\RequirePackage{fix-cm}
\documentclass[11pt]{article}
\usepackage[letterpaper,margin=1in]{geometry}
\usepackage[T1]{fontenc}
\usepackage{lmodern}
\usepackage[authoryear,round]{natbib}
\setcitestyle{citesep={;},aysep={,},yysep={;}}
\usepackage{amsmath,amssymb,amsthm,mathtools,mathrsfs}
\DeclareFontFamily{U}{rsfs}{\skewchar\font127}
\DeclareFontShape{U}{rsfs}{m}{n}{<-6.5> rsfs5 <6.5-8> rsfs7 <8-> rsfs10}{}
\usepackage{enumitem,graphicx,wrapfig,longtable,booktabs,tabularx,microtype,needspace}
\usepackage[colorlinks=true,allcolors=blue]{hyperref}
\usepackage{macros}
\usepackage{comment}
\usepackage{caption}
\usepackage{pgfplots,import,placeins,xfrac,xurl}
\pgfplotsset{compat=1.18}
\graphicspath{{figures/}}
\usepackage{url}

\setlength{\parindent}{0pt}
\setlength{\parskip}{0.45\baselineskip plus 1pt minus 1pt}
\setlength{\emergencystretch}{2em}
\renewcommand{\topfraction}{0.95}
\renewcommand{\textfraction}{0.05}
\raggedbottom

\title{Global Exponential Convergence\\of Two-Layer Linear Network Training}
\author{Gabriel Peyr\'e\\
  {\normalsize CNRS, ENS, PSL Univ.}\\
  {\small\texttt{gabriel.peyre@ens.fr}}
  \and
  Stephen Zhang\\
  {\normalsize Flatiron Institute \& NYU} \\ 
  {\small\texttt{syz@syz.id.au}}}
\date{}
\hypersetup{
  pdftitle={Global Exponential Convergence of Two-Layer Linear Network Training},
  pdfauthor={Gabriel Peyr\'e and Stephen Zhang}
}

\newcommand{\fix}{\marginpar{FIX}}
\newcommand{\new}{\marginpar{NEW}}

\begin{document}
\maketitle

\begin{abstract}
We prove global exponential (linear) convergence with an explicit rate in the rich scaling for wide two-layer linear networks trained with smooth
Polyak--\L{}ojasiewicz predictor losses. Gradient flow in the factors closes exactly in terms of a finite-dimensional Bures flow of the neuron law covariance, in which the predictor dynamics are preconditioned by hidden covariance blocks.
Mean-field conservation laws provide uniform spectral lower bounds on the hidden preconditioning blocks when the initial covariance satisfies a spectral support gap condition.
This condition encompasses positive definiteness while still allowing for singular initializations.
For an initial covariance $\Sigma_0 = \sigma^2 \Id$, the loss converges to the global minimum with linear rate at least $4\sigma^2\kappa$, where $\kappa$ is the PL constant. We establish stability of this rate under finite-width sampling, as well as global convergence of factor gradient descent for an explicit stepsize interval depending on smoothness, the initial loss, and conserved spectral margins.
Our argument extends layerwise to deep linear ResNets, subject to a residual-path bound.
In the case of heavy-ball momentum, training dynamics close instead over positions and velocities in terms of a lifted phase covariance.
Linear convergence holds under an explicit condition on the energy and damping, specifying a window of admissible dampings. For two-scale white initializations, this interval is nonempty for sufficiently large position scales, with a fixed initial loss gap and velocity covariance. Numerical experiments illustrate the covariance geometry and compare the predicted and observed rates.

%%% Local Variables:
%%% mode: LaTeX
%%% TeX-master: "../main"
%%% End:

\end{abstract}

% !TEX root = ../main.tex
% Section: Introduction

\section{Introduction}
\label{sec:intro}

Two-layer linear networks are nonconvex in their weight factors, even when the
loss is a simple function of the predictor matrix.  Training moves in the
lifted space of incoming and outgoing neuron weights, whose second moments act
as a time-dependent preconditioner on the predictor.  These moments may in
principle become degenerate and slow down the dynamics, and it is not clear that the dynamics necessarily converge to a globally optimal predictor.
We consider this question in the rich mean-field scaling, in which both sides of each rank-one neuron evolve, and
smooth predictor losses satisfying a Polyak--\L{}ojasiewicz (PL) inequality, which subsume the standard quadratic loss.
% Unlike lazy or neural tangent kernel scalings, the hidden representation is not
% frozen and its conditioning must be controlled globally in time.

Central to our analysis are width-independent conservation laws for
the weight covariance: the signed covariance spectrum, relative to the input-output
block split, supplies conserved supporting inequalities and lower bounds on each hidden-block. The sum of these lower bounds conditions the predictor dynamics: a positive sum
is sufficient to guarantee global convergence of the predictor loss with an explicit rate. This analysis covers every positive-definite initialization and some singular ones.
For initial variance $\sigma^2 \Id$, the loss decays with rate at least $4\sigma^2\kappa$, where $\kappa$ is the PL constant.

The covariance viewpoint also makes the result robust to finite-width
sampling and factor gradient descent. For deep linear ResNets, the same
conservation argument applied layerwise yields linear convergence
conditional on a uniform bound on the residual path. With heavy-ball
momentum, closure is lifted to the joint phase covariance of positions
and velocities. Conformally transported spectral laws describe its
geometry, while energy dissipation gives an explicit initialization and
damping certificate for linear convergence. Special values of the
conserved quantities further permit lower-dimensional intrinsic dynamics;
these reductions are described in Appendix~\ref{app:intrinsic}.

\paragraph{Related work}
% \paragraph{Mean-field flows and feature learning.}
Mean-field analyses describe the collective evolution of neuron weight
distributions in wide two-layer networks~\citep{ChizatBach2018,MeiMontanariNguyen2018,RotskoffVandenEijnden2018}.
These formulations treat broad nonlinear feature maps and often prove asymptotic global optimality or propagation-of-chaos type statements under support, homogeneity, noise, or regularity assumptions.
In the present linear setting, exact second-moment closure connects
training to Bures--Wasserstein covariance geometry~\citep{Takatsu2011,BhatiaJainLim2019},
making a sharp quantitative linear convergence rate available.
This closure distinguishes our analysis from the NTK~\citep{JacotGabrielHongler2018} and lazy training~\citep{ChizatOyallonBach2019} regime: both incoming and outgoing weights move, and the evolving representation is essential to the rate.

% \paragraph{Linear networks and conservation laws.}
Deep linear networks have served as a tractable model of the nonconvexity created by matrix products, with results on their landscapes, exact dynamics,
and convergence~\citep{SaxeMcClellandGanguli2014,Kawaguchi2016,LaurentVonBrecht2018,AroraCohenGolowichHu2019}.
Finite-width analyses of two-layer matrix factorization furnish convergence guarantees under small, asymmetric, or general initializations~\citep{YeDu2021,XuMinTarmounMalladaVidal2023}, while recent work treats discrete-time low-rank adaptation directly~\citep{MuKlabjan2025}.
Our result instead starts from a population covariance satisfying a positive \emph{support gap} condition, in particular any positive-definite covariance, and obtains an explicit global rate that survives independent finite-width sampling.
Related work identifies Gramian differences and balancedness as invariants of factor training~\citep{DuHuLee2018,TarmounFrancaHaeffeleVidal2021} and develops
a general differential framework for identifying conservation laws~\citep{MarcotteGribonvalPeyre2023}.
Our mean-field laws are functionals of the neuron distribution that induce
compatible particle laws at every width. The conserved spectrum of the
signed covariance yields hidden-block lower bounds without requiring
a special balancedness identity.

% \paragraph{Deep residual networks}
Residual networks~\citep{HeZhangRenSun2016Deep,HeZhangRenSun2016Identity} were introduced as architectures built from perturbations of the identity, and their infinite depth limit leads naturally to neural ODEs~\citep{HeZhangRenSun2016Deep,ChenRubanovaBettencourtDuvenaud2018}.
Mean-field formulations have established global optimality and convergence results under expressivity, regularity, or local PL conditions~\citep{LuMaLuLuYing2020,DingChenLiWright2021,Isobe2023,BarboniPeyreVialard2024COTResNets}.
Our linear analysis provides a complementary condition: uniform conditioning of the forward and backward residual products is sufficient for convergence.

% \paragraph{Heavy-ball momentum.}
Momentum turns first-order descent into an inertial, second-order iteration; the heavy-ball method of~\cite{Polyak1964} and Nesterov acceleration~\citep{Nesterov1983} have mature convergence theories, including continuous-time interpretations~\citep{SuBoydCandes2016} and robust-control analyses~\citep{LessardRechtPackard2016}.
For deep linear networks, \citet{WangLinAbernethy2021,LiuTaoPan2022}
establish acceleration results under initialization and overparameterization
conditions. Here the appropriate state variable is the phase covariance, over which the dynamics close.
We give an explicit energy-damping certificate for global convergence
with general position-velocity correlations. 

\paragraph{Our contributions}
Section~\ref{sec:laws} identifies width-independent spectral conservation
laws and derives uniform hidden-block bounds from their support margins
(Propositions~\ref{prop:covariance-closure-laws} and~\ref{lem:margins}).
Their Wasserstein interpretation and local completeness among smooth covariance
laws common to all split losses are developed in
Appendices~\ref{app:mean-field} and~\ref{app:law-counting}.
Section~\ref{sec:convergence} proves global exponential convergence with an explicit rate under a smooth PL predictor loss from every initial covariance with positive
support gap (Theorem~\ref{thm:gf}).
The rate survives independent finite-width sampling at width proportional to the population covariance rank, which can be
as small as $\min(d_u,d_v)$ (Corollary~\ref{cor:finite-width} and
Remark~\ref{rem:rank-adaptive-width}).
We establish global linear convergence of factor
gradient descent under an explicit stepsize bound by controlling the
accumulated conservation error (Theorem~\ref{thm:gd}).
Section~\ref{sec:resnets} extends the layerwise argument to finite- and
continuous-depth linear ResNets under residual-path control
(Theorems~\ref{thm:resnet-finite} and~\ref{thm:resnet-continuous}), and
discusses why that control is not automatic (Remark~\ref{rem:resnet-path-escape}
in Appendix~\ref{app:resnets}).
Section~\ref{sec:momentum} shows that heavy-ball momentum dynamics closes on the phase covariance and identifies conformal conservation laws (Proposition~\ref{prop:phase-laws}).
Global exponential convergence is proven under an initial-energy
and damping condition, including a two-scale white specialization
(Theorem~\ref{thm:momentum} and Corollary~\ref{cor:two-scale-white-window}).
A computational toolbox, code reproducing all six article figures, and
illustrative Jupyter notebooks are available at
\url{https://github.com/gpeyre/linear-networks-global}.

%%% Local Variables:
%%% mode: LaTeX
%%% TeX-master: "../main"
%%% End:

% !TEX root = ../main.tex
% Section: Finite-Width Setup and Predictor Loss

\label{sec:model}

\paragraph{Setup}
Writing training in terms of the predictor separates the loss geometry
from the nonconvex factorization.
A width-$n$ two-layer linear network has factors
$U:=(u_1,\ldots,u_n)\in\R^{d_u\times n}$ and
$V:=(v_1,\ldots,v_n)\in\R^{d_v\times n}$.
Write $m:=d_u+d_v$ and $W:=\binom UV$. Its predictor is
$\Svu:=VU^\top/n$, and its factor objective is
$\cL(U,V):=\ell(VU^\top/n)$ for a differentiable predictor loss
$\ell:\R^{d_v\times d_u}\to\R$.
For example, if a data distribution $\rho$ on pairs $(z,y)$ satisfies
$y=\Svu_\star z$ for a teacher $\Svu_\star$, squared regression loss is
\[
  \textstyle \ell_{\rm quad}(\Svu):=\frac12\int\norm{(\Svu-\Svu_\star)z}^2\,\mathrm{d}\rho(z,y).
\]
More generally, a self-adjoint positive semidefinite operator
$\quadH$ on $\R^{d_v\times d_u}$ defines a quadratic loss
$\ell(S):=\frac12\ip{\quadH[S-\Svu_\star]}{S-\Svu_\star}$.
The regression example corresponds to $\quadH[X]=XC_z$, where
$C_z:=\int zz^\top\,\dd\rho(z,y)$ is the input second moment.

We state the assumptions on the predictor loss: assume throughout that \(\ell:\R^{d_v\times d_u}\to\R\) is \(C^1\) and bounded below, and write \(\ell_\star:=\inf_S\ell(S)\).

\begin{assumption}[Smooth PL predictor loss]
\label{ass:predictor-loss}\label{ass:loss}
There exist constants \(L>0\) and \(\kappa>0\) such that the loss
\(\ell\) is \(L\)-smooth and satisfies the \(\kappa\)-Polyak--\L{}ojasiewicz
inequality: \(\norm{\nabla\ell(S)-\nabla\ell(S')}_{\F}\le L\norm{S-S'}_{\F}\) for all \(S,S'\),
and \(\norm{\nabla\ell(S)}_{\F}^2 \ge 2\kappa\bigl(\ell(S)-\ell_\star\bigr)\) for all \(S\).
\end{assumption}

We take \(0<\kappa\le L\), as holds automatically for any nonconstant loss
satisfying Assumption~\ref{ass:loss}.
For a nonzero quadratic Hessian $\quadH\succeq0$, one can take
$L=\lambda_{\max}(\quadH)$ and $\kappa=\lambda_{\min}^{+}(\quadH)$.
The operator $\quadH$ may have a nontrivial kernel.

%%% Local Variables:
%%% mode: LaTeX
%%% TeX-master: "../main.tex"
%%% End:

\section{Mean Field Conservation Laws}
\label{sec:laws}\label{sec:conservation-hyperbolic}

\paragraph{Gradient flow and covariance closure}
Although training evolves all neuron weights, its effect on the
predictor is governed by a fixed-size covariance matrix.
We use the width-normalized gradient flow
$(\dot U_t,\dot V_t)=-n\nabla_{(U,V)}\cL(U_t,V_t)$, namely
\begin{equation}\label{eq:mean-field-gradient-flow}
 \dot U_t=-E_t^\top V_t,\qquad
 \dot V_t=-E_tU_t,\qquad E_t:=\nabla\ell(V_tU_t^\top/n).
\end{equation}
The factor $n$ fixes a common time scale across widths; equivalently
this is Euclidean gradient flow of $n\cL$.
With $W_t:=\binom{U_t}{V_t}$, define the uncentered covariance
\[
 \Sigma_t:=\tfrac1nW_tW_t^\top
 =\begin{psmallmatrix}\Suu_t&\Suv_t\\\Svu_t&\Svv_t\end{psmallmatrix},
 \qquad
 \Suu_t:=\tfrac{U_tU_t^\top}{n},\quad
 \Svv_t:=\tfrac{V_tV_t^\top}{n},\quad
 \Svu_t:=\tfrac{V_tU_t^\top}{n}.
\]
The hidden blocks $\Suu_t,\Svv_t$ do not enter the loss, but govern
the predictor's motion. We use the same notation $\Sigma$ whether the
covariance is specified directly or realized by finite-width factors:
the closed equation depends only on this matrix.
Set
$G_E:=\left(\begin{smallmatrix}0&E^\top\\E&0\end{smallmatrix}\right)$,
$G_t:=G_{E_t}$, and $G(S):=G_{\nabla\ell(S)}$, so $\dot W_t=-G_tW_t$.
The split structure of the gradient is encoded by the signature matrix
\[
 J:=\diag(\Id_{d_u},-\Id_{d_v}),\qquad G_tJ+JG_t=0.
\]
Our first result combines covariance closure with width-independent spectral conservation laws.
\begin{proposition}[Covariance closure and conservation laws]
\label{prop:covariance-closure-laws}
Let $\mathcal H_k(\Sigma):=\tr\!\left((J\Sigma)^k\right)$.
On its interval of existence, the factor flow
\eqref{eq:mean-field-gradient-flow} closes on the covariance $\Sigma_t$ as a Bures equation, for which the characteristic polynomial of $J\Sigma$ is preserved:
\begin{equation}\label{eq:bures-G}
 \dot\Sigma_t=-(G_t\Sigma_t+\Sigma_tG_t), \qquad
 \tfrac{\dd}{\dd t}\mathcal H_k(\Sigma_t)=0, \quad k=1,\ldots,m,
\end{equation}
\end{proposition}

\begin{comment}
\emph{Proof sketch.}
Differentiate $W_tW_t^\top/n$ to obtain the closed equation.
The identity $JG_t=-G_tJ$ shows that $J\Sigma_t$ evolves following a Lax equation: $\frac{\dd}{\dd t}(J\Sigma_t)=[G_t,J\Sigma_t] = G_t J \Sigma_t - J \Sigma_t G_t$, whose trace powers
are constant by cyclicity of the trace. Newton identities give the
characteristic polynomial; see Appendix~\ref{app:covariance-laws} for a full proof.\qed
\end{comment}

The covariance equation \eqref{eq:bures-G} is the Bures gradient
flow~\citep{Takatsu2011,BhatiaJainLim2019} of
$F(\Sigma):=\ell(\Svu)$, since $\nabla_\Sigma F=G/2$.
Its fundamental matrix $\Phi_t$, defined by $\dot\Phi_t=-G_t\Phi_t$ and $\Phi_0=\Id$, satisfies
\begin{equation}\label{eq:orbit}
 \Phi_t^\top J\Phi_t=  \Phi_tJ\Phi_t^\top=J,\qquad
 \Sigma_t=\Phi_t\Sigma_0\Phi_t^\top.
\end{equation}
The functions $\mathcal H_k$ are conserved at every width $n$, defining mean-field laws.
Appendix~\ref{app:mean-field} gives their measure formulation, the corresponding Wasserstein PDE, and describes a larger family of laws.

\paragraph{Conserved spectral margins}
The spectral conservation laws \eqref{eq:bures-G} give uniform lower bounds on the spectrum of the hidden blocks $(\Suu, \Svv)$ through two margins
of the full covariance.
\begin{definition}[Covariance support margins and gap]\label{def:support-gap}
For $\Sigma\succeq0$, the support margins and their sum are
% Compact spacing local to this definition.
\setlength{\abovedisplayskip}{2pt}
\setlength{\belowdisplayskip}{2pt}
\setlength{\abovedisplayshortskip}{0pt}
\setlength{\belowdisplayshortskip}{2pt}
\begin{equation}\label{eq:margins}
 a_\pm(\Sigma):=\max\{a\ge0:\Sigma\mp aJ\succeq0\},\qquad
 \gap{\Sigma}:=\aplus{\Sigma}+\aminus{\Sigma}.
\end{equation}
\end{definition}
These maxima are finite and nonnegative. If $\lambda_1\le\cdots\le\lambda_m$
are the eigenvalues of $\Sigma^{1/2}J\Sigma^{1/2}$, \emph{including zeros}, then
\begin{equation}\label{eq:gap-spectrum}
 \aplus{\Sigma}=\lambda_{d_v+1},\qquad
 \aminus{\Sigma}=-\lambda_{d_v},\qquad
 \gap{\Sigma}=\lambda_{d_v+1}-\lambda_{d_v}.
\end{equation}
Thus $\gap{\Sigma}$ is an indexed spectral gap. Equivalently, it is the
length of the positive-semidefinite interval
$[-\aminus{\Sigma},\aplus{\Sigma}]$ of the pencil $\Sigma-tJ$;
such intervals are standard in matrix-pencil theory
\citep{NguyenNguyen2023}. Appendix~\ref{app:technical-details} proves
\eqref{eq:gap-spectrum}, including at singular covariances.
For $\Sigma\succ0$, $\aplus{\Sigma}$ and $-\aminus{\Sigma}$ are the
signed eigenvalues of $\Sigma^{1/2} J \Sigma^{1/2}$ nearest zero on their respective sides, and
$\gap{\Sigma}\ge2\lambda_{\min}(\Sigma)>0$.

\begin{proposition}[Transported lower bounds]\label{lem:margins}
On its interval of existence, every solution of \eqref{eq:bures-G} from
$\Sigma_0\succeq0$ preserves $a_\pm(\Sigma_t)=a_\pm(\Sigma_0)$ and satisfies
\[
  \Sigma_t \mp a_{\pm}(\Sigma_0) J \succeq 0, \qquad
  \Suu_t\succeq\aplus{\Sigma_0}\Id_{d_u},\qquad
  \Svv_t\succeq\aminus{\Sigma_0}\Id_{d_v}.
\]
\end{proposition}

\begin{comment}
\emph{Proof sketch.}
The key is to track matrix inequalities, rather than eigenvalues.
The supporting inequalities hold at $t = 0$ by definition. Congruence by
$\Phi_t$ preserves them because $\Phi_tJ\Phi_t^\top=J$; congruence by
$\Phi_t^{-1}$ proves that the maximal margins are unchanged.
Taking principal blocks then gives the hidden bounds. The complete proof is in
Appendix~\ref{app:technical-details}.\qed
\end{comment}

The margins $a_\pm(\Sigma)$ need not equal the minimum eigenvalues of the initial hidden blocks: they satisfy
$\aplus{\Sigma}\le\lambda_{\min}(\Suu)$ and
$\aminus{\Sigma}\le\lambda_{\min}(\Svv)$, with equality when $\Svu=0$.
For instance, letting $\Sigma=\left(\begin{smallmatrix}1&r\\r&1\end{smallmatrix}\right)$,
$|r|\le1$, both margins are $\sqrt{1-r^2}$, whereas both hidden variances are one.
The support gap $g(\Sigma)$ controls predictor dissipation and can remain positive even when the full covariance is degenerate.
\begin{wrapfigure}[22]{r}{0.32\textwidth}
  \centering
  \resizebox{\linewidth}{!}{%% Creator: Matplotlib, PGF backend
%%
%% To include the figure in your LaTeX document, write
%%   \input{<filename>.pgf}
%%
%% Make sure the required packages are loaded in your preamble
%%   \usepackage{pgf}
%%
%% Also ensure that all the required font packages are loaded; for instance,
%% the lmodern package is sometimes necessary when using math font.
%%   \usepackage{lmodern}
%%
%% Figures using additional raster images can only be included by \input if
%% they are in the same directory as the main LaTeX file. For loading figures
%% from other directories you can use the `import` package
%%   \usepackage{import}
%%
%% and then include the figures with
%%   \import{<path to file>}{<filename>.pgf}
%%
%% Matplotlib used the following preamble
%%   \def\mathdefault#1{#1}
%%   \everymath=\expandafter{\the\everymath\displaystyle}
%%   \IfFileExists{scrextend.sty}{
%%     \usepackage[fontsize=6.000000pt]{scrextend}
%%   }{
%%     \renewcommand{\normalsize}{\fontsize{6.000000}{7.200000}\selectfont}
%%     \normalsize
%%   }
%%   \usepackage{amsfonts}
%%   \usepackage{amssymb}
%%   \ifdefined\pdftexversion\else  % non-pdftex case.
%%     \usepackage{fontspec}
%%   \fi
%%   \makeatletter\@ifpackageloaded{underscore}{}{\usepackage[strings]{underscore}}\makeatother
%%
\begingroup%
\makeatletter%
\begin{pgfpicture}%
\pgfpathrectangle{\pgfpointorigin}{\pgfqpoint{2.150000in}{3.350000in}}%
\pgfusepath{use as bounding box, clip}%
\begin{pgfscope}%
\pgfsetbuttcap%
\pgfsetmiterjoin%
\definecolor{currentfill}{rgb}{1.000000,1.000000,1.000000}%
\pgfsetfillcolor{currentfill}%
\pgfsetlinewidth{0.000000pt}%
\definecolor{currentstroke}{rgb}{1.000000,1.000000,1.000000}%
\pgfsetstrokecolor{currentstroke}%
\pgfsetdash{}{0pt}%
\pgfpathmoveto{\pgfqpoint{0.000000in}{0.000000in}}%
\pgfpathlineto{\pgfqpoint{2.150000in}{0.000000in}}%
\pgfpathlineto{\pgfqpoint{2.150000in}{3.350000in}}%
\pgfpathlineto{\pgfqpoint{0.000000in}{3.350000in}}%
\pgfpathlineto{\pgfqpoint{0.000000in}{0.000000in}}%
\pgfpathclose%
\pgfusepath{fill}%
\end{pgfscope}%
\begin{pgfscope}%
\pgfsetbuttcap%
\pgfsetmiterjoin%
\definecolor{currentfill}{rgb}{1.000000,1.000000,1.000000}%
\pgfsetfillcolor{currentfill}%
\pgfsetlinewidth{0.000000pt}%
\definecolor{currentstroke}{rgb}{0.000000,0.000000,0.000000}%
\pgfsetstrokecolor{currentstroke}%
\pgfsetstrokeopacity{0.000000}%
\pgfsetdash{}{0pt}%
\pgfpathmoveto{\pgfqpoint{0.537500in}{2.059688in}}%
\pgfpathlineto{\pgfqpoint{2.042500in}{2.059688in}}%
\pgfpathlineto{\pgfqpoint{2.042500in}{3.216000in}}%
\pgfpathlineto{\pgfqpoint{0.537500in}{3.216000in}}%
\pgfpathlineto{\pgfqpoint{0.537500in}{2.059688in}}%
\pgfpathclose%
\pgfusepath{fill}%
\end{pgfscope}%
\begin{pgfscope}%
\pgfpathrectangle{\pgfqpoint{0.537500in}{2.059688in}}{\pgfqpoint{1.505000in}{1.156312in}}%
\pgfusepath{clip}%
\pgfsetbuttcap%
\pgfsetroundjoin%
\definecolor{currentfill}{rgb}{0.933333,0.941176,0.949020}%
\pgfsetfillcolor{currentfill}%
\pgfsetlinewidth{1.003750pt}%
\definecolor{currentstroke}{rgb}{0.933333,0.941176,0.949020}%
\pgfsetstrokecolor{currentstroke}%
\pgfsetdash{}{0pt}%
\pgfsys@defobject{currentmarker}{\pgfqpoint{0.537500in}{2.060627in}}{\pgfqpoint{2.042500in}{3.216000in}}{%
\pgfpathmoveto{\pgfqpoint{0.537500in}{3.216000in}}%
\pgfpathlineto{\pgfqpoint{0.537500in}{2.840980in}}%
\pgfpathlineto{\pgfqpoint{0.539653in}{2.838566in}}%
\pgfpathlineto{\pgfqpoint{0.541806in}{2.836151in}}%
\pgfpathlineto{\pgfqpoint{0.543959in}{2.833737in}}%
\pgfpathlineto{\pgfqpoint{0.546112in}{2.831323in}}%
\pgfpathlineto{\pgfqpoint{0.548265in}{2.828909in}}%
\pgfpathlineto{\pgfqpoint{0.550418in}{2.826494in}}%
\pgfpathlineto{\pgfqpoint{0.552572in}{2.824080in}}%
\pgfpathlineto{\pgfqpoint{0.554725in}{2.821666in}}%
\pgfpathlineto{\pgfqpoint{0.556878in}{2.819251in}}%
\pgfpathlineto{\pgfqpoint{0.559031in}{2.816837in}}%
\pgfpathlineto{\pgfqpoint{0.561184in}{2.814423in}}%
\pgfpathlineto{\pgfqpoint{0.563337in}{2.812009in}}%
\pgfpathlineto{\pgfqpoint{0.565490in}{2.809594in}}%
\pgfpathlineto{\pgfqpoint{0.567643in}{2.807180in}}%
\pgfpathlineto{\pgfqpoint{0.569796in}{2.804766in}}%
\pgfpathlineto{\pgfqpoint{0.571949in}{2.802351in}}%
\pgfpathlineto{\pgfqpoint{0.574102in}{2.799937in}}%
\pgfpathlineto{\pgfqpoint{0.576255in}{2.797523in}}%
\pgfpathlineto{\pgfqpoint{0.578408in}{2.795108in}}%
\pgfpathlineto{\pgfqpoint{0.580562in}{2.792694in}}%
\pgfpathlineto{\pgfqpoint{0.582715in}{2.790280in}}%
\pgfpathlineto{\pgfqpoint{0.584868in}{2.787866in}}%
\pgfpathlineto{\pgfqpoint{0.587021in}{2.785451in}}%
\pgfpathlineto{\pgfqpoint{0.589174in}{2.783037in}}%
\pgfpathlineto{\pgfqpoint{0.591327in}{2.780623in}}%
\pgfpathlineto{\pgfqpoint{0.593480in}{2.778208in}}%
\pgfpathlineto{\pgfqpoint{0.595633in}{2.775794in}}%
\pgfpathlineto{\pgfqpoint{0.597786in}{2.773380in}}%
\pgfpathlineto{\pgfqpoint{0.599939in}{2.770966in}}%
\pgfpathlineto{\pgfqpoint{0.602092in}{2.768551in}}%
\pgfpathlineto{\pgfqpoint{0.604245in}{2.766137in}}%
\pgfpathlineto{\pgfqpoint{0.606398in}{2.763723in}}%
\pgfpathlineto{\pgfqpoint{0.608552in}{2.761308in}}%
\pgfpathlineto{\pgfqpoint{0.610705in}{2.758894in}}%
\pgfpathlineto{\pgfqpoint{0.612858in}{2.756480in}}%
\pgfpathlineto{\pgfqpoint{0.615011in}{2.754066in}}%
\pgfpathlineto{\pgfqpoint{0.617164in}{2.751651in}}%
\pgfpathlineto{\pgfqpoint{0.619317in}{2.749237in}}%
\pgfpathlineto{\pgfqpoint{0.621470in}{2.746823in}}%
\pgfpathlineto{\pgfqpoint{0.623623in}{2.744408in}}%
\pgfpathlineto{\pgfqpoint{0.625776in}{2.741994in}}%
\pgfpathlineto{\pgfqpoint{0.627929in}{2.739580in}}%
\pgfpathlineto{\pgfqpoint{0.630082in}{2.737165in}}%
\pgfpathlineto{\pgfqpoint{0.632235in}{2.734751in}}%
\pgfpathlineto{\pgfqpoint{0.634388in}{2.732337in}}%
\pgfpathlineto{\pgfqpoint{0.636541in}{2.729923in}}%
\pgfpathlineto{\pgfqpoint{0.638695in}{2.727508in}}%
\pgfpathlineto{\pgfqpoint{0.640848in}{2.725094in}}%
\pgfpathlineto{\pgfqpoint{0.643001in}{2.722680in}}%
\pgfpathlineto{\pgfqpoint{0.645154in}{2.720265in}}%
\pgfpathlineto{\pgfqpoint{0.647307in}{2.717851in}}%
\pgfpathlineto{\pgfqpoint{0.649460in}{2.715437in}}%
\pgfpathlineto{\pgfqpoint{0.651613in}{2.713023in}}%
\pgfpathlineto{\pgfqpoint{0.653766in}{2.710608in}}%
\pgfpathlineto{\pgfqpoint{0.655919in}{2.708194in}}%
\pgfpathlineto{\pgfqpoint{0.658072in}{2.705780in}}%
\pgfpathlineto{\pgfqpoint{0.660225in}{2.703365in}}%
\pgfpathlineto{\pgfqpoint{0.662378in}{2.700951in}}%
\pgfpathlineto{\pgfqpoint{0.664531in}{2.698537in}}%
\pgfpathlineto{\pgfqpoint{0.666685in}{2.696123in}}%
\pgfpathlineto{\pgfqpoint{0.668838in}{2.693708in}}%
\pgfpathlineto{\pgfqpoint{0.670991in}{2.691294in}}%
\pgfpathlineto{\pgfqpoint{0.673144in}{2.688880in}}%
\pgfpathlineto{\pgfqpoint{0.675297in}{2.686465in}}%
\pgfpathlineto{\pgfqpoint{0.677450in}{2.684051in}}%
\pgfpathlineto{\pgfqpoint{0.679603in}{2.681637in}}%
\pgfpathlineto{\pgfqpoint{0.681756in}{2.679222in}}%
\pgfpathlineto{\pgfqpoint{0.683909in}{2.676808in}}%
\pgfpathlineto{\pgfqpoint{0.686062in}{2.674394in}}%
\pgfpathlineto{\pgfqpoint{0.688215in}{2.671980in}}%
\pgfpathlineto{\pgfqpoint{0.690368in}{2.669565in}}%
\pgfpathlineto{\pgfqpoint{0.692521in}{2.667151in}}%
\pgfpathlineto{\pgfqpoint{0.694675in}{2.664737in}}%
\pgfpathlineto{\pgfqpoint{0.696828in}{2.662322in}}%
\pgfpathlineto{\pgfqpoint{0.698981in}{2.659908in}}%
\pgfpathlineto{\pgfqpoint{0.701134in}{2.657494in}}%
\pgfpathlineto{\pgfqpoint{0.703287in}{2.655080in}}%
\pgfpathlineto{\pgfqpoint{0.705440in}{2.652665in}}%
\pgfpathlineto{\pgfqpoint{0.707593in}{2.650251in}}%
\pgfpathlineto{\pgfqpoint{0.709746in}{2.647837in}}%
\pgfpathlineto{\pgfqpoint{0.711899in}{2.645422in}}%
\pgfpathlineto{\pgfqpoint{0.714052in}{2.643008in}}%
\pgfpathlineto{\pgfqpoint{0.716205in}{2.640594in}}%
\pgfpathlineto{\pgfqpoint{0.718358in}{2.638179in}}%
\pgfpathlineto{\pgfqpoint{0.720511in}{2.635765in}}%
\pgfpathlineto{\pgfqpoint{0.722665in}{2.633351in}}%
\pgfpathlineto{\pgfqpoint{0.724818in}{2.630937in}}%
\pgfpathlineto{\pgfqpoint{0.726971in}{2.628522in}}%
\pgfpathlineto{\pgfqpoint{0.729124in}{2.626108in}}%
\pgfpathlineto{\pgfqpoint{0.731277in}{2.623694in}}%
\pgfpathlineto{\pgfqpoint{0.733430in}{2.621279in}}%
\pgfpathlineto{\pgfqpoint{0.735583in}{2.618865in}}%
\pgfpathlineto{\pgfqpoint{0.737736in}{2.616451in}}%
\pgfpathlineto{\pgfqpoint{0.739889in}{2.614037in}}%
\pgfpathlineto{\pgfqpoint{0.742042in}{2.611622in}}%
\pgfpathlineto{\pgfqpoint{0.744195in}{2.609208in}}%
\pgfpathlineto{\pgfqpoint{0.746348in}{2.606794in}}%
\pgfpathlineto{\pgfqpoint{0.748501in}{2.604379in}}%
\pgfpathlineto{\pgfqpoint{0.750655in}{2.601965in}}%
\pgfpathlineto{\pgfqpoint{0.752808in}{2.599551in}}%
\pgfpathlineto{\pgfqpoint{0.754961in}{2.597137in}}%
\pgfpathlineto{\pgfqpoint{0.757114in}{2.594722in}}%
\pgfpathlineto{\pgfqpoint{0.759267in}{2.592308in}}%
\pgfpathlineto{\pgfqpoint{0.761420in}{2.589894in}}%
\pgfpathlineto{\pgfqpoint{0.763573in}{2.587479in}}%
\pgfpathlineto{\pgfqpoint{0.765726in}{2.585065in}}%
\pgfpathlineto{\pgfqpoint{0.767879in}{2.582651in}}%
\pgfpathlineto{\pgfqpoint{0.770032in}{2.580236in}}%
\pgfpathlineto{\pgfqpoint{0.772185in}{2.577822in}}%
\pgfpathlineto{\pgfqpoint{0.774338in}{2.575408in}}%
\pgfpathlineto{\pgfqpoint{0.776491in}{2.572994in}}%
\pgfpathlineto{\pgfqpoint{0.778644in}{2.570579in}}%
\pgfpathlineto{\pgfqpoint{0.780798in}{2.568165in}}%
\pgfpathlineto{\pgfqpoint{0.782951in}{2.565751in}}%
\pgfpathlineto{\pgfqpoint{0.785104in}{2.563336in}}%
\pgfpathlineto{\pgfqpoint{0.787257in}{2.560922in}}%
\pgfpathlineto{\pgfqpoint{0.789410in}{2.558508in}}%
\pgfpathlineto{\pgfqpoint{0.791563in}{2.556094in}}%
\pgfpathlineto{\pgfqpoint{0.793716in}{2.553679in}}%
\pgfpathlineto{\pgfqpoint{0.795869in}{2.551265in}}%
\pgfpathlineto{\pgfqpoint{0.798022in}{2.548851in}}%
\pgfpathlineto{\pgfqpoint{0.800175in}{2.546436in}}%
\pgfpathlineto{\pgfqpoint{0.802328in}{2.544022in}}%
\pgfpathlineto{\pgfqpoint{0.804481in}{2.541608in}}%
\pgfpathlineto{\pgfqpoint{0.806634in}{2.539194in}}%
\pgfpathlineto{\pgfqpoint{0.808788in}{2.536779in}}%
\pgfpathlineto{\pgfqpoint{0.810941in}{2.534365in}}%
\pgfpathlineto{\pgfqpoint{0.813094in}{2.531951in}}%
\pgfpathlineto{\pgfqpoint{0.815247in}{2.529536in}}%
\pgfpathlineto{\pgfqpoint{0.817400in}{2.527122in}}%
\pgfpathlineto{\pgfqpoint{0.819553in}{2.524708in}}%
\pgfpathlineto{\pgfqpoint{0.821706in}{2.522293in}}%
\pgfpathlineto{\pgfqpoint{0.823859in}{2.519879in}}%
\pgfpathlineto{\pgfqpoint{0.826012in}{2.517465in}}%
\pgfpathlineto{\pgfqpoint{0.828165in}{2.515051in}}%
\pgfpathlineto{\pgfqpoint{0.830318in}{2.512636in}}%
\pgfpathlineto{\pgfqpoint{0.832471in}{2.510222in}}%
\pgfpathlineto{\pgfqpoint{0.834624in}{2.507808in}}%
\pgfpathlineto{\pgfqpoint{0.836778in}{2.505393in}}%
\pgfpathlineto{\pgfqpoint{0.838931in}{2.502979in}}%
\pgfpathlineto{\pgfqpoint{0.841084in}{2.500565in}}%
\pgfpathlineto{\pgfqpoint{0.843237in}{2.498151in}}%
\pgfpathlineto{\pgfqpoint{0.845390in}{2.495736in}}%
\pgfpathlineto{\pgfqpoint{0.847543in}{2.493322in}}%
\pgfpathlineto{\pgfqpoint{0.849696in}{2.490908in}}%
\pgfpathlineto{\pgfqpoint{0.851849in}{2.488493in}}%
\pgfpathlineto{\pgfqpoint{0.854002in}{2.486079in}}%
\pgfpathlineto{\pgfqpoint{0.856155in}{2.483665in}}%
\pgfpathlineto{\pgfqpoint{0.858308in}{2.481251in}}%
\pgfpathlineto{\pgfqpoint{0.860461in}{2.478836in}}%
\pgfpathlineto{\pgfqpoint{0.862614in}{2.476422in}}%
\pgfpathlineto{\pgfqpoint{0.864768in}{2.474008in}}%
\pgfpathlineto{\pgfqpoint{0.866921in}{2.471593in}}%
\pgfpathlineto{\pgfqpoint{0.869074in}{2.469179in}}%
\pgfpathlineto{\pgfqpoint{0.871227in}{2.466765in}}%
\pgfpathlineto{\pgfqpoint{0.873380in}{2.464350in}}%
\pgfpathlineto{\pgfqpoint{0.875533in}{2.461936in}}%
\pgfpathlineto{\pgfqpoint{0.877686in}{2.459522in}}%
\pgfpathlineto{\pgfqpoint{0.879839in}{2.457108in}}%
\pgfpathlineto{\pgfqpoint{0.881992in}{2.454693in}}%
\pgfpathlineto{\pgfqpoint{0.884145in}{2.452279in}}%
\pgfpathlineto{\pgfqpoint{0.886298in}{2.449865in}}%
\pgfpathlineto{\pgfqpoint{0.888451in}{2.447450in}}%
\pgfpathlineto{\pgfqpoint{0.890604in}{2.445036in}}%
\pgfpathlineto{\pgfqpoint{0.892758in}{2.442622in}}%
\pgfpathlineto{\pgfqpoint{0.894911in}{2.440208in}}%
\pgfpathlineto{\pgfqpoint{0.897064in}{2.437793in}}%
\pgfpathlineto{\pgfqpoint{0.899217in}{2.435379in}}%
\pgfpathlineto{\pgfqpoint{0.901370in}{2.432965in}}%
\pgfpathlineto{\pgfqpoint{0.903523in}{2.430550in}}%
\pgfpathlineto{\pgfqpoint{0.905676in}{2.428136in}}%
\pgfpathlineto{\pgfqpoint{0.907829in}{2.425722in}}%
\pgfpathlineto{\pgfqpoint{0.909982in}{2.423308in}}%
\pgfpathlineto{\pgfqpoint{0.912135in}{2.420893in}}%
\pgfpathlineto{\pgfqpoint{0.914288in}{2.418479in}}%
\pgfpathlineto{\pgfqpoint{0.916441in}{2.416065in}}%
\pgfpathlineto{\pgfqpoint{0.918594in}{2.413650in}}%
\pgfpathlineto{\pgfqpoint{0.920747in}{2.411236in}}%
\pgfpathlineto{\pgfqpoint{0.922901in}{2.408822in}}%
\pgfpathlineto{\pgfqpoint{0.925054in}{2.406407in}}%
\pgfpathlineto{\pgfqpoint{0.927207in}{2.403993in}}%
\pgfpathlineto{\pgfqpoint{0.929360in}{2.401579in}}%
\pgfpathlineto{\pgfqpoint{0.931513in}{2.399165in}}%
\pgfpathlineto{\pgfqpoint{0.933666in}{2.396750in}}%
\pgfpathlineto{\pgfqpoint{0.935819in}{2.394336in}}%
\pgfpathlineto{\pgfqpoint{0.937972in}{2.391922in}}%
\pgfpathlineto{\pgfqpoint{0.940125in}{2.389507in}}%
\pgfpathlineto{\pgfqpoint{0.942278in}{2.387093in}}%
\pgfpathlineto{\pgfqpoint{0.944431in}{2.384679in}}%
\pgfpathlineto{\pgfqpoint{0.946584in}{2.382265in}}%
\pgfpathlineto{\pgfqpoint{0.948737in}{2.379850in}}%
\pgfpathlineto{\pgfqpoint{0.950891in}{2.377436in}}%
\pgfpathlineto{\pgfqpoint{0.953044in}{2.375022in}}%
\pgfpathlineto{\pgfqpoint{0.955197in}{2.372607in}}%
\pgfpathlineto{\pgfqpoint{0.957350in}{2.370193in}}%
\pgfpathlineto{\pgfqpoint{0.959503in}{2.367779in}}%
\pgfpathlineto{\pgfqpoint{0.961656in}{2.365365in}}%
\pgfpathlineto{\pgfqpoint{0.963809in}{2.362950in}}%
\pgfpathlineto{\pgfqpoint{0.965962in}{2.360536in}}%
\pgfpathlineto{\pgfqpoint{0.968115in}{2.358122in}}%
\pgfpathlineto{\pgfqpoint{0.970268in}{2.355707in}}%
\pgfpathlineto{\pgfqpoint{0.972421in}{2.353293in}}%
\pgfpathlineto{\pgfqpoint{0.974574in}{2.350879in}}%
\pgfpathlineto{\pgfqpoint{0.976727in}{2.348464in}}%
\pgfpathlineto{\pgfqpoint{0.978881in}{2.346050in}}%
\pgfpathlineto{\pgfqpoint{0.981034in}{2.343636in}}%
\pgfpathlineto{\pgfqpoint{0.983187in}{2.341222in}}%
\pgfpathlineto{\pgfqpoint{0.985340in}{2.338807in}}%
\pgfpathlineto{\pgfqpoint{0.987493in}{2.336393in}}%
\pgfpathlineto{\pgfqpoint{0.989646in}{2.333979in}}%
\pgfpathlineto{\pgfqpoint{0.991799in}{2.331564in}}%
\pgfpathlineto{\pgfqpoint{0.993952in}{2.329150in}}%
\pgfpathlineto{\pgfqpoint{0.996105in}{2.326736in}}%
\pgfpathlineto{\pgfqpoint{0.998258in}{2.324322in}}%
\pgfpathlineto{\pgfqpoint{1.000411in}{2.321907in}}%
\pgfpathlineto{\pgfqpoint{1.002564in}{2.319493in}}%
\pgfpathlineto{\pgfqpoint{1.004717in}{2.317079in}}%
\pgfpathlineto{\pgfqpoint{1.006871in}{2.314664in}}%
\pgfpathlineto{\pgfqpoint{1.009024in}{2.312250in}}%
\pgfpathlineto{\pgfqpoint{1.011177in}{2.309836in}}%
\pgfpathlineto{\pgfqpoint{1.013330in}{2.307422in}}%
\pgfpathlineto{\pgfqpoint{1.015483in}{2.305007in}}%
\pgfpathlineto{\pgfqpoint{1.017636in}{2.302593in}}%
\pgfpathlineto{\pgfqpoint{1.019789in}{2.300179in}}%
\pgfpathlineto{\pgfqpoint{1.021942in}{2.297764in}}%
\pgfpathlineto{\pgfqpoint{1.024095in}{2.295350in}}%
\pgfpathlineto{\pgfqpoint{1.026248in}{2.292936in}}%
\pgfpathlineto{\pgfqpoint{1.028401in}{2.290521in}}%
\pgfpathlineto{\pgfqpoint{1.030554in}{2.288107in}}%
\pgfpathlineto{\pgfqpoint{1.032707in}{2.285693in}}%
\pgfpathlineto{\pgfqpoint{1.034861in}{2.283279in}}%
\pgfpathlineto{\pgfqpoint{1.037014in}{2.280864in}}%
\pgfpathlineto{\pgfqpoint{1.039167in}{2.278450in}}%
\pgfpathlineto{\pgfqpoint{1.041320in}{2.276036in}}%
\pgfpathlineto{\pgfqpoint{1.043473in}{2.273621in}}%
\pgfpathlineto{\pgfqpoint{1.045626in}{2.271207in}}%
\pgfpathlineto{\pgfqpoint{1.047779in}{2.268793in}}%
\pgfpathlineto{\pgfqpoint{1.049932in}{2.266379in}}%
\pgfpathlineto{\pgfqpoint{1.052085in}{2.263964in}}%
\pgfpathlineto{\pgfqpoint{1.054238in}{2.261550in}}%
\pgfpathlineto{\pgfqpoint{1.056391in}{2.259136in}}%
\pgfpathlineto{\pgfqpoint{1.058544in}{2.256721in}}%
\pgfpathlineto{\pgfqpoint{1.060697in}{2.254307in}}%
\pgfpathlineto{\pgfqpoint{1.062851in}{2.251893in}}%
\pgfpathlineto{\pgfqpoint{1.065004in}{2.249479in}}%
\pgfpathlineto{\pgfqpoint{1.067157in}{2.247064in}}%
\pgfpathlineto{\pgfqpoint{1.069310in}{2.244650in}}%
\pgfpathlineto{\pgfqpoint{1.071463in}{2.242236in}}%
\pgfpathlineto{\pgfqpoint{1.073616in}{2.239821in}}%
\pgfpathlineto{\pgfqpoint{1.075769in}{2.237407in}}%
\pgfpathlineto{\pgfqpoint{1.077922in}{2.234993in}}%
\pgfpathlineto{\pgfqpoint{1.080075in}{2.232578in}}%
\pgfpathlineto{\pgfqpoint{1.082228in}{2.230164in}}%
\pgfpathlineto{\pgfqpoint{1.084381in}{2.227750in}}%
\pgfpathlineto{\pgfqpoint{1.086534in}{2.225336in}}%
\pgfpathlineto{\pgfqpoint{1.088687in}{2.222921in}}%
\pgfpathlineto{\pgfqpoint{1.090840in}{2.220507in}}%
\pgfpathlineto{\pgfqpoint{1.092994in}{2.218093in}}%
\pgfpathlineto{\pgfqpoint{1.095147in}{2.215678in}}%
\pgfpathlineto{\pgfqpoint{1.097300in}{2.213264in}}%
\pgfpathlineto{\pgfqpoint{1.099453in}{2.210850in}}%
\pgfpathlineto{\pgfqpoint{1.101606in}{2.208436in}}%
\pgfpathlineto{\pgfqpoint{1.103759in}{2.206021in}}%
\pgfpathlineto{\pgfqpoint{1.105912in}{2.203607in}}%
\pgfpathlineto{\pgfqpoint{1.108065in}{2.201193in}}%
\pgfpathlineto{\pgfqpoint{1.110218in}{2.198778in}}%
\pgfpathlineto{\pgfqpoint{1.112371in}{2.196364in}}%
\pgfpathlineto{\pgfqpoint{1.114524in}{2.193950in}}%
\pgfpathlineto{\pgfqpoint{1.116677in}{2.191536in}}%
\pgfpathlineto{\pgfqpoint{1.118830in}{2.189121in}}%
\pgfpathlineto{\pgfqpoint{1.120984in}{2.186707in}}%
\pgfpathlineto{\pgfqpoint{1.123137in}{2.184293in}}%
\pgfpathlineto{\pgfqpoint{1.125290in}{2.181878in}}%
\pgfpathlineto{\pgfqpoint{1.127443in}{2.179464in}}%
\pgfpathlineto{\pgfqpoint{1.129596in}{2.177050in}}%
\pgfpathlineto{\pgfqpoint{1.131749in}{2.174635in}}%
\pgfpathlineto{\pgfqpoint{1.133902in}{2.172221in}}%
\pgfpathlineto{\pgfqpoint{1.136055in}{2.169807in}}%
\pgfpathlineto{\pgfqpoint{1.138208in}{2.167393in}}%
\pgfpathlineto{\pgfqpoint{1.140361in}{2.164978in}}%
\pgfpathlineto{\pgfqpoint{1.142514in}{2.162564in}}%
\pgfpathlineto{\pgfqpoint{1.144667in}{2.160150in}}%
\pgfpathlineto{\pgfqpoint{1.146820in}{2.157735in}}%
\pgfpathlineto{\pgfqpoint{1.148974in}{2.155321in}}%
\pgfpathlineto{\pgfqpoint{1.151127in}{2.152907in}}%
\pgfpathlineto{\pgfqpoint{1.153280in}{2.150493in}}%
\pgfpathlineto{\pgfqpoint{1.155433in}{2.148078in}}%
\pgfpathlineto{\pgfqpoint{1.157586in}{2.145664in}}%
\pgfpathlineto{\pgfqpoint{1.159739in}{2.143250in}}%
\pgfpathlineto{\pgfqpoint{1.161892in}{2.140835in}}%
\pgfpathlineto{\pgfqpoint{1.164045in}{2.138421in}}%
\pgfpathlineto{\pgfqpoint{1.166198in}{2.136007in}}%
\pgfpathlineto{\pgfqpoint{1.168351in}{2.133593in}}%
\pgfpathlineto{\pgfqpoint{1.170504in}{2.131178in}}%
\pgfpathlineto{\pgfqpoint{1.172657in}{2.128764in}}%
\pgfpathlineto{\pgfqpoint{1.174810in}{2.126350in}}%
\pgfpathlineto{\pgfqpoint{1.176964in}{2.123935in}}%
\pgfpathlineto{\pgfqpoint{1.179117in}{2.121521in}}%
\pgfpathlineto{\pgfqpoint{1.181270in}{2.119107in}}%
\pgfpathlineto{\pgfqpoint{1.183423in}{2.116692in}}%
\pgfpathlineto{\pgfqpoint{1.185576in}{2.114278in}}%
\pgfpathlineto{\pgfqpoint{1.187729in}{2.111864in}}%
\pgfpathlineto{\pgfqpoint{1.189882in}{2.109450in}}%
\pgfpathlineto{\pgfqpoint{1.192035in}{2.107035in}}%
\pgfpathlineto{\pgfqpoint{1.194188in}{2.104621in}}%
\pgfpathlineto{\pgfqpoint{1.196341in}{2.102207in}}%
\pgfpathlineto{\pgfqpoint{1.198494in}{2.099792in}}%
\pgfpathlineto{\pgfqpoint{1.200647in}{2.097378in}}%
\pgfpathlineto{\pgfqpoint{1.202800in}{2.094964in}}%
\pgfpathlineto{\pgfqpoint{1.204954in}{2.092550in}}%
\pgfpathlineto{\pgfqpoint{1.207107in}{2.090135in}}%
\pgfpathlineto{\pgfqpoint{1.209260in}{2.087721in}}%
\pgfpathlineto{\pgfqpoint{1.211413in}{2.085307in}}%
\pgfpathlineto{\pgfqpoint{1.213566in}{2.082892in}}%
\pgfpathlineto{\pgfqpoint{1.215719in}{2.080478in}}%
\pgfpathlineto{\pgfqpoint{1.217872in}{2.078064in}}%
\pgfpathlineto{\pgfqpoint{1.220025in}{2.075650in}}%
\pgfpathlineto{\pgfqpoint{1.222178in}{2.073235in}}%
\pgfpathlineto{\pgfqpoint{1.224331in}{2.070821in}}%
\pgfpathlineto{\pgfqpoint{1.226484in}{2.068407in}}%
\pgfpathlineto{\pgfqpoint{1.228637in}{2.065992in}}%
\pgfpathlineto{\pgfqpoint{1.230790in}{2.063578in}}%
\pgfpathlineto{\pgfqpoint{1.232943in}{2.061164in}}%
\pgfpathlineto{\pgfqpoint{1.235097in}{2.060627in}}%
\pgfpathlineto{\pgfqpoint{1.237250in}{2.063042in}}%
\pgfpathlineto{\pgfqpoint{1.239403in}{2.065456in}}%
\pgfpathlineto{\pgfqpoint{1.241556in}{2.067870in}}%
\pgfpathlineto{\pgfqpoint{1.243709in}{2.070284in}}%
\pgfpathlineto{\pgfqpoint{1.245862in}{2.072699in}}%
\pgfpathlineto{\pgfqpoint{1.248015in}{2.075113in}}%
\pgfpathlineto{\pgfqpoint{1.250168in}{2.077527in}}%
\pgfpathlineto{\pgfqpoint{1.252321in}{2.079942in}}%
\pgfpathlineto{\pgfqpoint{1.254474in}{2.082356in}}%
\pgfpathlineto{\pgfqpoint{1.256627in}{2.084770in}}%
\pgfpathlineto{\pgfqpoint{1.258780in}{2.087184in}}%
\pgfpathlineto{\pgfqpoint{1.260933in}{2.089599in}}%
\pgfpathlineto{\pgfqpoint{1.263087in}{2.092013in}}%
\pgfpathlineto{\pgfqpoint{1.265240in}{2.094427in}}%
\pgfpathlineto{\pgfqpoint{1.267393in}{2.096842in}}%
\pgfpathlineto{\pgfqpoint{1.269546in}{2.099256in}}%
\pgfpathlineto{\pgfqpoint{1.271699in}{2.101670in}}%
\pgfpathlineto{\pgfqpoint{1.273852in}{2.104085in}}%
\pgfpathlineto{\pgfqpoint{1.276005in}{2.106499in}}%
\pgfpathlineto{\pgfqpoint{1.278158in}{2.108913in}}%
\pgfpathlineto{\pgfqpoint{1.280311in}{2.111327in}}%
\pgfpathlineto{\pgfqpoint{1.282464in}{2.113742in}}%
\pgfpathlineto{\pgfqpoint{1.284617in}{2.116156in}}%
\pgfpathlineto{\pgfqpoint{1.286770in}{2.118570in}}%
\pgfpathlineto{\pgfqpoint{1.288923in}{2.120985in}}%
\pgfpathlineto{\pgfqpoint{1.291077in}{2.123399in}}%
\pgfpathlineto{\pgfqpoint{1.293230in}{2.125813in}}%
\pgfpathlineto{\pgfqpoint{1.295383in}{2.128227in}}%
\pgfpathlineto{\pgfqpoint{1.297536in}{2.130642in}}%
\pgfpathlineto{\pgfqpoint{1.299689in}{2.133056in}}%
\pgfpathlineto{\pgfqpoint{1.301842in}{2.135470in}}%
\pgfpathlineto{\pgfqpoint{1.303995in}{2.137885in}}%
\pgfpathlineto{\pgfqpoint{1.306148in}{2.140299in}}%
\pgfpathlineto{\pgfqpoint{1.308301in}{2.142713in}}%
\pgfpathlineto{\pgfqpoint{1.310454in}{2.145127in}}%
\pgfpathlineto{\pgfqpoint{1.312607in}{2.147542in}}%
\pgfpathlineto{\pgfqpoint{1.314760in}{2.149956in}}%
\pgfpathlineto{\pgfqpoint{1.316913in}{2.152370in}}%
\pgfpathlineto{\pgfqpoint{1.319067in}{2.154785in}}%
\pgfpathlineto{\pgfqpoint{1.321220in}{2.157199in}}%
\pgfpathlineto{\pgfqpoint{1.323373in}{2.159613in}}%
\pgfpathlineto{\pgfqpoint{1.325526in}{2.162028in}}%
\pgfpathlineto{\pgfqpoint{1.327679in}{2.164442in}}%
\pgfpathlineto{\pgfqpoint{1.329832in}{2.166856in}}%
\pgfpathlineto{\pgfqpoint{1.331985in}{2.169270in}}%
\pgfpathlineto{\pgfqpoint{1.334138in}{2.171685in}}%
\pgfpathlineto{\pgfqpoint{1.336291in}{2.174099in}}%
\pgfpathlineto{\pgfqpoint{1.338444in}{2.176513in}}%
\pgfpathlineto{\pgfqpoint{1.340597in}{2.178928in}}%
\pgfpathlineto{\pgfqpoint{1.342750in}{2.181342in}}%
\pgfpathlineto{\pgfqpoint{1.344903in}{2.183756in}}%
\pgfpathlineto{\pgfqpoint{1.347057in}{2.186170in}}%
\pgfpathlineto{\pgfqpoint{1.349210in}{2.188585in}}%
\pgfpathlineto{\pgfqpoint{1.351363in}{2.190999in}}%
\pgfpathlineto{\pgfqpoint{1.353516in}{2.193413in}}%
\pgfpathlineto{\pgfqpoint{1.355669in}{2.195828in}}%
\pgfpathlineto{\pgfqpoint{1.357822in}{2.198242in}}%
\pgfpathlineto{\pgfqpoint{1.359975in}{2.200656in}}%
\pgfpathlineto{\pgfqpoint{1.362128in}{2.203070in}}%
\pgfpathlineto{\pgfqpoint{1.364281in}{2.205485in}}%
\pgfpathlineto{\pgfqpoint{1.366434in}{2.207899in}}%
\pgfpathlineto{\pgfqpoint{1.368587in}{2.210313in}}%
\pgfpathlineto{\pgfqpoint{1.370740in}{2.212728in}}%
\pgfpathlineto{\pgfqpoint{1.372893in}{2.215142in}}%
\pgfpathlineto{\pgfqpoint{1.375046in}{2.217556in}}%
\pgfpathlineto{\pgfqpoint{1.377200in}{2.219971in}}%
\pgfpathlineto{\pgfqpoint{1.379353in}{2.222385in}}%
\pgfpathlineto{\pgfqpoint{1.381506in}{2.224799in}}%
\pgfpathlineto{\pgfqpoint{1.383659in}{2.227213in}}%
\pgfpathlineto{\pgfqpoint{1.385812in}{2.229628in}}%
\pgfpathlineto{\pgfqpoint{1.387965in}{2.232042in}}%
\pgfpathlineto{\pgfqpoint{1.390118in}{2.234456in}}%
\pgfpathlineto{\pgfqpoint{1.392271in}{2.236871in}}%
\pgfpathlineto{\pgfqpoint{1.394424in}{2.239285in}}%
\pgfpathlineto{\pgfqpoint{1.396577in}{2.241699in}}%
\pgfpathlineto{\pgfqpoint{1.398730in}{2.244113in}}%
\pgfpathlineto{\pgfqpoint{1.400883in}{2.246528in}}%
\pgfpathlineto{\pgfqpoint{1.403036in}{2.248942in}}%
\pgfpathlineto{\pgfqpoint{1.405190in}{2.251356in}}%
\pgfpathlineto{\pgfqpoint{1.407343in}{2.253771in}}%
\pgfpathlineto{\pgfqpoint{1.409496in}{2.256185in}}%
\pgfpathlineto{\pgfqpoint{1.411649in}{2.258599in}}%
\pgfpathlineto{\pgfqpoint{1.413802in}{2.261013in}}%
\pgfpathlineto{\pgfqpoint{1.415955in}{2.263428in}}%
\pgfpathlineto{\pgfqpoint{1.418108in}{2.265842in}}%
\pgfpathlineto{\pgfqpoint{1.420261in}{2.268256in}}%
\pgfpathlineto{\pgfqpoint{1.422414in}{2.270671in}}%
\pgfpathlineto{\pgfqpoint{1.424567in}{2.273085in}}%
\pgfpathlineto{\pgfqpoint{1.426720in}{2.275499in}}%
\pgfpathlineto{\pgfqpoint{1.428873in}{2.277914in}}%
\pgfpathlineto{\pgfqpoint{1.431026in}{2.280328in}}%
\pgfpathlineto{\pgfqpoint{1.433180in}{2.282742in}}%
\pgfpathlineto{\pgfqpoint{1.435333in}{2.285156in}}%
\pgfpathlineto{\pgfqpoint{1.437486in}{2.287571in}}%
\pgfpathlineto{\pgfqpoint{1.439639in}{2.289985in}}%
\pgfpathlineto{\pgfqpoint{1.441792in}{2.292399in}}%
\pgfpathlineto{\pgfqpoint{1.443945in}{2.294814in}}%
\pgfpathlineto{\pgfqpoint{1.446098in}{2.297228in}}%
\pgfpathlineto{\pgfqpoint{1.448251in}{2.299642in}}%
\pgfpathlineto{\pgfqpoint{1.450404in}{2.302056in}}%
\pgfpathlineto{\pgfqpoint{1.452557in}{2.304471in}}%
\pgfpathlineto{\pgfqpoint{1.454710in}{2.306885in}}%
\pgfpathlineto{\pgfqpoint{1.456863in}{2.309299in}}%
\pgfpathlineto{\pgfqpoint{1.459016in}{2.311714in}}%
\pgfpathlineto{\pgfqpoint{1.461170in}{2.314128in}}%
\pgfpathlineto{\pgfqpoint{1.463323in}{2.316542in}}%
\pgfpathlineto{\pgfqpoint{1.465476in}{2.318956in}}%
\pgfpathlineto{\pgfqpoint{1.467629in}{2.321371in}}%
\pgfpathlineto{\pgfqpoint{1.469782in}{2.323785in}}%
\pgfpathlineto{\pgfqpoint{1.471935in}{2.326199in}}%
\pgfpathlineto{\pgfqpoint{1.474088in}{2.328614in}}%
\pgfpathlineto{\pgfqpoint{1.476241in}{2.331028in}}%
\pgfpathlineto{\pgfqpoint{1.478394in}{2.333442in}}%
\pgfpathlineto{\pgfqpoint{1.480547in}{2.335857in}}%
\pgfpathlineto{\pgfqpoint{1.482700in}{2.338271in}}%
\pgfpathlineto{\pgfqpoint{1.484853in}{2.340685in}}%
\pgfpathlineto{\pgfqpoint{1.487006in}{2.343099in}}%
\pgfpathlineto{\pgfqpoint{1.489160in}{2.345514in}}%
\pgfpathlineto{\pgfqpoint{1.491313in}{2.347928in}}%
\pgfpathlineto{\pgfqpoint{1.493466in}{2.350342in}}%
\pgfpathlineto{\pgfqpoint{1.495619in}{2.352757in}}%
\pgfpathlineto{\pgfqpoint{1.497772in}{2.355171in}}%
\pgfpathlineto{\pgfqpoint{1.499925in}{2.357585in}}%
\pgfpathlineto{\pgfqpoint{1.502078in}{2.359999in}}%
\pgfpathlineto{\pgfqpoint{1.504231in}{2.362414in}}%
\pgfpathlineto{\pgfqpoint{1.506384in}{2.364828in}}%
\pgfpathlineto{\pgfqpoint{1.508537in}{2.367242in}}%
\pgfpathlineto{\pgfqpoint{1.510690in}{2.369657in}}%
\pgfpathlineto{\pgfqpoint{1.512843in}{2.372071in}}%
\pgfpathlineto{\pgfqpoint{1.514996in}{2.374485in}}%
\pgfpathlineto{\pgfqpoint{1.517149in}{2.376899in}}%
\pgfpathlineto{\pgfqpoint{1.519303in}{2.379314in}}%
\pgfpathlineto{\pgfqpoint{1.521456in}{2.381728in}}%
\pgfpathlineto{\pgfqpoint{1.523609in}{2.384142in}}%
\pgfpathlineto{\pgfqpoint{1.525762in}{2.386557in}}%
\pgfpathlineto{\pgfqpoint{1.527915in}{2.388971in}}%
\pgfpathlineto{\pgfqpoint{1.530068in}{2.391385in}}%
\pgfpathlineto{\pgfqpoint{1.532221in}{2.393800in}}%
\pgfpathlineto{\pgfqpoint{1.534374in}{2.396214in}}%
\pgfpathlineto{\pgfqpoint{1.536527in}{2.398628in}}%
\pgfpathlineto{\pgfqpoint{1.538680in}{2.401042in}}%
\pgfpathlineto{\pgfqpoint{1.540833in}{2.403457in}}%
\pgfpathlineto{\pgfqpoint{1.542986in}{2.405871in}}%
\pgfpathlineto{\pgfqpoint{1.545139in}{2.408285in}}%
\pgfpathlineto{\pgfqpoint{1.547293in}{2.410700in}}%
\pgfpathlineto{\pgfqpoint{1.549446in}{2.413114in}}%
\pgfpathlineto{\pgfqpoint{1.551599in}{2.415528in}}%
\pgfpathlineto{\pgfqpoint{1.553752in}{2.417942in}}%
\pgfpathlineto{\pgfqpoint{1.555905in}{2.420357in}}%
\pgfpathlineto{\pgfqpoint{1.558058in}{2.422771in}}%
\pgfpathlineto{\pgfqpoint{1.560211in}{2.425185in}}%
\pgfpathlineto{\pgfqpoint{1.562364in}{2.427600in}}%
\pgfpathlineto{\pgfqpoint{1.564517in}{2.430014in}}%
\pgfpathlineto{\pgfqpoint{1.566670in}{2.432428in}}%
\pgfpathlineto{\pgfqpoint{1.568823in}{2.434842in}}%
\pgfpathlineto{\pgfqpoint{1.570976in}{2.437257in}}%
\pgfpathlineto{\pgfqpoint{1.573129in}{2.439671in}}%
\pgfpathlineto{\pgfqpoint{1.575283in}{2.442085in}}%
\pgfpathlineto{\pgfqpoint{1.577436in}{2.444500in}}%
\pgfpathlineto{\pgfqpoint{1.579589in}{2.446914in}}%
\pgfpathlineto{\pgfqpoint{1.581742in}{2.449328in}}%
\pgfpathlineto{\pgfqpoint{1.583895in}{2.451743in}}%
\pgfpathlineto{\pgfqpoint{1.586048in}{2.454157in}}%
\pgfpathlineto{\pgfqpoint{1.588201in}{2.456571in}}%
\pgfpathlineto{\pgfqpoint{1.590354in}{2.458985in}}%
\pgfpathlineto{\pgfqpoint{1.592507in}{2.461400in}}%
\pgfpathlineto{\pgfqpoint{1.594660in}{2.463814in}}%
\pgfpathlineto{\pgfqpoint{1.596813in}{2.466228in}}%
\pgfpathlineto{\pgfqpoint{1.598966in}{2.468643in}}%
\pgfpathlineto{\pgfqpoint{1.601119in}{2.471057in}}%
\pgfpathlineto{\pgfqpoint{1.603273in}{2.473471in}}%
\pgfpathlineto{\pgfqpoint{1.605426in}{2.475885in}}%
\pgfpathlineto{\pgfqpoint{1.607579in}{2.478300in}}%
\pgfpathlineto{\pgfqpoint{1.609732in}{2.480714in}}%
\pgfpathlineto{\pgfqpoint{1.611885in}{2.483128in}}%
\pgfpathlineto{\pgfqpoint{1.614038in}{2.485543in}}%
\pgfpathlineto{\pgfqpoint{1.616191in}{2.487957in}}%
\pgfpathlineto{\pgfqpoint{1.618344in}{2.490371in}}%
\pgfpathlineto{\pgfqpoint{1.620497in}{2.492785in}}%
\pgfpathlineto{\pgfqpoint{1.622650in}{2.495200in}}%
\pgfpathlineto{\pgfqpoint{1.624803in}{2.497614in}}%
\pgfpathlineto{\pgfqpoint{1.626956in}{2.500028in}}%
\pgfpathlineto{\pgfqpoint{1.629109in}{2.502443in}}%
\pgfpathlineto{\pgfqpoint{1.631263in}{2.504857in}}%
\pgfpathlineto{\pgfqpoint{1.633416in}{2.507271in}}%
\pgfpathlineto{\pgfqpoint{1.635569in}{2.509686in}}%
\pgfpathlineto{\pgfqpoint{1.637722in}{2.512100in}}%
\pgfpathlineto{\pgfqpoint{1.639875in}{2.514514in}}%
\pgfpathlineto{\pgfqpoint{1.642028in}{2.516928in}}%
\pgfpathlineto{\pgfqpoint{1.644181in}{2.519343in}}%
\pgfpathlineto{\pgfqpoint{1.646334in}{2.521757in}}%
\pgfpathlineto{\pgfqpoint{1.648487in}{2.524171in}}%
\pgfpathlineto{\pgfqpoint{1.650640in}{2.526586in}}%
\pgfpathlineto{\pgfqpoint{1.652793in}{2.529000in}}%
\pgfpathlineto{\pgfqpoint{1.654946in}{2.531414in}}%
\pgfpathlineto{\pgfqpoint{1.657099in}{2.533828in}}%
\pgfpathlineto{\pgfqpoint{1.659253in}{2.536243in}}%
\pgfpathlineto{\pgfqpoint{1.661406in}{2.538657in}}%
\pgfpathlineto{\pgfqpoint{1.663559in}{2.541071in}}%
\pgfpathlineto{\pgfqpoint{1.665712in}{2.543486in}}%
\pgfpathlineto{\pgfqpoint{1.667865in}{2.545900in}}%
\pgfpathlineto{\pgfqpoint{1.670018in}{2.548314in}}%
\pgfpathlineto{\pgfqpoint{1.672171in}{2.550728in}}%
\pgfpathlineto{\pgfqpoint{1.674324in}{2.553143in}}%
\pgfpathlineto{\pgfqpoint{1.676477in}{2.555557in}}%
\pgfpathlineto{\pgfqpoint{1.678630in}{2.557971in}}%
\pgfpathlineto{\pgfqpoint{1.680783in}{2.560386in}}%
\pgfpathlineto{\pgfqpoint{1.682936in}{2.562800in}}%
\pgfpathlineto{\pgfqpoint{1.685089in}{2.565214in}}%
\pgfpathlineto{\pgfqpoint{1.687242in}{2.567629in}}%
\pgfpathlineto{\pgfqpoint{1.689396in}{2.570043in}}%
\pgfpathlineto{\pgfqpoint{1.691549in}{2.572457in}}%
\pgfpathlineto{\pgfqpoint{1.693702in}{2.574871in}}%
\pgfpathlineto{\pgfqpoint{1.695855in}{2.577286in}}%
\pgfpathlineto{\pgfqpoint{1.698008in}{2.579700in}}%
\pgfpathlineto{\pgfqpoint{1.700161in}{2.582114in}}%
\pgfpathlineto{\pgfqpoint{1.702314in}{2.584529in}}%
\pgfpathlineto{\pgfqpoint{1.704467in}{2.586943in}}%
\pgfpathlineto{\pgfqpoint{1.706620in}{2.589357in}}%
\pgfpathlineto{\pgfqpoint{1.708773in}{2.591771in}}%
\pgfpathlineto{\pgfqpoint{1.710926in}{2.594186in}}%
\pgfpathlineto{\pgfqpoint{1.713079in}{2.596600in}}%
\pgfpathlineto{\pgfqpoint{1.715232in}{2.599014in}}%
\pgfpathlineto{\pgfqpoint{1.717386in}{2.601429in}}%
\pgfpathlineto{\pgfqpoint{1.719539in}{2.603843in}}%
\pgfpathlineto{\pgfqpoint{1.721692in}{2.606257in}}%
\pgfpathlineto{\pgfqpoint{1.723845in}{2.608671in}}%
\pgfpathlineto{\pgfqpoint{1.725998in}{2.611086in}}%
\pgfpathlineto{\pgfqpoint{1.728151in}{2.613500in}}%
\pgfpathlineto{\pgfqpoint{1.730304in}{2.615914in}}%
\pgfpathlineto{\pgfqpoint{1.732457in}{2.618329in}}%
\pgfpathlineto{\pgfqpoint{1.734610in}{2.620743in}}%
\pgfpathlineto{\pgfqpoint{1.736763in}{2.623157in}}%
\pgfpathlineto{\pgfqpoint{1.738916in}{2.625572in}}%
\pgfpathlineto{\pgfqpoint{1.741069in}{2.627986in}}%
\pgfpathlineto{\pgfqpoint{1.743222in}{2.630400in}}%
\pgfpathlineto{\pgfqpoint{1.745376in}{2.632814in}}%
\pgfpathlineto{\pgfqpoint{1.747529in}{2.635229in}}%
\pgfpathlineto{\pgfqpoint{1.749682in}{2.637643in}}%
\pgfpathlineto{\pgfqpoint{1.751835in}{2.640057in}}%
\pgfpathlineto{\pgfqpoint{1.753988in}{2.642472in}}%
\pgfpathlineto{\pgfqpoint{1.756141in}{2.644886in}}%
\pgfpathlineto{\pgfqpoint{1.758294in}{2.647300in}}%
\pgfpathlineto{\pgfqpoint{1.760447in}{2.649714in}}%
\pgfpathlineto{\pgfqpoint{1.762600in}{2.652129in}}%
\pgfpathlineto{\pgfqpoint{1.764753in}{2.654543in}}%
\pgfpathlineto{\pgfqpoint{1.766906in}{2.656957in}}%
\pgfpathlineto{\pgfqpoint{1.769059in}{2.659372in}}%
\pgfpathlineto{\pgfqpoint{1.771212in}{2.661786in}}%
\pgfpathlineto{\pgfqpoint{1.773366in}{2.664200in}}%
\pgfpathlineto{\pgfqpoint{1.775519in}{2.666614in}}%
\pgfpathlineto{\pgfqpoint{1.777672in}{2.669029in}}%
\pgfpathlineto{\pgfqpoint{1.779825in}{2.671443in}}%
\pgfpathlineto{\pgfqpoint{1.781978in}{2.673857in}}%
\pgfpathlineto{\pgfqpoint{1.784131in}{2.676272in}}%
\pgfpathlineto{\pgfqpoint{1.786284in}{2.678686in}}%
\pgfpathlineto{\pgfqpoint{1.788437in}{2.681100in}}%
\pgfpathlineto{\pgfqpoint{1.790590in}{2.683515in}}%
\pgfpathlineto{\pgfqpoint{1.792743in}{2.685929in}}%
\pgfpathlineto{\pgfqpoint{1.794896in}{2.688343in}}%
\pgfpathlineto{\pgfqpoint{1.797049in}{2.690757in}}%
\pgfpathlineto{\pgfqpoint{1.799202in}{2.693172in}}%
\pgfpathlineto{\pgfqpoint{1.801356in}{2.695586in}}%
\pgfpathlineto{\pgfqpoint{1.803509in}{2.698000in}}%
\pgfpathlineto{\pgfqpoint{1.805662in}{2.700415in}}%
\pgfpathlineto{\pgfqpoint{1.807815in}{2.702829in}}%
\pgfpathlineto{\pgfqpoint{1.809968in}{2.705243in}}%
\pgfpathlineto{\pgfqpoint{1.812121in}{2.707657in}}%
\pgfpathlineto{\pgfqpoint{1.814274in}{2.710072in}}%
\pgfpathlineto{\pgfqpoint{1.816427in}{2.712486in}}%
\pgfpathlineto{\pgfqpoint{1.818580in}{2.714900in}}%
\pgfpathlineto{\pgfqpoint{1.820733in}{2.717315in}}%
\pgfpathlineto{\pgfqpoint{1.822886in}{2.719729in}}%
\pgfpathlineto{\pgfqpoint{1.825039in}{2.722143in}}%
\pgfpathlineto{\pgfqpoint{1.827192in}{2.724557in}}%
\pgfpathlineto{\pgfqpoint{1.829345in}{2.726972in}}%
\pgfpathlineto{\pgfqpoint{1.831499in}{2.729386in}}%
\pgfpathlineto{\pgfqpoint{1.833652in}{2.731800in}}%
\pgfpathlineto{\pgfqpoint{1.835805in}{2.734215in}}%
\pgfpathlineto{\pgfqpoint{1.837958in}{2.736629in}}%
\pgfpathlineto{\pgfqpoint{1.840111in}{2.739043in}}%
\pgfpathlineto{\pgfqpoint{1.842264in}{2.741458in}}%
\pgfpathlineto{\pgfqpoint{1.844417in}{2.743872in}}%
\pgfpathlineto{\pgfqpoint{1.846570in}{2.746286in}}%
\pgfpathlineto{\pgfqpoint{1.848723in}{2.748700in}}%
\pgfpathlineto{\pgfqpoint{1.850876in}{2.751115in}}%
\pgfpathlineto{\pgfqpoint{1.853029in}{2.753529in}}%
\pgfpathlineto{\pgfqpoint{1.855182in}{2.755943in}}%
\pgfpathlineto{\pgfqpoint{1.857335in}{2.758358in}}%
\pgfpathlineto{\pgfqpoint{1.859489in}{2.760772in}}%
\pgfpathlineto{\pgfqpoint{1.861642in}{2.763186in}}%
\pgfpathlineto{\pgfqpoint{1.863795in}{2.765600in}}%
\pgfpathlineto{\pgfqpoint{1.865948in}{2.768015in}}%
\pgfpathlineto{\pgfqpoint{1.868101in}{2.770429in}}%
\pgfpathlineto{\pgfqpoint{1.870254in}{2.772843in}}%
\pgfpathlineto{\pgfqpoint{1.872407in}{2.775258in}}%
\pgfpathlineto{\pgfqpoint{1.874560in}{2.777672in}}%
\pgfpathlineto{\pgfqpoint{1.876713in}{2.780086in}}%
\pgfpathlineto{\pgfqpoint{1.878866in}{2.782500in}}%
\pgfpathlineto{\pgfqpoint{1.881019in}{2.784915in}}%
\pgfpathlineto{\pgfqpoint{1.883172in}{2.787329in}}%
\pgfpathlineto{\pgfqpoint{1.885325in}{2.789743in}}%
\pgfpathlineto{\pgfqpoint{1.887479in}{2.792158in}}%
\pgfpathlineto{\pgfqpoint{1.889632in}{2.794572in}}%
\pgfpathlineto{\pgfqpoint{1.891785in}{2.796986in}}%
\pgfpathlineto{\pgfqpoint{1.893938in}{2.799401in}}%
\pgfpathlineto{\pgfqpoint{1.896091in}{2.801815in}}%
\pgfpathlineto{\pgfqpoint{1.898244in}{2.804229in}}%
\pgfpathlineto{\pgfqpoint{1.900397in}{2.806643in}}%
\pgfpathlineto{\pgfqpoint{1.902550in}{2.809058in}}%
\pgfpathlineto{\pgfqpoint{1.904703in}{2.811472in}}%
\pgfpathlineto{\pgfqpoint{1.906856in}{2.813886in}}%
\pgfpathlineto{\pgfqpoint{1.909009in}{2.816301in}}%
\pgfpathlineto{\pgfqpoint{1.911162in}{2.818715in}}%
\pgfpathlineto{\pgfqpoint{1.913315in}{2.821129in}}%
\pgfpathlineto{\pgfqpoint{1.915469in}{2.823543in}}%
\pgfpathlineto{\pgfqpoint{1.917622in}{2.825958in}}%
\pgfpathlineto{\pgfqpoint{1.919775in}{2.828372in}}%
\pgfpathlineto{\pgfqpoint{1.921928in}{2.830786in}}%
\pgfpathlineto{\pgfqpoint{1.924081in}{2.833201in}}%
\pgfpathlineto{\pgfqpoint{1.926234in}{2.835615in}}%
\pgfpathlineto{\pgfqpoint{1.928387in}{2.838029in}}%
\pgfpathlineto{\pgfqpoint{1.930540in}{2.840444in}}%
\pgfpathlineto{\pgfqpoint{1.932693in}{2.842858in}}%
\pgfpathlineto{\pgfqpoint{1.934846in}{2.845272in}}%
\pgfpathlineto{\pgfqpoint{1.936999in}{2.847686in}}%
\pgfpathlineto{\pgfqpoint{1.939152in}{2.850101in}}%
\pgfpathlineto{\pgfqpoint{1.941305in}{2.852515in}}%
\pgfpathlineto{\pgfqpoint{1.943459in}{2.854929in}}%
\pgfpathlineto{\pgfqpoint{1.945612in}{2.857344in}}%
\pgfpathlineto{\pgfqpoint{1.947765in}{2.859758in}}%
\pgfpathlineto{\pgfqpoint{1.949918in}{2.862172in}}%
\pgfpathlineto{\pgfqpoint{1.952071in}{2.864586in}}%
\pgfpathlineto{\pgfqpoint{1.954224in}{2.867001in}}%
\pgfpathlineto{\pgfqpoint{1.956377in}{2.869415in}}%
\pgfpathlineto{\pgfqpoint{1.958530in}{2.871829in}}%
\pgfpathlineto{\pgfqpoint{1.960683in}{2.874244in}}%
\pgfpathlineto{\pgfqpoint{1.962836in}{2.876658in}}%
\pgfpathlineto{\pgfqpoint{1.964989in}{2.879072in}}%
\pgfpathlineto{\pgfqpoint{1.967142in}{2.881486in}}%
\pgfpathlineto{\pgfqpoint{1.969295in}{2.883901in}}%
\pgfpathlineto{\pgfqpoint{1.971448in}{2.886315in}}%
\pgfpathlineto{\pgfqpoint{1.973602in}{2.888729in}}%
\pgfpathlineto{\pgfqpoint{1.975755in}{2.891144in}}%
\pgfpathlineto{\pgfqpoint{1.977908in}{2.893558in}}%
\pgfpathlineto{\pgfqpoint{1.980061in}{2.895972in}}%
\pgfpathlineto{\pgfqpoint{1.982214in}{2.898387in}}%
\pgfpathlineto{\pgfqpoint{1.984367in}{2.900801in}}%
\pgfpathlineto{\pgfqpoint{1.986520in}{2.903215in}}%
\pgfpathlineto{\pgfqpoint{1.988673in}{2.905629in}}%
\pgfpathlineto{\pgfqpoint{1.990826in}{2.908044in}}%
\pgfpathlineto{\pgfqpoint{1.992979in}{2.910458in}}%
\pgfpathlineto{\pgfqpoint{1.995132in}{2.912872in}}%
\pgfpathlineto{\pgfqpoint{1.997285in}{2.915287in}}%
\pgfpathlineto{\pgfqpoint{1.999438in}{2.917701in}}%
\pgfpathlineto{\pgfqpoint{2.001592in}{2.920115in}}%
\pgfpathlineto{\pgfqpoint{2.003745in}{2.922529in}}%
\pgfpathlineto{\pgfqpoint{2.005898in}{2.924944in}}%
\pgfpathlineto{\pgfqpoint{2.008051in}{2.927358in}}%
\pgfpathlineto{\pgfqpoint{2.010204in}{2.929772in}}%
\pgfpathlineto{\pgfqpoint{2.012357in}{2.932187in}}%
\pgfpathlineto{\pgfqpoint{2.014510in}{2.934601in}}%
\pgfpathlineto{\pgfqpoint{2.016663in}{2.937015in}}%
\pgfpathlineto{\pgfqpoint{2.018816in}{2.939429in}}%
\pgfpathlineto{\pgfqpoint{2.020969in}{2.941844in}}%
\pgfpathlineto{\pgfqpoint{2.023122in}{2.944258in}}%
\pgfpathlineto{\pgfqpoint{2.025275in}{2.946672in}}%
\pgfpathlineto{\pgfqpoint{2.027428in}{2.949087in}}%
\pgfpathlineto{\pgfqpoint{2.029582in}{2.951501in}}%
\pgfpathlineto{\pgfqpoint{2.031735in}{2.953915in}}%
\pgfpathlineto{\pgfqpoint{2.033888in}{2.956330in}}%
\pgfpathlineto{\pgfqpoint{2.036041in}{2.958744in}}%
\pgfpathlineto{\pgfqpoint{2.038194in}{2.961158in}}%
\pgfpathlineto{\pgfqpoint{2.040347in}{2.963572in}}%
\pgfpathlineto{\pgfqpoint{2.042500in}{2.965987in}}%
\pgfpathlineto{\pgfqpoint{2.042500in}{3.216000in}}%
\pgfpathlineto{\pgfqpoint{2.042500in}{3.216000in}}%
\pgfpathlineto{\pgfqpoint{2.040347in}{3.216000in}}%
\pgfpathlineto{\pgfqpoint{2.038194in}{3.216000in}}%
\pgfpathlineto{\pgfqpoint{2.036041in}{3.216000in}}%
\pgfpathlineto{\pgfqpoint{2.033888in}{3.216000in}}%
\pgfpathlineto{\pgfqpoint{2.031735in}{3.216000in}}%
\pgfpathlineto{\pgfqpoint{2.029582in}{3.216000in}}%
\pgfpathlineto{\pgfqpoint{2.027428in}{3.216000in}}%
\pgfpathlineto{\pgfqpoint{2.025275in}{3.216000in}}%
\pgfpathlineto{\pgfqpoint{2.023122in}{3.216000in}}%
\pgfpathlineto{\pgfqpoint{2.020969in}{3.216000in}}%
\pgfpathlineto{\pgfqpoint{2.018816in}{3.216000in}}%
\pgfpathlineto{\pgfqpoint{2.016663in}{3.216000in}}%
\pgfpathlineto{\pgfqpoint{2.014510in}{3.216000in}}%
\pgfpathlineto{\pgfqpoint{2.012357in}{3.216000in}}%
\pgfpathlineto{\pgfqpoint{2.010204in}{3.216000in}}%
\pgfpathlineto{\pgfqpoint{2.008051in}{3.216000in}}%
\pgfpathlineto{\pgfqpoint{2.005898in}{3.216000in}}%
\pgfpathlineto{\pgfqpoint{2.003745in}{3.216000in}}%
\pgfpathlineto{\pgfqpoint{2.001592in}{3.216000in}}%
\pgfpathlineto{\pgfqpoint{1.999438in}{3.216000in}}%
\pgfpathlineto{\pgfqpoint{1.997285in}{3.216000in}}%
\pgfpathlineto{\pgfqpoint{1.995132in}{3.216000in}}%
\pgfpathlineto{\pgfqpoint{1.992979in}{3.216000in}}%
\pgfpathlineto{\pgfqpoint{1.990826in}{3.216000in}}%
\pgfpathlineto{\pgfqpoint{1.988673in}{3.216000in}}%
\pgfpathlineto{\pgfqpoint{1.986520in}{3.216000in}}%
\pgfpathlineto{\pgfqpoint{1.984367in}{3.216000in}}%
\pgfpathlineto{\pgfqpoint{1.982214in}{3.216000in}}%
\pgfpathlineto{\pgfqpoint{1.980061in}{3.216000in}}%
\pgfpathlineto{\pgfqpoint{1.977908in}{3.216000in}}%
\pgfpathlineto{\pgfqpoint{1.975755in}{3.216000in}}%
\pgfpathlineto{\pgfqpoint{1.973602in}{3.216000in}}%
\pgfpathlineto{\pgfqpoint{1.971448in}{3.216000in}}%
\pgfpathlineto{\pgfqpoint{1.969295in}{3.216000in}}%
\pgfpathlineto{\pgfqpoint{1.967142in}{3.216000in}}%
\pgfpathlineto{\pgfqpoint{1.964989in}{3.216000in}}%
\pgfpathlineto{\pgfqpoint{1.962836in}{3.216000in}}%
\pgfpathlineto{\pgfqpoint{1.960683in}{3.216000in}}%
\pgfpathlineto{\pgfqpoint{1.958530in}{3.216000in}}%
\pgfpathlineto{\pgfqpoint{1.956377in}{3.216000in}}%
\pgfpathlineto{\pgfqpoint{1.954224in}{3.216000in}}%
\pgfpathlineto{\pgfqpoint{1.952071in}{3.216000in}}%
\pgfpathlineto{\pgfqpoint{1.949918in}{3.216000in}}%
\pgfpathlineto{\pgfqpoint{1.947765in}{3.216000in}}%
\pgfpathlineto{\pgfqpoint{1.945612in}{3.216000in}}%
\pgfpathlineto{\pgfqpoint{1.943459in}{3.216000in}}%
\pgfpathlineto{\pgfqpoint{1.941305in}{3.216000in}}%
\pgfpathlineto{\pgfqpoint{1.939152in}{3.216000in}}%
\pgfpathlineto{\pgfqpoint{1.936999in}{3.216000in}}%
\pgfpathlineto{\pgfqpoint{1.934846in}{3.216000in}}%
\pgfpathlineto{\pgfqpoint{1.932693in}{3.216000in}}%
\pgfpathlineto{\pgfqpoint{1.930540in}{3.216000in}}%
\pgfpathlineto{\pgfqpoint{1.928387in}{3.216000in}}%
\pgfpathlineto{\pgfqpoint{1.926234in}{3.216000in}}%
\pgfpathlineto{\pgfqpoint{1.924081in}{3.216000in}}%
\pgfpathlineto{\pgfqpoint{1.921928in}{3.216000in}}%
\pgfpathlineto{\pgfqpoint{1.919775in}{3.216000in}}%
\pgfpathlineto{\pgfqpoint{1.917622in}{3.216000in}}%
\pgfpathlineto{\pgfqpoint{1.915469in}{3.216000in}}%
\pgfpathlineto{\pgfqpoint{1.913315in}{3.216000in}}%
\pgfpathlineto{\pgfqpoint{1.911162in}{3.216000in}}%
\pgfpathlineto{\pgfqpoint{1.909009in}{3.216000in}}%
\pgfpathlineto{\pgfqpoint{1.906856in}{3.216000in}}%
\pgfpathlineto{\pgfqpoint{1.904703in}{3.216000in}}%
\pgfpathlineto{\pgfqpoint{1.902550in}{3.216000in}}%
\pgfpathlineto{\pgfqpoint{1.900397in}{3.216000in}}%
\pgfpathlineto{\pgfqpoint{1.898244in}{3.216000in}}%
\pgfpathlineto{\pgfqpoint{1.896091in}{3.216000in}}%
\pgfpathlineto{\pgfqpoint{1.893938in}{3.216000in}}%
\pgfpathlineto{\pgfqpoint{1.891785in}{3.216000in}}%
\pgfpathlineto{\pgfqpoint{1.889632in}{3.216000in}}%
\pgfpathlineto{\pgfqpoint{1.887479in}{3.216000in}}%
\pgfpathlineto{\pgfqpoint{1.885325in}{3.216000in}}%
\pgfpathlineto{\pgfqpoint{1.883172in}{3.216000in}}%
\pgfpathlineto{\pgfqpoint{1.881019in}{3.216000in}}%
\pgfpathlineto{\pgfqpoint{1.878866in}{3.216000in}}%
\pgfpathlineto{\pgfqpoint{1.876713in}{3.216000in}}%
\pgfpathlineto{\pgfqpoint{1.874560in}{3.216000in}}%
\pgfpathlineto{\pgfqpoint{1.872407in}{3.216000in}}%
\pgfpathlineto{\pgfqpoint{1.870254in}{3.216000in}}%
\pgfpathlineto{\pgfqpoint{1.868101in}{3.216000in}}%
\pgfpathlineto{\pgfqpoint{1.865948in}{3.216000in}}%
\pgfpathlineto{\pgfqpoint{1.863795in}{3.216000in}}%
\pgfpathlineto{\pgfqpoint{1.861642in}{3.216000in}}%
\pgfpathlineto{\pgfqpoint{1.859489in}{3.216000in}}%
\pgfpathlineto{\pgfqpoint{1.857335in}{3.216000in}}%
\pgfpathlineto{\pgfqpoint{1.855182in}{3.216000in}}%
\pgfpathlineto{\pgfqpoint{1.853029in}{3.216000in}}%
\pgfpathlineto{\pgfqpoint{1.850876in}{3.216000in}}%
\pgfpathlineto{\pgfqpoint{1.848723in}{3.216000in}}%
\pgfpathlineto{\pgfqpoint{1.846570in}{3.216000in}}%
\pgfpathlineto{\pgfqpoint{1.844417in}{3.216000in}}%
\pgfpathlineto{\pgfqpoint{1.842264in}{3.216000in}}%
\pgfpathlineto{\pgfqpoint{1.840111in}{3.216000in}}%
\pgfpathlineto{\pgfqpoint{1.837958in}{3.216000in}}%
\pgfpathlineto{\pgfqpoint{1.835805in}{3.216000in}}%
\pgfpathlineto{\pgfqpoint{1.833652in}{3.216000in}}%
\pgfpathlineto{\pgfqpoint{1.831499in}{3.216000in}}%
\pgfpathlineto{\pgfqpoint{1.829345in}{3.216000in}}%
\pgfpathlineto{\pgfqpoint{1.827192in}{3.216000in}}%
\pgfpathlineto{\pgfqpoint{1.825039in}{3.216000in}}%
\pgfpathlineto{\pgfqpoint{1.822886in}{3.216000in}}%
\pgfpathlineto{\pgfqpoint{1.820733in}{3.216000in}}%
\pgfpathlineto{\pgfqpoint{1.818580in}{3.216000in}}%
\pgfpathlineto{\pgfqpoint{1.816427in}{3.216000in}}%
\pgfpathlineto{\pgfqpoint{1.814274in}{3.216000in}}%
\pgfpathlineto{\pgfqpoint{1.812121in}{3.216000in}}%
\pgfpathlineto{\pgfqpoint{1.809968in}{3.216000in}}%
\pgfpathlineto{\pgfqpoint{1.807815in}{3.216000in}}%
\pgfpathlineto{\pgfqpoint{1.805662in}{3.216000in}}%
\pgfpathlineto{\pgfqpoint{1.803509in}{3.216000in}}%
\pgfpathlineto{\pgfqpoint{1.801356in}{3.216000in}}%
\pgfpathlineto{\pgfqpoint{1.799202in}{3.216000in}}%
\pgfpathlineto{\pgfqpoint{1.797049in}{3.216000in}}%
\pgfpathlineto{\pgfqpoint{1.794896in}{3.216000in}}%
\pgfpathlineto{\pgfqpoint{1.792743in}{3.216000in}}%
\pgfpathlineto{\pgfqpoint{1.790590in}{3.216000in}}%
\pgfpathlineto{\pgfqpoint{1.788437in}{3.216000in}}%
\pgfpathlineto{\pgfqpoint{1.786284in}{3.216000in}}%
\pgfpathlineto{\pgfqpoint{1.784131in}{3.216000in}}%
\pgfpathlineto{\pgfqpoint{1.781978in}{3.216000in}}%
\pgfpathlineto{\pgfqpoint{1.779825in}{3.216000in}}%
\pgfpathlineto{\pgfqpoint{1.777672in}{3.216000in}}%
\pgfpathlineto{\pgfqpoint{1.775519in}{3.216000in}}%
\pgfpathlineto{\pgfqpoint{1.773366in}{3.216000in}}%
\pgfpathlineto{\pgfqpoint{1.771212in}{3.216000in}}%
\pgfpathlineto{\pgfqpoint{1.769059in}{3.216000in}}%
\pgfpathlineto{\pgfqpoint{1.766906in}{3.216000in}}%
\pgfpathlineto{\pgfqpoint{1.764753in}{3.216000in}}%
\pgfpathlineto{\pgfqpoint{1.762600in}{3.216000in}}%
\pgfpathlineto{\pgfqpoint{1.760447in}{3.216000in}}%
\pgfpathlineto{\pgfqpoint{1.758294in}{3.216000in}}%
\pgfpathlineto{\pgfqpoint{1.756141in}{3.216000in}}%
\pgfpathlineto{\pgfqpoint{1.753988in}{3.216000in}}%
\pgfpathlineto{\pgfqpoint{1.751835in}{3.216000in}}%
\pgfpathlineto{\pgfqpoint{1.749682in}{3.216000in}}%
\pgfpathlineto{\pgfqpoint{1.747529in}{3.216000in}}%
\pgfpathlineto{\pgfqpoint{1.745376in}{3.216000in}}%
\pgfpathlineto{\pgfqpoint{1.743222in}{3.216000in}}%
\pgfpathlineto{\pgfqpoint{1.741069in}{3.216000in}}%
\pgfpathlineto{\pgfqpoint{1.738916in}{3.216000in}}%
\pgfpathlineto{\pgfqpoint{1.736763in}{3.216000in}}%
\pgfpathlineto{\pgfqpoint{1.734610in}{3.216000in}}%
\pgfpathlineto{\pgfqpoint{1.732457in}{3.216000in}}%
\pgfpathlineto{\pgfqpoint{1.730304in}{3.216000in}}%
\pgfpathlineto{\pgfqpoint{1.728151in}{3.216000in}}%
\pgfpathlineto{\pgfqpoint{1.725998in}{3.216000in}}%
\pgfpathlineto{\pgfqpoint{1.723845in}{3.216000in}}%
\pgfpathlineto{\pgfqpoint{1.721692in}{3.216000in}}%
\pgfpathlineto{\pgfqpoint{1.719539in}{3.216000in}}%
\pgfpathlineto{\pgfqpoint{1.717386in}{3.216000in}}%
\pgfpathlineto{\pgfqpoint{1.715232in}{3.216000in}}%
\pgfpathlineto{\pgfqpoint{1.713079in}{3.216000in}}%
\pgfpathlineto{\pgfqpoint{1.710926in}{3.216000in}}%
\pgfpathlineto{\pgfqpoint{1.708773in}{3.216000in}}%
\pgfpathlineto{\pgfqpoint{1.706620in}{3.216000in}}%
\pgfpathlineto{\pgfqpoint{1.704467in}{3.216000in}}%
\pgfpathlineto{\pgfqpoint{1.702314in}{3.216000in}}%
\pgfpathlineto{\pgfqpoint{1.700161in}{3.216000in}}%
\pgfpathlineto{\pgfqpoint{1.698008in}{3.216000in}}%
\pgfpathlineto{\pgfqpoint{1.695855in}{3.216000in}}%
\pgfpathlineto{\pgfqpoint{1.693702in}{3.216000in}}%
\pgfpathlineto{\pgfqpoint{1.691549in}{3.216000in}}%
\pgfpathlineto{\pgfqpoint{1.689396in}{3.216000in}}%
\pgfpathlineto{\pgfqpoint{1.687242in}{3.216000in}}%
\pgfpathlineto{\pgfqpoint{1.685089in}{3.216000in}}%
\pgfpathlineto{\pgfqpoint{1.682936in}{3.216000in}}%
\pgfpathlineto{\pgfqpoint{1.680783in}{3.216000in}}%
\pgfpathlineto{\pgfqpoint{1.678630in}{3.216000in}}%
\pgfpathlineto{\pgfqpoint{1.676477in}{3.216000in}}%
\pgfpathlineto{\pgfqpoint{1.674324in}{3.216000in}}%
\pgfpathlineto{\pgfqpoint{1.672171in}{3.216000in}}%
\pgfpathlineto{\pgfqpoint{1.670018in}{3.216000in}}%
\pgfpathlineto{\pgfqpoint{1.667865in}{3.216000in}}%
\pgfpathlineto{\pgfqpoint{1.665712in}{3.216000in}}%
\pgfpathlineto{\pgfqpoint{1.663559in}{3.216000in}}%
\pgfpathlineto{\pgfqpoint{1.661406in}{3.216000in}}%
\pgfpathlineto{\pgfqpoint{1.659253in}{3.216000in}}%
\pgfpathlineto{\pgfqpoint{1.657099in}{3.216000in}}%
\pgfpathlineto{\pgfqpoint{1.654946in}{3.216000in}}%
\pgfpathlineto{\pgfqpoint{1.652793in}{3.216000in}}%
\pgfpathlineto{\pgfqpoint{1.650640in}{3.216000in}}%
\pgfpathlineto{\pgfqpoint{1.648487in}{3.216000in}}%
\pgfpathlineto{\pgfqpoint{1.646334in}{3.216000in}}%
\pgfpathlineto{\pgfqpoint{1.644181in}{3.216000in}}%
\pgfpathlineto{\pgfqpoint{1.642028in}{3.216000in}}%
\pgfpathlineto{\pgfqpoint{1.639875in}{3.216000in}}%
\pgfpathlineto{\pgfqpoint{1.637722in}{3.216000in}}%
\pgfpathlineto{\pgfqpoint{1.635569in}{3.216000in}}%
\pgfpathlineto{\pgfqpoint{1.633416in}{3.216000in}}%
\pgfpathlineto{\pgfqpoint{1.631263in}{3.216000in}}%
\pgfpathlineto{\pgfqpoint{1.629109in}{3.216000in}}%
\pgfpathlineto{\pgfqpoint{1.626956in}{3.216000in}}%
\pgfpathlineto{\pgfqpoint{1.624803in}{3.216000in}}%
\pgfpathlineto{\pgfqpoint{1.622650in}{3.216000in}}%
\pgfpathlineto{\pgfqpoint{1.620497in}{3.216000in}}%
\pgfpathlineto{\pgfqpoint{1.618344in}{3.216000in}}%
\pgfpathlineto{\pgfqpoint{1.616191in}{3.216000in}}%
\pgfpathlineto{\pgfqpoint{1.614038in}{3.216000in}}%
\pgfpathlineto{\pgfqpoint{1.611885in}{3.216000in}}%
\pgfpathlineto{\pgfqpoint{1.609732in}{3.216000in}}%
\pgfpathlineto{\pgfqpoint{1.607579in}{3.216000in}}%
\pgfpathlineto{\pgfqpoint{1.605426in}{3.216000in}}%
\pgfpathlineto{\pgfqpoint{1.603273in}{3.216000in}}%
\pgfpathlineto{\pgfqpoint{1.601119in}{3.216000in}}%
\pgfpathlineto{\pgfqpoint{1.598966in}{3.216000in}}%
\pgfpathlineto{\pgfqpoint{1.596813in}{3.216000in}}%
\pgfpathlineto{\pgfqpoint{1.594660in}{3.216000in}}%
\pgfpathlineto{\pgfqpoint{1.592507in}{3.216000in}}%
\pgfpathlineto{\pgfqpoint{1.590354in}{3.216000in}}%
\pgfpathlineto{\pgfqpoint{1.588201in}{3.216000in}}%
\pgfpathlineto{\pgfqpoint{1.586048in}{3.216000in}}%
\pgfpathlineto{\pgfqpoint{1.583895in}{3.216000in}}%
\pgfpathlineto{\pgfqpoint{1.581742in}{3.216000in}}%
\pgfpathlineto{\pgfqpoint{1.579589in}{3.216000in}}%
\pgfpathlineto{\pgfqpoint{1.577436in}{3.216000in}}%
\pgfpathlineto{\pgfqpoint{1.575283in}{3.216000in}}%
\pgfpathlineto{\pgfqpoint{1.573129in}{3.216000in}}%
\pgfpathlineto{\pgfqpoint{1.570976in}{3.216000in}}%
\pgfpathlineto{\pgfqpoint{1.568823in}{3.216000in}}%
\pgfpathlineto{\pgfqpoint{1.566670in}{3.216000in}}%
\pgfpathlineto{\pgfqpoint{1.564517in}{3.216000in}}%
\pgfpathlineto{\pgfqpoint{1.562364in}{3.216000in}}%
\pgfpathlineto{\pgfqpoint{1.560211in}{3.216000in}}%
\pgfpathlineto{\pgfqpoint{1.558058in}{3.216000in}}%
\pgfpathlineto{\pgfqpoint{1.555905in}{3.216000in}}%
\pgfpathlineto{\pgfqpoint{1.553752in}{3.216000in}}%
\pgfpathlineto{\pgfqpoint{1.551599in}{3.216000in}}%
\pgfpathlineto{\pgfqpoint{1.549446in}{3.216000in}}%
\pgfpathlineto{\pgfqpoint{1.547293in}{3.216000in}}%
\pgfpathlineto{\pgfqpoint{1.545139in}{3.216000in}}%
\pgfpathlineto{\pgfqpoint{1.542986in}{3.216000in}}%
\pgfpathlineto{\pgfqpoint{1.540833in}{3.216000in}}%
\pgfpathlineto{\pgfqpoint{1.538680in}{3.216000in}}%
\pgfpathlineto{\pgfqpoint{1.536527in}{3.216000in}}%
\pgfpathlineto{\pgfqpoint{1.534374in}{3.216000in}}%
\pgfpathlineto{\pgfqpoint{1.532221in}{3.216000in}}%
\pgfpathlineto{\pgfqpoint{1.530068in}{3.216000in}}%
\pgfpathlineto{\pgfqpoint{1.527915in}{3.216000in}}%
\pgfpathlineto{\pgfqpoint{1.525762in}{3.216000in}}%
\pgfpathlineto{\pgfqpoint{1.523609in}{3.216000in}}%
\pgfpathlineto{\pgfqpoint{1.521456in}{3.216000in}}%
\pgfpathlineto{\pgfqpoint{1.519303in}{3.216000in}}%
\pgfpathlineto{\pgfqpoint{1.517149in}{3.216000in}}%
\pgfpathlineto{\pgfqpoint{1.514996in}{3.216000in}}%
\pgfpathlineto{\pgfqpoint{1.512843in}{3.216000in}}%
\pgfpathlineto{\pgfqpoint{1.510690in}{3.216000in}}%
\pgfpathlineto{\pgfqpoint{1.508537in}{3.216000in}}%
\pgfpathlineto{\pgfqpoint{1.506384in}{3.216000in}}%
\pgfpathlineto{\pgfqpoint{1.504231in}{3.216000in}}%
\pgfpathlineto{\pgfqpoint{1.502078in}{3.216000in}}%
\pgfpathlineto{\pgfqpoint{1.499925in}{3.216000in}}%
\pgfpathlineto{\pgfqpoint{1.497772in}{3.216000in}}%
\pgfpathlineto{\pgfqpoint{1.495619in}{3.216000in}}%
\pgfpathlineto{\pgfqpoint{1.493466in}{3.216000in}}%
\pgfpathlineto{\pgfqpoint{1.491313in}{3.216000in}}%
\pgfpathlineto{\pgfqpoint{1.489160in}{3.216000in}}%
\pgfpathlineto{\pgfqpoint{1.487006in}{3.216000in}}%
\pgfpathlineto{\pgfqpoint{1.484853in}{3.216000in}}%
\pgfpathlineto{\pgfqpoint{1.482700in}{3.216000in}}%
\pgfpathlineto{\pgfqpoint{1.480547in}{3.216000in}}%
\pgfpathlineto{\pgfqpoint{1.478394in}{3.216000in}}%
\pgfpathlineto{\pgfqpoint{1.476241in}{3.216000in}}%
\pgfpathlineto{\pgfqpoint{1.474088in}{3.216000in}}%
\pgfpathlineto{\pgfqpoint{1.471935in}{3.216000in}}%
\pgfpathlineto{\pgfqpoint{1.469782in}{3.216000in}}%
\pgfpathlineto{\pgfqpoint{1.467629in}{3.216000in}}%
\pgfpathlineto{\pgfqpoint{1.465476in}{3.216000in}}%
\pgfpathlineto{\pgfqpoint{1.463323in}{3.216000in}}%
\pgfpathlineto{\pgfqpoint{1.461170in}{3.216000in}}%
\pgfpathlineto{\pgfqpoint{1.459016in}{3.216000in}}%
\pgfpathlineto{\pgfqpoint{1.456863in}{3.216000in}}%
\pgfpathlineto{\pgfqpoint{1.454710in}{3.216000in}}%
\pgfpathlineto{\pgfqpoint{1.452557in}{3.216000in}}%
\pgfpathlineto{\pgfqpoint{1.450404in}{3.216000in}}%
\pgfpathlineto{\pgfqpoint{1.448251in}{3.216000in}}%
\pgfpathlineto{\pgfqpoint{1.446098in}{3.216000in}}%
\pgfpathlineto{\pgfqpoint{1.443945in}{3.216000in}}%
\pgfpathlineto{\pgfqpoint{1.441792in}{3.216000in}}%
\pgfpathlineto{\pgfqpoint{1.439639in}{3.216000in}}%
\pgfpathlineto{\pgfqpoint{1.437486in}{3.216000in}}%
\pgfpathlineto{\pgfqpoint{1.435333in}{3.216000in}}%
\pgfpathlineto{\pgfqpoint{1.433180in}{3.216000in}}%
\pgfpathlineto{\pgfqpoint{1.431026in}{3.216000in}}%
\pgfpathlineto{\pgfqpoint{1.428873in}{3.216000in}}%
\pgfpathlineto{\pgfqpoint{1.426720in}{3.216000in}}%
\pgfpathlineto{\pgfqpoint{1.424567in}{3.216000in}}%
\pgfpathlineto{\pgfqpoint{1.422414in}{3.216000in}}%
\pgfpathlineto{\pgfqpoint{1.420261in}{3.216000in}}%
\pgfpathlineto{\pgfqpoint{1.418108in}{3.216000in}}%
\pgfpathlineto{\pgfqpoint{1.415955in}{3.216000in}}%
\pgfpathlineto{\pgfqpoint{1.413802in}{3.216000in}}%
\pgfpathlineto{\pgfqpoint{1.411649in}{3.216000in}}%
\pgfpathlineto{\pgfqpoint{1.409496in}{3.216000in}}%
\pgfpathlineto{\pgfqpoint{1.407343in}{3.216000in}}%
\pgfpathlineto{\pgfqpoint{1.405190in}{3.216000in}}%
\pgfpathlineto{\pgfqpoint{1.403036in}{3.216000in}}%
\pgfpathlineto{\pgfqpoint{1.400883in}{3.216000in}}%
\pgfpathlineto{\pgfqpoint{1.398730in}{3.216000in}}%
\pgfpathlineto{\pgfqpoint{1.396577in}{3.216000in}}%
\pgfpathlineto{\pgfqpoint{1.394424in}{3.216000in}}%
\pgfpathlineto{\pgfqpoint{1.392271in}{3.216000in}}%
\pgfpathlineto{\pgfqpoint{1.390118in}{3.216000in}}%
\pgfpathlineto{\pgfqpoint{1.387965in}{3.216000in}}%
\pgfpathlineto{\pgfqpoint{1.385812in}{3.216000in}}%
\pgfpathlineto{\pgfqpoint{1.383659in}{3.216000in}}%
\pgfpathlineto{\pgfqpoint{1.381506in}{3.216000in}}%
\pgfpathlineto{\pgfqpoint{1.379353in}{3.216000in}}%
\pgfpathlineto{\pgfqpoint{1.377200in}{3.216000in}}%
\pgfpathlineto{\pgfqpoint{1.375046in}{3.216000in}}%
\pgfpathlineto{\pgfqpoint{1.372893in}{3.216000in}}%
\pgfpathlineto{\pgfqpoint{1.370740in}{3.216000in}}%
\pgfpathlineto{\pgfqpoint{1.368587in}{3.216000in}}%
\pgfpathlineto{\pgfqpoint{1.366434in}{3.216000in}}%
\pgfpathlineto{\pgfqpoint{1.364281in}{3.216000in}}%
\pgfpathlineto{\pgfqpoint{1.362128in}{3.216000in}}%
\pgfpathlineto{\pgfqpoint{1.359975in}{3.216000in}}%
\pgfpathlineto{\pgfqpoint{1.357822in}{3.216000in}}%
\pgfpathlineto{\pgfqpoint{1.355669in}{3.216000in}}%
\pgfpathlineto{\pgfqpoint{1.353516in}{3.216000in}}%
\pgfpathlineto{\pgfqpoint{1.351363in}{3.216000in}}%
\pgfpathlineto{\pgfqpoint{1.349210in}{3.216000in}}%
\pgfpathlineto{\pgfqpoint{1.347057in}{3.216000in}}%
\pgfpathlineto{\pgfqpoint{1.344903in}{3.216000in}}%
\pgfpathlineto{\pgfqpoint{1.342750in}{3.216000in}}%
\pgfpathlineto{\pgfqpoint{1.340597in}{3.216000in}}%
\pgfpathlineto{\pgfqpoint{1.338444in}{3.216000in}}%
\pgfpathlineto{\pgfqpoint{1.336291in}{3.216000in}}%
\pgfpathlineto{\pgfqpoint{1.334138in}{3.216000in}}%
\pgfpathlineto{\pgfqpoint{1.331985in}{3.216000in}}%
\pgfpathlineto{\pgfqpoint{1.329832in}{3.216000in}}%
\pgfpathlineto{\pgfqpoint{1.327679in}{3.216000in}}%
\pgfpathlineto{\pgfqpoint{1.325526in}{3.216000in}}%
\pgfpathlineto{\pgfqpoint{1.323373in}{3.216000in}}%
\pgfpathlineto{\pgfqpoint{1.321220in}{3.216000in}}%
\pgfpathlineto{\pgfqpoint{1.319067in}{3.216000in}}%
\pgfpathlineto{\pgfqpoint{1.316913in}{3.216000in}}%
\pgfpathlineto{\pgfqpoint{1.314760in}{3.216000in}}%
\pgfpathlineto{\pgfqpoint{1.312607in}{3.216000in}}%
\pgfpathlineto{\pgfqpoint{1.310454in}{3.216000in}}%
\pgfpathlineto{\pgfqpoint{1.308301in}{3.216000in}}%
\pgfpathlineto{\pgfqpoint{1.306148in}{3.216000in}}%
\pgfpathlineto{\pgfqpoint{1.303995in}{3.216000in}}%
\pgfpathlineto{\pgfqpoint{1.301842in}{3.216000in}}%
\pgfpathlineto{\pgfqpoint{1.299689in}{3.216000in}}%
\pgfpathlineto{\pgfqpoint{1.297536in}{3.216000in}}%
\pgfpathlineto{\pgfqpoint{1.295383in}{3.216000in}}%
\pgfpathlineto{\pgfqpoint{1.293230in}{3.216000in}}%
\pgfpathlineto{\pgfqpoint{1.291077in}{3.216000in}}%
\pgfpathlineto{\pgfqpoint{1.288923in}{3.216000in}}%
\pgfpathlineto{\pgfqpoint{1.286770in}{3.216000in}}%
\pgfpathlineto{\pgfqpoint{1.284617in}{3.216000in}}%
\pgfpathlineto{\pgfqpoint{1.282464in}{3.216000in}}%
\pgfpathlineto{\pgfqpoint{1.280311in}{3.216000in}}%
\pgfpathlineto{\pgfqpoint{1.278158in}{3.216000in}}%
\pgfpathlineto{\pgfqpoint{1.276005in}{3.216000in}}%
\pgfpathlineto{\pgfqpoint{1.273852in}{3.216000in}}%
\pgfpathlineto{\pgfqpoint{1.271699in}{3.216000in}}%
\pgfpathlineto{\pgfqpoint{1.269546in}{3.216000in}}%
\pgfpathlineto{\pgfqpoint{1.267393in}{3.216000in}}%
\pgfpathlineto{\pgfqpoint{1.265240in}{3.216000in}}%
\pgfpathlineto{\pgfqpoint{1.263087in}{3.216000in}}%
\pgfpathlineto{\pgfqpoint{1.260933in}{3.216000in}}%
\pgfpathlineto{\pgfqpoint{1.258780in}{3.216000in}}%
\pgfpathlineto{\pgfqpoint{1.256627in}{3.216000in}}%
\pgfpathlineto{\pgfqpoint{1.254474in}{3.216000in}}%
\pgfpathlineto{\pgfqpoint{1.252321in}{3.216000in}}%
\pgfpathlineto{\pgfqpoint{1.250168in}{3.216000in}}%
\pgfpathlineto{\pgfqpoint{1.248015in}{3.216000in}}%
\pgfpathlineto{\pgfqpoint{1.245862in}{3.216000in}}%
\pgfpathlineto{\pgfqpoint{1.243709in}{3.216000in}}%
\pgfpathlineto{\pgfqpoint{1.241556in}{3.216000in}}%
\pgfpathlineto{\pgfqpoint{1.239403in}{3.216000in}}%
\pgfpathlineto{\pgfqpoint{1.237250in}{3.216000in}}%
\pgfpathlineto{\pgfqpoint{1.235097in}{3.216000in}}%
\pgfpathlineto{\pgfqpoint{1.232943in}{3.216000in}}%
\pgfpathlineto{\pgfqpoint{1.230790in}{3.216000in}}%
\pgfpathlineto{\pgfqpoint{1.228637in}{3.216000in}}%
\pgfpathlineto{\pgfqpoint{1.226484in}{3.216000in}}%
\pgfpathlineto{\pgfqpoint{1.224331in}{3.216000in}}%
\pgfpathlineto{\pgfqpoint{1.222178in}{3.216000in}}%
\pgfpathlineto{\pgfqpoint{1.220025in}{3.216000in}}%
\pgfpathlineto{\pgfqpoint{1.217872in}{3.216000in}}%
\pgfpathlineto{\pgfqpoint{1.215719in}{3.216000in}}%
\pgfpathlineto{\pgfqpoint{1.213566in}{3.216000in}}%
\pgfpathlineto{\pgfqpoint{1.211413in}{3.216000in}}%
\pgfpathlineto{\pgfqpoint{1.209260in}{3.216000in}}%
\pgfpathlineto{\pgfqpoint{1.207107in}{3.216000in}}%
\pgfpathlineto{\pgfqpoint{1.204954in}{3.216000in}}%
\pgfpathlineto{\pgfqpoint{1.202800in}{3.216000in}}%
\pgfpathlineto{\pgfqpoint{1.200647in}{3.216000in}}%
\pgfpathlineto{\pgfqpoint{1.198494in}{3.216000in}}%
\pgfpathlineto{\pgfqpoint{1.196341in}{3.216000in}}%
\pgfpathlineto{\pgfqpoint{1.194188in}{3.216000in}}%
\pgfpathlineto{\pgfqpoint{1.192035in}{3.216000in}}%
\pgfpathlineto{\pgfqpoint{1.189882in}{3.216000in}}%
\pgfpathlineto{\pgfqpoint{1.187729in}{3.216000in}}%
\pgfpathlineto{\pgfqpoint{1.185576in}{3.216000in}}%
\pgfpathlineto{\pgfqpoint{1.183423in}{3.216000in}}%
\pgfpathlineto{\pgfqpoint{1.181270in}{3.216000in}}%
\pgfpathlineto{\pgfqpoint{1.179117in}{3.216000in}}%
\pgfpathlineto{\pgfqpoint{1.176964in}{3.216000in}}%
\pgfpathlineto{\pgfqpoint{1.174810in}{3.216000in}}%
\pgfpathlineto{\pgfqpoint{1.172657in}{3.216000in}}%
\pgfpathlineto{\pgfqpoint{1.170504in}{3.216000in}}%
\pgfpathlineto{\pgfqpoint{1.168351in}{3.216000in}}%
\pgfpathlineto{\pgfqpoint{1.166198in}{3.216000in}}%
\pgfpathlineto{\pgfqpoint{1.164045in}{3.216000in}}%
\pgfpathlineto{\pgfqpoint{1.161892in}{3.216000in}}%
\pgfpathlineto{\pgfqpoint{1.159739in}{3.216000in}}%
\pgfpathlineto{\pgfqpoint{1.157586in}{3.216000in}}%
\pgfpathlineto{\pgfqpoint{1.155433in}{3.216000in}}%
\pgfpathlineto{\pgfqpoint{1.153280in}{3.216000in}}%
\pgfpathlineto{\pgfqpoint{1.151127in}{3.216000in}}%
\pgfpathlineto{\pgfqpoint{1.148974in}{3.216000in}}%
\pgfpathlineto{\pgfqpoint{1.146820in}{3.216000in}}%
\pgfpathlineto{\pgfqpoint{1.144667in}{3.216000in}}%
\pgfpathlineto{\pgfqpoint{1.142514in}{3.216000in}}%
\pgfpathlineto{\pgfqpoint{1.140361in}{3.216000in}}%
\pgfpathlineto{\pgfqpoint{1.138208in}{3.216000in}}%
\pgfpathlineto{\pgfqpoint{1.136055in}{3.216000in}}%
\pgfpathlineto{\pgfqpoint{1.133902in}{3.216000in}}%
\pgfpathlineto{\pgfqpoint{1.131749in}{3.216000in}}%
\pgfpathlineto{\pgfqpoint{1.129596in}{3.216000in}}%
\pgfpathlineto{\pgfqpoint{1.127443in}{3.216000in}}%
\pgfpathlineto{\pgfqpoint{1.125290in}{3.216000in}}%
\pgfpathlineto{\pgfqpoint{1.123137in}{3.216000in}}%
\pgfpathlineto{\pgfqpoint{1.120984in}{3.216000in}}%
\pgfpathlineto{\pgfqpoint{1.118830in}{3.216000in}}%
\pgfpathlineto{\pgfqpoint{1.116677in}{3.216000in}}%
\pgfpathlineto{\pgfqpoint{1.114524in}{3.216000in}}%
\pgfpathlineto{\pgfqpoint{1.112371in}{3.216000in}}%
\pgfpathlineto{\pgfqpoint{1.110218in}{3.216000in}}%
\pgfpathlineto{\pgfqpoint{1.108065in}{3.216000in}}%
\pgfpathlineto{\pgfqpoint{1.105912in}{3.216000in}}%
\pgfpathlineto{\pgfqpoint{1.103759in}{3.216000in}}%
\pgfpathlineto{\pgfqpoint{1.101606in}{3.216000in}}%
\pgfpathlineto{\pgfqpoint{1.099453in}{3.216000in}}%
\pgfpathlineto{\pgfqpoint{1.097300in}{3.216000in}}%
\pgfpathlineto{\pgfqpoint{1.095147in}{3.216000in}}%
\pgfpathlineto{\pgfqpoint{1.092994in}{3.216000in}}%
\pgfpathlineto{\pgfqpoint{1.090840in}{3.216000in}}%
\pgfpathlineto{\pgfqpoint{1.088687in}{3.216000in}}%
\pgfpathlineto{\pgfqpoint{1.086534in}{3.216000in}}%
\pgfpathlineto{\pgfqpoint{1.084381in}{3.216000in}}%
\pgfpathlineto{\pgfqpoint{1.082228in}{3.216000in}}%
\pgfpathlineto{\pgfqpoint{1.080075in}{3.216000in}}%
\pgfpathlineto{\pgfqpoint{1.077922in}{3.216000in}}%
\pgfpathlineto{\pgfqpoint{1.075769in}{3.216000in}}%
\pgfpathlineto{\pgfqpoint{1.073616in}{3.216000in}}%
\pgfpathlineto{\pgfqpoint{1.071463in}{3.216000in}}%
\pgfpathlineto{\pgfqpoint{1.069310in}{3.216000in}}%
\pgfpathlineto{\pgfqpoint{1.067157in}{3.216000in}}%
\pgfpathlineto{\pgfqpoint{1.065004in}{3.216000in}}%
\pgfpathlineto{\pgfqpoint{1.062851in}{3.216000in}}%
\pgfpathlineto{\pgfqpoint{1.060697in}{3.216000in}}%
\pgfpathlineto{\pgfqpoint{1.058544in}{3.216000in}}%
\pgfpathlineto{\pgfqpoint{1.056391in}{3.216000in}}%
\pgfpathlineto{\pgfqpoint{1.054238in}{3.216000in}}%
\pgfpathlineto{\pgfqpoint{1.052085in}{3.216000in}}%
\pgfpathlineto{\pgfqpoint{1.049932in}{3.216000in}}%
\pgfpathlineto{\pgfqpoint{1.047779in}{3.216000in}}%
\pgfpathlineto{\pgfqpoint{1.045626in}{3.216000in}}%
\pgfpathlineto{\pgfqpoint{1.043473in}{3.216000in}}%
\pgfpathlineto{\pgfqpoint{1.041320in}{3.216000in}}%
\pgfpathlineto{\pgfqpoint{1.039167in}{3.216000in}}%
\pgfpathlineto{\pgfqpoint{1.037014in}{3.216000in}}%
\pgfpathlineto{\pgfqpoint{1.034861in}{3.216000in}}%
\pgfpathlineto{\pgfqpoint{1.032707in}{3.216000in}}%
\pgfpathlineto{\pgfqpoint{1.030554in}{3.216000in}}%
\pgfpathlineto{\pgfqpoint{1.028401in}{3.216000in}}%
\pgfpathlineto{\pgfqpoint{1.026248in}{3.216000in}}%
\pgfpathlineto{\pgfqpoint{1.024095in}{3.216000in}}%
\pgfpathlineto{\pgfqpoint{1.021942in}{3.216000in}}%
\pgfpathlineto{\pgfqpoint{1.019789in}{3.216000in}}%
\pgfpathlineto{\pgfqpoint{1.017636in}{3.216000in}}%
\pgfpathlineto{\pgfqpoint{1.015483in}{3.216000in}}%
\pgfpathlineto{\pgfqpoint{1.013330in}{3.216000in}}%
\pgfpathlineto{\pgfqpoint{1.011177in}{3.216000in}}%
\pgfpathlineto{\pgfqpoint{1.009024in}{3.216000in}}%
\pgfpathlineto{\pgfqpoint{1.006871in}{3.216000in}}%
\pgfpathlineto{\pgfqpoint{1.004717in}{3.216000in}}%
\pgfpathlineto{\pgfqpoint{1.002564in}{3.216000in}}%
\pgfpathlineto{\pgfqpoint{1.000411in}{3.216000in}}%
\pgfpathlineto{\pgfqpoint{0.998258in}{3.216000in}}%
\pgfpathlineto{\pgfqpoint{0.996105in}{3.216000in}}%
\pgfpathlineto{\pgfqpoint{0.993952in}{3.216000in}}%
\pgfpathlineto{\pgfqpoint{0.991799in}{3.216000in}}%
\pgfpathlineto{\pgfqpoint{0.989646in}{3.216000in}}%
\pgfpathlineto{\pgfqpoint{0.987493in}{3.216000in}}%
\pgfpathlineto{\pgfqpoint{0.985340in}{3.216000in}}%
\pgfpathlineto{\pgfqpoint{0.983187in}{3.216000in}}%
\pgfpathlineto{\pgfqpoint{0.981034in}{3.216000in}}%
\pgfpathlineto{\pgfqpoint{0.978881in}{3.216000in}}%
\pgfpathlineto{\pgfqpoint{0.976727in}{3.216000in}}%
\pgfpathlineto{\pgfqpoint{0.974574in}{3.216000in}}%
\pgfpathlineto{\pgfqpoint{0.972421in}{3.216000in}}%
\pgfpathlineto{\pgfqpoint{0.970268in}{3.216000in}}%
\pgfpathlineto{\pgfqpoint{0.968115in}{3.216000in}}%
\pgfpathlineto{\pgfqpoint{0.965962in}{3.216000in}}%
\pgfpathlineto{\pgfqpoint{0.963809in}{3.216000in}}%
\pgfpathlineto{\pgfqpoint{0.961656in}{3.216000in}}%
\pgfpathlineto{\pgfqpoint{0.959503in}{3.216000in}}%
\pgfpathlineto{\pgfqpoint{0.957350in}{3.216000in}}%
\pgfpathlineto{\pgfqpoint{0.955197in}{3.216000in}}%
\pgfpathlineto{\pgfqpoint{0.953044in}{3.216000in}}%
\pgfpathlineto{\pgfqpoint{0.950891in}{3.216000in}}%
\pgfpathlineto{\pgfqpoint{0.948737in}{3.216000in}}%
\pgfpathlineto{\pgfqpoint{0.946584in}{3.216000in}}%
\pgfpathlineto{\pgfqpoint{0.944431in}{3.216000in}}%
\pgfpathlineto{\pgfqpoint{0.942278in}{3.216000in}}%
\pgfpathlineto{\pgfqpoint{0.940125in}{3.216000in}}%
\pgfpathlineto{\pgfqpoint{0.937972in}{3.216000in}}%
\pgfpathlineto{\pgfqpoint{0.935819in}{3.216000in}}%
\pgfpathlineto{\pgfqpoint{0.933666in}{3.216000in}}%
\pgfpathlineto{\pgfqpoint{0.931513in}{3.216000in}}%
\pgfpathlineto{\pgfqpoint{0.929360in}{3.216000in}}%
\pgfpathlineto{\pgfqpoint{0.927207in}{3.216000in}}%
\pgfpathlineto{\pgfqpoint{0.925054in}{3.216000in}}%
\pgfpathlineto{\pgfqpoint{0.922901in}{3.216000in}}%
\pgfpathlineto{\pgfqpoint{0.920747in}{3.216000in}}%
\pgfpathlineto{\pgfqpoint{0.918594in}{3.216000in}}%
\pgfpathlineto{\pgfqpoint{0.916441in}{3.216000in}}%
\pgfpathlineto{\pgfqpoint{0.914288in}{3.216000in}}%
\pgfpathlineto{\pgfqpoint{0.912135in}{3.216000in}}%
\pgfpathlineto{\pgfqpoint{0.909982in}{3.216000in}}%
\pgfpathlineto{\pgfqpoint{0.907829in}{3.216000in}}%
\pgfpathlineto{\pgfqpoint{0.905676in}{3.216000in}}%
\pgfpathlineto{\pgfqpoint{0.903523in}{3.216000in}}%
\pgfpathlineto{\pgfqpoint{0.901370in}{3.216000in}}%
\pgfpathlineto{\pgfqpoint{0.899217in}{3.216000in}}%
\pgfpathlineto{\pgfqpoint{0.897064in}{3.216000in}}%
\pgfpathlineto{\pgfqpoint{0.894911in}{3.216000in}}%
\pgfpathlineto{\pgfqpoint{0.892758in}{3.216000in}}%
\pgfpathlineto{\pgfqpoint{0.890604in}{3.216000in}}%
\pgfpathlineto{\pgfqpoint{0.888451in}{3.216000in}}%
\pgfpathlineto{\pgfqpoint{0.886298in}{3.216000in}}%
\pgfpathlineto{\pgfqpoint{0.884145in}{3.216000in}}%
\pgfpathlineto{\pgfqpoint{0.881992in}{3.216000in}}%
\pgfpathlineto{\pgfqpoint{0.879839in}{3.216000in}}%
\pgfpathlineto{\pgfqpoint{0.877686in}{3.216000in}}%
\pgfpathlineto{\pgfqpoint{0.875533in}{3.216000in}}%
\pgfpathlineto{\pgfqpoint{0.873380in}{3.216000in}}%
\pgfpathlineto{\pgfqpoint{0.871227in}{3.216000in}}%
\pgfpathlineto{\pgfqpoint{0.869074in}{3.216000in}}%
\pgfpathlineto{\pgfqpoint{0.866921in}{3.216000in}}%
\pgfpathlineto{\pgfqpoint{0.864768in}{3.216000in}}%
\pgfpathlineto{\pgfqpoint{0.862614in}{3.216000in}}%
\pgfpathlineto{\pgfqpoint{0.860461in}{3.216000in}}%
\pgfpathlineto{\pgfqpoint{0.858308in}{3.216000in}}%
\pgfpathlineto{\pgfqpoint{0.856155in}{3.216000in}}%
\pgfpathlineto{\pgfqpoint{0.854002in}{3.216000in}}%
\pgfpathlineto{\pgfqpoint{0.851849in}{3.216000in}}%
\pgfpathlineto{\pgfqpoint{0.849696in}{3.216000in}}%
\pgfpathlineto{\pgfqpoint{0.847543in}{3.216000in}}%
\pgfpathlineto{\pgfqpoint{0.845390in}{3.216000in}}%
\pgfpathlineto{\pgfqpoint{0.843237in}{3.216000in}}%
\pgfpathlineto{\pgfqpoint{0.841084in}{3.216000in}}%
\pgfpathlineto{\pgfqpoint{0.838931in}{3.216000in}}%
\pgfpathlineto{\pgfqpoint{0.836778in}{3.216000in}}%
\pgfpathlineto{\pgfqpoint{0.834624in}{3.216000in}}%
\pgfpathlineto{\pgfqpoint{0.832471in}{3.216000in}}%
\pgfpathlineto{\pgfqpoint{0.830318in}{3.216000in}}%
\pgfpathlineto{\pgfqpoint{0.828165in}{3.216000in}}%
\pgfpathlineto{\pgfqpoint{0.826012in}{3.216000in}}%
\pgfpathlineto{\pgfqpoint{0.823859in}{3.216000in}}%
\pgfpathlineto{\pgfqpoint{0.821706in}{3.216000in}}%
\pgfpathlineto{\pgfqpoint{0.819553in}{3.216000in}}%
\pgfpathlineto{\pgfqpoint{0.817400in}{3.216000in}}%
\pgfpathlineto{\pgfqpoint{0.815247in}{3.216000in}}%
\pgfpathlineto{\pgfqpoint{0.813094in}{3.216000in}}%
\pgfpathlineto{\pgfqpoint{0.810941in}{3.216000in}}%
\pgfpathlineto{\pgfqpoint{0.808788in}{3.216000in}}%
\pgfpathlineto{\pgfqpoint{0.806634in}{3.216000in}}%
\pgfpathlineto{\pgfqpoint{0.804481in}{3.216000in}}%
\pgfpathlineto{\pgfqpoint{0.802328in}{3.216000in}}%
\pgfpathlineto{\pgfqpoint{0.800175in}{3.216000in}}%
\pgfpathlineto{\pgfqpoint{0.798022in}{3.216000in}}%
\pgfpathlineto{\pgfqpoint{0.795869in}{3.216000in}}%
\pgfpathlineto{\pgfqpoint{0.793716in}{3.216000in}}%
\pgfpathlineto{\pgfqpoint{0.791563in}{3.216000in}}%
\pgfpathlineto{\pgfqpoint{0.789410in}{3.216000in}}%
\pgfpathlineto{\pgfqpoint{0.787257in}{3.216000in}}%
\pgfpathlineto{\pgfqpoint{0.785104in}{3.216000in}}%
\pgfpathlineto{\pgfqpoint{0.782951in}{3.216000in}}%
\pgfpathlineto{\pgfqpoint{0.780798in}{3.216000in}}%
\pgfpathlineto{\pgfqpoint{0.778644in}{3.216000in}}%
\pgfpathlineto{\pgfqpoint{0.776491in}{3.216000in}}%
\pgfpathlineto{\pgfqpoint{0.774338in}{3.216000in}}%
\pgfpathlineto{\pgfqpoint{0.772185in}{3.216000in}}%
\pgfpathlineto{\pgfqpoint{0.770032in}{3.216000in}}%
\pgfpathlineto{\pgfqpoint{0.767879in}{3.216000in}}%
\pgfpathlineto{\pgfqpoint{0.765726in}{3.216000in}}%
\pgfpathlineto{\pgfqpoint{0.763573in}{3.216000in}}%
\pgfpathlineto{\pgfqpoint{0.761420in}{3.216000in}}%
\pgfpathlineto{\pgfqpoint{0.759267in}{3.216000in}}%
\pgfpathlineto{\pgfqpoint{0.757114in}{3.216000in}}%
\pgfpathlineto{\pgfqpoint{0.754961in}{3.216000in}}%
\pgfpathlineto{\pgfqpoint{0.752808in}{3.216000in}}%
\pgfpathlineto{\pgfqpoint{0.750655in}{3.216000in}}%
\pgfpathlineto{\pgfqpoint{0.748501in}{3.216000in}}%
\pgfpathlineto{\pgfqpoint{0.746348in}{3.216000in}}%
\pgfpathlineto{\pgfqpoint{0.744195in}{3.216000in}}%
\pgfpathlineto{\pgfqpoint{0.742042in}{3.216000in}}%
\pgfpathlineto{\pgfqpoint{0.739889in}{3.216000in}}%
\pgfpathlineto{\pgfqpoint{0.737736in}{3.216000in}}%
\pgfpathlineto{\pgfqpoint{0.735583in}{3.216000in}}%
\pgfpathlineto{\pgfqpoint{0.733430in}{3.216000in}}%
\pgfpathlineto{\pgfqpoint{0.731277in}{3.216000in}}%
\pgfpathlineto{\pgfqpoint{0.729124in}{3.216000in}}%
\pgfpathlineto{\pgfqpoint{0.726971in}{3.216000in}}%
\pgfpathlineto{\pgfqpoint{0.724818in}{3.216000in}}%
\pgfpathlineto{\pgfqpoint{0.722665in}{3.216000in}}%
\pgfpathlineto{\pgfqpoint{0.720511in}{3.216000in}}%
\pgfpathlineto{\pgfqpoint{0.718358in}{3.216000in}}%
\pgfpathlineto{\pgfqpoint{0.716205in}{3.216000in}}%
\pgfpathlineto{\pgfqpoint{0.714052in}{3.216000in}}%
\pgfpathlineto{\pgfqpoint{0.711899in}{3.216000in}}%
\pgfpathlineto{\pgfqpoint{0.709746in}{3.216000in}}%
\pgfpathlineto{\pgfqpoint{0.707593in}{3.216000in}}%
\pgfpathlineto{\pgfqpoint{0.705440in}{3.216000in}}%
\pgfpathlineto{\pgfqpoint{0.703287in}{3.216000in}}%
\pgfpathlineto{\pgfqpoint{0.701134in}{3.216000in}}%
\pgfpathlineto{\pgfqpoint{0.698981in}{3.216000in}}%
\pgfpathlineto{\pgfqpoint{0.696828in}{3.216000in}}%
\pgfpathlineto{\pgfqpoint{0.694675in}{3.216000in}}%
\pgfpathlineto{\pgfqpoint{0.692521in}{3.216000in}}%
\pgfpathlineto{\pgfqpoint{0.690368in}{3.216000in}}%
\pgfpathlineto{\pgfqpoint{0.688215in}{3.216000in}}%
\pgfpathlineto{\pgfqpoint{0.686062in}{3.216000in}}%
\pgfpathlineto{\pgfqpoint{0.683909in}{3.216000in}}%
\pgfpathlineto{\pgfqpoint{0.681756in}{3.216000in}}%
\pgfpathlineto{\pgfqpoint{0.679603in}{3.216000in}}%
\pgfpathlineto{\pgfqpoint{0.677450in}{3.216000in}}%
\pgfpathlineto{\pgfqpoint{0.675297in}{3.216000in}}%
\pgfpathlineto{\pgfqpoint{0.673144in}{3.216000in}}%
\pgfpathlineto{\pgfqpoint{0.670991in}{3.216000in}}%
\pgfpathlineto{\pgfqpoint{0.668838in}{3.216000in}}%
\pgfpathlineto{\pgfqpoint{0.666685in}{3.216000in}}%
\pgfpathlineto{\pgfqpoint{0.664531in}{3.216000in}}%
\pgfpathlineto{\pgfqpoint{0.662378in}{3.216000in}}%
\pgfpathlineto{\pgfqpoint{0.660225in}{3.216000in}}%
\pgfpathlineto{\pgfqpoint{0.658072in}{3.216000in}}%
\pgfpathlineto{\pgfqpoint{0.655919in}{3.216000in}}%
\pgfpathlineto{\pgfqpoint{0.653766in}{3.216000in}}%
\pgfpathlineto{\pgfqpoint{0.651613in}{3.216000in}}%
\pgfpathlineto{\pgfqpoint{0.649460in}{3.216000in}}%
\pgfpathlineto{\pgfqpoint{0.647307in}{3.216000in}}%
\pgfpathlineto{\pgfqpoint{0.645154in}{3.216000in}}%
\pgfpathlineto{\pgfqpoint{0.643001in}{3.216000in}}%
\pgfpathlineto{\pgfqpoint{0.640848in}{3.216000in}}%
\pgfpathlineto{\pgfqpoint{0.638695in}{3.216000in}}%
\pgfpathlineto{\pgfqpoint{0.636541in}{3.216000in}}%
\pgfpathlineto{\pgfqpoint{0.634388in}{3.216000in}}%
\pgfpathlineto{\pgfqpoint{0.632235in}{3.216000in}}%
\pgfpathlineto{\pgfqpoint{0.630082in}{3.216000in}}%
\pgfpathlineto{\pgfqpoint{0.627929in}{3.216000in}}%
\pgfpathlineto{\pgfqpoint{0.625776in}{3.216000in}}%
\pgfpathlineto{\pgfqpoint{0.623623in}{3.216000in}}%
\pgfpathlineto{\pgfqpoint{0.621470in}{3.216000in}}%
\pgfpathlineto{\pgfqpoint{0.619317in}{3.216000in}}%
\pgfpathlineto{\pgfqpoint{0.617164in}{3.216000in}}%
\pgfpathlineto{\pgfqpoint{0.615011in}{3.216000in}}%
\pgfpathlineto{\pgfqpoint{0.612858in}{3.216000in}}%
\pgfpathlineto{\pgfqpoint{0.610705in}{3.216000in}}%
\pgfpathlineto{\pgfqpoint{0.608552in}{3.216000in}}%
\pgfpathlineto{\pgfqpoint{0.606398in}{3.216000in}}%
\pgfpathlineto{\pgfqpoint{0.604245in}{3.216000in}}%
\pgfpathlineto{\pgfqpoint{0.602092in}{3.216000in}}%
\pgfpathlineto{\pgfqpoint{0.599939in}{3.216000in}}%
\pgfpathlineto{\pgfqpoint{0.597786in}{3.216000in}}%
\pgfpathlineto{\pgfqpoint{0.595633in}{3.216000in}}%
\pgfpathlineto{\pgfqpoint{0.593480in}{3.216000in}}%
\pgfpathlineto{\pgfqpoint{0.591327in}{3.216000in}}%
\pgfpathlineto{\pgfqpoint{0.589174in}{3.216000in}}%
\pgfpathlineto{\pgfqpoint{0.587021in}{3.216000in}}%
\pgfpathlineto{\pgfqpoint{0.584868in}{3.216000in}}%
\pgfpathlineto{\pgfqpoint{0.582715in}{3.216000in}}%
\pgfpathlineto{\pgfqpoint{0.580562in}{3.216000in}}%
\pgfpathlineto{\pgfqpoint{0.578408in}{3.216000in}}%
\pgfpathlineto{\pgfqpoint{0.576255in}{3.216000in}}%
\pgfpathlineto{\pgfqpoint{0.574102in}{3.216000in}}%
\pgfpathlineto{\pgfqpoint{0.571949in}{3.216000in}}%
\pgfpathlineto{\pgfqpoint{0.569796in}{3.216000in}}%
\pgfpathlineto{\pgfqpoint{0.567643in}{3.216000in}}%
\pgfpathlineto{\pgfqpoint{0.565490in}{3.216000in}}%
\pgfpathlineto{\pgfqpoint{0.563337in}{3.216000in}}%
\pgfpathlineto{\pgfqpoint{0.561184in}{3.216000in}}%
\pgfpathlineto{\pgfqpoint{0.559031in}{3.216000in}}%
\pgfpathlineto{\pgfqpoint{0.556878in}{3.216000in}}%
\pgfpathlineto{\pgfqpoint{0.554725in}{3.216000in}}%
\pgfpathlineto{\pgfqpoint{0.552572in}{3.216000in}}%
\pgfpathlineto{\pgfqpoint{0.550418in}{3.216000in}}%
\pgfpathlineto{\pgfqpoint{0.548265in}{3.216000in}}%
\pgfpathlineto{\pgfqpoint{0.546112in}{3.216000in}}%
\pgfpathlineto{\pgfqpoint{0.543959in}{3.216000in}}%
\pgfpathlineto{\pgfqpoint{0.541806in}{3.216000in}}%
\pgfpathlineto{\pgfqpoint{0.539653in}{3.216000in}}%
\pgfpathlineto{\pgfqpoint{0.537500in}{3.216000in}}%
\pgfpathlineto{\pgfqpoint{0.537500in}{3.216000in}}%
\pgfpathclose%
\pgfusepath{stroke,fill}%
}%
\begin{pgfscope}%
\pgfsys@transformshift{0.000000in}{0.000000in}%
\pgfsys@useobject{currentmarker}{}%
\end{pgfscope}%
\end{pgfscope}%
\begin{pgfscope}%
\pgfpathrectangle{\pgfqpoint{0.537500in}{2.059688in}}{\pgfqpoint{1.505000in}{1.156312in}}%
\pgfusepath{clip}%
\pgfsetrectcap%
\pgfsetroundjoin%
\pgfsetlinewidth{0.903375pt}%
\definecolor{currentstroke}{rgb}{0.717647,0.717647,0.717647}%
\pgfsetstrokecolor{currentstroke}%
\pgfsetstrokeopacity{0.900000}%
\pgfsetdash{}{0pt}%
\pgfpathmoveto{\pgfqpoint{0.537500in}{2.850917in}}%
\pgfpathlineto{\pgfqpoint{0.686062in}{2.686976in}}%
\pgfpathlineto{\pgfqpoint{0.791563in}{2.571591in}}%
\pgfpathlineto{\pgfqpoint{0.871227in}{2.485526in}}%
\pgfpathlineto{\pgfqpoint{0.931513in}{2.421449in}}%
\pgfpathlineto{\pgfqpoint{0.978881in}{2.372146in}}%
\pgfpathlineto{\pgfqpoint{1.017636in}{2.332872in}}%
\pgfpathlineto{\pgfqpoint{1.047779in}{2.303310in}}%
\pgfpathlineto{\pgfqpoint{1.073616in}{2.278947in}}%
\pgfpathlineto{\pgfqpoint{1.095147in}{2.259587in}}%
\pgfpathlineto{\pgfqpoint{1.114524in}{2.243135in}}%
\pgfpathlineto{\pgfqpoint{1.131749in}{2.229510in}}%
\pgfpathlineto{\pgfqpoint{1.146820in}{2.218559in}}%
\pgfpathlineto{\pgfqpoint{1.159739in}{2.210052in}}%
\pgfpathlineto{\pgfqpoint{1.172657in}{2.202510in}}%
\pgfpathlineto{\pgfqpoint{1.183423in}{2.197079in}}%
\pgfpathlineto{\pgfqpoint{1.194188in}{2.192525in}}%
\pgfpathlineto{\pgfqpoint{1.202800in}{2.189577in}}%
\pgfpathlineto{\pgfqpoint{1.211413in}{2.187293in}}%
\pgfpathlineto{\pgfqpoint{1.220025in}{2.185710in}}%
\pgfpathlineto{\pgfqpoint{1.228637in}{2.184854in}}%
\pgfpathlineto{\pgfqpoint{1.237250in}{2.184740in}}%
\pgfpathlineto{\pgfqpoint{1.245862in}{2.185370in}}%
\pgfpathlineto{\pgfqpoint{1.254474in}{2.186734in}}%
\pgfpathlineto{\pgfqpoint{1.263087in}{2.188807in}}%
\pgfpathlineto{\pgfqpoint{1.271699in}{2.191556in}}%
\pgfpathlineto{\pgfqpoint{1.282464in}{2.195881in}}%
\pgfpathlineto{\pgfqpoint{1.293230in}{2.201107in}}%
\pgfpathlineto{\pgfqpoint{1.303995in}{2.207138in}}%
\pgfpathlineto{\pgfqpoint{1.316913in}{2.215305in}}%
\pgfpathlineto{\pgfqpoint{1.329832in}{2.224344in}}%
\pgfpathlineto{\pgfqpoint{1.344903in}{2.235812in}}%
\pgfpathlineto{\pgfqpoint{1.362128in}{2.249912in}}%
\pgfpathlineto{\pgfqpoint{1.381506in}{2.266783in}}%
\pgfpathlineto{\pgfqpoint{1.403036in}{2.286500in}}%
\pgfpathlineto{\pgfqpoint{1.428873in}{2.311182in}}%
\pgfpathlineto{\pgfqpoint{1.459016in}{2.341013in}}%
\pgfpathlineto{\pgfqpoint{1.495619in}{2.378304in}}%
\pgfpathlineto{\pgfqpoint{1.540833in}{2.425480in}}%
\pgfpathlineto{\pgfqpoint{1.596813in}{2.485013in}}%
\pgfpathlineto{\pgfqpoint{1.667865in}{2.561713in}}%
\pgfpathlineto{\pgfqpoint{1.760447in}{2.662811in}}%
\pgfpathlineto{\pgfqpoint{1.883172in}{2.797989in}}%
\pgfpathlineto{\pgfqpoint{2.042500in}{2.974567in}}%
\pgfpathlineto{\pgfqpoint{2.042500in}{2.974567in}}%
\pgfusepath{stroke}%
\end{pgfscope}%
\begin{pgfscope}%
\pgfpathrectangle{\pgfqpoint{0.537500in}{2.059688in}}{\pgfqpoint{1.505000in}{1.156312in}}%
\pgfusepath{clip}%
\pgfsetrectcap%
\pgfsetroundjoin%
\pgfsetlinewidth{0.903375pt}%
\definecolor{currentstroke}{rgb}{0.717647,0.717647,0.717647}%
\pgfsetstrokecolor{currentstroke}%
\pgfsetstrokeopacity{0.900000}%
\pgfsetdash{}{0pt}%
\pgfpathmoveto{\pgfqpoint{0.537500in}{2.871029in}}%
\pgfpathlineto{\pgfqpoint{0.630082in}{2.771610in}}%
\pgfpathlineto{\pgfqpoint{0.705440in}{2.691731in}}%
\pgfpathlineto{\pgfqpoint{0.767879in}{2.626563in}}%
\pgfpathlineto{\pgfqpoint{0.821706in}{2.571411in}}%
\pgfpathlineto{\pgfqpoint{0.866921in}{2.526081in}}%
\pgfpathlineto{\pgfqpoint{0.907829in}{2.486112in}}%
\pgfpathlineto{\pgfqpoint{0.942278in}{2.453453in}}%
\pgfpathlineto{\pgfqpoint{0.972421in}{2.425831in}}%
\pgfpathlineto{\pgfqpoint{1.000411in}{2.401178in}}%
\pgfpathlineto{\pgfqpoint{1.024095in}{2.381236in}}%
\pgfpathlineto{\pgfqpoint{1.045626in}{2.363986in}}%
\pgfpathlineto{\pgfqpoint{1.065004in}{2.349304in}}%
\pgfpathlineto{\pgfqpoint{1.084381in}{2.335553in}}%
\pgfpathlineto{\pgfqpoint{1.101606in}{2.324230in}}%
\pgfpathlineto{\pgfqpoint{1.116677in}{2.315110in}}%
\pgfpathlineto{\pgfqpoint{1.131749in}{2.306811in}}%
\pgfpathlineto{\pgfqpoint{1.146820in}{2.299417in}}%
\pgfpathlineto{\pgfqpoint{1.159739in}{2.293866in}}%
\pgfpathlineto{\pgfqpoint{1.172657in}{2.289097in}}%
\pgfpathlineto{\pgfqpoint{1.185576in}{2.285158in}}%
\pgfpathlineto{\pgfqpoint{1.198494in}{2.282096in}}%
\pgfpathlineto{\pgfqpoint{1.209260in}{2.280239in}}%
\pgfpathlineto{\pgfqpoint{1.220025in}{2.279032in}}%
\pgfpathlineto{\pgfqpoint{1.230790in}{2.278485in}}%
\pgfpathlineto{\pgfqpoint{1.241556in}{2.278603in}}%
\pgfpathlineto{\pgfqpoint{1.252321in}{2.279386in}}%
\pgfpathlineto{\pgfqpoint{1.263087in}{2.280825in}}%
\pgfpathlineto{\pgfqpoint{1.273852in}{2.282910in}}%
\pgfpathlineto{\pgfqpoint{1.286770in}{2.286236in}}%
\pgfpathlineto{\pgfqpoint{1.299689in}{2.290425in}}%
\pgfpathlineto{\pgfqpoint{1.312607in}{2.295432in}}%
\pgfpathlineto{\pgfqpoint{1.325526in}{2.301204in}}%
\pgfpathlineto{\pgfqpoint{1.340597in}{2.308836in}}%
\pgfpathlineto{\pgfqpoint{1.355669in}{2.317352in}}%
\pgfpathlineto{\pgfqpoint{1.370740in}{2.326667in}}%
\pgfpathlineto{\pgfqpoint{1.387965in}{2.338189in}}%
\pgfpathlineto{\pgfqpoint{1.407343in}{2.352134in}}%
\pgfpathlineto{\pgfqpoint{1.426720in}{2.366985in}}%
\pgfpathlineto{\pgfqpoint{1.448251in}{2.384395in}}%
\pgfpathlineto{\pgfqpoint{1.471935in}{2.404485in}}%
\pgfpathlineto{\pgfqpoint{1.497772in}{2.427338in}}%
\pgfpathlineto{\pgfqpoint{1.527915in}{2.455016in}}%
\pgfpathlineto{\pgfqpoint{1.562364in}{2.487725in}}%
\pgfpathlineto{\pgfqpoint{1.601119in}{2.525608in}}%
\pgfpathlineto{\pgfqpoint{1.644181in}{2.568745in}}%
\pgfpathlineto{\pgfqpoint{1.695855in}{2.621617in}}%
\pgfpathlineto{\pgfqpoint{1.753988in}{2.682178in}}%
\pgfpathlineto{\pgfqpoint{1.822886in}{2.755037in}}%
\pgfpathlineto{\pgfqpoint{1.904703in}{2.842654in}}%
\pgfpathlineto{\pgfqpoint{2.003745in}{2.949830in}}%
\pgfpathlineto{\pgfqpoint{2.042500in}{2.992015in}}%
\pgfpathlineto{\pgfqpoint{2.042500in}{2.992015in}}%
\pgfusepath{stroke}%
\end{pgfscope}%
\begin{pgfscope}%
\pgfpathrectangle{\pgfqpoint{0.537500in}{2.059688in}}{\pgfqpoint{1.505000in}{1.156312in}}%
\pgfusepath{clip}%
\pgfsetrectcap%
\pgfsetroundjoin%
\pgfsetlinewidth{0.903375pt}%
\definecolor{currentstroke}{rgb}{0.717647,0.717647,0.717647}%
\pgfsetstrokecolor{currentstroke}%
\pgfsetstrokeopacity{0.900000}%
\pgfsetdash{}{0pt}%
\pgfpathmoveto{\pgfqpoint{0.537500in}{2.913265in}}%
\pgfpathlineto{\pgfqpoint{0.602092in}{2.847510in}}%
\pgfpathlineto{\pgfqpoint{0.658072in}{2.791543in}}%
\pgfpathlineto{\pgfqpoint{0.707593in}{2.743019in}}%
\pgfpathlineto{\pgfqpoint{0.752808in}{2.699710in}}%
\pgfpathlineto{\pgfqpoint{0.793716in}{2.661522in}}%
\pgfpathlineto{\pgfqpoint{0.830318in}{2.628317in}}%
\pgfpathlineto{\pgfqpoint{0.864768in}{2.598054in}}%
\pgfpathlineto{\pgfqpoint{0.894911in}{2.572496in}}%
\pgfpathlineto{\pgfqpoint{0.922901in}{2.549659in}}%
\pgfpathlineto{\pgfqpoint{0.948737in}{2.529455in}}%
\pgfpathlineto{\pgfqpoint{0.974574in}{2.510209in}}%
\pgfpathlineto{\pgfqpoint{0.998258in}{2.493517in}}%
\pgfpathlineto{\pgfqpoint{1.019789in}{2.479226in}}%
\pgfpathlineto{\pgfqpoint{1.041320in}{2.465869in}}%
\pgfpathlineto{\pgfqpoint{1.060697in}{2.454724in}}%
\pgfpathlineto{\pgfqpoint{1.080075in}{2.444484in}}%
\pgfpathlineto{\pgfqpoint{1.097300in}{2.436202in}}%
\pgfpathlineto{\pgfqpoint{1.114524in}{2.428745in}}%
\pgfpathlineto{\pgfqpoint{1.131749in}{2.422165in}}%
\pgfpathlineto{\pgfqpoint{1.146820in}{2.417165in}}%
\pgfpathlineto{\pgfqpoint{1.161892in}{2.412904in}}%
\pgfpathlineto{\pgfqpoint{1.176964in}{2.409409in}}%
\pgfpathlineto{\pgfqpoint{1.192035in}{2.406702in}}%
\pgfpathlineto{\pgfqpoint{1.207107in}{2.404802in}}%
\pgfpathlineto{\pgfqpoint{1.222178in}{2.403724in}}%
\pgfpathlineto{\pgfqpoint{1.237250in}{2.403473in}}%
\pgfpathlineto{\pgfqpoint{1.252321in}{2.404053in}}%
\pgfpathlineto{\pgfqpoint{1.267393in}{2.405459in}}%
\pgfpathlineto{\pgfqpoint{1.282464in}{2.407680in}}%
\pgfpathlineto{\pgfqpoint{1.297536in}{2.410703in}}%
\pgfpathlineto{\pgfqpoint{1.312607in}{2.414505in}}%
\pgfpathlineto{\pgfqpoint{1.327679in}{2.419063in}}%
\pgfpathlineto{\pgfqpoint{1.342750in}{2.424347in}}%
\pgfpathlineto{\pgfqpoint{1.359975in}{2.431237in}}%
\pgfpathlineto{\pgfqpoint{1.377200in}{2.438986in}}%
\pgfpathlineto{\pgfqpoint{1.394424in}{2.447543in}}%
\pgfpathlineto{\pgfqpoint{1.413802in}{2.458071in}}%
\pgfpathlineto{\pgfqpoint{1.433180in}{2.469481in}}%
\pgfpathlineto{\pgfqpoint{1.452557in}{2.481701in}}%
\pgfpathlineto{\pgfqpoint{1.474088in}{2.496149in}}%
\pgfpathlineto{\pgfqpoint{1.497772in}{2.512995in}}%
\pgfpathlineto{\pgfqpoint{1.521456in}{2.530736in}}%
\pgfpathlineto{\pgfqpoint{1.547293in}{2.550998in}}%
\pgfpathlineto{\pgfqpoint{1.575283in}{2.573891in}}%
\pgfpathlineto{\pgfqpoint{1.605426in}{2.599500in}}%
\pgfpathlineto{\pgfqpoint{1.637722in}{2.627890in}}%
\pgfpathlineto{\pgfqpoint{1.674324in}{2.661082in}}%
\pgfpathlineto{\pgfqpoint{1.713079in}{2.697223in}}%
\pgfpathlineto{\pgfqpoint{1.756141in}{2.738388in}}%
\pgfpathlineto{\pgfqpoint{1.803509in}{2.784684in}}%
\pgfpathlineto{\pgfqpoint{1.857335in}{2.838351in}}%
\pgfpathlineto{\pgfqpoint{1.917622in}{2.899537in}}%
\pgfpathlineto{\pgfqpoint{1.984367in}{2.968339in}}%
\pgfpathlineto{\pgfqpoint{2.042500in}{3.028994in}}%
\pgfpathlineto{\pgfqpoint{2.042500in}{3.028994in}}%
\pgfusepath{stroke}%
\end{pgfscope}%
\begin{pgfscope}%
\pgfpathrectangle{\pgfqpoint{0.537500in}{2.059688in}}{\pgfqpoint{1.505000in}{1.156312in}}%
\pgfusepath{clip}%
\pgfsetrectcap%
\pgfsetroundjoin%
\pgfsetlinewidth{0.903375pt}%
\definecolor{currentstroke}{rgb}{0.717647,0.717647,0.717647}%
\pgfsetstrokecolor{currentstroke}%
\pgfsetstrokeopacity{0.900000}%
\pgfsetdash{}{0pt}%
\pgfpathmoveto{\pgfqpoint{0.537500in}{2.987289in}}%
\pgfpathlineto{\pgfqpoint{0.587021in}{2.941027in}}%
\pgfpathlineto{\pgfqpoint{0.632235in}{2.899770in}}%
\pgfpathlineto{\pgfqpoint{0.675297in}{2.861484in}}%
\pgfpathlineto{\pgfqpoint{0.714052in}{2.827991in}}%
\pgfpathlineto{\pgfqpoint{0.750655in}{2.797313in}}%
\pgfpathlineto{\pgfqpoint{0.785104in}{2.769398in}}%
\pgfpathlineto{\pgfqpoint{0.817400in}{2.744174in}}%
\pgfpathlineto{\pgfqpoint{0.847543in}{2.721553in}}%
\pgfpathlineto{\pgfqpoint{0.875533in}{2.701428in}}%
\pgfpathlineto{\pgfqpoint{0.903523in}{2.682236in}}%
\pgfpathlineto{\pgfqpoint{0.929360in}{2.665424in}}%
\pgfpathlineto{\pgfqpoint{0.953044in}{2.650841in}}%
\pgfpathlineto{\pgfqpoint{0.976727in}{2.637112in}}%
\pgfpathlineto{\pgfqpoint{0.998258in}{2.625425in}}%
\pgfpathlineto{\pgfqpoint{1.019789in}{2.614542in}}%
\pgfpathlineto{\pgfqpoint{1.041320in}{2.604512in}}%
\pgfpathlineto{\pgfqpoint{1.060697in}{2.596254in}}%
\pgfpathlineto{\pgfqpoint{1.080075in}{2.588761in}}%
\pgfpathlineto{\pgfqpoint{1.099453in}{2.582064in}}%
\pgfpathlineto{\pgfqpoint{1.118830in}{2.576195in}}%
\pgfpathlineto{\pgfqpoint{1.136055in}{2.571697in}}%
\pgfpathlineto{\pgfqpoint{1.153280in}{2.567893in}}%
\pgfpathlineto{\pgfqpoint{1.170504in}{2.564800in}}%
\pgfpathlineto{\pgfqpoint{1.187729in}{2.562430in}}%
\pgfpathlineto{\pgfqpoint{1.204954in}{2.560794in}}%
\pgfpathlineto{\pgfqpoint{1.222178in}{2.559898in}}%
\pgfpathlineto{\pgfqpoint{1.239403in}{2.559748in}}%
\pgfpathlineto{\pgfqpoint{1.256627in}{2.560344in}}%
\pgfpathlineto{\pgfqpoint{1.273852in}{2.561682in}}%
\pgfpathlineto{\pgfqpoint{1.291077in}{2.563757in}}%
\pgfpathlineto{\pgfqpoint{1.308301in}{2.566561in}}%
\pgfpathlineto{\pgfqpoint{1.325526in}{2.570080in}}%
\pgfpathlineto{\pgfqpoint{1.342750in}{2.574301in}}%
\pgfpathlineto{\pgfqpoint{1.362128in}{2.579866in}}%
\pgfpathlineto{\pgfqpoint{1.381506in}{2.586270in}}%
\pgfpathlineto{\pgfqpoint{1.400883in}{2.593482in}}%
\pgfpathlineto{\pgfqpoint{1.420261in}{2.601470in}}%
\pgfpathlineto{\pgfqpoint{1.439639in}{2.610200in}}%
\pgfpathlineto{\pgfqpoint{1.461170in}{2.620729in}}%
\pgfpathlineto{\pgfqpoint{1.482700in}{2.632082in}}%
\pgfpathlineto{\pgfqpoint{1.506384in}{2.645467in}}%
\pgfpathlineto{\pgfqpoint{1.530068in}{2.659730in}}%
\pgfpathlineto{\pgfqpoint{1.555905in}{2.676218in}}%
\pgfpathlineto{\pgfqpoint{1.581742in}{2.693602in}}%
\pgfpathlineto{\pgfqpoint{1.609732in}{2.713362in}}%
\pgfpathlineto{\pgfqpoint{1.637722in}{2.734005in}}%
\pgfpathlineto{\pgfqpoint{1.667865in}{2.757132in}}%
\pgfpathlineto{\pgfqpoint{1.700161in}{2.782845in}}%
\pgfpathlineto{\pgfqpoint{1.734610in}{2.811226in}}%
\pgfpathlineto{\pgfqpoint{1.771212in}{2.842343in}}%
\pgfpathlineto{\pgfqpoint{1.809968in}{2.876247in}}%
\pgfpathlineto{\pgfqpoint{1.853029in}{2.914932in}}%
\pgfpathlineto{\pgfqpoint{1.898244in}{2.956554in}}%
\pgfpathlineto{\pgfqpoint{1.947765in}{3.003160in}}%
\pgfpathlineto{\pgfqpoint{2.001592in}{3.054857in}}%
\pgfpathlineto{\pgfqpoint{2.042500in}{3.094774in}}%
\pgfpathlineto{\pgfqpoint{2.042500in}{3.094774in}}%
\pgfusepath{stroke}%
\end{pgfscope}%
\begin{pgfscope}%
\pgfpathrectangle{\pgfqpoint{0.537500in}{2.059688in}}{\pgfqpoint{1.505000in}{1.156312in}}%
\pgfusepath{clip}%
\pgfsetrectcap%
\pgfsetroundjoin%
\pgfsetlinewidth{0.903375pt}%
\definecolor{currentstroke}{rgb}{0.717647,0.717647,0.717647}%
\pgfsetstrokecolor{currentstroke}%
\pgfsetstrokeopacity{0.900000}%
\pgfsetdash{}{0pt}%
\pgfpathmoveto{\pgfqpoint{0.537500in}{3.100420in}}%
\pgfpathlineto{\pgfqpoint{0.580562in}{3.064678in}}%
\pgfpathlineto{\pgfqpoint{0.619317in}{3.033436in}}%
\pgfpathlineto{\pgfqpoint{0.658072in}{3.003161in}}%
\pgfpathlineto{\pgfqpoint{0.694675in}{2.975543in}}%
\pgfpathlineto{\pgfqpoint{0.729124in}{2.950496in}}%
\pgfpathlineto{\pgfqpoint{0.761420in}{2.927919in}}%
\pgfpathlineto{\pgfqpoint{0.791563in}{2.907701in}}%
\pgfpathlineto{\pgfqpoint{0.821706in}{2.888368in}}%
\pgfpathlineto{\pgfqpoint{0.849696in}{2.871265in}}%
\pgfpathlineto{\pgfqpoint{0.877686in}{2.855033in}}%
\pgfpathlineto{\pgfqpoint{0.903523in}{2.840870in}}%
\pgfpathlineto{\pgfqpoint{0.929360in}{2.827540in}}%
\pgfpathlineto{\pgfqpoint{0.955197in}{2.815086in}}%
\pgfpathlineto{\pgfqpoint{0.978881in}{2.804477in}}%
\pgfpathlineto{\pgfqpoint{1.002564in}{2.794675in}}%
\pgfpathlineto{\pgfqpoint{1.026248in}{2.785712in}}%
\pgfpathlineto{\pgfqpoint{1.047779in}{2.778320in}}%
\pgfpathlineto{\pgfqpoint{1.069310in}{2.771670in}}%
\pgfpathlineto{\pgfqpoint{1.090840in}{2.765783in}}%
\pgfpathlineto{\pgfqpoint{1.112371in}{2.760678in}}%
\pgfpathlineto{\pgfqpoint{1.133902in}{2.756374in}}%
\pgfpathlineto{\pgfqpoint{1.155433in}{2.752883in}}%
\pgfpathlineto{\pgfqpoint{1.176964in}{2.750220in}}%
\pgfpathlineto{\pgfqpoint{1.196341in}{2.748538in}}%
\pgfpathlineto{\pgfqpoint{1.215719in}{2.747539in}}%
\pgfpathlineto{\pgfqpoint{1.235097in}{2.747226in}}%
\pgfpathlineto{\pgfqpoint{1.254474in}{2.747599in}}%
\pgfpathlineto{\pgfqpoint{1.273852in}{2.748657in}}%
\pgfpathlineto{\pgfqpoint{1.293230in}{2.750398in}}%
\pgfpathlineto{\pgfqpoint{1.312607in}{2.752815in}}%
\pgfpathlineto{\pgfqpoint{1.334138in}{2.756287in}}%
\pgfpathlineto{\pgfqpoint{1.355669in}{2.760574in}}%
\pgfpathlineto{\pgfqpoint{1.377200in}{2.765661in}}%
\pgfpathlineto{\pgfqpoint{1.398730in}{2.771531in}}%
\pgfpathlineto{\pgfqpoint{1.420261in}{2.778164in}}%
\pgfpathlineto{\pgfqpoint{1.441792in}{2.785540in}}%
\pgfpathlineto{\pgfqpoint{1.465476in}{2.794485in}}%
\pgfpathlineto{\pgfqpoint{1.489160in}{2.804271in}}%
\pgfpathlineto{\pgfqpoint{1.512843in}{2.814864in}}%
\pgfpathlineto{\pgfqpoint{1.536527in}{2.826230in}}%
\pgfpathlineto{\pgfqpoint{1.562364in}{2.839474in}}%
\pgfpathlineto{\pgfqpoint{1.588201in}{2.853554in}}%
\pgfpathlineto{\pgfqpoint{1.616191in}{2.869701in}}%
\pgfpathlineto{\pgfqpoint{1.644181in}{2.886724in}}%
\pgfpathlineto{\pgfqpoint{1.674324in}{2.905977in}}%
\pgfpathlineto{\pgfqpoint{1.704467in}{2.926120in}}%
\pgfpathlineto{\pgfqpoint{1.736763in}{2.948622in}}%
\pgfpathlineto{\pgfqpoint{1.769059in}{2.972008in}}%
\pgfpathlineto{\pgfqpoint{1.803509in}{2.997851in}}%
\pgfpathlineto{\pgfqpoint{1.840111in}{3.026244in}}%
\pgfpathlineto{\pgfqpoint{1.878866in}{3.057267in}}%
\pgfpathlineto{\pgfqpoint{1.919775in}{3.090989in}}%
\pgfpathlineto{\pgfqpoint{1.962836in}{3.127464in}}%
\pgfpathlineto{\pgfqpoint{2.008051in}{3.166737in}}%
\pgfpathlineto{\pgfqpoint{2.042500in}{3.197266in}}%
\pgfpathlineto{\pgfqpoint{2.042500in}{3.197266in}}%
\pgfusepath{stroke}%
\end{pgfscope}%
\begin{pgfscope}%
\pgfpathrectangle{\pgfqpoint{0.537500in}{2.059688in}}{\pgfqpoint{1.505000in}{1.156312in}}%
\pgfusepath{clip}%
\pgfsetrectcap%
\pgfsetroundjoin%
\pgfsetlinewidth{0.903375pt}%
\definecolor{currentstroke}{rgb}{0.717647,0.717647,0.717647}%
\pgfsetstrokecolor{currentstroke}%
\pgfsetstrokeopacity{0.900000}%
\pgfsetdash{}{0pt}%
\pgfpathmoveto{\pgfqpoint{0.555627in}{3.219333in}}%
\pgfpathlineto{\pgfqpoint{0.593480in}{3.191933in}}%
\pgfpathlineto{\pgfqpoint{0.630082in}{3.166342in}}%
\pgfpathlineto{\pgfqpoint{0.664531in}{3.143125in}}%
\pgfpathlineto{\pgfqpoint{0.698981in}{3.120806in}}%
\pgfpathlineto{\pgfqpoint{0.731277in}{3.100750in}}%
\pgfpathlineto{\pgfqpoint{0.763573in}{3.081584in}}%
\pgfpathlineto{\pgfqpoint{0.793716in}{3.064543in}}%
\pgfpathlineto{\pgfqpoint{0.823859in}{3.048365in}}%
\pgfpathlineto{\pgfqpoint{0.851849in}{3.034153in}}%
\pgfpathlineto{\pgfqpoint{0.879839in}{3.020755in}}%
\pgfpathlineto{\pgfqpoint{0.907829in}{3.008207in}}%
\pgfpathlineto{\pgfqpoint{0.933666in}{2.997408in}}%
\pgfpathlineto{\pgfqpoint{0.959503in}{2.987388in}}%
\pgfpathlineto{\pgfqpoint{0.985340in}{2.978173in}}%
\pgfpathlineto{\pgfqpoint{1.011177in}{2.969787in}}%
\pgfpathlineto{\pgfqpoint{1.034861in}{2.962849in}}%
\pgfpathlineto{\pgfqpoint{1.058544in}{2.956644in}}%
\pgfpathlineto{\pgfqpoint{1.082228in}{2.951186in}}%
\pgfpathlineto{\pgfqpoint{1.105912in}{2.946491in}}%
\pgfpathlineto{\pgfqpoint{1.129596in}{2.942570in}}%
\pgfpathlineto{\pgfqpoint{1.153280in}{2.939434in}}%
\pgfpathlineto{\pgfqpoint{1.176964in}{2.937090in}}%
\pgfpathlineto{\pgfqpoint{1.200647in}{2.935546in}}%
\pgfpathlineto{\pgfqpoint{1.224331in}{2.934806in}}%
\pgfpathlineto{\pgfqpoint{1.248015in}{2.934871in}}%
\pgfpathlineto{\pgfqpoint{1.271699in}{2.935742in}}%
\pgfpathlineto{\pgfqpoint{1.295383in}{2.937415in}}%
\pgfpathlineto{\pgfqpoint{1.319067in}{2.939887in}}%
\pgfpathlineto{\pgfqpoint{1.342750in}{2.943151in}}%
\pgfpathlineto{\pgfqpoint{1.366434in}{2.947198in}}%
\pgfpathlineto{\pgfqpoint{1.390118in}{2.952017in}}%
\pgfpathlineto{\pgfqpoint{1.413802in}{2.957596in}}%
\pgfpathlineto{\pgfqpoint{1.437486in}{2.963921in}}%
\pgfpathlineto{\pgfqpoint{1.461170in}{2.970977in}}%
\pgfpathlineto{\pgfqpoint{1.487006in}{2.979486in}}%
\pgfpathlineto{\pgfqpoint{1.512843in}{2.988822in}}%
\pgfpathlineto{\pgfqpoint{1.538680in}{2.998959in}}%
\pgfpathlineto{\pgfqpoint{1.564517in}{3.009871in}}%
\pgfpathlineto{\pgfqpoint{1.592507in}{3.022538in}}%
\pgfpathlineto{\pgfqpoint{1.620497in}{3.036049in}}%
\pgfpathlineto{\pgfqpoint{1.648487in}{3.050370in}}%
\pgfpathlineto{\pgfqpoint{1.678630in}{3.066660in}}%
\pgfpathlineto{\pgfqpoint{1.708773in}{3.083807in}}%
\pgfpathlineto{\pgfqpoint{1.741069in}{3.103081in}}%
\pgfpathlineto{\pgfqpoint{1.773366in}{3.123240in}}%
\pgfpathlineto{\pgfqpoint{1.807815in}{3.145661in}}%
\pgfpathlineto{\pgfqpoint{1.842264in}{3.168975in}}%
\pgfpathlineto{\pgfqpoint{1.878866in}{3.194661in}}%
\pgfpathlineto{\pgfqpoint{1.912892in}{3.219333in}}%
\pgfpathlineto{\pgfqpoint{1.912892in}{3.219333in}}%
\pgfusepath{stroke}%
\end{pgfscope}%
\begin{pgfscope}%
\pgfpathrectangle{\pgfqpoint{0.537500in}{2.059688in}}{\pgfqpoint{1.505000in}{1.156312in}}%
\pgfusepath{clip}%
\pgfsetrectcap%
\pgfsetroundjoin%
\pgfsetlinewidth{1.003750pt}%
\definecolor{currentstroke}{rgb}{1.000000,1.000000,1.000000}%
\pgfsetstrokecolor{currentstroke}%
\pgfsetdash{}{0pt}%
\pgfpathmoveto{\pgfqpoint{0.676852in}{2.059688in}}%
\pgfpathlineto{\pgfqpoint{0.676852in}{3.216000in}}%
\pgfusepath{stroke}%
\end{pgfscope}%
\begin{pgfscope}%
\pgfsetbuttcap%
\pgfsetroundjoin%
\definecolor{currentfill}{rgb}{0.000000,0.000000,0.000000}%
\pgfsetfillcolor{currentfill}%
\pgfsetlinewidth{0.803000pt}%
\definecolor{currentstroke}{rgb}{0.000000,0.000000,0.000000}%
\pgfsetstrokecolor{currentstroke}%
\pgfsetdash{}{0pt}%
\pgfsys@defobject{currentmarker}{\pgfqpoint{0.000000in}{-0.048611in}}{\pgfqpoint{0.000000in}{0.000000in}}{%
\pgfpathmoveto{\pgfqpoint{0.000000in}{0.000000in}}%
\pgfpathlineto{\pgfqpoint{0.000000in}{-0.048611in}}%
\pgfusepath{stroke,fill}%
}%
\begin{pgfscope}%
\pgfsys@transformshift{0.676852in}{2.059688in}%
\pgfsys@useobject{currentmarker}{}%
\end{pgfscope}%
\end{pgfscope}%
\begin{pgfscope}%
\definecolor{textcolor}{rgb}{0.000000,0.000000,0.000000}%
\pgfsetstrokecolor{textcolor}%
\pgfsetfillcolor{textcolor}%
\pgftext[x=0.676852in,y=1.962466in,,top]{\color{textcolor}{\rmfamily\fontsize{6.000000}{7.200000}\selectfont\catcode`\^=\active\def^{\ifmmode\sp\else\^{}\fi}\catcode`\%=\active\def%{\%}\ensuremath{-}1.0}}%
\end{pgfscope}%
\begin{pgfscope}%
\pgfpathrectangle{\pgfqpoint{0.537500in}{2.059688in}}{\pgfqpoint{1.505000in}{1.156312in}}%
\pgfusepath{clip}%
\pgfsetrectcap%
\pgfsetroundjoin%
\pgfsetlinewidth{1.003750pt}%
\definecolor{currentstroke}{rgb}{1.000000,1.000000,1.000000}%
\pgfsetstrokecolor{currentstroke}%
\pgfsetdash{}{0pt}%
\pgfpathmoveto{\pgfqpoint{0.955556in}{2.059688in}}%
\pgfpathlineto{\pgfqpoint{0.955556in}{3.216000in}}%
\pgfusepath{stroke}%
\end{pgfscope}%
\begin{pgfscope}%
\pgfsetbuttcap%
\pgfsetroundjoin%
\definecolor{currentfill}{rgb}{0.000000,0.000000,0.000000}%
\pgfsetfillcolor{currentfill}%
\pgfsetlinewidth{0.803000pt}%
\definecolor{currentstroke}{rgb}{0.000000,0.000000,0.000000}%
\pgfsetstrokecolor{currentstroke}%
\pgfsetdash{}{0pt}%
\pgfsys@defobject{currentmarker}{\pgfqpoint{0.000000in}{-0.048611in}}{\pgfqpoint{0.000000in}{0.000000in}}{%
\pgfpathmoveto{\pgfqpoint{0.000000in}{0.000000in}}%
\pgfpathlineto{\pgfqpoint{0.000000in}{-0.048611in}}%
\pgfusepath{stroke,fill}%
}%
\begin{pgfscope}%
\pgfsys@transformshift{0.955556in}{2.059688in}%
\pgfsys@useobject{currentmarker}{}%
\end{pgfscope}%
\end{pgfscope}%
\begin{pgfscope}%
\definecolor{textcolor}{rgb}{0.000000,0.000000,0.000000}%
\pgfsetstrokecolor{textcolor}%
\pgfsetfillcolor{textcolor}%
\pgftext[x=0.955556in,y=1.962466in,,top]{\color{textcolor}{\rmfamily\fontsize{6.000000}{7.200000}\selectfont\catcode`\^=\active\def^{\ifmmode\sp\else\^{}\fi}\catcode`\%=\active\def%{\%}\ensuremath{-}0.5}}%
\end{pgfscope}%
\begin{pgfscope}%
\pgfpathrectangle{\pgfqpoint{0.537500in}{2.059688in}}{\pgfqpoint{1.505000in}{1.156312in}}%
\pgfusepath{clip}%
\pgfsetrectcap%
\pgfsetroundjoin%
\pgfsetlinewidth{1.003750pt}%
\definecolor{currentstroke}{rgb}{1.000000,1.000000,1.000000}%
\pgfsetstrokecolor{currentstroke}%
\pgfsetdash{}{0pt}%
\pgfpathmoveto{\pgfqpoint{1.234259in}{2.059688in}}%
\pgfpathlineto{\pgfqpoint{1.234259in}{3.216000in}}%
\pgfusepath{stroke}%
\end{pgfscope}%
\begin{pgfscope}%
\pgfsetbuttcap%
\pgfsetroundjoin%
\definecolor{currentfill}{rgb}{0.000000,0.000000,0.000000}%
\pgfsetfillcolor{currentfill}%
\pgfsetlinewidth{0.803000pt}%
\definecolor{currentstroke}{rgb}{0.000000,0.000000,0.000000}%
\pgfsetstrokecolor{currentstroke}%
\pgfsetdash{}{0pt}%
\pgfsys@defobject{currentmarker}{\pgfqpoint{0.000000in}{-0.048611in}}{\pgfqpoint{0.000000in}{0.000000in}}{%
\pgfpathmoveto{\pgfqpoint{0.000000in}{0.000000in}}%
\pgfpathlineto{\pgfqpoint{0.000000in}{-0.048611in}}%
\pgfusepath{stroke,fill}%
}%
\begin{pgfscope}%
\pgfsys@transformshift{1.234259in}{2.059688in}%
\pgfsys@useobject{currentmarker}{}%
\end{pgfscope}%
\end{pgfscope}%
\begin{pgfscope}%
\definecolor{textcolor}{rgb}{0.000000,0.000000,0.000000}%
\pgfsetstrokecolor{textcolor}%
\pgfsetfillcolor{textcolor}%
\pgftext[x=1.234259in,y=1.962466in,,top]{\color{textcolor}{\rmfamily\fontsize{6.000000}{7.200000}\selectfont\catcode`\^=\active\def^{\ifmmode\sp\else\^{}\fi}\catcode`\%=\active\def%{\%}0.0}}%
\end{pgfscope}%
\begin{pgfscope}%
\pgfpathrectangle{\pgfqpoint{0.537500in}{2.059688in}}{\pgfqpoint{1.505000in}{1.156312in}}%
\pgfusepath{clip}%
\pgfsetrectcap%
\pgfsetroundjoin%
\pgfsetlinewidth{1.003750pt}%
\definecolor{currentstroke}{rgb}{1.000000,1.000000,1.000000}%
\pgfsetstrokecolor{currentstroke}%
\pgfsetdash{}{0pt}%
\pgfpathmoveto{\pgfqpoint{1.512963in}{2.059688in}}%
\pgfpathlineto{\pgfqpoint{1.512963in}{3.216000in}}%
\pgfusepath{stroke}%
\end{pgfscope}%
\begin{pgfscope}%
\pgfsetbuttcap%
\pgfsetroundjoin%
\definecolor{currentfill}{rgb}{0.000000,0.000000,0.000000}%
\pgfsetfillcolor{currentfill}%
\pgfsetlinewidth{0.803000pt}%
\definecolor{currentstroke}{rgb}{0.000000,0.000000,0.000000}%
\pgfsetstrokecolor{currentstroke}%
\pgfsetdash{}{0pt}%
\pgfsys@defobject{currentmarker}{\pgfqpoint{0.000000in}{-0.048611in}}{\pgfqpoint{0.000000in}{0.000000in}}{%
\pgfpathmoveto{\pgfqpoint{0.000000in}{0.000000in}}%
\pgfpathlineto{\pgfqpoint{0.000000in}{-0.048611in}}%
\pgfusepath{stroke,fill}%
}%
\begin{pgfscope}%
\pgfsys@transformshift{1.512963in}{2.059688in}%
\pgfsys@useobject{currentmarker}{}%
\end{pgfscope}%
\end{pgfscope}%
\begin{pgfscope}%
\definecolor{textcolor}{rgb}{0.000000,0.000000,0.000000}%
\pgfsetstrokecolor{textcolor}%
\pgfsetfillcolor{textcolor}%
\pgftext[x=1.512963in,y=1.962466in,,top]{\color{textcolor}{\rmfamily\fontsize{6.000000}{7.200000}\selectfont\catcode`\^=\active\def^{\ifmmode\sp\else\^{}\fi}\catcode`\%=\active\def%{\%}0.5}}%
\end{pgfscope}%
\begin{pgfscope}%
\pgfpathrectangle{\pgfqpoint{0.537500in}{2.059688in}}{\pgfqpoint{1.505000in}{1.156312in}}%
\pgfusepath{clip}%
\pgfsetrectcap%
\pgfsetroundjoin%
\pgfsetlinewidth{1.003750pt}%
\definecolor{currentstroke}{rgb}{1.000000,1.000000,1.000000}%
\pgfsetstrokecolor{currentstroke}%
\pgfsetdash{}{0pt}%
\pgfpathmoveto{\pgfqpoint{1.791667in}{2.059688in}}%
\pgfpathlineto{\pgfqpoint{1.791667in}{3.216000in}}%
\pgfusepath{stroke}%
\end{pgfscope}%
\begin{pgfscope}%
\pgfsetbuttcap%
\pgfsetroundjoin%
\definecolor{currentfill}{rgb}{0.000000,0.000000,0.000000}%
\pgfsetfillcolor{currentfill}%
\pgfsetlinewidth{0.803000pt}%
\definecolor{currentstroke}{rgb}{0.000000,0.000000,0.000000}%
\pgfsetstrokecolor{currentstroke}%
\pgfsetdash{}{0pt}%
\pgfsys@defobject{currentmarker}{\pgfqpoint{0.000000in}{-0.048611in}}{\pgfqpoint{0.000000in}{0.000000in}}{%
\pgfpathmoveto{\pgfqpoint{0.000000in}{0.000000in}}%
\pgfpathlineto{\pgfqpoint{0.000000in}{-0.048611in}}%
\pgfusepath{stroke,fill}%
}%
\begin{pgfscope}%
\pgfsys@transformshift{1.791667in}{2.059688in}%
\pgfsys@useobject{currentmarker}{}%
\end{pgfscope}%
\end{pgfscope}%
\begin{pgfscope}%
\definecolor{textcolor}{rgb}{0.000000,0.000000,0.000000}%
\pgfsetstrokecolor{textcolor}%
\pgfsetfillcolor{textcolor}%
\pgftext[x=1.791667in,y=1.962466in,,top]{\color{textcolor}{\rmfamily\fontsize{6.000000}{7.200000}\selectfont\catcode`\^=\active\def^{\ifmmode\sp\else\^{}\fi}\catcode`\%=\active\def%{\%}1.0}}%
\end{pgfscope}%
\begin{pgfscope}%
\definecolor{textcolor}{rgb}{0.000000,0.000000,0.000000}%
\pgfsetstrokecolor{textcolor}%
\pgfsetfillcolor{textcolor}%
\pgftext[x=1.290000in,y=1.823577in,,top]{\color{textcolor}{\rmfamily\fontsize{6.000000}{7.200000}\selectfont\catcode`\^=\active\def^{\ifmmode\sp\else\^{}\fi}\catcode`\%=\active\def%{\%}predictor $r$}}%
\end{pgfscope}%
\begin{pgfscope}%
\pgfpathrectangle{\pgfqpoint{0.537500in}{2.059688in}}{\pgfqpoint{1.505000in}{1.156312in}}%
\pgfusepath{clip}%
\pgfsetrectcap%
\pgfsetroundjoin%
\pgfsetlinewidth{1.003750pt}%
\definecolor{currentstroke}{rgb}{1.000000,1.000000,1.000000}%
\pgfsetstrokecolor{currentstroke}%
\pgfsetdash{}{0pt}%
\pgfpathmoveto{\pgfqpoint{0.537500in}{2.059688in}}%
\pgfpathlineto{\pgfqpoint{2.042500in}{2.059688in}}%
\pgfusepath{stroke}%
\end{pgfscope}%
\begin{pgfscope}%
\pgfsetbuttcap%
\pgfsetroundjoin%
\definecolor{currentfill}{rgb}{0.000000,0.000000,0.000000}%
\pgfsetfillcolor{currentfill}%
\pgfsetlinewidth{0.803000pt}%
\definecolor{currentstroke}{rgb}{0.000000,0.000000,0.000000}%
\pgfsetstrokecolor{currentstroke}%
\pgfsetdash{}{0pt}%
\pgfsys@defobject{currentmarker}{\pgfqpoint{-0.048611in}{0.000000in}}{\pgfqpoint{-0.000000in}{0.000000in}}{%
\pgfpathmoveto{\pgfqpoint{-0.000000in}{0.000000in}}%
\pgfpathlineto{\pgfqpoint{-0.048611in}{0.000000in}}%
\pgfusepath{stroke,fill}%
}%
\begin{pgfscope}%
\pgfsys@transformshift{0.537500in}{2.059688in}%
\pgfsys@useobject{currentmarker}{}%
\end{pgfscope}%
\end{pgfscope}%
\begin{pgfscope}%
\definecolor{textcolor}{rgb}{0.000000,0.000000,0.000000}%
\pgfsetstrokecolor{textcolor}%
\pgfsetfillcolor{textcolor}%
\pgftext[x=0.307750in, y=2.028031in, left, base]{\color{textcolor}{\rmfamily\fontsize{6.000000}{7.200000}\selectfont\catcode`\^=\active\def^{\ifmmode\sp\else\^{}\fi}\catcode`\%=\active\def%{\%}0.0}}%
\end{pgfscope}%
\begin{pgfscope}%
\pgfpathrectangle{\pgfqpoint{0.537500in}{2.059688in}}{\pgfqpoint{1.505000in}{1.156312in}}%
\pgfusepath{clip}%
\pgfsetrectcap%
\pgfsetroundjoin%
\pgfsetlinewidth{1.003750pt}%
\definecolor{currentstroke}{rgb}{1.000000,1.000000,1.000000}%
\pgfsetstrokecolor{currentstroke}%
\pgfsetdash{}{0pt}%
\pgfpathmoveto{\pgfqpoint{0.537500in}{2.372205in}}%
\pgfpathlineto{\pgfqpoint{2.042500in}{2.372205in}}%
\pgfusepath{stroke}%
\end{pgfscope}%
\begin{pgfscope}%
\pgfsetbuttcap%
\pgfsetroundjoin%
\definecolor{currentfill}{rgb}{0.000000,0.000000,0.000000}%
\pgfsetfillcolor{currentfill}%
\pgfsetlinewidth{0.803000pt}%
\definecolor{currentstroke}{rgb}{0.000000,0.000000,0.000000}%
\pgfsetstrokecolor{currentstroke}%
\pgfsetdash{}{0pt}%
\pgfsys@defobject{currentmarker}{\pgfqpoint{-0.048611in}{0.000000in}}{\pgfqpoint{-0.000000in}{0.000000in}}{%
\pgfpathmoveto{\pgfqpoint{-0.000000in}{0.000000in}}%
\pgfpathlineto{\pgfqpoint{-0.048611in}{0.000000in}}%
\pgfusepath{stroke,fill}%
}%
\begin{pgfscope}%
\pgfsys@transformshift{0.537500in}{2.372205in}%
\pgfsys@useobject{currentmarker}{}%
\end{pgfscope}%
\end{pgfscope}%
\begin{pgfscope}%
\definecolor{textcolor}{rgb}{0.000000,0.000000,0.000000}%
\pgfsetstrokecolor{textcolor}%
\pgfsetfillcolor{textcolor}%
\pgftext[x=0.307750in, y=2.340548in, left, base]{\color{textcolor}{\rmfamily\fontsize{6.000000}{7.200000}\selectfont\catcode`\^=\active\def^{\ifmmode\sp\else\^{}\fi}\catcode`\%=\active\def%{\%}0.5}}%
\end{pgfscope}%
\begin{pgfscope}%
\pgfpathrectangle{\pgfqpoint{0.537500in}{2.059688in}}{\pgfqpoint{1.505000in}{1.156312in}}%
\pgfusepath{clip}%
\pgfsetrectcap%
\pgfsetroundjoin%
\pgfsetlinewidth{1.003750pt}%
\definecolor{currentstroke}{rgb}{1.000000,1.000000,1.000000}%
\pgfsetstrokecolor{currentstroke}%
\pgfsetdash{}{0pt}%
\pgfpathmoveto{\pgfqpoint{0.537500in}{2.684722in}}%
\pgfpathlineto{\pgfqpoint{2.042500in}{2.684722in}}%
\pgfusepath{stroke}%
\end{pgfscope}%
\begin{pgfscope}%
\pgfsetbuttcap%
\pgfsetroundjoin%
\definecolor{currentfill}{rgb}{0.000000,0.000000,0.000000}%
\pgfsetfillcolor{currentfill}%
\pgfsetlinewidth{0.803000pt}%
\definecolor{currentstroke}{rgb}{0.000000,0.000000,0.000000}%
\pgfsetstrokecolor{currentstroke}%
\pgfsetdash{}{0pt}%
\pgfsys@defobject{currentmarker}{\pgfqpoint{-0.048611in}{0.000000in}}{\pgfqpoint{-0.000000in}{0.000000in}}{%
\pgfpathmoveto{\pgfqpoint{-0.000000in}{0.000000in}}%
\pgfpathlineto{\pgfqpoint{-0.048611in}{0.000000in}}%
\pgfusepath{stroke,fill}%
}%
\begin{pgfscope}%
\pgfsys@transformshift{0.537500in}{2.684722in}%
\pgfsys@useobject{currentmarker}{}%
\end{pgfscope}%
\end{pgfscope}%
\begin{pgfscope}%
\definecolor{textcolor}{rgb}{0.000000,0.000000,0.000000}%
\pgfsetstrokecolor{textcolor}%
\pgfsetfillcolor{textcolor}%
\pgftext[x=0.307750in, y=2.653065in, left, base]{\color{textcolor}{\rmfamily\fontsize{6.000000}{7.200000}\selectfont\catcode`\^=\active\def^{\ifmmode\sp\else\^{}\fi}\catcode`\%=\active\def%{\%}1.0}}%
\end{pgfscope}%
\begin{pgfscope}%
\pgfpathrectangle{\pgfqpoint{0.537500in}{2.059688in}}{\pgfqpoint{1.505000in}{1.156312in}}%
\pgfusepath{clip}%
\pgfsetrectcap%
\pgfsetroundjoin%
\pgfsetlinewidth{1.003750pt}%
\definecolor{currentstroke}{rgb}{1.000000,1.000000,1.000000}%
\pgfsetstrokecolor{currentstroke}%
\pgfsetdash{}{0pt}%
\pgfpathmoveto{\pgfqpoint{0.537500in}{2.997238in}}%
\pgfpathlineto{\pgfqpoint{2.042500in}{2.997238in}}%
\pgfusepath{stroke}%
\end{pgfscope}%
\begin{pgfscope}%
\pgfsetbuttcap%
\pgfsetroundjoin%
\definecolor{currentfill}{rgb}{0.000000,0.000000,0.000000}%
\pgfsetfillcolor{currentfill}%
\pgfsetlinewidth{0.803000pt}%
\definecolor{currentstroke}{rgb}{0.000000,0.000000,0.000000}%
\pgfsetstrokecolor{currentstroke}%
\pgfsetdash{}{0pt}%
\pgfsys@defobject{currentmarker}{\pgfqpoint{-0.048611in}{0.000000in}}{\pgfqpoint{-0.000000in}{0.000000in}}{%
\pgfpathmoveto{\pgfqpoint{-0.000000in}{0.000000in}}%
\pgfpathlineto{\pgfqpoint{-0.048611in}{0.000000in}}%
\pgfusepath{stroke,fill}%
}%
\begin{pgfscope}%
\pgfsys@transformshift{0.537500in}{2.997238in}%
\pgfsys@useobject{currentmarker}{}%
\end{pgfscope}%
\end{pgfscope}%
\begin{pgfscope}%
\definecolor{textcolor}{rgb}{0.000000,0.000000,0.000000}%
\pgfsetstrokecolor{textcolor}%
\pgfsetfillcolor{textcolor}%
\pgftext[x=0.307750in, y=2.965581in, left, base]{\color{textcolor}{\rmfamily\fontsize{6.000000}{7.200000}\selectfont\catcode`\^=\active\def^{\ifmmode\sp\else\^{}\fi}\catcode`\%=\active\def%{\%}1.5}}%
\end{pgfscope}%
\begin{pgfscope}%
\definecolor{textcolor}{rgb}{0.000000,0.000000,0.000000}%
\pgfsetstrokecolor{textcolor}%
\pgfsetfillcolor{textcolor}%
\pgftext[x=0.252195in,y=2.637844in,,bottom,rotate=90.000000]{\color{textcolor}{\rmfamily\fontsize{6.000000}{7.200000}\selectfont\catcode`\^=\active\def^{\ifmmode\sp\else\^{}\fi}\catcode`\%=\active\def%{\%}shared variance $s$}}%
\end{pgfscope}%
\begin{pgfscope}%
\pgfpathrectangle{\pgfqpoint{0.537500in}{2.059688in}}{\pgfqpoint{1.505000in}{1.156312in}}%
\pgfusepath{clip}%
\pgfsetrectcap%
\pgfsetroundjoin%
\pgfsetlinewidth{1.003750pt}%
\definecolor{currentstroke}{rgb}{0.447059,0.447059,0.447059}%
\pgfsetstrokecolor{currentstroke}%
\pgfsetdash{}{0pt}%
\pgfpathmoveto{\pgfqpoint{0.537500in}{2.840980in}}%
\pgfpathlineto{\pgfqpoint{1.232943in}{2.061164in}}%
\pgfpathlineto{\pgfqpoint{1.235097in}{2.060627in}}%
\pgfpathlineto{\pgfqpoint{2.042500in}{2.965987in}}%
\pgfpathlineto{\pgfqpoint{2.042500in}{2.965987in}}%
\pgfusepath{stroke}%
\end{pgfscope}%
\begin{pgfscope}%
\pgfpathrectangle{\pgfqpoint{0.537500in}{2.059688in}}{\pgfqpoint{1.505000in}{1.156312in}}%
\pgfusepath{clip}%
\pgfsetrectcap%
\pgfsetroundjoin%
\pgfsetlinewidth{2.007500pt}%
\definecolor{currentstroke}{rgb}{0.168627,0.611765,0.411765}%
\pgfsetstrokecolor{currentstroke}%
\pgfsetdash{}{0pt}%
\pgfpathmoveto{\pgfqpoint{1.791667in}{2.684722in}}%
\pgfpathlineto{\pgfqpoint{1.791667in}{3.216000in}}%
\pgfusepath{stroke}%
\end{pgfscope}%
\begin{pgfscope}%
\pgfsetrectcap%
\pgfsetmiterjoin%
\pgfsetlinewidth{0.803000pt}%
\definecolor{currentstroke}{rgb}{0.000000,0.000000,0.000000}%
\pgfsetstrokecolor{currentstroke}%
\pgfsetdash{}{0pt}%
\pgfpathmoveto{\pgfqpoint{0.537500in}{2.059688in}}%
\pgfpathlineto{\pgfqpoint{0.537500in}{3.216000in}}%
\pgfusepath{stroke}%
\end{pgfscope}%
\begin{pgfscope}%
\pgfsetrectcap%
\pgfsetmiterjoin%
\pgfsetlinewidth{0.803000pt}%
\definecolor{currentstroke}{rgb}{0.000000,0.000000,0.000000}%
\pgfsetstrokecolor{currentstroke}%
\pgfsetdash{}{0pt}%
\pgfpathmoveto{\pgfqpoint{0.537500in}{2.059688in}}%
\pgfpathlineto{\pgfqpoint{2.042500in}{2.059688in}}%
\pgfusepath{stroke}%
\end{pgfscope}%
\begin{pgfscope}%
\pgfpathrectangle{\pgfqpoint{0.537500in}{2.059688in}}{\pgfqpoint{1.505000in}{1.156312in}}%
\pgfusepath{clip}%
\pgfsetrectcap%
\pgfsetroundjoin%
\pgfsetlinewidth{2.158062pt}%
\definecolor{currentstroke}{rgb}{0.141176,0.345098,0.650980}%
\pgfsetstrokecolor{currentstroke}%
\pgfsetdash{}{0pt}%
\pgfpathmoveto{\pgfqpoint{0.648981in}{2.732899in}}%
\pgfpathlineto{\pgfqpoint{0.756547in}{2.615965in}}%
\pgfpathlineto{\pgfqpoint{0.826223in}{2.541192in}}%
\pgfpathlineto{\pgfqpoint{0.883040in}{2.481120in}}%
\pgfpathlineto{\pgfqpoint{0.930240in}{2.432137in}}%
\pgfpathlineto{\pgfqpoint{0.970077in}{2.391737in}}%
\pgfpathlineto{\pgfqpoint{1.004172in}{2.358130in}}%
\pgfpathlineto{\pgfqpoint{1.033717in}{2.330003in}}%
\pgfpathlineto{\pgfqpoint{1.059614in}{2.306373in}}%
\pgfpathlineto{\pgfqpoint{1.082556in}{2.286490in}}%
\pgfpathlineto{\pgfqpoint{1.103089in}{2.269774in}}%
\pgfpathlineto{\pgfqpoint{1.121646in}{2.255770in}}%
\pgfpathlineto{\pgfqpoint{1.135307in}{2.246273in}}%
\pgfpathlineto{\pgfqpoint{1.148079in}{2.238129in}}%
\pgfpathlineto{\pgfqpoint{1.160091in}{2.231208in}}%
\pgfpathlineto{\pgfqpoint{1.171454in}{2.225405in}}%
\pgfpathlineto{\pgfqpoint{1.182266in}{2.220628in}}%
\pgfpathlineto{\pgfqpoint{1.192612in}{2.216798in}}%
\pgfpathlineto{\pgfqpoint{1.202568in}{2.213848in}}%
\pgfpathlineto{\pgfqpoint{1.212200in}{2.211722in}}%
\pgfpathlineto{\pgfqpoint{1.221568in}{2.210370in}}%
\pgfpathlineto{\pgfqpoint{1.230725in}{2.209749in}}%
\pgfpathlineto{\pgfqpoint{1.239720in}{2.209821in}}%
\pgfpathlineto{\pgfqpoint{1.248593in}{2.210555in}}%
\pgfpathlineto{\pgfqpoint{1.257385in}{2.211921in}}%
\pgfpathlineto{\pgfqpoint{1.268311in}{2.214480in}}%
\pgfpathlineto{\pgfqpoint{1.279221in}{2.217942in}}%
\pgfpathlineto{\pgfqpoint{1.290168in}{2.222269in}}%
\pgfpathlineto{\pgfqpoint{1.301194in}{2.227425in}}%
\pgfpathlineto{\pgfqpoint{1.312336in}{2.233376in}}%
\pgfpathlineto{\pgfqpoint{1.325902in}{2.241519in}}%
\pgfpathlineto{\pgfqpoint{1.339718in}{2.250702in}}%
\pgfpathlineto{\pgfqpoint{1.356186in}{2.262652in}}%
\pgfpathlineto{\pgfqpoint{1.373047in}{2.275841in}}%
\pgfpathlineto{\pgfqpoint{1.392794in}{2.292291in}}%
\pgfpathlineto{\pgfqpoint{1.415564in}{2.312342in}}%
\pgfpathlineto{\pgfqpoint{1.441420in}{2.336208in}}%
\pgfpathlineto{\pgfqpoint{1.470280in}{2.363901in}}%
\pgfpathlineto{\pgfqpoint{1.504453in}{2.397765in}}%
\pgfpathlineto{\pgfqpoint{1.545628in}{2.439694in}}%
\pgfpathlineto{\pgfqpoint{1.593898in}{2.489956in}}%
\pgfpathlineto{\pgfqpoint{1.653933in}{2.553608in}}%
\pgfpathlineto{\pgfqpoint{1.727725in}{2.632995in}}%
\pgfpathlineto{\pgfqpoint{1.787975in}{2.698446in}}%
\pgfpathlineto{\pgfqpoint{1.787975in}{2.698446in}}%
\pgfusepath{stroke}%
\end{pgfscope}%
\begin{pgfscope}%
\pgfsetroundcap%
\pgfsetroundjoin%
\pgfsetlinewidth{1.254687pt}%
\definecolor{currentstroke}{rgb}{0.141176,0.345098,0.650980}%
\pgfsetstrokecolor{currentstroke}%
\pgfsetdash{}{0pt}%
\pgfpathmoveto{\pgfqpoint{1.577532in}{2.472805in}}%
\pgfpathquadraticcurveto{\pgfqpoint{1.578720in}{2.474047in}}{\pgfqpoint{1.566490in}{2.461263in}}%
\pgfusepath{stroke}%
\end{pgfscope}%
\begin{pgfscope}%
\pgfsetroundcap%
\pgfsetroundjoin%
\definecolor{currentfill}{rgb}{0.141176,0.345098,0.650980}%
\pgfsetfillcolor{currentfill}%
\pgfsetlinewidth{1.254687pt}%
\definecolor{currentstroke}{rgb}{0.141176,0.345098,0.650980}%
\pgfsetstrokecolor{currentstroke}%
\pgfsetdash{}{0pt}%
\pgfpathmoveto{\pgfqpoint{1.531404in}{2.448699in}}%
\pgfpathlineto{\pgfqpoint{1.566490in}{2.461263in}}%
\pgfpathlineto{\pgfqpoint{1.555490in}{2.425656in}}%
\pgfpathlineto{\pgfqpoint{1.531404in}{2.448699in}}%
\pgfpathclose%
\pgfusepath{stroke,fill}%
\end{pgfscope}%
\begin{pgfscope}%
\pgfpathrectangle{\pgfqpoint{0.537500in}{2.059688in}}{\pgfqpoint{1.505000in}{1.156312in}}%
\pgfusepath{clip}%
\pgfsetrectcap%
\pgfsetroundjoin%
\pgfsetlinewidth{2.158062pt}%
\definecolor{currentstroke}{rgb}{0.819608,0.286275,0.356863}%
\pgfsetstrokecolor{currentstroke}%
\pgfsetdash{}{0pt}%
\pgfpathmoveto{\pgfqpoint{0.721444in}{2.681840in}}%
\pgfpathlineto{\pgfqpoint{0.771731in}{2.630130in}}%
\pgfpathlineto{\pgfqpoint{0.815874in}{2.585530in}}%
\pgfpathlineto{\pgfqpoint{0.854966in}{2.546823in}}%
\pgfpathlineto{\pgfqpoint{0.889861in}{2.513063in}}%
\pgfpathlineto{\pgfqpoint{0.921236in}{2.483496in}}%
\pgfpathlineto{\pgfqpoint{0.949635in}{2.457523in}}%
\pgfpathlineto{\pgfqpoint{0.975502in}{2.434654in}}%
\pgfpathlineto{\pgfqpoint{0.999199in}{2.414492in}}%
\pgfpathlineto{\pgfqpoint{1.021029in}{2.396706in}}%
\pgfpathlineto{\pgfqpoint{1.041245in}{2.381021in}}%
\pgfpathlineto{\pgfqpoint{1.060061in}{2.367206in}}%
\pgfpathlineto{\pgfqpoint{1.077657in}{2.355066in}}%
\pgfpathlineto{\pgfqpoint{1.094189in}{2.344436in}}%
\pgfpathlineto{\pgfqpoint{1.109792in}{2.335173in}}%
\pgfpathlineto{\pgfqpoint{1.124580in}{2.327154in}}%
\pgfpathlineto{\pgfqpoint{1.138655in}{2.320274in}}%
\pgfpathlineto{\pgfqpoint{1.152105in}{2.314440in}}%
\pgfpathlineto{\pgfqpoint{1.165008in}{2.309573in}}%
\pgfpathlineto{\pgfqpoint{1.177431in}{2.305601in}}%
\pgfpathlineto{\pgfqpoint{1.189435in}{2.302461in}}%
\pgfpathlineto{\pgfqpoint{1.201074in}{2.300098in}}%
\pgfpathlineto{\pgfqpoint{1.212394in}{2.298463in}}%
\pgfpathlineto{\pgfqpoint{1.223438in}{2.297511in}}%
\pgfpathlineto{\pgfqpoint{1.234243in}{2.297201in}}%
\pgfpathlineto{\pgfqpoint{1.244843in}{2.297497in}}%
\pgfpathlineto{\pgfqpoint{1.255266in}{2.298366in}}%
\pgfpathlineto{\pgfqpoint{1.268935in}{2.300363in}}%
\pgfpathlineto{\pgfqpoint{1.282389in}{2.303255in}}%
\pgfpathlineto{\pgfqpoint{1.295675in}{2.306983in}}%
\pgfpathlineto{\pgfqpoint{1.308828in}{2.311489in}}%
\pgfpathlineto{\pgfqpoint{1.321879in}{2.316720in}}%
\pgfpathlineto{\pgfqpoint{1.334851in}{2.322624in}}%
\pgfpathlineto{\pgfqpoint{1.350983in}{2.330877in}}%
\pgfpathlineto{\pgfqpoint{1.367044in}{2.340013in}}%
\pgfpathlineto{\pgfqpoint{1.383044in}{2.349941in}}%
\pgfpathlineto{\pgfqpoint{1.402166in}{2.362774in}}%
\pgfpathlineto{\pgfqpoint{1.421189in}{2.376466in}}%
\pgfpathlineto{\pgfqpoint{1.443216in}{2.393324in}}%
\pgfpathlineto{\pgfqpoint{1.464997in}{2.410907in}}%
\pgfpathlineto{\pgfqpoint{1.489477in}{2.431592in}}%
\pgfpathlineto{\pgfqpoint{1.516315in}{2.455216in}}%
\pgfpathlineto{\pgfqpoint{1.547806in}{2.483982in}}%
\pgfpathlineto{\pgfqpoint{1.582688in}{2.516919in}}%
\pgfpathlineto{\pgfqpoint{1.621471in}{2.554595in}}%
\pgfpathlineto{\pgfqpoint{1.666455in}{2.599391in}}%
\pgfpathlineto{\pgfqpoint{1.717184in}{2.651000in}}%
\pgfpathlineto{\pgfqpoint{1.776258in}{2.712206in}}%
\pgfpathlineto{\pgfqpoint{1.790609in}{2.727219in}}%
\pgfpathlineto{\pgfqpoint{1.790609in}{2.727219in}}%
\pgfusepath{stroke}%
\end{pgfscope}%
\begin{pgfscope}%
\pgfsetroundcap%
\pgfsetroundjoin%
\pgfsetlinewidth{1.254687pt}%
\definecolor{currentstroke}{rgb}{0.819608,0.286275,0.356863}%
\pgfsetstrokecolor{currentstroke}%
\pgfsetdash{}{0pt}%
\pgfpathmoveto{\pgfqpoint{1.694524in}{2.627822in}}%
\pgfpathquadraticcurveto{\pgfqpoint{1.695262in}{2.628575in}}{\pgfqpoint{1.682405in}{2.615474in}}%
\pgfusepath{stroke}%
\end{pgfscope}%
\begin{pgfscope}%
\pgfsetroundcap%
\pgfsetroundjoin%
\definecolor{currentfill}{rgb}{0.819608,0.286275,0.356863}%
\pgfsetfillcolor{currentfill}%
\pgfsetlinewidth{1.254687pt}%
\definecolor{currentstroke}{rgb}{0.819608,0.286275,0.356863}%
\pgfsetstrokecolor{currentstroke}%
\pgfsetdash{}{0pt}%
\pgfpathmoveto{\pgfqpoint{1.647161in}{2.603358in}}%
\pgfpathlineto{\pgfqpoint{1.682405in}{2.615474in}}%
\pgfpathlineto{\pgfqpoint{1.670951in}{2.580010in}}%
\pgfpathlineto{\pgfqpoint{1.647161in}{2.603358in}}%
\pgfpathclose%
\pgfusepath{stroke,fill}%
\end{pgfscope}%
\begin{pgfscope}%
\pgfpathrectangle{\pgfqpoint{0.537500in}{2.059688in}}{\pgfqpoint{1.505000in}{1.156312in}}%
\pgfusepath{clip}%
\pgfsetrectcap%
\pgfsetroundjoin%
\pgfsetlinewidth{2.158062pt}%
\definecolor{currentstroke}{rgb}{0.000000,0.529412,0.423529}%
\pgfsetstrokecolor{currentstroke}%
\pgfsetdash{}{0pt}%
\pgfpathmoveto{\pgfqpoint{0.810630in}{2.657241in}}%
\pgfpathlineto{\pgfqpoint{0.841008in}{2.630536in}}%
\pgfpathlineto{\pgfqpoint{0.869187in}{2.606496in}}%
\pgfpathlineto{\pgfqpoint{0.895430in}{2.584829in}}%
\pgfpathlineto{\pgfqpoint{0.919959in}{2.565285in}}%
\pgfpathlineto{\pgfqpoint{0.942969in}{2.547651in}}%
\pgfpathlineto{\pgfqpoint{0.964625in}{2.531741in}}%
\pgfpathlineto{\pgfqpoint{0.985072in}{2.517395in}}%
\pgfpathlineto{\pgfqpoint{1.004437in}{2.504472in}}%
\pgfpathlineto{\pgfqpoint{1.022828in}{2.492849in}}%
\pgfpathlineto{\pgfqpoint{1.040344in}{2.482419in}}%
\pgfpathlineto{\pgfqpoint{1.057071in}{2.473084in}}%
\pgfpathlineto{\pgfqpoint{1.073083in}{2.464761in}}%
\pgfpathlineto{\pgfqpoint{1.088450in}{2.457372in}}%
\pgfpathlineto{\pgfqpoint{1.103230in}{2.450850in}}%
\pgfpathlineto{\pgfqpoint{1.117477in}{2.445134in}}%
\pgfpathlineto{\pgfqpoint{1.131240in}{2.440168in}}%
\pgfpathlineto{\pgfqpoint{1.144561in}{2.435902in}}%
\pgfpathlineto{\pgfqpoint{1.157479in}{2.432291in}}%
\pgfpathlineto{\pgfqpoint{1.176175in}{2.428012in}}%
\pgfpathlineto{\pgfqpoint{1.194144in}{2.424988in}}%
\pgfpathlineto{\pgfqpoint{1.211474in}{2.423107in}}%
\pgfpathlineto{\pgfqpoint{1.228240in}{2.422271in}}%
\pgfpathlineto{\pgfqpoint{1.244505in}{2.422390in}}%
\pgfpathlineto{\pgfqpoint{1.260326in}{2.423384in}}%
\pgfpathlineto{\pgfqpoint{1.275747in}{2.425180in}}%
\pgfpathlineto{\pgfqpoint{1.290809in}{2.427712in}}%
\pgfpathlineto{\pgfqpoint{1.305544in}{2.430915in}}%
\pgfpathlineto{\pgfqpoint{1.319979in}{2.434734in}}%
\pgfpathlineto{\pgfqpoint{1.334137in}{2.439113in}}%
\pgfpathlineto{\pgfqpoint{1.352615in}{2.445737in}}%
\pgfpathlineto{\pgfqpoint{1.370666in}{2.453155in}}%
\pgfpathlineto{\pgfqpoint{1.388312in}{2.461261in}}%
\pgfpathlineto{\pgfqpoint{1.405569in}{2.469957in}}%
\pgfpathlineto{\pgfqpoint{1.426602in}{2.481515in}}%
\pgfpathlineto{\pgfqpoint{1.447042in}{2.493681in}}%
\pgfpathlineto{\pgfqpoint{1.466880in}{2.506297in}}%
\pgfpathlineto{\pgfqpoint{1.489872in}{2.521828in}}%
\pgfpathlineto{\pgfqpoint{1.515531in}{2.540204in}}%
\pgfpathlineto{\pgfqpoint{1.539883in}{2.558553in}}%
\pgfpathlineto{\pgfqpoint{1.566046in}{2.579144in}}%
\pgfpathlineto{\pgfqpoint{1.596159in}{2.603838in}}%
\pgfpathlineto{\pgfqpoint{1.628427in}{2.631330in}}%
\pgfpathlineto{\pgfqpoint{1.662901in}{2.661719in}}%
\pgfpathlineto{\pgfqpoint{1.700993in}{2.696339in}}%
\pgfpathlineto{\pgfqpoint{1.742005in}{2.734652in}}%
\pgfpathlineto{\pgfqpoint{1.787371in}{2.778081in}}%
\pgfpathlineto{\pgfqpoint{1.791433in}{2.782018in}}%
\pgfpathlineto{\pgfqpoint{1.791433in}{2.782018in}}%
\pgfusepath{stroke}%
\end{pgfscope}%
\begin{pgfscope}%
\pgfsetroundcap%
\pgfsetroundjoin%
\pgfsetlinewidth{1.254687pt}%
\definecolor{currentstroke}{rgb}{0.000000,0.529412,0.423529}%
\pgfsetstrokecolor{currentstroke}%
\pgfsetdash{}{0pt}%
\pgfpathmoveto{\pgfqpoint{1.756640in}{2.748551in}}%
\pgfpathquadraticcurveto{\pgfqpoint{1.756962in}{2.748858in}}{\pgfqpoint{1.743237in}{2.735770in}}%
\pgfusepath{stroke}%
\end{pgfscope}%
\begin{pgfscope}%
\pgfsetroundcap%
\pgfsetroundjoin%
\definecolor{currentfill}{rgb}{0.000000,0.529412,0.423529}%
\pgfsetfillcolor{currentfill}%
\pgfsetlinewidth{1.254687pt}%
\definecolor{currentstroke}{rgb}{0.000000,0.529412,0.423529}%
\pgfsetstrokecolor{currentstroke}%
\pgfsetdash{}{0pt}%
\pgfpathmoveto{\pgfqpoint{1.707613in}{2.724826in}}%
\pgfpathlineto{\pgfqpoint{1.743237in}{2.735770in}}%
\pgfpathlineto{\pgfqpoint{1.730617in}{2.700704in}}%
\pgfpathlineto{\pgfqpoint{1.707613in}{2.724826in}}%
\pgfpathclose%
\pgfusepath{stroke,fill}%
\end{pgfscope}%
\begin{pgfscope}%
\pgfpathrectangle{\pgfqpoint{0.537500in}{2.059688in}}{\pgfqpoint{1.505000in}{1.156312in}}%
\pgfusepath{clip}%
\pgfsetrectcap%
\pgfsetroundjoin%
\pgfsetlinewidth{2.158062pt}%
\definecolor{currentstroke}{rgb}{0.478431,0.243137,0.615686}%
\pgfsetstrokecolor{currentstroke}%
\pgfsetdash{}{0pt}%
\pgfpathmoveto{\pgfqpoint{0.910963in}{2.687466in}}%
\pgfpathlineto{\pgfqpoint{0.939822in}{2.669352in}}%
\pgfpathlineto{\pgfqpoint{0.966980in}{2.653412in}}%
\pgfpathlineto{\pgfqpoint{0.992629in}{2.639426in}}%
\pgfpathlineto{\pgfqpoint{1.016935in}{2.627200in}}%
\pgfpathlineto{\pgfqpoint{1.040039in}{2.616567in}}%
\pgfpathlineto{\pgfqpoint{1.062065in}{2.607380in}}%
\pgfpathlineto{\pgfqpoint{1.083121in}{2.599509in}}%
\pgfpathlineto{\pgfqpoint{1.103300in}{2.592838in}}%
\pgfpathlineto{\pgfqpoint{1.122683in}{2.587265in}}%
\pgfpathlineto{\pgfqpoint{1.141344in}{2.582698in}}%
\pgfpathlineto{\pgfqpoint{1.159345in}{2.579054in}}%
\pgfpathlineto{\pgfqpoint{1.176741in}{2.576258in}}%
\pgfpathlineto{\pgfqpoint{1.193581in}{2.574241in}}%
\pgfpathlineto{\pgfqpoint{1.209907in}{2.572943in}}%
\pgfpathlineto{\pgfqpoint{1.225759in}{2.572304in}}%
\pgfpathlineto{\pgfqpoint{1.241170in}{2.572274in}}%
\pgfpathlineto{\pgfqpoint{1.256167in}{2.572804in}}%
\pgfpathlineto{\pgfqpoint{1.270779in}{2.573849in}}%
\pgfpathlineto{\pgfqpoint{1.292021in}{2.576292in}}%
\pgfpathlineto{\pgfqpoint{1.312509in}{2.579672in}}%
\pgfpathlineto{\pgfqpoint{1.332299in}{2.583873in}}%
\pgfpathlineto{\pgfqpoint{1.351434in}{2.588789in}}%
\pgfpathlineto{\pgfqpoint{1.369950in}{2.594323in}}%
\pgfpathlineto{\pgfqpoint{1.387876in}{2.600388in}}%
\pgfpathlineto{\pgfqpoint{1.405238in}{2.606901in}}%
\pgfpathlineto{\pgfqpoint{1.427539in}{2.616156in}}%
\pgfpathlineto{\pgfqpoint{1.448905in}{2.625917in}}%
\pgfpathlineto{\pgfqpoint{1.469364in}{2.636042in}}%
\pgfpathlineto{\pgfqpoint{1.493692in}{2.649020in}}%
\pgfpathlineto{\pgfqpoint{1.516669in}{2.662154in}}%
\pgfpathlineto{\pgfqpoint{1.542496in}{2.677868in}}%
\pgfpathlineto{\pgfqpoint{1.566466in}{2.693288in}}%
\pgfpathlineto{\pgfqpoint{1.592151in}{2.710638in}}%
\pgfpathlineto{\pgfqpoint{1.621712in}{2.731580in}}%
\pgfpathlineto{\pgfqpoint{1.650254in}{2.752706in}}%
\pgfpathlineto{\pgfqpoint{1.681196in}{2.776519in}}%
\pgfpathlineto{\pgfqpoint{1.714144in}{2.802818in}}%
\pgfpathlineto{\pgfqpoint{1.749686in}{2.832166in}}%
\pgfpathlineto{\pgfqpoint{1.787588in}{2.864459in}}%
\pgfpathlineto{\pgfqpoint{1.791626in}{2.867954in}}%
\pgfpathlineto{\pgfqpoint{1.791626in}{2.867954in}}%
\pgfusepath{stroke}%
\end{pgfscope}%
\begin{pgfscope}%
\pgfsetroundcap%
\pgfsetroundjoin%
\pgfsetlinewidth{1.254687pt}%
\definecolor{currentstroke}{rgb}{0.478431,0.243137,0.615686}%
\pgfsetstrokecolor{currentstroke}%
\pgfsetdash{}{0pt}%
\pgfpathmoveto{\pgfqpoint{1.780164in}{2.858058in}}%
\pgfpathquadraticcurveto{\pgfqpoint{1.780286in}{2.858163in}}{\pgfqpoint{1.765691in}{2.845613in}}%
\pgfusepath{stroke}%
\end{pgfscope}%
\begin{pgfscope}%
\pgfsetroundcap%
\pgfsetroundjoin%
\definecolor{currentfill}{rgb}{0.478431,0.243137,0.615686}%
\pgfsetfillcolor{currentfill}%
\pgfsetlinewidth{1.254687pt}%
\definecolor{currentstroke}{rgb}{0.478431,0.243137,0.615686}%
\pgfsetstrokecolor{currentstroke}%
\pgfsetdash{}{0pt}%
\pgfpathmoveto{\pgfqpoint{1.729550in}{2.836518in}}%
\pgfpathlineto{\pgfqpoint{1.765691in}{2.845613in}}%
\pgfpathlineto{\pgfqpoint{1.751282in}{2.811243in}}%
\pgfpathlineto{\pgfqpoint{1.729550in}{2.836518in}}%
\pgfpathclose%
\pgfusepath{stroke,fill}%
\end{pgfscope}%
\begin{pgfscope}%
\pgfpathrectangle{\pgfqpoint{0.537500in}{2.059688in}}{\pgfqpoint{1.505000in}{1.156312in}}%
\pgfusepath{clip}%
\pgfsetrectcap%
\pgfsetroundjoin%
\pgfsetlinewidth{2.158062pt}%
\definecolor{currentstroke}{rgb}{0.878431,0.552941,0.176471}%
\pgfsetstrokecolor{currentstroke}%
\pgfsetdash{}{0pt}%
\pgfpathmoveto{\pgfqpoint{1.044741in}{2.791271in}}%
\pgfpathlineto{\pgfqpoint{1.073428in}{2.782582in}}%
\pgfpathlineto{\pgfqpoint{1.100723in}{2.775561in}}%
\pgfpathlineto{\pgfqpoint{1.126755in}{2.770029in}}%
\pgfpathlineto{\pgfqpoint{1.151640in}{2.765829in}}%
\pgfpathlineto{\pgfqpoint{1.175474in}{2.762822in}}%
\pgfpathlineto{\pgfqpoint{1.198341in}{2.760883in}}%
\pgfpathlineto{\pgfqpoint{1.220316in}{2.759900in}}%
\pgfpathlineto{\pgfqpoint{1.241460in}{2.759772in}}%
\pgfpathlineto{\pgfqpoint{1.261829in}{2.760408in}}%
\pgfpathlineto{\pgfqpoint{1.281472in}{2.761725in}}%
\pgfpathlineto{\pgfqpoint{1.300429in}{2.763647in}}%
\pgfpathlineto{\pgfqpoint{1.318737in}{2.766106in}}%
\pgfpathlineto{\pgfqpoint{1.336430in}{2.769038in}}%
\pgfpathlineto{\pgfqpoint{1.361875in}{2.774202in}}%
\pgfpathlineto{\pgfqpoint{1.386079in}{2.780128in}}%
\pgfpathlineto{\pgfqpoint{1.409114in}{2.786665in}}%
\pgfpathlineto{\pgfqpoint{1.431038in}{2.793677in}}%
\pgfpathlineto{\pgfqpoint{1.451904in}{2.801047in}}%
\pgfpathlineto{\pgfqpoint{1.478156in}{2.811251in}}%
\pgfpathlineto{\pgfqpoint{1.502701in}{2.821698in}}%
\pgfpathlineto{\pgfqpoint{1.525626in}{2.832214in}}%
\pgfpathlineto{\pgfqpoint{1.552122in}{2.845240in}}%
\pgfpathlineto{\pgfqpoint{1.576353in}{2.857936in}}%
\pgfpathlineto{\pgfqpoint{1.602635in}{2.872509in}}%
\pgfpathlineto{\pgfqpoint{1.629737in}{2.888366in}}%
\pgfpathlineto{\pgfqpoint{1.659505in}{2.906698in}}%
\pgfpathlineto{\pgfqpoint{1.689135in}{2.925838in}}%
\pgfpathlineto{\pgfqpoint{1.719723in}{2.946470in}}%
\pgfpathlineto{\pgfqpoint{1.752399in}{2.969422in}}%
\pgfpathlineto{\pgfqpoint{1.786587in}{2.994371in}}%
\pgfpathlineto{\pgfqpoint{1.791662in}{2.998151in}}%
\pgfpathlineto{\pgfqpoint{1.791662in}{2.998151in}}%
\pgfusepath{stroke}%
\end{pgfscope}%
\begin{pgfscope}%
\pgfsetroundcap%
\pgfsetroundjoin%
\pgfsetlinewidth{1.254687pt}%
\definecolor{currentstroke}{rgb}{0.878431,0.552941,0.176471}%
\pgfsetstrokecolor{currentstroke}%
\pgfsetdash{}{0pt}%
\pgfpathmoveto{\pgfqpoint{1.788674in}{2.995923in}}%
\pgfpathquadraticcurveto{\pgfqpoint{1.788711in}{2.995951in}}{\pgfqpoint{1.773180in}{2.984386in}}%
\pgfusepath{stroke}%
\end{pgfscope}%
\begin{pgfscope}%
\pgfsetroundcap%
\pgfsetroundjoin%
\definecolor{currentfill}{rgb}{0.878431,0.552941,0.176471}%
\pgfsetfillcolor{currentfill}%
\pgfsetlinewidth{1.254687pt}%
\definecolor{currentstroke}{rgb}{0.878431,0.552941,0.176471}%
\pgfsetstrokecolor{currentstroke}%
\pgfsetdash{}{0pt}%
\pgfpathmoveto{\pgfqpoint{1.736491in}{2.977846in}}%
\pgfpathlineto{\pgfqpoint{1.773180in}{2.984386in}}%
\pgfpathlineto{\pgfqpoint{1.756398in}{2.951111in}}%
\pgfpathlineto{\pgfqpoint{1.736491in}{2.977846in}}%
\pgfpathclose%
\pgfusepath{stroke,fill}%
\end{pgfscope}%
\begin{pgfscope}%
\definecolor{textcolor}{rgb}{0.372549,0.372549,0.372549}%
\pgfsetstrokecolor{textcolor}%
\pgfsetfillcolor{textcolor}%
\pgftext[x=0.576519in,y=3.122245in,left,base]{\color{textcolor}{\rmfamily\fontsize{6.000000}{7.200000}\selectfont\catcode`\^=\active\def^{\ifmmode\sp\else\^{}\fi}\catcode`\%=\active\def%{\%}leaves $s^2-r^2=\alpha^2$}}%
\end{pgfscope}%
\begin{pgfscope}%
\definecolor{textcolor}{rgb}{0.125490,0.498039,0.321569}%
\pgfsetstrokecolor{textcolor}%
\pgfsetfillcolor{textcolor}%
\pgftext[x=1.811176in,y=3.028490in,left,base]{\color{textcolor}{\rmfamily\fontsize{6.000000}{7.200000}\selectfont\catcode`\^=\active\def^{\ifmmode\sp\else\^{}\fi}\catcode`\%=\active\def%{\%}$r=\beta$}}%
\end{pgfscope}%
\begin{pgfscope}%
\definecolor{textcolor}{rgb}{0.000000,0.000000,0.000000}%
\pgfsetstrokecolor{textcolor}%
\pgfsetfillcolor{textcolor}%
\pgftext[x=1.290000in,y=3.243778in,,base]{\color{textcolor}{\rmfamily\fontsize{8.000000}{9.600000}\selectfont\catcode`\^=\active\def^{\ifmmode\sp\else\^{}\fi}\catcode`\%=\active\def%{\%}(a) Equal variances}}%
\end{pgfscope}%
\begin{pgfscope}%
\pgfpathrectangle{\pgfqpoint{0.537500in}{2.059688in}}{\pgfqpoint{1.505000in}{1.156312in}}%
\pgfusepath{clip}%
\pgfsetbuttcap%
\pgfsetroundjoin%
\definecolor{currentfill}{rgb}{0.141176,0.345098,0.650980}%
\pgfsetfillcolor{currentfill}%
\pgfsetlinewidth{1.003750pt}%
\definecolor{currentstroke}{rgb}{0.141176,0.345098,0.650980}%
\pgfsetstrokecolor{currentstroke}%
\pgfsetdash{}{0pt}%
\pgfsys@defobject{currentmarker}{\pgfqpoint{-0.034021in}{-0.034021in}}{\pgfqpoint{0.034021in}{0.034021in}}{%
\pgfpathmoveto{\pgfqpoint{0.000000in}{-0.034021in}}%
\pgfpathcurveto{\pgfqpoint{0.009022in}{-0.034021in}}{\pgfqpoint{0.017676in}{-0.030436in}}{\pgfqpoint{0.024056in}{-0.024056in}}%
\pgfpathcurveto{\pgfqpoint{0.030436in}{-0.017676in}}{\pgfqpoint{0.034021in}{-0.009022in}}{\pgfqpoint{0.034021in}{0.000000in}}%
\pgfpathcurveto{\pgfqpoint{0.034021in}{0.009022in}}{\pgfqpoint{0.030436in}{0.017676in}}{\pgfqpoint{0.024056in}{0.024056in}}%
\pgfpathcurveto{\pgfqpoint{0.017676in}{0.030436in}}{\pgfqpoint{0.009022in}{0.034021in}}{\pgfqpoint{0.000000in}{0.034021in}}%
\pgfpathcurveto{\pgfqpoint{-0.009022in}{0.034021in}}{\pgfqpoint{-0.017676in}{0.030436in}}{\pgfqpoint{-0.024056in}{0.024056in}}%
\pgfpathcurveto{\pgfqpoint{-0.030436in}{0.017676in}}{\pgfqpoint{-0.034021in}{0.009022in}}{\pgfqpoint{-0.034021in}{0.000000in}}%
\pgfpathcurveto{\pgfqpoint{-0.034021in}{-0.009022in}}{\pgfqpoint{-0.030436in}{-0.017676in}}{\pgfqpoint{-0.024056in}{-0.024056in}}%
\pgfpathcurveto{\pgfqpoint{-0.017676in}{-0.030436in}}{\pgfqpoint{-0.009022in}{-0.034021in}}{\pgfqpoint{0.000000in}{-0.034021in}}%
\pgfpathlineto{\pgfqpoint{0.000000in}{-0.034021in}}%
\pgfpathclose%
\pgfusepath{stroke,fill}%
}%
\begin{pgfscope}%
\pgfsys@transformshift{0.648981in}{2.732899in}%
\pgfsys@useobject{currentmarker}{}%
\end{pgfscope}%
\end{pgfscope}%
\begin{pgfscope}%
\pgfpathrectangle{\pgfqpoint{0.537500in}{2.059688in}}{\pgfqpoint{1.505000in}{1.156312in}}%
\pgfusepath{clip}%
\pgfsetbuttcap%
\pgfsetroundjoin%
\definecolor{currentfill}{rgb}{0.141176,0.345098,0.650980}%
\pgfsetfillcolor{currentfill}%
\pgfsetlinewidth{1.003750pt}%
\definecolor{currentstroke}{rgb}{0.141176,0.345098,0.650980}%
\pgfsetstrokecolor{currentstroke}%
\pgfsetdash{}{0pt}%
\pgfsys@defobject{currentmarker}{\pgfqpoint{-0.036747in}{-0.036747in}}{\pgfqpoint{0.036747in}{0.036747in}}{%
\pgfpathmoveto{\pgfqpoint{-0.036747in}{-0.036747in}}%
\pgfpathlineto{\pgfqpoint{0.036747in}{-0.036747in}}%
\pgfpathlineto{\pgfqpoint{0.036747in}{0.036747in}}%
\pgfpathlineto{\pgfqpoint{-0.036747in}{0.036747in}}%
\pgfpathlineto{\pgfqpoint{-0.036747in}{-0.036747in}}%
\pgfpathclose%
\pgfusepath{stroke,fill}%
}%
\begin{pgfscope}%
\pgfsys@transformshift{1.787975in}{2.698446in}%
\pgfsys@useobject{currentmarker}{}%
\end{pgfscope}%
\end{pgfscope}%
\begin{pgfscope}%
\pgfpathrectangle{\pgfqpoint{0.537500in}{2.059688in}}{\pgfqpoint{1.505000in}{1.156312in}}%
\pgfusepath{clip}%
\pgfsetbuttcap%
\pgfsetroundjoin%
\definecolor{currentfill}{rgb}{0.819608,0.286275,0.356863}%
\pgfsetfillcolor{currentfill}%
\pgfsetlinewidth{1.003750pt}%
\definecolor{currentstroke}{rgb}{0.819608,0.286275,0.356863}%
\pgfsetstrokecolor{currentstroke}%
\pgfsetdash{}{0pt}%
\pgfsys@defobject{currentmarker}{\pgfqpoint{-0.034021in}{-0.034021in}}{\pgfqpoint{0.034021in}{0.034021in}}{%
\pgfpathmoveto{\pgfqpoint{0.000000in}{-0.034021in}}%
\pgfpathcurveto{\pgfqpoint{0.009022in}{-0.034021in}}{\pgfqpoint{0.017676in}{-0.030436in}}{\pgfqpoint{0.024056in}{-0.024056in}}%
\pgfpathcurveto{\pgfqpoint{0.030436in}{-0.017676in}}{\pgfqpoint{0.034021in}{-0.009022in}}{\pgfqpoint{0.034021in}{0.000000in}}%
\pgfpathcurveto{\pgfqpoint{0.034021in}{0.009022in}}{\pgfqpoint{0.030436in}{0.017676in}}{\pgfqpoint{0.024056in}{0.024056in}}%
\pgfpathcurveto{\pgfqpoint{0.017676in}{0.030436in}}{\pgfqpoint{0.009022in}{0.034021in}}{\pgfqpoint{0.000000in}{0.034021in}}%
\pgfpathcurveto{\pgfqpoint{-0.009022in}{0.034021in}}{\pgfqpoint{-0.017676in}{0.030436in}}{\pgfqpoint{-0.024056in}{0.024056in}}%
\pgfpathcurveto{\pgfqpoint{-0.030436in}{0.017676in}}{\pgfqpoint{-0.034021in}{0.009022in}}{\pgfqpoint{-0.034021in}{0.000000in}}%
\pgfpathcurveto{\pgfqpoint{-0.034021in}{-0.009022in}}{\pgfqpoint{-0.030436in}{-0.017676in}}{\pgfqpoint{-0.024056in}{-0.024056in}}%
\pgfpathcurveto{\pgfqpoint{-0.017676in}{-0.030436in}}{\pgfqpoint{-0.009022in}{-0.034021in}}{\pgfqpoint{0.000000in}{-0.034021in}}%
\pgfpathlineto{\pgfqpoint{0.000000in}{-0.034021in}}%
\pgfpathclose%
\pgfusepath{stroke,fill}%
}%
\begin{pgfscope}%
\pgfsys@transformshift{0.721444in}{2.681840in}%
\pgfsys@useobject{currentmarker}{}%
\end{pgfscope}%
\end{pgfscope}%
\begin{pgfscope}%
\pgfpathrectangle{\pgfqpoint{0.537500in}{2.059688in}}{\pgfqpoint{1.505000in}{1.156312in}}%
\pgfusepath{clip}%
\pgfsetbuttcap%
\pgfsetroundjoin%
\definecolor{currentfill}{rgb}{0.819608,0.286275,0.356863}%
\pgfsetfillcolor{currentfill}%
\pgfsetlinewidth{1.003750pt}%
\definecolor{currentstroke}{rgb}{0.819608,0.286275,0.356863}%
\pgfsetstrokecolor{currentstroke}%
\pgfsetdash{}{0pt}%
\pgfsys@defobject{currentmarker}{\pgfqpoint{-0.036747in}{-0.036747in}}{\pgfqpoint{0.036747in}{0.036747in}}{%
\pgfpathmoveto{\pgfqpoint{-0.036747in}{-0.036747in}}%
\pgfpathlineto{\pgfqpoint{0.036747in}{-0.036747in}}%
\pgfpathlineto{\pgfqpoint{0.036747in}{0.036747in}}%
\pgfpathlineto{\pgfqpoint{-0.036747in}{0.036747in}}%
\pgfpathlineto{\pgfqpoint{-0.036747in}{-0.036747in}}%
\pgfpathclose%
\pgfusepath{stroke,fill}%
}%
\begin{pgfscope}%
\pgfsys@transformshift{1.790609in}{2.727219in}%
\pgfsys@useobject{currentmarker}{}%
\end{pgfscope}%
\end{pgfscope}%
\begin{pgfscope}%
\pgfpathrectangle{\pgfqpoint{0.537500in}{2.059688in}}{\pgfqpoint{1.505000in}{1.156312in}}%
\pgfusepath{clip}%
\pgfsetbuttcap%
\pgfsetroundjoin%
\definecolor{currentfill}{rgb}{0.000000,0.529412,0.423529}%
\pgfsetfillcolor{currentfill}%
\pgfsetlinewidth{1.003750pt}%
\definecolor{currentstroke}{rgb}{0.000000,0.529412,0.423529}%
\pgfsetstrokecolor{currentstroke}%
\pgfsetdash{}{0pt}%
\pgfsys@defobject{currentmarker}{\pgfqpoint{-0.034021in}{-0.034021in}}{\pgfqpoint{0.034021in}{0.034021in}}{%
\pgfpathmoveto{\pgfqpoint{0.000000in}{-0.034021in}}%
\pgfpathcurveto{\pgfqpoint{0.009022in}{-0.034021in}}{\pgfqpoint{0.017676in}{-0.030436in}}{\pgfqpoint{0.024056in}{-0.024056in}}%
\pgfpathcurveto{\pgfqpoint{0.030436in}{-0.017676in}}{\pgfqpoint{0.034021in}{-0.009022in}}{\pgfqpoint{0.034021in}{0.000000in}}%
\pgfpathcurveto{\pgfqpoint{0.034021in}{0.009022in}}{\pgfqpoint{0.030436in}{0.017676in}}{\pgfqpoint{0.024056in}{0.024056in}}%
\pgfpathcurveto{\pgfqpoint{0.017676in}{0.030436in}}{\pgfqpoint{0.009022in}{0.034021in}}{\pgfqpoint{0.000000in}{0.034021in}}%
\pgfpathcurveto{\pgfqpoint{-0.009022in}{0.034021in}}{\pgfqpoint{-0.017676in}{0.030436in}}{\pgfqpoint{-0.024056in}{0.024056in}}%
\pgfpathcurveto{\pgfqpoint{-0.030436in}{0.017676in}}{\pgfqpoint{-0.034021in}{0.009022in}}{\pgfqpoint{-0.034021in}{0.000000in}}%
\pgfpathcurveto{\pgfqpoint{-0.034021in}{-0.009022in}}{\pgfqpoint{-0.030436in}{-0.017676in}}{\pgfqpoint{-0.024056in}{-0.024056in}}%
\pgfpathcurveto{\pgfqpoint{-0.017676in}{-0.030436in}}{\pgfqpoint{-0.009022in}{-0.034021in}}{\pgfqpoint{0.000000in}{-0.034021in}}%
\pgfpathlineto{\pgfqpoint{0.000000in}{-0.034021in}}%
\pgfpathclose%
\pgfusepath{stroke,fill}%
}%
\begin{pgfscope}%
\pgfsys@transformshift{0.810630in}{2.657241in}%
\pgfsys@useobject{currentmarker}{}%
\end{pgfscope}%
\end{pgfscope}%
\begin{pgfscope}%
\pgfpathrectangle{\pgfqpoint{0.537500in}{2.059688in}}{\pgfqpoint{1.505000in}{1.156312in}}%
\pgfusepath{clip}%
\pgfsetbuttcap%
\pgfsetroundjoin%
\definecolor{currentfill}{rgb}{0.000000,0.529412,0.423529}%
\pgfsetfillcolor{currentfill}%
\pgfsetlinewidth{1.003750pt}%
\definecolor{currentstroke}{rgb}{0.000000,0.529412,0.423529}%
\pgfsetstrokecolor{currentstroke}%
\pgfsetdash{}{0pt}%
\pgfsys@defobject{currentmarker}{\pgfqpoint{-0.036747in}{-0.036747in}}{\pgfqpoint{0.036747in}{0.036747in}}{%
\pgfpathmoveto{\pgfqpoint{-0.036747in}{-0.036747in}}%
\pgfpathlineto{\pgfqpoint{0.036747in}{-0.036747in}}%
\pgfpathlineto{\pgfqpoint{0.036747in}{0.036747in}}%
\pgfpathlineto{\pgfqpoint{-0.036747in}{0.036747in}}%
\pgfpathlineto{\pgfqpoint{-0.036747in}{-0.036747in}}%
\pgfpathclose%
\pgfusepath{stroke,fill}%
}%
\begin{pgfscope}%
\pgfsys@transformshift{1.791433in}{2.782018in}%
\pgfsys@useobject{currentmarker}{}%
\end{pgfscope}%
\end{pgfscope}%
\begin{pgfscope}%
\pgfpathrectangle{\pgfqpoint{0.537500in}{2.059688in}}{\pgfqpoint{1.505000in}{1.156312in}}%
\pgfusepath{clip}%
\pgfsetbuttcap%
\pgfsetroundjoin%
\definecolor{currentfill}{rgb}{0.478431,0.243137,0.615686}%
\pgfsetfillcolor{currentfill}%
\pgfsetlinewidth{1.003750pt}%
\definecolor{currentstroke}{rgb}{0.478431,0.243137,0.615686}%
\pgfsetstrokecolor{currentstroke}%
\pgfsetdash{}{0pt}%
\pgfsys@defobject{currentmarker}{\pgfqpoint{-0.034021in}{-0.034021in}}{\pgfqpoint{0.034021in}{0.034021in}}{%
\pgfpathmoveto{\pgfqpoint{0.000000in}{-0.034021in}}%
\pgfpathcurveto{\pgfqpoint{0.009022in}{-0.034021in}}{\pgfqpoint{0.017676in}{-0.030436in}}{\pgfqpoint{0.024056in}{-0.024056in}}%
\pgfpathcurveto{\pgfqpoint{0.030436in}{-0.017676in}}{\pgfqpoint{0.034021in}{-0.009022in}}{\pgfqpoint{0.034021in}{0.000000in}}%
\pgfpathcurveto{\pgfqpoint{0.034021in}{0.009022in}}{\pgfqpoint{0.030436in}{0.017676in}}{\pgfqpoint{0.024056in}{0.024056in}}%
\pgfpathcurveto{\pgfqpoint{0.017676in}{0.030436in}}{\pgfqpoint{0.009022in}{0.034021in}}{\pgfqpoint{0.000000in}{0.034021in}}%
\pgfpathcurveto{\pgfqpoint{-0.009022in}{0.034021in}}{\pgfqpoint{-0.017676in}{0.030436in}}{\pgfqpoint{-0.024056in}{0.024056in}}%
\pgfpathcurveto{\pgfqpoint{-0.030436in}{0.017676in}}{\pgfqpoint{-0.034021in}{0.009022in}}{\pgfqpoint{-0.034021in}{0.000000in}}%
\pgfpathcurveto{\pgfqpoint{-0.034021in}{-0.009022in}}{\pgfqpoint{-0.030436in}{-0.017676in}}{\pgfqpoint{-0.024056in}{-0.024056in}}%
\pgfpathcurveto{\pgfqpoint{-0.017676in}{-0.030436in}}{\pgfqpoint{-0.009022in}{-0.034021in}}{\pgfqpoint{0.000000in}{-0.034021in}}%
\pgfpathlineto{\pgfqpoint{0.000000in}{-0.034021in}}%
\pgfpathclose%
\pgfusepath{stroke,fill}%
}%
\begin{pgfscope}%
\pgfsys@transformshift{0.910963in}{2.687466in}%
\pgfsys@useobject{currentmarker}{}%
\end{pgfscope}%
\end{pgfscope}%
\begin{pgfscope}%
\pgfpathrectangle{\pgfqpoint{0.537500in}{2.059688in}}{\pgfqpoint{1.505000in}{1.156312in}}%
\pgfusepath{clip}%
\pgfsetbuttcap%
\pgfsetroundjoin%
\definecolor{currentfill}{rgb}{0.478431,0.243137,0.615686}%
\pgfsetfillcolor{currentfill}%
\pgfsetlinewidth{1.003750pt}%
\definecolor{currentstroke}{rgb}{0.478431,0.243137,0.615686}%
\pgfsetstrokecolor{currentstroke}%
\pgfsetdash{}{0pt}%
\pgfsys@defobject{currentmarker}{\pgfqpoint{-0.036747in}{-0.036747in}}{\pgfqpoint{0.036747in}{0.036747in}}{%
\pgfpathmoveto{\pgfqpoint{-0.036747in}{-0.036747in}}%
\pgfpathlineto{\pgfqpoint{0.036747in}{-0.036747in}}%
\pgfpathlineto{\pgfqpoint{0.036747in}{0.036747in}}%
\pgfpathlineto{\pgfqpoint{-0.036747in}{0.036747in}}%
\pgfpathlineto{\pgfqpoint{-0.036747in}{-0.036747in}}%
\pgfpathclose%
\pgfusepath{stroke,fill}%
}%
\begin{pgfscope}%
\pgfsys@transformshift{1.791626in}{2.867954in}%
\pgfsys@useobject{currentmarker}{}%
\end{pgfscope}%
\end{pgfscope}%
\begin{pgfscope}%
\pgfpathrectangle{\pgfqpoint{0.537500in}{2.059688in}}{\pgfqpoint{1.505000in}{1.156312in}}%
\pgfusepath{clip}%
\pgfsetbuttcap%
\pgfsetroundjoin%
\definecolor{currentfill}{rgb}{0.878431,0.552941,0.176471}%
\pgfsetfillcolor{currentfill}%
\pgfsetlinewidth{1.003750pt}%
\definecolor{currentstroke}{rgb}{0.878431,0.552941,0.176471}%
\pgfsetstrokecolor{currentstroke}%
\pgfsetdash{}{0pt}%
\pgfsys@defobject{currentmarker}{\pgfqpoint{-0.034021in}{-0.034021in}}{\pgfqpoint{0.034021in}{0.034021in}}{%
\pgfpathmoveto{\pgfqpoint{0.000000in}{-0.034021in}}%
\pgfpathcurveto{\pgfqpoint{0.009022in}{-0.034021in}}{\pgfqpoint{0.017676in}{-0.030436in}}{\pgfqpoint{0.024056in}{-0.024056in}}%
\pgfpathcurveto{\pgfqpoint{0.030436in}{-0.017676in}}{\pgfqpoint{0.034021in}{-0.009022in}}{\pgfqpoint{0.034021in}{0.000000in}}%
\pgfpathcurveto{\pgfqpoint{0.034021in}{0.009022in}}{\pgfqpoint{0.030436in}{0.017676in}}{\pgfqpoint{0.024056in}{0.024056in}}%
\pgfpathcurveto{\pgfqpoint{0.017676in}{0.030436in}}{\pgfqpoint{0.009022in}{0.034021in}}{\pgfqpoint{0.000000in}{0.034021in}}%
\pgfpathcurveto{\pgfqpoint{-0.009022in}{0.034021in}}{\pgfqpoint{-0.017676in}{0.030436in}}{\pgfqpoint{-0.024056in}{0.024056in}}%
\pgfpathcurveto{\pgfqpoint{-0.030436in}{0.017676in}}{\pgfqpoint{-0.034021in}{0.009022in}}{\pgfqpoint{-0.034021in}{0.000000in}}%
\pgfpathcurveto{\pgfqpoint{-0.034021in}{-0.009022in}}{\pgfqpoint{-0.030436in}{-0.017676in}}{\pgfqpoint{-0.024056in}{-0.024056in}}%
\pgfpathcurveto{\pgfqpoint{-0.017676in}{-0.030436in}}{\pgfqpoint{-0.009022in}{-0.034021in}}{\pgfqpoint{0.000000in}{-0.034021in}}%
\pgfpathlineto{\pgfqpoint{0.000000in}{-0.034021in}}%
\pgfpathclose%
\pgfusepath{stroke,fill}%
}%
\begin{pgfscope}%
\pgfsys@transformshift{1.044741in}{2.791271in}%
\pgfsys@useobject{currentmarker}{}%
\end{pgfscope}%
\end{pgfscope}%
\begin{pgfscope}%
\pgfpathrectangle{\pgfqpoint{0.537500in}{2.059688in}}{\pgfqpoint{1.505000in}{1.156312in}}%
\pgfusepath{clip}%
\pgfsetbuttcap%
\pgfsetroundjoin%
\definecolor{currentfill}{rgb}{0.878431,0.552941,0.176471}%
\pgfsetfillcolor{currentfill}%
\pgfsetlinewidth{1.003750pt}%
\definecolor{currentstroke}{rgb}{0.878431,0.552941,0.176471}%
\pgfsetstrokecolor{currentstroke}%
\pgfsetdash{}{0pt}%
\pgfsys@defobject{currentmarker}{\pgfqpoint{-0.036747in}{-0.036747in}}{\pgfqpoint{0.036747in}{0.036747in}}{%
\pgfpathmoveto{\pgfqpoint{-0.036747in}{-0.036747in}}%
\pgfpathlineto{\pgfqpoint{0.036747in}{-0.036747in}}%
\pgfpathlineto{\pgfqpoint{0.036747in}{0.036747in}}%
\pgfpathlineto{\pgfqpoint{-0.036747in}{0.036747in}}%
\pgfpathlineto{\pgfqpoint{-0.036747in}{-0.036747in}}%
\pgfpathclose%
\pgfusepath{stroke,fill}%
}%
\begin{pgfscope}%
\pgfsys@transformshift{1.791662in}{2.998151in}%
\pgfsys@useobject{currentmarker}{}%
\end{pgfscope}%
\end{pgfscope}%
\begin{pgfscope}%
\pgfpathrectangle{\pgfqpoint{0.537500in}{2.059688in}}{\pgfqpoint{1.505000in}{1.156312in}}%
\pgfusepath{clip}%
\pgfsetbuttcap%
\pgfsetroundjoin%
\definecolor{currentfill}{rgb}{0.168627,0.611765,0.411765}%
\pgfsetfillcolor{currentfill}%
\pgfsetlinewidth{1.003750pt}%
\definecolor{currentstroke}{rgb}{0.168627,0.611765,0.411765}%
\pgfsetstrokecolor{currentstroke}%
\pgfsetdash{}{0pt}%
\pgfsys@defobject{currentmarker}{\pgfqpoint{-0.034021in}{-0.034021in}}{\pgfqpoint{0.034021in}{0.034021in}}{%
\pgfpathmoveto{\pgfqpoint{0.000000in}{-0.034021in}}%
\pgfpathcurveto{\pgfqpoint{0.009022in}{-0.034021in}}{\pgfqpoint{0.017676in}{-0.030436in}}{\pgfqpoint{0.024056in}{-0.024056in}}%
\pgfpathcurveto{\pgfqpoint{0.030436in}{-0.017676in}}{\pgfqpoint{0.034021in}{-0.009022in}}{\pgfqpoint{0.034021in}{0.000000in}}%
\pgfpathcurveto{\pgfqpoint{0.034021in}{0.009022in}}{\pgfqpoint{0.030436in}{0.017676in}}{\pgfqpoint{0.024056in}{0.024056in}}%
\pgfpathcurveto{\pgfqpoint{0.017676in}{0.030436in}}{\pgfqpoint{0.009022in}{0.034021in}}{\pgfqpoint{0.000000in}{0.034021in}}%
\pgfpathcurveto{\pgfqpoint{-0.009022in}{0.034021in}}{\pgfqpoint{-0.017676in}{0.030436in}}{\pgfqpoint{-0.024056in}{0.024056in}}%
\pgfpathcurveto{\pgfqpoint{-0.030436in}{0.017676in}}{\pgfqpoint{-0.034021in}{0.009022in}}{\pgfqpoint{-0.034021in}{0.000000in}}%
\pgfpathcurveto{\pgfqpoint{-0.034021in}{-0.009022in}}{\pgfqpoint{-0.030436in}{-0.017676in}}{\pgfqpoint{-0.024056in}{-0.024056in}}%
\pgfpathcurveto{\pgfqpoint{-0.017676in}{-0.030436in}}{\pgfqpoint{-0.009022in}{-0.034021in}}{\pgfqpoint{0.000000in}{-0.034021in}}%
\pgfpathlineto{\pgfqpoint{0.000000in}{-0.034021in}}%
\pgfpathclose%
\pgfusepath{stroke,fill}%
}%
\begin{pgfscope}%
\pgfsys@transformshift{1.791667in}{2.684722in}%
\pgfsys@useobject{currentmarker}{}%
\end{pgfscope}%
\end{pgfscope}%
\begin{pgfscope}%
\pgfsetbuttcap%
\pgfsetmiterjoin%
\definecolor{currentfill}{rgb}{1.000000,1.000000,1.000000}%
\pgfsetfillcolor{currentfill}%
\pgfsetlinewidth{0.000000pt}%
\definecolor{currentstroke}{rgb}{0.000000,0.000000,0.000000}%
\pgfsetstrokecolor{currentstroke}%
\pgfsetstrokeopacity{0.000000}%
\pgfsetdash{}{0pt}%
\pgfpathmoveto{\pgfqpoint{0.407592in}{0.222904in}}%
\pgfpathlineto{\pgfqpoint{1.818754in}{0.222904in}}%
\pgfpathlineto{\pgfqpoint{1.818754in}{1.634067in}}%
\pgfpathlineto{\pgfqpoint{0.407592in}{1.634067in}}%
\pgfpathlineto{\pgfqpoint{0.407592in}{0.222904in}}%
\pgfpathclose%
\pgfusepath{fill}%
\end{pgfscope}%
\begin{pgfscope}%
\pgfsetbuttcap%
\pgfsetmiterjoin%
\pgfsetlinewidth{0.000000pt}%
\definecolor{currentstroke}{rgb}{0.950000,0.950000,0.950000}%
\pgfsetstrokecolor{currentstroke}%
\pgfsetstrokeopacity{0.000000}%
\pgfsetdash{}{0pt}%
\pgfpathmoveto{\pgfqpoint{0.418841in}{0.578914in}}%
\pgfpathlineto{\pgfqpoint{0.882898in}{0.862989in}}%
\pgfpathlineto{\pgfqpoint{0.873607in}{1.529400in}}%
\pgfpathlineto{\pgfqpoint{0.389895in}{1.293091in}}%
\pgfusepath{}%
\end{pgfscope}%
\begin{pgfscope}%
\pgfsetbuttcap%
\pgfsetmiterjoin%
\pgfsetlinewidth{0.000000pt}%
\definecolor{currentstroke}{rgb}{0.900000,0.900000,0.900000}%
\pgfsetstrokecolor{currentstroke}%
\pgfsetstrokeopacity{0.000000}%
\pgfsetdash{}{0pt}%
\pgfpathmoveto{\pgfqpoint{0.882898in}{0.862989in}}%
\pgfpathlineto{\pgfqpoint{1.837675in}{0.619202in}}%
\pgfpathlineto{\pgfqpoint{1.865965in}{1.326696in}}%
\pgfpathlineto{\pgfqpoint{0.873607in}{1.529400in}}%
\pgfusepath{}%
\end{pgfscope}%
\begin{pgfscope}%
\pgfsetbuttcap%
\pgfsetmiterjoin%
\pgfsetlinewidth{0.000000pt}%
\definecolor{currentstroke}{rgb}{0.925000,0.925000,0.925000}%
\pgfsetstrokecolor{currentstroke}%
\pgfsetstrokeopacity{0.000000}%
\pgfsetdash{}{0pt}%
\pgfpathmoveto{\pgfqpoint{0.418841in}{0.578914in}}%
\pgfpathlineto{\pgfqpoint{1.424602in}{0.289843in}}%
\pgfpathlineto{\pgfqpoint{1.837675in}{0.619202in}}%
\pgfpathlineto{\pgfqpoint{0.882898in}{0.862989in}}%
\pgfusepath{}%
\end{pgfscope}%
\begin{pgfscope}%
\pgfsetrectcap%
\pgfsetroundjoin%
\pgfsetlinewidth{0.803000pt}%
\definecolor{currentstroke}{rgb}{0.000000,0.000000,0.000000}%
\pgfsetstrokecolor{currentstroke}%
\pgfsetdash{}{0pt}%
\pgfpathmoveto{\pgfqpoint{0.418841in}{0.578914in}}%
\pgfpathlineto{\pgfqpoint{1.424602in}{0.289843in}}%
\pgfusepath{stroke}%
\end{pgfscope}%
\begin{pgfscope}%
\pgfsetrectcap%
\pgfsetroundjoin%
\pgfsetlinewidth{0.803000pt}%
\definecolor{currentstroke}{rgb}{0.000000,0.000000,0.000000}%
\pgfsetstrokecolor{currentstroke}%
\pgfsetdash{}{0pt}%
\pgfpathmoveto{\pgfqpoint{0.500034in}{0.559223in}}%
\pgfpathlineto{\pgfqpoint{0.488021in}{0.551726in}}%
\pgfusepath{stroke}%
\end{pgfscope}%
\begin{pgfscope}%
\definecolor{textcolor}{rgb}{0.000000,0.000000,0.000000}%
\pgfsetstrokecolor{textcolor}%
\pgfsetfillcolor{textcolor}%
\pgftext[x=0.432780in,y=0.423379in,,top]{\color{textcolor}{\rmfamily\fontsize{5.000000}{6.000000}\selectfont\catcode`\^=\active\def^{\ifmmode\sp\else\^{}\fi}\catcode`\%=\active\def%{\%}\ensuremath{-}1.0}}%
\end{pgfscope}%
\begin{pgfscope}%
\pgfsetrectcap%
\pgfsetroundjoin%
\pgfsetlinewidth{0.803000pt}%
\definecolor{currentstroke}{rgb}{0.000000,0.000000,0.000000}%
\pgfsetstrokecolor{currentstroke}%
\pgfsetdash{}{0pt}%
\pgfpathmoveto{\pgfqpoint{0.678989in}{0.507840in}}%
\pgfpathlineto{\pgfqpoint{0.667181in}{0.500129in}}%
\pgfusepath{stroke}%
\end{pgfscope}%
\begin{pgfscope}%
\definecolor{textcolor}{rgb}{0.000000,0.000000,0.000000}%
\pgfsetstrokecolor{textcolor}%
\pgfsetfillcolor{textcolor}%
\pgftext[x=0.612032in,y=0.369697in,,top]{\color{textcolor}{\rmfamily\fontsize{5.000000}{6.000000}\selectfont\catcode`\^=\active\def^{\ifmmode\sp\else\^{}\fi}\catcode`\%=\active\def%{\%}\ensuremath{-}0.5}}%
\end{pgfscope}%
\begin{pgfscope}%
\pgfsetrectcap%
\pgfsetroundjoin%
\pgfsetlinewidth{0.803000pt}%
\definecolor{currentstroke}{rgb}{0.000000,0.000000,0.000000}%
\pgfsetstrokecolor{currentstroke}%
\pgfsetdash{}{0pt}%
\pgfpathmoveto{\pgfqpoint{0.863089in}{0.454980in}}%
\pgfpathlineto{\pgfqpoint{0.851503in}{0.447045in}}%
\pgfusepath{stroke}%
\end{pgfscope}%
\begin{pgfscope}%
\definecolor{textcolor}{rgb}{0.000000,0.000000,0.000000}%
\pgfsetstrokecolor{textcolor}%
\pgfsetfillcolor{textcolor}%
\pgftext[x=0.796474in,y=0.314461in,,top]{\color{textcolor}{\rmfamily\fontsize{5.000000}{6.000000}\selectfont\catcode`\^=\active\def^{\ifmmode\sp\else\^{}\fi}\catcode`\%=\active\def%{\%}0.0}}%
\end{pgfscope}%
\begin{pgfscope}%
\pgfsetrectcap%
\pgfsetroundjoin%
\pgfsetlinewidth{0.803000pt}%
\definecolor{currentstroke}{rgb}{0.000000,0.000000,0.000000}%
\pgfsetstrokecolor{currentstroke}%
\pgfsetdash{}{0pt}%
\pgfpathmoveto{\pgfqpoint{1.052559in}{0.400579in}}%
\pgfpathlineto{\pgfqpoint{1.041212in}{0.392410in}}%
\pgfusepath{stroke}%
\end{pgfscope}%
\begin{pgfscope}%
\definecolor{textcolor}{rgb}{0.000000,0.000000,0.000000}%
\pgfsetstrokecolor{textcolor}%
\pgfsetfillcolor{textcolor}%
\pgftext[x=0.986334in,y=0.257602in,,top]{\color{textcolor}{\rmfamily\fontsize{5.000000}{6.000000}\selectfont\catcode`\^=\active\def^{\ifmmode\sp\else\^{}\fi}\catcode`\%=\active\def%{\%}0.5}}%
\end{pgfscope}%
\begin{pgfscope}%
\pgfsetrectcap%
\pgfsetroundjoin%
\pgfsetlinewidth{0.803000pt}%
\definecolor{currentstroke}{rgb}{0.000000,0.000000,0.000000}%
\pgfsetstrokecolor{currentstroke}%
\pgfsetdash{}{0pt}%
\pgfpathmoveto{\pgfqpoint{1.247636in}{0.344567in}}%
\pgfpathlineto{\pgfqpoint{1.236548in}{0.336154in}}%
\pgfusepath{stroke}%
\end{pgfscope}%
\begin{pgfscope}%
\definecolor{textcolor}{rgb}{0.000000,0.000000,0.000000}%
\pgfsetstrokecolor{textcolor}%
\pgfsetfillcolor{textcolor}%
\pgftext[x=1.181855in,y=0.199048in,,top]{\color{textcolor}{\rmfamily\fontsize{5.000000}{6.000000}\selectfont\catcode`\^=\active\def^{\ifmmode\sp\else\^{}\fi}\catcode`\%=\active\def%{\%}1.0}}%
\end{pgfscope}%
\begin{pgfscope}%
\definecolor{textcolor}{rgb}{0.000000,0.000000,0.000000}%
\pgfsetstrokecolor{textcolor}%
\pgfsetfillcolor{textcolor}%
\pgftext[x=0.757857in,y=0.120925in,,]{\color{textcolor}{\rmfamily\fontsize{6.000000}{7.200000}\selectfont\catcode`\^=\active\def^{\ifmmode\sp\else\^{}\fi}\catcode`\%=\active\def%{\%}$r$}}%
\end{pgfscope}%
\begin{pgfscope}%
\pgfsetrectcap%
\pgfsetroundjoin%
\pgfsetlinewidth{0.803000pt}%
\definecolor{currentstroke}{rgb}{0.000000,0.000000,0.000000}%
\pgfsetstrokecolor{currentstroke}%
\pgfsetdash{}{0pt}%
\pgfpathmoveto{\pgfqpoint{1.837675in}{0.619202in}}%
\pgfpathlineto{\pgfqpoint{1.424602in}{0.289843in}}%
\pgfusepath{stroke}%
\end{pgfscope}%
\begin{pgfscope}%
\pgfsetrectcap%
\pgfsetroundjoin%
\pgfsetlinewidth{0.803000pt}%
\definecolor{currentstroke}{rgb}{0.000000,0.000000,0.000000}%
\pgfsetstrokecolor{currentstroke}%
\pgfsetdash{}{0pt}%
\pgfpathmoveto{\pgfqpoint{1.516903in}{0.372665in}}%
\pgfpathlineto{\pgfqpoint{1.542643in}{0.365473in}}%
\pgfusepath{stroke}%
\end{pgfscope}%
\begin{pgfscope}%
\definecolor{textcolor}{rgb}{0.000000,0.000000,0.000000}%
\pgfsetstrokecolor{textcolor}%
\pgfsetfillcolor{textcolor}%
\pgftext[x=1.665803in,y=0.235089in,,top]{\color{textcolor}{\rmfamily\fontsize{5.000000}{6.000000}\selectfont\catcode`\^=\active\def^{\ifmmode\sp\else\^{}\fi}\catcode`\%=\active\def%{\%}\ensuremath{-}0.5}}%
\end{pgfscope}%
\begin{pgfscope}%
\pgfsetrectcap%
\pgfsetroundjoin%
\pgfsetlinewidth{0.803000pt}%
\definecolor{currentstroke}{rgb}{0.000000,0.000000,0.000000}%
\pgfsetstrokecolor{currentstroke}%
\pgfsetdash{}{0pt}%
\pgfpathmoveto{\pgfqpoint{1.631821in}{0.464072in}}%
\pgfpathlineto{\pgfqpoint{1.657155in}{0.457221in}}%
\pgfusepath{stroke}%
\end{pgfscope}%
\begin{pgfscope}%
\definecolor{textcolor}{rgb}{0.000000,0.000000,0.000000}%
\pgfsetstrokecolor{textcolor}%
\pgfsetfillcolor{textcolor}%
\pgftext[x=1.777758in,y=0.330311in,,top]{\color{textcolor}{\rmfamily\fontsize{5.000000}{6.000000}\selectfont\catcode`\^=\active\def^{\ifmmode\sp\else\^{}\fi}\catcode`\%=\active\def%{\%}0.0}}%
\end{pgfscope}%
\begin{pgfscope}%
\pgfsetrectcap%
\pgfsetroundjoin%
\pgfsetlinewidth{0.803000pt}%
\definecolor{currentstroke}{rgb}{0.000000,0.000000,0.000000}%
\pgfsetstrokecolor{currentstroke}%
\pgfsetdash{}{0pt}%
\pgfpathmoveto{\pgfqpoint{1.741358in}{0.551200in}}%
\pgfpathlineto{\pgfqpoint{1.766295in}{0.544666in}}%
\pgfusepath{stroke}%
\end{pgfscope}%
\begin{pgfscope}%
\definecolor{textcolor}{rgb}{0.000000,0.000000,0.000000}%
\pgfsetstrokecolor{textcolor}%
\pgfsetfillcolor{textcolor}%
\pgftext[x=1.884443in,y=0.421050in,,top]{\color{textcolor}{\rmfamily\fontsize{5.000000}{6.000000}\selectfont\catcode`\^=\active\def^{\ifmmode\sp\else\^{}\fi}\catcode`\%=\active\def%{\%}0.5}}%
\end{pgfscope}%
\begin{pgfscope}%
\definecolor{textcolor}{rgb}{0.000000,0.000000,0.000000}%
\pgfsetstrokecolor{textcolor}%
\pgfsetfillcolor{textcolor}%
\pgftext[x=1.976512in,y=0.140262in,,]{\color{textcolor}{\rmfamily\fontsize{6.000000}{7.200000}\selectfont\catcode`\^=\active\def^{\ifmmode\sp\else\^{}\fi}\catcode`\%=\active\def%{\%}$b$}}%
\end{pgfscope}%
\begin{pgfscope}%
\pgfsetrectcap%
\pgfsetroundjoin%
\pgfsetlinewidth{0.803000pt}%
\definecolor{currentstroke}{rgb}{0.000000,0.000000,0.000000}%
\pgfsetstrokecolor{currentstroke}%
\pgfsetdash{}{0pt}%
\pgfpathmoveto{\pgfqpoint{1.837675in}{0.619202in}}%
\pgfpathlineto{\pgfqpoint{1.865965in}{1.326696in}}%
\pgfusepath{stroke}%
\end{pgfscope}%
\begin{pgfscope}%
\pgfsetrectcap%
\pgfsetroundjoin%
\pgfsetlinewidth{0.803000pt}%
\definecolor{currentstroke}{rgb}{0.000000,0.000000,0.000000}%
\pgfsetstrokecolor{currentstroke}%
\pgfsetdash{}{0pt}%
\pgfpathmoveto{\pgfqpoint{1.829482in}{0.621294in}}%
\pgfpathlineto{\pgfqpoint{1.854092in}{0.615010in}}%
\pgfusepath{stroke}%
\end{pgfscope}%
\begin{pgfscope}%
\definecolor{textcolor}{rgb}{0.000000,0.000000,0.000000}%
\pgfsetstrokecolor{textcolor}%
\pgfsetfillcolor{textcolor}%
\pgftext[x=2.022987in,y=0.624464in,,top]{\color{textcolor}{\rmfamily\fontsize{5.000000}{6.000000}\selectfont\catcode`\^=\active\def^{\ifmmode\sp\else\^{}\fi}\catcode`\%=\active\def%{\%}0.0}}%
\end{pgfscope}%
\begin{pgfscope}%
\pgfsetrectcap%
\pgfsetroundjoin%
\pgfsetlinewidth{0.803000pt}%
\definecolor{currentstroke}{rgb}{0.000000,0.000000,0.000000}%
\pgfsetstrokecolor{currentstroke}%
\pgfsetdash{}{0pt}%
\pgfpathmoveto{\pgfqpoint{1.836942in}{0.810078in}}%
\pgfpathlineto{\pgfqpoint{1.861829in}{0.804063in}}%
\pgfusepath{stroke}%
\end{pgfscope}%
\begin{pgfscope}%
\definecolor{textcolor}{rgb}{0.000000,0.000000,0.000000}%
\pgfsetstrokecolor{textcolor}%
\pgfsetfillcolor{textcolor}%
\pgftext[x=2.032509in,y=0.813112in,,top]{\color{textcolor}{\rmfamily\fontsize{5.000000}{6.000000}\selectfont\catcode`\^=\active\def^{\ifmmode\sp\else\^{}\fi}\catcode`\%=\active\def%{\%}0.5}}%
\end{pgfscope}%
\begin{pgfscope}%
\pgfsetrectcap%
\pgfsetroundjoin%
\pgfsetlinewidth{0.803000pt}%
\definecolor{currentstroke}{rgb}{0.000000,0.000000,0.000000}%
\pgfsetstrokecolor{currentstroke}%
\pgfsetdash{}{0pt}%
\pgfpathmoveto{\pgfqpoint{1.844563in}{1.002944in}}%
\pgfpathlineto{\pgfqpoint{1.869734in}{0.997212in}}%
\pgfusepath{stroke}%
\end{pgfscope}%
\begin{pgfscope}%
\definecolor{textcolor}{rgb}{0.000000,0.000000,0.000000}%
\pgfsetstrokecolor{textcolor}%
\pgfsetfillcolor{textcolor}%
\pgftext[x=2.042236in,y=1.005836in,,top]{\color{textcolor}{\rmfamily\fontsize{5.000000}{6.000000}\selectfont\catcode`\^=\active\def^{\ifmmode\sp\else\^{}\fi}\catcode`\%=\active\def%{\%}1.0}}%
\end{pgfscope}%
\begin{pgfscope}%
\pgfsetrectcap%
\pgfsetroundjoin%
\pgfsetlinewidth{0.803000pt}%
\definecolor{currentstroke}{rgb}{0.000000,0.000000,0.000000}%
\pgfsetstrokecolor{currentstroke}%
\pgfsetdash{}{0pt}%
\pgfpathmoveto{\pgfqpoint{1.852351in}{1.200029in}}%
\pgfpathlineto{\pgfqpoint{1.877812in}{1.194592in}}%
\pgfusepath{stroke}%
\end{pgfscope}%
\begin{pgfscope}%
\definecolor{textcolor}{rgb}{0.000000,0.000000,0.000000}%
\pgfsetstrokecolor{textcolor}%
\pgfsetfillcolor{textcolor}%
\pgftext[x=2.052177in,y=1.202771in,,top]{\color{textcolor}{\rmfamily\fontsize{5.000000}{6.000000}\selectfont\catcode`\^=\active\def^{\ifmmode\sp\else\^{}\fi}\catcode`\%=\active\def%{\%}1.5}}%
\end{pgfscope}%
\begin{pgfscope}%
\definecolor{textcolor}{rgb}{0.000000,0.000000,0.000000}%
\pgfsetstrokecolor{textcolor}%
\pgfsetfillcolor{textcolor}%
\pgftext[x=2.228830in,y=0.975668in,,]{\color{textcolor}{\rmfamily\fontsize{6.000000}{7.200000}\selectfont\catcode`\^=\active\def^{\ifmmode\sp\else\^{}\fi}\catcode`\%=\active\def%{\%}$a$}}%
\end{pgfscope}%
\begin{pgfscope}%
\pgfpathrectangle{\pgfqpoint{0.407592in}{0.222904in}}{\pgfqpoint{1.411163in}{1.411163in}}%
\pgfusepath{clip}%
\pgfsetrectcap%
\pgfsetroundjoin%
\pgfsetlinewidth{0.451687pt}%
\definecolor{currentstroke}{rgb}{0.560784,0.588235,0.611765}%
\pgfsetstrokecolor{currentstroke}%
\pgfsetstrokeopacity{0.450000}%
\pgfsetdash{}{0pt}%
\pgfpathmoveto{\pgfqpoint{1.257376in}{0.737514in}}%
\pgfpathlineto{\pgfqpoint{1.265273in}{0.743397in}}%
\pgfpathlineto{\pgfqpoint{1.272040in}{0.749493in}}%
\pgfpathlineto{\pgfqpoint{1.277643in}{0.755763in}}%
\pgfpathlineto{\pgfqpoint{1.282050in}{0.762165in}}%
\pgfpathlineto{\pgfqpoint{1.285242in}{0.768659in}}%
\pgfpathlineto{\pgfqpoint{1.286688in}{0.773018in}}%
\pgfpathlineto{\pgfqpoint{1.287585in}{0.777387in}}%
\pgfpathlineto{\pgfqpoint{1.287934in}{0.781754in}}%
\pgfpathlineto{\pgfqpoint{1.287735in}{0.786108in}}%
\pgfpathlineto{\pgfqpoint{1.286991in}{0.790435in}}%
\pgfpathlineto{\pgfqpoint{1.285706in}{0.794724in}}%
\pgfpathlineto{\pgfqpoint{1.283886in}{0.798964in}}%
\pgfpathlineto{\pgfqpoint{1.281538in}{0.803143in}}%
\pgfpathlineto{\pgfqpoint{1.277043in}{0.809273in}}%
\pgfpathlineto{\pgfqpoint{1.271412in}{0.815205in}}%
\pgfpathlineto{\pgfqpoint{1.264686in}{0.820905in}}%
\pgfpathlineto{\pgfqpoint{1.256908in}{0.826340in}}%
\pgfpathlineto{\pgfqpoint{1.248132in}{0.831477in}}%
\pgfpathlineto{\pgfqpoint{1.238414in}{0.836289in}}%
\pgfpathlineto{\pgfqpoint{1.227816in}{0.840747in}}%
\pgfpathlineto{\pgfqpoint{1.216402in}{0.844828in}}%
\pgfpathlineto{\pgfqpoint{1.204244in}{0.848508in}}%
\pgfpathlineto{\pgfqpoint{1.191415in}{0.851768in}}%
\pgfpathlineto{\pgfqpoint{1.177991in}{0.854589in}}%
\pgfpathlineto{\pgfqpoint{1.159305in}{0.857643in}}%
\pgfpathlineto{\pgfqpoint{1.139899in}{0.859862in}}%
\pgfpathlineto{\pgfqpoint{1.119974in}{0.861224in}}%
\pgfpathlineto{\pgfqpoint{1.099735in}{0.861716in}}%
\pgfpathlineto{\pgfqpoint{1.079390in}{0.861334in}}%
\pgfpathlineto{\pgfqpoint{1.059147in}{0.860082in}}%
\pgfpathlineto{\pgfqpoint{1.039215in}{0.857971in}}%
\pgfpathlineto{\pgfqpoint{1.019800in}{0.855021in}}%
\pgfpathlineto{\pgfqpoint{1.001103in}{0.851261in}}%
\pgfpathlineto{\pgfqpoint{0.983321in}{0.846727in}}%
\pgfpathlineto{\pgfqpoint{0.970698in}{0.842844in}}%
\pgfpathlineto{\pgfqpoint{0.958773in}{0.838572in}}%
\pgfpathlineto{\pgfqpoint{0.947617in}{0.833935in}}%
\pgfpathlineto{\pgfqpoint{0.937302in}{0.828957in}}%
\pgfpathlineto{\pgfqpoint{0.927891in}{0.823668in}}%
\pgfpathlineto{\pgfqpoint{0.919445in}{0.818098in}}%
\pgfpathlineto{\pgfqpoint{0.912020in}{0.812278in}}%
\pgfpathlineto{\pgfqpoint{0.905666in}{0.806242in}}%
\pgfpathlineto{\pgfqpoint{0.900427in}{0.800027in}}%
\pgfpathlineto{\pgfqpoint{0.896342in}{0.793668in}}%
\pgfpathlineto{\pgfqpoint{0.893441in}{0.787205in}}%
\pgfpathlineto{\pgfqpoint{0.892179in}{0.782857in}}%
\pgfpathlineto{\pgfqpoint{0.891460in}{0.778492in}}%
\pgfpathlineto{\pgfqpoint{0.891288in}{0.774122in}}%
\pgfpathlineto{\pgfqpoint{0.891666in}{0.769759in}}%
\pgfpathlineto{\pgfqpoint{0.892596in}{0.765416in}}%
\pgfpathlineto{\pgfqpoint{0.894076in}{0.761103in}}%
\pgfpathlineto{\pgfqpoint{0.896104in}{0.756834in}}%
\pgfpathlineto{\pgfqpoint{0.898678in}{0.752621in}}%
\pgfpathlineto{\pgfqpoint{0.903548in}{0.746433in}}%
\pgfpathlineto{\pgfqpoint{0.909607in}{0.740438in}}%
\pgfpathlineto{\pgfqpoint{0.916822in}{0.734676in}}%
\pgfpathlineto{\pgfqpoint{0.925152in}{0.729186in}}%
\pgfpathlineto{\pgfqpoint{0.934549in}{0.724004in}}%
\pgfpathlineto{\pgfqpoint{0.944956in}{0.719165in}}%
\pgfpathlineto{\pgfqpoint{0.956310in}{0.714704in}}%
\pgfpathlineto{\pgfqpoint{0.968538in}{0.710650in}}%
\pgfpathlineto{\pgfqpoint{0.981563in}{0.707032in}}%
\pgfpathlineto{\pgfqpoint{0.995303in}{0.703875in}}%
\pgfpathlineto{\pgfqpoint{1.009666in}{0.701201in}}%
\pgfpathlineto{\pgfqpoint{1.024560in}{0.699029in}}%
\pgfpathlineto{\pgfqpoint{1.039886in}{0.697375in}}%
\pgfpathlineto{\pgfqpoint{1.060817in}{0.695995in}}%
\pgfpathlineto{\pgfqpoint{1.082086in}{0.695572in}}%
\pgfpathlineto{\pgfqpoint{1.103443in}{0.696114in}}%
\pgfpathlineto{\pgfqpoint{1.124633in}{0.697611in}}%
\pgfpathlineto{\pgfqpoint{1.145406in}{0.700047in}}%
\pgfpathlineto{\pgfqpoint{1.160563in}{0.702471in}}%
\pgfpathlineto{\pgfqpoint{1.175248in}{0.705388in}}%
\pgfpathlineto{\pgfqpoint{1.189364in}{0.708778in}}%
\pgfpathlineto{\pgfqpoint{1.202820in}{0.712616in}}%
\pgfpathlineto{\pgfqpoint{1.215529in}{0.716877in}}%
\pgfpathlineto{\pgfqpoint{1.227410in}{0.721530in}}%
\pgfpathlineto{\pgfqpoint{1.238388in}{0.726544in}}%
\pgfpathlineto{\pgfqpoint{1.248397in}{0.731884in}}%
\pgfpathlineto{\pgfqpoint{1.257376in}{0.737514in}}%
\pgfpathlineto{\pgfqpoint{1.257376in}{0.737514in}}%
\pgfusepath{stroke}%
\end{pgfscope}%
\begin{pgfscope}%
\pgfpathrectangle{\pgfqpoint{0.407592in}{0.222904in}}{\pgfqpoint{1.411163in}{1.411163in}}%
\pgfusepath{clip}%
\pgfsetrectcap%
\pgfsetroundjoin%
\pgfsetlinewidth{0.451687pt}%
\definecolor{currentstroke}{rgb}{0.560784,0.588235,0.611765}%
\pgfsetstrokecolor{currentstroke}%
\pgfsetstrokeopacity{0.450000}%
\pgfsetdash{}{0pt}%
\pgfpathmoveto{\pgfqpoint{1.432199in}{0.870327in}}%
\pgfpathlineto{\pgfqpoint{1.442805in}{0.877983in}}%
\pgfpathlineto{\pgfqpoint{1.452363in}{0.885817in}}%
\pgfpathlineto{\pgfqpoint{1.460854in}{0.893806in}}%
\pgfpathlineto{\pgfqpoint{1.468260in}{0.901926in}}%
\pgfpathlineto{\pgfqpoint{1.474568in}{0.910151in}}%
\pgfpathlineto{\pgfqpoint{1.479767in}{0.918458in}}%
\pgfpathlineto{\pgfqpoint{1.483853in}{0.926822in}}%
\pgfpathlineto{\pgfqpoint{1.486821in}{0.935219in}}%
\pgfpathlineto{\pgfqpoint{1.488671in}{0.943625in}}%
\pgfpathlineto{\pgfqpoint{1.489408in}{0.952017in}}%
\pgfpathlineto{\pgfqpoint{1.489038in}{0.960371in}}%
\pgfpathlineto{\pgfqpoint{1.487569in}{0.968665in}}%
\pgfpathlineto{\pgfqpoint{1.485016in}{0.976877in}}%
\pgfpathlineto{\pgfqpoint{1.481393in}{0.984984in}}%
\pgfpathlineto{\pgfqpoint{1.476717in}{0.992967in}}%
\pgfpathlineto{\pgfqpoint{1.471010in}{1.000805in}}%
\pgfpathlineto{\pgfqpoint{1.464293in}{1.008477in}}%
\pgfpathlineto{\pgfqpoint{1.456592in}{1.015966in}}%
\pgfpathlineto{\pgfqpoint{1.447934in}{1.023252in}}%
\pgfpathlineto{\pgfqpoint{1.438347in}{1.030319in}}%
\pgfpathlineto{\pgfqpoint{1.427863in}{1.037150in}}%
\pgfpathlineto{\pgfqpoint{1.416514in}{1.043729in}}%
\pgfpathlineto{\pgfqpoint{1.404334in}{1.050040in}}%
\pgfpathlineto{\pgfqpoint{1.391359in}{1.056070in}}%
\pgfpathlineto{\pgfqpoint{1.377625in}{1.061805in}}%
\pgfpathlineto{\pgfqpoint{1.363173in}{1.067232in}}%
\pgfpathlineto{\pgfqpoint{1.348040in}{1.072339in}}%
\pgfpathlineto{\pgfqpoint{1.332267in}{1.077116in}}%
\pgfpathlineto{\pgfqpoint{1.307500in}{1.083640in}}%
\pgfpathlineto{\pgfqpoint{1.281531in}{1.089367in}}%
\pgfpathlineto{\pgfqpoint{1.254509in}{1.094270in}}%
\pgfpathlineto{\pgfqpoint{1.226584in}{1.098326in}}%
\pgfpathlineto{\pgfqpoint{1.197911in}{1.101517in}}%
\pgfpathlineto{\pgfqpoint{1.168645in}{1.103829in}}%
\pgfpathlineto{\pgfqpoint{1.138946in}{1.105250in}}%
\pgfpathlineto{\pgfqpoint{1.108971in}{1.105775in}}%
\pgfpathlineto{\pgfqpoint{1.078883in}{1.105402in}}%
\pgfpathlineto{\pgfqpoint{1.048840in}{1.104131in}}%
\pgfpathlineto{\pgfqpoint{1.019005in}{1.101969in}}%
\pgfpathlineto{\pgfqpoint{0.989538in}{1.098925in}}%
\pgfpathlineto{\pgfqpoint{0.960598in}{1.095014in}}%
\pgfpathlineto{\pgfqpoint{0.932343in}{1.090252in}}%
\pgfpathlineto{\pgfqpoint{0.904932in}{1.084663in}}%
\pgfpathlineto{\pgfqpoint{0.878516in}{1.078271in}}%
\pgfpathlineto{\pgfqpoint{0.853249in}{1.071108in}}%
\pgfpathlineto{\pgfqpoint{0.829279in}{1.063207in}}%
\pgfpathlineto{\pgfqpoint{0.814089in}{1.057548in}}%
\pgfpathlineto{\pgfqpoint{0.799582in}{1.051592in}}%
\pgfpathlineto{\pgfqpoint{0.785796in}{1.045351in}}%
\pgfpathlineto{\pgfqpoint{0.772771in}{1.038838in}}%
\pgfpathlineto{\pgfqpoint{0.760545in}{1.032070in}}%
\pgfpathlineto{\pgfqpoint{0.749155in}{1.025061in}}%
\pgfpathlineto{\pgfqpoint{0.738635in}{1.017828in}}%
\pgfpathlineto{\pgfqpoint{0.729018in}{1.010389in}}%
\pgfpathlineto{\pgfqpoint{0.720338in}{1.002761in}}%
\pgfpathlineto{\pgfqpoint{0.712622in}{0.994964in}}%
\pgfpathlineto{\pgfqpoint{0.705899in}{0.987016in}}%
\pgfpathlineto{\pgfqpoint{0.700194in}{0.978938in}}%
\pgfpathlineto{\pgfqpoint{0.695530in}{0.970750in}}%
\pgfpathlineto{\pgfqpoint{0.691928in}{0.962475in}}%
\pgfpathlineto{\pgfqpoint{0.689406in}{0.954134in}}%
\pgfpathlineto{\pgfqpoint{0.687979in}{0.945750in}}%
\pgfpathlineto{\pgfqpoint{0.687661in}{0.937345in}}%
\pgfpathlineto{\pgfqpoint{0.688459in}{0.928943in}}%
\pgfpathlineto{\pgfqpoint{0.690382in}{0.920569in}}%
\pgfpathlineto{\pgfqpoint{0.693431in}{0.912245in}}%
\pgfpathlineto{\pgfqpoint{0.697608in}{0.903997in}}%
\pgfpathlineto{\pgfqpoint{0.702907in}{0.895848in}}%
\pgfpathlineto{\pgfqpoint{0.709324in}{0.887824in}}%
\pgfpathlineto{\pgfqpoint{0.716845in}{0.879948in}}%
\pgfpathlineto{\pgfqpoint{0.725458in}{0.872245in}}%
\pgfpathlineto{\pgfqpoint{0.735145in}{0.864739in}}%
\pgfpathlineto{\pgfqpoint{0.745883in}{0.857455in}}%
\pgfpathlineto{\pgfqpoint{0.757648in}{0.850414in}}%
\pgfpathlineto{\pgfqpoint{0.770409in}{0.843641in}}%
\pgfpathlineto{\pgfqpoint{0.784135in}{0.837157in}}%
\pgfpathlineto{\pgfqpoint{0.798789in}{0.830985in}}%
\pgfpathlineto{\pgfqpoint{0.814331in}{0.825144in}}%
\pgfpathlineto{\pgfqpoint{0.830717in}{0.819655in}}%
\pgfpathlineto{\pgfqpoint{0.847901in}{0.814536in}}%
\pgfpathlineto{\pgfqpoint{0.865832in}{0.809805in}}%
\pgfpathlineto{\pgfqpoint{0.884458in}{0.805478in}}%
\pgfpathlineto{\pgfqpoint{0.903723in}{0.801572in}}%
\pgfpathlineto{\pgfqpoint{0.923567in}{0.798099in}}%
\pgfpathlineto{\pgfqpoint{0.943931in}{0.795071in}}%
\pgfpathlineto{\pgfqpoint{0.964751in}{0.792500in}}%
\pgfpathlineto{\pgfqpoint{0.985961in}{0.790395in}}%
\pgfpathlineto{\pgfqpoint{1.007495in}{0.788764in}}%
\pgfpathlineto{\pgfqpoint{1.040254in}{0.787216in}}%
\pgfpathlineto{\pgfqpoint{1.073354in}{0.786759in}}%
\pgfpathlineto{\pgfqpoint{1.106558in}{0.787396in}}%
\pgfpathlineto{\pgfqpoint{1.139626in}{0.789122in}}%
\pgfpathlineto{\pgfqpoint{1.172322in}{0.791924in}}%
\pgfpathlineto{\pgfqpoint{1.204413in}{0.795777in}}%
\pgfpathlineto{\pgfqpoint{1.225359in}{0.798916in}}%
\pgfpathlineto{\pgfqpoint{1.245870in}{0.802497in}}%
\pgfpathlineto{\pgfqpoint{1.265883in}{0.806509in}}%
\pgfpathlineto{\pgfqpoint{1.285337in}{0.810937in}}%
\pgfpathlineto{\pgfqpoint{1.304173in}{0.815765in}}%
\pgfpathlineto{\pgfqpoint{1.322334in}{0.820977in}}%
\pgfpathlineto{\pgfqpoint{1.339768in}{0.826555in}}%
\pgfpathlineto{\pgfqpoint{1.356425in}{0.832480in}}%
\pgfpathlineto{\pgfqpoint{1.372257in}{0.838731in}}%
\pgfpathlineto{\pgfqpoint{1.387220in}{0.845289in}}%
\pgfpathlineto{\pgfqpoint{1.401274in}{0.852130in}}%
\pgfpathlineto{\pgfqpoint{1.414383in}{0.859233in}}%
\pgfpathlineto{\pgfqpoint{1.426512in}{0.866574in}}%
\pgfpathlineto{\pgfqpoint{1.432199in}{0.870327in}}%
\pgfpathlineto{\pgfqpoint{1.432199in}{0.870327in}}%
\pgfusepath{stroke}%
\end{pgfscope}%
\begin{pgfscope}%
\pgfpathrectangle{\pgfqpoint{0.407592in}{0.222904in}}{\pgfqpoint{1.411163in}{1.411163in}}%
\pgfusepath{clip}%
\pgfsetrectcap%
\pgfsetroundjoin%
\pgfsetlinewidth{0.451687pt}%
\definecolor{currentstroke}{rgb}{0.560784,0.588235,0.611765}%
\pgfsetstrokecolor{currentstroke}%
\pgfsetstrokeopacity{0.450000}%
\pgfsetdash{}{0pt}%
\pgfpathmoveto{\pgfqpoint{1.615239in}{1.009383in}}%
\pgfpathlineto{\pgfqpoint{1.630853in}{1.020628in}}%
\pgfpathlineto{\pgfqpoint{1.644808in}{1.032109in}}%
\pgfpathlineto{\pgfqpoint{1.657082in}{1.043786in}}%
\pgfpathlineto{\pgfqpoint{1.667655in}{1.055623in}}%
\pgfpathlineto{\pgfqpoint{1.676515in}{1.067583in}}%
\pgfpathlineto{\pgfqpoint{1.680301in}{1.073598in}}%
\pgfpathlineto{\pgfqpoint{1.683657in}{1.079630in}}%
\pgfpathlineto{\pgfqpoint{1.686582in}{1.085674in}}%
\pgfpathlineto{\pgfqpoint{1.689078in}{1.091727in}}%
\pgfpathlineto{\pgfqpoint{1.691145in}{1.097783in}}%
\pgfpathlineto{\pgfqpoint{1.692784in}{1.103839in}}%
\pgfpathlineto{\pgfqpoint{1.693996in}{1.109890in}}%
\pgfpathlineto{\pgfqpoint{1.694783in}{1.115932in}}%
\pgfpathlineto{\pgfqpoint{1.695148in}{1.121961in}}%
\pgfpathlineto{\pgfqpoint{1.695091in}{1.127972in}}%
\pgfpathlineto{\pgfqpoint{1.694616in}{1.133962in}}%
\pgfpathlineto{\pgfqpoint{1.693726in}{1.139926in}}%
\pgfpathlineto{\pgfqpoint{1.692424in}{1.145861in}}%
\pgfpathlineto{\pgfqpoint{1.690712in}{1.151763in}}%
\pgfpathlineto{\pgfqpoint{1.688595in}{1.157627in}}%
\pgfpathlineto{\pgfqpoint{1.686076in}{1.163451in}}%
\pgfpathlineto{\pgfqpoint{1.683159in}{1.169230in}}%
\pgfpathlineto{\pgfqpoint{1.676148in}{1.180641in}}%
\pgfpathlineto{\pgfqpoint{1.667600in}{1.191832in}}%
\pgfpathlineto{\pgfqpoint{1.657553in}{1.202776in}}%
\pgfpathlineto{\pgfqpoint{1.646050in}{1.213448in}}%
\pgfpathlineto{\pgfqpoint{1.633135in}{1.223823in}}%
\pgfpathlineto{\pgfqpoint{1.618855in}{1.233880in}}%
\pgfpathlineto{\pgfqpoint{1.603260in}{1.243595in}}%
\pgfpathlineto{\pgfqpoint{1.586402in}{1.252949in}}%
\pgfpathlineto{\pgfqpoint{1.568332in}{1.261922in}}%
\pgfpathlineto{\pgfqpoint{1.549106in}{1.270496in}}%
\pgfpathlineto{\pgfqpoint{1.528781in}{1.278654in}}%
\pgfpathlineto{\pgfqpoint{1.507412in}{1.286380in}}%
\pgfpathlineto{\pgfqpoint{1.485058in}{1.293659in}}%
\pgfpathlineto{\pgfqpoint{1.461779in}{1.300479in}}%
\pgfpathlineto{\pgfqpoint{1.437636in}{1.306826in}}%
\pgfpathlineto{\pgfqpoint{1.412688in}{1.312689in}}%
\pgfpathlineto{\pgfqpoint{1.386997in}{1.318057in}}%
\pgfpathlineto{\pgfqpoint{1.360627in}{1.322922in}}%
\pgfpathlineto{\pgfqpoint{1.333638in}{1.327275in}}%
\pgfpathlineto{\pgfqpoint{1.306095in}{1.331109in}}%
\pgfpathlineto{\pgfqpoint{1.278060in}{1.334416in}}%
\pgfpathlineto{\pgfqpoint{1.249598in}{1.337193in}}%
\pgfpathlineto{\pgfqpoint{1.220773in}{1.339433in}}%
\pgfpathlineto{\pgfqpoint{1.191648in}{1.341134in}}%
\pgfpathlineto{\pgfqpoint{1.162288in}{1.342292in}}%
\pgfpathlineto{\pgfqpoint{1.132758in}{1.342906in}}%
\pgfpathlineto{\pgfqpoint{1.103122in}{1.342975in}}%
\pgfpathlineto{\pgfqpoint{1.073446in}{1.342499in}}%
\pgfpathlineto{\pgfqpoint{1.043795in}{1.341478in}}%
\pgfpathlineto{\pgfqpoint{1.014232in}{1.339914in}}%
\pgfpathlineto{\pgfqpoint{0.984825in}{1.337810in}}%
\pgfpathlineto{\pgfqpoint{0.955636in}{1.335169in}}%
\pgfpathlineto{\pgfqpoint{0.926733in}{1.331995in}}%
\pgfpathlineto{\pgfqpoint{0.898179in}{1.328294in}}%
\pgfpathlineto{\pgfqpoint{0.870040in}{1.324072in}}%
\pgfpathlineto{\pgfqpoint{0.842381in}{1.319336in}}%
\pgfpathlineto{\pgfqpoint{0.815266in}{1.314094in}}%
\pgfpathlineto{\pgfqpoint{0.788760in}{1.308355in}}%
\pgfpathlineto{\pgfqpoint{0.762926in}{1.302129in}}%
\pgfpathlineto{\pgfqpoint{0.737829in}{1.295428in}}%
\pgfpathlineto{\pgfqpoint{0.713532in}{1.288264in}}%
\pgfpathlineto{\pgfqpoint{0.690098in}{1.280649in}}%
\pgfpathlineto{\pgfqpoint{0.667589in}{1.272599in}}%
\pgfpathlineto{\pgfqpoint{0.646065in}{1.264129in}}%
\pgfpathlineto{\pgfqpoint{0.625588in}{1.255255in}}%
\pgfpathlineto{\pgfqpoint{0.606216in}{1.245996in}}%
\pgfpathlineto{\pgfqpoint{0.588007in}{1.236370in}}%
\pgfpathlineto{\pgfqpoint{0.571019in}{1.226398in}}%
\pgfpathlineto{\pgfqpoint{0.555305in}{1.216101in}}%
\pgfpathlineto{\pgfqpoint{0.540920in}{1.205501in}}%
\pgfpathlineto{\pgfqpoint{0.527913in}{1.194624in}}%
\pgfpathlineto{\pgfqpoint{0.516335in}{1.183493in}}%
\pgfpathlineto{\pgfqpoint{0.506231in}{1.172135in}}%
\pgfpathlineto{\pgfqpoint{0.497646in}{1.160578in}}%
\pgfpathlineto{\pgfqpoint{0.490618in}{1.148850in}}%
\pgfpathlineto{\pgfqpoint{0.487701in}{1.142932in}}%
\pgfpathlineto{\pgfqpoint{0.485187in}{1.136982in}}%
\pgfpathlineto{\pgfqpoint{0.483080in}{1.131005in}}%
\pgfpathlineto{\pgfqpoint{0.481385in}{1.125004in}}%
\pgfpathlineto{\pgfqpoint{0.480104in}{1.118983in}}%
\pgfpathlineto{\pgfqpoint{0.479242in}{1.112948in}}%
\pgfpathlineto{\pgfqpoint{0.478801in}{1.106901in}}%
\pgfpathlineto{\pgfqpoint{0.478784in}{1.100847in}}%
\pgfpathlineto{\pgfqpoint{0.479194in}{1.094790in}}%
\pgfpathlineto{\pgfqpoint{0.480033in}{1.088735in}}%
\pgfpathlineto{\pgfqpoint{0.481302in}{1.082686in}}%
\pgfpathlineto{\pgfqpoint{0.483004in}{1.076647in}}%
\pgfpathlineto{\pgfqpoint{0.485140in}{1.070623in}}%
\pgfpathlineto{\pgfqpoint{0.487710in}{1.064619in}}%
\pgfpathlineto{\pgfqpoint{0.490715in}{1.058638in}}%
\pgfpathlineto{\pgfqpoint{0.494156in}{1.052685in}}%
\pgfpathlineto{\pgfqpoint{0.498033in}{1.046766in}}%
\pgfpathlineto{\pgfqpoint{0.502344in}{1.040884in}}%
\pgfpathlineto{\pgfqpoint{0.512269in}{1.029252in}}%
\pgfpathlineto{\pgfqpoint{0.523920in}{1.017826in}}%
\pgfpathlineto{\pgfqpoint{0.537279in}{1.006644in}}%
\pgfpathlineto{\pgfqpoint{0.552324in}{0.995743in}}%
\pgfpathlineto{\pgfqpoint{0.569025in}{0.985160in}}%
\pgfpathlineto{\pgfqpoint{0.587346in}{0.974931in}}%
\pgfpathlineto{\pgfqpoint{0.607243in}{0.965093in}}%
\pgfpathlineto{\pgfqpoint{0.628667in}{0.955681in}}%
\pgfpathlineto{\pgfqpoint{0.651562in}{0.946729in}}%
\pgfpathlineto{\pgfqpoint{0.675867in}{0.938270in}}%
\pgfpathlineto{\pgfqpoint{0.701511in}{0.930336in}}%
\pgfpathlineto{\pgfqpoint{0.728419in}{0.922957in}}%
\pgfpathlineto{\pgfqpoint{0.756512in}{0.916161in}}%
\pgfpathlineto{\pgfqpoint{0.785702in}{0.909975in}}%
\pgfpathlineto{\pgfqpoint{0.815897in}{0.904423in}}%
\pgfpathlineto{\pgfqpoint{0.847001in}{0.899526in}}%
\pgfpathlineto{\pgfqpoint{0.878911in}{0.895306in}}%
\pgfpathlineto{\pgfqpoint{0.911523in}{0.891777in}}%
\pgfpathlineto{\pgfqpoint{0.944727in}{0.888956in}}%
\pgfpathlineto{\pgfqpoint{0.978411in}{0.886852in}}%
\pgfpathlineto{\pgfqpoint{1.012462in}{0.885476in}}%
\pgfpathlineto{\pgfqpoint{1.046761in}{0.884831in}}%
\pgfpathlineto{\pgfqpoint{1.081192in}{0.884922in}}%
\pgfpathlineto{\pgfqpoint{1.115637in}{0.885748in}}%
\pgfpathlineto{\pgfqpoint{1.149978in}{0.887305in}}%
\pgfpathlineto{\pgfqpoint{1.184097in}{0.889586in}}%
\pgfpathlineto{\pgfqpoint{1.217878in}{0.892584in}}%
\pgfpathlineto{\pgfqpoint{1.251206in}{0.896285in}}%
\pgfpathlineto{\pgfqpoint{1.283971in}{0.900674in}}%
\pgfpathlineto{\pgfqpoint{1.316063in}{0.905734in}}%
\pgfpathlineto{\pgfqpoint{1.347377in}{0.911445in}}%
\pgfpathlineto{\pgfqpoint{1.377811in}{0.917785in}}%
\pgfpathlineto{\pgfqpoint{1.407270in}{0.924727in}}%
\pgfpathlineto{\pgfqpoint{1.435662in}{0.932246in}}%
\pgfpathlineto{\pgfqpoint{1.462900in}{0.940313in}}%
\pgfpathlineto{\pgfqpoint{1.488904in}{0.948897in}}%
\pgfpathlineto{\pgfqpoint{1.513598in}{0.957965in}}%
\pgfpathlineto{\pgfqpoint{1.536914in}{0.967486in}}%
\pgfpathlineto{\pgfqpoint{1.558789in}{0.977424in}}%
\pgfpathlineto{\pgfqpoint{1.579167in}{0.987743in}}%
\pgfpathlineto{\pgfqpoint{1.597998in}{0.998409in}}%
\pgfpathlineto{\pgfqpoint{1.615239in}{1.009383in}}%
\pgfpathlineto{\pgfqpoint{1.615239in}{1.009383in}}%
\pgfusepath{stroke}%
\end{pgfscope}%
\begin{pgfscope}%
\pgfpathrectangle{\pgfqpoint{0.407592in}{0.222904in}}{\pgfqpoint{1.411163in}{1.411163in}}%
\pgfusepath{clip}%
\pgfsetrectcap%
\pgfsetroundjoin%
\pgfsetlinewidth{0.451687pt}%
\definecolor{currentstroke}{rgb}{0.560784,0.588235,0.611765}%
\pgfsetstrokecolor{currentstroke}%
\pgfsetstrokeopacity{0.450000}%
\pgfsetdash{}{0pt}%
\pgfpathmoveto{\pgfqpoint{1.763662in}{1.122140in}}%
\pgfpathlineto{\pgfqpoint{1.782996in}{1.136039in}}%
\pgfpathlineto{\pgfqpoint{1.791849in}{1.143089in}}%
\pgfpathlineto{\pgfqpoint{1.800155in}{1.150199in}}%
\pgfpathlineto{\pgfqpoint{1.807911in}{1.157362in}}%
\pgfpathlineto{\pgfqpoint{1.815114in}{1.164572in}}%
\pgfpathlineto{\pgfqpoint{1.822088in}{1.172211in}}%
\pgfpathmoveto{\pgfqpoint{1.822088in}{1.340804in}}%
\pgfpathlineto{\pgfqpoint{1.818991in}{1.344252in}}%
\pgfpathlineto{\pgfqpoint{1.812535in}{1.350843in}}%
\pgfpathlineto{\pgfqpoint{1.805624in}{1.357350in}}%
\pgfpathlineto{\pgfqpoint{1.790464in}{1.370098in}}%
\pgfpathlineto{\pgfqpoint{1.773573in}{1.382469in}}%
\pgfpathlineto{\pgfqpoint{1.755016in}{1.394439in}}%
\pgfpathlineto{\pgfqpoint{1.734858in}{1.405982in}}%
\pgfpathlineto{\pgfqpoint{1.713167in}{1.417078in}}%
\pgfpathlineto{\pgfqpoint{1.690012in}{1.427704in}}%
\pgfpathlineto{\pgfqpoint{1.665464in}{1.437843in}}%
\pgfpathlineto{\pgfqpoint{1.639594in}{1.447475in}}%
\pgfpathlineto{\pgfqpoint{1.612475in}{1.456585in}}%
\pgfpathlineto{\pgfqpoint{1.584179in}{1.465157in}}%
\pgfpathlineto{\pgfqpoint{1.554782in}{1.473178in}}%
\pgfpathlineto{\pgfqpoint{1.524355in}{1.480635in}}%
\pgfpathlineto{\pgfqpoint{1.492975in}{1.487515in}}%
\pgfpathlineto{\pgfqpoint{1.460716in}{1.493810in}}%
\pgfpathlineto{\pgfqpoint{1.427652in}{1.499508in}}%
\pgfpathlineto{\pgfqpoint{1.393859in}{1.504603in}}%
\pgfpathlineto{\pgfqpoint{1.359412in}{1.509087in}}%
\pgfpathlineto{\pgfqpoint{1.324386in}{1.512953in}}%
\pgfpathlineto{\pgfqpoint{1.288857in}{1.516197in}}%
\pgfpathlineto{\pgfqpoint{1.252899in}{1.518813in}}%
\pgfpathlineto{\pgfqpoint{1.216588in}{1.520799in}}%
\pgfpathlineto{\pgfqpoint{1.179999in}{1.522151in}}%
\pgfpathlineto{\pgfqpoint{1.143208in}{1.522867in}}%
\pgfpathlineto{\pgfqpoint{1.106290in}{1.522948in}}%
\pgfpathlineto{\pgfqpoint{1.069322in}{1.522392in}}%
\pgfpathlineto{\pgfqpoint{1.032377in}{1.521200in}}%
\pgfpathlineto{\pgfqpoint{0.995534in}{1.519375in}}%
\pgfpathlineto{\pgfqpoint{0.958866in}{1.516918in}}%
\pgfpathlineto{\pgfqpoint{0.922452in}{1.513833in}}%
\pgfpathlineto{\pgfqpoint{0.886366in}{1.510123in}}%
\pgfpathlineto{\pgfqpoint{0.850686in}{1.505795in}}%
\pgfpathlineto{\pgfqpoint{0.815487in}{1.500854in}}%
\pgfpathlineto{\pgfqpoint{0.780849in}{1.495307in}}%
\pgfpathlineto{\pgfqpoint{0.746846in}{1.489163in}}%
\pgfpathlineto{\pgfqpoint{0.713557in}{1.482430in}}%
\pgfpathlineto{\pgfqpoint{0.681058in}{1.475118in}}%
\pgfpathlineto{\pgfqpoint{0.649427in}{1.467239in}}%
\pgfpathlineto{\pgfqpoint{0.618742in}{1.458804in}}%
\pgfpathlineto{\pgfqpoint{0.589080in}{1.449829in}}%
\pgfpathlineto{\pgfqpoint{0.560517in}{1.440327in}}%
\pgfpathlineto{\pgfqpoint{0.533130in}{1.430315in}}%
\pgfpathlineto{\pgfqpoint{0.506995in}{1.419810in}}%
\pgfpathlineto{\pgfqpoint{0.482188in}{1.408831in}}%
\pgfpathlineto{\pgfqpoint{0.458784in}{1.397399in}}%
\pgfpathlineto{\pgfqpoint{0.436856in}{1.385535in}}%
\pgfpathlineto{\pgfqpoint{0.416476in}{1.373263in}}%
\pgfpathlineto{\pgfqpoint{0.404258in}{1.365154in}}%
\pgfpathmoveto{\pgfqpoint{0.404258in}{1.102084in}}%
\pgfpathlineto{\pgfqpoint{0.409555in}{1.098627in}}%
\pgfpathlineto{\pgfqpoint{0.431448in}{1.085692in}}%
\pgfpathlineto{\pgfqpoint{0.455462in}{1.073197in}}%
\pgfpathlineto{\pgfqpoint{0.481542in}{1.061190in}}%
\pgfpathlineto{\pgfqpoint{0.509624in}{1.049715in}}%
\pgfpathlineto{\pgfqpoint{0.539635in}{1.038819in}}%
\pgfpathlineto{\pgfqpoint{0.571492in}{1.028544in}}%
\pgfpathlineto{\pgfqpoint{0.605103in}{1.018932in}}%
\pgfpathlineto{\pgfqpoint{0.640368in}{1.010023in}}%
\pgfpathlineto{\pgfqpoint{0.677177in}{1.001854in}}%
\pgfpathlineto{\pgfqpoint{0.715411in}{0.994459in}}%
\pgfpathlineto{\pgfqpoint{0.754945in}{0.987870in}}%
\pgfpathlineto{\pgfqpoint{0.795645in}{0.982116in}}%
\pgfpathlineto{\pgfqpoint{0.837373in}{0.977221in}}%
\pgfpathlineto{\pgfqpoint{0.879982in}{0.973208in}}%
\pgfpathlineto{\pgfqpoint{0.923322in}{0.970093in}}%
\pgfpathlineto{\pgfqpoint{0.967239in}{0.967892in}}%
\pgfpathlineto{\pgfqpoint{1.011573in}{0.966613in}}%
\pgfpathlineto{\pgfqpoint{1.056165in}{0.966263in}}%
\pgfpathlineto{\pgfqpoint{1.100854in}{0.966843in}}%
\pgfpathlineto{\pgfqpoint{1.145476in}{0.968350in}}%
\pgfpathlineto{\pgfqpoint{1.189870in}{0.970778in}}%
\pgfpathlineto{\pgfqpoint{1.233876in}{0.974116in}}%
\pgfpathlineto{\pgfqpoint{1.277338in}{0.978349in}}%
\pgfpathlineto{\pgfqpoint{1.320099in}{0.983459in}}%
\pgfpathlineto{\pgfqpoint{1.362012in}{0.989422in}}%
\pgfpathlineto{\pgfqpoint{1.402931in}{0.996213in}}%
\pgfpathlineto{\pgfqpoint{1.442717in}{1.003802in}}%
\pgfpathlineto{\pgfqpoint{1.481240in}{1.012157in}}%
\pgfpathlineto{\pgfqpoint{1.518373in}{1.021244in}}%
\pgfpathlineto{\pgfqpoint{1.553999in}{1.031023in}}%
\pgfpathlineto{\pgfqpoint{1.588011in}{1.041455in}}%
\pgfpathlineto{\pgfqpoint{1.620307in}{1.052498in}}%
\pgfpathlineto{\pgfqpoint{1.650797in}{1.064109in}}%
\pgfpathlineto{\pgfqpoint{1.679399in}{1.076241in}}%
\pgfpathlineto{\pgfqpoint{1.706039in}{1.088849in}}%
\pgfpathlineto{\pgfqpoint{1.730654in}{1.101886in}}%
\pgfpathlineto{\pgfqpoint{1.753189in}{1.115303in}}%
\pgfpathlineto{\pgfqpoint{1.763662in}{1.122140in}}%
\pgfpathlineto{\pgfqpoint{1.763662in}{1.122140in}}%
\pgfusepath{stroke}%
\end{pgfscope}%
\begin{pgfscope}%
\pgfpathrectangle{\pgfqpoint{0.407592in}{0.222904in}}{\pgfqpoint{1.411163in}{1.411163in}}%
\pgfusepath{clip}%
\pgfsetrectcap%
\pgfsetroundjoin%
\pgfsetlinewidth{0.401500pt}%
\definecolor{currentstroke}{rgb}{0.560784,0.588235,0.611765}%
\pgfsetstrokecolor{currentstroke}%
\pgfsetstrokeopacity{0.350000}%
\pgfsetdash{}{0pt}%
\pgfpathmoveto{\pgfqpoint{1.090230in}{0.610533in}}%
\pgfpathlineto{\pgfqpoint{1.098147in}{0.616548in}}%
\pgfpathlineto{\pgfqpoint{1.106082in}{0.622576in}}%
\pgfpathlineto{\pgfqpoint{1.114034in}{0.628617in}}%
\pgfpathlineto{\pgfqpoint{1.122002in}{0.634671in}}%
\pgfpathlineto{\pgfqpoint{1.129988in}{0.640737in}}%
\pgfpathlineto{\pgfqpoint{1.137991in}{0.646817in}}%
\pgfpathlineto{\pgfqpoint{1.146011in}{0.652910in}}%
\pgfpathlineto{\pgfqpoint{1.154048in}{0.659015in}}%
\pgfpathlineto{\pgfqpoint{1.162102in}{0.665134in}}%
\pgfpathlineto{\pgfqpoint{1.170173in}{0.671266in}}%
\pgfpathlineto{\pgfqpoint{1.178262in}{0.677411in}}%
\pgfpathlineto{\pgfqpoint{1.186368in}{0.683570in}}%
\pgfpathlineto{\pgfqpoint{1.194492in}{0.689741in}}%
\pgfpathlineto{\pgfqpoint{1.202633in}{0.695926in}}%
\pgfpathlineto{\pgfqpoint{1.210792in}{0.702124in}}%
\pgfpathlineto{\pgfqpoint{1.218969in}{0.708336in}}%
\pgfpathlineto{\pgfqpoint{1.227163in}{0.714561in}}%
\pgfpathlineto{\pgfqpoint{1.235375in}{0.720800in}}%
\pgfpathlineto{\pgfqpoint{1.243605in}{0.727052in}}%
\pgfpathlineto{\pgfqpoint{1.251852in}{0.733318in}}%
\pgfpathlineto{\pgfqpoint{1.260118in}{0.739597in}}%
\pgfpathlineto{\pgfqpoint{1.268401in}{0.745890in}}%
\pgfpathlineto{\pgfqpoint{1.276703in}{0.752196in}}%
\pgfpathlineto{\pgfqpoint{1.285022in}{0.758517in}}%
\pgfpathlineto{\pgfqpoint{1.293360in}{0.764851in}}%
\pgfpathlineto{\pgfqpoint{1.301716in}{0.771199in}}%
\pgfpathlineto{\pgfqpoint{1.310090in}{0.777561in}}%
\pgfpathlineto{\pgfqpoint{1.318483in}{0.783937in}}%
\pgfpathlineto{\pgfqpoint{1.326894in}{0.790327in}}%
\pgfpathlineto{\pgfqpoint{1.335323in}{0.796731in}}%
\pgfpathlineto{\pgfqpoint{1.343771in}{0.803148in}}%
\pgfpathlineto{\pgfqpoint{1.352238in}{0.809580in}}%
\pgfpathlineto{\pgfqpoint{1.360723in}{0.816027in}}%
\pgfpathlineto{\pgfqpoint{1.369227in}{0.822487in}}%
\pgfpathlineto{\pgfqpoint{1.377749in}{0.828962in}}%
\pgfpathlineto{\pgfqpoint{1.386291in}{0.835451in}}%
\pgfpathlineto{\pgfqpoint{1.394851in}{0.841954in}}%
\pgfpathlineto{\pgfqpoint{1.403430in}{0.848472in}}%
\pgfpathlineto{\pgfqpoint{1.412029in}{0.855004in}}%
\pgfpathlineto{\pgfqpoint{1.420646in}{0.861550in}}%
\pgfpathlineto{\pgfqpoint{1.429283in}{0.868112in}}%
\pgfpathlineto{\pgfqpoint{1.437938in}{0.874687in}}%
\pgfpathlineto{\pgfqpoint{1.446613in}{0.881278in}}%
\pgfpathlineto{\pgfqpoint{1.455308in}{0.887883in}}%
\pgfpathlineto{\pgfqpoint{1.464021in}{0.894503in}}%
\pgfpathlineto{\pgfqpoint{1.472755in}{0.901137in}}%
\pgfpathlineto{\pgfqpoint{1.481507in}{0.907787in}}%
\pgfpathlineto{\pgfqpoint{1.490280in}{0.914451in}}%
\pgfpathlineto{\pgfqpoint{1.499072in}{0.921130in}}%
\pgfpathlineto{\pgfqpoint{1.507883in}{0.927824in}}%
\pgfpathlineto{\pgfqpoint{1.516715in}{0.934534in}}%
\pgfpathlineto{\pgfqpoint{1.525566in}{0.941258in}}%
\pgfpathlineto{\pgfqpoint{1.534437in}{0.947998in}}%
\pgfpathlineto{\pgfqpoint{1.543329in}{0.954752in}}%
\pgfpathlineto{\pgfqpoint{1.552240in}{0.961522in}}%
\pgfpathlineto{\pgfqpoint{1.561171in}{0.968307in}}%
\pgfpathlineto{\pgfqpoint{1.570123in}{0.975108in}}%
\pgfpathlineto{\pgfqpoint{1.579095in}{0.981924in}}%
\pgfpathlineto{\pgfqpoint{1.588087in}{0.988755in}}%
\pgfpathlineto{\pgfqpoint{1.597100in}{0.995602in}}%
\pgfpathlineto{\pgfqpoint{1.606133in}{1.002465in}}%
\pgfpathlineto{\pgfqpoint{1.615186in}{1.009343in}}%
\pgfpathlineto{\pgfqpoint{1.624261in}{1.016236in}}%
\pgfpathlineto{\pgfqpoint{1.633356in}{1.023146in}}%
\pgfpathlineto{\pgfqpoint{1.642471in}{1.030071in}}%
\pgfpathlineto{\pgfqpoint{1.651608in}{1.037012in}}%
\pgfpathlineto{\pgfqpoint{1.660765in}{1.043969in}}%
\pgfpathlineto{\pgfqpoint{1.669944in}{1.050942in}}%
\pgfpathlineto{\pgfqpoint{1.679143in}{1.057931in}}%
\pgfpathlineto{\pgfqpoint{1.688364in}{1.064935in}}%
\pgfpathlineto{\pgfqpoint{1.697605in}{1.071956in}}%
\pgfpathlineto{\pgfqpoint{1.706868in}{1.078993in}}%
\pgfpathlineto{\pgfqpoint{1.716153in}{1.086047in}}%
\pgfpathlineto{\pgfqpoint{1.725458in}{1.093116in}}%
\pgfpathlineto{\pgfqpoint{1.734786in}{1.100202in}}%
\pgfpathlineto{\pgfqpoint{1.744134in}{1.107304in}}%
\pgfpathlineto{\pgfqpoint{1.753505in}{1.114423in}}%
\pgfpathlineto{\pgfqpoint{1.762897in}{1.121558in}}%
\pgfpathlineto{\pgfqpoint{1.772311in}{1.128710in}}%
\pgfusepath{stroke}%
\end{pgfscope}%
\begin{pgfscope}%
\pgfpathrectangle{\pgfqpoint{0.407592in}{0.222904in}}{\pgfqpoint{1.411163in}{1.411163in}}%
\pgfusepath{clip}%
\pgfsetrectcap%
\pgfsetroundjoin%
\pgfsetlinewidth{0.401500pt}%
\definecolor{currentstroke}{rgb}{0.560784,0.588235,0.611765}%
\pgfsetstrokecolor{currentstroke}%
\pgfsetstrokeopacity{0.350000}%
\pgfsetdash{}{0pt}%
\pgfpathmoveto{\pgfqpoint{1.090230in}{0.610533in}}%
\pgfpathlineto{\pgfqpoint{1.099520in}{0.619706in}}%
\pgfpathlineto{\pgfqpoint{1.108814in}{0.628883in}}%
\pgfpathlineto{\pgfqpoint{1.118111in}{0.638063in}}%
\pgfpathlineto{\pgfqpoint{1.127411in}{0.647246in}}%
\pgfpathlineto{\pgfqpoint{1.136715in}{0.656433in}}%
\pgfpathlineto{\pgfqpoint{1.146022in}{0.665623in}}%
\pgfpathlineto{\pgfqpoint{1.155332in}{0.674816in}}%
\pgfpathlineto{\pgfqpoint{1.164646in}{0.684012in}}%
\pgfpathlineto{\pgfqpoint{1.173963in}{0.693212in}}%
\pgfpathlineto{\pgfqpoint{1.183283in}{0.702414in}}%
\pgfpathlineto{\pgfqpoint{1.192606in}{0.711621in}}%
\pgfpathlineto{\pgfqpoint{1.201933in}{0.720830in}}%
\pgfpathlineto{\pgfqpoint{1.211263in}{0.730043in}}%
\pgfpathlineto{\pgfqpoint{1.220597in}{0.739259in}}%
\pgfpathlineto{\pgfqpoint{1.229934in}{0.748478in}}%
\pgfpathlineto{\pgfqpoint{1.239274in}{0.757700in}}%
\pgfpathlineto{\pgfqpoint{1.248617in}{0.766926in}}%
\pgfpathlineto{\pgfqpoint{1.257964in}{0.776155in}}%
\pgfpathlineto{\pgfqpoint{1.267314in}{0.785388in}}%
\pgfpathlineto{\pgfqpoint{1.276668in}{0.794623in}}%
\pgfpathlineto{\pgfqpoint{1.286024in}{0.803862in}}%
\pgfpathlineto{\pgfqpoint{1.295384in}{0.813104in}}%
\pgfpathlineto{\pgfqpoint{1.304748in}{0.822350in}}%
\pgfpathlineto{\pgfqpoint{1.314115in}{0.831599in}}%
\pgfpathlineto{\pgfqpoint{1.323485in}{0.840851in}}%
\pgfpathlineto{\pgfqpoint{1.332858in}{0.850106in}}%
\pgfpathlineto{\pgfqpoint{1.342235in}{0.859365in}}%
\pgfpathlineto{\pgfqpoint{1.351615in}{0.868627in}}%
\pgfpathlineto{\pgfqpoint{1.360998in}{0.877892in}}%
\pgfpathlineto{\pgfqpoint{1.370385in}{0.887161in}}%
\pgfpathlineto{\pgfqpoint{1.379775in}{0.896433in}}%
\pgfpathlineto{\pgfqpoint{1.389169in}{0.905708in}}%
\pgfpathlineto{\pgfqpoint{1.398566in}{0.914986in}}%
\pgfpathlineto{\pgfqpoint{1.407966in}{0.924268in}}%
\pgfpathlineto{\pgfqpoint{1.417369in}{0.933553in}}%
\pgfpathlineto{\pgfqpoint{1.426776in}{0.942842in}}%
\pgfpathlineto{\pgfqpoint{1.436187in}{0.952133in}}%
\pgfpathlineto{\pgfqpoint{1.445600in}{0.961429in}}%
\pgfpathlineto{\pgfqpoint{1.455017in}{0.970727in}}%
\pgfpathlineto{\pgfqpoint{1.464438in}{0.980029in}}%
\pgfpathlineto{\pgfqpoint{1.473861in}{0.989334in}}%
\pgfpathlineto{\pgfqpoint{1.483288in}{0.998642in}}%
\pgfpathlineto{\pgfqpoint{1.492719in}{1.007954in}}%
\pgfpathlineto{\pgfqpoint{1.502153in}{1.017269in}}%
\pgfpathlineto{\pgfqpoint{1.511590in}{1.026587in}}%
\pgfpathlineto{\pgfqpoint{1.521031in}{1.035909in}}%
\pgfpathlineto{\pgfqpoint{1.530475in}{1.045234in}}%
\pgfpathlineto{\pgfqpoint{1.539922in}{1.054562in}}%
\pgfpathlineto{\pgfqpoint{1.549373in}{1.063894in}}%
\pgfpathlineto{\pgfqpoint{1.558827in}{1.073229in}}%
\pgfpathlineto{\pgfqpoint{1.568284in}{1.082568in}}%
\pgfpathlineto{\pgfqpoint{1.577745in}{1.091909in}}%
\pgfpathlineto{\pgfqpoint{1.587209in}{1.101254in}}%
\pgfpathlineto{\pgfqpoint{1.596677in}{1.110603in}}%
\pgfpathlineto{\pgfqpoint{1.606148in}{1.119955in}}%
\pgfpathlineto{\pgfqpoint{1.615622in}{1.129310in}}%
\pgfpathlineto{\pgfqpoint{1.625100in}{1.138668in}}%
\pgfpathlineto{\pgfqpoint{1.634581in}{1.148030in}}%
\pgfpathlineto{\pgfqpoint{1.644066in}{1.157395in}}%
\pgfpathlineto{\pgfqpoint{1.653554in}{1.166764in}}%
\pgfpathlineto{\pgfqpoint{1.663046in}{1.176136in}}%
\pgfpathlineto{\pgfqpoint{1.672540in}{1.185511in}}%
\pgfpathlineto{\pgfqpoint{1.682039in}{1.194890in}}%
\pgfpathlineto{\pgfqpoint{1.691540in}{1.204272in}}%
\pgfpathlineto{\pgfqpoint{1.701045in}{1.213657in}}%
\pgfpathlineto{\pgfqpoint{1.710554in}{1.223046in}}%
\pgfpathlineto{\pgfqpoint{1.720066in}{1.232438in}}%
\pgfpathlineto{\pgfqpoint{1.729581in}{1.241833in}}%
\pgfpathlineto{\pgfqpoint{1.739100in}{1.251232in}}%
\pgfpathlineto{\pgfqpoint{1.748622in}{1.260634in}}%
\pgfpathlineto{\pgfqpoint{1.758147in}{1.270040in}}%
\pgfpathlineto{\pgfqpoint{1.767676in}{1.279449in}}%
\pgfpathlineto{\pgfqpoint{1.777209in}{1.288861in}}%
\pgfpathlineto{\pgfqpoint{1.786745in}{1.298277in}}%
\pgfpathlineto{\pgfqpoint{1.796284in}{1.307696in}}%
\pgfpathlineto{\pgfqpoint{1.805827in}{1.317119in}}%
\pgfpathlineto{\pgfqpoint{1.815373in}{1.326545in}}%
\pgfpathlineto{\pgfqpoint{1.822088in}{1.333175in}}%
\pgfusepath{stroke}%
\end{pgfscope}%
\begin{pgfscope}%
\pgfpathrectangle{\pgfqpoint{0.407592in}{0.222904in}}{\pgfqpoint{1.411163in}{1.411163in}}%
\pgfusepath{clip}%
\pgfsetrectcap%
\pgfsetroundjoin%
\pgfsetlinewidth{0.401500pt}%
\definecolor{currentstroke}{rgb}{0.560784,0.588235,0.611765}%
\pgfsetstrokecolor{currentstroke}%
\pgfsetstrokeopacity{0.350000}%
\pgfsetdash{}{0pt}%
\pgfpathmoveto{\pgfqpoint{1.090230in}{0.610533in}}%
\pgfpathlineto{\pgfqpoint{1.095442in}{0.622265in}}%
\pgfpathlineto{\pgfqpoint{1.100649in}{0.633984in}}%
\pgfpathlineto{\pgfqpoint{1.105851in}{0.645690in}}%
\pgfpathlineto{\pgfqpoint{1.111047in}{0.657384in}}%
\pgfpathlineto{\pgfqpoint{1.116237in}{0.669065in}}%
\pgfpathlineto{\pgfqpoint{1.121422in}{0.680734in}}%
\pgfpathlineto{\pgfqpoint{1.126601in}{0.692390in}}%
\pgfpathlineto{\pgfqpoint{1.131774in}{0.704033in}}%
\pgfpathlineto{\pgfqpoint{1.136942in}{0.715664in}}%
\pgfpathlineto{\pgfqpoint{1.142105in}{0.727283in}}%
\pgfpathlineto{\pgfqpoint{1.147262in}{0.738889in}}%
\pgfpathlineto{\pgfqpoint{1.152413in}{0.750482in}}%
\pgfpathlineto{\pgfqpoint{1.157559in}{0.762063in}}%
\pgfpathlineto{\pgfqpoint{1.162699in}{0.773632in}}%
\pgfpathlineto{\pgfqpoint{1.167834in}{0.785188in}}%
\pgfpathlineto{\pgfqpoint{1.172963in}{0.796732in}}%
\pgfpathlineto{\pgfqpoint{1.178087in}{0.808263in}}%
\pgfpathlineto{\pgfqpoint{1.183205in}{0.819782in}}%
\pgfpathlineto{\pgfqpoint{1.188318in}{0.831289in}}%
\pgfpathlineto{\pgfqpoint{1.193425in}{0.842784in}}%
\pgfpathlineto{\pgfqpoint{1.198527in}{0.854266in}}%
\pgfpathlineto{\pgfqpoint{1.203623in}{0.865736in}}%
\pgfpathlineto{\pgfqpoint{1.208714in}{0.877194in}}%
\pgfpathlineto{\pgfqpoint{1.213800in}{0.888639in}}%
\pgfpathlineto{\pgfqpoint{1.218880in}{0.900072in}}%
\pgfpathlineto{\pgfqpoint{1.223955in}{0.911493in}}%
\pgfpathlineto{\pgfqpoint{1.229024in}{0.922902in}}%
\pgfpathlineto{\pgfqpoint{1.234088in}{0.934299in}}%
\pgfpathlineto{\pgfqpoint{1.239146in}{0.945684in}}%
\pgfpathlineto{\pgfqpoint{1.244200in}{0.957056in}}%
\pgfpathlineto{\pgfqpoint{1.249247in}{0.968417in}}%
\pgfpathlineto{\pgfqpoint{1.254290in}{0.979765in}}%
\pgfpathlineto{\pgfqpoint{1.259327in}{0.991101in}}%
\pgfpathlineto{\pgfqpoint{1.264358in}{1.002426in}}%
\pgfpathlineto{\pgfqpoint{1.269385in}{1.013738in}}%
\pgfpathlineto{\pgfqpoint{1.274406in}{1.025038in}}%
\pgfpathlineto{\pgfqpoint{1.279422in}{1.036326in}}%
\pgfpathlineto{\pgfqpoint{1.284432in}{1.047603in}}%
\pgfpathlineto{\pgfqpoint{1.289437in}{1.058867in}}%
\pgfpathlineto{\pgfqpoint{1.294437in}{1.070120in}}%
\pgfpathlineto{\pgfqpoint{1.299431in}{1.081360in}}%
\pgfpathlineto{\pgfqpoint{1.304421in}{1.092589in}}%
\pgfpathlineto{\pgfqpoint{1.309405in}{1.103806in}}%
\pgfpathlineto{\pgfqpoint{1.314383in}{1.115011in}}%
\pgfpathlineto{\pgfqpoint{1.319357in}{1.126204in}}%
\pgfpathlineto{\pgfqpoint{1.324325in}{1.137385in}}%
\pgfpathlineto{\pgfqpoint{1.329288in}{1.148555in}}%
\pgfpathlineto{\pgfqpoint{1.334246in}{1.159713in}}%
\pgfpathlineto{\pgfqpoint{1.339198in}{1.170859in}}%
\pgfpathlineto{\pgfqpoint{1.344145in}{1.181993in}}%
\pgfpathlineto{\pgfqpoint{1.349087in}{1.193116in}}%
\pgfpathlineto{\pgfqpoint{1.354024in}{1.204227in}}%
\pgfpathlineto{\pgfqpoint{1.358956in}{1.215326in}}%
\pgfpathlineto{\pgfqpoint{1.363883in}{1.226413in}}%
\pgfpathlineto{\pgfqpoint{1.368804in}{1.237489in}}%
\pgfpathlineto{\pgfqpoint{1.373720in}{1.248554in}}%
\pgfpathlineto{\pgfqpoint{1.378631in}{1.259606in}}%
\pgfpathlineto{\pgfqpoint{1.383537in}{1.270647in}}%
\pgfpathlineto{\pgfqpoint{1.388438in}{1.281677in}}%
\pgfpathlineto{\pgfqpoint{1.393333in}{1.292695in}}%
\pgfpathlineto{\pgfqpoint{1.398224in}{1.303702in}}%
\pgfpathlineto{\pgfqpoint{1.403109in}{1.314697in}}%
\pgfpathlineto{\pgfqpoint{1.407989in}{1.325680in}}%
\pgfpathlineto{\pgfqpoint{1.412865in}{1.336652in}}%
\pgfpathlineto{\pgfqpoint{1.417735in}{1.347613in}}%
\pgfpathlineto{\pgfqpoint{1.422600in}{1.358562in}}%
\pgfpathlineto{\pgfqpoint{1.427460in}{1.369499in}}%
\pgfpathlineto{\pgfqpoint{1.432314in}{1.380426in}}%
\pgfpathlineto{\pgfqpoint{1.437164in}{1.391341in}}%
\pgfpathlineto{\pgfqpoint{1.442009in}{1.402244in}}%
\pgfpathlineto{\pgfqpoint{1.446849in}{1.413136in}}%
\pgfpathlineto{\pgfqpoint{1.451683in}{1.424017in}}%
\pgfpathlineto{\pgfqpoint{1.456513in}{1.434887in}}%
\pgfpathlineto{\pgfqpoint{1.461338in}{1.445745in}}%
\pgfpathlineto{\pgfqpoint{1.466157in}{1.456592in}}%
\pgfpathlineto{\pgfqpoint{1.470972in}{1.467428in}}%
\pgfpathlineto{\pgfqpoint{1.475781in}{1.478252in}}%
\pgfpathlineto{\pgfqpoint{1.480586in}{1.489065in}}%
\pgfpathlineto{\pgfqpoint{1.485386in}{1.499867in}}%
\pgfusepath{stroke}%
\end{pgfscope}%
\begin{pgfscope}%
\pgfpathrectangle{\pgfqpoint{0.407592in}{0.222904in}}{\pgfqpoint{1.411163in}{1.411163in}}%
\pgfusepath{clip}%
\pgfsetrectcap%
\pgfsetroundjoin%
\pgfsetlinewidth{0.401500pt}%
\definecolor{currentstroke}{rgb}{0.560784,0.588235,0.611765}%
\pgfsetstrokecolor{currentstroke}%
\pgfsetstrokeopacity{0.350000}%
\pgfsetdash{}{0pt}%
\pgfpathmoveto{\pgfqpoint{1.090230in}{0.610533in}}%
\pgfpathlineto{\pgfqpoint{1.088307in}{0.622729in}}%
\pgfpathlineto{\pgfqpoint{1.086386in}{0.634909in}}%
\pgfpathlineto{\pgfqpoint{1.084468in}{0.647072in}}%
\pgfpathlineto{\pgfqpoint{1.082553in}{0.659219in}}%
\pgfpathlineto{\pgfqpoint{1.080640in}{0.671350in}}%
\pgfpathlineto{\pgfqpoint{1.078730in}{0.683464in}}%
\pgfpathlineto{\pgfqpoint{1.076822in}{0.695563in}}%
\pgfpathlineto{\pgfqpoint{1.074917in}{0.707645in}}%
\pgfpathlineto{\pgfqpoint{1.073015in}{0.719711in}}%
\pgfpathlineto{\pgfqpoint{1.071115in}{0.731762in}}%
\pgfpathlineto{\pgfqpoint{1.069217in}{0.743796in}}%
\pgfpathlineto{\pgfqpoint{1.067322in}{0.755814in}}%
\pgfpathlineto{\pgfqpoint{1.065430in}{0.767816in}}%
\pgfpathlineto{\pgfqpoint{1.063540in}{0.779802in}}%
\pgfpathlineto{\pgfqpoint{1.061653in}{0.791772in}}%
\pgfpathlineto{\pgfqpoint{1.059768in}{0.803726in}}%
\pgfpathlineto{\pgfqpoint{1.057885in}{0.815665in}}%
\pgfpathlineto{\pgfqpoint{1.056005in}{0.827587in}}%
\pgfpathlineto{\pgfqpoint{1.054128in}{0.839494in}}%
\pgfpathlineto{\pgfqpoint{1.052253in}{0.851385in}}%
\pgfpathlineto{\pgfqpoint{1.050381in}{0.863260in}}%
\pgfpathlineto{\pgfqpoint{1.048511in}{0.875120in}}%
\pgfpathlineto{\pgfqpoint{1.046643in}{0.886964in}}%
\pgfpathlineto{\pgfqpoint{1.044778in}{0.898792in}}%
\pgfpathlineto{\pgfqpoint{1.042915in}{0.910605in}}%
\pgfpathlineto{\pgfqpoint{1.041055in}{0.922402in}}%
\pgfpathlineto{\pgfqpoint{1.039198in}{0.934184in}}%
\pgfpathlineto{\pgfqpoint{1.037342in}{0.945950in}}%
\pgfpathlineto{\pgfqpoint{1.035490in}{0.957700in}}%
\pgfpathlineto{\pgfqpoint{1.033639in}{0.969435in}}%
\pgfpathlineto{\pgfqpoint{1.031791in}{0.981155in}}%
\pgfpathlineto{\pgfqpoint{1.029946in}{0.992859in}}%
\pgfpathlineto{\pgfqpoint{1.028103in}{1.004548in}}%
\pgfpathlineto{\pgfqpoint{1.026262in}{1.016222in}}%
\pgfpathlineto{\pgfqpoint{1.024424in}{1.027880in}}%
\pgfpathlineto{\pgfqpoint{1.022588in}{1.039523in}}%
\pgfpathlineto{\pgfqpoint{1.020755in}{1.051151in}}%
\pgfpathlineto{\pgfqpoint{1.018924in}{1.062763in}}%
\pgfpathlineto{\pgfqpoint{1.017095in}{1.074361in}}%
\pgfpathlineto{\pgfqpoint{1.015269in}{1.085943in}}%
\pgfpathlineto{\pgfqpoint{1.013445in}{1.097510in}}%
\pgfpathlineto{\pgfqpoint{1.011624in}{1.109062in}}%
\pgfpathlineto{\pgfqpoint{1.009804in}{1.120599in}}%
\pgfpathlineto{\pgfqpoint{1.007988in}{1.132121in}}%
\pgfpathlineto{\pgfqpoint{1.006173in}{1.143628in}}%
\pgfpathlineto{\pgfqpoint{1.004361in}{1.155120in}}%
\pgfpathlineto{\pgfqpoint{1.002552in}{1.166597in}}%
\pgfpathlineto{\pgfqpoint{1.000744in}{1.178059in}}%
\pgfpathlineto{\pgfqpoint{0.998940in}{1.189506in}}%
\pgfpathlineto{\pgfqpoint{0.997137in}{1.200938in}}%
\pgfpathlineto{\pgfqpoint{0.995337in}{1.212355in}}%
\pgfpathlineto{\pgfqpoint{0.993539in}{1.223758in}}%
\pgfpathlineto{\pgfqpoint{0.991743in}{1.235146in}}%
\pgfpathlineto{\pgfqpoint{0.989950in}{1.246519in}}%
\pgfpathlineto{\pgfqpoint{0.988159in}{1.257878in}}%
\pgfpathlineto{\pgfqpoint{0.986370in}{1.269221in}}%
\pgfpathlineto{\pgfqpoint{0.984584in}{1.280550in}}%
\pgfpathlineto{\pgfqpoint{0.982800in}{1.291865in}}%
\pgfpathlineto{\pgfqpoint{0.981018in}{1.303165in}}%
\pgfpathlineto{\pgfqpoint{0.979239in}{1.314450in}}%
\pgfpathlineto{\pgfqpoint{0.977462in}{1.325721in}}%
\pgfpathlineto{\pgfqpoint{0.975687in}{1.336977in}}%
\pgfpathlineto{\pgfqpoint{0.973914in}{1.348219in}}%
\pgfpathlineto{\pgfqpoint{0.972144in}{1.359446in}}%
\pgfpathlineto{\pgfqpoint{0.970376in}{1.370659in}}%
\pgfpathlineto{\pgfqpoint{0.968610in}{1.381857in}}%
\pgfpathlineto{\pgfqpoint{0.966847in}{1.393041in}}%
\pgfpathlineto{\pgfqpoint{0.965086in}{1.404211in}}%
\pgfpathlineto{\pgfqpoint{0.963327in}{1.415366in}}%
\pgfpathlineto{\pgfqpoint{0.961570in}{1.426507in}}%
\pgfpathlineto{\pgfqpoint{0.959816in}{1.437634in}}%
\pgfpathlineto{\pgfqpoint{0.958064in}{1.448746in}}%
\pgfpathlineto{\pgfqpoint{0.956314in}{1.459845in}}%
\pgfpathlineto{\pgfqpoint{0.954566in}{1.470929in}}%
\pgfpathlineto{\pgfqpoint{0.952820in}{1.481999in}}%
\pgfpathlineto{\pgfqpoint{0.951077in}{1.493055in}}%
\pgfpathlineto{\pgfqpoint{0.949336in}{1.504096in}}%
\pgfpathlineto{\pgfqpoint{0.947597in}{1.515124in}}%
\pgfpathlineto{\pgfqpoint{0.945861in}{1.526138in}}%
\pgfusepath{stroke}%
\end{pgfscope}%
\begin{pgfscope}%
\pgfpathrectangle{\pgfqpoint{0.407592in}{0.222904in}}{\pgfqpoint{1.411163in}{1.411163in}}%
\pgfusepath{clip}%
\pgfsetrectcap%
\pgfsetroundjoin%
\pgfsetlinewidth{0.401500pt}%
\definecolor{currentstroke}{rgb}{0.560784,0.588235,0.611765}%
\pgfsetstrokecolor{currentstroke}%
\pgfsetstrokeopacity{0.350000}%
\pgfsetdash{}{0pt}%
\pgfpathmoveto{\pgfqpoint{1.090230in}{0.610533in}}%
\pgfpathlineto{\pgfqpoint{1.082282in}{0.620829in}}%
\pgfpathlineto{\pgfqpoint{1.074337in}{0.631122in}}%
\pgfpathlineto{\pgfqpoint{1.066394in}{0.641412in}}%
\pgfpathlineto{\pgfqpoint{1.058454in}{0.651699in}}%
\pgfpathlineto{\pgfqpoint{1.050515in}{0.661983in}}%
\pgfpathlineto{\pgfqpoint{1.042579in}{0.672265in}}%
\pgfpathlineto{\pgfqpoint{1.034645in}{0.682544in}}%
\pgfpathlineto{\pgfqpoint{1.026713in}{0.692820in}}%
\pgfpathlineto{\pgfqpoint{1.018783in}{0.703093in}}%
\pgfpathlineto{\pgfqpoint{1.010856in}{0.713363in}}%
\pgfpathlineto{\pgfqpoint{1.002930in}{0.723631in}}%
\pgfpathlineto{\pgfqpoint{0.995007in}{0.733896in}}%
\pgfpathlineto{\pgfqpoint{0.987086in}{0.744157in}}%
\pgfpathlineto{\pgfqpoint{0.979167in}{0.754417in}}%
\pgfpathlineto{\pgfqpoint{0.971250in}{0.764673in}}%
\pgfpathlineto{\pgfqpoint{0.963336in}{0.774926in}}%
\pgfpathlineto{\pgfqpoint{0.955423in}{0.785177in}}%
\pgfpathlineto{\pgfqpoint{0.947513in}{0.795425in}}%
\pgfpathlineto{\pgfqpoint{0.939605in}{0.805670in}}%
\pgfpathlineto{\pgfqpoint{0.931699in}{0.815912in}}%
\pgfpathlineto{\pgfqpoint{0.923795in}{0.826151in}}%
\pgfpathlineto{\pgfqpoint{0.915894in}{0.836388in}}%
\pgfpathlineto{\pgfqpoint{0.907994in}{0.846622in}}%
\pgfpathlineto{\pgfqpoint{0.900097in}{0.856853in}}%
\pgfpathlineto{\pgfqpoint{0.892202in}{0.867081in}}%
\pgfpathlineto{\pgfqpoint{0.884309in}{0.877306in}}%
\pgfpathlineto{\pgfqpoint{0.876418in}{0.887529in}}%
\pgfpathlineto{\pgfqpoint{0.868530in}{0.897749in}}%
\pgfpathlineto{\pgfqpoint{0.860643in}{0.907966in}}%
\pgfpathlineto{\pgfqpoint{0.852759in}{0.918180in}}%
\pgfpathlineto{\pgfqpoint{0.844877in}{0.928391in}}%
\pgfpathlineto{\pgfqpoint{0.836997in}{0.938600in}}%
\pgfpathlineto{\pgfqpoint{0.829119in}{0.948806in}}%
\pgfpathlineto{\pgfqpoint{0.821244in}{0.959009in}}%
\pgfpathlineto{\pgfqpoint{0.813370in}{0.969209in}}%
\pgfpathlineto{\pgfqpoint{0.805499in}{0.979406in}}%
\pgfpathlineto{\pgfqpoint{0.797630in}{0.989601in}}%
\pgfpathlineto{\pgfqpoint{0.789763in}{0.999793in}}%
\pgfpathlineto{\pgfqpoint{0.781898in}{1.009982in}}%
\pgfpathlineto{\pgfqpoint{0.774035in}{1.020168in}}%
\pgfpathlineto{\pgfqpoint{0.766175in}{1.030352in}}%
\pgfpathlineto{\pgfqpoint{0.758316in}{1.040532in}}%
\pgfpathlineto{\pgfqpoint{0.750460in}{1.050710in}}%
\pgfpathlineto{\pgfqpoint{0.742606in}{1.060886in}}%
\pgfpathlineto{\pgfqpoint{0.734754in}{1.071058in}}%
\pgfpathlineto{\pgfqpoint{0.726904in}{1.081227in}}%
\pgfpathlineto{\pgfqpoint{0.719056in}{1.091394in}}%
\pgfpathlineto{\pgfqpoint{0.711211in}{1.101558in}}%
\pgfpathlineto{\pgfqpoint{0.703367in}{1.111720in}}%
\pgfpathlineto{\pgfqpoint{0.695526in}{1.121878in}}%
\pgfpathlineto{\pgfqpoint{0.687687in}{1.132034in}}%
\pgfpathlineto{\pgfqpoint{0.679850in}{1.142187in}}%
\pgfpathlineto{\pgfqpoint{0.672015in}{1.152337in}}%
\pgfpathlineto{\pgfqpoint{0.664182in}{1.162484in}}%
\pgfpathlineto{\pgfqpoint{0.656352in}{1.172629in}}%
\pgfpathlineto{\pgfqpoint{0.648523in}{1.182771in}}%
\pgfpathlineto{\pgfqpoint{0.640697in}{1.192910in}}%
\pgfpathlineto{\pgfqpoint{0.632873in}{1.203046in}}%
\pgfpathlineto{\pgfqpoint{0.625051in}{1.213180in}}%
\pgfpathlineto{\pgfqpoint{0.617231in}{1.223311in}}%
\pgfpathlineto{\pgfqpoint{0.609413in}{1.233439in}}%
\pgfpathlineto{\pgfqpoint{0.601598in}{1.243564in}}%
\pgfpathlineto{\pgfqpoint{0.593784in}{1.253687in}}%
\pgfpathlineto{\pgfqpoint{0.585973in}{1.263807in}}%
\pgfpathlineto{\pgfqpoint{0.578163in}{1.273924in}}%
\pgfpathlineto{\pgfqpoint{0.570356in}{1.284038in}}%
\pgfpathlineto{\pgfqpoint{0.562551in}{1.294149in}}%
\pgfpathlineto{\pgfqpoint{0.554748in}{1.304258in}}%
\pgfpathlineto{\pgfqpoint{0.546948in}{1.314364in}}%
\pgfpathlineto{\pgfqpoint{0.539149in}{1.324467in}}%
\pgfpathlineto{\pgfqpoint{0.531352in}{1.334568in}}%
\pgfpathlineto{\pgfqpoint{0.523558in}{1.344666in}}%
\pgfpathlineto{\pgfqpoint{0.515766in}{1.354761in}}%
\pgfpathlineto{\pgfqpoint{0.507976in}{1.364853in}}%
\pgfpathlineto{\pgfqpoint{0.500188in}{1.374942in}}%
\pgfpathlineto{\pgfqpoint{0.492402in}{1.385029in}}%
\pgfpathlineto{\pgfqpoint{0.484618in}{1.395113in}}%
\pgfpathlineto{\pgfqpoint{0.476836in}{1.405195in}}%
\pgfpathlineto{\pgfqpoint{0.469057in}{1.415273in}}%
\pgfusepath{stroke}%
\end{pgfscope}%
\begin{pgfscope}%
\pgfpathrectangle{\pgfqpoint{0.407592in}{0.222904in}}{\pgfqpoint{1.411163in}{1.411163in}}%
\pgfusepath{clip}%
\pgfsetrectcap%
\pgfsetroundjoin%
\pgfsetlinewidth{0.401500pt}%
\definecolor{currentstroke}{rgb}{0.560784,0.588235,0.611765}%
\pgfsetstrokecolor{currentstroke}%
\pgfsetstrokeopacity{0.350000}%
\pgfsetdash{}{0pt}%
\pgfpathmoveto{\pgfqpoint{1.090230in}{0.610533in}}%
\pgfpathlineto{\pgfqpoint{1.080895in}{0.617672in}}%
\pgfpathlineto{\pgfqpoint{1.071546in}{0.624823in}}%
\pgfpathlineto{\pgfqpoint{1.062184in}{0.631983in}}%
\pgfpathlineto{\pgfqpoint{1.052807in}{0.639155in}}%
\pgfpathlineto{\pgfqpoint{1.043416in}{0.646337in}}%
\pgfpathlineto{\pgfqpoint{1.034011in}{0.653531in}}%
\pgfpathlineto{\pgfqpoint{1.024592in}{0.660734in}}%
\pgfpathlineto{\pgfqpoint{1.015159in}{0.667949in}}%
\pgfpathlineto{\pgfqpoint{1.005712in}{0.675175in}}%
\pgfpathlineto{\pgfqpoint{0.996250in}{0.682411in}}%
\pgfpathlineto{\pgfqpoint{0.986775in}{0.689659in}}%
\pgfpathlineto{\pgfqpoint{0.977284in}{0.696917in}}%
\pgfpathlineto{\pgfqpoint{0.967780in}{0.704187in}}%
\pgfpathlineto{\pgfqpoint{0.958261in}{0.711467in}}%
\pgfpathlineto{\pgfqpoint{0.948728in}{0.718758in}}%
\pgfpathlineto{\pgfqpoint{0.939180in}{0.726060in}}%
\pgfpathlineto{\pgfqpoint{0.929618in}{0.733374in}}%
\pgfpathlineto{\pgfqpoint{0.920042in}{0.740698in}}%
\pgfpathlineto{\pgfqpoint{0.910451in}{0.748034in}}%
\pgfpathlineto{\pgfqpoint{0.900845in}{0.755381in}}%
\pgfpathlineto{\pgfqpoint{0.891225in}{0.762739in}}%
\pgfpathlineto{\pgfqpoint{0.881590in}{0.770108in}}%
\pgfpathlineto{\pgfqpoint{0.871941in}{0.777488in}}%
\pgfpathlineto{\pgfqpoint{0.862276in}{0.784879in}}%
\pgfpathlineto{\pgfqpoint{0.852597in}{0.792282in}}%
\pgfpathlineto{\pgfqpoint{0.842904in}{0.799696in}}%
\pgfpathlineto{\pgfqpoint{0.833195in}{0.807121in}}%
\pgfpathlineto{\pgfqpoint{0.823472in}{0.814558in}}%
\pgfpathlineto{\pgfqpoint{0.813734in}{0.822006in}}%
\pgfpathlineto{\pgfqpoint{0.803981in}{0.829466in}}%
\pgfpathlineto{\pgfqpoint{0.794213in}{0.836936in}}%
\pgfpathlineto{\pgfqpoint{0.784430in}{0.844419in}}%
\pgfpathlineto{\pgfqpoint{0.774632in}{0.851912in}}%
\pgfpathlineto{\pgfqpoint{0.764819in}{0.859418in}}%
\pgfpathlineto{\pgfqpoint{0.754991in}{0.866934in}}%
\pgfpathlineto{\pgfqpoint{0.745148in}{0.874463in}}%
\pgfpathlineto{\pgfqpoint{0.735290in}{0.882003in}}%
\pgfpathlineto{\pgfqpoint{0.725417in}{0.889554in}}%
\pgfpathlineto{\pgfqpoint{0.715528in}{0.897117in}}%
\pgfpathlineto{\pgfqpoint{0.705624in}{0.904692in}}%
\pgfpathlineto{\pgfqpoint{0.695705in}{0.912278in}}%
\pgfpathlineto{\pgfqpoint{0.685771in}{0.919877in}}%
\pgfpathlineto{\pgfqpoint{0.675821in}{0.927487in}}%
\pgfpathlineto{\pgfqpoint{0.665856in}{0.935108in}}%
\pgfpathlineto{\pgfqpoint{0.655876in}{0.942742in}}%
\pgfpathlineto{\pgfqpoint{0.645880in}{0.950387in}}%
\pgfpathlineto{\pgfqpoint{0.635868in}{0.958044in}}%
\pgfpathlineto{\pgfqpoint{0.625841in}{0.965713in}}%
\pgfpathlineto{\pgfqpoint{0.615799in}{0.973394in}}%
\pgfpathlineto{\pgfqpoint{0.605741in}{0.981087in}}%
\pgfpathlineto{\pgfqpoint{0.595667in}{0.988791in}}%
\pgfpathlineto{\pgfqpoint{0.585578in}{0.996508in}}%
\pgfpathlineto{\pgfqpoint{0.575472in}{1.004237in}}%
\pgfpathlineto{\pgfqpoint{0.565352in}{1.011978in}}%
\pgfpathlineto{\pgfqpoint{0.555215in}{1.019730in}}%
\pgfpathlineto{\pgfqpoint{0.545062in}{1.027495in}}%
\pgfpathlineto{\pgfqpoint{0.534894in}{1.035272in}}%
\pgfpathlineto{\pgfqpoint{0.524710in}{1.043062in}}%
\pgfpathlineto{\pgfqpoint{0.514510in}{1.050863in}}%
\pgfpathlineto{\pgfqpoint{0.504294in}{1.058677in}}%
\pgfpathlineto{\pgfqpoint{0.494062in}{1.066502in}}%
\pgfpathlineto{\pgfqpoint{0.483814in}{1.074341in}}%
\pgfpathlineto{\pgfqpoint{0.473549in}{1.082191in}}%
\pgfpathlineto{\pgfqpoint{0.463269in}{1.090054in}}%
\pgfpathlineto{\pgfqpoint{0.452973in}{1.097929in}}%
\pgfpathlineto{\pgfqpoint{0.442660in}{1.105816in}}%
\pgfpathlineto{\pgfqpoint{0.432331in}{1.113716in}}%
\pgfpathlineto{\pgfqpoint{0.421986in}{1.121628in}}%
\pgfpathlineto{\pgfqpoint{0.411625in}{1.129553in}}%
\pgfpathlineto{\pgfqpoint{0.404258in}{1.135187in}}%
\pgfusepath{stroke}%
\end{pgfscope}%
\begin{pgfscope}%
\pgfpathrectangle{\pgfqpoint{0.407592in}{0.222904in}}{\pgfqpoint{1.411163in}{1.411163in}}%
\pgfusepath{clip}%
\pgfsetrectcap%
\pgfsetroundjoin%
\pgfsetlinewidth{0.401500pt}%
\definecolor{currentstroke}{rgb}{0.560784,0.588235,0.611765}%
\pgfsetstrokecolor{currentstroke}%
\pgfsetstrokeopacity{0.350000}%
\pgfsetdash{}{0pt}%
\pgfpathmoveto{\pgfqpoint{1.090230in}{0.610533in}}%
\pgfpathlineto{\pgfqpoint{1.084967in}{0.615108in}}%
\pgfpathlineto{\pgfqpoint{1.079690in}{0.619696in}}%
\pgfpathlineto{\pgfqpoint{1.074397in}{0.624297in}}%
\pgfpathlineto{\pgfqpoint{1.069088in}{0.628912in}}%
\pgfpathlineto{\pgfqpoint{1.063764in}{0.633541in}}%
\pgfpathlineto{\pgfqpoint{1.058424in}{0.638184in}}%
\pgfpathlineto{\pgfqpoint{1.053068in}{0.642840in}}%
\pgfpathlineto{\pgfqpoint{1.047696in}{0.647510in}}%
\pgfpathlineto{\pgfqpoint{1.042308in}{0.652194in}}%
\pgfpathlineto{\pgfqpoint{1.036904in}{0.656891in}}%
\pgfpathlineto{\pgfqpoint{1.031484in}{0.661603in}}%
\pgfpathlineto{\pgfqpoint{1.026048in}{0.666329in}}%
\pgfpathlineto{\pgfqpoint{1.020596in}{0.671069in}}%
\pgfpathlineto{\pgfqpoint{1.015127in}{0.675823in}}%
\pgfpathlineto{\pgfqpoint{1.009642in}{0.680592in}}%
\pgfpathlineto{\pgfqpoint{1.004141in}{0.685374in}}%
\pgfpathlineto{\pgfqpoint{0.998623in}{0.690171in}}%
\pgfpathlineto{\pgfqpoint{0.993088in}{0.694983in}}%
\pgfpathlineto{\pgfqpoint{0.987537in}{0.699809in}}%
\pgfpathlineto{\pgfqpoint{0.981969in}{0.704649in}}%
\pgfpathlineto{\pgfqpoint{0.976384in}{0.709505in}}%
\pgfpathlineto{\pgfqpoint{0.970782in}{0.714375in}}%
\pgfpathlineto{\pgfqpoint{0.965164in}{0.719259in}}%
\pgfpathlineto{\pgfqpoint{0.959528in}{0.724159in}}%
\pgfpathlineto{\pgfqpoint{0.953875in}{0.729073in}}%
\pgfpathlineto{\pgfqpoint{0.948205in}{0.734003in}}%
\pgfpathlineto{\pgfqpoint{0.942517in}{0.738947in}}%
\pgfpathlineto{\pgfqpoint{0.936812in}{0.743907in}}%
\pgfpathlineto{\pgfqpoint{0.931090in}{0.748881in}}%
\pgfpathlineto{\pgfqpoint{0.925350in}{0.753871in}}%
\pgfpathlineto{\pgfqpoint{0.919592in}{0.758877in}}%
\pgfpathlineto{\pgfqpoint{0.913817in}{0.763897in}}%
\pgfpathlineto{\pgfqpoint{0.908024in}{0.768933in}}%
\pgfpathlineto{\pgfqpoint{0.902213in}{0.773985in}}%
\pgfpathlineto{\pgfqpoint{0.896385in}{0.779052in}}%
\pgfpathlineto{\pgfqpoint{0.890538in}{0.784135in}}%
\pgfpathlineto{\pgfqpoint{0.884673in}{0.789234in}}%
\pgfpathlineto{\pgfqpoint{0.878790in}{0.794349in}}%
\pgfpathlineto{\pgfqpoint{0.872888in}{0.799479in}}%
\pgfpathlineto{\pgfqpoint{0.866968in}{0.804625in}}%
\pgfpathlineto{\pgfqpoint{0.861030in}{0.809788in}}%
\pgfpathlineto{\pgfqpoint{0.855073in}{0.814966in}}%
\pgfpathlineto{\pgfqpoint{0.849098in}{0.820161in}}%
\pgfpathlineto{\pgfqpoint{0.843104in}{0.825372in}}%
\pgfpathlineto{\pgfqpoint{0.837091in}{0.830600in}}%
\pgfpathlineto{\pgfqpoint{0.831059in}{0.835843in}}%
\pgfpathlineto{\pgfqpoint{0.825008in}{0.841104in}}%
\pgfpathlineto{\pgfqpoint{0.818938in}{0.846381in}}%
\pgfpathlineto{\pgfqpoint{0.812849in}{0.851674in}}%
\pgfpathlineto{\pgfqpoint{0.806741in}{0.856984in}}%
\pgfpathlineto{\pgfqpoint{0.800613in}{0.862311in}}%
\pgfpathlineto{\pgfqpoint{0.794466in}{0.867655in}}%
\pgfpathlineto{\pgfqpoint{0.788299in}{0.873016in}}%
\pgfpathlineto{\pgfqpoint{0.782113in}{0.878394in}}%
\pgfpathlineto{\pgfqpoint{0.775907in}{0.883790in}}%
\pgfpathlineto{\pgfqpoint{0.769681in}{0.889202in}}%
\pgfpathlineto{\pgfqpoint{0.763436in}{0.894632in}}%
\pgfpathlineto{\pgfqpoint{0.757170in}{0.900079in}}%
\pgfpathlineto{\pgfqpoint{0.750884in}{0.905543in}}%
\pgfpathlineto{\pgfqpoint{0.744578in}{0.911026in}}%
\pgfpathlineto{\pgfqpoint{0.738252in}{0.916525in}}%
\pgfpathlineto{\pgfqpoint{0.731905in}{0.922043in}}%
\pgfpathlineto{\pgfqpoint{0.725538in}{0.927578in}}%
\pgfpathlineto{\pgfqpoint{0.719150in}{0.933131in}}%
\pgfpathlineto{\pgfqpoint{0.712742in}{0.938703in}}%
\pgfpathlineto{\pgfqpoint{0.706312in}{0.944292in}}%
\pgfpathlineto{\pgfqpoint{0.699862in}{0.949899in}}%
\pgfpathlineto{\pgfqpoint{0.693391in}{0.955525in}}%
\pgfpathlineto{\pgfqpoint{0.686899in}{0.961169in}}%
\pgfpathlineto{\pgfqpoint{0.680385in}{0.966831in}}%
\pgfpathlineto{\pgfqpoint{0.673851in}{0.972512in}}%
\pgfpathlineto{\pgfqpoint{0.667295in}{0.978212in}}%
\pgfpathlineto{\pgfqpoint{0.660717in}{0.983930in}}%
\pgfpathlineto{\pgfqpoint{0.654118in}{0.989667in}}%
\pgfpathlineto{\pgfqpoint{0.647497in}{0.995423in}}%
\pgfpathlineto{\pgfqpoint{0.640854in}{1.001198in}}%
\pgfpathlineto{\pgfqpoint{0.634189in}{1.006992in}}%
\pgfpathlineto{\pgfqpoint{0.627502in}{1.012806in}}%
\pgfpathlineto{\pgfqpoint{0.620793in}{1.018638in}}%
\pgfusepath{stroke}%
\end{pgfscope}%
\begin{pgfscope}%
\pgfpathrectangle{\pgfqpoint{0.407592in}{0.222904in}}{\pgfqpoint{1.411163in}{1.411163in}}%
\pgfusepath{clip}%
\pgfsetrectcap%
\pgfsetroundjoin%
\pgfsetlinewidth{0.401500pt}%
\definecolor{currentstroke}{rgb}{0.560784,0.588235,0.611765}%
\pgfsetstrokecolor{currentstroke}%
\pgfsetstrokeopacity{0.350000}%
\pgfsetdash{}{0pt}%
\pgfpathmoveto{\pgfqpoint{1.090230in}{0.610533in}}%
\pgfpathlineto{\pgfqpoint{1.092118in}{0.614641in}}%
\pgfpathlineto{\pgfqpoint{1.094013in}{0.618763in}}%
\pgfpathlineto{\pgfqpoint{1.095913in}{0.622898in}}%
\pgfpathlineto{\pgfqpoint{1.097820in}{0.627046in}}%
\pgfpathlineto{\pgfqpoint{1.099733in}{0.631207in}}%
\pgfpathlineto{\pgfqpoint{1.101652in}{0.635382in}}%
\pgfpathlineto{\pgfqpoint{1.103577in}{0.639570in}}%
\pgfpathlineto{\pgfqpoint{1.105508in}{0.643772in}}%
\pgfpathlineto{\pgfqpoint{1.107446in}{0.647987in}}%
\pgfpathlineto{\pgfqpoint{1.109390in}{0.652216in}}%
\pgfpathlineto{\pgfqpoint{1.111340in}{0.656459in}}%
\pgfpathlineto{\pgfqpoint{1.113297in}{0.660715in}}%
\pgfpathlineto{\pgfqpoint{1.115260in}{0.664986in}}%
\pgfpathlineto{\pgfqpoint{1.117229in}{0.669270in}}%
\pgfpathlineto{\pgfqpoint{1.119205in}{0.673569in}}%
\pgfpathlineto{\pgfqpoint{1.121187in}{0.677881in}}%
\pgfpathlineto{\pgfqpoint{1.123176in}{0.682208in}}%
\pgfpathlineto{\pgfqpoint{1.125171in}{0.686549in}}%
\pgfpathlineto{\pgfqpoint{1.127173in}{0.690904in}}%
\pgfpathlineto{\pgfqpoint{1.129182in}{0.695273in}}%
\pgfpathlineto{\pgfqpoint{1.131197in}{0.699657in}}%
\pgfpathlineto{\pgfqpoint{1.133219in}{0.704056in}}%
\pgfpathlineto{\pgfqpoint{1.135247in}{0.708469in}}%
\pgfpathlineto{\pgfqpoint{1.137282in}{0.712896in}}%
\pgfpathlineto{\pgfqpoint{1.139324in}{0.717338in}}%
\pgfpathlineto{\pgfqpoint{1.141373in}{0.721796in}}%
\pgfpathlineto{\pgfqpoint{1.143429in}{0.726268in}}%
\pgfpathlineto{\pgfqpoint{1.145491in}{0.730755in}}%
\pgfpathlineto{\pgfqpoint{1.147561in}{0.735256in}}%
\pgfpathlineto{\pgfqpoint{1.149637in}{0.739774in}}%
\pgfpathlineto{\pgfqpoint{1.151720in}{0.744306in}}%
\pgfpathlineto{\pgfqpoint{1.153811in}{0.748853in}}%
\pgfpathlineto{\pgfqpoint{1.155908in}{0.753416in}}%
\pgfpathlineto{\pgfqpoint{1.158012in}{0.757994in}}%
\pgfpathlineto{\pgfqpoint{1.160124in}{0.762588in}}%
\pgfpathlineto{\pgfqpoint{1.162242in}{0.767197in}}%
\pgfpathlineto{\pgfqpoint{1.164368in}{0.771822in}}%
\pgfpathlineto{\pgfqpoint{1.166501in}{0.776462in}}%
\pgfpathlineto{\pgfqpoint{1.168642in}{0.781119in}}%
\pgfpathlineto{\pgfqpoint{1.170789in}{0.785791in}}%
\pgfpathlineto{\pgfqpoint{1.172944in}{0.790479in}}%
\pgfpathlineto{\pgfqpoint{1.175107in}{0.795183in}}%
\pgfpathlineto{\pgfqpoint{1.177276in}{0.799903in}}%
\pgfpathlineto{\pgfqpoint{1.179454in}{0.804640in}}%
\pgfpathlineto{\pgfqpoint{1.181638in}{0.809393in}}%
\pgfpathlineto{\pgfqpoint{1.183831in}{0.814162in}}%
\pgfpathlineto{\pgfqpoint{1.186030in}{0.818947in}}%
\pgfpathlineto{\pgfqpoint{1.188238in}{0.823750in}}%
\pgfpathlineto{\pgfqpoint{1.190453in}{0.828568in}}%
\pgfpathlineto{\pgfqpoint{1.192675in}{0.833404in}}%
\pgfpathlineto{\pgfqpoint{1.194906in}{0.838256in}}%
\pgfpathlineto{\pgfqpoint{1.197144in}{0.843125in}}%
\pgfpathlineto{\pgfqpoint{1.199390in}{0.848011in}}%
\pgfpathlineto{\pgfqpoint{1.201644in}{0.852915in}}%
\pgfpathlineto{\pgfqpoint{1.203906in}{0.857835in}}%
\pgfpathlineto{\pgfqpoint{1.206175in}{0.862773in}}%
\pgfpathlineto{\pgfqpoint{1.208453in}{0.867727in}}%
\pgfpathlineto{\pgfqpoint{1.210738in}{0.872700in}}%
\pgfpathlineto{\pgfqpoint{1.213032in}{0.877690in}}%
\pgfpathlineto{\pgfqpoint{1.215334in}{0.882697in}}%
\pgfpathlineto{\pgfqpoint{1.217644in}{0.887722in}}%
\pgfpathlineto{\pgfqpoint{1.219962in}{0.892765in}}%
\pgfpathlineto{\pgfqpoint{1.222288in}{0.897826in}}%
\pgfpathlineto{\pgfqpoint{1.224623in}{0.902905in}}%
\pgfpathlineto{\pgfqpoint{1.226966in}{0.908002in}}%
\pgfpathlineto{\pgfqpoint{1.229317in}{0.913117in}}%
\pgfpathlineto{\pgfqpoint{1.231676in}{0.918250in}}%
\pgfpathlineto{\pgfqpoint{1.234045in}{0.923402in}}%
\pgfpathlineto{\pgfqpoint{1.236421in}{0.928572in}}%
\pgfpathlineto{\pgfqpoint{1.238806in}{0.933761in}}%
\pgfpathlineto{\pgfqpoint{1.241200in}{0.938969in}}%
\pgfpathlineto{\pgfqpoint{1.243602in}{0.944195in}}%
\pgfpathlineto{\pgfqpoint{1.246013in}{0.949440in}}%
\pgfpathlineto{\pgfqpoint{1.248433in}{0.954704in}}%
\pgfpathlineto{\pgfqpoint{1.250862in}{0.959988in}}%
\pgfpathlineto{\pgfqpoint{1.253299in}{0.965290in}}%
\pgfpathlineto{\pgfqpoint{1.255745in}{0.970612in}}%
\pgfpathlineto{\pgfqpoint{1.258200in}{0.975953in}}%
\pgfpathlineto{\pgfqpoint{1.260664in}{0.981313in}}%
\pgfusepath{stroke}%
\end{pgfscope}%
\begin{pgfscope}%
\pgfpathrectangle{\pgfqpoint{0.407592in}{0.222904in}}{\pgfqpoint{1.411163in}{1.411163in}}%
\pgfusepath{clip}%
\pgfsetrectcap%
\pgfsetroundjoin%
\pgfsetlinewidth{1.806750pt}%
\definecolor{currentstroke}{rgb}{0.168627,0.611765,0.411765}%
\pgfsetstrokecolor{currentstroke}%
\pgfsetdash{}{0pt}%
\pgfpathmoveto{\pgfqpoint{1.247522in}{0.897771in}}%
\pgfpathlineto{\pgfqpoint{1.264082in}{0.890537in}}%
\pgfpathlineto{\pgfqpoint{1.280506in}{0.884135in}}%
\pgfpathlineto{\pgfqpoint{1.296799in}{0.878629in}}%
\pgfpathlineto{\pgfqpoint{1.310664in}{0.874677in}}%
\pgfpathlineto{\pgfqpoint{1.324439in}{0.871482in}}%
\pgfpathlineto{\pgfqpoint{1.338128in}{0.869093in}}%
\pgfpathlineto{\pgfqpoint{1.351734in}{0.867560in}}%
\pgfpathlineto{\pgfqpoint{1.363012in}{0.866973in}}%
\pgfpathlineto{\pgfqpoint{1.374236in}{0.867045in}}%
\pgfpathlineto{\pgfqpoint{1.385410in}{0.867806in}}%
\pgfpathlineto{\pgfqpoint{1.396535in}{0.869281in}}%
\pgfpathlineto{\pgfqpoint{1.407613in}{0.871498in}}%
\pgfpathlineto{\pgfqpoint{1.418647in}{0.874479in}}%
\pgfpathlineto{\pgfqpoint{1.429639in}{0.878244in}}%
\pgfpathlineto{\pgfqpoint{1.440591in}{0.882810in}}%
\pgfpathlineto{\pgfqpoint{1.451505in}{0.888189in}}%
\pgfpathlineto{\pgfqpoint{1.462383in}{0.894389in}}%
\pgfpathlineto{\pgfqpoint{1.473228in}{0.901413in}}%
\pgfpathlineto{\pgfqpoint{1.484040in}{0.909259in}}%
\pgfpathlineto{\pgfqpoint{1.494821in}{0.917918in}}%
\pgfpathlineto{\pgfqpoint{1.505574in}{0.927379in}}%
\pgfpathlineto{\pgfqpoint{1.518442in}{0.939765in}}%
\pgfpathlineto{\pgfqpoint{1.531271in}{0.953241in}}%
\pgfpathlineto{\pgfqpoint{1.544065in}{0.967763in}}%
\pgfpathlineto{\pgfqpoint{1.556825in}{0.983277in}}%
\pgfpathlineto{\pgfqpoint{1.571670in}{1.002559in}}%
\pgfpathlineto{\pgfqpoint{1.586472in}{1.023023in}}%
\pgfpathlineto{\pgfqpoint{1.603337in}{1.047737in}}%
\pgfpathlineto{\pgfqpoint{1.620147in}{1.073730in}}%
\pgfpathlineto{\pgfqpoint{1.638995in}{1.104332in}}%
\pgfpathlineto{\pgfqpoint{1.657776in}{1.136199in}}%
\pgfpathlineto{\pgfqpoint{1.678565in}{1.172889in}}%
\pgfpathlineto{\pgfqpoint{1.682713in}{1.180372in}}%
\pgfpathlineto{\pgfqpoint{1.682713in}{1.180372in}}%
\pgfusepath{stroke}%
\end{pgfscope}%
\begin{pgfscope}%
\pgfpathrectangle{\pgfqpoint{0.407592in}{0.222904in}}{\pgfqpoint{1.411163in}{1.411163in}}%
\pgfusepath{clip}%
\pgfsetrectcap%
\pgfsetroundjoin%
\pgfsetlinewidth{2.308625pt}%
\definecolor{currentstroke}{rgb}{0.141176,0.345098,0.650980}%
\pgfsetstrokecolor{currentstroke}%
\pgfsetdash{}{0pt}%
\pgfpathmoveto{\pgfqpoint{0.537673in}{1.095850in}}%
\pgfpathlineto{\pgfqpoint{0.583281in}{1.045325in}}%
\pgfpathlineto{\pgfqpoint{0.623058in}{1.002745in}}%
\pgfpathlineto{\pgfqpoint{0.658171in}{0.966597in}}%
\pgfpathlineto{\pgfqpoint{0.689504in}{0.935734in}}%
\pgfpathlineto{\pgfqpoint{0.717737in}{0.909269in}}%
\pgfpathlineto{\pgfqpoint{0.743404in}{0.886509in}}%
\pgfpathlineto{\pgfqpoint{0.766930in}{0.866900in}}%
\pgfpathlineto{\pgfqpoint{0.788654in}{0.849996in}}%
\pgfpathlineto{\pgfqpoint{0.808853in}{0.835433in}}%
\pgfpathlineto{\pgfqpoint{0.827752in}{0.822909in}}%
\pgfpathlineto{\pgfqpoint{0.845537in}{0.812174in}}%
\pgfpathlineto{\pgfqpoint{0.862362in}{0.803018in}}%
\pgfpathlineto{\pgfqpoint{0.878355in}{0.795260in}}%
\pgfpathlineto{\pgfqpoint{0.893624in}{0.788747in}}%
\pgfpathlineto{\pgfqpoint{0.908259in}{0.783345in}}%
\pgfpathlineto{\pgfqpoint{0.922333in}{0.778940in}}%
\pgfpathlineto{\pgfqpoint{0.935911in}{0.775428in}}%
\pgfpathlineto{\pgfqpoint{0.949044in}{0.772720in}}%
\pgfpathlineto{\pgfqpoint{0.961776in}{0.770736in}}%
\pgfpathlineto{\pgfqpoint{0.974145in}{0.769403in}}%
\pgfpathlineto{\pgfqpoint{0.990120in}{0.768527in}}%
\pgfpathlineto{\pgfqpoint{1.005557in}{0.768557in}}%
\pgfpathlineto{\pgfqpoint{1.020497in}{0.769369in}}%
\pgfpathlineto{\pgfqpoint{1.034971in}{0.770855in}}%
\pgfpathlineto{\pgfqpoint{1.052439in}{0.773510in}}%
\pgfpathlineto{\pgfqpoint{1.069237in}{0.776894in}}%
\pgfpathlineto{\pgfqpoint{1.085379in}{0.780854in}}%
\pgfpathlineto{\pgfqpoint{1.103888in}{0.786181in}}%
\pgfpathlineto{\pgfqpoint{1.124289in}{0.792937in}}%
\pgfpathlineto{\pgfqpoint{1.146015in}{0.801046in}}%
\pgfpathlineto{\pgfqpoint{1.168412in}{0.810280in}}%
\pgfpathlineto{\pgfqpoint{1.192850in}{0.821246in}}%
\pgfpathlineto{\pgfqpoint{1.217673in}{0.833215in}}%
\pgfpathlineto{\pgfqpoint{1.245306in}{0.847389in}}%
\pgfpathlineto{\pgfqpoint{1.275238in}{0.863608in}}%
\pgfpathlineto{\pgfqpoint{1.307959in}{0.882211in}}%
\pgfpathlineto{\pgfqpoint{1.318296in}{0.888253in}}%
\pgfpathlineto{\pgfqpoint{1.318296in}{0.888253in}}%
\pgfusepath{stroke}%
\end{pgfscope}%
\begin{pgfscope}%
\pgfpathrectangle{\pgfqpoint{0.407592in}{0.222904in}}{\pgfqpoint{1.411163in}{1.411163in}}%
\pgfusepath{clip}%
\pgfsetrectcap%
\pgfsetroundjoin%
\pgfsetlinewidth{2.308625pt}%
\definecolor{currentstroke}{rgb}{0.819608,0.286275,0.356863}%
\pgfsetstrokecolor{currentstroke}%
\pgfsetdash{}{0pt}%
\pgfpathmoveto{\pgfqpoint{0.672155in}{1.057312in}}%
\pgfpathlineto{\pgfqpoint{0.708311in}{1.016342in}}%
\pgfpathlineto{\pgfqpoint{0.740163in}{0.981221in}}%
\pgfpathlineto{\pgfqpoint{0.768505in}{0.950926in}}%
\pgfpathlineto{\pgfqpoint{0.793952in}{0.924665in}}%
\pgfpathlineto{\pgfqpoint{0.816990in}{0.901810in}}%
\pgfpathlineto{\pgfqpoint{0.838008in}{0.881864in}}%
\pgfpathlineto{\pgfqpoint{0.857320in}{0.864421in}}%
\pgfpathlineto{\pgfqpoint{0.875184in}{0.849151in}}%
\pgfpathlineto{\pgfqpoint{0.891812in}{0.835782in}}%
\pgfpathlineto{\pgfqpoint{0.907379in}{0.824087in}}%
\pgfpathlineto{\pgfqpoint{0.922035in}{0.813875in}}%
\pgfpathlineto{\pgfqpoint{0.935904in}{0.804984in}}%
\pgfpathlineto{\pgfqpoint{0.949090in}{0.797275in}}%
\pgfpathlineto{\pgfqpoint{0.961686in}{0.790631in}}%
\pgfpathlineto{\pgfqpoint{0.973766in}{0.784947in}}%
\pgfpathlineto{\pgfqpoint{0.985398in}{0.780135in}}%
\pgfpathlineto{\pgfqpoint{0.996636in}{0.776115in}}%
\pgfpathlineto{\pgfqpoint{1.007531in}{0.772818in}}%
\pgfpathlineto{\pgfqpoint{1.021591in}{0.769438in}}%
\pgfpathlineto{\pgfqpoint{1.035194in}{0.767104in}}%
\pgfpathlineto{\pgfqpoint{1.048406in}{0.765700in}}%
\pgfpathlineto{\pgfqpoint{1.061278in}{0.765128in}}%
\pgfpathlineto{\pgfqpoint{1.073852in}{0.765297in}}%
\pgfpathlineto{\pgfqpoint{1.086159in}{0.766126in}}%
\pgfpathlineto{\pgfqpoint{1.101203in}{0.767980in}}%
\pgfpathlineto{\pgfqpoint{1.115895in}{0.770622in}}%
\pgfpathlineto{\pgfqpoint{1.130252in}{0.773931in}}%
\pgfpathlineto{\pgfqpoint{1.147043in}{0.778635in}}%
\pgfpathlineto{\pgfqpoint{1.163349in}{0.783976in}}%
\pgfpathlineto{\pgfqpoint{1.181730in}{0.790809in}}%
\pgfpathlineto{\pgfqpoint{1.201832in}{0.799150in}}%
\pgfpathlineto{\pgfqpoint{1.223199in}{0.808874in}}%
\pgfpathlineto{\pgfqpoint{1.245264in}{0.819713in}}%
\pgfpathlineto{\pgfqpoint{1.271139in}{0.833297in}}%
\pgfpathlineto{\pgfqpoint{1.298486in}{0.848510in}}%
\pgfpathlineto{\pgfqpoint{1.329398in}{0.866578in}}%
\pgfpathlineto{\pgfqpoint{1.364170in}{0.887801in}}%
\pgfpathlineto{\pgfqpoint{1.396682in}{0.908341in}}%
\pgfpathlineto{\pgfqpoint{1.396682in}{0.908341in}}%
\pgfusepath{stroke}%
\end{pgfscope}%
\begin{pgfscope}%
\pgfpathrectangle{\pgfqpoint{0.407592in}{0.222904in}}{\pgfqpoint{1.411163in}{1.411163in}}%
\pgfusepath{clip}%
\pgfsetrectcap%
\pgfsetroundjoin%
\pgfsetlinewidth{2.308625pt}%
\definecolor{currentstroke}{rgb}{0.000000,0.529412,0.423529}%
\pgfsetstrokecolor{currentstroke}%
\pgfsetdash{}{0pt}%
\pgfpathmoveto{\pgfqpoint{0.812524in}{1.057517in}}%
\pgfpathlineto{\pgfqpoint{0.843065in}{1.026044in}}%
\pgfpathlineto{\pgfqpoint{0.870461in}{0.998966in}}%
\pgfpathlineto{\pgfqpoint{0.895259in}{0.975582in}}%
\pgfpathlineto{\pgfqpoint{0.917889in}{0.955335in}}%
\pgfpathlineto{\pgfqpoint{0.938697in}{0.937777in}}%
\pgfpathlineto{\pgfqpoint{0.957966in}{0.922543in}}%
\pgfpathlineto{\pgfqpoint{0.975926in}{0.909335in}}%
\pgfpathlineto{\pgfqpoint{0.992767in}{0.897903in}}%
\pgfpathlineto{\pgfqpoint{1.008646in}{0.888038in}}%
\pgfpathlineto{\pgfqpoint{1.023697in}{0.879564in}}%
\pgfpathlineto{\pgfqpoint{1.038029in}{0.872331in}}%
\pgfpathlineto{\pgfqpoint{1.051738in}{0.866208in}}%
\pgfpathlineto{\pgfqpoint{1.064901in}{0.861083in}}%
\pgfpathlineto{\pgfqpoint{1.077585in}{0.856858in}}%
\pgfpathlineto{\pgfqpoint{1.089847in}{0.853446in}}%
\pgfpathlineto{\pgfqpoint{1.101736in}{0.850771in}}%
\pgfpathlineto{\pgfqpoint{1.113290in}{0.848764in}}%
\pgfpathlineto{\pgfqpoint{1.128233in}{0.847022in}}%
\pgfpathlineto{\pgfqpoint{1.142706in}{0.846229in}}%
\pgfpathlineto{\pgfqpoint{1.156759in}{0.846271in}}%
\pgfpathlineto{\pgfqpoint{1.170429in}{0.847044in}}%
\pgfpathlineto{\pgfqpoint{1.183745in}{0.848456in}}%
\pgfpathlineto{\pgfqpoint{1.199918in}{0.850992in}}%
\pgfpathlineto{\pgfqpoint{1.215588in}{0.854249in}}%
\pgfpathlineto{\pgfqpoint{1.230765in}{0.858093in}}%
\pgfpathlineto{\pgfqpoint{1.248325in}{0.863316in}}%
\pgfpathlineto{\pgfqpoint{1.267896in}{0.870019in}}%
\pgfpathlineto{\pgfqpoint{1.286450in}{0.877142in}}%
\pgfpathlineto{\pgfqpoint{1.306365in}{0.885522in}}%
\pgfpathlineto{\pgfqpoint{1.329259in}{0.895981in}}%
\pgfpathlineto{\pgfqpoint{1.353738in}{0.908011in}}%
\pgfpathlineto{\pgfqpoint{1.379794in}{0.921636in}}%
\pgfpathlineto{\pgfqpoint{1.409527in}{0.938060in}}%
\pgfpathlineto{\pgfqpoint{1.441853in}{0.956797in}}%
\pgfpathlineto{\pgfqpoint{1.473472in}{0.975862in}}%
\pgfpathlineto{\pgfqpoint{1.473472in}{0.975862in}}%
\pgfusepath{stroke}%
\end{pgfscope}%
\begin{pgfscope}%
\pgfpathrectangle{\pgfqpoint{0.407592in}{0.222904in}}{\pgfqpoint{1.411163in}{1.411163in}}%
\pgfusepath{clip}%
\pgfsetrectcap%
\pgfsetroundjoin%
\pgfsetlinewidth{2.308625pt}%
\definecolor{currentstroke}{rgb}{0.478431,0.243137,0.615686}%
\pgfsetstrokecolor{currentstroke}%
\pgfsetdash{}{0pt}%
\pgfpathmoveto{\pgfqpoint{0.835551in}{1.147966in}}%
\pgfpathlineto{\pgfqpoint{0.870813in}{1.108269in}}%
\pgfpathlineto{\pgfqpoint{0.901885in}{1.074249in}}%
\pgfpathlineto{\pgfqpoint{0.929538in}{1.044913in}}%
\pgfpathlineto{\pgfqpoint{0.954375in}{1.019491in}}%
\pgfpathlineto{\pgfqpoint{0.976868in}{0.997376in}}%
\pgfpathlineto{\pgfqpoint{0.997395in}{0.978083in}}%
\pgfpathlineto{\pgfqpoint{1.016264in}{0.961221in}}%
\pgfpathlineto{\pgfqpoint{1.033725in}{0.946468in}}%
\pgfpathlineto{\pgfqpoint{1.049983in}{0.933561in}}%
\pgfpathlineto{\pgfqpoint{1.065212in}{0.922279in}}%
\pgfpathlineto{\pgfqpoint{1.079555in}{0.912438in}}%
\pgfpathlineto{\pgfqpoint{1.093134in}{0.903879in}}%
\pgfpathlineto{\pgfqpoint{1.106052in}{0.896470in}}%
\pgfpathlineto{\pgfqpoint{1.118395in}{0.890094in}}%
\pgfpathlineto{\pgfqpoint{1.130240in}{0.884652in}}%
\pgfpathlineto{\pgfqpoint{1.141649in}{0.880057in}}%
\pgfpathlineto{\pgfqpoint{1.152678in}{0.876233in}}%
\pgfpathlineto{\pgfqpoint{1.166872in}{0.872215in}}%
\pgfpathlineto{\pgfqpoint{1.180566in}{0.869306in}}%
\pgfpathlineto{\pgfqpoint{1.193835in}{0.867384in}}%
\pgfpathlineto{\pgfqpoint{1.206739in}{0.866342in}}%
\pgfpathlineto{\pgfqpoint{1.219327in}{0.866085in}}%
\pgfpathlineto{\pgfqpoint{1.231635in}{0.866529in}}%
\pgfpathlineto{\pgfqpoint{1.243692in}{0.867597in}}%
\pgfpathlineto{\pgfqpoint{1.258439in}{0.869702in}}%
\pgfpathlineto{\pgfqpoint{1.272850in}{0.872548in}}%
\pgfpathlineto{\pgfqpoint{1.286937in}{0.876022in}}%
\pgfpathlineto{\pgfqpoint{1.303416in}{0.880874in}}%
\pgfpathlineto{\pgfqpoint{1.319417in}{0.886317in}}%
\pgfpathlineto{\pgfqpoint{1.337451in}{0.893219in}}%
\pgfpathlineto{\pgfqpoint{1.357165in}{0.901581in}}%
\pgfpathlineto{\pgfqpoint{1.378102in}{0.911272in}}%
\pgfpathlineto{\pgfqpoint{1.401762in}{0.923080in}}%
\pgfpathlineto{\pgfqpoint{1.426813in}{0.936417in}}%
\pgfpathlineto{\pgfqpoint{1.454691in}{0.952099in}}%
\pgfpathlineto{\pgfqpoint{1.486132in}{0.970656in}}%
\pgfpathlineto{\pgfqpoint{1.521413in}{0.992369in}}%
\pgfpathlineto{\pgfqpoint{1.547118in}{1.008671in}}%
\pgfpathlineto{\pgfqpoint{1.547118in}{1.008671in}}%
\pgfusepath{stroke}%
\end{pgfscope}%
\begin{pgfscope}%
\pgfpathrectangle{\pgfqpoint{0.407592in}{0.222904in}}{\pgfqpoint{1.411163in}{1.411163in}}%
\pgfusepath{clip}%
\pgfsetrectcap%
\pgfsetroundjoin%
\pgfsetlinewidth{2.308625pt}%
\definecolor{currentstroke}{rgb}{0.878431,0.552941,0.176471}%
\pgfsetstrokecolor{currentstroke}%
\pgfsetdash{}{0pt}%
\pgfpathmoveto{\pgfqpoint{0.867618in}{1.272438in}}%
\pgfpathlineto{\pgfqpoint{0.911021in}{1.225022in}}%
\pgfpathlineto{\pgfqpoint{0.948898in}{1.185089in}}%
\pgfpathlineto{\pgfqpoint{0.982354in}{1.151214in}}%
\pgfpathlineto{\pgfqpoint{1.012228in}{1.122317in}}%
\pgfpathlineto{\pgfqpoint{1.039168in}{1.097564in}}%
\pgfpathlineto{\pgfqpoint{1.063680in}{1.076302in}}%
\pgfpathlineto{\pgfqpoint{1.086165in}{1.058008in}}%
\pgfpathlineto{\pgfqpoint{1.106947in}{1.042264in}}%
\pgfpathlineto{\pgfqpoint{1.126288in}{1.028726in}}%
\pgfpathlineto{\pgfqpoint{1.144400in}{1.017112in}}%
\pgfpathlineto{\pgfqpoint{1.161461in}{1.007185in}}%
\pgfpathlineto{\pgfqpoint{1.177617in}{0.998746in}}%
\pgfpathlineto{\pgfqpoint{1.192989in}{0.991626in}}%
\pgfpathlineto{\pgfqpoint{1.207677in}{0.985681in}}%
\pgfpathlineto{\pgfqpoint{1.221768in}{0.980784in}}%
\pgfpathlineto{\pgfqpoint{1.235331in}{0.976828in}}%
\pgfpathlineto{\pgfqpoint{1.248426in}{0.973714in}}%
\pgfpathlineto{\pgfqpoint{1.261103in}{0.971359in}}%
\pgfpathlineto{\pgfqpoint{1.273402in}{0.969687in}}%
\pgfpathlineto{\pgfqpoint{1.289272in}{0.968402in}}%
\pgfpathlineto{\pgfqpoint{1.304595in}{0.968066in}}%
\pgfpathlineto{\pgfqpoint{1.319420in}{0.968553in}}%
\pgfpathlineto{\pgfqpoint{1.333783in}{0.969750in}}%
\pgfpathlineto{\pgfqpoint{1.347708in}{0.971555in}}%
\pgfpathlineto{\pgfqpoint{1.364527in}{0.974530in}}%
\pgfpathlineto{\pgfqpoint{1.380708in}{0.978156in}}%
\pgfpathlineto{\pgfqpoint{1.399297in}{0.983169in}}%
\pgfpathlineto{\pgfqpoint{1.416982in}{0.988707in}}%
\pgfpathlineto{\pgfqpoint{1.436461in}{0.995588in}}%
\pgfpathlineto{\pgfqpoint{1.457181in}{1.003713in}}%
\pgfpathlineto{\pgfqpoint{1.478506in}{1.012845in}}%
\pgfpathlineto{\pgfqpoint{1.501727in}{1.023571in}}%
\pgfpathlineto{\pgfqpoint{1.526894in}{1.035997in}}%
\pgfpathlineto{\pgfqpoint{1.555052in}{1.050749in}}%
\pgfpathlineto{\pgfqpoint{1.586043in}{1.067868in}}%
\pgfpathlineto{\pgfqpoint{1.619859in}{1.087432in}}%
\pgfpathlineto{\pgfqpoint{1.619876in}{1.087442in}}%
\pgfpathlineto{\pgfqpoint{1.619876in}{1.087442in}}%
\pgfusepath{stroke}%
\end{pgfscope}%
\begin{pgfscope}%
\pgfpathrectangle{\pgfqpoint{0.407592in}{0.222904in}}{\pgfqpoint{1.411163in}{1.411163in}}%
\pgfusepath{clip}%
\pgfsetbuttcap%
\pgfsetroundjoin%
\definecolor{currentfill}{rgb}{0.878431,0.552941,0.176471}%
\pgfsetfillcolor{currentfill}%
\pgfsetlinewidth{1.003750pt}%
\definecolor{currentstroke}{rgb}{0.878431,0.552941,0.176471}%
\pgfsetstrokecolor{currentstroke}%
\pgfsetdash{}{0pt}%
\pgfsys@defobject{currentmarker}{\pgfqpoint{-0.034021in}{-0.034021in}}{\pgfqpoint{0.034021in}{0.034021in}}{%
\pgfpathmoveto{\pgfqpoint{0.000000in}{-0.034021in}}%
\pgfpathcurveto{\pgfqpoint{0.009022in}{-0.034021in}}{\pgfqpoint{0.017676in}{-0.030436in}}{\pgfqpoint{0.024056in}{-0.024056in}}%
\pgfpathcurveto{\pgfqpoint{0.030436in}{-0.017676in}}{\pgfqpoint{0.034021in}{-0.009022in}}{\pgfqpoint{0.034021in}{0.000000in}}%
\pgfpathcurveto{\pgfqpoint{0.034021in}{0.009022in}}{\pgfqpoint{0.030436in}{0.017676in}}{\pgfqpoint{0.024056in}{0.024056in}}%
\pgfpathcurveto{\pgfqpoint{0.017676in}{0.030436in}}{\pgfqpoint{0.009022in}{0.034021in}}{\pgfqpoint{0.000000in}{0.034021in}}%
\pgfpathcurveto{\pgfqpoint{-0.009022in}{0.034021in}}{\pgfqpoint{-0.017676in}{0.030436in}}{\pgfqpoint{-0.024056in}{0.024056in}}%
\pgfpathcurveto{\pgfqpoint{-0.030436in}{0.017676in}}{\pgfqpoint{-0.034021in}{0.009022in}}{\pgfqpoint{-0.034021in}{0.000000in}}%
\pgfpathcurveto{\pgfqpoint{-0.034021in}{-0.009022in}}{\pgfqpoint{-0.030436in}{-0.017676in}}{\pgfqpoint{-0.024056in}{-0.024056in}}%
\pgfpathcurveto{\pgfqpoint{-0.017676in}{-0.030436in}}{\pgfqpoint{-0.009022in}{-0.034021in}}{\pgfqpoint{0.000000in}{-0.034021in}}%
\pgfpathlineto{\pgfqpoint{0.000000in}{-0.034021in}}%
\pgfpathclose%
\pgfusepath{stroke,fill}%
}%
\begin{pgfscope}%
\pgfsys@transformshift{0.867618in}{1.272438in}%
\pgfsys@useobject{currentmarker}{}%
\end{pgfscope}%
\end{pgfscope}%
\begin{pgfscope}%
\definecolor{textcolor}{rgb}{0.349020,0.376471,0.403922}%
\pgfsetstrokecolor{textcolor}%
\pgfsetfillcolor{textcolor}%
\pgftext[x=0.450955in,y=1.358414in,left,base]{\color{textcolor}{\rmfamily\fontsize{6.000000}{7.200000}\selectfont\catcode`\^=\active\def^{\ifmmode\sp\else\^{}\fi}\catcode`\%=\active\def%{\%}cone $a^2\geq r^2+b^2$}}%
\end{pgfscope}%
\begin{pgfscope}%
\definecolor{textcolor}{rgb}{0.125490,0.498039,0.321569}%
\pgfsetstrokecolor{textcolor}%
\pgfsetfillcolor{textcolor}%
\pgftext[x=1.299191in,y=1.332664in,left,base]{\color{textcolor}{\rmfamily\fontsize{6.000000}{7.200000}\selectfont\catcode`\^=\active\def^{\ifmmode\sp\else\^{}\fi}\catcode`\%=\active\def%{\%}minima $r=\beta$}}%
\end{pgfscope}%
\begin{pgfscope}%
\definecolor{textcolor}{rgb}{0.000000,0.000000,0.000000}%
\pgfsetstrokecolor{textcolor}%
\pgfsetfillcolor{textcolor}%
\pgftext[x=1.113173in,y=1.660759in,,base]{\color{textcolor}{\rmfamily\fontsize{8.000000}{9.600000}\selectfont\catcode`\^=\active\def^{\ifmmode\sp\else\^{}\fi}\catcode`\%=\active\def%{\%}(b) Unequal variances}}%
\end{pgfscope}%
\begin{pgfscope}%
\pgfpathrectangle{\pgfqpoint{0.407592in}{0.222904in}}{\pgfqpoint{1.411163in}{1.411163in}}%
\pgfusepath{clip}%
\pgfsetbuttcap%
\pgfsetroundjoin%
\definecolor{currentfill}{rgb}{0.478431,0.243137,0.615686}%
\pgfsetfillcolor{currentfill}%
\pgfsetlinewidth{1.003750pt}%
\definecolor{currentstroke}{rgb}{0.478431,0.243137,0.615686}%
\pgfsetstrokecolor{currentstroke}%
\pgfsetdash{}{0pt}%
\pgfsys@defobject{currentmarker}{\pgfqpoint{-0.034021in}{-0.034021in}}{\pgfqpoint{0.034021in}{0.034021in}}{%
\pgfpathmoveto{\pgfqpoint{0.000000in}{-0.034021in}}%
\pgfpathcurveto{\pgfqpoint{0.009022in}{-0.034021in}}{\pgfqpoint{0.017676in}{-0.030436in}}{\pgfqpoint{0.024056in}{-0.024056in}}%
\pgfpathcurveto{\pgfqpoint{0.030436in}{-0.017676in}}{\pgfqpoint{0.034021in}{-0.009022in}}{\pgfqpoint{0.034021in}{0.000000in}}%
\pgfpathcurveto{\pgfqpoint{0.034021in}{0.009022in}}{\pgfqpoint{0.030436in}{0.017676in}}{\pgfqpoint{0.024056in}{0.024056in}}%
\pgfpathcurveto{\pgfqpoint{0.017676in}{0.030436in}}{\pgfqpoint{0.009022in}{0.034021in}}{\pgfqpoint{0.000000in}{0.034021in}}%
\pgfpathcurveto{\pgfqpoint{-0.009022in}{0.034021in}}{\pgfqpoint{-0.017676in}{0.030436in}}{\pgfqpoint{-0.024056in}{0.024056in}}%
\pgfpathcurveto{\pgfqpoint{-0.030436in}{0.017676in}}{\pgfqpoint{-0.034021in}{0.009022in}}{\pgfqpoint{-0.034021in}{0.000000in}}%
\pgfpathcurveto{\pgfqpoint{-0.034021in}{-0.009022in}}{\pgfqpoint{-0.030436in}{-0.017676in}}{\pgfqpoint{-0.024056in}{-0.024056in}}%
\pgfpathcurveto{\pgfqpoint{-0.017676in}{-0.030436in}}{\pgfqpoint{-0.009022in}{-0.034021in}}{\pgfqpoint{0.000000in}{-0.034021in}}%
\pgfpathlineto{\pgfqpoint{0.000000in}{-0.034021in}}%
\pgfpathclose%
\pgfusepath{stroke,fill}%
}%
\begin{pgfscope}%
\pgfsys@transformshift{0.835551in}{1.147966in}%
\pgfsys@useobject{currentmarker}{}%
\end{pgfscope}%
\end{pgfscope}%
\begin{pgfscope}%
\pgfpathrectangle{\pgfqpoint{0.407592in}{0.222904in}}{\pgfqpoint{1.411163in}{1.411163in}}%
\pgfusepath{clip}%
\pgfsetbuttcap%
\pgfsetroundjoin%
\definecolor{currentfill}{rgb}{0.000000,0.529412,0.423529}%
\pgfsetfillcolor{currentfill}%
\pgfsetlinewidth{1.003750pt}%
\definecolor{currentstroke}{rgb}{0.000000,0.529412,0.423529}%
\pgfsetstrokecolor{currentstroke}%
\pgfsetdash{}{0pt}%
\pgfsys@defobject{currentmarker}{\pgfqpoint{-0.034021in}{-0.034021in}}{\pgfqpoint{0.034021in}{0.034021in}}{%
\pgfpathmoveto{\pgfqpoint{0.000000in}{-0.034021in}}%
\pgfpathcurveto{\pgfqpoint{0.009022in}{-0.034021in}}{\pgfqpoint{0.017676in}{-0.030436in}}{\pgfqpoint{0.024056in}{-0.024056in}}%
\pgfpathcurveto{\pgfqpoint{0.030436in}{-0.017676in}}{\pgfqpoint{0.034021in}{-0.009022in}}{\pgfqpoint{0.034021in}{0.000000in}}%
\pgfpathcurveto{\pgfqpoint{0.034021in}{0.009022in}}{\pgfqpoint{0.030436in}{0.017676in}}{\pgfqpoint{0.024056in}{0.024056in}}%
\pgfpathcurveto{\pgfqpoint{0.017676in}{0.030436in}}{\pgfqpoint{0.009022in}{0.034021in}}{\pgfqpoint{0.000000in}{0.034021in}}%
\pgfpathcurveto{\pgfqpoint{-0.009022in}{0.034021in}}{\pgfqpoint{-0.017676in}{0.030436in}}{\pgfqpoint{-0.024056in}{0.024056in}}%
\pgfpathcurveto{\pgfqpoint{-0.030436in}{0.017676in}}{\pgfqpoint{-0.034021in}{0.009022in}}{\pgfqpoint{-0.034021in}{0.000000in}}%
\pgfpathcurveto{\pgfqpoint{-0.034021in}{-0.009022in}}{\pgfqpoint{-0.030436in}{-0.017676in}}{\pgfqpoint{-0.024056in}{-0.024056in}}%
\pgfpathcurveto{\pgfqpoint{-0.017676in}{-0.030436in}}{\pgfqpoint{-0.009022in}{-0.034021in}}{\pgfqpoint{0.000000in}{-0.034021in}}%
\pgfpathlineto{\pgfqpoint{0.000000in}{-0.034021in}}%
\pgfpathclose%
\pgfusepath{stroke,fill}%
}%
\begin{pgfscope}%
\pgfsys@transformshift{0.812524in}{1.057517in}%
\pgfsys@useobject{currentmarker}{}%
\end{pgfscope}%
\end{pgfscope}%
\begin{pgfscope}%
\pgfpathrectangle{\pgfqpoint{0.407592in}{0.222904in}}{\pgfqpoint{1.411163in}{1.411163in}}%
\pgfusepath{clip}%
\pgfsetbuttcap%
\pgfsetroundjoin%
\definecolor{currentfill}{rgb}{0.819608,0.286275,0.356863}%
\pgfsetfillcolor{currentfill}%
\pgfsetlinewidth{1.003750pt}%
\definecolor{currentstroke}{rgb}{0.819608,0.286275,0.356863}%
\pgfsetstrokecolor{currentstroke}%
\pgfsetdash{}{0pt}%
\pgfsys@defobject{currentmarker}{\pgfqpoint{-0.034021in}{-0.034021in}}{\pgfqpoint{0.034021in}{0.034021in}}{%
\pgfpathmoveto{\pgfqpoint{0.000000in}{-0.034021in}}%
\pgfpathcurveto{\pgfqpoint{0.009022in}{-0.034021in}}{\pgfqpoint{0.017676in}{-0.030436in}}{\pgfqpoint{0.024056in}{-0.024056in}}%
\pgfpathcurveto{\pgfqpoint{0.030436in}{-0.017676in}}{\pgfqpoint{0.034021in}{-0.009022in}}{\pgfqpoint{0.034021in}{0.000000in}}%
\pgfpathcurveto{\pgfqpoint{0.034021in}{0.009022in}}{\pgfqpoint{0.030436in}{0.017676in}}{\pgfqpoint{0.024056in}{0.024056in}}%
\pgfpathcurveto{\pgfqpoint{0.017676in}{0.030436in}}{\pgfqpoint{0.009022in}{0.034021in}}{\pgfqpoint{0.000000in}{0.034021in}}%
\pgfpathcurveto{\pgfqpoint{-0.009022in}{0.034021in}}{\pgfqpoint{-0.017676in}{0.030436in}}{\pgfqpoint{-0.024056in}{0.024056in}}%
\pgfpathcurveto{\pgfqpoint{-0.030436in}{0.017676in}}{\pgfqpoint{-0.034021in}{0.009022in}}{\pgfqpoint{-0.034021in}{0.000000in}}%
\pgfpathcurveto{\pgfqpoint{-0.034021in}{-0.009022in}}{\pgfqpoint{-0.030436in}{-0.017676in}}{\pgfqpoint{-0.024056in}{-0.024056in}}%
\pgfpathcurveto{\pgfqpoint{-0.017676in}{-0.030436in}}{\pgfqpoint{-0.009022in}{-0.034021in}}{\pgfqpoint{0.000000in}{-0.034021in}}%
\pgfpathlineto{\pgfqpoint{0.000000in}{-0.034021in}}%
\pgfpathclose%
\pgfusepath{stroke,fill}%
}%
\begin{pgfscope}%
\pgfsys@transformshift{0.672155in}{1.057312in}%
\pgfsys@useobject{currentmarker}{}%
\end{pgfscope}%
\end{pgfscope}%
\begin{pgfscope}%
\pgfpathrectangle{\pgfqpoint{0.407592in}{0.222904in}}{\pgfqpoint{1.411163in}{1.411163in}}%
\pgfusepath{clip}%
\pgfsetbuttcap%
\pgfsetroundjoin%
\definecolor{currentfill}{rgb}{0.878431,0.552941,0.176471}%
\pgfsetfillcolor{currentfill}%
\pgfsetlinewidth{1.003750pt}%
\definecolor{currentstroke}{rgb}{0.878431,0.552941,0.176471}%
\pgfsetstrokecolor{currentstroke}%
\pgfsetdash{}{0pt}%
\pgfsys@defobject{currentmarker}{\pgfqpoint{-0.038036in}{-0.038036in}}{\pgfqpoint{0.038036in}{0.038036in}}{%
\pgfpathmoveto{\pgfqpoint{-0.038036in}{-0.038036in}}%
\pgfpathlineto{\pgfqpoint{0.038036in}{-0.038036in}}%
\pgfpathlineto{\pgfqpoint{0.038036in}{0.038036in}}%
\pgfpathlineto{\pgfqpoint{-0.038036in}{0.038036in}}%
\pgfpathlineto{\pgfqpoint{-0.038036in}{-0.038036in}}%
\pgfpathclose%
\pgfusepath{stroke,fill}%
}%
\begin{pgfscope}%
\pgfsys@transformshift{1.619876in}{1.087442in}%
\pgfsys@useobject{currentmarker}{}%
\end{pgfscope}%
\end{pgfscope}%
\begin{pgfscope}%
\pgfpathrectangle{\pgfqpoint{0.407592in}{0.222904in}}{\pgfqpoint{1.411163in}{1.411163in}}%
\pgfusepath{clip}%
\pgfsetbuttcap%
\pgfsetroundjoin%
\definecolor{currentfill}{rgb}{0.141176,0.345098,0.650980}%
\pgfsetfillcolor{currentfill}%
\pgfsetlinewidth{1.003750pt}%
\definecolor{currentstroke}{rgb}{0.141176,0.345098,0.650980}%
\pgfsetstrokecolor{currentstroke}%
\pgfsetdash{}{0pt}%
\pgfsys@defobject{currentmarker}{\pgfqpoint{-0.034021in}{-0.034021in}}{\pgfqpoint{0.034021in}{0.034021in}}{%
\pgfpathmoveto{\pgfqpoint{0.000000in}{-0.034021in}}%
\pgfpathcurveto{\pgfqpoint{0.009022in}{-0.034021in}}{\pgfqpoint{0.017676in}{-0.030436in}}{\pgfqpoint{0.024056in}{-0.024056in}}%
\pgfpathcurveto{\pgfqpoint{0.030436in}{-0.017676in}}{\pgfqpoint{0.034021in}{-0.009022in}}{\pgfqpoint{0.034021in}{0.000000in}}%
\pgfpathcurveto{\pgfqpoint{0.034021in}{0.009022in}}{\pgfqpoint{0.030436in}{0.017676in}}{\pgfqpoint{0.024056in}{0.024056in}}%
\pgfpathcurveto{\pgfqpoint{0.017676in}{0.030436in}}{\pgfqpoint{0.009022in}{0.034021in}}{\pgfqpoint{0.000000in}{0.034021in}}%
\pgfpathcurveto{\pgfqpoint{-0.009022in}{0.034021in}}{\pgfqpoint{-0.017676in}{0.030436in}}{\pgfqpoint{-0.024056in}{0.024056in}}%
\pgfpathcurveto{\pgfqpoint{-0.030436in}{0.017676in}}{\pgfqpoint{-0.034021in}{0.009022in}}{\pgfqpoint{-0.034021in}{0.000000in}}%
\pgfpathcurveto{\pgfqpoint{-0.034021in}{-0.009022in}}{\pgfqpoint{-0.030436in}{-0.017676in}}{\pgfqpoint{-0.024056in}{-0.024056in}}%
\pgfpathcurveto{\pgfqpoint{-0.017676in}{-0.030436in}}{\pgfqpoint{-0.009022in}{-0.034021in}}{\pgfqpoint{0.000000in}{-0.034021in}}%
\pgfpathlineto{\pgfqpoint{0.000000in}{-0.034021in}}%
\pgfpathclose%
\pgfusepath{stroke,fill}%
}%
\begin{pgfscope}%
\pgfsys@transformshift{0.537673in}{1.095850in}%
\pgfsys@useobject{currentmarker}{}%
\end{pgfscope}%
\end{pgfscope}%
\begin{pgfscope}%
\pgfpathrectangle{\pgfqpoint{0.407592in}{0.222904in}}{\pgfqpoint{1.411163in}{1.411163in}}%
\pgfusepath{clip}%
\pgfsetbuttcap%
\pgfsetroundjoin%
\definecolor{currentfill}{rgb}{0.478431,0.243137,0.615686}%
\pgfsetfillcolor{currentfill}%
\pgfsetlinewidth{1.003750pt}%
\definecolor{currentstroke}{rgb}{0.478431,0.243137,0.615686}%
\pgfsetstrokecolor{currentstroke}%
\pgfsetdash{}{0pt}%
\pgfsys@defobject{currentmarker}{\pgfqpoint{-0.038036in}{-0.038036in}}{\pgfqpoint{0.038036in}{0.038036in}}{%
\pgfpathmoveto{\pgfqpoint{-0.038036in}{-0.038036in}}%
\pgfpathlineto{\pgfqpoint{0.038036in}{-0.038036in}}%
\pgfpathlineto{\pgfqpoint{0.038036in}{0.038036in}}%
\pgfpathlineto{\pgfqpoint{-0.038036in}{0.038036in}}%
\pgfpathlineto{\pgfqpoint{-0.038036in}{-0.038036in}}%
\pgfpathclose%
\pgfusepath{stroke,fill}%
}%
\begin{pgfscope}%
\pgfsys@transformshift{1.547118in}{1.008671in}%
\pgfsys@useobject{currentmarker}{}%
\end{pgfscope}%
\end{pgfscope}%
\begin{pgfscope}%
\pgfpathrectangle{\pgfqpoint{0.407592in}{0.222904in}}{\pgfqpoint{1.411163in}{1.411163in}}%
\pgfusepath{clip}%
\pgfsetbuttcap%
\pgfsetroundjoin%
\definecolor{currentfill}{rgb}{0.000000,0.529412,0.423529}%
\pgfsetfillcolor{currentfill}%
\pgfsetlinewidth{1.003750pt}%
\definecolor{currentstroke}{rgb}{0.000000,0.529412,0.423529}%
\pgfsetstrokecolor{currentstroke}%
\pgfsetdash{}{0pt}%
\pgfsys@defobject{currentmarker}{\pgfqpoint{-0.038036in}{-0.038036in}}{\pgfqpoint{0.038036in}{0.038036in}}{%
\pgfpathmoveto{\pgfqpoint{-0.038036in}{-0.038036in}}%
\pgfpathlineto{\pgfqpoint{0.038036in}{-0.038036in}}%
\pgfpathlineto{\pgfqpoint{0.038036in}{0.038036in}}%
\pgfpathlineto{\pgfqpoint{-0.038036in}{0.038036in}}%
\pgfpathlineto{\pgfqpoint{-0.038036in}{-0.038036in}}%
\pgfpathclose%
\pgfusepath{stroke,fill}%
}%
\begin{pgfscope}%
\pgfsys@transformshift{1.473472in}{0.975862in}%
\pgfsys@useobject{currentmarker}{}%
\end{pgfscope}%
\end{pgfscope}%
\begin{pgfscope}%
\pgfpathrectangle{\pgfqpoint{0.407592in}{0.222904in}}{\pgfqpoint{1.411163in}{1.411163in}}%
\pgfusepath{clip}%
\pgfsetbuttcap%
\pgfsetroundjoin%
\definecolor{currentfill}{rgb}{0.819608,0.286275,0.356863}%
\pgfsetfillcolor{currentfill}%
\pgfsetlinewidth{1.003750pt}%
\definecolor{currentstroke}{rgb}{0.819608,0.286275,0.356863}%
\pgfsetstrokecolor{currentstroke}%
\pgfsetdash{}{0pt}%
\pgfsys@defobject{currentmarker}{\pgfqpoint{-0.038036in}{-0.038036in}}{\pgfqpoint{0.038036in}{0.038036in}}{%
\pgfpathmoveto{\pgfqpoint{-0.038036in}{-0.038036in}}%
\pgfpathlineto{\pgfqpoint{0.038036in}{-0.038036in}}%
\pgfpathlineto{\pgfqpoint{0.038036in}{0.038036in}}%
\pgfpathlineto{\pgfqpoint{-0.038036in}{0.038036in}}%
\pgfpathlineto{\pgfqpoint{-0.038036in}{-0.038036in}}%
\pgfpathclose%
\pgfusepath{stroke,fill}%
}%
\begin{pgfscope}%
\pgfsys@transformshift{1.396682in}{0.908341in}%
\pgfsys@useobject{currentmarker}{}%
\end{pgfscope}%
\end{pgfscope}%
\begin{pgfscope}%
\pgfpathrectangle{\pgfqpoint{0.407592in}{0.222904in}}{\pgfqpoint{1.411163in}{1.411163in}}%
\pgfusepath{clip}%
\pgfsetbuttcap%
\pgfsetroundjoin%
\definecolor{currentfill}{rgb}{0.141176,0.345098,0.650980}%
\pgfsetfillcolor{currentfill}%
\pgfsetlinewidth{1.003750pt}%
\definecolor{currentstroke}{rgb}{0.141176,0.345098,0.650980}%
\pgfsetstrokecolor{currentstroke}%
\pgfsetdash{}{0pt}%
\pgfsys@defobject{currentmarker}{\pgfqpoint{-0.038036in}{-0.038036in}}{\pgfqpoint{0.038036in}{0.038036in}}{%
\pgfpathmoveto{\pgfqpoint{-0.038036in}{-0.038036in}}%
\pgfpathlineto{\pgfqpoint{0.038036in}{-0.038036in}}%
\pgfpathlineto{\pgfqpoint{0.038036in}{0.038036in}}%
\pgfpathlineto{\pgfqpoint{-0.038036in}{0.038036in}}%
\pgfpathlineto{\pgfqpoint{-0.038036in}{-0.038036in}}%
\pgfpathclose%
\pgfusepath{stroke,fill}%
}%
\begin{pgfscope}%
\pgfsys@transformshift{1.318296in}{0.888253in}%
\pgfsys@useobject{currentmarker}{}%
\end{pgfscope}%
\end{pgfscope}%
\begin{pgfscope}%
\pgfpathrectangle{\pgfqpoint{0.407592in}{0.222904in}}{\pgfqpoint{1.411163in}{1.411163in}}%
\pgfusepath{clip}%
\pgfsetbuttcap%
\pgfsetroundjoin%
\definecolor{currentfill}{rgb}{0.141176,0.345098,0.650980}%
\pgfsetfillcolor{currentfill}%
\pgfsetfillopacity{0.080000}%
\pgfsetlinewidth{0.000000pt}%
\definecolor{currentstroke}{rgb}{0.000000,0.000000,0.000000}%
\pgfsetstrokecolor{currentstroke}%
\pgfsetdash{}{0pt}%
\pgfpathmoveto{\pgfqpoint{0.554113in}{1.064739in}}%
\pgfusepath{fill}%
\end{pgfscope}%
\begin{pgfscope}%
\pgfpathrectangle{\pgfqpoint{0.407592in}{0.222904in}}{\pgfqpoint{1.411163in}{1.411163in}}%
\pgfusepath{clip}%
\pgfsetbuttcap%
\pgfsetroundjoin%
\definecolor{currentfill}{rgb}{0.141176,0.345098,0.650980}%
\pgfsetfillcolor{currentfill}%
\pgfsetfillopacity{0.080000}%
\pgfsetlinewidth{0.000000pt}%
\definecolor{currentstroke}{rgb}{0.000000,0.000000,0.000000}%
\pgfsetstrokecolor{currentstroke}%
\pgfsetdash{}{0pt}%
\pgfpathmoveto{\pgfqpoint{0.554113in}{1.064739in}}%
\pgfpathlineto{\pgfqpoint{0.552790in}{1.106160in}}%
\pgfusepath{fill}%
\end{pgfscope}%
\begin{pgfscope}%
\pgfpathrectangle{\pgfqpoint{0.407592in}{0.222904in}}{\pgfqpoint{1.411163in}{1.411163in}}%
\pgfusepath{clip}%
\pgfsetbuttcap%
\pgfsetroundjoin%
\definecolor{currentfill}{rgb}{0.141176,0.345098,0.650980}%
\pgfsetfillcolor{currentfill}%
\pgfsetfillopacity{0.080000}%
\pgfsetlinewidth{0.000000pt}%
\definecolor{currentstroke}{rgb}{0.000000,0.000000,0.000000}%
\pgfsetstrokecolor{currentstroke}%
\pgfsetdash{}{0pt}%
\pgfpathmoveto{\pgfqpoint{0.552790in}{1.106160in}}%
\pgfpathlineto{\pgfqpoint{0.551459in}{1.147771in}}%
\pgfpathlineto{\pgfqpoint{0.489239in}{1.162676in}}%
\pgfusepath{fill}%
\end{pgfscope}%
\begin{pgfscope}%
\pgfpathrectangle{\pgfqpoint{0.407592in}{0.222904in}}{\pgfqpoint{1.411163in}{1.411163in}}%
\pgfusepath{clip}%
\pgfsetbuttcap%
\pgfsetroundjoin%
\definecolor{currentfill}{rgb}{0.141176,0.345098,0.650980}%
\pgfsetfillcolor{currentfill}%
\pgfsetfillopacity{0.080000}%
\pgfsetlinewidth{0.000000pt}%
\definecolor{currentstroke}{rgb}{0.000000,0.000000,0.000000}%
\pgfsetstrokecolor{currentstroke}%
\pgfsetdash{}{0pt}%
\pgfpathmoveto{\pgfqpoint{0.617856in}{1.007946in}}%
\pgfpathlineto{\pgfqpoint{0.616675in}{1.049346in}}%
\pgfpathlineto{\pgfqpoint{0.554113in}{1.064739in}}%
\pgfusepath{fill}%
\end{pgfscope}%
\begin{pgfscope}%
\pgfpathrectangle{\pgfqpoint{0.407592in}{0.222904in}}{\pgfqpoint{1.411163in}{1.411163in}}%
\pgfusepath{clip}%
\pgfsetbuttcap%
\pgfsetroundjoin%
\definecolor{currentfill}{rgb}{0.141176,0.345098,0.650980}%
\pgfsetfillcolor{currentfill}%
\pgfsetfillopacity{0.080000}%
\pgfsetlinewidth{0.000000pt}%
\definecolor{currentstroke}{rgb}{0.000000,0.000000,0.000000}%
\pgfsetstrokecolor{currentstroke}%
\pgfsetdash{}{0pt}%
\pgfpathmoveto{\pgfqpoint{0.617856in}{1.007946in}}%
\pgfusepath{fill}%
\end{pgfscope}%
\begin{pgfscope}%
\pgfpathrectangle{\pgfqpoint{0.407592in}{0.222904in}}{\pgfqpoint{1.411163in}{1.411163in}}%
\pgfusepath{clip}%
\pgfsetbuttcap%
\pgfsetroundjoin%
\definecolor{currentfill}{rgb}{0.141176,0.345098,0.650980}%
\pgfsetfillcolor{currentfill}%
\pgfsetfillopacity{0.080000}%
\pgfsetlinewidth{0.000000pt}%
\definecolor{currentstroke}{rgb}{0.000000,0.000000,0.000000}%
\pgfsetstrokecolor{currentstroke}%
\pgfsetdash{}{0pt}%
\pgfpathmoveto{\pgfqpoint{0.489239in}{1.162676in}}%
\pgfpathlineto{\pgfqpoint{0.551459in}{1.147771in}}%
\pgfpathlineto{\pgfqpoint{0.550123in}{1.189574in}}%
\pgfpathlineto{\pgfqpoint{0.487767in}{1.204307in}}%
\pgfpathlineto{\pgfqpoint{0.489239in}{1.162676in}}%
\pgfpathclose%
\pgfusepath{fill}%
\end{pgfscope}%
\begin{pgfscope}%
\pgfpathrectangle{\pgfqpoint{0.407592in}{0.222904in}}{\pgfqpoint{1.411163in}{1.411163in}}%
\pgfusepath{clip}%
\pgfsetbuttcap%
\pgfsetroundjoin%
\definecolor{currentfill}{rgb}{0.141176,0.345098,0.650980}%
\pgfsetfillcolor{currentfill}%
\pgfsetfillopacity{0.080000}%
\pgfsetlinewidth{0.000000pt}%
\definecolor{currentstroke}{rgb}{0.000000,0.000000,0.000000}%
\pgfsetstrokecolor{currentstroke}%
\pgfsetdash{}{0pt}%
\pgfpathmoveto{\pgfqpoint{0.554113in}{1.064739in}}%
\pgfpathlineto{\pgfqpoint{0.616675in}{1.049346in}}%
\pgfpathlineto{\pgfqpoint{0.615489in}{1.090936in}}%
\pgfpathlineto{\pgfqpoint{0.552790in}{1.106160in}}%
\pgfpathlineto{\pgfqpoint{0.554113in}{1.064739in}}%
\pgfpathclose%
\pgfusepath{fill}%
\end{pgfscope}%
\begin{pgfscope}%
\pgfpathrectangle{\pgfqpoint{0.407592in}{0.222904in}}{\pgfqpoint{1.411163in}{1.411163in}}%
\pgfusepath{clip}%
\pgfsetbuttcap%
\pgfsetroundjoin%
\definecolor{currentfill}{rgb}{0.141176,0.345098,0.650980}%
\pgfsetfillcolor{currentfill}%
\pgfsetfillopacity{0.080000}%
\pgfsetlinewidth{0.000000pt}%
\definecolor{currentstroke}{rgb}{0.000000,0.000000,0.000000}%
\pgfsetstrokecolor{currentstroke}%
\pgfsetdash{}{0pt}%
\pgfpathmoveto{\pgfqpoint{0.681935in}{0.950853in}}%
\pgfpathlineto{\pgfqpoint{0.680899in}{0.992231in}}%
\pgfpathlineto{\pgfqpoint{0.617856in}{1.007946in}}%
\pgfusepath{fill}%
\end{pgfscope}%
\begin{pgfscope}%
\pgfpathrectangle{\pgfqpoint{0.407592in}{0.222904in}}{\pgfqpoint{1.411163in}{1.411163in}}%
\pgfusepath{clip}%
\pgfsetbuttcap%
\pgfsetroundjoin%
\definecolor{currentfill}{rgb}{0.141176,0.345098,0.650980}%
\pgfsetfillcolor{currentfill}%
\pgfsetfillopacity{0.080000}%
\pgfsetlinewidth{0.000000pt}%
\definecolor{currentstroke}{rgb}{0.000000,0.000000,0.000000}%
\pgfsetstrokecolor{currentstroke}%
\pgfsetdash{}{0pt}%
\pgfpathmoveto{\pgfqpoint{0.681935in}{0.950853in}}%
\pgfusepath{fill}%
\end{pgfscope}%
\begin{pgfscope}%
\pgfpathrectangle{\pgfqpoint{0.407592in}{0.222904in}}{\pgfqpoint{1.411163in}{1.411163in}}%
\pgfusepath{clip}%
\pgfsetbuttcap%
\pgfsetroundjoin%
\definecolor{currentfill}{rgb}{0.141176,0.345098,0.650980}%
\pgfsetfillcolor{currentfill}%
\pgfsetfillopacity{0.080000}%
\pgfsetlinewidth{0.000000pt}%
\definecolor{currentstroke}{rgb}{0.000000,0.000000,0.000000}%
\pgfsetstrokecolor{currentstroke}%
\pgfsetdash{}{0pt}%
\pgfpathmoveto{\pgfqpoint{0.487767in}{1.204307in}}%
\pgfpathlineto{\pgfqpoint{0.550123in}{1.189574in}}%
\pgfpathlineto{\pgfqpoint{0.548781in}{1.231569in}}%
\pgfpathlineto{\pgfqpoint{0.486288in}{1.246129in}}%
\pgfpathlineto{\pgfqpoint{0.487767in}{1.204307in}}%
\pgfpathclose%
\pgfusepath{fill}%
\end{pgfscope}%
\begin{pgfscope}%
\pgfpathrectangle{\pgfqpoint{0.407592in}{0.222904in}}{\pgfqpoint{1.411163in}{1.411163in}}%
\pgfusepath{clip}%
\pgfsetbuttcap%
\pgfsetroundjoin%
\definecolor{currentfill}{rgb}{0.141176,0.345098,0.650980}%
\pgfsetfillcolor{currentfill}%
\pgfsetfillopacity{0.080000}%
\pgfsetlinewidth{0.000000pt}%
\definecolor{currentstroke}{rgb}{0.000000,0.000000,0.000000}%
\pgfsetstrokecolor{currentstroke}%
\pgfsetdash{}{0pt}%
\pgfpathmoveto{\pgfqpoint{0.552790in}{1.106160in}}%
\pgfpathlineto{\pgfqpoint{0.615489in}{1.090936in}}%
\pgfpathlineto{\pgfqpoint{0.614297in}{1.132718in}}%
\pgfpathlineto{\pgfqpoint{0.551459in}{1.147771in}}%
\pgfpathlineto{\pgfqpoint{0.552790in}{1.106160in}}%
\pgfpathclose%
\pgfusepath{fill}%
\end{pgfscope}%
\begin{pgfscope}%
\pgfpathrectangle{\pgfqpoint{0.407592in}{0.222904in}}{\pgfqpoint{1.411163in}{1.411163in}}%
\pgfusepath{clip}%
\pgfsetbuttcap%
\pgfsetroundjoin%
\definecolor{currentfill}{rgb}{0.141176,0.345098,0.650980}%
\pgfsetfillcolor{currentfill}%
\pgfsetfillopacity{0.080000}%
\pgfsetlinewidth{0.000000pt}%
\definecolor{currentstroke}{rgb}{0.000000,0.000000,0.000000}%
\pgfsetstrokecolor{currentstroke}%
\pgfsetdash{}{0pt}%
\pgfpathmoveto{\pgfqpoint{0.617856in}{1.007946in}}%
\pgfpathlineto{\pgfqpoint{0.680899in}{0.992231in}}%
\pgfpathlineto{\pgfqpoint{0.679857in}{1.033800in}}%
\pgfpathlineto{\pgfqpoint{0.616675in}{1.049346in}}%
\pgfpathlineto{\pgfqpoint{0.617856in}{1.007946in}}%
\pgfpathclose%
\pgfusepath{fill}%
\end{pgfscope}%
\begin{pgfscope}%
\pgfpathrectangle{\pgfqpoint{0.407592in}{0.222904in}}{\pgfqpoint{1.411163in}{1.411163in}}%
\pgfusepath{clip}%
\pgfsetbuttcap%
\pgfsetroundjoin%
\definecolor{currentfill}{rgb}{0.141176,0.345098,0.650980}%
\pgfsetfillcolor{currentfill}%
\pgfsetfillopacity{0.080000}%
\pgfsetlinewidth{0.000000pt}%
\definecolor{currentstroke}{rgb}{0.000000,0.000000,0.000000}%
\pgfsetstrokecolor{currentstroke}%
\pgfsetdash{}{0pt}%
\pgfpathmoveto{\pgfqpoint{0.746353in}{0.893459in}}%
\pgfpathlineto{\pgfqpoint{0.745462in}{0.934814in}}%
\pgfpathlineto{\pgfqpoint{0.681935in}{0.950853in}}%
\pgfusepath{fill}%
\end{pgfscope}%
\begin{pgfscope}%
\pgfpathrectangle{\pgfqpoint{0.407592in}{0.222904in}}{\pgfqpoint{1.411163in}{1.411163in}}%
\pgfusepath{clip}%
\pgfsetbuttcap%
\pgfsetroundjoin%
\definecolor{currentfill}{rgb}{0.141176,0.345098,0.650980}%
\pgfsetfillcolor{currentfill}%
\pgfsetfillopacity{0.080000}%
\pgfsetlinewidth{0.000000pt}%
\definecolor{currentstroke}{rgb}{0.000000,0.000000,0.000000}%
\pgfsetstrokecolor{currentstroke}%
\pgfsetdash{}{0pt}%
\pgfpathmoveto{\pgfqpoint{0.486288in}{1.246129in}}%
\pgfpathlineto{\pgfqpoint{0.548781in}{1.231569in}}%
\pgfpathlineto{\pgfqpoint{0.547433in}{1.273759in}}%
\pgfpathlineto{\pgfqpoint{0.484803in}{1.288144in}}%
\pgfpathlineto{\pgfqpoint{0.486288in}{1.246129in}}%
\pgfpathclose%
\pgfusepath{fill}%
\end{pgfscope}%
\begin{pgfscope}%
\pgfpathrectangle{\pgfqpoint{0.407592in}{0.222904in}}{\pgfqpoint{1.411163in}{1.411163in}}%
\pgfusepath{clip}%
\pgfsetbuttcap%
\pgfsetroundjoin%
\definecolor{currentfill}{rgb}{0.141176,0.345098,0.650980}%
\pgfsetfillcolor{currentfill}%
\pgfsetfillopacity{0.080000}%
\pgfsetlinewidth{0.000000pt}%
\definecolor{currentstroke}{rgb}{0.000000,0.000000,0.000000}%
\pgfsetstrokecolor{currentstroke}%
\pgfsetdash{}{0pt}%
\pgfpathmoveto{\pgfqpoint{0.746353in}{0.893459in}}%
\pgfusepath{fill}%
\end{pgfscope}%
\begin{pgfscope}%
\pgfpathrectangle{\pgfqpoint{0.407592in}{0.222904in}}{\pgfqpoint{1.411163in}{1.411163in}}%
\pgfusepath{clip}%
\pgfsetbuttcap%
\pgfsetroundjoin%
\definecolor{currentfill}{rgb}{0.141176,0.345098,0.650980}%
\pgfsetfillcolor{currentfill}%
\pgfsetfillopacity{0.080000}%
\pgfsetlinewidth{0.000000pt}%
\definecolor{currentstroke}{rgb}{0.000000,0.000000,0.000000}%
\pgfsetstrokecolor{currentstroke}%
\pgfsetdash{}{0pt}%
\pgfpathmoveto{\pgfqpoint{0.551459in}{1.147771in}}%
\pgfpathlineto{\pgfqpoint{0.614297in}{1.132718in}}%
\pgfpathlineto{\pgfqpoint{0.613099in}{1.174694in}}%
\pgfpathlineto{\pgfqpoint{0.550123in}{1.189574in}}%
\pgfpathlineto{\pgfqpoint{0.551459in}{1.147771in}}%
\pgfpathclose%
\pgfusepath{fill}%
\end{pgfscope}%
\begin{pgfscope}%
\pgfpathrectangle{\pgfqpoint{0.407592in}{0.222904in}}{\pgfqpoint{1.411163in}{1.411163in}}%
\pgfusepath{clip}%
\pgfsetbuttcap%
\pgfsetroundjoin%
\definecolor{currentfill}{rgb}{0.141176,0.345098,0.650980}%
\pgfsetfillcolor{currentfill}%
\pgfsetfillopacity{0.080000}%
\pgfsetlinewidth{0.000000pt}%
\definecolor{currentstroke}{rgb}{0.000000,0.000000,0.000000}%
\pgfsetstrokecolor{currentstroke}%
\pgfsetdash{}{0pt}%
\pgfpathmoveto{\pgfqpoint{0.616675in}{1.049346in}}%
\pgfpathlineto{\pgfqpoint{0.679857in}{1.033800in}}%
\pgfpathlineto{\pgfqpoint{0.678811in}{1.075561in}}%
\pgfpathlineto{\pgfqpoint{0.615489in}{1.090936in}}%
\pgfpathlineto{\pgfqpoint{0.616675in}{1.049346in}}%
\pgfpathclose%
\pgfusepath{fill}%
\end{pgfscope}%
\begin{pgfscope}%
\pgfpathrectangle{\pgfqpoint{0.407592in}{0.222904in}}{\pgfqpoint{1.411163in}{1.411163in}}%
\pgfusepath{clip}%
\pgfsetbuttcap%
\pgfsetroundjoin%
\definecolor{currentfill}{rgb}{0.141176,0.345098,0.650980}%
\pgfsetfillcolor{currentfill}%
\pgfsetfillopacity{0.080000}%
\pgfsetlinewidth{0.000000pt}%
\definecolor{currentstroke}{rgb}{0.000000,0.000000,0.000000}%
\pgfsetstrokecolor{currentstroke}%
\pgfsetdash{}{0pt}%
\pgfpathmoveto{\pgfqpoint{0.681935in}{0.950853in}}%
\pgfpathlineto{\pgfqpoint{0.745462in}{0.934814in}}%
\pgfpathlineto{\pgfqpoint{0.744568in}{0.976360in}}%
\pgfpathlineto{\pgfqpoint{0.680899in}{0.992231in}}%
\pgfpathlineto{\pgfqpoint{0.681935in}{0.950853in}}%
\pgfpathclose%
\pgfusepath{fill}%
\end{pgfscope}%
\begin{pgfscope}%
\pgfpathrectangle{\pgfqpoint{0.407592in}{0.222904in}}{\pgfqpoint{1.411163in}{1.411163in}}%
\pgfusepath{clip}%
\pgfsetbuttcap%
\pgfsetroundjoin%
\definecolor{currentfill}{rgb}{0.141176,0.345098,0.650980}%
\pgfsetfillcolor{currentfill}%
\pgfsetfillopacity{0.080000}%
\pgfsetlinewidth{0.000000pt}%
\definecolor{currentstroke}{rgb}{0.000000,0.000000,0.000000}%
\pgfsetstrokecolor{currentstroke}%
\pgfsetdash{}{0pt}%
\pgfpathmoveto{\pgfqpoint{0.811113in}{0.835760in}}%
\pgfpathlineto{\pgfqpoint{0.810369in}{0.877091in}}%
\pgfpathlineto{\pgfqpoint{0.746353in}{0.893459in}}%
\pgfusepath{fill}%
\end{pgfscope}%
\begin{pgfscope}%
\pgfpathrectangle{\pgfqpoint{0.407592in}{0.222904in}}{\pgfqpoint{1.411163in}{1.411163in}}%
\pgfusepath{clip}%
\pgfsetbuttcap%
\pgfsetroundjoin%
\definecolor{currentfill}{rgb}{0.141176,0.345098,0.650980}%
\pgfsetfillcolor{currentfill}%
\pgfsetfillopacity{0.080000}%
\pgfsetlinewidth{0.000000pt}%
\definecolor{currentstroke}{rgb}{0.000000,0.000000,0.000000}%
\pgfsetstrokecolor{currentstroke}%
\pgfsetdash{}{0pt}%
\pgfpathmoveto{\pgfqpoint{0.550123in}{1.189574in}}%
\pgfpathlineto{\pgfqpoint{0.613099in}{1.174694in}}%
\pgfpathlineto{\pgfqpoint{0.611896in}{1.216864in}}%
\pgfpathlineto{\pgfqpoint{0.548781in}{1.231569in}}%
\pgfpathlineto{\pgfqpoint{0.550123in}{1.189574in}}%
\pgfpathclose%
\pgfusepath{fill}%
\end{pgfscope}%
\begin{pgfscope}%
\pgfpathrectangle{\pgfqpoint{0.407592in}{0.222904in}}{\pgfqpoint{1.411163in}{1.411163in}}%
\pgfusepath{clip}%
\pgfsetbuttcap%
\pgfsetroundjoin%
\definecolor{currentfill}{rgb}{0.141176,0.345098,0.650980}%
\pgfsetfillcolor{currentfill}%
\pgfsetfillopacity{0.080000}%
\pgfsetlinewidth{0.000000pt}%
\definecolor{currentstroke}{rgb}{0.000000,0.000000,0.000000}%
\pgfsetstrokecolor{currentstroke}%
\pgfsetdash{}{0pt}%
\pgfpathmoveto{\pgfqpoint{0.811113in}{0.835760in}}%
\pgfusepath{fill}%
\end{pgfscope}%
\begin{pgfscope}%
\pgfpathrectangle{\pgfqpoint{0.407592in}{0.222904in}}{\pgfqpoint{1.411163in}{1.411163in}}%
\pgfusepath{clip}%
\pgfsetbuttcap%
\pgfsetroundjoin%
\definecolor{currentfill}{rgb}{0.141176,0.345098,0.650980}%
\pgfsetfillcolor{currentfill}%
\pgfsetfillopacity{0.080000}%
\pgfsetlinewidth{0.000000pt}%
\definecolor{currentstroke}{rgb}{0.000000,0.000000,0.000000}%
\pgfsetstrokecolor{currentstroke}%
\pgfsetdash{}{0pt}%
\pgfpathmoveto{\pgfqpoint{0.615489in}{1.090936in}}%
\pgfpathlineto{\pgfqpoint{0.678811in}{1.075561in}}%
\pgfpathlineto{\pgfqpoint{0.677760in}{1.117516in}}%
\pgfpathlineto{\pgfqpoint{0.614297in}{1.132718in}}%
\pgfpathlineto{\pgfqpoint{0.615489in}{1.090936in}}%
\pgfpathclose%
\pgfusepath{fill}%
\end{pgfscope}%
\begin{pgfscope}%
\pgfpathrectangle{\pgfqpoint{0.407592in}{0.222904in}}{\pgfqpoint{1.411163in}{1.411163in}}%
\pgfusepath{clip}%
\pgfsetbuttcap%
\pgfsetroundjoin%
\definecolor{currentfill}{rgb}{0.141176,0.345098,0.650980}%
\pgfsetfillcolor{currentfill}%
\pgfsetfillopacity{0.080000}%
\pgfsetlinewidth{0.000000pt}%
\definecolor{currentstroke}{rgb}{0.000000,0.000000,0.000000}%
\pgfsetstrokecolor{currentstroke}%
\pgfsetdash{}{0pt}%
\pgfpathmoveto{\pgfqpoint{0.680899in}{0.992231in}}%
\pgfpathlineto{\pgfqpoint{0.744568in}{0.976360in}}%
\pgfpathlineto{\pgfqpoint{0.743669in}{1.018099in}}%
\pgfpathlineto{\pgfqpoint{0.679857in}{1.033800in}}%
\pgfpathlineto{\pgfqpoint{0.680899in}{0.992231in}}%
\pgfpathclose%
\pgfusepath{fill}%
\end{pgfscope}%
\begin{pgfscope}%
\pgfpathrectangle{\pgfqpoint{0.407592in}{0.222904in}}{\pgfqpoint{1.411163in}{1.411163in}}%
\pgfusepath{clip}%
\pgfsetbuttcap%
\pgfsetroundjoin%
\definecolor{currentfill}{rgb}{0.141176,0.345098,0.650980}%
\pgfsetfillcolor{currentfill}%
\pgfsetfillopacity{0.080000}%
\pgfsetlinewidth{0.000000pt}%
\definecolor{currentstroke}{rgb}{0.000000,0.000000,0.000000}%
\pgfsetstrokecolor{currentstroke}%
\pgfsetdash{}{0pt}%
\pgfpathmoveto{\pgfqpoint{0.746353in}{0.893459in}}%
\pgfpathlineto{\pgfqpoint{0.810369in}{0.877091in}}%
\pgfpathlineto{\pgfqpoint{0.809623in}{0.918615in}}%
\pgfpathlineto{\pgfqpoint{0.745462in}{0.934814in}}%
\pgfpathlineto{\pgfqpoint{0.746353in}{0.893459in}}%
\pgfpathclose%
\pgfusepath{fill}%
\end{pgfscope}%
\begin{pgfscope}%
\pgfpathrectangle{\pgfqpoint{0.407592in}{0.222904in}}{\pgfqpoint{1.411163in}{1.411163in}}%
\pgfusepath{clip}%
\pgfsetbuttcap%
\pgfsetroundjoin%
\definecolor{currentfill}{rgb}{0.141176,0.345098,0.650980}%
\pgfsetfillcolor{currentfill}%
\pgfsetfillopacity{0.080000}%
\pgfsetlinewidth{0.000000pt}%
\definecolor{currentstroke}{rgb}{0.000000,0.000000,0.000000}%
\pgfsetstrokecolor{currentstroke}%
\pgfsetdash{}{0pt}%
\pgfpathmoveto{\pgfqpoint{0.548781in}{1.231569in}}%
\pgfpathlineto{\pgfqpoint{0.611896in}{1.216864in}}%
\pgfpathlineto{\pgfqpoint{0.610687in}{1.259230in}}%
\pgfpathlineto{\pgfqpoint{0.547433in}{1.273759in}}%
\pgfpathlineto{\pgfqpoint{0.548781in}{1.231569in}}%
\pgfpathclose%
\pgfusepath{fill}%
\end{pgfscope}%
\begin{pgfscope}%
\pgfpathrectangle{\pgfqpoint{0.407592in}{0.222904in}}{\pgfqpoint{1.411163in}{1.411163in}}%
\pgfusepath{clip}%
\pgfsetbuttcap%
\pgfsetroundjoin%
\definecolor{currentfill}{rgb}{0.141176,0.345098,0.650980}%
\pgfsetfillcolor{currentfill}%
\pgfsetfillopacity{0.080000}%
\pgfsetlinewidth{0.000000pt}%
\definecolor{currentstroke}{rgb}{0.000000,0.000000,0.000000}%
\pgfsetstrokecolor{currentstroke}%
\pgfsetdash{}{0pt}%
\pgfpathmoveto{\pgfqpoint{0.876216in}{0.777754in}}%
\pgfpathlineto{\pgfqpoint{0.875622in}{0.819061in}}%
\pgfpathlineto{\pgfqpoint{0.811113in}{0.835760in}}%
\pgfusepath{fill}%
\end{pgfscope}%
\begin{pgfscope}%
\pgfpathrectangle{\pgfqpoint{0.407592in}{0.222904in}}{\pgfqpoint{1.411163in}{1.411163in}}%
\pgfusepath{clip}%
\pgfsetbuttcap%
\pgfsetroundjoin%
\definecolor{currentfill}{rgb}{0.141176,0.345098,0.650980}%
\pgfsetfillcolor{currentfill}%
\pgfsetfillopacity{0.080000}%
\pgfsetlinewidth{0.000000pt}%
\definecolor{currentstroke}{rgb}{0.000000,0.000000,0.000000}%
\pgfsetstrokecolor{currentstroke}%
\pgfsetdash{}{0pt}%
\pgfpathmoveto{\pgfqpoint{0.614297in}{1.132718in}}%
\pgfpathlineto{\pgfqpoint{0.677760in}{1.117516in}}%
\pgfpathlineto{\pgfqpoint{0.676704in}{1.159665in}}%
\pgfpathlineto{\pgfqpoint{0.613099in}{1.174694in}}%
\pgfpathlineto{\pgfqpoint{0.614297in}{1.132718in}}%
\pgfpathclose%
\pgfusepath{fill}%
\end{pgfscope}%
\begin{pgfscope}%
\pgfpathrectangle{\pgfqpoint{0.407592in}{0.222904in}}{\pgfqpoint{1.411163in}{1.411163in}}%
\pgfusepath{clip}%
\pgfsetbuttcap%
\pgfsetroundjoin%
\definecolor{currentfill}{rgb}{0.141176,0.345098,0.650980}%
\pgfsetfillcolor{currentfill}%
\pgfsetfillopacity{0.080000}%
\pgfsetlinewidth{0.000000pt}%
\definecolor{currentstroke}{rgb}{0.000000,0.000000,0.000000}%
\pgfsetstrokecolor{currentstroke}%
\pgfsetdash{}{0pt}%
\pgfpathmoveto{\pgfqpoint{0.876216in}{0.777754in}}%
\pgfusepath{fill}%
\end{pgfscope}%
\begin{pgfscope}%
\pgfpathrectangle{\pgfqpoint{0.407592in}{0.222904in}}{\pgfqpoint{1.411163in}{1.411163in}}%
\pgfusepath{clip}%
\pgfsetbuttcap%
\pgfsetroundjoin%
\definecolor{currentfill}{rgb}{0.141176,0.345098,0.650980}%
\pgfsetfillcolor{currentfill}%
\pgfsetfillopacity{0.080000}%
\pgfsetlinewidth{0.000000pt}%
\definecolor{currentstroke}{rgb}{0.000000,0.000000,0.000000}%
\pgfsetstrokecolor{currentstroke}%
\pgfsetdash{}{0pt}%
\pgfpathmoveto{\pgfqpoint{0.679857in}{1.033800in}}%
\pgfpathlineto{\pgfqpoint{0.743669in}{1.018099in}}%
\pgfpathlineto{\pgfqpoint{0.742765in}{1.060032in}}%
\pgfpathlineto{\pgfqpoint{0.678811in}{1.075561in}}%
\pgfpathlineto{\pgfqpoint{0.679857in}{1.033800in}}%
\pgfpathclose%
\pgfusepath{fill}%
\end{pgfscope}%
\begin{pgfscope}%
\pgfpathrectangle{\pgfqpoint{0.407592in}{0.222904in}}{\pgfqpoint{1.411163in}{1.411163in}}%
\pgfusepath{clip}%
\pgfsetbuttcap%
\pgfsetroundjoin%
\definecolor{currentfill}{rgb}{0.141176,0.345098,0.650980}%
\pgfsetfillcolor{currentfill}%
\pgfsetfillopacity{0.080000}%
\pgfsetlinewidth{0.000000pt}%
\definecolor{currentstroke}{rgb}{0.000000,0.000000,0.000000}%
\pgfsetstrokecolor{currentstroke}%
\pgfsetdash{}{0pt}%
\pgfpathmoveto{\pgfqpoint{0.745462in}{0.934814in}}%
\pgfpathlineto{\pgfqpoint{0.809623in}{0.918615in}}%
\pgfpathlineto{\pgfqpoint{0.808873in}{0.960331in}}%
\pgfpathlineto{\pgfqpoint{0.744568in}{0.976360in}}%
\pgfpathlineto{\pgfqpoint{0.745462in}{0.934814in}}%
\pgfpathclose%
\pgfusepath{fill}%
\end{pgfscope}%
\begin{pgfscope}%
\pgfpathrectangle{\pgfqpoint{0.407592in}{0.222904in}}{\pgfqpoint{1.411163in}{1.411163in}}%
\pgfusepath{clip}%
\pgfsetbuttcap%
\pgfsetroundjoin%
\definecolor{currentfill}{rgb}{0.141176,0.345098,0.650980}%
\pgfsetfillcolor{currentfill}%
\pgfsetfillopacity{0.080000}%
\pgfsetlinewidth{0.000000pt}%
\definecolor{currentstroke}{rgb}{0.000000,0.000000,0.000000}%
\pgfsetstrokecolor{currentstroke}%
\pgfsetdash{}{0pt}%
\pgfpathmoveto{\pgfqpoint{0.811113in}{0.835760in}}%
\pgfpathlineto{\pgfqpoint{0.875622in}{0.819061in}}%
\pgfpathlineto{\pgfqpoint{0.875025in}{0.860560in}}%
\pgfpathlineto{\pgfqpoint{0.810369in}{0.877091in}}%
\pgfpathlineto{\pgfqpoint{0.811113in}{0.835760in}}%
\pgfpathclose%
\pgfusepath{fill}%
\end{pgfscope}%
\begin{pgfscope}%
\pgfpathrectangle{\pgfqpoint{0.407592in}{0.222904in}}{\pgfqpoint{1.411163in}{1.411163in}}%
\pgfusepath{clip}%
\pgfsetbuttcap%
\pgfsetroundjoin%
\definecolor{currentfill}{rgb}{0.141176,0.345098,0.650980}%
\pgfsetfillcolor{currentfill}%
\pgfsetfillopacity{0.080000}%
\pgfsetlinewidth{0.000000pt}%
\definecolor{currentstroke}{rgb}{0.000000,0.000000,0.000000}%
\pgfsetstrokecolor{currentstroke}%
\pgfsetdash{}{0pt}%
\pgfpathmoveto{\pgfqpoint{0.613099in}{1.174694in}}%
\pgfpathlineto{\pgfqpoint{0.676704in}{1.159665in}}%
\pgfpathlineto{\pgfqpoint{0.675643in}{1.202012in}}%
\pgfpathlineto{\pgfqpoint{0.611896in}{1.216864in}}%
\pgfpathlineto{\pgfqpoint{0.613099in}{1.174694in}}%
\pgfpathclose%
\pgfusepath{fill}%
\end{pgfscope}%
\begin{pgfscope}%
\pgfpathrectangle{\pgfqpoint{0.407592in}{0.222904in}}{\pgfqpoint{1.411163in}{1.411163in}}%
\pgfusepath{clip}%
\pgfsetbuttcap%
\pgfsetroundjoin%
\definecolor{currentfill}{rgb}{0.141176,0.345098,0.650980}%
\pgfsetfillcolor{currentfill}%
\pgfsetfillopacity{0.080000}%
\pgfsetlinewidth{0.000000pt}%
\definecolor{currentstroke}{rgb}{0.000000,0.000000,0.000000}%
\pgfsetstrokecolor{currentstroke}%
\pgfsetdash{}{0pt}%
\pgfpathmoveto{\pgfqpoint{0.941224in}{0.760721in}}%
\pgfpathlineto{\pgfqpoint{0.876216in}{0.777754in}}%
\pgfusepath{fill}%
\end{pgfscope}%
\begin{pgfscope}%
\pgfpathrectangle{\pgfqpoint{0.407592in}{0.222904in}}{\pgfqpoint{1.411163in}{1.411163in}}%
\pgfusepath{clip}%
\pgfsetbuttcap%
\pgfsetroundjoin%
\definecolor{currentfill}{rgb}{0.141176,0.345098,0.650980}%
\pgfsetfillcolor{currentfill}%
\pgfsetfillopacity{0.080000}%
\pgfsetlinewidth{0.000000pt}%
\definecolor{currentstroke}{rgb}{0.000000,0.000000,0.000000}%
\pgfsetstrokecolor{currentstroke}%
\pgfsetdash{}{0pt}%
\pgfpathmoveto{\pgfqpoint{0.678811in}{1.075561in}}%
\pgfpathlineto{\pgfqpoint{0.742765in}{1.060032in}}%
\pgfpathlineto{\pgfqpoint{0.741858in}{1.102161in}}%
\pgfpathlineto{\pgfqpoint{0.677760in}{1.117516in}}%
\pgfpathlineto{\pgfqpoint{0.678811in}{1.075561in}}%
\pgfpathclose%
\pgfusepath{fill}%
\end{pgfscope}%
\begin{pgfscope}%
\pgfpathrectangle{\pgfqpoint{0.407592in}{0.222904in}}{\pgfqpoint{1.411163in}{1.411163in}}%
\pgfusepath{clip}%
\pgfsetbuttcap%
\pgfsetroundjoin%
\definecolor{currentfill}{rgb}{0.141176,0.345098,0.650980}%
\pgfsetfillcolor{currentfill}%
\pgfsetfillopacity{0.080000}%
\pgfsetlinewidth{0.000000pt}%
\definecolor{currentstroke}{rgb}{0.000000,0.000000,0.000000}%
\pgfsetstrokecolor{currentstroke}%
\pgfsetdash{}{0pt}%
\pgfpathmoveto{\pgfqpoint{0.744568in}{0.976360in}}%
\pgfpathlineto{\pgfqpoint{0.808873in}{0.960331in}}%
\pgfpathlineto{\pgfqpoint{0.808119in}{1.002242in}}%
\pgfpathlineto{\pgfqpoint{0.743669in}{1.018099in}}%
\pgfpathlineto{\pgfqpoint{0.744568in}{0.976360in}}%
\pgfpathclose%
\pgfusepath{fill}%
\end{pgfscope}%
\begin{pgfscope}%
\pgfpathrectangle{\pgfqpoint{0.407592in}{0.222904in}}{\pgfqpoint{1.411163in}{1.411163in}}%
\pgfusepath{clip}%
\pgfsetbuttcap%
\pgfsetroundjoin%
\definecolor{currentfill}{rgb}{0.141176,0.345098,0.650980}%
\pgfsetfillcolor{currentfill}%
\pgfsetfillopacity{0.080000}%
\pgfsetlinewidth{0.000000pt}%
\definecolor{currentstroke}{rgb}{0.000000,0.000000,0.000000}%
\pgfsetstrokecolor{currentstroke}%
\pgfsetdash{}{0pt}%
\pgfpathmoveto{\pgfqpoint{0.810369in}{0.877091in}}%
\pgfpathlineto{\pgfqpoint{0.875025in}{0.860560in}}%
\pgfpathlineto{\pgfqpoint{0.874426in}{0.902253in}}%
\pgfpathlineto{\pgfqpoint{0.809623in}{0.918615in}}%
\pgfpathlineto{\pgfqpoint{0.810369in}{0.877091in}}%
\pgfpathclose%
\pgfusepath{fill}%
\end{pgfscope}%
\begin{pgfscope}%
\pgfpathrectangle{\pgfqpoint{0.407592in}{0.222904in}}{\pgfqpoint{1.411163in}{1.411163in}}%
\pgfusepath{clip}%
\pgfsetbuttcap%
\pgfsetroundjoin%
\definecolor{currentfill}{rgb}{0.141176,0.345098,0.650980}%
\pgfsetfillcolor{currentfill}%
\pgfsetfillopacity{0.080000}%
\pgfsetlinewidth{0.000000pt}%
\definecolor{currentstroke}{rgb}{0.000000,0.000000,0.000000}%
\pgfsetstrokecolor{currentstroke}%
\pgfsetdash{}{0pt}%
\pgfpathmoveto{\pgfqpoint{0.876216in}{0.777754in}}%
\pgfpathlineto{\pgfqpoint{0.941224in}{0.760721in}}%
\pgfpathlineto{\pgfqpoint{0.940778in}{0.802195in}}%
\pgfpathlineto{\pgfqpoint{0.875622in}{0.819061in}}%
\pgfpathlineto{\pgfqpoint{0.876216in}{0.777754in}}%
\pgfpathclose%
\pgfusepath{fill}%
\end{pgfscope}%
\begin{pgfscope}%
\pgfpathrectangle{\pgfqpoint{0.407592in}{0.222904in}}{\pgfqpoint{1.411163in}{1.411163in}}%
\pgfusepath{clip}%
\pgfsetbuttcap%
\pgfsetroundjoin%
\definecolor{currentfill}{rgb}{0.141176,0.345098,0.650980}%
\pgfsetfillcolor{currentfill}%
\pgfsetfillopacity{0.080000}%
\pgfsetlinewidth{0.000000pt}%
\definecolor{currentstroke}{rgb}{0.000000,0.000000,0.000000}%
\pgfsetstrokecolor{currentstroke}%
\pgfsetdash{}{0pt}%
\pgfpathmoveto{\pgfqpoint{0.611896in}{1.216864in}}%
\pgfpathlineto{\pgfqpoint{0.675643in}{1.202012in}}%
\pgfpathlineto{\pgfqpoint{0.674577in}{1.244555in}}%
\pgfpathlineto{\pgfqpoint{0.610687in}{1.259230in}}%
\pgfpathlineto{\pgfqpoint{0.611896in}{1.216864in}}%
\pgfpathclose%
\pgfusepath{fill}%
\end{pgfscope}%
\begin{pgfscope}%
\pgfpathrectangle{\pgfqpoint{0.407592in}{0.222904in}}{\pgfqpoint{1.411163in}{1.411163in}}%
\pgfusepath{clip}%
\pgfsetbuttcap%
\pgfsetroundjoin%
\definecolor{currentfill}{rgb}{0.141176,0.345098,0.650980}%
\pgfsetfillcolor{currentfill}%
\pgfsetfillopacity{0.080000}%
\pgfsetlinewidth{0.000000pt}%
\definecolor{currentstroke}{rgb}{0.000000,0.000000,0.000000}%
\pgfsetstrokecolor{currentstroke}%
\pgfsetdash{}{0pt}%
\pgfpathmoveto{\pgfqpoint{0.677760in}{1.117516in}}%
\pgfpathlineto{\pgfqpoint{0.741858in}{1.102161in}}%
\pgfpathlineto{\pgfqpoint{0.740946in}{1.144486in}}%
\pgfpathlineto{\pgfqpoint{0.676704in}{1.159665in}}%
\pgfpathlineto{\pgfqpoint{0.677760in}{1.117516in}}%
\pgfpathclose%
\pgfusepath{fill}%
\end{pgfscope}%
\begin{pgfscope}%
\pgfpathrectangle{\pgfqpoint{0.407592in}{0.222904in}}{\pgfqpoint{1.411163in}{1.411163in}}%
\pgfusepath{clip}%
\pgfsetbuttcap%
\pgfsetroundjoin%
\definecolor{currentfill}{rgb}{0.141176,0.345098,0.650980}%
\pgfsetfillcolor{currentfill}%
\pgfsetfillopacity{0.080000}%
\pgfsetlinewidth{0.000000pt}%
\definecolor{currentstroke}{rgb}{0.000000,0.000000,0.000000}%
\pgfsetstrokecolor{currentstroke}%
\pgfsetdash{}{0pt}%
\pgfpathmoveto{\pgfqpoint{0.743669in}{1.018099in}}%
\pgfpathlineto{\pgfqpoint{0.808119in}{1.002242in}}%
\pgfpathlineto{\pgfqpoint{0.807362in}{1.044348in}}%
\pgfpathlineto{\pgfqpoint{0.742765in}{1.060032in}}%
\pgfpathlineto{\pgfqpoint{0.743669in}{1.018099in}}%
\pgfpathclose%
\pgfusepath{fill}%
\end{pgfscope}%
\begin{pgfscope}%
\pgfpathrectangle{\pgfqpoint{0.407592in}{0.222904in}}{\pgfqpoint{1.411163in}{1.411163in}}%
\pgfusepath{clip}%
\pgfsetbuttcap%
\pgfsetroundjoin%
\definecolor{currentfill}{rgb}{0.141176,0.345098,0.650980}%
\pgfsetfillcolor{currentfill}%
\pgfsetfillopacity{0.080000}%
\pgfsetlinewidth{0.000000pt}%
\definecolor{currentstroke}{rgb}{0.000000,0.000000,0.000000}%
\pgfsetstrokecolor{currentstroke}%
\pgfsetdash{}{0pt}%
\pgfpathmoveto{\pgfqpoint{0.809623in}{0.918615in}}%
\pgfpathlineto{\pgfqpoint{0.874426in}{0.902253in}}%
\pgfpathlineto{\pgfqpoint{0.873823in}{0.944141in}}%
\pgfpathlineto{\pgfqpoint{0.808873in}{0.960331in}}%
\pgfpathlineto{\pgfqpoint{0.809623in}{0.918615in}}%
\pgfpathclose%
\pgfusepath{fill}%
\end{pgfscope}%
\begin{pgfscope}%
\pgfpathrectangle{\pgfqpoint{0.407592in}{0.222904in}}{\pgfqpoint{1.411163in}{1.411163in}}%
\pgfusepath{clip}%
\pgfsetbuttcap%
\pgfsetroundjoin%
\definecolor{currentfill}{rgb}{0.141176,0.345098,0.650980}%
\pgfsetfillcolor{currentfill}%
\pgfsetfillopacity{0.080000}%
\pgfsetlinewidth{0.000000pt}%
\definecolor{currentstroke}{rgb}{0.000000,0.000000,0.000000}%
\pgfsetstrokecolor{currentstroke}%
\pgfsetdash{}{0pt}%
\pgfpathmoveto{\pgfqpoint{0.875622in}{0.819061in}}%
\pgfpathlineto{\pgfqpoint{0.940778in}{0.802195in}}%
\pgfpathlineto{\pgfqpoint{0.940331in}{0.843863in}}%
\pgfpathlineto{\pgfqpoint{0.875025in}{0.860560in}}%
\pgfpathlineto{\pgfqpoint{0.875622in}{0.819061in}}%
\pgfpathclose%
\pgfusepath{fill}%
\end{pgfscope}%
\begin{pgfscope}%
\pgfpathrectangle{\pgfqpoint{0.407592in}{0.222904in}}{\pgfqpoint{1.411163in}{1.411163in}}%
\pgfusepath{clip}%
\pgfsetbuttcap%
\pgfsetroundjoin%
\definecolor{currentfill}{rgb}{0.141176,0.345098,0.650980}%
\pgfsetfillcolor{currentfill}%
\pgfsetfillopacity{0.080000}%
\pgfsetlinewidth{0.000000pt}%
\definecolor{currentstroke}{rgb}{0.000000,0.000000,0.000000}%
\pgfsetstrokecolor{currentstroke}%
\pgfsetdash{}{0pt}%
\pgfpathmoveto{\pgfqpoint{0.676704in}{1.159665in}}%
\pgfpathlineto{\pgfqpoint{0.740946in}{1.144486in}}%
\pgfpathlineto{\pgfqpoint{0.740030in}{1.187010in}}%
\pgfpathlineto{\pgfqpoint{0.675643in}{1.202012in}}%
\pgfpathlineto{\pgfqpoint{0.676704in}{1.159665in}}%
\pgfpathclose%
\pgfusepath{fill}%
\end{pgfscope}%
\begin{pgfscope}%
\pgfpathrectangle{\pgfqpoint{0.407592in}{0.222904in}}{\pgfqpoint{1.411163in}{1.411163in}}%
\pgfusepath{clip}%
\pgfsetbuttcap%
\pgfsetroundjoin%
\definecolor{currentfill}{rgb}{0.141176,0.345098,0.650980}%
\pgfsetfillcolor{currentfill}%
\pgfsetfillopacity{0.080000}%
\pgfsetlinewidth{0.000000pt}%
\definecolor{currentstroke}{rgb}{0.000000,0.000000,0.000000}%
\pgfsetstrokecolor{currentstroke}%
\pgfsetdash{}{0pt}%
\pgfpathmoveto{\pgfqpoint{1.006884in}{0.743517in}}%
\pgfpathlineto{\pgfqpoint{0.941224in}{0.760721in}}%
\pgfusepath{fill}%
\end{pgfscope}%
\begin{pgfscope}%
\pgfpathrectangle{\pgfqpoint{0.407592in}{0.222904in}}{\pgfqpoint{1.411163in}{1.411163in}}%
\pgfusepath{clip}%
\pgfsetbuttcap%
\pgfsetroundjoin%
\definecolor{currentfill}{rgb}{0.141176,0.345098,0.650980}%
\pgfsetfillcolor{currentfill}%
\pgfsetfillopacity{0.080000}%
\pgfsetlinewidth{0.000000pt}%
\definecolor{currentstroke}{rgb}{0.000000,0.000000,0.000000}%
\pgfsetstrokecolor{currentstroke}%
\pgfsetdash{}{0pt}%
\pgfpathmoveto{\pgfqpoint{0.742765in}{1.060032in}}%
\pgfpathlineto{\pgfqpoint{0.807362in}{1.044348in}}%
\pgfpathlineto{\pgfqpoint{0.806601in}{1.086652in}}%
\pgfpathlineto{\pgfqpoint{0.741858in}{1.102161in}}%
\pgfpathlineto{\pgfqpoint{0.742765in}{1.060032in}}%
\pgfpathclose%
\pgfusepath{fill}%
\end{pgfscope}%
\begin{pgfscope}%
\pgfpathrectangle{\pgfqpoint{0.407592in}{0.222904in}}{\pgfqpoint{1.411163in}{1.411163in}}%
\pgfusepath{clip}%
\pgfsetbuttcap%
\pgfsetroundjoin%
\definecolor{currentfill}{rgb}{0.141176,0.345098,0.650980}%
\pgfsetfillcolor{currentfill}%
\pgfsetfillopacity{0.080000}%
\pgfsetlinewidth{0.000000pt}%
\definecolor{currentstroke}{rgb}{0.000000,0.000000,0.000000}%
\pgfsetstrokecolor{currentstroke}%
\pgfsetdash{}{0pt}%
\pgfpathmoveto{\pgfqpoint{0.808873in}{0.960331in}}%
\pgfpathlineto{\pgfqpoint{0.873823in}{0.944141in}}%
\pgfpathlineto{\pgfqpoint{0.873218in}{0.986224in}}%
\pgfpathlineto{\pgfqpoint{0.808119in}{1.002242in}}%
\pgfpathlineto{\pgfqpoint{0.808873in}{0.960331in}}%
\pgfpathclose%
\pgfusepath{fill}%
\end{pgfscope}%
\begin{pgfscope}%
\pgfpathrectangle{\pgfqpoint{0.407592in}{0.222904in}}{\pgfqpoint{1.411163in}{1.411163in}}%
\pgfusepath{clip}%
\pgfsetbuttcap%
\pgfsetroundjoin%
\definecolor{currentfill}{rgb}{0.141176,0.345098,0.650980}%
\pgfsetfillcolor{currentfill}%
\pgfsetfillopacity{0.080000}%
\pgfsetlinewidth{0.000000pt}%
\definecolor{currentstroke}{rgb}{0.000000,0.000000,0.000000}%
\pgfsetstrokecolor{currentstroke}%
\pgfsetdash{}{0pt}%
\pgfpathmoveto{\pgfqpoint{0.875025in}{0.860560in}}%
\pgfpathlineto{\pgfqpoint{0.940331in}{0.843863in}}%
\pgfpathlineto{\pgfqpoint{0.939881in}{0.885727in}}%
\pgfpathlineto{\pgfqpoint{0.874426in}{0.902253in}}%
\pgfpathlineto{\pgfqpoint{0.875025in}{0.860560in}}%
\pgfpathclose%
\pgfusepath{fill}%
\end{pgfscope}%
\begin{pgfscope}%
\pgfpathrectangle{\pgfqpoint{0.407592in}{0.222904in}}{\pgfqpoint{1.411163in}{1.411163in}}%
\pgfusepath{clip}%
\pgfsetbuttcap%
\pgfsetroundjoin%
\definecolor{currentfill}{rgb}{0.141176,0.345098,0.650980}%
\pgfsetfillcolor{currentfill}%
\pgfsetfillopacity{0.080000}%
\pgfsetlinewidth{0.000000pt}%
\definecolor{currentstroke}{rgb}{0.000000,0.000000,0.000000}%
\pgfsetstrokecolor{currentstroke}%
\pgfsetdash{}{0pt}%
\pgfpathmoveto{\pgfqpoint{0.941224in}{0.760721in}}%
\pgfpathlineto{\pgfqpoint{1.006884in}{0.743517in}}%
\pgfpathlineto{\pgfqpoint{1.006591in}{0.785160in}}%
\pgfpathlineto{\pgfqpoint{0.940778in}{0.802195in}}%
\pgfpathlineto{\pgfqpoint{0.941224in}{0.760721in}}%
\pgfpathclose%
\pgfusepath{fill}%
\end{pgfscope}%
\begin{pgfscope}%
\pgfpathrectangle{\pgfqpoint{0.407592in}{0.222904in}}{\pgfqpoint{1.411163in}{1.411163in}}%
\pgfusepath{clip}%
\pgfsetbuttcap%
\pgfsetroundjoin%
\definecolor{currentfill}{rgb}{0.141176,0.345098,0.650980}%
\pgfsetfillcolor{currentfill}%
\pgfsetfillopacity{0.080000}%
\pgfsetlinewidth{0.000000pt}%
\definecolor{currentstroke}{rgb}{0.000000,0.000000,0.000000}%
\pgfsetstrokecolor{currentstroke}%
\pgfsetdash{}{0pt}%
\pgfpathmoveto{\pgfqpoint{0.675643in}{1.202012in}}%
\pgfpathlineto{\pgfqpoint{0.740030in}{1.187010in}}%
\pgfpathlineto{\pgfqpoint{0.739110in}{1.229733in}}%
\pgfpathlineto{\pgfqpoint{0.674577in}{1.244555in}}%
\pgfpathlineto{\pgfqpoint{0.675643in}{1.202012in}}%
\pgfpathclose%
\pgfusepath{fill}%
\end{pgfscope}%
\begin{pgfscope}%
\pgfpathrectangle{\pgfqpoint{0.407592in}{0.222904in}}{\pgfqpoint{1.411163in}{1.411163in}}%
\pgfusepath{clip}%
\pgfsetbuttcap%
\pgfsetroundjoin%
\definecolor{currentfill}{rgb}{0.141176,0.345098,0.650980}%
\pgfsetfillcolor{currentfill}%
\pgfsetfillopacity{0.080000}%
\pgfsetlinewidth{0.000000pt}%
\definecolor{currentstroke}{rgb}{0.000000,0.000000,0.000000}%
\pgfsetstrokecolor{currentstroke}%
\pgfsetdash{}{0pt}%
\pgfpathmoveto{\pgfqpoint{0.741858in}{1.102161in}}%
\pgfpathlineto{\pgfqpoint{0.806601in}{1.086652in}}%
\pgfpathlineto{\pgfqpoint{0.805837in}{1.129154in}}%
\pgfpathlineto{\pgfqpoint{0.740946in}{1.144486in}}%
\pgfpathlineto{\pgfqpoint{0.741858in}{1.102161in}}%
\pgfpathclose%
\pgfusepath{fill}%
\end{pgfscope}%
\begin{pgfscope}%
\pgfpathrectangle{\pgfqpoint{0.407592in}{0.222904in}}{\pgfqpoint{1.411163in}{1.411163in}}%
\pgfusepath{clip}%
\pgfsetbuttcap%
\pgfsetroundjoin%
\definecolor{currentfill}{rgb}{0.141176,0.345098,0.650980}%
\pgfsetfillcolor{currentfill}%
\pgfsetfillopacity{0.080000}%
\pgfsetlinewidth{0.000000pt}%
\definecolor{currentstroke}{rgb}{0.000000,0.000000,0.000000}%
\pgfsetstrokecolor{currentstroke}%
\pgfsetdash{}{0pt}%
\pgfpathmoveto{\pgfqpoint{1.006884in}{0.743517in}}%
\pgfusepath{fill}%
\end{pgfscope}%
\begin{pgfscope}%
\pgfpathrectangle{\pgfqpoint{0.407592in}{0.222904in}}{\pgfqpoint{1.411163in}{1.411163in}}%
\pgfusepath{clip}%
\pgfsetbuttcap%
\pgfsetroundjoin%
\definecolor{currentfill}{rgb}{0.141176,0.345098,0.650980}%
\pgfsetfillcolor{currentfill}%
\pgfsetfillopacity{0.080000}%
\pgfsetlinewidth{0.000000pt}%
\definecolor{currentstroke}{rgb}{0.000000,0.000000,0.000000}%
\pgfsetstrokecolor{currentstroke}%
\pgfsetdash{}{0pt}%
\pgfpathmoveto{\pgfqpoint{0.808119in}{1.002242in}}%
\pgfpathlineto{\pgfqpoint{0.873218in}{0.986224in}}%
\pgfpathlineto{\pgfqpoint{0.872610in}{1.028506in}}%
\pgfpathlineto{\pgfqpoint{0.807362in}{1.044348in}}%
\pgfpathlineto{\pgfqpoint{0.808119in}{1.002242in}}%
\pgfpathclose%
\pgfusepath{fill}%
\end{pgfscope}%
\begin{pgfscope}%
\pgfpathrectangle{\pgfqpoint{0.407592in}{0.222904in}}{\pgfqpoint{1.411163in}{1.411163in}}%
\pgfusepath{clip}%
\pgfsetbuttcap%
\pgfsetroundjoin%
\definecolor{currentfill}{rgb}{0.141176,0.345098,0.650980}%
\pgfsetfillcolor{currentfill}%
\pgfsetfillopacity{0.080000}%
\pgfsetlinewidth{0.000000pt}%
\definecolor{currentstroke}{rgb}{0.000000,0.000000,0.000000}%
\pgfsetstrokecolor{currentstroke}%
\pgfsetdash{}{0pt}%
\pgfpathmoveto{\pgfqpoint{0.874426in}{0.902253in}}%
\pgfpathlineto{\pgfqpoint{0.939881in}{0.885727in}}%
\pgfpathlineto{\pgfqpoint{0.939429in}{0.927787in}}%
\pgfpathlineto{\pgfqpoint{0.873823in}{0.944141in}}%
\pgfpathlineto{\pgfqpoint{0.874426in}{0.902253in}}%
\pgfpathclose%
\pgfusepath{fill}%
\end{pgfscope}%
\begin{pgfscope}%
\pgfpathrectangle{\pgfqpoint{0.407592in}{0.222904in}}{\pgfqpoint{1.411163in}{1.411163in}}%
\pgfusepath{clip}%
\pgfsetbuttcap%
\pgfsetroundjoin%
\definecolor{currentfill}{rgb}{0.141176,0.345098,0.650980}%
\pgfsetfillcolor{currentfill}%
\pgfsetfillopacity{0.080000}%
\pgfsetlinewidth{0.000000pt}%
\definecolor{currentstroke}{rgb}{0.000000,0.000000,0.000000}%
\pgfsetstrokecolor{currentstroke}%
\pgfsetdash{}{0pt}%
\pgfpathmoveto{\pgfqpoint{0.940778in}{0.802195in}}%
\pgfpathlineto{\pgfqpoint{1.006591in}{0.785160in}}%
\pgfpathlineto{\pgfqpoint{1.006295in}{0.826998in}}%
\pgfpathlineto{\pgfqpoint{0.940331in}{0.843863in}}%
\pgfpathlineto{\pgfqpoint{0.940778in}{0.802195in}}%
\pgfpathclose%
\pgfusepath{fill}%
\end{pgfscope}%
\begin{pgfscope}%
\pgfpathrectangle{\pgfqpoint{0.407592in}{0.222904in}}{\pgfqpoint{1.411163in}{1.411163in}}%
\pgfusepath{clip}%
\pgfsetbuttcap%
\pgfsetroundjoin%
\definecolor{currentfill}{rgb}{0.141176,0.345098,0.650980}%
\pgfsetfillcolor{currentfill}%
\pgfsetfillopacity{0.080000}%
\pgfsetlinewidth{0.000000pt}%
\definecolor{currentstroke}{rgb}{0.000000,0.000000,0.000000}%
\pgfsetstrokecolor{currentstroke}%
\pgfsetdash{}{0pt}%
\pgfpathmoveto{\pgfqpoint{0.740946in}{1.144486in}}%
\pgfpathlineto{\pgfqpoint{0.805837in}{1.129154in}}%
\pgfpathlineto{\pgfqpoint{0.805069in}{1.171856in}}%
\pgfpathlineto{\pgfqpoint{0.740030in}{1.187010in}}%
\pgfpathlineto{\pgfqpoint{0.740946in}{1.144486in}}%
\pgfpathclose%
\pgfusepath{fill}%
\end{pgfscope}%
\begin{pgfscope}%
\pgfpathrectangle{\pgfqpoint{0.407592in}{0.222904in}}{\pgfqpoint{1.411163in}{1.411163in}}%
\pgfusepath{clip}%
\pgfsetbuttcap%
\pgfsetroundjoin%
\definecolor{currentfill}{rgb}{0.141176,0.345098,0.650980}%
\pgfsetfillcolor{currentfill}%
\pgfsetfillopacity{0.080000}%
\pgfsetlinewidth{0.000000pt}%
\definecolor{currentstroke}{rgb}{0.000000,0.000000,0.000000}%
\pgfsetstrokecolor{currentstroke}%
\pgfsetdash{}{0pt}%
\pgfpathmoveto{\pgfqpoint{0.807362in}{1.044348in}}%
\pgfpathlineto{\pgfqpoint{0.872610in}{1.028506in}}%
\pgfpathlineto{\pgfqpoint{0.871999in}{1.070986in}}%
\pgfpathlineto{\pgfqpoint{0.806601in}{1.086652in}}%
\pgfpathlineto{\pgfqpoint{0.807362in}{1.044348in}}%
\pgfpathclose%
\pgfusepath{fill}%
\end{pgfscope}%
\begin{pgfscope}%
\pgfpathrectangle{\pgfqpoint{0.407592in}{0.222904in}}{\pgfqpoint{1.411163in}{1.411163in}}%
\pgfusepath{clip}%
\pgfsetbuttcap%
\pgfsetroundjoin%
\definecolor{currentfill}{rgb}{0.141176,0.345098,0.650980}%
\pgfsetfillcolor{currentfill}%
\pgfsetfillopacity{0.080000}%
\pgfsetlinewidth{0.000000pt}%
\definecolor{currentstroke}{rgb}{0.000000,0.000000,0.000000}%
\pgfsetstrokecolor{currentstroke}%
\pgfsetdash{}{0pt}%
\pgfpathmoveto{\pgfqpoint{0.873823in}{0.944141in}}%
\pgfpathlineto{\pgfqpoint{0.939429in}{0.927787in}}%
\pgfpathlineto{\pgfqpoint{0.938975in}{0.970045in}}%
\pgfpathlineto{\pgfqpoint{0.873218in}{0.986224in}}%
\pgfpathlineto{\pgfqpoint{0.873823in}{0.944141in}}%
\pgfpathclose%
\pgfusepath{fill}%
\end{pgfscope}%
\begin{pgfscope}%
\pgfpathrectangle{\pgfqpoint{0.407592in}{0.222904in}}{\pgfqpoint{1.411163in}{1.411163in}}%
\pgfusepath{clip}%
\pgfsetbuttcap%
\pgfsetroundjoin%
\definecolor{currentfill}{rgb}{0.141176,0.345098,0.650980}%
\pgfsetfillcolor{currentfill}%
\pgfsetfillopacity{0.080000}%
\pgfsetlinewidth{0.000000pt}%
\definecolor{currentstroke}{rgb}{0.000000,0.000000,0.000000}%
\pgfsetstrokecolor{currentstroke}%
\pgfsetdash{}{0pt}%
\pgfpathmoveto{\pgfqpoint{1.006884in}{0.743517in}}%
\pgfpathlineto{\pgfqpoint{1.073069in}{0.767951in}}%
\pgfpathlineto{\pgfqpoint{1.006591in}{0.785160in}}%
\pgfusepath{fill}%
\end{pgfscope}%
\begin{pgfscope}%
\pgfpathrectangle{\pgfqpoint{0.407592in}{0.222904in}}{\pgfqpoint{1.411163in}{1.411163in}}%
\pgfusepath{clip}%
\pgfsetbuttcap%
\pgfsetroundjoin%
\definecolor{currentfill}{rgb}{0.141176,0.345098,0.650980}%
\pgfsetfillcolor{currentfill}%
\pgfsetfillopacity{0.080000}%
\pgfsetlinewidth{0.000000pt}%
\definecolor{currentstroke}{rgb}{0.000000,0.000000,0.000000}%
\pgfsetstrokecolor{currentstroke}%
\pgfsetdash{}{0pt}%
\pgfpathmoveto{\pgfqpoint{0.940331in}{0.843863in}}%
\pgfpathlineto{\pgfqpoint{1.006295in}{0.826998in}}%
\pgfpathlineto{\pgfqpoint{1.005999in}{0.869033in}}%
\pgfpathlineto{\pgfqpoint{0.939881in}{0.885727in}}%
\pgfpathlineto{\pgfqpoint{0.940331in}{0.843863in}}%
\pgfpathclose%
\pgfusepath{fill}%
\end{pgfscope}%
\begin{pgfscope}%
\pgfpathrectangle{\pgfqpoint{0.407592in}{0.222904in}}{\pgfqpoint{1.411163in}{1.411163in}}%
\pgfusepath{clip}%
\pgfsetbuttcap%
\pgfsetroundjoin%
\definecolor{currentfill}{rgb}{0.141176,0.345098,0.650980}%
\pgfsetfillcolor{currentfill}%
\pgfsetfillopacity{0.080000}%
\pgfsetlinewidth{0.000000pt}%
\definecolor{currentstroke}{rgb}{0.000000,0.000000,0.000000}%
\pgfsetstrokecolor{currentstroke}%
\pgfsetdash{}{0pt}%
\pgfpathmoveto{\pgfqpoint{0.740030in}{1.187010in}}%
\pgfpathlineto{\pgfqpoint{0.805069in}{1.171856in}}%
\pgfpathlineto{\pgfqpoint{0.804297in}{1.214760in}}%
\pgfpathlineto{\pgfqpoint{0.739110in}{1.229733in}}%
\pgfpathlineto{\pgfqpoint{0.740030in}{1.187010in}}%
\pgfpathclose%
\pgfusepath{fill}%
\end{pgfscope}%
\begin{pgfscope}%
\pgfpathrectangle{\pgfqpoint{0.407592in}{0.222904in}}{\pgfqpoint{1.411163in}{1.411163in}}%
\pgfusepath{clip}%
\pgfsetbuttcap%
\pgfsetroundjoin%
\definecolor{currentfill}{rgb}{0.141176,0.345098,0.650980}%
\pgfsetfillcolor{currentfill}%
\pgfsetfillopacity{0.080000}%
\pgfsetlinewidth{0.000000pt}%
\definecolor{currentstroke}{rgb}{0.000000,0.000000,0.000000}%
\pgfsetstrokecolor{currentstroke}%
\pgfsetdash{}{0pt}%
\pgfpathmoveto{\pgfqpoint{0.806601in}{1.086652in}}%
\pgfpathlineto{\pgfqpoint{0.871999in}{1.070986in}}%
\pgfpathlineto{\pgfqpoint{0.871385in}{1.113667in}}%
\pgfpathlineto{\pgfqpoint{0.805837in}{1.129154in}}%
\pgfpathlineto{\pgfqpoint{0.806601in}{1.086652in}}%
\pgfpathclose%
\pgfusepath{fill}%
\end{pgfscope}%
\begin{pgfscope}%
\pgfpathrectangle{\pgfqpoint{0.407592in}{0.222904in}}{\pgfqpoint{1.411163in}{1.411163in}}%
\pgfusepath{clip}%
\pgfsetbuttcap%
\pgfsetroundjoin%
\definecolor{currentfill}{rgb}{0.141176,0.345098,0.650980}%
\pgfsetfillcolor{currentfill}%
\pgfsetfillopacity{0.080000}%
\pgfsetlinewidth{0.000000pt}%
\definecolor{currentstroke}{rgb}{0.000000,0.000000,0.000000}%
\pgfsetstrokecolor{currentstroke}%
\pgfsetdash{}{0pt}%
\pgfpathmoveto{\pgfqpoint{0.873218in}{0.986224in}}%
\pgfpathlineto{\pgfqpoint{0.938975in}{0.970045in}}%
\pgfpathlineto{\pgfqpoint{0.938519in}{1.012502in}}%
\pgfpathlineto{\pgfqpoint{0.872610in}{1.028506in}}%
\pgfpathlineto{\pgfqpoint{0.873218in}{0.986224in}}%
\pgfpathclose%
\pgfusepath{fill}%
\end{pgfscope}%
\begin{pgfscope}%
\pgfpathrectangle{\pgfqpoint{0.407592in}{0.222904in}}{\pgfqpoint{1.411163in}{1.411163in}}%
\pgfusepath{clip}%
\pgfsetbuttcap%
\pgfsetroundjoin%
\definecolor{currentfill}{rgb}{0.141176,0.345098,0.650980}%
\pgfsetfillcolor{currentfill}%
\pgfsetfillopacity{0.080000}%
\pgfsetlinewidth{0.000000pt}%
\definecolor{currentstroke}{rgb}{0.000000,0.000000,0.000000}%
\pgfsetstrokecolor{currentstroke}%
\pgfsetdash{}{0pt}%
\pgfpathmoveto{\pgfqpoint{0.939881in}{0.885727in}}%
\pgfpathlineto{\pgfqpoint{1.005999in}{0.869033in}}%
\pgfpathlineto{\pgfqpoint{1.005701in}{0.911267in}}%
\pgfpathlineto{\pgfqpoint{0.939429in}{0.927787in}}%
\pgfpathlineto{\pgfqpoint{0.939881in}{0.885727in}}%
\pgfpathclose%
\pgfusepath{fill}%
\end{pgfscope}%
\begin{pgfscope}%
\pgfpathrectangle{\pgfqpoint{0.407592in}{0.222904in}}{\pgfqpoint{1.411163in}{1.411163in}}%
\pgfusepath{clip}%
\pgfsetbuttcap%
\pgfsetroundjoin%
\definecolor{currentfill}{rgb}{0.141176,0.345098,0.650980}%
\pgfsetfillcolor{currentfill}%
\pgfsetfillopacity{0.080000}%
\pgfsetlinewidth{0.000000pt}%
\definecolor{currentstroke}{rgb}{0.000000,0.000000,0.000000}%
\pgfsetstrokecolor{currentstroke}%
\pgfsetdash{}{0pt}%
\pgfpathmoveto{\pgfqpoint{1.006591in}{0.785160in}}%
\pgfpathlineto{\pgfqpoint{1.073069in}{0.767951in}}%
\pgfpathlineto{\pgfqpoint{1.072929in}{0.809961in}}%
\pgfpathlineto{\pgfqpoint{1.006295in}{0.826998in}}%
\pgfpathlineto{\pgfqpoint{1.006591in}{0.785160in}}%
\pgfpathclose%
\pgfusepath{fill}%
\end{pgfscope}%
\begin{pgfscope}%
\pgfpathrectangle{\pgfqpoint{0.407592in}{0.222904in}}{\pgfqpoint{1.411163in}{1.411163in}}%
\pgfusepath{clip}%
\pgfsetbuttcap%
\pgfsetroundjoin%
\definecolor{currentfill}{rgb}{0.141176,0.345098,0.650980}%
\pgfsetfillcolor{currentfill}%
\pgfsetfillopacity{0.080000}%
\pgfsetlinewidth{0.000000pt}%
\definecolor{currentstroke}{rgb}{0.000000,0.000000,0.000000}%
\pgfsetstrokecolor{currentstroke}%
\pgfsetdash{}{0pt}%
\pgfpathmoveto{\pgfqpoint{0.805837in}{1.129154in}}%
\pgfpathlineto{\pgfqpoint{0.871385in}{1.113667in}}%
\pgfpathlineto{\pgfqpoint{0.870768in}{1.156549in}}%
\pgfpathlineto{\pgfqpoint{0.805069in}{1.171856in}}%
\pgfpathlineto{\pgfqpoint{0.805837in}{1.129154in}}%
\pgfpathclose%
\pgfusepath{fill}%
\end{pgfscope}%
\begin{pgfscope}%
\pgfpathrectangle{\pgfqpoint{0.407592in}{0.222904in}}{\pgfqpoint{1.411163in}{1.411163in}}%
\pgfusepath{clip}%
\pgfsetbuttcap%
\pgfsetroundjoin%
\definecolor{currentfill}{rgb}{0.141176,0.345098,0.650980}%
\pgfsetfillcolor{currentfill}%
\pgfsetfillopacity{0.080000}%
\pgfsetlinewidth{0.000000pt}%
\definecolor{currentstroke}{rgb}{0.000000,0.000000,0.000000}%
\pgfsetstrokecolor{currentstroke}%
\pgfsetdash{}{0pt}%
\pgfpathmoveto{\pgfqpoint{1.073069in}{0.767951in}}%
\pgfusepath{fill}%
\end{pgfscope}%
\begin{pgfscope}%
\pgfpathrectangle{\pgfqpoint{0.407592in}{0.222904in}}{\pgfqpoint{1.411163in}{1.411163in}}%
\pgfusepath{clip}%
\pgfsetbuttcap%
\pgfsetroundjoin%
\definecolor{currentfill}{rgb}{0.141176,0.345098,0.650980}%
\pgfsetfillcolor{currentfill}%
\pgfsetfillopacity{0.080000}%
\pgfsetlinewidth{0.000000pt}%
\definecolor{currentstroke}{rgb}{0.000000,0.000000,0.000000}%
\pgfsetstrokecolor{currentstroke}%
\pgfsetdash{}{0pt}%
\pgfpathmoveto{\pgfqpoint{0.872610in}{1.028506in}}%
\pgfpathlineto{\pgfqpoint{0.938519in}{1.012502in}}%
\pgfpathlineto{\pgfqpoint{0.938061in}{1.055161in}}%
\pgfpathlineto{\pgfqpoint{0.871999in}{1.070986in}}%
\pgfpathlineto{\pgfqpoint{0.872610in}{1.028506in}}%
\pgfpathclose%
\pgfusepath{fill}%
\end{pgfscope}%
\begin{pgfscope}%
\pgfpathrectangle{\pgfqpoint{0.407592in}{0.222904in}}{\pgfqpoint{1.411163in}{1.411163in}}%
\pgfusepath{clip}%
\pgfsetbuttcap%
\pgfsetroundjoin%
\definecolor{currentfill}{rgb}{0.141176,0.345098,0.650980}%
\pgfsetfillcolor{currentfill}%
\pgfsetfillopacity{0.080000}%
\pgfsetlinewidth{0.000000pt}%
\definecolor{currentstroke}{rgb}{0.000000,0.000000,0.000000}%
\pgfsetstrokecolor{currentstroke}%
\pgfsetdash{}{0pt}%
\pgfpathmoveto{\pgfqpoint{0.939429in}{0.927787in}}%
\pgfpathlineto{\pgfqpoint{1.005701in}{0.911267in}}%
\pgfpathlineto{\pgfqpoint{1.005401in}{0.953701in}}%
\pgfpathlineto{\pgfqpoint{0.938975in}{0.970045in}}%
\pgfpathlineto{\pgfqpoint{0.939429in}{0.927787in}}%
\pgfpathclose%
\pgfusepath{fill}%
\end{pgfscope}%
\begin{pgfscope}%
\pgfpathrectangle{\pgfqpoint{0.407592in}{0.222904in}}{\pgfqpoint{1.411163in}{1.411163in}}%
\pgfusepath{clip}%
\pgfsetbuttcap%
\pgfsetroundjoin%
\definecolor{currentfill}{rgb}{0.141176,0.345098,0.650980}%
\pgfsetfillcolor{currentfill}%
\pgfsetfillopacity{0.080000}%
\pgfsetlinewidth{0.000000pt}%
\definecolor{currentstroke}{rgb}{0.000000,0.000000,0.000000}%
\pgfsetstrokecolor{currentstroke}%
\pgfsetdash{}{0pt}%
\pgfpathmoveto{\pgfqpoint{1.006295in}{0.826998in}}%
\pgfpathlineto{\pgfqpoint{1.072929in}{0.809961in}}%
\pgfpathlineto{\pgfqpoint{1.072789in}{0.852170in}}%
\pgfpathlineto{\pgfqpoint{1.005999in}{0.869033in}}%
\pgfpathlineto{\pgfqpoint{1.006295in}{0.826998in}}%
\pgfpathclose%
\pgfusepath{fill}%
\end{pgfscope}%
\begin{pgfscope}%
\pgfpathrectangle{\pgfqpoint{0.407592in}{0.222904in}}{\pgfqpoint{1.411163in}{1.411163in}}%
\pgfusepath{clip}%
\pgfsetbuttcap%
\pgfsetroundjoin%
\definecolor{currentfill}{rgb}{0.141176,0.345098,0.650980}%
\pgfsetfillcolor{currentfill}%
\pgfsetfillopacity{0.080000}%
\pgfsetlinewidth{0.000000pt}%
\definecolor{currentstroke}{rgb}{0.000000,0.000000,0.000000}%
\pgfsetstrokecolor{currentstroke}%
\pgfsetdash{}{0pt}%
\pgfpathmoveto{\pgfqpoint{0.805069in}{1.171856in}}%
\pgfpathlineto{\pgfqpoint{0.870768in}{1.156549in}}%
\pgfpathlineto{\pgfqpoint{0.870148in}{1.199635in}}%
\pgfpathlineto{\pgfqpoint{0.804297in}{1.214760in}}%
\pgfpathlineto{\pgfqpoint{0.805069in}{1.171856in}}%
\pgfpathclose%
\pgfusepath{fill}%
\end{pgfscope}%
\begin{pgfscope}%
\pgfpathrectangle{\pgfqpoint{0.407592in}{0.222904in}}{\pgfqpoint{1.411163in}{1.411163in}}%
\pgfusepath{clip}%
\pgfsetbuttcap%
\pgfsetroundjoin%
\definecolor{currentfill}{rgb}{0.141176,0.345098,0.650980}%
\pgfsetfillcolor{currentfill}%
\pgfsetfillopacity{0.080000}%
\pgfsetlinewidth{0.000000pt}%
\definecolor{currentstroke}{rgb}{0.000000,0.000000,0.000000}%
\pgfsetstrokecolor{currentstroke}%
\pgfsetdash{}{0pt}%
\pgfpathmoveto{\pgfqpoint{0.871999in}{1.070986in}}%
\pgfpathlineto{\pgfqpoint{0.938061in}{1.055161in}}%
\pgfpathlineto{\pgfqpoint{0.937601in}{1.098021in}}%
\pgfpathlineto{\pgfqpoint{0.871385in}{1.113667in}}%
\pgfpathlineto{\pgfqpoint{0.871999in}{1.070986in}}%
\pgfpathclose%
\pgfusepath{fill}%
\end{pgfscope}%
\begin{pgfscope}%
\pgfpathrectangle{\pgfqpoint{0.407592in}{0.222904in}}{\pgfqpoint{1.411163in}{1.411163in}}%
\pgfusepath{clip}%
\pgfsetbuttcap%
\pgfsetroundjoin%
\definecolor{currentfill}{rgb}{0.141176,0.345098,0.650980}%
\pgfsetfillcolor{currentfill}%
\pgfsetfillopacity{0.080000}%
\pgfsetlinewidth{0.000000pt}%
\definecolor{currentstroke}{rgb}{0.000000,0.000000,0.000000}%
\pgfsetstrokecolor{currentstroke}%
\pgfsetdash{}{0pt}%
\pgfpathmoveto{\pgfqpoint{0.938975in}{0.970045in}}%
\pgfpathlineto{\pgfqpoint{1.005401in}{0.953701in}}%
\pgfpathlineto{\pgfqpoint{1.005101in}{0.996336in}}%
\pgfpathlineto{\pgfqpoint{0.938519in}{1.012502in}}%
\pgfpathlineto{\pgfqpoint{0.938975in}{0.970045in}}%
\pgfpathclose%
\pgfusepath{fill}%
\end{pgfscope}%
\begin{pgfscope}%
\pgfpathrectangle{\pgfqpoint{0.407592in}{0.222904in}}{\pgfqpoint{1.411163in}{1.411163in}}%
\pgfusepath{clip}%
\pgfsetbuttcap%
\pgfsetroundjoin%
\definecolor{currentfill}{rgb}{0.141176,0.345098,0.650980}%
\pgfsetfillcolor{currentfill}%
\pgfsetfillopacity{0.080000}%
\pgfsetlinewidth{0.000000pt}%
\definecolor{currentstroke}{rgb}{0.000000,0.000000,0.000000}%
\pgfsetstrokecolor{currentstroke}%
\pgfsetdash{}{0pt}%
\pgfpathmoveto{\pgfqpoint{1.073069in}{0.767951in}}%
\pgfpathlineto{\pgfqpoint{1.140243in}{0.792751in}}%
\pgfpathlineto{\pgfqpoint{1.072929in}{0.809961in}}%
\pgfusepath{fill}%
\end{pgfscope}%
\begin{pgfscope}%
\pgfpathrectangle{\pgfqpoint{0.407592in}{0.222904in}}{\pgfqpoint{1.411163in}{1.411163in}}%
\pgfusepath{clip}%
\pgfsetbuttcap%
\pgfsetroundjoin%
\definecolor{currentfill}{rgb}{0.141176,0.345098,0.650980}%
\pgfsetfillcolor{currentfill}%
\pgfsetfillopacity{0.080000}%
\pgfsetlinewidth{0.000000pt}%
\definecolor{currentstroke}{rgb}{0.000000,0.000000,0.000000}%
\pgfsetstrokecolor{currentstroke}%
\pgfsetdash{}{0pt}%
\pgfpathmoveto{\pgfqpoint{1.005999in}{0.869033in}}%
\pgfpathlineto{\pgfqpoint{1.072789in}{0.852170in}}%
\pgfpathlineto{\pgfqpoint{1.072648in}{0.894579in}}%
\pgfpathlineto{\pgfqpoint{1.005701in}{0.911267in}}%
\pgfpathlineto{\pgfqpoint{1.005999in}{0.869033in}}%
\pgfpathclose%
\pgfusepath{fill}%
\end{pgfscope}%
\begin{pgfscope}%
\pgfpathrectangle{\pgfqpoint{0.407592in}{0.222904in}}{\pgfqpoint{1.411163in}{1.411163in}}%
\pgfusepath{clip}%
\pgfsetbuttcap%
\pgfsetroundjoin%
\definecolor{currentfill}{rgb}{0.141176,0.345098,0.650980}%
\pgfsetfillcolor{currentfill}%
\pgfsetfillopacity{0.080000}%
\pgfsetlinewidth{0.000000pt}%
\definecolor{currentstroke}{rgb}{0.000000,0.000000,0.000000}%
\pgfsetstrokecolor{currentstroke}%
\pgfsetdash{}{0pt}%
\pgfpathmoveto{\pgfqpoint{0.871385in}{1.113667in}}%
\pgfpathlineto{\pgfqpoint{0.937601in}{1.098021in}}%
\pgfpathlineto{\pgfqpoint{0.937138in}{1.141086in}}%
\pgfpathlineto{\pgfqpoint{0.870768in}{1.156549in}}%
\pgfpathlineto{\pgfqpoint{0.871385in}{1.113667in}}%
\pgfpathclose%
\pgfusepath{fill}%
\end{pgfscope}%
\begin{pgfscope}%
\pgfpathrectangle{\pgfqpoint{0.407592in}{0.222904in}}{\pgfqpoint{1.411163in}{1.411163in}}%
\pgfusepath{clip}%
\pgfsetbuttcap%
\pgfsetroundjoin%
\definecolor{currentfill}{rgb}{0.141176,0.345098,0.650980}%
\pgfsetfillcolor{currentfill}%
\pgfsetfillopacity{0.080000}%
\pgfsetlinewidth{0.000000pt}%
\definecolor{currentstroke}{rgb}{0.000000,0.000000,0.000000}%
\pgfsetstrokecolor{currentstroke}%
\pgfsetdash{}{0pt}%
\pgfpathmoveto{\pgfqpoint{0.938519in}{1.012502in}}%
\pgfpathlineto{\pgfqpoint{1.005101in}{0.996336in}}%
\pgfpathlineto{\pgfqpoint{1.004798in}{1.039174in}}%
\pgfpathlineto{\pgfqpoint{0.938061in}{1.055161in}}%
\pgfpathlineto{\pgfqpoint{0.938519in}{1.012502in}}%
\pgfpathclose%
\pgfusepath{fill}%
\end{pgfscope}%
\begin{pgfscope}%
\pgfpathrectangle{\pgfqpoint{0.407592in}{0.222904in}}{\pgfqpoint{1.411163in}{1.411163in}}%
\pgfusepath{clip}%
\pgfsetbuttcap%
\pgfsetroundjoin%
\definecolor{currentfill}{rgb}{0.141176,0.345098,0.650980}%
\pgfsetfillcolor{currentfill}%
\pgfsetfillopacity{0.080000}%
\pgfsetlinewidth{0.000000pt}%
\definecolor{currentstroke}{rgb}{0.000000,0.000000,0.000000}%
\pgfsetstrokecolor{currentstroke}%
\pgfsetdash{}{0pt}%
\pgfpathmoveto{\pgfqpoint{1.005701in}{0.911267in}}%
\pgfpathlineto{\pgfqpoint{1.072648in}{0.894579in}}%
\pgfpathlineto{\pgfqpoint{1.072506in}{0.937190in}}%
\pgfpathlineto{\pgfqpoint{1.005401in}{0.953701in}}%
\pgfpathlineto{\pgfqpoint{1.005701in}{0.911267in}}%
\pgfpathclose%
\pgfusepath{fill}%
\end{pgfscope}%
\begin{pgfscope}%
\pgfpathrectangle{\pgfqpoint{0.407592in}{0.222904in}}{\pgfqpoint{1.411163in}{1.411163in}}%
\pgfusepath{clip}%
\pgfsetbuttcap%
\pgfsetroundjoin%
\definecolor{currentfill}{rgb}{0.141176,0.345098,0.650980}%
\pgfsetfillcolor{currentfill}%
\pgfsetfillopacity{0.080000}%
\pgfsetlinewidth{0.000000pt}%
\definecolor{currentstroke}{rgb}{0.000000,0.000000,0.000000}%
\pgfsetstrokecolor{currentstroke}%
\pgfsetdash{}{0pt}%
\pgfpathmoveto{\pgfqpoint{1.072929in}{0.809961in}}%
\pgfpathlineto{\pgfqpoint{1.140243in}{0.792751in}}%
\pgfpathlineto{\pgfqpoint{1.140262in}{0.835135in}}%
\pgfpathlineto{\pgfqpoint{1.072789in}{0.852170in}}%
\pgfpathlineto{\pgfqpoint{1.072929in}{0.809961in}}%
\pgfpathclose%
\pgfusepath{fill}%
\end{pgfscope}%
\begin{pgfscope}%
\pgfpathrectangle{\pgfqpoint{0.407592in}{0.222904in}}{\pgfqpoint{1.411163in}{1.411163in}}%
\pgfusepath{clip}%
\pgfsetbuttcap%
\pgfsetroundjoin%
\definecolor{currentfill}{rgb}{0.141176,0.345098,0.650980}%
\pgfsetfillcolor{currentfill}%
\pgfsetfillopacity{0.080000}%
\pgfsetlinewidth{0.000000pt}%
\definecolor{currentstroke}{rgb}{0.000000,0.000000,0.000000}%
\pgfsetstrokecolor{currentstroke}%
\pgfsetdash{}{0pt}%
\pgfpathmoveto{\pgfqpoint{0.870768in}{1.156549in}}%
\pgfpathlineto{\pgfqpoint{0.937138in}{1.141086in}}%
\pgfpathlineto{\pgfqpoint{0.936673in}{1.184355in}}%
\pgfpathlineto{\pgfqpoint{0.870148in}{1.199635in}}%
\pgfpathlineto{\pgfqpoint{0.870768in}{1.156549in}}%
\pgfpathclose%
\pgfusepath{fill}%
\end{pgfscope}%
\begin{pgfscope}%
\pgfpathrectangle{\pgfqpoint{0.407592in}{0.222904in}}{\pgfqpoint{1.411163in}{1.411163in}}%
\pgfusepath{clip}%
\pgfsetbuttcap%
\pgfsetroundjoin%
\definecolor{currentfill}{rgb}{0.141176,0.345098,0.650980}%
\pgfsetfillcolor{currentfill}%
\pgfsetfillopacity{0.080000}%
\pgfsetlinewidth{0.000000pt}%
\definecolor{currentstroke}{rgb}{0.000000,0.000000,0.000000}%
\pgfsetstrokecolor{currentstroke}%
\pgfsetdash{}{0pt}%
\pgfpathmoveto{\pgfqpoint{1.140243in}{0.792751in}}%
\pgfusepath{fill}%
\end{pgfscope}%
\begin{pgfscope}%
\pgfpathrectangle{\pgfqpoint{0.407592in}{0.222904in}}{\pgfqpoint{1.411163in}{1.411163in}}%
\pgfusepath{clip}%
\pgfsetbuttcap%
\pgfsetroundjoin%
\definecolor{currentfill}{rgb}{0.141176,0.345098,0.650980}%
\pgfsetfillcolor{currentfill}%
\pgfsetfillopacity{0.080000}%
\pgfsetlinewidth{0.000000pt}%
\definecolor{currentstroke}{rgb}{0.000000,0.000000,0.000000}%
\pgfsetstrokecolor{currentstroke}%
\pgfsetdash{}{0pt}%
\pgfpathmoveto{\pgfqpoint{0.938061in}{1.055161in}}%
\pgfpathlineto{\pgfqpoint{1.004798in}{1.039174in}}%
\pgfpathlineto{\pgfqpoint{1.004495in}{1.082216in}}%
\pgfpathlineto{\pgfqpoint{0.937601in}{1.098021in}}%
\pgfpathlineto{\pgfqpoint{0.938061in}{1.055161in}}%
\pgfpathclose%
\pgfusepath{fill}%
\end{pgfscope}%
\begin{pgfscope}%
\pgfpathrectangle{\pgfqpoint{0.407592in}{0.222904in}}{\pgfqpoint{1.411163in}{1.411163in}}%
\pgfusepath{clip}%
\pgfsetbuttcap%
\pgfsetroundjoin%
\definecolor{currentfill}{rgb}{0.141176,0.345098,0.650980}%
\pgfsetfillcolor{currentfill}%
\pgfsetfillopacity{0.080000}%
\pgfsetlinewidth{0.000000pt}%
\definecolor{currentstroke}{rgb}{0.000000,0.000000,0.000000}%
\pgfsetstrokecolor{currentstroke}%
\pgfsetdash{}{0pt}%
\pgfpathmoveto{\pgfqpoint{1.005401in}{0.953701in}}%
\pgfpathlineto{\pgfqpoint{1.072506in}{0.937190in}}%
\pgfpathlineto{\pgfqpoint{1.072364in}{0.980004in}}%
\pgfpathlineto{\pgfqpoint{1.005101in}{0.996336in}}%
\pgfpathlineto{\pgfqpoint{1.005401in}{0.953701in}}%
\pgfpathclose%
\pgfusepath{fill}%
\end{pgfscope}%
\begin{pgfscope}%
\pgfpathrectangle{\pgfqpoint{0.407592in}{0.222904in}}{\pgfqpoint{1.411163in}{1.411163in}}%
\pgfusepath{clip}%
\pgfsetbuttcap%
\pgfsetroundjoin%
\definecolor{currentfill}{rgb}{0.141176,0.345098,0.650980}%
\pgfsetfillcolor{currentfill}%
\pgfsetfillopacity{0.080000}%
\pgfsetlinewidth{0.000000pt}%
\definecolor{currentstroke}{rgb}{0.000000,0.000000,0.000000}%
\pgfsetstrokecolor{currentstroke}%
\pgfsetdash{}{0pt}%
\pgfpathmoveto{\pgfqpoint{1.072789in}{0.852170in}}%
\pgfpathlineto{\pgfqpoint{1.140262in}{0.835135in}}%
\pgfpathlineto{\pgfqpoint{1.140281in}{0.877720in}}%
\pgfpathlineto{\pgfqpoint{1.072648in}{0.894579in}}%
\pgfpathlineto{\pgfqpoint{1.072789in}{0.852170in}}%
\pgfpathclose%
\pgfusepath{fill}%
\end{pgfscope}%
\begin{pgfscope}%
\pgfpathrectangle{\pgfqpoint{0.407592in}{0.222904in}}{\pgfqpoint{1.411163in}{1.411163in}}%
\pgfusepath{clip}%
\pgfsetbuttcap%
\pgfsetroundjoin%
\definecolor{currentfill}{rgb}{0.141176,0.345098,0.650980}%
\pgfsetfillcolor{currentfill}%
\pgfsetfillopacity{0.080000}%
\pgfsetlinewidth{0.000000pt}%
\definecolor{currentstroke}{rgb}{0.000000,0.000000,0.000000}%
\pgfsetstrokecolor{currentstroke}%
\pgfsetdash{}{0pt}%
\pgfpathmoveto{\pgfqpoint{0.937601in}{1.098021in}}%
\pgfpathlineto{\pgfqpoint{1.004495in}{1.082216in}}%
\pgfpathlineto{\pgfqpoint{1.004189in}{1.125463in}}%
\pgfpathlineto{\pgfqpoint{0.937138in}{1.141086in}}%
\pgfpathlineto{\pgfqpoint{0.937601in}{1.098021in}}%
\pgfpathclose%
\pgfusepath{fill}%
\end{pgfscope}%
\begin{pgfscope}%
\pgfpathrectangle{\pgfqpoint{0.407592in}{0.222904in}}{\pgfqpoint{1.411163in}{1.411163in}}%
\pgfusepath{clip}%
\pgfsetbuttcap%
\pgfsetroundjoin%
\definecolor{currentfill}{rgb}{0.141176,0.345098,0.650980}%
\pgfsetfillcolor{currentfill}%
\pgfsetfillopacity{0.080000}%
\pgfsetlinewidth{0.000000pt}%
\definecolor{currentstroke}{rgb}{0.000000,0.000000,0.000000}%
\pgfsetstrokecolor{currentstroke}%
\pgfsetdash{}{0pt}%
\pgfpathmoveto{\pgfqpoint{1.005101in}{0.996336in}}%
\pgfpathlineto{\pgfqpoint{1.072364in}{0.980004in}}%
\pgfpathlineto{\pgfqpoint{1.072221in}{1.023023in}}%
\pgfpathlineto{\pgfqpoint{1.004798in}{1.039174in}}%
\pgfpathlineto{\pgfqpoint{1.005101in}{0.996336in}}%
\pgfpathclose%
\pgfusepath{fill}%
\end{pgfscope}%
\begin{pgfscope}%
\pgfpathrectangle{\pgfqpoint{0.407592in}{0.222904in}}{\pgfqpoint{1.411163in}{1.411163in}}%
\pgfusepath{clip}%
\pgfsetbuttcap%
\pgfsetroundjoin%
\definecolor{currentfill}{rgb}{0.141176,0.345098,0.650980}%
\pgfsetfillcolor{currentfill}%
\pgfsetfillopacity{0.080000}%
\pgfsetlinewidth{0.000000pt}%
\definecolor{currentstroke}{rgb}{0.000000,0.000000,0.000000}%
\pgfsetstrokecolor{currentstroke}%
\pgfsetdash{}{0pt}%
\pgfpathmoveto{\pgfqpoint{1.140243in}{0.792751in}}%
\pgfpathlineto{\pgfqpoint{1.208428in}{0.817924in}}%
\pgfpathlineto{\pgfqpoint{1.140262in}{0.835135in}}%
\pgfusepath{fill}%
\end{pgfscope}%
\begin{pgfscope}%
\pgfpathrectangle{\pgfqpoint{0.407592in}{0.222904in}}{\pgfqpoint{1.411163in}{1.411163in}}%
\pgfusepath{clip}%
\pgfsetbuttcap%
\pgfsetroundjoin%
\definecolor{currentfill}{rgb}{0.141176,0.345098,0.650980}%
\pgfsetfillcolor{currentfill}%
\pgfsetfillopacity{0.080000}%
\pgfsetlinewidth{0.000000pt}%
\definecolor{currentstroke}{rgb}{0.000000,0.000000,0.000000}%
\pgfsetstrokecolor{currentstroke}%
\pgfsetdash{}{0pt}%
\pgfpathmoveto{\pgfqpoint{1.072648in}{0.894579in}}%
\pgfpathlineto{\pgfqpoint{1.140281in}{0.877720in}}%
\pgfpathlineto{\pgfqpoint{1.140300in}{0.920510in}}%
\pgfpathlineto{\pgfqpoint{1.072506in}{0.937190in}}%
\pgfpathlineto{\pgfqpoint{1.072648in}{0.894579in}}%
\pgfpathclose%
\pgfusepath{fill}%
\end{pgfscope}%
\begin{pgfscope}%
\pgfpathrectangle{\pgfqpoint{0.407592in}{0.222904in}}{\pgfqpoint{1.411163in}{1.411163in}}%
\pgfusepath{clip}%
\pgfsetbuttcap%
\pgfsetroundjoin%
\definecolor{currentfill}{rgb}{0.141176,0.345098,0.650980}%
\pgfsetfillcolor{currentfill}%
\pgfsetfillopacity{0.080000}%
\pgfsetlinewidth{0.000000pt}%
\definecolor{currentstroke}{rgb}{0.000000,0.000000,0.000000}%
\pgfsetstrokecolor{currentstroke}%
\pgfsetdash{}{0pt}%
\pgfpathmoveto{\pgfqpoint{0.937138in}{1.141086in}}%
\pgfpathlineto{\pgfqpoint{1.004189in}{1.125463in}}%
\pgfpathlineto{\pgfqpoint{1.003883in}{1.168918in}}%
\pgfpathlineto{\pgfqpoint{0.936673in}{1.184355in}}%
\pgfpathlineto{\pgfqpoint{0.937138in}{1.141086in}}%
\pgfpathclose%
\pgfusepath{fill}%
\end{pgfscope}%
\begin{pgfscope}%
\pgfpathrectangle{\pgfqpoint{0.407592in}{0.222904in}}{\pgfqpoint{1.411163in}{1.411163in}}%
\pgfusepath{clip}%
\pgfsetbuttcap%
\pgfsetroundjoin%
\definecolor{currentfill}{rgb}{0.141176,0.345098,0.650980}%
\pgfsetfillcolor{currentfill}%
\pgfsetfillopacity{0.080000}%
\pgfsetlinewidth{0.000000pt}%
\definecolor{currentstroke}{rgb}{0.000000,0.000000,0.000000}%
\pgfsetstrokecolor{currentstroke}%
\pgfsetdash{}{0pt}%
\pgfpathmoveto{\pgfqpoint{1.004798in}{1.039174in}}%
\pgfpathlineto{\pgfqpoint{1.072221in}{1.023023in}}%
\pgfpathlineto{\pgfqpoint{1.072077in}{1.066248in}}%
\pgfpathlineto{\pgfqpoint{1.004495in}{1.082216in}}%
\pgfpathlineto{\pgfqpoint{1.004798in}{1.039174in}}%
\pgfpathclose%
\pgfusepath{fill}%
\end{pgfscope}%
\begin{pgfscope}%
\pgfpathrectangle{\pgfqpoint{0.407592in}{0.222904in}}{\pgfqpoint{1.411163in}{1.411163in}}%
\pgfusepath{clip}%
\pgfsetbuttcap%
\pgfsetroundjoin%
\definecolor{currentfill}{rgb}{0.141176,0.345098,0.650980}%
\pgfsetfillcolor{currentfill}%
\pgfsetfillopacity{0.080000}%
\pgfsetlinewidth{0.000000pt}%
\definecolor{currentstroke}{rgb}{0.000000,0.000000,0.000000}%
\pgfsetstrokecolor{currentstroke}%
\pgfsetdash{}{0pt}%
\pgfpathmoveto{\pgfqpoint{1.072506in}{0.937190in}}%
\pgfpathlineto{\pgfqpoint{1.140300in}{0.920510in}}%
\pgfpathlineto{\pgfqpoint{1.140319in}{0.963504in}}%
\pgfpathlineto{\pgfqpoint{1.072364in}{0.980004in}}%
\pgfpathlineto{\pgfqpoint{1.072506in}{0.937190in}}%
\pgfpathclose%
\pgfusepath{fill}%
\end{pgfscope}%
\begin{pgfscope}%
\pgfpathrectangle{\pgfqpoint{0.407592in}{0.222904in}}{\pgfqpoint{1.411163in}{1.411163in}}%
\pgfusepath{clip}%
\pgfsetbuttcap%
\pgfsetroundjoin%
\definecolor{currentfill}{rgb}{0.141176,0.345098,0.650980}%
\pgfsetfillcolor{currentfill}%
\pgfsetfillopacity{0.080000}%
\pgfsetlinewidth{0.000000pt}%
\definecolor{currentstroke}{rgb}{0.000000,0.000000,0.000000}%
\pgfsetstrokecolor{currentstroke}%
\pgfsetdash{}{0pt}%
\pgfpathmoveto{\pgfqpoint{1.140262in}{0.835135in}}%
\pgfpathlineto{\pgfqpoint{1.208428in}{0.817924in}}%
\pgfpathlineto{\pgfqpoint{1.208610in}{0.860688in}}%
\pgfpathlineto{\pgfqpoint{1.140281in}{0.877720in}}%
\pgfpathlineto{\pgfqpoint{1.140262in}{0.835135in}}%
\pgfpathclose%
\pgfusepath{fill}%
\end{pgfscope}%
\begin{pgfscope}%
\pgfpathrectangle{\pgfqpoint{0.407592in}{0.222904in}}{\pgfqpoint{1.411163in}{1.411163in}}%
\pgfusepath{clip}%
\pgfsetbuttcap%
\pgfsetroundjoin%
\definecolor{currentfill}{rgb}{0.141176,0.345098,0.650980}%
\pgfsetfillcolor{currentfill}%
\pgfsetfillopacity{0.080000}%
\pgfsetlinewidth{0.000000pt}%
\definecolor{currentstroke}{rgb}{0.000000,0.000000,0.000000}%
\pgfsetstrokecolor{currentstroke}%
\pgfsetdash{}{0pt}%
\pgfpathmoveto{\pgfqpoint{1.208428in}{0.817924in}}%
\pgfusepath{fill}%
\end{pgfscope}%
\begin{pgfscope}%
\pgfpathrectangle{\pgfqpoint{0.407592in}{0.222904in}}{\pgfqpoint{1.411163in}{1.411163in}}%
\pgfusepath{clip}%
\pgfsetbuttcap%
\pgfsetroundjoin%
\definecolor{currentfill}{rgb}{0.141176,0.345098,0.650980}%
\pgfsetfillcolor{currentfill}%
\pgfsetfillopacity{0.080000}%
\pgfsetlinewidth{0.000000pt}%
\definecolor{currentstroke}{rgb}{0.000000,0.000000,0.000000}%
\pgfsetstrokecolor{currentstroke}%
\pgfsetdash{}{0pt}%
\pgfpathmoveto{\pgfqpoint{1.004495in}{1.082216in}}%
\pgfpathlineto{\pgfqpoint{1.072077in}{1.066248in}}%
\pgfpathlineto{\pgfqpoint{1.071932in}{1.109680in}}%
\pgfpathlineto{\pgfqpoint{1.004189in}{1.125463in}}%
\pgfpathlineto{\pgfqpoint{1.004495in}{1.082216in}}%
\pgfpathclose%
\pgfusepath{fill}%
\end{pgfscope}%
\begin{pgfscope}%
\pgfpathrectangle{\pgfqpoint{0.407592in}{0.222904in}}{\pgfqpoint{1.411163in}{1.411163in}}%
\pgfusepath{clip}%
\pgfsetbuttcap%
\pgfsetroundjoin%
\definecolor{currentfill}{rgb}{0.141176,0.345098,0.650980}%
\pgfsetfillcolor{currentfill}%
\pgfsetfillopacity{0.080000}%
\pgfsetlinewidth{0.000000pt}%
\definecolor{currentstroke}{rgb}{0.000000,0.000000,0.000000}%
\pgfsetstrokecolor{currentstroke}%
\pgfsetdash{}{0pt}%
\pgfpathmoveto{\pgfqpoint{1.072364in}{0.980004in}}%
\pgfpathlineto{\pgfqpoint{1.140319in}{0.963504in}}%
\pgfpathlineto{\pgfqpoint{1.140339in}{1.006705in}}%
\pgfpathlineto{\pgfqpoint{1.072221in}{1.023023in}}%
\pgfpathlineto{\pgfqpoint{1.072364in}{0.980004in}}%
\pgfpathclose%
\pgfusepath{fill}%
\end{pgfscope}%
\begin{pgfscope}%
\pgfpathrectangle{\pgfqpoint{0.407592in}{0.222904in}}{\pgfqpoint{1.411163in}{1.411163in}}%
\pgfusepath{clip}%
\pgfsetbuttcap%
\pgfsetroundjoin%
\definecolor{currentfill}{rgb}{0.141176,0.345098,0.650980}%
\pgfsetfillcolor{currentfill}%
\pgfsetfillopacity{0.080000}%
\pgfsetlinewidth{0.000000pt}%
\definecolor{currentstroke}{rgb}{0.000000,0.000000,0.000000}%
\pgfsetstrokecolor{currentstroke}%
\pgfsetdash{}{0pt}%
\pgfpathmoveto{\pgfqpoint{1.140281in}{0.877720in}}%
\pgfpathlineto{\pgfqpoint{1.208610in}{0.860688in}}%
\pgfpathlineto{\pgfqpoint{1.208794in}{0.903657in}}%
\pgfpathlineto{\pgfqpoint{1.140300in}{0.920510in}}%
\pgfpathlineto{\pgfqpoint{1.140281in}{0.877720in}}%
\pgfpathclose%
\pgfusepath{fill}%
\end{pgfscope}%
\begin{pgfscope}%
\pgfpathrectangle{\pgfqpoint{0.407592in}{0.222904in}}{\pgfqpoint{1.411163in}{1.411163in}}%
\pgfusepath{clip}%
\pgfsetbuttcap%
\pgfsetroundjoin%
\definecolor{currentfill}{rgb}{0.141176,0.345098,0.650980}%
\pgfsetfillcolor{currentfill}%
\pgfsetfillopacity{0.080000}%
\pgfsetlinewidth{0.000000pt}%
\definecolor{currentstroke}{rgb}{0.000000,0.000000,0.000000}%
\pgfsetstrokecolor{currentstroke}%
\pgfsetdash{}{0pt}%
\pgfpathmoveto{\pgfqpoint{1.208428in}{0.817924in}}%
\pgfpathlineto{\pgfqpoint{1.208610in}{0.860688in}}%
\pgfusepath{fill}%
\end{pgfscope}%
\begin{pgfscope}%
\pgfpathrectangle{\pgfqpoint{0.407592in}{0.222904in}}{\pgfqpoint{1.411163in}{1.411163in}}%
\pgfusepath{clip}%
\pgfsetbuttcap%
\pgfsetroundjoin%
\definecolor{currentfill}{rgb}{0.141176,0.345098,0.650980}%
\pgfsetfillcolor{currentfill}%
\pgfsetfillopacity{0.080000}%
\pgfsetlinewidth{0.000000pt}%
\definecolor{currentstroke}{rgb}{0.000000,0.000000,0.000000}%
\pgfsetstrokecolor{currentstroke}%
\pgfsetdash{}{0pt}%
\pgfpathmoveto{\pgfqpoint{1.004189in}{1.125463in}}%
\pgfpathlineto{\pgfqpoint{1.071932in}{1.109680in}}%
\pgfpathlineto{\pgfqpoint{1.071787in}{1.153321in}}%
\pgfpathlineto{\pgfqpoint{1.003883in}{1.168918in}}%
\pgfpathlineto{\pgfqpoint{1.004189in}{1.125463in}}%
\pgfpathclose%
\pgfusepath{fill}%
\end{pgfscope}%
\begin{pgfscope}%
\pgfpathrectangle{\pgfqpoint{0.407592in}{0.222904in}}{\pgfqpoint{1.411163in}{1.411163in}}%
\pgfusepath{clip}%
\pgfsetbuttcap%
\pgfsetroundjoin%
\definecolor{currentfill}{rgb}{0.141176,0.345098,0.650980}%
\pgfsetfillcolor{currentfill}%
\pgfsetfillopacity{0.080000}%
\pgfsetlinewidth{0.000000pt}%
\definecolor{currentstroke}{rgb}{0.000000,0.000000,0.000000}%
\pgfsetstrokecolor{currentstroke}%
\pgfsetdash{}{0pt}%
\pgfpathmoveto{\pgfqpoint{1.072221in}{1.023023in}}%
\pgfpathlineto{\pgfqpoint{1.140339in}{1.006705in}}%
\pgfpathlineto{\pgfqpoint{1.140358in}{1.050114in}}%
\pgfpathlineto{\pgfqpoint{1.072077in}{1.066248in}}%
\pgfpathlineto{\pgfqpoint{1.072221in}{1.023023in}}%
\pgfpathclose%
\pgfusepath{fill}%
\end{pgfscope}%
\begin{pgfscope}%
\pgfpathrectangle{\pgfqpoint{0.407592in}{0.222904in}}{\pgfqpoint{1.411163in}{1.411163in}}%
\pgfusepath{clip}%
\pgfsetbuttcap%
\pgfsetroundjoin%
\definecolor{currentfill}{rgb}{0.141176,0.345098,0.650980}%
\pgfsetfillcolor{currentfill}%
\pgfsetfillopacity{0.080000}%
\pgfsetlinewidth{0.000000pt}%
\definecolor{currentstroke}{rgb}{0.000000,0.000000,0.000000}%
\pgfsetstrokecolor{currentstroke}%
\pgfsetdash{}{0pt}%
\pgfpathmoveto{\pgfqpoint{1.140300in}{0.920510in}}%
\pgfpathlineto{\pgfqpoint{1.208794in}{0.903657in}}%
\pgfpathlineto{\pgfqpoint{1.208978in}{0.946833in}}%
\pgfpathlineto{\pgfqpoint{1.140319in}{0.963504in}}%
\pgfpathlineto{\pgfqpoint{1.140300in}{0.920510in}}%
\pgfpathclose%
\pgfusepath{fill}%
\end{pgfscope}%
\begin{pgfscope}%
\pgfpathrectangle{\pgfqpoint{0.407592in}{0.222904in}}{\pgfqpoint{1.411163in}{1.411163in}}%
\pgfusepath{clip}%
\pgfsetbuttcap%
\pgfsetroundjoin%
\definecolor{currentfill}{rgb}{0.141176,0.345098,0.650980}%
\pgfsetfillcolor{currentfill}%
\pgfsetfillopacity{0.080000}%
\pgfsetlinewidth{0.000000pt}%
\definecolor{currentstroke}{rgb}{0.000000,0.000000,0.000000}%
\pgfsetstrokecolor{currentstroke}%
\pgfsetdash{}{0pt}%
\pgfpathmoveto{\pgfqpoint{1.072077in}{1.066248in}}%
\pgfpathlineto{\pgfqpoint{1.140358in}{1.050114in}}%
\pgfpathlineto{\pgfqpoint{1.140378in}{1.093733in}}%
\pgfpathlineto{\pgfqpoint{1.071932in}{1.109680in}}%
\pgfpathlineto{\pgfqpoint{1.072077in}{1.066248in}}%
\pgfpathclose%
\pgfusepath{fill}%
\end{pgfscope}%
\begin{pgfscope}%
\pgfpathrectangle{\pgfqpoint{0.407592in}{0.222904in}}{\pgfqpoint{1.411163in}{1.411163in}}%
\pgfusepath{clip}%
\pgfsetbuttcap%
\pgfsetroundjoin%
\definecolor{currentfill}{rgb}{0.141176,0.345098,0.650980}%
\pgfsetfillcolor{currentfill}%
\pgfsetfillopacity{0.080000}%
\pgfsetlinewidth{0.000000pt}%
\definecolor{currentstroke}{rgb}{0.000000,0.000000,0.000000}%
\pgfsetstrokecolor{currentstroke}%
\pgfsetdash{}{0pt}%
\pgfpathmoveto{\pgfqpoint{1.208610in}{0.860688in}}%
\pgfpathlineto{\pgfqpoint{1.277998in}{0.886630in}}%
\pgfpathlineto{\pgfqpoint{1.208794in}{0.903657in}}%
\pgfusepath{fill}%
\end{pgfscope}%
\begin{pgfscope}%
\pgfpathrectangle{\pgfqpoint{0.407592in}{0.222904in}}{\pgfqpoint{1.411163in}{1.411163in}}%
\pgfusepath{clip}%
\pgfsetbuttcap%
\pgfsetroundjoin%
\definecolor{currentfill}{rgb}{0.141176,0.345098,0.650980}%
\pgfsetfillcolor{currentfill}%
\pgfsetfillopacity{0.080000}%
\pgfsetlinewidth{0.000000pt}%
\definecolor{currentstroke}{rgb}{0.000000,0.000000,0.000000}%
\pgfsetstrokecolor{currentstroke}%
\pgfsetdash{}{0pt}%
\pgfpathmoveto{\pgfqpoint{1.140319in}{0.963504in}}%
\pgfpathlineto{\pgfqpoint{1.208978in}{0.946833in}}%
\pgfpathlineto{\pgfqpoint{1.209163in}{0.990218in}}%
\pgfpathlineto{\pgfqpoint{1.140339in}{1.006705in}}%
\pgfpathlineto{\pgfqpoint{1.140319in}{0.963504in}}%
\pgfpathclose%
\pgfusepath{fill}%
\end{pgfscope}%
\begin{pgfscope}%
\pgfpathrectangle{\pgfqpoint{0.407592in}{0.222904in}}{\pgfqpoint{1.411163in}{1.411163in}}%
\pgfusepath{clip}%
\pgfsetbuttcap%
\pgfsetroundjoin%
\definecolor{currentfill}{rgb}{0.141176,0.345098,0.650980}%
\pgfsetfillcolor{currentfill}%
\pgfsetfillopacity{0.080000}%
\pgfsetlinewidth{0.000000pt}%
\definecolor{currentstroke}{rgb}{0.000000,0.000000,0.000000}%
\pgfsetstrokecolor{currentstroke}%
\pgfsetdash{}{0pt}%
\pgfpathmoveto{\pgfqpoint{1.071932in}{1.109680in}}%
\pgfpathlineto{\pgfqpoint{1.140378in}{1.093733in}}%
\pgfpathlineto{\pgfqpoint{1.140398in}{1.137562in}}%
\pgfpathlineto{\pgfqpoint{1.071787in}{1.153321in}}%
\pgfpathlineto{\pgfqpoint{1.071932in}{1.109680in}}%
\pgfpathclose%
\pgfusepath{fill}%
\end{pgfscope}%
\begin{pgfscope}%
\pgfpathrectangle{\pgfqpoint{0.407592in}{0.222904in}}{\pgfqpoint{1.411163in}{1.411163in}}%
\pgfusepath{clip}%
\pgfsetbuttcap%
\pgfsetroundjoin%
\definecolor{currentfill}{rgb}{0.141176,0.345098,0.650980}%
\pgfsetfillcolor{currentfill}%
\pgfsetfillopacity{0.080000}%
\pgfsetlinewidth{0.000000pt}%
\definecolor{currentstroke}{rgb}{0.000000,0.000000,0.000000}%
\pgfsetstrokecolor{currentstroke}%
\pgfsetdash{}{0pt}%
\pgfpathmoveto{\pgfqpoint{1.140339in}{1.006705in}}%
\pgfpathlineto{\pgfqpoint{1.209163in}{0.990218in}}%
\pgfpathlineto{\pgfqpoint{1.209350in}{1.033813in}}%
\pgfpathlineto{\pgfqpoint{1.140358in}{1.050114in}}%
\pgfpathlineto{\pgfqpoint{1.140339in}{1.006705in}}%
\pgfpathclose%
\pgfusepath{fill}%
\end{pgfscope}%
\begin{pgfscope}%
\pgfpathrectangle{\pgfqpoint{0.407592in}{0.222904in}}{\pgfqpoint{1.411163in}{1.411163in}}%
\pgfusepath{clip}%
\pgfsetbuttcap%
\pgfsetroundjoin%
\definecolor{currentfill}{rgb}{0.141176,0.345098,0.650980}%
\pgfsetfillcolor{currentfill}%
\pgfsetfillopacity{0.080000}%
\pgfsetlinewidth{0.000000pt}%
\definecolor{currentstroke}{rgb}{0.000000,0.000000,0.000000}%
\pgfsetstrokecolor{currentstroke}%
\pgfsetdash{}{0pt}%
\pgfpathmoveto{\pgfqpoint{1.208794in}{0.903657in}}%
\pgfpathlineto{\pgfqpoint{1.277998in}{0.886630in}}%
\pgfpathlineto{\pgfqpoint{1.278351in}{0.929989in}}%
\pgfpathlineto{\pgfqpoint{1.208978in}{0.946833in}}%
\pgfpathlineto{\pgfqpoint{1.208794in}{0.903657in}}%
\pgfpathclose%
\pgfusepath{fill}%
\end{pgfscope}%
\begin{pgfscope}%
\pgfpathrectangle{\pgfqpoint{0.407592in}{0.222904in}}{\pgfqpoint{1.411163in}{1.411163in}}%
\pgfusepath{clip}%
\pgfsetbuttcap%
\pgfsetroundjoin%
\definecolor{currentfill}{rgb}{0.141176,0.345098,0.650980}%
\pgfsetfillcolor{currentfill}%
\pgfsetfillopacity{0.080000}%
\pgfsetlinewidth{0.000000pt}%
\definecolor{currentstroke}{rgb}{0.000000,0.000000,0.000000}%
\pgfsetstrokecolor{currentstroke}%
\pgfsetdash{}{0pt}%
\pgfpathmoveto{\pgfqpoint{1.277998in}{0.886630in}}%
\pgfusepath{fill}%
\end{pgfscope}%
\begin{pgfscope}%
\pgfpathrectangle{\pgfqpoint{0.407592in}{0.222904in}}{\pgfqpoint{1.411163in}{1.411163in}}%
\pgfusepath{clip}%
\pgfsetbuttcap%
\pgfsetroundjoin%
\definecolor{currentfill}{rgb}{0.141176,0.345098,0.650980}%
\pgfsetfillcolor{currentfill}%
\pgfsetfillopacity{0.080000}%
\pgfsetlinewidth{0.000000pt}%
\definecolor{currentstroke}{rgb}{0.000000,0.000000,0.000000}%
\pgfsetstrokecolor{currentstroke}%
\pgfsetdash{}{0pt}%
\pgfpathmoveto{\pgfqpoint{1.140358in}{1.050114in}}%
\pgfpathlineto{\pgfqpoint{1.209350in}{1.033813in}}%
\pgfpathlineto{\pgfqpoint{1.209537in}{1.077619in}}%
\pgfpathlineto{\pgfqpoint{1.140378in}{1.093733in}}%
\pgfpathlineto{\pgfqpoint{1.140358in}{1.050114in}}%
\pgfpathclose%
\pgfusepath{fill}%
\end{pgfscope}%
\begin{pgfscope}%
\pgfpathrectangle{\pgfqpoint{0.407592in}{0.222904in}}{\pgfqpoint{1.411163in}{1.411163in}}%
\pgfusepath{clip}%
\pgfsetbuttcap%
\pgfsetroundjoin%
\definecolor{currentfill}{rgb}{0.141176,0.345098,0.650980}%
\pgfsetfillcolor{currentfill}%
\pgfsetfillopacity{0.080000}%
\pgfsetlinewidth{0.000000pt}%
\definecolor{currentstroke}{rgb}{0.000000,0.000000,0.000000}%
\pgfsetstrokecolor{currentstroke}%
\pgfsetdash{}{0pt}%
\pgfpathmoveto{\pgfqpoint{1.208978in}{0.946833in}}%
\pgfpathlineto{\pgfqpoint{1.278351in}{0.929989in}}%
\pgfpathlineto{\pgfqpoint{1.278706in}{0.973559in}}%
\pgfpathlineto{\pgfqpoint{1.209163in}{0.990218in}}%
\pgfpathlineto{\pgfqpoint{1.208978in}{0.946833in}}%
\pgfpathclose%
\pgfusepath{fill}%
\end{pgfscope}%
\begin{pgfscope}%
\pgfpathrectangle{\pgfqpoint{0.407592in}{0.222904in}}{\pgfqpoint{1.411163in}{1.411163in}}%
\pgfusepath{clip}%
\pgfsetbuttcap%
\pgfsetroundjoin%
\definecolor{currentfill}{rgb}{0.141176,0.345098,0.650980}%
\pgfsetfillcolor{currentfill}%
\pgfsetfillopacity{0.080000}%
\pgfsetlinewidth{0.000000pt}%
\definecolor{currentstroke}{rgb}{0.000000,0.000000,0.000000}%
\pgfsetstrokecolor{currentstroke}%
\pgfsetdash{}{0pt}%
\pgfpathmoveto{\pgfqpoint{1.140378in}{1.093733in}}%
\pgfpathlineto{\pgfqpoint{1.209537in}{1.077619in}}%
\pgfpathlineto{\pgfqpoint{1.209725in}{1.121639in}}%
\pgfpathlineto{\pgfqpoint{1.140398in}{1.137562in}}%
\pgfpathlineto{\pgfqpoint{1.140378in}{1.093733in}}%
\pgfpathclose%
\pgfusepath{fill}%
\end{pgfscope}%
\begin{pgfscope}%
\pgfpathrectangle{\pgfqpoint{0.407592in}{0.222904in}}{\pgfqpoint{1.411163in}{1.411163in}}%
\pgfusepath{clip}%
\pgfsetbuttcap%
\pgfsetroundjoin%
\definecolor{currentfill}{rgb}{0.141176,0.345098,0.650980}%
\pgfsetfillcolor{currentfill}%
\pgfsetfillopacity{0.080000}%
\pgfsetlinewidth{0.000000pt}%
\definecolor{currentstroke}{rgb}{0.000000,0.000000,0.000000}%
\pgfsetstrokecolor{currentstroke}%
\pgfsetdash{}{0pt}%
\pgfpathmoveto{\pgfqpoint{1.277998in}{0.886630in}}%
\pgfpathlineto{\pgfqpoint{1.348450in}{0.912969in}}%
\pgfpathlineto{\pgfqpoint{1.278351in}{0.929989in}}%
\pgfusepath{fill}%
\end{pgfscope}%
\begin{pgfscope}%
\pgfpathrectangle{\pgfqpoint{0.407592in}{0.222904in}}{\pgfqpoint{1.411163in}{1.411163in}}%
\pgfusepath{clip}%
\pgfsetbuttcap%
\pgfsetroundjoin%
\definecolor{currentfill}{rgb}{0.141176,0.345098,0.650980}%
\pgfsetfillcolor{currentfill}%
\pgfsetfillopacity{0.080000}%
\pgfsetlinewidth{0.000000pt}%
\definecolor{currentstroke}{rgb}{0.000000,0.000000,0.000000}%
\pgfsetstrokecolor{currentstroke}%
\pgfsetdash{}{0pt}%
\pgfpathmoveto{\pgfqpoint{1.209163in}{0.990218in}}%
\pgfpathlineto{\pgfqpoint{1.278706in}{0.973559in}}%
\pgfpathlineto{\pgfqpoint{1.279062in}{1.017342in}}%
\pgfpathlineto{\pgfqpoint{1.209350in}{1.033813in}}%
\pgfpathlineto{\pgfqpoint{1.209163in}{0.990218in}}%
\pgfpathclose%
\pgfusepath{fill}%
\end{pgfscope}%
\begin{pgfscope}%
\pgfpathrectangle{\pgfqpoint{0.407592in}{0.222904in}}{\pgfqpoint{1.411163in}{1.411163in}}%
\pgfusepath{clip}%
\pgfsetbuttcap%
\pgfsetroundjoin%
\definecolor{currentfill}{rgb}{0.141176,0.345098,0.650980}%
\pgfsetfillcolor{currentfill}%
\pgfsetfillopacity{0.080000}%
\pgfsetlinewidth{0.000000pt}%
\definecolor{currentstroke}{rgb}{0.000000,0.000000,0.000000}%
\pgfsetstrokecolor{currentstroke}%
\pgfsetdash{}{0pt}%
\pgfpathmoveto{\pgfqpoint{1.209350in}{1.033813in}}%
\pgfpathlineto{\pgfqpoint{1.279062in}{1.017342in}}%
\pgfpathlineto{\pgfqpoint{1.279420in}{1.061337in}}%
\pgfpathlineto{\pgfqpoint{1.209537in}{1.077619in}}%
\pgfpathlineto{\pgfqpoint{1.209350in}{1.033813in}}%
\pgfpathclose%
\pgfusepath{fill}%
\end{pgfscope}%
\begin{pgfscope}%
\pgfpathrectangle{\pgfqpoint{0.407592in}{0.222904in}}{\pgfqpoint{1.411163in}{1.411163in}}%
\pgfusepath{clip}%
\pgfsetbuttcap%
\pgfsetroundjoin%
\definecolor{currentfill}{rgb}{0.141176,0.345098,0.650980}%
\pgfsetfillcolor{currentfill}%
\pgfsetfillopacity{0.080000}%
\pgfsetlinewidth{0.000000pt}%
\definecolor{currentstroke}{rgb}{0.000000,0.000000,0.000000}%
\pgfsetstrokecolor{currentstroke}%
\pgfsetdash{}{0pt}%
\pgfpathmoveto{\pgfqpoint{1.278351in}{0.929989in}}%
\pgfpathlineto{\pgfqpoint{1.348450in}{0.912969in}}%
\pgfpathlineto{\pgfqpoint{1.348977in}{0.956726in}}%
\pgfpathlineto{\pgfqpoint{1.278706in}{0.973559in}}%
\pgfpathlineto{\pgfqpoint{1.278351in}{0.929989in}}%
\pgfpathclose%
\pgfusepath{fill}%
\end{pgfscope}%
\begin{pgfscope}%
\pgfpathrectangle{\pgfqpoint{0.407592in}{0.222904in}}{\pgfqpoint{1.411163in}{1.411163in}}%
\pgfusepath{clip}%
\pgfsetbuttcap%
\pgfsetroundjoin%
\definecolor{currentfill}{rgb}{0.141176,0.345098,0.650980}%
\pgfsetfillcolor{currentfill}%
\pgfsetfillopacity{0.080000}%
\pgfsetlinewidth{0.000000pt}%
\definecolor{currentstroke}{rgb}{0.000000,0.000000,0.000000}%
\pgfsetstrokecolor{currentstroke}%
\pgfsetdash{}{0pt}%
\pgfpathmoveto{\pgfqpoint{1.348450in}{0.912969in}}%
\pgfusepath{fill}%
\end{pgfscope}%
\begin{pgfscope}%
\pgfpathrectangle{\pgfqpoint{0.407592in}{0.222904in}}{\pgfqpoint{1.411163in}{1.411163in}}%
\pgfusepath{clip}%
\pgfsetbuttcap%
\pgfsetroundjoin%
\definecolor{currentfill}{rgb}{0.141176,0.345098,0.650980}%
\pgfsetfillcolor{currentfill}%
\pgfsetfillopacity{0.080000}%
\pgfsetlinewidth{0.000000pt}%
\definecolor{currentstroke}{rgb}{0.000000,0.000000,0.000000}%
\pgfsetstrokecolor{currentstroke}%
\pgfsetdash{}{0pt}%
\pgfpathmoveto{\pgfqpoint{1.209537in}{1.077619in}}%
\pgfpathlineto{\pgfqpoint{1.279420in}{1.061337in}}%
\pgfpathlineto{\pgfqpoint{1.279780in}{1.105548in}}%
\pgfpathlineto{\pgfqpoint{1.209725in}{1.121639in}}%
\pgfpathlineto{\pgfqpoint{1.209537in}{1.077619in}}%
\pgfpathclose%
\pgfusepath{fill}%
\end{pgfscope}%
\begin{pgfscope}%
\pgfpathrectangle{\pgfqpoint{0.407592in}{0.222904in}}{\pgfqpoint{1.411163in}{1.411163in}}%
\pgfusepath{clip}%
\pgfsetbuttcap%
\pgfsetroundjoin%
\definecolor{currentfill}{rgb}{0.141176,0.345098,0.650980}%
\pgfsetfillcolor{currentfill}%
\pgfsetfillopacity{0.080000}%
\pgfsetlinewidth{0.000000pt}%
\definecolor{currentstroke}{rgb}{0.000000,0.000000,0.000000}%
\pgfsetstrokecolor{currentstroke}%
\pgfsetdash{}{0pt}%
\pgfpathmoveto{\pgfqpoint{1.278706in}{0.973559in}}%
\pgfpathlineto{\pgfqpoint{1.348977in}{0.956726in}}%
\pgfpathlineto{\pgfqpoint{1.349507in}{1.000697in}}%
\pgfpathlineto{\pgfqpoint{1.279062in}{1.017342in}}%
\pgfpathlineto{\pgfqpoint{1.278706in}{0.973559in}}%
\pgfpathclose%
\pgfusepath{fill}%
\end{pgfscope}%
\begin{pgfscope}%
\pgfpathrectangle{\pgfqpoint{0.407592in}{0.222904in}}{\pgfqpoint{1.411163in}{1.411163in}}%
\pgfusepath{clip}%
\pgfsetbuttcap%
\pgfsetroundjoin%
\definecolor{currentfill}{rgb}{0.141176,0.345098,0.650980}%
\pgfsetfillcolor{currentfill}%
\pgfsetfillopacity{0.080000}%
\pgfsetlinewidth{0.000000pt}%
\definecolor{currentstroke}{rgb}{0.000000,0.000000,0.000000}%
\pgfsetstrokecolor{currentstroke}%
\pgfsetdash{}{0pt}%
\pgfpathmoveto{\pgfqpoint{1.348450in}{0.912969in}}%
\pgfpathlineto{\pgfqpoint{1.419989in}{0.939715in}}%
\pgfpathlineto{\pgfqpoint{1.348977in}{0.956726in}}%
\pgfusepath{fill}%
\end{pgfscope}%
\begin{pgfscope}%
\pgfpathrectangle{\pgfqpoint{0.407592in}{0.222904in}}{\pgfqpoint{1.411163in}{1.411163in}}%
\pgfusepath{clip}%
\pgfsetbuttcap%
\pgfsetroundjoin%
\definecolor{currentfill}{rgb}{0.141176,0.345098,0.650980}%
\pgfsetfillcolor{currentfill}%
\pgfsetfillopacity{0.080000}%
\pgfsetlinewidth{0.000000pt}%
\definecolor{currentstroke}{rgb}{0.000000,0.000000,0.000000}%
\pgfsetstrokecolor{currentstroke}%
\pgfsetdash{}{0pt}%
\pgfpathmoveto{\pgfqpoint{1.279062in}{1.017342in}}%
\pgfpathlineto{\pgfqpoint{1.349507in}{1.000697in}}%
\pgfpathlineto{\pgfqpoint{1.350040in}{1.044884in}}%
\pgfpathlineto{\pgfqpoint{1.279420in}{1.061337in}}%
\pgfpathlineto{\pgfqpoint{1.279062in}{1.017342in}}%
\pgfpathclose%
\pgfusepath{fill}%
\end{pgfscope}%
\begin{pgfscope}%
\pgfpathrectangle{\pgfqpoint{0.407592in}{0.222904in}}{\pgfqpoint{1.411163in}{1.411163in}}%
\pgfusepath{clip}%
\pgfsetbuttcap%
\pgfsetroundjoin%
\definecolor{currentfill}{rgb}{0.141176,0.345098,0.650980}%
\pgfsetfillcolor{currentfill}%
\pgfsetfillopacity{0.080000}%
\pgfsetlinewidth{0.000000pt}%
\definecolor{currentstroke}{rgb}{0.000000,0.000000,0.000000}%
\pgfsetstrokecolor{currentstroke}%
\pgfsetdash{}{0pt}%
\pgfpathmoveto{\pgfqpoint{1.279420in}{1.061337in}}%
\pgfpathlineto{\pgfqpoint{1.350040in}{1.044884in}}%
\pgfpathlineto{\pgfqpoint{1.350575in}{1.089287in}}%
\pgfpathlineto{\pgfqpoint{1.279780in}{1.105548in}}%
\pgfpathlineto{\pgfqpoint{1.279420in}{1.061337in}}%
\pgfpathclose%
\pgfusepath{fill}%
\end{pgfscope}%
\begin{pgfscope}%
\pgfpathrectangle{\pgfqpoint{0.407592in}{0.222904in}}{\pgfqpoint{1.411163in}{1.411163in}}%
\pgfusepath{clip}%
\pgfsetbuttcap%
\pgfsetroundjoin%
\definecolor{currentfill}{rgb}{0.141176,0.345098,0.650980}%
\pgfsetfillcolor{currentfill}%
\pgfsetfillopacity{0.080000}%
\pgfsetlinewidth{0.000000pt}%
\definecolor{currentstroke}{rgb}{0.000000,0.000000,0.000000}%
\pgfsetstrokecolor{currentstroke}%
\pgfsetdash{}{0pt}%
\pgfpathmoveto{\pgfqpoint{1.348977in}{0.956726in}}%
\pgfpathlineto{\pgfqpoint{1.419989in}{0.939715in}}%
\pgfpathlineto{\pgfqpoint{1.420696in}{0.983877in}}%
\pgfpathlineto{\pgfqpoint{1.349507in}{1.000697in}}%
\pgfpathlineto{\pgfqpoint{1.348977in}{0.956726in}}%
\pgfpathclose%
\pgfusepath{fill}%
\end{pgfscope}%
\begin{pgfscope}%
\pgfpathrectangle{\pgfqpoint{0.407592in}{0.222904in}}{\pgfqpoint{1.411163in}{1.411163in}}%
\pgfusepath{clip}%
\pgfsetbuttcap%
\pgfsetroundjoin%
\definecolor{currentfill}{rgb}{0.141176,0.345098,0.650980}%
\pgfsetfillcolor{currentfill}%
\pgfsetfillopacity{0.080000}%
\pgfsetlinewidth{0.000000pt}%
\definecolor{currentstroke}{rgb}{0.000000,0.000000,0.000000}%
\pgfsetstrokecolor{currentstroke}%
\pgfsetdash{}{0pt}%
\pgfpathmoveto{\pgfqpoint{1.419989in}{0.939715in}}%
\pgfusepath{fill}%
\end{pgfscope}%
\begin{pgfscope}%
\pgfpathrectangle{\pgfqpoint{0.407592in}{0.222904in}}{\pgfqpoint{1.411163in}{1.411163in}}%
\pgfusepath{clip}%
\pgfsetbuttcap%
\pgfsetroundjoin%
\definecolor{currentfill}{rgb}{0.141176,0.345098,0.650980}%
\pgfsetfillcolor{currentfill}%
\pgfsetfillopacity{0.080000}%
\pgfsetlinewidth{0.000000pt}%
\definecolor{currentstroke}{rgb}{0.000000,0.000000,0.000000}%
\pgfsetstrokecolor{currentstroke}%
\pgfsetdash{}{0pt}%
\pgfpathmoveto{\pgfqpoint{1.349507in}{1.000697in}}%
\pgfpathlineto{\pgfqpoint{1.420696in}{0.983877in}}%
\pgfpathlineto{\pgfqpoint{1.421407in}{1.028256in}}%
\pgfpathlineto{\pgfqpoint{1.350040in}{1.044884in}}%
\pgfpathlineto{\pgfqpoint{1.349507in}{1.000697in}}%
\pgfpathclose%
\pgfusepath{fill}%
\end{pgfscope}%
\begin{pgfscope}%
\pgfpathrectangle{\pgfqpoint{0.407592in}{0.222904in}}{\pgfqpoint{1.411163in}{1.411163in}}%
\pgfusepath{clip}%
\pgfsetbuttcap%
\pgfsetroundjoin%
\definecolor{currentfill}{rgb}{0.141176,0.345098,0.650980}%
\pgfsetfillcolor{currentfill}%
\pgfsetfillopacity{0.080000}%
\pgfsetlinewidth{0.000000pt}%
\definecolor{currentstroke}{rgb}{0.000000,0.000000,0.000000}%
\pgfsetstrokecolor{currentstroke}%
\pgfsetdash{}{0pt}%
\pgfpathmoveto{\pgfqpoint{1.419989in}{0.939715in}}%
\pgfpathlineto{\pgfqpoint{1.420696in}{0.983877in}}%
\pgfusepath{fill}%
\end{pgfscope}%
\begin{pgfscope}%
\pgfpathrectangle{\pgfqpoint{0.407592in}{0.222904in}}{\pgfqpoint{1.411163in}{1.411163in}}%
\pgfusepath{clip}%
\pgfsetbuttcap%
\pgfsetroundjoin%
\definecolor{currentfill}{rgb}{0.141176,0.345098,0.650980}%
\pgfsetfillcolor{currentfill}%
\pgfsetfillopacity{0.080000}%
\pgfsetlinewidth{0.000000pt}%
\definecolor{currentstroke}{rgb}{0.000000,0.000000,0.000000}%
\pgfsetstrokecolor{currentstroke}%
\pgfsetdash{}{0pt}%
\pgfpathmoveto{\pgfqpoint{1.350040in}{1.044884in}}%
\pgfpathlineto{\pgfqpoint{1.421407in}{1.028256in}}%
\pgfpathlineto{\pgfqpoint{1.422122in}{1.072854in}}%
\pgfpathlineto{\pgfqpoint{1.350575in}{1.089287in}}%
\pgfpathlineto{\pgfqpoint{1.350040in}{1.044884in}}%
\pgfpathclose%
\pgfusepath{fill}%
\end{pgfscope}%
\begin{pgfscope}%
\pgfpathrectangle{\pgfqpoint{0.407592in}{0.222904in}}{\pgfqpoint{1.411163in}{1.411163in}}%
\pgfusepath{clip}%
\pgfsetbuttcap%
\pgfsetroundjoin%
\definecolor{currentfill}{rgb}{0.141176,0.345098,0.650980}%
\pgfsetfillcolor{currentfill}%
\pgfsetfillopacity{0.080000}%
\pgfsetlinewidth{0.000000pt}%
\definecolor{currentstroke}{rgb}{0.000000,0.000000,0.000000}%
\pgfsetstrokecolor{currentstroke}%
\pgfsetdash{}{0pt}%
\pgfpathmoveto{\pgfqpoint{1.420696in}{0.983877in}}%
\pgfpathlineto{\pgfqpoint{1.493535in}{1.011451in}}%
\pgfpathlineto{\pgfqpoint{1.421407in}{1.028256in}}%
\pgfusepath{fill}%
\end{pgfscope}%
\begin{pgfscope}%
\pgfpathrectangle{\pgfqpoint{0.407592in}{0.222904in}}{\pgfqpoint{1.411163in}{1.411163in}}%
\pgfusepath{clip}%
\pgfsetbuttcap%
\pgfsetroundjoin%
\definecolor{currentfill}{rgb}{0.141176,0.345098,0.650980}%
\pgfsetfillcolor{currentfill}%
\pgfsetfillopacity{0.080000}%
\pgfsetlinewidth{0.000000pt}%
\definecolor{currentstroke}{rgb}{0.000000,0.000000,0.000000}%
\pgfsetstrokecolor{currentstroke}%
\pgfsetdash{}{0pt}%
\pgfpathmoveto{\pgfqpoint{1.421407in}{1.028256in}}%
\pgfpathlineto{\pgfqpoint{1.493535in}{1.011451in}}%
\pgfpathlineto{\pgfqpoint{1.494432in}{1.056245in}}%
\pgfpathlineto{\pgfqpoint{1.422122in}{1.072854in}}%
\pgfpathlineto{\pgfqpoint{1.421407in}{1.028256in}}%
\pgfpathclose%
\pgfusepath{fill}%
\end{pgfscope}%
\begin{pgfscope}%
\pgfpathrectangle{\pgfqpoint{0.407592in}{0.222904in}}{\pgfqpoint{1.411163in}{1.411163in}}%
\pgfusepath{clip}%
\pgfsetbuttcap%
\pgfsetroundjoin%
\definecolor{currentfill}{rgb}{0.168627,0.611765,0.411765}%
\pgfsetfillcolor{currentfill}%
\pgfsetfillopacity{0.180000}%
\pgfsetlinewidth{0.000000pt}%
\definecolor{currentstroke}{rgb}{0.000000,0.000000,0.000000}%
\pgfsetstrokecolor{currentstroke}%
\pgfsetdash{}{0pt}%
\pgfpathmoveto{\pgfqpoint{1.668468in}{1.154894in}}%
\pgfpathlineto{\pgfqpoint{1.682713in}{1.180372in}}%
\pgfpathlineto{\pgfqpoint{1.683370in}{1.202112in}}%
\pgfpathlineto{\pgfqpoint{1.669109in}{1.176649in}}%
\pgfpathlineto{\pgfqpoint{1.668468in}{1.154894in}}%
\pgfpathclose%
\pgfusepath{fill}%
\end{pgfscope}%
\begin{pgfscope}%
\pgfpathrectangle{\pgfqpoint{0.407592in}{0.222904in}}{\pgfqpoint{1.411163in}{1.411163in}}%
\pgfusepath{clip}%
\pgfsetbuttcap%
\pgfsetroundjoin%
\definecolor{currentfill}{rgb}{0.168627,0.611765,0.411765}%
\pgfsetfillcolor{currentfill}%
\pgfsetfillopacity{0.180000}%
\pgfsetlinewidth{0.000000pt}%
\definecolor{currentstroke}{rgb}{0.000000,0.000000,0.000000}%
\pgfsetstrokecolor{currentstroke}%
\pgfsetdash{}{0pt}%
\pgfpathmoveto{\pgfqpoint{1.669109in}{1.176649in}}%
\pgfpathlineto{\pgfqpoint{1.683370in}{1.202112in}}%
\pgfpathlineto{\pgfqpoint{1.684029in}{1.223905in}}%
\pgfpathlineto{\pgfqpoint{1.669753in}{1.198456in}}%
\pgfpathlineto{\pgfqpoint{1.669109in}{1.176649in}}%
\pgfpathclose%
\pgfusepath{fill}%
\end{pgfscope}%
\begin{pgfscope}%
\pgfpathrectangle{\pgfqpoint{0.407592in}{0.222904in}}{\pgfqpoint{1.411163in}{1.411163in}}%
\pgfusepath{clip}%
\pgfsetbuttcap%
\pgfsetroundjoin%
\definecolor{currentfill}{rgb}{0.168627,0.611765,0.411765}%
\pgfsetfillcolor{currentfill}%
\pgfsetfillopacity{0.180000}%
\pgfsetlinewidth{0.000000pt}%
\definecolor{currentstroke}{rgb}{0.000000,0.000000,0.000000}%
\pgfsetstrokecolor{currentstroke}%
\pgfsetdash{}{0pt}%
\pgfpathmoveto{\pgfqpoint{1.654183in}{1.130004in}}%
\pgfpathlineto{\pgfqpoint{1.668468in}{1.154894in}}%
\pgfpathlineto{\pgfqpoint{1.669109in}{1.176649in}}%
\pgfpathlineto{\pgfqpoint{1.654809in}{1.151775in}}%
\pgfpathlineto{\pgfqpoint{1.654183in}{1.130004in}}%
\pgfpathclose%
\pgfusepath{fill}%
\end{pgfscope}%
\begin{pgfscope}%
\pgfpathrectangle{\pgfqpoint{0.407592in}{0.222904in}}{\pgfqpoint{1.411163in}{1.411163in}}%
\pgfusepath{clip}%
\pgfsetbuttcap%
\pgfsetroundjoin%
\definecolor{currentfill}{rgb}{0.168627,0.611765,0.411765}%
\pgfsetfillcolor{currentfill}%
\pgfsetfillopacity{0.180000}%
\pgfsetlinewidth{0.000000pt}%
\definecolor{currentstroke}{rgb}{0.000000,0.000000,0.000000}%
\pgfsetstrokecolor{currentstroke}%
\pgfsetdash{}{0pt}%
\pgfpathmoveto{\pgfqpoint{1.669753in}{1.198456in}}%
\pgfpathlineto{\pgfqpoint{1.684029in}{1.223905in}}%
\pgfpathlineto{\pgfqpoint{1.684690in}{1.245750in}}%
\pgfpathlineto{\pgfqpoint{1.670397in}{1.220316in}}%
\pgfpathlineto{\pgfqpoint{1.669753in}{1.198456in}}%
\pgfpathclose%
\pgfusepath{fill}%
\end{pgfscope}%
\begin{pgfscope}%
\pgfpathrectangle{\pgfqpoint{0.407592in}{0.222904in}}{\pgfqpoint{1.411163in}{1.411163in}}%
\pgfusepath{clip}%
\pgfsetbuttcap%
\pgfsetroundjoin%
\definecolor{currentfill}{rgb}{0.168627,0.611765,0.411765}%
\pgfsetfillcolor{currentfill}%
\pgfsetfillopacity{0.180000}%
\pgfsetlinewidth{0.000000pt}%
\definecolor{currentstroke}{rgb}{0.000000,0.000000,0.000000}%
\pgfsetstrokecolor{currentstroke}%
\pgfsetdash{}{0pt}%
\pgfpathmoveto{\pgfqpoint{1.654809in}{1.151775in}}%
\pgfpathlineto{\pgfqpoint{1.669109in}{1.176649in}}%
\pgfpathlineto{\pgfqpoint{1.669753in}{1.198456in}}%
\pgfpathlineto{\pgfqpoint{1.655437in}{1.173599in}}%
\pgfpathlineto{\pgfqpoint{1.654809in}{1.151775in}}%
\pgfpathclose%
\pgfusepath{fill}%
\end{pgfscope}%
\begin{pgfscope}%
\pgfpathrectangle{\pgfqpoint{0.407592in}{0.222904in}}{\pgfqpoint{1.411163in}{1.411163in}}%
\pgfusepath{clip}%
\pgfsetbuttcap%
\pgfsetroundjoin%
\definecolor{currentfill}{rgb}{0.168627,0.611765,0.411765}%
\pgfsetfillcolor{currentfill}%
\pgfsetfillopacity{0.180000}%
\pgfsetlinewidth{0.000000pt}%
\definecolor{currentstroke}{rgb}{0.000000,0.000000,0.000000}%
\pgfsetstrokecolor{currentstroke}%
\pgfsetdash{}{0pt}%
\pgfpathmoveto{\pgfqpoint{1.670397in}{1.220316in}}%
\pgfpathlineto{\pgfqpoint{1.684690in}{1.245750in}}%
\pgfpathlineto{\pgfqpoint{1.685352in}{1.267647in}}%
\pgfpathlineto{\pgfqpoint{1.671044in}{1.242228in}}%
\pgfpathlineto{\pgfqpoint{1.670397in}{1.220316in}}%
\pgfpathclose%
\pgfusepath{fill}%
\end{pgfscope}%
\begin{pgfscope}%
\pgfpathrectangle{\pgfqpoint{0.407592in}{0.222904in}}{\pgfqpoint{1.411163in}{1.411163in}}%
\pgfusepath{clip}%
\pgfsetbuttcap%
\pgfsetroundjoin%
\definecolor{currentfill}{rgb}{0.168627,0.611765,0.411765}%
\pgfsetfillcolor{currentfill}%
\pgfsetfillopacity{0.180000}%
\pgfsetlinewidth{0.000000pt}%
\definecolor{currentstroke}{rgb}{0.000000,0.000000,0.000000}%
\pgfsetstrokecolor{currentstroke}%
\pgfsetdash{}{0pt}%
\pgfpathmoveto{\pgfqpoint{1.639860in}{1.105771in}}%
\pgfpathlineto{\pgfqpoint{1.654183in}{1.130004in}}%
\pgfpathlineto{\pgfqpoint{1.654809in}{1.151775in}}%
\pgfpathlineto{\pgfqpoint{1.640470in}{1.127559in}}%
\pgfpathlineto{\pgfqpoint{1.639860in}{1.105771in}}%
\pgfpathclose%
\pgfusepath{fill}%
\end{pgfscope}%
\begin{pgfscope}%
\pgfpathrectangle{\pgfqpoint{0.407592in}{0.222904in}}{\pgfqpoint{1.411163in}{1.411163in}}%
\pgfusepath{clip}%
\pgfsetbuttcap%
\pgfsetroundjoin%
\definecolor{currentfill}{rgb}{0.168627,0.611765,0.411765}%
\pgfsetfillcolor{currentfill}%
\pgfsetfillopacity{0.180000}%
\pgfsetlinewidth{0.000000pt}%
\definecolor{currentstroke}{rgb}{0.000000,0.000000,0.000000}%
\pgfsetstrokecolor{currentstroke}%
\pgfsetdash{}{0pt}%
\pgfpathmoveto{\pgfqpoint{1.655437in}{1.173599in}}%
\pgfpathlineto{\pgfqpoint{1.669753in}{1.198456in}}%
\pgfpathlineto{\pgfqpoint{1.670397in}{1.220316in}}%
\pgfpathlineto{\pgfqpoint{1.656066in}{1.195474in}}%
\pgfpathlineto{\pgfqpoint{1.655437in}{1.173599in}}%
\pgfpathclose%
\pgfusepath{fill}%
\end{pgfscope}%
\begin{pgfscope}%
\pgfpathrectangle{\pgfqpoint{0.407592in}{0.222904in}}{\pgfqpoint{1.411163in}{1.411163in}}%
\pgfusepath{clip}%
\pgfsetbuttcap%
\pgfsetroundjoin%
\definecolor{currentfill}{rgb}{0.168627,0.611765,0.411765}%
\pgfsetfillcolor{currentfill}%
\pgfsetfillopacity{0.180000}%
\pgfsetlinewidth{0.000000pt}%
\definecolor{currentstroke}{rgb}{0.000000,0.000000,0.000000}%
\pgfsetstrokecolor{currentstroke}%
\pgfsetdash{}{0pt}%
\pgfpathmoveto{\pgfqpoint{1.671044in}{1.242228in}}%
\pgfpathlineto{\pgfqpoint{1.685352in}{1.267647in}}%
\pgfpathlineto{\pgfqpoint{1.686016in}{1.289597in}}%
\pgfpathlineto{\pgfqpoint{1.671692in}{1.264192in}}%
\pgfpathlineto{\pgfqpoint{1.671044in}{1.242228in}}%
\pgfpathclose%
\pgfusepath{fill}%
\end{pgfscope}%
\begin{pgfscope}%
\pgfpathrectangle{\pgfqpoint{0.407592in}{0.222904in}}{\pgfqpoint{1.411163in}{1.411163in}}%
\pgfusepath{clip}%
\pgfsetbuttcap%
\pgfsetroundjoin%
\definecolor{currentfill}{rgb}{0.168627,0.611765,0.411765}%
\pgfsetfillcolor{currentfill}%
\pgfsetfillopacity{0.180000}%
\pgfsetlinewidth{0.000000pt}%
\definecolor{currentstroke}{rgb}{0.000000,0.000000,0.000000}%
\pgfsetstrokecolor{currentstroke}%
\pgfsetdash{}{0pt}%
\pgfpathmoveto{\pgfqpoint{1.640470in}{1.127559in}}%
\pgfpathlineto{\pgfqpoint{1.654809in}{1.151775in}}%
\pgfpathlineto{\pgfqpoint{1.655437in}{1.173599in}}%
\pgfpathlineto{\pgfqpoint{1.641082in}{1.149401in}}%
\pgfpathlineto{\pgfqpoint{1.640470in}{1.127559in}}%
\pgfpathclose%
\pgfusepath{fill}%
\end{pgfscope}%
\begin{pgfscope}%
\pgfpathrectangle{\pgfqpoint{0.407592in}{0.222904in}}{\pgfqpoint{1.411163in}{1.411163in}}%
\pgfusepath{clip}%
\pgfsetbuttcap%
\pgfsetroundjoin%
\definecolor{currentfill}{rgb}{0.168627,0.611765,0.411765}%
\pgfsetfillcolor{currentfill}%
\pgfsetfillopacity{0.180000}%
\pgfsetlinewidth{0.000000pt}%
\definecolor{currentstroke}{rgb}{0.000000,0.000000,0.000000}%
\pgfsetstrokecolor{currentstroke}%
\pgfsetdash{}{0pt}%
\pgfpathmoveto{\pgfqpoint{1.656066in}{1.195474in}}%
\pgfpathlineto{\pgfqpoint{1.670397in}{1.220316in}}%
\pgfpathlineto{\pgfqpoint{1.671044in}{1.242228in}}%
\pgfpathlineto{\pgfqpoint{1.656696in}{1.217403in}}%
\pgfpathlineto{\pgfqpoint{1.656066in}{1.195474in}}%
\pgfpathclose%
\pgfusepath{fill}%
\end{pgfscope}%
\begin{pgfscope}%
\pgfpathrectangle{\pgfqpoint{0.407592in}{0.222904in}}{\pgfqpoint{1.411163in}{1.411163in}}%
\pgfusepath{clip}%
\pgfsetbuttcap%
\pgfsetroundjoin%
\definecolor{currentfill}{rgb}{0.168627,0.611765,0.411765}%
\pgfsetfillcolor{currentfill}%
\pgfsetfillopacity{0.180000}%
\pgfsetlinewidth{0.000000pt}%
\definecolor{currentstroke}{rgb}{0.000000,0.000000,0.000000}%
\pgfsetstrokecolor{currentstroke}%
\pgfsetdash{}{0pt}%
\pgfpathmoveto{\pgfqpoint{1.671692in}{1.264192in}}%
\pgfpathlineto{\pgfqpoint{1.686016in}{1.289597in}}%
\pgfpathlineto{\pgfqpoint{1.686681in}{1.311599in}}%
\pgfpathlineto{\pgfqpoint{1.672341in}{1.286210in}}%
\pgfpathlineto{\pgfqpoint{1.671692in}{1.264192in}}%
\pgfpathclose%
\pgfusepath{fill}%
\end{pgfscope}%
\begin{pgfscope}%
\pgfpathrectangle{\pgfqpoint{0.407592in}{0.222904in}}{\pgfqpoint{1.411163in}{1.411163in}}%
\pgfusepath{clip}%
\pgfsetbuttcap%
\pgfsetroundjoin%
\definecolor{currentfill}{rgb}{0.168627,0.611765,0.411765}%
\pgfsetfillcolor{currentfill}%
\pgfsetfillopacity{0.180000}%
\pgfsetlinewidth{0.000000pt}%
\definecolor{currentstroke}{rgb}{0.000000,0.000000,0.000000}%
\pgfsetstrokecolor{currentstroke}%
\pgfsetdash{}{0pt}%
\pgfpathmoveto{\pgfqpoint{1.625498in}{1.082268in}}%
\pgfpathlineto{\pgfqpoint{1.639860in}{1.105771in}}%
\pgfpathlineto{\pgfqpoint{1.640470in}{1.127559in}}%
\pgfpathlineto{\pgfqpoint{1.626092in}{1.104076in}}%
\pgfpathlineto{\pgfqpoint{1.625498in}{1.082268in}}%
\pgfpathclose%
\pgfusepath{fill}%
\end{pgfscope}%
\begin{pgfscope}%
\pgfpathrectangle{\pgfqpoint{0.407592in}{0.222904in}}{\pgfqpoint{1.411163in}{1.411163in}}%
\pgfusepath{clip}%
\pgfsetbuttcap%
\pgfsetroundjoin%
\definecolor{currentfill}{rgb}{0.168627,0.611765,0.411765}%
\pgfsetfillcolor{currentfill}%
\pgfsetfillopacity{0.180000}%
\pgfsetlinewidth{0.000000pt}%
\definecolor{currentstroke}{rgb}{0.000000,0.000000,0.000000}%
\pgfsetstrokecolor{currentstroke}%
\pgfsetdash{}{0pt}%
\pgfpathmoveto{\pgfqpoint{1.641082in}{1.149401in}}%
\pgfpathlineto{\pgfqpoint{1.655437in}{1.173599in}}%
\pgfpathlineto{\pgfqpoint{1.656066in}{1.195474in}}%
\pgfpathlineto{\pgfqpoint{1.641695in}{1.171294in}}%
\pgfpathlineto{\pgfqpoint{1.641082in}{1.149401in}}%
\pgfpathclose%
\pgfusepath{fill}%
\end{pgfscope}%
\begin{pgfscope}%
\pgfpathrectangle{\pgfqpoint{0.407592in}{0.222904in}}{\pgfqpoint{1.411163in}{1.411163in}}%
\pgfusepath{clip}%
\pgfsetbuttcap%
\pgfsetroundjoin%
\definecolor{currentfill}{rgb}{0.168627,0.611765,0.411765}%
\pgfsetfillcolor{currentfill}%
\pgfsetfillopacity{0.180000}%
\pgfsetlinewidth{0.000000pt}%
\definecolor{currentstroke}{rgb}{0.000000,0.000000,0.000000}%
\pgfsetstrokecolor{currentstroke}%
\pgfsetdash{}{0pt}%
\pgfpathmoveto{\pgfqpoint{1.656696in}{1.217403in}}%
\pgfpathlineto{\pgfqpoint{1.671044in}{1.242228in}}%
\pgfpathlineto{\pgfqpoint{1.671692in}{1.264192in}}%
\pgfpathlineto{\pgfqpoint{1.657328in}{1.239384in}}%
\pgfpathlineto{\pgfqpoint{1.656696in}{1.217403in}}%
\pgfpathclose%
\pgfusepath{fill}%
\end{pgfscope}%
\begin{pgfscope}%
\pgfpathrectangle{\pgfqpoint{0.407592in}{0.222904in}}{\pgfqpoint{1.411163in}{1.411163in}}%
\pgfusepath{clip}%
\pgfsetbuttcap%
\pgfsetroundjoin%
\definecolor{currentfill}{rgb}{0.168627,0.611765,0.411765}%
\pgfsetfillcolor{currentfill}%
\pgfsetfillopacity{0.180000}%
\pgfsetlinewidth{0.000000pt}%
\definecolor{currentstroke}{rgb}{0.000000,0.000000,0.000000}%
\pgfsetstrokecolor{currentstroke}%
\pgfsetdash{}{0pt}%
\pgfpathmoveto{\pgfqpoint{1.672341in}{1.286210in}}%
\pgfpathlineto{\pgfqpoint{1.686681in}{1.311599in}}%
\pgfpathlineto{\pgfqpoint{1.687348in}{1.333655in}}%
\pgfpathlineto{\pgfqpoint{1.672992in}{1.308281in}}%
\pgfpathlineto{\pgfqpoint{1.672341in}{1.286210in}}%
\pgfpathclose%
\pgfusepath{fill}%
\end{pgfscope}%
\begin{pgfscope}%
\pgfpathrectangle{\pgfqpoint{0.407592in}{0.222904in}}{\pgfqpoint{1.411163in}{1.411163in}}%
\pgfusepath{clip}%
\pgfsetbuttcap%
\pgfsetroundjoin%
\definecolor{currentfill}{rgb}{0.168627,0.611765,0.411765}%
\pgfsetfillcolor{currentfill}%
\pgfsetfillopacity{0.180000}%
\pgfsetlinewidth{0.000000pt}%
\definecolor{currentstroke}{rgb}{0.000000,0.000000,0.000000}%
\pgfsetstrokecolor{currentstroke}%
\pgfsetdash{}{0pt}%
\pgfpathmoveto{\pgfqpoint{1.626092in}{1.104076in}}%
\pgfpathlineto{\pgfqpoint{1.640470in}{1.127559in}}%
\pgfpathlineto{\pgfqpoint{1.641082in}{1.149401in}}%
\pgfpathlineto{\pgfqpoint{1.626687in}{1.125937in}}%
\pgfpathlineto{\pgfqpoint{1.626092in}{1.104076in}}%
\pgfpathclose%
\pgfusepath{fill}%
\end{pgfscope}%
\begin{pgfscope}%
\pgfpathrectangle{\pgfqpoint{0.407592in}{0.222904in}}{\pgfqpoint{1.411163in}{1.411163in}}%
\pgfusepath{clip}%
\pgfsetbuttcap%
\pgfsetroundjoin%
\definecolor{currentfill}{rgb}{0.168627,0.611765,0.411765}%
\pgfsetfillcolor{currentfill}%
\pgfsetfillopacity{0.180000}%
\pgfsetlinewidth{0.000000pt}%
\definecolor{currentstroke}{rgb}{0.000000,0.000000,0.000000}%
\pgfsetstrokecolor{currentstroke}%
\pgfsetdash{}{0pt}%
\pgfpathmoveto{\pgfqpoint{1.641695in}{1.171294in}}%
\pgfpathlineto{\pgfqpoint{1.656066in}{1.195474in}}%
\pgfpathlineto{\pgfqpoint{1.656696in}{1.217403in}}%
\pgfpathlineto{\pgfqpoint{1.642309in}{1.193241in}}%
\pgfpathlineto{\pgfqpoint{1.641695in}{1.171294in}}%
\pgfpathclose%
\pgfusepath{fill}%
\end{pgfscope}%
\begin{pgfscope}%
\pgfpathrectangle{\pgfqpoint{0.407592in}{0.222904in}}{\pgfqpoint{1.411163in}{1.411163in}}%
\pgfusepath{clip}%
\pgfsetbuttcap%
\pgfsetroundjoin%
\definecolor{currentfill}{rgb}{0.168627,0.611765,0.411765}%
\pgfsetfillcolor{currentfill}%
\pgfsetfillopacity{0.180000}%
\pgfsetlinewidth{0.000000pt}%
\definecolor{currentstroke}{rgb}{0.000000,0.000000,0.000000}%
\pgfsetstrokecolor{currentstroke}%
\pgfsetdash{}{0pt}%
\pgfpathmoveto{\pgfqpoint{1.657328in}{1.239384in}}%
\pgfpathlineto{\pgfqpoint{1.671692in}{1.264192in}}%
\pgfpathlineto{\pgfqpoint{1.672341in}{1.286210in}}%
\pgfpathlineto{\pgfqpoint{1.657962in}{1.261418in}}%
\pgfpathlineto{\pgfqpoint{1.657328in}{1.239384in}}%
\pgfpathclose%
\pgfusepath{fill}%
\end{pgfscope}%
\begin{pgfscope}%
\pgfpathrectangle{\pgfqpoint{0.407592in}{0.222904in}}{\pgfqpoint{1.411163in}{1.411163in}}%
\pgfusepath{clip}%
\pgfsetbuttcap%
\pgfsetroundjoin%
\definecolor{currentfill}{rgb}{0.168627,0.611765,0.411765}%
\pgfsetfillcolor{currentfill}%
\pgfsetfillopacity{0.180000}%
\pgfsetlinewidth{0.000000pt}%
\definecolor{currentstroke}{rgb}{0.000000,0.000000,0.000000}%
\pgfsetstrokecolor{currentstroke}%
\pgfsetdash{}{0pt}%
\pgfpathmoveto{\pgfqpoint{1.672992in}{1.308281in}}%
\pgfpathlineto{\pgfqpoint{1.687348in}{1.333655in}}%
\pgfpathlineto{\pgfqpoint{1.688017in}{1.355763in}}%
\pgfpathlineto{\pgfqpoint{1.673645in}{1.330405in}}%
\pgfpathlineto{\pgfqpoint{1.672992in}{1.308281in}}%
\pgfpathclose%
\pgfusepath{fill}%
\end{pgfscope}%
\begin{pgfscope}%
\pgfpathrectangle{\pgfqpoint{0.407592in}{0.222904in}}{\pgfqpoint{1.411163in}{1.411163in}}%
\pgfusepath{clip}%
\pgfsetbuttcap%
\pgfsetroundjoin%
\definecolor{currentfill}{rgb}{0.168627,0.611765,0.411765}%
\pgfsetfillcolor{currentfill}%
\pgfsetfillopacity{0.180000}%
\pgfsetlinewidth{0.000000pt}%
\definecolor{currentstroke}{rgb}{0.000000,0.000000,0.000000}%
\pgfsetstrokecolor{currentstroke}%
\pgfsetdash{}{0pt}%
\pgfpathmoveto{\pgfqpoint{1.611096in}{1.059575in}}%
\pgfpathlineto{\pgfqpoint{1.625498in}{1.082268in}}%
\pgfpathlineto{\pgfqpoint{1.626092in}{1.104076in}}%
\pgfpathlineto{\pgfqpoint{1.611674in}{1.081405in}}%
\pgfpathlineto{\pgfqpoint{1.611096in}{1.059575in}}%
\pgfpathclose%
\pgfusepath{fill}%
\end{pgfscope}%
\begin{pgfscope}%
\pgfpathrectangle{\pgfqpoint{0.407592in}{0.222904in}}{\pgfqpoint{1.411163in}{1.411163in}}%
\pgfusepath{clip}%
\pgfsetbuttcap%
\pgfsetroundjoin%
\definecolor{currentfill}{rgb}{0.168627,0.611765,0.411765}%
\pgfsetfillcolor{currentfill}%
\pgfsetfillopacity{0.180000}%
\pgfsetlinewidth{0.000000pt}%
\definecolor{currentstroke}{rgb}{0.000000,0.000000,0.000000}%
\pgfsetstrokecolor{currentstroke}%
\pgfsetdash{}{0pt}%
\pgfpathmoveto{\pgfqpoint{1.626687in}{1.125937in}}%
\pgfpathlineto{\pgfqpoint{1.641082in}{1.149401in}}%
\pgfpathlineto{\pgfqpoint{1.641695in}{1.171294in}}%
\pgfpathlineto{\pgfqpoint{1.627284in}{1.147851in}}%
\pgfpathlineto{\pgfqpoint{1.626687in}{1.125937in}}%
\pgfpathclose%
\pgfusepath{fill}%
\end{pgfscope}%
\begin{pgfscope}%
\pgfpathrectangle{\pgfqpoint{0.407592in}{0.222904in}}{\pgfqpoint{1.411163in}{1.411163in}}%
\pgfusepath{clip}%
\pgfsetbuttcap%
\pgfsetroundjoin%
\definecolor{currentfill}{rgb}{0.168627,0.611765,0.411765}%
\pgfsetfillcolor{currentfill}%
\pgfsetfillopacity{0.180000}%
\pgfsetlinewidth{0.000000pt}%
\definecolor{currentstroke}{rgb}{0.000000,0.000000,0.000000}%
\pgfsetstrokecolor{currentstroke}%
\pgfsetdash{}{0pt}%
\pgfpathmoveto{\pgfqpoint{1.642309in}{1.193241in}}%
\pgfpathlineto{\pgfqpoint{1.656696in}{1.217403in}}%
\pgfpathlineto{\pgfqpoint{1.657328in}{1.239384in}}%
\pgfpathlineto{\pgfqpoint{1.642925in}{1.215241in}}%
\pgfpathlineto{\pgfqpoint{1.642309in}{1.193241in}}%
\pgfpathclose%
\pgfusepath{fill}%
\end{pgfscope}%
\begin{pgfscope}%
\pgfpathrectangle{\pgfqpoint{0.407592in}{0.222904in}}{\pgfqpoint{1.411163in}{1.411163in}}%
\pgfusepath{clip}%
\pgfsetbuttcap%
\pgfsetroundjoin%
\definecolor{currentfill}{rgb}{0.168627,0.611765,0.411765}%
\pgfsetfillcolor{currentfill}%
\pgfsetfillopacity{0.180000}%
\pgfsetlinewidth{0.000000pt}%
\definecolor{currentstroke}{rgb}{0.000000,0.000000,0.000000}%
\pgfsetstrokecolor{currentstroke}%
\pgfsetdash{}{0pt}%
\pgfpathmoveto{\pgfqpoint{1.657962in}{1.261418in}}%
\pgfpathlineto{\pgfqpoint{1.672341in}{1.286210in}}%
\pgfpathlineto{\pgfqpoint{1.672992in}{1.308281in}}%
\pgfpathlineto{\pgfqpoint{1.658597in}{1.283506in}}%
\pgfpathlineto{\pgfqpoint{1.657962in}{1.261418in}}%
\pgfpathclose%
\pgfusepath{fill}%
\end{pgfscope}%
\begin{pgfscope}%
\pgfpathrectangle{\pgfqpoint{0.407592in}{0.222904in}}{\pgfqpoint{1.411163in}{1.411163in}}%
\pgfusepath{clip}%
\pgfsetbuttcap%
\pgfsetroundjoin%
\definecolor{currentfill}{rgb}{0.168627,0.611765,0.411765}%
\pgfsetfillcolor{currentfill}%
\pgfsetfillopacity{0.180000}%
\pgfsetlinewidth{0.000000pt}%
\definecolor{currentstroke}{rgb}{0.000000,0.000000,0.000000}%
\pgfsetstrokecolor{currentstroke}%
\pgfsetdash{}{0pt}%
\pgfpathmoveto{\pgfqpoint{1.673645in}{1.330405in}}%
\pgfpathlineto{\pgfqpoint{1.688017in}{1.355763in}}%
\pgfpathlineto{\pgfqpoint{1.688687in}{1.377925in}}%
\pgfpathlineto{\pgfqpoint{1.674299in}{1.352582in}}%
\pgfpathlineto{\pgfqpoint{1.673645in}{1.330405in}}%
\pgfpathclose%
\pgfusepath{fill}%
\end{pgfscope}%
\begin{pgfscope}%
\pgfpathrectangle{\pgfqpoint{0.407592in}{0.222904in}}{\pgfqpoint{1.411163in}{1.411163in}}%
\pgfusepath{clip}%
\pgfsetbuttcap%
\pgfsetroundjoin%
\definecolor{currentfill}{rgb}{0.168627,0.611765,0.411765}%
\pgfsetfillcolor{currentfill}%
\pgfsetfillopacity{0.180000}%
\pgfsetlinewidth{0.000000pt}%
\definecolor{currentstroke}{rgb}{0.000000,0.000000,0.000000}%
\pgfsetstrokecolor{currentstroke}%
\pgfsetdash{}{0pt}%
\pgfpathmoveto{\pgfqpoint{1.611674in}{1.081405in}}%
\pgfpathlineto{\pgfqpoint{1.626092in}{1.104076in}}%
\pgfpathlineto{\pgfqpoint{1.626687in}{1.125937in}}%
\pgfpathlineto{\pgfqpoint{1.612254in}{1.103288in}}%
\pgfpathlineto{\pgfqpoint{1.611674in}{1.081405in}}%
\pgfpathclose%
\pgfusepath{fill}%
\end{pgfscope}%
\begin{pgfscope}%
\pgfpathrectangle{\pgfqpoint{0.407592in}{0.222904in}}{\pgfqpoint{1.411163in}{1.411163in}}%
\pgfusepath{clip}%
\pgfsetbuttcap%
\pgfsetroundjoin%
\definecolor{currentfill}{rgb}{0.168627,0.611765,0.411765}%
\pgfsetfillcolor{currentfill}%
\pgfsetfillopacity{0.180000}%
\pgfsetlinewidth{0.000000pt}%
\definecolor{currentstroke}{rgb}{0.000000,0.000000,0.000000}%
\pgfsetstrokecolor{currentstroke}%
\pgfsetdash{}{0pt}%
\pgfpathmoveto{\pgfqpoint{1.627284in}{1.147851in}}%
\pgfpathlineto{\pgfqpoint{1.641695in}{1.171294in}}%
\pgfpathlineto{\pgfqpoint{1.642309in}{1.193241in}}%
\pgfpathlineto{\pgfqpoint{1.627883in}{1.169818in}}%
\pgfpathlineto{\pgfqpoint{1.627284in}{1.147851in}}%
\pgfpathclose%
\pgfusepath{fill}%
\end{pgfscope}%
\begin{pgfscope}%
\pgfpathrectangle{\pgfqpoint{0.407592in}{0.222904in}}{\pgfqpoint{1.411163in}{1.411163in}}%
\pgfusepath{clip}%
\pgfsetbuttcap%
\pgfsetroundjoin%
\definecolor{currentfill}{rgb}{0.168627,0.611765,0.411765}%
\pgfsetfillcolor{currentfill}%
\pgfsetfillopacity{0.180000}%
\pgfsetlinewidth{0.000000pt}%
\definecolor{currentstroke}{rgb}{0.000000,0.000000,0.000000}%
\pgfsetstrokecolor{currentstroke}%
\pgfsetdash{}{0pt}%
\pgfpathmoveto{\pgfqpoint{1.642925in}{1.215241in}}%
\pgfpathlineto{\pgfqpoint{1.657328in}{1.239384in}}%
\pgfpathlineto{\pgfqpoint{1.657962in}{1.261418in}}%
\pgfpathlineto{\pgfqpoint{1.643543in}{1.237293in}}%
\pgfpathlineto{\pgfqpoint{1.642925in}{1.215241in}}%
\pgfpathclose%
\pgfusepath{fill}%
\end{pgfscope}%
\begin{pgfscope}%
\pgfpathrectangle{\pgfqpoint{0.407592in}{0.222904in}}{\pgfqpoint{1.411163in}{1.411163in}}%
\pgfusepath{clip}%
\pgfsetbuttcap%
\pgfsetroundjoin%
\definecolor{currentfill}{rgb}{0.168627,0.611765,0.411765}%
\pgfsetfillcolor{currentfill}%
\pgfsetfillopacity{0.180000}%
\pgfsetlinewidth{0.000000pt}%
\definecolor{currentstroke}{rgb}{0.000000,0.000000,0.000000}%
\pgfsetstrokecolor{currentstroke}%
\pgfsetdash{}{0pt}%
\pgfpathmoveto{\pgfqpoint{1.658597in}{1.283506in}}%
\pgfpathlineto{\pgfqpoint{1.672992in}{1.308281in}}%
\pgfpathlineto{\pgfqpoint{1.673645in}{1.330405in}}%
\pgfpathlineto{\pgfqpoint{1.659233in}{1.305647in}}%
\pgfpathlineto{\pgfqpoint{1.658597in}{1.283506in}}%
\pgfpathclose%
\pgfusepath{fill}%
\end{pgfscope}%
\begin{pgfscope}%
\pgfpathrectangle{\pgfqpoint{0.407592in}{0.222904in}}{\pgfqpoint{1.411163in}{1.411163in}}%
\pgfusepath{clip}%
\pgfsetbuttcap%
\pgfsetroundjoin%
\definecolor{currentfill}{rgb}{0.168627,0.611765,0.411765}%
\pgfsetfillcolor{currentfill}%
\pgfsetfillopacity{0.180000}%
\pgfsetlinewidth{0.000000pt}%
\definecolor{currentstroke}{rgb}{0.000000,0.000000,0.000000}%
\pgfsetstrokecolor{currentstroke}%
\pgfsetdash{}{0pt}%
\pgfpathmoveto{\pgfqpoint{1.674299in}{1.352582in}}%
\pgfpathlineto{\pgfqpoint{1.688687in}{1.377925in}}%
\pgfpathlineto{\pgfqpoint{1.689359in}{1.400140in}}%
\pgfpathlineto{\pgfqpoint{1.674955in}{1.374813in}}%
\pgfpathlineto{\pgfqpoint{1.674299in}{1.352582in}}%
\pgfpathclose%
\pgfusepath{fill}%
\end{pgfscope}%
\begin{pgfscope}%
\pgfpathrectangle{\pgfqpoint{0.407592in}{0.222904in}}{\pgfqpoint{1.411163in}{1.411163in}}%
\pgfusepath{clip}%
\pgfsetbuttcap%
\pgfsetroundjoin%
\definecolor{currentfill}{rgb}{0.168627,0.611765,0.411765}%
\pgfsetfillcolor{currentfill}%
\pgfsetfillopacity{0.180000}%
\pgfsetlinewidth{0.000000pt}%
\definecolor{currentstroke}{rgb}{0.000000,0.000000,0.000000}%
\pgfsetstrokecolor{currentstroke}%
\pgfsetdash{}{0pt}%
\pgfpathmoveto{\pgfqpoint{1.596655in}{1.037776in}}%
\pgfpathlineto{\pgfqpoint{1.611096in}{1.059575in}}%
\pgfpathlineto{\pgfqpoint{1.611674in}{1.081405in}}%
\pgfpathlineto{\pgfqpoint{1.597217in}{1.059631in}}%
\pgfpathlineto{\pgfqpoint{1.596655in}{1.037776in}}%
\pgfpathclose%
\pgfusepath{fill}%
\end{pgfscope}%
\begin{pgfscope}%
\pgfpathrectangle{\pgfqpoint{0.407592in}{0.222904in}}{\pgfqpoint{1.411163in}{1.411163in}}%
\pgfusepath{clip}%
\pgfsetbuttcap%
\pgfsetroundjoin%
\definecolor{currentfill}{rgb}{0.168627,0.611765,0.411765}%
\pgfsetfillcolor{currentfill}%
\pgfsetfillopacity{0.180000}%
\pgfsetlinewidth{0.000000pt}%
\definecolor{currentstroke}{rgb}{0.000000,0.000000,0.000000}%
\pgfsetstrokecolor{currentstroke}%
\pgfsetdash{}{0pt}%
\pgfpathmoveto{\pgfqpoint{1.612254in}{1.103288in}}%
\pgfpathlineto{\pgfqpoint{1.626687in}{1.125937in}}%
\pgfpathlineto{\pgfqpoint{1.627284in}{1.147851in}}%
\pgfpathlineto{\pgfqpoint{1.612835in}{1.125224in}}%
\pgfpathlineto{\pgfqpoint{1.612254in}{1.103288in}}%
\pgfpathclose%
\pgfusepath{fill}%
\end{pgfscope}%
\begin{pgfscope}%
\pgfpathrectangle{\pgfqpoint{0.407592in}{0.222904in}}{\pgfqpoint{1.411163in}{1.411163in}}%
\pgfusepath{clip}%
\pgfsetbuttcap%
\pgfsetroundjoin%
\definecolor{currentfill}{rgb}{0.168627,0.611765,0.411765}%
\pgfsetfillcolor{currentfill}%
\pgfsetfillopacity{0.180000}%
\pgfsetlinewidth{0.000000pt}%
\definecolor{currentstroke}{rgb}{0.000000,0.000000,0.000000}%
\pgfsetstrokecolor{currentstroke}%
\pgfsetdash{}{0pt}%
\pgfpathmoveto{\pgfqpoint{1.627883in}{1.169818in}}%
\pgfpathlineto{\pgfqpoint{1.642309in}{1.193241in}}%
\pgfpathlineto{\pgfqpoint{1.642925in}{1.215241in}}%
\pgfpathlineto{\pgfqpoint{1.628483in}{1.191838in}}%
\pgfpathlineto{\pgfqpoint{1.627883in}{1.169818in}}%
\pgfpathclose%
\pgfusepath{fill}%
\end{pgfscope}%
\begin{pgfscope}%
\pgfpathrectangle{\pgfqpoint{0.407592in}{0.222904in}}{\pgfqpoint{1.411163in}{1.411163in}}%
\pgfusepath{clip}%
\pgfsetbuttcap%
\pgfsetroundjoin%
\definecolor{currentfill}{rgb}{0.168627,0.611765,0.411765}%
\pgfsetfillcolor{currentfill}%
\pgfsetfillopacity{0.180000}%
\pgfsetlinewidth{0.000000pt}%
\definecolor{currentstroke}{rgb}{0.000000,0.000000,0.000000}%
\pgfsetstrokecolor{currentstroke}%
\pgfsetdash{}{0pt}%
\pgfpathmoveto{\pgfqpoint{1.643543in}{1.237293in}}%
\pgfpathlineto{\pgfqpoint{1.657962in}{1.261418in}}%
\pgfpathlineto{\pgfqpoint{1.658597in}{1.283506in}}%
\pgfpathlineto{\pgfqpoint{1.644162in}{1.259399in}}%
\pgfpathlineto{\pgfqpoint{1.643543in}{1.237293in}}%
\pgfpathclose%
\pgfusepath{fill}%
\end{pgfscope}%
\begin{pgfscope}%
\pgfpathrectangle{\pgfqpoint{0.407592in}{0.222904in}}{\pgfqpoint{1.411163in}{1.411163in}}%
\pgfusepath{clip}%
\pgfsetbuttcap%
\pgfsetroundjoin%
\definecolor{currentfill}{rgb}{0.168627,0.611765,0.411765}%
\pgfsetfillcolor{currentfill}%
\pgfsetfillopacity{0.180000}%
\pgfsetlinewidth{0.000000pt}%
\definecolor{currentstroke}{rgb}{0.000000,0.000000,0.000000}%
\pgfsetstrokecolor{currentstroke}%
\pgfsetdash{}{0pt}%
\pgfpathmoveto{\pgfqpoint{1.659233in}{1.305647in}}%
\pgfpathlineto{\pgfqpoint{1.673645in}{1.330405in}}%
\pgfpathlineto{\pgfqpoint{1.674299in}{1.352582in}}%
\pgfpathlineto{\pgfqpoint{1.659872in}{1.327841in}}%
\pgfpathlineto{\pgfqpoint{1.659233in}{1.305647in}}%
\pgfpathclose%
\pgfusepath{fill}%
\end{pgfscope}%
\begin{pgfscope}%
\pgfpathrectangle{\pgfqpoint{0.407592in}{0.222904in}}{\pgfqpoint{1.411163in}{1.411163in}}%
\pgfusepath{clip}%
\pgfsetbuttcap%
\pgfsetroundjoin%
\definecolor{currentfill}{rgb}{0.168627,0.611765,0.411765}%
\pgfsetfillcolor{currentfill}%
\pgfsetfillopacity{0.180000}%
\pgfsetlinewidth{0.000000pt}%
\definecolor{currentstroke}{rgb}{0.000000,0.000000,0.000000}%
\pgfsetstrokecolor{currentstroke}%
\pgfsetdash{}{0pt}%
\pgfpathmoveto{\pgfqpoint{1.674955in}{1.374813in}}%
\pgfpathlineto{\pgfqpoint{1.689359in}{1.400140in}}%
\pgfpathlineto{\pgfqpoint{1.690032in}{1.422409in}}%
\pgfpathlineto{\pgfqpoint{1.675612in}{1.397098in}}%
\pgfpathlineto{\pgfqpoint{1.674955in}{1.374813in}}%
\pgfpathclose%
\pgfusepath{fill}%
\end{pgfscope}%
\begin{pgfscope}%
\pgfpathrectangle{\pgfqpoint{0.407592in}{0.222904in}}{\pgfqpoint{1.411163in}{1.411163in}}%
\pgfusepath{clip}%
\pgfsetbuttcap%
\pgfsetroundjoin%
\definecolor{currentfill}{rgb}{0.168627,0.611765,0.411765}%
\pgfsetfillcolor{currentfill}%
\pgfsetfillopacity{0.180000}%
\pgfsetlinewidth{0.000000pt}%
\definecolor{currentstroke}{rgb}{0.000000,0.000000,0.000000}%
\pgfsetstrokecolor{currentstroke}%
\pgfsetdash{}{0pt}%
\pgfpathmoveto{\pgfqpoint{1.597217in}{1.059631in}}%
\pgfpathlineto{\pgfqpoint{1.611674in}{1.081405in}}%
\pgfpathlineto{\pgfqpoint{1.612254in}{1.103288in}}%
\pgfpathlineto{\pgfqpoint{1.597781in}{1.081538in}}%
\pgfpathlineto{\pgfqpoint{1.597217in}{1.059631in}}%
\pgfpathclose%
\pgfusepath{fill}%
\end{pgfscope}%
\begin{pgfscope}%
\pgfpathrectangle{\pgfqpoint{0.407592in}{0.222904in}}{\pgfqpoint{1.411163in}{1.411163in}}%
\pgfusepath{clip}%
\pgfsetbuttcap%
\pgfsetroundjoin%
\definecolor{currentfill}{rgb}{0.168627,0.611765,0.411765}%
\pgfsetfillcolor{currentfill}%
\pgfsetfillopacity{0.180000}%
\pgfsetlinewidth{0.000000pt}%
\definecolor{currentstroke}{rgb}{0.000000,0.000000,0.000000}%
\pgfsetstrokecolor{currentstroke}%
\pgfsetdash{}{0pt}%
\pgfpathmoveto{\pgfqpoint{1.612835in}{1.125224in}}%
\pgfpathlineto{\pgfqpoint{1.627284in}{1.147851in}}%
\pgfpathlineto{\pgfqpoint{1.627883in}{1.169818in}}%
\pgfpathlineto{\pgfqpoint{1.613417in}{1.147213in}}%
\pgfpathlineto{\pgfqpoint{1.612835in}{1.125224in}}%
\pgfpathclose%
\pgfusepath{fill}%
\end{pgfscope}%
\begin{pgfscope}%
\pgfpathrectangle{\pgfqpoint{0.407592in}{0.222904in}}{\pgfqpoint{1.411163in}{1.411163in}}%
\pgfusepath{clip}%
\pgfsetbuttcap%
\pgfsetroundjoin%
\definecolor{currentfill}{rgb}{0.168627,0.611765,0.411765}%
\pgfsetfillcolor{currentfill}%
\pgfsetfillopacity{0.180000}%
\pgfsetlinewidth{0.000000pt}%
\definecolor{currentstroke}{rgb}{0.000000,0.000000,0.000000}%
\pgfsetstrokecolor{currentstroke}%
\pgfsetdash{}{0pt}%
\pgfpathmoveto{\pgfqpoint{1.628483in}{1.191838in}}%
\pgfpathlineto{\pgfqpoint{1.642925in}{1.215241in}}%
\pgfpathlineto{\pgfqpoint{1.643543in}{1.237293in}}%
\pgfpathlineto{\pgfqpoint{1.629084in}{1.213911in}}%
\pgfpathlineto{\pgfqpoint{1.628483in}{1.191838in}}%
\pgfpathclose%
\pgfusepath{fill}%
\end{pgfscope}%
\begin{pgfscope}%
\pgfpathrectangle{\pgfqpoint{0.407592in}{0.222904in}}{\pgfqpoint{1.411163in}{1.411163in}}%
\pgfusepath{clip}%
\pgfsetbuttcap%
\pgfsetroundjoin%
\definecolor{currentfill}{rgb}{0.168627,0.611765,0.411765}%
\pgfsetfillcolor{currentfill}%
\pgfsetfillopacity{0.180000}%
\pgfsetlinewidth{0.000000pt}%
\definecolor{currentstroke}{rgb}{0.000000,0.000000,0.000000}%
\pgfsetstrokecolor{currentstroke}%
\pgfsetdash{}{0pt}%
\pgfpathmoveto{\pgfqpoint{1.644162in}{1.259399in}}%
\pgfpathlineto{\pgfqpoint{1.658597in}{1.283506in}}%
\pgfpathlineto{\pgfqpoint{1.659233in}{1.305647in}}%
\pgfpathlineto{\pgfqpoint{1.644782in}{1.281559in}}%
\pgfpathlineto{\pgfqpoint{1.644162in}{1.259399in}}%
\pgfpathclose%
\pgfusepath{fill}%
\end{pgfscope}%
\begin{pgfscope}%
\pgfpathrectangle{\pgfqpoint{0.407592in}{0.222904in}}{\pgfqpoint{1.411163in}{1.411163in}}%
\pgfusepath{clip}%
\pgfsetbuttcap%
\pgfsetroundjoin%
\definecolor{currentfill}{rgb}{0.168627,0.611765,0.411765}%
\pgfsetfillcolor{currentfill}%
\pgfsetfillopacity{0.180000}%
\pgfsetlinewidth{0.000000pt}%
\definecolor{currentstroke}{rgb}{0.000000,0.000000,0.000000}%
\pgfsetstrokecolor{currentstroke}%
\pgfsetdash{}{0pt}%
\pgfpathmoveto{\pgfqpoint{1.659872in}{1.327841in}}%
\pgfpathlineto{\pgfqpoint{1.674299in}{1.352582in}}%
\pgfpathlineto{\pgfqpoint{1.674955in}{1.374813in}}%
\pgfpathlineto{\pgfqpoint{1.660511in}{1.350090in}}%
\pgfpathlineto{\pgfqpoint{1.659872in}{1.327841in}}%
\pgfpathclose%
\pgfusepath{fill}%
\end{pgfscope}%
\begin{pgfscope}%
\pgfpathrectangle{\pgfqpoint{0.407592in}{0.222904in}}{\pgfqpoint{1.411163in}{1.411163in}}%
\pgfusepath{clip}%
\pgfsetbuttcap%
\pgfsetroundjoin%
\definecolor{currentfill}{rgb}{0.168627,0.611765,0.411765}%
\pgfsetfillcolor{currentfill}%
\pgfsetfillopacity{0.180000}%
\pgfsetlinewidth{0.000000pt}%
\definecolor{currentstroke}{rgb}{0.000000,0.000000,0.000000}%
\pgfsetstrokecolor{currentstroke}%
\pgfsetdash{}{0pt}%
\pgfpathmoveto{\pgfqpoint{1.675612in}{1.397098in}}%
\pgfpathlineto{\pgfqpoint{1.690032in}{1.422409in}}%
\pgfpathlineto{\pgfqpoint{1.690707in}{1.444732in}}%
\pgfpathlineto{\pgfqpoint{1.676271in}{1.419437in}}%
\pgfpathlineto{\pgfqpoint{1.675612in}{1.397098in}}%
\pgfpathclose%
\pgfusepath{fill}%
\end{pgfscope}%
\begin{pgfscope}%
\pgfpathrectangle{\pgfqpoint{0.407592in}{0.222904in}}{\pgfqpoint{1.411163in}{1.411163in}}%
\pgfusepath{clip}%
\pgfsetbuttcap%
\pgfsetroundjoin%
\definecolor{currentfill}{rgb}{0.168627,0.611765,0.411765}%
\pgfsetfillcolor{currentfill}%
\pgfsetfillopacity{0.180000}%
\pgfsetlinewidth{0.000000pt}%
\definecolor{currentstroke}{rgb}{0.000000,0.000000,0.000000}%
\pgfsetstrokecolor{currentstroke}%
\pgfsetdash{}{0pt}%
\pgfpathmoveto{\pgfqpoint{1.582174in}{1.016959in}}%
\pgfpathlineto{\pgfqpoint{1.596655in}{1.037776in}}%
\pgfpathlineto{\pgfqpoint{1.597217in}{1.059631in}}%
\pgfpathlineto{\pgfqpoint{1.582720in}{1.038841in}}%
\pgfpathlineto{\pgfqpoint{1.582174in}{1.016959in}}%
\pgfpathclose%
\pgfusepath{fill}%
\end{pgfscope}%
\begin{pgfscope}%
\pgfpathrectangle{\pgfqpoint{0.407592in}{0.222904in}}{\pgfqpoint{1.411163in}{1.411163in}}%
\pgfusepath{clip}%
\pgfsetbuttcap%
\pgfsetroundjoin%
\definecolor{currentfill}{rgb}{0.168627,0.611765,0.411765}%
\pgfsetfillcolor{currentfill}%
\pgfsetfillopacity{0.180000}%
\pgfsetlinewidth{0.000000pt}%
\definecolor{currentstroke}{rgb}{0.000000,0.000000,0.000000}%
\pgfsetstrokecolor{currentstroke}%
\pgfsetdash{}{0pt}%
\pgfpathmoveto{\pgfqpoint{1.597781in}{1.081538in}}%
\pgfpathlineto{\pgfqpoint{1.612254in}{1.103288in}}%
\pgfpathlineto{\pgfqpoint{1.612835in}{1.125224in}}%
\pgfpathlineto{\pgfqpoint{1.598345in}{1.103499in}}%
\pgfpathlineto{\pgfqpoint{1.597781in}{1.081538in}}%
\pgfpathclose%
\pgfusepath{fill}%
\end{pgfscope}%
\begin{pgfscope}%
\pgfpathrectangle{\pgfqpoint{0.407592in}{0.222904in}}{\pgfqpoint{1.411163in}{1.411163in}}%
\pgfusepath{clip}%
\pgfsetbuttcap%
\pgfsetroundjoin%
\definecolor{currentfill}{rgb}{0.168627,0.611765,0.411765}%
\pgfsetfillcolor{currentfill}%
\pgfsetfillopacity{0.180000}%
\pgfsetlinewidth{0.000000pt}%
\definecolor{currentstroke}{rgb}{0.000000,0.000000,0.000000}%
\pgfsetstrokecolor{currentstroke}%
\pgfsetdash{}{0pt}%
\pgfpathmoveto{\pgfqpoint{1.613417in}{1.147213in}}%
\pgfpathlineto{\pgfqpoint{1.627883in}{1.169818in}}%
\pgfpathlineto{\pgfqpoint{1.628483in}{1.191838in}}%
\pgfpathlineto{\pgfqpoint{1.614001in}{1.169256in}}%
\pgfpathlineto{\pgfqpoint{1.613417in}{1.147213in}}%
\pgfpathclose%
\pgfusepath{fill}%
\end{pgfscope}%
\begin{pgfscope}%
\pgfpathrectangle{\pgfqpoint{0.407592in}{0.222904in}}{\pgfqpoint{1.411163in}{1.411163in}}%
\pgfusepath{clip}%
\pgfsetbuttcap%
\pgfsetroundjoin%
\definecolor{currentfill}{rgb}{0.168627,0.611765,0.411765}%
\pgfsetfillcolor{currentfill}%
\pgfsetfillopacity{0.180000}%
\pgfsetlinewidth{0.000000pt}%
\definecolor{currentstroke}{rgb}{0.000000,0.000000,0.000000}%
\pgfsetstrokecolor{currentstroke}%
\pgfsetdash{}{0pt}%
\pgfpathmoveto{\pgfqpoint{1.629084in}{1.213911in}}%
\pgfpathlineto{\pgfqpoint{1.643543in}{1.237293in}}%
\pgfpathlineto{\pgfqpoint{1.644162in}{1.259399in}}%
\pgfpathlineto{\pgfqpoint{1.629687in}{1.236038in}}%
\pgfpathlineto{\pgfqpoint{1.629084in}{1.213911in}}%
\pgfpathclose%
\pgfusepath{fill}%
\end{pgfscope}%
\begin{pgfscope}%
\pgfpathrectangle{\pgfqpoint{0.407592in}{0.222904in}}{\pgfqpoint{1.411163in}{1.411163in}}%
\pgfusepath{clip}%
\pgfsetbuttcap%
\pgfsetroundjoin%
\definecolor{currentfill}{rgb}{0.168627,0.611765,0.411765}%
\pgfsetfillcolor{currentfill}%
\pgfsetfillopacity{0.180000}%
\pgfsetlinewidth{0.000000pt}%
\definecolor{currentstroke}{rgb}{0.000000,0.000000,0.000000}%
\pgfsetstrokecolor{currentstroke}%
\pgfsetdash{}{0pt}%
\pgfpathmoveto{\pgfqpoint{1.644782in}{1.281559in}}%
\pgfpathlineto{\pgfqpoint{1.659233in}{1.305647in}}%
\pgfpathlineto{\pgfqpoint{1.659872in}{1.327841in}}%
\pgfpathlineto{\pgfqpoint{1.645404in}{1.303773in}}%
\pgfpathlineto{\pgfqpoint{1.644782in}{1.281559in}}%
\pgfpathclose%
\pgfusepath{fill}%
\end{pgfscope}%
\begin{pgfscope}%
\pgfpathrectangle{\pgfqpoint{0.407592in}{0.222904in}}{\pgfqpoint{1.411163in}{1.411163in}}%
\pgfusepath{clip}%
\pgfsetbuttcap%
\pgfsetroundjoin%
\definecolor{currentfill}{rgb}{0.168627,0.611765,0.411765}%
\pgfsetfillcolor{currentfill}%
\pgfsetfillopacity{0.180000}%
\pgfsetlinewidth{0.000000pt}%
\definecolor{currentstroke}{rgb}{0.000000,0.000000,0.000000}%
\pgfsetstrokecolor{currentstroke}%
\pgfsetdash{}{0pt}%
\pgfpathmoveto{\pgfqpoint{1.660511in}{1.350090in}}%
\pgfpathlineto{\pgfqpoint{1.674955in}{1.374813in}}%
\pgfpathlineto{\pgfqpoint{1.675612in}{1.397098in}}%
\pgfpathlineto{\pgfqpoint{1.661153in}{1.372392in}}%
\pgfpathlineto{\pgfqpoint{1.660511in}{1.350090in}}%
\pgfpathclose%
\pgfusepath{fill}%
\end{pgfscope}%
\begin{pgfscope}%
\pgfpathrectangle{\pgfqpoint{0.407592in}{0.222904in}}{\pgfqpoint{1.411163in}{1.411163in}}%
\pgfusepath{clip}%
\pgfsetbuttcap%
\pgfsetroundjoin%
\definecolor{currentfill}{rgb}{0.168627,0.611765,0.411765}%
\pgfsetfillcolor{currentfill}%
\pgfsetfillopacity{0.180000}%
\pgfsetlinewidth{0.000000pt}%
\definecolor{currentstroke}{rgb}{0.000000,0.000000,0.000000}%
\pgfsetstrokecolor{currentstroke}%
\pgfsetdash{}{0pt}%
\pgfpathmoveto{\pgfqpoint{1.676271in}{1.419437in}}%
\pgfpathlineto{\pgfqpoint{1.690707in}{1.444732in}}%
\pgfpathlineto{\pgfqpoint{1.691384in}{1.467109in}}%
\pgfpathlineto{\pgfqpoint{1.676932in}{1.441830in}}%
\pgfpathlineto{\pgfqpoint{1.676271in}{1.419437in}}%
\pgfpathclose%
\pgfusepath{fill}%
\end{pgfscope}%
\begin{pgfscope}%
\pgfpathrectangle{\pgfqpoint{0.407592in}{0.222904in}}{\pgfqpoint{1.411163in}{1.411163in}}%
\pgfusepath{clip}%
\pgfsetbuttcap%
\pgfsetroundjoin%
\definecolor{currentfill}{rgb}{0.168627,0.611765,0.411765}%
\pgfsetfillcolor{currentfill}%
\pgfsetfillopacity{0.180000}%
\pgfsetlinewidth{0.000000pt}%
\definecolor{currentstroke}{rgb}{0.000000,0.000000,0.000000}%
\pgfsetstrokecolor{currentstroke}%
\pgfsetdash{}{0pt}%
\pgfpathmoveto{\pgfqpoint{1.582720in}{1.038841in}}%
\pgfpathlineto{\pgfqpoint{1.597217in}{1.059631in}}%
\pgfpathlineto{\pgfqpoint{1.597781in}{1.081538in}}%
\pgfpathlineto{\pgfqpoint{1.583267in}{1.060775in}}%
\pgfpathlineto{\pgfqpoint{1.582720in}{1.038841in}}%
\pgfpathclose%
\pgfusepath{fill}%
\end{pgfscope}%
\begin{pgfscope}%
\pgfpathrectangle{\pgfqpoint{0.407592in}{0.222904in}}{\pgfqpoint{1.411163in}{1.411163in}}%
\pgfusepath{clip}%
\pgfsetbuttcap%
\pgfsetroundjoin%
\definecolor{currentfill}{rgb}{0.168627,0.611765,0.411765}%
\pgfsetfillcolor{currentfill}%
\pgfsetfillopacity{0.180000}%
\pgfsetlinewidth{0.000000pt}%
\definecolor{currentstroke}{rgb}{0.000000,0.000000,0.000000}%
\pgfsetstrokecolor{currentstroke}%
\pgfsetdash{}{0pt}%
\pgfpathmoveto{\pgfqpoint{1.598345in}{1.103499in}}%
\pgfpathlineto{\pgfqpoint{1.612835in}{1.125224in}}%
\pgfpathlineto{\pgfqpoint{1.613417in}{1.147213in}}%
\pgfpathlineto{\pgfqpoint{1.598911in}{1.125513in}}%
\pgfpathlineto{\pgfqpoint{1.598345in}{1.103499in}}%
\pgfpathclose%
\pgfusepath{fill}%
\end{pgfscope}%
\begin{pgfscope}%
\pgfpathrectangle{\pgfqpoint{0.407592in}{0.222904in}}{\pgfqpoint{1.411163in}{1.411163in}}%
\pgfusepath{clip}%
\pgfsetbuttcap%
\pgfsetroundjoin%
\definecolor{currentfill}{rgb}{0.168627,0.611765,0.411765}%
\pgfsetfillcolor{currentfill}%
\pgfsetfillopacity{0.180000}%
\pgfsetlinewidth{0.000000pt}%
\definecolor{currentstroke}{rgb}{0.000000,0.000000,0.000000}%
\pgfsetstrokecolor{currentstroke}%
\pgfsetdash{}{0pt}%
\pgfpathmoveto{\pgfqpoint{1.614001in}{1.169256in}}%
\pgfpathlineto{\pgfqpoint{1.628483in}{1.191838in}}%
\pgfpathlineto{\pgfqpoint{1.629084in}{1.213911in}}%
\pgfpathlineto{\pgfqpoint{1.614586in}{1.191352in}}%
\pgfpathlineto{\pgfqpoint{1.614001in}{1.169256in}}%
\pgfpathclose%
\pgfusepath{fill}%
\end{pgfscope}%
\begin{pgfscope}%
\pgfpathrectangle{\pgfqpoint{0.407592in}{0.222904in}}{\pgfqpoint{1.411163in}{1.411163in}}%
\pgfusepath{clip}%
\pgfsetbuttcap%
\pgfsetroundjoin%
\definecolor{currentfill}{rgb}{0.168627,0.611765,0.411765}%
\pgfsetfillcolor{currentfill}%
\pgfsetfillopacity{0.180000}%
\pgfsetlinewidth{0.000000pt}%
\definecolor{currentstroke}{rgb}{0.000000,0.000000,0.000000}%
\pgfsetstrokecolor{currentstroke}%
\pgfsetdash{}{0pt}%
\pgfpathmoveto{\pgfqpoint{1.629687in}{1.236038in}}%
\pgfpathlineto{\pgfqpoint{1.644162in}{1.259399in}}%
\pgfpathlineto{\pgfqpoint{1.644782in}{1.281559in}}%
\pgfpathlineto{\pgfqpoint{1.630291in}{1.258218in}}%
\pgfpathlineto{\pgfqpoint{1.629687in}{1.236038in}}%
\pgfpathclose%
\pgfusepath{fill}%
\end{pgfscope}%
\begin{pgfscope}%
\pgfpathrectangle{\pgfqpoint{0.407592in}{0.222904in}}{\pgfqpoint{1.411163in}{1.411163in}}%
\pgfusepath{clip}%
\pgfsetbuttcap%
\pgfsetroundjoin%
\definecolor{currentfill}{rgb}{0.168627,0.611765,0.411765}%
\pgfsetfillcolor{currentfill}%
\pgfsetfillopacity{0.180000}%
\pgfsetlinewidth{0.000000pt}%
\definecolor{currentstroke}{rgb}{0.000000,0.000000,0.000000}%
\pgfsetstrokecolor{currentstroke}%
\pgfsetdash{}{0pt}%
\pgfpathmoveto{\pgfqpoint{1.645404in}{1.303773in}}%
\pgfpathlineto{\pgfqpoint{1.659872in}{1.327841in}}%
\pgfpathlineto{\pgfqpoint{1.660511in}{1.350090in}}%
\pgfpathlineto{\pgfqpoint{1.646028in}{1.326040in}}%
\pgfpathlineto{\pgfqpoint{1.645404in}{1.303773in}}%
\pgfpathclose%
\pgfusepath{fill}%
\end{pgfscope}%
\begin{pgfscope}%
\pgfpathrectangle{\pgfqpoint{0.407592in}{0.222904in}}{\pgfqpoint{1.411163in}{1.411163in}}%
\pgfusepath{clip}%
\pgfsetbuttcap%
\pgfsetroundjoin%
\definecolor{currentfill}{rgb}{0.168627,0.611765,0.411765}%
\pgfsetfillcolor{currentfill}%
\pgfsetfillopacity{0.180000}%
\pgfsetlinewidth{0.000000pt}%
\definecolor{currentstroke}{rgb}{0.000000,0.000000,0.000000}%
\pgfsetstrokecolor{currentstroke}%
\pgfsetdash{}{0pt}%
\pgfpathmoveto{\pgfqpoint{1.661153in}{1.372392in}}%
\pgfpathlineto{\pgfqpoint{1.675612in}{1.397098in}}%
\pgfpathlineto{\pgfqpoint{1.676271in}{1.419437in}}%
\pgfpathlineto{\pgfqpoint{1.661795in}{1.394749in}}%
\pgfpathlineto{\pgfqpoint{1.661153in}{1.372392in}}%
\pgfpathclose%
\pgfusepath{fill}%
\end{pgfscope}%
\begin{pgfscope}%
\pgfpathrectangle{\pgfqpoint{0.407592in}{0.222904in}}{\pgfqpoint{1.411163in}{1.411163in}}%
\pgfusepath{clip}%
\pgfsetbuttcap%
\pgfsetroundjoin%
\definecolor{currentfill}{rgb}{0.168627,0.611765,0.411765}%
\pgfsetfillcolor{currentfill}%
\pgfsetfillopacity{0.180000}%
\pgfsetlinewidth{0.000000pt}%
\definecolor{currentstroke}{rgb}{0.000000,0.000000,0.000000}%
\pgfsetstrokecolor{currentstroke}%
\pgfsetdash{}{0pt}%
\pgfpathmoveto{\pgfqpoint{1.567653in}{0.997214in}}%
\pgfpathlineto{\pgfqpoint{1.582174in}{1.016959in}}%
\pgfpathlineto{\pgfqpoint{1.582720in}{1.038841in}}%
\pgfpathlineto{\pgfqpoint{1.568182in}{1.019125in}}%
\pgfpathlineto{\pgfqpoint{1.567653in}{0.997214in}}%
\pgfpathclose%
\pgfusepath{fill}%
\end{pgfscope}%
\begin{pgfscope}%
\pgfpathrectangle{\pgfqpoint{0.407592in}{0.222904in}}{\pgfqpoint{1.411163in}{1.411163in}}%
\pgfusepath{clip}%
\pgfsetbuttcap%
\pgfsetroundjoin%
\definecolor{currentfill}{rgb}{0.168627,0.611765,0.411765}%
\pgfsetfillcolor{currentfill}%
\pgfsetfillopacity{0.180000}%
\pgfsetlinewidth{0.000000pt}%
\definecolor{currentstroke}{rgb}{0.000000,0.000000,0.000000}%
\pgfsetstrokecolor{currentstroke}%
\pgfsetdash{}{0pt}%
\pgfpathmoveto{\pgfqpoint{1.583267in}{1.060775in}}%
\pgfpathlineto{\pgfqpoint{1.597781in}{1.081538in}}%
\pgfpathlineto{\pgfqpoint{1.598345in}{1.103499in}}%
\pgfpathlineto{\pgfqpoint{1.583816in}{1.082763in}}%
\pgfpathlineto{\pgfqpoint{1.583267in}{1.060775in}}%
\pgfpathclose%
\pgfusepath{fill}%
\end{pgfscope}%
\begin{pgfscope}%
\pgfpathrectangle{\pgfqpoint{0.407592in}{0.222904in}}{\pgfqpoint{1.411163in}{1.411163in}}%
\pgfusepath{clip}%
\pgfsetbuttcap%
\pgfsetroundjoin%
\definecolor{currentfill}{rgb}{0.168627,0.611765,0.411765}%
\pgfsetfillcolor{currentfill}%
\pgfsetfillopacity{0.180000}%
\pgfsetlinewidth{0.000000pt}%
\definecolor{currentstroke}{rgb}{0.000000,0.000000,0.000000}%
\pgfsetstrokecolor{currentstroke}%
\pgfsetdash{}{0pt}%
\pgfpathmoveto{\pgfqpoint{1.598911in}{1.125513in}}%
\pgfpathlineto{\pgfqpoint{1.613417in}{1.147213in}}%
\pgfpathlineto{\pgfqpoint{1.614001in}{1.169256in}}%
\pgfpathlineto{\pgfqpoint{1.599479in}{1.147580in}}%
\pgfpathlineto{\pgfqpoint{1.598911in}{1.125513in}}%
\pgfpathclose%
\pgfusepath{fill}%
\end{pgfscope}%
\begin{pgfscope}%
\pgfpathrectangle{\pgfqpoint{0.407592in}{0.222904in}}{\pgfqpoint{1.411163in}{1.411163in}}%
\pgfusepath{clip}%
\pgfsetbuttcap%
\pgfsetroundjoin%
\definecolor{currentfill}{rgb}{0.168627,0.611765,0.411765}%
\pgfsetfillcolor{currentfill}%
\pgfsetfillopacity{0.180000}%
\pgfsetlinewidth{0.000000pt}%
\definecolor{currentstroke}{rgb}{0.000000,0.000000,0.000000}%
\pgfsetstrokecolor{currentstroke}%
\pgfsetdash{}{0pt}%
\pgfpathmoveto{\pgfqpoint{1.614586in}{1.191352in}}%
\pgfpathlineto{\pgfqpoint{1.629084in}{1.213911in}}%
\pgfpathlineto{\pgfqpoint{1.629687in}{1.236038in}}%
\pgfpathlineto{\pgfqpoint{1.615173in}{1.213502in}}%
\pgfpathlineto{\pgfqpoint{1.614586in}{1.191352in}}%
\pgfpathclose%
\pgfusepath{fill}%
\end{pgfscope}%
\begin{pgfscope}%
\pgfpathrectangle{\pgfqpoint{0.407592in}{0.222904in}}{\pgfqpoint{1.411163in}{1.411163in}}%
\pgfusepath{clip}%
\pgfsetbuttcap%
\pgfsetroundjoin%
\definecolor{currentfill}{rgb}{0.168627,0.611765,0.411765}%
\pgfsetfillcolor{currentfill}%
\pgfsetfillopacity{0.180000}%
\pgfsetlinewidth{0.000000pt}%
\definecolor{currentstroke}{rgb}{0.000000,0.000000,0.000000}%
\pgfsetstrokecolor{currentstroke}%
\pgfsetdash{}{0pt}%
\pgfpathmoveto{\pgfqpoint{1.630291in}{1.258218in}}%
\pgfpathlineto{\pgfqpoint{1.644782in}{1.281559in}}%
\pgfpathlineto{\pgfqpoint{1.645404in}{1.303773in}}%
\pgfpathlineto{\pgfqpoint{1.630897in}{1.280453in}}%
\pgfpathlineto{\pgfqpoint{1.630291in}{1.258218in}}%
\pgfpathclose%
\pgfusepath{fill}%
\end{pgfscope}%
\begin{pgfscope}%
\pgfpathrectangle{\pgfqpoint{0.407592in}{0.222904in}}{\pgfqpoint{1.411163in}{1.411163in}}%
\pgfusepath{clip}%
\pgfsetbuttcap%
\pgfsetroundjoin%
\definecolor{currentfill}{rgb}{0.168627,0.611765,0.411765}%
\pgfsetfillcolor{currentfill}%
\pgfsetfillopacity{0.180000}%
\pgfsetlinewidth{0.000000pt}%
\definecolor{currentstroke}{rgb}{0.000000,0.000000,0.000000}%
\pgfsetstrokecolor{currentstroke}%
\pgfsetdash{}{0pt}%
\pgfpathmoveto{\pgfqpoint{1.646028in}{1.326040in}}%
\pgfpathlineto{\pgfqpoint{1.660511in}{1.350090in}}%
\pgfpathlineto{\pgfqpoint{1.661153in}{1.372392in}}%
\pgfpathlineto{\pgfqpoint{1.646653in}{1.348362in}}%
\pgfpathlineto{\pgfqpoint{1.646028in}{1.326040in}}%
\pgfpathclose%
\pgfusepath{fill}%
\end{pgfscope}%
\begin{pgfscope}%
\pgfpathrectangle{\pgfqpoint{0.407592in}{0.222904in}}{\pgfqpoint{1.411163in}{1.411163in}}%
\pgfusepath{clip}%
\pgfsetbuttcap%
\pgfsetroundjoin%
\definecolor{currentfill}{rgb}{0.168627,0.611765,0.411765}%
\pgfsetfillcolor{currentfill}%
\pgfsetfillopacity{0.180000}%
\pgfsetlinewidth{0.000000pt}%
\definecolor{currentstroke}{rgb}{0.000000,0.000000,0.000000}%
\pgfsetstrokecolor{currentstroke}%
\pgfsetdash{}{0pt}%
\pgfpathmoveto{\pgfqpoint{1.661795in}{1.394749in}}%
\pgfpathlineto{\pgfqpoint{1.676271in}{1.419437in}}%
\pgfpathlineto{\pgfqpoint{1.676932in}{1.441830in}}%
\pgfpathlineto{\pgfqpoint{1.662440in}{1.417160in}}%
\pgfpathlineto{\pgfqpoint{1.661795in}{1.394749in}}%
\pgfpathclose%
\pgfusepath{fill}%
\end{pgfscope}%
\begin{pgfscope}%
\pgfpathrectangle{\pgfqpoint{0.407592in}{0.222904in}}{\pgfqpoint{1.411163in}{1.411163in}}%
\pgfusepath{clip}%
\pgfsetbuttcap%
\pgfsetroundjoin%
\definecolor{currentfill}{rgb}{0.168627,0.611765,0.411765}%
\pgfsetfillcolor{currentfill}%
\pgfsetfillopacity{0.180000}%
\pgfsetlinewidth{0.000000pt}%
\definecolor{currentstroke}{rgb}{0.000000,0.000000,0.000000}%
\pgfsetstrokecolor{currentstroke}%
\pgfsetdash{}{0pt}%
\pgfpathmoveto{\pgfqpoint{1.568182in}{1.019125in}}%
\pgfpathlineto{\pgfqpoint{1.582720in}{1.038841in}}%
\pgfpathlineto{\pgfqpoint{1.583267in}{1.060775in}}%
\pgfpathlineto{\pgfqpoint{1.568713in}{1.041090in}}%
\pgfpathlineto{\pgfqpoint{1.568182in}{1.019125in}}%
\pgfpathclose%
\pgfusepath{fill}%
\end{pgfscope}%
\begin{pgfscope}%
\pgfpathrectangle{\pgfqpoint{0.407592in}{0.222904in}}{\pgfqpoint{1.411163in}{1.411163in}}%
\pgfusepath{clip}%
\pgfsetbuttcap%
\pgfsetroundjoin%
\definecolor{currentfill}{rgb}{0.168627,0.611765,0.411765}%
\pgfsetfillcolor{currentfill}%
\pgfsetfillopacity{0.180000}%
\pgfsetlinewidth{0.000000pt}%
\definecolor{currentstroke}{rgb}{0.000000,0.000000,0.000000}%
\pgfsetstrokecolor{currentstroke}%
\pgfsetdash{}{0pt}%
\pgfpathmoveto{\pgfqpoint{1.583816in}{1.082763in}}%
\pgfpathlineto{\pgfqpoint{1.598345in}{1.103499in}}%
\pgfpathlineto{\pgfqpoint{1.598911in}{1.125513in}}%
\pgfpathlineto{\pgfqpoint{1.584365in}{1.104805in}}%
\pgfpathlineto{\pgfqpoint{1.583816in}{1.082763in}}%
\pgfpathclose%
\pgfusepath{fill}%
\end{pgfscope}%
\begin{pgfscope}%
\pgfpathrectangle{\pgfqpoint{0.407592in}{0.222904in}}{\pgfqpoint{1.411163in}{1.411163in}}%
\pgfusepath{clip}%
\pgfsetbuttcap%
\pgfsetroundjoin%
\definecolor{currentfill}{rgb}{0.168627,0.611765,0.411765}%
\pgfsetfillcolor{currentfill}%
\pgfsetfillopacity{0.180000}%
\pgfsetlinewidth{0.000000pt}%
\definecolor{currentstroke}{rgb}{0.000000,0.000000,0.000000}%
\pgfsetstrokecolor{currentstroke}%
\pgfsetdash{}{0pt}%
\pgfpathmoveto{\pgfqpoint{1.599479in}{1.147580in}}%
\pgfpathlineto{\pgfqpoint{1.614001in}{1.169256in}}%
\pgfpathlineto{\pgfqpoint{1.614586in}{1.191352in}}%
\pgfpathlineto{\pgfqpoint{1.600048in}{1.169701in}}%
\pgfpathlineto{\pgfqpoint{1.599479in}{1.147580in}}%
\pgfpathclose%
\pgfusepath{fill}%
\end{pgfscope}%
\begin{pgfscope}%
\pgfpathrectangle{\pgfqpoint{0.407592in}{0.222904in}}{\pgfqpoint{1.411163in}{1.411163in}}%
\pgfusepath{clip}%
\pgfsetbuttcap%
\pgfsetroundjoin%
\definecolor{currentfill}{rgb}{0.168627,0.611765,0.411765}%
\pgfsetfillcolor{currentfill}%
\pgfsetfillopacity{0.180000}%
\pgfsetlinewidth{0.000000pt}%
\definecolor{currentstroke}{rgb}{0.000000,0.000000,0.000000}%
\pgfsetstrokecolor{currentstroke}%
\pgfsetdash{}{0pt}%
\pgfpathmoveto{\pgfqpoint{1.615173in}{1.213502in}}%
\pgfpathlineto{\pgfqpoint{1.629687in}{1.236038in}}%
\pgfpathlineto{\pgfqpoint{1.630291in}{1.258218in}}%
\pgfpathlineto{\pgfqpoint{1.615761in}{1.235705in}}%
\pgfpathlineto{\pgfqpoint{1.615173in}{1.213502in}}%
\pgfpathclose%
\pgfusepath{fill}%
\end{pgfscope}%
\begin{pgfscope}%
\pgfpathrectangle{\pgfqpoint{0.407592in}{0.222904in}}{\pgfqpoint{1.411163in}{1.411163in}}%
\pgfusepath{clip}%
\pgfsetbuttcap%
\pgfsetroundjoin%
\definecolor{currentfill}{rgb}{0.168627,0.611765,0.411765}%
\pgfsetfillcolor{currentfill}%
\pgfsetfillopacity{0.180000}%
\pgfsetlinewidth{0.000000pt}%
\definecolor{currentstroke}{rgb}{0.000000,0.000000,0.000000}%
\pgfsetstrokecolor{currentstroke}%
\pgfsetdash{}{0pt}%
\pgfpathmoveto{\pgfqpoint{1.630897in}{1.280453in}}%
\pgfpathlineto{\pgfqpoint{1.645404in}{1.303773in}}%
\pgfpathlineto{\pgfqpoint{1.646028in}{1.326040in}}%
\pgfpathlineto{\pgfqpoint{1.631504in}{1.302742in}}%
\pgfpathlineto{\pgfqpoint{1.630897in}{1.280453in}}%
\pgfpathclose%
\pgfusepath{fill}%
\end{pgfscope}%
\begin{pgfscope}%
\pgfpathrectangle{\pgfqpoint{0.407592in}{0.222904in}}{\pgfqpoint{1.411163in}{1.411163in}}%
\pgfusepath{clip}%
\pgfsetbuttcap%
\pgfsetroundjoin%
\definecolor{currentfill}{rgb}{0.168627,0.611765,0.411765}%
\pgfsetfillcolor{currentfill}%
\pgfsetfillopacity{0.180000}%
\pgfsetlinewidth{0.000000pt}%
\definecolor{currentstroke}{rgb}{0.000000,0.000000,0.000000}%
\pgfsetstrokecolor{currentstroke}%
\pgfsetdash{}{0pt}%
\pgfpathmoveto{\pgfqpoint{1.646653in}{1.348362in}}%
\pgfpathlineto{\pgfqpoint{1.661153in}{1.372392in}}%
\pgfpathlineto{\pgfqpoint{1.661795in}{1.394749in}}%
\pgfpathlineto{\pgfqpoint{1.647279in}{1.370738in}}%
\pgfpathlineto{\pgfqpoint{1.646653in}{1.348362in}}%
\pgfpathclose%
\pgfusepath{fill}%
\end{pgfscope}%
\begin{pgfscope}%
\pgfpathrectangle{\pgfqpoint{0.407592in}{0.222904in}}{\pgfqpoint{1.411163in}{1.411163in}}%
\pgfusepath{clip}%
\pgfsetbuttcap%
\pgfsetroundjoin%
\definecolor{currentfill}{rgb}{0.168627,0.611765,0.411765}%
\pgfsetfillcolor{currentfill}%
\pgfsetfillopacity{0.180000}%
\pgfsetlinewidth{0.000000pt}%
\definecolor{currentstroke}{rgb}{0.000000,0.000000,0.000000}%
\pgfsetstrokecolor{currentstroke}%
\pgfsetdash{}{0pt}%
\pgfpathmoveto{\pgfqpoint{1.568713in}{1.041090in}}%
\pgfpathlineto{\pgfqpoint{1.583267in}{1.060775in}}%
\pgfpathlineto{\pgfqpoint{1.583816in}{1.082763in}}%
\pgfpathlineto{\pgfqpoint{1.569245in}{1.063108in}}%
\pgfpathlineto{\pgfqpoint{1.568713in}{1.041090in}}%
\pgfpathclose%
\pgfusepath{fill}%
\end{pgfscope}%
\begin{pgfscope}%
\pgfpathrectangle{\pgfqpoint{0.407592in}{0.222904in}}{\pgfqpoint{1.411163in}{1.411163in}}%
\pgfusepath{clip}%
\pgfsetbuttcap%
\pgfsetroundjoin%
\definecolor{currentfill}{rgb}{0.168627,0.611765,0.411765}%
\pgfsetfillcolor{currentfill}%
\pgfsetfillopacity{0.180000}%
\pgfsetlinewidth{0.000000pt}%
\definecolor{currentstroke}{rgb}{0.000000,0.000000,0.000000}%
\pgfsetstrokecolor{currentstroke}%
\pgfsetdash{}{0pt}%
\pgfpathmoveto{\pgfqpoint{1.584365in}{1.104805in}}%
\pgfpathlineto{\pgfqpoint{1.598911in}{1.125513in}}%
\pgfpathlineto{\pgfqpoint{1.599479in}{1.147580in}}%
\pgfpathlineto{\pgfqpoint{1.584916in}{1.126900in}}%
\pgfpathlineto{\pgfqpoint{1.584365in}{1.104805in}}%
\pgfpathclose%
\pgfusepath{fill}%
\end{pgfscope}%
\begin{pgfscope}%
\pgfpathrectangle{\pgfqpoint{0.407592in}{0.222904in}}{\pgfqpoint{1.411163in}{1.411163in}}%
\pgfusepath{clip}%
\pgfsetbuttcap%
\pgfsetroundjoin%
\definecolor{currentfill}{rgb}{0.168627,0.611765,0.411765}%
\pgfsetfillcolor{currentfill}%
\pgfsetfillopacity{0.180000}%
\pgfsetlinewidth{0.000000pt}%
\definecolor{currentstroke}{rgb}{0.000000,0.000000,0.000000}%
\pgfsetstrokecolor{currentstroke}%
\pgfsetdash{}{0pt}%
\pgfpathmoveto{\pgfqpoint{1.553089in}{0.978630in}}%
\pgfpathlineto{\pgfqpoint{1.567653in}{0.997214in}}%
\pgfpathlineto{\pgfqpoint{1.568182in}{1.019125in}}%
\pgfpathlineto{\pgfqpoint{1.553602in}{1.000574in}}%
\pgfpathlineto{\pgfqpoint{1.553089in}{0.978630in}}%
\pgfpathclose%
\pgfusepath{fill}%
\end{pgfscope}%
\begin{pgfscope}%
\pgfpathrectangle{\pgfqpoint{0.407592in}{0.222904in}}{\pgfqpoint{1.411163in}{1.411163in}}%
\pgfusepath{clip}%
\pgfsetbuttcap%
\pgfsetroundjoin%
\definecolor{currentfill}{rgb}{0.168627,0.611765,0.411765}%
\pgfsetfillcolor{currentfill}%
\pgfsetfillopacity{0.180000}%
\pgfsetlinewidth{0.000000pt}%
\definecolor{currentstroke}{rgb}{0.000000,0.000000,0.000000}%
\pgfsetstrokecolor{currentstroke}%
\pgfsetdash{}{0pt}%
\pgfpathmoveto{\pgfqpoint{1.600048in}{1.169701in}}%
\pgfpathlineto{\pgfqpoint{1.614586in}{1.191352in}}%
\pgfpathlineto{\pgfqpoint{1.615173in}{1.213502in}}%
\pgfpathlineto{\pgfqpoint{1.600618in}{1.191877in}}%
\pgfpathlineto{\pgfqpoint{1.600048in}{1.169701in}}%
\pgfpathclose%
\pgfusepath{fill}%
\end{pgfscope}%
\begin{pgfscope}%
\pgfpathrectangle{\pgfqpoint{0.407592in}{0.222904in}}{\pgfqpoint{1.411163in}{1.411163in}}%
\pgfusepath{clip}%
\pgfsetbuttcap%
\pgfsetroundjoin%
\definecolor{currentfill}{rgb}{0.168627,0.611765,0.411765}%
\pgfsetfillcolor{currentfill}%
\pgfsetfillopacity{0.180000}%
\pgfsetlinewidth{0.000000pt}%
\definecolor{currentstroke}{rgb}{0.000000,0.000000,0.000000}%
\pgfsetstrokecolor{currentstroke}%
\pgfsetdash{}{0pt}%
\pgfpathmoveto{\pgfqpoint{1.615761in}{1.235705in}}%
\pgfpathlineto{\pgfqpoint{1.630291in}{1.258218in}}%
\pgfpathlineto{\pgfqpoint{1.630897in}{1.280453in}}%
\pgfpathlineto{\pgfqpoint{1.616350in}{1.257963in}}%
\pgfpathlineto{\pgfqpoint{1.615761in}{1.235705in}}%
\pgfpathclose%
\pgfusepath{fill}%
\end{pgfscope}%
\begin{pgfscope}%
\pgfpathrectangle{\pgfqpoint{0.407592in}{0.222904in}}{\pgfqpoint{1.411163in}{1.411163in}}%
\pgfusepath{clip}%
\pgfsetbuttcap%
\pgfsetroundjoin%
\definecolor{currentfill}{rgb}{0.168627,0.611765,0.411765}%
\pgfsetfillcolor{currentfill}%
\pgfsetfillopacity{0.180000}%
\pgfsetlinewidth{0.000000pt}%
\definecolor{currentstroke}{rgb}{0.000000,0.000000,0.000000}%
\pgfsetstrokecolor{currentstroke}%
\pgfsetdash{}{0pt}%
\pgfpathmoveto{\pgfqpoint{1.631504in}{1.302742in}}%
\pgfpathlineto{\pgfqpoint{1.646028in}{1.326040in}}%
\pgfpathlineto{\pgfqpoint{1.646653in}{1.348362in}}%
\pgfpathlineto{\pgfqpoint{1.632113in}{1.325085in}}%
\pgfpathlineto{\pgfqpoint{1.631504in}{1.302742in}}%
\pgfpathclose%
\pgfusepath{fill}%
\end{pgfscope}%
\begin{pgfscope}%
\pgfpathrectangle{\pgfqpoint{0.407592in}{0.222904in}}{\pgfqpoint{1.411163in}{1.411163in}}%
\pgfusepath{clip}%
\pgfsetbuttcap%
\pgfsetroundjoin%
\definecolor{currentfill}{rgb}{0.168627,0.611765,0.411765}%
\pgfsetfillcolor{currentfill}%
\pgfsetfillopacity{0.180000}%
\pgfsetlinewidth{0.000000pt}%
\definecolor{currentstroke}{rgb}{0.000000,0.000000,0.000000}%
\pgfsetstrokecolor{currentstroke}%
\pgfsetdash{}{0pt}%
\pgfpathmoveto{\pgfqpoint{1.647279in}{1.370738in}}%
\pgfpathlineto{\pgfqpoint{1.661795in}{1.394749in}}%
\pgfpathlineto{\pgfqpoint{1.662440in}{1.417160in}}%
\pgfpathlineto{\pgfqpoint{1.647907in}{1.393169in}}%
\pgfpathlineto{\pgfqpoint{1.647279in}{1.370738in}}%
\pgfpathclose%
\pgfusepath{fill}%
\end{pgfscope}%
\begin{pgfscope}%
\pgfpathrectangle{\pgfqpoint{0.407592in}{0.222904in}}{\pgfqpoint{1.411163in}{1.411163in}}%
\pgfusepath{clip}%
\pgfsetbuttcap%
\pgfsetroundjoin%
\definecolor{currentfill}{rgb}{0.168627,0.611765,0.411765}%
\pgfsetfillcolor{currentfill}%
\pgfsetfillopacity{0.180000}%
\pgfsetlinewidth{0.000000pt}%
\definecolor{currentstroke}{rgb}{0.000000,0.000000,0.000000}%
\pgfsetstrokecolor{currentstroke}%
\pgfsetdash{}{0pt}%
\pgfpathmoveto{\pgfqpoint{1.569245in}{1.063108in}}%
\pgfpathlineto{\pgfqpoint{1.583816in}{1.082763in}}%
\pgfpathlineto{\pgfqpoint{1.584365in}{1.104805in}}%
\pgfpathlineto{\pgfqpoint{1.569778in}{1.085180in}}%
\pgfpathlineto{\pgfqpoint{1.569245in}{1.063108in}}%
\pgfpathclose%
\pgfusepath{fill}%
\end{pgfscope}%
\begin{pgfscope}%
\pgfpathrectangle{\pgfqpoint{0.407592in}{0.222904in}}{\pgfqpoint{1.411163in}{1.411163in}}%
\pgfusepath{clip}%
\pgfsetbuttcap%
\pgfsetroundjoin%
\definecolor{currentfill}{rgb}{0.168627,0.611765,0.411765}%
\pgfsetfillcolor{currentfill}%
\pgfsetfillopacity{0.180000}%
\pgfsetlinewidth{0.000000pt}%
\definecolor{currentstroke}{rgb}{0.000000,0.000000,0.000000}%
\pgfsetstrokecolor{currentstroke}%
\pgfsetdash{}{0pt}%
\pgfpathmoveto{\pgfqpoint{1.584916in}{1.126900in}}%
\pgfpathlineto{\pgfqpoint{1.599479in}{1.147580in}}%
\pgfpathlineto{\pgfqpoint{1.600048in}{1.169701in}}%
\pgfpathlineto{\pgfqpoint{1.585469in}{1.149049in}}%
\pgfpathlineto{\pgfqpoint{1.584916in}{1.126900in}}%
\pgfpathclose%
\pgfusepath{fill}%
\end{pgfscope}%
\begin{pgfscope}%
\pgfpathrectangle{\pgfqpoint{0.407592in}{0.222904in}}{\pgfqpoint{1.411163in}{1.411163in}}%
\pgfusepath{clip}%
\pgfsetbuttcap%
\pgfsetroundjoin%
\definecolor{currentfill}{rgb}{0.168627,0.611765,0.411765}%
\pgfsetfillcolor{currentfill}%
\pgfsetfillopacity{0.180000}%
\pgfsetlinewidth{0.000000pt}%
\definecolor{currentstroke}{rgb}{0.000000,0.000000,0.000000}%
\pgfsetstrokecolor{currentstroke}%
\pgfsetdash{}{0pt}%
\pgfpathmoveto{\pgfqpoint{1.553602in}{1.000574in}}%
\pgfpathlineto{\pgfqpoint{1.568182in}{1.019125in}}%
\pgfpathlineto{\pgfqpoint{1.568713in}{1.041090in}}%
\pgfpathlineto{\pgfqpoint{1.554116in}{1.022572in}}%
\pgfpathlineto{\pgfqpoint{1.553602in}{1.000574in}}%
\pgfpathclose%
\pgfusepath{fill}%
\end{pgfscope}%
\begin{pgfscope}%
\pgfpathrectangle{\pgfqpoint{0.407592in}{0.222904in}}{\pgfqpoint{1.411163in}{1.411163in}}%
\pgfusepath{clip}%
\pgfsetbuttcap%
\pgfsetroundjoin%
\definecolor{currentfill}{rgb}{0.168627,0.611765,0.411765}%
\pgfsetfillcolor{currentfill}%
\pgfsetfillopacity{0.180000}%
\pgfsetlinewidth{0.000000pt}%
\definecolor{currentstroke}{rgb}{0.000000,0.000000,0.000000}%
\pgfsetstrokecolor{currentstroke}%
\pgfsetdash{}{0pt}%
\pgfpathmoveto{\pgfqpoint{1.600618in}{1.191877in}}%
\pgfpathlineto{\pgfqpoint{1.615173in}{1.213502in}}%
\pgfpathlineto{\pgfqpoint{1.615761in}{1.235705in}}%
\pgfpathlineto{\pgfqpoint{1.601190in}{1.214106in}}%
\pgfpathlineto{\pgfqpoint{1.600618in}{1.191877in}}%
\pgfpathclose%
\pgfusepath{fill}%
\end{pgfscope}%
\begin{pgfscope}%
\pgfpathrectangle{\pgfqpoint{0.407592in}{0.222904in}}{\pgfqpoint{1.411163in}{1.411163in}}%
\pgfusepath{clip}%
\pgfsetbuttcap%
\pgfsetroundjoin%
\definecolor{currentfill}{rgb}{0.168627,0.611765,0.411765}%
\pgfsetfillcolor{currentfill}%
\pgfsetfillopacity{0.180000}%
\pgfsetlinewidth{0.000000pt}%
\definecolor{currentstroke}{rgb}{0.000000,0.000000,0.000000}%
\pgfsetstrokecolor{currentstroke}%
\pgfsetdash{}{0pt}%
\pgfpathmoveto{\pgfqpoint{1.616350in}{1.257963in}}%
\pgfpathlineto{\pgfqpoint{1.630897in}{1.280453in}}%
\pgfpathlineto{\pgfqpoint{1.631504in}{1.302742in}}%
\pgfpathlineto{\pgfqpoint{1.616941in}{1.280275in}}%
\pgfpathlineto{\pgfqpoint{1.616350in}{1.257963in}}%
\pgfpathclose%
\pgfusepath{fill}%
\end{pgfscope}%
\begin{pgfscope}%
\pgfpathrectangle{\pgfqpoint{0.407592in}{0.222904in}}{\pgfqpoint{1.411163in}{1.411163in}}%
\pgfusepath{clip}%
\pgfsetbuttcap%
\pgfsetroundjoin%
\definecolor{currentfill}{rgb}{0.168627,0.611765,0.411765}%
\pgfsetfillcolor{currentfill}%
\pgfsetfillopacity{0.180000}%
\pgfsetlinewidth{0.000000pt}%
\definecolor{currentstroke}{rgb}{0.000000,0.000000,0.000000}%
\pgfsetstrokecolor{currentstroke}%
\pgfsetdash{}{0pt}%
\pgfpathmoveto{\pgfqpoint{1.632113in}{1.325085in}}%
\pgfpathlineto{\pgfqpoint{1.646653in}{1.348362in}}%
\pgfpathlineto{\pgfqpoint{1.647279in}{1.370738in}}%
\pgfpathlineto{\pgfqpoint{1.632723in}{1.347483in}}%
\pgfpathlineto{\pgfqpoint{1.632113in}{1.325085in}}%
\pgfpathclose%
\pgfusepath{fill}%
\end{pgfscope}%
\begin{pgfscope}%
\pgfpathrectangle{\pgfqpoint{0.407592in}{0.222904in}}{\pgfqpoint{1.411163in}{1.411163in}}%
\pgfusepath{clip}%
\pgfsetbuttcap%
\pgfsetroundjoin%
\definecolor{currentfill}{rgb}{0.168627,0.611765,0.411765}%
\pgfsetfillcolor{currentfill}%
\pgfsetfillopacity{0.180000}%
\pgfsetlinewidth{0.000000pt}%
\definecolor{currentstroke}{rgb}{0.000000,0.000000,0.000000}%
\pgfsetstrokecolor{currentstroke}%
\pgfsetdash{}{0pt}%
\pgfpathmoveto{\pgfqpoint{1.569778in}{1.085180in}}%
\pgfpathlineto{\pgfqpoint{1.584365in}{1.104805in}}%
\pgfpathlineto{\pgfqpoint{1.584916in}{1.126900in}}%
\pgfpathlineto{\pgfqpoint{1.570313in}{1.107305in}}%
\pgfpathlineto{\pgfqpoint{1.569778in}{1.085180in}}%
\pgfpathclose%
\pgfusepath{fill}%
\end{pgfscope}%
\begin{pgfscope}%
\pgfpathrectangle{\pgfqpoint{0.407592in}{0.222904in}}{\pgfqpoint{1.411163in}{1.411163in}}%
\pgfusepath{clip}%
\pgfsetbuttcap%
\pgfsetroundjoin%
\definecolor{currentfill}{rgb}{0.168627,0.611765,0.411765}%
\pgfsetfillcolor{currentfill}%
\pgfsetfillopacity{0.180000}%
\pgfsetlinewidth{0.000000pt}%
\definecolor{currentstroke}{rgb}{0.000000,0.000000,0.000000}%
\pgfsetstrokecolor{currentstroke}%
\pgfsetdash{}{0pt}%
\pgfpathmoveto{\pgfqpoint{1.585469in}{1.149049in}}%
\pgfpathlineto{\pgfqpoint{1.600048in}{1.169701in}}%
\pgfpathlineto{\pgfqpoint{1.600618in}{1.191877in}}%
\pgfpathlineto{\pgfqpoint{1.586023in}{1.171252in}}%
\pgfpathlineto{\pgfqpoint{1.585469in}{1.149049in}}%
\pgfpathclose%
\pgfusepath{fill}%
\end{pgfscope}%
\begin{pgfscope}%
\pgfpathrectangle{\pgfqpoint{0.407592in}{0.222904in}}{\pgfqpoint{1.411163in}{1.411163in}}%
\pgfusepath{clip}%
\pgfsetbuttcap%
\pgfsetroundjoin%
\definecolor{currentfill}{rgb}{0.168627,0.611765,0.411765}%
\pgfsetfillcolor{currentfill}%
\pgfsetfillopacity{0.180000}%
\pgfsetlinewidth{0.000000pt}%
\definecolor{currentstroke}{rgb}{0.000000,0.000000,0.000000}%
\pgfsetstrokecolor{currentstroke}%
\pgfsetdash{}{0pt}%
\pgfpathmoveto{\pgfqpoint{1.554116in}{1.022572in}}%
\pgfpathlineto{\pgfqpoint{1.568713in}{1.041090in}}%
\pgfpathlineto{\pgfqpoint{1.569245in}{1.063108in}}%
\pgfpathlineto{\pgfqpoint{1.554632in}{1.044623in}}%
\pgfpathlineto{\pgfqpoint{1.554116in}{1.022572in}}%
\pgfpathclose%
\pgfusepath{fill}%
\end{pgfscope}%
\begin{pgfscope}%
\pgfpathrectangle{\pgfqpoint{0.407592in}{0.222904in}}{\pgfqpoint{1.411163in}{1.411163in}}%
\pgfusepath{clip}%
\pgfsetbuttcap%
\pgfsetroundjoin%
\definecolor{currentfill}{rgb}{0.168627,0.611765,0.411765}%
\pgfsetfillcolor{currentfill}%
\pgfsetfillopacity{0.180000}%
\pgfsetlinewidth{0.000000pt}%
\definecolor{currentstroke}{rgb}{0.000000,0.000000,0.000000}%
\pgfsetstrokecolor{currentstroke}%
\pgfsetdash{}{0pt}%
\pgfpathmoveto{\pgfqpoint{1.601190in}{1.214106in}}%
\pgfpathlineto{\pgfqpoint{1.615761in}{1.235705in}}%
\pgfpathlineto{\pgfqpoint{1.616350in}{1.257963in}}%
\pgfpathlineto{\pgfqpoint{1.601763in}{1.236389in}}%
\pgfpathlineto{\pgfqpoint{1.601190in}{1.214106in}}%
\pgfpathclose%
\pgfusepath{fill}%
\end{pgfscope}%
\begin{pgfscope}%
\pgfpathrectangle{\pgfqpoint{0.407592in}{0.222904in}}{\pgfqpoint{1.411163in}{1.411163in}}%
\pgfusepath{clip}%
\pgfsetbuttcap%
\pgfsetroundjoin%
\definecolor{currentfill}{rgb}{0.168627,0.611765,0.411765}%
\pgfsetfillcolor{currentfill}%
\pgfsetfillopacity{0.180000}%
\pgfsetlinewidth{0.000000pt}%
\definecolor{currentstroke}{rgb}{0.000000,0.000000,0.000000}%
\pgfsetstrokecolor{currentstroke}%
\pgfsetdash{}{0pt}%
\pgfpathmoveto{\pgfqpoint{1.538481in}{0.961294in}}%
\pgfpathlineto{\pgfqpoint{1.553089in}{0.978630in}}%
\pgfpathlineto{\pgfqpoint{1.553602in}{1.000574in}}%
\pgfpathlineto{\pgfqpoint{1.538978in}{0.983275in}}%
\pgfpathlineto{\pgfqpoint{1.538481in}{0.961294in}}%
\pgfpathclose%
\pgfusepath{fill}%
\end{pgfscope}%
\begin{pgfscope}%
\pgfpathrectangle{\pgfqpoint{0.407592in}{0.222904in}}{\pgfqpoint{1.411163in}{1.411163in}}%
\pgfusepath{clip}%
\pgfsetbuttcap%
\pgfsetroundjoin%
\definecolor{currentfill}{rgb}{0.168627,0.611765,0.411765}%
\pgfsetfillcolor{currentfill}%
\pgfsetfillopacity{0.180000}%
\pgfsetlinewidth{0.000000pt}%
\definecolor{currentstroke}{rgb}{0.000000,0.000000,0.000000}%
\pgfsetstrokecolor{currentstroke}%
\pgfsetdash{}{0pt}%
\pgfpathmoveto{\pgfqpoint{1.616941in}{1.280275in}}%
\pgfpathlineto{\pgfqpoint{1.631504in}{1.302742in}}%
\pgfpathlineto{\pgfqpoint{1.632113in}{1.325085in}}%
\pgfpathlineto{\pgfqpoint{1.617533in}{1.302642in}}%
\pgfpathlineto{\pgfqpoint{1.616941in}{1.280275in}}%
\pgfpathclose%
\pgfusepath{fill}%
\end{pgfscope}%
\begin{pgfscope}%
\pgfpathrectangle{\pgfqpoint{0.407592in}{0.222904in}}{\pgfqpoint{1.411163in}{1.411163in}}%
\pgfusepath{clip}%
\pgfsetbuttcap%
\pgfsetroundjoin%
\definecolor{currentfill}{rgb}{0.168627,0.611765,0.411765}%
\pgfsetfillcolor{currentfill}%
\pgfsetfillopacity{0.180000}%
\pgfsetlinewidth{0.000000pt}%
\definecolor{currentstroke}{rgb}{0.000000,0.000000,0.000000}%
\pgfsetstrokecolor{currentstroke}%
\pgfsetdash{}{0pt}%
\pgfpathmoveto{\pgfqpoint{1.632723in}{1.347483in}}%
\pgfpathlineto{\pgfqpoint{1.647279in}{1.370738in}}%
\pgfpathlineto{\pgfqpoint{1.647907in}{1.393169in}}%
\pgfpathlineto{\pgfqpoint{1.633335in}{1.369935in}}%
\pgfpathlineto{\pgfqpoint{1.632723in}{1.347483in}}%
\pgfpathclose%
\pgfusepath{fill}%
\end{pgfscope}%
\begin{pgfscope}%
\pgfpathrectangle{\pgfqpoint{0.407592in}{0.222904in}}{\pgfqpoint{1.411163in}{1.411163in}}%
\pgfusepath{clip}%
\pgfsetbuttcap%
\pgfsetroundjoin%
\definecolor{currentfill}{rgb}{0.168627,0.611765,0.411765}%
\pgfsetfillcolor{currentfill}%
\pgfsetfillopacity{0.180000}%
\pgfsetlinewidth{0.000000pt}%
\definecolor{currentstroke}{rgb}{0.000000,0.000000,0.000000}%
\pgfsetstrokecolor{currentstroke}%
\pgfsetdash{}{0pt}%
\pgfpathmoveto{\pgfqpoint{1.570313in}{1.107305in}}%
\pgfpathlineto{\pgfqpoint{1.584916in}{1.126900in}}%
\pgfpathlineto{\pgfqpoint{1.585469in}{1.149049in}}%
\pgfpathlineto{\pgfqpoint{1.570849in}{1.129485in}}%
\pgfpathlineto{\pgfqpoint{1.570313in}{1.107305in}}%
\pgfpathclose%
\pgfusepath{fill}%
\end{pgfscope}%
\begin{pgfscope}%
\pgfpathrectangle{\pgfqpoint{0.407592in}{0.222904in}}{\pgfqpoint{1.411163in}{1.411163in}}%
\pgfusepath{clip}%
\pgfsetbuttcap%
\pgfsetroundjoin%
\definecolor{currentfill}{rgb}{0.168627,0.611765,0.411765}%
\pgfsetfillcolor{currentfill}%
\pgfsetfillopacity{0.180000}%
\pgfsetlinewidth{0.000000pt}%
\definecolor{currentstroke}{rgb}{0.000000,0.000000,0.000000}%
\pgfsetstrokecolor{currentstroke}%
\pgfsetdash{}{0pt}%
\pgfpathmoveto{\pgfqpoint{1.586023in}{1.171252in}}%
\pgfpathlineto{\pgfqpoint{1.600618in}{1.191877in}}%
\pgfpathlineto{\pgfqpoint{1.601190in}{1.214106in}}%
\pgfpathlineto{\pgfqpoint{1.586578in}{1.193510in}}%
\pgfpathlineto{\pgfqpoint{1.586023in}{1.171252in}}%
\pgfpathclose%
\pgfusepath{fill}%
\end{pgfscope}%
\begin{pgfscope}%
\pgfpathrectangle{\pgfqpoint{0.407592in}{0.222904in}}{\pgfqpoint{1.411163in}{1.411163in}}%
\pgfusepath{clip}%
\pgfsetbuttcap%
\pgfsetroundjoin%
\definecolor{currentfill}{rgb}{0.168627,0.611765,0.411765}%
\pgfsetfillcolor{currentfill}%
\pgfsetfillopacity{0.180000}%
\pgfsetlinewidth{0.000000pt}%
\definecolor{currentstroke}{rgb}{0.000000,0.000000,0.000000}%
\pgfsetstrokecolor{currentstroke}%
\pgfsetdash{}{0pt}%
\pgfpathmoveto{\pgfqpoint{1.554632in}{1.044623in}}%
\pgfpathlineto{\pgfqpoint{1.569245in}{1.063108in}}%
\pgfpathlineto{\pgfqpoint{1.569778in}{1.085180in}}%
\pgfpathlineto{\pgfqpoint{1.555148in}{1.066728in}}%
\pgfpathlineto{\pgfqpoint{1.554632in}{1.044623in}}%
\pgfpathclose%
\pgfusepath{fill}%
\end{pgfscope}%
\begin{pgfscope}%
\pgfpathrectangle{\pgfqpoint{0.407592in}{0.222904in}}{\pgfqpoint{1.411163in}{1.411163in}}%
\pgfusepath{clip}%
\pgfsetbuttcap%
\pgfsetroundjoin%
\definecolor{currentfill}{rgb}{0.168627,0.611765,0.411765}%
\pgfsetfillcolor{currentfill}%
\pgfsetfillopacity{0.180000}%
\pgfsetlinewidth{0.000000pt}%
\definecolor{currentstroke}{rgb}{0.000000,0.000000,0.000000}%
\pgfsetstrokecolor{currentstroke}%
\pgfsetdash{}{0pt}%
\pgfpathmoveto{\pgfqpoint{1.601763in}{1.236389in}}%
\pgfpathlineto{\pgfqpoint{1.616350in}{1.257963in}}%
\pgfpathlineto{\pgfqpoint{1.616941in}{1.280275in}}%
\pgfpathlineto{\pgfqpoint{1.602337in}{1.258728in}}%
\pgfpathlineto{\pgfqpoint{1.601763in}{1.236389in}}%
\pgfpathclose%
\pgfusepath{fill}%
\end{pgfscope}%
\begin{pgfscope}%
\pgfpathrectangle{\pgfqpoint{0.407592in}{0.222904in}}{\pgfqpoint{1.411163in}{1.411163in}}%
\pgfusepath{clip}%
\pgfsetbuttcap%
\pgfsetroundjoin%
\definecolor{currentfill}{rgb}{0.168627,0.611765,0.411765}%
\pgfsetfillcolor{currentfill}%
\pgfsetfillopacity{0.180000}%
\pgfsetlinewidth{0.000000pt}%
\definecolor{currentstroke}{rgb}{0.000000,0.000000,0.000000}%
\pgfsetstrokecolor{currentstroke}%
\pgfsetdash{}{0pt}%
\pgfpathmoveto{\pgfqpoint{1.538978in}{0.983275in}}%
\pgfpathlineto{\pgfqpoint{1.553602in}{1.000574in}}%
\pgfpathlineto{\pgfqpoint{1.554116in}{1.022572in}}%
\pgfpathlineto{\pgfqpoint{1.539476in}{1.005309in}}%
\pgfpathlineto{\pgfqpoint{1.538978in}{0.983275in}}%
\pgfpathclose%
\pgfusepath{fill}%
\end{pgfscope}%
\begin{pgfscope}%
\pgfpathrectangle{\pgfqpoint{0.407592in}{0.222904in}}{\pgfqpoint{1.411163in}{1.411163in}}%
\pgfusepath{clip}%
\pgfsetbuttcap%
\pgfsetroundjoin%
\definecolor{currentfill}{rgb}{0.168627,0.611765,0.411765}%
\pgfsetfillcolor{currentfill}%
\pgfsetfillopacity{0.180000}%
\pgfsetlinewidth{0.000000pt}%
\definecolor{currentstroke}{rgb}{0.000000,0.000000,0.000000}%
\pgfsetstrokecolor{currentstroke}%
\pgfsetdash{}{0pt}%
\pgfpathmoveto{\pgfqpoint{1.617533in}{1.302642in}}%
\pgfpathlineto{\pgfqpoint{1.632113in}{1.325085in}}%
\pgfpathlineto{\pgfqpoint{1.632723in}{1.347483in}}%
\pgfpathlineto{\pgfqpoint{1.618127in}{1.325064in}}%
\pgfpathlineto{\pgfqpoint{1.617533in}{1.302642in}}%
\pgfpathclose%
\pgfusepath{fill}%
\end{pgfscope}%
\begin{pgfscope}%
\pgfpathrectangle{\pgfqpoint{0.407592in}{0.222904in}}{\pgfqpoint{1.411163in}{1.411163in}}%
\pgfusepath{clip}%
\pgfsetbuttcap%
\pgfsetroundjoin%
\definecolor{currentfill}{rgb}{0.168627,0.611765,0.411765}%
\pgfsetfillcolor{currentfill}%
\pgfsetfillopacity{0.180000}%
\pgfsetlinewidth{0.000000pt}%
\definecolor{currentstroke}{rgb}{0.000000,0.000000,0.000000}%
\pgfsetstrokecolor{currentstroke}%
\pgfsetdash{}{0pt}%
\pgfpathmoveto{\pgfqpoint{1.570849in}{1.129485in}}%
\pgfpathlineto{\pgfqpoint{1.585469in}{1.149049in}}%
\pgfpathlineto{\pgfqpoint{1.586023in}{1.171252in}}%
\pgfpathlineto{\pgfqpoint{1.571386in}{1.151719in}}%
\pgfpathlineto{\pgfqpoint{1.570849in}{1.129485in}}%
\pgfpathclose%
\pgfusepath{fill}%
\end{pgfscope}%
\begin{pgfscope}%
\pgfpathrectangle{\pgfqpoint{0.407592in}{0.222904in}}{\pgfqpoint{1.411163in}{1.411163in}}%
\pgfusepath{clip}%
\pgfsetbuttcap%
\pgfsetroundjoin%
\definecolor{currentfill}{rgb}{0.168627,0.611765,0.411765}%
\pgfsetfillcolor{currentfill}%
\pgfsetfillopacity{0.180000}%
\pgfsetlinewidth{0.000000pt}%
\definecolor{currentstroke}{rgb}{0.000000,0.000000,0.000000}%
\pgfsetstrokecolor{currentstroke}%
\pgfsetdash{}{0pt}%
\pgfpathmoveto{\pgfqpoint{1.586578in}{1.193510in}}%
\pgfpathlineto{\pgfqpoint{1.601190in}{1.214106in}}%
\pgfpathlineto{\pgfqpoint{1.601763in}{1.236389in}}%
\pgfpathlineto{\pgfqpoint{1.587134in}{1.215822in}}%
\pgfpathlineto{\pgfqpoint{1.586578in}{1.193510in}}%
\pgfpathclose%
\pgfusepath{fill}%
\end{pgfscope}%
\begin{pgfscope}%
\pgfpathrectangle{\pgfqpoint{0.407592in}{0.222904in}}{\pgfqpoint{1.411163in}{1.411163in}}%
\pgfusepath{clip}%
\pgfsetbuttcap%
\pgfsetroundjoin%
\definecolor{currentfill}{rgb}{0.168627,0.611765,0.411765}%
\pgfsetfillcolor{currentfill}%
\pgfsetfillopacity{0.180000}%
\pgfsetlinewidth{0.000000pt}%
\definecolor{currentstroke}{rgb}{0.000000,0.000000,0.000000}%
\pgfsetstrokecolor{currentstroke}%
\pgfsetdash{}{0pt}%
\pgfpathmoveto{\pgfqpoint{1.555148in}{1.066728in}}%
\pgfpathlineto{\pgfqpoint{1.569778in}{1.085180in}}%
\pgfpathlineto{\pgfqpoint{1.570313in}{1.107305in}}%
\pgfpathlineto{\pgfqpoint{1.555666in}{1.088888in}}%
\pgfpathlineto{\pgfqpoint{1.555148in}{1.066728in}}%
\pgfpathclose%
\pgfusepath{fill}%
\end{pgfscope}%
\begin{pgfscope}%
\pgfpathrectangle{\pgfqpoint{0.407592in}{0.222904in}}{\pgfqpoint{1.411163in}{1.411163in}}%
\pgfusepath{clip}%
\pgfsetbuttcap%
\pgfsetroundjoin%
\definecolor{currentfill}{rgb}{0.168627,0.611765,0.411765}%
\pgfsetfillcolor{currentfill}%
\pgfsetfillopacity{0.180000}%
\pgfsetlinewidth{0.000000pt}%
\definecolor{currentstroke}{rgb}{0.000000,0.000000,0.000000}%
\pgfsetstrokecolor{currentstroke}%
\pgfsetdash{}{0pt}%
\pgfpathmoveto{\pgfqpoint{1.602337in}{1.258728in}}%
\pgfpathlineto{\pgfqpoint{1.616941in}{1.280275in}}%
\pgfpathlineto{\pgfqpoint{1.617533in}{1.302642in}}%
\pgfpathlineto{\pgfqpoint{1.602913in}{1.281120in}}%
\pgfpathlineto{\pgfqpoint{1.602337in}{1.258728in}}%
\pgfpathclose%
\pgfusepath{fill}%
\end{pgfscope}%
\begin{pgfscope}%
\pgfpathrectangle{\pgfqpoint{0.407592in}{0.222904in}}{\pgfqpoint{1.411163in}{1.411163in}}%
\pgfusepath{clip}%
\pgfsetbuttcap%
\pgfsetroundjoin%
\definecolor{currentfill}{rgb}{0.168627,0.611765,0.411765}%
\pgfsetfillcolor{currentfill}%
\pgfsetfillopacity{0.180000}%
\pgfsetlinewidth{0.000000pt}%
\definecolor{currentstroke}{rgb}{0.000000,0.000000,0.000000}%
\pgfsetstrokecolor{currentstroke}%
\pgfsetdash{}{0pt}%
\pgfpathmoveto{\pgfqpoint{1.539476in}{1.005309in}}%
\pgfpathlineto{\pgfqpoint{1.554116in}{1.022572in}}%
\pgfpathlineto{\pgfqpoint{1.554632in}{1.044623in}}%
\pgfpathlineto{\pgfqpoint{1.539975in}{1.027396in}}%
\pgfpathlineto{\pgfqpoint{1.539476in}{1.005309in}}%
\pgfpathclose%
\pgfusepath{fill}%
\end{pgfscope}%
\begin{pgfscope}%
\pgfpathrectangle{\pgfqpoint{0.407592in}{0.222904in}}{\pgfqpoint{1.411163in}{1.411163in}}%
\pgfusepath{clip}%
\pgfsetbuttcap%
\pgfsetroundjoin%
\definecolor{currentfill}{rgb}{0.168627,0.611765,0.411765}%
\pgfsetfillcolor{currentfill}%
\pgfsetfillopacity{0.180000}%
\pgfsetlinewidth{0.000000pt}%
\definecolor{currentstroke}{rgb}{0.000000,0.000000,0.000000}%
\pgfsetstrokecolor{currentstroke}%
\pgfsetdash{}{0pt}%
\pgfpathmoveto{\pgfqpoint{1.618127in}{1.325064in}}%
\pgfpathlineto{\pgfqpoint{1.632723in}{1.347483in}}%
\pgfpathlineto{\pgfqpoint{1.633335in}{1.369935in}}%
\pgfpathlineto{\pgfqpoint{1.618722in}{1.347540in}}%
\pgfpathlineto{\pgfqpoint{1.618127in}{1.325064in}}%
\pgfpathclose%
\pgfusepath{fill}%
\end{pgfscope}%
\begin{pgfscope}%
\pgfpathrectangle{\pgfqpoint{0.407592in}{0.222904in}}{\pgfqpoint{1.411163in}{1.411163in}}%
\pgfusepath{clip}%
\pgfsetbuttcap%
\pgfsetroundjoin%
\definecolor{currentfill}{rgb}{0.168627,0.611765,0.411765}%
\pgfsetfillcolor{currentfill}%
\pgfsetfillopacity{0.180000}%
\pgfsetlinewidth{0.000000pt}%
\definecolor{currentstroke}{rgb}{0.000000,0.000000,0.000000}%
\pgfsetstrokecolor{currentstroke}%
\pgfsetdash{}{0pt}%
\pgfpathmoveto{\pgfqpoint{1.523829in}{0.945288in}}%
\pgfpathlineto{\pgfqpoint{1.538481in}{0.961294in}}%
\pgfpathlineto{\pgfqpoint{1.538978in}{0.983275in}}%
\pgfpathlineto{\pgfqpoint{1.524308in}{0.967308in}}%
\pgfpathlineto{\pgfqpoint{1.523829in}{0.945288in}}%
\pgfpathclose%
\pgfusepath{fill}%
\end{pgfscope}%
\begin{pgfscope}%
\pgfpathrectangle{\pgfqpoint{0.407592in}{0.222904in}}{\pgfqpoint{1.411163in}{1.411163in}}%
\pgfusepath{clip}%
\pgfsetbuttcap%
\pgfsetroundjoin%
\definecolor{currentfill}{rgb}{0.168627,0.611765,0.411765}%
\pgfsetfillcolor{currentfill}%
\pgfsetfillopacity{0.180000}%
\pgfsetlinewidth{0.000000pt}%
\definecolor{currentstroke}{rgb}{0.000000,0.000000,0.000000}%
\pgfsetstrokecolor{currentstroke}%
\pgfsetdash{}{0pt}%
\pgfpathmoveto{\pgfqpoint{1.571386in}{1.151719in}}%
\pgfpathlineto{\pgfqpoint{1.586023in}{1.171252in}}%
\pgfpathlineto{\pgfqpoint{1.586578in}{1.193510in}}%
\pgfpathlineto{\pgfqpoint{1.571924in}{1.174008in}}%
\pgfpathlineto{\pgfqpoint{1.571386in}{1.151719in}}%
\pgfpathclose%
\pgfusepath{fill}%
\end{pgfscope}%
\begin{pgfscope}%
\pgfpathrectangle{\pgfqpoint{0.407592in}{0.222904in}}{\pgfqpoint{1.411163in}{1.411163in}}%
\pgfusepath{clip}%
\pgfsetbuttcap%
\pgfsetroundjoin%
\definecolor{currentfill}{rgb}{0.168627,0.611765,0.411765}%
\pgfsetfillcolor{currentfill}%
\pgfsetfillopacity{0.180000}%
\pgfsetlinewidth{0.000000pt}%
\definecolor{currentstroke}{rgb}{0.000000,0.000000,0.000000}%
\pgfsetstrokecolor{currentstroke}%
\pgfsetdash{}{0pt}%
\pgfpathmoveto{\pgfqpoint{1.587134in}{1.215822in}}%
\pgfpathlineto{\pgfqpoint{1.601763in}{1.236389in}}%
\pgfpathlineto{\pgfqpoint{1.602337in}{1.258728in}}%
\pgfpathlineto{\pgfqpoint{1.587692in}{1.238188in}}%
\pgfpathlineto{\pgfqpoint{1.587134in}{1.215822in}}%
\pgfpathclose%
\pgfusepath{fill}%
\end{pgfscope}%
\begin{pgfscope}%
\pgfpathrectangle{\pgfqpoint{0.407592in}{0.222904in}}{\pgfqpoint{1.411163in}{1.411163in}}%
\pgfusepath{clip}%
\pgfsetbuttcap%
\pgfsetroundjoin%
\definecolor{currentfill}{rgb}{0.168627,0.611765,0.411765}%
\pgfsetfillcolor{currentfill}%
\pgfsetfillopacity{0.180000}%
\pgfsetlinewidth{0.000000pt}%
\definecolor{currentstroke}{rgb}{0.000000,0.000000,0.000000}%
\pgfsetstrokecolor{currentstroke}%
\pgfsetdash{}{0pt}%
\pgfpathmoveto{\pgfqpoint{1.555666in}{1.088888in}}%
\pgfpathlineto{\pgfqpoint{1.570313in}{1.107305in}}%
\pgfpathlineto{\pgfqpoint{1.570849in}{1.129485in}}%
\pgfpathlineto{\pgfqpoint{1.556186in}{1.111101in}}%
\pgfpathlineto{\pgfqpoint{1.555666in}{1.088888in}}%
\pgfpathclose%
\pgfusepath{fill}%
\end{pgfscope}%
\begin{pgfscope}%
\pgfpathrectangle{\pgfqpoint{0.407592in}{0.222904in}}{\pgfqpoint{1.411163in}{1.411163in}}%
\pgfusepath{clip}%
\pgfsetbuttcap%
\pgfsetroundjoin%
\definecolor{currentfill}{rgb}{0.168627,0.611765,0.411765}%
\pgfsetfillcolor{currentfill}%
\pgfsetfillopacity{0.180000}%
\pgfsetlinewidth{0.000000pt}%
\definecolor{currentstroke}{rgb}{0.000000,0.000000,0.000000}%
\pgfsetstrokecolor{currentstroke}%
\pgfsetdash{}{0pt}%
\pgfpathmoveto{\pgfqpoint{1.602913in}{1.281120in}}%
\pgfpathlineto{\pgfqpoint{1.617533in}{1.302642in}}%
\pgfpathlineto{\pgfqpoint{1.618127in}{1.325064in}}%
\pgfpathlineto{\pgfqpoint{1.603490in}{1.303568in}}%
\pgfpathlineto{\pgfqpoint{1.602913in}{1.281120in}}%
\pgfpathclose%
\pgfusepath{fill}%
\end{pgfscope}%
\begin{pgfscope}%
\pgfpathrectangle{\pgfqpoint{0.407592in}{0.222904in}}{\pgfqpoint{1.411163in}{1.411163in}}%
\pgfusepath{clip}%
\pgfsetbuttcap%
\pgfsetroundjoin%
\definecolor{currentfill}{rgb}{0.168627,0.611765,0.411765}%
\pgfsetfillcolor{currentfill}%
\pgfsetfillopacity{0.180000}%
\pgfsetlinewidth{0.000000pt}%
\definecolor{currentstroke}{rgb}{0.000000,0.000000,0.000000}%
\pgfsetstrokecolor{currentstroke}%
\pgfsetdash{}{0pt}%
\pgfpathmoveto{\pgfqpoint{1.539975in}{1.027396in}}%
\pgfpathlineto{\pgfqpoint{1.554632in}{1.044623in}}%
\pgfpathlineto{\pgfqpoint{1.555148in}{1.066728in}}%
\pgfpathlineto{\pgfqpoint{1.540475in}{1.049539in}}%
\pgfpathlineto{\pgfqpoint{1.539975in}{1.027396in}}%
\pgfpathclose%
\pgfusepath{fill}%
\end{pgfscope}%
\begin{pgfscope}%
\pgfpathrectangle{\pgfqpoint{0.407592in}{0.222904in}}{\pgfqpoint{1.411163in}{1.411163in}}%
\pgfusepath{clip}%
\pgfsetbuttcap%
\pgfsetroundjoin%
\definecolor{currentfill}{rgb}{0.168627,0.611765,0.411765}%
\pgfsetfillcolor{currentfill}%
\pgfsetfillopacity{0.180000}%
\pgfsetlinewidth{0.000000pt}%
\definecolor{currentstroke}{rgb}{0.000000,0.000000,0.000000}%
\pgfsetstrokecolor{currentstroke}%
\pgfsetdash{}{0pt}%
\pgfpathmoveto{\pgfqpoint{1.524308in}{0.967308in}}%
\pgfpathlineto{\pgfqpoint{1.538978in}{0.983275in}}%
\pgfpathlineto{\pgfqpoint{1.539476in}{1.005309in}}%
\pgfpathlineto{\pgfqpoint{1.524790in}{0.989382in}}%
\pgfpathlineto{\pgfqpoint{1.524308in}{0.967308in}}%
\pgfpathclose%
\pgfusepath{fill}%
\end{pgfscope}%
\begin{pgfscope}%
\pgfpathrectangle{\pgfqpoint{0.407592in}{0.222904in}}{\pgfqpoint{1.411163in}{1.411163in}}%
\pgfusepath{clip}%
\pgfsetbuttcap%
\pgfsetroundjoin%
\definecolor{currentfill}{rgb}{0.168627,0.611765,0.411765}%
\pgfsetfillcolor{currentfill}%
\pgfsetfillopacity{0.180000}%
\pgfsetlinewidth{0.000000pt}%
\definecolor{currentstroke}{rgb}{0.000000,0.000000,0.000000}%
\pgfsetstrokecolor{currentstroke}%
\pgfsetdash{}{0pt}%
\pgfpathmoveto{\pgfqpoint{1.571924in}{1.174008in}}%
\pgfpathlineto{\pgfqpoint{1.586578in}{1.193510in}}%
\pgfpathlineto{\pgfqpoint{1.587134in}{1.215822in}}%
\pgfpathlineto{\pgfqpoint{1.572464in}{1.196352in}}%
\pgfpathlineto{\pgfqpoint{1.571924in}{1.174008in}}%
\pgfpathclose%
\pgfusepath{fill}%
\end{pgfscope}%
\begin{pgfscope}%
\pgfpathrectangle{\pgfqpoint{0.407592in}{0.222904in}}{\pgfqpoint{1.411163in}{1.411163in}}%
\pgfusepath{clip}%
\pgfsetbuttcap%
\pgfsetroundjoin%
\definecolor{currentfill}{rgb}{0.168627,0.611765,0.411765}%
\pgfsetfillcolor{currentfill}%
\pgfsetfillopacity{0.180000}%
\pgfsetlinewidth{0.000000pt}%
\definecolor{currentstroke}{rgb}{0.000000,0.000000,0.000000}%
\pgfsetstrokecolor{currentstroke}%
\pgfsetdash{}{0pt}%
\pgfpathmoveto{\pgfqpoint{1.587692in}{1.238188in}}%
\pgfpathlineto{\pgfqpoint{1.602337in}{1.258728in}}%
\pgfpathlineto{\pgfqpoint{1.602913in}{1.281120in}}%
\pgfpathlineto{\pgfqpoint{1.588251in}{1.260610in}}%
\pgfpathlineto{\pgfqpoint{1.587692in}{1.238188in}}%
\pgfpathclose%
\pgfusepath{fill}%
\end{pgfscope}%
\begin{pgfscope}%
\pgfpathrectangle{\pgfqpoint{0.407592in}{0.222904in}}{\pgfqpoint{1.411163in}{1.411163in}}%
\pgfusepath{clip}%
\pgfsetbuttcap%
\pgfsetroundjoin%
\definecolor{currentfill}{rgb}{0.168627,0.611765,0.411765}%
\pgfsetfillcolor{currentfill}%
\pgfsetfillopacity{0.180000}%
\pgfsetlinewidth{0.000000pt}%
\definecolor{currentstroke}{rgb}{0.000000,0.000000,0.000000}%
\pgfsetstrokecolor{currentstroke}%
\pgfsetdash{}{0pt}%
\pgfpathmoveto{\pgfqpoint{1.556186in}{1.111101in}}%
\pgfpathlineto{\pgfqpoint{1.570849in}{1.129485in}}%
\pgfpathlineto{\pgfqpoint{1.571386in}{1.151719in}}%
\pgfpathlineto{\pgfqpoint{1.556706in}{1.133370in}}%
\pgfpathlineto{\pgfqpoint{1.556186in}{1.111101in}}%
\pgfpathclose%
\pgfusepath{fill}%
\end{pgfscope}%
\begin{pgfscope}%
\pgfpathrectangle{\pgfqpoint{0.407592in}{0.222904in}}{\pgfqpoint{1.411163in}{1.411163in}}%
\pgfusepath{clip}%
\pgfsetbuttcap%
\pgfsetroundjoin%
\definecolor{currentfill}{rgb}{0.168627,0.611765,0.411765}%
\pgfsetfillcolor{currentfill}%
\pgfsetfillopacity{0.180000}%
\pgfsetlinewidth{0.000000pt}%
\definecolor{currentstroke}{rgb}{0.000000,0.000000,0.000000}%
\pgfsetstrokecolor{currentstroke}%
\pgfsetdash{}{0pt}%
\pgfpathmoveto{\pgfqpoint{1.603490in}{1.303568in}}%
\pgfpathlineto{\pgfqpoint{1.618127in}{1.325064in}}%
\pgfpathlineto{\pgfqpoint{1.618722in}{1.347540in}}%
\pgfpathlineto{\pgfqpoint{1.604069in}{1.326071in}}%
\pgfpathlineto{\pgfqpoint{1.603490in}{1.303568in}}%
\pgfpathclose%
\pgfusepath{fill}%
\end{pgfscope}%
\begin{pgfscope}%
\pgfpathrectangle{\pgfqpoint{0.407592in}{0.222904in}}{\pgfqpoint{1.411163in}{1.411163in}}%
\pgfusepath{clip}%
\pgfsetbuttcap%
\pgfsetroundjoin%
\definecolor{currentfill}{rgb}{0.168627,0.611765,0.411765}%
\pgfsetfillcolor{currentfill}%
\pgfsetfillopacity{0.180000}%
\pgfsetlinewidth{0.000000pt}%
\definecolor{currentstroke}{rgb}{0.000000,0.000000,0.000000}%
\pgfsetstrokecolor{currentstroke}%
\pgfsetdash{}{0pt}%
\pgfpathmoveto{\pgfqpoint{1.540475in}{1.049539in}}%
\pgfpathlineto{\pgfqpoint{1.555148in}{1.066728in}}%
\pgfpathlineto{\pgfqpoint{1.555666in}{1.088888in}}%
\pgfpathlineto{\pgfqpoint{1.540976in}{1.071735in}}%
\pgfpathlineto{\pgfqpoint{1.540475in}{1.049539in}}%
\pgfpathclose%
\pgfusepath{fill}%
\end{pgfscope}%
\begin{pgfscope}%
\pgfpathrectangle{\pgfqpoint{0.407592in}{0.222904in}}{\pgfqpoint{1.411163in}{1.411163in}}%
\pgfusepath{clip}%
\pgfsetbuttcap%
\pgfsetroundjoin%
\definecolor{currentfill}{rgb}{0.168627,0.611765,0.411765}%
\pgfsetfillcolor{currentfill}%
\pgfsetfillopacity{0.180000}%
\pgfsetlinewidth{0.000000pt}%
\definecolor{currentstroke}{rgb}{0.000000,0.000000,0.000000}%
\pgfsetstrokecolor{currentstroke}%
\pgfsetdash{}{0pt}%
\pgfpathmoveto{\pgfqpoint{1.524790in}{0.989382in}}%
\pgfpathlineto{\pgfqpoint{1.539476in}{1.005309in}}%
\pgfpathlineto{\pgfqpoint{1.539975in}{1.027396in}}%
\pgfpathlineto{\pgfqpoint{1.525272in}{1.011510in}}%
\pgfpathlineto{\pgfqpoint{1.524790in}{0.989382in}}%
\pgfpathclose%
\pgfusepath{fill}%
\end{pgfscope}%
\begin{pgfscope}%
\pgfpathrectangle{\pgfqpoint{0.407592in}{0.222904in}}{\pgfqpoint{1.411163in}{1.411163in}}%
\pgfusepath{clip}%
\pgfsetbuttcap%
\pgfsetroundjoin%
\definecolor{currentfill}{rgb}{0.168627,0.611765,0.411765}%
\pgfsetfillcolor{currentfill}%
\pgfsetfillopacity{0.180000}%
\pgfsetlinewidth{0.000000pt}%
\definecolor{currentstroke}{rgb}{0.000000,0.000000,0.000000}%
\pgfsetstrokecolor{currentstroke}%
\pgfsetdash{}{0pt}%
\pgfpathmoveto{\pgfqpoint{1.509128in}{0.930684in}}%
\pgfpathlineto{\pgfqpoint{1.523829in}{0.945288in}}%
\pgfpathlineto{\pgfqpoint{1.524308in}{0.967308in}}%
\pgfpathlineto{\pgfqpoint{1.509591in}{0.952748in}}%
\pgfpathlineto{\pgfqpoint{1.509128in}{0.930684in}}%
\pgfpathclose%
\pgfusepath{fill}%
\end{pgfscope}%
\begin{pgfscope}%
\pgfpathrectangle{\pgfqpoint{0.407592in}{0.222904in}}{\pgfqpoint{1.411163in}{1.411163in}}%
\pgfusepath{clip}%
\pgfsetbuttcap%
\pgfsetroundjoin%
\definecolor{currentfill}{rgb}{0.168627,0.611765,0.411765}%
\pgfsetfillcolor{currentfill}%
\pgfsetfillopacity{0.180000}%
\pgfsetlinewidth{0.000000pt}%
\definecolor{currentstroke}{rgb}{0.000000,0.000000,0.000000}%
\pgfsetstrokecolor{currentstroke}%
\pgfsetdash{}{0pt}%
\pgfpathmoveto{\pgfqpoint{1.572464in}{1.196352in}}%
\pgfpathlineto{\pgfqpoint{1.587134in}{1.215822in}}%
\pgfpathlineto{\pgfqpoint{1.587692in}{1.238188in}}%
\pgfpathlineto{\pgfqpoint{1.573006in}{1.218750in}}%
\pgfpathlineto{\pgfqpoint{1.572464in}{1.196352in}}%
\pgfpathclose%
\pgfusepath{fill}%
\end{pgfscope}%
\begin{pgfscope}%
\pgfpathrectangle{\pgfqpoint{0.407592in}{0.222904in}}{\pgfqpoint{1.411163in}{1.411163in}}%
\pgfusepath{clip}%
\pgfsetbuttcap%
\pgfsetroundjoin%
\definecolor{currentfill}{rgb}{0.168627,0.611765,0.411765}%
\pgfsetfillcolor{currentfill}%
\pgfsetfillopacity{0.180000}%
\pgfsetlinewidth{0.000000pt}%
\definecolor{currentstroke}{rgb}{0.000000,0.000000,0.000000}%
\pgfsetstrokecolor{currentstroke}%
\pgfsetdash{}{0pt}%
\pgfpathmoveto{\pgfqpoint{1.588251in}{1.260610in}}%
\pgfpathlineto{\pgfqpoint{1.602913in}{1.281120in}}%
\pgfpathlineto{\pgfqpoint{1.603490in}{1.303568in}}%
\pgfpathlineto{\pgfqpoint{1.588812in}{1.283087in}}%
\pgfpathlineto{\pgfqpoint{1.588251in}{1.260610in}}%
\pgfpathclose%
\pgfusepath{fill}%
\end{pgfscope}%
\begin{pgfscope}%
\pgfpathrectangle{\pgfqpoint{0.407592in}{0.222904in}}{\pgfqpoint{1.411163in}{1.411163in}}%
\pgfusepath{clip}%
\pgfsetbuttcap%
\pgfsetroundjoin%
\definecolor{currentfill}{rgb}{0.168627,0.611765,0.411765}%
\pgfsetfillcolor{currentfill}%
\pgfsetfillopacity{0.180000}%
\pgfsetlinewidth{0.000000pt}%
\definecolor{currentstroke}{rgb}{0.000000,0.000000,0.000000}%
\pgfsetstrokecolor{currentstroke}%
\pgfsetdash{}{0pt}%
\pgfpathmoveto{\pgfqpoint{1.556706in}{1.133370in}}%
\pgfpathlineto{\pgfqpoint{1.571386in}{1.151719in}}%
\pgfpathlineto{\pgfqpoint{1.571924in}{1.174008in}}%
\pgfpathlineto{\pgfqpoint{1.557228in}{1.155693in}}%
\pgfpathlineto{\pgfqpoint{1.556706in}{1.133370in}}%
\pgfpathclose%
\pgfusepath{fill}%
\end{pgfscope}%
\begin{pgfscope}%
\pgfpathrectangle{\pgfqpoint{0.407592in}{0.222904in}}{\pgfqpoint{1.411163in}{1.411163in}}%
\pgfusepath{clip}%
\pgfsetbuttcap%
\pgfsetroundjoin%
\definecolor{currentfill}{rgb}{0.168627,0.611765,0.411765}%
\pgfsetfillcolor{currentfill}%
\pgfsetfillopacity{0.180000}%
\pgfsetlinewidth{0.000000pt}%
\definecolor{currentstroke}{rgb}{0.000000,0.000000,0.000000}%
\pgfsetstrokecolor{currentstroke}%
\pgfsetdash{}{0pt}%
\pgfpathmoveto{\pgfqpoint{1.540976in}{1.071735in}}%
\pgfpathlineto{\pgfqpoint{1.555666in}{1.088888in}}%
\pgfpathlineto{\pgfqpoint{1.556186in}{1.111101in}}%
\pgfpathlineto{\pgfqpoint{1.541479in}{1.093986in}}%
\pgfpathlineto{\pgfqpoint{1.540976in}{1.071735in}}%
\pgfpathclose%
\pgfusepath{fill}%
\end{pgfscope}%
\begin{pgfscope}%
\pgfpathrectangle{\pgfqpoint{0.407592in}{0.222904in}}{\pgfqpoint{1.411163in}{1.411163in}}%
\pgfusepath{clip}%
\pgfsetbuttcap%
\pgfsetroundjoin%
\definecolor{currentfill}{rgb}{0.168627,0.611765,0.411765}%
\pgfsetfillcolor{currentfill}%
\pgfsetfillopacity{0.180000}%
\pgfsetlinewidth{0.000000pt}%
\definecolor{currentstroke}{rgb}{0.000000,0.000000,0.000000}%
\pgfsetstrokecolor{currentstroke}%
\pgfsetdash{}{0pt}%
\pgfpathmoveto{\pgfqpoint{1.525272in}{1.011510in}}%
\pgfpathlineto{\pgfqpoint{1.539975in}{1.027396in}}%
\pgfpathlineto{\pgfqpoint{1.540475in}{1.049539in}}%
\pgfpathlineto{\pgfqpoint{1.525755in}{1.033693in}}%
\pgfpathlineto{\pgfqpoint{1.525272in}{1.011510in}}%
\pgfpathclose%
\pgfusepath{fill}%
\end{pgfscope}%
\begin{pgfscope}%
\pgfpathrectangle{\pgfqpoint{0.407592in}{0.222904in}}{\pgfqpoint{1.411163in}{1.411163in}}%
\pgfusepath{clip}%
\pgfsetbuttcap%
\pgfsetroundjoin%
\definecolor{currentfill}{rgb}{0.168627,0.611765,0.411765}%
\pgfsetfillcolor{currentfill}%
\pgfsetfillopacity{0.180000}%
\pgfsetlinewidth{0.000000pt}%
\definecolor{currentstroke}{rgb}{0.000000,0.000000,0.000000}%
\pgfsetstrokecolor{currentstroke}%
\pgfsetdash{}{0pt}%
\pgfpathmoveto{\pgfqpoint{1.509591in}{0.952748in}}%
\pgfpathlineto{\pgfqpoint{1.524308in}{0.967308in}}%
\pgfpathlineto{\pgfqpoint{1.524790in}{0.989382in}}%
\pgfpathlineto{\pgfqpoint{1.510055in}{0.974865in}}%
\pgfpathlineto{\pgfqpoint{1.509591in}{0.952748in}}%
\pgfpathclose%
\pgfusepath{fill}%
\end{pgfscope}%
\begin{pgfscope}%
\pgfpathrectangle{\pgfqpoint{0.407592in}{0.222904in}}{\pgfqpoint{1.411163in}{1.411163in}}%
\pgfusepath{clip}%
\pgfsetbuttcap%
\pgfsetroundjoin%
\definecolor{currentfill}{rgb}{0.168627,0.611765,0.411765}%
\pgfsetfillcolor{currentfill}%
\pgfsetfillopacity{0.180000}%
\pgfsetlinewidth{0.000000pt}%
\definecolor{currentstroke}{rgb}{0.000000,0.000000,0.000000}%
\pgfsetstrokecolor{currentstroke}%
\pgfsetdash{}{0pt}%
\pgfpathmoveto{\pgfqpoint{1.573006in}{1.218750in}}%
\pgfpathlineto{\pgfqpoint{1.587692in}{1.238188in}}%
\pgfpathlineto{\pgfqpoint{1.588251in}{1.260610in}}%
\pgfpathlineto{\pgfqpoint{1.573548in}{1.241203in}}%
\pgfpathlineto{\pgfqpoint{1.573006in}{1.218750in}}%
\pgfpathclose%
\pgfusepath{fill}%
\end{pgfscope}%
\begin{pgfscope}%
\pgfpathrectangle{\pgfqpoint{0.407592in}{0.222904in}}{\pgfqpoint{1.411163in}{1.411163in}}%
\pgfusepath{clip}%
\pgfsetbuttcap%
\pgfsetroundjoin%
\definecolor{currentfill}{rgb}{0.168627,0.611765,0.411765}%
\pgfsetfillcolor{currentfill}%
\pgfsetfillopacity{0.180000}%
\pgfsetlinewidth{0.000000pt}%
\definecolor{currentstroke}{rgb}{0.000000,0.000000,0.000000}%
\pgfsetstrokecolor{currentstroke}%
\pgfsetdash{}{0pt}%
\pgfpathmoveto{\pgfqpoint{1.588812in}{1.283087in}}%
\pgfpathlineto{\pgfqpoint{1.603490in}{1.303568in}}%
\pgfpathlineto{\pgfqpoint{1.604069in}{1.326071in}}%
\pgfpathlineto{\pgfqpoint{1.589374in}{1.305619in}}%
\pgfpathlineto{\pgfqpoint{1.588812in}{1.283087in}}%
\pgfpathclose%
\pgfusepath{fill}%
\end{pgfscope}%
\begin{pgfscope}%
\pgfpathrectangle{\pgfqpoint{0.407592in}{0.222904in}}{\pgfqpoint{1.411163in}{1.411163in}}%
\pgfusepath{clip}%
\pgfsetbuttcap%
\pgfsetroundjoin%
\definecolor{currentfill}{rgb}{0.168627,0.611765,0.411765}%
\pgfsetfillcolor{currentfill}%
\pgfsetfillopacity{0.180000}%
\pgfsetlinewidth{0.000000pt}%
\definecolor{currentstroke}{rgb}{0.000000,0.000000,0.000000}%
\pgfsetstrokecolor{currentstroke}%
\pgfsetdash{}{0pt}%
\pgfpathmoveto{\pgfqpoint{1.557228in}{1.155693in}}%
\pgfpathlineto{\pgfqpoint{1.571924in}{1.174008in}}%
\pgfpathlineto{\pgfqpoint{1.572464in}{1.196352in}}%
\pgfpathlineto{\pgfqpoint{1.557751in}{1.178071in}}%
\pgfpathlineto{\pgfqpoint{1.557228in}{1.155693in}}%
\pgfpathclose%
\pgfusepath{fill}%
\end{pgfscope}%
\begin{pgfscope}%
\pgfpathrectangle{\pgfqpoint{0.407592in}{0.222904in}}{\pgfqpoint{1.411163in}{1.411163in}}%
\pgfusepath{clip}%
\pgfsetbuttcap%
\pgfsetroundjoin%
\definecolor{currentfill}{rgb}{0.168627,0.611765,0.411765}%
\pgfsetfillcolor{currentfill}%
\pgfsetfillopacity{0.180000}%
\pgfsetlinewidth{0.000000pt}%
\definecolor{currentstroke}{rgb}{0.000000,0.000000,0.000000}%
\pgfsetstrokecolor{currentstroke}%
\pgfsetdash{}{0pt}%
\pgfpathmoveto{\pgfqpoint{1.541479in}{1.093986in}}%
\pgfpathlineto{\pgfqpoint{1.556186in}{1.111101in}}%
\pgfpathlineto{\pgfqpoint{1.556706in}{1.133370in}}%
\pgfpathlineto{\pgfqpoint{1.541983in}{1.116292in}}%
\pgfpathlineto{\pgfqpoint{1.541479in}{1.093986in}}%
\pgfpathclose%
\pgfusepath{fill}%
\end{pgfscope}%
\begin{pgfscope}%
\pgfpathrectangle{\pgfqpoint{0.407592in}{0.222904in}}{\pgfqpoint{1.411163in}{1.411163in}}%
\pgfusepath{clip}%
\pgfsetbuttcap%
\pgfsetroundjoin%
\definecolor{currentfill}{rgb}{0.168627,0.611765,0.411765}%
\pgfsetfillcolor{currentfill}%
\pgfsetfillopacity{0.180000}%
\pgfsetlinewidth{0.000000pt}%
\definecolor{currentstroke}{rgb}{0.000000,0.000000,0.000000}%
\pgfsetstrokecolor{currentstroke}%
\pgfsetdash{}{0pt}%
\pgfpathmoveto{\pgfqpoint{1.525755in}{1.033693in}}%
\pgfpathlineto{\pgfqpoint{1.540475in}{1.049539in}}%
\pgfpathlineto{\pgfqpoint{1.540976in}{1.071735in}}%
\pgfpathlineto{\pgfqpoint{1.526240in}{1.055930in}}%
\pgfpathlineto{\pgfqpoint{1.525755in}{1.033693in}}%
\pgfpathclose%
\pgfusepath{fill}%
\end{pgfscope}%
\begin{pgfscope}%
\pgfpathrectangle{\pgfqpoint{0.407592in}{0.222904in}}{\pgfqpoint{1.411163in}{1.411163in}}%
\pgfusepath{clip}%
\pgfsetbuttcap%
\pgfsetroundjoin%
\definecolor{currentfill}{rgb}{0.168627,0.611765,0.411765}%
\pgfsetfillcolor{currentfill}%
\pgfsetfillopacity{0.180000}%
\pgfsetlinewidth{0.000000pt}%
\definecolor{currentstroke}{rgb}{0.000000,0.000000,0.000000}%
\pgfsetstrokecolor{currentstroke}%
\pgfsetdash{}{0pt}%
\pgfpathmoveto{\pgfqpoint{1.510055in}{0.974865in}}%
\pgfpathlineto{\pgfqpoint{1.524790in}{0.989382in}}%
\pgfpathlineto{\pgfqpoint{1.525272in}{1.011510in}}%
\pgfpathlineto{\pgfqpoint{1.510521in}{0.997038in}}%
\pgfpathlineto{\pgfqpoint{1.510055in}{0.974865in}}%
\pgfpathclose%
\pgfusepath{fill}%
\end{pgfscope}%
\begin{pgfscope}%
\pgfpathrectangle{\pgfqpoint{0.407592in}{0.222904in}}{\pgfqpoint{1.411163in}{1.411163in}}%
\pgfusepath{clip}%
\pgfsetbuttcap%
\pgfsetroundjoin%
\definecolor{currentfill}{rgb}{0.168627,0.611765,0.411765}%
\pgfsetfillcolor{currentfill}%
\pgfsetfillopacity{0.180000}%
\pgfsetlinewidth{0.000000pt}%
\definecolor{currentstroke}{rgb}{0.000000,0.000000,0.000000}%
\pgfsetstrokecolor{currentstroke}%
\pgfsetdash{}{0pt}%
\pgfpathmoveto{\pgfqpoint{1.494376in}{0.917544in}}%
\pgfpathlineto{\pgfqpoint{1.509128in}{0.930684in}}%
\pgfpathlineto{\pgfqpoint{1.509591in}{0.952748in}}%
\pgfpathlineto{\pgfqpoint{1.494822in}{0.939655in}}%
\pgfpathlineto{\pgfqpoint{1.494376in}{0.917544in}}%
\pgfpathclose%
\pgfusepath{fill}%
\end{pgfscope}%
\begin{pgfscope}%
\pgfpathrectangle{\pgfqpoint{0.407592in}{0.222904in}}{\pgfqpoint{1.411163in}{1.411163in}}%
\pgfusepath{clip}%
\pgfsetbuttcap%
\pgfsetroundjoin%
\definecolor{currentfill}{rgb}{0.168627,0.611765,0.411765}%
\pgfsetfillcolor{currentfill}%
\pgfsetfillopacity{0.180000}%
\pgfsetlinewidth{0.000000pt}%
\definecolor{currentstroke}{rgb}{0.000000,0.000000,0.000000}%
\pgfsetstrokecolor{currentstroke}%
\pgfsetdash{}{0pt}%
\pgfpathmoveto{\pgfqpoint{1.573548in}{1.241203in}}%
\pgfpathlineto{\pgfqpoint{1.588251in}{1.260610in}}%
\pgfpathlineto{\pgfqpoint{1.588812in}{1.283087in}}%
\pgfpathlineto{\pgfqpoint{1.574092in}{1.263712in}}%
\pgfpathlineto{\pgfqpoint{1.573548in}{1.241203in}}%
\pgfpathclose%
\pgfusepath{fill}%
\end{pgfscope}%
\begin{pgfscope}%
\pgfpathrectangle{\pgfqpoint{0.407592in}{0.222904in}}{\pgfqpoint{1.411163in}{1.411163in}}%
\pgfusepath{clip}%
\pgfsetbuttcap%
\pgfsetroundjoin%
\definecolor{currentfill}{rgb}{0.168627,0.611765,0.411765}%
\pgfsetfillcolor{currentfill}%
\pgfsetfillopacity{0.180000}%
\pgfsetlinewidth{0.000000pt}%
\definecolor{currentstroke}{rgb}{0.000000,0.000000,0.000000}%
\pgfsetstrokecolor{currentstroke}%
\pgfsetdash{}{0pt}%
\pgfpathmoveto{\pgfqpoint{1.557751in}{1.178071in}}%
\pgfpathlineto{\pgfqpoint{1.572464in}{1.196352in}}%
\pgfpathlineto{\pgfqpoint{1.573006in}{1.218750in}}%
\pgfpathlineto{\pgfqpoint{1.558276in}{1.200504in}}%
\pgfpathlineto{\pgfqpoint{1.557751in}{1.178071in}}%
\pgfpathclose%
\pgfusepath{fill}%
\end{pgfscope}%
\begin{pgfscope}%
\pgfpathrectangle{\pgfqpoint{0.407592in}{0.222904in}}{\pgfqpoint{1.411163in}{1.411163in}}%
\pgfusepath{clip}%
\pgfsetbuttcap%
\pgfsetroundjoin%
\definecolor{currentfill}{rgb}{0.168627,0.611765,0.411765}%
\pgfsetfillcolor{currentfill}%
\pgfsetfillopacity{0.180000}%
\pgfsetlinewidth{0.000000pt}%
\definecolor{currentstroke}{rgb}{0.000000,0.000000,0.000000}%
\pgfsetstrokecolor{currentstroke}%
\pgfsetdash{}{0pt}%
\pgfpathmoveto{\pgfqpoint{1.541983in}{1.116292in}}%
\pgfpathlineto{\pgfqpoint{1.556706in}{1.133370in}}%
\pgfpathlineto{\pgfqpoint{1.557228in}{1.155693in}}%
\pgfpathlineto{\pgfqpoint{1.542488in}{1.138653in}}%
\pgfpathlineto{\pgfqpoint{1.541983in}{1.116292in}}%
\pgfpathclose%
\pgfusepath{fill}%
\end{pgfscope}%
\begin{pgfscope}%
\pgfpathrectangle{\pgfqpoint{0.407592in}{0.222904in}}{\pgfqpoint{1.411163in}{1.411163in}}%
\pgfusepath{clip}%
\pgfsetbuttcap%
\pgfsetroundjoin%
\definecolor{currentfill}{rgb}{0.168627,0.611765,0.411765}%
\pgfsetfillcolor{currentfill}%
\pgfsetfillopacity{0.180000}%
\pgfsetlinewidth{0.000000pt}%
\definecolor{currentstroke}{rgb}{0.000000,0.000000,0.000000}%
\pgfsetstrokecolor{currentstroke}%
\pgfsetdash{}{0pt}%
\pgfpathmoveto{\pgfqpoint{1.526240in}{1.055930in}}%
\pgfpathlineto{\pgfqpoint{1.540976in}{1.071735in}}%
\pgfpathlineto{\pgfqpoint{1.541479in}{1.093986in}}%
\pgfpathlineto{\pgfqpoint{1.526726in}{1.078222in}}%
\pgfpathlineto{\pgfqpoint{1.526240in}{1.055930in}}%
\pgfpathclose%
\pgfusepath{fill}%
\end{pgfscope}%
\begin{pgfscope}%
\pgfpathrectangle{\pgfqpoint{0.407592in}{0.222904in}}{\pgfqpoint{1.411163in}{1.411163in}}%
\pgfusepath{clip}%
\pgfsetbuttcap%
\pgfsetroundjoin%
\definecolor{currentfill}{rgb}{0.168627,0.611765,0.411765}%
\pgfsetfillcolor{currentfill}%
\pgfsetfillopacity{0.180000}%
\pgfsetlinewidth{0.000000pt}%
\definecolor{currentstroke}{rgb}{0.000000,0.000000,0.000000}%
\pgfsetstrokecolor{currentstroke}%
\pgfsetdash{}{0pt}%
\pgfpathmoveto{\pgfqpoint{1.510521in}{0.997038in}}%
\pgfpathlineto{\pgfqpoint{1.525272in}{1.011510in}}%
\pgfpathlineto{\pgfqpoint{1.525755in}{1.033693in}}%
\pgfpathlineto{\pgfqpoint{1.510987in}{1.019264in}}%
\pgfpathlineto{\pgfqpoint{1.510521in}{0.997038in}}%
\pgfpathclose%
\pgfusepath{fill}%
\end{pgfscope}%
\begin{pgfscope}%
\pgfpathrectangle{\pgfqpoint{0.407592in}{0.222904in}}{\pgfqpoint{1.411163in}{1.411163in}}%
\pgfusepath{clip}%
\pgfsetbuttcap%
\pgfsetroundjoin%
\definecolor{currentfill}{rgb}{0.168627,0.611765,0.411765}%
\pgfsetfillcolor{currentfill}%
\pgfsetfillopacity{0.180000}%
\pgfsetlinewidth{0.000000pt}%
\definecolor{currentstroke}{rgb}{0.000000,0.000000,0.000000}%
\pgfsetstrokecolor{currentstroke}%
\pgfsetdash{}{0pt}%
\pgfpathmoveto{\pgfqpoint{1.494822in}{0.939655in}}%
\pgfpathlineto{\pgfqpoint{1.509591in}{0.952748in}}%
\pgfpathlineto{\pgfqpoint{1.510055in}{0.974865in}}%
\pgfpathlineto{\pgfqpoint{1.495270in}{0.961820in}}%
\pgfpathlineto{\pgfqpoint{1.494822in}{0.939655in}}%
\pgfpathclose%
\pgfusepath{fill}%
\end{pgfscope}%
\begin{pgfscope}%
\pgfpathrectangle{\pgfqpoint{0.407592in}{0.222904in}}{\pgfqpoint{1.411163in}{1.411163in}}%
\pgfusepath{clip}%
\pgfsetbuttcap%
\pgfsetroundjoin%
\definecolor{currentfill}{rgb}{0.168627,0.611765,0.411765}%
\pgfsetfillcolor{currentfill}%
\pgfsetfillopacity{0.180000}%
\pgfsetlinewidth{0.000000pt}%
\definecolor{currentstroke}{rgb}{0.000000,0.000000,0.000000}%
\pgfsetstrokecolor{currentstroke}%
\pgfsetdash{}{0pt}%
\pgfpathmoveto{\pgfqpoint{1.574092in}{1.263712in}}%
\pgfpathlineto{\pgfqpoint{1.588812in}{1.283087in}}%
\pgfpathlineto{\pgfqpoint{1.589374in}{1.305619in}}%
\pgfpathlineto{\pgfqpoint{1.574637in}{1.286277in}}%
\pgfpathlineto{\pgfqpoint{1.574092in}{1.263712in}}%
\pgfpathclose%
\pgfusepath{fill}%
\end{pgfscope}%
\begin{pgfscope}%
\pgfpathrectangle{\pgfqpoint{0.407592in}{0.222904in}}{\pgfqpoint{1.411163in}{1.411163in}}%
\pgfusepath{clip}%
\pgfsetbuttcap%
\pgfsetroundjoin%
\definecolor{currentfill}{rgb}{0.168627,0.611765,0.411765}%
\pgfsetfillcolor{currentfill}%
\pgfsetfillopacity{0.180000}%
\pgfsetlinewidth{0.000000pt}%
\definecolor{currentstroke}{rgb}{0.000000,0.000000,0.000000}%
\pgfsetstrokecolor{currentstroke}%
\pgfsetdash{}{0pt}%
\pgfpathmoveto{\pgfqpoint{1.558276in}{1.200504in}}%
\pgfpathlineto{\pgfqpoint{1.573006in}{1.218750in}}%
\pgfpathlineto{\pgfqpoint{1.573548in}{1.241203in}}%
\pgfpathlineto{\pgfqpoint{1.558802in}{1.222993in}}%
\pgfpathlineto{\pgfqpoint{1.558276in}{1.200504in}}%
\pgfpathclose%
\pgfusepath{fill}%
\end{pgfscope}%
\begin{pgfscope}%
\pgfpathrectangle{\pgfqpoint{0.407592in}{0.222904in}}{\pgfqpoint{1.411163in}{1.411163in}}%
\pgfusepath{clip}%
\pgfsetbuttcap%
\pgfsetroundjoin%
\definecolor{currentfill}{rgb}{0.168627,0.611765,0.411765}%
\pgfsetfillcolor{currentfill}%
\pgfsetfillopacity{0.180000}%
\pgfsetlinewidth{0.000000pt}%
\definecolor{currentstroke}{rgb}{0.000000,0.000000,0.000000}%
\pgfsetstrokecolor{currentstroke}%
\pgfsetdash{}{0pt}%
\pgfpathmoveto{\pgfqpoint{1.542488in}{1.138653in}}%
\pgfpathlineto{\pgfqpoint{1.557228in}{1.155693in}}%
\pgfpathlineto{\pgfqpoint{1.557751in}{1.178071in}}%
\pgfpathlineto{\pgfqpoint{1.542994in}{1.161069in}}%
\pgfpathlineto{\pgfqpoint{1.542488in}{1.138653in}}%
\pgfpathclose%
\pgfusepath{fill}%
\end{pgfscope}%
\begin{pgfscope}%
\pgfpathrectangle{\pgfqpoint{0.407592in}{0.222904in}}{\pgfqpoint{1.411163in}{1.411163in}}%
\pgfusepath{clip}%
\pgfsetbuttcap%
\pgfsetroundjoin%
\definecolor{currentfill}{rgb}{0.168627,0.611765,0.411765}%
\pgfsetfillcolor{currentfill}%
\pgfsetfillopacity{0.180000}%
\pgfsetlinewidth{0.000000pt}%
\definecolor{currentstroke}{rgb}{0.000000,0.000000,0.000000}%
\pgfsetstrokecolor{currentstroke}%
\pgfsetdash{}{0pt}%
\pgfpathmoveto{\pgfqpoint{1.526726in}{1.078222in}}%
\pgfpathlineto{\pgfqpoint{1.541479in}{1.093986in}}%
\pgfpathlineto{\pgfqpoint{1.541983in}{1.116292in}}%
\pgfpathlineto{\pgfqpoint{1.527213in}{1.100569in}}%
\pgfpathlineto{\pgfqpoint{1.526726in}{1.078222in}}%
\pgfpathclose%
\pgfusepath{fill}%
\end{pgfscope}%
\begin{pgfscope}%
\pgfpathrectangle{\pgfqpoint{0.407592in}{0.222904in}}{\pgfqpoint{1.411163in}{1.411163in}}%
\pgfusepath{clip}%
\pgfsetbuttcap%
\pgfsetroundjoin%
\definecolor{currentfill}{rgb}{0.168627,0.611765,0.411765}%
\pgfsetfillcolor{currentfill}%
\pgfsetfillopacity{0.180000}%
\pgfsetlinewidth{0.000000pt}%
\definecolor{currentstroke}{rgb}{0.000000,0.000000,0.000000}%
\pgfsetstrokecolor{currentstroke}%
\pgfsetdash{}{0pt}%
\pgfpathmoveto{\pgfqpoint{1.510987in}{1.019264in}}%
\pgfpathlineto{\pgfqpoint{1.525755in}{1.033693in}}%
\pgfpathlineto{\pgfqpoint{1.526240in}{1.055930in}}%
\pgfpathlineto{\pgfqpoint{1.511455in}{1.041546in}}%
\pgfpathlineto{\pgfqpoint{1.510987in}{1.019264in}}%
\pgfpathclose%
\pgfusepath{fill}%
\end{pgfscope}%
\begin{pgfscope}%
\pgfpathrectangle{\pgfqpoint{0.407592in}{0.222904in}}{\pgfqpoint{1.411163in}{1.411163in}}%
\pgfusepath{clip}%
\pgfsetbuttcap%
\pgfsetroundjoin%
\definecolor{currentfill}{rgb}{0.168627,0.611765,0.411765}%
\pgfsetfillcolor{currentfill}%
\pgfsetfillopacity{0.180000}%
\pgfsetlinewidth{0.000000pt}%
\definecolor{currentstroke}{rgb}{0.000000,0.000000,0.000000}%
\pgfsetstrokecolor{currentstroke}%
\pgfsetdash{}{0pt}%
\pgfpathmoveto{\pgfqpoint{1.495270in}{0.961820in}}%
\pgfpathlineto{\pgfqpoint{1.510055in}{0.974865in}}%
\pgfpathlineto{\pgfqpoint{1.510521in}{0.997038in}}%
\pgfpathlineto{\pgfqpoint{1.495718in}{0.984040in}}%
\pgfpathlineto{\pgfqpoint{1.495270in}{0.961820in}}%
\pgfpathclose%
\pgfusepath{fill}%
\end{pgfscope}%
\begin{pgfscope}%
\pgfpathrectangle{\pgfqpoint{0.407592in}{0.222904in}}{\pgfqpoint{1.411163in}{1.411163in}}%
\pgfusepath{clip}%
\pgfsetbuttcap%
\pgfsetroundjoin%
\definecolor{currentfill}{rgb}{0.168627,0.611765,0.411765}%
\pgfsetfillcolor{currentfill}%
\pgfsetfillopacity{0.180000}%
\pgfsetlinewidth{0.000000pt}%
\definecolor{currentstroke}{rgb}{0.000000,0.000000,0.000000}%
\pgfsetstrokecolor{currentstroke}%
\pgfsetdash{}{0pt}%
\pgfpathmoveto{\pgfqpoint{1.558802in}{1.222993in}}%
\pgfpathlineto{\pgfqpoint{1.573548in}{1.241203in}}%
\pgfpathlineto{\pgfqpoint{1.574092in}{1.263712in}}%
\pgfpathlineto{\pgfqpoint{1.559329in}{1.245537in}}%
\pgfpathlineto{\pgfqpoint{1.558802in}{1.222993in}}%
\pgfpathclose%
\pgfusepath{fill}%
\end{pgfscope}%
\begin{pgfscope}%
\pgfpathrectangle{\pgfqpoint{0.407592in}{0.222904in}}{\pgfqpoint{1.411163in}{1.411163in}}%
\pgfusepath{clip}%
\pgfsetbuttcap%
\pgfsetroundjoin%
\definecolor{currentfill}{rgb}{0.168627,0.611765,0.411765}%
\pgfsetfillcolor{currentfill}%
\pgfsetfillopacity{0.180000}%
\pgfsetlinewidth{0.000000pt}%
\definecolor{currentstroke}{rgb}{0.000000,0.000000,0.000000}%
\pgfsetstrokecolor{currentstroke}%
\pgfsetdash{}{0pt}%
\pgfpathmoveto{\pgfqpoint{1.542994in}{1.161069in}}%
\pgfpathlineto{\pgfqpoint{1.557751in}{1.178071in}}%
\pgfpathlineto{\pgfqpoint{1.558276in}{1.200504in}}%
\pgfpathlineto{\pgfqpoint{1.543502in}{1.183540in}}%
\pgfpathlineto{\pgfqpoint{1.542994in}{1.161069in}}%
\pgfpathclose%
\pgfusepath{fill}%
\end{pgfscope}%
\begin{pgfscope}%
\pgfpathrectangle{\pgfqpoint{0.407592in}{0.222904in}}{\pgfqpoint{1.411163in}{1.411163in}}%
\pgfusepath{clip}%
\pgfsetbuttcap%
\pgfsetroundjoin%
\definecolor{currentfill}{rgb}{0.168627,0.611765,0.411765}%
\pgfsetfillcolor{currentfill}%
\pgfsetfillopacity{0.180000}%
\pgfsetlinewidth{0.000000pt}%
\definecolor{currentstroke}{rgb}{0.000000,0.000000,0.000000}%
\pgfsetstrokecolor{currentstroke}%
\pgfsetdash{}{0pt}%
\pgfpathmoveto{\pgfqpoint{1.479570in}{0.905913in}}%
\pgfpathlineto{\pgfqpoint{1.494376in}{0.917544in}}%
\pgfpathlineto{\pgfqpoint{1.494822in}{0.939655in}}%
\pgfpathlineto{\pgfqpoint{1.479999in}{0.928075in}}%
\pgfpathlineto{\pgfqpoint{1.479570in}{0.905913in}}%
\pgfpathclose%
\pgfusepath{fill}%
\end{pgfscope}%
\begin{pgfscope}%
\pgfpathrectangle{\pgfqpoint{0.407592in}{0.222904in}}{\pgfqpoint{1.411163in}{1.411163in}}%
\pgfusepath{clip}%
\pgfsetbuttcap%
\pgfsetroundjoin%
\definecolor{currentfill}{rgb}{0.168627,0.611765,0.411765}%
\pgfsetfillcolor{currentfill}%
\pgfsetfillopacity{0.180000}%
\pgfsetlinewidth{0.000000pt}%
\definecolor{currentstroke}{rgb}{0.000000,0.000000,0.000000}%
\pgfsetstrokecolor{currentstroke}%
\pgfsetdash{}{0pt}%
\pgfpathmoveto{\pgfqpoint{1.527213in}{1.100569in}}%
\pgfpathlineto{\pgfqpoint{1.541983in}{1.116292in}}%
\pgfpathlineto{\pgfqpoint{1.542488in}{1.138653in}}%
\pgfpathlineto{\pgfqpoint{1.527701in}{1.122971in}}%
\pgfpathlineto{\pgfqpoint{1.527213in}{1.100569in}}%
\pgfpathclose%
\pgfusepath{fill}%
\end{pgfscope}%
\begin{pgfscope}%
\pgfpathrectangle{\pgfqpoint{0.407592in}{0.222904in}}{\pgfqpoint{1.411163in}{1.411163in}}%
\pgfusepath{clip}%
\pgfsetbuttcap%
\pgfsetroundjoin%
\definecolor{currentfill}{rgb}{0.168627,0.611765,0.411765}%
\pgfsetfillcolor{currentfill}%
\pgfsetfillopacity{0.180000}%
\pgfsetlinewidth{0.000000pt}%
\definecolor{currentstroke}{rgb}{0.000000,0.000000,0.000000}%
\pgfsetstrokecolor{currentstroke}%
\pgfsetdash{}{0pt}%
\pgfpathmoveto{\pgfqpoint{1.511455in}{1.041546in}}%
\pgfpathlineto{\pgfqpoint{1.526240in}{1.055930in}}%
\pgfpathlineto{\pgfqpoint{1.526726in}{1.078222in}}%
\pgfpathlineto{\pgfqpoint{1.511924in}{1.063883in}}%
\pgfpathlineto{\pgfqpoint{1.511455in}{1.041546in}}%
\pgfpathclose%
\pgfusepath{fill}%
\end{pgfscope}%
\begin{pgfscope}%
\pgfpathrectangle{\pgfqpoint{0.407592in}{0.222904in}}{\pgfqpoint{1.411163in}{1.411163in}}%
\pgfusepath{clip}%
\pgfsetbuttcap%
\pgfsetroundjoin%
\definecolor{currentfill}{rgb}{0.168627,0.611765,0.411765}%
\pgfsetfillcolor{currentfill}%
\pgfsetfillopacity{0.180000}%
\pgfsetlinewidth{0.000000pt}%
\definecolor{currentstroke}{rgb}{0.000000,0.000000,0.000000}%
\pgfsetstrokecolor{currentstroke}%
\pgfsetdash{}{0pt}%
\pgfpathmoveto{\pgfqpoint{1.495718in}{0.984040in}}%
\pgfpathlineto{\pgfqpoint{1.510521in}{0.997038in}}%
\pgfpathlineto{\pgfqpoint{1.510987in}{1.019264in}}%
\pgfpathlineto{\pgfqpoint{1.496168in}{1.006315in}}%
\pgfpathlineto{\pgfqpoint{1.495718in}{0.984040in}}%
\pgfpathclose%
\pgfusepath{fill}%
\end{pgfscope}%
\begin{pgfscope}%
\pgfpathrectangle{\pgfqpoint{0.407592in}{0.222904in}}{\pgfqpoint{1.411163in}{1.411163in}}%
\pgfusepath{clip}%
\pgfsetbuttcap%
\pgfsetroundjoin%
\definecolor{currentfill}{rgb}{0.168627,0.611765,0.411765}%
\pgfsetfillcolor{currentfill}%
\pgfsetfillopacity{0.180000}%
\pgfsetlinewidth{0.000000pt}%
\definecolor{currentstroke}{rgb}{0.000000,0.000000,0.000000}%
\pgfsetstrokecolor{currentstroke}%
\pgfsetdash{}{0pt}%
\pgfpathmoveto{\pgfqpoint{1.559329in}{1.245537in}}%
\pgfpathlineto{\pgfqpoint{1.574092in}{1.263712in}}%
\pgfpathlineto{\pgfqpoint{1.574637in}{1.286277in}}%
\pgfpathlineto{\pgfqpoint{1.559857in}{1.268137in}}%
\pgfpathlineto{\pgfqpoint{1.559329in}{1.245537in}}%
\pgfpathclose%
\pgfusepath{fill}%
\end{pgfscope}%
\begin{pgfscope}%
\pgfpathrectangle{\pgfqpoint{0.407592in}{0.222904in}}{\pgfqpoint{1.411163in}{1.411163in}}%
\pgfusepath{clip}%
\pgfsetbuttcap%
\pgfsetroundjoin%
\definecolor{currentfill}{rgb}{0.168627,0.611765,0.411765}%
\pgfsetfillcolor{currentfill}%
\pgfsetfillopacity{0.180000}%
\pgfsetlinewidth{0.000000pt}%
\definecolor{currentstroke}{rgb}{0.000000,0.000000,0.000000}%
\pgfsetstrokecolor{currentstroke}%
\pgfsetdash{}{0pt}%
\pgfpathmoveto{\pgfqpoint{1.543502in}{1.183540in}}%
\pgfpathlineto{\pgfqpoint{1.558276in}{1.200504in}}%
\pgfpathlineto{\pgfqpoint{1.558802in}{1.222993in}}%
\pgfpathlineto{\pgfqpoint{1.544011in}{1.206067in}}%
\pgfpathlineto{\pgfqpoint{1.543502in}{1.183540in}}%
\pgfpathclose%
\pgfusepath{fill}%
\end{pgfscope}%
\begin{pgfscope}%
\pgfpathrectangle{\pgfqpoint{0.407592in}{0.222904in}}{\pgfqpoint{1.411163in}{1.411163in}}%
\pgfusepath{clip}%
\pgfsetbuttcap%
\pgfsetroundjoin%
\definecolor{currentfill}{rgb}{0.168627,0.611765,0.411765}%
\pgfsetfillcolor{currentfill}%
\pgfsetfillopacity{0.180000}%
\pgfsetlinewidth{0.000000pt}%
\definecolor{currentstroke}{rgb}{0.000000,0.000000,0.000000}%
\pgfsetstrokecolor{currentstroke}%
\pgfsetdash{}{0pt}%
\pgfpathmoveto{\pgfqpoint{1.479999in}{0.928075in}}%
\pgfpathlineto{\pgfqpoint{1.494822in}{0.939655in}}%
\pgfpathlineto{\pgfqpoint{1.495270in}{0.961820in}}%
\pgfpathlineto{\pgfqpoint{1.480429in}{0.950293in}}%
\pgfpathlineto{\pgfqpoint{1.479999in}{0.928075in}}%
\pgfpathclose%
\pgfusepath{fill}%
\end{pgfscope}%
\begin{pgfscope}%
\pgfpathrectangle{\pgfqpoint{0.407592in}{0.222904in}}{\pgfqpoint{1.411163in}{1.411163in}}%
\pgfusepath{clip}%
\pgfsetbuttcap%
\pgfsetroundjoin%
\definecolor{currentfill}{rgb}{0.168627,0.611765,0.411765}%
\pgfsetfillcolor{currentfill}%
\pgfsetfillopacity{0.180000}%
\pgfsetlinewidth{0.000000pt}%
\definecolor{currentstroke}{rgb}{0.000000,0.000000,0.000000}%
\pgfsetstrokecolor{currentstroke}%
\pgfsetdash{}{0pt}%
\pgfpathmoveto{\pgfqpoint{1.527701in}{1.122971in}}%
\pgfpathlineto{\pgfqpoint{1.542488in}{1.138653in}}%
\pgfpathlineto{\pgfqpoint{1.542994in}{1.161069in}}%
\pgfpathlineto{\pgfqpoint{1.528190in}{1.145429in}}%
\pgfpathlineto{\pgfqpoint{1.527701in}{1.122971in}}%
\pgfpathclose%
\pgfusepath{fill}%
\end{pgfscope}%
\begin{pgfscope}%
\pgfpathrectangle{\pgfqpoint{0.407592in}{0.222904in}}{\pgfqpoint{1.411163in}{1.411163in}}%
\pgfusepath{clip}%
\pgfsetbuttcap%
\pgfsetroundjoin%
\definecolor{currentfill}{rgb}{0.168627,0.611765,0.411765}%
\pgfsetfillcolor{currentfill}%
\pgfsetfillopacity{0.180000}%
\pgfsetlinewidth{0.000000pt}%
\definecolor{currentstroke}{rgb}{0.000000,0.000000,0.000000}%
\pgfsetstrokecolor{currentstroke}%
\pgfsetdash{}{0pt}%
\pgfpathmoveto{\pgfqpoint{1.511924in}{1.063883in}}%
\pgfpathlineto{\pgfqpoint{1.526726in}{1.078222in}}%
\pgfpathlineto{\pgfqpoint{1.527213in}{1.100569in}}%
\pgfpathlineto{\pgfqpoint{1.512394in}{1.086275in}}%
\pgfpathlineto{\pgfqpoint{1.511924in}{1.063883in}}%
\pgfpathclose%
\pgfusepath{fill}%
\end{pgfscope}%
\begin{pgfscope}%
\pgfpathrectangle{\pgfqpoint{0.407592in}{0.222904in}}{\pgfqpoint{1.411163in}{1.411163in}}%
\pgfusepath{clip}%
\pgfsetbuttcap%
\pgfsetroundjoin%
\definecolor{currentfill}{rgb}{0.168627,0.611765,0.411765}%
\pgfsetfillcolor{currentfill}%
\pgfsetfillopacity{0.180000}%
\pgfsetlinewidth{0.000000pt}%
\definecolor{currentstroke}{rgb}{0.000000,0.000000,0.000000}%
\pgfsetstrokecolor{currentstroke}%
\pgfsetdash{}{0pt}%
\pgfpathmoveto{\pgfqpoint{1.496168in}{1.006315in}}%
\pgfpathlineto{\pgfqpoint{1.510987in}{1.019264in}}%
\pgfpathlineto{\pgfqpoint{1.511455in}{1.041546in}}%
\pgfpathlineto{\pgfqpoint{1.496619in}{1.028646in}}%
\pgfpathlineto{\pgfqpoint{1.496168in}{1.006315in}}%
\pgfpathclose%
\pgfusepath{fill}%
\end{pgfscope}%
\begin{pgfscope}%
\pgfpathrectangle{\pgfqpoint{0.407592in}{0.222904in}}{\pgfqpoint{1.411163in}{1.411163in}}%
\pgfusepath{clip}%
\pgfsetbuttcap%
\pgfsetroundjoin%
\definecolor{currentfill}{rgb}{0.168627,0.611765,0.411765}%
\pgfsetfillcolor{currentfill}%
\pgfsetfillopacity{0.180000}%
\pgfsetlinewidth{0.000000pt}%
\definecolor{currentstroke}{rgb}{0.000000,0.000000,0.000000}%
\pgfsetstrokecolor{currentstroke}%
\pgfsetdash{}{0pt}%
\pgfpathmoveto{\pgfqpoint{1.544011in}{1.206067in}}%
\pgfpathlineto{\pgfqpoint{1.558802in}{1.222993in}}%
\pgfpathlineto{\pgfqpoint{1.559329in}{1.245537in}}%
\pgfpathlineto{\pgfqpoint{1.544521in}{1.228650in}}%
\pgfpathlineto{\pgfqpoint{1.544011in}{1.206067in}}%
\pgfpathclose%
\pgfusepath{fill}%
\end{pgfscope}%
\begin{pgfscope}%
\pgfpathrectangle{\pgfqpoint{0.407592in}{0.222904in}}{\pgfqpoint{1.411163in}{1.411163in}}%
\pgfusepath{clip}%
\pgfsetbuttcap%
\pgfsetroundjoin%
\definecolor{currentfill}{rgb}{0.168627,0.611765,0.411765}%
\pgfsetfillcolor{currentfill}%
\pgfsetfillopacity{0.180000}%
\pgfsetlinewidth{0.000000pt}%
\definecolor{currentstroke}{rgb}{0.000000,0.000000,0.000000}%
\pgfsetstrokecolor{currentstroke}%
\pgfsetdash{}{0pt}%
\pgfpathmoveto{\pgfqpoint{1.480429in}{0.950293in}}%
\pgfpathlineto{\pgfqpoint{1.495270in}{0.961820in}}%
\pgfpathlineto{\pgfqpoint{1.495718in}{0.984040in}}%
\pgfpathlineto{\pgfqpoint{1.480861in}{0.972565in}}%
\pgfpathlineto{\pgfqpoint{1.480429in}{0.950293in}}%
\pgfpathclose%
\pgfusepath{fill}%
\end{pgfscope}%
\begin{pgfscope}%
\pgfpathrectangle{\pgfqpoint{0.407592in}{0.222904in}}{\pgfqpoint{1.411163in}{1.411163in}}%
\pgfusepath{clip}%
\pgfsetbuttcap%
\pgfsetroundjoin%
\definecolor{currentfill}{rgb}{0.168627,0.611765,0.411765}%
\pgfsetfillcolor{currentfill}%
\pgfsetfillopacity{0.180000}%
\pgfsetlinewidth{0.000000pt}%
\definecolor{currentstroke}{rgb}{0.000000,0.000000,0.000000}%
\pgfsetstrokecolor{currentstroke}%
\pgfsetdash{}{0pt}%
\pgfpathmoveto{\pgfqpoint{1.528190in}{1.145429in}}%
\pgfpathlineto{\pgfqpoint{1.542994in}{1.161069in}}%
\pgfpathlineto{\pgfqpoint{1.543502in}{1.183540in}}%
\pgfpathlineto{\pgfqpoint{1.528681in}{1.167943in}}%
\pgfpathlineto{\pgfqpoint{1.528190in}{1.145429in}}%
\pgfpathclose%
\pgfusepath{fill}%
\end{pgfscope}%
\begin{pgfscope}%
\pgfpathrectangle{\pgfqpoint{0.407592in}{0.222904in}}{\pgfqpoint{1.411163in}{1.411163in}}%
\pgfusepath{clip}%
\pgfsetbuttcap%
\pgfsetroundjoin%
\definecolor{currentfill}{rgb}{0.168627,0.611765,0.411765}%
\pgfsetfillcolor{currentfill}%
\pgfsetfillopacity{0.180000}%
\pgfsetlinewidth{0.000000pt}%
\definecolor{currentstroke}{rgb}{0.000000,0.000000,0.000000}%
\pgfsetstrokecolor{currentstroke}%
\pgfsetdash{}{0pt}%
\pgfpathmoveto{\pgfqpoint{1.512394in}{1.086275in}}%
\pgfpathlineto{\pgfqpoint{1.527213in}{1.100569in}}%
\pgfpathlineto{\pgfqpoint{1.527701in}{1.122971in}}%
\pgfpathlineto{\pgfqpoint{1.512865in}{1.108723in}}%
\pgfpathlineto{\pgfqpoint{1.512394in}{1.086275in}}%
\pgfpathclose%
\pgfusepath{fill}%
\end{pgfscope}%
\begin{pgfscope}%
\pgfpathrectangle{\pgfqpoint{0.407592in}{0.222904in}}{\pgfqpoint{1.411163in}{1.411163in}}%
\pgfusepath{clip}%
\pgfsetbuttcap%
\pgfsetroundjoin%
\definecolor{currentfill}{rgb}{0.168627,0.611765,0.411765}%
\pgfsetfillcolor{currentfill}%
\pgfsetfillopacity{0.180000}%
\pgfsetlinewidth{0.000000pt}%
\definecolor{currentstroke}{rgb}{0.000000,0.000000,0.000000}%
\pgfsetstrokecolor{currentstroke}%
\pgfsetdash{}{0pt}%
\pgfpathmoveto{\pgfqpoint{1.464705in}{0.895822in}}%
\pgfpathlineto{\pgfqpoint{1.479570in}{0.905913in}}%
\pgfpathlineto{\pgfqpoint{1.479999in}{0.928075in}}%
\pgfpathlineto{\pgfqpoint{1.465117in}{0.918040in}}%
\pgfpathlineto{\pgfqpoint{1.464705in}{0.895822in}}%
\pgfpathclose%
\pgfusepath{fill}%
\end{pgfscope}%
\begin{pgfscope}%
\pgfpathrectangle{\pgfqpoint{0.407592in}{0.222904in}}{\pgfqpoint{1.411163in}{1.411163in}}%
\pgfusepath{clip}%
\pgfsetbuttcap%
\pgfsetroundjoin%
\definecolor{currentfill}{rgb}{0.168627,0.611765,0.411765}%
\pgfsetfillcolor{currentfill}%
\pgfsetfillopacity{0.180000}%
\pgfsetlinewidth{0.000000pt}%
\definecolor{currentstroke}{rgb}{0.000000,0.000000,0.000000}%
\pgfsetstrokecolor{currentstroke}%
\pgfsetdash{}{0pt}%
\pgfpathmoveto{\pgfqpoint{1.496619in}{1.028646in}}%
\pgfpathlineto{\pgfqpoint{1.511455in}{1.041546in}}%
\pgfpathlineto{\pgfqpoint{1.511924in}{1.063883in}}%
\pgfpathlineto{\pgfqpoint{1.497070in}{1.051031in}}%
\pgfpathlineto{\pgfqpoint{1.496619in}{1.028646in}}%
\pgfpathclose%
\pgfusepath{fill}%
\end{pgfscope}%
\begin{pgfscope}%
\pgfpathrectangle{\pgfqpoint{0.407592in}{0.222904in}}{\pgfqpoint{1.411163in}{1.411163in}}%
\pgfusepath{clip}%
\pgfsetbuttcap%
\pgfsetroundjoin%
\definecolor{currentfill}{rgb}{0.168627,0.611765,0.411765}%
\pgfsetfillcolor{currentfill}%
\pgfsetfillopacity{0.180000}%
\pgfsetlinewidth{0.000000pt}%
\definecolor{currentstroke}{rgb}{0.000000,0.000000,0.000000}%
\pgfsetstrokecolor{currentstroke}%
\pgfsetdash{}{0pt}%
\pgfpathmoveto{\pgfqpoint{1.544521in}{1.228650in}}%
\pgfpathlineto{\pgfqpoint{1.559329in}{1.245537in}}%
\pgfpathlineto{\pgfqpoint{1.559857in}{1.268137in}}%
\pgfpathlineto{\pgfqpoint{1.545032in}{1.251289in}}%
\pgfpathlineto{\pgfqpoint{1.544521in}{1.228650in}}%
\pgfpathclose%
\pgfusepath{fill}%
\end{pgfscope}%
\begin{pgfscope}%
\pgfpathrectangle{\pgfqpoint{0.407592in}{0.222904in}}{\pgfqpoint{1.411163in}{1.411163in}}%
\pgfusepath{clip}%
\pgfsetbuttcap%
\pgfsetroundjoin%
\definecolor{currentfill}{rgb}{0.168627,0.611765,0.411765}%
\pgfsetfillcolor{currentfill}%
\pgfsetfillopacity{0.180000}%
\pgfsetlinewidth{0.000000pt}%
\definecolor{currentstroke}{rgb}{0.000000,0.000000,0.000000}%
\pgfsetstrokecolor{currentstroke}%
\pgfsetdash{}{0pt}%
\pgfpathmoveto{\pgfqpoint{1.480861in}{0.972565in}}%
\pgfpathlineto{\pgfqpoint{1.495718in}{0.984040in}}%
\pgfpathlineto{\pgfqpoint{1.496168in}{1.006315in}}%
\pgfpathlineto{\pgfqpoint{1.481293in}{0.994893in}}%
\pgfpathlineto{\pgfqpoint{1.480861in}{0.972565in}}%
\pgfpathclose%
\pgfusepath{fill}%
\end{pgfscope}%
\begin{pgfscope}%
\pgfpathrectangle{\pgfqpoint{0.407592in}{0.222904in}}{\pgfqpoint{1.411163in}{1.411163in}}%
\pgfusepath{clip}%
\pgfsetbuttcap%
\pgfsetroundjoin%
\definecolor{currentfill}{rgb}{0.168627,0.611765,0.411765}%
\pgfsetfillcolor{currentfill}%
\pgfsetfillopacity{0.180000}%
\pgfsetlinewidth{0.000000pt}%
\definecolor{currentstroke}{rgb}{0.000000,0.000000,0.000000}%
\pgfsetstrokecolor{currentstroke}%
\pgfsetdash{}{0pt}%
\pgfpathmoveto{\pgfqpoint{1.528681in}{1.167943in}}%
\pgfpathlineto{\pgfqpoint{1.543502in}{1.183540in}}%
\pgfpathlineto{\pgfqpoint{1.544011in}{1.206067in}}%
\pgfpathlineto{\pgfqpoint{1.529173in}{1.190512in}}%
\pgfpathlineto{\pgfqpoint{1.528681in}{1.167943in}}%
\pgfpathclose%
\pgfusepath{fill}%
\end{pgfscope}%
\begin{pgfscope}%
\pgfpathrectangle{\pgfqpoint{0.407592in}{0.222904in}}{\pgfqpoint{1.411163in}{1.411163in}}%
\pgfusepath{clip}%
\pgfsetbuttcap%
\pgfsetroundjoin%
\definecolor{currentfill}{rgb}{0.168627,0.611765,0.411765}%
\pgfsetfillcolor{currentfill}%
\pgfsetfillopacity{0.180000}%
\pgfsetlinewidth{0.000000pt}%
\definecolor{currentstroke}{rgb}{0.000000,0.000000,0.000000}%
\pgfsetstrokecolor{currentstroke}%
\pgfsetdash{}{0pt}%
\pgfpathmoveto{\pgfqpoint{1.512865in}{1.108723in}}%
\pgfpathlineto{\pgfqpoint{1.527701in}{1.122971in}}%
\pgfpathlineto{\pgfqpoint{1.528190in}{1.145429in}}%
\pgfpathlineto{\pgfqpoint{1.513338in}{1.131226in}}%
\pgfpathlineto{\pgfqpoint{1.512865in}{1.108723in}}%
\pgfpathclose%
\pgfusepath{fill}%
\end{pgfscope}%
\begin{pgfscope}%
\pgfpathrectangle{\pgfqpoint{0.407592in}{0.222904in}}{\pgfqpoint{1.411163in}{1.411163in}}%
\pgfusepath{clip}%
\pgfsetbuttcap%
\pgfsetroundjoin%
\definecolor{currentfill}{rgb}{0.168627,0.611765,0.411765}%
\pgfsetfillcolor{currentfill}%
\pgfsetfillopacity{0.180000}%
\pgfsetlinewidth{0.000000pt}%
\definecolor{currentstroke}{rgb}{0.000000,0.000000,0.000000}%
\pgfsetstrokecolor{currentstroke}%
\pgfsetdash{}{0pt}%
\pgfpathmoveto{\pgfqpoint{1.465117in}{0.918040in}}%
\pgfpathlineto{\pgfqpoint{1.479999in}{0.928075in}}%
\pgfpathlineto{\pgfqpoint{1.480429in}{0.950293in}}%
\pgfpathlineto{\pgfqpoint{1.465530in}{0.940313in}}%
\pgfpathlineto{\pgfqpoint{1.465117in}{0.918040in}}%
\pgfpathclose%
\pgfusepath{fill}%
\end{pgfscope}%
\begin{pgfscope}%
\pgfpathrectangle{\pgfqpoint{0.407592in}{0.222904in}}{\pgfqpoint{1.411163in}{1.411163in}}%
\pgfusepath{clip}%
\pgfsetbuttcap%
\pgfsetroundjoin%
\definecolor{currentfill}{rgb}{0.168627,0.611765,0.411765}%
\pgfsetfillcolor{currentfill}%
\pgfsetfillopacity{0.180000}%
\pgfsetlinewidth{0.000000pt}%
\definecolor{currentstroke}{rgb}{0.000000,0.000000,0.000000}%
\pgfsetstrokecolor{currentstroke}%
\pgfsetdash{}{0pt}%
\pgfpathmoveto{\pgfqpoint{1.497070in}{1.051031in}}%
\pgfpathlineto{\pgfqpoint{1.511924in}{1.063883in}}%
\pgfpathlineto{\pgfqpoint{1.512394in}{1.086275in}}%
\pgfpathlineto{\pgfqpoint{1.497523in}{1.073473in}}%
\pgfpathlineto{\pgfqpoint{1.497070in}{1.051031in}}%
\pgfpathclose%
\pgfusepath{fill}%
\end{pgfscope}%
\begin{pgfscope}%
\pgfpathrectangle{\pgfqpoint{0.407592in}{0.222904in}}{\pgfqpoint{1.411163in}{1.411163in}}%
\pgfusepath{clip}%
\pgfsetbuttcap%
\pgfsetroundjoin%
\definecolor{currentfill}{rgb}{0.168627,0.611765,0.411765}%
\pgfsetfillcolor{currentfill}%
\pgfsetfillopacity{0.180000}%
\pgfsetlinewidth{0.000000pt}%
\definecolor{currentstroke}{rgb}{0.000000,0.000000,0.000000}%
\pgfsetstrokecolor{currentstroke}%
\pgfsetdash{}{0pt}%
\pgfpathmoveto{\pgfqpoint{1.481293in}{0.994893in}}%
\pgfpathlineto{\pgfqpoint{1.496168in}{1.006315in}}%
\pgfpathlineto{\pgfqpoint{1.496619in}{1.028646in}}%
\pgfpathlineto{\pgfqpoint{1.481727in}{1.017276in}}%
\pgfpathlineto{\pgfqpoint{1.481293in}{0.994893in}}%
\pgfpathclose%
\pgfusepath{fill}%
\end{pgfscope}%
\begin{pgfscope}%
\pgfpathrectangle{\pgfqpoint{0.407592in}{0.222904in}}{\pgfqpoint{1.411163in}{1.411163in}}%
\pgfusepath{clip}%
\pgfsetbuttcap%
\pgfsetroundjoin%
\definecolor{currentfill}{rgb}{0.168627,0.611765,0.411765}%
\pgfsetfillcolor{currentfill}%
\pgfsetfillopacity{0.180000}%
\pgfsetlinewidth{0.000000pt}%
\definecolor{currentstroke}{rgb}{0.000000,0.000000,0.000000}%
\pgfsetstrokecolor{currentstroke}%
\pgfsetdash{}{0pt}%
\pgfpathmoveto{\pgfqpoint{1.529173in}{1.190512in}}%
\pgfpathlineto{\pgfqpoint{1.544011in}{1.206067in}}%
\pgfpathlineto{\pgfqpoint{1.544521in}{1.228650in}}%
\pgfpathlineto{\pgfqpoint{1.529666in}{1.213137in}}%
\pgfpathlineto{\pgfqpoint{1.529173in}{1.190512in}}%
\pgfpathclose%
\pgfusepath{fill}%
\end{pgfscope}%
\begin{pgfscope}%
\pgfpathrectangle{\pgfqpoint{0.407592in}{0.222904in}}{\pgfqpoint{1.411163in}{1.411163in}}%
\pgfusepath{clip}%
\pgfsetbuttcap%
\pgfsetroundjoin%
\definecolor{currentfill}{rgb}{0.168627,0.611765,0.411765}%
\pgfsetfillcolor{currentfill}%
\pgfsetfillopacity{0.180000}%
\pgfsetlinewidth{0.000000pt}%
\definecolor{currentstroke}{rgb}{0.000000,0.000000,0.000000}%
\pgfsetstrokecolor{currentstroke}%
\pgfsetdash{}{0pt}%
\pgfpathmoveto{\pgfqpoint{1.513338in}{1.131226in}}%
\pgfpathlineto{\pgfqpoint{1.528190in}{1.145429in}}%
\pgfpathlineto{\pgfqpoint{1.528681in}{1.167943in}}%
\pgfpathlineto{\pgfqpoint{1.513811in}{1.153785in}}%
\pgfpathlineto{\pgfqpoint{1.513338in}{1.131226in}}%
\pgfpathclose%
\pgfusepath{fill}%
\end{pgfscope}%
\begin{pgfscope}%
\pgfpathrectangle{\pgfqpoint{0.407592in}{0.222904in}}{\pgfqpoint{1.411163in}{1.411163in}}%
\pgfusepath{clip}%
\pgfsetbuttcap%
\pgfsetroundjoin%
\definecolor{currentfill}{rgb}{0.168627,0.611765,0.411765}%
\pgfsetfillcolor{currentfill}%
\pgfsetfillopacity{0.180000}%
\pgfsetlinewidth{0.000000pt}%
\definecolor{currentstroke}{rgb}{0.000000,0.000000,0.000000}%
\pgfsetstrokecolor{currentstroke}%
\pgfsetdash{}{0pt}%
\pgfpathmoveto{\pgfqpoint{1.465530in}{0.940313in}}%
\pgfpathlineto{\pgfqpoint{1.480429in}{0.950293in}}%
\pgfpathlineto{\pgfqpoint{1.480861in}{0.972565in}}%
\pgfpathlineto{\pgfqpoint{1.465944in}{0.962642in}}%
\pgfpathlineto{\pgfqpoint{1.465530in}{0.940313in}}%
\pgfpathclose%
\pgfusepath{fill}%
\end{pgfscope}%
\begin{pgfscope}%
\pgfpathrectangle{\pgfqpoint{0.407592in}{0.222904in}}{\pgfqpoint{1.411163in}{1.411163in}}%
\pgfusepath{clip}%
\pgfsetbuttcap%
\pgfsetroundjoin%
\definecolor{currentfill}{rgb}{0.168627,0.611765,0.411765}%
\pgfsetfillcolor{currentfill}%
\pgfsetfillopacity{0.180000}%
\pgfsetlinewidth{0.000000pt}%
\definecolor{currentstroke}{rgb}{0.000000,0.000000,0.000000}%
\pgfsetstrokecolor{currentstroke}%
\pgfsetdash{}{0pt}%
\pgfpathmoveto{\pgfqpoint{1.497523in}{1.073473in}}%
\pgfpathlineto{\pgfqpoint{1.512394in}{1.086275in}}%
\pgfpathlineto{\pgfqpoint{1.512865in}{1.108723in}}%
\pgfpathlineto{\pgfqpoint{1.497978in}{1.095970in}}%
\pgfpathlineto{\pgfqpoint{1.497523in}{1.073473in}}%
\pgfpathclose%
\pgfusepath{fill}%
\end{pgfscope}%
\begin{pgfscope}%
\pgfpathrectangle{\pgfqpoint{0.407592in}{0.222904in}}{\pgfqpoint{1.411163in}{1.411163in}}%
\pgfusepath{clip}%
\pgfsetbuttcap%
\pgfsetroundjoin%
\definecolor{currentfill}{rgb}{0.168627,0.611765,0.411765}%
\pgfsetfillcolor{currentfill}%
\pgfsetfillopacity{0.180000}%
\pgfsetlinewidth{0.000000pt}%
\definecolor{currentstroke}{rgb}{0.000000,0.000000,0.000000}%
\pgfsetstrokecolor{currentstroke}%
\pgfsetdash{}{0pt}%
\pgfpathmoveto{\pgfqpoint{1.449776in}{0.887281in}}%
\pgfpathlineto{\pgfqpoint{1.464705in}{0.895822in}}%
\pgfpathlineto{\pgfqpoint{1.465117in}{0.918040in}}%
\pgfpathlineto{\pgfqpoint{1.450171in}{0.909559in}}%
\pgfpathlineto{\pgfqpoint{1.449776in}{0.887281in}}%
\pgfpathclose%
\pgfusepath{fill}%
\end{pgfscope}%
\begin{pgfscope}%
\pgfpathrectangle{\pgfqpoint{0.407592in}{0.222904in}}{\pgfqpoint{1.411163in}{1.411163in}}%
\pgfusepath{clip}%
\pgfsetbuttcap%
\pgfsetroundjoin%
\definecolor{currentfill}{rgb}{0.168627,0.611765,0.411765}%
\pgfsetfillcolor{currentfill}%
\pgfsetfillopacity{0.180000}%
\pgfsetlinewidth{0.000000pt}%
\definecolor{currentstroke}{rgb}{0.000000,0.000000,0.000000}%
\pgfsetstrokecolor{currentstroke}%
\pgfsetdash{}{0pt}%
\pgfpathmoveto{\pgfqpoint{1.481727in}{1.017276in}}%
\pgfpathlineto{\pgfqpoint{1.496619in}{1.028646in}}%
\pgfpathlineto{\pgfqpoint{1.497070in}{1.051031in}}%
\pgfpathlineto{\pgfqpoint{1.482162in}{1.039714in}}%
\pgfpathlineto{\pgfqpoint{1.481727in}{1.017276in}}%
\pgfpathclose%
\pgfusepath{fill}%
\end{pgfscope}%
\begin{pgfscope}%
\pgfpathrectangle{\pgfqpoint{0.407592in}{0.222904in}}{\pgfqpoint{1.411163in}{1.411163in}}%
\pgfusepath{clip}%
\pgfsetbuttcap%
\pgfsetroundjoin%
\definecolor{currentfill}{rgb}{0.168627,0.611765,0.411765}%
\pgfsetfillcolor{currentfill}%
\pgfsetfillopacity{0.180000}%
\pgfsetlinewidth{0.000000pt}%
\definecolor{currentstroke}{rgb}{0.000000,0.000000,0.000000}%
\pgfsetstrokecolor{currentstroke}%
\pgfsetdash{}{0pt}%
\pgfpathmoveto{\pgfqpoint{1.529666in}{1.213137in}}%
\pgfpathlineto{\pgfqpoint{1.544521in}{1.228650in}}%
\pgfpathlineto{\pgfqpoint{1.545032in}{1.251289in}}%
\pgfpathlineto{\pgfqpoint{1.530160in}{1.235819in}}%
\pgfpathlineto{\pgfqpoint{1.529666in}{1.213137in}}%
\pgfpathclose%
\pgfusepath{fill}%
\end{pgfscope}%
\begin{pgfscope}%
\pgfpathrectangle{\pgfqpoint{0.407592in}{0.222904in}}{\pgfqpoint{1.411163in}{1.411163in}}%
\pgfusepath{clip}%
\pgfsetbuttcap%
\pgfsetroundjoin%
\definecolor{currentfill}{rgb}{0.168627,0.611765,0.411765}%
\pgfsetfillcolor{currentfill}%
\pgfsetfillopacity{0.180000}%
\pgfsetlinewidth{0.000000pt}%
\definecolor{currentstroke}{rgb}{0.000000,0.000000,0.000000}%
\pgfsetstrokecolor{currentstroke}%
\pgfsetdash{}{0pt}%
\pgfpathmoveto{\pgfqpoint{1.513811in}{1.153785in}}%
\pgfpathlineto{\pgfqpoint{1.528681in}{1.167943in}}%
\pgfpathlineto{\pgfqpoint{1.529173in}{1.190512in}}%
\pgfpathlineto{\pgfqpoint{1.514286in}{1.176401in}}%
\pgfpathlineto{\pgfqpoint{1.513811in}{1.153785in}}%
\pgfpathclose%
\pgfusepath{fill}%
\end{pgfscope}%
\begin{pgfscope}%
\pgfpathrectangle{\pgfqpoint{0.407592in}{0.222904in}}{\pgfqpoint{1.411163in}{1.411163in}}%
\pgfusepath{clip}%
\pgfsetbuttcap%
\pgfsetroundjoin%
\definecolor{currentfill}{rgb}{0.168627,0.611765,0.411765}%
\pgfsetfillcolor{currentfill}%
\pgfsetfillopacity{0.180000}%
\pgfsetlinewidth{0.000000pt}%
\definecolor{currentstroke}{rgb}{0.000000,0.000000,0.000000}%
\pgfsetstrokecolor{currentstroke}%
\pgfsetdash{}{0pt}%
\pgfpathmoveto{\pgfqpoint{1.465944in}{0.962642in}}%
\pgfpathlineto{\pgfqpoint{1.480861in}{0.972565in}}%
\pgfpathlineto{\pgfqpoint{1.481293in}{0.994893in}}%
\pgfpathlineto{\pgfqpoint{1.466360in}{0.985026in}}%
\pgfpathlineto{\pgfqpoint{1.465944in}{0.962642in}}%
\pgfpathclose%
\pgfusepath{fill}%
\end{pgfscope}%
\begin{pgfscope}%
\pgfpathrectangle{\pgfqpoint{0.407592in}{0.222904in}}{\pgfqpoint{1.411163in}{1.411163in}}%
\pgfusepath{clip}%
\pgfsetbuttcap%
\pgfsetroundjoin%
\definecolor{currentfill}{rgb}{0.168627,0.611765,0.411765}%
\pgfsetfillcolor{currentfill}%
\pgfsetfillopacity{0.180000}%
\pgfsetlinewidth{0.000000pt}%
\definecolor{currentstroke}{rgb}{0.000000,0.000000,0.000000}%
\pgfsetstrokecolor{currentstroke}%
\pgfsetdash{}{0pt}%
\pgfpathmoveto{\pgfqpoint{1.497978in}{1.095970in}}%
\pgfpathlineto{\pgfqpoint{1.512865in}{1.108723in}}%
\pgfpathlineto{\pgfqpoint{1.513338in}{1.131226in}}%
\pgfpathlineto{\pgfqpoint{1.498433in}{1.118523in}}%
\pgfpathlineto{\pgfqpoint{1.497978in}{1.095970in}}%
\pgfpathclose%
\pgfusepath{fill}%
\end{pgfscope}%
\begin{pgfscope}%
\pgfpathrectangle{\pgfqpoint{0.407592in}{0.222904in}}{\pgfqpoint{1.411163in}{1.411163in}}%
\pgfusepath{clip}%
\pgfsetbuttcap%
\pgfsetroundjoin%
\definecolor{currentfill}{rgb}{0.168627,0.611765,0.411765}%
\pgfsetfillcolor{currentfill}%
\pgfsetfillopacity{0.180000}%
\pgfsetlinewidth{0.000000pt}%
\definecolor{currentstroke}{rgb}{0.000000,0.000000,0.000000}%
\pgfsetstrokecolor{currentstroke}%
\pgfsetdash{}{0pt}%
\pgfpathmoveto{\pgfqpoint{1.450171in}{0.909559in}}%
\pgfpathlineto{\pgfqpoint{1.465117in}{0.918040in}}%
\pgfpathlineto{\pgfqpoint{1.465530in}{0.940313in}}%
\pgfpathlineto{\pgfqpoint{1.450567in}{0.931893in}}%
\pgfpathlineto{\pgfqpoint{1.450171in}{0.909559in}}%
\pgfpathclose%
\pgfusepath{fill}%
\end{pgfscope}%
\begin{pgfscope}%
\pgfpathrectangle{\pgfqpoint{0.407592in}{0.222904in}}{\pgfqpoint{1.411163in}{1.411163in}}%
\pgfusepath{clip}%
\pgfsetbuttcap%
\pgfsetroundjoin%
\definecolor{currentfill}{rgb}{0.168627,0.611765,0.411765}%
\pgfsetfillcolor{currentfill}%
\pgfsetfillopacity{0.180000}%
\pgfsetlinewidth{0.000000pt}%
\definecolor{currentstroke}{rgb}{0.000000,0.000000,0.000000}%
\pgfsetstrokecolor{currentstroke}%
\pgfsetdash{}{0pt}%
\pgfpathmoveto{\pgfqpoint{1.482162in}{1.039714in}}%
\pgfpathlineto{\pgfqpoint{1.497070in}{1.051031in}}%
\pgfpathlineto{\pgfqpoint{1.497523in}{1.073473in}}%
\pgfpathlineto{\pgfqpoint{1.482597in}{1.062209in}}%
\pgfpathlineto{\pgfqpoint{1.482162in}{1.039714in}}%
\pgfpathclose%
\pgfusepath{fill}%
\end{pgfscope}%
\begin{pgfscope}%
\pgfpathrectangle{\pgfqpoint{0.407592in}{0.222904in}}{\pgfqpoint{1.411163in}{1.411163in}}%
\pgfusepath{clip}%
\pgfsetbuttcap%
\pgfsetroundjoin%
\definecolor{currentfill}{rgb}{0.168627,0.611765,0.411765}%
\pgfsetfillcolor{currentfill}%
\pgfsetfillopacity{0.180000}%
\pgfsetlinewidth{0.000000pt}%
\definecolor{currentstroke}{rgb}{0.000000,0.000000,0.000000}%
\pgfsetstrokecolor{currentstroke}%
\pgfsetdash{}{0pt}%
\pgfpathmoveto{\pgfqpoint{1.514286in}{1.176401in}}%
\pgfpathlineto{\pgfqpoint{1.529173in}{1.190512in}}%
\pgfpathlineto{\pgfqpoint{1.529666in}{1.213137in}}%
\pgfpathlineto{\pgfqpoint{1.514762in}{1.199073in}}%
\pgfpathlineto{\pgfqpoint{1.514286in}{1.176401in}}%
\pgfpathclose%
\pgfusepath{fill}%
\end{pgfscope}%
\begin{pgfscope}%
\pgfpathrectangle{\pgfqpoint{0.407592in}{0.222904in}}{\pgfqpoint{1.411163in}{1.411163in}}%
\pgfusepath{clip}%
\pgfsetbuttcap%
\pgfsetroundjoin%
\definecolor{currentfill}{rgb}{0.168627,0.611765,0.411765}%
\pgfsetfillcolor{currentfill}%
\pgfsetfillopacity{0.180000}%
\pgfsetlinewidth{0.000000pt}%
\definecolor{currentstroke}{rgb}{0.000000,0.000000,0.000000}%
\pgfsetstrokecolor{currentstroke}%
\pgfsetdash{}{0pt}%
\pgfpathmoveto{\pgfqpoint{1.466360in}{0.985026in}}%
\pgfpathlineto{\pgfqpoint{1.481293in}{0.994893in}}%
\pgfpathlineto{\pgfqpoint{1.481727in}{1.017276in}}%
\pgfpathlineto{\pgfqpoint{1.466776in}{1.007466in}}%
\pgfpathlineto{\pgfqpoint{1.466360in}{0.985026in}}%
\pgfpathclose%
\pgfusepath{fill}%
\end{pgfscope}%
\begin{pgfscope}%
\pgfpathrectangle{\pgfqpoint{0.407592in}{0.222904in}}{\pgfqpoint{1.411163in}{1.411163in}}%
\pgfusepath{clip}%
\pgfsetbuttcap%
\pgfsetroundjoin%
\definecolor{currentfill}{rgb}{0.168627,0.611765,0.411765}%
\pgfsetfillcolor{currentfill}%
\pgfsetfillopacity{0.180000}%
\pgfsetlinewidth{0.000000pt}%
\definecolor{currentstroke}{rgb}{0.000000,0.000000,0.000000}%
\pgfsetstrokecolor{currentstroke}%
\pgfsetdash{}{0pt}%
\pgfpathmoveto{\pgfqpoint{1.498433in}{1.118523in}}%
\pgfpathlineto{\pgfqpoint{1.513338in}{1.131226in}}%
\pgfpathlineto{\pgfqpoint{1.513811in}{1.153785in}}%
\pgfpathlineto{\pgfqpoint{1.498889in}{1.141132in}}%
\pgfpathlineto{\pgfqpoint{1.498433in}{1.118523in}}%
\pgfpathclose%
\pgfusepath{fill}%
\end{pgfscope}%
\begin{pgfscope}%
\pgfpathrectangle{\pgfqpoint{0.407592in}{0.222904in}}{\pgfqpoint{1.411163in}{1.411163in}}%
\pgfusepath{clip}%
\pgfsetbuttcap%
\pgfsetroundjoin%
\definecolor{currentfill}{rgb}{0.168627,0.611765,0.411765}%
\pgfsetfillcolor{currentfill}%
\pgfsetfillopacity{0.180000}%
\pgfsetlinewidth{0.000000pt}%
\definecolor{currentstroke}{rgb}{0.000000,0.000000,0.000000}%
\pgfsetstrokecolor{currentstroke}%
\pgfsetdash{}{0pt}%
\pgfpathmoveto{\pgfqpoint{1.450567in}{0.931893in}}%
\pgfpathlineto{\pgfqpoint{1.465530in}{0.940313in}}%
\pgfpathlineto{\pgfqpoint{1.465944in}{0.962642in}}%
\pgfpathlineto{\pgfqpoint{1.450964in}{0.954282in}}%
\pgfpathlineto{\pgfqpoint{1.450567in}{0.931893in}}%
\pgfpathclose%
\pgfusepath{fill}%
\end{pgfscope}%
\begin{pgfscope}%
\pgfpathrectangle{\pgfqpoint{0.407592in}{0.222904in}}{\pgfqpoint{1.411163in}{1.411163in}}%
\pgfusepath{clip}%
\pgfsetbuttcap%
\pgfsetroundjoin%
\definecolor{currentfill}{rgb}{0.168627,0.611765,0.411765}%
\pgfsetfillcolor{currentfill}%
\pgfsetfillopacity{0.180000}%
\pgfsetlinewidth{0.000000pt}%
\definecolor{currentstroke}{rgb}{0.000000,0.000000,0.000000}%
\pgfsetstrokecolor{currentstroke}%
\pgfsetdash{}{0pt}%
\pgfpathmoveto{\pgfqpoint{1.482597in}{1.062209in}}%
\pgfpathlineto{\pgfqpoint{1.497523in}{1.073473in}}%
\pgfpathlineto{\pgfqpoint{1.497978in}{1.095970in}}%
\pgfpathlineto{\pgfqpoint{1.483034in}{1.084760in}}%
\pgfpathlineto{\pgfqpoint{1.482597in}{1.062209in}}%
\pgfpathclose%
\pgfusepath{fill}%
\end{pgfscope}%
\begin{pgfscope}%
\pgfpathrectangle{\pgfqpoint{0.407592in}{0.222904in}}{\pgfqpoint{1.411163in}{1.411163in}}%
\pgfusepath{clip}%
\pgfsetbuttcap%
\pgfsetroundjoin%
\definecolor{currentfill}{rgb}{0.168627,0.611765,0.411765}%
\pgfsetfillcolor{currentfill}%
\pgfsetfillopacity{0.180000}%
\pgfsetlinewidth{0.000000pt}%
\definecolor{currentstroke}{rgb}{0.000000,0.000000,0.000000}%
\pgfsetstrokecolor{currentstroke}%
\pgfsetdash{}{0pt}%
\pgfpathmoveto{\pgfqpoint{1.514762in}{1.199073in}}%
\pgfpathlineto{\pgfqpoint{1.529666in}{1.213137in}}%
\pgfpathlineto{\pgfqpoint{1.530160in}{1.235819in}}%
\pgfpathlineto{\pgfqpoint{1.515239in}{1.221801in}}%
\pgfpathlineto{\pgfqpoint{1.514762in}{1.199073in}}%
\pgfpathclose%
\pgfusepath{fill}%
\end{pgfscope}%
\begin{pgfscope}%
\pgfpathrectangle{\pgfqpoint{0.407592in}{0.222904in}}{\pgfqpoint{1.411163in}{1.411163in}}%
\pgfusepath{clip}%
\pgfsetbuttcap%
\pgfsetroundjoin%
\definecolor{currentfill}{rgb}{0.168627,0.611765,0.411765}%
\pgfsetfillcolor{currentfill}%
\pgfsetfillopacity{0.180000}%
\pgfsetlinewidth{0.000000pt}%
\definecolor{currentstroke}{rgb}{0.000000,0.000000,0.000000}%
\pgfsetstrokecolor{currentstroke}%
\pgfsetdash{}{0pt}%
\pgfpathmoveto{\pgfqpoint{1.466776in}{1.007466in}}%
\pgfpathlineto{\pgfqpoint{1.481727in}{1.017276in}}%
\pgfpathlineto{\pgfqpoint{1.482162in}{1.039714in}}%
\pgfpathlineto{\pgfqpoint{1.467193in}{1.029962in}}%
\pgfpathlineto{\pgfqpoint{1.466776in}{1.007466in}}%
\pgfpathclose%
\pgfusepath{fill}%
\end{pgfscope}%
\begin{pgfscope}%
\pgfpathrectangle{\pgfqpoint{0.407592in}{0.222904in}}{\pgfqpoint{1.411163in}{1.411163in}}%
\pgfusepath{clip}%
\pgfsetbuttcap%
\pgfsetroundjoin%
\definecolor{currentfill}{rgb}{0.168627,0.611765,0.411765}%
\pgfsetfillcolor{currentfill}%
\pgfsetfillopacity{0.180000}%
\pgfsetlinewidth{0.000000pt}%
\definecolor{currentstroke}{rgb}{0.000000,0.000000,0.000000}%
\pgfsetstrokecolor{currentstroke}%
\pgfsetdash{}{0pt}%
\pgfpathmoveto{\pgfqpoint{1.434780in}{0.880284in}}%
\pgfpathlineto{\pgfqpoint{1.449776in}{0.887281in}}%
\pgfpathlineto{\pgfqpoint{1.450171in}{0.909559in}}%
\pgfpathlineto{\pgfqpoint{1.435158in}{0.902627in}}%
\pgfpathlineto{\pgfqpoint{1.434780in}{0.880284in}}%
\pgfpathclose%
\pgfusepath{fill}%
\end{pgfscope}%
\begin{pgfscope}%
\pgfpathrectangle{\pgfqpoint{0.407592in}{0.222904in}}{\pgfqpoint{1.411163in}{1.411163in}}%
\pgfusepath{clip}%
\pgfsetbuttcap%
\pgfsetroundjoin%
\definecolor{currentfill}{rgb}{0.168627,0.611765,0.411765}%
\pgfsetfillcolor{currentfill}%
\pgfsetfillopacity{0.180000}%
\pgfsetlinewidth{0.000000pt}%
\definecolor{currentstroke}{rgb}{0.000000,0.000000,0.000000}%
\pgfsetstrokecolor{currentstroke}%
\pgfsetdash{}{0pt}%
\pgfpathmoveto{\pgfqpoint{1.498889in}{1.141132in}}%
\pgfpathlineto{\pgfqpoint{1.513811in}{1.153785in}}%
\pgfpathlineto{\pgfqpoint{1.514286in}{1.176401in}}%
\pgfpathlineto{\pgfqpoint{1.499347in}{1.163798in}}%
\pgfpathlineto{\pgfqpoint{1.498889in}{1.141132in}}%
\pgfpathclose%
\pgfusepath{fill}%
\end{pgfscope}%
\begin{pgfscope}%
\pgfpathrectangle{\pgfqpoint{0.407592in}{0.222904in}}{\pgfqpoint{1.411163in}{1.411163in}}%
\pgfusepath{clip}%
\pgfsetbuttcap%
\pgfsetroundjoin%
\definecolor{currentfill}{rgb}{0.168627,0.611765,0.411765}%
\pgfsetfillcolor{currentfill}%
\pgfsetfillopacity{0.180000}%
\pgfsetlinewidth{0.000000pt}%
\definecolor{currentstroke}{rgb}{0.000000,0.000000,0.000000}%
\pgfsetstrokecolor{currentstroke}%
\pgfsetdash{}{0pt}%
\pgfpathmoveto{\pgfqpoint{1.450964in}{0.954282in}}%
\pgfpathlineto{\pgfqpoint{1.465944in}{0.962642in}}%
\pgfpathlineto{\pgfqpoint{1.466360in}{0.985026in}}%
\pgfpathlineto{\pgfqpoint{1.451362in}{0.976727in}}%
\pgfpathlineto{\pgfqpoint{1.450964in}{0.954282in}}%
\pgfpathclose%
\pgfusepath{fill}%
\end{pgfscope}%
\begin{pgfscope}%
\pgfpathrectangle{\pgfqpoint{0.407592in}{0.222904in}}{\pgfqpoint{1.411163in}{1.411163in}}%
\pgfusepath{clip}%
\pgfsetbuttcap%
\pgfsetroundjoin%
\definecolor{currentfill}{rgb}{0.168627,0.611765,0.411765}%
\pgfsetfillcolor{currentfill}%
\pgfsetfillopacity{0.180000}%
\pgfsetlinewidth{0.000000pt}%
\definecolor{currentstroke}{rgb}{0.000000,0.000000,0.000000}%
\pgfsetstrokecolor{currentstroke}%
\pgfsetdash{}{0pt}%
\pgfpathmoveto{\pgfqpoint{1.483034in}{1.084760in}}%
\pgfpathlineto{\pgfqpoint{1.497978in}{1.095970in}}%
\pgfpathlineto{\pgfqpoint{1.498433in}{1.118523in}}%
\pgfpathlineto{\pgfqpoint{1.483472in}{1.107367in}}%
\pgfpathlineto{\pgfqpoint{1.483034in}{1.084760in}}%
\pgfpathclose%
\pgfusepath{fill}%
\end{pgfscope}%
\begin{pgfscope}%
\pgfpathrectangle{\pgfqpoint{0.407592in}{0.222904in}}{\pgfqpoint{1.411163in}{1.411163in}}%
\pgfusepath{clip}%
\pgfsetbuttcap%
\pgfsetroundjoin%
\definecolor{currentfill}{rgb}{0.168627,0.611765,0.411765}%
\pgfsetfillcolor{currentfill}%
\pgfsetfillopacity{0.180000}%
\pgfsetlinewidth{0.000000pt}%
\definecolor{currentstroke}{rgb}{0.000000,0.000000,0.000000}%
\pgfsetstrokecolor{currentstroke}%
\pgfsetdash{}{0pt}%
\pgfpathmoveto{\pgfqpoint{1.467193in}{1.029962in}}%
\pgfpathlineto{\pgfqpoint{1.482162in}{1.039714in}}%
\pgfpathlineto{\pgfqpoint{1.482597in}{1.062209in}}%
\pgfpathlineto{\pgfqpoint{1.467612in}{1.052514in}}%
\pgfpathlineto{\pgfqpoint{1.467193in}{1.029962in}}%
\pgfpathclose%
\pgfusepath{fill}%
\end{pgfscope}%
\begin{pgfscope}%
\pgfpathrectangle{\pgfqpoint{0.407592in}{0.222904in}}{\pgfqpoint{1.411163in}{1.411163in}}%
\pgfusepath{clip}%
\pgfsetbuttcap%
\pgfsetroundjoin%
\definecolor{currentfill}{rgb}{0.168627,0.611765,0.411765}%
\pgfsetfillcolor{currentfill}%
\pgfsetfillopacity{0.180000}%
\pgfsetlinewidth{0.000000pt}%
\definecolor{currentstroke}{rgb}{0.000000,0.000000,0.000000}%
\pgfsetstrokecolor{currentstroke}%
\pgfsetdash{}{0pt}%
\pgfpathmoveto{\pgfqpoint{1.435158in}{0.902627in}}%
\pgfpathlineto{\pgfqpoint{1.450171in}{0.909559in}}%
\pgfpathlineto{\pgfqpoint{1.450567in}{0.931893in}}%
\pgfpathlineto{\pgfqpoint{1.435536in}{0.925025in}}%
\pgfpathlineto{\pgfqpoint{1.435158in}{0.902627in}}%
\pgfpathclose%
\pgfusepath{fill}%
\end{pgfscope}%
\begin{pgfscope}%
\pgfpathrectangle{\pgfqpoint{0.407592in}{0.222904in}}{\pgfqpoint{1.411163in}{1.411163in}}%
\pgfusepath{clip}%
\pgfsetbuttcap%
\pgfsetroundjoin%
\definecolor{currentfill}{rgb}{0.168627,0.611765,0.411765}%
\pgfsetfillcolor{currentfill}%
\pgfsetfillopacity{0.180000}%
\pgfsetlinewidth{0.000000pt}%
\definecolor{currentstroke}{rgb}{0.000000,0.000000,0.000000}%
\pgfsetstrokecolor{currentstroke}%
\pgfsetdash{}{0pt}%
\pgfpathmoveto{\pgfqpoint{1.499347in}{1.163798in}}%
\pgfpathlineto{\pgfqpoint{1.514286in}{1.176401in}}%
\pgfpathlineto{\pgfqpoint{1.514762in}{1.199073in}}%
\pgfpathlineto{\pgfqpoint{1.499805in}{1.186520in}}%
\pgfpathlineto{\pgfqpoint{1.499347in}{1.163798in}}%
\pgfpathclose%
\pgfusepath{fill}%
\end{pgfscope}%
\begin{pgfscope}%
\pgfpathrectangle{\pgfqpoint{0.407592in}{0.222904in}}{\pgfqpoint{1.411163in}{1.411163in}}%
\pgfusepath{clip}%
\pgfsetbuttcap%
\pgfsetroundjoin%
\definecolor{currentfill}{rgb}{0.168627,0.611765,0.411765}%
\pgfsetfillcolor{currentfill}%
\pgfsetfillopacity{0.180000}%
\pgfsetlinewidth{0.000000pt}%
\definecolor{currentstroke}{rgb}{0.000000,0.000000,0.000000}%
\pgfsetstrokecolor{currentstroke}%
\pgfsetdash{}{0pt}%
\pgfpathmoveto{\pgfqpoint{1.451362in}{0.976727in}}%
\pgfpathlineto{\pgfqpoint{1.466360in}{0.985026in}}%
\pgfpathlineto{\pgfqpoint{1.466776in}{1.007466in}}%
\pgfpathlineto{\pgfqpoint{1.451761in}{0.999229in}}%
\pgfpathlineto{\pgfqpoint{1.451362in}{0.976727in}}%
\pgfpathclose%
\pgfusepath{fill}%
\end{pgfscope}%
\begin{pgfscope}%
\pgfpathrectangle{\pgfqpoint{0.407592in}{0.222904in}}{\pgfqpoint{1.411163in}{1.411163in}}%
\pgfusepath{clip}%
\pgfsetbuttcap%
\pgfsetroundjoin%
\definecolor{currentfill}{rgb}{0.168627,0.611765,0.411765}%
\pgfsetfillcolor{currentfill}%
\pgfsetfillopacity{0.180000}%
\pgfsetlinewidth{0.000000pt}%
\definecolor{currentstroke}{rgb}{0.000000,0.000000,0.000000}%
\pgfsetstrokecolor{currentstroke}%
\pgfsetdash{}{0pt}%
\pgfpathmoveto{\pgfqpoint{1.483472in}{1.107367in}}%
\pgfpathlineto{\pgfqpoint{1.498433in}{1.118523in}}%
\pgfpathlineto{\pgfqpoint{1.498889in}{1.141132in}}%
\pgfpathlineto{\pgfqpoint{1.483911in}{1.130030in}}%
\pgfpathlineto{\pgfqpoint{1.483472in}{1.107367in}}%
\pgfpathclose%
\pgfusepath{fill}%
\end{pgfscope}%
\begin{pgfscope}%
\pgfpathrectangle{\pgfqpoint{0.407592in}{0.222904in}}{\pgfqpoint{1.411163in}{1.411163in}}%
\pgfusepath{clip}%
\pgfsetbuttcap%
\pgfsetroundjoin%
\definecolor{currentfill}{rgb}{0.168627,0.611765,0.411765}%
\pgfsetfillcolor{currentfill}%
\pgfsetfillopacity{0.180000}%
\pgfsetlinewidth{0.000000pt}%
\definecolor{currentstroke}{rgb}{0.000000,0.000000,0.000000}%
\pgfsetstrokecolor{currentstroke}%
\pgfsetdash{}{0pt}%
\pgfpathmoveto{\pgfqpoint{1.467612in}{1.052514in}}%
\pgfpathlineto{\pgfqpoint{1.482597in}{1.062209in}}%
\pgfpathlineto{\pgfqpoint{1.483034in}{1.084760in}}%
\pgfpathlineto{\pgfqpoint{1.468031in}{1.075123in}}%
\pgfpathlineto{\pgfqpoint{1.467612in}{1.052514in}}%
\pgfpathclose%
\pgfusepath{fill}%
\end{pgfscope}%
\begin{pgfscope}%
\pgfpathrectangle{\pgfqpoint{0.407592in}{0.222904in}}{\pgfqpoint{1.411163in}{1.411163in}}%
\pgfusepath{clip}%
\pgfsetbuttcap%
\pgfsetroundjoin%
\definecolor{currentfill}{rgb}{0.168627,0.611765,0.411765}%
\pgfsetfillcolor{currentfill}%
\pgfsetfillopacity{0.180000}%
\pgfsetlinewidth{0.000000pt}%
\definecolor{currentstroke}{rgb}{0.000000,0.000000,0.000000}%
\pgfsetstrokecolor{currentstroke}%
\pgfsetdash{}{0pt}%
\pgfpathmoveto{\pgfqpoint{1.435536in}{0.925025in}}%
\pgfpathlineto{\pgfqpoint{1.450567in}{0.931893in}}%
\pgfpathlineto{\pgfqpoint{1.450964in}{0.954282in}}%
\pgfpathlineto{\pgfqpoint{1.435916in}{0.947480in}}%
\pgfpathlineto{\pgfqpoint{1.435536in}{0.925025in}}%
\pgfpathclose%
\pgfusepath{fill}%
\end{pgfscope}%
\begin{pgfscope}%
\pgfpathrectangle{\pgfqpoint{0.407592in}{0.222904in}}{\pgfqpoint{1.411163in}{1.411163in}}%
\pgfusepath{clip}%
\pgfsetbuttcap%
\pgfsetroundjoin%
\definecolor{currentfill}{rgb}{0.168627,0.611765,0.411765}%
\pgfsetfillcolor{currentfill}%
\pgfsetfillopacity{0.180000}%
\pgfsetlinewidth{0.000000pt}%
\definecolor{currentstroke}{rgb}{0.000000,0.000000,0.000000}%
\pgfsetstrokecolor{currentstroke}%
\pgfsetdash{}{0pt}%
\pgfpathmoveto{\pgfqpoint{1.499805in}{1.186520in}}%
\pgfpathlineto{\pgfqpoint{1.514762in}{1.199073in}}%
\pgfpathlineto{\pgfqpoint{1.515239in}{1.221801in}}%
\pgfpathlineto{\pgfqpoint{1.500265in}{1.209299in}}%
\pgfpathlineto{\pgfqpoint{1.499805in}{1.186520in}}%
\pgfpathclose%
\pgfusepath{fill}%
\end{pgfscope}%
\begin{pgfscope}%
\pgfpathrectangle{\pgfqpoint{0.407592in}{0.222904in}}{\pgfqpoint{1.411163in}{1.411163in}}%
\pgfusepath{clip}%
\pgfsetbuttcap%
\pgfsetroundjoin%
\definecolor{currentfill}{rgb}{0.168627,0.611765,0.411765}%
\pgfsetfillcolor{currentfill}%
\pgfsetfillopacity{0.180000}%
\pgfsetlinewidth{0.000000pt}%
\definecolor{currentstroke}{rgb}{0.000000,0.000000,0.000000}%
\pgfsetstrokecolor{currentstroke}%
\pgfsetdash{}{0pt}%
\pgfpathmoveto{\pgfqpoint{1.451761in}{0.999229in}}%
\pgfpathlineto{\pgfqpoint{1.466776in}{1.007466in}}%
\pgfpathlineto{\pgfqpoint{1.467193in}{1.029962in}}%
\pgfpathlineto{\pgfqpoint{1.452161in}{1.021787in}}%
\pgfpathlineto{\pgfqpoint{1.451761in}{0.999229in}}%
\pgfpathclose%
\pgfusepath{fill}%
\end{pgfscope}%
\begin{pgfscope}%
\pgfpathrectangle{\pgfqpoint{0.407592in}{0.222904in}}{\pgfqpoint{1.411163in}{1.411163in}}%
\pgfusepath{clip}%
\pgfsetbuttcap%
\pgfsetroundjoin%
\definecolor{currentfill}{rgb}{0.168627,0.611765,0.411765}%
\pgfsetfillcolor{currentfill}%
\pgfsetfillopacity{0.180000}%
\pgfsetlinewidth{0.000000pt}%
\definecolor{currentstroke}{rgb}{0.000000,0.000000,0.000000}%
\pgfsetstrokecolor{currentstroke}%
\pgfsetdash{}{0pt}%
\pgfpathmoveto{\pgfqpoint{1.483911in}{1.130030in}}%
\pgfpathlineto{\pgfqpoint{1.498889in}{1.141132in}}%
\pgfpathlineto{\pgfqpoint{1.499347in}{1.163798in}}%
\pgfpathlineto{\pgfqpoint{1.484352in}{1.152750in}}%
\pgfpathlineto{\pgfqpoint{1.483911in}{1.130030in}}%
\pgfpathclose%
\pgfusepath{fill}%
\end{pgfscope}%
\begin{pgfscope}%
\pgfpathrectangle{\pgfqpoint{0.407592in}{0.222904in}}{\pgfqpoint{1.411163in}{1.411163in}}%
\pgfusepath{clip}%
\pgfsetbuttcap%
\pgfsetroundjoin%
\definecolor{currentfill}{rgb}{0.168627,0.611765,0.411765}%
\pgfsetfillcolor{currentfill}%
\pgfsetfillopacity{0.180000}%
\pgfsetlinewidth{0.000000pt}%
\definecolor{currentstroke}{rgb}{0.000000,0.000000,0.000000}%
\pgfsetstrokecolor{currentstroke}%
\pgfsetdash{}{0pt}%
\pgfpathmoveto{\pgfqpoint{1.419710in}{0.874808in}}%
\pgfpathlineto{\pgfqpoint{1.434780in}{0.880284in}}%
\pgfpathlineto{\pgfqpoint{1.435158in}{0.902627in}}%
\pgfpathlineto{\pgfqpoint{1.420070in}{0.897219in}}%
\pgfpathlineto{\pgfqpoint{1.419710in}{0.874808in}}%
\pgfpathclose%
\pgfusepath{fill}%
\end{pgfscope}%
\begin{pgfscope}%
\pgfpathrectangle{\pgfqpoint{0.407592in}{0.222904in}}{\pgfqpoint{1.411163in}{1.411163in}}%
\pgfusepath{clip}%
\pgfsetbuttcap%
\pgfsetroundjoin%
\definecolor{currentfill}{rgb}{0.168627,0.611765,0.411765}%
\pgfsetfillcolor{currentfill}%
\pgfsetfillopacity{0.180000}%
\pgfsetlinewidth{0.000000pt}%
\definecolor{currentstroke}{rgb}{0.000000,0.000000,0.000000}%
\pgfsetstrokecolor{currentstroke}%
\pgfsetdash{}{0pt}%
\pgfpathmoveto{\pgfqpoint{1.468031in}{1.075123in}}%
\pgfpathlineto{\pgfqpoint{1.483034in}{1.084760in}}%
\pgfpathlineto{\pgfqpoint{1.483472in}{1.107367in}}%
\pgfpathlineto{\pgfqpoint{1.468452in}{1.097788in}}%
\pgfpathlineto{\pgfqpoint{1.468031in}{1.075123in}}%
\pgfpathclose%
\pgfusepath{fill}%
\end{pgfscope}%
\begin{pgfscope}%
\pgfpathrectangle{\pgfqpoint{0.407592in}{0.222904in}}{\pgfqpoint{1.411163in}{1.411163in}}%
\pgfusepath{clip}%
\pgfsetbuttcap%
\pgfsetroundjoin%
\definecolor{currentfill}{rgb}{0.168627,0.611765,0.411765}%
\pgfsetfillcolor{currentfill}%
\pgfsetfillopacity{0.180000}%
\pgfsetlinewidth{0.000000pt}%
\definecolor{currentstroke}{rgb}{0.000000,0.000000,0.000000}%
\pgfsetstrokecolor{currentstroke}%
\pgfsetdash{}{0pt}%
\pgfpathmoveto{\pgfqpoint{1.435916in}{0.947480in}}%
\pgfpathlineto{\pgfqpoint{1.450964in}{0.954282in}}%
\pgfpathlineto{\pgfqpoint{1.451362in}{0.976727in}}%
\pgfpathlineto{\pgfqpoint{1.436296in}{0.969990in}}%
\pgfpathlineto{\pgfqpoint{1.435916in}{0.947480in}}%
\pgfpathclose%
\pgfusepath{fill}%
\end{pgfscope}%
\begin{pgfscope}%
\pgfpathrectangle{\pgfqpoint{0.407592in}{0.222904in}}{\pgfqpoint{1.411163in}{1.411163in}}%
\pgfusepath{clip}%
\pgfsetbuttcap%
\pgfsetroundjoin%
\definecolor{currentfill}{rgb}{0.168627,0.611765,0.411765}%
\pgfsetfillcolor{currentfill}%
\pgfsetfillopacity{0.180000}%
\pgfsetlinewidth{0.000000pt}%
\definecolor{currentstroke}{rgb}{0.000000,0.000000,0.000000}%
\pgfsetstrokecolor{currentstroke}%
\pgfsetdash{}{0pt}%
\pgfpathmoveto{\pgfqpoint{1.452161in}{1.021787in}}%
\pgfpathlineto{\pgfqpoint{1.467193in}{1.029962in}}%
\pgfpathlineto{\pgfqpoint{1.467612in}{1.052514in}}%
\pgfpathlineto{\pgfqpoint{1.452562in}{1.044401in}}%
\pgfpathlineto{\pgfqpoint{1.452161in}{1.021787in}}%
\pgfpathclose%
\pgfusepath{fill}%
\end{pgfscope}%
\begin{pgfscope}%
\pgfpathrectangle{\pgfqpoint{0.407592in}{0.222904in}}{\pgfqpoint{1.411163in}{1.411163in}}%
\pgfusepath{clip}%
\pgfsetbuttcap%
\pgfsetroundjoin%
\definecolor{currentfill}{rgb}{0.168627,0.611765,0.411765}%
\pgfsetfillcolor{currentfill}%
\pgfsetfillopacity{0.180000}%
\pgfsetlinewidth{0.000000pt}%
\definecolor{currentstroke}{rgb}{0.000000,0.000000,0.000000}%
\pgfsetstrokecolor{currentstroke}%
\pgfsetdash{}{0pt}%
\pgfpathmoveto{\pgfqpoint{1.484352in}{1.152750in}}%
\pgfpathlineto{\pgfqpoint{1.499347in}{1.163798in}}%
\pgfpathlineto{\pgfqpoint{1.499805in}{1.186520in}}%
\pgfpathlineto{\pgfqpoint{1.484793in}{1.175528in}}%
\pgfpathlineto{\pgfqpoint{1.484352in}{1.152750in}}%
\pgfpathclose%
\pgfusepath{fill}%
\end{pgfscope}%
\begin{pgfscope}%
\pgfpathrectangle{\pgfqpoint{0.407592in}{0.222904in}}{\pgfqpoint{1.411163in}{1.411163in}}%
\pgfusepath{clip}%
\pgfsetbuttcap%
\pgfsetroundjoin%
\definecolor{currentfill}{rgb}{0.168627,0.611765,0.411765}%
\pgfsetfillcolor{currentfill}%
\pgfsetfillopacity{0.180000}%
\pgfsetlinewidth{0.000000pt}%
\definecolor{currentstroke}{rgb}{0.000000,0.000000,0.000000}%
\pgfsetstrokecolor{currentstroke}%
\pgfsetdash{}{0pt}%
\pgfpathmoveto{\pgfqpoint{1.420070in}{0.897219in}}%
\pgfpathlineto{\pgfqpoint{1.435158in}{0.902627in}}%
\pgfpathlineto{\pgfqpoint{1.435536in}{0.925025in}}%
\pgfpathlineto{\pgfqpoint{1.420431in}{0.919686in}}%
\pgfpathlineto{\pgfqpoint{1.420070in}{0.897219in}}%
\pgfpathclose%
\pgfusepath{fill}%
\end{pgfscope}%
\begin{pgfscope}%
\pgfpathrectangle{\pgfqpoint{0.407592in}{0.222904in}}{\pgfqpoint{1.411163in}{1.411163in}}%
\pgfusepath{clip}%
\pgfsetbuttcap%
\pgfsetroundjoin%
\definecolor{currentfill}{rgb}{0.168627,0.611765,0.411765}%
\pgfsetfillcolor{currentfill}%
\pgfsetfillopacity{0.180000}%
\pgfsetlinewidth{0.000000pt}%
\definecolor{currentstroke}{rgb}{0.000000,0.000000,0.000000}%
\pgfsetstrokecolor{currentstroke}%
\pgfsetdash{}{0pt}%
\pgfpathmoveto{\pgfqpoint{1.468452in}{1.097788in}}%
\pgfpathlineto{\pgfqpoint{1.483472in}{1.107367in}}%
\pgfpathlineto{\pgfqpoint{1.483911in}{1.130030in}}%
\pgfpathlineto{\pgfqpoint{1.468874in}{1.120510in}}%
\pgfpathlineto{\pgfqpoint{1.468452in}{1.097788in}}%
\pgfpathclose%
\pgfusepath{fill}%
\end{pgfscope}%
\begin{pgfscope}%
\pgfpathrectangle{\pgfqpoint{0.407592in}{0.222904in}}{\pgfqpoint{1.411163in}{1.411163in}}%
\pgfusepath{clip}%
\pgfsetbuttcap%
\pgfsetroundjoin%
\definecolor{currentfill}{rgb}{0.168627,0.611765,0.411765}%
\pgfsetfillcolor{currentfill}%
\pgfsetfillopacity{0.180000}%
\pgfsetlinewidth{0.000000pt}%
\definecolor{currentstroke}{rgb}{0.000000,0.000000,0.000000}%
\pgfsetstrokecolor{currentstroke}%
\pgfsetdash{}{0pt}%
\pgfpathmoveto{\pgfqpoint{1.436296in}{0.969990in}}%
\pgfpathlineto{\pgfqpoint{1.451362in}{0.976727in}}%
\pgfpathlineto{\pgfqpoint{1.451761in}{0.999229in}}%
\pgfpathlineto{\pgfqpoint{1.436677in}{0.992557in}}%
\pgfpathlineto{\pgfqpoint{1.436296in}{0.969990in}}%
\pgfpathclose%
\pgfusepath{fill}%
\end{pgfscope}%
\begin{pgfscope}%
\pgfpathrectangle{\pgfqpoint{0.407592in}{0.222904in}}{\pgfqpoint{1.411163in}{1.411163in}}%
\pgfusepath{clip}%
\pgfsetbuttcap%
\pgfsetroundjoin%
\definecolor{currentfill}{rgb}{0.168627,0.611765,0.411765}%
\pgfsetfillcolor{currentfill}%
\pgfsetfillopacity{0.180000}%
\pgfsetlinewidth{0.000000pt}%
\definecolor{currentstroke}{rgb}{0.000000,0.000000,0.000000}%
\pgfsetstrokecolor{currentstroke}%
\pgfsetdash{}{0pt}%
\pgfpathmoveto{\pgfqpoint{1.452562in}{1.044401in}}%
\pgfpathlineto{\pgfqpoint{1.467612in}{1.052514in}}%
\pgfpathlineto{\pgfqpoint{1.468031in}{1.075123in}}%
\pgfpathlineto{\pgfqpoint{1.452964in}{1.067072in}}%
\pgfpathlineto{\pgfqpoint{1.452562in}{1.044401in}}%
\pgfpathclose%
\pgfusepath{fill}%
\end{pgfscope}%
\begin{pgfscope}%
\pgfpathrectangle{\pgfqpoint{0.407592in}{0.222904in}}{\pgfqpoint{1.411163in}{1.411163in}}%
\pgfusepath{clip}%
\pgfsetbuttcap%
\pgfsetroundjoin%
\definecolor{currentfill}{rgb}{0.168627,0.611765,0.411765}%
\pgfsetfillcolor{currentfill}%
\pgfsetfillopacity{0.180000}%
\pgfsetlinewidth{0.000000pt}%
\definecolor{currentstroke}{rgb}{0.000000,0.000000,0.000000}%
\pgfsetstrokecolor{currentstroke}%
\pgfsetdash{}{0pt}%
\pgfpathmoveto{\pgfqpoint{1.484793in}{1.175528in}}%
\pgfpathlineto{\pgfqpoint{1.499805in}{1.186520in}}%
\pgfpathlineto{\pgfqpoint{1.500265in}{1.209299in}}%
\pgfpathlineto{\pgfqpoint{1.485235in}{1.198362in}}%
\pgfpathlineto{\pgfqpoint{1.484793in}{1.175528in}}%
\pgfpathclose%
\pgfusepath{fill}%
\end{pgfscope}%
\begin{pgfscope}%
\pgfpathrectangle{\pgfqpoint{0.407592in}{0.222904in}}{\pgfqpoint{1.411163in}{1.411163in}}%
\pgfusepath{clip}%
\pgfsetbuttcap%
\pgfsetroundjoin%
\definecolor{currentfill}{rgb}{0.168627,0.611765,0.411765}%
\pgfsetfillcolor{currentfill}%
\pgfsetfillopacity{0.180000}%
\pgfsetlinewidth{0.000000pt}%
\definecolor{currentstroke}{rgb}{0.000000,0.000000,0.000000}%
\pgfsetstrokecolor{currentstroke}%
\pgfsetdash{}{0pt}%
\pgfpathmoveto{\pgfqpoint{1.420431in}{0.919686in}}%
\pgfpathlineto{\pgfqpoint{1.435536in}{0.925025in}}%
\pgfpathlineto{\pgfqpoint{1.435916in}{0.947480in}}%
\pgfpathlineto{\pgfqpoint{1.420793in}{0.942210in}}%
\pgfpathlineto{\pgfqpoint{1.420431in}{0.919686in}}%
\pgfpathclose%
\pgfusepath{fill}%
\end{pgfscope}%
\begin{pgfscope}%
\pgfpathrectangle{\pgfqpoint{0.407592in}{0.222904in}}{\pgfqpoint{1.411163in}{1.411163in}}%
\pgfusepath{clip}%
\pgfsetbuttcap%
\pgfsetroundjoin%
\definecolor{currentfill}{rgb}{0.168627,0.611765,0.411765}%
\pgfsetfillcolor{currentfill}%
\pgfsetfillopacity{0.180000}%
\pgfsetlinewidth{0.000000pt}%
\definecolor{currentstroke}{rgb}{0.000000,0.000000,0.000000}%
\pgfsetstrokecolor{currentstroke}%
\pgfsetdash{}{0pt}%
\pgfpathmoveto{\pgfqpoint{1.468874in}{1.120510in}}%
\pgfpathlineto{\pgfqpoint{1.483911in}{1.130030in}}%
\pgfpathlineto{\pgfqpoint{1.484352in}{1.152750in}}%
\pgfpathlineto{\pgfqpoint{1.469296in}{1.143289in}}%
\pgfpathlineto{\pgfqpoint{1.468874in}{1.120510in}}%
\pgfpathclose%
\pgfusepath{fill}%
\end{pgfscope}%
\begin{pgfscope}%
\pgfpathrectangle{\pgfqpoint{0.407592in}{0.222904in}}{\pgfqpoint{1.411163in}{1.411163in}}%
\pgfusepath{clip}%
\pgfsetbuttcap%
\pgfsetroundjoin%
\definecolor{currentfill}{rgb}{0.168627,0.611765,0.411765}%
\pgfsetfillcolor{currentfill}%
\pgfsetfillopacity{0.180000}%
\pgfsetlinewidth{0.000000pt}%
\definecolor{currentstroke}{rgb}{0.000000,0.000000,0.000000}%
\pgfsetstrokecolor{currentstroke}%
\pgfsetdash{}{0pt}%
\pgfpathmoveto{\pgfqpoint{1.436677in}{0.992557in}}%
\pgfpathlineto{\pgfqpoint{1.451761in}{0.999229in}}%
\pgfpathlineto{\pgfqpoint{1.452161in}{1.021787in}}%
\pgfpathlineto{\pgfqpoint{1.437060in}{1.015181in}}%
\pgfpathlineto{\pgfqpoint{1.436677in}{0.992557in}}%
\pgfpathclose%
\pgfusepath{fill}%
\end{pgfscope}%
\begin{pgfscope}%
\pgfpathrectangle{\pgfqpoint{0.407592in}{0.222904in}}{\pgfqpoint{1.411163in}{1.411163in}}%
\pgfusepath{clip}%
\pgfsetbuttcap%
\pgfsetroundjoin%
\definecolor{currentfill}{rgb}{0.168627,0.611765,0.411765}%
\pgfsetfillcolor{currentfill}%
\pgfsetfillopacity{0.180000}%
\pgfsetlinewidth{0.000000pt}%
\definecolor{currentstroke}{rgb}{0.000000,0.000000,0.000000}%
\pgfsetstrokecolor{currentstroke}%
\pgfsetdash{}{0pt}%
\pgfpathmoveto{\pgfqpoint{1.404562in}{0.870811in}}%
\pgfpathlineto{\pgfqpoint{1.419710in}{0.874808in}}%
\pgfpathlineto{\pgfqpoint{1.420070in}{0.897219in}}%
\pgfpathlineto{\pgfqpoint{1.404904in}{0.893295in}}%
\pgfpathlineto{\pgfqpoint{1.404562in}{0.870811in}}%
\pgfpathclose%
\pgfusepath{fill}%
\end{pgfscope}%
\begin{pgfscope}%
\pgfpathrectangle{\pgfqpoint{0.407592in}{0.222904in}}{\pgfqpoint{1.411163in}{1.411163in}}%
\pgfusepath{clip}%
\pgfsetbuttcap%
\pgfsetroundjoin%
\definecolor{currentfill}{rgb}{0.168627,0.611765,0.411765}%
\pgfsetfillcolor{currentfill}%
\pgfsetfillopacity{0.180000}%
\pgfsetlinewidth{0.000000pt}%
\definecolor{currentstroke}{rgb}{0.000000,0.000000,0.000000}%
\pgfsetstrokecolor{currentstroke}%
\pgfsetdash{}{0pt}%
\pgfpathmoveto{\pgfqpoint{1.452964in}{1.067072in}}%
\pgfpathlineto{\pgfqpoint{1.468031in}{1.075123in}}%
\pgfpathlineto{\pgfqpoint{1.468452in}{1.097788in}}%
\pgfpathlineto{\pgfqpoint{1.453367in}{1.089800in}}%
\pgfpathlineto{\pgfqpoint{1.452964in}{1.067072in}}%
\pgfpathclose%
\pgfusepath{fill}%
\end{pgfscope}%
\begin{pgfscope}%
\pgfpathrectangle{\pgfqpoint{0.407592in}{0.222904in}}{\pgfqpoint{1.411163in}{1.411163in}}%
\pgfusepath{clip}%
\pgfsetbuttcap%
\pgfsetroundjoin%
\definecolor{currentfill}{rgb}{0.168627,0.611765,0.411765}%
\pgfsetfillcolor{currentfill}%
\pgfsetfillopacity{0.180000}%
\pgfsetlinewidth{0.000000pt}%
\definecolor{currentstroke}{rgb}{0.000000,0.000000,0.000000}%
\pgfsetstrokecolor{currentstroke}%
\pgfsetdash{}{0pt}%
\pgfpathmoveto{\pgfqpoint{1.420793in}{0.942210in}}%
\pgfpathlineto{\pgfqpoint{1.435916in}{0.947480in}}%
\pgfpathlineto{\pgfqpoint{1.436296in}{0.969990in}}%
\pgfpathlineto{\pgfqpoint{1.421156in}{0.964790in}}%
\pgfpathlineto{\pgfqpoint{1.420793in}{0.942210in}}%
\pgfpathclose%
\pgfusepath{fill}%
\end{pgfscope}%
\begin{pgfscope}%
\pgfpathrectangle{\pgfqpoint{0.407592in}{0.222904in}}{\pgfqpoint{1.411163in}{1.411163in}}%
\pgfusepath{clip}%
\pgfsetbuttcap%
\pgfsetroundjoin%
\definecolor{currentfill}{rgb}{0.168627,0.611765,0.411765}%
\pgfsetfillcolor{currentfill}%
\pgfsetfillopacity{0.180000}%
\pgfsetlinewidth{0.000000pt}%
\definecolor{currentstroke}{rgb}{0.000000,0.000000,0.000000}%
\pgfsetstrokecolor{currentstroke}%
\pgfsetdash{}{0pt}%
\pgfpathmoveto{\pgfqpoint{1.469296in}{1.143289in}}%
\pgfpathlineto{\pgfqpoint{1.484352in}{1.152750in}}%
\pgfpathlineto{\pgfqpoint{1.484793in}{1.175528in}}%
\pgfpathlineto{\pgfqpoint{1.469720in}{1.166126in}}%
\pgfpathlineto{\pgfqpoint{1.469296in}{1.143289in}}%
\pgfpathclose%
\pgfusepath{fill}%
\end{pgfscope}%
\begin{pgfscope}%
\pgfpathrectangle{\pgfqpoint{0.407592in}{0.222904in}}{\pgfqpoint{1.411163in}{1.411163in}}%
\pgfusepath{clip}%
\pgfsetbuttcap%
\pgfsetroundjoin%
\definecolor{currentfill}{rgb}{0.168627,0.611765,0.411765}%
\pgfsetfillcolor{currentfill}%
\pgfsetfillopacity{0.180000}%
\pgfsetlinewidth{0.000000pt}%
\definecolor{currentstroke}{rgb}{0.000000,0.000000,0.000000}%
\pgfsetstrokecolor{currentstroke}%
\pgfsetdash{}{0pt}%
\pgfpathmoveto{\pgfqpoint{1.437060in}{1.015181in}}%
\pgfpathlineto{\pgfqpoint{1.452161in}{1.021787in}}%
\pgfpathlineto{\pgfqpoint{1.452562in}{1.044401in}}%
\pgfpathlineto{\pgfqpoint{1.437443in}{1.037862in}}%
\pgfpathlineto{\pgfqpoint{1.437060in}{1.015181in}}%
\pgfpathclose%
\pgfusepath{fill}%
\end{pgfscope}%
\begin{pgfscope}%
\pgfpathrectangle{\pgfqpoint{0.407592in}{0.222904in}}{\pgfqpoint{1.411163in}{1.411163in}}%
\pgfusepath{clip}%
\pgfsetbuttcap%
\pgfsetroundjoin%
\definecolor{currentfill}{rgb}{0.168627,0.611765,0.411765}%
\pgfsetfillcolor{currentfill}%
\pgfsetfillopacity{0.180000}%
\pgfsetlinewidth{0.000000pt}%
\definecolor{currentstroke}{rgb}{0.000000,0.000000,0.000000}%
\pgfsetstrokecolor{currentstroke}%
\pgfsetdash{}{0pt}%
\pgfpathmoveto{\pgfqpoint{1.404904in}{0.893295in}}%
\pgfpathlineto{\pgfqpoint{1.420070in}{0.897219in}}%
\pgfpathlineto{\pgfqpoint{1.420431in}{0.919686in}}%
\pgfpathlineto{\pgfqpoint{1.405247in}{0.915836in}}%
\pgfpathlineto{\pgfqpoint{1.404904in}{0.893295in}}%
\pgfpathclose%
\pgfusepath{fill}%
\end{pgfscope}%
\begin{pgfscope}%
\pgfpathrectangle{\pgfqpoint{0.407592in}{0.222904in}}{\pgfqpoint{1.411163in}{1.411163in}}%
\pgfusepath{clip}%
\pgfsetbuttcap%
\pgfsetroundjoin%
\definecolor{currentfill}{rgb}{0.168627,0.611765,0.411765}%
\pgfsetfillcolor{currentfill}%
\pgfsetfillopacity{0.180000}%
\pgfsetlinewidth{0.000000pt}%
\definecolor{currentstroke}{rgb}{0.000000,0.000000,0.000000}%
\pgfsetstrokecolor{currentstroke}%
\pgfsetdash{}{0pt}%
\pgfpathmoveto{\pgfqpoint{1.453367in}{1.089800in}}%
\pgfpathlineto{\pgfqpoint{1.468452in}{1.097788in}}%
\pgfpathlineto{\pgfqpoint{1.468874in}{1.120510in}}%
\pgfpathlineto{\pgfqpoint{1.453771in}{1.112585in}}%
\pgfpathlineto{\pgfqpoint{1.453367in}{1.089800in}}%
\pgfpathclose%
\pgfusepath{fill}%
\end{pgfscope}%
\begin{pgfscope}%
\pgfpathrectangle{\pgfqpoint{0.407592in}{0.222904in}}{\pgfqpoint{1.411163in}{1.411163in}}%
\pgfusepath{clip}%
\pgfsetbuttcap%
\pgfsetroundjoin%
\definecolor{currentfill}{rgb}{0.168627,0.611765,0.411765}%
\pgfsetfillcolor{currentfill}%
\pgfsetfillopacity{0.180000}%
\pgfsetlinewidth{0.000000pt}%
\definecolor{currentstroke}{rgb}{0.000000,0.000000,0.000000}%
\pgfsetstrokecolor{currentstroke}%
\pgfsetdash{}{0pt}%
\pgfpathmoveto{\pgfqpoint{1.421156in}{0.964790in}}%
\pgfpathlineto{\pgfqpoint{1.436296in}{0.969990in}}%
\pgfpathlineto{\pgfqpoint{1.436677in}{0.992557in}}%
\pgfpathlineto{\pgfqpoint{1.421519in}{0.987428in}}%
\pgfpathlineto{\pgfqpoint{1.421156in}{0.964790in}}%
\pgfpathclose%
\pgfusepath{fill}%
\end{pgfscope}%
\begin{pgfscope}%
\pgfpathrectangle{\pgfqpoint{0.407592in}{0.222904in}}{\pgfqpoint{1.411163in}{1.411163in}}%
\pgfusepath{clip}%
\pgfsetbuttcap%
\pgfsetroundjoin%
\definecolor{currentfill}{rgb}{0.168627,0.611765,0.411765}%
\pgfsetfillcolor{currentfill}%
\pgfsetfillopacity{0.180000}%
\pgfsetlinewidth{0.000000pt}%
\definecolor{currentstroke}{rgb}{0.000000,0.000000,0.000000}%
\pgfsetstrokecolor{currentstroke}%
\pgfsetdash{}{0pt}%
\pgfpathmoveto{\pgfqpoint{1.469720in}{1.166126in}}%
\pgfpathlineto{\pgfqpoint{1.484793in}{1.175528in}}%
\pgfpathlineto{\pgfqpoint{1.485235in}{1.198362in}}%
\pgfpathlineto{\pgfqpoint{1.470145in}{1.189020in}}%
\pgfpathlineto{\pgfqpoint{1.469720in}{1.166126in}}%
\pgfpathclose%
\pgfusepath{fill}%
\end{pgfscope}%
\begin{pgfscope}%
\pgfpathrectangle{\pgfqpoint{0.407592in}{0.222904in}}{\pgfqpoint{1.411163in}{1.411163in}}%
\pgfusepath{clip}%
\pgfsetbuttcap%
\pgfsetroundjoin%
\definecolor{currentfill}{rgb}{0.168627,0.611765,0.411765}%
\pgfsetfillcolor{currentfill}%
\pgfsetfillopacity{0.180000}%
\pgfsetlinewidth{0.000000pt}%
\definecolor{currentstroke}{rgb}{0.000000,0.000000,0.000000}%
\pgfsetstrokecolor{currentstroke}%
\pgfsetdash{}{0pt}%
\pgfpathmoveto{\pgfqpoint{1.437443in}{1.037862in}}%
\pgfpathlineto{\pgfqpoint{1.452562in}{1.044401in}}%
\pgfpathlineto{\pgfqpoint{1.452964in}{1.067072in}}%
\pgfpathlineto{\pgfqpoint{1.437827in}{1.060600in}}%
\pgfpathlineto{\pgfqpoint{1.437443in}{1.037862in}}%
\pgfpathclose%
\pgfusepath{fill}%
\end{pgfscope}%
\begin{pgfscope}%
\pgfpathrectangle{\pgfqpoint{0.407592in}{0.222904in}}{\pgfqpoint{1.411163in}{1.411163in}}%
\pgfusepath{clip}%
\pgfsetbuttcap%
\pgfsetroundjoin%
\definecolor{currentfill}{rgb}{0.168627,0.611765,0.411765}%
\pgfsetfillcolor{currentfill}%
\pgfsetfillopacity{0.180000}%
\pgfsetlinewidth{0.000000pt}%
\definecolor{currentstroke}{rgb}{0.000000,0.000000,0.000000}%
\pgfsetstrokecolor{currentstroke}%
\pgfsetdash{}{0pt}%
\pgfpathmoveto{\pgfqpoint{1.405247in}{0.915836in}}%
\pgfpathlineto{\pgfqpoint{1.420431in}{0.919686in}}%
\pgfpathlineto{\pgfqpoint{1.420793in}{0.942210in}}%
\pgfpathlineto{\pgfqpoint{1.405591in}{0.938433in}}%
\pgfpathlineto{\pgfqpoint{1.405247in}{0.915836in}}%
\pgfpathclose%
\pgfusepath{fill}%
\end{pgfscope}%
\begin{pgfscope}%
\pgfpathrectangle{\pgfqpoint{0.407592in}{0.222904in}}{\pgfqpoint{1.411163in}{1.411163in}}%
\pgfusepath{clip}%
\pgfsetbuttcap%
\pgfsetroundjoin%
\definecolor{currentfill}{rgb}{0.168627,0.611765,0.411765}%
\pgfsetfillcolor{currentfill}%
\pgfsetfillopacity{0.180000}%
\pgfsetlinewidth{0.000000pt}%
\definecolor{currentstroke}{rgb}{0.000000,0.000000,0.000000}%
\pgfsetstrokecolor{currentstroke}%
\pgfsetdash{}{0pt}%
\pgfpathmoveto{\pgfqpoint{1.453771in}{1.112585in}}%
\pgfpathlineto{\pgfqpoint{1.468874in}{1.120510in}}%
\pgfpathlineto{\pgfqpoint{1.469296in}{1.143289in}}%
\pgfpathlineto{\pgfqpoint{1.454176in}{1.135428in}}%
\pgfpathlineto{\pgfqpoint{1.453771in}{1.112585in}}%
\pgfpathclose%
\pgfusepath{fill}%
\end{pgfscope}%
\begin{pgfscope}%
\pgfpathrectangle{\pgfqpoint{0.407592in}{0.222904in}}{\pgfqpoint{1.411163in}{1.411163in}}%
\pgfusepath{clip}%
\pgfsetbuttcap%
\pgfsetroundjoin%
\definecolor{currentfill}{rgb}{0.168627,0.611765,0.411765}%
\pgfsetfillcolor{currentfill}%
\pgfsetfillopacity{0.180000}%
\pgfsetlinewidth{0.000000pt}%
\definecolor{currentstroke}{rgb}{0.000000,0.000000,0.000000}%
\pgfsetstrokecolor{currentstroke}%
\pgfsetdash{}{0pt}%
\pgfpathmoveto{\pgfqpoint{1.421519in}{0.987428in}}%
\pgfpathlineto{\pgfqpoint{1.436677in}{0.992557in}}%
\pgfpathlineto{\pgfqpoint{1.437060in}{1.015181in}}%
\pgfpathlineto{\pgfqpoint{1.421884in}{1.010122in}}%
\pgfpathlineto{\pgfqpoint{1.421519in}{0.987428in}}%
\pgfpathclose%
\pgfusepath{fill}%
\end{pgfscope}%
\begin{pgfscope}%
\pgfpathrectangle{\pgfqpoint{0.407592in}{0.222904in}}{\pgfqpoint{1.411163in}{1.411163in}}%
\pgfusepath{clip}%
\pgfsetbuttcap%
\pgfsetroundjoin%
\definecolor{currentfill}{rgb}{0.168627,0.611765,0.411765}%
\pgfsetfillcolor{currentfill}%
\pgfsetfillopacity{0.180000}%
\pgfsetlinewidth{0.000000pt}%
\definecolor{currentstroke}{rgb}{0.000000,0.000000,0.000000}%
\pgfsetstrokecolor{currentstroke}%
\pgfsetdash{}{0pt}%
\pgfpathmoveto{\pgfqpoint{1.437827in}{1.060600in}}%
\pgfpathlineto{\pgfqpoint{1.452964in}{1.067072in}}%
\pgfpathlineto{\pgfqpoint{1.453367in}{1.089800in}}%
\pgfpathlineto{\pgfqpoint{1.438213in}{1.083395in}}%
\pgfpathlineto{\pgfqpoint{1.437827in}{1.060600in}}%
\pgfpathclose%
\pgfusepath{fill}%
\end{pgfscope}%
\begin{pgfscope}%
\pgfpathrectangle{\pgfqpoint{0.407592in}{0.222904in}}{\pgfqpoint{1.411163in}{1.411163in}}%
\pgfusepath{clip}%
\pgfsetbuttcap%
\pgfsetroundjoin%
\definecolor{currentfill}{rgb}{0.168627,0.611765,0.411765}%
\pgfsetfillcolor{currentfill}%
\pgfsetfillopacity{0.180000}%
\pgfsetlinewidth{0.000000pt}%
\definecolor{currentstroke}{rgb}{0.000000,0.000000,0.000000}%
\pgfsetstrokecolor{currentstroke}%
\pgfsetdash{}{0pt}%
\pgfpathmoveto{\pgfqpoint{1.389328in}{0.868242in}}%
\pgfpathlineto{\pgfqpoint{1.404562in}{0.870811in}}%
\pgfpathlineto{\pgfqpoint{1.404904in}{0.893295in}}%
\pgfpathlineto{\pgfqpoint{1.389653in}{0.890803in}}%
\pgfpathlineto{\pgfqpoint{1.389328in}{0.868242in}}%
\pgfpathclose%
\pgfusepath{fill}%
\end{pgfscope}%
\begin{pgfscope}%
\pgfpathrectangle{\pgfqpoint{0.407592in}{0.222904in}}{\pgfqpoint{1.411163in}{1.411163in}}%
\pgfusepath{clip}%
\pgfsetbuttcap%
\pgfsetroundjoin%
\definecolor{currentfill}{rgb}{0.168627,0.611765,0.411765}%
\pgfsetfillcolor{currentfill}%
\pgfsetfillopacity{0.180000}%
\pgfsetlinewidth{0.000000pt}%
\definecolor{currentstroke}{rgb}{0.000000,0.000000,0.000000}%
\pgfsetstrokecolor{currentstroke}%
\pgfsetdash{}{0pt}%
\pgfpathmoveto{\pgfqpoint{1.405591in}{0.938433in}}%
\pgfpathlineto{\pgfqpoint{1.420793in}{0.942210in}}%
\pgfpathlineto{\pgfqpoint{1.421156in}{0.964790in}}%
\pgfpathlineto{\pgfqpoint{1.405936in}{0.961088in}}%
\pgfpathlineto{\pgfqpoint{1.405591in}{0.938433in}}%
\pgfpathclose%
\pgfusepath{fill}%
\end{pgfscope}%
\begin{pgfscope}%
\pgfpathrectangle{\pgfqpoint{0.407592in}{0.222904in}}{\pgfqpoint{1.411163in}{1.411163in}}%
\pgfusepath{clip}%
\pgfsetbuttcap%
\pgfsetroundjoin%
\definecolor{currentfill}{rgb}{0.168627,0.611765,0.411765}%
\pgfsetfillcolor{currentfill}%
\pgfsetfillopacity{0.180000}%
\pgfsetlinewidth{0.000000pt}%
\definecolor{currentstroke}{rgb}{0.000000,0.000000,0.000000}%
\pgfsetstrokecolor{currentstroke}%
\pgfsetdash{}{0pt}%
\pgfpathmoveto{\pgfqpoint{1.454176in}{1.135428in}}%
\pgfpathlineto{\pgfqpoint{1.469296in}{1.143289in}}%
\pgfpathlineto{\pgfqpoint{1.469720in}{1.166126in}}%
\pgfpathlineto{\pgfqpoint{1.454582in}{1.158328in}}%
\pgfpathlineto{\pgfqpoint{1.454176in}{1.135428in}}%
\pgfpathclose%
\pgfusepath{fill}%
\end{pgfscope}%
\begin{pgfscope}%
\pgfpathrectangle{\pgfqpoint{0.407592in}{0.222904in}}{\pgfqpoint{1.411163in}{1.411163in}}%
\pgfusepath{clip}%
\pgfsetbuttcap%
\pgfsetroundjoin%
\definecolor{currentfill}{rgb}{0.168627,0.611765,0.411765}%
\pgfsetfillcolor{currentfill}%
\pgfsetfillopacity{0.180000}%
\pgfsetlinewidth{0.000000pt}%
\definecolor{currentstroke}{rgb}{0.000000,0.000000,0.000000}%
\pgfsetstrokecolor{currentstroke}%
\pgfsetdash{}{0pt}%
\pgfpathmoveto{\pgfqpoint{1.421884in}{1.010122in}}%
\pgfpathlineto{\pgfqpoint{1.437060in}{1.015181in}}%
\pgfpathlineto{\pgfqpoint{1.437443in}{1.037862in}}%
\pgfpathlineto{\pgfqpoint{1.422249in}{1.032873in}}%
\pgfpathlineto{\pgfqpoint{1.421884in}{1.010122in}}%
\pgfpathclose%
\pgfusepath{fill}%
\end{pgfscope}%
\begin{pgfscope}%
\pgfpathrectangle{\pgfqpoint{0.407592in}{0.222904in}}{\pgfqpoint{1.411163in}{1.411163in}}%
\pgfusepath{clip}%
\pgfsetbuttcap%
\pgfsetroundjoin%
\definecolor{currentfill}{rgb}{0.168627,0.611765,0.411765}%
\pgfsetfillcolor{currentfill}%
\pgfsetfillopacity{0.180000}%
\pgfsetlinewidth{0.000000pt}%
\definecolor{currentstroke}{rgb}{0.000000,0.000000,0.000000}%
\pgfsetstrokecolor{currentstroke}%
\pgfsetdash{}{0pt}%
\pgfpathmoveto{\pgfqpoint{1.438213in}{1.083395in}}%
\pgfpathlineto{\pgfqpoint{1.453367in}{1.089800in}}%
\pgfpathlineto{\pgfqpoint{1.453771in}{1.112585in}}%
\pgfpathlineto{\pgfqpoint{1.438599in}{1.106248in}}%
\pgfpathlineto{\pgfqpoint{1.438213in}{1.083395in}}%
\pgfpathclose%
\pgfusepath{fill}%
\end{pgfscope}%
\begin{pgfscope}%
\pgfpathrectangle{\pgfqpoint{0.407592in}{0.222904in}}{\pgfqpoint{1.411163in}{1.411163in}}%
\pgfusepath{clip}%
\pgfsetbuttcap%
\pgfsetroundjoin%
\definecolor{currentfill}{rgb}{0.168627,0.611765,0.411765}%
\pgfsetfillcolor{currentfill}%
\pgfsetfillopacity{0.180000}%
\pgfsetlinewidth{0.000000pt}%
\definecolor{currentstroke}{rgb}{0.000000,0.000000,0.000000}%
\pgfsetstrokecolor{currentstroke}%
\pgfsetdash{}{0pt}%
\pgfpathmoveto{\pgfqpoint{1.389653in}{0.890803in}}%
\pgfpathlineto{\pgfqpoint{1.404904in}{0.893295in}}%
\pgfpathlineto{\pgfqpoint{1.405247in}{0.915836in}}%
\pgfpathlineto{\pgfqpoint{1.389978in}{0.913421in}}%
\pgfpathlineto{\pgfqpoint{1.389653in}{0.890803in}}%
\pgfpathclose%
\pgfusepath{fill}%
\end{pgfscope}%
\begin{pgfscope}%
\pgfpathrectangle{\pgfqpoint{0.407592in}{0.222904in}}{\pgfqpoint{1.411163in}{1.411163in}}%
\pgfusepath{clip}%
\pgfsetbuttcap%
\pgfsetroundjoin%
\definecolor{currentfill}{rgb}{0.168627,0.611765,0.411765}%
\pgfsetfillcolor{currentfill}%
\pgfsetfillopacity{0.180000}%
\pgfsetlinewidth{0.000000pt}%
\definecolor{currentstroke}{rgb}{0.000000,0.000000,0.000000}%
\pgfsetstrokecolor{currentstroke}%
\pgfsetdash{}{0pt}%
\pgfpathmoveto{\pgfqpoint{1.405936in}{0.961088in}}%
\pgfpathlineto{\pgfqpoint{1.421156in}{0.964790in}}%
\pgfpathlineto{\pgfqpoint{1.421519in}{0.987428in}}%
\pgfpathlineto{\pgfqpoint{1.406281in}{0.983799in}}%
\pgfpathlineto{\pgfqpoint{1.405936in}{0.961088in}}%
\pgfpathclose%
\pgfusepath{fill}%
\end{pgfscope}%
\begin{pgfscope}%
\pgfpathrectangle{\pgfqpoint{0.407592in}{0.222904in}}{\pgfqpoint{1.411163in}{1.411163in}}%
\pgfusepath{clip}%
\pgfsetbuttcap%
\pgfsetroundjoin%
\definecolor{currentfill}{rgb}{0.168627,0.611765,0.411765}%
\pgfsetfillcolor{currentfill}%
\pgfsetfillopacity{0.180000}%
\pgfsetlinewidth{0.000000pt}%
\definecolor{currentstroke}{rgb}{0.000000,0.000000,0.000000}%
\pgfsetstrokecolor{currentstroke}%
\pgfsetdash{}{0pt}%
\pgfpathmoveto{\pgfqpoint{1.454582in}{1.158328in}}%
\pgfpathlineto{\pgfqpoint{1.469720in}{1.166126in}}%
\pgfpathlineto{\pgfqpoint{1.470145in}{1.189020in}}%
\pgfpathlineto{\pgfqpoint{1.454989in}{1.181286in}}%
\pgfpathlineto{\pgfqpoint{1.454582in}{1.158328in}}%
\pgfpathclose%
\pgfusepath{fill}%
\end{pgfscope}%
\begin{pgfscope}%
\pgfpathrectangle{\pgfqpoint{0.407592in}{0.222904in}}{\pgfqpoint{1.411163in}{1.411163in}}%
\pgfusepath{clip}%
\pgfsetbuttcap%
\pgfsetroundjoin%
\definecolor{currentfill}{rgb}{0.168627,0.611765,0.411765}%
\pgfsetfillcolor{currentfill}%
\pgfsetfillopacity{0.180000}%
\pgfsetlinewidth{0.000000pt}%
\definecolor{currentstroke}{rgb}{0.000000,0.000000,0.000000}%
\pgfsetstrokecolor{currentstroke}%
\pgfsetdash{}{0pt}%
\pgfpathmoveto{\pgfqpoint{1.422249in}{1.032873in}}%
\pgfpathlineto{\pgfqpoint{1.437443in}{1.037862in}}%
\pgfpathlineto{\pgfqpoint{1.437827in}{1.060600in}}%
\pgfpathlineto{\pgfqpoint{1.422616in}{1.055682in}}%
\pgfpathlineto{\pgfqpoint{1.422249in}{1.032873in}}%
\pgfpathclose%
\pgfusepath{fill}%
\end{pgfscope}%
\begin{pgfscope}%
\pgfpathrectangle{\pgfqpoint{0.407592in}{0.222904in}}{\pgfqpoint{1.411163in}{1.411163in}}%
\pgfusepath{clip}%
\pgfsetbuttcap%
\pgfsetroundjoin%
\definecolor{currentfill}{rgb}{0.168627,0.611765,0.411765}%
\pgfsetfillcolor{currentfill}%
\pgfsetfillopacity{0.180000}%
\pgfsetlinewidth{0.000000pt}%
\definecolor{currentstroke}{rgb}{0.000000,0.000000,0.000000}%
\pgfsetstrokecolor{currentstroke}%
\pgfsetdash{}{0pt}%
\pgfpathmoveto{\pgfqpoint{1.438599in}{1.106248in}}%
\pgfpathlineto{\pgfqpoint{1.453771in}{1.112585in}}%
\pgfpathlineto{\pgfqpoint{1.454176in}{1.135428in}}%
\pgfpathlineto{\pgfqpoint{1.438986in}{1.129158in}}%
\pgfpathlineto{\pgfqpoint{1.438599in}{1.106248in}}%
\pgfpathclose%
\pgfusepath{fill}%
\end{pgfscope}%
\begin{pgfscope}%
\pgfpathrectangle{\pgfqpoint{0.407592in}{0.222904in}}{\pgfqpoint{1.411163in}{1.411163in}}%
\pgfusepath{clip}%
\pgfsetbuttcap%
\pgfsetroundjoin%
\definecolor{currentfill}{rgb}{0.168627,0.611765,0.411765}%
\pgfsetfillcolor{currentfill}%
\pgfsetfillopacity{0.180000}%
\pgfsetlinewidth{0.000000pt}%
\definecolor{currentstroke}{rgb}{0.000000,0.000000,0.000000}%
\pgfsetstrokecolor{currentstroke}%
\pgfsetdash{}{0pt}%
\pgfpathmoveto{\pgfqpoint{1.389978in}{0.913421in}}%
\pgfpathlineto{\pgfqpoint{1.405247in}{0.915836in}}%
\pgfpathlineto{\pgfqpoint{1.405591in}{0.938433in}}%
\pgfpathlineto{\pgfqpoint{1.390303in}{0.936096in}}%
\pgfpathlineto{\pgfqpoint{1.389978in}{0.913421in}}%
\pgfpathclose%
\pgfusepath{fill}%
\end{pgfscope}%
\begin{pgfscope}%
\pgfpathrectangle{\pgfqpoint{0.407592in}{0.222904in}}{\pgfqpoint{1.411163in}{1.411163in}}%
\pgfusepath{clip}%
\pgfsetbuttcap%
\pgfsetroundjoin%
\definecolor{currentfill}{rgb}{0.168627,0.611765,0.411765}%
\pgfsetfillcolor{currentfill}%
\pgfsetfillopacity{0.180000}%
\pgfsetlinewidth{0.000000pt}%
\definecolor{currentstroke}{rgb}{0.000000,0.000000,0.000000}%
\pgfsetstrokecolor{currentstroke}%
\pgfsetdash{}{0pt}%
\pgfpathmoveto{\pgfqpoint{1.406281in}{0.983799in}}%
\pgfpathlineto{\pgfqpoint{1.421519in}{0.987428in}}%
\pgfpathlineto{\pgfqpoint{1.421884in}{1.010122in}}%
\pgfpathlineto{\pgfqpoint{1.406628in}{1.006568in}}%
\pgfpathlineto{\pgfqpoint{1.406281in}{0.983799in}}%
\pgfpathclose%
\pgfusepath{fill}%
\end{pgfscope}%
\begin{pgfscope}%
\pgfpathrectangle{\pgfqpoint{0.407592in}{0.222904in}}{\pgfqpoint{1.411163in}{1.411163in}}%
\pgfusepath{clip}%
\pgfsetbuttcap%
\pgfsetroundjoin%
\definecolor{currentfill}{rgb}{0.168627,0.611765,0.411765}%
\pgfsetfillcolor{currentfill}%
\pgfsetfillopacity{0.180000}%
\pgfsetlinewidth{0.000000pt}%
\definecolor{currentstroke}{rgb}{0.000000,0.000000,0.000000}%
\pgfsetstrokecolor{currentstroke}%
\pgfsetdash{}{0pt}%
\pgfpathmoveto{\pgfqpoint{1.422616in}{1.055682in}}%
\pgfpathlineto{\pgfqpoint{1.437827in}{1.060600in}}%
\pgfpathlineto{\pgfqpoint{1.438213in}{1.083395in}}%
\pgfpathlineto{\pgfqpoint{1.422983in}{1.078549in}}%
\pgfpathlineto{\pgfqpoint{1.422616in}{1.055682in}}%
\pgfpathclose%
\pgfusepath{fill}%
\end{pgfscope}%
\begin{pgfscope}%
\pgfpathrectangle{\pgfqpoint{0.407592in}{0.222904in}}{\pgfqpoint{1.411163in}{1.411163in}}%
\pgfusepath{clip}%
\pgfsetbuttcap%
\pgfsetroundjoin%
\definecolor{currentfill}{rgb}{0.168627,0.611765,0.411765}%
\pgfsetfillcolor{currentfill}%
\pgfsetfillopacity{0.180000}%
\pgfsetlinewidth{0.000000pt}%
\definecolor{currentstroke}{rgb}{0.000000,0.000000,0.000000}%
\pgfsetstrokecolor{currentstroke}%
\pgfsetdash{}{0pt}%
\pgfpathmoveto{\pgfqpoint{1.438986in}{1.129158in}}%
\pgfpathlineto{\pgfqpoint{1.454176in}{1.135428in}}%
\pgfpathlineto{\pgfqpoint{1.454582in}{1.158328in}}%
\pgfpathlineto{\pgfqpoint{1.439374in}{1.152126in}}%
\pgfpathlineto{\pgfqpoint{1.438986in}{1.129158in}}%
\pgfpathclose%
\pgfusepath{fill}%
\end{pgfscope}%
\begin{pgfscope}%
\pgfpathrectangle{\pgfqpoint{0.407592in}{0.222904in}}{\pgfqpoint{1.411163in}{1.411163in}}%
\pgfusepath{clip}%
\pgfsetbuttcap%
\pgfsetroundjoin%
\definecolor{currentfill}{rgb}{0.168627,0.611765,0.411765}%
\pgfsetfillcolor{currentfill}%
\pgfsetfillopacity{0.180000}%
\pgfsetlinewidth{0.000000pt}%
\definecolor{currentstroke}{rgb}{0.000000,0.000000,0.000000}%
\pgfsetstrokecolor{currentstroke}%
\pgfsetdash{}{0pt}%
\pgfpathmoveto{\pgfqpoint{1.374005in}{0.867037in}}%
\pgfpathlineto{\pgfqpoint{1.389328in}{0.868242in}}%
\pgfpathlineto{\pgfqpoint{1.389653in}{0.890803in}}%
\pgfpathlineto{\pgfqpoint{1.374311in}{0.889679in}}%
\pgfpathlineto{\pgfqpoint{1.374005in}{0.867037in}}%
\pgfpathclose%
\pgfusepath{fill}%
\end{pgfscope}%
\begin{pgfscope}%
\pgfpathrectangle{\pgfqpoint{0.407592in}{0.222904in}}{\pgfqpoint{1.411163in}{1.411163in}}%
\pgfusepath{clip}%
\pgfsetbuttcap%
\pgfsetroundjoin%
\definecolor{currentfill}{rgb}{0.168627,0.611765,0.411765}%
\pgfsetfillcolor{currentfill}%
\pgfsetfillopacity{0.180000}%
\pgfsetlinewidth{0.000000pt}%
\definecolor{currentstroke}{rgb}{0.000000,0.000000,0.000000}%
\pgfsetstrokecolor{currentstroke}%
\pgfsetdash{}{0pt}%
\pgfpathmoveto{\pgfqpoint{1.390303in}{0.936096in}}%
\pgfpathlineto{\pgfqpoint{1.405591in}{0.938433in}}%
\pgfpathlineto{\pgfqpoint{1.405936in}{0.961088in}}%
\pgfpathlineto{\pgfqpoint{1.390630in}{0.958829in}}%
\pgfpathlineto{\pgfqpoint{1.390303in}{0.936096in}}%
\pgfpathclose%
\pgfusepath{fill}%
\end{pgfscope}%
\begin{pgfscope}%
\pgfpathrectangle{\pgfqpoint{0.407592in}{0.222904in}}{\pgfqpoint{1.411163in}{1.411163in}}%
\pgfusepath{clip}%
\pgfsetbuttcap%
\pgfsetroundjoin%
\definecolor{currentfill}{rgb}{0.168627,0.611765,0.411765}%
\pgfsetfillcolor{currentfill}%
\pgfsetfillopacity{0.180000}%
\pgfsetlinewidth{0.000000pt}%
\definecolor{currentstroke}{rgb}{0.000000,0.000000,0.000000}%
\pgfsetstrokecolor{currentstroke}%
\pgfsetdash{}{0pt}%
\pgfpathmoveto{\pgfqpoint{1.406628in}{1.006568in}}%
\pgfpathlineto{\pgfqpoint{1.421884in}{1.010122in}}%
\pgfpathlineto{\pgfqpoint{1.422249in}{1.032873in}}%
\pgfpathlineto{\pgfqpoint{1.406975in}{1.029395in}}%
\pgfpathlineto{\pgfqpoint{1.406628in}{1.006568in}}%
\pgfpathclose%
\pgfusepath{fill}%
\end{pgfscope}%
\begin{pgfscope}%
\pgfpathrectangle{\pgfqpoint{0.407592in}{0.222904in}}{\pgfqpoint{1.411163in}{1.411163in}}%
\pgfusepath{clip}%
\pgfsetbuttcap%
\pgfsetroundjoin%
\definecolor{currentfill}{rgb}{0.168627,0.611765,0.411765}%
\pgfsetfillcolor{currentfill}%
\pgfsetfillopacity{0.180000}%
\pgfsetlinewidth{0.000000pt}%
\definecolor{currentstroke}{rgb}{0.000000,0.000000,0.000000}%
\pgfsetstrokecolor{currentstroke}%
\pgfsetdash{}{0pt}%
\pgfpathmoveto{\pgfqpoint{1.422983in}{1.078549in}}%
\pgfpathlineto{\pgfqpoint{1.438213in}{1.083395in}}%
\pgfpathlineto{\pgfqpoint{1.438599in}{1.106248in}}%
\pgfpathlineto{\pgfqpoint{1.423351in}{1.101474in}}%
\pgfpathlineto{\pgfqpoint{1.422983in}{1.078549in}}%
\pgfpathclose%
\pgfusepath{fill}%
\end{pgfscope}%
\begin{pgfscope}%
\pgfpathrectangle{\pgfqpoint{0.407592in}{0.222904in}}{\pgfqpoint{1.411163in}{1.411163in}}%
\pgfusepath{clip}%
\pgfsetbuttcap%
\pgfsetroundjoin%
\definecolor{currentfill}{rgb}{0.168627,0.611765,0.411765}%
\pgfsetfillcolor{currentfill}%
\pgfsetfillopacity{0.180000}%
\pgfsetlinewidth{0.000000pt}%
\definecolor{currentstroke}{rgb}{0.000000,0.000000,0.000000}%
\pgfsetstrokecolor{currentstroke}%
\pgfsetdash{}{0pt}%
\pgfpathmoveto{\pgfqpoint{1.439374in}{1.152126in}}%
\pgfpathlineto{\pgfqpoint{1.454582in}{1.158328in}}%
\pgfpathlineto{\pgfqpoint{1.454989in}{1.181286in}}%
\pgfpathlineto{\pgfqpoint{1.439763in}{1.175153in}}%
\pgfpathlineto{\pgfqpoint{1.439374in}{1.152126in}}%
\pgfpathclose%
\pgfusepath{fill}%
\end{pgfscope}%
\begin{pgfscope}%
\pgfpathrectangle{\pgfqpoint{0.407592in}{0.222904in}}{\pgfqpoint{1.411163in}{1.411163in}}%
\pgfusepath{clip}%
\pgfsetbuttcap%
\pgfsetroundjoin%
\definecolor{currentfill}{rgb}{0.168627,0.611765,0.411765}%
\pgfsetfillcolor{currentfill}%
\pgfsetfillopacity{0.180000}%
\pgfsetlinewidth{0.000000pt}%
\definecolor{currentstroke}{rgb}{0.000000,0.000000,0.000000}%
\pgfsetstrokecolor{currentstroke}%
\pgfsetdash{}{0pt}%
\pgfpathmoveto{\pgfqpoint{1.374311in}{0.889679in}}%
\pgfpathlineto{\pgfqpoint{1.389653in}{0.890803in}}%
\pgfpathlineto{\pgfqpoint{1.389978in}{0.913421in}}%
\pgfpathlineto{\pgfqpoint{1.374618in}{0.912379in}}%
\pgfpathlineto{\pgfqpoint{1.374311in}{0.889679in}}%
\pgfpathclose%
\pgfusepath{fill}%
\end{pgfscope}%
\begin{pgfscope}%
\pgfpathrectangle{\pgfqpoint{0.407592in}{0.222904in}}{\pgfqpoint{1.411163in}{1.411163in}}%
\pgfusepath{clip}%
\pgfsetbuttcap%
\pgfsetroundjoin%
\definecolor{currentfill}{rgb}{0.168627,0.611765,0.411765}%
\pgfsetfillcolor{currentfill}%
\pgfsetfillopacity{0.180000}%
\pgfsetlinewidth{0.000000pt}%
\definecolor{currentstroke}{rgb}{0.000000,0.000000,0.000000}%
\pgfsetstrokecolor{currentstroke}%
\pgfsetdash{}{0pt}%
\pgfpathmoveto{\pgfqpoint{1.390630in}{0.958829in}}%
\pgfpathlineto{\pgfqpoint{1.405936in}{0.961088in}}%
\pgfpathlineto{\pgfqpoint{1.406281in}{0.983799in}}%
\pgfpathlineto{\pgfqpoint{1.390958in}{0.981619in}}%
\pgfpathlineto{\pgfqpoint{1.390630in}{0.958829in}}%
\pgfpathclose%
\pgfusepath{fill}%
\end{pgfscope}%
\begin{pgfscope}%
\pgfpathrectangle{\pgfqpoint{0.407592in}{0.222904in}}{\pgfqpoint{1.411163in}{1.411163in}}%
\pgfusepath{clip}%
\pgfsetbuttcap%
\pgfsetroundjoin%
\definecolor{currentfill}{rgb}{0.168627,0.611765,0.411765}%
\pgfsetfillcolor{currentfill}%
\pgfsetfillopacity{0.180000}%
\pgfsetlinewidth{0.000000pt}%
\definecolor{currentstroke}{rgb}{0.000000,0.000000,0.000000}%
\pgfsetstrokecolor{currentstroke}%
\pgfsetdash{}{0pt}%
\pgfpathmoveto{\pgfqpoint{1.406975in}{1.029395in}}%
\pgfpathlineto{\pgfqpoint{1.422249in}{1.032873in}}%
\pgfpathlineto{\pgfqpoint{1.422616in}{1.055682in}}%
\pgfpathlineto{\pgfqpoint{1.407324in}{1.052280in}}%
\pgfpathlineto{\pgfqpoint{1.406975in}{1.029395in}}%
\pgfpathclose%
\pgfusepath{fill}%
\end{pgfscope}%
\begin{pgfscope}%
\pgfpathrectangle{\pgfqpoint{0.407592in}{0.222904in}}{\pgfqpoint{1.411163in}{1.411163in}}%
\pgfusepath{clip}%
\pgfsetbuttcap%
\pgfsetroundjoin%
\definecolor{currentfill}{rgb}{0.168627,0.611765,0.411765}%
\pgfsetfillcolor{currentfill}%
\pgfsetfillopacity{0.180000}%
\pgfsetlinewidth{0.000000pt}%
\definecolor{currentstroke}{rgb}{0.000000,0.000000,0.000000}%
\pgfsetstrokecolor{currentstroke}%
\pgfsetdash{}{0pt}%
\pgfpathmoveto{\pgfqpoint{1.423351in}{1.101474in}}%
\pgfpathlineto{\pgfqpoint{1.438599in}{1.106248in}}%
\pgfpathlineto{\pgfqpoint{1.438986in}{1.129158in}}%
\pgfpathlineto{\pgfqpoint{1.423720in}{1.124457in}}%
\pgfpathlineto{\pgfqpoint{1.423351in}{1.101474in}}%
\pgfpathclose%
\pgfusepath{fill}%
\end{pgfscope}%
\begin{pgfscope}%
\pgfpathrectangle{\pgfqpoint{0.407592in}{0.222904in}}{\pgfqpoint{1.411163in}{1.411163in}}%
\pgfusepath{clip}%
\pgfsetbuttcap%
\pgfsetroundjoin%
\definecolor{currentfill}{rgb}{0.168627,0.611765,0.411765}%
\pgfsetfillcolor{currentfill}%
\pgfsetfillopacity{0.180000}%
\pgfsetlinewidth{0.000000pt}%
\definecolor{currentstroke}{rgb}{0.000000,0.000000,0.000000}%
\pgfsetstrokecolor{currentstroke}%
\pgfsetdash{}{0pt}%
\pgfpathmoveto{\pgfqpoint{1.374618in}{0.912379in}}%
\pgfpathlineto{\pgfqpoint{1.389978in}{0.913421in}}%
\pgfpathlineto{\pgfqpoint{1.390303in}{0.936096in}}%
\pgfpathlineto{\pgfqpoint{1.374925in}{0.935136in}}%
\pgfpathlineto{\pgfqpoint{1.374618in}{0.912379in}}%
\pgfpathclose%
\pgfusepath{fill}%
\end{pgfscope}%
\begin{pgfscope}%
\pgfpathrectangle{\pgfqpoint{0.407592in}{0.222904in}}{\pgfqpoint{1.411163in}{1.411163in}}%
\pgfusepath{clip}%
\pgfsetbuttcap%
\pgfsetroundjoin%
\definecolor{currentfill}{rgb}{0.168627,0.611765,0.411765}%
\pgfsetfillcolor{currentfill}%
\pgfsetfillopacity{0.180000}%
\pgfsetlinewidth{0.000000pt}%
\definecolor{currentstroke}{rgb}{0.000000,0.000000,0.000000}%
\pgfsetstrokecolor{currentstroke}%
\pgfsetdash{}{0pt}%
\pgfpathmoveto{\pgfqpoint{1.390958in}{0.981619in}}%
\pgfpathlineto{\pgfqpoint{1.406281in}{0.983799in}}%
\pgfpathlineto{\pgfqpoint{1.406628in}{1.006568in}}%
\pgfpathlineto{\pgfqpoint{1.391286in}{1.004467in}}%
\pgfpathlineto{\pgfqpoint{1.390958in}{0.981619in}}%
\pgfpathclose%
\pgfusepath{fill}%
\end{pgfscope}%
\begin{pgfscope}%
\pgfpathrectangle{\pgfqpoint{0.407592in}{0.222904in}}{\pgfqpoint{1.411163in}{1.411163in}}%
\pgfusepath{clip}%
\pgfsetbuttcap%
\pgfsetroundjoin%
\definecolor{currentfill}{rgb}{0.168627,0.611765,0.411765}%
\pgfsetfillcolor{currentfill}%
\pgfsetfillopacity{0.180000}%
\pgfsetlinewidth{0.000000pt}%
\definecolor{currentstroke}{rgb}{0.000000,0.000000,0.000000}%
\pgfsetstrokecolor{currentstroke}%
\pgfsetdash{}{0pt}%
\pgfpathmoveto{\pgfqpoint{1.407324in}{1.052280in}}%
\pgfpathlineto{\pgfqpoint{1.422616in}{1.055682in}}%
\pgfpathlineto{\pgfqpoint{1.422983in}{1.078549in}}%
\pgfpathlineto{\pgfqpoint{1.407673in}{1.075222in}}%
\pgfpathlineto{\pgfqpoint{1.407324in}{1.052280in}}%
\pgfpathclose%
\pgfusepath{fill}%
\end{pgfscope}%
\begin{pgfscope}%
\pgfpathrectangle{\pgfqpoint{0.407592in}{0.222904in}}{\pgfqpoint{1.411163in}{1.411163in}}%
\pgfusepath{clip}%
\pgfsetbuttcap%
\pgfsetroundjoin%
\definecolor{currentfill}{rgb}{0.168627,0.611765,0.411765}%
\pgfsetfillcolor{currentfill}%
\pgfsetfillopacity{0.180000}%
\pgfsetlinewidth{0.000000pt}%
\definecolor{currentstroke}{rgb}{0.000000,0.000000,0.000000}%
\pgfsetstrokecolor{currentstroke}%
\pgfsetdash{}{0pt}%
\pgfpathmoveto{\pgfqpoint{1.423720in}{1.124457in}}%
\pgfpathlineto{\pgfqpoint{1.438986in}{1.129158in}}%
\pgfpathlineto{\pgfqpoint{1.439374in}{1.152126in}}%
\pgfpathlineto{\pgfqpoint{1.424090in}{1.147498in}}%
\pgfpathlineto{\pgfqpoint{1.423720in}{1.124457in}}%
\pgfpathclose%
\pgfusepath{fill}%
\end{pgfscope}%
\begin{pgfscope}%
\pgfpathrectangle{\pgfqpoint{0.407592in}{0.222904in}}{\pgfqpoint{1.411163in}{1.411163in}}%
\pgfusepath{clip}%
\pgfsetbuttcap%
\pgfsetroundjoin%
\definecolor{currentfill}{rgb}{0.168627,0.611765,0.411765}%
\pgfsetfillcolor{currentfill}%
\pgfsetfillopacity{0.180000}%
\pgfsetlinewidth{0.000000pt}%
\definecolor{currentstroke}{rgb}{0.000000,0.000000,0.000000}%
\pgfsetstrokecolor{currentstroke}%
\pgfsetdash{}{0pt}%
\pgfpathmoveto{\pgfqpoint{1.358585in}{0.867127in}}%
\pgfpathlineto{\pgfqpoint{1.374005in}{0.867037in}}%
\pgfpathlineto{\pgfqpoint{1.374311in}{0.889679in}}%
\pgfpathlineto{\pgfqpoint{1.358873in}{0.889854in}}%
\pgfpathlineto{\pgfqpoint{1.358585in}{0.867127in}}%
\pgfpathclose%
\pgfusepath{fill}%
\end{pgfscope}%
\begin{pgfscope}%
\pgfpathrectangle{\pgfqpoint{0.407592in}{0.222904in}}{\pgfqpoint{1.411163in}{1.411163in}}%
\pgfusepath{clip}%
\pgfsetbuttcap%
\pgfsetroundjoin%
\definecolor{currentfill}{rgb}{0.168627,0.611765,0.411765}%
\pgfsetfillcolor{currentfill}%
\pgfsetfillopacity{0.180000}%
\pgfsetlinewidth{0.000000pt}%
\definecolor{currentstroke}{rgb}{0.000000,0.000000,0.000000}%
\pgfsetstrokecolor{currentstroke}%
\pgfsetdash{}{0pt}%
\pgfpathmoveto{\pgfqpoint{1.374925in}{0.935136in}}%
\pgfpathlineto{\pgfqpoint{1.390303in}{0.936096in}}%
\pgfpathlineto{\pgfqpoint{1.390630in}{0.958829in}}%
\pgfpathlineto{\pgfqpoint{1.375233in}{0.957951in}}%
\pgfpathlineto{\pgfqpoint{1.374925in}{0.935136in}}%
\pgfpathclose%
\pgfusepath{fill}%
\end{pgfscope}%
\begin{pgfscope}%
\pgfpathrectangle{\pgfqpoint{0.407592in}{0.222904in}}{\pgfqpoint{1.411163in}{1.411163in}}%
\pgfusepath{clip}%
\pgfsetbuttcap%
\pgfsetroundjoin%
\definecolor{currentfill}{rgb}{0.168627,0.611765,0.411765}%
\pgfsetfillcolor{currentfill}%
\pgfsetfillopacity{0.180000}%
\pgfsetlinewidth{0.000000pt}%
\definecolor{currentstroke}{rgb}{0.000000,0.000000,0.000000}%
\pgfsetstrokecolor{currentstroke}%
\pgfsetdash{}{0pt}%
\pgfpathmoveto{\pgfqpoint{1.391286in}{1.004467in}}%
\pgfpathlineto{\pgfqpoint{1.406628in}{1.006568in}}%
\pgfpathlineto{\pgfqpoint{1.406975in}{1.029395in}}%
\pgfpathlineto{\pgfqpoint{1.391615in}{1.027374in}}%
\pgfpathlineto{\pgfqpoint{1.391286in}{1.004467in}}%
\pgfpathclose%
\pgfusepath{fill}%
\end{pgfscope}%
\begin{pgfscope}%
\pgfpathrectangle{\pgfqpoint{0.407592in}{0.222904in}}{\pgfqpoint{1.411163in}{1.411163in}}%
\pgfusepath{clip}%
\pgfsetbuttcap%
\pgfsetroundjoin%
\definecolor{currentfill}{rgb}{0.168627,0.611765,0.411765}%
\pgfsetfillcolor{currentfill}%
\pgfsetfillopacity{0.180000}%
\pgfsetlinewidth{0.000000pt}%
\definecolor{currentstroke}{rgb}{0.000000,0.000000,0.000000}%
\pgfsetstrokecolor{currentstroke}%
\pgfsetdash{}{0pt}%
\pgfpathmoveto{\pgfqpoint{1.407673in}{1.075222in}}%
\pgfpathlineto{\pgfqpoint{1.422983in}{1.078549in}}%
\pgfpathlineto{\pgfqpoint{1.423351in}{1.101474in}}%
\pgfpathlineto{\pgfqpoint{1.408023in}{1.098223in}}%
\pgfpathlineto{\pgfqpoint{1.407673in}{1.075222in}}%
\pgfpathclose%
\pgfusepath{fill}%
\end{pgfscope}%
\begin{pgfscope}%
\pgfpathrectangle{\pgfqpoint{0.407592in}{0.222904in}}{\pgfqpoint{1.411163in}{1.411163in}}%
\pgfusepath{clip}%
\pgfsetbuttcap%
\pgfsetroundjoin%
\definecolor{currentfill}{rgb}{0.168627,0.611765,0.411765}%
\pgfsetfillcolor{currentfill}%
\pgfsetfillopacity{0.180000}%
\pgfsetlinewidth{0.000000pt}%
\definecolor{currentstroke}{rgb}{0.000000,0.000000,0.000000}%
\pgfsetstrokecolor{currentstroke}%
\pgfsetdash{}{0pt}%
\pgfpathmoveto{\pgfqpoint{1.424090in}{1.147498in}}%
\pgfpathlineto{\pgfqpoint{1.439374in}{1.152126in}}%
\pgfpathlineto{\pgfqpoint{1.439763in}{1.175153in}}%
\pgfpathlineto{\pgfqpoint{1.424461in}{1.170598in}}%
\pgfpathlineto{\pgfqpoint{1.424090in}{1.147498in}}%
\pgfpathclose%
\pgfusepath{fill}%
\end{pgfscope}%
\begin{pgfscope}%
\pgfpathrectangle{\pgfqpoint{0.407592in}{0.222904in}}{\pgfqpoint{1.411163in}{1.411163in}}%
\pgfusepath{clip}%
\pgfsetbuttcap%
\pgfsetroundjoin%
\definecolor{currentfill}{rgb}{0.168627,0.611765,0.411765}%
\pgfsetfillcolor{currentfill}%
\pgfsetfillopacity{0.180000}%
\pgfsetlinewidth{0.000000pt}%
\definecolor{currentstroke}{rgb}{0.000000,0.000000,0.000000}%
\pgfsetstrokecolor{currentstroke}%
\pgfsetdash{}{0pt}%
\pgfpathmoveto{\pgfqpoint{1.358873in}{0.889854in}}%
\pgfpathlineto{\pgfqpoint{1.374311in}{0.889679in}}%
\pgfpathlineto{\pgfqpoint{1.374618in}{0.912379in}}%
\pgfpathlineto{\pgfqpoint{1.359161in}{0.912639in}}%
\pgfpathlineto{\pgfqpoint{1.358873in}{0.889854in}}%
\pgfpathclose%
\pgfusepath{fill}%
\end{pgfscope}%
\begin{pgfscope}%
\pgfpathrectangle{\pgfqpoint{0.407592in}{0.222904in}}{\pgfqpoint{1.411163in}{1.411163in}}%
\pgfusepath{clip}%
\pgfsetbuttcap%
\pgfsetroundjoin%
\definecolor{currentfill}{rgb}{0.168627,0.611765,0.411765}%
\pgfsetfillcolor{currentfill}%
\pgfsetfillopacity{0.180000}%
\pgfsetlinewidth{0.000000pt}%
\definecolor{currentstroke}{rgb}{0.000000,0.000000,0.000000}%
\pgfsetstrokecolor{currentstroke}%
\pgfsetdash{}{0pt}%
\pgfpathmoveto{\pgfqpoint{1.375233in}{0.957951in}}%
\pgfpathlineto{\pgfqpoint{1.390630in}{0.958829in}}%
\pgfpathlineto{\pgfqpoint{1.390958in}{0.981619in}}%
\pgfpathlineto{\pgfqpoint{1.375542in}{0.980824in}}%
\pgfpathlineto{\pgfqpoint{1.375233in}{0.957951in}}%
\pgfpathclose%
\pgfusepath{fill}%
\end{pgfscope}%
\begin{pgfscope}%
\pgfpathrectangle{\pgfqpoint{0.407592in}{0.222904in}}{\pgfqpoint{1.411163in}{1.411163in}}%
\pgfusepath{clip}%
\pgfsetbuttcap%
\pgfsetroundjoin%
\definecolor{currentfill}{rgb}{0.168627,0.611765,0.411765}%
\pgfsetfillcolor{currentfill}%
\pgfsetfillopacity{0.180000}%
\pgfsetlinewidth{0.000000pt}%
\definecolor{currentstroke}{rgb}{0.000000,0.000000,0.000000}%
\pgfsetstrokecolor{currentstroke}%
\pgfsetdash{}{0pt}%
\pgfpathmoveto{\pgfqpoint{1.391615in}{1.027374in}}%
\pgfpathlineto{\pgfqpoint{1.406975in}{1.029395in}}%
\pgfpathlineto{\pgfqpoint{1.407324in}{1.052280in}}%
\pgfpathlineto{\pgfqpoint{1.391945in}{1.050338in}}%
\pgfpathlineto{\pgfqpoint{1.391615in}{1.027374in}}%
\pgfpathclose%
\pgfusepath{fill}%
\end{pgfscope}%
\begin{pgfscope}%
\pgfpathrectangle{\pgfqpoint{0.407592in}{0.222904in}}{\pgfqpoint{1.411163in}{1.411163in}}%
\pgfusepath{clip}%
\pgfsetbuttcap%
\pgfsetroundjoin%
\definecolor{currentfill}{rgb}{0.168627,0.611765,0.411765}%
\pgfsetfillcolor{currentfill}%
\pgfsetfillopacity{0.180000}%
\pgfsetlinewidth{0.000000pt}%
\definecolor{currentstroke}{rgb}{0.000000,0.000000,0.000000}%
\pgfsetstrokecolor{currentstroke}%
\pgfsetdash{}{0pt}%
\pgfpathmoveto{\pgfqpoint{1.408023in}{1.098223in}}%
\pgfpathlineto{\pgfqpoint{1.423351in}{1.101474in}}%
\pgfpathlineto{\pgfqpoint{1.423720in}{1.124457in}}%
\pgfpathlineto{\pgfqpoint{1.408374in}{1.121283in}}%
\pgfpathlineto{\pgfqpoint{1.408023in}{1.098223in}}%
\pgfpathclose%
\pgfusepath{fill}%
\end{pgfscope}%
\begin{pgfscope}%
\pgfpathrectangle{\pgfqpoint{0.407592in}{0.222904in}}{\pgfqpoint{1.411163in}{1.411163in}}%
\pgfusepath{clip}%
\pgfsetbuttcap%
\pgfsetroundjoin%
\definecolor{currentfill}{rgb}{0.168627,0.611765,0.411765}%
\pgfsetfillcolor{currentfill}%
\pgfsetfillopacity{0.180000}%
\pgfsetlinewidth{0.000000pt}%
\definecolor{currentstroke}{rgb}{0.000000,0.000000,0.000000}%
\pgfsetstrokecolor{currentstroke}%
\pgfsetdash{}{0pt}%
\pgfpathmoveto{\pgfqpoint{1.359161in}{0.912639in}}%
\pgfpathlineto{\pgfqpoint{1.374618in}{0.912379in}}%
\pgfpathlineto{\pgfqpoint{1.374925in}{0.935136in}}%
\pgfpathlineto{\pgfqpoint{1.359450in}{0.935482in}}%
\pgfpathlineto{\pgfqpoint{1.359161in}{0.912639in}}%
\pgfpathclose%
\pgfusepath{fill}%
\end{pgfscope}%
\begin{pgfscope}%
\pgfpathrectangle{\pgfqpoint{0.407592in}{0.222904in}}{\pgfqpoint{1.411163in}{1.411163in}}%
\pgfusepath{clip}%
\pgfsetbuttcap%
\pgfsetroundjoin%
\definecolor{currentfill}{rgb}{0.168627,0.611765,0.411765}%
\pgfsetfillcolor{currentfill}%
\pgfsetfillopacity{0.180000}%
\pgfsetlinewidth{0.000000pt}%
\definecolor{currentstroke}{rgb}{0.000000,0.000000,0.000000}%
\pgfsetstrokecolor{currentstroke}%
\pgfsetdash{}{0pt}%
\pgfpathmoveto{\pgfqpoint{1.375542in}{0.980824in}}%
\pgfpathlineto{\pgfqpoint{1.390958in}{0.981619in}}%
\pgfpathlineto{\pgfqpoint{1.391286in}{1.004467in}}%
\pgfpathlineto{\pgfqpoint{1.375852in}{1.003755in}}%
\pgfpathlineto{\pgfqpoint{1.375542in}{0.980824in}}%
\pgfpathclose%
\pgfusepath{fill}%
\end{pgfscope}%
\begin{pgfscope}%
\pgfpathrectangle{\pgfqpoint{0.407592in}{0.222904in}}{\pgfqpoint{1.411163in}{1.411163in}}%
\pgfusepath{clip}%
\pgfsetbuttcap%
\pgfsetroundjoin%
\definecolor{currentfill}{rgb}{0.168627,0.611765,0.411765}%
\pgfsetfillcolor{currentfill}%
\pgfsetfillopacity{0.180000}%
\pgfsetlinewidth{0.000000pt}%
\definecolor{currentstroke}{rgb}{0.000000,0.000000,0.000000}%
\pgfsetstrokecolor{currentstroke}%
\pgfsetdash{}{0pt}%
\pgfpathmoveto{\pgfqpoint{1.391945in}{1.050338in}}%
\pgfpathlineto{\pgfqpoint{1.407324in}{1.052280in}}%
\pgfpathlineto{\pgfqpoint{1.407673in}{1.075222in}}%
\pgfpathlineto{\pgfqpoint{1.392276in}{1.073361in}}%
\pgfpathlineto{\pgfqpoint{1.391945in}{1.050338in}}%
\pgfpathclose%
\pgfusepath{fill}%
\end{pgfscope}%
\begin{pgfscope}%
\pgfpathrectangle{\pgfqpoint{0.407592in}{0.222904in}}{\pgfqpoint{1.411163in}{1.411163in}}%
\pgfusepath{clip}%
\pgfsetbuttcap%
\pgfsetroundjoin%
\definecolor{currentfill}{rgb}{0.168627,0.611765,0.411765}%
\pgfsetfillcolor{currentfill}%
\pgfsetfillopacity{0.180000}%
\pgfsetlinewidth{0.000000pt}%
\definecolor{currentstroke}{rgb}{0.000000,0.000000,0.000000}%
\pgfsetstrokecolor{currentstroke}%
\pgfsetdash{}{0pt}%
\pgfpathmoveto{\pgfqpoint{1.408374in}{1.121283in}}%
\pgfpathlineto{\pgfqpoint{1.423720in}{1.124457in}}%
\pgfpathlineto{\pgfqpoint{1.424090in}{1.147498in}}%
\pgfpathlineto{\pgfqpoint{1.408726in}{1.144401in}}%
\pgfpathlineto{\pgfqpoint{1.408374in}{1.121283in}}%
\pgfpathclose%
\pgfusepath{fill}%
\end{pgfscope}%
\begin{pgfscope}%
\pgfpathrectangle{\pgfqpoint{0.407592in}{0.222904in}}{\pgfqpoint{1.411163in}{1.411163in}}%
\pgfusepath{clip}%
\pgfsetbuttcap%
\pgfsetroundjoin%
\definecolor{currentfill}{rgb}{0.168627,0.611765,0.411765}%
\pgfsetfillcolor{currentfill}%
\pgfsetfillopacity{0.180000}%
\pgfsetlinewidth{0.000000pt}%
\definecolor{currentstroke}{rgb}{0.000000,0.000000,0.000000}%
\pgfsetstrokecolor{currentstroke}%
\pgfsetdash{}{0pt}%
\pgfpathmoveto{\pgfqpoint{1.343064in}{0.868436in}}%
\pgfpathlineto{\pgfqpoint{1.358585in}{0.867127in}}%
\pgfpathlineto{\pgfqpoint{1.358873in}{0.889854in}}%
\pgfpathlineto{\pgfqpoint{1.343333in}{0.891252in}}%
\pgfpathlineto{\pgfqpoint{1.343064in}{0.868436in}}%
\pgfpathclose%
\pgfusepath{fill}%
\end{pgfscope}%
\begin{pgfscope}%
\pgfpathrectangle{\pgfqpoint{0.407592in}{0.222904in}}{\pgfqpoint{1.411163in}{1.411163in}}%
\pgfusepath{clip}%
\pgfsetbuttcap%
\pgfsetroundjoin%
\definecolor{currentfill}{rgb}{0.168627,0.611765,0.411765}%
\pgfsetfillcolor{currentfill}%
\pgfsetfillopacity{0.180000}%
\pgfsetlinewidth{0.000000pt}%
\definecolor{currentstroke}{rgb}{0.000000,0.000000,0.000000}%
\pgfsetstrokecolor{currentstroke}%
\pgfsetdash{}{0pt}%
\pgfpathmoveto{\pgfqpoint{1.359450in}{0.935482in}}%
\pgfpathlineto{\pgfqpoint{1.374925in}{0.935136in}}%
\pgfpathlineto{\pgfqpoint{1.375233in}{0.957951in}}%
\pgfpathlineto{\pgfqpoint{1.359740in}{0.958383in}}%
\pgfpathlineto{\pgfqpoint{1.359450in}{0.935482in}}%
\pgfpathclose%
\pgfusepath{fill}%
\end{pgfscope}%
\begin{pgfscope}%
\pgfpathrectangle{\pgfqpoint{0.407592in}{0.222904in}}{\pgfqpoint{1.411163in}{1.411163in}}%
\pgfusepath{clip}%
\pgfsetbuttcap%
\pgfsetroundjoin%
\definecolor{currentfill}{rgb}{0.168627,0.611765,0.411765}%
\pgfsetfillcolor{currentfill}%
\pgfsetfillopacity{0.180000}%
\pgfsetlinewidth{0.000000pt}%
\definecolor{currentstroke}{rgb}{0.000000,0.000000,0.000000}%
\pgfsetstrokecolor{currentstroke}%
\pgfsetdash{}{0pt}%
\pgfpathmoveto{\pgfqpoint{1.375852in}{1.003755in}}%
\pgfpathlineto{\pgfqpoint{1.391286in}{1.004467in}}%
\pgfpathlineto{\pgfqpoint{1.391615in}{1.027374in}}%
\pgfpathlineto{\pgfqpoint{1.376163in}{1.026745in}}%
\pgfpathlineto{\pgfqpoint{1.375852in}{1.003755in}}%
\pgfpathclose%
\pgfusepath{fill}%
\end{pgfscope}%
\begin{pgfscope}%
\pgfpathrectangle{\pgfqpoint{0.407592in}{0.222904in}}{\pgfqpoint{1.411163in}{1.411163in}}%
\pgfusepath{clip}%
\pgfsetbuttcap%
\pgfsetroundjoin%
\definecolor{currentfill}{rgb}{0.168627,0.611765,0.411765}%
\pgfsetfillcolor{currentfill}%
\pgfsetfillopacity{0.180000}%
\pgfsetlinewidth{0.000000pt}%
\definecolor{currentstroke}{rgb}{0.000000,0.000000,0.000000}%
\pgfsetstrokecolor{currentstroke}%
\pgfsetdash{}{0pt}%
\pgfpathmoveto{\pgfqpoint{1.392276in}{1.073361in}}%
\pgfpathlineto{\pgfqpoint{1.407673in}{1.075222in}}%
\pgfpathlineto{\pgfqpoint{1.408023in}{1.098223in}}%
\pgfpathlineto{\pgfqpoint{1.392608in}{1.096443in}}%
\pgfpathlineto{\pgfqpoint{1.392276in}{1.073361in}}%
\pgfpathclose%
\pgfusepath{fill}%
\end{pgfscope}%
\begin{pgfscope}%
\pgfpathrectangle{\pgfqpoint{0.407592in}{0.222904in}}{\pgfqpoint{1.411163in}{1.411163in}}%
\pgfusepath{clip}%
\pgfsetbuttcap%
\pgfsetroundjoin%
\definecolor{currentfill}{rgb}{0.168627,0.611765,0.411765}%
\pgfsetfillcolor{currentfill}%
\pgfsetfillopacity{0.180000}%
\pgfsetlinewidth{0.000000pt}%
\definecolor{currentstroke}{rgb}{0.000000,0.000000,0.000000}%
\pgfsetstrokecolor{currentstroke}%
\pgfsetdash{}{0pt}%
\pgfpathmoveto{\pgfqpoint{1.408726in}{1.144401in}}%
\pgfpathlineto{\pgfqpoint{1.424090in}{1.147498in}}%
\pgfpathlineto{\pgfqpoint{1.424461in}{1.170598in}}%
\pgfpathlineto{\pgfqpoint{1.409078in}{1.167579in}}%
\pgfpathlineto{\pgfqpoint{1.408726in}{1.144401in}}%
\pgfpathclose%
\pgfusepath{fill}%
\end{pgfscope}%
\begin{pgfscope}%
\pgfpathrectangle{\pgfqpoint{0.407592in}{0.222904in}}{\pgfqpoint{1.411163in}{1.411163in}}%
\pgfusepath{clip}%
\pgfsetbuttcap%
\pgfsetroundjoin%
\definecolor{currentfill}{rgb}{0.168627,0.611765,0.411765}%
\pgfsetfillcolor{currentfill}%
\pgfsetfillopacity{0.180000}%
\pgfsetlinewidth{0.000000pt}%
\definecolor{currentstroke}{rgb}{0.000000,0.000000,0.000000}%
\pgfsetstrokecolor{currentstroke}%
\pgfsetdash{}{0pt}%
\pgfpathmoveto{\pgfqpoint{1.343333in}{0.891252in}}%
\pgfpathlineto{\pgfqpoint{1.358873in}{0.889854in}}%
\pgfpathlineto{\pgfqpoint{1.359161in}{0.912639in}}%
\pgfpathlineto{\pgfqpoint{1.343602in}{0.914126in}}%
\pgfpathlineto{\pgfqpoint{1.343333in}{0.891252in}}%
\pgfpathclose%
\pgfusepath{fill}%
\end{pgfscope}%
\begin{pgfscope}%
\pgfpathrectangle{\pgfqpoint{0.407592in}{0.222904in}}{\pgfqpoint{1.411163in}{1.411163in}}%
\pgfusepath{clip}%
\pgfsetbuttcap%
\pgfsetroundjoin%
\definecolor{currentfill}{rgb}{0.168627,0.611765,0.411765}%
\pgfsetfillcolor{currentfill}%
\pgfsetfillopacity{0.180000}%
\pgfsetlinewidth{0.000000pt}%
\definecolor{currentstroke}{rgb}{0.000000,0.000000,0.000000}%
\pgfsetstrokecolor{currentstroke}%
\pgfsetdash{}{0pt}%
\pgfpathmoveto{\pgfqpoint{1.359740in}{0.958383in}}%
\pgfpathlineto{\pgfqpoint{1.375233in}{0.957951in}}%
\pgfpathlineto{\pgfqpoint{1.375542in}{0.980824in}}%
\pgfpathlineto{\pgfqpoint{1.360030in}{0.981342in}}%
\pgfpathlineto{\pgfqpoint{1.359740in}{0.958383in}}%
\pgfpathclose%
\pgfusepath{fill}%
\end{pgfscope}%
\begin{pgfscope}%
\pgfpathrectangle{\pgfqpoint{0.407592in}{0.222904in}}{\pgfqpoint{1.411163in}{1.411163in}}%
\pgfusepath{clip}%
\pgfsetbuttcap%
\pgfsetroundjoin%
\definecolor{currentfill}{rgb}{0.168627,0.611765,0.411765}%
\pgfsetfillcolor{currentfill}%
\pgfsetfillopacity{0.180000}%
\pgfsetlinewidth{0.000000pt}%
\definecolor{currentstroke}{rgb}{0.000000,0.000000,0.000000}%
\pgfsetstrokecolor{currentstroke}%
\pgfsetdash{}{0pt}%
\pgfpathmoveto{\pgfqpoint{1.376163in}{1.026745in}}%
\pgfpathlineto{\pgfqpoint{1.391615in}{1.027374in}}%
\pgfpathlineto{\pgfqpoint{1.391945in}{1.050338in}}%
\pgfpathlineto{\pgfqpoint{1.376475in}{1.049793in}}%
\pgfpathlineto{\pgfqpoint{1.376163in}{1.026745in}}%
\pgfpathclose%
\pgfusepath{fill}%
\end{pgfscope}%
\begin{pgfscope}%
\pgfpathrectangle{\pgfqpoint{0.407592in}{0.222904in}}{\pgfqpoint{1.411163in}{1.411163in}}%
\pgfusepath{clip}%
\pgfsetbuttcap%
\pgfsetroundjoin%
\definecolor{currentfill}{rgb}{0.168627,0.611765,0.411765}%
\pgfsetfillcolor{currentfill}%
\pgfsetfillopacity{0.180000}%
\pgfsetlinewidth{0.000000pt}%
\definecolor{currentstroke}{rgb}{0.000000,0.000000,0.000000}%
\pgfsetstrokecolor{currentstroke}%
\pgfsetdash{}{0pt}%
\pgfpathmoveto{\pgfqpoint{1.392608in}{1.096443in}}%
\pgfpathlineto{\pgfqpoint{1.408023in}{1.098223in}}%
\pgfpathlineto{\pgfqpoint{1.408374in}{1.121283in}}%
\pgfpathlineto{\pgfqpoint{1.392940in}{1.119583in}}%
\pgfpathlineto{\pgfqpoint{1.392608in}{1.096443in}}%
\pgfpathclose%
\pgfusepath{fill}%
\end{pgfscope}%
\begin{pgfscope}%
\pgfpathrectangle{\pgfqpoint{0.407592in}{0.222904in}}{\pgfqpoint{1.411163in}{1.411163in}}%
\pgfusepath{clip}%
\pgfsetbuttcap%
\pgfsetroundjoin%
\definecolor{currentfill}{rgb}{0.168627,0.611765,0.411765}%
\pgfsetfillcolor{currentfill}%
\pgfsetfillopacity{0.180000}%
\pgfsetlinewidth{0.000000pt}%
\definecolor{currentstroke}{rgb}{0.000000,0.000000,0.000000}%
\pgfsetstrokecolor{currentstroke}%
\pgfsetdash{}{0pt}%
\pgfpathmoveto{\pgfqpoint{1.343602in}{0.914126in}}%
\pgfpathlineto{\pgfqpoint{1.359161in}{0.912639in}}%
\pgfpathlineto{\pgfqpoint{1.359450in}{0.935482in}}%
\pgfpathlineto{\pgfqpoint{1.343873in}{0.937058in}}%
\pgfpathlineto{\pgfqpoint{1.343602in}{0.914126in}}%
\pgfpathclose%
\pgfusepath{fill}%
\end{pgfscope}%
\begin{pgfscope}%
\pgfpathrectangle{\pgfqpoint{0.407592in}{0.222904in}}{\pgfqpoint{1.411163in}{1.411163in}}%
\pgfusepath{clip}%
\pgfsetbuttcap%
\pgfsetroundjoin%
\definecolor{currentfill}{rgb}{0.168627,0.611765,0.411765}%
\pgfsetfillcolor{currentfill}%
\pgfsetfillopacity{0.180000}%
\pgfsetlinewidth{0.000000pt}%
\definecolor{currentstroke}{rgb}{0.000000,0.000000,0.000000}%
\pgfsetstrokecolor{currentstroke}%
\pgfsetdash{}{0pt}%
\pgfpathmoveto{\pgfqpoint{1.360030in}{0.981342in}}%
\pgfpathlineto{\pgfqpoint{1.375542in}{0.980824in}}%
\pgfpathlineto{\pgfqpoint{1.375852in}{1.003755in}}%
\pgfpathlineto{\pgfqpoint{1.360321in}{1.004361in}}%
\pgfpathlineto{\pgfqpoint{1.360030in}{0.981342in}}%
\pgfpathclose%
\pgfusepath{fill}%
\end{pgfscope}%
\begin{pgfscope}%
\pgfpathrectangle{\pgfqpoint{0.407592in}{0.222904in}}{\pgfqpoint{1.411163in}{1.411163in}}%
\pgfusepath{clip}%
\pgfsetbuttcap%
\pgfsetroundjoin%
\definecolor{currentfill}{rgb}{0.168627,0.611765,0.411765}%
\pgfsetfillcolor{currentfill}%
\pgfsetfillopacity{0.180000}%
\pgfsetlinewidth{0.000000pt}%
\definecolor{currentstroke}{rgb}{0.000000,0.000000,0.000000}%
\pgfsetstrokecolor{currentstroke}%
\pgfsetdash{}{0pt}%
\pgfpathmoveto{\pgfqpoint{1.376475in}{1.049793in}}%
\pgfpathlineto{\pgfqpoint{1.391945in}{1.050338in}}%
\pgfpathlineto{\pgfqpoint{1.392276in}{1.073361in}}%
\pgfpathlineto{\pgfqpoint{1.376787in}{1.072901in}}%
\pgfpathlineto{\pgfqpoint{1.376475in}{1.049793in}}%
\pgfpathclose%
\pgfusepath{fill}%
\end{pgfscope}%
\begin{pgfscope}%
\pgfpathrectangle{\pgfqpoint{0.407592in}{0.222904in}}{\pgfqpoint{1.411163in}{1.411163in}}%
\pgfusepath{clip}%
\pgfsetbuttcap%
\pgfsetroundjoin%
\definecolor{currentfill}{rgb}{0.168627,0.611765,0.411765}%
\pgfsetfillcolor{currentfill}%
\pgfsetfillopacity{0.180000}%
\pgfsetlinewidth{0.000000pt}%
\definecolor{currentstroke}{rgb}{0.000000,0.000000,0.000000}%
\pgfsetstrokecolor{currentstroke}%
\pgfsetdash{}{0pt}%
\pgfpathmoveto{\pgfqpoint{1.392940in}{1.119583in}}%
\pgfpathlineto{\pgfqpoint{1.408374in}{1.121283in}}%
\pgfpathlineto{\pgfqpoint{1.408726in}{1.144401in}}%
\pgfpathlineto{\pgfqpoint{1.393273in}{1.142783in}}%
\pgfpathlineto{\pgfqpoint{1.392940in}{1.119583in}}%
\pgfpathclose%
\pgfusepath{fill}%
\end{pgfscope}%
\begin{pgfscope}%
\pgfpathrectangle{\pgfqpoint{0.407592in}{0.222904in}}{\pgfqpoint{1.411163in}{1.411163in}}%
\pgfusepath{clip}%
\pgfsetbuttcap%
\pgfsetroundjoin%
\definecolor{currentfill}{rgb}{0.168627,0.611765,0.411765}%
\pgfsetfillcolor{currentfill}%
\pgfsetfillopacity{0.180000}%
\pgfsetlinewidth{0.000000pt}%
\definecolor{currentstroke}{rgb}{0.000000,0.000000,0.000000}%
\pgfsetstrokecolor{currentstroke}%
\pgfsetdash{}{0pt}%
\pgfpathmoveto{\pgfqpoint{1.343873in}{0.937058in}}%
\pgfpathlineto{\pgfqpoint{1.359450in}{0.935482in}}%
\pgfpathlineto{\pgfqpoint{1.359740in}{0.958383in}}%
\pgfpathlineto{\pgfqpoint{1.344144in}{0.960050in}}%
\pgfpathlineto{\pgfqpoint{1.343873in}{0.937058in}}%
\pgfpathclose%
\pgfusepath{fill}%
\end{pgfscope}%
\begin{pgfscope}%
\pgfpathrectangle{\pgfqpoint{0.407592in}{0.222904in}}{\pgfqpoint{1.411163in}{1.411163in}}%
\pgfusepath{clip}%
\pgfsetbuttcap%
\pgfsetroundjoin%
\definecolor{currentfill}{rgb}{0.168627,0.611765,0.411765}%
\pgfsetfillcolor{currentfill}%
\pgfsetfillopacity{0.180000}%
\pgfsetlinewidth{0.000000pt}%
\definecolor{currentstroke}{rgb}{0.000000,0.000000,0.000000}%
\pgfsetstrokecolor{currentstroke}%
\pgfsetdash{}{0pt}%
\pgfpathmoveto{\pgfqpoint{1.360321in}{1.004361in}}%
\pgfpathlineto{\pgfqpoint{1.375852in}{1.003755in}}%
\pgfpathlineto{\pgfqpoint{1.376163in}{1.026745in}}%
\pgfpathlineto{\pgfqpoint{1.360613in}{1.027438in}}%
\pgfpathlineto{\pgfqpoint{1.360321in}{1.004361in}}%
\pgfpathclose%
\pgfusepath{fill}%
\end{pgfscope}%
\begin{pgfscope}%
\pgfpathrectangle{\pgfqpoint{0.407592in}{0.222904in}}{\pgfqpoint{1.411163in}{1.411163in}}%
\pgfusepath{clip}%
\pgfsetbuttcap%
\pgfsetroundjoin%
\definecolor{currentfill}{rgb}{0.168627,0.611765,0.411765}%
\pgfsetfillcolor{currentfill}%
\pgfsetfillopacity{0.180000}%
\pgfsetlinewidth{0.000000pt}%
\definecolor{currentstroke}{rgb}{0.000000,0.000000,0.000000}%
\pgfsetstrokecolor{currentstroke}%
\pgfsetdash{}{0pt}%
\pgfpathmoveto{\pgfqpoint{1.327436in}{0.870889in}}%
\pgfpathlineto{\pgfqpoint{1.343064in}{0.868436in}}%
\pgfpathlineto{\pgfqpoint{1.343333in}{0.891252in}}%
\pgfpathlineto{\pgfqpoint{1.327686in}{0.893797in}}%
\pgfpathlineto{\pgfqpoint{1.327436in}{0.870889in}}%
\pgfpathclose%
\pgfusepath{fill}%
\end{pgfscope}%
\begin{pgfscope}%
\pgfpathrectangle{\pgfqpoint{0.407592in}{0.222904in}}{\pgfqpoint{1.411163in}{1.411163in}}%
\pgfusepath{clip}%
\pgfsetbuttcap%
\pgfsetroundjoin%
\definecolor{currentfill}{rgb}{0.168627,0.611765,0.411765}%
\pgfsetfillcolor{currentfill}%
\pgfsetfillopacity{0.180000}%
\pgfsetlinewidth{0.000000pt}%
\definecolor{currentstroke}{rgb}{0.000000,0.000000,0.000000}%
\pgfsetstrokecolor{currentstroke}%
\pgfsetdash{}{0pt}%
\pgfpathmoveto{\pgfqpoint{1.376787in}{1.072901in}}%
\pgfpathlineto{\pgfqpoint{1.392276in}{1.073361in}}%
\pgfpathlineto{\pgfqpoint{1.392608in}{1.096443in}}%
\pgfpathlineto{\pgfqpoint{1.377100in}{1.096067in}}%
\pgfpathlineto{\pgfqpoint{1.376787in}{1.072901in}}%
\pgfpathclose%
\pgfusepath{fill}%
\end{pgfscope}%
\begin{pgfscope}%
\pgfpathrectangle{\pgfqpoint{0.407592in}{0.222904in}}{\pgfqpoint{1.411163in}{1.411163in}}%
\pgfusepath{clip}%
\pgfsetbuttcap%
\pgfsetroundjoin%
\definecolor{currentfill}{rgb}{0.168627,0.611765,0.411765}%
\pgfsetfillcolor{currentfill}%
\pgfsetfillopacity{0.180000}%
\pgfsetlinewidth{0.000000pt}%
\definecolor{currentstroke}{rgb}{0.000000,0.000000,0.000000}%
\pgfsetstrokecolor{currentstroke}%
\pgfsetdash{}{0pt}%
\pgfpathmoveto{\pgfqpoint{1.393273in}{1.142783in}}%
\pgfpathlineto{\pgfqpoint{1.408726in}{1.144401in}}%
\pgfpathlineto{\pgfqpoint{1.409078in}{1.167579in}}%
\pgfpathlineto{\pgfqpoint{1.393608in}{1.166043in}}%
\pgfpathlineto{\pgfqpoint{1.393273in}{1.142783in}}%
\pgfpathclose%
\pgfusepath{fill}%
\end{pgfscope}%
\begin{pgfscope}%
\pgfpathrectangle{\pgfqpoint{0.407592in}{0.222904in}}{\pgfqpoint{1.411163in}{1.411163in}}%
\pgfusepath{clip}%
\pgfsetbuttcap%
\pgfsetroundjoin%
\definecolor{currentfill}{rgb}{0.168627,0.611765,0.411765}%
\pgfsetfillcolor{currentfill}%
\pgfsetfillopacity{0.180000}%
\pgfsetlinewidth{0.000000pt}%
\definecolor{currentstroke}{rgb}{0.000000,0.000000,0.000000}%
\pgfsetstrokecolor{currentstroke}%
\pgfsetdash{}{0pt}%
\pgfpathmoveto{\pgfqpoint{1.344144in}{0.960050in}}%
\pgfpathlineto{\pgfqpoint{1.359740in}{0.958383in}}%
\pgfpathlineto{\pgfqpoint{1.360030in}{0.981342in}}%
\pgfpathlineto{\pgfqpoint{1.344415in}{0.983100in}}%
\pgfpathlineto{\pgfqpoint{1.344144in}{0.960050in}}%
\pgfpathclose%
\pgfusepath{fill}%
\end{pgfscope}%
\begin{pgfscope}%
\pgfpathrectangle{\pgfqpoint{0.407592in}{0.222904in}}{\pgfqpoint{1.411163in}{1.411163in}}%
\pgfusepath{clip}%
\pgfsetbuttcap%
\pgfsetroundjoin%
\definecolor{currentfill}{rgb}{0.168627,0.611765,0.411765}%
\pgfsetfillcolor{currentfill}%
\pgfsetfillopacity{0.180000}%
\pgfsetlinewidth{0.000000pt}%
\definecolor{currentstroke}{rgb}{0.000000,0.000000,0.000000}%
\pgfsetstrokecolor{currentstroke}%
\pgfsetdash{}{0pt}%
\pgfpathmoveto{\pgfqpoint{1.360613in}{1.027438in}}%
\pgfpathlineto{\pgfqpoint{1.376163in}{1.026745in}}%
\pgfpathlineto{\pgfqpoint{1.376475in}{1.049793in}}%
\pgfpathlineto{\pgfqpoint{1.360906in}{1.050575in}}%
\pgfpathlineto{\pgfqpoint{1.360613in}{1.027438in}}%
\pgfpathclose%
\pgfusepath{fill}%
\end{pgfscope}%
\begin{pgfscope}%
\pgfpathrectangle{\pgfqpoint{0.407592in}{0.222904in}}{\pgfqpoint{1.411163in}{1.411163in}}%
\pgfusepath{clip}%
\pgfsetbuttcap%
\pgfsetroundjoin%
\definecolor{currentfill}{rgb}{0.168627,0.611765,0.411765}%
\pgfsetfillcolor{currentfill}%
\pgfsetfillopacity{0.180000}%
\pgfsetlinewidth{0.000000pt}%
\definecolor{currentstroke}{rgb}{0.000000,0.000000,0.000000}%
\pgfsetstrokecolor{currentstroke}%
\pgfsetdash{}{0pt}%
\pgfpathmoveto{\pgfqpoint{1.327686in}{0.893797in}}%
\pgfpathlineto{\pgfqpoint{1.343333in}{0.891252in}}%
\pgfpathlineto{\pgfqpoint{1.343602in}{0.914126in}}%
\pgfpathlineto{\pgfqpoint{1.327936in}{0.916764in}}%
\pgfpathlineto{\pgfqpoint{1.327686in}{0.893797in}}%
\pgfpathclose%
\pgfusepath{fill}%
\end{pgfscope}%
\begin{pgfscope}%
\pgfpathrectangle{\pgfqpoint{0.407592in}{0.222904in}}{\pgfqpoint{1.411163in}{1.411163in}}%
\pgfusepath{clip}%
\pgfsetbuttcap%
\pgfsetroundjoin%
\definecolor{currentfill}{rgb}{0.168627,0.611765,0.411765}%
\pgfsetfillcolor{currentfill}%
\pgfsetfillopacity{0.180000}%
\pgfsetlinewidth{0.000000pt}%
\definecolor{currentstroke}{rgb}{0.000000,0.000000,0.000000}%
\pgfsetstrokecolor{currentstroke}%
\pgfsetdash{}{0pt}%
\pgfpathmoveto{\pgfqpoint{1.377100in}{1.096067in}}%
\pgfpathlineto{\pgfqpoint{1.392608in}{1.096443in}}%
\pgfpathlineto{\pgfqpoint{1.392940in}{1.119583in}}%
\pgfpathlineto{\pgfqpoint{1.377414in}{1.119293in}}%
\pgfpathlineto{\pgfqpoint{1.377100in}{1.096067in}}%
\pgfpathclose%
\pgfusepath{fill}%
\end{pgfscope}%
\begin{pgfscope}%
\pgfpathrectangle{\pgfqpoint{0.407592in}{0.222904in}}{\pgfqpoint{1.411163in}{1.411163in}}%
\pgfusepath{clip}%
\pgfsetbuttcap%
\pgfsetroundjoin%
\definecolor{currentfill}{rgb}{0.168627,0.611765,0.411765}%
\pgfsetfillcolor{currentfill}%
\pgfsetfillopacity{0.180000}%
\pgfsetlinewidth{0.000000pt}%
\definecolor{currentstroke}{rgb}{0.000000,0.000000,0.000000}%
\pgfsetstrokecolor{currentstroke}%
\pgfsetdash{}{0pt}%
\pgfpathmoveto{\pgfqpoint{1.344415in}{0.983100in}}%
\pgfpathlineto{\pgfqpoint{1.360030in}{0.981342in}}%
\pgfpathlineto{\pgfqpoint{1.360321in}{1.004361in}}%
\pgfpathlineto{\pgfqpoint{1.344687in}{1.006209in}}%
\pgfpathlineto{\pgfqpoint{1.344415in}{0.983100in}}%
\pgfpathclose%
\pgfusepath{fill}%
\end{pgfscope}%
\begin{pgfscope}%
\pgfpathrectangle{\pgfqpoint{0.407592in}{0.222904in}}{\pgfqpoint{1.411163in}{1.411163in}}%
\pgfusepath{clip}%
\pgfsetbuttcap%
\pgfsetroundjoin%
\definecolor{currentfill}{rgb}{0.168627,0.611765,0.411765}%
\pgfsetfillcolor{currentfill}%
\pgfsetfillopacity{0.180000}%
\pgfsetlinewidth{0.000000pt}%
\definecolor{currentstroke}{rgb}{0.000000,0.000000,0.000000}%
\pgfsetstrokecolor{currentstroke}%
\pgfsetdash{}{0pt}%
\pgfpathmoveto{\pgfqpoint{1.360906in}{1.050575in}}%
\pgfpathlineto{\pgfqpoint{1.376475in}{1.049793in}}%
\pgfpathlineto{\pgfqpoint{1.376787in}{1.072901in}}%
\pgfpathlineto{\pgfqpoint{1.361200in}{1.073771in}}%
\pgfpathlineto{\pgfqpoint{1.360906in}{1.050575in}}%
\pgfpathclose%
\pgfusepath{fill}%
\end{pgfscope}%
\begin{pgfscope}%
\pgfpathrectangle{\pgfqpoint{0.407592in}{0.222904in}}{\pgfqpoint{1.411163in}{1.411163in}}%
\pgfusepath{clip}%
\pgfsetbuttcap%
\pgfsetroundjoin%
\definecolor{currentfill}{rgb}{0.168627,0.611765,0.411765}%
\pgfsetfillcolor{currentfill}%
\pgfsetfillopacity{0.180000}%
\pgfsetlinewidth{0.000000pt}%
\definecolor{currentstroke}{rgb}{0.000000,0.000000,0.000000}%
\pgfsetstrokecolor{currentstroke}%
\pgfsetdash{}{0pt}%
\pgfpathmoveto{\pgfqpoint{1.327936in}{0.916764in}}%
\pgfpathlineto{\pgfqpoint{1.343602in}{0.914126in}}%
\pgfpathlineto{\pgfqpoint{1.343873in}{0.937058in}}%
\pgfpathlineto{\pgfqpoint{1.328188in}{0.939790in}}%
\pgfpathlineto{\pgfqpoint{1.327936in}{0.916764in}}%
\pgfpathclose%
\pgfusepath{fill}%
\end{pgfscope}%
\begin{pgfscope}%
\pgfpathrectangle{\pgfqpoint{0.407592in}{0.222904in}}{\pgfqpoint{1.411163in}{1.411163in}}%
\pgfusepath{clip}%
\pgfsetbuttcap%
\pgfsetroundjoin%
\definecolor{currentfill}{rgb}{0.168627,0.611765,0.411765}%
\pgfsetfillcolor{currentfill}%
\pgfsetfillopacity{0.180000}%
\pgfsetlinewidth{0.000000pt}%
\definecolor{currentstroke}{rgb}{0.000000,0.000000,0.000000}%
\pgfsetstrokecolor{currentstroke}%
\pgfsetdash{}{0pt}%
\pgfpathmoveto{\pgfqpoint{1.377414in}{1.119293in}}%
\pgfpathlineto{\pgfqpoint{1.392940in}{1.119583in}}%
\pgfpathlineto{\pgfqpoint{1.393273in}{1.142783in}}%
\pgfpathlineto{\pgfqpoint{1.377728in}{1.142579in}}%
\pgfpathlineto{\pgfqpoint{1.377414in}{1.119293in}}%
\pgfpathclose%
\pgfusepath{fill}%
\end{pgfscope}%
\begin{pgfscope}%
\pgfpathrectangle{\pgfqpoint{0.407592in}{0.222904in}}{\pgfqpoint{1.411163in}{1.411163in}}%
\pgfusepath{clip}%
\pgfsetbuttcap%
\pgfsetroundjoin%
\definecolor{currentfill}{rgb}{0.168627,0.611765,0.411765}%
\pgfsetfillcolor{currentfill}%
\pgfsetfillopacity{0.180000}%
\pgfsetlinewidth{0.000000pt}%
\definecolor{currentstroke}{rgb}{0.000000,0.000000,0.000000}%
\pgfsetstrokecolor{currentstroke}%
\pgfsetdash{}{0pt}%
\pgfpathmoveto{\pgfqpoint{1.344687in}{1.006209in}}%
\pgfpathlineto{\pgfqpoint{1.360321in}{1.004361in}}%
\pgfpathlineto{\pgfqpoint{1.360613in}{1.027438in}}%
\pgfpathlineto{\pgfqpoint{1.344960in}{1.029378in}}%
\pgfpathlineto{\pgfqpoint{1.344687in}{1.006209in}}%
\pgfpathclose%
\pgfusepath{fill}%
\end{pgfscope}%
\begin{pgfscope}%
\pgfpathrectangle{\pgfqpoint{0.407592in}{0.222904in}}{\pgfqpoint{1.411163in}{1.411163in}}%
\pgfusepath{clip}%
\pgfsetbuttcap%
\pgfsetroundjoin%
\definecolor{currentfill}{rgb}{0.168627,0.611765,0.411765}%
\pgfsetfillcolor{currentfill}%
\pgfsetfillopacity{0.180000}%
\pgfsetlinewidth{0.000000pt}%
\definecolor{currentstroke}{rgb}{0.000000,0.000000,0.000000}%
\pgfsetstrokecolor{currentstroke}%
\pgfsetdash{}{0pt}%
\pgfpathmoveto{\pgfqpoint{1.361200in}{1.073771in}}%
\pgfpathlineto{\pgfqpoint{1.376787in}{1.072901in}}%
\pgfpathlineto{\pgfqpoint{1.377100in}{1.096067in}}%
\pgfpathlineto{\pgfqpoint{1.361494in}{1.097026in}}%
\pgfpathlineto{\pgfqpoint{1.361200in}{1.073771in}}%
\pgfpathclose%
\pgfusepath{fill}%
\end{pgfscope}%
\begin{pgfscope}%
\pgfpathrectangle{\pgfqpoint{0.407592in}{0.222904in}}{\pgfqpoint{1.411163in}{1.411163in}}%
\pgfusepath{clip}%
\pgfsetbuttcap%
\pgfsetroundjoin%
\definecolor{currentfill}{rgb}{0.168627,0.611765,0.411765}%
\pgfsetfillcolor{currentfill}%
\pgfsetfillopacity{0.180000}%
\pgfsetlinewidth{0.000000pt}%
\definecolor{currentstroke}{rgb}{0.000000,0.000000,0.000000}%
\pgfsetstrokecolor{currentstroke}%
\pgfsetdash{}{0pt}%
\pgfpathmoveto{\pgfqpoint{1.328188in}{0.939790in}}%
\pgfpathlineto{\pgfqpoint{1.343873in}{0.937058in}}%
\pgfpathlineto{\pgfqpoint{1.344144in}{0.960050in}}%
\pgfpathlineto{\pgfqpoint{1.328439in}{0.962875in}}%
\pgfpathlineto{\pgfqpoint{1.328188in}{0.939790in}}%
\pgfpathclose%
\pgfusepath{fill}%
\end{pgfscope}%
\begin{pgfscope}%
\pgfpathrectangle{\pgfqpoint{0.407592in}{0.222904in}}{\pgfqpoint{1.411163in}{1.411163in}}%
\pgfusepath{clip}%
\pgfsetbuttcap%
\pgfsetroundjoin%
\definecolor{currentfill}{rgb}{0.168627,0.611765,0.411765}%
\pgfsetfillcolor{currentfill}%
\pgfsetfillopacity{0.180000}%
\pgfsetlinewidth{0.000000pt}%
\definecolor{currentstroke}{rgb}{0.000000,0.000000,0.000000}%
\pgfsetstrokecolor{currentstroke}%
\pgfsetdash{}{0pt}%
\pgfpathmoveto{\pgfqpoint{1.377728in}{1.142579in}}%
\pgfpathlineto{\pgfqpoint{1.393273in}{1.142783in}}%
\pgfpathlineto{\pgfqpoint{1.393608in}{1.166043in}}%
\pgfpathlineto{\pgfqpoint{1.378044in}{1.165924in}}%
\pgfpathlineto{\pgfqpoint{1.377728in}{1.142579in}}%
\pgfpathclose%
\pgfusepath{fill}%
\end{pgfscope}%
\begin{pgfscope}%
\pgfpathrectangle{\pgfqpoint{0.407592in}{0.222904in}}{\pgfqpoint{1.411163in}{1.411163in}}%
\pgfusepath{clip}%
\pgfsetbuttcap%
\pgfsetroundjoin%
\definecolor{currentfill}{rgb}{0.168627,0.611765,0.411765}%
\pgfsetfillcolor{currentfill}%
\pgfsetfillopacity{0.180000}%
\pgfsetlinewidth{0.000000pt}%
\definecolor{currentstroke}{rgb}{0.000000,0.000000,0.000000}%
\pgfsetstrokecolor{currentstroke}%
\pgfsetdash{}{0pt}%
\pgfpathmoveto{\pgfqpoint{1.311696in}{0.874411in}}%
\pgfpathlineto{\pgfqpoint{1.327436in}{0.870889in}}%
\pgfpathlineto{\pgfqpoint{1.327686in}{0.893797in}}%
\pgfpathlineto{\pgfqpoint{1.311927in}{0.897415in}}%
\pgfpathlineto{\pgfqpoint{1.311696in}{0.874411in}}%
\pgfpathclose%
\pgfusepath{fill}%
\end{pgfscope}%
\begin{pgfscope}%
\pgfpathrectangle{\pgfqpoint{0.407592in}{0.222904in}}{\pgfqpoint{1.411163in}{1.411163in}}%
\pgfusepath{clip}%
\pgfsetbuttcap%
\pgfsetroundjoin%
\definecolor{currentfill}{rgb}{0.168627,0.611765,0.411765}%
\pgfsetfillcolor{currentfill}%
\pgfsetfillopacity{0.180000}%
\pgfsetlinewidth{0.000000pt}%
\definecolor{currentstroke}{rgb}{0.000000,0.000000,0.000000}%
\pgfsetstrokecolor{currentstroke}%
\pgfsetdash{}{0pt}%
\pgfpathmoveto{\pgfqpoint{1.344960in}{1.029378in}}%
\pgfpathlineto{\pgfqpoint{1.360613in}{1.027438in}}%
\pgfpathlineto{\pgfqpoint{1.360906in}{1.050575in}}%
\pgfpathlineto{\pgfqpoint{1.345234in}{1.052607in}}%
\pgfpathlineto{\pgfqpoint{1.344960in}{1.029378in}}%
\pgfpathclose%
\pgfusepath{fill}%
\end{pgfscope}%
\begin{pgfscope}%
\pgfpathrectangle{\pgfqpoint{0.407592in}{0.222904in}}{\pgfqpoint{1.411163in}{1.411163in}}%
\pgfusepath{clip}%
\pgfsetbuttcap%
\pgfsetroundjoin%
\definecolor{currentfill}{rgb}{0.168627,0.611765,0.411765}%
\pgfsetfillcolor{currentfill}%
\pgfsetfillopacity{0.180000}%
\pgfsetlinewidth{0.000000pt}%
\definecolor{currentstroke}{rgb}{0.000000,0.000000,0.000000}%
\pgfsetstrokecolor{currentstroke}%
\pgfsetdash{}{0pt}%
\pgfpathmoveto{\pgfqpoint{1.361494in}{1.097026in}}%
\pgfpathlineto{\pgfqpoint{1.377100in}{1.096067in}}%
\pgfpathlineto{\pgfqpoint{1.377414in}{1.119293in}}%
\pgfpathlineto{\pgfqpoint{1.361789in}{1.120342in}}%
\pgfpathlineto{\pgfqpoint{1.361494in}{1.097026in}}%
\pgfpathclose%
\pgfusepath{fill}%
\end{pgfscope}%
\begin{pgfscope}%
\pgfpathrectangle{\pgfqpoint{0.407592in}{0.222904in}}{\pgfqpoint{1.411163in}{1.411163in}}%
\pgfusepath{clip}%
\pgfsetbuttcap%
\pgfsetroundjoin%
\definecolor{currentfill}{rgb}{0.168627,0.611765,0.411765}%
\pgfsetfillcolor{currentfill}%
\pgfsetfillopacity{0.180000}%
\pgfsetlinewidth{0.000000pt}%
\definecolor{currentstroke}{rgb}{0.000000,0.000000,0.000000}%
\pgfsetstrokecolor{currentstroke}%
\pgfsetdash{}{0pt}%
\pgfpathmoveto{\pgfqpoint{1.328439in}{0.962875in}}%
\pgfpathlineto{\pgfqpoint{1.344144in}{0.960050in}}%
\pgfpathlineto{\pgfqpoint{1.344415in}{0.983100in}}%
\pgfpathlineto{\pgfqpoint{1.328692in}{0.986020in}}%
\pgfpathlineto{\pgfqpoint{1.328439in}{0.962875in}}%
\pgfpathclose%
\pgfusepath{fill}%
\end{pgfscope}%
\begin{pgfscope}%
\pgfpathrectangle{\pgfqpoint{0.407592in}{0.222904in}}{\pgfqpoint{1.411163in}{1.411163in}}%
\pgfusepath{clip}%
\pgfsetbuttcap%
\pgfsetroundjoin%
\definecolor{currentfill}{rgb}{0.168627,0.611765,0.411765}%
\pgfsetfillcolor{currentfill}%
\pgfsetfillopacity{0.180000}%
\pgfsetlinewidth{0.000000pt}%
\definecolor{currentstroke}{rgb}{0.000000,0.000000,0.000000}%
\pgfsetstrokecolor{currentstroke}%
\pgfsetdash{}{0pt}%
\pgfpathmoveto{\pgfqpoint{1.311927in}{0.897415in}}%
\pgfpathlineto{\pgfqpoint{1.327686in}{0.893797in}}%
\pgfpathlineto{\pgfqpoint{1.327936in}{0.916764in}}%
\pgfpathlineto{\pgfqpoint{1.312158in}{0.920478in}}%
\pgfpathlineto{\pgfqpoint{1.311927in}{0.897415in}}%
\pgfpathclose%
\pgfusepath{fill}%
\end{pgfscope}%
\begin{pgfscope}%
\pgfpathrectangle{\pgfqpoint{0.407592in}{0.222904in}}{\pgfqpoint{1.411163in}{1.411163in}}%
\pgfusepath{clip}%
\pgfsetbuttcap%
\pgfsetroundjoin%
\definecolor{currentfill}{rgb}{0.168627,0.611765,0.411765}%
\pgfsetfillcolor{currentfill}%
\pgfsetfillopacity{0.180000}%
\pgfsetlinewidth{0.000000pt}%
\definecolor{currentstroke}{rgb}{0.000000,0.000000,0.000000}%
\pgfsetstrokecolor{currentstroke}%
\pgfsetdash{}{0pt}%
\pgfpathmoveto{\pgfqpoint{1.345234in}{1.052607in}}%
\pgfpathlineto{\pgfqpoint{1.360906in}{1.050575in}}%
\pgfpathlineto{\pgfqpoint{1.361200in}{1.073771in}}%
\pgfpathlineto{\pgfqpoint{1.345509in}{1.075895in}}%
\pgfpathlineto{\pgfqpoint{1.345234in}{1.052607in}}%
\pgfpathclose%
\pgfusepath{fill}%
\end{pgfscope}%
\begin{pgfscope}%
\pgfpathrectangle{\pgfqpoint{0.407592in}{0.222904in}}{\pgfqpoint{1.411163in}{1.411163in}}%
\pgfusepath{clip}%
\pgfsetbuttcap%
\pgfsetroundjoin%
\definecolor{currentfill}{rgb}{0.168627,0.611765,0.411765}%
\pgfsetfillcolor{currentfill}%
\pgfsetfillopacity{0.180000}%
\pgfsetlinewidth{0.000000pt}%
\definecolor{currentstroke}{rgb}{0.000000,0.000000,0.000000}%
\pgfsetstrokecolor{currentstroke}%
\pgfsetdash{}{0pt}%
\pgfpathmoveto{\pgfqpoint{1.361789in}{1.120342in}}%
\pgfpathlineto{\pgfqpoint{1.377414in}{1.119293in}}%
\pgfpathlineto{\pgfqpoint{1.377728in}{1.142579in}}%
\pgfpathlineto{\pgfqpoint{1.362085in}{1.143717in}}%
\pgfpathlineto{\pgfqpoint{1.361789in}{1.120342in}}%
\pgfpathclose%
\pgfusepath{fill}%
\end{pgfscope}%
\begin{pgfscope}%
\pgfpathrectangle{\pgfqpoint{0.407592in}{0.222904in}}{\pgfqpoint{1.411163in}{1.411163in}}%
\pgfusepath{clip}%
\pgfsetbuttcap%
\pgfsetroundjoin%
\definecolor{currentfill}{rgb}{0.168627,0.611765,0.411765}%
\pgfsetfillcolor{currentfill}%
\pgfsetfillopacity{0.180000}%
\pgfsetlinewidth{0.000000pt}%
\definecolor{currentstroke}{rgb}{0.000000,0.000000,0.000000}%
\pgfsetstrokecolor{currentstroke}%
\pgfsetdash{}{0pt}%
\pgfpathmoveto{\pgfqpoint{1.328692in}{0.986020in}}%
\pgfpathlineto{\pgfqpoint{1.344415in}{0.983100in}}%
\pgfpathlineto{\pgfqpoint{1.344687in}{1.006209in}}%
\pgfpathlineto{\pgfqpoint{1.328945in}{1.009224in}}%
\pgfpathlineto{\pgfqpoint{1.328692in}{0.986020in}}%
\pgfpathclose%
\pgfusepath{fill}%
\end{pgfscope}%
\begin{pgfscope}%
\pgfpathrectangle{\pgfqpoint{0.407592in}{0.222904in}}{\pgfqpoint{1.411163in}{1.411163in}}%
\pgfusepath{clip}%
\pgfsetbuttcap%
\pgfsetroundjoin%
\definecolor{currentfill}{rgb}{0.168627,0.611765,0.411765}%
\pgfsetfillcolor{currentfill}%
\pgfsetfillopacity{0.180000}%
\pgfsetlinewidth{0.000000pt}%
\definecolor{currentstroke}{rgb}{0.000000,0.000000,0.000000}%
\pgfsetstrokecolor{currentstroke}%
\pgfsetdash{}{0pt}%
\pgfpathmoveto{\pgfqpoint{1.312158in}{0.920478in}}%
\pgfpathlineto{\pgfqpoint{1.327936in}{0.916764in}}%
\pgfpathlineto{\pgfqpoint{1.328188in}{0.939790in}}%
\pgfpathlineto{\pgfqpoint{1.312390in}{0.943601in}}%
\pgfpathlineto{\pgfqpoint{1.312158in}{0.920478in}}%
\pgfpathclose%
\pgfusepath{fill}%
\end{pgfscope}%
\begin{pgfscope}%
\pgfpathrectangle{\pgfqpoint{0.407592in}{0.222904in}}{\pgfqpoint{1.411163in}{1.411163in}}%
\pgfusepath{clip}%
\pgfsetbuttcap%
\pgfsetroundjoin%
\definecolor{currentfill}{rgb}{0.168627,0.611765,0.411765}%
\pgfsetfillcolor{currentfill}%
\pgfsetfillopacity{0.180000}%
\pgfsetlinewidth{0.000000pt}%
\definecolor{currentstroke}{rgb}{0.000000,0.000000,0.000000}%
\pgfsetstrokecolor{currentstroke}%
\pgfsetdash{}{0pt}%
\pgfpathmoveto{\pgfqpoint{1.345509in}{1.075895in}}%
\pgfpathlineto{\pgfqpoint{1.361200in}{1.073771in}}%
\pgfpathlineto{\pgfqpoint{1.361494in}{1.097026in}}%
\pgfpathlineto{\pgfqpoint{1.345784in}{1.099243in}}%
\pgfpathlineto{\pgfqpoint{1.345509in}{1.075895in}}%
\pgfpathclose%
\pgfusepath{fill}%
\end{pgfscope}%
\begin{pgfscope}%
\pgfpathrectangle{\pgfqpoint{0.407592in}{0.222904in}}{\pgfqpoint{1.411163in}{1.411163in}}%
\pgfusepath{clip}%
\pgfsetbuttcap%
\pgfsetroundjoin%
\definecolor{currentfill}{rgb}{0.168627,0.611765,0.411765}%
\pgfsetfillcolor{currentfill}%
\pgfsetfillopacity{0.180000}%
\pgfsetlinewidth{0.000000pt}%
\definecolor{currentstroke}{rgb}{0.000000,0.000000,0.000000}%
\pgfsetstrokecolor{currentstroke}%
\pgfsetdash{}{0pt}%
\pgfpathmoveto{\pgfqpoint{1.362085in}{1.143717in}}%
\pgfpathlineto{\pgfqpoint{1.377728in}{1.142579in}}%
\pgfpathlineto{\pgfqpoint{1.378044in}{1.165924in}}%
\pgfpathlineto{\pgfqpoint{1.362381in}{1.167153in}}%
\pgfpathlineto{\pgfqpoint{1.362085in}{1.143717in}}%
\pgfpathclose%
\pgfusepath{fill}%
\end{pgfscope}%
\begin{pgfscope}%
\pgfpathrectangle{\pgfqpoint{0.407592in}{0.222904in}}{\pgfqpoint{1.411163in}{1.411163in}}%
\pgfusepath{clip}%
\pgfsetbuttcap%
\pgfsetroundjoin%
\definecolor{currentfill}{rgb}{0.168627,0.611765,0.411765}%
\pgfsetfillcolor{currentfill}%
\pgfsetfillopacity{0.180000}%
\pgfsetlinewidth{0.000000pt}%
\definecolor{currentstroke}{rgb}{0.000000,0.000000,0.000000}%
\pgfsetstrokecolor{currentstroke}%
\pgfsetdash{}{0pt}%
\pgfpathmoveto{\pgfqpoint{1.328945in}{1.009224in}}%
\pgfpathlineto{\pgfqpoint{1.344687in}{1.006209in}}%
\pgfpathlineto{\pgfqpoint{1.344960in}{1.029378in}}%
\pgfpathlineto{\pgfqpoint{1.329199in}{1.032488in}}%
\pgfpathlineto{\pgfqpoint{1.328945in}{1.009224in}}%
\pgfpathclose%
\pgfusepath{fill}%
\end{pgfscope}%
\begin{pgfscope}%
\pgfpathrectangle{\pgfqpoint{0.407592in}{0.222904in}}{\pgfqpoint{1.411163in}{1.411163in}}%
\pgfusepath{clip}%
\pgfsetbuttcap%
\pgfsetroundjoin%
\definecolor{currentfill}{rgb}{0.168627,0.611765,0.411765}%
\pgfsetfillcolor{currentfill}%
\pgfsetfillopacity{0.180000}%
\pgfsetlinewidth{0.000000pt}%
\definecolor{currentstroke}{rgb}{0.000000,0.000000,0.000000}%
\pgfsetstrokecolor{currentstroke}%
\pgfsetdash{}{0pt}%
\pgfpathmoveto{\pgfqpoint{1.312390in}{0.943601in}}%
\pgfpathlineto{\pgfqpoint{1.328188in}{0.939790in}}%
\pgfpathlineto{\pgfqpoint{1.328439in}{0.962875in}}%
\pgfpathlineto{\pgfqpoint{1.312622in}{0.966783in}}%
\pgfpathlineto{\pgfqpoint{1.312390in}{0.943601in}}%
\pgfpathclose%
\pgfusepath{fill}%
\end{pgfscope}%
\begin{pgfscope}%
\pgfpathrectangle{\pgfqpoint{0.407592in}{0.222904in}}{\pgfqpoint{1.411163in}{1.411163in}}%
\pgfusepath{clip}%
\pgfsetbuttcap%
\pgfsetroundjoin%
\definecolor{currentfill}{rgb}{0.168627,0.611765,0.411765}%
\pgfsetfillcolor{currentfill}%
\pgfsetfillopacity{0.180000}%
\pgfsetlinewidth{0.000000pt}%
\definecolor{currentstroke}{rgb}{0.000000,0.000000,0.000000}%
\pgfsetstrokecolor{currentstroke}%
\pgfsetdash{}{0pt}%
\pgfpathmoveto{\pgfqpoint{1.295839in}{0.878928in}}%
\pgfpathlineto{\pgfqpoint{1.311696in}{0.874411in}}%
\pgfpathlineto{\pgfqpoint{1.311927in}{0.897415in}}%
\pgfpathlineto{\pgfqpoint{1.296051in}{0.902031in}}%
\pgfpathlineto{\pgfqpoint{1.295839in}{0.878928in}}%
\pgfpathclose%
\pgfusepath{fill}%
\end{pgfscope}%
\begin{pgfscope}%
\pgfpathrectangle{\pgfqpoint{0.407592in}{0.222904in}}{\pgfqpoint{1.411163in}{1.411163in}}%
\pgfusepath{clip}%
\pgfsetbuttcap%
\pgfsetroundjoin%
\definecolor{currentfill}{rgb}{0.168627,0.611765,0.411765}%
\pgfsetfillcolor{currentfill}%
\pgfsetfillopacity{0.180000}%
\pgfsetlinewidth{0.000000pt}%
\definecolor{currentstroke}{rgb}{0.000000,0.000000,0.000000}%
\pgfsetstrokecolor{currentstroke}%
\pgfsetdash{}{0pt}%
\pgfpathmoveto{\pgfqpoint{1.345784in}{1.099243in}}%
\pgfpathlineto{\pgfqpoint{1.361494in}{1.097026in}}%
\pgfpathlineto{\pgfqpoint{1.361789in}{1.120342in}}%
\pgfpathlineto{\pgfqpoint{1.346060in}{1.122652in}}%
\pgfpathlineto{\pgfqpoint{1.345784in}{1.099243in}}%
\pgfpathclose%
\pgfusepath{fill}%
\end{pgfscope}%
\begin{pgfscope}%
\pgfpathrectangle{\pgfqpoint{0.407592in}{0.222904in}}{\pgfqpoint{1.411163in}{1.411163in}}%
\pgfusepath{clip}%
\pgfsetbuttcap%
\pgfsetroundjoin%
\definecolor{currentfill}{rgb}{0.168627,0.611765,0.411765}%
\pgfsetfillcolor{currentfill}%
\pgfsetfillopacity{0.180000}%
\pgfsetlinewidth{0.000000pt}%
\definecolor{currentstroke}{rgb}{0.000000,0.000000,0.000000}%
\pgfsetstrokecolor{currentstroke}%
\pgfsetdash{}{0pt}%
\pgfpathmoveto{\pgfqpoint{1.329199in}{1.032488in}}%
\pgfpathlineto{\pgfqpoint{1.344960in}{1.029378in}}%
\pgfpathlineto{\pgfqpoint{1.345234in}{1.052607in}}%
\pgfpathlineto{\pgfqpoint{1.329453in}{1.055812in}}%
\pgfpathlineto{\pgfqpoint{1.329199in}{1.032488in}}%
\pgfpathclose%
\pgfusepath{fill}%
\end{pgfscope}%
\begin{pgfscope}%
\pgfpathrectangle{\pgfqpoint{0.407592in}{0.222904in}}{\pgfqpoint{1.411163in}{1.411163in}}%
\pgfusepath{clip}%
\pgfsetbuttcap%
\pgfsetroundjoin%
\definecolor{currentfill}{rgb}{0.168627,0.611765,0.411765}%
\pgfsetfillcolor{currentfill}%
\pgfsetfillopacity{0.180000}%
\pgfsetlinewidth{0.000000pt}%
\definecolor{currentstroke}{rgb}{0.000000,0.000000,0.000000}%
\pgfsetstrokecolor{currentstroke}%
\pgfsetdash{}{0pt}%
\pgfpathmoveto{\pgfqpoint{1.312622in}{0.966783in}}%
\pgfpathlineto{\pgfqpoint{1.328439in}{0.962875in}}%
\pgfpathlineto{\pgfqpoint{1.328692in}{0.986020in}}%
\pgfpathlineto{\pgfqpoint{1.312855in}{0.990025in}}%
\pgfpathlineto{\pgfqpoint{1.312622in}{0.966783in}}%
\pgfpathclose%
\pgfusepath{fill}%
\end{pgfscope}%
\begin{pgfscope}%
\pgfpathrectangle{\pgfqpoint{0.407592in}{0.222904in}}{\pgfqpoint{1.411163in}{1.411163in}}%
\pgfusepath{clip}%
\pgfsetbuttcap%
\pgfsetroundjoin%
\definecolor{currentfill}{rgb}{0.168627,0.611765,0.411765}%
\pgfsetfillcolor{currentfill}%
\pgfsetfillopacity{0.180000}%
\pgfsetlinewidth{0.000000pt}%
\definecolor{currentstroke}{rgb}{0.000000,0.000000,0.000000}%
\pgfsetstrokecolor{currentstroke}%
\pgfsetdash{}{0pt}%
\pgfpathmoveto{\pgfqpoint{1.296051in}{0.902031in}}%
\pgfpathlineto{\pgfqpoint{1.311927in}{0.897415in}}%
\pgfpathlineto{\pgfqpoint{1.312158in}{0.920478in}}%
\pgfpathlineto{\pgfqpoint{1.296262in}{0.925194in}}%
\pgfpathlineto{\pgfqpoint{1.296051in}{0.902031in}}%
\pgfpathclose%
\pgfusepath{fill}%
\end{pgfscope}%
\begin{pgfscope}%
\pgfpathrectangle{\pgfqpoint{0.407592in}{0.222904in}}{\pgfqpoint{1.411163in}{1.411163in}}%
\pgfusepath{clip}%
\pgfsetbuttcap%
\pgfsetroundjoin%
\definecolor{currentfill}{rgb}{0.168627,0.611765,0.411765}%
\pgfsetfillcolor{currentfill}%
\pgfsetfillopacity{0.180000}%
\pgfsetlinewidth{0.000000pt}%
\definecolor{currentstroke}{rgb}{0.000000,0.000000,0.000000}%
\pgfsetstrokecolor{currentstroke}%
\pgfsetdash{}{0pt}%
\pgfpathmoveto{\pgfqpoint{1.346060in}{1.122652in}}%
\pgfpathlineto{\pgfqpoint{1.361789in}{1.120342in}}%
\pgfpathlineto{\pgfqpoint{1.362085in}{1.143717in}}%
\pgfpathlineto{\pgfqpoint{1.346336in}{1.146121in}}%
\pgfpathlineto{\pgfqpoint{1.346060in}{1.122652in}}%
\pgfpathclose%
\pgfusepath{fill}%
\end{pgfscope}%
\begin{pgfscope}%
\pgfpathrectangle{\pgfqpoint{0.407592in}{0.222904in}}{\pgfqpoint{1.411163in}{1.411163in}}%
\pgfusepath{clip}%
\pgfsetbuttcap%
\pgfsetroundjoin%
\definecolor{currentfill}{rgb}{0.168627,0.611765,0.411765}%
\pgfsetfillcolor{currentfill}%
\pgfsetfillopacity{0.180000}%
\pgfsetlinewidth{0.000000pt}%
\definecolor{currentstroke}{rgb}{0.000000,0.000000,0.000000}%
\pgfsetstrokecolor{currentstroke}%
\pgfsetdash{}{0pt}%
\pgfpathmoveto{\pgfqpoint{1.329453in}{1.055812in}}%
\pgfpathlineto{\pgfqpoint{1.345234in}{1.052607in}}%
\pgfpathlineto{\pgfqpoint{1.345509in}{1.075895in}}%
\pgfpathlineto{\pgfqpoint{1.329708in}{1.079196in}}%
\pgfpathlineto{\pgfqpoint{1.329453in}{1.055812in}}%
\pgfpathclose%
\pgfusepath{fill}%
\end{pgfscope}%
\begin{pgfscope}%
\pgfpathrectangle{\pgfqpoint{0.407592in}{0.222904in}}{\pgfqpoint{1.411163in}{1.411163in}}%
\pgfusepath{clip}%
\pgfsetbuttcap%
\pgfsetroundjoin%
\definecolor{currentfill}{rgb}{0.168627,0.611765,0.411765}%
\pgfsetfillcolor{currentfill}%
\pgfsetfillopacity{0.180000}%
\pgfsetlinewidth{0.000000pt}%
\definecolor{currentstroke}{rgb}{0.000000,0.000000,0.000000}%
\pgfsetstrokecolor{currentstroke}%
\pgfsetdash{}{0pt}%
\pgfpathmoveto{\pgfqpoint{1.312855in}{0.990025in}}%
\pgfpathlineto{\pgfqpoint{1.328692in}{0.986020in}}%
\pgfpathlineto{\pgfqpoint{1.328945in}{1.009224in}}%
\pgfpathlineto{\pgfqpoint{1.313089in}{1.013327in}}%
\pgfpathlineto{\pgfqpoint{1.312855in}{0.990025in}}%
\pgfpathclose%
\pgfusepath{fill}%
\end{pgfscope}%
\begin{pgfscope}%
\pgfpathrectangle{\pgfqpoint{0.407592in}{0.222904in}}{\pgfqpoint{1.411163in}{1.411163in}}%
\pgfusepath{clip}%
\pgfsetbuttcap%
\pgfsetroundjoin%
\definecolor{currentfill}{rgb}{0.168627,0.611765,0.411765}%
\pgfsetfillcolor{currentfill}%
\pgfsetfillopacity{0.180000}%
\pgfsetlinewidth{0.000000pt}%
\definecolor{currentstroke}{rgb}{0.000000,0.000000,0.000000}%
\pgfsetstrokecolor{currentstroke}%
\pgfsetdash{}{0pt}%
\pgfpathmoveto{\pgfqpoint{1.296262in}{0.925194in}}%
\pgfpathlineto{\pgfqpoint{1.312158in}{0.920478in}}%
\pgfpathlineto{\pgfqpoint{1.312390in}{0.943601in}}%
\pgfpathlineto{\pgfqpoint{1.296474in}{0.948416in}}%
\pgfpathlineto{\pgfqpoint{1.296262in}{0.925194in}}%
\pgfpathclose%
\pgfusepath{fill}%
\end{pgfscope}%
\begin{pgfscope}%
\pgfpathrectangle{\pgfqpoint{0.407592in}{0.222904in}}{\pgfqpoint{1.411163in}{1.411163in}}%
\pgfusepath{clip}%
\pgfsetbuttcap%
\pgfsetroundjoin%
\definecolor{currentfill}{rgb}{0.168627,0.611765,0.411765}%
\pgfsetfillcolor{currentfill}%
\pgfsetfillopacity{0.180000}%
\pgfsetlinewidth{0.000000pt}%
\definecolor{currentstroke}{rgb}{0.000000,0.000000,0.000000}%
\pgfsetstrokecolor{currentstroke}%
\pgfsetdash{}{0pt}%
\pgfpathmoveto{\pgfqpoint{1.346336in}{1.146121in}}%
\pgfpathlineto{\pgfqpoint{1.362085in}{1.143717in}}%
\pgfpathlineto{\pgfqpoint{1.362381in}{1.167153in}}%
\pgfpathlineto{\pgfqpoint{1.346613in}{1.169651in}}%
\pgfpathlineto{\pgfqpoint{1.346336in}{1.146121in}}%
\pgfpathclose%
\pgfusepath{fill}%
\end{pgfscope}%
\begin{pgfscope}%
\pgfpathrectangle{\pgfqpoint{0.407592in}{0.222904in}}{\pgfqpoint{1.411163in}{1.411163in}}%
\pgfusepath{clip}%
\pgfsetbuttcap%
\pgfsetroundjoin%
\definecolor{currentfill}{rgb}{0.168627,0.611765,0.411765}%
\pgfsetfillcolor{currentfill}%
\pgfsetfillopacity{0.180000}%
\pgfsetlinewidth{0.000000pt}%
\definecolor{currentstroke}{rgb}{0.000000,0.000000,0.000000}%
\pgfsetstrokecolor{currentstroke}%
\pgfsetdash{}{0pt}%
\pgfpathmoveto{\pgfqpoint{1.329708in}{1.079196in}}%
\pgfpathlineto{\pgfqpoint{1.345509in}{1.075895in}}%
\pgfpathlineto{\pgfqpoint{1.345784in}{1.099243in}}%
\pgfpathlineto{\pgfqpoint{1.329964in}{1.102641in}}%
\pgfpathlineto{\pgfqpoint{1.329708in}{1.079196in}}%
\pgfpathclose%
\pgfusepath{fill}%
\end{pgfscope}%
\begin{pgfscope}%
\pgfpathrectangle{\pgfqpoint{0.407592in}{0.222904in}}{\pgfqpoint{1.411163in}{1.411163in}}%
\pgfusepath{clip}%
\pgfsetbuttcap%
\pgfsetroundjoin%
\definecolor{currentfill}{rgb}{0.168627,0.611765,0.411765}%
\pgfsetfillcolor{currentfill}%
\pgfsetfillopacity{0.180000}%
\pgfsetlinewidth{0.000000pt}%
\definecolor{currentstroke}{rgb}{0.000000,0.000000,0.000000}%
\pgfsetstrokecolor{currentstroke}%
\pgfsetdash{}{0pt}%
\pgfpathmoveto{\pgfqpoint{1.313089in}{1.013327in}}%
\pgfpathlineto{\pgfqpoint{1.328945in}{1.009224in}}%
\pgfpathlineto{\pgfqpoint{1.329199in}{1.032488in}}%
\pgfpathlineto{\pgfqpoint{1.313323in}{1.036690in}}%
\pgfpathlineto{\pgfqpoint{1.313089in}{1.013327in}}%
\pgfpathclose%
\pgfusepath{fill}%
\end{pgfscope}%
\begin{pgfscope}%
\pgfpathrectangle{\pgfqpoint{0.407592in}{0.222904in}}{\pgfqpoint{1.411163in}{1.411163in}}%
\pgfusepath{clip}%
\pgfsetbuttcap%
\pgfsetroundjoin%
\definecolor{currentfill}{rgb}{0.168627,0.611765,0.411765}%
\pgfsetfillcolor{currentfill}%
\pgfsetfillopacity{0.180000}%
\pgfsetlinewidth{0.000000pt}%
\definecolor{currentstroke}{rgb}{0.000000,0.000000,0.000000}%
\pgfsetstrokecolor{currentstroke}%
\pgfsetdash{}{0pt}%
\pgfpathmoveto{\pgfqpoint{1.296474in}{0.948416in}}%
\pgfpathlineto{\pgfqpoint{1.312390in}{0.943601in}}%
\pgfpathlineto{\pgfqpoint{1.312622in}{0.966783in}}%
\pgfpathlineto{\pgfqpoint{1.296687in}{0.971699in}}%
\pgfpathlineto{\pgfqpoint{1.296474in}{0.948416in}}%
\pgfpathclose%
\pgfusepath{fill}%
\end{pgfscope}%
\begin{pgfscope}%
\pgfpathrectangle{\pgfqpoint{0.407592in}{0.222904in}}{\pgfqpoint{1.411163in}{1.411163in}}%
\pgfusepath{clip}%
\pgfsetbuttcap%
\pgfsetroundjoin%
\definecolor{currentfill}{rgb}{0.168627,0.611765,0.411765}%
\pgfsetfillcolor{currentfill}%
\pgfsetfillopacity{0.180000}%
\pgfsetlinewidth{0.000000pt}%
\definecolor{currentstroke}{rgb}{0.000000,0.000000,0.000000}%
\pgfsetstrokecolor{currentstroke}%
\pgfsetdash{}{0pt}%
\pgfpathmoveto{\pgfqpoint{1.279861in}{0.884370in}}%
\pgfpathlineto{\pgfqpoint{1.295839in}{0.878928in}}%
\pgfpathlineto{\pgfqpoint{1.296051in}{0.902031in}}%
\pgfpathlineto{\pgfqpoint{1.280052in}{0.907576in}}%
\pgfpathlineto{\pgfqpoint{1.279861in}{0.884370in}}%
\pgfpathclose%
\pgfusepath{fill}%
\end{pgfscope}%
\begin{pgfscope}%
\pgfpathrectangle{\pgfqpoint{0.407592in}{0.222904in}}{\pgfqpoint{1.411163in}{1.411163in}}%
\pgfusepath{clip}%
\pgfsetbuttcap%
\pgfsetroundjoin%
\definecolor{currentfill}{rgb}{0.168627,0.611765,0.411765}%
\pgfsetfillcolor{currentfill}%
\pgfsetfillopacity{0.180000}%
\pgfsetlinewidth{0.000000pt}%
\definecolor{currentstroke}{rgb}{0.000000,0.000000,0.000000}%
\pgfsetstrokecolor{currentstroke}%
\pgfsetdash{}{0pt}%
\pgfpathmoveto{\pgfqpoint{1.329964in}{1.102641in}}%
\pgfpathlineto{\pgfqpoint{1.345784in}{1.099243in}}%
\pgfpathlineto{\pgfqpoint{1.346060in}{1.122652in}}%
\pgfpathlineto{\pgfqpoint{1.330221in}{1.126147in}}%
\pgfpathlineto{\pgfqpoint{1.329964in}{1.102641in}}%
\pgfpathclose%
\pgfusepath{fill}%
\end{pgfscope}%
\begin{pgfscope}%
\pgfpathrectangle{\pgfqpoint{0.407592in}{0.222904in}}{\pgfqpoint{1.411163in}{1.411163in}}%
\pgfusepath{clip}%
\pgfsetbuttcap%
\pgfsetroundjoin%
\definecolor{currentfill}{rgb}{0.168627,0.611765,0.411765}%
\pgfsetfillcolor{currentfill}%
\pgfsetfillopacity{0.180000}%
\pgfsetlinewidth{0.000000pt}%
\definecolor{currentstroke}{rgb}{0.000000,0.000000,0.000000}%
\pgfsetstrokecolor{currentstroke}%
\pgfsetdash{}{0pt}%
\pgfpathmoveto{\pgfqpoint{1.313323in}{1.036690in}}%
\pgfpathlineto{\pgfqpoint{1.329199in}{1.032488in}}%
\pgfpathlineto{\pgfqpoint{1.329453in}{1.055812in}}%
\pgfpathlineto{\pgfqpoint{1.313558in}{1.060113in}}%
\pgfpathlineto{\pgfqpoint{1.313323in}{1.036690in}}%
\pgfpathclose%
\pgfusepath{fill}%
\end{pgfscope}%
\begin{pgfscope}%
\pgfpathrectangle{\pgfqpoint{0.407592in}{0.222904in}}{\pgfqpoint{1.411163in}{1.411163in}}%
\pgfusepath{clip}%
\pgfsetbuttcap%
\pgfsetroundjoin%
\definecolor{currentfill}{rgb}{0.168627,0.611765,0.411765}%
\pgfsetfillcolor{currentfill}%
\pgfsetfillopacity{0.180000}%
\pgfsetlinewidth{0.000000pt}%
\definecolor{currentstroke}{rgb}{0.000000,0.000000,0.000000}%
\pgfsetstrokecolor{currentstroke}%
\pgfsetdash{}{0pt}%
\pgfpathmoveto{\pgfqpoint{1.296687in}{0.971699in}}%
\pgfpathlineto{\pgfqpoint{1.312622in}{0.966783in}}%
\pgfpathlineto{\pgfqpoint{1.312855in}{0.990025in}}%
\pgfpathlineto{\pgfqpoint{1.296901in}{0.995042in}}%
\pgfpathlineto{\pgfqpoint{1.296687in}{0.971699in}}%
\pgfpathclose%
\pgfusepath{fill}%
\end{pgfscope}%
\begin{pgfscope}%
\pgfpathrectangle{\pgfqpoint{0.407592in}{0.222904in}}{\pgfqpoint{1.411163in}{1.411163in}}%
\pgfusepath{clip}%
\pgfsetbuttcap%
\pgfsetroundjoin%
\definecolor{currentfill}{rgb}{0.168627,0.611765,0.411765}%
\pgfsetfillcolor{currentfill}%
\pgfsetfillopacity{0.180000}%
\pgfsetlinewidth{0.000000pt}%
\definecolor{currentstroke}{rgb}{0.000000,0.000000,0.000000}%
\pgfsetstrokecolor{currentstroke}%
\pgfsetdash{}{0pt}%
\pgfpathmoveto{\pgfqpoint{1.280052in}{0.907576in}}%
\pgfpathlineto{\pgfqpoint{1.296051in}{0.902031in}}%
\pgfpathlineto{\pgfqpoint{1.296262in}{0.925194in}}%
\pgfpathlineto{\pgfqpoint{1.280244in}{0.930841in}}%
\pgfpathlineto{\pgfqpoint{1.280052in}{0.907576in}}%
\pgfpathclose%
\pgfusepath{fill}%
\end{pgfscope}%
\begin{pgfscope}%
\pgfpathrectangle{\pgfqpoint{0.407592in}{0.222904in}}{\pgfqpoint{1.411163in}{1.411163in}}%
\pgfusepath{clip}%
\pgfsetbuttcap%
\pgfsetroundjoin%
\definecolor{currentfill}{rgb}{0.168627,0.611765,0.411765}%
\pgfsetfillcolor{currentfill}%
\pgfsetfillopacity{0.180000}%
\pgfsetlinewidth{0.000000pt}%
\definecolor{currentstroke}{rgb}{0.000000,0.000000,0.000000}%
\pgfsetstrokecolor{currentstroke}%
\pgfsetdash{}{0pt}%
\pgfpathmoveto{\pgfqpoint{1.330221in}{1.126147in}}%
\pgfpathlineto{\pgfqpoint{1.346060in}{1.122652in}}%
\pgfpathlineto{\pgfqpoint{1.346336in}{1.146121in}}%
\pgfpathlineto{\pgfqpoint{1.330478in}{1.149713in}}%
\pgfpathlineto{\pgfqpoint{1.330221in}{1.126147in}}%
\pgfpathclose%
\pgfusepath{fill}%
\end{pgfscope}%
\begin{pgfscope}%
\pgfpathrectangle{\pgfqpoint{0.407592in}{0.222904in}}{\pgfqpoint{1.411163in}{1.411163in}}%
\pgfusepath{clip}%
\pgfsetbuttcap%
\pgfsetroundjoin%
\definecolor{currentfill}{rgb}{0.168627,0.611765,0.411765}%
\pgfsetfillcolor{currentfill}%
\pgfsetfillopacity{0.180000}%
\pgfsetlinewidth{0.000000pt}%
\definecolor{currentstroke}{rgb}{0.000000,0.000000,0.000000}%
\pgfsetstrokecolor{currentstroke}%
\pgfsetdash{}{0pt}%
\pgfpathmoveto{\pgfqpoint{1.313558in}{1.060113in}}%
\pgfpathlineto{\pgfqpoint{1.329453in}{1.055812in}}%
\pgfpathlineto{\pgfqpoint{1.329708in}{1.079196in}}%
\pgfpathlineto{\pgfqpoint{1.313794in}{1.083597in}}%
\pgfpathlineto{\pgfqpoint{1.313558in}{1.060113in}}%
\pgfpathclose%
\pgfusepath{fill}%
\end{pgfscope}%
\begin{pgfscope}%
\pgfpathrectangle{\pgfqpoint{0.407592in}{0.222904in}}{\pgfqpoint{1.411163in}{1.411163in}}%
\pgfusepath{clip}%
\pgfsetbuttcap%
\pgfsetroundjoin%
\definecolor{currentfill}{rgb}{0.168627,0.611765,0.411765}%
\pgfsetfillcolor{currentfill}%
\pgfsetfillopacity{0.180000}%
\pgfsetlinewidth{0.000000pt}%
\definecolor{currentstroke}{rgb}{0.000000,0.000000,0.000000}%
\pgfsetstrokecolor{currentstroke}%
\pgfsetdash{}{0pt}%
\pgfpathmoveto{\pgfqpoint{1.296901in}{0.995042in}}%
\pgfpathlineto{\pgfqpoint{1.312855in}{0.990025in}}%
\pgfpathlineto{\pgfqpoint{1.313089in}{1.013327in}}%
\pgfpathlineto{\pgfqpoint{1.297114in}{1.018446in}}%
\pgfpathlineto{\pgfqpoint{1.296901in}{0.995042in}}%
\pgfpathclose%
\pgfusepath{fill}%
\end{pgfscope}%
\begin{pgfscope}%
\pgfpathrectangle{\pgfqpoint{0.407592in}{0.222904in}}{\pgfqpoint{1.411163in}{1.411163in}}%
\pgfusepath{clip}%
\pgfsetbuttcap%
\pgfsetroundjoin%
\definecolor{currentfill}{rgb}{0.168627,0.611765,0.411765}%
\pgfsetfillcolor{currentfill}%
\pgfsetfillopacity{0.180000}%
\pgfsetlinewidth{0.000000pt}%
\definecolor{currentstroke}{rgb}{0.000000,0.000000,0.000000}%
\pgfsetstrokecolor{currentstroke}%
\pgfsetdash{}{0pt}%
\pgfpathmoveto{\pgfqpoint{1.280244in}{0.930841in}}%
\pgfpathlineto{\pgfqpoint{1.296262in}{0.925194in}}%
\pgfpathlineto{\pgfqpoint{1.296474in}{0.948416in}}%
\pgfpathlineto{\pgfqpoint{1.280436in}{0.954167in}}%
\pgfpathlineto{\pgfqpoint{1.280244in}{0.930841in}}%
\pgfpathclose%
\pgfusepath{fill}%
\end{pgfscope}%
\begin{pgfscope}%
\pgfpathrectangle{\pgfqpoint{0.407592in}{0.222904in}}{\pgfqpoint{1.411163in}{1.411163in}}%
\pgfusepath{clip}%
\pgfsetbuttcap%
\pgfsetroundjoin%
\definecolor{currentfill}{rgb}{0.168627,0.611765,0.411765}%
\pgfsetfillcolor{currentfill}%
\pgfsetfillopacity{0.180000}%
\pgfsetlinewidth{0.000000pt}%
\definecolor{currentstroke}{rgb}{0.000000,0.000000,0.000000}%
\pgfsetstrokecolor{currentstroke}%
\pgfsetdash{}{0pt}%
\pgfpathmoveto{\pgfqpoint{1.330478in}{1.149713in}}%
\pgfpathlineto{\pgfqpoint{1.346336in}{1.146121in}}%
\pgfpathlineto{\pgfqpoint{1.346613in}{1.169651in}}%
\pgfpathlineto{\pgfqpoint{1.330735in}{1.173342in}}%
\pgfpathlineto{\pgfqpoint{1.330478in}{1.149713in}}%
\pgfpathclose%
\pgfusepath{fill}%
\end{pgfscope}%
\begin{pgfscope}%
\pgfpathrectangle{\pgfqpoint{0.407592in}{0.222904in}}{\pgfqpoint{1.411163in}{1.411163in}}%
\pgfusepath{clip}%
\pgfsetbuttcap%
\pgfsetroundjoin%
\definecolor{currentfill}{rgb}{0.168627,0.611765,0.411765}%
\pgfsetfillcolor{currentfill}%
\pgfsetfillopacity{0.180000}%
\pgfsetlinewidth{0.000000pt}%
\definecolor{currentstroke}{rgb}{0.000000,0.000000,0.000000}%
\pgfsetstrokecolor{currentstroke}%
\pgfsetdash{}{0pt}%
\pgfpathmoveto{\pgfqpoint{1.313794in}{1.083597in}}%
\pgfpathlineto{\pgfqpoint{1.329708in}{1.079196in}}%
\pgfpathlineto{\pgfqpoint{1.329964in}{1.102641in}}%
\pgfpathlineto{\pgfqpoint{1.314030in}{1.107142in}}%
\pgfpathlineto{\pgfqpoint{1.313794in}{1.083597in}}%
\pgfpathclose%
\pgfusepath{fill}%
\end{pgfscope}%
\begin{pgfscope}%
\pgfpathrectangle{\pgfqpoint{0.407592in}{0.222904in}}{\pgfqpoint{1.411163in}{1.411163in}}%
\pgfusepath{clip}%
\pgfsetbuttcap%
\pgfsetroundjoin%
\definecolor{currentfill}{rgb}{0.168627,0.611765,0.411765}%
\pgfsetfillcolor{currentfill}%
\pgfsetfillopacity{0.180000}%
\pgfsetlinewidth{0.000000pt}%
\definecolor{currentstroke}{rgb}{0.000000,0.000000,0.000000}%
\pgfsetstrokecolor{currentstroke}%
\pgfsetdash{}{0pt}%
\pgfpathmoveto{\pgfqpoint{1.297114in}{1.018446in}}%
\pgfpathlineto{\pgfqpoint{1.313089in}{1.013327in}}%
\pgfpathlineto{\pgfqpoint{1.313323in}{1.036690in}}%
\pgfpathlineto{\pgfqpoint{1.297329in}{1.041911in}}%
\pgfpathlineto{\pgfqpoint{1.297114in}{1.018446in}}%
\pgfpathclose%
\pgfusepath{fill}%
\end{pgfscope}%
\begin{pgfscope}%
\pgfpathrectangle{\pgfqpoint{0.407592in}{0.222904in}}{\pgfqpoint{1.411163in}{1.411163in}}%
\pgfusepath{clip}%
\pgfsetbuttcap%
\pgfsetroundjoin%
\definecolor{currentfill}{rgb}{0.168627,0.611765,0.411765}%
\pgfsetfillcolor{currentfill}%
\pgfsetfillopacity{0.180000}%
\pgfsetlinewidth{0.000000pt}%
\definecolor{currentstroke}{rgb}{0.000000,0.000000,0.000000}%
\pgfsetstrokecolor{currentstroke}%
\pgfsetdash{}{0pt}%
\pgfpathmoveto{\pgfqpoint{1.280436in}{0.954167in}}%
\pgfpathlineto{\pgfqpoint{1.296474in}{0.948416in}}%
\pgfpathlineto{\pgfqpoint{1.296687in}{0.971699in}}%
\pgfpathlineto{\pgfqpoint{1.280629in}{0.977553in}}%
\pgfpathlineto{\pgfqpoint{1.280436in}{0.954167in}}%
\pgfpathclose%
\pgfusepath{fill}%
\end{pgfscope}%
\begin{pgfscope}%
\pgfpathrectangle{\pgfqpoint{0.407592in}{0.222904in}}{\pgfqpoint{1.411163in}{1.411163in}}%
\pgfusepath{clip}%
\pgfsetbuttcap%
\pgfsetroundjoin%
\definecolor{currentfill}{rgb}{0.168627,0.611765,0.411765}%
\pgfsetfillcolor{currentfill}%
\pgfsetfillopacity{0.180000}%
\pgfsetlinewidth{0.000000pt}%
\definecolor{currentstroke}{rgb}{0.000000,0.000000,0.000000}%
\pgfsetstrokecolor{currentstroke}%
\pgfsetdash{}{0pt}%
\pgfpathmoveto{\pgfqpoint{1.263757in}{0.890672in}}%
\pgfpathlineto{\pgfqpoint{1.279861in}{0.884370in}}%
\pgfpathlineto{\pgfqpoint{1.280052in}{0.907576in}}%
\pgfpathlineto{\pgfqpoint{1.263928in}{0.913982in}}%
\pgfpathlineto{\pgfqpoint{1.263757in}{0.890672in}}%
\pgfpathclose%
\pgfusepath{fill}%
\end{pgfscope}%
\begin{pgfscope}%
\pgfpathrectangle{\pgfqpoint{0.407592in}{0.222904in}}{\pgfqpoint{1.411163in}{1.411163in}}%
\pgfusepath{clip}%
\pgfsetbuttcap%
\pgfsetroundjoin%
\definecolor{currentfill}{rgb}{0.168627,0.611765,0.411765}%
\pgfsetfillcolor{currentfill}%
\pgfsetfillopacity{0.180000}%
\pgfsetlinewidth{0.000000pt}%
\definecolor{currentstroke}{rgb}{0.000000,0.000000,0.000000}%
\pgfsetstrokecolor{currentstroke}%
\pgfsetdash{}{0pt}%
\pgfpathmoveto{\pgfqpoint{1.314030in}{1.107142in}}%
\pgfpathlineto{\pgfqpoint{1.329964in}{1.102641in}}%
\pgfpathlineto{\pgfqpoint{1.330221in}{1.126147in}}%
\pgfpathlineto{\pgfqpoint{1.314266in}{1.130749in}}%
\pgfpathlineto{\pgfqpoint{1.314030in}{1.107142in}}%
\pgfpathclose%
\pgfusepath{fill}%
\end{pgfscope}%
\begin{pgfscope}%
\pgfpathrectangle{\pgfqpoint{0.407592in}{0.222904in}}{\pgfqpoint{1.411163in}{1.411163in}}%
\pgfusepath{clip}%
\pgfsetbuttcap%
\pgfsetroundjoin%
\definecolor{currentfill}{rgb}{0.168627,0.611765,0.411765}%
\pgfsetfillcolor{currentfill}%
\pgfsetfillopacity{0.180000}%
\pgfsetlinewidth{0.000000pt}%
\definecolor{currentstroke}{rgb}{0.000000,0.000000,0.000000}%
\pgfsetstrokecolor{currentstroke}%
\pgfsetdash{}{0pt}%
\pgfpathmoveto{\pgfqpoint{1.297329in}{1.041911in}}%
\pgfpathlineto{\pgfqpoint{1.313323in}{1.036690in}}%
\pgfpathlineto{\pgfqpoint{1.313558in}{1.060113in}}%
\pgfpathlineto{\pgfqpoint{1.297544in}{1.065437in}}%
\pgfpathlineto{\pgfqpoint{1.297329in}{1.041911in}}%
\pgfpathclose%
\pgfusepath{fill}%
\end{pgfscope}%
\begin{pgfscope}%
\pgfpathrectangle{\pgfqpoint{0.407592in}{0.222904in}}{\pgfqpoint{1.411163in}{1.411163in}}%
\pgfusepath{clip}%
\pgfsetbuttcap%
\pgfsetroundjoin%
\definecolor{currentfill}{rgb}{0.168627,0.611765,0.411765}%
\pgfsetfillcolor{currentfill}%
\pgfsetfillopacity{0.180000}%
\pgfsetlinewidth{0.000000pt}%
\definecolor{currentstroke}{rgb}{0.000000,0.000000,0.000000}%
\pgfsetstrokecolor{currentstroke}%
\pgfsetdash{}{0pt}%
\pgfpathmoveto{\pgfqpoint{1.280629in}{0.977553in}}%
\pgfpathlineto{\pgfqpoint{1.296687in}{0.971699in}}%
\pgfpathlineto{\pgfqpoint{1.296901in}{0.995042in}}%
\pgfpathlineto{\pgfqpoint{1.280823in}{1.001001in}}%
\pgfpathlineto{\pgfqpoint{1.280629in}{0.977553in}}%
\pgfpathclose%
\pgfusepath{fill}%
\end{pgfscope}%
\begin{pgfscope}%
\pgfpathrectangle{\pgfqpoint{0.407592in}{0.222904in}}{\pgfqpoint{1.411163in}{1.411163in}}%
\pgfusepath{clip}%
\pgfsetbuttcap%
\pgfsetroundjoin%
\definecolor{currentfill}{rgb}{0.168627,0.611765,0.411765}%
\pgfsetfillcolor{currentfill}%
\pgfsetfillopacity{0.180000}%
\pgfsetlinewidth{0.000000pt}%
\definecolor{currentstroke}{rgb}{0.000000,0.000000,0.000000}%
\pgfsetstrokecolor{currentstroke}%
\pgfsetdash{}{0pt}%
\pgfpathmoveto{\pgfqpoint{1.263928in}{0.913982in}}%
\pgfpathlineto{\pgfqpoint{1.280052in}{0.907576in}}%
\pgfpathlineto{\pgfqpoint{1.280244in}{0.930841in}}%
\pgfpathlineto{\pgfqpoint{1.264099in}{0.937353in}}%
\pgfpathlineto{\pgfqpoint{1.263928in}{0.913982in}}%
\pgfpathclose%
\pgfusepath{fill}%
\end{pgfscope}%
\begin{pgfscope}%
\pgfpathrectangle{\pgfqpoint{0.407592in}{0.222904in}}{\pgfqpoint{1.411163in}{1.411163in}}%
\pgfusepath{clip}%
\pgfsetbuttcap%
\pgfsetroundjoin%
\definecolor{currentfill}{rgb}{0.168627,0.611765,0.411765}%
\pgfsetfillcolor{currentfill}%
\pgfsetfillopacity{0.180000}%
\pgfsetlinewidth{0.000000pt}%
\definecolor{currentstroke}{rgb}{0.000000,0.000000,0.000000}%
\pgfsetstrokecolor{currentstroke}%
\pgfsetdash{}{0pt}%
\pgfpathmoveto{\pgfqpoint{1.314266in}{1.130749in}}%
\pgfpathlineto{\pgfqpoint{1.330221in}{1.126147in}}%
\pgfpathlineto{\pgfqpoint{1.330478in}{1.149713in}}%
\pgfpathlineto{\pgfqpoint{1.314504in}{1.154417in}}%
\pgfpathlineto{\pgfqpoint{1.314266in}{1.130749in}}%
\pgfpathclose%
\pgfusepath{fill}%
\end{pgfscope}%
\begin{pgfscope}%
\pgfpathrectangle{\pgfqpoint{0.407592in}{0.222904in}}{\pgfqpoint{1.411163in}{1.411163in}}%
\pgfusepath{clip}%
\pgfsetbuttcap%
\pgfsetroundjoin%
\definecolor{currentfill}{rgb}{0.168627,0.611765,0.411765}%
\pgfsetfillcolor{currentfill}%
\pgfsetfillopacity{0.180000}%
\pgfsetlinewidth{0.000000pt}%
\definecolor{currentstroke}{rgb}{0.000000,0.000000,0.000000}%
\pgfsetstrokecolor{currentstroke}%
\pgfsetdash{}{0pt}%
\pgfpathmoveto{\pgfqpoint{1.297544in}{1.065437in}}%
\pgfpathlineto{\pgfqpoint{1.313558in}{1.060113in}}%
\pgfpathlineto{\pgfqpoint{1.313794in}{1.083597in}}%
\pgfpathlineto{\pgfqpoint{1.297759in}{1.089024in}}%
\pgfpathlineto{\pgfqpoint{1.297544in}{1.065437in}}%
\pgfpathclose%
\pgfusepath{fill}%
\end{pgfscope}%
\begin{pgfscope}%
\pgfpathrectangle{\pgfqpoint{0.407592in}{0.222904in}}{\pgfqpoint{1.411163in}{1.411163in}}%
\pgfusepath{clip}%
\pgfsetbuttcap%
\pgfsetroundjoin%
\definecolor{currentfill}{rgb}{0.168627,0.611765,0.411765}%
\pgfsetfillcolor{currentfill}%
\pgfsetfillopacity{0.180000}%
\pgfsetlinewidth{0.000000pt}%
\definecolor{currentstroke}{rgb}{0.000000,0.000000,0.000000}%
\pgfsetstrokecolor{currentstroke}%
\pgfsetdash{}{0pt}%
\pgfpathmoveto{\pgfqpoint{1.280823in}{1.001001in}}%
\pgfpathlineto{\pgfqpoint{1.296901in}{0.995042in}}%
\pgfpathlineto{\pgfqpoint{1.297114in}{1.018446in}}%
\pgfpathlineto{\pgfqpoint{1.281016in}{1.024509in}}%
\pgfpathlineto{\pgfqpoint{1.280823in}{1.001001in}}%
\pgfpathclose%
\pgfusepath{fill}%
\end{pgfscope}%
\begin{pgfscope}%
\pgfpathrectangle{\pgfqpoint{0.407592in}{0.222904in}}{\pgfqpoint{1.411163in}{1.411163in}}%
\pgfusepath{clip}%
\pgfsetbuttcap%
\pgfsetroundjoin%
\definecolor{currentfill}{rgb}{0.168627,0.611765,0.411765}%
\pgfsetfillcolor{currentfill}%
\pgfsetfillopacity{0.180000}%
\pgfsetlinewidth{0.000000pt}%
\definecolor{currentstroke}{rgb}{0.000000,0.000000,0.000000}%
\pgfsetstrokecolor{currentstroke}%
\pgfsetdash{}{0pt}%
\pgfpathmoveto{\pgfqpoint{1.264099in}{0.937353in}}%
\pgfpathlineto{\pgfqpoint{1.280244in}{0.930841in}}%
\pgfpathlineto{\pgfqpoint{1.280436in}{0.954167in}}%
\pgfpathlineto{\pgfqpoint{1.264271in}{0.960785in}}%
\pgfpathlineto{\pgfqpoint{1.264099in}{0.937353in}}%
\pgfpathclose%
\pgfusepath{fill}%
\end{pgfscope}%
\begin{pgfscope}%
\pgfpathrectangle{\pgfqpoint{0.407592in}{0.222904in}}{\pgfqpoint{1.411163in}{1.411163in}}%
\pgfusepath{clip}%
\pgfsetbuttcap%
\pgfsetroundjoin%
\definecolor{currentfill}{rgb}{0.168627,0.611765,0.411765}%
\pgfsetfillcolor{currentfill}%
\pgfsetfillopacity{0.180000}%
\pgfsetlinewidth{0.000000pt}%
\definecolor{currentstroke}{rgb}{0.000000,0.000000,0.000000}%
\pgfsetstrokecolor{currentstroke}%
\pgfsetdash{}{0pt}%
\pgfpathmoveto{\pgfqpoint{1.314504in}{1.154417in}}%
\pgfpathlineto{\pgfqpoint{1.330478in}{1.149713in}}%
\pgfpathlineto{\pgfqpoint{1.330735in}{1.173342in}}%
\pgfpathlineto{\pgfqpoint{1.314742in}{1.178147in}}%
\pgfpathlineto{\pgfqpoint{1.314504in}{1.154417in}}%
\pgfpathclose%
\pgfusepath{fill}%
\end{pgfscope}%
\begin{pgfscope}%
\pgfpathrectangle{\pgfqpoint{0.407592in}{0.222904in}}{\pgfqpoint{1.411163in}{1.411163in}}%
\pgfusepath{clip}%
\pgfsetbuttcap%
\pgfsetroundjoin%
\definecolor{currentfill}{rgb}{0.168627,0.611765,0.411765}%
\pgfsetfillcolor{currentfill}%
\pgfsetfillopacity{0.180000}%
\pgfsetlinewidth{0.000000pt}%
\definecolor{currentstroke}{rgb}{0.000000,0.000000,0.000000}%
\pgfsetstrokecolor{currentstroke}%
\pgfsetdash{}{0pt}%
\pgfpathmoveto{\pgfqpoint{1.297759in}{1.089024in}}%
\pgfpathlineto{\pgfqpoint{1.313794in}{1.083597in}}%
\pgfpathlineto{\pgfqpoint{1.314030in}{1.107142in}}%
\pgfpathlineto{\pgfqpoint{1.297975in}{1.112673in}}%
\pgfpathlineto{\pgfqpoint{1.297759in}{1.089024in}}%
\pgfpathclose%
\pgfusepath{fill}%
\end{pgfscope}%
\begin{pgfscope}%
\pgfpathrectangle{\pgfqpoint{0.407592in}{0.222904in}}{\pgfqpoint{1.411163in}{1.411163in}}%
\pgfusepath{clip}%
\pgfsetbuttcap%
\pgfsetroundjoin%
\definecolor{currentfill}{rgb}{0.168627,0.611765,0.411765}%
\pgfsetfillcolor{currentfill}%
\pgfsetfillopacity{0.180000}%
\pgfsetlinewidth{0.000000pt}%
\definecolor{currentstroke}{rgb}{0.000000,0.000000,0.000000}%
\pgfsetstrokecolor{currentstroke}%
\pgfsetdash{}{0pt}%
\pgfpathmoveto{\pgfqpoint{1.281016in}{1.024509in}}%
\pgfpathlineto{\pgfqpoint{1.297114in}{1.018446in}}%
\pgfpathlineto{\pgfqpoint{1.297329in}{1.041911in}}%
\pgfpathlineto{\pgfqpoint{1.281211in}{1.048079in}}%
\pgfpathlineto{\pgfqpoint{1.281016in}{1.024509in}}%
\pgfpathclose%
\pgfusepath{fill}%
\end{pgfscope}%
\begin{pgfscope}%
\pgfpathrectangle{\pgfqpoint{0.407592in}{0.222904in}}{\pgfqpoint{1.411163in}{1.411163in}}%
\pgfusepath{clip}%
\pgfsetbuttcap%
\pgfsetroundjoin%
\definecolor{currentfill}{rgb}{0.168627,0.611765,0.411765}%
\pgfsetfillcolor{currentfill}%
\pgfsetfillopacity{0.180000}%
\pgfsetlinewidth{0.000000pt}%
\definecolor{currentstroke}{rgb}{0.000000,0.000000,0.000000}%
\pgfsetstrokecolor{currentstroke}%
\pgfsetdash{}{0pt}%
\pgfpathmoveto{\pgfqpoint{1.264271in}{0.960785in}}%
\pgfpathlineto{\pgfqpoint{1.280436in}{0.954167in}}%
\pgfpathlineto{\pgfqpoint{1.280629in}{0.977553in}}%
\pgfpathlineto{\pgfqpoint{1.264444in}{0.984279in}}%
\pgfpathlineto{\pgfqpoint{1.264271in}{0.960785in}}%
\pgfpathclose%
\pgfusepath{fill}%
\end{pgfscope}%
\begin{pgfscope}%
\pgfpathrectangle{\pgfqpoint{0.407592in}{0.222904in}}{\pgfqpoint{1.411163in}{1.411163in}}%
\pgfusepath{clip}%
\pgfsetbuttcap%
\pgfsetroundjoin%
\definecolor{currentfill}{rgb}{0.168627,0.611765,0.411765}%
\pgfsetfillcolor{currentfill}%
\pgfsetfillopacity{0.180000}%
\pgfsetlinewidth{0.000000pt}%
\definecolor{currentstroke}{rgb}{0.000000,0.000000,0.000000}%
\pgfsetstrokecolor{currentstroke}%
\pgfsetdash{}{0pt}%
\pgfpathmoveto{\pgfqpoint{1.247522in}{0.897771in}}%
\pgfpathlineto{\pgfqpoint{1.263757in}{0.890672in}}%
\pgfpathlineto{\pgfqpoint{1.263928in}{0.913982in}}%
\pgfpathlineto{\pgfqpoint{1.247673in}{0.921190in}}%
\pgfpathlineto{\pgfqpoint{1.247522in}{0.897771in}}%
\pgfpathclose%
\pgfusepath{fill}%
\end{pgfscope}%
\begin{pgfscope}%
\pgfpathrectangle{\pgfqpoint{0.407592in}{0.222904in}}{\pgfqpoint{1.411163in}{1.411163in}}%
\pgfusepath{clip}%
\pgfsetbuttcap%
\pgfsetroundjoin%
\definecolor{currentfill}{rgb}{0.168627,0.611765,0.411765}%
\pgfsetfillcolor{currentfill}%
\pgfsetfillopacity{0.180000}%
\pgfsetlinewidth{0.000000pt}%
\definecolor{currentstroke}{rgb}{0.000000,0.000000,0.000000}%
\pgfsetstrokecolor{currentstroke}%
\pgfsetdash{}{0pt}%
\pgfpathmoveto{\pgfqpoint{1.297975in}{1.112673in}}%
\pgfpathlineto{\pgfqpoint{1.314030in}{1.107142in}}%
\pgfpathlineto{\pgfqpoint{1.314266in}{1.130749in}}%
\pgfpathlineto{\pgfqpoint{1.298192in}{1.136383in}}%
\pgfpathlineto{\pgfqpoint{1.297975in}{1.112673in}}%
\pgfpathclose%
\pgfusepath{fill}%
\end{pgfscope}%
\begin{pgfscope}%
\pgfpathrectangle{\pgfqpoint{0.407592in}{0.222904in}}{\pgfqpoint{1.411163in}{1.411163in}}%
\pgfusepath{clip}%
\pgfsetbuttcap%
\pgfsetroundjoin%
\definecolor{currentfill}{rgb}{0.168627,0.611765,0.411765}%
\pgfsetfillcolor{currentfill}%
\pgfsetfillopacity{0.180000}%
\pgfsetlinewidth{0.000000pt}%
\definecolor{currentstroke}{rgb}{0.000000,0.000000,0.000000}%
\pgfsetstrokecolor{currentstroke}%
\pgfsetdash{}{0pt}%
\pgfpathmoveto{\pgfqpoint{1.281211in}{1.048079in}}%
\pgfpathlineto{\pgfqpoint{1.297329in}{1.041911in}}%
\pgfpathlineto{\pgfqpoint{1.297544in}{1.065437in}}%
\pgfpathlineto{\pgfqpoint{1.281405in}{1.071711in}}%
\pgfpathlineto{\pgfqpoint{1.281211in}{1.048079in}}%
\pgfpathclose%
\pgfusepath{fill}%
\end{pgfscope}%
\begin{pgfscope}%
\pgfpathrectangle{\pgfqpoint{0.407592in}{0.222904in}}{\pgfqpoint{1.411163in}{1.411163in}}%
\pgfusepath{clip}%
\pgfsetbuttcap%
\pgfsetroundjoin%
\definecolor{currentfill}{rgb}{0.168627,0.611765,0.411765}%
\pgfsetfillcolor{currentfill}%
\pgfsetfillopacity{0.180000}%
\pgfsetlinewidth{0.000000pt}%
\definecolor{currentstroke}{rgb}{0.000000,0.000000,0.000000}%
\pgfsetstrokecolor{currentstroke}%
\pgfsetdash{}{0pt}%
\pgfpathmoveto{\pgfqpoint{1.264444in}{0.984279in}}%
\pgfpathlineto{\pgfqpoint{1.280629in}{0.977553in}}%
\pgfpathlineto{\pgfqpoint{1.280823in}{1.001001in}}%
\pgfpathlineto{\pgfqpoint{1.264617in}{1.007834in}}%
\pgfpathlineto{\pgfqpoint{1.264444in}{0.984279in}}%
\pgfpathclose%
\pgfusepath{fill}%
\end{pgfscope}%
\begin{pgfscope}%
\pgfpathrectangle{\pgfqpoint{0.407592in}{0.222904in}}{\pgfqpoint{1.411163in}{1.411163in}}%
\pgfusepath{clip}%
\pgfsetbuttcap%
\pgfsetroundjoin%
\definecolor{currentfill}{rgb}{0.168627,0.611765,0.411765}%
\pgfsetfillcolor{currentfill}%
\pgfsetfillopacity{0.180000}%
\pgfsetlinewidth{0.000000pt}%
\definecolor{currentstroke}{rgb}{0.000000,0.000000,0.000000}%
\pgfsetstrokecolor{currentstroke}%
\pgfsetdash{}{0pt}%
\pgfpathmoveto{\pgfqpoint{1.247673in}{0.921190in}}%
\pgfpathlineto{\pgfqpoint{1.263928in}{0.913982in}}%
\pgfpathlineto{\pgfqpoint{1.264099in}{0.937353in}}%
\pgfpathlineto{\pgfqpoint{1.247824in}{0.944669in}}%
\pgfpathlineto{\pgfqpoint{1.247673in}{0.921190in}}%
\pgfpathclose%
\pgfusepath{fill}%
\end{pgfscope}%
\begin{pgfscope}%
\pgfpathrectangle{\pgfqpoint{0.407592in}{0.222904in}}{\pgfqpoint{1.411163in}{1.411163in}}%
\pgfusepath{clip}%
\pgfsetbuttcap%
\pgfsetroundjoin%
\definecolor{currentfill}{rgb}{0.168627,0.611765,0.411765}%
\pgfsetfillcolor{currentfill}%
\pgfsetfillopacity{0.180000}%
\pgfsetlinewidth{0.000000pt}%
\definecolor{currentstroke}{rgb}{0.000000,0.000000,0.000000}%
\pgfsetstrokecolor{currentstroke}%
\pgfsetdash{}{0pt}%
\pgfpathmoveto{\pgfqpoint{1.298192in}{1.136383in}}%
\pgfpathlineto{\pgfqpoint{1.314266in}{1.130749in}}%
\pgfpathlineto{\pgfqpoint{1.314504in}{1.154417in}}%
\pgfpathlineto{\pgfqpoint{1.298409in}{1.160156in}}%
\pgfpathlineto{\pgfqpoint{1.298192in}{1.136383in}}%
\pgfpathclose%
\pgfusepath{fill}%
\end{pgfscope}%
\begin{pgfscope}%
\pgfpathrectangle{\pgfqpoint{0.407592in}{0.222904in}}{\pgfqpoint{1.411163in}{1.411163in}}%
\pgfusepath{clip}%
\pgfsetbuttcap%
\pgfsetroundjoin%
\definecolor{currentfill}{rgb}{0.168627,0.611765,0.411765}%
\pgfsetfillcolor{currentfill}%
\pgfsetfillopacity{0.180000}%
\pgfsetlinewidth{0.000000pt}%
\definecolor{currentstroke}{rgb}{0.000000,0.000000,0.000000}%
\pgfsetstrokecolor{currentstroke}%
\pgfsetdash{}{0pt}%
\pgfpathmoveto{\pgfqpoint{1.281405in}{1.071711in}}%
\pgfpathlineto{\pgfqpoint{1.297544in}{1.065437in}}%
\pgfpathlineto{\pgfqpoint{1.297759in}{1.089024in}}%
\pgfpathlineto{\pgfqpoint{1.281601in}{1.095405in}}%
\pgfpathlineto{\pgfqpoint{1.281405in}{1.071711in}}%
\pgfpathclose%
\pgfusepath{fill}%
\end{pgfscope}%
\begin{pgfscope}%
\pgfpathrectangle{\pgfqpoint{0.407592in}{0.222904in}}{\pgfqpoint{1.411163in}{1.411163in}}%
\pgfusepath{clip}%
\pgfsetbuttcap%
\pgfsetroundjoin%
\definecolor{currentfill}{rgb}{0.168627,0.611765,0.411765}%
\pgfsetfillcolor{currentfill}%
\pgfsetfillopacity{0.180000}%
\pgfsetlinewidth{0.000000pt}%
\definecolor{currentstroke}{rgb}{0.000000,0.000000,0.000000}%
\pgfsetstrokecolor{currentstroke}%
\pgfsetdash{}{0pt}%
\pgfpathmoveto{\pgfqpoint{1.264617in}{1.007834in}}%
\pgfpathlineto{\pgfqpoint{1.280823in}{1.001001in}}%
\pgfpathlineto{\pgfqpoint{1.281016in}{1.024509in}}%
\pgfpathlineto{\pgfqpoint{1.264790in}{1.031450in}}%
\pgfpathlineto{\pgfqpoint{1.264617in}{1.007834in}}%
\pgfpathclose%
\pgfusepath{fill}%
\end{pgfscope}%
\begin{pgfscope}%
\pgfpathrectangle{\pgfqpoint{0.407592in}{0.222904in}}{\pgfqpoint{1.411163in}{1.411163in}}%
\pgfusepath{clip}%
\pgfsetbuttcap%
\pgfsetroundjoin%
\definecolor{currentfill}{rgb}{0.168627,0.611765,0.411765}%
\pgfsetfillcolor{currentfill}%
\pgfsetfillopacity{0.180000}%
\pgfsetlinewidth{0.000000pt}%
\definecolor{currentstroke}{rgb}{0.000000,0.000000,0.000000}%
\pgfsetstrokecolor{currentstroke}%
\pgfsetdash{}{0pt}%
\pgfpathmoveto{\pgfqpoint{1.247824in}{0.944669in}}%
\pgfpathlineto{\pgfqpoint{1.264099in}{0.937353in}}%
\pgfpathlineto{\pgfqpoint{1.264271in}{0.960785in}}%
\pgfpathlineto{\pgfqpoint{1.247975in}{0.968211in}}%
\pgfpathlineto{\pgfqpoint{1.247824in}{0.944669in}}%
\pgfpathclose%
\pgfusepath{fill}%
\end{pgfscope}%
\begin{pgfscope}%
\pgfpathrectangle{\pgfqpoint{0.407592in}{0.222904in}}{\pgfqpoint{1.411163in}{1.411163in}}%
\pgfusepath{clip}%
\pgfsetbuttcap%
\pgfsetroundjoin%
\definecolor{currentfill}{rgb}{0.168627,0.611765,0.411765}%
\pgfsetfillcolor{currentfill}%
\pgfsetfillopacity{0.180000}%
\pgfsetlinewidth{0.000000pt}%
\definecolor{currentstroke}{rgb}{0.000000,0.000000,0.000000}%
\pgfsetstrokecolor{currentstroke}%
\pgfsetdash{}{0pt}%
\pgfpathmoveto{\pgfqpoint{1.298409in}{1.160156in}}%
\pgfpathlineto{\pgfqpoint{1.314504in}{1.154417in}}%
\pgfpathlineto{\pgfqpoint{1.314742in}{1.178147in}}%
\pgfpathlineto{\pgfqpoint{1.298627in}{1.183991in}}%
\pgfpathlineto{\pgfqpoint{1.298409in}{1.160156in}}%
\pgfpathclose%
\pgfusepath{fill}%
\end{pgfscope}%
\begin{pgfscope}%
\pgfpathrectangle{\pgfqpoint{0.407592in}{0.222904in}}{\pgfqpoint{1.411163in}{1.411163in}}%
\pgfusepath{clip}%
\pgfsetbuttcap%
\pgfsetroundjoin%
\definecolor{currentfill}{rgb}{0.168627,0.611765,0.411765}%
\pgfsetfillcolor{currentfill}%
\pgfsetfillopacity{0.180000}%
\pgfsetlinewidth{0.000000pt}%
\definecolor{currentstroke}{rgb}{0.000000,0.000000,0.000000}%
\pgfsetstrokecolor{currentstroke}%
\pgfsetdash{}{0pt}%
\pgfpathmoveto{\pgfqpoint{1.281601in}{1.095405in}}%
\pgfpathlineto{\pgfqpoint{1.297759in}{1.089024in}}%
\pgfpathlineto{\pgfqpoint{1.297975in}{1.112673in}}%
\pgfpathlineto{\pgfqpoint{1.281797in}{1.119160in}}%
\pgfpathlineto{\pgfqpoint{1.281601in}{1.095405in}}%
\pgfpathclose%
\pgfusepath{fill}%
\end{pgfscope}%
\begin{pgfscope}%
\pgfpathrectangle{\pgfqpoint{0.407592in}{0.222904in}}{\pgfqpoint{1.411163in}{1.411163in}}%
\pgfusepath{clip}%
\pgfsetbuttcap%
\pgfsetroundjoin%
\definecolor{currentfill}{rgb}{0.168627,0.611765,0.411765}%
\pgfsetfillcolor{currentfill}%
\pgfsetfillopacity{0.180000}%
\pgfsetlinewidth{0.000000pt}%
\definecolor{currentstroke}{rgb}{0.000000,0.000000,0.000000}%
\pgfsetstrokecolor{currentstroke}%
\pgfsetdash{}{0pt}%
\pgfpathmoveto{\pgfqpoint{1.264790in}{1.031450in}}%
\pgfpathlineto{\pgfqpoint{1.281016in}{1.024509in}}%
\pgfpathlineto{\pgfqpoint{1.281211in}{1.048079in}}%
\pgfpathlineto{\pgfqpoint{1.264964in}{1.055128in}}%
\pgfpathlineto{\pgfqpoint{1.264790in}{1.031450in}}%
\pgfpathclose%
\pgfusepath{fill}%
\end{pgfscope}%
\begin{pgfscope}%
\pgfpathrectangle{\pgfqpoint{0.407592in}{0.222904in}}{\pgfqpoint{1.411163in}{1.411163in}}%
\pgfusepath{clip}%
\pgfsetbuttcap%
\pgfsetroundjoin%
\definecolor{currentfill}{rgb}{0.168627,0.611765,0.411765}%
\pgfsetfillcolor{currentfill}%
\pgfsetfillopacity{0.180000}%
\pgfsetlinewidth{0.000000pt}%
\definecolor{currentstroke}{rgb}{0.000000,0.000000,0.000000}%
\pgfsetstrokecolor{currentstroke}%
\pgfsetdash{}{0pt}%
\pgfpathmoveto{\pgfqpoint{1.247975in}{0.968211in}}%
\pgfpathlineto{\pgfqpoint{1.264271in}{0.960785in}}%
\pgfpathlineto{\pgfqpoint{1.264444in}{0.984279in}}%
\pgfpathlineto{\pgfqpoint{1.248127in}{0.991814in}}%
\pgfpathlineto{\pgfqpoint{1.247975in}{0.968211in}}%
\pgfpathclose%
\pgfusepath{fill}%
\end{pgfscope}%
\begin{pgfscope}%
\pgfpathrectangle{\pgfqpoint{0.407592in}{0.222904in}}{\pgfqpoint{1.411163in}{1.411163in}}%
\pgfusepath{clip}%
\pgfsetbuttcap%
\pgfsetroundjoin%
\definecolor{currentfill}{rgb}{0.168627,0.611765,0.411765}%
\pgfsetfillcolor{currentfill}%
\pgfsetfillopacity{0.180000}%
\pgfsetlinewidth{0.000000pt}%
\definecolor{currentstroke}{rgb}{0.000000,0.000000,0.000000}%
\pgfsetstrokecolor{currentstroke}%
\pgfsetdash{}{0pt}%
\pgfpathmoveto{\pgfqpoint{1.281797in}{1.119160in}}%
\pgfpathlineto{\pgfqpoint{1.297975in}{1.112673in}}%
\pgfpathlineto{\pgfqpoint{1.298192in}{1.136383in}}%
\pgfpathlineto{\pgfqpoint{1.281993in}{1.142978in}}%
\pgfpathlineto{\pgfqpoint{1.281797in}{1.119160in}}%
\pgfpathclose%
\pgfusepath{fill}%
\end{pgfscope}%
\begin{pgfscope}%
\pgfpathrectangle{\pgfqpoint{0.407592in}{0.222904in}}{\pgfqpoint{1.411163in}{1.411163in}}%
\pgfusepath{clip}%
\pgfsetbuttcap%
\pgfsetroundjoin%
\definecolor{currentfill}{rgb}{0.168627,0.611765,0.411765}%
\pgfsetfillcolor{currentfill}%
\pgfsetfillopacity{0.180000}%
\pgfsetlinewidth{0.000000pt}%
\definecolor{currentstroke}{rgb}{0.000000,0.000000,0.000000}%
\pgfsetstrokecolor{currentstroke}%
\pgfsetdash{}{0pt}%
\pgfpathmoveto{\pgfqpoint{1.264964in}{1.055128in}}%
\pgfpathlineto{\pgfqpoint{1.281211in}{1.048079in}}%
\pgfpathlineto{\pgfqpoint{1.281405in}{1.071711in}}%
\pgfpathlineto{\pgfqpoint{1.265138in}{1.078869in}}%
\pgfpathlineto{\pgfqpoint{1.264964in}{1.055128in}}%
\pgfpathclose%
\pgfusepath{fill}%
\end{pgfscope}%
\begin{pgfscope}%
\pgfpathrectangle{\pgfqpoint{0.407592in}{0.222904in}}{\pgfqpoint{1.411163in}{1.411163in}}%
\pgfusepath{clip}%
\pgfsetbuttcap%
\pgfsetroundjoin%
\definecolor{currentfill}{rgb}{0.168627,0.611765,0.411765}%
\pgfsetfillcolor{currentfill}%
\pgfsetfillopacity{0.180000}%
\pgfsetlinewidth{0.000000pt}%
\definecolor{currentstroke}{rgb}{0.000000,0.000000,0.000000}%
\pgfsetstrokecolor{currentstroke}%
\pgfsetdash{}{0pt}%
\pgfpathmoveto{\pgfqpoint{1.248127in}{0.991814in}}%
\pgfpathlineto{\pgfqpoint{1.264444in}{0.984279in}}%
\pgfpathlineto{\pgfqpoint{1.264617in}{1.007834in}}%
\pgfpathlineto{\pgfqpoint{1.248279in}{1.015479in}}%
\pgfpathlineto{\pgfqpoint{1.248127in}{0.991814in}}%
\pgfpathclose%
\pgfusepath{fill}%
\end{pgfscope}%
\begin{pgfscope}%
\pgfpathrectangle{\pgfqpoint{0.407592in}{0.222904in}}{\pgfqpoint{1.411163in}{1.411163in}}%
\pgfusepath{clip}%
\pgfsetbuttcap%
\pgfsetroundjoin%
\definecolor{currentfill}{rgb}{0.168627,0.611765,0.411765}%
\pgfsetfillcolor{currentfill}%
\pgfsetfillopacity{0.180000}%
\pgfsetlinewidth{0.000000pt}%
\definecolor{currentstroke}{rgb}{0.000000,0.000000,0.000000}%
\pgfsetstrokecolor{currentstroke}%
\pgfsetdash{}{0pt}%
\pgfpathmoveto{\pgfqpoint{1.281993in}{1.142978in}}%
\pgfpathlineto{\pgfqpoint{1.298192in}{1.136383in}}%
\pgfpathlineto{\pgfqpoint{1.298409in}{1.160156in}}%
\pgfpathlineto{\pgfqpoint{1.282190in}{1.166859in}}%
\pgfpathlineto{\pgfqpoint{1.281993in}{1.142978in}}%
\pgfpathclose%
\pgfusepath{fill}%
\end{pgfscope}%
\begin{pgfscope}%
\pgfpathrectangle{\pgfqpoint{0.407592in}{0.222904in}}{\pgfqpoint{1.411163in}{1.411163in}}%
\pgfusepath{clip}%
\pgfsetbuttcap%
\pgfsetroundjoin%
\definecolor{currentfill}{rgb}{0.168627,0.611765,0.411765}%
\pgfsetfillcolor{currentfill}%
\pgfsetfillopacity{0.180000}%
\pgfsetlinewidth{0.000000pt}%
\definecolor{currentstroke}{rgb}{0.000000,0.000000,0.000000}%
\pgfsetstrokecolor{currentstroke}%
\pgfsetdash{}{0pt}%
\pgfpathmoveto{\pgfqpoint{1.265138in}{1.078869in}}%
\pgfpathlineto{\pgfqpoint{1.281405in}{1.071711in}}%
\pgfpathlineto{\pgfqpoint{1.281601in}{1.095405in}}%
\pgfpathlineto{\pgfqpoint{1.265313in}{1.102672in}}%
\pgfpathlineto{\pgfqpoint{1.265138in}{1.078869in}}%
\pgfpathclose%
\pgfusepath{fill}%
\end{pgfscope}%
\begin{pgfscope}%
\pgfpathrectangle{\pgfqpoint{0.407592in}{0.222904in}}{\pgfqpoint{1.411163in}{1.411163in}}%
\pgfusepath{clip}%
\pgfsetbuttcap%
\pgfsetroundjoin%
\definecolor{currentfill}{rgb}{0.168627,0.611765,0.411765}%
\pgfsetfillcolor{currentfill}%
\pgfsetfillopacity{0.180000}%
\pgfsetlinewidth{0.000000pt}%
\definecolor{currentstroke}{rgb}{0.000000,0.000000,0.000000}%
\pgfsetstrokecolor{currentstroke}%
\pgfsetdash{}{0pt}%
\pgfpathmoveto{\pgfqpoint{1.248279in}{1.015479in}}%
\pgfpathlineto{\pgfqpoint{1.264617in}{1.007834in}}%
\pgfpathlineto{\pgfqpoint{1.264790in}{1.031450in}}%
\pgfpathlineto{\pgfqpoint{1.248432in}{1.039206in}}%
\pgfpathlineto{\pgfqpoint{1.248279in}{1.015479in}}%
\pgfpathclose%
\pgfusepath{fill}%
\end{pgfscope}%
\begin{pgfscope}%
\pgfpathrectangle{\pgfqpoint{0.407592in}{0.222904in}}{\pgfqpoint{1.411163in}{1.411163in}}%
\pgfusepath{clip}%
\pgfsetbuttcap%
\pgfsetroundjoin%
\definecolor{currentfill}{rgb}{0.168627,0.611765,0.411765}%
\pgfsetfillcolor{currentfill}%
\pgfsetfillopacity{0.180000}%
\pgfsetlinewidth{0.000000pt}%
\definecolor{currentstroke}{rgb}{0.000000,0.000000,0.000000}%
\pgfsetstrokecolor{currentstroke}%
\pgfsetdash{}{0pt}%
\pgfpathmoveto{\pgfqpoint{1.282190in}{1.166859in}}%
\pgfpathlineto{\pgfqpoint{1.298409in}{1.160156in}}%
\pgfpathlineto{\pgfqpoint{1.298627in}{1.183991in}}%
\pgfpathlineto{\pgfqpoint{1.282387in}{1.190802in}}%
\pgfpathlineto{\pgfqpoint{1.282190in}{1.166859in}}%
\pgfpathclose%
\pgfusepath{fill}%
\end{pgfscope}%
\begin{pgfscope}%
\pgfpathrectangle{\pgfqpoint{0.407592in}{0.222904in}}{\pgfqpoint{1.411163in}{1.411163in}}%
\pgfusepath{clip}%
\pgfsetbuttcap%
\pgfsetroundjoin%
\definecolor{currentfill}{rgb}{0.168627,0.611765,0.411765}%
\pgfsetfillcolor{currentfill}%
\pgfsetfillopacity{0.180000}%
\pgfsetlinewidth{0.000000pt}%
\definecolor{currentstroke}{rgb}{0.000000,0.000000,0.000000}%
\pgfsetstrokecolor{currentstroke}%
\pgfsetdash{}{0pt}%
\pgfpathmoveto{\pgfqpoint{1.265313in}{1.102672in}}%
\pgfpathlineto{\pgfqpoint{1.281601in}{1.095405in}}%
\pgfpathlineto{\pgfqpoint{1.281797in}{1.119160in}}%
\pgfpathlineto{\pgfqpoint{1.265488in}{1.126538in}}%
\pgfpathlineto{\pgfqpoint{1.265313in}{1.102672in}}%
\pgfpathclose%
\pgfusepath{fill}%
\end{pgfscope}%
\begin{pgfscope}%
\pgfpathrectangle{\pgfqpoint{0.407592in}{0.222904in}}{\pgfqpoint{1.411163in}{1.411163in}}%
\pgfusepath{clip}%
\pgfsetbuttcap%
\pgfsetroundjoin%
\definecolor{currentfill}{rgb}{0.168627,0.611765,0.411765}%
\pgfsetfillcolor{currentfill}%
\pgfsetfillopacity{0.180000}%
\pgfsetlinewidth{0.000000pt}%
\definecolor{currentstroke}{rgb}{0.000000,0.000000,0.000000}%
\pgfsetstrokecolor{currentstroke}%
\pgfsetdash{}{0pt}%
\pgfpathmoveto{\pgfqpoint{1.248432in}{1.039206in}}%
\pgfpathlineto{\pgfqpoint{1.264790in}{1.031450in}}%
\pgfpathlineto{\pgfqpoint{1.264964in}{1.055128in}}%
\pgfpathlineto{\pgfqpoint{1.248585in}{1.062996in}}%
\pgfpathlineto{\pgfqpoint{1.248432in}{1.039206in}}%
\pgfpathclose%
\pgfusepath{fill}%
\end{pgfscope}%
\begin{pgfscope}%
\pgfpathrectangle{\pgfqpoint{0.407592in}{0.222904in}}{\pgfqpoint{1.411163in}{1.411163in}}%
\pgfusepath{clip}%
\pgfsetbuttcap%
\pgfsetroundjoin%
\definecolor{currentfill}{rgb}{0.168627,0.611765,0.411765}%
\pgfsetfillcolor{currentfill}%
\pgfsetfillopacity{0.180000}%
\pgfsetlinewidth{0.000000pt}%
\definecolor{currentstroke}{rgb}{0.000000,0.000000,0.000000}%
\pgfsetstrokecolor{currentstroke}%
\pgfsetdash{}{0pt}%
\pgfpathmoveto{\pgfqpoint{1.265488in}{1.126538in}}%
\pgfpathlineto{\pgfqpoint{1.281797in}{1.119160in}}%
\pgfpathlineto{\pgfqpoint{1.281993in}{1.142978in}}%
\pgfpathlineto{\pgfqpoint{1.265664in}{1.150467in}}%
\pgfpathlineto{\pgfqpoint{1.265488in}{1.126538in}}%
\pgfpathclose%
\pgfusepath{fill}%
\end{pgfscope}%
\begin{pgfscope}%
\pgfpathrectangle{\pgfqpoint{0.407592in}{0.222904in}}{\pgfqpoint{1.411163in}{1.411163in}}%
\pgfusepath{clip}%
\pgfsetbuttcap%
\pgfsetroundjoin%
\definecolor{currentfill}{rgb}{0.168627,0.611765,0.411765}%
\pgfsetfillcolor{currentfill}%
\pgfsetfillopacity{0.180000}%
\pgfsetlinewidth{0.000000pt}%
\definecolor{currentstroke}{rgb}{0.000000,0.000000,0.000000}%
\pgfsetstrokecolor{currentstroke}%
\pgfsetdash{}{0pt}%
\pgfpathmoveto{\pgfqpoint{1.248585in}{1.062996in}}%
\pgfpathlineto{\pgfqpoint{1.264964in}{1.055128in}}%
\pgfpathlineto{\pgfqpoint{1.265138in}{1.078869in}}%
\pgfpathlineto{\pgfqpoint{1.248738in}{1.086849in}}%
\pgfpathlineto{\pgfqpoint{1.248585in}{1.062996in}}%
\pgfpathclose%
\pgfusepath{fill}%
\end{pgfscope}%
\begin{pgfscope}%
\pgfpathrectangle{\pgfqpoint{0.407592in}{0.222904in}}{\pgfqpoint{1.411163in}{1.411163in}}%
\pgfusepath{clip}%
\pgfsetbuttcap%
\pgfsetroundjoin%
\definecolor{currentfill}{rgb}{0.168627,0.611765,0.411765}%
\pgfsetfillcolor{currentfill}%
\pgfsetfillopacity{0.180000}%
\pgfsetlinewidth{0.000000pt}%
\definecolor{currentstroke}{rgb}{0.000000,0.000000,0.000000}%
\pgfsetstrokecolor{currentstroke}%
\pgfsetdash{}{0pt}%
\pgfpathmoveto{\pgfqpoint{1.265664in}{1.150467in}}%
\pgfpathlineto{\pgfqpoint{1.281993in}{1.142978in}}%
\pgfpathlineto{\pgfqpoint{1.282190in}{1.166859in}}%
\pgfpathlineto{\pgfqpoint{1.265840in}{1.174459in}}%
\pgfpathlineto{\pgfqpoint{1.265664in}{1.150467in}}%
\pgfpathclose%
\pgfusepath{fill}%
\end{pgfscope}%
\begin{pgfscope}%
\pgfpathrectangle{\pgfqpoint{0.407592in}{0.222904in}}{\pgfqpoint{1.411163in}{1.411163in}}%
\pgfusepath{clip}%
\pgfsetbuttcap%
\pgfsetroundjoin%
\definecolor{currentfill}{rgb}{0.168627,0.611765,0.411765}%
\pgfsetfillcolor{currentfill}%
\pgfsetfillopacity{0.180000}%
\pgfsetlinewidth{0.000000pt}%
\definecolor{currentstroke}{rgb}{0.000000,0.000000,0.000000}%
\pgfsetstrokecolor{currentstroke}%
\pgfsetdash{}{0pt}%
\pgfpathmoveto{\pgfqpoint{1.248738in}{1.086849in}}%
\pgfpathlineto{\pgfqpoint{1.265138in}{1.078869in}}%
\pgfpathlineto{\pgfqpoint{1.265313in}{1.102672in}}%
\pgfpathlineto{\pgfqpoint{1.248892in}{1.110765in}}%
\pgfpathlineto{\pgfqpoint{1.248738in}{1.086849in}}%
\pgfpathclose%
\pgfusepath{fill}%
\end{pgfscope}%
\begin{pgfscope}%
\pgfpathrectangle{\pgfqpoint{0.407592in}{0.222904in}}{\pgfqpoint{1.411163in}{1.411163in}}%
\pgfusepath{clip}%
\pgfsetbuttcap%
\pgfsetroundjoin%
\definecolor{currentfill}{rgb}{0.168627,0.611765,0.411765}%
\pgfsetfillcolor{currentfill}%
\pgfsetfillopacity{0.180000}%
\pgfsetlinewidth{0.000000pt}%
\definecolor{currentstroke}{rgb}{0.000000,0.000000,0.000000}%
\pgfsetstrokecolor{currentstroke}%
\pgfsetdash{}{0pt}%
\pgfpathmoveto{\pgfqpoint{1.265840in}{1.174459in}}%
\pgfpathlineto{\pgfqpoint{1.282190in}{1.166859in}}%
\pgfpathlineto{\pgfqpoint{1.282387in}{1.190802in}}%
\pgfpathlineto{\pgfqpoint{1.266017in}{1.198514in}}%
\pgfpathlineto{\pgfqpoint{1.265840in}{1.174459in}}%
\pgfpathclose%
\pgfusepath{fill}%
\end{pgfscope}%
\begin{pgfscope}%
\pgfpathrectangle{\pgfqpoint{0.407592in}{0.222904in}}{\pgfqpoint{1.411163in}{1.411163in}}%
\pgfusepath{clip}%
\pgfsetbuttcap%
\pgfsetroundjoin%
\definecolor{currentfill}{rgb}{0.168627,0.611765,0.411765}%
\pgfsetfillcolor{currentfill}%
\pgfsetfillopacity{0.180000}%
\pgfsetlinewidth{0.000000pt}%
\definecolor{currentstroke}{rgb}{0.000000,0.000000,0.000000}%
\pgfsetstrokecolor{currentstroke}%
\pgfsetdash{}{0pt}%
\pgfpathmoveto{\pgfqpoint{1.248892in}{1.110765in}}%
\pgfpathlineto{\pgfqpoint{1.265313in}{1.102672in}}%
\pgfpathlineto{\pgfqpoint{1.265488in}{1.126538in}}%
\pgfpathlineto{\pgfqpoint{1.249047in}{1.134743in}}%
\pgfpathlineto{\pgfqpoint{1.248892in}{1.110765in}}%
\pgfpathclose%
\pgfusepath{fill}%
\end{pgfscope}%
\begin{pgfscope}%
\pgfpathrectangle{\pgfqpoint{0.407592in}{0.222904in}}{\pgfqpoint{1.411163in}{1.411163in}}%
\pgfusepath{clip}%
\pgfsetbuttcap%
\pgfsetroundjoin%
\definecolor{currentfill}{rgb}{0.168627,0.611765,0.411765}%
\pgfsetfillcolor{currentfill}%
\pgfsetfillopacity{0.180000}%
\pgfsetlinewidth{0.000000pt}%
\definecolor{currentstroke}{rgb}{0.000000,0.000000,0.000000}%
\pgfsetstrokecolor{currentstroke}%
\pgfsetdash{}{0pt}%
\pgfpathmoveto{\pgfqpoint{1.249047in}{1.134743in}}%
\pgfpathlineto{\pgfqpoint{1.265488in}{1.126538in}}%
\pgfpathlineto{\pgfqpoint{1.265664in}{1.150467in}}%
\pgfpathlineto{\pgfqpoint{1.249201in}{1.158786in}}%
\pgfpathlineto{\pgfqpoint{1.249047in}{1.134743in}}%
\pgfpathclose%
\pgfusepath{fill}%
\end{pgfscope}%
\begin{pgfscope}%
\pgfpathrectangle{\pgfqpoint{0.407592in}{0.222904in}}{\pgfqpoint{1.411163in}{1.411163in}}%
\pgfusepath{clip}%
\pgfsetbuttcap%
\pgfsetroundjoin%
\definecolor{currentfill}{rgb}{0.168627,0.611765,0.411765}%
\pgfsetfillcolor{currentfill}%
\pgfsetfillopacity{0.180000}%
\pgfsetlinewidth{0.000000pt}%
\definecolor{currentstroke}{rgb}{0.000000,0.000000,0.000000}%
\pgfsetstrokecolor{currentstroke}%
\pgfsetdash{}{0pt}%
\pgfpathmoveto{\pgfqpoint{1.249201in}{1.158786in}}%
\pgfpathlineto{\pgfqpoint{1.265664in}{1.150467in}}%
\pgfpathlineto{\pgfqpoint{1.265840in}{1.174459in}}%
\pgfpathlineto{\pgfqpoint{1.249356in}{1.182892in}}%
\pgfpathlineto{\pgfqpoint{1.249201in}{1.158786in}}%
\pgfpathclose%
\pgfusepath{fill}%
\end{pgfscope}%
\begin{pgfscope}%
\pgfpathrectangle{\pgfqpoint{0.407592in}{0.222904in}}{\pgfqpoint{1.411163in}{1.411163in}}%
\pgfusepath{clip}%
\pgfsetbuttcap%
\pgfsetroundjoin%
\definecolor{currentfill}{rgb}{0.168627,0.611765,0.411765}%
\pgfsetfillcolor{currentfill}%
\pgfsetfillopacity{0.180000}%
\pgfsetlinewidth{0.000000pt}%
\definecolor{currentstroke}{rgb}{0.000000,0.000000,0.000000}%
\pgfsetstrokecolor{currentstroke}%
\pgfsetdash{}{0pt}%
\pgfpathmoveto{\pgfqpoint{1.249356in}{1.182892in}}%
\pgfpathlineto{\pgfqpoint{1.265840in}{1.174459in}}%
\pgfpathlineto{\pgfqpoint{1.266017in}{1.198514in}}%
\pgfpathlineto{\pgfqpoint{1.249512in}{1.207062in}}%
\pgfpathlineto{\pgfqpoint{1.249356in}{1.182892in}}%
\pgfpathclose%
\pgfusepath{fill}%
\end{pgfscope}%
\begin{pgfscope}%
\pgfpathrectangle{\pgfqpoint{0.407592in}{0.222904in}}{\pgfqpoint{1.411163in}{1.411163in}}%
\pgfusepath{clip}%
\pgfsetbuttcap%
\pgfsetroundjoin%
\definecolor{currentfill}{rgb}{0.556863,0.654902,0.847059}%
\pgfsetfillcolor{currentfill}%
\pgfsetfillopacity{0.130000}%
\pgfsetlinewidth{0.000000pt}%
\definecolor{currentstroke}{rgb}{0.000000,0.000000,0.000000}%
\pgfsetstrokecolor{currentstroke}%
\pgfsetdash{}{0pt}%
\pgfpathmoveto{\pgfqpoint{0.871478in}{1.394437in}}%
\pgfpathlineto{\pgfqpoint{0.886848in}{1.378173in}}%
\pgfpathlineto{\pgfqpoint{0.902230in}{1.362091in}}%
\pgfpathlineto{\pgfqpoint{0.910091in}{1.373743in}}%
\pgfpathlineto{\pgfqpoint{0.894724in}{1.389630in}}%
\pgfpathlineto{\pgfqpoint{0.879369in}{1.405704in}}%
\pgfusepath{fill}%
\end{pgfscope}%
\begin{pgfscope}%
\pgfpathrectangle{\pgfqpoint{0.407592in}{0.222904in}}{\pgfqpoint{1.411163in}{1.411163in}}%
\pgfusepath{clip}%
\pgfsetbuttcap%
\pgfsetroundjoin%
\definecolor{currentfill}{rgb}{0.556863,0.654902,0.847059}%
\pgfsetfillcolor{currentfill}%
\pgfsetfillopacity{0.130000}%
\pgfsetlinewidth{0.000000pt}%
\definecolor{currentstroke}{rgb}{0.000000,0.000000,0.000000}%
\pgfsetstrokecolor{currentstroke}%
\pgfsetdash{}{0pt}%
\pgfpathmoveto{\pgfqpoint{0.902230in}{1.362091in}}%
\pgfpathlineto{\pgfqpoint{0.917624in}{1.346205in}}%
\pgfpathlineto{\pgfqpoint{0.933031in}{1.330531in}}%
\pgfpathlineto{\pgfqpoint{0.940863in}{1.342588in}}%
\pgfpathlineto{\pgfqpoint{0.925470in}{1.358057in}}%
\pgfpathlineto{\pgfqpoint{0.910091in}{1.373743in}}%
\pgfusepath{fill}%
\end{pgfscope}%
\begin{pgfscope}%
\pgfpathrectangle{\pgfqpoint{0.407592in}{0.222904in}}{\pgfqpoint{1.411163in}{1.411163in}}%
\pgfusepath{clip}%
\pgfsetbuttcap%
\pgfsetroundjoin%
\definecolor{currentfill}{rgb}{0.556863,0.654902,0.847059}%
\pgfsetfillcolor{currentfill}%
\pgfsetfillopacity{0.130000}%
\pgfsetlinewidth{0.000000pt}%
\definecolor{currentstroke}{rgb}{0.000000,0.000000,0.000000}%
\pgfsetstrokecolor{currentstroke}%
\pgfsetdash{}{0pt}%
\pgfpathmoveto{\pgfqpoint{0.855637in}{1.372404in}}%
\pgfpathlineto{\pgfqpoint{0.871038in}{1.355759in}}%
\pgfpathlineto{\pgfqpoint{0.886451in}{1.339285in}}%
\pgfpathlineto{\pgfqpoint{0.894350in}{1.350602in}}%
\pgfpathlineto{\pgfqpoint{0.902230in}{1.362091in}}%
\pgfpathlineto{\pgfqpoint{0.886848in}{1.378173in}}%
\pgfpathlineto{\pgfqpoint{0.871478in}{1.394437in}}%
\pgfpathlineto{\pgfqpoint{0.863568in}{1.383335in}}%
\pgfpathlineto{\pgfqpoint{0.855637in}{1.372404in}}%
\pgfpathclose%
\pgfusepath{fill}%
\end{pgfscope}%
\begin{pgfscope}%
\pgfpathrectangle{\pgfqpoint{0.407592in}{0.222904in}}{\pgfqpoint{1.411163in}{1.411163in}}%
\pgfusepath{clip}%
\pgfsetbuttcap%
\pgfsetroundjoin%
\definecolor{currentfill}{rgb}{0.556863,0.654902,0.847059}%
\pgfsetfillcolor{currentfill}%
\pgfsetfillopacity{0.130000}%
\pgfsetlinewidth{0.000000pt}%
\definecolor{currentstroke}{rgb}{0.000000,0.000000,0.000000}%
\pgfsetstrokecolor{currentstroke}%
\pgfsetdash{}{0pt}%
\pgfpathmoveto{\pgfqpoint{0.933031in}{1.330531in}}%
\pgfpathlineto{\pgfqpoint{0.948450in}{1.315085in}}%
\pgfpathlineto{\pgfqpoint{0.963882in}{1.299885in}}%
\pgfpathlineto{\pgfqpoint{0.971687in}{1.312368in}}%
\pgfpathlineto{\pgfqpoint{0.956268in}{1.327353in}}%
\pgfpathlineto{\pgfqpoint{0.940863in}{1.342588in}}%
\pgfusepath{fill}%
\end{pgfscope}%
\begin{pgfscope}%
\pgfpathrectangle{\pgfqpoint{0.407592in}{0.222904in}}{\pgfqpoint{1.411163in}{1.411163in}}%
\pgfusepath{clip}%
\pgfsetbuttcap%
\pgfsetroundjoin%
\definecolor{currentfill}{rgb}{0.556863,0.654902,0.847059}%
\pgfsetfillcolor{currentfill}%
\pgfsetfillopacity{0.130000}%
\pgfsetlinewidth{0.000000pt}%
\definecolor{currentstroke}{rgb}{0.000000,0.000000,0.000000}%
\pgfsetstrokecolor{currentstroke}%
\pgfsetdash{}{0pt}%
\pgfpathmoveto{\pgfqpoint{0.886451in}{1.339285in}}%
\pgfpathlineto{\pgfqpoint{0.901875in}{1.322996in}}%
\pgfpathlineto{\pgfqpoint{0.917311in}{1.306908in}}%
\pgfpathlineto{\pgfqpoint{0.925180in}{1.318634in}}%
\pgfpathlineto{\pgfqpoint{0.933031in}{1.330531in}}%
\pgfpathlineto{\pgfqpoint{0.917624in}{1.346205in}}%
\pgfpathlineto{\pgfqpoint{0.902230in}{1.362091in}}%
\pgfpathlineto{\pgfqpoint{0.894350in}{1.350602in}}%
\pgfpathlineto{\pgfqpoint{0.886451in}{1.339285in}}%
\pgfpathclose%
\pgfusepath{fill}%
\end{pgfscope}%
\begin{pgfscope}%
\pgfpathrectangle{\pgfqpoint{0.407592in}{0.222904in}}{\pgfqpoint{1.411163in}{1.411163in}}%
\pgfusepath{clip}%
\pgfsetbuttcap%
\pgfsetroundjoin%
\definecolor{currentfill}{rgb}{0.556863,0.654902,0.847059}%
\pgfsetfillcolor{currentfill}%
\pgfsetfillopacity{0.130000}%
\pgfsetlinewidth{0.000000pt}%
\definecolor{currentstroke}{rgb}{0.000000,0.000000,0.000000}%
\pgfsetstrokecolor{currentstroke}%
\pgfsetdash{}{0pt}%
\pgfpathmoveto{\pgfqpoint{0.963882in}{1.299885in}}%
\pgfpathlineto{\pgfqpoint{0.979329in}{1.284951in}}%
\pgfpathlineto{\pgfqpoint{0.994790in}{1.270304in}}%
\pgfpathlineto{\pgfqpoint{1.002568in}{1.283228in}}%
\pgfpathlineto{\pgfqpoint{0.987120in}{1.297653in}}%
\pgfpathlineto{\pgfqpoint{0.971687in}{1.312368in}}%
\pgfusepath{fill}%
\end{pgfscope}%
\begin{pgfscope}%
\pgfpathrectangle{\pgfqpoint{0.407592in}{0.222904in}}{\pgfqpoint{1.411163in}{1.411163in}}%
\pgfusepath{clip}%
\pgfsetbuttcap%
\pgfsetroundjoin%
\definecolor{currentfill}{rgb}{0.556863,0.654902,0.847059}%
\pgfsetfillcolor{currentfill}%
\pgfsetfillopacity{0.130000}%
\pgfsetlinewidth{0.000000pt}%
\definecolor{currentstroke}{rgb}{0.000000,0.000000,0.000000}%
\pgfsetstrokecolor{currentstroke}%
\pgfsetdash{}{0pt}%
\pgfpathmoveto{\pgfqpoint{0.917311in}{1.306908in}}%
\pgfpathlineto{\pgfqpoint{0.932759in}{1.291037in}}%
\pgfpathlineto{\pgfqpoint{0.948220in}{1.275402in}}%
\pgfpathlineto{\pgfqpoint{0.956060in}{1.287560in}}%
\pgfpathlineto{\pgfqpoint{0.963882in}{1.299885in}}%
\pgfpathlineto{\pgfqpoint{0.948450in}{1.315085in}}%
\pgfpathlineto{\pgfqpoint{0.933031in}{1.330531in}}%
\pgfpathlineto{\pgfqpoint{0.925180in}{1.318634in}}%
\pgfpathlineto{\pgfqpoint{0.917311in}{1.306908in}}%
\pgfpathclose%
\pgfusepath{fill}%
\end{pgfscope}%
\begin{pgfscope}%
\pgfpathrectangle{\pgfqpoint{0.407592in}{0.222904in}}{\pgfqpoint{1.411163in}{1.411163in}}%
\pgfusepath{clip}%
\pgfsetbuttcap%
\pgfsetroundjoin%
\definecolor{currentfill}{rgb}{0.556863,0.654902,0.847059}%
\pgfsetfillcolor{currentfill}%
\pgfsetfillopacity{0.130000}%
\pgfsetlinewidth{0.000000pt}%
\definecolor{currentstroke}{rgb}{0.000000,0.000000,0.000000}%
\pgfsetstrokecolor{currentstroke}%
\pgfsetdash{}{0pt}%
\pgfpathmoveto{\pgfqpoint{0.839713in}{1.351088in}}%
\pgfpathlineto{\pgfqpoint{0.855147in}{1.334064in}}%
\pgfpathlineto{\pgfqpoint{0.870591in}{1.317199in}}%
\pgfpathlineto{\pgfqpoint{0.878531in}{1.328148in}}%
\pgfpathlineto{\pgfqpoint{0.886451in}{1.339285in}}%
\pgfpathlineto{\pgfqpoint{0.871038in}{1.355759in}}%
\pgfpathlineto{\pgfqpoint{0.855637in}{1.372404in}}%
\pgfpathlineto{\pgfqpoint{0.847685in}{1.361653in}}%
\pgfpathlineto{\pgfqpoint{0.839713in}{1.351088in}}%
\pgfpathclose%
\pgfusepath{fill}%
\end{pgfscope}%
\begin{pgfscope}%
\pgfpathrectangle{\pgfqpoint{0.407592in}{0.222904in}}{\pgfqpoint{1.411163in}{1.411163in}}%
\pgfusepath{clip}%
\pgfsetbuttcap%
\pgfsetroundjoin%
\definecolor{currentfill}{rgb}{0.556863,0.654902,0.847059}%
\pgfsetfillcolor{currentfill}%
\pgfsetfillopacity{0.130000}%
\pgfsetlinewidth{0.000000pt}%
\definecolor{currentstroke}{rgb}{0.000000,0.000000,0.000000}%
\pgfsetstrokecolor{currentstroke}%
\pgfsetdash{}{0pt}%
\pgfpathmoveto{\pgfqpoint{0.994790in}{1.270304in}}%
\pgfpathlineto{\pgfqpoint{1.010266in}{1.255965in}}%
\pgfpathlineto{\pgfqpoint{1.025758in}{1.241959in}}%
\pgfpathlineto{\pgfqpoint{1.033512in}{1.255338in}}%
\pgfpathlineto{\pgfqpoint{1.018032in}{1.269115in}}%
\pgfpathlineto{\pgfqpoint{1.002568in}{1.283228in}}%
\pgfusepath{fill}%
\end{pgfscope}%
\begin{pgfscope}%
\pgfpathrectangle{\pgfqpoint{0.407592in}{0.222904in}}{\pgfqpoint{1.411163in}{1.411163in}}%
\pgfusepath{clip}%
\pgfsetbuttcap%
\pgfsetroundjoin%
\definecolor{currentfill}{rgb}{0.556863,0.654902,0.847059}%
\pgfsetfillcolor{currentfill}%
\pgfsetfillopacity{0.130000}%
\pgfsetlinewidth{0.000000pt}%
\definecolor{currentstroke}{rgb}{0.000000,0.000000,0.000000}%
\pgfsetstrokecolor{currentstroke}%
\pgfsetdash{}{0pt}%
\pgfpathmoveto{\pgfqpoint{0.948220in}{1.275402in}}%
\pgfpathlineto{\pgfqpoint{0.963694in}{1.260023in}}%
\pgfpathlineto{\pgfqpoint{0.979181in}{1.244921in}}%
\pgfpathlineto{\pgfqpoint{0.986994in}{1.257531in}}%
\pgfpathlineto{\pgfqpoint{0.994790in}{1.270304in}}%
\pgfpathlineto{\pgfqpoint{0.979329in}{1.284951in}}%
\pgfpathlineto{\pgfqpoint{0.963882in}{1.299885in}}%
\pgfpathlineto{\pgfqpoint{0.956060in}{1.287560in}}%
\pgfpathlineto{\pgfqpoint{0.948220in}{1.275402in}}%
\pgfpathclose%
\pgfusepath{fill}%
\end{pgfscope}%
\begin{pgfscope}%
\pgfpathrectangle{\pgfqpoint{0.407592in}{0.222904in}}{\pgfqpoint{1.411163in}{1.411163in}}%
\pgfusepath{clip}%
\pgfsetbuttcap%
\pgfsetroundjoin%
\definecolor{currentfill}{rgb}{0.556863,0.654902,0.847059}%
\pgfsetfillcolor{currentfill}%
\pgfsetfillopacity{0.130000}%
\pgfsetlinewidth{0.000000pt}%
\definecolor{currentstroke}{rgb}{0.000000,0.000000,0.000000}%
\pgfsetstrokecolor{currentstroke}%
\pgfsetdash{}{0pt}%
\pgfpathmoveto{\pgfqpoint{0.870591in}{1.317199in}}%
\pgfpathlineto{\pgfqpoint{0.886047in}{1.300507in}}%
\pgfpathlineto{\pgfqpoint{0.901514in}{1.284003in}}%
\pgfpathlineto{\pgfqpoint{0.909422in}{1.295361in}}%
\pgfpathlineto{\pgfqpoint{0.917311in}{1.306908in}}%
\pgfpathlineto{\pgfqpoint{0.901875in}{1.322996in}}%
\pgfpathlineto{\pgfqpoint{0.886451in}{1.339285in}}%
\pgfpathlineto{\pgfqpoint{0.878531in}{1.328148in}}%
\pgfpathlineto{\pgfqpoint{0.870591in}{1.317199in}}%
\pgfpathclose%
\pgfusepath{fill}%
\end{pgfscope}%
\begin{pgfscope}%
\pgfpathrectangle{\pgfqpoint{0.407592in}{0.222904in}}{\pgfqpoint{1.411163in}{1.411163in}}%
\pgfusepath{clip}%
\pgfsetbuttcap%
\pgfsetroundjoin%
\definecolor{currentfill}{rgb}{0.556863,0.654902,0.847059}%
\pgfsetfillcolor{currentfill}%
\pgfsetfillopacity{0.130000}%
\pgfsetlinewidth{0.000000pt}%
\definecolor{currentstroke}{rgb}{0.000000,0.000000,0.000000}%
\pgfsetstrokecolor{currentstroke}%
\pgfsetdash{}{0pt}%
\pgfpathmoveto{\pgfqpoint{1.025758in}{1.241959in}}%
\pgfpathlineto{\pgfqpoint{1.041266in}{1.228311in}}%
\pgfpathlineto{\pgfqpoint{1.056792in}{1.215048in}}%
\pgfpathlineto{\pgfqpoint{1.064524in}{1.228888in}}%
\pgfpathlineto{\pgfqpoint{1.049009in}{1.241920in}}%
\pgfpathlineto{\pgfqpoint{1.033512in}{1.255338in}}%
\pgfusepath{fill}%
\end{pgfscope}%
\begin{pgfscope}%
\pgfpathrectangle{\pgfqpoint{0.407592in}{0.222904in}}{\pgfqpoint{1.411163in}{1.411163in}}%
\pgfusepath{clip}%
\pgfsetbuttcap%
\pgfsetroundjoin%
\definecolor{currentfill}{rgb}{0.556863,0.654902,0.847059}%
\pgfsetfillcolor{currentfill}%
\pgfsetfillopacity{0.130000}%
\pgfsetlinewidth{0.000000pt}%
\definecolor{currentstroke}{rgb}{0.000000,0.000000,0.000000}%
\pgfsetstrokecolor{currentstroke}%
\pgfsetdash{}{0pt}%
\pgfpathmoveto{\pgfqpoint{0.901514in}{1.284003in}}%
\pgfpathlineto{\pgfqpoint{0.916992in}{1.267704in}}%
\pgfpathlineto{\pgfqpoint{0.932483in}{1.251628in}}%
\pgfpathlineto{\pgfqpoint{0.940361in}{1.263421in}}%
\pgfpathlineto{\pgfqpoint{0.948220in}{1.275402in}}%
\pgfpathlineto{\pgfqpoint{0.932759in}{1.291037in}}%
\pgfpathlineto{\pgfqpoint{0.917311in}{1.306908in}}%
\pgfpathlineto{\pgfqpoint{0.909422in}{1.295361in}}%
\pgfpathlineto{\pgfqpoint{0.901514in}{1.284003in}}%
\pgfpathclose%
\pgfusepath{fill}%
\end{pgfscope}%
\begin{pgfscope}%
\pgfpathrectangle{\pgfqpoint{0.407592in}{0.222904in}}{\pgfqpoint{1.411163in}{1.411163in}}%
\pgfusepath{clip}%
\pgfsetbuttcap%
\pgfsetroundjoin%
\definecolor{currentfill}{rgb}{0.556863,0.654902,0.847059}%
\pgfsetfillcolor{currentfill}%
\pgfsetfillopacity{0.130000}%
\pgfsetlinewidth{0.000000pt}%
\definecolor{currentstroke}{rgb}{0.000000,0.000000,0.000000}%
\pgfsetstrokecolor{currentstroke}%
\pgfsetdash{}{0pt}%
\pgfpathmoveto{\pgfqpoint{0.979181in}{1.244921in}}%
\pgfpathlineto{\pgfqpoint{0.994683in}{1.230119in}}%
\pgfpathlineto{\pgfqpoint{1.010200in}{1.215642in}}%
\pgfpathlineto{\pgfqpoint{1.017987in}{1.228723in}}%
\pgfpathlineto{\pgfqpoint{1.025758in}{1.241959in}}%
\pgfpathlineto{\pgfqpoint{1.010266in}{1.255965in}}%
\pgfpathlineto{\pgfqpoint{0.994790in}{1.270304in}}%
\pgfpathlineto{\pgfqpoint{0.986994in}{1.257531in}}%
\pgfpathlineto{\pgfqpoint{0.979181in}{1.244921in}}%
\pgfpathclose%
\pgfusepath{fill}%
\end{pgfscope}%
\begin{pgfscope}%
\pgfpathrectangle{\pgfqpoint{0.407592in}{0.222904in}}{\pgfqpoint{1.411163in}{1.411163in}}%
\pgfusepath{clip}%
\pgfsetbuttcap%
\pgfsetroundjoin%
\definecolor{currentfill}{rgb}{0.556863,0.654902,0.847059}%
\pgfsetfillcolor{currentfill}%
\pgfsetfillopacity{0.130000}%
\pgfsetlinewidth{0.000000pt}%
\definecolor{currentstroke}{rgb}{0.000000,0.000000,0.000000}%
\pgfsetstrokecolor{currentstroke}%
\pgfsetdash{}{0pt}%
\pgfpathmoveto{\pgfqpoint{0.823702in}{1.330551in}}%
\pgfpathlineto{\pgfqpoint{0.839170in}{1.313154in}}%
\pgfpathlineto{\pgfqpoint{0.854648in}{1.295902in}}%
\pgfpathlineto{\pgfqpoint{0.862630in}{1.306447in}}%
\pgfpathlineto{\pgfqpoint{0.870591in}{1.317199in}}%
\pgfpathlineto{\pgfqpoint{0.855147in}{1.334064in}}%
\pgfpathlineto{\pgfqpoint{0.839713in}{1.351088in}}%
\pgfpathlineto{\pgfqpoint{0.831718in}{1.340718in}}%
\pgfpathlineto{\pgfqpoint{0.823702in}{1.330551in}}%
\pgfpathclose%
\pgfusepath{fill}%
\end{pgfscope}%
\begin{pgfscope}%
\pgfpathrectangle{\pgfqpoint{0.407592in}{0.222904in}}{\pgfqpoint{1.411163in}{1.411163in}}%
\pgfusepath{clip}%
\pgfsetbuttcap%
\pgfsetroundjoin%
\definecolor{currentfill}{rgb}{0.556863,0.654902,0.847059}%
\pgfsetfillcolor{currentfill}%
\pgfsetfillopacity{0.130000}%
\pgfsetlinewidth{0.000000pt}%
\definecolor{currentstroke}{rgb}{0.000000,0.000000,0.000000}%
\pgfsetstrokecolor{currentstroke}%
\pgfsetdash{}{0pt}%
\pgfpathmoveto{\pgfqpoint{0.932483in}{1.251628in}}%
\pgfpathlineto{\pgfqpoint{0.947986in}{1.235795in}}%
\pgfpathlineto{\pgfqpoint{0.963502in}{1.220228in}}%
\pgfpathlineto{\pgfqpoint{0.971351in}{1.232482in}}%
\pgfpathlineto{\pgfqpoint{0.979181in}{1.244921in}}%
\pgfpathlineto{\pgfqpoint{0.963694in}{1.260023in}}%
\pgfpathlineto{\pgfqpoint{0.948220in}{1.275402in}}%
\pgfpathlineto{\pgfqpoint{0.940361in}{1.263421in}}%
\pgfpathlineto{\pgfqpoint{0.932483in}{1.251628in}}%
\pgfpathclose%
\pgfusepath{fill}%
\end{pgfscope}%
\begin{pgfscope}%
\pgfpathrectangle{\pgfqpoint{0.407592in}{0.222904in}}{\pgfqpoint{1.411163in}{1.411163in}}%
\pgfusepath{clip}%
\pgfsetbuttcap%
\pgfsetroundjoin%
\definecolor{currentfill}{rgb}{0.556863,0.654902,0.847059}%
\pgfsetfillcolor{currentfill}%
\pgfsetfillopacity{0.130000}%
\pgfsetlinewidth{0.000000pt}%
\definecolor{currentstroke}{rgb}{0.000000,0.000000,0.000000}%
\pgfsetstrokecolor{currentstroke}%
\pgfsetdash{}{0pt}%
\pgfpathmoveto{\pgfqpoint{1.056792in}{1.215048in}}%
\pgfpathlineto{\pgfqpoint{1.072337in}{1.202198in}}%
\pgfpathlineto{\pgfqpoint{1.087900in}{1.189791in}}%
\pgfpathlineto{\pgfqpoint{1.095612in}{1.204093in}}%
\pgfpathlineto{\pgfqpoint{1.080058in}{1.216270in}}%
\pgfpathlineto{\pgfqpoint{1.064524in}{1.228888in}}%
\pgfusepath{fill}%
\end{pgfscope}%
\begin{pgfscope}%
\pgfpathrectangle{\pgfqpoint{0.407592in}{0.222904in}}{\pgfqpoint{1.411163in}{1.411163in}}%
\pgfusepath{clip}%
\pgfsetbuttcap%
\pgfsetroundjoin%
\definecolor{currentfill}{rgb}{0.556863,0.654902,0.847059}%
\pgfsetfillcolor{currentfill}%
\pgfsetfillopacity{0.130000}%
\pgfsetlinewidth{0.000000pt}%
\definecolor{currentstroke}{rgb}{0.000000,0.000000,0.000000}%
\pgfsetstrokecolor{currentstroke}%
\pgfsetdash{}{0pt}%
\pgfpathmoveto{\pgfqpoint{1.010200in}{1.215642in}}%
\pgfpathlineto{\pgfqpoint{1.025732in}{1.201517in}}%
\pgfpathlineto{\pgfqpoint{1.041281in}{1.187774in}}%
\pgfpathlineto{\pgfqpoint{1.049045in}{1.201339in}}%
\pgfpathlineto{\pgfqpoint{1.056792in}{1.215048in}}%
\pgfpathlineto{\pgfqpoint{1.041266in}{1.228311in}}%
\pgfpathlineto{\pgfqpoint{1.025758in}{1.241959in}}%
\pgfpathlineto{\pgfqpoint{1.017987in}{1.228723in}}%
\pgfpathlineto{\pgfqpoint{1.010200in}{1.215642in}}%
\pgfpathclose%
\pgfusepath{fill}%
\end{pgfscope}%
\begin{pgfscope}%
\pgfpathrectangle{\pgfqpoint{0.407592in}{0.222904in}}{\pgfqpoint{1.411163in}{1.411163in}}%
\pgfusepath{clip}%
\pgfsetbuttcap%
\pgfsetroundjoin%
\definecolor{currentfill}{rgb}{0.556863,0.654902,0.847059}%
\pgfsetfillcolor{currentfill}%
\pgfsetfillopacity{0.130000}%
\pgfsetlinewidth{0.000000pt}%
\definecolor{currentstroke}{rgb}{0.000000,0.000000,0.000000}%
\pgfsetstrokecolor{currentstroke}%
\pgfsetdash{}{0pt}%
\pgfpathmoveto{\pgfqpoint{0.854648in}{1.295902in}}%
\pgfpathlineto{\pgfqpoint{0.870136in}{1.278810in}}%
\pgfpathlineto{\pgfqpoint{0.885636in}{1.261892in}}%
\pgfpathlineto{\pgfqpoint{0.893585in}{1.272843in}}%
\pgfpathlineto{\pgfqpoint{0.901514in}{1.284003in}}%
\pgfpathlineto{\pgfqpoint{0.886047in}{1.300507in}}%
\pgfpathlineto{\pgfqpoint{0.870591in}{1.317199in}}%
\pgfpathlineto{\pgfqpoint{0.862630in}{1.306447in}}%
\pgfpathlineto{\pgfqpoint{0.854648in}{1.295902in}}%
\pgfpathclose%
\pgfusepath{fill}%
\end{pgfscope}%
\begin{pgfscope}%
\pgfpathrectangle{\pgfqpoint{0.407592in}{0.222904in}}{\pgfqpoint{1.411163in}{1.411163in}}%
\pgfusepath{clip}%
\pgfsetbuttcap%
\pgfsetroundjoin%
\definecolor{currentfill}{rgb}{0.556863,0.654902,0.847059}%
\pgfsetfillcolor{currentfill}%
\pgfsetfillopacity{0.130000}%
\pgfsetlinewidth{0.000000pt}%
\definecolor{currentstroke}{rgb}{0.000000,0.000000,0.000000}%
\pgfsetstrokecolor{currentstroke}%
\pgfsetdash{}{0pt}%
\pgfpathmoveto{\pgfqpoint{0.963502in}{1.220228in}}%
\pgfpathlineto{\pgfqpoint{0.979031in}{1.204950in}}%
\pgfpathlineto{\pgfqpoint{0.994574in}{1.189987in}}%
\pgfpathlineto{\pgfqpoint{1.002396in}{1.202726in}}%
\pgfpathlineto{\pgfqpoint{1.010200in}{1.215642in}}%
\pgfpathlineto{\pgfqpoint{0.994683in}{1.230119in}}%
\pgfpathlineto{\pgfqpoint{0.979181in}{1.244921in}}%
\pgfpathlineto{\pgfqpoint{0.971351in}{1.232482in}}%
\pgfpathlineto{\pgfqpoint{0.963502in}{1.220228in}}%
\pgfpathclose%
\pgfusepath{fill}%
\end{pgfscope}%
\begin{pgfscope}%
\pgfpathrectangle{\pgfqpoint{0.407592in}{0.222904in}}{\pgfqpoint{1.411163in}{1.411163in}}%
\pgfusepath{clip}%
\pgfsetbuttcap%
\pgfsetroundjoin%
\definecolor{currentfill}{rgb}{0.556863,0.654902,0.847059}%
\pgfsetfillcolor{currentfill}%
\pgfsetfillopacity{0.130000}%
\pgfsetlinewidth{0.000000pt}%
\definecolor{currentstroke}{rgb}{0.000000,0.000000,0.000000}%
\pgfsetstrokecolor{currentstroke}%
\pgfsetdash{}{0pt}%
\pgfpathmoveto{\pgfqpoint{0.885636in}{1.261892in}}%
\pgfpathlineto{\pgfqpoint{0.901146in}{1.245164in}}%
\pgfpathlineto{\pgfqpoint{0.916668in}{1.228646in}}%
\pgfpathlineto{\pgfqpoint{0.924585in}{1.240032in}}%
\pgfpathlineto{\pgfqpoint{0.932483in}{1.251628in}}%
\pgfpathlineto{\pgfqpoint{0.916992in}{1.267704in}}%
\pgfpathlineto{\pgfqpoint{0.901514in}{1.284003in}}%
\pgfpathlineto{\pgfqpoint{0.893585in}{1.272843in}}%
\pgfpathlineto{\pgfqpoint{0.885636in}{1.261892in}}%
\pgfpathclose%
\pgfusepath{fill}%
\end{pgfscope}%
\begin{pgfscope}%
\pgfpathrectangle{\pgfqpoint{0.407592in}{0.222904in}}{\pgfqpoint{1.411163in}{1.411163in}}%
\pgfusepath{clip}%
\pgfsetbuttcap%
\pgfsetroundjoin%
\definecolor{currentfill}{rgb}{0.556863,0.654902,0.847059}%
\pgfsetfillcolor{currentfill}%
\pgfsetfillopacity{0.130000}%
\pgfsetlinewidth{0.000000pt}%
\definecolor{currentstroke}{rgb}{0.000000,0.000000,0.000000}%
\pgfsetstrokecolor{currentstroke}%
\pgfsetdash{}{0pt}%
\pgfpathmoveto{\pgfqpoint{1.087900in}{1.189791in}}%
\pgfpathlineto{\pgfqpoint{1.103485in}{1.177859in}}%
\pgfpathlineto{\pgfqpoint{1.119091in}{1.166434in}}%
\pgfpathlineto{\pgfqpoint{1.126785in}{1.181183in}}%
\pgfpathlineto{\pgfqpoint{1.111187in}{1.192387in}}%
\pgfpathlineto{\pgfqpoint{1.095612in}{1.204093in}}%
\pgfusepath{fill}%
\end{pgfscope}%
\begin{pgfscope}%
\pgfpathrectangle{\pgfqpoint{0.407592in}{0.222904in}}{\pgfqpoint{1.411163in}{1.411163in}}%
\pgfusepath{clip}%
\pgfsetbuttcap%
\pgfsetroundjoin%
\definecolor{currentfill}{rgb}{0.556863,0.654902,0.847059}%
\pgfsetfillcolor{currentfill}%
\pgfsetfillopacity{0.130000}%
\pgfsetlinewidth{0.000000pt}%
\definecolor{currentstroke}{rgb}{0.000000,0.000000,0.000000}%
\pgfsetstrokecolor{currentstroke}%
\pgfsetdash{}{0pt}%
\pgfpathmoveto{\pgfqpoint{1.041281in}{1.187774in}}%
\pgfpathlineto{\pgfqpoint{1.056848in}{1.174443in}}%
\pgfpathlineto{\pgfqpoint{1.072432in}{1.161557in}}%
\pgfpathlineto{\pgfqpoint{1.080174in}{1.175609in}}%
\pgfpathlineto{\pgfqpoint{1.087900in}{1.189791in}}%
\pgfpathlineto{\pgfqpoint{1.072337in}{1.202198in}}%
\pgfpathlineto{\pgfqpoint{1.056792in}{1.215048in}}%
\pgfpathlineto{\pgfqpoint{1.049045in}{1.201339in}}%
\pgfpathlineto{\pgfqpoint{1.041281in}{1.187774in}}%
\pgfpathclose%
\pgfusepath{fill}%
\end{pgfscope}%
\begin{pgfscope}%
\pgfpathrectangle{\pgfqpoint{0.407592in}{0.222904in}}{\pgfqpoint{1.411163in}{1.411163in}}%
\pgfusepath{clip}%
\pgfsetbuttcap%
\pgfsetroundjoin%
\definecolor{currentfill}{rgb}{0.556863,0.654902,0.847059}%
\pgfsetfillcolor{currentfill}%
\pgfsetfillopacity{0.130000}%
\pgfsetlinewidth{0.000000pt}%
\definecolor{currentstroke}{rgb}{0.000000,0.000000,0.000000}%
\pgfsetstrokecolor{currentstroke}%
\pgfsetdash{}{0pt}%
\pgfpathmoveto{\pgfqpoint{0.807601in}{1.310858in}}%
\pgfpathlineto{\pgfqpoint{0.823104in}{1.293097in}}%
\pgfpathlineto{\pgfqpoint{0.838617in}{1.275467in}}%
\pgfpathlineto{\pgfqpoint{0.846643in}{1.285572in}}%
\pgfpathlineto{\pgfqpoint{0.854648in}{1.295902in}}%
\pgfpathlineto{\pgfqpoint{0.839170in}{1.313154in}}%
\pgfpathlineto{\pgfqpoint{0.823702in}{1.330551in}}%
\pgfpathlineto{\pgfqpoint{0.815663in}{1.320595in}}%
\pgfpathlineto{\pgfqpoint{0.807601in}{1.310858in}}%
\pgfpathclose%
\pgfusepath{fill}%
\end{pgfscope}%
\begin{pgfscope}%
\pgfpathrectangle{\pgfqpoint{0.407592in}{0.222904in}}{\pgfqpoint{1.411163in}{1.411163in}}%
\pgfusepath{clip}%
\pgfsetbuttcap%
\pgfsetroundjoin%
\definecolor{currentfill}{rgb}{0.556863,0.654902,0.847059}%
\pgfsetfillcolor{currentfill}%
\pgfsetfillopacity{0.130000}%
\pgfsetlinewidth{0.000000pt}%
\definecolor{currentstroke}{rgb}{0.000000,0.000000,0.000000}%
\pgfsetstrokecolor{currentstroke}%
\pgfsetdash{}{0pt}%
\pgfpathmoveto{\pgfqpoint{0.916668in}{1.228646in}}%
\pgfpathlineto{\pgfqpoint{0.932201in}{1.212356in}}%
\pgfpathlineto{\pgfqpoint{0.947747in}{1.196318in}}%
\pgfpathlineto{\pgfqpoint{0.955634in}{1.208169in}}%
\pgfpathlineto{\pgfqpoint{0.963502in}{1.220228in}}%
\pgfpathlineto{\pgfqpoint{0.947986in}{1.235795in}}%
\pgfpathlineto{\pgfqpoint{0.932483in}{1.251628in}}%
\pgfpathlineto{\pgfqpoint{0.924585in}{1.240032in}}%
\pgfpathlineto{\pgfqpoint{0.916668in}{1.228646in}}%
\pgfpathclose%
\pgfusepath{fill}%
\end{pgfscope}%
\begin{pgfscope}%
\pgfpathrectangle{\pgfqpoint{0.407592in}{0.222904in}}{\pgfqpoint{1.411163in}{1.411163in}}%
\pgfusepath{clip}%
\pgfsetbuttcap%
\pgfsetroundjoin%
\definecolor{currentfill}{rgb}{0.556863,0.654902,0.847059}%
\pgfsetfillcolor{currentfill}%
\pgfsetfillopacity{0.130000}%
\pgfsetlinewidth{0.000000pt}%
\definecolor{currentstroke}{rgb}{0.000000,0.000000,0.000000}%
\pgfsetstrokecolor{currentstroke}%
\pgfsetdash{}{0pt}%
\pgfpathmoveto{\pgfqpoint{0.994574in}{1.189987in}}%
\pgfpathlineto{\pgfqpoint{1.010132in}{1.175368in}}%
\pgfpathlineto{\pgfqpoint{1.025705in}{1.161124in}}%
\pgfpathlineto{\pgfqpoint{1.033502in}{1.174365in}}%
\pgfpathlineto{\pgfqpoint{1.041281in}{1.187774in}}%
\pgfpathlineto{\pgfqpoint{1.025732in}{1.201517in}}%
\pgfpathlineto{\pgfqpoint{1.010200in}{1.215642in}}%
\pgfpathlineto{\pgfqpoint{1.002396in}{1.202726in}}%
\pgfpathlineto{\pgfqpoint{0.994574in}{1.189987in}}%
\pgfpathclose%
\pgfusepath{fill}%
\end{pgfscope}%
\begin{pgfscope}%
\pgfpathrectangle{\pgfqpoint{0.407592in}{0.222904in}}{\pgfqpoint{1.411163in}{1.411163in}}%
\pgfusepath{clip}%
\pgfsetbuttcap%
\pgfsetroundjoin%
\definecolor{currentfill}{rgb}{0.556863,0.654902,0.847059}%
\pgfsetfillcolor{currentfill}%
\pgfsetfillopacity{0.130000}%
\pgfsetlinewidth{0.000000pt}%
\definecolor{currentstroke}{rgb}{0.000000,0.000000,0.000000}%
\pgfsetstrokecolor{currentstroke}%
\pgfsetdash{}{0pt}%
\pgfpathmoveto{\pgfqpoint{0.838617in}{1.275467in}}%
\pgfpathlineto{\pgfqpoint{0.854140in}{1.257983in}}%
\pgfpathlineto{\pgfqpoint{0.869673in}{1.240657in}}%
\pgfpathlineto{\pgfqpoint{0.877665in}{1.251160in}}%
\pgfpathlineto{\pgfqpoint{0.885636in}{1.261892in}}%
\pgfpathlineto{\pgfqpoint{0.870136in}{1.278810in}}%
\pgfpathlineto{\pgfqpoint{0.854648in}{1.295902in}}%
\pgfpathlineto{\pgfqpoint{0.846643in}{1.285572in}}%
\pgfpathlineto{\pgfqpoint{0.838617in}{1.275467in}}%
\pgfpathclose%
\pgfusepath{fill}%
\end{pgfscope}%
\begin{pgfscope}%
\pgfpathrectangle{\pgfqpoint{0.407592in}{0.222904in}}{\pgfqpoint{1.411163in}{1.411163in}}%
\pgfusepath{clip}%
\pgfsetbuttcap%
\pgfsetroundjoin%
\definecolor{currentfill}{rgb}{0.556863,0.654902,0.847059}%
\pgfsetfillcolor{currentfill}%
\pgfsetfillopacity{0.130000}%
\pgfsetlinewidth{0.000000pt}%
\definecolor{currentstroke}{rgb}{0.000000,0.000000,0.000000}%
\pgfsetstrokecolor{currentstroke}%
\pgfsetdash{}{0pt}%
\pgfpathmoveto{\pgfqpoint{1.072432in}{1.161557in}}%
\pgfpathlineto{\pgfqpoint{1.088036in}{1.149150in}}%
\pgfpathlineto{\pgfqpoint{1.103662in}{1.137258in}}%
\pgfpathlineto{\pgfqpoint{1.111384in}{1.151788in}}%
\pgfpathlineto{\pgfqpoint{1.119091in}{1.166434in}}%
\pgfpathlineto{\pgfqpoint{1.103485in}{1.177859in}}%
\pgfpathlineto{\pgfqpoint{1.087900in}{1.189791in}}%
\pgfpathlineto{\pgfqpoint{1.080174in}{1.175609in}}%
\pgfpathlineto{\pgfqpoint{1.072432in}{1.161557in}}%
\pgfpathclose%
\pgfusepath{fill}%
\end{pgfscope}%
\begin{pgfscope}%
\pgfpathrectangle{\pgfqpoint{0.407592in}{0.222904in}}{\pgfqpoint{1.411163in}{1.411163in}}%
\pgfusepath{clip}%
\pgfsetbuttcap%
\pgfsetroundjoin%
\definecolor{currentfill}{rgb}{0.556863,0.654902,0.847059}%
\pgfsetfillcolor{currentfill}%
\pgfsetfillopacity{0.130000}%
\pgfsetlinewidth{0.000000pt}%
\definecolor{currentstroke}{rgb}{0.000000,0.000000,0.000000}%
\pgfsetstrokecolor{currentstroke}%
\pgfsetdash{}{0pt}%
\pgfpathmoveto{\pgfqpoint{1.119091in}{1.166434in}}%
\pgfpathlineto{\pgfqpoint{1.134721in}{1.155547in}}%
\pgfpathlineto{\pgfqpoint{1.150376in}{1.145233in}}%
\pgfpathlineto{\pgfqpoint{1.158053in}{1.160402in}}%
\pgfpathlineto{\pgfqpoint{1.142406in}{1.170511in}}%
\pgfpathlineto{\pgfqpoint{1.126785in}{1.181183in}}%
\pgfusepath{fill}%
\end{pgfscope}%
\begin{pgfscope}%
\pgfpathrectangle{\pgfqpoint{0.407592in}{0.222904in}}{\pgfqpoint{1.411163in}{1.411163in}}%
\pgfusepath{clip}%
\pgfsetbuttcap%
\pgfsetroundjoin%
\definecolor{currentfill}{rgb}{0.556863,0.654902,0.847059}%
\pgfsetfillcolor{currentfill}%
\pgfsetfillopacity{0.130000}%
\pgfsetlinewidth{0.000000pt}%
\definecolor{currentstroke}{rgb}{0.000000,0.000000,0.000000}%
\pgfsetstrokecolor{currentstroke}%
\pgfsetdash{}{0pt}%
\pgfpathmoveto{\pgfqpoint{0.947747in}{1.196318in}}%
\pgfpathlineto{\pgfqpoint{0.963305in}{1.180555in}}%
\pgfpathlineto{\pgfqpoint{0.978877in}{1.165094in}}%
\pgfpathlineto{\pgfqpoint{0.986734in}{1.177439in}}%
\pgfpathlineto{\pgfqpoint{0.994574in}{1.189987in}}%
\pgfpathlineto{\pgfqpoint{0.979031in}{1.204950in}}%
\pgfpathlineto{\pgfqpoint{0.963502in}{1.220228in}}%
\pgfpathlineto{\pgfqpoint{0.955634in}{1.208169in}}%
\pgfpathlineto{\pgfqpoint{0.947747in}{1.196318in}}%
\pgfpathclose%
\pgfusepath{fill}%
\end{pgfscope}%
\begin{pgfscope}%
\pgfpathrectangle{\pgfqpoint{0.407592in}{0.222904in}}{\pgfqpoint{1.411163in}{1.411163in}}%
\pgfusepath{clip}%
\pgfsetbuttcap%
\pgfsetroundjoin%
\definecolor{currentfill}{rgb}{0.556863,0.654902,0.847059}%
\pgfsetfillcolor{currentfill}%
\pgfsetfillopacity{0.130000}%
\pgfsetlinewidth{0.000000pt}%
\definecolor{currentstroke}{rgb}{0.000000,0.000000,0.000000}%
\pgfsetstrokecolor{currentstroke}%
\pgfsetdash{}{0pt}%
\pgfpathmoveto{\pgfqpoint{0.869673in}{1.240657in}}%
\pgfpathlineto{\pgfqpoint{0.885216in}{1.223506in}}%
\pgfpathlineto{\pgfqpoint{0.900771in}{1.206548in}}%
\pgfpathlineto{\pgfqpoint{0.908730in}{1.217481in}}%
\pgfpathlineto{\pgfqpoint{0.916668in}{1.228646in}}%
\pgfpathlineto{\pgfqpoint{0.901146in}{1.245164in}}%
\pgfpathlineto{\pgfqpoint{0.885636in}{1.261892in}}%
\pgfpathlineto{\pgfqpoint{0.877665in}{1.251160in}}%
\pgfpathlineto{\pgfqpoint{0.869673in}{1.240657in}}%
\pgfpathclose%
\pgfusepath{fill}%
\end{pgfscope}%
\begin{pgfscope}%
\pgfpathrectangle{\pgfqpoint{0.407592in}{0.222904in}}{\pgfqpoint{1.411163in}{1.411163in}}%
\pgfusepath{clip}%
\pgfsetbuttcap%
\pgfsetroundjoin%
\definecolor{currentfill}{rgb}{0.556863,0.654902,0.847059}%
\pgfsetfillcolor{currentfill}%
\pgfsetfillopacity{0.130000}%
\pgfsetlinewidth{0.000000pt}%
\definecolor{currentstroke}{rgb}{0.000000,0.000000,0.000000}%
\pgfsetstrokecolor{currentstroke}%
\pgfsetdash{}{0pt}%
\pgfpathmoveto{\pgfqpoint{1.025705in}{1.161124in}}%
\pgfpathlineto{\pgfqpoint{1.041295in}{1.147287in}}%
\pgfpathlineto{\pgfqpoint{1.056902in}{1.133893in}}%
\pgfpathlineto{\pgfqpoint{1.064675in}{1.147647in}}%
\pgfpathlineto{\pgfqpoint{1.072432in}{1.161557in}}%
\pgfpathlineto{\pgfqpoint{1.056848in}{1.174443in}}%
\pgfpathlineto{\pgfqpoint{1.041281in}{1.187774in}}%
\pgfpathlineto{\pgfqpoint{1.033502in}{1.174365in}}%
\pgfpathlineto{\pgfqpoint{1.025705in}{1.161124in}}%
\pgfpathclose%
\pgfusepath{fill}%
\end{pgfscope}%
\begin{pgfscope}%
\pgfpathrectangle{\pgfqpoint{0.407592in}{0.222904in}}{\pgfqpoint{1.411163in}{1.411163in}}%
\pgfusepath{clip}%
\pgfsetbuttcap%
\pgfsetroundjoin%
\definecolor{currentfill}{rgb}{0.556863,0.654902,0.847059}%
\pgfsetfillcolor{currentfill}%
\pgfsetfillopacity{0.130000}%
\pgfsetlinewidth{0.000000pt}%
\definecolor{currentstroke}{rgb}{0.000000,0.000000,0.000000}%
\pgfsetstrokecolor{currentstroke}%
\pgfsetdash{}{0pt}%
\pgfpathmoveto{\pgfqpoint{0.791405in}{1.292074in}}%
\pgfpathlineto{\pgfqpoint{0.806944in}{1.273963in}}%
\pgfpathlineto{\pgfqpoint{0.822493in}{1.255969in}}%
\pgfpathlineto{\pgfqpoint{0.830567in}{1.265596in}}%
\pgfpathlineto{\pgfqpoint{0.838617in}{1.275467in}}%
\pgfpathlineto{\pgfqpoint{0.823104in}{1.293097in}}%
\pgfpathlineto{\pgfqpoint{0.807601in}{1.310858in}}%
\pgfpathlineto{\pgfqpoint{0.799515in}{1.301348in}}%
\pgfpathlineto{\pgfqpoint{0.791405in}{1.292074in}}%
\pgfpathclose%
\pgfusepath{fill}%
\end{pgfscope}%
\begin{pgfscope}%
\pgfpathrectangle{\pgfqpoint{0.407592in}{0.222904in}}{\pgfqpoint{1.411163in}{1.411163in}}%
\pgfusepath{clip}%
\pgfsetbuttcap%
\pgfsetroundjoin%
\definecolor{currentfill}{rgb}{0.556863,0.654902,0.847059}%
\pgfsetfillcolor{currentfill}%
\pgfsetfillopacity{0.130000}%
\pgfsetlinewidth{0.000000pt}%
\definecolor{currentstroke}{rgb}{0.000000,0.000000,0.000000}%
\pgfsetstrokecolor{currentstroke}%
\pgfsetdash{}{0pt}%
\pgfpathmoveto{\pgfqpoint{0.978877in}{1.165094in}}%
\pgfpathlineto{\pgfqpoint{0.994462in}{1.149965in}}%
\pgfpathlineto{\pgfqpoint{1.010061in}{1.135201in}}%
\pgfpathlineto{\pgfqpoint{1.017892in}{1.148064in}}%
\pgfpathlineto{\pgfqpoint{1.025705in}{1.161124in}}%
\pgfpathlineto{\pgfqpoint{1.010132in}{1.175368in}}%
\pgfpathlineto{\pgfqpoint{0.994574in}{1.189987in}}%
\pgfpathlineto{\pgfqpoint{0.986734in}{1.177439in}}%
\pgfpathlineto{\pgfqpoint{0.978877in}{1.165094in}}%
\pgfpathclose%
\pgfusepath{fill}%
\end{pgfscope}%
\begin{pgfscope}%
\pgfpathrectangle{\pgfqpoint{0.407592in}{0.222904in}}{\pgfqpoint{1.411163in}{1.411163in}}%
\pgfusepath{clip}%
\pgfsetbuttcap%
\pgfsetroundjoin%
\definecolor{currentfill}{rgb}{0.556863,0.654902,0.847059}%
\pgfsetfillcolor{currentfill}%
\pgfsetfillopacity{0.130000}%
\pgfsetlinewidth{0.000000pt}%
\definecolor{currentstroke}{rgb}{0.000000,0.000000,0.000000}%
\pgfsetstrokecolor{currentstroke}%
\pgfsetdash{}{0pt}%
\pgfpathmoveto{\pgfqpoint{0.900771in}{1.206548in}}%
\pgfpathlineto{\pgfqpoint{0.916337in}{1.189803in}}%
\pgfpathlineto{\pgfqpoint{0.931914in}{1.173292in}}%
\pgfpathlineto{\pgfqpoint{0.939840in}{1.184688in}}%
\pgfpathlineto{\pgfqpoint{0.947747in}{1.196318in}}%
\pgfpathlineto{\pgfqpoint{0.932201in}{1.212356in}}%
\pgfpathlineto{\pgfqpoint{0.916668in}{1.228646in}}%
\pgfpathlineto{\pgfqpoint{0.908730in}{1.217481in}}%
\pgfpathlineto{\pgfqpoint{0.900771in}{1.206548in}}%
\pgfpathclose%
\pgfusepath{fill}%
\end{pgfscope}%
\begin{pgfscope}%
\pgfpathrectangle{\pgfqpoint{0.407592in}{0.222904in}}{\pgfqpoint{1.411163in}{1.411163in}}%
\pgfusepath{clip}%
\pgfsetbuttcap%
\pgfsetroundjoin%
\definecolor{currentfill}{rgb}{0.556863,0.654902,0.847059}%
\pgfsetfillcolor{currentfill}%
\pgfsetfillopacity{0.130000}%
\pgfsetlinewidth{0.000000pt}%
\definecolor{currentstroke}{rgb}{0.000000,0.000000,0.000000}%
\pgfsetstrokecolor{currentstroke}%
\pgfsetdash{}{0pt}%
\pgfpathmoveto{\pgfqpoint{1.103662in}{1.137258in}}%
\pgfpathlineto{\pgfqpoint{1.119309in}{1.125917in}}%
\pgfpathlineto{\pgfqpoint{1.134980in}{1.115167in}}%
\pgfpathlineto{\pgfqpoint{1.142685in}{1.130151in}}%
\pgfpathlineto{\pgfqpoint{1.150376in}{1.145233in}}%
\pgfpathlineto{\pgfqpoint{1.134721in}{1.155547in}}%
\pgfpathlineto{\pgfqpoint{1.119091in}{1.166434in}}%
\pgfpathlineto{\pgfqpoint{1.111384in}{1.151788in}}%
\pgfpathlineto{\pgfqpoint{1.103662in}{1.137258in}}%
\pgfpathclose%
\pgfusepath{fill}%
\end{pgfscope}%
\begin{pgfscope}%
\pgfpathrectangle{\pgfqpoint{0.407592in}{0.222904in}}{\pgfqpoint{1.411163in}{1.411163in}}%
\pgfusepath{clip}%
\pgfsetbuttcap%
\pgfsetroundjoin%
\definecolor{currentfill}{rgb}{0.556863,0.654902,0.847059}%
\pgfsetfillcolor{currentfill}%
\pgfsetfillopacity{0.130000}%
\pgfsetlinewidth{0.000000pt}%
\definecolor{currentstroke}{rgb}{0.000000,0.000000,0.000000}%
\pgfsetstrokecolor{currentstroke}%
\pgfsetdash{}{0pt}%
\pgfpathmoveto{\pgfqpoint{1.150376in}{1.145233in}}%
\pgfpathlineto{\pgfqpoint{1.166057in}{1.135524in}}%
\pgfpathlineto{\pgfqpoint{1.181767in}{1.126452in}}%
\pgfpathlineto{\pgfqpoint{1.189430in}{1.141996in}}%
\pgfpathlineto{\pgfqpoint{1.173727in}{1.150887in}}%
\pgfpathlineto{\pgfqpoint{1.158053in}{1.160402in}}%
\pgfusepath{fill}%
\end{pgfscope}%
\begin{pgfscope}%
\pgfpathrectangle{\pgfqpoint{0.407592in}{0.222904in}}{\pgfqpoint{1.411163in}{1.411163in}}%
\pgfusepath{clip}%
\pgfsetbuttcap%
\pgfsetroundjoin%
\definecolor{currentfill}{rgb}{0.556863,0.654902,0.847059}%
\pgfsetfillcolor{currentfill}%
\pgfsetfillopacity{0.130000}%
\pgfsetlinewidth{0.000000pt}%
\definecolor{currentstroke}{rgb}{0.000000,0.000000,0.000000}%
\pgfsetstrokecolor{currentstroke}%
\pgfsetdash{}{0pt}%
\pgfpathmoveto{\pgfqpoint{1.056902in}{1.133893in}}%
\pgfpathlineto{\pgfqpoint{1.072527in}{1.120981in}}%
\pgfpathlineto{\pgfqpoint{1.088173in}{1.108589in}}%
\pgfpathlineto{\pgfqpoint{1.095925in}{1.122853in}}%
\pgfpathlineto{\pgfqpoint{1.103662in}{1.137258in}}%
\pgfpathlineto{\pgfqpoint{1.088036in}{1.149150in}}%
\pgfpathlineto{\pgfqpoint{1.072432in}{1.161557in}}%
\pgfpathlineto{\pgfqpoint{1.064675in}{1.147647in}}%
\pgfpathlineto{\pgfqpoint{1.056902in}{1.133893in}}%
\pgfpathclose%
\pgfusepath{fill}%
\end{pgfscope}%
\begin{pgfscope}%
\pgfpathrectangle{\pgfqpoint{0.407592in}{0.222904in}}{\pgfqpoint{1.411163in}{1.411163in}}%
\pgfusepath{clip}%
\pgfsetbuttcap%
\pgfsetroundjoin%
\definecolor{currentfill}{rgb}{0.556863,0.654902,0.847059}%
\pgfsetfillcolor{currentfill}%
\pgfsetfillopacity{0.130000}%
\pgfsetlinewidth{0.000000pt}%
\definecolor{currentstroke}{rgb}{0.000000,0.000000,0.000000}%
\pgfsetstrokecolor{currentstroke}%
\pgfsetdash{}{0pt}%
\pgfpathmoveto{\pgfqpoint{0.822493in}{1.255969in}}%
\pgfpathlineto{\pgfqpoint{0.838052in}{1.238104in}}%
\pgfpathlineto{\pgfqpoint{0.853620in}{1.220383in}}%
\pgfpathlineto{\pgfqpoint{0.861658in}{1.230394in}}%
\pgfpathlineto{\pgfqpoint{0.869673in}{1.240657in}}%
\pgfpathlineto{\pgfqpoint{0.854140in}{1.257983in}}%
\pgfpathlineto{\pgfqpoint{0.838617in}{1.275467in}}%
\pgfpathlineto{\pgfqpoint{0.830567in}{1.265596in}}%
\pgfpathlineto{\pgfqpoint{0.822493in}{1.255969in}}%
\pgfpathclose%
\pgfusepath{fill}%
\end{pgfscope}%
\begin{pgfscope}%
\pgfpathrectangle{\pgfqpoint{0.407592in}{0.222904in}}{\pgfqpoint{1.411163in}{1.411163in}}%
\pgfusepath{clip}%
\pgfsetbuttcap%
\pgfsetroundjoin%
\definecolor{currentfill}{rgb}{0.556863,0.654902,0.847059}%
\pgfsetfillcolor{currentfill}%
\pgfsetfillopacity{0.130000}%
\pgfsetlinewidth{0.000000pt}%
\definecolor{currentstroke}{rgb}{0.000000,0.000000,0.000000}%
\pgfsetstrokecolor{currentstroke}%
\pgfsetdash{}{0pt}%
\pgfpathmoveto{\pgfqpoint{0.931914in}{1.173292in}}%
\pgfpathlineto{\pgfqpoint{0.947503in}{1.157040in}}%
\pgfpathlineto{\pgfqpoint{0.963104in}{1.141075in}}%
\pgfpathlineto{\pgfqpoint{0.971000in}{1.152968in}}%
\pgfpathlineto{\pgfqpoint{0.978877in}{1.165094in}}%
\pgfpathlineto{\pgfqpoint{0.963305in}{1.180555in}}%
\pgfpathlineto{\pgfqpoint{0.947747in}{1.196318in}}%
\pgfpathlineto{\pgfqpoint{0.939840in}{1.184688in}}%
\pgfpathlineto{\pgfqpoint{0.931914in}{1.173292in}}%
\pgfpathclose%
\pgfusepath{fill}%
\end{pgfscope}%
\begin{pgfscope}%
\pgfpathrectangle{\pgfqpoint{0.407592in}{0.222904in}}{\pgfqpoint{1.411163in}{1.411163in}}%
\pgfusepath{clip}%
\pgfsetbuttcap%
\pgfsetroundjoin%
\definecolor{currentfill}{rgb}{0.556863,0.654902,0.847059}%
\pgfsetfillcolor{currentfill}%
\pgfsetfillopacity{0.130000}%
\pgfsetlinewidth{0.000000pt}%
\definecolor{currentstroke}{rgb}{0.000000,0.000000,0.000000}%
\pgfsetstrokecolor{currentstroke}%
\pgfsetdash{}{0pt}%
\pgfpathmoveto{\pgfqpoint{1.010061in}{1.135201in}}%
\pgfpathlineto{\pgfqpoint{1.025676in}{1.120837in}}%
\pgfpathlineto{\pgfqpoint{1.041307in}{1.106910in}}%
\pgfpathlineto{\pgfqpoint{1.049113in}{1.120309in}}%
\pgfpathlineto{\pgfqpoint{1.056902in}{1.133893in}}%
\pgfpathlineto{\pgfqpoint{1.041295in}{1.147287in}}%
\pgfpathlineto{\pgfqpoint{1.025705in}{1.161124in}}%
\pgfpathlineto{\pgfqpoint{1.017892in}{1.148064in}}%
\pgfpathlineto{\pgfqpoint{1.010061in}{1.135201in}}%
\pgfpathclose%
\pgfusepath{fill}%
\end{pgfscope}%
\begin{pgfscope}%
\pgfpathrectangle{\pgfqpoint{0.407592in}{0.222904in}}{\pgfqpoint{1.411163in}{1.411163in}}%
\pgfusepath{clip}%
\pgfsetbuttcap%
\pgfsetroundjoin%
\definecolor{currentfill}{rgb}{0.556863,0.654902,0.847059}%
\pgfsetfillcolor{currentfill}%
\pgfsetfillopacity{0.130000}%
\pgfsetlinewidth{0.000000pt}%
\definecolor{currentstroke}{rgb}{0.000000,0.000000,0.000000}%
\pgfsetstrokecolor{currentstroke}%
\pgfsetdash{}{0pt}%
\pgfpathmoveto{\pgfqpoint{0.853620in}{1.220383in}}%
\pgfpathlineto{\pgfqpoint{0.869199in}{1.202819in}}%
\pgfpathlineto{\pgfqpoint{0.884788in}{1.185431in}}%
\pgfpathlineto{\pgfqpoint{0.892791in}{1.195861in}}%
\pgfpathlineto{\pgfqpoint{0.900771in}{1.206548in}}%
\pgfpathlineto{\pgfqpoint{0.885216in}{1.223506in}}%
\pgfpathlineto{\pgfqpoint{0.869673in}{1.240657in}}%
\pgfpathlineto{\pgfqpoint{0.861658in}{1.230394in}}%
\pgfpathlineto{\pgfqpoint{0.853620in}{1.220383in}}%
\pgfpathclose%
\pgfusepath{fill}%
\end{pgfscope}%
\begin{pgfscope}%
\pgfpathrectangle{\pgfqpoint{0.407592in}{0.222904in}}{\pgfqpoint{1.411163in}{1.411163in}}%
\pgfusepath{clip}%
\pgfsetbuttcap%
\pgfsetroundjoin%
\definecolor{currentfill}{rgb}{0.556863,0.654902,0.847059}%
\pgfsetfillcolor{currentfill}%
\pgfsetfillopacity{0.130000}%
\pgfsetlinewidth{0.000000pt}%
\definecolor{currentstroke}{rgb}{0.000000,0.000000,0.000000}%
\pgfsetstrokecolor{currentstroke}%
\pgfsetdash{}{0pt}%
\pgfpathmoveto{\pgfqpoint{1.134980in}{1.115167in}}%
\pgfpathlineto{\pgfqpoint{1.150676in}{1.105044in}}%
\pgfpathlineto{\pgfqpoint{1.166400in}{1.095585in}}%
\pgfpathlineto{\pgfqpoint{1.174090in}{1.110979in}}%
\pgfpathlineto{\pgfqpoint{1.181767in}{1.126452in}}%
\pgfpathlineto{\pgfqpoint{1.166057in}{1.135524in}}%
\pgfpathlineto{\pgfqpoint{1.150376in}{1.145233in}}%
\pgfpathlineto{\pgfqpoint{1.142685in}{1.130151in}}%
\pgfpathlineto{\pgfqpoint{1.134980in}{1.115167in}}%
\pgfpathclose%
\pgfusepath{fill}%
\end{pgfscope}%
\begin{pgfscope}%
\pgfpathrectangle{\pgfqpoint{0.407592in}{0.222904in}}{\pgfqpoint{1.411163in}{1.411163in}}%
\pgfusepath{clip}%
\pgfsetbuttcap%
\pgfsetroundjoin%
\definecolor{currentfill}{rgb}{0.556863,0.654902,0.847059}%
\pgfsetfillcolor{currentfill}%
\pgfsetfillopacity{0.130000}%
\pgfsetlinewidth{0.000000pt}%
\definecolor{currentstroke}{rgb}{0.000000,0.000000,0.000000}%
\pgfsetstrokecolor{currentstroke}%
\pgfsetdash{}{0pt}%
\pgfpathmoveto{\pgfqpoint{1.088173in}{1.108589in}}%
\pgfpathlineto{\pgfqpoint{1.103839in}{1.096760in}}%
\pgfpathlineto{\pgfqpoint{1.119527in}{1.085537in}}%
\pgfpathlineto{\pgfqpoint{1.127261in}{1.100291in}}%
\pgfpathlineto{\pgfqpoint{1.134980in}{1.115167in}}%
\pgfpathlineto{\pgfqpoint{1.119309in}{1.125917in}}%
\pgfpathlineto{\pgfqpoint{1.103662in}{1.137258in}}%
\pgfpathlineto{\pgfqpoint{1.095925in}{1.122853in}}%
\pgfpathlineto{\pgfqpoint{1.088173in}{1.108589in}}%
\pgfpathclose%
\pgfusepath{fill}%
\end{pgfscope}%
\begin{pgfscope}%
\pgfpathrectangle{\pgfqpoint{0.407592in}{0.222904in}}{\pgfqpoint{1.411163in}{1.411163in}}%
\pgfusepath{clip}%
\pgfsetbuttcap%
\pgfsetroundjoin%
\definecolor{currentfill}{rgb}{0.556863,0.654902,0.847059}%
\pgfsetfillcolor{currentfill}%
\pgfsetfillopacity{0.130000}%
\pgfsetlinewidth{0.000000pt}%
\definecolor{currentstroke}{rgb}{0.000000,0.000000,0.000000}%
\pgfsetstrokecolor{currentstroke}%
\pgfsetdash{}{0pt}%
\pgfpathmoveto{\pgfqpoint{0.963104in}{1.141075in}}%
\pgfpathlineto{\pgfqpoint{0.978718in}{1.125426in}}%
\pgfpathlineto{\pgfqpoint{0.994346in}{1.110128in}}%
\pgfpathlineto{\pgfqpoint{1.002213in}{1.122550in}}%
\pgfpathlineto{\pgfqpoint{1.010061in}{1.135201in}}%
\pgfpathlineto{\pgfqpoint{0.994462in}{1.149965in}}%
\pgfpathlineto{\pgfqpoint{0.978877in}{1.165094in}}%
\pgfpathlineto{\pgfqpoint{0.971000in}{1.152968in}}%
\pgfpathlineto{\pgfqpoint{0.963104in}{1.141075in}}%
\pgfpathclose%
\pgfusepath{fill}%
\end{pgfscope}%
\begin{pgfscope}%
\pgfpathrectangle{\pgfqpoint{0.407592in}{0.222904in}}{\pgfqpoint{1.411163in}{1.411163in}}%
\pgfusepath{clip}%
\pgfsetbuttcap%
\pgfsetroundjoin%
\definecolor{currentfill}{rgb}{0.556863,0.654902,0.847059}%
\pgfsetfillcolor{currentfill}%
\pgfsetfillopacity{0.130000}%
\pgfsetlinewidth{0.000000pt}%
\definecolor{currentstroke}{rgb}{0.000000,0.000000,0.000000}%
\pgfsetstrokecolor{currentstroke}%
\pgfsetdash{}{0pt}%
\pgfpathmoveto{\pgfqpoint{1.181767in}{1.126452in}}%
\pgfpathlineto{\pgfqpoint{1.197506in}{1.118048in}}%
\pgfpathlineto{\pgfqpoint{1.213278in}{1.110341in}}%
\pgfpathlineto{\pgfqpoint{1.220930in}{1.126195in}}%
\pgfpathlineto{\pgfqpoint{1.205163in}{1.133756in}}%
\pgfpathlineto{\pgfqpoint{1.189430in}{1.141996in}}%
\pgfusepath{fill}%
\end{pgfscope}%
\begin{pgfscope}%
\pgfpathrectangle{\pgfqpoint{0.407592in}{0.222904in}}{\pgfqpoint{1.411163in}{1.411163in}}%
\pgfusepath{clip}%
\pgfsetbuttcap%
\pgfsetroundjoin%
\definecolor{currentfill}{rgb}{0.556863,0.654902,0.847059}%
\pgfsetfillcolor{currentfill}%
\pgfsetfillopacity{0.130000}%
\pgfsetlinewidth{0.000000pt}%
\definecolor{currentstroke}{rgb}{0.000000,0.000000,0.000000}%
\pgfsetstrokecolor{currentstroke}%
\pgfsetdash{}{0pt}%
\pgfpathmoveto{\pgfqpoint{0.775110in}{1.274265in}}%
\pgfpathlineto{\pgfqpoint{0.790687in}{1.255823in}}%
\pgfpathlineto{\pgfqpoint{0.806273in}{1.237483in}}%
\pgfpathlineto{\pgfqpoint{0.814396in}{1.246595in}}%
\pgfpathlineto{\pgfqpoint{0.822493in}{1.255969in}}%
\pgfpathlineto{\pgfqpoint{0.806944in}{1.273963in}}%
\pgfpathlineto{\pgfqpoint{0.791405in}{1.292074in}}%
\pgfpathlineto{\pgfqpoint{0.783270in}{1.283044in}}%
\pgfpathlineto{\pgfqpoint{0.775110in}{1.274265in}}%
\pgfpathclose%
\pgfusepath{fill}%
\end{pgfscope}%
\begin{pgfscope}%
\pgfpathrectangle{\pgfqpoint{0.407592in}{0.222904in}}{\pgfqpoint{1.411163in}{1.411163in}}%
\pgfusepath{clip}%
\pgfsetbuttcap%
\pgfsetroundjoin%
\definecolor{currentfill}{rgb}{0.556863,0.654902,0.847059}%
\pgfsetfillcolor{currentfill}%
\pgfsetfillopacity{0.130000}%
\pgfsetlinewidth{0.000000pt}%
\definecolor{currentstroke}{rgb}{0.000000,0.000000,0.000000}%
\pgfsetstrokecolor{currentstroke}%
\pgfsetdash{}{0pt}%
\pgfpathmoveto{\pgfqpoint{0.884788in}{1.185431in}}%
\pgfpathlineto{\pgfqpoint{0.900387in}{1.168237in}}%
\pgfpathlineto{\pgfqpoint{0.915998in}{1.151260in}}%
\pgfpathlineto{\pgfqpoint{0.923966in}{1.162144in}}%
\pgfpathlineto{\pgfqpoint{0.931914in}{1.173292in}}%
\pgfpathlineto{\pgfqpoint{0.916337in}{1.189803in}}%
\pgfpathlineto{\pgfqpoint{0.900771in}{1.206548in}}%
\pgfpathlineto{\pgfqpoint{0.892791in}{1.195861in}}%
\pgfpathlineto{\pgfqpoint{0.884788in}{1.185431in}}%
\pgfpathclose%
\pgfusepath{fill}%
\end{pgfscope}%
\begin{pgfscope}%
\pgfpathrectangle{\pgfqpoint{0.407592in}{0.222904in}}{\pgfqpoint{1.411163in}{1.411163in}}%
\pgfusepath{clip}%
\pgfsetbuttcap%
\pgfsetroundjoin%
\definecolor{currentfill}{rgb}{0.556863,0.654902,0.847059}%
\pgfsetfillcolor{currentfill}%
\pgfsetfillopacity{0.130000}%
\pgfsetlinewidth{0.000000pt}%
\definecolor{currentstroke}{rgb}{0.000000,0.000000,0.000000}%
\pgfsetstrokecolor{currentstroke}%
\pgfsetdash{}{0pt}%
\pgfpathmoveto{\pgfqpoint{1.041307in}{1.106910in}}%
\pgfpathlineto{\pgfqpoint{1.056955in}{1.093462in}}%
\pgfpathlineto{\pgfqpoint{1.072622in}{1.080538in}}%
\pgfpathlineto{\pgfqpoint{1.080405in}{1.094479in}}%
\pgfpathlineto{\pgfqpoint{1.088173in}{1.108589in}}%
\pgfpathlineto{\pgfqpoint{1.072527in}{1.120981in}}%
\pgfpathlineto{\pgfqpoint{1.056902in}{1.133893in}}%
\pgfpathlineto{\pgfqpoint{1.049113in}{1.120309in}}%
\pgfpathlineto{\pgfqpoint{1.041307in}{1.106910in}}%
\pgfpathclose%
\pgfusepath{fill}%
\end{pgfscope}%
\begin{pgfscope}%
\pgfpathrectangle{\pgfqpoint{0.407592in}{0.222904in}}{\pgfqpoint{1.411163in}{1.411163in}}%
\pgfusepath{clip}%
\pgfsetbuttcap%
\pgfsetroundjoin%
\definecolor{currentfill}{rgb}{0.556863,0.654902,0.847059}%
\pgfsetfillcolor{currentfill}%
\pgfsetfillopacity{0.130000}%
\pgfsetlinewidth{0.000000pt}%
\definecolor{currentstroke}{rgb}{0.000000,0.000000,0.000000}%
\pgfsetstrokecolor{currentstroke}%
\pgfsetdash{}{0pt}%
\pgfpathmoveto{\pgfqpoint{0.806273in}{1.237483in}}%
\pgfpathlineto{\pgfqpoint{0.821869in}{1.219256in}}%
\pgfpathlineto{\pgfqpoint{0.837474in}{1.201157in}}%
\pgfpathlineto{\pgfqpoint{0.845559in}{1.210633in}}%
\pgfpathlineto{\pgfqpoint{0.853620in}{1.220383in}}%
\pgfpathlineto{\pgfqpoint{0.838052in}{1.238104in}}%
\pgfpathlineto{\pgfqpoint{0.822493in}{1.255969in}}%
\pgfpathlineto{\pgfqpoint{0.814396in}{1.246595in}}%
\pgfpathlineto{\pgfqpoint{0.806273in}{1.237483in}}%
\pgfpathclose%
\pgfusepath{fill}%
\end{pgfscope}%
\begin{pgfscope}%
\pgfpathrectangle{\pgfqpoint{0.407592in}{0.222904in}}{\pgfqpoint{1.411163in}{1.411163in}}%
\pgfusepath{clip}%
\pgfsetbuttcap%
\pgfsetroundjoin%
\definecolor{currentfill}{rgb}{0.556863,0.654902,0.847059}%
\pgfsetfillcolor{currentfill}%
\pgfsetfillopacity{0.130000}%
\pgfsetlinewidth{0.000000pt}%
\definecolor{currentstroke}{rgb}{0.000000,0.000000,0.000000}%
\pgfsetstrokecolor{currentstroke}%
\pgfsetdash{}{0pt}%
\pgfpathmoveto{\pgfqpoint{0.994346in}{1.110128in}}%
\pgfpathlineto{\pgfqpoint{1.009988in}{1.095217in}}%
\pgfpathlineto{\pgfqpoint{1.025645in}{1.080735in}}%
\pgfpathlineto{\pgfqpoint{1.033485in}{1.093713in}}%
\pgfpathlineto{\pgfqpoint{1.041307in}{1.106910in}}%
\pgfpathlineto{\pgfqpoint{1.025676in}{1.120837in}}%
\pgfpathlineto{\pgfqpoint{1.010061in}{1.135201in}}%
\pgfpathlineto{\pgfqpoint{1.002213in}{1.122550in}}%
\pgfpathlineto{\pgfqpoint{0.994346in}{1.110128in}}%
\pgfpathclose%
\pgfusepath{fill}%
\end{pgfscope}%
\begin{pgfscope}%
\pgfpathrectangle{\pgfqpoint{0.407592in}{0.222904in}}{\pgfqpoint{1.411163in}{1.411163in}}%
\pgfusepath{clip}%
\pgfsetbuttcap%
\pgfsetroundjoin%
\definecolor{currentfill}{rgb}{0.556863,0.654902,0.847059}%
\pgfsetfillcolor{currentfill}%
\pgfsetfillopacity{0.130000}%
\pgfsetlinewidth{0.000000pt}%
\definecolor{currentstroke}{rgb}{0.000000,0.000000,0.000000}%
\pgfsetstrokecolor{currentstroke}%
\pgfsetdash{}{0pt}%
\pgfpathmoveto{\pgfqpoint{0.915998in}{1.151260in}}%
\pgfpathlineto{\pgfqpoint{0.931619in}{1.134522in}}%
\pgfpathlineto{\pgfqpoint{0.947252in}{1.118052in}}%
\pgfpathlineto{\pgfqpoint{0.955188in}{1.129431in}}%
\pgfpathlineto{\pgfqpoint{0.963104in}{1.141075in}}%
\pgfpathlineto{\pgfqpoint{0.947503in}{1.157040in}}%
\pgfpathlineto{\pgfqpoint{0.931914in}{1.173292in}}%
\pgfpathlineto{\pgfqpoint{0.923966in}{1.162144in}}%
\pgfpathlineto{\pgfqpoint{0.915998in}{1.151260in}}%
\pgfpathclose%
\pgfusepath{fill}%
\end{pgfscope}%
\begin{pgfscope}%
\pgfpathrectangle{\pgfqpoint{0.407592in}{0.222904in}}{\pgfqpoint{1.411163in}{1.411163in}}%
\pgfusepath{clip}%
\pgfsetbuttcap%
\pgfsetroundjoin%
\definecolor{currentfill}{rgb}{0.556863,0.654902,0.847059}%
\pgfsetfillcolor{currentfill}%
\pgfsetfillopacity{0.130000}%
\pgfsetlinewidth{0.000000pt}%
\definecolor{currentstroke}{rgb}{0.000000,0.000000,0.000000}%
\pgfsetstrokecolor{currentstroke}%
\pgfsetdash{}{0pt}%
\pgfpathmoveto{\pgfqpoint{1.166400in}{1.095585in}}%
\pgfpathlineto{\pgfqpoint{1.182153in}{1.086828in}}%
\pgfpathlineto{\pgfqpoint{1.197937in}{1.078806in}}%
\pgfpathlineto{\pgfqpoint{1.205614in}{1.094541in}}%
\pgfpathlineto{\pgfqpoint{1.213278in}{1.110341in}}%
\pgfpathlineto{\pgfqpoint{1.197506in}{1.118048in}}%
\pgfpathlineto{\pgfqpoint{1.181767in}{1.126452in}}%
\pgfpathlineto{\pgfqpoint{1.174090in}{1.110979in}}%
\pgfpathlineto{\pgfqpoint{1.166400in}{1.095585in}}%
\pgfpathclose%
\pgfusepath{fill}%
\end{pgfscope}%
\begin{pgfscope}%
\pgfpathrectangle{\pgfqpoint{0.407592in}{0.222904in}}{\pgfqpoint{1.411163in}{1.411163in}}%
\pgfusepath{clip}%
\pgfsetbuttcap%
\pgfsetroundjoin%
\definecolor{currentfill}{rgb}{0.556863,0.654902,0.847059}%
\pgfsetfillcolor{currentfill}%
\pgfsetfillopacity{0.130000}%
\pgfsetlinewidth{0.000000pt}%
\definecolor{currentstroke}{rgb}{0.000000,0.000000,0.000000}%
\pgfsetstrokecolor{currentstroke}%
\pgfsetdash{}{0pt}%
\pgfpathmoveto{\pgfqpoint{1.119527in}{1.085537in}}%
\pgfpathlineto{\pgfqpoint{1.135240in}{1.074963in}}%
\pgfpathlineto{\pgfqpoint{1.150979in}{1.065082in}}%
\pgfpathlineto{\pgfqpoint{1.158696in}{1.080282in}}%
\pgfpathlineto{\pgfqpoint{1.166400in}{1.095585in}}%
\pgfpathlineto{\pgfqpoint{1.150676in}{1.105044in}}%
\pgfpathlineto{\pgfqpoint{1.134980in}{1.115167in}}%
\pgfpathlineto{\pgfqpoint{1.127261in}{1.100291in}}%
\pgfpathlineto{\pgfqpoint{1.119527in}{1.085537in}}%
\pgfpathclose%
\pgfusepath{fill}%
\end{pgfscope}%
\begin{pgfscope}%
\pgfpathrectangle{\pgfqpoint{0.407592in}{0.222904in}}{\pgfqpoint{1.411163in}{1.411163in}}%
\pgfusepath{clip}%
\pgfsetbuttcap%
\pgfsetroundjoin%
\definecolor{currentfill}{rgb}{0.556863,0.654902,0.847059}%
\pgfsetfillcolor{currentfill}%
\pgfsetfillopacity{0.130000}%
\pgfsetlinewidth{0.000000pt}%
\definecolor{currentstroke}{rgb}{0.000000,0.000000,0.000000}%
\pgfsetstrokecolor{currentstroke}%
\pgfsetdash{}{0pt}%
\pgfpathmoveto{\pgfqpoint{0.837474in}{1.201157in}}%
\pgfpathlineto{\pgfqpoint{0.853089in}{1.183198in}}%
\pgfpathlineto{\pgfqpoint{0.868714in}{1.165395in}}%
\pgfpathlineto{\pgfqpoint{0.876762in}{1.175272in}}%
\pgfpathlineto{\pgfqpoint{0.884788in}{1.185431in}}%
\pgfpathlineto{\pgfqpoint{0.869199in}{1.202819in}}%
\pgfpathlineto{\pgfqpoint{0.853620in}{1.220383in}}%
\pgfpathlineto{\pgfqpoint{0.845559in}{1.210633in}}%
\pgfpathlineto{\pgfqpoint{0.837474in}{1.201157in}}%
\pgfpathclose%
\pgfusepath{fill}%
\end{pgfscope}%
\begin{pgfscope}%
\pgfpathrectangle{\pgfqpoint{0.407592in}{0.222904in}}{\pgfqpoint{1.411163in}{1.411163in}}%
\pgfusepath{clip}%
\pgfsetbuttcap%
\pgfsetroundjoin%
\definecolor{currentfill}{rgb}{0.556863,0.654902,0.847059}%
\pgfsetfillcolor{currentfill}%
\pgfsetfillopacity{0.130000}%
\pgfsetlinewidth{0.000000pt}%
\definecolor{currentstroke}{rgb}{0.000000,0.000000,0.000000}%
\pgfsetstrokecolor{currentstroke}%
\pgfsetdash{}{0pt}%
\pgfpathmoveto{\pgfqpoint{1.072622in}{1.080538in}}%
\pgfpathlineto{\pgfqpoint{1.088308in}{1.068185in}}%
\pgfpathlineto{\pgfqpoint{1.104016in}{1.056450in}}%
\pgfpathlineto{\pgfqpoint{1.111779in}{1.070918in}}%
\pgfpathlineto{\pgfqpoint{1.119527in}{1.085537in}}%
\pgfpathlineto{\pgfqpoint{1.103839in}{1.096760in}}%
\pgfpathlineto{\pgfqpoint{1.088173in}{1.108589in}}%
\pgfpathlineto{\pgfqpoint{1.080405in}{1.094479in}}%
\pgfpathlineto{\pgfqpoint{1.072622in}{1.080538in}}%
\pgfpathclose%
\pgfusepath{fill}%
\end{pgfscope}%
\begin{pgfscope}%
\pgfpathrectangle{\pgfqpoint{0.407592in}{0.222904in}}{\pgfqpoint{1.411163in}{1.411163in}}%
\pgfusepath{clip}%
\pgfsetbuttcap%
\pgfsetroundjoin%
\definecolor{currentfill}{rgb}{0.556863,0.654902,0.847059}%
\pgfsetfillcolor{currentfill}%
\pgfsetfillopacity{0.130000}%
\pgfsetlinewidth{0.000000pt}%
\definecolor{currentstroke}{rgb}{0.000000,0.000000,0.000000}%
\pgfsetstrokecolor{currentstroke}%
\pgfsetdash{}{0pt}%
\pgfpathmoveto{\pgfqpoint{1.213278in}{1.110341in}}%
\pgfpathlineto{\pgfqpoint{1.229085in}{1.103356in}}%
\pgfpathlineto{\pgfqpoint{1.244927in}{1.097116in}}%
\pgfpathlineto{\pgfqpoint{1.252568in}{1.113201in}}%
\pgfpathlineto{\pgfqpoint{1.236730in}{1.119337in}}%
\pgfpathlineto{\pgfqpoint{1.220930in}{1.126195in}}%
\pgfusepath{fill}%
\end{pgfscope}%
\begin{pgfscope}%
\pgfpathrectangle{\pgfqpoint{0.407592in}{0.222904in}}{\pgfqpoint{1.411163in}{1.411163in}}%
\pgfusepath{clip}%
\pgfsetbuttcap%
\pgfsetroundjoin%
\definecolor{currentfill}{rgb}{0.556863,0.654902,0.847059}%
\pgfsetfillcolor{currentfill}%
\pgfsetfillopacity{0.130000}%
\pgfsetlinewidth{0.000000pt}%
\definecolor{currentstroke}{rgb}{0.000000,0.000000,0.000000}%
\pgfsetstrokecolor{currentstroke}%
\pgfsetdash{}{0pt}%
\pgfpathmoveto{\pgfqpoint{0.947252in}{1.118052in}}%
\pgfpathlineto{\pgfqpoint{0.962897in}{1.101880in}}%
\pgfpathlineto{\pgfqpoint{0.978555in}{1.086041in}}%
\pgfpathlineto{\pgfqpoint{0.986460in}{1.097952in}}%
\pgfpathlineto{\pgfqpoint{0.994346in}{1.110128in}}%
\pgfpathlineto{\pgfqpoint{0.978718in}{1.125426in}}%
\pgfpathlineto{\pgfqpoint{0.963104in}{1.141075in}}%
\pgfpathlineto{\pgfqpoint{0.955188in}{1.129431in}}%
\pgfpathlineto{\pgfqpoint{0.947252in}{1.118052in}}%
\pgfpathclose%
\pgfusepath{fill}%
\end{pgfscope}%
\begin{pgfscope}%
\pgfpathrectangle{\pgfqpoint{0.407592in}{0.222904in}}{\pgfqpoint{1.411163in}{1.411163in}}%
\pgfusepath{clip}%
\pgfsetbuttcap%
\pgfsetroundjoin%
\definecolor{currentfill}{rgb}{0.556863,0.654902,0.847059}%
\pgfsetfillcolor{currentfill}%
\pgfsetfillopacity{0.130000}%
\pgfsetlinewidth{0.000000pt}%
\definecolor{currentstroke}{rgb}{0.000000,0.000000,0.000000}%
\pgfsetstrokecolor{currentstroke}%
\pgfsetdash{}{0pt}%
\pgfpathmoveto{\pgfqpoint{1.025645in}{1.080735in}}%
\pgfpathlineto{\pgfqpoint{1.041317in}{1.066726in}}%
\pgfpathlineto{\pgfqpoint{1.057007in}{1.053238in}}%
\pgfpathlineto{\pgfqpoint{1.064823in}{1.066785in}}%
\pgfpathlineto{\pgfqpoint{1.072622in}{1.080538in}}%
\pgfpathlineto{\pgfqpoint{1.056955in}{1.093462in}}%
\pgfpathlineto{\pgfqpoint{1.041307in}{1.106910in}}%
\pgfpathlineto{\pgfqpoint{1.033485in}{1.093713in}}%
\pgfpathlineto{\pgfqpoint{1.025645in}{1.080735in}}%
\pgfpathclose%
\pgfusepath{fill}%
\end{pgfscope}%
\begin{pgfscope}%
\pgfpathrectangle{\pgfqpoint{0.407592in}{0.222904in}}{\pgfqpoint{1.411163in}{1.411163in}}%
\pgfusepath{clip}%
\pgfsetbuttcap%
\pgfsetroundjoin%
\definecolor{currentfill}{rgb}{0.556863,0.654902,0.847059}%
\pgfsetfillcolor{currentfill}%
\pgfsetfillopacity{0.130000}%
\pgfsetlinewidth{0.000000pt}%
\definecolor{currentstroke}{rgb}{0.000000,0.000000,0.000000}%
\pgfsetstrokecolor{currentstroke}%
\pgfsetdash{}{0pt}%
\pgfpathmoveto{\pgfqpoint{0.868714in}{1.165395in}}%
\pgfpathlineto{\pgfqpoint{0.884348in}{1.147768in}}%
\pgfpathlineto{\pgfqpoint{0.899994in}{1.130337in}}%
\pgfpathlineto{\pgfqpoint{0.908007in}{1.140652in}}%
\pgfpathlineto{\pgfqpoint{0.915998in}{1.151260in}}%
\pgfpathlineto{\pgfqpoint{0.900387in}{1.168237in}}%
\pgfpathlineto{\pgfqpoint{0.884788in}{1.185431in}}%
\pgfpathlineto{\pgfqpoint{0.876762in}{1.175272in}}%
\pgfpathlineto{\pgfqpoint{0.868714in}{1.165395in}}%
\pgfpathclose%
\pgfusepath{fill}%
\end{pgfscope}%
\begin{pgfscope}%
\pgfpathrectangle{\pgfqpoint{0.407592in}{0.222904in}}{\pgfqpoint{1.411163in}{1.411163in}}%
\pgfusepath{clip}%
\pgfsetbuttcap%
\pgfsetroundjoin%
\definecolor{currentfill}{rgb}{0.556863,0.654902,0.847059}%
\pgfsetfillcolor{currentfill}%
\pgfsetfillopacity{0.130000}%
\pgfsetlinewidth{0.000000pt}%
\definecolor{currentstroke}{rgb}{0.000000,0.000000,0.000000}%
\pgfsetstrokecolor{currentstroke}%
\pgfsetdash{}{0pt}%
\pgfpathmoveto{\pgfqpoint{0.758710in}{1.257496in}}%
\pgfpathlineto{\pgfqpoint{0.774326in}{1.238746in}}%
\pgfpathlineto{\pgfqpoint{0.789951in}{1.220083in}}%
\pgfpathlineto{\pgfqpoint{0.798125in}{1.228643in}}%
\pgfpathlineto{\pgfqpoint{0.806273in}{1.237483in}}%
\pgfpathlineto{\pgfqpoint{0.790687in}{1.255823in}}%
\pgfpathlineto{\pgfqpoint{0.775110in}{1.274265in}}%
\pgfpathlineto{\pgfqpoint{0.766923in}{1.265747in}}%
\pgfpathlineto{\pgfqpoint{0.758710in}{1.257496in}}%
\pgfpathclose%
\pgfusepath{fill}%
\end{pgfscope}%
\begin{pgfscope}%
\pgfpathrectangle{\pgfqpoint{0.407592in}{0.222904in}}{\pgfqpoint{1.411163in}{1.411163in}}%
\pgfusepath{clip}%
\pgfsetbuttcap%
\pgfsetroundjoin%
\definecolor{currentfill}{rgb}{0.556863,0.654902,0.847059}%
\pgfsetfillcolor{currentfill}%
\pgfsetfillopacity{0.130000}%
\pgfsetlinewidth{0.000000pt}%
\definecolor{currentstroke}{rgb}{0.000000,0.000000,0.000000}%
\pgfsetstrokecolor{currentstroke}%
\pgfsetdash{}{0pt}%
\pgfpathmoveto{\pgfqpoint{1.150979in}{1.065082in}}%
\pgfpathlineto{\pgfqpoint{1.166746in}{1.055936in}}%
\pgfpathlineto{\pgfqpoint{1.182543in}{1.047567in}}%
\pgfpathlineto{\pgfqpoint{1.190247in}{1.063144in}}%
\pgfpathlineto{\pgfqpoint{1.197937in}{1.078806in}}%
\pgfpathlineto{\pgfqpoint{1.182153in}{1.086828in}}%
\pgfpathlineto{\pgfqpoint{1.166400in}{1.095585in}}%
\pgfpathlineto{\pgfqpoint{1.158696in}{1.080282in}}%
\pgfpathlineto{\pgfqpoint{1.150979in}{1.065082in}}%
\pgfpathclose%
\pgfusepath{fill}%
\end{pgfscope}%
\begin{pgfscope}%
\pgfpathrectangle{\pgfqpoint{0.407592in}{0.222904in}}{\pgfqpoint{1.411163in}{1.411163in}}%
\pgfusepath{clip}%
\pgfsetbuttcap%
\pgfsetroundjoin%
\definecolor{currentfill}{rgb}{0.556863,0.654902,0.847059}%
\pgfsetfillcolor{currentfill}%
\pgfsetfillopacity{0.130000}%
\pgfsetlinewidth{0.000000pt}%
\definecolor{currentstroke}{rgb}{0.000000,0.000000,0.000000}%
\pgfsetstrokecolor{currentstroke}%
\pgfsetdash{}{0pt}%
\pgfpathmoveto{\pgfqpoint{0.978555in}{1.086041in}}%
\pgfpathlineto{\pgfqpoint{0.994226in}{1.070572in}}%
\pgfpathlineto{\pgfqpoint{1.009910in}{1.055518in}}%
\pgfpathlineto{\pgfqpoint{1.017787in}{1.067997in}}%
\pgfpathlineto{\pgfqpoint{1.025645in}{1.080735in}}%
\pgfpathlineto{\pgfqpoint{1.009988in}{1.095217in}}%
\pgfpathlineto{\pgfqpoint{0.994346in}{1.110128in}}%
\pgfpathlineto{\pgfqpoint{0.986460in}{1.097952in}}%
\pgfpathlineto{\pgfqpoint{0.978555in}{1.086041in}}%
\pgfpathclose%
\pgfusepath{fill}%
\end{pgfscope}%
\begin{pgfscope}%
\pgfpathrectangle{\pgfqpoint{0.407592in}{0.222904in}}{\pgfqpoint{1.411163in}{1.411163in}}%
\pgfusepath{clip}%
\pgfsetbuttcap%
\pgfsetroundjoin%
\definecolor{currentfill}{rgb}{0.556863,0.654902,0.847059}%
\pgfsetfillcolor{currentfill}%
\pgfsetfillopacity{0.130000}%
\pgfsetlinewidth{0.000000pt}%
\definecolor{currentstroke}{rgb}{0.000000,0.000000,0.000000}%
\pgfsetstrokecolor{currentstroke}%
\pgfsetdash{}{0pt}%
\pgfpathmoveto{\pgfqpoint{1.104016in}{1.056450in}}%
\pgfpathlineto{\pgfqpoint{1.119746in}{1.045385in}}%
\pgfpathlineto{\pgfqpoint{1.135502in}{1.035041in}}%
\pgfpathlineto{\pgfqpoint{1.143247in}{1.049996in}}%
\pgfpathlineto{\pgfqpoint{1.150979in}{1.065082in}}%
\pgfpathlineto{\pgfqpoint{1.135240in}{1.074963in}}%
\pgfpathlineto{\pgfqpoint{1.119527in}{1.085537in}}%
\pgfpathlineto{\pgfqpoint{1.111779in}{1.070918in}}%
\pgfpathlineto{\pgfqpoint{1.104016in}{1.056450in}}%
\pgfpathclose%
\pgfusepath{fill}%
\end{pgfscope}%
\begin{pgfscope}%
\pgfpathrectangle{\pgfqpoint{0.407592in}{0.222904in}}{\pgfqpoint{1.411163in}{1.411163in}}%
\pgfusepath{clip}%
\pgfsetbuttcap%
\pgfsetroundjoin%
\definecolor{currentfill}{rgb}{0.556863,0.654902,0.847059}%
\pgfsetfillcolor{currentfill}%
\pgfsetfillopacity{0.130000}%
\pgfsetlinewidth{0.000000pt}%
\definecolor{currentstroke}{rgb}{0.000000,0.000000,0.000000}%
\pgfsetstrokecolor{currentstroke}%
\pgfsetdash{}{0pt}%
\pgfpathmoveto{\pgfqpoint{0.789951in}{1.220083in}}%
\pgfpathlineto{\pgfqpoint{0.805585in}{1.201520in}}%
\pgfpathlineto{\pgfqpoint{0.821228in}{1.183066in}}%
\pgfpathlineto{\pgfqpoint{0.829364in}{1.191964in}}%
\pgfpathlineto{\pgfqpoint{0.837474in}{1.201157in}}%
\pgfpathlineto{\pgfqpoint{0.821869in}{1.219256in}}%
\pgfpathlineto{\pgfqpoint{0.806273in}{1.237483in}}%
\pgfpathlineto{\pgfqpoint{0.798125in}{1.228643in}}%
\pgfpathlineto{\pgfqpoint{0.789951in}{1.220083in}}%
\pgfpathclose%
\pgfusepath{fill}%
\end{pgfscope}%
\begin{pgfscope}%
\pgfpathrectangle{\pgfqpoint{0.407592in}{0.222904in}}{\pgfqpoint{1.411163in}{1.411163in}}%
\pgfusepath{clip}%
\pgfsetbuttcap%
\pgfsetroundjoin%
\definecolor{currentfill}{rgb}{0.556863,0.654902,0.847059}%
\pgfsetfillcolor{currentfill}%
\pgfsetfillopacity{0.130000}%
\pgfsetlinewidth{0.000000pt}%
\definecolor{currentstroke}{rgb}{0.000000,0.000000,0.000000}%
\pgfsetstrokecolor{currentstroke}%
\pgfsetdash{}{0pt}%
\pgfpathmoveto{\pgfqpoint{1.197937in}{1.078806in}}%
\pgfpathlineto{\pgfqpoint{1.213755in}{1.071550in}}%
\pgfpathlineto{\pgfqpoint{1.229608in}{1.065087in}}%
\pgfpathlineto{\pgfqpoint{1.237274in}{1.081075in}}%
\pgfpathlineto{\pgfqpoint{1.244927in}{1.097116in}}%
\pgfpathlineto{\pgfqpoint{1.229085in}{1.103356in}}%
\pgfpathlineto{\pgfqpoint{1.213278in}{1.110341in}}%
\pgfpathlineto{\pgfqpoint{1.205614in}{1.094541in}}%
\pgfpathlineto{\pgfqpoint{1.197937in}{1.078806in}}%
\pgfpathclose%
\pgfusepath{fill}%
\end{pgfscope}%
\begin{pgfscope}%
\pgfpathrectangle{\pgfqpoint{0.407592in}{0.222904in}}{\pgfqpoint{1.411163in}{1.411163in}}%
\pgfusepath{clip}%
\pgfsetbuttcap%
\pgfsetroundjoin%
\definecolor{currentfill}{rgb}{0.556863,0.654902,0.847059}%
\pgfsetfillcolor{currentfill}%
\pgfsetfillopacity{0.130000}%
\pgfsetlinewidth{0.000000pt}%
\definecolor{currentstroke}{rgb}{0.000000,0.000000,0.000000}%
\pgfsetstrokecolor{currentstroke}%
\pgfsetdash{}{0pt}%
\pgfpathmoveto{\pgfqpoint{0.899994in}{1.130337in}}%
\pgfpathlineto{\pgfqpoint{0.915649in}{1.113125in}}%
\pgfpathlineto{\pgfqpoint{0.931316in}{1.096159in}}%
\pgfpathlineto{\pgfqpoint{0.939295in}{1.106956in}}%
\pgfpathlineto{\pgfqpoint{0.947252in}{1.118052in}}%
\pgfpathlineto{\pgfqpoint{0.931619in}{1.134522in}}%
\pgfpathlineto{\pgfqpoint{0.915998in}{1.151260in}}%
\pgfpathlineto{\pgfqpoint{0.908007in}{1.140652in}}%
\pgfpathlineto{\pgfqpoint{0.899994in}{1.130337in}}%
\pgfpathclose%
\pgfusepath{fill}%
\end{pgfscope}%
\begin{pgfscope}%
\pgfpathrectangle{\pgfqpoint{0.407592in}{0.222904in}}{\pgfqpoint{1.411163in}{1.411163in}}%
\pgfusepath{clip}%
\pgfsetbuttcap%
\pgfsetroundjoin%
\definecolor{currentfill}{rgb}{0.556863,0.654902,0.847059}%
\pgfsetfillcolor{currentfill}%
\pgfsetfillopacity{0.130000}%
\pgfsetlinewidth{0.000000pt}%
\definecolor{currentstroke}{rgb}{0.000000,0.000000,0.000000}%
\pgfsetstrokecolor{currentstroke}%
\pgfsetdash{}{0pt}%
\pgfpathmoveto{\pgfqpoint{1.057007in}{1.053238in}}%
\pgfpathlineto{\pgfqpoint{1.072716in}{1.040324in}}%
\pgfpathlineto{\pgfqpoint{1.088444in}{1.028039in}}%
\pgfpathlineto{\pgfqpoint{1.096238in}{1.042151in}}%
\pgfpathlineto{\pgfqpoint{1.104016in}{1.056450in}}%
\pgfpathlineto{\pgfqpoint{1.088308in}{1.068185in}}%
\pgfpathlineto{\pgfqpoint{1.072622in}{1.080538in}}%
\pgfpathlineto{\pgfqpoint{1.064823in}{1.066785in}}%
\pgfpathlineto{\pgfqpoint{1.057007in}{1.053238in}}%
\pgfpathclose%
\pgfusepath{fill}%
\end{pgfscope}%
\begin{pgfscope}%
\pgfpathrectangle{\pgfqpoint{0.407592in}{0.222904in}}{\pgfqpoint{1.411163in}{1.411163in}}%
\pgfusepath{clip}%
\pgfsetbuttcap%
\pgfsetroundjoin%
\definecolor{currentfill}{rgb}{0.556863,0.654902,0.847059}%
\pgfsetfillcolor{currentfill}%
\pgfsetfillopacity{0.130000}%
\pgfsetlinewidth{0.000000pt}%
\definecolor{currentstroke}{rgb}{0.000000,0.000000,0.000000}%
\pgfsetstrokecolor{currentstroke}%
\pgfsetdash{}{0pt}%
\pgfpathmoveto{\pgfqpoint{0.821228in}{1.183066in}}%
\pgfpathlineto{\pgfqpoint{0.836881in}{1.164736in}}%
\pgfpathlineto{\pgfqpoint{0.852543in}{1.146543in}}%
\pgfpathlineto{\pgfqpoint{0.860641in}{1.155815in}}%
\pgfpathlineto{\pgfqpoint{0.868714in}{1.165395in}}%
\pgfpathlineto{\pgfqpoint{0.853089in}{1.183198in}}%
\pgfpathlineto{\pgfqpoint{0.837474in}{1.201157in}}%
\pgfpathlineto{\pgfqpoint{0.829364in}{1.191964in}}%
\pgfpathlineto{\pgfqpoint{0.821228in}{1.183066in}}%
\pgfpathclose%
\pgfusepath{fill}%
\end{pgfscope}%
\begin{pgfscope}%
\pgfpathrectangle{\pgfqpoint{0.407592in}{0.222904in}}{\pgfqpoint{1.411163in}{1.411163in}}%
\pgfusepath{clip}%
\pgfsetbuttcap%
\pgfsetroundjoin%
\definecolor{currentfill}{rgb}{0.556863,0.654902,0.847059}%
\pgfsetfillcolor{currentfill}%
\pgfsetfillopacity{0.130000}%
\pgfsetlinewidth{0.000000pt}%
\definecolor{currentstroke}{rgb}{0.000000,0.000000,0.000000}%
\pgfsetstrokecolor{currentstroke}%
\pgfsetdash{}{0pt}%
\pgfpathmoveto{\pgfqpoint{1.244927in}{1.097116in}}%
\pgfpathlineto{\pgfqpoint{1.260808in}{1.091640in}}%
\pgfpathlineto{\pgfqpoint{1.276730in}{1.086941in}}%
\pgfpathlineto{\pgfqpoint{1.284361in}{1.103162in}}%
\pgfpathlineto{\pgfqpoint{1.268444in}{1.107805in}}%
\pgfpathlineto{\pgfqpoint{1.252568in}{1.113201in}}%
\pgfusepath{fill}%
\end{pgfscope}%
\begin{pgfscope}%
\pgfpathrectangle{\pgfqpoint{0.407592in}{0.222904in}}{\pgfqpoint{1.411163in}{1.411163in}}%
\pgfusepath{clip}%
\pgfsetbuttcap%
\pgfsetroundjoin%
\definecolor{currentfill}{rgb}{0.556863,0.654902,0.847059}%
\pgfsetfillcolor{currentfill}%
\pgfsetfillopacity{0.130000}%
\pgfsetlinewidth{0.000000pt}%
\definecolor{currentstroke}{rgb}{0.000000,0.000000,0.000000}%
\pgfsetstrokecolor{currentstroke}%
\pgfsetdash{}{0pt}%
\pgfpathmoveto{\pgfqpoint{0.931316in}{1.096159in}}%
\pgfpathlineto{\pgfqpoint{0.946994in}{1.079469in}}%
\pgfpathlineto{\pgfqpoint{0.962683in}{1.063091in}}%
\pgfpathlineto{\pgfqpoint{0.970629in}{1.074413in}}%
\pgfpathlineto{\pgfqpoint{0.978555in}{1.086041in}}%
\pgfpathlineto{\pgfqpoint{0.962897in}{1.101880in}}%
\pgfpathlineto{\pgfqpoint{0.947252in}{1.118052in}}%
\pgfpathlineto{\pgfqpoint{0.939295in}{1.106956in}}%
\pgfpathlineto{\pgfqpoint{0.931316in}{1.096159in}}%
\pgfpathclose%
\pgfusepath{fill}%
\end{pgfscope}%
\begin{pgfscope}%
\pgfpathrectangle{\pgfqpoint{0.407592in}{0.222904in}}{\pgfqpoint{1.411163in}{1.411163in}}%
\pgfusepath{clip}%
\pgfsetbuttcap%
\pgfsetroundjoin%
\definecolor{currentfill}{rgb}{0.556863,0.654902,0.847059}%
\pgfsetfillcolor{currentfill}%
\pgfsetfillopacity{0.130000}%
\pgfsetlinewidth{0.000000pt}%
\definecolor{currentstroke}{rgb}{0.000000,0.000000,0.000000}%
\pgfsetstrokecolor{currentstroke}%
\pgfsetdash{}{0pt}%
\pgfpathmoveto{\pgfqpoint{1.009910in}{1.055518in}}%
\pgfpathlineto{\pgfqpoint{1.025610in}{1.040925in}}%
\pgfpathlineto{\pgfqpoint{1.041325in}{1.026847in}}%
\pgfpathlineto{\pgfqpoint{1.049175in}{1.039918in}}%
\pgfpathlineto{\pgfqpoint{1.057007in}{1.053238in}}%
\pgfpathlineto{\pgfqpoint{1.041317in}{1.066726in}}%
\pgfpathlineto{\pgfqpoint{1.025645in}{1.080735in}}%
\pgfpathlineto{\pgfqpoint{1.017787in}{1.067997in}}%
\pgfpathlineto{\pgfqpoint{1.009910in}{1.055518in}}%
\pgfpathclose%
\pgfusepath{fill}%
\end{pgfscope}%
\begin{pgfscope}%
\pgfpathrectangle{\pgfqpoint{0.407592in}{0.222904in}}{\pgfqpoint{1.411163in}{1.411163in}}%
\pgfusepath{clip}%
\pgfsetbuttcap%
\pgfsetroundjoin%
\definecolor{currentfill}{rgb}{0.556863,0.654902,0.847059}%
\pgfsetfillcolor{currentfill}%
\pgfsetfillopacity{0.130000}%
\pgfsetlinewidth{0.000000pt}%
\definecolor{currentstroke}{rgb}{0.000000,0.000000,0.000000}%
\pgfsetstrokecolor{currentstroke}%
\pgfsetdash{}{0pt}%
\pgfpathmoveto{\pgfqpoint{0.852543in}{1.146543in}}%
\pgfpathlineto{\pgfqpoint{0.868215in}{1.128507in}}%
\pgfpathlineto{\pgfqpoint{0.883896in}{1.110644in}}%
\pgfpathlineto{\pgfqpoint{0.891957in}{1.120329in}}%
\pgfpathlineto{\pgfqpoint{0.899994in}{1.130337in}}%
\pgfpathlineto{\pgfqpoint{0.884348in}{1.147768in}}%
\pgfpathlineto{\pgfqpoint{0.868714in}{1.165395in}}%
\pgfpathlineto{\pgfqpoint{0.860641in}{1.155815in}}%
\pgfpathlineto{\pgfqpoint{0.852543in}{1.146543in}}%
\pgfpathclose%
\pgfusepath{fill}%
\end{pgfscope}%
\begin{pgfscope}%
\pgfpathrectangle{\pgfqpoint{0.407592in}{0.222904in}}{\pgfqpoint{1.411163in}{1.411163in}}%
\pgfusepath{clip}%
\pgfsetbuttcap%
\pgfsetroundjoin%
\definecolor{currentfill}{rgb}{0.556863,0.654902,0.847059}%
\pgfsetfillcolor{currentfill}%
\pgfsetfillopacity{0.130000}%
\pgfsetlinewidth{0.000000pt}%
\definecolor{currentstroke}{rgb}{0.000000,0.000000,0.000000}%
\pgfsetstrokecolor{currentstroke}%
\pgfsetdash{}{0pt}%
\pgfpathmoveto{\pgfqpoint{1.135502in}{1.035041in}}%
\pgfpathlineto{\pgfqpoint{1.151283in}{1.025468in}}%
\pgfpathlineto{\pgfqpoint{1.167094in}{1.016715in}}%
\pgfpathlineto{\pgfqpoint{1.174825in}{1.032086in}}%
\pgfpathlineto{\pgfqpoint{1.182543in}{1.047567in}}%
\pgfpathlineto{\pgfqpoint{1.166746in}{1.055936in}}%
\pgfpathlineto{\pgfqpoint{1.150979in}{1.065082in}}%
\pgfpathlineto{\pgfqpoint{1.143247in}{1.049996in}}%
\pgfpathlineto{\pgfqpoint{1.135502in}{1.035041in}}%
\pgfpathclose%
\pgfusepath{fill}%
\end{pgfscope}%
\begin{pgfscope}%
\pgfpathrectangle{\pgfqpoint{0.407592in}{0.222904in}}{\pgfqpoint{1.411163in}{1.411163in}}%
\pgfusepath{clip}%
\pgfsetbuttcap%
\pgfsetroundjoin%
\definecolor{currentfill}{rgb}{0.556863,0.654902,0.847059}%
\pgfsetfillcolor{currentfill}%
\pgfsetfillopacity{0.130000}%
\pgfsetlinewidth{0.000000pt}%
\definecolor{currentstroke}{rgb}{0.000000,0.000000,0.000000}%
\pgfsetstrokecolor{currentstroke}%
\pgfsetdash{}{0pt}%
\pgfpathmoveto{\pgfqpoint{1.182543in}{1.047567in}}%
\pgfpathlineto{\pgfqpoint{1.198372in}{1.040011in}}%
\pgfpathlineto{\pgfqpoint{1.214237in}{1.033302in}}%
\pgfpathlineto{\pgfqpoint{1.221929in}{1.049159in}}%
\pgfpathlineto{\pgfqpoint{1.229608in}{1.065087in}}%
\pgfpathlineto{\pgfqpoint{1.213755in}{1.071550in}}%
\pgfpathlineto{\pgfqpoint{1.197937in}{1.078806in}}%
\pgfpathlineto{\pgfqpoint{1.190247in}{1.063144in}}%
\pgfpathlineto{\pgfqpoint{1.182543in}{1.047567in}}%
\pgfpathclose%
\pgfusepath{fill}%
\end{pgfscope}%
\begin{pgfscope}%
\pgfpathrectangle{\pgfqpoint{0.407592in}{0.222904in}}{\pgfqpoint{1.411163in}{1.411163in}}%
\pgfusepath{clip}%
\pgfsetbuttcap%
\pgfsetroundjoin%
\definecolor{currentfill}{rgb}{0.556863,0.654902,0.847059}%
\pgfsetfillcolor{currentfill}%
\pgfsetfillopacity{0.130000}%
\pgfsetlinewidth{0.000000pt}%
\definecolor{currentstroke}{rgb}{0.000000,0.000000,0.000000}%
\pgfsetstrokecolor{currentstroke}%
\pgfsetdash{}{0pt}%
\pgfpathmoveto{\pgfqpoint{0.742202in}{1.241828in}}%
\pgfpathlineto{\pgfqpoint{0.757857in}{1.222798in}}%
\pgfpathlineto{\pgfqpoint{0.773522in}{1.203843in}}%
\pgfpathlineto{\pgfqpoint{0.781750in}{1.211814in}}%
\pgfpathlineto{\pgfqpoint{0.789951in}{1.220083in}}%
\pgfpathlineto{\pgfqpoint{0.774326in}{1.238746in}}%
\pgfpathlineto{\pgfqpoint{0.758710in}{1.257496in}}%
\pgfpathlineto{\pgfqpoint{0.750470in}{1.249521in}}%
\pgfpathlineto{\pgfqpoint{0.742202in}{1.241828in}}%
\pgfpathclose%
\pgfusepath{fill}%
\end{pgfscope}%
\begin{pgfscope}%
\pgfpathrectangle{\pgfqpoint{0.407592in}{0.222904in}}{\pgfqpoint{1.411163in}{1.411163in}}%
\pgfusepath{clip}%
\pgfsetbuttcap%
\pgfsetroundjoin%
\definecolor{currentfill}{rgb}{0.556863,0.654902,0.847059}%
\pgfsetfillcolor{currentfill}%
\pgfsetfillopacity{0.130000}%
\pgfsetlinewidth{0.000000pt}%
\definecolor{currentstroke}{rgb}{0.000000,0.000000,0.000000}%
\pgfsetstrokecolor{currentstroke}%
\pgfsetdash{}{0pt}%
\pgfpathmoveto{\pgfqpoint{1.088444in}{1.028039in}}%
\pgfpathlineto{\pgfqpoint{1.104193in}{1.016442in}}%
\pgfpathlineto{\pgfqpoint{1.119966in}{1.005592in}}%
\pgfpathlineto{\pgfqpoint{1.127741in}{1.020233in}}%
\pgfpathlineto{\pgfqpoint{1.135502in}{1.035041in}}%
\pgfpathlineto{\pgfqpoint{1.119746in}{1.045385in}}%
\pgfpathlineto{\pgfqpoint{1.104016in}{1.056450in}}%
\pgfpathlineto{\pgfqpoint{1.096238in}{1.042151in}}%
\pgfpathlineto{\pgfqpoint{1.088444in}{1.028039in}}%
\pgfpathclose%
\pgfusepath{fill}%
\end{pgfscope}%
\begin{pgfscope}%
\pgfpathrectangle{\pgfqpoint{0.407592in}{0.222904in}}{\pgfqpoint{1.411163in}{1.411163in}}%
\pgfusepath{clip}%
\pgfsetbuttcap%
\pgfsetroundjoin%
\definecolor{currentfill}{rgb}{0.556863,0.654902,0.847059}%
\pgfsetfillcolor{currentfill}%
\pgfsetfillopacity{0.130000}%
\pgfsetlinewidth{0.000000pt}%
\definecolor{currentstroke}{rgb}{0.000000,0.000000,0.000000}%
\pgfsetstrokecolor{currentstroke}%
\pgfsetdash{}{0pt}%
\pgfpathmoveto{\pgfqpoint{0.962683in}{1.063091in}}%
\pgfpathlineto{\pgfqpoint{0.978385in}{1.047062in}}%
\pgfpathlineto{\pgfqpoint{0.994100in}{1.031429in}}%
\pgfpathlineto{\pgfqpoint{1.002015in}{1.043321in}}%
\pgfpathlineto{\pgfqpoint{1.009910in}{1.055518in}}%
\pgfpathlineto{\pgfqpoint{0.994226in}{1.070572in}}%
\pgfpathlineto{\pgfqpoint{0.978555in}{1.086041in}}%
\pgfpathlineto{\pgfqpoint{0.970629in}{1.074413in}}%
\pgfpathlineto{\pgfqpoint{0.962683in}{1.063091in}}%
\pgfpathclose%
\pgfusepath{fill}%
\end{pgfscope}%
\begin{pgfscope}%
\pgfpathrectangle{\pgfqpoint{0.407592in}{0.222904in}}{\pgfqpoint{1.411163in}{1.411163in}}%
\pgfusepath{clip}%
\pgfsetbuttcap%
\pgfsetroundjoin%
\definecolor{currentfill}{rgb}{0.556863,0.654902,0.847059}%
\pgfsetfillcolor{currentfill}%
\pgfsetfillopacity{0.130000}%
\pgfsetlinewidth{0.000000pt}%
\definecolor{currentstroke}{rgb}{0.000000,0.000000,0.000000}%
\pgfsetstrokecolor{currentstroke}%
\pgfsetdash{}{0pt}%
\pgfpathmoveto{\pgfqpoint{1.229608in}{1.065087in}}%
\pgfpathlineto{\pgfqpoint{1.245499in}{1.059439in}}%
\pgfpathlineto{\pgfqpoint{1.261431in}{1.054624in}}%
\pgfpathlineto{\pgfqpoint{1.269087in}{1.070760in}}%
\pgfpathlineto{\pgfqpoint{1.276730in}{1.086941in}}%
\pgfpathlineto{\pgfqpoint{1.260808in}{1.091640in}}%
\pgfpathlineto{\pgfqpoint{1.244927in}{1.097116in}}%
\pgfpathlineto{\pgfqpoint{1.237274in}{1.081075in}}%
\pgfpathlineto{\pgfqpoint{1.229608in}{1.065087in}}%
\pgfpathclose%
\pgfusepath{fill}%
\end{pgfscope}%
\begin{pgfscope}%
\pgfpathrectangle{\pgfqpoint{0.407592in}{0.222904in}}{\pgfqpoint{1.411163in}{1.411163in}}%
\pgfusepath{clip}%
\pgfsetbuttcap%
\pgfsetroundjoin%
\definecolor{currentfill}{rgb}{0.556863,0.654902,0.847059}%
\pgfsetfillcolor{currentfill}%
\pgfsetfillopacity{0.130000}%
\pgfsetlinewidth{0.000000pt}%
\definecolor{currentstroke}{rgb}{0.000000,0.000000,0.000000}%
\pgfsetstrokecolor{currentstroke}%
\pgfsetdash{}{0pt}%
\pgfpathmoveto{\pgfqpoint{1.041325in}{1.026847in}}%
\pgfpathlineto{\pgfqpoint{1.057057in}{1.013340in}}%
\pgfpathlineto{\pgfqpoint{1.072808in}{1.000468in}}%
\pgfpathlineto{\pgfqpoint{1.080634in}{1.014137in}}%
\pgfpathlineto{\pgfqpoint{1.088444in}{1.028039in}}%
\pgfpathlineto{\pgfqpoint{1.072716in}{1.040324in}}%
\pgfpathlineto{\pgfqpoint{1.057007in}{1.053238in}}%
\pgfpathlineto{\pgfqpoint{1.049175in}{1.039918in}}%
\pgfpathlineto{\pgfqpoint{1.041325in}{1.026847in}}%
\pgfpathclose%
\pgfusepath{fill}%
\end{pgfscope}%
\begin{pgfscope}%
\pgfpathrectangle{\pgfqpoint{0.407592in}{0.222904in}}{\pgfqpoint{1.411163in}{1.411163in}}%
\pgfusepath{clip}%
\pgfsetbuttcap%
\pgfsetroundjoin%
\definecolor{currentfill}{rgb}{0.556863,0.654902,0.847059}%
\pgfsetfillcolor{currentfill}%
\pgfsetfillopacity{0.130000}%
\pgfsetlinewidth{0.000000pt}%
\definecolor{currentstroke}{rgb}{0.000000,0.000000,0.000000}%
\pgfsetstrokecolor{currentstroke}%
\pgfsetdash{}{0pt}%
\pgfpathmoveto{\pgfqpoint{0.883896in}{1.110644in}}%
\pgfpathlineto{\pgfqpoint{0.899588in}{1.092980in}}%
\pgfpathlineto{\pgfqpoint{0.915290in}{1.075537in}}%
\pgfpathlineto{\pgfqpoint{0.923314in}{1.085680in}}%
\pgfpathlineto{\pgfqpoint{0.931316in}{1.096159in}}%
\pgfpathlineto{\pgfqpoint{0.915649in}{1.113125in}}%
\pgfpathlineto{\pgfqpoint{0.899994in}{1.130337in}}%
\pgfpathlineto{\pgfqpoint{0.891957in}{1.120329in}}%
\pgfpathlineto{\pgfqpoint{0.883896in}{1.110644in}}%
\pgfpathclose%
\pgfusepath{fill}%
\end{pgfscope}%
\begin{pgfscope}%
\pgfpathrectangle{\pgfqpoint{0.407592in}{0.222904in}}{\pgfqpoint{1.411163in}{1.411163in}}%
\pgfusepath{clip}%
\pgfsetbuttcap%
\pgfsetroundjoin%
\definecolor{currentfill}{rgb}{0.556863,0.654902,0.847059}%
\pgfsetfillcolor{currentfill}%
\pgfsetfillopacity{0.130000}%
\pgfsetlinewidth{0.000000pt}%
\definecolor{currentstroke}{rgb}{0.000000,0.000000,0.000000}%
\pgfsetstrokecolor{currentstroke}%
\pgfsetdash{}{0pt}%
\pgfpathmoveto{\pgfqpoint{0.773522in}{1.203843in}}%
\pgfpathlineto{\pgfqpoint{0.789195in}{1.184972in}}%
\pgfpathlineto{\pgfqpoint{0.804877in}{1.166196in}}%
\pgfpathlineto{\pgfqpoint{0.813066in}{1.174473in}}%
\pgfpathlineto{\pgfqpoint{0.821228in}{1.183066in}}%
\pgfpathlineto{\pgfqpoint{0.805585in}{1.201520in}}%
\pgfpathlineto{\pgfqpoint{0.789951in}{1.220083in}}%
\pgfpathlineto{\pgfqpoint{0.781750in}{1.211814in}}%
\pgfpathlineto{\pgfqpoint{0.773522in}{1.203843in}}%
\pgfpathclose%
\pgfusepath{fill}%
\end{pgfscope}%
\begin{pgfscope}%
\pgfpathrectangle{\pgfqpoint{0.407592in}{0.222904in}}{\pgfqpoint{1.411163in}{1.411163in}}%
\pgfusepath{clip}%
\pgfsetbuttcap%
\pgfsetroundjoin%
\definecolor{currentfill}{rgb}{0.556863,0.654902,0.847059}%
\pgfsetfillcolor{currentfill}%
\pgfsetfillopacity{0.130000}%
\pgfsetlinewidth{0.000000pt}%
\definecolor{currentstroke}{rgb}{0.000000,0.000000,0.000000}%
\pgfsetstrokecolor{currentstroke}%
\pgfsetdash{}{0pt}%
\pgfpathmoveto{\pgfqpoint{1.276730in}{1.086941in}}%
\pgfpathlineto{\pgfqpoint{1.292696in}{1.083029in}}%
\pgfpathlineto{\pgfqpoint{1.308706in}{1.079908in}}%
\pgfpathlineto{\pgfqpoint{1.316328in}{1.096162in}}%
\pgfpathlineto{\pgfqpoint{1.300322in}{1.099279in}}%
\pgfpathlineto{\pgfqpoint{1.284361in}{1.103162in}}%
\pgfusepath{fill}%
\end{pgfscope}%
\begin{pgfscope}%
\pgfpathrectangle{\pgfqpoint{0.407592in}{0.222904in}}{\pgfqpoint{1.411163in}{1.411163in}}%
\pgfusepath{clip}%
\pgfsetbuttcap%
\pgfsetroundjoin%
\definecolor{currentfill}{rgb}{0.556863,0.654902,0.847059}%
\pgfsetfillcolor{currentfill}%
\pgfsetfillopacity{0.130000}%
\pgfsetlinewidth{0.000000pt}%
\definecolor{currentstroke}{rgb}{0.000000,0.000000,0.000000}%
\pgfsetstrokecolor{currentstroke}%
\pgfsetdash{}{0pt}%
\pgfpathmoveto{\pgfqpoint{0.994100in}{1.031429in}}%
\pgfpathlineto{\pgfqpoint{1.009828in}{1.016239in}}%
\pgfpathlineto{\pgfqpoint{1.025571in}{1.001551in}}%
\pgfpathlineto{\pgfqpoint{1.033457in}{1.014049in}}%
\pgfpathlineto{\pgfqpoint{1.041325in}{1.026847in}}%
\pgfpathlineto{\pgfqpoint{1.025610in}{1.040925in}}%
\pgfpathlineto{\pgfqpoint{1.009910in}{1.055518in}}%
\pgfpathlineto{\pgfqpoint{1.002015in}{1.043321in}}%
\pgfpathlineto{\pgfqpoint{0.994100in}{1.031429in}}%
\pgfpathclose%
\pgfusepath{fill}%
\end{pgfscope}%
\begin{pgfscope}%
\pgfpathrectangle{\pgfqpoint{0.407592in}{0.222904in}}{\pgfqpoint{1.411163in}{1.411163in}}%
\pgfusepath{clip}%
\pgfsetbuttcap%
\pgfsetroundjoin%
\definecolor{currentfill}{rgb}{0.556863,0.654902,0.847059}%
\pgfsetfillcolor{currentfill}%
\pgfsetfillopacity{0.130000}%
\pgfsetlinewidth{0.000000pt}%
\definecolor{currentstroke}{rgb}{0.000000,0.000000,0.000000}%
\pgfsetstrokecolor{currentstroke}%
\pgfsetdash{}{0pt}%
\pgfpathmoveto{\pgfqpoint{0.804877in}{1.166196in}}%
\pgfpathlineto{\pgfqpoint{0.820569in}{1.147526in}}%
\pgfpathlineto{\pgfqpoint{0.836270in}{1.128977in}}%
\pgfpathlineto{\pgfqpoint{0.844419in}{1.137593in}}%
\pgfpathlineto{\pgfqpoint{0.852543in}{1.146543in}}%
\pgfpathlineto{\pgfqpoint{0.836881in}{1.164736in}}%
\pgfpathlineto{\pgfqpoint{0.821228in}{1.183066in}}%
\pgfpathlineto{\pgfqpoint{0.813066in}{1.174473in}}%
\pgfpathlineto{\pgfqpoint{0.804877in}{1.166196in}}%
\pgfpathclose%
\pgfusepath{fill}%
\end{pgfscope}%
\begin{pgfscope}%
\pgfpathrectangle{\pgfqpoint{0.407592in}{0.222904in}}{\pgfqpoint{1.411163in}{1.411163in}}%
\pgfusepath{clip}%
\pgfsetbuttcap%
\pgfsetroundjoin%
\definecolor{currentfill}{rgb}{0.556863,0.654902,0.847059}%
\pgfsetfillcolor{currentfill}%
\pgfsetfillopacity{0.130000}%
\pgfsetlinewidth{0.000000pt}%
\definecolor{currentstroke}{rgb}{0.000000,0.000000,0.000000}%
\pgfsetstrokecolor{currentstroke}%
\pgfsetdash{}{0pt}%
\pgfpathmoveto{\pgfqpoint{0.915290in}{1.075537in}}%
\pgfpathlineto{\pgfqpoint{0.931002in}{1.058348in}}%
\pgfpathlineto{\pgfqpoint{0.946726in}{1.041444in}}%
\pgfpathlineto{\pgfqpoint{0.954716in}{1.052094in}}%
\pgfpathlineto{\pgfqpoint{0.962683in}{1.063091in}}%
\pgfpathlineto{\pgfqpoint{0.946994in}{1.079469in}}%
\pgfpathlineto{\pgfqpoint{0.931316in}{1.096159in}}%
\pgfpathlineto{\pgfqpoint{0.923314in}{1.085680in}}%
\pgfpathlineto{\pgfqpoint{0.915290in}{1.075537in}}%
\pgfpathclose%
\pgfusepath{fill}%
\end{pgfscope}%
\begin{pgfscope}%
\pgfpathrectangle{\pgfqpoint{0.407592in}{0.222904in}}{\pgfqpoint{1.411163in}{1.411163in}}%
\pgfusepath{clip}%
\pgfsetbuttcap%
\pgfsetroundjoin%
\definecolor{currentfill}{rgb}{0.556863,0.654902,0.847059}%
\pgfsetfillcolor{currentfill}%
\pgfsetfillopacity{0.130000}%
\pgfsetlinewidth{0.000000pt}%
\definecolor{currentstroke}{rgb}{0.000000,0.000000,0.000000}%
\pgfsetstrokecolor{currentstroke}%
\pgfsetdash{}{0pt}%
\pgfpathmoveto{\pgfqpoint{1.119966in}{1.005592in}}%
\pgfpathlineto{\pgfqpoint{1.135765in}{0.995548in}}%
\pgfpathlineto{\pgfqpoint{1.151591in}{0.986371in}}%
\pgfpathlineto{\pgfqpoint{1.159350in}{1.001471in}}%
\pgfpathlineto{\pgfqpoint{1.167094in}{1.016715in}}%
\pgfpathlineto{\pgfqpoint{1.151283in}{1.025468in}}%
\pgfpathlineto{\pgfqpoint{1.135502in}{1.035041in}}%
\pgfpathlineto{\pgfqpoint{1.127741in}{1.020233in}}%
\pgfpathlineto{\pgfqpoint{1.119966in}{1.005592in}}%
\pgfpathclose%
\pgfusepath{fill}%
\end{pgfscope}%
\begin{pgfscope}%
\pgfpathrectangle{\pgfqpoint{0.407592in}{0.222904in}}{\pgfqpoint{1.411163in}{1.411163in}}%
\pgfusepath{clip}%
\pgfsetbuttcap%
\pgfsetroundjoin%
\definecolor{currentfill}{rgb}{0.556863,0.654902,0.847059}%
\pgfsetfillcolor{currentfill}%
\pgfsetfillopacity{0.130000}%
\pgfsetlinewidth{0.000000pt}%
\definecolor{currentstroke}{rgb}{0.000000,0.000000,0.000000}%
\pgfsetstrokecolor{currentstroke}%
\pgfsetdash{}{0pt}%
\pgfpathmoveto{\pgfqpoint{1.167094in}{1.016715in}}%
\pgfpathlineto{\pgfqpoint{1.182937in}{1.008828in}}%
\pgfpathlineto{\pgfqpoint{1.198813in}{1.001845in}}%
\pgfpathlineto{\pgfqpoint{1.206532in}{1.017526in}}%
\pgfpathlineto{\pgfqpoint{1.214237in}{1.033302in}}%
\pgfpathlineto{\pgfqpoint{1.198372in}{1.040011in}}%
\pgfpathlineto{\pgfqpoint{1.182543in}{1.047567in}}%
\pgfpathlineto{\pgfqpoint{1.174825in}{1.032086in}}%
\pgfpathlineto{\pgfqpoint{1.167094in}{1.016715in}}%
\pgfpathclose%
\pgfusepath{fill}%
\end{pgfscope}%
\begin{pgfscope}%
\pgfpathrectangle{\pgfqpoint{0.407592in}{0.222904in}}{\pgfqpoint{1.411163in}{1.411163in}}%
\pgfusepath{clip}%
\pgfsetbuttcap%
\pgfsetroundjoin%
\definecolor{currentfill}{rgb}{0.556863,0.654902,0.847059}%
\pgfsetfillcolor{currentfill}%
\pgfsetfillopacity{0.130000}%
\pgfsetlinewidth{0.000000pt}%
\definecolor{currentstroke}{rgb}{0.000000,0.000000,0.000000}%
\pgfsetstrokecolor{currentstroke}%
\pgfsetdash{}{0pt}%
\pgfpathmoveto{\pgfqpoint{1.072808in}{1.000468in}}%
\pgfpathlineto{\pgfqpoint{1.088579in}{0.988297in}}%
\pgfpathlineto{\pgfqpoint{1.104371in}{0.976896in}}%
\pgfpathlineto{\pgfqpoint{1.112176in}{0.991138in}}%
\pgfpathlineto{\pgfqpoint{1.119966in}{1.005592in}}%
\pgfpathlineto{\pgfqpoint{1.104193in}{1.016442in}}%
\pgfpathlineto{\pgfqpoint{1.088444in}{1.028039in}}%
\pgfpathlineto{\pgfqpoint{1.080634in}{1.014137in}}%
\pgfpathlineto{\pgfqpoint{1.072808in}{1.000468in}}%
\pgfpathclose%
\pgfusepath{fill}%
\end{pgfscope}%
\begin{pgfscope}%
\pgfpathrectangle{\pgfqpoint{0.407592in}{0.222904in}}{\pgfqpoint{1.411163in}{1.411163in}}%
\pgfusepath{clip}%
\pgfsetbuttcap%
\pgfsetroundjoin%
\definecolor{currentfill}{rgb}{0.556863,0.654902,0.847059}%
\pgfsetfillcolor{currentfill}%
\pgfsetfillopacity{0.130000}%
\pgfsetlinewidth{0.000000pt}%
\definecolor{currentstroke}{rgb}{0.000000,0.000000,0.000000}%
\pgfsetstrokecolor{currentstroke}%
\pgfsetdash{}{0pt}%
\pgfpathmoveto{\pgfqpoint{1.214237in}{1.033302in}}%
\pgfpathlineto{\pgfqpoint{1.230139in}{1.027466in}}%
\pgfpathlineto{\pgfqpoint{1.246081in}{1.022524in}}%
\pgfpathlineto{\pgfqpoint{1.253762in}{1.038542in}}%
\pgfpathlineto{\pgfqpoint{1.261431in}{1.054624in}}%
\pgfpathlineto{\pgfqpoint{1.245499in}{1.059439in}}%
\pgfpathlineto{\pgfqpoint{1.229608in}{1.065087in}}%
\pgfpathlineto{\pgfqpoint{1.221929in}{1.049159in}}%
\pgfpathlineto{\pgfqpoint{1.214237in}{1.033302in}}%
\pgfpathclose%
\pgfusepath{fill}%
\end{pgfscope}%
\begin{pgfscope}%
\pgfpathrectangle{\pgfqpoint{0.407592in}{0.222904in}}{\pgfqpoint{1.411163in}{1.411163in}}%
\pgfusepath{clip}%
\pgfsetbuttcap%
\pgfsetroundjoin%
\definecolor{currentfill}{rgb}{0.556863,0.654902,0.847059}%
\pgfsetfillcolor{currentfill}%
\pgfsetfillopacity{0.130000}%
\pgfsetlinewidth{0.000000pt}%
\definecolor{currentstroke}{rgb}{0.000000,0.000000,0.000000}%
\pgfsetstrokecolor{currentstroke}%
\pgfsetdash{}{0pt}%
\pgfpathmoveto{\pgfqpoint{0.836270in}{1.128977in}}%
\pgfpathlineto{\pgfqpoint{0.851980in}{1.110563in}}%
\pgfpathlineto{\pgfqpoint{0.867700in}{1.092303in}}%
\pgfpathlineto{\pgfqpoint{0.875811in}{1.101297in}}%
\pgfpathlineto{\pgfqpoint{0.883896in}{1.110644in}}%
\pgfpathlineto{\pgfqpoint{0.868215in}{1.128507in}}%
\pgfpathlineto{\pgfqpoint{0.852543in}{1.146543in}}%
\pgfpathlineto{\pgfqpoint{0.844419in}{1.137593in}}%
\pgfpathlineto{\pgfqpoint{0.836270in}{1.128977in}}%
\pgfpathclose%
\pgfusepath{fill}%
\end{pgfscope}%
\begin{pgfscope}%
\pgfpathrectangle{\pgfqpoint{0.407592in}{0.222904in}}{\pgfqpoint{1.411163in}{1.411163in}}%
\pgfusepath{clip}%
\pgfsetbuttcap%
\pgfsetroundjoin%
\definecolor{currentfill}{rgb}{0.556863,0.654902,0.847059}%
\pgfsetfillcolor{currentfill}%
\pgfsetfillopacity{0.130000}%
\pgfsetlinewidth{0.000000pt}%
\definecolor{currentstroke}{rgb}{0.000000,0.000000,0.000000}%
\pgfsetstrokecolor{currentstroke}%
\pgfsetdash{}{0pt}%
\pgfpathmoveto{\pgfqpoint{0.946726in}{1.041444in}}%
\pgfpathlineto{\pgfqpoint{0.962461in}{1.024867in}}%
\pgfpathlineto{\pgfqpoint{0.978208in}{1.008659in}}%
\pgfpathlineto{\pgfqpoint{0.986164in}{1.019866in}}%
\pgfpathlineto{\pgfqpoint{0.994100in}{1.031429in}}%
\pgfpathlineto{\pgfqpoint{0.978385in}{1.047062in}}%
\pgfpathlineto{\pgfqpoint{0.962683in}{1.063091in}}%
\pgfpathlineto{\pgfqpoint{0.954716in}{1.052094in}}%
\pgfpathlineto{\pgfqpoint{0.946726in}{1.041444in}}%
\pgfpathclose%
\pgfusepath{fill}%
\end{pgfscope}%
\begin{pgfscope}%
\pgfpathrectangle{\pgfqpoint{0.407592in}{0.222904in}}{\pgfqpoint{1.411163in}{1.411163in}}%
\pgfusepath{clip}%
\pgfsetbuttcap%
\pgfsetroundjoin%
\definecolor{currentfill}{rgb}{0.556863,0.654902,0.847059}%
\pgfsetfillcolor{currentfill}%
\pgfsetfillopacity{0.130000}%
\pgfsetlinewidth{0.000000pt}%
\definecolor{currentstroke}{rgb}{0.000000,0.000000,0.000000}%
\pgfsetstrokecolor{currentstroke}%
\pgfsetdash{}{0pt}%
\pgfpathmoveto{\pgfqpoint{1.025571in}{1.001551in}}%
\pgfpathlineto{\pgfqpoint{1.041329in}{0.987427in}}%
\pgfpathlineto{\pgfqpoint{1.057105in}{0.973936in}}%
\pgfpathlineto{\pgfqpoint{1.064965in}{0.987058in}}%
\pgfpathlineto{\pgfqpoint{1.072808in}{1.000468in}}%
\pgfpathlineto{\pgfqpoint{1.057057in}{1.013340in}}%
\pgfpathlineto{\pgfqpoint{1.041325in}{1.026847in}}%
\pgfpathlineto{\pgfqpoint{1.033457in}{1.014049in}}%
\pgfpathlineto{\pgfqpoint{1.025571in}{1.001551in}}%
\pgfpathclose%
\pgfusepath{fill}%
\end{pgfscope}%
\begin{pgfscope}%
\pgfpathrectangle{\pgfqpoint{0.407592in}{0.222904in}}{\pgfqpoint{1.411163in}{1.411163in}}%
\pgfusepath{clip}%
\pgfsetbuttcap%
\pgfsetroundjoin%
\definecolor{currentfill}{rgb}{0.556863,0.654902,0.847059}%
\pgfsetfillcolor{currentfill}%
\pgfsetfillopacity{0.130000}%
\pgfsetlinewidth{0.000000pt}%
\definecolor{currentstroke}{rgb}{0.000000,0.000000,0.000000}%
\pgfsetstrokecolor{currentstroke}%
\pgfsetdash{}{0pt}%
\pgfpathmoveto{\pgfqpoint{1.261431in}{1.054624in}}%
\pgfpathlineto{\pgfqpoint{1.277406in}{1.050651in}}%
\pgfpathlineto{\pgfqpoint{1.293426in}{1.047526in}}%
\pgfpathlineto{\pgfqpoint{1.301073in}{1.063694in}}%
\pgfpathlineto{\pgfqpoint{1.308706in}{1.079908in}}%
\pgfpathlineto{\pgfqpoint{1.292696in}{1.083029in}}%
\pgfpathlineto{\pgfqpoint{1.276730in}{1.086941in}}%
\pgfpathlineto{\pgfqpoint{1.269087in}{1.070760in}}%
\pgfpathlineto{\pgfqpoint{1.261431in}{1.054624in}}%
\pgfpathclose%
\pgfusepath{fill}%
\end{pgfscope}%
\begin{pgfscope}%
\pgfpathrectangle{\pgfqpoint{0.407592in}{0.222904in}}{\pgfqpoint{1.411163in}{1.411163in}}%
\pgfusepath{clip}%
\pgfsetbuttcap%
\pgfsetroundjoin%
\definecolor{currentfill}{rgb}{0.556863,0.654902,0.847059}%
\pgfsetfillcolor{currentfill}%
\pgfsetfillopacity{0.130000}%
\pgfsetlinewidth{0.000000pt}%
\definecolor{currentstroke}{rgb}{0.000000,0.000000,0.000000}%
\pgfsetstrokecolor{currentstroke}%
\pgfsetdash{}{0pt}%
\pgfpathmoveto{\pgfqpoint{0.725579in}{1.227318in}}%
\pgfpathlineto{\pgfqpoint{0.741275in}{1.208042in}}%
\pgfpathlineto{\pgfqpoint{0.756980in}{1.188829in}}%
\pgfpathlineto{\pgfqpoint{0.765265in}{1.196178in}}%
\pgfpathlineto{\pgfqpoint{0.773522in}{1.203843in}}%
\pgfpathlineto{\pgfqpoint{0.757857in}{1.222798in}}%
\pgfpathlineto{\pgfqpoint{0.742202in}{1.241828in}}%
\pgfpathlineto{\pgfqpoint{0.733905in}{1.234425in}}%
\pgfpathlineto{\pgfqpoint{0.725579in}{1.227318in}}%
\pgfpathclose%
\pgfusepath{fill}%
\end{pgfscope}%
\begin{pgfscope}%
\pgfpathrectangle{\pgfqpoint{0.407592in}{0.222904in}}{\pgfqpoint{1.411163in}{1.411163in}}%
\pgfusepath{clip}%
\pgfsetbuttcap%
\pgfsetroundjoin%
\definecolor{currentfill}{rgb}{0.556863,0.654902,0.847059}%
\pgfsetfillcolor{currentfill}%
\pgfsetfillopacity{0.130000}%
\pgfsetlinewidth{0.000000pt}%
\definecolor{currentstroke}{rgb}{0.000000,0.000000,0.000000}%
\pgfsetstrokecolor{currentstroke}%
\pgfsetdash{}{0pt}%
\pgfpathmoveto{\pgfqpoint{0.867700in}{1.092303in}}%
\pgfpathlineto{\pgfqpoint{0.883429in}{1.074219in}}%
\pgfpathlineto{\pgfqpoint{0.899168in}{1.056333in}}%
\pgfpathlineto{\pgfqpoint{0.907241in}{1.065749in}}%
\pgfpathlineto{\pgfqpoint{0.915290in}{1.075537in}}%
\pgfpathlineto{\pgfqpoint{0.899588in}{1.092980in}}%
\pgfpathlineto{\pgfqpoint{0.883896in}{1.110644in}}%
\pgfpathlineto{\pgfqpoint{0.875811in}{1.101297in}}%
\pgfpathlineto{\pgfqpoint{0.867700in}{1.092303in}}%
\pgfpathclose%
\pgfusepath{fill}%
\end{pgfscope}%
\begin{pgfscope}%
\pgfpathrectangle{\pgfqpoint{0.407592in}{0.222904in}}{\pgfqpoint{1.411163in}{1.411163in}}%
\pgfusepath{clip}%
\pgfsetbuttcap%
\pgfsetroundjoin%
\definecolor{currentfill}{rgb}{0.556863,0.654902,0.847059}%
\pgfsetfillcolor{currentfill}%
\pgfsetfillopacity{0.130000}%
\pgfsetlinewidth{0.000000pt}%
\definecolor{currentstroke}{rgb}{0.000000,0.000000,0.000000}%
\pgfsetstrokecolor{currentstroke}%
\pgfsetdash{}{0pt}%
\pgfpathmoveto{\pgfqpoint{0.756980in}{1.188829in}}%
\pgfpathlineto{\pgfqpoint{0.772693in}{1.169687in}}%
\pgfpathlineto{\pgfqpoint{0.788416in}{1.150626in}}%
\pgfpathlineto{\pgfqpoint{0.796661in}{1.158244in}}%
\pgfpathlineto{\pgfqpoint{0.804877in}{1.166196in}}%
\pgfpathlineto{\pgfqpoint{0.789195in}{1.184972in}}%
\pgfpathlineto{\pgfqpoint{0.773522in}{1.203843in}}%
\pgfpathlineto{\pgfqpoint{0.765265in}{1.196178in}}%
\pgfpathlineto{\pgfqpoint{0.756980in}{1.188829in}}%
\pgfpathclose%
\pgfusepath{fill}%
\end{pgfscope}%
\begin{pgfscope}%
\pgfpathrectangle{\pgfqpoint{0.407592in}{0.222904in}}{\pgfqpoint{1.411163in}{1.411163in}}%
\pgfusepath{clip}%
\pgfsetbuttcap%
\pgfsetroundjoin%
\definecolor{currentfill}{rgb}{0.556863,0.654902,0.847059}%
\pgfsetfillcolor{currentfill}%
\pgfsetfillopacity{0.130000}%
\pgfsetlinewidth{0.000000pt}%
\definecolor{currentstroke}{rgb}{0.000000,0.000000,0.000000}%
\pgfsetstrokecolor{currentstroke}%
\pgfsetdash{}{0pt}%
\pgfpathmoveto{\pgfqpoint{0.978208in}{1.008659in}}%
\pgfpathlineto{\pgfqpoint{0.993967in}{0.992874in}}%
\pgfpathlineto{\pgfqpoint{1.009740in}{0.977571in}}%
\pgfpathlineto{\pgfqpoint{1.017665in}{0.989382in}}%
\pgfpathlineto{\pgfqpoint{1.025571in}{1.001551in}}%
\pgfpathlineto{\pgfqpoint{1.009828in}{1.016239in}}%
\pgfpathlineto{\pgfqpoint{0.994100in}{1.031429in}}%
\pgfpathlineto{\pgfqpoint{0.986164in}{1.019866in}}%
\pgfpathlineto{\pgfqpoint{0.978208in}{1.008659in}}%
\pgfpathclose%
\pgfusepath{fill}%
\end{pgfscope}%
\begin{pgfscope}%
\pgfpathrectangle{\pgfqpoint{0.407592in}{0.222904in}}{\pgfqpoint{1.411163in}{1.411163in}}%
\pgfusepath{clip}%
\pgfsetbuttcap%
\pgfsetroundjoin%
\definecolor{currentfill}{rgb}{0.556863,0.654902,0.847059}%
\pgfsetfillcolor{currentfill}%
\pgfsetfillopacity{0.130000}%
\pgfsetlinewidth{0.000000pt}%
\definecolor{currentstroke}{rgb}{0.000000,0.000000,0.000000}%
\pgfsetstrokecolor{currentstroke}%
\pgfsetdash{}{0pt}%
\pgfpathmoveto{\pgfqpoint{1.151591in}{0.986371in}}%
\pgfpathlineto{\pgfqpoint{1.167447in}{0.978114in}}%
\pgfpathlineto{\pgfqpoint{1.183336in}{0.970828in}}%
\pgfpathlineto{\pgfqpoint{1.191081in}{0.986273in}}%
\pgfpathlineto{\pgfqpoint{1.198813in}{1.001845in}}%
\pgfpathlineto{\pgfqpoint{1.182937in}{1.008828in}}%
\pgfpathlineto{\pgfqpoint{1.167094in}{1.016715in}}%
\pgfpathlineto{\pgfqpoint{1.159350in}{1.001471in}}%
\pgfpathlineto{\pgfqpoint{1.151591in}{0.986371in}}%
\pgfpathclose%
\pgfusepath{fill}%
\end{pgfscope}%
\begin{pgfscope}%
\pgfpathrectangle{\pgfqpoint{0.407592in}{0.222904in}}{\pgfqpoint{1.411163in}{1.411163in}}%
\pgfusepath{clip}%
\pgfsetbuttcap%
\pgfsetroundjoin%
\definecolor{currentfill}{rgb}{0.556863,0.654902,0.847059}%
\pgfsetfillcolor{currentfill}%
\pgfsetfillopacity{0.130000}%
\pgfsetlinewidth{0.000000pt}%
\definecolor{currentstroke}{rgb}{0.000000,0.000000,0.000000}%
\pgfsetstrokecolor{currentstroke}%
\pgfsetdash{}{0pt}%
\pgfpathmoveto{\pgfqpoint{1.104371in}{0.976896in}}%
\pgfpathlineto{\pgfqpoint{1.120187in}{0.966336in}}%
\pgfpathlineto{\pgfqpoint{1.136029in}{0.956689in}}%
\pgfpathlineto{\pgfqpoint{1.143817in}{0.971436in}}%
\pgfpathlineto{\pgfqpoint{1.151591in}{0.986371in}}%
\pgfpathlineto{\pgfqpoint{1.135765in}{0.995548in}}%
\pgfpathlineto{\pgfqpoint{1.119966in}{1.005592in}}%
\pgfpathlineto{\pgfqpoint{1.112176in}{0.991138in}}%
\pgfpathlineto{\pgfqpoint{1.104371in}{0.976896in}}%
\pgfpathclose%
\pgfusepath{fill}%
\end{pgfscope}%
\begin{pgfscope}%
\pgfpathrectangle{\pgfqpoint{0.407592in}{0.222904in}}{\pgfqpoint{1.411163in}{1.411163in}}%
\pgfusepath{clip}%
\pgfsetbuttcap%
\pgfsetroundjoin%
\definecolor{currentfill}{rgb}{0.556863,0.654902,0.847059}%
\pgfsetfillcolor{currentfill}%
\pgfsetfillopacity{0.130000}%
\pgfsetlinewidth{0.000000pt}%
\definecolor{currentstroke}{rgb}{0.000000,0.000000,0.000000}%
\pgfsetstrokecolor{currentstroke}%
\pgfsetdash{}{0pt}%
\pgfpathmoveto{\pgfqpoint{1.308706in}{1.079908in}}%
\pgfpathlineto{\pgfqpoint{1.324765in}{1.077576in}}%
\pgfpathlineto{\pgfqpoint{1.340873in}{1.076026in}}%
\pgfpathlineto{\pgfqpoint{1.348485in}{1.092211in}}%
\pgfpathlineto{\pgfqpoint{1.332382in}{1.093808in}}%
\pgfpathlineto{\pgfqpoint{1.316328in}{1.096162in}}%
\pgfusepath{fill}%
\end{pgfscope}%
\begin{pgfscope}%
\pgfpathrectangle{\pgfqpoint{0.407592in}{0.222904in}}{\pgfqpoint{1.411163in}{1.411163in}}%
\pgfusepath{clip}%
\pgfsetbuttcap%
\pgfsetroundjoin%
\definecolor{currentfill}{rgb}{0.556863,0.654902,0.847059}%
\pgfsetfillcolor{currentfill}%
\pgfsetfillopacity{0.130000}%
\pgfsetlinewidth{0.000000pt}%
\definecolor{currentstroke}{rgb}{0.000000,0.000000,0.000000}%
\pgfsetstrokecolor{currentstroke}%
\pgfsetdash{}{0pt}%
\pgfpathmoveto{\pgfqpoint{1.198813in}{1.001845in}}%
\pgfpathlineto{\pgfqpoint{1.214726in}{0.995801in}}%
\pgfpathlineto{\pgfqpoint{1.230679in}{0.990721in}}%
\pgfpathlineto{\pgfqpoint{1.238386in}{1.006579in}}%
\pgfpathlineto{\pgfqpoint{1.246081in}{1.022524in}}%
\pgfpathlineto{\pgfqpoint{1.230139in}{1.027466in}}%
\pgfpathlineto{\pgfqpoint{1.214237in}{1.033302in}}%
\pgfpathlineto{\pgfqpoint{1.206532in}{1.017526in}}%
\pgfpathlineto{\pgfqpoint{1.198813in}{1.001845in}}%
\pgfpathclose%
\pgfusepath{fill}%
\end{pgfscope}%
\begin{pgfscope}%
\pgfpathrectangle{\pgfqpoint{0.407592in}{0.222904in}}{\pgfqpoint{1.411163in}{1.411163in}}%
\pgfusepath{clip}%
\pgfsetbuttcap%
\pgfsetroundjoin%
\definecolor{currentfill}{rgb}{0.556863,0.654902,0.847059}%
\pgfsetfillcolor{currentfill}%
\pgfsetfillopacity{0.130000}%
\pgfsetlinewidth{0.000000pt}%
\definecolor{currentstroke}{rgb}{0.000000,0.000000,0.000000}%
\pgfsetstrokecolor{currentstroke}%
\pgfsetdash{}{0pt}%
\pgfpathmoveto{\pgfqpoint{0.899168in}{1.056333in}}%
\pgfpathlineto{\pgfqpoint{0.914917in}{1.038674in}}%
\pgfpathlineto{\pgfqpoint{0.930677in}{1.021276in}}%
\pgfpathlineto{\pgfqpoint{0.938713in}{1.031164in}}%
\pgfpathlineto{\pgfqpoint{0.946726in}{1.041444in}}%
\pgfpathlineto{\pgfqpoint{0.931002in}{1.058348in}}%
\pgfpathlineto{\pgfqpoint{0.915290in}{1.075537in}}%
\pgfpathlineto{\pgfqpoint{0.907241in}{1.065749in}}%
\pgfpathlineto{\pgfqpoint{0.899168in}{1.056333in}}%
\pgfpathclose%
\pgfusepath{fill}%
\end{pgfscope}%
\begin{pgfscope}%
\pgfpathrectangle{\pgfqpoint{0.407592in}{0.222904in}}{\pgfqpoint{1.411163in}{1.411163in}}%
\pgfusepath{clip}%
\pgfsetbuttcap%
\pgfsetroundjoin%
\definecolor{currentfill}{rgb}{0.556863,0.654902,0.847059}%
\pgfsetfillcolor{currentfill}%
\pgfsetfillopacity{0.130000}%
\pgfsetlinewidth{0.000000pt}%
\definecolor{currentstroke}{rgb}{0.000000,0.000000,0.000000}%
\pgfsetstrokecolor{currentstroke}%
\pgfsetdash{}{0pt}%
\pgfpathmoveto{\pgfqpoint{1.057105in}{0.973936in}}%
\pgfpathlineto{\pgfqpoint{1.072898in}{0.961153in}}%
\pgfpathlineto{\pgfqpoint{1.088712in}{0.949159in}}%
\pgfpathlineto{\pgfqpoint{1.096550in}{0.962892in}}%
\pgfpathlineto{\pgfqpoint{1.104371in}{0.976896in}}%
\pgfpathlineto{\pgfqpoint{1.088579in}{0.988297in}}%
\pgfpathlineto{\pgfqpoint{1.072808in}{1.000468in}}%
\pgfpathlineto{\pgfqpoint{1.064965in}{0.987058in}}%
\pgfpathlineto{\pgfqpoint{1.057105in}{0.973936in}}%
\pgfpathclose%
\pgfusepath{fill}%
\end{pgfscope}%
\begin{pgfscope}%
\pgfpathrectangle{\pgfqpoint{0.407592in}{0.222904in}}{\pgfqpoint{1.411163in}{1.411163in}}%
\pgfusepath{clip}%
\pgfsetbuttcap%
\pgfsetroundjoin%
\definecolor{currentfill}{rgb}{0.556863,0.654902,0.847059}%
\pgfsetfillcolor{currentfill}%
\pgfsetfillopacity{0.130000}%
\pgfsetlinewidth{0.000000pt}%
\definecolor{currentstroke}{rgb}{0.000000,0.000000,0.000000}%
\pgfsetstrokecolor{currentstroke}%
\pgfsetdash{}{0pt}%
\pgfpathmoveto{\pgfqpoint{0.788416in}{1.150626in}}%
\pgfpathlineto{\pgfqpoint{0.804148in}{1.131657in}}%
\pgfpathlineto{\pgfqpoint{0.819889in}{1.112792in}}%
\pgfpathlineto{\pgfqpoint{0.828093in}{1.120706in}}%
\pgfpathlineto{\pgfqpoint{0.836270in}{1.128977in}}%
\pgfpathlineto{\pgfqpoint{0.820569in}{1.147526in}}%
\pgfpathlineto{\pgfqpoint{0.804877in}{1.166196in}}%
\pgfpathlineto{\pgfqpoint{0.796661in}{1.158244in}}%
\pgfpathlineto{\pgfqpoint{0.788416in}{1.150626in}}%
\pgfpathclose%
\pgfusepath{fill}%
\end{pgfscope}%
\begin{pgfscope}%
\pgfpathrectangle{\pgfqpoint{0.407592in}{0.222904in}}{\pgfqpoint{1.411163in}{1.411163in}}%
\pgfusepath{clip}%
\pgfsetbuttcap%
\pgfsetroundjoin%
\definecolor{currentfill}{rgb}{0.556863,0.654902,0.847059}%
\pgfsetfillcolor{currentfill}%
\pgfsetfillopacity{0.130000}%
\pgfsetlinewidth{0.000000pt}%
\definecolor{currentstroke}{rgb}{0.000000,0.000000,0.000000}%
\pgfsetstrokecolor{currentstroke}%
\pgfsetdash{}{0pt}%
\pgfpathmoveto{\pgfqpoint{1.246081in}{1.022524in}}%
\pgfpathlineto{\pgfqpoint{1.262065in}{1.018489in}}%
\pgfpathlineto{\pgfqpoint{1.278095in}{1.015365in}}%
\pgfpathlineto{\pgfqpoint{1.285767in}{1.031413in}}%
\pgfpathlineto{\pgfqpoint{1.293426in}{1.047526in}}%
\pgfpathlineto{\pgfqpoint{1.277406in}{1.050651in}}%
\pgfpathlineto{\pgfqpoint{1.261431in}{1.054624in}}%
\pgfpathlineto{\pgfqpoint{1.253762in}{1.038542in}}%
\pgfpathlineto{\pgfqpoint{1.246081in}{1.022524in}}%
\pgfpathclose%
\pgfusepath{fill}%
\end{pgfscope}%
\begin{pgfscope}%
\pgfpathrectangle{\pgfqpoint{0.407592in}{0.222904in}}{\pgfqpoint{1.411163in}{1.411163in}}%
\pgfusepath{clip}%
\pgfsetbuttcap%
\pgfsetroundjoin%
\definecolor{currentfill}{rgb}{0.556863,0.654902,0.847059}%
\pgfsetfillcolor{currentfill}%
\pgfsetfillopacity{0.130000}%
\pgfsetlinewidth{0.000000pt}%
\definecolor{currentstroke}{rgb}{0.000000,0.000000,0.000000}%
\pgfsetstrokecolor{currentstroke}%
\pgfsetdash{}{0pt}%
\pgfpathmoveto{\pgfqpoint{1.009740in}{0.977571in}}%
\pgfpathlineto{\pgfqpoint{1.025527in}{0.962816in}}%
\pgfpathlineto{\pgfqpoint{1.041329in}{0.948685in}}%
\pgfpathlineto{\pgfqpoint{1.049226in}{0.961133in}}%
\pgfpathlineto{\pgfqpoint{1.057105in}{0.973936in}}%
\pgfpathlineto{\pgfqpoint{1.041329in}{0.987427in}}%
\pgfpathlineto{\pgfqpoint{1.025571in}{1.001551in}}%
\pgfpathlineto{\pgfqpoint{1.017665in}{0.989382in}}%
\pgfpathlineto{\pgfqpoint{1.009740in}{0.977571in}}%
\pgfpathclose%
\pgfusepath{fill}%
\end{pgfscope}%
\begin{pgfscope}%
\pgfpathrectangle{\pgfqpoint{0.407592in}{0.222904in}}{\pgfqpoint{1.411163in}{1.411163in}}%
\pgfusepath{clip}%
\pgfsetbuttcap%
\pgfsetroundjoin%
\definecolor{currentfill}{rgb}{0.556863,0.654902,0.847059}%
\pgfsetfillcolor{currentfill}%
\pgfsetfillopacity{0.130000}%
\pgfsetlinewidth{0.000000pt}%
\definecolor{currentstroke}{rgb}{0.000000,0.000000,0.000000}%
\pgfsetstrokecolor{currentstroke}%
\pgfsetdash{}{0pt}%
\pgfpathmoveto{\pgfqpoint{0.819889in}{1.112792in}}%
\pgfpathlineto{\pgfqpoint{0.835638in}{1.094044in}}%
\pgfpathlineto{\pgfqpoint{0.851397in}{1.075431in}}%
\pgfpathlineto{\pgfqpoint{0.859562in}{1.083677in}}%
\pgfpathlineto{\pgfqpoint{0.867700in}{1.092303in}}%
\pgfpathlineto{\pgfqpoint{0.851980in}{1.110563in}}%
\pgfpathlineto{\pgfqpoint{0.836270in}{1.128977in}}%
\pgfpathlineto{\pgfqpoint{0.828093in}{1.120706in}}%
\pgfpathlineto{\pgfqpoint{0.819889in}{1.112792in}}%
\pgfpathclose%
\pgfusepath{fill}%
\end{pgfscope}%
\begin{pgfscope}%
\pgfpathrectangle{\pgfqpoint{0.407592in}{0.222904in}}{\pgfqpoint{1.411163in}{1.411163in}}%
\pgfusepath{clip}%
\pgfsetbuttcap%
\pgfsetroundjoin%
\definecolor{currentfill}{rgb}{0.556863,0.654902,0.847059}%
\pgfsetfillcolor{currentfill}%
\pgfsetfillopacity{0.130000}%
\pgfsetlinewidth{0.000000pt}%
\definecolor{currentstroke}{rgb}{0.000000,0.000000,0.000000}%
\pgfsetstrokecolor{currentstroke}%
\pgfsetdash{}{0pt}%
\pgfpathmoveto{\pgfqpoint{0.930677in}{1.021276in}}%
\pgfpathlineto{\pgfqpoint{0.946447in}{1.004176in}}%
\pgfpathlineto{\pgfqpoint{0.962228in}{0.987419in}}%
\pgfpathlineto{\pgfqpoint{0.970229in}{0.997834in}}%
\pgfpathlineto{\pgfqpoint{0.978208in}{1.008659in}}%
\pgfpathlineto{\pgfqpoint{0.962461in}{1.024867in}}%
\pgfpathlineto{\pgfqpoint{0.946726in}{1.041444in}}%
\pgfpathlineto{\pgfqpoint{0.938713in}{1.031164in}}%
\pgfpathlineto{\pgfqpoint{0.930677in}{1.021276in}}%
\pgfpathclose%
\pgfusepath{fill}%
\end{pgfscope}%
\begin{pgfscope}%
\pgfpathrectangle{\pgfqpoint{0.407592in}{0.222904in}}{\pgfqpoint{1.411163in}{1.411163in}}%
\pgfusepath{clip}%
\pgfsetbuttcap%
\pgfsetroundjoin%
\definecolor{currentfill}{rgb}{0.556863,0.654902,0.847059}%
\pgfsetfillcolor{currentfill}%
\pgfsetfillopacity{0.130000}%
\pgfsetlinewidth{0.000000pt}%
\definecolor{currentstroke}{rgb}{0.000000,0.000000,0.000000}%
\pgfsetstrokecolor{currentstroke}%
\pgfsetdash{}{0pt}%
\pgfpathmoveto{\pgfqpoint{1.293426in}{1.047526in}}%
\pgfpathlineto{\pgfqpoint{1.309494in}{1.045246in}}%
\pgfpathlineto{\pgfqpoint{1.325612in}{1.043802in}}%
\pgfpathlineto{\pgfqpoint{1.333249in}{1.059887in}}%
\pgfpathlineto{\pgfqpoint{1.340873in}{1.076026in}}%
\pgfpathlineto{\pgfqpoint{1.324765in}{1.077576in}}%
\pgfpathlineto{\pgfqpoint{1.308706in}{1.079908in}}%
\pgfpathlineto{\pgfqpoint{1.301073in}{1.063694in}}%
\pgfpathlineto{\pgfqpoint{1.293426in}{1.047526in}}%
\pgfpathclose%
\pgfusepath{fill}%
\end{pgfscope}%
\begin{pgfscope}%
\pgfpathrectangle{\pgfqpoint{0.407592in}{0.222904in}}{\pgfqpoint{1.411163in}{1.411163in}}%
\pgfusepath{clip}%
\pgfsetbuttcap%
\pgfsetroundjoin%
\definecolor{currentfill}{rgb}{0.556863,0.654902,0.847059}%
\pgfsetfillcolor{currentfill}%
\pgfsetfillopacity{0.130000}%
\pgfsetlinewidth{0.000000pt}%
\definecolor{currentstroke}{rgb}{0.000000,0.000000,0.000000}%
\pgfsetstrokecolor{currentstroke}%
\pgfsetdash{}{0pt}%
\pgfpathmoveto{\pgfqpoint{0.708836in}{1.214017in}}%
\pgfpathlineto{\pgfqpoint{0.724573in}{1.194532in}}%
\pgfpathlineto{\pgfqpoint{0.740319in}{1.175101in}}%
\pgfpathlineto{\pgfqpoint{0.748665in}{1.181801in}}%
\pgfpathlineto{\pgfqpoint{0.756980in}{1.188829in}}%
\pgfpathlineto{\pgfqpoint{0.741275in}{1.208042in}}%
\pgfpathlineto{\pgfqpoint{0.725579in}{1.227318in}}%
\pgfpathlineto{\pgfqpoint{0.717222in}{1.220514in}}%
\pgfpathlineto{\pgfqpoint{0.708836in}{1.214017in}}%
\pgfpathclose%
\pgfusepath{fill}%
\end{pgfscope}%
\begin{pgfscope}%
\pgfpathrectangle{\pgfqpoint{0.407592in}{0.222904in}}{\pgfqpoint{1.411163in}{1.411163in}}%
\pgfusepath{clip}%
\pgfsetbuttcap%
\pgfsetroundjoin%
\definecolor{currentfill}{rgb}{0.556863,0.654902,0.847059}%
\pgfsetfillcolor{currentfill}%
\pgfsetfillopacity{0.130000}%
\pgfsetlinewidth{0.000000pt}%
\definecolor{currentstroke}{rgb}{0.000000,0.000000,0.000000}%
\pgfsetstrokecolor{currentstroke}%
\pgfsetdash{}{0pt}%
\pgfpathmoveto{\pgfqpoint{1.136029in}{0.956689in}}%
\pgfpathlineto{\pgfqpoint{1.151900in}{0.948023in}}%
\pgfpathlineto{\pgfqpoint{1.167803in}{0.940398in}}%
\pgfpathlineto{\pgfqpoint{1.175576in}{0.955528in}}%
\pgfpathlineto{\pgfqpoint{1.183336in}{0.970828in}}%
\pgfpathlineto{\pgfqpoint{1.167447in}{0.978114in}}%
\pgfpathlineto{\pgfqpoint{1.151591in}{0.986371in}}%
\pgfpathlineto{\pgfqpoint{1.143817in}{0.971436in}}%
\pgfpathlineto{\pgfqpoint{1.136029in}{0.956689in}}%
\pgfpathclose%
\pgfusepath{fill}%
\end{pgfscope}%
\begin{pgfscope}%
\pgfpathrectangle{\pgfqpoint{0.407592in}{0.222904in}}{\pgfqpoint{1.411163in}{1.411163in}}%
\pgfusepath{clip}%
\pgfsetbuttcap%
\pgfsetroundjoin%
\definecolor{currentfill}{rgb}{0.556863,0.654902,0.847059}%
\pgfsetfillcolor{currentfill}%
\pgfsetfillopacity{0.130000}%
\pgfsetlinewidth{0.000000pt}%
\definecolor{currentstroke}{rgb}{0.000000,0.000000,0.000000}%
\pgfsetstrokecolor{currentstroke}%
\pgfsetdash{}{0pt}%
\pgfpathmoveto{\pgfqpoint{1.088712in}{0.949159in}}%
\pgfpathlineto{\pgfqpoint{1.104548in}{0.938038in}}%
\pgfpathlineto{\pgfqpoint{1.120408in}{0.927876in}}%
\pgfpathlineto{\pgfqpoint{1.128227in}{0.942159in}}%
\pgfpathlineto{\pgfqpoint{1.136029in}{0.956689in}}%
\pgfpathlineto{\pgfqpoint{1.120187in}{0.966336in}}%
\pgfpathlineto{\pgfqpoint{1.104371in}{0.976896in}}%
\pgfpathlineto{\pgfqpoint{1.096550in}{0.962892in}}%
\pgfpathlineto{\pgfqpoint{1.088712in}{0.949159in}}%
\pgfpathclose%
\pgfusepath{fill}%
\end{pgfscope}%
\begin{pgfscope}%
\pgfpathrectangle{\pgfqpoint{0.407592in}{0.222904in}}{\pgfqpoint{1.411163in}{1.411163in}}%
\pgfusepath{clip}%
\pgfsetbuttcap%
\pgfsetroundjoin%
\definecolor{currentfill}{rgb}{0.556863,0.654902,0.847059}%
\pgfsetfillcolor{currentfill}%
\pgfsetfillopacity{0.130000}%
\pgfsetlinewidth{0.000000pt}%
\definecolor{currentstroke}{rgb}{0.000000,0.000000,0.000000}%
\pgfsetstrokecolor{currentstroke}%
\pgfsetdash{}{0pt}%
\pgfpathmoveto{\pgfqpoint{0.851397in}{1.075431in}}%
\pgfpathlineto{\pgfqpoint{0.867166in}{1.056972in}}%
\pgfpathlineto{\pgfqpoint{0.882944in}{1.038690in}}%
\pgfpathlineto{\pgfqpoint{0.891069in}{1.047307in}}%
\pgfpathlineto{\pgfqpoint{0.899168in}{1.056333in}}%
\pgfpathlineto{\pgfqpoint{0.883429in}{1.074219in}}%
\pgfpathlineto{\pgfqpoint{0.867700in}{1.092303in}}%
\pgfpathlineto{\pgfqpoint{0.859562in}{1.083677in}}%
\pgfpathlineto{\pgfqpoint{0.851397in}{1.075431in}}%
\pgfpathclose%
\pgfusepath{fill}%
\end{pgfscope}%
\begin{pgfscope}%
\pgfpathrectangle{\pgfqpoint{0.407592in}{0.222904in}}{\pgfqpoint{1.411163in}{1.411163in}}%
\pgfusepath{clip}%
\pgfsetbuttcap%
\pgfsetroundjoin%
\definecolor{currentfill}{rgb}{0.556863,0.654902,0.847059}%
\pgfsetfillcolor{currentfill}%
\pgfsetfillopacity{0.130000}%
\pgfsetlinewidth{0.000000pt}%
\definecolor{currentstroke}{rgb}{0.000000,0.000000,0.000000}%
\pgfsetstrokecolor{currentstroke}%
\pgfsetdash{}{0pt}%
\pgfpathmoveto{\pgfqpoint{1.183336in}{0.970828in}}%
\pgfpathlineto{\pgfqpoint{1.199261in}{0.964553in}}%
\pgfpathlineto{\pgfqpoint{1.215224in}{0.959322in}}%
\pgfpathlineto{\pgfqpoint{1.222958in}{0.974963in}}%
\pgfpathlineto{\pgfqpoint{1.230679in}{0.990721in}}%
\pgfpathlineto{\pgfqpoint{1.214726in}{0.995801in}}%
\pgfpathlineto{\pgfqpoint{1.198813in}{1.001845in}}%
\pgfpathlineto{\pgfqpoint{1.191081in}{0.986273in}}%
\pgfpathlineto{\pgfqpoint{1.183336in}{0.970828in}}%
\pgfpathclose%
\pgfusepath{fill}%
\end{pgfscope}%
\begin{pgfscope}%
\pgfpathrectangle{\pgfqpoint{0.407592in}{0.222904in}}{\pgfqpoint{1.411163in}{1.411163in}}%
\pgfusepath{clip}%
\pgfsetbuttcap%
\pgfsetroundjoin%
\definecolor{currentfill}{rgb}{0.556863,0.654902,0.847059}%
\pgfsetfillcolor{currentfill}%
\pgfsetfillopacity{0.130000}%
\pgfsetlinewidth{0.000000pt}%
\definecolor{currentstroke}{rgb}{0.000000,0.000000,0.000000}%
\pgfsetstrokecolor{currentstroke}%
\pgfsetdash{}{0pt}%
\pgfpathmoveto{\pgfqpoint{0.962228in}{0.987419in}}%
\pgfpathlineto{\pgfqpoint{0.978021in}{0.971057in}}%
\pgfpathlineto{\pgfqpoint{0.993826in}{0.955152in}}%
\pgfpathlineto{\pgfqpoint{1.001794in}{0.966150in}}%
\pgfpathlineto{\pgfqpoint{1.009740in}{0.977571in}}%
\pgfpathlineto{\pgfqpoint{0.993967in}{0.992874in}}%
\pgfpathlineto{\pgfqpoint{0.978208in}{1.008659in}}%
\pgfpathlineto{\pgfqpoint{0.970229in}{0.997834in}}%
\pgfpathlineto{\pgfqpoint{0.962228in}{0.987419in}}%
\pgfpathclose%
\pgfusepath{fill}%
\end{pgfscope}%
\begin{pgfscope}%
\pgfpathrectangle{\pgfqpoint{0.407592in}{0.222904in}}{\pgfqpoint{1.411163in}{1.411163in}}%
\pgfusepath{clip}%
\pgfsetbuttcap%
\pgfsetroundjoin%
\definecolor{currentfill}{rgb}{0.556863,0.654902,0.847059}%
\pgfsetfillcolor{currentfill}%
\pgfsetfillopacity{0.130000}%
\pgfsetlinewidth{0.000000pt}%
\definecolor{currentstroke}{rgb}{0.000000,0.000000,0.000000}%
\pgfsetstrokecolor{currentstroke}%
\pgfsetdash{}{0pt}%
\pgfpathmoveto{\pgfqpoint{1.041329in}{0.948685in}}%
\pgfpathlineto{\pgfqpoint{1.057148in}{0.935262in}}%
\pgfpathlineto{\pgfqpoint{1.072986in}{0.922640in}}%
\pgfpathlineto{\pgfqpoint{1.080858in}{0.935729in}}%
\pgfpathlineto{\pgfqpoint{1.088712in}{0.949159in}}%
\pgfpathlineto{\pgfqpoint{1.072898in}{0.961153in}}%
\pgfpathlineto{\pgfqpoint{1.057105in}{0.973936in}}%
\pgfpathlineto{\pgfqpoint{1.049226in}{0.961133in}}%
\pgfpathlineto{\pgfqpoint{1.041329in}{0.948685in}}%
\pgfpathclose%
\pgfusepath{fill}%
\end{pgfscope}%
\begin{pgfscope}%
\pgfpathrectangle{\pgfqpoint{0.407592in}{0.222904in}}{\pgfqpoint{1.411163in}{1.411163in}}%
\pgfusepath{clip}%
\pgfsetbuttcap%
\pgfsetroundjoin%
\definecolor{currentfill}{rgb}{0.556863,0.654902,0.847059}%
\pgfsetfillcolor{currentfill}%
\pgfsetfillopacity{0.130000}%
\pgfsetlinewidth{0.000000pt}%
\definecolor{currentstroke}{rgb}{0.000000,0.000000,0.000000}%
\pgfsetstrokecolor{currentstroke}%
\pgfsetdash{}{0pt}%
\pgfpathmoveto{\pgfqpoint{0.740319in}{1.175101in}}%
\pgfpathlineto{\pgfqpoint{0.756074in}{1.155731in}}%
\pgfpathlineto{\pgfqpoint{0.771838in}{1.136431in}}%
\pgfpathlineto{\pgfqpoint{0.780142in}{1.143353in}}%
\pgfpathlineto{\pgfqpoint{0.788416in}{1.150626in}}%
\pgfpathlineto{\pgfqpoint{0.772693in}{1.169687in}}%
\pgfpathlineto{\pgfqpoint{0.756980in}{1.188829in}}%
\pgfpathlineto{\pgfqpoint{0.748665in}{1.181801in}}%
\pgfpathlineto{\pgfqpoint{0.740319in}{1.175101in}}%
\pgfpathclose%
\pgfusepath{fill}%
\end{pgfscope}%
\begin{pgfscope}%
\pgfpathrectangle{\pgfqpoint{0.407592in}{0.222904in}}{\pgfqpoint{1.411163in}{1.411163in}}%
\pgfusepath{clip}%
\pgfsetbuttcap%
\pgfsetroundjoin%
\definecolor{currentfill}{rgb}{0.556863,0.654902,0.847059}%
\pgfsetfillcolor{currentfill}%
\pgfsetfillopacity{0.130000}%
\pgfsetlinewidth{0.000000pt}%
\definecolor{currentstroke}{rgb}{0.000000,0.000000,0.000000}%
\pgfsetstrokecolor{currentstroke}%
\pgfsetdash{}{0pt}%
\pgfpathmoveto{\pgfqpoint{1.230679in}{0.990721in}}%
\pgfpathlineto{\pgfqpoint{1.246673in}{0.986620in}}%
\pgfpathlineto{\pgfqpoint{1.262713in}{0.983505in}}%
\pgfpathlineto{\pgfqpoint{1.270411in}{0.999391in}}%
\pgfpathlineto{\pgfqpoint{1.278095in}{1.015365in}}%
\pgfpathlineto{\pgfqpoint{1.262065in}{1.018489in}}%
\pgfpathlineto{\pgfqpoint{1.246081in}{1.022524in}}%
\pgfpathlineto{\pgfqpoint{1.238386in}{1.006579in}}%
\pgfpathlineto{\pgfqpoint{1.230679in}{0.990721in}}%
\pgfpathclose%
\pgfusepath{fill}%
\end{pgfscope}%
\begin{pgfscope}%
\pgfpathrectangle{\pgfqpoint{0.407592in}{0.222904in}}{\pgfqpoint{1.411163in}{1.411163in}}%
\pgfusepath{clip}%
\pgfsetbuttcap%
\pgfsetroundjoin%
\definecolor{currentfill}{rgb}{0.556863,0.654902,0.847059}%
\pgfsetfillcolor{currentfill}%
\pgfsetfillopacity{0.130000}%
\pgfsetlinewidth{0.000000pt}%
\definecolor{currentstroke}{rgb}{0.000000,0.000000,0.000000}%
\pgfsetstrokecolor{currentstroke}%
\pgfsetdash{}{0pt}%
\pgfpathmoveto{\pgfqpoint{0.882944in}{1.038690in}}%
\pgfpathlineto{\pgfqpoint{0.898731in}{1.020610in}}%
\pgfpathlineto{\pgfqpoint{0.914528in}{1.002763in}}%
\pgfpathlineto{\pgfqpoint{0.922615in}{1.011802in}}%
\pgfpathlineto{\pgfqpoint{0.930677in}{1.021276in}}%
\pgfpathlineto{\pgfqpoint{0.914917in}{1.038674in}}%
\pgfpathlineto{\pgfqpoint{0.899168in}{1.056333in}}%
\pgfpathlineto{\pgfqpoint{0.891069in}{1.047307in}}%
\pgfpathlineto{\pgfqpoint{0.882944in}{1.038690in}}%
\pgfpathclose%
\pgfusepath{fill}%
\end{pgfscope}%
\begin{pgfscope}%
\pgfpathrectangle{\pgfqpoint{0.407592in}{0.222904in}}{\pgfqpoint{1.411163in}{1.411163in}}%
\pgfusepath{clip}%
\pgfsetbuttcap%
\pgfsetroundjoin%
\definecolor{currentfill}{rgb}{0.556863,0.654902,0.847059}%
\pgfsetfillcolor{currentfill}%
\pgfsetfillopacity{0.130000}%
\pgfsetlinewidth{0.000000pt}%
\definecolor{currentstroke}{rgb}{0.000000,0.000000,0.000000}%
\pgfsetstrokecolor{currentstroke}%
\pgfsetdash{}{0pt}%
\pgfpathmoveto{\pgfqpoint{1.340873in}{1.076026in}}%
\pgfpathlineto{\pgfqpoint{1.357034in}{1.075246in}}%
\pgfpathlineto{\pgfqpoint{1.373249in}{1.075222in}}%
\pgfpathlineto{\pgfqpoint{1.380851in}{1.091245in}}%
\pgfpathlineto{\pgfqpoint{1.364641in}{1.091362in}}%
\pgfpathlineto{\pgfqpoint{1.348485in}{1.092211in}}%
\pgfusepath{fill}%
\end{pgfscope}%
\begin{pgfscope}%
\pgfpathrectangle{\pgfqpoint{0.407592in}{0.222904in}}{\pgfqpoint{1.411163in}{1.411163in}}%
\pgfusepath{clip}%
\pgfsetbuttcap%
\pgfsetroundjoin%
\definecolor{currentfill}{rgb}{0.556863,0.654902,0.847059}%
\pgfsetfillcolor{currentfill}%
\pgfsetfillopacity{0.130000}%
\pgfsetlinewidth{0.000000pt}%
\definecolor{currentstroke}{rgb}{0.000000,0.000000,0.000000}%
\pgfsetstrokecolor{currentstroke}%
\pgfsetdash{}{0pt}%
\pgfpathmoveto{\pgfqpoint{0.771838in}{1.136431in}}%
\pgfpathlineto{\pgfqpoint{0.787611in}{1.117209in}}%
\pgfpathlineto{\pgfqpoint{0.803393in}{1.098077in}}%
\pgfpathlineto{\pgfqpoint{0.811655in}{1.105246in}}%
\pgfpathlineto{\pgfqpoint{0.819889in}{1.112792in}}%
\pgfpathlineto{\pgfqpoint{0.804148in}{1.131657in}}%
\pgfpathlineto{\pgfqpoint{0.788416in}{1.150626in}}%
\pgfpathlineto{\pgfqpoint{0.780142in}{1.143353in}}%
\pgfpathlineto{\pgfqpoint{0.771838in}{1.136431in}}%
\pgfpathclose%
\pgfusepath{fill}%
\end{pgfscope}%
\begin{pgfscope}%
\pgfpathrectangle{\pgfqpoint{0.407592in}{0.222904in}}{\pgfqpoint{1.411163in}{1.411163in}}%
\pgfusepath{clip}%
\pgfsetbuttcap%
\pgfsetroundjoin%
\definecolor{currentfill}{rgb}{0.556863,0.654902,0.847059}%
\pgfsetfillcolor{currentfill}%
\pgfsetfillopacity{0.130000}%
\pgfsetlinewidth{0.000000pt}%
\definecolor{currentstroke}{rgb}{0.000000,0.000000,0.000000}%
\pgfsetstrokecolor{currentstroke}%
\pgfsetdash{}{0pt}%
\pgfpathmoveto{\pgfqpoint{1.278095in}{1.015365in}}%
\pgfpathlineto{\pgfqpoint{1.294173in}{1.013150in}}%
\pgfpathlineto{\pgfqpoint{1.310301in}{1.011832in}}%
\pgfpathlineto{\pgfqpoint{1.317963in}{1.027780in}}%
\pgfpathlineto{\pgfqpoint{1.325612in}{1.043802in}}%
\pgfpathlineto{\pgfqpoint{1.309494in}{1.045246in}}%
\pgfpathlineto{\pgfqpoint{1.293426in}{1.047526in}}%
\pgfpathlineto{\pgfqpoint{1.285767in}{1.031413in}}%
\pgfpathlineto{\pgfqpoint{1.278095in}{1.015365in}}%
\pgfpathclose%
\pgfusepath{fill}%
\end{pgfscope}%
\begin{pgfscope}%
\pgfpathrectangle{\pgfqpoint{0.407592in}{0.222904in}}{\pgfqpoint{1.411163in}{1.411163in}}%
\pgfusepath{clip}%
\pgfsetbuttcap%
\pgfsetroundjoin%
\definecolor{currentfill}{rgb}{0.556863,0.654902,0.847059}%
\pgfsetfillcolor{currentfill}%
\pgfsetfillopacity{0.130000}%
\pgfsetlinewidth{0.000000pt}%
\definecolor{currentstroke}{rgb}{0.000000,0.000000,0.000000}%
\pgfsetstrokecolor{currentstroke}%
\pgfsetdash{}{0pt}%
\pgfpathmoveto{\pgfqpoint{0.993826in}{0.955152in}}%
\pgfpathlineto{\pgfqpoint{1.009644in}{0.939772in}}%
\pgfpathlineto{\pgfqpoint{1.025476in}{0.925000in}}%
\pgfpathlineto{\pgfqpoint{1.033413in}{0.936627in}}%
\pgfpathlineto{\pgfqpoint{1.041329in}{0.948685in}}%
\pgfpathlineto{\pgfqpoint{1.025527in}{0.962816in}}%
\pgfpathlineto{\pgfqpoint{1.009740in}{0.977571in}}%
\pgfpathlineto{\pgfqpoint{1.001794in}{0.966150in}}%
\pgfpathlineto{\pgfqpoint{0.993826in}{0.955152in}}%
\pgfpathclose%
\pgfusepath{fill}%
\end{pgfscope}%
\begin{pgfscope}%
\pgfpathrectangle{\pgfqpoint{0.407592in}{0.222904in}}{\pgfqpoint{1.411163in}{1.411163in}}%
\pgfusepath{clip}%
\pgfsetbuttcap%
\pgfsetroundjoin%
\definecolor{currentfill}{rgb}{0.556863,0.654902,0.847059}%
\pgfsetfillcolor{currentfill}%
\pgfsetfillopacity{0.130000}%
\pgfsetlinewidth{0.000000pt}%
\definecolor{currentstroke}{rgb}{0.000000,0.000000,0.000000}%
\pgfsetstrokecolor{currentstroke}%
\pgfsetdash{}{0pt}%
\pgfpathmoveto{\pgfqpoint{0.914528in}{1.002763in}}%
\pgfpathlineto{\pgfqpoint{0.930335in}{0.985188in}}%
\pgfpathlineto{\pgfqpoint{0.946153in}{0.967927in}}%
\pgfpathlineto{\pgfqpoint{0.954203in}{0.977441in}}%
\pgfpathlineto{\pgfqpoint{0.962228in}{0.987419in}}%
\pgfpathlineto{\pgfqpoint{0.946447in}{1.004176in}}%
\pgfpathlineto{\pgfqpoint{0.930677in}{1.021276in}}%
\pgfpathlineto{\pgfqpoint{0.922615in}{1.011802in}}%
\pgfpathlineto{\pgfqpoint{0.914528in}{1.002763in}}%
\pgfpathclose%
\pgfusepath{fill}%
\end{pgfscope}%
\begin{pgfscope}%
\pgfpathrectangle{\pgfqpoint{0.407592in}{0.222904in}}{\pgfqpoint{1.411163in}{1.411163in}}%
\pgfusepath{clip}%
\pgfsetbuttcap%
\pgfsetroundjoin%
\definecolor{currentfill}{rgb}{0.556863,0.654902,0.847059}%
\pgfsetfillcolor{currentfill}%
\pgfsetfillopacity{0.130000}%
\pgfsetlinewidth{0.000000pt}%
\definecolor{currentstroke}{rgb}{0.000000,0.000000,0.000000}%
\pgfsetstrokecolor{currentstroke}%
\pgfsetdash{}{0pt}%
\pgfpathmoveto{\pgfqpoint{1.120408in}{0.927876in}}%
\pgfpathlineto{\pgfqpoint{1.136296in}{0.918755in}}%
\pgfpathlineto{\pgfqpoint{1.152214in}{0.910754in}}%
\pgfpathlineto{\pgfqpoint{1.160016in}{0.925463in}}%
\pgfpathlineto{\pgfqpoint{1.167803in}{0.940398in}}%
\pgfpathlineto{\pgfqpoint{1.151900in}{0.948023in}}%
\pgfpathlineto{\pgfqpoint{1.136029in}{0.956689in}}%
\pgfpathlineto{\pgfqpoint{1.128227in}{0.942159in}}%
\pgfpathlineto{\pgfqpoint{1.120408in}{0.927876in}}%
\pgfpathclose%
\pgfusepath{fill}%
\end{pgfscope}%
\begin{pgfscope}%
\pgfpathrectangle{\pgfqpoint{0.407592in}{0.222904in}}{\pgfqpoint{1.411163in}{1.411163in}}%
\pgfusepath{clip}%
\pgfsetbuttcap%
\pgfsetroundjoin%
\definecolor{currentfill}{rgb}{0.556863,0.654902,0.847059}%
\pgfsetfillcolor{currentfill}%
\pgfsetfillopacity{0.130000}%
\pgfsetlinewidth{0.000000pt}%
\definecolor{currentstroke}{rgb}{0.000000,0.000000,0.000000}%
\pgfsetstrokecolor{currentstroke}%
\pgfsetdash{}{0pt}%
\pgfpathmoveto{\pgfqpoint{0.803393in}{1.098077in}}%
\pgfpathlineto{\pgfqpoint{0.819183in}{1.079049in}}%
\pgfpathlineto{\pgfqpoint{0.834983in}{1.060138in}}%
\pgfpathlineto{\pgfqpoint{0.843205in}{1.067581in}}%
\pgfpathlineto{\pgfqpoint{0.851397in}{1.075431in}}%
\pgfpathlineto{\pgfqpoint{0.835638in}{1.094044in}}%
\pgfpathlineto{\pgfqpoint{0.819889in}{1.112792in}}%
\pgfpathlineto{\pgfqpoint{0.811655in}{1.105246in}}%
\pgfpathlineto{\pgfqpoint{0.803393in}{1.098077in}}%
\pgfpathclose%
\pgfusepath{fill}%
\end{pgfscope}%
\begin{pgfscope}%
\pgfpathrectangle{\pgfqpoint{0.407592in}{0.222904in}}{\pgfqpoint{1.411163in}{1.411163in}}%
\pgfusepath{clip}%
\pgfsetbuttcap%
\pgfsetroundjoin%
\definecolor{currentfill}{rgb}{0.556863,0.654902,0.847059}%
\pgfsetfillcolor{currentfill}%
\pgfsetfillopacity{0.130000}%
\pgfsetlinewidth{0.000000pt}%
\definecolor{currentstroke}{rgb}{0.000000,0.000000,0.000000}%
\pgfsetstrokecolor{currentstroke}%
\pgfsetdash{}{0pt}%
\pgfpathmoveto{\pgfqpoint{1.167803in}{0.940398in}}%
\pgfpathlineto{\pgfqpoint{1.183741in}{0.933866in}}%
\pgfpathlineto{\pgfqpoint{1.199716in}{0.928468in}}%
\pgfpathlineto{\pgfqpoint{1.207477in}{0.943816in}}%
\pgfpathlineto{\pgfqpoint{1.215224in}{0.959322in}}%
\pgfpathlineto{\pgfqpoint{1.199261in}{0.964553in}}%
\pgfpathlineto{\pgfqpoint{1.183336in}{0.970828in}}%
\pgfpathlineto{\pgfqpoint{1.175576in}{0.955528in}}%
\pgfpathlineto{\pgfqpoint{1.167803in}{0.940398in}}%
\pgfpathclose%
\pgfusepath{fill}%
\end{pgfscope}%
\begin{pgfscope}%
\pgfpathrectangle{\pgfqpoint{0.407592in}{0.222904in}}{\pgfqpoint{1.411163in}{1.411163in}}%
\pgfusepath{clip}%
\pgfsetbuttcap%
\pgfsetroundjoin%
\definecolor{currentfill}{rgb}{0.556863,0.654902,0.847059}%
\pgfsetfillcolor{currentfill}%
\pgfsetfillopacity{0.130000}%
\pgfsetlinewidth{0.000000pt}%
\definecolor{currentstroke}{rgb}{0.000000,0.000000,0.000000}%
\pgfsetstrokecolor{currentstroke}%
\pgfsetdash{}{0pt}%
\pgfpathmoveto{\pgfqpoint{1.072986in}{0.922640in}}%
\pgfpathlineto{\pgfqpoint{1.088844in}{0.910918in}}%
\pgfpathlineto{\pgfqpoint{1.104724in}{0.900197in}}%
\pgfpathlineto{\pgfqpoint{1.112574in}{0.913875in}}%
\pgfpathlineto{\pgfqpoint{1.120408in}{0.927876in}}%
\pgfpathlineto{\pgfqpoint{1.104548in}{0.938038in}}%
\pgfpathlineto{\pgfqpoint{1.088712in}{0.949159in}}%
\pgfpathlineto{\pgfqpoint{1.080858in}{0.935729in}}%
\pgfpathlineto{\pgfqpoint{1.072986in}{0.922640in}}%
\pgfpathclose%
\pgfusepath{fill}%
\end{pgfscope}%
\begin{pgfscope}%
\pgfpathrectangle{\pgfqpoint{0.407592in}{0.222904in}}{\pgfqpoint{1.411163in}{1.411163in}}%
\pgfusepath{clip}%
\pgfsetbuttcap%
\pgfsetroundjoin%
\definecolor{currentfill}{rgb}{0.556863,0.654902,0.847059}%
\pgfsetfillcolor{currentfill}%
\pgfsetfillopacity{0.130000}%
\pgfsetlinewidth{0.000000pt}%
\definecolor{currentstroke}{rgb}{0.000000,0.000000,0.000000}%
\pgfsetstrokecolor{currentstroke}%
\pgfsetdash{}{0pt}%
\pgfpathmoveto{\pgfqpoint{1.325612in}{1.043802in}}%
\pgfpathlineto{\pgfqpoint{1.341783in}{1.043180in}}%
\pgfpathlineto{\pgfqpoint{1.358008in}{1.043360in}}%
\pgfpathlineto{\pgfqpoint{1.365634in}{1.059257in}}%
\pgfpathlineto{\pgfqpoint{1.373249in}{1.075222in}}%
\pgfpathlineto{\pgfqpoint{1.357034in}{1.075246in}}%
\pgfpathlineto{\pgfqpoint{1.340873in}{1.076026in}}%
\pgfpathlineto{\pgfqpoint{1.333249in}{1.059887in}}%
\pgfpathlineto{\pgfqpoint{1.325612in}{1.043802in}}%
\pgfpathclose%
\pgfusepath{fill}%
\end{pgfscope}%
\begin{pgfscope}%
\pgfpathrectangle{\pgfqpoint{0.407592in}{0.222904in}}{\pgfqpoint{1.411163in}{1.411163in}}%
\pgfusepath{clip}%
\pgfsetbuttcap%
\pgfsetroundjoin%
\definecolor{currentfill}{rgb}{0.556863,0.654902,0.847059}%
\pgfsetfillcolor{currentfill}%
\pgfsetfillopacity{0.130000}%
\pgfsetlinewidth{0.000000pt}%
\definecolor{currentstroke}{rgb}{0.000000,0.000000,0.000000}%
\pgfsetstrokecolor{currentstroke}%
\pgfsetdash{}{0pt}%
\pgfpathmoveto{\pgfqpoint{1.215224in}{0.959322in}}%
\pgfpathlineto{\pgfqpoint{1.231229in}{0.955153in}}%
\pgfpathlineto{\pgfqpoint{1.247279in}{0.952054in}}%
\pgfpathlineto{\pgfqpoint{1.255003in}{0.967721in}}%
\pgfpathlineto{\pgfqpoint{1.262713in}{0.983505in}}%
\pgfpathlineto{\pgfqpoint{1.246673in}{0.986620in}}%
\pgfpathlineto{\pgfqpoint{1.230679in}{0.990721in}}%
\pgfpathlineto{\pgfqpoint{1.222958in}{0.974963in}}%
\pgfpathlineto{\pgfqpoint{1.215224in}{0.959322in}}%
\pgfpathclose%
\pgfusepath{fill}%
\end{pgfscope}%
\begin{pgfscope}%
\pgfpathrectangle{\pgfqpoint{0.407592in}{0.222904in}}{\pgfqpoint{1.411163in}{1.411163in}}%
\pgfusepath{clip}%
\pgfsetbuttcap%
\pgfsetroundjoin%
\definecolor{currentfill}{rgb}{0.556863,0.654902,0.847059}%
\pgfsetfillcolor{currentfill}%
\pgfsetfillopacity{0.130000}%
\pgfsetlinewidth{0.000000pt}%
\definecolor{currentstroke}{rgb}{0.000000,0.000000,0.000000}%
\pgfsetstrokecolor{currentstroke}%
\pgfsetdash{}{0pt}%
\pgfpathmoveto{\pgfqpoint{0.691968in}{1.201968in}}%
\pgfpathlineto{\pgfqpoint{0.707747in}{1.182317in}}%
\pgfpathlineto{\pgfqpoint{0.723535in}{1.162712in}}%
\pgfpathlineto{\pgfqpoint{0.731943in}{1.168737in}}%
\pgfpathlineto{\pgfqpoint{0.740319in}{1.175101in}}%
\pgfpathlineto{\pgfqpoint{0.724573in}{1.194532in}}%
\pgfpathlineto{\pgfqpoint{0.708836in}{1.214017in}}%
\pgfpathlineto{\pgfqpoint{0.700418in}{1.207834in}}%
\pgfpathlineto{\pgfqpoint{0.691968in}{1.201968in}}%
\pgfpathclose%
\pgfusepath{fill}%
\end{pgfscope}%
\begin{pgfscope}%
\pgfpathrectangle{\pgfqpoint{0.407592in}{0.222904in}}{\pgfqpoint{1.411163in}{1.411163in}}%
\pgfusepath{clip}%
\pgfsetbuttcap%
\pgfsetroundjoin%
\definecolor{currentfill}{rgb}{0.556863,0.654902,0.847059}%
\pgfsetfillcolor{currentfill}%
\pgfsetfillopacity{0.130000}%
\pgfsetlinewidth{0.000000pt}%
\definecolor{currentstroke}{rgb}{0.000000,0.000000,0.000000}%
\pgfsetstrokecolor{currentstroke}%
\pgfsetdash{}{0pt}%
\pgfpathmoveto{\pgfqpoint{0.834983in}{1.060138in}}%
\pgfpathlineto{\pgfqpoint{0.850791in}{1.041362in}}%
\pgfpathlineto{\pgfqpoint{0.866609in}{1.022743in}}%
\pgfpathlineto{\pgfqpoint{0.874791in}{1.030496in}}%
\pgfpathlineto{\pgfqpoint{0.882944in}{1.038690in}}%
\pgfpathlineto{\pgfqpoint{0.867166in}{1.056972in}}%
\pgfpathlineto{\pgfqpoint{0.851397in}{1.075431in}}%
\pgfpathlineto{\pgfqpoint{0.843205in}{1.067581in}}%
\pgfpathlineto{\pgfqpoint{0.834983in}{1.060138in}}%
\pgfpathclose%
\pgfusepath{fill}%
\end{pgfscope}%
\begin{pgfscope}%
\pgfpathrectangle{\pgfqpoint{0.407592in}{0.222904in}}{\pgfqpoint{1.411163in}{1.411163in}}%
\pgfusepath{clip}%
\pgfsetbuttcap%
\pgfsetroundjoin%
\definecolor{currentfill}{rgb}{0.556863,0.654902,0.847059}%
\pgfsetfillcolor{currentfill}%
\pgfsetfillopacity{0.130000}%
\pgfsetlinewidth{0.000000pt}%
\definecolor{currentstroke}{rgb}{0.000000,0.000000,0.000000}%
\pgfsetstrokecolor{currentstroke}%
\pgfsetdash{}{0pt}%
\pgfpathmoveto{\pgfqpoint{1.025476in}{0.925000in}}%
\pgfpathlineto{\pgfqpoint{1.041323in}{0.910928in}}%
\pgfpathlineto{\pgfqpoint{1.057187in}{0.897660in}}%
\pgfpathlineto{\pgfqpoint{1.065096in}{0.909935in}}%
\pgfpathlineto{\pgfqpoint{1.072986in}{0.922640in}}%
\pgfpathlineto{\pgfqpoint{1.057148in}{0.935262in}}%
\pgfpathlineto{\pgfqpoint{1.041329in}{0.948685in}}%
\pgfpathlineto{\pgfqpoint{1.033413in}{0.936627in}}%
\pgfpathlineto{\pgfqpoint{1.025476in}{0.925000in}}%
\pgfpathclose%
\pgfusepath{fill}%
\end{pgfscope}%
\begin{pgfscope}%
\pgfpathrectangle{\pgfqpoint{0.407592in}{0.222904in}}{\pgfqpoint{1.411163in}{1.411163in}}%
\pgfusepath{clip}%
\pgfsetbuttcap%
\pgfsetroundjoin%
\definecolor{currentfill}{rgb}{0.556863,0.654902,0.847059}%
\pgfsetfillcolor{currentfill}%
\pgfsetfillopacity{0.130000}%
\pgfsetlinewidth{0.000000pt}%
\definecolor{currentstroke}{rgb}{0.000000,0.000000,0.000000}%
\pgfsetstrokecolor{currentstroke}%
\pgfsetdash{}{0pt}%
\pgfpathmoveto{\pgfqpoint{0.946153in}{0.967927in}}%
\pgfpathlineto{\pgfqpoint{0.961982in}{0.951031in}}%
\pgfpathlineto{\pgfqpoint{0.977821in}{0.934563in}}%
\pgfpathlineto{\pgfqpoint{0.985835in}{0.944611in}}%
\pgfpathlineto{\pgfqpoint{0.993826in}{0.955152in}}%
\pgfpathlineto{\pgfqpoint{0.978021in}{0.971057in}}%
\pgfpathlineto{\pgfqpoint{0.962228in}{0.987419in}}%
\pgfpathlineto{\pgfqpoint{0.954203in}{0.977441in}}%
\pgfpathlineto{\pgfqpoint{0.946153in}{0.967927in}}%
\pgfpathclose%
\pgfusepath{fill}%
\end{pgfscope}%
\begin{pgfscope}%
\pgfpathrectangle{\pgfqpoint{0.407592in}{0.222904in}}{\pgfqpoint{1.411163in}{1.411163in}}%
\pgfusepath{clip}%
\pgfsetbuttcap%
\pgfsetroundjoin%
\definecolor{currentfill}{rgb}{0.556863,0.654902,0.847059}%
\pgfsetfillcolor{currentfill}%
\pgfsetfillopacity{0.130000}%
\pgfsetlinewidth{0.000000pt}%
\definecolor{currentstroke}{rgb}{0.000000,0.000000,0.000000}%
\pgfsetstrokecolor{currentstroke}%
\pgfsetdash{}{0pt}%
\pgfpathmoveto{\pgfqpoint{1.262713in}{0.983505in}}%
\pgfpathlineto{\pgfqpoint{1.278801in}{0.981371in}}%
\pgfpathlineto{\pgfqpoint{1.294939in}{0.980202in}}%
\pgfpathlineto{\pgfqpoint{1.302626in}{0.995968in}}%
\pgfpathlineto{\pgfqpoint{1.310301in}{1.011832in}}%
\pgfpathlineto{\pgfqpoint{1.294173in}{1.013150in}}%
\pgfpathlineto{\pgfqpoint{1.278095in}{1.015365in}}%
\pgfpathlineto{\pgfqpoint{1.270411in}{0.999391in}}%
\pgfpathlineto{\pgfqpoint{1.262713in}{0.983505in}}%
\pgfpathclose%
\pgfusepath{fill}%
\end{pgfscope}%
\begin{pgfscope}%
\pgfpathrectangle{\pgfqpoint{0.407592in}{0.222904in}}{\pgfqpoint{1.411163in}{1.411163in}}%
\pgfusepath{clip}%
\pgfsetbuttcap%
\pgfsetroundjoin%
\definecolor{currentfill}{rgb}{0.556863,0.654902,0.847059}%
\pgfsetfillcolor{currentfill}%
\pgfsetfillopacity{0.130000}%
\pgfsetlinewidth{0.000000pt}%
\definecolor{currentstroke}{rgb}{0.000000,0.000000,0.000000}%
\pgfsetstrokecolor{currentstroke}%
\pgfsetdash{}{0pt}%
\pgfpathmoveto{\pgfqpoint{0.723535in}{1.162712in}}%
\pgfpathlineto{\pgfqpoint{0.739332in}{1.143161in}}%
\pgfpathlineto{\pgfqpoint{0.755138in}{1.123671in}}%
\pgfpathlineto{\pgfqpoint{0.763504in}{1.129868in}}%
\pgfpathlineto{\pgfqpoint{0.771838in}{1.136431in}}%
\pgfpathlineto{\pgfqpoint{0.756074in}{1.155731in}}%
\pgfpathlineto{\pgfqpoint{0.740319in}{1.175101in}}%
\pgfpathlineto{\pgfqpoint{0.731943in}{1.168737in}}%
\pgfpathlineto{\pgfqpoint{0.723535in}{1.162712in}}%
\pgfpathclose%
\pgfusepath{fill}%
\end{pgfscope}%
\begin{pgfscope}%
\pgfpathrectangle{\pgfqpoint{0.407592in}{0.222904in}}{\pgfqpoint{1.411163in}{1.411163in}}%
\pgfusepath{clip}%
\pgfsetbuttcap%
\pgfsetroundjoin%
\definecolor{currentfill}{rgb}{0.556863,0.654902,0.847059}%
\pgfsetfillcolor{currentfill}%
\pgfsetfillopacity{0.130000}%
\pgfsetlinewidth{0.000000pt}%
\definecolor{currentstroke}{rgb}{0.000000,0.000000,0.000000}%
\pgfsetstrokecolor{currentstroke}%
\pgfsetdash{}{0pt}%
\pgfpathmoveto{\pgfqpoint{0.866609in}{1.022743in}}%
\pgfpathlineto{\pgfqpoint{0.882436in}{1.004305in}}%
\pgfpathlineto{\pgfqpoint{0.898273in}{0.986078in}}%
\pgfpathlineto{\pgfqpoint{0.906414in}{0.994182in}}%
\pgfpathlineto{\pgfqpoint{0.914528in}{1.002763in}}%
\pgfpathlineto{\pgfqpoint{0.898731in}{1.020610in}}%
\pgfpathlineto{\pgfqpoint{0.882944in}{1.038690in}}%
\pgfpathlineto{\pgfqpoint{0.874791in}{1.030496in}}%
\pgfpathlineto{\pgfqpoint{0.866609in}{1.022743in}}%
\pgfpathclose%
\pgfusepath{fill}%
\end{pgfscope}%
\begin{pgfscope}%
\pgfpathrectangle{\pgfqpoint{0.407592in}{0.222904in}}{\pgfqpoint{1.411163in}{1.411163in}}%
\pgfusepath{clip}%
\pgfsetbuttcap%
\pgfsetroundjoin%
\definecolor{currentfill}{rgb}{0.556863,0.654902,0.847059}%
\pgfsetfillcolor{currentfill}%
\pgfsetfillopacity{0.130000}%
\pgfsetlinewidth{0.000000pt}%
\definecolor{currentstroke}{rgb}{0.000000,0.000000,0.000000}%
\pgfsetstrokecolor{currentstroke}%
\pgfsetdash{}{0pt}%
\pgfpathmoveto{\pgfqpoint{0.977821in}{0.934563in}}%
\pgfpathlineto{\pgfqpoint{0.993673in}{0.918594in}}%
\pgfpathlineto{\pgfqpoint{1.009537in}{0.903208in}}%
\pgfpathlineto{\pgfqpoint{1.017518in}{0.913846in}}%
\pgfpathlineto{\pgfqpoint{1.025476in}{0.925000in}}%
\pgfpathlineto{\pgfqpoint{1.009644in}{0.939772in}}%
\pgfpathlineto{\pgfqpoint{0.993826in}{0.955152in}}%
\pgfpathlineto{\pgfqpoint{0.985835in}{0.944611in}}%
\pgfpathlineto{\pgfqpoint{0.977821in}{0.934563in}}%
\pgfpathclose%
\pgfusepath{fill}%
\end{pgfscope}%
\begin{pgfscope}%
\pgfpathrectangle{\pgfqpoint{0.407592in}{0.222904in}}{\pgfqpoint{1.411163in}{1.411163in}}%
\pgfusepath{clip}%
\pgfsetbuttcap%
\pgfsetroundjoin%
\definecolor{currentfill}{rgb}{0.556863,0.654902,0.847059}%
\pgfsetfillcolor{currentfill}%
\pgfsetfillopacity{0.130000}%
\pgfsetlinewidth{0.000000pt}%
\definecolor{currentstroke}{rgb}{0.000000,0.000000,0.000000}%
\pgfsetstrokecolor{currentstroke}%
\pgfsetdash{}{0pt}%
\pgfpathmoveto{\pgfqpoint{1.104724in}{0.900197in}}%
\pgfpathlineto{\pgfqpoint{1.120630in}{0.890580in}}%
\pgfpathlineto{\pgfqpoint{1.136564in}{0.882163in}}%
\pgfpathlineto{\pgfqpoint{1.144397in}{0.896307in}}%
\pgfpathlineto{\pgfqpoint{1.152214in}{0.910754in}}%
\pgfpathlineto{\pgfqpoint{1.136296in}{0.918755in}}%
\pgfpathlineto{\pgfqpoint{1.120408in}{0.927876in}}%
\pgfpathlineto{\pgfqpoint{1.112574in}{0.913875in}}%
\pgfpathlineto{\pgfqpoint{1.104724in}{0.900197in}}%
\pgfpathclose%
\pgfusepath{fill}%
\end{pgfscope}%
\begin{pgfscope}%
\pgfpathrectangle{\pgfqpoint{0.407592in}{0.222904in}}{\pgfqpoint{1.411163in}{1.411163in}}%
\pgfusepath{clip}%
\pgfsetbuttcap%
\pgfsetroundjoin%
\definecolor{currentfill}{rgb}{0.556863,0.654902,0.847059}%
\pgfsetfillcolor{currentfill}%
\pgfsetfillopacity{0.130000}%
\pgfsetlinewidth{0.000000pt}%
\definecolor{currentstroke}{rgb}{0.000000,0.000000,0.000000}%
\pgfsetstrokecolor{currentstroke}%
\pgfsetdash{}{0pt}%
\pgfpathmoveto{\pgfqpoint{1.373249in}{1.075222in}}%
\pgfpathlineto{\pgfqpoint{1.389519in}{1.075932in}}%
\pgfpathlineto{\pgfqpoint{1.405848in}{1.077353in}}%
\pgfpathlineto{\pgfqpoint{1.413439in}{1.093134in}}%
\pgfpathlineto{\pgfqpoint{1.397116in}{1.091843in}}%
\pgfpathlineto{\pgfqpoint{1.380851in}{1.091245in}}%
\pgfusepath{fill}%
\end{pgfscope}%
\begin{pgfscope}%
\pgfpathrectangle{\pgfqpoint{0.407592in}{0.222904in}}{\pgfqpoint{1.411163in}{1.411163in}}%
\pgfusepath{clip}%
\pgfsetbuttcap%
\pgfsetroundjoin%
\definecolor{currentfill}{rgb}{0.556863,0.654902,0.847059}%
\pgfsetfillcolor{currentfill}%
\pgfsetfillopacity{0.130000}%
\pgfsetlinewidth{0.000000pt}%
\definecolor{currentstroke}{rgb}{0.000000,0.000000,0.000000}%
\pgfsetstrokecolor{currentstroke}%
\pgfsetdash{}{0pt}%
\pgfpathmoveto{\pgfqpoint{1.152214in}{0.910754in}}%
\pgfpathlineto{\pgfqpoint{1.168165in}{0.903937in}}%
\pgfpathlineto{\pgfqpoint{1.184153in}{0.898356in}}%
\pgfpathlineto{\pgfqpoint{1.191941in}{0.913305in}}%
\pgfpathlineto{\pgfqpoint{1.199716in}{0.928468in}}%
\pgfpathlineto{\pgfqpoint{1.183741in}{0.933866in}}%
\pgfpathlineto{\pgfqpoint{1.167803in}{0.940398in}}%
\pgfpathlineto{\pgfqpoint{1.160016in}{0.925463in}}%
\pgfpathlineto{\pgfqpoint{1.152214in}{0.910754in}}%
\pgfpathclose%
\pgfusepath{fill}%
\end{pgfscope}%
\begin{pgfscope}%
\pgfpathrectangle{\pgfqpoint{0.407592in}{0.222904in}}{\pgfqpoint{1.411163in}{1.411163in}}%
\pgfusepath{clip}%
\pgfsetbuttcap%
\pgfsetroundjoin%
\definecolor{currentfill}{rgb}{0.556863,0.654902,0.847059}%
\pgfsetfillcolor{currentfill}%
\pgfsetfillopacity{0.130000}%
\pgfsetlinewidth{0.000000pt}%
\definecolor{currentstroke}{rgb}{0.000000,0.000000,0.000000}%
\pgfsetstrokecolor{currentstroke}%
\pgfsetdash{}{0pt}%
\pgfpathmoveto{\pgfqpoint{0.755138in}{1.123671in}}%
\pgfpathlineto{\pgfqpoint{0.770952in}{1.104251in}}%
\pgfpathlineto{\pgfqpoint{0.786776in}{1.084910in}}%
\pgfpathlineto{\pgfqpoint{0.795100in}{1.091296in}}%
\pgfpathlineto{\pgfqpoint{0.803393in}{1.098077in}}%
\pgfpathlineto{\pgfqpoint{0.787611in}{1.117209in}}%
\pgfpathlineto{\pgfqpoint{0.771838in}{1.136431in}}%
\pgfpathlineto{\pgfqpoint{0.763504in}{1.129868in}}%
\pgfpathlineto{\pgfqpoint{0.755138in}{1.123671in}}%
\pgfpathclose%
\pgfusepath{fill}%
\end{pgfscope}%
\begin{pgfscope}%
\pgfpathrectangle{\pgfqpoint{0.407592in}{0.222904in}}{\pgfqpoint{1.411163in}{1.411163in}}%
\pgfusepath{clip}%
\pgfsetbuttcap%
\pgfsetroundjoin%
\definecolor{currentfill}{rgb}{0.556863,0.654902,0.847059}%
\pgfsetfillcolor{currentfill}%
\pgfsetfillopacity{0.130000}%
\pgfsetlinewidth{0.000000pt}%
\definecolor{currentstroke}{rgb}{0.000000,0.000000,0.000000}%
\pgfsetstrokecolor{currentstroke}%
\pgfsetdash{}{0pt}%
\pgfpathmoveto{\pgfqpoint{1.310301in}{1.011832in}}%
\pgfpathlineto{\pgfqpoint{1.326481in}{1.011393in}}%
\pgfpathlineto{\pgfqpoint{1.342717in}{1.011809in}}%
\pgfpathlineto{\pgfqpoint{1.350369in}{1.027540in}}%
\pgfpathlineto{\pgfqpoint{1.358008in}{1.043360in}}%
\pgfpathlineto{\pgfqpoint{1.341783in}{1.043180in}}%
\pgfpathlineto{\pgfqpoint{1.325612in}{1.043802in}}%
\pgfpathlineto{\pgfqpoint{1.317963in}{1.027780in}}%
\pgfpathlineto{\pgfqpoint{1.310301in}{1.011832in}}%
\pgfpathclose%
\pgfusepath{fill}%
\end{pgfscope}%
\begin{pgfscope}%
\pgfpathrectangle{\pgfqpoint{0.407592in}{0.222904in}}{\pgfqpoint{1.411163in}{1.411163in}}%
\pgfusepath{clip}%
\pgfsetbuttcap%
\pgfsetroundjoin%
\definecolor{currentfill}{rgb}{0.556863,0.654902,0.847059}%
\pgfsetfillcolor{currentfill}%
\pgfsetfillopacity{0.130000}%
\pgfsetlinewidth{0.000000pt}%
\definecolor{currentstroke}{rgb}{0.000000,0.000000,0.000000}%
\pgfsetstrokecolor{currentstroke}%
\pgfsetdash{}{0pt}%
\pgfpathmoveto{\pgfqpoint{1.057187in}{0.897660in}}%
\pgfpathlineto{\pgfqpoint{1.073069in}{0.885309in}}%
\pgfpathlineto{\pgfqpoint{1.088972in}{0.873997in}}%
\pgfpathlineto{\pgfqpoint{1.096857in}{0.886887in}}%
\pgfpathlineto{\pgfqpoint{1.104724in}{0.900197in}}%
\pgfpathlineto{\pgfqpoint{1.088844in}{0.910918in}}%
\pgfpathlineto{\pgfqpoint{1.072986in}{0.922640in}}%
\pgfpathlineto{\pgfqpoint{1.065096in}{0.909935in}}%
\pgfpathlineto{\pgfqpoint{1.057187in}{0.897660in}}%
\pgfpathclose%
\pgfusepath{fill}%
\end{pgfscope}%
\begin{pgfscope}%
\pgfpathrectangle{\pgfqpoint{0.407592in}{0.222904in}}{\pgfqpoint{1.411163in}{1.411163in}}%
\pgfusepath{clip}%
\pgfsetbuttcap%
\pgfsetroundjoin%
\definecolor{currentfill}{rgb}{0.556863,0.654902,0.847059}%
\pgfsetfillcolor{currentfill}%
\pgfsetfillopacity{0.130000}%
\pgfsetlinewidth{0.000000pt}%
\definecolor{currentstroke}{rgb}{0.000000,0.000000,0.000000}%
\pgfsetstrokecolor{currentstroke}%
\pgfsetdash{}{0pt}%
\pgfpathmoveto{\pgfqpoint{1.199716in}{0.928468in}}%
\pgfpathlineto{\pgfqpoint{1.215732in}{0.924229in}}%
\pgfpathlineto{\pgfqpoint{1.231793in}{0.921155in}}%
\pgfpathlineto{\pgfqpoint{1.239543in}{0.936525in}}%
\pgfpathlineto{\pgfqpoint{1.247279in}{0.952054in}}%
\pgfpathlineto{\pgfqpoint{1.231229in}{0.955153in}}%
\pgfpathlineto{\pgfqpoint{1.215224in}{0.959322in}}%
\pgfpathlineto{\pgfqpoint{1.207477in}{0.943816in}}%
\pgfpathlineto{\pgfqpoint{1.199716in}{0.928468in}}%
\pgfpathclose%
\pgfusepath{fill}%
\end{pgfscope}%
\begin{pgfscope}%
\pgfpathrectangle{\pgfqpoint{0.407592in}{0.222904in}}{\pgfqpoint{1.411163in}{1.411163in}}%
\pgfusepath{clip}%
\pgfsetbuttcap%
\pgfsetroundjoin%
\definecolor{currentfill}{rgb}{0.556863,0.654902,0.847059}%
\pgfsetfillcolor{currentfill}%
\pgfsetfillopacity{0.130000}%
\pgfsetlinewidth{0.000000pt}%
\definecolor{currentstroke}{rgb}{0.000000,0.000000,0.000000}%
\pgfsetstrokecolor{currentstroke}%
\pgfsetdash{}{0pt}%
\pgfpathmoveto{\pgfqpoint{0.898273in}{0.986078in}}%
\pgfpathlineto{\pgfqpoint{0.914119in}{0.968095in}}%
\pgfpathlineto{\pgfqpoint{0.929975in}{0.950400in}}%
\pgfpathlineto{\pgfqpoint{0.938078in}{0.958905in}}%
\pgfpathlineto{\pgfqpoint{0.946153in}{0.967927in}}%
\pgfpathlineto{\pgfqpoint{0.930335in}{0.985188in}}%
\pgfpathlineto{\pgfqpoint{0.914528in}{1.002763in}}%
\pgfpathlineto{\pgfqpoint{0.906414in}{0.994182in}}%
\pgfpathlineto{\pgfqpoint{0.898273in}{0.986078in}}%
\pgfpathclose%
\pgfusepath{fill}%
\end{pgfscope}%
\begin{pgfscope}%
\pgfpathrectangle{\pgfqpoint{0.407592in}{0.222904in}}{\pgfqpoint{1.411163in}{1.411163in}}%
\pgfusepath{clip}%
\pgfsetbuttcap%
\pgfsetroundjoin%
\definecolor{currentfill}{rgb}{0.556863,0.654902,0.847059}%
\pgfsetfillcolor{currentfill}%
\pgfsetfillopacity{0.130000}%
\pgfsetlinewidth{0.000000pt}%
\definecolor{currentstroke}{rgb}{0.000000,0.000000,0.000000}%
\pgfsetstrokecolor{currentstroke}%
\pgfsetdash{}{0pt}%
\pgfpathmoveto{\pgfqpoint{0.786776in}{1.084910in}}%
\pgfpathlineto{\pgfqpoint{0.802608in}{1.065660in}}%
\pgfpathlineto{\pgfqpoint{0.818449in}{1.046516in}}%
\pgfpathlineto{\pgfqpoint{0.826731in}{1.053113in}}%
\pgfpathlineto{\pgfqpoint{0.834983in}{1.060138in}}%
\pgfpathlineto{\pgfqpoint{0.819183in}{1.079049in}}%
\pgfpathlineto{\pgfqpoint{0.803393in}{1.098077in}}%
\pgfpathlineto{\pgfqpoint{0.795100in}{1.091296in}}%
\pgfpathlineto{\pgfqpoint{0.786776in}{1.084910in}}%
\pgfpathclose%
\pgfusepath{fill}%
\end{pgfscope}%
\begin{pgfscope}%
\pgfpathrectangle{\pgfqpoint{0.407592in}{0.222904in}}{\pgfqpoint{1.411163in}{1.411163in}}%
\pgfusepath{clip}%
\pgfsetbuttcap%
\pgfsetroundjoin%
\definecolor{currentfill}{rgb}{0.556863,0.654902,0.847059}%
\pgfsetfillcolor{currentfill}%
\pgfsetfillopacity{0.130000}%
\pgfsetlinewidth{0.000000pt}%
\definecolor{currentstroke}{rgb}{0.000000,0.000000,0.000000}%
\pgfsetstrokecolor{currentstroke}%
\pgfsetdash{}{0pt}%
\pgfpathmoveto{\pgfqpoint{1.247279in}{0.952054in}}%
\pgfpathlineto{\pgfqpoint{1.263377in}{0.950018in}}%
\pgfpathlineto{\pgfqpoint{1.279526in}{0.949027in}}%
\pgfpathlineto{\pgfqpoint{1.287239in}{0.964549in}}%
\pgfpathlineto{\pgfqpoint{1.294939in}{0.980202in}}%
\pgfpathlineto{\pgfqpoint{1.278801in}{0.981371in}}%
\pgfpathlineto{\pgfqpoint{1.262713in}{0.983505in}}%
\pgfpathlineto{\pgfqpoint{1.255003in}{0.967721in}}%
\pgfpathlineto{\pgfqpoint{1.247279in}{0.952054in}}%
\pgfpathclose%
\pgfusepath{fill}%
\end{pgfscope}%
\begin{pgfscope}%
\pgfpathrectangle{\pgfqpoint{0.407592in}{0.222904in}}{\pgfqpoint{1.411163in}{1.411163in}}%
\pgfusepath{clip}%
\pgfsetbuttcap%
\pgfsetroundjoin%
\definecolor{currentfill}{rgb}{0.556863,0.654902,0.847059}%
\pgfsetfillcolor{currentfill}%
\pgfsetfillopacity{0.130000}%
\pgfsetlinewidth{0.000000pt}%
\definecolor{currentstroke}{rgb}{0.000000,0.000000,0.000000}%
\pgfsetstrokecolor{currentstroke}%
\pgfsetdash{}{0pt}%
\pgfpathmoveto{\pgfqpoint{1.009537in}{0.903208in}}%
\pgfpathlineto{\pgfqpoint{1.025415in}{0.888506in}}%
\pgfpathlineto{\pgfqpoint{1.041308in}{0.874601in}}%
\pgfpathlineto{\pgfqpoint{1.049258in}{0.885864in}}%
\pgfpathlineto{\pgfqpoint{1.057187in}{0.897660in}}%
\pgfpathlineto{\pgfqpoint{1.041323in}{0.910928in}}%
\pgfpathlineto{\pgfqpoint{1.025476in}{0.925000in}}%
\pgfpathlineto{\pgfqpoint{1.017518in}{0.913846in}}%
\pgfpathlineto{\pgfqpoint{1.009537in}{0.903208in}}%
\pgfpathclose%
\pgfusepath{fill}%
\end{pgfscope}%
\begin{pgfscope}%
\pgfpathrectangle{\pgfqpoint{0.407592in}{0.222904in}}{\pgfqpoint{1.411163in}{1.411163in}}%
\pgfusepath{clip}%
\pgfsetbuttcap%
\pgfsetroundjoin%
\definecolor{currentfill}{rgb}{0.556863,0.654902,0.847059}%
\pgfsetfillcolor{currentfill}%
\pgfsetfillopacity{0.130000}%
\pgfsetlinewidth{0.000000pt}%
\definecolor{currentstroke}{rgb}{0.000000,0.000000,0.000000}%
\pgfsetstrokecolor{currentstroke}%
\pgfsetdash{}{0pt}%
\pgfpathmoveto{\pgfqpoint{1.358008in}{1.043360in}}%
\pgfpathlineto{\pgfqpoint{1.374289in}{1.044317in}}%
\pgfpathlineto{\pgfqpoint{1.390629in}{1.046024in}}%
\pgfpathlineto{\pgfqpoint{1.398245in}{1.061646in}}%
\pgfpathlineto{\pgfqpoint{1.405848in}{1.077353in}}%
\pgfpathlineto{\pgfqpoint{1.389519in}{1.075932in}}%
\pgfpathlineto{\pgfqpoint{1.373249in}{1.075222in}}%
\pgfpathlineto{\pgfqpoint{1.365634in}{1.059257in}}%
\pgfpathlineto{\pgfqpoint{1.358008in}{1.043360in}}%
\pgfpathclose%
\pgfusepath{fill}%
\end{pgfscope}%
\begin{pgfscope}%
\pgfpathrectangle{\pgfqpoint{0.407592in}{0.222904in}}{\pgfqpoint{1.411163in}{1.411163in}}%
\pgfusepath{clip}%
\pgfsetbuttcap%
\pgfsetroundjoin%
\definecolor{currentfill}{rgb}{0.556863,0.654902,0.847059}%
\pgfsetfillcolor{currentfill}%
\pgfsetfillopacity{0.130000}%
\pgfsetlinewidth{0.000000pt}%
\definecolor{currentstroke}{rgb}{0.000000,0.000000,0.000000}%
\pgfsetstrokecolor{currentstroke}%
\pgfsetdash{}{0pt}%
\pgfpathmoveto{\pgfqpoint{0.929975in}{0.950400in}}%
\pgfpathlineto{\pgfqpoint{0.945841in}{0.933042in}}%
\pgfpathlineto{\pgfqpoint{0.961718in}{0.916081in}}%
\pgfpathlineto{\pgfqpoint{0.969783in}{0.925042in}}%
\pgfpathlineto{\pgfqpoint{0.977821in}{0.934563in}}%
\pgfpathlineto{\pgfqpoint{0.961982in}{0.951031in}}%
\pgfpathlineto{\pgfqpoint{0.946153in}{0.967927in}}%
\pgfpathlineto{\pgfqpoint{0.938078in}{0.958905in}}%
\pgfpathlineto{\pgfqpoint{0.929975in}{0.950400in}}%
\pgfpathclose%
\pgfusepath{fill}%
\end{pgfscope}%
\begin{pgfscope}%
\pgfpathrectangle{\pgfqpoint{0.407592in}{0.222904in}}{\pgfqpoint{1.411163in}{1.411163in}}%
\pgfusepath{clip}%
\pgfsetbuttcap%
\pgfsetroundjoin%
\definecolor{currentfill}{rgb}{0.556863,0.654902,0.847059}%
\pgfsetfillcolor{currentfill}%
\pgfsetfillopacity{0.130000}%
\pgfsetlinewidth{0.000000pt}%
\definecolor{currentstroke}{rgb}{0.000000,0.000000,0.000000}%
\pgfsetstrokecolor{currentstroke}%
\pgfsetdash{}{0pt}%
\pgfpathmoveto{\pgfqpoint{0.818449in}{1.046516in}}%
\pgfpathlineto{\pgfqpoint{0.834299in}{1.027494in}}%
\pgfpathlineto{\pgfqpoint{0.850158in}{1.008613in}}%
\pgfpathlineto{\pgfqpoint{0.858399in}{1.015445in}}%
\pgfpathlineto{\pgfqpoint{0.866609in}{1.022743in}}%
\pgfpathlineto{\pgfqpoint{0.850791in}{1.041362in}}%
\pgfpathlineto{\pgfqpoint{0.834983in}{1.060138in}}%
\pgfpathlineto{\pgfqpoint{0.826731in}{1.053113in}}%
\pgfpathlineto{\pgfqpoint{0.818449in}{1.046516in}}%
\pgfpathclose%
\pgfusepath{fill}%
\end{pgfscope}%
\begin{pgfscope}%
\pgfpathrectangle{\pgfqpoint{0.407592in}{0.222904in}}{\pgfqpoint{1.411163in}{1.411163in}}%
\pgfusepath{clip}%
\pgfsetbuttcap%
\pgfsetroundjoin%
\definecolor{currentfill}{rgb}{0.556863,0.654902,0.847059}%
\pgfsetfillcolor{currentfill}%
\pgfsetfillopacity{0.130000}%
\pgfsetlinewidth{0.000000pt}%
\definecolor{currentstroke}{rgb}{0.000000,0.000000,0.000000}%
\pgfsetstrokecolor{currentstroke}%
\pgfsetdash{}{0pt}%
\pgfpathmoveto{\pgfqpoint{0.674969in}{1.191205in}}%
\pgfpathlineto{\pgfqpoint{0.690791in}{1.171432in}}%
\pgfpathlineto{\pgfqpoint{0.706621in}{1.151702in}}%
\pgfpathlineto{\pgfqpoint{0.715095in}{1.157033in}}%
\pgfpathlineto{\pgfqpoint{0.723535in}{1.162712in}}%
\pgfpathlineto{\pgfqpoint{0.707747in}{1.182317in}}%
\pgfpathlineto{\pgfqpoint{0.691968in}{1.201968in}}%
\pgfpathlineto{\pgfqpoint{0.683485in}{1.196424in}}%
\pgfpathlineto{\pgfqpoint{0.674969in}{1.191205in}}%
\pgfpathclose%
\pgfusepath{fill}%
\end{pgfscope}%
\begin{pgfscope}%
\pgfpathrectangle{\pgfqpoint{0.407592in}{0.222904in}}{\pgfqpoint{1.411163in}{1.411163in}}%
\pgfusepath{clip}%
\pgfsetbuttcap%
\pgfsetroundjoin%
\definecolor{currentfill}{rgb}{0.556863,0.654902,0.847059}%
\pgfsetfillcolor{currentfill}%
\pgfsetfillopacity{0.130000}%
\pgfsetlinewidth{0.000000pt}%
\definecolor{currentstroke}{rgb}{0.000000,0.000000,0.000000}%
\pgfsetstrokecolor{currentstroke}%
\pgfsetdash{}{0pt}%
\pgfpathmoveto{\pgfqpoint{1.136564in}{0.882163in}}%
\pgfpathlineto{\pgfqpoint{1.152530in}{0.875031in}}%
\pgfpathlineto{\pgfqpoint{1.168532in}{0.869248in}}%
\pgfpathlineto{\pgfqpoint{1.176350in}{0.883656in}}%
\pgfpathlineto{\pgfqpoint{1.184153in}{0.898356in}}%
\pgfpathlineto{\pgfqpoint{1.168165in}{0.903937in}}%
\pgfpathlineto{\pgfqpoint{1.152214in}{0.910754in}}%
\pgfpathlineto{\pgfqpoint{1.144397in}{0.896307in}}%
\pgfpathlineto{\pgfqpoint{1.136564in}{0.882163in}}%
\pgfpathclose%
\pgfusepath{fill}%
\end{pgfscope}%
\begin{pgfscope}%
\pgfpathrectangle{\pgfqpoint{0.407592in}{0.222904in}}{\pgfqpoint{1.411163in}{1.411163in}}%
\pgfusepath{clip}%
\pgfsetbuttcap%
\pgfsetroundjoin%
\definecolor{currentfill}{rgb}{0.556863,0.654902,0.847059}%
\pgfsetfillcolor{currentfill}%
\pgfsetfillopacity{0.130000}%
\pgfsetlinewidth{0.000000pt}%
\definecolor{currentstroke}{rgb}{0.000000,0.000000,0.000000}%
\pgfsetstrokecolor{currentstroke}%
\pgfsetdash{}{0pt}%
\pgfpathmoveto{\pgfqpoint{1.088972in}{0.873997in}}%
\pgfpathlineto{\pgfqpoint{1.104899in}{0.863848in}}%
\pgfpathlineto{\pgfqpoint{1.120852in}{0.854983in}}%
\pgfpathlineto{\pgfqpoint{1.128716in}{0.868370in}}%
\pgfpathlineto{\pgfqpoint{1.136564in}{0.882163in}}%
\pgfpathlineto{\pgfqpoint{1.120630in}{0.890580in}}%
\pgfpathlineto{\pgfqpoint{1.104724in}{0.900197in}}%
\pgfpathlineto{\pgfqpoint{1.096857in}{0.886887in}}%
\pgfpathlineto{\pgfqpoint{1.088972in}{0.873997in}}%
\pgfpathclose%
\pgfusepath{fill}%
\end{pgfscope}%
\begin{pgfscope}%
\pgfpathrectangle{\pgfqpoint{0.407592in}{0.222904in}}{\pgfqpoint{1.411163in}{1.411163in}}%
\pgfusepath{clip}%
\pgfsetbuttcap%
\pgfsetroundjoin%
\definecolor{currentfill}{rgb}{0.556863,0.654902,0.847059}%
\pgfsetfillcolor{currentfill}%
\pgfsetfillopacity{0.130000}%
\pgfsetlinewidth{0.000000pt}%
\definecolor{currentstroke}{rgb}{0.000000,0.000000,0.000000}%
\pgfsetstrokecolor{currentstroke}%
\pgfsetdash{}{0pt}%
\pgfpathmoveto{\pgfqpoint{1.294939in}{0.980202in}}%
\pgfpathlineto{\pgfqpoint{1.311130in}{0.979977in}}%
\pgfpathlineto{\pgfqpoint{1.327377in}{0.980663in}}%
\pgfpathlineto{\pgfqpoint{1.335053in}{0.996178in}}%
\pgfpathlineto{\pgfqpoint{1.342717in}{1.011809in}}%
\pgfpathlineto{\pgfqpoint{1.326481in}{1.011393in}}%
\pgfpathlineto{\pgfqpoint{1.310301in}{1.011832in}}%
\pgfpathlineto{\pgfqpoint{1.302626in}{0.995968in}}%
\pgfpathlineto{\pgfqpoint{1.294939in}{0.980202in}}%
\pgfpathclose%
\pgfusepath{fill}%
\end{pgfscope}%
\begin{pgfscope}%
\pgfpathrectangle{\pgfqpoint{0.407592in}{0.222904in}}{\pgfqpoint{1.411163in}{1.411163in}}%
\pgfusepath{clip}%
\pgfsetbuttcap%
\pgfsetroundjoin%
\definecolor{currentfill}{rgb}{0.556863,0.654902,0.847059}%
\pgfsetfillcolor{currentfill}%
\pgfsetfillopacity{0.130000}%
\pgfsetlinewidth{0.000000pt}%
\definecolor{currentstroke}{rgb}{0.000000,0.000000,0.000000}%
\pgfsetstrokecolor{currentstroke}%
\pgfsetdash{}{0pt}%
\pgfpathmoveto{\pgfqpoint{1.184153in}{0.898356in}}%
\pgfpathlineto{\pgfqpoint{1.200181in}{0.894041in}}%
\pgfpathlineto{\pgfqpoint{1.216253in}{0.891004in}}%
\pgfpathlineto{\pgfqpoint{1.224030in}{0.905972in}}%
\pgfpathlineto{\pgfqpoint{1.231793in}{0.921155in}}%
\pgfpathlineto{\pgfqpoint{1.215732in}{0.924229in}}%
\pgfpathlineto{\pgfqpoint{1.199716in}{0.928468in}}%
\pgfpathlineto{\pgfqpoint{1.191941in}{0.913305in}}%
\pgfpathlineto{\pgfqpoint{1.184153in}{0.898356in}}%
\pgfpathclose%
\pgfusepath{fill}%
\end{pgfscope}%
\begin{pgfscope}%
\pgfpathrectangle{\pgfqpoint{0.407592in}{0.222904in}}{\pgfqpoint{1.411163in}{1.411163in}}%
\pgfusepath{clip}%
\pgfsetbuttcap%
\pgfsetroundjoin%
\definecolor{currentfill}{rgb}{0.556863,0.654902,0.847059}%
\pgfsetfillcolor{currentfill}%
\pgfsetfillopacity{0.130000}%
\pgfsetlinewidth{0.000000pt}%
\definecolor{currentstroke}{rgb}{0.000000,0.000000,0.000000}%
\pgfsetstrokecolor{currentstroke}%
\pgfsetdash{}{0pt}%
\pgfpathmoveto{\pgfqpoint{0.850158in}{1.008613in}}%
\pgfpathlineto{\pgfqpoint{0.866026in}{0.989897in}}%
\pgfpathlineto{\pgfqpoint{0.881903in}{0.971372in}}%
\pgfpathlineto{\pgfqpoint{0.890103in}{0.978469in}}%
\pgfpathlineto{\pgfqpoint{0.898273in}{0.986078in}}%
\pgfpathlineto{\pgfqpoint{0.882436in}{1.004305in}}%
\pgfpathlineto{\pgfqpoint{0.866609in}{1.022743in}}%
\pgfpathlineto{\pgfqpoint{0.858399in}{1.015445in}}%
\pgfpathlineto{\pgfqpoint{0.850158in}{1.008613in}}%
\pgfpathclose%
\pgfusepath{fill}%
\end{pgfscope}%
\begin{pgfscope}%
\pgfpathrectangle{\pgfqpoint{0.407592in}{0.222904in}}{\pgfqpoint{1.411163in}{1.411163in}}%
\pgfusepath{clip}%
\pgfsetbuttcap%
\pgfsetroundjoin%
\definecolor{currentfill}{rgb}{0.556863,0.654902,0.847059}%
\pgfsetfillcolor{currentfill}%
\pgfsetfillopacity{0.130000}%
\pgfsetlinewidth{0.000000pt}%
\definecolor{currentstroke}{rgb}{0.000000,0.000000,0.000000}%
\pgfsetstrokecolor{currentstroke}%
\pgfsetdash{}{0pt}%
\pgfpathmoveto{\pgfqpoint{0.706621in}{1.151702in}}%
\pgfpathlineto{\pgfqpoint{0.722461in}{1.132021in}}%
\pgfpathlineto{\pgfqpoint{0.738309in}{1.112396in}}%
\pgfpathlineto{\pgfqpoint{0.746740in}{1.117846in}}%
\pgfpathlineto{\pgfqpoint{0.755138in}{1.123671in}}%
\pgfpathlineto{\pgfqpoint{0.739332in}{1.143161in}}%
\pgfpathlineto{\pgfqpoint{0.723535in}{1.162712in}}%
\pgfpathlineto{\pgfqpoint{0.715095in}{1.157033in}}%
\pgfpathlineto{\pgfqpoint{0.706621in}{1.151702in}}%
\pgfpathclose%
\pgfusepath{fill}%
\end{pgfscope}%
\begin{pgfscope}%
\pgfpathrectangle{\pgfqpoint{0.407592in}{0.222904in}}{\pgfqpoint{1.411163in}{1.411163in}}%
\pgfusepath{clip}%
\pgfsetbuttcap%
\pgfsetroundjoin%
\definecolor{currentfill}{rgb}{0.556863,0.654902,0.847059}%
\pgfsetfillcolor{currentfill}%
\pgfsetfillopacity{0.130000}%
\pgfsetlinewidth{0.000000pt}%
\definecolor{currentstroke}{rgb}{0.000000,0.000000,0.000000}%
\pgfsetstrokecolor{currentstroke}%
\pgfsetdash{}{0pt}%
\pgfpathmoveto{\pgfqpoint{0.961718in}{0.916081in}}%
\pgfpathlineto{\pgfqpoint{0.977606in}{0.899591in}}%
\pgfpathlineto{\pgfqpoint{0.993505in}{0.883658in}}%
\pgfpathlineto{\pgfqpoint{1.001533in}{0.893131in}}%
\pgfpathlineto{\pgfqpoint{1.009537in}{0.903208in}}%
\pgfpathlineto{\pgfqpoint{0.993673in}{0.918594in}}%
\pgfpathlineto{\pgfqpoint{0.977821in}{0.934563in}}%
\pgfpathlineto{\pgfqpoint{0.969783in}{0.925042in}}%
\pgfpathlineto{\pgfqpoint{0.961718in}{0.916081in}}%
\pgfpathclose%
\pgfusepath{fill}%
\end{pgfscope}%
\begin{pgfscope}%
\pgfpathrectangle{\pgfqpoint{0.407592in}{0.222904in}}{\pgfqpoint{1.411163in}{1.411163in}}%
\pgfusepath{clip}%
\pgfsetbuttcap%
\pgfsetroundjoin%
\definecolor{currentfill}{rgb}{0.556863,0.654902,0.847059}%
\pgfsetfillcolor{currentfill}%
\pgfsetfillopacity{0.130000}%
\pgfsetlinewidth{0.000000pt}%
\definecolor{currentstroke}{rgb}{0.000000,0.000000,0.000000}%
\pgfsetstrokecolor{currentstroke}%
\pgfsetdash{}{0pt}%
\pgfpathmoveto{\pgfqpoint{1.041308in}{0.874601in}}%
\pgfpathlineto{\pgfqpoint{1.057218in}{0.861622in}}%
\pgfpathlineto{\pgfqpoint{1.073146in}{0.849711in}}%
\pgfpathlineto{\pgfqpoint{1.081069in}{0.861584in}}%
\pgfpathlineto{\pgfqpoint{1.088972in}{0.873997in}}%
\pgfpathlineto{\pgfqpoint{1.073069in}{0.885309in}}%
\pgfpathlineto{\pgfqpoint{1.057187in}{0.897660in}}%
\pgfpathlineto{\pgfqpoint{1.049258in}{0.885864in}}%
\pgfpathlineto{\pgfqpoint{1.041308in}{0.874601in}}%
\pgfpathclose%
\pgfusepath{fill}%
\end{pgfscope}%
\begin{pgfscope}%
\pgfpathrectangle{\pgfqpoint{0.407592in}{0.222904in}}{\pgfqpoint{1.411163in}{1.411163in}}%
\pgfusepath{clip}%
\pgfsetbuttcap%
\pgfsetroundjoin%
\definecolor{currentfill}{rgb}{0.556863,0.654902,0.847059}%
\pgfsetfillcolor{currentfill}%
\pgfsetfillopacity{0.130000}%
\pgfsetlinewidth{0.000000pt}%
\definecolor{currentstroke}{rgb}{0.000000,0.000000,0.000000}%
\pgfsetstrokecolor{currentstroke}%
\pgfsetdash{}{0pt}%
\pgfpathmoveto{\pgfqpoint{1.231793in}{0.921155in}}%
\pgfpathlineto{\pgfqpoint{1.247902in}{0.919239in}}%
\pgfpathlineto{\pgfqpoint{1.264061in}{0.918456in}}%
\pgfpathlineto{\pgfqpoint{1.271800in}{0.933655in}}%
\pgfpathlineto{\pgfqpoint{1.279526in}{0.949027in}}%
\pgfpathlineto{\pgfqpoint{1.263377in}{0.950018in}}%
\pgfpathlineto{\pgfqpoint{1.247279in}{0.952054in}}%
\pgfpathlineto{\pgfqpoint{1.239543in}{0.936525in}}%
\pgfpathlineto{\pgfqpoint{1.231793in}{0.921155in}}%
\pgfpathclose%
\pgfusepath{fill}%
\end{pgfscope}%
\begin{pgfscope}%
\pgfpathrectangle{\pgfqpoint{0.407592in}{0.222904in}}{\pgfqpoint{1.411163in}{1.411163in}}%
\pgfusepath{clip}%
\pgfsetbuttcap%
\pgfsetroundjoin%
\definecolor{currentfill}{rgb}{0.556863,0.654902,0.847059}%
\pgfsetfillcolor{currentfill}%
\pgfsetfillopacity{0.130000}%
\pgfsetlinewidth{0.000000pt}%
\definecolor{currentstroke}{rgb}{0.000000,0.000000,0.000000}%
\pgfsetstrokecolor{currentstroke}%
\pgfsetdash{}{0pt}%
\pgfpathmoveto{\pgfqpoint{1.405848in}{1.077353in}}%
\pgfpathlineto{\pgfqpoint{1.422236in}{1.079459in}}%
\pgfpathlineto{\pgfqpoint{1.438685in}{1.082223in}}%
\pgfpathlineto{\pgfqpoint{1.446264in}{1.097699in}}%
\pgfpathlineto{\pgfqpoint{1.429821in}{1.095095in}}%
\pgfpathlineto{\pgfqpoint{1.413439in}{1.093134in}}%
\pgfusepath{fill}%
\end{pgfscope}%
\begin{pgfscope}%
\pgfpathrectangle{\pgfqpoint{0.407592in}{0.222904in}}{\pgfqpoint{1.411163in}{1.411163in}}%
\pgfusepath{clip}%
\pgfsetbuttcap%
\pgfsetroundjoin%
\definecolor{currentfill}{rgb}{0.556863,0.654902,0.847059}%
\pgfsetfillcolor{currentfill}%
\pgfsetfillopacity{0.130000}%
\pgfsetlinewidth{0.000000pt}%
\definecolor{currentstroke}{rgb}{0.000000,0.000000,0.000000}%
\pgfsetstrokecolor{currentstroke}%
\pgfsetdash{}{0pt}%
\pgfpathmoveto{\pgfqpoint{0.738309in}{1.112396in}}%
\pgfpathlineto{\pgfqpoint{0.754166in}{1.092835in}}%
\pgfpathlineto{\pgfqpoint{0.770031in}{1.073348in}}%
\pgfpathlineto{\pgfqpoint{0.778420in}{1.078925in}}%
\pgfpathlineto{\pgfqpoint{0.786776in}{1.084910in}}%
\pgfpathlineto{\pgfqpoint{0.770952in}{1.104251in}}%
\pgfpathlineto{\pgfqpoint{0.755138in}{1.123671in}}%
\pgfpathlineto{\pgfqpoint{0.746740in}{1.117846in}}%
\pgfpathlineto{\pgfqpoint{0.738309in}{1.112396in}}%
\pgfpathclose%
\pgfusepath{fill}%
\end{pgfscope}%
\begin{pgfscope}%
\pgfpathrectangle{\pgfqpoint{0.407592in}{0.222904in}}{\pgfqpoint{1.411163in}{1.411163in}}%
\pgfusepath{clip}%
\pgfsetbuttcap%
\pgfsetroundjoin%
\definecolor{currentfill}{rgb}{0.556863,0.654902,0.847059}%
\pgfsetfillcolor{currentfill}%
\pgfsetfillopacity{0.130000}%
\pgfsetlinewidth{0.000000pt}%
\definecolor{currentstroke}{rgb}{0.000000,0.000000,0.000000}%
\pgfsetstrokecolor{currentstroke}%
\pgfsetdash{}{0pt}%
\pgfpathmoveto{\pgfqpoint{1.342717in}{1.011809in}}%
\pgfpathlineto{\pgfqpoint{1.359010in}{1.013048in}}%
\pgfpathlineto{\pgfqpoint{1.375362in}{1.015076in}}%
\pgfpathlineto{\pgfqpoint{1.383002in}{1.030496in}}%
\pgfpathlineto{\pgfqpoint{1.390629in}{1.046024in}}%
\pgfpathlineto{\pgfqpoint{1.374289in}{1.044317in}}%
\pgfpathlineto{\pgfqpoint{1.358008in}{1.043360in}}%
\pgfpathlineto{\pgfqpoint{1.350369in}{1.027540in}}%
\pgfpathlineto{\pgfqpoint{1.342717in}{1.011809in}}%
\pgfpathclose%
\pgfusepath{fill}%
\end{pgfscope}%
\begin{pgfscope}%
\pgfpathrectangle{\pgfqpoint{0.407592in}{0.222904in}}{\pgfqpoint{1.411163in}{1.411163in}}%
\pgfusepath{clip}%
\pgfsetbuttcap%
\pgfsetroundjoin%
\definecolor{currentfill}{rgb}{0.556863,0.654902,0.847059}%
\pgfsetfillcolor{currentfill}%
\pgfsetfillopacity{0.130000}%
\pgfsetlinewidth{0.000000pt}%
\definecolor{currentstroke}{rgb}{0.000000,0.000000,0.000000}%
\pgfsetstrokecolor{currentstroke}%
\pgfsetdash{}{0pt}%
\pgfpathmoveto{\pgfqpoint{0.881903in}{0.971372in}}%
\pgfpathlineto{\pgfqpoint{0.897790in}{0.953072in}}%
\pgfpathlineto{\pgfqpoint{0.913686in}{0.935037in}}%
\pgfpathlineto{\pgfqpoint{0.921845in}{0.942437in}}%
\pgfpathlineto{\pgfqpoint{0.929975in}{0.950400in}}%
\pgfpathlineto{\pgfqpoint{0.914119in}{0.968095in}}%
\pgfpathlineto{\pgfqpoint{0.898273in}{0.986078in}}%
\pgfpathlineto{\pgfqpoint{0.890103in}{0.978469in}}%
\pgfpathlineto{\pgfqpoint{0.881903in}{0.971372in}}%
\pgfpathclose%
\pgfusepath{fill}%
\end{pgfscope}%
\begin{pgfscope}%
\pgfpathrectangle{\pgfqpoint{0.407592in}{0.222904in}}{\pgfqpoint{1.411163in}{1.411163in}}%
\pgfusepath{clip}%
\pgfsetbuttcap%
\pgfsetroundjoin%
\definecolor{currentfill}{rgb}{0.556863,0.654902,0.847059}%
\pgfsetfillcolor{currentfill}%
\pgfsetfillopacity{0.130000}%
\pgfsetlinewidth{0.000000pt}%
\definecolor{currentstroke}{rgb}{0.000000,0.000000,0.000000}%
\pgfsetstrokecolor{currentstroke}%
\pgfsetdash{}{0pt}%
\pgfpathmoveto{\pgfqpoint{0.993505in}{0.883658in}}%
\pgfpathlineto{\pgfqpoint{1.009416in}{0.868386in}}%
\pgfpathlineto{\pgfqpoint{1.025342in}{0.853898in}}%
\pgfpathlineto{\pgfqpoint{1.033336in}{0.863927in}}%
\pgfpathlineto{\pgfqpoint{1.041308in}{0.874601in}}%
\pgfpathlineto{\pgfqpoint{1.025415in}{0.888506in}}%
\pgfpathlineto{\pgfqpoint{1.009537in}{0.903208in}}%
\pgfpathlineto{\pgfqpoint{1.001533in}{0.893131in}}%
\pgfpathlineto{\pgfqpoint{0.993505in}{0.883658in}}%
\pgfpathclose%
\pgfusepath{fill}%
\end{pgfscope}%
\begin{pgfscope}%
\pgfpathrectangle{\pgfqpoint{0.407592in}{0.222904in}}{\pgfqpoint{1.411163in}{1.411163in}}%
\pgfusepath{clip}%
\pgfsetbuttcap%
\pgfsetroundjoin%
\definecolor{currentfill}{rgb}{0.556863,0.654902,0.847059}%
\pgfsetfillcolor{currentfill}%
\pgfsetfillopacity{0.130000}%
\pgfsetlinewidth{0.000000pt}%
\definecolor{currentstroke}{rgb}{0.000000,0.000000,0.000000}%
\pgfsetstrokecolor{currentstroke}%
\pgfsetdash{}{0pt}%
\pgfpathmoveto{\pgfqpoint{0.770031in}{1.073348in}}%
\pgfpathlineto{\pgfqpoint{0.785906in}{1.053946in}}%
\pgfpathlineto{\pgfqpoint{0.801789in}{1.034642in}}%
\pgfpathlineto{\pgfqpoint{0.810135in}{1.040357in}}%
\pgfpathlineto{\pgfqpoint{0.818449in}{1.046516in}}%
\pgfpathlineto{\pgfqpoint{0.802608in}{1.065660in}}%
\pgfpathlineto{\pgfqpoint{0.786776in}{1.084910in}}%
\pgfpathlineto{\pgfqpoint{0.778420in}{1.078925in}}%
\pgfpathlineto{\pgfqpoint{0.770031in}{1.073348in}}%
\pgfpathclose%
\pgfusepath{fill}%
\end{pgfscope}%
\begin{pgfscope}%
\pgfpathrectangle{\pgfqpoint{0.407592in}{0.222904in}}{\pgfqpoint{1.411163in}{1.411163in}}%
\pgfusepath{clip}%
\pgfsetbuttcap%
\pgfsetroundjoin%
\definecolor{currentfill}{rgb}{0.556863,0.654902,0.847059}%
\pgfsetfillcolor{currentfill}%
\pgfsetfillopacity{0.130000}%
\pgfsetlinewidth{0.000000pt}%
\definecolor{currentstroke}{rgb}{0.000000,0.000000,0.000000}%
\pgfsetstrokecolor{currentstroke}%
\pgfsetdash{}{0pt}%
\pgfpathmoveto{\pgfqpoint{1.279526in}{0.949027in}}%
\pgfpathlineto{\pgfqpoint{1.295728in}{0.949049in}}%
\pgfpathlineto{\pgfqpoint{1.311986in}{0.950046in}}%
\pgfpathlineto{\pgfqpoint{1.319688in}{0.965280in}}%
\pgfpathlineto{\pgfqpoint{1.327377in}{0.980663in}}%
\pgfpathlineto{\pgfqpoint{1.311130in}{0.979977in}}%
\pgfpathlineto{\pgfqpoint{1.294939in}{0.980202in}}%
\pgfpathlineto{\pgfqpoint{1.287239in}{0.964549in}}%
\pgfpathlineto{\pgfqpoint{1.279526in}{0.949027in}}%
\pgfpathclose%
\pgfusepath{fill}%
\end{pgfscope}%
\begin{pgfscope}%
\pgfpathrectangle{\pgfqpoint{0.407592in}{0.222904in}}{\pgfqpoint{1.411163in}{1.411163in}}%
\pgfusepath{clip}%
\pgfsetbuttcap%
\pgfsetroundjoin%
\definecolor{currentfill}{rgb}{0.556863,0.654902,0.847059}%
\pgfsetfillcolor{currentfill}%
\pgfsetfillopacity{0.130000}%
\pgfsetlinewidth{0.000000pt}%
\definecolor{currentstroke}{rgb}{0.000000,0.000000,0.000000}%
\pgfsetstrokecolor{currentstroke}%
\pgfsetdash{}{0pt}%
\pgfpathmoveto{\pgfqpoint{1.120852in}{0.854983in}}%
\pgfpathlineto{\pgfqpoint{1.136835in}{0.847507in}}%
\pgfpathlineto{\pgfqpoint{1.152852in}{0.841507in}}%
\pgfpathlineto{\pgfqpoint{1.160700in}{0.855179in}}%
\pgfpathlineto{\pgfqpoint{1.168532in}{0.869248in}}%
\pgfpathlineto{\pgfqpoint{1.152530in}{0.875031in}}%
\pgfpathlineto{\pgfqpoint{1.136564in}{0.882163in}}%
\pgfpathlineto{\pgfqpoint{1.128716in}{0.868370in}}%
\pgfpathlineto{\pgfqpoint{1.120852in}{0.854983in}}%
\pgfpathclose%
\pgfusepath{fill}%
\end{pgfscope}%
\begin{pgfscope}%
\pgfpathrectangle{\pgfqpoint{0.407592in}{0.222904in}}{\pgfqpoint{1.411163in}{1.411163in}}%
\pgfusepath{clip}%
\pgfsetbuttcap%
\pgfsetroundjoin%
\definecolor{currentfill}{rgb}{0.556863,0.654902,0.847059}%
\pgfsetfillcolor{currentfill}%
\pgfsetfillopacity{0.130000}%
\pgfsetlinewidth{0.000000pt}%
\definecolor{currentstroke}{rgb}{0.000000,0.000000,0.000000}%
\pgfsetstrokecolor{currentstroke}%
\pgfsetdash{}{0pt}%
\pgfpathmoveto{\pgfqpoint{0.913686in}{0.935037in}}%
\pgfpathlineto{\pgfqpoint{0.929591in}{0.917316in}}%
\pgfpathlineto{\pgfqpoint{0.945507in}{0.899967in}}%
\pgfpathlineto{\pgfqpoint{0.953626in}{0.907713in}}%
\pgfpathlineto{\pgfqpoint{0.961718in}{0.916081in}}%
\pgfpathlineto{\pgfqpoint{0.945841in}{0.933042in}}%
\pgfpathlineto{\pgfqpoint{0.929975in}{0.950400in}}%
\pgfpathlineto{\pgfqpoint{0.921845in}{0.942437in}}%
\pgfpathlineto{\pgfqpoint{0.913686in}{0.935037in}}%
\pgfpathclose%
\pgfusepath{fill}%
\end{pgfscope}%
\begin{pgfscope}%
\pgfpathrectangle{\pgfqpoint{0.407592in}{0.222904in}}{\pgfqpoint{1.411163in}{1.411163in}}%
\pgfusepath{clip}%
\pgfsetbuttcap%
\pgfsetroundjoin%
\definecolor{currentfill}{rgb}{0.556863,0.654902,0.847059}%
\pgfsetfillcolor{currentfill}%
\pgfsetfillopacity{0.130000}%
\pgfsetlinewidth{0.000000pt}%
\definecolor{currentstroke}{rgb}{0.000000,0.000000,0.000000}%
\pgfsetstrokecolor{currentstroke}%
\pgfsetdash{}{0pt}%
\pgfpathmoveto{\pgfqpoint{1.073146in}{0.849711in}}%
\pgfpathlineto{\pgfqpoint{1.089096in}{0.839014in}}%
\pgfpathlineto{\pgfqpoint{1.105070in}{0.829681in}}%
\pgfpathlineto{\pgfqpoint{1.112970in}{0.842063in}}%
\pgfpathlineto{\pgfqpoint{1.120852in}{0.854983in}}%
\pgfpathlineto{\pgfqpoint{1.104899in}{0.863848in}}%
\pgfpathlineto{\pgfqpoint{1.088972in}{0.873997in}}%
\pgfpathlineto{\pgfqpoint{1.081069in}{0.861584in}}%
\pgfpathlineto{\pgfqpoint{1.073146in}{0.849711in}}%
\pgfpathclose%
\pgfusepath{fill}%
\end{pgfscope}%
\begin{pgfscope}%
\pgfpathrectangle{\pgfqpoint{0.407592in}{0.222904in}}{\pgfqpoint{1.411163in}{1.411163in}}%
\pgfusepath{clip}%
\pgfsetbuttcap%
\pgfsetroundjoin%
\definecolor{currentfill}{rgb}{0.556863,0.654902,0.847059}%
\pgfsetfillcolor{currentfill}%
\pgfsetfillopacity{0.130000}%
\pgfsetlinewidth{0.000000pt}%
\definecolor{currentstroke}{rgb}{0.000000,0.000000,0.000000}%
\pgfsetstrokecolor{currentstroke}%
\pgfsetdash{}{0pt}%
\pgfpathmoveto{\pgfqpoint{1.168532in}{0.869248in}}%
\pgfpathlineto{\pgfqpoint{1.184574in}{0.864856in}}%
\pgfpathlineto{\pgfqpoint{1.200658in}{0.861867in}}%
\pgfpathlineto{\pgfqpoint{1.208463in}{0.876289in}}%
\pgfpathlineto{\pgfqpoint{1.216253in}{0.891004in}}%
\pgfpathlineto{\pgfqpoint{1.200181in}{0.894041in}}%
\pgfpathlineto{\pgfqpoint{1.184153in}{0.898356in}}%
\pgfpathlineto{\pgfqpoint{1.176350in}{0.883656in}}%
\pgfpathlineto{\pgfqpoint{1.168532in}{0.869248in}}%
\pgfpathclose%
\pgfusepath{fill}%
\end{pgfscope}%
\begin{pgfscope}%
\pgfpathrectangle{\pgfqpoint{0.407592in}{0.222904in}}{\pgfqpoint{1.411163in}{1.411163in}}%
\pgfusepath{clip}%
\pgfsetbuttcap%
\pgfsetroundjoin%
\definecolor{currentfill}{rgb}{0.556863,0.654902,0.847059}%
\pgfsetfillcolor{currentfill}%
\pgfsetfillopacity{0.130000}%
\pgfsetlinewidth{0.000000pt}%
\definecolor{currentstroke}{rgb}{0.000000,0.000000,0.000000}%
\pgfsetstrokecolor{currentstroke}%
\pgfsetdash{}{0pt}%
\pgfpathmoveto{\pgfqpoint{1.390629in}{1.046024in}}%
\pgfpathlineto{\pgfqpoint{1.407029in}{1.048448in}}%
\pgfpathlineto{\pgfqpoint{1.423491in}{1.051556in}}%
\pgfpathlineto{\pgfqpoint{1.431094in}{1.066839in}}%
\pgfpathlineto{\pgfqpoint{1.438685in}{1.082223in}}%
\pgfpathlineto{\pgfqpoint{1.422236in}{1.079459in}}%
\pgfpathlineto{\pgfqpoint{1.405848in}{1.077353in}}%
\pgfpathlineto{\pgfqpoint{1.398245in}{1.061646in}}%
\pgfpathlineto{\pgfqpoint{1.390629in}{1.046024in}}%
\pgfpathclose%
\pgfusepath{fill}%
\end{pgfscope}%
\begin{pgfscope}%
\pgfpathrectangle{\pgfqpoint{0.407592in}{0.222904in}}{\pgfqpoint{1.411163in}{1.411163in}}%
\pgfusepath{clip}%
\pgfsetbuttcap%
\pgfsetroundjoin%
\definecolor{currentfill}{rgb}{0.556863,0.654902,0.847059}%
\pgfsetfillcolor{currentfill}%
\pgfsetfillopacity{0.130000}%
\pgfsetlinewidth{0.000000pt}%
\definecolor{currentstroke}{rgb}{0.000000,0.000000,0.000000}%
\pgfsetstrokecolor{currentstroke}%
\pgfsetdash{}{0pt}%
\pgfpathmoveto{\pgfqpoint{0.801789in}{1.034642in}}%
\pgfpathlineto{\pgfqpoint{0.817681in}{1.015451in}}%
\pgfpathlineto{\pgfqpoint{0.833582in}{0.996393in}}%
\pgfpathlineto{\pgfqpoint{0.841886in}{1.002260in}}%
\pgfpathlineto{\pgfqpoint{0.850158in}{1.008613in}}%
\pgfpathlineto{\pgfqpoint{0.834299in}{1.027494in}}%
\pgfpathlineto{\pgfqpoint{0.818449in}{1.046516in}}%
\pgfpathlineto{\pgfqpoint{0.810135in}{1.040357in}}%
\pgfpathlineto{\pgfqpoint{0.801789in}{1.034642in}}%
\pgfpathclose%
\pgfusepath{fill}%
\end{pgfscope}%
\begin{pgfscope}%
\pgfpathrectangle{\pgfqpoint{0.407592in}{0.222904in}}{\pgfqpoint{1.411163in}{1.411163in}}%
\pgfusepath{clip}%
\pgfsetbuttcap%
\pgfsetroundjoin%
\definecolor{currentfill}{rgb}{0.556863,0.654902,0.847059}%
\pgfsetfillcolor{currentfill}%
\pgfsetfillopacity{0.130000}%
\pgfsetlinewidth{0.000000pt}%
\definecolor{currentstroke}{rgb}{0.000000,0.000000,0.000000}%
\pgfsetstrokecolor{currentstroke}%
\pgfsetdash{}{0pt}%
\pgfpathmoveto{\pgfqpoint{0.657834in}{1.181752in}}%
\pgfpathlineto{\pgfqpoint{0.673699in}{1.161904in}}%
\pgfpathlineto{\pgfqpoint{0.689572in}{1.142098in}}%
\pgfpathlineto{\pgfqpoint{0.698114in}{1.146723in}}%
\pgfpathlineto{\pgfqpoint{0.706621in}{1.151702in}}%
\pgfpathlineto{\pgfqpoint{0.690791in}{1.171432in}}%
\pgfpathlineto{\pgfqpoint{0.674969in}{1.191205in}}%
\pgfpathlineto{\pgfqpoint{0.666419in}{1.186314in}}%
\pgfpathlineto{\pgfqpoint{0.657834in}{1.181752in}}%
\pgfpathclose%
\pgfusepath{fill}%
\end{pgfscope}%
\begin{pgfscope}%
\pgfpathrectangle{\pgfqpoint{0.407592in}{0.222904in}}{\pgfqpoint{1.411163in}{1.411163in}}%
\pgfusepath{clip}%
\pgfsetbuttcap%
\pgfsetroundjoin%
\definecolor{currentfill}{rgb}{0.556863,0.654902,0.847059}%
\pgfsetfillcolor{currentfill}%
\pgfsetfillopacity{0.130000}%
\pgfsetlinewidth{0.000000pt}%
\definecolor{currentstroke}{rgb}{0.000000,0.000000,0.000000}%
\pgfsetstrokecolor{currentstroke}%
\pgfsetdash{}{0pt}%
\pgfpathmoveto{\pgfqpoint{1.025342in}{0.853898in}}%
\pgfpathlineto{\pgfqpoint{1.041282in}{0.840334in}}%
\pgfpathlineto{\pgfqpoint{1.057239in}{0.827855in}}%
\pgfpathlineto{\pgfqpoint{1.065203in}{0.838444in}}%
\pgfpathlineto{\pgfqpoint{1.073146in}{0.849711in}}%
\pgfpathlineto{\pgfqpoint{1.057218in}{0.861622in}}%
\pgfpathlineto{\pgfqpoint{1.041308in}{0.874601in}}%
\pgfpathlineto{\pgfqpoint{1.033336in}{0.863927in}}%
\pgfpathlineto{\pgfqpoint{1.025342in}{0.853898in}}%
\pgfpathclose%
\pgfusepath{fill}%
\end{pgfscope}%
\begin{pgfscope}%
\pgfpathrectangle{\pgfqpoint{0.407592in}{0.222904in}}{\pgfqpoint{1.411163in}{1.411163in}}%
\pgfusepath{clip}%
\pgfsetbuttcap%
\pgfsetroundjoin%
\definecolor{currentfill}{rgb}{0.556863,0.654902,0.847059}%
\pgfsetfillcolor{currentfill}%
\pgfsetfillopacity{0.130000}%
\pgfsetlinewidth{0.000000pt}%
\definecolor{currentstroke}{rgb}{0.000000,0.000000,0.000000}%
\pgfsetstrokecolor{currentstroke}%
\pgfsetdash{}{0pt}%
\pgfpathmoveto{\pgfqpoint{1.216253in}{0.891004in}}%
\pgfpathlineto{\pgfqpoint{1.232373in}{0.889232in}}%
\pgfpathlineto{\pgfqpoint{1.248544in}{0.888691in}}%
\pgfpathlineto{\pgfqpoint{1.256310in}{0.903458in}}%
\pgfpathlineto{\pgfqpoint{1.264061in}{0.918456in}}%
\pgfpathlineto{\pgfqpoint{1.247902in}{0.919239in}}%
\pgfpathlineto{\pgfqpoint{1.231793in}{0.921155in}}%
\pgfpathlineto{\pgfqpoint{1.224030in}{0.905972in}}%
\pgfpathlineto{\pgfqpoint{1.216253in}{0.891004in}}%
\pgfpathclose%
\pgfusepath{fill}%
\end{pgfscope}%
\begin{pgfscope}%
\pgfpathrectangle{\pgfqpoint{0.407592in}{0.222904in}}{\pgfqpoint{1.411163in}{1.411163in}}%
\pgfusepath{clip}%
\pgfsetbuttcap%
\pgfsetroundjoin%
\definecolor{currentfill}{rgb}{0.556863,0.654902,0.847059}%
\pgfsetfillcolor{currentfill}%
\pgfsetfillopacity{0.130000}%
\pgfsetlinewidth{0.000000pt}%
\definecolor{currentstroke}{rgb}{0.000000,0.000000,0.000000}%
\pgfsetstrokecolor{currentstroke}%
\pgfsetdash{}{0pt}%
\pgfpathmoveto{\pgfqpoint{0.945507in}{0.899967in}}%
\pgfpathlineto{\pgfqpoint{0.961432in}{0.883064in}}%
\pgfpathlineto{\pgfqpoint{0.977369in}{0.866693in}}%
\pgfpathlineto{\pgfqpoint{0.985450in}{0.874832in}}%
\pgfpathlineto{\pgfqpoint{0.993505in}{0.883658in}}%
\pgfpathlineto{\pgfqpoint{0.977606in}{0.899591in}}%
\pgfpathlineto{\pgfqpoint{0.961718in}{0.916081in}}%
\pgfpathlineto{\pgfqpoint{0.953626in}{0.907713in}}%
\pgfpathlineto{\pgfqpoint{0.945507in}{0.899967in}}%
\pgfpathclose%
\pgfusepath{fill}%
\end{pgfscope}%
\begin{pgfscope}%
\pgfpathrectangle{\pgfqpoint{0.407592in}{0.222904in}}{\pgfqpoint{1.411163in}{1.411163in}}%
\pgfusepath{clip}%
\pgfsetbuttcap%
\pgfsetroundjoin%
\definecolor{currentfill}{rgb}{0.556863,0.654902,0.847059}%
\pgfsetfillcolor{currentfill}%
\pgfsetfillopacity{0.130000}%
\pgfsetlinewidth{0.000000pt}%
\definecolor{currentstroke}{rgb}{0.000000,0.000000,0.000000}%
\pgfsetstrokecolor{currentstroke}%
\pgfsetdash{}{0pt}%
\pgfpathmoveto{\pgfqpoint{1.327377in}{0.980663in}}%
\pgfpathlineto{\pgfqpoint{1.343681in}{0.982223in}}%
\pgfpathlineto{\pgfqpoint{1.360045in}{0.984614in}}%
\pgfpathlineto{\pgfqpoint{1.367710in}{0.999777in}}%
\pgfpathlineto{\pgfqpoint{1.375362in}{1.015076in}}%
\pgfpathlineto{\pgfqpoint{1.359010in}{1.013048in}}%
\pgfpathlineto{\pgfqpoint{1.342717in}{1.011809in}}%
\pgfpathlineto{\pgfqpoint{1.335053in}{0.996178in}}%
\pgfpathlineto{\pgfqpoint{1.327377in}{0.980663in}}%
\pgfpathclose%
\pgfusepath{fill}%
\end{pgfscope}%
\begin{pgfscope}%
\pgfpathrectangle{\pgfqpoint{0.407592in}{0.222904in}}{\pgfqpoint{1.411163in}{1.411163in}}%
\pgfusepath{clip}%
\pgfsetbuttcap%
\pgfsetroundjoin%
\definecolor{currentfill}{rgb}{0.556863,0.654902,0.847059}%
\pgfsetfillcolor{currentfill}%
\pgfsetfillopacity{0.130000}%
\pgfsetlinewidth{0.000000pt}%
\definecolor{currentstroke}{rgb}{0.000000,0.000000,0.000000}%
\pgfsetstrokecolor{currentstroke}%
\pgfsetdash{}{0pt}%
\pgfpathmoveto{\pgfqpoint{0.833582in}{0.996393in}}%
\pgfpathlineto{\pgfqpoint{0.849492in}{0.977490in}}%
\pgfpathlineto{\pgfqpoint{0.865410in}{0.958767in}}%
\pgfpathlineto{\pgfqpoint{0.873673in}{0.964800in}}%
\pgfpathlineto{\pgfqpoint{0.881903in}{0.971372in}}%
\pgfpathlineto{\pgfqpoint{0.866026in}{0.989897in}}%
\pgfpathlineto{\pgfqpoint{0.850158in}{1.008613in}}%
\pgfpathlineto{\pgfqpoint{0.841886in}{1.002260in}}%
\pgfpathlineto{\pgfqpoint{0.833582in}{0.996393in}}%
\pgfpathclose%
\pgfusepath{fill}%
\end{pgfscope}%
\begin{pgfscope}%
\pgfpathrectangle{\pgfqpoint{0.407592in}{0.222904in}}{\pgfqpoint{1.411163in}{1.411163in}}%
\pgfusepath{clip}%
\pgfsetbuttcap%
\pgfsetroundjoin%
\definecolor{currentfill}{rgb}{0.556863,0.654902,0.847059}%
\pgfsetfillcolor{currentfill}%
\pgfsetfillopacity{0.130000}%
\pgfsetlinewidth{0.000000pt}%
\definecolor{currentstroke}{rgb}{0.000000,0.000000,0.000000}%
\pgfsetstrokecolor{currentstroke}%
\pgfsetdash{}{0pt}%
\pgfpathmoveto{\pgfqpoint{0.689572in}{1.142098in}}%
\pgfpathlineto{\pgfqpoint{0.705454in}{1.122340in}}%
\pgfpathlineto{\pgfqpoint{0.721345in}{1.102638in}}%
\pgfpathlineto{\pgfqpoint{0.729844in}{1.107326in}}%
\pgfpathlineto{\pgfqpoint{0.738309in}{1.112396in}}%
\pgfpathlineto{\pgfqpoint{0.722461in}{1.132021in}}%
\pgfpathlineto{\pgfqpoint{0.706621in}{1.151702in}}%
\pgfpathlineto{\pgfqpoint{0.698114in}{1.146723in}}%
\pgfpathlineto{\pgfqpoint{0.689572in}{1.142098in}}%
\pgfpathclose%
\pgfusepath{fill}%
\end{pgfscope}%
\begin{pgfscope}%
\pgfpathrectangle{\pgfqpoint{0.407592in}{0.222904in}}{\pgfqpoint{1.411163in}{1.411163in}}%
\pgfusepath{clip}%
\pgfsetbuttcap%
\pgfsetroundjoin%
\definecolor{currentfill}{rgb}{0.556863,0.654902,0.847059}%
\pgfsetfillcolor{currentfill}%
\pgfsetfillopacity{0.130000}%
\pgfsetlinewidth{0.000000pt}%
\definecolor{currentstroke}{rgb}{0.000000,0.000000,0.000000}%
\pgfsetstrokecolor{currentstroke}%
\pgfsetdash{}{0pt}%
\pgfpathmoveto{\pgfqpoint{1.264061in}{0.918456in}}%
\pgfpathlineto{\pgfqpoint{1.280275in}{0.918765in}}%
\pgfpathlineto{\pgfqpoint{1.296545in}{0.920117in}}%
\pgfpathlineto{\pgfqpoint{1.304272in}{0.934984in}}%
\pgfpathlineto{\pgfqpoint{1.311986in}{0.950046in}}%
\pgfpathlineto{\pgfqpoint{1.295728in}{0.949049in}}%
\pgfpathlineto{\pgfqpoint{1.279526in}{0.949027in}}%
\pgfpathlineto{\pgfqpoint{1.271800in}{0.933655in}}%
\pgfpathlineto{\pgfqpoint{1.264061in}{0.918456in}}%
\pgfpathclose%
\pgfusepath{fill}%
\end{pgfscope}%
\begin{pgfscope}%
\pgfpathrectangle{\pgfqpoint{0.407592in}{0.222904in}}{\pgfqpoint{1.411163in}{1.411163in}}%
\pgfusepath{clip}%
\pgfsetbuttcap%
\pgfsetroundjoin%
\definecolor{currentfill}{rgb}{0.556863,0.654902,0.847059}%
\pgfsetfillcolor{currentfill}%
\pgfsetfillopacity{0.130000}%
\pgfsetlinewidth{0.000000pt}%
\definecolor{currentstroke}{rgb}{0.000000,0.000000,0.000000}%
\pgfsetstrokecolor{currentstroke}%
\pgfsetdash{}{0pt}%
\pgfpathmoveto{\pgfqpoint{0.865410in}{0.958767in}}%
\pgfpathlineto{\pgfqpoint{0.881338in}{0.940256in}}%
\pgfpathlineto{\pgfqpoint{0.897275in}{0.921998in}}%
\pgfpathlineto{\pgfqpoint{0.905496in}{0.928219in}}%
\pgfpathlineto{\pgfqpoint{0.913686in}{0.935037in}}%
\pgfpathlineto{\pgfqpoint{0.897790in}{0.953072in}}%
\pgfpathlineto{\pgfqpoint{0.881903in}{0.971372in}}%
\pgfpathlineto{\pgfqpoint{0.873673in}{0.964800in}}%
\pgfpathlineto{\pgfqpoint{0.865410in}{0.958767in}}%
\pgfpathclose%
\pgfusepath{fill}%
\end{pgfscope}%
\begin{pgfscope}%
\pgfpathrectangle{\pgfqpoint{0.407592in}{0.222904in}}{\pgfqpoint{1.411163in}{1.411163in}}%
\pgfusepath{clip}%
\pgfsetbuttcap%
\pgfsetroundjoin%
\definecolor{currentfill}{rgb}{0.556863,0.654902,0.847059}%
\pgfsetfillcolor{currentfill}%
\pgfsetfillopacity{0.130000}%
\pgfsetlinewidth{0.000000pt}%
\definecolor{currentstroke}{rgb}{0.000000,0.000000,0.000000}%
\pgfsetstrokecolor{currentstroke}%
\pgfsetdash{}{0pt}%
\pgfpathmoveto{\pgfqpoint{0.721345in}{1.102638in}}%
\pgfpathlineto{\pgfqpoint{0.737244in}{1.082998in}}%
\pgfpathlineto{\pgfqpoint{0.753153in}{1.063432in}}%
\pgfpathlineto{\pgfqpoint{0.761609in}{1.068182in}}%
\pgfpathlineto{\pgfqpoint{0.770031in}{1.073348in}}%
\pgfpathlineto{\pgfqpoint{0.754166in}{1.092835in}}%
\pgfpathlineto{\pgfqpoint{0.738309in}{1.112396in}}%
\pgfpathlineto{\pgfqpoint{0.729844in}{1.107326in}}%
\pgfpathlineto{\pgfqpoint{0.721345in}{1.102638in}}%
\pgfpathclose%
\pgfusepath{fill}%
\end{pgfscope}%
\begin{pgfscope}%
\pgfpathrectangle{\pgfqpoint{0.407592in}{0.222904in}}{\pgfqpoint{1.411163in}{1.411163in}}%
\pgfusepath{clip}%
\pgfsetbuttcap%
\pgfsetroundjoin%
\definecolor{currentfill}{rgb}{0.556863,0.654902,0.847059}%
\pgfsetfillcolor{currentfill}%
\pgfsetfillopacity{0.130000}%
\pgfsetlinewidth{0.000000pt}%
\definecolor{currentstroke}{rgb}{0.000000,0.000000,0.000000}%
\pgfsetstrokecolor{currentstroke}%
\pgfsetdash{}{0pt}%
\pgfpathmoveto{\pgfqpoint{1.105070in}{0.829681in}}%
\pgfpathlineto{\pgfqpoint{1.121073in}{0.821845in}}%
\pgfpathlineto{\pgfqpoint{1.137108in}{0.815616in}}%
\pgfpathlineto{\pgfqpoint{1.144988in}{0.828294in}}%
\pgfpathlineto{\pgfqpoint{1.152852in}{0.841507in}}%
\pgfpathlineto{\pgfqpoint{1.136835in}{0.847507in}}%
\pgfpathlineto{\pgfqpoint{1.120852in}{0.854983in}}%
\pgfpathlineto{\pgfqpoint{1.112970in}{0.842063in}}%
\pgfpathlineto{\pgfqpoint{1.105070in}{0.829681in}}%
\pgfpathclose%
\pgfusepath{fill}%
\end{pgfscope}%
\begin{pgfscope}%
\pgfpathrectangle{\pgfqpoint{0.407592in}{0.222904in}}{\pgfqpoint{1.411163in}{1.411163in}}%
\pgfusepath{clip}%
\pgfsetbuttcap%
\pgfsetroundjoin%
\definecolor{currentfill}{rgb}{0.556863,0.654902,0.847059}%
\pgfsetfillcolor{currentfill}%
\pgfsetfillopacity{0.130000}%
\pgfsetlinewidth{0.000000pt}%
\definecolor{currentstroke}{rgb}{0.000000,0.000000,0.000000}%
\pgfsetstrokecolor{currentstroke}%
\pgfsetdash{}{0pt}%
\pgfpathmoveto{\pgfqpoint{0.977369in}{0.866693in}}%
\pgfpathlineto{\pgfqpoint{0.993317in}{0.850961in}}%
\pgfpathlineto{\pgfqpoint{1.009277in}{0.835997in}}%
\pgfpathlineto{\pgfqpoint{1.017322in}{0.844570in}}%
\pgfpathlineto{\pgfqpoint{1.025342in}{0.853898in}}%
\pgfpathlineto{\pgfqpoint{1.009416in}{0.868386in}}%
\pgfpathlineto{\pgfqpoint{0.993505in}{0.883658in}}%
\pgfpathlineto{\pgfqpoint{0.985450in}{0.874832in}}%
\pgfpathlineto{\pgfqpoint{0.977369in}{0.866693in}}%
\pgfpathclose%
\pgfusepath{fill}%
\end{pgfscope}%
\begin{pgfscope}%
\pgfpathrectangle{\pgfqpoint{0.407592in}{0.222904in}}{\pgfqpoint{1.411163in}{1.411163in}}%
\pgfusepath{clip}%
\pgfsetbuttcap%
\pgfsetroundjoin%
\definecolor{currentfill}{rgb}{0.556863,0.654902,0.847059}%
\pgfsetfillcolor{currentfill}%
\pgfsetfillopacity{0.130000}%
\pgfsetlinewidth{0.000000pt}%
\definecolor{currentstroke}{rgb}{0.000000,0.000000,0.000000}%
\pgfsetstrokecolor{currentstroke}%
\pgfsetdash{}{0pt}%
\pgfpathmoveto{\pgfqpoint{1.438685in}{1.082223in}}%
\pgfpathlineto{\pgfqpoint{1.455197in}{1.085615in}}%
\pgfpathlineto{\pgfqpoint{1.471773in}{1.089605in}}%
\pgfpathlineto{\pgfqpoint{1.479338in}{1.104731in}}%
\pgfpathlineto{\pgfqpoint{1.462769in}{1.100921in}}%
\pgfpathlineto{\pgfqpoint{1.446264in}{1.097699in}}%
\pgfusepath{fill}%
\end{pgfscope}%
\begin{pgfscope}%
\pgfpathrectangle{\pgfqpoint{0.407592in}{0.222904in}}{\pgfqpoint{1.411163in}{1.411163in}}%
\pgfusepath{clip}%
\pgfsetbuttcap%
\pgfsetroundjoin%
\definecolor{currentfill}{rgb}{0.556863,0.654902,0.847059}%
\pgfsetfillcolor{currentfill}%
\pgfsetfillopacity{0.130000}%
\pgfsetlinewidth{0.000000pt}%
\definecolor{currentstroke}{rgb}{0.000000,0.000000,0.000000}%
\pgfsetstrokecolor{currentstroke}%
\pgfsetdash{}{0pt}%
\pgfpathmoveto{\pgfqpoint{1.375362in}{1.015076in}}%
\pgfpathlineto{\pgfqpoint{1.391774in}{1.017856in}}%
\pgfpathlineto{\pgfqpoint{1.408250in}{1.021348in}}%
\pgfpathlineto{\pgfqpoint{1.415877in}{1.036388in}}%
\pgfpathlineto{\pgfqpoint{1.423491in}{1.051556in}}%
\pgfpathlineto{\pgfqpoint{1.407029in}{1.048448in}}%
\pgfpathlineto{\pgfqpoint{1.390629in}{1.046024in}}%
\pgfpathlineto{\pgfqpoint{1.383002in}{1.030496in}}%
\pgfpathlineto{\pgfqpoint{1.375362in}{1.015076in}}%
\pgfpathclose%
\pgfusepath{fill}%
\end{pgfscope}%
\begin{pgfscope}%
\pgfpathrectangle{\pgfqpoint{0.407592in}{0.222904in}}{\pgfqpoint{1.411163in}{1.411163in}}%
\pgfusepath{clip}%
\pgfsetbuttcap%
\pgfsetroundjoin%
\definecolor{currentfill}{rgb}{0.556863,0.654902,0.847059}%
\pgfsetfillcolor{currentfill}%
\pgfsetfillopacity{0.130000}%
\pgfsetlinewidth{0.000000pt}%
\definecolor{currentstroke}{rgb}{0.000000,0.000000,0.000000}%
\pgfsetstrokecolor{currentstroke}%
\pgfsetdash{}{0pt}%
\pgfpathmoveto{\pgfqpoint{1.152852in}{0.841507in}}%
\pgfpathlineto{\pgfqpoint{1.168908in}{0.837034in}}%
\pgfpathlineto{\pgfqpoint{1.185006in}{0.834104in}}%
\pgfpathlineto{\pgfqpoint{1.192840in}{0.847786in}}%
\pgfpathlineto{\pgfqpoint{1.200658in}{0.861867in}}%
\pgfpathlineto{\pgfqpoint{1.184574in}{0.864856in}}%
\pgfpathlineto{\pgfqpoint{1.168532in}{0.869248in}}%
\pgfpathlineto{\pgfqpoint{1.160700in}{0.855179in}}%
\pgfpathlineto{\pgfqpoint{1.152852in}{0.841507in}}%
\pgfpathclose%
\pgfusepath{fill}%
\end{pgfscope}%
\begin{pgfscope}%
\pgfpathrectangle{\pgfqpoint{0.407592in}{0.222904in}}{\pgfqpoint{1.411163in}{1.411163in}}%
\pgfusepath{clip}%
\pgfsetbuttcap%
\pgfsetroundjoin%
\definecolor{currentfill}{rgb}{0.556863,0.654902,0.847059}%
\pgfsetfillcolor{currentfill}%
\pgfsetfillopacity{0.130000}%
\pgfsetlinewidth{0.000000pt}%
\definecolor{currentstroke}{rgb}{0.000000,0.000000,0.000000}%
\pgfsetstrokecolor{currentstroke}%
\pgfsetdash{}{0pt}%
\pgfpathmoveto{\pgfqpoint{1.057239in}{0.827855in}}%
\pgfpathlineto{\pgfqpoint{1.073215in}{0.816632in}}%
\pgfpathlineto{\pgfqpoint{1.089214in}{0.806844in}}%
\pgfpathlineto{\pgfqpoint{1.097152in}{0.817914in}}%
\pgfpathlineto{\pgfqpoint{1.105070in}{0.829681in}}%
\pgfpathlineto{\pgfqpoint{1.089096in}{0.839014in}}%
\pgfpathlineto{\pgfqpoint{1.073146in}{0.849711in}}%
\pgfpathlineto{\pgfqpoint{1.065203in}{0.838444in}}%
\pgfpathlineto{\pgfqpoint{1.057239in}{0.827855in}}%
\pgfpathclose%
\pgfusepath{fill}%
\end{pgfscope}%
\begin{pgfscope}%
\pgfpathrectangle{\pgfqpoint{0.407592in}{0.222904in}}{\pgfqpoint{1.411163in}{1.411163in}}%
\pgfusepath{clip}%
\pgfsetbuttcap%
\pgfsetroundjoin%
\definecolor{currentfill}{rgb}{0.556863,0.654902,0.847059}%
\pgfsetfillcolor{currentfill}%
\pgfsetfillopacity{0.130000}%
\pgfsetlinewidth{0.000000pt}%
\definecolor{currentstroke}{rgb}{0.000000,0.000000,0.000000}%
\pgfsetstrokecolor{currentstroke}%
\pgfsetdash{}{0pt}%
\pgfpathmoveto{\pgfqpoint{0.753153in}{1.063432in}}%
\pgfpathlineto{\pgfqpoint{0.769070in}{1.043949in}}%
\pgfpathlineto{\pgfqpoint{0.784995in}{1.024562in}}%
\pgfpathlineto{\pgfqpoint{0.793409in}{1.029376in}}%
\pgfpathlineto{\pgfqpoint{0.801789in}{1.034642in}}%
\pgfpathlineto{\pgfqpoint{0.785906in}{1.053946in}}%
\pgfpathlineto{\pgfqpoint{0.770031in}{1.073348in}}%
\pgfpathlineto{\pgfqpoint{0.761609in}{1.068182in}}%
\pgfpathlineto{\pgfqpoint{0.753153in}{1.063432in}}%
\pgfpathclose%
\pgfusepath{fill}%
\end{pgfscope}%
\begin{pgfscope}%
\pgfpathrectangle{\pgfqpoint{0.407592in}{0.222904in}}{\pgfqpoint{1.411163in}{1.411163in}}%
\pgfusepath{clip}%
\pgfsetbuttcap%
\pgfsetroundjoin%
\definecolor{currentfill}{rgb}{0.556863,0.654902,0.847059}%
\pgfsetfillcolor{currentfill}%
\pgfsetfillopacity{0.130000}%
\pgfsetlinewidth{0.000000pt}%
\definecolor{currentstroke}{rgb}{0.000000,0.000000,0.000000}%
\pgfsetstrokecolor{currentstroke}%
\pgfsetdash{}{0pt}%
\pgfpathmoveto{\pgfqpoint{0.897275in}{0.921998in}}%
\pgfpathlineto{\pgfqpoint{0.913221in}{0.904039in}}%
\pgfpathlineto{\pgfqpoint{0.929177in}{0.886437in}}%
\pgfpathlineto{\pgfqpoint{0.937357in}{0.892868in}}%
\pgfpathlineto{\pgfqpoint{0.945507in}{0.899967in}}%
\pgfpathlineto{\pgfqpoint{0.929591in}{0.917316in}}%
\pgfpathlineto{\pgfqpoint{0.913686in}{0.935037in}}%
\pgfpathlineto{\pgfqpoint{0.905496in}{0.928219in}}%
\pgfpathlineto{\pgfqpoint{0.897275in}{0.921998in}}%
\pgfpathclose%
\pgfusepath{fill}%
\end{pgfscope}%
\begin{pgfscope}%
\pgfpathrectangle{\pgfqpoint{0.407592in}{0.222904in}}{\pgfqpoint{1.411163in}{1.411163in}}%
\pgfusepath{clip}%
\pgfsetbuttcap%
\pgfsetroundjoin%
\definecolor{currentfill}{rgb}{0.556863,0.654902,0.847059}%
\pgfsetfillcolor{currentfill}%
\pgfsetfillopacity{0.130000}%
\pgfsetlinewidth{0.000000pt}%
\definecolor{currentstroke}{rgb}{0.000000,0.000000,0.000000}%
\pgfsetstrokecolor{currentstroke}%
\pgfsetdash{}{0pt}%
\pgfpathmoveto{\pgfqpoint{1.200658in}{0.861867in}}%
\pgfpathlineto{\pgfqpoint{1.216791in}{0.860263in}}%
\pgfpathlineto{\pgfqpoint{1.232974in}{0.860001in}}%
\pgfpathlineto{\pgfqpoint{1.240766in}{0.874191in}}%
\pgfpathlineto{\pgfqpoint{1.248544in}{0.888691in}}%
\pgfpathlineto{\pgfqpoint{1.232373in}{0.889232in}}%
\pgfpathlineto{\pgfqpoint{1.216253in}{0.891004in}}%
\pgfpathlineto{\pgfqpoint{1.208463in}{0.876289in}}%
\pgfpathlineto{\pgfqpoint{1.200658in}{0.861867in}}%
\pgfpathclose%
\pgfusepath{fill}%
\end{pgfscope}%
\begin{pgfscope}%
\pgfpathrectangle{\pgfqpoint{0.407592in}{0.222904in}}{\pgfqpoint{1.411163in}{1.411163in}}%
\pgfusepath{clip}%
\pgfsetbuttcap%
\pgfsetroundjoin%
\definecolor{currentfill}{rgb}{0.556863,0.654902,0.847059}%
\pgfsetfillcolor{currentfill}%
\pgfsetfillopacity{0.130000}%
\pgfsetlinewidth{0.000000pt}%
\definecolor{currentstroke}{rgb}{0.000000,0.000000,0.000000}%
\pgfsetstrokecolor{currentstroke}%
\pgfsetdash{}{0pt}%
\pgfpathmoveto{\pgfqpoint{1.311986in}{0.950046in}}%
\pgfpathlineto{\pgfqpoint{1.328302in}{0.951970in}}%
\pgfpathlineto{\pgfqpoint{1.344679in}{0.954769in}}%
\pgfpathlineto{\pgfqpoint{1.352368in}{0.969605in}}%
\pgfpathlineto{\pgfqpoint{1.360045in}{0.984614in}}%
\pgfpathlineto{\pgfqpoint{1.343681in}{0.982223in}}%
\pgfpathlineto{\pgfqpoint{1.327377in}{0.980663in}}%
\pgfpathlineto{\pgfqpoint{1.319688in}{0.965280in}}%
\pgfpathlineto{\pgfqpoint{1.311986in}{0.950046in}}%
\pgfpathclose%
\pgfusepath{fill}%
\end{pgfscope}%
\begin{pgfscope}%
\pgfpathrectangle{\pgfqpoint{0.407592in}{0.222904in}}{\pgfqpoint{1.411163in}{1.411163in}}%
\pgfusepath{clip}%
\pgfsetbuttcap%
\pgfsetroundjoin%
\definecolor{currentfill}{rgb}{0.556863,0.654902,0.847059}%
\pgfsetfillcolor{currentfill}%
\pgfsetfillopacity{0.130000}%
\pgfsetlinewidth{0.000000pt}%
\definecolor{currentstroke}{rgb}{0.000000,0.000000,0.000000}%
\pgfsetstrokecolor{currentstroke}%
\pgfsetdash{}{0pt}%
\pgfpathmoveto{\pgfqpoint{1.009277in}{0.835997in}}%
\pgfpathlineto{\pgfqpoint{1.025250in}{0.821950in}}%
\pgfpathlineto{\pgfqpoint{1.041239in}{0.808995in}}%
\pgfpathlineto{\pgfqpoint{1.049251in}{0.818015in}}%
\pgfpathlineto{\pgfqpoint{1.057239in}{0.827855in}}%
\pgfpathlineto{\pgfqpoint{1.041282in}{0.840334in}}%
\pgfpathlineto{\pgfqpoint{1.025342in}{0.853898in}}%
\pgfpathlineto{\pgfqpoint{1.017322in}{0.844570in}}%
\pgfpathlineto{\pgfqpoint{1.009277in}{0.835997in}}%
\pgfpathclose%
\pgfusepath{fill}%
\end{pgfscope}%
\begin{pgfscope}%
\pgfpathrectangle{\pgfqpoint{0.407592in}{0.222904in}}{\pgfqpoint{1.411163in}{1.411163in}}%
\pgfusepath{clip}%
\pgfsetbuttcap%
\pgfsetroundjoin%
\definecolor{currentfill}{rgb}{0.556863,0.654902,0.847059}%
\pgfsetfillcolor{currentfill}%
\pgfsetfillopacity{0.130000}%
\pgfsetlinewidth{0.000000pt}%
\definecolor{currentstroke}{rgb}{0.000000,0.000000,0.000000}%
\pgfsetstrokecolor{currentstroke}%
\pgfsetdash{}{0pt}%
\pgfpathmoveto{\pgfqpoint{0.784995in}{1.024562in}}%
\pgfpathlineto{\pgfqpoint{0.800930in}{1.005288in}}%
\pgfpathlineto{\pgfqpoint{0.816873in}{0.986145in}}%
\pgfpathlineto{\pgfqpoint{0.825244in}{0.991020in}}%
\pgfpathlineto{\pgfqpoint{0.833582in}{0.996393in}}%
\pgfpathlineto{\pgfqpoint{0.817681in}{1.015451in}}%
\pgfpathlineto{\pgfqpoint{0.801789in}{1.034642in}}%
\pgfpathlineto{\pgfqpoint{0.793409in}{1.029376in}}%
\pgfpathlineto{\pgfqpoint{0.784995in}{1.024562in}}%
\pgfpathclose%
\pgfusepath{fill}%
\end{pgfscope}%
\begin{pgfscope}%
\pgfpathrectangle{\pgfqpoint{0.407592in}{0.222904in}}{\pgfqpoint{1.411163in}{1.411163in}}%
\pgfusepath{clip}%
\pgfsetbuttcap%
\pgfsetroundjoin%
\definecolor{currentfill}{rgb}{0.556863,0.654902,0.847059}%
\pgfsetfillcolor{currentfill}%
\pgfsetfillopacity{0.130000}%
\pgfsetlinewidth{0.000000pt}%
\definecolor{currentstroke}{rgb}{0.000000,0.000000,0.000000}%
\pgfsetstrokecolor{currentstroke}%
\pgfsetdash{}{0pt}%
\pgfpathmoveto{\pgfqpoint{1.423491in}{1.051556in}}%
\pgfpathlineto{\pgfqpoint{1.440017in}{1.055315in}}%
\pgfpathlineto{\pgfqpoint{1.456607in}{1.059690in}}%
\pgfpathlineto{\pgfqpoint{1.464196in}{1.074587in}}%
\pgfpathlineto{\pgfqpoint{1.471773in}{1.089605in}}%
\pgfpathlineto{\pgfqpoint{1.455197in}{1.085615in}}%
\pgfpathlineto{\pgfqpoint{1.438685in}{1.082223in}}%
\pgfpathlineto{\pgfqpoint{1.431094in}{1.066839in}}%
\pgfpathlineto{\pgfqpoint{1.423491in}{1.051556in}}%
\pgfpathclose%
\pgfusepath{fill}%
\end{pgfscope}%
\begin{pgfscope}%
\pgfpathrectangle{\pgfqpoint{0.407592in}{0.222904in}}{\pgfqpoint{1.411163in}{1.411163in}}%
\pgfusepath{clip}%
\pgfsetbuttcap%
\pgfsetroundjoin%
\definecolor{currentfill}{rgb}{0.556863,0.654902,0.847059}%
\pgfsetfillcolor{currentfill}%
\pgfsetfillopacity{0.130000}%
\pgfsetlinewidth{0.000000pt}%
\definecolor{currentstroke}{rgb}{0.000000,0.000000,0.000000}%
\pgfsetstrokecolor{currentstroke}%
\pgfsetdash{}{0pt}%
\pgfpathmoveto{\pgfqpoint{1.248544in}{0.888691in}}%
\pgfpathlineto{\pgfqpoint{1.264770in}{0.889330in}}%
\pgfpathlineto{\pgfqpoint{1.281052in}{0.891084in}}%
\pgfpathlineto{\pgfqpoint{1.288805in}{0.905473in}}%
\pgfpathlineto{\pgfqpoint{1.296545in}{0.920117in}}%
\pgfpathlineto{\pgfqpoint{1.280275in}{0.918765in}}%
\pgfpathlineto{\pgfqpoint{1.264061in}{0.918456in}}%
\pgfpathlineto{\pgfqpoint{1.256310in}{0.903458in}}%
\pgfpathlineto{\pgfqpoint{1.248544in}{0.888691in}}%
\pgfpathclose%
\pgfusepath{fill}%
\end{pgfscope}%
\begin{pgfscope}%
\pgfpathrectangle{\pgfqpoint{0.407592in}{0.222904in}}{\pgfqpoint{1.411163in}{1.411163in}}%
\pgfusepath{clip}%
\pgfsetbuttcap%
\pgfsetroundjoin%
\definecolor{currentfill}{rgb}{0.556863,0.654902,0.847059}%
\pgfsetfillcolor{currentfill}%
\pgfsetfillopacity{0.130000}%
\pgfsetlinewidth{0.000000pt}%
\definecolor{currentstroke}{rgb}{0.000000,0.000000,0.000000}%
\pgfsetstrokecolor{currentstroke}%
\pgfsetdash{}{0pt}%
\pgfpathmoveto{\pgfqpoint{0.929177in}{0.886437in}}%
\pgfpathlineto{\pgfqpoint{0.945143in}{0.869264in}}%
\pgfpathlineto{\pgfqpoint{0.961118in}{0.852608in}}%
\pgfpathlineto{\pgfqpoint{0.969259in}{0.859276in}}%
\pgfpathlineto{\pgfqpoint{0.977369in}{0.866693in}}%
\pgfpathlineto{\pgfqpoint{0.961432in}{0.883064in}}%
\pgfpathlineto{\pgfqpoint{0.945507in}{0.899967in}}%
\pgfpathlineto{\pgfqpoint{0.937357in}{0.892868in}}%
\pgfpathlineto{\pgfqpoint{0.929177in}{0.886437in}}%
\pgfpathclose%
\pgfusepath{fill}%
\end{pgfscope}%
\begin{pgfscope}%
\pgfpathrectangle{\pgfqpoint{0.407592in}{0.222904in}}{\pgfqpoint{1.411163in}{1.411163in}}%
\pgfusepath{clip}%
\pgfsetbuttcap%
\pgfsetroundjoin%
\definecolor{currentfill}{rgb}{0.556863,0.654902,0.847059}%
\pgfsetfillcolor{currentfill}%
\pgfsetfillopacity{0.130000}%
\pgfsetlinewidth{0.000000pt}%
\definecolor{currentstroke}{rgb}{0.000000,0.000000,0.000000}%
\pgfsetstrokecolor{currentstroke}%
\pgfsetdash{}{0pt}%
\pgfpathmoveto{\pgfqpoint{0.640558in}{1.173621in}}%
\pgfpathlineto{\pgfqpoint{0.656466in}{1.153746in}}%
\pgfpathlineto{\pgfqpoint{0.672382in}{1.133915in}}%
\pgfpathlineto{\pgfqpoint{0.680995in}{1.137829in}}%
\pgfpathlineto{\pgfqpoint{0.689572in}{1.142098in}}%
\pgfpathlineto{\pgfqpoint{0.673699in}{1.161904in}}%
\pgfpathlineto{\pgfqpoint{0.657834in}{1.181752in}}%
\pgfpathlineto{\pgfqpoint{0.649214in}{1.177521in}}%
\pgfpathlineto{\pgfqpoint{0.640558in}{1.173621in}}%
\pgfpathclose%
\pgfusepath{fill}%
\end{pgfscope}%
\begin{pgfscope}%
\pgfpathrectangle{\pgfqpoint{0.407592in}{0.222904in}}{\pgfqpoint{1.411163in}{1.411163in}}%
\pgfusepath{clip}%
\pgfsetbuttcap%
\pgfsetroundjoin%
\definecolor{currentfill}{rgb}{0.556863,0.654902,0.847059}%
\pgfsetfillcolor{currentfill}%
\pgfsetfillopacity{0.130000}%
\pgfsetlinewidth{0.000000pt}%
\definecolor{currentstroke}{rgb}{0.000000,0.000000,0.000000}%
\pgfsetstrokecolor{currentstroke}%
\pgfsetdash{}{0pt}%
\pgfpathmoveto{\pgfqpoint{1.360045in}{0.984614in}}%
\pgfpathlineto{\pgfqpoint{1.376471in}{0.987791in}}%
\pgfpathlineto{\pgfqpoint{1.392960in}{0.991706in}}%
\pgfpathlineto{\pgfqpoint{1.400611in}{1.006448in}}%
\pgfpathlineto{\pgfqpoint{1.408250in}{1.021348in}}%
\pgfpathlineto{\pgfqpoint{1.391774in}{1.017856in}}%
\pgfpathlineto{\pgfqpoint{1.375362in}{1.015076in}}%
\pgfpathlineto{\pgfqpoint{1.367710in}{0.999777in}}%
\pgfpathlineto{\pgfqpoint{1.360045in}{0.984614in}}%
\pgfpathclose%
\pgfusepath{fill}%
\end{pgfscope}%
\begin{pgfscope}%
\pgfpathrectangle{\pgfqpoint{0.407592in}{0.222904in}}{\pgfqpoint{1.411163in}{1.411163in}}%
\pgfusepath{clip}%
\pgfsetbuttcap%
\pgfsetroundjoin%
\definecolor{currentfill}{rgb}{0.556863,0.654902,0.847059}%
\pgfsetfillcolor{currentfill}%
\pgfsetfillopacity{0.130000}%
\pgfsetlinewidth{0.000000pt}%
\definecolor{currentstroke}{rgb}{0.000000,0.000000,0.000000}%
\pgfsetstrokecolor{currentstroke}%
\pgfsetdash{}{0pt}%
\pgfpathmoveto{\pgfqpoint{0.816873in}{0.986145in}}%
\pgfpathlineto{\pgfqpoint{0.832825in}{0.967154in}}%
\pgfpathlineto{\pgfqpoint{0.848786in}{0.948342in}}%
\pgfpathlineto{\pgfqpoint{0.857115in}{0.953279in}}%
\pgfpathlineto{\pgfqpoint{0.865410in}{0.958767in}}%
\pgfpathlineto{\pgfqpoint{0.849492in}{0.977490in}}%
\pgfpathlineto{\pgfqpoint{0.833582in}{0.996393in}}%
\pgfpathlineto{\pgfqpoint{0.825244in}{0.991020in}}%
\pgfpathlineto{\pgfqpoint{0.816873in}{0.986145in}}%
\pgfpathclose%
\pgfusepath{fill}%
\end{pgfscope}%
\begin{pgfscope}%
\pgfpathrectangle{\pgfqpoint{0.407592in}{0.222904in}}{\pgfqpoint{1.411163in}{1.411163in}}%
\pgfusepath{clip}%
\pgfsetbuttcap%
\pgfsetroundjoin%
\definecolor{currentfill}{rgb}{0.556863,0.654902,0.847059}%
\pgfsetfillcolor{currentfill}%
\pgfsetfillopacity{0.130000}%
\pgfsetlinewidth{0.000000pt}%
\definecolor{currentstroke}{rgb}{0.000000,0.000000,0.000000}%
\pgfsetstrokecolor{currentstroke}%
\pgfsetdash{}{0pt}%
\pgfpathmoveto{\pgfqpoint{1.089214in}{0.806844in}}%
\pgfpathlineto{\pgfqpoint{1.105238in}{0.798654in}}%
\pgfpathlineto{\pgfqpoint{1.121293in}{0.792199in}}%
\pgfpathlineto{\pgfqpoint{1.129210in}{0.803554in}}%
\pgfpathlineto{\pgfqpoint{1.137108in}{0.815616in}}%
\pgfpathlineto{\pgfqpoint{1.121073in}{0.821845in}}%
\pgfpathlineto{\pgfqpoint{1.105070in}{0.829681in}}%
\pgfpathlineto{\pgfqpoint{1.097152in}{0.817914in}}%
\pgfpathlineto{\pgfqpoint{1.089214in}{0.806844in}}%
\pgfpathclose%
\pgfusepath{fill}%
\end{pgfscope}%
\begin{pgfscope}%
\pgfpathrectangle{\pgfqpoint{0.407592in}{0.222904in}}{\pgfqpoint{1.411163in}{1.411163in}}%
\pgfusepath{clip}%
\pgfsetbuttcap%
\pgfsetroundjoin%
\definecolor{currentfill}{rgb}{0.556863,0.654902,0.847059}%
\pgfsetfillcolor{currentfill}%
\pgfsetfillopacity{0.130000}%
\pgfsetlinewidth{0.000000pt}%
\definecolor{currentstroke}{rgb}{0.000000,0.000000,0.000000}%
\pgfsetstrokecolor{currentstroke}%
\pgfsetdash{}{0pt}%
\pgfpathmoveto{\pgfqpoint{1.137108in}{0.815616in}}%
\pgfpathlineto{\pgfqpoint{1.153180in}{0.811061in}}%
\pgfpathlineto{\pgfqpoint{1.169293in}{0.808200in}}%
\pgfpathlineto{\pgfqpoint{1.177157in}{0.820884in}}%
\pgfpathlineto{\pgfqpoint{1.185006in}{0.834104in}}%
\pgfpathlineto{\pgfqpoint{1.168908in}{0.837034in}}%
\pgfpathlineto{\pgfqpoint{1.152852in}{0.841507in}}%
\pgfpathlineto{\pgfqpoint{1.144988in}{0.828294in}}%
\pgfpathlineto{\pgfqpoint{1.137108in}{0.815616in}}%
\pgfpathclose%
\pgfusepath{fill}%
\end{pgfscope}%
\begin{pgfscope}%
\pgfpathrectangle{\pgfqpoint{0.407592in}{0.222904in}}{\pgfqpoint{1.411163in}{1.411163in}}%
\pgfusepath{clip}%
\pgfsetbuttcap%
\pgfsetroundjoin%
\definecolor{currentfill}{rgb}{0.556863,0.654902,0.847059}%
\pgfsetfillcolor{currentfill}%
\pgfsetfillopacity{0.130000}%
\pgfsetlinewidth{0.000000pt}%
\definecolor{currentstroke}{rgb}{0.000000,0.000000,0.000000}%
\pgfsetstrokecolor{currentstroke}%
\pgfsetdash{}{0pt}%
\pgfpathmoveto{\pgfqpoint{0.672382in}{1.133915in}}%
\pgfpathlineto{\pgfqpoint{0.688307in}{1.114135in}}%
\pgfpathlineto{\pgfqpoint{0.704240in}{1.094413in}}%
\pgfpathlineto{\pgfqpoint{0.712811in}{1.098333in}}%
\pgfpathlineto{\pgfqpoint{0.721345in}{1.102638in}}%
\pgfpathlineto{\pgfqpoint{0.705454in}{1.122340in}}%
\pgfpathlineto{\pgfqpoint{0.689572in}{1.142098in}}%
\pgfpathlineto{\pgfqpoint{0.680995in}{1.137829in}}%
\pgfpathlineto{\pgfqpoint{0.672382in}{1.133915in}}%
\pgfpathclose%
\pgfusepath{fill}%
\end{pgfscope}%
\begin{pgfscope}%
\pgfpathrectangle{\pgfqpoint{0.407592in}{0.222904in}}{\pgfqpoint{1.411163in}{1.411163in}}%
\pgfusepath{clip}%
\pgfsetbuttcap%
\pgfsetroundjoin%
\definecolor{currentfill}{rgb}{0.556863,0.654902,0.847059}%
\pgfsetfillcolor{currentfill}%
\pgfsetfillopacity{0.130000}%
\pgfsetlinewidth{0.000000pt}%
\definecolor{currentstroke}{rgb}{0.000000,0.000000,0.000000}%
\pgfsetstrokecolor{currentstroke}%
\pgfsetdash{}{0pt}%
\pgfpathmoveto{\pgfqpoint{0.961118in}{0.852608in}}%
\pgfpathlineto{\pgfqpoint{0.977104in}{0.836576in}}%
\pgfpathlineto{\pgfqpoint{0.993102in}{0.821299in}}%
\pgfpathlineto{\pgfqpoint{1.001204in}{0.828226in}}%
\pgfpathlineto{\pgfqpoint{1.009277in}{0.835997in}}%
\pgfpathlineto{\pgfqpoint{0.993317in}{0.850961in}}%
\pgfpathlineto{\pgfqpoint{0.977369in}{0.866693in}}%
\pgfpathlineto{\pgfqpoint{0.969259in}{0.859276in}}%
\pgfpathlineto{\pgfqpoint{0.961118in}{0.852608in}}%
\pgfpathclose%
\pgfusepath{fill}%
\end{pgfscope}%
\begin{pgfscope}%
\pgfpathrectangle{\pgfqpoint{0.407592in}{0.222904in}}{\pgfqpoint{1.411163in}{1.411163in}}%
\pgfusepath{clip}%
\pgfsetbuttcap%
\pgfsetroundjoin%
\definecolor{currentfill}{rgb}{0.556863,0.654902,0.847059}%
\pgfsetfillcolor{currentfill}%
\pgfsetfillopacity{0.130000}%
\pgfsetlinewidth{0.000000pt}%
\definecolor{currentstroke}{rgb}{0.000000,0.000000,0.000000}%
\pgfsetstrokecolor{currentstroke}%
\pgfsetdash{}{0pt}%
\pgfpathmoveto{\pgfqpoint{1.296545in}{0.920117in}}%
\pgfpathlineto{\pgfqpoint{1.312873in}{0.922452in}}%
\pgfpathlineto{\pgfqpoint{1.329263in}{0.925705in}}%
\pgfpathlineto{\pgfqpoint{1.336977in}{0.940127in}}%
\pgfpathlineto{\pgfqpoint{1.344679in}{0.954769in}}%
\pgfpathlineto{\pgfqpoint{1.328302in}{0.951970in}}%
\pgfpathlineto{\pgfqpoint{1.311986in}{0.950046in}}%
\pgfpathlineto{\pgfqpoint{1.304272in}{0.934984in}}%
\pgfpathlineto{\pgfqpoint{1.296545in}{0.920117in}}%
\pgfpathclose%
\pgfusepath{fill}%
\end{pgfscope}%
\begin{pgfscope}%
\pgfpathrectangle{\pgfqpoint{0.407592in}{0.222904in}}{\pgfqpoint{1.411163in}{1.411163in}}%
\pgfusepath{clip}%
\pgfsetbuttcap%
\pgfsetroundjoin%
\definecolor{currentfill}{rgb}{0.556863,0.654902,0.847059}%
\pgfsetfillcolor{currentfill}%
\pgfsetfillopacity{0.130000}%
\pgfsetlinewidth{0.000000pt}%
\definecolor{currentstroke}{rgb}{0.000000,0.000000,0.000000}%
\pgfsetstrokecolor{currentstroke}%
\pgfsetdash{}{0pt}%
\pgfpathmoveto{\pgfqpoint{1.185006in}{0.834104in}}%
\pgfpathlineto{\pgfqpoint{1.201151in}{0.832694in}}%
\pgfpathlineto{\pgfqpoint{1.217348in}{0.832744in}}%
\pgfpathlineto{\pgfqpoint{1.225168in}{0.846167in}}%
\pgfpathlineto{\pgfqpoint{1.232974in}{0.860001in}}%
\pgfpathlineto{\pgfqpoint{1.216791in}{0.860263in}}%
\pgfpathlineto{\pgfqpoint{1.200658in}{0.861867in}}%
\pgfpathlineto{\pgfqpoint{1.192840in}{0.847786in}}%
\pgfpathlineto{\pgfqpoint{1.185006in}{0.834104in}}%
\pgfpathclose%
\pgfusepath{fill}%
\end{pgfscope}%
\begin{pgfscope}%
\pgfpathrectangle{\pgfqpoint{0.407592in}{0.222904in}}{\pgfqpoint{1.411163in}{1.411163in}}%
\pgfusepath{clip}%
\pgfsetbuttcap%
\pgfsetroundjoin%
\definecolor{currentfill}{rgb}{0.556863,0.654902,0.847059}%
\pgfsetfillcolor{currentfill}%
\pgfsetfillopacity{0.130000}%
\pgfsetlinewidth{0.000000pt}%
\definecolor{currentstroke}{rgb}{0.000000,0.000000,0.000000}%
\pgfsetstrokecolor{currentstroke}%
\pgfsetdash{}{0pt}%
\pgfpathmoveto{\pgfqpoint{1.041239in}{0.808995in}}%
\pgfpathlineto{\pgfqpoint{1.057245in}{0.797326in}}%
\pgfpathlineto{\pgfqpoint{1.073272in}{0.787145in}}%
\pgfpathlineto{\pgfqpoint{1.081254in}{0.796559in}}%
\pgfpathlineto{\pgfqpoint{1.089214in}{0.806844in}}%
\pgfpathlineto{\pgfqpoint{1.073215in}{0.816632in}}%
\pgfpathlineto{\pgfqpoint{1.057239in}{0.827855in}}%
\pgfpathlineto{\pgfqpoint{1.049251in}{0.818015in}}%
\pgfpathlineto{\pgfqpoint{1.041239in}{0.808995in}}%
\pgfpathclose%
\pgfusepath{fill}%
\end{pgfscope}%
\begin{pgfscope}%
\pgfpathrectangle{\pgfqpoint{0.407592in}{0.222904in}}{\pgfqpoint{1.411163in}{1.411163in}}%
\pgfusepath{clip}%
\pgfsetbuttcap%
\pgfsetroundjoin%
\definecolor{currentfill}{rgb}{0.556863,0.654902,0.847059}%
\pgfsetfillcolor{currentfill}%
\pgfsetfillopacity{0.130000}%
\pgfsetlinewidth{0.000000pt}%
\definecolor{currentstroke}{rgb}{0.000000,0.000000,0.000000}%
\pgfsetstrokecolor{currentstroke}%
\pgfsetdash{}{0pt}%
\pgfpathmoveto{\pgfqpoint{0.848786in}{0.948342in}}%
\pgfpathlineto{\pgfqpoint{0.864756in}{0.929741in}}%
\pgfpathlineto{\pgfqpoint{0.880735in}{0.911389in}}%
\pgfpathlineto{\pgfqpoint{0.889022in}{0.916386in}}%
\pgfpathlineto{\pgfqpoint{0.897275in}{0.921998in}}%
\pgfpathlineto{\pgfqpoint{0.881338in}{0.940256in}}%
\pgfpathlineto{\pgfqpoint{0.865410in}{0.958767in}}%
\pgfpathlineto{\pgfqpoint{0.857115in}{0.953279in}}%
\pgfpathlineto{\pgfqpoint{0.848786in}{0.948342in}}%
\pgfpathclose%
\pgfusepath{fill}%
\end{pgfscope}%
\begin{pgfscope}%
\pgfpathrectangle{\pgfqpoint{0.407592in}{0.222904in}}{\pgfqpoint{1.411163in}{1.411163in}}%
\pgfusepath{clip}%
\pgfsetbuttcap%
\pgfsetroundjoin%
\definecolor{currentfill}{rgb}{0.556863,0.654902,0.847059}%
\pgfsetfillcolor{currentfill}%
\pgfsetfillopacity{0.130000}%
\pgfsetlinewidth{0.000000pt}%
\definecolor{currentstroke}{rgb}{0.000000,0.000000,0.000000}%
\pgfsetstrokecolor{currentstroke}%
\pgfsetdash{}{0pt}%
\pgfpathmoveto{\pgfqpoint{0.704240in}{1.094413in}}%
\pgfpathlineto{\pgfqpoint{0.720183in}{1.074758in}}%
\pgfpathlineto{\pgfqpoint{0.736134in}{1.055179in}}%
\pgfpathlineto{\pgfqpoint{0.744661in}{1.059097in}}%
\pgfpathlineto{\pgfqpoint{0.753153in}{1.063432in}}%
\pgfpathlineto{\pgfqpoint{0.737244in}{1.082998in}}%
\pgfpathlineto{\pgfqpoint{0.721345in}{1.102638in}}%
\pgfpathlineto{\pgfqpoint{0.712811in}{1.098333in}}%
\pgfpathlineto{\pgfqpoint{0.704240in}{1.094413in}}%
\pgfpathclose%
\pgfusepath{fill}%
\end{pgfscope}%
\begin{pgfscope}%
\pgfpathrectangle{\pgfqpoint{0.407592in}{0.222904in}}{\pgfqpoint{1.411163in}{1.411163in}}%
\pgfusepath{clip}%
\pgfsetbuttcap%
\pgfsetroundjoin%
\definecolor{currentfill}{rgb}{0.556863,0.654902,0.847059}%
\pgfsetfillcolor{currentfill}%
\pgfsetfillopacity{0.130000}%
\pgfsetlinewidth{0.000000pt}%
\definecolor{currentstroke}{rgb}{0.000000,0.000000,0.000000}%
\pgfsetstrokecolor{currentstroke}%
\pgfsetdash{}{0pt}%
\pgfpathmoveto{\pgfqpoint{1.471773in}{1.089605in}}%
\pgfpathlineto{\pgfqpoint{1.488414in}{1.094163in}}%
\pgfpathlineto{\pgfqpoint{1.505121in}{1.099259in}}%
\pgfpathlineto{\pgfqpoint{1.512671in}{1.114007in}}%
\pgfpathlineto{\pgfqpoint{1.495971in}{1.109103in}}%
\pgfpathlineto{\pgfqpoint{1.479338in}{1.104731in}}%
\pgfusepath{fill}%
\end{pgfscope}%
\begin{pgfscope}%
\pgfpathrectangle{\pgfqpoint{0.407592in}{0.222904in}}{\pgfqpoint{1.411163in}{1.411163in}}%
\pgfusepath{clip}%
\pgfsetbuttcap%
\pgfsetroundjoin%
\definecolor{currentfill}{rgb}{0.556863,0.654902,0.847059}%
\pgfsetfillcolor{currentfill}%
\pgfsetfillopacity{0.130000}%
\pgfsetlinewidth{0.000000pt}%
\definecolor{currentstroke}{rgb}{0.000000,0.000000,0.000000}%
\pgfsetstrokecolor{currentstroke}%
\pgfsetdash{}{0pt}%
\pgfpathmoveto{\pgfqpoint{1.408250in}{1.021348in}}%
\pgfpathlineto{\pgfqpoint{1.424789in}{1.025511in}}%
\pgfpathlineto{\pgfqpoint{1.441394in}{1.030305in}}%
\pgfpathlineto{\pgfqpoint{1.449006in}{1.044924in}}%
\pgfpathlineto{\pgfqpoint{1.456607in}{1.059690in}}%
\pgfpathlineto{\pgfqpoint{1.440017in}{1.055315in}}%
\pgfpathlineto{\pgfqpoint{1.423491in}{1.051556in}}%
\pgfpathlineto{\pgfqpoint{1.415877in}{1.036388in}}%
\pgfpathlineto{\pgfqpoint{1.408250in}{1.021348in}}%
\pgfpathclose%
\pgfusepath{fill}%
\end{pgfscope}%
\begin{pgfscope}%
\pgfpathrectangle{\pgfqpoint{0.407592in}{0.222904in}}{\pgfqpoint{1.411163in}{1.411163in}}%
\pgfusepath{clip}%
\pgfsetbuttcap%
\pgfsetroundjoin%
\definecolor{currentfill}{rgb}{0.556863,0.654902,0.847059}%
\pgfsetfillcolor{currentfill}%
\pgfsetfillopacity{0.130000}%
\pgfsetlinewidth{0.000000pt}%
\definecolor{currentstroke}{rgb}{0.000000,0.000000,0.000000}%
\pgfsetstrokecolor{currentstroke}%
\pgfsetdash{}{0pt}%
\pgfpathmoveto{\pgfqpoint{1.232974in}{0.860001in}}%
\pgfpathlineto{\pgfqpoint{1.249211in}{0.861012in}}%
\pgfpathlineto{\pgfqpoint{1.265506in}{0.863214in}}%
\pgfpathlineto{\pgfqpoint{1.273286in}{0.876984in}}%
\pgfpathlineto{\pgfqpoint{1.281052in}{0.891084in}}%
\pgfpathlineto{\pgfqpoint{1.264770in}{0.889330in}}%
\pgfpathlineto{\pgfqpoint{1.248544in}{0.888691in}}%
\pgfpathlineto{\pgfqpoint{1.240766in}{0.874191in}}%
\pgfpathlineto{\pgfqpoint{1.232974in}{0.860001in}}%
\pgfpathclose%
\pgfusepath{fill}%
\end{pgfscope}%
\begin{pgfscope}%
\pgfpathrectangle{\pgfqpoint{0.407592in}{0.222904in}}{\pgfqpoint{1.411163in}{1.411163in}}%
\pgfusepath{clip}%
\pgfsetbuttcap%
\pgfsetroundjoin%
\definecolor{currentfill}{rgb}{0.556863,0.654902,0.847059}%
\pgfsetfillcolor{currentfill}%
\pgfsetfillopacity{0.130000}%
\pgfsetlinewidth{0.000000pt}%
\definecolor{currentstroke}{rgb}{0.000000,0.000000,0.000000}%
\pgfsetstrokecolor{currentstroke}%
\pgfsetdash{}{0pt}%
\pgfpathmoveto{\pgfqpoint{0.880735in}{0.911389in}}%
\pgfpathlineto{\pgfqpoint{0.896723in}{0.893334in}}%
\pgfpathlineto{\pgfqpoint{0.912720in}{0.875633in}}%
\pgfpathlineto{\pgfqpoint{0.920965in}{0.880689in}}%
\pgfpathlineto{\pgfqpoint{0.929177in}{0.886437in}}%
\pgfpathlineto{\pgfqpoint{0.913221in}{0.904039in}}%
\pgfpathlineto{\pgfqpoint{0.897275in}{0.921998in}}%
\pgfpathlineto{\pgfqpoint{0.889022in}{0.916386in}}%
\pgfpathlineto{\pgfqpoint{0.880735in}{0.911389in}}%
\pgfpathclose%
\pgfusepath{fill}%
\end{pgfscope}%
\begin{pgfscope}%
\pgfpathrectangle{\pgfqpoint{0.407592in}{0.222904in}}{\pgfqpoint{1.411163in}{1.411163in}}%
\pgfusepath{clip}%
\pgfsetbuttcap%
\pgfsetroundjoin%
\definecolor{currentfill}{rgb}{0.556863,0.654902,0.847059}%
\pgfsetfillcolor{currentfill}%
\pgfsetfillopacity{0.130000}%
\pgfsetlinewidth{0.000000pt}%
\definecolor{currentstroke}{rgb}{0.000000,0.000000,0.000000}%
\pgfsetstrokecolor{currentstroke}%
\pgfsetdash{}{0pt}%
\pgfpathmoveto{\pgfqpoint{0.736134in}{1.055179in}}%
\pgfpathlineto{\pgfqpoint{0.752094in}{1.035688in}}%
\pgfpathlineto{\pgfqpoint{0.768062in}{1.016299in}}%
\pgfpathlineto{\pgfqpoint{0.776547in}{1.020204in}}%
\pgfpathlineto{\pgfqpoint{0.784995in}{1.024562in}}%
\pgfpathlineto{\pgfqpoint{0.769070in}{1.043949in}}%
\pgfpathlineto{\pgfqpoint{0.753153in}{1.063432in}}%
\pgfpathlineto{\pgfqpoint{0.744661in}{1.059097in}}%
\pgfpathlineto{\pgfqpoint{0.736134in}{1.055179in}}%
\pgfpathclose%
\pgfusepath{fill}%
\end{pgfscope}%
\begin{pgfscope}%
\pgfpathrectangle{\pgfqpoint{0.407592in}{0.222904in}}{\pgfqpoint{1.411163in}{1.411163in}}%
\pgfusepath{clip}%
\pgfsetbuttcap%
\pgfsetroundjoin%
\definecolor{currentfill}{rgb}{0.556863,0.654902,0.847059}%
\pgfsetfillcolor{currentfill}%
\pgfsetfillopacity{0.130000}%
\pgfsetlinewidth{0.000000pt}%
\definecolor{currentstroke}{rgb}{0.000000,0.000000,0.000000}%
\pgfsetstrokecolor{currentstroke}%
\pgfsetdash{}{0pt}%
\pgfpathmoveto{\pgfqpoint{0.993102in}{0.821299in}}%
\pgfpathlineto{\pgfqpoint{1.009112in}{0.806933in}}%
\pgfpathlineto{\pgfqpoint{1.025135in}{0.793661in}}%
\pgfpathlineto{\pgfqpoint{1.033201in}{0.800859in}}%
\pgfpathlineto{\pgfqpoint{1.041239in}{0.808995in}}%
\pgfpathlineto{\pgfqpoint{1.025250in}{0.821950in}}%
\pgfpathlineto{\pgfqpoint{1.009277in}{0.835997in}}%
\pgfpathlineto{\pgfqpoint{1.001204in}{0.828226in}}%
\pgfpathlineto{\pgfqpoint{0.993102in}{0.821299in}}%
\pgfpathclose%
\pgfusepath{fill}%
\end{pgfscope}%
\begin{pgfscope}%
\pgfpathrectangle{\pgfqpoint{0.407592in}{0.222904in}}{\pgfqpoint{1.411163in}{1.411163in}}%
\pgfusepath{clip}%
\pgfsetbuttcap%
\pgfsetroundjoin%
\definecolor{currentfill}{rgb}{0.556863,0.654902,0.847059}%
\pgfsetfillcolor{currentfill}%
\pgfsetfillopacity{0.130000}%
\pgfsetlinewidth{0.000000pt}%
\definecolor{currentstroke}{rgb}{0.000000,0.000000,0.000000}%
\pgfsetstrokecolor{currentstroke}%
\pgfsetdash{}{0pt}%
\pgfpathmoveto{\pgfqpoint{1.344679in}{0.954769in}}%
\pgfpathlineto{\pgfqpoint{1.361118in}{0.958387in}}%
\pgfpathlineto{\pgfqpoint{1.377621in}{0.962768in}}%
\pgfpathlineto{\pgfqpoint{1.385296in}{0.977140in}}%
\pgfpathlineto{\pgfqpoint{1.392960in}{0.991706in}}%
\pgfpathlineto{\pgfqpoint{1.376471in}{0.987791in}}%
\pgfpathlineto{\pgfqpoint{1.360045in}{0.984614in}}%
\pgfpathlineto{\pgfqpoint{1.352368in}{0.969605in}}%
\pgfpathlineto{\pgfqpoint{1.344679in}{0.954769in}}%
\pgfpathclose%
\pgfusepath{fill}%
\end{pgfscope}%
\begin{pgfscope}%
\pgfpathrectangle{\pgfqpoint{0.407592in}{0.222904in}}{\pgfqpoint{1.411163in}{1.411163in}}%
\pgfusepath{clip}%
\pgfsetbuttcap%
\pgfsetroundjoin%
\definecolor{currentfill}{rgb}{0.556863,0.654902,0.847059}%
\pgfsetfillcolor{currentfill}%
\pgfsetfillopacity{0.130000}%
\pgfsetlinewidth{0.000000pt}%
\definecolor{currentstroke}{rgb}{0.000000,0.000000,0.000000}%
\pgfsetstrokecolor{currentstroke}%
\pgfsetdash{}{0pt}%
\pgfpathmoveto{\pgfqpoint{0.768062in}{1.016299in}}%
\pgfpathlineto{\pgfqpoint{0.784039in}{0.997028in}}%
\pgfpathlineto{\pgfqpoint{0.800025in}{0.977893in}}%
\pgfpathlineto{\pgfqpoint{0.808467in}{0.981770in}}%
\pgfpathlineto{\pgfqpoint{0.816873in}{0.986145in}}%
\pgfpathlineto{\pgfqpoint{0.800930in}{1.005288in}}%
\pgfpathlineto{\pgfqpoint{0.784995in}{1.024562in}}%
\pgfpathlineto{\pgfqpoint{0.776547in}{1.020204in}}%
\pgfpathlineto{\pgfqpoint{0.768062in}{1.016299in}}%
\pgfpathclose%
\pgfusepath{fill}%
\end{pgfscope}%
\begin{pgfscope}%
\pgfpathrectangle{\pgfqpoint{0.407592in}{0.222904in}}{\pgfqpoint{1.411163in}{1.411163in}}%
\pgfusepath{clip}%
\pgfsetbuttcap%
\pgfsetroundjoin%
\definecolor{currentfill}{rgb}{0.556863,0.654902,0.847059}%
\pgfsetfillcolor{currentfill}%
\pgfsetfillopacity{0.130000}%
\pgfsetlinewidth{0.000000pt}%
\definecolor{currentstroke}{rgb}{0.000000,0.000000,0.000000}%
\pgfsetstrokecolor{currentstroke}%
\pgfsetdash{}{0pt}%
\pgfpathmoveto{\pgfqpoint{1.121293in}{0.792199in}}%
\pgfpathlineto{\pgfqpoint{1.137384in}{0.787564in}}%
\pgfpathlineto{\pgfqpoint{1.153515in}{0.784774in}}%
\pgfpathlineto{\pgfqpoint{1.161412in}{0.796133in}}%
\pgfpathlineto{\pgfqpoint{1.169293in}{0.808200in}}%
\pgfpathlineto{\pgfqpoint{1.153180in}{0.811061in}}%
\pgfpathlineto{\pgfqpoint{1.137108in}{0.815616in}}%
\pgfpathlineto{\pgfqpoint{1.129210in}{0.803554in}}%
\pgfpathlineto{\pgfqpoint{1.121293in}{0.792199in}}%
\pgfpathclose%
\pgfusepath{fill}%
\end{pgfscope}%
\begin{pgfscope}%
\pgfpathrectangle{\pgfqpoint{0.407592in}{0.222904in}}{\pgfqpoint{1.411163in}{1.411163in}}%
\pgfusepath{clip}%
\pgfsetbuttcap%
\pgfsetroundjoin%
\definecolor{currentfill}{rgb}{0.556863,0.654902,0.847059}%
\pgfsetfillcolor{currentfill}%
\pgfsetfillopacity{0.130000}%
\pgfsetlinewidth{0.000000pt}%
\definecolor{currentstroke}{rgb}{0.000000,0.000000,0.000000}%
\pgfsetstrokecolor{currentstroke}%
\pgfsetdash{}{0pt}%
\pgfpathmoveto{\pgfqpoint{0.912720in}{0.875633in}}%
\pgfpathlineto{\pgfqpoint{0.928726in}{0.858359in}}%
\pgfpathlineto{\pgfqpoint{0.944742in}{0.841599in}}%
\pgfpathlineto{\pgfqpoint{0.952947in}{0.846712in}}%
\pgfpathlineto{\pgfqpoint{0.961118in}{0.852608in}}%
\pgfpathlineto{\pgfqpoint{0.945143in}{0.869264in}}%
\pgfpathlineto{\pgfqpoint{0.929177in}{0.886437in}}%
\pgfpathlineto{\pgfqpoint{0.920965in}{0.880689in}}%
\pgfpathlineto{\pgfqpoint{0.912720in}{0.875633in}}%
\pgfpathclose%
\pgfusepath{fill}%
\end{pgfscope}%
\begin{pgfscope}%
\pgfpathrectangle{\pgfqpoint{0.407592in}{0.222904in}}{\pgfqpoint{1.411163in}{1.411163in}}%
\pgfusepath{clip}%
\pgfsetbuttcap%
\pgfsetroundjoin%
\definecolor{currentfill}{rgb}{0.556863,0.654902,0.847059}%
\pgfsetfillcolor{currentfill}%
\pgfsetfillopacity{0.130000}%
\pgfsetlinewidth{0.000000pt}%
\definecolor{currentstroke}{rgb}{0.000000,0.000000,0.000000}%
\pgfsetstrokecolor{currentstroke}%
\pgfsetdash{}{0pt}%
\pgfpathmoveto{\pgfqpoint{1.073272in}{0.787145in}}%
\pgfpathlineto{\pgfqpoint{1.089322in}{0.778647in}}%
\pgfpathlineto{\pgfqpoint{1.105400in}{0.771992in}}%
\pgfpathlineto{\pgfqpoint{1.113357in}{0.781646in}}%
\pgfpathlineto{\pgfqpoint{1.121293in}{0.792199in}}%
\pgfpathlineto{\pgfqpoint{1.105238in}{0.798654in}}%
\pgfpathlineto{\pgfqpoint{1.089214in}{0.806844in}}%
\pgfpathlineto{\pgfqpoint{1.081254in}{0.796559in}}%
\pgfpathlineto{\pgfqpoint{1.073272in}{0.787145in}}%
\pgfpathclose%
\pgfusepath{fill}%
\end{pgfscope}%
\begin{pgfscope}%
\pgfpathrectangle{\pgfqpoint{0.407592in}{0.222904in}}{\pgfqpoint{1.411163in}{1.411163in}}%
\pgfusepath{clip}%
\pgfsetbuttcap%
\pgfsetroundjoin%
\definecolor{currentfill}{rgb}{0.556863,0.654902,0.847059}%
\pgfsetfillcolor{currentfill}%
\pgfsetfillopacity{0.130000}%
\pgfsetlinewidth{0.000000pt}%
\definecolor{currentstroke}{rgb}{0.000000,0.000000,0.000000}%
\pgfsetstrokecolor{currentstroke}%
\pgfsetdash{}{0pt}%
\pgfpathmoveto{\pgfqpoint{1.281052in}{0.891084in}}%
\pgfpathlineto{\pgfqpoint{1.297393in}{0.893877in}}%
\pgfpathlineto{\pgfqpoint{1.313797in}{0.897632in}}%
\pgfpathlineto{\pgfqpoint{1.321536in}{0.911530in}}%
\pgfpathlineto{\pgfqpoint{1.329263in}{0.925705in}}%
\pgfpathlineto{\pgfqpoint{1.312873in}{0.922452in}}%
\pgfpathlineto{\pgfqpoint{1.296545in}{0.920117in}}%
\pgfpathlineto{\pgfqpoint{1.288805in}{0.905473in}}%
\pgfpathlineto{\pgfqpoint{1.281052in}{0.891084in}}%
\pgfpathclose%
\pgfusepath{fill}%
\end{pgfscope}%
\begin{pgfscope}%
\pgfpathrectangle{\pgfqpoint{0.407592in}{0.222904in}}{\pgfqpoint{1.411163in}{1.411163in}}%
\pgfusepath{clip}%
\pgfsetbuttcap%
\pgfsetroundjoin%
\definecolor{currentfill}{rgb}{0.556863,0.654902,0.847059}%
\pgfsetfillcolor{currentfill}%
\pgfsetfillopacity{0.130000}%
\pgfsetlinewidth{0.000000pt}%
\definecolor{currentstroke}{rgb}{0.000000,0.000000,0.000000}%
\pgfsetstrokecolor{currentstroke}%
\pgfsetdash{}{0pt}%
\pgfpathmoveto{\pgfqpoint{1.456607in}{1.059690in}}%
\pgfpathlineto{\pgfqpoint{1.473263in}{1.064645in}}%
\pgfpathlineto{\pgfqpoint{1.489986in}{1.070146in}}%
\pgfpathlineto{\pgfqpoint{1.497560in}{1.084635in}}%
\pgfpathlineto{\pgfqpoint{1.505121in}{1.099259in}}%
\pgfpathlineto{\pgfqpoint{1.488414in}{1.094163in}}%
\pgfpathlineto{\pgfqpoint{1.471773in}{1.089605in}}%
\pgfpathlineto{\pgfqpoint{1.464196in}{1.074587in}}%
\pgfpathlineto{\pgfqpoint{1.456607in}{1.059690in}}%
\pgfpathclose%
\pgfusepath{fill}%
\end{pgfscope}%
\begin{pgfscope}%
\pgfpathrectangle{\pgfqpoint{0.407592in}{0.222904in}}{\pgfqpoint{1.411163in}{1.411163in}}%
\pgfusepath{clip}%
\pgfsetbuttcap%
\pgfsetroundjoin%
\definecolor{currentfill}{rgb}{0.556863,0.654902,0.847059}%
\pgfsetfillcolor{currentfill}%
\pgfsetfillopacity{0.130000}%
\pgfsetlinewidth{0.000000pt}%
\definecolor{currentstroke}{rgb}{0.000000,0.000000,0.000000}%
\pgfsetstrokecolor{currentstroke}%
\pgfsetdash{}{0pt}%
\pgfpathmoveto{\pgfqpoint{0.623136in}{1.166815in}}%
\pgfpathlineto{\pgfqpoint{0.639086in}{1.146958in}}%
\pgfpathlineto{\pgfqpoint{0.655045in}{1.127152in}}%
\pgfpathlineto{\pgfqpoint{0.663732in}{1.130356in}}%
\pgfpathlineto{\pgfqpoint{0.672382in}{1.133915in}}%
\pgfpathlineto{\pgfqpoint{0.656466in}{1.153746in}}%
\pgfpathlineto{\pgfqpoint{0.640558in}{1.173621in}}%
\pgfpathlineto{\pgfqpoint{0.631866in}{1.170053in}}%
\pgfpathlineto{\pgfqpoint{0.623136in}{1.166815in}}%
\pgfpathclose%
\pgfusepath{fill}%
\end{pgfscope}%
\begin{pgfscope}%
\pgfpathrectangle{\pgfqpoint{0.407592in}{0.222904in}}{\pgfqpoint{1.411163in}{1.411163in}}%
\pgfusepath{clip}%
\pgfsetbuttcap%
\pgfsetroundjoin%
\definecolor{currentfill}{rgb}{0.556863,0.654902,0.847059}%
\pgfsetfillcolor{currentfill}%
\pgfsetfillopacity{0.130000}%
\pgfsetlinewidth{0.000000pt}%
\definecolor{currentstroke}{rgb}{0.000000,0.000000,0.000000}%
\pgfsetstrokecolor{currentstroke}%
\pgfsetdash{}{0pt}%
\pgfpathmoveto{\pgfqpoint{1.169293in}{0.808200in}}%
\pgfpathlineto{\pgfqpoint{1.185453in}{0.807000in}}%
\pgfpathlineto{\pgfqpoint{1.201664in}{0.807384in}}%
\pgfpathlineto{\pgfqpoint{1.209513in}{0.819793in}}%
\pgfpathlineto{\pgfqpoint{1.217348in}{0.832744in}}%
\pgfpathlineto{\pgfqpoint{1.201151in}{0.832694in}}%
\pgfpathlineto{\pgfqpoint{1.185006in}{0.834104in}}%
\pgfpathlineto{\pgfqpoint{1.177157in}{0.820884in}}%
\pgfpathlineto{\pgfqpoint{1.169293in}{0.808200in}}%
\pgfpathclose%
\pgfusepath{fill}%
\end{pgfscope}%
\begin{pgfscope}%
\pgfpathrectangle{\pgfqpoint{0.407592in}{0.222904in}}{\pgfqpoint{1.411163in}{1.411163in}}%
\pgfusepath{clip}%
\pgfsetbuttcap%
\pgfsetroundjoin%
\definecolor{currentfill}{rgb}{0.556863,0.654902,0.847059}%
\pgfsetfillcolor{currentfill}%
\pgfsetfillopacity{0.130000}%
\pgfsetlinewidth{0.000000pt}%
\definecolor{currentstroke}{rgb}{0.000000,0.000000,0.000000}%
\pgfsetstrokecolor{currentstroke}%
\pgfsetdash{}{0pt}%
\pgfpathmoveto{\pgfqpoint{0.800025in}{0.977893in}}%
\pgfpathlineto{\pgfqpoint{0.816020in}{0.958917in}}%
\pgfpathlineto{\pgfqpoint{0.832024in}{0.940128in}}%
\pgfpathlineto{\pgfqpoint{0.840423in}{0.943959in}}%
\pgfpathlineto{\pgfqpoint{0.848786in}{0.948342in}}%
\pgfpathlineto{\pgfqpoint{0.832825in}{0.967154in}}%
\pgfpathlineto{\pgfqpoint{0.816873in}{0.986145in}}%
\pgfpathlineto{\pgfqpoint{0.808467in}{0.981770in}}%
\pgfpathlineto{\pgfqpoint{0.800025in}{0.977893in}}%
\pgfpathclose%
\pgfusepath{fill}%
\end{pgfscope}%
\begin{pgfscope}%
\pgfpathrectangle{\pgfqpoint{0.407592in}{0.222904in}}{\pgfqpoint{1.411163in}{1.411163in}}%
\pgfusepath{clip}%
\pgfsetbuttcap%
\pgfsetroundjoin%
\definecolor{currentfill}{rgb}{0.556863,0.654902,0.847059}%
\pgfsetfillcolor{currentfill}%
\pgfsetfillopacity{0.130000}%
\pgfsetlinewidth{0.000000pt}%
\definecolor{currentstroke}{rgb}{0.000000,0.000000,0.000000}%
\pgfsetstrokecolor{currentstroke}%
\pgfsetdash{}{0pt}%
\pgfpathmoveto{\pgfqpoint{1.392960in}{0.991706in}}%
\pgfpathlineto{\pgfqpoint{1.409513in}{0.996312in}}%
\pgfpathlineto{\pgfqpoint{1.426133in}{1.001563in}}%
\pgfpathlineto{\pgfqpoint{1.433770in}{1.015845in}}%
\pgfpathlineto{\pgfqpoint{1.441394in}{1.030305in}}%
\pgfpathlineto{\pgfqpoint{1.424789in}{1.025511in}}%
\pgfpathlineto{\pgfqpoint{1.408250in}{1.021348in}}%
\pgfpathlineto{\pgfqpoint{1.400611in}{1.006448in}}%
\pgfpathlineto{\pgfqpoint{1.392960in}{0.991706in}}%
\pgfpathclose%
\pgfusepath{fill}%
\end{pgfscope}%
\begin{pgfscope}%
\pgfpathrectangle{\pgfqpoint{0.407592in}{0.222904in}}{\pgfqpoint{1.411163in}{1.411163in}}%
\pgfusepath{clip}%
\pgfsetbuttcap%
\pgfsetroundjoin%
\definecolor{currentfill}{rgb}{0.556863,0.654902,0.847059}%
\pgfsetfillcolor{currentfill}%
\pgfsetfillopacity{0.130000}%
\pgfsetlinewidth{0.000000pt}%
\definecolor{currentstroke}{rgb}{0.000000,0.000000,0.000000}%
\pgfsetstrokecolor{currentstroke}%
\pgfsetdash{}{0pt}%
\pgfpathmoveto{\pgfqpoint{1.025135in}{0.793661in}}%
\pgfpathlineto{\pgfqpoint{1.041175in}{0.781692in}}%
\pgfpathlineto{\pgfqpoint{1.057232in}{0.771245in}}%
\pgfpathlineto{\pgfqpoint{1.065265in}{0.778684in}}%
\pgfpathlineto{\pgfqpoint{1.073272in}{0.787145in}}%
\pgfpathlineto{\pgfqpoint{1.057245in}{0.797326in}}%
\pgfpathlineto{\pgfqpoint{1.041239in}{0.808995in}}%
\pgfpathlineto{\pgfqpoint{1.033201in}{0.800859in}}%
\pgfpathlineto{\pgfqpoint{1.025135in}{0.793661in}}%
\pgfpathclose%
\pgfusepath{fill}%
\end{pgfscope}%
\begin{pgfscope}%
\pgfpathrectangle{\pgfqpoint{0.407592in}{0.222904in}}{\pgfqpoint{1.411163in}{1.411163in}}%
\pgfusepath{clip}%
\pgfsetbuttcap%
\pgfsetroundjoin%
\definecolor{currentfill}{rgb}{0.556863,0.654902,0.847059}%
\pgfsetfillcolor{currentfill}%
\pgfsetfillopacity{0.130000}%
\pgfsetlinewidth{0.000000pt}%
\definecolor{currentstroke}{rgb}{0.000000,0.000000,0.000000}%
\pgfsetstrokecolor{currentstroke}%
\pgfsetdash{}{0pt}%
\pgfpathmoveto{\pgfqpoint{0.655045in}{1.127152in}}%
\pgfpathlineto{\pgfqpoint{0.671013in}{1.107403in}}%
\pgfpathlineto{\pgfqpoint{0.686990in}{1.087719in}}%
\pgfpathlineto{\pgfqpoint{0.695634in}{1.090875in}}%
\pgfpathlineto{\pgfqpoint{0.704240in}{1.094413in}}%
\pgfpathlineto{\pgfqpoint{0.688307in}{1.114135in}}%
\pgfpathlineto{\pgfqpoint{0.672382in}{1.133915in}}%
\pgfpathlineto{\pgfqpoint{0.663732in}{1.130356in}}%
\pgfpathlineto{\pgfqpoint{0.655045in}{1.127152in}}%
\pgfpathclose%
\pgfusepath{fill}%
\end{pgfscope}%
\begin{pgfscope}%
\pgfpathrectangle{\pgfqpoint{0.407592in}{0.222904in}}{\pgfqpoint{1.411163in}{1.411163in}}%
\pgfusepath{clip}%
\pgfsetbuttcap%
\pgfsetroundjoin%
\definecolor{currentfill}{rgb}{0.556863,0.654902,0.847059}%
\pgfsetfillcolor{currentfill}%
\pgfsetfillopacity{0.130000}%
\pgfsetlinewidth{0.000000pt}%
\definecolor{currentstroke}{rgb}{0.000000,0.000000,0.000000}%
\pgfsetstrokecolor{currentstroke}%
\pgfsetdash{}{0pt}%
\pgfpathmoveto{\pgfqpoint{0.944742in}{0.841599in}}%
\pgfpathlineto{\pgfqpoint{0.960768in}{0.825461in}}%
\pgfpathlineto{\pgfqpoint{0.976805in}{0.810076in}}%
\pgfpathlineto{\pgfqpoint{0.984970in}{0.815243in}}%
\pgfpathlineto{\pgfqpoint{0.993102in}{0.821299in}}%
\pgfpathlineto{\pgfqpoint{0.977104in}{0.836576in}}%
\pgfpathlineto{\pgfqpoint{0.961118in}{0.852608in}}%
\pgfpathlineto{\pgfqpoint{0.952947in}{0.846712in}}%
\pgfpathlineto{\pgfqpoint{0.944742in}{0.841599in}}%
\pgfpathclose%
\pgfusepath{fill}%
\end{pgfscope}%
\begin{pgfscope}%
\pgfpathrectangle{\pgfqpoint{0.407592in}{0.222904in}}{\pgfqpoint{1.411163in}{1.411163in}}%
\pgfusepath{clip}%
\pgfsetbuttcap%
\pgfsetroundjoin%
\definecolor{currentfill}{rgb}{0.556863,0.654902,0.847059}%
\pgfsetfillcolor{currentfill}%
\pgfsetfillopacity{0.130000}%
\pgfsetlinewidth{0.000000pt}%
\definecolor{currentstroke}{rgb}{0.000000,0.000000,0.000000}%
\pgfsetstrokecolor{currentstroke}%
\pgfsetdash{}{0pt}%
\pgfpathmoveto{\pgfqpoint{1.217348in}{0.832744in}}%
\pgfpathlineto{\pgfqpoint{1.233599in}{0.834165in}}%
\pgfpathlineto{\pgfqpoint{1.249907in}{0.836852in}}%
\pgfpathlineto{\pgfqpoint{1.257714in}{0.849820in}}%
\pgfpathlineto{\pgfqpoint{1.265506in}{0.863214in}}%
\pgfpathlineto{\pgfqpoint{1.249211in}{0.861012in}}%
\pgfpathlineto{\pgfqpoint{1.232974in}{0.860001in}}%
\pgfpathlineto{\pgfqpoint{1.225168in}{0.846167in}}%
\pgfpathlineto{\pgfqpoint{1.217348in}{0.832744in}}%
\pgfpathclose%
\pgfusepath{fill}%
\end{pgfscope}%
\begin{pgfscope}%
\pgfpathrectangle{\pgfqpoint{0.407592in}{0.222904in}}{\pgfqpoint{1.411163in}{1.411163in}}%
\pgfusepath{clip}%
\pgfsetbuttcap%
\pgfsetroundjoin%
\definecolor{currentfill}{rgb}{0.556863,0.654902,0.847059}%
\pgfsetfillcolor{currentfill}%
\pgfsetfillopacity{0.130000}%
\pgfsetlinewidth{0.000000pt}%
\definecolor{currentstroke}{rgb}{0.000000,0.000000,0.000000}%
\pgfsetstrokecolor{currentstroke}%
\pgfsetdash{}{0pt}%
\pgfpathmoveto{\pgfqpoint{0.832024in}{0.940128in}}%
\pgfpathlineto{\pgfqpoint{0.848036in}{0.921557in}}%
\pgfpathlineto{\pgfqpoint{0.864057in}{0.903246in}}%
\pgfpathlineto{\pgfqpoint{0.872414in}{0.907010in}}%
\pgfpathlineto{\pgfqpoint{0.880735in}{0.911389in}}%
\pgfpathlineto{\pgfqpoint{0.864756in}{0.929741in}}%
\pgfpathlineto{\pgfqpoint{0.848786in}{0.948342in}}%
\pgfpathlineto{\pgfqpoint{0.840423in}{0.943959in}}%
\pgfpathlineto{\pgfqpoint{0.832024in}{0.940128in}}%
\pgfpathclose%
\pgfusepath{fill}%
\end{pgfscope}%
\begin{pgfscope}%
\pgfpathrectangle{\pgfqpoint{0.407592in}{0.222904in}}{\pgfqpoint{1.411163in}{1.411163in}}%
\pgfusepath{clip}%
\pgfsetbuttcap%
\pgfsetroundjoin%
\definecolor{currentfill}{rgb}{0.556863,0.654902,0.847059}%
\pgfsetfillcolor{currentfill}%
\pgfsetfillopacity{0.130000}%
\pgfsetlinewidth{0.000000pt}%
\definecolor{currentstroke}{rgb}{0.000000,0.000000,0.000000}%
\pgfsetstrokecolor{currentstroke}%
\pgfsetdash{}{0pt}%
\pgfpathmoveto{\pgfqpoint{1.329263in}{0.925705in}}%
\pgfpathlineto{\pgfqpoint{1.345716in}{0.929810in}}%
\pgfpathlineto{\pgfqpoint{1.362233in}{0.934699in}}%
\pgfpathlineto{\pgfqpoint{1.369933in}{0.948613in}}%
\pgfpathlineto{\pgfqpoint{1.377621in}{0.962768in}}%
\pgfpathlineto{\pgfqpoint{1.361118in}{0.958387in}}%
\pgfpathlineto{\pgfqpoint{1.344679in}{0.954769in}}%
\pgfpathlineto{\pgfqpoint{1.336977in}{0.940127in}}%
\pgfpathlineto{\pgfqpoint{1.329263in}{0.925705in}}%
\pgfpathclose%
\pgfusepath{fill}%
\end{pgfscope}%
\begin{pgfscope}%
\pgfpathrectangle{\pgfqpoint{0.407592in}{0.222904in}}{\pgfqpoint{1.411163in}{1.411163in}}%
\pgfusepath{clip}%
\pgfsetbuttcap%
\pgfsetroundjoin%
\definecolor{currentfill}{rgb}{0.556863,0.654902,0.847059}%
\pgfsetfillcolor{currentfill}%
\pgfsetfillopacity{0.130000}%
\pgfsetlinewidth{0.000000pt}%
\definecolor{currentstroke}{rgb}{0.000000,0.000000,0.000000}%
\pgfsetstrokecolor{currentstroke}%
\pgfsetdash{}{0pt}%
\pgfpathmoveto{\pgfqpoint{0.686990in}{1.087719in}}%
\pgfpathlineto{\pgfqpoint{0.702975in}{1.068110in}}%
\pgfpathlineto{\pgfqpoint{0.718969in}{1.048586in}}%
\pgfpathlineto{\pgfqpoint{0.727570in}{1.051676in}}%
\pgfpathlineto{\pgfqpoint{0.736134in}{1.055179in}}%
\pgfpathlineto{\pgfqpoint{0.720183in}{1.074758in}}%
\pgfpathlineto{\pgfqpoint{0.704240in}{1.094413in}}%
\pgfpathlineto{\pgfqpoint{0.695634in}{1.090875in}}%
\pgfpathlineto{\pgfqpoint{0.686990in}{1.087719in}}%
\pgfpathclose%
\pgfusepath{fill}%
\end{pgfscope}%
\begin{pgfscope}%
\pgfpathrectangle{\pgfqpoint{0.407592in}{0.222904in}}{\pgfqpoint{1.411163in}{1.411163in}}%
\pgfusepath{clip}%
\pgfsetbuttcap%
\pgfsetroundjoin%
\definecolor{currentfill}{rgb}{0.556863,0.654902,0.847059}%
\pgfsetfillcolor{currentfill}%
\pgfsetfillopacity{0.130000}%
\pgfsetlinewidth{0.000000pt}%
\definecolor{currentstroke}{rgb}{0.000000,0.000000,0.000000}%
\pgfsetstrokecolor{currentstroke}%
\pgfsetdash{}{0pt}%
\pgfpathmoveto{\pgfqpoint{1.505121in}{1.099259in}}%
\pgfpathlineto{\pgfqpoint{1.521896in}{1.104863in}}%
\pgfpathlineto{\pgfqpoint{1.538740in}{1.110948in}}%
\pgfpathlineto{\pgfqpoint{1.546274in}{1.125305in}}%
\pgfpathlineto{\pgfqpoint{1.529438in}{1.119417in}}%
\pgfpathlineto{\pgfqpoint{1.512671in}{1.114007in}}%
\pgfusepath{fill}%
\end{pgfscope}%
\begin{pgfscope}%
\pgfpathrectangle{\pgfqpoint{0.407592in}{0.222904in}}{\pgfqpoint{1.411163in}{1.411163in}}%
\pgfusepath{clip}%
\pgfsetbuttcap%
\pgfsetroundjoin%
\definecolor{currentfill}{rgb}{0.556863,0.654902,0.847059}%
\pgfsetfillcolor{currentfill}%
\pgfsetfillopacity{0.130000}%
\pgfsetlinewidth{0.000000pt}%
\definecolor{currentstroke}{rgb}{0.000000,0.000000,0.000000}%
\pgfsetstrokecolor{currentstroke}%
\pgfsetdash{}{0pt}%
\pgfpathmoveto{\pgfqpoint{0.864057in}{0.903246in}}%
\pgfpathlineto{\pgfqpoint{0.880087in}{0.885241in}}%
\pgfpathlineto{\pgfqpoint{0.896126in}{0.867602in}}%
\pgfpathlineto{\pgfqpoint{0.904440in}{0.871272in}}%
\pgfpathlineto{\pgfqpoint{0.912720in}{0.875633in}}%
\pgfpathlineto{\pgfqpoint{0.896723in}{0.893334in}}%
\pgfpathlineto{\pgfqpoint{0.880735in}{0.911389in}}%
\pgfpathlineto{\pgfqpoint{0.872414in}{0.907010in}}%
\pgfpathlineto{\pgfqpoint{0.864057in}{0.903246in}}%
\pgfpathclose%
\pgfusepath{fill}%
\end{pgfscope}%
\begin{pgfscope}%
\pgfpathrectangle{\pgfqpoint{0.407592in}{0.222904in}}{\pgfqpoint{1.411163in}{1.411163in}}%
\pgfusepath{clip}%
\pgfsetbuttcap%
\pgfsetroundjoin%
\definecolor{currentfill}{rgb}{0.556863,0.654902,0.847059}%
\pgfsetfillcolor{currentfill}%
\pgfsetfillopacity{0.130000}%
\pgfsetlinewidth{0.000000pt}%
\definecolor{currentstroke}{rgb}{0.000000,0.000000,0.000000}%
\pgfsetstrokecolor{currentstroke}%
\pgfsetdash{}{0pt}%
\pgfpathmoveto{\pgfqpoint{1.441394in}{1.030305in}}%
\pgfpathlineto{\pgfqpoint{1.458065in}{1.035690in}}%
\pgfpathlineto{\pgfqpoint{1.474805in}{1.041628in}}%
\pgfpathlineto{\pgfqpoint{1.482401in}{1.055806in}}%
\pgfpathlineto{\pgfqpoint{1.489986in}{1.070146in}}%
\pgfpathlineto{\pgfqpoint{1.473263in}{1.064645in}}%
\pgfpathlineto{\pgfqpoint{1.456607in}{1.059690in}}%
\pgfpathlineto{\pgfqpoint{1.449006in}{1.044924in}}%
\pgfpathlineto{\pgfqpoint{1.441394in}{1.030305in}}%
\pgfpathclose%
\pgfusepath{fill}%
\end{pgfscope}%
\begin{pgfscope}%
\pgfpathrectangle{\pgfqpoint{0.407592in}{0.222904in}}{\pgfqpoint{1.411163in}{1.411163in}}%
\pgfusepath{clip}%
\pgfsetbuttcap%
\pgfsetroundjoin%
\definecolor{currentfill}{rgb}{0.556863,0.654902,0.847059}%
\pgfsetfillcolor{currentfill}%
\pgfsetfillopacity{0.130000}%
\pgfsetlinewidth{0.000000pt}%
\definecolor{currentstroke}{rgb}{0.000000,0.000000,0.000000}%
\pgfsetstrokecolor{currentstroke}%
\pgfsetdash{}{0pt}%
\pgfpathmoveto{\pgfqpoint{1.265506in}{0.863214in}}%
\pgfpathlineto{\pgfqpoint{1.281862in}{0.866511in}}%
\pgfpathlineto{\pgfqpoint{1.298279in}{0.870807in}}%
\pgfpathlineto{\pgfqpoint{1.306044in}{0.884045in}}%
\pgfpathlineto{\pgfqpoint{1.313797in}{0.897632in}}%
\pgfpathlineto{\pgfqpoint{1.297393in}{0.893877in}}%
\pgfpathlineto{\pgfqpoint{1.281052in}{0.891084in}}%
\pgfpathlineto{\pgfqpoint{1.273286in}{0.876984in}}%
\pgfpathlineto{\pgfqpoint{1.265506in}{0.863214in}}%
\pgfpathclose%
\pgfusepath{fill}%
\end{pgfscope}%
\begin{pgfscope}%
\pgfpathrectangle{\pgfqpoint{0.407592in}{0.222904in}}{\pgfqpoint{1.411163in}{1.411163in}}%
\pgfusepath{clip}%
\pgfsetbuttcap%
\pgfsetroundjoin%
\definecolor{currentfill}{rgb}{0.556863,0.654902,0.847059}%
\pgfsetfillcolor{currentfill}%
\pgfsetfillopacity{0.130000}%
\pgfsetlinewidth{0.000000pt}%
\definecolor{currentstroke}{rgb}{0.000000,0.000000,0.000000}%
\pgfsetstrokecolor{currentstroke}%
\pgfsetdash{}{0pt}%
\pgfpathmoveto{\pgfqpoint{0.718969in}{1.048586in}}%
\pgfpathlineto{\pgfqpoint{0.734972in}{1.029159in}}%
\pgfpathlineto{\pgfqpoint{0.750983in}{1.009845in}}%
\pgfpathlineto{\pgfqpoint{0.759541in}{1.012848in}}%
\pgfpathlineto{\pgfqpoint{0.768062in}{1.016299in}}%
\pgfpathlineto{\pgfqpoint{0.752094in}{1.035688in}}%
\pgfpathlineto{\pgfqpoint{0.736134in}{1.055179in}}%
\pgfpathlineto{\pgfqpoint{0.727570in}{1.051676in}}%
\pgfpathlineto{\pgfqpoint{0.718969in}{1.048586in}}%
\pgfpathclose%
\pgfusepath{fill}%
\end{pgfscope}%
\begin{pgfscope}%
\pgfpathrectangle{\pgfqpoint{0.407592in}{0.222904in}}{\pgfqpoint{1.411163in}{1.411163in}}%
\pgfusepath{clip}%
\pgfsetbuttcap%
\pgfsetroundjoin%
\definecolor{currentfill}{rgb}{0.556863,0.654902,0.847059}%
\pgfsetfillcolor{currentfill}%
\pgfsetfillopacity{0.130000}%
\pgfsetlinewidth{0.000000pt}%
\definecolor{currentstroke}{rgb}{0.000000,0.000000,0.000000}%
\pgfsetstrokecolor{currentstroke}%
\pgfsetdash{}{0pt}%
\pgfpathmoveto{\pgfqpoint{0.976805in}{0.810076in}}%
\pgfpathlineto{\pgfqpoint{0.992853in}{0.795603in}}%
\pgfpathlineto{\pgfqpoint{1.008914in}{0.782227in}}%
\pgfpathlineto{\pgfqpoint{1.017040in}{0.787443in}}%
\pgfpathlineto{\pgfqpoint{1.025135in}{0.793661in}}%
\pgfpathlineto{\pgfqpoint{1.009112in}{0.806933in}}%
\pgfpathlineto{\pgfqpoint{0.993102in}{0.821299in}}%
\pgfpathlineto{\pgfqpoint{0.984970in}{0.815243in}}%
\pgfpathlineto{\pgfqpoint{0.976805in}{0.810076in}}%
\pgfpathclose%
\pgfusepath{fill}%
\end{pgfscope}%
\begin{pgfscope}%
\pgfpathrectangle{\pgfqpoint{0.407592in}{0.222904in}}{\pgfqpoint{1.411163in}{1.411163in}}%
\pgfusepath{clip}%
\pgfsetbuttcap%
\pgfsetroundjoin%
\definecolor{currentfill}{rgb}{0.556863,0.654902,0.847059}%
\pgfsetfillcolor{currentfill}%
\pgfsetfillopacity{0.130000}%
\pgfsetlinewidth{0.000000pt}%
\definecolor{currentstroke}{rgb}{0.000000,0.000000,0.000000}%
\pgfsetstrokecolor{currentstroke}%
\pgfsetdash{}{0pt}%
\pgfpathmoveto{\pgfqpoint{1.105400in}{0.771992in}}%
\pgfpathlineto{\pgfqpoint{1.121512in}{0.767285in}}%
\pgfpathlineto{\pgfqpoint{1.137663in}{0.764555in}}%
\pgfpathlineto{\pgfqpoint{1.145599in}{0.774216in}}%
\pgfpathlineto{\pgfqpoint{1.153515in}{0.784774in}}%
\pgfpathlineto{\pgfqpoint{1.137384in}{0.787564in}}%
\pgfpathlineto{\pgfqpoint{1.121293in}{0.792199in}}%
\pgfpathlineto{\pgfqpoint{1.113357in}{0.781646in}}%
\pgfpathlineto{\pgfqpoint{1.105400in}{0.771992in}}%
\pgfpathclose%
\pgfusepath{fill}%
\end{pgfscope}%
\begin{pgfscope}%
\pgfpathrectangle{\pgfqpoint{0.407592in}{0.222904in}}{\pgfqpoint{1.411163in}{1.411163in}}%
\pgfusepath{clip}%
\pgfsetbuttcap%
\pgfsetroundjoin%
\definecolor{currentfill}{rgb}{0.556863,0.654902,0.847059}%
\pgfsetfillcolor{currentfill}%
\pgfsetfillopacity{0.130000}%
\pgfsetlinewidth{0.000000pt}%
\definecolor{currentstroke}{rgb}{0.000000,0.000000,0.000000}%
\pgfsetstrokecolor{currentstroke}%
\pgfsetdash{}{0pt}%
\pgfpathmoveto{\pgfqpoint{1.153515in}{0.784774in}}%
\pgfpathlineto{\pgfqpoint{1.169691in}{0.783786in}}%
\pgfpathlineto{\pgfqpoint{1.185918in}{0.784499in}}%
\pgfpathlineto{\pgfqpoint{1.193799in}{0.795592in}}%
\pgfpathlineto{\pgfqpoint{1.201664in}{0.807384in}}%
\pgfpathlineto{\pgfqpoint{1.185453in}{0.807000in}}%
\pgfpathlineto{\pgfqpoint{1.169293in}{0.808200in}}%
\pgfpathlineto{\pgfqpoint{1.161412in}{0.796133in}}%
\pgfpathlineto{\pgfqpoint{1.153515in}{0.784774in}}%
\pgfpathclose%
\pgfusepath{fill}%
\end{pgfscope}%
\begin{pgfscope}%
\pgfpathrectangle{\pgfqpoint{0.407592in}{0.222904in}}{\pgfqpoint{1.411163in}{1.411163in}}%
\pgfusepath{clip}%
\pgfsetbuttcap%
\pgfsetroundjoin%
\definecolor{currentfill}{rgb}{0.556863,0.654902,0.847059}%
\pgfsetfillcolor{currentfill}%
\pgfsetfillopacity{0.130000}%
\pgfsetlinewidth{0.000000pt}%
\definecolor{currentstroke}{rgb}{0.000000,0.000000,0.000000}%
\pgfsetstrokecolor{currentstroke}%
\pgfsetdash{}{0pt}%
\pgfpathmoveto{\pgfqpoint{1.057232in}{0.771245in}}%
\pgfpathlineto{\pgfqpoint{1.073312in}{0.762532in}}%
\pgfpathlineto{\pgfqpoint{1.089418in}{0.755736in}}%
\pgfpathlineto{\pgfqpoint{1.097421in}{0.763328in}}%
\pgfpathlineto{\pgfqpoint{1.105400in}{0.771992in}}%
\pgfpathlineto{\pgfqpoint{1.089322in}{0.778647in}}%
\pgfpathlineto{\pgfqpoint{1.073272in}{0.787145in}}%
\pgfpathlineto{\pgfqpoint{1.065265in}{0.778684in}}%
\pgfpathlineto{\pgfqpoint{1.057232in}{0.771245in}}%
\pgfpathclose%
\pgfusepath{fill}%
\end{pgfscope}%
\begin{pgfscope}%
\pgfpathrectangle{\pgfqpoint{0.407592in}{0.222904in}}{\pgfqpoint{1.411163in}{1.411163in}}%
\pgfusepath{clip}%
\pgfsetbuttcap%
\pgfsetroundjoin%
\definecolor{currentfill}{rgb}{0.556863,0.654902,0.847059}%
\pgfsetfillcolor{currentfill}%
\pgfsetfillopacity{0.130000}%
\pgfsetlinewidth{0.000000pt}%
\definecolor{currentstroke}{rgb}{0.000000,0.000000,0.000000}%
\pgfsetstrokecolor{currentstroke}%
\pgfsetdash{}{0pt}%
\pgfpathmoveto{\pgfqpoint{1.377621in}{0.962768in}}%
\pgfpathlineto{\pgfqpoint{1.394190in}{0.967857in}}%
\pgfpathlineto{\pgfqpoint{1.410825in}{0.973598in}}%
\pgfpathlineto{\pgfqpoint{1.418485in}{0.987474in}}%
\pgfpathlineto{\pgfqpoint{1.426133in}{1.001563in}}%
\pgfpathlineto{\pgfqpoint{1.409513in}{0.996312in}}%
\pgfpathlineto{\pgfqpoint{1.392960in}{0.991706in}}%
\pgfpathlineto{\pgfqpoint{1.385296in}{0.977140in}}%
\pgfpathlineto{\pgfqpoint{1.377621in}{0.962768in}}%
\pgfpathclose%
\pgfusepath{fill}%
\end{pgfscope}%
\begin{pgfscope}%
\pgfpathrectangle{\pgfqpoint{0.407592in}{0.222904in}}{\pgfqpoint{1.411163in}{1.411163in}}%
\pgfusepath{clip}%
\pgfsetbuttcap%
\pgfsetroundjoin%
\definecolor{currentfill}{rgb}{0.556863,0.654902,0.847059}%
\pgfsetfillcolor{currentfill}%
\pgfsetfillopacity{0.130000}%
\pgfsetlinewidth{0.000000pt}%
\definecolor{currentstroke}{rgb}{0.000000,0.000000,0.000000}%
\pgfsetstrokecolor{currentstroke}%
\pgfsetdash{}{0pt}%
\pgfpathmoveto{\pgfqpoint{0.750983in}{1.009845in}}%
\pgfpathlineto{\pgfqpoint{0.767003in}{0.990659in}}%
\pgfpathlineto{\pgfqpoint{0.783032in}{0.971623in}}%
\pgfpathlineto{\pgfqpoint{0.791547in}{0.974513in}}%
\pgfpathlineto{\pgfqpoint{0.800025in}{0.977893in}}%
\pgfpathlineto{\pgfqpoint{0.784039in}{0.997028in}}%
\pgfpathlineto{\pgfqpoint{0.768062in}{1.016299in}}%
\pgfpathlineto{\pgfqpoint{0.759541in}{1.012848in}}%
\pgfpathlineto{\pgfqpoint{0.750983in}{1.009845in}}%
\pgfpathclose%
\pgfusepath{fill}%
\end{pgfscope}%
\begin{pgfscope}%
\pgfpathrectangle{\pgfqpoint{0.407592in}{0.222904in}}{\pgfqpoint{1.411163in}{1.411163in}}%
\pgfusepath{clip}%
\pgfsetbuttcap%
\pgfsetroundjoin%
\definecolor{currentfill}{rgb}{0.556863,0.654902,0.847059}%
\pgfsetfillcolor{currentfill}%
\pgfsetfillopacity{0.130000}%
\pgfsetlinewidth{0.000000pt}%
\definecolor{currentstroke}{rgb}{0.000000,0.000000,0.000000}%
\pgfsetstrokecolor{currentstroke}%
\pgfsetdash{}{0pt}%
\pgfpathmoveto{\pgfqpoint{0.896126in}{0.867602in}}%
\pgfpathlineto{\pgfqpoint{0.912173in}{0.850400in}}%
\pgfpathlineto{\pgfqpoint{0.928231in}{0.833726in}}%
\pgfpathlineto{\pgfqpoint{0.936504in}{0.837272in}}%
\pgfpathlineto{\pgfqpoint{0.944742in}{0.841599in}}%
\pgfpathlineto{\pgfqpoint{0.928726in}{0.858359in}}%
\pgfpathlineto{\pgfqpoint{0.912720in}{0.875633in}}%
\pgfpathlineto{\pgfqpoint{0.904440in}{0.871272in}}%
\pgfpathlineto{\pgfqpoint{0.896126in}{0.867602in}}%
\pgfpathclose%
\pgfusepath{fill}%
\end{pgfscope}%
\begin{pgfscope}%
\pgfpathrectangle{\pgfqpoint{0.407592in}{0.222904in}}{\pgfqpoint{1.411163in}{1.411163in}}%
\pgfusepath{clip}%
\pgfsetbuttcap%
\pgfsetroundjoin%
\definecolor{currentfill}{rgb}{0.556863,0.654902,0.847059}%
\pgfsetfillcolor{currentfill}%
\pgfsetfillopacity{0.130000}%
\pgfsetlinewidth{0.000000pt}%
\definecolor{currentstroke}{rgb}{0.000000,0.000000,0.000000}%
\pgfsetstrokecolor{currentstroke}%
\pgfsetdash{}{0pt}%
\pgfpathmoveto{\pgfqpoint{1.201664in}{0.807384in}}%
\pgfpathlineto{\pgfqpoint{1.217929in}{0.809237in}}%
\pgfpathlineto{\pgfqpoint{1.234253in}{0.812426in}}%
\pgfpathlineto{\pgfqpoint{1.242087in}{0.824367in}}%
\pgfpathlineto{\pgfqpoint{1.249907in}{0.836852in}}%
\pgfpathlineto{\pgfqpoint{1.233599in}{0.834165in}}%
\pgfpathlineto{\pgfqpoint{1.217348in}{0.832744in}}%
\pgfpathlineto{\pgfqpoint{1.209513in}{0.819793in}}%
\pgfpathlineto{\pgfqpoint{1.201664in}{0.807384in}}%
\pgfpathclose%
\pgfusepath{fill}%
\end{pgfscope}%
\begin{pgfscope}%
\pgfpathrectangle{\pgfqpoint{0.407592in}{0.222904in}}{\pgfqpoint{1.411163in}{1.411163in}}%
\pgfusepath{clip}%
\pgfsetbuttcap%
\pgfsetroundjoin%
\definecolor{currentfill}{rgb}{0.556863,0.654902,0.847059}%
\pgfsetfillcolor{currentfill}%
\pgfsetfillopacity{0.130000}%
\pgfsetlinewidth{0.000000pt}%
\definecolor{currentstroke}{rgb}{0.000000,0.000000,0.000000}%
\pgfsetstrokecolor{currentstroke}%
\pgfsetdash{}{0pt}%
\pgfpathmoveto{\pgfqpoint{1.313797in}{0.897632in}}%
\pgfpathlineto{\pgfqpoint{1.330264in}{0.902265in}}%
\pgfpathlineto{\pgfqpoint{1.346797in}{0.907699in}}%
\pgfpathlineto{\pgfqpoint{1.354521in}{0.921051in}}%
\pgfpathlineto{\pgfqpoint{1.362233in}{0.934699in}}%
\pgfpathlineto{\pgfqpoint{1.345716in}{0.929810in}}%
\pgfpathlineto{\pgfqpoint{1.329263in}{0.925705in}}%
\pgfpathlineto{\pgfqpoint{1.321536in}{0.911530in}}%
\pgfpathlineto{\pgfqpoint{1.313797in}{0.897632in}}%
\pgfpathclose%
\pgfusepath{fill}%
\end{pgfscope}%
\begin{pgfscope}%
\pgfpathrectangle{\pgfqpoint{0.407592in}{0.222904in}}{\pgfqpoint{1.411163in}{1.411163in}}%
\pgfusepath{clip}%
\pgfsetbuttcap%
\pgfsetroundjoin%
\definecolor{currentfill}{rgb}{0.556863,0.654902,0.847059}%
\pgfsetfillcolor{currentfill}%
\pgfsetfillopacity{0.130000}%
\pgfsetlinewidth{0.000000pt}%
\definecolor{currentstroke}{rgb}{0.000000,0.000000,0.000000}%
\pgfsetstrokecolor{currentstroke}%
\pgfsetdash{}{0pt}%
\pgfpathmoveto{\pgfqpoint{0.605563in}{1.161323in}}%
\pgfpathlineto{\pgfqpoint{0.621556in}{1.141530in}}%
\pgfpathlineto{\pgfqpoint{0.637558in}{1.121797in}}%
\pgfpathlineto{\pgfqpoint{0.646321in}{1.124300in}}%
\pgfpathlineto{\pgfqpoint{0.655045in}{1.127152in}}%
\pgfpathlineto{\pgfqpoint{0.639086in}{1.146958in}}%
\pgfpathlineto{\pgfqpoint{0.623136in}{1.166815in}}%
\pgfpathlineto{\pgfqpoint{0.614368in}{1.163906in}}%
\pgfpathlineto{\pgfqpoint{0.605563in}{1.161323in}}%
\pgfpathclose%
\pgfusepath{fill}%
\end{pgfscope}%
\begin{pgfscope}%
\pgfpathrectangle{\pgfqpoint{0.407592in}{0.222904in}}{\pgfqpoint{1.411163in}{1.411163in}}%
\pgfusepath{clip}%
\pgfsetbuttcap%
\pgfsetroundjoin%
\definecolor{currentfill}{rgb}{0.556863,0.654902,0.847059}%
\pgfsetfillcolor{currentfill}%
\pgfsetfillopacity{0.130000}%
\pgfsetlinewidth{0.000000pt}%
\definecolor{currentstroke}{rgb}{0.000000,0.000000,0.000000}%
\pgfsetstrokecolor{currentstroke}%
\pgfsetdash{}{0pt}%
\pgfpathmoveto{\pgfqpoint{1.489986in}{1.070146in}}%
\pgfpathlineto{\pgfqpoint{1.506778in}{1.076162in}}%
\pgfpathlineto{\pgfqpoint{1.523638in}{1.082659in}}%
\pgfpathlineto{\pgfqpoint{1.531195in}{1.096729in}}%
\pgfpathlineto{\pgfqpoint{1.538740in}{1.110948in}}%
\pgfpathlineto{\pgfqpoint{1.521896in}{1.104863in}}%
\pgfpathlineto{\pgfqpoint{1.505121in}{1.099259in}}%
\pgfpathlineto{\pgfqpoint{1.497560in}{1.084635in}}%
\pgfpathlineto{\pgfqpoint{1.489986in}{1.070146in}}%
\pgfpathclose%
\pgfusepath{fill}%
\end{pgfscope}%
\begin{pgfscope}%
\pgfpathrectangle{\pgfqpoint{0.407592in}{0.222904in}}{\pgfqpoint{1.411163in}{1.411163in}}%
\pgfusepath{clip}%
\pgfsetbuttcap%
\pgfsetroundjoin%
\definecolor{currentfill}{rgb}{0.556863,0.654902,0.847059}%
\pgfsetfillcolor{currentfill}%
\pgfsetfillopacity{0.130000}%
\pgfsetlinewidth{0.000000pt}%
\definecolor{currentstroke}{rgb}{0.000000,0.000000,0.000000}%
\pgfsetstrokecolor{currentstroke}%
\pgfsetdash{}{0pt}%
\pgfpathmoveto{\pgfqpoint{0.783032in}{0.971623in}}%
\pgfpathlineto{\pgfqpoint{0.799069in}{0.952761in}}%
\pgfpathlineto{\pgfqpoint{0.815115in}{0.934101in}}%
\pgfpathlineto{\pgfqpoint{0.823588in}{0.936844in}}%
\pgfpathlineto{\pgfqpoint{0.832024in}{0.940128in}}%
\pgfpathlineto{\pgfqpoint{0.816020in}{0.958917in}}%
\pgfpathlineto{\pgfqpoint{0.800025in}{0.977893in}}%
\pgfpathlineto{\pgfqpoint{0.791547in}{0.974513in}}%
\pgfpathlineto{\pgfqpoint{0.783032in}{0.971623in}}%
\pgfpathclose%
\pgfusepath{fill}%
\end{pgfscope}%
\begin{pgfscope}%
\pgfpathrectangle{\pgfqpoint{0.407592in}{0.222904in}}{\pgfqpoint{1.411163in}{1.411163in}}%
\pgfusepath{clip}%
\pgfsetbuttcap%
\pgfsetroundjoin%
\definecolor{currentfill}{rgb}{0.556863,0.654902,0.847059}%
\pgfsetfillcolor{currentfill}%
\pgfsetfillopacity{0.130000}%
\pgfsetlinewidth{0.000000pt}%
\definecolor{currentstroke}{rgb}{0.000000,0.000000,0.000000}%
\pgfsetstrokecolor{currentstroke}%
\pgfsetdash{}{0pt}%
\pgfpathmoveto{\pgfqpoint{1.008914in}{0.782227in}}%
\pgfpathlineto{\pgfqpoint{1.024989in}{0.770158in}}%
\pgfpathlineto{\pgfqpoint{1.041082in}{0.759620in}}%
\pgfpathlineto{\pgfqpoint{1.049172in}{0.764880in}}%
\pgfpathlineto{\pgfqpoint{1.057232in}{0.771245in}}%
\pgfpathlineto{\pgfqpoint{1.041175in}{0.781692in}}%
\pgfpathlineto{\pgfqpoint{1.025135in}{0.793661in}}%
\pgfpathlineto{\pgfqpoint{1.017040in}{0.787443in}}%
\pgfpathlineto{\pgfqpoint{1.008914in}{0.782227in}}%
\pgfpathclose%
\pgfusepath{fill}%
\end{pgfscope}%
\begin{pgfscope}%
\pgfpathrectangle{\pgfqpoint{0.407592in}{0.222904in}}{\pgfqpoint{1.411163in}{1.411163in}}%
\pgfusepath{clip}%
\pgfsetbuttcap%
\pgfsetroundjoin%
\definecolor{currentfill}{rgb}{0.556863,0.654902,0.847059}%
\pgfsetfillcolor{currentfill}%
\pgfsetfillopacity{0.130000}%
\pgfsetlinewidth{0.000000pt}%
\definecolor{currentstroke}{rgb}{0.000000,0.000000,0.000000}%
\pgfsetstrokecolor{currentstroke}%
\pgfsetdash{}{0pt}%
\pgfpathmoveto{\pgfqpoint{0.637558in}{1.121797in}}%
\pgfpathlineto{\pgfqpoint{0.653569in}{1.102130in}}%
\pgfpathlineto{\pgfqpoint{0.669588in}{1.082539in}}%
\pgfpathlineto{\pgfqpoint{0.678308in}{1.084942in}}%
\pgfpathlineto{\pgfqpoint{0.686990in}{1.087719in}}%
\pgfpathlineto{\pgfqpoint{0.671013in}{1.107403in}}%
\pgfpathlineto{\pgfqpoint{0.655045in}{1.127152in}}%
\pgfpathlineto{\pgfqpoint{0.646321in}{1.124300in}}%
\pgfpathlineto{\pgfqpoint{0.637558in}{1.121797in}}%
\pgfpathclose%
\pgfusepath{fill}%
\end{pgfscope}%
\begin{pgfscope}%
\pgfpathrectangle{\pgfqpoint{0.407592in}{0.222904in}}{\pgfqpoint{1.411163in}{1.411163in}}%
\pgfusepath{clip}%
\pgfsetbuttcap%
\pgfsetroundjoin%
\definecolor{currentfill}{rgb}{0.556863,0.654902,0.847059}%
\pgfsetfillcolor{currentfill}%
\pgfsetfillopacity{0.130000}%
\pgfsetlinewidth{0.000000pt}%
\definecolor{currentstroke}{rgb}{0.000000,0.000000,0.000000}%
\pgfsetstrokecolor{currentstroke}%
\pgfsetdash{}{0pt}%
\pgfpathmoveto{\pgfqpoint{0.928231in}{0.833726in}}%
\pgfpathlineto{\pgfqpoint{0.944298in}{0.817685in}}%
\pgfpathlineto{\pgfqpoint{0.960375in}{0.802408in}}%
\pgfpathlineto{\pgfqpoint{0.968607in}{0.805801in}}%
\pgfpathlineto{\pgfqpoint{0.976805in}{0.810076in}}%
\pgfpathlineto{\pgfqpoint{0.960768in}{0.825461in}}%
\pgfpathlineto{\pgfqpoint{0.944742in}{0.841599in}}%
\pgfpathlineto{\pgfqpoint{0.936504in}{0.837272in}}%
\pgfpathlineto{\pgfqpoint{0.928231in}{0.833726in}}%
\pgfpathclose%
\pgfusepath{fill}%
\end{pgfscope}%
\begin{pgfscope}%
\pgfpathrectangle{\pgfqpoint{0.407592in}{0.222904in}}{\pgfqpoint{1.411163in}{1.411163in}}%
\pgfusepath{clip}%
\pgfsetbuttcap%
\pgfsetroundjoin%
\definecolor{currentfill}{rgb}{0.556863,0.654902,0.847059}%
\pgfsetfillcolor{currentfill}%
\pgfsetfillopacity{0.130000}%
\pgfsetlinewidth{0.000000pt}%
\definecolor{currentstroke}{rgb}{0.000000,0.000000,0.000000}%
\pgfsetstrokecolor{currentstroke}%
\pgfsetdash{}{0pt}%
\pgfpathmoveto{\pgfqpoint{1.426133in}{1.001563in}}%
\pgfpathlineto{\pgfqpoint{1.442821in}{1.007411in}}%
\pgfpathlineto{\pgfqpoint{1.459576in}{1.013815in}}%
\pgfpathlineto{\pgfqpoint{1.467196in}{1.027626in}}%
\pgfpathlineto{\pgfqpoint{1.474805in}{1.041628in}}%
\pgfpathlineto{\pgfqpoint{1.458065in}{1.035690in}}%
\pgfpathlineto{\pgfqpoint{1.441394in}{1.030305in}}%
\pgfpathlineto{\pgfqpoint{1.433770in}{1.015845in}}%
\pgfpathlineto{\pgfqpoint{1.426133in}{1.001563in}}%
\pgfpathclose%
\pgfusepath{fill}%
\end{pgfscope}%
\begin{pgfscope}%
\pgfpathrectangle{\pgfqpoint{0.407592in}{0.222904in}}{\pgfqpoint{1.411163in}{1.411163in}}%
\pgfusepath{clip}%
\pgfsetbuttcap%
\pgfsetroundjoin%
\definecolor{currentfill}{rgb}{0.556863,0.654902,0.847059}%
\pgfsetfillcolor{currentfill}%
\pgfsetfillopacity{0.130000}%
\pgfsetlinewidth{0.000000pt}%
\definecolor{currentstroke}{rgb}{0.000000,0.000000,0.000000}%
\pgfsetstrokecolor{currentstroke}%
\pgfsetdash{}{0pt}%
\pgfpathmoveto{\pgfqpoint{1.249907in}{0.836852in}}%
\pgfpathlineto{\pgfqpoint{1.266277in}{0.840686in}}%
\pgfpathlineto{\pgfqpoint{1.282710in}{0.845550in}}%
\pgfpathlineto{\pgfqpoint{1.290501in}{0.857961in}}%
\pgfpathlineto{\pgfqpoint{1.298279in}{0.870807in}}%
\pgfpathlineto{\pgfqpoint{1.281862in}{0.866511in}}%
\pgfpathlineto{\pgfqpoint{1.265506in}{0.863214in}}%
\pgfpathlineto{\pgfqpoint{1.257714in}{0.849820in}}%
\pgfpathlineto{\pgfqpoint{1.249907in}{0.836852in}}%
\pgfpathclose%
\pgfusepath{fill}%
\end{pgfscope}%
\begin{pgfscope}%
\pgfpathrectangle{\pgfqpoint{0.407592in}{0.222904in}}{\pgfqpoint{1.411163in}{1.411163in}}%
\pgfusepath{clip}%
\pgfsetbuttcap%
\pgfsetroundjoin%
\definecolor{currentfill}{rgb}{0.556863,0.654902,0.847059}%
\pgfsetfillcolor{currentfill}%
\pgfsetfillopacity{0.130000}%
\pgfsetlinewidth{0.000000pt}%
\definecolor{currentstroke}{rgb}{0.000000,0.000000,0.000000}%
\pgfsetstrokecolor{currentstroke}%
\pgfsetdash{}{0pt}%
\pgfpathmoveto{\pgfqpoint{0.815115in}{0.934101in}}%
\pgfpathlineto{\pgfqpoint{0.831170in}{0.915678in}}%
\pgfpathlineto{\pgfqpoint{0.847233in}{0.897533in}}%
\pgfpathlineto{\pgfqpoint{0.855664in}{0.900090in}}%
\pgfpathlineto{\pgfqpoint{0.864057in}{0.903246in}}%
\pgfpathlineto{\pgfqpoint{0.848036in}{0.921557in}}%
\pgfpathlineto{\pgfqpoint{0.832024in}{0.940128in}}%
\pgfpathlineto{\pgfqpoint{0.823588in}{0.936844in}}%
\pgfpathlineto{\pgfqpoint{0.815115in}{0.934101in}}%
\pgfpathclose%
\pgfusepath{fill}%
\end{pgfscope}%
\begin{pgfscope}%
\pgfpathrectangle{\pgfqpoint{0.407592in}{0.222904in}}{\pgfqpoint{1.411163in}{1.411163in}}%
\pgfusepath{clip}%
\pgfsetbuttcap%
\pgfsetroundjoin%
\definecolor{currentfill}{rgb}{0.556863,0.654902,0.847059}%
\pgfsetfillcolor{currentfill}%
\pgfsetfillopacity{0.130000}%
\pgfsetlinewidth{0.000000pt}%
\definecolor{currentstroke}{rgb}{0.000000,0.000000,0.000000}%
\pgfsetstrokecolor{currentstroke}%
\pgfsetdash{}{0pt}%
\pgfpathmoveto{\pgfqpoint{1.362233in}{0.934699in}}%
\pgfpathlineto{\pgfqpoint{1.378817in}{0.940307in}}%
\pgfpathlineto{\pgfqpoint{1.395469in}{0.946572in}}%
\pgfpathlineto{\pgfqpoint{1.403153in}{0.959957in}}%
\pgfpathlineto{\pgfqpoint{1.410825in}{0.973598in}}%
\pgfpathlineto{\pgfqpoint{1.394190in}{0.967857in}}%
\pgfpathlineto{\pgfqpoint{1.377621in}{0.962768in}}%
\pgfpathlineto{\pgfqpoint{1.369933in}{0.948613in}}%
\pgfpathlineto{\pgfqpoint{1.362233in}{0.934699in}}%
\pgfpathclose%
\pgfusepath{fill}%
\end{pgfscope}%
\begin{pgfscope}%
\pgfpathrectangle{\pgfqpoint{0.407592in}{0.222904in}}{\pgfqpoint{1.411163in}{1.411163in}}%
\pgfusepath{clip}%
\pgfsetbuttcap%
\pgfsetroundjoin%
\definecolor{currentfill}{rgb}{0.556863,0.654902,0.847059}%
\pgfsetfillcolor{currentfill}%
\pgfsetfillopacity{0.130000}%
\pgfsetlinewidth{0.000000pt}%
\definecolor{currentstroke}{rgb}{0.000000,0.000000,0.000000}%
\pgfsetstrokecolor{currentstroke}%
\pgfsetdash{}{0pt}%
\pgfpathmoveto{\pgfqpoint{0.669588in}{1.082539in}}%
\pgfpathlineto{\pgfqpoint{0.685616in}{1.063034in}}%
\pgfpathlineto{\pgfqpoint{0.701653in}{1.043626in}}%
\pgfpathlineto{\pgfqpoint{0.710330in}{1.045904in}}%
\pgfpathlineto{\pgfqpoint{0.718969in}{1.048586in}}%
\pgfpathlineto{\pgfqpoint{0.702975in}{1.068110in}}%
\pgfpathlineto{\pgfqpoint{0.686990in}{1.087719in}}%
\pgfpathlineto{\pgfqpoint{0.678308in}{1.084942in}}%
\pgfpathlineto{\pgfqpoint{0.669588in}{1.082539in}}%
\pgfpathclose%
\pgfusepath{fill}%
\end{pgfscope}%
\begin{pgfscope}%
\pgfpathrectangle{\pgfqpoint{0.407592in}{0.222904in}}{\pgfqpoint{1.411163in}{1.411163in}}%
\pgfusepath{clip}%
\pgfsetbuttcap%
\pgfsetroundjoin%
\definecolor{currentfill}{rgb}{0.556863,0.654902,0.847059}%
\pgfsetfillcolor{currentfill}%
\pgfsetfillopacity{0.130000}%
\pgfsetlinewidth{0.000000pt}%
\definecolor{currentstroke}{rgb}{0.000000,0.000000,0.000000}%
\pgfsetstrokecolor{currentstroke}%
\pgfsetdash{}{0pt}%
\pgfpathmoveto{\pgfqpoint{1.089418in}{0.755736in}}%
\pgfpathlineto{\pgfqpoint{1.105555in}{0.750973in}}%
\pgfpathlineto{\pgfqpoint{1.121729in}{0.748276in}}%
\pgfpathlineto{\pgfqpoint{1.129707in}{0.755882in}}%
\pgfpathlineto{\pgfqpoint{1.137663in}{0.764555in}}%
\pgfpathlineto{\pgfqpoint{1.121512in}{0.767285in}}%
\pgfpathlineto{\pgfqpoint{1.105400in}{0.771992in}}%
\pgfpathlineto{\pgfqpoint{1.097421in}{0.763328in}}%
\pgfpathlineto{\pgfqpoint{1.089418in}{0.755736in}}%
\pgfpathclose%
\pgfusepath{fill}%
\end{pgfscope}%
\begin{pgfscope}%
\pgfpathrectangle{\pgfqpoint{0.407592in}{0.222904in}}{\pgfqpoint{1.411163in}{1.411163in}}%
\pgfusepath{clip}%
\pgfsetbuttcap%
\pgfsetroundjoin%
\definecolor{currentfill}{rgb}{0.556863,0.654902,0.847059}%
\pgfsetfillcolor{currentfill}%
\pgfsetfillopacity{0.130000}%
\pgfsetlinewidth{0.000000pt}%
\definecolor{currentstroke}{rgb}{0.000000,0.000000,0.000000}%
\pgfsetstrokecolor{currentstroke}%
\pgfsetdash{}{0pt}%
\pgfpathmoveto{\pgfqpoint{1.137663in}{0.764555in}}%
\pgfpathlineto{\pgfqpoint{1.153859in}{0.763749in}}%
\pgfpathlineto{\pgfqpoint{1.170103in}{0.764745in}}%
\pgfpathlineto{\pgfqpoint{1.178019in}{0.774189in}}%
\pgfpathlineto{\pgfqpoint{1.185918in}{0.784499in}}%
\pgfpathlineto{\pgfqpoint{1.169691in}{0.783786in}}%
\pgfpathlineto{\pgfqpoint{1.153515in}{0.784774in}}%
\pgfpathlineto{\pgfqpoint{1.145599in}{0.774216in}}%
\pgfpathlineto{\pgfqpoint{1.137663in}{0.764555in}}%
\pgfpathclose%
\pgfusepath{fill}%
\end{pgfscope}%
\begin{pgfscope}%
\pgfpathrectangle{\pgfqpoint{0.407592in}{0.222904in}}{\pgfqpoint{1.411163in}{1.411163in}}%
\pgfusepath{clip}%
\pgfsetbuttcap%
\pgfsetroundjoin%
\definecolor{currentfill}{rgb}{0.556863,0.654902,0.847059}%
\pgfsetfillcolor{currentfill}%
\pgfsetfillopacity{0.130000}%
\pgfsetlinewidth{0.000000pt}%
\definecolor{currentstroke}{rgb}{0.000000,0.000000,0.000000}%
\pgfsetstrokecolor{currentstroke}%
\pgfsetdash{}{0pt}%
\pgfpathmoveto{\pgfqpoint{0.960375in}{0.802408in}}%
\pgfpathlineto{\pgfqpoint{0.976463in}{0.788051in}}%
\pgfpathlineto{\pgfqpoint{0.992563in}{0.774795in}}%
\pgfpathlineto{\pgfqpoint{1.000755in}{0.778016in}}%
\pgfpathlineto{\pgfqpoint{1.008914in}{0.782227in}}%
\pgfpathlineto{\pgfqpoint{0.992853in}{0.795603in}}%
\pgfpathlineto{\pgfqpoint{0.976805in}{0.810076in}}%
\pgfpathlineto{\pgfqpoint{0.968607in}{0.805801in}}%
\pgfpathlineto{\pgfqpoint{0.960375in}{0.802408in}}%
\pgfpathclose%
\pgfusepath{fill}%
\end{pgfscope}%
\begin{pgfscope}%
\pgfpathrectangle{\pgfqpoint{0.407592in}{0.222904in}}{\pgfqpoint{1.411163in}{1.411163in}}%
\pgfusepath{clip}%
\pgfsetbuttcap%
\pgfsetroundjoin%
\definecolor{currentfill}{rgb}{0.556863,0.654902,0.847059}%
\pgfsetfillcolor{currentfill}%
\pgfsetfillopacity{0.130000}%
\pgfsetlinewidth{0.000000pt}%
\definecolor{currentstroke}{rgb}{0.000000,0.000000,0.000000}%
\pgfsetstrokecolor{currentstroke}%
\pgfsetdash{}{0pt}%
\pgfpathmoveto{\pgfqpoint{0.847233in}{0.897533in}}%
\pgfpathlineto{\pgfqpoint{0.863305in}{0.879715in}}%
\pgfpathlineto{\pgfqpoint{0.879386in}{0.862286in}}%
\pgfpathlineto{\pgfqpoint{0.887774in}{0.864612in}}%
\pgfpathlineto{\pgfqpoint{0.896126in}{0.867602in}}%
\pgfpathlineto{\pgfqpoint{0.880087in}{0.885241in}}%
\pgfpathlineto{\pgfqpoint{0.864057in}{0.903246in}}%
\pgfpathlineto{\pgfqpoint{0.855664in}{0.900090in}}%
\pgfpathlineto{\pgfqpoint{0.847233in}{0.897533in}}%
\pgfpathclose%
\pgfusepath{fill}%
\end{pgfscope}%
\begin{pgfscope}%
\pgfpathrectangle{\pgfqpoint{0.407592in}{0.222904in}}{\pgfqpoint{1.411163in}{1.411163in}}%
\pgfusepath{clip}%
\pgfsetbuttcap%
\pgfsetroundjoin%
\definecolor{currentfill}{rgb}{0.556863,0.654902,0.847059}%
\pgfsetfillcolor{currentfill}%
\pgfsetfillopacity{0.130000}%
\pgfsetlinewidth{0.000000pt}%
\definecolor{currentstroke}{rgb}{0.000000,0.000000,0.000000}%
\pgfsetstrokecolor{currentstroke}%
\pgfsetdash{}{0pt}%
\pgfpathmoveto{\pgfqpoint{0.701653in}{1.043626in}}%
\pgfpathlineto{\pgfqpoint{0.717698in}{1.024330in}}%
\pgfpathlineto{\pgfqpoint{0.733752in}{1.005162in}}%
\pgfpathlineto{\pgfqpoint{0.742387in}{1.007285in}}%
\pgfpathlineto{\pgfqpoint{0.750983in}{1.009845in}}%
\pgfpathlineto{\pgfqpoint{0.734972in}{1.029159in}}%
\pgfpathlineto{\pgfqpoint{0.718969in}{1.048586in}}%
\pgfpathlineto{\pgfqpoint{0.710330in}{1.045904in}}%
\pgfpathlineto{\pgfqpoint{0.701653in}{1.043626in}}%
\pgfpathclose%
\pgfusepath{fill}%
\end{pgfscope}%
\begin{pgfscope}%
\pgfpathrectangle{\pgfqpoint{0.407592in}{0.222904in}}{\pgfqpoint{1.411163in}{1.411163in}}%
\pgfusepath{clip}%
\pgfsetbuttcap%
\pgfsetroundjoin%
\definecolor{currentfill}{rgb}{0.556863,0.654902,0.847059}%
\pgfsetfillcolor{currentfill}%
\pgfsetfillopacity{0.130000}%
\pgfsetlinewidth{0.000000pt}%
\definecolor{currentstroke}{rgb}{0.000000,0.000000,0.000000}%
\pgfsetstrokecolor{currentstroke}%
\pgfsetdash{}{0pt}%
\pgfpathmoveto{\pgfqpoint{1.298279in}{0.870807in}}%
\pgfpathlineto{\pgfqpoint{1.314761in}{0.876005in}}%
\pgfpathlineto{\pgfqpoint{1.331310in}{0.882011in}}%
\pgfpathlineto{\pgfqpoint{1.339060in}{0.894674in}}%
\pgfpathlineto{\pgfqpoint{1.346797in}{0.907699in}}%
\pgfpathlineto{\pgfqpoint{1.330264in}{0.902265in}}%
\pgfpathlineto{\pgfqpoint{1.313797in}{0.897632in}}%
\pgfpathlineto{\pgfqpoint{1.306044in}{0.884045in}}%
\pgfpathlineto{\pgfqpoint{1.298279in}{0.870807in}}%
\pgfpathclose%
\pgfusepath{fill}%
\end{pgfscope}%
\begin{pgfscope}%
\pgfpathrectangle{\pgfqpoint{0.407592in}{0.222904in}}{\pgfqpoint{1.411163in}{1.411163in}}%
\pgfusepath{clip}%
\pgfsetbuttcap%
\pgfsetroundjoin%
\definecolor{currentfill}{rgb}{0.556863,0.654902,0.847059}%
\pgfsetfillcolor{currentfill}%
\pgfsetfillopacity{0.130000}%
\pgfsetlinewidth{0.000000pt}%
\definecolor{currentstroke}{rgb}{0.000000,0.000000,0.000000}%
\pgfsetstrokecolor{currentstroke}%
\pgfsetdash{}{0pt}%
\pgfpathmoveto{\pgfqpoint{1.041082in}{0.759620in}}%
\pgfpathlineto{\pgfqpoint{1.057194in}{0.750833in}}%
\pgfpathlineto{\pgfqpoint{1.073331in}{0.743981in}}%
\pgfpathlineto{\pgfqpoint{1.081388in}{0.749276in}}%
\pgfpathlineto{\pgfqpoint{1.089418in}{0.755736in}}%
\pgfpathlineto{\pgfqpoint{1.073312in}{0.762532in}}%
\pgfpathlineto{\pgfqpoint{1.057232in}{0.771245in}}%
\pgfpathlineto{\pgfqpoint{1.049172in}{0.764880in}}%
\pgfpathlineto{\pgfqpoint{1.041082in}{0.759620in}}%
\pgfpathclose%
\pgfusepath{fill}%
\end{pgfscope}%
\begin{pgfscope}%
\pgfpathrectangle{\pgfqpoint{0.407592in}{0.222904in}}{\pgfqpoint{1.411163in}{1.411163in}}%
\pgfusepath{clip}%
\pgfsetbuttcap%
\pgfsetroundjoin%
\definecolor{currentfill}{rgb}{0.556863,0.654902,0.847059}%
\pgfsetfillcolor{currentfill}%
\pgfsetfillopacity{0.130000}%
\pgfsetlinewidth{0.000000pt}%
\definecolor{currentstroke}{rgb}{0.000000,0.000000,0.000000}%
\pgfsetstrokecolor{currentstroke}%
\pgfsetdash{}{0pt}%
\pgfpathmoveto{\pgfqpoint{1.185918in}{0.784499in}}%
\pgfpathlineto{\pgfqpoint{1.202199in}{0.786772in}}%
\pgfpathlineto{\pgfqpoint{1.218538in}{0.790440in}}%
\pgfpathlineto{\pgfqpoint{1.226403in}{0.801094in}}%
\pgfpathlineto{\pgfqpoint{1.234253in}{0.812426in}}%
\pgfpathlineto{\pgfqpoint{1.217929in}{0.809237in}}%
\pgfpathlineto{\pgfqpoint{1.201664in}{0.807384in}}%
\pgfpathlineto{\pgfqpoint{1.193799in}{0.795592in}}%
\pgfpathlineto{\pgfqpoint{1.185918in}{0.784499in}}%
\pgfpathclose%
\pgfusepath{fill}%
\end{pgfscope}%
\begin{pgfscope}%
\pgfpathrectangle{\pgfqpoint{0.407592in}{0.222904in}}{\pgfqpoint{1.411163in}{1.411163in}}%
\pgfusepath{clip}%
\pgfsetbuttcap%
\pgfsetroundjoin%
\definecolor{currentfill}{rgb}{0.556863,0.654902,0.847059}%
\pgfsetfillcolor{currentfill}%
\pgfsetfillopacity{0.130000}%
\pgfsetlinewidth{0.000000pt}%
\definecolor{currentstroke}{rgb}{0.000000,0.000000,0.000000}%
\pgfsetstrokecolor{currentstroke}%
\pgfsetdash{}{0pt}%
\pgfpathmoveto{\pgfqpoint{1.474805in}{1.041628in}}%
\pgfpathlineto{\pgfqpoint{1.491613in}{1.048082in}}%
\pgfpathlineto{\pgfqpoint{1.508491in}{1.055016in}}%
\pgfpathlineto{\pgfqpoint{1.516070in}{1.068750in}}%
\pgfpathlineto{\pgfqpoint{1.523638in}{1.082659in}}%
\pgfpathlineto{\pgfqpoint{1.506778in}{1.076162in}}%
\pgfpathlineto{\pgfqpoint{1.489986in}{1.070146in}}%
\pgfpathlineto{\pgfqpoint{1.482401in}{1.055806in}}%
\pgfpathlineto{\pgfqpoint{1.474805in}{1.041628in}}%
\pgfpathclose%
\pgfusepath{fill}%
\end{pgfscope}%
\begin{pgfscope}%
\pgfpathrectangle{\pgfqpoint{0.407592in}{0.222904in}}{\pgfqpoint{1.411163in}{1.411163in}}%
\pgfusepath{clip}%
\pgfsetbuttcap%
\pgfsetroundjoin%
\definecolor{currentfill}{rgb}{0.556863,0.654902,0.847059}%
\pgfsetfillcolor{currentfill}%
\pgfsetfillopacity{0.130000}%
\pgfsetlinewidth{0.000000pt}%
\definecolor{currentstroke}{rgb}{0.000000,0.000000,0.000000}%
\pgfsetstrokecolor{currentstroke}%
\pgfsetdash{}{0pt}%
\pgfpathmoveto{\pgfqpoint{1.538740in}{1.110948in}}%
\pgfpathlineto{\pgfqpoint{1.555653in}{1.117485in}}%
\pgfpathlineto{\pgfqpoint{1.572637in}{1.124449in}}%
\pgfpathlineto{\pgfqpoint{1.580153in}{1.138411in}}%
\pgfpathlineto{\pgfqpoint{1.563178in}{1.131644in}}%
\pgfpathlineto{\pgfqpoint{1.546274in}{1.125305in}}%
\pgfusepath{fill}%
\end{pgfscope}%
\begin{pgfscope}%
\pgfpathrectangle{\pgfqpoint{0.407592in}{0.222904in}}{\pgfqpoint{1.411163in}{1.411163in}}%
\pgfusepath{clip}%
\pgfsetbuttcap%
\pgfsetroundjoin%
\definecolor{currentfill}{rgb}{0.556863,0.654902,0.847059}%
\pgfsetfillcolor{currentfill}%
\pgfsetfillopacity{0.130000}%
\pgfsetlinewidth{0.000000pt}%
\definecolor{currentstroke}{rgb}{0.000000,0.000000,0.000000}%
\pgfsetstrokecolor{currentstroke}%
\pgfsetdash{}{0pt}%
\pgfpathmoveto{\pgfqpoint{1.410825in}{0.973598in}}%
\pgfpathlineto{\pgfqpoint{1.427529in}{0.979942in}}%
\pgfpathlineto{\pgfqpoint{1.444301in}{0.986838in}}%
\pgfpathlineto{\pgfqpoint{1.451945in}{1.000213in}}%
\pgfpathlineto{\pgfqpoint{1.459576in}{1.013815in}}%
\pgfpathlineto{\pgfqpoint{1.442821in}{1.007411in}}%
\pgfpathlineto{\pgfqpoint{1.426133in}{1.001563in}}%
\pgfpathlineto{\pgfqpoint{1.418485in}{0.987474in}}%
\pgfpathlineto{\pgfqpoint{1.410825in}{0.973598in}}%
\pgfpathclose%
\pgfusepath{fill}%
\end{pgfscope}%
\begin{pgfscope}%
\pgfpathrectangle{\pgfqpoint{0.407592in}{0.222904in}}{\pgfqpoint{1.411163in}{1.411163in}}%
\pgfusepath{clip}%
\pgfsetbuttcap%
\pgfsetroundjoin%
\definecolor{currentfill}{rgb}{0.556863,0.654902,0.847059}%
\pgfsetfillcolor{currentfill}%
\pgfsetfillopacity{0.130000}%
\pgfsetlinewidth{0.000000pt}%
\definecolor{currentstroke}{rgb}{0.000000,0.000000,0.000000}%
\pgfsetstrokecolor{currentstroke}%
\pgfsetdash{}{0pt}%
\pgfpathmoveto{\pgfqpoint{0.733752in}{1.005162in}}%
\pgfpathlineto{\pgfqpoint{0.749815in}{0.986140in}}%
\pgfpathlineto{\pgfqpoint{0.765886in}{0.967286in}}%
\pgfpathlineto{\pgfqpoint{0.774478in}{0.969217in}}%
\pgfpathlineto{\pgfqpoint{0.783032in}{0.971623in}}%
\pgfpathlineto{\pgfqpoint{0.767003in}{0.990659in}}%
\pgfpathlineto{\pgfqpoint{0.750983in}{1.009845in}}%
\pgfpathlineto{\pgfqpoint{0.742387in}{1.007285in}}%
\pgfpathlineto{\pgfqpoint{0.733752in}{1.005162in}}%
\pgfpathclose%
\pgfusepath{fill}%
\end{pgfscope}%
\begin{pgfscope}%
\pgfpathrectangle{\pgfqpoint{0.407592in}{0.222904in}}{\pgfqpoint{1.411163in}{1.411163in}}%
\pgfusepath{clip}%
\pgfsetbuttcap%
\pgfsetroundjoin%
\definecolor{currentfill}{rgb}{0.556863,0.654902,0.847059}%
\pgfsetfillcolor{currentfill}%
\pgfsetfillopacity{0.130000}%
\pgfsetlinewidth{0.000000pt}%
\definecolor{currentstroke}{rgb}{0.000000,0.000000,0.000000}%
\pgfsetstrokecolor{currentstroke}%
\pgfsetdash{}{0pt}%
\pgfpathmoveto{\pgfqpoint{0.879386in}{0.862286in}}%
\pgfpathlineto{\pgfqpoint{0.895476in}{0.845319in}}%
\pgfpathlineto{\pgfqpoint{0.911574in}{0.828902in}}%
\pgfpathlineto{\pgfqpoint{0.919921in}{0.830943in}}%
\pgfpathlineto{\pgfqpoint{0.928231in}{0.833726in}}%
\pgfpathlineto{\pgfqpoint{0.912173in}{0.850400in}}%
\pgfpathlineto{\pgfqpoint{0.896126in}{0.867602in}}%
\pgfpathlineto{\pgfqpoint{0.887774in}{0.864612in}}%
\pgfpathlineto{\pgfqpoint{0.879386in}{0.862286in}}%
\pgfpathclose%
\pgfusepath{fill}%
\end{pgfscope}%
\begin{pgfscope}%
\pgfpathrectangle{\pgfqpoint{0.407592in}{0.222904in}}{\pgfqpoint{1.411163in}{1.411163in}}%
\pgfusepath{clip}%
\pgfsetbuttcap%
\pgfsetroundjoin%
\definecolor{currentfill}{rgb}{0.556863,0.654902,0.847059}%
\pgfsetfillcolor{currentfill}%
\pgfsetfillopacity{0.130000}%
\pgfsetlinewidth{0.000000pt}%
\definecolor{currentstroke}{rgb}{0.000000,0.000000,0.000000}%
\pgfsetstrokecolor{currentstroke}%
\pgfsetdash{}{0pt}%
\pgfpathmoveto{\pgfqpoint{1.234253in}{0.812426in}}%
\pgfpathlineto{\pgfqpoint{1.250637in}{0.816806in}}%
\pgfpathlineto{\pgfqpoint{1.267086in}{0.822236in}}%
\pgfpathlineto{\pgfqpoint{1.274905in}{0.833625in}}%
\pgfpathlineto{\pgfqpoint{1.282710in}{0.845550in}}%
\pgfpathlineto{\pgfqpoint{1.266277in}{0.840686in}}%
\pgfpathlineto{\pgfqpoint{1.249907in}{0.836852in}}%
\pgfpathlineto{\pgfqpoint{1.242087in}{0.824367in}}%
\pgfpathlineto{\pgfqpoint{1.234253in}{0.812426in}}%
\pgfpathclose%
\pgfusepath{fill}%
\end{pgfscope}%
\begin{pgfscope}%
\pgfpathrectangle{\pgfqpoint{0.407592in}{0.222904in}}{\pgfqpoint{1.411163in}{1.411163in}}%
\pgfusepath{clip}%
\pgfsetbuttcap%
\pgfsetroundjoin%
\definecolor{currentfill}{rgb}{0.556863,0.654902,0.847059}%
\pgfsetfillcolor{currentfill}%
\pgfsetfillopacity{0.130000}%
\pgfsetlinewidth{0.000000pt}%
\definecolor{currentstroke}{rgb}{0.000000,0.000000,0.000000}%
\pgfsetstrokecolor{currentstroke}%
\pgfsetdash{}{0pt}%
\pgfpathmoveto{\pgfqpoint{0.587834in}{1.157126in}}%
\pgfpathlineto{\pgfqpoint{0.603870in}{1.137440in}}%
\pgfpathlineto{\pgfqpoint{0.619915in}{1.117823in}}%
\pgfpathlineto{\pgfqpoint{0.628756in}{1.119639in}}%
\pgfpathlineto{\pgfqpoint{0.637558in}{1.121797in}}%
\pgfpathlineto{\pgfqpoint{0.621556in}{1.141530in}}%
\pgfpathlineto{\pgfqpoint{0.605563in}{1.161323in}}%
\pgfpathlineto{\pgfqpoint{0.596718in}{1.159064in}}%
\pgfpathlineto{\pgfqpoint{0.587834in}{1.157126in}}%
\pgfpathclose%
\pgfusepath{fill}%
\end{pgfscope}%
\begin{pgfscope}%
\pgfpathrectangle{\pgfqpoint{0.407592in}{0.222904in}}{\pgfqpoint{1.411163in}{1.411163in}}%
\pgfusepath{clip}%
\pgfsetbuttcap%
\pgfsetroundjoin%
\definecolor{currentfill}{rgb}{0.556863,0.654902,0.847059}%
\pgfsetfillcolor{currentfill}%
\pgfsetfillopacity{0.130000}%
\pgfsetlinewidth{0.000000pt}%
\definecolor{currentstroke}{rgb}{0.000000,0.000000,0.000000}%
\pgfsetstrokecolor{currentstroke}%
\pgfsetdash{}{0pt}%
\pgfpathmoveto{\pgfqpoint{0.992563in}{0.774795in}}%
\pgfpathlineto{\pgfqpoint{1.008677in}{0.762842in}}%
\pgfpathlineto{\pgfqpoint{1.024806in}{0.752405in}}%
\pgfpathlineto{\pgfqpoint{1.032960in}{0.755469in}}%
\pgfpathlineto{\pgfqpoint{1.041082in}{0.759620in}}%
\pgfpathlineto{\pgfqpoint{1.024989in}{0.770158in}}%
\pgfpathlineto{\pgfqpoint{1.008914in}{0.782227in}}%
\pgfpathlineto{\pgfqpoint{1.000755in}{0.778016in}}%
\pgfpathlineto{\pgfqpoint{0.992563in}{0.774795in}}%
\pgfpathclose%
\pgfusepath{fill}%
\end{pgfscope}%
\begin{pgfscope}%
\pgfpathrectangle{\pgfqpoint{0.407592in}{0.222904in}}{\pgfqpoint{1.411163in}{1.411163in}}%
\pgfusepath{clip}%
\pgfsetbuttcap%
\pgfsetroundjoin%
\definecolor{currentfill}{rgb}{0.556863,0.654902,0.847059}%
\pgfsetfillcolor{currentfill}%
\pgfsetfillopacity{0.130000}%
\pgfsetlinewidth{0.000000pt}%
\definecolor{currentstroke}{rgb}{0.000000,0.000000,0.000000}%
\pgfsetstrokecolor{currentstroke}%
\pgfsetdash{}{0pt}%
\pgfpathmoveto{\pgfqpoint{1.346797in}{0.907699in}}%
\pgfpathlineto{\pgfqpoint{1.363396in}{0.913858in}}%
\pgfpathlineto{\pgfqpoint{1.380064in}{0.920672in}}%
\pgfpathlineto{\pgfqpoint{1.387773in}{0.933468in}}%
\pgfpathlineto{\pgfqpoint{1.395469in}{0.946572in}}%
\pgfpathlineto{\pgfqpoint{1.378817in}{0.940307in}}%
\pgfpathlineto{\pgfqpoint{1.362233in}{0.934699in}}%
\pgfpathlineto{\pgfqpoint{1.354521in}{0.921051in}}%
\pgfpathlineto{\pgfqpoint{1.346797in}{0.907699in}}%
\pgfpathclose%
\pgfusepath{fill}%
\end{pgfscope}%
\begin{pgfscope}%
\pgfpathrectangle{\pgfqpoint{0.407592in}{0.222904in}}{\pgfqpoint{1.411163in}{1.411163in}}%
\pgfusepath{clip}%
\pgfsetbuttcap%
\pgfsetroundjoin%
\definecolor{currentfill}{rgb}{0.556863,0.654902,0.847059}%
\pgfsetfillcolor{currentfill}%
\pgfsetfillopacity{0.130000}%
\pgfsetlinewidth{0.000000pt}%
\definecolor{currentstroke}{rgb}{0.000000,0.000000,0.000000}%
\pgfsetstrokecolor{currentstroke}%
\pgfsetdash{}{0pt}%
\pgfpathmoveto{\pgfqpoint{0.765886in}{0.967286in}}%
\pgfpathlineto{\pgfqpoint{0.781966in}{0.948626in}}%
\pgfpathlineto{\pgfqpoint{0.798055in}{0.930191in}}%
\pgfpathlineto{\pgfqpoint{0.806604in}{0.931888in}}%
\pgfpathlineto{\pgfqpoint{0.815115in}{0.934101in}}%
\pgfpathlineto{\pgfqpoint{0.799069in}{0.952761in}}%
\pgfpathlineto{\pgfqpoint{0.783032in}{0.971623in}}%
\pgfpathlineto{\pgfqpoint{0.774478in}{0.969217in}}%
\pgfpathlineto{\pgfqpoint{0.765886in}{0.967286in}}%
\pgfpathclose%
\pgfusepath{fill}%
\end{pgfscope}%
\begin{pgfscope}%
\pgfpathrectangle{\pgfqpoint{0.407592in}{0.222904in}}{\pgfqpoint{1.411163in}{1.411163in}}%
\pgfusepath{clip}%
\pgfsetbuttcap%
\pgfsetroundjoin%
\definecolor{currentfill}{rgb}{0.556863,0.654902,0.847059}%
\pgfsetfillcolor{currentfill}%
\pgfsetfillopacity{0.130000}%
\pgfsetlinewidth{0.000000pt}%
\definecolor{currentstroke}{rgb}{0.000000,0.000000,0.000000}%
\pgfsetstrokecolor{currentstroke}%
\pgfsetdash{}{0pt}%
\pgfpathmoveto{\pgfqpoint{1.523638in}{1.082659in}}%
\pgfpathlineto{\pgfqpoint{1.540569in}{1.089608in}}%
\pgfpathlineto{\pgfqpoint{1.557571in}{1.096981in}}%
\pgfpathlineto{\pgfqpoint{1.565110in}{1.110635in}}%
\pgfpathlineto{\pgfqpoint{1.572637in}{1.124449in}}%
\pgfpathlineto{\pgfqpoint{1.555653in}{1.117485in}}%
\pgfpathlineto{\pgfqpoint{1.538740in}{1.110948in}}%
\pgfpathlineto{\pgfqpoint{1.531195in}{1.096729in}}%
\pgfpathlineto{\pgfqpoint{1.523638in}{1.082659in}}%
\pgfpathclose%
\pgfusepath{fill}%
\end{pgfscope}%
\begin{pgfscope}%
\pgfpathrectangle{\pgfqpoint{0.407592in}{0.222904in}}{\pgfqpoint{1.411163in}{1.411163in}}%
\pgfusepath{clip}%
\pgfsetbuttcap%
\pgfsetroundjoin%
\definecolor{currentfill}{rgb}{0.556863,0.654902,0.847059}%
\pgfsetfillcolor{currentfill}%
\pgfsetfillopacity{0.130000}%
\pgfsetlinewidth{0.000000pt}%
\definecolor{currentstroke}{rgb}{0.000000,0.000000,0.000000}%
\pgfsetstrokecolor{currentstroke}%
\pgfsetdash{}{0pt}%
\pgfpathmoveto{\pgfqpoint{0.911574in}{0.828902in}}%
\pgfpathlineto{\pgfqpoint{0.927683in}{0.813142in}}%
\pgfpathlineto{\pgfqpoint{0.943801in}{0.798165in}}%
\pgfpathlineto{\pgfqpoint{0.952106in}{0.799874in}}%
\pgfpathlineto{\pgfqpoint{0.960375in}{0.802408in}}%
\pgfpathlineto{\pgfqpoint{0.944298in}{0.817685in}}%
\pgfpathlineto{\pgfqpoint{0.928231in}{0.833726in}}%
\pgfpathlineto{\pgfqpoint{0.919921in}{0.830943in}}%
\pgfpathlineto{\pgfqpoint{0.911574in}{0.828902in}}%
\pgfpathclose%
\pgfusepath{fill}%
\end{pgfscope}%
\begin{pgfscope}%
\pgfpathrectangle{\pgfqpoint{0.407592in}{0.222904in}}{\pgfqpoint{1.411163in}{1.411163in}}%
\pgfusepath{clip}%
\pgfsetbuttcap%
\pgfsetroundjoin%
\definecolor{currentfill}{rgb}{0.556863,0.654902,0.847059}%
\pgfsetfillcolor{currentfill}%
\pgfsetfillopacity{0.130000}%
\pgfsetlinewidth{0.000000pt}%
\definecolor{currentstroke}{rgb}{0.000000,0.000000,0.000000}%
\pgfsetstrokecolor{currentstroke}%
\pgfsetdash{}{0pt}%
\pgfpathmoveto{\pgfqpoint{0.619915in}{1.117823in}}%
\pgfpathlineto{\pgfqpoint{0.635969in}{1.098286in}}%
\pgfpathlineto{\pgfqpoint{0.652031in}{1.078837in}}%
\pgfpathlineto{\pgfqpoint{0.660829in}{1.080506in}}%
\pgfpathlineto{\pgfqpoint{0.669588in}{1.082539in}}%
\pgfpathlineto{\pgfqpoint{0.653569in}{1.102130in}}%
\pgfpathlineto{\pgfqpoint{0.637558in}{1.121797in}}%
\pgfpathlineto{\pgfqpoint{0.628756in}{1.119639in}}%
\pgfpathlineto{\pgfqpoint{0.619915in}{1.117823in}}%
\pgfpathclose%
\pgfusepath{fill}%
\end{pgfscope}%
\begin{pgfscope}%
\pgfpathrectangle{\pgfqpoint{0.407592in}{0.222904in}}{\pgfqpoint{1.411163in}{1.411163in}}%
\pgfusepath{clip}%
\pgfsetbuttcap%
\pgfsetroundjoin%
\definecolor{currentfill}{rgb}{0.556863,0.654902,0.847059}%
\pgfsetfillcolor{currentfill}%
\pgfsetfillopacity{0.130000}%
\pgfsetlinewidth{0.000000pt}%
\definecolor{currentstroke}{rgb}{0.000000,0.000000,0.000000}%
\pgfsetstrokecolor{currentstroke}%
\pgfsetdash{}{0pt}%
\pgfpathmoveto{\pgfqpoint{1.121729in}{0.748276in}}%
\pgfpathlineto{\pgfqpoint{1.137947in}{0.747584in}}%
\pgfpathlineto{\pgfqpoint{1.154213in}{0.748759in}}%
\pgfpathlineto{\pgfqpoint{1.162168in}{0.756246in}}%
\pgfpathlineto{\pgfqpoint{1.170103in}{0.764745in}}%
\pgfpathlineto{\pgfqpoint{1.153859in}{0.763749in}}%
\pgfpathlineto{\pgfqpoint{1.137663in}{0.764555in}}%
\pgfpathlineto{\pgfqpoint{1.129707in}{0.755882in}}%
\pgfpathlineto{\pgfqpoint{1.121729in}{0.748276in}}%
\pgfpathclose%
\pgfusepath{fill}%
\end{pgfscope}%
\begin{pgfscope}%
\pgfpathrectangle{\pgfqpoint{0.407592in}{0.222904in}}{\pgfqpoint{1.411163in}{1.411163in}}%
\pgfusepath{clip}%
\pgfsetbuttcap%
\pgfsetroundjoin%
\definecolor{currentfill}{rgb}{0.556863,0.654902,0.847059}%
\pgfsetfillcolor{currentfill}%
\pgfsetfillopacity{0.130000}%
\pgfsetlinewidth{0.000000pt}%
\definecolor{currentstroke}{rgb}{0.000000,0.000000,0.000000}%
\pgfsetstrokecolor{currentstroke}%
\pgfsetdash{}{0pt}%
\pgfpathmoveto{\pgfqpoint{1.459576in}{1.013815in}}%
\pgfpathlineto{\pgfqpoint{1.476402in}{1.020733in}}%
\pgfpathlineto{\pgfqpoint{1.493298in}{1.028128in}}%
\pgfpathlineto{\pgfqpoint{1.500900in}{1.041471in}}%
\pgfpathlineto{\pgfqpoint{1.508491in}{1.055016in}}%
\pgfpathlineto{\pgfqpoint{1.491613in}{1.048082in}}%
\pgfpathlineto{\pgfqpoint{1.474805in}{1.041628in}}%
\pgfpathlineto{\pgfqpoint{1.467196in}{1.027626in}}%
\pgfpathlineto{\pgfqpoint{1.459576in}{1.013815in}}%
\pgfpathclose%
\pgfusepath{fill}%
\end{pgfscope}%
\begin{pgfscope}%
\pgfpathrectangle{\pgfqpoint{0.407592in}{0.222904in}}{\pgfqpoint{1.411163in}{1.411163in}}%
\pgfusepath{clip}%
\pgfsetbuttcap%
\pgfsetroundjoin%
\definecolor{currentfill}{rgb}{0.556863,0.654902,0.847059}%
\pgfsetfillcolor{currentfill}%
\pgfsetfillopacity{0.130000}%
\pgfsetlinewidth{0.000000pt}%
\definecolor{currentstroke}{rgb}{0.000000,0.000000,0.000000}%
\pgfsetstrokecolor{currentstroke}%
\pgfsetdash{}{0pt}%
\pgfpathmoveto{\pgfqpoint{1.282710in}{0.845550in}}%
\pgfpathlineto{\pgfqpoint{1.299207in}{0.851329in}}%
\pgfpathlineto{\pgfqpoint{1.315772in}{0.857916in}}%
\pgfpathlineto{\pgfqpoint{1.323547in}{0.869745in}}%
\pgfpathlineto{\pgfqpoint{1.331310in}{0.882011in}}%
\pgfpathlineto{\pgfqpoint{1.314761in}{0.876005in}}%
\pgfpathlineto{\pgfqpoint{1.298279in}{0.870807in}}%
\pgfpathlineto{\pgfqpoint{1.290501in}{0.857961in}}%
\pgfpathlineto{\pgfqpoint{1.282710in}{0.845550in}}%
\pgfpathclose%
\pgfusepath{fill}%
\end{pgfscope}%
\begin{pgfscope}%
\pgfpathrectangle{\pgfqpoint{0.407592in}{0.222904in}}{\pgfqpoint{1.411163in}{1.411163in}}%
\pgfusepath{clip}%
\pgfsetbuttcap%
\pgfsetroundjoin%
\definecolor{currentfill}{rgb}{0.556863,0.654902,0.847059}%
\pgfsetfillcolor{currentfill}%
\pgfsetfillopacity{0.130000}%
\pgfsetlinewidth{0.000000pt}%
\definecolor{currentstroke}{rgb}{0.000000,0.000000,0.000000}%
\pgfsetstrokecolor{currentstroke}%
\pgfsetdash{}{0pt}%
\pgfpathmoveto{\pgfqpoint{1.073331in}{0.743981in}}%
\pgfpathlineto{\pgfqpoint{1.089498in}{0.739185in}}%
\pgfpathlineto{\pgfqpoint{1.105701in}{0.736477in}}%
\pgfpathlineto{\pgfqpoint{1.113728in}{0.741796in}}%
\pgfpathlineto{\pgfqpoint{1.121729in}{0.748276in}}%
\pgfpathlineto{\pgfqpoint{1.105555in}{0.750973in}}%
\pgfpathlineto{\pgfqpoint{1.089418in}{0.755736in}}%
\pgfpathlineto{\pgfqpoint{1.081388in}{0.749276in}}%
\pgfpathlineto{\pgfqpoint{1.073331in}{0.743981in}}%
\pgfpathclose%
\pgfusepath{fill}%
\end{pgfscope}%
\begin{pgfscope}%
\pgfpathrectangle{\pgfqpoint{0.407592in}{0.222904in}}{\pgfqpoint{1.411163in}{1.411163in}}%
\pgfusepath{clip}%
\pgfsetbuttcap%
\pgfsetroundjoin%
\definecolor{currentfill}{rgb}{0.556863,0.654902,0.847059}%
\pgfsetfillcolor{currentfill}%
\pgfsetfillopacity{0.130000}%
\pgfsetlinewidth{0.000000pt}%
\definecolor{currentstroke}{rgb}{0.000000,0.000000,0.000000}%
\pgfsetstrokecolor{currentstroke}%
\pgfsetdash{}{0pt}%
\pgfpathmoveto{\pgfqpoint{0.798055in}{0.930191in}}%
\pgfpathlineto{\pgfqpoint{0.814151in}{0.912018in}}%
\pgfpathlineto{\pgfqpoint{0.830257in}{0.894149in}}%
\pgfpathlineto{\pgfqpoint{0.838764in}{0.895558in}}%
\pgfpathlineto{\pgfqpoint{0.847233in}{0.897533in}}%
\pgfpathlineto{\pgfqpoint{0.831170in}{0.915678in}}%
\pgfpathlineto{\pgfqpoint{0.815115in}{0.934101in}}%
\pgfpathlineto{\pgfqpoint{0.806604in}{0.931888in}}%
\pgfpathlineto{\pgfqpoint{0.798055in}{0.930191in}}%
\pgfpathclose%
\pgfusepath{fill}%
\end{pgfscope}%
\begin{pgfscope}%
\pgfpathrectangle{\pgfqpoint{0.407592in}{0.222904in}}{\pgfqpoint{1.411163in}{1.411163in}}%
\pgfusepath{clip}%
\pgfsetbuttcap%
\pgfsetroundjoin%
\definecolor{currentfill}{rgb}{0.556863,0.654902,0.847059}%
\pgfsetfillcolor{currentfill}%
\pgfsetfillopacity{0.130000}%
\pgfsetlinewidth{0.000000pt}%
\definecolor{currentstroke}{rgb}{0.000000,0.000000,0.000000}%
\pgfsetstrokecolor{currentstroke}%
\pgfsetdash{}{0pt}%
\pgfpathmoveto{\pgfqpoint{1.170103in}{0.764745in}}%
\pgfpathlineto{\pgfqpoint{1.186402in}{0.767373in}}%
\pgfpathlineto{\pgfqpoint{1.202760in}{0.771439in}}%
\pgfpathlineto{\pgfqpoint{1.210657in}{0.780533in}}%
\pgfpathlineto{\pgfqpoint{1.218538in}{0.790440in}}%
\pgfpathlineto{\pgfqpoint{1.202199in}{0.786772in}}%
\pgfpathlineto{\pgfqpoint{1.185918in}{0.784499in}}%
\pgfpathlineto{\pgfqpoint{1.178019in}{0.774189in}}%
\pgfpathlineto{\pgfqpoint{1.170103in}{0.764745in}}%
\pgfpathclose%
\pgfusepath{fill}%
\end{pgfscope}%
\begin{pgfscope}%
\pgfpathrectangle{\pgfqpoint{0.407592in}{0.222904in}}{\pgfqpoint{1.411163in}{1.411163in}}%
\pgfusepath{clip}%
\pgfsetbuttcap%
\pgfsetroundjoin%
\definecolor{currentfill}{rgb}{0.556863,0.654902,0.847059}%
\pgfsetfillcolor{currentfill}%
\pgfsetfillopacity{0.130000}%
\pgfsetlinewidth{0.000000pt}%
\definecolor{currentstroke}{rgb}{0.000000,0.000000,0.000000}%
\pgfsetstrokecolor{currentstroke}%
\pgfsetdash{}{0pt}%
\pgfpathmoveto{\pgfqpoint{0.652031in}{1.078837in}}%
\pgfpathlineto{\pgfqpoint{0.668101in}{1.059489in}}%
\pgfpathlineto{\pgfqpoint{0.684181in}{1.040255in}}%
\pgfpathlineto{\pgfqpoint{0.692937in}{1.041746in}}%
\pgfpathlineto{\pgfqpoint{0.701653in}{1.043626in}}%
\pgfpathlineto{\pgfqpoint{0.685616in}{1.063034in}}%
\pgfpathlineto{\pgfqpoint{0.669588in}{1.082539in}}%
\pgfpathlineto{\pgfqpoint{0.660829in}{1.080506in}}%
\pgfpathlineto{\pgfqpoint{0.652031in}{1.078837in}}%
\pgfpathclose%
\pgfusepath{fill}%
\end{pgfscope}%
\begin{pgfscope}%
\pgfpathrectangle{\pgfqpoint{0.407592in}{0.222904in}}{\pgfqpoint{1.411163in}{1.411163in}}%
\pgfusepath{clip}%
\pgfsetbuttcap%
\pgfsetroundjoin%
\definecolor{currentfill}{rgb}{0.556863,0.654902,0.847059}%
\pgfsetfillcolor{currentfill}%
\pgfsetfillopacity{0.130000}%
\pgfsetlinewidth{0.000000pt}%
\definecolor{currentstroke}{rgb}{0.000000,0.000000,0.000000}%
\pgfsetstrokecolor{currentstroke}%
\pgfsetdash{}{0pt}%
\pgfpathmoveto{\pgfqpoint{1.395469in}{0.946572in}}%
\pgfpathlineto{\pgfqpoint{1.412189in}{0.953437in}}%
\pgfpathlineto{\pgfqpoint{1.428979in}{0.960848in}}%
\pgfpathlineto{\pgfqpoint{1.436646in}{0.973709in}}%
\pgfpathlineto{\pgfqpoint{1.444301in}{0.986838in}}%
\pgfpathlineto{\pgfqpoint{1.427529in}{0.979942in}}%
\pgfpathlineto{\pgfqpoint{1.410825in}{0.973598in}}%
\pgfpathlineto{\pgfqpoint{1.403153in}{0.959957in}}%
\pgfpathlineto{\pgfqpoint{1.395469in}{0.946572in}}%
\pgfpathclose%
\pgfusepath{fill}%
\end{pgfscope}%
\begin{pgfscope}%
\pgfpathrectangle{\pgfqpoint{0.407592in}{0.222904in}}{\pgfqpoint{1.411163in}{1.411163in}}%
\pgfusepath{clip}%
\pgfsetbuttcap%
\pgfsetroundjoin%
\definecolor{currentfill}{rgb}{0.556863,0.654902,0.847059}%
\pgfsetfillcolor{currentfill}%
\pgfsetfillopacity{0.130000}%
\pgfsetlinewidth{0.000000pt}%
\definecolor{currentstroke}{rgb}{0.000000,0.000000,0.000000}%
\pgfsetstrokecolor{currentstroke}%
\pgfsetdash{}{0pt}%
\pgfpathmoveto{\pgfqpoint{1.024806in}{0.752405in}}%
\pgfpathlineto{\pgfqpoint{1.040956in}{0.743693in}}%
\pgfpathlineto{\pgfqpoint{1.057128in}{0.736877in}}%
\pgfpathlineto{\pgfqpoint{1.065245in}{0.739857in}}%
\pgfpathlineto{\pgfqpoint{1.073331in}{0.743981in}}%
\pgfpathlineto{\pgfqpoint{1.057194in}{0.750833in}}%
\pgfpathlineto{\pgfqpoint{1.041082in}{0.759620in}}%
\pgfpathlineto{\pgfqpoint{1.032960in}{0.755469in}}%
\pgfpathlineto{\pgfqpoint{1.024806in}{0.752405in}}%
\pgfpathclose%
\pgfusepath{fill}%
\end{pgfscope}%
\begin{pgfscope}%
\pgfpathrectangle{\pgfqpoint{0.407592in}{0.222904in}}{\pgfqpoint{1.411163in}{1.411163in}}%
\pgfusepath{clip}%
\pgfsetbuttcap%
\pgfsetroundjoin%
\definecolor{currentfill}{rgb}{0.556863,0.654902,0.847059}%
\pgfsetfillcolor{currentfill}%
\pgfsetfillopacity{0.130000}%
\pgfsetlinewidth{0.000000pt}%
\definecolor{currentstroke}{rgb}{0.000000,0.000000,0.000000}%
\pgfsetstrokecolor{currentstroke}%
\pgfsetdash{}{0pt}%
\pgfpathmoveto{\pgfqpoint{0.943801in}{0.798165in}}%
\pgfpathlineto{\pgfqpoint{0.959930in}{0.784119in}}%
\pgfpathlineto{\pgfqpoint{0.976071in}{0.771175in}}%
\pgfpathlineto{\pgfqpoint{0.984335in}{0.772531in}}%
\pgfpathlineto{\pgfqpoint{0.992563in}{0.774795in}}%
\pgfpathlineto{\pgfqpoint{0.976463in}{0.788051in}}%
\pgfpathlineto{\pgfqpoint{0.960375in}{0.802408in}}%
\pgfpathlineto{\pgfqpoint{0.952106in}{0.799874in}}%
\pgfpathlineto{\pgfqpoint{0.943801in}{0.798165in}}%
\pgfpathclose%
\pgfusepath{fill}%
\end{pgfscope}%
\begin{pgfscope}%
\pgfpathrectangle{\pgfqpoint{0.407592in}{0.222904in}}{\pgfqpoint{1.411163in}{1.411163in}}%
\pgfusepath{clip}%
\pgfsetbuttcap%
\pgfsetroundjoin%
\definecolor{currentfill}{rgb}{0.556863,0.654902,0.847059}%
\pgfsetfillcolor{currentfill}%
\pgfsetfillopacity{0.130000}%
\pgfsetlinewidth{0.000000pt}%
\definecolor{currentstroke}{rgb}{0.000000,0.000000,0.000000}%
\pgfsetstrokecolor{currentstroke}%
\pgfsetdash{}{0pt}%
\pgfpathmoveto{\pgfqpoint{0.830257in}{0.894149in}}%
\pgfpathlineto{\pgfqpoint{0.846371in}{0.876637in}}%
\pgfpathlineto{\pgfqpoint{0.862494in}{0.859543in}}%
\pgfpathlineto{\pgfqpoint{0.870959in}{0.860605in}}%
\pgfpathlineto{\pgfqpoint{0.879386in}{0.862286in}}%
\pgfpathlineto{\pgfqpoint{0.863305in}{0.879715in}}%
\pgfpathlineto{\pgfqpoint{0.847233in}{0.897533in}}%
\pgfpathlineto{\pgfqpoint{0.838764in}{0.895558in}}%
\pgfpathlineto{\pgfqpoint{0.830257in}{0.894149in}}%
\pgfpathclose%
\pgfusepath{fill}%
\end{pgfscope}%
\begin{pgfscope}%
\pgfpathrectangle{\pgfqpoint{0.407592in}{0.222904in}}{\pgfqpoint{1.411163in}{1.411163in}}%
\pgfusepath{clip}%
\pgfsetbuttcap%
\pgfsetroundjoin%
\definecolor{currentfill}{rgb}{0.556863,0.654902,0.847059}%
\pgfsetfillcolor{currentfill}%
\pgfsetfillopacity{0.130000}%
\pgfsetlinewidth{0.000000pt}%
\definecolor{currentstroke}{rgb}{0.000000,0.000000,0.000000}%
\pgfsetstrokecolor{currentstroke}%
\pgfsetdash{}{0pt}%
\pgfpathmoveto{\pgfqpoint{0.684181in}{1.040255in}}%
\pgfpathlineto{\pgfqpoint{0.700269in}{1.021150in}}%
\pgfpathlineto{\pgfqpoint{0.716365in}{1.002191in}}%
\pgfpathlineto{\pgfqpoint{0.725079in}{1.003467in}}%
\pgfpathlineto{\pgfqpoint{0.733752in}{1.005162in}}%
\pgfpathlineto{\pgfqpoint{0.717698in}{1.024330in}}%
\pgfpathlineto{\pgfqpoint{0.701653in}{1.043626in}}%
\pgfpathlineto{\pgfqpoint{0.692937in}{1.041746in}}%
\pgfpathlineto{\pgfqpoint{0.684181in}{1.040255in}}%
\pgfpathclose%
\pgfusepath{fill}%
\end{pgfscope}%
\begin{pgfscope}%
\pgfpathrectangle{\pgfqpoint{0.407592in}{0.222904in}}{\pgfqpoint{1.411163in}{1.411163in}}%
\pgfusepath{clip}%
\pgfsetbuttcap%
\pgfsetroundjoin%
\definecolor{currentfill}{rgb}{0.556863,0.654902,0.847059}%
\pgfsetfillcolor{currentfill}%
\pgfsetfillopacity{0.130000}%
\pgfsetlinewidth{0.000000pt}%
\definecolor{currentstroke}{rgb}{0.000000,0.000000,0.000000}%
\pgfsetstrokecolor{currentstroke}%
\pgfsetdash{}{0pt}%
\pgfpathmoveto{\pgfqpoint{1.218538in}{0.790440in}}%
\pgfpathlineto{\pgfqpoint{1.234939in}{0.795332in}}%
\pgfpathlineto{\pgfqpoint{1.251404in}{0.801285in}}%
\pgfpathlineto{\pgfqpoint{1.259252in}{0.811438in}}%
\pgfpathlineto{\pgfqpoint{1.267086in}{0.822236in}}%
\pgfpathlineto{\pgfqpoint{1.250637in}{0.816806in}}%
\pgfpathlineto{\pgfqpoint{1.234253in}{0.812426in}}%
\pgfpathlineto{\pgfqpoint{1.226403in}{0.801094in}}%
\pgfpathlineto{\pgfqpoint{1.218538in}{0.790440in}}%
\pgfpathclose%
\pgfusepath{fill}%
\end{pgfscope}%
\begin{pgfscope}%
\pgfpathrectangle{\pgfqpoint{0.407592in}{0.222904in}}{\pgfqpoint{1.411163in}{1.411163in}}%
\pgfusepath{clip}%
\pgfsetbuttcap%
\pgfsetroundjoin%
\definecolor{currentfill}{rgb}{0.556863,0.654902,0.847059}%
\pgfsetfillcolor{currentfill}%
\pgfsetfillopacity{0.130000}%
\pgfsetlinewidth{0.000000pt}%
\definecolor{currentstroke}{rgb}{0.000000,0.000000,0.000000}%
\pgfsetstrokecolor{currentstroke}%
\pgfsetdash{}{0pt}%
\pgfpathmoveto{\pgfqpoint{1.331310in}{0.882011in}}%
\pgfpathlineto{\pgfqpoint{1.347926in}{0.888740in}}%
\pgfpathlineto{\pgfqpoint{1.364611in}{0.896115in}}%
\pgfpathlineto{\pgfqpoint{1.372344in}{0.908210in}}%
\pgfpathlineto{\pgfqpoint{1.380064in}{0.920672in}}%
\pgfpathlineto{\pgfqpoint{1.363396in}{0.913858in}}%
\pgfpathlineto{\pgfqpoint{1.346797in}{0.907699in}}%
\pgfpathlineto{\pgfqpoint{1.339060in}{0.894674in}}%
\pgfpathlineto{\pgfqpoint{1.331310in}{0.882011in}}%
\pgfpathclose%
\pgfusepath{fill}%
\end{pgfscope}%
\begin{pgfscope}%
\pgfpathrectangle{\pgfqpoint{0.407592in}{0.222904in}}{\pgfqpoint{1.411163in}{1.411163in}}%
\pgfusepath{clip}%
\pgfsetbuttcap%
\pgfsetroundjoin%
\definecolor{currentfill}{rgb}{0.556863,0.654902,0.847059}%
\pgfsetfillcolor{currentfill}%
\pgfsetfillopacity{0.130000}%
\pgfsetlinewidth{0.000000pt}%
\definecolor{currentstroke}{rgb}{0.000000,0.000000,0.000000}%
\pgfsetstrokecolor{currentstroke}%
\pgfsetdash{}{0pt}%
\pgfpathmoveto{\pgfqpoint{1.508491in}{1.055016in}}%
\pgfpathlineto{\pgfqpoint{1.525440in}{1.062400in}}%
\pgfpathlineto{\pgfqpoint{1.542460in}{1.070200in}}%
\pgfpathlineto{\pgfqpoint{1.550021in}{1.083499in}}%
\pgfpathlineto{\pgfqpoint{1.557571in}{1.096981in}}%
\pgfpathlineto{\pgfqpoint{1.540569in}{1.089608in}}%
\pgfpathlineto{\pgfqpoint{1.523638in}{1.082659in}}%
\pgfpathlineto{\pgfqpoint{1.516070in}{1.068750in}}%
\pgfpathlineto{\pgfqpoint{1.508491in}{1.055016in}}%
\pgfpathclose%
\pgfusepath{fill}%
\end{pgfscope}%
\begin{pgfscope}%
\pgfpathrectangle{\pgfqpoint{0.407592in}{0.222904in}}{\pgfqpoint{1.411163in}{1.411163in}}%
\pgfusepath{clip}%
\pgfsetbuttcap%
\pgfsetroundjoin%
\definecolor{currentfill}{rgb}{0.556863,0.654902,0.847059}%
\pgfsetfillcolor{currentfill}%
\pgfsetfillopacity{0.130000}%
\pgfsetlinewidth{0.000000pt}%
\definecolor{currentstroke}{rgb}{0.000000,0.000000,0.000000}%
\pgfsetstrokecolor{currentstroke}%
\pgfsetdash{}{0pt}%
\pgfpathmoveto{\pgfqpoint{1.572637in}{1.124449in}}%
\pgfpathlineto{\pgfqpoint{1.589692in}{1.131815in}}%
\pgfpathlineto{\pgfqpoint{1.606818in}{1.139560in}}%
\pgfpathlineto{\pgfqpoint{1.614316in}{1.153133in}}%
\pgfpathlineto{\pgfqpoint{1.597198in}{1.145582in}}%
\pgfpathlineto{\pgfqpoint{1.580153in}{1.138411in}}%
\pgfusepath{fill}%
\end{pgfscope}%
\begin{pgfscope}%
\pgfpathrectangle{\pgfqpoint{0.407592in}{0.222904in}}{\pgfqpoint{1.411163in}{1.411163in}}%
\pgfusepath{clip}%
\pgfsetbuttcap%
\pgfsetroundjoin%
\definecolor{currentfill}{rgb}{0.556863,0.654902,0.847059}%
\pgfsetfillcolor{currentfill}%
\pgfsetfillopacity{0.130000}%
\pgfsetlinewidth{0.000000pt}%
\definecolor{currentstroke}{rgb}{0.000000,0.000000,0.000000}%
\pgfsetstrokecolor{currentstroke}%
\pgfsetdash{}{0pt}%
\pgfpathmoveto{\pgfqpoint{0.716365in}{1.002191in}}%
\pgfpathlineto{\pgfqpoint{0.732470in}{0.983400in}}%
\pgfpathlineto{\pgfqpoint{0.748584in}{0.964799in}}%
\pgfpathlineto{\pgfqpoint{0.757255in}{0.965817in}}%
\pgfpathlineto{\pgfqpoint{0.765886in}{0.967286in}}%
\pgfpathlineto{\pgfqpoint{0.749815in}{0.986140in}}%
\pgfpathlineto{\pgfqpoint{0.733752in}{1.005162in}}%
\pgfpathlineto{\pgfqpoint{0.725079in}{1.003467in}}%
\pgfpathlineto{\pgfqpoint{0.716365in}{1.002191in}}%
\pgfpathclose%
\pgfusepath{fill}%
\end{pgfscope}%
\begin{pgfscope}%
\pgfpathrectangle{\pgfqpoint{0.407592in}{0.222904in}}{\pgfqpoint{1.411163in}{1.411163in}}%
\pgfusepath{clip}%
\pgfsetbuttcap%
\pgfsetroundjoin%
\definecolor{currentfill}{rgb}{0.556863,0.654902,0.847059}%
\pgfsetfillcolor{currentfill}%
\pgfsetfillopacity{0.130000}%
\pgfsetlinewidth{0.000000pt}%
\definecolor{currentstroke}{rgb}{0.000000,0.000000,0.000000}%
\pgfsetstrokecolor{currentstroke}%
\pgfsetdash{}{0pt}%
\pgfpathmoveto{\pgfqpoint{0.862494in}{0.859543in}}%
\pgfpathlineto{\pgfqpoint{0.878626in}{0.842941in}}%
\pgfpathlineto{\pgfqpoint{0.894767in}{0.826919in}}%
\pgfpathlineto{\pgfqpoint{0.903190in}{0.827572in}}%
\pgfpathlineto{\pgfqpoint{0.911574in}{0.828902in}}%
\pgfpathlineto{\pgfqpoint{0.895476in}{0.845319in}}%
\pgfpathlineto{\pgfqpoint{0.879386in}{0.862286in}}%
\pgfpathlineto{\pgfqpoint{0.870959in}{0.860605in}}%
\pgfpathlineto{\pgfqpoint{0.862494in}{0.859543in}}%
\pgfpathclose%
\pgfusepath{fill}%
\end{pgfscope}%
\begin{pgfscope}%
\pgfpathrectangle{\pgfqpoint{0.407592in}{0.222904in}}{\pgfqpoint{1.411163in}{1.411163in}}%
\pgfusepath{clip}%
\pgfsetbuttcap%
\pgfsetroundjoin%
\definecolor{currentfill}{rgb}{0.556863,0.654902,0.847059}%
\pgfsetfillcolor{currentfill}%
\pgfsetfillopacity{0.130000}%
\pgfsetlinewidth{0.000000pt}%
\definecolor{currentstroke}{rgb}{0.000000,0.000000,0.000000}%
\pgfsetstrokecolor{currentstroke}%
\pgfsetdash{}{0pt}%
\pgfpathmoveto{\pgfqpoint{1.444301in}{0.986838in}}%
\pgfpathlineto{\pgfqpoint{1.461144in}{0.994243in}}%
\pgfpathlineto{\pgfqpoint{1.478058in}{1.002116in}}%
\pgfpathlineto{\pgfqpoint{1.485684in}{1.015004in}}%
\pgfpathlineto{\pgfqpoint{1.493298in}{1.028128in}}%
\pgfpathlineto{\pgfqpoint{1.476402in}{1.020733in}}%
\pgfpathlineto{\pgfqpoint{1.459576in}{1.013815in}}%
\pgfpathlineto{\pgfqpoint{1.451945in}{1.000213in}}%
\pgfpathlineto{\pgfqpoint{1.444301in}{0.986838in}}%
\pgfpathclose%
\pgfusepath{fill}%
\end{pgfscope}%
\begin{pgfscope}%
\pgfpathrectangle{\pgfqpoint{0.407592in}{0.222904in}}{\pgfqpoint{1.411163in}{1.411163in}}%
\pgfusepath{clip}%
\pgfsetbuttcap%
\pgfsetroundjoin%
\definecolor{currentfill}{rgb}{0.556863,0.654902,0.847059}%
\pgfsetfillcolor{currentfill}%
\pgfsetfillopacity{0.130000}%
\pgfsetlinewidth{0.000000pt}%
\definecolor{currentstroke}{rgb}{0.000000,0.000000,0.000000}%
\pgfsetstrokecolor{currentstroke}%
\pgfsetdash{}{0pt}%
\pgfpathmoveto{\pgfqpoint{1.267086in}{0.822236in}}%
\pgfpathlineto{\pgfqpoint{1.283600in}{0.828585in}}%
\pgfpathlineto{\pgfqpoint{1.300181in}{0.835734in}}%
\pgfpathlineto{\pgfqpoint{1.307983in}{0.846565in}}%
\pgfpathlineto{\pgfqpoint{1.315772in}{0.857916in}}%
\pgfpathlineto{\pgfqpoint{1.299207in}{0.851329in}}%
\pgfpathlineto{\pgfqpoint{1.282710in}{0.845550in}}%
\pgfpathlineto{\pgfqpoint{1.274905in}{0.833625in}}%
\pgfpathlineto{\pgfqpoint{1.267086in}{0.822236in}}%
\pgfpathclose%
\pgfusepath{fill}%
\end{pgfscope}%
\begin{pgfscope}%
\pgfpathrectangle{\pgfqpoint{0.407592in}{0.222904in}}{\pgfqpoint{1.411163in}{1.411163in}}%
\pgfusepath{clip}%
\pgfsetbuttcap%
\pgfsetroundjoin%
\definecolor{currentfill}{rgb}{0.556863,0.654902,0.847059}%
\pgfsetfillcolor{currentfill}%
\pgfsetfillopacity{0.130000}%
\pgfsetlinewidth{0.000000pt}%
\definecolor{currentstroke}{rgb}{0.000000,0.000000,0.000000}%
\pgfsetstrokecolor{currentstroke}%
\pgfsetdash{}{0pt}%
\pgfpathmoveto{\pgfqpoint{0.569946in}{1.154195in}}%
\pgfpathlineto{\pgfqpoint{0.586025in}{1.134653in}}%
\pgfpathlineto{\pgfqpoint{0.602112in}{1.115194in}}%
\pgfpathlineto{\pgfqpoint{0.611034in}{1.116343in}}%
\pgfpathlineto{\pgfqpoint{0.619915in}{1.117823in}}%
\pgfpathlineto{\pgfqpoint{0.603870in}{1.137440in}}%
\pgfpathlineto{\pgfqpoint{0.587834in}{1.157126in}}%
\pgfpathlineto{\pgfqpoint{0.578910in}{1.155504in}}%
\pgfpathlineto{\pgfqpoint{0.569946in}{1.154195in}}%
\pgfpathclose%
\pgfusepath{fill}%
\end{pgfscope}%
\begin{pgfscope}%
\pgfpathrectangle{\pgfqpoint{0.407592in}{0.222904in}}{\pgfqpoint{1.411163in}{1.411163in}}%
\pgfusepath{clip}%
\pgfsetbuttcap%
\pgfsetroundjoin%
\definecolor{currentfill}{rgb}{0.556863,0.654902,0.847059}%
\pgfsetfillcolor{currentfill}%
\pgfsetfillopacity{0.130000}%
\pgfsetlinewidth{0.000000pt}%
\definecolor{currentstroke}{rgb}{0.000000,0.000000,0.000000}%
\pgfsetstrokecolor{currentstroke}%
\pgfsetdash{}{0pt}%
\pgfpathmoveto{\pgfqpoint{0.976071in}{0.771175in}}%
\pgfpathlineto{\pgfqpoint{0.992225in}{0.759518in}}%
\pgfpathlineto{\pgfqpoint{1.008395in}{0.749341in}}%
\pgfpathlineto{\pgfqpoint{1.016618in}{0.750384in}}%
\pgfpathlineto{\pgfqpoint{1.024806in}{0.752405in}}%
\pgfpathlineto{\pgfqpoint{1.008677in}{0.762842in}}%
\pgfpathlineto{\pgfqpoint{0.992563in}{0.774795in}}%
\pgfpathlineto{\pgfqpoint{0.984335in}{0.772531in}}%
\pgfpathlineto{\pgfqpoint{0.976071in}{0.771175in}}%
\pgfpathclose%
\pgfusepath{fill}%
\end{pgfscope}%
\begin{pgfscope}%
\pgfpathrectangle{\pgfqpoint{0.407592in}{0.222904in}}{\pgfqpoint{1.411163in}{1.411163in}}%
\pgfusepath{clip}%
\pgfsetbuttcap%
\pgfsetroundjoin%
\definecolor{currentfill}{rgb}{0.556863,0.654902,0.847059}%
\pgfsetfillcolor{currentfill}%
\pgfsetfillopacity{0.130000}%
\pgfsetlinewidth{0.000000pt}%
\definecolor{currentstroke}{rgb}{0.000000,0.000000,0.000000}%
\pgfsetstrokecolor{currentstroke}%
\pgfsetdash{}{0pt}%
\pgfpathmoveto{\pgfqpoint{0.748584in}{0.964799in}}%
\pgfpathlineto{\pgfqpoint{0.764706in}{0.946418in}}%
\pgfpathlineto{\pgfqpoint{0.780837in}{0.928288in}}%
\pgfpathlineto{\pgfqpoint{0.789466in}{0.928997in}}%
\pgfpathlineto{\pgfqpoint{0.798055in}{0.930191in}}%
\pgfpathlineto{\pgfqpoint{0.781966in}{0.948626in}}%
\pgfpathlineto{\pgfqpoint{0.765886in}{0.967286in}}%
\pgfpathlineto{\pgfqpoint{0.757255in}{0.965817in}}%
\pgfpathlineto{\pgfqpoint{0.748584in}{0.964799in}}%
\pgfpathclose%
\pgfusepath{fill}%
\end{pgfscope}%
\begin{pgfscope}%
\pgfpathrectangle{\pgfqpoint{0.407592in}{0.222904in}}{\pgfqpoint{1.411163in}{1.411163in}}%
\pgfusepath{clip}%
\pgfsetbuttcap%
\pgfsetroundjoin%
\definecolor{currentfill}{rgb}{0.556863,0.654902,0.847059}%
\pgfsetfillcolor{currentfill}%
\pgfsetfillopacity{0.130000}%
\pgfsetlinewidth{0.000000pt}%
\definecolor{currentstroke}{rgb}{0.000000,0.000000,0.000000}%
\pgfsetstrokecolor{currentstroke}%
\pgfsetdash{}{0pt}%
\pgfpathmoveto{\pgfqpoint{1.105701in}{0.736477in}}%
\pgfpathlineto{\pgfqpoint{1.121944in}{0.735796in}}%
\pgfpathlineto{\pgfqpoint{1.138235in}{0.737000in}}%
\pgfpathlineto{\pgfqpoint{1.146236in}{0.742334in}}%
\pgfpathlineto{\pgfqpoint{1.154213in}{0.748759in}}%
\pgfpathlineto{\pgfqpoint{1.137947in}{0.747584in}}%
\pgfpathlineto{\pgfqpoint{1.121729in}{0.748276in}}%
\pgfpathlineto{\pgfqpoint{1.113728in}{0.741796in}}%
\pgfpathlineto{\pgfqpoint{1.105701in}{0.736477in}}%
\pgfpathclose%
\pgfusepath{fill}%
\end{pgfscope}%
\begin{pgfscope}%
\pgfpathrectangle{\pgfqpoint{0.407592in}{0.222904in}}{\pgfqpoint{1.411163in}{1.411163in}}%
\pgfusepath{clip}%
\pgfsetbuttcap%
\pgfsetroundjoin%
\definecolor{currentfill}{rgb}{0.556863,0.654902,0.847059}%
\pgfsetfillcolor{currentfill}%
\pgfsetfillopacity{0.130000}%
\pgfsetlinewidth{0.000000pt}%
\definecolor{currentstroke}{rgb}{0.000000,0.000000,0.000000}%
\pgfsetstrokecolor{currentstroke}%
\pgfsetdash{}{0pt}%
\pgfpathmoveto{\pgfqpoint{1.380064in}{0.920672in}}%
\pgfpathlineto{\pgfqpoint{1.396802in}{0.928076in}}%
\pgfpathlineto{\pgfqpoint{1.413609in}{0.936014in}}%
\pgfpathlineto{\pgfqpoint{1.421300in}{0.948275in}}%
\pgfpathlineto{\pgfqpoint{1.428979in}{0.960848in}}%
\pgfpathlineto{\pgfqpoint{1.412189in}{0.953437in}}%
\pgfpathlineto{\pgfqpoint{1.395469in}{0.946572in}}%
\pgfpathlineto{\pgfqpoint{1.387773in}{0.933468in}}%
\pgfpathlineto{\pgfqpoint{1.380064in}{0.920672in}}%
\pgfpathclose%
\pgfusepath{fill}%
\end{pgfscope}%
\begin{pgfscope}%
\pgfpathrectangle{\pgfqpoint{0.407592in}{0.222904in}}{\pgfqpoint{1.411163in}{1.411163in}}%
\pgfusepath{clip}%
\pgfsetbuttcap%
\pgfsetroundjoin%
\definecolor{currentfill}{rgb}{0.556863,0.654902,0.847059}%
\pgfsetfillcolor{currentfill}%
\pgfsetfillopacity{0.130000}%
\pgfsetlinewidth{0.000000pt}%
\definecolor{currentstroke}{rgb}{0.000000,0.000000,0.000000}%
\pgfsetstrokecolor{currentstroke}%
\pgfsetdash{}{0pt}%
\pgfpathmoveto{\pgfqpoint{1.154213in}{0.748759in}}%
\pgfpathlineto{\pgfqpoint{1.170532in}{0.751611in}}%
\pgfpathlineto{\pgfqpoint{1.186909in}{0.755929in}}%
\pgfpathlineto{\pgfqpoint{1.194844in}{0.763221in}}%
\pgfpathlineto{\pgfqpoint{1.202760in}{0.771439in}}%
\pgfpathlineto{\pgfqpoint{1.186402in}{0.767373in}}%
\pgfpathlineto{\pgfqpoint{1.170103in}{0.764745in}}%
\pgfpathlineto{\pgfqpoint{1.162168in}{0.756246in}}%
\pgfpathlineto{\pgfqpoint{1.154213in}{0.748759in}}%
\pgfpathclose%
\pgfusepath{fill}%
\end{pgfscope}%
\begin{pgfscope}%
\pgfpathrectangle{\pgfqpoint{0.407592in}{0.222904in}}{\pgfqpoint{1.411163in}{1.411163in}}%
\pgfusepath{clip}%
\pgfsetbuttcap%
\pgfsetroundjoin%
\definecolor{currentfill}{rgb}{0.556863,0.654902,0.847059}%
\pgfsetfillcolor{currentfill}%
\pgfsetfillopacity{0.130000}%
\pgfsetlinewidth{0.000000pt}%
\definecolor{currentstroke}{rgb}{0.000000,0.000000,0.000000}%
\pgfsetstrokecolor{currentstroke}%
\pgfsetdash{}{0pt}%
\pgfpathmoveto{\pgfqpoint{1.557571in}{1.096981in}}%
\pgfpathlineto{\pgfqpoint{1.574644in}{1.104752in}}%
\pgfpathlineto{\pgfqpoint{1.591791in}{1.112895in}}%
\pgfpathlineto{\pgfqpoint{1.599310in}{1.126143in}}%
\pgfpathlineto{\pgfqpoint{1.606818in}{1.139560in}}%
\pgfpathlineto{\pgfqpoint{1.589692in}{1.131815in}}%
\pgfpathlineto{\pgfqpoint{1.572637in}{1.124449in}}%
\pgfpathlineto{\pgfqpoint{1.565110in}{1.110635in}}%
\pgfpathlineto{\pgfqpoint{1.557571in}{1.096981in}}%
\pgfpathclose%
\pgfusepath{fill}%
\end{pgfscope}%
\begin{pgfscope}%
\pgfpathrectangle{\pgfqpoint{0.407592in}{0.222904in}}{\pgfqpoint{1.411163in}{1.411163in}}%
\pgfusepath{clip}%
\pgfsetbuttcap%
\pgfsetroundjoin%
\definecolor{currentfill}{rgb}{0.556863,0.654902,0.847059}%
\pgfsetfillcolor{currentfill}%
\pgfsetfillopacity{0.130000}%
\pgfsetlinewidth{0.000000pt}%
\definecolor{currentstroke}{rgb}{0.000000,0.000000,0.000000}%
\pgfsetstrokecolor{currentstroke}%
\pgfsetdash{}{0pt}%
\pgfpathmoveto{\pgfqpoint{0.602112in}{1.115194in}}%
\pgfpathlineto{\pgfqpoint{0.618208in}{1.095828in}}%
\pgfpathlineto{\pgfqpoint{0.634313in}{1.076567in}}%
\pgfpathlineto{\pgfqpoint{0.643192in}{1.077527in}}%
\pgfpathlineto{\pgfqpoint{0.652031in}{1.078837in}}%
\pgfpathlineto{\pgfqpoint{0.635969in}{1.098286in}}%
\pgfpathlineto{\pgfqpoint{0.619915in}{1.117823in}}%
\pgfpathlineto{\pgfqpoint{0.611034in}{1.116343in}}%
\pgfpathlineto{\pgfqpoint{0.602112in}{1.115194in}}%
\pgfpathclose%
\pgfusepath{fill}%
\end{pgfscope}%
\begin{pgfscope}%
\pgfpathrectangle{\pgfqpoint{0.407592in}{0.222904in}}{\pgfqpoint{1.411163in}{1.411163in}}%
\pgfusepath{clip}%
\pgfsetbuttcap%
\pgfsetroundjoin%
\definecolor{currentfill}{rgb}{0.556863,0.654902,0.847059}%
\pgfsetfillcolor{currentfill}%
\pgfsetfillopacity{0.130000}%
\pgfsetlinewidth{0.000000pt}%
\definecolor{currentstroke}{rgb}{0.000000,0.000000,0.000000}%
\pgfsetstrokecolor{currentstroke}%
\pgfsetdash{}{0pt}%
\pgfpathmoveto{\pgfqpoint{0.894767in}{0.826919in}}%
\pgfpathlineto{\pgfqpoint{0.910917in}{0.811580in}}%
\pgfpathlineto{\pgfqpoint{0.927077in}{0.797044in}}%
\pgfpathlineto{\pgfqpoint{0.935458in}{0.797237in}}%
\pgfpathlineto{\pgfqpoint{0.943801in}{0.798165in}}%
\pgfpathlineto{\pgfqpoint{0.927683in}{0.813142in}}%
\pgfpathlineto{\pgfqpoint{0.911574in}{0.828902in}}%
\pgfpathlineto{\pgfqpoint{0.903190in}{0.827572in}}%
\pgfpathlineto{\pgfqpoint{0.894767in}{0.826919in}}%
\pgfpathclose%
\pgfusepath{fill}%
\end{pgfscope}%
\begin{pgfscope}%
\pgfpathrectangle{\pgfqpoint{0.407592in}{0.222904in}}{\pgfqpoint{1.411163in}{1.411163in}}%
\pgfusepath{clip}%
\pgfsetbuttcap%
\pgfsetroundjoin%
\definecolor{currentfill}{rgb}{0.556863,0.654902,0.847059}%
\pgfsetfillcolor{currentfill}%
\pgfsetfillopacity{0.130000}%
\pgfsetlinewidth{0.000000pt}%
\definecolor{currentstroke}{rgb}{0.000000,0.000000,0.000000}%
\pgfsetstrokecolor{currentstroke}%
\pgfsetdash{}{0pt}%
\pgfpathmoveto{\pgfqpoint{1.057128in}{0.736877in}}%
\pgfpathlineto{\pgfqpoint{1.073329in}{0.732072in}}%
\pgfpathlineto{\pgfqpoint{1.089563in}{0.729307in}}%
\pgfpathlineto{\pgfqpoint{1.097646in}{0.732322in}}%
\pgfpathlineto{\pgfqpoint{1.105701in}{0.736477in}}%
\pgfpathlineto{\pgfqpoint{1.089498in}{0.739185in}}%
\pgfpathlineto{\pgfqpoint{1.073331in}{0.743981in}}%
\pgfpathlineto{\pgfqpoint{1.065245in}{0.739857in}}%
\pgfpathlineto{\pgfqpoint{1.057128in}{0.736877in}}%
\pgfpathclose%
\pgfusepath{fill}%
\end{pgfscope}%
\begin{pgfscope}%
\pgfpathrectangle{\pgfqpoint{0.407592in}{0.222904in}}{\pgfqpoint{1.411163in}{1.411163in}}%
\pgfusepath{clip}%
\pgfsetbuttcap%
\pgfsetroundjoin%
\definecolor{currentfill}{rgb}{0.556863,0.654902,0.847059}%
\pgfsetfillcolor{currentfill}%
\pgfsetfillopacity{0.130000}%
\pgfsetlinewidth{0.000000pt}%
\definecolor{currentstroke}{rgb}{0.000000,0.000000,0.000000}%
\pgfsetstrokecolor{currentstroke}%
\pgfsetdash{}{0pt}%
\pgfpathmoveto{\pgfqpoint{0.780837in}{0.928288in}}%
\pgfpathlineto{\pgfqpoint{0.796976in}{0.910448in}}%
\pgfpathlineto{\pgfqpoint{0.813124in}{0.892943in}}%
\pgfpathlineto{\pgfqpoint{0.821710in}{0.893285in}}%
\pgfpathlineto{\pgfqpoint{0.830257in}{0.894149in}}%
\pgfpathlineto{\pgfqpoint{0.814151in}{0.912018in}}%
\pgfpathlineto{\pgfqpoint{0.798055in}{0.930191in}}%
\pgfpathlineto{\pgfqpoint{0.789466in}{0.928997in}}%
\pgfpathlineto{\pgfqpoint{0.780837in}{0.928288in}}%
\pgfpathclose%
\pgfusepath{fill}%
\end{pgfscope}%
\begin{pgfscope}%
\pgfpathrectangle{\pgfqpoint{0.407592in}{0.222904in}}{\pgfqpoint{1.411163in}{1.411163in}}%
\pgfusepath{clip}%
\pgfsetbuttcap%
\pgfsetroundjoin%
\definecolor{currentfill}{rgb}{0.556863,0.654902,0.847059}%
\pgfsetfillcolor{currentfill}%
\pgfsetfillopacity{0.130000}%
\pgfsetlinewidth{0.000000pt}%
\definecolor{currentstroke}{rgb}{0.000000,0.000000,0.000000}%
\pgfsetstrokecolor{currentstroke}%
\pgfsetdash{}{0pt}%
\pgfpathmoveto{\pgfqpoint{1.315772in}{0.857916in}}%
\pgfpathlineto{\pgfqpoint{1.332405in}{0.865216in}}%
\pgfpathlineto{\pgfqpoint{1.349107in}{0.873144in}}%
\pgfpathlineto{\pgfqpoint{1.356865in}{0.884415in}}%
\pgfpathlineto{\pgfqpoint{1.364611in}{0.896115in}}%
\pgfpathlineto{\pgfqpoint{1.347926in}{0.888740in}}%
\pgfpathlineto{\pgfqpoint{1.331310in}{0.882011in}}%
\pgfpathlineto{\pgfqpoint{1.323547in}{0.869745in}}%
\pgfpathlineto{\pgfqpoint{1.315772in}{0.857916in}}%
\pgfpathclose%
\pgfusepath{fill}%
\end{pgfscope}%
\begin{pgfscope}%
\pgfpathrectangle{\pgfqpoint{0.407592in}{0.222904in}}{\pgfqpoint{1.411163in}{1.411163in}}%
\pgfusepath{clip}%
\pgfsetbuttcap%
\pgfsetroundjoin%
\definecolor{currentfill}{rgb}{0.556863,0.654902,0.847059}%
\pgfsetfillcolor{currentfill}%
\pgfsetfillopacity{0.130000}%
\pgfsetlinewidth{0.000000pt}%
\definecolor{currentstroke}{rgb}{0.000000,0.000000,0.000000}%
\pgfsetstrokecolor{currentstroke}%
\pgfsetdash{}{0pt}%
\pgfpathmoveto{\pgfqpoint{1.493298in}{1.028128in}}%
\pgfpathlineto{\pgfqpoint{1.510265in}{1.035964in}}%
\pgfpathlineto{\pgfqpoint{1.527304in}{1.044208in}}%
\pgfpathlineto{\pgfqpoint{1.534888in}{1.057099in}}%
\pgfpathlineto{\pgfqpoint{1.542460in}{1.070200in}}%
\pgfpathlineto{\pgfqpoint{1.525440in}{1.062400in}}%
\pgfpathlineto{\pgfqpoint{1.508491in}{1.055016in}}%
\pgfpathlineto{\pgfqpoint{1.500900in}{1.041471in}}%
\pgfpathlineto{\pgfqpoint{1.493298in}{1.028128in}}%
\pgfpathclose%
\pgfusepath{fill}%
\end{pgfscope}%
\begin{pgfscope}%
\pgfpathrectangle{\pgfqpoint{0.407592in}{0.222904in}}{\pgfqpoint{1.411163in}{1.411163in}}%
\pgfusepath{clip}%
\pgfsetbuttcap%
\pgfsetroundjoin%
\definecolor{currentfill}{rgb}{0.556863,0.654902,0.847059}%
\pgfsetfillcolor{currentfill}%
\pgfsetfillopacity{0.130000}%
\pgfsetlinewidth{0.000000pt}%
\definecolor{currentstroke}{rgb}{0.000000,0.000000,0.000000}%
\pgfsetstrokecolor{currentstroke}%
\pgfsetdash{}{0pt}%
\pgfpathmoveto{\pgfqpoint{1.202760in}{0.771439in}}%
\pgfpathlineto{\pgfqpoint{1.219178in}{0.776751in}}%
\pgfpathlineto{\pgfqpoint{1.235661in}{0.783126in}}%
\pgfpathlineto{\pgfqpoint{1.243541in}{0.791830in}}%
\pgfpathlineto{\pgfqpoint{1.251404in}{0.801285in}}%
\pgfpathlineto{\pgfqpoint{1.234939in}{0.795332in}}%
\pgfpathlineto{\pgfqpoint{1.218538in}{0.790440in}}%
\pgfpathlineto{\pgfqpoint{1.210657in}{0.780533in}}%
\pgfpathlineto{\pgfqpoint{1.202760in}{0.771439in}}%
\pgfpathclose%
\pgfusepath{fill}%
\end{pgfscope}%
\begin{pgfscope}%
\pgfpathrectangle{\pgfqpoint{0.407592in}{0.222904in}}{\pgfqpoint{1.411163in}{1.411163in}}%
\pgfusepath{clip}%
\pgfsetbuttcap%
\pgfsetroundjoin%
\definecolor{currentfill}{rgb}{0.556863,0.654902,0.847059}%
\pgfsetfillcolor{currentfill}%
\pgfsetfillopacity{0.130000}%
\pgfsetlinewidth{0.000000pt}%
\definecolor{currentstroke}{rgb}{0.000000,0.000000,0.000000}%
\pgfsetstrokecolor{currentstroke}%
\pgfsetdash{}{0pt}%
\pgfpathmoveto{\pgfqpoint{0.634313in}{1.076567in}}%
\pgfpathlineto{\pgfqpoint{0.650426in}{1.057423in}}%
\pgfpathlineto{\pgfqpoint{0.666548in}{1.038411in}}%
\pgfpathlineto{\pgfqpoint{0.675385in}{1.039146in}}%
\pgfpathlineto{\pgfqpoint{0.684181in}{1.040255in}}%
\pgfpathlineto{\pgfqpoint{0.668101in}{1.059489in}}%
\pgfpathlineto{\pgfqpoint{0.652031in}{1.078837in}}%
\pgfpathlineto{\pgfqpoint{0.643192in}{1.077527in}}%
\pgfpathlineto{\pgfqpoint{0.634313in}{1.076567in}}%
\pgfpathclose%
\pgfusepath{fill}%
\end{pgfscope}%
\begin{pgfscope}%
\pgfpathrectangle{\pgfqpoint{0.407592in}{0.222904in}}{\pgfqpoint{1.411163in}{1.411163in}}%
\pgfusepath{clip}%
\pgfsetbuttcap%
\pgfsetroundjoin%
\definecolor{currentfill}{rgb}{0.556863,0.654902,0.847059}%
\pgfsetfillcolor{currentfill}%
\pgfsetfillopacity{0.130000}%
\pgfsetlinewidth{0.000000pt}%
\definecolor{currentstroke}{rgb}{0.000000,0.000000,0.000000}%
\pgfsetstrokecolor{currentstroke}%
\pgfsetdash{}{0pt}%
\pgfpathmoveto{\pgfqpoint{1.428979in}{0.960848in}}%
\pgfpathlineto{\pgfqpoint{1.445840in}{0.968756in}}%
\pgfpathlineto{\pgfqpoint{1.462772in}{0.977118in}}%
\pgfpathlineto{\pgfqpoint{1.470421in}{0.989481in}}%
\pgfpathlineto{\pgfqpoint{1.478058in}{1.002116in}}%
\pgfpathlineto{\pgfqpoint{1.461144in}{0.994243in}}%
\pgfpathlineto{\pgfqpoint{1.444301in}{0.986838in}}%
\pgfpathlineto{\pgfqpoint{1.436646in}{0.973709in}}%
\pgfpathlineto{\pgfqpoint{1.428979in}{0.960848in}}%
\pgfpathclose%
\pgfusepath{fill}%
\end{pgfscope}%
\begin{pgfscope}%
\pgfpathrectangle{\pgfqpoint{0.407592in}{0.222904in}}{\pgfqpoint{1.411163in}{1.411163in}}%
\pgfusepath{clip}%
\pgfsetbuttcap%
\pgfsetroundjoin%
\definecolor{currentfill}{rgb}{0.556863,0.654902,0.847059}%
\pgfsetfillcolor{currentfill}%
\pgfsetfillopacity{0.130000}%
\pgfsetlinewidth{0.000000pt}%
\definecolor{currentstroke}{rgb}{0.000000,0.000000,0.000000}%
\pgfsetstrokecolor{currentstroke}%
\pgfsetdash{}{0pt}%
\pgfpathmoveto{\pgfqpoint{1.008395in}{0.749341in}}%
\pgfpathlineto{\pgfqpoint{1.024584in}{0.740825in}}%
\pgfpathlineto{\pgfqpoint{1.040795in}{0.734121in}}%
\pgfpathlineto{\pgfqpoint{1.048978in}{0.734990in}}%
\pgfpathlineto{\pgfqpoint{1.057128in}{0.736877in}}%
\pgfpathlineto{\pgfqpoint{1.040956in}{0.743693in}}%
\pgfpathlineto{\pgfqpoint{1.024806in}{0.752405in}}%
\pgfpathlineto{\pgfqpoint{1.016618in}{0.750384in}}%
\pgfpathlineto{\pgfqpoint{1.008395in}{0.749341in}}%
\pgfpathclose%
\pgfusepath{fill}%
\end{pgfscope}%
\begin{pgfscope}%
\pgfpathrectangle{\pgfqpoint{0.407592in}{0.222904in}}{\pgfqpoint{1.411163in}{1.411163in}}%
\pgfusepath{clip}%
\pgfsetbuttcap%
\pgfsetroundjoin%
\definecolor{currentfill}{rgb}{0.556863,0.654902,0.847059}%
\pgfsetfillcolor{currentfill}%
\pgfsetfillopacity{0.130000}%
\pgfsetlinewidth{0.000000pt}%
\definecolor{currentstroke}{rgb}{0.000000,0.000000,0.000000}%
\pgfsetstrokecolor{currentstroke}%
\pgfsetdash{}{0pt}%
\pgfpathmoveto{\pgfqpoint{0.927077in}{0.797044in}}%
\pgfpathlineto{\pgfqpoint{0.943248in}{0.783450in}}%
\pgfpathlineto{\pgfqpoint{0.959431in}{0.770949in}}%
\pgfpathlineto{\pgfqpoint{0.967770in}{0.770670in}}%
\pgfpathlineto{\pgfqpoint{0.976071in}{0.771175in}}%
\pgfpathlineto{\pgfqpoint{0.959930in}{0.784119in}}%
\pgfpathlineto{\pgfqpoint{0.943801in}{0.798165in}}%
\pgfpathlineto{\pgfqpoint{0.935458in}{0.797237in}}%
\pgfpathlineto{\pgfqpoint{0.927077in}{0.797044in}}%
\pgfpathclose%
\pgfusepath{fill}%
\end{pgfscope}%
\begin{pgfscope}%
\pgfpathrectangle{\pgfqpoint{0.407592in}{0.222904in}}{\pgfqpoint{1.411163in}{1.411163in}}%
\pgfusepath{clip}%
\pgfsetbuttcap%
\pgfsetroundjoin%
\definecolor{currentfill}{rgb}{0.556863,0.654902,0.847059}%
\pgfsetfillcolor{currentfill}%
\pgfsetfillopacity{0.130000}%
\pgfsetlinewidth{0.000000pt}%
\definecolor{currentstroke}{rgb}{0.000000,0.000000,0.000000}%
\pgfsetstrokecolor{currentstroke}%
\pgfsetdash{}{0pt}%
\pgfpathmoveto{\pgfqpoint{0.813124in}{0.892943in}}%
\pgfpathlineto{\pgfqpoint{0.829280in}{0.875827in}}%
\pgfpathlineto{\pgfqpoint{0.845445in}{0.859161in}}%
\pgfpathlineto{\pgfqpoint{0.853990in}{0.859072in}}%
\pgfpathlineto{\pgfqpoint{0.862494in}{0.859543in}}%
\pgfpathlineto{\pgfqpoint{0.846371in}{0.876637in}}%
\pgfpathlineto{\pgfqpoint{0.830257in}{0.894149in}}%
\pgfpathlineto{\pgfqpoint{0.821710in}{0.893285in}}%
\pgfpathlineto{\pgfqpoint{0.813124in}{0.892943in}}%
\pgfpathclose%
\pgfusepath{fill}%
\end{pgfscope}%
\begin{pgfscope}%
\pgfpathrectangle{\pgfqpoint{0.407592in}{0.222904in}}{\pgfqpoint{1.411163in}{1.411163in}}%
\pgfusepath{clip}%
\pgfsetbuttcap%
\pgfsetroundjoin%
\definecolor{currentfill}{rgb}{0.556863,0.654902,0.847059}%
\pgfsetfillcolor{currentfill}%
\pgfsetfillopacity{0.130000}%
\pgfsetlinewidth{0.000000pt}%
\definecolor{currentstroke}{rgb}{0.000000,0.000000,0.000000}%
\pgfsetstrokecolor{currentstroke}%
\pgfsetdash{}{0pt}%
\pgfpathmoveto{\pgfqpoint{0.666548in}{1.038411in}}%
\pgfpathlineto{\pgfqpoint{0.682678in}{1.019548in}}%
\pgfpathlineto{\pgfqpoint{0.698817in}{1.000853in}}%
\pgfpathlineto{\pgfqpoint{0.707612in}{1.001324in}}%
\pgfpathlineto{\pgfqpoint{0.716365in}{1.002191in}}%
\pgfpathlineto{\pgfqpoint{0.700269in}{1.021150in}}%
\pgfpathlineto{\pgfqpoint{0.684181in}{1.040255in}}%
\pgfpathlineto{\pgfqpoint{0.675385in}{1.039146in}}%
\pgfpathlineto{\pgfqpoint{0.666548in}{1.038411in}}%
\pgfpathclose%
\pgfusepath{fill}%
\end{pgfscope}%
\begin{pgfscope}%
\pgfpathrectangle{\pgfqpoint{0.407592in}{0.222904in}}{\pgfqpoint{1.411163in}{1.411163in}}%
\pgfusepath{clip}%
\pgfsetbuttcap%
\pgfsetroundjoin%
\definecolor{currentfill}{rgb}{0.556863,0.654902,0.847059}%
\pgfsetfillcolor{currentfill}%
\pgfsetfillopacity{0.130000}%
\pgfsetlinewidth{0.000000pt}%
\definecolor{currentstroke}{rgb}{0.000000,0.000000,0.000000}%
\pgfsetstrokecolor{currentstroke}%
\pgfsetdash{}{0pt}%
\pgfpathmoveto{\pgfqpoint{1.251404in}{0.801285in}}%
\pgfpathlineto{\pgfqpoint{1.267935in}{0.808151in}}%
\pgfpathlineto{\pgfqpoint{1.284534in}{0.815801in}}%
\pgfpathlineto{\pgfqpoint{1.292365in}{0.825465in}}%
\pgfpathlineto{\pgfqpoint{1.300181in}{0.835734in}}%
\pgfpathlineto{\pgfqpoint{1.283600in}{0.828585in}}%
\pgfpathlineto{\pgfqpoint{1.267086in}{0.822236in}}%
\pgfpathlineto{\pgfqpoint{1.259252in}{0.811438in}}%
\pgfpathlineto{\pgfqpoint{1.251404in}{0.801285in}}%
\pgfpathclose%
\pgfusepath{fill}%
\end{pgfscope}%
\begin{pgfscope}%
\pgfpathrectangle{\pgfqpoint{0.407592in}{0.222904in}}{\pgfqpoint{1.411163in}{1.411163in}}%
\pgfusepath{clip}%
\pgfsetbuttcap%
\pgfsetroundjoin%
\definecolor{currentfill}{rgb}{0.556863,0.654902,0.847059}%
\pgfsetfillcolor{currentfill}%
\pgfsetfillopacity{0.130000}%
\pgfsetlinewidth{0.000000pt}%
\definecolor{currentstroke}{rgb}{0.000000,0.000000,0.000000}%
\pgfsetstrokecolor{currentstroke}%
\pgfsetdash{}{0pt}%
\pgfpathmoveto{\pgfqpoint{1.364611in}{0.896115in}}%
\pgfpathlineto{\pgfqpoint{1.381366in}{0.904064in}}%
\pgfpathlineto{\pgfqpoint{1.398191in}{0.912528in}}%
\pgfpathlineto{\pgfqpoint{1.405906in}{0.924090in}}%
\pgfpathlineto{\pgfqpoint{1.413609in}{0.936014in}}%
\pgfpathlineto{\pgfqpoint{1.396802in}{0.928076in}}%
\pgfpathlineto{\pgfqpoint{1.380064in}{0.920672in}}%
\pgfpathlineto{\pgfqpoint{1.372344in}{0.908210in}}%
\pgfpathlineto{\pgfqpoint{1.364611in}{0.896115in}}%
\pgfpathclose%
\pgfusepath{fill}%
\end{pgfscope}%
\begin{pgfscope}%
\pgfpathrectangle{\pgfqpoint{0.407592in}{0.222904in}}{\pgfqpoint{1.411163in}{1.411163in}}%
\pgfusepath{clip}%
\pgfsetbuttcap%
\pgfsetroundjoin%
\definecolor{currentfill}{rgb}{0.556863,0.654902,0.847059}%
\pgfsetfillcolor{currentfill}%
\pgfsetfillopacity{0.130000}%
\pgfsetlinewidth{0.000000pt}%
\definecolor{currentstroke}{rgb}{0.000000,0.000000,0.000000}%
\pgfsetstrokecolor{currentstroke}%
\pgfsetdash{}{0pt}%
\pgfpathmoveto{\pgfqpoint{1.606818in}{1.139560in}}%
\pgfpathlineto{\pgfqpoint{1.624018in}{1.147661in}}%
\pgfpathlineto{\pgfqpoint{1.641291in}{1.156100in}}%
\pgfpathlineto{\pgfqpoint{1.648768in}{1.169296in}}%
\pgfpathlineto{\pgfqpoint{1.631505in}{1.161045in}}%
\pgfpathlineto{\pgfqpoint{1.614316in}{1.153133in}}%
\pgfusepath{fill}%
\end{pgfscope}%
\begin{pgfscope}%
\pgfpathrectangle{\pgfqpoint{0.407592in}{0.222904in}}{\pgfqpoint{1.411163in}{1.411163in}}%
\pgfusepath{clip}%
\pgfsetbuttcap%
\pgfsetroundjoin%
\definecolor{currentfill}{rgb}{0.556863,0.654902,0.847059}%
\pgfsetfillcolor{currentfill}%
\pgfsetfillopacity{0.130000}%
\pgfsetlinewidth{0.000000pt}%
\definecolor{currentstroke}{rgb}{0.000000,0.000000,0.000000}%
\pgfsetstrokecolor{currentstroke}%
\pgfsetdash{}{0pt}%
\pgfpathmoveto{\pgfqpoint{1.542460in}{1.070200in}}%
\pgfpathlineto{\pgfqpoint{1.559553in}{1.078391in}}%
\pgfpathlineto{\pgfqpoint{1.576718in}{1.086945in}}%
\pgfpathlineto{\pgfqpoint{1.584260in}{1.099825in}}%
\pgfpathlineto{\pgfqpoint{1.591791in}{1.112895in}}%
\pgfpathlineto{\pgfqpoint{1.574644in}{1.104752in}}%
\pgfpathlineto{\pgfqpoint{1.557571in}{1.096981in}}%
\pgfpathlineto{\pgfqpoint{1.550021in}{1.083499in}}%
\pgfpathlineto{\pgfqpoint{1.542460in}{1.070200in}}%
\pgfpathclose%
\pgfusepath{fill}%
\end{pgfscope}%
\begin{pgfscope}%
\pgfpathrectangle{\pgfqpoint{0.407592in}{0.222904in}}{\pgfqpoint{1.411163in}{1.411163in}}%
\pgfusepath{clip}%
\pgfsetbuttcap%
\pgfsetroundjoin%
\definecolor{currentfill}{rgb}{0.556863,0.654902,0.847059}%
\pgfsetfillcolor{currentfill}%
\pgfsetfillopacity{0.130000}%
\pgfsetlinewidth{0.000000pt}%
\definecolor{currentstroke}{rgb}{0.000000,0.000000,0.000000}%
\pgfsetstrokecolor{currentstroke}%
\pgfsetdash{}{0pt}%
\pgfpathmoveto{\pgfqpoint{0.698817in}{1.000853in}}%
\pgfpathlineto{\pgfqpoint{0.714965in}{0.982348in}}%
\pgfpathlineto{\pgfqpoint{0.731121in}{0.964059in}}%
\pgfpathlineto{\pgfqpoint{0.739873in}{0.964218in}}%
\pgfpathlineto{\pgfqpoint{0.748584in}{0.964799in}}%
\pgfpathlineto{\pgfqpoint{0.732470in}{0.983400in}}%
\pgfpathlineto{\pgfqpoint{0.716365in}{1.002191in}}%
\pgfpathlineto{\pgfqpoint{0.707612in}{1.001324in}}%
\pgfpathlineto{\pgfqpoint{0.698817in}{1.000853in}}%
\pgfpathclose%
\pgfusepath{fill}%
\end{pgfscope}%
\begin{pgfscope}%
\pgfpathrectangle{\pgfqpoint{0.407592in}{0.222904in}}{\pgfqpoint{1.411163in}{1.411163in}}%
\pgfusepath{clip}%
\pgfsetbuttcap%
\pgfsetroundjoin%
\definecolor{currentfill}{rgb}{0.556863,0.654902,0.847059}%
\pgfsetfillcolor{currentfill}%
\pgfsetfillopacity{0.130000}%
\pgfsetlinewidth{0.000000pt}%
\definecolor{currentstroke}{rgb}{0.000000,0.000000,0.000000}%
\pgfsetstrokecolor{currentstroke}%
\pgfsetdash{}{0pt}%
\pgfpathmoveto{\pgfqpoint{0.845445in}{0.859161in}}%
\pgfpathlineto{\pgfqpoint{0.861619in}{0.843019in}}%
\pgfpathlineto{\pgfqpoint{0.877802in}{0.827486in}}%
\pgfpathlineto{\pgfqpoint{0.886304in}{0.826904in}}%
\pgfpathlineto{\pgfqpoint{0.894767in}{0.826919in}}%
\pgfpathlineto{\pgfqpoint{0.878626in}{0.842941in}}%
\pgfpathlineto{\pgfqpoint{0.862494in}{0.859543in}}%
\pgfpathlineto{\pgfqpoint{0.853990in}{0.859072in}}%
\pgfpathlineto{\pgfqpoint{0.845445in}{0.859161in}}%
\pgfpathclose%
\pgfusepath{fill}%
\end{pgfscope}%
\begin{pgfscope}%
\pgfpathrectangle{\pgfqpoint{0.407592in}{0.222904in}}{\pgfqpoint{1.411163in}{1.411163in}}%
\pgfusepath{clip}%
\pgfsetbuttcap%
\pgfsetroundjoin%
\definecolor{currentfill}{rgb}{0.556863,0.654902,0.847059}%
\pgfsetfillcolor{currentfill}%
\pgfsetfillopacity{0.130000}%
\pgfsetlinewidth{0.000000pt}%
\definecolor{currentstroke}{rgb}{0.000000,0.000000,0.000000}%
\pgfsetstrokecolor{currentstroke}%
\pgfsetdash{}{0pt}%
\pgfpathmoveto{\pgfqpoint{1.478058in}{1.002116in}}%
\pgfpathlineto{\pgfqpoint{1.495044in}{1.010418in}}%
\pgfpathlineto{\pgfqpoint{1.512102in}{1.019116in}}%
\pgfpathlineto{\pgfqpoint{1.519709in}{1.031542in}}%
\pgfpathlineto{\pgfqpoint{1.527304in}{1.044208in}}%
\pgfpathlineto{\pgfqpoint{1.510265in}{1.035964in}}%
\pgfpathlineto{\pgfqpoint{1.493298in}{1.028128in}}%
\pgfpathlineto{\pgfqpoint{1.485684in}{1.015004in}}%
\pgfpathlineto{\pgfqpoint{1.478058in}{1.002116in}}%
\pgfpathclose%
\pgfusepath{fill}%
\end{pgfscope}%
\begin{pgfscope}%
\pgfpathrectangle{\pgfqpoint{0.407592in}{0.222904in}}{\pgfqpoint{1.411163in}{1.411163in}}%
\pgfusepath{clip}%
\pgfsetbuttcap%
\pgfsetroundjoin%
\definecolor{currentfill}{rgb}{0.556863,0.654902,0.847059}%
\pgfsetfillcolor{currentfill}%
\pgfsetfillopacity{0.130000}%
\pgfsetlinewidth{0.000000pt}%
\definecolor{currentstroke}{rgb}{0.000000,0.000000,0.000000}%
\pgfsetstrokecolor{currentstroke}%
\pgfsetdash{}{0pt}%
\pgfpathmoveto{\pgfqpoint{0.551894in}{1.152493in}}%
\pgfpathlineto{\pgfqpoint{0.568015in}{1.133129in}}%
\pgfpathlineto{\pgfqpoint{0.584145in}{1.113863in}}%
\pgfpathlineto{\pgfqpoint{0.593149in}{1.114369in}}%
\pgfpathlineto{\pgfqpoint{0.602112in}{1.115194in}}%
\pgfpathlineto{\pgfqpoint{0.586025in}{1.134653in}}%
\pgfpathlineto{\pgfqpoint{0.569946in}{1.154195in}}%
\pgfpathlineto{\pgfqpoint{0.560940in}{1.153193in}}%
\pgfpathlineto{\pgfqpoint{0.551894in}{1.152493in}}%
\pgfpathclose%
\pgfusepath{fill}%
\end{pgfscope}%
\begin{pgfscope}%
\pgfpathrectangle{\pgfqpoint{0.407592in}{0.222904in}}{\pgfqpoint{1.411163in}{1.411163in}}%
\pgfusepath{clip}%
\pgfsetbuttcap%
\pgfsetroundjoin%
\definecolor{currentfill}{rgb}{0.556863,0.654902,0.847059}%
\pgfsetfillcolor{currentfill}%
\pgfsetfillopacity{0.130000}%
\pgfsetlinewidth{0.000000pt}%
\definecolor{currentstroke}{rgb}{0.000000,0.000000,0.000000}%
\pgfsetstrokecolor{currentstroke}%
\pgfsetdash{}{0pt}%
\pgfpathmoveto{\pgfqpoint{1.138235in}{0.737000in}}%
\pgfpathlineto{\pgfqpoint{1.154578in}{0.739893in}}%
\pgfpathlineto{\pgfqpoint{1.170978in}{0.744261in}}%
\pgfpathlineto{\pgfqpoint{1.178955in}{0.749601in}}%
\pgfpathlineto{\pgfqpoint{1.186909in}{0.755929in}}%
\pgfpathlineto{\pgfqpoint{1.170532in}{0.751611in}}%
\pgfpathlineto{\pgfqpoint{1.154213in}{0.748759in}}%
\pgfpathlineto{\pgfqpoint{1.146236in}{0.742334in}}%
\pgfpathlineto{\pgfqpoint{1.138235in}{0.737000in}}%
\pgfpathclose%
\pgfusepath{fill}%
\end{pgfscope}%
\begin{pgfscope}%
\pgfpathrectangle{\pgfqpoint{0.407592in}{0.222904in}}{\pgfqpoint{1.411163in}{1.411163in}}%
\pgfusepath{clip}%
\pgfsetbuttcap%
\pgfsetroundjoin%
\definecolor{currentfill}{rgb}{0.556863,0.654902,0.847059}%
\pgfsetfillcolor{currentfill}%
\pgfsetfillopacity{0.130000}%
\pgfsetlinewidth{0.000000pt}%
\definecolor{currentstroke}{rgb}{0.000000,0.000000,0.000000}%
\pgfsetstrokecolor{currentstroke}%
\pgfsetdash{}{0pt}%
\pgfpathmoveto{\pgfqpoint{1.089563in}{0.729307in}}%
\pgfpathlineto{\pgfqpoint{1.105838in}{0.728526in}}%
\pgfpathlineto{\pgfqpoint{1.122158in}{0.729595in}}%
\pgfpathlineto{\pgfqpoint{1.130210in}{0.732760in}}%
\pgfpathlineto{\pgfqpoint{1.138235in}{0.737000in}}%
\pgfpathlineto{\pgfqpoint{1.121944in}{0.735796in}}%
\pgfpathlineto{\pgfqpoint{1.105701in}{0.736477in}}%
\pgfpathlineto{\pgfqpoint{1.097646in}{0.732322in}}%
\pgfpathlineto{\pgfqpoint{1.089563in}{0.729307in}}%
\pgfpathclose%
\pgfusepath{fill}%
\end{pgfscope}%
\begin{pgfscope}%
\pgfpathrectangle{\pgfqpoint{0.407592in}{0.222904in}}{\pgfqpoint{1.411163in}{1.411163in}}%
\pgfusepath{clip}%
\pgfsetbuttcap%
\pgfsetroundjoin%
\definecolor{currentfill}{rgb}{0.556863,0.654902,0.847059}%
\pgfsetfillcolor{currentfill}%
\pgfsetfillopacity{0.130000}%
\pgfsetlinewidth{0.000000pt}%
\definecolor{currentstroke}{rgb}{0.000000,0.000000,0.000000}%
\pgfsetstrokecolor{currentstroke}%
\pgfsetdash{}{0pt}%
\pgfpathmoveto{\pgfqpoint{1.300181in}{0.835734in}}%
\pgfpathlineto{\pgfqpoint{1.316831in}{0.843577in}}%
\pgfpathlineto{\pgfqpoint{1.333551in}{0.852024in}}%
\pgfpathlineto{\pgfqpoint{1.341335in}{0.862336in}}%
\pgfpathlineto{\pgfqpoint{1.349107in}{0.873144in}}%
\pgfpathlineto{\pgfqpoint{1.332405in}{0.865216in}}%
\pgfpathlineto{\pgfqpoint{1.315772in}{0.857916in}}%
\pgfpathlineto{\pgfqpoint{1.307983in}{0.846565in}}%
\pgfpathlineto{\pgfqpoint{1.300181in}{0.835734in}}%
\pgfpathclose%
\pgfusepath{fill}%
\end{pgfscope}%
\begin{pgfscope}%
\pgfpathrectangle{\pgfqpoint{0.407592in}{0.222904in}}{\pgfqpoint{1.411163in}{1.411163in}}%
\pgfusepath{clip}%
\pgfsetbuttcap%
\pgfsetroundjoin%
\definecolor{currentfill}{rgb}{0.556863,0.654902,0.847059}%
\pgfsetfillcolor{currentfill}%
\pgfsetfillopacity{0.130000}%
\pgfsetlinewidth{0.000000pt}%
\definecolor{currentstroke}{rgb}{0.000000,0.000000,0.000000}%
\pgfsetstrokecolor{currentstroke}%
\pgfsetdash{}{0pt}%
\pgfpathmoveto{\pgfqpoint{0.959431in}{0.770949in}}%
\pgfpathlineto{\pgfqpoint{0.975627in}{0.759706in}}%
\pgfpathlineto{\pgfqpoint{0.991839in}{0.749886in}}%
\pgfpathlineto{\pgfqpoint{1.000136in}{0.749202in}}%
\pgfpathlineto{\pgfqpoint{1.008395in}{0.749341in}}%
\pgfpathlineto{\pgfqpoint{0.992225in}{0.759518in}}%
\pgfpathlineto{\pgfqpoint{0.976071in}{0.771175in}}%
\pgfpathlineto{\pgfqpoint{0.967770in}{0.770670in}}%
\pgfpathlineto{\pgfqpoint{0.959431in}{0.770949in}}%
\pgfpathclose%
\pgfusepath{fill}%
\end{pgfscope}%
\begin{pgfscope}%
\pgfpathrectangle{\pgfqpoint{0.407592in}{0.222904in}}{\pgfqpoint{1.411163in}{1.411163in}}%
\pgfusepath{clip}%
\pgfsetbuttcap%
\pgfsetroundjoin%
\definecolor{currentfill}{rgb}{0.556863,0.654902,0.847059}%
\pgfsetfillcolor{currentfill}%
\pgfsetfillopacity{0.130000}%
\pgfsetlinewidth{0.000000pt}%
\definecolor{currentstroke}{rgb}{0.000000,0.000000,0.000000}%
\pgfsetstrokecolor{currentstroke}%
\pgfsetdash{}{0pt}%
\pgfpathmoveto{\pgfqpoint{0.731121in}{0.964059in}}%
\pgfpathlineto{\pgfqpoint{0.747285in}{0.946016in}}%
\pgfpathlineto{\pgfqpoint{0.763458in}{0.928253in}}%
\pgfpathlineto{\pgfqpoint{0.772168in}{0.928046in}}%
\pgfpathlineto{\pgfqpoint{0.780837in}{0.928288in}}%
\pgfpathlineto{\pgfqpoint{0.764706in}{0.946418in}}%
\pgfpathlineto{\pgfqpoint{0.748584in}{0.964799in}}%
\pgfpathlineto{\pgfqpoint{0.739873in}{0.964218in}}%
\pgfpathlineto{\pgfqpoint{0.731121in}{0.964059in}}%
\pgfpathclose%
\pgfusepath{fill}%
\end{pgfscope}%
\begin{pgfscope}%
\pgfpathrectangle{\pgfqpoint{0.407592in}{0.222904in}}{\pgfqpoint{1.411163in}{1.411163in}}%
\pgfusepath{clip}%
\pgfsetbuttcap%
\pgfsetroundjoin%
\definecolor{currentfill}{rgb}{0.556863,0.654902,0.847059}%
\pgfsetfillcolor{currentfill}%
\pgfsetfillopacity{0.130000}%
\pgfsetlinewidth{0.000000pt}%
\definecolor{currentstroke}{rgb}{0.000000,0.000000,0.000000}%
\pgfsetstrokecolor{currentstroke}%
\pgfsetdash{}{0pt}%
\pgfpathmoveto{\pgfqpoint{1.186909in}{0.755929in}}%
\pgfpathlineto{\pgfqpoint{1.203348in}{0.761502in}}%
\pgfpathlineto{\pgfqpoint{1.219850in}{0.768138in}}%
\pgfpathlineto{\pgfqpoint{1.227764in}{0.775216in}}%
\pgfpathlineto{\pgfqpoint{1.235661in}{0.783126in}}%
\pgfpathlineto{\pgfqpoint{1.219178in}{0.776751in}}%
\pgfpathlineto{\pgfqpoint{1.202760in}{0.771439in}}%
\pgfpathlineto{\pgfqpoint{1.194844in}{0.763221in}}%
\pgfpathlineto{\pgfqpoint{1.186909in}{0.755929in}}%
\pgfpathclose%
\pgfusepath{fill}%
\end{pgfscope}%
\begin{pgfscope}%
\pgfpathrectangle{\pgfqpoint{0.407592in}{0.222904in}}{\pgfqpoint{1.411163in}{1.411163in}}%
\pgfusepath{clip}%
\pgfsetbuttcap%
\pgfsetroundjoin%
\definecolor{currentfill}{rgb}{0.556863,0.654902,0.847059}%
\pgfsetfillcolor{currentfill}%
\pgfsetfillopacity{0.130000}%
\pgfsetlinewidth{0.000000pt}%
\definecolor{currentstroke}{rgb}{0.000000,0.000000,0.000000}%
\pgfsetstrokecolor{currentstroke}%
\pgfsetdash{}{0pt}%
\pgfpathmoveto{\pgfqpoint{1.413609in}{0.936014in}}%
\pgfpathlineto{\pgfqpoint{1.430489in}{0.944433in}}%
\pgfpathlineto{\pgfqpoint{1.447440in}{0.953287in}}%
\pgfpathlineto{\pgfqpoint{1.455112in}{0.965047in}}%
\pgfpathlineto{\pgfqpoint{1.462772in}{0.977118in}}%
\pgfpathlineto{\pgfqpoint{1.445840in}{0.968756in}}%
\pgfpathlineto{\pgfqpoint{1.428979in}{0.960848in}}%
\pgfpathlineto{\pgfqpoint{1.421300in}{0.948275in}}%
\pgfpathlineto{\pgfqpoint{1.413609in}{0.936014in}}%
\pgfpathclose%
\pgfusepath{fill}%
\end{pgfscope}%
\begin{pgfscope}%
\pgfpathrectangle{\pgfqpoint{0.407592in}{0.222904in}}{\pgfqpoint{1.411163in}{1.411163in}}%
\pgfusepath{clip}%
\pgfsetbuttcap%
\pgfsetroundjoin%
\definecolor{currentfill}{rgb}{0.556863,0.654902,0.847059}%
\pgfsetfillcolor{currentfill}%
\pgfsetfillopacity{0.130000}%
\pgfsetlinewidth{0.000000pt}%
\definecolor{currentstroke}{rgb}{0.000000,0.000000,0.000000}%
\pgfsetstrokecolor{currentstroke}%
\pgfsetdash{}{0pt}%
\pgfpathmoveto{\pgfqpoint{1.040795in}{0.734121in}}%
\pgfpathlineto{\pgfqpoint{1.057034in}{0.729326in}}%
\pgfpathlineto{\pgfqpoint{1.073305in}{0.726467in}}%
\pgfpathlineto{\pgfqpoint{1.081450in}{0.727379in}}%
\pgfpathlineto{\pgfqpoint{1.089563in}{0.729307in}}%
\pgfpathlineto{\pgfqpoint{1.073329in}{0.732072in}}%
\pgfpathlineto{\pgfqpoint{1.057128in}{0.736877in}}%
\pgfpathlineto{\pgfqpoint{1.048978in}{0.734990in}}%
\pgfpathlineto{\pgfqpoint{1.040795in}{0.734121in}}%
\pgfpathclose%
\pgfusepath{fill}%
\end{pgfscope}%
\begin{pgfscope}%
\pgfpathrectangle{\pgfqpoint{0.407592in}{0.222904in}}{\pgfqpoint{1.411163in}{1.411163in}}%
\pgfusepath{clip}%
\pgfsetbuttcap%
\pgfsetroundjoin%
\definecolor{currentfill}{rgb}{0.556863,0.654902,0.847059}%
\pgfsetfillcolor{currentfill}%
\pgfsetfillopacity{0.130000}%
\pgfsetlinewidth{0.000000pt}%
\definecolor{currentstroke}{rgb}{0.000000,0.000000,0.000000}%
\pgfsetstrokecolor{currentstroke}%
\pgfsetdash{}{0pt}%
\pgfpathmoveto{\pgfqpoint{0.584145in}{1.113863in}}%
\pgfpathlineto{\pgfqpoint{0.600284in}{1.094706in}}%
\pgfpathlineto{\pgfqpoint{0.616431in}{1.075670in}}%
\pgfpathlineto{\pgfqpoint{0.625393in}{1.075951in}}%
\pgfpathlineto{\pgfqpoint{0.634313in}{1.076567in}}%
\pgfpathlineto{\pgfqpoint{0.618208in}{1.095828in}}%
\pgfpathlineto{\pgfqpoint{0.602112in}{1.115194in}}%
\pgfpathlineto{\pgfqpoint{0.593149in}{1.114369in}}%
\pgfpathlineto{\pgfqpoint{0.584145in}{1.113863in}}%
\pgfpathclose%
\pgfusepath{fill}%
\end{pgfscope}%
\begin{pgfscope}%
\pgfpathrectangle{\pgfqpoint{0.407592in}{0.222904in}}{\pgfqpoint{1.411163in}{1.411163in}}%
\pgfusepath{clip}%
\pgfsetbuttcap%
\pgfsetroundjoin%
\definecolor{currentfill}{rgb}{0.556863,0.654902,0.847059}%
\pgfsetfillcolor{currentfill}%
\pgfsetfillopacity{0.130000}%
\pgfsetlinewidth{0.000000pt}%
\definecolor{currentstroke}{rgb}{0.000000,0.000000,0.000000}%
\pgfsetstrokecolor{currentstroke}%
\pgfsetdash{}{0pt}%
\pgfpathmoveto{\pgfqpoint{0.877802in}{0.827486in}}%
\pgfpathlineto{\pgfqpoint{0.893995in}{0.812659in}}%
\pgfpathlineto{\pgfqpoint{0.910197in}{0.798651in}}%
\pgfpathlineto{\pgfqpoint{0.918657in}{0.797533in}}%
\pgfpathlineto{\pgfqpoint{0.927077in}{0.797044in}}%
\pgfpathlineto{\pgfqpoint{0.910917in}{0.811580in}}%
\pgfpathlineto{\pgfqpoint{0.894767in}{0.826919in}}%
\pgfpathlineto{\pgfqpoint{0.886304in}{0.826904in}}%
\pgfpathlineto{\pgfqpoint{0.877802in}{0.827486in}}%
\pgfpathclose%
\pgfusepath{fill}%
\end{pgfscope}%
\begin{pgfscope}%
\pgfpathrectangle{\pgfqpoint{0.407592in}{0.222904in}}{\pgfqpoint{1.411163in}{1.411163in}}%
\pgfusepath{clip}%
\pgfsetbuttcap%
\pgfsetroundjoin%
\definecolor{currentfill}{rgb}{0.556863,0.654902,0.847059}%
\pgfsetfillcolor{currentfill}%
\pgfsetfillopacity{0.130000}%
\pgfsetlinewidth{0.000000pt}%
\definecolor{currentstroke}{rgb}{0.000000,0.000000,0.000000}%
\pgfsetstrokecolor{currentstroke}%
\pgfsetdash{}{0pt}%
\pgfpathmoveto{\pgfqpoint{1.591791in}{1.112895in}}%
\pgfpathlineto{\pgfqpoint{1.609010in}{1.121387in}}%
\pgfpathlineto{\pgfqpoint{1.626303in}{1.130208in}}%
\pgfpathlineto{\pgfqpoint{1.633802in}{1.143067in}}%
\pgfpathlineto{\pgfqpoint{1.641291in}{1.156100in}}%
\pgfpathlineto{\pgfqpoint{1.624018in}{1.147661in}}%
\pgfpathlineto{\pgfqpoint{1.606818in}{1.139560in}}%
\pgfpathlineto{\pgfqpoint{1.599310in}{1.126143in}}%
\pgfpathlineto{\pgfqpoint{1.591791in}{1.112895in}}%
\pgfpathclose%
\pgfusepath{fill}%
\end{pgfscope}%
\begin{pgfscope}%
\pgfpathrectangle{\pgfqpoint{0.407592in}{0.222904in}}{\pgfqpoint{1.411163in}{1.411163in}}%
\pgfusepath{clip}%
\pgfsetbuttcap%
\pgfsetroundjoin%
\definecolor{currentfill}{rgb}{0.556863,0.654902,0.847059}%
\pgfsetfillcolor{currentfill}%
\pgfsetfillopacity{0.130000}%
\pgfsetlinewidth{0.000000pt}%
\definecolor{currentstroke}{rgb}{0.000000,0.000000,0.000000}%
\pgfsetstrokecolor{currentstroke}%
\pgfsetdash{}{0pt}%
\pgfpathmoveto{\pgfqpoint{0.763458in}{0.928253in}}%
\pgfpathlineto{\pgfqpoint{0.779639in}{0.910810in}}%
\pgfpathlineto{\pgfqpoint{0.795829in}{0.893733in}}%
\pgfpathlineto{\pgfqpoint{0.804497in}{0.893101in}}%
\pgfpathlineto{\pgfqpoint{0.813124in}{0.892943in}}%
\pgfpathlineto{\pgfqpoint{0.796976in}{0.910448in}}%
\pgfpathlineto{\pgfqpoint{0.780837in}{0.928288in}}%
\pgfpathlineto{\pgfqpoint{0.772168in}{0.928046in}}%
\pgfpathlineto{\pgfqpoint{0.763458in}{0.928253in}}%
\pgfpathclose%
\pgfusepath{fill}%
\end{pgfscope}%
\begin{pgfscope}%
\pgfpathrectangle{\pgfqpoint{0.407592in}{0.222904in}}{\pgfqpoint{1.411163in}{1.411163in}}%
\pgfusepath{clip}%
\pgfsetbuttcap%
\pgfsetroundjoin%
\definecolor{currentfill}{rgb}{0.556863,0.654902,0.847059}%
\pgfsetfillcolor{currentfill}%
\pgfsetfillopacity{0.130000}%
\pgfsetlinewidth{0.000000pt}%
\definecolor{currentstroke}{rgb}{0.000000,0.000000,0.000000}%
\pgfsetstrokecolor{currentstroke}%
\pgfsetdash{}{0pt}%
\pgfpathmoveto{\pgfqpoint{1.235661in}{0.783126in}}%
\pgfpathlineto{\pgfqpoint{1.252210in}{0.790404in}}%
\pgfpathlineto{\pgfqpoint{1.268827in}{0.798448in}}%
\pgfpathlineto{\pgfqpoint{1.276688in}{0.806783in}}%
\pgfpathlineto{\pgfqpoint{1.284534in}{0.815801in}}%
\pgfpathlineto{\pgfqpoint{1.267935in}{0.808151in}}%
\pgfpathlineto{\pgfqpoint{1.251404in}{0.801285in}}%
\pgfpathlineto{\pgfqpoint{1.243541in}{0.791830in}}%
\pgfpathlineto{\pgfqpoint{1.235661in}{0.783126in}}%
\pgfpathclose%
\pgfusepath{fill}%
\end{pgfscope}%
\begin{pgfscope}%
\pgfpathrectangle{\pgfqpoint{0.407592in}{0.222904in}}{\pgfqpoint{1.411163in}{1.411163in}}%
\pgfusepath{clip}%
\pgfsetbuttcap%
\pgfsetroundjoin%
\definecolor{currentfill}{rgb}{0.556863,0.654902,0.847059}%
\pgfsetfillcolor{currentfill}%
\pgfsetfillopacity{0.130000}%
\pgfsetlinewidth{0.000000pt}%
\definecolor{currentstroke}{rgb}{0.000000,0.000000,0.000000}%
\pgfsetstrokecolor{currentstroke}%
\pgfsetdash{}{0pt}%
\pgfpathmoveto{\pgfqpoint{1.349107in}{0.873144in}}%
\pgfpathlineto{\pgfqpoint{1.365879in}{0.881626in}}%
\pgfpathlineto{\pgfqpoint{1.382724in}{0.890597in}}%
\pgfpathlineto{\pgfqpoint{1.390464in}{0.901355in}}%
\pgfpathlineto{\pgfqpoint{1.398191in}{0.912528in}}%
\pgfpathlineto{\pgfqpoint{1.381366in}{0.904064in}}%
\pgfpathlineto{\pgfqpoint{1.364611in}{0.896115in}}%
\pgfpathlineto{\pgfqpoint{1.356865in}{0.884415in}}%
\pgfpathlineto{\pgfqpoint{1.349107in}{0.873144in}}%
\pgfpathclose%
\pgfusepath{fill}%
\end{pgfscope}%
\begin{pgfscope}%
\pgfpathrectangle{\pgfqpoint{0.407592in}{0.222904in}}{\pgfqpoint{1.411163in}{1.411163in}}%
\pgfusepath{clip}%
\pgfsetbuttcap%
\pgfsetroundjoin%
\definecolor{currentfill}{rgb}{0.556863,0.654902,0.847059}%
\pgfsetfillcolor{currentfill}%
\pgfsetfillopacity{0.130000}%
\pgfsetlinewidth{0.000000pt}%
\definecolor{currentstroke}{rgb}{0.000000,0.000000,0.000000}%
\pgfsetstrokecolor{currentstroke}%
\pgfsetdash{}{0pt}%
\pgfpathmoveto{\pgfqpoint{0.616431in}{1.075670in}}%
\pgfpathlineto{\pgfqpoint{0.632587in}{1.056770in}}%
\pgfpathlineto{\pgfqpoint{0.648751in}{1.038021in}}%
\pgfpathlineto{\pgfqpoint{0.657670in}{1.038039in}}%
\pgfpathlineto{\pgfqpoint{0.666548in}{1.038411in}}%
\pgfpathlineto{\pgfqpoint{0.650426in}{1.057423in}}%
\pgfpathlineto{\pgfqpoint{0.634313in}{1.076567in}}%
\pgfpathlineto{\pgfqpoint{0.625393in}{1.075951in}}%
\pgfpathlineto{\pgfqpoint{0.616431in}{1.075670in}}%
\pgfpathclose%
\pgfusepath{fill}%
\end{pgfscope}%
\begin{pgfscope}%
\pgfpathrectangle{\pgfqpoint{0.407592in}{0.222904in}}{\pgfqpoint{1.411163in}{1.411163in}}%
\pgfusepath{clip}%
\pgfsetbuttcap%
\pgfsetroundjoin%
\definecolor{currentfill}{rgb}{0.556863,0.654902,0.847059}%
\pgfsetfillcolor{currentfill}%
\pgfsetfillopacity{0.130000}%
\pgfsetlinewidth{0.000000pt}%
\definecolor{currentstroke}{rgb}{0.000000,0.000000,0.000000}%
\pgfsetstrokecolor{currentstroke}%
\pgfsetdash{}{0pt}%
\pgfpathmoveto{\pgfqpoint{1.527304in}{1.044208in}}%
\pgfpathlineto{\pgfqpoint{1.544416in}{1.052830in}}%
\pgfpathlineto{\pgfqpoint{1.561602in}{1.061805in}}%
\pgfpathlineto{\pgfqpoint{1.569166in}{1.074267in}}%
\pgfpathlineto{\pgfqpoint{1.576718in}{1.086945in}}%
\pgfpathlineto{\pgfqpoint{1.559553in}{1.078391in}}%
\pgfpathlineto{\pgfqpoint{1.542460in}{1.070200in}}%
\pgfpathlineto{\pgfqpoint{1.534888in}{1.057099in}}%
\pgfpathlineto{\pgfqpoint{1.527304in}{1.044208in}}%
\pgfpathclose%
\pgfusepath{fill}%
\end{pgfscope}%
\begin{pgfscope}%
\pgfpathrectangle{\pgfqpoint{0.407592in}{0.222904in}}{\pgfqpoint{1.411163in}{1.411163in}}%
\pgfusepath{clip}%
\pgfsetbuttcap%
\pgfsetroundjoin%
\definecolor{currentfill}{rgb}{0.556863,0.654902,0.847059}%
\pgfsetfillcolor{currentfill}%
\pgfsetfillopacity{0.130000}%
\pgfsetlinewidth{0.000000pt}%
\definecolor{currentstroke}{rgb}{0.000000,0.000000,0.000000}%
\pgfsetstrokecolor{currentstroke}%
\pgfsetdash{}{0pt}%
\pgfpathmoveto{\pgfqpoint{0.991839in}{0.749886in}}%
\pgfpathlineto{\pgfqpoint{1.008070in}{0.741640in}}%
\pgfpathlineto{\pgfqpoint{1.024323in}{0.735092in}}%
\pgfpathlineto{\pgfqpoint{1.032577in}{0.734186in}}%
\pgfpathlineto{\pgfqpoint{1.040795in}{0.734121in}}%
\pgfpathlineto{\pgfqpoint{1.024584in}{0.740825in}}%
\pgfpathlineto{\pgfqpoint{1.008395in}{0.749341in}}%
\pgfpathlineto{\pgfqpoint{1.000136in}{0.749202in}}%
\pgfpathlineto{\pgfqpoint{0.991839in}{0.749886in}}%
\pgfpathclose%
\pgfusepath{fill}%
\end{pgfscope}%
\begin{pgfscope}%
\pgfpathrectangle{\pgfqpoint{0.407592in}{0.222904in}}{\pgfqpoint{1.411163in}{1.411163in}}%
\pgfusepath{clip}%
\pgfsetbuttcap%
\pgfsetroundjoin%
\definecolor{currentfill}{rgb}{0.556863,0.654902,0.847059}%
\pgfsetfillcolor{currentfill}%
\pgfsetfillopacity{0.130000}%
\pgfsetlinewidth{0.000000pt}%
\definecolor{currentstroke}{rgb}{0.000000,0.000000,0.000000}%
\pgfsetstrokecolor{currentstroke}%
\pgfsetdash{}{0pt}%
\pgfpathmoveto{\pgfqpoint{1.462772in}{0.977118in}}%
\pgfpathlineto{\pgfqpoint{1.479777in}{0.985894in}}%
\pgfpathlineto{\pgfqpoint{1.496855in}{0.995049in}}%
\pgfpathlineto{\pgfqpoint{1.504485in}{1.006946in}}%
\pgfpathlineto{\pgfqpoint{1.512102in}{1.019116in}}%
\pgfpathlineto{\pgfqpoint{1.495044in}{1.010418in}}%
\pgfpathlineto{\pgfqpoint{1.478058in}{1.002116in}}%
\pgfpathlineto{\pgfqpoint{1.470421in}{0.989481in}}%
\pgfpathlineto{\pgfqpoint{1.462772in}{0.977118in}}%
\pgfpathclose%
\pgfusepath{fill}%
\end{pgfscope}%
\begin{pgfscope}%
\pgfpathrectangle{\pgfqpoint{0.407592in}{0.222904in}}{\pgfqpoint{1.411163in}{1.411163in}}%
\pgfusepath{clip}%
\pgfsetbuttcap%
\pgfsetroundjoin%
\definecolor{currentfill}{rgb}{0.556863,0.654902,0.847059}%
\pgfsetfillcolor{currentfill}%
\pgfsetfillopacity{0.130000}%
\pgfsetlinewidth{0.000000pt}%
\definecolor{currentstroke}{rgb}{0.000000,0.000000,0.000000}%
\pgfsetstrokecolor{currentstroke}%
\pgfsetdash{}{0pt}%
\pgfpathmoveto{\pgfqpoint{0.910197in}{0.798651in}}%
\pgfpathlineto{\pgfqpoint{0.926411in}{0.785584in}}%
\pgfpathlineto{\pgfqpoint{0.942637in}{0.773594in}}%
\pgfpathlineto{\pgfqpoint{0.951053in}{0.771946in}}%
\pgfpathlineto{\pgfqpoint{0.959431in}{0.770949in}}%
\pgfpathlineto{\pgfqpoint{0.943248in}{0.783450in}}%
\pgfpathlineto{\pgfqpoint{0.927077in}{0.797044in}}%
\pgfpathlineto{\pgfqpoint{0.918657in}{0.797533in}}%
\pgfpathlineto{\pgfqpoint{0.910197in}{0.798651in}}%
\pgfpathclose%
\pgfusepath{fill}%
\end{pgfscope}%
\begin{pgfscope}%
\pgfpathrectangle{\pgfqpoint{0.407592in}{0.222904in}}{\pgfqpoint{1.411163in}{1.411163in}}%
\pgfusepath{clip}%
\pgfsetbuttcap%
\pgfsetroundjoin%
\definecolor{currentfill}{rgb}{0.556863,0.654902,0.847059}%
\pgfsetfillcolor{currentfill}%
\pgfsetfillopacity{0.130000}%
\pgfsetlinewidth{0.000000pt}%
\definecolor{currentstroke}{rgb}{0.000000,0.000000,0.000000}%
\pgfsetstrokecolor{currentstroke}%
\pgfsetdash{}{0pt}%
\pgfpathmoveto{\pgfqpoint{0.648751in}{1.038021in}}%
\pgfpathlineto{\pgfqpoint{0.664924in}{1.019442in}}%
\pgfpathlineto{\pgfqpoint{0.681105in}{1.001054in}}%
\pgfpathlineto{\pgfqpoint{0.689982in}{1.000767in}}%
\pgfpathlineto{\pgfqpoint{0.698817in}{1.000853in}}%
\pgfpathlineto{\pgfqpoint{0.682678in}{1.019548in}}%
\pgfpathlineto{\pgfqpoint{0.666548in}{1.038411in}}%
\pgfpathlineto{\pgfqpoint{0.657670in}{1.038039in}}%
\pgfpathlineto{\pgfqpoint{0.648751in}{1.038021in}}%
\pgfpathclose%
\pgfusepath{fill}%
\end{pgfscope}%
\begin{pgfscope}%
\pgfpathrectangle{\pgfqpoint{0.407592in}{0.222904in}}{\pgfqpoint{1.411163in}{1.411163in}}%
\pgfusepath{clip}%
\pgfsetbuttcap%
\pgfsetroundjoin%
\definecolor{currentfill}{rgb}{0.556863,0.654902,0.847059}%
\pgfsetfillcolor{currentfill}%
\pgfsetfillopacity{0.130000}%
\pgfsetlinewidth{0.000000pt}%
\definecolor{currentstroke}{rgb}{0.000000,0.000000,0.000000}%
\pgfsetstrokecolor{currentstroke}%
\pgfsetdash{}{0pt}%
\pgfpathmoveto{\pgfqpoint{0.795829in}{0.893733in}}%
\pgfpathlineto{\pgfqpoint{0.812028in}{0.877075in}}%
\pgfpathlineto{\pgfqpoint{0.828235in}{0.860898in}}%
\pgfpathlineto{\pgfqpoint{0.836860in}{0.859781in}}%
\pgfpathlineto{\pgfqpoint{0.845445in}{0.859161in}}%
\pgfpathlineto{\pgfqpoint{0.829280in}{0.875827in}}%
\pgfpathlineto{\pgfqpoint{0.813124in}{0.892943in}}%
\pgfpathlineto{\pgfqpoint{0.804497in}{0.893101in}}%
\pgfpathlineto{\pgfqpoint{0.795829in}{0.893733in}}%
\pgfpathclose%
\pgfusepath{fill}%
\end{pgfscope}%
\begin{pgfscope}%
\pgfpathrectangle{\pgfqpoint{0.407592in}{0.222904in}}{\pgfqpoint{1.411163in}{1.411163in}}%
\pgfusepath{clip}%
\pgfsetbuttcap%
\pgfsetroundjoin%
\definecolor{currentfill}{rgb}{0.556863,0.654902,0.847059}%
\pgfsetfillcolor{currentfill}%
\pgfsetfillopacity{0.130000}%
\pgfsetlinewidth{0.000000pt}%
\definecolor{currentstroke}{rgb}{0.000000,0.000000,0.000000}%
\pgfsetstrokecolor{currentstroke}%
\pgfsetdash{}{0pt}%
\pgfpathmoveto{\pgfqpoint{1.284534in}{0.815801in}}%
\pgfpathlineto{\pgfqpoint{1.301202in}{0.824123in}}%
\pgfpathlineto{\pgfqpoint{1.317940in}{0.833022in}}%
\pgfpathlineto{\pgfqpoint{1.325752in}{0.842242in}}%
\pgfpathlineto{\pgfqpoint{1.333551in}{0.852024in}}%
\pgfpathlineto{\pgfqpoint{1.316831in}{0.843577in}}%
\pgfpathlineto{\pgfqpoint{1.300181in}{0.835734in}}%
\pgfpathlineto{\pgfqpoint{1.292365in}{0.825465in}}%
\pgfpathlineto{\pgfqpoint{1.284534in}{0.815801in}}%
\pgfpathclose%
\pgfusepath{fill}%
\end{pgfscope}%
\begin{pgfscope}%
\pgfpathrectangle{\pgfqpoint{0.407592in}{0.222904in}}{\pgfqpoint{1.411163in}{1.411163in}}%
\pgfusepath{clip}%
\pgfsetbuttcap%
\pgfsetroundjoin%
\definecolor{currentfill}{rgb}{0.556863,0.654902,0.847059}%
\pgfsetfillcolor{currentfill}%
\pgfsetfillopacity{0.130000}%
\pgfsetlinewidth{0.000000pt}%
\definecolor{currentstroke}{rgb}{0.000000,0.000000,0.000000}%
\pgfsetstrokecolor{currentstroke}%
\pgfsetdash{}{0pt}%
\pgfpathmoveto{\pgfqpoint{1.398191in}{0.912528in}}%
\pgfpathlineto{\pgfqpoint{1.415089in}{0.921451in}}%
\pgfpathlineto{\pgfqpoint{1.432059in}{0.930787in}}%
\pgfpathlineto{\pgfqpoint{1.439755in}{0.941860in}}%
\pgfpathlineto{\pgfqpoint{1.447440in}{0.953287in}}%
\pgfpathlineto{\pgfqpoint{1.430489in}{0.944433in}}%
\pgfpathlineto{\pgfqpoint{1.413609in}{0.936014in}}%
\pgfpathlineto{\pgfqpoint{1.405906in}{0.924090in}}%
\pgfpathlineto{\pgfqpoint{1.398191in}{0.912528in}}%
\pgfpathclose%
\pgfusepath{fill}%
\end{pgfscope}%
\begin{pgfscope}%
\pgfpathrectangle{\pgfqpoint{0.407592in}{0.222904in}}{\pgfqpoint{1.411163in}{1.411163in}}%
\pgfusepath{clip}%
\pgfsetbuttcap%
\pgfsetroundjoin%
\definecolor{currentfill}{rgb}{0.556863,0.654902,0.847059}%
\pgfsetfillcolor{currentfill}%
\pgfsetfillopacity{0.130000}%
\pgfsetlinewidth{0.000000pt}%
\definecolor{currentstroke}{rgb}{0.000000,0.000000,0.000000}%
\pgfsetstrokecolor{currentstroke}%
\pgfsetdash{}{0pt}%
\pgfpathmoveto{\pgfqpoint{1.641291in}{1.156100in}}%
\pgfpathlineto{\pgfqpoint{1.658638in}{1.164855in}}%
\pgfpathlineto{\pgfqpoint{1.676059in}{1.173911in}}%
\pgfpathlineto{\pgfqpoint{1.683516in}{1.186745in}}%
\pgfpathlineto{\pgfqpoint{1.666105in}{1.177868in}}%
\pgfpathlineto{\pgfqpoint{1.648768in}{1.169296in}}%
\pgfusepath{fill}%
\end{pgfscope}%
\begin{pgfscope}%
\pgfpathrectangle{\pgfqpoint{0.407592in}{0.222904in}}{\pgfqpoint{1.411163in}{1.411163in}}%
\pgfusepath{clip}%
\pgfsetbuttcap%
\pgfsetroundjoin%
\definecolor{currentfill}{rgb}{0.556863,0.654902,0.847059}%
\pgfsetfillcolor{currentfill}%
\pgfsetfillopacity{0.130000}%
\pgfsetlinewidth{0.000000pt}%
\definecolor{currentstroke}{rgb}{0.000000,0.000000,0.000000}%
\pgfsetstrokecolor{currentstroke}%
\pgfsetdash{}{0pt}%
\pgfpathmoveto{\pgfqpoint{1.576718in}{1.086945in}}%
\pgfpathlineto{\pgfqpoint{1.593958in}{1.095838in}}%
\pgfpathlineto{\pgfqpoint{1.611272in}{1.105049in}}%
\pgfpathlineto{\pgfqpoint{1.618793in}{1.117531in}}%
\pgfpathlineto{\pgfqpoint{1.626303in}{1.130208in}}%
\pgfpathlineto{\pgfqpoint{1.609010in}{1.121387in}}%
\pgfpathlineto{\pgfqpoint{1.591791in}{1.112895in}}%
\pgfpathlineto{\pgfqpoint{1.584260in}{1.099825in}}%
\pgfpathlineto{\pgfqpoint{1.576718in}{1.086945in}}%
\pgfpathclose%
\pgfusepath{fill}%
\end{pgfscope}%
\begin{pgfscope}%
\pgfpathrectangle{\pgfqpoint{0.407592in}{0.222904in}}{\pgfqpoint{1.411163in}{1.411163in}}%
\pgfusepath{clip}%
\pgfsetbuttcap%
\pgfsetroundjoin%
\definecolor{currentfill}{rgb}{0.556863,0.654902,0.847059}%
\pgfsetfillcolor{currentfill}%
\pgfsetfillopacity{0.130000}%
\pgfsetlinewidth{0.000000pt}%
\definecolor{currentstroke}{rgb}{0.000000,0.000000,0.000000}%
\pgfsetstrokecolor{currentstroke}%
\pgfsetdash{}{0pt}%
\pgfpathmoveto{\pgfqpoint{1.122158in}{0.729595in}}%
\pgfpathlineto{\pgfqpoint{1.138528in}{0.732333in}}%
\pgfpathlineto{\pgfqpoint{1.154954in}{0.736535in}}%
\pgfpathlineto{\pgfqpoint{1.162978in}{0.739910in}}%
\pgfpathlineto{\pgfqpoint{1.170978in}{0.744261in}}%
\pgfpathlineto{\pgfqpoint{1.154578in}{0.739893in}}%
\pgfpathlineto{\pgfqpoint{1.138235in}{0.737000in}}%
\pgfpathlineto{\pgfqpoint{1.130210in}{0.732760in}}%
\pgfpathlineto{\pgfqpoint{1.122158in}{0.729595in}}%
\pgfpathclose%
\pgfusepath{fill}%
\end{pgfscope}%
\begin{pgfscope}%
\pgfpathrectangle{\pgfqpoint{0.407592in}{0.222904in}}{\pgfqpoint{1.411163in}{1.411163in}}%
\pgfusepath{clip}%
\pgfsetbuttcap%
\pgfsetroundjoin%
\definecolor{currentfill}{rgb}{0.556863,0.654902,0.847059}%
\pgfsetfillcolor{currentfill}%
\pgfsetfillopacity{0.130000}%
\pgfsetlinewidth{0.000000pt}%
\definecolor{currentstroke}{rgb}{0.000000,0.000000,0.000000}%
\pgfsetstrokecolor{currentstroke}%
\pgfsetdash{}{0pt}%
\pgfpathmoveto{\pgfqpoint{0.681105in}{1.001054in}}%
\pgfpathlineto{\pgfqpoint{0.697294in}{0.982880in}}%
\pgfpathlineto{\pgfqpoint{0.713492in}{0.964947in}}%
\pgfpathlineto{\pgfqpoint{0.722327in}{0.964308in}}%
\pgfpathlineto{\pgfqpoint{0.731121in}{0.964059in}}%
\pgfpathlineto{\pgfqpoint{0.714965in}{0.982348in}}%
\pgfpathlineto{\pgfqpoint{0.698817in}{1.000853in}}%
\pgfpathlineto{\pgfqpoint{0.689982in}{1.000767in}}%
\pgfpathlineto{\pgfqpoint{0.681105in}{1.001054in}}%
\pgfpathclose%
\pgfusepath{fill}%
\end{pgfscope}%
\begin{pgfscope}%
\pgfpathrectangle{\pgfqpoint{0.407592in}{0.222904in}}{\pgfqpoint{1.411163in}{1.411163in}}%
\pgfusepath{clip}%
\pgfsetbuttcap%
\pgfsetroundjoin%
\definecolor{currentfill}{rgb}{0.556863,0.654902,0.847059}%
\pgfsetfillcolor{currentfill}%
\pgfsetfillopacity{0.130000}%
\pgfsetlinewidth{0.000000pt}%
\definecolor{currentstroke}{rgb}{0.000000,0.000000,0.000000}%
\pgfsetstrokecolor{currentstroke}%
\pgfsetdash{}{0pt}%
\pgfpathmoveto{\pgfqpoint{1.073305in}{0.726467in}}%
\pgfpathlineto{\pgfqpoint{1.089614in}{0.725496in}}%
\pgfpathlineto{\pgfqpoint{1.105968in}{0.726299in}}%
\pgfpathlineto{\pgfqpoint{1.114077in}{0.727461in}}%
\pgfpathlineto{\pgfqpoint{1.122158in}{0.729595in}}%
\pgfpathlineto{\pgfqpoint{1.105838in}{0.728526in}}%
\pgfpathlineto{\pgfqpoint{1.089563in}{0.729307in}}%
\pgfpathlineto{\pgfqpoint{1.081450in}{0.727379in}}%
\pgfpathlineto{\pgfqpoint{1.073305in}{0.726467in}}%
\pgfpathclose%
\pgfusepath{fill}%
\end{pgfscope}%
\begin{pgfscope}%
\pgfpathrectangle{\pgfqpoint{0.407592in}{0.222904in}}{\pgfqpoint{1.411163in}{1.411163in}}%
\pgfusepath{clip}%
\pgfsetbuttcap%
\pgfsetroundjoin%
\definecolor{currentfill}{rgb}{0.556863,0.654902,0.847059}%
\pgfsetfillcolor{currentfill}%
\pgfsetfillopacity{0.130000}%
\pgfsetlinewidth{0.000000pt}%
\definecolor{currentstroke}{rgb}{0.000000,0.000000,0.000000}%
\pgfsetstrokecolor{currentstroke}%
\pgfsetdash{}{0pt}%
\pgfpathmoveto{\pgfqpoint{0.828235in}{0.860898in}}%
\pgfpathlineto{\pgfqpoint{0.844451in}{0.845273in}}%
\pgfpathlineto{\pgfqpoint{0.860677in}{0.830281in}}%
\pgfpathlineto{\pgfqpoint{0.869260in}{0.828625in}}%
\pgfpathlineto{\pgfqpoint{0.877802in}{0.827486in}}%
\pgfpathlineto{\pgfqpoint{0.861619in}{0.843019in}}%
\pgfpathlineto{\pgfqpoint{0.845445in}{0.859161in}}%
\pgfpathlineto{\pgfqpoint{0.836860in}{0.859781in}}%
\pgfpathlineto{\pgfqpoint{0.828235in}{0.860898in}}%
\pgfpathclose%
\pgfusepath{fill}%
\end{pgfscope}%
\begin{pgfscope}%
\pgfpathrectangle{\pgfqpoint{0.407592in}{0.222904in}}{\pgfqpoint{1.411163in}{1.411163in}}%
\pgfusepath{clip}%
\pgfsetbuttcap%
\pgfsetroundjoin%
\definecolor{currentfill}{rgb}{0.556863,0.654902,0.847059}%
\pgfsetfillcolor{currentfill}%
\pgfsetfillopacity{0.130000}%
\pgfsetlinewidth{0.000000pt}%
\definecolor{currentstroke}{rgb}{0.000000,0.000000,0.000000}%
\pgfsetstrokecolor{currentstroke}%
\pgfsetdash{}{0pt}%
\pgfpathmoveto{\pgfqpoint{1.170978in}{0.744261in}}%
\pgfpathlineto{\pgfqpoint{1.187438in}{0.749888in}}%
\pgfpathlineto{\pgfqpoint{1.203961in}{0.756580in}}%
\pgfpathlineto{\pgfqpoint{1.211916in}{0.761921in}}%
\pgfpathlineto{\pgfqpoint{1.219850in}{0.768138in}}%
\pgfpathlineto{\pgfqpoint{1.203348in}{0.761502in}}%
\pgfpathlineto{\pgfqpoint{1.186909in}{0.755929in}}%
\pgfpathlineto{\pgfqpoint{1.178955in}{0.749601in}}%
\pgfpathlineto{\pgfqpoint{1.170978in}{0.744261in}}%
\pgfpathclose%
\pgfusepath{fill}%
\end{pgfscope}%
\begin{pgfscope}%
\pgfpathrectangle{\pgfqpoint{0.407592in}{0.222904in}}{\pgfqpoint{1.411163in}{1.411163in}}%
\pgfusepath{clip}%
\pgfsetbuttcap%
\pgfsetroundjoin%
\definecolor{currentfill}{rgb}{0.556863,0.654902,0.847059}%
\pgfsetfillcolor{currentfill}%
\pgfsetfillopacity{0.130000}%
\pgfsetlinewidth{0.000000pt}%
\definecolor{currentstroke}{rgb}{0.000000,0.000000,0.000000}%
\pgfsetstrokecolor{currentstroke}%
\pgfsetdash{}{0pt}%
\pgfpathmoveto{\pgfqpoint{0.533675in}{1.151976in}}%
\pgfpathlineto{\pgfqpoint{0.549838in}{1.132820in}}%
\pgfpathlineto{\pgfqpoint{0.566011in}{1.113777in}}%
\pgfpathlineto{\pgfqpoint{0.575099in}{1.113668in}}%
\pgfpathlineto{\pgfqpoint{0.584145in}{1.113863in}}%
\pgfpathlineto{\pgfqpoint{0.568015in}{1.133129in}}%
\pgfpathlineto{\pgfqpoint{0.551894in}{1.152493in}}%
\pgfpathlineto{\pgfqpoint{0.542805in}{1.152089in}}%
\pgfpathlineto{\pgfqpoint{0.533675in}{1.151976in}}%
\pgfpathclose%
\pgfusepath{fill}%
\end{pgfscope}%
\begin{pgfscope}%
\pgfpathrectangle{\pgfqpoint{0.407592in}{0.222904in}}{\pgfqpoint{1.411163in}{1.411163in}}%
\pgfusepath{clip}%
\pgfsetbuttcap%
\pgfsetroundjoin%
\definecolor{currentfill}{rgb}{0.556863,0.654902,0.847059}%
\pgfsetfillcolor{currentfill}%
\pgfsetfillopacity{0.130000}%
\pgfsetlinewidth{0.000000pt}%
\definecolor{currentstroke}{rgb}{0.000000,0.000000,0.000000}%
\pgfsetstrokecolor{currentstroke}%
\pgfsetdash{}{0pt}%
\pgfpathmoveto{\pgfqpoint{1.512102in}{1.019116in}}%
\pgfpathlineto{\pgfqpoint{1.529234in}{1.028178in}}%
\pgfpathlineto{\pgfqpoint{1.546440in}{1.037577in}}%
\pgfpathlineto{\pgfqpoint{1.554027in}{1.049570in}}%
\pgfpathlineto{\pgfqpoint{1.561602in}{1.061805in}}%
\pgfpathlineto{\pgfqpoint{1.544416in}{1.052830in}}%
\pgfpathlineto{\pgfqpoint{1.527304in}{1.044208in}}%
\pgfpathlineto{\pgfqpoint{1.519709in}{1.031542in}}%
\pgfpathlineto{\pgfqpoint{1.512102in}{1.019116in}}%
\pgfpathclose%
\pgfusepath{fill}%
\end{pgfscope}%
\begin{pgfscope}%
\pgfpathrectangle{\pgfqpoint{0.407592in}{0.222904in}}{\pgfqpoint{1.411163in}{1.411163in}}%
\pgfusepath{clip}%
\pgfsetbuttcap%
\pgfsetroundjoin%
\definecolor{currentfill}{rgb}{0.556863,0.654902,0.847059}%
\pgfsetfillcolor{currentfill}%
\pgfsetfillopacity{0.130000}%
\pgfsetlinewidth{0.000000pt}%
\definecolor{currentstroke}{rgb}{0.000000,0.000000,0.000000}%
\pgfsetstrokecolor{currentstroke}%
\pgfsetdash{}{0pt}%
\pgfpathmoveto{\pgfqpoint{0.942637in}{0.773594in}}%
\pgfpathlineto{\pgfqpoint{0.958877in}{0.762819in}}%
\pgfpathlineto{\pgfqpoint{0.975132in}{0.753395in}}%
\pgfpathlineto{\pgfqpoint{0.983505in}{0.751310in}}%
\pgfpathlineto{\pgfqpoint{0.991839in}{0.749886in}}%
\pgfpathlineto{\pgfqpoint{0.975627in}{0.759706in}}%
\pgfpathlineto{\pgfqpoint{0.959431in}{0.770949in}}%
\pgfpathlineto{\pgfqpoint{0.951053in}{0.771946in}}%
\pgfpathlineto{\pgfqpoint{0.942637in}{0.773594in}}%
\pgfpathclose%
\pgfusepath{fill}%
\end{pgfscope}%
\begin{pgfscope}%
\pgfpathrectangle{\pgfqpoint{0.407592in}{0.222904in}}{\pgfqpoint{1.411163in}{1.411163in}}%
\pgfusepath{clip}%
\pgfsetbuttcap%
\pgfsetroundjoin%
\definecolor{currentfill}{rgb}{0.556863,0.654902,0.847059}%
\pgfsetfillcolor{currentfill}%
\pgfsetfillopacity{0.130000}%
\pgfsetlinewidth{0.000000pt}%
\definecolor{currentstroke}{rgb}{0.000000,0.000000,0.000000}%
\pgfsetstrokecolor{currentstroke}%
\pgfsetdash{}{0pt}%
\pgfpathmoveto{\pgfqpoint{1.333551in}{0.852024in}}%
\pgfpathlineto{\pgfqpoint{1.350342in}{0.860999in}}%
\pgfpathlineto{\pgfqpoint{1.367204in}{0.870434in}}%
\pgfpathlineto{\pgfqpoint{1.374970in}{0.880281in}}%
\pgfpathlineto{\pgfqpoint{1.382724in}{0.890597in}}%
\pgfpathlineto{\pgfqpoint{1.365879in}{0.881626in}}%
\pgfpathlineto{\pgfqpoint{1.349107in}{0.873144in}}%
\pgfpathlineto{\pgfqpoint{1.341335in}{0.862336in}}%
\pgfpathlineto{\pgfqpoint{1.333551in}{0.852024in}}%
\pgfpathclose%
\pgfusepath{fill}%
\end{pgfscope}%
\begin{pgfscope}%
\pgfpathrectangle{\pgfqpoint{0.407592in}{0.222904in}}{\pgfqpoint{1.411163in}{1.411163in}}%
\pgfusepath{clip}%
\pgfsetbuttcap%
\pgfsetroundjoin%
\definecolor{currentfill}{rgb}{0.556863,0.654902,0.847059}%
\pgfsetfillcolor{currentfill}%
\pgfsetfillopacity{0.130000}%
\pgfsetlinewidth{0.000000pt}%
\definecolor{currentstroke}{rgb}{0.000000,0.000000,0.000000}%
\pgfsetstrokecolor{currentstroke}%
\pgfsetdash{}{0pt}%
\pgfpathmoveto{\pgfqpoint{0.713492in}{0.964947in}}%
\pgfpathlineto{\pgfqpoint{0.729699in}{0.947286in}}%
\pgfpathlineto{\pgfqpoint{0.745914in}{0.929934in}}%
\pgfpathlineto{\pgfqpoint{0.754707in}{0.928889in}}%
\pgfpathlineto{\pgfqpoint{0.763458in}{0.928253in}}%
\pgfpathlineto{\pgfqpoint{0.747285in}{0.946016in}}%
\pgfpathlineto{\pgfqpoint{0.731121in}{0.964059in}}%
\pgfpathlineto{\pgfqpoint{0.722327in}{0.964308in}}%
\pgfpathlineto{\pgfqpoint{0.713492in}{0.964947in}}%
\pgfpathclose%
\pgfusepath{fill}%
\end{pgfscope}%
\begin{pgfscope}%
\pgfpathrectangle{\pgfqpoint{0.407592in}{0.222904in}}{\pgfqpoint{1.411163in}{1.411163in}}%
\pgfusepath{clip}%
\pgfsetbuttcap%
\pgfsetroundjoin%
\definecolor{currentfill}{rgb}{0.556863,0.654902,0.847059}%
\pgfsetfillcolor{currentfill}%
\pgfsetfillopacity{0.130000}%
\pgfsetlinewidth{0.000000pt}%
\definecolor{currentstroke}{rgb}{0.000000,0.000000,0.000000}%
\pgfsetstrokecolor{currentstroke}%
\pgfsetdash{}{0pt}%
\pgfpathmoveto{\pgfqpoint{1.219850in}{0.768138in}}%
\pgfpathlineto{\pgfqpoint{1.236418in}{0.775670in}}%
\pgfpathlineto{\pgfqpoint{1.253054in}{0.783954in}}%
\pgfpathlineto{\pgfqpoint{1.260949in}{0.790829in}}%
\pgfpathlineto{\pgfqpoint{1.268827in}{0.798448in}}%
\pgfpathlineto{\pgfqpoint{1.252210in}{0.790404in}}%
\pgfpathlineto{\pgfqpoint{1.235661in}{0.783126in}}%
\pgfpathlineto{\pgfqpoint{1.227764in}{0.775216in}}%
\pgfpathlineto{\pgfqpoint{1.219850in}{0.768138in}}%
\pgfpathclose%
\pgfusepath{fill}%
\end{pgfscope}%
\begin{pgfscope}%
\pgfpathrectangle{\pgfqpoint{0.407592in}{0.222904in}}{\pgfqpoint{1.411163in}{1.411163in}}%
\pgfusepath{clip}%
\pgfsetbuttcap%
\pgfsetroundjoin%
\definecolor{currentfill}{rgb}{0.556863,0.654902,0.847059}%
\pgfsetfillcolor{currentfill}%
\pgfsetfillopacity{0.130000}%
\pgfsetlinewidth{0.000000pt}%
\definecolor{currentstroke}{rgb}{0.000000,0.000000,0.000000}%
\pgfsetstrokecolor{currentstroke}%
\pgfsetdash{}{0pt}%
\pgfpathmoveto{\pgfqpoint{0.566011in}{1.113777in}}%
\pgfpathlineto{\pgfqpoint{0.582192in}{1.094859in}}%
\pgfpathlineto{\pgfqpoint{0.598381in}{1.076080in}}%
\pgfpathlineto{\pgfqpoint{0.607427in}{1.075716in}}%
\pgfpathlineto{\pgfqpoint{0.616431in}{1.075670in}}%
\pgfpathlineto{\pgfqpoint{0.600284in}{1.094706in}}%
\pgfpathlineto{\pgfqpoint{0.584145in}{1.113863in}}%
\pgfpathlineto{\pgfqpoint{0.575099in}{1.113668in}}%
\pgfpathlineto{\pgfqpoint{0.566011in}{1.113777in}}%
\pgfpathclose%
\pgfusepath{fill}%
\end{pgfscope}%
\begin{pgfscope}%
\pgfpathrectangle{\pgfqpoint{0.407592in}{0.222904in}}{\pgfqpoint{1.411163in}{1.411163in}}%
\pgfusepath{clip}%
\pgfsetbuttcap%
\pgfsetroundjoin%
\definecolor{currentfill}{rgb}{0.556863,0.654902,0.847059}%
\pgfsetfillcolor{currentfill}%
\pgfsetfillopacity{0.130000}%
\pgfsetlinewidth{0.000000pt}%
\definecolor{currentstroke}{rgb}{0.000000,0.000000,0.000000}%
\pgfsetstrokecolor{currentstroke}%
\pgfsetdash{}{0pt}%
\pgfpathmoveto{\pgfqpoint{1.024323in}{0.735092in}}%
\pgfpathlineto{\pgfqpoint{1.040603in}{0.730320in}}%
\pgfpathlineto{\pgfqpoint{1.056914in}{0.727347in}}%
\pgfpathlineto{\pgfqpoint{1.065127in}{0.726487in}}%
\pgfpathlineto{\pgfqpoint{1.073305in}{0.726467in}}%
\pgfpathlineto{\pgfqpoint{1.057034in}{0.729326in}}%
\pgfpathlineto{\pgfqpoint{1.040795in}{0.734121in}}%
\pgfpathlineto{\pgfqpoint{1.032577in}{0.734186in}}%
\pgfpathlineto{\pgfqpoint{1.024323in}{0.735092in}}%
\pgfpathclose%
\pgfusepath{fill}%
\end{pgfscope}%
\begin{pgfscope}%
\pgfpathrectangle{\pgfqpoint{0.407592in}{0.222904in}}{\pgfqpoint{1.411163in}{1.411163in}}%
\pgfusepath{clip}%
\pgfsetbuttcap%
\pgfsetroundjoin%
\definecolor{currentfill}{rgb}{0.556863,0.654902,0.847059}%
\pgfsetfillcolor{currentfill}%
\pgfsetfillopacity{0.130000}%
\pgfsetlinewidth{0.000000pt}%
\definecolor{currentstroke}{rgb}{0.000000,0.000000,0.000000}%
\pgfsetstrokecolor{currentstroke}%
\pgfsetdash{}{0pt}%
\pgfpathmoveto{\pgfqpoint{1.447440in}{0.953287in}}%
\pgfpathlineto{\pgfqpoint{1.464464in}{0.962535in}}%
\pgfpathlineto{\pgfqpoint{1.481561in}{0.972142in}}%
\pgfpathlineto{\pgfqpoint{1.489214in}{0.983442in}}%
\pgfpathlineto{\pgfqpoint{1.496855in}{0.995049in}}%
\pgfpathlineto{\pgfqpoint{1.479777in}{0.985894in}}%
\pgfpathlineto{\pgfqpoint{1.462772in}{0.977118in}}%
\pgfpathlineto{\pgfqpoint{1.455112in}{0.965047in}}%
\pgfpathlineto{\pgfqpoint{1.447440in}{0.953287in}}%
\pgfpathclose%
\pgfusepath{fill}%
\end{pgfscope}%
\begin{pgfscope}%
\pgfpathrectangle{\pgfqpoint{0.407592in}{0.222904in}}{\pgfqpoint{1.411163in}{1.411163in}}%
\pgfusepath{clip}%
\pgfsetbuttcap%
\pgfsetroundjoin%
\definecolor{currentfill}{rgb}{0.556863,0.654902,0.847059}%
\pgfsetfillcolor{currentfill}%
\pgfsetfillopacity{0.130000}%
\pgfsetlinewidth{0.000000pt}%
\definecolor{currentstroke}{rgb}{0.000000,0.000000,0.000000}%
\pgfsetstrokecolor{currentstroke}%
\pgfsetdash{}{0pt}%
\pgfpathmoveto{\pgfqpoint{0.860677in}{0.830281in}}%
\pgfpathlineto{\pgfqpoint{0.876912in}{0.816011in}}%
\pgfpathlineto{\pgfqpoint{0.893158in}{0.802566in}}%
\pgfpathlineto{\pgfqpoint{0.901698in}{0.800345in}}%
\pgfpathlineto{\pgfqpoint{0.910197in}{0.798651in}}%
\pgfpathlineto{\pgfqpoint{0.893995in}{0.812659in}}%
\pgfpathlineto{\pgfqpoint{0.877802in}{0.827486in}}%
\pgfpathlineto{\pgfqpoint{0.869260in}{0.828625in}}%
\pgfpathlineto{\pgfqpoint{0.860677in}{0.830281in}}%
\pgfpathclose%
\pgfusepath{fill}%
\end{pgfscope}%
\begin{pgfscope}%
\pgfpathrectangle{\pgfqpoint{0.407592in}{0.222904in}}{\pgfqpoint{1.411163in}{1.411163in}}%
\pgfusepath{clip}%
\pgfsetbuttcap%
\pgfsetroundjoin%
\definecolor{currentfill}{rgb}{0.556863,0.654902,0.847059}%
\pgfsetfillcolor{currentfill}%
\pgfsetfillopacity{0.130000}%
\pgfsetlinewidth{0.000000pt}%
\definecolor{currentstroke}{rgb}{0.000000,0.000000,0.000000}%
\pgfsetstrokecolor{currentstroke}%
\pgfsetdash{}{0pt}%
\pgfpathmoveto{\pgfqpoint{1.626303in}{1.130208in}}%
\pgfpathlineto{\pgfqpoint{1.643671in}{1.139337in}}%
\pgfpathlineto{\pgfqpoint{1.661114in}{1.148756in}}%
\pgfpathlineto{\pgfqpoint{1.668592in}{1.161245in}}%
\pgfpathlineto{\pgfqpoint{1.676059in}{1.173911in}}%
\pgfpathlineto{\pgfqpoint{1.658638in}{1.164855in}}%
\pgfpathlineto{\pgfqpoint{1.641291in}{1.156100in}}%
\pgfpathlineto{\pgfqpoint{1.633802in}{1.143067in}}%
\pgfpathlineto{\pgfqpoint{1.626303in}{1.130208in}}%
\pgfpathclose%
\pgfusepath{fill}%
\end{pgfscope}%
\begin{pgfscope}%
\pgfpathrectangle{\pgfqpoint{0.407592in}{0.222904in}}{\pgfqpoint{1.411163in}{1.411163in}}%
\pgfusepath{clip}%
\pgfsetbuttcap%
\pgfsetroundjoin%
\definecolor{currentfill}{rgb}{0.556863,0.654902,0.847059}%
\pgfsetfillcolor{currentfill}%
\pgfsetfillopacity{0.130000}%
\pgfsetlinewidth{0.000000pt}%
\definecolor{currentstroke}{rgb}{0.000000,0.000000,0.000000}%
\pgfsetstrokecolor{currentstroke}%
\pgfsetdash{}{0pt}%
\pgfpathmoveto{\pgfqpoint{0.745914in}{0.929934in}}%
\pgfpathlineto{\pgfqpoint{0.762138in}{0.912929in}}%
\pgfpathlineto{\pgfqpoint{0.778370in}{0.896319in}}%
\pgfpathlineto{\pgfqpoint{0.787120in}{0.894814in}}%
\pgfpathlineto{\pgfqpoint{0.795829in}{0.893733in}}%
\pgfpathlineto{\pgfqpoint{0.779639in}{0.910810in}}%
\pgfpathlineto{\pgfqpoint{0.763458in}{0.928253in}}%
\pgfpathlineto{\pgfqpoint{0.754707in}{0.928889in}}%
\pgfpathlineto{\pgfqpoint{0.745914in}{0.929934in}}%
\pgfpathclose%
\pgfusepath{fill}%
\end{pgfscope}%
\begin{pgfscope}%
\pgfpathrectangle{\pgfqpoint{0.407592in}{0.222904in}}{\pgfqpoint{1.411163in}{1.411163in}}%
\pgfusepath{clip}%
\pgfsetbuttcap%
\pgfsetroundjoin%
\definecolor{currentfill}{rgb}{0.556863,0.654902,0.847059}%
\pgfsetfillcolor{currentfill}%
\pgfsetfillopacity{0.130000}%
\pgfsetlinewidth{0.000000pt}%
\definecolor{currentstroke}{rgb}{0.000000,0.000000,0.000000}%
\pgfsetstrokecolor{currentstroke}%
\pgfsetdash{}{0pt}%
\pgfpathmoveto{\pgfqpoint{0.598381in}{1.076080in}}%
\pgfpathlineto{\pgfqpoint{0.614579in}{1.057456in}}%
\pgfpathlineto{\pgfqpoint{0.630786in}{1.039003in}}%
\pgfpathlineto{\pgfqpoint{0.639789in}{1.038346in}}%
\pgfpathlineto{\pgfqpoint{0.648751in}{1.038021in}}%
\pgfpathlineto{\pgfqpoint{0.632587in}{1.056770in}}%
\pgfpathlineto{\pgfqpoint{0.616431in}{1.075670in}}%
\pgfpathlineto{\pgfqpoint{0.607427in}{1.075716in}}%
\pgfpathlineto{\pgfqpoint{0.598381in}{1.076080in}}%
\pgfpathclose%
\pgfusepath{fill}%
\end{pgfscope}%
\begin{pgfscope}%
\pgfpathrectangle{\pgfqpoint{0.407592in}{0.222904in}}{\pgfqpoint{1.411163in}{1.411163in}}%
\pgfusepath{clip}%
\pgfsetbuttcap%
\pgfsetroundjoin%
\definecolor{currentfill}{rgb}{0.556863,0.654902,0.847059}%
\pgfsetfillcolor{currentfill}%
\pgfsetfillopacity{0.130000}%
\pgfsetlinewidth{0.000000pt}%
\definecolor{currentstroke}{rgb}{0.000000,0.000000,0.000000}%
\pgfsetstrokecolor{currentstroke}%
\pgfsetdash{}{0pt}%
\pgfpathmoveto{\pgfqpoint{1.268827in}{0.798448in}}%
\pgfpathlineto{\pgfqpoint{1.285513in}{0.807141in}}%
\pgfpathlineto{\pgfqpoint{1.302270in}{0.816388in}}%
\pgfpathlineto{\pgfqpoint{1.310112in}{0.824395in}}%
\pgfpathlineto{\pgfqpoint{1.317940in}{0.833022in}}%
\pgfpathlineto{\pgfqpoint{1.301202in}{0.824123in}}%
\pgfpathlineto{\pgfqpoint{1.284534in}{0.815801in}}%
\pgfpathlineto{\pgfqpoint{1.276688in}{0.806783in}}%
\pgfpathlineto{\pgfqpoint{1.268827in}{0.798448in}}%
\pgfpathclose%
\pgfusepath{fill}%
\end{pgfscope}%
\begin{pgfscope}%
\pgfpathrectangle{\pgfqpoint{0.407592in}{0.222904in}}{\pgfqpoint{1.411163in}{1.411163in}}%
\pgfusepath{clip}%
\pgfsetbuttcap%
\pgfsetroundjoin%
\definecolor{currentfill}{rgb}{0.556863,0.654902,0.847059}%
\pgfsetfillcolor{currentfill}%
\pgfsetfillopacity{0.130000}%
\pgfsetlinewidth{0.000000pt}%
\definecolor{currentstroke}{rgb}{0.000000,0.000000,0.000000}%
\pgfsetstrokecolor{currentstroke}%
\pgfsetdash{}{0pt}%
\pgfpathmoveto{\pgfqpoint{1.561602in}{1.061805in}}%
\pgfpathlineto{\pgfqpoint{1.578862in}{1.071106in}}%
\pgfpathlineto{\pgfqpoint{1.596196in}{1.080712in}}%
\pgfpathlineto{\pgfqpoint{1.603740in}{1.092772in}}%
\pgfpathlineto{\pgfqpoint{1.611272in}{1.105049in}}%
\pgfpathlineto{\pgfqpoint{1.593958in}{1.095838in}}%
\pgfpathlineto{\pgfqpoint{1.576718in}{1.086945in}}%
\pgfpathlineto{\pgfqpoint{1.569166in}{1.074267in}}%
\pgfpathlineto{\pgfqpoint{1.561602in}{1.061805in}}%
\pgfpathclose%
\pgfusepath{fill}%
\end{pgfscope}%
\begin{pgfscope}%
\pgfpathrectangle{\pgfqpoint{0.407592in}{0.222904in}}{\pgfqpoint{1.411163in}{1.411163in}}%
\pgfusepath{clip}%
\pgfsetbuttcap%
\pgfsetroundjoin%
\definecolor{currentfill}{rgb}{0.556863,0.654902,0.847059}%
\pgfsetfillcolor{currentfill}%
\pgfsetfillopacity{0.130000}%
\pgfsetlinewidth{0.000000pt}%
\definecolor{currentstroke}{rgb}{0.000000,0.000000,0.000000}%
\pgfsetstrokecolor{currentstroke}%
\pgfsetdash{}{0pt}%
\pgfpathmoveto{\pgfqpoint{1.382724in}{0.890597in}}%
\pgfpathlineto{\pgfqpoint{1.399640in}{0.900001in}}%
\pgfpathlineto{\pgfqpoint{1.416629in}{0.909791in}}%
\pgfpathlineto{\pgfqpoint{1.424350in}{0.920090in}}%
\pgfpathlineto{\pgfqpoint{1.432059in}{0.930787in}}%
\pgfpathlineto{\pgfqpoint{1.415089in}{0.921451in}}%
\pgfpathlineto{\pgfqpoint{1.398191in}{0.912528in}}%
\pgfpathlineto{\pgfqpoint{1.390464in}{0.901355in}}%
\pgfpathlineto{\pgfqpoint{1.382724in}{0.890597in}}%
\pgfpathclose%
\pgfusepath{fill}%
\end{pgfscope}%
\begin{pgfscope}%
\pgfpathrectangle{\pgfqpoint{0.407592in}{0.222904in}}{\pgfqpoint{1.411163in}{1.411163in}}%
\pgfusepath{clip}%
\pgfsetbuttcap%
\pgfsetroundjoin%
\definecolor{currentfill}{rgb}{0.556863,0.654902,0.847059}%
\pgfsetfillcolor{currentfill}%
\pgfsetfillopacity{0.130000}%
\pgfsetlinewidth{0.000000pt}%
\definecolor{currentstroke}{rgb}{0.000000,0.000000,0.000000}%
\pgfsetstrokecolor{currentstroke}%
\pgfsetdash{}{0pt}%
\pgfpathmoveto{\pgfqpoint{0.975132in}{0.753395in}}%
\pgfpathlineto{\pgfqpoint{0.991407in}{0.745446in}}%
\pgfpathlineto{\pgfqpoint{1.007705in}{0.739069in}}%
\pgfpathlineto{\pgfqpoint{1.016032in}{0.736748in}}%
\pgfpathlineto{\pgfqpoint{1.024323in}{0.735092in}}%
\pgfpathlineto{\pgfqpoint{1.008070in}{0.741640in}}%
\pgfpathlineto{\pgfqpoint{0.991839in}{0.749886in}}%
\pgfpathlineto{\pgfqpoint{0.983505in}{0.751310in}}%
\pgfpathlineto{\pgfqpoint{0.975132in}{0.753395in}}%
\pgfpathclose%
\pgfusepath{fill}%
\end{pgfscope}%
\begin{pgfscope}%
\pgfpathrectangle{\pgfqpoint{0.407592in}{0.222904in}}{\pgfqpoint{1.411163in}{1.411163in}}%
\pgfusepath{clip}%
\pgfsetbuttcap%
\pgfsetroundjoin%
\definecolor{currentfill}{rgb}{0.556863,0.654902,0.847059}%
\pgfsetfillcolor{currentfill}%
\pgfsetfillopacity{0.130000}%
\pgfsetlinewidth{0.000000pt}%
\definecolor{currentstroke}{rgb}{0.000000,0.000000,0.000000}%
\pgfsetstrokecolor{currentstroke}%
\pgfsetdash{}{0pt}%
\pgfpathmoveto{\pgfqpoint{0.630786in}{1.039003in}}%
\pgfpathlineto{\pgfqpoint{0.647001in}{1.020741in}}%
\pgfpathlineto{\pgfqpoint{0.663224in}{1.002692in}}%
\pgfpathlineto{\pgfqpoint{0.672186in}{1.001700in}}%
\pgfpathlineto{\pgfqpoint{0.681105in}{1.001054in}}%
\pgfpathlineto{\pgfqpoint{0.664924in}{1.019442in}}%
\pgfpathlineto{\pgfqpoint{0.648751in}{1.038021in}}%
\pgfpathlineto{\pgfqpoint{0.639789in}{1.038346in}}%
\pgfpathlineto{\pgfqpoint{0.630786in}{1.039003in}}%
\pgfpathclose%
\pgfusepath{fill}%
\end{pgfscope}%
\begin{pgfscope}%
\pgfpathrectangle{\pgfqpoint{0.407592in}{0.222904in}}{\pgfqpoint{1.411163in}{1.411163in}}%
\pgfusepath{clip}%
\pgfsetbuttcap%
\pgfsetroundjoin%
\definecolor{currentfill}{rgb}{0.556863,0.654902,0.847059}%
\pgfsetfillcolor{currentfill}%
\pgfsetfillopacity{0.130000}%
\pgfsetlinewidth{0.000000pt}%
\definecolor{currentstroke}{rgb}{0.000000,0.000000,0.000000}%
\pgfsetstrokecolor{currentstroke}%
\pgfsetdash{}{0pt}%
\pgfpathmoveto{\pgfqpoint{1.496855in}{0.995049in}}%
\pgfpathlineto{\pgfqpoint{1.514007in}{1.004551in}}%
\pgfpathlineto{\pgfqpoint{1.531233in}{1.014372in}}%
\pgfpathlineto{\pgfqpoint{1.538842in}{1.025840in}}%
\pgfpathlineto{\pgfqpoint{1.546440in}{1.037577in}}%
\pgfpathlineto{\pgfqpoint{1.529234in}{1.028178in}}%
\pgfpathlineto{\pgfqpoint{1.512102in}{1.019116in}}%
\pgfpathlineto{\pgfqpoint{1.504485in}{1.006946in}}%
\pgfpathlineto{\pgfqpoint{1.496855in}{0.995049in}}%
\pgfpathclose%
\pgfusepath{fill}%
\end{pgfscope}%
\begin{pgfscope}%
\pgfpathrectangle{\pgfqpoint{0.407592in}{0.222904in}}{\pgfqpoint{1.411163in}{1.411163in}}%
\pgfusepath{clip}%
\pgfsetbuttcap%
\pgfsetroundjoin%
\definecolor{currentfill}{rgb}{0.556863,0.654902,0.847059}%
\pgfsetfillcolor{currentfill}%
\pgfsetfillopacity{0.130000}%
\pgfsetlinewidth{0.000000pt}%
\definecolor{currentstroke}{rgb}{0.000000,0.000000,0.000000}%
\pgfsetstrokecolor{currentstroke}%
\pgfsetdash{}{0pt}%
\pgfpathmoveto{\pgfqpoint{0.778370in}{0.896319in}}%
\pgfpathlineto{\pgfqpoint{0.794611in}{0.880157in}}%
\pgfpathlineto{\pgfqpoint{0.810860in}{0.864500in}}%
\pgfpathlineto{\pgfqpoint{0.819568in}{0.862481in}}%
\pgfpathlineto{\pgfqpoint{0.828235in}{0.860898in}}%
\pgfpathlineto{\pgfqpoint{0.812028in}{0.877075in}}%
\pgfpathlineto{\pgfqpoint{0.795829in}{0.893733in}}%
\pgfpathlineto{\pgfqpoint{0.787120in}{0.894814in}}%
\pgfpathlineto{\pgfqpoint{0.778370in}{0.896319in}}%
\pgfpathclose%
\pgfusepath{fill}%
\end{pgfscope}%
\begin{pgfscope}%
\pgfpathrectangle{\pgfqpoint{0.407592in}{0.222904in}}{\pgfqpoint{1.411163in}{1.411163in}}%
\pgfusepath{clip}%
\pgfsetbuttcap%
\pgfsetroundjoin%
\definecolor{currentfill}{rgb}{0.556863,0.654902,0.847059}%
\pgfsetfillcolor{currentfill}%
\pgfsetfillopacity{0.130000}%
\pgfsetlinewidth{0.000000pt}%
\definecolor{currentstroke}{rgb}{0.000000,0.000000,0.000000}%
\pgfsetstrokecolor{currentstroke}%
\pgfsetdash{}{0pt}%
\pgfpathmoveto{\pgfqpoint{0.893158in}{0.802566in}}%
\pgfpathlineto{\pgfqpoint{0.909415in}{0.790053in}}%
\pgfpathlineto{\pgfqpoint{0.925686in}{0.778588in}}%
\pgfpathlineto{\pgfqpoint{0.934181in}{0.775828in}}%
\pgfpathlineto{\pgfqpoint{0.942637in}{0.773594in}}%
\pgfpathlineto{\pgfqpoint{0.926411in}{0.785584in}}%
\pgfpathlineto{\pgfqpoint{0.910197in}{0.798651in}}%
\pgfpathlineto{\pgfqpoint{0.901698in}{0.800345in}}%
\pgfpathlineto{\pgfqpoint{0.893158in}{0.802566in}}%
\pgfpathclose%
\pgfusepath{fill}%
\end{pgfscope}%
\begin{pgfscope}%
\pgfpathrectangle{\pgfqpoint{0.407592in}{0.222904in}}{\pgfqpoint{1.411163in}{1.411163in}}%
\pgfusepath{clip}%
\pgfsetbuttcap%
\pgfsetroundjoin%
\definecolor{currentfill}{rgb}{0.556863,0.654902,0.847059}%
\pgfsetfillcolor{currentfill}%
\pgfsetfillopacity{0.130000}%
\pgfsetlinewidth{0.000000pt}%
\definecolor{currentstroke}{rgb}{0.000000,0.000000,0.000000}%
\pgfsetstrokecolor{currentstroke}%
\pgfsetdash{}{0pt}%
\pgfpathmoveto{\pgfqpoint{1.105968in}{0.726299in}}%
\pgfpathlineto{\pgfqpoint{1.122370in}{0.728719in}}%
\pgfpathlineto{\pgfqpoint{1.138826in}{0.732573in}}%
\pgfpathlineto{\pgfqpoint{1.146904in}{0.734104in}}%
\pgfpathlineto{\pgfqpoint{1.154954in}{0.736535in}}%
\pgfpathlineto{\pgfqpoint{1.138528in}{0.732333in}}%
\pgfpathlineto{\pgfqpoint{1.122158in}{0.729595in}}%
\pgfpathlineto{\pgfqpoint{1.114077in}{0.727461in}}%
\pgfpathlineto{\pgfqpoint{1.105968in}{0.726299in}}%
\pgfpathclose%
\pgfusepath{fill}%
\end{pgfscope}%
\begin{pgfscope}%
\pgfpathrectangle{\pgfqpoint{0.407592in}{0.222904in}}{\pgfqpoint{1.411163in}{1.411163in}}%
\pgfusepath{clip}%
\pgfsetbuttcap%
\pgfsetroundjoin%
\definecolor{currentfill}{rgb}{0.556863,0.654902,0.847059}%
\pgfsetfillcolor{currentfill}%
\pgfsetfillopacity{0.130000}%
\pgfsetlinewidth{0.000000pt}%
\definecolor{currentstroke}{rgb}{0.000000,0.000000,0.000000}%
\pgfsetstrokecolor{currentstroke}%
\pgfsetdash{}{0pt}%
\pgfpathmoveto{\pgfqpoint{1.317940in}{0.833022in}}%
\pgfpathlineto{\pgfqpoint{1.334749in}{0.842420in}}%
\pgfpathlineto{\pgfqpoint{1.351630in}{0.852252in}}%
\pgfpathlineto{\pgfqpoint{1.359424in}{0.861083in}}%
\pgfpathlineto{\pgfqpoint{1.367204in}{0.870434in}}%
\pgfpathlineto{\pgfqpoint{1.350342in}{0.860999in}}%
\pgfpathlineto{\pgfqpoint{1.333551in}{0.852024in}}%
\pgfpathlineto{\pgfqpoint{1.325752in}{0.842242in}}%
\pgfpathlineto{\pgfqpoint{1.317940in}{0.833022in}}%
\pgfpathclose%
\pgfusepath{fill}%
\end{pgfscope}%
\begin{pgfscope}%
\pgfpathrectangle{\pgfqpoint{0.407592in}{0.222904in}}{\pgfqpoint{1.411163in}{1.411163in}}%
\pgfusepath{clip}%
\pgfsetbuttcap%
\pgfsetroundjoin%
\definecolor{currentfill}{rgb}{0.556863,0.654902,0.847059}%
\pgfsetfillcolor{currentfill}%
\pgfsetfillopacity{0.130000}%
\pgfsetlinewidth{0.000000pt}%
\definecolor{currentstroke}{rgb}{0.000000,0.000000,0.000000}%
\pgfsetstrokecolor{currentstroke}%
\pgfsetdash{}{0pt}%
\pgfpathmoveto{\pgfqpoint{1.154954in}{0.736535in}}%
\pgfpathlineto{\pgfqpoint{1.171439in}{0.741996in}}%
\pgfpathlineto{\pgfqpoint{1.187987in}{0.748527in}}%
\pgfpathlineto{\pgfqpoint{1.195986in}{0.752118in}}%
\pgfpathlineto{\pgfqpoint{1.203961in}{0.756580in}}%
\pgfpathlineto{\pgfqpoint{1.187438in}{0.749888in}}%
\pgfpathlineto{\pgfqpoint{1.170978in}{0.744261in}}%
\pgfpathlineto{\pgfqpoint{1.162978in}{0.739910in}}%
\pgfpathlineto{\pgfqpoint{1.154954in}{0.736535in}}%
\pgfpathclose%
\pgfusepath{fill}%
\end{pgfscope}%
\begin{pgfscope}%
\pgfpathrectangle{\pgfqpoint{0.407592in}{0.222904in}}{\pgfqpoint{1.411163in}{1.411163in}}%
\pgfusepath{clip}%
\pgfsetbuttcap%
\pgfsetroundjoin%
\definecolor{currentfill}{rgb}{0.556863,0.654902,0.847059}%
\pgfsetfillcolor{currentfill}%
\pgfsetfillopacity{0.130000}%
\pgfsetlinewidth{0.000000pt}%
\definecolor{currentstroke}{rgb}{0.000000,0.000000,0.000000}%
\pgfsetstrokecolor{currentstroke}%
\pgfsetdash{}{0pt}%
\pgfpathmoveto{\pgfqpoint{0.663224in}{1.002692in}}%
\pgfpathlineto{\pgfqpoint{0.679456in}{0.984880in}}%
\pgfpathlineto{\pgfqpoint{0.695696in}{0.967335in}}%
\pgfpathlineto{\pgfqpoint{0.704615in}{0.965961in}}%
\pgfpathlineto{\pgfqpoint{0.713492in}{0.964947in}}%
\pgfpathlineto{\pgfqpoint{0.697294in}{0.982880in}}%
\pgfpathlineto{\pgfqpoint{0.681105in}{1.001054in}}%
\pgfpathlineto{\pgfqpoint{0.672186in}{1.001700in}}%
\pgfpathlineto{\pgfqpoint{0.663224in}{1.002692in}}%
\pgfpathclose%
\pgfusepath{fill}%
\end{pgfscope}%
\begin{pgfscope}%
\pgfpathrectangle{\pgfqpoint{0.407592in}{0.222904in}}{\pgfqpoint{1.411163in}{1.411163in}}%
\pgfusepath{clip}%
\pgfsetbuttcap%
\pgfsetroundjoin%
\definecolor{currentfill}{rgb}{0.556863,0.654902,0.847059}%
\pgfsetfillcolor{currentfill}%
\pgfsetfillopacity{0.130000}%
\pgfsetlinewidth{0.000000pt}%
\definecolor{currentstroke}{rgb}{0.000000,0.000000,0.000000}%
\pgfsetstrokecolor{currentstroke}%
\pgfsetdash{}{0pt}%
\pgfpathmoveto{\pgfqpoint{1.432059in}{0.930787in}}%
\pgfpathlineto{\pgfqpoint{1.449103in}{0.940493in}}%
\pgfpathlineto{\pgfqpoint{1.466220in}{0.950535in}}%
\pgfpathlineto{\pgfqpoint{1.473897in}{0.961167in}}%
\pgfpathlineto{\pgfqpoint{1.481561in}{0.972142in}}%
\pgfpathlineto{\pgfqpoint{1.464464in}{0.962535in}}%
\pgfpathlineto{\pgfqpoint{1.447440in}{0.953287in}}%
\pgfpathlineto{\pgfqpoint{1.439755in}{0.941860in}}%
\pgfpathlineto{\pgfqpoint{1.432059in}{0.930787in}}%
\pgfpathclose%
\pgfusepath{fill}%
\end{pgfscope}%
\begin{pgfscope}%
\pgfpathrectangle{\pgfqpoint{0.407592in}{0.222904in}}{\pgfqpoint{1.411163in}{1.411163in}}%
\pgfusepath{clip}%
\pgfsetbuttcap%
\pgfsetroundjoin%
\definecolor{currentfill}{rgb}{0.556863,0.654902,0.847059}%
\pgfsetfillcolor{currentfill}%
\pgfsetfillopacity{0.130000}%
\pgfsetlinewidth{0.000000pt}%
\definecolor{currentstroke}{rgb}{0.000000,0.000000,0.000000}%
\pgfsetstrokecolor{currentstroke}%
\pgfsetdash{}{0pt}%
\pgfpathmoveto{\pgfqpoint{1.676059in}{1.173911in}}%
\pgfpathlineto{\pgfqpoint{1.693557in}{1.183250in}}%
\pgfpathlineto{\pgfqpoint{1.711130in}{1.192858in}}%
\pgfpathlineto{\pgfqpoint{1.718564in}{1.205346in}}%
\pgfpathlineto{\pgfqpoint{1.701002in}{1.195909in}}%
\pgfpathlineto{\pgfqpoint{1.683516in}{1.186745in}}%
\pgfusepath{fill}%
\end{pgfscope}%
\begin{pgfscope}%
\pgfpathrectangle{\pgfqpoint{0.407592in}{0.222904in}}{\pgfqpoint{1.411163in}{1.411163in}}%
\pgfusepath{clip}%
\pgfsetbuttcap%
\pgfsetroundjoin%
\definecolor{currentfill}{rgb}{0.556863,0.654902,0.847059}%
\pgfsetfillcolor{currentfill}%
\pgfsetfillopacity{0.130000}%
\pgfsetlinewidth{0.000000pt}%
\definecolor{currentstroke}{rgb}{0.000000,0.000000,0.000000}%
\pgfsetstrokecolor{currentstroke}%
\pgfsetdash{}{0pt}%
\pgfpathmoveto{\pgfqpoint{1.611272in}{1.105049in}}%
\pgfpathlineto{\pgfqpoint{1.628661in}{1.114558in}}%
\pgfpathlineto{\pgfqpoint{1.646125in}{1.124345in}}%
\pgfpathlineto{\pgfqpoint{1.653625in}{1.136453in}}%
\pgfpathlineto{\pgfqpoint{1.661114in}{1.148756in}}%
\pgfpathlineto{\pgfqpoint{1.643671in}{1.139337in}}%
\pgfpathlineto{\pgfqpoint{1.626303in}{1.130208in}}%
\pgfpathlineto{\pgfqpoint{1.618793in}{1.117531in}}%
\pgfpathlineto{\pgfqpoint{1.611272in}{1.105049in}}%
\pgfpathclose%
\pgfusepath{fill}%
\end{pgfscope}%
\begin{pgfscope}%
\pgfpathrectangle{\pgfqpoint{0.407592in}{0.222904in}}{\pgfqpoint{1.411163in}{1.411163in}}%
\pgfusepath{clip}%
\pgfsetbuttcap%
\pgfsetroundjoin%
\definecolor{currentfill}{rgb}{0.556863,0.654902,0.847059}%
\pgfsetfillcolor{currentfill}%
\pgfsetfillopacity{0.130000}%
\pgfsetlinewidth{0.000000pt}%
\definecolor{currentstroke}{rgb}{0.000000,0.000000,0.000000}%
\pgfsetstrokecolor{currentstroke}%
\pgfsetdash{}{0pt}%
\pgfpathmoveto{\pgfqpoint{1.056914in}{0.727347in}}%
\pgfpathlineto{\pgfqpoint{1.073263in}{0.726135in}}%
\pgfpathlineto{\pgfqpoint{1.089654in}{0.726594in}}%
\pgfpathlineto{\pgfqpoint{1.097827in}{0.726036in}}%
\pgfpathlineto{\pgfqpoint{1.105968in}{0.726299in}}%
\pgfpathlineto{\pgfqpoint{1.089614in}{0.725496in}}%
\pgfpathlineto{\pgfqpoint{1.073305in}{0.726467in}}%
\pgfpathlineto{\pgfqpoint{1.065127in}{0.726487in}}%
\pgfpathlineto{\pgfqpoint{1.056914in}{0.727347in}}%
\pgfpathclose%
\pgfusepath{fill}%
\end{pgfscope}%
\begin{pgfscope}%
\pgfpathrectangle{\pgfqpoint{0.407592in}{0.222904in}}{\pgfqpoint{1.411163in}{1.411163in}}%
\pgfusepath{clip}%
\pgfsetbuttcap%
\pgfsetroundjoin%
\definecolor{currentfill}{rgb}{0.556863,0.654902,0.847059}%
\pgfsetfillcolor{currentfill}%
\pgfsetfillopacity{0.130000}%
\pgfsetlinewidth{0.000000pt}%
\definecolor{currentstroke}{rgb}{0.000000,0.000000,0.000000}%
\pgfsetstrokecolor{currentstroke}%
\pgfsetdash{}{0pt}%
\pgfpathmoveto{\pgfqpoint{0.515285in}{1.152597in}}%
\pgfpathlineto{\pgfqpoint{0.531491in}{1.133673in}}%
\pgfpathlineto{\pgfqpoint{0.547706in}{1.114878in}}%
\pgfpathlineto{\pgfqpoint{0.556880in}{1.114183in}}%
\pgfpathlineto{\pgfqpoint{0.566011in}{1.113777in}}%
\pgfpathlineto{\pgfqpoint{0.549838in}{1.132820in}}%
\pgfpathlineto{\pgfqpoint{0.533675in}{1.151976in}}%
\pgfpathlineto{\pgfqpoint{0.524501in}{1.152148in}}%
\pgfpathlineto{\pgfqpoint{0.515285in}{1.152597in}}%
\pgfpathclose%
\pgfusepath{fill}%
\end{pgfscope}%
\begin{pgfscope}%
\pgfpathrectangle{\pgfqpoint{0.407592in}{0.222904in}}{\pgfqpoint{1.411163in}{1.411163in}}%
\pgfusepath{clip}%
\pgfsetbuttcap%
\pgfsetroundjoin%
\definecolor{currentfill}{rgb}{0.556863,0.654902,0.847059}%
\pgfsetfillcolor{currentfill}%
\pgfsetfillopacity{0.130000}%
\pgfsetlinewidth{0.000000pt}%
\definecolor{currentstroke}{rgb}{0.000000,0.000000,0.000000}%
\pgfsetstrokecolor{currentstroke}%
\pgfsetdash{}{0pt}%
\pgfpathmoveto{\pgfqpoint{1.203961in}{0.756580in}}%
\pgfpathlineto{\pgfqpoint{1.220551in}{0.764167in}}%
\pgfpathlineto{\pgfqpoint{1.237207in}{0.772505in}}%
\pgfpathlineto{\pgfqpoint{1.245140in}{0.777842in}}%
\pgfpathlineto{\pgfqpoint{1.253054in}{0.783954in}}%
\pgfpathlineto{\pgfqpoint{1.236418in}{0.775670in}}%
\pgfpathlineto{\pgfqpoint{1.219850in}{0.768138in}}%
\pgfpathlineto{\pgfqpoint{1.211916in}{0.761921in}}%
\pgfpathlineto{\pgfqpoint{1.203961in}{0.756580in}}%
\pgfpathclose%
\pgfusepath{fill}%
\end{pgfscope}%
\begin{pgfscope}%
\pgfpathrectangle{\pgfqpoint{0.407592in}{0.222904in}}{\pgfqpoint{1.411163in}{1.411163in}}%
\pgfusepath{clip}%
\pgfsetbuttcap%
\pgfsetroundjoin%
\definecolor{currentfill}{rgb}{0.556863,0.654902,0.847059}%
\pgfsetfillcolor{currentfill}%
\pgfsetfillopacity{0.130000}%
\pgfsetlinewidth{0.000000pt}%
\definecolor{currentstroke}{rgb}{0.000000,0.000000,0.000000}%
\pgfsetstrokecolor{currentstroke}%
\pgfsetdash{}{0pt}%
\pgfpathmoveto{\pgfqpoint{0.810860in}{0.864500in}}%
\pgfpathlineto{\pgfqpoint{0.827120in}{0.849418in}}%
\pgfpathlineto{\pgfqpoint{0.843388in}{0.834985in}}%
\pgfpathlineto{\pgfqpoint{0.852053in}{0.832413in}}%
\pgfpathlineto{\pgfqpoint{0.860677in}{0.830281in}}%
\pgfpathlineto{\pgfqpoint{0.844451in}{0.845273in}}%
\pgfpathlineto{\pgfqpoint{0.828235in}{0.860898in}}%
\pgfpathlineto{\pgfqpoint{0.819568in}{0.862481in}}%
\pgfpathlineto{\pgfqpoint{0.810860in}{0.864500in}}%
\pgfpathclose%
\pgfusepath{fill}%
\end{pgfscope}%
\begin{pgfscope}%
\pgfpathrectangle{\pgfqpoint{0.407592in}{0.222904in}}{\pgfqpoint{1.411163in}{1.411163in}}%
\pgfusepath{clip}%
\pgfsetbuttcap%
\pgfsetroundjoin%
\definecolor{currentfill}{rgb}{0.556863,0.654902,0.847059}%
\pgfsetfillcolor{currentfill}%
\pgfsetfillopacity{0.130000}%
\pgfsetlinewidth{0.000000pt}%
\definecolor{currentstroke}{rgb}{0.000000,0.000000,0.000000}%
\pgfsetstrokecolor{currentstroke}%
\pgfsetdash{}{0pt}%
\pgfpathmoveto{\pgfqpoint{1.546440in}{1.037577in}}%
\pgfpathlineto{\pgfqpoint{1.563721in}{1.047288in}}%
\pgfpathlineto{\pgfqpoint{1.581076in}{1.057287in}}%
\pgfpathlineto{\pgfqpoint{1.588642in}{1.068879in}}%
\pgfpathlineto{\pgfqpoint{1.596196in}{1.080712in}}%
\pgfpathlineto{\pgfqpoint{1.578862in}{1.071106in}}%
\pgfpathlineto{\pgfqpoint{1.561602in}{1.061805in}}%
\pgfpathlineto{\pgfqpoint{1.554027in}{1.049570in}}%
\pgfpathlineto{\pgfqpoint{1.546440in}{1.037577in}}%
\pgfpathclose%
\pgfusepath{fill}%
\end{pgfscope}%
\begin{pgfscope}%
\pgfpathrectangle{\pgfqpoint{0.407592in}{0.222904in}}{\pgfqpoint{1.411163in}{1.411163in}}%
\pgfusepath{clip}%
\pgfsetbuttcap%
\pgfsetroundjoin%
\definecolor{currentfill}{rgb}{0.556863,0.654902,0.847059}%
\pgfsetfillcolor{currentfill}%
\pgfsetfillopacity{0.130000}%
\pgfsetlinewidth{0.000000pt}%
\definecolor{currentstroke}{rgb}{0.000000,0.000000,0.000000}%
\pgfsetstrokecolor{currentstroke}%
\pgfsetdash{}{0pt}%
\pgfpathmoveto{\pgfqpoint{0.695696in}{0.967335in}}%
\pgfpathlineto{\pgfqpoint{0.711945in}{0.950086in}}%
\pgfpathlineto{\pgfqpoint{0.728202in}{0.933171in}}%
\pgfpathlineto{\pgfqpoint{0.737079in}{0.931368in}}%
\pgfpathlineto{\pgfqpoint{0.745914in}{0.929934in}}%
\pgfpathlineto{\pgfqpoint{0.729699in}{0.947286in}}%
\pgfpathlineto{\pgfqpoint{0.713492in}{0.964947in}}%
\pgfpathlineto{\pgfqpoint{0.704615in}{0.965961in}}%
\pgfpathlineto{\pgfqpoint{0.695696in}{0.967335in}}%
\pgfpathclose%
\pgfusepath{fill}%
\end{pgfscope}%
\begin{pgfscope}%
\pgfpathrectangle{\pgfqpoint{0.407592in}{0.222904in}}{\pgfqpoint{1.411163in}{1.411163in}}%
\pgfusepath{clip}%
\pgfsetbuttcap%
\pgfsetroundjoin%
\definecolor{currentfill}{rgb}{0.556863,0.654902,0.847059}%
\pgfsetfillcolor{currentfill}%
\pgfsetfillopacity{0.130000}%
\pgfsetlinewidth{0.000000pt}%
\definecolor{currentstroke}{rgb}{0.000000,0.000000,0.000000}%
\pgfsetstrokecolor{currentstroke}%
\pgfsetdash{}{0pt}%
\pgfpathmoveto{\pgfqpoint{0.925686in}{0.778588in}}%
\pgfpathlineto{\pgfqpoint{0.941970in}{0.768287in}}%
\pgfpathlineto{\pgfqpoint{0.958272in}{0.759259in}}%
\pgfpathlineto{\pgfqpoint{0.966722in}{0.756067in}}%
\pgfpathlineto{\pgfqpoint{0.975132in}{0.753395in}}%
\pgfpathlineto{\pgfqpoint{0.958877in}{0.762819in}}%
\pgfpathlineto{\pgfqpoint{0.942637in}{0.773594in}}%
\pgfpathlineto{\pgfqpoint{0.934181in}{0.775828in}}%
\pgfpathlineto{\pgfqpoint{0.925686in}{0.778588in}}%
\pgfpathclose%
\pgfusepath{fill}%
\end{pgfscope}%
\begin{pgfscope}%
\pgfpathrectangle{\pgfqpoint{0.407592in}{0.222904in}}{\pgfqpoint{1.411163in}{1.411163in}}%
\pgfusepath{clip}%
\pgfsetbuttcap%
\pgfsetroundjoin%
\definecolor{currentfill}{rgb}{0.556863,0.654902,0.847059}%
\pgfsetfillcolor{currentfill}%
\pgfsetfillopacity{0.130000}%
\pgfsetlinewidth{0.000000pt}%
\definecolor{currentstroke}{rgb}{0.000000,0.000000,0.000000}%
\pgfsetstrokecolor{currentstroke}%
\pgfsetdash{}{0pt}%
\pgfpathmoveto{\pgfqpoint{1.367204in}{0.870434in}}%
\pgfpathlineto{\pgfqpoint{1.384140in}{0.880275in}}%
\pgfpathlineto{\pgfqpoint{1.401148in}{0.890475in}}%
\pgfpathlineto{\pgfqpoint{1.408895in}{0.899912in}}%
\pgfpathlineto{\pgfqpoint{1.416629in}{0.909791in}}%
\pgfpathlineto{\pgfqpoint{1.399640in}{0.900001in}}%
\pgfpathlineto{\pgfqpoint{1.382724in}{0.890597in}}%
\pgfpathlineto{\pgfqpoint{1.374970in}{0.880281in}}%
\pgfpathlineto{\pgfqpoint{1.367204in}{0.870434in}}%
\pgfpathclose%
\pgfusepath{fill}%
\end{pgfscope}%
\begin{pgfscope}%
\pgfpathrectangle{\pgfqpoint{0.407592in}{0.222904in}}{\pgfqpoint{1.411163in}{1.411163in}}%
\pgfusepath{clip}%
\pgfsetbuttcap%
\pgfsetroundjoin%
\definecolor{currentfill}{rgb}{0.556863,0.654902,0.847059}%
\pgfsetfillcolor{currentfill}%
\pgfsetfillopacity{0.130000}%
\pgfsetlinewidth{0.000000pt}%
\definecolor{currentstroke}{rgb}{0.000000,0.000000,0.000000}%
\pgfsetstrokecolor{currentstroke}%
\pgfsetdash{}{0pt}%
\pgfpathmoveto{\pgfqpoint{0.547706in}{1.114878in}}%
\pgfpathlineto{\pgfqpoint{0.563929in}{1.096224in}}%
\pgfpathlineto{\pgfqpoint{0.580161in}{1.077728in}}%
\pgfpathlineto{\pgfqpoint{0.589292in}{1.076754in}}%
\pgfpathlineto{\pgfqpoint{0.598381in}{1.076080in}}%
\pgfpathlineto{\pgfqpoint{0.582192in}{1.094859in}}%
\pgfpathlineto{\pgfqpoint{0.566011in}{1.113777in}}%
\pgfpathlineto{\pgfqpoint{0.556880in}{1.114183in}}%
\pgfpathlineto{\pgfqpoint{0.547706in}{1.114878in}}%
\pgfpathclose%
\pgfusepath{fill}%
\end{pgfscope}%
\begin{pgfscope}%
\pgfpathrectangle{\pgfqpoint{0.407592in}{0.222904in}}{\pgfqpoint{1.411163in}{1.411163in}}%
\pgfusepath{clip}%
\pgfsetbuttcap%
\pgfsetroundjoin%
\definecolor{currentfill}{rgb}{0.556863,0.654902,0.847059}%
\pgfsetfillcolor{currentfill}%
\pgfsetfillopacity{0.130000}%
\pgfsetlinewidth{0.000000pt}%
\definecolor{currentstroke}{rgb}{0.000000,0.000000,0.000000}%
\pgfsetstrokecolor{currentstroke}%
\pgfsetdash{}{0pt}%
\pgfpathmoveto{\pgfqpoint{1.253054in}{0.783954in}}%
\pgfpathlineto{\pgfqpoint{1.269759in}{0.792873in}}%
\pgfpathlineto{\pgfqpoint{1.286535in}{0.802331in}}%
\pgfpathlineto{\pgfqpoint{1.294411in}{0.809027in}}%
\pgfpathlineto{\pgfqpoint{1.302270in}{0.816388in}}%
\pgfpathlineto{\pgfqpoint{1.285513in}{0.807141in}}%
\pgfpathlineto{\pgfqpoint{1.268827in}{0.798448in}}%
\pgfpathlineto{\pgfqpoint{1.260949in}{0.790829in}}%
\pgfpathlineto{\pgfqpoint{1.253054in}{0.783954in}}%
\pgfpathclose%
\pgfusepath{fill}%
\end{pgfscope}%
\begin{pgfscope}%
\pgfpathrectangle{\pgfqpoint{0.407592in}{0.222904in}}{\pgfqpoint{1.411163in}{1.411163in}}%
\pgfusepath{clip}%
\pgfsetbuttcap%
\pgfsetroundjoin%
\definecolor{currentfill}{rgb}{0.556863,0.654902,0.847059}%
\pgfsetfillcolor{currentfill}%
\pgfsetfillopacity{0.130000}%
\pgfsetlinewidth{0.000000pt}%
\definecolor{currentstroke}{rgb}{0.000000,0.000000,0.000000}%
\pgfsetstrokecolor{currentstroke}%
\pgfsetdash{}{0pt}%
\pgfpathmoveto{\pgfqpoint{1.007705in}{0.739069in}}%
\pgfpathlineto{\pgfqpoint{1.024029in}{0.734327in}}%
\pgfpathlineto{\pgfqpoint{1.040384in}{0.731236in}}%
\pgfpathlineto{\pgfqpoint{1.048667in}{0.728959in}}%
\pgfpathlineto{\pgfqpoint{1.056914in}{0.727347in}}%
\pgfpathlineto{\pgfqpoint{1.040603in}{0.730320in}}%
\pgfpathlineto{\pgfqpoint{1.024323in}{0.735092in}}%
\pgfpathlineto{\pgfqpoint{1.016032in}{0.736748in}}%
\pgfpathlineto{\pgfqpoint{1.007705in}{0.739069in}}%
\pgfpathclose%
\pgfusepath{fill}%
\end{pgfscope}%
\begin{pgfscope}%
\pgfpathrectangle{\pgfqpoint{0.407592in}{0.222904in}}{\pgfqpoint{1.411163in}{1.411163in}}%
\pgfusepath{clip}%
\pgfsetbuttcap%
\pgfsetroundjoin%
\definecolor{currentfill}{rgb}{0.556863,0.654902,0.847059}%
\pgfsetfillcolor{currentfill}%
\pgfsetfillopacity{0.130000}%
\pgfsetlinewidth{0.000000pt}%
\definecolor{currentstroke}{rgb}{0.000000,0.000000,0.000000}%
\pgfsetstrokecolor{currentstroke}%
\pgfsetdash{}{0pt}%
\pgfpathmoveto{\pgfqpoint{1.481561in}{0.972142in}}%
\pgfpathlineto{\pgfqpoint{1.498733in}{0.982075in}}%
\pgfpathlineto{\pgfqpoint{1.515980in}{0.992308in}}%
\pgfpathlineto{\pgfqpoint{1.523612in}{1.003190in}}%
\pgfpathlineto{\pgfqpoint{1.531233in}{1.014372in}}%
\pgfpathlineto{\pgfqpoint{1.514007in}{1.004551in}}%
\pgfpathlineto{\pgfqpoint{1.496855in}{0.995049in}}%
\pgfpathlineto{\pgfqpoint{1.489214in}{0.983442in}}%
\pgfpathlineto{\pgfqpoint{1.481561in}{0.972142in}}%
\pgfpathclose%
\pgfusepath{fill}%
\end{pgfscope}%
\begin{pgfscope}%
\pgfpathrectangle{\pgfqpoint{0.407592in}{0.222904in}}{\pgfqpoint{1.411163in}{1.411163in}}%
\pgfusepath{clip}%
\pgfsetbuttcap%
\pgfsetroundjoin%
\definecolor{currentfill}{rgb}{0.556863,0.654902,0.847059}%
\pgfsetfillcolor{currentfill}%
\pgfsetfillopacity{0.130000}%
\pgfsetlinewidth{0.000000pt}%
\definecolor{currentstroke}{rgb}{0.000000,0.000000,0.000000}%
\pgfsetstrokecolor{currentstroke}%
\pgfsetdash{}{0pt}%
\pgfpathmoveto{\pgfqpoint{0.843388in}{0.834985in}}%
\pgfpathlineto{\pgfqpoint{0.859667in}{0.821282in}}%
\pgfpathlineto{\pgfqpoint{0.875957in}{0.808400in}}%
\pgfpathlineto{\pgfqpoint{0.884578in}{0.805265in}}%
\pgfpathlineto{\pgfqpoint{0.893158in}{0.802566in}}%
\pgfpathlineto{\pgfqpoint{0.876912in}{0.816011in}}%
\pgfpathlineto{\pgfqpoint{0.860677in}{0.830281in}}%
\pgfpathlineto{\pgfqpoint{0.852053in}{0.832413in}}%
\pgfpathlineto{\pgfqpoint{0.843388in}{0.834985in}}%
\pgfpathclose%
\pgfusepath{fill}%
\end{pgfscope}%
\begin{pgfscope}%
\pgfpathrectangle{\pgfqpoint{0.407592in}{0.222904in}}{\pgfqpoint{1.411163in}{1.411163in}}%
\pgfusepath{clip}%
\pgfsetbuttcap%
\pgfsetroundjoin%
\definecolor{currentfill}{rgb}{0.556863,0.654902,0.847059}%
\pgfsetfillcolor{currentfill}%
\pgfsetfillopacity{0.130000}%
\pgfsetlinewidth{0.000000pt}%
\definecolor{currentstroke}{rgb}{0.000000,0.000000,0.000000}%
\pgfsetstrokecolor{currentstroke}%
\pgfsetdash{}{0pt}%
\pgfpathmoveto{\pgfqpoint{1.661114in}{1.148756in}}%
\pgfpathlineto{\pgfqpoint{1.678633in}{1.158449in}}%
\pgfpathlineto{\pgfqpoint{1.696228in}{1.168401in}}%
\pgfpathlineto{\pgfqpoint{1.703684in}{1.180540in}}%
\pgfpathlineto{\pgfqpoint{1.711130in}{1.192858in}}%
\pgfpathlineto{\pgfqpoint{1.693557in}{1.183250in}}%
\pgfpathlineto{\pgfqpoint{1.676059in}{1.173911in}}%
\pgfpathlineto{\pgfqpoint{1.668592in}{1.161245in}}%
\pgfpathlineto{\pgfqpoint{1.661114in}{1.148756in}}%
\pgfpathclose%
\pgfusepath{fill}%
\end{pgfscope}%
\begin{pgfscope}%
\pgfpathrectangle{\pgfqpoint{0.407592in}{0.222904in}}{\pgfqpoint{1.411163in}{1.411163in}}%
\pgfusepath{clip}%
\pgfsetbuttcap%
\pgfsetroundjoin%
\definecolor{currentfill}{rgb}{0.556863,0.654902,0.847059}%
\pgfsetfillcolor{currentfill}%
\pgfsetfillopacity{0.130000}%
\pgfsetlinewidth{0.000000pt}%
\definecolor{currentstroke}{rgb}{0.000000,0.000000,0.000000}%
\pgfsetstrokecolor{currentstroke}%
\pgfsetdash{}{0pt}%
\pgfpathmoveto{\pgfqpoint{0.580161in}{1.077728in}}%
\pgfpathlineto{\pgfqpoint{0.596401in}{1.059404in}}%
\pgfpathlineto{\pgfqpoint{0.612650in}{1.041271in}}%
\pgfpathlineto{\pgfqpoint{0.621739in}{1.039982in}}%
\pgfpathlineto{\pgfqpoint{0.630786in}{1.039003in}}%
\pgfpathlineto{\pgfqpoint{0.614579in}{1.057456in}}%
\pgfpathlineto{\pgfqpoint{0.598381in}{1.076080in}}%
\pgfpathlineto{\pgfqpoint{0.589292in}{1.076754in}}%
\pgfpathlineto{\pgfqpoint{0.580161in}{1.077728in}}%
\pgfpathclose%
\pgfusepath{fill}%
\end{pgfscope}%
\begin{pgfscope}%
\pgfpathrectangle{\pgfqpoint{0.407592in}{0.222904in}}{\pgfqpoint{1.411163in}{1.411163in}}%
\pgfusepath{clip}%
\pgfsetbuttcap%
\pgfsetroundjoin%
\definecolor{currentfill}{rgb}{0.556863,0.654902,0.847059}%
\pgfsetfillcolor{currentfill}%
\pgfsetfillopacity{0.130000}%
\pgfsetlinewidth{0.000000pt}%
\definecolor{currentstroke}{rgb}{0.000000,0.000000,0.000000}%
\pgfsetstrokecolor{currentstroke}%
\pgfsetdash{}{0pt}%
\pgfpathmoveto{\pgfqpoint{0.728202in}{0.933171in}}%
\pgfpathlineto{\pgfqpoint{0.744468in}{0.916630in}}%
\pgfpathlineto{\pgfqpoint{0.760743in}{0.900507in}}%
\pgfpathlineto{\pgfqpoint{0.769577in}{0.898225in}}%
\pgfpathlineto{\pgfqpoint{0.778370in}{0.896319in}}%
\pgfpathlineto{\pgfqpoint{0.762138in}{0.912929in}}%
\pgfpathlineto{\pgfqpoint{0.745914in}{0.929934in}}%
\pgfpathlineto{\pgfqpoint{0.737079in}{0.931368in}}%
\pgfpathlineto{\pgfqpoint{0.728202in}{0.933171in}}%
\pgfpathclose%
\pgfusepath{fill}%
\end{pgfscope}%
\begin{pgfscope}%
\pgfpathrectangle{\pgfqpoint{0.407592in}{0.222904in}}{\pgfqpoint{1.411163in}{1.411163in}}%
\pgfusepath{clip}%
\pgfsetbuttcap%
\pgfsetroundjoin%
\definecolor{currentfill}{rgb}{0.556863,0.654902,0.847059}%
\pgfsetfillcolor{currentfill}%
\pgfsetfillopacity{0.130000}%
\pgfsetlinewidth{0.000000pt}%
\definecolor{currentstroke}{rgb}{0.000000,0.000000,0.000000}%
\pgfsetstrokecolor{currentstroke}%
\pgfsetdash{}{0pt}%
\pgfpathmoveto{\pgfqpoint{1.302270in}{0.816388in}}%
\pgfpathlineto{\pgfqpoint{1.319098in}{0.826110in}}%
\pgfpathlineto{\pgfqpoint{1.335998in}{0.836242in}}%
\pgfpathlineto{\pgfqpoint{1.343822in}{0.843965in}}%
\pgfpathlineto{\pgfqpoint{1.351630in}{0.852252in}}%
\pgfpathlineto{\pgfqpoint{1.334749in}{0.842420in}}%
\pgfpathlineto{\pgfqpoint{1.317940in}{0.833022in}}%
\pgfpathlineto{\pgfqpoint{1.310112in}{0.824395in}}%
\pgfpathlineto{\pgfqpoint{1.302270in}{0.816388in}}%
\pgfpathclose%
\pgfusepath{fill}%
\end{pgfscope}%
\begin{pgfscope}%
\pgfpathrectangle{\pgfqpoint{0.407592in}{0.222904in}}{\pgfqpoint{1.411163in}{1.411163in}}%
\pgfusepath{clip}%
\pgfsetbuttcap%
\pgfsetroundjoin%
\definecolor{currentfill}{rgb}{0.556863,0.654902,0.847059}%
\pgfsetfillcolor{currentfill}%
\pgfsetfillopacity{0.130000}%
\pgfsetlinewidth{0.000000pt}%
\definecolor{currentstroke}{rgb}{0.000000,0.000000,0.000000}%
\pgfsetstrokecolor{currentstroke}%
\pgfsetdash{}{0pt}%
\pgfpathmoveto{\pgfqpoint{1.596196in}{1.080712in}}%
\pgfpathlineto{\pgfqpoint{1.613607in}{1.090602in}}%
\pgfpathlineto{\pgfqpoint{1.631093in}{1.100757in}}%
\pgfpathlineto{\pgfqpoint{1.638615in}{1.112443in}}%
\pgfpathlineto{\pgfqpoint{1.646125in}{1.124345in}}%
\pgfpathlineto{\pgfqpoint{1.628661in}{1.114558in}}%
\pgfpathlineto{\pgfqpoint{1.611272in}{1.105049in}}%
\pgfpathlineto{\pgfqpoint{1.603740in}{1.092772in}}%
\pgfpathlineto{\pgfqpoint{1.596196in}{1.080712in}}%
\pgfpathclose%
\pgfusepath{fill}%
\end{pgfscope}%
\begin{pgfscope}%
\pgfpathrectangle{\pgfqpoint{0.407592in}{0.222904in}}{\pgfqpoint{1.411163in}{1.411163in}}%
\pgfusepath{clip}%
\pgfsetbuttcap%
\pgfsetroundjoin%
\definecolor{currentfill}{rgb}{0.556863,0.654902,0.847059}%
\pgfsetfillcolor{currentfill}%
\pgfsetfillopacity{0.130000}%
\pgfsetlinewidth{0.000000pt}%
\definecolor{currentstroke}{rgb}{0.000000,0.000000,0.000000}%
\pgfsetstrokecolor{currentstroke}%
\pgfsetdash{}{0pt}%
\pgfpathmoveto{\pgfqpoint{1.416629in}{0.909791in}}%
\pgfpathlineto{\pgfqpoint{1.433692in}{0.919927in}}%
\pgfpathlineto{\pgfqpoint{1.450830in}{0.930373in}}%
\pgfpathlineto{\pgfqpoint{1.458531in}{0.940264in}}%
\pgfpathlineto{\pgfqpoint{1.466220in}{0.950535in}}%
\pgfpathlineto{\pgfqpoint{1.449103in}{0.940493in}}%
\pgfpathlineto{\pgfqpoint{1.432059in}{0.930787in}}%
\pgfpathlineto{\pgfqpoint{1.424350in}{0.920090in}}%
\pgfpathlineto{\pgfqpoint{1.416629in}{0.909791in}}%
\pgfpathclose%
\pgfusepath{fill}%
\end{pgfscope}%
\begin{pgfscope}%
\pgfpathrectangle{\pgfqpoint{0.407592in}{0.222904in}}{\pgfqpoint{1.411163in}{1.411163in}}%
\pgfusepath{clip}%
\pgfsetbuttcap%
\pgfsetroundjoin%
\definecolor{currentfill}{rgb}{0.556863,0.654902,0.847059}%
\pgfsetfillcolor{currentfill}%
\pgfsetfillopacity{0.130000}%
\pgfsetlinewidth{0.000000pt}%
\definecolor{currentstroke}{rgb}{0.000000,0.000000,0.000000}%
\pgfsetstrokecolor{currentstroke}%
\pgfsetdash{}{0pt}%
\pgfpathmoveto{\pgfqpoint{0.612650in}{1.041271in}}%
\pgfpathlineto{\pgfqpoint{0.628907in}{1.023350in}}%
\pgfpathlineto{\pgfqpoint{0.645172in}{1.005662in}}%
\pgfpathlineto{\pgfqpoint{0.654220in}{1.004017in}}%
\pgfpathlineto{\pgfqpoint{0.663224in}{1.002692in}}%
\pgfpathlineto{\pgfqpoint{0.647001in}{1.020741in}}%
\pgfpathlineto{\pgfqpoint{0.630786in}{1.039003in}}%
\pgfpathlineto{\pgfqpoint{0.621739in}{1.039982in}}%
\pgfpathlineto{\pgfqpoint{0.612650in}{1.041271in}}%
\pgfpathclose%
\pgfusepath{fill}%
\end{pgfscope}%
\begin{pgfscope}%
\pgfpathrectangle{\pgfqpoint{0.407592in}{0.222904in}}{\pgfqpoint{1.411163in}{1.411163in}}%
\pgfusepath{clip}%
\pgfsetbuttcap%
\pgfsetroundjoin%
\definecolor{currentfill}{rgb}{0.556863,0.654902,0.847059}%
\pgfsetfillcolor{currentfill}%
\pgfsetfillopacity{0.130000}%
\pgfsetlinewidth{0.000000pt}%
\definecolor{currentstroke}{rgb}{0.000000,0.000000,0.000000}%
\pgfsetstrokecolor{currentstroke}%
\pgfsetdash{}{0pt}%
\pgfpathmoveto{\pgfqpoint{0.958272in}{0.759259in}}%
\pgfpathlineto{\pgfqpoint{0.974593in}{0.751603in}}%
\pgfpathlineto{\pgfqpoint{0.990936in}{0.745394in}}%
\pgfpathlineto{\pgfqpoint{0.999340in}{0.741975in}}%
\pgfpathlineto{\pgfqpoint{1.007705in}{0.739069in}}%
\pgfpathlineto{\pgfqpoint{0.991407in}{0.745446in}}%
\pgfpathlineto{\pgfqpoint{0.975132in}{0.753395in}}%
\pgfpathlineto{\pgfqpoint{0.966722in}{0.756067in}}%
\pgfpathlineto{\pgfqpoint{0.958272in}{0.759259in}}%
\pgfpathclose%
\pgfusepath{fill}%
\end{pgfscope}%
\begin{pgfscope}%
\pgfpathrectangle{\pgfqpoint{0.407592in}{0.222904in}}{\pgfqpoint{1.411163in}{1.411163in}}%
\pgfusepath{clip}%
\pgfsetbuttcap%
\pgfsetroundjoin%
\definecolor{currentfill}{rgb}{0.556863,0.654902,0.847059}%
\pgfsetfillcolor{currentfill}%
\pgfsetfillopacity{0.130000}%
\pgfsetlinewidth{0.000000pt}%
\definecolor{currentstroke}{rgb}{0.000000,0.000000,0.000000}%
\pgfsetstrokecolor{currentstroke}%
\pgfsetdash{}{0pt}%
\pgfpathmoveto{\pgfqpoint{1.138826in}{0.732573in}}%
\pgfpathlineto{\pgfqpoint{1.155340in}{0.737677in}}%
\pgfpathlineto{\pgfqpoint{1.171915in}{0.743857in}}%
\pgfpathlineto{\pgfqpoint{1.179964in}{0.745784in}}%
\pgfpathlineto{\pgfqpoint{1.187987in}{0.748527in}}%
\pgfpathlineto{\pgfqpoint{1.171439in}{0.741996in}}%
\pgfpathlineto{\pgfqpoint{1.154954in}{0.736535in}}%
\pgfpathlineto{\pgfqpoint{1.146904in}{0.734104in}}%
\pgfpathlineto{\pgfqpoint{1.138826in}{0.732573in}}%
\pgfpathclose%
\pgfusepath{fill}%
\end{pgfscope}%
\begin{pgfscope}%
\pgfpathrectangle{\pgfqpoint{0.407592in}{0.222904in}}{\pgfqpoint{1.411163in}{1.411163in}}%
\pgfusepath{clip}%
\pgfsetbuttcap%
\pgfsetroundjoin%
\definecolor{currentfill}{rgb}{0.556863,0.654902,0.847059}%
\pgfsetfillcolor{currentfill}%
\pgfsetfillopacity{0.130000}%
\pgfsetlinewidth{0.000000pt}%
\definecolor{currentstroke}{rgb}{0.000000,0.000000,0.000000}%
\pgfsetstrokecolor{currentstroke}%
\pgfsetdash{}{0pt}%
\pgfpathmoveto{\pgfqpoint{0.760743in}{0.900507in}}%
\pgfpathlineto{\pgfqpoint{0.777027in}{0.884854in}}%
\pgfpathlineto{\pgfqpoint{0.793319in}{0.869728in}}%
\pgfpathlineto{\pgfqpoint{0.802111in}{0.866925in}}%
\pgfpathlineto{\pgfqpoint{0.810860in}{0.864500in}}%
\pgfpathlineto{\pgfqpoint{0.794611in}{0.880157in}}%
\pgfpathlineto{\pgfqpoint{0.778370in}{0.896319in}}%
\pgfpathlineto{\pgfqpoint{0.769577in}{0.898225in}}%
\pgfpathlineto{\pgfqpoint{0.760743in}{0.900507in}}%
\pgfpathclose%
\pgfusepath{fill}%
\end{pgfscope}%
\begin{pgfscope}%
\pgfpathrectangle{\pgfqpoint{0.407592in}{0.222904in}}{\pgfqpoint{1.411163in}{1.411163in}}%
\pgfusepath{clip}%
\pgfsetbuttcap%
\pgfsetroundjoin%
\definecolor{currentfill}{rgb}{0.556863,0.654902,0.847059}%
\pgfsetfillcolor{currentfill}%
\pgfsetfillopacity{0.130000}%
\pgfsetlinewidth{0.000000pt}%
\definecolor{currentstroke}{rgb}{0.000000,0.000000,0.000000}%
\pgfsetstrokecolor{currentstroke}%
\pgfsetdash{}{0pt}%
\pgfpathmoveto{\pgfqpoint{1.531233in}{1.014372in}}%
\pgfpathlineto{\pgfqpoint{1.548535in}{1.024487in}}%
\pgfpathlineto{\pgfqpoint{1.565912in}{1.034874in}}%
\pgfpathlineto{\pgfqpoint{1.573500in}{1.045948in}}%
\pgfpathlineto{\pgfqpoint{1.581076in}{1.057287in}}%
\pgfpathlineto{\pgfqpoint{1.563721in}{1.047288in}}%
\pgfpathlineto{\pgfqpoint{1.546440in}{1.037577in}}%
\pgfpathlineto{\pgfqpoint{1.538842in}{1.025840in}}%
\pgfpathlineto{\pgfqpoint{1.531233in}{1.014372in}}%
\pgfpathclose%
\pgfusepath{fill}%
\end{pgfscope}%
\begin{pgfscope}%
\pgfpathrectangle{\pgfqpoint{0.407592in}{0.222904in}}{\pgfqpoint{1.411163in}{1.411163in}}%
\pgfusepath{clip}%
\pgfsetbuttcap%
\pgfsetroundjoin%
\definecolor{currentfill}{rgb}{0.556863,0.654902,0.847059}%
\pgfsetfillcolor{currentfill}%
\pgfsetfillopacity{0.130000}%
\pgfsetlinewidth{0.000000pt}%
\definecolor{currentstroke}{rgb}{0.000000,0.000000,0.000000}%
\pgfsetstrokecolor{currentstroke}%
\pgfsetdash{}{0pt}%
\pgfpathmoveto{\pgfqpoint{1.089654in}{0.726594in}}%
\pgfpathlineto{\pgfqpoint{1.106092in}{0.728593in}}%
\pgfpathlineto{\pgfqpoint{1.122582in}{0.731981in}}%
\pgfpathlineto{\pgfqpoint{1.130719in}{0.731886in}}%
\pgfpathlineto{\pgfqpoint{1.138826in}{0.732573in}}%
\pgfpathlineto{\pgfqpoint{1.122370in}{0.728719in}}%
\pgfpathlineto{\pgfqpoint{1.105968in}{0.726299in}}%
\pgfpathlineto{\pgfqpoint{1.097827in}{0.726036in}}%
\pgfpathlineto{\pgfqpoint{1.089654in}{0.726594in}}%
\pgfpathclose%
\pgfusepath{fill}%
\end{pgfscope}%
\begin{pgfscope}%
\pgfpathrectangle{\pgfqpoint{0.407592in}{0.222904in}}{\pgfqpoint{1.411163in}{1.411163in}}%
\pgfusepath{clip}%
\pgfsetbuttcap%
\pgfsetroundjoin%
\definecolor{currentfill}{rgb}{0.556863,0.654902,0.847059}%
\pgfsetfillcolor{currentfill}%
\pgfsetfillopacity{0.130000}%
\pgfsetlinewidth{0.000000pt}%
\definecolor{currentstroke}{rgb}{0.000000,0.000000,0.000000}%
\pgfsetstrokecolor{currentstroke}%
\pgfsetdash{}{0pt}%
\pgfpathmoveto{\pgfqpoint{0.875957in}{0.808400in}}%
\pgfpathlineto{\pgfqpoint{0.892259in}{0.796432in}}%
\pgfpathlineto{\pgfqpoint{0.908575in}{0.785475in}}%
\pgfpathlineto{\pgfqpoint{0.917150in}{0.781820in}}%
\pgfpathlineto{\pgfqpoint{0.925686in}{0.778588in}}%
\pgfpathlineto{\pgfqpoint{0.909415in}{0.790053in}}%
\pgfpathlineto{\pgfqpoint{0.893158in}{0.802566in}}%
\pgfpathlineto{\pgfqpoint{0.884578in}{0.805265in}}%
\pgfpathlineto{\pgfqpoint{0.875957in}{0.808400in}}%
\pgfpathclose%
\pgfusepath{fill}%
\end{pgfscope}%
\begin{pgfscope}%
\pgfpathrectangle{\pgfqpoint{0.407592in}{0.222904in}}{\pgfqpoint{1.411163in}{1.411163in}}%
\pgfusepath{clip}%
\pgfsetbuttcap%
\pgfsetroundjoin%
\definecolor{currentfill}{rgb}{0.556863,0.654902,0.847059}%
\pgfsetfillcolor{currentfill}%
\pgfsetfillopacity{0.130000}%
\pgfsetlinewidth{0.000000pt}%
\definecolor{currentstroke}{rgb}{0.000000,0.000000,0.000000}%
\pgfsetstrokecolor{currentstroke}%
\pgfsetdash{}{0pt}%
\pgfpathmoveto{\pgfqpoint{1.187987in}{0.748527in}}%
\pgfpathlineto{\pgfqpoint{1.204599in}{0.755962in}}%
\pgfpathlineto{\pgfqpoint{1.221279in}{0.764160in}}%
\pgfpathlineto{\pgfqpoint{1.229254in}{0.767947in}}%
\pgfpathlineto{\pgfqpoint{1.237207in}{0.772505in}}%
\pgfpathlineto{\pgfqpoint{1.220551in}{0.764167in}}%
\pgfpathlineto{\pgfqpoint{1.203961in}{0.756580in}}%
\pgfpathlineto{\pgfqpoint{1.195986in}{0.752118in}}%
\pgfpathlineto{\pgfqpoint{1.187987in}{0.748527in}}%
\pgfpathclose%
\pgfusepath{fill}%
\end{pgfscope}%
\begin{pgfscope}%
\pgfpathrectangle{\pgfqpoint{0.407592in}{0.222904in}}{\pgfqpoint{1.411163in}{1.411163in}}%
\pgfusepath{clip}%
\pgfsetbuttcap%
\pgfsetroundjoin%
\definecolor{currentfill}{rgb}{0.556863,0.654902,0.847059}%
\pgfsetfillcolor{currentfill}%
\pgfsetfillopacity{0.130000}%
\pgfsetlinewidth{0.000000pt}%
\definecolor{currentstroke}{rgb}{0.000000,0.000000,0.000000}%
\pgfsetstrokecolor{currentstroke}%
\pgfsetdash{}{0pt}%
\pgfpathmoveto{\pgfqpoint{1.351630in}{0.852252in}}%
\pgfpathlineto{\pgfqpoint{1.368585in}{0.862462in}}%
\pgfpathlineto{\pgfqpoint{1.385613in}{0.873006in}}%
\pgfpathlineto{\pgfqpoint{1.393388in}{0.881500in}}%
\pgfpathlineto{\pgfqpoint{1.401148in}{0.890475in}}%
\pgfpathlineto{\pgfqpoint{1.384140in}{0.880275in}}%
\pgfpathlineto{\pgfqpoint{1.367204in}{0.870434in}}%
\pgfpathlineto{\pgfqpoint{1.359424in}{0.861083in}}%
\pgfpathlineto{\pgfqpoint{1.351630in}{0.852252in}}%
\pgfpathclose%
\pgfusepath{fill}%
\end{pgfscope}%
\begin{pgfscope}%
\pgfpathrectangle{\pgfqpoint{0.407592in}{0.222904in}}{\pgfqpoint{1.411163in}{1.411163in}}%
\pgfusepath{clip}%
\pgfsetbuttcap%
\pgfsetroundjoin%
\definecolor{currentfill}{rgb}{0.556863,0.654902,0.847059}%
\pgfsetfillcolor{currentfill}%
\pgfsetfillopacity{0.130000}%
\pgfsetlinewidth{0.000000pt}%
\definecolor{currentstroke}{rgb}{0.000000,0.000000,0.000000}%
\pgfsetstrokecolor{currentstroke}%
\pgfsetdash{}{0pt}%
\pgfpathmoveto{\pgfqpoint{0.645172in}{1.005662in}}%
\pgfpathlineto{\pgfqpoint{0.661447in}{0.988235in}}%
\pgfpathlineto{\pgfqpoint{0.677729in}{0.971095in}}%
\pgfpathlineto{\pgfqpoint{0.686734in}{0.969051in}}%
\pgfpathlineto{\pgfqpoint{0.695696in}{0.967335in}}%
\pgfpathlineto{\pgfqpoint{0.679456in}{0.984880in}}%
\pgfpathlineto{\pgfqpoint{0.663224in}{1.002692in}}%
\pgfpathlineto{\pgfqpoint{0.654220in}{1.004017in}}%
\pgfpathlineto{\pgfqpoint{0.645172in}{1.005662in}}%
\pgfpathclose%
\pgfusepath{fill}%
\end{pgfscope}%
\begin{pgfscope}%
\pgfpathrectangle{\pgfqpoint{0.407592in}{0.222904in}}{\pgfqpoint{1.411163in}{1.411163in}}%
\pgfusepath{clip}%
\pgfsetbuttcap%
\pgfsetroundjoin%
\definecolor{currentfill}{rgb}{0.556863,0.654902,0.847059}%
\pgfsetfillcolor{currentfill}%
\pgfsetfillopacity{0.130000}%
\pgfsetlinewidth{0.000000pt}%
\definecolor{currentstroke}{rgb}{0.000000,0.000000,0.000000}%
\pgfsetstrokecolor{currentstroke}%
\pgfsetdash{}{0pt}%
\pgfpathmoveto{\pgfqpoint{1.040384in}{0.731236in}}%
\pgfpathlineto{\pgfqpoint{1.056775in}{0.729769in}}%
\pgfpathlineto{\pgfqpoint{1.073207in}{0.729855in}}%
\pgfpathlineto{\pgfqpoint{1.081448in}{0.727893in}}%
\pgfpathlineto{\pgfqpoint{1.089654in}{0.726594in}}%
\pgfpathlineto{\pgfqpoint{1.073263in}{0.726135in}}%
\pgfpathlineto{\pgfqpoint{1.056914in}{0.727347in}}%
\pgfpathlineto{\pgfqpoint{1.048667in}{0.728959in}}%
\pgfpathlineto{\pgfqpoint{1.040384in}{0.731236in}}%
\pgfpathclose%
\pgfusepath{fill}%
\end{pgfscope}%
\begin{pgfscope}%
\pgfpathrectangle{\pgfqpoint{0.407592in}{0.222904in}}{\pgfqpoint{1.411163in}{1.411163in}}%
\pgfusepath{clip}%
\pgfsetbuttcap%
\pgfsetroundjoin%
\definecolor{currentfill}{rgb}{0.556863,0.654902,0.847059}%
\pgfsetfillcolor{currentfill}%
\pgfsetfillopacity{0.130000}%
\pgfsetlinewidth{0.000000pt}%
\definecolor{currentstroke}{rgb}{0.000000,0.000000,0.000000}%
\pgfsetstrokecolor{currentstroke}%
\pgfsetdash{}{0pt}%
\pgfpathmoveto{\pgfqpoint{0.496721in}{1.154304in}}%
\pgfpathlineto{\pgfqpoint{0.512970in}{1.135632in}}%
\pgfpathlineto{\pgfqpoint{0.529227in}{1.117104in}}%
\pgfpathlineto{\pgfqpoint{0.538488in}{1.115854in}}%
\pgfpathlineto{\pgfqpoint{0.547706in}{1.114878in}}%
\pgfpathlineto{\pgfqpoint{0.531491in}{1.133673in}}%
\pgfpathlineto{\pgfqpoint{0.515285in}{1.152597in}}%
\pgfpathlineto{\pgfqpoint{0.506025in}{1.153318in}}%
\pgfpathlineto{\pgfqpoint{0.496721in}{1.154304in}}%
\pgfpathclose%
\pgfusepath{fill}%
\end{pgfscope}%
\begin{pgfscope}%
\pgfpathrectangle{\pgfqpoint{0.407592in}{0.222904in}}{\pgfqpoint{1.411163in}{1.411163in}}%
\pgfusepath{clip}%
\pgfsetbuttcap%
\pgfsetroundjoin%
\definecolor{currentfill}{rgb}{0.556863,0.654902,0.847059}%
\pgfsetfillcolor{currentfill}%
\pgfsetfillopacity{0.130000}%
\pgfsetlinewidth{0.000000pt}%
\definecolor{currentstroke}{rgb}{0.000000,0.000000,0.000000}%
\pgfsetstrokecolor{currentstroke}%
\pgfsetdash{}{0pt}%
\pgfpathmoveto{\pgfqpoint{1.711130in}{1.192858in}}%
\pgfpathlineto{\pgfqpoint{1.728779in}{1.202720in}}%
\pgfpathlineto{\pgfqpoint{1.746506in}{1.212824in}}%
\pgfpathlineto{\pgfqpoint{1.753918in}{1.224984in}}%
\pgfpathlineto{\pgfqpoint{1.736203in}{1.215042in}}%
\pgfpathlineto{\pgfqpoint{1.718564in}{1.205346in}}%
\pgfusepath{fill}%
\end{pgfscope}%
\begin{pgfscope}%
\pgfpathrectangle{\pgfqpoint{0.407592in}{0.222904in}}{\pgfqpoint{1.411163in}{1.411163in}}%
\pgfusepath{clip}%
\pgfsetbuttcap%
\pgfsetroundjoin%
\definecolor{currentfill}{rgb}{0.556863,0.654902,0.847059}%
\pgfsetfillcolor{currentfill}%
\pgfsetfillopacity{0.130000}%
\pgfsetlinewidth{0.000000pt}%
\definecolor{currentstroke}{rgb}{0.000000,0.000000,0.000000}%
\pgfsetstrokecolor{currentstroke}%
\pgfsetdash{}{0pt}%
\pgfpathmoveto{\pgfqpoint{1.466220in}{0.950535in}}%
\pgfpathlineto{\pgfqpoint{1.483412in}{0.960881in}}%
\pgfpathlineto{\pgfqpoint{1.500680in}{0.971504in}}%
\pgfpathlineto{\pgfqpoint{1.508336in}{0.981741in}}%
\pgfpathlineto{\pgfqpoint{1.515980in}{0.992308in}}%
\pgfpathlineto{\pgfqpoint{1.498733in}{0.982075in}}%
\pgfpathlineto{\pgfqpoint{1.481561in}{0.972142in}}%
\pgfpathlineto{\pgfqpoint{1.473897in}{0.961167in}}%
\pgfpathlineto{\pgfqpoint{1.466220in}{0.950535in}}%
\pgfpathclose%
\pgfusepath{fill}%
\end{pgfscope}%
\begin{pgfscope}%
\pgfpathrectangle{\pgfqpoint{0.407592in}{0.222904in}}{\pgfqpoint{1.411163in}{1.411163in}}%
\pgfusepath{clip}%
\pgfsetbuttcap%
\pgfsetroundjoin%
\definecolor{currentfill}{rgb}{0.556863,0.654902,0.847059}%
\pgfsetfillcolor{currentfill}%
\pgfsetfillopacity{0.130000}%
\pgfsetlinewidth{0.000000pt}%
\definecolor{currentstroke}{rgb}{0.000000,0.000000,0.000000}%
\pgfsetstrokecolor{currentstroke}%
\pgfsetdash{}{0pt}%
\pgfpathmoveto{\pgfqpoint{1.646125in}{1.124345in}}%
\pgfpathlineto{\pgfqpoint{1.663666in}{1.134395in}}%
\pgfpathlineto{\pgfqpoint{1.681284in}{1.144691in}}%
\pgfpathlineto{\pgfqpoint{1.688761in}{1.156448in}}%
\pgfpathlineto{\pgfqpoint{1.696228in}{1.168401in}}%
\pgfpathlineto{\pgfqpoint{1.678633in}{1.158449in}}%
\pgfpathlineto{\pgfqpoint{1.661114in}{1.148756in}}%
\pgfpathlineto{\pgfqpoint{1.653625in}{1.136453in}}%
\pgfpathlineto{\pgfqpoint{1.646125in}{1.124345in}}%
\pgfpathclose%
\pgfusepath{fill}%
\end{pgfscope}%
\begin{pgfscope}%
\pgfpathrectangle{\pgfqpoint{0.407592in}{0.222904in}}{\pgfqpoint{1.411163in}{1.411163in}}%
\pgfusepath{clip}%
\pgfsetbuttcap%
\pgfsetroundjoin%
\definecolor{currentfill}{rgb}{0.556863,0.654902,0.847059}%
\pgfsetfillcolor{currentfill}%
\pgfsetfillopacity{0.130000}%
\pgfsetlinewidth{0.000000pt}%
\definecolor{currentstroke}{rgb}{0.000000,0.000000,0.000000}%
\pgfsetstrokecolor{currentstroke}%
\pgfsetdash{}{0pt}%
\pgfpathmoveto{\pgfqpoint{1.237207in}{0.772505in}}%
\pgfpathlineto{\pgfqpoint{1.253933in}{0.781477in}}%
\pgfpathlineto{\pgfqpoint{1.270730in}{0.790984in}}%
\pgfpathlineto{\pgfqpoint{1.278642in}{0.796313in}}%
\pgfpathlineto{\pgfqpoint{1.286535in}{0.802331in}}%
\pgfpathlineto{\pgfqpoint{1.269759in}{0.792873in}}%
\pgfpathlineto{\pgfqpoint{1.253054in}{0.783954in}}%
\pgfpathlineto{\pgfqpoint{1.245140in}{0.777842in}}%
\pgfpathlineto{\pgfqpoint{1.237207in}{0.772505in}}%
\pgfpathclose%
\pgfusepath{fill}%
\end{pgfscope}%
\begin{pgfscope}%
\pgfpathrectangle{\pgfqpoint{0.407592in}{0.222904in}}{\pgfqpoint{1.411163in}{1.411163in}}%
\pgfusepath{clip}%
\pgfsetbuttcap%
\pgfsetroundjoin%
\definecolor{currentfill}{rgb}{0.556863,0.654902,0.847059}%
\pgfsetfillcolor{currentfill}%
\pgfsetfillopacity{0.130000}%
\pgfsetlinewidth{0.000000pt}%
\definecolor{currentstroke}{rgb}{0.000000,0.000000,0.000000}%
\pgfsetstrokecolor{currentstroke}%
\pgfsetdash{}{0pt}%
\pgfpathmoveto{\pgfqpoint{0.793319in}{0.869728in}}%
\pgfpathlineto{\pgfqpoint{0.809622in}{0.855190in}}%
\pgfpathlineto{\pgfqpoint{0.825934in}{0.841309in}}%
\pgfpathlineto{\pgfqpoint{0.834682in}{0.837961in}}%
\pgfpathlineto{\pgfqpoint{0.843388in}{0.834985in}}%
\pgfpathlineto{\pgfqpoint{0.827120in}{0.849418in}}%
\pgfpathlineto{\pgfqpoint{0.810860in}{0.864500in}}%
\pgfpathlineto{\pgfqpoint{0.802111in}{0.866925in}}%
\pgfpathlineto{\pgfqpoint{0.793319in}{0.869728in}}%
\pgfpathclose%
\pgfusepath{fill}%
\end{pgfscope}%
\begin{pgfscope}%
\pgfpathrectangle{\pgfqpoint{0.407592in}{0.222904in}}{\pgfqpoint{1.411163in}{1.411163in}}%
\pgfusepath{clip}%
\pgfsetbuttcap%
\pgfsetroundjoin%
\definecolor{currentfill}{rgb}{0.556863,0.654902,0.847059}%
\pgfsetfillcolor{currentfill}%
\pgfsetfillopacity{0.130000}%
\pgfsetlinewidth{0.000000pt}%
\definecolor{currentstroke}{rgb}{0.000000,0.000000,0.000000}%
\pgfsetstrokecolor{currentstroke}%
\pgfsetdash{}{0pt}%
\pgfpathmoveto{\pgfqpoint{1.581076in}{1.057287in}}%
\pgfpathlineto{\pgfqpoint{1.598508in}{1.067556in}}%
\pgfpathlineto{\pgfqpoint{1.616016in}{1.078076in}}%
\pgfpathlineto{\pgfqpoint{1.623560in}{1.089298in}}%
\pgfpathlineto{\pgfqpoint{1.631093in}{1.100757in}}%
\pgfpathlineto{\pgfqpoint{1.613607in}{1.090602in}}%
\pgfpathlineto{\pgfqpoint{1.596196in}{1.080712in}}%
\pgfpathlineto{\pgfqpoint{1.588642in}{1.068879in}}%
\pgfpathlineto{\pgfqpoint{1.581076in}{1.057287in}}%
\pgfpathclose%
\pgfusepath{fill}%
\end{pgfscope}%
\begin{pgfscope}%
\pgfpathrectangle{\pgfqpoint{0.407592in}{0.222904in}}{\pgfqpoint{1.411163in}{1.411163in}}%
\pgfusepath{clip}%
\pgfsetbuttcap%
\pgfsetroundjoin%
\definecolor{currentfill}{rgb}{0.556863,0.654902,0.847059}%
\pgfsetfillcolor{currentfill}%
\pgfsetfillopacity{0.130000}%
\pgfsetlinewidth{0.000000pt}%
\definecolor{currentstroke}{rgb}{0.000000,0.000000,0.000000}%
\pgfsetstrokecolor{currentstroke}%
\pgfsetdash{}{0pt}%
\pgfpathmoveto{\pgfqpoint{0.677729in}{0.971095in}}%
\pgfpathlineto{\pgfqpoint{0.694020in}{0.954275in}}%
\pgfpathlineto{\pgfqpoint{0.710320in}{0.937811in}}%
\pgfpathlineto{\pgfqpoint{0.719283in}{0.935325in}}%
\pgfpathlineto{\pgfqpoint{0.728202in}{0.933171in}}%
\pgfpathlineto{\pgfqpoint{0.711945in}{0.950086in}}%
\pgfpathlineto{\pgfqpoint{0.695696in}{0.967335in}}%
\pgfpathlineto{\pgfqpoint{0.686734in}{0.969051in}}%
\pgfpathlineto{\pgfqpoint{0.677729in}{0.971095in}}%
\pgfpathclose%
\pgfusepath{fill}%
\end{pgfscope}%
\begin{pgfscope}%
\pgfpathrectangle{\pgfqpoint{0.407592in}{0.222904in}}{\pgfqpoint{1.411163in}{1.411163in}}%
\pgfusepath{clip}%
\pgfsetbuttcap%
\pgfsetroundjoin%
\definecolor{currentfill}{rgb}{0.556863,0.654902,0.847059}%
\pgfsetfillcolor{currentfill}%
\pgfsetfillopacity{0.130000}%
\pgfsetlinewidth{0.000000pt}%
\definecolor{currentstroke}{rgb}{0.000000,0.000000,0.000000}%
\pgfsetstrokecolor{currentstroke}%
\pgfsetdash{}{0pt}%
\pgfpathmoveto{\pgfqpoint{0.529227in}{1.117104in}}%
\pgfpathlineto{\pgfqpoint{0.545492in}{1.098734in}}%
\pgfpathlineto{\pgfqpoint{0.561766in}{1.080538in}}%
\pgfpathlineto{\pgfqpoint{0.570985in}{1.078992in}}%
\pgfpathlineto{\pgfqpoint{0.580161in}{1.077728in}}%
\pgfpathlineto{\pgfqpoint{0.563929in}{1.096224in}}%
\pgfpathlineto{\pgfqpoint{0.547706in}{1.114878in}}%
\pgfpathlineto{\pgfqpoint{0.538488in}{1.115854in}}%
\pgfpathlineto{\pgfqpoint{0.529227in}{1.117104in}}%
\pgfpathclose%
\pgfusepath{fill}%
\end{pgfscope}%
\begin{pgfscope}%
\pgfpathrectangle{\pgfqpoint{0.407592in}{0.222904in}}{\pgfqpoint{1.411163in}{1.411163in}}%
\pgfusepath{clip}%
\pgfsetbuttcap%
\pgfsetroundjoin%
\definecolor{currentfill}{rgb}{0.556863,0.654902,0.847059}%
\pgfsetfillcolor{currentfill}%
\pgfsetfillopacity{0.130000}%
\pgfsetlinewidth{0.000000pt}%
\definecolor{currentstroke}{rgb}{0.000000,0.000000,0.000000}%
\pgfsetstrokecolor{currentstroke}%
\pgfsetdash{}{0pt}%
\pgfpathmoveto{\pgfqpoint{0.908575in}{0.785475in}}%
\pgfpathlineto{\pgfqpoint{0.924906in}{0.775624in}}%
\pgfpathlineto{\pgfqpoint{0.941254in}{0.766967in}}%
\pgfpathlineto{\pgfqpoint{0.949783in}{0.762910in}}%
\pgfpathlineto{\pgfqpoint{0.958272in}{0.759259in}}%
\pgfpathlineto{\pgfqpoint{0.941970in}{0.768287in}}%
\pgfpathlineto{\pgfqpoint{0.925686in}{0.778588in}}%
\pgfpathlineto{\pgfqpoint{0.917150in}{0.781820in}}%
\pgfpathlineto{\pgfqpoint{0.908575in}{0.785475in}}%
\pgfpathclose%
\pgfusepath{fill}%
\end{pgfscope}%
\begin{pgfscope}%
\pgfpathrectangle{\pgfqpoint{0.407592in}{0.222904in}}{\pgfqpoint{1.411163in}{1.411163in}}%
\pgfusepath{clip}%
\pgfsetbuttcap%
\pgfsetroundjoin%
\definecolor{currentfill}{rgb}{0.556863,0.654902,0.847059}%
\pgfsetfillcolor{currentfill}%
\pgfsetfillopacity{0.130000}%
\pgfsetlinewidth{0.000000pt}%
\definecolor{currentstroke}{rgb}{0.000000,0.000000,0.000000}%
\pgfsetstrokecolor{currentstroke}%
\pgfsetdash{}{0pt}%
\pgfpathmoveto{\pgfqpoint{1.401148in}{0.890475in}}%
\pgfpathlineto{\pgfqpoint{1.418231in}{0.900994in}}%
\pgfpathlineto{\pgfqpoint{1.435389in}{0.911799in}}%
\pgfpathlineto{\pgfqpoint{1.443116in}{0.920879in}}%
\pgfpathlineto{\pgfqpoint{1.450830in}{0.930373in}}%
\pgfpathlineto{\pgfqpoint{1.433692in}{0.919927in}}%
\pgfpathlineto{\pgfqpoint{1.416629in}{0.909791in}}%
\pgfpathlineto{\pgfqpoint{1.408895in}{0.899912in}}%
\pgfpathlineto{\pgfqpoint{1.401148in}{0.890475in}}%
\pgfpathclose%
\pgfusepath{fill}%
\end{pgfscope}%
\begin{pgfscope}%
\pgfpathrectangle{\pgfqpoint{0.407592in}{0.222904in}}{\pgfqpoint{1.411163in}{1.411163in}}%
\pgfusepath{clip}%
\pgfsetbuttcap%
\pgfsetroundjoin%
\definecolor{currentfill}{rgb}{0.556863,0.654902,0.847059}%
\pgfsetfillcolor{currentfill}%
\pgfsetfillopacity{0.130000}%
\pgfsetlinewidth{0.000000pt}%
\definecolor{currentstroke}{rgb}{0.000000,0.000000,0.000000}%
\pgfsetstrokecolor{currentstroke}%
\pgfsetdash{}{0pt}%
\pgfpathmoveto{\pgfqpoint{1.286535in}{0.802331in}}%
\pgfpathlineto{\pgfqpoint{1.303383in}{0.812248in}}%
\pgfpathlineto{\pgfqpoint{1.320303in}{0.822560in}}%
\pgfpathlineto{\pgfqpoint{1.328159in}{0.829102in}}%
\pgfpathlineto{\pgfqpoint{1.335998in}{0.836242in}}%
\pgfpathlineto{\pgfqpoint{1.319098in}{0.826110in}}%
\pgfpathlineto{\pgfqpoint{1.302270in}{0.816388in}}%
\pgfpathlineto{\pgfqpoint{1.294411in}{0.809027in}}%
\pgfpathlineto{\pgfqpoint{1.286535in}{0.802331in}}%
\pgfpathclose%
\pgfusepath{fill}%
\end{pgfscope}%
\begin{pgfscope}%
\pgfpathrectangle{\pgfqpoint{0.407592in}{0.222904in}}{\pgfqpoint{1.411163in}{1.411163in}}%
\pgfusepath{clip}%
\pgfsetbuttcap%
\pgfsetroundjoin%
\definecolor{currentfill}{rgb}{0.556863,0.654902,0.847059}%
\pgfsetfillcolor{currentfill}%
\pgfsetfillopacity{0.130000}%
\pgfsetlinewidth{0.000000pt}%
\definecolor{currentstroke}{rgb}{0.000000,0.000000,0.000000}%
\pgfsetstrokecolor{currentstroke}%
\pgfsetdash{}{0pt}%
\pgfpathmoveto{\pgfqpoint{0.990936in}{0.745394in}}%
\pgfpathlineto{\pgfqpoint{1.007307in}{0.740681in}}%
\pgfpathlineto{\pgfqpoint{1.023708in}{0.737478in}}%
\pgfpathlineto{\pgfqpoint{1.032064in}{0.734099in}}%
\pgfpathlineto{\pgfqpoint{1.040384in}{0.731236in}}%
\pgfpathlineto{\pgfqpoint{1.024029in}{0.734327in}}%
\pgfpathlineto{\pgfqpoint{1.007705in}{0.739069in}}%
\pgfpathlineto{\pgfqpoint{0.999340in}{0.741975in}}%
\pgfpathlineto{\pgfqpoint{0.990936in}{0.745394in}}%
\pgfpathclose%
\pgfusepath{fill}%
\end{pgfscope}%
\begin{pgfscope}%
\pgfpathrectangle{\pgfqpoint{0.407592in}{0.222904in}}{\pgfqpoint{1.411163in}{1.411163in}}%
\pgfusepath{clip}%
\pgfsetbuttcap%
\pgfsetroundjoin%
\definecolor{currentfill}{rgb}{0.556863,0.654902,0.847059}%
\pgfsetfillcolor{currentfill}%
\pgfsetfillopacity{0.130000}%
\pgfsetlinewidth{0.000000pt}%
\definecolor{currentstroke}{rgb}{0.000000,0.000000,0.000000}%
\pgfsetstrokecolor{currentstroke}%
\pgfsetdash{}{0pt}%
\pgfpathmoveto{\pgfqpoint{1.515980in}{0.992308in}}%
\pgfpathlineto{\pgfqpoint{1.533302in}{1.002815in}}%
\pgfpathlineto{\pgfqpoint{1.550701in}{1.013574in}}%
\pgfpathlineto{\pgfqpoint{1.558312in}{1.024079in}}%
\pgfpathlineto{\pgfqpoint{1.565912in}{1.034874in}}%
\pgfpathlineto{\pgfqpoint{1.548535in}{1.024487in}}%
\pgfpathlineto{\pgfqpoint{1.531233in}{1.014372in}}%
\pgfpathlineto{\pgfqpoint{1.523612in}{1.003190in}}%
\pgfpathlineto{\pgfqpoint{1.515980in}{0.992308in}}%
\pgfpathclose%
\pgfusepath{fill}%
\end{pgfscope}%
\begin{pgfscope}%
\pgfpathrectangle{\pgfqpoint{0.407592in}{0.222904in}}{\pgfqpoint{1.411163in}{1.411163in}}%
\pgfusepath{clip}%
\pgfsetbuttcap%
\pgfsetroundjoin%
\definecolor{currentfill}{rgb}{0.556863,0.654902,0.847059}%
\pgfsetfillcolor{currentfill}%
\pgfsetfillopacity{0.130000}%
\pgfsetlinewidth{0.000000pt}%
\definecolor{currentstroke}{rgb}{0.000000,0.000000,0.000000}%
\pgfsetstrokecolor{currentstroke}%
\pgfsetdash{}{0pt}%
\pgfpathmoveto{\pgfqpoint{1.696228in}{1.168401in}}%
\pgfpathlineto{\pgfqpoint{1.713901in}{1.178597in}}%
\pgfpathlineto{\pgfqpoint{1.731650in}{1.189024in}}%
\pgfpathlineto{\pgfqpoint{1.739084in}{1.200835in}}%
\pgfpathlineto{\pgfqpoint{1.746506in}{1.212824in}}%
\pgfpathlineto{\pgfqpoint{1.728779in}{1.202720in}}%
\pgfpathlineto{\pgfqpoint{1.711130in}{1.192858in}}%
\pgfpathlineto{\pgfqpoint{1.703684in}{1.180540in}}%
\pgfpathlineto{\pgfqpoint{1.696228in}{1.168401in}}%
\pgfpathclose%
\pgfusepath{fill}%
\end{pgfscope}%
\begin{pgfscope}%
\pgfpathrectangle{\pgfqpoint{0.407592in}{0.222904in}}{\pgfqpoint{1.411163in}{1.411163in}}%
\pgfusepath{clip}%
\pgfsetbuttcap%
\pgfsetroundjoin%
\definecolor{currentfill}{rgb}{0.556863,0.654902,0.847059}%
\pgfsetfillcolor{currentfill}%
\pgfsetfillopacity{0.130000}%
\pgfsetlinewidth{0.000000pt}%
\definecolor{currentstroke}{rgb}{0.000000,0.000000,0.000000}%
\pgfsetstrokecolor{currentstroke}%
\pgfsetdash{}{0pt}%
\pgfpathmoveto{\pgfqpoint{0.561766in}{1.080538in}}%
\pgfpathlineto{\pgfqpoint{0.578049in}{1.062533in}}%
\pgfpathlineto{\pgfqpoint{0.594340in}{1.044738in}}%
\pgfpathlineto{\pgfqpoint{0.603517in}{1.042860in}}%
\pgfpathlineto{\pgfqpoint{0.612650in}{1.041271in}}%
\pgfpathlineto{\pgfqpoint{0.596401in}{1.059404in}}%
\pgfpathlineto{\pgfqpoint{0.580161in}{1.077728in}}%
\pgfpathlineto{\pgfqpoint{0.570985in}{1.078992in}}%
\pgfpathlineto{\pgfqpoint{0.561766in}{1.080538in}}%
\pgfpathclose%
\pgfusepath{fill}%
\end{pgfscope}%
\begin{pgfscope}%
\pgfpathrectangle{\pgfqpoint{0.407592in}{0.222904in}}{\pgfqpoint{1.411163in}{1.411163in}}%
\pgfusepath{clip}%
\pgfsetbuttcap%
\pgfsetroundjoin%
\definecolor{currentfill}{rgb}{0.556863,0.654902,0.847059}%
\pgfsetfillcolor{currentfill}%
\pgfsetfillopacity{0.130000}%
\pgfsetlinewidth{0.000000pt}%
\definecolor{currentstroke}{rgb}{0.000000,0.000000,0.000000}%
\pgfsetstrokecolor{currentstroke}%
\pgfsetdash{}{0pt}%
\pgfpathmoveto{\pgfqpoint{0.825934in}{0.841309in}}%
\pgfpathlineto{\pgfqpoint{0.842258in}{0.828159in}}%
\pgfpathlineto{\pgfqpoint{0.858593in}{0.815816in}}%
\pgfpathlineto{\pgfqpoint{0.867296in}{0.811929in}}%
\pgfpathlineto{\pgfqpoint{0.875957in}{0.808400in}}%
\pgfpathlineto{\pgfqpoint{0.859667in}{0.821282in}}%
\pgfpathlineto{\pgfqpoint{0.843388in}{0.834985in}}%
\pgfpathlineto{\pgfqpoint{0.834682in}{0.837961in}}%
\pgfpathlineto{\pgfqpoint{0.825934in}{0.841309in}}%
\pgfpathclose%
\pgfusepath{fill}%
\end{pgfscope}%
\begin{pgfscope}%
\pgfpathrectangle{\pgfqpoint{0.407592in}{0.222904in}}{\pgfqpoint{1.411163in}{1.411163in}}%
\pgfusepath{clip}%
\pgfsetbuttcap%
\pgfsetroundjoin%
\definecolor{currentfill}{rgb}{0.556863,0.654902,0.847059}%
\pgfsetfillcolor{currentfill}%
\pgfsetfillopacity{0.130000}%
\pgfsetlinewidth{0.000000pt}%
\definecolor{currentstroke}{rgb}{0.000000,0.000000,0.000000}%
\pgfsetstrokecolor{currentstroke}%
\pgfsetdash{}{0pt}%
\pgfpathmoveto{\pgfqpoint{0.710320in}{0.937811in}}%
\pgfpathlineto{\pgfqpoint{0.726629in}{0.921741in}}%
\pgfpathlineto{\pgfqpoint{0.742947in}{0.906110in}}%
\pgfpathlineto{\pgfqpoint{0.751866in}{0.903143in}}%
\pgfpathlineto{\pgfqpoint{0.760743in}{0.900507in}}%
\pgfpathlineto{\pgfqpoint{0.744468in}{0.916630in}}%
\pgfpathlineto{\pgfqpoint{0.728202in}{0.933171in}}%
\pgfpathlineto{\pgfqpoint{0.719283in}{0.935325in}}%
\pgfpathlineto{\pgfqpoint{0.710320in}{0.937811in}}%
\pgfpathclose%
\pgfusepath{fill}%
\end{pgfscope}%
\begin{pgfscope}%
\pgfpathrectangle{\pgfqpoint{0.407592in}{0.222904in}}{\pgfqpoint{1.411163in}{1.411163in}}%
\pgfusepath{clip}%
\pgfsetbuttcap%
\pgfsetroundjoin%
\definecolor{currentfill}{rgb}{0.556863,0.654902,0.847059}%
\pgfsetfillcolor{currentfill}%
\pgfsetfillopacity{0.130000}%
\pgfsetlinewidth{0.000000pt}%
\definecolor{currentstroke}{rgb}{0.000000,0.000000,0.000000}%
\pgfsetstrokecolor{currentstroke}%
\pgfsetdash{}{0pt}%
\pgfpathmoveto{\pgfqpoint{1.631093in}{1.100757in}}%
\pgfpathlineto{\pgfqpoint{1.648656in}{1.111162in}}%
\pgfpathlineto{\pgfqpoint{1.666296in}{1.121801in}}%
\pgfpathlineto{\pgfqpoint{1.673795in}{1.133139in}}%
\pgfpathlineto{\pgfqpoint{1.681284in}{1.144691in}}%
\pgfpathlineto{\pgfqpoint{1.663666in}{1.134395in}}%
\pgfpathlineto{\pgfqpoint{1.646125in}{1.124345in}}%
\pgfpathlineto{\pgfqpoint{1.638615in}{1.112443in}}%
\pgfpathlineto{\pgfqpoint{1.631093in}{1.100757in}}%
\pgfpathclose%
\pgfusepath{fill}%
\end{pgfscope}%
\begin{pgfscope}%
\pgfpathrectangle{\pgfqpoint{0.407592in}{0.222904in}}{\pgfqpoint{1.411163in}{1.411163in}}%
\pgfusepath{clip}%
\pgfsetbuttcap%
\pgfsetroundjoin%
\definecolor{currentfill}{rgb}{0.556863,0.654902,0.847059}%
\pgfsetfillcolor{currentfill}%
\pgfsetfillopacity{0.130000}%
\pgfsetlinewidth{0.000000pt}%
\definecolor{currentstroke}{rgb}{0.000000,0.000000,0.000000}%
\pgfsetstrokecolor{currentstroke}%
\pgfsetdash{}{0pt}%
\pgfpathmoveto{\pgfqpoint{1.335998in}{0.836242in}}%
\pgfpathlineto{\pgfqpoint{1.352972in}{0.846730in}}%
\pgfpathlineto{\pgfqpoint{1.370020in}{0.857532in}}%
\pgfpathlineto{\pgfqpoint{1.377824in}{0.865011in}}%
\pgfpathlineto{\pgfqpoint{1.385613in}{0.873006in}}%
\pgfpathlineto{\pgfqpoint{1.368585in}{0.862462in}}%
\pgfpathlineto{\pgfqpoint{1.351630in}{0.852252in}}%
\pgfpathlineto{\pgfqpoint{1.343822in}{0.843965in}}%
\pgfpathlineto{\pgfqpoint{1.335998in}{0.836242in}}%
\pgfpathclose%
\pgfusepath{fill}%
\end{pgfscope}%
\begin{pgfscope}%
\pgfpathrectangle{\pgfqpoint{0.407592in}{0.222904in}}{\pgfqpoint{1.411163in}{1.411163in}}%
\pgfusepath{clip}%
\pgfsetbuttcap%
\pgfsetroundjoin%
\definecolor{currentfill}{rgb}{0.556863,0.654902,0.847059}%
\pgfsetfillcolor{currentfill}%
\pgfsetfillopacity{0.130000}%
\pgfsetlinewidth{0.000000pt}%
\definecolor{currentstroke}{rgb}{0.000000,0.000000,0.000000}%
\pgfsetstrokecolor{currentstroke}%
\pgfsetdash{}{0pt}%
\pgfpathmoveto{\pgfqpoint{1.450830in}{0.930373in}}%
\pgfpathlineto{\pgfqpoint{1.468042in}{0.941100in}}%
\pgfpathlineto{\pgfqpoint{1.485330in}{0.952082in}}%
\pgfpathlineto{\pgfqpoint{1.493011in}{0.961613in}}%
\pgfpathlineto{\pgfqpoint{1.500680in}{0.971504in}}%
\pgfpathlineto{\pgfqpoint{1.483412in}{0.960881in}}%
\pgfpathlineto{\pgfqpoint{1.466220in}{0.950535in}}%
\pgfpathlineto{\pgfqpoint{1.458531in}{0.940264in}}%
\pgfpathlineto{\pgfqpoint{1.450830in}{0.930373in}}%
\pgfpathclose%
\pgfusepath{fill}%
\end{pgfscope}%
\begin{pgfscope}%
\pgfpathrectangle{\pgfqpoint{0.407592in}{0.222904in}}{\pgfqpoint{1.411163in}{1.411163in}}%
\pgfusepath{clip}%
\pgfsetbuttcap%
\pgfsetroundjoin%
\definecolor{currentfill}{rgb}{0.556863,0.654902,0.847059}%
\pgfsetfillcolor{currentfill}%
\pgfsetfillopacity{0.130000}%
\pgfsetlinewidth{0.000000pt}%
\definecolor{currentstroke}{rgb}{0.000000,0.000000,0.000000}%
\pgfsetstrokecolor{currentstroke}%
\pgfsetdash{}{0pt}%
\pgfpathmoveto{\pgfqpoint{1.122582in}{0.731981in}}%
\pgfpathlineto{\pgfqpoint{1.139128in}{0.736598in}}%
\pgfpathlineto{\pgfqpoint{1.155734in}{0.742290in}}%
\pgfpathlineto{\pgfqpoint{1.163838in}{0.742708in}}%
\pgfpathlineto{\pgfqpoint{1.171915in}{0.743857in}}%
\pgfpathlineto{\pgfqpoint{1.155340in}{0.737677in}}%
\pgfpathlineto{\pgfqpoint{1.138826in}{0.732573in}}%
\pgfpathlineto{\pgfqpoint{1.130719in}{0.731886in}}%
\pgfpathlineto{\pgfqpoint{1.122582in}{0.731981in}}%
\pgfpathclose%
\pgfusepath{fill}%
\end{pgfscope}%
\begin{pgfscope}%
\pgfpathrectangle{\pgfqpoint{0.407592in}{0.222904in}}{\pgfqpoint{1.411163in}{1.411163in}}%
\pgfusepath{clip}%
\pgfsetbuttcap%
\pgfsetroundjoin%
\definecolor{currentfill}{rgb}{0.556863,0.654902,0.847059}%
\pgfsetfillcolor{currentfill}%
\pgfsetfillopacity{0.130000}%
\pgfsetlinewidth{0.000000pt}%
\definecolor{currentstroke}{rgb}{0.000000,0.000000,0.000000}%
\pgfsetstrokecolor{currentstroke}%
\pgfsetdash{}{0pt}%
\pgfpathmoveto{\pgfqpoint{0.594340in}{1.044738in}}%
\pgfpathlineto{\pgfqpoint{0.610639in}{1.027172in}}%
\pgfpathlineto{\pgfqpoint{0.626947in}{1.009861in}}%
\pgfpathlineto{\pgfqpoint{0.636082in}{1.007614in}}%
\pgfpathlineto{\pgfqpoint{0.645172in}{1.005662in}}%
\pgfpathlineto{\pgfqpoint{0.628907in}{1.023350in}}%
\pgfpathlineto{\pgfqpoint{0.612650in}{1.041271in}}%
\pgfpathlineto{\pgfqpoint{0.603517in}{1.042860in}}%
\pgfpathlineto{\pgfqpoint{0.594340in}{1.044738in}}%
\pgfpathclose%
\pgfusepath{fill}%
\end{pgfscope}%
\begin{pgfscope}%
\pgfpathrectangle{\pgfqpoint{0.407592in}{0.222904in}}{\pgfqpoint{1.411163in}{1.411163in}}%
\pgfusepath{clip}%
\pgfsetbuttcap%
\pgfsetroundjoin%
\definecolor{currentfill}{rgb}{0.556863,0.654902,0.847059}%
\pgfsetfillcolor{currentfill}%
\pgfsetfillopacity{0.130000}%
\pgfsetlinewidth{0.000000pt}%
\definecolor{currentstroke}{rgb}{0.000000,0.000000,0.000000}%
\pgfsetstrokecolor{currentstroke}%
\pgfsetdash{}{0pt}%
\pgfpathmoveto{\pgfqpoint{1.171915in}{0.743857in}}%
\pgfpathlineto{\pgfqpoint{1.188553in}{0.750956in}}%
\pgfpathlineto{\pgfqpoint{1.205258in}{0.758838in}}%
\pgfpathlineto{\pgfqpoint{1.213280in}{0.761131in}}%
\pgfpathlineto{\pgfqpoint{1.221279in}{0.764160in}}%
\pgfpathlineto{\pgfqpoint{1.204599in}{0.755962in}}%
\pgfpathlineto{\pgfqpoint{1.187987in}{0.748527in}}%
\pgfpathlineto{\pgfqpoint{1.179964in}{0.745784in}}%
\pgfpathlineto{\pgfqpoint{1.171915in}{0.743857in}}%
\pgfpathclose%
\pgfusepath{fill}%
\end{pgfscope}%
\begin{pgfscope}%
\pgfpathrectangle{\pgfqpoint{0.407592in}{0.222904in}}{\pgfqpoint{1.411163in}{1.411163in}}%
\pgfusepath{clip}%
\pgfsetbuttcap%
\pgfsetroundjoin%
\definecolor{currentfill}{rgb}{0.556863,0.654902,0.847059}%
\pgfsetfillcolor{currentfill}%
\pgfsetfillopacity{0.130000}%
\pgfsetlinewidth{0.000000pt}%
\definecolor{currentstroke}{rgb}{0.000000,0.000000,0.000000}%
\pgfsetstrokecolor{currentstroke}%
\pgfsetdash{}{0pt}%
\pgfpathmoveto{\pgfqpoint{1.073207in}{0.729855in}}%
\pgfpathlineto{\pgfqpoint{1.089685in}{0.731394in}}%
\pgfpathlineto{\pgfqpoint{1.106213in}{0.734261in}}%
\pgfpathlineto{\pgfqpoint{1.114414in}{0.732794in}}%
\pgfpathlineto{\pgfqpoint{1.122582in}{0.731981in}}%
\pgfpathlineto{\pgfqpoint{1.106092in}{0.728593in}}%
\pgfpathlineto{\pgfqpoint{1.089654in}{0.726594in}}%
\pgfpathlineto{\pgfqpoint{1.081448in}{0.727893in}}%
\pgfpathlineto{\pgfqpoint{1.073207in}{0.729855in}}%
\pgfpathclose%
\pgfusepath{fill}%
\end{pgfscope}%
\begin{pgfscope}%
\pgfpathrectangle{\pgfqpoint{0.407592in}{0.222904in}}{\pgfqpoint{1.411163in}{1.411163in}}%
\pgfusepath{clip}%
\pgfsetbuttcap%
\pgfsetroundjoin%
\definecolor{currentfill}{rgb}{0.556863,0.654902,0.847059}%
\pgfsetfillcolor{currentfill}%
\pgfsetfillopacity{0.130000}%
\pgfsetlinewidth{0.000000pt}%
\definecolor{currentstroke}{rgb}{0.000000,0.000000,0.000000}%
\pgfsetstrokecolor{currentstroke}%
\pgfsetdash{}{0pt}%
\pgfpathmoveto{\pgfqpoint{0.941254in}{0.766967in}}%
\pgfpathlineto{\pgfqpoint{0.957622in}{0.759583in}}%
\pgfpathlineto{\pgfqpoint{0.974014in}{0.753530in}}%
\pgfpathlineto{\pgfqpoint{0.982494in}{0.749264in}}%
\pgfpathlineto{\pgfqpoint{0.990936in}{0.745394in}}%
\pgfpathlineto{\pgfqpoint{0.974593in}{0.751603in}}%
\pgfpathlineto{\pgfqpoint{0.958272in}{0.759259in}}%
\pgfpathlineto{\pgfqpoint{0.949783in}{0.762910in}}%
\pgfpathlineto{\pgfqpoint{0.941254in}{0.766967in}}%
\pgfpathclose%
\pgfusepath{fill}%
\end{pgfscope}%
\begin{pgfscope}%
\pgfpathrectangle{\pgfqpoint{0.407592in}{0.222904in}}{\pgfqpoint{1.411163in}{1.411163in}}%
\pgfusepath{clip}%
\pgfsetbuttcap%
\pgfsetroundjoin%
\definecolor{currentfill}{rgb}{0.556863,0.654902,0.847059}%
\pgfsetfillcolor{currentfill}%
\pgfsetfillopacity{0.130000}%
\pgfsetlinewidth{0.000000pt}%
\definecolor{currentstroke}{rgb}{0.000000,0.000000,0.000000}%
\pgfsetstrokecolor{currentstroke}%
\pgfsetdash{}{0pt}%
\pgfpathmoveto{\pgfqpoint{0.742947in}{0.906110in}}%
\pgfpathlineto{\pgfqpoint{0.759273in}{0.890966in}}%
\pgfpathlineto{\pgfqpoint{0.775610in}{0.876361in}}%
\pgfpathlineto{\pgfqpoint{0.784486in}{0.872881in}}%
\pgfpathlineto{\pgfqpoint{0.793319in}{0.869728in}}%
\pgfpathlineto{\pgfqpoint{0.777027in}{0.884854in}}%
\pgfpathlineto{\pgfqpoint{0.760743in}{0.900507in}}%
\pgfpathlineto{\pgfqpoint{0.751866in}{0.903143in}}%
\pgfpathlineto{\pgfqpoint{0.742947in}{0.906110in}}%
\pgfpathclose%
\pgfusepath{fill}%
\end{pgfscope}%
\begin{pgfscope}%
\pgfpathrectangle{\pgfqpoint{0.407592in}{0.222904in}}{\pgfqpoint{1.411163in}{1.411163in}}%
\pgfusepath{clip}%
\pgfsetbuttcap%
\pgfsetroundjoin%
\definecolor{currentfill}{rgb}{0.556863,0.654902,0.847059}%
\pgfsetfillcolor{currentfill}%
\pgfsetfillopacity{0.130000}%
\pgfsetlinewidth{0.000000pt}%
\definecolor{currentstroke}{rgb}{0.000000,0.000000,0.000000}%
\pgfsetstrokecolor{currentstroke}%
\pgfsetdash{}{0pt}%
\pgfpathmoveto{\pgfqpoint{1.565912in}{1.034874in}}%
\pgfpathlineto{\pgfqpoint{1.583365in}{1.045513in}}%
\pgfpathlineto{\pgfqpoint{1.600895in}{1.056387in}}%
\pgfpathlineto{\pgfqpoint{1.608461in}{1.067102in}}%
\pgfpathlineto{\pgfqpoint{1.616016in}{1.078076in}}%
\pgfpathlineto{\pgfqpoint{1.598508in}{1.067556in}}%
\pgfpathlineto{\pgfqpoint{1.581076in}{1.057287in}}%
\pgfpathlineto{\pgfqpoint{1.573500in}{1.045948in}}%
\pgfpathlineto{\pgfqpoint{1.565912in}{1.034874in}}%
\pgfpathclose%
\pgfusepath{fill}%
\end{pgfscope}%
\begin{pgfscope}%
\pgfpathrectangle{\pgfqpoint{0.407592in}{0.222904in}}{\pgfqpoint{1.411163in}{1.411163in}}%
\pgfusepath{clip}%
\pgfsetbuttcap%
\pgfsetroundjoin%
\definecolor{currentfill}{rgb}{0.556863,0.654902,0.847059}%
\pgfsetfillcolor{currentfill}%
\pgfsetfillopacity{0.130000}%
\pgfsetlinewidth{0.000000pt}%
\definecolor{currentstroke}{rgb}{0.000000,0.000000,0.000000}%
\pgfsetstrokecolor{currentstroke}%
\pgfsetdash{}{0pt}%
\pgfpathmoveto{\pgfqpoint{1.221279in}{0.764160in}}%
\pgfpathlineto{\pgfqpoint{1.238027in}{0.773003in}}%
\pgfpathlineto{\pgfqpoint{1.254845in}{0.782394in}}%
\pgfpathlineto{\pgfqpoint{1.262798in}{0.786346in}}%
\pgfpathlineto{\pgfqpoint{1.270730in}{0.790984in}}%
\pgfpathlineto{\pgfqpoint{1.253933in}{0.781477in}}%
\pgfpathlineto{\pgfqpoint{1.237207in}{0.772505in}}%
\pgfpathlineto{\pgfqpoint{1.229254in}{0.767947in}}%
\pgfpathlineto{\pgfqpoint{1.221279in}{0.764160in}}%
\pgfpathclose%
\pgfusepath{fill}%
\end{pgfscope}%
\begin{pgfscope}%
\pgfpathrectangle{\pgfqpoint{0.407592in}{0.222904in}}{\pgfqpoint{1.411163in}{1.411163in}}%
\pgfusepath{clip}%
\pgfsetbuttcap%
\pgfsetroundjoin%
\definecolor{currentfill}{rgb}{0.556863,0.654902,0.847059}%
\pgfsetfillcolor{currentfill}%
\pgfsetfillopacity{0.130000}%
\pgfsetlinewidth{0.000000pt}%
\definecolor{currentstroke}{rgb}{0.000000,0.000000,0.000000}%
\pgfsetstrokecolor{currentstroke}%
\pgfsetdash{}{0pt}%
\pgfpathmoveto{\pgfqpoint{0.858593in}{0.815816in}}%
\pgfpathlineto{\pgfqpoint{0.874940in}{0.804363in}}%
\pgfpathlineto{\pgfqpoint{0.891302in}{0.793879in}}%
\pgfpathlineto{\pgfqpoint{0.899959in}{0.789507in}}%
\pgfpathlineto{\pgfqpoint{0.908575in}{0.785475in}}%
\pgfpathlineto{\pgfqpoint{0.892259in}{0.796432in}}%
\pgfpathlineto{\pgfqpoint{0.875957in}{0.808400in}}%
\pgfpathlineto{\pgfqpoint{0.867296in}{0.811929in}}%
\pgfpathlineto{\pgfqpoint{0.858593in}{0.815816in}}%
\pgfpathclose%
\pgfusepath{fill}%
\end{pgfscope}%
\begin{pgfscope}%
\pgfpathrectangle{\pgfqpoint{0.407592in}{0.222904in}}{\pgfqpoint{1.411163in}{1.411163in}}%
\pgfusepath{clip}%
\pgfsetbuttcap%
\pgfsetroundjoin%
\definecolor{currentfill}{rgb}{0.556863,0.654902,0.847059}%
\pgfsetfillcolor{currentfill}%
\pgfsetfillopacity{0.130000}%
\pgfsetlinewidth{0.000000pt}%
\definecolor{currentstroke}{rgb}{0.000000,0.000000,0.000000}%
\pgfsetstrokecolor{currentstroke}%
\pgfsetdash{}{0pt}%
\pgfpathmoveto{\pgfqpoint{0.626947in}{1.009861in}}%
\pgfpathlineto{\pgfqpoint{0.643264in}{0.992828in}}%
\pgfpathlineto{\pgfqpoint{0.659589in}{0.976103in}}%
\pgfpathlineto{\pgfqpoint{0.668681in}{0.973451in}}%
\pgfpathlineto{\pgfqpoint{0.677729in}{0.971095in}}%
\pgfpathlineto{\pgfqpoint{0.661447in}{0.988235in}}%
\pgfpathlineto{\pgfqpoint{0.645172in}{1.005662in}}%
\pgfpathlineto{\pgfqpoint{0.636082in}{1.007614in}}%
\pgfpathlineto{\pgfqpoint{0.626947in}{1.009861in}}%
\pgfpathclose%
\pgfusepath{fill}%
\end{pgfscope}%
\begin{pgfscope}%
\pgfpathrectangle{\pgfqpoint{0.407592in}{0.222904in}}{\pgfqpoint{1.411163in}{1.411163in}}%
\pgfusepath{clip}%
\pgfsetbuttcap%
\pgfsetroundjoin%
\definecolor{currentfill}{rgb}{0.556863,0.654902,0.847059}%
\pgfsetfillcolor{currentfill}%
\pgfsetfillopacity{0.130000}%
\pgfsetlinewidth{0.000000pt}%
\definecolor{currentstroke}{rgb}{0.000000,0.000000,0.000000}%
\pgfsetstrokecolor{currentstroke}%
\pgfsetdash{}{0pt}%
\pgfpathmoveto{\pgfqpoint{1.385613in}{0.873006in}}%
\pgfpathlineto{\pgfqpoint{1.402716in}{0.883845in}}%
\pgfpathlineto{\pgfqpoint{1.419893in}{0.894947in}}%
\pgfpathlineto{\pgfqpoint{1.427648in}{0.903150in}}%
\pgfpathlineto{\pgfqpoint{1.435389in}{0.911799in}}%
\pgfpathlineto{\pgfqpoint{1.418231in}{0.900994in}}%
\pgfpathlineto{\pgfqpoint{1.401148in}{0.890475in}}%
\pgfpathlineto{\pgfqpoint{1.393388in}{0.881500in}}%
\pgfpathlineto{\pgfqpoint{1.385613in}{0.873006in}}%
\pgfpathclose%
\pgfusepath{fill}%
\end{pgfscope}%
\begin{pgfscope}%
\pgfpathrectangle{\pgfqpoint{0.407592in}{0.222904in}}{\pgfqpoint{1.411163in}{1.411163in}}%
\pgfusepath{clip}%
\pgfsetbuttcap%
\pgfsetroundjoin%
\definecolor{currentfill}{rgb}{0.556863,0.654902,0.847059}%
\pgfsetfillcolor{currentfill}%
\pgfsetfillopacity{0.130000}%
\pgfsetlinewidth{0.000000pt}%
\definecolor{currentstroke}{rgb}{0.000000,0.000000,0.000000}%
\pgfsetstrokecolor{currentstroke}%
\pgfsetdash{}{0pt}%
\pgfpathmoveto{\pgfqpoint{0.477981in}{1.157044in}}%
\pgfpathlineto{\pgfqpoint{0.494272in}{1.138639in}}%
\pgfpathlineto{\pgfqpoint{0.510571in}{1.120393in}}%
\pgfpathlineto{\pgfqpoint{0.519921in}{1.118620in}}%
\pgfpathlineto{\pgfqpoint{0.529227in}{1.117104in}}%
\pgfpathlineto{\pgfqpoint{0.512970in}{1.135632in}}%
\pgfpathlineto{\pgfqpoint{0.496721in}{1.154304in}}%
\pgfpathlineto{\pgfqpoint{0.487373in}{1.155548in}}%
\pgfpathlineto{\pgfqpoint{0.477981in}{1.157044in}}%
\pgfpathclose%
\pgfusepath{fill}%
\end{pgfscope}%
\begin{pgfscope}%
\pgfpathrectangle{\pgfqpoint{0.407592in}{0.222904in}}{\pgfqpoint{1.411163in}{1.411163in}}%
\pgfusepath{clip}%
\pgfsetbuttcap%
\pgfsetroundjoin%
\definecolor{currentfill}{rgb}{0.556863,0.654902,0.847059}%
\pgfsetfillcolor{currentfill}%
\pgfsetfillopacity{0.130000}%
\pgfsetlinewidth{0.000000pt}%
\definecolor{currentstroke}{rgb}{0.000000,0.000000,0.000000}%
\pgfsetstrokecolor{currentstroke}%
\pgfsetdash{}{0pt}%
\pgfpathmoveto{\pgfqpoint{1.023708in}{0.737478in}}%
\pgfpathlineto{\pgfqpoint{1.040144in}{0.735765in}}%
\pgfpathlineto{\pgfqpoint{1.056620in}{0.735487in}}%
\pgfpathlineto{\pgfqpoint{1.064932in}{0.732409in}}%
\pgfpathlineto{\pgfqpoint{1.073207in}{0.729855in}}%
\pgfpathlineto{\pgfqpoint{1.056775in}{0.729769in}}%
\pgfpathlineto{\pgfqpoint{1.040384in}{0.731236in}}%
\pgfpathlineto{\pgfqpoint{1.032064in}{0.734099in}}%
\pgfpathlineto{\pgfqpoint{1.023708in}{0.737478in}}%
\pgfpathclose%
\pgfusepath{fill}%
\end{pgfscope}%
\begin{pgfscope}%
\pgfpathrectangle{\pgfqpoint{0.407592in}{0.222904in}}{\pgfqpoint{1.411163in}{1.411163in}}%
\pgfusepath{clip}%
\pgfsetbuttcap%
\pgfsetroundjoin%
\definecolor{currentfill}{rgb}{0.556863,0.654902,0.847059}%
\pgfsetfillcolor{currentfill}%
\pgfsetfillopacity{0.130000}%
\pgfsetlinewidth{0.000000pt}%
\definecolor{currentstroke}{rgb}{0.000000,0.000000,0.000000}%
\pgfsetstrokecolor{currentstroke}%
\pgfsetdash{}{0pt}%
\pgfpathmoveto{\pgfqpoint{1.746506in}{1.212824in}}%
\pgfpathlineto{\pgfqpoint{1.764311in}{1.223157in}}%
\pgfpathlineto{\pgfqpoint{1.782194in}{1.233709in}}%
\pgfpathlineto{\pgfqpoint{1.789583in}{1.245560in}}%
\pgfpathlineto{\pgfqpoint{1.771711in}{1.235160in}}%
\pgfpathlineto{\pgfqpoint{1.753918in}{1.224984in}}%
\pgfusepath{fill}%
\end{pgfscope}%
\begin{pgfscope}%
\pgfpathrectangle{\pgfqpoint{0.407592in}{0.222904in}}{\pgfqpoint{1.411163in}{1.411163in}}%
\pgfusepath{clip}%
\pgfsetbuttcap%
\pgfsetroundjoin%
\definecolor{currentfill}{rgb}{0.556863,0.654902,0.847059}%
\pgfsetfillcolor{currentfill}%
\pgfsetfillopacity{0.130000}%
\pgfsetlinewidth{0.000000pt}%
\definecolor{currentstroke}{rgb}{0.000000,0.000000,0.000000}%
\pgfsetstrokecolor{currentstroke}%
\pgfsetdash{}{0pt}%
\pgfpathmoveto{\pgfqpoint{1.681284in}{1.144691in}}%
\pgfpathlineto{\pgfqpoint{1.698979in}{1.155221in}}%
\pgfpathlineto{\pgfqpoint{1.716752in}{1.165971in}}%
\pgfpathlineto{\pgfqpoint{1.724207in}{1.177400in}}%
\pgfpathlineto{\pgfqpoint{1.731650in}{1.189024in}}%
\pgfpathlineto{\pgfqpoint{1.713901in}{1.178597in}}%
\pgfpathlineto{\pgfqpoint{1.696228in}{1.168401in}}%
\pgfpathlineto{\pgfqpoint{1.688761in}{1.156448in}}%
\pgfpathlineto{\pgfqpoint{1.681284in}{1.144691in}}%
\pgfpathclose%
\pgfusepath{fill}%
\end{pgfscope}%
\begin{pgfscope}%
\pgfpathrectangle{\pgfqpoint{0.407592in}{0.222904in}}{\pgfqpoint{1.411163in}{1.411163in}}%
\pgfusepath{clip}%
\pgfsetbuttcap%
\pgfsetroundjoin%
\definecolor{currentfill}{rgb}{0.556863,0.654902,0.847059}%
\pgfsetfillcolor{currentfill}%
\pgfsetfillopacity{0.130000}%
\pgfsetlinewidth{0.000000pt}%
\definecolor{currentstroke}{rgb}{0.000000,0.000000,0.000000}%
\pgfsetstrokecolor{currentstroke}%
\pgfsetdash{}{0pt}%
\pgfpathmoveto{\pgfqpoint{1.500680in}{0.971504in}}%
\pgfpathlineto{\pgfqpoint{1.518023in}{0.982381in}}%
\pgfpathlineto{\pgfqpoint{1.535443in}{0.993491in}}%
\pgfpathlineto{\pgfqpoint{1.543078in}{1.003374in}}%
\pgfpathlineto{\pgfqpoint{1.550701in}{1.013574in}}%
\pgfpathlineto{\pgfqpoint{1.533302in}{1.002815in}}%
\pgfpathlineto{\pgfqpoint{1.515980in}{0.992308in}}%
\pgfpathlineto{\pgfqpoint{1.508336in}{0.981741in}}%
\pgfpathlineto{\pgfqpoint{1.500680in}{0.971504in}}%
\pgfpathclose%
\pgfusepath{fill}%
\end{pgfscope}%
\begin{pgfscope}%
\pgfpathrectangle{\pgfqpoint{0.407592in}{0.222904in}}{\pgfqpoint{1.411163in}{1.411163in}}%
\pgfusepath{clip}%
\pgfsetbuttcap%
\pgfsetroundjoin%
\definecolor{currentfill}{rgb}{0.556863,0.654902,0.847059}%
\pgfsetfillcolor{currentfill}%
\pgfsetfillopacity{0.130000}%
\pgfsetlinewidth{0.000000pt}%
\definecolor{currentstroke}{rgb}{0.000000,0.000000,0.000000}%
\pgfsetstrokecolor{currentstroke}%
\pgfsetdash{}{0pt}%
\pgfpathmoveto{\pgfqpoint{1.270730in}{0.790984in}}%
\pgfpathlineto{\pgfqpoint{1.287598in}{0.800949in}}%
\pgfpathlineto{\pgfqpoint{1.304539in}{0.811307in}}%
\pgfpathlineto{\pgfqpoint{1.312430in}{0.816626in}}%
\pgfpathlineto{\pgfqpoint{1.320303in}{0.822560in}}%
\pgfpathlineto{\pgfqpoint{1.303383in}{0.812248in}}%
\pgfpathlineto{\pgfqpoint{1.286535in}{0.802331in}}%
\pgfpathlineto{\pgfqpoint{1.278642in}{0.796313in}}%
\pgfpathlineto{\pgfqpoint{1.270730in}{0.790984in}}%
\pgfpathclose%
\pgfusepath{fill}%
\end{pgfscope}%
\begin{pgfscope}%
\pgfpathrectangle{\pgfqpoint{0.407592in}{0.222904in}}{\pgfqpoint{1.411163in}{1.411163in}}%
\pgfusepath{clip}%
\pgfsetbuttcap%
\pgfsetroundjoin%
\definecolor{currentfill}{rgb}{0.556863,0.654902,0.847059}%
\pgfsetfillcolor{currentfill}%
\pgfsetfillopacity{0.130000}%
\pgfsetlinewidth{0.000000pt}%
\definecolor{currentstroke}{rgb}{0.000000,0.000000,0.000000}%
\pgfsetstrokecolor{currentstroke}%
\pgfsetdash{}{0pt}%
\pgfpathmoveto{\pgfqpoint{0.775610in}{0.876361in}}%
\pgfpathlineto{\pgfqpoint{0.791956in}{0.862353in}}%
\pgfpathlineto{\pgfqpoint{0.808313in}{0.849002in}}%
\pgfpathlineto{\pgfqpoint{0.817145in}{0.844999in}}%
\pgfpathlineto{\pgfqpoint{0.825934in}{0.841309in}}%
\pgfpathlineto{\pgfqpoint{0.809622in}{0.855190in}}%
\pgfpathlineto{\pgfqpoint{0.793319in}{0.869728in}}%
\pgfpathlineto{\pgfqpoint{0.784486in}{0.872881in}}%
\pgfpathlineto{\pgfqpoint{0.775610in}{0.876361in}}%
\pgfpathclose%
\pgfusepath{fill}%
\end{pgfscope}%
\begin{pgfscope}%
\pgfpathrectangle{\pgfqpoint{0.407592in}{0.222904in}}{\pgfqpoint{1.411163in}{1.411163in}}%
\pgfusepath{clip}%
\pgfsetbuttcap%
\pgfsetroundjoin%
\definecolor{currentfill}{rgb}{0.556863,0.654902,0.847059}%
\pgfsetfillcolor{currentfill}%
\pgfsetfillopacity{0.130000}%
\pgfsetlinewidth{0.000000pt}%
\definecolor{currentstroke}{rgb}{0.000000,0.000000,0.000000}%
\pgfsetstrokecolor{currentstroke}%
\pgfsetdash{}{0pt}%
\pgfpathmoveto{\pgfqpoint{0.510571in}{1.120393in}}%
\pgfpathlineto{\pgfqpoint{0.526879in}{1.102321in}}%
\pgfpathlineto{\pgfqpoint{0.543195in}{1.084439in}}%
\pgfpathlineto{\pgfqpoint{0.552503in}{1.082357in}}%
\pgfpathlineto{\pgfqpoint{0.561766in}{1.080538in}}%
\pgfpathlineto{\pgfqpoint{0.545492in}{1.098734in}}%
\pgfpathlineto{\pgfqpoint{0.529227in}{1.117104in}}%
\pgfpathlineto{\pgfqpoint{0.519921in}{1.118620in}}%
\pgfpathlineto{\pgfqpoint{0.510571in}{1.120393in}}%
\pgfpathclose%
\pgfusepath{fill}%
\end{pgfscope}%
\begin{pgfscope}%
\pgfpathrectangle{\pgfqpoint{0.407592in}{0.222904in}}{\pgfqpoint{1.411163in}{1.411163in}}%
\pgfusepath{clip}%
\pgfsetbuttcap%
\pgfsetroundjoin%
\definecolor{currentfill}{rgb}{0.556863,0.654902,0.847059}%
\pgfsetfillcolor{currentfill}%
\pgfsetfillopacity{0.130000}%
\pgfsetlinewidth{0.000000pt}%
\definecolor{currentstroke}{rgb}{0.000000,0.000000,0.000000}%
\pgfsetstrokecolor{currentstroke}%
\pgfsetdash{}{0pt}%
\pgfpathmoveto{\pgfqpoint{0.659589in}{0.976103in}}%
\pgfpathlineto{\pgfqpoint{0.675923in}{0.959718in}}%
\pgfpathlineto{\pgfqpoint{0.692266in}{0.943706in}}%
\pgfpathlineto{\pgfqpoint{0.701315in}{0.940610in}}%
\pgfpathlineto{\pgfqpoint{0.710320in}{0.937811in}}%
\pgfpathlineto{\pgfqpoint{0.694020in}{0.954275in}}%
\pgfpathlineto{\pgfqpoint{0.677729in}{0.971095in}}%
\pgfpathlineto{\pgfqpoint{0.668681in}{0.973451in}}%
\pgfpathlineto{\pgfqpoint{0.659589in}{0.976103in}}%
\pgfpathclose%
\pgfusepath{fill}%
\end{pgfscope}%
\begin{pgfscope}%
\pgfpathrectangle{\pgfqpoint{0.407592in}{0.222904in}}{\pgfqpoint{1.411163in}{1.411163in}}%
\pgfusepath{clip}%
\pgfsetbuttcap%
\pgfsetroundjoin%
\definecolor{currentfill}{rgb}{0.556863,0.654902,0.847059}%
\pgfsetfillcolor{currentfill}%
\pgfsetfillopacity{0.130000}%
\pgfsetlinewidth{0.000000pt}%
\definecolor{currentstroke}{rgb}{0.000000,0.000000,0.000000}%
\pgfsetstrokecolor{currentstroke}%
\pgfsetdash{}{0pt}%
\pgfpathmoveto{\pgfqpoint{1.616016in}{1.078076in}}%
\pgfpathlineto{\pgfqpoint{1.633602in}{1.088830in}}%
\pgfpathlineto{\pgfqpoint{1.651264in}{1.099804in}}%
\pgfpathlineto{\pgfqpoint{1.658786in}{1.110686in}}%
\pgfpathlineto{\pgfqpoint{1.666296in}{1.121801in}}%
\pgfpathlineto{\pgfqpoint{1.648656in}{1.111162in}}%
\pgfpathlineto{\pgfqpoint{1.631093in}{1.100757in}}%
\pgfpathlineto{\pgfqpoint{1.623560in}{1.089298in}}%
\pgfpathlineto{\pgfqpoint{1.616016in}{1.078076in}}%
\pgfpathclose%
\pgfusepath{fill}%
\end{pgfscope}%
\begin{pgfscope}%
\pgfpathrectangle{\pgfqpoint{0.407592in}{0.222904in}}{\pgfqpoint{1.411163in}{1.411163in}}%
\pgfusepath{clip}%
\pgfsetbuttcap%
\pgfsetroundjoin%
\definecolor{currentfill}{rgb}{0.556863,0.654902,0.847059}%
\pgfsetfillcolor{currentfill}%
\pgfsetfillopacity{0.130000}%
\pgfsetlinewidth{0.000000pt}%
\definecolor{currentstroke}{rgb}{0.000000,0.000000,0.000000}%
\pgfsetstrokecolor{currentstroke}%
\pgfsetdash{}{0pt}%
\pgfpathmoveto{\pgfqpoint{0.891302in}{0.793879in}}%
\pgfpathlineto{\pgfqpoint{0.907681in}{0.784441in}}%
\pgfpathlineto{\pgfqpoint{0.924077in}{0.776122in}}%
\pgfpathlineto{\pgfqpoint{0.932685in}{0.771385in}}%
\pgfpathlineto{\pgfqpoint{0.941254in}{0.766967in}}%
\pgfpathlineto{\pgfqpoint{0.924906in}{0.775624in}}%
\pgfpathlineto{\pgfqpoint{0.908575in}{0.785475in}}%
\pgfpathlineto{\pgfqpoint{0.899959in}{0.789507in}}%
\pgfpathlineto{\pgfqpoint{0.891302in}{0.793879in}}%
\pgfpathclose%
\pgfusepath{fill}%
\end{pgfscope}%
\begin{pgfscope}%
\pgfpathrectangle{\pgfqpoint{0.407592in}{0.222904in}}{\pgfqpoint{1.411163in}{1.411163in}}%
\pgfusepath{clip}%
\pgfsetbuttcap%
\pgfsetroundjoin%
\definecolor{currentfill}{rgb}{0.556863,0.654902,0.847059}%
\pgfsetfillcolor{currentfill}%
\pgfsetfillopacity{0.130000}%
\pgfsetlinewidth{0.000000pt}%
\definecolor{currentstroke}{rgb}{0.000000,0.000000,0.000000}%
\pgfsetstrokecolor{currentstroke}%
\pgfsetdash{}{0pt}%
\pgfpathmoveto{\pgfqpoint{1.435389in}{0.911799in}}%
\pgfpathlineto{\pgfqpoint{1.452621in}{0.922862in}}%
\pgfpathlineto{\pgfqpoint{1.469930in}{0.934159in}}%
\pgfpathlineto{\pgfqpoint{1.477637in}{0.942926in}}%
\pgfpathlineto{\pgfqpoint{1.485330in}{0.952082in}}%
\pgfpathlineto{\pgfqpoint{1.468042in}{0.941100in}}%
\pgfpathlineto{\pgfqpoint{1.450830in}{0.930373in}}%
\pgfpathlineto{\pgfqpoint{1.443116in}{0.920879in}}%
\pgfpathlineto{\pgfqpoint{1.435389in}{0.911799in}}%
\pgfpathclose%
\pgfusepath{fill}%
\end{pgfscope}%
\begin{pgfscope}%
\pgfpathrectangle{\pgfqpoint{0.407592in}{0.222904in}}{\pgfqpoint{1.411163in}{1.411163in}}%
\pgfusepath{clip}%
\pgfsetbuttcap%
\pgfsetroundjoin%
\definecolor{currentfill}{rgb}{0.556863,0.654902,0.847059}%
\pgfsetfillcolor{currentfill}%
\pgfsetfillopacity{0.130000}%
\pgfsetlinewidth{0.000000pt}%
\definecolor{currentstroke}{rgb}{0.000000,0.000000,0.000000}%
\pgfsetstrokecolor{currentstroke}%
\pgfsetdash{}{0pt}%
\pgfpathmoveto{\pgfqpoint{1.320303in}{0.822560in}}%
\pgfpathlineto{\pgfqpoint{1.337297in}{0.833216in}}%
\pgfpathlineto{\pgfqpoint{1.354365in}{0.844172in}}%
\pgfpathlineto{\pgfqpoint{1.362201in}{0.850582in}}%
\pgfpathlineto{\pgfqpoint{1.370020in}{0.857532in}}%
\pgfpathlineto{\pgfqpoint{1.352972in}{0.846730in}}%
\pgfpathlineto{\pgfqpoint{1.335998in}{0.836242in}}%
\pgfpathlineto{\pgfqpoint{1.328159in}{0.829102in}}%
\pgfpathlineto{\pgfqpoint{1.320303in}{0.822560in}}%
\pgfpathclose%
\pgfusepath{fill}%
\end{pgfscope}%
\begin{pgfscope}%
\pgfpathrectangle{\pgfqpoint{0.407592in}{0.222904in}}{\pgfqpoint{1.411163in}{1.411163in}}%
\pgfusepath{clip}%
\pgfsetbuttcap%
\pgfsetroundjoin%
\definecolor{currentfill}{rgb}{0.556863,0.654902,0.847059}%
\pgfsetfillcolor{currentfill}%
\pgfsetfillopacity{0.130000}%
\pgfsetlinewidth{0.000000pt}%
\definecolor{currentstroke}{rgb}{0.000000,0.000000,0.000000}%
\pgfsetstrokecolor{currentstroke}%
\pgfsetdash{}{0pt}%
\pgfpathmoveto{\pgfqpoint{0.974014in}{0.753530in}}%
\pgfpathlineto{\pgfqpoint{0.990432in}{0.748845in}}%
\pgfpathlineto{\pgfqpoint{1.006881in}{0.745540in}}%
\pgfpathlineto{\pgfqpoint{1.015313in}{0.741310in}}%
\pgfpathlineto{\pgfqpoint{1.023708in}{0.737478in}}%
\pgfpathlineto{\pgfqpoint{1.007307in}{0.740681in}}%
\pgfpathlineto{\pgfqpoint{0.990936in}{0.745394in}}%
\pgfpathlineto{\pgfqpoint{0.982494in}{0.749264in}}%
\pgfpathlineto{\pgfqpoint{0.974014in}{0.753530in}}%
\pgfpathclose%
\pgfusepath{fill}%
\end{pgfscope}%
\begin{pgfscope}%
\pgfpathrectangle{\pgfqpoint{0.407592in}{0.222904in}}{\pgfqpoint{1.411163in}{1.411163in}}%
\pgfusepath{clip}%
\pgfsetbuttcap%
\pgfsetroundjoin%
\definecolor{currentfill}{rgb}{0.556863,0.654902,0.847059}%
\pgfsetfillcolor{currentfill}%
\pgfsetfillopacity{0.130000}%
\pgfsetlinewidth{0.000000pt}%
\definecolor{currentstroke}{rgb}{0.000000,0.000000,0.000000}%
\pgfsetstrokecolor{currentstroke}%
\pgfsetdash{}{0pt}%
\pgfpathmoveto{\pgfqpoint{0.543195in}{1.084439in}}%
\pgfpathlineto{\pgfqpoint{0.559520in}{1.066764in}}%
\pgfpathlineto{\pgfqpoint{0.575854in}{1.049316in}}%
\pgfpathlineto{\pgfqpoint{0.585119in}{1.046893in}}%
\pgfpathlineto{\pgfqpoint{0.594340in}{1.044738in}}%
\pgfpathlineto{\pgfqpoint{0.578049in}{1.062533in}}%
\pgfpathlineto{\pgfqpoint{0.561766in}{1.080538in}}%
\pgfpathlineto{\pgfqpoint{0.552503in}{1.082357in}}%
\pgfpathlineto{\pgfqpoint{0.543195in}{1.084439in}}%
\pgfpathclose%
\pgfusepath{fill}%
\end{pgfscope}%
\begin{pgfscope}%
\pgfpathrectangle{\pgfqpoint{0.407592in}{0.222904in}}{\pgfqpoint{1.411163in}{1.411163in}}%
\pgfusepath{clip}%
\pgfsetbuttcap%
\pgfsetroundjoin%
\definecolor{currentfill}{rgb}{0.556863,0.654902,0.847059}%
\pgfsetfillcolor{currentfill}%
\pgfsetfillopacity{0.130000}%
\pgfsetlinewidth{0.000000pt}%
\definecolor{currentstroke}{rgb}{0.000000,0.000000,0.000000}%
\pgfsetstrokecolor{currentstroke}%
\pgfsetdash{}{0pt}%
\pgfpathmoveto{\pgfqpoint{1.550701in}{1.013574in}}%
\pgfpathlineto{\pgfqpoint{1.568176in}{1.024568in}}%
\pgfpathlineto{\pgfqpoint{1.585728in}{1.035780in}}%
\pgfpathlineto{\pgfqpoint{1.593317in}{1.045943in}}%
\pgfpathlineto{\pgfqpoint{1.600895in}{1.056387in}}%
\pgfpathlineto{\pgfqpoint{1.583365in}{1.045513in}}%
\pgfpathlineto{\pgfqpoint{1.565912in}{1.034874in}}%
\pgfpathlineto{\pgfqpoint{1.558312in}{1.024079in}}%
\pgfpathlineto{\pgfqpoint{1.550701in}{1.013574in}}%
\pgfpathclose%
\pgfusepath{fill}%
\end{pgfscope}%
\begin{pgfscope}%
\pgfpathrectangle{\pgfqpoint{0.407592in}{0.222904in}}{\pgfqpoint{1.411163in}{1.411163in}}%
\pgfusepath{clip}%
\pgfsetbuttcap%
\pgfsetroundjoin%
\definecolor{currentfill}{rgb}{0.556863,0.654902,0.847059}%
\pgfsetfillcolor{currentfill}%
\pgfsetfillopacity{0.130000}%
\pgfsetlinewidth{0.000000pt}%
\definecolor{currentstroke}{rgb}{0.000000,0.000000,0.000000}%
\pgfsetstrokecolor{currentstroke}%
\pgfsetdash{}{0pt}%
\pgfpathmoveto{\pgfqpoint{1.731650in}{1.189024in}}%
\pgfpathlineto{\pgfqpoint{1.749479in}{1.199672in}}%
\pgfpathlineto{\pgfqpoint{1.767385in}{1.210528in}}%
\pgfpathlineto{\pgfqpoint{1.774795in}{1.222030in}}%
\pgfpathlineto{\pgfqpoint{1.782194in}{1.233709in}}%
\pgfpathlineto{\pgfqpoint{1.764311in}{1.223157in}}%
\pgfpathlineto{\pgfqpoint{1.746506in}{1.212824in}}%
\pgfpathlineto{\pgfqpoint{1.739084in}{1.200835in}}%
\pgfpathlineto{\pgfqpoint{1.731650in}{1.189024in}}%
\pgfpathclose%
\pgfusepath{fill}%
\end{pgfscope}%
\begin{pgfscope}%
\pgfpathrectangle{\pgfqpoint{0.407592in}{0.222904in}}{\pgfqpoint{1.411163in}{1.411163in}}%
\pgfusepath{clip}%
\pgfsetbuttcap%
\pgfsetroundjoin%
\definecolor{currentfill}{rgb}{0.556863,0.654902,0.847059}%
\pgfsetfillcolor{currentfill}%
\pgfsetfillopacity{0.130000}%
\pgfsetlinewidth{0.000000pt}%
\definecolor{currentstroke}{rgb}{0.000000,0.000000,0.000000}%
\pgfsetstrokecolor{currentstroke}%
\pgfsetdash{}{0pt}%
\pgfpathmoveto{\pgfqpoint{0.808313in}{0.849002in}}%
\pgfpathlineto{\pgfqpoint{0.824682in}{0.836375in}}%
\pgfpathlineto{\pgfqpoint{0.841063in}{0.824538in}}%
\pgfpathlineto{\pgfqpoint{0.849849in}{0.820029in}}%
\pgfpathlineto{\pgfqpoint{0.858593in}{0.815816in}}%
\pgfpathlineto{\pgfqpoint{0.842258in}{0.828159in}}%
\pgfpathlineto{\pgfqpoint{0.825934in}{0.841309in}}%
\pgfpathlineto{\pgfqpoint{0.817145in}{0.844999in}}%
\pgfpathlineto{\pgfqpoint{0.808313in}{0.849002in}}%
\pgfpathclose%
\pgfusepath{fill}%
\end{pgfscope}%
\begin{pgfscope}%
\pgfpathrectangle{\pgfqpoint{0.407592in}{0.222904in}}{\pgfqpoint{1.411163in}{1.411163in}}%
\pgfusepath{clip}%
\pgfsetbuttcap%
\pgfsetroundjoin%
\definecolor{currentfill}{rgb}{0.556863,0.654902,0.847059}%
\pgfsetfillcolor{currentfill}%
\pgfsetfillopacity{0.130000}%
\pgfsetlinewidth{0.000000pt}%
\definecolor{currentstroke}{rgb}{0.000000,0.000000,0.000000}%
\pgfsetstrokecolor{currentstroke}%
\pgfsetdash{}{0pt}%
\pgfpathmoveto{\pgfqpoint{0.692266in}{0.943706in}}%
\pgfpathlineto{\pgfqpoint{0.708618in}{0.928106in}}%
\pgfpathlineto{\pgfqpoint{0.724979in}{0.912959in}}%
\pgfpathlineto{\pgfqpoint{0.733984in}{0.909389in}}%
\pgfpathlineto{\pgfqpoint{0.742947in}{0.906110in}}%
\pgfpathlineto{\pgfqpoint{0.726629in}{0.921741in}}%
\pgfpathlineto{\pgfqpoint{0.710320in}{0.937811in}}%
\pgfpathlineto{\pgfqpoint{0.701315in}{0.940610in}}%
\pgfpathlineto{\pgfqpoint{0.692266in}{0.943706in}}%
\pgfpathclose%
\pgfusepath{fill}%
\end{pgfscope}%
\begin{pgfscope}%
\pgfpathrectangle{\pgfqpoint{0.407592in}{0.222904in}}{\pgfqpoint{1.411163in}{1.411163in}}%
\pgfusepath{clip}%
\pgfsetbuttcap%
\pgfsetroundjoin%
\definecolor{currentfill}{rgb}{0.556863,0.654902,0.847059}%
\pgfsetfillcolor{currentfill}%
\pgfsetfillopacity{0.130000}%
\pgfsetlinewidth{0.000000pt}%
\definecolor{currentstroke}{rgb}{0.000000,0.000000,0.000000}%
\pgfsetstrokecolor{currentstroke}%
\pgfsetdash{}{0pt}%
\pgfpathmoveto{\pgfqpoint{1.155734in}{0.742290in}}%
\pgfpathlineto{\pgfqpoint{1.172402in}{0.748915in}}%
\pgfpathlineto{\pgfqpoint{1.189135in}{0.756344in}}%
\pgfpathlineto{\pgfqpoint{1.197210in}{0.757253in}}%
\pgfpathlineto{\pgfqpoint{1.205258in}{0.758838in}}%
\pgfpathlineto{\pgfqpoint{1.188553in}{0.750956in}}%
\pgfpathlineto{\pgfqpoint{1.171915in}{0.743857in}}%
\pgfpathlineto{\pgfqpoint{1.163838in}{0.742708in}}%
\pgfpathlineto{\pgfqpoint{1.155734in}{0.742290in}}%
\pgfpathclose%
\pgfusepath{fill}%
\end{pgfscope}%
\begin{pgfscope}%
\pgfpathrectangle{\pgfqpoint{0.407592in}{0.222904in}}{\pgfqpoint{1.411163in}{1.411163in}}%
\pgfusepath{clip}%
\pgfsetbuttcap%
\pgfsetroundjoin%
\definecolor{currentfill}{rgb}{0.556863,0.654902,0.847059}%
\pgfsetfillcolor{currentfill}%
\pgfsetfillopacity{0.130000}%
\pgfsetlinewidth{0.000000pt}%
\definecolor{currentstroke}{rgb}{0.000000,0.000000,0.000000}%
\pgfsetstrokecolor{currentstroke}%
\pgfsetdash{}{0pt}%
\pgfpathmoveto{\pgfqpoint{1.106213in}{0.734261in}}%
\pgfpathlineto{\pgfqpoint{1.122795in}{0.738324in}}%
\pgfpathlineto{\pgfqpoint{1.139435in}{0.743451in}}%
\pgfpathlineto{\pgfqpoint{1.147600in}{0.742555in}}%
\pgfpathlineto{\pgfqpoint{1.155734in}{0.742290in}}%
\pgfpathlineto{\pgfqpoint{1.139128in}{0.736598in}}%
\pgfpathlineto{\pgfqpoint{1.122582in}{0.731981in}}%
\pgfpathlineto{\pgfqpoint{1.114414in}{0.732794in}}%
\pgfpathlineto{\pgfqpoint{1.106213in}{0.734261in}}%
\pgfpathclose%
\pgfusepath{fill}%
\end{pgfscope}%
\begin{pgfscope}%
\pgfpathrectangle{\pgfqpoint{0.407592in}{0.222904in}}{\pgfqpoint{1.411163in}{1.411163in}}%
\pgfusepath{clip}%
\pgfsetbuttcap%
\pgfsetroundjoin%
\definecolor{currentfill}{rgb}{0.556863,0.654902,0.847059}%
\pgfsetfillcolor{currentfill}%
\pgfsetfillopacity{0.130000}%
\pgfsetlinewidth{0.000000pt}%
\definecolor{currentstroke}{rgb}{0.000000,0.000000,0.000000}%
\pgfsetstrokecolor{currentstroke}%
\pgfsetdash{}{0pt}%
\pgfpathmoveto{\pgfqpoint{1.666296in}{1.121801in}}%
\pgfpathlineto{\pgfqpoint{1.684014in}{1.132660in}}%
\pgfpathlineto{\pgfqpoint{1.701811in}{1.143728in}}%
\pgfpathlineto{\pgfqpoint{1.709287in}{1.154744in}}%
\pgfpathlineto{\pgfqpoint{1.716752in}{1.165971in}}%
\pgfpathlineto{\pgfqpoint{1.698979in}{1.155221in}}%
\pgfpathlineto{\pgfqpoint{1.681284in}{1.144691in}}%
\pgfpathlineto{\pgfqpoint{1.673795in}{1.133139in}}%
\pgfpathlineto{\pgfqpoint{1.666296in}{1.121801in}}%
\pgfpathclose%
\pgfusepath{fill}%
\end{pgfscope}%
\begin{pgfscope}%
\pgfpathrectangle{\pgfqpoint{0.407592in}{0.222904in}}{\pgfqpoint{1.411163in}{1.411163in}}%
\pgfusepath{clip}%
\pgfsetbuttcap%
\pgfsetroundjoin%
\definecolor{currentfill}{rgb}{0.556863,0.654902,0.847059}%
\pgfsetfillcolor{currentfill}%
\pgfsetfillopacity{0.130000}%
\pgfsetlinewidth{0.000000pt}%
\definecolor{currentstroke}{rgb}{0.000000,0.000000,0.000000}%
\pgfsetstrokecolor{currentstroke}%
\pgfsetdash{}{0pt}%
\pgfpathmoveto{\pgfqpoint{1.370020in}{0.857532in}}%
\pgfpathlineto{\pgfqpoint{1.387143in}{0.868610in}}%
\pgfpathlineto{\pgfqpoint{1.404341in}{0.879934in}}%
\pgfpathlineto{\pgfqpoint{1.412124in}{0.887205in}}%
\pgfpathlineto{\pgfqpoint{1.419893in}{0.894947in}}%
\pgfpathlineto{\pgfqpoint{1.402716in}{0.883845in}}%
\pgfpathlineto{\pgfqpoint{1.385613in}{0.873006in}}%
\pgfpathlineto{\pgfqpoint{1.377824in}{0.865011in}}%
\pgfpathlineto{\pgfqpoint{1.370020in}{0.857532in}}%
\pgfpathclose%
\pgfusepath{fill}%
\end{pgfscope}%
\begin{pgfscope}%
\pgfpathrectangle{\pgfqpoint{0.407592in}{0.222904in}}{\pgfqpoint{1.411163in}{1.411163in}}%
\pgfusepath{clip}%
\pgfsetbuttcap%
\pgfsetroundjoin%
\definecolor{currentfill}{rgb}{0.556863,0.654902,0.847059}%
\pgfsetfillcolor{currentfill}%
\pgfsetfillopacity{0.130000}%
\pgfsetlinewidth{0.000000pt}%
\definecolor{currentstroke}{rgb}{0.000000,0.000000,0.000000}%
\pgfsetstrokecolor{currentstroke}%
\pgfsetdash{}{0pt}%
\pgfpathmoveto{\pgfqpoint{0.575854in}{1.049316in}}%
\pgfpathlineto{\pgfqpoint{0.592196in}{1.032116in}}%
\pgfpathlineto{\pgfqpoint{0.608546in}{1.015186in}}%
\pgfpathlineto{\pgfqpoint{0.617769in}{1.012388in}}%
\pgfpathlineto{\pgfqpoint{0.626947in}{1.009861in}}%
\pgfpathlineto{\pgfqpoint{0.610639in}{1.027172in}}%
\pgfpathlineto{\pgfqpoint{0.594340in}{1.044738in}}%
\pgfpathlineto{\pgfqpoint{0.585119in}{1.046893in}}%
\pgfpathlineto{\pgfqpoint{0.575854in}{1.049316in}}%
\pgfpathclose%
\pgfusepath{fill}%
\end{pgfscope}%
\begin{pgfscope}%
\pgfpathrectangle{\pgfqpoint{0.407592in}{0.222904in}}{\pgfqpoint{1.411163in}{1.411163in}}%
\pgfusepath{clip}%
\pgfsetbuttcap%
\pgfsetroundjoin%
\definecolor{currentfill}{rgb}{0.556863,0.654902,0.847059}%
\pgfsetfillcolor{currentfill}%
\pgfsetfillopacity{0.130000}%
\pgfsetlinewidth{0.000000pt}%
\definecolor{currentstroke}{rgb}{0.000000,0.000000,0.000000}%
\pgfsetstrokecolor{currentstroke}%
\pgfsetdash{}{0pt}%
\pgfpathmoveto{\pgfqpoint{1.485330in}{0.952082in}}%
\pgfpathlineto{\pgfqpoint{1.502695in}{0.963296in}}%
\pgfpathlineto{\pgfqpoint{1.520136in}{0.974725in}}%
\pgfpathlineto{\pgfqpoint{1.527795in}{0.983937in}}%
\pgfpathlineto{\pgfqpoint{1.535443in}{0.993491in}}%
\pgfpathlineto{\pgfqpoint{1.518023in}{0.982381in}}%
\pgfpathlineto{\pgfqpoint{1.500680in}{0.971504in}}%
\pgfpathlineto{\pgfqpoint{1.493011in}{0.961613in}}%
\pgfpathlineto{\pgfqpoint{1.485330in}{0.952082in}}%
\pgfpathclose%
\pgfusepath{fill}%
\end{pgfscope}%
\begin{pgfscope}%
\pgfpathrectangle{\pgfqpoint{0.407592in}{0.222904in}}{\pgfqpoint{1.411163in}{1.411163in}}%
\pgfusepath{clip}%
\pgfsetbuttcap%
\pgfsetroundjoin%
\definecolor{currentfill}{rgb}{0.556863,0.654902,0.847059}%
\pgfsetfillcolor{currentfill}%
\pgfsetfillopacity{0.130000}%
\pgfsetlinewidth{0.000000pt}%
\definecolor{currentstroke}{rgb}{0.000000,0.000000,0.000000}%
\pgfsetstrokecolor{currentstroke}%
\pgfsetdash{}{0pt}%
\pgfpathmoveto{\pgfqpoint{1.205258in}{0.758838in}}%
\pgfpathlineto{\pgfqpoint{1.222031in}{0.767388in}}%
\pgfpathlineto{\pgfqpoint{1.238873in}{0.776510in}}%
\pgfpathlineto{\pgfqpoint{1.246871in}{0.779121in}}%
\pgfpathlineto{\pgfqpoint{1.254845in}{0.782394in}}%
\pgfpathlineto{\pgfqpoint{1.238027in}{0.773003in}}%
\pgfpathlineto{\pgfqpoint{1.221279in}{0.764160in}}%
\pgfpathlineto{\pgfqpoint{1.213280in}{0.761131in}}%
\pgfpathlineto{\pgfqpoint{1.205258in}{0.758838in}}%
\pgfpathclose%
\pgfusepath{fill}%
\end{pgfscope}%
\begin{pgfscope}%
\pgfpathrectangle{\pgfqpoint{0.407592in}{0.222904in}}{\pgfqpoint{1.411163in}{1.411163in}}%
\pgfusepath{clip}%
\pgfsetbuttcap%
\pgfsetroundjoin%
\definecolor{currentfill}{rgb}{0.556863,0.654902,0.847059}%
\pgfsetfillcolor{currentfill}%
\pgfsetfillopacity{0.130000}%
\pgfsetlinewidth{0.000000pt}%
\definecolor{currentstroke}{rgb}{0.000000,0.000000,0.000000}%
\pgfsetstrokecolor{currentstroke}%
\pgfsetdash{}{0pt}%
\pgfpathmoveto{\pgfqpoint{1.056620in}{0.735487in}}%
\pgfpathlineto{\pgfqpoint{1.073141in}{0.736568in}}%
\pgfpathlineto{\pgfqpoint{1.089710in}{0.738908in}}%
\pgfpathlineto{\pgfqpoint{1.097979in}{0.736318in}}%
\pgfpathlineto{\pgfqpoint{1.106213in}{0.734261in}}%
\pgfpathlineto{\pgfqpoint{1.089685in}{0.731394in}}%
\pgfpathlineto{\pgfqpoint{1.073207in}{0.729855in}}%
\pgfpathlineto{\pgfqpoint{1.064932in}{0.732409in}}%
\pgfpathlineto{\pgfqpoint{1.056620in}{0.735487in}}%
\pgfpathclose%
\pgfusepath{fill}%
\end{pgfscope}%
\begin{pgfscope}%
\pgfpathrectangle{\pgfqpoint{0.407592in}{0.222904in}}{\pgfqpoint{1.411163in}{1.411163in}}%
\pgfusepath{clip}%
\pgfsetbuttcap%
\pgfsetroundjoin%
\definecolor{currentfill}{rgb}{0.556863,0.654902,0.847059}%
\pgfsetfillcolor{currentfill}%
\pgfsetfillopacity{0.130000}%
\pgfsetlinewidth{0.000000pt}%
\definecolor{currentstroke}{rgb}{0.000000,0.000000,0.000000}%
\pgfsetstrokecolor{currentstroke}%
\pgfsetdash{}{0pt}%
\pgfpathmoveto{\pgfqpoint{0.724979in}{0.912959in}}%
\pgfpathlineto{\pgfqpoint{0.741349in}{0.898310in}}%
\pgfpathlineto{\pgfqpoint{0.757730in}{0.884208in}}%
\pgfpathlineto{\pgfqpoint{0.766691in}{0.880144in}}%
\pgfpathlineto{\pgfqpoint{0.775610in}{0.876361in}}%
\pgfpathlineto{\pgfqpoint{0.759273in}{0.890966in}}%
\pgfpathlineto{\pgfqpoint{0.742947in}{0.906110in}}%
\pgfpathlineto{\pgfqpoint{0.733984in}{0.909389in}}%
\pgfpathlineto{\pgfqpoint{0.724979in}{0.912959in}}%
\pgfpathclose%
\pgfusepath{fill}%
\end{pgfscope}%
\begin{pgfscope}%
\pgfpathrectangle{\pgfqpoint{0.407592in}{0.222904in}}{\pgfqpoint{1.411163in}{1.411163in}}%
\pgfusepath{clip}%
\pgfsetbuttcap%
\pgfsetroundjoin%
\definecolor{currentfill}{rgb}{0.556863,0.654902,0.847059}%
\pgfsetfillcolor{currentfill}%
\pgfsetfillopacity{0.130000}%
\pgfsetlinewidth{0.000000pt}%
\definecolor{currentstroke}{rgb}{0.000000,0.000000,0.000000}%
\pgfsetstrokecolor{currentstroke}%
\pgfsetdash{}{0pt}%
\pgfpathmoveto{\pgfqpoint{0.924077in}{0.776122in}}%
\pgfpathlineto{\pgfqpoint{0.940494in}{0.768983in}}%
\pgfpathlineto{\pgfqpoint{0.956935in}{0.763069in}}%
\pgfpathlineto{\pgfqpoint{0.965494in}{0.758145in}}%
\pgfpathlineto{\pgfqpoint{0.974014in}{0.753530in}}%
\pgfpathlineto{\pgfqpoint{0.957622in}{0.759583in}}%
\pgfpathlineto{\pgfqpoint{0.941254in}{0.766967in}}%
\pgfpathlineto{\pgfqpoint{0.932685in}{0.771385in}}%
\pgfpathlineto{\pgfqpoint{0.924077in}{0.776122in}}%
\pgfpathclose%
\pgfusepath{fill}%
\end{pgfscope}%
\begin{pgfscope}%
\pgfpathrectangle{\pgfqpoint{0.407592in}{0.222904in}}{\pgfqpoint{1.411163in}{1.411163in}}%
\pgfusepath{clip}%
\pgfsetbuttcap%
\pgfsetroundjoin%
\definecolor{currentfill}{rgb}{0.556863,0.654902,0.847059}%
\pgfsetfillcolor{currentfill}%
\pgfsetfillopacity{0.130000}%
\pgfsetlinewidth{0.000000pt}%
\definecolor{currentstroke}{rgb}{0.000000,0.000000,0.000000}%
\pgfsetstrokecolor{currentstroke}%
\pgfsetdash{}{0pt}%
\pgfpathmoveto{\pgfqpoint{1.600895in}{1.056387in}}%
\pgfpathlineto{\pgfqpoint{1.618503in}{1.067480in}}%
\pgfpathlineto{\pgfqpoint{1.636188in}{1.078777in}}%
\pgfpathlineto{\pgfqpoint{1.643732in}{1.089165in}}%
\pgfpathlineto{\pgfqpoint{1.651264in}{1.099804in}}%
\pgfpathlineto{\pgfqpoint{1.633602in}{1.088830in}}%
\pgfpathlineto{\pgfqpoint{1.616016in}{1.078076in}}%
\pgfpathlineto{\pgfqpoint{1.608461in}{1.067102in}}%
\pgfpathlineto{\pgfqpoint{1.600895in}{1.056387in}}%
\pgfpathclose%
\pgfusepath{fill}%
\end{pgfscope}%
\begin{pgfscope}%
\pgfpathrectangle{\pgfqpoint{0.407592in}{0.222904in}}{\pgfqpoint{1.411163in}{1.411163in}}%
\pgfusepath{clip}%
\pgfsetbuttcap%
\pgfsetroundjoin%
\definecolor{currentfill}{rgb}{0.556863,0.654902,0.847059}%
\pgfsetfillcolor{currentfill}%
\pgfsetfillopacity{0.130000}%
\pgfsetlinewidth{0.000000pt}%
\definecolor{currentstroke}{rgb}{0.000000,0.000000,0.000000}%
\pgfsetstrokecolor{currentstroke}%
\pgfsetdash{}{0pt}%
\pgfpathmoveto{\pgfqpoint{1.254845in}{0.782394in}}%
\pgfpathlineto{\pgfqpoint{1.271735in}{0.792254in}}%
\pgfpathlineto{\pgfqpoint{1.288698in}{0.802519in}}%
\pgfpathlineto{\pgfqpoint{1.296628in}{0.806605in}}%
\pgfpathlineto{\pgfqpoint{1.304539in}{0.811307in}}%
\pgfpathlineto{\pgfqpoint{1.287598in}{0.800949in}}%
\pgfpathlineto{\pgfqpoint{1.270730in}{0.790984in}}%
\pgfpathlineto{\pgfqpoint{1.262798in}{0.786346in}}%
\pgfpathlineto{\pgfqpoint{1.254845in}{0.782394in}}%
\pgfpathclose%
\pgfusepath{fill}%
\end{pgfscope}%
\begin{pgfscope}%
\pgfpathrectangle{\pgfqpoint{0.407592in}{0.222904in}}{\pgfqpoint{1.411163in}{1.411163in}}%
\pgfusepath{clip}%
\pgfsetbuttcap%
\pgfsetroundjoin%
\definecolor{currentfill}{rgb}{0.556863,0.654902,0.847059}%
\pgfsetfillcolor{currentfill}%
\pgfsetfillopacity{0.130000}%
\pgfsetlinewidth{0.000000pt}%
\definecolor{currentstroke}{rgb}{0.000000,0.000000,0.000000}%
\pgfsetstrokecolor{currentstroke}%
\pgfsetdash{}{0pt}%
\pgfpathmoveto{\pgfqpoint{0.841063in}{0.824538in}}%
\pgfpathlineto{\pgfqpoint{0.857457in}{0.813559in}}%
\pgfpathlineto{\pgfqpoint{0.873867in}{0.803505in}}%
\pgfpathlineto{\pgfqpoint{0.882605in}{0.798555in}}%
\pgfpathlineto{\pgfqpoint{0.891302in}{0.793879in}}%
\pgfpathlineto{\pgfqpoint{0.874940in}{0.804363in}}%
\pgfpathlineto{\pgfqpoint{0.858593in}{0.815816in}}%
\pgfpathlineto{\pgfqpoint{0.849849in}{0.820029in}}%
\pgfpathlineto{\pgfqpoint{0.841063in}{0.824538in}}%
\pgfpathclose%
\pgfusepath{fill}%
\end{pgfscope}%
\begin{pgfscope}%
\pgfpathrectangle{\pgfqpoint{0.407592in}{0.222904in}}{\pgfqpoint{1.411163in}{1.411163in}}%
\pgfusepath{clip}%
\pgfsetbuttcap%
\pgfsetroundjoin%
\definecolor{currentfill}{rgb}{0.556863,0.654902,0.847059}%
\pgfsetfillcolor{currentfill}%
\pgfsetfillopacity{0.130000}%
\pgfsetlinewidth{0.000000pt}%
\definecolor{currentstroke}{rgb}{0.000000,0.000000,0.000000}%
\pgfsetstrokecolor{currentstroke}%
\pgfsetdash{}{0pt}%
\pgfpathmoveto{\pgfqpoint{0.459061in}{1.160764in}}%
\pgfpathlineto{\pgfqpoint{0.475395in}{1.142637in}}%
\pgfpathlineto{\pgfqpoint{0.491736in}{1.124683in}}%
\pgfpathlineto{\pgfqpoint{0.501176in}{1.122417in}}%
\pgfpathlineto{\pgfqpoint{0.510571in}{1.120393in}}%
\pgfpathlineto{\pgfqpoint{0.494272in}{1.138639in}}%
\pgfpathlineto{\pgfqpoint{0.477981in}{1.157044in}}%
\pgfpathlineto{\pgfqpoint{0.468544in}{1.158785in}}%
\pgfpathlineto{\pgfqpoint{0.459061in}{1.160764in}}%
\pgfpathclose%
\pgfusepath{fill}%
\end{pgfscope}%
\begin{pgfscope}%
\pgfpathrectangle{\pgfqpoint{0.407592in}{0.222904in}}{\pgfqpoint{1.411163in}{1.411163in}}%
\pgfusepath{clip}%
\pgfsetbuttcap%
\pgfsetroundjoin%
\definecolor{currentfill}{rgb}{0.556863,0.654902,0.847059}%
\pgfsetfillcolor{currentfill}%
\pgfsetfillopacity{0.130000}%
\pgfsetlinewidth{0.000000pt}%
\definecolor{currentstroke}{rgb}{0.000000,0.000000,0.000000}%
\pgfsetstrokecolor{currentstroke}%
\pgfsetdash{}{0pt}%
\pgfpathmoveto{\pgfqpoint{0.608546in}{1.015186in}}%
\pgfpathlineto{\pgfqpoint{0.624906in}{0.998553in}}%
\pgfpathlineto{\pgfqpoint{0.641274in}{0.982244in}}%
\pgfpathlineto{\pgfqpoint{0.650453in}{0.979039in}}%
\pgfpathlineto{\pgfqpoint{0.659589in}{0.976103in}}%
\pgfpathlineto{\pgfqpoint{0.643264in}{0.992828in}}%
\pgfpathlineto{\pgfqpoint{0.626947in}{1.009861in}}%
\pgfpathlineto{\pgfqpoint{0.617769in}{1.012388in}}%
\pgfpathlineto{\pgfqpoint{0.608546in}{1.015186in}}%
\pgfpathclose%
\pgfusepath{fill}%
\end{pgfscope}%
\begin{pgfscope}%
\pgfpathrectangle{\pgfqpoint{0.407592in}{0.222904in}}{\pgfqpoint{1.411163in}{1.411163in}}%
\pgfusepath{clip}%
\pgfsetbuttcap%
\pgfsetroundjoin%
\definecolor{currentfill}{rgb}{0.556863,0.654902,0.847059}%
\pgfsetfillcolor{currentfill}%
\pgfsetfillopacity{0.130000}%
\pgfsetlinewidth{0.000000pt}%
\definecolor{currentstroke}{rgb}{0.000000,0.000000,0.000000}%
\pgfsetstrokecolor{currentstroke}%
\pgfsetdash{}{0pt}%
\pgfpathmoveto{\pgfqpoint{1.419893in}{0.894947in}}%
\pgfpathlineto{\pgfqpoint{1.437147in}{0.906288in}}%
\pgfpathlineto{\pgfqpoint{1.454476in}{0.917843in}}%
\pgfpathlineto{\pgfqpoint{1.462210in}{0.925794in}}%
\pgfpathlineto{\pgfqpoint{1.469930in}{0.934159in}}%
\pgfpathlineto{\pgfqpoint{1.452621in}{0.922862in}}%
\pgfpathlineto{\pgfqpoint{1.435389in}{0.911799in}}%
\pgfpathlineto{\pgfqpoint{1.427648in}{0.903150in}}%
\pgfpathlineto{\pgfqpoint{1.419893in}{0.894947in}}%
\pgfpathclose%
\pgfusepath{fill}%
\end{pgfscope}%
\begin{pgfscope}%
\pgfpathrectangle{\pgfqpoint{0.407592in}{0.222904in}}{\pgfqpoint{1.411163in}{1.411163in}}%
\pgfusepath{clip}%
\pgfsetbuttcap%
\pgfsetroundjoin%
\definecolor{currentfill}{rgb}{0.556863,0.654902,0.847059}%
\pgfsetfillcolor{currentfill}%
\pgfsetfillopacity{0.130000}%
\pgfsetlinewidth{0.000000pt}%
\definecolor{currentstroke}{rgb}{0.000000,0.000000,0.000000}%
\pgfsetstrokecolor{currentstroke}%
\pgfsetdash{}{0pt}%
\pgfpathmoveto{\pgfqpoint{1.782194in}{1.233709in}}%
\pgfpathlineto{\pgfqpoint{1.800156in}{1.244470in}}%
\pgfpathlineto{\pgfqpoint{1.818197in}{1.255429in}}%
\pgfpathlineto{\pgfqpoint{1.825562in}{1.266988in}}%
\pgfpathlineto{\pgfqpoint{1.807533in}{1.256172in}}%
\pgfpathlineto{\pgfqpoint{1.789583in}{1.245560in}}%
\pgfusepath{fill}%
\end{pgfscope}%
\begin{pgfscope}%
\pgfpathrectangle{\pgfqpoint{0.407592in}{0.222904in}}{\pgfqpoint{1.411163in}{1.411163in}}%
\pgfusepath{clip}%
\pgfsetbuttcap%
\pgfsetroundjoin%
\definecolor{currentfill}{rgb}{0.556863,0.654902,0.847059}%
\pgfsetfillcolor{currentfill}%
\pgfsetfillopacity{0.130000}%
\pgfsetlinewidth{0.000000pt}%
\definecolor{currentstroke}{rgb}{0.000000,0.000000,0.000000}%
\pgfsetstrokecolor{currentstroke}%
\pgfsetdash{}{0pt}%
\pgfpathmoveto{\pgfqpoint{1.006881in}{0.745540in}}%
\pgfpathlineto{\pgfqpoint{1.023365in}{0.743601in}}%
\pgfpathlineto{\pgfqpoint{1.039888in}{0.742986in}}%
\pgfpathlineto{\pgfqpoint{1.048272in}{0.739031in}}%
\pgfpathlineto{\pgfqpoint{1.056620in}{0.735487in}}%
\pgfpathlineto{\pgfqpoint{1.040144in}{0.735765in}}%
\pgfpathlineto{\pgfqpoint{1.023708in}{0.737478in}}%
\pgfpathlineto{\pgfqpoint{1.015313in}{0.741310in}}%
\pgfpathlineto{\pgfqpoint{1.006881in}{0.745540in}}%
\pgfpathclose%
\pgfusepath{fill}%
\end{pgfscope}%
\begin{pgfscope}%
\pgfpathrectangle{\pgfqpoint{0.407592in}{0.222904in}}{\pgfqpoint{1.411163in}{1.411163in}}%
\pgfusepath{clip}%
\pgfsetbuttcap%
\pgfsetroundjoin%
\definecolor{currentfill}{rgb}{0.556863,0.654902,0.847059}%
\pgfsetfillcolor{currentfill}%
\pgfsetfillopacity{0.130000}%
\pgfsetlinewidth{0.000000pt}%
\definecolor{currentstroke}{rgb}{0.000000,0.000000,0.000000}%
\pgfsetstrokecolor{currentstroke}%
\pgfsetdash{}{0pt}%
\pgfpathmoveto{\pgfqpoint{1.716752in}{1.165971in}}%
\pgfpathlineto{\pgfqpoint{1.734604in}{1.176930in}}%
\pgfpathlineto{\pgfqpoint{1.752534in}{1.188087in}}%
\pgfpathlineto{\pgfqpoint{1.759965in}{1.199211in}}%
\pgfpathlineto{\pgfqpoint{1.767385in}{1.210528in}}%
\pgfpathlineto{\pgfqpoint{1.749479in}{1.199672in}}%
\pgfpathlineto{\pgfqpoint{1.731650in}{1.189024in}}%
\pgfpathlineto{\pgfqpoint{1.724207in}{1.177400in}}%
\pgfpathlineto{\pgfqpoint{1.716752in}{1.165971in}}%
\pgfpathclose%
\pgfusepath{fill}%
\end{pgfscope}%
\begin{pgfscope}%
\pgfpathrectangle{\pgfqpoint{0.407592in}{0.222904in}}{\pgfqpoint{1.411163in}{1.411163in}}%
\pgfusepath{clip}%
\pgfsetbuttcap%
\pgfsetroundjoin%
\definecolor{currentfill}{rgb}{0.556863,0.654902,0.847059}%
\pgfsetfillcolor{currentfill}%
\pgfsetfillopacity{0.130000}%
\pgfsetlinewidth{0.000000pt}%
\definecolor{currentstroke}{rgb}{0.000000,0.000000,0.000000}%
\pgfsetstrokecolor{currentstroke}%
\pgfsetdash{}{0pt}%
\pgfpathmoveto{\pgfqpoint{1.535443in}{0.993491in}}%
\pgfpathlineto{\pgfqpoint{1.552939in}{1.004816in}}%
\pgfpathlineto{\pgfqpoint{1.570514in}{1.016342in}}%
\pgfpathlineto{\pgfqpoint{1.578127in}{1.025909in}}%
\pgfpathlineto{\pgfqpoint{1.585728in}{1.035780in}}%
\pgfpathlineto{\pgfqpoint{1.568176in}{1.024568in}}%
\pgfpathlineto{\pgfqpoint{1.550701in}{1.013574in}}%
\pgfpathlineto{\pgfqpoint{1.543078in}{1.003374in}}%
\pgfpathlineto{\pgfqpoint{1.535443in}{0.993491in}}%
\pgfpathclose%
\pgfusepath{fill}%
\end{pgfscope}%
\begin{pgfscope}%
\pgfpathrectangle{\pgfqpoint{0.407592in}{0.222904in}}{\pgfqpoint{1.411163in}{1.411163in}}%
\pgfusepath{clip}%
\pgfsetbuttcap%
\pgfsetroundjoin%
\definecolor{currentfill}{rgb}{0.556863,0.654902,0.847059}%
\pgfsetfillcolor{currentfill}%
\pgfsetfillopacity{0.130000}%
\pgfsetlinewidth{0.000000pt}%
\definecolor{currentstroke}{rgb}{0.000000,0.000000,0.000000}%
\pgfsetstrokecolor{currentstroke}%
\pgfsetdash{}{0pt}%
\pgfpathmoveto{\pgfqpoint{1.304539in}{0.811307in}}%
\pgfpathlineto{\pgfqpoint{1.321553in}{0.822006in}}%
\pgfpathlineto{\pgfqpoint{1.338642in}{0.833004in}}%
\pgfpathlineto{\pgfqpoint{1.346512in}{0.838311in}}%
\pgfpathlineto{\pgfqpoint{1.354365in}{0.844172in}}%
\pgfpathlineto{\pgfqpoint{1.337297in}{0.833216in}}%
\pgfpathlineto{\pgfqpoint{1.320303in}{0.822560in}}%
\pgfpathlineto{\pgfqpoint{1.312430in}{0.816626in}}%
\pgfpathlineto{\pgfqpoint{1.304539in}{0.811307in}}%
\pgfpathclose%
\pgfusepath{fill}%
\end{pgfscope}%
\begin{pgfscope}%
\pgfpathrectangle{\pgfqpoint{0.407592in}{0.222904in}}{\pgfqpoint{1.411163in}{1.411163in}}%
\pgfusepath{clip}%
\pgfsetbuttcap%
\pgfsetroundjoin%
\definecolor{currentfill}{rgb}{0.556863,0.654902,0.847059}%
\pgfsetfillcolor{currentfill}%
\pgfsetfillopacity{0.130000}%
\pgfsetlinewidth{0.000000pt}%
\definecolor{currentstroke}{rgb}{0.000000,0.000000,0.000000}%
\pgfsetstrokecolor{currentstroke}%
\pgfsetdash{}{0pt}%
\pgfpathmoveto{\pgfqpoint{0.491736in}{1.124683in}}%
\pgfpathlineto{\pgfqpoint{0.508087in}{1.106918in}}%
\pgfpathlineto{\pgfqpoint{0.524446in}{1.089358in}}%
\pgfpathlineto{\pgfqpoint{0.533843in}{1.086776in}}%
\pgfpathlineto{\pgfqpoint{0.543195in}{1.084439in}}%
\pgfpathlineto{\pgfqpoint{0.526879in}{1.102321in}}%
\pgfpathlineto{\pgfqpoint{0.510571in}{1.120393in}}%
\pgfpathlineto{\pgfqpoint{0.501176in}{1.122417in}}%
\pgfpathlineto{\pgfqpoint{0.491736in}{1.124683in}}%
\pgfpathclose%
\pgfusepath{fill}%
\end{pgfscope}%
\begin{pgfscope}%
\pgfpathrectangle{\pgfqpoint{0.407592in}{0.222904in}}{\pgfqpoint{1.411163in}{1.411163in}}%
\pgfusepath{clip}%
\pgfsetbuttcap%
\pgfsetroundjoin%
\definecolor{currentfill}{rgb}{0.556863,0.654902,0.847059}%
\pgfsetfillcolor{currentfill}%
\pgfsetfillopacity{0.130000}%
\pgfsetlinewidth{0.000000pt}%
\definecolor{currentstroke}{rgb}{0.000000,0.000000,0.000000}%
\pgfsetstrokecolor{currentstroke}%
\pgfsetdash{}{0pt}%
\pgfpathmoveto{\pgfqpoint{0.757730in}{0.884208in}}%
\pgfpathlineto{\pgfqpoint{0.774121in}{0.870705in}}%
\pgfpathlineto{\pgfqpoint{0.790524in}{0.857854in}}%
\pgfpathlineto{\pgfqpoint{0.799440in}{0.853295in}}%
\pgfpathlineto{\pgfqpoint{0.808313in}{0.849002in}}%
\pgfpathlineto{\pgfqpoint{0.791956in}{0.862353in}}%
\pgfpathlineto{\pgfqpoint{0.775610in}{0.876361in}}%
\pgfpathlineto{\pgfqpoint{0.766691in}{0.880144in}}%
\pgfpathlineto{\pgfqpoint{0.757730in}{0.884208in}}%
\pgfpathclose%
\pgfusepath{fill}%
\end{pgfscope}%
\begin{pgfscope}%
\pgfpathrectangle{\pgfqpoint{0.407592in}{0.222904in}}{\pgfqpoint{1.411163in}{1.411163in}}%
\pgfusepath{clip}%
\pgfsetbuttcap%
\pgfsetroundjoin%
\definecolor{currentfill}{rgb}{0.556863,0.654902,0.847059}%
\pgfsetfillcolor{currentfill}%
\pgfsetfillopacity{0.130000}%
\pgfsetlinewidth{0.000000pt}%
\definecolor{currentstroke}{rgb}{0.000000,0.000000,0.000000}%
\pgfsetstrokecolor{currentstroke}%
\pgfsetdash{}{0pt}%
\pgfpathmoveto{\pgfqpoint{0.641274in}{0.982244in}}%
\pgfpathlineto{\pgfqpoint{0.657651in}{0.966290in}}%
\pgfpathlineto{\pgfqpoint{0.674037in}{0.950724in}}%
\pgfpathlineto{\pgfqpoint{0.683173in}{0.947083in}}%
\pgfpathlineto{\pgfqpoint{0.692266in}{0.943706in}}%
\pgfpathlineto{\pgfqpoint{0.675923in}{0.959718in}}%
\pgfpathlineto{\pgfqpoint{0.659589in}{0.976103in}}%
\pgfpathlineto{\pgfqpoint{0.650453in}{0.979039in}}%
\pgfpathlineto{\pgfqpoint{0.641274in}{0.982244in}}%
\pgfpathclose%
\pgfusepath{fill}%
\end{pgfscope}%
\begin{pgfscope}%
\pgfpathrectangle{\pgfqpoint{0.407592in}{0.222904in}}{\pgfqpoint{1.411163in}{1.411163in}}%
\pgfusepath{clip}%
\pgfsetbuttcap%
\pgfsetroundjoin%
\definecolor{currentfill}{rgb}{0.556863,0.654902,0.847059}%
\pgfsetfillcolor{currentfill}%
\pgfsetfillopacity{0.130000}%
\pgfsetlinewidth{0.000000pt}%
\definecolor{currentstroke}{rgb}{0.000000,0.000000,0.000000}%
\pgfsetstrokecolor{currentstroke}%
\pgfsetdash{}{0pt}%
\pgfpathmoveto{\pgfqpoint{1.651264in}{1.099804in}}%
\pgfpathlineto{\pgfqpoint{1.669006in}{1.110985in}}%
\pgfpathlineto{\pgfqpoint{1.686825in}{1.122362in}}%
\pgfpathlineto{\pgfqpoint{1.694323in}{1.132931in}}%
\pgfpathlineto{\pgfqpoint{1.701811in}{1.143728in}}%
\pgfpathlineto{\pgfqpoint{1.684014in}{1.132660in}}%
\pgfpathlineto{\pgfqpoint{1.666296in}{1.121801in}}%
\pgfpathlineto{\pgfqpoint{1.658786in}{1.110686in}}%
\pgfpathlineto{\pgfqpoint{1.651264in}{1.099804in}}%
\pgfpathclose%
\pgfusepath{fill}%
\end{pgfscope}%
\begin{pgfscope}%
\pgfpathrectangle{\pgfqpoint{0.407592in}{0.222904in}}{\pgfqpoint{1.411163in}{1.411163in}}%
\pgfusepath{clip}%
\pgfsetbuttcap%
\pgfsetroundjoin%
\definecolor{currentfill}{rgb}{0.556863,0.654902,0.847059}%
\pgfsetfillcolor{currentfill}%
\pgfsetfillopacity{0.130000}%
\pgfsetlinewidth{0.000000pt}%
\definecolor{currentstroke}{rgb}{0.000000,0.000000,0.000000}%
\pgfsetstrokecolor{currentstroke}%
\pgfsetdash{}{0pt}%
\pgfpathmoveto{\pgfqpoint{1.469930in}{0.934159in}}%
\pgfpathlineto{\pgfqpoint{1.487315in}{0.945669in}}%
\pgfpathlineto{\pgfqpoint{1.504778in}{0.957375in}}%
\pgfpathlineto{\pgfqpoint{1.512463in}{0.965867in}}%
\pgfpathlineto{\pgfqpoint{1.520136in}{0.974725in}}%
\pgfpathlineto{\pgfqpoint{1.502695in}{0.963296in}}%
\pgfpathlineto{\pgfqpoint{1.485330in}{0.952082in}}%
\pgfpathlineto{\pgfqpoint{1.477637in}{0.942926in}}%
\pgfpathlineto{\pgfqpoint{1.469930in}{0.934159in}}%
\pgfpathclose%
\pgfusepath{fill}%
\end{pgfscope}%
\begin{pgfscope}%
\pgfpathrectangle{\pgfqpoint{0.407592in}{0.222904in}}{\pgfqpoint{1.411163in}{1.411163in}}%
\pgfusepath{clip}%
\pgfsetbuttcap%
\pgfsetroundjoin%
\definecolor{currentfill}{rgb}{0.556863,0.654902,0.847059}%
\pgfsetfillcolor{currentfill}%
\pgfsetfillopacity{0.130000}%
\pgfsetlinewidth{0.000000pt}%
\definecolor{currentstroke}{rgb}{0.000000,0.000000,0.000000}%
\pgfsetstrokecolor{currentstroke}%
\pgfsetdash{}{0pt}%
\pgfpathmoveto{\pgfqpoint{0.873867in}{0.803505in}}%
\pgfpathlineto{\pgfqpoint{0.890293in}{0.794441in}}%
\pgfpathlineto{\pgfqpoint{0.906739in}{0.786422in}}%
\pgfpathlineto{\pgfqpoint{0.915428in}{0.781145in}}%
\pgfpathlineto{\pgfqpoint{0.924077in}{0.776122in}}%
\pgfpathlineto{\pgfqpoint{0.907681in}{0.784441in}}%
\pgfpathlineto{\pgfqpoint{0.891302in}{0.793879in}}%
\pgfpathlineto{\pgfqpoint{0.882605in}{0.798555in}}%
\pgfpathlineto{\pgfqpoint{0.873867in}{0.803505in}}%
\pgfpathclose%
\pgfusepath{fill}%
\end{pgfscope}%
\begin{pgfscope}%
\pgfpathrectangle{\pgfqpoint{0.407592in}{0.222904in}}{\pgfqpoint{1.411163in}{1.411163in}}%
\pgfusepath{clip}%
\pgfsetbuttcap%
\pgfsetroundjoin%
\definecolor{currentfill}{rgb}{0.556863,0.654902,0.847059}%
\pgfsetfillcolor{currentfill}%
\pgfsetfillopacity{0.130000}%
\pgfsetlinewidth{0.000000pt}%
\definecolor{currentstroke}{rgb}{0.000000,0.000000,0.000000}%
\pgfsetstrokecolor{currentstroke}%
\pgfsetdash{}{0pt}%
\pgfpathmoveto{\pgfqpoint{1.354365in}{0.844172in}}%
\pgfpathlineto{\pgfqpoint{1.371508in}{0.855394in}}%
\pgfpathlineto{\pgfqpoint{1.388726in}{0.866852in}}%
\pgfpathlineto{\pgfqpoint{1.396542in}{0.873147in}}%
\pgfpathlineto{\pgfqpoint{1.404341in}{0.879934in}}%
\pgfpathlineto{\pgfqpoint{1.387143in}{0.868610in}}%
\pgfpathlineto{\pgfqpoint{1.370020in}{0.857532in}}%
\pgfpathlineto{\pgfqpoint{1.362201in}{0.850582in}}%
\pgfpathlineto{\pgfqpoint{1.354365in}{0.844172in}}%
\pgfpathclose%
\pgfusepath{fill}%
\end{pgfscope}%
\begin{pgfscope}%
\pgfpathrectangle{\pgfqpoint{0.407592in}{0.222904in}}{\pgfqpoint{1.411163in}{1.411163in}}%
\pgfusepath{clip}%
\pgfsetbuttcap%
\pgfsetroundjoin%
\definecolor{currentfill}{rgb}{0.556863,0.654902,0.847059}%
\pgfsetfillcolor{currentfill}%
\pgfsetfillopacity{0.130000}%
\pgfsetlinewidth{0.000000pt}%
\definecolor{currentstroke}{rgb}{0.000000,0.000000,0.000000}%
\pgfsetstrokecolor{currentstroke}%
\pgfsetdash{}{0pt}%
\pgfpathmoveto{\pgfqpoint{0.524446in}{1.089358in}}%
\pgfpathlineto{\pgfqpoint{0.540813in}{1.072020in}}%
\pgfpathlineto{\pgfqpoint{0.557189in}{1.054924in}}%
\pgfpathlineto{\pgfqpoint{0.566544in}{1.051996in}}%
\pgfpathlineto{\pgfqpoint{0.575854in}{1.049316in}}%
\pgfpathlineto{\pgfqpoint{0.559520in}{1.066764in}}%
\pgfpathlineto{\pgfqpoint{0.543195in}{1.084439in}}%
\pgfpathlineto{\pgfqpoint{0.533843in}{1.086776in}}%
\pgfpathlineto{\pgfqpoint{0.524446in}{1.089358in}}%
\pgfpathclose%
\pgfusepath{fill}%
\end{pgfscope}%
\begin{pgfscope}%
\pgfpathrectangle{\pgfqpoint{0.407592in}{0.222904in}}{\pgfqpoint{1.411163in}{1.411163in}}%
\pgfusepath{clip}%
\pgfsetbuttcap%
\pgfsetroundjoin%
\definecolor{currentfill}{rgb}{0.556863,0.654902,0.847059}%
\pgfsetfillcolor{currentfill}%
\pgfsetfillopacity{0.130000}%
\pgfsetlinewidth{0.000000pt}%
\definecolor{currentstroke}{rgb}{0.000000,0.000000,0.000000}%
\pgfsetstrokecolor{currentstroke}%
\pgfsetdash{}{0pt}%
\pgfpathmoveto{\pgfqpoint{0.956935in}{0.763069in}}%
\pgfpathlineto{\pgfqpoint{0.973403in}{0.758409in}}%
\pgfpathlineto{\pgfqpoint{0.989901in}{0.755014in}}%
\pgfpathlineto{\pgfqpoint{0.998411in}{0.750122in}}%
\pgfpathlineto{\pgfqpoint{1.006881in}{0.745540in}}%
\pgfpathlineto{\pgfqpoint{0.990432in}{0.748845in}}%
\pgfpathlineto{\pgfqpoint{0.974014in}{0.753530in}}%
\pgfpathlineto{\pgfqpoint{0.965494in}{0.758145in}}%
\pgfpathlineto{\pgfqpoint{0.956935in}{0.763069in}}%
\pgfpathclose%
\pgfusepath{fill}%
\end{pgfscope}%
\begin{pgfscope}%
\pgfpathrectangle{\pgfqpoint{0.407592in}{0.222904in}}{\pgfqpoint{1.411163in}{1.411163in}}%
\pgfusepath{clip}%
\pgfsetbuttcap%
\pgfsetroundjoin%
\definecolor{currentfill}{rgb}{0.556863,0.654902,0.847059}%
\pgfsetfillcolor{currentfill}%
\pgfsetfillopacity{0.130000}%
\pgfsetlinewidth{0.000000pt}%
\definecolor{currentstroke}{rgb}{0.000000,0.000000,0.000000}%
\pgfsetstrokecolor{currentstroke}%
\pgfsetdash{}{0pt}%
\pgfpathmoveto{\pgfqpoint{1.767385in}{1.210528in}}%
\pgfpathlineto{\pgfqpoint{1.785372in}{1.221583in}}%
\pgfpathlineto{\pgfqpoint{1.803438in}{1.232829in}}%
\pgfpathlineto{\pgfqpoint{1.810823in}{1.244041in}}%
\pgfpathlineto{\pgfqpoint{1.818197in}{1.255429in}}%
\pgfpathlineto{\pgfqpoint{1.800156in}{1.244470in}}%
\pgfpathlineto{\pgfqpoint{1.782194in}{1.233709in}}%
\pgfpathlineto{\pgfqpoint{1.774795in}{1.222030in}}%
\pgfpathlineto{\pgfqpoint{1.767385in}{1.210528in}}%
\pgfpathclose%
\pgfusepath{fill}%
\end{pgfscope}%
\begin{pgfscope}%
\pgfpathrectangle{\pgfqpoint{0.407592in}{0.222904in}}{\pgfqpoint{1.411163in}{1.411163in}}%
\pgfusepath{clip}%
\pgfsetbuttcap%
\pgfsetroundjoin%
\definecolor{currentfill}{rgb}{0.556863,0.654902,0.847059}%
\pgfsetfillcolor{currentfill}%
\pgfsetfillopacity{0.130000}%
\pgfsetlinewidth{0.000000pt}%
\definecolor{currentstroke}{rgb}{0.000000,0.000000,0.000000}%
\pgfsetstrokecolor{currentstroke}%
\pgfsetdash{}{0pt}%
\pgfpathmoveto{\pgfqpoint{1.585728in}{1.035780in}}%
\pgfpathlineto{\pgfqpoint{1.603358in}{1.047194in}}%
\pgfpathlineto{\pgfqpoint{1.621066in}{1.058798in}}%
\pgfpathlineto{\pgfqpoint{1.628633in}{1.068652in}}%
\pgfpathlineto{\pgfqpoint{1.636188in}{1.078777in}}%
\pgfpathlineto{\pgfqpoint{1.618503in}{1.067480in}}%
\pgfpathlineto{\pgfqpoint{1.600895in}{1.056387in}}%
\pgfpathlineto{\pgfqpoint{1.593317in}{1.045943in}}%
\pgfpathlineto{\pgfqpoint{1.585728in}{1.035780in}}%
\pgfpathclose%
\pgfusepath{fill}%
\end{pgfscope}%
\begin{pgfscope}%
\pgfpathrectangle{\pgfqpoint{0.407592in}{0.222904in}}{\pgfqpoint{1.411163in}{1.411163in}}%
\pgfusepath{clip}%
\pgfsetbuttcap%
\pgfsetroundjoin%
\definecolor{currentfill}{rgb}{0.556863,0.654902,0.847059}%
\pgfsetfillcolor{currentfill}%
\pgfsetfillopacity{0.130000}%
\pgfsetlinewidth{0.000000pt}%
\definecolor{currentstroke}{rgb}{0.000000,0.000000,0.000000}%
\pgfsetstrokecolor{currentstroke}%
\pgfsetdash{}{0pt}%
\pgfpathmoveto{\pgfqpoint{0.674037in}{0.950724in}}%
\pgfpathlineto{\pgfqpoint{0.690432in}{0.935583in}}%
\pgfpathlineto{\pgfqpoint{0.706837in}{0.920904in}}%
\pgfpathlineto{\pgfqpoint{0.715930in}{0.916803in}}%
\pgfpathlineto{\pgfqpoint{0.724979in}{0.912959in}}%
\pgfpathlineto{\pgfqpoint{0.708618in}{0.928106in}}%
\pgfpathlineto{\pgfqpoint{0.692266in}{0.943706in}}%
\pgfpathlineto{\pgfqpoint{0.683173in}{0.947083in}}%
\pgfpathlineto{\pgfqpoint{0.674037in}{0.950724in}}%
\pgfpathclose%
\pgfusepath{fill}%
\end{pgfscope}%
\begin{pgfscope}%
\pgfpathrectangle{\pgfqpoint{0.407592in}{0.222904in}}{\pgfqpoint{1.411163in}{1.411163in}}%
\pgfusepath{clip}%
\pgfsetbuttcap%
\pgfsetroundjoin%
\definecolor{currentfill}{rgb}{0.556863,0.654902,0.847059}%
\pgfsetfillcolor{currentfill}%
\pgfsetfillopacity{0.130000}%
\pgfsetlinewidth{0.000000pt}%
\definecolor{currentstroke}{rgb}{0.000000,0.000000,0.000000}%
\pgfsetstrokecolor{currentstroke}%
\pgfsetdash{}{0pt}%
\pgfpathmoveto{\pgfqpoint{1.139435in}{0.743451in}}%
\pgfpathlineto{\pgfqpoint{1.156136in}{0.749515in}}%
\pgfpathlineto{\pgfqpoint{1.172900in}{0.756403in}}%
\pgfpathlineto{\pgfqpoint{1.181032in}{0.756073in}}%
\pgfpathlineto{\pgfqpoint{1.189135in}{0.756344in}}%
\pgfpathlineto{\pgfqpoint{1.172402in}{0.748915in}}%
\pgfpathlineto{\pgfqpoint{1.155734in}{0.742290in}}%
\pgfpathlineto{\pgfqpoint{1.147600in}{0.742555in}}%
\pgfpathlineto{\pgfqpoint{1.139435in}{0.743451in}}%
\pgfpathclose%
\pgfusepath{fill}%
\end{pgfscope}%
\begin{pgfscope}%
\pgfpathrectangle{\pgfqpoint{0.407592in}{0.222904in}}{\pgfqpoint{1.411163in}{1.411163in}}%
\pgfusepath{clip}%
\pgfsetbuttcap%
\pgfsetroundjoin%
\definecolor{currentfill}{rgb}{0.556863,0.654902,0.847059}%
\pgfsetfillcolor{currentfill}%
\pgfsetfillopacity{0.130000}%
\pgfsetlinewidth{0.000000pt}%
\definecolor{currentstroke}{rgb}{0.000000,0.000000,0.000000}%
\pgfsetstrokecolor{currentstroke}%
\pgfsetdash{}{0pt}%
\pgfpathmoveto{\pgfqpoint{0.790524in}{0.857854in}}%
\pgfpathlineto{\pgfqpoint{0.806938in}{0.845714in}}%
\pgfpathlineto{\pgfqpoint{0.823366in}{0.834341in}}%
\pgfpathlineto{\pgfqpoint{0.832235in}{0.829316in}}%
\pgfpathlineto{\pgfqpoint{0.841063in}{0.824538in}}%
\pgfpathlineto{\pgfqpoint{0.824682in}{0.836375in}}%
\pgfpathlineto{\pgfqpoint{0.808313in}{0.849002in}}%
\pgfpathlineto{\pgfqpoint{0.799440in}{0.853295in}}%
\pgfpathlineto{\pgfqpoint{0.790524in}{0.857854in}}%
\pgfpathclose%
\pgfusepath{fill}%
\end{pgfscope}%
\begin{pgfscope}%
\pgfpathrectangle{\pgfqpoint{0.407592in}{0.222904in}}{\pgfqpoint{1.411163in}{1.411163in}}%
\pgfusepath{clip}%
\pgfsetbuttcap%
\pgfsetroundjoin%
\definecolor{currentfill}{rgb}{0.556863,0.654902,0.847059}%
\pgfsetfillcolor{currentfill}%
\pgfsetfillopacity{0.130000}%
\pgfsetlinewidth{0.000000pt}%
\definecolor{currentstroke}{rgb}{0.000000,0.000000,0.000000}%
\pgfsetstrokecolor{currentstroke}%
\pgfsetdash{}{0pt}%
\pgfpathmoveto{\pgfqpoint{1.189135in}{0.756344in}}%
\pgfpathlineto{\pgfqpoint{1.205935in}{0.764469in}}%
\pgfpathlineto{\pgfqpoint{1.222804in}{0.773194in}}%
\pgfpathlineto{\pgfqpoint{1.230851in}{0.774542in}}%
\pgfpathlineto{\pgfqpoint{1.238873in}{0.776510in}}%
\pgfpathlineto{\pgfqpoint{1.222031in}{0.767388in}}%
\pgfpathlineto{\pgfqpoint{1.205258in}{0.758838in}}%
\pgfpathlineto{\pgfqpoint{1.197210in}{0.757253in}}%
\pgfpathlineto{\pgfqpoint{1.189135in}{0.756344in}}%
\pgfpathclose%
\pgfusepath{fill}%
\end{pgfscope}%
\begin{pgfscope}%
\pgfpathrectangle{\pgfqpoint{0.407592in}{0.222904in}}{\pgfqpoint{1.411163in}{1.411163in}}%
\pgfusepath{clip}%
\pgfsetbuttcap%
\pgfsetroundjoin%
\definecolor{currentfill}{rgb}{0.556863,0.654902,0.847059}%
\pgfsetfillcolor{currentfill}%
\pgfsetfillopacity{0.130000}%
\pgfsetlinewidth{0.000000pt}%
\definecolor{currentstroke}{rgb}{0.000000,0.000000,0.000000}%
\pgfsetstrokecolor{currentstroke}%
\pgfsetdash{}{0pt}%
\pgfpathmoveto{\pgfqpoint{1.089710in}{0.738908in}}%
\pgfpathlineto{\pgfqpoint{1.106332in}{0.742401in}}%
\pgfpathlineto{\pgfqpoint{1.123009in}{0.746937in}}%
\pgfpathlineto{\pgfqpoint{1.131239in}{0.744928in}}%
\pgfpathlineto{\pgfqpoint{1.139435in}{0.743451in}}%
\pgfpathlineto{\pgfqpoint{1.122795in}{0.738324in}}%
\pgfpathlineto{\pgfqpoint{1.106213in}{0.734261in}}%
\pgfpathlineto{\pgfqpoint{1.097979in}{0.736318in}}%
\pgfpathlineto{\pgfqpoint{1.089710in}{0.738908in}}%
\pgfpathclose%
\pgfusepath{fill}%
\end{pgfscope}%
\begin{pgfscope}%
\pgfpathrectangle{\pgfqpoint{0.407592in}{0.222904in}}{\pgfqpoint{1.411163in}{1.411163in}}%
\pgfusepath{clip}%
\pgfsetbuttcap%
\pgfsetroundjoin%
\definecolor{currentfill}{rgb}{0.556863,0.654902,0.847059}%
\pgfsetfillcolor{currentfill}%
\pgfsetfillopacity{0.130000}%
\pgfsetlinewidth{0.000000pt}%
\definecolor{currentstroke}{rgb}{0.000000,0.000000,0.000000}%
\pgfsetstrokecolor{currentstroke}%
\pgfsetdash{}{0pt}%
\pgfpathmoveto{\pgfqpoint{1.701811in}{1.143728in}}%
\pgfpathlineto{\pgfqpoint{1.719686in}{1.154992in}}%
\pgfpathlineto{\pgfqpoint{1.737641in}{1.166444in}}%
\pgfpathlineto{\pgfqpoint{1.745093in}{1.177162in}}%
\pgfpathlineto{\pgfqpoint{1.752534in}{1.188087in}}%
\pgfpathlineto{\pgfqpoint{1.734604in}{1.176930in}}%
\pgfpathlineto{\pgfqpoint{1.716752in}{1.165971in}}%
\pgfpathlineto{\pgfqpoint{1.709287in}{1.154744in}}%
\pgfpathlineto{\pgfqpoint{1.701811in}{1.143728in}}%
\pgfpathclose%
\pgfusepath{fill}%
\end{pgfscope}%
\begin{pgfscope}%
\pgfpathrectangle{\pgfqpoint{0.407592in}{0.222904in}}{\pgfqpoint{1.411163in}{1.411163in}}%
\pgfusepath{clip}%
\pgfsetbuttcap%
\pgfsetroundjoin%
\definecolor{currentfill}{rgb}{0.556863,0.654902,0.847059}%
\pgfsetfillcolor{currentfill}%
\pgfsetfillopacity{0.130000}%
\pgfsetlinewidth{0.000000pt}%
\definecolor{currentstroke}{rgb}{0.000000,0.000000,0.000000}%
\pgfsetstrokecolor{currentstroke}%
\pgfsetdash{}{0pt}%
\pgfpathmoveto{\pgfqpoint{0.557189in}{1.054924in}}%
\pgfpathlineto{\pgfqpoint{0.573574in}{1.038091in}}%
\pgfpathlineto{\pgfqpoint{0.589967in}{1.021543in}}%
\pgfpathlineto{\pgfqpoint{0.599279in}{1.018241in}}%
\pgfpathlineto{\pgfqpoint{0.608546in}{1.015186in}}%
\pgfpathlineto{\pgfqpoint{0.592196in}{1.032116in}}%
\pgfpathlineto{\pgfqpoint{0.575854in}{1.049316in}}%
\pgfpathlineto{\pgfqpoint{0.566544in}{1.051996in}}%
\pgfpathlineto{\pgfqpoint{0.557189in}{1.054924in}}%
\pgfpathclose%
\pgfusepath{fill}%
\end{pgfscope}%
\begin{pgfscope}%
\pgfpathrectangle{\pgfqpoint{0.407592in}{0.222904in}}{\pgfqpoint{1.411163in}{1.411163in}}%
\pgfusepath{clip}%
\pgfsetbuttcap%
\pgfsetroundjoin%
\definecolor{currentfill}{rgb}{0.556863,0.654902,0.847059}%
\pgfsetfillcolor{currentfill}%
\pgfsetfillopacity{0.130000}%
\pgfsetlinewidth{0.000000pt}%
\definecolor{currentstroke}{rgb}{0.000000,0.000000,0.000000}%
\pgfsetstrokecolor{currentstroke}%
\pgfsetdash{}{0pt}%
\pgfpathmoveto{\pgfqpoint{1.238873in}{0.776510in}}%
\pgfpathlineto{\pgfqpoint{1.255787in}{0.786123in}}%
\pgfpathlineto{\pgfqpoint{1.272772in}{0.796163in}}%
\pgfpathlineto{\pgfqpoint{1.280746in}{0.799041in}}%
\pgfpathlineto{\pgfqpoint{1.288698in}{0.802519in}}%
\pgfpathlineto{\pgfqpoint{1.271735in}{0.792254in}}%
\pgfpathlineto{\pgfqpoint{1.254845in}{0.782394in}}%
\pgfpathlineto{\pgfqpoint{1.246871in}{0.779121in}}%
\pgfpathlineto{\pgfqpoint{1.238873in}{0.776510in}}%
\pgfpathclose%
\pgfusepath{fill}%
\end{pgfscope}%
\begin{pgfscope}%
\pgfpathrectangle{\pgfqpoint{0.407592in}{0.222904in}}{\pgfqpoint{1.411163in}{1.411163in}}%
\pgfusepath{clip}%
\pgfsetbuttcap%
\pgfsetroundjoin%
\definecolor{currentfill}{rgb}{0.556863,0.654902,0.847059}%
\pgfsetfillcolor{currentfill}%
\pgfsetfillopacity{0.130000}%
\pgfsetlinewidth{0.000000pt}%
\definecolor{currentstroke}{rgb}{0.000000,0.000000,0.000000}%
\pgfsetstrokecolor{currentstroke}%
\pgfsetdash{}{0pt}%
\pgfpathmoveto{\pgfqpoint{1.404341in}{0.879934in}}%
\pgfpathlineto{\pgfqpoint{1.421615in}{0.891481in}}%
\pgfpathlineto{\pgfqpoint{1.438965in}{0.903229in}}%
\pgfpathlineto{\pgfqpoint{1.446728in}{0.910318in}}%
\pgfpathlineto{\pgfqpoint{1.454476in}{0.917843in}}%
\pgfpathlineto{\pgfqpoint{1.437147in}{0.906288in}}%
\pgfpathlineto{\pgfqpoint{1.419893in}{0.894947in}}%
\pgfpathlineto{\pgfqpoint{1.412124in}{0.887205in}}%
\pgfpathlineto{\pgfqpoint{1.404341in}{0.879934in}}%
\pgfpathclose%
\pgfusepath{fill}%
\end{pgfscope}%
\begin{pgfscope}%
\pgfpathrectangle{\pgfqpoint{0.407592in}{0.222904in}}{\pgfqpoint{1.411163in}{1.411163in}}%
\pgfusepath{clip}%
\pgfsetbuttcap%
\pgfsetroundjoin%
\definecolor{currentfill}{rgb}{0.556863,0.654902,0.847059}%
\pgfsetfillcolor{currentfill}%
\pgfsetfillopacity{0.130000}%
\pgfsetlinewidth{0.000000pt}%
\definecolor{currentstroke}{rgb}{0.000000,0.000000,0.000000}%
\pgfsetstrokecolor{currentstroke}%
\pgfsetdash{}{0pt}%
\pgfpathmoveto{\pgfqpoint{1.520136in}{0.974725in}}%
\pgfpathlineto{\pgfqpoint{1.537654in}{0.986351in}}%
\pgfpathlineto{\pgfqpoint{1.555250in}{0.998161in}}%
\pgfpathlineto{\pgfqpoint{1.562888in}{1.007089in}}%
\pgfpathlineto{\pgfqpoint{1.570514in}{1.016342in}}%
\pgfpathlineto{\pgfqpoint{1.552939in}{1.004816in}}%
\pgfpathlineto{\pgfqpoint{1.535443in}{0.993491in}}%
\pgfpathlineto{\pgfqpoint{1.527795in}{0.983937in}}%
\pgfpathlineto{\pgfqpoint{1.520136in}{0.974725in}}%
\pgfpathclose%
\pgfusepath{fill}%
\end{pgfscope}%
\begin{pgfscope}%
\pgfpathrectangle{\pgfqpoint{0.407592in}{0.222904in}}{\pgfqpoint{1.411163in}{1.411163in}}%
\pgfusepath{clip}%
\pgfsetbuttcap%
\pgfsetroundjoin%
\definecolor{currentfill}{rgb}{0.556863,0.654902,0.847059}%
\pgfsetfillcolor{currentfill}%
\pgfsetfillopacity{0.130000}%
\pgfsetlinewidth{0.000000pt}%
\definecolor{currentstroke}{rgb}{0.000000,0.000000,0.000000}%
\pgfsetstrokecolor{currentstroke}%
\pgfsetdash{}{0pt}%
\pgfpathmoveto{\pgfqpoint{1.818197in}{1.255429in}}%
\pgfpathlineto{\pgfqpoint{1.836319in}{1.266580in}}%
\pgfpathlineto{\pgfqpoint{1.843671in}{1.277999in}}%
\pgfpathlineto{\pgfqpoint{1.825562in}{1.266988in}}%
\pgfusepath{fill}%
\end{pgfscope}%
\begin{pgfscope}%
\pgfpathrectangle{\pgfqpoint{0.407592in}{0.222904in}}{\pgfqpoint{1.411163in}{1.411163in}}%
\pgfusepath{clip}%
\pgfsetbuttcap%
\pgfsetroundjoin%
\definecolor{currentfill}{rgb}{0.556863,0.654902,0.847059}%
\pgfsetfillcolor{currentfill}%
\pgfsetfillopacity{0.130000}%
\pgfsetlinewidth{0.000000pt}%
\definecolor{currentstroke}{rgb}{0.000000,0.000000,0.000000}%
\pgfsetstrokecolor{currentstroke}%
\pgfsetdash{}{0pt}%
\pgfpathmoveto{\pgfqpoint{1.039888in}{0.742986in}}%
\pgfpathlineto{\pgfqpoint{1.056454in}{0.743637in}}%
\pgfpathlineto{\pgfqpoint{1.073067in}{0.745476in}}%
\pgfpathlineto{\pgfqpoint{1.081406in}{0.741977in}}%
\pgfpathlineto{\pgfqpoint{1.089710in}{0.738908in}}%
\pgfpathlineto{\pgfqpoint{1.073141in}{0.736568in}}%
\pgfpathlineto{\pgfqpoint{1.056620in}{0.735487in}}%
\pgfpathlineto{\pgfqpoint{1.048272in}{0.739031in}}%
\pgfpathlineto{\pgfqpoint{1.039888in}{0.742986in}}%
\pgfpathclose%
\pgfusepath{fill}%
\end{pgfscope}%
\begin{pgfscope}%
\pgfpathrectangle{\pgfqpoint{0.407592in}{0.222904in}}{\pgfqpoint{1.411163in}{1.411163in}}%
\pgfusepath{clip}%
\pgfsetbuttcap%
\pgfsetroundjoin%
\definecolor{currentfill}{rgb}{0.556863,0.654902,0.847059}%
\pgfsetfillcolor{currentfill}%
\pgfsetfillopacity{0.130000}%
\pgfsetlinewidth{0.000000pt}%
\definecolor{currentstroke}{rgb}{0.000000,0.000000,0.000000}%
\pgfsetstrokecolor{currentstroke}%
\pgfsetdash{}{0pt}%
\pgfpathmoveto{\pgfqpoint{0.706837in}{0.920904in}}%
\pgfpathlineto{\pgfqpoint{0.723253in}{0.906731in}}%
\pgfpathlineto{\pgfqpoint{0.739678in}{0.893106in}}%
\pgfpathlineto{\pgfqpoint{0.748726in}{0.888535in}}%
\pgfpathlineto{\pgfqpoint{0.757730in}{0.884208in}}%
\pgfpathlineto{\pgfqpoint{0.741349in}{0.898310in}}%
\pgfpathlineto{\pgfqpoint{0.724979in}{0.912959in}}%
\pgfpathlineto{\pgfqpoint{0.715930in}{0.916803in}}%
\pgfpathlineto{\pgfqpoint{0.706837in}{0.920904in}}%
\pgfpathclose%
\pgfusepath{fill}%
\end{pgfscope}%
\begin{pgfscope}%
\pgfpathrectangle{\pgfqpoint{0.407592in}{0.222904in}}{\pgfqpoint{1.411163in}{1.411163in}}%
\pgfusepath{clip}%
\pgfsetbuttcap%
\pgfsetroundjoin%
\definecolor{currentfill}{rgb}{0.556863,0.654902,0.847059}%
\pgfsetfillcolor{currentfill}%
\pgfsetfillopacity{0.130000}%
\pgfsetlinewidth{0.000000pt}%
\definecolor{currentstroke}{rgb}{0.000000,0.000000,0.000000}%
\pgfsetstrokecolor{currentstroke}%
\pgfsetdash{}{0pt}%
\pgfpathmoveto{\pgfqpoint{1.636188in}{1.078777in}}%
\pgfpathlineto{\pgfqpoint{1.653952in}{1.090268in}}%
\pgfpathlineto{\pgfqpoint{1.671795in}{1.101941in}}%
\pgfpathlineto{\pgfqpoint{1.679316in}{1.112029in}}%
\pgfpathlineto{\pgfqpoint{1.686825in}{1.122362in}}%
\pgfpathlineto{\pgfqpoint{1.669006in}{1.110985in}}%
\pgfpathlineto{\pgfqpoint{1.651264in}{1.099804in}}%
\pgfpathlineto{\pgfqpoint{1.643732in}{1.089165in}}%
\pgfpathlineto{\pgfqpoint{1.636188in}{1.078777in}}%
\pgfpathclose%
\pgfusepath{fill}%
\end{pgfscope}%
\begin{pgfscope}%
\pgfpathrectangle{\pgfqpoint{0.407592in}{0.222904in}}{\pgfqpoint{1.411163in}{1.411163in}}%
\pgfusepath{clip}%
\pgfsetbuttcap%
\pgfsetroundjoin%
\definecolor{currentfill}{rgb}{0.556863,0.654902,0.847059}%
\pgfsetfillcolor{currentfill}%
\pgfsetfillopacity{0.130000}%
\pgfsetlinewidth{0.000000pt}%
\definecolor{currentstroke}{rgb}{0.000000,0.000000,0.000000}%
\pgfsetstrokecolor{currentstroke}%
\pgfsetdash{}{0pt}%
\pgfpathmoveto{\pgfqpoint{0.439960in}{1.165411in}}%
\pgfpathlineto{\pgfqpoint{0.456335in}{1.147569in}}%
\pgfpathlineto{\pgfqpoint{0.472720in}{1.129914in}}%
\pgfpathlineto{\pgfqpoint{0.482251in}{1.127185in}}%
\pgfpathlineto{\pgfqpoint{0.491736in}{1.124683in}}%
\pgfpathlineto{\pgfqpoint{0.475395in}{1.142637in}}%
\pgfpathlineto{\pgfqpoint{0.459061in}{1.160764in}}%
\pgfpathlineto{\pgfqpoint{0.449533in}{1.162975in}}%
\pgfpathlineto{\pgfqpoint{0.439960in}{1.165411in}}%
\pgfpathclose%
\pgfusepath{fill}%
\end{pgfscope}%
\begin{pgfscope}%
\pgfpathrectangle{\pgfqpoint{0.407592in}{0.222904in}}{\pgfqpoint{1.411163in}{1.411163in}}%
\pgfusepath{clip}%
\pgfsetbuttcap%
\pgfsetroundjoin%
\definecolor{currentfill}{rgb}{0.556863,0.654902,0.847059}%
\pgfsetfillcolor{currentfill}%
\pgfsetfillopacity{0.130000}%
\pgfsetlinewidth{0.000000pt}%
\definecolor{currentstroke}{rgb}{0.000000,0.000000,0.000000}%
\pgfsetstrokecolor{currentstroke}%
\pgfsetdash{}{0pt}%
\pgfpathmoveto{\pgfqpoint{1.288698in}{0.802519in}}%
\pgfpathlineto{\pgfqpoint{1.305734in}{0.813134in}}%
\pgfpathlineto{\pgfqpoint{1.322844in}{0.824057in}}%
\pgfpathlineto{\pgfqpoint{1.330753in}{0.828253in}}%
\pgfpathlineto{\pgfqpoint{1.338642in}{0.833004in}}%
\pgfpathlineto{\pgfqpoint{1.321553in}{0.822006in}}%
\pgfpathlineto{\pgfqpoint{1.304539in}{0.811307in}}%
\pgfpathlineto{\pgfqpoint{1.296628in}{0.806605in}}%
\pgfpathlineto{\pgfqpoint{1.288698in}{0.802519in}}%
\pgfpathclose%
\pgfusepath{fill}%
\end{pgfscope}%
\begin{pgfscope}%
\pgfpathrectangle{\pgfqpoint{0.407592in}{0.222904in}}{\pgfqpoint{1.411163in}{1.411163in}}%
\pgfusepath{clip}%
\pgfsetbuttcap%
\pgfsetroundjoin%
\definecolor{currentfill}{rgb}{0.556863,0.654902,0.847059}%
\pgfsetfillcolor{currentfill}%
\pgfsetfillopacity{0.130000}%
\pgfsetlinewidth{0.000000pt}%
\definecolor{currentstroke}{rgb}{0.000000,0.000000,0.000000}%
\pgfsetstrokecolor{currentstroke}%
\pgfsetdash{}{0pt}%
\pgfpathmoveto{\pgfqpoint{0.906739in}{0.786422in}}%
\pgfpathlineto{\pgfqpoint{0.923206in}{0.779499in}}%
\pgfpathlineto{\pgfqpoint{0.939697in}{0.773708in}}%
\pgfpathlineto{\pgfqpoint{0.948336in}{0.768266in}}%
\pgfpathlineto{\pgfqpoint{0.956935in}{0.763069in}}%
\pgfpathlineto{\pgfqpoint{0.940494in}{0.768983in}}%
\pgfpathlineto{\pgfqpoint{0.924077in}{0.776122in}}%
\pgfpathlineto{\pgfqpoint{0.915428in}{0.781145in}}%
\pgfpathlineto{\pgfqpoint{0.906739in}{0.786422in}}%
\pgfpathclose%
\pgfusepath{fill}%
\end{pgfscope}%
\begin{pgfscope}%
\pgfpathrectangle{\pgfqpoint{0.407592in}{0.222904in}}{\pgfqpoint{1.411163in}{1.411163in}}%
\pgfusepath{clip}%
\pgfsetbuttcap%
\pgfsetroundjoin%
\definecolor{currentfill}{rgb}{0.556863,0.654902,0.847059}%
\pgfsetfillcolor{currentfill}%
\pgfsetfillopacity{0.130000}%
\pgfsetlinewidth{0.000000pt}%
\definecolor{currentstroke}{rgb}{0.000000,0.000000,0.000000}%
\pgfsetstrokecolor{currentstroke}%
\pgfsetdash{}{0pt}%
\pgfpathmoveto{\pgfqpoint{0.589967in}{1.021543in}}%
\pgfpathlineto{\pgfqpoint{0.606370in}{1.005307in}}%
\pgfpathlineto{\pgfqpoint{0.622781in}{0.989409in}}%
\pgfpathlineto{\pgfqpoint{0.632050in}{0.985705in}}%
\pgfpathlineto{\pgfqpoint{0.641274in}{0.982244in}}%
\pgfpathlineto{\pgfqpoint{0.624906in}{0.998553in}}%
\pgfpathlineto{\pgfqpoint{0.608546in}{1.015186in}}%
\pgfpathlineto{\pgfqpoint{0.599279in}{1.018241in}}%
\pgfpathlineto{\pgfqpoint{0.589967in}{1.021543in}}%
\pgfpathclose%
\pgfusepath{fill}%
\end{pgfscope}%
\begin{pgfscope}%
\pgfpathrectangle{\pgfqpoint{0.407592in}{0.222904in}}{\pgfqpoint{1.411163in}{1.411163in}}%
\pgfusepath{clip}%
\pgfsetbuttcap%
\pgfsetroundjoin%
\definecolor{currentfill}{rgb}{0.556863,0.654902,0.847059}%
\pgfsetfillcolor{currentfill}%
\pgfsetfillopacity{0.130000}%
\pgfsetlinewidth{0.000000pt}%
\definecolor{currentstroke}{rgb}{0.000000,0.000000,0.000000}%
\pgfsetstrokecolor{currentstroke}%
\pgfsetdash{}{0pt}%
\pgfpathmoveto{\pgfqpoint{0.823366in}{0.834341in}}%
\pgfpathlineto{\pgfqpoint{0.839808in}{0.823794in}}%
\pgfpathlineto{\pgfqpoint{0.856266in}{0.814128in}}%
\pgfpathlineto{\pgfqpoint{0.865087in}{0.808704in}}%
\pgfpathlineto{\pgfqpoint{0.873867in}{0.803505in}}%
\pgfpathlineto{\pgfqpoint{0.857457in}{0.813559in}}%
\pgfpathlineto{\pgfqpoint{0.841063in}{0.824538in}}%
\pgfpathlineto{\pgfqpoint{0.832235in}{0.829316in}}%
\pgfpathlineto{\pgfqpoint{0.823366in}{0.834341in}}%
\pgfpathclose%
\pgfusepath{fill}%
\end{pgfscope}%
\begin{pgfscope}%
\pgfpathrectangle{\pgfqpoint{0.407592in}{0.222904in}}{\pgfqpoint{1.411163in}{1.411163in}}%
\pgfusepath{clip}%
\pgfsetbuttcap%
\pgfsetroundjoin%
\definecolor{currentfill}{rgb}{0.556863,0.654902,0.847059}%
\pgfsetfillcolor{currentfill}%
\pgfsetfillopacity{0.130000}%
\pgfsetlinewidth{0.000000pt}%
\definecolor{currentstroke}{rgb}{0.000000,0.000000,0.000000}%
\pgfsetstrokecolor{currentstroke}%
\pgfsetdash{}{0pt}%
\pgfpathmoveto{\pgfqpoint{1.454476in}{0.917843in}}%
\pgfpathlineto{\pgfqpoint{1.471883in}{0.929595in}}%
\pgfpathlineto{\pgfqpoint{1.489366in}{0.941528in}}%
\pgfpathlineto{\pgfqpoint{1.497079in}{0.949258in}}%
\pgfpathlineto{\pgfqpoint{1.504778in}{0.957375in}}%
\pgfpathlineto{\pgfqpoint{1.487315in}{0.945669in}}%
\pgfpathlineto{\pgfqpoint{1.469930in}{0.934159in}}%
\pgfpathlineto{\pgfqpoint{1.462210in}{0.925794in}}%
\pgfpathlineto{\pgfqpoint{1.454476in}{0.917843in}}%
\pgfpathclose%
\pgfusepath{fill}%
\end{pgfscope}%
\begin{pgfscope}%
\pgfpathrectangle{\pgfqpoint{0.407592in}{0.222904in}}{\pgfqpoint{1.411163in}{1.411163in}}%
\pgfusepath{clip}%
\pgfsetbuttcap%
\pgfsetroundjoin%
\definecolor{currentfill}{rgb}{0.556863,0.654902,0.847059}%
\pgfsetfillcolor{currentfill}%
\pgfsetfillopacity{0.130000}%
\pgfsetlinewidth{0.000000pt}%
\definecolor{currentstroke}{rgb}{0.000000,0.000000,0.000000}%
\pgfsetstrokecolor{currentstroke}%
\pgfsetdash{}{0pt}%
\pgfpathmoveto{\pgfqpoint{1.752534in}{1.188087in}}%
\pgfpathlineto{\pgfqpoint{1.770545in}{1.199433in}}%
\pgfpathlineto{\pgfqpoint{1.788636in}{1.210959in}}%
\pgfpathlineto{\pgfqpoint{1.796042in}{1.221799in}}%
\pgfpathlineto{\pgfqpoint{1.803438in}{1.232829in}}%
\pgfpathlineto{\pgfqpoint{1.785372in}{1.221583in}}%
\pgfpathlineto{\pgfqpoint{1.767385in}{1.210528in}}%
\pgfpathlineto{\pgfqpoint{1.759965in}{1.199211in}}%
\pgfpathlineto{\pgfqpoint{1.752534in}{1.188087in}}%
\pgfpathclose%
\pgfusepath{fill}%
\end{pgfscope}%
\begin{pgfscope}%
\pgfpathrectangle{\pgfqpoint{0.407592in}{0.222904in}}{\pgfqpoint{1.411163in}{1.411163in}}%
\pgfusepath{clip}%
\pgfsetbuttcap%
\pgfsetroundjoin%
\definecolor{currentfill}{rgb}{0.556863,0.654902,0.847059}%
\pgfsetfillcolor{currentfill}%
\pgfsetfillopacity{0.130000}%
\pgfsetlinewidth{0.000000pt}%
\definecolor{currentstroke}{rgb}{0.000000,0.000000,0.000000}%
\pgfsetstrokecolor{currentstroke}%
\pgfsetdash{}{0pt}%
\pgfpathmoveto{\pgfqpoint{0.989901in}{0.755014in}}%
\pgfpathlineto{\pgfqpoint{1.006434in}{0.752873in}}%
\pgfpathlineto{\pgfqpoint{1.023005in}{0.751955in}}%
\pgfpathlineto{\pgfqpoint{1.031465in}{0.747308in}}%
\pgfpathlineto{\pgfqpoint{1.039888in}{0.742986in}}%
\pgfpathlineto{\pgfqpoint{1.023365in}{0.743601in}}%
\pgfpathlineto{\pgfqpoint{1.006881in}{0.745540in}}%
\pgfpathlineto{\pgfqpoint{0.998411in}{0.750122in}}%
\pgfpathlineto{\pgfqpoint{0.989901in}{0.755014in}}%
\pgfpathclose%
\pgfusepath{fill}%
\end{pgfscope}%
\begin{pgfscope}%
\pgfpathrectangle{\pgfqpoint{0.407592in}{0.222904in}}{\pgfqpoint{1.411163in}{1.411163in}}%
\pgfusepath{clip}%
\pgfsetbuttcap%
\pgfsetroundjoin%
\definecolor{currentfill}{rgb}{0.556863,0.654902,0.847059}%
\pgfsetfillcolor{currentfill}%
\pgfsetfillopacity{0.130000}%
\pgfsetlinewidth{0.000000pt}%
\definecolor{currentstroke}{rgb}{0.000000,0.000000,0.000000}%
\pgfsetstrokecolor{currentstroke}%
\pgfsetdash{}{0pt}%
\pgfpathmoveto{\pgfqpoint{1.570514in}{1.016342in}}%
\pgfpathlineto{\pgfqpoint{1.588166in}{1.028055in}}%
\pgfpathlineto{\pgfqpoint{1.605897in}{1.039942in}}%
\pgfpathlineto{\pgfqpoint{1.613488in}{1.049225in}}%
\pgfpathlineto{\pgfqpoint{1.621066in}{1.058798in}}%
\pgfpathlineto{\pgfqpoint{1.603358in}{1.047194in}}%
\pgfpathlineto{\pgfqpoint{1.585728in}{1.035780in}}%
\pgfpathlineto{\pgfqpoint{1.578127in}{1.025909in}}%
\pgfpathlineto{\pgfqpoint{1.570514in}{1.016342in}}%
\pgfpathclose%
\pgfusepath{fill}%
\end{pgfscope}%
\begin{pgfscope}%
\pgfpathrectangle{\pgfqpoint{0.407592in}{0.222904in}}{\pgfqpoint{1.411163in}{1.411163in}}%
\pgfusepath{clip}%
\pgfsetbuttcap%
\pgfsetroundjoin%
\definecolor{currentfill}{rgb}{0.556863,0.654902,0.847059}%
\pgfsetfillcolor{currentfill}%
\pgfsetfillopacity{0.130000}%
\pgfsetlinewidth{0.000000pt}%
\definecolor{currentstroke}{rgb}{0.000000,0.000000,0.000000}%
\pgfsetstrokecolor{currentstroke}%
\pgfsetdash{}{0pt}%
\pgfpathmoveto{\pgfqpoint{0.472720in}{1.129914in}}%
\pgfpathlineto{\pgfqpoint{0.489113in}{1.112460in}}%
\pgfpathlineto{\pgfqpoint{0.505514in}{1.095225in}}%
\pgfpathlineto{\pgfqpoint{0.515003in}{1.092177in}}%
\pgfpathlineto{\pgfqpoint{0.524446in}{1.089358in}}%
\pgfpathlineto{\pgfqpoint{0.508087in}{1.106918in}}%
\pgfpathlineto{\pgfqpoint{0.491736in}{1.124683in}}%
\pgfpathlineto{\pgfqpoint{0.482251in}{1.127185in}}%
\pgfpathlineto{\pgfqpoint{0.472720in}{1.129914in}}%
\pgfpathclose%
\pgfusepath{fill}%
\end{pgfscope}%
\begin{pgfscope}%
\pgfpathrectangle{\pgfqpoint{0.407592in}{0.222904in}}{\pgfqpoint{1.411163in}{1.411163in}}%
\pgfusepath{clip}%
\pgfsetbuttcap%
\pgfsetroundjoin%
\definecolor{currentfill}{rgb}{0.556863,0.654902,0.847059}%
\pgfsetfillcolor{currentfill}%
\pgfsetfillopacity{0.130000}%
\pgfsetlinewidth{0.000000pt}%
\definecolor{currentstroke}{rgb}{0.000000,0.000000,0.000000}%
\pgfsetstrokecolor{currentstroke}%
\pgfsetdash{}{0pt}%
\pgfpathmoveto{\pgfqpoint{1.338642in}{0.833004in}}%
\pgfpathlineto{\pgfqpoint{1.355806in}{0.844265in}}%
\pgfpathlineto{\pgfqpoint{1.373045in}{0.855762in}}%
\pgfpathlineto{\pgfqpoint{1.380894in}{0.861056in}}%
\pgfpathlineto{\pgfqpoint{1.388726in}{0.866852in}}%
\pgfpathlineto{\pgfqpoint{1.371508in}{0.855394in}}%
\pgfpathlineto{\pgfqpoint{1.354365in}{0.844172in}}%
\pgfpathlineto{\pgfqpoint{1.346512in}{0.838311in}}%
\pgfpathlineto{\pgfqpoint{1.338642in}{0.833004in}}%
\pgfpathclose%
\pgfusepath{fill}%
\end{pgfscope}%
\begin{pgfscope}%
\pgfpathrectangle{\pgfqpoint{0.407592in}{0.222904in}}{\pgfqpoint{1.411163in}{1.411163in}}%
\pgfusepath{clip}%
\pgfsetbuttcap%
\pgfsetroundjoin%
\definecolor{currentfill}{rgb}{0.556863,0.654902,0.847059}%
\pgfsetfillcolor{currentfill}%
\pgfsetfillopacity{0.130000}%
\pgfsetlinewidth{0.000000pt}%
\definecolor{currentstroke}{rgb}{0.000000,0.000000,0.000000}%
\pgfsetstrokecolor{currentstroke}%
\pgfsetdash{}{0pt}%
\pgfpathmoveto{\pgfqpoint{0.739678in}{0.893106in}}%
\pgfpathlineto{\pgfqpoint{0.756115in}{0.880077in}}%
\pgfpathlineto{\pgfqpoint{0.772564in}{0.867691in}}%
\pgfpathlineto{\pgfqpoint{0.781565in}{0.862659in}}%
\pgfpathlineto{\pgfqpoint{0.790524in}{0.857854in}}%
\pgfpathlineto{\pgfqpoint{0.774121in}{0.870705in}}%
\pgfpathlineto{\pgfqpoint{0.757730in}{0.884208in}}%
\pgfpathlineto{\pgfqpoint{0.748726in}{0.888535in}}%
\pgfpathlineto{\pgfqpoint{0.739678in}{0.893106in}}%
\pgfpathclose%
\pgfusepath{fill}%
\end{pgfscope}%
\begin{pgfscope}%
\pgfpathrectangle{\pgfqpoint{0.407592in}{0.222904in}}{\pgfqpoint{1.411163in}{1.411163in}}%
\pgfusepath{clip}%
\pgfsetbuttcap%
\pgfsetroundjoin%
\definecolor{currentfill}{rgb}{0.556863,0.654902,0.847059}%
\pgfsetfillcolor{currentfill}%
\pgfsetfillopacity{0.130000}%
\pgfsetlinewidth{0.000000pt}%
\definecolor{currentstroke}{rgb}{0.000000,0.000000,0.000000}%
\pgfsetstrokecolor{currentstroke}%
\pgfsetdash{}{0pt}%
\pgfpathmoveto{\pgfqpoint{0.622781in}{0.989409in}}%
\pgfpathlineto{\pgfqpoint{0.639202in}{0.973878in}}%
\pgfpathlineto{\pgfqpoint{0.655632in}{0.958746in}}%
\pgfpathlineto{\pgfqpoint{0.664856in}{0.954617in}}%
\pgfpathlineto{\pgfqpoint{0.674037in}{0.950724in}}%
\pgfpathlineto{\pgfqpoint{0.657651in}{0.966290in}}%
\pgfpathlineto{\pgfqpoint{0.641274in}{0.982244in}}%
\pgfpathlineto{\pgfqpoint{0.632050in}{0.985705in}}%
\pgfpathlineto{\pgfqpoint{0.622781in}{0.989409in}}%
\pgfpathclose%
\pgfusepath{fill}%
\end{pgfscope}%
\begin{pgfscope}%
\pgfpathrectangle{\pgfqpoint{0.407592in}{0.222904in}}{\pgfqpoint{1.411163in}{1.411163in}}%
\pgfusepath{clip}%
\pgfsetbuttcap%
\pgfsetroundjoin%
\definecolor{currentfill}{rgb}{0.556863,0.654902,0.847059}%
\pgfsetfillcolor{currentfill}%
\pgfsetfillopacity{0.130000}%
\pgfsetlinewidth{0.000000pt}%
\definecolor{currentstroke}{rgb}{0.000000,0.000000,0.000000}%
\pgfsetstrokecolor{currentstroke}%
\pgfsetdash{}{0pt}%
\pgfpathmoveto{\pgfqpoint{1.803438in}{1.232829in}}%
\pgfpathlineto{\pgfqpoint{1.821584in}{1.244256in}}%
\pgfpathlineto{\pgfqpoint{1.828957in}{1.255330in}}%
\pgfpathlineto{\pgfqpoint{1.836319in}{1.266580in}}%
\pgfpathlineto{\pgfqpoint{1.818197in}{1.255429in}}%
\pgfpathlineto{\pgfqpoint{1.810823in}{1.244041in}}%
\pgfusepath{fill}%
\end{pgfscope}%
\begin{pgfscope}%
\pgfpathrectangle{\pgfqpoint{0.407592in}{0.222904in}}{\pgfqpoint{1.411163in}{1.411163in}}%
\pgfusepath{clip}%
\pgfsetbuttcap%
\pgfsetroundjoin%
\definecolor{currentfill}{rgb}{0.556863,0.654902,0.847059}%
\pgfsetfillcolor{currentfill}%
\pgfsetfillopacity{0.130000}%
\pgfsetlinewidth{0.000000pt}%
\definecolor{currentstroke}{rgb}{0.000000,0.000000,0.000000}%
\pgfsetstrokecolor{currentstroke}%
\pgfsetdash{}{0pt}%
\pgfpathmoveto{\pgfqpoint{1.686825in}{1.122362in}}%
\pgfpathlineto{\pgfqpoint{1.704724in}{1.133923in}}%
\pgfpathlineto{\pgfqpoint{1.722703in}{1.145660in}}%
\pgfpathlineto{\pgfqpoint{1.730178in}{1.155941in}}%
\pgfpathlineto{\pgfqpoint{1.737641in}{1.166444in}}%
\pgfpathlineto{\pgfqpoint{1.719686in}{1.154992in}}%
\pgfpathlineto{\pgfqpoint{1.701811in}{1.143728in}}%
\pgfpathlineto{\pgfqpoint{1.694323in}{1.132931in}}%
\pgfpathlineto{\pgfqpoint{1.686825in}{1.122362in}}%
\pgfpathclose%
\pgfusepath{fill}%
\end{pgfscope}%
\begin{pgfscope}%
\pgfpathrectangle{\pgfqpoint{0.407592in}{0.222904in}}{\pgfqpoint{1.411163in}{1.411163in}}%
\pgfusepath{clip}%
\pgfsetbuttcap%
\pgfsetroundjoin%
\definecolor{currentfill}{rgb}{0.556863,0.654902,0.847059}%
\pgfsetfillcolor{currentfill}%
\pgfsetfillopacity{0.130000}%
\pgfsetlinewidth{0.000000pt}%
\definecolor{currentstroke}{rgb}{0.000000,0.000000,0.000000}%
\pgfsetstrokecolor{currentstroke}%
\pgfsetdash{}{0pt}%
\pgfpathmoveto{\pgfqpoint{0.505514in}{1.095225in}}%
\pgfpathlineto{\pgfqpoint{0.521925in}{1.078226in}}%
\pgfpathlineto{\pgfqpoint{0.538344in}{1.061482in}}%
\pgfpathlineto{\pgfqpoint{0.547789in}{1.058089in}}%
\pgfpathlineto{\pgfqpoint{0.557189in}{1.054924in}}%
\pgfpathlineto{\pgfqpoint{0.540813in}{1.072020in}}%
\pgfpathlineto{\pgfqpoint{0.524446in}{1.089358in}}%
\pgfpathlineto{\pgfqpoint{0.515003in}{1.092177in}}%
\pgfpathlineto{\pgfqpoint{0.505514in}{1.095225in}}%
\pgfpathclose%
\pgfusepath{fill}%
\end{pgfscope}%
\begin{pgfscope}%
\pgfpathrectangle{\pgfqpoint{0.407592in}{0.222904in}}{\pgfqpoint{1.411163in}{1.411163in}}%
\pgfusepath{clip}%
\pgfsetbuttcap%
\pgfsetroundjoin%
\definecolor{currentfill}{rgb}{0.556863,0.654902,0.847059}%
\pgfsetfillcolor{currentfill}%
\pgfsetfillopacity{0.130000}%
\pgfsetlinewidth{0.000000pt}%
\definecolor{currentstroke}{rgb}{0.000000,0.000000,0.000000}%
\pgfsetstrokecolor{currentstroke}%
\pgfsetdash{}{0pt}%
\pgfpathmoveto{\pgfqpoint{1.504778in}{0.957375in}}%
\pgfpathlineto{\pgfqpoint{1.522318in}{0.969261in}}%
\pgfpathlineto{\pgfqpoint{1.539936in}{0.981316in}}%
\pgfpathlineto{\pgfqpoint{1.547600in}{0.989566in}}%
\pgfpathlineto{\pgfqpoint{1.555250in}{0.998161in}}%
\pgfpathlineto{\pgfqpoint{1.537654in}{0.986351in}}%
\pgfpathlineto{\pgfqpoint{1.520136in}{0.974725in}}%
\pgfpathlineto{\pgfqpoint{1.512463in}{0.965867in}}%
\pgfpathlineto{\pgfqpoint{1.504778in}{0.957375in}}%
\pgfpathclose%
\pgfusepath{fill}%
\end{pgfscope}%
\begin{pgfscope}%
\pgfpathrectangle{\pgfqpoint{0.407592in}{0.222904in}}{\pgfqpoint{1.411163in}{1.411163in}}%
\pgfusepath{clip}%
\pgfsetbuttcap%
\pgfsetroundjoin%
\definecolor{currentfill}{rgb}{0.556863,0.654902,0.847059}%
\pgfsetfillcolor{currentfill}%
\pgfsetfillopacity{0.130000}%
\pgfsetlinewidth{0.000000pt}%
\definecolor{currentstroke}{rgb}{0.000000,0.000000,0.000000}%
\pgfsetstrokecolor{currentstroke}%
\pgfsetdash{}{0pt}%
\pgfpathmoveto{\pgfqpoint{1.388726in}{0.866852in}}%
\pgfpathlineto{\pgfqpoint{1.406021in}{0.878524in}}%
\pgfpathlineto{\pgfqpoint{1.423392in}{0.890388in}}%
\pgfpathlineto{\pgfqpoint{1.431187in}{0.896583in}}%
\pgfpathlineto{\pgfqpoint{1.438965in}{0.903229in}}%
\pgfpathlineto{\pgfqpoint{1.421615in}{0.891481in}}%
\pgfpathlineto{\pgfqpoint{1.404341in}{0.879934in}}%
\pgfpathlineto{\pgfqpoint{1.396542in}{0.873147in}}%
\pgfpathlineto{\pgfqpoint{1.388726in}{0.866852in}}%
\pgfpathclose%
\pgfusepath{fill}%
\end{pgfscope}%
\begin{pgfscope}%
\pgfpathrectangle{\pgfqpoint{0.407592in}{0.222904in}}{\pgfqpoint{1.411163in}{1.411163in}}%
\pgfusepath{clip}%
\pgfsetbuttcap%
\pgfsetroundjoin%
\definecolor{currentfill}{rgb}{0.556863,0.654902,0.847059}%
\pgfsetfillcolor{currentfill}%
\pgfsetfillopacity{0.130000}%
\pgfsetlinewidth{0.000000pt}%
\definecolor{currentstroke}{rgb}{0.000000,0.000000,0.000000}%
\pgfsetstrokecolor{currentstroke}%
\pgfsetdash{}{0pt}%
\pgfpathmoveto{\pgfqpoint{0.856266in}{0.814128in}}%
\pgfpathlineto{\pgfqpoint{0.872742in}{0.805395in}}%
\pgfpathlineto{\pgfqpoint{0.889237in}{0.797644in}}%
\pgfpathlineto{\pgfqpoint{0.898008in}{0.791929in}}%
\pgfpathlineto{\pgfqpoint{0.906739in}{0.786422in}}%
\pgfpathlineto{\pgfqpoint{0.890293in}{0.794441in}}%
\pgfpathlineto{\pgfqpoint{0.873867in}{0.803505in}}%
\pgfpathlineto{\pgfqpoint{0.865087in}{0.808704in}}%
\pgfpathlineto{\pgfqpoint{0.856266in}{0.814128in}}%
\pgfpathclose%
\pgfusepath{fill}%
\end{pgfscope}%
\begin{pgfscope}%
\pgfpathrectangle{\pgfqpoint{0.407592in}{0.222904in}}{\pgfqpoint{1.411163in}{1.411163in}}%
\pgfusepath{clip}%
\pgfsetbuttcap%
\pgfsetroundjoin%
\definecolor{currentfill}{rgb}{0.556863,0.654902,0.847059}%
\pgfsetfillcolor{currentfill}%
\pgfsetfillopacity{0.130000}%
\pgfsetlinewidth{0.000000pt}%
\definecolor{currentstroke}{rgb}{0.000000,0.000000,0.000000}%
\pgfsetstrokecolor{currentstroke}%
\pgfsetdash{}{0pt}%
\pgfpathmoveto{\pgfqpoint{1.172900in}{0.756403in}}%
\pgfpathlineto{\pgfqpoint{1.189731in}{0.764011in}}%
\pgfpathlineto{\pgfqpoint{1.206629in}{0.772249in}}%
\pgfpathlineto{\pgfqpoint{1.214730in}{0.772440in}}%
\pgfpathlineto{\pgfqpoint{1.222804in}{0.773194in}}%
\pgfpathlineto{\pgfqpoint{1.205935in}{0.764469in}}%
\pgfpathlineto{\pgfqpoint{1.189135in}{0.756344in}}%
\pgfpathlineto{\pgfqpoint{1.181032in}{0.756073in}}%
\pgfpathlineto{\pgfqpoint{1.172900in}{0.756403in}}%
\pgfpathclose%
\pgfusepath{fill}%
\end{pgfscope}%
\begin{pgfscope}%
\pgfpathrectangle{\pgfqpoint{0.407592in}{0.222904in}}{\pgfqpoint{1.411163in}{1.411163in}}%
\pgfusepath{clip}%
\pgfsetbuttcap%
\pgfsetroundjoin%
\definecolor{currentfill}{rgb}{0.556863,0.654902,0.847059}%
\pgfsetfillcolor{currentfill}%
\pgfsetfillopacity{0.130000}%
\pgfsetlinewidth{0.000000pt}%
\definecolor{currentstroke}{rgb}{0.000000,0.000000,0.000000}%
\pgfsetstrokecolor{currentstroke}%
\pgfsetdash{}{0pt}%
\pgfpathmoveto{\pgfqpoint{1.123009in}{0.746937in}}%
\pgfpathlineto{\pgfqpoint{1.139746in}{0.752406in}}%
\pgfpathlineto{\pgfqpoint{1.156545in}{0.758707in}}%
\pgfpathlineto{\pgfqpoint{1.164738in}{0.757294in}}%
\pgfpathlineto{\pgfqpoint{1.172900in}{0.756403in}}%
\pgfpathlineto{\pgfqpoint{1.156136in}{0.749515in}}%
\pgfpathlineto{\pgfqpoint{1.139435in}{0.743451in}}%
\pgfpathlineto{\pgfqpoint{1.131239in}{0.744928in}}%
\pgfpathlineto{\pgfqpoint{1.123009in}{0.746937in}}%
\pgfpathclose%
\pgfusepath{fill}%
\end{pgfscope}%
\begin{pgfscope}%
\pgfpathrectangle{\pgfqpoint{0.407592in}{0.222904in}}{\pgfqpoint{1.411163in}{1.411163in}}%
\pgfusepath{clip}%
\pgfsetbuttcap%
\pgfsetroundjoin%
\definecolor{currentfill}{rgb}{0.556863,0.654902,0.847059}%
\pgfsetfillcolor{currentfill}%
\pgfsetfillopacity{0.130000}%
\pgfsetlinewidth{0.000000pt}%
\definecolor{currentstroke}{rgb}{0.000000,0.000000,0.000000}%
\pgfsetstrokecolor{currentstroke}%
\pgfsetdash{}{0pt}%
\pgfpathmoveto{\pgfqpoint{0.939697in}{0.773708in}}%
\pgfpathlineto{\pgfqpoint{0.956215in}{0.769072in}}%
\pgfpathlineto{\pgfqpoint{0.972764in}{0.765598in}}%
\pgfpathlineto{\pgfqpoint{0.981353in}{0.760183in}}%
\pgfpathlineto{\pgfqpoint{0.989901in}{0.755014in}}%
\pgfpathlineto{\pgfqpoint{0.973403in}{0.758409in}}%
\pgfpathlineto{\pgfqpoint{0.956935in}{0.763069in}}%
\pgfpathlineto{\pgfqpoint{0.948336in}{0.768266in}}%
\pgfpathlineto{\pgfqpoint{0.939697in}{0.773708in}}%
\pgfpathclose%
\pgfusepath{fill}%
\end{pgfscope}%
\begin{pgfscope}%
\pgfpathrectangle{\pgfqpoint{0.407592in}{0.222904in}}{\pgfqpoint{1.411163in}{1.411163in}}%
\pgfusepath{clip}%
\pgfsetbuttcap%
\pgfsetroundjoin%
\definecolor{currentfill}{rgb}{0.556863,0.654902,0.847059}%
\pgfsetfillcolor{currentfill}%
\pgfsetfillopacity{0.130000}%
\pgfsetlinewidth{0.000000pt}%
\definecolor{currentstroke}{rgb}{0.000000,0.000000,0.000000}%
\pgfsetstrokecolor{currentstroke}%
\pgfsetdash{}{0pt}%
\pgfpathmoveto{\pgfqpoint{1.621066in}{1.058798in}}%
\pgfpathlineto{\pgfqpoint{1.638853in}{1.070580in}}%
\pgfpathlineto{\pgfqpoint{1.656720in}{1.082531in}}%
\pgfpathlineto{\pgfqpoint{1.664264in}{1.092105in}}%
\pgfpathlineto{\pgfqpoint{1.671795in}{1.101941in}}%
\pgfpathlineto{\pgfqpoint{1.653952in}{1.090268in}}%
\pgfpathlineto{\pgfqpoint{1.636188in}{1.078777in}}%
\pgfpathlineto{\pgfqpoint{1.628633in}{1.068652in}}%
\pgfpathlineto{\pgfqpoint{1.621066in}{1.058798in}}%
\pgfpathclose%
\pgfusepath{fill}%
\end{pgfscope}%
\begin{pgfscope}%
\pgfpathrectangle{\pgfqpoint{0.407592in}{0.222904in}}{\pgfqpoint{1.411163in}{1.411163in}}%
\pgfusepath{clip}%
\pgfsetbuttcap%
\pgfsetroundjoin%
\definecolor{currentfill}{rgb}{0.556863,0.654902,0.847059}%
\pgfsetfillcolor{currentfill}%
\pgfsetfillopacity{0.130000}%
\pgfsetlinewidth{0.000000pt}%
\definecolor{currentstroke}{rgb}{0.000000,0.000000,0.000000}%
\pgfsetstrokecolor{currentstroke}%
\pgfsetdash{}{0pt}%
\pgfpathmoveto{\pgfqpoint{0.655632in}{0.958746in}}%
\pgfpathlineto{\pgfqpoint{0.672071in}{0.944046in}}%
\pgfpathlineto{\pgfqpoint{0.688521in}{0.929816in}}%
\pgfpathlineto{\pgfqpoint{0.697701in}{0.925247in}}%
\pgfpathlineto{\pgfqpoint{0.706837in}{0.920904in}}%
\pgfpathlineto{\pgfqpoint{0.690432in}{0.935583in}}%
\pgfpathlineto{\pgfqpoint{0.674037in}{0.950724in}}%
\pgfpathlineto{\pgfqpoint{0.664856in}{0.954617in}}%
\pgfpathlineto{\pgfqpoint{0.655632in}{0.958746in}}%
\pgfpathclose%
\pgfusepath{fill}%
\end{pgfscope}%
\begin{pgfscope}%
\pgfpathrectangle{\pgfqpoint{0.407592in}{0.222904in}}{\pgfqpoint{1.411163in}{1.411163in}}%
\pgfusepath{clip}%
\pgfsetbuttcap%
\pgfsetroundjoin%
\definecolor{currentfill}{rgb}{0.556863,0.654902,0.847059}%
\pgfsetfillcolor{currentfill}%
\pgfsetfillopacity{0.130000}%
\pgfsetlinewidth{0.000000pt}%
\definecolor{currentstroke}{rgb}{0.000000,0.000000,0.000000}%
\pgfsetstrokecolor{currentstroke}%
\pgfsetdash{}{0pt}%
\pgfpathmoveto{\pgfqpoint{1.222804in}{0.773194in}}%
\pgfpathlineto{\pgfqpoint{1.239743in}{0.782442in}}%
\pgfpathlineto{\pgfqpoint{1.256754in}{0.792144in}}%
\pgfpathlineto{\pgfqpoint{1.264775in}{0.793869in}}%
\pgfpathlineto{\pgfqpoint{1.272772in}{0.796163in}}%
\pgfpathlineto{\pgfqpoint{1.255787in}{0.786123in}}%
\pgfpathlineto{\pgfqpoint{1.238873in}{0.776510in}}%
\pgfpathlineto{\pgfqpoint{1.230851in}{0.774542in}}%
\pgfpathlineto{\pgfqpoint{1.222804in}{0.773194in}}%
\pgfpathclose%
\pgfusepath{fill}%
\end{pgfscope}%
\begin{pgfscope}%
\pgfpathrectangle{\pgfqpoint{0.407592in}{0.222904in}}{\pgfqpoint{1.411163in}{1.411163in}}%
\pgfusepath{clip}%
\pgfsetbuttcap%
\pgfsetroundjoin%
\definecolor{currentfill}{rgb}{0.556863,0.654902,0.847059}%
\pgfsetfillcolor{currentfill}%
\pgfsetfillopacity{0.130000}%
\pgfsetlinewidth{0.000000pt}%
\definecolor{currentstroke}{rgb}{0.000000,0.000000,0.000000}%
\pgfsetstrokecolor{currentstroke}%
\pgfsetdash{}{0pt}%
\pgfpathmoveto{\pgfqpoint{1.073067in}{0.745476in}}%
\pgfpathlineto{\pgfqpoint{1.089731in}{0.748417in}}%
\pgfpathlineto{\pgfqpoint{1.106449in}{0.752370in}}%
\pgfpathlineto{\pgfqpoint{1.114747in}{0.749432in}}%
\pgfpathlineto{\pgfqpoint{1.123009in}{0.746937in}}%
\pgfpathlineto{\pgfqpoint{1.106332in}{0.742401in}}%
\pgfpathlineto{\pgfqpoint{1.089710in}{0.738908in}}%
\pgfpathlineto{\pgfqpoint{1.081406in}{0.741977in}}%
\pgfpathlineto{\pgfqpoint{1.073067in}{0.745476in}}%
\pgfpathclose%
\pgfusepath{fill}%
\end{pgfscope}%
\begin{pgfscope}%
\pgfpathrectangle{\pgfqpoint{0.407592in}{0.222904in}}{\pgfqpoint{1.411163in}{1.411163in}}%
\pgfusepath{clip}%
\pgfsetbuttcap%
\pgfsetroundjoin%
\definecolor{currentfill}{rgb}{0.556863,0.654902,0.847059}%
\pgfsetfillcolor{currentfill}%
\pgfsetfillopacity{0.130000}%
\pgfsetlinewidth{0.000000pt}%
\definecolor{currentstroke}{rgb}{0.000000,0.000000,0.000000}%
\pgfsetstrokecolor{currentstroke}%
\pgfsetdash{}{0pt}%
\pgfpathmoveto{\pgfqpoint{0.772564in}{0.867691in}}%
\pgfpathlineto{\pgfqpoint{0.789025in}{0.855999in}}%
\pgfpathlineto{\pgfqpoint{0.805501in}{0.845049in}}%
\pgfpathlineto{\pgfqpoint{0.814454in}{0.839591in}}%
\pgfpathlineto{\pgfqpoint{0.823366in}{0.834341in}}%
\pgfpathlineto{\pgfqpoint{0.806938in}{0.845714in}}%
\pgfpathlineto{\pgfqpoint{0.790524in}{0.857854in}}%
\pgfpathlineto{\pgfqpoint{0.781565in}{0.862659in}}%
\pgfpathlineto{\pgfqpoint{0.772564in}{0.867691in}}%
\pgfpathclose%
\pgfusepath{fill}%
\end{pgfscope}%
\begin{pgfscope}%
\pgfpathrectangle{\pgfqpoint{0.407592in}{0.222904in}}{\pgfqpoint{1.411163in}{1.411163in}}%
\pgfusepath{clip}%
\pgfsetbuttcap%
\pgfsetroundjoin%
\definecolor{currentfill}{rgb}{0.556863,0.654902,0.847059}%
\pgfsetfillcolor{currentfill}%
\pgfsetfillopacity{0.130000}%
\pgfsetlinewidth{0.000000pt}%
\definecolor{currentstroke}{rgb}{0.000000,0.000000,0.000000}%
\pgfsetstrokecolor{currentstroke}%
\pgfsetdash{}{0pt}%
\pgfpathmoveto{\pgfqpoint{0.538344in}{1.061482in}}%
\pgfpathlineto{\pgfqpoint{0.554771in}{1.045015in}}%
\pgfpathlineto{\pgfqpoint{0.571208in}{1.028845in}}%
\pgfpathlineto{\pgfqpoint{0.580610in}{1.025081in}}%
\pgfpathlineto{\pgfqpoint{0.589967in}{1.021543in}}%
\pgfpathlineto{\pgfqpoint{0.573574in}{1.038091in}}%
\pgfpathlineto{\pgfqpoint{0.557189in}{1.054924in}}%
\pgfpathlineto{\pgfqpoint{0.547789in}{1.058089in}}%
\pgfpathlineto{\pgfqpoint{0.538344in}{1.061482in}}%
\pgfpathclose%
\pgfusepath{fill}%
\end{pgfscope}%
\begin{pgfscope}%
\pgfpathrectangle{\pgfqpoint{0.407592in}{0.222904in}}{\pgfqpoint{1.411163in}{1.411163in}}%
\pgfusepath{clip}%
\pgfsetbuttcap%
\pgfsetroundjoin%
\definecolor{currentfill}{rgb}{0.556863,0.654902,0.847059}%
\pgfsetfillcolor{currentfill}%
\pgfsetfillopacity{0.130000}%
\pgfsetlinewidth{0.000000pt}%
\definecolor{currentstroke}{rgb}{0.000000,0.000000,0.000000}%
\pgfsetstrokecolor{currentstroke}%
\pgfsetdash{}{0pt}%
\pgfpathmoveto{\pgfqpoint{1.737641in}{1.166444in}}%
\pgfpathlineto{\pgfqpoint{1.755676in}{1.178074in}}%
\pgfpathlineto{\pgfqpoint{1.773791in}{1.189874in}}%
\pgfpathlineto{\pgfqpoint{1.781219in}{1.200315in}}%
\pgfpathlineto{\pgfqpoint{1.788636in}{1.210959in}}%
\pgfpathlineto{\pgfqpoint{1.770545in}{1.199433in}}%
\pgfpathlineto{\pgfqpoint{1.752534in}{1.188087in}}%
\pgfpathlineto{\pgfqpoint{1.745093in}{1.177162in}}%
\pgfpathlineto{\pgfqpoint{1.737641in}{1.166444in}}%
\pgfpathclose%
\pgfusepath{fill}%
\end{pgfscope}%
\begin{pgfscope}%
\pgfpathrectangle{\pgfqpoint{0.407592in}{0.222904in}}{\pgfqpoint{1.411163in}{1.411163in}}%
\pgfusepath{clip}%
\pgfsetbuttcap%
\pgfsetroundjoin%
\definecolor{currentfill}{rgb}{0.556863,0.654902,0.847059}%
\pgfsetfillcolor{currentfill}%
\pgfsetfillopacity{0.130000}%
\pgfsetlinewidth{0.000000pt}%
\definecolor{currentstroke}{rgb}{0.000000,0.000000,0.000000}%
\pgfsetstrokecolor{currentstroke}%
\pgfsetdash{}{0pt}%
\pgfpathmoveto{\pgfqpoint{1.272772in}{0.796163in}}%
\pgfpathlineto{\pgfqpoint{1.289831in}{0.806573in}}%
\pgfpathlineto{\pgfqpoint{1.306964in}{0.817310in}}%
\pgfpathlineto{\pgfqpoint{1.314915in}{0.820412in}}%
\pgfpathlineto{\pgfqpoint{1.322844in}{0.824057in}}%
\pgfpathlineto{\pgfqpoint{1.305734in}{0.813134in}}%
\pgfpathlineto{\pgfqpoint{1.288698in}{0.802519in}}%
\pgfpathlineto{\pgfqpoint{1.280746in}{0.799041in}}%
\pgfpathlineto{\pgfqpoint{1.272772in}{0.796163in}}%
\pgfpathclose%
\pgfusepath{fill}%
\end{pgfscope}%
\begin{pgfscope}%
\pgfpathrectangle{\pgfqpoint{0.407592in}{0.222904in}}{\pgfqpoint{1.411163in}{1.411163in}}%
\pgfusepath{clip}%
\pgfsetbuttcap%
\pgfsetroundjoin%
\definecolor{currentfill}{rgb}{0.556863,0.654902,0.847059}%
\pgfsetfillcolor{currentfill}%
\pgfsetfillopacity{0.130000}%
\pgfsetlinewidth{0.000000pt}%
\definecolor{currentstroke}{rgb}{0.000000,0.000000,0.000000}%
\pgfsetstrokecolor{currentstroke}%
\pgfsetdash{}{0pt}%
\pgfpathmoveto{\pgfqpoint{1.438965in}{0.903229in}}%
\pgfpathlineto{\pgfqpoint{1.456392in}{0.915160in}}%
\pgfpathlineto{\pgfqpoint{1.473898in}{0.927261in}}%
\pgfpathlineto{\pgfqpoint{1.481639in}{0.934193in}}%
\pgfpathlineto{\pgfqpoint{1.489366in}{0.941528in}}%
\pgfpathlineto{\pgfqpoint{1.471883in}{0.929595in}}%
\pgfpathlineto{\pgfqpoint{1.454476in}{0.917843in}}%
\pgfpathlineto{\pgfqpoint{1.446728in}{0.910318in}}%
\pgfpathlineto{\pgfqpoint{1.438965in}{0.903229in}}%
\pgfpathclose%
\pgfusepath{fill}%
\end{pgfscope}%
\begin{pgfscope}%
\pgfpathrectangle{\pgfqpoint{0.407592in}{0.222904in}}{\pgfqpoint{1.411163in}{1.411163in}}%
\pgfusepath{clip}%
\pgfsetbuttcap%
\pgfsetroundjoin%
\definecolor{currentfill}{rgb}{0.556863,0.654902,0.847059}%
\pgfsetfillcolor{currentfill}%
\pgfsetfillopacity{0.130000}%
\pgfsetlinewidth{0.000000pt}%
\definecolor{currentstroke}{rgb}{0.000000,0.000000,0.000000}%
\pgfsetstrokecolor{currentstroke}%
\pgfsetdash{}{0pt}%
\pgfpathmoveto{\pgfqpoint{1.555250in}{0.998161in}}%
\pgfpathlineto{\pgfqpoint{1.572925in}{1.010141in}}%
\pgfpathlineto{\pgfqpoint{1.590679in}{1.022282in}}%
\pgfpathlineto{\pgfqpoint{1.598294in}{1.030958in}}%
\pgfpathlineto{\pgfqpoint{1.605897in}{1.039942in}}%
\pgfpathlineto{\pgfqpoint{1.588166in}{1.028055in}}%
\pgfpathlineto{\pgfqpoint{1.570514in}{1.016342in}}%
\pgfpathlineto{\pgfqpoint{1.562888in}{1.007089in}}%
\pgfpathlineto{\pgfqpoint{1.555250in}{0.998161in}}%
\pgfpathclose%
\pgfusepath{fill}%
\end{pgfscope}%
\begin{pgfscope}%
\pgfpathrectangle{\pgfqpoint{0.407592in}{0.222904in}}{\pgfqpoint{1.411163in}{1.411163in}}%
\pgfusepath{clip}%
\pgfsetbuttcap%
\pgfsetroundjoin%
\definecolor{currentfill}{rgb}{0.556863,0.654902,0.847059}%
\pgfsetfillcolor{currentfill}%
\pgfsetfillopacity{0.130000}%
\pgfsetlinewidth{0.000000pt}%
\definecolor{currentstroke}{rgb}{0.000000,0.000000,0.000000}%
\pgfsetstrokecolor{currentstroke}%
\pgfsetdash{}{0pt}%
\pgfpathmoveto{\pgfqpoint{1.023005in}{0.751955in}}%
\pgfpathlineto{\pgfqpoint{1.039619in}{0.752215in}}%
\pgfpathlineto{\pgfqpoint{1.056278in}{0.753592in}}%
\pgfpathlineto{\pgfqpoint{1.064691in}{0.749360in}}%
\pgfpathlineto{\pgfqpoint{1.073067in}{0.745476in}}%
\pgfpathlineto{\pgfqpoint{1.056454in}{0.743637in}}%
\pgfpathlineto{\pgfqpoint{1.039888in}{0.742986in}}%
\pgfpathlineto{\pgfqpoint{1.031465in}{0.747308in}}%
\pgfpathlineto{\pgfqpoint{1.023005in}{0.751955in}}%
\pgfpathclose%
\pgfusepath{fill}%
\end{pgfscope}%
\begin{pgfscope}%
\pgfpathrectangle{\pgfqpoint{0.407592in}{0.222904in}}{\pgfqpoint{1.411163in}{1.411163in}}%
\pgfusepath{clip}%
\pgfsetbuttcap%
\pgfsetroundjoin%
\definecolor{currentfill}{rgb}{0.556863,0.654902,0.847059}%
\pgfsetfillcolor{currentfill}%
\pgfsetfillopacity{0.130000}%
\pgfsetlinewidth{0.000000pt}%
\definecolor{currentstroke}{rgb}{0.000000,0.000000,0.000000}%
\pgfsetstrokecolor{currentstroke}%
\pgfsetdash{}{0pt}%
\pgfpathmoveto{\pgfqpoint{1.788636in}{1.210959in}}%
\pgfpathlineto{\pgfqpoint{1.806807in}{1.222657in}}%
\pgfpathlineto{\pgfqpoint{1.814201in}{1.233363in}}%
\pgfpathlineto{\pgfqpoint{1.821584in}{1.244256in}}%
\pgfpathlineto{\pgfqpoint{1.803438in}{1.232829in}}%
\pgfpathlineto{\pgfqpoint{1.796042in}{1.221799in}}%
\pgfusepath{fill}%
\end{pgfscope}%
\begin{pgfscope}%
\pgfpathrectangle{\pgfqpoint{0.407592in}{0.222904in}}{\pgfqpoint{1.411163in}{1.411163in}}%
\pgfusepath{clip}%
\pgfsetbuttcap%
\pgfsetroundjoin%
\definecolor{currentfill}{rgb}{0.556863,0.654902,0.847059}%
\pgfsetfillcolor{currentfill}%
\pgfsetfillopacity{0.130000}%
\pgfsetlinewidth{0.000000pt}%
\definecolor{currentstroke}{rgb}{0.000000,0.000000,0.000000}%
\pgfsetstrokecolor{currentstroke}%
\pgfsetdash{}{0pt}%
\pgfpathmoveto{\pgfqpoint{0.420674in}{1.170931in}}%
\pgfpathlineto{\pgfqpoint{0.437092in}{1.153379in}}%
\pgfpathlineto{\pgfqpoint{0.453519in}{1.136025in}}%
\pgfpathlineto{\pgfqpoint{0.463143in}{1.132863in}}%
\pgfpathlineto{\pgfqpoint{0.472720in}{1.129914in}}%
\pgfpathlineto{\pgfqpoint{0.456335in}{1.147569in}}%
\pgfpathlineto{\pgfqpoint{0.439960in}{1.165411in}}%
\pgfpathlineto{\pgfqpoint{0.430340in}{1.168065in}}%
\pgfpathlineto{\pgfqpoint{0.420674in}{1.170931in}}%
\pgfpathclose%
\pgfusepath{fill}%
\end{pgfscope}%
\begin{pgfscope}%
\pgfpathrectangle{\pgfqpoint{0.407592in}{0.222904in}}{\pgfqpoint{1.411163in}{1.411163in}}%
\pgfusepath{clip}%
\pgfsetbuttcap%
\pgfsetroundjoin%
\definecolor{currentfill}{rgb}{0.556863,0.654902,0.847059}%
\pgfsetfillcolor{currentfill}%
\pgfsetfillopacity{0.130000}%
\pgfsetlinewidth{0.000000pt}%
\definecolor{currentstroke}{rgb}{0.000000,0.000000,0.000000}%
\pgfsetstrokecolor{currentstroke}%
\pgfsetdash{}{0pt}%
\pgfpathmoveto{\pgfqpoint{0.688521in}{0.929816in}}%
\pgfpathlineto{\pgfqpoint{0.704981in}{0.916092in}}%
\pgfpathlineto{\pgfqpoint{0.721453in}{0.902916in}}%
\pgfpathlineto{\pgfqpoint{0.730587in}{0.897905in}}%
\pgfpathlineto{\pgfqpoint{0.739678in}{0.893106in}}%
\pgfpathlineto{\pgfqpoint{0.723253in}{0.906731in}}%
\pgfpathlineto{\pgfqpoint{0.706837in}{0.920904in}}%
\pgfpathlineto{\pgfqpoint{0.697701in}{0.925247in}}%
\pgfpathlineto{\pgfqpoint{0.688521in}{0.929816in}}%
\pgfpathclose%
\pgfusepath{fill}%
\end{pgfscope}%
\begin{pgfscope}%
\pgfpathrectangle{\pgfqpoint{0.407592in}{0.222904in}}{\pgfqpoint{1.411163in}{1.411163in}}%
\pgfusepath{clip}%
\pgfsetbuttcap%
\pgfsetroundjoin%
\definecolor{currentfill}{rgb}{0.556863,0.654902,0.847059}%
\pgfsetfillcolor{currentfill}%
\pgfsetfillopacity{0.130000}%
\pgfsetlinewidth{0.000000pt}%
\definecolor{currentstroke}{rgb}{0.000000,0.000000,0.000000}%
\pgfsetstrokecolor{currentstroke}%
\pgfsetdash{}{0pt}%
\pgfpathmoveto{\pgfqpoint{1.671795in}{1.101941in}}%
\pgfpathlineto{\pgfqpoint{1.689718in}{1.113785in}}%
\pgfpathlineto{\pgfqpoint{1.707721in}{1.125793in}}%
\pgfpathlineto{\pgfqpoint{1.715218in}{1.135608in}}%
\pgfpathlineto{\pgfqpoint{1.722703in}{1.145660in}}%
\pgfpathlineto{\pgfqpoint{1.704724in}{1.133923in}}%
\pgfpathlineto{\pgfqpoint{1.686825in}{1.122362in}}%
\pgfpathlineto{\pgfqpoint{1.679316in}{1.112029in}}%
\pgfpathlineto{\pgfqpoint{1.671795in}{1.101941in}}%
\pgfpathclose%
\pgfusepath{fill}%
\end{pgfscope}%
\begin{pgfscope}%
\pgfpathrectangle{\pgfqpoint{0.407592in}{0.222904in}}{\pgfqpoint{1.411163in}{1.411163in}}%
\pgfusepath{clip}%
\pgfsetbuttcap%
\pgfsetroundjoin%
\definecolor{currentfill}{rgb}{0.556863,0.654902,0.847059}%
\pgfsetfillcolor{currentfill}%
\pgfsetfillopacity{0.130000}%
\pgfsetlinewidth{0.000000pt}%
\definecolor{currentstroke}{rgb}{0.000000,0.000000,0.000000}%
\pgfsetstrokecolor{currentstroke}%
\pgfsetdash{}{0pt}%
\pgfpathmoveto{\pgfqpoint{1.322844in}{0.824057in}}%
\pgfpathlineto{\pgfqpoint{1.340029in}{0.835252in}}%
\pgfpathlineto{\pgfqpoint{1.357291in}{0.846689in}}%
\pgfpathlineto{\pgfqpoint{1.365177in}{0.850974in}}%
\pgfpathlineto{\pgfqpoint{1.373045in}{0.855762in}}%
\pgfpathlineto{\pgfqpoint{1.355806in}{0.844265in}}%
\pgfpathlineto{\pgfqpoint{1.338642in}{0.833004in}}%
\pgfpathlineto{\pgfqpoint{1.330753in}{0.828253in}}%
\pgfpathlineto{\pgfqpoint{1.322844in}{0.824057in}}%
\pgfpathclose%
\pgfusepath{fill}%
\end{pgfscope}%
\begin{pgfscope}%
\pgfpathrectangle{\pgfqpoint{0.407592in}{0.222904in}}{\pgfqpoint{1.411163in}{1.411163in}}%
\pgfusepath{clip}%
\pgfsetbuttcap%
\pgfsetroundjoin%
\definecolor{currentfill}{rgb}{0.556863,0.654902,0.847059}%
\pgfsetfillcolor{currentfill}%
\pgfsetfillopacity{0.130000}%
\pgfsetlinewidth{0.000000pt}%
\definecolor{currentstroke}{rgb}{0.000000,0.000000,0.000000}%
\pgfsetstrokecolor{currentstroke}%
\pgfsetdash{}{0pt}%
\pgfpathmoveto{\pgfqpoint{0.571208in}{1.028845in}}%
\pgfpathlineto{\pgfqpoint{0.587654in}{1.012999in}}%
\pgfpathlineto{\pgfqpoint{0.604109in}{0.997501in}}%
\pgfpathlineto{\pgfqpoint{0.613468in}{0.993344in}}%
\pgfpathlineto{\pgfqpoint{0.622781in}{0.989409in}}%
\pgfpathlineto{\pgfqpoint{0.606370in}{1.005307in}}%
\pgfpathlineto{\pgfqpoint{0.589967in}{1.021543in}}%
\pgfpathlineto{\pgfqpoint{0.580610in}{1.025081in}}%
\pgfpathlineto{\pgfqpoint{0.571208in}{1.028845in}}%
\pgfpathclose%
\pgfusepath{fill}%
\end{pgfscope}%
\begin{pgfscope}%
\pgfpathrectangle{\pgfqpoint{0.407592in}{0.222904in}}{\pgfqpoint{1.411163in}{1.411163in}}%
\pgfusepath{clip}%
\pgfsetbuttcap%
\pgfsetroundjoin%
\definecolor{currentfill}{rgb}{0.556863,0.654902,0.847059}%
\pgfsetfillcolor{currentfill}%
\pgfsetfillopacity{0.130000}%
\pgfsetlinewidth{0.000000pt}%
\definecolor{currentstroke}{rgb}{0.000000,0.000000,0.000000}%
\pgfsetstrokecolor{currentstroke}%
\pgfsetdash{}{0pt}%
\pgfpathmoveto{\pgfqpoint{0.889237in}{0.797644in}}%
\pgfpathlineto{\pgfqpoint{0.905755in}{0.790911in}}%
\pgfpathlineto{\pgfqpoint{0.922297in}{0.785227in}}%
\pgfpathlineto{\pgfqpoint{0.931017in}{0.779369in}}%
\pgfpathlineto{\pgfqpoint{0.939697in}{0.773708in}}%
\pgfpathlineto{\pgfqpoint{0.923206in}{0.779499in}}%
\pgfpathlineto{\pgfqpoint{0.906739in}{0.786422in}}%
\pgfpathlineto{\pgfqpoint{0.898008in}{0.791929in}}%
\pgfpathlineto{\pgfqpoint{0.889237in}{0.797644in}}%
\pgfpathclose%
\pgfusepath{fill}%
\end{pgfscope}%
\begin{pgfscope}%
\pgfpathrectangle{\pgfqpoint{0.407592in}{0.222904in}}{\pgfqpoint{1.411163in}{1.411163in}}%
\pgfusepath{clip}%
\pgfsetbuttcap%
\pgfsetroundjoin%
\definecolor{currentfill}{rgb}{0.556863,0.654902,0.847059}%
\pgfsetfillcolor{currentfill}%
\pgfsetfillopacity{0.130000}%
\pgfsetlinewidth{0.000000pt}%
\definecolor{currentstroke}{rgb}{0.000000,0.000000,0.000000}%
\pgfsetstrokecolor{currentstroke}%
\pgfsetdash{}{0pt}%
\pgfpathmoveto{\pgfqpoint{0.805501in}{0.845049in}}%
\pgfpathlineto{\pgfqpoint{0.821991in}{0.834892in}}%
\pgfpathlineto{\pgfqpoint{0.838498in}{0.825573in}}%
\pgfpathlineto{\pgfqpoint{0.847403in}{0.819756in}}%
\pgfpathlineto{\pgfqpoint{0.856266in}{0.814128in}}%
\pgfpathlineto{\pgfqpoint{0.839808in}{0.823794in}}%
\pgfpathlineto{\pgfqpoint{0.823366in}{0.834341in}}%
\pgfpathlineto{\pgfqpoint{0.814454in}{0.839591in}}%
\pgfpathlineto{\pgfqpoint{0.805501in}{0.845049in}}%
\pgfpathclose%
\pgfusepath{fill}%
\end{pgfscope}%
\begin{pgfscope}%
\pgfpathrectangle{\pgfqpoint{0.407592in}{0.222904in}}{\pgfqpoint{1.411163in}{1.411163in}}%
\pgfusepath{clip}%
\pgfsetbuttcap%
\pgfsetroundjoin%
\definecolor{currentfill}{rgb}{0.556863,0.654902,0.847059}%
\pgfsetfillcolor{currentfill}%
\pgfsetfillopacity{0.130000}%
\pgfsetlinewidth{0.000000pt}%
\definecolor{currentstroke}{rgb}{0.000000,0.000000,0.000000}%
\pgfsetstrokecolor{currentstroke}%
\pgfsetdash{}{0pt}%
\pgfpathmoveto{\pgfqpoint{1.489366in}{0.941528in}}%
\pgfpathlineto{\pgfqpoint{1.506928in}{0.953627in}}%
\pgfpathlineto{\pgfqpoint{1.524568in}{0.965882in}}%
\pgfpathlineto{\pgfqpoint{1.532259in}{0.973418in}}%
\pgfpathlineto{\pgfqpoint{1.539936in}{0.981316in}}%
\pgfpathlineto{\pgfqpoint{1.522318in}{0.969261in}}%
\pgfpathlineto{\pgfqpoint{1.504778in}{0.957375in}}%
\pgfpathlineto{\pgfqpoint{1.497079in}{0.949258in}}%
\pgfpathlineto{\pgfqpoint{1.489366in}{0.941528in}}%
\pgfpathclose%
\pgfusepath{fill}%
\end{pgfscope}%
\begin{pgfscope}%
\pgfpathrectangle{\pgfqpoint{0.407592in}{0.222904in}}{\pgfqpoint{1.411163in}{1.411163in}}%
\pgfusepath{clip}%
\pgfsetbuttcap%
\pgfsetroundjoin%
\definecolor{currentfill}{rgb}{0.556863,0.654902,0.847059}%
\pgfsetfillcolor{currentfill}%
\pgfsetfillopacity{0.130000}%
\pgfsetlinewidth{0.000000pt}%
\definecolor{currentstroke}{rgb}{0.000000,0.000000,0.000000}%
\pgfsetstrokecolor{currentstroke}%
\pgfsetdash{}{0pt}%
\pgfpathmoveto{\pgfqpoint{0.453519in}{1.136025in}}%
\pgfpathlineto{\pgfqpoint{0.469955in}{1.118886in}}%
\pgfpathlineto{\pgfqpoint{0.486399in}{1.101976in}}%
\pgfpathlineto{\pgfqpoint{0.495980in}{1.098494in}}%
\pgfpathlineto{\pgfqpoint{0.505514in}{1.095225in}}%
\pgfpathlineto{\pgfqpoint{0.489113in}{1.112460in}}%
\pgfpathlineto{\pgfqpoint{0.472720in}{1.129914in}}%
\pgfpathlineto{\pgfqpoint{0.463143in}{1.132863in}}%
\pgfpathlineto{\pgfqpoint{0.453519in}{1.136025in}}%
\pgfpathclose%
\pgfusepath{fill}%
\end{pgfscope}%
\begin{pgfscope}%
\pgfpathrectangle{\pgfqpoint{0.407592in}{0.222904in}}{\pgfqpoint{1.411163in}{1.411163in}}%
\pgfusepath{clip}%
\pgfsetbuttcap%
\pgfsetroundjoin%
\definecolor{currentfill}{rgb}{0.556863,0.654902,0.847059}%
\pgfsetfillcolor{currentfill}%
\pgfsetfillopacity{0.130000}%
\pgfsetlinewidth{0.000000pt}%
\definecolor{currentstroke}{rgb}{0.000000,0.000000,0.000000}%
\pgfsetstrokecolor{currentstroke}%
\pgfsetdash{}{0pt}%
\pgfpathmoveto{\pgfqpoint{0.972764in}{0.765598in}}%
\pgfpathlineto{\pgfqpoint{0.989348in}{0.763279in}}%
\pgfpathlineto{\pgfqpoint{1.005969in}{0.762093in}}%
\pgfpathlineto{\pgfqpoint{1.014506in}{0.756894in}}%
\pgfpathlineto{\pgfqpoint{1.023005in}{0.751955in}}%
\pgfpathlineto{\pgfqpoint{1.006434in}{0.752873in}}%
\pgfpathlineto{\pgfqpoint{0.989901in}{0.755014in}}%
\pgfpathlineto{\pgfqpoint{0.981353in}{0.760183in}}%
\pgfpathlineto{\pgfqpoint{0.972764in}{0.765598in}}%
\pgfpathclose%
\pgfusepath{fill}%
\end{pgfscope}%
\begin{pgfscope}%
\pgfpathrectangle{\pgfqpoint{0.407592in}{0.222904in}}{\pgfqpoint{1.411163in}{1.411163in}}%
\pgfusepath{clip}%
\pgfsetbuttcap%
\pgfsetroundjoin%
\definecolor{currentfill}{rgb}{0.556863,0.654902,0.847059}%
\pgfsetfillcolor{currentfill}%
\pgfsetfillopacity{0.130000}%
\pgfsetlinewidth{0.000000pt}%
\definecolor{currentstroke}{rgb}{0.000000,0.000000,0.000000}%
\pgfsetstrokecolor{currentstroke}%
\pgfsetdash{}{0pt}%
\pgfpathmoveto{\pgfqpoint{1.605897in}{1.039942in}}%
\pgfpathlineto{\pgfqpoint{1.623707in}{1.051993in}}%
\pgfpathlineto{\pgfqpoint{1.641597in}{1.064199in}}%
\pgfpathlineto{\pgfqpoint{1.649165in}{1.073227in}}%
\pgfpathlineto{\pgfqpoint{1.656720in}{1.082531in}}%
\pgfpathlineto{\pgfqpoint{1.638853in}{1.070580in}}%
\pgfpathlineto{\pgfqpoint{1.621066in}{1.058798in}}%
\pgfpathlineto{\pgfqpoint{1.613488in}{1.049225in}}%
\pgfpathlineto{\pgfqpoint{1.605897in}{1.039942in}}%
\pgfpathclose%
\pgfusepath{fill}%
\end{pgfscope}%
\begin{pgfscope}%
\pgfpathrectangle{\pgfqpoint{0.407592in}{0.222904in}}{\pgfqpoint{1.411163in}{1.411163in}}%
\pgfusepath{clip}%
\pgfsetbuttcap%
\pgfsetroundjoin%
\definecolor{currentfill}{rgb}{0.556863,0.654902,0.847059}%
\pgfsetfillcolor{currentfill}%
\pgfsetfillopacity{0.130000}%
\pgfsetlinewidth{0.000000pt}%
\definecolor{currentstroke}{rgb}{0.000000,0.000000,0.000000}%
\pgfsetstrokecolor{currentstroke}%
\pgfsetdash{}{0pt}%
\pgfpathmoveto{\pgfqpoint{1.373045in}{0.855762in}}%
\pgfpathlineto{\pgfqpoint{1.390361in}{0.867471in}}%
\pgfpathlineto{\pgfqpoint{1.407754in}{0.879371in}}%
\pgfpathlineto{\pgfqpoint{1.415582in}{0.884650in}}%
\pgfpathlineto{\pgfqpoint{1.423392in}{0.890388in}}%
\pgfpathlineto{\pgfqpoint{1.406021in}{0.878524in}}%
\pgfpathlineto{\pgfqpoint{1.388726in}{0.866852in}}%
\pgfpathlineto{\pgfqpoint{1.380894in}{0.861056in}}%
\pgfpathlineto{\pgfqpoint{1.373045in}{0.855762in}}%
\pgfpathclose%
\pgfusepath{fill}%
\end{pgfscope}%
\begin{pgfscope}%
\pgfpathrectangle{\pgfqpoint{0.407592in}{0.222904in}}{\pgfqpoint{1.411163in}{1.411163in}}%
\pgfusepath{clip}%
\pgfsetbuttcap%
\pgfsetroundjoin%
\definecolor{currentfill}{rgb}{0.556863,0.654902,0.847059}%
\pgfsetfillcolor{currentfill}%
\pgfsetfillopacity{0.130000}%
\pgfsetlinewidth{0.000000pt}%
\definecolor{currentstroke}{rgb}{0.000000,0.000000,0.000000}%
\pgfsetstrokecolor{currentstroke}%
\pgfsetdash{}{0pt}%
\pgfpathmoveto{\pgfqpoint{0.604109in}{0.997501in}}%
\pgfpathlineto{\pgfqpoint{0.620574in}{0.982379in}}%
\pgfpathlineto{\pgfqpoint{0.637048in}{0.967665in}}%
\pgfpathlineto{\pgfqpoint{0.646362in}{0.963099in}}%
\pgfpathlineto{\pgfqpoint{0.655632in}{0.958746in}}%
\pgfpathlineto{\pgfqpoint{0.639202in}{0.973878in}}%
\pgfpathlineto{\pgfqpoint{0.622781in}{0.989409in}}%
\pgfpathlineto{\pgfqpoint{0.613468in}{0.993344in}}%
\pgfpathlineto{\pgfqpoint{0.604109in}{0.997501in}}%
\pgfpathclose%
\pgfusepath{fill}%
\end{pgfscope}%
\begin{pgfscope}%
\pgfpathrectangle{\pgfqpoint{0.407592in}{0.222904in}}{\pgfqpoint{1.411163in}{1.411163in}}%
\pgfusepath{clip}%
\pgfsetbuttcap%
\pgfsetroundjoin%
\definecolor{currentfill}{rgb}{0.556863,0.654902,0.847059}%
\pgfsetfillcolor{currentfill}%
\pgfsetfillopacity{0.130000}%
\pgfsetlinewidth{0.000000pt}%
\definecolor{currentstroke}{rgb}{0.000000,0.000000,0.000000}%
\pgfsetstrokecolor{currentstroke}%
\pgfsetdash{}{0pt}%
\pgfpathmoveto{\pgfqpoint{0.721453in}{0.902916in}}%
\pgfpathlineto{\pgfqpoint{0.737936in}{0.890328in}}%
\pgfpathlineto{\pgfqpoint{0.754432in}{0.878371in}}%
\pgfpathlineto{\pgfqpoint{0.763519in}{0.872933in}}%
\pgfpathlineto{\pgfqpoint{0.772564in}{0.867691in}}%
\pgfpathlineto{\pgfqpoint{0.756115in}{0.880077in}}%
\pgfpathlineto{\pgfqpoint{0.739678in}{0.893106in}}%
\pgfpathlineto{\pgfqpoint{0.730587in}{0.897905in}}%
\pgfpathlineto{\pgfqpoint{0.721453in}{0.902916in}}%
\pgfpathclose%
\pgfusepath{fill}%
\end{pgfscope}%
\begin{pgfscope}%
\pgfpathrectangle{\pgfqpoint{0.407592in}{0.222904in}}{\pgfqpoint{1.411163in}{1.411163in}}%
\pgfusepath{clip}%
\pgfsetbuttcap%
\pgfsetroundjoin%
\definecolor{currentfill}{rgb}{0.556863,0.654902,0.847059}%
\pgfsetfillcolor{currentfill}%
\pgfsetfillopacity{0.130000}%
\pgfsetlinewidth{0.000000pt}%
\definecolor{currentstroke}{rgb}{0.000000,0.000000,0.000000}%
\pgfsetstrokecolor{currentstroke}%
\pgfsetdash{}{0pt}%
\pgfpathmoveto{\pgfqpoint{1.722703in}{1.145660in}}%
\pgfpathlineto{\pgfqpoint{1.740763in}{1.157563in}}%
\pgfpathlineto{\pgfqpoint{1.758903in}{1.169625in}}%
\pgfpathlineto{\pgfqpoint{1.766353in}{1.179642in}}%
\pgfpathlineto{\pgfqpoint{1.773791in}{1.189874in}}%
\pgfpathlineto{\pgfqpoint{1.755676in}{1.178074in}}%
\pgfpathlineto{\pgfqpoint{1.737641in}{1.166444in}}%
\pgfpathlineto{\pgfqpoint{1.730178in}{1.155941in}}%
\pgfpathlineto{\pgfqpoint{1.722703in}{1.145660in}}%
\pgfpathclose%
\pgfusepath{fill}%
\end{pgfscope}%
\begin{pgfscope}%
\pgfpathrectangle{\pgfqpoint{0.407592in}{0.222904in}}{\pgfqpoint{1.411163in}{1.411163in}}%
\pgfusepath{clip}%
\pgfsetbuttcap%
\pgfsetroundjoin%
\definecolor{currentfill}{rgb}{0.556863,0.654902,0.847059}%
\pgfsetfillcolor{currentfill}%
\pgfsetfillopacity{0.130000}%
\pgfsetlinewidth{0.000000pt}%
\definecolor{currentstroke}{rgb}{0.000000,0.000000,0.000000}%
\pgfsetstrokecolor{currentstroke}%
\pgfsetdash{}{0pt}%
\pgfpathmoveto{\pgfqpoint{0.486399in}{1.101976in}}%
\pgfpathlineto{\pgfqpoint{0.502853in}{1.085314in}}%
\pgfpathlineto{\pgfqpoint{0.519315in}{1.068919in}}%
\pgfpathlineto{\pgfqpoint{0.528852in}{1.065095in}}%
\pgfpathlineto{\pgfqpoint{0.538344in}{1.061482in}}%
\pgfpathlineto{\pgfqpoint{0.521925in}{1.078226in}}%
\pgfpathlineto{\pgfqpoint{0.505514in}{1.095225in}}%
\pgfpathlineto{\pgfqpoint{0.495980in}{1.098494in}}%
\pgfpathlineto{\pgfqpoint{0.486399in}{1.101976in}}%
\pgfpathclose%
\pgfusepath{fill}%
\end{pgfscope}%
\begin{pgfscope}%
\pgfpathrectangle{\pgfqpoint{0.407592in}{0.222904in}}{\pgfqpoint{1.411163in}{1.411163in}}%
\pgfusepath{clip}%
\pgfsetbuttcap%
\pgfsetroundjoin%
\definecolor{currentfill}{rgb}{0.556863,0.654902,0.847059}%
\pgfsetfillcolor{currentfill}%
\pgfsetfillopacity{0.130000}%
\pgfsetlinewidth{0.000000pt}%
\definecolor{currentstroke}{rgb}{0.000000,0.000000,0.000000}%
\pgfsetstrokecolor{currentstroke}%
\pgfsetdash{}{0pt}%
\pgfpathmoveto{\pgfqpoint{1.156545in}{0.758707in}}%
\pgfpathlineto{\pgfqpoint{1.173409in}{0.765747in}}%
\pgfpathlineto{\pgfqpoint{1.190339in}{0.773441in}}%
\pgfpathlineto{\pgfqpoint{1.198499in}{0.772594in}}%
\pgfpathlineto{\pgfqpoint{1.206629in}{0.772249in}}%
\pgfpathlineto{\pgfqpoint{1.189731in}{0.764011in}}%
\pgfpathlineto{\pgfqpoint{1.172900in}{0.756403in}}%
\pgfpathlineto{\pgfqpoint{1.164738in}{0.757294in}}%
\pgfpathlineto{\pgfqpoint{1.156545in}{0.758707in}}%
\pgfpathclose%
\pgfusepath{fill}%
\end{pgfscope}%
\begin{pgfscope}%
\pgfpathrectangle{\pgfqpoint{0.407592in}{0.222904in}}{\pgfqpoint{1.411163in}{1.411163in}}%
\pgfusepath{clip}%
\pgfsetbuttcap%
\pgfsetroundjoin%
\definecolor{currentfill}{rgb}{0.556863,0.654902,0.847059}%
\pgfsetfillcolor{currentfill}%
\pgfsetfillopacity{0.130000}%
\pgfsetlinewidth{0.000000pt}%
\definecolor{currentstroke}{rgb}{0.000000,0.000000,0.000000}%
\pgfsetstrokecolor{currentstroke}%
\pgfsetdash{}{0pt}%
\pgfpathmoveto{\pgfqpoint{1.206629in}{0.772249in}}%
\pgfpathlineto{\pgfqpoint{1.223596in}{0.781041in}}%
\pgfpathlineto{\pgfqpoint{1.240634in}{0.790318in}}%
\pgfpathlineto{\pgfqpoint{1.248707in}{0.790967in}}%
\pgfpathlineto{\pgfqpoint{1.256754in}{0.792144in}}%
\pgfpathlineto{\pgfqpoint{1.239743in}{0.782442in}}%
\pgfpathlineto{\pgfqpoint{1.222804in}{0.773194in}}%
\pgfpathlineto{\pgfqpoint{1.214730in}{0.772440in}}%
\pgfpathlineto{\pgfqpoint{1.206629in}{0.772249in}}%
\pgfpathclose%
\pgfusepath{fill}%
\end{pgfscope}%
\begin{pgfscope}%
\pgfpathrectangle{\pgfqpoint{0.407592in}{0.222904in}}{\pgfqpoint{1.411163in}{1.411163in}}%
\pgfusepath{clip}%
\pgfsetbuttcap%
\pgfsetroundjoin%
\definecolor{currentfill}{rgb}{0.556863,0.654902,0.847059}%
\pgfsetfillcolor{currentfill}%
\pgfsetfillopacity{0.130000}%
\pgfsetlinewidth{0.000000pt}%
\definecolor{currentstroke}{rgb}{0.000000,0.000000,0.000000}%
\pgfsetstrokecolor{currentstroke}%
\pgfsetdash{}{0pt}%
\pgfpathmoveto{\pgfqpoint{1.773791in}{1.189874in}}%
\pgfpathlineto{\pgfqpoint{1.791988in}{1.201836in}}%
\pgfpathlineto{\pgfqpoint{1.799403in}{1.212146in}}%
\pgfpathlineto{\pgfqpoint{1.806807in}{1.222657in}}%
\pgfpathlineto{\pgfqpoint{1.788636in}{1.210959in}}%
\pgfpathlineto{\pgfqpoint{1.781219in}{1.200315in}}%
\pgfusepath{fill}%
\end{pgfscope}%
\begin{pgfscope}%
\pgfpathrectangle{\pgfqpoint{0.407592in}{0.222904in}}{\pgfqpoint{1.411163in}{1.411163in}}%
\pgfusepath{clip}%
\pgfsetbuttcap%
\pgfsetroundjoin%
\definecolor{currentfill}{rgb}{0.556863,0.654902,0.847059}%
\pgfsetfillcolor{currentfill}%
\pgfsetfillopacity{0.130000}%
\pgfsetlinewidth{0.000000pt}%
\definecolor{currentstroke}{rgb}{0.000000,0.000000,0.000000}%
\pgfsetstrokecolor{currentstroke}%
\pgfsetdash{}{0pt}%
\pgfpathmoveto{\pgfqpoint{1.539936in}{0.981316in}}%
\pgfpathlineto{\pgfqpoint{1.557633in}{0.993528in}}%
\pgfpathlineto{\pgfqpoint{1.575410in}{1.005887in}}%
\pgfpathlineto{\pgfqpoint{1.583051in}{1.013922in}}%
\pgfpathlineto{\pgfqpoint{1.590679in}{1.022282in}}%
\pgfpathlineto{\pgfqpoint{1.572925in}{1.010141in}}%
\pgfpathlineto{\pgfqpoint{1.555250in}{0.998161in}}%
\pgfpathlineto{\pgfqpoint{1.547600in}{0.989566in}}%
\pgfpathlineto{\pgfqpoint{1.539936in}{0.981316in}}%
\pgfpathclose%
\pgfusepath{fill}%
\end{pgfscope}%
\begin{pgfscope}%
\pgfpathrectangle{\pgfqpoint{0.407592in}{0.222904in}}{\pgfqpoint{1.411163in}{1.411163in}}%
\pgfusepath{clip}%
\pgfsetbuttcap%
\pgfsetroundjoin%
\definecolor{currentfill}{rgb}{0.556863,0.654902,0.847059}%
\pgfsetfillcolor{currentfill}%
\pgfsetfillopacity{0.130000}%
\pgfsetlinewidth{0.000000pt}%
\definecolor{currentstroke}{rgb}{0.000000,0.000000,0.000000}%
\pgfsetstrokecolor{currentstroke}%
\pgfsetdash{}{0pt}%
\pgfpathmoveto{\pgfqpoint{1.106449in}{0.752370in}}%
\pgfpathlineto{\pgfqpoint{1.123225in}{0.757243in}}%
\pgfpathlineto{\pgfqpoint{1.140062in}{0.762947in}}%
\pgfpathlineto{\pgfqpoint{1.148320in}{0.760603in}}%
\pgfpathlineto{\pgfqpoint{1.156545in}{0.758707in}}%
\pgfpathlineto{\pgfqpoint{1.139746in}{0.752406in}}%
\pgfpathlineto{\pgfqpoint{1.123009in}{0.746937in}}%
\pgfpathlineto{\pgfqpoint{1.114747in}{0.749432in}}%
\pgfpathlineto{\pgfqpoint{1.106449in}{0.752370in}}%
\pgfpathclose%
\pgfusepath{fill}%
\end{pgfscope}%
\begin{pgfscope}%
\pgfpathrectangle{\pgfqpoint{0.407592in}{0.222904in}}{\pgfqpoint{1.411163in}{1.411163in}}%
\pgfusepath{clip}%
\pgfsetbuttcap%
\pgfsetroundjoin%
\definecolor{currentfill}{rgb}{0.556863,0.654902,0.847059}%
\pgfsetfillcolor{currentfill}%
\pgfsetfillopacity{0.130000}%
\pgfsetlinewidth{0.000000pt}%
\definecolor{currentstroke}{rgb}{0.000000,0.000000,0.000000}%
\pgfsetstrokecolor{currentstroke}%
\pgfsetdash{}{0pt}%
\pgfpathmoveto{\pgfqpoint{1.423392in}{0.890388in}}%
\pgfpathlineto{\pgfqpoint{1.440841in}{0.902429in}}%
\pgfpathlineto{\pgfqpoint{1.458368in}{0.914632in}}%
\pgfpathlineto{\pgfqpoint{1.466141in}{0.920739in}}%
\pgfpathlineto{\pgfqpoint{1.473898in}{0.927261in}}%
\pgfpathlineto{\pgfqpoint{1.456392in}{0.915160in}}%
\pgfpathlineto{\pgfqpoint{1.438965in}{0.903229in}}%
\pgfpathlineto{\pgfqpoint{1.431187in}{0.896583in}}%
\pgfpathlineto{\pgfqpoint{1.423392in}{0.890388in}}%
\pgfpathclose%
\pgfusepath{fill}%
\end{pgfscope}%
\begin{pgfscope}%
\pgfpathrectangle{\pgfqpoint{0.407592in}{0.222904in}}{\pgfqpoint{1.411163in}{1.411163in}}%
\pgfusepath{clip}%
\pgfsetbuttcap%
\pgfsetroundjoin%
\definecolor{currentfill}{rgb}{0.556863,0.654902,0.847059}%
\pgfsetfillcolor{currentfill}%
\pgfsetfillopacity{0.130000}%
\pgfsetlinewidth{0.000000pt}%
\definecolor{currentstroke}{rgb}{0.000000,0.000000,0.000000}%
\pgfsetstrokecolor{currentstroke}%
\pgfsetdash{}{0pt}%
\pgfpathmoveto{\pgfqpoint{0.838498in}{0.825573in}}%
\pgfpathlineto{\pgfqpoint{0.855024in}{0.817136in}}%
\pgfpathlineto{\pgfqpoint{0.871571in}{0.809619in}}%
\pgfpathlineto{\pgfqpoint{0.880425in}{0.803545in}}%
\pgfpathlineto{\pgfqpoint{0.889237in}{0.797644in}}%
\pgfpathlineto{\pgfqpoint{0.872742in}{0.805395in}}%
\pgfpathlineto{\pgfqpoint{0.856266in}{0.814128in}}%
\pgfpathlineto{\pgfqpoint{0.847403in}{0.819756in}}%
\pgfpathlineto{\pgfqpoint{0.838498in}{0.825573in}}%
\pgfpathclose%
\pgfusepath{fill}%
\end{pgfscope}%
\begin{pgfscope}%
\pgfpathrectangle{\pgfqpoint{0.407592in}{0.222904in}}{\pgfqpoint{1.411163in}{1.411163in}}%
\pgfusepath{clip}%
\pgfsetbuttcap%
\pgfsetroundjoin%
\definecolor{currentfill}{rgb}{0.556863,0.654902,0.847059}%
\pgfsetfillcolor{currentfill}%
\pgfsetfillopacity{0.130000}%
\pgfsetlinewidth{0.000000pt}%
\definecolor{currentstroke}{rgb}{0.000000,0.000000,0.000000}%
\pgfsetstrokecolor{currentstroke}%
\pgfsetdash{}{0pt}%
\pgfpathmoveto{\pgfqpoint{1.656720in}{1.082531in}}%
\pgfpathlineto{\pgfqpoint{1.674667in}{1.094641in}}%
\pgfpathlineto{\pgfqpoint{1.692694in}{1.106903in}}%
\pgfpathlineto{\pgfqpoint{1.700213in}{1.116222in}}%
\pgfpathlineto{\pgfqpoint{1.707721in}{1.125793in}}%
\pgfpathlineto{\pgfqpoint{1.689718in}{1.113785in}}%
\pgfpathlineto{\pgfqpoint{1.671795in}{1.101941in}}%
\pgfpathlineto{\pgfqpoint{1.664264in}{1.092105in}}%
\pgfpathlineto{\pgfqpoint{1.656720in}{1.082531in}}%
\pgfpathclose%
\pgfusepath{fill}%
\end{pgfscope}%
\begin{pgfscope}%
\pgfpathrectangle{\pgfqpoint{0.407592in}{0.222904in}}{\pgfqpoint{1.411163in}{1.411163in}}%
\pgfusepath{clip}%
\pgfsetbuttcap%
\pgfsetroundjoin%
\definecolor{currentfill}{rgb}{0.556863,0.654902,0.847059}%
\pgfsetfillcolor{currentfill}%
\pgfsetfillopacity{0.130000}%
\pgfsetlinewidth{0.000000pt}%
\definecolor{currentstroke}{rgb}{0.000000,0.000000,0.000000}%
\pgfsetstrokecolor{currentstroke}%
\pgfsetdash{}{0pt}%
\pgfpathmoveto{\pgfqpoint{1.256754in}{0.792144in}}%
\pgfpathlineto{\pgfqpoint{1.273838in}{0.802245in}}%
\pgfpathlineto{\pgfqpoint{1.290995in}{0.812698in}}%
\pgfpathlineto{\pgfqpoint{1.298992in}{0.814743in}}%
\pgfpathlineto{\pgfqpoint{1.306964in}{0.817310in}}%
\pgfpathlineto{\pgfqpoint{1.289831in}{0.806573in}}%
\pgfpathlineto{\pgfqpoint{1.272772in}{0.796163in}}%
\pgfpathlineto{\pgfqpoint{1.264775in}{0.793869in}}%
\pgfpathlineto{\pgfqpoint{1.256754in}{0.792144in}}%
\pgfpathclose%
\pgfusepath{fill}%
\end{pgfscope}%
\begin{pgfscope}%
\pgfpathrectangle{\pgfqpoint{0.407592in}{0.222904in}}{\pgfqpoint{1.411163in}{1.411163in}}%
\pgfusepath{clip}%
\pgfsetbuttcap%
\pgfsetroundjoin%
\definecolor{currentfill}{rgb}{0.556863,0.654902,0.847059}%
\pgfsetfillcolor{currentfill}%
\pgfsetfillopacity{0.130000}%
\pgfsetlinewidth{0.000000pt}%
\definecolor{currentstroke}{rgb}{0.000000,0.000000,0.000000}%
\pgfsetstrokecolor{currentstroke}%
\pgfsetdash{}{0pt}%
\pgfpathmoveto{\pgfqpoint{0.922297in}{0.785227in}}%
\pgfpathlineto{\pgfqpoint{0.938867in}{0.780611in}}%
\pgfpathlineto{\pgfqpoint{0.955468in}{0.777068in}}%
\pgfpathlineto{\pgfqpoint{0.964136in}{0.771233in}}%
\pgfpathlineto{\pgfqpoint{0.972764in}{0.765598in}}%
\pgfpathlineto{\pgfqpoint{0.956215in}{0.769072in}}%
\pgfpathlineto{\pgfqpoint{0.939697in}{0.773708in}}%
\pgfpathlineto{\pgfqpoint{0.931017in}{0.779369in}}%
\pgfpathlineto{\pgfqpoint{0.922297in}{0.785227in}}%
\pgfpathclose%
\pgfusepath{fill}%
\end{pgfscope}%
\begin{pgfscope}%
\pgfpathrectangle{\pgfqpoint{0.407592in}{0.222904in}}{\pgfqpoint{1.411163in}{1.411163in}}%
\pgfusepath{clip}%
\pgfsetbuttcap%
\pgfsetroundjoin%
\definecolor{currentfill}{rgb}{0.556863,0.654902,0.847059}%
\pgfsetfillcolor{currentfill}%
\pgfsetfillopacity{0.130000}%
\pgfsetlinewidth{0.000000pt}%
\definecolor{currentstroke}{rgb}{0.000000,0.000000,0.000000}%
\pgfsetstrokecolor{currentstroke}%
\pgfsetdash{}{0pt}%
\pgfpathmoveto{\pgfqpoint{0.637048in}{0.967665in}}%
\pgfpathlineto{\pgfqpoint{0.653533in}{0.953388in}}%
\pgfpathlineto{\pgfqpoint{0.670028in}{0.939582in}}%
\pgfpathlineto{\pgfqpoint{0.679297in}{0.934598in}}%
\pgfpathlineto{\pgfqpoint{0.688521in}{0.929816in}}%
\pgfpathlineto{\pgfqpoint{0.672071in}{0.944046in}}%
\pgfpathlineto{\pgfqpoint{0.655632in}{0.958746in}}%
\pgfpathlineto{\pgfqpoint{0.646362in}{0.963099in}}%
\pgfpathlineto{\pgfqpoint{0.637048in}{0.967665in}}%
\pgfpathclose%
\pgfusepath{fill}%
\end{pgfscope}%
\begin{pgfscope}%
\pgfpathrectangle{\pgfqpoint{0.407592in}{0.222904in}}{\pgfqpoint{1.411163in}{1.411163in}}%
\pgfusepath{clip}%
\pgfsetbuttcap%
\pgfsetroundjoin%
\definecolor{currentfill}{rgb}{0.556863,0.654902,0.847059}%
\pgfsetfillcolor{currentfill}%
\pgfsetfillopacity{0.130000}%
\pgfsetlinewidth{0.000000pt}%
\definecolor{currentstroke}{rgb}{0.000000,0.000000,0.000000}%
\pgfsetstrokecolor{currentstroke}%
\pgfsetdash{}{0pt}%
\pgfpathmoveto{\pgfqpoint{1.056278in}{0.753592in}}%
\pgfpathlineto{\pgfqpoint{1.072987in}{0.756019in}}%
\pgfpathlineto{\pgfqpoint{1.089748in}{0.759419in}}%
\pgfpathlineto{\pgfqpoint{1.098117in}{0.755711in}}%
\pgfpathlineto{\pgfqpoint{1.106449in}{0.752370in}}%
\pgfpathlineto{\pgfqpoint{1.089731in}{0.748417in}}%
\pgfpathlineto{\pgfqpoint{1.073067in}{0.745476in}}%
\pgfpathlineto{\pgfqpoint{1.064691in}{0.749360in}}%
\pgfpathlineto{\pgfqpoint{1.056278in}{0.753592in}}%
\pgfpathclose%
\pgfusepath{fill}%
\end{pgfscope}%
\begin{pgfscope}%
\pgfpathrectangle{\pgfqpoint{0.407592in}{0.222904in}}{\pgfqpoint{1.411163in}{1.411163in}}%
\pgfusepath{clip}%
\pgfsetbuttcap%
\pgfsetroundjoin%
\definecolor{currentfill}{rgb}{0.556863,0.654902,0.847059}%
\pgfsetfillcolor{currentfill}%
\pgfsetfillopacity{0.130000}%
\pgfsetlinewidth{0.000000pt}%
\definecolor{currentstroke}{rgb}{0.000000,0.000000,0.000000}%
\pgfsetstrokecolor{currentstroke}%
\pgfsetdash{}{0pt}%
\pgfpathmoveto{\pgfqpoint{0.754432in}{0.878371in}}%
\pgfpathlineto{\pgfqpoint{0.770941in}{0.867089in}}%
\pgfpathlineto{\pgfqpoint{0.787465in}{0.856524in}}%
\pgfpathlineto{\pgfqpoint{0.796504in}{0.850698in}}%
\pgfpathlineto{\pgfqpoint{0.805501in}{0.845049in}}%
\pgfpathlineto{\pgfqpoint{0.789025in}{0.855999in}}%
\pgfpathlineto{\pgfqpoint{0.772564in}{0.867691in}}%
\pgfpathlineto{\pgfqpoint{0.763519in}{0.872933in}}%
\pgfpathlineto{\pgfqpoint{0.754432in}{0.878371in}}%
\pgfpathclose%
\pgfusepath{fill}%
\end{pgfscope}%
\begin{pgfscope}%
\pgfpathrectangle{\pgfqpoint{0.407592in}{0.222904in}}{\pgfqpoint{1.411163in}{1.411163in}}%
\pgfusepath{clip}%
\pgfsetbuttcap%
\pgfsetroundjoin%
\definecolor{currentfill}{rgb}{0.556863,0.654902,0.847059}%
\pgfsetfillcolor{currentfill}%
\pgfsetfillopacity{0.130000}%
\pgfsetlinewidth{0.000000pt}%
\definecolor{currentstroke}{rgb}{0.000000,0.000000,0.000000}%
\pgfsetstrokecolor{currentstroke}%
\pgfsetdash{}{0pt}%
\pgfpathmoveto{\pgfqpoint{0.519315in}{1.068919in}}%
\pgfpathlineto{\pgfqpoint{0.535786in}{1.052811in}}%
\pgfpathlineto{\pgfqpoint{0.552267in}{1.037011in}}%
\pgfpathlineto{\pgfqpoint{0.561760in}{1.032825in}}%
\pgfpathlineto{\pgfqpoint{0.571208in}{1.028845in}}%
\pgfpathlineto{\pgfqpoint{0.554771in}{1.045015in}}%
\pgfpathlineto{\pgfqpoint{0.538344in}{1.061482in}}%
\pgfpathlineto{\pgfqpoint{0.528852in}{1.065095in}}%
\pgfpathlineto{\pgfqpoint{0.519315in}{1.068919in}}%
\pgfpathclose%
\pgfusepath{fill}%
\end{pgfscope}%
\begin{pgfscope}%
\pgfpathrectangle{\pgfqpoint{0.407592in}{0.222904in}}{\pgfqpoint{1.411163in}{1.411163in}}%
\pgfusepath{clip}%
\pgfsetbuttcap%
\pgfsetroundjoin%
\definecolor{currentfill}{rgb}{0.556863,0.654902,0.847059}%
\pgfsetfillcolor{currentfill}%
\pgfsetfillopacity{0.130000}%
\pgfsetlinewidth{0.000000pt}%
\definecolor{currentstroke}{rgb}{0.000000,0.000000,0.000000}%
\pgfsetstrokecolor{currentstroke}%
\pgfsetdash{}{0pt}%
\pgfpathmoveto{\pgfqpoint{1.306964in}{0.817310in}}%
\pgfpathlineto{\pgfqpoint{1.324173in}{0.828336in}}%
\pgfpathlineto{\pgfqpoint{1.341457in}{0.839620in}}%
\pgfpathlineto{\pgfqpoint{1.349384in}{0.842906in}}%
\pgfpathlineto{\pgfqpoint{1.357291in}{0.846689in}}%
\pgfpathlineto{\pgfqpoint{1.340029in}{0.835252in}}%
\pgfpathlineto{\pgfqpoint{1.322844in}{0.824057in}}%
\pgfpathlineto{\pgfqpoint{1.314915in}{0.820412in}}%
\pgfpathlineto{\pgfqpoint{1.306964in}{0.817310in}}%
\pgfpathclose%
\pgfusepath{fill}%
\end{pgfscope}%
\begin{pgfscope}%
\pgfpathrectangle{\pgfqpoint{0.407592in}{0.222904in}}{\pgfqpoint{1.411163in}{1.411163in}}%
\pgfusepath{clip}%
\pgfsetbuttcap%
\pgfsetroundjoin%
\definecolor{currentfill}{rgb}{0.556863,0.654902,0.847059}%
\pgfsetfillcolor{currentfill}%
\pgfsetfillopacity{0.130000}%
\pgfsetlinewidth{0.000000pt}%
\definecolor{currentstroke}{rgb}{0.000000,0.000000,0.000000}%
\pgfsetstrokecolor{currentstroke}%
\pgfsetdash{}{0pt}%
\pgfpathmoveto{\pgfqpoint{1.473898in}{0.927261in}}%
\pgfpathlineto{\pgfqpoint{1.491481in}{0.939518in}}%
\pgfpathlineto{\pgfqpoint{1.509144in}{0.951921in}}%
\pgfpathlineto{\pgfqpoint{1.516863in}{0.958714in}}%
\pgfpathlineto{\pgfqpoint{1.524568in}{0.965882in}}%
\pgfpathlineto{\pgfqpoint{1.506928in}{0.953627in}}%
\pgfpathlineto{\pgfqpoint{1.489366in}{0.941528in}}%
\pgfpathlineto{\pgfqpoint{1.481639in}{0.934193in}}%
\pgfpathlineto{\pgfqpoint{1.473898in}{0.927261in}}%
\pgfpathclose%
\pgfusepath{fill}%
\end{pgfscope}%
\begin{pgfscope}%
\pgfpathrectangle{\pgfqpoint{0.407592in}{0.222904in}}{\pgfqpoint{1.411163in}{1.411163in}}%
\pgfusepath{clip}%
\pgfsetbuttcap%
\pgfsetroundjoin%
\definecolor{currentfill}{rgb}{0.556863,0.654902,0.847059}%
\pgfsetfillcolor{currentfill}%
\pgfsetfillopacity{0.130000}%
\pgfsetlinewidth{0.000000pt}%
\definecolor{currentstroke}{rgb}{0.000000,0.000000,0.000000}%
\pgfsetstrokecolor{currentstroke}%
\pgfsetdash{}{0pt}%
\pgfpathmoveto{\pgfqpoint{1.590679in}{1.022282in}}%
\pgfpathlineto{\pgfqpoint{1.608512in}{1.034574in}}%
\pgfpathlineto{\pgfqpoint{1.626426in}{1.047008in}}%
\pgfpathlineto{\pgfqpoint{1.634018in}{1.055458in}}%
\pgfpathlineto{\pgfqpoint{1.641597in}{1.064199in}}%
\pgfpathlineto{\pgfqpoint{1.623707in}{1.051993in}}%
\pgfpathlineto{\pgfqpoint{1.605897in}{1.039942in}}%
\pgfpathlineto{\pgfqpoint{1.598294in}{1.030958in}}%
\pgfpathlineto{\pgfqpoint{1.590679in}{1.022282in}}%
\pgfpathclose%
\pgfusepath{fill}%
\end{pgfscope}%
\begin{pgfscope}%
\pgfpathrectangle{\pgfqpoint{0.407592in}{0.222904in}}{\pgfqpoint{1.411163in}{1.411163in}}%
\pgfusepath{clip}%
\pgfsetbuttcap%
\pgfsetroundjoin%
\definecolor{currentfill}{rgb}{0.556863,0.654902,0.847059}%
\pgfsetfillcolor{currentfill}%
\pgfsetfillopacity{0.130000}%
\pgfsetlinewidth{0.000000pt}%
\definecolor{currentstroke}{rgb}{0.000000,0.000000,0.000000}%
\pgfsetstrokecolor{currentstroke}%
\pgfsetdash{}{0pt}%
\pgfpathmoveto{\pgfqpoint{1.005969in}{0.762093in}}%
\pgfpathlineto{\pgfqpoint{1.022631in}{0.762003in}}%
\pgfpathlineto{\pgfqpoint{1.039339in}{0.762963in}}%
\pgfpathlineto{\pgfqpoint{1.047827in}{0.758137in}}%
\pgfpathlineto{\pgfqpoint{1.056278in}{0.753592in}}%
\pgfpathlineto{\pgfqpoint{1.039619in}{0.752215in}}%
\pgfpathlineto{\pgfqpoint{1.023005in}{0.751955in}}%
\pgfpathlineto{\pgfqpoint{1.014506in}{0.756894in}}%
\pgfpathlineto{\pgfqpoint{1.005969in}{0.762093in}}%
\pgfpathclose%
\pgfusepath{fill}%
\end{pgfscope}%
\begin{pgfscope}%
\pgfpathrectangle{\pgfqpoint{0.407592in}{0.222904in}}{\pgfqpoint{1.411163in}{1.411163in}}%
\pgfusepath{clip}%
\pgfsetbuttcap%
\pgfsetroundjoin%
\definecolor{currentfill}{rgb}{0.556863,0.654902,0.847059}%
\pgfsetfillcolor{currentfill}%
\pgfsetfillopacity{0.130000}%
\pgfsetlinewidth{0.000000pt}%
\definecolor{currentstroke}{rgb}{0.000000,0.000000,0.000000}%
\pgfsetstrokecolor{currentstroke}%
\pgfsetdash{}{0pt}%
\pgfpathmoveto{\pgfqpoint{0.670028in}{0.939582in}}%
\pgfpathlineto{\pgfqpoint{0.686534in}{0.926282in}}%
\pgfpathlineto{\pgfqpoint{0.703052in}{0.913523in}}%
\pgfpathlineto{\pgfqpoint{0.712274in}{0.908126in}}%
\pgfpathlineto{\pgfqpoint{0.721453in}{0.902916in}}%
\pgfpathlineto{\pgfqpoint{0.704981in}{0.916092in}}%
\pgfpathlineto{\pgfqpoint{0.688521in}{0.929816in}}%
\pgfpathlineto{\pgfqpoint{0.679297in}{0.934598in}}%
\pgfpathlineto{\pgfqpoint{0.670028in}{0.939582in}}%
\pgfpathclose%
\pgfusepath{fill}%
\end{pgfscope}%
\begin{pgfscope}%
\pgfpathrectangle{\pgfqpoint{0.407592in}{0.222904in}}{\pgfqpoint{1.411163in}{1.411163in}}%
\pgfusepath{clip}%
\pgfsetbuttcap%
\pgfsetroundjoin%
\definecolor{currentfill}{rgb}{0.556863,0.654902,0.847059}%
\pgfsetfillcolor{currentfill}%
\pgfsetfillopacity{0.130000}%
\pgfsetlinewidth{0.000000pt}%
\definecolor{currentstroke}{rgb}{0.000000,0.000000,0.000000}%
\pgfsetstrokecolor{currentstroke}%
\pgfsetdash{}{0pt}%
\pgfpathmoveto{\pgfqpoint{1.707721in}{1.125793in}}%
\pgfpathlineto{\pgfqpoint{1.725805in}{1.137956in}}%
\pgfpathlineto{\pgfqpoint{1.743970in}{1.150268in}}%
\pgfpathlineto{\pgfqpoint{1.751442in}{1.159832in}}%
\pgfpathlineto{\pgfqpoint{1.758903in}{1.169625in}}%
\pgfpathlineto{\pgfqpoint{1.740763in}{1.157563in}}%
\pgfpathlineto{\pgfqpoint{1.722703in}{1.145660in}}%
\pgfpathlineto{\pgfqpoint{1.715218in}{1.135608in}}%
\pgfpathlineto{\pgfqpoint{1.707721in}{1.125793in}}%
\pgfpathclose%
\pgfusepath{fill}%
\end{pgfscope}%
\begin{pgfscope}%
\pgfpathrectangle{\pgfqpoint{0.407592in}{0.222904in}}{\pgfqpoint{1.411163in}{1.411163in}}%
\pgfusepath{clip}%
\pgfsetbuttcap%
\pgfsetroundjoin%
\definecolor{currentfill}{rgb}{0.556863,0.654902,0.847059}%
\pgfsetfillcolor{currentfill}%
\pgfsetfillopacity{0.130000}%
\pgfsetlinewidth{0.000000pt}%
\definecolor{currentstroke}{rgb}{0.000000,0.000000,0.000000}%
\pgfsetstrokecolor{currentstroke}%
\pgfsetdash{}{0pt}%
\pgfpathmoveto{\pgfqpoint{0.552267in}{1.037011in}}%
\pgfpathlineto{\pgfqpoint{0.568757in}{1.021544in}}%
\pgfpathlineto{\pgfqpoint{0.585256in}{1.006433in}}%
\pgfpathlineto{\pgfqpoint{0.594705in}{1.001867in}}%
\pgfpathlineto{\pgfqpoint{0.604109in}{0.997501in}}%
\pgfpathlineto{\pgfqpoint{0.587654in}{1.012999in}}%
\pgfpathlineto{\pgfqpoint{0.571208in}{1.028845in}}%
\pgfpathlineto{\pgfqpoint{0.561760in}{1.032825in}}%
\pgfpathlineto{\pgfqpoint{0.552267in}{1.037011in}}%
\pgfpathclose%
\pgfusepath{fill}%
\end{pgfscope}%
\begin{pgfscope}%
\pgfpathrectangle{\pgfqpoint{0.407592in}{0.222904in}}{\pgfqpoint{1.411163in}{1.411163in}}%
\pgfusepath{clip}%
\pgfsetbuttcap%
\pgfsetroundjoin%
\definecolor{currentfill}{rgb}{0.556863,0.654902,0.847059}%
\pgfsetfillcolor{currentfill}%
\pgfsetfillopacity{0.130000}%
\pgfsetlinewidth{0.000000pt}%
\definecolor{currentstroke}{rgb}{0.000000,0.000000,0.000000}%
\pgfsetstrokecolor{currentstroke}%
\pgfsetdash{}{0pt}%
\pgfpathmoveto{\pgfqpoint{1.357291in}{0.846689in}}%
\pgfpathlineto{\pgfqpoint{1.374628in}{0.858345in}}%
\pgfpathlineto{\pgfqpoint{1.392043in}{0.870198in}}%
\pgfpathlineto{\pgfqpoint{1.399908in}{0.874554in}}%
\pgfpathlineto{\pgfqpoint{1.407754in}{0.879371in}}%
\pgfpathlineto{\pgfqpoint{1.390361in}{0.867471in}}%
\pgfpathlineto{\pgfqpoint{1.373045in}{0.855762in}}%
\pgfpathlineto{\pgfqpoint{1.365177in}{0.850974in}}%
\pgfpathlineto{\pgfqpoint{1.357291in}{0.846689in}}%
\pgfpathclose%
\pgfusepath{fill}%
\end{pgfscope}%
\begin{pgfscope}%
\pgfpathrectangle{\pgfqpoint{0.407592in}{0.222904in}}{\pgfqpoint{1.411163in}{1.411163in}}%
\pgfusepath{clip}%
\pgfsetbuttcap%
\pgfsetroundjoin%
\definecolor{currentfill}{rgb}{0.556863,0.654902,0.847059}%
\pgfsetfillcolor{currentfill}%
\pgfsetfillopacity{0.130000}%
\pgfsetlinewidth{0.000000pt}%
\definecolor{currentstroke}{rgb}{0.000000,0.000000,0.000000}%
\pgfsetstrokecolor{currentstroke}%
\pgfsetdash{}{0pt}%
\pgfpathmoveto{\pgfqpoint{0.871571in}{0.809619in}}%
\pgfpathlineto{\pgfqpoint{0.888140in}{0.803053in}}%
\pgfpathlineto{\pgfqpoint{0.904734in}{0.797463in}}%
\pgfpathlineto{\pgfqpoint{0.913536in}{0.791264in}}%
\pgfpathlineto{\pgfqpoint{0.922297in}{0.785227in}}%
\pgfpathlineto{\pgfqpoint{0.905755in}{0.790911in}}%
\pgfpathlineto{\pgfqpoint{0.889237in}{0.797644in}}%
\pgfpathlineto{\pgfqpoint{0.880425in}{0.803545in}}%
\pgfpathlineto{\pgfqpoint{0.871571in}{0.809619in}}%
\pgfpathclose%
\pgfusepath{fill}%
\end{pgfscope}%
\begin{pgfscope}%
\pgfpathrectangle{\pgfqpoint{0.407592in}{0.222904in}}{\pgfqpoint{1.411163in}{1.411163in}}%
\pgfusepath{clip}%
\pgfsetbuttcap%
\pgfsetroundjoin%
\definecolor{currentfill}{rgb}{0.556863,0.654902,0.847059}%
\pgfsetfillcolor{currentfill}%
\pgfsetfillopacity{0.130000}%
\pgfsetlinewidth{0.000000pt}%
\definecolor{currentstroke}{rgb}{0.000000,0.000000,0.000000}%
\pgfsetstrokecolor{currentstroke}%
\pgfsetdash{}{0pt}%
\pgfpathmoveto{\pgfqpoint{1.758903in}{1.169625in}}%
\pgfpathlineto{\pgfqpoint{1.777125in}{1.181840in}}%
\pgfpathlineto{\pgfqpoint{1.784562in}{1.191731in}}%
\pgfpathlineto{\pgfqpoint{1.791988in}{1.201836in}}%
\pgfpathlineto{\pgfqpoint{1.773791in}{1.189874in}}%
\pgfpathlineto{\pgfqpoint{1.766353in}{1.179642in}}%
\pgfusepath{fill}%
\end{pgfscope}%
\begin{pgfscope}%
\pgfpathrectangle{\pgfqpoint{0.407592in}{0.222904in}}{\pgfqpoint{1.411163in}{1.411163in}}%
\pgfusepath{clip}%
\pgfsetbuttcap%
\pgfsetroundjoin%
\definecolor{currentfill}{rgb}{0.556863,0.654902,0.847059}%
\pgfsetfillcolor{currentfill}%
\pgfsetfillopacity{0.130000}%
\pgfsetlinewidth{0.000000pt}%
\definecolor{currentstroke}{rgb}{0.000000,0.000000,0.000000}%
\pgfsetstrokecolor{currentstroke}%
\pgfsetdash{}{0pt}%
\pgfpathmoveto{\pgfqpoint{0.787465in}{0.856524in}}%
\pgfpathlineto{\pgfqpoint{0.804005in}{0.846717in}}%
\pgfpathlineto{\pgfqpoint{0.820562in}{0.837708in}}%
\pgfpathlineto{\pgfqpoint{0.829551in}{0.831561in}}%
\pgfpathlineto{\pgfqpoint{0.838498in}{0.825573in}}%
\pgfpathlineto{\pgfqpoint{0.821991in}{0.834892in}}%
\pgfpathlineto{\pgfqpoint{0.805501in}{0.845049in}}%
\pgfpathlineto{\pgfqpoint{0.796504in}{0.850698in}}%
\pgfpathlineto{\pgfqpoint{0.787465in}{0.856524in}}%
\pgfpathclose%
\pgfusepath{fill}%
\end{pgfscope}%
\begin{pgfscope}%
\pgfpathrectangle{\pgfqpoint{0.407592in}{0.222904in}}{\pgfqpoint{1.411163in}{1.411163in}}%
\pgfusepath{clip}%
\pgfsetbuttcap%
\pgfsetroundjoin%
\definecolor{currentfill}{rgb}{0.556863,0.654902,0.847059}%
\pgfsetfillcolor{currentfill}%
\pgfsetfillopacity{0.130000}%
\pgfsetlinewidth{0.000000pt}%
\definecolor{currentstroke}{rgb}{0.000000,0.000000,0.000000}%
\pgfsetstrokecolor{currentstroke}%
\pgfsetdash{}{0pt}%
\pgfpathmoveto{\pgfqpoint{1.524568in}{0.965882in}}%
\pgfpathlineto{\pgfqpoint{1.542288in}{0.978282in}}%
\pgfpathlineto{\pgfqpoint{1.560087in}{0.990819in}}%
\pgfpathlineto{\pgfqpoint{1.567755in}{0.998184in}}%
\pgfpathlineto{\pgfqpoint{1.575410in}{1.005887in}}%
\pgfpathlineto{\pgfqpoint{1.557633in}{0.993528in}}%
\pgfpathlineto{\pgfqpoint{1.539936in}{0.981316in}}%
\pgfpathlineto{\pgfqpoint{1.532259in}{0.973418in}}%
\pgfpathlineto{\pgfqpoint{1.524568in}{0.965882in}}%
\pgfpathclose%
\pgfusepath{fill}%
\end{pgfscope}%
\begin{pgfscope}%
\pgfpathrectangle{\pgfqpoint{0.407592in}{0.222904in}}{\pgfqpoint{1.411163in}{1.411163in}}%
\pgfusepath{clip}%
\pgfsetbuttcap%
\pgfsetroundjoin%
\definecolor{currentfill}{rgb}{0.556863,0.654902,0.847059}%
\pgfsetfillcolor{currentfill}%
\pgfsetfillopacity{0.130000}%
\pgfsetlinewidth{0.000000pt}%
\definecolor{currentstroke}{rgb}{0.000000,0.000000,0.000000}%
\pgfsetstrokecolor{currentstroke}%
\pgfsetdash{}{0pt}%
\pgfpathmoveto{\pgfqpoint{1.641597in}{1.064199in}}%
\pgfpathlineto{\pgfqpoint{1.659568in}{1.076553in}}%
\pgfpathlineto{\pgfqpoint{1.677619in}{1.089046in}}%
\pgfpathlineto{\pgfqpoint{1.685162in}{1.097842in}}%
\pgfpathlineto{\pgfqpoint{1.692694in}{1.106903in}}%
\pgfpathlineto{\pgfqpoint{1.674667in}{1.094641in}}%
\pgfpathlineto{\pgfqpoint{1.656720in}{1.082531in}}%
\pgfpathlineto{\pgfqpoint{1.649165in}{1.073227in}}%
\pgfpathlineto{\pgfqpoint{1.641597in}{1.064199in}}%
\pgfpathclose%
\pgfusepath{fill}%
\end{pgfscope}%
\begin{pgfscope}%
\pgfpathrectangle{\pgfqpoint{0.407592in}{0.222904in}}{\pgfqpoint{1.411163in}{1.411163in}}%
\pgfusepath{clip}%
\pgfsetbuttcap%
\pgfsetroundjoin%
\definecolor{currentfill}{rgb}{0.556863,0.654902,0.847059}%
\pgfsetfillcolor{currentfill}%
\pgfsetfillopacity{0.130000}%
\pgfsetlinewidth{0.000000pt}%
\definecolor{currentstroke}{rgb}{0.000000,0.000000,0.000000}%
\pgfsetstrokecolor{currentstroke}%
\pgfsetdash{}{0pt}%
\pgfpathmoveto{\pgfqpoint{0.955468in}{0.777068in}}%
\pgfpathlineto{\pgfqpoint{0.972103in}{0.774594in}}%
\pgfpathlineto{\pgfqpoint{0.988776in}{0.773171in}}%
\pgfpathlineto{\pgfqpoint{0.997392in}{0.767526in}}%
\pgfpathlineto{\pgfqpoint{1.005969in}{0.762093in}}%
\pgfpathlineto{\pgfqpoint{0.989348in}{0.763279in}}%
\pgfpathlineto{\pgfqpoint{0.972764in}{0.765598in}}%
\pgfpathlineto{\pgfqpoint{0.964136in}{0.771233in}}%
\pgfpathlineto{\pgfqpoint{0.955468in}{0.777068in}}%
\pgfpathclose%
\pgfusepath{fill}%
\end{pgfscope}%
\begin{pgfscope}%
\pgfpathrectangle{\pgfqpoint{0.407592in}{0.222904in}}{\pgfqpoint{1.411163in}{1.411163in}}%
\pgfusepath{clip}%
\pgfsetbuttcap%
\pgfsetroundjoin%
\definecolor{currentfill}{rgb}{0.556863,0.654902,0.847059}%
\pgfsetfillcolor{currentfill}%
\pgfsetfillopacity{0.130000}%
\pgfsetlinewidth{0.000000pt}%
\definecolor{currentstroke}{rgb}{0.000000,0.000000,0.000000}%
\pgfsetstrokecolor{currentstroke}%
\pgfsetdash{}{0pt}%
\pgfpathmoveto{\pgfqpoint{1.407754in}{0.879371in}}%
\pgfpathlineto{\pgfqpoint{1.425224in}{0.891446in}}%
\pgfpathlineto{\pgfqpoint{1.442773in}{0.903684in}}%
\pgfpathlineto{\pgfqpoint{1.450579in}{0.908946in}}%
\pgfpathlineto{\pgfqpoint{1.458368in}{0.914632in}}%
\pgfpathlineto{\pgfqpoint{1.440841in}{0.902429in}}%
\pgfpathlineto{\pgfqpoint{1.423392in}{0.890388in}}%
\pgfpathlineto{\pgfqpoint{1.415582in}{0.884650in}}%
\pgfpathlineto{\pgfqpoint{1.407754in}{0.879371in}}%
\pgfpathclose%
\pgfusepath{fill}%
\end{pgfscope}%
\begin{pgfscope}%
\pgfpathrectangle{\pgfqpoint{0.407592in}{0.222904in}}{\pgfqpoint{1.411163in}{1.411163in}}%
\pgfusepath{clip}%
\pgfsetbuttcap%
\pgfsetroundjoin%
\definecolor{currentfill}{rgb}{0.556863,0.654902,0.847059}%
\pgfsetfillcolor{currentfill}%
\pgfsetfillopacity{0.130000}%
\pgfsetlinewidth{0.000000pt}%
\definecolor{currentstroke}{rgb}{0.000000,0.000000,0.000000}%
\pgfsetstrokecolor{currentstroke}%
\pgfsetdash{}{0pt}%
\pgfpathmoveto{\pgfqpoint{0.585256in}{1.006433in}}%
\pgfpathlineto{\pgfqpoint{0.601765in}{0.991706in}}%
\pgfpathlineto{\pgfqpoint{0.618285in}{0.977389in}}%
\pgfpathlineto{\pgfqpoint{0.627689in}{0.972431in}}%
\pgfpathlineto{\pgfqpoint{0.637048in}{0.967665in}}%
\pgfpathlineto{\pgfqpoint{0.620574in}{0.982379in}}%
\pgfpathlineto{\pgfqpoint{0.604109in}{0.997501in}}%
\pgfpathlineto{\pgfqpoint{0.594705in}{1.001867in}}%
\pgfpathlineto{\pgfqpoint{0.585256in}{1.006433in}}%
\pgfpathclose%
\pgfusepath{fill}%
\end{pgfscope}%
\begin{pgfscope}%
\pgfpathrectangle{\pgfqpoint{0.407592in}{0.222904in}}{\pgfqpoint{1.411163in}{1.411163in}}%
\pgfusepath{clip}%
\pgfsetbuttcap%
\pgfsetroundjoin%
\definecolor{currentfill}{rgb}{0.556863,0.654902,0.847059}%
\pgfsetfillcolor{currentfill}%
\pgfsetfillopacity{0.130000}%
\pgfsetlinewidth{0.000000pt}%
\definecolor{currentstroke}{rgb}{0.000000,0.000000,0.000000}%
\pgfsetstrokecolor{currentstroke}%
\pgfsetdash{}{0pt}%
\pgfpathmoveto{\pgfqpoint{0.703052in}{0.913523in}}%
\pgfpathlineto{\pgfqpoint{0.719582in}{0.901344in}}%
\pgfpathlineto{\pgfqpoint{0.736126in}{0.889781in}}%
\pgfpathlineto{\pgfqpoint{0.745300in}{0.883991in}}%
\pgfpathlineto{\pgfqpoint{0.754432in}{0.878371in}}%
\pgfpathlineto{\pgfqpoint{0.737936in}{0.890328in}}%
\pgfpathlineto{\pgfqpoint{0.721453in}{0.902916in}}%
\pgfpathlineto{\pgfqpoint{0.712274in}{0.908126in}}%
\pgfpathlineto{\pgfqpoint{0.703052in}{0.913523in}}%
\pgfpathclose%
\pgfusepath{fill}%
\end{pgfscope}%
\begin{pgfscope}%
\pgfpathrectangle{\pgfqpoint{0.407592in}{0.222904in}}{\pgfqpoint{1.411163in}{1.411163in}}%
\pgfusepath{clip}%
\pgfsetbuttcap%
\pgfsetroundjoin%
\definecolor{currentfill}{rgb}{0.556863,0.654902,0.847059}%
\pgfsetfillcolor{currentfill}%
\pgfsetfillopacity{0.130000}%
\pgfsetlinewidth{0.000000pt}%
\definecolor{currentstroke}{rgb}{0.000000,0.000000,0.000000}%
\pgfsetstrokecolor{currentstroke}%
\pgfsetdash{}{0pt}%
\pgfpathmoveto{\pgfqpoint{1.190339in}{0.773441in}}%
\pgfpathlineto{\pgfqpoint{1.207337in}{0.781717in}}%
\pgfpathlineto{\pgfqpoint{1.224405in}{0.790509in}}%
\pgfpathlineto{\pgfqpoint{1.232534in}{0.790173in}}%
\pgfpathlineto{\pgfqpoint{1.240634in}{0.790318in}}%
\pgfpathlineto{\pgfqpoint{1.223596in}{0.781041in}}%
\pgfpathlineto{\pgfqpoint{1.206629in}{0.772249in}}%
\pgfpathlineto{\pgfqpoint{1.198499in}{0.772594in}}%
\pgfpathlineto{\pgfqpoint{1.190339in}{0.773441in}}%
\pgfpathclose%
\pgfusepath{fill}%
\end{pgfscope}%
\begin{pgfscope}%
\pgfpathrectangle{\pgfqpoint{0.407592in}{0.222904in}}{\pgfqpoint{1.411163in}{1.411163in}}%
\pgfusepath{clip}%
\pgfsetbuttcap%
\pgfsetroundjoin%
\definecolor{currentfill}{rgb}{0.556863,0.654902,0.847059}%
\pgfsetfillcolor{currentfill}%
\pgfsetfillopacity{0.130000}%
\pgfsetlinewidth{0.000000pt}%
\definecolor{currentstroke}{rgb}{0.000000,0.000000,0.000000}%
\pgfsetstrokecolor{currentstroke}%
\pgfsetdash{}{0pt}%
\pgfpathmoveto{\pgfqpoint{1.140062in}{0.762947in}}%
\pgfpathlineto{\pgfqpoint{1.156961in}{0.769400in}}%
\pgfpathlineto{\pgfqpoint{1.173926in}{0.776526in}}%
\pgfpathlineto{\pgfqpoint{1.182148in}{0.774762in}}%
\pgfpathlineto{\pgfqpoint{1.190339in}{0.773441in}}%
\pgfpathlineto{\pgfqpoint{1.173409in}{0.765747in}}%
\pgfpathlineto{\pgfqpoint{1.156545in}{0.758707in}}%
\pgfpathlineto{\pgfqpoint{1.148320in}{0.760603in}}%
\pgfpathlineto{\pgfqpoint{1.140062in}{0.762947in}}%
\pgfpathclose%
\pgfusepath{fill}%
\end{pgfscope}%
\begin{pgfscope}%
\pgfpathrectangle{\pgfqpoint{0.407592in}{0.222904in}}{\pgfqpoint{1.411163in}{1.411163in}}%
\pgfusepath{clip}%
\pgfsetbuttcap%
\pgfsetroundjoin%
\definecolor{currentfill}{rgb}{0.556863,0.654902,0.847059}%
\pgfsetfillcolor{currentfill}%
\pgfsetfillopacity{0.130000}%
\pgfsetlinewidth{0.000000pt}%
\definecolor{currentstroke}{rgb}{0.000000,0.000000,0.000000}%
\pgfsetstrokecolor{currentstroke}%
\pgfsetdash{}{0pt}%
\pgfpathmoveto{\pgfqpoint{1.240634in}{0.790318in}}%
\pgfpathlineto{\pgfqpoint{1.257745in}{0.800025in}}%
\pgfpathlineto{\pgfqpoint{1.274929in}{0.810113in}}%
\pgfpathlineto{\pgfqpoint{1.282975in}{0.811160in}}%
\pgfpathlineto{\pgfqpoint{1.290995in}{0.812698in}}%
\pgfpathlineto{\pgfqpoint{1.273838in}{0.802245in}}%
\pgfpathlineto{\pgfqpoint{1.256754in}{0.792144in}}%
\pgfpathlineto{\pgfqpoint{1.248707in}{0.790967in}}%
\pgfpathlineto{\pgfqpoint{1.240634in}{0.790318in}}%
\pgfpathclose%
\pgfusepath{fill}%
\end{pgfscope}%
\begin{pgfscope}%
\pgfpathrectangle{\pgfqpoint{0.407592in}{0.222904in}}{\pgfqpoint{1.411163in}{1.411163in}}%
\pgfusepath{clip}%
\pgfsetbuttcap%
\pgfsetroundjoin%
\definecolor{currentfill}{rgb}{0.556863,0.654902,0.847059}%
\pgfsetfillcolor{currentfill}%
\pgfsetfillopacity{0.130000}%
\pgfsetlinewidth{0.000000pt}%
\definecolor{currentstroke}{rgb}{0.000000,0.000000,0.000000}%
\pgfsetstrokecolor{currentstroke}%
\pgfsetdash{}{0pt}%
\pgfpathmoveto{\pgfqpoint{1.089748in}{0.759419in}}%
\pgfpathlineto{\pgfqpoint{1.106566in}{0.763719in}}%
\pgfpathlineto{\pgfqpoint{1.123443in}{0.768841in}}%
\pgfpathlineto{\pgfqpoint{1.131770in}{0.765704in}}%
\pgfpathlineto{\pgfqpoint{1.140062in}{0.762947in}}%
\pgfpathlineto{\pgfqpoint{1.123225in}{0.757243in}}%
\pgfpathlineto{\pgfqpoint{1.106449in}{0.752370in}}%
\pgfpathlineto{\pgfqpoint{1.098117in}{0.755711in}}%
\pgfpathlineto{\pgfqpoint{1.089748in}{0.759419in}}%
\pgfpathclose%
\pgfusepath{fill}%
\end{pgfscope}%
\begin{pgfscope}%
\pgfpathrectangle{\pgfqpoint{0.407592in}{0.222904in}}{\pgfqpoint{1.411163in}{1.411163in}}%
\pgfusepath{clip}%
\pgfsetbuttcap%
\pgfsetroundjoin%
\definecolor{currentfill}{rgb}{0.556863,0.654902,0.847059}%
\pgfsetfillcolor{currentfill}%
\pgfsetfillopacity{0.130000}%
\pgfsetlinewidth{0.000000pt}%
\definecolor{currentstroke}{rgb}{0.000000,0.000000,0.000000}%
\pgfsetstrokecolor{currentstroke}%
\pgfsetdash{}{0pt}%
\pgfpathmoveto{\pgfqpoint{1.575410in}{1.005887in}}%
\pgfpathlineto{\pgfqpoint{1.593266in}{1.018386in}}%
\pgfpathlineto{\pgfqpoint{1.611203in}{1.031016in}}%
\pgfpathlineto{\pgfqpoint{1.618821in}{1.038859in}}%
\pgfpathlineto{\pgfqpoint{1.626426in}{1.047008in}}%
\pgfpathlineto{\pgfqpoint{1.608512in}{1.034574in}}%
\pgfpathlineto{\pgfqpoint{1.590679in}{1.022282in}}%
\pgfpathlineto{\pgfqpoint{1.583051in}{1.013922in}}%
\pgfpathlineto{\pgfqpoint{1.575410in}{1.005887in}}%
\pgfpathclose%
\pgfusepath{fill}%
\end{pgfscope}%
\begin{pgfscope}%
\pgfpathrectangle{\pgfqpoint{0.407592in}{0.222904in}}{\pgfqpoint{1.411163in}{1.411163in}}%
\pgfusepath{clip}%
\pgfsetbuttcap%
\pgfsetroundjoin%
\definecolor{currentfill}{rgb}{0.556863,0.654902,0.847059}%
\pgfsetfillcolor{currentfill}%
\pgfsetfillopacity{0.130000}%
\pgfsetlinewidth{0.000000pt}%
\definecolor{currentstroke}{rgb}{0.000000,0.000000,0.000000}%
\pgfsetstrokecolor{currentstroke}%
\pgfsetdash{}{0pt}%
\pgfpathmoveto{\pgfqpoint{1.458368in}{0.914632in}}%
\pgfpathlineto{\pgfqpoint{1.475974in}{0.926987in}}%
\pgfpathlineto{\pgfqpoint{1.493658in}{0.939481in}}%
\pgfpathlineto{\pgfqpoint{1.501409in}{0.945508in}}%
\pgfpathlineto{\pgfqpoint{1.509144in}{0.951921in}}%
\pgfpathlineto{\pgfqpoint{1.491481in}{0.939518in}}%
\pgfpathlineto{\pgfqpoint{1.473898in}{0.927261in}}%
\pgfpathlineto{\pgfqpoint{1.466141in}{0.920739in}}%
\pgfpathlineto{\pgfqpoint{1.458368in}{0.914632in}}%
\pgfpathclose%
\pgfusepath{fill}%
\end{pgfscope}%
\begin{pgfscope}%
\pgfpathrectangle{\pgfqpoint{0.407592in}{0.222904in}}{\pgfqpoint{1.411163in}{1.411163in}}%
\pgfusepath{clip}%
\pgfsetbuttcap%
\pgfsetroundjoin%
\definecolor{currentfill}{rgb}{0.556863,0.654902,0.847059}%
\pgfsetfillcolor{currentfill}%
\pgfsetfillopacity{0.130000}%
\pgfsetlinewidth{0.000000pt}%
\definecolor{currentstroke}{rgb}{0.000000,0.000000,0.000000}%
\pgfsetstrokecolor{currentstroke}%
\pgfsetdash{}{0pt}%
\pgfpathmoveto{\pgfqpoint{1.290995in}{0.812698in}}%
\pgfpathlineto{\pgfqpoint{1.308228in}{0.823463in}}%
\pgfpathlineto{\pgfqpoint{1.325536in}{0.834508in}}%
\pgfpathlineto{\pgfqpoint{1.333508in}{0.836824in}}%
\pgfpathlineto{\pgfqpoint{1.341457in}{0.839620in}}%
\pgfpathlineto{\pgfqpoint{1.324173in}{0.828336in}}%
\pgfpathlineto{\pgfqpoint{1.306964in}{0.817310in}}%
\pgfpathlineto{\pgfqpoint{1.298992in}{0.814743in}}%
\pgfpathlineto{\pgfqpoint{1.290995in}{0.812698in}}%
\pgfpathclose%
\pgfusepath{fill}%
\end{pgfscope}%
\begin{pgfscope}%
\pgfpathrectangle{\pgfqpoint{0.407592in}{0.222904in}}{\pgfqpoint{1.411163in}{1.411163in}}%
\pgfusepath{clip}%
\pgfsetbuttcap%
\pgfsetroundjoin%
\definecolor{currentfill}{rgb}{0.556863,0.654902,0.847059}%
\pgfsetfillcolor{currentfill}%
\pgfsetfillopacity{0.130000}%
\pgfsetlinewidth{0.000000pt}%
\definecolor{currentstroke}{rgb}{0.000000,0.000000,0.000000}%
\pgfsetstrokecolor{currentstroke}%
\pgfsetdash{}{0pt}%
\pgfpathmoveto{\pgfqpoint{1.692694in}{1.106903in}}%
\pgfpathlineto{\pgfqpoint{1.710802in}{1.119308in}}%
\pgfpathlineto{\pgfqpoint{1.728992in}{1.131852in}}%
\pgfpathlineto{\pgfqpoint{1.736487in}{1.140939in}}%
\pgfpathlineto{\pgfqpoint{1.743970in}{1.150268in}}%
\pgfpathlineto{\pgfqpoint{1.725805in}{1.137956in}}%
\pgfpathlineto{\pgfqpoint{1.707721in}{1.125793in}}%
\pgfpathlineto{\pgfqpoint{1.700213in}{1.116222in}}%
\pgfpathlineto{\pgfqpoint{1.692694in}{1.106903in}}%
\pgfpathclose%
\pgfusepath{fill}%
\end{pgfscope}%
\begin{pgfscope}%
\pgfpathrectangle{\pgfqpoint{0.407592in}{0.222904in}}{\pgfqpoint{1.411163in}{1.411163in}}%
\pgfusepath{clip}%
\pgfsetbuttcap%
\pgfsetroundjoin%
\definecolor{currentfill}{rgb}{0.556863,0.654902,0.847059}%
\pgfsetfillcolor{currentfill}%
\pgfsetfillopacity{0.130000}%
\pgfsetlinewidth{0.000000pt}%
\definecolor{currentstroke}{rgb}{0.000000,0.000000,0.000000}%
\pgfsetstrokecolor{currentstroke}%
\pgfsetdash{}{0pt}%
\pgfpathmoveto{\pgfqpoint{0.820562in}{0.837708in}}%
\pgfpathlineto{\pgfqpoint{0.837139in}{0.829532in}}%
\pgfpathlineto{\pgfqpoint{0.853737in}{0.822223in}}%
\pgfpathlineto{\pgfqpoint{0.862675in}{0.815848in}}%
\pgfpathlineto{\pgfqpoint{0.871571in}{0.809619in}}%
\pgfpathlineto{\pgfqpoint{0.855024in}{0.817136in}}%
\pgfpathlineto{\pgfqpoint{0.838498in}{0.825573in}}%
\pgfpathlineto{\pgfqpoint{0.829551in}{0.831561in}}%
\pgfpathlineto{\pgfqpoint{0.820562in}{0.837708in}}%
\pgfpathclose%
\pgfusepath{fill}%
\end{pgfscope}%
\begin{pgfscope}%
\pgfpathrectangle{\pgfqpoint{0.407592in}{0.222904in}}{\pgfqpoint{1.411163in}{1.411163in}}%
\pgfusepath{clip}%
\pgfsetbuttcap%
\pgfsetroundjoin%
\definecolor{currentfill}{rgb}{0.556863,0.654902,0.847059}%
\pgfsetfillcolor{currentfill}%
\pgfsetfillopacity{0.130000}%
\pgfsetlinewidth{0.000000pt}%
\definecolor{currentstroke}{rgb}{0.000000,0.000000,0.000000}%
\pgfsetstrokecolor{currentstroke}%
\pgfsetdash{}{0pt}%
\pgfpathmoveto{\pgfqpoint{0.618285in}{0.977389in}}%
\pgfpathlineto{\pgfqpoint{0.634815in}{0.963512in}}%
\pgfpathlineto{\pgfqpoint{0.651356in}{0.950106in}}%
\pgfpathlineto{\pgfqpoint{0.660714in}{0.944754in}}%
\pgfpathlineto{\pgfqpoint{0.670028in}{0.939582in}}%
\pgfpathlineto{\pgfqpoint{0.653533in}{0.953388in}}%
\pgfpathlineto{\pgfqpoint{0.637048in}{0.967665in}}%
\pgfpathlineto{\pgfqpoint{0.627689in}{0.972431in}}%
\pgfpathlineto{\pgfqpoint{0.618285in}{0.977389in}}%
\pgfpathclose%
\pgfusepath{fill}%
\end{pgfscope}%
\begin{pgfscope}%
\pgfpathrectangle{\pgfqpoint{0.407592in}{0.222904in}}{\pgfqpoint{1.411163in}{1.411163in}}%
\pgfusepath{clip}%
\pgfsetbuttcap%
\pgfsetroundjoin%
\definecolor{currentfill}{rgb}{0.556863,0.654902,0.847059}%
\pgfsetfillcolor{currentfill}%
\pgfsetfillopacity{0.130000}%
\pgfsetlinewidth{0.000000pt}%
\definecolor{currentstroke}{rgb}{0.000000,0.000000,0.000000}%
\pgfsetstrokecolor{currentstroke}%
\pgfsetdash{}{0pt}%
\pgfpathmoveto{\pgfqpoint{1.039339in}{0.762963in}}%
\pgfpathlineto{\pgfqpoint{1.056095in}{0.764919in}}%
\pgfpathlineto{\pgfqpoint{1.072902in}{0.767809in}}%
\pgfpathlineto{\pgfqpoint{1.081344in}{0.763462in}}%
\pgfpathlineto{\pgfqpoint{1.089748in}{0.759419in}}%
\pgfpathlineto{\pgfqpoint{1.072987in}{0.756019in}}%
\pgfpathlineto{\pgfqpoint{1.056278in}{0.753592in}}%
\pgfpathlineto{\pgfqpoint{1.047827in}{0.758137in}}%
\pgfpathlineto{\pgfqpoint{1.039339in}{0.762963in}}%
\pgfpathclose%
\pgfusepath{fill}%
\end{pgfscope}%
\begin{pgfscope}%
\pgfpathrectangle{\pgfqpoint{0.407592in}{0.222904in}}{\pgfqpoint{1.411163in}{1.411163in}}%
\pgfusepath{clip}%
\pgfsetbuttcap%
\pgfsetroundjoin%
\definecolor{currentfill}{rgb}{0.556863,0.654902,0.847059}%
\pgfsetfillcolor{currentfill}%
\pgfsetfillopacity{0.130000}%
\pgfsetlinewidth{0.000000pt}%
\definecolor{currentstroke}{rgb}{0.000000,0.000000,0.000000}%
\pgfsetstrokecolor{currentstroke}%
\pgfsetdash{}{0pt}%
\pgfpathmoveto{\pgfqpoint{0.904734in}{0.797463in}}%
\pgfpathlineto{\pgfqpoint{0.921357in}{0.792865in}}%
\pgfpathlineto{\pgfqpoint{0.938010in}{0.789262in}}%
\pgfpathlineto{\pgfqpoint{0.946760in}{0.783083in}}%
\pgfpathlineto{\pgfqpoint{0.955468in}{0.777068in}}%
\pgfpathlineto{\pgfqpoint{0.938867in}{0.780611in}}%
\pgfpathlineto{\pgfqpoint{0.922297in}{0.785227in}}%
\pgfpathlineto{\pgfqpoint{0.913536in}{0.791264in}}%
\pgfpathlineto{\pgfqpoint{0.904734in}{0.797463in}}%
\pgfpathclose%
\pgfusepath{fill}%
\end{pgfscope}%
\begin{pgfscope}%
\pgfpathrectangle{\pgfqpoint{0.407592in}{0.222904in}}{\pgfqpoint{1.411163in}{1.411163in}}%
\pgfusepath{clip}%
\pgfsetbuttcap%
\pgfsetroundjoin%
\definecolor{currentfill}{rgb}{0.556863,0.654902,0.847059}%
\pgfsetfillcolor{currentfill}%
\pgfsetfillopacity{0.130000}%
\pgfsetlinewidth{0.000000pt}%
\definecolor{currentstroke}{rgb}{0.000000,0.000000,0.000000}%
\pgfsetstrokecolor{currentstroke}%
\pgfsetdash{}{0pt}%
\pgfpathmoveto{\pgfqpoint{1.743970in}{1.150268in}}%
\pgfpathlineto{\pgfqpoint{1.762217in}{1.162721in}}%
\pgfpathlineto{\pgfqpoint{1.769677in}{1.172168in}}%
\pgfpathlineto{\pgfqpoint{1.777125in}{1.181840in}}%
\pgfpathlineto{\pgfqpoint{1.758903in}{1.169625in}}%
\pgfpathlineto{\pgfqpoint{1.751442in}{1.159832in}}%
\pgfusepath{fill}%
\end{pgfscope}%
\begin{pgfscope}%
\pgfpathrectangle{\pgfqpoint{0.407592in}{0.222904in}}{\pgfqpoint{1.411163in}{1.411163in}}%
\pgfusepath{clip}%
\pgfsetbuttcap%
\pgfsetroundjoin%
\definecolor{currentfill}{rgb}{0.556863,0.654902,0.847059}%
\pgfsetfillcolor{currentfill}%
\pgfsetfillopacity{0.130000}%
\pgfsetlinewidth{0.000000pt}%
\definecolor{currentstroke}{rgb}{0.000000,0.000000,0.000000}%
\pgfsetstrokecolor{currentstroke}%
\pgfsetdash{}{0pt}%
\pgfpathmoveto{\pgfqpoint{0.736126in}{0.889781in}}%
\pgfpathlineto{\pgfqpoint{0.752683in}{0.878873in}}%
\pgfpathlineto{\pgfqpoint{0.769257in}{0.868655in}}%
\pgfpathlineto{\pgfqpoint{0.778382in}{0.862513in}}%
\pgfpathlineto{\pgfqpoint{0.787465in}{0.856524in}}%
\pgfpathlineto{\pgfqpoint{0.770941in}{0.867089in}}%
\pgfpathlineto{\pgfqpoint{0.754432in}{0.878371in}}%
\pgfpathlineto{\pgfqpoint{0.745300in}{0.883991in}}%
\pgfpathlineto{\pgfqpoint{0.736126in}{0.889781in}}%
\pgfpathclose%
\pgfusepath{fill}%
\end{pgfscope}%
\begin{pgfscope}%
\pgfpathrectangle{\pgfqpoint{0.407592in}{0.222904in}}{\pgfqpoint{1.411163in}{1.411163in}}%
\pgfusepath{clip}%
\pgfsetbuttcap%
\pgfsetroundjoin%
\definecolor{currentfill}{rgb}{0.556863,0.654902,0.847059}%
\pgfsetfillcolor{currentfill}%
\pgfsetfillopacity{0.130000}%
\pgfsetlinewidth{0.000000pt}%
\definecolor{currentstroke}{rgb}{0.000000,0.000000,0.000000}%
\pgfsetstrokecolor{currentstroke}%
\pgfsetdash{}{0pt}%
\pgfpathmoveto{\pgfqpoint{1.341457in}{0.839620in}}%
\pgfpathlineto{\pgfqpoint{1.358817in}{0.851136in}}%
\pgfpathlineto{\pgfqpoint{1.376255in}{0.862863in}}%
\pgfpathlineto{\pgfqpoint{1.384159in}{0.866302in}}%
\pgfpathlineto{\pgfqpoint{1.392043in}{0.870198in}}%
\pgfpathlineto{\pgfqpoint{1.374628in}{0.858345in}}%
\pgfpathlineto{\pgfqpoint{1.357291in}{0.846689in}}%
\pgfpathlineto{\pgfqpoint{1.349384in}{0.842906in}}%
\pgfpathlineto{\pgfqpoint{1.341457in}{0.839620in}}%
\pgfpathclose%
\pgfusepath{fill}%
\end{pgfscope}%
\begin{pgfscope}%
\pgfpathrectangle{\pgfqpoint{0.407592in}{0.222904in}}{\pgfqpoint{1.411163in}{1.411163in}}%
\pgfusepath{clip}%
\pgfsetbuttcap%
\pgfsetroundjoin%
\definecolor{currentfill}{rgb}{0.556863,0.654902,0.847059}%
\pgfsetfillcolor{currentfill}%
\pgfsetfillopacity{0.130000}%
\pgfsetlinewidth{0.000000pt}%
\definecolor{currentstroke}{rgb}{0.000000,0.000000,0.000000}%
\pgfsetstrokecolor{currentstroke}%
\pgfsetdash{}{0pt}%
\pgfpathmoveto{\pgfqpoint{1.509144in}{0.951921in}}%
\pgfpathlineto{\pgfqpoint{1.526886in}{0.964460in}}%
\pgfpathlineto{\pgfqpoint{1.544708in}{0.977130in}}%
\pgfpathlineto{\pgfqpoint{1.552405in}{0.983799in}}%
\pgfpathlineto{\pgfqpoint{1.560087in}{0.990819in}}%
\pgfpathlineto{\pgfqpoint{1.542288in}{0.978282in}}%
\pgfpathlineto{\pgfqpoint{1.524568in}{0.965882in}}%
\pgfpathlineto{\pgfqpoint{1.516863in}{0.958714in}}%
\pgfpathlineto{\pgfqpoint{1.509144in}{0.951921in}}%
\pgfpathclose%
\pgfusepath{fill}%
\end{pgfscope}%
\begin{pgfscope}%
\pgfpathrectangle{\pgfqpoint{0.407592in}{0.222904in}}{\pgfqpoint{1.411163in}{1.411163in}}%
\pgfusepath{clip}%
\pgfsetbuttcap%
\pgfsetroundjoin%
\definecolor{currentfill}{rgb}{0.556863,0.654902,0.847059}%
\pgfsetfillcolor{currentfill}%
\pgfsetfillopacity{0.130000}%
\pgfsetlinewidth{0.000000pt}%
\definecolor{currentstroke}{rgb}{0.000000,0.000000,0.000000}%
\pgfsetstrokecolor{currentstroke}%
\pgfsetdash{}{0pt}%
\pgfpathmoveto{\pgfqpoint{1.626426in}{1.047008in}}%
\pgfpathlineto{\pgfqpoint{1.644420in}{1.059579in}}%
\pgfpathlineto{\pgfqpoint{1.662495in}{1.072279in}}%
\pgfpathlineto{\pgfqpoint{1.670063in}{1.080523in}}%
\pgfpathlineto{\pgfqpoint{1.677619in}{1.089046in}}%
\pgfpathlineto{\pgfqpoint{1.659568in}{1.076553in}}%
\pgfpathlineto{\pgfqpoint{1.641597in}{1.064199in}}%
\pgfpathlineto{\pgfqpoint{1.634018in}{1.055458in}}%
\pgfpathlineto{\pgfqpoint{1.626426in}{1.047008in}}%
\pgfpathclose%
\pgfusepath{fill}%
\end{pgfscope}%
\begin{pgfscope}%
\pgfpathrectangle{\pgfqpoint{0.407592in}{0.222904in}}{\pgfqpoint{1.411163in}{1.411163in}}%
\pgfusepath{clip}%
\pgfsetbuttcap%
\pgfsetroundjoin%
\definecolor{currentfill}{rgb}{0.556863,0.654902,0.847059}%
\pgfsetfillcolor{currentfill}%
\pgfsetfillopacity{0.130000}%
\pgfsetlinewidth{0.000000pt}%
\definecolor{currentstroke}{rgb}{0.000000,0.000000,0.000000}%
\pgfsetstrokecolor{currentstroke}%
\pgfsetdash{}{0pt}%
\pgfpathmoveto{\pgfqpoint{0.988776in}{0.773171in}}%
\pgfpathlineto{\pgfqpoint{1.005489in}{0.772772in}}%
\pgfpathlineto{\pgfqpoint{1.022246in}{0.773360in}}%
\pgfpathlineto{\pgfqpoint{1.030812in}{0.768046in}}%
\pgfpathlineto{\pgfqpoint{1.039339in}{0.762963in}}%
\pgfpathlineto{\pgfqpoint{1.022631in}{0.762003in}}%
\pgfpathlineto{\pgfqpoint{1.005969in}{0.762093in}}%
\pgfpathlineto{\pgfqpoint{0.997392in}{0.767526in}}%
\pgfpathlineto{\pgfqpoint{0.988776in}{0.773171in}}%
\pgfpathclose%
\pgfusepath{fill}%
\end{pgfscope}%
\begin{pgfscope}%
\pgfpathrectangle{\pgfqpoint{0.407592in}{0.222904in}}{\pgfqpoint{1.411163in}{1.411163in}}%
\pgfusepath{clip}%
\pgfsetbuttcap%
\pgfsetroundjoin%
\definecolor{currentfill}{rgb}{0.556863,0.654902,0.847059}%
\pgfsetfillcolor{currentfill}%
\pgfsetfillopacity{0.130000}%
\pgfsetlinewidth{0.000000pt}%
\definecolor{currentstroke}{rgb}{0.000000,0.000000,0.000000}%
\pgfsetstrokecolor{currentstroke}%
\pgfsetdash{}{0pt}%
\pgfpathmoveto{\pgfqpoint{0.651356in}{0.950106in}}%
\pgfpathlineto{\pgfqpoint{0.667909in}{0.937202in}}%
\pgfpathlineto{\pgfqpoint{0.684474in}{0.924832in}}%
\pgfpathlineto{\pgfqpoint{0.693785in}{0.919095in}}%
\pgfpathlineto{\pgfqpoint{0.703052in}{0.913523in}}%
\pgfpathlineto{\pgfqpoint{0.686534in}{0.926282in}}%
\pgfpathlineto{\pgfqpoint{0.670028in}{0.939582in}}%
\pgfpathlineto{\pgfqpoint{0.660714in}{0.944754in}}%
\pgfpathlineto{\pgfqpoint{0.651356in}{0.950106in}}%
\pgfpathclose%
\pgfusepath{fill}%
\end{pgfscope}%
\begin{pgfscope}%
\pgfpathrectangle{\pgfqpoint{0.407592in}{0.222904in}}{\pgfqpoint{1.411163in}{1.411163in}}%
\pgfusepath{clip}%
\pgfsetbuttcap%
\pgfsetroundjoin%
\definecolor{currentfill}{rgb}{0.556863,0.654902,0.847059}%
\pgfsetfillcolor{currentfill}%
\pgfsetfillopacity{0.130000}%
\pgfsetlinewidth{0.000000pt}%
\definecolor{currentstroke}{rgb}{0.000000,0.000000,0.000000}%
\pgfsetstrokecolor{currentstroke}%
\pgfsetdash{}{0pt}%
\pgfpathmoveto{\pgfqpoint{1.392043in}{0.870198in}}%
\pgfpathlineto{\pgfqpoint{1.409536in}{0.882232in}}%
\pgfpathlineto{\pgfqpoint{1.427107in}{0.894433in}}%
\pgfpathlineto{\pgfqpoint{1.434950in}{0.898846in}}%
\pgfpathlineto{\pgfqpoint{1.442773in}{0.903684in}}%
\pgfpathlineto{\pgfqpoint{1.425224in}{0.891446in}}%
\pgfpathlineto{\pgfqpoint{1.407754in}{0.879371in}}%
\pgfpathlineto{\pgfqpoint{1.399908in}{0.874554in}}%
\pgfpathlineto{\pgfqpoint{1.392043in}{0.870198in}}%
\pgfpathclose%
\pgfusepath{fill}%
\end{pgfscope}%
\begin{pgfscope}%
\pgfpathrectangle{\pgfqpoint{0.407592in}{0.222904in}}{\pgfqpoint{1.411163in}{1.411163in}}%
\pgfusepath{clip}%
\pgfsetbuttcap%
\pgfsetroundjoin%
\definecolor{currentfill}{rgb}{0.556863,0.654902,0.847059}%
\pgfsetfillcolor{currentfill}%
\pgfsetfillopacity{0.130000}%
\pgfsetlinewidth{0.000000pt}%
\definecolor{currentstroke}{rgb}{0.000000,0.000000,0.000000}%
\pgfsetstrokecolor{currentstroke}%
\pgfsetdash{}{0pt}%
\pgfpathmoveto{\pgfqpoint{0.853737in}{0.822223in}}%
\pgfpathlineto{\pgfqpoint{0.870358in}{0.815804in}}%
\pgfpathlineto{\pgfqpoint{0.887006in}{0.810297in}}%
\pgfpathlineto{\pgfqpoint{0.895891in}{0.803812in}}%
\pgfpathlineto{\pgfqpoint{0.904734in}{0.797463in}}%
\pgfpathlineto{\pgfqpoint{0.888140in}{0.803053in}}%
\pgfpathlineto{\pgfqpoint{0.871571in}{0.809619in}}%
\pgfpathlineto{\pgfqpoint{0.862675in}{0.815848in}}%
\pgfpathlineto{\pgfqpoint{0.853737in}{0.822223in}}%
\pgfpathclose%
\pgfusepath{fill}%
\end{pgfscope}%
\begin{pgfscope}%
\pgfpathrectangle{\pgfqpoint{0.407592in}{0.222904in}}{\pgfqpoint{1.411163in}{1.411163in}}%
\pgfusepath{clip}%
\pgfsetbuttcap%
\pgfsetroundjoin%
\definecolor{currentfill}{rgb}{0.556863,0.654902,0.847059}%
\pgfsetfillcolor{currentfill}%
\pgfsetfillopacity{0.130000}%
\pgfsetlinewidth{0.000000pt}%
\definecolor{currentstroke}{rgb}{0.000000,0.000000,0.000000}%
\pgfsetstrokecolor{currentstroke}%
\pgfsetdash{}{0pt}%
\pgfpathmoveto{\pgfqpoint{0.769257in}{0.868655in}}%
\pgfpathlineto{\pgfqpoint{0.785847in}{0.859163in}}%
\pgfpathlineto{\pgfqpoint{0.802455in}{0.850431in}}%
\pgfpathlineto{\pgfqpoint{0.811530in}{0.844001in}}%
\pgfpathlineto{\pgfqpoint{0.820562in}{0.837708in}}%
\pgfpathlineto{\pgfqpoint{0.804005in}{0.846717in}}%
\pgfpathlineto{\pgfqpoint{0.787465in}{0.856524in}}%
\pgfpathlineto{\pgfqpoint{0.778382in}{0.862513in}}%
\pgfpathlineto{\pgfqpoint{0.769257in}{0.868655in}}%
\pgfpathclose%
\pgfusepath{fill}%
\end{pgfscope}%
\begin{pgfscope}%
\pgfpathrectangle{\pgfqpoint{0.407592in}{0.222904in}}{\pgfqpoint{1.411163in}{1.411163in}}%
\pgfusepath{clip}%
\pgfsetbuttcap%
\pgfsetroundjoin%
\definecolor{currentfill}{rgb}{0.556863,0.654902,0.847059}%
\pgfsetfillcolor{currentfill}%
\pgfsetfillopacity{0.130000}%
\pgfsetlinewidth{0.000000pt}%
\definecolor{currentstroke}{rgb}{0.000000,0.000000,0.000000}%
\pgfsetstrokecolor{currentstroke}%
\pgfsetdash{}{0pt}%
\pgfpathmoveto{\pgfqpoint{1.560087in}{0.990819in}}%
\pgfpathlineto{\pgfqpoint{1.577966in}{1.003486in}}%
\pgfpathlineto{\pgfqpoint{1.595927in}{1.016276in}}%
\pgfpathlineto{\pgfqpoint{1.603572in}{1.023487in}}%
\pgfpathlineto{\pgfqpoint{1.611203in}{1.031016in}}%
\pgfpathlineto{\pgfqpoint{1.593266in}{1.018386in}}%
\pgfpathlineto{\pgfqpoint{1.575410in}{1.005887in}}%
\pgfpathlineto{\pgfqpoint{1.567755in}{0.998184in}}%
\pgfpathlineto{\pgfqpoint{1.560087in}{0.990819in}}%
\pgfpathclose%
\pgfusepath{fill}%
\end{pgfscope}%
\begin{pgfscope}%
\pgfpathrectangle{\pgfqpoint{0.407592in}{0.222904in}}{\pgfqpoint{1.411163in}{1.411163in}}%
\pgfusepath{clip}%
\pgfsetbuttcap%
\pgfsetroundjoin%
\definecolor{currentfill}{rgb}{0.556863,0.654902,0.847059}%
\pgfsetfillcolor{currentfill}%
\pgfsetfillopacity{0.130000}%
\pgfsetlinewidth{0.000000pt}%
\definecolor{currentstroke}{rgb}{0.000000,0.000000,0.000000}%
\pgfsetstrokecolor{currentstroke}%
\pgfsetdash{}{0pt}%
\pgfpathmoveto{\pgfqpoint{1.677619in}{1.089046in}}%
\pgfpathlineto{\pgfqpoint{1.695752in}{1.101674in}}%
\pgfpathlineto{\pgfqpoint{1.713967in}{1.114429in}}%
\pgfpathlineto{\pgfqpoint{1.721485in}{1.123013in}}%
\pgfpathlineto{\pgfqpoint{1.728992in}{1.131852in}}%
\pgfpathlineto{\pgfqpoint{1.710802in}{1.119308in}}%
\pgfpathlineto{\pgfqpoint{1.692694in}{1.106903in}}%
\pgfpathlineto{\pgfqpoint{1.685162in}{1.097842in}}%
\pgfpathlineto{\pgfqpoint{1.677619in}{1.089046in}}%
\pgfpathclose%
\pgfusepath{fill}%
\end{pgfscope}%
\begin{pgfscope}%
\pgfpathrectangle{\pgfqpoint{0.407592in}{0.222904in}}{\pgfqpoint{1.411163in}{1.411163in}}%
\pgfusepath{clip}%
\pgfsetbuttcap%
\pgfsetroundjoin%
\definecolor{currentfill}{rgb}{0.556863,0.654902,0.847059}%
\pgfsetfillcolor{currentfill}%
\pgfsetfillopacity{0.130000}%
\pgfsetlinewidth{0.000000pt}%
\definecolor{currentstroke}{rgb}{0.000000,0.000000,0.000000}%
\pgfsetstrokecolor{currentstroke}%
\pgfsetdash{}{0pt}%
\pgfpathmoveto{\pgfqpoint{1.442773in}{0.903684in}}%
\pgfpathlineto{\pgfqpoint{1.460401in}{0.916071in}}%
\pgfpathlineto{\pgfqpoint{1.478108in}{0.928598in}}%
\pgfpathlineto{\pgfqpoint{1.485892in}{0.933844in}}%
\pgfpathlineto{\pgfqpoint{1.493658in}{0.939481in}}%
\pgfpathlineto{\pgfqpoint{1.475974in}{0.926987in}}%
\pgfpathlineto{\pgfqpoint{1.458368in}{0.914632in}}%
\pgfpathlineto{\pgfqpoint{1.450579in}{0.908946in}}%
\pgfpathlineto{\pgfqpoint{1.442773in}{0.903684in}}%
\pgfpathclose%
\pgfusepath{fill}%
\end{pgfscope}%
\begin{pgfscope}%
\pgfpathrectangle{\pgfqpoint{0.407592in}{0.222904in}}{\pgfqpoint{1.411163in}{1.411163in}}%
\pgfusepath{clip}%
\pgfsetbuttcap%
\pgfsetroundjoin%
\definecolor{currentfill}{rgb}{0.556863,0.654902,0.847059}%
\pgfsetfillcolor{currentfill}%
\pgfsetfillopacity{0.130000}%
\pgfsetlinewidth{0.000000pt}%
\definecolor{currentstroke}{rgb}{0.000000,0.000000,0.000000}%
\pgfsetstrokecolor{currentstroke}%
\pgfsetdash{}{0pt}%
\pgfpathmoveto{\pgfqpoint{1.173926in}{0.776526in}}%
\pgfpathlineto{\pgfqpoint{1.190958in}{0.784255in}}%
\pgfpathlineto{\pgfqpoint{1.208059in}{0.792526in}}%
\pgfpathlineto{\pgfqpoint{1.216247in}{0.791301in}}%
\pgfpathlineto{\pgfqpoint{1.224405in}{0.790509in}}%
\pgfpathlineto{\pgfqpoint{1.207337in}{0.781717in}}%
\pgfpathlineto{\pgfqpoint{1.190339in}{0.773441in}}%
\pgfpathlineto{\pgfqpoint{1.182148in}{0.774762in}}%
\pgfpathlineto{\pgfqpoint{1.173926in}{0.776526in}}%
\pgfpathclose%
\pgfusepath{fill}%
\end{pgfscope}%
\begin{pgfscope}%
\pgfpathrectangle{\pgfqpoint{0.407592in}{0.222904in}}{\pgfqpoint{1.411163in}{1.411163in}}%
\pgfusepath{clip}%
\pgfsetbuttcap%
\pgfsetroundjoin%
\definecolor{currentfill}{rgb}{0.556863,0.654902,0.847059}%
\pgfsetfillcolor{currentfill}%
\pgfsetfillopacity{0.130000}%
\pgfsetlinewidth{0.000000pt}%
\definecolor{currentstroke}{rgb}{0.000000,0.000000,0.000000}%
\pgfsetstrokecolor{currentstroke}%
\pgfsetdash{}{0pt}%
\pgfpathmoveto{\pgfqpoint{0.938010in}{0.789262in}}%
\pgfpathlineto{\pgfqpoint{0.954698in}{0.786652in}}%
\pgfpathlineto{\pgfqpoint{0.971422in}{0.785022in}}%
\pgfpathlineto{\pgfqpoint{0.980119in}{0.779009in}}%
\pgfpathlineto{\pgfqpoint{0.988776in}{0.773171in}}%
\pgfpathlineto{\pgfqpoint{0.972103in}{0.774594in}}%
\pgfpathlineto{\pgfqpoint{0.955468in}{0.777068in}}%
\pgfpathlineto{\pgfqpoint{0.946760in}{0.783083in}}%
\pgfpathlineto{\pgfqpoint{0.938010in}{0.789262in}}%
\pgfpathclose%
\pgfusepath{fill}%
\end{pgfscope}%
\begin{pgfscope}%
\pgfpathrectangle{\pgfqpoint{0.407592in}{0.222904in}}{\pgfqpoint{1.411163in}{1.411163in}}%
\pgfusepath{clip}%
\pgfsetbuttcap%
\pgfsetroundjoin%
\definecolor{currentfill}{rgb}{0.556863,0.654902,0.847059}%
\pgfsetfillcolor{currentfill}%
\pgfsetfillopacity{0.130000}%
\pgfsetlinewidth{0.000000pt}%
\definecolor{currentstroke}{rgb}{0.000000,0.000000,0.000000}%
\pgfsetstrokecolor{currentstroke}%
\pgfsetdash{}{0pt}%
\pgfpathmoveto{\pgfqpoint{1.224405in}{0.790509in}}%
\pgfpathlineto{\pgfqpoint{1.241545in}{0.799761in}}%
\pgfpathlineto{\pgfqpoint{1.258758in}{0.809424in}}%
\pgfpathlineto{\pgfqpoint{1.266857in}{0.809540in}}%
\pgfpathlineto{\pgfqpoint{1.274929in}{0.810113in}}%
\pgfpathlineto{\pgfqpoint{1.257745in}{0.800025in}}%
\pgfpathlineto{\pgfqpoint{1.240634in}{0.790318in}}%
\pgfpathlineto{\pgfqpoint{1.232534in}{0.790173in}}%
\pgfpathlineto{\pgfqpoint{1.224405in}{0.790509in}}%
\pgfpathclose%
\pgfusepath{fill}%
\end{pgfscope}%
\begin{pgfscope}%
\pgfpathrectangle{\pgfqpoint{0.407592in}{0.222904in}}{\pgfqpoint{1.411163in}{1.411163in}}%
\pgfusepath{clip}%
\pgfsetbuttcap%
\pgfsetroundjoin%
\definecolor{currentfill}{rgb}{0.556863,0.654902,0.847059}%
\pgfsetfillcolor{currentfill}%
\pgfsetfillopacity{0.130000}%
\pgfsetlinewidth{0.000000pt}%
\definecolor{currentstroke}{rgb}{0.000000,0.000000,0.000000}%
\pgfsetstrokecolor{currentstroke}%
\pgfsetdash{}{0pt}%
\pgfpathmoveto{\pgfqpoint{1.123443in}{0.768841in}}%
\pgfpathlineto{\pgfqpoint{1.140382in}{0.774713in}}%
\pgfpathlineto{\pgfqpoint{1.157384in}{0.781268in}}%
\pgfpathlineto{\pgfqpoint{1.165672in}{0.778704in}}%
\pgfpathlineto{\pgfqpoint{1.173926in}{0.776526in}}%
\pgfpathlineto{\pgfqpoint{1.156961in}{0.769400in}}%
\pgfpathlineto{\pgfqpoint{1.140062in}{0.762947in}}%
\pgfpathlineto{\pgfqpoint{1.131770in}{0.765704in}}%
\pgfpathlineto{\pgfqpoint{1.123443in}{0.768841in}}%
\pgfpathclose%
\pgfusepath{fill}%
\end{pgfscope}%
\begin{pgfscope}%
\pgfpathrectangle{\pgfqpoint{0.407592in}{0.222904in}}{\pgfqpoint{1.411163in}{1.411163in}}%
\pgfusepath{clip}%
\pgfsetbuttcap%
\pgfsetroundjoin%
\definecolor{currentfill}{rgb}{0.556863,0.654902,0.847059}%
\pgfsetfillcolor{currentfill}%
\pgfsetfillopacity{0.130000}%
\pgfsetlinewidth{0.000000pt}%
\definecolor{currentstroke}{rgb}{0.000000,0.000000,0.000000}%
\pgfsetstrokecolor{currentstroke}%
\pgfsetdash{}{0pt}%
\pgfpathmoveto{\pgfqpoint{1.728992in}{1.131852in}}%
\pgfpathlineto{\pgfqpoint{1.747264in}{1.144528in}}%
\pgfpathlineto{\pgfqpoint{1.754747in}{1.153506in}}%
\pgfpathlineto{\pgfqpoint{1.762217in}{1.162721in}}%
\pgfpathlineto{\pgfqpoint{1.743970in}{1.150268in}}%
\pgfpathlineto{\pgfqpoint{1.736487in}{1.140939in}}%
\pgfusepath{fill}%
\end{pgfscope}%
\begin{pgfscope}%
\pgfpathrectangle{\pgfqpoint{0.407592in}{0.222904in}}{\pgfqpoint{1.411163in}{1.411163in}}%
\pgfusepath{clip}%
\pgfsetbuttcap%
\pgfsetroundjoin%
\definecolor{currentfill}{rgb}{0.556863,0.654902,0.847059}%
\pgfsetfillcolor{currentfill}%
\pgfsetfillopacity{0.130000}%
\pgfsetlinewidth{0.000000pt}%
\definecolor{currentstroke}{rgb}{0.000000,0.000000,0.000000}%
\pgfsetstrokecolor{currentstroke}%
\pgfsetdash{}{0pt}%
\pgfpathmoveto{\pgfqpoint{0.684474in}{0.924832in}}%
\pgfpathlineto{\pgfqpoint{0.701052in}{0.913030in}}%
\pgfpathlineto{\pgfqpoint{0.717644in}{0.901828in}}%
\pgfpathlineto{\pgfqpoint{0.726907in}{0.895730in}}%
\pgfpathlineto{\pgfqpoint{0.736126in}{0.889781in}}%
\pgfpathlineto{\pgfqpoint{0.719582in}{0.901344in}}%
\pgfpathlineto{\pgfqpoint{0.703052in}{0.913523in}}%
\pgfpathlineto{\pgfqpoint{0.693785in}{0.919095in}}%
\pgfpathlineto{\pgfqpoint{0.684474in}{0.924832in}}%
\pgfpathclose%
\pgfusepath{fill}%
\end{pgfscope}%
\begin{pgfscope}%
\pgfpathrectangle{\pgfqpoint{0.407592in}{0.222904in}}{\pgfqpoint{1.411163in}{1.411163in}}%
\pgfusepath{clip}%
\pgfsetbuttcap%
\pgfsetroundjoin%
\definecolor{currentfill}{rgb}{0.556863,0.654902,0.847059}%
\pgfsetfillcolor{currentfill}%
\pgfsetfillopacity{0.130000}%
\pgfsetlinewidth{0.000000pt}%
\definecolor{currentstroke}{rgb}{0.000000,0.000000,0.000000}%
\pgfsetstrokecolor{currentstroke}%
\pgfsetdash{}{0pt}%
\pgfpathmoveto{\pgfqpoint{1.274929in}{0.810113in}}%
\pgfpathlineto{\pgfqpoint{1.292188in}{0.820541in}}%
\pgfpathlineto{\pgfqpoint{1.309521in}{0.831275in}}%
\pgfpathlineto{\pgfqpoint{1.317541in}{0.832663in}}%
\pgfpathlineto{\pgfqpoint{1.325536in}{0.834508in}}%
\pgfpathlineto{\pgfqpoint{1.308228in}{0.823463in}}%
\pgfpathlineto{\pgfqpoint{1.290995in}{0.812698in}}%
\pgfpathlineto{\pgfqpoint{1.282975in}{0.811160in}}%
\pgfpathlineto{\pgfqpoint{1.274929in}{0.810113in}}%
\pgfpathclose%
\pgfusepath{fill}%
\end{pgfscope}%
\begin{pgfscope}%
\pgfpathrectangle{\pgfqpoint{0.407592in}{0.222904in}}{\pgfqpoint{1.411163in}{1.411163in}}%
\pgfusepath{clip}%
\pgfsetbuttcap%
\pgfsetroundjoin%
\definecolor{currentfill}{rgb}{0.556863,0.654902,0.847059}%
\pgfsetfillcolor{currentfill}%
\pgfsetfillopacity{0.130000}%
\pgfsetlinewidth{0.000000pt}%
\definecolor{currentstroke}{rgb}{0.000000,0.000000,0.000000}%
\pgfsetstrokecolor{currentstroke}%
\pgfsetdash{}{0pt}%
\pgfpathmoveto{\pgfqpoint{1.072902in}{0.767809in}}%
\pgfpathlineto{\pgfqpoint{1.089764in}{0.771570in}}%
\pgfpathlineto{\pgfqpoint{1.106684in}{0.776140in}}%
\pgfpathlineto{\pgfqpoint{1.115081in}{0.772329in}}%
\pgfpathlineto{\pgfqpoint{1.123443in}{0.768841in}}%
\pgfpathlineto{\pgfqpoint{1.106566in}{0.763719in}}%
\pgfpathlineto{\pgfqpoint{1.089748in}{0.759419in}}%
\pgfpathlineto{\pgfqpoint{1.081344in}{0.763462in}}%
\pgfpathlineto{\pgfqpoint{1.072902in}{0.767809in}}%
\pgfpathclose%
\pgfusepath{fill}%
\end{pgfscope}%
\begin{pgfscope}%
\pgfpathrectangle{\pgfqpoint{0.407592in}{0.222904in}}{\pgfqpoint{1.411163in}{1.411163in}}%
\pgfusepath{clip}%
\pgfsetbuttcap%
\pgfsetroundjoin%
\definecolor{currentfill}{rgb}{0.556863,0.654902,0.847059}%
\pgfsetfillcolor{currentfill}%
\pgfsetfillopacity{0.130000}%
\pgfsetlinewidth{0.000000pt}%
\definecolor{currentstroke}{rgb}{0.000000,0.000000,0.000000}%
\pgfsetstrokecolor{currentstroke}%
\pgfsetdash{}{0pt}%
\pgfpathmoveto{\pgfqpoint{1.611203in}{1.031016in}}%
\pgfpathlineto{\pgfqpoint{1.629221in}{1.043773in}}%
\pgfpathlineto{\pgfqpoint{1.647320in}{1.056650in}}%
\pgfpathlineto{\pgfqpoint{1.654914in}{1.064319in}}%
\pgfpathlineto{\pgfqpoint{1.662495in}{1.072279in}}%
\pgfpathlineto{\pgfqpoint{1.644420in}{1.059579in}}%
\pgfpathlineto{\pgfqpoint{1.626426in}{1.047008in}}%
\pgfpathlineto{\pgfqpoint{1.618821in}{1.038859in}}%
\pgfpathlineto{\pgfqpoint{1.611203in}{1.031016in}}%
\pgfpathclose%
\pgfusepath{fill}%
\end{pgfscope}%
\begin{pgfscope}%
\pgfpathrectangle{\pgfqpoint{0.407592in}{0.222904in}}{\pgfqpoint{1.411163in}{1.411163in}}%
\pgfusepath{clip}%
\pgfsetbuttcap%
\pgfsetroundjoin%
\definecolor{currentfill}{rgb}{0.556863,0.654902,0.847059}%
\pgfsetfillcolor{currentfill}%
\pgfsetfillopacity{0.130000}%
\pgfsetlinewidth{0.000000pt}%
\definecolor{currentstroke}{rgb}{0.000000,0.000000,0.000000}%
\pgfsetstrokecolor{currentstroke}%
\pgfsetdash{}{0pt}%
\pgfpathmoveto{\pgfqpoint{1.493658in}{0.939481in}}%
\pgfpathlineto{\pgfqpoint{1.511423in}{0.952108in}}%
\pgfpathlineto{\pgfqpoint{1.529268in}{0.964860in}}%
\pgfpathlineto{\pgfqpoint{1.536995in}{0.970815in}}%
\pgfpathlineto{\pgfqpoint{1.544708in}{0.977130in}}%
\pgfpathlineto{\pgfqpoint{1.526886in}{0.964460in}}%
\pgfpathlineto{\pgfqpoint{1.509144in}{0.951921in}}%
\pgfpathlineto{\pgfqpoint{1.501409in}{0.945508in}}%
\pgfpathlineto{\pgfqpoint{1.493658in}{0.939481in}}%
\pgfpathclose%
\pgfusepath{fill}%
\end{pgfscope}%
\begin{pgfscope}%
\pgfpathrectangle{\pgfqpoint{0.407592in}{0.222904in}}{\pgfqpoint{1.411163in}{1.411163in}}%
\pgfusepath{clip}%
\pgfsetbuttcap%
\pgfsetroundjoin%
\definecolor{currentfill}{rgb}{0.556863,0.654902,0.847059}%
\pgfsetfillcolor{currentfill}%
\pgfsetfillopacity{0.130000}%
\pgfsetlinewidth{0.000000pt}%
\definecolor{currentstroke}{rgb}{0.000000,0.000000,0.000000}%
\pgfsetstrokecolor{currentstroke}%
\pgfsetdash{}{0pt}%
\pgfpathmoveto{\pgfqpoint{1.325536in}{0.834508in}}%
\pgfpathlineto{\pgfqpoint{1.342921in}{0.845805in}}%
\pgfpathlineto{\pgfqpoint{1.360383in}{0.857330in}}%
\pgfpathlineto{\pgfqpoint{1.368330in}{0.859874in}}%
\pgfpathlineto{\pgfqpoint{1.376255in}{0.862863in}}%
\pgfpathlineto{\pgfqpoint{1.358817in}{0.851136in}}%
\pgfpathlineto{\pgfqpoint{1.341457in}{0.839620in}}%
\pgfpathlineto{\pgfqpoint{1.333508in}{0.836824in}}%
\pgfpathlineto{\pgfqpoint{1.325536in}{0.834508in}}%
\pgfpathclose%
\pgfusepath{fill}%
\end{pgfscope}%
\begin{pgfscope}%
\pgfpathrectangle{\pgfqpoint{0.407592in}{0.222904in}}{\pgfqpoint{1.411163in}{1.411163in}}%
\pgfusepath{clip}%
\pgfsetbuttcap%
\pgfsetroundjoin%
\definecolor{currentfill}{rgb}{0.556863,0.654902,0.847059}%
\pgfsetfillcolor{currentfill}%
\pgfsetfillopacity{0.130000}%
\pgfsetlinewidth{0.000000pt}%
\definecolor{currentstroke}{rgb}{0.000000,0.000000,0.000000}%
\pgfsetstrokecolor{currentstroke}%
\pgfsetdash{}{0pt}%
\pgfpathmoveto{\pgfqpoint{0.802455in}{0.850431in}}%
\pgfpathlineto{\pgfqpoint{0.819083in}{0.842488in}}%
\pgfpathlineto{\pgfqpoint{0.835734in}{0.835362in}}%
\pgfpathlineto{\pgfqpoint{0.844757in}{0.828730in}}%
\pgfpathlineto{\pgfqpoint{0.853737in}{0.822223in}}%
\pgfpathlineto{\pgfqpoint{0.837139in}{0.829532in}}%
\pgfpathlineto{\pgfqpoint{0.820562in}{0.837708in}}%
\pgfpathlineto{\pgfqpoint{0.811530in}{0.844001in}}%
\pgfpathlineto{\pgfqpoint{0.802455in}{0.850431in}}%
\pgfpathclose%
\pgfusepath{fill}%
\end{pgfscope}%
\begin{pgfscope}%
\pgfpathrectangle{\pgfqpoint{0.407592in}{0.222904in}}{\pgfqpoint{1.411163in}{1.411163in}}%
\pgfusepath{clip}%
\pgfsetbuttcap%
\pgfsetroundjoin%
\definecolor{currentfill}{rgb}{0.556863,0.654902,0.847059}%
\pgfsetfillcolor{currentfill}%
\pgfsetfillopacity{0.130000}%
\pgfsetlinewidth{0.000000pt}%
\definecolor{currentstroke}{rgb}{0.000000,0.000000,0.000000}%
\pgfsetstrokecolor{currentstroke}%
\pgfsetdash{}{0pt}%
\pgfpathmoveto{\pgfqpoint{1.022246in}{0.773360in}}%
\pgfpathlineto{\pgfqpoint{1.039050in}{0.774890in}}%
\pgfpathlineto{\pgfqpoint{1.055905in}{0.777314in}}%
\pgfpathlineto{\pgfqpoint{1.064423in}{0.772434in}}%
\pgfpathlineto{\pgfqpoint{1.072902in}{0.767809in}}%
\pgfpathlineto{\pgfqpoint{1.056095in}{0.764919in}}%
\pgfpathlineto{\pgfqpoint{1.039339in}{0.762963in}}%
\pgfpathlineto{\pgfqpoint{1.030812in}{0.768046in}}%
\pgfpathlineto{\pgfqpoint{1.022246in}{0.773360in}}%
\pgfpathclose%
\pgfusepath{fill}%
\end{pgfscope}%
\begin{pgfscope}%
\pgfpathrectangle{\pgfqpoint{0.407592in}{0.222904in}}{\pgfqpoint{1.411163in}{1.411163in}}%
\pgfusepath{clip}%
\pgfsetbuttcap%
\pgfsetroundjoin%
\definecolor{currentfill}{rgb}{0.556863,0.654902,0.847059}%
\pgfsetfillcolor{currentfill}%
\pgfsetfillopacity{0.130000}%
\pgfsetlinewidth{0.000000pt}%
\definecolor{currentstroke}{rgb}{0.000000,0.000000,0.000000}%
\pgfsetstrokecolor{currentstroke}%
\pgfsetdash{}{0pt}%
\pgfpathmoveto{\pgfqpoint{0.887006in}{0.810297in}}%
\pgfpathlineto{\pgfqpoint{0.903681in}{0.805714in}}%
\pgfpathlineto{\pgfqpoint{0.920388in}{0.802059in}}%
\pgfpathlineto{\pgfqpoint{0.929220in}{0.795592in}}%
\pgfpathlineto{\pgfqpoint{0.938010in}{0.789262in}}%
\pgfpathlineto{\pgfqpoint{0.921357in}{0.792865in}}%
\pgfpathlineto{\pgfqpoint{0.904734in}{0.797463in}}%
\pgfpathlineto{\pgfqpoint{0.895891in}{0.803812in}}%
\pgfpathlineto{\pgfqpoint{0.887006in}{0.810297in}}%
\pgfpathclose%
\pgfusepath{fill}%
\end{pgfscope}%
\begin{pgfscope}%
\pgfpathrectangle{\pgfqpoint{0.407592in}{0.222904in}}{\pgfqpoint{1.411163in}{1.411163in}}%
\pgfusepath{clip}%
\pgfsetbuttcap%
\pgfsetroundjoin%
\definecolor{currentfill}{rgb}{0.556863,0.654902,0.847059}%
\pgfsetfillcolor{currentfill}%
\pgfsetfillopacity{0.130000}%
\pgfsetlinewidth{0.000000pt}%
\definecolor{currentstroke}{rgb}{0.000000,0.000000,0.000000}%
\pgfsetstrokecolor{currentstroke}%
\pgfsetdash{}{0pt}%
\pgfpathmoveto{\pgfqpoint{0.717644in}{0.901828in}}%
\pgfpathlineto{\pgfqpoint{0.734251in}{0.891260in}}%
\pgfpathlineto{\pgfqpoint{0.750874in}{0.881356in}}%
\pgfpathlineto{\pgfqpoint{0.760087in}{0.874939in}}%
\pgfpathlineto{\pgfqpoint{0.769257in}{0.868655in}}%
\pgfpathlineto{\pgfqpoint{0.752683in}{0.878873in}}%
\pgfpathlineto{\pgfqpoint{0.736126in}{0.889781in}}%
\pgfpathlineto{\pgfqpoint{0.726907in}{0.895730in}}%
\pgfpathlineto{\pgfqpoint{0.717644in}{0.901828in}}%
\pgfpathclose%
\pgfusepath{fill}%
\end{pgfscope}%
\begin{pgfscope}%
\pgfpathrectangle{\pgfqpoint{0.407592in}{0.222904in}}{\pgfqpoint{1.411163in}{1.411163in}}%
\pgfusepath{clip}%
\pgfsetbuttcap%
\pgfsetroundjoin%
\definecolor{currentfill}{rgb}{0.556863,0.654902,0.847059}%
\pgfsetfillcolor{currentfill}%
\pgfsetfillopacity{0.130000}%
\pgfsetlinewidth{0.000000pt}%
\definecolor{currentstroke}{rgb}{0.000000,0.000000,0.000000}%
\pgfsetstrokecolor{currentstroke}%
\pgfsetdash{}{0pt}%
\pgfpathmoveto{\pgfqpoint{1.376255in}{0.862863in}}%
\pgfpathlineto{\pgfqpoint{1.393771in}{0.874781in}}%
\pgfpathlineto{\pgfqpoint{1.411366in}{0.886876in}}%
\pgfpathlineto{\pgfqpoint{1.419246in}{0.890444in}}%
\pgfpathlineto{\pgfqpoint{1.427107in}{0.894433in}}%
\pgfpathlineto{\pgfqpoint{1.409536in}{0.882232in}}%
\pgfpathlineto{\pgfqpoint{1.392043in}{0.870198in}}%
\pgfpathlineto{\pgfqpoint{1.384159in}{0.866302in}}%
\pgfpathlineto{\pgfqpoint{1.376255in}{0.862863in}}%
\pgfpathclose%
\pgfusepath{fill}%
\end{pgfscope}%
\begin{pgfscope}%
\pgfpathrectangle{\pgfqpoint{0.407592in}{0.222904in}}{\pgfqpoint{1.411163in}{1.411163in}}%
\pgfusepath{clip}%
\pgfsetbuttcap%
\pgfsetroundjoin%
\definecolor{currentfill}{rgb}{0.556863,0.654902,0.847059}%
\pgfsetfillcolor{currentfill}%
\pgfsetfillopacity{0.130000}%
\pgfsetlinewidth{0.000000pt}%
\definecolor{currentstroke}{rgb}{0.000000,0.000000,0.000000}%
\pgfsetstrokecolor{currentstroke}%
\pgfsetdash{}{0pt}%
\pgfpathmoveto{\pgfqpoint{1.544708in}{0.977130in}}%
\pgfpathlineto{\pgfqpoint{1.562610in}{0.989921in}}%
\pgfpathlineto{\pgfqpoint{1.580594in}{1.002830in}}%
\pgfpathlineto{\pgfqpoint{1.588267in}{1.009389in}}%
\pgfpathlineto{\pgfqpoint{1.595927in}{1.016276in}}%
\pgfpathlineto{\pgfqpoint{1.577966in}{1.003486in}}%
\pgfpathlineto{\pgfqpoint{1.560087in}{0.990819in}}%
\pgfpathlineto{\pgfqpoint{1.552405in}{0.983799in}}%
\pgfpathlineto{\pgfqpoint{1.544708in}{0.977130in}}%
\pgfpathclose%
\pgfusepath{fill}%
\end{pgfscope}%
\begin{pgfscope}%
\pgfpathrectangle{\pgfqpoint{0.407592in}{0.222904in}}{\pgfqpoint{1.411163in}{1.411163in}}%
\pgfusepath{clip}%
\pgfsetbuttcap%
\pgfsetroundjoin%
\definecolor{currentfill}{rgb}{0.556863,0.654902,0.847059}%
\pgfsetfillcolor{currentfill}%
\pgfsetfillopacity{0.130000}%
\pgfsetlinewidth{0.000000pt}%
\definecolor{currentstroke}{rgb}{0.000000,0.000000,0.000000}%
\pgfsetstrokecolor{currentstroke}%
\pgfsetdash{}{0pt}%
\pgfpathmoveto{\pgfqpoint{1.662495in}{1.072279in}}%
\pgfpathlineto{\pgfqpoint{1.680652in}{1.085103in}}%
\pgfpathlineto{\pgfqpoint{1.698892in}{1.098045in}}%
\pgfpathlineto{\pgfqpoint{1.706436in}{1.106104in}}%
\pgfpathlineto{\pgfqpoint{1.713967in}{1.114429in}}%
\pgfpathlineto{\pgfqpoint{1.695752in}{1.101674in}}%
\pgfpathlineto{\pgfqpoint{1.677619in}{1.089046in}}%
\pgfpathlineto{\pgfqpoint{1.670063in}{1.080523in}}%
\pgfpathlineto{\pgfqpoint{1.662495in}{1.072279in}}%
\pgfpathclose%
\pgfusepath{fill}%
\end{pgfscope}%
\begin{pgfscope}%
\pgfpathrectangle{\pgfqpoint{0.407592in}{0.222904in}}{\pgfqpoint{1.411163in}{1.411163in}}%
\pgfusepath{clip}%
\pgfsetbuttcap%
\pgfsetroundjoin%
\definecolor{currentfill}{rgb}{0.556863,0.654902,0.847059}%
\pgfsetfillcolor{currentfill}%
\pgfsetfillopacity{0.130000}%
\pgfsetlinewidth{0.000000pt}%
\definecolor{currentstroke}{rgb}{0.000000,0.000000,0.000000}%
\pgfsetstrokecolor{currentstroke}%
\pgfsetdash{}{0pt}%
\pgfpathmoveto{\pgfqpoint{0.971422in}{0.785022in}}%
\pgfpathlineto{\pgfqpoint{0.988187in}{0.784350in}}%
\pgfpathlineto{\pgfqpoint{1.004995in}{0.784606in}}%
\pgfpathlineto{\pgfqpoint{1.013640in}{0.778886in}}%
\pgfpathlineto{\pgfqpoint{1.022246in}{0.773360in}}%
\pgfpathlineto{\pgfqpoint{1.005489in}{0.772772in}}%
\pgfpathlineto{\pgfqpoint{0.988776in}{0.773171in}}%
\pgfpathlineto{\pgfqpoint{0.980119in}{0.779009in}}%
\pgfpathlineto{\pgfqpoint{0.971422in}{0.785022in}}%
\pgfpathclose%
\pgfusepath{fill}%
\end{pgfscope}%
\begin{pgfscope}%
\pgfpathrectangle{\pgfqpoint{0.407592in}{0.222904in}}{\pgfqpoint{1.411163in}{1.411163in}}%
\pgfusepath{clip}%
\pgfsetbuttcap%
\pgfsetroundjoin%
\definecolor{currentfill}{rgb}{0.556863,0.654902,0.847059}%
\pgfsetfillcolor{currentfill}%
\pgfsetfillopacity{0.130000}%
\pgfsetlinewidth{0.000000pt}%
\definecolor{currentstroke}{rgb}{0.000000,0.000000,0.000000}%
\pgfsetstrokecolor{currentstroke}%
\pgfsetdash{}{0pt}%
\pgfpathmoveto{\pgfqpoint{1.713967in}{1.114429in}}%
\pgfpathlineto{\pgfqpoint{1.732264in}{1.127307in}}%
\pgfpathlineto{\pgfqpoint{1.739770in}{1.135793in}}%
\pgfpathlineto{\pgfqpoint{1.747264in}{1.144528in}}%
\pgfpathlineto{\pgfqpoint{1.728992in}{1.131852in}}%
\pgfpathlineto{\pgfqpoint{1.721485in}{1.123013in}}%
\pgfusepath{fill}%
\end{pgfscope}%
\begin{pgfscope}%
\pgfpathrectangle{\pgfqpoint{0.407592in}{0.222904in}}{\pgfqpoint{1.411163in}{1.411163in}}%
\pgfusepath{clip}%
\pgfsetbuttcap%
\pgfsetroundjoin%
\definecolor{currentfill}{rgb}{0.556863,0.654902,0.847059}%
\pgfsetfillcolor{currentfill}%
\pgfsetfillopacity{0.130000}%
\pgfsetlinewidth{0.000000pt}%
\definecolor{currentstroke}{rgb}{0.000000,0.000000,0.000000}%
\pgfsetstrokecolor{currentstroke}%
\pgfsetdash{}{0pt}%
\pgfpathmoveto{\pgfqpoint{1.427107in}{0.894433in}}%
\pgfpathlineto{\pgfqpoint{1.444758in}{0.906787in}}%
\pgfpathlineto{\pgfqpoint{1.462488in}{0.919286in}}%
\pgfpathlineto{\pgfqpoint{1.470307in}{0.923745in}}%
\pgfpathlineto{\pgfqpoint{1.478108in}{0.928598in}}%
\pgfpathlineto{\pgfqpoint{1.460401in}{0.916071in}}%
\pgfpathlineto{\pgfqpoint{1.442773in}{0.903684in}}%
\pgfpathlineto{\pgfqpoint{1.434950in}{0.898846in}}%
\pgfpathlineto{\pgfqpoint{1.427107in}{0.894433in}}%
\pgfpathclose%
\pgfusepath{fill}%
\end{pgfscope}%
\begin{pgfscope}%
\pgfpathrectangle{\pgfqpoint{0.407592in}{0.222904in}}{\pgfqpoint{1.411163in}{1.411163in}}%
\pgfusepath{clip}%
\pgfsetbuttcap%
\pgfsetroundjoin%
\definecolor{currentfill}{rgb}{0.556863,0.654902,0.847059}%
\pgfsetfillcolor{currentfill}%
\pgfsetfillopacity{0.130000}%
\pgfsetlinewidth{0.000000pt}%
\definecolor{currentstroke}{rgb}{0.000000,0.000000,0.000000}%
\pgfsetstrokecolor{currentstroke}%
\pgfsetdash{}{0pt}%
\pgfpathmoveto{\pgfqpoint{0.835734in}{0.835362in}}%
\pgfpathlineto{\pgfqpoint{0.852408in}{0.829074in}}%
\pgfpathlineto{\pgfqpoint{0.869109in}{0.823639in}}%
\pgfpathlineto{\pgfqpoint{0.878078in}{0.816909in}}%
\pgfpathlineto{\pgfqpoint{0.887006in}{0.810297in}}%
\pgfpathlineto{\pgfqpoint{0.870358in}{0.815804in}}%
\pgfpathlineto{\pgfqpoint{0.853737in}{0.822223in}}%
\pgfpathlineto{\pgfqpoint{0.844757in}{0.828730in}}%
\pgfpathlineto{\pgfqpoint{0.835734in}{0.835362in}}%
\pgfpathclose%
\pgfusepath{fill}%
\end{pgfscope}%
\begin{pgfscope}%
\pgfpathrectangle{\pgfqpoint{0.407592in}{0.222904in}}{\pgfqpoint{1.411163in}{1.411163in}}%
\pgfusepath{clip}%
\pgfsetbuttcap%
\pgfsetroundjoin%
\definecolor{currentfill}{rgb}{0.556863,0.654902,0.847059}%
\pgfsetfillcolor{currentfill}%
\pgfsetfillopacity{0.130000}%
\pgfsetlinewidth{0.000000pt}%
\definecolor{currentstroke}{rgb}{0.000000,0.000000,0.000000}%
\pgfsetstrokecolor{currentstroke}%
\pgfsetdash{}{0pt}%
\pgfpathmoveto{\pgfqpoint{1.208059in}{0.792526in}}%
\pgfpathlineto{\pgfqpoint{1.225230in}{0.801284in}}%
\pgfpathlineto{\pgfqpoint{1.242474in}{0.810480in}}%
\pgfpathlineto{\pgfqpoint{1.250630in}{0.809743in}}%
\pgfpathlineto{\pgfqpoint{1.258758in}{0.809424in}}%
\pgfpathlineto{\pgfqpoint{1.241545in}{0.799761in}}%
\pgfpathlineto{\pgfqpoint{1.224405in}{0.790509in}}%
\pgfpathlineto{\pgfqpoint{1.216247in}{0.791301in}}%
\pgfpathlineto{\pgfqpoint{1.208059in}{0.792526in}}%
\pgfpathclose%
\pgfusepath{fill}%
\end{pgfscope}%
\begin{pgfscope}%
\pgfpathrectangle{\pgfqpoint{0.407592in}{0.222904in}}{\pgfqpoint{1.411163in}{1.411163in}}%
\pgfusepath{clip}%
\pgfsetbuttcap%
\pgfsetroundjoin%
\definecolor{currentfill}{rgb}{0.556863,0.654902,0.847059}%
\pgfsetfillcolor{currentfill}%
\pgfsetfillopacity{0.130000}%
\pgfsetlinewidth{0.000000pt}%
\definecolor{currentstroke}{rgb}{0.000000,0.000000,0.000000}%
\pgfsetstrokecolor{currentstroke}%
\pgfsetdash{}{0pt}%
\pgfpathmoveto{\pgfqpoint{1.157384in}{0.781268in}}%
\pgfpathlineto{\pgfqpoint{1.174452in}{0.788442in}}%
\pgfpathlineto{\pgfqpoint{1.191588in}{0.796177in}}%
\pgfpathlineto{\pgfqpoint{1.199839in}{0.794159in}}%
\pgfpathlineto{\pgfqpoint{1.208059in}{0.792526in}}%
\pgfpathlineto{\pgfqpoint{1.190958in}{0.784255in}}%
\pgfpathlineto{\pgfqpoint{1.173926in}{0.776526in}}%
\pgfpathlineto{\pgfqpoint{1.165672in}{0.778704in}}%
\pgfpathlineto{\pgfqpoint{1.157384in}{0.781268in}}%
\pgfpathclose%
\pgfusepath{fill}%
\end{pgfscope}%
\begin{pgfscope}%
\pgfpathrectangle{\pgfqpoint{0.407592in}{0.222904in}}{\pgfqpoint{1.411163in}{1.411163in}}%
\pgfusepath{clip}%
\pgfsetbuttcap%
\pgfsetroundjoin%
\definecolor{currentfill}{rgb}{0.556863,0.654902,0.847059}%
\pgfsetfillcolor{currentfill}%
\pgfsetfillopacity{0.130000}%
\pgfsetlinewidth{0.000000pt}%
\definecolor{currentstroke}{rgb}{0.000000,0.000000,0.000000}%
\pgfsetstrokecolor{currentstroke}%
\pgfsetdash{}{0pt}%
\pgfpathmoveto{\pgfqpoint{1.595927in}{1.016276in}}%
\pgfpathlineto{\pgfqpoint{1.613968in}{1.029183in}}%
\pgfpathlineto{\pgfqpoint{1.632092in}{1.042204in}}%
\pgfpathlineto{\pgfqpoint{1.639713in}{1.049276in}}%
\pgfpathlineto{\pgfqpoint{1.647320in}{1.056650in}}%
\pgfpathlineto{\pgfqpoint{1.629221in}{1.043773in}}%
\pgfpathlineto{\pgfqpoint{1.611203in}{1.031016in}}%
\pgfpathlineto{\pgfqpoint{1.603572in}{1.023487in}}%
\pgfpathlineto{\pgfqpoint{1.595927in}{1.016276in}}%
\pgfpathclose%
\pgfusepath{fill}%
\end{pgfscope}%
\begin{pgfscope}%
\pgfpathrectangle{\pgfqpoint{0.407592in}{0.222904in}}{\pgfqpoint{1.411163in}{1.411163in}}%
\pgfusepath{clip}%
\pgfsetbuttcap%
\pgfsetroundjoin%
\definecolor{currentfill}{rgb}{0.556863,0.654902,0.847059}%
\pgfsetfillcolor{currentfill}%
\pgfsetfillopacity{0.130000}%
\pgfsetlinewidth{0.000000pt}%
\definecolor{currentstroke}{rgb}{0.000000,0.000000,0.000000}%
\pgfsetstrokecolor{currentstroke}%
\pgfsetdash{}{0pt}%
\pgfpathmoveto{\pgfqpoint{0.750874in}{0.881356in}}%
\pgfpathlineto{\pgfqpoint{0.767515in}{0.872148in}}%
\pgfpathlineto{\pgfqpoint{0.784175in}{0.863664in}}%
\pgfpathlineto{\pgfqpoint{0.793337in}{0.856988in}}%
\pgfpathlineto{\pgfqpoint{0.802455in}{0.850431in}}%
\pgfpathlineto{\pgfqpoint{0.785847in}{0.859163in}}%
\pgfpathlineto{\pgfqpoint{0.769257in}{0.868655in}}%
\pgfpathlineto{\pgfqpoint{0.760087in}{0.874939in}}%
\pgfpathlineto{\pgfqpoint{0.750874in}{0.881356in}}%
\pgfpathclose%
\pgfusepath{fill}%
\end{pgfscope}%
\begin{pgfscope}%
\pgfpathrectangle{\pgfqpoint{0.407592in}{0.222904in}}{\pgfqpoint{1.411163in}{1.411163in}}%
\pgfusepath{clip}%
\pgfsetbuttcap%
\pgfsetroundjoin%
\definecolor{currentfill}{rgb}{0.556863,0.654902,0.847059}%
\pgfsetfillcolor{currentfill}%
\pgfsetfillopacity{0.130000}%
\pgfsetlinewidth{0.000000pt}%
\definecolor{currentstroke}{rgb}{0.000000,0.000000,0.000000}%
\pgfsetstrokecolor{currentstroke}%
\pgfsetdash{}{0pt}%
\pgfpathmoveto{\pgfqpoint{1.258758in}{0.809424in}}%
\pgfpathlineto{\pgfqpoint{1.276044in}{0.819455in}}%
\pgfpathlineto{\pgfqpoint{1.293406in}{0.829819in}}%
\pgfpathlineto{\pgfqpoint{1.301477in}{0.830332in}}%
\pgfpathlineto{\pgfqpoint{1.309521in}{0.831275in}}%
\pgfpathlineto{\pgfqpoint{1.292188in}{0.820541in}}%
\pgfpathlineto{\pgfqpoint{1.274929in}{0.810113in}}%
\pgfpathlineto{\pgfqpoint{1.266857in}{0.809540in}}%
\pgfpathlineto{\pgfqpoint{1.258758in}{0.809424in}}%
\pgfpathclose%
\pgfusepath{fill}%
\end{pgfscope}%
\begin{pgfscope}%
\pgfpathrectangle{\pgfqpoint{0.407592in}{0.222904in}}{\pgfqpoint{1.411163in}{1.411163in}}%
\pgfusepath{clip}%
\pgfsetbuttcap%
\pgfsetroundjoin%
\definecolor{currentfill}{rgb}{0.556863,0.654902,0.847059}%
\pgfsetfillcolor{currentfill}%
\pgfsetfillopacity{0.130000}%
\pgfsetlinewidth{0.000000pt}%
\definecolor{currentstroke}{rgb}{0.000000,0.000000,0.000000}%
\pgfsetstrokecolor{currentstroke}%
\pgfsetdash{}{0pt}%
\pgfpathmoveto{\pgfqpoint{1.106684in}{0.776140in}}%
\pgfpathlineto{\pgfqpoint{1.123663in}{0.781453in}}%
\pgfpathlineto{\pgfqpoint{1.140706in}{0.787452in}}%
\pgfpathlineto{\pgfqpoint{1.149062in}{0.784192in}}%
\pgfpathlineto{\pgfqpoint{1.157384in}{0.781268in}}%
\pgfpathlineto{\pgfqpoint{1.140382in}{0.774713in}}%
\pgfpathlineto{\pgfqpoint{1.123443in}{0.768841in}}%
\pgfpathlineto{\pgfqpoint{1.115081in}{0.772329in}}%
\pgfpathlineto{\pgfqpoint{1.106684in}{0.776140in}}%
\pgfpathclose%
\pgfusepath{fill}%
\end{pgfscope}%
\begin{pgfscope}%
\pgfpathrectangle{\pgfqpoint{0.407592in}{0.222904in}}{\pgfqpoint{1.411163in}{1.411163in}}%
\pgfusepath{clip}%
\pgfsetbuttcap%
\pgfsetroundjoin%
\definecolor{currentfill}{rgb}{0.556863,0.654902,0.847059}%
\pgfsetfillcolor{currentfill}%
\pgfsetfillopacity{0.130000}%
\pgfsetlinewidth{0.000000pt}%
\definecolor{currentstroke}{rgb}{0.000000,0.000000,0.000000}%
\pgfsetstrokecolor{currentstroke}%
\pgfsetdash{}{0pt}%
\pgfpathmoveto{\pgfqpoint{1.478108in}{0.928598in}}%
\pgfpathlineto{\pgfqpoint{1.495895in}{0.941256in}}%
\pgfpathlineto{\pgfqpoint{1.513763in}{0.954039in}}%
\pgfpathlineto{\pgfqpoint{1.521524in}{0.959267in}}%
\pgfpathlineto{\pgfqpoint{1.529268in}{0.964860in}}%
\pgfpathlineto{\pgfqpoint{1.511423in}{0.952108in}}%
\pgfpathlineto{\pgfqpoint{1.493658in}{0.939481in}}%
\pgfpathlineto{\pgfqpoint{1.485892in}{0.933844in}}%
\pgfpathlineto{\pgfqpoint{1.478108in}{0.928598in}}%
\pgfpathclose%
\pgfusepath{fill}%
\end{pgfscope}%
\begin{pgfscope}%
\pgfpathrectangle{\pgfqpoint{0.407592in}{0.222904in}}{\pgfqpoint{1.411163in}{1.411163in}}%
\pgfusepath{clip}%
\pgfsetbuttcap%
\pgfsetroundjoin%
\definecolor{currentfill}{rgb}{0.556863,0.654902,0.847059}%
\pgfsetfillcolor{currentfill}%
\pgfsetfillopacity{0.130000}%
\pgfsetlinewidth{0.000000pt}%
\definecolor{currentstroke}{rgb}{0.000000,0.000000,0.000000}%
\pgfsetstrokecolor{currentstroke}%
\pgfsetdash{}{0pt}%
\pgfpathmoveto{\pgfqpoint{0.920388in}{0.802059in}}%
\pgfpathlineto{\pgfqpoint{0.937129in}{0.799331in}}%
\pgfpathlineto{\pgfqpoint{0.953907in}{0.797518in}}%
\pgfpathlineto{\pgfqpoint{0.962685in}{0.791196in}}%
\pgfpathlineto{\pgfqpoint{0.971422in}{0.785022in}}%
\pgfpathlineto{\pgfqpoint{0.954698in}{0.786652in}}%
\pgfpathlineto{\pgfqpoint{0.938010in}{0.789262in}}%
\pgfpathlineto{\pgfqpoint{0.929220in}{0.795592in}}%
\pgfpathlineto{\pgfqpoint{0.920388in}{0.802059in}}%
\pgfpathclose%
\pgfusepath{fill}%
\end{pgfscope}%
\begin{pgfscope}%
\pgfpathrectangle{\pgfqpoint{0.407592in}{0.222904in}}{\pgfqpoint{1.411163in}{1.411163in}}%
\pgfusepath{clip}%
\pgfsetbuttcap%
\pgfsetroundjoin%
\definecolor{currentfill}{rgb}{0.556863,0.654902,0.847059}%
\pgfsetfillcolor{currentfill}%
\pgfsetfillopacity{0.130000}%
\pgfsetlinewidth{0.000000pt}%
\definecolor{currentstroke}{rgb}{0.000000,0.000000,0.000000}%
\pgfsetstrokecolor{currentstroke}%
\pgfsetdash{}{0pt}%
\pgfpathmoveto{\pgfqpoint{1.309521in}{0.831275in}}%
\pgfpathlineto{\pgfqpoint{1.326932in}{0.842284in}}%
\pgfpathlineto{\pgfqpoint{1.344419in}{0.853544in}}%
\pgfpathlineto{\pgfqpoint{1.352413in}{0.855223in}}%
\pgfpathlineto{\pgfqpoint{1.360383in}{0.857330in}}%
\pgfpathlineto{\pgfqpoint{1.342921in}{0.845805in}}%
\pgfpathlineto{\pgfqpoint{1.325536in}{0.834508in}}%
\pgfpathlineto{\pgfqpoint{1.317541in}{0.832663in}}%
\pgfpathlineto{\pgfqpoint{1.309521in}{0.831275in}}%
\pgfpathclose%
\pgfusepath{fill}%
\end{pgfscope}%
\begin{pgfscope}%
\pgfpathrectangle{\pgfqpoint{0.407592in}{0.222904in}}{\pgfqpoint{1.411163in}{1.411163in}}%
\pgfusepath{clip}%
\pgfsetbuttcap%
\pgfsetroundjoin%
\definecolor{currentfill}{rgb}{0.556863,0.654902,0.847059}%
\pgfsetfillcolor{currentfill}%
\pgfsetfillopacity{0.130000}%
\pgfsetlinewidth{0.000000pt}%
\definecolor{currentstroke}{rgb}{0.000000,0.000000,0.000000}%
\pgfsetstrokecolor{currentstroke}%
\pgfsetdash{}{0pt}%
\pgfpathmoveto{\pgfqpoint{1.055905in}{0.777314in}}%
\pgfpathlineto{\pgfqpoint{1.072813in}{0.780580in}}%
\pgfpathlineto{\pgfqpoint{1.089778in}{0.784632in}}%
\pgfpathlineto{\pgfqpoint{1.098249in}{0.780248in}}%
\pgfpathlineto{\pgfqpoint{1.106684in}{0.776140in}}%
\pgfpathlineto{\pgfqpoint{1.089764in}{0.771570in}}%
\pgfpathlineto{\pgfqpoint{1.072902in}{0.767809in}}%
\pgfpathlineto{\pgfqpoint{1.064423in}{0.772434in}}%
\pgfpathlineto{\pgfqpoint{1.055905in}{0.777314in}}%
\pgfpathclose%
\pgfusepath{fill}%
\end{pgfscope}%
\begin{pgfscope}%
\pgfpathrectangle{\pgfqpoint{0.407592in}{0.222904in}}{\pgfqpoint{1.411163in}{1.411163in}}%
\pgfusepath{clip}%
\pgfsetbuttcap%
\pgfsetroundjoin%
\definecolor{currentfill}{rgb}{0.556863,0.654902,0.847059}%
\pgfsetfillcolor{currentfill}%
\pgfsetfillopacity{0.130000}%
\pgfsetlinewidth{0.000000pt}%
\definecolor{currentstroke}{rgb}{0.000000,0.000000,0.000000}%
\pgfsetstrokecolor{currentstroke}%
\pgfsetdash{}{0pt}%
\pgfpathmoveto{\pgfqpoint{1.647320in}{1.056650in}}%
\pgfpathlineto{\pgfqpoint{1.665502in}{1.069642in}}%
\pgfpathlineto{\pgfqpoint{1.683766in}{1.082745in}}%
\pgfpathlineto{\pgfqpoint{1.691336in}{1.090258in}}%
\pgfpathlineto{\pgfqpoint{1.698892in}{1.098045in}}%
\pgfpathlineto{\pgfqpoint{1.680652in}{1.085103in}}%
\pgfpathlineto{\pgfqpoint{1.662495in}{1.072279in}}%
\pgfpathlineto{\pgfqpoint{1.654914in}{1.064319in}}%
\pgfpathlineto{\pgfqpoint{1.647320in}{1.056650in}}%
\pgfpathclose%
\pgfusepath{fill}%
\end{pgfscope}%
\begin{pgfscope}%
\pgfpathrectangle{\pgfqpoint{0.407592in}{0.222904in}}{\pgfqpoint{1.411163in}{1.411163in}}%
\pgfusepath{clip}%
\pgfsetbuttcap%
\pgfsetroundjoin%
\definecolor{currentfill}{rgb}{0.556863,0.654902,0.847059}%
\pgfsetfillcolor{currentfill}%
\pgfsetfillopacity{0.130000}%
\pgfsetlinewidth{0.000000pt}%
\definecolor{currentstroke}{rgb}{0.000000,0.000000,0.000000}%
\pgfsetstrokecolor{currentstroke}%
\pgfsetdash{}{0pt}%
\pgfpathmoveto{\pgfqpoint{1.360383in}{0.857330in}}%
\pgfpathlineto{\pgfqpoint{1.377923in}{0.869064in}}%
\pgfpathlineto{\pgfqpoint{1.395541in}{0.880990in}}%
\pgfpathlineto{\pgfqpoint{1.403464in}{0.883727in}}%
\pgfpathlineto{\pgfqpoint{1.411366in}{0.886876in}}%
\pgfpathlineto{\pgfqpoint{1.393771in}{0.874781in}}%
\pgfpathlineto{\pgfqpoint{1.376255in}{0.862863in}}%
\pgfpathlineto{\pgfqpoint{1.368330in}{0.859874in}}%
\pgfpathlineto{\pgfqpoint{1.360383in}{0.857330in}}%
\pgfpathclose%
\pgfusepath{fill}%
\end{pgfscope}%
\begin{pgfscope}%
\pgfpathrectangle{\pgfqpoint{0.407592in}{0.222904in}}{\pgfqpoint{1.411163in}{1.411163in}}%
\pgfusepath{clip}%
\pgfsetbuttcap%
\pgfsetroundjoin%
\definecolor{currentfill}{rgb}{0.556863,0.654902,0.847059}%
\pgfsetfillcolor{currentfill}%
\pgfsetfillopacity{0.130000}%
\pgfsetlinewidth{0.000000pt}%
\definecolor{currentstroke}{rgb}{0.000000,0.000000,0.000000}%
\pgfsetstrokecolor{currentstroke}%
\pgfsetdash{}{0pt}%
\pgfpathmoveto{\pgfqpoint{1.529268in}{0.964860in}}%
\pgfpathlineto{\pgfqpoint{1.547193in}{0.977731in}}%
\pgfpathlineto{\pgfqpoint{1.565200in}{0.990715in}}%
\pgfpathlineto{\pgfqpoint{1.572905in}{0.996605in}}%
\pgfpathlineto{\pgfqpoint{1.580594in}{1.002830in}}%
\pgfpathlineto{\pgfqpoint{1.562610in}{0.989921in}}%
\pgfpathlineto{\pgfqpoint{1.544708in}{0.977130in}}%
\pgfpathlineto{\pgfqpoint{1.536995in}{0.970815in}}%
\pgfpathlineto{\pgfqpoint{1.529268in}{0.964860in}}%
\pgfpathclose%
\pgfusepath{fill}%
\end{pgfscope}%
\begin{pgfscope}%
\pgfpathrectangle{\pgfqpoint{0.407592in}{0.222904in}}{\pgfqpoint{1.411163in}{1.411163in}}%
\pgfusepath{clip}%
\pgfsetbuttcap%
\pgfsetroundjoin%
\definecolor{currentfill}{rgb}{0.556863,0.654902,0.847059}%
\pgfsetfillcolor{currentfill}%
\pgfsetfillopacity{0.130000}%
\pgfsetlinewidth{0.000000pt}%
\definecolor{currentstroke}{rgb}{0.000000,0.000000,0.000000}%
\pgfsetstrokecolor{currentstroke}%
\pgfsetdash{}{0pt}%
\pgfpathmoveto{\pgfqpoint{1.698892in}{1.098045in}}%
\pgfpathlineto{\pgfqpoint{1.717215in}{1.111103in}}%
\pgfpathlineto{\pgfqpoint{1.724746in}{1.119075in}}%
\pgfpathlineto{\pgfqpoint{1.732264in}{1.127307in}}%
\pgfpathlineto{\pgfqpoint{1.713967in}{1.114429in}}%
\pgfpathlineto{\pgfqpoint{1.706436in}{1.106104in}}%
\pgfusepath{fill}%
\end{pgfscope}%
\begin{pgfscope}%
\pgfpathrectangle{\pgfqpoint{0.407592in}{0.222904in}}{\pgfqpoint{1.411163in}{1.411163in}}%
\pgfusepath{clip}%
\pgfsetbuttcap%
\pgfsetroundjoin%
\definecolor{currentfill}{rgb}{0.556863,0.654902,0.847059}%
\pgfsetfillcolor{currentfill}%
\pgfsetfillopacity{0.130000}%
\pgfsetlinewidth{0.000000pt}%
\definecolor{currentstroke}{rgb}{0.000000,0.000000,0.000000}%
\pgfsetstrokecolor{currentstroke}%
\pgfsetdash{}{0pt}%
\pgfpathmoveto{\pgfqpoint{1.004995in}{0.784606in}}%
\pgfpathlineto{\pgfqpoint{1.021850in}{0.785755in}}%
\pgfpathlineto{\pgfqpoint{1.038754in}{0.787757in}}%
\pgfpathlineto{\pgfqpoint{1.047349in}{0.782428in}}%
\pgfpathlineto{\pgfqpoint{1.055905in}{0.777314in}}%
\pgfpathlineto{\pgfqpoint{1.039050in}{0.774890in}}%
\pgfpathlineto{\pgfqpoint{1.022246in}{0.773360in}}%
\pgfpathlineto{\pgfqpoint{1.013640in}{0.778886in}}%
\pgfpathlineto{\pgfqpoint{1.004995in}{0.784606in}}%
\pgfpathclose%
\pgfusepath{fill}%
\end{pgfscope}%
\begin{pgfscope}%
\pgfpathrectangle{\pgfqpoint{0.407592in}{0.222904in}}{\pgfqpoint{1.411163in}{1.411163in}}%
\pgfusepath{clip}%
\pgfsetbuttcap%
\pgfsetroundjoin%
\definecolor{currentfill}{rgb}{0.556863,0.654902,0.847059}%
\pgfsetfillcolor{currentfill}%
\pgfsetfillopacity{0.130000}%
\pgfsetlinewidth{0.000000pt}%
\definecolor{currentstroke}{rgb}{0.000000,0.000000,0.000000}%
\pgfsetstrokecolor{currentstroke}%
\pgfsetdash{}{0pt}%
\pgfpathmoveto{\pgfqpoint{0.784175in}{0.863664in}}%
\pgfpathlineto{\pgfqpoint{0.800856in}{0.855929in}}%
\pgfpathlineto{\pgfqpoint{0.817559in}{0.848965in}}%
\pgfpathlineto{\pgfqpoint{0.826668in}{0.842110in}}%
\pgfpathlineto{\pgfqpoint{0.835734in}{0.835362in}}%
\pgfpathlineto{\pgfqpoint{0.819083in}{0.842488in}}%
\pgfpathlineto{\pgfqpoint{0.802455in}{0.850431in}}%
\pgfpathlineto{\pgfqpoint{0.793337in}{0.856988in}}%
\pgfpathlineto{\pgfqpoint{0.784175in}{0.863664in}}%
\pgfpathclose%
\pgfusepath{fill}%
\end{pgfscope}%
\begin{pgfscope}%
\pgfpathrectangle{\pgfqpoint{0.407592in}{0.222904in}}{\pgfqpoint{1.411163in}{1.411163in}}%
\pgfusepath{clip}%
\pgfsetbuttcap%
\pgfsetroundjoin%
\definecolor{currentfill}{rgb}{0.556863,0.654902,0.847059}%
\pgfsetfillcolor{currentfill}%
\pgfsetfillopacity{0.130000}%
\pgfsetlinewidth{0.000000pt}%
\definecolor{currentstroke}{rgb}{0.000000,0.000000,0.000000}%
\pgfsetstrokecolor{currentstroke}%
\pgfsetdash{}{0pt}%
\pgfpathmoveto{\pgfqpoint{0.869109in}{0.823639in}}%
\pgfpathlineto{\pgfqpoint{0.885838in}{0.819069in}}%
\pgfpathlineto{\pgfqpoint{0.902599in}{0.815369in}}%
\pgfpathlineto{\pgfqpoint{0.911515in}{0.808655in}}%
\pgfpathlineto{\pgfqpoint{0.920388in}{0.802059in}}%
\pgfpathlineto{\pgfqpoint{0.903681in}{0.805714in}}%
\pgfpathlineto{\pgfqpoint{0.887006in}{0.810297in}}%
\pgfpathlineto{\pgfqpoint{0.878078in}{0.816909in}}%
\pgfpathlineto{\pgfqpoint{0.869109in}{0.823639in}}%
\pgfpathclose%
\pgfusepath{fill}%
\end{pgfscope}%
\begin{pgfscope}%
\pgfpathrectangle{\pgfqpoint{0.407592in}{0.222904in}}{\pgfqpoint{1.411163in}{1.411163in}}%
\pgfusepath{clip}%
\pgfsetbuttcap%
\pgfsetroundjoin%
\definecolor{currentfill}{rgb}{0.556863,0.654902,0.847059}%
\pgfsetfillcolor{currentfill}%
\pgfsetfillopacity{0.130000}%
\pgfsetlinewidth{0.000000pt}%
\definecolor{currentstroke}{rgb}{0.000000,0.000000,0.000000}%
\pgfsetstrokecolor{currentstroke}%
\pgfsetdash{}{0pt}%
\pgfpathmoveto{\pgfqpoint{1.411366in}{0.886876in}}%
\pgfpathlineto{\pgfqpoint{1.429039in}{0.899135in}}%
\pgfpathlineto{\pgfqpoint{1.446793in}{0.911546in}}%
\pgfpathlineto{\pgfqpoint{1.454650in}{0.915220in}}%
\pgfpathlineto{\pgfqpoint{1.462488in}{0.919286in}}%
\pgfpathlineto{\pgfqpoint{1.444758in}{0.906787in}}%
\pgfpathlineto{\pgfqpoint{1.427107in}{0.894433in}}%
\pgfpathlineto{\pgfqpoint{1.419246in}{0.890444in}}%
\pgfpathlineto{\pgfqpoint{1.411366in}{0.886876in}}%
\pgfpathclose%
\pgfusepath{fill}%
\end{pgfscope}%
\begin{pgfscope}%
\pgfpathrectangle{\pgfqpoint{0.407592in}{0.222904in}}{\pgfqpoint{1.411163in}{1.411163in}}%
\pgfusepath{clip}%
\pgfsetbuttcap%
\pgfsetroundjoin%
\definecolor{currentfill}{rgb}{0.556863,0.654902,0.847059}%
\pgfsetfillcolor{currentfill}%
\pgfsetfillopacity{0.130000}%
\pgfsetlinewidth{0.000000pt}%
\definecolor{currentstroke}{rgb}{0.000000,0.000000,0.000000}%
\pgfsetstrokecolor{currentstroke}%
\pgfsetdash{}{0pt}%
\pgfpathmoveto{\pgfqpoint{1.580594in}{1.002830in}}%
\pgfpathlineto{\pgfqpoint{1.598659in}{1.015851in}}%
\pgfpathlineto{\pgfqpoint{1.616807in}{1.028979in}}%
\pgfpathlineto{\pgfqpoint{1.624456in}{1.035437in}}%
\pgfpathlineto{\pgfqpoint{1.632092in}{1.042204in}}%
\pgfpathlineto{\pgfqpoint{1.613968in}{1.029183in}}%
\pgfpathlineto{\pgfqpoint{1.595927in}{1.016276in}}%
\pgfpathlineto{\pgfqpoint{1.588267in}{1.009389in}}%
\pgfpathlineto{\pgfqpoint{1.580594in}{1.002830in}}%
\pgfpathclose%
\pgfusepath{fill}%
\end{pgfscope}%
\begin{pgfscope}%
\pgfpathrectangle{\pgfqpoint{0.407592in}{0.222904in}}{\pgfqpoint{1.411163in}{1.411163in}}%
\pgfusepath{clip}%
\pgfsetbuttcap%
\pgfsetroundjoin%
\definecolor{currentfill}{rgb}{0.556863,0.654902,0.847059}%
\pgfsetfillcolor{currentfill}%
\pgfsetfillopacity{0.130000}%
\pgfsetlinewidth{0.000000pt}%
\definecolor{currentstroke}{rgb}{0.000000,0.000000,0.000000}%
\pgfsetstrokecolor{currentstroke}%
\pgfsetdash{}{0pt}%
\pgfpathmoveto{\pgfqpoint{0.953907in}{0.797518in}}%
\pgfpathlineto{\pgfqpoint{0.970724in}{0.796604in}}%
\pgfpathlineto{\pgfqpoint{0.987584in}{0.796565in}}%
\pgfpathlineto{\pgfqpoint{0.996310in}{0.790503in}}%
\pgfpathlineto{\pgfqpoint{1.004995in}{0.784606in}}%
\pgfpathlineto{\pgfqpoint{0.988187in}{0.784350in}}%
\pgfpathlineto{\pgfqpoint{0.971422in}{0.785022in}}%
\pgfpathlineto{\pgfqpoint{0.962685in}{0.791196in}}%
\pgfpathlineto{\pgfqpoint{0.953907in}{0.797518in}}%
\pgfpathclose%
\pgfusepath{fill}%
\end{pgfscope}%
\begin{pgfscope}%
\pgfpathrectangle{\pgfqpoint{0.407592in}{0.222904in}}{\pgfqpoint{1.411163in}{1.411163in}}%
\pgfusepath{clip}%
\pgfsetbuttcap%
\pgfsetroundjoin%
\definecolor{currentfill}{rgb}{0.556863,0.654902,0.847059}%
\pgfsetfillcolor{currentfill}%
\pgfsetfillopacity{0.130000}%
\pgfsetlinewidth{0.000000pt}%
\definecolor{currentstroke}{rgb}{0.000000,0.000000,0.000000}%
\pgfsetstrokecolor{currentstroke}%
\pgfsetdash{}{0pt}%
\pgfpathmoveto{\pgfqpoint{1.462488in}{0.919286in}}%
\pgfpathlineto{\pgfqpoint{1.480299in}{0.931920in}}%
\pgfpathlineto{\pgfqpoint{1.498190in}{0.944682in}}%
\pgfpathlineto{\pgfqpoint{1.505986in}{0.949177in}}%
\pgfpathlineto{\pgfqpoint{1.513763in}{0.954039in}}%
\pgfpathlineto{\pgfqpoint{1.495895in}{0.941256in}}%
\pgfpathlineto{\pgfqpoint{1.478108in}{0.928598in}}%
\pgfpathlineto{\pgfqpoint{1.470307in}{0.923745in}}%
\pgfpathlineto{\pgfqpoint{1.462488in}{0.919286in}}%
\pgfpathclose%
\pgfusepath{fill}%
\end{pgfscope}%
\begin{pgfscope}%
\pgfpathrectangle{\pgfqpoint{0.407592in}{0.222904in}}{\pgfqpoint{1.411163in}{1.411163in}}%
\pgfusepath{clip}%
\pgfsetbuttcap%
\pgfsetroundjoin%
\definecolor{currentfill}{rgb}{0.556863,0.654902,0.847059}%
\pgfsetfillcolor{currentfill}%
\pgfsetfillopacity{0.130000}%
\pgfsetlinewidth{0.000000pt}%
\definecolor{currentstroke}{rgb}{0.000000,0.000000,0.000000}%
\pgfsetstrokecolor{currentstroke}%
\pgfsetdash{}{0pt}%
\pgfpathmoveto{\pgfqpoint{1.191588in}{0.796177in}}%
\pgfpathlineto{\pgfqpoint{1.208794in}{0.804421in}}%
\pgfpathlineto{\pgfqpoint{1.226070in}{0.813128in}}%
\pgfpathlineto{\pgfqpoint{1.234287in}{0.811615in}}%
\pgfpathlineto{\pgfqpoint{1.242474in}{0.810480in}}%
\pgfpathlineto{\pgfqpoint{1.225230in}{0.801284in}}%
\pgfpathlineto{\pgfqpoint{1.208059in}{0.792526in}}%
\pgfpathlineto{\pgfqpoint{1.199839in}{0.794159in}}%
\pgfpathlineto{\pgfqpoint{1.191588in}{0.796177in}}%
\pgfpathclose%
\pgfusepath{fill}%
\end{pgfscope}%
\begin{pgfscope}%
\pgfpathrectangle{\pgfqpoint{0.407592in}{0.222904in}}{\pgfqpoint{1.411163in}{1.411163in}}%
\pgfusepath{clip}%
\pgfsetbuttcap%
\pgfsetroundjoin%
\definecolor{currentfill}{rgb}{0.556863,0.654902,0.847059}%
\pgfsetfillcolor{currentfill}%
\pgfsetfillopacity{0.130000}%
\pgfsetlinewidth{0.000000pt}%
\definecolor{currentstroke}{rgb}{0.000000,0.000000,0.000000}%
\pgfsetstrokecolor{currentstroke}%
\pgfsetdash{}{0pt}%
\pgfpathmoveto{\pgfqpoint{1.140706in}{0.787452in}}%
\pgfpathlineto{\pgfqpoint{1.157813in}{0.794077in}}%
\pgfpathlineto{\pgfqpoint{1.174986in}{0.801278in}}%
\pgfpathlineto{\pgfqpoint{1.183304in}{0.798557in}}%
\pgfpathlineto{\pgfqpoint{1.191588in}{0.796177in}}%
\pgfpathlineto{\pgfqpoint{1.174452in}{0.788442in}}%
\pgfpathlineto{\pgfqpoint{1.157384in}{0.781268in}}%
\pgfpathlineto{\pgfqpoint{1.149062in}{0.784192in}}%
\pgfpathlineto{\pgfqpoint{1.140706in}{0.787452in}}%
\pgfpathclose%
\pgfusepath{fill}%
\end{pgfscope}%
\begin{pgfscope}%
\pgfpathrectangle{\pgfqpoint{0.407592in}{0.222904in}}{\pgfqpoint{1.411163in}{1.411163in}}%
\pgfusepath{clip}%
\pgfsetbuttcap%
\pgfsetroundjoin%
\definecolor{currentfill}{rgb}{0.556863,0.654902,0.847059}%
\pgfsetfillcolor{currentfill}%
\pgfsetfillopacity{0.130000}%
\pgfsetlinewidth{0.000000pt}%
\definecolor{currentstroke}{rgb}{0.000000,0.000000,0.000000}%
\pgfsetstrokecolor{currentstroke}%
\pgfsetdash{}{0pt}%
\pgfpathmoveto{\pgfqpoint{1.242474in}{0.810480in}}%
\pgfpathlineto{\pgfqpoint{1.259790in}{0.820073in}}%
\pgfpathlineto{\pgfqpoint{1.277182in}{0.830025in}}%
\pgfpathlineto{\pgfqpoint{1.285308in}{0.829722in}}%
\pgfpathlineto{\pgfqpoint{1.293406in}{0.829819in}}%
\pgfpathlineto{\pgfqpoint{1.276044in}{0.819455in}}%
\pgfpathlineto{\pgfqpoint{1.258758in}{0.809424in}}%
\pgfpathlineto{\pgfqpoint{1.250630in}{0.809743in}}%
\pgfpathlineto{\pgfqpoint{1.242474in}{0.810480in}}%
\pgfpathclose%
\pgfusepath{fill}%
\end{pgfscope}%
\begin{pgfscope}%
\pgfpathrectangle{\pgfqpoint{0.407592in}{0.222904in}}{\pgfqpoint{1.411163in}{1.411163in}}%
\pgfusepath{clip}%
\pgfsetbuttcap%
\pgfsetroundjoin%
\definecolor{currentfill}{rgb}{0.556863,0.654902,0.847059}%
\pgfsetfillcolor{currentfill}%
\pgfsetfillopacity{0.130000}%
\pgfsetlinewidth{0.000000pt}%
\definecolor{currentstroke}{rgb}{0.000000,0.000000,0.000000}%
\pgfsetstrokecolor{currentstroke}%
\pgfsetdash{}{0pt}%
\pgfpathmoveto{\pgfqpoint{1.632092in}{1.042204in}}%
\pgfpathlineto{\pgfqpoint{1.650298in}{1.055333in}}%
\pgfpathlineto{\pgfqpoint{1.668587in}{1.068568in}}%
\pgfpathlineto{\pgfqpoint{1.676183in}{1.075514in}}%
\pgfpathlineto{\pgfqpoint{1.683766in}{1.082745in}}%
\pgfpathlineto{\pgfqpoint{1.665502in}{1.069642in}}%
\pgfpathlineto{\pgfqpoint{1.647320in}{1.056650in}}%
\pgfpathlineto{\pgfqpoint{1.639713in}{1.049276in}}%
\pgfpathlineto{\pgfqpoint{1.632092in}{1.042204in}}%
\pgfpathclose%
\pgfusepath{fill}%
\end{pgfscope}%
\begin{pgfscope}%
\pgfpathrectangle{\pgfqpoint{0.407592in}{0.222904in}}{\pgfqpoint{1.411163in}{1.411163in}}%
\pgfusepath{clip}%
\pgfsetbuttcap%
\pgfsetroundjoin%
\definecolor{currentfill}{rgb}{0.556863,0.654902,0.847059}%
\pgfsetfillcolor{currentfill}%
\pgfsetfillopacity{0.130000}%
\pgfsetlinewidth{0.000000pt}%
\definecolor{currentstroke}{rgb}{0.000000,0.000000,0.000000}%
\pgfsetstrokecolor{currentstroke}%
\pgfsetdash{}{0pt}%
\pgfpathmoveto{\pgfqpoint{0.817559in}{0.848965in}}%
\pgfpathlineto{\pgfqpoint{0.834287in}{0.842792in}}%
\pgfpathlineto{\pgfqpoint{0.851042in}{0.837422in}}%
\pgfpathlineto{\pgfqpoint{0.860097in}{0.830479in}}%
\pgfpathlineto{\pgfqpoint{0.869109in}{0.823639in}}%
\pgfpathlineto{\pgfqpoint{0.852408in}{0.829074in}}%
\pgfpathlineto{\pgfqpoint{0.835734in}{0.835362in}}%
\pgfpathlineto{\pgfqpoint{0.826668in}{0.842110in}}%
\pgfpathlineto{\pgfqpoint{0.817559in}{0.848965in}}%
\pgfpathclose%
\pgfusepath{fill}%
\end{pgfscope}%
\begin{pgfscope}%
\pgfpathrectangle{\pgfqpoint{0.407592in}{0.222904in}}{\pgfqpoint{1.411163in}{1.411163in}}%
\pgfusepath{clip}%
\pgfsetbuttcap%
\pgfsetroundjoin%
\definecolor{currentfill}{rgb}{0.556863,0.654902,0.847059}%
\pgfsetfillcolor{currentfill}%
\pgfsetfillopacity{0.130000}%
\pgfsetlinewidth{0.000000pt}%
\definecolor{currentstroke}{rgb}{0.000000,0.000000,0.000000}%
\pgfsetstrokecolor{currentstroke}%
\pgfsetdash{}{0pt}%
\pgfpathmoveto{\pgfqpoint{1.293406in}{0.829819in}}%
\pgfpathlineto{\pgfqpoint{1.310843in}{0.840485in}}%
\pgfpathlineto{\pgfqpoint{1.328358in}{0.851425in}}%
\pgfpathlineto{\pgfqpoint{1.336401in}{0.852281in}}%
\pgfpathlineto{\pgfqpoint{1.344419in}{0.853544in}}%
\pgfpathlineto{\pgfqpoint{1.326932in}{0.842284in}}%
\pgfpathlineto{\pgfqpoint{1.309521in}{0.831275in}}%
\pgfpathlineto{\pgfqpoint{1.301477in}{0.830332in}}%
\pgfpathlineto{\pgfqpoint{1.293406in}{0.829819in}}%
\pgfpathclose%
\pgfusepath{fill}%
\end{pgfscope}%
\begin{pgfscope}%
\pgfpathrectangle{\pgfqpoint{0.407592in}{0.222904in}}{\pgfqpoint{1.411163in}{1.411163in}}%
\pgfusepath{clip}%
\pgfsetbuttcap%
\pgfsetroundjoin%
\definecolor{currentfill}{rgb}{0.556863,0.654902,0.847059}%
\pgfsetfillcolor{currentfill}%
\pgfsetfillopacity{0.130000}%
\pgfsetlinewidth{0.000000pt}%
\definecolor{currentstroke}{rgb}{0.000000,0.000000,0.000000}%
\pgfsetstrokecolor{currentstroke}%
\pgfsetdash{}{0pt}%
\pgfpathmoveto{\pgfqpoint{1.089778in}{0.784632in}}%
\pgfpathlineto{\pgfqpoint{1.106801in}{0.789419in}}%
\pgfpathlineto{\pgfqpoint{1.123886in}{0.794886in}}%
\pgfpathlineto{\pgfqpoint{1.132314in}{0.791023in}}%
\pgfpathlineto{\pgfqpoint{1.140706in}{0.787452in}}%
\pgfpathlineto{\pgfqpoint{1.123663in}{0.781453in}}%
\pgfpathlineto{\pgfqpoint{1.106684in}{0.776140in}}%
\pgfpathlineto{\pgfqpoint{1.098249in}{0.780248in}}%
\pgfpathlineto{\pgfqpoint{1.089778in}{0.784632in}}%
\pgfpathclose%
\pgfusepath{fill}%
\end{pgfscope}%
\begin{pgfscope}%
\pgfpathrectangle{\pgfqpoint{0.407592in}{0.222904in}}{\pgfqpoint{1.411163in}{1.411163in}}%
\pgfusepath{clip}%
\pgfsetbuttcap%
\pgfsetroundjoin%
\definecolor{currentfill}{rgb}{0.556863,0.654902,0.847059}%
\pgfsetfillcolor{currentfill}%
\pgfsetfillopacity{0.130000}%
\pgfsetlinewidth{0.000000pt}%
\definecolor{currentstroke}{rgb}{0.000000,0.000000,0.000000}%
\pgfsetstrokecolor{currentstroke}%
\pgfsetdash{}{0pt}%
\pgfpathmoveto{\pgfqpoint{1.683766in}{1.082745in}}%
\pgfpathlineto{\pgfqpoint{1.702114in}{1.095957in}}%
\pgfpathlineto{\pgfqpoint{1.709671in}{1.103395in}}%
\pgfpathlineto{\pgfqpoint{1.717215in}{1.111103in}}%
\pgfpathlineto{\pgfqpoint{1.698892in}{1.098045in}}%
\pgfpathlineto{\pgfqpoint{1.691336in}{1.090258in}}%
\pgfusepath{fill}%
\end{pgfscope}%
\begin{pgfscope}%
\pgfpathrectangle{\pgfqpoint{0.407592in}{0.222904in}}{\pgfqpoint{1.411163in}{1.411163in}}%
\pgfusepath{clip}%
\pgfsetbuttcap%
\pgfsetroundjoin%
\definecolor{currentfill}{rgb}{0.556863,0.654902,0.847059}%
\pgfsetfillcolor{currentfill}%
\pgfsetfillopacity{0.130000}%
\pgfsetlinewidth{0.000000pt}%
\definecolor{currentstroke}{rgb}{0.000000,0.000000,0.000000}%
\pgfsetstrokecolor{currentstroke}%
\pgfsetdash{}{0pt}%
\pgfpathmoveto{\pgfqpoint{1.513763in}{0.954039in}}%
\pgfpathlineto{\pgfqpoint{1.531712in}{0.966941in}}%
\pgfpathlineto{\pgfqpoint{1.549743in}{0.979955in}}%
\pgfpathlineto{\pgfqpoint{1.557480in}{0.985165in}}%
\pgfpathlineto{\pgfqpoint{1.565200in}{0.990715in}}%
\pgfpathlineto{\pgfqpoint{1.547193in}{0.977731in}}%
\pgfpathlineto{\pgfqpoint{1.529268in}{0.964860in}}%
\pgfpathlineto{\pgfqpoint{1.521524in}{0.959267in}}%
\pgfpathlineto{\pgfqpoint{1.513763in}{0.954039in}}%
\pgfpathclose%
\pgfusepath{fill}%
\end{pgfscope}%
\begin{pgfscope}%
\pgfpathrectangle{\pgfqpoint{0.407592in}{0.222904in}}{\pgfqpoint{1.411163in}{1.411163in}}%
\pgfusepath{clip}%
\pgfsetbuttcap%
\pgfsetroundjoin%
\definecolor{currentfill}{rgb}{0.556863,0.654902,0.847059}%
\pgfsetfillcolor{currentfill}%
\pgfsetfillopacity{0.130000}%
\pgfsetlinewidth{0.000000pt}%
\definecolor{currentstroke}{rgb}{0.000000,0.000000,0.000000}%
\pgfsetstrokecolor{currentstroke}%
\pgfsetdash{}{0pt}%
\pgfpathmoveto{\pgfqpoint{0.902599in}{0.815369in}}%
\pgfpathlineto{\pgfqpoint{0.919394in}{0.812536in}}%
\pgfpathlineto{\pgfqpoint{0.936226in}{0.810562in}}%
\pgfpathlineto{\pgfqpoint{0.945087in}{0.803976in}}%
\pgfpathlineto{\pgfqpoint{0.953907in}{0.797518in}}%
\pgfpathlineto{\pgfqpoint{0.937129in}{0.799331in}}%
\pgfpathlineto{\pgfqpoint{0.920388in}{0.802059in}}%
\pgfpathlineto{\pgfqpoint{0.911515in}{0.808655in}}%
\pgfpathlineto{\pgfqpoint{0.902599in}{0.815369in}}%
\pgfpathclose%
\pgfusepath{fill}%
\end{pgfscope}%
\begin{pgfscope}%
\pgfpathrectangle{\pgfqpoint{0.407592in}{0.222904in}}{\pgfqpoint{1.411163in}{1.411163in}}%
\pgfusepath{clip}%
\pgfsetbuttcap%
\pgfsetroundjoin%
\definecolor{currentfill}{rgb}{0.556863,0.654902,0.847059}%
\pgfsetfillcolor{currentfill}%
\pgfsetfillopacity{0.130000}%
\pgfsetlinewidth{0.000000pt}%
\definecolor{currentstroke}{rgb}{0.000000,0.000000,0.000000}%
\pgfsetstrokecolor{currentstroke}%
\pgfsetdash{}{0pt}%
\pgfpathmoveto{\pgfqpoint{1.344419in}{0.853544in}}%
\pgfpathlineto{\pgfqpoint{1.361984in}{0.865032in}}%
\pgfpathlineto{\pgfqpoint{1.379628in}{0.876730in}}%
\pgfpathlineto{\pgfqpoint{1.387596in}{0.878660in}}%
\pgfpathlineto{\pgfqpoint{1.395541in}{0.880990in}}%
\pgfpathlineto{\pgfqpoint{1.377923in}{0.869064in}}%
\pgfpathlineto{\pgfqpoint{1.360383in}{0.857330in}}%
\pgfpathlineto{\pgfqpoint{1.352413in}{0.855223in}}%
\pgfpathlineto{\pgfqpoint{1.344419in}{0.853544in}}%
\pgfpathclose%
\pgfusepath{fill}%
\end{pgfscope}%
\begin{pgfscope}%
\pgfpathrectangle{\pgfqpoint{0.407592in}{0.222904in}}{\pgfqpoint{1.411163in}{1.411163in}}%
\pgfusepath{clip}%
\pgfsetbuttcap%
\pgfsetroundjoin%
\definecolor{currentfill}{rgb}{0.556863,0.654902,0.847059}%
\pgfsetfillcolor{currentfill}%
\pgfsetfillopacity{0.130000}%
\pgfsetlinewidth{0.000000pt}%
\definecolor{currentstroke}{rgb}{0.000000,0.000000,0.000000}%
\pgfsetstrokecolor{currentstroke}%
\pgfsetdash{}{0pt}%
\pgfpathmoveto{\pgfqpoint{1.038754in}{0.787757in}}%
\pgfpathlineto{\pgfqpoint{1.055710in}{0.790569in}}%
\pgfpathlineto{\pgfqpoint{1.072722in}{0.794145in}}%
\pgfpathlineto{\pgfqpoint{1.081269in}{0.789271in}}%
\pgfpathlineto{\pgfqpoint{1.089778in}{0.784632in}}%
\pgfpathlineto{\pgfqpoint{1.072813in}{0.780580in}}%
\pgfpathlineto{\pgfqpoint{1.055905in}{0.777314in}}%
\pgfpathlineto{\pgfqpoint{1.047349in}{0.782428in}}%
\pgfpathlineto{\pgfqpoint{1.038754in}{0.787757in}}%
\pgfpathclose%
\pgfusepath{fill}%
\end{pgfscope}%
\begin{pgfscope}%
\pgfpathrectangle{\pgfqpoint{0.407592in}{0.222904in}}{\pgfqpoint{1.411163in}{1.411163in}}%
\pgfusepath{clip}%
\pgfsetbuttcap%
\pgfsetroundjoin%
\definecolor{currentfill}{rgb}{0.556863,0.654902,0.847059}%
\pgfsetfillcolor{currentfill}%
\pgfsetfillopacity{0.130000}%
\pgfsetlinewidth{0.000000pt}%
\definecolor{currentstroke}{rgb}{0.000000,0.000000,0.000000}%
\pgfsetstrokecolor{currentstroke}%
\pgfsetdash{}{0pt}%
\pgfpathmoveto{\pgfqpoint{1.395541in}{0.880990in}}%
\pgfpathlineto{\pgfqpoint{1.413239in}{0.893092in}}%
\pgfpathlineto{\pgfqpoint{1.431018in}{0.905360in}}%
\pgfpathlineto{\pgfqpoint{1.438916in}{0.908260in}}%
\pgfpathlineto{\pgfqpoint{1.446793in}{0.911546in}}%
\pgfpathlineto{\pgfqpoint{1.429039in}{0.899135in}}%
\pgfpathlineto{\pgfqpoint{1.411366in}{0.886876in}}%
\pgfpathlineto{\pgfqpoint{1.403464in}{0.883727in}}%
\pgfpathlineto{\pgfqpoint{1.395541in}{0.880990in}}%
\pgfpathclose%
\pgfusepath{fill}%
\end{pgfscope}%
\begin{pgfscope}%
\pgfpathrectangle{\pgfqpoint{0.407592in}{0.222904in}}{\pgfqpoint{1.411163in}{1.411163in}}%
\pgfusepath{clip}%
\pgfsetbuttcap%
\pgfsetroundjoin%
\definecolor{currentfill}{rgb}{0.556863,0.654902,0.847059}%
\pgfsetfillcolor{currentfill}%
\pgfsetfillopacity{0.130000}%
\pgfsetlinewidth{0.000000pt}%
\definecolor{currentstroke}{rgb}{0.000000,0.000000,0.000000}%
\pgfsetstrokecolor{currentstroke}%
\pgfsetdash{}{0pt}%
\pgfpathmoveto{\pgfqpoint{1.565200in}{0.990715in}}%
\pgfpathlineto{\pgfqpoint{1.583290in}{1.003808in}}%
\pgfpathlineto{\pgfqpoint{1.601462in}{1.017006in}}%
\pgfpathlineto{\pgfqpoint{1.609142in}{1.022835in}}%
\pgfpathlineto{\pgfqpoint{1.616807in}{1.028979in}}%
\pgfpathlineto{\pgfqpoint{1.598659in}{1.015851in}}%
\pgfpathlineto{\pgfqpoint{1.580594in}{1.002830in}}%
\pgfpathlineto{\pgfqpoint{1.572905in}{0.996605in}}%
\pgfpathlineto{\pgfqpoint{1.565200in}{0.990715in}}%
\pgfpathclose%
\pgfusepath{fill}%
\end{pgfscope}%
\begin{pgfscope}%
\pgfpathrectangle{\pgfqpoint{0.407592in}{0.222904in}}{\pgfqpoint{1.411163in}{1.411163in}}%
\pgfusepath{clip}%
\pgfsetbuttcap%
\pgfsetroundjoin%
\definecolor{currentfill}{rgb}{0.556863,0.654902,0.847059}%
\pgfsetfillcolor{currentfill}%
\pgfsetfillopacity{0.130000}%
\pgfsetlinewidth{0.000000pt}%
\definecolor{currentstroke}{rgb}{0.000000,0.000000,0.000000}%
\pgfsetstrokecolor{currentstroke}%
\pgfsetdash{}{0pt}%
\pgfpathmoveto{\pgfqpoint{0.987584in}{0.796565in}}%
\pgfpathlineto{\pgfqpoint{1.004490in}{0.797373in}}%
\pgfpathlineto{\pgfqpoint{1.021445in}{0.798995in}}%
\pgfpathlineto{\pgfqpoint{1.030119in}{0.793284in}}%
\pgfpathlineto{\pgfqpoint{1.038754in}{0.787757in}}%
\pgfpathlineto{\pgfqpoint{1.021850in}{0.785755in}}%
\pgfpathlineto{\pgfqpoint{1.004995in}{0.784606in}}%
\pgfpathlineto{\pgfqpoint{0.996310in}{0.790503in}}%
\pgfpathlineto{\pgfqpoint{0.987584in}{0.796565in}}%
\pgfpathclose%
\pgfusepath{fill}%
\end{pgfscope}%
\begin{pgfscope}%
\pgfpathrectangle{\pgfqpoint{0.407592in}{0.222904in}}{\pgfqpoint{1.411163in}{1.411163in}}%
\pgfusepath{clip}%
\pgfsetbuttcap%
\pgfsetroundjoin%
\definecolor{currentfill}{rgb}{0.556863,0.654902,0.847059}%
\pgfsetfillcolor{currentfill}%
\pgfsetfillopacity{0.130000}%
\pgfsetlinewidth{0.000000pt}%
\definecolor{currentstroke}{rgb}{0.000000,0.000000,0.000000}%
\pgfsetstrokecolor{currentstroke}%
\pgfsetdash{}{0pt}%
\pgfpathmoveto{\pgfqpoint{0.851042in}{0.837422in}}%
\pgfpathlineto{\pgfqpoint{0.867826in}{0.832864in}}%
\pgfpathlineto{\pgfqpoint{0.884641in}{0.829124in}}%
\pgfpathlineto{\pgfqpoint{0.893641in}{0.822194in}}%
\pgfpathlineto{\pgfqpoint{0.902599in}{0.815369in}}%
\pgfpathlineto{\pgfqpoint{0.885838in}{0.819069in}}%
\pgfpathlineto{\pgfqpoint{0.869109in}{0.823639in}}%
\pgfpathlineto{\pgfqpoint{0.860097in}{0.830479in}}%
\pgfpathlineto{\pgfqpoint{0.851042in}{0.837422in}}%
\pgfpathclose%
\pgfusepath{fill}%
\end{pgfscope}%
\begin{pgfscope}%
\pgfpathrectangle{\pgfqpoint{0.407592in}{0.222904in}}{\pgfqpoint{1.411163in}{1.411163in}}%
\pgfusepath{clip}%
\pgfsetbuttcap%
\pgfsetroundjoin%
\definecolor{currentfill}{rgb}{0.556863,0.654902,0.847059}%
\pgfsetfillcolor{currentfill}%
\pgfsetfillopacity{0.130000}%
\pgfsetlinewidth{0.000000pt}%
\definecolor{currentstroke}{rgb}{0.000000,0.000000,0.000000}%
\pgfsetstrokecolor{currentstroke}%
\pgfsetdash{}{0pt}%
\pgfpathmoveto{\pgfqpoint{1.446793in}{0.911546in}}%
\pgfpathlineto{\pgfqpoint{1.464627in}{0.924100in}}%
\pgfpathlineto{\pgfqpoint{1.482543in}{0.936789in}}%
\pgfpathlineto{\pgfqpoint{1.490376in}{0.940553in}}%
\pgfpathlineto{\pgfqpoint{1.498190in}{0.944682in}}%
\pgfpathlineto{\pgfqpoint{1.480299in}{0.931920in}}%
\pgfpathlineto{\pgfqpoint{1.462488in}{0.919286in}}%
\pgfpathlineto{\pgfqpoint{1.454650in}{0.915220in}}%
\pgfpathlineto{\pgfqpoint{1.446793in}{0.911546in}}%
\pgfpathclose%
\pgfusepath{fill}%
\end{pgfscope}%
\begin{pgfscope}%
\pgfpathrectangle{\pgfqpoint{0.407592in}{0.222904in}}{\pgfqpoint{1.411163in}{1.411163in}}%
\pgfusepath{clip}%
\pgfsetbuttcap%
\pgfsetroundjoin%
\definecolor{currentfill}{rgb}{0.556863,0.654902,0.847059}%
\pgfsetfillcolor{currentfill}%
\pgfsetfillopacity{0.130000}%
\pgfsetlinewidth{0.000000pt}%
\definecolor{currentstroke}{rgb}{0.000000,0.000000,0.000000}%
\pgfsetstrokecolor{currentstroke}%
\pgfsetdash{}{0pt}%
\pgfpathmoveto{\pgfqpoint{1.616807in}{1.028979in}}%
\pgfpathlineto{\pgfqpoint{1.635037in}{1.042211in}}%
\pgfpathlineto{\pgfqpoint{1.653351in}{1.055544in}}%
\pgfpathlineto{\pgfqpoint{1.660976in}{1.061910in}}%
\pgfpathlineto{\pgfqpoint{1.668587in}{1.068568in}}%
\pgfpathlineto{\pgfqpoint{1.650298in}{1.055333in}}%
\pgfpathlineto{\pgfqpoint{1.632092in}{1.042204in}}%
\pgfpathlineto{\pgfqpoint{1.624456in}{1.035437in}}%
\pgfpathlineto{\pgfqpoint{1.616807in}{1.028979in}}%
\pgfpathclose%
\pgfusepath{fill}%
\end{pgfscope}%
\begin{pgfscope}%
\pgfpathrectangle{\pgfqpoint{0.407592in}{0.222904in}}{\pgfqpoint{1.411163in}{1.411163in}}%
\pgfusepath{clip}%
\pgfsetbuttcap%
\pgfsetroundjoin%
\definecolor{currentfill}{rgb}{0.556863,0.654902,0.847059}%
\pgfsetfillcolor{currentfill}%
\pgfsetfillopacity{0.130000}%
\pgfsetlinewidth{0.000000pt}%
\definecolor{currentstroke}{rgb}{0.000000,0.000000,0.000000}%
\pgfsetstrokecolor{currentstroke}%
\pgfsetdash{}{0pt}%
\pgfpathmoveto{\pgfqpoint{1.668587in}{1.068568in}}%
\pgfpathlineto{\pgfqpoint{1.686960in}{1.081904in}}%
\pgfpathlineto{\pgfqpoint{1.694544in}{1.088792in}}%
\pgfpathlineto{\pgfqpoint{1.702114in}{1.095957in}}%
\pgfpathlineto{\pgfqpoint{1.683766in}{1.082745in}}%
\pgfpathlineto{\pgfqpoint{1.676183in}{1.075514in}}%
\pgfusepath{fill}%
\end{pgfscope}%
\begin{pgfscope}%
\pgfpathrectangle{\pgfqpoint{0.407592in}{0.222904in}}{\pgfqpoint{1.411163in}{1.411163in}}%
\pgfusepath{clip}%
\pgfsetbuttcap%
\pgfsetroundjoin%
\definecolor{currentfill}{rgb}{0.556863,0.654902,0.847059}%
\pgfsetfillcolor{currentfill}%
\pgfsetfillopacity{0.130000}%
\pgfsetlinewidth{0.000000pt}%
\definecolor{currentstroke}{rgb}{0.000000,0.000000,0.000000}%
\pgfsetstrokecolor{currentstroke}%
\pgfsetdash{}{0pt}%
\pgfpathmoveto{\pgfqpoint{0.936226in}{0.810562in}}%
\pgfpathlineto{\pgfqpoint{0.953097in}{0.809434in}}%
\pgfpathlineto{\pgfqpoint{0.970010in}{0.809134in}}%
\pgfpathlineto{\pgfqpoint{0.978818in}{0.802779in}}%
\pgfpathlineto{\pgfqpoint{0.987584in}{0.796565in}}%
\pgfpathlineto{\pgfqpoint{0.970724in}{0.796604in}}%
\pgfpathlineto{\pgfqpoint{0.953907in}{0.797518in}}%
\pgfpathlineto{\pgfqpoint{0.945087in}{0.803976in}}%
\pgfpathlineto{\pgfqpoint{0.936226in}{0.810562in}}%
\pgfpathclose%
\pgfusepath{fill}%
\end{pgfscope}%
\begin{pgfscope}%
\pgfpathrectangle{\pgfqpoint{0.407592in}{0.222904in}}{\pgfqpoint{1.411163in}{1.411163in}}%
\pgfusepath{clip}%
\pgfsetbuttcap%
\pgfsetroundjoin%
\definecolor{currentfill}{rgb}{0.556863,0.654902,0.847059}%
\pgfsetfillcolor{currentfill}%
\pgfsetfillopacity{0.130000}%
\pgfsetlinewidth{0.000000pt}%
\definecolor{currentstroke}{rgb}{0.000000,0.000000,0.000000}%
\pgfsetstrokecolor{currentstroke}%
\pgfsetdash{}{0pt}%
\pgfpathmoveto{\pgfqpoint{1.174986in}{0.801278in}}%
\pgfpathlineto{\pgfqpoint{1.192228in}{0.809005in}}%
\pgfpathlineto{\pgfqpoint{1.209540in}{0.817214in}}%
\pgfpathlineto{\pgfqpoint{1.217822in}{0.815000in}}%
\pgfpathlineto{\pgfqpoint{1.226070in}{0.813128in}}%
\pgfpathlineto{\pgfqpoint{1.208794in}{0.804421in}}%
\pgfpathlineto{\pgfqpoint{1.191588in}{0.796177in}}%
\pgfpathlineto{\pgfqpoint{1.183304in}{0.798557in}}%
\pgfpathlineto{\pgfqpoint{1.174986in}{0.801278in}}%
\pgfpathclose%
\pgfusepath{fill}%
\end{pgfscope}%
\begin{pgfscope}%
\pgfpathrectangle{\pgfqpoint{0.407592in}{0.222904in}}{\pgfqpoint{1.411163in}{1.411163in}}%
\pgfusepath{clip}%
\pgfsetbuttcap%
\pgfsetroundjoin%
\definecolor{currentfill}{rgb}{0.556863,0.654902,0.847059}%
\pgfsetfillcolor{currentfill}%
\pgfsetfillopacity{0.130000}%
\pgfsetlinewidth{0.000000pt}%
\definecolor{currentstroke}{rgb}{0.000000,0.000000,0.000000}%
\pgfsetstrokecolor{currentstroke}%
\pgfsetdash{}{0pt}%
\pgfpathmoveto{\pgfqpoint{1.226070in}{0.813128in}}%
\pgfpathlineto{\pgfqpoint{1.243419in}{0.822256in}}%
\pgfpathlineto{\pgfqpoint{1.260842in}{0.831767in}}%
\pgfpathlineto{\pgfqpoint{1.269027in}{0.830712in}}%
\pgfpathlineto{\pgfqpoint{1.277182in}{0.830025in}}%
\pgfpathlineto{\pgfqpoint{1.259790in}{0.820073in}}%
\pgfpathlineto{\pgfqpoint{1.242474in}{0.810480in}}%
\pgfpathlineto{\pgfqpoint{1.234287in}{0.811615in}}%
\pgfpathlineto{\pgfqpoint{1.226070in}{0.813128in}}%
\pgfpathclose%
\pgfusepath{fill}%
\end{pgfscope}%
\begin{pgfscope}%
\pgfpathrectangle{\pgfqpoint{0.407592in}{0.222904in}}{\pgfqpoint{1.411163in}{1.411163in}}%
\pgfusepath{clip}%
\pgfsetbuttcap%
\pgfsetroundjoin%
\definecolor{currentfill}{rgb}{0.556863,0.654902,0.847059}%
\pgfsetfillcolor{currentfill}%
\pgfsetfillopacity{0.130000}%
\pgfsetlinewidth{0.000000pt}%
\definecolor{currentstroke}{rgb}{0.000000,0.000000,0.000000}%
\pgfsetstrokecolor{currentstroke}%
\pgfsetdash{}{0pt}%
\pgfpathmoveto{\pgfqpoint{1.498190in}{0.944682in}}%
\pgfpathlineto{\pgfqpoint{1.516163in}{0.957565in}}%
\pgfpathlineto{\pgfqpoint{1.534218in}{0.970564in}}%
\pgfpathlineto{\pgfqpoint{1.541989in}{0.975088in}}%
\pgfpathlineto{\pgfqpoint{1.549743in}{0.979955in}}%
\pgfpathlineto{\pgfqpoint{1.531712in}{0.966941in}}%
\pgfpathlineto{\pgfqpoint{1.513763in}{0.954039in}}%
\pgfpathlineto{\pgfqpoint{1.505986in}{0.949177in}}%
\pgfpathlineto{\pgfqpoint{1.498190in}{0.944682in}}%
\pgfpathclose%
\pgfusepath{fill}%
\end{pgfscope}%
\begin{pgfscope}%
\pgfpathrectangle{\pgfqpoint{0.407592in}{0.222904in}}{\pgfqpoint{1.411163in}{1.411163in}}%
\pgfusepath{clip}%
\pgfsetbuttcap%
\pgfsetroundjoin%
\definecolor{currentfill}{rgb}{0.556863,0.654902,0.847059}%
\pgfsetfillcolor{currentfill}%
\pgfsetfillopacity{0.130000}%
\pgfsetlinewidth{0.000000pt}%
\definecolor{currentstroke}{rgb}{0.000000,0.000000,0.000000}%
\pgfsetstrokecolor{currentstroke}%
\pgfsetdash{}{0pt}%
\pgfpathmoveto{\pgfqpoint{1.123886in}{0.794886in}}%
\pgfpathlineto{\pgfqpoint{1.141034in}{0.800983in}}%
\pgfpathlineto{\pgfqpoint{1.158247in}{0.807663in}}%
\pgfpathlineto{\pgfqpoint{1.166634in}{0.804320in}}%
\pgfpathlineto{\pgfqpoint{1.174986in}{0.801278in}}%
\pgfpathlineto{\pgfqpoint{1.157813in}{0.794077in}}%
\pgfpathlineto{\pgfqpoint{1.140706in}{0.787452in}}%
\pgfpathlineto{\pgfqpoint{1.132314in}{0.791023in}}%
\pgfpathlineto{\pgfqpoint{1.123886in}{0.794886in}}%
\pgfpathclose%
\pgfusepath{fill}%
\end{pgfscope}%
\begin{pgfscope}%
\pgfpathrectangle{\pgfqpoint{0.407592in}{0.222904in}}{\pgfqpoint{1.411163in}{1.411163in}}%
\pgfusepath{clip}%
\pgfsetbuttcap%
\pgfsetroundjoin%
\definecolor{currentfill}{rgb}{0.556863,0.654902,0.847059}%
\pgfsetfillcolor{currentfill}%
\pgfsetfillopacity{0.130000}%
\pgfsetlinewidth{0.000000pt}%
\definecolor{currentstroke}{rgb}{0.000000,0.000000,0.000000}%
\pgfsetstrokecolor{currentstroke}%
\pgfsetdash{}{0pt}%
\pgfpathmoveto{\pgfqpoint{1.277182in}{0.830025in}}%
\pgfpathlineto{\pgfqpoint{1.294648in}{0.840304in}}%
\pgfpathlineto{\pgfqpoint{1.312192in}{0.850882in}}%
\pgfpathlineto{\pgfqpoint{1.320288in}{0.850962in}}%
\pgfpathlineto{\pgfqpoint{1.328358in}{0.851425in}}%
\pgfpathlineto{\pgfqpoint{1.310843in}{0.840485in}}%
\pgfpathlineto{\pgfqpoint{1.293406in}{0.829819in}}%
\pgfpathlineto{\pgfqpoint{1.285308in}{0.829722in}}%
\pgfpathlineto{\pgfqpoint{1.277182in}{0.830025in}}%
\pgfpathclose%
\pgfusepath{fill}%
\end{pgfscope}%
\begin{pgfscope}%
\pgfpathrectangle{\pgfqpoint{0.407592in}{0.222904in}}{\pgfqpoint{1.411163in}{1.411163in}}%
\pgfusepath{clip}%
\pgfsetbuttcap%
\pgfsetroundjoin%
\definecolor{currentfill}{rgb}{0.556863,0.654902,0.847059}%
\pgfsetfillcolor{currentfill}%
\pgfsetfillopacity{0.130000}%
\pgfsetlinewidth{0.000000pt}%
\definecolor{currentstroke}{rgb}{0.000000,0.000000,0.000000}%
\pgfsetstrokecolor{currentstroke}%
\pgfsetdash{}{0pt}%
\pgfpathmoveto{\pgfqpoint{1.328358in}{0.851425in}}%
\pgfpathlineto{\pgfqpoint{1.345950in}{0.862615in}}%
\pgfpathlineto{\pgfqpoint{1.363621in}{0.874036in}}%
\pgfpathlineto{\pgfqpoint{1.371637in}{0.875192in}}%
\pgfpathlineto{\pgfqpoint{1.379628in}{0.876730in}}%
\pgfpathlineto{\pgfqpoint{1.361984in}{0.865032in}}%
\pgfpathlineto{\pgfqpoint{1.344419in}{0.853544in}}%
\pgfpathlineto{\pgfqpoint{1.336401in}{0.852281in}}%
\pgfpathlineto{\pgfqpoint{1.328358in}{0.851425in}}%
\pgfpathclose%
\pgfusepath{fill}%
\end{pgfscope}%
\begin{pgfscope}%
\pgfpathrectangle{\pgfqpoint{0.407592in}{0.222904in}}{\pgfqpoint{1.411163in}{1.411163in}}%
\pgfusepath{clip}%
\pgfsetbuttcap%
\pgfsetroundjoin%
\definecolor{currentfill}{rgb}{0.556863,0.654902,0.847059}%
\pgfsetfillcolor{currentfill}%
\pgfsetfillopacity{0.130000}%
\pgfsetlinewidth{0.000000pt}%
\definecolor{currentstroke}{rgb}{0.000000,0.000000,0.000000}%
\pgfsetstrokecolor{currentstroke}%
\pgfsetdash{}{0pt}%
\pgfpathmoveto{\pgfqpoint{1.072722in}{0.794145in}}%
\pgfpathlineto{\pgfqpoint{1.089791in}{0.798439in}}%
\pgfpathlineto{\pgfqpoint{1.106920in}{0.803406in}}%
\pgfpathlineto{\pgfqpoint{1.115421in}{0.799019in}}%
\pgfpathlineto{\pgfqpoint{1.123886in}{0.794886in}}%
\pgfpathlineto{\pgfqpoint{1.106801in}{0.789419in}}%
\pgfpathlineto{\pgfqpoint{1.089778in}{0.784632in}}%
\pgfpathlineto{\pgfqpoint{1.081269in}{0.789271in}}%
\pgfpathlineto{\pgfqpoint{1.072722in}{0.794145in}}%
\pgfpathclose%
\pgfusepath{fill}%
\end{pgfscope}%
\begin{pgfscope}%
\pgfpathrectangle{\pgfqpoint{0.407592in}{0.222904in}}{\pgfqpoint{1.411163in}{1.411163in}}%
\pgfusepath{clip}%
\pgfsetbuttcap%
\pgfsetroundjoin%
\definecolor{currentfill}{rgb}{0.556863,0.654902,0.847059}%
\pgfsetfillcolor{currentfill}%
\pgfsetfillopacity{0.130000}%
\pgfsetlinewidth{0.000000pt}%
\definecolor{currentstroke}{rgb}{0.000000,0.000000,0.000000}%
\pgfsetstrokecolor{currentstroke}%
\pgfsetdash{}{0pt}%
\pgfpathmoveto{\pgfqpoint{1.549743in}{0.979955in}}%
\pgfpathlineto{\pgfqpoint{1.567857in}{0.993078in}}%
\pgfpathlineto{\pgfqpoint{1.586053in}{1.006306in}}%
\pgfpathlineto{\pgfqpoint{1.593765in}{1.011496in}}%
\pgfpathlineto{\pgfqpoint{1.601462in}{1.017006in}}%
\pgfpathlineto{\pgfqpoint{1.583290in}{1.003808in}}%
\pgfpathlineto{\pgfqpoint{1.565200in}{0.990715in}}%
\pgfpathlineto{\pgfqpoint{1.557480in}{0.985165in}}%
\pgfpathlineto{\pgfqpoint{1.549743in}{0.979955in}}%
\pgfpathclose%
\pgfusepath{fill}%
\end{pgfscope}%
\begin{pgfscope}%
\pgfpathrectangle{\pgfqpoint{0.407592in}{0.222904in}}{\pgfqpoint{1.411163in}{1.411163in}}%
\pgfusepath{clip}%
\pgfsetbuttcap%
\pgfsetroundjoin%
\definecolor{currentfill}{rgb}{0.556863,0.654902,0.847059}%
\pgfsetfillcolor{currentfill}%
\pgfsetfillopacity{0.130000}%
\pgfsetlinewidth{0.000000pt}%
\definecolor{currentstroke}{rgb}{0.000000,0.000000,0.000000}%
\pgfsetstrokecolor{currentstroke}%
\pgfsetdash{}{0pt}%
\pgfpathmoveto{\pgfqpoint{1.379628in}{0.876730in}}%
\pgfpathlineto{\pgfqpoint{1.397352in}{0.888622in}}%
\pgfpathlineto{\pgfqpoint{1.415156in}{0.900695in}}%
\pgfpathlineto{\pgfqpoint{1.423098in}{0.902840in}}%
\pgfpathlineto{\pgfqpoint{1.431018in}{0.905360in}}%
\pgfpathlineto{\pgfqpoint{1.413239in}{0.893092in}}%
\pgfpathlineto{\pgfqpoint{1.395541in}{0.880990in}}%
\pgfpathlineto{\pgfqpoint{1.387596in}{0.878660in}}%
\pgfpathlineto{\pgfqpoint{1.379628in}{0.876730in}}%
\pgfpathclose%
\pgfusepath{fill}%
\end{pgfscope}%
\begin{pgfscope}%
\pgfpathrectangle{\pgfqpoint{0.407592in}{0.222904in}}{\pgfqpoint{1.411163in}{1.411163in}}%
\pgfusepath{clip}%
\pgfsetbuttcap%
\pgfsetroundjoin%
\definecolor{currentfill}{rgb}{0.556863,0.654902,0.847059}%
\pgfsetfillcolor{currentfill}%
\pgfsetfillopacity{0.130000}%
\pgfsetlinewidth{0.000000pt}%
\definecolor{currentstroke}{rgb}{0.000000,0.000000,0.000000}%
\pgfsetstrokecolor{currentstroke}%
\pgfsetdash{}{0pt}%
\pgfpathmoveto{\pgfqpoint{0.884641in}{0.829124in}}%
\pgfpathlineto{\pgfqpoint{0.901491in}{0.826198in}}%
\pgfpathlineto{\pgfqpoint{0.918377in}{0.824082in}}%
\pgfpathlineto{\pgfqpoint{0.927322in}{0.817267in}}%
\pgfpathlineto{\pgfqpoint{0.936226in}{0.810562in}}%
\pgfpathlineto{\pgfqpoint{0.919394in}{0.812536in}}%
\pgfpathlineto{\pgfqpoint{0.902599in}{0.815369in}}%
\pgfpathlineto{\pgfqpoint{0.893641in}{0.822194in}}%
\pgfpathlineto{\pgfqpoint{0.884641in}{0.829124in}}%
\pgfpathclose%
\pgfusepath{fill}%
\end{pgfscope}%
\begin{pgfscope}%
\pgfpathrectangle{\pgfqpoint{0.407592in}{0.222904in}}{\pgfqpoint{1.411163in}{1.411163in}}%
\pgfusepath{clip}%
\pgfsetbuttcap%
\pgfsetroundjoin%
\definecolor{currentfill}{rgb}{0.556863,0.654902,0.847059}%
\pgfsetfillcolor{currentfill}%
\pgfsetfillopacity{0.130000}%
\pgfsetlinewidth{0.000000pt}%
\definecolor{currentstroke}{rgb}{0.000000,0.000000,0.000000}%
\pgfsetstrokecolor{currentstroke}%
\pgfsetdash{}{0pt}%
\pgfpathmoveto{\pgfqpoint{1.021445in}{0.798995in}}%
\pgfpathlineto{\pgfqpoint{1.038451in}{0.801395in}}%
\pgfpathlineto{\pgfqpoint{1.055511in}{0.804534in}}%
\pgfpathlineto{\pgfqpoint{1.064136in}{0.799238in}}%
\pgfpathlineto{\pgfqpoint{1.072722in}{0.794145in}}%
\pgfpathlineto{\pgfqpoint{1.055710in}{0.790569in}}%
\pgfpathlineto{\pgfqpoint{1.038754in}{0.787757in}}%
\pgfpathlineto{\pgfqpoint{1.030119in}{0.793284in}}%
\pgfpathlineto{\pgfqpoint{1.021445in}{0.798995in}}%
\pgfpathclose%
\pgfusepath{fill}%
\end{pgfscope}%
\begin{pgfscope}%
\pgfpathrectangle{\pgfqpoint{0.407592in}{0.222904in}}{\pgfqpoint{1.411163in}{1.411163in}}%
\pgfusepath{clip}%
\pgfsetbuttcap%
\pgfsetroundjoin%
\definecolor{currentfill}{rgb}{0.556863,0.654902,0.847059}%
\pgfsetfillcolor{currentfill}%
\pgfsetfillopacity{0.130000}%
\pgfsetlinewidth{0.000000pt}%
\definecolor{currentstroke}{rgb}{0.000000,0.000000,0.000000}%
\pgfsetstrokecolor{currentstroke}%
\pgfsetdash{}{0pt}%
\pgfpathmoveto{\pgfqpoint{1.431018in}{0.905360in}}%
\pgfpathlineto{\pgfqpoint{1.448876in}{0.917782in}}%
\pgfpathlineto{\pgfqpoint{1.466817in}{0.930348in}}%
\pgfpathlineto{\pgfqpoint{1.474690in}{0.933389in}}%
\pgfpathlineto{\pgfqpoint{1.482543in}{0.936789in}}%
\pgfpathlineto{\pgfqpoint{1.464627in}{0.924100in}}%
\pgfpathlineto{\pgfqpoint{1.446793in}{0.911546in}}%
\pgfpathlineto{\pgfqpoint{1.438916in}{0.908260in}}%
\pgfpathlineto{\pgfqpoint{1.431018in}{0.905360in}}%
\pgfpathclose%
\pgfusepath{fill}%
\end{pgfscope}%
\begin{pgfscope}%
\pgfpathrectangle{\pgfqpoint{0.407592in}{0.222904in}}{\pgfqpoint{1.411163in}{1.411163in}}%
\pgfusepath{clip}%
\pgfsetbuttcap%
\pgfsetroundjoin%
\definecolor{currentfill}{rgb}{0.556863,0.654902,0.847059}%
\pgfsetfillcolor{currentfill}%
\pgfsetfillopacity{0.130000}%
\pgfsetlinewidth{0.000000pt}%
\definecolor{currentstroke}{rgb}{0.000000,0.000000,0.000000}%
\pgfsetstrokecolor{currentstroke}%
\pgfsetdash{}{0pt}%
\pgfpathmoveto{\pgfqpoint{1.601462in}{1.017006in}}%
\pgfpathlineto{\pgfqpoint{1.619717in}{1.030305in}}%
\pgfpathlineto{\pgfqpoint{1.638056in}{1.043703in}}%
\pgfpathlineto{\pgfqpoint{1.645711in}{1.049475in}}%
\pgfpathlineto{\pgfqpoint{1.653351in}{1.055544in}}%
\pgfpathlineto{\pgfqpoint{1.635037in}{1.042211in}}%
\pgfpathlineto{\pgfqpoint{1.616807in}{1.028979in}}%
\pgfpathlineto{\pgfqpoint{1.609142in}{1.022835in}}%
\pgfpathlineto{\pgfqpoint{1.601462in}{1.017006in}}%
\pgfpathclose%
\pgfusepath{fill}%
\end{pgfscope}%
\begin{pgfscope}%
\pgfpathrectangle{\pgfqpoint{0.407592in}{0.222904in}}{\pgfqpoint{1.411163in}{1.411163in}}%
\pgfusepath{clip}%
\pgfsetbuttcap%
\pgfsetroundjoin%
\definecolor{currentfill}{rgb}{0.556863,0.654902,0.847059}%
\pgfsetfillcolor{currentfill}%
\pgfsetfillopacity{0.130000}%
\pgfsetlinewidth{0.000000pt}%
\definecolor{currentstroke}{rgb}{0.000000,0.000000,0.000000}%
\pgfsetstrokecolor{currentstroke}%
\pgfsetdash{}{0pt}%
\pgfpathmoveto{\pgfqpoint{1.653351in}{1.055544in}}%
\pgfpathlineto{\pgfqpoint{1.671749in}{1.068975in}}%
\pgfpathlineto{\pgfqpoint{1.679362in}{1.075297in}}%
\pgfpathlineto{\pgfqpoint{1.686960in}{1.081904in}}%
\pgfpathlineto{\pgfqpoint{1.668587in}{1.068568in}}%
\pgfpathlineto{\pgfqpoint{1.660976in}{1.061910in}}%
\pgfusepath{fill}%
\end{pgfscope}%
\begin{pgfscope}%
\pgfpathrectangle{\pgfqpoint{0.407592in}{0.222904in}}{\pgfqpoint{1.411163in}{1.411163in}}%
\pgfusepath{clip}%
\pgfsetbuttcap%
\pgfsetroundjoin%
\definecolor{currentfill}{rgb}{0.556863,0.654902,0.847059}%
\pgfsetfillcolor{currentfill}%
\pgfsetfillopacity{0.130000}%
\pgfsetlinewidth{0.000000pt}%
\definecolor{currentstroke}{rgb}{0.000000,0.000000,0.000000}%
\pgfsetstrokecolor{currentstroke}%
\pgfsetdash{}{0pt}%
\pgfpathmoveto{\pgfqpoint{0.970010in}{0.809134in}}%
\pgfpathlineto{\pgfqpoint{0.986968in}{0.809636in}}%
\pgfpathlineto{\pgfqpoint{1.003975in}{0.810916in}}%
\pgfpathlineto{\pgfqpoint{1.012730in}{0.804876in}}%
\pgfpathlineto{\pgfqpoint{1.021445in}{0.798995in}}%
\pgfpathlineto{\pgfqpoint{1.004490in}{0.797373in}}%
\pgfpathlineto{\pgfqpoint{0.987584in}{0.796565in}}%
\pgfpathlineto{\pgfqpoint{0.978818in}{0.802779in}}%
\pgfpathlineto{\pgfqpoint{0.970010in}{0.809134in}}%
\pgfpathclose%
\pgfusepath{fill}%
\end{pgfscope}%
\begin{pgfscope}%
\pgfpathrectangle{\pgfqpoint{0.407592in}{0.222904in}}{\pgfqpoint{1.411163in}{1.411163in}}%
\pgfusepath{clip}%
\pgfsetbuttcap%
\pgfsetroundjoin%
\definecolor{currentfill}{rgb}{0.556863,0.654902,0.847059}%
\pgfsetfillcolor{currentfill}%
\pgfsetfillopacity{0.130000}%
\pgfsetlinewidth{0.000000pt}%
\definecolor{currentstroke}{rgb}{0.000000,0.000000,0.000000}%
\pgfsetstrokecolor{currentstroke}%
\pgfsetdash{}{0pt}%
\pgfpathmoveto{\pgfqpoint{1.482543in}{0.936789in}}%
\pgfpathlineto{\pgfqpoint{1.500540in}{0.949606in}}%
\pgfpathlineto{\pgfqpoint{1.518620in}{0.962544in}}%
\pgfpathlineto{\pgfqpoint{1.526428in}{0.966383in}}%
\pgfpathlineto{\pgfqpoint{1.534218in}{0.970564in}}%
\pgfpathlineto{\pgfqpoint{1.516163in}{0.957565in}}%
\pgfpathlineto{\pgfqpoint{1.498190in}{0.944682in}}%
\pgfpathlineto{\pgfqpoint{1.490376in}{0.940553in}}%
\pgfpathlineto{\pgfqpoint{1.482543in}{0.936789in}}%
\pgfpathclose%
\pgfusepath{fill}%
\end{pgfscope}%
\begin{pgfscope}%
\pgfpathrectangle{\pgfqpoint{0.407592in}{0.222904in}}{\pgfqpoint{1.411163in}{1.411163in}}%
\pgfusepath{clip}%
\pgfsetbuttcap%
\pgfsetroundjoin%
\definecolor{currentfill}{rgb}{0.556863,0.654902,0.847059}%
\pgfsetfillcolor{currentfill}%
\pgfsetfillopacity{0.130000}%
\pgfsetlinewidth{0.000000pt}%
\definecolor{currentstroke}{rgb}{0.000000,0.000000,0.000000}%
\pgfsetstrokecolor{currentstroke}%
\pgfsetdash{}{0pt}%
\pgfpathmoveto{\pgfqpoint{1.209540in}{0.817214in}}%
\pgfpathlineto{\pgfqpoint{1.226924in}{0.825864in}}%
\pgfpathlineto{\pgfqpoint{1.244381in}{0.834920in}}%
\pgfpathlineto{\pgfqpoint{1.252627in}{0.833175in}}%
\pgfpathlineto{\pgfqpoint{1.260842in}{0.831767in}}%
\pgfpathlineto{\pgfqpoint{1.243419in}{0.822256in}}%
\pgfpathlineto{\pgfqpoint{1.226070in}{0.813128in}}%
\pgfpathlineto{\pgfqpoint{1.217822in}{0.815000in}}%
\pgfpathlineto{\pgfqpoint{1.209540in}{0.817214in}}%
\pgfpathclose%
\pgfusepath{fill}%
\end{pgfscope}%
\begin{pgfscope}%
\pgfpathrectangle{\pgfqpoint{0.407592in}{0.222904in}}{\pgfqpoint{1.411163in}{1.411163in}}%
\pgfusepath{clip}%
\pgfsetbuttcap%
\pgfsetroundjoin%
\definecolor{currentfill}{rgb}{0.556863,0.654902,0.847059}%
\pgfsetfillcolor{currentfill}%
\pgfsetfillopacity{0.130000}%
\pgfsetlinewidth{0.000000pt}%
\definecolor{currentstroke}{rgb}{0.000000,0.000000,0.000000}%
\pgfsetstrokecolor{currentstroke}%
\pgfsetdash{}{0pt}%
\pgfpathmoveto{\pgfqpoint{1.158247in}{0.807663in}}%
\pgfpathlineto{\pgfqpoint{1.175528in}{0.814880in}}%
\pgfpathlineto{\pgfqpoint{1.192878in}{0.822594in}}%
\pgfpathlineto{\pgfqpoint{1.201226in}{0.819751in}}%
\pgfpathlineto{\pgfqpoint{1.209540in}{0.817214in}}%
\pgfpathlineto{\pgfqpoint{1.192228in}{0.809005in}}%
\pgfpathlineto{\pgfqpoint{1.174986in}{0.801278in}}%
\pgfpathlineto{\pgfqpoint{1.166634in}{0.804320in}}%
\pgfpathlineto{\pgfqpoint{1.158247in}{0.807663in}}%
\pgfpathclose%
\pgfusepath{fill}%
\end{pgfscope}%
\begin{pgfscope}%
\pgfpathrectangle{\pgfqpoint{0.407592in}{0.222904in}}{\pgfqpoint{1.411163in}{1.411163in}}%
\pgfusepath{clip}%
\pgfsetbuttcap%
\pgfsetroundjoin%
\definecolor{currentfill}{rgb}{0.556863,0.654902,0.847059}%
\pgfsetfillcolor{currentfill}%
\pgfsetfillopacity{0.130000}%
\pgfsetlinewidth{0.000000pt}%
\definecolor{currentstroke}{rgb}{0.000000,0.000000,0.000000}%
\pgfsetstrokecolor{currentstroke}%
\pgfsetdash{}{0pt}%
\pgfpathmoveto{\pgfqpoint{1.260842in}{0.831767in}}%
\pgfpathlineto{\pgfqpoint{1.278340in}{0.841630in}}%
\pgfpathlineto{\pgfqpoint{1.295914in}{0.851815in}}%
\pgfpathlineto{\pgfqpoint{1.304067in}{0.851171in}}%
\pgfpathlineto{\pgfqpoint{1.312192in}{0.850882in}}%
\pgfpathlineto{\pgfqpoint{1.294648in}{0.840304in}}%
\pgfpathlineto{\pgfqpoint{1.277182in}{0.830025in}}%
\pgfpathlineto{\pgfqpoint{1.269027in}{0.830712in}}%
\pgfpathlineto{\pgfqpoint{1.260842in}{0.831767in}}%
\pgfpathclose%
\pgfusepath{fill}%
\end{pgfscope}%
\begin{pgfscope}%
\pgfpathrectangle{\pgfqpoint{0.407592in}{0.222904in}}{\pgfqpoint{1.411163in}{1.411163in}}%
\pgfusepath{clip}%
\pgfsetbuttcap%
\pgfsetroundjoin%
\definecolor{currentfill}{rgb}{0.556863,0.654902,0.847059}%
\pgfsetfillcolor{currentfill}%
\pgfsetfillopacity{0.130000}%
\pgfsetlinewidth{0.000000pt}%
\definecolor{currentstroke}{rgb}{0.000000,0.000000,0.000000}%
\pgfsetstrokecolor{currentstroke}%
\pgfsetdash{}{0pt}%
\pgfpathmoveto{\pgfqpoint{0.918377in}{0.824082in}}%
\pgfpathlineto{\pgfqpoint{0.935302in}{0.822765in}}%
\pgfpathlineto{\pgfqpoint{0.952269in}{0.822230in}}%
\pgfpathlineto{\pgfqpoint{0.961161in}{0.815620in}}%
\pgfpathlineto{\pgfqpoint{0.970010in}{0.809134in}}%
\pgfpathlineto{\pgfqpoint{0.953097in}{0.809434in}}%
\pgfpathlineto{\pgfqpoint{0.936226in}{0.810562in}}%
\pgfpathlineto{\pgfqpoint{0.927322in}{0.817267in}}%
\pgfpathlineto{\pgfqpoint{0.918377in}{0.824082in}}%
\pgfpathclose%
\pgfusepath{fill}%
\end{pgfscope}%
\begin{pgfscope}%
\pgfpathrectangle{\pgfqpoint{0.407592in}{0.222904in}}{\pgfqpoint{1.411163in}{1.411163in}}%
\pgfusepath{clip}%
\pgfsetbuttcap%
\pgfsetroundjoin%
\definecolor{currentfill}{rgb}{0.556863,0.654902,0.847059}%
\pgfsetfillcolor{currentfill}%
\pgfsetfillopacity{0.130000}%
\pgfsetlinewidth{0.000000pt}%
\definecolor{currentstroke}{rgb}{0.000000,0.000000,0.000000}%
\pgfsetstrokecolor{currentstroke}%
\pgfsetdash{}{0pt}%
\pgfpathmoveto{\pgfqpoint{1.534218in}{0.970564in}}%
\pgfpathlineto{\pgfqpoint{1.552356in}{0.983674in}}%
\pgfpathlineto{\pgfqpoint{1.570577in}{0.996891in}}%
\pgfpathlineto{\pgfqpoint{1.578324in}{1.001437in}}%
\pgfpathlineto{\pgfqpoint{1.586053in}{1.006306in}}%
\pgfpathlineto{\pgfqpoint{1.567857in}{0.993078in}}%
\pgfpathlineto{\pgfqpoint{1.549743in}{0.979955in}}%
\pgfpathlineto{\pgfqpoint{1.541989in}{0.975088in}}%
\pgfpathlineto{\pgfqpoint{1.534218in}{0.970564in}}%
\pgfpathclose%
\pgfusepath{fill}%
\end{pgfscope}%
\begin{pgfscope}%
\pgfpathrectangle{\pgfqpoint{0.407592in}{0.222904in}}{\pgfqpoint{1.411163in}{1.411163in}}%
\pgfusepath{clip}%
\pgfsetbuttcap%
\pgfsetroundjoin%
\definecolor{currentfill}{rgb}{0.556863,0.654902,0.847059}%
\pgfsetfillcolor{currentfill}%
\pgfsetfillopacity{0.130000}%
\pgfsetlinewidth{0.000000pt}%
\definecolor{currentstroke}{rgb}{0.000000,0.000000,0.000000}%
\pgfsetstrokecolor{currentstroke}%
\pgfsetdash{}{0pt}%
\pgfpathmoveto{\pgfqpoint{1.312192in}{0.850882in}}%
\pgfpathlineto{\pgfqpoint{1.329812in}{0.861734in}}%
\pgfpathlineto{\pgfqpoint{1.347511in}{0.872838in}}%
\pgfpathlineto{\pgfqpoint{1.355579in}{0.873255in}}%
\pgfpathlineto{\pgfqpoint{1.363621in}{0.874036in}}%
\pgfpathlineto{\pgfqpoint{1.345950in}{0.862615in}}%
\pgfpathlineto{\pgfqpoint{1.328358in}{0.851425in}}%
\pgfpathlineto{\pgfqpoint{1.320288in}{0.850962in}}%
\pgfpathlineto{\pgfqpoint{1.312192in}{0.850882in}}%
\pgfpathclose%
\pgfusepath{fill}%
\end{pgfscope}%
\begin{pgfscope}%
\pgfpathrectangle{\pgfqpoint{0.407592in}{0.222904in}}{\pgfqpoint{1.411163in}{1.411163in}}%
\pgfusepath{clip}%
\pgfsetbuttcap%
\pgfsetroundjoin%
\definecolor{currentfill}{rgb}{0.556863,0.654902,0.847059}%
\pgfsetfillcolor{currentfill}%
\pgfsetfillopacity{0.130000}%
\pgfsetlinewidth{0.000000pt}%
\definecolor{currentstroke}{rgb}{0.000000,0.000000,0.000000}%
\pgfsetstrokecolor{currentstroke}%
\pgfsetdash{}{0pt}%
\pgfpathmoveto{\pgfqpoint{1.106920in}{0.803406in}}%
\pgfpathlineto{\pgfqpoint{1.124111in}{0.809000in}}%
\pgfpathlineto{\pgfqpoint{1.141367in}{0.815180in}}%
\pgfpathlineto{\pgfqpoint{1.149825in}{0.811288in}}%
\pgfpathlineto{\pgfqpoint{1.158247in}{0.807663in}}%
\pgfpathlineto{\pgfqpoint{1.141034in}{0.800983in}}%
\pgfpathlineto{\pgfqpoint{1.123886in}{0.794886in}}%
\pgfpathlineto{\pgfqpoint{1.115421in}{0.799019in}}%
\pgfpathlineto{\pgfqpoint{1.106920in}{0.803406in}}%
\pgfpathclose%
\pgfusepath{fill}%
\end{pgfscope}%
\begin{pgfscope}%
\pgfpathrectangle{\pgfqpoint{0.407592in}{0.222904in}}{\pgfqpoint{1.411163in}{1.411163in}}%
\pgfusepath{clip}%
\pgfsetbuttcap%
\pgfsetroundjoin%
\definecolor{currentfill}{rgb}{0.556863,0.654902,0.847059}%
\pgfsetfillcolor{currentfill}%
\pgfsetfillopacity{0.130000}%
\pgfsetlinewidth{0.000000pt}%
\definecolor{currentstroke}{rgb}{0.000000,0.000000,0.000000}%
\pgfsetstrokecolor{currentstroke}%
\pgfsetdash{}{0pt}%
\pgfpathmoveto{\pgfqpoint{1.363621in}{0.874036in}}%
\pgfpathlineto{\pgfqpoint{1.381371in}{0.885671in}}%
\pgfpathlineto{\pgfqpoint{1.399201in}{0.897504in}}%
\pgfpathlineto{\pgfqpoint{1.407190in}{0.898918in}}%
\pgfpathlineto{\pgfqpoint{1.415156in}{0.900695in}}%
\pgfpathlineto{\pgfqpoint{1.397352in}{0.888622in}}%
\pgfpathlineto{\pgfqpoint{1.379628in}{0.876730in}}%
\pgfpathlineto{\pgfqpoint{1.371637in}{0.875192in}}%
\pgfpathlineto{\pgfqpoint{1.363621in}{0.874036in}}%
\pgfpathclose%
\pgfusepath{fill}%
\end{pgfscope}%
\begin{pgfscope}%
\pgfpathrectangle{\pgfqpoint{0.407592in}{0.222904in}}{\pgfqpoint{1.411163in}{1.411163in}}%
\pgfusepath{clip}%
\pgfsetbuttcap%
\pgfsetroundjoin%
\definecolor{currentfill}{rgb}{0.556863,0.654902,0.847059}%
\pgfsetfillcolor{currentfill}%
\pgfsetfillopacity{0.130000}%
\pgfsetlinewidth{0.000000pt}%
\definecolor{currentstroke}{rgb}{0.000000,0.000000,0.000000}%
\pgfsetstrokecolor{currentstroke}%
\pgfsetdash{}{0pt}%
\pgfpathmoveto{\pgfqpoint{1.055511in}{0.804534in}}%
\pgfpathlineto{\pgfqpoint{1.072627in}{0.808373in}}%
\pgfpathlineto{\pgfqpoint{1.089802in}{0.812873in}}%
\pgfpathlineto{\pgfqpoint{1.098380in}{0.808029in}}%
\pgfpathlineto{\pgfqpoint{1.106920in}{0.803406in}}%
\pgfpathlineto{\pgfqpoint{1.089791in}{0.798439in}}%
\pgfpathlineto{\pgfqpoint{1.072722in}{0.794145in}}%
\pgfpathlineto{\pgfqpoint{1.064136in}{0.799238in}}%
\pgfpathlineto{\pgfqpoint{1.055511in}{0.804534in}}%
\pgfpathclose%
\pgfusepath{fill}%
\end{pgfscope}%
\begin{pgfscope}%
\pgfpathrectangle{\pgfqpoint{0.407592in}{0.222904in}}{\pgfqpoint{1.411163in}{1.411163in}}%
\pgfusepath{clip}%
\pgfsetbuttcap%
\pgfsetroundjoin%
\definecolor{currentfill}{rgb}{0.556863,0.654902,0.847059}%
\pgfsetfillcolor{currentfill}%
\pgfsetfillopacity{0.130000}%
\pgfsetlinewidth{0.000000pt}%
\definecolor{currentstroke}{rgb}{0.000000,0.000000,0.000000}%
\pgfsetstrokecolor{currentstroke}%
\pgfsetdash{}{0pt}%
\pgfpathmoveto{\pgfqpoint{1.586053in}{1.006306in}}%
\pgfpathlineto{\pgfqpoint{1.604333in}{1.019634in}}%
\pgfpathlineto{\pgfqpoint{1.622697in}{1.033061in}}%
\pgfpathlineto{\pgfqpoint{1.630384in}{1.038231in}}%
\pgfpathlineto{\pgfqpoint{1.638056in}{1.043703in}}%
\pgfpathlineto{\pgfqpoint{1.619717in}{1.030305in}}%
\pgfpathlineto{\pgfqpoint{1.601462in}{1.017006in}}%
\pgfpathlineto{\pgfqpoint{1.593765in}{1.011496in}}%
\pgfpathlineto{\pgfqpoint{1.586053in}{1.006306in}}%
\pgfpathclose%
\pgfusepath{fill}%
\end{pgfscope}%
\begin{pgfscope}%
\pgfpathrectangle{\pgfqpoint{0.407592in}{0.222904in}}{\pgfqpoint{1.411163in}{1.411163in}}%
\pgfusepath{clip}%
\pgfsetbuttcap%
\pgfsetroundjoin%
\definecolor{currentfill}{rgb}{0.556863,0.654902,0.847059}%
\pgfsetfillcolor{currentfill}%
\pgfsetfillopacity{0.130000}%
\pgfsetlinewidth{0.000000pt}%
\definecolor{currentstroke}{rgb}{0.000000,0.000000,0.000000}%
\pgfsetstrokecolor{currentstroke}%
\pgfsetdash{}{0pt}%
\pgfpathmoveto{\pgfqpoint{1.415156in}{0.900695in}}%
\pgfpathlineto{\pgfqpoint{1.433040in}{0.912935in}}%
\pgfpathlineto{\pgfqpoint{1.451006in}{0.925334in}}%
\pgfpathlineto{\pgfqpoint{1.458922in}{0.927665in}}%
\pgfpathlineto{\pgfqpoint{1.466817in}{0.930348in}}%
\pgfpathlineto{\pgfqpoint{1.448876in}{0.917782in}}%
\pgfpathlineto{\pgfqpoint{1.431018in}{0.905360in}}%
\pgfpathlineto{\pgfqpoint{1.423098in}{0.902840in}}%
\pgfpathlineto{\pgfqpoint{1.415156in}{0.900695in}}%
\pgfpathclose%
\pgfusepath{fill}%
\end{pgfscope}%
\begin{pgfscope}%
\pgfpathrectangle{\pgfqpoint{0.407592in}{0.222904in}}{\pgfqpoint{1.411163in}{1.411163in}}%
\pgfusepath{clip}%
\pgfsetbuttcap%
\pgfsetroundjoin%
\definecolor{currentfill}{rgb}{0.556863,0.654902,0.847059}%
\pgfsetfillcolor{currentfill}%
\pgfsetfillopacity{0.130000}%
\pgfsetlinewidth{0.000000pt}%
\definecolor{currentstroke}{rgb}{0.000000,0.000000,0.000000}%
\pgfsetstrokecolor{currentstroke}%
\pgfsetdash{}{0pt}%
\pgfpathmoveto{\pgfqpoint{1.638056in}{1.043703in}}%
\pgfpathlineto{\pgfqpoint{1.656479in}{1.057196in}}%
\pgfpathlineto{\pgfqpoint{1.664122in}{1.062941in}}%
\pgfpathlineto{\pgfqpoint{1.671749in}{1.068975in}}%
\pgfpathlineto{\pgfqpoint{1.653351in}{1.055544in}}%
\pgfpathlineto{\pgfqpoint{1.645711in}{1.049475in}}%
\pgfusepath{fill}%
\end{pgfscope}%
\begin{pgfscope}%
\pgfpathrectangle{\pgfqpoint{0.407592in}{0.222904in}}{\pgfqpoint{1.411163in}{1.411163in}}%
\pgfusepath{clip}%
\pgfsetbuttcap%
\pgfsetroundjoin%
\definecolor{currentfill}{rgb}{0.556863,0.654902,0.847059}%
\pgfsetfillcolor{currentfill}%
\pgfsetfillopacity{0.130000}%
\pgfsetlinewidth{0.000000pt}%
\definecolor{currentstroke}{rgb}{0.000000,0.000000,0.000000}%
\pgfsetstrokecolor{currentstroke}%
\pgfsetdash{}{0pt}%
\pgfpathmoveto{\pgfqpoint{1.003975in}{0.810916in}}%
\pgfpathlineto{\pgfqpoint{1.021031in}{0.812942in}}%
\pgfpathlineto{\pgfqpoint{1.038141in}{0.815682in}}%
\pgfpathlineto{\pgfqpoint{1.046846in}{0.810019in}}%
\pgfpathlineto{\pgfqpoint{1.055511in}{0.804534in}}%
\pgfpathlineto{\pgfqpoint{1.038451in}{0.801395in}}%
\pgfpathlineto{\pgfqpoint{1.021445in}{0.798995in}}%
\pgfpathlineto{\pgfqpoint{1.012730in}{0.804876in}}%
\pgfpathlineto{\pgfqpoint{1.003975in}{0.810916in}}%
\pgfpathclose%
\pgfusepath{fill}%
\end{pgfscope}%
\begin{pgfscope}%
\pgfpathrectangle{\pgfqpoint{0.407592in}{0.222904in}}{\pgfqpoint{1.411163in}{1.411163in}}%
\pgfusepath{clip}%
\pgfsetbuttcap%
\pgfsetroundjoin%
\definecolor{currentfill}{rgb}{0.556863,0.654902,0.847059}%
\pgfsetfillcolor{currentfill}%
\pgfsetfillopacity{0.130000}%
\pgfsetlinewidth{0.000000pt}%
\definecolor{currentstroke}{rgb}{0.000000,0.000000,0.000000}%
\pgfsetstrokecolor{currentstroke}%
\pgfsetdash{}{0pt}%
\pgfpathmoveto{\pgfqpoint{1.466817in}{0.930348in}}%
\pgfpathlineto{\pgfqpoint{1.484839in}{0.943052in}}%
\pgfpathlineto{\pgfqpoint{1.502944in}{0.955886in}}%
\pgfpathlineto{\pgfqpoint{1.510792in}{0.959046in}}%
\pgfpathlineto{\pgfqpoint{1.518620in}{0.962544in}}%
\pgfpathlineto{\pgfqpoint{1.500540in}{0.949606in}}%
\pgfpathlineto{\pgfqpoint{1.482543in}{0.936789in}}%
\pgfpathlineto{\pgfqpoint{1.474690in}{0.933389in}}%
\pgfpathlineto{\pgfqpoint{1.466817in}{0.930348in}}%
\pgfpathclose%
\pgfusepath{fill}%
\end{pgfscope}%
\begin{pgfscope}%
\pgfpathrectangle{\pgfqpoint{0.407592in}{0.222904in}}{\pgfqpoint{1.411163in}{1.411163in}}%
\pgfusepath{clip}%
\pgfsetbuttcap%
\pgfsetroundjoin%
\definecolor{currentfill}{rgb}{0.556863,0.654902,0.847059}%
\pgfsetfillcolor{currentfill}%
\pgfsetfillopacity{0.130000}%
\pgfsetlinewidth{0.000000pt}%
\definecolor{currentstroke}{rgb}{0.000000,0.000000,0.000000}%
\pgfsetstrokecolor{currentstroke}%
\pgfsetdash{}{0pt}%
\pgfpathmoveto{\pgfqpoint{0.952269in}{0.822230in}}%
\pgfpathlineto{\pgfqpoint{0.969281in}{0.822460in}}%
\pgfpathlineto{\pgfqpoint{0.986340in}{0.823431in}}%
\pgfpathlineto{\pgfqpoint{0.995178in}{0.817104in}}%
\pgfpathlineto{\pgfqpoint{1.003975in}{0.810916in}}%
\pgfpathlineto{\pgfqpoint{0.986968in}{0.809636in}}%
\pgfpathlineto{\pgfqpoint{0.970010in}{0.809134in}}%
\pgfpathlineto{\pgfqpoint{0.961161in}{0.815620in}}%
\pgfpathlineto{\pgfqpoint{0.952269in}{0.822230in}}%
\pgfpathclose%
\pgfusepath{fill}%
\end{pgfscope}%
\begin{pgfscope}%
\pgfpathrectangle{\pgfqpoint{0.407592in}{0.222904in}}{\pgfqpoint{1.411163in}{1.411163in}}%
\pgfusepath{clip}%
\pgfsetbuttcap%
\pgfsetroundjoin%
\definecolor{currentfill}{rgb}{0.556863,0.654902,0.847059}%
\pgfsetfillcolor{currentfill}%
\pgfsetfillopacity{0.130000}%
\pgfsetlinewidth{0.000000pt}%
\definecolor{currentstroke}{rgb}{0.000000,0.000000,0.000000}%
\pgfsetstrokecolor{currentstroke}%
\pgfsetdash{}{0pt}%
\pgfpathmoveto{\pgfqpoint{1.518620in}{0.962544in}}%
\pgfpathlineto{\pgfqpoint{1.536783in}{0.975599in}}%
\pgfpathlineto{\pgfqpoint{1.555029in}{0.988766in}}%
\pgfpathlineto{\pgfqpoint{1.562812in}{0.992667in}}%
\pgfpathlineto{\pgfqpoint{1.570577in}{0.996891in}}%
\pgfpathlineto{\pgfqpoint{1.552356in}{0.983674in}}%
\pgfpathlineto{\pgfqpoint{1.534218in}{0.970564in}}%
\pgfpathlineto{\pgfqpoint{1.526428in}{0.966383in}}%
\pgfpathlineto{\pgfqpoint{1.518620in}{0.962544in}}%
\pgfpathclose%
\pgfusepath{fill}%
\end{pgfscope}%
\begin{pgfscope}%
\pgfpathrectangle{\pgfqpoint{0.407592in}{0.222904in}}{\pgfqpoint{1.411163in}{1.411163in}}%
\pgfusepath{clip}%
\pgfsetbuttcap%
\pgfsetroundjoin%
\definecolor{currentfill}{rgb}{0.556863,0.654902,0.847059}%
\pgfsetfillcolor{currentfill}%
\pgfsetfillopacity{0.130000}%
\pgfsetlinewidth{0.000000pt}%
\definecolor{currentstroke}{rgb}{0.000000,0.000000,0.000000}%
\pgfsetstrokecolor{currentstroke}%
\pgfsetdash{}{0pt}%
\pgfpathmoveto{\pgfqpoint{1.192878in}{0.822594in}}%
\pgfpathlineto{\pgfqpoint{1.210299in}{0.830765in}}%
\pgfpathlineto{\pgfqpoint{1.227792in}{0.839360in}}%
\pgfpathlineto{\pgfqpoint{1.236103in}{0.836987in}}%
\pgfpathlineto{\pgfqpoint{1.244381in}{0.834920in}}%
\pgfpathlineto{\pgfqpoint{1.226924in}{0.825864in}}%
\pgfpathlineto{\pgfqpoint{1.209540in}{0.817214in}}%
\pgfpathlineto{\pgfqpoint{1.201226in}{0.819751in}}%
\pgfpathlineto{\pgfqpoint{1.192878in}{0.822594in}}%
\pgfpathclose%
\pgfusepath{fill}%
\end{pgfscope}%
\begin{pgfscope}%
\pgfpathrectangle{\pgfqpoint{0.407592in}{0.222904in}}{\pgfqpoint{1.411163in}{1.411163in}}%
\pgfusepath{clip}%
\pgfsetbuttcap%
\pgfsetroundjoin%
\definecolor{currentfill}{rgb}{0.556863,0.654902,0.847059}%
\pgfsetfillcolor{currentfill}%
\pgfsetfillopacity{0.130000}%
\pgfsetlinewidth{0.000000pt}%
\definecolor{currentstroke}{rgb}{0.000000,0.000000,0.000000}%
\pgfsetstrokecolor{currentstroke}%
\pgfsetdash{}{0pt}%
\pgfpathmoveto{\pgfqpoint{1.244381in}{0.834920in}}%
\pgfpathlineto{\pgfqpoint{1.261912in}{0.844349in}}%
\pgfpathlineto{\pgfqpoint{1.279519in}{0.854121in}}%
\pgfpathlineto{\pgfqpoint{1.287732in}{0.852803in}}%
\pgfpathlineto{\pgfqpoint{1.295914in}{0.851815in}}%
\pgfpathlineto{\pgfqpoint{1.278340in}{0.841630in}}%
\pgfpathlineto{\pgfqpoint{1.260842in}{0.831767in}}%
\pgfpathlineto{\pgfqpoint{1.252627in}{0.833175in}}%
\pgfpathlineto{\pgfqpoint{1.244381in}{0.834920in}}%
\pgfpathclose%
\pgfusepath{fill}%
\end{pgfscope}%
\begin{pgfscope}%
\pgfpathrectangle{\pgfqpoint{0.407592in}{0.222904in}}{\pgfqpoint{1.411163in}{1.411163in}}%
\pgfusepath{clip}%
\pgfsetbuttcap%
\pgfsetroundjoin%
\definecolor{currentfill}{rgb}{0.556863,0.654902,0.847059}%
\pgfsetfillcolor{currentfill}%
\pgfsetfillopacity{0.130000}%
\pgfsetlinewidth{0.000000pt}%
\definecolor{currentstroke}{rgb}{0.000000,0.000000,0.000000}%
\pgfsetstrokecolor{currentstroke}%
\pgfsetdash{}{0pt}%
\pgfpathmoveto{\pgfqpoint{1.141367in}{0.815180in}}%
\pgfpathlineto{\pgfqpoint{1.158688in}{0.821904in}}%
\pgfpathlineto{\pgfqpoint{1.176079in}{0.829133in}}%
\pgfpathlineto{\pgfqpoint{1.184496in}{0.825726in}}%
\pgfpathlineto{\pgfqpoint{1.192878in}{0.822594in}}%
\pgfpathlineto{\pgfqpoint{1.175528in}{0.814880in}}%
\pgfpathlineto{\pgfqpoint{1.158247in}{0.807663in}}%
\pgfpathlineto{\pgfqpoint{1.149825in}{0.811288in}}%
\pgfpathlineto{\pgfqpoint{1.141367in}{0.815180in}}%
\pgfpathclose%
\pgfusepath{fill}%
\end{pgfscope}%
\begin{pgfscope}%
\pgfpathrectangle{\pgfqpoint{0.407592in}{0.222904in}}{\pgfqpoint{1.411163in}{1.411163in}}%
\pgfusepath{clip}%
\pgfsetbuttcap%
\pgfsetroundjoin%
\definecolor{currentfill}{rgb}{0.556863,0.654902,0.847059}%
\pgfsetfillcolor{currentfill}%
\pgfsetfillopacity{0.130000}%
\pgfsetlinewidth{0.000000pt}%
\definecolor{currentstroke}{rgb}{0.000000,0.000000,0.000000}%
\pgfsetstrokecolor{currentstroke}%
\pgfsetdash{}{0pt}%
\pgfpathmoveto{\pgfqpoint{1.295914in}{0.851815in}}%
\pgfpathlineto{\pgfqpoint{1.313565in}{0.862297in}}%
\pgfpathlineto{\pgfqpoint{1.331294in}{0.873053in}}%
\pgfpathlineto{\pgfqpoint{1.339417in}{0.872774in}}%
\pgfpathlineto{\pgfqpoint{1.347511in}{0.872838in}}%
\pgfpathlineto{\pgfqpoint{1.329812in}{0.861734in}}%
\pgfpathlineto{\pgfqpoint{1.312192in}{0.850882in}}%
\pgfpathlineto{\pgfqpoint{1.304067in}{0.851171in}}%
\pgfpathlineto{\pgfqpoint{1.295914in}{0.851815in}}%
\pgfpathclose%
\pgfusepath{fill}%
\end{pgfscope}%
\begin{pgfscope}%
\pgfpathrectangle{\pgfqpoint{0.407592in}{0.222904in}}{\pgfqpoint{1.411163in}{1.411163in}}%
\pgfusepath{clip}%
\pgfsetbuttcap%
\pgfsetroundjoin%
\definecolor{currentfill}{rgb}{0.556863,0.654902,0.847059}%
\pgfsetfillcolor{currentfill}%
\pgfsetfillopacity{0.130000}%
\pgfsetlinewidth{0.000000pt}%
\definecolor{currentstroke}{rgb}{0.000000,0.000000,0.000000}%
\pgfsetstrokecolor{currentstroke}%
\pgfsetdash{}{0pt}%
\pgfpathmoveto{\pgfqpoint{1.347511in}{0.872838in}}%
\pgfpathlineto{\pgfqpoint{1.365290in}{0.884175in}}%
\pgfpathlineto{\pgfqpoint{1.383148in}{0.895730in}}%
\pgfpathlineto{\pgfqpoint{1.391187in}{0.896443in}}%
\pgfpathlineto{\pgfqpoint{1.399201in}{0.897504in}}%
\pgfpathlineto{\pgfqpoint{1.381371in}{0.885671in}}%
\pgfpathlineto{\pgfqpoint{1.363621in}{0.874036in}}%
\pgfpathlineto{\pgfqpoint{1.355579in}{0.873255in}}%
\pgfpathlineto{\pgfqpoint{1.347511in}{0.872838in}}%
\pgfpathclose%
\pgfusepath{fill}%
\end{pgfscope}%
\begin{pgfscope}%
\pgfpathrectangle{\pgfqpoint{0.407592in}{0.222904in}}{\pgfqpoint{1.411163in}{1.411163in}}%
\pgfusepath{clip}%
\pgfsetbuttcap%
\pgfsetroundjoin%
\definecolor{currentfill}{rgb}{0.556863,0.654902,0.847059}%
\pgfsetfillcolor{currentfill}%
\pgfsetfillopacity{0.130000}%
\pgfsetlinewidth{0.000000pt}%
\definecolor{currentstroke}{rgb}{0.000000,0.000000,0.000000}%
\pgfsetstrokecolor{currentstroke}%
\pgfsetdash{}{0pt}%
\pgfpathmoveto{\pgfqpoint{1.089802in}{0.812873in}}%
\pgfpathlineto{\pgfqpoint{1.107039in}{0.817994in}}%
\pgfpathlineto{\pgfqpoint{1.124339in}{0.823699in}}%
\pgfpathlineto{\pgfqpoint{1.132871in}{0.819322in}}%
\pgfpathlineto{\pgfqpoint{1.141367in}{0.815180in}}%
\pgfpathlineto{\pgfqpoint{1.124111in}{0.809000in}}%
\pgfpathlineto{\pgfqpoint{1.106920in}{0.803406in}}%
\pgfpathlineto{\pgfqpoint{1.098380in}{0.808029in}}%
\pgfpathlineto{\pgfqpoint{1.089802in}{0.812873in}}%
\pgfpathclose%
\pgfusepath{fill}%
\end{pgfscope}%
\begin{pgfscope}%
\pgfpathrectangle{\pgfqpoint{0.407592in}{0.222904in}}{\pgfqpoint{1.411163in}{1.411163in}}%
\pgfusepath{clip}%
\pgfsetbuttcap%
\pgfsetroundjoin%
\definecolor{currentfill}{rgb}{0.556863,0.654902,0.847059}%
\pgfsetfillcolor{currentfill}%
\pgfsetfillopacity{0.130000}%
\pgfsetlinewidth{0.000000pt}%
\definecolor{currentstroke}{rgb}{0.000000,0.000000,0.000000}%
\pgfsetstrokecolor{currentstroke}%
\pgfsetdash{}{0pt}%
\pgfpathmoveto{\pgfqpoint{1.570577in}{0.996891in}}%
\pgfpathlineto{\pgfqpoint{1.588882in}{1.010212in}}%
\pgfpathlineto{\pgfqpoint{1.607272in}{1.023632in}}%
\pgfpathlineto{\pgfqpoint{1.614993in}{1.028195in}}%
\pgfpathlineto{\pgfqpoint{1.622697in}{1.033061in}}%
\pgfpathlineto{\pgfqpoint{1.604333in}{1.019634in}}%
\pgfpathlineto{\pgfqpoint{1.586053in}{1.006306in}}%
\pgfpathlineto{\pgfqpoint{1.578324in}{1.001437in}}%
\pgfpathlineto{\pgfqpoint{1.570577in}{0.996891in}}%
\pgfpathclose%
\pgfusepath{fill}%
\end{pgfscope}%
\begin{pgfscope}%
\pgfpathrectangle{\pgfqpoint{0.407592in}{0.222904in}}{\pgfqpoint{1.411163in}{1.411163in}}%
\pgfusepath{clip}%
\pgfsetbuttcap%
\pgfsetroundjoin%
\definecolor{currentfill}{rgb}{0.556863,0.654902,0.847059}%
\pgfsetfillcolor{currentfill}%
\pgfsetfillopacity{0.130000}%
\pgfsetlinewidth{0.000000pt}%
\definecolor{currentstroke}{rgb}{0.000000,0.000000,0.000000}%
\pgfsetstrokecolor{currentstroke}%
\pgfsetdash{}{0pt}%
\pgfpathmoveto{\pgfqpoint{1.622697in}{1.033061in}}%
\pgfpathlineto{\pgfqpoint{1.641146in}{1.046584in}}%
\pgfpathlineto{\pgfqpoint{1.648821in}{1.051743in}}%
\pgfpathlineto{\pgfqpoint{1.656479in}{1.057196in}}%
\pgfpathlineto{\pgfqpoint{1.638056in}{1.043703in}}%
\pgfpathlineto{\pgfqpoint{1.630384in}{1.038231in}}%
\pgfusepath{fill}%
\end{pgfscope}%
\begin{pgfscope}%
\pgfpathrectangle{\pgfqpoint{0.407592in}{0.222904in}}{\pgfqpoint{1.411163in}{1.411163in}}%
\pgfusepath{clip}%
\pgfsetbuttcap%
\pgfsetroundjoin%
\definecolor{currentfill}{rgb}{0.556863,0.654902,0.847059}%
\pgfsetfillcolor{currentfill}%
\pgfsetfillopacity{0.130000}%
\pgfsetlinewidth{0.000000pt}%
\definecolor{currentstroke}{rgb}{0.000000,0.000000,0.000000}%
\pgfsetstrokecolor{currentstroke}%
\pgfsetdash{}{0pt}%
\pgfpathmoveto{\pgfqpoint{1.399201in}{0.897504in}}%
\pgfpathlineto{\pgfqpoint{1.417112in}{0.909521in}}%
\pgfpathlineto{\pgfqpoint{1.435105in}{0.921711in}}%
\pgfpathlineto{\pgfqpoint{1.443067in}{0.923352in}}%
\pgfpathlineto{\pgfqpoint{1.451006in}{0.925334in}}%
\pgfpathlineto{\pgfqpoint{1.433040in}{0.912935in}}%
\pgfpathlineto{\pgfqpoint{1.415156in}{0.900695in}}%
\pgfpathlineto{\pgfqpoint{1.407190in}{0.898918in}}%
\pgfpathlineto{\pgfqpoint{1.399201in}{0.897504in}}%
\pgfpathclose%
\pgfusepath{fill}%
\end{pgfscope}%
\begin{pgfscope}%
\pgfpathrectangle{\pgfqpoint{0.407592in}{0.222904in}}{\pgfqpoint{1.411163in}{1.411163in}}%
\pgfusepath{clip}%
\pgfsetbuttcap%
\pgfsetroundjoin%
\definecolor{currentfill}{rgb}{0.556863,0.654902,0.847059}%
\pgfsetfillcolor{currentfill}%
\pgfsetfillopacity{0.130000}%
\pgfsetlinewidth{0.000000pt}%
\definecolor{currentstroke}{rgb}{0.000000,0.000000,0.000000}%
\pgfsetstrokecolor{currentstroke}%
\pgfsetdash{}{0pt}%
\pgfpathmoveto{\pgfqpoint{1.038141in}{0.815682in}}%
\pgfpathlineto{\pgfqpoint{1.055307in}{0.819102in}}%
\pgfpathlineto{\pgfqpoint{1.072530in}{0.823169in}}%
\pgfpathlineto{\pgfqpoint{1.081186in}{0.817924in}}%
\pgfpathlineto{\pgfqpoint{1.089802in}{0.812873in}}%
\pgfpathlineto{\pgfqpoint{1.072627in}{0.808373in}}%
\pgfpathlineto{\pgfqpoint{1.055511in}{0.804534in}}%
\pgfpathlineto{\pgfqpoint{1.046846in}{0.810019in}}%
\pgfpathlineto{\pgfqpoint{1.038141in}{0.815682in}}%
\pgfpathclose%
\pgfusepath{fill}%
\end{pgfscope}%
\begin{pgfscope}%
\pgfpathrectangle{\pgfqpoint{0.407592in}{0.222904in}}{\pgfqpoint{1.411163in}{1.411163in}}%
\pgfusepath{clip}%
\pgfsetbuttcap%
\pgfsetroundjoin%
\definecolor{currentfill}{rgb}{0.556863,0.654902,0.847059}%
\pgfsetfillcolor{currentfill}%
\pgfsetfillopacity{0.130000}%
\pgfsetlinewidth{0.000000pt}%
\definecolor{currentstroke}{rgb}{0.000000,0.000000,0.000000}%
\pgfsetstrokecolor{currentstroke}%
\pgfsetdash{}{0pt}%
\pgfpathmoveto{\pgfqpoint{1.451006in}{0.925334in}}%
\pgfpathlineto{\pgfqpoint{1.469054in}{0.937883in}}%
\pgfpathlineto{\pgfqpoint{1.487185in}{0.950572in}}%
\pgfpathlineto{\pgfqpoint{1.495075in}{0.953063in}}%
\pgfpathlineto{\pgfqpoint{1.502944in}{0.955886in}}%
\pgfpathlineto{\pgfqpoint{1.484839in}{0.943052in}}%
\pgfpathlineto{\pgfqpoint{1.466817in}{0.930348in}}%
\pgfpathlineto{\pgfqpoint{1.458922in}{0.927665in}}%
\pgfpathlineto{\pgfqpoint{1.451006in}{0.925334in}}%
\pgfpathclose%
\pgfusepath{fill}%
\end{pgfscope}%
\begin{pgfscope}%
\pgfpathrectangle{\pgfqpoint{0.407592in}{0.222904in}}{\pgfqpoint{1.411163in}{1.411163in}}%
\pgfusepath{clip}%
\pgfsetbuttcap%
\pgfsetroundjoin%
\definecolor{currentfill}{rgb}{0.556863,0.654902,0.847059}%
\pgfsetfillcolor{currentfill}%
\pgfsetfillopacity{0.130000}%
\pgfsetlinewidth{0.000000pt}%
\definecolor{currentstroke}{rgb}{0.000000,0.000000,0.000000}%
\pgfsetstrokecolor{currentstroke}%
\pgfsetdash{}{0pt}%
\pgfpathmoveto{\pgfqpoint{0.986340in}{0.823431in}}%
\pgfpathlineto{\pgfqpoint{1.003449in}{0.825118in}}%
\pgfpathlineto{\pgfqpoint{1.020610in}{0.827494in}}%
\pgfpathlineto{\pgfqpoint{1.029396in}{0.821510in}}%
\pgfpathlineto{\pgfqpoint{1.038141in}{0.815682in}}%
\pgfpathlineto{\pgfqpoint{1.021031in}{0.812942in}}%
\pgfpathlineto{\pgfqpoint{1.003975in}{0.810916in}}%
\pgfpathlineto{\pgfqpoint{0.995178in}{0.817104in}}%
\pgfpathlineto{\pgfqpoint{0.986340in}{0.823431in}}%
\pgfpathclose%
\pgfusepath{fill}%
\end{pgfscope}%
\begin{pgfscope}%
\pgfpathrectangle{\pgfqpoint{0.407592in}{0.222904in}}{\pgfqpoint{1.411163in}{1.411163in}}%
\pgfusepath{clip}%
\pgfsetbuttcap%
\pgfsetroundjoin%
\definecolor{currentfill}{rgb}{0.556863,0.654902,0.847059}%
\pgfsetfillcolor{currentfill}%
\pgfsetfillopacity{0.130000}%
\pgfsetlinewidth{0.000000pt}%
\definecolor{currentstroke}{rgb}{0.000000,0.000000,0.000000}%
\pgfsetstrokecolor{currentstroke}%
\pgfsetdash{}{0pt}%
\pgfpathmoveto{\pgfqpoint{1.502944in}{0.955886in}}%
\pgfpathlineto{\pgfqpoint{1.521132in}{0.968845in}}%
\pgfpathlineto{\pgfqpoint{1.539404in}{0.981923in}}%
\pgfpathlineto{\pgfqpoint{1.547227in}{0.985185in}}%
\pgfpathlineto{\pgfqpoint{1.555029in}{0.988766in}}%
\pgfpathlineto{\pgfqpoint{1.536783in}{0.975599in}}%
\pgfpathlineto{\pgfqpoint{1.518620in}{0.962544in}}%
\pgfpathlineto{\pgfqpoint{1.510792in}{0.959046in}}%
\pgfpathlineto{\pgfqpoint{1.502944in}{0.955886in}}%
\pgfpathclose%
\pgfusepath{fill}%
\end{pgfscope}%
\begin{pgfscope}%
\pgfpathrectangle{\pgfqpoint{0.407592in}{0.222904in}}{\pgfqpoint{1.411163in}{1.411163in}}%
\pgfusepath{clip}%
\pgfsetbuttcap%
\pgfsetroundjoin%
\definecolor{currentfill}{rgb}{0.556863,0.654902,0.847059}%
\pgfsetfillcolor{currentfill}%
\pgfsetfillopacity{0.130000}%
\pgfsetlinewidth{0.000000pt}%
\definecolor{currentstroke}{rgb}{0.000000,0.000000,0.000000}%
\pgfsetstrokecolor{currentstroke}%
\pgfsetdash{}{0pt}%
\pgfpathmoveto{\pgfqpoint{1.227792in}{0.839360in}}%
\pgfpathlineto{\pgfqpoint{1.245359in}{0.848346in}}%
\pgfpathlineto{\pgfqpoint{1.263000in}{0.857695in}}%
\pgfpathlineto{\pgfqpoint{1.271275in}{0.855756in}}%
\pgfpathlineto{\pgfqpoint{1.279519in}{0.854121in}}%
\pgfpathlineto{\pgfqpoint{1.261912in}{0.844349in}}%
\pgfpathlineto{\pgfqpoint{1.244381in}{0.834920in}}%
\pgfpathlineto{\pgfqpoint{1.236103in}{0.836987in}}%
\pgfpathlineto{\pgfqpoint{1.227792in}{0.839360in}}%
\pgfpathclose%
\pgfusepath{fill}%
\end{pgfscope}%
\begin{pgfscope}%
\pgfpathrectangle{\pgfqpoint{0.407592in}{0.222904in}}{\pgfqpoint{1.411163in}{1.411163in}}%
\pgfusepath{clip}%
\pgfsetbuttcap%
\pgfsetroundjoin%
\definecolor{currentfill}{rgb}{0.556863,0.654902,0.847059}%
\pgfsetfillcolor{currentfill}%
\pgfsetfillopacity{0.130000}%
\pgfsetlinewidth{0.000000pt}%
\definecolor{currentstroke}{rgb}{0.000000,0.000000,0.000000}%
\pgfsetstrokecolor{currentstroke}%
\pgfsetdash{}{0pt}%
\pgfpathmoveto{\pgfqpoint{1.176079in}{0.829133in}}%
\pgfpathlineto{\pgfqpoint{1.193538in}{0.836833in}}%
\pgfpathlineto{\pgfqpoint{1.211070in}{0.844971in}}%
\pgfpathlineto{\pgfqpoint{1.219448in}{0.842026in}}%
\pgfpathlineto{\pgfqpoint{1.227792in}{0.839360in}}%
\pgfpathlineto{\pgfqpoint{1.210299in}{0.830765in}}%
\pgfpathlineto{\pgfqpoint{1.192878in}{0.822594in}}%
\pgfpathlineto{\pgfqpoint{1.184496in}{0.825726in}}%
\pgfpathlineto{\pgfqpoint{1.176079in}{0.829133in}}%
\pgfpathclose%
\pgfusepath{fill}%
\end{pgfscope}%
\begin{pgfscope}%
\pgfpathrectangle{\pgfqpoint{0.407592in}{0.222904in}}{\pgfqpoint{1.411163in}{1.411163in}}%
\pgfusepath{clip}%
\pgfsetbuttcap%
\pgfsetroundjoin%
\definecolor{currentfill}{rgb}{0.556863,0.654902,0.847059}%
\pgfsetfillcolor{currentfill}%
\pgfsetfillopacity{0.130000}%
\pgfsetlinewidth{0.000000pt}%
\definecolor{currentstroke}{rgb}{0.000000,0.000000,0.000000}%
\pgfsetstrokecolor{currentstroke}%
\pgfsetdash{}{0pt}%
\pgfpathmoveto{\pgfqpoint{1.279519in}{0.854121in}}%
\pgfpathlineto{\pgfqpoint{1.297202in}{0.864211in}}%
\pgfpathlineto{\pgfqpoint{1.314963in}{0.874596in}}%
\pgfpathlineto{\pgfqpoint{1.323143in}{0.873664in}}%
\pgfpathlineto{\pgfqpoint{1.331294in}{0.873053in}}%
\pgfpathlineto{\pgfqpoint{1.313565in}{0.862297in}}%
\pgfpathlineto{\pgfqpoint{1.295914in}{0.851815in}}%
\pgfpathlineto{\pgfqpoint{1.287732in}{0.852803in}}%
\pgfpathlineto{\pgfqpoint{1.279519in}{0.854121in}}%
\pgfpathclose%
\pgfusepath{fill}%
\end{pgfscope}%
\begin{pgfscope}%
\pgfpathrectangle{\pgfqpoint{0.407592in}{0.222904in}}{\pgfqpoint{1.411163in}{1.411163in}}%
\pgfusepath{clip}%
\pgfsetbuttcap%
\pgfsetroundjoin%
\definecolor{currentfill}{rgb}{0.556863,0.654902,0.847059}%
\pgfsetfillcolor{currentfill}%
\pgfsetfillopacity{0.130000}%
\pgfsetlinewidth{0.000000pt}%
\definecolor{currentstroke}{rgb}{0.000000,0.000000,0.000000}%
\pgfsetstrokecolor{currentstroke}%
\pgfsetdash{}{0pt}%
\pgfpathmoveto{\pgfqpoint{1.555029in}{0.988766in}}%
\pgfpathlineto{\pgfqpoint{1.573360in}{1.002040in}}%
\pgfpathlineto{\pgfqpoint{1.591775in}{1.015419in}}%
\pgfpathlineto{\pgfqpoint{1.599532in}{1.019374in}}%
\pgfpathlineto{\pgfqpoint{1.607272in}{1.023632in}}%
\pgfpathlineto{\pgfqpoint{1.588882in}{1.010212in}}%
\pgfpathlineto{\pgfqpoint{1.570577in}{0.996891in}}%
\pgfpathlineto{\pgfqpoint{1.562812in}{0.992667in}}%
\pgfpathlineto{\pgfqpoint{1.555029in}{0.988766in}}%
\pgfpathclose%
\pgfusepath{fill}%
\end{pgfscope}%
\begin{pgfscope}%
\pgfpathrectangle{\pgfqpoint{0.407592in}{0.222904in}}{\pgfqpoint{1.411163in}{1.411163in}}%
\pgfusepath{clip}%
\pgfsetbuttcap%
\pgfsetroundjoin%
\definecolor{currentfill}{rgb}{0.556863,0.654902,0.847059}%
\pgfsetfillcolor{currentfill}%
\pgfsetfillopacity{0.130000}%
\pgfsetlinewidth{0.000000pt}%
\definecolor{currentstroke}{rgb}{0.000000,0.000000,0.000000}%
\pgfsetstrokecolor{currentstroke}%
\pgfsetdash{}{0pt}%
\pgfpathmoveto{\pgfqpoint{1.607272in}{1.023632in}}%
\pgfpathlineto{\pgfqpoint{1.625746in}{1.037151in}}%
\pgfpathlineto{\pgfqpoint{1.633455in}{1.041720in}}%
\pgfpathlineto{\pgfqpoint{1.641146in}{1.046584in}}%
\pgfpathlineto{\pgfqpoint{1.622697in}{1.033061in}}%
\pgfpathlineto{\pgfqpoint{1.614993in}{1.028195in}}%
\pgfusepath{fill}%
\end{pgfscope}%
\begin{pgfscope}%
\pgfpathrectangle{\pgfqpoint{0.407592in}{0.222904in}}{\pgfqpoint{1.411163in}{1.411163in}}%
\pgfusepath{clip}%
\pgfsetbuttcap%
\pgfsetroundjoin%
\definecolor{currentfill}{rgb}{0.556863,0.654902,0.847059}%
\pgfsetfillcolor{currentfill}%
\pgfsetfillopacity{0.130000}%
\pgfsetlinewidth{0.000000pt}%
\definecolor{currentstroke}{rgb}{0.000000,0.000000,0.000000}%
\pgfsetstrokecolor{currentstroke}%
\pgfsetdash{}{0pt}%
\pgfpathmoveto{\pgfqpoint{1.124339in}{0.823699in}}%
\pgfpathlineto{\pgfqpoint{1.141704in}{0.829950in}}%
\pgfpathlineto{\pgfqpoint{1.159136in}{0.836713in}}%
\pgfpathlineto{\pgfqpoint{1.167625in}{0.832800in}}%
\pgfpathlineto{\pgfqpoint{1.176079in}{0.829133in}}%
\pgfpathlineto{\pgfqpoint{1.158688in}{0.821904in}}%
\pgfpathlineto{\pgfqpoint{1.141367in}{0.815180in}}%
\pgfpathlineto{\pgfqpoint{1.132871in}{0.819322in}}%
\pgfpathlineto{\pgfqpoint{1.124339in}{0.823699in}}%
\pgfpathclose%
\pgfusepath{fill}%
\end{pgfscope}%
\begin{pgfscope}%
\pgfpathrectangle{\pgfqpoint{0.407592in}{0.222904in}}{\pgfqpoint{1.411163in}{1.411163in}}%
\pgfusepath{clip}%
\pgfsetbuttcap%
\pgfsetroundjoin%
\definecolor{currentfill}{rgb}{0.556863,0.654902,0.847059}%
\pgfsetfillcolor{currentfill}%
\pgfsetfillopacity{0.130000}%
\pgfsetlinewidth{0.000000pt}%
\definecolor{currentstroke}{rgb}{0.000000,0.000000,0.000000}%
\pgfsetstrokecolor{currentstroke}%
\pgfsetdash{}{0pt}%
\pgfpathmoveto{\pgfqpoint{1.331294in}{0.873053in}}%
\pgfpathlineto{\pgfqpoint{1.349102in}{0.884062in}}%
\pgfpathlineto{\pgfqpoint{1.366990in}{0.895308in}}%
\pgfpathlineto{\pgfqpoint{1.375083in}{0.895354in}}%
\pgfpathlineto{\pgfqpoint{1.383148in}{0.895730in}}%
\pgfpathlineto{\pgfqpoint{1.365290in}{0.884175in}}%
\pgfpathlineto{\pgfqpoint{1.347511in}{0.872838in}}%
\pgfpathlineto{\pgfqpoint{1.339417in}{0.872774in}}%
\pgfpathlineto{\pgfqpoint{1.331294in}{0.873053in}}%
\pgfpathclose%
\pgfusepath{fill}%
\end{pgfscope}%
\begin{pgfscope}%
\pgfpathrectangle{\pgfqpoint{0.407592in}{0.222904in}}{\pgfqpoint{1.411163in}{1.411163in}}%
\pgfusepath{clip}%
\pgfsetbuttcap%
\pgfsetroundjoin%
\definecolor{currentfill}{rgb}{0.556863,0.654902,0.847059}%
\pgfsetfillcolor{currentfill}%
\pgfsetfillopacity{0.130000}%
\pgfsetlinewidth{0.000000pt}%
\definecolor{currentstroke}{rgb}{0.000000,0.000000,0.000000}%
\pgfsetstrokecolor{currentstroke}%
\pgfsetdash{}{0pt}%
\pgfpathmoveto{\pgfqpoint{1.383148in}{0.895730in}}%
\pgfpathlineto{\pgfqpoint{1.401087in}{0.907486in}}%
\pgfpathlineto{\pgfqpoint{1.419108in}{0.919433in}}%
\pgfpathlineto{\pgfqpoint{1.427119in}{0.920407in}}%
\pgfpathlineto{\pgfqpoint{1.435105in}{0.921711in}}%
\pgfpathlineto{\pgfqpoint{1.417112in}{0.909521in}}%
\pgfpathlineto{\pgfqpoint{1.399201in}{0.897504in}}%
\pgfpathlineto{\pgfqpoint{1.391187in}{0.896443in}}%
\pgfpathlineto{\pgfqpoint{1.383148in}{0.895730in}}%
\pgfpathclose%
\pgfusepath{fill}%
\end{pgfscope}%
\begin{pgfscope}%
\pgfpathrectangle{\pgfqpoint{0.407592in}{0.222904in}}{\pgfqpoint{1.411163in}{1.411163in}}%
\pgfusepath{clip}%
\pgfsetbuttcap%
\pgfsetroundjoin%
\definecolor{currentfill}{rgb}{0.556863,0.654902,0.847059}%
\pgfsetfillcolor{currentfill}%
\pgfsetfillopacity{0.130000}%
\pgfsetlinewidth{0.000000pt}%
\definecolor{currentstroke}{rgb}{0.000000,0.000000,0.000000}%
\pgfsetstrokecolor{currentstroke}%
\pgfsetdash{}{0pt}%
\pgfpathmoveto{\pgfqpoint{1.072530in}{0.823169in}}%
\pgfpathlineto{\pgfqpoint{1.089813in}{0.827849in}}%
\pgfpathlineto{\pgfqpoint{1.107159in}{0.833106in}}%
\pgfpathlineto{\pgfqpoint{1.115768in}{0.828298in}}%
\pgfpathlineto{\pgfqpoint{1.124339in}{0.823699in}}%
\pgfpathlineto{\pgfqpoint{1.107039in}{0.817994in}}%
\pgfpathlineto{\pgfqpoint{1.089802in}{0.812873in}}%
\pgfpathlineto{\pgfqpoint{1.081186in}{0.817924in}}%
\pgfpathlineto{\pgfqpoint{1.072530in}{0.823169in}}%
\pgfpathclose%
\pgfusepath{fill}%
\end{pgfscope}%
\begin{pgfscope}%
\pgfpathrectangle{\pgfqpoint{0.407592in}{0.222904in}}{\pgfqpoint{1.411163in}{1.411163in}}%
\pgfusepath{clip}%
\pgfsetbuttcap%
\pgfsetroundjoin%
\definecolor{currentfill}{rgb}{0.556863,0.654902,0.847059}%
\pgfsetfillcolor{currentfill}%
\pgfsetfillopacity{0.130000}%
\pgfsetlinewidth{0.000000pt}%
\definecolor{currentstroke}{rgb}{0.000000,0.000000,0.000000}%
\pgfsetstrokecolor{currentstroke}%
\pgfsetdash{}{0pt}%
\pgfpathmoveto{\pgfqpoint{1.435105in}{0.921711in}}%
\pgfpathlineto{\pgfqpoint{1.453180in}{0.934065in}}%
\pgfpathlineto{\pgfqpoint{1.471338in}{0.946572in}}%
\pgfpathlineto{\pgfqpoint{1.479273in}{0.948410in}}%
\pgfpathlineto{\pgfqpoint{1.487185in}{0.950572in}}%
\pgfpathlineto{\pgfqpoint{1.469054in}{0.937883in}}%
\pgfpathlineto{\pgfqpoint{1.451006in}{0.925334in}}%
\pgfpathlineto{\pgfqpoint{1.443067in}{0.923352in}}%
\pgfpathlineto{\pgfqpoint{1.435105in}{0.921711in}}%
\pgfpathclose%
\pgfusepath{fill}%
\end{pgfscope}%
\begin{pgfscope}%
\pgfpathrectangle{\pgfqpoint{0.407592in}{0.222904in}}{\pgfqpoint{1.411163in}{1.411163in}}%
\pgfusepath{clip}%
\pgfsetbuttcap%
\pgfsetroundjoin%
\definecolor{currentfill}{rgb}{0.556863,0.654902,0.847059}%
\pgfsetfillcolor{currentfill}%
\pgfsetfillopacity{0.130000}%
\pgfsetlinewidth{0.000000pt}%
\definecolor{currentstroke}{rgb}{0.000000,0.000000,0.000000}%
\pgfsetstrokecolor{currentstroke}%
\pgfsetdash{}{0pt}%
\pgfpathmoveto{\pgfqpoint{1.020610in}{0.827494in}}%
\pgfpathlineto{\pgfqpoint{1.037826in}{0.830530in}}%
\pgfpathlineto{\pgfqpoint{1.055099in}{0.834198in}}%
\pgfpathlineto{\pgfqpoint{1.063834in}{0.828598in}}%
\pgfpathlineto{\pgfqpoint{1.072530in}{0.823169in}}%
\pgfpathlineto{\pgfqpoint{1.055307in}{0.819102in}}%
\pgfpathlineto{\pgfqpoint{1.038141in}{0.815682in}}%
\pgfpathlineto{\pgfqpoint{1.029396in}{0.821510in}}%
\pgfpathlineto{\pgfqpoint{1.020610in}{0.827494in}}%
\pgfpathclose%
\pgfusepath{fill}%
\end{pgfscope}%
\begin{pgfscope}%
\pgfpathrectangle{\pgfqpoint{0.407592in}{0.222904in}}{\pgfqpoint{1.411163in}{1.411163in}}%
\pgfusepath{clip}%
\pgfsetbuttcap%
\pgfsetroundjoin%
\definecolor{currentfill}{rgb}{0.556863,0.654902,0.847059}%
\pgfsetfillcolor{currentfill}%
\pgfsetfillopacity{0.130000}%
\pgfsetlinewidth{0.000000pt}%
\definecolor{currentstroke}{rgb}{0.000000,0.000000,0.000000}%
\pgfsetstrokecolor{currentstroke}%
\pgfsetdash{}{0pt}%
\pgfpathmoveto{\pgfqpoint{1.487185in}{0.950572in}}%
\pgfpathlineto{\pgfqpoint{1.505400in}{0.963396in}}%
\pgfpathlineto{\pgfqpoint{1.523698in}{0.976349in}}%
\pgfpathlineto{\pgfqpoint{1.531562in}{0.978978in}}%
\pgfpathlineto{\pgfqpoint{1.539404in}{0.981923in}}%
\pgfpathlineto{\pgfqpoint{1.521132in}{0.968845in}}%
\pgfpathlineto{\pgfqpoint{1.502944in}{0.955886in}}%
\pgfpathlineto{\pgfqpoint{1.495075in}{0.953063in}}%
\pgfpathlineto{\pgfqpoint{1.487185in}{0.950572in}}%
\pgfpathclose%
\pgfusepath{fill}%
\end{pgfscope}%
\begin{pgfscope}%
\pgfpathrectangle{\pgfqpoint{0.407592in}{0.222904in}}{\pgfqpoint{1.411163in}{1.411163in}}%
\pgfusepath{clip}%
\pgfsetbuttcap%
\pgfsetroundjoin%
\definecolor{currentfill}{rgb}{0.556863,0.654902,0.847059}%
\pgfsetfillcolor{currentfill}%
\pgfsetfillopacity{0.130000}%
\pgfsetlinewidth{0.000000pt}%
\definecolor{currentstroke}{rgb}{0.000000,0.000000,0.000000}%
\pgfsetstrokecolor{currentstroke}%
\pgfsetdash{}{0pt}%
\pgfpathmoveto{\pgfqpoint{1.591775in}{1.015419in}}%
\pgfpathlineto{\pgfqpoint{1.610276in}{1.028900in}}%
\pgfpathlineto{\pgfqpoint{1.618020in}{1.032878in}}%
\pgfpathlineto{\pgfqpoint{1.625746in}{1.037151in}}%
\pgfpathlineto{\pgfqpoint{1.607272in}{1.023632in}}%
\pgfpathlineto{\pgfqpoint{1.599532in}{1.019374in}}%
\pgfusepath{fill}%
\end{pgfscope}%
\begin{pgfscope}%
\pgfpathrectangle{\pgfqpoint{0.407592in}{0.222904in}}{\pgfqpoint{1.411163in}{1.411163in}}%
\pgfusepath{clip}%
\pgfsetbuttcap%
\pgfsetroundjoin%
\definecolor{currentfill}{rgb}{0.556863,0.654902,0.847059}%
\pgfsetfillcolor{currentfill}%
\pgfsetfillopacity{0.130000}%
\pgfsetlinewidth{0.000000pt}%
\definecolor{currentstroke}{rgb}{0.000000,0.000000,0.000000}%
\pgfsetstrokecolor{currentstroke}%
\pgfsetdash{}{0pt}%
\pgfpathmoveto{\pgfqpoint{1.539404in}{0.981923in}}%
\pgfpathlineto{\pgfqpoint{1.557761in}{0.995115in}}%
\pgfpathlineto{\pgfqpoint{1.576203in}{1.008419in}}%
\pgfpathlineto{\pgfqpoint{1.583999in}{1.011768in}}%
\pgfpathlineto{\pgfqpoint{1.591775in}{1.015419in}}%
\pgfpathlineto{\pgfqpoint{1.573360in}{1.002040in}}%
\pgfpathlineto{\pgfqpoint{1.555029in}{0.988766in}}%
\pgfpathlineto{\pgfqpoint{1.547227in}{0.985185in}}%
\pgfpathlineto{\pgfqpoint{1.539404in}{0.981923in}}%
\pgfpathclose%
\pgfusepath{fill}%
\end{pgfscope}%
\begin{pgfscope}%
\pgfpathrectangle{\pgfqpoint{0.407592in}{0.222904in}}{\pgfqpoint{1.411163in}{1.411163in}}%
\pgfusepath{clip}%
\pgfsetbuttcap%
\pgfsetroundjoin%
\definecolor{currentfill}{rgb}{0.556863,0.654902,0.847059}%
\pgfsetfillcolor{currentfill}%
\pgfsetfillopacity{0.130000}%
\pgfsetlinewidth{0.000000pt}%
\definecolor{currentstroke}{rgb}{0.000000,0.000000,0.000000}%
\pgfsetstrokecolor{currentstroke}%
\pgfsetdash{}{0pt}%
\pgfpathmoveto{\pgfqpoint{1.211070in}{0.844971in}}%
\pgfpathlineto{\pgfqpoint{1.228674in}{0.853515in}}%
\pgfpathlineto{\pgfqpoint{1.246352in}{0.862437in}}%
\pgfpathlineto{\pgfqpoint{1.254693in}{0.859926in}}%
\pgfpathlineto{\pgfqpoint{1.263000in}{0.857695in}}%
\pgfpathlineto{\pgfqpoint{1.245359in}{0.848346in}}%
\pgfpathlineto{\pgfqpoint{1.227792in}{0.839360in}}%
\pgfpathlineto{\pgfqpoint{1.219448in}{0.842026in}}%
\pgfpathlineto{\pgfqpoint{1.211070in}{0.844971in}}%
\pgfpathclose%
\pgfusepath{fill}%
\end{pgfscope}%
\begin{pgfscope}%
\pgfpathrectangle{\pgfqpoint{0.407592in}{0.222904in}}{\pgfqpoint{1.411163in}{1.411163in}}%
\pgfusepath{clip}%
\pgfsetbuttcap%
\pgfsetroundjoin%
\definecolor{currentfill}{rgb}{0.556863,0.654902,0.847059}%
\pgfsetfillcolor{currentfill}%
\pgfsetfillopacity{0.130000}%
\pgfsetlinewidth{0.000000pt}%
\definecolor{currentstroke}{rgb}{0.000000,0.000000,0.000000}%
\pgfsetstrokecolor{currentstroke}%
\pgfsetdash{}{0pt}%
\pgfpathmoveto{\pgfqpoint{1.263000in}{0.857695in}}%
\pgfpathlineto{\pgfqpoint{1.280718in}{0.867381in}}%
\pgfpathlineto{\pgfqpoint{1.298512in}{0.877380in}}%
\pgfpathlineto{\pgfqpoint{1.306753in}{0.875838in}}%
\pgfpathlineto{\pgfqpoint{1.314963in}{0.874596in}}%
\pgfpathlineto{\pgfqpoint{1.297202in}{0.864211in}}%
\pgfpathlineto{\pgfqpoint{1.279519in}{0.854121in}}%
\pgfpathlineto{\pgfqpoint{1.271275in}{0.855756in}}%
\pgfpathlineto{\pgfqpoint{1.263000in}{0.857695in}}%
\pgfpathclose%
\pgfusepath{fill}%
\end{pgfscope}%
\begin{pgfscope}%
\pgfpathrectangle{\pgfqpoint{0.407592in}{0.222904in}}{\pgfqpoint{1.411163in}{1.411163in}}%
\pgfusepath{clip}%
\pgfsetbuttcap%
\pgfsetroundjoin%
\definecolor{currentfill}{rgb}{0.556863,0.654902,0.847059}%
\pgfsetfillcolor{currentfill}%
\pgfsetfillopacity{0.130000}%
\pgfsetlinewidth{0.000000pt}%
\definecolor{currentstroke}{rgb}{0.000000,0.000000,0.000000}%
\pgfsetstrokecolor{currentstroke}%
\pgfsetdash{}{0pt}%
\pgfpathmoveto{\pgfqpoint{1.159136in}{0.836713in}}%
\pgfpathlineto{\pgfqpoint{1.176637in}{0.843955in}}%
\pgfpathlineto{\pgfqpoint{1.194208in}{0.851644in}}%
\pgfpathlineto{\pgfqpoint{1.202657in}{0.848181in}}%
\pgfpathlineto{\pgfqpoint{1.211070in}{0.844971in}}%
\pgfpathlineto{\pgfqpoint{1.193538in}{0.836833in}}%
\pgfpathlineto{\pgfqpoint{1.176079in}{0.829133in}}%
\pgfpathlineto{\pgfqpoint{1.167625in}{0.832800in}}%
\pgfpathlineto{\pgfqpoint{1.159136in}{0.836713in}}%
\pgfpathclose%
\pgfusepath{fill}%
\end{pgfscope}%
\begin{pgfscope}%
\pgfpathrectangle{\pgfqpoint{0.407592in}{0.222904in}}{\pgfqpoint{1.411163in}{1.411163in}}%
\pgfusepath{clip}%
\pgfsetbuttcap%
\pgfsetroundjoin%
\definecolor{currentfill}{rgb}{0.556863,0.654902,0.847059}%
\pgfsetfillcolor{currentfill}%
\pgfsetfillopacity{0.130000}%
\pgfsetlinewidth{0.000000pt}%
\definecolor{currentstroke}{rgb}{0.000000,0.000000,0.000000}%
\pgfsetstrokecolor{currentstroke}%
\pgfsetdash{}{0pt}%
\pgfpathmoveto{\pgfqpoint{1.314963in}{0.874596in}}%
\pgfpathlineto{\pgfqpoint{1.332803in}{0.885254in}}%
\pgfpathlineto{\pgfqpoint{1.350722in}{0.896167in}}%
\pgfpathlineto{\pgfqpoint{1.358870in}{0.895582in}}%
\pgfpathlineto{\pgfqpoint{1.366990in}{0.895308in}}%
\pgfpathlineto{\pgfqpoint{1.349102in}{0.884062in}}%
\pgfpathlineto{\pgfqpoint{1.331294in}{0.873053in}}%
\pgfpathlineto{\pgfqpoint{1.323143in}{0.873664in}}%
\pgfpathlineto{\pgfqpoint{1.314963in}{0.874596in}}%
\pgfpathclose%
\pgfusepath{fill}%
\end{pgfscope}%
\begin{pgfscope}%
\pgfpathrectangle{\pgfqpoint{0.407592in}{0.222904in}}{\pgfqpoint{1.411163in}{1.411163in}}%
\pgfusepath{clip}%
\pgfsetbuttcap%
\pgfsetroundjoin%
\definecolor{currentfill}{rgb}{0.556863,0.654902,0.847059}%
\pgfsetfillcolor{currentfill}%
\pgfsetfillopacity{0.130000}%
\pgfsetlinewidth{0.000000pt}%
\definecolor{currentstroke}{rgb}{0.000000,0.000000,0.000000}%
\pgfsetstrokecolor{currentstroke}%
\pgfsetdash{}{0pt}%
\pgfpathmoveto{\pgfqpoint{1.107159in}{0.833106in}}%
\pgfpathlineto{\pgfqpoint{1.124569in}{0.838910in}}%
\pgfpathlineto{\pgfqpoint{1.142045in}{0.845227in}}%
\pgfpathlineto{\pgfqpoint{1.150609in}{0.840859in}}%
\pgfpathlineto{\pgfqpoint{1.159136in}{0.836713in}}%
\pgfpathlineto{\pgfqpoint{1.141704in}{0.829950in}}%
\pgfpathlineto{\pgfqpoint{1.124339in}{0.823699in}}%
\pgfpathlineto{\pgfqpoint{1.115768in}{0.828298in}}%
\pgfpathlineto{\pgfqpoint{1.107159in}{0.833106in}}%
\pgfpathclose%
\pgfusepath{fill}%
\end{pgfscope}%
\begin{pgfscope}%
\pgfpathrectangle{\pgfqpoint{0.407592in}{0.222904in}}{\pgfqpoint{1.411163in}{1.411163in}}%
\pgfusepath{clip}%
\pgfsetbuttcap%
\pgfsetroundjoin%
\definecolor{currentfill}{rgb}{0.556863,0.654902,0.847059}%
\pgfsetfillcolor{currentfill}%
\pgfsetfillopacity{0.130000}%
\pgfsetlinewidth{0.000000pt}%
\definecolor{currentstroke}{rgb}{0.000000,0.000000,0.000000}%
\pgfsetstrokecolor{currentstroke}%
\pgfsetdash{}{0pt}%
\pgfpathmoveto{\pgfqpoint{1.366990in}{0.895308in}}%
\pgfpathlineto{\pgfqpoint{1.384959in}{0.906774in}}%
\pgfpathlineto{\pgfqpoint{1.403009in}{0.918447in}}%
\pgfpathlineto{\pgfqpoint{1.411072in}{0.918782in}}%
\pgfpathlineto{\pgfqpoint{1.419108in}{0.919433in}}%
\pgfpathlineto{\pgfqpoint{1.401087in}{0.907486in}}%
\pgfpathlineto{\pgfqpoint{1.383148in}{0.895730in}}%
\pgfpathlineto{\pgfqpoint{1.375083in}{0.895354in}}%
\pgfpathlineto{\pgfqpoint{1.366990in}{0.895308in}}%
\pgfpathclose%
\pgfusepath{fill}%
\end{pgfscope}%
\begin{pgfscope}%
\pgfpathrectangle{\pgfqpoint{0.407592in}{0.222904in}}{\pgfqpoint{1.411163in}{1.411163in}}%
\pgfusepath{clip}%
\pgfsetbuttcap%
\pgfsetroundjoin%
\definecolor{currentfill}{rgb}{0.556863,0.654902,0.847059}%
\pgfsetfillcolor{currentfill}%
\pgfsetfillopacity{0.130000}%
\pgfsetlinewidth{0.000000pt}%
\definecolor{currentstroke}{rgb}{0.000000,0.000000,0.000000}%
\pgfsetstrokecolor{currentstroke}%
\pgfsetdash{}{0pt}%
\pgfpathmoveto{\pgfqpoint{1.419108in}{0.919433in}}%
\pgfpathlineto{\pgfqpoint{1.437211in}{0.931558in}}%
\pgfpathlineto{\pgfqpoint{1.455397in}{0.943851in}}%
\pgfpathlineto{\pgfqpoint{1.463380in}{0.945054in}}%
\pgfpathlineto{\pgfqpoint{1.471338in}{0.946572in}}%
\pgfpathlineto{\pgfqpoint{1.453180in}{0.934065in}}%
\pgfpathlineto{\pgfqpoint{1.435105in}{0.921711in}}%
\pgfpathlineto{\pgfqpoint{1.427119in}{0.920407in}}%
\pgfpathlineto{\pgfqpoint{1.419108in}{0.919433in}}%
\pgfpathclose%
\pgfusepath{fill}%
\end{pgfscope}%
\begin{pgfscope}%
\pgfpathrectangle{\pgfqpoint{0.407592in}{0.222904in}}{\pgfqpoint{1.411163in}{1.411163in}}%
\pgfusepath{clip}%
\pgfsetbuttcap%
\pgfsetroundjoin%
\definecolor{currentfill}{rgb}{0.556863,0.654902,0.847059}%
\pgfsetfillcolor{currentfill}%
\pgfsetfillopacity{0.130000}%
\pgfsetlinewidth{0.000000pt}%
\definecolor{currentstroke}{rgb}{0.000000,0.000000,0.000000}%
\pgfsetstrokecolor{currentstroke}%
\pgfsetdash{}{0pt}%
\pgfpathmoveto{\pgfqpoint{1.055099in}{0.834198in}}%
\pgfpathlineto{\pgfqpoint{1.072431in}{0.838466in}}%
\pgfpathlineto{\pgfqpoint{1.089824in}{0.843305in}}%
\pgfpathlineto{\pgfqpoint{1.098511in}{0.838112in}}%
\pgfpathlineto{\pgfqpoint{1.107159in}{0.833106in}}%
\pgfpathlineto{\pgfqpoint{1.089813in}{0.827849in}}%
\pgfpathlineto{\pgfqpoint{1.072530in}{0.823169in}}%
\pgfpathlineto{\pgfqpoint{1.063834in}{0.828598in}}%
\pgfpathlineto{\pgfqpoint{1.055099in}{0.834198in}}%
\pgfpathclose%
\pgfusepath{fill}%
\end{pgfscope}%
\begin{pgfscope}%
\pgfpathrectangle{\pgfqpoint{0.407592in}{0.222904in}}{\pgfqpoint{1.411163in}{1.411163in}}%
\pgfusepath{clip}%
\pgfsetbuttcap%
\pgfsetroundjoin%
\definecolor{currentfill}{rgb}{0.556863,0.654902,0.847059}%
\pgfsetfillcolor{currentfill}%
\pgfsetfillopacity{0.130000}%
\pgfsetlinewidth{0.000000pt}%
\definecolor{currentstroke}{rgb}{0.000000,0.000000,0.000000}%
\pgfsetstrokecolor{currentstroke}%
\pgfsetdash{}{0pt}%
\pgfpathmoveto{\pgfqpoint{1.471338in}{0.946572in}}%
\pgfpathlineto{\pgfqpoint{1.489580in}{0.959227in}}%
\pgfpathlineto{\pgfqpoint{1.507906in}{0.972020in}}%
\pgfpathlineto{\pgfqpoint{1.515813in}{0.974031in}}%
\pgfpathlineto{\pgfqpoint{1.523698in}{0.976349in}}%
\pgfpathlineto{\pgfqpoint{1.505400in}{0.963396in}}%
\pgfpathlineto{\pgfqpoint{1.487185in}{0.950572in}}%
\pgfpathlineto{\pgfqpoint{1.479273in}{0.948410in}}%
\pgfpathlineto{\pgfqpoint{1.471338in}{0.946572in}}%
\pgfpathclose%
\pgfusepath{fill}%
\end{pgfscope}%
\begin{pgfscope}%
\pgfpathrectangle{\pgfqpoint{0.407592in}{0.222904in}}{\pgfqpoint{1.411163in}{1.411163in}}%
\pgfusepath{clip}%
\pgfsetbuttcap%
\pgfsetroundjoin%
\definecolor{currentfill}{rgb}{0.556863,0.654902,0.847059}%
\pgfsetfillcolor{currentfill}%
\pgfsetfillopacity{0.130000}%
\pgfsetlinewidth{0.000000pt}%
\definecolor{currentstroke}{rgb}{0.000000,0.000000,0.000000}%
\pgfsetstrokecolor{currentstroke}%
\pgfsetdash{}{0pt}%
\pgfpathmoveto{\pgfqpoint{1.576203in}{1.008419in}}%
\pgfpathlineto{\pgfqpoint{1.594731in}{1.021829in}}%
\pgfpathlineto{\pgfqpoint{1.602513in}{1.025218in}}%
\pgfpathlineto{\pgfqpoint{1.610276in}{1.028900in}}%
\pgfpathlineto{\pgfqpoint{1.591775in}{1.015419in}}%
\pgfpathlineto{\pgfqpoint{1.583999in}{1.011768in}}%
\pgfusepath{fill}%
\end{pgfscope}%
\begin{pgfscope}%
\pgfpathrectangle{\pgfqpoint{0.407592in}{0.222904in}}{\pgfqpoint{1.411163in}{1.411163in}}%
\pgfusepath{clip}%
\pgfsetbuttcap%
\pgfsetroundjoin%
\definecolor{currentfill}{rgb}{0.556863,0.654902,0.847059}%
\pgfsetfillcolor{currentfill}%
\pgfsetfillopacity{0.130000}%
\pgfsetlinewidth{0.000000pt}%
\definecolor{currentstroke}{rgb}{0.000000,0.000000,0.000000}%
\pgfsetstrokecolor{currentstroke}%
\pgfsetdash{}{0pt}%
\pgfpathmoveto{\pgfqpoint{1.523698in}{0.976349in}}%
\pgfpathlineto{\pgfqpoint{1.542082in}{0.989425in}}%
\pgfpathlineto{\pgfqpoint{1.560551in}{1.002619in}}%
\pgfpathlineto{\pgfqpoint{1.568387in}{1.005370in}}%
\pgfpathlineto{\pgfqpoint{1.576203in}{1.008419in}}%
\pgfpathlineto{\pgfqpoint{1.557761in}{0.995115in}}%
\pgfpathlineto{\pgfqpoint{1.539404in}{0.981923in}}%
\pgfpathlineto{\pgfqpoint{1.531562in}{0.978978in}}%
\pgfpathlineto{\pgfqpoint{1.523698in}{0.976349in}}%
\pgfpathclose%
\pgfusepath{fill}%
\end{pgfscope}%
\begin{pgfscope}%
\pgfpathrectangle{\pgfqpoint{0.407592in}{0.222904in}}{\pgfqpoint{1.411163in}{1.411163in}}%
\pgfusepath{clip}%
\pgfsetbuttcap%
\pgfsetroundjoin%
\definecolor{currentfill}{rgb}{0.556863,0.654902,0.847059}%
\pgfsetfillcolor{currentfill}%
\pgfsetfillopacity{0.130000}%
\pgfsetlinewidth{0.000000pt}%
\definecolor{currentstroke}{rgb}{0.000000,0.000000,0.000000}%
\pgfsetstrokecolor{currentstroke}%
\pgfsetdash{}{0pt}%
\pgfpathmoveto{\pgfqpoint{1.246352in}{0.862437in}}%
\pgfpathlineto{\pgfqpoint{1.264105in}{0.871713in}}%
\pgfpathlineto{\pgfqpoint{1.281936in}{0.881319in}}%
\pgfpathlineto{\pgfqpoint{1.290240in}{0.879210in}}%
\pgfpathlineto{\pgfqpoint{1.298512in}{0.877380in}}%
\pgfpathlineto{\pgfqpoint{1.280718in}{0.867381in}}%
\pgfpathlineto{\pgfqpoint{1.263000in}{0.857695in}}%
\pgfpathlineto{\pgfqpoint{1.254693in}{0.859926in}}%
\pgfpathlineto{\pgfqpoint{1.246352in}{0.862437in}}%
\pgfpathclose%
\pgfusepath{fill}%
\end{pgfscope}%
\begin{pgfscope}%
\pgfpathrectangle{\pgfqpoint{0.407592in}{0.222904in}}{\pgfqpoint{1.411163in}{1.411163in}}%
\pgfusepath{clip}%
\pgfsetbuttcap%
\pgfsetroundjoin%
\definecolor{currentfill}{rgb}{0.556863,0.654902,0.847059}%
\pgfsetfillcolor{currentfill}%
\pgfsetfillopacity{0.130000}%
\pgfsetlinewidth{0.000000pt}%
\definecolor{currentstroke}{rgb}{0.000000,0.000000,0.000000}%
\pgfsetstrokecolor{currentstroke}%
\pgfsetdash{}{0pt}%
\pgfpathmoveto{\pgfqpoint{1.194208in}{0.851644in}}%
\pgfpathlineto{\pgfqpoint{1.211852in}{0.859752in}}%
\pgfpathlineto{\pgfqpoint{1.229568in}{0.868251in}}%
\pgfpathlineto{\pgfqpoint{1.237977in}{0.865216in}}%
\pgfpathlineto{\pgfqpoint{1.246352in}{0.862437in}}%
\pgfpathlineto{\pgfqpoint{1.228674in}{0.853515in}}%
\pgfpathlineto{\pgfqpoint{1.211070in}{0.844971in}}%
\pgfpathlineto{\pgfqpoint{1.202657in}{0.848181in}}%
\pgfpathlineto{\pgfqpoint{1.194208in}{0.851644in}}%
\pgfpathclose%
\pgfusepath{fill}%
\end{pgfscope}%
\begin{pgfscope}%
\pgfpathrectangle{\pgfqpoint{0.407592in}{0.222904in}}{\pgfqpoint{1.411163in}{1.411163in}}%
\pgfusepath{clip}%
\pgfsetbuttcap%
\pgfsetroundjoin%
\definecolor{currentfill}{rgb}{0.556863,0.654902,0.847059}%
\pgfsetfillcolor{currentfill}%
\pgfsetfillopacity{0.130000}%
\pgfsetlinewidth{0.000000pt}%
\definecolor{currentstroke}{rgb}{0.000000,0.000000,0.000000}%
\pgfsetstrokecolor{currentstroke}%
\pgfsetdash{}{0pt}%
\pgfpathmoveto{\pgfqpoint{1.298512in}{0.877380in}}%
\pgfpathlineto{\pgfqpoint{1.316385in}{0.887671in}}%
\pgfpathlineto{\pgfqpoint{1.334337in}{0.898235in}}%
\pgfpathlineto{\pgfqpoint{1.342544in}{0.897054in}}%
\pgfpathlineto{\pgfqpoint{1.350722in}{0.896167in}}%
\pgfpathlineto{\pgfqpoint{1.332803in}{0.885254in}}%
\pgfpathlineto{\pgfqpoint{1.314963in}{0.874596in}}%
\pgfpathlineto{\pgfqpoint{1.306753in}{0.875838in}}%
\pgfpathlineto{\pgfqpoint{1.298512in}{0.877380in}}%
\pgfpathclose%
\pgfusepath{fill}%
\end{pgfscope}%
\begin{pgfscope}%
\pgfpathrectangle{\pgfqpoint{0.407592in}{0.222904in}}{\pgfqpoint{1.411163in}{1.411163in}}%
\pgfusepath{clip}%
\pgfsetbuttcap%
\pgfsetroundjoin%
\definecolor{currentfill}{rgb}{0.556863,0.654902,0.847059}%
\pgfsetfillcolor{currentfill}%
\pgfsetfillopacity{0.130000}%
\pgfsetlinewidth{0.000000pt}%
\definecolor{currentstroke}{rgb}{0.000000,0.000000,0.000000}%
\pgfsetstrokecolor{currentstroke}%
\pgfsetdash{}{0pt}%
\pgfpathmoveto{\pgfqpoint{1.142045in}{0.845227in}}%
\pgfpathlineto{\pgfqpoint{1.159589in}{0.852027in}}%
\pgfpathlineto{\pgfqpoint{1.177203in}{0.859282in}}%
\pgfpathlineto{\pgfqpoint{1.185724in}{0.855348in}}%
\pgfpathlineto{\pgfqpoint{1.194208in}{0.851644in}}%
\pgfpathlineto{\pgfqpoint{1.176637in}{0.843955in}}%
\pgfpathlineto{\pgfqpoint{1.159136in}{0.836713in}}%
\pgfpathlineto{\pgfqpoint{1.150609in}{0.840859in}}%
\pgfpathlineto{\pgfqpoint{1.142045in}{0.845227in}}%
\pgfpathclose%
\pgfusepath{fill}%
\end{pgfscope}%
\begin{pgfscope}%
\pgfpathrectangle{\pgfqpoint{0.407592in}{0.222904in}}{\pgfqpoint{1.411163in}{1.411163in}}%
\pgfusepath{clip}%
\pgfsetbuttcap%
\pgfsetroundjoin%
\definecolor{currentfill}{rgb}{0.556863,0.654902,0.847059}%
\pgfsetfillcolor{currentfill}%
\pgfsetfillopacity{0.130000}%
\pgfsetlinewidth{0.000000pt}%
\definecolor{currentstroke}{rgb}{0.000000,0.000000,0.000000}%
\pgfsetstrokecolor{currentstroke}%
\pgfsetdash{}{0pt}%
\pgfpathmoveto{\pgfqpoint{1.350722in}{0.896167in}}%
\pgfpathlineto{\pgfqpoint{1.368722in}{0.907319in}}%
\pgfpathlineto{\pgfqpoint{1.386803in}{0.918694in}}%
\pgfpathlineto{\pgfqpoint{1.394920in}{0.918420in}}%
\pgfpathlineto{\pgfqpoint{1.403009in}{0.918447in}}%
\pgfpathlineto{\pgfqpoint{1.384959in}{0.906774in}}%
\pgfpathlineto{\pgfqpoint{1.366990in}{0.895308in}}%
\pgfpathlineto{\pgfqpoint{1.358870in}{0.895582in}}%
\pgfpathlineto{\pgfqpoint{1.350722in}{0.896167in}}%
\pgfpathclose%
\pgfusepath{fill}%
\end{pgfscope}%
\begin{pgfscope}%
\pgfpathrectangle{\pgfqpoint{0.407592in}{0.222904in}}{\pgfqpoint{1.411163in}{1.411163in}}%
\pgfusepath{clip}%
\pgfsetbuttcap%
\pgfsetroundjoin%
\definecolor{currentfill}{rgb}{0.556863,0.654902,0.847059}%
\pgfsetfillcolor{currentfill}%
\pgfsetfillopacity{0.130000}%
\pgfsetlinewidth{0.000000pt}%
\definecolor{currentstroke}{rgb}{0.000000,0.000000,0.000000}%
\pgfsetstrokecolor{currentstroke}%
\pgfsetdash{}{0pt}%
\pgfpathmoveto{\pgfqpoint{1.403009in}{0.918447in}}%
\pgfpathlineto{\pgfqpoint{1.421141in}{0.930314in}}%
\pgfpathlineto{\pgfqpoint{1.439357in}{0.942364in}}%
\pgfpathlineto{\pgfqpoint{1.447390in}{0.942956in}}%
\pgfpathlineto{\pgfqpoint{1.455397in}{0.943851in}}%
\pgfpathlineto{\pgfqpoint{1.437211in}{0.931558in}}%
\pgfpathlineto{\pgfqpoint{1.419108in}{0.919433in}}%
\pgfpathlineto{\pgfqpoint{1.411072in}{0.918782in}}%
\pgfpathlineto{\pgfqpoint{1.403009in}{0.918447in}}%
\pgfpathclose%
\pgfusepath{fill}%
\end{pgfscope}%
\begin{pgfscope}%
\pgfpathrectangle{\pgfqpoint{0.407592in}{0.222904in}}{\pgfqpoint{1.411163in}{1.411163in}}%
\pgfusepath{clip}%
\pgfsetbuttcap%
\pgfsetroundjoin%
\definecolor{currentfill}{rgb}{0.556863,0.654902,0.847059}%
\pgfsetfillcolor{currentfill}%
\pgfsetfillopacity{0.130000}%
\pgfsetlinewidth{0.000000pt}%
\definecolor{currentstroke}{rgb}{0.000000,0.000000,0.000000}%
\pgfsetstrokecolor{currentstroke}%
\pgfsetdash{}{0pt}%
\pgfpathmoveto{\pgfqpoint{1.089824in}{0.843305in}}%
\pgfpathlineto{\pgfqpoint{1.107280in}{0.848686in}}%
\pgfpathlineto{\pgfqpoint{1.124802in}{0.854579in}}%
\pgfpathlineto{\pgfqpoint{1.133443in}{0.849803in}}%
\pgfpathlineto{\pgfqpoint{1.142045in}{0.845227in}}%
\pgfpathlineto{\pgfqpoint{1.124569in}{0.838910in}}%
\pgfpathlineto{\pgfqpoint{1.107159in}{0.833106in}}%
\pgfpathlineto{\pgfqpoint{1.098511in}{0.838112in}}%
\pgfpathlineto{\pgfqpoint{1.089824in}{0.843305in}}%
\pgfpathclose%
\pgfusepath{fill}%
\end{pgfscope}%
\begin{pgfscope}%
\pgfpathrectangle{\pgfqpoint{0.407592in}{0.222904in}}{\pgfqpoint{1.411163in}{1.411163in}}%
\pgfusepath{clip}%
\pgfsetbuttcap%
\pgfsetroundjoin%
\definecolor{currentfill}{rgb}{0.556863,0.654902,0.847059}%
\pgfsetfillcolor{currentfill}%
\pgfsetfillopacity{0.130000}%
\pgfsetlinewidth{0.000000pt}%
\definecolor{currentstroke}{rgb}{0.000000,0.000000,0.000000}%
\pgfsetstrokecolor{currentstroke}%
\pgfsetdash{}{0pt}%
\pgfpathmoveto{\pgfqpoint{1.455397in}{0.943851in}}%
\pgfpathlineto{\pgfqpoint{1.473667in}{0.956304in}}%
\pgfpathlineto{\pgfqpoint{1.492021in}{0.968909in}}%
\pgfpathlineto{\pgfqpoint{1.499975in}{0.970315in}}%
\pgfpathlineto{\pgfqpoint{1.507906in}{0.972020in}}%
\pgfpathlineto{\pgfqpoint{1.489580in}{0.959227in}}%
\pgfpathlineto{\pgfqpoint{1.471338in}{0.946572in}}%
\pgfpathlineto{\pgfqpoint{1.463380in}{0.945054in}}%
\pgfpathlineto{\pgfqpoint{1.455397in}{0.943851in}}%
\pgfpathclose%
\pgfusepath{fill}%
\end{pgfscope}%
\begin{pgfscope}%
\pgfpathrectangle{\pgfqpoint{0.407592in}{0.222904in}}{\pgfqpoint{1.411163in}{1.411163in}}%
\pgfusepath{clip}%
\pgfsetbuttcap%
\pgfsetroundjoin%
\definecolor{currentfill}{rgb}{0.556863,0.654902,0.847059}%
\pgfsetfillcolor{currentfill}%
\pgfsetfillopacity{0.130000}%
\pgfsetlinewidth{0.000000pt}%
\definecolor{currentstroke}{rgb}{0.000000,0.000000,0.000000}%
\pgfsetstrokecolor{currentstroke}%
\pgfsetdash{}{0pt}%
\pgfpathmoveto{\pgfqpoint{1.560551in}{1.002619in}}%
\pgfpathlineto{\pgfqpoint{1.579106in}{1.015929in}}%
\pgfpathlineto{\pgfqpoint{1.586928in}{1.018733in}}%
\pgfpathlineto{\pgfqpoint{1.594731in}{1.021829in}}%
\pgfpathlineto{\pgfqpoint{1.576203in}{1.008419in}}%
\pgfpathlineto{\pgfqpoint{1.568387in}{1.005370in}}%
\pgfusepath{fill}%
\end{pgfscope}%
\begin{pgfscope}%
\pgfpathrectangle{\pgfqpoint{0.407592in}{0.222904in}}{\pgfqpoint{1.411163in}{1.411163in}}%
\pgfusepath{clip}%
\pgfsetbuttcap%
\pgfsetroundjoin%
\definecolor{currentfill}{rgb}{0.556863,0.654902,0.847059}%
\pgfsetfillcolor{currentfill}%
\pgfsetfillopacity{0.130000}%
\pgfsetlinewidth{0.000000pt}%
\definecolor{currentstroke}{rgb}{0.000000,0.000000,0.000000}%
\pgfsetstrokecolor{currentstroke}%
\pgfsetdash{}{0pt}%
\pgfpathmoveto{\pgfqpoint{1.507906in}{0.972020in}}%
\pgfpathlineto{\pgfqpoint{1.526317in}{0.984948in}}%
\pgfpathlineto{\pgfqpoint{1.544813in}{0.998003in}}%
\pgfpathlineto{\pgfqpoint{1.552693in}{1.000165in}}%
\pgfpathlineto{\pgfqpoint{1.560551in}{1.002619in}}%
\pgfpathlineto{\pgfqpoint{1.542082in}{0.989425in}}%
\pgfpathlineto{\pgfqpoint{1.523698in}{0.976349in}}%
\pgfpathlineto{\pgfqpoint{1.515813in}{0.974031in}}%
\pgfpathlineto{\pgfqpoint{1.507906in}{0.972020in}}%
\pgfpathclose%
\pgfusepath{fill}%
\end{pgfscope}%
\begin{pgfscope}%
\pgfpathrectangle{\pgfqpoint{0.407592in}{0.222904in}}{\pgfqpoint{1.411163in}{1.411163in}}%
\pgfusepath{clip}%
\pgfsetbuttcap%
\pgfsetroundjoin%
\definecolor{currentfill}{rgb}{0.556863,0.654902,0.847059}%
\pgfsetfillcolor{currentfill}%
\pgfsetfillopacity{0.130000}%
\pgfsetlinewidth{0.000000pt}%
\definecolor{currentstroke}{rgb}{0.000000,0.000000,0.000000}%
\pgfsetstrokecolor{currentstroke}%
\pgfsetdash{}{0pt}%
\pgfpathmoveto{\pgfqpoint{1.229568in}{0.868251in}}%
\pgfpathlineto{\pgfqpoint{1.247360in}{0.877118in}}%
\pgfpathlineto{\pgfqpoint{1.265228in}{0.886328in}}%
\pgfpathlineto{\pgfqpoint{1.273598in}{0.883695in}}%
\pgfpathlineto{\pgfqpoint{1.281936in}{0.881319in}}%
\pgfpathlineto{\pgfqpoint{1.264105in}{0.871713in}}%
\pgfpathlineto{\pgfqpoint{1.246352in}{0.862437in}}%
\pgfpathlineto{\pgfqpoint{1.237977in}{0.865216in}}%
\pgfpathlineto{\pgfqpoint{1.229568in}{0.868251in}}%
\pgfpathclose%
\pgfusepath{fill}%
\end{pgfscope}%
\begin{pgfscope}%
\pgfpathrectangle{\pgfqpoint{0.407592in}{0.222904in}}{\pgfqpoint{1.411163in}{1.411163in}}%
\pgfusepath{clip}%
\pgfsetbuttcap%
\pgfsetroundjoin%
\definecolor{currentfill}{rgb}{0.556863,0.654902,0.847059}%
\pgfsetfillcolor{currentfill}%
\pgfsetfillopacity{0.130000}%
\pgfsetlinewidth{0.000000pt}%
\definecolor{currentstroke}{rgb}{0.000000,0.000000,0.000000}%
\pgfsetstrokecolor{currentstroke}%
\pgfsetdash{}{0pt}%
\pgfpathmoveto{\pgfqpoint{1.281936in}{0.881319in}}%
\pgfpathlineto{\pgfqpoint{1.299843in}{0.891233in}}%
\pgfpathlineto{\pgfqpoint{1.317830in}{0.901437in}}%
\pgfpathlineto{\pgfqpoint{1.326099in}{0.899698in}}%
\pgfpathlineto{\pgfqpoint{1.334337in}{0.898235in}}%
\pgfpathlineto{\pgfqpoint{1.316385in}{0.887671in}}%
\pgfpathlineto{\pgfqpoint{1.298512in}{0.877380in}}%
\pgfpathlineto{\pgfqpoint{1.290240in}{0.879210in}}%
\pgfpathlineto{\pgfqpoint{1.281936in}{0.881319in}}%
\pgfpathclose%
\pgfusepath{fill}%
\end{pgfscope}%
\begin{pgfscope}%
\pgfpathrectangle{\pgfqpoint{0.407592in}{0.222904in}}{\pgfqpoint{1.411163in}{1.411163in}}%
\pgfusepath{clip}%
\pgfsetbuttcap%
\pgfsetroundjoin%
\definecolor{currentfill}{rgb}{0.556863,0.654902,0.847059}%
\pgfsetfillcolor{currentfill}%
\pgfsetfillopacity{0.130000}%
\pgfsetlinewidth{0.000000pt}%
\definecolor{currentstroke}{rgb}{0.000000,0.000000,0.000000}%
\pgfsetstrokecolor{currentstroke}%
\pgfsetdash{}{0pt}%
\pgfpathmoveto{\pgfqpoint{1.177203in}{0.859282in}}%
\pgfpathlineto{\pgfqpoint{1.194887in}{0.866964in}}%
\pgfpathlineto{\pgfqpoint{1.212645in}{0.875048in}}%
\pgfpathlineto{\pgfqpoint{1.221125in}{0.871532in}}%
\pgfpathlineto{\pgfqpoint{1.229568in}{0.868251in}}%
\pgfpathlineto{\pgfqpoint{1.211852in}{0.859752in}}%
\pgfpathlineto{\pgfqpoint{1.194208in}{0.851644in}}%
\pgfpathlineto{\pgfqpoint{1.185724in}{0.855348in}}%
\pgfpathlineto{\pgfqpoint{1.177203in}{0.859282in}}%
\pgfpathclose%
\pgfusepath{fill}%
\end{pgfscope}%
\begin{pgfscope}%
\pgfpathrectangle{\pgfqpoint{0.407592in}{0.222904in}}{\pgfqpoint{1.411163in}{1.411163in}}%
\pgfusepath{clip}%
\pgfsetbuttcap%
\pgfsetroundjoin%
\definecolor{currentfill}{rgb}{0.556863,0.654902,0.847059}%
\pgfsetfillcolor{currentfill}%
\pgfsetfillopacity{0.130000}%
\pgfsetlinewidth{0.000000pt}%
\definecolor{currentstroke}{rgb}{0.000000,0.000000,0.000000}%
\pgfsetstrokecolor{currentstroke}%
\pgfsetdash{}{0pt}%
\pgfpathmoveto{\pgfqpoint{1.334337in}{0.898235in}}%
\pgfpathlineto{\pgfqpoint{1.352369in}{0.909054in}}%
\pgfpathlineto{\pgfqpoint{1.370483in}{0.920114in}}%
\pgfpathlineto{\pgfqpoint{1.378657in}{0.919261in}}%
\pgfpathlineto{\pgfqpoint{1.386803in}{0.918694in}}%
\pgfpathlineto{\pgfqpoint{1.368722in}{0.907319in}}%
\pgfpathlineto{\pgfqpoint{1.350722in}{0.896167in}}%
\pgfpathlineto{\pgfqpoint{1.342544in}{0.897054in}}%
\pgfpathlineto{\pgfqpoint{1.334337in}{0.898235in}}%
\pgfpathclose%
\pgfusepath{fill}%
\end{pgfscope}%
\begin{pgfscope}%
\pgfpathrectangle{\pgfqpoint{0.407592in}{0.222904in}}{\pgfqpoint{1.411163in}{1.411163in}}%
\pgfusepath{clip}%
\pgfsetbuttcap%
\pgfsetroundjoin%
\definecolor{currentfill}{rgb}{0.556863,0.654902,0.847059}%
\pgfsetfillcolor{currentfill}%
\pgfsetfillopacity{0.130000}%
\pgfsetlinewidth{0.000000pt}%
\definecolor{currentstroke}{rgb}{0.000000,0.000000,0.000000}%
\pgfsetstrokecolor{currentstroke}%
\pgfsetdash{}{0pt}%
\pgfpathmoveto{\pgfqpoint{1.124802in}{0.854579in}}%
\pgfpathlineto{\pgfqpoint{1.142391in}{0.860958in}}%
\pgfpathlineto{\pgfqpoint{1.160049in}{0.867796in}}%
\pgfpathlineto{\pgfqpoint{1.168645in}{0.863434in}}%
\pgfpathlineto{\pgfqpoint{1.177203in}{0.859282in}}%
\pgfpathlineto{\pgfqpoint{1.159589in}{0.852027in}}%
\pgfpathlineto{\pgfqpoint{1.142045in}{0.845227in}}%
\pgfpathlineto{\pgfqpoint{1.133443in}{0.849803in}}%
\pgfpathlineto{\pgfqpoint{1.124802in}{0.854579in}}%
\pgfpathclose%
\pgfusepath{fill}%
\end{pgfscope}%
\begin{pgfscope}%
\pgfpathrectangle{\pgfqpoint{0.407592in}{0.222904in}}{\pgfqpoint{1.411163in}{1.411163in}}%
\pgfusepath{clip}%
\pgfsetbuttcap%
\pgfsetroundjoin%
\definecolor{currentfill}{rgb}{0.556863,0.654902,0.847059}%
\pgfsetfillcolor{currentfill}%
\pgfsetfillopacity{0.130000}%
\pgfsetlinewidth{0.000000pt}%
\definecolor{currentstroke}{rgb}{0.000000,0.000000,0.000000}%
\pgfsetstrokecolor{currentstroke}%
\pgfsetdash{}{0pt}%
\pgfpathmoveto{\pgfqpoint{1.386803in}{0.918694in}}%
\pgfpathlineto{\pgfqpoint{1.404966in}{0.930280in}}%
\pgfpathlineto{\pgfqpoint{1.423212in}{0.942065in}}%
\pgfpathlineto{\pgfqpoint{1.431298in}{0.942069in}}%
\pgfpathlineto{\pgfqpoint{1.439357in}{0.942364in}}%
\pgfpathlineto{\pgfqpoint{1.421141in}{0.930314in}}%
\pgfpathlineto{\pgfqpoint{1.403009in}{0.918447in}}%
\pgfpathlineto{\pgfqpoint{1.394920in}{0.918420in}}%
\pgfpathlineto{\pgfqpoint{1.386803in}{0.918694in}}%
\pgfpathclose%
\pgfusepath{fill}%
\end{pgfscope}%
\begin{pgfscope}%
\pgfpathrectangle{\pgfqpoint{0.407592in}{0.222904in}}{\pgfqpoint{1.411163in}{1.411163in}}%
\pgfusepath{clip}%
\pgfsetbuttcap%
\pgfsetroundjoin%
\definecolor{currentfill}{rgb}{0.556863,0.654902,0.847059}%
\pgfsetfillcolor{currentfill}%
\pgfsetfillopacity{0.130000}%
\pgfsetlinewidth{0.000000pt}%
\definecolor{currentstroke}{rgb}{0.000000,0.000000,0.000000}%
\pgfsetstrokecolor{currentstroke}%
\pgfsetdash{}{0pt}%
\pgfpathmoveto{\pgfqpoint{1.439357in}{0.942364in}}%
\pgfpathlineto{\pgfqpoint{1.457656in}{0.954589in}}%
\pgfpathlineto{\pgfqpoint{1.476040in}{0.966978in}}%
\pgfpathlineto{\pgfqpoint{1.484043in}{0.967798in}}%
\pgfpathlineto{\pgfqpoint{1.492021in}{0.968909in}}%
\pgfpathlineto{\pgfqpoint{1.473667in}{0.956304in}}%
\pgfpathlineto{\pgfqpoint{1.455397in}{0.943851in}}%
\pgfpathlineto{\pgfqpoint{1.447390in}{0.942956in}}%
\pgfpathlineto{\pgfqpoint{1.439357in}{0.942364in}}%
\pgfpathclose%
\pgfusepath{fill}%
\end{pgfscope}%
\begin{pgfscope}%
\pgfpathrectangle{\pgfqpoint{0.407592in}{0.222904in}}{\pgfqpoint{1.411163in}{1.411163in}}%
\pgfusepath{clip}%
\pgfsetbuttcap%
\pgfsetroundjoin%
\definecolor{currentfill}{rgb}{0.556863,0.654902,0.847059}%
\pgfsetfillcolor{currentfill}%
\pgfsetfillopacity{0.130000}%
\pgfsetlinewidth{0.000000pt}%
\definecolor{currentstroke}{rgb}{0.000000,0.000000,0.000000}%
\pgfsetstrokecolor{currentstroke}%
\pgfsetdash{}{0pt}%
\pgfpathmoveto{\pgfqpoint{1.544813in}{0.998003in}}%
\pgfpathlineto{\pgfqpoint{1.563397in}{1.011183in}}%
\pgfpathlineto{\pgfqpoint{1.571262in}{1.013413in}}%
\pgfpathlineto{\pgfqpoint{1.579106in}{1.015929in}}%
\pgfpathlineto{\pgfqpoint{1.560551in}{1.002619in}}%
\pgfpathlineto{\pgfqpoint{1.552693in}{1.000165in}}%
\pgfusepath{fill}%
\end{pgfscope}%
\begin{pgfscope}%
\pgfpathrectangle{\pgfqpoint{0.407592in}{0.222904in}}{\pgfqpoint{1.411163in}{1.411163in}}%
\pgfusepath{clip}%
\pgfsetbuttcap%
\pgfsetroundjoin%
\definecolor{currentfill}{rgb}{0.556863,0.654902,0.847059}%
\pgfsetfillcolor{currentfill}%
\pgfsetfillopacity{0.130000}%
\pgfsetlinewidth{0.000000pt}%
\definecolor{currentstroke}{rgb}{0.000000,0.000000,0.000000}%
\pgfsetstrokecolor{currentstroke}%
\pgfsetdash{}{0pt}%
\pgfpathmoveto{\pgfqpoint{1.492021in}{0.968909in}}%
\pgfpathlineto{\pgfqpoint{1.510461in}{0.981659in}}%
\pgfpathlineto{\pgfqpoint{1.528986in}{0.994547in}}%
\pgfpathlineto{\pgfqpoint{1.536912in}{0.996132in}}%
\pgfpathlineto{\pgfqpoint{1.544813in}{0.998003in}}%
\pgfpathlineto{\pgfqpoint{1.526317in}{0.984948in}}%
\pgfpathlineto{\pgfqpoint{1.507906in}{0.972020in}}%
\pgfpathlineto{\pgfqpoint{1.499975in}{0.970315in}}%
\pgfpathlineto{\pgfqpoint{1.492021in}{0.968909in}}%
\pgfpathclose%
\pgfusepath{fill}%
\end{pgfscope}%
\begin{pgfscope}%
\pgfpathrectangle{\pgfqpoint{0.407592in}{0.222904in}}{\pgfqpoint{1.411163in}{1.411163in}}%
\pgfusepath{clip}%
\pgfsetbuttcap%
\pgfsetroundjoin%
\definecolor{currentfill}{rgb}{0.556863,0.654902,0.847059}%
\pgfsetfillcolor{currentfill}%
\pgfsetfillopacity{0.130000}%
\pgfsetlinewidth{0.000000pt}%
\definecolor{currentstroke}{rgb}{0.000000,0.000000,0.000000}%
\pgfsetstrokecolor{currentstroke}%
\pgfsetdash{}{0pt}%
\pgfpathmoveto{\pgfqpoint{1.265228in}{0.886328in}}%
\pgfpathlineto{\pgfqpoint{1.283173in}{0.895863in}}%
\pgfpathlineto{\pgfqpoint{1.301196in}{0.905701in}}%
\pgfpathlineto{\pgfqpoint{1.309529in}{0.903440in}}%
\pgfpathlineto{\pgfqpoint{1.317830in}{0.901437in}}%
\pgfpathlineto{\pgfqpoint{1.299843in}{0.891233in}}%
\pgfpathlineto{\pgfqpoint{1.281936in}{0.881319in}}%
\pgfpathlineto{\pgfqpoint{1.273598in}{0.883695in}}%
\pgfpathlineto{\pgfqpoint{1.265228in}{0.886328in}}%
\pgfpathclose%
\pgfusepath{fill}%
\end{pgfscope}%
\begin{pgfscope}%
\pgfpathrectangle{\pgfqpoint{0.407592in}{0.222904in}}{\pgfqpoint{1.411163in}{1.411163in}}%
\pgfusepath{clip}%
\pgfsetbuttcap%
\pgfsetroundjoin%
\definecolor{currentfill}{rgb}{0.556863,0.654902,0.847059}%
\pgfsetfillcolor{currentfill}%
\pgfsetfillopacity{0.130000}%
\pgfsetlinewidth{0.000000pt}%
\definecolor{currentstroke}{rgb}{0.000000,0.000000,0.000000}%
\pgfsetstrokecolor{currentstroke}%
\pgfsetdash{}{0pt}%
\pgfpathmoveto{\pgfqpoint{1.212645in}{0.875048in}}%
\pgfpathlineto{\pgfqpoint{1.230476in}{0.883511in}}%
\pgfpathlineto{\pgfqpoint{1.248383in}{0.892330in}}%
\pgfpathlineto{\pgfqpoint{1.256823in}{0.889210in}}%
\pgfpathlineto{\pgfqpoint{1.265228in}{0.886328in}}%
\pgfpathlineto{\pgfqpoint{1.247360in}{0.877118in}}%
\pgfpathlineto{\pgfqpoint{1.229568in}{0.868251in}}%
\pgfpathlineto{\pgfqpoint{1.221125in}{0.871532in}}%
\pgfpathlineto{\pgfqpoint{1.212645in}{0.875048in}}%
\pgfpathclose%
\pgfusepath{fill}%
\end{pgfscope}%
\begin{pgfscope}%
\pgfpathrectangle{\pgfqpoint{0.407592in}{0.222904in}}{\pgfqpoint{1.411163in}{1.411163in}}%
\pgfusepath{clip}%
\pgfsetbuttcap%
\pgfsetroundjoin%
\definecolor{currentfill}{rgb}{0.556863,0.654902,0.847059}%
\pgfsetfillcolor{currentfill}%
\pgfsetfillopacity{0.130000}%
\pgfsetlinewidth{0.000000pt}%
\definecolor{currentstroke}{rgb}{0.000000,0.000000,0.000000}%
\pgfsetstrokecolor{currentstroke}%
\pgfsetdash{}{0pt}%
\pgfpathmoveto{\pgfqpoint{1.317830in}{0.901437in}}%
\pgfpathlineto{\pgfqpoint{1.335897in}{0.911912in}}%
\pgfpathlineto{\pgfqpoint{1.354044in}{0.922644in}}%
\pgfpathlineto{\pgfqpoint{1.362279in}{0.921244in}}%
\pgfpathlineto{\pgfqpoint{1.370483in}{0.920114in}}%
\pgfpathlineto{\pgfqpoint{1.352369in}{0.909054in}}%
\pgfpathlineto{\pgfqpoint{1.334337in}{0.898235in}}%
\pgfpathlineto{\pgfqpoint{1.326099in}{0.899698in}}%
\pgfpathlineto{\pgfqpoint{1.317830in}{0.901437in}}%
\pgfpathclose%
\pgfusepath{fill}%
\end{pgfscope}%
\begin{pgfscope}%
\pgfpathrectangle{\pgfqpoint{0.407592in}{0.222904in}}{\pgfqpoint{1.411163in}{1.411163in}}%
\pgfusepath{clip}%
\pgfsetbuttcap%
\pgfsetroundjoin%
\definecolor{currentfill}{rgb}{0.556863,0.654902,0.847059}%
\pgfsetfillcolor{currentfill}%
\pgfsetfillopacity{0.130000}%
\pgfsetlinewidth{0.000000pt}%
\definecolor{currentstroke}{rgb}{0.000000,0.000000,0.000000}%
\pgfsetstrokecolor{currentstroke}%
\pgfsetdash{}{0pt}%
\pgfpathmoveto{\pgfqpoint{1.160049in}{0.867796in}}%
\pgfpathlineto{\pgfqpoint{1.177777in}{0.875067in}}%
\pgfpathlineto{\pgfqpoint{1.195577in}{0.882748in}}%
\pgfpathlineto{\pgfqpoint{1.204129in}{0.878790in}}%
\pgfpathlineto{\pgfqpoint{1.212645in}{0.875048in}}%
\pgfpathlineto{\pgfqpoint{1.194887in}{0.866964in}}%
\pgfpathlineto{\pgfqpoint{1.177203in}{0.859282in}}%
\pgfpathlineto{\pgfqpoint{1.168645in}{0.863434in}}%
\pgfpathlineto{\pgfqpoint{1.160049in}{0.867796in}}%
\pgfpathclose%
\pgfusepath{fill}%
\end{pgfscope}%
\begin{pgfscope}%
\pgfpathrectangle{\pgfqpoint{0.407592in}{0.222904in}}{\pgfqpoint{1.411163in}{1.411163in}}%
\pgfusepath{clip}%
\pgfsetbuttcap%
\pgfsetroundjoin%
\definecolor{currentfill}{rgb}{0.556863,0.654902,0.847059}%
\pgfsetfillcolor{currentfill}%
\pgfsetfillopacity{0.130000}%
\pgfsetlinewidth{0.000000pt}%
\definecolor{currentstroke}{rgb}{0.000000,0.000000,0.000000}%
\pgfsetstrokecolor{currentstroke}%
\pgfsetdash{}{0pt}%
\pgfpathmoveto{\pgfqpoint{1.370483in}{0.920114in}}%
\pgfpathlineto{\pgfqpoint{1.388678in}{0.931400in}}%
\pgfpathlineto{\pgfqpoint{1.406957in}{0.942900in}}%
\pgfpathlineto{\pgfqpoint{1.415098in}{0.942344in}}%
\pgfpathlineto{\pgfqpoint{1.423212in}{0.942065in}}%
\pgfpathlineto{\pgfqpoint{1.404966in}{0.930280in}}%
\pgfpathlineto{\pgfqpoint{1.386803in}{0.918694in}}%
\pgfpathlineto{\pgfqpoint{1.378657in}{0.919261in}}%
\pgfpathlineto{\pgfqpoint{1.370483in}{0.920114in}}%
\pgfpathclose%
\pgfusepath{fill}%
\end{pgfscope}%
\begin{pgfscope}%
\pgfpathrectangle{\pgfqpoint{0.407592in}{0.222904in}}{\pgfqpoint{1.411163in}{1.411163in}}%
\pgfusepath{clip}%
\pgfsetbuttcap%
\pgfsetroundjoin%
\definecolor{currentfill}{rgb}{0.556863,0.654902,0.847059}%
\pgfsetfillcolor{currentfill}%
\pgfsetfillopacity{0.130000}%
\pgfsetlinewidth{0.000000pt}%
\definecolor{currentstroke}{rgb}{0.000000,0.000000,0.000000}%
\pgfsetstrokecolor{currentstroke}%
\pgfsetdash{}{0pt}%
\pgfpathmoveto{\pgfqpoint{1.423212in}{0.942065in}}%
\pgfpathlineto{\pgfqpoint{1.441542in}{0.954037in}}%
\pgfpathlineto{\pgfqpoint{1.459956in}{0.966189in}}%
\pgfpathlineto{\pgfqpoint{1.468011in}{0.966444in}}%
\pgfpathlineto{\pgfqpoint{1.476040in}{0.966978in}}%
\pgfpathlineto{\pgfqpoint{1.457656in}{0.954589in}}%
\pgfpathlineto{\pgfqpoint{1.439357in}{0.942364in}}%
\pgfpathlineto{\pgfqpoint{1.431298in}{0.942069in}}%
\pgfpathlineto{\pgfqpoint{1.423212in}{0.942065in}}%
\pgfpathclose%
\pgfusepath{fill}%
\end{pgfscope}%
\begin{pgfscope}%
\pgfpathrectangle{\pgfqpoint{0.407592in}{0.222904in}}{\pgfqpoint{1.411163in}{1.411163in}}%
\pgfusepath{clip}%
\pgfsetbuttcap%
\pgfsetroundjoin%
\definecolor{currentfill}{rgb}{0.556863,0.654902,0.847059}%
\pgfsetfillcolor{currentfill}%
\pgfsetfillopacity{0.130000}%
\pgfsetlinewidth{0.000000pt}%
\definecolor{currentstroke}{rgb}{0.000000,0.000000,0.000000}%
\pgfsetstrokecolor{currentstroke}%
\pgfsetdash{}{0pt}%
\pgfpathmoveto{\pgfqpoint{1.528986in}{0.994547in}}%
\pgfpathlineto{\pgfqpoint{1.547598in}{1.007570in}}%
\pgfpathlineto{\pgfqpoint{1.555509in}{1.009236in}}%
\pgfpathlineto{\pgfqpoint{1.563397in}{1.011183in}}%
\pgfpathlineto{\pgfqpoint{1.544813in}{0.998003in}}%
\pgfpathlineto{\pgfqpoint{1.536912in}{0.996132in}}%
\pgfusepath{fill}%
\end{pgfscope}%
\begin{pgfscope}%
\pgfpathrectangle{\pgfqpoint{0.407592in}{0.222904in}}{\pgfqpoint{1.411163in}{1.411163in}}%
\pgfusepath{clip}%
\pgfsetbuttcap%
\pgfsetroundjoin%
\definecolor{currentfill}{rgb}{0.556863,0.654902,0.847059}%
\pgfsetfillcolor{currentfill}%
\pgfsetfillopacity{0.130000}%
\pgfsetlinewidth{0.000000pt}%
\definecolor{currentstroke}{rgb}{0.000000,0.000000,0.000000}%
\pgfsetstrokecolor{currentstroke}%
\pgfsetdash{}{0pt}%
\pgfpathmoveto{\pgfqpoint{1.476040in}{0.966978in}}%
\pgfpathlineto{\pgfqpoint{1.494509in}{0.979525in}}%
\pgfpathlineto{\pgfqpoint{1.513064in}{0.992222in}}%
\pgfpathlineto{\pgfqpoint{1.521038in}{0.993246in}}%
\pgfpathlineto{\pgfqpoint{1.528986in}{0.994547in}}%
\pgfpathlineto{\pgfqpoint{1.510461in}{0.981659in}}%
\pgfpathlineto{\pgfqpoint{1.492021in}{0.968909in}}%
\pgfpathlineto{\pgfqpoint{1.484043in}{0.967798in}}%
\pgfpathlineto{\pgfqpoint{1.476040in}{0.966978in}}%
\pgfpathclose%
\pgfusepath{fill}%
\end{pgfscope}%
\begin{pgfscope}%
\pgfpathrectangle{\pgfqpoint{0.407592in}{0.222904in}}{\pgfqpoint{1.411163in}{1.411163in}}%
\pgfusepath{clip}%
\pgfsetbuttcap%
\pgfsetroundjoin%
\definecolor{currentfill}{rgb}{0.556863,0.654902,0.847059}%
\pgfsetfillcolor{currentfill}%
\pgfsetfillopacity{0.130000}%
\pgfsetlinewidth{0.000000pt}%
\definecolor{currentstroke}{rgb}{0.000000,0.000000,0.000000}%
\pgfsetstrokecolor{currentstroke}%
\pgfsetdash{}{0pt}%
\pgfpathmoveto{\pgfqpoint{1.248383in}{0.892330in}}%
\pgfpathlineto{\pgfqpoint{1.266367in}{0.901485in}}%
\pgfpathlineto{\pgfqpoint{1.284428in}{0.910956in}}%
\pgfpathlineto{\pgfqpoint{1.292829in}{0.908209in}}%
\pgfpathlineto{\pgfqpoint{1.301196in}{0.905701in}}%
\pgfpathlineto{\pgfqpoint{1.283173in}{0.895863in}}%
\pgfpathlineto{\pgfqpoint{1.265228in}{0.886328in}}%
\pgfpathlineto{\pgfqpoint{1.256823in}{0.889210in}}%
\pgfpathlineto{\pgfqpoint{1.248383in}{0.892330in}}%
\pgfpathclose%
\pgfusepath{fill}%
\end{pgfscope}%
\begin{pgfscope}%
\pgfpathrectangle{\pgfqpoint{0.407592in}{0.222904in}}{\pgfqpoint{1.411163in}{1.411163in}}%
\pgfusepath{clip}%
\pgfsetbuttcap%
\pgfsetroundjoin%
\definecolor{currentfill}{rgb}{0.556863,0.654902,0.847059}%
\pgfsetfillcolor{currentfill}%
\pgfsetfillopacity{0.130000}%
\pgfsetlinewidth{0.000000pt}%
\definecolor{currentstroke}{rgb}{0.000000,0.000000,0.000000}%
\pgfsetstrokecolor{currentstroke}%
\pgfsetdash{}{0pt}%
\pgfpathmoveto{\pgfqpoint{1.301196in}{0.905701in}}%
\pgfpathlineto{\pgfqpoint{1.319298in}{0.915825in}}%
\pgfpathlineto{\pgfqpoint{1.337482in}{0.926220in}}%
\pgfpathlineto{\pgfqpoint{1.345779in}{0.924305in}}%
\pgfpathlineto{\pgfqpoint{1.354044in}{0.922644in}}%
\pgfpathlineto{\pgfqpoint{1.335897in}{0.911912in}}%
\pgfpathlineto{\pgfqpoint{1.317830in}{0.901437in}}%
\pgfpathlineto{\pgfqpoint{1.309529in}{0.903440in}}%
\pgfpathlineto{\pgfqpoint{1.301196in}{0.905701in}}%
\pgfpathclose%
\pgfusepath{fill}%
\end{pgfscope}%
\begin{pgfscope}%
\pgfpathrectangle{\pgfqpoint{0.407592in}{0.222904in}}{\pgfqpoint{1.411163in}{1.411163in}}%
\pgfusepath{clip}%
\pgfsetbuttcap%
\pgfsetroundjoin%
\definecolor{currentfill}{rgb}{0.556863,0.654902,0.847059}%
\pgfsetfillcolor{currentfill}%
\pgfsetfillopacity{0.130000}%
\pgfsetlinewidth{0.000000pt}%
\definecolor{currentstroke}{rgb}{0.000000,0.000000,0.000000}%
\pgfsetstrokecolor{currentstroke}%
\pgfsetdash{}{0pt}%
\pgfpathmoveto{\pgfqpoint{1.195577in}{0.882748in}}%
\pgfpathlineto{\pgfqpoint{1.213450in}{0.890816in}}%
\pgfpathlineto{\pgfqpoint{1.231398in}{0.899249in}}%
\pgfpathlineto{\pgfqpoint{1.239909in}{0.895679in}}%
\pgfpathlineto{\pgfqpoint{1.248383in}{0.892330in}}%
\pgfpathlineto{\pgfqpoint{1.230476in}{0.883511in}}%
\pgfpathlineto{\pgfqpoint{1.212645in}{0.875048in}}%
\pgfpathlineto{\pgfqpoint{1.204129in}{0.878790in}}%
\pgfpathlineto{\pgfqpoint{1.195577in}{0.882748in}}%
\pgfpathclose%
\pgfusepath{fill}%
\end{pgfscope}%
\begin{pgfscope}%
\pgfpathrectangle{\pgfqpoint{0.407592in}{0.222904in}}{\pgfqpoint{1.411163in}{1.411163in}}%
\pgfusepath{clip}%
\pgfsetbuttcap%
\pgfsetroundjoin%
\definecolor{currentfill}{rgb}{0.556863,0.654902,0.847059}%
\pgfsetfillcolor{currentfill}%
\pgfsetfillopacity{0.130000}%
\pgfsetlinewidth{0.000000pt}%
\definecolor{currentstroke}{rgb}{0.000000,0.000000,0.000000}%
\pgfsetstrokecolor{currentstroke}%
\pgfsetdash{}{0pt}%
\pgfpathmoveto{\pgfqpoint{1.354044in}{0.922644in}}%
\pgfpathlineto{\pgfqpoint{1.372274in}{0.933616in}}%
\pgfpathlineto{\pgfqpoint{1.390585in}{0.944818in}}%
\pgfpathlineto{\pgfqpoint{1.398786in}{0.943727in}}%
\pgfpathlineto{\pgfqpoint{1.406957in}{0.942900in}}%
\pgfpathlineto{\pgfqpoint{1.388678in}{0.931400in}}%
\pgfpathlineto{\pgfqpoint{1.370483in}{0.920114in}}%
\pgfpathlineto{\pgfqpoint{1.362279in}{0.921244in}}%
\pgfpathlineto{\pgfqpoint{1.354044in}{0.922644in}}%
\pgfpathclose%
\pgfusepath{fill}%
\end{pgfscope}%
\begin{pgfscope}%
\pgfpathrectangle{\pgfqpoint{0.407592in}{0.222904in}}{\pgfqpoint{1.411163in}{1.411163in}}%
\pgfusepath{clip}%
\pgfsetbuttcap%
\pgfsetroundjoin%
\definecolor{currentfill}{rgb}{0.556863,0.654902,0.847059}%
\pgfsetfillcolor{currentfill}%
\pgfsetfillopacity{0.130000}%
\pgfsetlinewidth{0.000000pt}%
\definecolor{currentstroke}{rgb}{0.000000,0.000000,0.000000}%
\pgfsetstrokecolor{currentstroke}%
\pgfsetdash{}{0pt}%
\pgfpathmoveto{\pgfqpoint{1.406957in}{0.942900in}}%
\pgfpathlineto{\pgfqpoint{1.425319in}{0.954603in}}%
\pgfpathlineto{\pgfqpoint{1.443765in}{0.966498in}}%
\pgfpathlineto{\pgfqpoint{1.451875in}{0.966209in}}%
\pgfpathlineto{\pgfqpoint{1.459956in}{0.966189in}}%
\pgfpathlineto{\pgfqpoint{1.441542in}{0.954037in}}%
\pgfpathlineto{\pgfqpoint{1.423212in}{0.942065in}}%
\pgfpathlineto{\pgfqpoint{1.415098in}{0.942344in}}%
\pgfpathlineto{\pgfqpoint{1.406957in}{0.942900in}}%
\pgfpathclose%
\pgfusepath{fill}%
\end{pgfscope}%
\begin{pgfscope}%
\pgfpathrectangle{\pgfqpoint{0.407592in}{0.222904in}}{\pgfqpoint{1.411163in}{1.411163in}}%
\pgfusepath{clip}%
\pgfsetbuttcap%
\pgfsetroundjoin%
\definecolor{currentfill}{rgb}{0.556863,0.654902,0.847059}%
\pgfsetfillcolor{currentfill}%
\pgfsetfillopacity{0.130000}%
\pgfsetlinewidth{0.000000pt}%
\definecolor{currentstroke}{rgb}{0.000000,0.000000,0.000000}%
\pgfsetstrokecolor{currentstroke}%
\pgfsetdash{}{0pt}%
\pgfpathmoveto{\pgfqpoint{1.513064in}{0.992222in}}%
\pgfpathlineto{\pgfqpoint{1.531706in}{1.005064in}}%
\pgfpathlineto{\pgfqpoint{1.539664in}{1.006181in}}%
\pgfpathlineto{\pgfqpoint{1.547598in}{1.007570in}}%
\pgfpathlineto{\pgfqpoint{1.528986in}{0.994547in}}%
\pgfpathlineto{\pgfqpoint{1.521038in}{0.993246in}}%
\pgfusepath{fill}%
\end{pgfscope}%
\begin{pgfscope}%
\pgfpathrectangle{\pgfqpoint{0.407592in}{0.222904in}}{\pgfqpoint{1.411163in}{1.411163in}}%
\pgfusepath{clip}%
\pgfsetbuttcap%
\pgfsetroundjoin%
\definecolor{currentfill}{rgb}{0.556863,0.654902,0.847059}%
\pgfsetfillcolor{currentfill}%
\pgfsetfillopacity{0.130000}%
\pgfsetlinewidth{0.000000pt}%
\definecolor{currentstroke}{rgb}{0.000000,0.000000,0.000000}%
\pgfsetstrokecolor{currentstroke}%
\pgfsetdash{}{0pt}%
\pgfpathmoveto{\pgfqpoint{1.459956in}{0.966189in}}%
\pgfpathlineto{\pgfqpoint{1.478456in}{0.978511in}}%
\pgfpathlineto{\pgfqpoint{1.497042in}{0.990995in}}%
\pgfpathlineto{\pgfqpoint{1.505066in}{0.991474in}}%
\pgfpathlineto{\pgfqpoint{1.513064in}{0.992222in}}%
\pgfpathlineto{\pgfqpoint{1.494509in}{0.979525in}}%
\pgfpathlineto{\pgfqpoint{1.476040in}{0.966978in}}%
\pgfpathlineto{\pgfqpoint{1.468011in}{0.966444in}}%
\pgfpathlineto{\pgfqpoint{1.459956in}{0.966189in}}%
\pgfpathclose%
\pgfusepath{fill}%
\end{pgfscope}%
\begin{pgfscope}%
\pgfpathrectangle{\pgfqpoint{0.407592in}{0.222904in}}{\pgfqpoint{1.411163in}{1.411163in}}%
\pgfusepath{clip}%
\pgfsetbuttcap%
\pgfsetroundjoin%
\definecolor{currentfill}{rgb}{0.556863,0.654902,0.847059}%
\pgfsetfillcolor{currentfill}%
\pgfsetfillopacity{0.130000}%
\pgfsetlinewidth{0.000000pt}%
\definecolor{currentstroke}{rgb}{0.000000,0.000000,0.000000}%
\pgfsetstrokecolor{currentstroke}%
\pgfsetdash{}{0pt}%
\pgfpathmoveto{\pgfqpoint{1.284428in}{0.910956in}}%
\pgfpathlineto{\pgfqpoint{1.302569in}{0.920727in}}%
\pgfpathlineto{\pgfqpoint{1.320789in}{0.930781in}}%
\pgfpathlineto{\pgfqpoint{1.329152in}{0.928381in}}%
\pgfpathlineto{\pgfqpoint{1.337482in}{0.926220in}}%
\pgfpathlineto{\pgfqpoint{1.319298in}{0.915825in}}%
\pgfpathlineto{\pgfqpoint{1.301196in}{0.905701in}}%
\pgfpathlineto{\pgfqpoint{1.292829in}{0.908209in}}%
\pgfpathlineto{\pgfqpoint{1.284428in}{0.910956in}}%
\pgfpathclose%
\pgfusepath{fill}%
\end{pgfscope}%
\begin{pgfscope}%
\pgfpathrectangle{\pgfqpoint{0.407592in}{0.222904in}}{\pgfqpoint{1.411163in}{1.411163in}}%
\pgfusepath{clip}%
\pgfsetbuttcap%
\pgfsetroundjoin%
\definecolor{currentfill}{rgb}{0.556863,0.654902,0.847059}%
\pgfsetfillcolor{currentfill}%
\pgfsetfillopacity{0.130000}%
\pgfsetlinewidth{0.000000pt}%
\definecolor{currentstroke}{rgb}{0.000000,0.000000,0.000000}%
\pgfsetstrokecolor{currentstroke}%
\pgfsetdash{}{0pt}%
\pgfpathmoveto{\pgfqpoint{1.231398in}{0.899249in}}%
\pgfpathlineto{\pgfqpoint{1.249422in}{0.908029in}}%
\pgfpathlineto{\pgfqpoint{1.267523in}{0.917137in}}%
\pgfpathlineto{\pgfqpoint{1.275993in}{0.913935in}}%
\pgfpathlineto{\pgfqpoint{1.284428in}{0.910956in}}%
\pgfpathlineto{\pgfqpoint{1.266367in}{0.901485in}}%
\pgfpathlineto{\pgfqpoint{1.248383in}{0.892330in}}%
\pgfpathlineto{\pgfqpoint{1.239909in}{0.895679in}}%
\pgfpathlineto{\pgfqpoint{1.231398in}{0.899249in}}%
\pgfpathclose%
\pgfusepath{fill}%
\end{pgfscope}%
\begin{pgfscope}%
\pgfpathrectangle{\pgfqpoint{0.407592in}{0.222904in}}{\pgfqpoint{1.411163in}{1.411163in}}%
\pgfusepath{clip}%
\pgfsetbuttcap%
\pgfsetroundjoin%
\definecolor{currentfill}{rgb}{0.556863,0.654902,0.847059}%
\pgfsetfillcolor{currentfill}%
\pgfsetfillopacity{0.130000}%
\pgfsetlinewidth{0.000000pt}%
\definecolor{currentstroke}{rgb}{0.000000,0.000000,0.000000}%
\pgfsetstrokecolor{currentstroke}%
\pgfsetdash{}{0pt}%
\pgfpathmoveto{\pgfqpoint{1.337482in}{0.926220in}}%
\pgfpathlineto{\pgfqpoint{1.355746in}{0.936870in}}%
\pgfpathlineto{\pgfqpoint{1.374093in}{0.947763in}}%
\pgfpathlineto{\pgfqpoint{1.382355in}{0.946165in}}%
\pgfpathlineto{\pgfqpoint{1.390585in}{0.944818in}}%
\pgfpathlineto{\pgfqpoint{1.372274in}{0.933616in}}%
\pgfpathlineto{\pgfqpoint{1.354044in}{0.922644in}}%
\pgfpathlineto{\pgfqpoint{1.345779in}{0.924305in}}%
\pgfpathlineto{\pgfqpoint{1.337482in}{0.926220in}}%
\pgfpathclose%
\pgfusepath{fill}%
\end{pgfscope}%
\begin{pgfscope}%
\pgfpathrectangle{\pgfqpoint{0.407592in}{0.222904in}}{\pgfqpoint{1.411163in}{1.411163in}}%
\pgfusepath{clip}%
\pgfsetbuttcap%
\pgfsetroundjoin%
\definecolor{currentfill}{rgb}{0.556863,0.654902,0.847059}%
\pgfsetfillcolor{currentfill}%
\pgfsetfillopacity{0.130000}%
\pgfsetlinewidth{0.000000pt}%
\definecolor{currentstroke}{rgb}{0.000000,0.000000,0.000000}%
\pgfsetstrokecolor{currentstroke}%
\pgfsetdash{}{0pt}%
\pgfpathmoveto{\pgfqpoint{1.497042in}{0.990995in}}%
\pgfpathlineto{\pgfqpoint{1.515715in}{1.003635in}}%
\pgfpathlineto{\pgfqpoint{1.523724in}{1.004217in}}%
\pgfpathlineto{\pgfqpoint{1.531706in}{1.005064in}}%
\pgfpathlineto{\pgfqpoint{1.513064in}{0.992222in}}%
\pgfpathlineto{\pgfqpoint{1.505066in}{0.991474in}}%
\pgfusepath{fill}%
\end{pgfscope}%
\begin{pgfscope}%
\pgfpathrectangle{\pgfqpoint{0.407592in}{0.222904in}}{\pgfqpoint{1.411163in}{1.411163in}}%
\pgfusepath{clip}%
\pgfsetbuttcap%
\pgfsetroundjoin%
\definecolor{currentfill}{rgb}{0.556863,0.654902,0.847059}%
\pgfsetfillcolor{currentfill}%
\pgfsetfillopacity{0.130000}%
\pgfsetlinewidth{0.000000pt}%
\definecolor{currentstroke}{rgb}{0.000000,0.000000,0.000000}%
\pgfsetstrokecolor{currentstroke}%
\pgfsetdash{}{0pt}%
\pgfpathmoveto{\pgfqpoint{1.390585in}{0.944818in}}%
\pgfpathlineto{\pgfqpoint{1.408981in}{0.956235in}}%
\pgfpathlineto{\pgfqpoint{1.427461in}{0.967859in}}%
\pgfpathlineto{\pgfqpoint{1.435628in}{0.967050in}}%
\pgfpathlineto{\pgfqpoint{1.443765in}{0.966498in}}%
\pgfpathlineto{\pgfqpoint{1.425319in}{0.954603in}}%
\pgfpathlineto{\pgfqpoint{1.406957in}{0.942900in}}%
\pgfpathlineto{\pgfqpoint{1.398786in}{0.943727in}}%
\pgfpathlineto{\pgfqpoint{1.390585in}{0.944818in}}%
\pgfpathclose%
\pgfusepath{fill}%
\end{pgfscope}%
\begin{pgfscope}%
\pgfpathrectangle{\pgfqpoint{0.407592in}{0.222904in}}{\pgfqpoint{1.411163in}{1.411163in}}%
\pgfusepath{clip}%
\pgfsetbuttcap%
\pgfsetroundjoin%
\definecolor{currentfill}{rgb}{0.556863,0.654902,0.847059}%
\pgfsetfillcolor{currentfill}%
\pgfsetfillopacity{0.130000}%
\pgfsetlinewidth{0.000000pt}%
\definecolor{currentstroke}{rgb}{0.000000,0.000000,0.000000}%
\pgfsetstrokecolor{currentstroke}%
\pgfsetdash{}{0pt}%
\pgfpathmoveto{\pgfqpoint{1.443765in}{0.966498in}}%
\pgfpathlineto{\pgfqpoint{1.462297in}{0.978576in}}%
\pgfpathlineto{\pgfqpoint{1.480915in}{0.990829in}}%
\pgfpathlineto{\pgfqpoint{1.488992in}{0.990782in}}%
\pgfpathlineto{\pgfqpoint{1.497042in}{0.990995in}}%
\pgfpathlineto{\pgfqpoint{1.478456in}{0.978511in}}%
\pgfpathlineto{\pgfqpoint{1.459956in}{0.966189in}}%
\pgfpathlineto{\pgfqpoint{1.451875in}{0.966209in}}%
\pgfpathlineto{\pgfqpoint{1.443765in}{0.966498in}}%
\pgfpathclose%
\pgfusepath{fill}%
\end{pgfscope}%
\begin{pgfscope}%
\pgfpathrectangle{\pgfqpoint{0.407592in}{0.222904in}}{\pgfqpoint{1.411163in}{1.411163in}}%
\pgfusepath{clip}%
\pgfsetbuttcap%
\pgfsetroundjoin%
\definecolor{currentfill}{rgb}{0.556863,0.654902,0.847059}%
\pgfsetfillcolor{currentfill}%
\pgfsetfillopacity{0.130000}%
\pgfsetlinewidth{0.000000pt}%
\definecolor{currentstroke}{rgb}{0.000000,0.000000,0.000000}%
\pgfsetstrokecolor{currentstroke}%
\pgfsetdash{}{0pt}%
\pgfpathmoveto{\pgfqpoint{1.480915in}{0.990829in}}%
\pgfpathlineto{\pgfqpoint{1.499620in}{1.003250in}}%
\pgfpathlineto{\pgfqpoint{1.507681in}{1.003314in}}%
\pgfpathlineto{\pgfqpoint{1.515715in}{1.003635in}}%
\pgfpathlineto{\pgfqpoint{1.497042in}{0.990995in}}%
\pgfpathlineto{\pgfqpoint{1.488992in}{0.990782in}}%
\pgfusepath{fill}%
\end{pgfscope}%
\begin{pgfscope}%
\pgfpathrectangle{\pgfqpoint{0.407592in}{0.222904in}}{\pgfqpoint{1.411163in}{1.411163in}}%
\pgfusepath{clip}%
\pgfsetbuttcap%
\pgfsetroundjoin%
\definecolor{currentfill}{rgb}{0.556863,0.654902,0.847059}%
\pgfsetfillcolor{currentfill}%
\pgfsetfillopacity{0.130000}%
\pgfsetlinewidth{0.000000pt}%
\definecolor{currentstroke}{rgb}{0.000000,0.000000,0.000000}%
\pgfsetstrokecolor{currentstroke}%
\pgfsetdash{}{0pt}%
\pgfpathmoveto{\pgfqpoint{1.267523in}{0.917137in}}%
\pgfpathlineto{\pgfqpoint{1.285703in}{0.926555in}}%
\pgfpathlineto{\pgfqpoint{1.303963in}{0.936268in}}%
\pgfpathlineto{\pgfqpoint{1.312393in}{0.933413in}}%
\pgfpathlineto{\pgfqpoint{1.320789in}{0.930781in}}%
\pgfpathlineto{\pgfqpoint{1.302569in}{0.920727in}}%
\pgfpathlineto{\pgfqpoint{1.284428in}{0.910956in}}%
\pgfpathlineto{\pgfqpoint{1.275993in}{0.913935in}}%
\pgfpathlineto{\pgfqpoint{1.267523in}{0.917137in}}%
\pgfpathclose%
\pgfusepath{fill}%
\end{pgfscope}%
\begin{pgfscope}%
\pgfpathrectangle{\pgfqpoint{0.407592in}{0.222904in}}{\pgfqpoint{1.411163in}{1.411163in}}%
\pgfusepath{clip}%
\pgfsetbuttcap%
\pgfsetroundjoin%
\definecolor{currentfill}{rgb}{0.556863,0.654902,0.847059}%
\pgfsetfillcolor{currentfill}%
\pgfsetfillopacity{0.130000}%
\pgfsetlinewidth{0.000000pt}%
\definecolor{currentstroke}{rgb}{0.000000,0.000000,0.000000}%
\pgfsetstrokecolor{currentstroke}%
\pgfsetdash{}{0pt}%
\pgfpathmoveto{\pgfqpoint{1.320789in}{0.930781in}}%
\pgfpathlineto{\pgfqpoint{1.339091in}{0.941104in}}%
\pgfpathlineto{\pgfqpoint{1.357475in}{0.951681in}}%
\pgfpathlineto{\pgfqpoint{1.365800in}{0.949604in}}%
\pgfpathlineto{\pgfqpoint{1.374093in}{0.947763in}}%
\pgfpathlineto{\pgfqpoint{1.355746in}{0.936870in}}%
\pgfpathlineto{\pgfqpoint{1.337482in}{0.926220in}}%
\pgfpathlineto{\pgfqpoint{1.329152in}{0.928381in}}%
\pgfpathlineto{\pgfqpoint{1.320789in}{0.930781in}}%
\pgfpathclose%
\pgfusepath{fill}%
\end{pgfscope}%
\begin{pgfscope}%
\pgfpathrectangle{\pgfqpoint{0.407592in}{0.222904in}}{\pgfqpoint{1.411163in}{1.411163in}}%
\pgfusepath{clip}%
\pgfsetbuttcap%
\pgfsetroundjoin%
\definecolor{currentfill}{rgb}{0.556863,0.654902,0.847059}%
\pgfsetfillcolor{currentfill}%
\pgfsetfillopacity{0.130000}%
\pgfsetlinewidth{0.000000pt}%
\definecolor{currentstroke}{rgb}{0.000000,0.000000,0.000000}%
\pgfsetstrokecolor{currentstroke}%
\pgfsetdash{}{0pt}%
\pgfpathmoveto{\pgfqpoint{1.374093in}{0.947763in}}%
\pgfpathlineto{\pgfqpoint{1.392524in}{0.958885in}}%
\pgfpathlineto{\pgfqpoint{1.411039in}{0.970226in}}%
\pgfpathlineto{\pgfqpoint{1.419265in}{0.968920in}}%
\pgfpathlineto{\pgfqpoint{1.427461in}{0.967859in}}%
\pgfpathlineto{\pgfqpoint{1.408981in}{0.956235in}}%
\pgfpathlineto{\pgfqpoint{1.390585in}{0.944818in}}%
\pgfpathlineto{\pgfqpoint{1.382355in}{0.946165in}}%
\pgfpathlineto{\pgfqpoint{1.374093in}{0.947763in}}%
\pgfpathclose%
\pgfusepath{fill}%
\end{pgfscope}%
\begin{pgfscope}%
\pgfpathrectangle{\pgfqpoint{0.407592in}{0.222904in}}{\pgfqpoint{1.411163in}{1.411163in}}%
\pgfusepath{clip}%
\pgfsetbuttcap%
\pgfsetroundjoin%
\definecolor{currentfill}{rgb}{0.556863,0.654902,0.847059}%
\pgfsetfillcolor{currentfill}%
\pgfsetfillopacity{0.130000}%
\pgfsetlinewidth{0.000000pt}%
\definecolor{currentstroke}{rgb}{0.000000,0.000000,0.000000}%
\pgfsetstrokecolor{currentstroke}%
\pgfsetdash{}{0pt}%
\pgfpathmoveto{\pgfqpoint{1.427461in}{0.967859in}}%
\pgfpathlineto{\pgfqpoint{1.446026in}{0.979679in}}%
\pgfpathlineto{\pgfqpoint{1.464678in}{0.991685in}}%
\pgfpathlineto{\pgfqpoint{1.472811in}{0.991132in}}%
\pgfpathlineto{\pgfqpoint{1.480915in}{0.990829in}}%
\pgfpathlineto{\pgfqpoint{1.462297in}{0.978576in}}%
\pgfpathlineto{\pgfqpoint{1.443765in}{0.966498in}}%
\pgfpathlineto{\pgfqpoint{1.435628in}{0.967050in}}%
\pgfpathlineto{\pgfqpoint{1.427461in}{0.967859in}}%
\pgfpathclose%
\pgfusepath{fill}%
\end{pgfscope}%
\begin{pgfscope}%
\pgfpathrectangle{\pgfqpoint{0.407592in}{0.222904in}}{\pgfqpoint{1.411163in}{1.411163in}}%
\pgfusepath{clip}%
\pgfsetbuttcap%
\pgfsetroundjoin%
\definecolor{currentfill}{rgb}{0.556863,0.654902,0.847059}%
\pgfsetfillcolor{currentfill}%
\pgfsetfillopacity{0.130000}%
\pgfsetlinewidth{0.000000pt}%
\definecolor{currentstroke}{rgb}{0.000000,0.000000,0.000000}%
\pgfsetstrokecolor{currentstroke}%
\pgfsetdash{}{0pt}%
\pgfpathmoveto{\pgfqpoint{1.464678in}{0.991685in}}%
\pgfpathlineto{\pgfqpoint{1.483416in}{1.003871in}}%
\pgfpathlineto{\pgfqpoint{1.491532in}{1.003437in}}%
\pgfpathlineto{\pgfqpoint{1.499620in}{1.003250in}}%
\pgfpathlineto{\pgfqpoint{1.480915in}{0.990829in}}%
\pgfpathlineto{\pgfqpoint{1.472811in}{0.991132in}}%
\pgfusepath{fill}%
\end{pgfscope}%
\begin{pgfscope}%
\pgfpathrectangle{\pgfqpoint{0.407592in}{0.222904in}}{\pgfqpoint{1.411163in}{1.411163in}}%
\pgfusepath{clip}%
\pgfsetbuttcap%
\pgfsetroundjoin%
\definecolor{currentfill}{rgb}{0.556863,0.654902,0.847059}%
\pgfsetfillcolor{currentfill}%
\pgfsetfillopacity{0.130000}%
\pgfsetlinewidth{0.000000pt}%
\definecolor{currentstroke}{rgb}{0.000000,0.000000,0.000000}%
\pgfsetstrokecolor{currentstroke}%
\pgfsetdash{}{0pt}%
\pgfpathmoveto{\pgfqpoint{1.303963in}{0.936268in}}%
\pgfpathlineto{\pgfqpoint{1.322303in}{0.946261in}}%
\pgfpathlineto{\pgfqpoint{1.340726in}{0.956521in}}%
\pgfpathlineto{\pgfqpoint{1.349117in}{0.953989in}}%
\pgfpathlineto{\pgfqpoint{1.357475in}{0.951681in}}%
\pgfpathlineto{\pgfqpoint{1.339091in}{0.941104in}}%
\pgfpathlineto{\pgfqpoint{1.320789in}{0.930781in}}%
\pgfpathlineto{\pgfqpoint{1.312393in}{0.933413in}}%
\pgfpathlineto{\pgfqpoint{1.303963in}{0.936268in}}%
\pgfpathclose%
\pgfusepath{fill}%
\end{pgfscope}%
\begin{pgfscope}%
\pgfpathrectangle{\pgfqpoint{0.407592in}{0.222904in}}{\pgfqpoint{1.411163in}{1.411163in}}%
\pgfusepath{clip}%
\pgfsetbuttcap%
\pgfsetroundjoin%
\definecolor{currentfill}{rgb}{0.556863,0.654902,0.847059}%
\pgfsetfillcolor{currentfill}%
\pgfsetfillopacity{0.130000}%
\pgfsetlinewidth{0.000000pt}%
\definecolor{currentstroke}{rgb}{0.000000,0.000000,0.000000}%
\pgfsetstrokecolor{currentstroke}%
\pgfsetdash{}{0pt}%
\pgfpathmoveto{\pgfqpoint{1.357475in}{0.951681in}}%
\pgfpathlineto{\pgfqpoint{1.375942in}{0.962501in}}%
\pgfpathlineto{\pgfqpoint{1.394493in}{0.973552in}}%
\pgfpathlineto{\pgfqpoint{1.402782in}{0.971772in}}%
\pgfpathlineto{\pgfqpoint{1.411039in}{0.970226in}}%
\pgfpathlineto{\pgfqpoint{1.392524in}{0.958885in}}%
\pgfpathlineto{\pgfqpoint{1.374093in}{0.947763in}}%
\pgfpathlineto{\pgfqpoint{1.365800in}{0.949604in}}%
\pgfpathlineto{\pgfqpoint{1.357475in}{0.951681in}}%
\pgfpathclose%
\pgfusepath{fill}%
\end{pgfscope}%
\begin{pgfscope}%
\pgfpathrectangle{\pgfqpoint{0.407592in}{0.222904in}}{\pgfqpoint{1.411163in}{1.411163in}}%
\pgfusepath{clip}%
\pgfsetbuttcap%
\pgfsetroundjoin%
\definecolor{currentfill}{rgb}{0.556863,0.654902,0.847059}%
\pgfsetfillcolor{currentfill}%
\pgfsetfillopacity{0.130000}%
\pgfsetlinewidth{0.000000pt}%
\definecolor{currentstroke}{rgb}{0.000000,0.000000,0.000000}%
\pgfsetstrokecolor{currentstroke}%
\pgfsetdash{}{0pt}%
\pgfpathmoveto{\pgfqpoint{1.411039in}{0.970226in}}%
\pgfpathlineto{\pgfqpoint{1.429639in}{0.981775in}}%
\pgfpathlineto{\pgfqpoint{1.448325in}{0.993524in}}%
\pgfpathlineto{\pgfqpoint{1.456516in}{0.992484in}}%
\pgfpathlineto{\pgfqpoint{1.464678in}{0.991685in}}%
\pgfpathlineto{\pgfqpoint{1.446026in}{0.979679in}}%
\pgfpathlineto{\pgfqpoint{1.427461in}{0.967859in}}%
\pgfpathlineto{\pgfqpoint{1.419265in}{0.968920in}}%
\pgfpathlineto{\pgfqpoint{1.411039in}{0.970226in}}%
\pgfpathclose%
\pgfusepath{fill}%
\end{pgfscope}%
\begin{pgfscope}%
\pgfpathrectangle{\pgfqpoint{0.407592in}{0.222904in}}{\pgfqpoint{1.411163in}{1.411163in}}%
\pgfusepath{clip}%
\pgfsetbuttcap%
\pgfsetroundjoin%
\definecolor{currentfill}{rgb}{0.556863,0.654902,0.847059}%
\pgfsetfillcolor{currentfill}%
\pgfsetfillopacity{0.130000}%
\pgfsetlinewidth{0.000000pt}%
\definecolor{currentstroke}{rgb}{0.000000,0.000000,0.000000}%
\pgfsetstrokecolor{currentstroke}%
\pgfsetdash{}{0pt}%
\pgfpathmoveto{\pgfqpoint{1.448325in}{0.993524in}}%
\pgfpathlineto{\pgfqpoint{1.467098in}{1.005463in}}%
\pgfpathlineto{\pgfqpoint{1.475272in}{1.004548in}}%
\pgfpathlineto{\pgfqpoint{1.483416in}{1.003871in}}%
\pgfpathlineto{\pgfqpoint{1.464678in}{0.991685in}}%
\pgfpathlineto{\pgfqpoint{1.456516in}{0.992484in}}%
\pgfusepath{fill}%
\end{pgfscope}%
\begin{pgfscope}%
\pgfpathrectangle{\pgfqpoint{0.407592in}{0.222904in}}{\pgfqpoint{1.411163in}{1.411163in}}%
\pgfusepath{clip}%
\pgfsetbuttcap%
\pgfsetroundjoin%
\definecolor{currentfill}{rgb}{0.556863,0.654902,0.847059}%
\pgfsetfillcolor{currentfill}%
\pgfsetfillopacity{0.130000}%
\pgfsetlinewidth{0.000000pt}%
\definecolor{currentstroke}{rgb}{0.000000,0.000000,0.000000}%
\pgfsetstrokecolor{currentstroke}%
\pgfsetdash{}{0pt}%
\pgfpathmoveto{\pgfqpoint{1.340726in}{0.956521in}}%
\pgfpathlineto{\pgfqpoint{1.359231in}{0.967034in}}%
\pgfpathlineto{\pgfqpoint{1.377820in}{0.977789in}}%
\pgfpathlineto{\pgfqpoint{1.386173in}{0.975559in}}%
\pgfpathlineto{\pgfqpoint{1.394493in}{0.973552in}}%
\pgfpathlineto{\pgfqpoint{1.375942in}{0.962501in}}%
\pgfpathlineto{\pgfqpoint{1.357475in}{0.951681in}}%
\pgfpathlineto{\pgfqpoint{1.349117in}{0.953989in}}%
\pgfpathlineto{\pgfqpoint{1.340726in}{0.956521in}}%
\pgfpathclose%
\pgfusepath{fill}%
\end{pgfscope}%
\begin{pgfscope}%
\pgfpathrectangle{\pgfqpoint{0.407592in}{0.222904in}}{\pgfqpoint{1.411163in}{1.411163in}}%
\pgfusepath{clip}%
\pgfsetbuttcap%
\pgfsetroundjoin%
\definecolor{currentfill}{rgb}{0.556863,0.654902,0.847059}%
\pgfsetfillcolor{currentfill}%
\pgfsetfillopacity{0.130000}%
\pgfsetlinewidth{0.000000pt}%
\definecolor{currentstroke}{rgb}{0.000000,0.000000,0.000000}%
\pgfsetstrokecolor{currentstroke}%
\pgfsetdash{}{0pt}%
\pgfpathmoveto{\pgfqpoint{1.394493in}{0.973552in}}%
\pgfpathlineto{\pgfqpoint{1.413130in}{0.984822in}}%
\pgfpathlineto{\pgfqpoint{1.431852in}{0.996304in}}%
\pgfpathlineto{\pgfqpoint{1.440104in}{0.994799in}}%
\pgfpathlineto{\pgfqpoint{1.448325in}{0.993524in}}%
\pgfpathlineto{\pgfqpoint{1.429639in}{0.981775in}}%
\pgfpathlineto{\pgfqpoint{1.411039in}{0.970226in}}%
\pgfpathlineto{\pgfqpoint{1.402782in}{0.971772in}}%
\pgfpathlineto{\pgfqpoint{1.394493in}{0.973552in}}%
\pgfpathclose%
\pgfusepath{fill}%
\end{pgfscope}%
\begin{pgfscope}%
\pgfpathrectangle{\pgfqpoint{0.407592in}{0.222904in}}{\pgfqpoint{1.411163in}{1.411163in}}%
\pgfusepath{clip}%
\pgfsetbuttcap%
\pgfsetroundjoin%
\definecolor{currentfill}{rgb}{0.556863,0.654902,0.847059}%
\pgfsetfillcolor{currentfill}%
\pgfsetfillopacity{0.130000}%
\pgfsetlinewidth{0.000000pt}%
\definecolor{currentstroke}{rgb}{0.000000,0.000000,0.000000}%
\pgfsetstrokecolor{currentstroke}%
\pgfsetdash{}{0pt}%
\pgfpathmoveto{\pgfqpoint{1.431852in}{0.996304in}}%
\pgfpathlineto{\pgfqpoint{1.450661in}{1.007987in}}%
\pgfpathlineto{\pgfqpoint{1.458895in}{1.006611in}}%
\pgfpathlineto{\pgfqpoint{1.467098in}{1.005463in}}%
\pgfpathlineto{\pgfqpoint{1.448325in}{0.993524in}}%
\pgfpathlineto{\pgfqpoint{1.440104in}{0.994799in}}%
\pgfusepath{fill}%
\end{pgfscope}%
\begin{pgfscope}%
\pgfpathrectangle{\pgfqpoint{0.407592in}{0.222904in}}{\pgfqpoint{1.411163in}{1.411163in}}%
\pgfusepath{clip}%
\pgfsetbuttcap%
\pgfsetroundjoin%
\definecolor{currentfill}{rgb}{0.556863,0.654902,0.847059}%
\pgfsetfillcolor{currentfill}%
\pgfsetfillopacity{0.130000}%
\pgfsetlinewidth{0.000000pt}%
\definecolor{currentstroke}{rgb}{0.000000,0.000000,0.000000}%
\pgfsetstrokecolor{currentstroke}%
\pgfsetdash{}{0pt}%
\pgfpathmoveto{\pgfqpoint{1.377820in}{0.977789in}}%
\pgfpathlineto{\pgfqpoint{1.396494in}{0.988775in}}%
\pgfpathlineto{\pgfqpoint{1.415254in}{0.999983in}}%
\pgfpathlineto{\pgfqpoint{1.423569in}{0.998033in}}%
\pgfpathlineto{\pgfqpoint{1.431852in}{0.996304in}}%
\pgfpathlineto{\pgfqpoint{1.413130in}{0.984822in}}%
\pgfpathlineto{\pgfqpoint{1.394493in}{0.973552in}}%
\pgfpathlineto{\pgfqpoint{1.386173in}{0.975559in}}%
\pgfpathlineto{\pgfqpoint{1.377820in}{0.977789in}}%
\pgfpathclose%
\pgfusepath{fill}%
\end{pgfscope}%
\begin{pgfscope}%
\pgfpathrectangle{\pgfqpoint{0.407592in}{0.222904in}}{\pgfqpoint{1.411163in}{1.411163in}}%
\pgfusepath{clip}%
\pgfsetbuttcap%
\pgfsetroundjoin%
\definecolor{currentfill}{rgb}{0.556863,0.654902,0.847059}%
\pgfsetfillcolor{currentfill}%
\pgfsetfillopacity{0.130000}%
\pgfsetlinewidth{0.000000pt}%
\definecolor{currentstroke}{rgb}{0.000000,0.000000,0.000000}%
\pgfsetstrokecolor{currentstroke}%
\pgfsetdash{}{0pt}%
\pgfpathmoveto{\pgfqpoint{1.415254in}{0.999983in}}%
\pgfpathlineto{\pgfqpoint{1.434100in}{1.011403in}}%
\pgfpathlineto{\pgfqpoint{1.442396in}{1.009586in}}%
\pgfpathlineto{\pgfqpoint{1.450661in}{1.007987in}}%
\pgfpathlineto{\pgfqpoint{1.431852in}{0.996304in}}%
\pgfpathlineto{\pgfqpoint{1.423569in}{0.998033in}}%
\pgfusepath{fill}%
\end{pgfscope}%
\begin{pgfscope}%
\pgfpathrectangle{\pgfqpoint{0.407592in}{0.222904in}}{\pgfqpoint{1.411163in}{1.411163in}}%
\pgfusepath{clip}%
\pgfsetbuttcap%
\pgfsetroundjoin%
\definecolor{currentfill}{rgb}{0.909804,0.698039,0.423529}%
\pgfsetfillcolor{currentfill}%
\pgfsetfillopacity{0.120000}%
\pgfsetlinewidth{0.000000pt}%
\definecolor{currentstroke}{rgb}{0.000000,0.000000,0.000000}%
\pgfsetstrokecolor{currentstroke}%
\pgfsetdash{}{0pt}%
\pgfpathmoveto{\pgfqpoint{0.870800in}{1.442368in}}%
\pgfpathlineto{\pgfqpoint{0.886196in}{1.427129in}}%
\pgfpathlineto{\pgfqpoint{0.901605in}{1.412100in}}%
\pgfpathlineto{\pgfqpoint{0.909495in}{1.423095in}}%
\pgfpathlineto{\pgfqpoint{0.894100in}{1.437968in}}%
\pgfpathlineto{\pgfqpoint{0.878719in}{1.453053in}}%
\pgfusepath{fill}%
\end{pgfscope}%
\begin{pgfscope}%
\pgfpathrectangle{\pgfqpoint{0.407592in}{0.222904in}}{\pgfqpoint{1.411163in}{1.411163in}}%
\pgfusepath{clip}%
\pgfsetbuttcap%
\pgfsetroundjoin%
\definecolor{currentfill}{rgb}{0.909804,0.698039,0.423529}%
\pgfsetfillcolor{currentfill}%
\pgfsetfillopacity{0.120000}%
\pgfsetlinewidth{0.000000pt}%
\definecolor{currentstroke}{rgb}{0.000000,0.000000,0.000000}%
\pgfsetstrokecolor{currentstroke}%
\pgfsetdash{}{0pt}%
\pgfpathmoveto{\pgfqpoint{0.901605in}{1.412100in}}%
\pgfpathlineto{\pgfqpoint{0.917027in}{1.397294in}}%
\pgfpathlineto{\pgfqpoint{0.932464in}{1.382727in}}%
\pgfpathlineto{\pgfqpoint{0.940326in}{1.394042in}}%
\pgfpathlineto{\pgfqpoint{0.924904in}{1.408449in}}%
\pgfpathlineto{\pgfqpoint{0.909495in}{1.423095in}}%
\pgfusepath{fill}%
\end{pgfscope}%
\begin{pgfscope}%
\pgfpathrectangle{\pgfqpoint{0.407592in}{0.222904in}}{\pgfqpoint{1.411163in}{1.411163in}}%
\pgfusepath{clip}%
\pgfsetbuttcap%
\pgfsetroundjoin%
\definecolor{currentfill}{rgb}{0.909804,0.698039,0.423529}%
\pgfsetfillcolor{currentfill}%
\pgfsetfillopacity{0.120000}%
\pgfsetlinewidth{0.000000pt}%
\definecolor{currentstroke}{rgb}{0.000000,0.000000,0.000000}%
\pgfsetstrokecolor{currentstroke}%
\pgfsetdash{}{0pt}%
\pgfpathmoveto{\pgfqpoint{0.854899in}{1.421494in}}%
\pgfpathlineto{\pgfqpoint{0.870325in}{1.405949in}}%
\pgfpathlineto{\pgfqpoint{0.885764in}{1.390608in}}%
\pgfpathlineto{\pgfqpoint{0.893694in}{1.401268in}}%
\pgfpathlineto{\pgfqpoint{0.901605in}{1.412100in}}%
\pgfpathlineto{\pgfqpoint{0.886196in}{1.427129in}}%
\pgfpathlineto{\pgfqpoint{0.870800in}{1.442368in}}%
\pgfpathlineto{\pgfqpoint{0.862860in}{1.431846in}}%
\pgfpathlineto{\pgfqpoint{0.854899in}{1.421494in}}%
\pgfpathclose%
\pgfusepath{fill}%
\end{pgfscope}%
\begin{pgfscope}%
\pgfpathrectangle{\pgfqpoint{0.407592in}{0.222904in}}{\pgfqpoint{1.411163in}{1.411163in}}%
\pgfusepath{clip}%
\pgfsetbuttcap%
\pgfsetroundjoin%
\definecolor{currentfill}{rgb}{0.909804,0.698039,0.423529}%
\pgfsetfillcolor{currentfill}%
\pgfsetfillopacity{0.120000}%
\pgfsetlinewidth{0.000000pt}%
\definecolor{currentstroke}{rgb}{0.000000,0.000000,0.000000}%
\pgfsetstrokecolor{currentstroke}%
\pgfsetdash{}{0pt}%
\pgfpathmoveto{\pgfqpoint{0.932464in}{1.382727in}}%
\pgfpathlineto{\pgfqpoint{0.947915in}{1.368413in}}%
\pgfpathlineto{\pgfqpoint{0.963382in}{1.354369in}}%
\pgfpathlineto{\pgfqpoint{0.971217in}{1.366012in}}%
\pgfpathlineto{\pgfqpoint{0.955764in}{1.379892in}}%
\pgfpathlineto{\pgfqpoint{0.940326in}{1.394042in}}%
\pgfusepath{fill}%
\end{pgfscope}%
\begin{pgfscope}%
\pgfpathrectangle{\pgfqpoint{0.407592in}{0.222904in}}{\pgfqpoint{1.411163in}{1.411163in}}%
\pgfusepath{clip}%
\pgfsetbuttcap%
\pgfsetroundjoin%
\definecolor{currentfill}{rgb}{0.909804,0.698039,0.423529}%
\pgfsetfillcolor{currentfill}%
\pgfsetfillopacity{0.120000}%
\pgfsetlinewidth{0.000000pt}%
\definecolor{currentstroke}{rgb}{0.000000,0.000000,0.000000}%
\pgfsetstrokecolor{currentstroke}%
\pgfsetdash{}{0pt}%
\pgfpathmoveto{\pgfqpoint{0.885764in}{1.390608in}}%
\pgfpathlineto{\pgfqpoint{0.901216in}{1.375486in}}%
\pgfpathlineto{\pgfqpoint{0.916681in}{1.360597in}}%
\pgfpathlineto{\pgfqpoint{0.924582in}{1.371576in}}%
\pgfpathlineto{\pgfqpoint{0.932464in}{1.382727in}}%
\pgfpathlineto{\pgfqpoint{0.917027in}{1.397294in}}%
\pgfpathlineto{\pgfqpoint{0.901605in}{1.412100in}}%
\pgfpathlineto{\pgfqpoint{0.893694in}{1.401268in}}%
\pgfpathlineto{\pgfqpoint{0.885764in}{1.390608in}}%
\pgfpathclose%
\pgfusepath{fill}%
\end{pgfscope}%
\begin{pgfscope}%
\pgfpathrectangle{\pgfqpoint{0.407592in}{0.222904in}}{\pgfqpoint{1.411163in}{1.411163in}}%
\pgfusepath{clip}%
\pgfsetbuttcap%
\pgfsetroundjoin%
\definecolor{currentfill}{rgb}{0.909804,0.698039,0.423529}%
\pgfsetfillcolor{currentfill}%
\pgfsetfillopacity{0.120000}%
\pgfsetlinewidth{0.000000pt}%
\definecolor{currentstroke}{rgb}{0.000000,0.000000,0.000000}%
\pgfsetstrokecolor{currentstroke}%
\pgfsetdash{}{0pt}%
\pgfpathmoveto{\pgfqpoint{0.963382in}{1.354369in}}%
\pgfpathlineto{\pgfqpoint{0.978864in}{1.340612in}}%
\pgfpathlineto{\pgfqpoint{0.994363in}{1.327159in}}%
\pgfpathlineto{\pgfqpoint{1.002172in}{1.339134in}}%
\pgfpathlineto{\pgfqpoint{0.986686in}{1.352420in}}%
\pgfpathlineto{\pgfqpoint{0.971217in}{1.366012in}}%
\pgfusepath{fill}%
\end{pgfscope}%
\begin{pgfscope}%
\pgfpathrectangle{\pgfqpoint{0.407592in}{0.222904in}}{\pgfqpoint{1.411163in}{1.411163in}}%
\pgfusepath{clip}%
\pgfsetbuttcap%
\pgfsetroundjoin%
\definecolor{currentfill}{rgb}{0.909804,0.698039,0.423529}%
\pgfsetfillcolor{currentfill}%
\pgfsetfillopacity{0.120000}%
\pgfsetlinewidth{0.000000pt}%
\definecolor{currentstroke}{rgb}{0.000000,0.000000,0.000000}%
\pgfsetstrokecolor{currentstroke}%
\pgfsetdash{}{0pt}%
\pgfpathmoveto{\pgfqpoint{0.916681in}{1.360597in}}%
\pgfpathlineto{\pgfqpoint{0.932161in}{1.345958in}}%
\pgfpathlineto{\pgfqpoint{0.947655in}{1.331584in}}%
\pgfpathlineto{\pgfqpoint{0.955528in}{1.342890in}}%
\pgfpathlineto{\pgfqpoint{0.963382in}{1.354369in}}%
\pgfpathlineto{\pgfqpoint{0.947915in}{1.368413in}}%
\pgfpathlineto{\pgfqpoint{0.932464in}{1.382727in}}%
\pgfpathlineto{\pgfqpoint{0.924582in}{1.371576in}}%
\pgfpathlineto{\pgfqpoint{0.916681in}{1.360597in}}%
\pgfpathclose%
\pgfusepath{fill}%
\end{pgfscope}%
\begin{pgfscope}%
\pgfpathrectangle{\pgfqpoint{0.407592in}{0.222904in}}{\pgfqpoint{1.411163in}{1.411163in}}%
\pgfusepath{clip}%
\pgfsetbuttcap%
\pgfsetroundjoin%
\definecolor{currentfill}{rgb}{0.909804,0.698039,0.423529}%
\pgfsetfillcolor{currentfill}%
\pgfsetfillopacity{0.120000}%
\pgfsetlinewidth{0.000000pt}%
\definecolor{currentstroke}{rgb}{0.000000,0.000000,0.000000}%
\pgfsetstrokecolor{currentstroke}%
\pgfsetdash{}{0pt}%
\pgfpathmoveto{\pgfqpoint{0.838913in}{1.401324in}}%
\pgfpathlineto{\pgfqpoint{0.854371in}{1.385477in}}%
\pgfpathlineto{\pgfqpoint{0.869840in}{1.369829in}}%
\pgfpathlineto{\pgfqpoint{0.877813in}{1.380126in}}%
\pgfpathlineto{\pgfqpoint{0.885764in}{1.390608in}}%
\pgfpathlineto{\pgfqpoint{0.870325in}{1.405949in}}%
\pgfpathlineto{\pgfqpoint{0.854899in}{1.421494in}}%
\pgfpathlineto{\pgfqpoint{0.846917in}{1.411317in}}%
\pgfpathlineto{\pgfqpoint{0.838913in}{1.401324in}}%
\pgfpathclose%
\pgfusepath{fill}%
\end{pgfscope}%
\begin{pgfscope}%
\pgfpathrectangle{\pgfqpoint{0.407592in}{0.222904in}}{\pgfqpoint{1.411163in}{1.411163in}}%
\pgfusepath{clip}%
\pgfsetbuttcap%
\pgfsetroundjoin%
\definecolor{currentfill}{rgb}{0.909804,0.698039,0.423529}%
\pgfsetfillcolor{currentfill}%
\pgfsetfillopacity{0.120000}%
\pgfsetlinewidth{0.000000pt}%
\definecolor{currentstroke}{rgb}{0.000000,0.000000,0.000000}%
\pgfsetstrokecolor{currentstroke}%
\pgfsetdash{}{0pt}%
\pgfpathmoveto{\pgfqpoint{0.994363in}{1.327159in}}%
\pgfpathlineto{\pgfqpoint{1.009879in}{1.314030in}}%
\pgfpathlineto{\pgfqpoint{1.025412in}{1.301244in}}%
\pgfpathlineto{\pgfqpoint{1.033198in}{1.313552in}}%
\pgfpathlineto{\pgfqpoint{1.017676in}{1.326171in}}%
\pgfpathlineto{\pgfqpoint{1.002172in}{1.339134in}}%
\pgfusepath{fill}%
\end{pgfscope}%
\begin{pgfscope}%
\pgfpathrectangle{\pgfqpoint{0.407592in}{0.222904in}}{\pgfqpoint{1.411163in}{1.411163in}}%
\pgfusepath{clip}%
\pgfsetbuttcap%
\pgfsetroundjoin%
\definecolor{currentfill}{rgb}{0.909804,0.698039,0.423529}%
\pgfsetfillcolor{currentfill}%
\pgfsetfillopacity{0.120000}%
\pgfsetlinewidth{0.000000pt}%
\definecolor{currentstroke}{rgb}{0.000000,0.000000,0.000000}%
\pgfsetstrokecolor{currentstroke}%
\pgfsetdash{}{0pt}%
\pgfpathmoveto{\pgfqpoint{0.869840in}{1.369829in}}%
\pgfpathlineto{\pgfqpoint{0.885323in}{1.354394in}}%
\pgfpathlineto{\pgfqpoint{0.900818in}{1.339187in}}%
\pgfpathlineto{\pgfqpoint{0.908760in}{1.349798in}}%
\pgfpathlineto{\pgfqpoint{0.916681in}{1.360597in}}%
\pgfpathlineto{\pgfqpoint{0.901216in}{1.375486in}}%
\pgfpathlineto{\pgfqpoint{0.885764in}{1.390608in}}%
\pgfpathlineto{\pgfqpoint{0.877813in}{1.380126in}}%
\pgfpathlineto{\pgfqpoint{0.869840in}{1.369829in}}%
\pgfpathclose%
\pgfusepath{fill}%
\end{pgfscope}%
\begin{pgfscope}%
\pgfpathrectangle{\pgfqpoint{0.407592in}{0.222904in}}{\pgfqpoint{1.411163in}{1.411163in}}%
\pgfusepath{clip}%
\pgfsetbuttcap%
\pgfsetroundjoin%
\definecolor{currentfill}{rgb}{0.909804,0.698039,0.423529}%
\pgfsetfillcolor{currentfill}%
\pgfsetfillopacity{0.120000}%
\pgfsetlinewidth{0.000000pt}%
\definecolor{currentstroke}{rgb}{0.000000,0.000000,0.000000}%
\pgfsetstrokecolor{currentstroke}%
\pgfsetdash{}{0pt}%
\pgfpathmoveto{\pgfqpoint{0.947655in}{1.331584in}}%
\pgfpathlineto{\pgfqpoint{0.963164in}{1.317494in}}%
\pgfpathlineto{\pgfqpoint{0.978689in}{1.303706in}}%
\pgfpathlineto{\pgfqpoint{0.986535in}{1.315347in}}%
\pgfpathlineto{\pgfqpoint{0.994363in}{1.327159in}}%
\pgfpathlineto{\pgfqpoint{0.978864in}{1.340612in}}%
\pgfpathlineto{\pgfqpoint{0.963382in}{1.354369in}}%
\pgfpathlineto{\pgfqpoint{0.955528in}{1.342890in}}%
\pgfpathlineto{\pgfqpoint{0.947655in}{1.331584in}}%
\pgfpathclose%
\pgfusepath{fill}%
\end{pgfscope}%
\begin{pgfscope}%
\pgfpathrectangle{\pgfqpoint{0.407592in}{0.222904in}}{\pgfqpoint{1.411163in}{1.411163in}}%
\pgfusepath{clip}%
\pgfsetbuttcap%
\pgfsetroundjoin%
\definecolor{currentfill}{rgb}{0.909804,0.698039,0.423529}%
\pgfsetfillcolor{currentfill}%
\pgfsetfillopacity{0.120000}%
\pgfsetlinewidth{0.000000pt}%
\definecolor{currentstroke}{rgb}{0.000000,0.000000,0.000000}%
\pgfsetstrokecolor{currentstroke}%
\pgfsetdash{}{0pt}%
\pgfpathmoveto{\pgfqpoint{1.025412in}{1.301244in}}%
\pgfpathlineto{\pgfqpoint{1.040965in}{1.288822in}}%
\pgfpathlineto{\pgfqpoint{1.056537in}{1.276784in}}%
\pgfpathlineto{\pgfqpoint{1.064300in}{1.289422in}}%
\pgfpathlineto{\pgfqpoint{1.048739in}{1.301295in}}%
\pgfpathlineto{\pgfqpoint{1.033198in}{1.313552in}}%
\pgfusepath{fill}%
\end{pgfscope}%
\begin{pgfscope}%
\pgfpathrectangle{\pgfqpoint{0.407592in}{0.222904in}}{\pgfqpoint{1.411163in}{1.411163in}}%
\pgfusepath{clip}%
\pgfsetbuttcap%
\pgfsetroundjoin%
\definecolor{currentfill}{rgb}{0.909804,0.698039,0.423529}%
\pgfsetfillcolor{currentfill}%
\pgfsetfillopacity{0.120000}%
\pgfsetlinewidth{0.000000pt}%
\definecolor{currentstroke}{rgb}{0.000000,0.000000,0.000000}%
\pgfsetstrokecolor{currentstroke}%
\pgfsetdash{}{0pt}%
\pgfpathmoveto{\pgfqpoint{0.900818in}{1.339187in}}%
\pgfpathlineto{\pgfqpoint{0.916327in}{1.324223in}}%
\pgfpathlineto{\pgfqpoint{0.931850in}{1.309521in}}%
\pgfpathlineto{\pgfqpoint{0.939762in}{1.320458in}}%
\pgfpathlineto{\pgfqpoint{0.947655in}{1.331584in}}%
\pgfpathlineto{\pgfqpoint{0.932161in}{1.345958in}}%
\pgfpathlineto{\pgfqpoint{0.916681in}{1.360597in}}%
\pgfpathlineto{\pgfqpoint{0.908760in}{1.349798in}}%
\pgfpathlineto{\pgfqpoint{0.900818in}{1.339187in}}%
\pgfpathclose%
\pgfusepath{fill}%
\end{pgfscope}%
\begin{pgfscope}%
\pgfpathrectangle{\pgfqpoint{0.407592in}{0.222904in}}{\pgfqpoint{1.411163in}{1.411163in}}%
\pgfusepath{clip}%
\pgfsetbuttcap%
\pgfsetroundjoin%
\definecolor{currentfill}{rgb}{0.909804,0.698039,0.423529}%
\pgfsetfillcolor{currentfill}%
\pgfsetfillopacity{0.120000}%
\pgfsetlinewidth{0.000000pt}%
\definecolor{currentstroke}{rgb}{0.000000,0.000000,0.000000}%
\pgfsetstrokecolor{currentstroke}%
\pgfsetdash{}{0pt}%
\pgfpathmoveto{\pgfqpoint{0.822839in}{1.381909in}}%
\pgfpathlineto{\pgfqpoint{0.838329in}{1.365768in}}%
\pgfpathlineto{\pgfqpoint{0.853831in}{1.349820in}}%
\pgfpathlineto{\pgfqpoint{0.861847in}{1.359724in}}%
\pgfpathlineto{\pgfqpoint{0.869840in}{1.369829in}}%
\pgfpathlineto{\pgfqpoint{0.854371in}{1.385477in}}%
\pgfpathlineto{\pgfqpoint{0.838913in}{1.401324in}}%
\pgfpathlineto{\pgfqpoint{0.830887in}{1.391518in}}%
\pgfpathlineto{\pgfqpoint{0.822839in}{1.381909in}}%
\pgfpathclose%
\pgfusepath{fill}%
\end{pgfscope}%
\begin{pgfscope}%
\pgfpathrectangle{\pgfqpoint{0.407592in}{0.222904in}}{\pgfqpoint{1.411163in}{1.411163in}}%
\pgfusepath{clip}%
\pgfsetbuttcap%
\pgfsetroundjoin%
\definecolor{currentfill}{rgb}{0.909804,0.698039,0.423529}%
\pgfsetfillcolor{currentfill}%
\pgfsetfillopacity{0.120000}%
\pgfsetlinewidth{0.000000pt}%
\definecolor{currentstroke}{rgb}{0.000000,0.000000,0.000000}%
\pgfsetstrokecolor{currentstroke}%
\pgfsetdash{}{0pt}%
\pgfpathmoveto{\pgfqpoint{0.978689in}{1.303706in}}%
\pgfpathlineto{\pgfqpoint{0.994230in}{1.290241in}}%
\pgfpathlineto{\pgfqpoint{1.009789in}{1.277119in}}%
\pgfpathlineto{\pgfqpoint{1.017609in}{1.289096in}}%
\pgfpathlineto{\pgfqpoint{1.025412in}{1.301244in}}%
\pgfpathlineto{\pgfqpoint{1.009879in}{1.314030in}}%
\pgfpathlineto{\pgfqpoint{0.994363in}{1.327159in}}%
\pgfpathlineto{\pgfqpoint{0.986535in}{1.315347in}}%
\pgfpathlineto{\pgfqpoint{0.978689in}{1.303706in}}%
\pgfpathclose%
\pgfusepath{fill}%
\end{pgfscope}%
\begin{pgfscope}%
\pgfpathrectangle{\pgfqpoint{0.407592in}{0.222904in}}{\pgfqpoint{1.411163in}{1.411163in}}%
\pgfusepath{clip}%
\pgfsetbuttcap%
\pgfsetroundjoin%
\definecolor{currentfill}{rgb}{0.909804,0.698039,0.423529}%
\pgfsetfillcolor{currentfill}%
\pgfsetfillopacity{0.120000}%
\pgfsetlinewidth{0.000000pt}%
\definecolor{currentstroke}{rgb}{0.000000,0.000000,0.000000}%
\pgfsetstrokecolor{currentstroke}%
\pgfsetdash{}{0pt}%
\pgfpathmoveto{\pgfqpoint{0.931850in}{1.309521in}}%
\pgfpathlineto{\pgfqpoint{0.947387in}{1.295099in}}%
\pgfpathlineto{\pgfqpoint{0.962940in}{1.280976in}}%
\pgfpathlineto{\pgfqpoint{0.970824in}{1.292246in}}%
\pgfpathlineto{\pgfqpoint{0.978689in}{1.303706in}}%
\pgfpathlineto{\pgfqpoint{0.963164in}{1.317494in}}%
\pgfpathlineto{\pgfqpoint{0.947655in}{1.331584in}}%
\pgfpathlineto{\pgfqpoint{0.939762in}{1.320458in}}%
\pgfpathlineto{\pgfqpoint{0.931850in}{1.309521in}}%
\pgfpathclose%
\pgfusepath{fill}%
\end{pgfscope}%
\begin{pgfscope}%
\pgfpathrectangle{\pgfqpoint{0.407592in}{0.222904in}}{\pgfqpoint{1.411163in}{1.411163in}}%
\pgfusepath{clip}%
\pgfsetbuttcap%
\pgfsetroundjoin%
\definecolor{currentfill}{rgb}{0.909804,0.698039,0.423529}%
\pgfsetfillcolor{currentfill}%
\pgfsetfillopacity{0.120000}%
\pgfsetlinewidth{0.000000pt}%
\definecolor{currentstroke}{rgb}{0.000000,0.000000,0.000000}%
\pgfsetstrokecolor{currentstroke}%
\pgfsetdash{}{0pt}%
\pgfpathmoveto{\pgfqpoint{0.853831in}{1.349820in}}%
\pgfpathlineto{\pgfqpoint{0.869345in}{1.334077in}}%
\pgfpathlineto{\pgfqpoint{0.884871in}{1.318557in}}%
\pgfpathlineto{\pgfqpoint{0.892856in}{1.328770in}}%
\pgfpathlineto{\pgfqpoint{0.900818in}{1.339187in}}%
\pgfpathlineto{\pgfqpoint{0.885323in}{1.354394in}}%
\pgfpathlineto{\pgfqpoint{0.869840in}{1.369829in}}%
\pgfpathlineto{\pgfqpoint{0.861847in}{1.359724in}}%
\pgfpathlineto{\pgfqpoint{0.853831in}{1.349820in}}%
\pgfpathclose%
\pgfusepath{fill}%
\end{pgfscope}%
\begin{pgfscope}%
\pgfpathrectangle{\pgfqpoint{0.407592in}{0.222904in}}{\pgfqpoint{1.411163in}{1.411163in}}%
\pgfusepath{clip}%
\pgfsetbuttcap%
\pgfsetroundjoin%
\definecolor{currentfill}{rgb}{0.909804,0.698039,0.423529}%
\pgfsetfillcolor{currentfill}%
\pgfsetfillopacity{0.120000}%
\pgfsetlinewidth{0.000000pt}%
\definecolor{currentstroke}{rgb}{0.000000,0.000000,0.000000}%
\pgfsetstrokecolor{currentstroke}%
\pgfsetdash{}{0pt}%
\pgfpathmoveto{\pgfqpoint{1.056537in}{1.276784in}}%
\pgfpathlineto{\pgfqpoint{1.072130in}{1.265154in}}%
\pgfpathlineto{\pgfqpoint{1.087744in}{1.253952in}}%
\pgfpathlineto{\pgfqpoint{1.095486in}{1.266909in}}%
\pgfpathlineto{\pgfqpoint{1.079882in}{1.277953in}}%
\pgfpathlineto{\pgfqpoint{1.064300in}{1.289422in}}%
\pgfusepath{fill}%
\end{pgfscope}%
\begin{pgfscope}%
\pgfpathrectangle{\pgfqpoint{0.407592in}{0.222904in}}{\pgfqpoint{1.411163in}{1.411163in}}%
\pgfusepath{clip}%
\pgfsetbuttcap%
\pgfsetroundjoin%
\definecolor{currentfill}{rgb}{0.909804,0.698039,0.423529}%
\pgfsetfillcolor{currentfill}%
\pgfsetfillopacity{0.120000}%
\pgfsetlinewidth{0.000000pt}%
\definecolor{currentstroke}{rgb}{0.000000,0.000000,0.000000}%
\pgfsetstrokecolor{currentstroke}%
\pgfsetdash{}{0pt}%
\pgfpathmoveto{\pgfqpoint{1.009789in}{1.277119in}}%
\pgfpathlineto{\pgfqpoint{1.025365in}{1.264361in}}%
\pgfpathlineto{\pgfqpoint{1.040961in}{1.251991in}}%
\pgfpathlineto{\pgfqpoint{1.048757in}{1.264304in}}%
\pgfpathlineto{\pgfqpoint{1.056537in}{1.276784in}}%
\pgfpathlineto{\pgfqpoint{1.040965in}{1.288822in}}%
\pgfpathlineto{\pgfqpoint{1.025412in}{1.301244in}}%
\pgfpathlineto{\pgfqpoint{1.017609in}{1.289096in}}%
\pgfpathlineto{\pgfqpoint{1.009789in}{1.277119in}}%
\pgfpathclose%
\pgfusepath{fill}%
\end{pgfscope}%
\begin{pgfscope}%
\pgfpathrectangle{\pgfqpoint{0.407592in}{0.222904in}}{\pgfqpoint{1.411163in}{1.411163in}}%
\pgfusepath{clip}%
\pgfsetbuttcap%
\pgfsetroundjoin%
\definecolor{currentfill}{rgb}{0.909804,0.698039,0.423529}%
\pgfsetfillcolor{currentfill}%
\pgfsetfillopacity{0.120000}%
\pgfsetlinewidth{0.000000pt}%
\definecolor{currentstroke}{rgb}{0.000000,0.000000,0.000000}%
\pgfsetstrokecolor{currentstroke}%
\pgfsetdash{}{0pt}%
\pgfpathmoveto{\pgfqpoint{0.884871in}{1.318557in}}%
\pgfpathlineto{\pgfqpoint{0.900411in}{1.303275in}}%
\pgfpathlineto{\pgfqpoint{0.915964in}{1.288249in}}%
\pgfpathlineto{\pgfqpoint{0.923917in}{1.298782in}}%
\pgfpathlineto{\pgfqpoint{0.931850in}{1.309521in}}%
\pgfpathlineto{\pgfqpoint{0.916327in}{1.324223in}}%
\pgfpathlineto{\pgfqpoint{0.900818in}{1.339187in}}%
\pgfpathlineto{\pgfqpoint{0.892856in}{1.328770in}}%
\pgfpathlineto{\pgfqpoint{0.884871in}{1.318557in}}%
\pgfpathclose%
\pgfusepath{fill}%
\end{pgfscope}%
\begin{pgfscope}%
\pgfpathrectangle{\pgfqpoint{0.407592in}{0.222904in}}{\pgfqpoint{1.411163in}{1.411163in}}%
\pgfusepath{clip}%
\pgfsetbuttcap%
\pgfsetroundjoin%
\definecolor{currentfill}{rgb}{0.909804,0.698039,0.423529}%
\pgfsetfillcolor{currentfill}%
\pgfsetfillopacity{0.120000}%
\pgfsetlinewidth{0.000000pt}%
\definecolor{currentstroke}{rgb}{0.000000,0.000000,0.000000}%
\pgfsetstrokecolor{currentstroke}%
\pgfsetdash{}{0pt}%
\pgfpathmoveto{\pgfqpoint{0.962940in}{1.280976in}}%
\pgfpathlineto{\pgfqpoint{0.978508in}{1.267173in}}%
\pgfpathlineto{\pgfqpoint{0.994092in}{1.253712in}}%
\pgfpathlineto{\pgfqpoint{1.001950in}{1.265321in}}%
\pgfpathlineto{\pgfqpoint{1.009789in}{1.277119in}}%
\pgfpathlineto{\pgfqpoint{0.994230in}{1.290241in}}%
\pgfpathlineto{\pgfqpoint{0.978689in}{1.303706in}}%
\pgfpathlineto{\pgfqpoint{0.970824in}{1.292246in}}%
\pgfpathlineto{\pgfqpoint{0.962940in}{1.280976in}}%
\pgfpathclose%
\pgfusepath{fill}%
\end{pgfscope}%
\begin{pgfscope}%
\pgfpathrectangle{\pgfqpoint{0.407592in}{0.222904in}}{\pgfqpoint{1.411163in}{1.411163in}}%
\pgfusepath{clip}%
\pgfsetbuttcap%
\pgfsetroundjoin%
\definecolor{currentfill}{rgb}{0.909804,0.698039,0.423529}%
\pgfsetfillcolor{currentfill}%
\pgfsetfillopacity{0.120000}%
\pgfsetlinewidth{0.000000pt}%
\definecolor{currentstroke}{rgb}{0.000000,0.000000,0.000000}%
\pgfsetstrokecolor{currentstroke}%
\pgfsetdash{}{0pt}%
\pgfpathmoveto{\pgfqpoint{0.806672in}{1.363303in}}%
\pgfpathlineto{\pgfqpoint{0.822196in}{1.346878in}}%
\pgfpathlineto{\pgfqpoint{0.837731in}{1.330638in}}%
\pgfpathlineto{\pgfqpoint{0.845793in}{1.340122in}}%
\pgfpathlineto{\pgfqpoint{0.853831in}{1.349820in}}%
\pgfpathlineto{\pgfqpoint{0.838329in}{1.365768in}}%
\pgfpathlineto{\pgfqpoint{0.822839in}{1.381909in}}%
\pgfpathlineto{\pgfqpoint{0.814767in}{1.372501in}}%
\pgfpathlineto{\pgfqpoint{0.806672in}{1.363303in}}%
\pgfpathclose%
\pgfusepath{fill}%
\end{pgfscope}%
\begin{pgfscope}%
\pgfpathrectangle{\pgfqpoint{0.407592in}{0.222904in}}{\pgfqpoint{1.411163in}{1.411163in}}%
\pgfusepath{clip}%
\pgfsetbuttcap%
\pgfsetroundjoin%
\definecolor{currentfill}{rgb}{0.909804,0.698039,0.423529}%
\pgfsetfillcolor{currentfill}%
\pgfsetfillopacity{0.120000}%
\pgfsetlinewidth{0.000000pt}%
\definecolor{currentstroke}{rgb}{0.000000,0.000000,0.000000}%
\pgfsetstrokecolor{currentstroke}%
\pgfsetdash{}{0pt}%
\pgfpathmoveto{\pgfqpoint{1.087744in}{1.253952in}}%
\pgfpathlineto{\pgfqpoint{1.103382in}{1.243201in}}%
\pgfpathlineto{\pgfqpoint{1.119043in}{1.232925in}}%
\pgfpathlineto{\pgfqpoint{1.126765in}{1.246184in}}%
\pgfpathlineto{\pgfqpoint{1.111113in}{1.256312in}}%
\pgfpathlineto{\pgfqpoint{1.095486in}{1.266909in}}%
\pgfusepath{fill}%
\end{pgfscope}%
\begin{pgfscope}%
\pgfpathrectangle{\pgfqpoint{0.407592in}{0.222904in}}{\pgfqpoint{1.411163in}{1.411163in}}%
\pgfusepath{clip}%
\pgfsetbuttcap%
\pgfsetroundjoin%
\definecolor{currentfill}{rgb}{0.909804,0.698039,0.423529}%
\pgfsetfillcolor{currentfill}%
\pgfsetfillopacity{0.120000}%
\pgfsetlinewidth{0.000000pt}%
\definecolor{currentstroke}{rgb}{0.000000,0.000000,0.000000}%
\pgfsetstrokecolor{currentstroke}%
\pgfsetdash{}{0pt}%
\pgfpathmoveto{\pgfqpoint{1.040961in}{1.251991in}}%
\pgfpathlineto{\pgfqpoint{1.056576in}{1.240032in}}%
\pgfpathlineto{\pgfqpoint{1.072212in}{1.228508in}}%
\pgfpathlineto{\pgfqpoint{1.079986in}{1.241148in}}%
\pgfpathlineto{\pgfqpoint{1.087744in}{1.253952in}}%
\pgfpathlineto{\pgfqpoint{1.072130in}{1.265154in}}%
\pgfpathlineto{\pgfqpoint{1.056537in}{1.276784in}}%
\pgfpathlineto{\pgfqpoint{1.048757in}{1.264304in}}%
\pgfpathlineto{\pgfqpoint{1.040961in}{1.251991in}}%
\pgfpathclose%
\pgfusepath{fill}%
\end{pgfscope}%
\begin{pgfscope}%
\pgfpathrectangle{\pgfqpoint{0.407592in}{0.222904in}}{\pgfqpoint{1.411163in}{1.411163in}}%
\pgfusepath{clip}%
\pgfsetbuttcap%
\pgfsetroundjoin%
\definecolor{currentfill}{rgb}{0.909804,0.698039,0.423529}%
\pgfsetfillcolor{currentfill}%
\pgfsetfillopacity{0.120000}%
\pgfsetlinewidth{0.000000pt}%
\definecolor{currentstroke}{rgb}{0.000000,0.000000,0.000000}%
\pgfsetstrokecolor{currentstroke}%
\pgfsetdash{}{0pt}%
\pgfpathmoveto{\pgfqpoint{0.915964in}{1.288249in}}%
\pgfpathlineto{\pgfqpoint{0.931531in}{1.273497in}}%
\pgfpathlineto{\pgfqpoint{0.947112in}{1.259041in}}%
\pgfpathlineto{\pgfqpoint{0.955036in}{1.269904in}}%
\pgfpathlineto{\pgfqpoint{0.962940in}{1.280976in}}%
\pgfpathlineto{\pgfqpoint{0.947387in}{1.295099in}}%
\pgfpathlineto{\pgfqpoint{0.931850in}{1.309521in}}%
\pgfpathlineto{\pgfqpoint{0.923917in}{1.298782in}}%
\pgfpathlineto{\pgfqpoint{0.915964in}{1.288249in}}%
\pgfpathclose%
\pgfusepath{fill}%
\end{pgfscope}%
\begin{pgfscope}%
\pgfpathrectangle{\pgfqpoint{0.407592in}{0.222904in}}{\pgfqpoint{1.411163in}{1.411163in}}%
\pgfusepath{clip}%
\pgfsetbuttcap%
\pgfsetroundjoin%
\definecolor{currentfill}{rgb}{0.909804,0.698039,0.423529}%
\pgfsetfillcolor{currentfill}%
\pgfsetfillopacity{0.120000}%
\pgfsetlinewidth{0.000000pt}%
\definecolor{currentstroke}{rgb}{0.000000,0.000000,0.000000}%
\pgfsetstrokecolor{currentstroke}%
\pgfsetdash{}{0pt}%
\pgfpathmoveto{\pgfqpoint{0.837731in}{1.330638in}}%
\pgfpathlineto{\pgfqpoint{0.853278in}{1.314598in}}%
\pgfpathlineto{\pgfqpoint{0.868837in}{1.298773in}}%
\pgfpathlineto{\pgfqpoint{0.876866in}{1.308555in}}%
\pgfpathlineto{\pgfqpoint{0.884871in}{1.318557in}}%
\pgfpathlineto{\pgfqpoint{0.869345in}{1.334077in}}%
\pgfpathlineto{\pgfqpoint{0.853831in}{1.349820in}}%
\pgfpathlineto{\pgfqpoint{0.845793in}{1.340122in}}%
\pgfpathlineto{\pgfqpoint{0.837731in}{1.330638in}}%
\pgfpathclose%
\pgfusepath{fill}%
\end{pgfscope}%
\begin{pgfscope}%
\pgfpathrectangle{\pgfqpoint{0.407592in}{0.222904in}}{\pgfqpoint{1.411163in}{1.411163in}}%
\pgfusepath{clip}%
\pgfsetbuttcap%
\pgfsetroundjoin%
\definecolor{currentfill}{rgb}{0.909804,0.698039,0.423529}%
\pgfsetfillcolor{currentfill}%
\pgfsetfillopacity{0.120000}%
\pgfsetlinewidth{0.000000pt}%
\definecolor{currentstroke}{rgb}{0.000000,0.000000,0.000000}%
\pgfsetstrokecolor{currentstroke}%
\pgfsetdash{}{0pt}%
\pgfpathmoveto{\pgfqpoint{0.994092in}{1.253712in}}%
\pgfpathlineto{\pgfqpoint{1.009694in}{1.240615in}}%
\pgfpathlineto{\pgfqpoint{1.025314in}{1.227908in}}%
\pgfpathlineto{\pgfqpoint{1.033146in}{1.239856in}}%
\pgfpathlineto{\pgfqpoint{1.040961in}{1.251991in}}%
\pgfpathlineto{\pgfqpoint{1.025365in}{1.264361in}}%
\pgfpathlineto{\pgfqpoint{1.009789in}{1.277119in}}%
\pgfpathlineto{\pgfqpoint{1.001950in}{1.265321in}}%
\pgfpathlineto{\pgfqpoint{0.994092in}{1.253712in}}%
\pgfpathclose%
\pgfusepath{fill}%
\end{pgfscope}%
\begin{pgfscope}%
\pgfpathrectangle{\pgfqpoint{0.407592in}{0.222904in}}{\pgfqpoint{1.411163in}{1.411163in}}%
\pgfusepath{clip}%
\pgfsetbuttcap%
\pgfsetroundjoin%
\definecolor{currentfill}{rgb}{0.909804,0.698039,0.423529}%
\pgfsetfillcolor{currentfill}%
\pgfsetfillopacity{0.120000}%
\pgfsetlinewidth{0.000000pt}%
\definecolor{currentstroke}{rgb}{0.000000,0.000000,0.000000}%
\pgfsetstrokecolor{currentstroke}%
\pgfsetdash{}{0pt}%
\pgfpathmoveto{\pgfqpoint{1.072212in}{1.228508in}}%
\pgfpathlineto{\pgfqpoint{1.087870in}{1.217443in}}%
\pgfpathlineto{\pgfqpoint{1.103552in}{1.206863in}}%
\pgfpathlineto{\pgfqpoint{1.111306in}{1.219814in}}%
\pgfpathlineto{\pgfqpoint{1.119043in}{1.232925in}}%
\pgfpathlineto{\pgfqpoint{1.103382in}{1.243201in}}%
\pgfpathlineto{\pgfqpoint{1.087744in}{1.253952in}}%
\pgfpathlineto{\pgfqpoint{1.079986in}{1.241148in}}%
\pgfpathlineto{\pgfqpoint{1.072212in}{1.228508in}}%
\pgfpathclose%
\pgfusepath{fill}%
\end{pgfscope}%
\begin{pgfscope}%
\pgfpathrectangle{\pgfqpoint{0.407592in}{0.222904in}}{\pgfqpoint{1.411163in}{1.411163in}}%
\pgfusepath{clip}%
\pgfsetbuttcap%
\pgfsetroundjoin%
\definecolor{currentfill}{rgb}{0.909804,0.698039,0.423529}%
\pgfsetfillcolor{currentfill}%
\pgfsetfillopacity{0.120000}%
\pgfsetlinewidth{0.000000pt}%
\definecolor{currentstroke}{rgb}{0.000000,0.000000,0.000000}%
\pgfsetstrokecolor{currentstroke}%
\pgfsetdash{}{0pt}%
\pgfpathmoveto{\pgfqpoint{0.947112in}{1.259041in}}%
\pgfpathlineto{\pgfqpoint{0.962708in}{1.244902in}}%
\pgfpathlineto{\pgfqpoint{0.978320in}{1.231102in}}%
\pgfpathlineto{\pgfqpoint{0.986216in}{1.242302in}}%
\pgfpathlineto{\pgfqpoint{0.994092in}{1.253712in}}%
\pgfpathlineto{\pgfqpoint{0.978508in}{1.267173in}}%
\pgfpathlineto{\pgfqpoint{0.962940in}{1.280976in}}%
\pgfpathlineto{\pgfqpoint{0.955036in}{1.269904in}}%
\pgfpathlineto{\pgfqpoint{0.947112in}{1.259041in}}%
\pgfpathclose%
\pgfusepath{fill}%
\end{pgfscope}%
\begin{pgfscope}%
\pgfpathrectangle{\pgfqpoint{0.407592in}{0.222904in}}{\pgfqpoint{1.411163in}{1.411163in}}%
\pgfusepath{clip}%
\pgfsetbuttcap%
\pgfsetroundjoin%
\definecolor{currentfill}{rgb}{0.909804,0.698039,0.423529}%
\pgfsetfillcolor{currentfill}%
\pgfsetfillopacity{0.120000}%
\pgfsetlinewidth{0.000000pt}%
\definecolor{currentstroke}{rgb}{0.000000,0.000000,0.000000}%
\pgfsetstrokecolor{currentstroke}%
\pgfsetdash{}{0pt}%
\pgfpathmoveto{\pgfqpoint{0.868837in}{1.298773in}}%
\pgfpathlineto{\pgfqpoint{0.884409in}{1.283180in}}%
\pgfpathlineto{\pgfqpoint{0.899993in}{1.267837in}}%
\pgfpathlineto{\pgfqpoint{0.907989in}{1.277931in}}%
\pgfpathlineto{\pgfqpoint{0.915964in}{1.288249in}}%
\pgfpathlineto{\pgfqpoint{0.900411in}{1.303275in}}%
\pgfpathlineto{\pgfqpoint{0.884871in}{1.318557in}}%
\pgfpathlineto{\pgfqpoint{0.876866in}{1.308555in}}%
\pgfpathlineto{\pgfqpoint{0.868837in}{1.298773in}}%
\pgfpathclose%
\pgfusepath{fill}%
\end{pgfscope}%
\begin{pgfscope}%
\pgfpathrectangle{\pgfqpoint{0.407592in}{0.222904in}}{\pgfqpoint{1.411163in}{1.411163in}}%
\pgfusepath{clip}%
\pgfsetbuttcap%
\pgfsetroundjoin%
\definecolor{currentfill}{rgb}{0.909804,0.698039,0.423529}%
\pgfsetfillcolor{currentfill}%
\pgfsetfillopacity{0.120000}%
\pgfsetlinewidth{0.000000pt}%
\definecolor{currentstroke}{rgb}{0.000000,0.000000,0.000000}%
\pgfsetstrokecolor{currentstroke}%
\pgfsetdash{}{0pt}%
\pgfpathmoveto{\pgfqpoint{1.119043in}{1.232925in}}%
\pgfpathlineto{\pgfqpoint{1.134730in}{1.223145in}}%
\pgfpathlineto{\pgfqpoint{1.150444in}{1.213884in}}%
\pgfpathlineto{\pgfqpoint{1.158148in}{1.227420in}}%
\pgfpathlineto{\pgfqpoint{1.142443in}{1.236547in}}%
\pgfpathlineto{\pgfqpoint{1.126765in}{1.246184in}}%
\pgfusepath{fill}%
\end{pgfscope}%
\begin{pgfscope}%
\pgfpathrectangle{\pgfqpoint{0.407592in}{0.222904in}}{\pgfqpoint{1.411163in}{1.411163in}}%
\pgfusepath{clip}%
\pgfsetbuttcap%
\pgfsetroundjoin%
\definecolor{currentfill}{rgb}{0.909804,0.698039,0.423529}%
\pgfsetfillcolor{currentfill}%
\pgfsetfillopacity{0.120000}%
\pgfsetlinewidth{0.000000pt}%
\definecolor{currentstroke}{rgb}{0.000000,0.000000,0.000000}%
\pgfsetstrokecolor{currentstroke}%
\pgfsetdash{}{0pt}%
\pgfpathmoveto{\pgfqpoint{0.790410in}{1.345557in}}%
\pgfpathlineto{\pgfqpoint{0.805968in}{1.328861in}}%
\pgfpathlineto{\pgfqpoint{0.821538in}{1.312343in}}%
\pgfpathlineto{\pgfqpoint{0.829647in}{1.321376in}}%
\pgfpathlineto{\pgfqpoint{0.837731in}{1.330638in}}%
\pgfpathlineto{\pgfqpoint{0.822196in}{1.346878in}}%
\pgfpathlineto{\pgfqpoint{0.806672in}{1.363303in}}%
\pgfpathlineto{\pgfqpoint{0.798553in}{1.354319in}}%
\pgfpathlineto{\pgfqpoint{0.790410in}{1.345557in}}%
\pgfpathclose%
\pgfusepath{fill}%
\end{pgfscope}%
\begin{pgfscope}%
\pgfpathrectangle{\pgfqpoint{0.407592in}{0.222904in}}{\pgfqpoint{1.411163in}{1.411163in}}%
\pgfusepath{clip}%
\pgfsetbuttcap%
\pgfsetroundjoin%
\definecolor{currentfill}{rgb}{0.909804,0.698039,0.423529}%
\pgfsetfillcolor{currentfill}%
\pgfsetfillopacity{0.120000}%
\pgfsetlinewidth{0.000000pt}%
\definecolor{currentstroke}{rgb}{0.000000,0.000000,0.000000}%
\pgfsetstrokecolor{currentstroke}%
\pgfsetdash{}{0pt}%
\pgfpathmoveto{\pgfqpoint{1.025314in}{1.227908in}}%
\pgfpathlineto{\pgfqpoint{1.040953in}{1.215615in}}%
\pgfpathlineto{\pgfqpoint{1.056612in}{1.203762in}}%
\pgfpathlineto{\pgfqpoint{1.064421in}{1.216042in}}%
\pgfpathlineto{\pgfqpoint{1.072212in}{1.228508in}}%
\pgfpathlineto{\pgfqpoint{1.056576in}{1.240032in}}%
\pgfpathlineto{\pgfqpoint{1.040961in}{1.251991in}}%
\pgfpathlineto{\pgfqpoint{1.033146in}{1.239856in}}%
\pgfpathlineto{\pgfqpoint{1.025314in}{1.227908in}}%
\pgfpathclose%
\pgfusepath{fill}%
\end{pgfscope}%
\begin{pgfscope}%
\pgfpathrectangle{\pgfqpoint{0.407592in}{0.222904in}}{\pgfqpoint{1.411163in}{1.411163in}}%
\pgfusepath{clip}%
\pgfsetbuttcap%
\pgfsetroundjoin%
\definecolor{currentfill}{rgb}{0.909804,0.698039,0.423529}%
\pgfsetfillcolor{currentfill}%
\pgfsetfillopacity{0.120000}%
\pgfsetlinewidth{0.000000pt}%
\definecolor{currentstroke}{rgb}{0.000000,0.000000,0.000000}%
\pgfsetstrokecolor{currentstroke}%
\pgfsetdash{}{0pt}%
\pgfpathmoveto{\pgfqpoint{0.899993in}{1.267837in}}%
\pgfpathlineto{\pgfqpoint{0.915590in}{1.252764in}}%
\pgfpathlineto{\pgfqpoint{0.931202in}{1.237981in}}%
\pgfpathlineto{\pgfqpoint{0.939167in}{1.248397in}}%
\pgfpathlineto{\pgfqpoint{0.947112in}{1.259041in}}%
\pgfpathlineto{\pgfqpoint{0.931531in}{1.273497in}}%
\pgfpathlineto{\pgfqpoint{0.915964in}{1.288249in}}%
\pgfpathlineto{\pgfqpoint{0.907989in}{1.277931in}}%
\pgfpathlineto{\pgfqpoint{0.899993in}{1.267837in}}%
\pgfpathclose%
\pgfusepath{fill}%
\end{pgfscope}%
\begin{pgfscope}%
\pgfpathrectangle{\pgfqpoint{0.407592in}{0.222904in}}{\pgfqpoint{1.411163in}{1.411163in}}%
\pgfusepath{clip}%
\pgfsetbuttcap%
\pgfsetroundjoin%
\definecolor{currentfill}{rgb}{0.909804,0.698039,0.423529}%
\pgfsetfillcolor{currentfill}%
\pgfsetfillopacity{0.120000}%
\pgfsetlinewidth{0.000000pt}%
\definecolor{currentstroke}{rgb}{0.000000,0.000000,0.000000}%
\pgfsetstrokecolor{currentstroke}%
\pgfsetdash{}{0pt}%
\pgfpathmoveto{\pgfqpoint{0.978320in}{1.231102in}}%
\pgfpathlineto{\pgfqpoint{0.993948in}{1.217667in}}%
\pgfpathlineto{\pgfqpoint{1.009594in}{1.204621in}}%
\pgfpathlineto{\pgfqpoint{1.017463in}{1.216159in}}%
\pgfpathlineto{\pgfqpoint{1.025314in}{1.227908in}}%
\pgfpathlineto{\pgfqpoint{1.009694in}{1.240615in}}%
\pgfpathlineto{\pgfqpoint{0.994092in}{1.253712in}}%
\pgfpathlineto{\pgfqpoint{0.986216in}{1.242302in}}%
\pgfpathlineto{\pgfqpoint{0.978320in}{1.231102in}}%
\pgfpathclose%
\pgfusepath{fill}%
\end{pgfscope}%
\begin{pgfscope}%
\pgfpathrectangle{\pgfqpoint{0.407592in}{0.222904in}}{\pgfqpoint{1.411163in}{1.411163in}}%
\pgfusepath{clip}%
\pgfsetbuttcap%
\pgfsetroundjoin%
\definecolor{currentfill}{rgb}{0.909804,0.698039,0.423529}%
\pgfsetfillcolor{currentfill}%
\pgfsetfillopacity{0.120000}%
\pgfsetlinewidth{0.000000pt}%
\definecolor{currentstroke}{rgb}{0.000000,0.000000,0.000000}%
\pgfsetstrokecolor{currentstroke}%
\pgfsetdash{}{0pt}%
\pgfpathmoveto{\pgfqpoint{0.821538in}{1.312343in}}%
\pgfpathlineto{\pgfqpoint{0.837119in}{1.296017in}}%
\pgfpathlineto{\pgfqpoint{0.852711in}{1.279900in}}%
\pgfpathlineto{\pgfqpoint{0.860786in}{1.289218in}}%
\pgfpathlineto{\pgfqpoint{0.868837in}{1.298773in}}%
\pgfpathlineto{\pgfqpoint{0.853278in}{1.314598in}}%
\pgfpathlineto{\pgfqpoint{0.837731in}{1.330638in}}%
\pgfpathlineto{\pgfqpoint{0.829647in}{1.321376in}}%
\pgfpathlineto{\pgfqpoint{0.821538in}{1.312343in}}%
\pgfpathclose%
\pgfusepath{fill}%
\end{pgfscope}%
\begin{pgfscope}%
\pgfpathrectangle{\pgfqpoint{0.407592in}{0.222904in}}{\pgfqpoint{1.411163in}{1.411163in}}%
\pgfusepath{clip}%
\pgfsetbuttcap%
\pgfsetroundjoin%
\definecolor{currentfill}{rgb}{0.909804,0.698039,0.423529}%
\pgfsetfillcolor{currentfill}%
\pgfsetfillopacity{0.120000}%
\pgfsetlinewidth{0.000000pt}%
\definecolor{currentstroke}{rgb}{0.000000,0.000000,0.000000}%
\pgfsetstrokecolor{currentstroke}%
\pgfsetdash{}{0pt}%
\pgfpathmoveto{\pgfqpoint{1.103552in}{1.206863in}}%
\pgfpathlineto{\pgfqpoint{1.119258in}{1.196791in}}%
\pgfpathlineto{\pgfqpoint{1.134991in}{1.187253in}}%
\pgfpathlineto{\pgfqpoint{1.142725in}{1.200491in}}%
\pgfpathlineto{\pgfqpoint{1.150444in}{1.213884in}}%
\pgfpathlineto{\pgfqpoint{1.134730in}{1.223145in}}%
\pgfpathlineto{\pgfqpoint{1.119043in}{1.232925in}}%
\pgfpathlineto{\pgfqpoint{1.111306in}{1.219814in}}%
\pgfpathlineto{\pgfqpoint{1.103552in}{1.206863in}}%
\pgfpathclose%
\pgfusepath{fill}%
\end{pgfscope}%
\begin{pgfscope}%
\pgfpathrectangle{\pgfqpoint{0.407592in}{0.222904in}}{\pgfqpoint{1.411163in}{1.411163in}}%
\pgfusepath{clip}%
\pgfsetbuttcap%
\pgfsetroundjoin%
\definecolor{currentfill}{rgb}{0.909804,0.698039,0.423529}%
\pgfsetfillcolor{currentfill}%
\pgfsetfillopacity{0.120000}%
\pgfsetlinewidth{0.000000pt}%
\definecolor{currentstroke}{rgb}{0.000000,0.000000,0.000000}%
\pgfsetstrokecolor{currentstroke}%
\pgfsetdash{}{0pt}%
\pgfpathmoveto{\pgfqpoint{1.150444in}{1.213884in}}%
\pgfpathlineto{\pgfqpoint{1.166187in}{1.205163in}}%
\pgfpathlineto{\pgfqpoint{1.181959in}{1.197003in}}%
\pgfpathlineto{\pgfqpoint{1.189647in}{1.210783in}}%
\pgfpathlineto{\pgfqpoint{1.173882in}{1.218826in}}%
\pgfpathlineto{\pgfqpoint{1.158148in}{1.227420in}}%
\pgfusepath{fill}%
\end{pgfscope}%
\begin{pgfscope}%
\pgfpathrectangle{\pgfqpoint{0.407592in}{0.222904in}}{\pgfqpoint{1.411163in}{1.411163in}}%
\pgfusepath{clip}%
\pgfsetbuttcap%
\pgfsetroundjoin%
\definecolor{currentfill}{rgb}{0.909804,0.698039,0.423529}%
\pgfsetfillcolor{currentfill}%
\pgfsetfillopacity{0.120000}%
\pgfsetlinewidth{0.000000pt}%
\definecolor{currentstroke}{rgb}{0.000000,0.000000,0.000000}%
\pgfsetstrokecolor{currentstroke}%
\pgfsetdash{}{0pt}%
\pgfpathmoveto{\pgfqpoint{1.056612in}{1.203762in}}%
\pgfpathlineto{\pgfqpoint{1.072293in}{1.192376in}}%
\pgfpathlineto{\pgfqpoint{1.087995in}{1.181483in}}%
\pgfpathlineto{\pgfqpoint{1.095782in}{1.194082in}}%
\pgfpathlineto{\pgfqpoint{1.103552in}{1.206863in}}%
\pgfpathlineto{\pgfqpoint{1.087870in}{1.217443in}}%
\pgfpathlineto{\pgfqpoint{1.072212in}{1.228508in}}%
\pgfpathlineto{\pgfqpoint{1.064421in}{1.216042in}}%
\pgfpathlineto{\pgfqpoint{1.056612in}{1.203762in}}%
\pgfpathclose%
\pgfusepath{fill}%
\end{pgfscope}%
\begin{pgfscope}%
\pgfpathrectangle{\pgfqpoint{0.407592in}{0.222904in}}{\pgfqpoint{1.411163in}{1.411163in}}%
\pgfusepath{clip}%
\pgfsetbuttcap%
\pgfsetroundjoin%
\definecolor{currentfill}{rgb}{0.909804,0.698039,0.423529}%
\pgfsetfillcolor{currentfill}%
\pgfsetfillopacity{0.120000}%
\pgfsetlinewidth{0.000000pt}%
\definecolor{currentstroke}{rgb}{0.000000,0.000000,0.000000}%
\pgfsetstrokecolor{currentstroke}%
\pgfsetdash{}{0pt}%
\pgfpathmoveto{\pgfqpoint{0.931202in}{1.237981in}}%
\pgfpathlineto{\pgfqpoint{0.946827in}{1.223510in}}%
\pgfpathlineto{\pgfqpoint{0.962468in}{1.209377in}}%
\pgfpathlineto{\pgfqpoint{0.970404in}{1.220124in}}%
\pgfpathlineto{\pgfqpoint{0.978320in}{1.231102in}}%
\pgfpathlineto{\pgfqpoint{0.962708in}{1.244902in}}%
\pgfpathlineto{\pgfqpoint{0.947112in}{1.259041in}}%
\pgfpathlineto{\pgfqpoint{0.939167in}{1.248397in}}%
\pgfpathlineto{\pgfqpoint{0.931202in}{1.237981in}}%
\pgfpathclose%
\pgfusepath{fill}%
\end{pgfscope}%
\begin{pgfscope}%
\pgfpathrectangle{\pgfqpoint{0.407592in}{0.222904in}}{\pgfqpoint{1.411163in}{1.411163in}}%
\pgfusepath{clip}%
\pgfsetbuttcap%
\pgfsetroundjoin%
\definecolor{currentfill}{rgb}{0.909804,0.698039,0.423529}%
\pgfsetfillcolor{currentfill}%
\pgfsetfillopacity{0.120000}%
\pgfsetlinewidth{0.000000pt}%
\definecolor{currentstroke}{rgb}{0.000000,0.000000,0.000000}%
\pgfsetstrokecolor{currentstroke}%
\pgfsetdash{}{0pt}%
\pgfpathmoveto{\pgfqpoint{0.852711in}{1.279900in}}%
\pgfpathlineto{\pgfqpoint{0.868316in}{1.264008in}}%
\pgfpathlineto{\pgfqpoint{0.883933in}{1.248360in}}%
\pgfpathlineto{\pgfqpoint{0.891974in}{1.257977in}}%
\pgfpathlineto{\pgfqpoint{0.899993in}{1.267837in}}%
\pgfpathlineto{\pgfqpoint{0.884409in}{1.283180in}}%
\pgfpathlineto{\pgfqpoint{0.868837in}{1.298773in}}%
\pgfpathlineto{\pgfqpoint{0.860786in}{1.289218in}}%
\pgfpathlineto{\pgfqpoint{0.852711in}{1.279900in}}%
\pgfpathclose%
\pgfusepath{fill}%
\end{pgfscope}%
\begin{pgfscope}%
\pgfpathrectangle{\pgfqpoint{0.407592in}{0.222904in}}{\pgfqpoint{1.411163in}{1.411163in}}%
\pgfusepath{clip}%
\pgfsetbuttcap%
\pgfsetroundjoin%
\definecolor{currentfill}{rgb}{0.909804,0.698039,0.423529}%
\pgfsetfillcolor{currentfill}%
\pgfsetfillopacity{0.120000}%
\pgfsetlinewidth{0.000000pt}%
\definecolor{currentstroke}{rgb}{0.000000,0.000000,0.000000}%
\pgfsetstrokecolor{currentstroke}%
\pgfsetdash{}{0pt}%
\pgfpathmoveto{\pgfqpoint{1.009594in}{1.204621in}}%
\pgfpathlineto{\pgfqpoint{1.025258in}{1.191991in}}%
\pgfpathlineto{\pgfqpoint{1.040942in}{1.179806in}}%
\pgfpathlineto{\pgfqpoint{1.048786in}{1.191679in}}%
\pgfpathlineto{\pgfqpoint{1.056612in}{1.203762in}}%
\pgfpathlineto{\pgfqpoint{1.040953in}{1.215615in}}%
\pgfpathlineto{\pgfqpoint{1.025314in}{1.227908in}}%
\pgfpathlineto{\pgfqpoint{1.017463in}{1.216159in}}%
\pgfpathlineto{\pgfqpoint{1.009594in}{1.204621in}}%
\pgfpathclose%
\pgfusepath{fill}%
\end{pgfscope}%
\begin{pgfscope}%
\pgfpathrectangle{\pgfqpoint{0.407592in}{0.222904in}}{\pgfqpoint{1.411163in}{1.411163in}}%
\pgfusepath{clip}%
\pgfsetbuttcap%
\pgfsetroundjoin%
\definecolor{currentfill}{rgb}{0.909804,0.698039,0.423529}%
\pgfsetfillcolor{currentfill}%
\pgfsetfillopacity{0.120000}%
\pgfsetlinewidth{0.000000pt}%
\definecolor{currentstroke}{rgb}{0.000000,0.000000,0.000000}%
\pgfsetstrokecolor{currentstroke}%
\pgfsetdash{}{0pt}%
\pgfpathmoveto{\pgfqpoint{0.774047in}{1.328725in}}%
\pgfpathlineto{\pgfqpoint{0.789641in}{1.311773in}}%
\pgfpathlineto{\pgfqpoint{0.805246in}{1.294992in}}%
\pgfpathlineto{\pgfqpoint{0.813405in}{1.303546in}}%
\pgfpathlineto{\pgfqpoint{0.821538in}{1.312343in}}%
\pgfpathlineto{\pgfqpoint{0.805968in}{1.328861in}}%
\pgfpathlineto{\pgfqpoint{0.790410in}{1.345557in}}%
\pgfpathlineto{\pgfqpoint{0.782241in}{1.337024in}}%
\pgfpathlineto{\pgfqpoint{0.774047in}{1.328725in}}%
\pgfpathclose%
\pgfusepath{fill}%
\end{pgfscope}%
\begin{pgfscope}%
\pgfpathrectangle{\pgfqpoint{0.407592in}{0.222904in}}{\pgfqpoint{1.411163in}{1.411163in}}%
\pgfusepath{clip}%
\pgfsetbuttcap%
\pgfsetroundjoin%
\definecolor{currentfill}{rgb}{0.909804,0.698039,0.423529}%
\pgfsetfillcolor{currentfill}%
\pgfsetfillopacity{0.120000}%
\pgfsetlinewidth{0.000000pt}%
\definecolor{currentstroke}{rgb}{0.000000,0.000000,0.000000}%
\pgfsetstrokecolor{currentstroke}%
\pgfsetdash{}{0pt}%
\pgfpathmoveto{\pgfqpoint{1.134991in}{1.187253in}}%
\pgfpathlineto{\pgfqpoint{1.150751in}{1.178272in}}%
\pgfpathlineto{\pgfqpoint{1.166540in}{1.169872in}}%
\pgfpathlineto{\pgfqpoint{1.174257in}{1.183362in}}%
\pgfpathlineto{\pgfqpoint{1.181959in}{1.197003in}}%
\pgfpathlineto{\pgfqpoint{1.166187in}{1.205163in}}%
\pgfpathlineto{\pgfqpoint{1.150444in}{1.213884in}}%
\pgfpathlineto{\pgfqpoint{1.142725in}{1.200491in}}%
\pgfpathlineto{\pgfqpoint{1.134991in}{1.187253in}}%
\pgfpathclose%
\pgfusepath{fill}%
\end{pgfscope}%
\begin{pgfscope}%
\pgfpathrectangle{\pgfqpoint{0.407592in}{0.222904in}}{\pgfqpoint{1.411163in}{1.411163in}}%
\pgfusepath{clip}%
\pgfsetbuttcap%
\pgfsetroundjoin%
\definecolor{currentfill}{rgb}{0.909804,0.698039,0.423529}%
\pgfsetfillcolor{currentfill}%
\pgfsetfillopacity{0.120000}%
\pgfsetlinewidth{0.000000pt}%
\definecolor{currentstroke}{rgb}{0.000000,0.000000,0.000000}%
\pgfsetstrokecolor{currentstroke}%
\pgfsetdash{}{0pt}%
\pgfpathmoveto{\pgfqpoint{0.883933in}{1.248360in}}%
\pgfpathlineto{\pgfqpoint{0.899562in}{1.232975in}}%
\pgfpathlineto{\pgfqpoint{0.915205in}{1.217875in}}%
\pgfpathlineto{\pgfqpoint{0.923214in}{1.227803in}}%
\pgfpathlineto{\pgfqpoint{0.931202in}{1.237981in}}%
\pgfpathlineto{\pgfqpoint{0.915590in}{1.252764in}}%
\pgfpathlineto{\pgfqpoint{0.899993in}{1.267837in}}%
\pgfpathlineto{\pgfqpoint{0.891974in}{1.257977in}}%
\pgfpathlineto{\pgfqpoint{0.883933in}{1.248360in}}%
\pgfpathclose%
\pgfusepath{fill}%
\end{pgfscope}%
\begin{pgfscope}%
\pgfpathrectangle{\pgfqpoint{0.407592in}{0.222904in}}{\pgfqpoint{1.411163in}{1.411163in}}%
\pgfusepath{clip}%
\pgfsetbuttcap%
\pgfsetroundjoin%
\definecolor{currentfill}{rgb}{0.909804,0.698039,0.423529}%
\pgfsetfillcolor{currentfill}%
\pgfsetfillopacity{0.120000}%
\pgfsetlinewidth{0.000000pt}%
\definecolor{currentstroke}{rgb}{0.000000,0.000000,0.000000}%
\pgfsetstrokecolor{currentstroke}%
\pgfsetdash{}{0pt}%
\pgfpathmoveto{\pgfqpoint{0.962468in}{1.209377in}}%
\pgfpathlineto{\pgfqpoint{0.978125in}{1.195605in}}%
\pgfpathlineto{\pgfqpoint{0.993798in}{1.182222in}}%
\pgfpathlineto{\pgfqpoint{1.001706in}{1.193304in}}%
\pgfpathlineto{\pgfqpoint{1.009594in}{1.204621in}}%
\pgfpathlineto{\pgfqpoint{0.993948in}{1.217667in}}%
\pgfpathlineto{\pgfqpoint{0.978320in}{1.231102in}}%
\pgfpathlineto{\pgfqpoint{0.970404in}{1.220124in}}%
\pgfpathlineto{\pgfqpoint{0.962468in}{1.209377in}}%
\pgfpathclose%
\pgfusepath{fill}%
\end{pgfscope}%
\begin{pgfscope}%
\pgfpathrectangle{\pgfqpoint{0.407592in}{0.222904in}}{\pgfqpoint{1.411163in}{1.411163in}}%
\pgfusepath{clip}%
\pgfsetbuttcap%
\pgfsetroundjoin%
\definecolor{currentfill}{rgb}{0.909804,0.698039,0.423529}%
\pgfsetfillcolor{currentfill}%
\pgfsetfillopacity{0.120000}%
\pgfsetlinewidth{0.000000pt}%
\definecolor{currentstroke}{rgb}{0.000000,0.000000,0.000000}%
\pgfsetstrokecolor{currentstroke}%
\pgfsetdash{}{0pt}%
\pgfpathmoveto{\pgfqpoint{1.087995in}{1.181483in}}%
\pgfpathlineto{\pgfqpoint{1.103722in}{1.171111in}}%
\pgfpathlineto{\pgfqpoint{1.119474in}{1.161288in}}%
\pgfpathlineto{\pgfqpoint{1.127240in}{1.174181in}}%
\pgfpathlineto{\pgfqpoint{1.134991in}{1.187253in}}%
\pgfpathlineto{\pgfqpoint{1.119258in}{1.196791in}}%
\pgfpathlineto{\pgfqpoint{1.103552in}{1.206863in}}%
\pgfpathlineto{\pgfqpoint{1.095782in}{1.194082in}}%
\pgfpathlineto{\pgfqpoint{1.087995in}{1.181483in}}%
\pgfpathclose%
\pgfusepath{fill}%
\end{pgfscope}%
\begin{pgfscope}%
\pgfpathrectangle{\pgfqpoint{0.407592in}{0.222904in}}{\pgfqpoint{1.411163in}{1.411163in}}%
\pgfusepath{clip}%
\pgfsetbuttcap%
\pgfsetroundjoin%
\definecolor{currentfill}{rgb}{0.909804,0.698039,0.423529}%
\pgfsetfillcolor{currentfill}%
\pgfsetfillopacity{0.120000}%
\pgfsetlinewidth{0.000000pt}%
\definecolor{currentstroke}{rgb}{0.000000,0.000000,0.000000}%
\pgfsetstrokecolor{currentstroke}%
\pgfsetdash{}{0pt}%
\pgfpathmoveto{\pgfqpoint{1.181959in}{1.197003in}}%
\pgfpathlineto{\pgfqpoint{1.197763in}{1.189423in}}%
\pgfpathlineto{\pgfqpoint{1.213601in}{1.182442in}}%
\pgfpathlineto{\pgfqpoint{1.221273in}{1.196422in}}%
\pgfpathlineto{\pgfqpoint{1.205443in}{1.203309in}}%
\pgfpathlineto{\pgfqpoint{1.189647in}{1.210783in}}%
\pgfusepath{fill}%
\end{pgfscope}%
\begin{pgfscope}%
\pgfpathrectangle{\pgfqpoint{0.407592in}{0.222904in}}{\pgfqpoint{1.411163in}{1.411163in}}%
\pgfusepath{clip}%
\pgfsetbuttcap%
\pgfsetroundjoin%
\definecolor{currentfill}{rgb}{0.909804,0.698039,0.423529}%
\pgfsetfillcolor{currentfill}%
\pgfsetfillopacity{0.120000}%
\pgfsetlinewidth{0.000000pt}%
\definecolor{currentstroke}{rgb}{0.000000,0.000000,0.000000}%
\pgfsetstrokecolor{currentstroke}%
\pgfsetdash{}{0pt}%
\pgfpathmoveto{\pgfqpoint{0.805246in}{1.294992in}}%
\pgfpathlineto{\pgfqpoint{0.820862in}{1.278397in}}%
\pgfpathlineto{\pgfqpoint{0.836490in}{1.262004in}}%
\pgfpathlineto{\pgfqpoint{0.844613in}{1.270826in}}%
\pgfpathlineto{\pgfqpoint{0.852711in}{1.279900in}}%
\pgfpathlineto{\pgfqpoint{0.837119in}{1.296017in}}%
\pgfpathlineto{\pgfqpoint{0.821538in}{1.312343in}}%
\pgfpathlineto{\pgfqpoint{0.813405in}{1.303546in}}%
\pgfpathlineto{\pgfqpoint{0.805246in}{1.294992in}}%
\pgfpathclose%
\pgfusepath{fill}%
\end{pgfscope}%
\begin{pgfscope}%
\pgfpathrectangle{\pgfqpoint{0.407592in}{0.222904in}}{\pgfqpoint{1.411163in}{1.411163in}}%
\pgfusepath{clip}%
\pgfsetbuttcap%
\pgfsetroundjoin%
\definecolor{currentfill}{rgb}{0.909804,0.698039,0.423529}%
\pgfsetfillcolor{currentfill}%
\pgfsetfillopacity{0.120000}%
\pgfsetlinewidth{0.000000pt}%
\definecolor{currentstroke}{rgb}{0.000000,0.000000,0.000000}%
\pgfsetstrokecolor{currentstroke}%
\pgfsetdash{}{0pt}%
\pgfpathmoveto{\pgfqpoint{1.040942in}{1.179806in}}%
\pgfpathlineto{\pgfqpoint{1.056645in}{1.168093in}}%
\pgfpathlineto{\pgfqpoint{1.072371in}{1.156883in}}%
\pgfpathlineto{\pgfqpoint{1.080192in}{1.169079in}}%
\pgfpathlineto{\pgfqpoint{1.087995in}{1.181483in}}%
\pgfpathlineto{\pgfqpoint{1.072293in}{1.192376in}}%
\pgfpathlineto{\pgfqpoint{1.056612in}{1.203762in}}%
\pgfpathlineto{\pgfqpoint{1.048786in}{1.191679in}}%
\pgfpathlineto{\pgfqpoint{1.040942in}{1.179806in}}%
\pgfpathclose%
\pgfusepath{fill}%
\end{pgfscope}%
\begin{pgfscope}%
\pgfpathrectangle{\pgfqpoint{0.407592in}{0.222904in}}{\pgfqpoint{1.411163in}{1.411163in}}%
\pgfusepath{clip}%
\pgfsetbuttcap%
\pgfsetroundjoin%
\definecolor{currentfill}{rgb}{0.909804,0.698039,0.423529}%
\pgfsetfillcolor{currentfill}%
\pgfsetfillopacity{0.120000}%
\pgfsetlinewidth{0.000000pt}%
\definecolor{currentstroke}{rgb}{0.000000,0.000000,0.000000}%
\pgfsetstrokecolor{currentstroke}%
\pgfsetdash{}{0pt}%
\pgfpathmoveto{\pgfqpoint{0.915205in}{1.217875in}}%
\pgfpathlineto{\pgfqpoint{0.930862in}{1.203083in}}%
\pgfpathlineto{\pgfqpoint{0.946532in}{1.188624in}}%
\pgfpathlineto{\pgfqpoint{0.954511in}{1.198873in}}%
\pgfpathlineto{\pgfqpoint{0.962468in}{1.209377in}}%
\pgfpathlineto{\pgfqpoint{0.946827in}{1.223510in}}%
\pgfpathlineto{\pgfqpoint{0.931202in}{1.237981in}}%
\pgfpathlineto{\pgfqpoint{0.923214in}{1.227803in}}%
\pgfpathlineto{\pgfqpoint{0.915205in}{1.217875in}}%
\pgfpathclose%
\pgfusepath{fill}%
\end{pgfscope}%
\begin{pgfscope}%
\pgfpathrectangle{\pgfqpoint{0.407592in}{0.222904in}}{\pgfqpoint{1.411163in}{1.411163in}}%
\pgfusepath{clip}%
\pgfsetbuttcap%
\pgfsetroundjoin%
\definecolor{currentfill}{rgb}{0.909804,0.698039,0.423529}%
\pgfsetfillcolor{currentfill}%
\pgfsetfillopacity{0.120000}%
\pgfsetlinewidth{0.000000pt}%
\definecolor{currentstroke}{rgb}{0.000000,0.000000,0.000000}%
\pgfsetstrokecolor{currentstroke}%
\pgfsetdash{}{0pt}%
\pgfpathmoveto{\pgfqpoint{0.836490in}{1.262004in}}%
\pgfpathlineto{\pgfqpoint{0.852128in}{1.245829in}}%
\pgfpathlineto{\pgfqpoint{0.867779in}{1.229891in}}%
\pgfpathlineto{\pgfqpoint{0.875868in}{1.238995in}}%
\pgfpathlineto{\pgfqpoint{0.883933in}{1.248360in}}%
\pgfpathlineto{\pgfqpoint{0.868316in}{1.264008in}}%
\pgfpathlineto{\pgfqpoint{0.852711in}{1.279900in}}%
\pgfpathlineto{\pgfqpoint{0.844613in}{1.270826in}}%
\pgfpathlineto{\pgfqpoint{0.836490in}{1.262004in}}%
\pgfpathclose%
\pgfusepath{fill}%
\end{pgfscope}%
\begin{pgfscope}%
\pgfpathrectangle{\pgfqpoint{0.407592in}{0.222904in}}{\pgfqpoint{1.411163in}{1.411163in}}%
\pgfusepath{clip}%
\pgfsetbuttcap%
\pgfsetroundjoin%
\definecolor{currentfill}{rgb}{0.909804,0.698039,0.423529}%
\pgfsetfillcolor{currentfill}%
\pgfsetfillopacity{0.120000}%
\pgfsetlinewidth{0.000000pt}%
\definecolor{currentstroke}{rgb}{0.000000,0.000000,0.000000}%
\pgfsetstrokecolor{currentstroke}%
\pgfsetdash{}{0pt}%
\pgfpathmoveto{\pgfqpoint{0.993798in}{1.182222in}}%
\pgfpathlineto{\pgfqpoint{1.009488in}{1.169257in}}%
\pgfpathlineto{\pgfqpoint{1.025197in}{1.156739in}}%
\pgfpathlineto{\pgfqpoint{1.033079in}{1.168155in}}%
\pgfpathlineto{\pgfqpoint{1.040942in}{1.179806in}}%
\pgfpathlineto{\pgfqpoint{1.025258in}{1.191991in}}%
\pgfpathlineto{\pgfqpoint{1.009594in}{1.204621in}}%
\pgfpathlineto{\pgfqpoint{1.001706in}{1.193304in}}%
\pgfpathlineto{\pgfqpoint{0.993798in}{1.182222in}}%
\pgfpathclose%
\pgfusepath{fill}%
\end{pgfscope}%
\begin{pgfscope}%
\pgfpathrectangle{\pgfqpoint{0.407592in}{0.222904in}}{\pgfqpoint{1.411163in}{1.411163in}}%
\pgfusepath{clip}%
\pgfsetbuttcap%
\pgfsetroundjoin%
\definecolor{currentfill}{rgb}{0.909804,0.698039,0.423529}%
\pgfsetfillcolor{currentfill}%
\pgfsetfillopacity{0.120000}%
\pgfsetlinewidth{0.000000pt}%
\definecolor{currentstroke}{rgb}{0.000000,0.000000,0.000000}%
\pgfsetstrokecolor{currentstroke}%
\pgfsetdash{}{0pt}%
\pgfpathmoveto{\pgfqpoint{1.166540in}{1.169872in}}%
\pgfpathlineto{\pgfqpoint{1.182360in}{1.162074in}}%
\pgfpathlineto{\pgfqpoint{1.198213in}{1.154898in}}%
\pgfpathlineto{\pgfqpoint{1.205914in}{1.168597in}}%
\pgfpathlineto{\pgfqpoint{1.213601in}{1.182442in}}%
\pgfpathlineto{\pgfqpoint{1.197763in}{1.189423in}}%
\pgfpathlineto{\pgfqpoint{1.181959in}{1.197003in}}%
\pgfpathlineto{\pgfqpoint{1.174257in}{1.183362in}}%
\pgfpathlineto{\pgfqpoint{1.166540in}{1.169872in}}%
\pgfpathclose%
\pgfusepath{fill}%
\end{pgfscope}%
\begin{pgfscope}%
\pgfpathrectangle{\pgfqpoint{0.407592in}{0.222904in}}{\pgfqpoint{1.411163in}{1.411163in}}%
\pgfusepath{clip}%
\pgfsetbuttcap%
\pgfsetroundjoin%
\definecolor{currentfill}{rgb}{0.909804,0.698039,0.423529}%
\pgfsetfillcolor{currentfill}%
\pgfsetfillopacity{0.120000}%
\pgfsetlinewidth{0.000000pt}%
\definecolor{currentstroke}{rgb}{0.000000,0.000000,0.000000}%
\pgfsetstrokecolor{currentstroke}%
\pgfsetdash{}{0pt}%
\pgfpathmoveto{\pgfqpoint{1.119474in}{1.161288in}}%
\pgfpathlineto{\pgfqpoint{1.135253in}{1.152039in}}%
\pgfpathlineto{\pgfqpoint{1.151060in}{1.143390in}}%
\pgfpathlineto{\pgfqpoint{1.158808in}{1.156544in}}%
\pgfpathlineto{\pgfqpoint{1.166540in}{1.169872in}}%
\pgfpathlineto{\pgfqpoint{1.150751in}{1.178272in}}%
\pgfpathlineto{\pgfqpoint{1.134991in}{1.187253in}}%
\pgfpathlineto{\pgfqpoint{1.127240in}{1.174181in}}%
\pgfpathlineto{\pgfqpoint{1.119474in}{1.161288in}}%
\pgfpathclose%
\pgfusepath{fill}%
\end{pgfscope}%
\begin{pgfscope}%
\pgfpathrectangle{\pgfqpoint{0.407592in}{0.222904in}}{\pgfqpoint{1.411163in}{1.411163in}}%
\pgfusepath{clip}%
\pgfsetbuttcap%
\pgfsetroundjoin%
\definecolor{currentfill}{rgb}{0.909804,0.698039,0.423529}%
\pgfsetfillcolor{currentfill}%
\pgfsetfillopacity{0.120000}%
\pgfsetlinewidth{0.000000pt}%
\definecolor{currentstroke}{rgb}{0.000000,0.000000,0.000000}%
\pgfsetstrokecolor{currentstroke}%
\pgfsetdash{}{0pt}%
\pgfpathmoveto{\pgfqpoint{0.757579in}{1.312855in}}%
\pgfpathlineto{\pgfqpoint{0.773210in}{1.295666in}}%
\pgfpathlineto{\pgfqpoint{0.788852in}{1.278642in}}%
\pgfpathlineto{\pgfqpoint{0.797062in}{1.286689in}}%
\pgfpathlineto{\pgfqpoint{0.805246in}{1.294992in}}%
\pgfpathlineto{\pgfqpoint{0.789641in}{1.311773in}}%
\pgfpathlineto{\pgfqpoint{0.774047in}{1.328725in}}%
\pgfpathlineto{\pgfqpoint{0.765826in}{1.320667in}}%
\pgfpathlineto{\pgfqpoint{0.757579in}{1.312855in}}%
\pgfpathclose%
\pgfusepath{fill}%
\end{pgfscope}%
\begin{pgfscope}%
\pgfpathrectangle{\pgfqpoint{0.407592in}{0.222904in}}{\pgfqpoint{1.411163in}{1.411163in}}%
\pgfusepath{clip}%
\pgfsetbuttcap%
\pgfsetroundjoin%
\definecolor{currentfill}{rgb}{0.909804,0.698039,0.423529}%
\pgfsetfillcolor{currentfill}%
\pgfsetfillopacity{0.120000}%
\pgfsetlinewidth{0.000000pt}%
\definecolor{currentstroke}{rgb}{0.000000,0.000000,0.000000}%
\pgfsetstrokecolor{currentstroke}%
\pgfsetdash{}{0pt}%
\pgfpathmoveto{\pgfqpoint{1.072371in}{1.156883in}}%
\pgfpathlineto{\pgfqpoint{1.088119in}{1.146205in}}%
\pgfpathlineto{\pgfqpoint{1.103892in}{1.136089in}}%
\pgfpathlineto{\pgfqpoint{1.111691in}{1.148586in}}%
\pgfpathlineto{\pgfqpoint{1.119474in}{1.161288in}}%
\pgfpathlineto{\pgfqpoint{1.103722in}{1.171111in}}%
\pgfpathlineto{\pgfqpoint{1.087995in}{1.181483in}}%
\pgfpathlineto{\pgfqpoint{1.080192in}{1.169079in}}%
\pgfpathlineto{\pgfqpoint{1.072371in}{1.156883in}}%
\pgfpathclose%
\pgfusepath{fill}%
\end{pgfscope}%
\begin{pgfscope}%
\pgfpathrectangle{\pgfqpoint{0.407592in}{0.222904in}}{\pgfqpoint{1.411163in}{1.411163in}}%
\pgfusepath{clip}%
\pgfsetbuttcap%
\pgfsetroundjoin%
\definecolor{currentfill}{rgb}{0.909804,0.698039,0.423529}%
\pgfsetfillcolor{currentfill}%
\pgfsetfillopacity{0.120000}%
\pgfsetlinewidth{0.000000pt}%
\definecolor{currentstroke}{rgb}{0.000000,0.000000,0.000000}%
\pgfsetstrokecolor{currentstroke}%
\pgfsetdash{}{0pt}%
\pgfpathmoveto{\pgfqpoint{1.213601in}{1.182442in}}%
\pgfpathlineto{\pgfqpoint{1.229473in}{1.176074in}}%
\pgfpathlineto{\pgfqpoint{1.245382in}{1.170334in}}%
\pgfpathlineto{\pgfqpoint{1.253040in}{1.184463in}}%
\pgfpathlineto{\pgfqpoint{1.237138in}{1.190135in}}%
\pgfpathlineto{\pgfqpoint{1.221273in}{1.196422in}}%
\pgfusepath{fill}%
\end{pgfscope}%
\begin{pgfscope}%
\pgfpathrectangle{\pgfqpoint{0.407592in}{0.222904in}}{\pgfqpoint{1.411163in}{1.411163in}}%
\pgfusepath{clip}%
\pgfsetbuttcap%
\pgfsetroundjoin%
\definecolor{currentfill}{rgb}{0.909804,0.698039,0.423529}%
\pgfsetfillcolor{currentfill}%
\pgfsetfillopacity{0.120000}%
\pgfsetlinewidth{0.000000pt}%
\definecolor{currentstroke}{rgb}{0.000000,0.000000,0.000000}%
\pgfsetstrokecolor{currentstroke}%
\pgfsetdash{}{0pt}%
\pgfpathmoveto{\pgfqpoint{0.946532in}{1.188624in}}%
\pgfpathlineto{\pgfqpoint{0.962218in}{1.174523in}}%
\pgfpathlineto{\pgfqpoint{0.977920in}{1.160811in}}%
\pgfpathlineto{\pgfqpoint{0.985869in}{1.171387in}}%
\pgfpathlineto{\pgfqpoint{0.993798in}{1.182222in}}%
\pgfpathlineto{\pgfqpoint{0.978125in}{1.195605in}}%
\pgfpathlineto{\pgfqpoint{0.962468in}{1.209377in}}%
\pgfpathlineto{\pgfqpoint{0.954511in}{1.198873in}}%
\pgfpathlineto{\pgfqpoint{0.946532in}{1.188624in}}%
\pgfpathclose%
\pgfusepath{fill}%
\end{pgfscope}%
\begin{pgfscope}%
\pgfpathrectangle{\pgfqpoint{0.407592in}{0.222904in}}{\pgfqpoint{1.411163in}{1.411163in}}%
\pgfusepath{clip}%
\pgfsetbuttcap%
\pgfsetroundjoin%
\definecolor{currentfill}{rgb}{0.909804,0.698039,0.423529}%
\pgfsetfillcolor{currentfill}%
\pgfsetfillopacity{0.120000}%
\pgfsetlinewidth{0.000000pt}%
\definecolor{currentstroke}{rgb}{0.000000,0.000000,0.000000}%
\pgfsetstrokecolor{currentstroke}%
\pgfsetdash{}{0pt}%
\pgfpathmoveto{\pgfqpoint{0.867779in}{1.229891in}}%
\pgfpathlineto{\pgfqpoint{0.883442in}{1.214210in}}%
\pgfpathlineto{\pgfqpoint{0.899118in}{1.198808in}}%
\pgfpathlineto{\pgfqpoint{0.907173in}{1.208207in}}%
\pgfpathlineto{\pgfqpoint{0.915205in}{1.217875in}}%
\pgfpathlineto{\pgfqpoint{0.899562in}{1.232975in}}%
\pgfpathlineto{\pgfqpoint{0.883933in}{1.248360in}}%
\pgfpathlineto{\pgfqpoint{0.875868in}{1.238995in}}%
\pgfpathlineto{\pgfqpoint{0.867779in}{1.229891in}}%
\pgfpathclose%
\pgfusepath{fill}%
\end{pgfscope}%
\begin{pgfscope}%
\pgfpathrectangle{\pgfqpoint{0.407592in}{0.222904in}}{\pgfqpoint{1.411163in}{1.411163in}}%
\pgfusepath{clip}%
\pgfsetbuttcap%
\pgfsetroundjoin%
\definecolor{currentfill}{rgb}{0.909804,0.698039,0.423529}%
\pgfsetfillcolor{currentfill}%
\pgfsetfillopacity{0.120000}%
\pgfsetlinewidth{0.000000pt}%
\definecolor{currentstroke}{rgb}{0.000000,0.000000,0.000000}%
\pgfsetstrokecolor{currentstroke}%
\pgfsetdash{}{0pt}%
\pgfpathmoveto{\pgfqpoint{1.025197in}{1.156739in}}%
\pgfpathlineto{\pgfqpoint{1.040926in}{1.144699in}}%
\pgfpathlineto{\pgfqpoint{1.056675in}{1.133169in}}%
\pgfpathlineto{\pgfqpoint{1.064532in}{1.144908in}}%
\pgfpathlineto{\pgfqpoint{1.072371in}{1.156883in}}%
\pgfpathlineto{\pgfqpoint{1.056645in}{1.168093in}}%
\pgfpathlineto{\pgfqpoint{1.040942in}{1.179806in}}%
\pgfpathlineto{\pgfqpoint{1.033079in}{1.168155in}}%
\pgfpathlineto{\pgfqpoint{1.025197in}{1.156739in}}%
\pgfpathclose%
\pgfusepath{fill}%
\end{pgfscope}%
\begin{pgfscope}%
\pgfpathrectangle{\pgfqpoint{0.407592in}{0.222904in}}{\pgfqpoint{1.411163in}{1.411163in}}%
\pgfusepath{clip}%
\pgfsetbuttcap%
\pgfsetroundjoin%
\definecolor{currentfill}{rgb}{0.909804,0.698039,0.423529}%
\pgfsetfillcolor{currentfill}%
\pgfsetfillopacity{0.120000}%
\pgfsetlinewidth{0.000000pt}%
\definecolor{currentstroke}{rgb}{0.000000,0.000000,0.000000}%
\pgfsetstrokecolor{currentstroke}%
\pgfsetdash{}{0pt}%
\pgfpathmoveto{\pgfqpoint{0.788852in}{1.278642in}}%
\pgfpathlineto{\pgfqpoint{0.804504in}{1.261798in}}%
\pgfpathlineto{\pgfqpoint{0.820167in}{1.245148in}}%
\pgfpathlineto{\pgfqpoint{0.828341in}{1.253442in}}%
\pgfpathlineto{\pgfqpoint{0.836490in}{1.262004in}}%
\pgfpathlineto{\pgfqpoint{0.820862in}{1.278397in}}%
\pgfpathlineto{\pgfqpoint{0.805246in}{1.294992in}}%
\pgfpathlineto{\pgfqpoint{0.797062in}{1.286689in}}%
\pgfpathlineto{\pgfqpoint{0.788852in}{1.278642in}}%
\pgfpathclose%
\pgfusepath{fill}%
\end{pgfscope}%
\begin{pgfscope}%
\pgfpathrectangle{\pgfqpoint{0.407592in}{0.222904in}}{\pgfqpoint{1.411163in}{1.411163in}}%
\pgfusepath{clip}%
\pgfsetbuttcap%
\pgfsetroundjoin%
\definecolor{currentfill}{rgb}{0.909804,0.698039,0.423529}%
\pgfsetfillcolor{currentfill}%
\pgfsetfillopacity{0.120000}%
\pgfsetlinewidth{0.000000pt}%
\definecolor{currentstroke}{rgb}{0.000000,0.000000,0.000000}%
\pgfsetstrokecolor{currentstroke}%
\pgfsetdash{}{0pt}%
\pgfpathmoveto{\pgfqpoint{0.899118in}{1.198808in}}%
\pgfpathlineto{\pgfqpoint{0.914806in}{1.183708in}}%
\pgfpathlineto{\pgfqpoint{0.930509in}{1.168936in}}%
\pgfpathlineto{\pgfqpoint{0.938532in}{1.178641in}}%
\pgfpathlineto{\pgfqpoint{0.946532in}{1.188624in}}%
\pgfpathlineto{\pgfqpoint{0.930862in}{1.203083in}}%
\pgfpathlineto{\pgfqpoint{0.915205in}{1.217875in}}%
\pgfpathlineto{\pgfqpoint{0.907173in}{1.208207in}}%
\pgfpathlineto{\pgfqpoint{0.899118in}{1.198808in}}%
\pgfpathclose%
\pgfusepath{fill}%
\end{pgfscope}%
\begin{pgfscope}%
\pgfpathrectangle{\pgfqpoint{0.407592in}{0.222904in}}{\pgfqpoint{1.411163in}{1.411163in}}%
\pgfusepath{clip}%
\pgfsetbuttcap%
\pgfsetroundjoin%
\definecolor{currentfill}{rgb}{0.909804,0.698039,0.423529}%
\pgfsetfillcolor{currentfill}%
\pgfsetfillopacity{0.120000}%
\pgfsetlinewidth{0.000000pt}%
\definecolor{currentstroke}{rgb}{0.000000,0.000000,0.000000}%
\pgfsetstrokecolor{currentstroke}%
\pgfsetdash{}{0pt}%
\pgfpathmoveto{\pgfqpoint{0.977920in}{1.160811in}}%
\pgfpathlineto{\pgfqpoint{0.993639in}{1.147515in}}%
\pgfpathlineto{\pgfqpoint{1.009375in}{1.134669in}}%
\pgfpathlineto{\pgfqpoint{1.017296in}{1.145572in}}%
\pgfpathlineto{\pgfqpoint{1.025197in}{1.156739in}}%
\pgfpathlineto{\pgfqpoint{1.009488in}{1.169257in}}%
\pgfpathlineto{\pgfqpoint{0.993798in}{1.182222in}}%
\pgfpathlineto{\pgfqpoint{0.985869in}{1.171387in}}%
\pgfpathlineto{\pgfqpoint{0.977920in}{1.160811in}}%
\pgfpathclose%
\pgfusepath{fill}%
\end{pgfscope}%
\begin{pgfscope}%
\pgfpathrectangle{\pgfqpoint{0.407592in}{0.222904in}}{\pgfqpoint{1.411163in}{1.411163in}}%
\pgfusepath{clip}%
\pgfsetbuttcap%
\pgfsetroundjoin%
\definecolor{currentfill}{rgb}{0.909804,0.698039,0.423529}%
\pgfsetfillcolor{currentfill}%
\pgfsetfillopacity{0.120000}%
\pgfsetlinewidth{0.000000pt}%
\definecolor{currentstroke}{rgb}{0.000000,0.000000,0.000000}%
\pgfsetstrokecolor{currentstroke}%
\pgfsetdash{}{0pt}%
\pgfpathmoveto{\pgfqpoint{1.151060in}{1.143390in}}%
\pgfpathlineto{\pgfqpoint{1.166897in}{1.135365in}}%
\pgfpathlineto{\pgfqpoint{1.182767in}{1.127988in}}%
\pgfpathlineto{\pgfqpoint{1.190497in}{1.141358in}}%
\pgfpathlineto{\pgfqpoint{1.198213in}{1.154898in}}%
\pgfpathlineto{\pgfqpoint{1.182360in}{1.162074in}}%
\pgfpathlineto{\pgfqpoint{1.166540in}{1.169872in}}%
\pgfpathlineto{\pgfqpoint{1.158808in}{1.156544in}}%
\pgfpathlineto{\pgfqpoint{1.151060in}{1.143390in}}%
\pgfpathclose%
\pgfusepath{fill}%
\end{pgfscope}%
\begin{pgfscope}%
\pgfpathrectangle{\pgfqpoint{0.407592in}{0.222904in}}{\pgfqpoint{1.411163in}{1.411163in}}%
\pgfusepath{clip}%
\pgfsetbuttcap%
\pgfsetroundjoin%
\definecolor{currentfill}{rgb}{0.909804,0.698039,0.423529}%
\pgfsetfillcolor{currentfill}%
\pgfsetfillopacity{0.120000}%
\pgfsetlinewidth{0.000000pt}%
\definecolor{currentstroke}{rgb}{0.000000,0.000000,0.000000}%
\pgfsetstrokecolor{currentstroke}%
\pgfsetdash{}{0pt}%
\pgfpathmoveto{\pgfqpoint{1.198213in}{1.154898in}}%
\pgfpathlineto{\pgfqpoint{1.214101in}{1.148362in}}%
\pgfpathlineto{\pgfqpoint{1.230024in}{1.142482in}}%
\pgfpathlineto{\pgfqpoint{1.237711in}{1.156336in}}%
\pgfpathlineto{\pgfqpoint{1.245382in}{1.170334in}}%
\pgfpathlineto{\pgfqpoint{1.229473in}{1.176074in}}%
\pgfpathlineto{\pgfqpoint{1.213601in}{1.182442in}}%
\pgfpathlineto{\pgfqpoint{1.205914in}{1.168597in}}%
\pgfpathlineto{\pgfqpoint{1.198213in}{1.154898in}}%
\pgfpathclose%
\pgfusepath{fill}%
\end{pgfscope}%
\begin{pgfscope}%
\pgfpathrectangle{\pgfqpoint{0.407592in}{0.222904in}}{\pgfqpoint{1.411163in}{1.411163in}}%
\pgfusepath{clip}%
\pgfsetbuttcap%
\pgfsetroundjoin%
\definecolor{currentfill}{rgb}{0.909804,0.698039,0.423529}%
\pgfsetfillcolor{currentfill}%
\pgfsetfillopacity{0.120000}%
\pgfsetlinewidth{0.000000pt}%
\definecolor{currentstroke}{rgb}{0.000000,0.000000,0.000000}%
\pgfsetstrokecolor{currentstroke}%
\pgfsetdash{}{0pt}%
\pgfpathmoveto{\pgfqpoint{0.820167in}{1.245148in}}%
\pgfpathlineto{\pgfqpoint{0.835842in}{1.228710in}}%
\pgfpathlineto{\pgfqpoint{0.851527in}{1.212502in}}%
\pgfpathlineto{\pgfqpoint{0.859666in}{1.221057in}}%
\pgfpathlineto{\pgfqpoint{0.867779in}{1.229891in}}%
\pgfpathlineto{\pgfqpoint{0.852128in}{1.245829in}}%
\pgfpathlineto{\pgfqpoint{0.836490in}{1.262004in}}%
\pgfpathlineto{\pgfqpoint{0.828341in}{1.253442in}}%
\pgfpathlineto{\pgfqpoint{0.820167in}{1.245148in}}%
\pgfpathclose%
\pgfusepath{fill}%
\end{pgfscope}%
\begin{pgfscope}%
\pgfpathrectangle{\pgfqpoint{0.407592in}{0.222904in}}{\pgfqpoint{1.411163in}{1.411163in}}%
\pgfusepath{clip}%
\pgfsetbuttcap%
\pgfsetroundjoin%
\definecolor{currentfill}{rgb}{0.909804,0.698039,0.423529}%
\pgfsetfillcolor{currentfill}%
\pgfsetfillopacity{0.120000}%
\pgfsetlinewidth{0.000000pt}%
\definecolor{currentstroke}{rgb}{0.000000,0.000000,0.000000}%
\pgfsetstrokecolor{currentstroke}%
\pgfsetdash{}{0pt}%
\pgfpathmoveto{\pgfqpoint{1.103892in}{1.136089in}}%
\pgfpathlineto{\pgfqpoint{1.119690in}{1.126565in}}%
\pgfpathlineto{\pgfqpoint{1.135516in}{1.117660in}}%
\pgfpathlineto{\pgfqpoint{1.143296in}{1.130424in}}%
\pgfpathlineto{\pgfqpoint{1.151060in}{1.143390in}}%
\pgfpathlineto{\pgfqpoint{1.135253in}{1.152039in}}%
\pgfpathlineto{\pgfqpoint{1.119474in}{1.161288in}}%
\pgfpathlineto{\pgfqpoint{1.111691in}{1.148586in}}%
\pgfpathlineto{\pgfqpoint{1.103892in}{1.136089in}}%
\pgfpathclose%
\pgfusepath{fill}%
\end{pgfscope}%
\begin{pgfscope}%
\pgfpathrectangle{\pgfqpoint{0.407592in}{0.222904in}}{\pgfqpoint{1.411163in}{1.411163in}}%
\pgfusepath{clip}%
\pgfsetbuttcap%
\pgfsetroundjoin%
\definecolor{currentfill}{rgb}{0.909804,0.698039,0.423529}%
\pgfsetfillcolor{currentfill}%
\pgfsetfillopacity{0.120000}%
\pgfsetlinewidth{0.000000pt}%
\definecolor{currentstroke}{rgb}{0.000000,0.000000,0.000000}%
\pgfsetstrokecolor{currentstroke}%
\pgfsetdash{}{0pt}%
\pgfpathmoveto{\pgfqpoint{1.056675in}{1.133169in}}%
\pgfpathlineto{\pgfqpoint{1.072446in}{1.122182in}}%
\pgfpathlineto{\pgfqpoint{1.088241in}{1.111769in}}%
\pgfpathlineto{\pgfqpoint{1.096075in}{1.123812in}}%
\pgfpathlineto{\pgfqpoint{1.103892in}{1.136089in}}%
\pgfpathlineto{\pgfqpoint{1.088119in}{1.146205in}}%
\pgfpathlineto{\pgfqpoint{1.072371in}{1.156883in}}%
\pgfpathlineto{\pgfqpoint{1.064532in}{1.144908in}}%
\pgfpathlineto{\pgfqpoint{1.056675in}{1.133169in}}%
\pgfpathclose%
\pgfusepath{fill}%
\end{pgfscope}%
\begin{pgfscope}%
\pgfpathrectangle{\pgfqpoint{0.407592in}{0.222904in}}{\pgfqpoint{1.411163in}{1.411163in}}%
\pgfusepath{clip}%
\pgfsetbuttcap%
\pgfsetroundjoin%
\definecolor{currentfill}{rgb}{0.909804,0.698039,0.423529}%
\pgfsetfillcolor{currentfill}%
\pgfsetfillopacity{0.120000}%
\pgfsetlinewidth{0.000000pt}%
\definecolor{currentstroke}{rgb}{0.000000,0.000000,0.000000}%
\pgfsetstrokecolor{currentstroke}%
\pgfsetdash{}{0pt}%
\pgfpathmoveto{\pgfqpoint{1.245382in}{1.170334in}}%
\pgfpathlineto{\pgfqpoint{1.261331in}{1.165232in}}%
\pgfpathlineto{\pgfqpoint{1.277319in}{1.160778in}}%
\pgfpathlineto{\pgfqpoint{1.284964in}{1.175001in}}%
\pgfpathlineto{\pgfqpoint{1.268982in}{1.179416in}}%
\pgfpathlineto{\pgfqpoint{1.253040in}{1.184463in}}%
\pgfusepath{fill}%
\end{pgfscope}%
\begin{pgfscope}%
\pgfpathrectangle{\pgfqpoint{0.407592in}{0.222904in}}{\pgfqpoint{1.411163in}{1.411163in}}%
\pgfusepath{clip}%
\pgfsetbuttcap%
\pgfsetroundjoin%
\definecolor{currentfill}{rgb}{0.909804,0.698039,0.423529}%
\pgfsetfillcolor{currentfill}%
\pgfsetfillopacity{0.120000}%
\pgfsetlinewidth{0.000000pt}%
\definecolor{currentstroke}{rgb}{0.000000,0.000000,0.000000}%
\pgfsetstrokecolor{currentstroke}%
\pgfsetdash{}{0pt}%
\pgfpathmoveto{\pgfqpoint{0.741002in}{1.297996in}}%
\pgfpathlineto{\pgfqpoint{0.756671in}{1.280592in}}%
\pgfpathlineto{\pgfqpoint{0.772350in}{1.263348in}}%
\pgfpathlineto{\pgfqpoint{0.780615in}{1.270860in}}%
\pgfpathlineto{\pgfqpoint{0.788852in}{1.278642in}}%
\pgfpathlineto{\pgfqpoint{0.773210in}{1.295666in}}%
\pgfpathlineto{\pgfqpoint{0.757579in}{1.312855in}}%
\pgfpathlineto{\pgfqpoint{0.749305in}{1.305297in}}%
\pgfpathlineto{\pgfqpoint{0.741002in}{1.297996in}}%
\pgfpathclose%
\pgfusepath{fill}%
\end{pgfscope}%
\begin{pgfscope}%
\pgfpathrectangle{\pgfqpoint{0.407592in}{0.222904in}}{\pgfqpoint{1.411163in}{1.411163in}}%
\pgfusepath{clip}%
\pgfsetbuttcap%
\pgfsetroundjoin%
\definecolor{currentfill}{rgb}{0.909804,0.698039,0.423529}%
\pgfsetfillcolor{currentfill}%
\pgfsetfillopacity{0.120000}%
\pgfsetlinewidth{0.000000pt}%
\definecolor{currentstroke}{rgb}{0.000000,0.000000,0.000000}%
\pgfsetstrokecolor{currentstroke}%
\pgfsetdash{}{0pt}%
\pgfpathmoveto{\pgfqpoint{0.930509in}{1.168936in}}%
\pgfpathlineto{\pgfqpoint{0.946226in}{1.154520in}}%
\pgfpathlineto{\pgfqpoint{0.961958in}{1.140489in}}%
\pgfpathlineto{\pgfqpoint{0.969950in}{1.150507in}}%
\pgfpathlineto{\pgfqpoint{0.977920in}{1.160811in}}%
\pgfpathlineto{\pgfqpoint{0.962218in}{1.174523in}}%
\pgfpathlineto{\pgfqpoint{0.946532in}{1.188624in}}%
\pgfpathlineto{\pgfqpoint{0.938532in}{1.178641in}}%
\pgfpathlineto{\pgfqpoint{0.930509in}{1.168936in}}%
\pgfpathclose%
\pgfusepath{fill}%
\end{pgfscope}%
\begin{pgfscope}%
\pgfpathrectangle{\pgfqpoint{0.407592in}{0.222904in}}{\pgfqpoint{1.411163in}{1.411163in}}%
\pgfusepath{clip}%
\pgfsetbuttcap%
\pgfsetroundjoin%
\definecolor{currentfill}{rgb}{0.909804,0.698039,0.423529}%
\pgfsetfillcolor{currentfill}%
\pgfsetfillopacity{0.120000}%
\pgfsetlinewidth{0.000000pt}%
\definecolor{currentstroke}{rgb}{0.000000,0.000000,0.000000}%
\pgfsetstrokecolor{currentstroke}%
\pgfsetdash{}{0pt}%
\pgfpathmoveto{\pgfqpoint{0.851527in}{1.212502in}}%
\pgfpathlineto{\pgfqpoint{0.867225in}{1.196545in}}%
\pgfpathlineto{\pgfqpoint{0.882935in}{1.180860in}}%
\pgfpathlineto{\pgfqpoint{0.891038in}{1.189689in}}%
\pgfpathlineto{\pgfqpoint{0.899118in}{1.198808in}}%
\pgfpathlineto{\pgfqpoint{0.883442in}{1.214210in}}%
\pgfpathlineto{\pgfqpoint{0.867779in}{1.229891in}}%
\pgfpathlineto{\pgfqpoint{0.859666in}{1.221057in}}%
\pgfpathlineto{\pgfqpoint{0.851527in}{1.212502in}}%
\pgfpathclose%
\pgfusepath{fill}%
\end{pgfscope}%
\begin{pgfscope}%
\pgfpathrectangle{\pgfqpoint{0.407592in}{0.222904in}}{\pgfqpoint{1.411163in}{1.411163in}}%
\pgfusepath{clip}%
\pgfsetbuttcap%
\pgfsetroundjoin%
\definecolor{currentfill}{rgb}{0.909804,0.698039,0.423529}%
\pgfsetfillcolor{currentfill}%
\pgfsetfillopacity{0.120000}%
\pgfsetlinewidth{0.000000pt}%
\definecolor{currentstroke}{rgb}{0.000000,0.000000,0.000000}%
\pgfsetstrokecolor{currentstroke}%
\pgfsetdash{}{0pt}%
\pgfpathmoveto{\pgfqpoint{1.009375in}{1.134669in}}%
\pgfpathlineto{\pgfqpoint{1.025129in}{1.122305in}}%
\pgfpathlineto{\pgfqpoint{1.040904in}{1.110457in}}%
\pgfpathlineto{\pgfqpoint{1.048799in}{1.121680in}}%
\pgfpathlineto{\pgfqpoint{1.056675in}{1.133169in}}%
\pgfpathlineto{\pgfqpoint{1.040926in}{1.144699in}}%
\pgfpathlineto{\pgfqpoint{1.025197in}{1.156739in}}%
\pgfpathlineto{\pgfqpoint{1.017296in}{1.145572in}}%
\pgfpathlineto{\pgfqpoint{1.009375in}{1.134669in}}%
\pgfpathclose%
\pgfusepath{fill}%
\end{pgfscope}%
\begin{pgfscope}%
\pgfpathrectangle{\pgfqpoint{0.407592in}{0.222904in}}{\pgfqpoint{1.411163in}{1.411163in}}%
\pgfusepath{clip}%
\pgfsetbuttcap%
\pgfsetroundjoin%
\definecolor{currentfill}{rgb}{0.909804,0.698039,0.423529}%
\pgfsetfillcolor{currentfill}%
\pgfsetfillopacity{0.120000}%
\pgfsetlinewidth{0.000000pt}%
\definecolor{currentstroke}{rgb}{0.000000,0.000000,0.000000}%
\pgfsetstrokecolor{currentstroke}%
\pgfsetdash{}{0pt}%
\pgfpathmoveto{\pgfqpoint{0.772350in}{1.263348in}}%
\pgfpathlineto{\pgfqpoint{0.788040in}{1.246277in}}%
\pgfpathlineto{\pgfqpoint{0.803740in}{1.229394in}}%
\pgfpathlineto{\pgfqpoint{0.811967in}{1.237130in}}%
\pgfpathlineto{\pgfqpoint{0.820167in}{1.245148in}}%
\pgfpathlineto{\pgfqpoint{0.804504in}{1.261798in}}%
\pgfpathlineto{\pgfqpoint{0.788852in}{1.278642in}}%
\pgfpathlineto{\pgfqpoint{0.780615in}{1.270860in}}%
\pgfpathlineto{\pgfqpoint{0.772350in}{1.263348in}}%
\pgfpathclose%
\pgfusepath{fill}%
\end{pgfscope}%
\begin{pgfscope}%
\pgfpathrectangle{\pgfqpoint{0.407592in}{0.222904in}}{\pgfqpoint{1.411163in}{1.411163in}}%
\pgfusepath{clip}%
\pgfsetbuttcap%
\pgfsetroundjoin%
\definecolor{currentfill}{rgb}{0.909804,0.698039,0.423529}%
\pgfsetfillcolor{currentfill}%
\pgfsetfillopacity{0.120000}%
\pgfsetlinewidth{0.000000pt}%
\definecolor{currentstroke}{rgb}{0.000000,0.000000,0.000000}%
\pgfsetstrokecolor{currentstroke}%
\pgfsetdash{}{0pt}%
\pgfpathmoveto{\pgfqpoint{1.182767in}{1.127988in}}%
\pgfpathlineto{\pgfqpoint{1.198670in}{1.121278in}}%
\pgfpathlineto{\pgfqpoint{1.214609in}{1.115252in}}%
\pgfpathlineto{\pgfqpoint{1.222324in}{1.128783in}}%
\pgfpathlineto{\pgfqpoint{1.230024in}{1.142482in}}%
\pgfpathlineto{\pgfqpoint{1.214101in}{1.148362in}}%
\pgfpathlineto{\pgfqpoint{1.198213in}{1.154898in}}%
\pgfpathlineto{\pgfqpoint{1.190497in}{1.141358in}}%
\pgfpathlineto{\pgfqpoint{1.182767in}{1.127988in}}%
\pgfpathclose%
\pgfusepath{fill}%
\end{pgfscope}%
\begin{pgfscope}%
\pgfpathrectangle{\pgfqpoint{0.407592in}{0.222904in}}{\pgfqpoint{1.411163in}{1.411163in}}%
\pgfusepath{clip}%
\pgfsetbuttcap%
\pgfsetroundjoin%
\definecolor{currentfill}{rgb}{0.909804,0.698039,0.423529}%
\pgfsetfillcolor{currentfill}%
\pgfsetfillopacity{0.120000}%
\pgfsetlinewidth{0.000000pt}%
\definecolor{currentstroke}{rgb}{0.000000,0.000000,0.000000}%
\pgfsetstrokecolor{currentstroke}%
\pgfsetdash{}{0pt}%
\pgfpathmoveto{\pgfqpoint{1.135516in}{1.117660in}}%
\pgfpathlineto{\pgfqpoint{1.151372in}{1.109402in}}%
\pgfpathlineto{\pgfqpoint{1.167259in}{1.101817in}}%
\pgfpathlineto{\pgfqpoint{1.175021in}{1.114803in}}%
\pgfpathlineto{\pgfqpoint{1.182767in}{1.127988in}}%
\pgfpathlineto{\pgfqpoint{1.166897in}{1.135365in}}%
\pgfpathlineto{\pgfqpoint{1.151060in}{1.143390in}}%
\pgfpathlineto{\pgfqpoint{1.143296in}{1.130424in}}%
\pgfpathlineto{\pgfqpoint{1.135516in}{1.117660in}}%
\pgfpathclose%
\pgfusepath{fill}%
\end{pgfscope}%
\begin{pgfscope}%
\pgfpathrectangle{\pgfqpoint{0.407592in}{0.222904in}}{\pgfqpoint{1.411163in}{1.411163in}}%
\pgfusepath{clip}%
\pgfsetbuttcap%
\pgfsetroundjoin%
\definecolor{currentfill}{rgb}{0.909804,0.698039,0.423529}%
\pgfsetfillcolor{currentfill}%
\pgfsetfillopacity{0.120000}%
\pgfsetlinewidth{0.000000pt}%
\definecolor{currentstroke}{rgb}{0.000000,0.000000,0.000000}%
\pgfsetstrokecolor{currentstroke}%
\pgfsetdash{}{0pt}%
\pgfpathmoveto{\pgfqpoint{1.230024in}{1.142482in}}%
\pgfpathlineto{\pgfqpoint{1.245987in}{1.137269in}}%
\pgfpathlineto{\pgfqpoint{1.261989in}{1.132735in}}%
\pgfpathlineto{\pgfqpoint{1.269661in}{1.146685in}}%
\pgfpathlineto{\pgfqpoint{1.277319in}{1.160778in}}%
\pgfpathlineto{\pgfqpoint{1.261331in}{1.165232in}}%
\pgfpathlineto{\pgfqpoint{1.245382in}{1.170334in}}%
\pgfpathlineto{\pgfqpoint{1.237711in}{1.156336in}}%
\pgfpathlineto{\pgfqpoint{1.230024in}{1.142482in}}%
\pgfpathclose%
\pgfusepath{fill}%
\end{pgfscope}%
\begin{pgfscope}%
\pgfpathrectangle{\pgfqpoint{0.407592in}{0.222904in}}{\pgfqpoint{1.411163in}{1.411163in}}%
\pgfusepath{clip}%
\pgfsetbuttcap%
\pgfsetroundjoin%
\definecolor{currentfill}{rgb}{0.909804,0.698039,0.423529}%
\pgfsetfillcolor{currentfill}%
\pgfsetfillopacity{0.120000}%
\pgfsetlinewidth{0.000000pt}%
\definecolor{currentstroke}{rgb}{0.000000,0.000000,0.000000}%
\pgfsetstrokecolor{currentstroke}%
\pgfsetdash{}{0pt}%
\pgfpathmoveto{\pgfqpoint{0.882935in}{1.180860in}}%
\pgfpathlineto{\pgfqpoint{0.898657in}{1.165472in}}%
\pgfpathlineto{\pgfqpoint{0.914392in}{1.150408in}}%
\pgfpathlineto{\pgfqpoint{0.922462in}{1.159521in}}%
\pgfpathlineto{\pgfqpoint{0.930509in}{1.168936in}}%
\pgfpathlineto{\pgfqpoint{0.914806in}{1.183708in}}%
\pgfpathlineto{\pgfqpoint{0.899118in}{1.198808in}}%
\pgfpathlineto{\pgfqpoint{0.891038in}{1.189689in}}%
\pgfpathlineto{\pgfqpoint{0.882935in}{1.180860in}}%
\pgfpathclose%
\pgfusepath{fill}%
\end{pgfscope}%
\begin{pgfscope}%
\pgfpathrectangle{\pgfqpoint{0.407592in}{0.222904in}}{\pgfqpoint{1.411163in}{1.411163in}}%
\pgfusepath{clip}%
\pgfsetbuttcap%
\pgfsetroundjoin%
\definecolor{currentfill}{rgb}{0.909804,0.698039,0.423529}%
\pgfsetfillcolor{currentfill}%
\pgfsetfillopacity{0.120000}%
\pgfsetlinewidth{0.000000pt}%
\definecolor{currentstroke}{rgb}{0.000000,0.000000,0.000000}%
\pgfsetstrokecolor{currentstroke}%
\pgfsetdash{}{0pt}%
\pgfpathmoveto{\pgfqpoint{0.961958in}{1.140489in}}%
\pgfpathlineto{\pgfqpoint{0.977706in}{1.126874in}}%
\pgfpathlineto{\pgfqpoint{0.993470in}{1.113709in}}%
\pgfpathlineto{\pgfqpoint{1.001433in}{1.124043in}}%
\pgfpathlineto{\pgfqpoint{1.009375in}{1.134669in}}%
\pgfpathlineto{\pgfqpoint{0.993639in}{1.147515in}}%
\pgfpathlineto{\pgfqpoint{0.977920in}{1.160811in}}%
\pgfpathlineto{\pgfqpoint{0.969950in}{1.150507in}}%
\pgfpathlineto{\pgfqpoint{0.961958in}{1.140489in}}%
\pgfpathclose%
\pgfusepath{fill}%
\end{pgfscope}%
\begin{pgfscope}%
\pgfpathrectangle{\pgfqpoint{0.407592in}{0.222904in}}{\pgfqpoint{1.411163in}{1.411163in}}%
\pgfusepath{clip}%
\pgfsetbuttcap%
\pgfsetroundjoin%
\definecolor{currentfill}{rgb}{0.909804,0.698039,0.423529}%
\pgfsetfillcolor{currentfill}%
\pgfsetfillopacity{0.120000}%
\pgfsetlinewidth{0.000000pt}%
\definecolor{currentstroke}{rgb}{0.000000,0.000000,0.000000}%
\pgfsetstrokecolor{currentstroke}%
\pgfsetdash{}{0pt}%
\pgfpathmoveto{\pgfqpoint{1.088241in}{1.111769in}}%
\pgfpathlineto{\pgfqpoint{1.104061in}{1.101965in}}%
\pgfpathlineto{\pgfqpoint{1.119907in}{1.092800in}}%
\pgfpathlineto{\pgfqpoint{1.127720in}{1.105113in}}%
\pgfpathlineto{\pgfqpoint{1.135516in}{1.117660in}}%
\pgfpathlineto{\pgfqpoint{1.119690in}{1.126565in}}%
\pgfpathlineto{\pgfqpoint{1.103892in}{1.136089in}}%
\pgfpathlineto{\pgfqpoint{1.096075in}{1.123812in}}%
\pgfpathlineto{\pgfqpoint{1.088241in}{1.111769in}}%
\pgfpathclose%
\pgfusepath{fill}%
\end{pgfscope}%
\begin{pgfscope}%
\pgfpathrectangle{\pgfqpoint{0.407592in}{0.222904in}}{\pgfqpoint{1.411163in}{1.411163in}}%
\pgfusepath{clip}%
\pgfsetbuttcap%
\pgfsetroundjoin%
\definecolor{currentfill}{rgb}{0.909804,0.698039,0.423529}%
\pgfsetfillcolor{currentfill}%
\pgfsetfillopacity{0.120000}%
\pgfsetlinewidth{0.000000pt}%
\definecolor{currentstroke}{rgb}{0.000000,0.000000,0.000000}%
\pgfsetstrokecolor{currentstroke}%
\pgfsetdash{}{0pt}%
\pgfpathmoveto{\pgfqpoint{0.803740in}{1.229394in}}%
\pgfpathlineto{\pgfqpoint{0.819451in}{1.212717in}}%
\pgfpathlineto{\pgfqpoint{0.835173in}{1.196263in}}%
\pgfpathlineto{\pgfqpoint{0.843363in}{1.204235in}}%
\pgfpathlineto{\pgfqpoint{0.851527in}{1.212502in}}%
\pgfpathlineto{\pgfqpoint{0.835842in}{1.228710in}}%
\pgfpathlineto{\pgfqpoint{0.820167in}{1.245148in}}%
\pgfpathlineto{\pgfqpoint{0.811967in}{1.237130in}}%
\pgfpathlineto{\pgfqpoint{0.803740in}{1.229394in}}%
\pgfpathclose%
\pgfusepath{fill}%
\end{pgfscope}%
\begin{pgfscope}%
\pgfpathrectangle{\pgfqpoint{0.407592in}{0.222904in}}{\pgfqpoint{1.411163in}{1.411163in}}%
\pgfusepath{clip}%
\pgfsetbuttcap%
\pgfsetroundjoin%
\definecolor{currentfill}{rgb}{0.909804,0.698039,0.423529}%
\pgfsetfillcolor{currentfill}%
\pgfsetfillopacity{0.120000}%
\pgfsetlinewidth{0.000000pt}%
\definecolor{currentstroke}{rgb}{0.000000,0.000000,0.000000}%
\pgfsetstrokecolor{currentstroke}%
\pgfsetdash{}{0pt}%
\pgfpathmoveto{\pgfqpoint{1.040904in}{1.110457in}}%
\pgfpathlineto{\pgfqpoint{1.056700in}{1.099161in}}%
\pgfpathlineto{\pgfqpoint{1.072518in}{1.088453in}}%
\pgfpathlineto{\pgfqpoint{1.080389in}{1.099977in}}%
\pgfpathlineto{\pgfqpoint{1.088241in}{1.111769in}}%
\pgfpathlineto{\pgfqpoint{1.072446in}{1.122182in}}%
\pgfpathlineto{\pgfqpoint{1.056675in}{1.133169in}}%
\pgfpathlineto{\pgfqpoint{1.048799in}{1.121680in}}%
\pgfpathlineto{\pgfqpoint{1.040904in}{1.110457in}}%
\pgfpathclose%
\pgfusepath{fill}%
\end{pgfscope}%
\begin{pgfscope}%
\pgfpathrectangle{\pgfqpoint{0.407592in}{0.222904in}}{\pgfqpoint{1.411163in}{1.411163in}}%
\pgfusepath{clip}%
\pgfsetbuttcap%
\pgfsetroundjoin%
\definecolor{currentfill}{rgb}{0.909804,0.698039,0.423529}%
\pgfsetfillcolor{currentfill}%
\pgfsetfillopacity{0.120000}%
\pgfsetlinewidth{0.000000pt}%
\definecolor{currentstroke}{rgb}{0.000000,0.000000,0.000000}%
\pgfsetstrokecolor{currentstroke}%
\pgfsetdash{}{0pt}%
\pgfpathmoveto{\pgfqpoint{1.277319in}{1.160778in}}%
\pgfpathlineto{\pgfqpoint{1.293351in}{1.156977in}}%
\pgfpathlineto{\pgfqpoint{1.309427in}{1.153831in}}%
\pgfpathlineto{\pgfqpoint{1.317059in}{1.168089in}}%
\pgfpathlineto{\pgfqpoint{1.300989in}{1.171225in}}%
\pgfpathlineto{\pgfqpoint{1.284964in}{1.175001in}}%
\pgfusepath{fill}%
\end{pgfscope}%
\begin{pgfscope}%
\pgfpathrectangle{\pgfqpoint{0.407592in}{0.222904in}}{\pgfqpoint{1.411163in}{1.411163in}}%
\pgfusepath{clip}%
\pgfsetbuttcap%
\pgfsetroundjoin%
\definecolor{currentfill}{rgb}{0.909804,0.698039,0.423529}%
\pgfsetfillcolor{currentfill}%
\pgfsetfillopacity{0.120000}%
\pgfsetlinewidth{0.000000pt}%
\definecolor{currentstroke}{rgb}{0.000000,0.000000,0.000000}%
\pgfsetstrokecolor{currentstroke}%
\pgfsetdash{}{0pt}%
\pgfpathmoveto{\pgfqpoint{0.914392in}{1.150408in}}%
\pgfpathlineto{\pgfqpoint{0.930142in}{1.135694in}}%
\pgfpathlineto{\pgfqpoint{0.945905in}{1.121363in}}%
\pgfpathlineto{\pgfqpoint{0.953943in}{1.130770in}}%
\pgfpathlineto{\pgfqpoint{0.961958in}{1.140489in}}%
\pgfpathlineto{\pgfqpoint{0.946226in}{1.154520in}}%
\pgfpathlineto{\pgfqpoint{0.930509in}{1.168936in}}%
\pgfpathlineto{\pgfqpoint{0.922462in}{1.159521in}}%
\pgfpathlineto{\pgfqpoint{0.914392in}{1.150408in}}%
\pgfpathclose%
\pgfusepath{fill}%
\end{pgfscope}%
\begin{pgfscope}%
\pgfpathrectangle{\pgfqpoint{0.407592in}{0.222904in}}{\pgfqpoint{1.411163in}{1.411163in}}%
\pgfusepath{clip}%
\pgfsetbuttcap%
\pgfsetroundjoin%
\definecolor{currentfill}{rgb}{0.909804,0.698039,0.423529}%
\pgfsetfillcolor{currentfill}%
\pgfsetfillopacity{0.120000}%
\pgfsetlinewidth{0.000000pt}%
\definecolor{currentstroke}{rgb}{0.000000,0.000000,0.000000}%
\pgfsetstrokecolor{currentstroke}%
\pgfsetdash{}{0pt}%
\pgfpathmoveto{\pgfqpoint{0.724311in}{1.284190in}}%
\pgfpathlineto{\pgfqpoint{0.740019in}{1.266597in}}%
\pgfpathlineto{\pgfqpoint{0.755736in}{1.249158in}}%
\pgfpathlineto{\pgfqpoint{0.764058in}{1.256112in}}%
\pgfpathlineto{\pgfqpoint{0.772350in}{1.263348in}}%
\pgfpathlineto{\pgfqpoint{0.756671in}{1.280592in}}%
\pgfpathlineto{\pgfqpoint{0.741002in}{1.297996in}}%
\pgfpathlineto{\pgfqpoint{0.732671in}{1.290959in}}%
\pgfpathlineto{\pgfqpoint{0.724311in}{1.284190in}}%
\pgfpathclose%
\pgfusepath{fill}%
\end{pgfscope}%
\begin{pgfscope}%
\pgfpathrectangle{\pgfqpoint{0.407592in}{0.222904in}}{\pgfqpoint{1.411163in}{1.411163in}}%
\pgfusepath{clip}%
\pgfsetbuttcap%
\pgfsetroundjoin%
\definecolor{currentfill}{rgb}{0.909804,0.698039,0.423529}%
\pgfsetfillcolor{currentfill}%
\pgfsetfillopacity{0.120000}%
\pgfsetlinewidth{0.000000pt}%
\definecolor{currentstroke}{rgb}{0.000000,0.000000,0.000000}%
\pgfsetstrokecolor{currentstroke}%
\pgfsetdash{}{0pt}%
\pgfpathmoveto{\pgfqpoint{0.993470in}{1.113709in}}%
\pgfpathlineto{\pgfqpoint{1.009253in}{1.101030in}}%
\pgfpathlineto{\pgfqpoint{1.025054in}{1.088872in}}%
\pgfpathlineto{\pgfqpoint{1.032989in}{1.099515in}}%
\pgfpathlineto{\pgfqpoint{1.040904in}{1.110457in}}%
\pgfpathlineto{\pgfqpoint{1.025129in}{1.122305in}}%
\pgfpathlineto{\pgfqpoint{1.009375in}{1.134669in}}%
\pgfpathlineto{\pgfqpoint{1.001433in}{1.124043in}}%
\pgfpathlineto{\pgfqpoint{0.993470in}{1.113709in}}%
\pgfpathclose%
\pgfusepath{fill}%
\end{pgfscope}%
\begin{pgfscope}%
\pgfpathrectangle{\pgfqpoint{0.407592in}{0.222904in}}{\pgfqpoint{1.411163in}{1.411163in}}%
\pgfusepath{clip}%
\pgfsetbuttcap%
\pgfsetroundjoin%
\definecolor{currentfill}{rgb}{0.909804,0.698039,0.423529}%
\pgfsetfillcolor{currentfill}%
\pgfsetfillopacity{0.120000}%
\pgfsetlinewidth{0.000000pt}%
\definecolor{currentstroke}{rgb}{0.000000,0.000000,0.000000}%
\pgfsetstrokecolor{currentstroke}%
\pgfsetdash{}{0pt}%
\pgfpathmoveto{\pgfqpoint{0.835173in}{1.196263in}}%
\pgfpathlineto{\pgfqpoint{0.850906in}{1.180055in}}%
\pgfpathlineto{\pgfqpoint{0.866651in}{1.164113in}}%
\pgfpathlineto{\pgfqpoint{0.874806in}{1.172332in}}%
\pgfpathlineto{\pgfqpoint{0.882935in}{1.180860in}}%
\pgfpathlineto{\pgfqpoint{0.867225in}{1.196545in}}%
\pgfpathlineto{\pgfqpoint{0.851527in}{1.212502in}}%
\pgfpathlineto{\pgfqpoint{0.843363in}{1.204235in}}%
\pgfpathlineto{\pgfqpoint{0.835173in}{1.196263in}}%
\pgfpathclose%
\pgfusepath{fill}%
\end{pgfscope}%
\begin{pgfscope}%
\pgfpathrectangle{\pgfqpoint{0.407592in}{0.222904in}}{\pgfqpoint{1.411163in}{1.411163in}}%
\pgfusepath{clip}%
\pgfsetbuttcap%
\pgfsetroundjoin%
\definecolor{currentfill}{rgb}{0.909804,0.698039,0.423529}%
\pgfsetfillcolor{currentfill}%
\pgfsetfillopacity{0.120000}%
\pgfsetlinewidth{0.000000pt}%
\definecolor{currentstroke}{rgb}{0.000000,0.000000,0.000000}%
\pgfsetstrokecolor{currentstroke}%
\pgfsetdash{}{0pt}%
\pgfpathmoveto{\pgfqpoint{1.167259in}{1.101817in}}%
\pgfpathlineto{\pgfqpoint{1.183179in}{1.094928in}}%
\pgfpathlineto{\pgfqpoint{1.199134in}{1.088752in}}%
\pgfpathlineto{\pgfqpoint{1.206879in}{1.101904in}}%
\pgfpathlineto{\pgfqpoint{1.214609in}{1.115252in}}%
\pgfpathlineto{\pgfqpoint{1.198670in}{1.121278in}}%
\pgfpathlineto{\pgfqpoint{1.182767in}{1.127988in}}%
\pgfpathlineto{\pgfqpoint{1.175021in}{1.114803in}}%
\pgfpathlineto{\pgfqpoint{1.167259in}{1.101817in}}%
\pgfpathclose%
\pgfusepath{fill}%
\end{pgfscope}%
\begin{pgfscope}%
\pgfpathrectangle{\pgfqpoint{0.407592in}{0.222904in}}{\pgfqpoint{1.411163in}{1.411163in}}%
\pgfusepath{clip}%
\pgfsetbuttcap%
\pgfsetroundjoin%
\definecolor{currentfill}{rgb}{0.909804,0.698039,0.423529}%
\pgfsetfillcolor{currentfill}%
\pgfsetfillopacity{0.120000}%
\pgfsetlinewidth{0.000000pt}%
\definecolor{currentstroke}{rgb}{0.000000,0.000000,0.000000}%
\pgfsetstrokecolor{currentstroke}%
\pgfsetdash{}{0pt}%
\pgfpathmoveto{\pgfqpoint{1.119907in}{1.092800in}}%
\pgfpathlineto{\pgfqpoint{1.135782in}{1.084305in}}%
\pgfpathlineto{\pgfqpoint{1.151687in}{1.076507in}}%
\pgfpathlineto{\pgfqpoint{1.159481in}{1.089047in}}%
\pgfpathlineto{\pgfqpoint{1.167259in}{1.101817in}}%
\pgfpathlineto{\pgfqpoint{1.151372in}{1.109402in}}%
\pgfpathlineto{\pgfqpoint{1.135516in}{1.117660in}}%
\pgfpathlineto{\pgfqpoint{1.127720in}{1.105113in}}%
\pgfpathlineto{\pgfqpoint{1.119907in}{1.092800in}}%
\pgfpathclose%
\pgfusepath{fill}%
\end{pgfscope}%
\begin{pgfscope}%
\pgfpathrectangle{\pgfqpoint{0.407592in}{0.222904in}}{\pgfqpoint{1.411163in}{1.411163in}}%
\pgfusepath{clip}%
\pgfsetbuttcap%
\pgfsetroundjoin%
\definecolor{currentfill}{rgb}{0.909804,0.698039,0.423529}%
\pgfsetfillcolor{currentfill}%
\pgfsetfillopacity{0.120000}%
\pgfsetlinewidth{0.000000pt}%
\definecolor{currentstroke}{rgb}{0.000000,0.000000,0.000000}%
\pgfsetstrokecolor{currentstroke}%
\pgfsetdash{}{0pt}%
\pgfpathmoveto{\pgfqpoint{1.214609in}{1.115252in}}%
\pgfpathlineto{\pgfqpoint{1.230586in}{1.109925in}}%
\pgfpathlineto{\pgfqpoint{1.246603in}{1.105308in}}%
\pgfpathlineto{\pgfqpoint{1.254303in}{1.118938in}}%
\pgfpathlineto{\pgfqpoint{1.261989in}{1.132735in}}%
\pgfpathlineto{\pgfqpoint{1.245987in}{1.137269in}}%
\pgfpathlineto{\pgfqpoint{1.230024in}{1.142482in}}%
\pgfpathlineto{\pgfqpoint{1.222324in}{1.128783in}}%
\pgfpathlineto{\pgfqpoint{1.214609in}{1.115252in}}%
\pgfpathclose%
\pgfusepath{fill}%
\end{pgfscope}%
\begin{pgfscope}%
\pgfpathrectangle{\pgfqpoint{0.407592in}{0.222904in}}{\pgfqpoint{1.411163in}{1.411163in}}%
\pgfusepath{clip}%
\pgfsetbuttcap%
\pgfsetroundjoin%
\definecolor{currentfill}{rgb}{0.909804,0.698039,0.423529}%
\pgfsetfillcolor{currentfill}%
\pgfsetfillopacity{0.120000}%
\pgfsetlinewidth{0.000000pt}%
\definecolor{currentstroke}{rgb}{0.000000,0.000000,0.000000}%
\pgfsetstrokecolor{currentstroke}%
\pgfsetdash{}{0pt}%
\pgfpathmoveto{\pgfqpoint{0.755736in}{1.249158in}}%
\pgfpathlineto{\pgfqpoint{0.771464in}{1.231887in}}%
\pgfpathlineto{\pgfqpoint{0.787202in}{1.214799in}}%
\pgfpathlineto{\pgfqpoint{0.795485in}{1.221948in}}%
\pgfpathlineto{\pgfqpoint{0.803740in}{1.229394in}}%
\pgfpathlineto{\pgfqpoint{0.788040in}{1.246277in}}%
\pgfpathlineto{\pgfqpoint{0.772350in}{1.263348in}}%
\pgfpathlineto{\pgfqpoint{0.764058in}{1.256112in}}%
\pgfpathlineto{\pgfqpoint{0.755736in}{1.249158in}}%
\pgfpathclose%
\pgfusepath{fill}%
\end{pgfscope}%
\begin{pgfscope}%
\pgfpathrectangle{\pgfqpoint{0.407592in}{0.222904in}}{\pgfqpoint{1.411163in}{1.411163in}}%
\pgfusepath{clip}%
\pgfsetbuttcap%
\pgfsetroundjoin%
\definecolor{currentfill}{rgb}{0.909804,0.698039,0.423529}%
\pgfsetfillcolor{currentfill}%
\pgfsetfillopacity{0.120000}%
\pgfsetlinewidth{0.000000pt}%
\definecolor{currentstroke}{rgb}{0.000000,0.000000,0.000000}%
\pgfsetstrokecolor{currentstroke}%
\pgfsetdash{}{0pt}%
\pgfpathmoveto{\pgfqpoint{1.261989in}{1.132735in}}%
\pgfpathlineto{\pgfqpoint{1.278034in}{1.128884in}}%
\pgfpathlineto{\pgfqpoint{1.294124in}{1.125720in}}%
\pgfpathlineto{\pgfqpoint{1.301782in}{1.139705in}}%
\pgfpathlineto{\pgfqpoint{1.309427in}{1.153831in}}%
\pgfpathlineto{\pgfqpoint{1.293351in}{1.156977in}}%
\pgfpathlineto{\pgfqpoint{1.277319in}{1.160778in}}%
\pgfpathlineto{\pgfqpoint{1.269661in}{1.146685in}}%
\pgfpathlineto{\pgfqpoint{1.261989in}{1.132735in}}%
\pgfpathclose%
\pgfusepath{fill}%
\end{pgfscope}%
\begin{pgfscope}%
\pgfpathrectangle{\pgfqpoint{0.407592in}{0.222904in}}{\pgfqpoint{1.411163in}{1.411163in}}%
\pgfusepath{clip}%
\pgfsetbuttcap%
\pgfsetroundjoin%
\definecolor{currentfill}{rgb}{0.909804,0.698039,0.423529}%
\pgfsetfillcolor{currentfill}%
\pgfsetfillopacity{0.120000}%
\pgfsetlinewidth{0.000000pt}%
\definecolor{currentstroke}{rgb}{0.000000,0.000000,0.000000}%
\pgfsetstrokecolor{currentstroke}%
\pgfsetdash{}{0pt}%
\pgfpathmoveto{\pgfqpoint{1.072518in}{1.088453in}}%
\pgfpathlineto{\pgfqpoint{1.088360in}{1.078367in}}%
\pgfpathlineto{\pgfqpoint{1.104228in}{1.068940in}}%
\pgfpathlineto{\pgfqpoint{1.112077in}{1.080736in}}%
\pgfpathlineto{\pgfqpoint{1.119907in}{1.092800in}}%
\pgfpathlineto{\pgfqpoint{1.104061in}{1.101965in}}%
\pgfpathlineto{\pgfqpoint{1.088241in}{1.111769in}}%
\pgfpathlineto{\pgfqpoint{1.080389in}{1.099977in}}%
\pgfpathlineto{\pgfqpoint{1.072518in}{1.088453in}}%
\pgfpathclose%
\pgfusepath{fill}%
\end{pgfscope}%
\begin{pgfscope}%
\pgfpathrectangle{\pgfqpoint{0.407592in}{0.222904in}}{\pgfqpoint{1.411163in}{1.411163in}}%
\pgfusepath{clip}%
\pgfsetbuttcap%
\pgfsetroundjoin%
\definecolor{currentfill}{rgb}{0.909804,0.698039,0.423529}%
\pgfsetfillcolor{currentfill}%
\pgfsetfillopacity{0.120000}%
\pgfsetlinewidth{0.000000pt}%
\definecolor{currentstroke}{rgb}{0.000000,0.000000,0.000000}%
\pgfsetstrokecolor{currentstroke}%
\pgfsetdash{}{0pt}%
\pgfpathmoveto{\pgfqpoint{0.945905in}{1.121363in}}%
\pgfpathlineto{\pgfqpoint{0.961684in}{1.107445in}}%
\pgfpathlineto{\pgfqpoint{0.977479in}{1.093978in}}%
\pgfpathlineto{\pgfqpoint{0.985486in}{1.103683in}}%
\pgfpathlineto{\pgfqpoint{0.993470in}{1.113709in}}%
\pgfpathlineto{\pgfqpoint{0.977706in}{1.126874in}}%
\pgfpathlineto{\pgfqpoint{0.961958in}{1.140489in}}%
\pgfpathlineto{\pgfqpoint{0.953943in}{1.130770in}}%
\pgfpathlineto{\pgfqpoint{0.945905in}{1.121363in}}%
\pgfpathclose%
\pgfusepath{fill}%
\end{pgfscope}%
\begin{pgfscope}%
\pgfpathrectangle{\pgfqpoint{0.407592in}{0.222904in}}{\pgfqpoint{1.411163in}{1.411163in}}%
\pgfusepath{clip}%
\pgfsetbuttcap%
\pgfsetroundjoin%
\definecolor{currentfill}{rgb}{0.909804,0.698039,0.423529}%
\pgfsetfillcolor{currentfill}%
\pgfsetfillopacity{0.120000}%
\pgfsetlinewidth{0.000000pt}%
\definecolor{currentstroke}{rgb}{0.000000,0.000000,0.000000}%
\pgfsetstrokecolor{currentstroke}%
\pgfsetdash{}{0pt}%
\pgfpathmoveto{\pgfqpoint{0.866651in}{1.164113in}}%
\pgfpathlineto{\pgfqpoint{0.882408in}{1.148462in}}%
\pgfpathlineto{\pgfqpoint{0.898178in}{1.133129in}}%
\pgfpathlineto{\pgfqpoint{0.906298in}{1.141606in}}%
\pgfpathlineto{\pgfqpoint{0.914392in}{1.150408in}}%
\pgfpathlineto{\pgfqpoint{0.898657in}{1.165472in}}%
\pgfpathlineto{\pgfqpoint{0.882935in}{1.180860in}}%
\pgfpathlineto{\pgfqpoint{0.874806in}{1.172332in}}%
\pgfpathlineto{\pgfqpoint{0.866651in}{1.164113in}}%
\pgfpathclose%
\pgfusepath{fill}%
\end{pgfscope}%
\begin{pgfscope}%
\pgfpathrectangle{\pgfqpoint{0.407592in}{0.222904in}}{\pgfqpoint{1.411163in}{1.411163in}}%
\pgfusepath{clip}%
\pgfsetbuttcap%
\pgfsetroundjoin%
\definecolor{currentfill}{rgb}{0.909804,0.698039,0.423529}%
\pgfsetfillcolor{currentfill}%
\pgfsetfillopacity{0.120000}%
\pgfsetlinewidth{0.000000pt}%
\definecolor{currentstroke}{rgb}{0.000000,0.000000,0.000000}%
\pgfsetstrokecolor{currentstroke}%
\pgfsetdash{}{0pt}%
\pgfpathmoveto{\pgfqpoint{1.025054in}{1.088872in}}%
\pgfpathlineto{\pgfqpoint{1.040876in}{1.077273in}}%
\pgfpathlineto{\pgfqpoint{1.056719in}{1.066273in}}%
\pgfpathlineto{\pgfqpoint{1.064628in}{1.077212in}}%
\pgfpathlineto{\pgfqpoint{1.072518in}{1.088453in}}%
\pgfpathlineto{\pgfqpoint{1.056700in}{1.099161in}}%
\pgfpathlineto{\pgfqpoint{1.040904in}{1.110457in}}%
\pgfpathlineto{\pgfqpoint{1.032989in}{1.099515in}}%
\pgfpathlineto{\pgfqpoint{1.025054in}{1.088872in}}%
\pgfpathclose%
\pgfusepath{fill}%
\end{pgfscope}%
\begin{pgfscope}%
\pgfpathrectangle{\pgfqpoint{0.407592in}{0.222904in}}{\pgfqpoint{1.411163in}{1.411163in}}%
\pgfusepath{clip}%
\pgfsetbuttcap%
\pgfsetroundjoin%
\definecolor{currentfill}{rgb}{0.909804,0.698039,0.423529}%
\pgfsetfillcolor{currentfill}%
\pgfsetfillopacity{0.120000}%
\pgfsetlinewidth{0.000000pt}%
\definecolor{currentstroke}{rgb}{0.000000,0.000000,0.000000}%
\pgfsetstrokecolor{currentstroke}%
\pgfsetdash{}{0pt}%
\pgfpathmoveto{\pgfqpoint{0.787202in}{1.214799in}}%
\pgfpathlineto{\pgfqpoint{0.802950in}{1.197910in}}%
\pgfpathlineto{\pgfqpoint{0.818710in}{1.181241in}}%
\pgfpathlineto{\pgfqpoint{0.826955in}{1.188596in}}%
\pgfpathlineto{\pgfqpoint{0.835173in}{1.196263in}}%
\pgfpathlineto{\pgfqpoint{0.819451in}{1.212717in}}%
\pgfpathlineto{\pgfqpoint{0.803740in}{1.229394in}}%
\pgfpathlineto{\pgfqpoint{0.795485in}{1.221948in}}%
\pgfpathlineto{\pgfqpoint{0.787202in}{1.214799in}}%
\pgfpathclose%
\pgfusepath{fill}%
\end{pgfscope}%
\begin{pgfscope}%
\pgfpathrectangle{\pgfqpoint{0.407592in}{0.222904in}}{\pgfqpoint{1.411163in}{1.411163in}}%
\pgfusepath{clip}%
\pgfsetbuttcap%
\pgfsetroundjoin%
\definecolor{currentfill}{rgb}{0.909804,0.698039,0.423529}%
\pgfsetfillcolor{currentfill}%
\pgfsetfillopacity{0.120000}%
\pgfsetlinewidth{0.000000pt}%
\definecolor{currentstroke}{rgb}{0.000000,0.000000,0.000000}%
\pgfsetstrokecolor{currentstroke}%
\pgfsetdash{}{0pt}%
\pgfpathmoveto{\pgfqpoint{1.309427in}{1.153831in}}%
\pgfpathlineto{\pgfqpoint{1.325550in}{1.151341in}}%
\pgfpathlineto{\pgfqpoint{1.341720in}{1.149504in}}%
\pgfpathlineto{\pgfqpoint{1.349339in}{1.163737in}}%
\pgfpathlineto{\pgfqpoint{1.333175in}{1.165594in}}%
\pgfpathlineto{\pgfqpoint{1.317059in}{1.168089in}}%
\pgfusepath{fill}%
\end{pgfscope}%
\begin{pgfscope}%
\pgfpathrectangle{\pgfqpoint{0.407592in}{0.222904in}}{\pgfqpoint{1.411163in}{1.411163in}}%
\pgfusepath{clip}%
\pgfsetbuttcap%
\pgfsetroundjoin%
\definecolor{currentfill}{rgb}{0.909804,0.698039,0.423529}%
\pgfsetfillcolor{currentfill}%
\pgfsetfillopacity{0.120000}%
\pgfsetlinewidth{0.000000pt}%
\definecolor{currentstroke}{rgb}{0.000000,0.000000,0.000000}%
\pgfsetstrokecolor{currentstroke}%
\pgfsetdash{}{0pt}%
\pgfpathmoveto{\pgfqpoint{0.898178in}{1.133129in}}%
\pgfpathlineto{\pgfqpoint{0.913961in}{1.118143in}}%
\pgfpathlineto{\pgfqpoint{0.929758in}{1.103535in}}%
\pgfpathlineto{\pgfqpoint{0.937844in}{1.112280in}}%
\pgfpathlineto{\pgfqpoint{0.945905in}{1.121363in}}%
\pgfpathlineto{\pgfqpoint{0.930142in}{1.135694in}}%
\pgfpathlineto{\pgfqpoint{0.914392in}{1.150408in}}%
\pgfpathlineto{\pgfqpoint{0.906298in}{1.141606in}}%
\pgfpathlineto{\pgfqpoint{0.898178in}{1.133129in}}%
\pgfpathclose%
\pgfusepath{fill}%
\end{pgfscope}%
\begin{pgfscope}%
\pgfpathrectangle{\pgfqpoint{0.407592in}{0.222904in}}{\pgfqpoint{1.411163in}{1.411163in}}%
\pgfusepath{clip}%
\pgfsetbuttcap%
\pgfsetroundjoin%
\definecolor{currentfill}{rgb}{0.909804,0.698039,0.423529}%
\pgfsetfillcolor{currentfill}%
\pgfsetfillopacity{0.120000}%
\pgfsetlinewidth{0.000000pt}%
\definecolor{currentstroke}{rgb}{0.000000,0.000000,0.000000}%
\pgfsetstrokecolor{currentstroke}%
\pgfsetdash{}{0pt}%
\pgfpathmoveto{\pgfqpoint{0.977479in}{1.093978in}}%
\pgfpathlineto{\pgfqpoint{0.993291in}{1.080998in}}%
\pgfpathlineto{\pgfqpoint{1.009121in}{1.068543in}}%
\pgfpathlineto{\pgfqpoint{1.017098in}{1.078542in}}%
\pgfpathlineto{\pgfqpoint{1.025054in}{1.088872in}}%
\pgfpathlineto{\pgfqpoint{1.009253in}{1.101030in}}%
\pgfpathlineto{\pgfqpoint{0.993470in}{1.113709in}}%
\pgfpathlineto{\pgfqpoint{0.985486in}{1.103683in}}%
\pgfpathlineto{\pgfqpoint{0.977479in}{1.093978in}}%
\pgfpathclose%
\pgfusepath{fill}%
\end{pgfscope}%
\begin{pgfscope}%
\pgfpathrectangle{\pgfqpoint{0.407592in}{0.222904in}}{\pgfqpoint{1.411163in}{1.411163in}}%
\pgfusepath{clip}%
\pgfsetbuttcap%
\pgfsetroundjoin%
\definecolor{currentfill}{rgb}{0.909804,0.698039,0.423529}%
\pgfsetfillcolor{currentfill}%
\pgfsetfillopacity{0.120000}%
\pgfsetlinewidth{0.000000pt}%
\definecolor{currentstroke}{rgb}{0.000000,0.000000,0.000000}%
\pgfsetstrokecolor{currentstroke}%
\pgfsetdash{}{0pt}%
\pgfpathmoveto{\pgfqpoint{1.151687in}{1.076507in}}%
\pgfpathlineto{\pgfqpoint{1.167625in}{1.069434in}}%
\pgfpathlineto{\pgfqpoint{1.183598in}{1.063106in}}%
\pgfpathlineto{\pgfqpoint{1.191374in}{1.075814in}}%
\pgfpathlineto{\pgfqpoint{1.199134in}{1.088752in}}%
\pgfpathlineto{\pgfqpoint{1.183179in}{1.094928in}}%
\pgfpathlineto{\pgfqpoint{1.167259in}{1.101817in}}%
\pgfpathlineto{\pgfqpoint{1.159481in}{1.089047in}}%
\pgfpathlineto{\pgfqpoint{1.151687in}{1.076507in}}%
\pgfpathclose%
\pgfusepath{fill}%
\end{pgfscope}%
\begin{pgfscope}%
\pgfpathrectangle{\pgfqpoint{0.407592in}{0.222904in}}{\pgfqpoint{1.411163in}{1.411163in}}%
\pgfusepath{clip}%
\pgfsetbuttcap%
\pgfsetroundjoin%
\definecolor{currentfill}{rgb}{0.909804,0.698039,0.423529}%
\pgfsetfillcolor{currentfill}%
\pgfsetfillopacity{0.120000}%
\pgfsetlinewidth{0.000000pt}%
\definecolor{currentstroke}{rgb}{0.000000,0.000000,0.000000}%
\pgfsetstrokecolor{currentstroke}%
\pgfsetdash{}{0pt}%
\pgfpathmoveto{\pgfqpoint{1.199134in}{1.088752in}}%
\pgfpathlineto{\pgfqpoint{1.215127in}{1.083308in}}%
\pgfpathlineto{\pgfqpoint{1.231159in}{1.078608in}}%
\pgfpathlineto{\pgfqpoint{1.238888in}{1.091860in}}%
\pgfpathlineto{\pgfqpoint{1.246603in}{1.105308in}}%
\pgfpathlineto{\pgfqpoint{1.230586in}{1.109925in}}%
\pgfpathlineto{\pgfqpoint{1.214609in}{1.115252in}}%
\pgfpathlineto{\pgfqpoint{1.206879in}{1.101904in}}%
\pgfpathlineto{\pgfqpoint{1.199134in}{1.088752in}}%
\pgfpathclose%
\pgfusepath{fill}%
\end{pgfscope}%
\begin{pgfscope}%
\pgfpathrectangle{\pgfqpoint{0.407592in}{0.222904in}}{\pgfqpoint{1.411163in}{1.411163in}}%
\pgfusepath{clip}%
\pgfsetbuttcap%
\pgfsetroundjoin%
\definecolor{currentfill}{rgb}{0.909804,0.698039,0.423529}%
\pgfsetfillcolor{currentfill}%
\pgfsetfillopacity{0.120000}%
\pgfsetlinewidth{0.000000pt}%
\definecolor{currentstroke}{rgb}{0.000000,0.000000,0.000000}%
\pgfsetstrokecolor{currentstroke}%
\pgfsetdash{}{0pt}%
\pgfpathmoveto{\pgfqpoint{0.707502in}{1.271476in}}%
\pgfpathlineto{\pgfqpoint{0.723249in}{1.253721in}}%
\pgfpathlineto{\pgfqpoint{0.739005in}{1.236117in}}%
\pgfpathlineto{\pgfqpoint{0.747386in}{1.242491in}}%
\pgfpathlineto{\pgfqpoint{0.755736in}{1.249158in}}%
\pgfpathlineto{\pgfqpoint{0.740019in}{1.266597in}}%
\pgfpathlineto{\pgfqpoint{0.724311in}{1.284190in}}%
\pgfpathlineto{\pgfqpoint{0.715922in}{1.277694in}}%
\pgfpathlineto{\pgfqpoint{0.707502in}{1.271476in}}%
\pgfpathclose%
\pgfusepath{fill}%
\end{pgfscope}%
\begin{pgfscope}%
\pgfpathrectangle{\pgfqpoint{0.407592in}{0.222904in}}{\pgfqpoint{1.411163in}{1.411163in}}%
\pgfusepath{clip}%
\pgfsetbuttcap%
\pgfsetroundjoin%
\definecolor{currentfill}{rgb}{0.909804,0.698039,0.423529}%
\pgfsetfillcolor{currentfill}%
\pgfsetfillopacity{0.120000}%
\pgfsetlinewidth{0.000000pt}%
\definecolor{currentstroke}{rgb}{0.000000,0.000000,0.000000}%
\pgfsetstrokecolor{currentstroke}%
\pgfsetdash{}{0pt}%
\pgfpathmoveto{\pgfqpoint{1.104228in}{1.068940in}}%
\pgfpathlineto{\pgfqpoint{1.120124in}{1.060205in}}%
\pgfpathlineto{\pgfqpoint{1.136050in}{1.052193in}}%
\pgfpathlineto{\pgfqpoint{1.143877in}{1.064217in}}%
\pgfpathlineto{\pgfqpoint{1.151687in}{1.076507in}}%
\pgfpathlineto{\pgfqpoint{1.135782in}{1.084305in}}%
\pgfpathlineto{\pgfqpoint{1.119907in}{1.092800in}}%
\pgfpathlineto{\pgfqpoint{1.112077in}{1.080736in}}%
\pgfpathlineto{\pgfqpoint{1.104228in}{1.068940in}}%
\pgfpathclose%
\pgfusepath{fill}%
\end{pgfscope}%
\begin{pgfscope}%
\pgfpathrectangle{\pgfqpoint{0.407592in}{0.222904in}}{\pgfqpoint{1.411163in}{1.411163in}}%
\pgfusepath{clip}%
\pgfsetbuttcap%
\pgfsetroundjoin%
\definecolor{currentfill}{rgb}{0.909804,0.698039,0.423529}%
\pgfsetfillcolor{currentfill}%
\pgfsetfillopacity{0.120000}%
\pgfsetlinewidth{0.000000pt}%
\definecolor{currentstroke}{rgb}{0.000000,0.000000,0.000000}%
\pgfsetstrokecolor{currentstroke}%
\pgfsetdash{}{0pt}%
\pgfpathmoveto{\pgfqpoint{0.818710in}{1.181241in}}%
\pgfpathlineto{\pgfqpoint{0.834480in}{1.164810in}}%
\pgfpathlineto{\pgfqpoint{0.850261in}{1.148640in}}%
\pgfpathlineto{\pgfqpoint{0.858470in}{1.156213in}}%
\pgfpathlineto{\pgfqpoint{0.866651in}{1.164113in}}%
\pgfpathlineto{\pgfqpoint{0.850906in}{1.180055in}}%
\pgfpathlineto{\pgfqpoint{0.835173in}{1.196263in}}%
\pgfpathlineto{\pgfqpoint{0.826955in}{1.188596in}}%
\pgfpathlineto{\pgfqpoint{0.818710in}{1.181241in}}%
\pgfpathclose%
\pgfusepath{fill}%
\end{pgfscope}%
\begin{pgfscope}%
\pgfpathrectangle{\pgfqpoint{0.407592in}{0.222904in}}{\pgfqpoint{1.411163in}{1.411163in}}%
\pgfusepath{clip}%
\pgfsetbuttcap%
\pgfsetroundjoin%
\definecolor{currentfill}{rgb}{0.909804,0.698039,0.423529}%
\pgfsetfillcolor{currentfill}%
\pgfsetfillopacity{0.120000}%
\pgfsetlinewidth{0.000000pt}%
\definecolor{currentstroke}{rgb}{0.000000,0.000000,0.000000}%
\pgfsetstrokecolor{currentstroke}%
\pgfsetdash{}{0pt}%
\pgfpathmoveto{\pgfqpoint{1.246603in}{1.105308in}}%
\pgfpathlineto{\pgfqpoint{1.262662in}{1.101408in}}%
\pgfpathlineto{\pgfqpoint{1.278765in}{1.098228in}}%
\pgfpathlineto{\pgfqpoint{1.286451in}{1.111890in}}%
\pgfpathlineto{\pgfqpoint{1.294124in}{1.125720in}}%
\pgfpathlineto{\pgfqpoint{1.278034in}{1.128884in}}%
\pgfpathlineto{\pgfqpoint{1.261989in}{1.132735in}}%
\pgfpathlineto{\pgfqpoint{1.254303in}{1.118938in}}%
\pgfpathlineto{\pgfqpoint{1.246603in}{1.105308in}}%
\pgfpathclose%
\pgfusepath{fill}%
\end{pgfscope}%
\begin{pgfscope}%
\pgfpathrectangle{\pgfqpoint{0.407592in}{0.222904in}}{\pgfqpoint{1.411163in}{1.411163in}}%
\pgfusepath{clip}%
\pgfsetbuttcap%
\pgfsetroundjoin%
\definecolor{currentfill}{rgb}{0.909804,0.698039,0.423529}%
\pgfsetfillcolor{currentfill}%
\pgfsetfillopacity{0.120000}%
\pgfsetlinewidth{0.000000pt}%
\definecolor{currentstroke}{rgb}{0.000000,0.000000,0.000000}%
\pgfsetstrokecolor{currentstroke}%
\pgfsetdash{}{0pt}%
\pgfpathmoveto{\pgfqpoint{1.056719in}{1.066273in}}%
\pgfpathlineto{\pgfqpoint{1.072586in}{1.055911in}}%
\pgfpathlineto{\pgfqpoint{1.088477in}{1.046225in}}%
\pgfpathlineto{\pgfqpoint{1.096362in}{1.057430in}}%
\pgfpathlineto{\pgfqpoint{1.104228in}{1.068940in}}%
\pgfpathlineto{\pgfqpoint{1.088360in}{1.078367in}}%
\pgfpathlineto{\pgfqpoint{1.072518in}{1.088453in}}%
\pgfpathlineto{\pgfqpoint{1.064628in}{1.077212in}}%
\pgfpathlineto{\pgfqpoint{1.056719in}{1.066273in}}%
\pgfpathclose%
\pgfusepath{fill}%
\end{pgfscope}%
\begin{pgfscope}%
\pgfpathrectangle{\pgfqpoint{0.407592in}{0.222904in}}{\pgfqpoint{1.411163in}{1.411163in}}%
\pgfusepath{clip}%
\pgfsetbuttcap%
\pgfsetroundjoin%
\definecolor{currentfill}{rgb}{0.909804,0.698039,0.423529}%
\pgfsetfillcolor{currentfill}%
\pgfsetfillopacity{0.120000}%
\pgfsetlinewidth{0.000000pt}%
\definecolor{currentstroke}{rgb}{0.000000,0.000000,0.000000}%
\pgfsetstrokecolor{currentstroke}%
\pgfsetdash{}{0pt}%
\pgfpathmoveto{\pgfqpoint{1.294124in}{1.125720in}}%
\pgfpathlineto{\pgfqpoint{1.310259in}{1.123243in}}%
\pgfpathlineto{\pgfqpoint{1.326443in}{1.121449in}}%
\pgfpathlineto{\pgfqpoint{1.334089in}{1.135404in}}%
\pgfpathlineto{\pgfqpoint{1.341720in}{1.149504in}}%
\pgfpathlineto{\pgfqpoint{1.325550in}{1.151341in}}%
\pgfpathlineto{\pgfqpoint{1.309427in}{1.153831in}}%
\pgfpathlineto{\pgfqpoint{1.301782in}{1.139705in}}%
\pgfpathlineto{\pgfqpoint{1.294124in}{1.125720in}}%
\pgfpathclose%
\pgfusepath{fill}%
\end{pgfscope}%
\begin{pgfscope}%
\pgfpathrectangle{\pgfqpoint{0.407592in}{0.222904in}}{\pgfqpoint{1.411163in}{1.411163in}}%
\pgfusepath{clip}%
\pgfsetbuttcap%
\pgfsetroundjoin%
\definecolor{currentfill}{rgb}{0.909804,0.698039,0.423529}%
\pgfsetfillcolor{currentfill}%
\pgfsetfillopacity{0.120000}%
\pgfsetlinewidth{0.000000pt}%
\definecolor{currentstroke}{rgb}{0.000000,0.000000,0.000000}%
\pgfsetstrokecolor{currentstroke}%
\pgfsetdash{}{0pt}%
\pgfpathmoveto{\pgfqpoint{0.929758in}{1.103535in}}%
\pgfpathlineto{\pgfqpoint{0.945569in}{1.089340in}}%
\pgfpathlineto{\pgfqpoint{0.961396in}{1.075593in}}%
\pgfpathlineto{\pgfqpoint{0.969449in}{1.084610in}}%
\pgfpathlineto{\pgfqpoint{0.977479in}{1.093978in}}%
\pgfpathlineto{\pgfqpoint{0.961684in}{1.107445in}}%
\pgfpathlineto{\pgfqpoint{0.945905in}{1.121363in}}%
\pgfpathlineto{\pgfqpoint{0.937844in}{1.112280in}}%
\pgfpathlineto{\pgfqpoint{0.929758in}{1.103535in}}%
\pgfpathclose%
\pgfusepath{fill}%
\end{pgfscope}%
\begin{pgfscope}%
\pgfpathrectangle{\pgfqpoint{0.407592in}{0.222904in}}{\pgfqpoint{1.411163in}{1.411163in}}%
\pgfusepath{clip}%
\pgfsetbuttcap%
\pgfsetroundjoin%
\definecolor{currentfill}{rgb}{0.909804,0.698039,0.423529}%
\pgfsetfillcolor{currentfill}%
\pgfsetfillopacity{0.120000}%
\pgfsetlinewidth{0.000000pt}%
\definecolor{currentstroke}{rgb}{0.000000,0.000000,0.000000}%
\pgfsetstrokecolor{currentstroke}%
\pgfsetdash{}{0pt}%
\pgfpathmoveto{\pgfqpoint{0.739005in}{1.236117in}}%
\pgfpathlineto{\pgfqpoint{0.754772in}{1.218675in}}%
\pgfpathlineto{\pgfqpoint{0.770549in}{1.201413in}}%
\pgfpathlineto{\pgfqpoint{0.778890in}{1.207952in}}%
\pgfpathlineto{\pgfqpoint{0.787202in}{1.214799in}}%
\pgfpathlineto{\pgfqpoint{0.771464in}{1.231887in}}%
\pgfpathlineto{\pgfqpoint{0.755736in}{1.249158in}}%
\pgfpathlineto{\pgfqpoint{0.747386in}{1.242491in}}%
\pgfpathlineto{\pgfqpoint{0.739005in}{1.236117in}}%
\pgfpathclose%
\pgfusepath{fill}%
\end{pgfscope}%
\begin{pgfscope}%
\pgfpathrectangle{\pgfqpoint{0.407592in}{0.222904in}}{\pgfqpoint{1.411163in}{1.411163in}}%
\pgfusepath{clip}%
\pgfsetbuttcap%
\pgfsetroundjoin%
\definecolor{currentfill}{rgb}{0.909804,0.698039,0.423529}%
\pgfsetfillcolor{currentfill}%
\pgfsetfillopacity{0.120000}%
\pgfsetlinewidth{0.000000pt}%
\definecolor{currentstroke}{rgb}{0.000000,0.000000,0.000000}%
\pgfsetstrokecolor{currentstroke}%
\pgfsetdash{}{0pt}%
\pgfpathmoveto{\pgfqpoint{0.850261in}{1.148640in}}%
\pgfpathlineto{\pgfqpoint{0.866055in}{1.132757in}}%
\pgfpathlineto{\pgfqpoint{0.881860in}{1.117187in}}%
\pgfpathlineto{\pgfqpoint{0.890032in}{1.124986in}}%
\pgfpathlineto{\pgfqpoint{0.898178in}{1.133129in}}%
\pgfpathlineto{\pgfqpoint{0.882408in}{1.148462in}}%
\pgfpathlineto{\pgfqpoint{0.866651in}{1.164113in}}%
\pgfpathlineto{\pgfqpoint{0.858470in}{1.156213in}}%
\pgfpathlineto{\pgfqpoint{0.850261in}{1.148640in}}%
\pgfpathclose%
\pgfusepath{fill}%
\end{pgfscope}%
\begin{pgfscope}%
\pgfpathrectangle{\pgfqpoint{0.407592in}{0.222904in}}{\pgfqpoint{1.411163in}{1.411163in}}%
\pgfusepath{clip}%
\pgfsetbuttcap%
\pgfsetroundjoin%
\definecolor{currentfill}{rgb}{0.909804,0.698039,0.423529}%
\pgfsetfillcolor{currentfill}%
\pgfsetfillopacity{0.120000}%
\pgfsetlinewidth{0.000000pt}%
\definecolor{currentstroke}{rgb}{0.000000,0.000000,0.000000}%
\pgfsetstrokecolor{currentstroke}%
\pgfsetdash{}{0pt}%
\pgfpathmoveto{\pgfqpoint{1.009121in}{1.068543in}}%
\pgfpathlineto{\pgfqpoint{1.024970in}{1.056655in}}%
\pgfpathlineto{\pgfqpoint{1.040840in}{1.045374in}}%
\pgfpathlineto{\pgfqpoint{1.048790in}{1.055655in}}%
\pgfpathlineto{\pgfqpoint{1.056719in}{1.066273in}}%
\pgfpathlineto{\pgfqpoint{1.040876in}{1.077273in}}%
\pgfpathlineto{\pgfqpoint{1.025054in}{1.088872in}}%
\pgfpathlineto{\pgfqpoint{1.017098in}{1.078542in}}%
\pgfpathlineto{\pgfqpoint{1.009121in}{1.068543in}}%
\pgfpathclose%
\pgfusepath{fill}%
\end{pgfscope}%
\begin{pgfscope}%
\pgfpathrectangle{\pgfqpoint{0.407592in}{0.222904in}}{\pgfqpoint{1.411163in}{1.411163in}}%
\pgfusepath{clip}%
\pgfsetbuttcap%
\pgfsetroundjoin%
\definecolor{currentfill}{rgb}{0.909804,0.698039,0.423529}%
\pgfsetfillcolor{currentfill}%
\pgfsetfillopacity{0.120000}%
\pgfsetlinewidth{0.000000pt}%
\definecolor{currentstroke}{rgb}{0.000000,0.000000,0.000000}%
\pgfsetstrokecolor{currentstroke}%
\pgfsetdash{}{0pt}%
\pgfpathmoveto{\pgfqpoint{0.770549in}{1.201413in}}%
\pgfpathlineto{\pgfqpoint{0.786336in}{1.184345in}}%
\pgfpathlineto{\pgfqpoint{0.802133in}{1.167492in}}%
\pgfpathlineto{\pgfqpoint{0.810436in}{1.174204in}}%
\pgfpathlineto{\pgfqpoint{0.818710in}{1.181241in}}%
\pgfpathlineto{\pgfqpoint{0.802950in}{1.197910in}}%
\pgfpathlineto{\pgfqpoint{0.787202in}{1.214799in}}%
\pgfpathlineto{\pgfqpoint{0.778890in}{1.207952in}}%
\pgfpathlineto{\pgfqpoint{0.770549in}{1.201413in}}%
\pgfpathclose%
\pgfusepath{fill}%
\end{pgfscope}%
\begin{pgfscope}%
\pgfpathrectangle{\pgfqpoint{0.407592in}{0.222904in}}{\pgfqpoint{1.411163in}{1.411163in}}%
\pgfusepath{clip}%
\pgfsetbuttcap%
\pgfsetroundjoin%
\definecolor{currentfill}{rgb}{0.909804,0.698039,0.423529}%
\pgfsetfillcolor{currentfill}%
\pgfsetfillopacity{0.120000}%
\pgfsetlinewidth{0.000000pt}%
\definecolor{currentstroke}{rgb}{0.000000,0.000000,0.000000}%
\pgfsetstrokecolor{currentstroke}%
\pgfsetdash{}{0pt}%
\pgfpathmoveto{\pgfqpoint{1.341720in}{1.149504in}}%
\pgfpathlineto{\pgfqpoint{1.357942in}{1.148313in}}%
\pgfpathlineto{\pgfqpoint{1.374215in}{1.147760in}}%
\pgfpathlineto{\pgfqpoint{1.381820in}{1.161912in}}%
\pgfpathlineto{\pgfqpoint{1.365554in}{1.162512in}}%
\pgfpathlineto{\pgfqpoint{1.349339in}{1.163737in}}%
\pgfusepath{fill}%
\end{pgfscope}%
\begin{pgfscope}%
\pgfpathrectangle{\pgfqpoint{0.407592in}{0.222904in}}{\pgfqpoint{1.411163in}{1.411163in}}%
\pgfusepath{clip}%
\pgfsetbuttcap%
\pgfsetroundjoin%
\definecolor{currentfill}{rgb}{0.909804,0.698039,0.423529}%
\pgfsetfillcolor{currentfill}%
\pgfsetfillopacity{0.120000}%
\pgfsetlinewidth{0.000000pt}%
\definecolor{currentstroke}{rgb}{0.000000,0.000000,0.000000}%
\pgfsetstrokecolor{currentstroke}%
\pgfsetdash{}{0pt}%
\pgfpathmoveto{\pgfqpoint{1.183598in}{1.063106in}}%
\pgfpathlineto{\pgfqpoint{1.199607in}{1.057542in}}%
\pgfpathlineto{\pgfqpoint{1.215655in}{1.052757in}}%
\pgfpathlineto{\pgfqpoint{1.223415in}{1.065568in}}%
\pgfpathlineto{\pgfqpoint{1.231159in}{1.078608in}}%
\pgfpathlineto{\pgfqpoint{1.215127in}{1.083308in}}%
\pgfpathlineto{\pgfqpoint{1.199134in}{1.088752in}}%
\pgfpathlineto{\pgfqpoint{1.191374in}{1.075814in}}%
\pgfpathlineto{\pgfqpoint{1.183598in}{1.063106in}}%
\pgfpathclose%
\pgfusepath{fill}%
\end{pgfscope}%
\begin{pgfscope}%
\pgfpathrectangle{\pgfqpoint{0.407592in}{0.222904in}}{\pgfqpoint{1.411163in}{1.411163in}}%
\pgfusepath{clip}%
\pgfsetbuttcap%
\pgfsetroundjoin%
\definecolor{currentfill}{rgb}{0.909804,0.698039,0.423529}%
\pgfsetfillcolor{currentfill}%
\pgfsetfillopacity{0.120000}%
\pgfsetlinewidth{0.000000pt}%
\definecolor{currentstroke}{rgb}{0.000000,0.000000,0.000000}%
\pgfsetstrokecolor{currentstroke}%
\pgfsetdash{}{0pt}%
\pgfpathmoveto{\pgfqpoint{1.136050in}{1.052193in}}%
\pgfpathlineto{\pgfqpoint{1.152007in}{1.044934in}}%
\pgfpathlineto{\pgfqpoint{1.167997in}{1.038452in}}%
\pgfpathlineto{\pgfqpoint{1.175806in}{1.050646in}}%
\pgfpathlineto{\pgfqpoint{1.183598in}{1.063106in}}%
\pgfpathlineto{\pgfqpoint{1.167625in}{1.069434in}}%
\pgfpathlineto{\pgfqpoint{1.151687in}{1.076507in}}%
\pgfpathlineto{\pgfqpoint{1.143877in}{1.064217in}}%
\pgfpathlineto{\pgfqpoint{1.136050in}{1.052193in}}%
\pgfpathclose%
\pgfusepath{fill}%
\end{pgfscope}%
\begin{pgfscope}%
\pgfpathrectangle{\pgfqpoint{0.407592in}{0.222904in}}{\pgfqpoint{1.411163in}{1.411163in}}%
\pgfusepath{clip}%
\pgfsetbuttcap%
\pgfsetroundjoin%
\definecolor{currentfill}{rgb}{0.909804,0.698039,0.423529}%
\pgfsetfillcolor{currentfill}%
\pgfsetfillopacity{0.120000}%
\pgfsetlinewidth{0.000000pt}%
\definecolor{currentstroke}{rgb}{0.000000,0.000000,0.000000}%
\pgfsetstrokecolor{currentstroke}%
\pgfsetdash{}{0pt}%
\pgfpathmoveto{\pgfqpoint{0.961396in}{1.075593in}}%
\pgfpathlineto{\pgfqpoint{0.977238in}{1.062334in}}%
\pgfpathlineto{\pgfqpoint{0.993099in}{1.049603in}}%
\pgfpathlineto{\pgfqpoint{1.001121in}{1.058891in}}%
\pgfpathlineto{\pgfqpoint{1.009121in}{1.068543in}}%
\pgfpathlineto{\pgfqpoint{0.993291in}{1.080998in}}%
\pgfpathlineto{\pgfqpoint{0.977479in}{1.093978in}}%
\pgfpathlineto{\pgfqpoint{0.969449in}{1.084610in}}%
\pgfpathlineto{\pgfqpoint{0.961396in}{1.075593in}}%
\pgfpathclose%
\pgfusepath{fill}%
\end{pgfscope}%
\begin{pgfscope}%
\pgfpathrectangle{\pgfqpoint{0.407592in}{0.222904in}}{\pgfqpoint{1.411163in}{1.411163in}}%
\pgfusepath{clip}%
\pgfsetbuttcap%
\pgfsetroundjoin%
\definecolor{currentfill}{rgb}{0.909804,0.698039,0.423529}%
\pgfsetfillcolor{currentfill}%
\pgfsetfillopacity{0.120000}%
\pgfsetlinewidth{0.000000pt}%
\definecolor{currentstroke}{rgb}{0.000000,0.000000,0.000000}%
\pgfsetstrokecolor{currentstroke}%
\pgfsetdash{}{0pt}%
\pgfpathmoveto{\pgfqpoint{0.881860in}{1.117187in}}%
\pgfpathlineto{\pgfqpoint{0.897678in}{1.101959in}}%
\pgfpathlineto{\pgfqpoint{0.913509in}{1.087106in}}%
\pgfpathlineto{\pgfqpoint{0.921647in}{1.095140in}}%
\pgfpathlineto{\pgfqpoint{0.929758in}{1.103535in}}%
\pgfpathlineto{\pgfqpoint{0.913961in}{1.118143in}}%
\pgfpathlineto{\pgfqpoint{0.898178in}{1.133129in}}%
\pgfpathlineto{\pgfqpoint{0.890032in}{1.124986in}}%
\pgfpathlineto{\pgfqpoint{0.881860in}{1.117187in}}%
\pgfpathclose%
\pgfusepath{fill}%
\end{pgfscope}%
\begin{pgfscope}%
\pgfpathrectangle{\pgfqpoint{0.407592in}{0.222904in}}{\pgfqpoint{1.411163in}{1.411163in}}%
\pgfusepath{clip}%
\pgfsetbuttcap%
\pgfsetroundjoin%
\definecolor{currentfill}{rgb}{0.909804,0.698039,0.423529}%
\pgfsetfillcolor{currentfill}%
\pgfsetfillopacity{0.120000}%
\pgfsetlinewidth{0.000000pt}%
\definecolor{currentstroke}{rgb}{0.000000,0.000000,0.000000}%
\pgfsetstrokecolor{currentstroke}%
\pgfsetdash{}{0pt}%
\pgfpathmoveto{\pgfqpoint{1.231159in}{1.078608in}}%
\pgfpathlineto{\pgfqpoint{1.247233in}{1.074658in}}%
\pgfpathlineto{\pgfqpoint{1.263350in}{1.071463in}}%
\pgfpathlineto{\pgfqpoint{1.271065in}{1.084747in}}%
\pgfpathlineto{\pgfqpoint{1.278765in}{1.098228in}}%
\pgfpathlineto{\pgfqpoint{1.262662in}{1.101408in}}%
\pgfpathlineto{\pgfqpoint{1.246603in}{1.105308in}}%
\pgfpathlineto{\pgfqpoint{1.238888in}{1.091860in}}%
\pgfpathlineto{\pgfqpoint{1.231159in}{1.078608in}}%
\pgfpathclose%
\pgfusepath{fill}%
\end{pgfscope}%
\begin{pgfscope}%
\pgfpathrectangle{\pgfqpoint{0.407592in}{0.222904in}}{\pgfqpoint{1.411163in}{1.411163in}}%
\pgfusepath{clip}%
\pgfsetbuttcap%
\pgfsetroundjoin%
\definecolor{currentfill}{rgb}{0.909804,0.698039,0.423529}%
\pgfsetfillcolor{currentfill}%
\pgfsetfillopacity{0.120000}%
\pgfsetlinewidth{0.000000pt}%
\definecolor{currentstroke}{rgb}{0.000000,0.000000,0.000000}%
\pgfsetstrokecolor{currentstroke}%
\pgfsetdash{}{0pt}%
\pgfpathmoveto{\pgfqpoint{1.088477in}{1.046225in}}%
\pgfpathlineto{\pgfqpoint{1.104395in}{1.037252in}}%
\pgfpathlineto{\pgfqpoint{1.120342in}{1.029027in}}%
\pgfpathlineto{\pgfqpoint{1.128205in}{1.040457in}}%
\pgfpathlineto{\pgfqpoint{1.136050in}{1.052193in}}%
\pgfpathlineto{\pgfqpoint{1.120124in}{1.060205in}}%
\pgfpathlineto{\pgfqpoint{1.104228in}{1.068940in}}%
\pgfpathlineto{\pgfqpoint{1.096362in}{1.057430in}}%
\pgfpathlineto{\pgfqpoint{1.088477in}{1.046225in}}%
\pgfpathclose%
\pgfusepath{fill}%
\end{pgfscope}%
\begin{pgfscope}%
\pgfpathrectangle{\pgfqpoint{0.407592in}{0.222904in}}{\pgfqpoint{1.411163in}{1.411163in}}%
\pgfusepath{clip}%
\pgfsetbuttcap%
\pgfsetroundjoin%
\definecolor{currentfill}{rgb}{0.909804,0.698039,0.423529}%
\pgfsetfillcolor{currentfill}%
\pgfsetfillopacity{0.120000}%
\pgfsetlinewidth{0.000000pt}%
\definecolor{currentstroke}{rgb}{0.000000,0.000000,0.000000}%
\pgfsetstrokecolor{currentstroke}%
\pgfsetdash{}{0pt}%
\pgfpathmoveto{\pgfqpoint{1.278765in}{1.098228in}}%
\pgfpathlineto{\pgfqpoint{1.294915in}{1.095767in}}%
\pgfpathlineto{\pgfqpoint{1.311112in}{1.094022in}}%
\pgfpathlineto{\pgfqpoint{1.318785in}{1.107651in}}%
\pgfpathlineto{\pgfqpoint{1.326443in}{1.121449in}}%
\pgfpathlineto{\pgfqpoint{1.310259in}{1.123243in}}%
\pgfpathlineto{\pgfqpoint{1.294124in}{1.125720in}}%
\pgfpathlineto{\pgfqpoint{1.286451in}{1.111890in}}%
\pgfpathlineto{\pgfqpoint{1.278765in}{1.098228in}}%
\pgfpathclose%
\pgfusepath{fill}%
\end{pgfscope}%
\begin{pgfscope}%
\pgfpathrectangle{\pgfqpoint{0.407592in}{0.222904in}}{\pgfqpoint{1.411163in}{1.411163in}}%
\pgfusepath{clip}%
\pgfsetbuttcap%
\pgfsetroundjoin%
\definecolor{currentfill}{rgb}{0.909804,0.698039,0.423529}%
\pgfsetfillcolor{currentfill}%
\pgfsetfillopacity{0.120000}%
\pgfsetlinewidth{0.000000pt}%
\definecolor{currentstroke}{rgb}{0.000000,0.000000,0.000000}%
\pgfsetstrokecolor{currentstroke}%
\pgfsetdash{}{0pt}%
\pgfpathmoveto{\pgfqpoint{0.802133in}{1.167492in}}%
\pgfpathlineto{\pgfqpoint{0.817941in}{1.150873in}}%
\pgfpathlineto{\pgfqpoint{0.833760in}{1.134511in}}%
\pgfpathlineto{\pgfqpoint{0.842025in}{1.141404in}}%
\pgfpathlineto{\pgfqpoint{0.850261in}{1.148640in}}%
\pgfpathlineto{\pgfqpoint{0.834480in}{1.164810in}}%
\pgfpathlineto{\pgfqpoint{0.818710in}{1.181241in}}%
\pgfpathlineto{\pgfqpoint{0.810436in}{1.174204in}}%
\pgfpathlineto{\pgfqpoint{0.802133in}{1.167492in}}%
\pgfpathclose%
\pgfusepath{fill}%
\end{pgfscope}%
\begin{pgfscope}%
\pgfpathrectangle{\pgfqpoint{0.407592in}{0.222904in}}{\pgfqpoint{1.411163in}{1.411163in}}%
\pgfusepath{clip}%
\pgfsetbuttcap%
\pgfsetroundjoin%
\definecolor{currentfill}{rgb}{0.909804,0.698039,0.423529}%
\pgfsetfillcolor{currentfill}%
\pgfsetfillopacity{0.120000}%
\pgfsetlinewidth{0.000000pt}%
\definecolor{currentstroke}{rgb}{0.000000,0.000000,0.000000}%
\pgfsetstrokecolor{currentstroke}%
\pgfsetdash{}{0pt}%
\pgfpathmoveto{\pgfqpoint{1.040840in}{1.045374in}}%
\pgfpathlineto{\pgfqpoint{1.056732in}{1.034745in}}%
\pgfpathlineto{\pgfqpoint{1.072648in}{1.024807in}}%
\pgfpathlineto{\pgfqpoint{1.080573in}{1.035344in}}%
\pgfpathlineto{\pgfqpoint{1.088477in}{1.046225in}}%
\pgfpathlineto{\pgfqpoint{1.072586in}{1.055911in}}%
\pgfpathlineto{\pgfqpoint{1.056719in}{1.066273in}}%
\pgfpathlineto{\pgfqpoint{1.048790in}{1.055655in}}%
\pgfpathlineto{\pgfqpoint{1.040840in}{1.045374in}}%
\pgfpathclose%
\pgfusepath{fill}%
\end{pgfscope}%
\begin{pgfscope}%
\pgfpathrectangle{\pgfqpoint{0.407592in}{0.222904in}}{\pgfqpoint{1.411163in}{1.411163in}}%
\pgfusepath{clip}%
\pgfsetbuttcap%
\pgfsetroundjoin%
\definecolor{currentfill}{rgb}{0.909804,0.698039,0.423529}%
\pgfsetfillcolor{currentfill}%
\pgfsetfillopacity{0.120000}%
\pgfsetlinewidth{0.000000pt}%
\definecolor{currentstroke}{rgb}{0.000000,0.000000,0.000000}%
\pgfsetstrokecolor{currentstroke}%
\pgfsetdash{}{0pt}%
\pgfpathmoveto{\pgfqpoint{0.690568in}{1.259886in}}%
\pgfpathlineto{\pgfqpoint{0.706355in}{1.242001in}}%
\pgfpathlineto{\pgfqpoint{0.722152in}{1.224261in}}%
\pgfpathlineto{\pgfqpoint{0.730594in}{1.230039in}}%
\pgfpathlineto{\pgfqpoint{0.739005in}{1.236117in}}%
\pgfpathlineto{\pgfqpoint{0.723249in}{1.253721in}}%
\pgfpathlineto{\pgfqpoint{0.707502in}{1.271476in}}%
\pgfpathlineto{\pgfqpoint{0.699051in}{1.265539in}}%
\pgfpathlineto{\pgfqpoint{0.690568in}{1.259886in}}%
\pgfpathclose%
\pgfusepath{fill}%
\end{pgfscope}%
\begin{pgfscope}%
\pgfpathrectangle{\pgfqpoint{0.407592in}{0.222904in}}{\pgfqpoint{1.411163in}{1.411163in}}%
\pgfusepath{clip}%
\pgfsetbuttcap%
\pgfsetroundjoin%
\definecolor{currentfill}{rgb}{0.909804,0.698039,0.423529}%
\pgfsetfillcolor{currentfill}%
\pgfsetfillopacity{0.120000}%
\pgfsetlinewidth{0.000000pt}%
\definecolor{currentstroke}{rgb}{0.000000,0.000000,0.000000}%
\pgfsetstrokecolor{currentstroke}%
\pgfsetdash{}{0pt}%
\pgfpathmoveto{\pgfqpoint{1.326443in}{1.121449in}}%
\pgfpathlineto{\pgfqpoint{1.342677in}{1.120331in}}%
\pgfpathlineto{\pgfqpoint{1.358964in}{1.119878in}}%
\pgfpathlineto{\pgfqpoint{1.366596in}{1.133745in}}%
\pgfpathlineto{\pgfqpoint{1.374215in}{1.147760in}}%
\pgfpathlineto{\pgfqpoint{1.357942in}{1.148313in}}%
\pgfpathlineto{\pgfqpoint{1.341720in}{1.149504in}}%
\pgfpathlineto{\pgfqpoint{1.334089in}{1.135404in}}%
\pgfpathlineto{\pgfqpoint{1.326443in}{1.121449in}}%
\pgfpathclose%
\pgfusepath{fill}%
\end{pgfscope}%
\begin{pgfscope}%
\pgfpathrectangle{\pgfqpoint{0.407592in}{0.222904in}}{\pgfqpoint{1.411163in}{1.411163in}}%
\pgfusepath{clip}%
\pgfsetbuttcap%
\pgfsetroundjoin%
\definecolor{currentfill}{rgb}{0.909804,0.698039,0.423529}%
\pgfsetfillcolor{currentfill}%
\pgfsetfillopacity{0.120000}%
\pgfsetlinewidth{0.000000pt}%
\definecolor{currentstroke}{rgb}{0.000000,0.000000,0.000000}%
\pgfsetstrokecolor{currentstroke}%
\pgfsetdash{}{0pt}%
\pgfpathmoveto{\pgfqpoint{0.913509in}{1.087106in}}%
\pgfpathlineto{\pgfqpoint{0.929355in}{1.072663in}}%
\pgfpathlineto{\pgfqpoint{0.945215in}{1.058668in}}%
\pgfpathlineto{\pgfqpoint{0.953318in}{1.066941in}}%
\pgfpathlineto{\pgfqpoint{0.961396in}{1.075593in}}%
\pgfpathlineto{\pgfqpoint{0.945569in}{1.089340in}}%
\pgfpathlineto{\pgfqpoint{0.929758in}{1.103535in}}%
\pgfpathlineto{\pgfqpoint{0.921647in}{1.095140in}}%
\pgfpathlineto{\pgfqpoint{0.913509in}{1.087106in}}%
\pgfpathclose%
\pgfusepath{fill}%
\end{pgfscope}%
\begin{pgfscope}%
\pgfpathrectangle{\pgfqpoint{0.407592in}{0.222904in}}{\pgfqpoint{1.411163in}{1.411163in}}%
\pgfusepath{clip}%
\pgfsetbuttcap%
\pgfsetroundjoin%
\definecolor{currentfill}{rgb}{0.909804,0.698039,0.423529}%
\pgfsetfillcolor{currentfill}%
\pgfsetfillopacity{0.120000}%
\pgfsetlinewidth{0.000000pt}%
\definecolor{currentstroke}{rgb}{0.000000,0.000000,0.000000}%
\pgfsetstrokecolor{currentstroke}%
\pgfsetdash{}{0pt}%
\pgfpathmoveto{\pgfqpoint{0.993099in}{1.049603in}}%
\pgfpathlineto{\pgfqpoint{1.008977in}{1.037445in}}%
\pgfpathlineto{\pgfqpoint{1.024876in}{1.025903in}}%
\pgfpathlineto{\pgfqpoint{1.032869in}{1.035451in}}%
\pgfpathlineto{\pgfqpoint{1.040840in}{1.045374in}}%
\pgfpathlineto{\pgfqpoint{1.024970in}{1.056655in}}%
\pgfpathlineto{\pgfqpoint{1.009121in}{1.068543in}}%
\pgfpathlineto{\pgfqpoint{1.001121in}{1.058891in}}%
\pgfpathlineto{\pgfqpoint{0.993099in}{1.049603in}}%
\pgfpathclose%
\pgfusepath{fill}%
\end{pgfscope}%
\begin{pgfscope}%
\pgfpathrectangle{\pgfqpoint{0.407592in}{0.222904in}}{\pgfqpoint{1.411163in}{1.411163in}}%
\pgfusepath{clip}%
\pgfsetbuttcap%
\pgfsetroundjoin%
\definecolor{currentfill}{rgb}{0.909804,0.698039,0.423529}%
\pgfsetfillcolor{currentfill}%
\pgfsetfillopacity{0.120000}%
\pgfsetlinewidth{0.000000pt}%
\definecolor{currentstroke}{rgb}{0.000000,0.000000,0.000000}%
\pgfsetstrokecolor{currentstroke}%
\pgfsetdash{}{0pt}%
\pgfpathmoveto{\pgfqpoint{0.722152in}{1.224261in}}%
\pgfpathlineto{\pgfqpoint{0.737958in}{1.206682in}}%
\pgfpathlineto{\pgfqpoint{0.753774in}{1.189279in}}%
\pgfpathlineto{\pgfqpoint{0.762177in}{1.195187in}}%
\pgfpathlineto{\pgfqpoint{0.770549in}{1.201413in}}%
\pgfpathlineto{\pgfqpoint{0.754772in}{1.218675in}}%
\pgfpathlineto{\pgfqpoint{0.739005in}{1.236117in}}%
\pgfpathlineto{\pgfqpoint{0.730594in}{1.230039in}}%
\pgfpathlineto{\pgfqpoint{0.722152in}{1.224261in}}%
\pgfpathclose%
\pgfusepath{fill}%
\end{pgfscope}%
\begin{pgfscope}%
\pgfpathrectangle{\pgfqpoint{0.407592in}{0.222904in}}{\pgfqpoint{1.411163in}{1.411163in}}%
\pgfusepath{clip}%
\pgfsetbuttcap%
\pgfsetroundjoin%
\definecolor{currentfill}{rgb}{0.909804,0.698039,0.423529}%
\pgfsetfillcolor{currentfill}%
\pgfsetfillopacity{0.120000}%
\pgfsetlinewidth{0.000000pt}%
\definecolor{currentstroke}{rgb}{0.000000,0.000000,0.000000}%
\pgfsetstrokecolor{currentstroke}%
\pgfsetdash{}{0pt}%
\pgfpathmoveto{\pgfqpoint{0.833760in}{1.134511in}}%
\pgfpathlineto{\pgfqpoint{0.849591in}{1.118431in}}%
\pgfpathlineto{\pgfqpoint{0.865433in}{1.102660in}}%
\pgfpathlineto{\pgfqpoint{0.873660in}{1.109742in}}%
\pgfpathlineto{\pgfqpoint{0.881860in}{1.117187in}}%
\pgfpathlineto{\pgfqpoint{0.866055in}{1.132757in}}%
\pgfpathlineto{\pgfqpoint{0.850261in}{1.148640in}}%
\pgfpathlineto{\pgfqpoint{0.842025in}{1.141404in}}%
\pgfpathlineto{\pgfqpoint{0.833760in}{1.134511in}}%
\pgfpathclose%
\pgfusepath{fill}%
\end{pgfscope}%
\begin{pgfscope}%
\pgfpathrectangle{\pgfqpoint{0.407592in}{0.222904in}}{\pgfqpoint{1.411163in}{1.411163in}}%
\pgfusepath{clip}%
\pgfsetbuttcap%
\pgfsetroundjoin%
\definecolor{currentfill}{rgb}{0.909804,0.698039,0.423529}%
\pgfsetfillcolor{currentfill}%
\pgfsetfillopacity{0.120000}%
\pgfsetlinewidth{0.000000pt}%
\definecolor{currentstroke}{rgb}{0.000000,0.000000,0.000000}%
\pgfsetstrokecolor{currentstroke}%
\pgfsetdash{}{0pt}%
\pgfpathmoveto{\pgfqpoint{1.167997in}{1.038452in}}%
\pgfpathlineto{\pgfqpoint{1.184024in}{1.032768in}}%
\pgfpathlineto{\pgfqpoint{1.200090in}{1.027897in}}%
\pgfpathlineto{\pgfqpoint{1.207880in}{1.040194in}}%
\pgfpathlineto{\pgfqpoint{1.215655in}{1.052757in}}%
\pgfpathlineto{\pgfqpoint{1.199607in}{1.057542in}}%
\pgfpathlineto{\pgfqpoint{1.183598in}{1.063106in}}%
\pgfpathlineto{\pgfqpoint{1.175806in}{1.050646in}}%
\pgfpathlineto{\pgfqpoint{1.167997in}{1.038452in}}%
\pgfpathclose%
\pgfusepath{fill}%
\end{pgfscope}%
\begin{pgfscope}%
\pgfpathrectangle{\pgfqpoint{0.407592in}{0.222904in}}{\pgfqpoint{1.411163in}{1.411163in}}%
\pgfusepath{clip}%
\pgfsetbuttcap%
\pgfsetroundjoin%
\definecolor{currentfill}{rgb}{0.909804,0.698039,0.423529}%
\pgfsetfillcolor{currentfill}%
\pgfsetfillopacity{0.120000}%
\pgfsetlinewidth{0.000000pt}%
\definecolor{currentstroke}{rgb}{0.000000,0.000000,0.000000}%
\pgfsetstrokecolor{currentstroke}%
\pgfsetdash{}{0pt}%
\pgfpathmoveto{\pgfqpoint{1.120342in}{1.029027in}}%
\pgfpathlineto{\pgfqpoint{1.136319in}{1.021583in}}%
\pgfpathlineto{\pgfqpoint{1.152330in}{1.014947in}}%
\pgfpathlineto{\pgfqpoint{1.160172in}{1.026546in}}%
\pgfpathlineto{\pgfqpoint{1.167997in}{1.038452in}}%
\pgfpathlineto{\pgfqpoint{1.152007in}{1.044934in}}%
\pgfpathlineto{\pgfqpoint{1.136050in}{1.052193in}}%
\pgfpathlineto{\pgfqpoint{1.128205in}{1.040457in}}%
\pgfpathlineto{\pgfqpoint{1.120342in}{1.029027in}}%
\pgfpathclose%
\pgfusepath{fill}%
\end{pgfscope}%
\begin{pgfscope}%
\pgfpathrectangle{\pgfqpoint{0.407592in}{0.222904in}}{\pgfqpoint{1.411163in}{1.411163in}}%
\pgfusepath{clip}%
\pgfsetbuttcap%
\pgfsetroundjoin%
\definecolor{currentfill}{rgb}{0.909804,0.698039,0.423529}%
\pgfsetfillcolor{currentfill}%
\pgfsetfillopacity{0.120000}%
\pgfsetlinewidth{0.000000pt}%
\definecolor{currentstroke}{rgb}{0.000000,0.000000,0.000000}%
\pgfsetstrokecolor{currentstroke}%
\pgfsetdash{}{0pt}%
\pgfpathmoveto{\pgfqpoint{1.215655in}{1.052757in}}%
\pgfpathlineto{\pgfqpoint{1.231745in}{1.048758in}}%
\pgfpathlineto{\pgfqpoint{1.247878in}{1.045550in}}%
\pgfpathlineto{\pgfqpoint{1.255621in}{1.058392in}}%
\pgfpathlineto{\pgfqpoint{1.263350in}{1.071463in}}%
\pgfpathlineto{\pgfqpoint{1.247233in}{1.074658in}}%
\pgfpathlineto{\pgfqpoint{1.231159in}{1.078608in}}%
\pgfpathlineto{\pgfqpoint{1.223415in}{1.065568in}}%
\pgfpathlineto{\pgfqpoint{1.215655in}{1.052757in}}%
\pgfpathclose%
\pgfusepath{fill}%
\end{pgfscope}%
\begin{pgfscope}%
\pgfpathrectangle{\pgfqpoint{0.407592in}{0.222904in}}{\pgfqpoint{1.411163in}{1.411163in}}%
\pgfusepath{clip}%
\pgfsetbuttcap%
\pgfsetroundjoin%
\definecolor{currentfill}{rgb}{0.909804,0.698039,0.423529}%
\pgfsetfillcolor{currentfill}%
\pgfsetfillopacity{0.120000}%
\pgfsetlinewidth{0.000000pt}%
\definecolor{currentstroke}{rgb}{0.000000,0.000000,0.000000}%
\pgfsetstrokecolor{currentstroke}%
\pgfsetdash{}{0pt}%
\pgfpathmoveto{\pgfqpoint{1.072648in}{1.024807in}}%
\pgfpathlineto{\pgfqpoint{1.088590in}{1.015603in}}%
\pgfpathlineto{\pgfqpoint{1.104560in}{1.007171in}}%
\pgfpathlineto{\pgfqpoint{1.112460in}{1.017925in}}%
\pgfpathlineto{\pgfqpoint{1.120342in}{1.029027in}}%
\pgfpathlineto{\pgfqpoint{1.104395in}{1.037252in}}%
\pgfpathlineto{\pgfqpoint{1.088477in}{1.046225in}}%
\pgfpathlineto{\pgfqpoint{1.080573in}{1.035344in}}%
\pgfpathlineto{\pgfqpoint{1.072648in}{1.024807in}}%
\pgfpathclose%
\pgfusepath{fill}%
\end{pgfscope}%
\begin{pgfscope}%
\pgfpathrectangle{\pgfqpoint{0.407592in}{0.222904in}}{\pgfqpoint{1.411163in}{1.411163in}}%
\pgfusepath{clip}%
\pgfsetbuttcap%
\pgfsetroundjoin%
\definecolor{currentfill}{rgb}{0.909804,0.698039,0.423529}%
\pgfsetfillcolor{currentfill}%
\pgfsetfillopacity{0.120000}%
\pgfsetlinewidth{0.000000pt}%
\definecolor{currentstroke}{rgb}{0.000000,0.000000,0.000000}%
\pgfsetstrokecolor{currentstroke}%
\pgfsetdash{}{0pt}%
\pgfpathmoveto{\pgfqpoint{0.753774in}{1.189279in}}%
\pgfpathlineto{\pgfqpoint{0.769601in}{1.172068in}}%
\pgfpathlineto{\pgfqpoint{0.785437in}{1.155067in}}%
\pgfpathlineto{\pgfqpoint{0.793800in}{1.161112in}}%
\pgfpathlineto{\pgfqpoint{0.802133in}{1.167492in}}%
\pgfpathlineto{\pgfqpoint{0.786336in}{1.184345in}}%
\pgfpathlineto{\pgfqpoint{0.770549in}{1.201413in}}%
\pgfpathlineto{\pgfqpoint{0.762177in}{1.195187in}}%
\pgfpathlineto{\pgfqpoint{0.753774in}{1.189279in}}%
\pgfpathclose%
\pgfusepath{fill}%
\end{pgfscope}%
\begin{pgfscope}%
\pgfpathrectangle{\pgfqpoint{0.407592in}{0.222904in}}{\pgfqpoint{1.411163in}{1.411163in}}%
\pgfusepath{clip}%
\pgfsetbuttcap%
\pgfsetroundjoin%
\definecolor{currentfill}{rgb}{0.909804,0.698039,0.423529}%
\pgfsetfillcolor{currentfill}%
\pgfsetfillopacity{0.120000}%
\pgfsetlinewidth{0.000000pt}%
\definecolor{currentstroke}{rgb}{0.000000,0.000000,0.000000}%
\pgfsetstrokecolor{currentstroke}%
\pgfsetdash{}{0pt}%
\pgfpathmoveto{\pgfqpoint{1.374215in}{1.147760in}}%
\pgfpathlineto{\pgfqpoint{1.390541in}{1.147834in}}%
\pgfpathlineto{\pgfqpoint{1.406923in}{1.148521in}}%
\pgfpathlineto{\pgfqpoint{1.414516in}{1.162542in}}%
\pgfpathlineto{\pgfqpoint{1.398141in}{1.161926in}}%
\pgfpathlineto{\pgfqpoint{1.381820in}{1.161912in}}%
\pgfusepath{fill}%
\end{pgfscope}%
\begin{pgfscope}%
\pgfpathrectangle{\pgfqpoint{0.407592in}{0.222904in}}{\pgfqpoint{1.411163in}{1.411163in}}%
\pgfusepath{clip}%
\pgfsetbuttcap%
\pgfsetroundjoin%
\definecolor{currentfill}{rgb}{0.909804,0.698039,0.423529}%
\pgfsetfillcolor{currentfill}%
\pgfsetfillopacity{0.120000}%
\pgfsetlinewidth{0.000000pt}%
\definecolor{currentstroke}{rgb}{0.000000,0.000000,0.000000}%
\pgfsetstrokecolor{currentstroke}%
\pgfsetdash{}{0pt}%
\pgfpathmoveto{\pgfqpoint{0.945215in}{1.058668in}}%
\pgfpathlineto{\pgfqpoint{0.961090in}{1.045160in}}%
\pgfpathlineto{\pgfqpoint{0.976982in}{1.032183in}}%
\pgfpathlineto{\pgfqpoint{0.985052in}{1.040695in}}%
\pgfpathlineto{\pgfqpoint{0.993099in}{1.049603in}}%
\pgfpathlineto{\pgfqpoint{0.977238in}{1.062334in}}%
\pgfpathlineto{\pgfqpoint{0.961396in}{1.075593in}}%
\pgfpathlineto{\pgfqpoint{0.953318in}{1.066941in}}%
\pgfpathlineto{\pgfqpoint{0.945215in}{1.058668in}}%
\pgfpathclose%
\pgfusepath{fill}%
\end{pgfscope}%
\begin{pgfscope}%
\pgfpathrectangle{\pgfqpoint{0.407592in}{0.222904in}}{\pgfqpoint{1.411163in}{1.411163in}}%
\pgfusepath{clip}%
\pgfsetbuttcap%
\pgfsetroundjoin%
\definecolor{currentfill}{rgb}{0.909804,0.698039,0.423529}%
\pgfsetfillcolor{currentfill}%
\pgfsetfillopacity{0.120000}%
\pgfsetlinewidth{0.000000pt}%
\definecolor{currentstroke}{rgb}{0.000000,0.000000,0.000000}%
\pgfsetstrokecolor{currentstroke}%
\pgfsetdash{}{0pt}%
\pgfpathmoveto{\pgfqpoint{1.263350in}{1.071463in}}%
\pgfpathlineto{\pgfqpoint{1.279514in}{1.069022in}}%
\pgfpathlineto{\pgfqpoint{1.295726in}{1.067329in}}%
\pgfpathlineto{\pgfqpoint{1.303426in}{1.080576in}}%
\pgfpathlineto{\pgfqpoint{1.311112in}{1.094022in}}%
\pgfpathlineto{\pgfqpoint{1.294915in}{1.095767in}}%
\pgfpathlineto{\pgfqpoint{1.278765in}{1.098228in}}%
\pgfpathlineto{\pgfqpoint{1.271065in}{1.084747in}}%
\pgfpathlineto{\pgfqpoint{1.263350in}{1.071463in}}%
\pgfpathclose%
\pgfusepath{fill}%
\end{pgfscope}%
\begin{pgfscope}%
\pgfpathrectangle{\pgfqpoint{0.407592in}{0.222904in}}{\pgfqpoint{1.411163in}{1.411163in}}%
\pgfusepath{clip}%
\pgfsetbuttcap%
\pgfsetroundjoin%
\definecolor{currentfill}{rgb}{0.909804,0.698039,0.423529}%
\pgfsetfillcolor{currentfill}%
\pgfsetfillopacity{0.120000}%
\pgfsetlinewidth{0.000000pt}%
\definecolor{currentstroke}{rgb}{0.000000,0.000000,0.000000}%
\pgfsetstrokecolor{currentstroke}%
\pgfsetdash{}{0pt}%
\pgfpathmoveto{\pgfqpoint{0.865433in}{1.102660in}}%
\pgfpathlineto{\pgfqpoint{0.881287in}{1.087228in}}%
\pgfpathlineto{\pgfqpoint{0.897155in}{1.072168in}}%
\pgfpathlineto{\pgfqpoint{0.905346in}{1.079446in}}%
\pgfpathlineto{\pgfqpoint{0.913509in}{1.087106in}}%
\pgfpathlineto{\pgfqpoint{0.897678in}{1.101959in}}%
\pgfpathlineto{\pgfqpoint{0.881860in}{1.117187in}}%
\pgfpathlineto{\pgfqpoint{0.873660in}{1.109742in}}%
\pgfpathlineto{\pgfqpoint{0.865433in}{1.102660in}}%
\pgfpathclose%
\pgfusepath{fill}%
\end{pgfscope}%
\begin{pgfscope}%
\pgfpathrectangle{\pgfqpoint{0.407592in}{0.222904in}}{\pgfqpoint{1.411163in}{1.411163in}}%
\pgfusepath{clip}%
\pgfsetbuttcap%
\pgfsetroundjoin%
\definecolor{currentfill}{rgb}{0.909804,0.698039,0.423529}%
\pgfsetfillcolor{currentfill}%
\pgfsetfillopacity{0.120000}%
\pgfsetlinewidth{0.000000pt}%
\definecolor{currentstroke}{rgb}{0.000000,0.000000,0.000000}%
\pgfsetstrokecolor{currentstroke}%
\pgfsetdash{}{0pt}%
\pgfpathmoveto{\pgfqpoint{1.024876in}{1.025903in}}%
\pgfpathlineto{\pgfqpoint{1.040795in}{1.015022in}}%
\pgfpathlineto{\pgfqpoint{1.056738in}{1.004849in}}%
\pgfpathlineto{\pgfqpoint{1.064704in}{1.014636in}}%
\pgfpathlineto{\pgfqpoint{1.072648in}{1.024807in}}%
\pgfpathlineto{\pgfqpoint{1.056732in}{1.034745in}}%
\pgfpathlineto{\pgfqpoint{1.040840in}{1.045374in}}%
\pgfpathlineto{\pgfqpoint{1.032869in}{1.035451in}}%
\pgfpathlineto{\pgfqpoint{1.024876in}{1.025903in}}%
\pgfpathclose%
\pgfusepath{fill}%
\end{pgfscope}%
\begin{pgfscope}%
\pgfpathrectangle{\pgfqpoint{0.407592in}{0.222904in}}{\pgfqpoint{1.411163in}{1.411163in}}%
\pgfusepath{clip}%
\pgfsetbuttcap%
\pgfsetroundjoin%
\definecolor{currentfill}{rgb}{0.909804,0.698039,0.423529}%
\pgfsetfillcolor{currentfill}%
\pgfsetfillopacity{0.120000}%
\pgfsetlinewidth{0.000000pt}%
\definecolor{currentstroke}{rgb}{0.000000,0.000000,0.000000}%
\pgfsetstrokecolor{currentstroke}%
\pgfsetdash{}{0pt}%
\pgfpathmoveto{\pgfqpoint{1.311112in}{1.094022in}}%
\pgfpathlineto{\pgfqpoint{1.327360in}{1.092982in}}%
\pgfpathlineto{\pgfqpoint{1.343660in}{1.092638in}}%
\pgfpathlineto{\pgfqpoint{1.351319in}{1.106172in}}%
\pgfpathlineto{\pgfqpoint{1.358964in}{1.119878in}}%
\pgfpathlineto{\pgfqpoint{1.342677in}{1.120331in}}%
\pgfpathlineto{\pgfqpoint{1.326443in}{1.121449in}}%
\pgfpathlineto{\pgfqpoint{1.318785in}{1.107651in}}%
\pgfpathlineto{\pgfqpoint{1.311112in}{1.094022in}}%
\pgfpathclose%
\pgfusepath{fill}%
\end{pgfscope}%
\begin{pgfscope}%
\pgfpathrectangle{\pgfqpoint{0.407592in}{0.222904in}}{\pgfqpoint{1.411163in}{1.411163in}}%
\pgfusepath{clip}%
\pgfsetbuttcap%
\pgfsetroundjoin%
\definecolor{currentfill}{rgb}{0.909804,0.698039,0.423529}%
\pgfsetfillcolor{currentfill}%
\pgfsetfillopacity{0.120000}%
\pgfsetlinewidth{0.000000pt}%
\definecolor{currentstroke}{rgb}{0.000000,0.000000,0.000000}%
\pgfsetstrokecolor{currentstroke}%
\pgfsetdash{}{0pt}%
\pgfpathmoveto{\pgfqpoint{0.785437in}{1.155067in}}%
\pgfpathlineto{\pgfqpoint{0.801284in}{1.138298in}}%
\pgfpathlineto{\pgfqpoint{0.817142in}{1.121783in}}%
\pgfpathlineto{\pgfqpoint{0.825466in}{1.127969in}}%
\pgfpathlineto{\pgfqpoint{0.833760in}{1.134511in}}%
\pgfpathlineto{\pgfqpoint{0.817941in}{1.150873in}}%
\pgfpathlineto{\pgfqpoint{0.802133in}{1.167492in}}%
\pgfpathlineto{\pgfqpoint{0.793800in}{1.161112in}}%
\pgfpathlineto{\pgfqpoint{0.785437in}{1.155067in}}%
\pgfpathclose%
\pgfusepath{fill}%
\end{pgfscope}%
\begin{pgfscope}%
\pgfpathrectangle{\pgfqpoint{0.407592in}{0.222904in}}{\pgfqpoint{1.411163in}{1.411163in}}%
\pgfusepath{clip}%
\pgfsetbuttcap%
\pgfsetroundjoin%
\definecolor{currentfill}{rgb}{0.909804,0.698039,0.423529}%
\pgfsetfillcolor{currentfill}%
\pgfsetfillopacity{0.120000}%
\pgfsetlinewidth{0.000000pt}%
\definecolor{currentstroke}{rgb}{0.000000,0.000000,0.000000}%
\pgfsetstrokecolor{currentstroke}%
\pgfsetdash{}{0pt}%
\pgfpathmoveto{\pgfqpoint{0.673507in}{1.249448in}}%
\pgfpathlineto{\pgfqpoint{0.689334in}{1.231462in}}%
\pgfpathlineto{\pgfqpoint{0.705171in}{1.213622in}}%
\pgfpathlineto{\pgfqpoint{0.713678in}{1.218788in}}%
\pgfpathlineto{\pgfqpoint{0.722152in}{1.224261in}}%
\pgfpathlineto{\pgfqpoint{0.706355in}{1.242001in}}%
\pgfpathlineto{\pgfqpoint{0.690568in}{1.259886in}}%
\pgfpathlineto{\pgfqpoint{0.682054in}{1.254522in}}%
\pgfpathlineto{\pgfqpoint{0.673507in}{1.249448in}}%
\pgfpathclose%
\pgfusepath{fill}%
\end{pgfscope}%
\begin{pgfscope}%
\pgfpathrectangle{\pgfqpoint{0.407592in}{0.222904in}}{\pgfqpoint{1.411163in}{1.411163in}}%
\pgfusepath{clip}%
\pgfsetbuttcap%
\pgfsetroundjoin%
\definecolor{currentfill}{rgb}{0.909804,0.698039,0.423529}%
\pgfsetfillcolor{currentfill}%
\pgfsetfillopacity{0.120000}%
\pgfsetlinewidth{0.000000pt}%
\definecolor{currentstroke}{rgb}{0.000000,0.000000,0.000000}%
\pgfsetstrokecolor{currentstroke}%
\pgfsetdash{}{0pt}%
\pgfpathmoveto{\pgfqpoint{0.897155in}{1.072168in}}%
\pgfpathlineto{\pgfqpoint{0.913035in}{1.057515in}}%
\pgfpathlineto{\pgfqpoint{0.928930in}{1.043308in}}%
\pgfpathlineto{\pgfqpoint{0.937085in}{1.050786in}}%
\pgfpathlineto{\pgfqpoint{0.945215in}{1.058668in}}%
\pgfpathlineto{\pgfqpoint{0.929355in}{1.072663in}}%
\pgfpathlineto{\pgfqpoint{0.913509in}{1.087106in}}%
\pgfpathlineto{\pgfqpoint{0.905346in}{1.079446in}}%
\pgfpathlineto{\pgfqpoint{0.897155in}{1.072168in}}%
\pgfpathclose%
\pgfusepath{fill}%
\end{pgfscope}%
\begin{pgfscope}%
\pgfpathrectangle{\pgfqpoint{0.407592in}{0.222904in}}{\pgfqpoint{1.411163in}{1.411163in}}%
\pgfusepath{clip}%
\pgfsetbuttcap%
\pgfsetroundjoin%
\definecolor{currentfill}{rgb}{0.909804,0.698039,0.423529}%
\pgfsetfillcolor{currentfill}%
\pgfsetfillopacity{0.120000}%
\pgfsetlinewidth{0.000000pt}%
\definecolor{currentstroke}{rgb}{0.000000,0.000000,0.000000}%
\pgfsetstrokecolor{currentstroke}%
\pgfsetdash{}{0pt}%
\pgfpathmoveto{\pgfqpoint{1.358964in}{1.119878in}}%
\pgfpathlineto{\pgfqpoint{1.375304in}{1.120079in}}%
\pgfpathlineto{\pgfqpoint{1.391700in}{1.120917in}}%
\pgfpathlineto{\pgfqpoint{1.399318in}{1.134643in}}%
\pgfpathlineto{\pgfqpoint{1.406923in}{1.148521in}}%
\pgfpathlineto{\pgfqpoint{1.390541in}{1.147834in}}%
\pgfpathlineto{\pgfqpoint{1.374215in}{1.147760in}}%
\pgfpathlineto{\pgfqpoint{1.366596in}{1.133745in}}%
\pgfpathlineto{\pgfqpoint{1.358964in}{1.119878in}}%
\pgfpathclose%
\pgfusepath{fill}%
\end{pgfscope}%
\begin{pgfscope}%
\pgfpathrectangle{\pgfqpoint{0.407592in}{0.222904in}}{\pgfqpoint{1.411163in}{1.411163in}}%
\pgfusepath{clip}%
\pgfsetbuttcap%
\pgfsetroundjoin%
\definecolor{currentfill}{rgb}{0.909804,0.698039,0.423529}%
\pgfsetfillcolor{currentfill}%
\pgfsetfillopacity{0.120000}%
\pgfsetlinewidth{0.000000pt}%
\definecolor{currentstroke}{rgb}{0.000000,0.000000,0.000000}%
\pgfsetstrokecolor{currentstroke}%
\pgfsetdash{}{0pt}%
\pgfpathmoveto{\pgfqpoint{0.976982in}{1.032183in}}%
\pgfpathlineto{\pgfqpoint{0.992891in}{1.019782in}}%
\pgfpathlineto{\pgfqpoint{1.008820in}{1.008004in}}%
\pgfpathlineto{\pgfqpoint{1.016859in}{1.016748in}}%
\pgfpathlineto{\pgfqpoint{1.024876in}{1.025903in}}%
\pgfpathlineto{\pgfqpoint{1.008977in}{1.037445in}}%
\pgfpathlineto{\pgfqpoint{0.993099in}{1.049603in}}%
\pgfpathlineto{\pgfqpoint{0.985052in}{1.040695in}}%
\pgfpathlineto{\pgfqpoint{0.976982in}{1.032183in}}%
\pgfpathclose%
\pgfusepath{fill}%
\end{pgfscope}%
\begin{pgfscope}%
\pgfpathrectangle{\pgfqpoint{0.407592in}{0.222904in}}{\pgfqpoint{1.411163in}{1.411163in}}%
\pgfusepath{clip}%
\pgfsetbuttcap%
\pgfsetroundjoin%
\definecolor{currentfill}{rgb}{0.909804,0.698039,0.423529}%
\pgfsetfillcolor{currentfill}%
\pgfsetfillopacity{0.120000}%
\pgfsetlinewidth{0.000000pt}%
\definecolor{currentstroke}{rgb}{0.000000,0.000000,0.000000}%
\pgfsetstrokecolor{currentstroke}%
\pgfsetdash{}{0pt}%
\pgfpathmoveto{\pgfqpoint{1.152330in}{1.014947in}}%
\pgfpathlineto{\pgfqpoint{1.168376in}{1.009143in}}%
\pgfpathlineto{\pgfqpoint{1.184459in}{1.004188in}}%
\pgfpathlineto{\pgfqpoint{1.192283in}{1.015888in}}%
\pgfpathlineto{\pgfqpoint{1.200090in}{1.027897in}}%
\pgfpathlineto{\pgfqpoint{1.184024in}{1.032768in}}%
\pgfpathlineto{\pgfqpoint{1.167997in}{1.038452in}}%
\pgfpathlineto{\pgfqpoint{1.160172in}{1.026546in}}%
\pgfpathlineto{\pgfqpoint{1.152330in}{1.014947in}}%
\pgfpathclose%
\pgfusepath{fill}%
\end{pgfscope}%
\begin{pgfscope}%
\pgfpathrectangle{\pgfqpoint{0.407592in}{0.222904in}}{\pgfqpoint{1.411163in}{1.411163in}}%
\pgfusepath{clip}%
\pgfsetbuttcap%
\pgfsetroundjoin%
\definecolor{currentfill}{rgb}{0.909804,0.698039,0.423529}%
\pgfsetfillcolor{currentfill}%
\pgfsetfillopacity{0.120000}%
\pgfsetlinewidth{0.000000pt}%
\definecolor{currentstroke}{rgb}{0.000000,0.000000,0.000000}%
\pgfsetstrokecolor{currentstroke}%
\pgfsetdash{}{0pt}%
\pgfpathmoveto{\pgfqpoint{1.200090in}{1.027897in}}%
\pgfpathlineto{\pgfqpoint{1.216196in}{1.023850in}}%
\pgfpathlineto{\pgfqpoint{1.232345in}{1.020630in}}%
\pgfpathlineto{\pgfqpoint{1.240119in}{1.032956in}}%
\pgfpathlineto{\pgfqpoint{1.247878in}{1.045550in}}%
\pgfpathlineto{\pgfqpoint{1.231745in}{1.048758in}}%
\pgfpathlineto{\pgfqpoint{1.215655in}{1.052757in}}%
\pgfpathlineto{\pgfqpoint{1.207880in}{1.040194in}}%
\pgfpathlineto{\pgfqpoint{1.200090in}{1.027897in}}%
\pgfpathclose%
\pgfusepath{fill}%
\end{pgfscope}%
\begin{pgfscope}%
\pgfpathrectangle{\pgfqpoint{0.407592in}{0.222904in}}{\pgfqpoint{1.411163in}{1.411163in}}%
\pgfusepath{clip}%
\pgfsetbuttcap%
\pgfsetroundjoin%
\definecolor{currentfill}{rgb}{0.909804,0.698039,0.423529}%
\pgfsetfillcolor{currentfill}%
\pgfsetfillopacity{0.120000}%
\pgfsetlinewidth{0.000000pt}%
\definecolor{currentstroke}{rgb}{0.000000,0.000000,0.000000}%
\pgfsetstrokecolor{currentstroke}%
\pgfsetdash{}{0pt}%
\pgfpathmoveto{\pgfqpoint{1.104560in}{1.007171in}}%
\pgfpathlineto{\pgfqpoint{1.120559in}{0.999547in}}%
\pgfpathlineto{\pgfqpoint{1.136591in}{0.992761in}}%
\pgfpathlineto{\pgfqpoint{1.144469in}{1.003678in}}%
\pgfpathlineto{\pgfqpoint{1.152330in}{1.014947in}}%
\pgfpathlineto{\pgfqpoint{1.136319in}{1.021583in}}%
\pgfpathlineto{\pgfqpoint{1.120342in}{1.029027in}}%
\pgfpathlineto{\pgfqpoint{1.112460in}{1.017925in}}%
\pgfpathlineto{\pgfqpoint{1.104560in}{1.007171in}}%
\pgfpathclose%
\pgfusepath{fill}%
\end{pgfscope}%
\begin{pgfscope}%
\pgfpathrectangle{\pgfqpoint{0.407592in}{0.222904in}}{\pgfqpoint{1.411163in}{1.411163in}}%
\pgfusepath{clip}%
\pgfsetbuttcap%
\pgfsetroundjoin%
\definecolor{currentfill}{rgb}{0.909804,0.698039,0.423529}%
\pgfsetfillcolor{currentfill}%
\pgfsetfillopacity{0.120000}%
\pgfsetlinewidth{0.000000pt}%
\definecolor{currentstroke}{rgb}{0.000000,0.000000,0.000000}%
\pgfsetstrokecolor{currentstroke}%
\pgfsetdash{}{0pt}%
\pgfpathmoveto{\pgfqpoint{0.817142in}{1.121783in}}%
\pgfpathlineto{\pgfqpoint{0.833011in}{1.105546in}}%
\pgfpathlineto{\pgfqpoint{0.848891in}{1.089615in}}%
\pgfpathlineto{\pgfqpoint{0.857177in}{1.095949in}}%
\pgfpathlineto{\pgfqpoint{0.865433in}{1.102660in}}%
\pgfpathlineto{\pgfqpoint{0.849591in}{1.118431in}}%
\pgfpathlineto{\pgfqpoint{0.833760in}{1.134511in}}%
\pgfpathlineto{\pgfqpoint{0.825466in}{1.127969in}}%
\pgfpathlineto{\pgfqpoint{0.817142in}{1.121783in}}%
\pgfpathclose%
\pgfusepath{fill}%
\end{pgfscope}%
\begin{pgfscope}%
\pgfpathrectangle{\pgfqpoint{0.407592in}{0.222904in}}{\pgfqpoint{1.411163in}{1.411163in}}%
\pgfusepath{clip}%
\pgfsetbuttcap%
\pgfsetroundjoin%
\definecolor{currentfill}{rgb}{0.909804,0.698039,0.423529}%
\pgfsetfillcolor{currentfill}%
\pgfsetfillopacity{0.120000}%
\pgfsetlinewidth{0.000000pt}%
\definecolor{currentstroke}{rgb}{0.000000,0.000000,0.000000}%
\pgfsetstrokecolor{currentstroke}%
\pgfsetdash{}{0pt}%
\pgfpathmoveto{\pgfqpoint{0.705171in}{1.213622in}}%
\pgfpathlineto{\pgfqpoint{0.721018in}{1.195940in}}%
\pgfpathlineto{\pgfqpoint{0.736874in}{1.178432in}}%
\pgfpathlineto{\pgfqpoint{0.745340in}{1.183693in}}%
\pgfpathlineto{\pgfqpoint{0.753774in}{1.189279in}}%
\pgfpathlineto{\pgfqpoint{0.737958in}{1.206682in}}%
\pgfpathlineto{\pgfqpoint{0.722152in}{1.224261in}}%
\pgfpathlineto{\pgfqpoint{0.713678in}{1.218788in}}%
\pgfpathlineto{\pgfqpoint{0.705171in}{1.213622in}}%
\pgfpathclose%
\pgfusepath{fill}%
\end{pgfscope}%
\begin{pgfscope}%
\pgfpathrectangle{\pgfqpoint{0.407592in}{0.222904in}}{\pgfqpoint{1.411163in}{1.411163in}}%
\pgfusepath{clip}%
\pgfsetbuttcap%
\pgfsetroundjoin%
\definecolor{currentfill}{rgb}{0.909804,0.698039,0.423529}%
\pgfsetfillcolor{currentfill}%
\pgfsetfillopacity{0.120000}%
\pgfsetlinewidth{0.000000pt}%
\definecolor{currentstroke}{rgb}{0.000000,0.000000,0.000000}%
\pgfsetstrokecolor{currentstroke}%
\pgfsetdash{}{0pt}%
\pgfpathmoveto{\pgfqpoint{1.247878in}{1.045550in}}%
\pgfpathlineto{\pgfqpoint{1.264057in}{1.043131in}}%
\pgfpathlineto{\pgfqpoint{1.280284in}{1.041495in}}%
\pgfpathlineto{\pgfqpoint{1.288012in}{1.054297in}}%
\pgfpathlineto{\pgfqpoint{1.295726in}{1.067329in}}%
\pgfpathlineto{\pgfqpoint{1.279514in}{1.069022in}}%
\pgfpathlineto{\pgfqpoint{1.263350in}{1.071463in}}%
\pgfpathlineto{\pgfqpoint{1.255621in}{1.058392in}}%
\pgfpathlineto{\pgfqpoint{1.247878in}{1.045550in}}%
\pgfpathclose%
\pgfusepath{fill}%
\end{pgfscope}%
\begin{pgfscope}%
\pgfpathrectangle{\pgfqpoint{0.407592in}{0.222904in}}{\pgfqpoint{1.411163in}{1.411163in}}%
\pgfusepath{clip}%
\pgfsetbuttcap%
\pgfsetroundjoin%
\definecolor{currentfill}{rgb}{0.909804,0.698039,0.423529}%
\pgfsetfillcolor{currentfill}%
\pgfsetfillopacity{0.120000}%
\pgfsetlinewidth{0.000000pt}%
\definecolor{currentstroke}{rgb}{0.000000,0.000000,0.000000}%
\pgfsetstrokecolor{currentstroke}%
\pgfsetdash{}{0pt}%
\pgfpathmoveto{\pgfqpoint{1.056738in}{1.004849in}}%
\pgfpathlineto{\pgfqpoint{1.072705in}{0.995427in}}%
\pgfpathlineto{\pgfqpoint{1.088699in}{0.986799in}}%
\pgfpathlineto{\pgfqpoint{1.096640in}{0.996789in}}%
\pgfpathlineto{\pgfqpoint{1.104560in}{1.007171in}}%
\pgfpathlineto{\pgfqpoint{1.088590in}{1.015603in}}%
\pgfpathlineto{\pgfqpoint{1.072648in}{1.024807in}}%
\pgfpathlineto{\pgfqpoint{1.064704in}{1.014636in}}%
\pgfpathlineto{\pgfqpoint{1.056738in}{1.004849in}}%
\pgfpathclose%
\pgfusepath{fill}%
\end{pgfscope}%
\begin{pgfscope}%
\pgfpathrectangle{\pgfqpoint{0.407592in}{0.222904in}}{\pgfqpoint{1.411163in}{1.411163in}}%
\pgfusepath{clip}%
\pgfsetbuttcap%
\pgfsetroundjoin%
\definecolor{currentfill}{rgb}{0.909804,0.698039,0.423529}%
\pgfsetfillcolor{currentfill}%
\pgfsetfillopacity{0.120000}%
\pgfsetlinewidth{0.000000pt}%
\definecolor{currentstroke}{rgb}{0.000000,0.000000,0.000000}%
\pgfsetstrokecolor{currentstroke}%
\pgfsetdash{}{0pt}%
\pgfpathmoveto{\pgfqpoint{0.928930in}{1.043308in}}%
\pgfpathlineto{\pgfqpoint{0.944839in}{1.029589in}}%
\pgfpathlineto{\pgfqpoint{0.960764in}{1.016402in}}%
\pgfpathlineto{\pgfqpoint{0.968886in}{1.024081in}}%
\pgfpathlineto{\pgfqpoint{0.976982in}{1.032183in}}%
\pgfpathlineto{\pgfqpoint{0.961090in}{1.045160in}}%
\pgfpathlineto{\pgfqpoint{0.945215in}{1.058668in}}%
\pgfpathlineto{\pgfqpoint{0.937085in}{1.050786in}}%
\pgfpathlineto{\pgfqpoint{0.928930in}{1.043308in}}%
\pgfpathclose%
\pgfusepath{fill}%
\end{pgfscope}%
\begin{pgfscope}%
\pgfpathrectangle{\pgfqpoint{0.407592in}{0.222904in}}{\pgfqpoint{1.411163in}{1.411163in}}%
\pgfusepath{clip}%
\pgfsetbuttcap%
\pgfsetroundjoin%
\definecolor{currentfill}{rgb}{0.909804,0.698039,0.423529}%
\pgfsetfillcolor{currentfill}%
\pgfsetfillopacity{0.120000}%
\pgfsetlinewidth{0.000000pt}%
\definecolor{currentstroke}{rgb}{0.000000,0.000000,0.000000}%
\pgfsetstrokecolor{currentstroke}%
\pgfsetdash{}{0pt}%
\pgfpathmoveto{\pgfqpoint{1.295726in}{1.067329in}}%
\pgfpathlineto{\pgfqpoint{1.311989in}{1.066375in}}%
\pgfpathlineto{\pgfqpoint{1.328303in}{1.066145in}}%
\pgfpathlineto{\pgfqpoint{1.335989in}{1.079291in}}%
\pgfpathlineto{\pgfqpoint{1.343660in}{1.092638in}}%
\pgfpathlineto{\pgfqpoint{1.327360in}{1.092982in}}%
\pgfpathlineto{\pgfqpoint{1.311112in}{1.094022in}}%
\pgfpathlineto{\pgfqpoint{1.303426in}{1.080576in}}%
\pgfpathlineto{\pgfqpoint{1.295726in}{1.067329in}}%
\pgfpathclose%
\pgfusepath{fill}%
\end{pgfscope}%
\begin{pgfscope}%
\pgfpathrectangle{\pgfqpoint{0.407592in}{0.222904in}}{\pgfqpoint{1.411163in}{1.411163in}}%
\pgfusepath{clip}%
\pgfsetbuttcap%
\pgfsetroundjoin%
\definecolor{currentfill}{rgb}{0.909804,0.698039,0.423529}%
\pgfsetfillcolor{currentfill}%
\pgfsetfillopacity{0.120000}%
\pgfsetlinewidth{0.000000pt}%
\definecolor{currentstroke}{rgb}{0.000000,0.000000,0.000000}%
\pgfsetstrokecolor{currentstroke}%
\pgfsetdash{}{0pt}%
\pgfpathmoveto{\pgfqpoint{0.736874in}{1.178432in}}%
\pgfpathlineto{\pgfqpoint{0.752741in}{1.161114in}}%
\pgfpathlineto{\pgfqpoint{0.768617in}{1.144006in}}%
\pgfpathlineto{\pgfqpoint{0.777043in}{1.149364in}}%
\pgfpathlineto{\pgfqpoint{0.785437in}{1.155067in}}%
\pgfpathlineto{\pgfqpoint{0.769601in}{1.172068in}}%
\pgfpathlineto{\pgfqpoint{0.753774in}{1.189279in}}%
\pgfpathlineto{\pgfqpoint{0.745340in}{1.183693in}}%
\pgfpathlineto{\pgfqpoint{0.736874in}{1.178432in}}%
\pgfpathclose%
\pgfusepath{fill}%
\end{pgfscope}%
\begin{pgfscope}%
\pgfpathrectangle{\pgfqpoint{0.407592in}{0.222904in}}{\pgfqpoint{1.411163in}{1.411163in}}%
\pgfusepath{clip}%
\pgfsetbuttcap%
\pgfsetroundjoin%
\definecolor{currentfill}{rgb}{0.909804,0.698039,0.423529}%
\pgfsetfillcolor{currentfill}%
\pgfsetfillopacity{0.120000}%
\pgfsetlinewidth{0.000000pt}%
\definecolor{currentstroke}{rgb}{0.000000,0.000000,0.000000}%
\pgfsetstrokecolor{currentstroke}%
\pgfsetdash{}{0pt}%
\pgfpathmoveto{\pgfqpoint{1.406923in}{1.148521in}}%
\pgfpathlineto{\pgfqpoint{1.423363in}{1.149807in}}%
\pgfpathlineto{\pgfqpoint{1.439860in}{1.151675in}}%
\pgfpathlineto{\pgfqpoint{1.447439in}{1.165521in}}%
\pgfpathlineto{\pgfqpoint{1.430948in}{1.163745in}}%
\pgfpathlineto{\pgfqpoint{1.414516in}{1.162542in}}%
\pgfusepath{fill}%
\end{pgfscope}%
\begin{pgfscope}%
\pgfpathrectangle{\pgfqpoint{0.407592in}{0.222904in}}{\pgfqpoint{1.411163in}{1.411163in}}%
\pgfusepath{clip}%
\pgfsetbuttcap%
\pgfsetroundjoin%
\definecolor{currentfill}{rgb}{0.909804,0.698039,0.423529}%
\pgfsetfillcolor{currentfill}%
\pgfsetfillopacity{0.120000}%
\pgfsetlinewidth{0.000000pt}%
\definecolor{currentstroke}{rgb}{0.000000,0.000000,0.000000}%
\pgfsetstrokecolor{currentstroke}%
\pgfsetdash{}{0pt}%
\pgfpathmoveto{\pgfqpoint{0.848891in}{1.089615in}}%
\pgfpathlineto{\pgfqpoint{0.864783in}{1.074021in}}%
\pgfpathlineto{\pgfqpoint{0.880687in}{1.058797in}}%
\pgfpathlineto{\pgfqpoint{0.888935in}{1.065282in}}%
\pgfpathlineto{\pgfqpoint{0.897155in}{1.072168in}}%
\pgfpathlineto{\pgfqpoint{0.881287in}{1.087228in}}%
\pgfpathlineto{\pgfqpoint{0.865433in}{1.102660in}}%
\pgfpathlineto{\pgfqpoint{0.857177in}{1.095949in}}%
\pgfpathlineto{\pgfqpoint{0.848891in}{1.089615in}}%
\pgfpathclose%
\pgfusepath{fill}%
\end{pgfscope}%
\begin{pgfscope}%
\pgfpathrectangle{\pgfqpoint{0.407592in}{0.222904in}}{\pgfqpoint{1.411163in}{1.411163in}}%
\pgfusepath{clip}%
\pgfsetbuttcap%
\pgfsetroundjoin%
\definecolor{currentfill}{rgb}{0.909804,0.698039,0.423529}%
\pgfsetfillcolor{currentfill}%
\pgfsetfillopacity{0.120000}%
\pgfsetlinewidth{0.000000pt}%
\definecolor{currentstroke}{rgb}{0.000000,0.000000,0.000000}%
\pgfsetstrokecolor{currentstroke}%
\pgfsetdash{}{0pt}%
\pgfpathmoveto{\pgfqpoint{1.008820in}{1.008004in}}%
\pgfpathlineto{\pgfqpoint{1.024769in}{0.996898in}}%
\pgfpathlineto{\pgfqpoint{1.040740in}{0.986511in}}%
\pgfpathlineto{\pgfqpoint{1.048750in}{0.995467in}}%
\pgfpathlineto{\pgfqpoint{1.056738in}{1.004849in}}%
\pgfpathlineto{\pgfqpoint{1.040795in}{1.015022in}}%
\pgfpathlineto{\pgfqpoint{1.024876in}{1.025903in}}%
\pgfpathlineto{\pgfqpoint{1.016859in}{1.016748in}}%
\pgfpathlineto{\pgfqpoint{1.008820in}{1.008004in}}%
\pgfpathclose%
\pgfusepath{fill}%
\end{pgfscope}%
\begin{pgfscope}%
\pgfpathrectangle{\pgfqpoint{0.407592in}{0.222904in}}{\pgfqpoint{1.411163in}{1.411163in}}%
\pgfusepath{clip}%
\pgfsetbuttcap%
\pgfsetroundjoin%
\definecolor{currentfill}{rgb}{0.909804,0.698039,0.423529}%
\pgfsetfillcolor{currentfill}%
\pgfsetfillopacity{0.120000}%
\pgfsetlinewidth{0.000000pt}%
\definecolor{currentstroke}{rgb}{0.000000,0.000000,0.000000}%
\pgfsetstrokecolor{currentstroke}%
\pgfsetdash{}{0pt}%
\pgfpathmoveto{\pgfqpoint{1.343660in}{1.092638in}}%
\pgfpathlineto{\pgfqpoint{1.360015in}{1.092973in}}%
\pgfpathlineto{\pgfqpoint{1.376425in}{1.093971in}}%
\pgfpathlineto{\pgfqpoint{1.384069in}{1.107355in}}%
\pgfpathlineto{\pgfqpoint{1.391700in}{1.120917in}}%
\pgfpathlineto{\pgfqpoint{1.375304in}{1.120079in}}%
\pgfpathlineto{\pgfqpoint{1.358964in}{1.119878in}}%
\pgfpathlineto{\pgfqpoint{1.351319in}{1.106172in}}%
\pgfpathlineto{\pgfqpoint{1.343660in}{1.092638in}}%
\pgfpathclose%
\pgfusepath{fill}%
\end{pgfscope}%
\begin{pgfscope}%
\pgfpathrectangle{\pgfqpoint{0.407592in}{0.222904in}}{\pgfqpoint{1.411163in}{1.411163in}}%
\pgfusepath{clip}%
\pgfsetbuttcap%
\pgfsetroundjoin%
\definecolor{currentfill}{rgb}{0.909804,0.698039,0.423529}%
\pgfsetfillcolor{currentfill}%
\pgfsetfillopacity{0.120000}%
\pgfsetlinewidth{0.000000pt}%
\definecolor{currentstroke}{rgb}{0.000000,0.000000,0.000000}%
\pgfsetstrokecolor{currentstroke}%
\pgfsetdash{}{0pt}%
\pgfpathmoveto{\pgfqpoint{0.768617in}{1.144006in}}%
\pgfpathlineto{\pgfqpoint{0.784503in}{1.127127in}}%
\pgfpathlineto{\pgfqpoint{0.800401in}{1.110500in}}%
\pgfpathlineto{\pgfqpoint{0.808787in}{1.115958in}}%
\pgfpathlineto{\pgfqpoint{0.817142in}{1.121783in}}%
\pgfpathlineto{\pgfqpoint{0.801284in}{1.138298in}}%
\pgfpathlineto{\pgfqpoint{0.785437in}{1.155067in}}%
\pgfpathlineto{\pgfqpoint{0.777043in}{1.149364in}}%
\pgfpathlineto{\pgfqpoint{0.768617in}{1.144006in}}%
\pgfpathclose%
\pgfusepath{fill}%
\end{pgfscope}%
\begin{pgfscope}%
\pgfpathrectangle{\pgfqpoint{0.407592in}{0.222904in}}{\pgfqpoint{1.411163in}{1.411163in}}%
\pgfusepath{clip}%
\pgfsetbuttcap%
\pgfsetroundjoin%
\definecolor{currentfill}{rgb}{0.909804,0.698039,0.423529}%
\pgfsetfillcolor{currentfill}%
\pgfsetfillopacity{0.120000}%
\pgfsetlinewidth{0.000000pt}%
\definecolor{currentstroke}{rgb}{0.000000,0.000000,0.000000}%
\pgfsetstrokecolor{currentstroke}%
\pgfsetdash{}{0pt}%
\pgfpathmoveto{\pgfqpoint{1.136591in}{0.992761in}}%
\pgfpathlineto{\pgfqpoint{1.152657in}{0.986840in}}%
\pgfpathlineto{\pgfqpoint{1.168761in}{0.981802in}}%
\pgfpathlineto{\pgfqpoint{1.176619in}{0.992818in}}%
\pgfpathlineto{\pgfqpoint{1.184459in}{1.004188in}}%
\pgfpathlineto{\pgfqpoint{1.168376in}{1.009143in}}%
\pgfpathlineto{\pgfqpoint{1.152330in}{1.014947in}}%
\pgfpathlineto{\pgfqpoint{1.144469in}{1.003678in}}%
\pgfpathlineto{\pgfqpoint{1.136591in}{0.992761in}}%
\pgfpathclose%
\pgfusepath{fill}%
\end{pgfscope}%
\begin{pgfscope}%
\pgfpathrectangle{\pgfqpoint{0.407592in}{0.222904in}}{\pgfqpoint{1.411163in}{1.411163in}}%
\pgfusepath{clip}%
\pgfsetbuttcap%
\pgfsetroundjoin%
\definecolor{currentfill}{rgb}{0.909804,0.698039,0.423529}%
\pgfsetfillcolor{currentfill}%
\pgfsetfillopacity{0.120000}%
\pgfsetlinewidth{0.000000pt}%
\definecolor{currentstroke}{rgb}{0.000000,0.000000,0.000000}%
\pgfsetstrokecolor{currentstroke}%
\pgfsetdash{}{0pt}%
\pgfpathmoveto{\pgfqpoint{1.184459in}{1.004188in}}%
\pgfpathlineto{\pgfqpoint{1.200583in}{1.000093in}}%
\pgfpathlineto{\pgfqpoint{1.216750in}{0.996863in}}%
\pgfpathlineto{\pgfqpoint{1.224556in}{1.008592in}}%
\pgfpathlineto{\pgfqpoint{1.232345in}{1.020630in}}%
\pgfpathlineto{\pgfqpoint{1.216196in}{1.023850in}}%
\pgfpathlineto{\pgfqpoint{1.200090in}{1.027897in}}%
\pgfpathlineto{\pgfqpoint{1.192283in}{1.015888in}}%
\pgfpathlineto{\pgfqpoint{1.184459in}{1.004188in}}%
\pgfpathclose%
\pgfusepath{fill}%
\end{pgfscope}%
\begin{pgfscope}%
\pgfpathrectangle{\pgfqpoint{0.407592in}{0.222904in}}{\pgfqpoint{1.411163in}{1.411163in}}%
\pgfusepath{clip}%
\pgfsetbuttcap%
\pgfsetroundjoin%
\definecolor{currentfill}{rgb}{0.909804,0.698039,0.423529}%
\pgfsetfillcolor{currentfill}%
\pgfsetfillopacity{0.120000}%
\pgfsetlinewidth{0.000000pt}%
\definecolor{currentstroke}{rgb}{0.000000,0.000000,0.000000}%
\pgfsetstrokecolor{currentstroke}%
\pgfsetdash{}{0pt}%
\pgfpathmoveto{\pgfqpoint{0.656311in}{1.240179in}}%
\pgfpathlineto{\pgfqpoint{0.672180in}{1.222126in}}%
\pgfpathlineto{\pgfqpoint{0.688058in}{1.204219in}}%
\pgfpathlineto{\pgfqpoint{0.696632in}{1.208765in}}%
\pgfpathlineto{\pgfqpoint{0.705171in}{1.213622in}}%
\pgfpathlineto{\pgfqpoint{0.689334in}{1.231462in}}%
\pgfpathlineto{\pgfqpoint{0.673507in}{1.249448in}}%
\pgfpathlineto{\pgfqpoint{0.664926in}{1.244667in}}%
\pgfpathlineto{\pgfqpoint{0.656311in}{1.240179in}}%
\pgfpathclose%
\pgfusepath{fill}%
\end{pgfscope}%
\begin{pgfscope}%
\pgfpathrectangle{\pgfqpoint{0.407592in}{0.222904in}}{\pgfqpoint{1.411163in}{1.411163in}}%
\pgfusepath{clip}%
\pgfsetbuttcap%
\pgfsetroundjoin%
\definecolor{currentfill}{rgb}{0.909804,0.698039,0.423529}%
\pgfsetfillcolor{currentfill}%
\pgfsetfillopacity{0.120000}%
\pgfsetlinewidth{0.000000pt}%
\definecolor{currentstroke}{rgb}{0.000000,0.000000,0.000000}%
\pgfsetstrokecolor{currentstroke}%
\pgfsetdash{}{0pt}%
\pgfpathmoveto{\pgfqpoint{0.960764in}{1.016402in}}%
\pgfpathlineto{\pgfqpoint{0.976707in}{1.003795in}}%
\pgfpathlineto{\pgfqpoint{0.992667in}{0.991816in}}%
\pgfpathlineto{\pgfqpoint{1.000756in}{0.999688in}}%
\pgfpathlineto{\pgfqpoint{1.008820in}{1.008004in}}%
\pgfpathlineto{\pgfqpoint{0.992891in}{1.019782in}}%
\pgfpathlineto{\pgfqpoint{0.976982in}{1.032183in}}%
\pgfpathlineto{\pgfqpoint{0.968886in}{1.024081in}}%
\pgfpathlineto{\pgfqpoint{0.960764in}{1.016402in}}%
\pgfpathclose%
\pgfusepath{fill}%
\end{pgfscope}%
\begin{pgfscope}%
\pgfpathrectangle{\pgfqpoint{0.407592in}{0.222904in}}{\pgfqpoint{1.411163in}{1.411163in}}%
\pgfusepath{clip}%
\pgfsetbuttcap%
\pgfsetroundjoin%
\definecolor{currentfill}{rgb}{0.909804,0.698039,0.423529}%
\pgfsetfillcolor{currentfill}%
\pgfsetfillopacity{0.120000}%
\pgfsetlinewidth{0.000000pt}%
\definecolor{currentstroke}{rgb}{0.000000,0.000000,0.000000}%
\pgfsetstrokecolor{currentstroke}%
\pgfsetdash{}{0pt}%
\pgfpathmoveto{\pgfqpoint{0.880687in}{1.058797in}}%
\pgfpathlineto{\pgfqpoint{0.896604in}{1.043978in}}%
\pgfpathlineto{\pgfqpoint{0.912535in}{1.029605in}}%
\pgfpathlineto{\pgfqpoint{0.920746in}{1.036245in}}%
\pgfpathlineto{\pgfqpoint{0.928930in}{1.043308in}}%
\pgfpathlineto{\pgfqpoint{0.913035in}{1.057515in}}%
\pgfpathlineto{\pgfqpoint{0.897155in}{1.072168in}}%
\pgfpathlineto{\pgfqpoint{0.888935in}{1.065282in}}%
\pgfpathlineto{\pgfqpoint{0.880687in}{1.058797in}}%
\pgfpathclose%
\pgfusepath{fill}%
\end{pgfscope}%
\begin{pgfscope}%
\pgfpathrectangle{\pgfqpoint{0.407592in}{0.222904in}}{\pgfqpoint{1.411163in}{1.411163in}}%
\pgfusepath{clip}%
\pgfsetbuttcap%
\pgfsetroundjoin%
\definecolor{currentfill}{rgb}{0.909804,0.698039,0.423529}%
\pgfsetfillcolor{currentfill}%
\pgfsetfillopacity{0.120000}%
\pgfsetlinewidth{0.000000pt}%
\definecolor{currentstroke}{rgb}{0.000000,0.000000,0.000000}%
\pgfsetstrokecolor{currentstroke}%
\pgfsetdash{}{0pt}%
\pgfpathmoveto{\pgfqpoint{1.088699in}{0.986799in}}%
\pgfpathlineto{\pgfqpoint{1.104722in}{0.979004in}}%
\pgfpathlineto{\pgfqpoint{1.120777in}{0.972075in}}%
\pgfpathlineto{\pgfqpoint{1.128694in}{0.982219in}}%
\pgfpathlineto{\pgfqpoint{1.136591in}{0.992761in}}%
\pgfpathlineto{\pgfqpoint{1.120559in}{0.999547in}}%
\pgfpathlineto{\pgfqpoint{1.104560in}{1.007171in}}%
\pgfpathlineto{\pgfqpoint{1.096640in}{0.996789in}}%
\pgfpathlineto{\pgfqpoint{1.088699in}{0.986799in}}%
\pgfpathclose%
\pgfusepath{fill}%
\end{pgfscope}%
\begin{pgfscope}%
\pgfpathrectangle{\pgfqpoint{0.407592in}{0.222904in}}{\pgfqpoint{1.411163in}{1.411163in}}%
\pgfusepath{clip}%
\pgfsetbuttcap%
\pgfsetroundjoin%
\definecolor{currentfill}{rgb}{0.909804,0.698039,0.423529}%
\pgfsetfillcolor{currentfill}%
\pgfsetfillopacity{0.120000}%
\pgfsetlinewidth{0.000000pt}%
\definecolor{currentstroke}{rgb}{0.000000,0.000000,0.000000}%
\pgfsetstrokecolor{currentstroke}%
\pgfsetdash{}{0pt}%
\pgfpathmoveto{\pgfqpoint{1.232345in}{1.020630in}}%
\pgfpathlineto{\pgfqpoint{1.248540in}{1.018236in}}%
\pgfpathlineto{\pgfqpoint{1.264783in}{1.016660in}}%
\pgfpathlineto{\pgfqpoint{1.272541in}{1.028944in}}%
\pgfpathlineto{\pgfqpoint{1.280284in}{1.041495in}}%
\pgfpathlineto{\pgfqpoint{1.264057in}{1.043131in}}%
\pgfpathlineto{\pgfqpoint{1.247878in}{1.045550in}}%
\pgfpathlineto{\pgfqpoint{1.240119in}{1.032956in}}%
\pgfpathlineto{\pgfqpoint{1.232345in}{1.020630in}}%
\pgfpathclose%
\pgfusepath{fill}%
\end{pgfscope}%
\begin{pgfscope}%
\pgfpathrectangle{\pgfqpoint{0.407592in}{0.222904in}}{\pgfqpoint{1.411163in}{1.411163in}}%
\pgfusepath{clip}%
\pgfsetbuttcap%
\pgfsetroundjoin%
\definecolor{currentfill}{rgb}{0.909804,0.698039,0.423529}%
\pgfsetfillcolor{currentfill}%
\pgfsetfillopacity{0.120000}%
\pgfsetlinewidth{0.000000pt}%
\definecolor{currentstroke}{rgb}{0.000000,0.000000,0.000000}%
\pgfsetstrokecolor{currentstroke}%
\pgfsetdash{}{0pt}%
\pgfpathmoveto{\pgfqpoint{1.391700in}{1.120917in}}%
\pgfpathlineto{\pgfqpoint{1.408154in}{1.122375in}}%
\pgfpathlineto{\pgfqpoint{1.424666in}{1.124435in}}%
\pgfpathlineto{\pgfqpoint{1.432269in}{1.137976in}}%
\pgfpathlineto{\pgfqpoint{1.439860in}{1.151675in}}%
\pgfpathlineto{\pgfqpoint{1.423363in}{1.149807in}}%
\pgfpathlineto{\pgfqpoint{1.406923in}{1.148521in}}%
\pgfpathlineto{\pgfqpoint{1.399318in}{1.134643in}}%
\pgfpathlineto{\pgfqpoint{1.391700in}{1.120917in}}%
\pgfpathclose%
\pgfusepath{fill}%
\end{pgfscope}%
\begin{pgfscope}%
\pgfpathrectangle{\pgfqpoint{0.407592in}{0.222904in}}{\pgfqpoint{1.411163in}{1.411163in}}%
\pgfusepath{clip}%
\pgfsetbuttcap%
\pgfsetroundjoin%
\definecolor{currentfill}{rgb}{0.909804,0.698039,0.423529}%
\pgfsetfillcolor{currentfill}%
\pgfsetfillopacity{0.120000}%
\pgfsetlinewidth{0.000000pt}%
\definecolor{currentstroke}{rgb}{0.000000,0.000000,0.000000}%
\pgfsetstrokecolor{currentstroke}%
\pgfsetdash{}{0pt}%
\pgfpathmoveto{\pgfqpoint{0.800401in}{1.110500in}}%
\pgfpathlineto{\pgfqpoint{0.816308in}{1.094150in}}%
\pgfpathlineto{\pgfqpoint{0.832227in}{1.078106in}}%
\pgfpathlineto{\pgfqpoint{0.840575in}{1.083666in}}%
\pgfpathlineto{\pgfqpoint{0.848891in}{1.089615in}}%
\pgfpathlineto{\pgfqpoint{0.833011in}{1.105546in}}%
\pgfpathlineto{\pgfqpoint{0.817142in}{1.121783in}}%
\pgfpathlineto{\pgfqpoint{0.808787in}{1.115958in}}%
\pgfpathlineto{\pgfqpoint{0.800401in}{1.110500in}}%
\pgfpathclose%
\pgfusepath{fill}%
\end{pgfscope}%
\begin{pgfscope}%
\pgfpathrectangle{\pgfqpoint{0.407592in}{0.222904in}}{\pgfqpoint{1.411163in}{1.411163in}}%
\pgfusepath{clip}%
\pgfsetbuttcap%
\pgfsetroundjoin%
\definecolor{currentfill}{rgb}{0.909804,0.698039,0.423529}%
\pgfsetfillcolor{currentfill}%
\pgfsetfillopacity{0.120000}%
\pgfsetlinewidth{0.000000pt}%
\definecolor{currentstroke}{rgb}{0.000000,0.000000,0.000000}%
\pgfsetstrokecolor{currentstroke}%
\pgfsetdash{}{0pt}%
\pgfpathmoveto{\pgfqpoint{1.040740in}{0.986511in}}%
\pgfpathlineto{\pgfqpoint{1.056734in}{0.976891in}}%
\pgfpathlineto{\pgfqpoint{1.072755in}{0.968084in}}%
\pgfpathlineto{\pgfqpoint{1.080738in}{0.977224in}}%
\pgfpathlineto{\pgfqpoint{1.088699in}{0.986799in}}%
\pgfpathlineto{\pgfqpoint{1.072705in}{0.995427in}}%
\pgfpathlineto{\pgfqpoint{1.056738in}{1.004849in}}%
\pgfpathlineto{\pgfqpoint{1.048750in}{0.995467in}}%
\pgfpathlineto{\pgfqpoint{1.040740in}{0.986511in}}%
\pgfpathclose%
\pgfusepath{fill}%
\end{pgfscope}%
\begin{pgfscope}%
\pgfpathrectangle{\pgfqpoint{0.407592in}{0.222904in}}{\pgfqpoint{1.411163in}{1.411163in}}%
\pgfusepath{clip}%
\pgfsetbuttcap%
\pgfsetroundjoin%
\definecolor{currentfill}{rgb}{0.909804,0.698039,0.423529}%
\pgfsetfillcolor{currentfill}%
\pgfsetfillopacity{0.120000}%
\pgfsetlinewidth{0.000000pt}%
\definecolor{currentstroke}{rgb}{0.000000,0.000000,0.000000}%
\pgfsetstrokecolor{currentstroke}%
\pgfsetdash{}{0pt}%
\pgfpathmoveto{\pgfqpoint{1.280284in}{1.041495in}}%
\pgfpathlineto{\pgfqpoint{1.296561in}{1.040631in}}%
\pgfpathlineto{\pgfqpoint{1.312891in}{1.040522in}}%
\pgfpathlineto{\pgfqpoint{1.320604in}{1.053217in}}%
\pgfpathlineto{\pgfqpoint{1.328303in}{1.066145in}}%
\pgfpathlineto{\pgfqpoint{1.311989in}{1.066375in}}%
\pgfpathlineto{\pgfqpoint{1.295726in}{1.067329in}}%
\pgfpathlineto{\pgfqpoint{1.288012in}{1.054297in}}%
\pgfpathlineto{\pgfqpoint{1.280284in}{1.041495in}}%
\pgfpathclose%
\pgfusepath{fill}%
\end{pgfscope}%
\begin{pgfscope}%
\pgfpathrectangle{\pgfqpoint{0.407592in}{0.222904in}}{\pgfqpoint{1.411163in}{1.411163in}}%
\pgfusepath{clip}%
\pgfsetbuttcap%
\pgfsetroundjoin%
\definecolor{currentfill}{rgb}{0.909804,0.698039,0.423529}%
\pgfsetfillcolor{currentfill}%
\pgfsetfillopacity{0.120000}%
\pgfsetlinewidth{0.000000pt}%
\definecolor{currentstroke}{rgb}{0.000000,0.000000,0.000000}%
\pgfsetstrokecolor{currentstroke}%
\pgfsetdash{}{0pt}%
\pgfpathmoveto{\pgfqpoint{0.688058in}{1.204219in}}%
\pgfpathlineto{\pgfqpoint{0.703946in}{1.186469in}}%
\pgfpathlineto{\pgfqpoint{0.719843in}{1.168894in}}%
\pgfpathlineto{\pgfqpoint{0.728375in}{1.173498in}}%
\pgfpathlineto{\pgfqpoint{0.736874in}{1.178432in}}%
\pgfpathlineto{\pgfqpoint{0.721018in}{1.195940in}}%
\pgfpathlineto{\pgfqpoint{0.705171in}{1.213622in}}%
\pgfpathlineto{\pgfqpoint{0.696632in}{1.208765in}}%
\pgfpathlineto{\pgfqpoint{0.688058in}{1.204219in}}%
\pgfpathclose%
\pgfusepath{fill}%
\end{pgfscope}%
\begin{pgfscope}%
\pgfpathrectangle{\pgfqpoint{0.407592in}{0.222904in}}{\pgfqpoint{1.411163in}{1.411163in}}%
\pgfusepath{clip}%
\pgfsetbuttcap%
\pgfsetroundjoin%
\definecolor{currentfill}{rgb}{0.909804,0.698039,0.423529}%
\pgfsetfillcolor{currentfill}%
\pgfsetfillopacity{0.120000}%
\pgfsetlinewidth{0.000000pt}%
\definecolor{currentstroke}{rgb}{0.000000,0.000000,0.000000}%
\pgfsetstrokecolor{currentstroke}%
\pgfsetdash{}{0pt}%
\pgfpathmoveto{\pgfqpoint{0.912535in}{1.029605in}}%
\pgfpathlineto{\pgfqpoint{0.928480in}{1.015719in}}%
\pgfpathlineto{\pgfqpoint{0.944440in}{1.002367in}}%
\pgfpathlineto{\pgfqpoint{0.952616in}{1.009161in}}%
\pgfpathlineto{\pgfqpoint{0.960764in}{1.016402in}}%
\pgfpathlineto{\pgfqpoint{0.944839in}{1.029589in}}%
\pgfpathlineto{\pgfqpoint{0.928930in}{1.043308in}}%
\pgfpathlineto{\pgfqpoint{0.920746in}{1.036245in}}%
\pgfpathlineto{\pgfqpoint{0.912535in}{1.029605in}}%
\pgfpathclose%
\pgfusepath{fill}%
\end{pgfscope}%
\begin{pgfscope}%
\pgfpathrectangle{\pgfqpoint{0.407592in}{0.222904in}}{\pgfqpoint{1.411163in}{1.411163in}}%
\pgfusepath{clip}%
\pgfsetbuttcap%
\pgfsetroundjoin%
\definecolor{currentfill}{rgb}{0.909804,0.698039,0.423529}%
\pgfsetfillcolor{currentfill}%
\pgfsetfillopacity{0.120000}%
\pgfsetlinewidth{0.000000pt}%
\definecolor{currentstroke}{rgb}{0.000000,0.000000,0.000000}%
\pgfsetstrokecolor{currentstroke}%
\pgfsetdash{}{0pt}%
\pgfpathmoveto{\pgfqpoint{1.328303in}{1.066145in}}%
\pgfpathlineto{\pgfqpoint{1.344672in}{1.066624in}}%
\pgfpathlineto{\pgfqpoint{1.361097in}{1.067790in}}%
\pgfpathlineto{\pgfqpoint{1.368768in}{1.080778in}}%
\pgfpathlineto{\pgfqpoint{1.376425in}{1.093971in}}%
\pgfpathlineto{\pgfqpoint{1.360015in}{1.092973in}}%
\pgfpathlineto{\pgfqpoint{1.343660in}{1.092638in}}%
\pgfpathlineto{\pgfqpoint{1.335989in}{1.079291in}}%
\pgfpathlineto{\pgfqpoint{1.328303in}{1.066145in}}%
\pgfpathclose%
\pgfusepath{fill}%
\end{pgfscope}%
\begin{pgfscope}%
\pgfpathrectangle{\pgfqpoint{0.407592in}{0.222904in}}{\pgfqpoint{1.411163in}{1.411163in}}%
\pgfusepath{clip}%
\pgfsetbuttcap%
\pgfsetroundjoin%
\definecolor{currentfill}{rgb}{0.909804,0.698039,0.423529}%
\pgfsetfillcolor{currentfill}%
\pgfsetfillopacity{0.120000}%
\pgfsetlinewidth{0.000000pt}%
\definecolor{currentstroke}{rgb}{0.000000,0.000000,0.000000}%
\pgfsetstrokecolor{currentstroke}%
\pgfsetdash{}{0pt}%
\pgfpathmoveto{\pgfqpoint{0.992667in}{0.991816in}}%
\pgfpathlineto{\pgfqpoint{1.008647in}{0.980516in}}%
\pgfpathlineto{\pgfqpoint{1.024648in}{0.969945in}}%
\pgfpathlineto{\pgfqpoint{1.032706in}{0.977997in}}%
\pgfpathlineto{\pgfqpoint{1.040740in}{0.986511in}}%
\pgfpathlineto{\pgfqpoint{1.024769in}{0.996898in}}%
\pgfpathlineto{\pgfqpoint{1.008820in}{1.008004in}}%
\pgfpathlineto{\pgfqpoint{1.000756in}{0.999688in}}%
\pgfpathlineto{\pgfqpoint{0.992667in}{0.991816in}}%
\pgfpathclose%
\pgfusepath{fill}%
\end{pgfscope}%
\begin{pgfscope}%
\pgfpathrectangle{\pgfqpoint{0.407592in}{0.222904in}}{\pgfqpoint{1.411163in}{1.411163in}}%
\pgfusepath{clip}%
\pgfsetbuttcap%
\pgfsetroundjoin%
\definecolor{currentfill}{rgb}{0.909804,0.698039,0.423529}%
\pgfsetfillcolor{currentfill}%
\pgfsetfillopacity{0.120000}%
\pgfsetlinewidth{0.000000pt}%
\definecolor{currentstroke}{rgb}{0.000000,0.000000,0.000000}%
\pgfsetstrokecolor{currentstroke}%
\pgfsetdash{}{0pt}%
\pgfpathmoveto{\pgfqpoint{0.719843in}{1.168894in}}%
\pgfpathlineto{\pgfqpoint{0.735750in}{1.151509in}}%
\pgfpathlineto{\pgfqpoint{0.751666in}{1.134334in}}%
\pgfpathlineto{\pgfqpoint{0.760158in}{1.138995in}}%
\pgfpathlineto{\pgfqpoint{0.768617in}{1.144006in}}%
\pgfpathlineto{\pgfqpoint{0.752741in}{1.161114in}}%
\pgfpathlineto{\pgfqpoint{0.736874in}{1.178432in}}%
\pgfpathlineto{\pgfqpoint{0.728375in}{1.173498in}}%
\pgfpathlineto{\pgfqpoint{0.719843in}{1.168894in}}%
\pgfpathclose%
\pgfusepath{fill}%
\end{pgfscope}%
\begin{pgfscope}%
\pgfpathrectangle{\pgfqpoint{0.407592in}{0.222904in}}{\pgfqpoint{1.411163in}{1.411163in}}%
\pgfusepath{clip}%
\pgfsetbuttcap%
\pgfsetroundjoin%
\definecolor{currentfill}{rgb}{0.909804,0.698039,0.423529}%
\pgfsetfillcolor{currentfill}%
\pgfsetfillopacity{0.120000}%
\pgfsetlinewidth{0.000000pt}%
\definecolor{currentstroke}{rgb}{0.000000,0.000000,0.000000}%
\pgfsetstrokecolor{currentstroke}%
\pgfsetdash{}{0pt}%
\pgfpathmoveto{\pgfqpoint{0.832227in}{1.078106in}}%
\pgfpathlineto{\pgfqpoint{0.848158in}{1.062396in}}%
\pgfpathlineto{\pgfqpoint{0.864100in}{1.047056in}}%
\pgfpathlineto{\pgfqpoint{0.872409in}{1.052719in}}%
\pgfpathlineto{\pgfqpoint{0.880687in}{1.058797in}}%
\pgfpathlineto{\pgfqpoint{0.864783in}{1.074021in}}%
\pgfpathlineto{\pgfqpoint{0.848891in}{1.089615in}}%
\pgfpathlineto{\pgfqpoint{0.840575in}{1.083666in}}%
\pgfpathlineto{\pgfqpoint{0.832227in}{1.078106in}}%
\pgfpathclose%
\pgfusepath{fill}%
\end{pgfscope}%
\begin{pgfscope}%
\pgfpathrectangle{\pgfqpoint{0.407592in}{0.222904in}}{\pgfqpoint{1.411163in}{1.411163in}}%
\pgfusepath{clip}%
\pgfsetbuttcap%
\pgfsetroundjoin%
\definecolor{currentfill}{rgb}{0.909804,0.698039,0.423529}%
\pgfsetfillcolor{currentfill}%
\pgfsetfillopacity{0.120000}%
\pgfsetlinewidth{0.000000pt}%
\definecolor{currentstroke}{rgb}{0.000000,0.000000,0.000000}%
\pgfsetstrokecolor{currentstroke}%
\pgfsetdash{}{0pt}%
\pgfpathmoveto{\pgfqpoint{1.439860in}{1.151675in}}%
\pgfpathlineto{\pgfqpoint{1.456418in}{1.154106in}}%
\pgfpathlineto{\pgfqpoint{1.473037in}{1.157082in}}%
\pgfpathlineto{\pgfqpoint{1.480601in}{1.170718in}}%
\pgfpathlineto{\pgfqpoint{1.463989in}{1.167851in}}%
\pgfpathlineto{\pgfqpoint{1.447439in}{1.165521in}}%
\pgfusepath{fill}%
\end{pgfscope}%
\begin{pgfscope}%
\pgfpathrectangle{\pgfqpoint{0.407592in}{0.222904in}}{\pgfqpoint{1.411163in}{1.411163in}}%
\pgfusepath{clip}%
\pgfsetbuttcap%
\pgfsetroundjoin%
\definecolor{currentfill}{rgb}{0.909804,0.698039,0.423529}%
\pgfsetfillcolor{currentfill}%
\pgfsetfillopacity{0.120000}%
\pgfsetlinewidth{0.000000pt}%
\definecolor{currentstroke}{rgb}{0.000000,0.000000,0.000000}%
\pgfsetstrokecolor{currentstroke}%
\pgfsetdash{}{0pt}%
\pgfpathmoveto{\pgfqpoint{1.168761in}{0.981802in}}%
\pgfpathlineto{\pgfqpoint{1.184904in}{0.977661in}}%
\pgfpathlineto{\pgfqpoint{1.201089in}{0.974421in}}%
\pgfpathlineto{\pgfqpoint{1.208928in}{0.985465in}}%
\pgfpathlineto{\pgfqpoint{1.216750in}{0.996863in}}%
\pgfpathlineto{\pgfqpoint{1.200583in}{1.000093in}}%
\pgfpathlineto{\pgfqpoint{1.184459in}{1.004188in}}%
\pgfpathlineto{\pgfqpoint{1.176619in}{0.992818in}}%
\pgfpathlineto{\pgfqpoint{1.168761in}{0.981802in}}%
\pgfpathclose%
\pgfusepath{fill}%
\end{pgfscope}%
\begin{pgfscope}%
\pgfpathrectangle{\pgfqpoint{0.407592in}{0.222904in}}{\pgfqpoint{1.411163in}{1.411163in}}%
\pgfusepath{clip}%
\pgfsetbuttcap%
\pgfsetroundjoin%
\definecolor{currentfill}{rgb}{0.909804,0.698039,0.423529}%
\pgfsetfillcolor{currentfill}%
\pgfsetfillopacity{0.120000}%
\pgfsetlinewidth{0.000000pt}%
\definecolor{currentstroke}{rgb}{0.000000,0.000000,0.000000}%
\pgfsetstrokecolor{currentstroke}%
\pgfsetdash{}{0pt}%
\pgfpathmoveto{\pgfqpoint{1.120777in}{0.972075in}}%
\pgfpathlineto{\pgfqpoint{1.136865in}{0.966042in}}%
\pgfpathlineto{\pgfqpoint{1.152990in}{0.960926in}}%
\pgfpathlineto{\pgfqpoint{1.160885in}{0.971164in}}%
\pgfpathlineto{\pgfqpoint{1.168761in}{0.981802in}}%
\pgfpathlineto{\pgfqpoint{1.152657in}{0.986840in}}%
\pgfpathlineto{\pgfqpoint{1.136591in}{0.992761in}}%
\pgfpathlineto{\pgfqpoint{1.128694in}{0.982219in}}%
\pgfpathlineto{\pgfqpoint{1.120777in}{0.972075in}}%
\pgfpathclose%
\pgfusepath{fill}%
\end{pgfscope}%
\begin{pgfscope}%
\pgfpathrectangle{\pgfqpoint{0.407592in}{0.222904in}}{\pgfqpoint{1.411163in}{1.411163in}}%
\pgfusepath{clip}%
\pgfsetbuttcap%
\pgfsetroundjoin%
\definecolor{currentfill}{rgb}{0.909804,0.698039,0.423529}%
\pgfsetfillcolor{currentfill}%
\pgfsetfillopacity{0.120000}%
\pgfsetlinewidth{0.000000pt}%
\definecolor{currentstroke}{rgb}{0.000000,0.000000,0.000000}%
\pgfsetstrokecolor{currentstroke}%
\pgfsetdash{}{0pt}%
\pgfpathmoveto{\pgfqpoint{1.376425in}{1.093971in}}%
\pgfpathlineto{\pgfqpoint{1.392893in}{1.095612in}}%
\pgfpathlineto{\pgfqpoint{1.409420in}{1.097874in}}%
\pgfpathlineto{\pgfqpoint{1.417049in}{1.111064in}}%
\pgfpathlineto{\pgfqpoint{1.424666in}{1.124435in}}%
\pgfpathlineto{\pgfqpoint{1.408154in}{1.122375in}}%
\pgfpathlineto{\pgfqpoint{1.391700in}{1.120917in}}%
\pgfpathlineto{\pgfqpoint{1.384069in}{1.107355in}}%
\pgfpathlineto{\pgfqpoint{1.376425in}{1.093971in}}%
\pgfpathclose%
\pgfusepath{fill}%
\end{pgfscope}%
\begin{pgfscope}%
\pgfpathrectangle{\pgfqpoint{0.407592in}{0.222904in}}{\pgfqpoint{1.411163in}{1.411163in}}%
\pgfusepath{clip}%
\pgfsetbuttcap%
\pgfsetroundjoin%
\definecolor{currentfill}{rgb}{0.909804,0.698039,0.423529}%
\pgfsetfillcolor{currentfill}%
\pgfsetfillopacity{0.120000}%
\pgfsetlinewidth{0.000000pt}%
\definecolor{currentstroke}{rgb}{0.000000,0.000000,0.000000}%
\pgfsetstrokecolor{currentstroke}%
\pgfsetdash{}{0pt}%
\pgfpathmoveto{\pgfqpoint{1.216750in}{0.996863in}}%
\pgfpathlineto{\pgfqpoint{1.232962in}{0.994494in}}%
\pgfpathlineto{\pgfqpoint{1.249222in}{0.992980in}}%
\pgfpathlineto{\pgfqpoint{1.257010in}{1.004665in}}%
\pgfpathlineto{\pgfqpoint{1.264783in}{1.016660in}}%
\pgfpathlineto{\pgfqpoint{1.248540in}{1.018236in}}%
\pgfpathlineto{\pgfqpoint{1.232345in}{1.020630in}}%
\pgfpathlineto{\pgfqpoint{1.224556in}{1.008592in}}%
\pgfpathlineto{\pgfqpoint{1.216750in}{0.996863in}}%
\pgfpathclose%
\pgfusepath{fill}%
\end{pgfscope}%
\begin{pgfscope}%
\pgfpathrectangle{\pgfqpoint{0.407592in}{0.222904in}}{\pgfqpoint{1.411163in}{1.411163in}}%
\pgfusepath{clip}%
\pgfsetbuttcap%
\pgfsetroundjoin%
\definecolor{currentfill}{rgb}{0.909804,0.698039,0.423529}%
\pgfsetfillcolor{currentfill}%
\pgfsetfillopacity{0.120000}%
\pgfsetlinewidth{0.000000pt}%
\definecolor{currentstroke}{rgb}{0.000000,0.000000,0.000000}%
\pgfsetstrokecolor{currentstroke}%
\pgfsetdash{}{0pt}%
\pgfpathmoveto{\pgfqpoint{1.072755in}{0.968084in}}%
\pgfpathlineto{\pgfqpoint{1.088804in}{0.960132in}}%
\pgfpathlineto{\pgfqpoint{1.104883in}{0.953073in}}%
\pgfpathlineto{\pgfqpoint{1.112840in}{0.962352in}}%
\pgfpathlineto{\pgfqpoint{1.120777in}{0.972075in}}%
\pgfpathlineto{\pgfqpoint{1.104722in}{0.979004in}}%
\pgfpathlineto{\pgfqpoint{1.088699in}{0.986799in}}%
\pgfpathlineto{\pgfqpoint{1.080738in}{0.977224in}}%
\pgfpathlineto{\pgfqpoint{1.072755in}{0.968084in}}%
\pgfpathclose%
\pgfusepath{fill}%
\end{pgfscope}%
\begin{pgfscope}%
\pgfpathrectangle{\pgfqpoint{0.407592in}{0.222904in}}{\pgfqpoint{1.411163in}{1.411163in}}%
\pgfusepath{clip}%
\pgfsetbuttcap%
\pgfsetroundjoin%
\definecolor{currentfill}{rgb}{0.909804,0.698039,0.423529}%
\pgfsetfillcolor{currentfill}%
\pgfsetfillopacity{0.120000}%
\pgfsetlinewidth{0.000000pt}%
\definecolor{currentstroke}{rgb}{0.000000,0.000000,0.000000}%
\pgfsetstrokecolor{currentstroke}%
\pgfsetdash{}{0pt}%
\pgfpathmoveto{\pgfqpoint{0.751666in}{1.134334in}}%
\pgfpathlineto{\pgfqpoint{0.767593in}{1.117388in}}%
\pgfpathlineto{\pgfqpoint{0.783531in}{1.100694in}}%
\pgfpathlineto{\pgfqpoint{0.791982in}{1.105411in}}%
\pgfpathlineto{\pgfqpoint{0.800401in}{1.110500in}}%
\pgfpathlineto{\pgfqpoint{0.784503in}{1.127127in}}%
\pgfpathlineto{\pgfqpoint{0.768617in}{1.144006in}}%
\pgfpathlineto{\pgfqpoint{0.760158in}{1.138995in}}%
\pgfpathlineto{\pgfqpoint{0.751666in}{1.134334in}}%
\pgfpathclose%
\pgfusepath{fill}%
\end{pgfscope}%
\begin{pgfscope}%
\pgfpathrectangle{\pgfqpoint{0.407592in}{0.222904in}}{\pgfqpoint{1.411163in}{1.411163in}}%
\pgfusepath{clip}%
\pgfsetbuttcap%
\pgfsetroundjoin%
\definecolor{currentfill}{rgb}{0.909804,0.698039,0.423529}%
\pgfsetfillcolor{currentfill}%
\pgfsetfillopacity{0.120000}%
\pgfsetlinewidth{0.000000pt}%
\definecolor{currentstroke}{rgb}{0.000000,0.000000,0.000000}%
\pgfsetstrokecolor{currentstroke}%
\pgfsetdash{}{0pt}%
\pgfpathmoveto{\pgfqpoint{0.944440in}{1.002367in}}%
\pgfpathlineto{\pgfqpoint{0.960416in}{0.989596in}}%
\pgfpathlineto{\pgfqpoint{0.976410in}{0.977458in}}%
\pgfpathlineto{\pgfqpoint{0.984552in}{0.984402in}}%
\pgfpathlineto{\pgfqpoint{0.992667in}{0.991816in}}%
\pgfpathlineto{\pgfqpoint{0.976707in}{1.003795in}}%
\pgfpathlineto{\pgfqpoint{0.960764in}{1.016402in}}%
\pgfpathlineto{\pgfqpoint{0.952616in}{1.009161in}}%
\pgfpathlineto{\pgfqpoint{0.944440in}{1.002367in}}%
\pgfpathclose%
\pgfusepath{fill}%
\end{pgfscope}%
\begin{pgfscope}%
\pgfpathrectangle{\pgfqpoint{0.407592in}{0.222904in}}{\pgfqpoint{1.411163in}{1.411163in}}%
\pgfusepath{clip}%
\pgfsetbuttcap%
\pgfsetroundjoin%
\definecolor{currentfill}{rgb}{0.909804,0.698039,0.423529}%
\pgfsetfillcolor{currentfill}%
\pgfsetfillopacity{0.120000}%
\pgfsetlinewidth{0.000000pt}%
\definecolor{currentstroke}{rgb}{0.000000,0.000000,0.000000}%
\pgfsetstrokecolor{currentstroke}%
\pgfsetdash{}{0pt}%
\pgfpathmoveto{\pgfqpoint{1.264783in}{1.016660in}}%
\pgfpathlineto{\pgfqpoint{1.281076in}{1.015889in}}%
\pgfpathlineto{\pgfqpoint{1.297422in}{1.015905in}}%
\pgfpathlineto{\pgfqpoint{1.305163in}{1.028078in}}%
\pgfpathlineto{\pgfqpoint{1.312891in}{1.040522in}}%
\pgfpathlineto{\pgfqpoint{1.296561in}{1.040631in}}%
\pgfpathlineto{\pgfqpoint{1.280284in}{1.041495in}}%
\pgfpathlineto{\pgfqpoint{1.272541in}{1.028944in}}%
\pgfpathlineto{\pgfqpoint{1.264783in}{1.016660in}}%
\pgfpathclose%
\pgfusepath{fill}%
\end{pgfscope}%
\begin{pgfscope}%
\pgfpathrectangle{\pgfqpoint{0.407592in}{0.222904in}}{\pgfqpoint{1.411163in}{1.411163in}}%
\pgfusepath{clip}%
\pgfsetbuttcap%
\pgfsetroundjoin%
\definecolor{currentfill}{rgb}{0.909804,0.698039,0.423529}%
\pgfsetfillcolor{currentfill}%
\pgfsetfillopacity{0.120000}%
\pgfsetlinewidth{0.000000pt}%
\definecolor{currentstroke}{rgb}{0.000000,0.000000,0.000000}%
\pgfsetstrokecolor{currentstroke}%
\pgfsetdash{}{0pt}%
\pgfpathmoveto{\pgfqpoint{0.864100in}{1.047056in}}%
\pgfpathlineto{\pgfqpoint{0.880055in}{1.032120in}}%
\pgfpathlineto{\pgfqpoint{0.896023in}{1.017630in}}%
\pgfpathlineto{\pgfqpoint{0.904294in}{1.023398in}}%
\pgfpathlineto{\pgfqpoint{0.912535in}{1.029605in}}%
\pgfpathlineto{\pgfqpoint{0.896604in}{1.043978in}}%
\pgfpathlineto{\pgfqpoint{0.880687in}{1.058797in}}%
\pgfpathlineto{\pgfqpoint{0.872409in}{1.052719in}}%
\pgfpathlineto{\pgfqpoint{0.864100in}{1.047056in}}%
\pgfpathclose%
\pgfusepath{fill}%
\end{pgfscope}%
\begin{pgfscope}%
\pgfpathrectangle{\pgfqpoint{0.407592in}{0.222904in}}{\pgfqpoint{1.411163in}{1.411163in}}%
\pgfusepath{clip}%
\pgfsetbuttcap%
\pgfsetroundjoin%
\definecolor{currentfill}{rgb}{0.909804,0.698039,0.423529}%
\pgfsetfillcolor{currentfill}%
\pgfsetfillopacity{0.120000}%
\pgfsetlinewidth{0.000000pt}%
\definecolor{currentstroke}{rgb}{0.000000,0.000000,0.000000}%
\pgfsetstrokecolor{currentstroke}%
\pgfsetdash{}{0pt}%
\pgfpathmoveto{\pgfqpoint{0.638978in}{1.232091in}}%
\pgfpathlineto{\pgfqpoint{0.654888in}{1.214004in}}%
\pgfpathlineto{\pgfqpoint{0.670808in}{1.196064in}}%
\pgfpathlineto{\pgfqpoint{0.679450in}{1.199985in}}%
\pgfpathlineto{\pgfqpoint{0.688058in}{1.204219in}}%
\pgfpathlineto{\pgfqpoint{0.672180in}{1.222126in}}%
\pgfpathlineto{\pgfqpoint{0.656311in}{1.240179in}}%
\pgfpathlineto{\pgfqpoint{0.647662in}{1.235987in}}%
\pgfpathlineto{\pgfqpoint{0.638978in}{1.232091in}}%
\pgfpathclose%
\pgfusepath{fill}%
\end{pgfscope}%
\begin{pgfscope}%
\pgfpathrectangle{\pgfqpoint{0.407592in}{0.222904in}}{\pgfqpoint{1.411163in}{1.411163in}}%
\pgfusepath{clip}%
\pgfsetbuttcap%
\pgfsetroundjoin%
\definecolor{currentfill}{rgb}{0.909804,0.698039,0.423529}%
\pgfsetfillcolor{currentfill}%
\pgfsetfillopacity{0.120000}%
\pgfsetlinewidth{0.000000pt}%
\definecolor{currentstroke}{rgb}{0.000000,0.000000,0.000000}%
\pgfsetstrokecolor{currentstroke}%
\pgfsetdash{}{0pt}%
\pgfpathmoveto{\pgfqpoint{1.024648in}{0.969945in}}%
\pgfpathlineto{\pgfqpoint{1.040672in}{0.960155in}}%
\pgfpathlineto{\pgfqpoint{1.056720in}{0.951194in}}%
\pgfpathlineto{\pgfqpoint{1.064749in}{0.959401in}}%
\pgfpathlineto{\pgfqpoint{1.072755in}{0.968084in}}%
\pgfpathlineto{\pgfqpoint{1.056734in}{0.976891in}}%
\pgfpathlineto{\pgfqpoint{1.040740in}{0.986511in}}%
\pgfpathlineto{\pgfqpoint{1.032706in}{0.977997in}}%
\pgfpathlineto{\pgfqpoint{1.024648in}{0.969945in}}%
\pgfpathclose%
\pgfusepath{fill}%
\end{pgfscope}%
\begin{pgfscope}%
\pgfpathrectangle{\pgfqpoint{0.407592in}{0.222904in}}{\pgfqpoint{1.411163in}{1.411163in}}%
\pgfusepath{clip}%
\pgfsetbuttcap%
\pgfsetroundjoin%
\definecolor{currentfill}{rgb}{0.909804,0.698039,0.423529}%
\pgfsetfillcolor{currentfill}%
\pgfsetfillopacity{0.120000}%
\pgfsetlinewidth{0.000000pt}%
\definecolor{currentstroke}{rgb}{0.000000,0.000000,0.000000}%
\pgfsetstrokecolor{currentstroke}%
\pgfsetdash{}{0pt}%
\pgfpathmoveto{\pgfqpoint{1.424666in}{1.124435in}}%
\pgfpathlineto{\pgfqpoint{1.441239in}{1.127076in}}%
\pgfpathlineto{\pgfqpoint{1.457873in}{1.130277in}}%
\pgfpathlineto{\pgfqpoint{1.465462in}{1.143598in}}%
\pgfpathlineto{\pgfqpoint{1.473037in}{1.157082in}}%
\pgfpathlineto{\pgfqpoint{1.456418in}{1.154106in}}%
\pgfpathlineto{\pgfqpoint{1.439860in}{1.151675in}}%
\pgfpathlineto{\pgfqpoint{1.432269in}{1.137976in}}%
\pgfpathlineto{\pgfqpoint{1.424666in}{1.124435in}}%
\pgfpathclose%
\pgfusepath{fill}%
\end{pgfscope}%
\begin{pgfscope}%
\pgfpathrectangle{\pgfqpoint{0.407592in}{0.222904in}}{\pgfqpoint{1.411163in}{1.411163in}}%
\pgfusepath{clip}%
\pgfsetbuttcap%
\pgfsetroundjoin%
\definecolor{currentfill}{rgb}{0.909804,0.698039,0.423529}%
\pgfsetfillcolor{currentfill}%
\pgfsetfillopacity{0.120000}%
\pgfsetlinewidth{0.000000pt}%
\definecolor{currentstroke}{rgb}{0.000000,0.000000,0.000000}%
\pgfsetstrokecolor{currentstroke}%
\pgfsetdash{}{0pt}%
\pgfpathmoveto{\pgfqpoint{1.312891in}{1.040522in}}%
\pgfpathlineto{\pgfqpoint{1.329275in}{1.041149in}}%
\pgfpathlineto{\pgfqpoint{1.345716in}{1.042490in}}%
\pgfpathlineto{\pgfqpoint{1.353413in}{1.055022in}}%
\pgfpathlineto{\pgfqpoint{1.361097in}{1.067790in}}%
\pgfpathlineto{\pgfqpoint{1.344672in}{1.066624in}}%
\pgfpathlineto{\pgfqpoint{1.328303in}{1.066145in}}%
\pgfpathlineto{\pgfqpoint{1.320604in}{1.053217in}}%
\pgfpathlineto{\pgfqpoint{1.312891in}{1.040522in}}%
\pgfpathclose%
\pgfusepath{fill}%
\end{pgfscope}%
\begin{pgfscope}%
\pgfpathrectangle{\pgfqpoint{0.407592in}{0.222904in}}{\pgfqpoint{1.411163in}{1.411163in}}%
\pgfusepath{clip}%
\pgfsetbuttcap%
\pgfsetroundjoin%
\definecolor{currentfill}{rgb}{0.909804,0.698039,0.423529}%
\pgfsetfillcolor{currentfill}%
\pgfsetfillopacity{0.120000}%
\pgfsetlinewidth{0.000000pt}%
\definecolor{currentstroke}{rgb}{0.000000,0.000000,0.000000}%
\pgfsetstrokecolor{currentstroke}%
\pgfsetdash{}{0pt}%
\pgfpathmoveto{\pgfqpoint{0.783531in}{1.100694in}}%
\pgfpathlineto{\pgfqpoint{0.799478in}{1.084278in}}%
\pgfpathlineto{\pgfqpoint{0.815437in}{1.068166in}}%
\pgfpathlineto{\pgfqpoint{0.823848in}{1.072938in}}%
\pgfpathlineto{\pgfqpoint{0.832227in}{1.078106in}}%
\pgfpathlineto{\pgfqpoint{0.816308in}{1.094150in}}%
\pgfpathlineto{\pgfqpoint{0.800401in}{1.110500in}}%
\pgfpathlineto{\pgfqpoint{0.791982in}{1.105411in}}%
\pgfpathlineto{\pgfqpoint{0.783531in}{1.100694in}}%
\pgfpathclose%
\pgfusepath{fill}%
\end{pgfscope}%
\begin{pgfscope}%
\pgfpathrectangle{\pgfqpoint{0.407592in}{0.222904in}}{\pgfqpoint{1.411163in}{1.411163in}}%
\pgfusepath{clip}%
\pgfsetbuttcap%
\pgfsetroundjoin%
\definecolor{currentfill}{rgb}{0.909804,0.698039,0.423529}%
\pgfsetfillcolor{currentfill}%
\pgfsetfillopacity{0.120000}%
\pgfsetlinewidth{0.000000pt}%
\definecolor{currentstroke}{rgb}{0.000000,0.000000,0.000000}%
\pgfsetstrokecolor{currentstroke}%
\pgfsetdash{}{0pt}%
\pgfpathmoveto{\pgfqpoint{0.670808in}{1.196064in}}%
\pgfpathlineto{\pgfqpoint{0.686737in}{1.178284in}}%
\pgfpathlineto{\pgfqpoint{0.702675in}{1.160679in}}%
\pgfpathlineto{\pgfqpoint{0.711276in}{1.164621in}}%
\pgfpathlineto{\pgfqpoint{0.719843in}{1.168894in}}%
\pgfpathlineto{\pgfqpoint{0.703946in}{1.186469in}}%
\pgfpathlineto{\pgfqpoint{0.688058in}{1.204219in}}%
\pgfpathlineto{\pgfqpoint{0.679450in}{1.199985in}}%
\pgfpathlineto{\pgfqpoint{0.670808in}{1.196064in}}%
\pgfpathclose%
\pgfusepath{fill}%
\end{pgfscope}%
\begin{pgfscope}%
\pgfpathrectangle{\pgfqpoint{0.407592in}{0.222904in}}{\pgfqpoint{1.411163in}{1.411163in}}%
\pgfusepath{clip}%
\pgfsetbuttcap%
\pgfsetroundjoin%
\definecolor{currentfill}{rgb}{0.909804,0.698039,0.423529}%
\pgfsetfillcolor{currentfill}%
\pgfsetfillopacity{0.120000}%
\pgfsetlinewidth{0.000000pt}%
\definecolor{currentstroke}{rgb}{0.000000,0.000000,0.000000}%
\pgfsetstrokecolor{currentstroke}%
\pgfsetdash{}{0pt}%
\pgfpathmoveto{\pgfqpoint{0.896023in}{1.017630in}}%
\pgfpathlineto{\pgfqpoint{0.912005in}{1.003627in}}%
\pgfpathlineto{\pgfqpoint{0.928001in}{0.990159in}}%
\pgfpathlineto{\pgfqpoint{0.936235in}{0.996031in}}%
\pgfpathlineto{\pgfqpoint{0.944440in}{1.002367in}}%
\pgfpathlineto{\pgfqpoint{0.928480in}{1.015719in}}%
\pgfpathlineto{\pgfqpoint{0.912535in}{1.029605in}}%
\pgfpathlineto{\pgfqpoint{0.904294in}{1.023398in}}%
\pgfpathlineto{\pgfqpoint{0.896023in}{1.017630in}}%
\pgfpathclose%
\pgfusepath{fill}%
\end{pgfscope}%
\begin{pgfscope}%
\pgfpathrectangle{\pgfqpoint{0.407592in}{0.222904in}}{\pgfqpoint{1.411163in}{1.411163in}}%
\pgfusepath{clip}%
\pgfsetbuttcap%
\pgfsetroundjoin%
\definecolor{currentfill}{rgb}{0.909804,0.698039,0.423529}%
\pgfsetfillcolor{currentfill}%
\pgfsetfillopacity{0.120000}%
\pgfsetlinewidth{0.000000pt}%
\definecolor{currentstroke}{rgb}{0.000000,0.000000,0.000000}%
\pgfsetstrokecolor{currentstroke}%
\pgfsetdash{}{0pt}%
\pgfpathmoveto{\pgfqpoint{0.976410in}{0.977458in}}%
\pgfpathlineto{\pgfqpoint{0.992422in}{0.966004in}}%
\pgfpathlineto{\pgfqpoint{1.008455in}{0.955288in}}%
\pgfpathlineto{\pgfqpoint{1.016564in}{0.962371in}}%
\pgfpathlineto{\pgfqpoint{1.024648in}{0.969945in}}%
\pgfpathlineto{\pgfqpoint{1.008647in}{0.980516in}}%
\pgfpathlineto{\pgfqpoint{0.992667in}{0.991816in}}%
\pgfpathlineto{\pgfqpoint{0.984552in}{0.984402in}}%
\pgfpathlineto{\pgfqpoint{0.976410in}{0.977458in}}%
\pgfpathclose%
\pgfusepath{fill}%
\end{pgfscope}%
\begin{pgfscope}%
\pgfpathrectangle{\pgfqpoint{0.407592in}{0.222904in}}{\pgfqpoint{1.411163in}{1.411163in}}%
\pgfusepath{clip}%
\pgfsetbuttcap%
\pgfsetroundjoin%
\definecolor{currentfill}{rgb}{0.909804,0.698039,0.423529}%
\pgfsetfillcolor{currentfill}%
\pgfsetfillopacity{0.120000}%
\pgfsetlinewidth{0.000000pt}%
\definecolor{currentstroke}{rgb}{0.000000,0.000000,0.000000}%
\pgfsetstrokecolor{currentstroke}%
\pgfsetdash{}{0pt}%
\pgfpathmoveto{\pgfqpoint{1.361097in}{1.067790in}}%
\pgfpathlineto{\pgfqpoint{1.377580in}{1.069621in}}%
\pgfpathlineto{\pgfqpoint{1.394123in}{1.072094in}}%
\pgfpathlineto{\pgfqpoint{1.401778in}{1.084880in}}%
\pgfpathlineto{\pgfqpoint{1.409420in}{1.097874in}}%
\pgfpathlineto{\pgfqpoint{1.392893in}{1.095612in}}%
\pgfpathlineto{\pgfqpoint{1.376425in}{1.093971in}}%
\pgfpathlineto{\pgfqpoint{1.368768in}{1.080778in}}%
\pgfpathlineto{\pgfqpoint{1.361097in}{1.067790in}}%
\pgfpathclose%
\pgfusepath{fill}%
\end{pgfscope}%
\begin{pgfscope}%
\pgfpathrectangle{\pgfqpoint{0.407592in}{0.222904in}}{\pgfqpoint{1.411163in}{1.411163in}}%
\pgfusepath{clip}%
\pgfsetbuttcap%
\pgfsetroundjoin%
\definecolor{currentfill}{rgb}{0.909804,0.698039,0.423529}%
\pgfsetfillcolor{currentfill}%
\pgfsetfillopacity{0.120000}%
\pgfsetlinewidth{0.000000pt}%
\definecolor{currentstroke}{rgb}{0.000000,0.000000,0.000000}%
\pgfsetstrokecolor{currentstroke}%
\pgfsetdash{}{0pt}%
\pgfpathmoveto{\pgfqpoint{1.152990in}{0.960926in}}%
\pgfpathlineto{\pgfqpoint{1.169154in}{0.956740in}}%
\pgfpathlineto{\pgfqpoint{1.185359in}{0.953490in}}%
\pgfpathlineto{\pgfqpoint{1.193233in}{0.963755in}}%
\pgfpathlineto{\pgfqpoint{1.201089in}{0.974421in}}%
\pgfpathlineto{\pgfqpoint{1.184904in}{0.977661in}}%
\pgfpathlineto{\pgfqpoint{1.168761in}{0.981802in}}%
\pgfpathlineto{\pgfqpoint{1.160885in}{0.971164in}}%
\pgfpathlineto{\pgfqpoint{1.152990in}{0.960926in}}%
\pgfpathclose%
\pgfusepath{fill}%
\end{pgfscope}%
\begin{pgfscope}%
\pgfpathrectangle{\pgfqpoint{0.407592in}{0.222904in}}{\pgfqpoint{1.411163in}{1.411163in}}%
\pgfusepath{clip}%
\pgfsetbuttcap%
\pgfsetroundjoin%
\definecolor{currentfill}{rgb}{0.909804,0.698039,0.423529}%
\pgfsetfillcolor{currentfill}%
\pgfsetfillopacity{0.120000}%
\pgfsetlinewidth{0.000000pt}%
\definecolor{currentstroke}{rgb}{0.000000,0.000000,0.000000}%
\pgfsetstrokecolor{currentstroke}%
\pgfsetdash{}{0pt}%
\pgfpathmoveto{\pgfqpoint{0.815437in}{1.068166in}}%
\pgfpathlineto{\pgfqpoint{0.831407in}{1.052391in}}%
\pgfpathlineto{\pgfqpoint{0.847388in}{1.036985in}}%
\pgfpathlineto{\pgfqpoint{0.855760in}{1.041810in}}%
\pgfpathlineto{\pgfqpoint{0.864100in}{1.047056in}}%
\pgfpathlineto{\pgfqpoint{0.848158in}{1.062396in}}%
\pgfpathlineto{\pgfqpoint{0.832227in}{1.078106in}}%
\pgfpathlineto{\pgfqpoint{0.823848in}{1.072938in}}%
\pgfpathlineto{\pgfqpoint{0.815437in}{1.068166in}}%
\pgfpathclose%
\pgfusepath{fill}%
\end{pgfscope}%
\begin{pgfscope}%
\pgfpathrectangle{\pgfqpoint{0.407592in}{0.222904in}}{\pgfqpoint{1.411163in}{1.411163in}}%
\pgfusepath{clip}%
\pgfsetbuttcap%
\pgfsetroundjoin%
\definecolor{currentfill}{rgb}{0.909804,0.698039,0.423529}%
\pgfsetfillcolor{currentfill}%
\pgfsetfillopacity{0.120000}%
\pgfsetlinewidth{0.000000pt}%
\definecolor{currentstroke}{rgb}{0.000000,0.000000,0.000000}%
\pgfsetstrokecolor{currentstroke}%
\pgfsetdash{}{0pt}%
\pgfpathmoveto{\pgfqpoint{0.702675in}{1.160679in}}%
\pgfpathlineto{\pgfqpoint{0.718623in}{1.143267in}}%
\pgfpathlineto{\pgfqpoint{0.734581in}{1.126066in}}%
\pgfpathlineto{\pgfqpoint{0.743141in}{1.130024in}}%
\pgfpathlineto{\pgfqpoint{0.751666in}{1.134334in}}%
\pgfpathlineto{\pgfqpoint{0.735750in}{1.151509in}}%
\pgfpathlineto{\pgfqpoint{0.719843in}{1.168894in}}%
\pgfpathlineto{\pgfqpoint{0.711276in}{1.164621in}}%
\pgfpathlineto{\pgfqpoint{0.702675in}{1.160679in}}%
\pgfpathclose%
\pgfusepath{fill}%
\end{pgfscope}%
\begin{pgfscope}%
\pgfpathrectangle{\pgfqpoint{0.407592in}{0.222904in}}{\pgfqpoint{1.411163in}{1.411163in}}%
\pgfusepath{clip}%
\pgfsetbuttcap%
\pgfsetroundjoin%
\definecolor{currentfill}{rgb}{0.909804,0.698039,0.423529}%
\pgfsetfillcolor{currentfill}%
\pgfsetfillopacity{0.120000}%
\pgfsetlinewidth{0.000000pt}%
\definecolor{currentstroke}{rgb}{0.000000,0.000000,0.000000}%
\pgfsetstrokecolor{currentstroke}%
\pgfsetdash{}{0pt}%
\pgfpathmoveto{\pgfqpoint{1.201089in}{0.974421in}}%
\pgfpathlineto{\pgfqpoint{1.217319in}{0.972078in}}%
\pgfpathlineto{\pgfqpoint{1.233597in}{0.970624in}}%
\pgfpathlineto{\pgfqpoint{1.241417in}{0.981625in}}%
\pgfpathlineto{\pgfqpoint{1.249222in}{0.992980in}}%
\pgfpathlineto{\pgfqpoint{1.232962in}{0.994494in}}%
\pgfpathlineto{\pgfqpoint{1.216750in}{0.996863in}}%
\pgfpathlineto{\pgfqpoint{1.208928in}{0.985465in}}%
\pgfpathlineto{\pgfqpoint{1.201089in}{0.974421in}}%
\pgfpathclose%
\pgfusepath{fill}%
\end{pgfscope}%
\begin{pgfscope}%
\pgfpathrectangle{\pgfqpoint{0.407592in}{0.222904in}}{\pgfqpoint{1.411163in}{1.411163in}}%
\pgfusepath{clip}%
\pgfsetbuttcap%
\pgfsetroundjoin%
\definecolor{currentfill}{rgb}{0.909804,0.698039,0.423529}%
\pgfsetfillcolor{currentfill}%
\pgfsetfillopacity{0.120000}%
\pgfsetlinewidth{0.000000pt}%
\definecolor{currentstroke}{rgb}{0.000000,0.000000,0.000000}%
\pgfsetstrokecolor{currentstroke}%
\pgfsetdash{}{0pt}%
\pgfpathmoveto{\pgfqpoint{1.104883in}{0.953073in}}%
\pgfpathlineto{\pgfqpoint{1.120995in}{0.946937in}}%
\pgfpathlineto{\pgfqpoint{1.137142in}{0.941748in}}%
\pgfpathlineto{\pgfqpoint{1.145076in}{0.951113in}}%
\pgfpathlineto{\pgfqpoint{1.152990in}{0.960926in}}%
\pgfpathlineto{\pgfqpoint{1.136865in}{0.966042in}}%
\pgfpathlineto{\pgfqpoint{1.120777in}{0.972075in}}%
\pgfpathlineto{\pgfqpoint{1.112840in}{0.962352in}}%
\pgfpathlineto{\pgfqpoint{1.104883in}{0.953073in}}%
\pgfpathclose%
\pgfusepath{fill}%
\end{pgfscope}%
\begin{pgfscope}%
\pgfpathrectangle{\pgfqpoint{0.407592in}{0.222904in}}{\pgfqpoint{1.411163in}{1.411163in}}%
\pgfusepath{clip}%
\pgfsetbuttcap%
\pgfsetroundjoin%
\definecolor{currentfill}{rgb}{0.909804,0.698039,0.423529}%
\pgfsetfillcolor{currentfill}%
\pgfsetfillopacity{0.120000}%
\pgfsetlinewidth{0.000000pt}%
\definecolor{currentstroke}{rgb}{0.000000,0.000000,0.000000}%
\pgfsetstrokecolor{currentstroke}%
\pgfsetdash{}{0pt}%
\pgfpathmoveto{\pgfqpoint{1.473037in}{1.157082in}}%
\pgfpathlineto{\pgfqpoint{1.489719in}{1.160582in}}%
\pgfpathlineto{\pgfqpoint{1.506465in}{1.164588in}}%
\pgfpathlineto{\pgfqpoint{1.514013in}{1.177989in}}%
\pgfpathlineto{\pgfqpoint{1.497275in}{1.174104in}}%
\pgfpathlineto{\pgfqpoint{1.480601in}{1.170718in}}%
\pgfusepath{fill}%
\end{pgfscope}%
\begin{pgfscope}%
\pgfpathrectangle{\pgfqpoint{0.407592in}{0.222904in}}{\pgfqpoint{1.411163in}{1.411163in}}%
\pgfusepath{clip}%
\pgfsetbuttcap%
\pgfsetroundjoin%
\definecolor{currentfill}{rgb}{0.909804,0.698039,0.423529}%
\pgfsetfillcolor{currentfill}%
\pgfsetfillopacity{0.120000}%
\pgfsetlinewidth{0.000000pt}%
\definecolor{currentstroke}{rgb}{0.000000,0.000000,0.000000}%
\pgfsetstrokecolor{currentstroke}%
\pgfsetdash{}{0pt}%
\pgfpathmoveto{\pgfqpoint{1.249222in}{0.992980in}}%
\pgfpathlineto{\pgfqpoint{1.265531in}{0.992303in}}%
\pgfpathlineto{\pgfqpoint{1.281894in}{0.992444in}}%
\pgfpathlineto{\pgfqpoint{1.289665in}{1.004020in}}%
\pgfpathlineto{\pgfqpoint{1.297422in}{1.015905in}}%
\pgfpathlineto{\pgfqpoint{1.281076in}{1.015889in}}%
\pgfpathlineto{\pgfqpoint{1.264783in}{1.016660in}}%
\pgfpathlineto{\pgfqpoint{1.257010in}{1.004665in}}%
\pgfpathlineto{\pgfqpoint{1.249222in}{0.992980in}}%
\pgfpathclose%
\pgfusepath{fill}%
\end{pgfscope}%
\begin{pgfscope}%
\pgfpathrectangle{\pgfqpoint{0.407592in}{0.222904in}}{\pgfqpoint{1.411163in}{1.411163in}}%
\pgfusepath{clip}%
\pgfsetbuttcap%
\pgfsetroundjoin%
\definecolor{currentfill}{rgb}{0.909804,0.698039,0.423529}%
\pgfsetfillcolor{currentfill}%
\pgfsetfillopacity{0.120000}%
\pgfsetlinewidth{0.000000pt}%
\definecolor{currentstroke}{rgb}{0.000000,0.000000,0.000000}%
\pgfsetstrokecolor{currentstroke}%
\pgfsetdash{}{0pt}%
\pgfpathmoveto{\pgfqpoint{1.056720in}{0.951194in}}%
\pgfpathlineto{\pgfqpoint{1.072797in}{0.943107in}}%
\pgfpathlineto{\pgfqpoint{1.088902in}{0.935933in}}%
\pgfpathlineto{\pgfqpoint{1.096904in}{0.944260in}}%
\pgfpathlineto{\pgfqpoint{1.104883in}{0.953073in}}%
\pgfpathlineto{\pgfqpoint{1.088804in}{0.960132in}}%
\pgfpathlineto{\pgfqpoint{1.072755in}{0.968084in}}%
\pgfpathlineto{\pgfqpoint{1.064749in}{0.959401in}}%
\pgfpathlineto{\pgfqpoint{1.056720in}{0.951194in}}%
\pgfpathclose%
\pgfusepath{fill}%
\end{pgfscope}%
\begin{pgfscope}%
\pgfpathrectangle{\pgfqpoint{0.407592in}{0.222904in}}{\pgfqpoint{1.411163in}{1.411163in}}%
\pgfusepath{clip}%
\pgfsetbuttcap%
\pgfsetroundjoin%
\definecolor{currentfill}{rgb}{0.909804,0.698039,0.423529}%
\pgfsetfillcolor{currentfill}%
\pgfsetfillopacity{0.120000}%
\pgfsetlinewidth{0.000000pt}%
\definecolor{currentstroke}{rgb}{0.000000,0.000000,0.000000}%
\pgfsetstrokecolor{currentstroke}%
\pgfsetdash{}{0pt}%
\pgfpathmoveto{\pgfqpoint{1.409420in}{1.097874in}}%
\pgfpathlineto{\pgfqpoint{1.426008in}{1.100735in}}%
\pgfpathlineto{\pgfqpoint{1.442659in}{1.104169in}}%
\pgfpathlineto{\pgfqpoint{1.450272in}{1.117130in}}%
\pgfpathlineto{\pgfqpoint{1.457873in}{1.130277in}}%
\pgfpathlineto{\pgfqpoint{1.441239in}{1.127076in}}%
\pgfpathlineto{\pgfqpoint{1.424666in}{1.124435in}}%
\pgfpathlineto{\pgfqpoint{1.417049in}{1.111064in}}%
\pgfpathlineto{\pgfqpoint{1.409420in}{1.097874in}}%
\pgfpathclose%
\pgfusepath{fill}%
\end{pgfscope}%
\begin{pgfscope}%
\pgfpathrectangle{\pgfqpoint{0.407592in}{0.222904in}}{\pgfqpoint{1.411163in}{1.411163in}}%
\pgfusepath{clip}%
\pgfsetbuttcap%
\pgfsetroundjoin%
\definecolor{currentfill}{rgb}{0.909804,0.698039,0.423529}%
\pgfsetfillcolor{currentfill}%
\pgfsetfillopacity{0.120000}%
\pgfsetlinewidth{0.000000pt}%
\definecolor{currentstroke}{rgb}{0.000000,0.000000,0.000000}%
\pgfsetstrokecolor{currentstroke}%
\pgfsetdash{}{0pt}%
\pgfpathmoveto{\pgfqpoint{0.928001in}{0.990159in}}%
\pgfpathlineto{\pgfqpoint{0.944013in}{0.977275in}}%
\pgfpathlineto{\pgfqpoint{0.960042in}{0.965026in}}%
\pgfpathlineto{\pgfqpoint{0.968240in}{0.970996in}}%
\pgfpathlineto{\pgfqpoint{0.976410in}{0.977458in}}%
\pgfpathlineto{\pgfqpoint{0.960416in}{0.989596in}}%
\pgfpathlineto{\pgfqpoint{0.944440in}{1.002367in}}%
\pgfpathlineto{\pgfqpoint{0.936235in}{0.996031in}}%
\pgfpathlineto{\pgfqpoint{0.928001in}{0.990159in}}%
\pgfpathclose%
\pgfusepath{fill}%
\end{pgfscope}%
\begin{pgfscope}%
\pgfpathrectangle{\pgfqpoint{0.407592in}{0.222904in}}{\pgfqpoint{1.411163in}{1.411163in}}%
\pgfusepath{clip}%
\pgfsetbuttcap%
\pgfsetroundjoin%
\definecolor{currentfill}{rgb}{0.909804,0.698039,0.423529}%
\pgfsetfillcolor{currentfill}%
\pgfsetfillopacity{0.120000}%
\pgfsetlinewidth{0.000000pt}%
\definecolor{currentstroke}{rgb}{0.000000,0.000000,0.000000}%
\pgfsetstrokecolor{currentstroke}%
\pgfsetdash{}{0pt}%
\pgfpathmoveto{\pgfqpoint{0.734581in}{1.126066in}}%
\pgfpathlineto{\pgfqpoint{0.750548in}{1.109096in}}%
\pgfpathlineto{\pgfqpoint{0.766526in}{1.092380in}}%
\pgfpathlineto{\pgfqpoint{0.775045in}{1.096350in}}%
\pgfpathlineto{\pgfqpoint{0.783531in}{1.100694in}}%
\pgfpathlineto{\pgfqpoint{0.767593in}{1.117388in}}%
\pgfpathlineto{\pgfqpoint{0.751666in}{1.134334in}}%
\pgfpathlineto{\pgfqpoint{0.743141in}{1.130024in}}%
\pgfpathlineto{\pgfqpoint{0.734581in}{1.126066in}}%
\pgfpathclose%
\pgfusepath{fill}%
\end{pgfscope}%
\begin{pgfscope}%
\pgfpathrectangle{\pgfqpoint{0.407592in}{0.222904in}}{\pgfqpoint{1.411163in}{1.411163in}}%
\pgfusepath{clip}%
\pgfsetbuttcap%
\pgfsetroundjoin%
\definecolor{currentfill}{rgb}{0.909804,0.698039,0.423529}%
\pgfsetfillcolor{currentfill}%
\pgfsetfillopacity{0.120000}%
\pgfsetlinewidth{0.000000pt}%
\definecolor{currentstroke}{rgb}{0.000000,0.000000,0.000000}%
\pgfsetstrokecolor{currentstroke}%
\pgfsetdash{}{0pt}%
\pgfpathmoveto{\pgfqpoint{1.297422in}{1.015905in}}%
\pgfpathlineto{\pgfqpoint{1.313822in}{1.016685in}}%
\pgfpathlineto{\pgfqpoint{1.330278in}{1.018203in}}%
\pgfpathlineto{\pgfqpoint{1.338004in}{1.030211in}}%
\pgfpathlineto{\pgfqpoint{1.345716in}{1.042490in}}%
\pgfpathlineto{\pgfqpoint{1.329275in}{1.041149in}}%
\pgfpathlineto{\pgfqpoint{1.312891in}{1.040522in}}%
\pgfpathlineto{\pgfqpoint{1.305163in}{1.028078in}}%
\pgfpathlineto{\pgfqpoint{1.297422in}{1.015905in}}%
\pgfpathclose%
\pgfusepath{fill}%
\end{pgfscope}%
\begin{pgfscope}%
\pgfpathrectangle{\pgfqpoint{0.407592in}{0.222904in}}{\pgfqpoint{1.411163in}{1.411163in}}%
\pgfusepath{clip}%
\pgfsetbuttcap%
\pgfsetroundjoin%
\definecolor{currentfill}{rgb}{0.909804,0.698039,0.423529}%
\pgfsetfillcolor{currentfill}%
\pgfsetfillopacity{0.120000}%
\pgfsetlinewidth{0.000000pt}%
\definecolor{currentstroke}{rgb}{0.000000,0.000000,0.000000}%
\pgfsetstrokecolor{currentstroke}%
\pgfsetdash{}{0pt}%
\pgfpathmoveto{\pgfqpoint{0.847388in}{1.036985in}}%
\pgfpathlineto{\pgfqpoint{0.863382in}{1.021984in}}%
\pgfpathlineto{\pgfqpoint{0.879388in}{1.007429in}}%
\pgfpathlineto{\pgfqpoint{0.887721in}{1.012306in}}%
\pgfpathlineto{\pgfqpoint{0.896023in}{1.017630in}}%
\pgfpathlineto{\pgfqpoint{0.880055in}{1.032120in}}%
\pgfpathlineto{\pgfqpoint{0.864100in}{1.047056in}}%
\pgfpathlineto{\pgfqpoint{0.855760in}{1.041810in}}%
\pgfpathlineto{\pgfqpoint{0.847388in}{1.036985in}}%
\pgfpathclose%
\pgfusepath{fill}%
\end{pgfscope}%
\begin{pgfscope}%
\pgfpathrectangle{\pgfqpoint{0.407592in}{0.222904in}}{\pgfqpoint{1.411163in}{1.411163in}}%
\pgfusepath{clip}%
\pgfsetbuttcap%
\pgfsetroundjoin%
\definecolor{currentfill}{rgb}{0.909804,0.698039,0.423529}%
\pgfsetfillcolor{currentfill}%
\pgfsetfillopacity{0.120000}%
\pgfsetlinewidth{0.000000pt}%
\definecolor{currentstroke}{rgb}{0.000000,0.000000,0.000000}%
\pgfsetstrokecolor{currentstroke}%
\pgfsetdash{}{0pt}%
\pgfpathmoveto{\pgfqpoint{1.008455in}{0.955288in}}%
\pgfpathlineto{\pgfqpoint{1.024510in}{0.945362in}}%
\pgfpathlineto{\pgfqpoint{1.040590in}{0.936278in}}%
\pgfpathlineto{\pgfqpoint{1.048668in}{0.943481in}}%
\pgfpathlineto{\pgfqpoint{1.056720in}{0.951194in}}%
\pgfpathlineto{\pgfqpoint{1.040672in}{0.960155in}}%
\pgfpathlineto{\pgfqpoint{1.024648in}{0.969945in}}%
\pgfpathlineto{\pgfqpoint{1.016564in}{0.962371in}}%
\pgfpathlineto{\pgfqpoint{1.008455in}{0.955288in}}%
\pgfpathclose%
\pgfusepath{fill}%
\end{pgfscope}%
\begin{pgfscope}%
\pgfpathrectangle{\pgfqpoint{0.407592in}{0.222904in}}{\pgfqpoint{1.411163in}{1.411163in}}%
\pgfusepath{clip}%
\pgfsetbuttcap%
\pgfsetroundjoin%
\definecolor{currentfill}{rgb}{0.909804,0.698039,0.423529}%
\pgfsetfillcolor{currentfill}%
\pgfsetfillopacity{0.120000}%
\pgfsetlinewidth{0.000000pt}%
\definecolor{currentstroke}{rgb}{0.000000,0.000000,0.000000}%
\pgfsetstrokecolor{currentstroke}%
\pgfsetdash{}{0pt}%
\pgfpathmoveto{\pgfqpoint{0.621502in}{1.225186in}}%
\pgfpathlineto{\pgfqpoint{0.637454in}{1.207098in}}%
\pgfpathlineto{\pgfqpoint{0.653415in}{1.189160in}}%
\pgfpathlineto{\pgfqpoint{0.662129in}{1.192456in}}%
\pgfpathlineto{\pgfqpoint{0.670808in}{1.196064in}}%
\pgfpathlineto{\pgfqpoint{0.654888in}{1.214004in}}%
\pgfpathlineto{\pgfqpoint{0.638978in}{1.232091in}}%
\pgfpathlineto{\pgfqpoint{0.630258in}{1.228490in}}%
\pgfpathlineto{\pgfqpoint{0.621502in}{1.225186in}}%
\pgfpathclose%
\pgfusepath{fill}%
\end{pgfscope}%
\begin{pgfscope}%
\pgfpathrectangle{\pgfqpoint{0.407592in}{0.222904in}}{\pgfqpoint{1.411163in}{1.411163in}}%
\pgfusepath{clip}%
\pgfsetbuttcap%
\pgfsetroundjoin%
\definecolor{currentfill}{rgb}{0.909804,0.698039,0.423529}%
\pgfsetfillcolor{currentfill}%
\pgfsetfillopacity{0.120000}%
\pgfsetlinewidth{0.000000pt}%
\definecolor{currentstroke}{rgb}{0.000000,0.000000,0.000000}%
\pgfsetstrokecolor{currentstroke}%
\pgfsetdash{}{0pt}%
\pgfpathmoveto{\pgfqpoint{1.457873in}{1.130277in}}%
\pgfpathlineto{\pgfqpoint{1.474571in}{1.134015in}}%
\pgfpathlineto{\pgfqpoint{1.491333in}{1.138269in}}%
\pgfpathlineto{\pgfqpoint{1.498905in}{1.151344in}}%
\pgfpathlineto{\pgfqpoint{1.506465in}{1.164588in}}%
\pgfpathlineto{\pgfqpoint{1.489719in}{1.160582in}}%
\pgfpathlineto{\pgfqpoint{1.473037in}{1.157082in}}%
\pgfpathlineto{\pgfqpoint{1.465462in}{1.143598in}}%
\pgfpathlineto{\pgfqpoint{1.457873in}{1.130277in}}%
\pgfpathclose%
\pgfusepath{fill}%
\end{pgfscope}%
\begin{pgfscope}%
\pgfpathrectangle{\pgfqpoint{0.407592in}{0.222904in}}{\pgfqpoint{1.411163in}{1.411163in}}%
\pgfusepath{clip}%
\pgfsetbuttcap%
\pgfsetroundjoin%
\definecolor{currentfill}{rgb}{0.909804,0.698039,0.423529}%
\pgfsetfillcolor{currentfill}%
\pgfsetfillopacity{0.120000}%
\pgfsetlinewidth{0.000000pt}%
\definecolor{currentstroke}{rgb}{0.000000,0.000000,0.000000}%
\pgfsetstrokecolor{currentstroke}%
\pgfsetdash{}{0pt}%
\pgfpathmoveto{\pgfqpoint{1.345716in}{1.042490in}}%
\pgfpathlineto{\pgfqpoint{1.362214in}{1.044518in}}%
\pgfpathlineto{\pgfqpoint{1.378773in}{1.047207in}}%
\pgfpathlineto{\pgfqpoint{1.386455in}{1.059531in}}%
\pgfpathlineto{\pgfqpoint{1.394123in}{1.072094in}}%
\pgfpathlineto{\pgfqpoint{1.377580in}{1.069621in}}%
\pgfpathlineto{\pgfqpoint{1.361097in}{1.067790in}}%
\pgfpathlineto{\pgfqpoint{1.353413in}{1.055022in}}%
\pgfpathlineto{\pgfqpoint{1.345716in}{1.042490in}}%
\pgfpathclose%
\pgfusepath{fill}%
\end{pgfscope}%
\begin{pgfscope}%
\pgfpathrectangle{\pgfqpoint{0.407592in}{0.222904in}}{\pgfqpoint{1.411163in}{1.411163in}}%
\pgfusepath{clip}%
\pgfsetbuttcap%
\pgfsetroundjoin%
\definecolor{currentfill}{rgb}{0.909804,0.698039,0.423529}%
\pgfsetfillcolor{currentfill}%
\pgfsetfillopacity{0.120000}%
\pgfsetlinewidth{0.000000pt}%
\definecolor{currentstroke}{rgb}{0.000000,0.000000,0.000000}%
\pgfsetstrokecolor{currentstroke}%
\pgfsetdash{}{0pt}%
\pgfpathmoveto{\pgfqpoint{0.766526in}{1.092380in}}%
\pgfpathlineto{\pgfqpoint{0.782515in}{1.075944in}}%
\pgfpathlineto{\pgfqpoint{0.798513in}{1.059814in}}%
\pgfpathlineto{\pgfqpoint{0.806992in}{1.063792in}}%
\pgfpathlineto{\pgfqpoint{0.815437in}{1.068166in}}%
\pgfpathlineto{\pgfqpoint{0.799478in}{1.084278in}}%
\pgfpathlineto{\pgfqpoint{0.783531in}{1.100694in}}%
\pgfpathlineto{\pgfqpoint{0.775045in}{1.096350in}}%
\pgfpathlineto{\pgfqpoint{0.766526in}{1.092380in}}%
\pgfpathclose%
\pgfusepath{fill}%
\end{pgfscope}%
\begin{pgfscope}%
\pgfpathrectangle{\pgfqpoint{0.407592in}{0.222904in}}{\pgfqpoint{1.411163in}{1.411163in}}%
\pgfusepath{clip}%
\pgfsetbuttcap%
\pgfsetroundjoin%
\definecolor{currentfill}{rgb}{0.909804,0.698039,0.423529}%
\pgfsetfillcolor{currentfill}%
\pgfsetfillopacity{0.120000}%
\pgfsetlinewidth{0.000000pt}%
\definecolor{currentstroke}{rgb}{0.000000,0.000000,0.000000}%
\pgfsetstrokecolor{currentstroke}%
\pgfsetdash{}{0pt}%
\pgfpathmoveto{\pgfqpoint{0.653415in}{1.189160in}}%
\pgfpathlineto{\pgfqpoint{0.669386in}{1.171385in}}%
\pgfpathlineto{\pgfqpoint{0.685366in}{1.153789in}}%
\pgfpathlineto{\pgfqpoint{0.694038in}{1.157069in}}%
\pgfpathlineto{\pgfqpoint{0.702675in}{1.160679in}}%
\pgfpathlineto{\pgfqpoint{0.686737in}{1.178284in}}%
\pgfpathlineto{\pgfqpoint{0.670808in}{1.196064in}}%
\pgfpathlineto{\pgfqpoint{0.662129in}{1.192456in}}%
\pgfpathlineto{\pgfqpoint{0.653415in}{1.189160in}}%
\pgfpathclose%
\pgfusepath{fill}%
\end{pgfscope}%
\begin{pgfscope}%
\pgfpathrectangle{\pgfqpoint{0.407592in}{0.222904in}}{\pgfqpoint{1.411163in}{1.411163in}}%
\pgfusepath{clip}%
\pgfsetbuttcap%
\pgfsetroundjoin%
\definecolor{currentfill}{rgb}{0.909804,0.698039,0.423529}%
\pgfsetfillcolor{currentfill}%
\pgfsetfillopacity{0.120000}%
\pgfsetlinewidth{0.000000pt}%
\definecolor{currentstroke}{rgb}{0.000000,0.000000,0.000000}%
\pgfsetstrokecolor{currentstroke}%
\pgfsetdash{}{0pt}%
\pgfpathmoveto{\pgfqpoint{1.137142in}{0.941748in}}%
\pgfpathlineto{\pgfqpoint{1.153328in}{0.937520in}}%
\pgfpathlineto{\pgfqpoint{1.169555in}{0.934259in}}%
\pgfpathlineto{\pgfqpoint{1.177467in}{0.943650in}}%
\pgfpathlineto{\pgfqpoint{1.185359in}{0.953490in}}%
\pgfpathlineto{\pgfqpoint{1.169154in}{0.956740in}}%
\pgfpathlineto{\pgfqpoint{1.152990in}{0.960926in}}%
\pgfpathlineto{\pgfqpoint{1.145076in}{0.951113in}}%
\pgfpathlineto{\pgfqpoint{1.137142in}{0.941748in}}%
\pgfpathclose%
\pgfusepath{fill}%
\end{pgfscope}%
\begin{pgfscope}%
\pgfpathrectangle{\pgfqpoint{0.407592in}{0.222904in}}{\pgfqpoint{1.411163in}{1.411163in}}%
\pgfusepath{clip}%
\pgfsetbuttcap%
\pgfsetroundjoin%
\definecolor{currentfill}{rgb}{0.909804,0.698039,0.423529}%
\pgfsetfillcolor{currentfill}%
\pgfsetfillopacity{0.120000}%
\pgfsetlinewidth{0.000000pt}%
\definecolor{currentstroke}{rgb}{0.000000,0.000000,0.000000}%
\pgfsetstrokecolor{currentstroke}%
\pgfsetdash{}{0pt}%
\pgfpathmoveto{\pgfqpoint{1.185359in}{0.953490in}}%
\pgfpathlineto{\pgfqpoint{1.201609in}{0.951171in}}%
\pgfpathlineto{\pgfqpoint{1.217906in}{0.949774in}}%
\pgfpathlineto{\pgfqpoint{1.225760in}{0.959999in}}%
\pgfpathlineto{\pgfqpoint{1.233597in}{0.970624in}}%
\pgfpathlineto{\pgfqpoint{1.217319in}{0.972078in}}%
\pgfpathlineto{\pgfqpoint{1.201089in}{0.974421in}}%
\pgfpathlineto{\pgfqpoint{1.193233in}{0.963755in}}%
\pgfpathlineto{\pgfqpoint{1.185359in}{0.953490in}}%
\pgfpathclose%
\pgfusepath{fill}%
\end{pgfscope}%
\begin{pgfscope}%
\pgfpathrectangle{\pgfqpoint{0.407592in}{0.222904in}}{\pgfqpoint{1.411163in}{1.411163in}}%
\pgfusepath{clip}%
\pgfsetbuttcap%
\pgfsetroundjoin%
\definecolor{currentfill}{rgb}{0.909804,0.698039,0.423529}%
\pgfsetfillcolor{currentfill}%
\pgfsetfillopacity{0.120000}%
\pgfsetlinewidth{0.000000pt}%
\definecolor{currentstroke}{rgb}{0.000000,0.000000,0.000000}%
\pgfsetstrokecolor{currentstroke}%
\pgfsetdash{}{0pt}%
\pgfpathmoveto{\pgfqpoint{0.960042in}{0.965026in}}%
\pgfpathlineto{\pgfqpoint{0.976089in}{0.953465in}}%
\pgfpathlineto{\pgfqpoint{0.992156in}{0.942647in}}%
\pgfpathlineto{\pgfqpoint{1.000319in}{0.948710in}}%
\pgfpathlineto{\pgfqpoint{1.008455in}{0.955288in}}%
\pgfpathlineto{\pgfqpoint{0.992422in}{0.966004in}}%
\pgfpathlineto{\pgfqpoint{0.976410in}{0.977458in}}%
\pgfpathlineto{\pgfqpoint{0.968240in}{0.970996in}}%
\pgfpathlineto{\pgfqpoint{0.960042in}{0.965026in}}%
\pgfpathclose%
\pgfusepath{fill}%
\end{pgfscope}%
\begin{pgfscope}%
\pgfpathrectangle{\pgfqpoint{0.407592in}{0.222904in}}{\pgfqpoint{1.411163in}{1.411163in}}%
\pgfusepath{clip}%
\pgfsetbuttcap%
\pgfsetroundjoin%
\definecolor{currentfill}{rgb}{0.909804,0.698039,0.423529}%
\pgfsetfillcolor{currentfill}%
\pgfsetfillopacity{0.120000}%
\pgfsetlinewidth{0.000000pt}%
\definecolor{currentstroke}{rgb}{0.000000,0.000000,0.000000}%
\pgfsetstrokecolor{currentstroke}%
\pgfsetdash{}{0pt}%
\pgfpathmoveto{\pgfqpoint{0.879388in}{1.007429in}}%
\pgfpathlineto{\pgfqpoint{0.895408in}{0.993364in}}%
\pgfpathlineto{\pgfqpoint{0.911442in}{0.979834in}}%
\pgfpathlineto{\pgfqpoint{0.919737in}{0.984759in}}%
\pgfpathlineto{\pgfqpoint{0.928001in}{0.990159in}}%
\pgfpathlineto{\pgfqpoint{0.912005in}{1.003627in}}%
\pgfpathlineto{\pgfqpoint{0.896023in}{1.017630in}}%
\pgfpathlineto{\pgfqpoint{0.887721in}{1.012306in}}%
\pgfpathlineto{\pgfqpoint{0.879388in}{1.007429in}}%
\pgfpathclose%
\pgfusepath{fill}%
\end{pgfscope}%
\begin{pgfscope}%
\pgfpathrectangle{\pgfqpoint{0.407592in}{0.222904in}}{\pgfqpoint{1.411163in}{1.411163in}}%
\pgfusepath{clip}%
\pgfsetbuttcap%
\pgfsetroundjoin%
\definecolor{currentfill}{rgb}{0.909804,0.698039,0.423529}%
\pgfsetfillcolor{currentfill}%
\pgfsetfillopacity{0.120000}%
\pgfsetlinewidth{0.000000pt}%
\definecolor{currentstroke}{rgb}{0.000000,0.000000,0.000000}%
\pgfsetstrokecolor{currentstroke}%
\pgfsetdash{}{0pt}%
\pgfpathmoveto{\pgfqpoint{1.088902in}{0.935933in}}%
\pgfpathlineto{\pgfqpoint{1.105040in}{0.929706in}}%
\pgfpathlineto{\pgfqpoint{1.121212in}{0.924452in}}%
\pgfpathlineto{\pgfqpoint{1.129188in}{0.932854in}}%
\pgfpathlineto{\pgfqpoint{1.137142in}{0.941748in}}%
\pgfpathlineto{\pgfqpoint{1.120995in}{0.946937in}}%
\pgfpathlineto{\pgfqpoint{1.104883in}{0.953073in}}%
\pgfpathlineto{\pgfqpoint{1.096904in}{0.944260in}}%
\pgfpathlineto{\pgfqpoint{1.088902in}{0.935933in}}%
\pgfpathclose%
\pgfusepath{fill}%
\end{pgfscope}%
\begin{pgfscope}%
\pgfpathrectangle{\pgfqpoint{0.407592in}{0.222904in}}{\pgfqpoint{1.411163in}{1.411163in}}%
\pgfusepath{clip}%
\pgfsetbuttcap%
\pgfsetroundjoin%
\definecolor{currentfill}{rgb}{0.909804,0.698039,0.423529}%
\pgfsetfillcolor{currentfill}%
\pgfsetfillopacity{0.120000}%
\pgfsetlinewidth{0.000000pt}%
\definecolor{currentstroke}{rgb}{0.000000,0.000000,0.000000}%
\pgfsetstrokecolor{currentstroke}%
\pgfsetdash{}{0pt}%
\pgfpathmoveto{\pgfqpoint{1.233597in}{0.970624in}}%
\pgfpathlineto{\pgfqpoint{1.249924in}{0.970040in}}%
\pgfpathlineto{\pgfqpoint{1.266304in}{0.970304in}}%
\pgfpathlineto{\pgfqpoint{1.274107in}{0.981198in}}%
\pgfpathlineto{\pgfqpoint{1.281894in}{0.992444in}}%
\pgfpathlineto{\pgfqpoint{1.265531in}{0.992303in}}%
\pgfpathlineto{\pgfqpoint{1.249222in}{0.992980in}}%
\pgfpathlineto{\pgfqpoint{1.241417in}{0.981625in}}%
\pgfpathlineto{\pgfqpoint{1.233597in}{0.970624in}}%
\pgfpathclose%
\pgfusepath{fill}%
\end{pgfscope}%
\begin{pgfscope}%
\pgfpathrectangle{\pgfqpoint{0.407592in}{0.222904in}}{\pgfqpoint{1.411163in}{1.411163in}}%
\pgfusepath{clip}%
\pgfsetbuttcap%
\pgfsetroundjoin%
\definecolor{currentfill}{rgb}{0.909804,0.698039,0.423529}%
\pgfsetfillcolor{currentfill}%
\pgfsetfillopacity{0.120000}%
\pgfsetlinewidth{0.000000pt}%
\definecolor{currentstroke}{rgb}{0.000000,0.000000,0.000000}%
\pgfsetstrokecolor{currentstroke}%
\pgfsetdash{}{0pt}%
\pgfpathmoveto{\pgfqpoint{1.394123in}{1.072094in}}%
\pgfpathlineto{\pgfqpoint{1.410727in}{1.075181in}}%
\pgfpathlineto{\pgfqpoint{1.427394in}{1.078857in}}%
\pgfpathlineto{\pgfqpoint{1.435033in}{1.091407in}}%
\pgfpathlineto{\pgfqpoint{1.442659in}{1.104169in}}%
\pgfpathlineto{\pgfqpoint{1.426008in}{1.100735in}}%
\pgfpathlineto{\pgfqpoint{1.409420in}{1.097874in}}%
\pgfpathlineto{\pgfqpoint{1.401778in}{1.084880in}}%
\pgfpathlineto{\pgfqpoint{1.394123in}{1.072094in}}%
\pgfpathclose%
\pgfusepath{fill}%
\end{pgfscope}%
\begin{pgfscope}%
\pgfpathrectangle{\pgfqpoint{0.407592in}{0.222904in}}{\pgfqpoint{1.411163in}{1.411163in}}%
\pgfusepath{clip}%
\pgfsetbuttcap%
\pgfsetroundjoin%
\definecolor{currentfill}{rgb}{0.909804,0.698039,0.423529}%
\pgfsetfillcolor{currentfill}%
\pgfsetfillopacity{0.120000}%
\pgfsetlinewidth{0.000000pt}%
\definecolor{currentstroke}{rgb}{0.000000,0.000000,0.000000}%
\pgfsetstrokecolor{currentstroke}%
\pgfsetdash{}{0pt}%
\pgfpathmoveto{\pgfqpoint{0.798513in}{1.059814in}}%
\pgfpathlineto{\pgfqpoint{0.814523in}{1.044023in}}%
\pgfpathlineto{\pgfqpoint{0.830545in}{1.028602in}}%
\pgfpathlineto{\pgfqpoint{0.838983in}{1.032582in}}%
\pgfpathlineto{\pgfqpoint{0.847388in}{1.036985in}}%
\pgfpathlineto{\pgfqpoint{0.831407in}{1.052391in}}%
\pgfpathlineto{\pgfqpoint{0.815437in}{1.068166in}}%
\pgfpathlineto{\pgfqpoint{0.806992in}{1.063792in}}%
\pgfpathlineto{\pgfqpoint{0.798513in}{1.059814in}}%
\pgfpathclose%
\pgfusepath{fill}%
\end{pgfscope}%
\begin{pgfscope}%
\pgfpathrectangle{\pgfqpoint{0.407592in}{0.222904in}}{\pgfqpoint{1.411163in}{1.411163in}}%
\pgfusepath{clip}%
\pgfsetbuttcap%
\pgfsetroundjoin%
\definecolor{currentfill}{rgb}{0.909804,0.698039,0.423529}%
\pgfsetfillcolor{currentfill}%
\pgfsetfillopacity{0.120000}%
\pgfsetlinewidth{0.000000pt}%
\definecolor{currentstroke}{rgb}{0.000000,0.000000,0.000000}%
\pgfsetstrokecolor{currentstroke}%
\pgfsetdash{}{0pt}%
\pgfpathmoveto{\pgfqpoint{0.685366in}{1.153789in}}%
\pgfpathlineto{\pgfqpoint{0.701355in}{1.136388in}}%
\pgfpathlineto{\pgfqpoint{0.717354in}{1.119202in}}%
\pgfpathlineto{\pgfqpoint{0.725985in}{1.122459in}}%
\pgfpathlineto{\pgfqpoint{0.734581in}{1.126066in}}%
\pgfpathlineto{\pgfqpoint{0.718623in}{1.143267in}}%
\pgfpathlineto{\pgfqpoint{0.702675in}{1.160679in}}%
\pgfpathlineto{\pgfqpoint{0.694038in}{1.157069in}}%
\pgfpathlineto{\pgfqpoint{0.685366in}{1.153789in}}%
\pgfpathclose%
\pgfusepath{fill}%
\end{pgfscope}%
\begin{pgfscope}%
\pgfpathrectangle{\pgfqpoint{0.407592in}{0.222904in}}{\pgfqpoint{1.411163in}{1.411163in}}%
\pgfusepath{clip}%
\pgfsetbuttcap%
\pgfsetroundjoin%
\definecolor{currentfill}{rgb}{0.909804,0.698039,0.423529}%
\pgfsetfillcolor{currentfill}%
\pgfsetfillopacity{0.120000}%
\pgfsetlinewidth{0.000000pt}%
\definecolor{currentstroke}{rgb}{0.000000,0.000000,0.000000}%
\pgfsetstrokecolor{currentstroke}%
\pgfsetdash{}{0pt}%
\pgfpathmoveto{\pgfqpoint{1.040590in}{0.936278in}}%
\pgfpathlineto{\pgfqpoint{1.056695in}{0.928081in}}%
\pgfpathlineto{\pgfqpoint{1.072829in}{0.920814in}}%
\pgfpathlineto{\pgfqpoint{1.080878in}{0.928112in}}%
\pgfpathlineto{\pgfqpoint{1.088902in}{0.935933in}}%
\pgfpathlineto{\pgfqpoint{1.072797in}{0.943107in}}%
\pgfpathlineto{\pgfqpoint{1.056720in}{0.951194in}}%
\pgfpathlineto{\pgfqpoint{1.048668in}{0.943481in}}%
\pgfpathlineto{\pgfqpoint{1.040590in}{0.936278in}}%
\pgfpathclose%
\pgfusepath{fill}%
\end{pgfscope}%
\begin{pgfscope}%
\pgfpathrectangle{\pgfqpoint{0.407592in}{0.222904in}}{\pgfqpoint{1.411163in}{1.411163in}}%
\pgfusepath{clip}%
\pgfsetbuttcap%
\pgfsetroundjoin%
\definecolor{currentfill}{rgb}{0.909804,0.698039,0.423529}%
\pgfsetfillcolor{currentfill}%
\pgfsetfillopacity{0.120000}%
\pgfsetlinewidth{0.000000pt}%
\definecolor{currentstroke}{rgb}{0.000000,0.000000,0.000000}%
\pgfsetstrokecolor{currentstroke}%
\pgfsetdash{}{0pt}%
\pgfpathmoveto{\pgfqpoint{1.281894in}{0.992444in}}%
\pgfpathlineto{\pgfqpoint{1.298310in}{0.993377in}}%
\pgfpathlineto{\pgfqpoint{1.314784in}{0.995073in}}%
\pgfpathlineto{\pgfqpoint{1.322538in}{1.006484in}}%
\pgfpathlineto{\pgfqpoint{1.330278in}{1.018203in}}%
\pgfpathlineto{\pgfqpoint{1.313822in}{1.016685in}}%
\pgfpathlineto{\pgfqpoint{1.297422in}{1.015905in}}%
\pgfpathlineto{\pgfqpoint{1.289665in}{1.004020in}}%
\pgfpathlineto{\pgfqpoint{1.281894in}{0.992444in}}%
\pgfpathclose%
\pgfusepath{fill}%
\end{pgfscope}%
\begin{pgfscope}%
\pgfpathrectangle{\pgfqpoint{0.407592in}{0.222904in}}{\pgfqpoint{1.411163in}{1.411163in}}%
\pgfusepath{clip}%
\pgfsetbuttcap%
\pgfsetroundjoin%
\definecolor{currentfill}{rgb}{0.909804,0.698039,0.423529}%
\pgfsetfillcolor{currentfill}%
\pgfsetfillopacity{0.120000}%
\pgfsetlinewidth{0.000000pt}%
\definecolor{currentstroke}{rgb}{0.000000,0.000000,0.000000}%
\pgfsetstrokecolor{currentstroke}%
\pgfsetdash{}{0pt}%
\pgfpathmoveto{\pgfqpoint{1.506465in}{1.164588in}}%
\pgfpathlineto{\pgfqpoint{1.523277in}{1.169078in}}%
\pgfpathlineto{\pgfqpoint{1.540154in}{1.174033in}}%
\pgfpathlineto{\pgfqpoint{1.547685in}{1.187181in}}%
\pgfpathlineto{\pgfqpoint{1.530816in}{1.182354in}}%
\pgfpathlineto{\pgfqpoint{1.514013in}{1.177989in}}%
\pgfusepath{fill}%
\end{pgfscope}%
\begin{pgfscope}%
\pgfpathrectangle{\pgfqpoint{0.407592in}{0.222904in}}{\pgfqpoint{1.411163in}{1.411163in}}%
\pgfusepath{clip}%
\pgfsetbuttcap%
\pgfsetroundjoin%
\definecolor{currentfill}{rgb}{0.909804,0.698039,0.423529}%
\pgfsetfillcolor{currentfill}%
\pgfsetfillopacity{0.120000}%
\pgfsetlinewidth{0.000000pt}%
\definecolor{currentstroke}{rgb}{0.000000,0.000000,0.000000}%
\pgfsetstrokecolor{currentstroke}%
\pgfsetdash{}{0pt}%
\pgfpathmoveto{\pgfqpoint{0.911442in}{0.979834in}}%
\pgfpathlineto{\pgfqpoint{0.927491in}{0.966889in}}%
\pgfpathlineto{\pgfqpoint{0.943557in}{0.954581in}}%
\pgfpathlineto{\pgfqpoint{0.951814in}{0.959552in}}%
\pgfpathlineto{\pgfqpoint{0.960042in}{0.965026in}}%
\pgfpathlineto{\pgfqpoint{0.944013in}{0.977275in}}%
\pgfpathlineto{\pgfqpoint{0.928001in}{0.990159in}}%
\pgfpathlineto{\pgfqpoint{0.919737in}{0.984759in}}%
\pgfpathlineto{\pgfqpoint{0.911442in}{0.979834in}}%
\pgfpathclose%
\pgfusepath{fill}%
\end{pgfscope}%
\begin{pgfscope}%
\pgfpathrectangle{\pgfqpoint{0.407592in}{0.222904in}}{\pgfqpoint{1.411163in}{1.411163in}}%
\pgfusepath{clip}%
\pgfsetbuttcap%
\pgfsetroundjoin%
\definecolor{currentfill}{rgb}{0.909804,0.698039,0.423529}%
\pgfsetfillcolor{currentfill}%
\pgfsetfillopacity{0.120000}%
\pgfsetlinewidth{0.000000pt}%
\definecolor{currentstroke}{rgb}{0.000000,0.000000,0.000000}%
\pgfsetstrokecolor{currentstroke}%
\pgfsetdash{}{0pt}%
\pgfpathmoveto{\pgfqpoint{1.442659in}{1.104169in}}%
\pgfpathlineto{\pgfqpoint{1.459373in}{1.108154in}}%
\pgfpathlineto{\pgfqpoint{1.476152in}{1.112664in}}%
\pgfpathlineto{\pgfqpoint{1.483749in}{1.125371in}}%
\pgfpathlineto{\pgfqpoint{1.491333in}{1.138269in}}%
\pgfpathlineto{\pgfqpoint{1.474571in}{1.134015in}}%
\pgfpathlineto{\pgfqpoint{1.457873in}{1.130277in}}%
\pgfpathlineto{\pgfqpoint{1.450272in}{1.117130in}}%
\pgfpathlineto{\pgfqpoint{1.442659in}{1.104169in}}%
\pgfpathclose%
\pgfusepath{fill}%
\end{pgfscope}%
\begin{pgfscope}%
\pgfpathrectangle{\pgfqpoint{0.407592in}{0.222904in}}{\pgfqpoint{1.411163in}{1.411163in}}%
\pgfusepath{clip}%
\pgfsetbuttcap%
\pgfsetroundjoin%
\definecolor{currentfill}{rgb}{0.909804,0.698039,0.423529}%
\pgfsetfillcolor{currentfill}%
\pgfsetfillopacity{0.120000}%
\pgfsetlinewidth{0.000000pt}%
\definecolor{currentstroke}{rgb}{0.000000,0.000000,0.000000}%
\pgfsetstrokecolor{currentstroke}%
\pgfsetdash{}{0pt}%
\pgfpathmoveto{\pgfqpoint{1.330278in}{1.018203in}}%
\pgfpathlineto{\pgfqpoint{1.346793in}{1.020431in}}%
\pgfpathlineto{\pgfqpoint{1.363369in}{1.023338in}}%
\pgfpathlineto{\pgfqpoint{1.371078in}{1.035137in}}%
\pgfpathlineto{\pgfqpoint{1.378773in}{1.047207in}}%
\pgfpathlineto{\pgfqpoint{1.362214in}{1.044518in}}%
\pgfpathlineto{\pgfqpoint{1.345716in}{1.042490in}}%
\pgfpathlineto{\pgfqpoint{1.338004in}{1.030211in}}%
\pgfpathlineto{\pgfqpoint{1.330278in}{1.018203in}}%
\pgfpathclose%
\pgfusepath{fill}%
\end{pgfscope}%
\begin{pgfscope}%
\pgfpathrectangle{\pgfqpoint{0.407592in}{0.222904in}}{\pgfqpoint{1.411163in}{1.411163in}}%
\pgfusepath{clip}%
\pgfsetbuttcap%
\pgfsetroundjoin%
\definecolor{currentfill}{rgb}{0.909804,0.698039,0.423529}%
\pgfsetfillcolor{currentfill}%
\pgfsetfillopacity{0.120000}%
\pgfsetlinewidth{0.000000pt}%
\definecolor{currentstroke}{rgb}{0.000000,0.000000,0.000000}%
\pgfsetstrokecolor{currentstroke}%
\pgfsetdash{}{0pt}%
\pgfpathmoveto{\pgfqpoint{0.717354in}{1.119202in}}%
\pgfpathlineto{\pgfqpoint{0.733363in}{1.102251in}}%
\pgfpathlineto{\pgfqpoint{0.749382in}{1.085557in}}%
\pgfpathlineto{\pgfqpoint{0.757972in}{1.088783in}}%
\pgfpathlineto{\pgfqpoint{0.766526in}{1.092380in}}%
\pgfpathlineto{\pgfqpoint{0.750548in}{1.109096in}}%
\pgfpathlineto{\pgfqpoint{0.734581in}{1.126066in}}%
\pgfpathlineto{\pgfqpoint{0.725985in}{1.122459in}}%
\pgfpathlineto{\pgfqpoint{0.717354in}{1.119202in}}%
\pgfpathclose%
\pgfusepath{fill}%
\end{pgfscope}%
\begin{pgfscope}%
\pgfpathrectangle{\pgfqpoint{0.407592in}{0.222904in}}{\pgfqpoint{1.411163in}{1.411163in}}%
\pgfusepath{clip}%
\pgfsetbuttcap%
\pgfsetroundjoin%
\definecolor{currentfill}{rgb}{0.909804,0.698039,0.423529}%
\pgfsetfillcolor{currentfill}%
\pgfsetfillopacity{0.120000}%
\pgfsetlinewidth{0.000000pt}%
\definecolor{currentstroke}{rgb}{0.000000,0.000000,0.000000}%
\pgfsetstrokecolor{currentstroke}%
\pgfsetdash{}{0pt}%
\pgfpathmoveto{\pgfqpoint{0.992156in}{0.942647in}}%
\pgfpathlineto{\pgfqpoint{1.008244in}{0.932626in}}%
\pgfpathlineto{\pgfqpoint{1.024355in}{0.923453in}}%
\pgfpathlineto{\pgfqpoint{1.032486in}{0.929598in}}%
\pgfpathlineto{\pgfqpoint{1.040590in}{0.936278in}}%
\pgfpathlineto{\pgfqpoint{1.024510in}{0.945362in}}%
\pgfpathlineto{\pgfqpoint{1.008455in}{0.955288in}}%
\pgfpathlineto{\pgfqpoint{1.000319in}{0.948710in}}%
\pgfpathlineto{\pgfqpoint{0.992156in}{0.942647in}}%
\pgfpathclose%
\pgfusepath{fill}%
\end{pgfscope}%
\begin{pgfscope}%
\pgfpathrectangle{\pgfqpoint{0.407592in}{0.222904in}}{\pgfqpoint{1.411163in}{1.411163in}}%
\pgfusepath{clip}%
\pgfsetbuttcap%
\pgfsetroundjoin%
\definecolor{currentfill}{rgb}{0.909804,0.698039,0.423529}%
\pgfsetfillcolor{currentfill}%
\pgfsetfillopacity{0.120000}%
\pgfsetlinewidth{0.000000pt}%
\definecolor{currentstroke}{rgb}{0.000000,0.000000,0.000000}%
\pgfsetstrokecolor{currentstroke}%
\pgfsetdash{}{0pt}%
\pgfpathmoveto{\pgfqpoint{0.830545in}{1.028602in}}%
\pgfpathlineto{\pgfqpoint{0.846578in}{1.013590in}}%
\pgfpathlineto{\pgfqpoint{0.862624in}{0.999025in}}%
\pgfpathlineto{\pgfqpoint{0.871023in}{1.003002in}}%
\pgfpathlineto{\pgfqpoint{0.879388in}{1.007429in}}%
\pgfpathlineto{\pgfqpoint{0.863382in}{1.021984in}}%
\pgfpathlineto{\pgfqpoint{0.847388in}{1.036985in}}%
\pgfpathlineto{\pgfqpoint{0.838983in}{1.032582in}}%
\pgfpathlineto{\pgfqpoint{0.830545in}{1.028602in}}%
\pgfpathclose%
\pgfusepath{fill}%
\end{pgfscope}%
\begin{pgfscope}%
\pgfpathrectangle{\pgfqpoint{0.407592in}{0.222904in}}{\pgfqpoint{1.411163in}{1.411163in}}%
\pgfusepath{clip}%
\pgfsetbuttcap%
\pgfsetroundjoin%
\definecolor{currentfill}{rgb}{0.909804,0.698039,0.423529}%
\pgfsetfillcolor{currentfill}%
\pgfsetfillopacity{0.120000}%
\pgfsetlinewidth{0.000000pt}%
\definecolor{currentstroke}{rgb}{0.000000,0.000000,0.000000}%
\pgfsetstrokecolor{currentstroke}%
\pgfsetdash{}{0pt}%
\pgfpathmoveto{\pgfqpoint{0.603879in}{1.219459in}}%
\pgfpathlineto{\pgfqpoint{0.619873in}{1.201403in}}%
\pgfpathlineto{\pgfqpoint{0.635876in}{1.183500in}}%
\pgfpathlineto{\pgfqpoint{0.644664in}{1.186175in}}%
\pgfpathlineto{\pgfqpoint{0.653415in}{1.189160in}}%
\pgfpathlineto{\pgfqpoint{0.637454in}{1.207098in}}%
\pgfpathlineto{\pgfqpoint{0.621502in}{1.225186in}}%
\pgfpathlineto{\pgfqpoint{0.612709in}{1.222176in}}%
\pgfpathlineto{\pgfqpoint{0.603879in}{1.219459in}}%
\pgfpathclose%
\pgfusepath{fill}%
\end{pgfscope}%
\begin{pgfscope}%
\pgfpathrectangle{\pgfqpoint{0.407592in}{0.222904in}}{\pgfqpoint{1.411163in}{1.411163in}}%
\pgfusepath{clip}%
\pgfsetbuttcap%
\pgfsetroundjoin%
\definecolor{currentfill}{rgb}{0.909804,0.698039,0.423529}%
\pgfsetfillcolor{currentfill}%
\pgfsetfillopacity{0.120000}%
\pgfsetlinewidth{0.000000pt}%
\definecolor{currentstroke}{rgb}{0.000000,0.000000,0.000000}%
\pgfsetstrokecolor{currentstroke}%
\pgfsetdash{}{0pt}%
\pgfpathmoveto{\pgfqpoint{1.169555in}{0.934259in}}%
\pgfpathlineto{\pgfqpoint{1.185826in}{0.931960in}}%
\pgfpathlineto{\pgfqpoint{1.202143in}{0.930612in}}%
\pgfpathlineto{\pgfqpoint{1.210034in}{0.939970in}}%
\pgfpathlineto{\pgfqpoint{1.217906in}{0.949774in}}%
\pgfpathlineto{\pgfqpoint{1.201609in}{0.951171in}}%
\pgfpathlineto{\pgfqpoint{1.185359in}{0.953490in}}%
\pgfpathlineto{\pgfqpoint{1.177467in}{0.943650in}}%
\pgfpathlineto{\pgfqpoint{1.169555in}{0.934259in}}%
\pgfpathclose%
\pgfusepath{fill}%
\end{pgfscope}%
\begin{pgfscope}%
\pgfpathrectangle{\pgfqpoint{0.407592in}{0.222904in}}{\pgfqpoint{1.411163in}{1.411163in}}%
\pgfusepath{clip}%
\pgfsetbuttcap%
\pgfsetroundjoin%
\definecolor{currentfill}{rgb}{0.909804,0.698039,0.423529}%
\pgfsetfillcolor{currentfill}%
\pgfsetfillopacity{0.120000}%
\pgfsetlinewidth{0.000000pt}%
\definecolor{currentstroke}{rgb}{0.000000,0.000000,0.000000}%
\pgfsetstrokecolor{currentstroke}%
\pgfsetdash{}{0pt}%
\pgfpathmoveto{\pgfqpoint{1.121212in}{0.924452in}}%
\pgfpathlineto{\pgfqpoint{1.137422in}{0.920185in}}%
\pgfpathlineto{\pgfqpoint{1.153672in}{0.916911in}}%
\pgfpathlineto{\pgfqpoint{1.161624in}{0.925338in}}%
\pgfpathlineto{\pgfqpoint{1.169555in}{0.934259in}}%
\pgfpathlineto{\pgfqpoint{1.153328in}{0.937520in}}%
\pgfpathlineto{\pgfqpoint{1.137142in}{0.941748in}}%
\pgfpathlineto{\pgfqpoint{1.129188in}{0.932854in}}%
\pgfpathlineto{\pgfqpoint{1.121212in}{0.924452in}}%
\pgfpathclose%
\pgfusepath{fill}%
\end{pgfscope}%
\begin{pgfscope}%
\pgfpathrectangle{\pgfqpoint{0.407592in}{0.222904in}}{\pgfqpoint{1.411163in}{1.411163in}}%
\pgfusepath{clip}%
\pgfsetbuttcap%
\pgfsetroundjoin%
\definecolor{currentfill}{rgb}{0.909804,0.698039,0.423529}%
\pgfsetfillcolor{currentfill}%
\pgfsetfillopacity{0.120000}%
\pgfsetlinewidth{0.000000pt}%
\definecolor{currentstroke}{rgb}{0.000000,0.000000,0.000000}%
\pgfsetstrokecolor{currentstroke}%
\pgfsetdash{}{0pt}%
\pgfpathmoveto{\pgfqpoint{1.378773in}{1.047207in}}%
\pgfpathlineto{\pgfqpoint{1.395393in}{1.050527in}}%
\pgfpathlineto{\pgfqpoint{1.412077in}{1.054448in}}%
\pgfpathlineto{\pgfqpoint{1.419742in}{1.066533in}}%
\pgfpathlineto{\pgfqpoint{1.427394in}{1.078857in}}%
\pgfpathlineto{\pgfqpoint{1.410727in}{1.075181in}}%
\pgfpathlineto{\pgfqpoint{1.394123in}{1.072094in}}%
\pgfpathlineto{\pgfqpoint{1.386455in}{1.059531in}}%
\pgfpathlineto{\pgfqpoint{1.378773in}{1.047207in}}%
\pgfpathclose%
\pgfusepath{fill}%
\end{pgfscope}%
\begin{pgfscope}%
\pgfpathrectangle{\pgfqpoint{0.407592in}{0.222904in}}{\pgfqpoint{1.411163in}{1.411163in}}%
\pgfusepath{clip}%
\pgfsetbuttcap%
\pgfsetroundjoin%
\definecolor{currentfill}{rgb}{0.909804,0.698039,0.423529}%
\pgfsetfillcolor{currentfill}%
\pgfsetfillopacity{0.120000}%
\pgfsetlinewidth{0.000000pt}%
\definecolor{currentstroke}{rgb}{0.000000,0.000000,0.000000}%
\pgfsetstrokecolor{currentstroke}%
\pgfsetdash{}{0pt}%
\pgfpathmoveto{\pgfqpoint{1.491333in}{1.138269in}}%
\pgfpathlineto{\pgfqpoint{1.508161in}{1.143015in}}%
\pgfpathlineto{\pgfqpoint{1.525056in}{1.148233in}}%
\pgfpathlineto{\pgfqpoint{1.532611in}{1.161047in}}%
\pgfpathlineto{\pgfqpoint{1.540154in}{1.174033in}}%
\pgfpathlineto{\pgfqpoint{1.523277in}{1.169078in}}%
\pgfpathlineto{\pgfqpoint{1.506465in}{1.164588in}}%
\pgfpathlineto{\pgfqpoint{1.498905in}{1.151344in}}%
\pgfpathlineto{\pgfqpoint{1.491333in}{1.138269in}}%
\pgfpathclose%
\pgfusepath{fill}%
\end{pgfscope}%
\begin{pgfscope}%
\pgfpathrectangle{\pgfqpoint{0.407592in}{0.222904in}}{\pgfqpoint{1.411163in}{1.411163in}}%
\pgfusepath{clip}%
\pgfsetbuttcap%
\pgfsetroundjoin%
\definecolor{currentfill}{rgb}{0.909804,0.698039,0.423529}%
\pgfsetfillcolor{currentfill}%
\pgfsetfillopacity{0.120000}%
\pgfsetlinewidth{0.000000pt}%
\definecolor{currentstroke}{rgb}{0.000000,0.000000,0.000000}%
\pgfsetstrokecolor{currentstroke}%
\pgfsetdash{}{0pt}%
\pgfpathmoveto{\pgfqpoint{0.749382in}{1.085557in}}%
\pgfpathlineto{\pgfqpoint{0.765412in}{1.069147in}}%
\pgfpathlineto{\pgfqpoint{0.781452in}{1.053047in}}%
\pgfpathlineto{\pgfqpoint{0.790000in}{1.056233in}}%
\pgfpathlineto{\pgfqpoint{0.798513in}{1.059814in}}%
\pgfpathlineto{\pgfqpoint{0.782515in}{1.075944in}}%
\pgfpathlineto{\pgfqpoint{0.766526in}{1.092380in}}%
\pgfpathlineto{\pgfqpoint{0.757972in}{1.088783in}}%
\pgfpathlineto{\pgfqpoint{0.749382in}{1.085557in}}%
\pgfpathclose%
\pgfusepath{fill}%
\end{pgfscope}%
\begin{pgfscope}%
\pgfpathrectangle{\pgfqpoint{0.407592in}{0.222904in}}{\pgfqpoint{1.411163in}{1.411163in}}%
\pgfusepath{clip}%
\pgfsetbuttcap%
\pgfsetroundjoin%
\definecolor{currentfill}{rgb}{0.909804,0.698039,0.423529}%
\pgfsetfillcolor{currentfill}%
\pgfsetfillopacity{0.120000}%
\pgfsetlinewidth{0.000000pt}%
\definecolor{currentstroke}{rgb}{0.000000,0.000000,0.000000}%
\pgfsetstrokecolor{currentstroke}%
\pgfsetdash{}{0pt}%
\pgfpathmoveto{\pgfqpoint{1.217906in}{0.949774in}}%
\pgfpathlineto{\pgfqpoint{1.234252in}{0.949277in}}%
\pgfpathlineto{\pgfqpoint{1.250651in}{0.949655in}}%
\pgfpathlineto{\pgfqpoint{1.258486in}{0.959782in}}%
\pgfpathlineto{\pgfqpoint{1.266304in}{0.970304in}}%
\pgfpathlineto{\pgfqpoint{1.249924in}{0.970040in}}%
\pgfpathlineto{\pgfqpoint{1.233597in}{0.970624in}}%
\pgfpathlineto{\pgfqpoint{1.225760in}{0.959999in}}%
\pgfpathlineto{\pgfqpoint{1.217906in}{0.949774in}}%
\pgfpathclose%
\pgfusepath{fill}%
\end{pgfscope}%
\begin{pgfscope}%
\pgfpathrectangle{\pgfqpoint{0.407592in}{0.222904in}}{\pgfqpoint{1.411163in}{1.411163in}}%
\pgfusepath{clip}%
\pgfsetbuttcap%
\pgfsetroundjoin%
\definecolor{currentfill}{rgb}{0.909804,0.698039,0.423529}%
\pgfsetfillcolor{currentfill}%
\pgfsetfillopacity{0.120000}%
\pgfsetlinewidth{0.000000pt}%
\definecolor{currentstroke}{rgb}{0.000000,0.000000,0.000000}%
\pgfsetstrokecolor{currentstroke}%
\pgfsetdash{}{0pt}%
\pgfpathmoveto{\pgfqpoint{1.072829in}{0.920814in}}%
\pgfpathlineto{\pgfqpoint{1.088994in}{0.914513in}}%
\pgfpathlineto{\pgfqpoint{1.105193in}{0.909202in}}%
\pgfpathlineto{\pgfqpoint{1.113214in}{0.916562in}}%
\pgfpathlineto{\pgfqpoint{1.121212in}{0.924452in}}%
\pgfpathlineto{\pgfqpoint{1.105040in}{0.929706in}}%
\pgfpathlineto{\pgfqpoint{1.088902in}{0.935933in}}%
\pgfpathlineto{\pgfqpoint{1.080878in}{0.928112in}}%
\pgfpathlineto{\pgfqpoint{1.072829in}{0.920814in}}%
\pgfpathclose%
\pgfusepath{fill}%
\end{pgfscope}%
\begin{pgfscope}%
\pgfpathrectangle{\pgfqpoint{0.407592in}{0.222904in}}{\pgfqpoint{1.411163in}{1.411163in}}%
\pgfusepath{clip}%
\pgfsetbuttcap%
\pgfsetroundjoin%
\definecolor{currentfill}{rgb}{0.909804,0.698039,0.423529}%
\pgfsetfillcolor{currentfill}%
\pgfsetfillopacity{0.120000}%
\pgfsetlinewidth{0.000000pt}%
\definecolor{currentstroke}{rgb}{0.000000,0.000000,0.000000}%
\pgfsetstrokecolor{currentstroke}%
\pgfsetdash{}{0pt}%
\pgfpathmoveto{\pgfqpoint{0.943557in}{0.954581in}}%
\pgfpathlineto{\pgfqpoint{0.959640in}{0.942964in}}%
\pgfpathlineto{\pgfqpoint{0.975742in}{0.932092in}}%
\pgfpathlineto{\pgfqpoint{0.983963in}{0.937106in}}%
\pgfpathlineto{\pgfqpoint{0.992156in}{0.942647in}}%
\pgfpathlineto{\pgfqpoint{0.976089in}{0.953465in}}%
\pgfpathlineto{\pgfqpoint{0.960042in}{0.965026in}}%
\pgfpathlineto{\pgfqpoint{0.951814in}{0.959552in}}%
\pgfpathlineto{\pgfqpoint{0.943557in}{0.954581in}}%
\pgfpathclose%
\pgfusepath{fill}%
\end{pgfscope}%
\begin{pgfscope}%
\pgfpathrectangle{\pgfqpoint{0.407592in}{0.222904in}}{\pgfqpoint{1.411163in}{1.411163in}}%
\pgfusepath{clip}%
\pgfsetbuttcap%
\pgfsetroundjoin%
\definecolor{currentfill}{rgb}{0.909804,0.698039,0.423529}%
\pgfsetfillcolor{currentfill}%
\pgfsetfillopacity{0.120000}%
\pgfsetlinewidth{0.000000pt}%
\definecolor{currentstroke}{rgb}{0.000000,0.000000,0.000000}%
\pgfsetstrokecolor{currentstroke}%
\pgfsetdash{}{0pt}%
\pgfpathmoveto{\pgfqpoint{0.635876in}{1.183500in}}%
\pgfpathlineto{\pgfqpoint{0.651888in}{1.165765in}}%
\pgfpathlineto{\pgfqpoint{0.667910in}{1.148213in}}%
\pgfpathlineto{\pgfqpoint{0.676656in}{1.150838in}}%
\pgfpathlineto{\pgfqpoint{0.685366in}{1.153789in}}%
\pgfpathlineto{\pgfqpoint{0.669386in}{1.171385in}}%
\pgfpathlineto{\pgfqpoint{0.653415in}{1.189160in}}%
\pgfpathlineto{\pgfqpoint{0.644664in}{1.186175in}}%
\pgfpathlineto{\pgfqpoint{0.635876in}{1.183500in}}%
\pgfpathclose%
\pgfusepath{fill}%
\end{pgfscope}%
\begin{pgfscope}%
\pgfpathrectangle{\pgfqpoint{0.407592in}{0.222904in}}{\pgfqpoint{1.411163in}{1.411163in}}%
\pgfusepath{clip}%
\pgfsetbuttcap%
\pgfsetroundjoin%
\definecolor{currentfill}{rgb}{0.909804,0.698039,0.423529}%
\pgfsetfillcolor{currentfill}%
\pgfsetfillopacity{0.120000}%
\pgfsetlinewidth{0.000000pt}%
\definecolor{currentstroke}{rgb}{0.000000,0.000000,0.000000}%
\pgfsetstrokecolor{currentstroke}%
\pgfsetdash{}{0pt}%
\pgfpathmoveto{\pgfqpoint{0.862624in}{0.999025in}}%
\pgfpathlineto{\pgfqpoint{0.878683in}{0.984951in}}%
\pgfpathlineto{\pgfqpoint{0.894755in}{0.971414in}}%
\pgfpathlineto{\pgfqpoint{0.903115in}{0.975386in}}%
\pgfpathlineto{\pgfqpoint{0.911442in}{0.979834in}}%
\pgfpathlineto{\pgfqpoint{0.895408in}{0.993364in}}%
\pgfpathlineto{\pgfqpoint{0.879388in}{1.007429in}}%
\pgfpathlineto{\pgfqpoint{0.871023in}{1.003002in}}%
\pgfpathlineto{\pgfqpoint{0.862624in}{0.999025in}}%
\pgfpathclose%
\pgfusepath{fill}%
\end{pgfscope}%
\begin{pgfscope}%
\pgfpathrectangle{\pgfqpoint{0.407592in}{0.222904in}}{\pgfqpoint{1.411163in}{1.411163in}}%
\pgfusepath{clip}%
\pgfsetbuttcap%
\pgfsetroundjoin%
\definecolor{currentfill}{rgb}{0.909804,0.698039,0.423529}%
\pgfsetfillcolor{currentfill}%
\pgfsetfillopacity{0.120000}%
\pgfsetlinewidth{0.000000pt}%
\definecolor{currentstroke}{rgb}{0.000000,0.000000,0.000000}%
\pgfsetstrokecolor{currentstroke}%
\pgfsetdash{}{0pt}%
\pgfpathmoveto{\pgfqpoint{1.266304in}{0.970304in}}%
\pgfpathlineto{\pgfqpoint{1.282739in}{0.971386in}}%
\pgfpathlineto{\pgfqpoint{1.299230in}{0.973254in}}%
\pgfpathlineto{\pgfqpoint{1.307015in}{0.983990in}}%
\pgfpathlineto{\pgfqpoint{1.314784in}{0.995073in}}%
\pgfpathlineto{\pgfqpoint{1.298310in}{0.993377in}}%
\pgfpathlineto{\pgfqpoint{1.281894in}{0.992444in}}%
\pgfpathlineto{\pgfqpoint{1.274107in}{0.981198in}}%
\pgfpathlineto{\pgfqpoint{1.266304in}{0.970304in}}%
\pgfpathclose%
\pgfusepath{fill}%
\end{pgfscope}%
\begin{pgfscope}%
\pgfpathrectangle{\pgfqpoint{0.407592in}{0.222904in}}{\pgfqpoint{1.411163in}{1.411163in}}%
\pgfusepath{clip}%
\pgfsetbuttcap%
\pgfsetroundjoin%
\definecolor{currentfill}{rgb}{0.909804,0.698039,0.423529}%
\pgfsetfillcolor{currentfill}%
\pgfsetfillopacity{0.120000}%
\pgfsetlinewidth{0.000000pt}%
\definecolor{currentstroke}{rgb}{0.000000,0.000000,0.000000}%
\pgfsetstrokecolor{currentstroke}%
\pgfsetdash{}{0pt}%
\pgfpathmoveto{\pgfqpoint{1.024355in}{0.923453in}}%
\pgfpathlineto{\pgfqpoint{1.040492in}{0.915178in}}%
\pgfpathlineto{\pgfqpoint{1.056656in}{0.907844in}}%
\pgfpathlineto{\pgfqpoint{1.064756in}{0.914054in}}%
\pgfpathlineto{\pgfqpoint{1.072829in}{0.920814in}}%
\pgfpathlineto{\pgfqpoint{1.056695in}{0.928081in}}%
\pgfpathlineto{\pgfqpoint{1.040590in}{0.936278in}}%
\pgfpathlineto{\pgfqpoint{1.032486in}{0.929598in}}%
\pgfpathlineto{\pgfqpoint{1.024355in}{0.923453in}}%
\pgfpathclose%
\pgfusepath{fill}%
\end{pgfscope}%
\begin{pgfscope}%
\pgfpathrectangle{\pgfqpoint{0.407592in}{0.222904in}}{\pgfqpoint{1.411163in}{1.411163in}}%
\pgfusepath{clip}%
\pgfsetbuttcap%
\pgfsetroundjoin%
\definecolor{currentfill}{rgb}{0.909804,0.698039,0.423529}%
\pgfsetfillcolor{currentfill}%
\pgfsetfillopacity{0.120000}%
\pgfsetlinewidth{0.000000pt}%
\definecolor{currentstroke}{rgb}{0.000000,0.000000,0.000000}%
\pgfsetstrokecolor{currentstroke}%
\pgfsetdash{}{0pt}%
\pgfpathmoveto{\pgfqpoint{1.427394in}{1.078857in}}%
\pgfpathlineto{\pgfqpoint{1.444124in}{1.083094in}}%
\pgfpathlineto{\pgfqpoint{1.460920in}{1.087866in}}%
\pgfpathlineto{\pgfqpoint{1.468542in}{1.100158in}}%
\pgfpathlineto{\pgfqpoint{1.476152in}{1.112664in}}%
\pgfpathlineto{\pgfqpoint{1.459373in}{1.108154in}}%
\pgfpathlineto{\pgfqpoint{1.442659in}{1.104169in}}%
\pgfpathlineto{\pgfqpoint{1.435033in}{1.091407in}}%
\pgfpathlineto{\pgfqpoint{1.427394in}{1.078857in}}%
\pgfpathclose%
\pgfusepath{fill}%
\end{pgfscope}%
\begin{pgfscope}%
\pgfpathrectangle{\pgfqpoint{0.407592in}{0.222904in}}{\pgfqpoint{1.411163in}{1.411163in}}%
\pgfusepath{clip}%
\pgfsetbuttcap%
\pgfsetroundjoin%
\definecolor{currentfill}{rgb}{0.909804,0.698039,0.423529}%
\pgfsetfillcolor{currentfill}%
\pgfsetfillopacity{0.120000}%
\pgfsetlinewidth{0.000000pt}%
\definecolor{currentstroke}{rgb}{0.000000,0.000000,0.000000}%
\pgfsetstrokecolor{currentstroke}%
\pgfsetdash{}{0pt}%
\pgfpathmoveto{\pgfqpoint{0.781452in}{1.053047in}}%
\pgfpathlineto{\pgfqpoint{0.797502in}{1.037288in}}%
\pgfpathlineto{\pgfqpoint{0.813564in}{1.021904in}}%
\pgfpathlineto{\pgfqpoint{0.822072in}{1.025044in}}%
\pgfpathlineto{\pgfqpoint{0.830545in}{1.028602in}}%
\pgfpathlineto{\pgfqpoint{0.814523in}{1.044023in}}%
\pgfpathlineto{\pgfqpoint{0.798513in}{1.059814in}}%
\pgfpathlineto{\pgfqpoint{0.790000in}{1.056233in}}%
\pgfpathlineto{\pgfqpoint{0.781452in}{1.053047in}}%
\pgfpathclose%
\pgfusepath{fill}%
\end{pgfscope}%
\begin{pgfscope}%
\pgfpathrectangle{\pgfqpoint{0.407592in}{0.222904in}}{\pgfqpoint{1.411163in}{1.411163in}}%
\pgfusepath{clip}%
\pgfsetbuttcap%
\pgfsetroundjoin%
\definecolor{currentfill}{rgb}{0.909804,0.698039,0.423529}%
\pgfsetfillcolor{currentfill}%
\pgfsetfillopacity{0.120000}%
\pgfsetlinewidth{0.000000pt}%
\definecolor{currentstroke}{rgb}{0.000000,0.000000,0.000000}%
\pgfsetstrokecolor{currentstroke}%
\pgfsetdash{}{0pt}%
\pgfpathmoveto{\pgfqpoint{0.667910in}{1.148213in}}%
\pgfpathlineto{\pgfqpoint{0.683941in}{1.130862in}}%
\pgfpathlineto{\pgfqpoint{0.699982in}{1.113730in}}%
\pgfpathlineto{\pgfqpoint{0.708687in}{1.116293in}}%
\pgfpathlineto{\pgfqpoint{0.717354in}{1.119202in}}%
\pgfpathlineto{\pgfqpoint{0.701355in}{1.136388in}}%
\pgfpathlineto{\pgfqpoint{0.685366in}{1.153789in}}%
\pgfpathlineto{\pgfqpoint{0.676656in}{1.150838in}}%
\pgfpathlineto{\pgfqpoint{0.667910in}{1.148213in}}%
\pgfpathclose%
\pgfusepath{fill}%
\end{pgfscope}%
\begin{pgfscope}%
\pgfpathrectangle{\pgfqpoint{0.407592in}{0.222904in}}{\pgfqpoint{1.411163in}{1.411163in}}%
\pgfusepath{clip}%
\pgfsetbuttcap%
\pgfsetroundjoin%
\definecolor{currentfill}{rgb}{0.909804,0.698039,0.423529}%
\pgfsetfillcolor{currentfill}%
\pgfsetfillopacity{0.120000}%
\pgfsetlinewidth{0.000000pt}%
\definecolor{currentstroke}{rgb}{0.000000,0.000000,0.000000}%
\pgfsetstrokecolor{currentstroke}%
\pgfsetdash{}{0pt}%
\pgfpathmoveto{\pgfqpoint{1.314784in}{0.995073in}}%
\pgfpathlineto{\pgfqpoint{1.331316in}{0.997499in}}%
\pgfpathlineto{\pgfqpoint{1.347909in}{1.000622in}}%
\pgfpathlineto{\pgfqpoint{1.355646in}{1.011827in}}%
\pgfpathlineto{\pgfqpoint{1.363369in}{1.023338in}}%
\pgfpathlineto{\pgfqpoint{1.346793in}{1.020431in}}%
\pgfpathlineto{\pgfqpoint{1.330278in}{1.018203in}}%
\pgfpathlineto{\pgfqpoint{1.322538in}{1.006484in}}%
\pgfpathlineto{\pgfqpoint{1.314784in}{0.995073in}}%
\pgfpathclose%
\pgfusepath{fill}%
\end{pgfscope}%
\begin{pgfscope}%
\pgfpathrectangle{\pgfqpoint{0.407592in}{0.222904in}}{\pgfqpoint{1.411163in}{1.411163in}}%
\pgfusepath{clip}%
\pgfsetbuttcap%
\pgfsetroundjoin%
\definecolor{currentfill}{rgb}{0.909804,0.698039,0.423529}%
\pgfsetfillcolor{currentfill}%
\pgfsetfillopacity{0.120000}%
\pgfsetlinewidth{0.000000pt}%
\definecolor{currentstroke}{rgb}{0.000000,0.000000,0.000000}%
\pgfsetstrokecolor{currentstroke}%
\pgfsetdash{}{0pt}%
\pgfpathmoveto{\pgfqpoint{1.540154in}{1.174033in}}%
\pgfpathlineto{\pgfqpoint{1.557099in}{1.179432in}}%
\pgfpathlineto{\pgfqpoint{1.574112in}{1.185256in}}%
\pgfpathlineto{\pgfqpoint{1.581626in}{1.198139in}}%
\pgfpathlineto{\pgfqpoint{1.564622in}{1.192448in}}%
\pgfpathlineto{\pgfqpoint{1.547685in}{1.187181in}}%
\pgfusepath{fill}%
\end{pgfscope}%
\begin{pgfscope}%
\pgfpathrectangle{\pgfqpoint{0.407592in}{0.222904in}}{\pgfqpoint{1.411163in}{1.411163in}}%
\pgfusepath{clip}%
\pgfsetbuttcap%
\pgfsetroundjoin%
\definecolor{currentfill}{rgb}{0.909804,0.698039,0.423529}%
\pgfsetfillcolor{currentfill}%
\pgfsetfillopacity{0.120000}%
\pgfsetlinewidth{0.000000pt}%
\definecolor{currentstroke}{rgb}{0.000000,0.000000,0.000000}%
\pgfsetstrokecolor{currentstroke}%
\pgfsetdash{}{0pt}%
\pgfpathmoveto{\pgfqpoint{0.894755in}{0.971414in}}%
\pgfpathlineto{\pgfqpoint{0.910843in}{0.958463in}}%
\pgfpathlineto{\pgfqpoint{0.926946in}{0.946149in}}%
\pgfpathlineto{\pgfqpoint{0.935267in}{0.950114in}}%
\pgfpathlineto{\pgfqpoint{0.943557in}{0.954581in}}%
\pgfpathlineto{\pgfqpoint{0.927491in}{0.966889in}}%
\pgfpathlineto{\pgfqpoint{0.911442in}{0.979834in}}%
\pgfpathlineto{\pgfqpoint{0.903115in}{0.975386in}}%
\pgfpathlineto{\pgfqpoint{0.894755in}{0.971414in}}%
\pgfpathclose%
\pgfusepath{fill}%
\end{pgfscope}%
\begin{pgfscope}%
\pgfpathrectangle{\pgfqpoint{0.407592in}{0.222904in}}{\pgfqpoint{1.411163in}{1.411163in}}%
\pgfusepath{clip}%
\pgfsetbuttcap%
\pgfsetroundjoin%
\definecolor{currentfill}{rgb}{0.909804,0.698039,0.423529}%
\pgfsetfillcolor{currentfill}%
\pgfsetfillopacity{0.120000}%
\pgfsetlinewidth{0.000000pt}%
\definecolor{currentstroke}{rgb}{0.000000,0.000000,0.000000}%
\pgfsetstrokecolor{currentstroke}%
\pgfsetdash{}{0pt}%
\pgfpathmoveto{\pgfqpoint{0.975742in}{0.932092in}}%
\pgfpathlineto{\pgfqpoint{0.991864in}{0.922019in}}%
\pgfpathlineto{\pgfqpoint{1.008010in}{0.912799in}}%
\pgfpathlineto{\pgfqpoint{1.016197in}{0.917852in}}%
\pgfpathlineto{\pgfqpoint{1.024355in}{0.923453in}}%
\pgfpathlineto{\pgfqpoint{1.008244in}{0.932626in}}%
\pgfpathlineto{\pgfqpoint{0.992156in}{0.942647in}}%
\pgfpathlineto{\pgfqpoint{0.983963in}{0.937106in}}%
\pgfpathlineto{\pgfqpoint{0.975742in}{0.932092in}}%
\pgfpathclose%
\pgfusepath{fill}%
\end{pgfscope}%
\begin{pgfscope}%
\pgfpathrectangle{\pgfqpoint{0.407592in}{0.222904in}}{\pgfqpoint{1.411163in}{1.411163in}}%
\pgfusepath{clip}%
\pgfsetbuttcap%
\pgfsetroundjoin%
\definecolor{currentfill}{rgb}{0.909804,0.698039,0.423529}%
\pgfsetfillcolor{currentfill}%
\pgfsetfillopacity{0.120000}%
\pgfsetlinewidth{0.000000pt}%
\definecolor{currentstroke}{rgb}{0.000000,0.000000,0.000000}%
\pgfsetstrokecolor{currentstroke}%
\pgfsetdash{}{0pt}%
\pgfpathmoveto{\pgfqpoint{1.476152in}{1.112664in}}%
\pgfpathlineto{\pgfqpoint{1.492997in}{1.117674in}}%
\pgfpathlineto{\pgfqpoint{1.509909in}{1.123162in}}%
\pgfpathlineto{\pgfqpoint{1.517489in}{1.135601in}}%
\pgfpathlineto{\pgfqpoint{1.525056in}{1.148233in}}%
\pgfpathlineto{\pgfqpoint{1.508161in}{1.143015in}}%
\pgfpathlineto{\pgfqpoint{1.491333in}{1.138269in}}%
\pgfpathlineto{\pgfqpoint{1.483749in}{1.125371in}}%
\pgfpathlineto{\pgfqpoint{1.476152in}{1.112664in}}%
\pgfpathclose%
\pgfusepath{fill}%
\end{pgfscope}%
\begin{pgfscope}%
\pgfpathrectangle{\pgfqpoint{0.407592in}{0.222904in}}{\pgfqpoint{1.411163in}{1.411163in}}%
\pgfusepath{clip}%
\pgfsetbuttcap%
\pgfsetroundjoin%
\definecolor{currentfill}{rgb}{0.909804,0.698039,0.423529}%
\pgfsetfillcolor{currentfill}%
\pgfsetfillopacity{0.120000}%
\pgfsetlinewidth{0.000000pt}%
\definecolor{currentstroke}{rgb}{0.000000,0.000000,0.000000}%
\pgfsetstrokecolor{currentstroke}%
\pgfsetdash{}{0pt}%
\pgfpathmoveto{\pgfqpoint{1.363369in}{1.023338in}}%
\pgfpathlineto{\pgfqpoint{1.380006in}{1.026891in}}%
\pgfpathlineto{\pgfqpoint{1.396706in}{1.031059in}}%
\pgfpathlineto{\pgfqpoint{1.404398in}{1.042618in}}%
\pgfpathlineto{\pgfqpoint{1.412077in}{1.054448in}}%
\pgfpathlineto{\pgfqpoint{1.395393in}{1.050527in}}%
\pgfpathlineto{\pgfqpoint{1.378773in}{1.047207in}}%
\pgfpathlineto{\pgfqpoint{1.371078in}{1.035137in}}%
\pgfpathlineto{\pgfqpoint{1.363369in}{1.023338in}}%
\pgfpathclose%
\pgfusepath{fill}%
\end{pgfscope}%
\begin{pgfscope}%
\pgfpathrectangle{\pgfqpoint{0.407592in}{0.222904in}}{\pgfqpoint{1.411163in}{1.411163in}}%
\pgfusepath{clip}%
\pgfsetbuttcap%
\pgfsetroundjoin%
\definecolor{currentfill}{rgb}{0.909804,0.698039,0.423529}%
\pgfsetfillcolor{currentfill}%
\pgfsetfillopacity{0.120000}%
\pgfsetlinewidth{0.000000pt}%
\definecolor{currentstroke}{rgb}{0.000000,0.000000,0.000000}%
\pgfsetstrokecolor{currentstroke}%
\pgfsetdash{}{0pt}%
\pgfpathmoveto{\pgfqpoint{0.699982in}{1.113730in}}%
\pgfpathlineto{\pgfqpoint{0.716033in}{1.096838in}}%
\pgfpathlineto{\pgfqpoint{0.732094in}{1.080209in}}%
\pgfpathlineto{\pgfqpoint{0.740756in}{1.082700in}}%
\pgfpathlineto{\pgfqpoint{0.749382in}{1.085557in}}%
\pgfpathlineto{\pgfqpoint{0.733363in}{1.102251in}}%
\pgfpathlineto{\pgfqpoint{0.717354in}{1.119202in}}%
\pgfpathlineto{\pgfqpoint{0.708687in}{1.116293in}}%
\pgfpathlineto{\pgfqpoint{0.699982in}{1.113730in}}%
\pgfpathclose%
\pgfusepath{fill}%
\end{pgfscope}%
\begin{pgfscope}%
\pgfpathrectangle{\pgfqpoint{0.407592in}{0.222904in}}{\pgfqpoint{1.411163in}{1.411163in}}%
\pgfusepath{clip}%
\pgfsetbuttcap%
\pgfsetroundjoin%
\definecolor{currentfill}{rgb}{0.909804,0.698039,0.423529}%
\pgfsetfillcolor{currentfill}%
\pgfsetfillopacity{0.120000}%
\pgfsetlinewidth{0.000000pt}%
\definecolor{currentstroke}{rgb}{0.000000,0.000000,0.000000}%
\pgfsetstrokecolor{currentstroke}%
\pgfsetdash{}{0pt}%
\pgfpathmoveto{\pgfqpoint{0.813564in}{1.021904in}}%
\pgfpathlineto{\pgfqpoint{0.829638in}{1.006931in}}%
\pgfpathlineto{\pgfqpoint{0.845724in}{0.992408in}}%
\pgfpathlineto{\pgfqpoint{0.854191in}{0.995495in}}%
\pgfpathlineto{\pgfqpoint{0.862624in}{0.999025in}}%
\pgfpathlineto{\pgfqpoint{0.846578in}{1.013590in}}%
\pgfpathlineto{\pgfqpoint{0.830545in}{1.028602in}}%
\pgfpathlineto{\pgfqpoint{0.822072in}{1.025044in}}%
\pgfpathlineto{\pgfqpoint{0.813564in}{1.021904in}}%
\pgfpathclose%
\pgfusepath{fill}%
\end{pgfscope}%
\begin{pgfscope}%
\pgfpathrectangle{\pgfqpoint{0.407592in}{0.222904in}}{\pgfqpoint{1.411163in}{1.411163in}}%
\pgfusepath{clip}%
\pgfsetbuttcap%
\pgfsetroundjoin%
\definecolor{currentfill}{rgb}{0.909804,0.698039,0.423529}%
\pgfsetfillcolor{currentfill}%
\pgfsetfillopacity{0.120000}%
\pgfsetlinewidth{0.000000pt}%
\definecolor{currentstroke}{rgb}{0.000000,0.000000,0.000000}%
\pgfsetstrokecolor{currentstroke}%
\pgfsetdash{}{0pt}%
\pgfpathmoveto{\pgfqpoint{1.153672in}{0.916911in}}%
\pgfpathlineto{\pgfqpoint{1.169966in}{0.914626in}}%
\pgfpathlineto{\pgfqpoint{1.186306in}{0.913316in}}%
\pgfpathlineto{\pgfqpoint{1.194234in}{0.921720in}}%
\pgfpathlineto{\pgfqpoint{1.202143in}{0.930612in}}%
\pgfpathlineto{\pgfqpoint{1.185826in}{0.931960in}}%
\pgfpathlineto{\pgfqpoint{1.169555in}{0.934259in}}%
\pgfpathlineto{\pgfqpoint{1.161624in}{0.925338in}}%
\pgfpathlineto{\pgfqpoint{1.153672in}{0.916911in}}%
\pgfpathclose%
\pgfusepath{fill}%
\end{pgfscope}%
\begin{pgfscope}%
\pgfpathrectangle{\pgfqpoint{0.407592in}{0.222904in}}{\pgfqpoint{1.411163in}{1.411163in}}%
\pgfusepath{clip}%
\pgfsetbuttcap%
\pgfsetroundjoin%
\definecolor{currentfill}{rgb}{0.909804,0.698039,0.423529}%
\pgfsetfillcolor{currentfill}%
\pgfsetfillopacity{0.120000}%
\pgfsetlinewidth{0.000000pt}%
\definecolor{currentstroke}{rgb}{0.000000,0.000000,0.000000}%
\pgfsetstrokecolor{currentstroke}%
\pgfsetdash{}{0pt}%
\pgfpathmoveto{\pgfqpoint{1.105193in}{0.909202in}}%
\pgfpathlineto{\pgfqpoint{1.121429in}{0.904899in}}%
\pgfpathlineto{\pgfqpoint{1.137705in}{0.901610in}}%
\pgfpathlineto{\pgfqpoint{1.145700in}{0.908995in}}%
\pgfpathlineto{\pgfqpoint{1.153672in}{0.916911in}}%
\pgfpathlineto{\pgfqpoint{1.137422in}{0.920185in}}%
\pgfpathlineto{\pgfqpoint{1.121212in}{0.924452in}}%
\pgfpathlineto{\pgfqpoint{1.113214in}{0.916562in}}%
\pgfpathlineto{\pgfqpoint{1.105193in}{0.909202in}}%
\pgfpathclose%
\pgfusepath{fill}%
\end{pgfscope}%
\begin{pgfscope}%
\pgfpathrectangle{\pgfqpoint{0.407592in}{0.222904in}}{\pgfqpoint{1.411163in}{1.411163in}}%
\pgfusepath{clip}%
\pgfsetbuttcap%
\pgfsetroundjoin%
\definecolor{currentfill}{rgb}{0.909804,0.698039,0.423529}%
\pgfsetfillcolor{currentfill}%
\pgfsetfillopacity{0.120000}%
\pgfsetlinewidth{0.000000pt}%
\definecolor{currentstroke}{rgb}{0.000000,0.000000,0.000000}%
\pgfsetstrokecolor{currentstroke}%
\pgfsetdash{}{0pt}%
\pgfpathmoveto{\pgfqpoint{1.202143in}{0.930612in}}%
\pgfpathlineto{\pgfqpoint{1.218510in}{0.930192in}}%
\pgfpathlineto{\pgfqpoint{1.234929in}{0.930672in}}%
\pgfpathlineto{\pgfqpoint{1.242799in}{0.939945in}}%
\pgfpathlineto{\pgfqpoint{1.250651in}{0.949655in}}%
\pgfpathlineto{\pgfqpoint{1.234252in}{0.949277in}}%
\pgfpathlineto{\pgfqpoint{1.217906in}{0.949774in}}%
\pgfpathlineto{\pgfqpoint{1.210034in}{0.939970in}}%
\pgfpathlineto{\pgfqpoint{1.202143in}{0.930612in}}%
\pgfpathclose%
\pgfusepath{fill}%
\end{pgfscope}%
\begin{pgfscope}%
\pgfpathrectangle{\pgfqpoint{0.407592in}{0.222904in}}{\pgfqpoint{1.411163in}{1.411163in}}%
\pgfusepath{clip}%
\pgfsetbuttcap%
\pgfsetroundjoin%
\definecolor{currentfill}{rgb}{0.909804,0.698039,0.423529}%
\pgfsetfillcolor{currentfill}%
\pgfsetfillopacity{0.120000}%
\pgfsetlinewidth{0.000000pt}%
\definecolor{currentstroke}{rgb}{0.000000,0.000000,0.000000}%
\pgfsetstrokecolor{currentstroke}%
\pgfsetdash{}{0pt}%
\pgfpathmoveto{\pgfqpoint{0.586104in}{1.214899in}}%
\pgfpathlineto{\pgfqpoint{0.602140in}{1.196904in}}%
\pgfpathlineto{\pgfqpoint{0.618185in}{1.179069in}}%
\pgfpathlineto{\pgfqpoint{0.627050in}{1.181132in}}%
\pgfpathlineto{\pgfqpoint{0.635876in}{1.183500in}}%
\pgfpathlineto{\pgfqpoint{0.619873in}{1.201403in}}%
\pgfpathlineto{\pgfqpoint{0.603879in}{1.219459in}}%
\pgfpathlineto{\pgfqpoint{0.595011in}{1.217035in}}%
\pgfpathlineto{\pgfqpoint{0.586104in}{1.214899in}}%
\pgfpathclose%
\pgfusepath{fill}%
\end{pgfscope}%
\begin{pgfscope}%
\pgfpathrectangle{\pgfqpoint{0.407592in}{0.222904in}}{\pgfqpoint{1.411163in}{1.411163in}}%
\pgfusepath{clip}%
\pgfsetbuttcap%
\pgfsetroundjoin%
\definecolor{currentfill}{rgb}{0.909804,0.698039,0.423529}%
\pgfsetfillcolor{currentfill}%
\pgfsetfillopacity{0.120000}%
\pgfsetlinewidth{0.000000pt}%
\definecolor{currentstroke}{rgb}{0.000000,0.000000,0.000000}%
\pgfsetstrokecolor{currentstroke}%
\pgfsetdash{}{0pt}%
\pgfpathmoveto{\pgfqpoint{1.250651in}{0.949655in}}%
\pgfpathlineto{\pgfqpoint{1.267104in}{0.950877in}}%
\pgfpathlineto{\pgfqpoint{1.283614in}{0.952905in}}%
\pgfpathlineto{\pgfqpoint{1.291430in}{0.962886in}}%
\pgfpathlineto{\pgfqpoint{1.299230in}{0.973254in}}%
\pgfpathlineto{\pgfqpoint{1.282739in}{0.971386in}}%
\pgfpathlineto{\pgfqpoint{1.266304in}{0.970304in}}%
\pgfpathlineto{\pgfqpoint{1.258486in}{0.959782in}}%
\pgfpathlineto{\pgfqpoint{1.250651in}{0.949655in}}%
\pgfpathclose%
\pgfusepath{fill}%
\end{pgfscope}%
\begin{pgfscope}%
\pgfpathrectangle{\pgfqpoint{0.407592in}{0.222904in}}{\pgfqpoint{1.411163in}{1.411163in}}%
\pgfusepath{clip}%
\pgfsetbuttcap%
\pgfsetroundjoin%
\definecolor{currentfill}{rgb}{0.909804,0.698039,0.423529}%
\pgfsetfillcolor{currentfill}%
\pgfsetfillopacity{0.120000}%
\pgfsetlinewidth{0.000000pt}%
\definecolor{currentstroke}{rgb}{0.000000,0.000000,0.000000}%
\pgfsetstrokecolor{currentstroke}%
\pgfsetdash{}{0pt}%
\pgfpathmoveto{\pgfqpoint{1.056656in}{0.907844in}}%
\pgfpathlineto{\pgfqpoint{1.072851in}{0.901485in}}%
\pgfpathlineto{\pgfqpoint{1.089079in}{0.896130in}}%
\pgfpathlineto{\pgfqpoint{1.097149in}{0.902387in}}%
\pgfpathlineto{\pgfqpoint{1.105193in}{0.909202in}}%
\pgfpathlineto{\pgfqpoint{1.088994in}{0.914513in}}%
\pgfpathlineto{\pgfqpoint{1.072829in}{0.920814in}}%
\pgfpathlineto{\pgfqpoint{1.064756in}{0.914054in}}%
\pgfpathlineto{\pgfqpoint{1.056656in}{0.907844in}}%
\pgfpathclose%
\pgfusepath{fill}%
\end{pgfscope}%
\begin{pgfscope}%
\pgfpathrectangle{\pgfqpoint{0.407592in}{0.222904in}}{\pgfqpoint{1.411163in}{1.411163in}}%
\pgfusepath{clip}%
\pgfsetbuttcap%
\pgfsetroundjoin%
\definecolor{currentfill}{rgb}{0.909804,0.698039,0.423529}%
\pgfsetfillcolor{currentfill}%
\pgfsetfillopacity{0.120000}%
\pgfsetlinewidth{0.000000pt}%
\definecolor{currentstroke}{rgb}{0.000000,0.000000,0.000000}%
\pgfsetstrokecolor{currentstroke}%
\pgfsetdash{}{0pt}%
\pgfpathmoveto{\pgfqpoint{1.412077in}{1.054448in}}%
\pgfpathlineto{\pgfqpoint{1.428824in}{1.058942in}}%
\pgfpathlineto{\pgfqpoint{1.445638in}{1.063977in}}%
\pgfpathlineto{\pgfqpoint{1.453286in}{1.075801in}}%
\pgfpathlineto{\pgfqpoint{1.460920in}{1.087866in}}%
\pgfpathlineto{\pgfqpoint{1.444124in}{1.083094in}}%
\pgfpathlineto{\pgfqpoint{1.427394in}{1.078857in}}%
\pgfpathlineto{\pgfqpoint{1.419742in}{1.066533in}}%
\pgfpathlineto{\pgfqpoint{1.412077in}{1.054448in}}%
\pgfpathclose%
\pgfusepath{fill}%
\end{pgfscope}%
\begin{pgfscope}%
\pgfpathrectangle{\pgfqpoint{0.407592in}{0.222904in}}{\pgfqpoint{1.411163in}{1.411163in}}%
\pgfusepath{clip}%
\pgfsetbuttcap%
\pgfsetroundjoin%
\definecolor{currentfill}{rgb}{0.909804,0.698039,0.423529}%
\pgfsetfillcolor{currentfill}%
\pgfsetfillopacity{0.120000}%
\pgfsetlinewidth{0.000000pt}%
\definecolor{currentstroke}{rgb}{0.000000,0.000000,0.000000}%
\pgfsetstrokecolor{currentstroke}%
\pgfsetdash{}{0pt}%
\pgfpathmoveto{\pgfqpoint{0.732094in}{1.080209in}}%
\pgfpathlineto{\pgfqpoint{0.748165in}{1.063868in}}%
\pgfpathlineto{\pgfqpoint{0.764246in}{1.047843in}}%
\pgfpathlineto{\pgfqpoint{0.772867in}{1.050252in}}%
\pgfpathlineto{\pgfqpoint{0.781452in}{1.053047in}}%
\pgfpathlineto{\pgfqpoint{0.765412in}{1.069147in}}%
\pgfpathlineto{\pgfqpoint{0.749382in}{1.085557in}}%
\pgfpathlineto{\pgfqpoint{0.740756in}{1.082700in}}%
\pgfpathlineto{\pgfqpoint{0.732094in}{1.080209in}}%
\pgfpathclose%
\pgfusepath{fill}%
\end{pgfscope}%
\begin{pgfscope}%
\pgfpathrectangle{\pgfqpoint{0.407592in}{0.222904in}}{\pgfqpoint{1.411163in}{1.411163in}}%
\pgfusepath{clip}%
\pgfsetbuttcap%
\pgfsetroundjoin%
\definecolor{currentfill}{rgb}{0.909804,0.698039,0.423529}%
\pgfsetfillcolor{currentfill}%
\pgfsetfillopacity{0.120000}%
\pgfsetlinewidth{0.000000pt}%
\definecolor{currentstroke}{rgb}{0.000000,0.000000,0.000000}%
\pgfsetstrokecolor{currentstroke}%
\pgfsetdash{}{0pt}%
\pgfpathmoveto{\pgfqpoint{0.926946in}{0.946149in}}%
\pgfpathlineto{\pgfqpoint{0.943067in}{0.934527in}}%
\pgfpathlineto{\pgfqpoint{0.959206in}{0.923650in}}%
\pgfpathlineto{\pgfqpoint{0.967489in}{0.927607in}}%
\pgfpathlineto{\pgfqpoint{0.975742in}{0.932092in}}%
\pgfpathlineto{\pgfqpoint{0.959640in}{0.942964in}}%
\pgfpathlineto{\pgfqpoint{0.943557in}{0.954581in}}%
\pgfpathlineto{\pgfqpoint{0.935267in}{0.950114in}}%
\pgfpathlineto{\pgfqpoint{0.926946in}{0.946149in}}%
\pgfpathclose%
\pgfusepath{fill}%
\end{pgfscope}%
\begin{pgfscope}%
\pgfpathrectangle{\pgfqpoint{0.407592in}{0.222904in}}{\pgfqpoint{1.411163in}{1.411163in}}%
\pgfusepath{clip}%
\pgfsetbuttcap%
\pgfsetroundjoin%
\definecolor{currentfill}{rgb}{0.909804,0.698039,0.423529}%
\pgfsetfillcolor{currentfill}%
\pgfsetfillopacity{0.120000}%
\pgfsetlinewidth{0.000000pt}%
\definecolor{currentstroke}{rgb}{0.000000,0.000000,0.000000}%
\pgfsetstrokecolor{currentstroke}%
\pgfsetdash{}{0pt}%
\pgfpathmoveto{\pgfqpoint{1.525056in}{1.148233in}}%
\pgfpathlineto{\pgfqpoint{1.542019in}{1.153900in}}%
\pgfpathlineto{\pgfqpoint{1.559050in}{1.159996in}}%
\pgfpathlineto{\pgfqpoint{1.566587in}{1.172539in}}%
\pgfpathlineto{\pgfqpoint{1.574112in}{1.185256in}}%
\pgfpathlineto{\pgfqpoint{1.557099in}{1.179432in}}%
\pgfpathlineto{\pgfqpoint{1.540154in}{1.174033in}}%
\pgfpathlineto{\pgfqpoint{1.532611in}{1.161047in}}%
\pgfpathlineto{\pgfqpoint{1.525056in}{1.148233in}}%
\pgfpathclose%
\pgfusepath{fill}%
\end{pgfscope}%
\begin{pgfscope}%
\pgfpathrectangle{\pgfqpoint{0.407592in}{0.222904in}}{\pgfqpoint{1.411163in}{1.411163in}}%
\pgfusepath{clip}%
\pgfsetbuttcap%
\pgfsetroundjoin%
\definecolor{currentfill}{rgb}{0.909804,0.698039,0.423529}%
\pgfsetfillcolor{currentfill}%
\pgfsetfillopacity{0.120000}%
\pgfsetlinewidth{0.000000pt}%
\definecolor{currentstroke}{rgb}{0.000000,0.000000,0.000000}%
\pgfsetstrokecolor{currentstroke}%
\pgfsetdash{}{0pt}%
\pgfpathmoveto{\pgfqpoint{0.845724in}{0.992408in}}%
\pgfpathlineto{\pgfqpoint{0.861823in}{0.978379in}}%
\pgfpathlineto{\pgfqpoint{0.877935in}{0.964887in}}%
\pgfpathlineto{\pgfqpoint{0.886362in}{0.967916in}}%
\pgfpathlineto{\pgfqpoint{0.894755in}{0.971414in}}%
\pgfpathlineto{\pgfqpoint{0.878683in}{0.984951in}}%
\pgfpathlineto{\pgfqpoint{0.862624in}{0.999025in}}%
\pgfpathlineto{\pgfqpoint{0.854191in}{0.995495in}}%
\pgfpathlineto{\pgfqpoint{0.845724in}{0.992408in}}%
\pgfpathclose%
\pgfusepath{fill}%
\end{pgfscope}%
\begin{pgfscope}%
\pgfpathrectangle{\pgfqpoint{0.407592in}{0.222904in}}{\pgfqpoint{1.411163in}{1.411163in}}%
\pgfusepath{clip}%
\pgfsetbuttcap%
\pgfsetroundjoin%
\definecolor{currentfill}{rgb}{0.909804,0.698039,0.423529}%
\pgfsetfillcolor{currentfill}%
\pgfsetfillopacity{0.120000}%
\pgfsetlinewidth{0.000000pt}%
\definecolor{currentstroke}{rgb}{0.000000,0.000000,0.000000}%
\pgfsetstrokecolor{currentstroke}%
\pgfsetdash{}{0pt}%
\pgfpathmoveto{\pgfqpoint{0.618185in}{1.179069in}}%
\pgfpathlineto{\pgfqpoint{0.634240in}{1.161406in}}%
\pgfpathlineto{\pgfqpoint{0.650304in}{1.143933in}}%
\pgfpathlineto{\pgfqpoint{0.659126in}{1.145913in}}%
\pgfpathlineto{\pgfqpoint{0.667910in}{1.148213in}}%
\pgfpathlineto{\pgfqpoint{0.651888in}{1.165765in}}%
\pgfpathlineto{\pgfqpoint{0.635876in}{1.183500in}}%
\pgfpathlineto{\pgfqpoint{0.627050in}{1.181132in}}%
\pgfpathlineto{\pgfqpoint{0.618185in}{1.179069in}}%
\pgfpathclose%
\pgfusepath{fill}%
\end{pgfscope}%
\begin{pgfscope}%
\pgfpathrectangle{\pgfqpoint{0.407592in}{0.222904in}}{\pgfqpoint{1.411163in}{1.411163in}}%
\pgfusepath{clip}%
\pgfsetbuttcap%
\pgfsetroundjoin%
\definecolor{currentfill}{rgb}{0.909804,0.698039,0.423529}%
\pgfsetfillcolor{currentfill}%
\pgfsetfillopacity{0.120000}%
\pgfsetlinewidth{0.000000pt}%
\definecolor{currentstroke}{rgb}{0.000000,0.000000,0.000000}%
\pgfsetstrokecolor{currentstroke}%
\pgfsetdash{}{0pt}%
\pgfpathmoveto{\pgfqpoint{1.299230in}{0.973254in}}%
\pgfpathlineto{\pgfqpoint{1.315780in}{0.975872in}}%
\pgfpathlineto{\pgfqpoint{1.332390in}{0.979202in}}%
\pgfpathlineto{\pgfqpoint{1.340157in}{0.989741in}}%
\pgfpathlineto{\pgfqpoint{1.347909in}{1.000622in}}%
\pgfpathlineto{\pgfqpoint{1.331316in}{0.997499in}}%
\pgfpathlineto{\pgfqpoint{1.314784in}{0.995073in}}%
\pgfpathlineto{\pgfqpoint{1.307015in}{0.983990in}}%
\pgfpathlineto{\pgfqpoint{1.299230in}{0.973254in}}%
\pgfpathclose%
\pgfusepath{fill}%
\end{pgfscope}%
\begin{pgfscope}%
\pgfpathrectangle{\pgfqpoint{0.407592in}{0.222904in}}{\pgfqpoint{1.411163in}{1.411163in}}%
\pgfusepath{clip}%
\pgfsetbuttcap%
\pgfsetroundjoin%
\definecolor{currentfill}{rgb}{0.909804,0.698039,0.423529}%
\pgfsetfillcolor{currentfill}%
\pgfsetfillopacity{0.120000}%
\pgfsetlinewidth{0.000000pt}%
\definecolor{currentstroke}{rgb}{0.000000,0.000000,0.000000}%
\pgfsetstrokecolor{currentstroke}%
\pgfsetdash{}{0pt}%
\pgfpathmoveto{\pgfqpoint{1.008010in}{0.912799in}}%
\pgfpathlineto{\pgfqpoint{1.024180in}{0.904479in}}%
\pgfpathlineto{\pgfqpoint{1.040377in}{0.897104in}}%
\pgfpathlineto{\pgfqpoint{1.048530in}{0.902192in}}%
\pgfpathlineto{\pgfqpoint{1.056656in}{0.907844in}}%
\pgfpathlineto{\pgfqpoint{1.040492in}{0.915178in}}%
\pgfpathlineto{\pgfqpoint{1.024355in}{0.923453in}}%
\pgfpathlineto{\pgfqpoint{1.016197in}{0.917852in}}%
\pgfpathlineto{\pgfqpoint{1.008010in}{0.912799in}}%
\pgfpathclose%
\pgfusepath{fill}%
\end{pgfscope}%
\begin{pgfscope}%
\pgfpathrectangle{\pgfqpoint{0.407592in}{0.222904in}}{\pgfqpoint{1.411163in}{1.411163in}}%
\pgfusepath{clip}%
\pgfsetbuttcap%
\pgfsetroundjoin%
\definecolor{currentfill}{rgb}{0.909804,0.698039,0.423529}%
\pgfsetfillcolor{currentfill}%
\pgfsetfillopacity{0.120000}%
\pgfsetlinewidth{0.000000pt}%
\definecolor{currentstroke}{rgb}{0.000000,0.000000,0.000000}%
\pgfsetstrokecolor{currentstroke}%
\pgfsetdash{}{0pt}%
\pgfpathmoveto{\pgfqpoint{1.460920in}{1.087866in}}%
\pgfpathlineto{\pgfqpoint{1.477783in}{1.093145in}}%
\pgfpathlineto{\pgfqpoint{1.494713in}{1.098906in}}%
\pgfpathlineto{\pgfqpoint{1.502317in}{1.110926in}}%
\pgfpathlineto{\pgfqpoint{1.509909in}{1.123162in}}%
\pgfpathlineto{\pgfqpoint{1.492997in}{1.117674in}}%
\pgfpathlineto{\pgfqpoint{1.476152in}{1.112664in}}%
\pgfpathlineto{\pgfqpoint{1.468542in}{1.100158in}}%
\pgfpathlineto{\pgfqpoint{1.460920in}{1.087866in}}%
\pgfpathclose%
\pgfusepath{fill}%
\end{pgfscope}%
\begin{pgfscope}%
\pgfpathrectangle{\pgfqpoint{0.407592in}{0.222904in}}{\pgfqpoint{1.411163in}{1.411163in}}%
\pgfusepath{clip}%
\pgfsetbuttcap%
\pgfsetroundjoin%
\definecolor{currentfill}{rgb}{0.909804,0.698039,0.423529}%
\pgfsetfillcolor{currentfill}%
\pgfsetfillopacity{0.120000}%
\pgfsetlinewidth{0.000000pt}%
\definecolor{currentstroke}{rgb}{0.000000,0.000000,0.000000}%
\pgfsetstrokecolor{currentstroke}%
\pgfsetdash{}{0pt}%
\pgfpathmoveto{\pgfqpoint{0.764246in}{1.047843in}}%
\pgfpathlineto{\pgfqpoint{0.780338in}{1.032164in}}%
\pgfpathlineto{\pgfqpoint{0.796441in}{1.016864in}}%
\pgfpathlineto{\pgfqpoint{0.805021in}{1.019179in}}%
\pgfpathlineto{\pgfqpoint{0.813564in}{1.021904in}}%
\pgfpathlineto{\pgfqpoint{0.797502in}{1.037288in}}%
\pgfpathlineto{\pgfqpoint{0.781452in}{1.053047in}}%
\pgfpathlineto{\pgfqpoint{0.772867in}{1.050252in}}%
\pgfpathlineto{\pgfqpoint{0.764246in}{1.047843in}}%
\pgfpathclose%
\pgfusepath{fill}%
\end{pgfscope}%
\begin{pgfscope}%
\pgfpathrectangle{\pgfqpoint{0.407592in}{0.222904in}}{\pgfqpoint{1.411163in}{1.411163in}}%
\pgfusepath{clip}%
\pgfsetbuttcap%
\pgfsetroundjoin%
\definecolor{currentfill}{rgb}{0.909804,0.698039,0.423529}%
\pgfsetfillcolor{currentfill}%
\pgfsetfillopacity{0.120000}%
\pgfsetlinewidth{0.000000pt}%
\definecolor{currentstroke}{rgb}{0.000000,0.000000,0.000000}%
\pgfsetstrokecolor{currentstroke}%
\pgfsetdash{}{0pt}%
\pgfpathmoveto{\pgfqpoint{1.347909in}{1.000622in}}%
\pgfpathlineto{\pgfqpoint{1.364563in}{1.004406in}}%
\pgfpathlineto{\pgfqpoint{1.381282in}{1.008814in}}%
\pgfpathlineto{\pgfqpoint{1.389001in}{1.019785in}}%
\pgfpathlineto{\pgfqpoint{1.396706in}{1.031059in}}%
\pgfpathlineto{\pgfqpoint{1.380006in}{1.026891in}}%
\pgfpathlineto{\pgfqpoint{1.363369in}{1.023338in}}%
\pgfpathlineto{\pgfqpoint{1.355646in}{1.011827in}}%
\pgfpathlineto{\pgfqpoint{1.347909in}{1.000622in}}%
\pgfpathclose%
\pgfusepath{fill}%
\end{pgfscope}%
\begin{pgfscope}%
\pgfpathrectangle{\pgfqpoint{0.407592in}{0.222904in}}{\pgfqpoint{1.411163in}{1.411163in}}%
\pgfusepath{clip}%
\pgfsetbuttcap%
\pgfsetroundjoin%
\definecolor{currentfill}{rgb}{0.909804,0.698039,0.423529}%
\pgfsetfillcolor{currentfill}%
\pgfsetfillopacity{0.120000}%
\pgfsetlinewidth{0.000000pt}%
\definecolor{currentstroke}{rgb}{0.000000,0.000000,0.000000}%
\pgfsetstrokecolor{currentstroke}%
\pgfsetdash{}{0pt}%
\pgfpathmoveto{\pgfqpoint{0.650304in}{1.143933in}}%
\pgfpathlineto{\pgfqpoint{0.666377in}{1.126667in}}%
\pgfpathlineto{\pgfqpoint{0.682460in}{1.109626in}}%
\pgfpathlineto{\pgfqpoint{0.691240in}{1.111509in}}%
\pgfpathlineto{\pgfqpoint{0.699982in}{1.113730in}}%
\pgfpathlineto{\pgfqpoint{0.683941in}{1.130862in}}%
\pgfpathlineto{\pgfqpoint{0.667910in}{1.148213in}}%
\pgfpathlineto{\pgfqpoint{0.659126in}{1.145913in}}%
\pgfpathlineto{\pgfqpoint{0.650304in}{1.143933in}}%
\pgfpathclose%
\pgfusepath{fill}%
\end{pgfscope}%
\begin{pgfscope}%
\pgfpathrectangle{\pgfqpoint{0.407592in}{0.222904in}}{\pgfqpoint{1.411163in}{1.411163in}}%
\pgfusepath{clip}%
\pgfsetbuttcap%
\pgfsetroundjoin%
\definecolor{currentfill}{rgb}{0.909804,0.698039,0.423529}%
\pgfsetfillcolor{currentfill}%
\pgfsetfillopacity{0.120000}%
\pgfsetlinewidth{0.000000pt}%
\definecolor{currentstroke}{rgb}{0.000000,0.000000,0.000000}%
\pgfsetstrokecolor{currentstroke}%
\pgfsetdash{}{0pt}%
\pgfpathmoveto{\pgfqpoint{0.877935in}{0.964887in}}%
\pgfpathlineto{\pgfqpoint{0.894062in}{0.951983in}}%
\pgfpathlineto{\pgfqpoint{0.910205in}{0.939716in}}%
\pgfpathlineto{\pgfqpoint{0.918592in}{0.942685in}}%
\pgfpathlineto{\pgfqpoint{0.926946in}{0.946149in}}%
\pgfpathlineto{\pgfqpoint{0.910843in}{0.958463in}}%
\pgfpathlineto{\pgfqpoint{0.894755in}{0.971414in}}%
\pgfpathlineto{\pgfqpoint{0.886362in}{0.967916in}}%
\pgfpathlineto{\pgfqpoint{0.877935in}{0.964887in}}%
\pgfpathclose%
\pgfusepath{fill}%
\end{pgfscope}%
\begin{pgfscope}%
\pgfpathrectangle{\pgfqpoint{0.407592in}{0.222904in}}{\pgfqpoint{1.411163in}{1.411163in}}%
\pgfusepath{clip}%
\pgfsetbuttcap%
\pgfsetroundjoin%
\definecolor{currentfill}{rgb}{0.909804,0.698039,0.423529}%
\pgfsetfillcolor{currentfill}%
\pgfsetfillopacity{0.120000}%
\pgfsetlinewidth{0.000000pt}%
\definecolor{currentstroke}{rgb}{0.000000,0.000000,0.000000}%
\pgfsetstrokecolor{currentstroke}%
\pgfsetdash{}{0pt}%
\pgfpathmoveto{\pgfqpoint{0.959206in}{0.923650in}}%
\pgfpathlineto{\pgfqpoint{0.975365in}{0.913571in}}%
\pgfpathlineto{\pgfqpoint{0.991546in}{0.904343in}}%
\pgfpathlineto{\pgfqpoint{0.999793in}{0.908296in}}%
\pgfpathlineto{\pgfqpoint{1.008010in}{0.912799in}}%
\pgfpathlineto{\pgfqpoint{0.991864in}{0.922019in}}%
\pgfpathlineto{\pgfqpoint{0.975742in}{0.932092in}}%
\pgfpathlineto{\pgfqpoint{0.967489in}{0.927607in}}%
\pgfpathlineto{\pgfqpoint{0.959206in}{0.923650in}}%
\pgfpathclose%
\pgfusepath{fill}%
\end{pgfscope}%
\begin{pgfscope}%
\pgfpathrectangle{\pgfqpoint{0.407592in}{0.222904in}}{\pgfqpoint{1.411163in}{1.411163in}}%
\pgfusepath{clip}%
\pgfsetbuttcap%
\pgfsetroundjoin%
\definecolor{currentfill}{rgb}{0.909804,0.698039,0.423529}%
\pgfsetfillcolor{currentfill}%
\pgfsetfillopacity{0.120000}%
\pgfsetlinewidth{0.000000pt}%
\definecolor{currentstroke}{rgb}{0.000000,0.000000,0.000000}%
\pgfsetstrokecolor{currentstroke}%
\pgfsetdash{}{0pt}%
\pgfpathmoveto{\pgfqpoint{1.137705in}{0.901610in}}%
\pgfpathlineto{\pgfqpoint{1.154023in}{0.899330in}}%
\pgfpathlineto{\pgfqpoint{1.170387in}{0.898045in}}%
\pgfpathlineto{\pgfqpoint{1.178357in}{0.905419in}}%
\pgfpathlineto{\pgfqpoint{1.186306in}{0.913316in}}%
\pgfpathlineto{\pgfqpoint{1.169966in}{0.914626in}}%
\pgfpathlineto{\pgfqpoint{1.153672in}{0.916911in}}%
\pgfpathlineto{\pgfqpoint{1.145700in}{0.908995in}}%
\pgfpathlineto{\pgfqpoint{1.137705in}{0.901610in}}%
\pgfpathclose%
\pgfusepath{fill}%
\end{pgfscope}%
\begin{pgfscope}%
\pgfpathrectangle{\pgfqpoint{0.407592in}{0.222904in}}{\pgfqpoint{1.411163in}{1.411163in}}%
\pgfusepath{clip}%
\pgfsetbuttcap%
\pgfsetroundjoin%
\definecolor{currentfill}{rgb}{0.909804,0.698039,0.423529}%
\pgfsetfillcolor{currentfill}%
\pgfsetfillopacity{0.120000}%
\pgfsetlinewidth{0.000000pt}%
\definecolor{currentstroke}{rgb}{0.000000,0.000000,0.000000}%
\pgfsetstrokecolor{currentstroke}%
\pgfsetdash{}{0pt}%
\pgfpathmoveto{\pgfqpoint{1.186306in}{0.913316in}}%
\pgfpathlineto{\pgfqpoint{1.202694in}{0.912958in}}%
\pgfpathlineto{\pgfqpoint{1.219135in}{0.913520in}}%
\pgfpathlineto{\pgfqpoint{1.227041in}{0.921858in}}%
\pgfpathlineto{\pgfqpoint{1.234929in}{0.930672in}}%
\pgfpathlineto{\pgfqpoint{1.218510in}{0.930192in}}%
\pgfpathlineto{\pgfqpoint{1.202143in}{0.930612in}}%
\pgfpathlineto{\pgfqpoint{1.194234in}{0.921720in}}%
\pgfpathlineto{\pgfqpoint{1.186306in}{0.913316in}}%
\pgfpathclose%
\pgfusepath{fill}%
\end{pgfscope}%
\begin{pgfscope}%
\pgfpathrectangle{\pgfqpoint{0.407592in}{0.222904in}}{\pgfqpoint{1.411163in}{1.411163in}}%
\pgfusepath{clip}%
\pgfsetbuttcap%
\pgfsetroundjoin%
\definecolor{currentfill}{rgb}{0.909804,0.698039,0.423529}%
\pgfsetfillcolor{currentfill}%
\pgfsetfillopacity{0.120000}%
\pgfsetlinewidth{0.000000pt}%
\definecolor{currentstroke}{rgb}{0.000000,0.000000,0.000000}%
\pgfsetstrokecolor{currentstroke}%
\pgfsetdash{}{0pt}%
\pgfpathmoveto{\pgfqpoint{1.574112in}{1.185256in}}%
\pgfpathlineto{\pgfqpoint{1.591195in}{1.191486in}}%
\pgfpathlineto{\pgfqpoint{1.608348in}{1.198103in}}%
\pgfpathlineto{\pgfqpoint{1.615843in}{1.210716in}}%
\pgfpathlineto{\pgfqpoint{1.598699in}{1.204234in}}%
\pgfpathlineto{\pgfqpoint{1.581626in}{1.198139in}}%
\pgfusepath{fill}%
\end{pgfscope}%
\begin{pgfscope}%
\pgfpathrectangle{\pgfqpoint{0.407592in}{0.222904in}}{\pgfqpoint{1.411163in}{1.411163in}}%
\pgfusepath{clip}%
\pgfsetbuttcap%
\pgfsetroundjoin%
\definecolor{currentfill}{rgb}{0.909804,0.698039,0.423529}%
\pgfsetfillcolor{currentfill}%
\pgfsetfillopacity{0.120000}%
\pgfsetlinewidth{0.000000pt}%
\definecolor{currentstroke}{rgb}{0.000000,0.000000,0.000000}%
\pgfsetstrokecolor{currentstroke}%
\pgfsetdash{}{0pt}%
\pgfpathmoveto{\pgfqpoint{1.089079in}{0.896130in}}%
\pgfpathlineto{\pgfqpoint{1.105343in}{0.891796in}}%
\pgfpathlineto{\pgfqpoint{1.121647in}{0.888489in}}%
\pgfpathlineto{\pgfqpoint{1.129688in}{0.894771in}}%
\pgfpathlineto{\pgfqpoint{1.137705in}{0.901610in}}%
\pgfpathlineto{\pgfqpoint{1.121429in}{0.904899in}}%
\pgfpathlineto{\pgfqpoint{1.105193in}{0.909202in}}%
\pgfpathlineto{\pgfqpoint{1.097149in}{0.902387in}}%
\pgfpathlineto{\pgfqpoint{1.089079in}{0.896130in}}%
\pgfpathclose%
\pgfusepath{fill}%
\end{pgfscope}%
\begin{pgfscope}%
\pgfpathrectangle{\pgfqpoint{0.407592in}{0.222904in}}{\pgfqpoint{1.411163in}{1.411163in}}%
\pgfusepath{clip}%
\pgfsetbuttcap%
\pgfsetroundjoin%
\definecolor{currentfill}{rgb}{0.909804,0.698039,0.423529}%
\pgfsetfillcolor{currentfill}%
\pgfsetfillopacity{0.120000}%
\pgfsetlinewidth{0.000000pt}%
\definecolor{currentstroke}{rgb}{0.000000,0.000000,0.000000}%
\pgfsetstrokecolor{currentstroke}%
\pgfsetdash{}{0pt}%
\pgfpathmoveto{\pgfqpoint{1.396706in}{1.031059in}}%
\pgfpathlineto{\pgfqpoint{1.413472in}{1.035808in}}%
\pgfpathlineto{\pgfqpoint{1.430303in}{1.041105in}}%
\pgfpathlineto{\pgfqpoint{1.437977in}{1.052407in}}%
\pgfpathlineto{\pgfqpoint{1.445638in}{1.063977in}}%
\pgfpathlineto{\pgfqpoint{1.428824in}{1.058942in}}%
\pgfpathlineto{\pgfqpoint{1.412077in}{1.054448in}}%
\pgfpathlineto{\pgfqpoint{1.404398in}{1.042618in}}%
\pgfpathlineto{\pgfqpoint{1.396706in}{1.031059in}}%
\pgfpathclose%
\pgfusepath{fill}%
\end{pgfscope}%
\begin{pgfscope}%
\pgfpathrectangle{\pgfqpoint{0.407592in}{0.222904in}}{\pgfqpoint{1.411163in}{1.411163in}}%
\pgfusepath{clip}%
\pgfsetbuttcap%
\pgfsetroundjoin%
\definecolor{currentfill}{rgb}{0.909804,0.698039,0.423529}%
\pgfsetfillcolor{currentfill}%
\pgfsetfillopacity{0.120000}%
\pgfsetlinewidth{0.000000pt}%
\definecolor{currentstroke}{rgb}{0.000000,0.000000,0.000000}%
\pgfsetstrokecolor{currentstroke}%
\pgfsetdash{}{0pt}%
\pgfpathmoveto{\pgfqpoint{1.234929in}{0.930672in}}%
\pgfpathlineto{\pgfqpoint{1.251402in}{0.932017in}}%
\pgfpathlineto{\pgfqpoint{1.267932in}{0.934188in}}%
\pgfpathlineto{\pgfqpoint{1.275781in}{0.943333in}}%
\pgfpathlineto{\pgfqpoint{1.283614in}{0.952905in}}%
\pgfpathlineto{\pgfqpoint{1.267104in}{0.950877in}}%
\pgfpathlineto{\pgfqpoint{1.250651in}{0.949655in}}%
\pgfpathlineto{\pgfqpoint{1.242799in}{0.939945in}}%
\pgfpathlineto{\pgfqpoint{1.234929in}{0.930672in}}%
\pgfpathclose%
\pgfusepath{fill}%
\end{pgfscope}%
\begin{pgfscope}%
\pgfpathrectangle{\pgfqpoint{0.407592in}{0.222904in}}{\pgfqpoint{1.411163in}{1.411163in}}%
\pgfusepath{clip}%
\pgfsetbuttcap%
\pgfsetroundjoin%
\definecolor{currentfill}{rgb}{0.909804,0.698039,0.423529}%
\pgfsetfillcolor{currentfill}%
\pgfsetfillopacity{0.120000}%
\pgfsetlinewidth{0.000000pt}%
\definecolor{currentstroke}{rgb}{0.000000,0.000000,0.000000}%
\pgfsetstrokecolor{currentstroke}%
\pgfsetdash{}{0pt}%
\pgfpathmoveto{\pgfqpoint{1.509909in}{1.123162in}}%
\pgfpathlineto{\pgfqpoint{1.526890in}{1.129102in}}%
\pgfpathlineto{\pgfqpoint{1.543940in}{1.135473in}}%
\pgfpathlineto{\pgfqpoint{1.551501in}{1.147637in}}%
\pgfpathlineto{\pgfqpoint{1.559050in}{1.159996in}}%
\pgfpathlineto{\pgfqpoint{1.542019in}{1.153900in}}%
\pgfpathlineto{\pgfqpoint{1.525056in}{1.148233in}}%
\pgfpathlineto{\pgfqpoint{1.517489in}{1.135601in}}%
\pgfpathlineto{\pgfqpoint{1.509909in}{1.123162in}}%
\pgfpathclose%
\pgfusepath{fill}%
\end{pgfscope}%
\begin{pgfscope}%
\pgfpathrectangle{\pgfqpoint{0.407592in}{0.222904in}}{\pgfqpoint{1.411163in}{1.411163in}}%
\pgfusepath{clip}%
\pgfsetbuttcap%
\pgfsetroundjoin%
\definecolor{currentfill}{rgb}{0.909804,0.698039,0.423529}%
\pgfsetfillcolor{currentfill}%
\pgfsetfillopacity{0.120000}%
\pgfsetlinewidth{0.000000pt}%
\definecolor{currentstroke}{rgb}{0.000000,0.000000,0.000000}%
\pgfsetstrokecolor{currentstroke}%
\pgfsetdash{}{0pt}%
\pgfpathmoveto{\pgfqpoint{0.682460in}{1.109626in}}%
\pgfpathlineto{\pgfqpoint{0.698553in}{1.092831in}}%
\pgfpathlineto{\pgfqpoint{0.714655in}{1.076306in}}%
\pgfpathlineto{\pgfqpoint{0.723393in}{1.078079in}}%
\pgfpathlineto{\pgfqpoint{0.732094in}{1.080209in}}%
\pgfpathlineto{\pgfqpoint{0.716033in}{1.096838in}}%
\pgfpathlineto{\pgfqpoint{0.699982in}{1.113730in}}%
\pgfpathlineto{\pgfqpoint{0.691240in}{1.111509in}}%
\pgfpathlineto{\pgfqpoint{0.682460in}{1.109626in}}%
\pgfpathclose%
\pgfusepath{fill}%
\end{pgfscope}%
\begin{pgfscope}%
\pgfpathrectangle{\pgfqpoint{0.407592in}{0.222904in}}{\pgfqpoint{1.411163in}{1.411163in}}%
\pgfusepath{clip}%
\pgfsetbuttcap%
\pgfsetroundjoin%
\definecolor{currentfill}{rgb}{0.909804,0.698039,0.423529}%
\pgfsetfillcolor{currentfill}%
\pgfsetfillopacity{0.120000}%
\pgfsetlinewidth{0.000000pt}%
\definecolor{currentstroke}{rgb}{0.000000,0.000000,0.000000}%
\pgfsetstrokecolor{currentstroke}%
\pgfsetdash{}{0pt}%
\pgfpathmoveto{\pgfqpoint{0.796441in}{1.016864in}}%
\pgfpathlineto{\pgfqpoint{0.812556in}{1.001978in}}%
\pgfpathlineto{\pgfqpoint{0.828683in}{0.987546in}}%
\pgfpathlineto{\pgfqpoint{0.837221in}{0.989761in}}%
\pgfpathlineto{\pgfqpoint{0.845724in}{0.992408in}}%
\pgfpathlineto{\pgfqpoint{0.829638in}{1.006931in}}%
\pgfpathlineto{\pgfqpoint{0.813564in}{1.021904in}}%
\pgfpathlineto{\pgfqpoint{0.805021in}{1.019179in}}%
\pgfpathlineto{\pgfqpoint{0.796441in}{1.016864in}}%
\pgfpathclose%
\pgfusepath{fill}%
\end{pgfscope}%
\begin{pgfscope}%
\pgfpathrectangle{\pgfqpoint{0.407592in}{0.222904in}}{\pgfqpoint{1.411163in}{1.411163in}}%
\pgfusepath{clip}%
\pgfsetbuttcap%
\pgfsetroundjoin%
\definecolor{currentfill}{rgb}{0.909804,0.698039,0.423529}%
\pgfsetfillcolor{currentfill}%
\pgfsetfillopacity{0.120000}%
\pgfsetlinewidth{0.000000pt}%
\definecolor{currentstroke}{rgb}{0.000000,0.000000,0.000000}%
\pgfsetstrokecolor{currentstroke}%
\pgfsetdash{}{0pt}%
\pgfpathmoveto{\pgfqpoint{1.040377in}{0.897104in}}%
\pgfpathlineto{\pgfqpoint{1.056604in}{0.890710in}}%
\pgfpathlineto{\pgfqpoint{1.072863in}{0.885323in}}%
\pgfpathlineto{\pgfqpoint{1.080984in}{0.890440in}}%
\pgfpathlineto{\pgfqpoint{1.089079in}{0.896130in}}%
\pgfpathlineto{\pgfqpoint{1.072851in}{0.901485in}}%
\pgfpathlineto{\pgfqpoint{1.056656in}{0.907844in}}%
\pgfpathlineto{\pgfqpoint{1.048530in}{0.902192in}}%
\pgfpathlineto{\pgfqpoint{1.040377in}{0.897104in}}%
\pgfpathclose%
\pgfusepath{fill}%
\end{pgfscope}%
\begin{pgfscope}%
\pgfpathrectangle{\pgfqpoint{0.407592in}{0.222904in}}{\pgfqpoint{1.411163in}{1.411163in}}%
\pgfusepath{clip}%
\pgfsetbuttcap%
\pgfsetroundjoin%
\definecolor{currentfill}{rgb}{0.909804,0.698039,0.423529}%
\pgfsetfillcolor{currentfill}%
\pgfsetfillopacity{0.120000}%
\pgfsetlinewidth{0.000000pt}%
\definecolor{currentstroke}{rgb}{0.000000,0.000000,0.000000}%
\pgfsetstrokecolor{currentstroke}%
\pgfsetdash{}{0pt}%
\pgfpathmoveto{\pgfqpoint{0.568173in}{1.211487in}}%
\pgfpathlineto{\pgfqpoint{0.584252in}{1.193582in}}%
\pgfpathlineto{\pgfqpoint{0.600339in}{1.175842in}}%
\pgfpathlineto{\pgfqpoint{0.609282in}{1.177306in}}%
\pgfpathlineto{\pgfqpoint{0.618185in}{1.179069in}}%
\pgfpathlineto{\pgfqpoint{0.602140in}{1.196904in}}%
\pgfpathlineto{\pgfqpoint{0.586104in}{1.214899in}}%
\pgfpathlineto{\pgfqpoint{0.577158in}{1.213051in}}%
\pgfpathlineto{\pgfqpoint{0.568173in}{1.211487in}}%
\pgfpathclose%
\pgfusepath{fill}%
\end{pgfscope}%
\begin{pgfscope}%
\pgfpathrectangle{\pgfqpoint{0.407592in}{0.222904in}}{\pgfqpoint{1.411163in}{1.411163in}}%
\pgfusepath{clip}%
\pgfsetbuttcap%
\pgfsetroundjoin%
\definecolor{currentfill}{rgb}{0.909804,0.698039,0.423529}%
\pgfsetfillcolor{currentfill}%
\pgfsetfillopacity{0.120000}%
\pgfsetlinewidth{0.000000pt}%
\definecolor{currentstroke}{rgb}{0.000000,0.000000,0.000000}%
\pgfsetstrokecolor{currentstroke}%
\pgfsetdash{}{0pt}%
\pgfpathmoveto{\pgfqpoint{1.283614in}{0.952905in}}%
\pgfpathlineto{\pgfqpoint{1.300182in}{0.955702in}}%
\pgfpathlineto{\pgfqpoint{1.316812in}{0.959225in}}%
\pgfpathlineto{\pgfqpoint{1.324609in}{0.969024in}}%
\pgfpathlineto{\pgfqpoint{1.332390in}{0.979202in}}%
\pgfpathlineto{\pgfqpoint{1.315780in}{0.975872in}}%
\pgfpathlineto{\pgfqpoint{1.299230in}{0.973254in}}%
\pgfpathlineto{\pgfqpoint{1.291430in}{0.962886in}}%
\pgfpathlineto{\pgfqpoint{1.283614in}{0.952905in}}%
\pgfpathclose%
\pgfusepath{fill}%
\end{pgfscope}%
\begin{pgfscope}%
\pgfpathrectangle{\pgfqpoint{0.407592in}{0.222904in}}{\pgfqpoint{1.411163in}{1.411163in}}%
\pgfusepath{clip}%
\pgfsetbuttcap%
\pgfsetroundjoin%
\definecolor{currentfill}{rgb}{0.909804,0.698039,0.423529}%
\pgfsetfillcolor{currentfill}%
\pgfsetfillopacity{0.120000}%
\pgfsetlinewidth{0.000000pt}%
\definecolor{currentstroke}{rgb}{0.000000,0.000000,0.000000}%
\pgfsetstrokecolor{currentstroke}%
\pgfsetdash{}{0pt}%
\pgfpathmoveto{\pgfqpoint{0.910205in}{0.939716in}}%
\pgfpathlineto{\pgfqpoint{0.926364in}{0.928138in}}%
\pgfpathlineto{\pgfqpoint{0.942541in}{0.917303in}}%
\pgfpathlineto{\pgfqpoint{0.950890in}{0.920217in}}%
\pgfpathlineto{\pgfqpoint{0.959206in}{0.923650in}}%
\pgfpathlineto{\pgfqpoint{0.943067in}{0.934527in}}%
\pgfpathlineto{\pgfqpoint{0.926946in}{0.946149in}}%
\pgfpathlineto{\pgfqpoint{0.918592in}{0.942685in}}%
\pgfpathlineto{\pgfqpoint{0.910205in}{0.939716in}}%
\pgfpathclose%
\pgfusepath{fill}%
\end{pgfscope}%
\begin{pgfscope}%
\pgfpathrectangle{\pgfqpoint{0.407592in}{0.222904in}}{\pgfqpoint{1.411163in}{1.411163in}}%
\pgfusepath{clip}%
\pgfsetbuttcap%
\pgfsetroundjoin%
\definecolor{currentfill}{rgb}{0.909804,0.698039,0.423529}%
\pgfsetfillcolor{currentfill}%
\pgfsetfillopacity{0.120000}%
\pgfsetlinewidth{0.000000pt}%
\definecolor{currentstroke}{rgb}{0.000000,0.000000,0.000000}%
\pgfsetstrokecolor{currentstroke}%
\pgfsetdash{}{0pt}%
\pgfpathmoveto{\pgfqpoint{1.445638in}{1.063977in}}%
\pgfpathlineto{\pgfqpoint{1.462519in}{1.069526in}}%
\pgfpathlineto{\pgfqpoint{1.479467in}{1.075559in}}%
\pgfpathlineto{\pgfqpoint{1.487097in}{1.087113in}}%
\pgfpathlineto{\pgfqpoint{1.494713in}{1.098906in}}%
\pgfpathlineto{\pgfqpoint{1.477783in}{1.093145in}}%
\pgfpathlineto{\pgfqpoint{1.460920in}{1.087866in}}%
\pgfpathlineto{\pgfqpoint{1.453286in}{1.075801in}}%
\pgfpathlineto{\pgfqpoint{1.445638in}{1.063977in}}%
\pgfpathclose%
\pgfusepath{fill}%
\end{pgfscope}%
\begin{pgfscope}%
\pgfpathrectangle{\pgfqpoint{0.407592in}{0.222904in}}{\pgfqpoint{1.411163in}{1.411163in}}%
\pgfusepath{clip}%
\pgfsetbuttcap%
\pgfsetroundjoin%
\definecolor{currentfill}{rgb}{0.909804,0.698039,0.423529}%
\pgfsetfillcolor{currentfill}%
\pgfsetfillopacity{0.120000}%
\pgfsetlinewidth{0.000000pt}%
\definecolor{currentstroke}{rgb}{0.000000,0.000000,0.000000}%
\pgfsetstrokecolor{currentstroke}%
\pgfsetdash{}{0pt}%
\pgfpathmoveto{\pgfqpoint{0.714655in}{1.076306in}}%
\pgfpathlineto{\pgfqpoint{0.730768in}{1.060074in}}%
\pgfpathlineto{\pgfqpoint{0.746891in}{1.044165in}}%
\pgfpathlineto{\pgfqpoint{0.755587in}{1.045816in}}%
\pgfpathlineto{\pgfqpoint{0.764246in}{1.047843in}}%
\pgfpathlineto{\pgfqpoint{0.748165in}{1.063868in}}%
\pgfpathlineto{\pgfqpoint{0.732094in}{1.080209in}}%
\pgfpathlineto{\pgfqpoint{0.723393in}{1.078079in}}%
\pgfpathlineto{\pgfqpoint{0.714655in}{1.076306in}}%
\pgfpathclose%
\pgfusepath{fill}%
\end{pgfscope}%
\begin{pgfscope}%
\pgfpathrectangle{\pgfqpoint{0.407592in}{0.222904in}}{\pgfqpoint{1.411163in}{1.411163in}}%
\pgfusepath{clip}%
\pgfsetbuttcap%
\pgfsetroundjoin%
\definecolor{currentfill}{rgb}{0.909804,0.698039,0.423529}%
\pgfsetfillcolor{currentfill}%
\pgfsetfillopacity{0.120000}%
\pgfsetlinewidth{0.000000pt}%
\definecolor{currentstroke}{rgb}{0.000000,0.000000,0.000000}%
\pgfsetstrokecolor{currentstroke}%
\pgfsetdash{}{0pt}%
\pgfpathmoveto{\pgfqpoint{1.559050in}{1.159996in}}%
\pgfpathlineto{\pgfqpoint{1.576151in}{1.166498in}}%
\pgfpathlineto{\pgfqpoint{1.593323in}{1.173389in}}%
\pgfpathlineto{\pgfqpoint{1.600842in}{1.185658in}}%
\pgfpathlineto{\pgfqpoint{1.608348in}{1.198103in}}%
\pgfpathlineto{\pgfqpoint{1.591195in}{1.191486in}}%
\pgfpathlineto{\pgfqpoint{1.574112in}{1.185256in}}%
\pgfpathlineto{\pgfqpoint{1.566587in}{1.172539in}}%
\pgfpathlineto{\pgfqpoint{1.559050in}{1.159996in}}%
\pgfpathclose%
\pgfusepath{fill}%
\end{pgfscope}%
\begin{pgfscope}%
\pgfpathrectangle{\pgfqpoint{0.407592in}{0.222904in}}{\pgfqpoint{1.411163in}{1.411163in}}%
\pgfusepath{clip}%
\pgfsetbuttcap%
\pgfsetroundjoin%
\definecolor{currentfill}{rgb}{0.909804,0.698039,0.423529}%
\pgfsetfillcolor{currentfill}%
\pgfsetfillopacity{0.120000}%
\pgfsetlinewidth{0.000000pt}%
\definecolor{currentstroke}{rgb}{0.000000,0.000000,0.000000}%
\pgfsetstrokecolor{currentstroke}%
\pgfsetdash{}{0pt}%
\pgfpathmoveto{\pgfqpoint{0.828683in}{0.987546in}}%
\pgfpathlineto{\pgfqpoint{0.844822in}{0.973610in}}%
\pgfpathlineto{\pgfqpoint{0.860975in}{0.960214in}}%
\pgfpathlineto{\pgfqpoint{0.869473in}{0.962323in}}%
\pgfpathlineto{\pgfqpoint{0.877935in}{0.964887in}}%
\pgfpathlineto{\pgfqpoint{0.861823in}{0.978379in}}%
\pgfpathlineto{\pgfqpoint{0.845724in}{0.992408in}}%
\pgfpathlineto{\pgfqpoint{0.837221in}{0.989761in}}%
\pgfpathlineto{\pgfqpoint{0.828683in}{0.987546in}}%
\pgfpathclose%
\pgfusepath{fill}%
\end{pgfscope}%
\begin{pgfscope}%
\pgfpathrectangle{\pgfqpoint{0.407592in}{0.222904in}}{\pgfqpoint{1.411163in}{1.411163in}}%
\pgfusepath{clip}%
\pgfsetbuttcap%
\pgfsetroundjoin%
\definecolor{currentfill}{rgb}{0.909804,0.698039,0.423529}%
\pgfsetfillcolor{currentfill}%
\pgfsetfillopacity{0.120000}%
\pgfsetlinewidth{0.000000pt}%
\definecolor{currentstroke}{rgb}{0.000000,0.000000,0.000000}%
\pgfsetstrokecolor{currentstroke}%
\pgfsetdash{}{0pt}%
\pgfpathmoveto{\pgfqpoint{1.332390in}{0.979202in}}%
\pgfpathlineto{\pgfqpoint{1.349063in}{0.983206in}}%
\pgfpathlineto{\pgfqpoint{1.365800in}{0.987845in}}%
\pgfpathlineto{\pgfqpoint{1.373548in}{0.998162in}}%
\pgfpathlineto{\pgfqpoint{1.381282in}{1.008814in}}%
\pgfpathlineto{\pgfqpoint{1.364563in}{1.004406in}}%
\pgfpathlineto{\pgfqpoint{1.347909in}{1.000622in}}%
\pgfpathlineto{\pgfqpoint{1.340157in}{0.989741in}}%
\pgfpathlineto{\pgfqpoint{1.332390in}{0.979202in}}%
\pgfpathclose%
\pgfusepath{fill}%
\end{pgfscope}%
\begin{pgfscope}%
\pgfpathrectangle{\pgfqpoint{0.407592in}{0.222904in}}{\pgfqpoint{1.411163in}{1.411163in}}%
\pgfusepath{clip}%
\pgfsetbuttcap%
\pgfsetroundjoin%
\definecolor{currentfill}{rgb}{0.909804,0.698039,0.423529}%
\pgfsetfillcolor{currentfill}%
\pgfsetfillopacity{0.120000}%
\pgfsetlinewidth{0.000000pt}%
\definecolor{currentstroke}{rgb}{0.000000,0.000000,0.000000}%
\pgfsetstrokecolor{currentstroke}%
\pgfsetdash{}{0pt}%
\pgfpathmoveto{\pgfqpoint{0.991546in}{0.904343in}}%
\pgfpathlineto{\pgfqpoint{1.007751in}{0.896014in}}%
\pgfpathlineto{\pgfqpoint{1.023983in}{0.888627in}}%
\pgfpathlineto{\pgfqpoint{1.032195in}{0.892583in}}%
\pgfpathlineto{\pgfqpoint{1.040377in}{0.897104in}}%
\pgfpathlineto{\pgfqpoint{1.024180in}{0.904479in}}%
\pgfpathlineto{\pgfqpoint{1.008010in}{0.912799in}}%
\pgfpathlineto{\pgfqpoint{0.999793in}{0.908296in}}%
\pgfpathlineto{\pgfqpoint{0.991546in}{0.904343in}}%
\pgfpathclose%
\pgfusepath{fill}%
\end{pgfscope}%
\begin{pgfscope}%
\pgfpathrectangle{\pgfqpoint{0.407592in}{0.222904in}}{\pgfqpoint{1.411163in}{1.411163in}}%
\pgfusepath{clip}%
\pgfsetbuttcap%
\pgfsetroundjoin%
\definecolor{currentfill}{rgb}{0.909804,0.698039,0.423529}%
\pgfsetfillcolor{currentfill}%
\pgfsetfillopacity{0.120000}%
\pgfsetlinewidth{0.000000pt}%
\definecolor{currentstroke}{rgb}{0.000000,0.000000,0.000000}%
\pgfsetstrokecolor{currentstroke}%
\pgfsetdash{}{0pt}%
\pgfpathmoveto{\pgfqpoint{0.600339in}{1.175842in}}%
\pgfpathlineto{\pgfqpoint{0.616436in}{1.158283in}}%
\pgfpathlineto{\pgfqpoint{0.632542in}{1.140920in}}%
\pgfpathlineto{\pgfqpoint{0.641443in}{1.142270in}}%
\pgfpathlineto{\pgfqpoint{0.650304in}{1.143933in}}%
\pgfpathlineto{\pgfqpoint{0.634240in}{1.161406in}}%
\pgfpathlineto{\pgfqpoint{0.618185in}{1.179069in}}%
\pgfpathlineto{\pgfqpoint{0.609282in}{1.177306in}}%
\pgfpathlineto{\pgfqpoint{0.600339in}{1.175842in}}%
\pgfpathclose%
\pgfusepath{fill}%
\end{pgfscope}%
\begin{pgfscope}%
\pgfpathrectangle{\pgfqpoint{0.407592in}{0.222904in}}{\pgfqpoint{1.411163in}{1.411163in}}%
\pgfusepath{clip}%
\pgfsetbuttcap%
\pgfsetroundjoin%
\definecolor{currentfill}{rgb}{0.909804,0.698039,0.423529}%
\pgfsetfillcolor{currentfill}%
\pgfsetfillopacity{0.120000}%
\pgfsetlinewidth{0.000000pt}%
\definecolor{currentstroke}{rgb}{0.000000,0.000000,0.000000}%
\pgfsetstrokecolor{currentstroke}%
\pgfsetdash{}{0pt}%
\pgfpathmoveto{\pgfqpoint{0.746891in}{1.044165in}}%
\pgfpathlineto{\pgfqpoint{0.763025in}{1.028607in}}%
\pgfpathlineto{\pgfqpoint{0.779170in}{1.013433in}}%
\pgfpathlineto{\pgfqpoint{0.787824in}{1.014951in}}%
\pgfpathlineto{\pgfqpoint{0.796441in}{1.016864in}}%
\pgfpathlineto{\pgfqpoint{0.780338in}{1.032164in}}%
\pgfpathlineto{\pgfqpoint{0.764246in}{1.047843in}}%
\pgfpathlineto{\pgfqpoint{0.755587in}{1.045816in}}%
\pgfpathlineto{\pgfqpoint{0.746891in}{1.044165in}}%
\pgfpathclose%
\pgfusepath{fill}%
\end{pgfscope}%
\begin{pgfscope}%
\pgfpathrectangle{\pgfqpoint{0.407592in}{0.222904in}}{\pgfqpoint{1.411163in}{1.411163in}}%
\pgfusepath{clip}%
\pgfsetbuttcap%
\pgfsetroundjoin%
\definecolor{currentfill}{rgb}{0.909804,0.698039,0.423529}%
\pgfsetfillcolor{currentfill}%
\pgfsetfillopacity{0.120000}%
\pgfsetlinewidth{0.000000pt}%
\definecolor{currentstroke}{rgb}{0.000000,0.000000,0.000000}%
\pgfsetstrokecolor{currentstroke}%
\pgfsetdash{}{0pt}%
\pgfpathmoveto{\pgfqpoint{1.494713in}{1.098906in}}%
\pgfpathlineto{\pgfqpoint{1.511712in}{1.105122in}}%
\pgfpathlineto{\pgfqpoint{1.528781in}{1.111770in}}%
\pgfpathlineto{\pgfqpoint{1.536367in}{1.123513in}}%
\pgfpathlineto{\pgfqpoint{1.543940in}{1.135473in}}%
\pgfpathlineto{\pgfqpoint{1.526890in}{1.129102in}}%
\pgfpathlineto{\pgfqpoint{1.509909in}{1.123162in}}%
\pgfpathlineto{\pgfqpoint{1.502317in}{1.110926in}}%
\pgfpathlineto{\pgfqpoint{1.494713in}{1.098906in}}%
\pgfpathclose%
\pgfusepath{fill}%
\end{pgfscope}%
\begin{pgfscope}%
\pgfpathrectangle{\pgfqpoint{0.407592in}{0.222904in}}{\pgfqpoint{1.411163in}{1.411163in}}%
\pgfusepath{clip}%
\pgfsetbuttcap%
\pgfsetroundjoin%
\definecolor{currentfill}{rgb}{0.909804,0.698039,0.423529}%
\pgfsetfillcolor{currentfill}%
\pgfsetfillopacity{0.120000}%
\pgfsetlinewidth{0.000000pt}%
\definecolor{currentstroke}{rgb}{0.000000,0.000000,0.000000}%
\pgfsetstrokecolor{currentstroke}%
\pgfsetdash{}{0pt}%
\pgfpathmoveto{\pgfqpoint{1.381282in}{1.008814in}}%
\pgfpathlineto{\pgfqpoint{1.398065in}{1.013812in}}%
\pgfpathlineto{\pgfqpoint{1.414915in}{1.019364in}}%
\pgfpathlineto{\pgfqpoint{1.422616in}{1.030086in}}%
\pgfpathlineto{\pgfqpoint{1.430303in}{1.041105in}}%
\pgfpathlineto{\pgfqpoint{1.413472in}{1.035808in}}%
\pgfpathlineto{\pgfqpoint{1.396706in}{1.031059in}}%
\pgfpathlineto{\pgfqpoint{1.389001in}{1.019785in}}%
\pgfpathlineto{\pgfqpoint{1.381282in}{1.008814in}}%
\pgfpathclose%
\pgfusepath{fill}%
\end{pgfscope}%
\begin{pgfscope}%
\pgfpathrectangle{\pgfqpoint{0.407592in}{0.222904in}}{\pgfqpoint{1.411163in}{1.411163in}}%
\pgfusepath{clip}%
\pgfsetbuttcap%
\pgfsetroundjoin%
\definecolor{currentfill}{rgb}{0.909804,0.698039,0.423529}%
\pgfsetfillcolor{currentfill}%
\pgfsetfillopacity{0.120000}%
\pgfsetlinewidth{0.000000pt}%
\definecolor{currentstroke}{rgb}{0.000000,0.000000,0.000000}%
\pgfsetstrokecolor{currentstroke}%
\pgfsetdash{}{0pt}%
\pgfpathmoveto{\pgfqpoint{1.170387in}{0.898045in}}%
\pgfpathlineto{\pgfqpoint{1.186799in}{0.897728in}}%
\pgfpathlineto{\pgfqpoint{1.203263in}{0.898348in}}%
\pgfpathlineto{\pgfqpoint{1.211209in}{0.905679in}}%
\pgfpathlineto{\pgfqpoint{1.219135in}{0.913520in}}%
\pgfpathlineto{\pgfqpoint{1.202694in}{0.912958in}}%
\pgfpathlineto{\pgfqpoint{1.186306in}{0.913316in}}%
\pgfpathlineto{\pgfqpoint{1.178357in}{0.905419in}}%
\pgfpathlineto{\pgfqpoint{1.170387in}{0.898045in}}%
\pgfpathclose%
\pgfusepath{fill}%
\end{pgfscope}%
\begin{pgfscope}%
\pgfpathrectangle{\pgfqpoint{0.407592in}{0.222904in}}{\pgfqpoint{1.411163in}{1.411163in}}%
\pgfusepath{clip}%
\pgfsetbuttcap%
\pgfsetroundjoin%
\definecolor{currentfill}{rgb}{0.909804,0.698039,0.423529}%
\pgfsetfillcolor{currentfill}%
\pgfsetfillopacity{0.120000}%
\pgfsetlinewidth{0.000000pt}%
\definecolor{currentstroke}{rgb}{0.000000,0.000000,0.000000}%
\pgfsetstrokecolor{currentstroke}%
\pgfsetdash{}{0pt}%
\pgfpathmoveto{\pgfqpoint{1.121647in}{0.888489in}}%
\pgfpathlineto{\pgfqpoint{1.137991in}{0.886203in}}%
\pgfpathlineto{\pgfqpoint{1.154381in}{0.884923in}}%
\pgfpathlineto{\pgfqpoint{1.162396in}{0.891209in}}%
\pgfpathlineto{\pgfqpoint{1.170387in}{0.898045in}}%
\pgfpathlineto{\pgfqpoint{1.154023in}{0.899330in}}%
\pgfpathlineto{\pgfqpoint{1.137705in}{0.901610in}}%
\pgfpathlineto{\pgfqpoint{1.129688in}{0.894771in}}%
\pgfpathlineto{\pgfqpoint{1.121647in}{0.888489in}}%
\pgfpathclose%
\pgfusepath{fill}%
\end{pgfscope}%
\begin{pgfscope}%
\pgfpathrectangle{\pgfqpoint{0.407592in}{0.222904in}}{\pgfqpoint{1.411163in}{1.411163in}}%
\pgfusepath{clip}%
\pgfsetbuttcap%
\pgfsetroundjoin%
\definecolor{currentfill}{rgb}{0.909804,0.698039,0.423529}%
\pgfsetfillcolor{currentfill}%
\pgfsetfillopacity{0.120000}%
\pgfsetlinewidth{0.000000pt}%
\definecolor{currentstroke}{rgb}{0.000000,0.000000,0.000000}%
\pgfsetstrokecolor{currentstroke}%
\pgfsetdash{}{0pt}%
\pgfpathmoveto{\pgfqpoint{0.632542in}{1.140920in}}%
\pgfpathlineto{\pgfqpoint{0.648658in}{1.123770in}}%
\pgfpathlineto{\pgfqpoint{0.664783in}{1.106854in}}%
\pgfpathlineto{\pgfqpoint{0.673641in}{1.108076in}}%
\pgfpathlineto{\pgfqpoint{0.682460in}{1.109626in}}%
\pgfpathlineto{\pgfqpoint{0.666377in}{1.126667in}}%
\pgfpathlineto{\pgfqpoint{0.650304in}{1.143933in}}%
\pgfpathlineto{\pgfqpoint{0.641443in}{1.142270in}}%
\pgfpathlineto{\pgfqpoint{0.632542in}{1.140920in}}%
\pgfpathclose%
\pgfusepath{fill}%
\end{pgfscope}%
\begin{pgfscope}%
\pgfpathrectangle{\pgfqpoint{0.407592in}{0.222904in}}{\pgfqpoint{1.411163in}{1.411163in}}%
\pgfusepath{clip}%
\pgfsetbuttcap%
\pgfsetroundjoin%
\definecolor{currentfill}{rgb}{0.909804,0.698039,0.423529}%
\pgfsetfillcolor{currentfill}%
\pgfsetfillopacity{0.120000}%
\pgfsetlinewidth{0.000000pt}%
\definecolor{currentstroke}{rgb}{0.000000,0.000000,0.000000}%
\pgfsetstrokecolor{currentstroke}%
\pgfsetdash{}{0pt}%
\pgfpathmoveto{\pgfqpoint{1.219135in}{0.913520in}}%
\pgfpathlineto{\pgfqpoint{1.235629in}{0.914966in}}%
\pgfpathlineto{\pgfqpoint{1.252180in}{0.917253in}}%
\pgfpathlineto{\pgfqpoint{1.260065in}{0.925489in}}%
\pgfpathlineto{\pgfqpoint{1.267932in}{0.934188in}}%
\pgfpathlineto{\pgfqpoint{1.251402in}{0.932017in}}%
\pgfpathlineto{\pgfqpoint{1.234929in}{0.930672in}}%
\pgfpathlineto{\pgfqpoint{1.227041in}{0.921858in}}%
\pgfpathlineto{\pgfqpoint{1.219135in}{0.913520in}}%
\pgfpathclose%
\pgfusepath{fill}%
\end{pgfscope}%
\begin{pgfscope}%
\pgfpathrectangle{\pgfqpoint{0.407592in}{0.222904in}}{\pgfqpoint{1.411163in}{1.411163in}}%
\pgfusepath{clip}%
\pgfsetbuttcap%
\pgfsetroundjoin%
\definecolor{currentfill}{rgb}{0.909804,0.698039,0.423529}%
\pgfsetfillcolor{currentfill}%
\pgfsetfillopacity{0.120000}%
\pgfsetlinewidth{0.000000pt}%
\definecolor{currentstroke}{rgb}{0.000000,0.000000,0.000000}%
\pgfsetstrokecolor{currentstroke}%
\pgfsetdash{}{0pt}%
\pgfpathmoveto{\pgfqpoint{0.942541in}{0.917303in}}%
\pgfpathlineto{\pgfqpoint{0.958738in}{0.907261in}}%
\pgfpathlineto{\pgfqpoint{0.974957in}{0.898065in}}%
\pgfpathlineto{\pgfqpoint{0.983268in}{0.900935in}}%
\pgfpathlineto{\pgfqpoint{0.991546in}{0.904343in}}%
\pgfpathlineto{\pgfqpoint{0.975365in}{0.913571in}}%
\pgfpathlineto{\pgfqpoint{0.959206in}{0.923650in}}%
\pgfpathlineto{\pgfqpoint{0.950890in}{0.920217in}}%
\pgfpathlineto{\pgfqpoint{0.942541in}{0.917303in}}%
\pgfpathclose%
\pgfusepath{fill}%
\end{pgfscope}%
\begin{pgfscope}%
\pgfpathrectangle{\pgfqpoint{0.407592in}{0.222904in}}{\pgfqpoint{1.411163in}{1.411163in}}%
\pgfusepath{clip}%
\pgfsetbuttcap%
\pgfsetroundjoin%
\definecolor{currentfill}{rgb}{0.909804,0.698039,0.423529}%
\pgfsetfillcolor{currentfill}%
\pgfsetfillopacity{0.120000}%
\pgfsetlinewidth{0.000000pt}%
\definecolor{currentstroke}{rgb}{0.000000,0.000000,0.000000}%
\pgfsetstrokecolor{currentstroke}%
\pgfsetdash{}{0pt}%
\pgfpathmoveto{\pgfqpoint{0.860975in}{0.960214in}}%
\pgfpathlineto{\pgfqpoint{0.877143in}{0.947404in}}%
\pgfpathlineto{\pgfqpoint{0.893326in}{0.935230in}}%
\pgfpathlineto{\pgfqpoint{0.901783in}{0.937234in}}%
\pgfpathlineto{\pgfqpoint{0.910205in}{0.939716in}}%
\pgfpathlineto{\pgfqpoint{0.894062in}{0.951983in}}%
\pgfpathlineto{\pgfqpoint{0.877935in}{0.964887in}}%
\pgfpathlineto{\pgfqpoint{0.869473in}{0.962323in}}%
\pgfpathlineto{\pgfqpoint{0.860975in}{0.960214in}}%
\pgfpathclose%
\pgfusepath{fill}%
\end{pgfscope}%
\begin{pgfscope}%
\pgfpathrectangle{\pgfqpoint{0.407592in}{0.222904in}}{\pgfqpoint{1.411163in}{1.411163in}}%
\pgfusepath{clip}%
\pgfsetbuttcap%
\pgfsetroundjoin%
\definecolor{currentfill}{rgb}{0.909804,0.698039,0.423529}%
\pgfsetfillcolor{currentfill}%
\pgfsetfillopacity{0.120000}%
\pgfsetlinewidth{0.000000pt}%
\definecolor{currentstroke}{rgb}{0.000000,0.000000,0.000000}%
\pgfsetstrokecolor{currentstroke}%
\pgfsetdash{}{0pt}%
\pgfpathmoveto{\pgfqpoint{1.072863in}{0.885323in}}%
\pgfpathlineto{\pgfqpoint{1.089157in}{0.880962in}}%
\pgfpathlineto{\pgfqpoint{1.105490in}{0.877632in}}%
\pgfpathlineto{\pgfqpoint{1.113581in}{0.882774in}}%
\pgfpathlineto{\pgfqpoint{1.121647in}{0.888489in}}%
\pgfpathlineto{\pgfqpoint{1.105343in}{0.891796in}}%
\pgfpathlineto{\pgfqpoint{1.089079in}{0.896130in}}%
\pgfpathlineto{\pgfqpoint{1.080984in}{0.890440in}}%
\pgfpathlineto{\pgfqpoint{1.072863in}{0.885323in}}%
\pgfpathclose%
\pgfusepath{fill}%
\end{pgfscope}%
\begin{pgfscope}%
\pgfpathrectangle{\pgfqpoint{0.407592in}{0.222904in}}{\pgfqpoint{1.411163in}{1.411163in}}%
\pgfusepath{clip}%
\pgfsetbuttcap%
\pgfsetroundjoin%
\definecolor{currentfill}{rgb}{0.909804,0.698039,0.423529}%
\pgfsetfillcolor{currentfill}%
\pgfsetfillopacity{0.120000}%
\pgfsetlinewidth{0.000000pt}%
\definecolor{currentstroke}{rgb}{0.000000,0.000000,0.000000}%
\pgfsetstrokecolor{currentstroke}%
\pgfsetdash{}{0pt}%
\pgfpathmoveto{\pgfqpoint{1.267932in}{0.934188in}}%
\pgfpathlineto{\pgfqpoint{1.284520in}{0.937142in}}%
\pgfpathlineto{\pgfqpoint{1.301169in}{0.940834in}}%
\pgfpathlineto{\pgfqpoint{1.308999in}{0.949822in}}%
\pgfpathlineto{\pgfqpoint{1.316812in}{0.959225in}}%
\pgfpathlineto{\pgfqpoint{1.300182in}{0.955702in}}%
\pgfpathlineto{\pgfqpoint{1.283614in}{0.952905in}}%
\pgfpathlineto{\pgfqpoint{1.275781in}{0.943333in}}%
\pgfpathlineto{\pgfqpoint{1.267932in}{0.934188in}}%
\pgfpathclose%
\pgfusepath{fill}%
\end{pgfscope}%
\begin{pgfscope}%
\pgfpathrectangle{\pgfqpoint{0.407592in}{0.222904in}}{\pgfqpoint{1.411163in}{1.411163in}}%
\pgfusepath{clip}%
\pgfsetbuttcap%
\pgfsetroundjoin%
\definecolor{currentfill}{rgb}{0.909804,0.698039,0.423529}%
\pgfsetfillcolor{currentfill}%
\pgfsetfillopacity{0.120000}%
\pgfsetlinewidth{0.000000pt}%
\definecolor{currentstroke}{rgb}{0.000000,0.000000,0.000000}%
\pgfsetstrokecolor{currentstroke}%
\pgfsetdash{}{0pt}%
\pgfpathmoveto{\pgfqpoint{1.608348in}{1.198103in}}%
\pgfpathlineto{\pgfqpoint{1.625572in}{1.205090in}}%
\pgfpathlineto{\pgfqpoint{1.642868in}{1.212429in}}%
\pgfpathlineto{\pgfqpoint{1.650343in}{1.224772in}}%
\pgfpathlineto{\pgfqpoint{1.633057in}{1.217568in}}%
\pgfpathlineto{\pgfqpoint{1.615843in}{1.210716in}}%
\pgfusepath{fill}%
\end{pgfscope}%
\begin{pgfscope}%
\pgfpathrectangle{\pgfqpoint{0.407592in}{0.222904in}}{\pgfqpoint{1.411163in}{1.411163in}}%
\pgfusepath{clip}%
\pgfsetbuttcap%
\pgfsetroundjoin%
\definecolor{currentfill}{rgb}{0.909804,0.698039,0.423529}%
\pgfsetfillcolor{currentfill}%
\pgfsetfillopacity{0.120000}%
\pgfsetlinewidth{0.000000pt}%
\definecolor{currentstroke}{rgb}{0.000000,0.000000,0.000000}%
\pgfsetstrokecolor{currentstroke}%
\pgfsetdash{}{0pt}%
\pgfpathmoveto{\pgfqpoint{1.430303in}{1.041105in}}%
\pgfpathlineto{\pgfqpoint{1.447202in}{1.046920in}}%
\pgfpathlineto{\pgfqpoint{1.464170in}{1.053222in}}%
\pgfpathlineto{\pgfqpoint{1.471825in}{1.064258in}}%
\pgfpathlineto{\pgfqpoint{1.479467in}{1.075559in}}%
\pgfpathlineto{\pgfqpoint{1.462519in}{1.069526in}}%
\pgfpathlineto{\pgfqpoint{1.445638in}{1.063977in}}%
\pgfpathlineto{\pgfqpoint{1.437977in}{1.052407in}}%
\pgfpathlineto{\pgfqpoint{1.430303in}{1.041105in}}%
\pgfpathclose%
\pgfusepath{fill}%
\end{pgfscope}%
\begin{pgfscope}%
\pgfpathrectangle{\pgfqpoint{0.407592in}{0.222904in}}{\pgfqpoint{1.411163in}{1.411163in}}%
\pgfusepath{clip}%
\pgfsetbuttcap%
\pgfsetroundjoin%
\definecolor{currentfill}{rgb}{0.909804,0.698039,0.423529}%
\pgfsetfillcolor{currentfill}%
\pgfsetfillopacity{0.120000}%
\pgfsetlinewidth{0.000000pt}%
\definecolor{currentstroke}{rgb}{0.000000,0.000000,0.000000}%
\pgfsetstrokecolor{currentstroke}%
\pgfsetdash{}{0pt}%
\pgfpathmoveto{\pgfqpoint{0.779170in}{1.013433in}}%
\pgfpathlineto{\pgfqpoint{0.795327in}{0.998679in}}%
\pgfpathlineto{\pgfqpoint{0.811495in}{0.984382in}}%
\pgfpathlineto{\pgfqpoint{0.820108in}{0.985756in}}%
\pgfpathlineto{\pgfqpoint{0.828683in}{0.987546in}}%
\pgfpathlineto{\pgfqpoint{0.812556in}{1.001978in}}%
\pgfpathlineto{\pgfqpoint{0.796441in}{1.016864in}}%
\pgfpathlineto{\pgfqpoint{0.787824in}{1.014951in}}%
\pgfpathlineto{\pgfqpoint{0.779170in}{1.013433in}}%
\pgfpathclose%
\pgfusepath{fill}%
\end{pgfscope}%
\begin{pgfscope}%
\pgfpathrectangle{\pgfqpoint{0.407592in}{0.222904in}}{\pgfqpoint{1.411163in}{1.411163in}}%
\pgfusepath{clip}%
\pgfsetbuttcap%
\pgfsetroundjoin%
\definecolor{currentfill}{rgb}{0.909804,0.698039,0.423529}%
\pgfsetfillcolor{currentfill}%
\pgfsetfillopacity{0.120000}%
\pgfsetlinewidth{0.000000pt}%
\definecolor{currentstroke}{rgb}{0.000000,0.000000,0.000000}%
\pgfsetstrokecolor{currentstroke}%
\pgfsetdash{}{0pt}%
\pgfpathmoveto{\pgfqpoint{0.664783in}{1.106854in}}%
\pgfpathlineto{\pgfqpoint{0.680918in}{1.090190in}}%
\pgfpathlineto{\pgfqpoint{0.697063in}{1.073803in}}%
\pgfpathlineto{\pgfqpoint{0.705879in}{1.074882in}}%
\pgfpathlineto{\pgfqpoint{0.714655in}{1.076306in}}%
\pgfpathlineto{\pgfqpoint{0.698553in}{1.092831in}}%
\pgfpathlineto{\pgfqpoint{0.682460in}{1.109626in}}%
\pgfpathlineto{\pgfqpoint{0.673641in}{1.108076in}}%
\pgfpathlineto{\pgfqpoint{0.664783in}{1.106854in}}%
\pgfpathclose%
\pgfusepath{fill}%
\end{pgfscope}%
\begin{pgfscope}%
\pgfpathrectangle{\pgfqpoint{0.407592in}{0.222904in}}{\pgfqpoint{1.411163in}{1.411163in}}%
\pgfusepath{clip}%
\pgfsetbuttcap%
\pgfsetroundjoin%
\definecolor{currentfill}{rgb}{0.909804,0.698039,0.423529}%
\pgfsetfillcolor{currentfill}%
\pgfsetfillopacity{0.120000}%
\pgfsetlinewidth{0.000000pt}%
\definecolor{currentstroke}{rgb}{0.000000,0.000000,0.000000}%
\pgfsetstrokecolor{currentstroke}%
\pgfsetdash{}{0pt}%
\pgfpathmoveto{\pgfqpoint{1.543940in}{1.135473in}}%
\pgfpathlineto{\pgfqpoint{1.561060in}{1.142252in}}%
\pgfpathlineto{\pgfqpoint{1.578252in}{1.149419in}}%
\pgfpathlineto{\pgfqpoint{1.585794in}{1.161307in}}%
\pgfpathlineto{\pgfqpoint{1.593323in}{1.173389in}}%
\pgfpathlineto{\pgfqpoint{1.576151in}{1.166498in}}%
\pgfpathlineto{\pgfqpoint{1.559050in}{1.159996in}}%
\pgfpathlineto{\pgfqpoint{1.551501in}{1.147637in}}%
\pgfpathlineto{\pgfqpoint{1.543940in}{1.135473in}}%
\pgfpathclose%
\pgfusepath{fill}%
\end{pgfscope}%
\begin{pgfscope}%
\pgfpathrectangle{\pgfqpoint{0.407592in}{0.222904in}}{\pgfqpoint{1.411163in}{1.411163in}}%
\pgfusepath{clip}%
\pgfsetbuttcap%
\pgfsetroundjoin%
\definecolor{currentfill}{rgb}{0.909804,0.698039,0.423529}%
\pgfsetfillcolor{currentfill}%
\pgfsetfillopacity{0.120000}%
\pgfsetlinewidth{0.000000pt}%
\definecolor{currentstroke}{rgb}{0.000000,0.000000,0.000000}%
\pgfsetstrokecolor{currentstroke}%
\pgfsetdash{}{0pt}%
\pgfpathmoveto{\pgfqpoint{1.023983in}{0.888627in}}%
\pgfpathlineto{\pgfqpoint{1.040244in}{0.882217in}}%
\pgfpathlineto{\pgfqpoint{1.056536in}{0.876811in}}%
\pgfpathlineto{\pgfqpoint{1.064714in}{0.880781in}}%
\pgfpathlineto{\pgfqpoint{1.072863in}{0.885323in}}%
\pgfpathlineto{\pgfqpoint{1.056604in}{0.890710in}}%
\pgfpathlineto{\pgfqpoint{1.040377in}{0.897104in}}%
\pgfpathlineto{\pgfqpoint{1.032195in}{0.892583in}}%
\pgfpathlineto{\pgfqpoint{1.023983in}{0.888627in}}%
\pgfpathclose%
\pgfusepath{fill}%
\end{pgfscope}%
\begin{pgfscope}%
\pgfpathrectangle{\pgfqpoint{0.407592in}{0.222904in}}{\pgfqpoint{1.411163in}{1.411163in}}%
\pgfusepath{clip}%
\pgfsetbuttcap%
\pgfsetroundjoin%
\definecolor{currentfill}{rgb}{0.909804,0.698039,0.423529}%
\pgfsetfillcolor{currentfill}%
\pgfsetfillopacity{0.120000}%
\pgfsetlinewidth{0.000000pt}%
\definecolor{currentstroke}{rgb}{0.000000,0.000000,0.000000}%
\pgfsetstrokecolor{currentstroke}%
\pgfsetdash{}{0pt}%
\pgfpathmoveto{\pgfqpoint{1.316812in}{0.959225in}}%
\pgfpathlineto{\pgfqpoint{1.333504in}{0.963432in}}%
\pgfpathlineto{\pgfqpoint{1.350260in}{0.968282in}}%
\pgfpathlineto{\pgfqpoint{1.358038in}{0.977879in}}%
\pgfpathlineto{\pgfqpoint{1.365800in}{0.987845in}}%
\pgfpathlineto{\pgfqpoint{1.349063in}{0.983206in}}%
\pgfpathlineto{\pgfqpoint{1.332390in}{0.979202in}}%
\pgfpathlineto{\pgfqpoint{1.324609in}{0.969024in}}%
\pgfpathlineto{\pgfqpoint{1.316812in}{0.959225in}}%
\pgfpathclose%
\pgfusepath{fill}%
\end{pgfscope}%
\begin{pgfscope}%
\pgfpathrectangle{\pgfqpoint{0.407592in}{0.222904in}}{\pgfqpoint{1.411163in}{1.411163in}}%
\pgfusepath{clip}%
\pgfsetbuttcap%
\pgfsetroundjoin%
\definecolor{currentfill}{rgb}{0.909804,0.698039,0.423529}%
\pgfsetfillcolor{currentfill}%
\pgfsetfillopacity{0.120000}%
\pgfsetlinewidth{0.000000pt}%
\definecolor{currentstroke}{rgb}{0.000000,0.000000,0.000000}%
\pgfsetstrokecolor{currentstroke}%
\pgfsetdash{}{0pt}%
\pgfpathmoveto{\pgfqpoint{0.550083in}{1.209196in}}%
\pgfpathlineto{\pgfqpoint{0.566204in}{1.191407in}}%
\pgfpathlineto{\pgfqpoint{0.582334in}{1.173790in}}%
\pgfpathlineto{\pgfqpoint{0.591357in}{1.174672in}}%
\pgfpathlineto{\pgfqpoint{0.600339in}{1.175842in}}%
\pgfpathlineto{\pgfqpoint{0.584252in}{1.193582in}}%
\pgfpathlineto{\pgfqpoint{0.568173in}{1.211487in}}%
\pgfpathlineto{\pgfqpoint{0.559148in}{1.210203in}}%
\pgfpathlineto{\pgfqpoint{0.550083in}{1.209196in}}%
\pgfpathclose%
\pgfusepath{fill}%
\end{pgfscope}%
\begin{pgfscope}%
\pgfpathrectangle{\pgfqpoint{0.407592in}{0.222904in}}{\pgfqpoint{1.411163in}{1.411163in}}%
\pgfusepath{clip}%
\pgfsetbuttcap%
\pgfsetroundjoin%
\definecolor{currentfill}{rgb}{0.909804,0.698039,0.423529}%
\pgfsetfillcolor{currentfill}%
\pgfsetfillopacity{0.120000}%
\pgfsetlinewidth{0.000000pt}%
\definecolor{currentstroke}{rgb}{0.000000,0.000000,0.000000}%
\pgfsetstrokecolor{currentstroke}%
\pgfsetdash{}{0pt}%
\pgfpathmoveto{\pgfqpoint{0.893326in}{0.935230in}}%
\pgfpathlineto{\pgfqpoint{0.909525in}{0.923743in}}%
\pgfpathlineto{\pgfqpoint{0.925742in}{0.912991in}}%
\pgfpathlineto{\pgfqpoint{0.934159in}{0.914898in}}%
\pgfpathlineto{\pgfqpoint{0.942541in}{0.917303in}}%
\pgfpathlineto{\pgfqpoint{0.926364in}{0.928138in}}%
\pgfpathlineto{\pgfqpoint{0.910205in}{0.939716in}}%
\pgfpathlineto{\pgfqpoint{0.901783in}{0.937234in}}%
\pgfpathlineto{\pgfqpoint{0.893326in}{0.935230in}}%
\pgfpathclose%
\pgfusepath{fill}%
\end{pgfscope}%
\begin{pgfscope}%
\pgfpathrectangle{\pgfqpoint{0.407592in}{0.222904in}}{\pgfqpoint{1.411163in}{1.411163in}}%
\pgfusepath{clip}%
\pgfsetbuttcap%
\pgfsetroundjoin%
\definecolor{currentfill}{rgb}{0.909804,0.698039,0.423529}%
\pgfsetfillcolor{currentfill}%
\pgfsetfillopacity{0.120000}%
\pgfsetlinewidth{0.000000pt}%
\definecolor{currentstroke}{rgb}{0.000000,0.000000,0.000000}%
\pgfsetstrokecolor{currentstroke}%
\pgfsetdash{}{0pt}%
\pgfpathmoveto{\pgfqpoint{0.697063in}{1.073803in}}%
\pgfpathlineto{\pgfqpoint{0.713218in}{1.057717in}}%
\pgfpathlineto{\pgfqpoint{0.729383in}{1.041960in}}%
\pgfpathlineto{\pgfqpoint{0.738157in}{1.042882in}}%
\pgfpathlineto{\pgfqpoint{0.746891in}{1.044165in}}%
\pgfpathlineto{\pgfqpoint{0.730768in}{1.060074in}}%
\pgfpathlineto{\pgfqpoint{0.714655in}{1.076306in}}%
\pgfpathlineto{\pgfqpoint{0.705879in}{1.074882in}}%
\pgfpathlineto{\pgfqpoint{0.697063in}{1.073803in}}%
\pgfpathclose%
\pgfusepath{fill}%
\end{pgfscope}%
\begin{pgfscope}%
\pgfpathrectangle{\pgfqpoint{0.407592in}{0.222904in}}{\pgfqpoint{1.411163in}{1.411163in}}%
\pgfusepath{clip}%
\pgfsetbuttcap%
\pgfsetroundjoin%
\definecolor{currentfill}{rgb}{0.909804,0.698039,0.423529}%
\pgfsetfillcolor{currentfill}%
\pgfsetfillopacity{0.120000}%
\pgfsetlinewidth{0.000000pt}%
\definecolor{currentstroke}{rgb}{0.000000,0.000000,0.000000}%
\pgfsetstrokecolor{currentstroke}%
\pgfsetdash{}{0pt}%
\pgfpathmoveto{\pgfqpoint{1.479467in}{1.075559in}}%
\pgfpathlineto{\pgfqpoint{1.496485in}{1.082050in}}%
\pgfpathlineto{\pgfqpoint{1.513573in}{1.088973in}}%
\pgfpathlineto{\pgfqpoint{1.521184in}{1.100252in}}%
\pgfpathlineto{\pgfqpoint{1.528781in}{1.111770in}}%
\pgfpathlineto{\pgfqpoint{1.511712in}{1.105122in}}%
\pgfpathlineto{\pgfqpoint{1.494713in}{1.098906in}}%
\pgfpathlineto{\pgfqpoint{1.487097in}{1.087113in}}%
\pgfpathlineto{\pgfqpoint{1.479467in}{1.075559in}}%
\pgfpathclose%
\pgfusepath{fill}%
\end{pgfscope}%
\begin{pgfscope}%
\pgfpathrectangle{\pgfqpoint{0.407592in}{0.222904in}}{\pgfqpoint{1.411163in}{1.411163in}}%
\pgfusepath{clip}%
\pgfsetbuttcap%
\pgfsetroundjoin%
\definecolor{currentfill}{rgb}{0.909804,0.698039,0.423529}%
\pgfsetfillcolor{currentfill}%
\pgfsetfillopacity{0.120000}%
\pgfsetlinewidth{0.000000pt}%
\definecolor{currentstroke}{rgb}{0.000000,0.000000,0.000000}%
\pgfsetstrokecolor{currentstroke}%
\pgfsetdash{}{0pt}%
\pgfpathmoveto{\pgfqpoint{0.974957in}{0.898065in}}%
\pgfpathlineto{\pgfqpoint{0.991200in}{0.889760in}}%
\pgfpathlineto{\pgfqpoint{1.007468in}{0.882387in}}%
\pgfpathlineto{\pgfqpoint{1.015741in}{0.885231in}}%
\pgfpathlineto{\pgfqpoint{1.023983in}{0.888627in}}%
\pgfpathlineto{\pgfqpoint{1.007751in}{0.896014in}}%
\pgfpathlineto{\pgfqpoint{0.991546in}{0.904343in}}%
\pgfpathlineto{\pgfqpoint{0.983268in}{0.900935in}}%
\pgfpathlineto{\pgfqpoint{0.974957in}{0.898065in}}%
\pgfpathclose%
\pgfusepath{fill}%
\end{pgfscope}%
\begin{pgfscope}%
\pgfpathrectangle{\pgfqpoint{0.407592in}{0.222904in}}{\pgfqpoint{1.411163in}{1.411163in}}%
\pgfusepath{clip}%
\pgfsetbuttcap%
\pgfsetroundjoin%
\definecolor{currentfill}{rgb}{0.909804,0.698039,0.423529}%
\pgfsetfillcolor{currentfill}%
\pgfsetfillopacity{0.120000}%
\pgfsetlinewidth{0.000000pt}%
\definecolor{currentstroke}{rgb}{0.000000,0.000000,0.000000}%
\pgfsetstrokecolor{currentstroke}%
\pgfsetdash{}{0pt}%
\pgfpathmoveto{\pgfqpoint{1.365800in}{0.987845in}}%
\pgfpathlineto{\pgfqpoint{1.382603in}{0.993079in}}%
\pgfpathlineto{\pgfqpoint{1.399472in}{0.998872in}}%
\pgfpathlineto{\pgfqpoint{1.407201in}{1.008955in}}%
\pgfpathlineto{\pgfqpoint{1.414915in}{1.019364in}}%
\pgfpathlineto{\pgfqpoint{1.398065in}{1.013812in}}%
\pgfpathlineto{\pgfqpoint{1.381282in}{1.008814in}}%
\pgfpathlineto{\pgfqpoint{1.373548in}{0.998162in}}%
\pgfpathlineto{\pgfqpoint{1.365800in}{0.987845in}}%
\pgfpathclose%
\pgfusepath{fill}%
\end{pgfscope}%
\begin{pgfscope}%
\pgfpathrectangle{\pgfqpoint{0.407592in}{0.222904in}}{\pgfqpoint{1.411163in}{1.411163in}}%
\pgfusepath{clip}%
\pgfsetbuttcap%
\pgfsetroundjoin%
\definecolor{currentfill}{rgb}{0.909804,0.698039,0.423529}%
\pgfsetfillcolor{currentfill}%
\pgfsetfillopacity{0.120000}%
\pgfsetlinewidth{0.000000pt}%
\definecolor{currentstroke}{rgb}{0.000000,0.000000,0.000000}%
\pgfsetstrokecolor{currentstroke}%
\pgfsetdash{}{0pt}%
\pgfpathmoveto{\pgfqpoint{0.811495in}{0.984382in}}%
\pgfpathlineto{\pgfqpoint{0.827676in}{0.970582in}}%
\pgfpathlineto{\pgfqpoint{0.843871in}{0.957324in}}%
\pgfpathlineto{\pgfqpoint{0.852442in}{0.958551in}}%
\pgfpathlineto{\pgfqpoint{0.860975in}{0.960214in}}%
\pgfpathlineto{\pgfqpoint{0.844822in}{0.973610in}}%
\pgfpathlineto{\pgfqpoint{0.828683in}{0.987546in}}%
\pgfpathlineto{\pgfqpoint{0.820108in}{0.985756in}}%
\pgfpathlineto{\pgfqpoint{0.811495in}{0.984382in}}%
\pgfpathclose%
\pgfusepath{fill}%
\end{pgfscope}%
\begin{pgfscope}%
\pgfpathrectangle{\pgfqpoint{0.407592in}{0.222904in}}{\pgfqpoint{1.411163in}{1.411163in}}%
\pgfusepath{clip}%
\pgfsetbuttcap%
\pgfsetroundjoin%
\definecolor{currentfill}{rgb}{0.909804,0.698039,0.423529}%
\pgfsetfillcolor{currentfill}%
\pgfsetfillopacity{0.120000}%
\pgfsetlinewidth{0.000000pt}%
\definecolor{currentstroke}{rgb}{0.000000,0.000000,0.000000}%
\pgfsetstrokecolor{currentstroke}%
\pgfsetdash{}{0pt}%
\pgfpathmoveto{\pgfqpoint{0.582334in}{1.173790in}}%
\pgfpathlineto{\pgfqpoint{0.598473in}{1.156362in}}%
\pgfpathlineto{\pgfqpoint{0.614622in}{1.139137in}}%
\pgfpathlineto{\pgfqpoint{0.623602in}{1.139877in}}%
\pgfpathlineto{\pgfqpoint{0.632542in}{1.140920in}}%
\pgfpathlineto{\pgfqpoint{0.616436in}{1.158283in}}%
\pgfpathlineto{\pgfqpoint{0.600339in}{1.175842in}}%
\pgfpathlineto{\pgfqpoint{0.591357in}{1.174672in}}%
\pgfpathlineto{\pgfqpoint{0.582334in}{1.173790in}}%
\pgfpathclose%
\pgfusepath{fill}%
\end{pgfscope}%
\begin{pgfscope}%
\pgfpathrectangle{\pgfqpoint{0.407592in}{0.222904in}}{\pgfqpoint{1.411163in}{1.411163in}}%
\pgfusepath{clip}%
\pgfsetbuttcap%
\pgfsetroundjoin%
\definecolor{currentfill}{rgb}{0.909804,0.698039,0.423529}%
\pgfsetfillcolor{currentfill}%
\pgfsetfillopacity{0.120000}%
\pgfsetlinewidth{0.000000pt}%
\definecolor{currentstroke}{rgb}{0.000000,0.000000,0.000000}%
\pgfsetstrokecolor{currentstroke}%
\pgfsetdash{}{0pt}%
\pgfpathmoveto{\pgfqpoint{1.593323in}{1.173389in}}%
\pgfpathlineto{\pgfqpoint{1.610567in}{1.180649in}}%
\pgfpathlineto{\pgfqpoint{1.627883in}{1.188260in}}%
\pgfpathlineto{\pgfqpoint{1.635381in}{1.200256in}}%
\pgfpathlineto{\pgfqpoint{1.642868in}{1.212429in}}%
\pgfpathlineto{\pgfqpoint{1.625572in}{1.205090in}}%
\pgfpathlineto{\pgfqpoint{1.608348in}{1.198103in}}%
\pgfpathlineto{\pgfqpoint{1.600842in}{1.185658in}}%
\pgfpathlineto{\pgfqpoint{1.593323in}{1.173389in}}%
\pgfpathclose%
\pgfusepath{fill}%
\end{pgfscope}%
\begin{pgfscope}%
\pgfpathrectangle{\pgfqpoint{0.407592in}{0.222904in}}{\pgfqpoint{1.411163in}{1.411163in}}%
\pgfusepath{clip}%
\pgfsetbuttcap%
\pgfsetroundjoin%
\definecolor{currentfill}{rgb}{0.909804,0.698039,0.423529}%
\pgfsetfillcolor{currentfill}%
\pgfsetfillopacity{0.120000}%
\pgfsetlinewidth{0.000000pt}%
\definecolor{currentstroke}{rgb}{0.000000,0.000000,0.000000}%
\pgfsetstrokecolor{currentstroke}%
\pgfsetdash{}{0pt}%
\pgfpathmoveto{\pgfqpoint{1.154381in}{0.884923in}}%
\pgfpathlineto{\pgfqpoint{1.170819in}{0.884624in}}%
\pgfpathlineto{\pgfqpoint{1.187307in}{0.885272in}}%
\pgfpathlineto{\pgfqpoint{1.195296in}{0.891542in}}%
\pgfpathlineto{\pgfqpoint{1.203263in}{0.898348in}}%
\pgfpathlineto{\pgfqpoint{1.186799in}{0.897728in}}%
\pgfpathlineto{\pgfqpoint{1.170387in}{0.898045in}}%
\pgfpathlineto{\pgfqpoint{1.162396in}{0.891209in}}%
\pgfpathlineto{\pgfqpoint{1.154381in}{0.884923in}}%
\pgfpathclose%
\pgfusepath{fill}%
\end{pgfscope}%
\begin{pgfscope}%
\pgfpathrectangle{\pgfqpoint{0.407592in}{0.222904in}}{\pgfqpoint{1.411163in}{1.411163in}}%
\pgfusepath{clip}%
\pgfsetbuttcap%
\pgfsetroundjoin%
\definecolor{currentfill}{rgb}{0.909804,0.698039,0.423529}%
\pgfsetfillcolor{currentfill}%
\pgfsetfillopacity{0.120000}%
\pgfsetlinewidth{0.000000pt}%
\definecolor{currentstroke}{rgb}{0.000000,0.000000,0.000000}%
\pgfsetstrokecolor{currentstroke}%
\pgfsetdash{}{0pt}%
\pgfpathmoveto{\pgfqpoint{1.203263in}{0.898348in}}%
\pgfpathlineto{\pgfqpoint{1.219780in}{0.899865in}}%
\pgfpathlineto{\pgfqpoint{1.236353in}{0.902235in}}%
\pgfpathlineto{\pgfqpoint{1.244276in}{0.909497in}}%
\pgfpathlineto{\pgfqpoint{1.252180in}{0.917253in}}%
\pgfpathlineto{\pgfqpoint{1.235629in}{0.914966in}}%
\pgfpathlineto{\pgfqpoint{1.219135in}{0.913520in}}%
\pgfpathlineto{\pgfqpoint{1.211209in}{0.905679in}}%
\pgfpathlineto{\pgfqpoint{1.203263in}{0.898348in}}%
\pgfpathclose%
\pgfusepath{fill}%
\end{pgfscope}%
\begin{pgfscope}%
\pgfpathrectangle{\pgfqpoint{0.407592in}{0.222904in}}{\pgfqpoint{1.411163in}{1.411163in}}%
\pgfusepath{clip}%
\pgfsetbuttcap%
\pgfsetroundjoin%
\definecolor{currentfill}{rgb}{0.909804,0.698039,0.423529}%
\pgfsetfillcolor{currentfill}%
\pgfsetfillopacity{0.120000}%
\pgfsetlinewidth{0.000000pt}%
\definecolor{currentstroke}{rgb}{0.000000,0.000000,0.000000}%
\pgfsetstrokecolor{currentstroke}%
\pgfsetdash{}{0pt}%
\pgfpathmoveto{\pgfqpoint{1.105490in}{0.877632in}}%
\pgfpathlineto{\pgfqpoint{1.121863in}{0.875328in}}%
\pgfpathlineto{\pgfqpoint{1.138281in}{0.874035in}}%
\pgfpathlineto{\pgfqpoint{1.146344in}{0.879197in}}%
\pgfpathlineto{\pgfqpoint{1.154381in}{0.884923in}}%
\pgfpathlineto{\pgfqpoint{1.137991in}{0.886203in}}%
\pgfpathlineto{\pgfqpoint{1.121647in}{0.888489in}}%
\pgfpathlineto{\pgfqpoint{1.113581in}{0.882774in}}%
\pgfpathlineto{\pgfqpoint{1.105490in}{0.877632in}}%
\pgfpathclose%
\pgfusepath{fill}%
\end{pgfscope}%
\begin{pgfscope}%
\pgfpathrectangle{\pgfqpoint{0.407592in}{0.222904in}}{\pgfqpoint{1.411163in}{1.411163in}}%
\pgfusepath{clip}%
\pgfsetbuttcap%
\pgfsetroundjoin%
\definecolor{currentfill}{rgb}{0.909804,0.698039,0.423529}%
\pgfsetfillcolor{currentfill}%
\pgfsetfillopacity{0.120000}%
\pgfsetlinewidth{0.000000pt}%
\definecolor{currentstroke}{rgb}{0.000000,0.000000,0.000000}%
\pgfsetstrokecolor{currentstroke}%
\pgfsetdash{}{0pt}%
\pgfpathmoveto{\pgfqpoint{0.729383in}{1.041960in}}%
\pgfpathlineto{\pgfqpoint{0.745560in}{1.026560in}}%
\pgfpathlineto{\pgfqpoint{0.761747in}{1.011550in}}%
\pgfpathlineto{\pgfqpoint{0.770478in}{1.012303in}}%
\pgfpathlineto{\pgfqpoint{0.779170in}{1.013433in}}%
\pgfpathlineto{\pgfqpoint{0.763025in}{1.028607in}}%
\pgfpathlineto{\pgfqpoint{0.746891in}{1.044165in}}%
\pgfpathlineto{\pgfqpoint{0.738157in}{1.042882in}}%
\pgfpathlineto{\pgfqpoint{0.729383in}{1.041960in}}%
\pgfpathclose%
\pgfusepath{fill}%
\end{pgfscope}%
\begin{pgfscope}%
\pgfpathrectangle{\pgfqpoint{0.407592in}{0.222904in}}{\pgfqpoint{1.411163in}{1.411163in}}%
\pgfusepath{clip}%
\pgfsetbuttcap%
\pgfsetroundjoin%
\definecolor{currentfill}{rgb}{0.909804,0.698039,0.423529}%
\pgfsetfillcolor{currentfill}%
\pgfsetfillopacity{0.120000}%
\pgfsetlinewidth{0.000000pt}%
\definecolor{currentstroke}{rgb}{0.000000,0.000000,0.000000}%
\pgfsetstrokecolor{currentstroke}%
\pgfsetdash{}{0pt}%
\pgfpathmoveto{\pgfqpoint{1.252180in}{0.917253in}}%
\pgfpathlineto{\pgfqpoint{1.268789in}{0.920336in}}%
\pgfpathlineto{\pgfqpoint{1.285459in}{0.924167in}}%
\pgfpathlineto{\pgfqpoint{1.293323in}{0.932277in}}%
\pgfpathlineto{\pgfqpoint{1.301169in}{0.940834in}}%
\pgfpathlineto{\pgfqpoint{1.284520in}{0.937142in}}%
\pgfpathlineto{\pgfqpoint{1.267932in}{0.934188in}}%
\pgfpathlineto{\pgfqpoint{1.260065in}{0.925489in}}%
\pgfpathlineto{\pgfqpoint{1.252180in}{0.917253in}}%
\pgfpathclose%
\pgfusepath{fill}%
\end{pgfscope}%
\begin{pgfscope}%
\pgfpathrectangle{\pgfqpoint{0.407592in}{0.222904in}}{\pgfqpoint{1.411163in}{1.411163in}}%
\pgfusepath{clip}%
\pgfsetbuttcap%
\pgfsetroundjoin%
\definecolor{currentfill}{rgb}{0.909804,0.698039,0.423529}%
\pgfsetfillcolor{currentfill}%
\pgfsetfillopacity{0.120000}%
\pgfsetlinewidth{0.000000pt}%
\definecolor{currentstroke}{rgb}{0.000000,0.000000,0.000000}%
\pgfsetstrokecolor{currentstroke}%
\pgfsetdash{}{0pt}%
\pgfpathmoveto{\pgfqpoint{1.414915in}{1.019364in}}%
\pgfpathlineto{\pgfqpoint{1.431833in}{1.025437in}}%
\pgfpathlineto{\pgfqpoint{1.448820in}{1.031998in}}%
\pgfpathlineto{\pgfqpoint{1.456502in}{1.042465in}}%
\pgfpathlineto{\pgfqpoint{1.464170in}{1.053222in}}%
\pgfpathlineto{\pgfqpoint{1.447202in}{1.046920in}}%
\pgfpathlineto{\pgfqpoint{1.430303in}{1.041105in}}%
\pgfpathlineto{\pgfqpoint{1.422616in}{1.030086in}}%
\pgfpathlineto{\pgfqpoint{1.414915in}{1.019364in}}%
\pgfpathclose%
\pgfusepath{fill}%
\end{pgfscope}%
\begin{pgfscope}%
\pgfpathrectangle{\pgfqpoint{0.407592in}{0.222904in}}{\pgfqpoint{1.411163in}{1.411163in}}%
\pgfusepath{clip}%
\pgfsetbuttcap%
\pgfsetroundjoin%
\definecolor{currentfill}{rgb}{0.909804,0.698039,0.423529}%
\pgfsetfillcolor{currentfill}%
\pgfsetfillopacity{0.120000}%
\pgfsetlinewidth{0.000000pt}%
\definecolor{currentstroke}{rgb}{0.000000,0.000000,0.000000}%
\pgfsetstrokecolor{currentstroke}%
\pgfsetdash{}{0pt}%
\pgfpathmoveto{\pgfqpoint{1.528781in}{1.111770in}}%
\pgfpathlineto{\pgfqpoint{1.545921in}{1.118826in}}%
\pgfpathlineto{\pgfqpoint{1.563132in}{1.126268in}}%
\pgfpathlineto{\pgfqpoint{1.570698in}{1.137736in}}%
\pgfpathlineto{\pgfqpoint{1.578252in}{1.149419in}}%
\pgfpathlineto{\pgfqpoint{1.561060in}{1.142252in}}%
\pgfpathlineto{\pgfqpoint{1.543940in}{1.135473in}}%
\pgfpathlineto{\pgfqpoint{1.536367in}{1.123513in}}%
\pgfpathlineto{\pgfqpoint{1.528781in}{1.111770in}}%
\pgfpathclose%
\pgfusepath{fill}%
\end{pgfscope}%
\begin{pgfscope}%
\pgfpathrectangle{\pgfqpoint{0.407592in}{0.222904in}}{\pgfqpoint{1.411163in}{1.411163in}}%
\pgfusepath{clip}%
\pgfsetbuttcap%
\pgfsetroundjoin%
\definecolor{currentfill}{rgb}{0.909804,0.698039,0.423529}%
\pgfsetfillcolor{currentfill}%
\pgfsetfillopacity{0.120000}%
\pgfsetlinewidth{0.000000pt}%
\definecolor{currentstroke}{rgb}{0.000000,0.000000,0.000000}%
\pgfsetstrokecolor{currentstroke}%
\pgfsetdash{}{0pt}%
\pgfpathmoveto{\pgfqpoint{0.614622in}{1.139137in}}%
\pgfpathlineto{\pgfqpoint{0.630780in}{1.122133in}}%
\pgfpathlineto{\pgfqpoint{0.646948in}{1.105370in}}%
\pgfpathlineto{\pgfqpoint{0.655886in}{1.105954in}}%
\pgfpathlineto{\pgfqpoint{0.664783in}{1.106854in}}%
\pgfpathlineto{\pgfqpoint{0.648658in}{1.123770in}}%
\pgfpathlineto{\pgfqpoint{0.632542in}{1.140920in}}%
\pgfpathlineto{\pgfqpoint{0.623602in}{1.139877in}}%
\pgfpathlineto{\pgfqpoint{0.614622in}{1.139137in}}%
\pgfpathclose%
\pgfusepath{fill}%
\end{pgfscope}%
\begin{pgfscope}%
\pgfpathrectangle{\pgfqpoint{0.407592in}{0.222904in}}{\pgfqpoint{1.411163in}{1.411163in}}%
\pgfusepath{clip}%
\pgfsetbuttcap%
\pgfsetroundjoin%
\definecolor{currentfill}{rgb}{0.909804,0.698039,0.423529}%
\pgfsetfillcolor{currentfill}%
\pgfsetfillopacity{0.120000}%
\pgfsetlinewidth{0.000000pt}%
\definecolor{currentstroke}{rgb}{0.000000,0.000000,0.000000}%
\pgfsetstrokecolor{currentstroke}%
\pgfsetdash{}{0pt}%
\pgfpathmoveto{\pgfqpoint{1.056536in}{0.876811in}}%
\pgfpathlineto{\pgfqpoint{1.072863in}{0.872427in}}%
\pgfpathlineto{\pgfqpoint{1.089227in}{0.869071in}}%
\pgfpathlineto{\pgfqpoint{1.097372in}{0.873065in}}%
\pgfpathlineto{\pgfqpoint{1.105490in}{0.877632in}}%
\pgfpathlineto{\pgfqpoint{1.089157in}{0.880962in}}%
\pgfpathlineto{\pgfqpoint{1.072863in}{0.885323in}}%
\pgfpathlineto{\pgfqpoint{1.064714in}{0.880781in}}%
\pgfpathlineto{\pgfqpoint{1.056536in}{0.876811in}}%
\pgfpathclose%
\pgfusepath{fill}%
\end{pgfscope}%
\begin{pgfscope}%
\pgfpathrectangle{\pgfqpoint{0.407592in}{0.222904in}}{\pgfqpoint{1.411163in}{1.411163in}}%
\pgfusepath{clip}%
\pgfsetbuttcap%
\pgfsetroundjoin%
\definecolor{currentfill}{rgb}{0.909804,0.698039,0.423529}%
\pgfsetfillcolor{currentfill}%
\pgfsetfillopacity{0.120000}%
\pgfsetlinewidth{0.000000pt}%
\definecolor{currentstroke}{rgb}{0.000000,0.000000,0.000000}%
\pgfsetstrokecolor{currentstroke}%
\pgfsetdash{}{0pt}%
\pgfpathmoveto{\pgfqpoint{0.925742in}{0.912991in}}%
\pgfpathlineto{\pgfqpoint{0.941979in}{0.903026in}}%
\pgfpathlineto{\pgfqpoint{0.958237in}{0.893896in}}%
\pgfpathlineto{\pgfqpoint{0.966614in}{0.895723in}}%
\pgfpathlineto{\pgfqpoint{0.974957in}{0.898065in}}%
\pgfpathlineto{\pgfqpoint{0.958738in}{0.907261in}}%
\pgfpathlineto{\pgfqpoint{0.942541in}{0.917303in}}%
\pgfpathlineto{\pgfqpoint{0.934159in}{0.914898in}}%
\pgfpathlineto{\pgfqpoint{0.925742in}{0.912991in}}%
\pgfpathclose%
\pgfusepath{fill}%
\end{pgfscope}%
\begin{pgfscope}%
\pgfpathrectangle{\pgfqpoint{0.407592in}{0.222904in}}{\pgfqpoint{1.411163in}{1.411163in}}%
\pgfusepath{clip}%
\pgfsetbuttcap%
\pgfsetroundjoin%
\definecolor{currentfill}{rgb}{0.909804,0.698039,0.423529}%
\pgfsetfillcolor{currentfill}%
\pgfsetfillopacity{0.120000}%
\pgfsetlinewidth{0.000000pt}%
\definecolor{currentstroke}{rgb}{0.000000,0.000000,0.000000}%
\pgfsetstrokecolor{currentstroke}%
\pgfsetdash{}{0pt}%
\pgfpathmoveto{\pgfqpoint{0.843871in}{0.957324in}}%
\pgfpathlineto{\pgfqpoint{0.860079in}{0.944651in}}%
\pgfpathlineto{\pgfqpoint{0.876303in}{0.932611in}}%
\pgfpathlineto{\pgfqpoint{0.884833in}{0.933694in}}%
\pgfpathlineto{\pgfqpoint{0.893326in}{0.935230in}}%
\pgfpathlineto{\pgfqpoint{0.877143in}{0.947404in}}%
\pgfpathlineto{\pgfqpoint{0.860975in}{0.960214in}}%
\pgfpathlineto{\pgfqpoint{0.852442in}{0.958551in}}%
\pgfpathlineto{\pgfqpoint{0.843871in}{0.957324in}}%
\pgfpathclose%
\pgfusepath{fill}%
\end{pgfscope}%
\begin{pgfscope}%
\pgfpathrectangle{\pgfqpoint{0.407592in}{0.222904in}}{\pgfqpoint{1.411163in}{1.411163in}}%
\pgfusepath{clip}%
\pgfsetbuttcap%
\pgfsetroundjoin%
\definecolor{currentfill}{rgb}{0.909804,0.698039,0.423529}%
\pgfsetfillcolor{currentfill}%
\pgfsetfillopacity{0.120000}%
\pgfsetlinewidth{0.000000pt}%
\definecolor{currentstroke}{rgb}{0.000000,0.000000,0.000000}%
\pgfsetstrokecolor{currentstroke}%
\pgfsetdash{}{0pt}%
\pgfpathmoveto{\pgfqpoint{1.301169in}{0.940834in}}%
\pgfpathlineto{\pgfqpoint{1.317881in}{0.945221in}}%
\pgfpathlineto{\pgfqpoint{1.334657in}{0.950256in}}%
\pgfpathlineto{\pgfqpoint{1.342466in}{0.959069in}}%
\pgfpathlineto{\pgfqpoint{1.350260in}{0.968282in}}%
\pgfpathlineto{\pgfqpoint{1.333504in}{0.963432in}}%
\pgfpathlineto{\pgfqpoint{1.316812in}{0.959225in}}%
\pgfpathlineto{\pgfqpoint{1.308999in}{0.949822in}}%
\pgfpathlineto{\pgfqpoint{1.301169in}{0.940834in}}%
\pgfpathclose%
\pgfusepath{fill}%
\end{pgfscope}%
\begin{pgfscope}%
\pgfpathrectangle{\pgfqpoint{0.407592in}{0.222904in}}{\pgfqpoint{1.411163in}{1.411163in}}%
\pgfusepath{clip}%
\pgfsetbuttcap%
\pgfsetroundjoin%
\definecolor{currentfill}{rgb}{0.909804,0.698039,0.423529}%
\pgfsetfillcolor{currentfill}%
\pgfsetfillopacity{0.120000}%
\pgfsetlinewidth{0.000000pt}%
\definecolor{currentstroke}{rgb}{0.000000,0.000000,0.000000}%
\pgfsetstrokecolor{currentstroke}%
\pgfsetdash{}{0pt}%
\pgfpathmoveto{\pgfqpoint{1.007468in}{0.882387in}}%
\pgfpathlineto{\pgfqpoint{1.023765in}{0.875982in}}%
\pgfpathlineto{\pgfqpoint{1.040092in}{0.870570in}}%
\pgfpathlineto{\pgfqpoint{1.048329in}{0.873411in}}%
\pgfpathlineto{\pgfqpoint{1.056536in}{0.876811in}}%
\pgfpathlineto{\pgfqpoint{1.040244in}{0.882217in}}%
\pgfpathlineto{\pgfqpoint{1.023983in}{0.888627in}}%
\pgfpathlineto{\pgfqpoint{1.015741in}{0.885231in}}%
\pgfpathlineto{\pgfqpoint{1.007468in}{0.882387in}}%
\pgfpathclose%
\pgfusepath{fill}%
\end{pgfscope}%
\begin{pgfscope}%
\pgfpathrectangle{\pgfqpoint{0.407592in}{0.222904in}}{\pgfqpoint{1.411163in}{1.411163in}}%
\pgfusepath{clip}%
\pgfsetbuttcap%
\pgfsetroundjoin%
\definecolor{currentfill}{rgb}{0.909804,0.698039,0.423529}%
\pgfsetfillcolor{currentfill}%
\pgfsetfillopacity{0.120000}%
\pgfsetlinewidth{0.000000pt}%
\definecolor{currentstroke}{rgb}{0.000000,0.000000,0.000000}%
\pgfsetstrokecolor{currentstroke}%
\pgfsetdash{}{0pt}%
\pgfpathmoveto{\pgfqpoint{1.642868in}{1.212429in}}%
\pgfpathlineto{\pgfqpoint{1.660237in}{1.220105in}}%
\pgfpathlineto{\pgfqpoint{1.677679in}{1.228102in}}%
\pgfpathlineto{\pgfqpoint{1.685134in}{1.240178in}}%
\pgfpathlineto{\pgfqpoint{1.667702in}{1.232314in}}%
\pgfpathlineto{\pgfqpoint{1.650343in}{1.224772in}}%
\pgfusepath{fill}%
\end{pgfscope}%
\begin{pgfscope}%
\pgfpathrectangle{\pgfqpoint{0.407592in}{0.222904in}}{\pgfqpoint{1.411163in}{1.411163in}}%
\pgfusepath{clip}%
\pgfsetbuttcap%
\pgfsetroundjoin%
\definecolor{currentfill}{rgb}{0.909804,0.698039,0.423529}%
\pgfsetfillcolor{currentfill}%
\pgfsetfillopacity{0.120000}%
\pgfsetlinewidth{0.000000pt}%
\definecolor{currentstroke}{rgb}{0.000000,0.000000,0.000000}%
\pgfsetstrokecolor{currentstroke}%
\pgfsetdash{}{0pt}%
\pgfpathmoveto{\pgfqpoint{0.761747in}{1.011550in}}%
\pgfpathlineto{\pgfqpoint{0.777945in}{0.996964in}}%
\pgfpathlineto{\pgfqpoint{0.794156in}{0.982838in}}%
\pgfpathlineto{\pgfqpoint{0.802845in}{0.983413in}}%
\pgfpathlineto{\pgfqpoint{0.811495in}{0.984382in}}%
\pgfpathlineto{\pgfqpoint{0.795327in}{0.998679in}}%
\pgfpathlineto{\pgfqpoint{0.779170in}{1.013433in}}%
\pgfpathlineto{\pgfqpoint{0.770478in}{1.012303in}}%
\pgfpathlineto{\pgfqpoint{0.761747in}{1.011550in}}%
\pgfpathclose%
\pgfusepath{fill}%
\end{pgfscope}%
\begin{pgfscope}%
\pgfpathrectangle{\pgfqpoint{0.407592in}{0.222904in}}{\pgfqpoint{1.411163in}{1.411163in}}%
\pgfusepath{clip}%
\pgfsetbuttcap%
\pgfsetroundjoin%
\definecolor{currentfill}{rgb}{0.909804,0.698039,0.423529}%
\pgfsetfillcolor{currentfill}%
\pgfsetfillopacity{0.120000}%
\pgfsetlinewidth{0.000000pt}%
\definecolor{currentstroke}{rgb}{0.000000,0.000000,0.000000}%
\pgfsetstrokecolor{currentstroke}%
\pgfsetdash{}{0pt}%
\pgfpathmoveto{\pgfqpoint{0.646948in}{1.105370in}}%
\pgfpathlineto{\pgfqpoint{0.663125in}{1.088867in}}%
\pgfpathlineto{\pgfqpoint{0.679313in}{1.072649in}}%
\pgfpathlineto{\pgfqpoint{0.688208in}{1.073061in}}%
\pgfpathlineto{\pgfqpoint{0.697063in}{1.073803in}}%
\pgfpathlineto{\pgfqpoint{0.680918in}{1.090190in}}%
\pgfpathlineto{\pgfqpoint{0.664783in}{1.106854in}}%
\pgfpathlineto{\pgfqpoint{0.655886in}{1.105954in}}%
\pgfpathlineto{\pgfqpoint{0.646948in}{1.105370in}}%
\pgfpathclose%
\pgfusepath{fill}%
\end{pgfscope}%
\begin{pgfscope}%
\pgfpathrectangle{\pgfqpoint{0.407592in}{0.222904in}}{\pgfqpoint{1.411163in}{1.411163in}}%
\pgfusepath{clip}%
\pgfsetbuttcap%
\pgfsetroundjoin%
\definecolor{currentfill}{rgb}{0.909804,0.698039,0.423529}%
\pgfsetfillcolor{currentfill}%
\pgfsetfillopacity{0.120000}%
\pgfsetlinewidth{0.000000pt}%
\definecolor{currentstroke}{rgb}{0.000000,0.000000,0.000000}%
\pgfsetstrokecolor{currentstroke}%
\pgfsetdash{}{0pt}%
\pgfpathmoveto{\pgfqpoint{1.464170in}{1.053222in}}%
\pgfpathlineto{\pgfqpoint{1.481207in}{1.059983in}}%
\pgfpathlineto{\pgfqpoint{1.498315in}{1.067173in}}%
\pgfpathlineto{\pgfqpoint{1.505950in}{1.077943in}}%
\pgfpathlineto{\pgfqpoint{1.513573in}{1.088973in}}%
\pgfpathlineto{\pgfqpoint{1.496485in}{1.082050in}}%
\pgfpathlineto{\pgfqpoint{1.479467in}{1.075559in}}%
\pgfpathlineto{\pgfqpoint{1.471825in}{1.064258in}}%
\pgfpathlineto{\pgfqpoint{1.464170in}{1.053222in}}%
\pgfpathclose%
\pgfusepath{fill}%
\end{pgfscope}%
\begin{pgfscope}%
\pgfpathrectangle{\pgfqpoint{0.407592in}{0.222904in}}{\pgfqpoint{1.411163in}{1.411163in}}%
\pgfusepath{clip}%
\pgfsetbuttcap%
\pgfsetroundjoin%
\definecolor{currentfill}{rgb}{0.909804,0.698039,0.423529}%
\pgfsetfillcolor{currentfill}%
\pgfsetfillopacity{0.120000}%
\pgfsetlinewidth{0.000000pt}%
\definecolor{currentstroke}{rgb}{0.000000,0.000000,0.000000}%
\pgfsetstrokecolor{currentstroke}%
\pgfsetdash{}{0pt}%
\pgfpathmoveto{\pgfqpoint{1.578252in}{1.149419in}}%
\pgfpathlineto{\pgfqpoint{1.595515in}{1.156953in}}%
\pgfpathlineto{\pgfqpoint{1.612852in}{1.164837in}}%
\pgfpathlineto{\pgfqpoint{1.620373in}{1.176451in}}%
\pgfpathlineto{\pgfqpoint{1.627883in}{1.188260in}}%
\pgfpathlineto{\pgfqpoint{1.610567in}{1.180649in}}%
\pgfpathlineto{\pgfqpoint{1.593323in}{1.173389in}}%
\pgfpathlineto{\pgfqpoint{1.585794in}{1.161307in}}%
\pgfpathlineto{\pgfqpoint{1.578252in}{1.149419in}}%
\pgfpathclose%
\pgfusepath{fill}%
\end{pgfscope}%
\begin{pgfscope}%
\pgfpathrectangle{\pgfqpoint{0.407592in}{0.222904in}}{\pgfqpoint{1.411163in}{1.411163in}}%
\pgfusepath{clip}%
\pgfsetbuttcap%
\pgfsetroundjoin%
\definecolor{currentfill}{rgb}{0.909804,0.698039,0.423529}%
\pgfsetfillcolor{currentfill}%
\pgfsetfillopacity{0.120000}%
\pgfsetlinewidth{0.000000pt}%
\definecolor{currentstroke}{rgb}{0.000000,0.000000,0.000000}%
\pgfsetstrokecolor{currentstroke}%
\pgfsetdash{}{0pt}%
\pgfpathmoveto{\pgfqpoint{1.350260in}{0.968282in}}%
\pgfpathlineto{\pgfqpoint{1.367081in}{0.973734in}}%
\pgfpathlineto{\pgfqpoint{1.383970in}{0.979747in}}%
\pgfpathlineto{\pgfqpoint{1.391728in}{0.989132in}}%
\pgfpathlineto{\pgfqpoint{1.399472in}{0.998872in}}%
\pgfpathlineto{\pgfqpoint{1.382603in}{0.993079in}}%
\pgfpathlineto{\pgfqpoint{1.365800in}{0.987845in}}%
\pgfpathlineto{\pgfqpoint{1.358038in}{0.977879in}}%
\pgfpathlineto{\pgfqpoint{1.350260in}{0.968282in}}%
\pgfpathclose%
\pgfusepath{fill}%
\end{pgfscope}%
\begin{pgfscope}%
\pgfpathrectangle{\pgfqpoint{0.407592in}{0.222904in}}{\pgfqpoint{1.411163in}{1.411163in}}%
\pgfusepath{clip}%
\pgfsetbuttcap%
\pgfsetroundjoin%
\definecolor{currentfill}{rgb}{0.909804,0.698039,0.423529}%
\pgfsetfillcolor{currentfill}%
\pgfsetfillopacity{0.120000}%
\pgfsetlinewidth{0.000000pt}%
\definecolor{currentstroke}{rgb}{0.000000,0.000000,0.000000}%
\pgfsetstrokecolor{currentstroke}%
\pgfsetdash{}{0pt}%
\pgfpathmoveto{\pgfqpoint{0.531830in}{1.207998in}}%
\pgfpathlineto{\pgfqpoint{0.547993in}{1.190347in}}%
\pgfpathlineto{\pgfqpoint{0.564166in}{1.172877in}}%
\pgfpathlineto{\pgfqpoint{0.573270in}{1.173194in}}%
\pgfpathlineto{\pgfqpoint{0.582334in}{1.173790in}}%
\pgfpathlineto{\pgfqpoint{0.566204in}{1.191407in}}%
\pgfpathlineto{\pgfqpoint{0.550083in}{1.209196in}}%
\pgfpathlineto{\pgfqpoint{0.540977in}{1.208462in}}%
\pgfpathlineto{\pgfqpoint{0.531830in}{1.207998in}}%
\pgfpathclose%
\pgfusepath{fill}%
\end{pgfscope}%
\begin{pgfscope}%
\pgfpathrectangle{\pgfqpoint{0.407592in}{0.222904in}}{\pgfqpoint{1.411163in}{1.411163in}}%
\pgfusepath{clip}%
\pgfsetbuttcap%
\pgfsetroundjoin%
\definecolor{currentfill}{rgb}{0.909804,0.698039,0.423529}%
\pgfsetfillcolor{currentfill}%
\pgfsetfillopacity{0.120000}%
\pgfsetlinewidth{0.000000pt}%
\definecolor{currentstroke}{rgb}{0.000000,0.000000,0.000000}%
\pgfsetstrokecolor{currentstroke}%
\pgfsetdash{}{0pt}%
\pgfpathmoveto{\pgfqpoint{0.876303in}{0.932611in}}%
\pgfpathlineto{\pgfqpoint{0.892544in}{0.921252in}}%
\pgfpathlineto{\pgfqpoint{0.908802in}{0.910621in}}%
\pgfpathlineto{\pgfqpoint{0.917290in}{0.911570in}}%
\pgfpathlineto{\pgfqpoint{0.925742in}{0.912991in}}%
\pgfpathlineto{\pgfqpoint{0.909525in}{0.923743in}}%
\pgfpathlineto{\pgfqpoint{0.893326in}{0.935230in}}%
\pgfpathlineto{\pgfqpoint{0.884833in}{0.933694in}}%
\pgfpathlineto{\pgfqpoint{0.876303in}{0.932611in}}%
\pgfpathclose%
\pgfusepath{fill}%
\end{pgfscope}%
\begin{pgfscope}%
\pgfpathrectangle{\pgfqpoint{0.407592in}{0.222904in}}{\pgfqpoint{1.411163in}{1.411163in}}%
\pgfusepath{clip}%
\pgfsetbuttcap%
\pgfsetroundjoin%
\definecolor{currentfill}{rgb}{0.909804,0.698039,0.423529}%
\pgfsetfillcolor{currentfill}%
\pgfsetfillopacity{0.120000}%
\pgfsetlinewidth{0.000000pt}%
\definecolor{currentstroke}{rgb}{0.000000,0.000000,0.000000}%
\pgfsetstrokecolor{currentstroke}%
\pgfsetdash{}{0pt}%
\pgfpathmoveto{\pgfqpoint{0.958237in}{0.893896in}}%
\pgfpathlineto{\pgfqpoint{0.974518in}{0.885645in}}%
\pgfpathlineto{\pgfqpoint{0.990824in}{0.878311in}}%
\pgfpathlineto{\pgfqpoint{0.999163in}{0.880085in}}%
\pgfpathlineto{\pgfqpoint{1.007468in}{0.882387in}}%
\pgfpathlineto{\pgfqpoint{0.991200in}{0.889760in}}%
\pgfpathlineto{\pgfqpoint{0.974957in}{0.898065in}}%
\pgfpathlineto{\pgfqpoint{0.966614in}{0.895723in}}%
\pgfpathlineto{\pgfqpoint{0.958237in}{0.893896in}}%
\pgfpathclose%
\pgfusepath{fill}%
\end{pgfscope}%
\begin{pgfscope}%
\pgfpathrectangle{\pgfqpoint{0.407592in}{0.222904in}}{\pgfqpoint{1.411163in}{1.411163in}}%
\pgfusepath{clip}%
\pgfsetbuttcap%
\pgfsetroundjoin%
\definecolor{currentfill}{rgb}{0.909804,0.698039,0.423529}%
\pgfsetfillcolor{currentfill}%
\pgfsetfillopacity{0.120000}%
\pgfsetlinewidth{0.000000pt}%
\definecolor{currentstroke}{rgb}{0.000000,0.000000,0.000000}%
\pgfsetstrokecolor{currentstroke}%
\pgfsetdash{}{0pt}%
\pgfpathmoveto{\pgfqpoint{0.679313in}{1.072649in}}%
\pgfpathlineto{\pgfqpoint{0.695510in}{1.056739in}}%
\pgfpathlineto{\pgfqpoint{0.711718in}{1.041165in}}%
\pgfpathlineto{\pgfqpoint{0.720571in}{1.041390in}}%
\pgfpathlineto{\pgfqpoint{0.729383in}{1.041960in}}%
\pgfpathlineto{\pgfqpoint{0.713218in}{1.057717in}}%
\pgfpathlineto{\pgfqpoint{0.697063in}{1.073803in}}%
\pgfpathlineto{\pgfqpoint{0.688208in}{1.073061in}}%
\pgfpathlineto{\pgfqpoint{0.679313in}{1.072649in}}%
\pgfpathclose%
\pgfusepath{fill}%
\end{pgfscope}%
\begin{pgfscope}%
\pgfpathrectangle{\pgfqpoint{0.407592in}{0.222904in}}{\pgfqpoint{1.411163in}{1.411163in}}%
\pgfusepath{clip}%
\pgfsetbuttcap%
\pgfsetroundjoin%
\definecolor{currentfill}{rgb}{0.909804,0.698039,0.423529}%
\pgfsetfillcolor{currentfill}%
\pgfsetfillopacity{0.120000}%
\pgfsetlinewidth{0.000000pt}%
\definecolor{currentstroke}{rgb}{0.000000,0.000000,0.000000}%
\pgfsetstrokecolor{currentstroke}%
\pgfsetdash{}{0pt}%
\pgfpathmoveto{\pgfqpoint{1.138281in}{0.874035in}}%
\pgfpathlineto{\pgfqpoint{1.154747in}{0.873727in}}%
\pgfpathlineto{\pgfqpoint{1.171262in}{0.874370in}}%
\pgfpathlineto{\pgfqpoint{1.179296in}{0.879546in}}%
\pgfpathlineto{\pgfqpoint{1.187307in}{0.885272in}}%
\pgfpathlineto{\pgfqpoint{1.170819in}{0.884624in}}%
\pgfpathlineto{\pgfqpoint{1.154381in}{0.884923in}}%
\pgfpathlineto{\pgfqpoint{1.146344in}{0.879197in}}%
\pgfpathlineto{\pgfqpoint{1.138281in}{0.874035in}}%
\pgfpathclose%
\pgfusepath{fill}%
\end{pgfscope}%
\begin{pgfscope}%
\pgfpathrectangle{\pgfqpoint{0.407592in}{0.222904in}}{\pgfqpoint{1.411163in}{1.411163in}}%
\pgfusepath{clip}%
\pgfsetbuttcap%
\pgfsetroundjoin%
\definecolor{currentfill}{rgb}{0.909804,0.698039,0.423529}%
\pgfsetfillcolor{currentfill}%
\pgfsetfillopacity{0.120000}%
\pgfsetlinewidth{0.000000pt}%
\definecolor{currentstroke}{rgb}{0.000000,0.000000,0.000000}%
\pgfsetstrokecolor{currentstroke}%
\pgfsetdash{}{0pt}%
\pgfpathmoveto{\pgfqpoint{1.513573in}{1.088973in}}%
\pgfpathlineto{\pgfqpoint{1.530733in}{1.096303in}}%
\pgfpathlineto{\pgfqpoint{1.547964in}{1.104016in}}%
\pgfpathlineto{\pgfqpoint{1.555554in}{1.115024in}}%
\pgfpathlineto{\pgfqpoint{1.563132in}{1.126268in}}%
\pgfpathlineto{\pgfqpoint{1.545921in}{1.118826in}}%
\pgfpathlineto{\pgfqpoint{1.528781in}{1.111770in}}%
\pgfpathlineto{\pgfqpoint{1.521184in}{1.100252in}}%
\pgfpathlineto{\pgfqpoint{1.513573in}{1.088973in}}%
\pgfpathclose%
\pgfusepath{fill}%
\end{pgfscope}%
\begin{pgfscope}%
\pgfpathrectangle{\pgfqpoint{0.407592in}{0.222904in}}{\pgfqpoint{1.411163in}{1.411163in}}%
\pgfusepath{clip}%
\pgfsetbuttcap%
\pgfsetroundjoin%
\definecolor{currentfill}{rgb}{0.909804,0.698039,0.423529}%
\pgfsetfillcolor{currentfill}%
\pgfsetfillopacity{0.120000}%
\pgfsetlinewidth{0.000000pt}%
\definecolor{currentstroke}{rgb}{0.000000,0.000000,0.000000}%
\pgfsetstrokecolor{currentstroke}%
\pgfsetdash{}{0pt}%
\pgfpathmoveto{\pgfqpoint{1.187307in}{0.885272in}}%
\pgfpathlineto{\pgfqpoint{1.203849in}{0.886826in}}%
\pgfpathlineto{\pgfqpoint{1.220446in}{0.889240in}}%
\pgfpathlineto{\pgfqpoint{1.228410in}{0.895480in}}%
\pgfpathlineto{\pgfqpoint{1.236353in}{0.902235in}}%
\pgfpathlineto{\pgfqpoint{1.219780in}{0.899865in}}%
\pgfpathlineto{\pgfqpoint{1.203263in}{0.898348in}}%
\pgfpathlineto{\pgfqpoint{1.195296in}{0.891542in}}%
\pgfpathlineto{\pgfqpoint{1.187307in}{0.885272in}}%
\pgfpathclose%
\pgfusepath{fill}%
\end{pgfscope}%
\begin{pgfscope}%
\pgfpathrectangle{\pgfqpoint{0.407592in}{0.222904in}}{\pgfqpoint{1.411163in}{1.411163in}}%
\pgfusepath{clip}%
\pgfsetbuttcap%
\pgfsetroundjoin%
\definecolor{currentfill}{rgb}{0.909804,0.698039,0.423529}%
\pgfsetfillcolor{currentfill}%
\pgfsetfillopacity{0.120000}%
\pgfsetlinewidth{0.000000pt}%
\definecolor{currentstroke}{rgb}{0.000000,0.000000,0.000000}%
\pgfsetstrokecolor{currentstroke}%
\pgfsetdash{}{0pt}%
\pgfpathmoveto{\pgfqpoint{0.794156in}{0.982838in}}%
\pgfpathlineto{\pgfqpoint{0.810379in}{0.969212in}}%
\pgfpathlineto{\pgfqpoint{0.826616in}{0.956127in}}%
\pgfpathlineto{\pgfqpoint{0.835262in}{0.956520in}}%
\pgfpathlineto{\pgfqpoint{0.843871in}{0.957324in}}%
\pgfpathlineto{\pgfqpoint{0.827676in}{0.970582in}}%
\pgfpathlineto{\pgfqpoint{0.811495in}{0.984382in}}%
\pgfpathlineto{\pgfqpoint{0.802845in}{0.983413in}}%
\pgfpathlineto{\pgfqpoint{0.794156in}{0.982838in}}%
\pgfpathclose%
\pgfusepath{fill}%
\end{pgfscope}%
\begin{pgfscope}%
\pgfpathrectangle{\pgfqpoint{0.407592in}{0.222904in}}{\pgfqpoint{1.411163in}{1.411163in}}%
\pgfusepath{clip}%
\pgfsetbuttcap%
\pgfsetroundjoin%
\definecolor{currentfill}{rgb}{0.909804,0.698039,0.423529}%
\pgfsetfillcolor{currentfill}%
\pgfsetfillopacity{0.120000}%
\pgfsetlinewidth{0.000000pt}%
\definecolor{currentstroke}{rgb}{0.000000,0.000000,0.000000}%
\pgfsetstrokecolor{currentstroke}%
\pgfsetdash{}{0pt}%
\pgfpathmoveto{\pgfqpoint{1.399472in}{0.998872in}}%
\pgfpathlineto{\pgfqpoint{1.416409in}{1.005188in}}%
\pgfpathlineto{\pgfqpoint{1.433415in}{1.011992in}}%
\pgfpathlineto{\pgfqpoint{1.441124in}{1.021836in}}%
\pgfpathlineto{\pgfqpoint{1.448820in}{1.031998in}}%
\pgfpathlineto{\pgfqpoint{1.431833in}{1.025437in}}%
\pgfpathlineto{\pgfqpoint{1.414915in}{1.019364in}}%
\pgfpathlineto{\pgfqpoint{1.407201in}{1.008955in}}%
\pgfpathlineto{\pgfqpoint{1.399472in}{0.998872in}}%
\pgfpathclose%
\pgfusepath{fill}%
\end{pgfscope}%
\begin{pgfscope}%
\pgfpathrectangle{\pgfqpoint{0.407592in}{0.222904in}}{\pgfqpoint{1.411163in}{1.411163in}}%
\pgfusepath{clip}%
\pgfsetbuttcap%
\pgfsetroundjoin%
\definecolor{currentfill}{rgb}{0.909804,0.698039,0.423529}%
\pgfsetfillcolor{currentfill}%
\pgfsetfillopacity{0.120000}%
\pgfsetlinewidth{0.000000pt}%
\definecolor{currentstroke}{rgb}{0.000000,0.000000,0.000000}%
\pgfsetstrokecolor{currentstroke}%
\pgfsetdash{}{0pt}%
\pgfpathmoveto{\pgfqpoint{0.564166in}{1.172877in}}%
\pgfpathlineto{\pgfqpoint{0.580348in}{1.155603in}}%
\pgfpathlineto{\pgfqpoint{0.596539in}{1.138540in}}%
\pgfpathlineto{\pgfqpoint{0.605601in}{1.138693in}}%
\pgfpathlineto{\pgfqpoint{0.614622in}{1.139137in}}%
\pgfpathlineto{\pgfqpoint{0.598473in}{1.156362in}}%
\pgfpathlineto{\pgfqpoint{0.582334in}{1.173790in}}%
\pgfpathlineto{\pgfqpoint{0.573270in}{1.173194in}}%
\pgfpathlineto{\pgfqpoint{0.564166in}{1.172877in}}%
\pgfpathclose%
\pgfusepath{fill}%
\end{pgfscope}%
\begin{pgfscope}%
\pgfpathrectangle{\pgfqpoint{0.407592in}{0.222904in}}{\pgfqpoint{1.411163in}{1.411163in}}%
\pgfusepath{clip}%
\pgfsetbuttcap%
\pgfsetroundjoin%
\definecolor{currentfill}{rgb}{0.909804,0.698039,0.423529}%
\pgfsetfillcolor{currentfill}%
\pgfsetfillopacity{0.120000}%
\pgfsetlinewidth{0.000000pt}%
\definecolor{currentstroke}{rgb}{0.000000,0.000000,0.000000}%
\pgfsetstrokecolor{currentstroke}%
\pgfsetdash{}{0pt}%
\pgfpathmoveto{\pgfqpoint{1.089227in}{0.869071in}}%
\pgfpathlineto{\pgfqpoint{1.105632in}{0.866737in}}%
\pgfpathlineto{\pgfqpoint{1.122080in}{0.865411in}}%
\pgfpathlineto{\pgfqpoint{1.130194in}{0.869440in}}%
\pgfpathlineto{\pgfqpoint{1.138281in}{0.874035in}}%
\pgfpathlineto{\pgfqpoint{1.121863in}{0.875328in}}%
\pgfpathlineto{\pgfqpoint{1.105490in}{0.877632in}}%
\pgfpathlineto{\pgfqpoint{1.097372in}{0.873065in}}%
\pgfpathlineto{\pgfqpoint{1.089227in}{0.869071in}}%
\pgfpathclose%
\pgfusepath{fill}%
\end{pgfscope}%
\begin{pgfscope}%
\pgfpathrectangle{\pgfqpoint{0.407592in}{0.222904in}}{\pgfqpoint{1.411163in}{1.411163in}}%
\pgfusepath{clip}%
\pgfsetbuttcap%
\pgfsetroundjoin%
\definecolor{currentfill}{rgb}{0.909804,0.698039,0.423529}%
\pgfsetfillcolor{currentfill}%
\pgfsetfillopacity{0.120000}%
\pgfsetlinewidth{0.000000pt}%
\definecolor{currentstroke}{rgb}{0.000000,0.000000,0.000000}%
\pgfsetstrokecolor{currentstroke}%
\pgfsetdash{}{0pt}%
\pgfpathmoveto{\pgfqpoint{1.236353in}{0.902235in}}%
\pgfpathlineto{\pgfqpoint{1.252984in}{0.905410in}}%
\pgfpathlineto{\pgfqpoint{1.269677in}{0.909342in}}%
\pgfpathlineto{\pgfqpoint{1.277577in}{0.916518in}}%
\pgfpathlineto{\pgfqpoint{1.285459in}{0.924167in}}%
\pgfpathlineto{\pgfqpoint{1.268789in}{0.920336in}}%
\pgfpathlineto{\pgfqpoint{1.252180in}{0.917253in}}%
\pgfpathlineto{\pgfqpoint{1.244276in}{0.909497in}}%
\pgfpathlineto{\pgfqpoint{1.236353in}{0.902235in}}%
\pgfpathclose%
\pgfusepath{fill}%
\end{pgfscope}%
\begin{pgfscope}%
\pgfpathrectangle{\pgfqpoint{0.407592in}{0.222904in}}{\pgfqpoint{1.411163in}{1.411163in}}%
\pgfusepath{clip}%
\pgfsetbuttcap%
\pgfsetroundjoin%
\definecolor{currentfill}{rgb}{0.909804,0.698039,0.423529}%
\pgfsetfillcolor{currentfill}%
\pgfsetfillopacity{0.120000}%
\pgfsetlinewidth{0.000000pt}%
\definecolor{currentstroke}{rgb}{0.000000,0.000000,0.000000}%
\pgfsetstrokecolor{currentstroke}%
\pgfsetdash{}{0pt}%
\pgfpathmoveto{\pgfqpoint{1.627883in}{1.188260in}}%
\pgfpathlineto{\pgfqpoint{1.645272in}{1.196206in}}%
\pgfpathlineto{\pgfqpoint{1.662735in}{1.204469in}}%
\pgfpathlineto{\pgfqpoint{1.670213in}{1.216196in}}%
\pgfpathlineto{\pgfqpoint{1.677679in}{1.228102in}}%
\pgfpathlineto{\pgfqpoint{1.660237in}{1.220105in}}%
\pgfpathlineto{\pgfqpoint{1.642868in}{1.212429in}}%
\pgfpathlineto{\pgfqpoint{1.635381in}{1.200256in}}%
\pgfpathlineto{\pgfqpoint{1.627883in}{1.188260in}}%
\pgfpathclose%
\pgfusepath{fill}%
\end{pgfscope}%
\begin{pgfscope}%
\pgfpathrectangle{\pgfqpoint{0.407592in}{0.222904in}}{\pgfqpoint{1.411163in}{1.411163in}}%
\pgfusepath{clip}%
\pgfsetbuttcap%
\pgfsetroundjoin%
\definecolor{currentfill}{rgb}{0.909804,0.698039,0.423529}%
\pgfsetfillcolor{currentfill}%
\pgfsetfillopacity{0.120000}%
\pgfsetlinewidth{0.000000pt}%
\definecolor{currentstroke}{rgb}{0.000000,0.000000,0.000000}%
\pgfsetstrokecolor{currentstroke}%
\pgfsetdash{}{0pt}%
\pgfpathmoveto{\pgfqpoint{1.285459in}{0.924167in}}%
\pgfpathlineto{\pgfqpoint{1.302192in}{0.928700in}}%
\pgfpathlineto{\pgfqpoint{1.318989in}{0.933887in}}%
\pgfpathlineto{\pgfqpoint{1.326831in}{0.941857in}}%
\pgfpathlineto{\pgfqpoint{1.334657in}{0.950256in}}%
\pgfpathlineto{\pgfqpoint{1.317881in}{0.945221in}}%
\pgfpathlineto{\pgfqpoint{1.301169in}{0.940834in}}%
\pgfpathlineto{\pgfqpoint{1.293323in}{0.932277in}}%
\pgfpathlineto{\pgfqpoint{1.285459in}{0.924167in}}%
\pgfpathclose%
\pgfusepath{fill}%
\end{pgfscope}%
\begin{pgfscope}%
\pgfpathrectangle{\pgfqpoint{0.407592in}{0.222904in}}{\pgfqpoint{1.411163in}{1.411163in}}%
\pgfusepath{clip}%
\pgfsetbuttcap%
\pgfsetroundjoin%
\definecolor{currentfill}{rgb}{0.909804,0.698039,0.423529}%
\pgfsetfillcolor{currentfill}%
\pgfsetfillopacity{0.120000}%
\pgfsetlinewidth{0.000000pt}%
\definecolor{currentstroke}{rgb}{0.000000,0.000000,0.000000}%
\pgfsetstrokecolor{currentstroke}%
\pgfsetdash{}{0pt}%
\pgfpathmoveto{\pgfqpoint{1.040092in}{0.870570in}}%
\pgfpathlineto{\pgfqpoint{1.056454in}{0.866168in}}%
\pgfpathlineto{\pgfqpoint{1.072852in}{0.862781in}}%
\pgfpathlineto{\pgfqpoint{1.081054in}{0.865646in}}%
\pgfpathlineto{\pgfqpoint{1.089227in}{0.869071in}}%
\pgfpathlineto{\pgfqpoint{1.072863in}{0.872427in}}%
\pgfpathlineto{\pgfqpoint{1.056536in}{0.876811in}}%
\pgfpathlineto{\pgfqpoint{1.048329in}{0.873411in}}%
\pgfpathlineto{\pgfqpoint{1.040092in}{0.870570in}}%
\pgfpathclose%
\pgfusepath{fill}%
\end{pgfscope}%
\begin{pgfscope}%
\pgfpathrectangle{\pgfqpoint{0.407592in}{0.222904in}}{\pgfqpoint{1.411163in}{1.411163in}}%
\pgfusepath{clip}%
\pgfsetbuttcap%
\pgfsetroundjoin%
\definecolor{currentfill}{rgb}{0.909804,0.698039,0.423529}%
\pgfsetfillcolor{currentfill}%
\pgfsetfillopacity{0.120000}%
\pgfsetlinewidth{0.000000pt}%
\definecolor{currentstroke}{rgb}{0.000000,0.000000,0.000000}%
\pgfsetstrokecolor{currentstroke}%
\pgfsetdash{}{0pt}%
\pgfpathmoveto{\pgfqpoint{0.711718in}{1.041165in}}%
\pgfpathlineto{\pgfqpoint{0.727937in}{1.025954in}}%
\pgfpathlineto{\pgfqpoint{0.744167in}{1.011138in}}%
\pgfpathlineto{\pgfqpoint{0.752976in}{1.011165in}}%
\pgfpathlineto{\pgfqpoint{0.761747in}{1.011550in}}%
\pgfpathlineto{\pgfqpoint{0.745560in}{1.026560in}}%
\pgfpathlineto{\pgfqpoint{0.729383in}{1.041960in}}%
\pgfpathlineto{\pgfqpoint{0.720571in}{1.041390in}}%
\pgfpathlineto{\pgfqpoint{0.711718in}{1.041165in}}%
\pgfpathclose%
\pgfusepath{fill}%
\end{pgfscope}%
\begin{pgfscope}%
\pgfpathrectangle{\pgfqpoint{0.407592in}{0.222904in}}{\pgfqpoint{1.411163in}{1.411163in}}%
\pgfusepath{clip}%
\pgfsetbuttcap%
\pgfsetroundjoin%
\definecolor{currentfill}{rgb}{0.909804,0.698039,0.423529}%
\pgfsetfillcolor{currentfill}%
\pgfsetfillopacity{0.120000}%
\pgfsetlinewidth{0.000000pt}%
\definecolor{currentstroke}{rgb}{0.000000,0.000000,0.000000}%
\pgfsetstrokecolor{currentstroke}%
\pgfsetdash{}{0pt}%
\pgfpathmoveto{\pgfqpoint{1.448820in}{1.031998in}}%
\pgfpathlineto{\pgfqpoint{1.465877in}{1.039017in}}%
\pgfpathlineto{\pgfqpoint{1.483004in}{1.046464in}}%
\pgfpathlineto{\pgfqpoint{1.490666in}{1.056677in}}%
\pgfpathlineto{\pgfqpoint{1.498315in}{1.067173in}}%
\pgfpathlineto{\pgfqpoint{1.481207in}{1.059983in}}%
\pgfpathlineto{\pgfqpoint{1.464170in}{1.053222in}}%
\pgfpathlineto{\pgfqpoint{1.456502in}{1.042465in}}%
\pgfpathlineto{\pgfqpoint{1.448820in}{1.031998in}}%
\pgfpathclose%
\pgfusepath{fill}%
\end{pgfscope}%
\begin{pgfscope}%
\pgfpathrectangle{\pgfqpoint{0.407592in}{0.222904in}}{\pgfqpoint{1.411163in}{1.411163in}}%
\pgfusepath{clip}%
\pgfsetbuttcap%
\pgfsetroundjoin%
\definecolor{currentfill}{rgb}{0.909804,0.698039,0.423529}%
\pgfsetfillcolor{currentfill}%
\pgfsetfillopacity{0.120000}%
\pgfsetlinewidth{0.000000pt}%
\definecolor{currentstroke}{rgb}{0.000000,0.000000,0.000000}%
\pgfsetstrokecolor{currentstroke}%
\pgfsetdash{}{0pt}%
\pgfpathmoveto{\pgfqpoint{0.596539in}{1.138540in}}%
\pgfpathlineto{\pgfqpoint{0.612740in}{1.121707in}}%
\pgfpathlineto{\pgfqpoint{0.628950in}{1.105123in}}%
\pgfpathlineto{\pgfqpoint{0.637969in}{1.105095in}}%
\pgfpathlineto{\pgfqpoint{0.646948in}{1.105370in}}%
\pgfpathlineto{\pgfqpoint{0.630780in}{1.122133in}}%
\pgfpathlineto{\pgfqpoint{0.614622in}{1.139137in}}%
\pgfpathlineto{\pgfqpoint{0.605601in}{1.138693in}}%
\pgfpathlineto{\pgfqpoint{0.596539in}{1.138540in}}%
\pgfpathclose%
\pgfusepath{fill}%
\end{pgfscope}%
\begin{pgfscope}%
\pgfpathrectangle{\pgfqpoint{0.407592in}{0.222904in}}{\pgfqpoint{1.411163in}{1.411163in}}%
\pgfusepath{clip}%
\pgfsetbuttcap%
\pgfsetroundjoin%
\definecolor{currentfill}{rgb}{0.909804,0.698039,0.423529}%
\pgfsetfillcolor{currentfill}%
\pgfsetfillopacity{0.120000}%
\pgfsetlinewidth{0.000000pt}%
\definecolor{currentstroke}{rgb}{0.000000,0.000000,0.000000}%
\pgfsetstrokecolor{currentstroke}%
\pgfsetdash{}{0pt}%
\pgfpathmoveto{\pgfqpoint{1.563132in}{1.126268in}}%
\pgfpathlineto{\pgfqpoint{1.580416in}{1.134075in}}%
\pgfpathlineto{\pgfqpoint{1.597773in}{1.142228in}}%
\pgfpathlineto{\pgfqpoint{1.605318in}{1.153426in}}%
\pgfpathlineto{\pgfqpoint{1.612852in}{1.164837in}}%
\pgfpathlineto{\pgfqpoint{1.595515in}{1.156953in}}%
\pgfpathlineto{\pgfqpoint{1.578252in}{1.149419in}}%
\pgfpathlineto{\pgfqpoint{1.570698in}{1.137736in}}%
\pgfpathlineto{\pgfqpoint{1.563132in}{1.126268in}}%
\pgfpathclose%
\pgfusepath{fill}%
\end{pgfscope}%
\begin{pgfscope}%
\pgfpathrectangle{\pgfqpoint{0.407592in}{0.222904in}}{\pgfqpoint{1.411163in}{1.411163in}}%
\pgfusepath{clip}%
\pgfsetbuttcap%
\pgfsetroundjoin%
\definecolor{currentfill}{rgb}{0.909804,0.698039,0.423529}%
\pgfsetfillcolor{currentfill}%
\pgfsetfillopacity{0.120000}%
\pgfsetlinewidth{0.000000pt}%
\definecolor{currentstroke}{rgb}{0.000000,0.000000,0.000000}%
\pgfsetstrokecolor{currentstroke}%
\pgfsetdash{}{0pt}%
\pgfpathmoveto{\pgfqpoint{0.908802in}{0.910621in}}%
\pgfpathlineto{\pgfqpoint{0.925080in}{0.900766in}}%
\pgfpathlineto{\pgfqpoint{0.941378in}{0.891731in}}%
\pgfpathlineto{\pgfqpoint{0.949825in}{0.892571in}}%
\pgfpathlineto{\pgfqpoint{0.958237in}{0.893896in}}%
\pgfpathlineto{\pgfqpoint{0.941979in}{0.903026in}}%
\pgfpathlineto{\pgfqpoint{0.925742in}{0.912991in}}%
\pgfpathlineto{\pgfqpoint{0.917290in}{0.911570in}}%
\pgfpathlineto{\pgfqpoint{0.908802in}{0.910621in}}%
\pgfpathclose%
\pgfusepath{fill}%
\end{pgfscope}%
\begin{pgfscope}%
\pgfpathrectangle{\pgfqpoint{0.407592in}{0.222904in}}{\pgfqpoint{1.411163in}{1.411163in}}%
\pgfusepath{clip}%
\pgfsetbuttcap%
\pgfsetroundjoin%
\definecolor{currentfill}{rgb}{0.909804,0.698039,0.423529}%
\pgfsetfillcolor{currentfill}%
\pgfsetfillopacity{0.120000}%
\pgfsetlinewidth{0.000000pt}%
\definecolor{currentstroke}{rgb}{0.000000,0.000000,0.000000}%
\pgfsetstrokecolor{currentstroke}%
\pgfsetdash{}{0pt}%
\pgfpathmoveto{\pgfqpoint{0.826616in}{0.956127in}}%
\pgfpathlineto{\pgfqpoint{0.842867in}{0.943626in}}%
\pgfpathlineto{\pgfqpoint{0.859133in}{0.931753in}}%
\pgfpathlineto{\pgfqpoint{0.867737in}{0.931969in}}%
\pgfpathlineto{\pgfqpoint{0.876303in}{0.932611in}}%
\pgfpathlineto{\pgfqpoint{0.860079in}{0.944651in}}%
\pgfpathlineto{\pgfqpoint{0.843871in}{0.957324in}}%
\pgfpathlineto{\pgfqpoint{0.835262in}{0.956520in}}%
\pgfpathlineto{\pgfqpoint{0.826616in}{0.956127in}}%
\pgfpathclose%
\pgfusepath{fill}%
\end{pgfscope}%
\begin{pgfscope}%
\pgfpathrectangle{\pgfqpoint{0.407592in}{0.222904in}}{\pgfqpoint{1.411163in}{1.411163in}}%
\pgfusepath{clip}%
\pgfsetbuttcap%
\pgfsetroundjoin%
\definecolor{currentfill}{rgb}{0.909804,0.698039,0.423529}%
\pgfsetfillcolor{currentfill}%
\pgfsetfillopacity{0.120000}%
\pgfsetlinewidth{0.000000pt}%
\definecolor{currentstroke}{rgb}{0.000000,0.000000,0.000000}%
\pgfsetstrokecolor{currentstroke}%
\pgfsetdash{}{0pt}%
\pgfpathmoveto{\pgfqpoint{1.334657in}{0.950256in}}%
\pgfpathlineto{\pgfqpoint{1.351499in}{0.955898in}}%
\pgfpathlineto{\pgfqpoint{1.368408in}{0.962103in}}%
\pgfpathlineto{\pgfqpoint{1.376197in}{0.970733in}}%
\pgfpathlineto{\pgfqpoint{1.383970in}{0.979747in}}%
\pgfpathlineto{\pgfqpoint{1.367081in}{0.973734in}}%
\pgfpathlineto{\pgfqpoint{1.350260in}{0.968282in}}%
\pgfpathlineto{\pgfqpoint{1.342466in}{0.959069in}}%
\pgfpathlineto{\pgfqpoint{1.334657in}{0.950256in}}%
\pgfpathclose%
\pgfusepath{fill}%
\end{pgfscope}%
\begin{pgfscope}%
\pgfpathrectangle{\pgfqpoint{0.407592in}{0.222904in}}{\pgfqpoint{1.411163in}{1.411163in}}%
\pgfusepath{clip}%
\pgfsetbuttcap%
\pgfsetroundjoin%
\definecolor{currentfill}{rgb}{0.909804,0.698039,0.423529}%
\pgfsetfillcolor{currentfill}%
\pgfsetfillopacity{0.120000}%
\pgfsetlinewidth{0.000000pt}%
\definecolor{currentstroke}{rgb}{0.000000,0.000000,0.000000}%
\pgfsetstrokecolor{currentstroke}%
\pgfsetdash{}{0pt}%
\pgfpathmoveto{\pgfqpoint{0.990824in}{0.878311in}}%
\pgfpathlineto{\pgfqpoint{1.007159in}{0.871928in}}%
\pgfpathlineto{\pgfqpoint{1.023524in}{0.866521in}}%
\pgfpathlineto{\pgfqpoint{1.031824in}{0.868278in}}%
\pgfpathlineto{\pgfqpoint{1.040092in}{0.870570in}}%
\pgfpathlineto{\pgfqpoint{1.023765in}{0.875982in}}%
\pgfpathlineto{\pgfqpoint{1.007468in}{0.882387in}}%
\pgfpathlineto{\pgfqpoint{0.999163in}{0.880085in}}%
\pgfpathlineto{\pgfqpoint{0.990824in}{0.878311in}}%
\pgfpathclose%
\pgfusepath{fill}%
\end{pgfscope}%
\begin{pgfscope}%
\pgfpathrectangle{\pgfqpoint{0.407592in}{0.222904in}}{\pgfqpoint{1.411163in}{1.411163in}}%
\pgfusepath{clip}%
\pgfsetbuttcap%
\pgfsetroundjoin%
\definecolor{currentfill}{rgb}{0.909804,0.698039,0.423529}%
\pgfsetfillcolor{currentfill}%
\pgfsetfillopacity{0.120000}%
\pgfsetlinewidth{0.000000pt}%
\definecolor{currentstroke}{rgb}{0.000000,0.000000,0.000000}%
\pgfsetstrokecolor{currentstroke}%
\pgfsetdash{}{0pt}%
\pgfpathmoveto{\pgfqpoint{0.744167in}{1.011138in}}%
\pgfpathlineto{\pgfqpoint{0.760408in}{0.996751in}}%
\pgfpathlineto{\pgfqpoint{0.776661in}{0.982827in}}%
\pgfpathlineto{\pgfqpoint{0.785428in}{0.982647in}}%
\pgfpathlineto{\pgfqpoint{0.794156in}{0.982838in}}%
\pgfpathlineto{\pgfqpoint{0.777945in}{0.996964in}}%
\pgfpathlineto{\pgfqpoint{0.761747in}{1.011550in}}%
\pgfpathlineto{\pgfqpoint{0.752976in}{1.011165in}}%
\pgfpathlineto{\pgfqpoint{0.744167in}{1.011138in}}%
\pgfpathclose%
\pgfusepath{fill}%
\end{pgfscope}%
\begin{pgfscope}%
\pgfpathrectangle{\pgfqpoint{0.407592in}{0.222904in}}{\pgfqpoint{1.411163in}{1.411163in}}%
\pgfusepath{clip}%
\pgfsetbuttcap%
\pgfsetroundjoin%
\definecolor{currentfill}{rgb}{0.909804,0.698039,0.423529}%
\pgfsetfillcolor{currentfill}%
\pgfsetfillopacity{0.120000}%
\pgfsetlinewidth{0.000000pt}%
\definecolor{currentstroke}{rgb}{0.000000,0.000000,0.000000}%
\pgfsetstrokecolor{currentstroke}%
\pgfsetdash{}{0pt}%
\pgfpathmoveto{\pgfqpoint{0.628950in}{1.105123in}}%
\pgfpathlineto{\pgfqpoint{0.645170in}{1.088807in}}%
\pgfpathlineto{\pgfqpoint{0.661400in}{1.072783in}}%
\pgfpathlineto{\pgfqpoint{0.670377in}{1.072559in}}%
\pgfpathlineto{\pgfqpoint{0.679313in}{1.072649in}}%
\pgfpathlineto{\pgfqpoint{0.663125in}{1.088867in}}%
\pgfpathlineto{\pgfqpoint{0.646948in}{1.105370in}}%
\pgfpathlineto{\pgfqpoint{0.637969in}{1.105095in}}%
\pgfpathlineto{\pgfqpoint{0.628950in}{1.105123in}}%
\pgfpathclose%
\pgfusepath{fill}%
\end{pgfscope}%
\begin{pgfscope}%
\pgfpathrectangle{\pgfqpoint{0.407592in}{0.222904in}}{\pgfqpoint{1.411163in}{1.411163in}}%
\pgfusepath{clip}%
\pgfsetbuttcap%
\pgfsetroundjoin%
\definecolor{currentfill}{rgb}{0.909804,0.698039,0.423529}%
\pgfsetfillcolor{currentfill}%
\pgfsetfillopacity{0.120000}%
\pgfsetlinewidth{0.000000pt}%
\definecolor{currentstroke}{rgb}{0.000000,0.000000,0.000000}%
\pgfsetstrokecolor{currentstroke}%
\pgfsetdash{}{0pt}%
\pgfpathmoveto{\pgfqpoint{1.677679in}{1.228102in}}%
\pgfpathlineto{\pgfqpoint{1.695195in}{1.236405in}}%
\pgfpathlineto{\pgfqpoint{1.712787in}{1.245000in}}%
\pgfpathlineto{\pgfqpoint{1.720220in}{1.256815in}}%
\pgfpathlineto{\pgfqpoint{1.702639in}{1.248349in}}%
\pgfpathlineto{\pgfqpoint{1.685134in}{1.240178in}}%
\pgfusepath{fill}%
\end{pgfscope}%
\begin{pgfscope}%
\pgfpathrectangle{\pgfqpoint{0.407592in}{0.222904in}}{\pgfqpoint{1.411163in}{1.411163in}}%
\pgfusepath{clip}%
\pgfsetbuttcap%
\pgfsetroundjoin%
\definecolor{currentfill}{rgb}{0.909804,0.698039,0.423529}%
\pgfsetfillcolor{currentfill}%
\pgfsetfillopacity{0.120000}%
\pgfsetlinewidth{0.000000pt}%
\definecolor{currentstroke}{rgb}{0.000000,0.000000,0.000000}%
\pgfsetstrokecolor{currentstroke}%
\pgfsetdash{}{0pt}%
\pgfpathmoveto{\pgfqpoint{1.498315in}{1.067173in}}%
\pgfpathlineto{\pgfqpoint{1.515494in}{1.074769in}}%
\pgfpathlineto{\pgfqpoint{1.532746in}{1.082745in}}%
\pgfpathlineto{\pgfqpoint{1.540361in}{1.093253in}}%
\pgfpathlineto{\pgfqpoint{1.547964in}{1.104016in}}%
\pgfpathlineto{\pgfqpoint{1.530733in}{1.096303in}}%
\pgfpathlineto{\pgfqpoint{1.513573in}{1.088973in}}%
\pgfpathlineto{\pgfqpoint{1.505950in}{1.077943in}}%
\pgfpathlineto{\pgfqpoint{1.498315in}{1.067173in}}%
\pgfpathclose%
\pgfusepath{fill}%
\end{pgfscope}%
\begin{pgfscope}%
\pgfpathrectangle{\pgfqpoint{0.407592in}{0.222904in}}{\pgfqpoint{1.411163in}{1.411163in}}%
\pgfusepath{clip}%
\pgfsetbuttcap%
\pgfsetroundjoin%
\definecolor{currentfill}{rgb}{0.909804,0.698039,0.423529}%
\pgfsetfillcolor{currentfill}%
\pgfsetfillopacity{0.120000}%
\pgfsetlinewidth{0.000000pt}%
\definecolor{currentstroke}{rgb}{0.000000,0.000000,0.000000}%
\pgfsetstrokecolor{currentstroke}%
\pgfsetdash{}{0pt}%
\pgfpathmoveto{\pgfqpoint{1.383970in}{0.979747in}}%
\pgfpathlineto{\pgfqpoint{1.400927in}{0.986284in}}%
\pgfpathlineto{\pgfqpoint{1.417954in}{0.993309in}}%
\pgfpathlineto{\pgfqpoint{1.425692in}{1.002479in}}%
\pgfpathlineto{\pgfqpoint{1.433415in}{1.011992in}}%
\pgfpathlineto{\pgfqpoint{1.416409in}{1.005188in}}%
\pgfpathlineto{\pgfqpoint{1.399472in}{0.998872in}}%
\pgfpathlineto{\pgfqpoint{1.391728in}{0.989132in}}%
\pgfpathlineto{\pgfqpoint{1.383970in}{0.979747in}}%
\pgfpathclose%
\pgfusepath{fill}%
\end{pgfscope}%
\begin{pgfscope}%
\pgfpathrectangle{\pgfqpoint{0.407592in}{0.222904in}}{\pgfqpoint{1.411163in}{1.411163in}}%
\pgfusepath{clip}%
\pgfsetbuttcap%
\pgfsetroundjoin%
\definecolor{currentfill}{rgb}{0.909804,0.698039,0.423529}%
\pgfsetfillcolor{currentfill}%
\pgfsetfillopacity{0.120000}%
\pgfsetlinewidth{0.000000pt}%
\definecolor{currentstroke}{rgb}{0.000000,0.000000,0.000000}%
\pgfsetstrokecolor{currentstroke}%
\pgfsetdash{}{0pt}%
\pgfpathmoveto{\pgfqpoint{1.612852in}{1.164837in}}%
\pgfpathlineto{\pgfqpoint{1.630262in}{1.173051in}}%
\pgfpathlineto{\pgfqpoint{1.647746in}{1.181579in}}%
\pgfpathlineto{\pgfqpoint{1.655247in}{1.192927in}}%
\pgfpathlineto{\pgfqpoint{1.662735in}{1.204469in}}%
\pgfpathlineto{\pgfqpoint{1.645272in}{1.196206in}}%
\pgfpathlineto{\pgfqpoint{1.627883in}{1.188260in}}%
\pgfpathlineto{\pgfqpoint{1.620373in}{1.176451in}}%
\pgfpathlineto{\pgfqpoint{1.612852in}{1.164837in}}%
\pgfpathclose%
\pgfusepath{fill}%
\end{pgfscope}%
\begin{pgfscope}%
\pgfpathrectangle{\pgfqpoint{0.407592in}{0.222904in}}{\pgfqpoint{1.411163in}{1.411163in}}%
\pgfusepath{clip}%
\pgfsetbuttcap%
\pgfsetroundjoin%
\definecolor{currentfill}{rgb}{0.909804,0.698039,0.423529}%
\pgfsetfillcolor{currentfill}%
\pgfsetfillopacity{0.120000}%
\pgfsetlinewidth{0.000000pt}%
\definecolor{currentstroke}{rgb}{0.000000,0.000000,0.000000}%
\pgfsetstrokecolor{currentstroke}%
\pgfsetdash{}{0pt}%
\pgfpathmoveto{\pgfqpoint{0.513409in}{1.207856in}}%
\pgfpathlineto{\pgfqpoint{0.529615in}{1.190365in}}%
\pgfpathlineto{\pgfqpoint{0.545830in}{1.173062in}}%
\pgfpathlineto{\pgfqpoint{0.555019in}{1.172835in}}%
\pgfpathlineto{\pgfqpoint{0.564166in}{1.172877in}}%
\pgfpathlineto{\pgfqpoint{0.547993in}{1.190347in}}%
\pgfpathlineto{\pgfqpoint{0.531830in}{1.207998in}}%
\pgfpathlineto{\pgfqpoint{0.522640in}{1.207797in}}%
\pgfpathlineto{\pgfqpoint{0.513409in}{1.207856in}}%
\pgfpathclose%
\pgfusepath{fill}%
\end{pgfscope}%
\begin{pgfscope}%
\pgfpathrectangle{\pgfqpoint{0.407592in}{0.222904in}}{\pgfqpoint{1.411163in}{1.411163in}}%
\pgfusepath{clip}%
\pgfsetbuttcap%
\pgfsetroundjoin%
\definecolor{currentfill}{rgb}{0.909804,0.698039,0.423529}%
\pgfsetfillcolor{currentfill}%
\pgfsetfillopacity{0.120000}%
\pgfsetlinewidth{0.000000pt}%
\definecolor{currentstroke}{rgb}{0.000000,0.000000,0.000000}%
\pgfsetstrokecolor{currentstroke}%
\pgfsetdash{}{0pt}%
\pgfpathmoveto{\pgfqpoint{1.171262in}{0.874370in}}%
\pgfpathlineto{\pgfqpoint{1.187830in}{0.875922in}}%
\pgfpathlineto{\pgfqpoint{1.204453in}{0.878338in}}%
\pgfpathlineto{\pgfqpoint{1.212460in}{0.883525in}}%
\pgfpathlineto{\pgfqpoint{1.220446in}{0.889240in}}%
\pgfpathlineto{\pgfqpoint{1.203849in}{0.886826in}}%
\pgfpathlineto{\pgfqpoint{1.187307in}{0.885272in}}%
\pgfpathlineto{\pgfqpoint{1.179296in}{0.879546in}}%
\pgfpathlineto{\pgfqpoint{1.171262in}{0.874370in}}%
\pgfpathclose%
\pgfusepath{fill}%
\end{pgfscope}%
\begin{pgfscope}%
\pgfpathrectangle{\pgfqpoint{0.407592in}{0.222904in}}{\pgfqpoint{1.411163in}{1.411163in}}%
\pgfusepath{clip}%
\pgfsetbuttcap%
\pgfsetroundjoin%
\definecolor{currentfill}{rgb}{0.909804,0.698039,0.423529}%
\pgfsetfillcolor{currentfill}%
\pgfsetfillopacity{0.120000}%
\pgfsetlinewidth{0.000000pt}%
\definecolor{currentstroke}{rgb}{0.000000,0.000000,0.000000}%
\pgfsetstrokecolor{currentstroke}%
\pgfsetdash{}{0pt}%
\pgfpathmoveto{\pgfqpoint{1.122080in}{0.865411in}}%
\pgfpathlineto{\pgfqpoint{1.138575in}{0.865066in}}%
\pgfpathlineto{\pgfqpoint{1.155119in}{0.865670in}}%
\pgfpathlineto{\pgfqpoint{1.163203in}{0.869745in}}%
\pgfpathlineto{\pgfqpoint{1.171262in}{0.874370in}}%
\pgfpathlineto{\pgfqpoint{1.154747in}{0.873727in}}%
\pgfpathlineto{\pgfqpoint{1.138281in}{0.874035in}}%
\pgfpathlineto{\pgfqpoint{1.130194in}{0.869440in}}%
\pgfpathlineto{\pgfqpoint{1.122080in}{0.865411in}}%
\pgfpathclose%
\pgfusepath{fill}%
\end{pgfscope}%
\begin{pgfscope}%
\pgfpathrectangle{\pgfqpoint{0.407592in}{0.222904in}}{\pgfqpoint{1.411163in}{1.411163in}}%
\pgfusepath{clip}%
\pgfsetbuttcap%
\pgfsetroundjoin%
\definecolor{currentfill}{rgb}{0.909804,0.698039,0.423529}%
\pgfsetfillcolor{currentfill}%
\pgfsetfillopacity{0.120000}%
\pgfsetlinewidth{0.000000pt}%
\definecolor{currentstroke}{rgb}{0.000000,0.000000,0.000000}%
\pgfsetstrokecolor{currentstroke}%
\pgfsetdash{}{0pt}%
\pgfpathmoveto{\pgfqpoint{0.859133in}{0.931753in}}%
\pgfpathlineto{\pgfqpoint{0.875416in}{0.920553in}}%
\pgfpathlineto{\pgfqpoint{0.891716in}{0.910071in}}%
\pgfpathlineto{\pgfqpoint{0.900278in}{0.910127in}}%
\pgfpathlineto{\pgfqpoint{0.908802in}{0.910621in}}%
\pgfpathlineto{\pgfqpoint{0.892544in}{0.921252in}}%
\pgfpathlineto{\pgfqpoint{0.876303in}{0.932611in}}%
\pgfpathlineto{\pgfqpoint{0.867737in}{0.931969in}}%
\pgfpathlineto{\pgfqpoint{0.859133in}{0.931753in}}%
\pgfpathclose%
\pgfusepath{fill}%
\end{pgfscope}%
\begin{pgfscope}%
\pgfpathrectangle{\pgfqpoint{0.407592in}{0.222904in}}{\pgfqpoint{1.411163in}{1.411163in}}%
\pgfusepath{clip}%
\pgfsetbuttcap%
\pgfsetroundjoin%
\definecolor{currentfill}{rgb}{0.909804,0.698039,0.423529}%
\pgfsetfillcolor{currentfill}%
\pgfsetfillopacity{0.120000}%
\pgfsetlinewidth{0.000000pt}%
\definecolor{currentstroke}{rgb}{0.000000,0.000000,0.000000}%
\pgfsetstrokecolor{currentstroke}%
\pgfsetdash{}{0pt}%
\pgfpathmoveto{\pgfqpoint{1.220446in}{0.889240in}}%
\pgfpathlineto{\pgfqpoint{1.237101in}{0.892466in}}%
\pgfpathlineto{\pgfqpoint{1.253817in}{0.896454in}}%
\pgfpathlineto{\pgfqpoint{1.261757in}{0.902651in}}%
\pgfpathlineto{\pgfqpoint{1.269677in}{0.909342in}}%
\pgfpathlineto{\pgfqpoint{1.252984in}{0.905410in}}%
\pgfpathlineto{\pgfqpoint{1.236353in}{0.902235in}}%
\pgfpathlineto{\pgfqpoint{1.228410in}{0.895480in}}%
\pgfpathlineto{\pgfqpoint{1.220446in}{0.889240in}}%
\pgfpathclose%
\pgfusepath{fill}%
\end{pgfscope}%
\begin{pgfscope}%
\pgfpathrectangle{\pgfqpoint{0.407592in}{0.222904in}}{\pgfqpoint{1.411163in}{1.411163in}}%
\pgfusepath{clip}%
\pgfsetbuttcap%
\pgfsetroundjoin%
\definecolor{currentfill}{rgb}{0.909804,0.698039,0.423529}%
\pgfsetfillcolor{currentfill}%
\pgfsetfillopacity{0.120000}%
\pgfsetlinewidth{0.000000pt}%
\definecolor{currentstroke}{rgb}{0.000000,0.000000,0.000000}%
\pgfsetstrokecolor{currentstroke}%
\pgfsetdash{}{0pt}%
\pgfpathmoveto{\pgfqpoint{0.941378in}{0.891731in}}%
\pgfpathlineto{\pgfqpoint{0.957700in}{0.883558in}}%
\pgfpathlineto{\pgfqpoint{0.974046in}{0.876284in}}%
\pgfpathlineto{\pgfqpoint{0.982453in}{0.877050in}}%
\pgfpathlineto{\pgfqpoint{0.990824in}{0.878311in}}%
\pgfpathlineto{\pgfqpoint{0.974518in}{0.885645in}}%
\pgfpathlineto{\pgfqpoint{0.958237in}{0.893896in}}%
\pgfpathlineto{\pgfqpoint{0.949825in}{0.892571in}}%
\pgfpathlineto{\pgfqpoint{0.941378in}{0.891731in}}%
\pgfpathclose%
\pgfusepath{fill}%
\end{pgfscope}%
\begin{pgfscope}%
\pgfpathrectangle{\pgfqpoint{0.407592in}{0.222904in}}{\pgfqpoint{1.411163in}{1.411163in}}%
\pgfusepath{clip}%
\pgfsetbuttcap%
\pgfsetroundjoin%
\definecolor{currentfill}{rgb}{0.909804,0.698039,0.423529}%
\pgfsetfillcolor{currentfill}%
\pgfsetfillopacity{0.120000}%
\pgfsetlinewidth{0.000000pt}%
\definecolor{currentstroke}{rgb}{0.000000,0.000000,0.000000}%
\pgfsetstrokecolor{currentstroke}%
\pgfsetdash{}{0pt}%
\pgfpathmoveto{\pgfqpoint{0.661400in}{1.072783in}}%
\pgfpathlineto{\pgfqpoint{0.677641in}{1.057074in}}%
\pgfpathlineto{\pgfqpoint{0.693891in}{1.041708in}}%
\pgfpathlineto{\pgfqpoint{0.702825in}{1.041274in}}%
\pgfpathlineto{\pgfqpoint{0.711718in}{1.041165in}}%
\pgfpathlineto{\pgfqpoint{0.695510in}{1.056739in}}%
\pgfpathlineto{\pgfqpoint{0.679313in}{1.072649in}}%
\pgfpathlineto{\pgfqpoint{0.670377in}{1.072559in}}%
\pgfpathlineto{\pgfqpoint{0.661400in}{1.072783in}}%
\pgfpathclose%
\pgfusepath{fill}%
\end{pgfscope}%
\begin{pgfscope}%
\pgfpathrectangle{\pgfqpoint{0.407592in}{0.222904in}}{\pgfqpoint{1.411163in}{1.411163in}}%
\pgfusepath{clip}%
\pgfsetbuttcap%
\pgfsetroundjoin%
\definecolor{currentfill}{rgb}{0.909804,0.698039,0.423529}%
\pgfsetfillcolor{currentfill}%
\pgfsetfillopacity{0.120000}%
\pgfsetlinewidth{0.000000pt}%
\definecolor{currentstroke}{rgb}{0.000000,0.000000,0.000000}%
\pgfsetstrokecolor{currentstroke}%
\pgfsetdash{}{0pt}%
\pgfpathmoveto{\pgfqpoint{1.072852in}{0.862781in}}%
\pgfpathlineto{\pgfqpoint{1.089290in}{0.860406in}}%
\pgfpathlineto{\pgfqpoint{1.105771in}{0.859027in}}%
\pgfpathlineto{\pgfqpoint{1.113940in}{0.861942in}}%
\pgfpathlineto{\pgfqpoint{1.122080in}{0.865411in}}%
\pgfpathlineto{\pgfqpoint{1.105632in}{0.866737in}}%
\pgfpathlineto{\pgfqpoint{1.089227in}{0.869071in}}%
\pgfpathlineto{\pgfqpoint{1.081054in}{0.865646in}}%
\pgfpathlineto{\pgfqpoint{1.072852in}{0.862781in}}%
\pgfpathclose%
\pgfusepath{fill}%
\end{pgfscope}%
\begin{pgfscope}%
\pgfpathrectangle{\pgfqpoint{0.407592in}{0.222904in}}{\pgfqpoint{1.411163in}{1.411163in}}%
\pgfusepath{clip}%
\pgfsetbuttcap%
\pgfsetroundjoin%
\definecolor{currentfill}{rgb}{0.909804,0.698039,0.423529}%
\pgfsetfillcolor{currentfill}%
\pgfsetfillopacity{0.120000}%
\pgfsetlinewidth{0.000000pt}%
\definecolor{currentstroke}{rgb}{0.000000,0.000000,0.000000}%
\pgfsetstrokecolor{currentstroke}%
\pgfsetdash{}{0pt}%
\pgfpathmoveto{\pgfqpoint{0.776661in}{0.982827in}}%
\pgfpathlineto{\pgfqpoint{0.792928in}{0.969403in}}%
\pgfpathlineto{\pgfqpoint{0.809207in}{0.956520in}}%
\pgfpathlineto{\pgfqpoint{0.817931in}{0.956132in}}%
\pgfpathlineto{\pgfqpoint{0.826616in}{0.956127in}}%
\pgfpathlineto{\pgfqpoint{0.810379in}{0.969212in}}%
\pgfpathlineto{\pgfqpoint{0.794156in}{0.982838in}}%
\pgfpathlineto{\pgfqpoint{0.785428in}{0.982647in}}%
\pgfpathlineto{\pgfqpoint{0.776661in}{0.982827in}}%
\pgfpathclose%
\pgfusepath{fill}%
\end{pgfscope}%
\begin{pgfscope}%
\pgfpathrectangle{\pgfqpoint{0.407592in}{0.222904in}}{\pgfqpoint{1.411163in}{1.411163in}}%
\pgfusepath{clip}%
\pgfsetbuttcap%
\pgfsetroundjoin%
\definecolor{currentfill}{rgb}{0.909804,0.698039,0.423529}%
\pgfsetfillcolor{currentfill}%
\pgfsetfillopacity{0.120000}%
\pgfsetlinewidth{0.000000pt}%
\definecolor{currentstroke}{rgb}{0.000000,0.000000,0.000000}%
\pgfsetstrokecolor{currentstroke}%
\pgfsetdash{}{0pt}%
\pgfpathmoveto{\pgfqpoint{1.269677in}{0.909342in}}%
\pgfpathlineto{\pgfqpoint{1.286431in}{0.913981in}}%
\pgfpathlineto{\pgfqpoint{1.303250in}{0.919278in}}%
\pgfpathlineto{\pgfqpoint{1.311128in}{0.926356in}}%
\pgfpathlineto{\pgfqpoint{1.318989in}{0.933887in}}%
\pgfpathlineto{\pgfqpoint{1.302192in}{0.928700in}}%
\pgfpathlineto{\pgfqpoint{1.285459in}{0.924167in}}%
\pgfpathlineto{\pgfqpoint{1.277577in}{0.916518in}}%
\pgfpathlineto{\pgfqpoint{1.269677in}{0.909342in}}%
\pgfpathclose%
\pgfusepath{fill}%
\end{pgfscope}%
\begin{pgfscope}%
\pgfpathrectangle{\pgfqpoint{0.407592in}{0.222904in}}{\pgfqpoint{1.411163in}{1.411163in}}%
\pgfusepath{clip}%
\pgfsetbuttcap%
\pgfsetroundjoin%
\definecolor{currentfill}{rgb}{0.909804,0.698039,0.423529}%
\pgfsetfillcolor{currentfill}%
\pgfsetfillopacity{0.120000}%
\pgfsetlinewidth{0.000000pt}%
\definecolor{currentstroke}{rgb}{0.000000,0.000000,0.000000}%
\pgfsetstrokecolor{currentstroke}%
\pgfsetdash{}{0pt}%
\pgfpathmoveto{\pgfqpoint{1.433415in}{1.011992in}}%
\pgfpathlineto{\pgfqpoint{1.450492in}{1.019253in}}%
\pgfpathlineto{\pgfqpoint{1.467640in}{1.026940in}}%
\pgfpathlineto{\pgfqpoint{1.475329in}{1.036548in}}%
\pgfpathlineto{\pgfqpoint{1.483004in}{1.046464in}}%
\pgfpathlineto{\pgfqpoint{1.465877in}{1.039017in}}%
\pgfpathlineto{\pgfqpoint{1.448820in}{1.031998in}}%
\pgfpathlineto{\pgfqpoint{1.441124in}{1.021836in}}%
\pgfpathlineto{\pgfqpoint{1.433415in}{1.011992in}}%
\pgfpathclose%
\pgfusepath{fill}%
\end{pgfscope}%
\begin{pgfscope}%
\pgfpathrectangle{\pgfqpoint{0.407592in}{0.222904in}}{\pgfqpoint{1.411163in}{1.411163in}}%
\pgfusepath{clip}%
\pgfsetbuttcap%
\pgfsetroundjoin%
\definecolor{currentfill}{rgb}{0.909804,0.698039,0.423529}%
\pgfsetfillcolor{currentfill}%
\pgfsetfillopacity{0.120000}%
\pgfsetlinewidth{0.000000pt}%
\definecolor{currentstroke}{rgb}{0.000000,0.000000,0.000000}%
\pgfsetstrokecolor{currentstroke}%
\pgfsetdash{}{0pt}%
\pgfpathmoveto{\pgfqpoint{1.547964in}{1.104016in}}%
\pgfpathlineto{\pgfqpoint{1.565268in}{1.112091in}}%
\pgfpathlineto{\pgfqpoint{1.582646in}{1.120508in}}%
\pgfpathlineto{\pgfqpoint{1.590216in}{1.131252in}}%
\pgfpathlineto{\pgfqpoint{1.597773in}{1.142228in}}%
\pgfpathlineto{\pgfqpoint{1.580416in}{1.134075in}}%
\pgfpathlineto{\pgfqpoint{1.563132in}{1.126268in}}%
\pgfpathlineto{\pgfqpoint{1.555554in}{1.115024in}}%
\pgfpathlineto{\pgfqpoint{1.547964in}{1.104016in}}%
\pgfpathclose%
\pgfusepath{fill}%
\end{pgfscope}%
\begin{pgfscope}%
\pgfpathrectangle{\pgfqpoint{0.407592in}{0.222904in}}{\pgfqpoint{1.411163in}{1.411163in}}%
\pgfusepath{clip}%
\pgfsetbuttcap%
\pgfsetroundjoin%
\definecolor{currentfill}{rgb}{0.909804,0.698039,0.423529}%
\pgfsetfillcolor{currentfill}%
\pgfsetfillopacity{0.120000}%
\pgfsetlinewidth{0.000000pt}%
\definecolor{currentstroke}{rgb}{0.000000,0.000000,0.000000}%
\pgfsetstrokecolor{currentstroke}%
\pgfsetdash{}{0pt}%
\pgfpathmoveto{\pgfqpoint{0.545830in}{1.173062in}}%
\pgfpathlineto{\pgfqpoint{0.562055in}{1.155963in}}%
\pgfpathlineto{\pgfqpoint{0.578289in}{1.139084in}}%
\pgfpathlineto{\pgfqpoint{0.587435in}{1.138673in}}%
\pgfpathlineto{\pgfqpoint{0.596539in}{1.138540in}}%
\pgfpathlineto{\pgfqpoint{0.580348in}{1.155603in}}%
\pgfpathlineto{\pgfqpoint{0.564166in}{1.172877in}}%
\pgfpathlineto{\pgfqpoint{0.555019in}{1.172835in}}%
\pgfpathlineto{\pgfqpoint{0.545830in}{1.173062in}}%
\pgfpathclose%
\pgfusepath{fill}%
\end{pgfscope}%
\begin{pgfscope}%
\pgfpathrectangle{\pgfqpoint{0.407592in}{0.222904in}}{\pgfqpoint{1.411163in}{1.411163in}}%
\pgfusepath{clip}%
\pgfsetbuttcap%
\pgfsetroundjoin%
\definecolor{currentfill}{rgb}{0.909804,0.698039,0.423529}%
\pgfsetfillcolor{currentfill}%
\pgfsetfillopacity{0.120000}%
\pgfsetlinewidth{0.000000pt}%
\definecolor{currentstroke}{rgb}{0.000000,0.000000,0.000000}%
\pgfsetstrokecolor{currentstroke}%
\pgfsetdash{}{0pt}%
\pgfpathmoveto{\pgfqpoint{1.662735in}{1.204469in}}%
\pgfpathlineto{\pgfqpoint{1.680273in}{1.213035in}}%
\pgfpathlineto{\pgfqpoint{1.697887in}{1.221889in}}%
\pgfpathlineto{\pgfqpoint{1.705342in}{1.233356in}}%
\pgfpathlineto{\pgfqpoint{1.712787in}{1.245000in}}%
\pgfpathlineto{\pgfqpoint{1.695195in}{1.236405in}}%
\pgfpathlineto{\pgfqpoint{1.677679in}{1.228102in}}%
\pgfpathlineto{\pgfqpoint{1.670213in}{1.216196in}}%
\pgfpathlineto{\pgfqpoint{1.662735in}{1.204469in}}%
\pgfpathclose%
\pgfusepath{fill}%
\end{pgfscope}%
\begin{pgfscope}%
\pgfpathrectangle{\pgfqpoint{0.407592in}{0.222904in}}{\pgfqpoint{1.411163in}{1.411163in}}%
\pgfusepath{clip}%
\pgfsetbuttcap%
\pgfsetroundjoin%
\definecolor{currentfill}{rgb}{0.909804,0.698039,0.423529}%
\pgfsetfillcolor{currentfill}%
\pgfsetfillopacity{0.120000}%
\pgfsetlinewidth{0.000000pt}%
\definecolor{currentstroke}{rgb}{0.000000,0.000000,0.000000}%
\pgfsetstrokecolor{currentstroke}%
\pgfsetdash{}{0pt}%
\pgfpathmoveto{\pgfqpoint{1.023524in}{0.866521in}}%
\pgfpathlineto{\pgfqpoint{1.039922in}{0.862105in}}%
\pgfpathlineto{\pgfqpoint{1.056357in}{0.858685in}}%
\pgfpathlineto{\pgfqpoint{1.064620in}{0.860466in}}%
\pgfpathlineto{\pgfqpoint{1.072852in}{0.862781in}}%
\pgfpathlineto{\pgfqpoint{1.056454in}{0.866168in}}%
\pgfpathlineto{\pgfqpoint{1.040092in}{0.870570in}}%
\pgfpathlineto{\pgfqpoint{1.031824in}{0.868278in}}%
\pgfpathlineto{\pgfqpoint{1.023524in}{0.866521in}}%
\pgfpathclose%
\pgfusepath{fill}%
\end{pgfscope}%
\begin{pgfscope}%
\pgfpathrectangle{\pgfqpoint{0.407592in}{0.222904in}}{\pgfqpoint{1.411163in}{1.411163in}}%
\pgfusepath{clip}%
\pgfsetbuttcap%
\pgfsetroundjoin%
\definecolor{currentfill}{rgb}{0.909804,0.698039,0.423529}%
\pgfsetfillcolor{currentfill}%
\pgfsetfillopacity{0.120000}%
\pgfsetlinewidth{0.000000pt}%
\definecolor{currentstroke}{rgb}{0.000000,0.000000,0.000000}%
\pgfsetstrokecolor{currentstroke}%
\pgfsetdash{}{0pt}%
\pgfpathmoveto{\pgfqpoint{1.318989in}{0.933887in}}%
\pgfpathlineto{\pgfqpoint{1.335851in}{0.939682in}}%
\pgfpathlineto{\pgfqpoint{1.352782in}{0.946044in}}%
\pgfpathlineto{\pgfqpoint{1.360603in}{0.953869in}}%
\pgfpathlineto{\pgfqpoint{1.368408in}{0.962103in}}%
\pgfpathlineto{\pgfqpoint{1.351499in}{0.955898in}}%
\pgfpathlineto{\pgfqpoint{1.334657in}{0.950256in}}%
\pgfpathlineto{\pgfqpoint{1.326831in}{0.941857in}}%
\pgfpathlineto{\pgfqpoint{1.318989in}{0.933887in}}%
\pgfpathclose%
\pgfusepath{fill}%
\end{pgfscope}%
\begin{pgfscope}%
\pgfpathrectangle{\pgfqpoint{0.407592in}{0.222904in}}{\pgfqpoint{1.411163in}{1.411163in}}%
\pgfusepath{clip}%
\pgfsetbuttcap%
\pgfsetroundjoin%
\definecolor{currentfill}{rgb}{0.909804,0.698039,0.423529}%
\pgfsetfillcolor{currentfill}%
\pgfsetfillopacity{0.120000}%
\pgfsetlinewidth{0.000000pt}%
\definecolor{currentstroke}{rgb}{0.000000,0.000000,0.000000}%
\pgfsetstrokecolor{currentstroke}%
\pgfsetdash{}{0pt}%
\pgfpathmoveto{\pgfqpoint{0.693891in}{1.041708in}}%
\pgfpathlineto{\pgfqpoint{0.710153in}{1.026711in}}%
\pgfpathlineto{\pgfqpoint{0.726426in}{1.012114in}}%
\pgfpathlineto{\pgfqpoint{0.735317in}{1.011458in}}%
\pgfpathlineto{\pgfqpoint{0.744167in}{1.011138in}}%
\pgfpathlineto{\pgfqpoint{0.727937in}{1.025954in}}%
\pgfpathlineto{\pgfqpoint{0.711718in}{1.041165in}}%
\pgfpathlineto{\pgfqpoint{0.702825in}{1.041274in}}%
\pgfpathlineto{\pgfqpoint{0.693891in}{1.041708in}}%
\pgfpathclose%
\pgfusepath{fill}%
\end{pgfscope}%
\begin{pgfscope}%
\pgfpathrectangle{\pgfqpoint{0.407592in}{0.222904in}}{\pgfqpoint{1.411163in}{1.411163in}}%
\pgfusepath{clip}%
\pgfsetbuttcap%
\pgfsetroundjoin%
\definecolor{currentfill}{rgb}{0.909804,0.698039,0.423529}%
\pgfsetfillcolor{currentfill}%
\pgfsetfillopacity{0.120000}%
\pgfsetlinewidth{0.000000pt}%
\definecolor{currentstroke}{rgb}{0.000000,0.000000,0.000000}%
\pgfsetstrokecolor{currentstroke}%
\pgfsetdash{}{0pt}%
\pgfpathmoveto{\pgfqpoint{0.578289in}{1.139084in}}%
\pgfpathlineto{\pgfqpoint{0.594533in}{1.122442in}}%
\pgfpathlineto{\pgfqpoint{0.610786in}{1.106057in}}%
\pgfpathlineto{\pgfqpoint{0.619889in}{1.105446in}}%
\pgfpathlineto{\pgfqpoint{0.628950in}{1.105123in}}%
\pgfpathlineto{\pgfqpoint{0.612740in}{1.121707in}}%
\pgfpathlineto{\pgfqpoint{0.596539in}{1.138540in}}%
\pgfpathlineto{\pgfqpoint{0.587435in}{1.138673in}}%
\pgfpathlineto{\pgfqpoint{0.578289in}{1.139084in}}%
\pgfpathclose%
\pgfusepath{fill}%
\end{pgfscope}%
\begin{pgfscope}%
\pgfpathrectangle{\pgfqpoint{0.407592in}{0.222904in}}{\pgfqpoint{1.411163in}{1.411163in}}%
\pgfusepath{clip}%
\pgfsetbuttcap%
\pgfsetroundjoin%
\definecolor{currentfill}{rgb}{0.909804,0.698039,0.423529}%
\pgfsetfillcolor{currentfill}%
\pgfsetfillopacity{0.120000}%
\pgfsetlinewidth{0.000000pt}%
\definecolor{currentstroke}{rgb}{0.000000,0.000000,0.000000}%
\pgfsetstrokecolor{currentstroke}%
\pgfsetdash{}{0pt}%
\pgfpathmoveto{\pgfqpoint{1.483004in}{1.046464in}}%
\pgfpathlineto{\pgfqpoint{1.500204in}{1.054313in}}%
\pgfpathlineto{\pgfqpoint{1.517476in}{1.062539in}}%
\pgfpathlineto{\pgfqpoint{1.525117in}{1.072504in}}%
\pgfpathlineto{\pgfqpoint{1.532746in}{1.082745in}}%
\pgfpathlineto{\pgfqpoint{1.515494in}{1.074769in}}%
\pgfpathlineto{\pgfqpoint{1.498315in}{1.067173in}}%
\pgfpathlineto{\pgfqpoint{1.490666in}{1.056677in}}%
\pgfpathlineto{\pgfqpoint{1.483004in}{1.046464in}}%
\pgfpathclose%
\pgfusepath{fill}%
\end{pgfscope}%
\begin{pgfscope}%
\pgfpathrectangle{\pgfqpoint{0.407592in}{0.222904in}}{\pgfqpoint{1.411163in}{1.411163in}}%
\pgfusepath{clip}%
\pgfsetbuttcap%
\pgfsetroundjoin%
\definecolor{currentfill}{rgb}{0.909804,0.698039,0.423529}%
\pgfsetfillcolor{currentfill}%
\pgfsetfillopacity{0.120000}%
\pgfsetlinewidth{0.000000pt}%
\definecolor{currentstroke}{rgb}{0.000000,0.000000,0.000000}%
\pgfsetstrokecolor{currentstroke}%
\pgfsetdash{}{0pt}%
\pgfpathmoveto{\pgfqpoint{0.891716in}{0.910071in}}%
\pgfpathlineto{\pgfqpoint{0.908036in}{0.900351in}}%
\pgfpathlineto{\pgfqpoint{0.924376in}{0.891435in}}%
\pgfpathlineto{\pgfqpoint{0.932895in}{0.891358in}}%
\pgfpathlineto{\pgfqpoint{0.941378in}{0.891731in}}%
\pgfpathlineto{\pgfqpoint{0.925080in}{0.900766in}}%
\pgfpathlineto{\pgfqpoint{0.908802in}{0.910621in}}%
\pgfpathlineto{\pgfqpoint{0.900278in}{0.910127in}}%
\pgfpathlineto{\pgfqpoint{0.891716in}{0.910071in}}%
\pgfpathclose%
\pgfusepath{fill}%
\end{pgfscope}%
\begin{pgfscope}%
\pgfpathrectangle{\pgfqpoint{0.407592in}{0.222904in}}{\pgfqpoint{1.411163in}{1.411163in}}%
\pgfusepath{clip}%
\pgfsetbuttcap%
\pgfsetroundjoin%
\definecolor{currentfill}{rgb}{0.909804,0.698039,0.423529}%
\pgfsetfillcolor{currentfill}%
\pgfsetfillopacity{0.120000}%
\pgfsetlinewidth{0.000000pt}%
\definecolor{currentstroke}{rgb}{0.000000,0.000000,0.000000}%
\pgfsetstrokecolor{currentstroke}%
\pgfsetdash{}{0pt}%
\pgfpathmoveto{\pgfqpoint{1.597773in}{1.142228in}}%
\pgfpathlineto{\pgfqpoint{1.615204in}{1.150708in}}%
\pgfpathlineto{\pgfqpoint{1.632710in}{1.159498in}}%
\pgfpathlineto{\pgfqpoint{1.640234in}{1.170433in}}%
\pgfpathlineto{\pgfqpoint{1.647746in}{1.181579in}}%
\pgfpathlineto{\pgfqpoint{1.630262in}{1.173051in}}%
\pgfpathlineto{\pgfqpoint{1.612852in}{1.164837in}}%
\pgfpathlineto{\pgfqpoint{1.605318in}{1.153426in}}%
\pgfpathlineto{\pgfqpoint{1.597773in}{1.142228in}}%
\pgfpathclose%
\pgfusepath{fill}%
\end{pgfscope}%
\begin{pgfscope}%
\pgfpathrectangle{\pgfqpoint{0.407592in}{0.222904in}}{\pgfqpoint{1.411163in}{1.411163in}}%
\pgfusepath{clip}%
\pgfsetbuttcap%
\pgfsetroundjoin%
\definecolor{currentfill}{rgb}{0.909804,0.698039,0.423529}%
\pgfsetfillcolor{currentfill}%
\pgfsetfillopacity{0.120000}%
\pgfsetlinewidth{0.000000pt}%
\definecolor{currentstroke}{rgb}{0.000000,0.000000,0.000000}%
\pgfsetstrokecolor{currentstroke}%
\pgfsetdash{}{0pt}%
\pgfpathmoveto{\pgfqpoint{0.809207in}{0.956520in}}%
\pgfpathlineto{\pgfqpoint{0.825501in}{0.944217in}}%
\pgfpathlineto{\pgfqpoint{0.841810in}{0.932535in}}%
\pgfpathlineto{\pgfqpoint{0.850491in}{0.931947in}}%
\pgfpathlineto{\pgfqpoint{0.859133in}{0.931753in}}%
\pgfpathlineto{\pgfqpoint{0.842867in}{0.943626in}}%
\pgfpathlineto{\pgfqpoint{0.826616in}{0.956127in}}%
\pgfpathlineto{\pgfqpoint{0.817931in}{0.956132in}}%
\pgfpathlineto{\pgfqpoint{0.809207in}{0.956520in}}%
\pgfpathclose%
\pgfusepath{fill}%
\end{pgfscope}%
\begin{pgfscope}%
\pgfpathrectangle{\pgfqpoint{0.407592in}{0.222904in}}{\pgfqpoint{1.411163in}{1.411163in}}%
\pgfusepath{clip}%
\pgfsetbuttcap%
\pgfsetroundjoin%
\definecolor{currentfill}{rgb}{0.909804,0.698039,0.423529}%
\pgfsetfillcolor{currentfill}%
\pgfsetfillopacity{0.120000}%
\pgfsetlinewidth{0.000000pt}%
\definecolor{currentstroke}{rgb}{0.000000,0.000000,0.000000}%
\pgfsetstrokecolor{currentstroke}%
\pgfsetdash{}{0pt}%
\pgfpathmoveto{\pgfqpoint{1.368408in}{0.962103in}}%
\pgfpathlineto{\pgfqpoint{1.385386in}{0.968832in}}%
\pgfpathlineto{\pgfqpoint{1.402433in}{0.976048in}}%
\pgfpathlineto{\pgfqpoint{1.410201in}{0.984495in}}%
\pgfpathlineto{\pgfqpoint{1.417954in}{0.993309in}}%
\pgfpathlineto{\pgfqpoint{1.400927in}{0.986284in}}%
\pgfpathlineto{\pgfqpoint{1.383970in}{0.979747in}}%
\pgfpathlineto{\pgfqpoint{1.376197in}{0.970733in}}%
\pgfpathlineto{\pgfqpoint{1.368408in}{0.962103in}}%
\pgfpathclose%
\pgfusepath{fill}%
\end{pgfscope}%
\begin{pgfscope}%
\pgfpathrectangle{\pgfqpoint{0.407592in}{0.222904in}}{\pgfqpoint{1.411163in}{1.411163in}}%
\pgfusepath{clip}%
\pgfsetbuttcap%
\pgfsetroundjoin%
\definecolor{currentfill}{rgb}{0.909804,0.698039,0.423529}%
\pgfsetfillcolor{currentfill}%
\pgfsetfillopacity{0.120000}%
\pgfsetlinewidth{0.000000pt}%
\definecolor{currentstroke}{rgb}{0.000000,0.000000,0.000000}%
\pgfsetstrokecolor{currentstroke}%
\pgfsetdash{}{0pt}%
\pgfpathmoveto{\pgfqpoint{0.974046in}{0.876284in}}%
\pgfpathlineto{\pgfqpoint{0.990421in}{0.869938in}}%
\pgfpathlineto{\pgfqpoint{1.006825in}{0.864545in}}%
\pgfpathlineto{\pgfqpoint{1.015191in}{0.865283in}}%
\pgfpathlineto{\pgfqpoint{1.023524in}{0.866521in}}%
\pgfpathlineto{\pgfqpoint{1.007159in}{0.871928in}}%
\pgfpathlineto{\pgfqpoint{0.990824in}{0.878311in}}%
\pgfpathlineto{\pgfqpoint{0.982453in}{0.877050in}}%
\pgfpathlineto{\pgfqpoint{0.974046in}{0.876284in}}%
\pgfpathclose%
\pgfusepath{fill}%
\end{pgfscope}%
\begin{pgfscope}%
\pgfpathrectangle{\pgfqpoint{0.407592in}{0.222904in}}{\pgfqpoint{1.411163in}{1.411163in}}%
\pgfusepath{clip}%
\pgfsetbuttcap%
\pgfsetroundjoin%
\definecolor{currentfill}{rgb}{0.909804,0.698039,0.423529}%
\pgfsetfillcolor{currentfill}%
\pgfsetfillopacity{0.120000}%
\pgfsetlinewidth{0.000000pt}%
\definecolor{currentstroke}{rgb}{0.000000,0.000000,0.000000}%
\pgfsetstrokecolor{currentstroke}%
\pgfsetdash{}{0pt}%
\pgfpathmoveto{\pgfqpoint{0.610786in}{1.106057in}}%
\pgfpathlineto{\pgfqpoint{0.627049in}{1.089948in}}%
\pgfpathlineto{\pgfqpoint{0.643322in}{1.074138in}}%
\pgfpathlineto{\pgfqpoint{0.652382in}{1.073312in}}%
\pgfpathlineto{\pgfqpoint{0.661400in}{1.072783in}}%
\pgfpathlineto{\pgfqpoint{0.645170in}{1.088807in}}%
\pgfpathlineto{\pgfqpoint{0.628950in}{1.105123in}}%
\pgfpathlineto{\pgfqpoint{0.619889in}{1.105446in}}%
\pgfpathlineto{\pgfqpoint{0.610786in}{1.106057in}}%
\pgfpathclose%
\pgfusepath{fill}%
\end{pgfscope}%
\begin{pgfscope}%
\pgfpathrectangle{\pgfqpoint{0.407592in}{0.222904in}}{\pgfqpoint{1.411163in}{1.411163in}}%
\pgfusepath{clip}%
\pgfsetbuttcap%
\pgfsetroundjoin%
\definecolor{currentfill}{rgb}{0.909804,0.698039,0.423529}%
\pgfsetfillcolor{currentfill}%
\pgfsetfillopacity{0.120000}%
\pgfsetlinewidth{0.000000pt}%
\definecolor{currentstroke}{rgb}{0.000000,0.000000,0.000000}%
\pgfsetstrokecolor{currentstroke}%
\pgfsetdash{}{0pt}%
\pgfpathmoveto{\pgfqpoint{0.726426in}{1.012114in}}%
\pgfpathlineto{\pgfqpoint{0.742710in}{0.997948in}}%
\pgfpathlineto{\pgfqpoint{0.759007in}{0.984248in}}%
\pgfpathlineto{\pgfqpoint{0.767854in}{0.983365in}}%
\pgfpathlineto{\pgfqpoint{0.776661in}{0.982827in}}%
\pgfpathlineto{\pgfqpoint{0.760408in}{0.996751in}}%
\pgfpathlineto{\pgfqpoint{0.744167in}{1.011138in}}%
\pgfpathlineto{\pgfqpoint{0.735317in}{1.011458in}}%
\pgfpathlineto{\pgfqpoint{0.726426in}{1.012114in}}%
\pgfpathclose%
\pgfusepath{fill}%
\end{pgfscope}%
\begin{pgfscope}%
\pgfpathrectangle{\pgfqpoint{0.407592in}{0.222904in}}{\pgfqpoint{1.411163in}{1.411163in}}%
\pgfusepath{clip}%
\pgfsetbuttcap%
\pgfsetroundjoin%
\definecolor{currentfill}{rgb}{0.909804,0.698039,0.423529}%
\pgfsetfillcolor{currentfill}%
\pgfsetfillopacity{0.120000}%
\pgfsetlinewidth{0.000000pt}%
\definecolor{currentstroke}{rgb}{0.000000,0.000000,0.000000}%
\pgfsetstrokecolor{currentstroke}%
\pgfsetdash{}{0pt}%
\pgfpathmoveto{\pgfqpoint{1.155119in}{0.865670in}}%
\pgfpathlineto{\pgfqpoint{1.171716in}{0.867182in}}%
\pgfpathlineto{\pgfqpoint{1.188366in}{0.869556in}}%
\pgfpathlineto{\pgfqpoint{1.196422in}{0.873682in}}%
\pgfpathlineto{\pgfqpoint{1.204453in}{0.878338in}}%
\pgfpathlineto{\pgfqpoint{1.187830in}{0.875922in}}%
\pgfpathlineto{\pgfqpoint{1.171262in}{0.874370in}}%
\pgfpathlineto{\pgfqpoint{1.163203in}{0.869745in}}%
\pgfpathlineto{\pgfqpoint{1.155119in}{0.865670in}}%
\pgfpathclose%
\pgfusepath{fill}%
\end{pgfscope}%
\begin{pgfscope}%
\pgfpathrectangle{\pgfqpoint{0.407592in}{0.222904in}}{\pgfqpoint{1.411163in}{1.411163in}}%
\pgfusepath{clip}%
\pgfsetbuttcap%
\pgfsetroundjoin%
\definecolor{currentfill}{rgb}{0.909804,0.698039,0.423529}%
\pgfsetfillcolor{currentfill}%
\pgfsetfillopacity{0.120000}%
\pgfsetlinewidth{0.000000pt}%
\definecolor{currentstroke}{rgb}{0.000000,0.000000,0.000000}%
\pgfsetstrokecolor{currentstroke}%
\pgfsetdash{}{0pt}%
\pgfpathmoveto{\pgfqpoint{1.532746in}{1.082745in}}%
\pgfpathlineto{\pgfqpoint{1.550071in}{1.091079in}}%
\pgfpathlineto{\pgfqpoint{1.567470in}{1.099749in}}%
\pgfpathlineto{\pgfqpoint{1.575064in}{1.110003in}}%
\pgfpathlineto{\pgfqpoint{1.582646in}{1.120508in}}%
\pgfpathlineto{\pgfqpoint{1.565268in}{1.112091in}}%
\pgfpathlineto{\pgfqpoint{1.547964in}{1.104016in}}%
\pgfpathlineto{\pgfqpoint{1.540361in}{1.093253in}}%
\pgfpathlineto{\pgfqpoint{1.532746in}{1.082745in}}%
\pgfpathclose%
\pgfusepath{fill}%
\end{pgfscope}%
\begin{pgfscope}%
\pgfpathrectangle{\pgfqpoint{0.407592in}{0.222904in}}{\pgfqpoint{1.411163in}{1.411163in}}%
\pgfusepath{clip}%
\pgfsetbuttcap%
\pgfsetroundjoin%
\definecolor{currentfill}{rgb}{0.909804,0.698039,0.423529}%
\pgfsetfillcolor{currentfill}%
\pgfsetfillopacity{0.120000}%
\pgfsetlinewidth{0.000000pt}%
\definecolor{currentstroke}{rgb}{0.000000,0.000000,0.000000}%
\pgfsetstrokecolor{currentstroke}%
\pgfsetdash{}{0pt}%
\pgfpathmoveto{\pgfqpoint{1.712787in}{1.245000in}}%
\pgfpathlineto{\pgfqpoint{1.730454in}{1.253875in}}%
\pgfpathlineto{\pgfqpoint{1.748197in}{1.263018in}}%
\pgfpathlineto{\pgfqpoint{1.755608in}{1.274579in}}%
\pgfpathlineto{\pgfqpoint{1.737876in}{1.265562in}}%
\pgfpathlineto{\pgfqpoint{1.720220in}{1.256815in}}%
\pgfusepath{fill}%
\end{pgfscope}%
\begin{pgfscope}%
\pgfpathrectangle{\pgfqpoint{0.407592in}{0.222904in}}{\pgfqpoint{1.411163in}{1.411163in}}%
\pgfusepath{clip}%
\pgfsetbuttcap%
\pgfsetroundjoin%
\definecolor{currentfill}{rgb}{0.909804,0.698039,0.423529}%
\pgfsetfillcolor{currentfill}%
\pgfsetfillopacity{0.120000}%
\pgfsetlinewidth{0.000000pt}%
\definecolor{currentstroke}{rgb}{0.000000,0.000000,0.000000}%
\pgfsetstrokecolor{currentstroke}%
\pgfsetdash{}{0pt}%
\pgfpathmoveto{\pgfqpoint{1.204453in}{0.878338in}}%
\pgfpathlineto{\pgfqpoint{1.221133in}{0.881569in}}%
\pgfpathlineto{\pgfqpoint{1.237874in}{0.885564in}}%
\pgfpathlineto{\pgfqpoint{1.245856in}{0.890757in}}%
\pgfpathlineto{\pgfqpoint{1.253817in}{0.896454in}}%
\pgfpathlineto{\pgfqpoint{1.237101in}{0.892466in}}%
\pgfpathlineto{\pgfqpoint{1.220446in}{0.889240in}}%
\pgfpathlineto{\pgfqpoint{1.212460in}{0.883525in}}%
\pgfpathlineto{\pgfqpoint{1.204453in}{0.878338in}}%
\pgfpathclose%
\pgfusepath{fill}%
\end{pgfscope}%
\begin{pgfscope}%
\pgfpathrectangle{\pgfqpoint{0.407592in}{0.222904in}}{\pgfqpoint{1.411163in}{1.411163in}}%
\pgfusepath{clip}%
\pgfsetbuttcap%
\pgfsetroundjoin%
\definecolor{currentfill}{rgb}{0.909804,0.698039,0.423529}%
\pgfsetfillcolor{currentfill}%
\pgfsetfillopacity{0.120000}%
\pgfsetlinewidth{0.000000pt}%
\definecolor{currentstroke}{rgb}{0.000000,0.000000,0.000000}%
\pgfsetstrokecolor{currentstroke}%
\pgfsetdash{}{0pt}%
\pgfpathmoveto{\pgfqpoint{1.417954in}{0.993309in}}%
\pgfpathlineto{\pgfqpoint{1.435051in}{1.000789in}}%
\pgfpathlineto{\pgfqpoint{1.452220in}{1.008691in}}%
\pgfpathlineto{\pgfqpoint{1.459937in}{1.017650in}}%
\pgfpathlineto{\pgfqpoint{1.467640in}{1.026940in}}%
\pgfpathlineto{\pgfqpoint{1.450492in}{1.019253in}}%
\pgfpathlineto{\pgfqpoint{1.433415in}{1.011992in}}%
\pgfpathlineto{\pgfqpoint{1.425692in}{1.002479in}}%
\pgfpathlineto{\pgfqpoint{1.417954in}{0.993309in}}%
\pgfpathclose%
\pgfusepath{fill}%
\end{pgfscope}%
\begin{pgfscope}%
\pgfpathrectangle{\pgfqpoint{0.407592in}{0.222904in}}{\pgfqpoint{1.411163in}{1.411163in}}%
\pgfusepath{clip}%
\pgfsetbuttcap%
\pgfsetroundjoin%
\definecolor{currentfill}{rgb}{0.909804,0.698039,0.423529}%
\pgfsetfillcolor{currentfill}%
\pgfsetfillopacity{0.120000}%
\pgfsetlinewidth{0.000000pt}%
\definecolor{currentstroke}{rgb}{0.000000,0.000000,0.000000}%
\pgfsetstrokecolor{currentstroke}%
\pgfsetdash{}{0pt}%
\pgfpathmoveto{\pgfqpoint{1.105771in}{0.859027in}}%
\pgfpathlineto{\pgfqpoint{1.122298in}{0.858620in}}%
\pgfpathlineto{\pgfqpoint{1.138873in}{0.859153in}}%
\pgfpathlineto{\pgfqpoint{1.147010in}{0.862142in}}%
\pgfpathlineto{\pgfqpoint{1.155119in}{0.865670in}}%
\pgfpathlineto{\pgfqpoint{1.138575in}{0.865066in}}%
\pgfpathlineto{\pgfqpoint{1.122080in}{0.865411in}}%
\pgfpathlineto{\pgfqpoint{1.113940in}{0.861942in}}%
\pgfpathlineto{\pgfqpoint{1.105771in}{0.859027in}}%
\pgfpathclose%
\pgfusepath{fill}%
\end{pgfscope}%
\begin{pgfscope}%
\pgfpathrectangle{\pgfqpoint{0.407592in}{0.222904in}}{\pgfqpoint{1.411163in}{1.411163in}}%
\pgfusepath{clip}%
\pgfsetbuttcap%
\pgfsetroundjoin%
\definecolor{currentfill}{rgb}{0.909804,0.698039,0.423529}%
\pgfsetfillcolor{currentfill}%
\pgfsetfillopacity{0.120000}%
\pgfsetlinewidth{0.000000pt}%
\definecolor{currentstroke}{rgb}{0.000000,0.000000,0.000000}%
\pgfsetstrokecolor{currentstroke}%
\pgfsetdash{}{0pt}%
\pgfpathmoveto{\pgfqpoint{0.494818in}{1.208735in}}%
\pgfpathlineto{\pgfqpoint{0.511067in}{1.191420in}}%
\pgfpathlineto{\pgfqpoint{0.527325in}{1.174301in}}%
\pgfpathlineto{\pgfqpoint{0.536599in}{1.173553in}}%
\pgfpathlineto{\pgfqpoint{0.545830in}{1.173062in}}%
\pgfpathlineto{\pgfqpoint{0.529615in}{1.190365in}}%
\pgfpathlineto{\pgfqpoint{0.513409in}{1.207856in}}%
\pgfpathlineto{\pgfqpoint{0.504135in}{1.208171in}}%
\pgfpathlineto{\pgfqpoint{0.494818in}{1.208735in}}%
\pgfpathclose%
\pgfusepath{fill}%
\end{pgfscope}%
\begin{pgfscope}%
\pgfpathrectangle{\pgfqpoint{0.407592in}{0.222904in}}{\pgfqpoint{1.411163in}{1.411163in}}%
\pgfusepath{clip}%
\pgfsetbuttcap%
\pgfsetroundjoin%
\definecolor{currentfill}{rgb}{0.909804,0.698039,0.423529}%
\pgfsetfillcolor{currentfill}%
\pgfsetfillopacity{0.120000}%
\pgfsetlinewidth{0.000000pt}%
\definecolor{currentstroke}{rgb}{0.000000,0.000000,0.000000}%
\pgfsetstrokecolor{currentstroke}%
\pgfsetdash{}{0pt}%
\pgfpathmoveto{\pgfqpoint{1.647746in}{1.181579in}}%
\pgfpathlineto{\pgfqpoint{1.665305in}{1.190406in}}%
\pgfpathlineto{\pgfqpoint{1.682941in}{1.199517in}}%
\pgfpathlineto{\pgfqpoint{1.690419in}{1.210607in}}%
\pgfpathlineto{\pgfqpoint{1.697887in}{1.221889in}}%
\pgfpathlineto{\pgfqpoint{1.680273in}{1.213035in}}%
\pgfpathlineto{\pgfqpoint{1.662735in}{1.204469in}}%
\pgfpathlineto{\pgfqpoint{1.655247in}{1.192927in}}%
\pgfpathlineto{\pgfqpoint{1.647746in}{1.181579in}}%
\pgfpathclose%
\pgfusepath{fill}%
\end{pgfscope}%
\begin{pgfscope}%
\pgfpathrectangle{\pgfqpoint{0.407592in}{0.222904in}}{\pgfqpoint{1.411163in}{1.411163in}}%
\pgfusepath{clip}%
\pgfsetbuttcap%
\pgfsetroundjoin%
\definecolor{currentfill}{rgb}{0.909804,0.698039,0.423529}%
\pgfsetfillcolor{currentfill}%
\pgfsetfillopacity{0.120000}%
\pgfsetlinewidth{0.000000pt}%
\definecolor{currentstroke}{rgb}{0.000000,0.000000,0.000000}%
\pgfsetstrokecolor{currentstroke}%
\pgfsetdash{}{0pt}%
\pgfpathmoveto{\pgfqpoint{1.253817in}{0.896454in}}%
\pgfpathlineto{\pgfqpoint{1.270595in}{0.901152in}}%
\pgfpathlineto{\pgfqpoint{1.287437in}{0.906512in}}%
\pgfpathlineto{\pgfqpoint{1.295353in}{0.912660in}}%
\pgfpathlineto{\pgfqpoint{1.303250in}{0.919278in}}%
\pgfpathlineto{\pgfqpoint{1.286431in}{0.913981in}}%
\pgfpathlineto{\pgfqpoint{1.269677in}{0.909342in}}%
\pgfpathlineto{\pgfqpoint{1.261757in}{0.902651in}}%
\pgfpathlineto{\pgfqpoint{1.253817in}{0.896454in}}%
\pgfpathclose%
\pgfusepath{fill}%
\end{pgfscope}%
\begin{pgfscope}%
\pgfpathrectangle{\pgfqpoint{0.407592in}{0.222904in}}{\pgfqpoint{1.411163in}{1.411163in}}%
\pgfusepath{clip}%
\pgfsetbuttcap%
\pgfsetroundjoin%
\definecolor{currentfill}{rgb}{0.909804,0.698039,0.423529}%
\pgfsetfillcolor{currentfill}%
\pgfsetfillopacity{0.120000}%
\pgfsetlinewidth{0.000000pt}%
\definecolor{currentstroke}{rgb}{0.000000,0.000000,0.000000}%
\pgfsetstrokecolor{currentstroke}%
\pgfsetdash{}{0pt}%
\pgfpathmoveto{\pgfqpoint{0.841810in}{0.932535in}}%
\pgfpathlineto{\pgfqpoint{0.858136in}{0.921518in}}%
\pgfpathlineto{\pgfqpoint{0.874480in}{0.911207in}}%
\pgfpathlineto{\pgfqpoint{0.883117in}{0.910437in}}%
\pgfpathlineto{\pgfqpoint{0.891716in}{0.910071in}}%
\pgfpathlineto{\pgfqpoint{0.875416in}{0.920553in}}%
\pgfpathlineto{\pgfqpoint{0.859133in}{0.931753in}}%
\pgfpathlineto{\pgfqpoint{0.850491in}{0.931947in}}%
\pgfpathlineto{\pgfqpoint{0.841810in}{0.932535in}}%
\pgfpathclose%
\pgfusepath{fill}%
\end{pgfscope}%
\begin{pgfscope}%
\pgfpathrectangle{\pgfqpoint{0.407592in}{0.222904in}}{\pgfqpoint{1.411163in}{1.411163in}}%
\pgfusepath{clip}%
\pgfsetbuttcap%
\pgfsetroundjoin%
\definecolor{currentfill}{rgb}{0.909804,0.698039,0.423529}%
\pgfsetfillcolor{currentfill}%
\pgfsetfillopacity{0.120000}%
\pgfsetlinewidth{0.000000pt}%
\definecolor{currentstroke}{rgb}{0.000000,0.000000,0.000000}%
\pgfsetstrokecolor{currentstroke}%
\pgfsetdash{}{0pt}%
\pgfpathmoveto{\pgfqpoint{0.924376in}{0.891435in}}%
\pgfpathlineto{\pgfqpoint{0.940740in}{0.883360in}}%
\pgfpathlineto{\pgfqpoint{0.957128in}{0.876160in}}%
\pgfpathlineto{\pgfqpoint{0.965605in}{0.875994in}}%
\pgfpathlineto{\pgfqpoint{0.974046in}{0.876284in}}%
\pgfpathlineto{\pgfqpoint{0.957700in}{0.883558in}}%
\pgfpathlineto{\pgfqpoint{0.941378in}{0.891731in}}%
\pgfpathlineto{\pgfqpoint{0.932895in}{0.891358in}}%
\pgfpathlineto{\pgfqpoint{0.924376in}{0.891435in}}%
\pgfpathclose%
\pgfusepath{fill}%
\end{pgfscope}%
\begin{pgfscope}%
\pgfpathrectangle{\pgfqpoint{0.407592in}{0.222904in}}{\pgfqpoint{1.411163in}{1.411163in}}%
\pgfusepath{clip}%
\pgfsetbuttcap%
\pgfsetroundjoin%
\definecolor{currentfill}{rgb}{0.909804,0.698039,0.423529}%
\pgfsetfillcolor{currentfill}%
\pgfsetfillopacity{0.120000}%
\pgfsetlinewidth{0.000000pt}%
\definecolor{currentstroke}{rgb}{0.000000,0.000000,0.000000}%
\pgfsetstrokecolor{currentstroke}%
\pgfsetdash{}{0pt}%
\pgfpathmoveto{\pgfqpoint{1.056357in}{0.858685in}}%
\pgfpathlineto{\pgfqpoint{1.072831in}{0.856258in}}%
\pgfpathlineto{\pgfqpoint{1.089346in}{0.854809in}}%
\pgfpathlineto{\pgfqpoint{1.097574in}{0.856654in}}%
\pgfpathlineto{\pgfqpoint{1.105771in}{0.859027in}}%
\pgfpathlineto{\pgfqpoint{1.089290in}{0.860406in}}%
\pgfpathlineto{\pgfqpoint{1.072852in}{0.862781in}}%
\pgfpathlineto{\pgfqpoint{1.064620in}{0.860466in}}%
\pgfpathlineto{\pgfqpoint{1.056357in}{0.858685in}}%
\pgfpathclose%
\pgfusepath{fill}%
\end{pgfscope}%
\begin{pgfscope}%
\pgfpathrectangle{\pgfqpoint{0.407592in}{0.222904in}}{\pgfqpoint{1.411163in}{1.411163in}}%
\pgfusepath{clip}%
\pgfsetbuttcap%
\pgfsetroundjoin%
\definecolor{currentfill}{rgb}{0.909804,0.698039,0.423529}%
\pgfsetfillcolor{currentfill}%
\pgfsetfillopacity{0.120000}%
\pgfsetlinewidth{0.000000pt}%
\definecolor{currentstroke}{rgb}{0.000000,0.000000,0.000000}%
\pgfsetstrokecolor{currentstroke}%
\pgfsetdash{}{0pt}%
\pgfpathmoveto{\pgfqpoint{0.643322in}{1.074138in}}%
\pgfpathlineto{\pgfqpoint{0.659606in}{1.058651in}}%
\pgfpathlineto{\pgfqpoint{0.675900in}{1.043512in}}%
\pgfpathlineto{\pgfqpoint{0.684916in}{1.042457in}}%
\pgfpathlineto{\pgfqpoint{0.693891in}{1.041708in}}%
\pgfpathlineto{\pgfqpoint{0.677641in}{1.057074in}}%
\pgfpathlineto{\pgfqpoint{0.661400in}{1.072783in}}%
\pgfpathlineto{\pgfqpoint{0.652382in}{1.073312in}}%
\pgfpathlineto{\pgfqpoint{0.643322in}{1.074138in}}%
\pgfpathclose%
\pgfusepath{fill}%
\end{pgfscope}%
\begin{pgfscope}%
\pgfpathrectangle{\pgfqpoint{0.407592in}{0.222904in}}{\pgfqpoint{1.411163in}{1.411163in}}%
\pgfusepath{clip}%
\pgfsetbuttcap%
\pgfsetroundjoin%
\definecolor{currentfill}{rgb}{0.909804,0.698039,0.423529}%
\pgfsetfillcolor{currentfill}%
\pgfsetfillopacity{0.120000}%
\pgfsetlinewidth{0.000000pt}%
\definecolor{currentstroke}{rgb}{0.000000,0.000000,0.000000}%
\pgfsetstrokecolor{currentstroke}%
\pgfsetdash{}{0pt}%
\pgfpathmoveto{\pgfqpoint{1.303250in}{0.919278in}}%
\pgfpathlineto{\pgfqpoint{1.320135in}{0.925186in}}%
\pgfpathlineto{\pgfqpoint{1.337087in}{0.931661in}}%
\pgfpathlineto{\pgfqpoint{1.344943in}{0.938638in}}%
\pgfpathlineto{\pgfqpoint{1.352782in}{0.946044in}}%
\pgfpathlineto{\pgfqpoint{1.335851in}{0.939682in}}%
\pgfpathlineto{\pgfqpoint{1.318989in}{0.933887in}}%
\pgfpathlineto{\pgfqpoint{1.311128in}{0.926356in}}%
\pgfpathlineto{\pgfqpoint{1.303250in}{0.919278in}}%
\pgfpathclose%
\pgfusepath{fill}%
\end{pgfscope}%
\begin{pgfscope}%
\pgfpathrectangle{\pgfqpoint{0.407592in}{0.222904in}}{\pgfqpoint{1.411163in}{1.411163in}}%
\pgfusepath{clip}%
\pgfsetbuttcap%
\pgfsetroundjoin%
\definecolor{currentfill}{rgb}{0.909804,0.698039,0.423529}%
\pgfsetfillcolor{currentfill}%
\pgfsetfillopacity{0.120000}%
\pgfsetlinewidth{0.000000pt}%
\definecolor{currentstroke}{rgb}{0.000000,0.000000,0.000000}%
\pgfsetstrokecolor{currentstroke}%
\pgfsetdash{}{0pt}%
\pgfpathmoveto{\pgfqpoint{0.759007in}{0.984248in}}%
\pgfpathlineto{\pgfqpoint{0.775317in}{0.971049in}}%
\pgfpathlineto{\pgfqpoint{0.791640in}{0.958388in}}%
\pgfpathlineto{\pgfqpoint{0.800444in}{0.957276in}}%
\pgfpathlineto{\pgfqpoint{0.809207in}{0.956520in}}%
\pgfpathlineto{\pgfqpoint{0.792928in}{0.969403in}}%
\pgfpathlineto{\pgfqpoint{0.776661in}{0.982827in}}%
\pgfpathlineto{\pgfqpoint{0.767854in}{0.983365in}}%
\pgfpathlineto{\pgfqpoint{0.759007in}{0.984248in}}%
\pgfpathclose%
\pgfusepath{fill}%
\end{pgfscope}%
\begin{pgfscope}%
\pgfpathrectangle{\pgfqpoint{0.407592in}{0.222904in}}{\pgfqpoint{1.411163in}{1.411163in}}%
\pgfusepath{clip}%
\pgfsetbuttcap%
\pgfsetroundjoin%
\definecolor{currentfill}{rgb}{0.909804,0.698039,0.423529}%
\pgfsetfillcolor{currentfill}%
\pgfsetfillopacity{0.120000}%
\pgfsetlinewidth{0.000000pt}%
\definecolor{currentstroke}{rgb}{0.000000,0.000000,0.000000}%
\pgfsetstrokecolor{currentstroke}%
\pgfsetdash{}{0pt}%
\pgfpathmoveto{\pgfqpoint{1.467640in}{1.026940in}}%
\pgfpathlineto{\pgfqpoint{1.484860in}{1.035024in}}%
\pgfpathlineto{\pgfqpoint{1.502153in}{1.043481in}}%
\pgfpathlineto{\pgfqpoint{1.509822in}{1.052861in}}%
\pgfpathlineto{\pgfqpoint{1.517476in}{1.062539in}}%
\pgfpathlineto{\pgfqpoint{1.500204in}{1.054313in}}%
\pgfpathlineto{\pgfqpoint{1.483004in}{1.046464in}}%
\pgfpathlineto{\pgfqpoint{1.475329in}{1.036548in}}%
\pgfpathlineto{\pgfqpoint{1.467640in}{1.026940in}}%
\pgfpathclose%
\pgfusepath{fill}%
\end{pgfscope}%
\begin{pgfscope}%
\pgfpathrectangle{\pgfqpoint{0.407592in}{0.222904in}}{\pgfqpoint{1.411163in}{1.411163in}}%
\pgfusepath{clip}%
\pgfsetbuttcap%
\pgfsetroundjoin%
\definecolor{currentfill}{rgb}{0.909804,0.698039,0.423529}%
\pgfsetfillcolor{currentfill}%
\pgfsetfillopacity{0.120000}%
\pgfsetlinewidth{0.000000pt}%
\definecolor{currentstroke}{rgb}{0.000000,0.000000,0.000000}%
\pgfsetstrokecolor{currentstroke}%
\pgfsetdash{}{0pt}%
\pgfpathmoveto{\pgfqpoint{0.527325in}{1.174301in}}%
\pgfpathlineto{\pgfqpoint{0.543593in}{1.157394in}}%
\pgfpathlineto{\pgfqpoint{0.559870in}{1.140715in}}%
\pgfpathlineto{\pgfqpoint{0.569101in}{1.139767in}}%
\pgfpathlineto{\pgfqpoint{0.578289in}{1.139084in}}%
\pgfpathlineto{\pgfqpoint{0.562055in}{1.155963in}}%
\pgfpathlineto{\pgfqpoint{0.545830in}{1.173062in}}%
\pgfpathlineto{\pgfqpoint{0.536599in}{1.173553in}}%
\pgfpathlineto{\pgfqpoint{0.527325in}{1.174301in}}%
\pgfpathclose%
\pgfusepath{fill}%
\end{pgfscope}%
\begin{pgfscope}%
\pgfpathrectangle{\pgfqpoint{0.407592in}{0.222904in}}{\pgfqpoint{1.411163in}{1.411163in}}%
\pgfusepath{clip}%
\pgfsetbuttcap%
\pgfsetroundjoin%
\definecolor{currentfill}{rgb}{0.909804,0.698039,0.423529}%
\pgfsetfillcolor{currentfill}%
\pgfsetfillopacity{0.120000}%
\pgfsetlinewidth{0.000000pt}%
\definecolor{currentstroke}{rgb}{0.000000,0.000000,0.000000}%
\pgfsetstrokecolor{currentstroke}%
\pgfsetdash{}{0pt}%
\pgfpathmoveto{\pgfqpoint{1.582646in}{1.120508in}}%
\pgfpathlineto{\pgfqpoint{1.600098in}{1.129246in}}%
\pgfpathlineto{\pgfqpoint{1.617626in}{1.138290in}}%
\pgfpathlineto{\pgfqpoint{1.625174in}{1.148780in}}%
\pgfpathlineto{\pgfqpoint{1.632710in}{1.159498in}}%
\pgfpathlineto{\pgfqpoint{1.615204in}{1.150708in}}%
\pgfpathlineto{\pgfqpoint{1.597773in}{1.142228in}}%
\pgfpathlineto{\pgfqpoint{1.590216in}{1.131252in}}%
\pgfpathlineto{\pgfqpoint{1.582646in}{1.120508in}}%
\pgfpathclose%
\pgfusepath{fill}%
\end{pgfscope}%
\begin{pgfscope}%
\pgfpathrectangle{\pgfqpoint{0.407592in}{0.222904in}}{\pgfqpoint{1.411163in}{1.411163in}}%
\pgfusepath{clip}%
\pgfsetbuttcap%
\pgfsetroundjoin%
\definecolor{currentfill}{rgb}{0.909804,0.698039,0.423529}%
\pgfsetfillcolor{currentfill}%
\pgfsetfillopacity{0.120000}%
\pgfsetlinewidth{0.000000pt}%
\definecolor{currentstroke}{rgb}{0.000000,0.000000,0.000000}%
\pgfsetstrokecolor{currentstroke}%
\pgfsetdash{}{0pt}%
\pgfpathmoveto{\pgfqpoint{1.006825in}{0.864545in}}%
\pgfpathlineto{\pgfqpoint{1.023262in}{0.860118in}}%
\pgfpathlineto{\pgfqpoint{1.039735in}{0.856663in}}%
\pgfpathlineto{\pgfqpoint{1.048062in}{0.857423in}}%
\pgfpathlineto{\pgfqpoint{1.056357in}{0.858685in}}%
\pgfpathlineto{\pgfqpoint{1.039922in}{0.862105in}}%
\pgfpathlineto{\pgfqpoint{1.023524in}{0.866521in}}%
\pgfpathlineto{\pgfqpoint{1.015191in}{0.865283in}}%
\pgfpathlineto{\pgfqpoint{1.006825in}{0.864545in}}%
\pgfpathclose%
\pgfusepath{fill}%
\end{pgfscope}%
\begin{pgfscope}%
\pgfpathrectangle{\pgfqpoint{0.407592in}{0.222904in}}{\pgfqpoint{1.411163in}{1.411163in}}%
\pgfusepath{clip}%
\pgfsetbuttcap%
\pgfsetroundjoin%
\definecolor{currentfill}{rgb}{0.909804,0.698039,0.423529}%
\pgfsetfillcolor{currentfill}%
\pgfsetfillopacity{0.120000}%
\pgfsetlinewidth{0.000000pt}%
\definecolor{currentstroke}{rgb}{0.000000,0.000000,0.000000}%
\pgfsetstrokecolor{currentstroke}%
\pgfsetdash{}{0pt}%
\pgfpathmoveto{\pgfqpoint{1.697887in}{1.221889in}}%
\pgfpathlineto{\pgfqpoint{1.715576in}{1.231019in}}%
\pgfpathlineto{\pgfqpoint{1.733342in}{1.240411in}}%
\pgfpathlineto{\pgfqpoint{1.740775in}{1.251626in}}%
\pgfpathlineto{\pgfqpoint{1.748197in}{1.263018in}}%
\pgfpathlineto{\pgfqpoint{1.730454in}{1.253875in}}%
\pgfpathlineto{\pgfqpoint{1.712787in}{1.245000in}}%
\pgfpathlineto{\pgfqpoint{1.705342in}{1.233356in}}%
\pgfpathlineto{\pgfqpoint{1.697887in}{1.221889in}}%
\pgfpathclose%
\pgfusepath{fill}%
\end{pgfscope}%
\begin{pgfscope}%
\pgfpathrectangle{\pgfqpoint{0.407592in}{0.222904in}}{\pgfqpoint{1.411163in}{1.411163in}}%
\pgfusepath{clip}%
\pgfsetbuttcap%
\pgfsetroundjoin%
\definecolor{currentfill}{rgb}{0.909804,0.698039,0.423529}%
\pgfsetfillcolor{currentfill}%
\pgfsetfillopacity{0.120000}%
\pgfsetlinewidth{0.000000pt}%
\definecolor{currentstroke}{rgb}{0.000000,0.000000,0.000000}%
\pgfsetstrokecolor{currentstroke}%
\pgfsetdash{}{0pt}%
\pgfpathmoveto{\pgfqpoint{1.352782in}{0.946044in}}%
\pgfpathlineto{\pgfqpoint{1.369780in}{0.952930in}}%
\pgfpathlineto{\pgfqpoint{1.386849in}{0.960301in}}%
\pgfpathlineto{\pgfqpoint{1.394649in}{0.967980in}}%
\pgfpathlineto{\pgfqpoint{1.402433in}{0.976048in}}%
\pgfpathlineto{\pgfqpoint{1.385386in}{0.968832in}}%
\pgfpathlineto{\pgfqpoint{1.368408in}{0.962103in}}%
\pgfpathlineto{\pgfqpoint{1.360603in}{0.953869in}}%
\pgfpathlineto{\pgfqpoint{1.352782in}{0.946044in}}%
\pgfpathclose%
\pgfusepath{fill}%
\end{pgfscope}%
\begin{pgfscope}%
\pgfpathrectangle{\pgfqpoint{0.407592in}{0.222904in}}{\pgfqpoint{1.411163in}{1.411163in}}%
\pgfusepath{clip}%
\pgfsetbuttcap%
\pgfsetroundjoin%
\definecolor{currentfill}{rgb}{0.909804,0.698039,0.423529}%
\pgfsetfillcolor{currentfill}%
\pgfsetfillopacity{0.120000}%
\pgfsetlinewidth{0.000000pt}%
\definecolor{currentstroke}{rgb}{0.000000,0.000000,0.000000}%
\pgfsetstrokecolor{currentstroke}%
\pgfsetdash{}{0pt}%
\pgfpathmoveto{\pgfqpoint{0.675900in}{1.043512in}}%
\pgfpathlineto{\pgfqpoint{0.692205in}{1.028748in}}%
\pgfpathlineto{\pgfqpoint{0.708521in}{1.014388in}}%
\pgfpathlineto{\pgfqpoint{0.717494in}{1.013094in}}%
\pgfpathlineto{\pgfqpoint{0.726426in}{1.012114in}}%
\pgfpathlineto{\pgfqpoint{0.710153in}{1.026711in}}%
\pgfpathlineto{\pgfqpoint{0.693891in}{1.041708in}}%
\pgfpathlineto{\pgfqpoint{0.684916in}{1.042457in}}%
\pgfpathlineto{\pgfqpoint{0.675900in}{1.043512in}}%
\pgfpathclose%
\pgfusepath{fill}%
\end{pgfscope}%
\begin{pgfscope}%
\pgfpathrectangle{\pgfqpoint{0.407592in}{0.222904in}}{\pgfqpoint{1.411163in}{1.411163in}}%
\pgfusepath{clip}%
\pgfsetbuttcap%
\pgfsetroundjoin%
\definecolor{currentfill}{rgb}{0.909804,0.698039,0.423529}%
\pgfsetfillcolor{currentfill}%
\pgfsetfillopacity{0.120000}%
\pgfsetlinewidth{0.000000pt}%
\definecolor{currentstroke}{rgb}{0.000000,0.000000,0.000000}%
\pgfsetstrokecolor{currentstroke}%
\pgfsetdash{}{0pt}%
\pgfpathmoveto{\pgfqpoint{0.559870in}{1.140715in}}%
\pgfpathlineto{\pgfqpoint{0.576157in}{1.124282in}}%
\pgfpathlineto{\pgfqpoint{0.592453in}{1.108112in}}%
\pgfpathlineto{\pgfqpoint{0.601641in}{1.106948in}}%
\pgfpathlineto{\pgfqpoint{0.610786in}{1.106057in}}%
\pgfpathlineto{\pgfqpoint{0.594533in}{1.122442in}}%
\pgfpathlineto{\pgfqpoint{0.578289in}{1.139084in}}%
\pgfpathlineto{\pgfqpoint{0.569101in}{1.139767in}}%
\pgfpathlineto{\pgfqpoint{0.559870in}{1.140715in}}%
\pgfpathclose%
\pgfusepath{fill}%
\end{pgfscope}%
\begin{pgfscope}%
\pgfpathrectangle{\pgfqpoint{0.407592in}{0.222904in}}{\pgfqpoint{1.411163in}{1.411163in}}%
\pgfusepath{clip}%
\pgfsetbuttcap%
\pgfsetroundjoin%
\definecolor{currentfill}{rgb}{0.909804,0.698039,0.423529}%
\pgfsetfillcolor{currentfill}%
\pgfsetfillopacity{0.120000}%
\pgfsetlinewidth{0.000000pt}%
\definecolor{currentstroke}{rgb}{0.000000,0.000000,0.000000}%
\pgfsetstrokecolor{currentstroke}%
\pgfsetdash{}{0pt}%
\pgfpathmoveto{\pgfqpoint{0.874480in}{0.911207in}}%
\pgfpathlineto{\pgfqpoint{0.890842in}{0.901641in}}%
\pgfpathlineto{\pgfqpoint{0.907226in}{0.892859in}}%
\pgfpathlineto{\pgfqpoint{0.915820in}{0.891941in}}%
\pgfpathlineto{\pgfqpoint{0.924376in}{0.891435in}}%
\pgfpathlineto{\pgfqpoint{0.908036in}{0.900351in}}%
\pgfpathlineto{\pgfqpoint{0.891716in}{0.910071in}}%
\pgfpathlineto{\pgfqpoint{0.883117in}{0.910437in}}%
\pgfpathlineto{\pgfqpoint{0.874480in}{0.911207in}}%
\pgfpathclose%
\pgfusepath{fill}%
\end{pgfscope}%
\begin{pgfscope}%
\pgfpathrectangle{\pgfqpoint{0.407592in}{0.222904in}}{\pgfqpoint{1.411163in}{1.411163in}}%
\pgfusepath{clip}%
\pgfsetbuttcap%
\pgfsetroundjoin%
\definecolor{currentfill}{rgb}{0.909804,0.698039,0.423529}%
\pgfsetfillcolor{currentfill}%
\pgfsetfillopacity{0.120000}%
\pgfsetlinewidth{0.000000pt}%
\definecolor{currentstroke}{rgb}{0.000000,0.000000,0.000000}%
\pgfsetstrokecolor{currentstroke}%
\pgfsetdash{}{0pt}%
\pgfpathmoveto{\pgfqpoint{1.517476in}{1.062539in}}%
\pgfpathlineto{\pgfqpoint{1.534822in}{1.071118in}}%
\pgfpathlineto{\pgfqpoint{1.552242in}{1.080028in}}%
\pgfpathlineto{\pgfqpoint{1.559862in}{1.089754in}}%
\pgfpathlineto{\pgfqpoint{1.567470in}{1.099749in}}%
\pgfpathlineto{\pgfqpoint{1.550071in}{1.091079in}}%
\pgfpathlineto{\pgfqpoint{1.532746in}{1.082745in}}%
\pgfpathlineto{\pgfqpoint{1.525117in}{1.072504in}}%
\pgfpathlineto{\pgfqpoint{1.517476in}{1.062539in}}%
\pgfpathclose%
\pgfusepath{fill}%
\end{pgfscope}%
\begin{pgfscope}%
\pgfpathrectangle{\pgfqpoint{0.407592in}{0.222904in}}{\pgfqpoint{1.411163in}{1.411163in}}%
\pgfusepath{clip}%
\pgfsetbuttcap%
\pgfsetroundjoin%
\definecolor{currentfill}{rgb}{0.909804,0.698039,0.423529}%
\pgfsetfillcolor{currentfill}%
\pgfsetfillopacity{0.120000}%
\pgfsetlinewidth{0.000000pt}%
\definecolor{currentstroke}{rgb}{0.000000,0.000000,0.000000}%
\pgfsetstrokecolor{currentstroke}%
\pgfsetdash{}{0pt}%
\pgfpathmoveto{\pgfqpoint{1.632710in}{1.159498in}}%
\pgfpathlineto{\pgfqpoint{1.650291in}{1.168581in}}%
\pgfpathlineto{\pgfqpoint{1.667949in}{1.177943in}}%
\pgfpathlineto{\pgfqpoint{1.675450in}{1.188627in}}%
\pgfpathlineto{\pgfqpoint{1.682941in}{1.199517in}}%
\pgfpathlineto{\pgfqpoint{1.665305in}{1.190406in}}%
\pgfpathlineto{\pgfqpoint{1.647746in}{1.181579in}}%
\pgfpathlineto{\pgfqpoint{1.640234in}{1.170433in}}%
\pgfpathlineto{\pgfqpoint{1.632710in}{1.159498in}}%
\pgfpathclose%
\pgfusepath{fill}%
\end{pgfscope}%
\begin{pgfscope}%
\pgfpathrectangle{\pgfqpoint{0.407592in}{0.222904in}}{\pgfqpoint{1.411163in}{1.411163in}}%
\pgfusepath{clip}%
\pgfsetbuttcap%
\pgfsetroundjoin%
\definecolor{currentfill}{rgb}{0.909804,0.698039,0.423529}%
\pgfsetfillcolor{currentfill}%
\pgfsetfillopacity{0.120000}%
\pgfsetlinewidth{0.000000pt}%
\definecolor{currentstroke}{rgb}{0.000000,0.000000,0.000000}%
\pgfsetstrokecolor{currentstroke}%
\pgfsetdash{}{0pt}%
\pgfpathmoveto{\pgfqpoint{0.791640in}{0.958388in}}%
\pgfpathlineto{\pgfqpoint{0.807978in}{0.946302in}}%
\pgfpathlineto{\pgfqpoint{0.824331in}{0.934830in}}%
\pgfpathlineto{\pgfqpoint{0.833090in}{0.933502in}}%
\pgfpathlineto{\pgfqpoint{0.841810in}{0.932535in}}%
\pgfpathlineto{\pgfqpoint{0.825501in}{0.944217in}}%
\pgfpathlineto{\pgfqpoint{0.809207in}{0.956520in}}%
\pgfpathlineto{\pgfqpoint{0.800444in}{0.957276in}}%
\pgfpathlineto{\pgfqpoint{0.791640in}{0.958388in}}%
\pgfpathclose%
\pgfusepath{fill}%
\end{pgfscope}%
\begin{pgfscope}%
\pgfpathrectangle{\pgfqpoint{0.407592in}{0.222904in}}{\pgfqpoint{1.411163in}{1.411163in}}%
\pgfusepath{clip}%
\pgfsetbuttcap%
\pgfsetroundjoin%
\definecolor{currentfill}{rgb}{0.909804,0.698039,0.423529}%
\pgfsetfillcolor{currentfill}%
\pgfsetfillopacity{0.120000}%
\pgfsetlinewidth{0.000000pt}%
\definecolor{currentstroke}{rgb}{0.000000,0.000000,0.000000}%
\pgfsetstrokecolor{currentstroke}%
\pgfsetdash{}{0pt}%
\pgfpathmoveto{\pgfqpoint{1.402433in}{0.976048in}}%
\pgfpathlineto{\pgfqpoint{1.419551in}{0.983717in}}%
\pgfpathlineto{\pgfqpoint{1.436741in}{0.991806in}}%
\pgfpathlineto{\pgfqpoint{1.444488in}{1.000073in}}%
\pgfpathlineto{\pgfqpoint{1.452220in}{1.008691in}}%
\pgfpathlineto{\pgfqpoint{1.435051in}{1.000789in}}%
\pgfpathlineto{\pgfqpoint{1.417954in}{0.993309in}}%
\pgfpathlineto{\pgfqpoint{1.410201in}{0.984495in}}%
\pgfpathlineto{\pgfqpoint{1.402433in}{0.976048in}}%
\pgfpathclose%
\pgfusepath{fill}%
\end{pgfscope}%
\begin{pgfscope}%
\pgfpathrectangle{\pgfqpoint{0.407592in}{0.222904in}}{\pgfqpoint{1.411163in}{1.411163in}}%
\pgfusepath{clip}%
\pgfsetbuttcap%
\pgfsetroundjoin%
\definecolor{currentfill}{rgb}{0.909804,0.698039,0.423529}%
\pgfsetfillcolor{currentfill}%
\pgfsetfillopacity{0.120000}%
\pgfsetlinewidth{0.000000pt}%
\definecolor{currentstroke}{rgb}{0.000000,0.000000,0.000000}%
\pgfsetstrokecolor{currentstroke}%
\pgfsetdash{}{0pt}%
\pgfpathmoveto{\pgfqpoint{0.957128in}{0.876160in}}%
\pgfpathlineto{\pgfqpoint{0.973544in}{0.869864in}}%
\pgfpathlineto{\pgfqpoint{0.989989in}{0.864492in}}%
\pgfpathlineto{\pgfqpoint{0.998424in}{0.864288in}}%
\pgfpathlineto{\pgfqpoint{1.006825in}{0.864545in}}%
\pgfpathlineto{\pgfqpoint{0.990421in}{0.869938in}}%
\pgfpathlineto{\pgfqpoint{0.974046in}{0.876284in}}%
\pgfpathlineto{\pgfqpoint{0.965605in}{0.875994in}}%
\pgfpathlineto{\pgfqpoint{0.957128in}{0.876160in}}%
\pgfpathclose%
\pgfusepath{fill}%
\end{pgfscope}%
\begin{pgfscope}%
\pgfpathrectangle{\pgfqpoint{0.407592in}{0.222904in}}{\pgfqpoint{1.411163in}{1.411163in}}%
\pgfusepath{clip}%
\pgfsetbuttcap%
\pgfsetroundjoin%
\definecolor{currentfill}{rgb}{0.909804,0.698039,0.423529}%
\pgfsetfillcolor{currentfill}%
\pgfsetfillopacity{0.120000}%
\pgfsetlinewidth{0.000000pt}%
\definecolor{currentstroke}{rgb}{0.000000,0.000000,0.000000}%
\pgfsetstrokecolor{currentstroke}%
\pgfsetdash{}{0pt}%
\pgfpathmoveto{\pgfqpoint{1.188366in}{0.869556in}}%
\pgfpathlineto{\pgfqpoint{1.205074in}{0.872745in}}%
\pgfpathlineto{\pgfqpoint{1.221842in}{0.876697in}}%
\pgfpathlineto{\pgfqpoint{1.229870in}{0.880878in}}%
\pgfpathlineto{\pgfqpoint{1.237874in}{0.885564in}}%
\pgfpathlineto{\pgfqpoint{1.221133in}{0.881569in}}%
\pgfpathlineto{\pgfqpoint{1.204453in}{0.878338in}}%
\pgfpathlineto{\pgfqpoint{1.196422in}{0.873682in}}%
\pgfpathlineto{\pgfqpoint{1.188366in}{0.869556in}}%
\pgfpathclose%
\pgfusepath{fill}%
\end{pgfscope}%
\begin{pgfscope}%
\pgfpathrectangle{\pgfqpoint{0.407592in}{0.222904in}}{\pgfqpoint{1.411163in}{1.411163in}}%
\pgfusepath{clip}%
\pgfsetbuttcap%
\pgfsetroundjoin%
\definecolor{currentfill}{rgb}{0.909804,0.698039,0.423529}%
\pgfsetfillcolor{currentfill}%
\pgfsetfillopacity{0.120000}%
\pgfsetlinewidth{0.000000pt}%
\definecolor{currentstroke}{rgb}{0.000000,0.000000,0.000000}%
\pgfsetstrokecolor{currentstroke}%
\pgfsetdash{}{0pt}%
\pgfpathmoveto{\pgfqpoint{1.138873in}{0.859153in}}%
\pgfpathlineto{\pgfqpoint{1.155500in}{0.860586in}}%
\pgfpathlineto{\pgfqpoint{1.172180in}{0.862877in}}%
\pgfpathlineto{\pgfqpoint{1.180286in}{0.865957in}}%
\pgfpathlineto{\pgfqpoint{1.188366in}{0.869556in}}%
\pgfpathlineto{\pgfqpoint{1.171716in}{0.867182in}}%
\pgfpathlineto{\pgfqpoint{1.155119in}{0.865670in}}%
\pgfpathlineto{\pgfqpoint{1.147010in}{0.862142in}}%
\pgfpathlineto{\pgfqpoint{1.138873in}{0.859153in}}%
\pgfpathclose%
\pgfusepath{fill}%
\end{pgfscope}%
\begin{pgfscope}%
\pgfpathrectangle{\pgfqpoint{0.407592in}{0.222904in}}{\pgfqpoint{1.411163in}{1.411163in}}%
\pgfusepath{clip}%
\pgfsetbuttcap%
\pgfsetroundjoin%
\definecolor{currentfill}{rgb}{0.909804,0.698039,0.423529}%
\pgfsetfillcolor{currentfill}%
\pgfsetfillopacity{0.120000}%
\pgfsetlinewidth{0.000000pt}%
\definecolor{currentstroke}{rgb}{0.000000,0.000000,0.000000}%
\pgfsetstrokecolor{currentstroke}%
\pgfsetdash{}{0pt}%
\pgfpathmoveto{\pgfqpoint{1.237874in}{0.885564in}}%
\pgfpathlineto{\pgfqpoint{1.254677in}{0.890272in}}%
\pgfpathlineto{\pgfqpoint{1.271543in}{0.895643in}}%
\pgfpathlineto{\pgfqpoint{1.279501in}{0.900838in}}%
\pgfpathlineto{\pgfqpoint{1.287437in}{0.906512in}}%
\pgfpathlineto{\pgfqpoint{1.270595in}{0.901152in}}%
\pgfpathlineto{\pgfqpoint{1.253817in}{0.896454in}}%
\pgfpathlineto{\pgfqpoint{1.245856in}{0.890757in}}%
\pgfpathlineto{\pgfqpoint{1.237874in}{0.885564in}}%
\pgfpathclose%
\pgfusepath{fill}%
\end{pgfscope}%
\begin{pgfscope}%
\pgfpathrectangle{\pgfqpoint{0.407592in}{0.222904in}}{\pgfqpoint{1.411163in}{1.411163in}}%
\pgfusepath{clip}%
\pgfsetbuttcap%
\pgfsetroundjoin%
\definecolor{currentfill}{rgb}{0.909804,0.698039,0.423529}%
\pgfsetfillcolor{currentfill}%
\pgfsetfillopacity{0.120000}%
\pgfsetlinewidth{0.000000pt}%
\definecolor{currentstroke}{rgb}{0.000000,0.000000,0.000000}%
\pgfsetstrokecolor{currentstroke}%
\pgfsetdash{}{0pt}%
\pgfpathmoveto{\pgfqpoint{0.592453in}{1.108112in}}%
\pgfpathlineto{\pgfqpoint{0.608759in}{1.092227in}}%
\pgfpathlineto{\pgfqpoint{0.625076in}{1.076648in}}%
\pgfpathlineto{\pgfqpoint{0.634220in}{1.075253in}}%
\pgfpathlineto{\pgfqpoint{0.643322in}{1.074138in}}%
\pgfpathlineto{\pgfqpoint{0.627049in}{1.089948in}}%
\pgfpathlineto{\pgfqpoint{0.610786in}{1.106057in}}%
\pgfpathlineto{\pgfqpoint{0.601641in}{1.106948in}}%
\pgfpathlineto{\pgfqpoint{0.592453in}{1.108112in}}%
\pgfpathclose%
\pgfusepath{fill}%
\end{pgfscope}%
\begin{pgfscope}%
\pgfpathrectangle{\pgfqpoint{0.407592in}{0.222904in}}{\pgfqpoint{1.411163in}{1.411163in}}%
\pgfusepath{clip}%
\pgfsetbuttcap%
\pgfsetroundjoin%
\definecolor{currentfill}{rgb}{0.909804,0.698039,0.423529}%
\pgfsetfillcolor{currentfill}%
\pgfsetfillopacity{0.120000}%
\pgfsetlinewidth{0.000000pt}%
\definecolor{currentstroke}{rgb}{0.000000,0.000000,0.000000}%
\pgfsetstrokecolor{currentstroke}%
\pgfsetdash{}{0pt}%
\pgfpathmoveto{\pgfqpoint{0.708521in}{1.014388in}}%
\pgfpathlineto{\pgfqpoint{0.724850in}{1.000462in}}%
\pgfpathlineto{\pgfqpoint{0.741190in}{0.987002in}}%
\pgfpathlineto{\pgfqpoint{0.750119in}{0.985465in}}%
\pgfpathlineto{\pgfqpoint{0.759007in}{0.984248in}}%
\pgfpathlineto{\pgfqpoint{0.742710in}{0.997948in}}%
\pgfpathlineto{\pgfqpoint{0.726426in}{1.012114in}}%
\pgfpathlineto{\pgfqpoint{0.717494in}{1.013094in}}%
\pgfpathlineto{\pgfqpoint{0.708521in}{1.014388in}}%
\pgfpathclose%
\pgfusepath{fill}%
\end{pgfscope}%
\begin{pgfscope}%
\pgfpathrectangle{\pgfqpoint{0.407592in}{0.222904in}}{\pgfqpoint{1.411163in}{1.411163in}}%
\pgfusepath{clip}%
\pgfsetbuttcap%
\pgfsetroundjoin%
\definecolor{currentfill}{rgb}{0.909804,0.698039,0.423529}%
\pgfsetfillcolor{currentfill}%
\pgfsetfillopacity{0.120000}%
\pgfsetlinewidth{0.000000pt}%
\definecolor{currentstroke}{rgb}{0.000000,0.000000,0.000000}%
\pgfsetstrokecolor{currentstroke}%
\pgfsetdash{}{0pt}%
\pgfpathmoveto{\pgfqpoint{1.089346in}{0.854809in}}%
\pgfpathlineto{\pgfqpoint{1.105907in}{0.854317in}}%
\pgfpathlineto{\pgfqpoint{1.122515in}{0.854750in}}%
\pgfpathlineto{\pgfqpoint{1.130709in}{0.856693in}}%
\pgfpathlineto{\pgfqpoint{1.138873in}{0.859153in}}%
\pgfpathlineto{\pgfqpoint{1.122298in}{0.858620in}}%
\pgfpathlineto{\pgfqpoint{1.105771in}{0.859027in}}%
\pgfpathlineto{\pgfqpoint{1.097574in}{0.856654in}}%
\pgfpathlineto{\pgfqpoint{1.089346in}{0.854809in}}%
\pgfpathclose%
\pgfusepath{fill}%
\end{pgfscope}%
\begin{pgfscope}%
\pgfpathrectangle{\pgfqpoint{0.407592in}{0.222904in}}{\pgfqpoint{1.411163in}{1.411163in}}%
\pgfusepath{clip}%
\pgfsetbuttcap%
\pgfsetroundjoin%
\definecolor{currentfill}{rgb}{0.909804,0.698039,0.423529}%
\pgfsetfillcolor{currentfill}%
\pgfsetfillopacity{0.120000}%
\pgfsetlinewidth{0.000000pt}%
\definecolor{currentstroke}{rgb}{0.000000,0.000000,0.000000}%
\pgfsetstrokecolor{currentstroke}%
\pgfsetdash{}{0pt}%
\pgfpathmoveto{\pgfqpoint{1.567470in}{1.099749in}}%
\pgfpathlineto{\pgfqpoint{1.584943in}{1.108737in}}%
\pgfpathlineto{\pgfqpoint{1.602493in}{1.118023in}}%
\pgfpathlineto{\pgfqpoint{1.610065in}{1.128035in}}%
\pgfpathlineto{\pgfqpoint{1.617626in}{1.138290in}}%
\pgfpathlineto{\pgfqpoint{1.600098in}{1.129246in}}%
\pgfpathlineto{\pgfqpoint{1.582646in}{1.120508in}}%
\pgfpathlineto{\pgfqpoint{1.575064in}{1.110003in}}%
\pgfpathlineto{\pgfqpoint{1.567470in}{1.099749in}}%
\pgfpathclose%
\pgfusepath{fill}%
\end{pgfscope}%
\begin{pgfscope}%
\pgfpathrectangle{\pgfqpoint{0.407592in}{0.222904in}}{\pgfqpoint{1.411163in}{1.411163in}}%
\pgfusepath{clip}%
\pgfsetbuttcap%
\pgfsetroundjoin%
\definecolor{currentfill}{rgb}{0.909804,0.698039,0.423529}%
\pgfsetfillcolor{currentfill}%
\pgfsetfillopacity{0.120000}%
\pgfsetlinewidth{0.000000pt}%
\definecolor{currentstroke}{rgb}{0.000000,0.000000,0.000000}%
\pgfsetstrokecolor{currentstroke}%
\pgfsetdash{}{0pt}%
\pgfpathmoveto{\pgfqpoint{1.748197in}{1.263018in}}%
\pgfpathlineto{\pgfqpoint{1.766017in}{1.272415in}}%
\pgfpathlineto{\pgfqpoint{1.783915in}{1.282058in}}%
\pgfpathlineto{\pgfqpoint{1.791303in}{1.293374in}}%
\pgfpathlineto{\pgfqpoint{1.773417in}{1.283853in}}%
\pgfpathlineto{\pgfqpoint{1.755608in}{1.274579in}}%
\pgfusepath{fill}%
\end{pgfscope}%
\begin{pgfscope}%
\pgfpathrectangle{\pgfqpoint{0.407592in}{0.222904in}}{\pgfqpoint{1.411163in}{1.411163in}}%
\pgfusepath{clip}%
\pgfsetbuttcap%
\pgfsetroundjoin%
\definecolor{currentfill}{rgb}{0.909804,0.698039,0.423529}%
\pgfsetfillcolor{currentfill}%
\pgfsetfillopacity{0.120000}%
\pgfsetlinewidth{0.000000pt}%
\definecolor{currentstroke}{rgb}{0.000000,0.000000,0.000000}%
\pgfsetstrokecolor{currentstroke}%
\pgfsetdash{}{0pt}%
\pgfpathmoveto{\pgfqpoint{1.452220in}{1.008691in}}%
\pgfpathlineto{\pgfqpoint{1.469461in}{1.016988in}}%
\pgfpathlineto{\pgfqpoint{1.486776in}{1.025652in}}%
\pgfpathlineto{\pgfqpoint{1.494472in}{1.034408in}}%
\pgfpathlineto{\pgfqpoint{1.502153in}{1.043481in}}%
\pgfpathlineto{\pgfqpoint{1.484860in}{1.035024in}}%
\pgfpathlineto{\pgfqpoint{1.467640in}{1.026940in}}%
\pgfpathlineto{\pgfqpoint{1.459937in}{1.017650in}}%
\pgfpathlineto{\pgfqpoint{1.452220in}{1.008691in}}%
\pgfpathclose%
\pgfusepath{fill}%
\end{pgfscope}%
\begin{pgfscope}%
\pgfpathrectangle{\pgfqpoint{0.407592in}{0.222904in}}{\pgfqpoint{1.411163in}{1.411163in}}%
\pgfusepath{clip}%
\pgfsetbuttcap%
\pgfsetroundjoin%
\definecolor{currentfill}{rgb}{0.909804,0.698039,0.423529}%
\pgfsetfillcolor{currentfill}%
\pgfsetfillopacity{0.120000}%
\pgfsetlinewidth{0.000000pt}%
\definecolor{currentstroke}{rgb}{0.000000,0.000000,0.000000}%
\pgfsetstrokecolor{currentstroke}%
\pgfsetdash{}{0pt}%
\pgfpathmoveto{\pgfqpoint{1.287437in}{0.906512in}}%
\pgfpathlineto{\pgfqpoint{1.304345in}{0.912484in}}%
\pgfpathlineto{\pgfqpoint{1.321320in}{0.919025in}}%
\pgfpathlineto{\pgfqpoint{1.329213in}{0.925121in}}%
\pgfpathlineto{\pgfqpoint{1.337087in}{0.931661in}}%
\pgfpathlineto{\pgfqpoint{1.320135in}{0.925186in}}%
\pgfpathlineto{\pgfqpoint{1.303250in}{0.919278in}}%
\pgfpathlineto{\pgfqpoint{1.295353in}{0.912660in}}%
\pgfpathlineto{\pgfqpoint{1.287437in}{0.906512in}}%
\pgfpathclose%
\pgfusepath{fill}%
\end{pgfscope}%
\begin{pgfscope}%
\pgfpathrectangle{\pgfqpoint{0.407592in}{0.222904in}}{\pgfqpoint{1.411163in}{1.411163in}}%
\pgfusepath{clip}%
\pgfsetbuttcap%
\pgfsetroundjoin%
\definecolor{currentfill}{rgb}{0.909804,0.698039,0.423529}%
\pgfsetfillcolor{currentfill}%
\pgfsetfillopacity{0.120000}%
\pgfsetlinewidth{0.000000pt}%
\definecolor{currentstroke}{rgb}{0.000000,0.000000,0.000000}%
\pgfsetstrokecolor{currentstroke}%
\pgfsetdash{}{0pt}%
\pgfpathmoveto{\pgfqpoint{0.476054in}{1.210594in}}%
\pgfpathlineto{\pgfqpoint{0.492346in}{1.193469in}}%
\pgfpathlineto{\pgfqpoint{0.508647in}{1.176548in}}%
\pgfpathlineto{\pgfqpoint{0.518008in}{1.175302in}}%
\pgfpathlineto{\pgfqpoint{0.527325in}{1.174301in}}%
\pgfpathlineto{\pgfqpoint{0.511067in}{1.191420in}}%
\pgfpathlineto{\pgfqpoint{0.494818in}{1.208735in}}%
\pgfpathlineto{\pgfqpoint{0.485458in}{1.209544in}}%
\pgfpathlineto{\pgfqpoint{0.476054in}{1.210594in}}%
\pgfpathclose%
\pgfusepath{fill}%
\end{pgfscope}%
\begin{pgfscope}%
\pgfpathrectangle{\pgfqpoint{0.407592in}{0.222904in}}{\pgfqpoint{1.411163in}{1.411163in}}%
\pgfusepath{clip}%
\pgfsetbuttcap%
\pgfsetroundjoin%
\definecolor{currentfill}{rgb}{0.909804,0.698039,0.423529}%
\pgfsetfillcolor{currentfill}%
\pgfsetfillopacity{0.120000}%
\pgfsetlinewidth{0.000000pt}%
\definecolor{currentstroke}{rgb}{0.000000,0.000000,0.000000}%
\pgfsetstrokecolor{currentstroke}%
\pgfsetdash{}{0pt}%
\pgfpathmoveto{\pgfqpoint{1.682941in}{1.199517in}}%
\pgfpathlineto{\pgfqpoint{1.700652in}{1.208899in}}%
\pgfpathlineto{\pgfqpoint{1.718441in}{1.218537in}}%
\pgfpathlineto{\pgfqpoint{1.725897in}{1.229379in}}%
\pgfpathlineto{\pgfqpoint{1.733342in}{1.240411in}}%
\pgfpathlineto{\pgfqpoint{1.715576in}{1.231019in}}%
\pgfpathlineto{\pgfqpoint{1.697887in}{1.221889in}}%
\pgfpathlineto{\pgfqpoint{1.690419in}{1.210607in}}%
\pgfpathlineto{\pgfqpoint{1.682941in}{1.199517in}}%
\pgfpathclose%
\pgfusepath{fill}%
\end{pgfscope}%
\begin{pgfscope}%
\pgfpathrectangle{\pgfqpoint{0.407592in}{0.222904in}}{\pgfqpoint{1.411163in}{1.411163in}}%
\pgfusepath{clip}%
\pgfsetbuttcap%
\pgfsetroundjoin%
\definecolor{currentfill}{rgb}{0.909804,0.698039,0.423529}%
\pgfsetfillcolor{currentfill}%
\pgfsetfillopacity{0.120000}%
\pgfsetlinewidth{0.000000pt}%
\definecolor{currentstroke}{rgb}{0.000000,0.000000,0.000000}%
\pgfsetstrokecolor{currentstroke}%
\pgfsetdash{}{0pt}%
\pgfpathmoveto{\pgfqpoint{1.039735in}{0.856663in}}%
\pgfpathlineto{\pgfqpoint{1.056246in}{0.854176in}}%
\pgfpathlineto{\pgfqpoint{1.072799in}{0.852644in}}%
\pgfpathlineto{\pgfqpoint{1.081088in}{0.853479in}}%
\pgfpathlineto{\pgfqpoint{1.089346in}{0.854809in}}%
\pgfpathlineto{\pgfqpoint{1.072831in}{0.856258in}}%
\pgfpathlineto{\pgfqpoint{1.056357in}{0.858685in}}%
\pgfpathlineto{\pgfqpoint{1.048062in}{0.857423in}}%
\pgfpathlineto{\pgfqpoint{1.039735in}{0.856663in}}%
\pgfpathclose%
\pgfusepath{fill}%
\end{pgfscope}%
\begin{pgfscope}%
\pgfpathrectangle{\pgfqpoint{0.407592in}{0.222904in}}{\pgfqpoint{1.411163in}{1.411163in}}%
\pgfusepath{clip}%
\pgfsetbuttcap%
\pgfsetroundjoin%
\definecolor{currentfill}{rgb}{0.909804,0.698039,0.423529}%
\pgfsetfillcolor{currentfill}%
\pgfsetfillopacity{0.120000}%
\pgfsetlinewidth{0.000000pt}%
\definecolor{currentstroke}{rgb}{0.000000,0.000000,0.000000}%
\pgfsetstrokecolor{currentstroke}%
\pgfsetdash{}{0pt}%
\pgfpathmoveto{\pgfqpoint{0.824331in}{0.934830in}}%
\pgfpathlineto{\pgfqpoint{0.840701in}{0.924011in}}%
\pgfpathlineto{\pgfqpoint{0.857088in}{0.913884in}}%
\pgfpathlineto{\pgfqpoint{0.865804in}{0.912361in}}%
\pgfpathlineto{\pgfqpoint{0.874480in}{0.911207in}}%
\pgfpathlineto{\pgfqpoint{0.858136in}{0.921518in}}%
\pgfpathlineto{\pgfqpoint{0.841810in}{0.932535in}}%
\pgfpathlineto{\pgfqpoint{0.833090in}{0.933502in}}%
\pgfpathlineto{\pgfqpoint{0.824331in}{0.934830in}}%
\pgfpathclose%
\pgfusepath{fill}%
\end{pgfscope}%
\begin{pgfscope}%
\pgfpathrectangle{\pgfqpoint{0.407592in}{0.222904in}}{\pgfqpoint{1.411163in}{1.411163in}}%
\pgfusepath{clip}%
\pgfsetbuttcap%
\pgfsetroundjoin%
\definecolor{currentfill}{rgb}{0.909804,0.698039,0.423529}%
\pgfsetfillcolor{currentfill}%
\pgfsetfillopacity{0.120000}%
\pgfsetlinewidth{0.000000pt}%
\definecolor{currentstroke}{rgb}{0.000000,0.000000,0.000000}%
\pgfsetstrokecolor{currentstroke}%
\pgfsetdash{}{0pt}%
\pgfpathmoveto{\pgfqpoint{0.907226in}{0.892859in}}%
\pgfpathlineto{\pgfqpoint{0.923633in}{0.884895in}}%
\pgfpathlineto{\pgfqpoint{0.940064in}{0.877782in}}%
\pgfpathlineto{\pgfqpoint{0.948615in}{0.876763in}}%
\pgfpathlineto{\pgfqpoint{0.957128in}{0.876160in}}%
\pgfpathlineto{\pgfqpoint{0.940740in}{0.883360in}}%
\pgfpathlineto{\pgfqpoint{0.924376in}{0.891435in}}%
\pgfpathlineto{\pgfqpoint{0.915820in}{0.891941in}}%
\pgfpathlineto{\pgfqpoint{0.907226in}{0.892859in}}%
\pgfpathclose%
\pgfusepath{fill}%
\end{pgfscope}%
\begin{pgfscope}%
\pgfpathrectangle{\pgfqpoint{0.407592in}{0.222904in}}{\pgfqpoint{1.411163in}{1.411163in}}%
\pgfusepath{clip}%
\pgfsetbuttcap%
\pgfsetroundjoin%
\definecolor{currentfill}{rgb}{0.909804,0.698039,0.423529}%
\pgfsetfillcolor{currentfill}%
\pgfsetfillopacity{0.120000}%
\pgfsetlinewidth{0.000000pt}%
\definecolor{currentstroke}{rgb}{0.000000,0.000000,0.000000}%
\pgfsetstrokecolor{currentstroke}%
\pgfsetdash{}{0pt}%
\pgfpathmoveto{\pgfqpoint{1.337087in}{0.931661in}}%
\pgfpathlineto{\pgfqpoint{1.354108in}{0.938661in}}%
\pgfpathlineto{\pgfqpoint{1.371199in}{0.946146in}}%
\pgfpathlineto{\pgfqpoint{1.379032in}{0.953020in}}%
\pgfpathlineto{\pgfqpoint{1.386849in}{0.960301in}}%
\pgfpathlineto{\pgfqpoint{1.369780in}{0.952930in}}%
\pgfpathlineto{\pgfqpoint{1.352782in}{0.946044in}}%
\pgfpathlineto{\pgfqpoint{1.344943in}{0.938638in}}%
\pgfpathlineto{\pgfqpoint{1.337087in}{0.931661in}}%
\pgfpathclose%
\pgfusepath{fill}%
\end{pgfscope}%
\begin{pgfscope}%
\pgfpathrectangle{\pgfqpoint{0.407592in}{0.222904in}}{\pgfqpoint{1.411163in}{1.411163in}}%
\pgfusepath{clip}%
\pgfsetbuttcap%
\pgfsetroundjoin%
\definecolor{currentfill}{rgb}{0.909804,0.698039,0.423529}%
\pgfsetfillcolor{currentfill}%
\pgfsetfillopacity{0.120000}%
\pgfsetlinewidth{0.000000pt}%
\definecolor{currentstroke}{rgb}{0.000000,0.000000,0.000000}%
\pgfsetstrokecolor{currentstroke}%
\pgfsetdash{}{0pt}%
\pgfpathmoveto{\pgfqpoint{0.625076in}{1.076648in}}%
\pgfpathlineto{\pgfqpoint{0.641403in}{1.061397in}}%
\pgfpathlineto{\pgfqpoint{0.657740in}{1.046500in}}%
\pgfpathlineto{\pgfqpoint{0.666841in}{1.044863in}}%
\pgfpathlineto{\pgfqpoint{0.675900in}{1.043512in}}%
\pgfpathlineto{\pgfqpoint{0.659606in}{1.058651in}}%
\pgfpathlineto{\pgfqpoint{0.643322in}{1.074138in}}%
\pgfpathlineto{\pgfqpoint{0.634220in}{1.075253in}}%
\pgfpathlineto{\pgfqpoint{0.625076in}{1.076648in}}%
\pgfpathclose%
\pgfusepath{fill}%
\end{pgfscope}%
\begin{pgfscope}%
\pgfpathrectangle{\pgfqpoint{0.407592in}{0.222904in}}{\pgfqpoint{1.411163in}{1.411163in}}%
\pgfusepath{clip}%
\pgfsetbuttcap%
\pgfsetroundjoin%
\definecolor{currentfill}{rgb}{0.909804,0.698039,0.423529}%
\pgfsetfillcolor{currentfill}%
\pgfsetfillopacity{0.120000}%
\pgfsetlinewidth{0.000000pt}%
\definecolor{currentstroke}{rgb}{0.000000,0.000000,0.000000}%
\pgfsetstrokecolor{currentstroke}%
\pgfsetdash{}{0pt}%
\pgfpathmoveto{\pgfqpoint{0.741190in}{0.987002in}}%
\pgfpathlineto{\pgfqpoint{0.757544in}{0.974042in}}%
\pgfpathlineto{\pgfqpoint{0.773911in}{0.961616in}}%
\pgfpathlineto{\pgfqpoint{0.782796in}{0.959839in}}%
\pgfpathlineto{\pgfqpoint{0.791640in}{0.958388in}}%
\pgfpathlineto{\pgfqpoint{0.775317in}{0.971049in}}%
\pgfpathlineto{\pgfqpoint{0.759007in}{0.984248in}}%
\pgfpathlineto{\pgfqpoint{0.750119in}{0.985465in}}%
\pgfpathlineto{\pgfqpoint{0.741190in}{0.987002in}}%
\pgfpathclose%
\pgfusepath{fill}%
\end{pgfscope}%
\begin{pgfscope}%
\pgfpathrectangle{\pgfqpoint{0.407592in}{0.222904in}}{\pgfqpoint{1.411163in}{1.411163in}}%
\pgfusepath{clip}%
\pgfsetbuttcap%
\pgfsetroundjoin%
\definecolor{currentfill}{rgb}{0.909804,0.698039,0.423529}%
\pgfsetfillcolor{currentfill}%
\pgfsetfillopacity{0.120000}%
\pgfsetlinewidth{0.000000pt}%
\definecolor{currentstroke}{rgb}{0.000000,0.000000,0.000000}%
\pgfsetstrokecolor{currentstroke}%
\pgfsetdash{}{0pt}%
\pgfpathmoveto{\pgfqpoint{0.508647in}{1.176548in}}%
\pgfpathlineto{\pgfqpoint{0.524958in}{1.159848in}}%
\pgfpathlineto{\pgfqpoint{0.541278in}{1.143383in}}%
\pgfpathlineto{\pgfqpoint{0.550596in}{1.141923in}}%
\pgfpathlineto{\pgfqpoint{0.559870in}{1.140715in}}%
\pgfpathlineto{\pgfqpoint{0.543593in}{1.157394in}}%
\pgfpathlineto{\pgfqpoint{0.527325in}{1.174301in}}%
\pgfpathlineto{\pgfqpoint{0.518008in}{1.175302in}}%
\pgfpathlineto{\pgfqpoint{0.508647in}{1.176548in}}%
\pgfpathclose%
\pgfusepath{fill}%
\end{pgfscope}%
\begin{pgfscope}%
\pgfpathrectangle{\pgfqpoint{0.407592in}{0.222904in}}{\pgfqpoint{1.411163in}{1.411163in}}%
\pgfusepath{clip}%
\pgfsetbuttcap%
\pgfsetroundjoin%
\definecolor{currentfill}{rgb}{0.909804,0.698039,0.423529}%
\pgfsetfillcolor{currentfill}%
\pgfsetfillopacity{0.120000}%
\pgfsetlinewidth{0.000000pt}%
\definecolor{currentstroke}{rgb}{0.000000,0.000000,0.000000}%
\pgfsetstrokecolor{currentstroke}%
\pgfsetdash{}{0pt}%
\pgfpathmoveto{\pgfqpoint{1.502153in}{1.043481in}}%
\pgfpathlineto{\pgfqpoint{1.519521in}{1.052286in}}%
\pgfpathlineto{\pgfqpoint{1.536963in}{1.061417in}}%
\pgfpathlineto{\pgfqpoint{1.544609in}{1.070579in}}%
\pgfpathlineto{\pgfqpoint{1.552242in}{1.080028in}}%
\pgfpathlineto{\pgfqpoint{1.534822in}{1.071118in}}%
\pgfpathlineto{\pgfqpoint{1.517476in}{1.062539in}}%
\pgfpathlineto{\pgfqpoint{1.509822in}{1.052861in}}%
\pgfpathlineto{\pgfqpoint{1.502153in}{1.043481in}}%
\pgfpathclose%
\pgfusepath{fill}%
\end{pgfscope}%
\begin{pgfscope}%
\pgfpathrectangle{\pgfqpoint{0.407592in}{0.222904in}}{\pgfqpoint{1.411163in}{1.411163in}}%
\pgfusepath{clip}%
\pgfsetbuttcap%
\pgfsetroundjoin%
\definecolor{currentfill}{rgb}{0.909804,0.698039,0.423529}%
\pgfsetfillcolor{currentfill}%
\pgfsetfillopacity{0.120000}%
\pgfsetlinewidth{0.000000pt}%
\definecolor{currentstroke}{rgb}{0.000000,0.000000,0.000000}%
\pgfsetstrokecolor{currentstroke}%
\pgfsetdash{}{0pt}%
\pgfpathmoveto{\pgfqpoint{1.617626in}{1.138290in}}%
\pgfpathlineto{\pgfqpoint{1.635229in}{1.147622in}}%
\pgfpathlineto{\pgfqpoint{1.652909in}{1.157226in}}%
\pgfpathlineto{\pgfqpoint{1.660435in}{1.167474in}}%
\pgfpathlineto{\pgfqpoint{1.667949in}{1.177943in}}%
\pgfpathlineto{\pgfqpoint{1.650291in}{1.168581in}}%
\pgfpathlineto{\pgfqpoint{1.632710in}{1.159498in}}%
\pgfpathlineto{\pgfqpoint{1.625174in}{1.148780in}}%
\pgfpathlineto{\pgfqpoint{1.617626in}{1.138290in}}%
\pgfpathclose%
\pgfusepath{fill}%
\end{pgfscope}%
\begin{pgfscope}%
\pgfpathrectangle{\pgfqpoint{0.407592in}{0.222904in}}{\pgfqpoint{1.411163in}{1.411163in}}%
\pgfusepath{clip}%
\pgfsetbuttcap%
\pgfsetroundjoin%
\definecolor{currentfill}{rgb}{0.909804,0.698039,0.423529}%
\pgfsetfillcolor{currentfill}%
\pgfsetfillopacity{0.120000}%
\pgfsetlinewidth{0.000000pt}%
\definecolor{currentstroke}{rgb}{0.000000,0.000000,0.000000}%
\pgfsetstrokecolor{currentstroke}%
\pgfsetdash{}{0pt}%
\pgfpathmoveto{\pgfqpoint{0.989989in}{0.864492in}}%
\pgfpathlineto{\pgfqpoint{1.006467in}{0.860057in}}%
\pgfpathlineto{\pgfqpoint{1.022980in}{0.856565in}}%
\pgfpathlineto{\pgfqpoint{1.031374in}{0.856383in}}%
\pgfpathlineto{\pgfqpoint{1.039735in}{0.856663in}}%
\pgfpathlineto{\pgfqpoint{1.023262in}{0.860118in}}%
\pgfpathlineto{\pgfqpoint{1.006825in}{0.864545in}}%
\pgfpathlineto{\pgfqpoint{0.998424in}{0.864288in}}%
\pgfpathlineto{\pgfqpoint{0.989989in}{0.864492in}}%
\pgfpathclose%
\pgfusepath{fill}%
\end{pgfscope}%
\begin{pgfscope}%
\pgfpathrectangle{\pgfqpoint{0.407592in}{0.222904in}}{\pgfqpoint{1.411163in}{1.411163in}}%
\pgfusepath{clip}%
\pgfsetbuttcap%
\pgfsetroundjoin%
\definecolor{currentfill}{rgb}{0.909804,0.698039,0.423529}%
\pgfsetfillcolor{currentfill}%
\pgfsetfillopacity{0.120000}%
\pgfsetlinewidth{0.000000pt}%
\definecolor{currentstroke}{rgb}{0.000000,0.000000,0.000000}%
\pgfsetstrokecolor{currentstroke}%
\pgfsetdash{}{0pt}%
\pgfpathmoveto{\pgfqpoint{1.386849in}{0.960301in}}%
\pgfpathlineto{\pgfqpoint{1.403989in}{0.968123in}}%
\pgfpathlineto{\pgfqpoint{1.421200in}{0.976363in}}%
\pgfpathlineto{\pgfqpoint{1.428979in}{0.983900in}}%
\pgfpathlineto{\pgfqpoint{1.436741in}{0.991806in}}%
\pgfpathlineto{\pgfqpoint{1.419551in}{0.983717in}}%
\pgfpathlineto{\pgfqpoint{1.402433in}{0.976048in}}%
\pgfpathlineto{\pgfqpoint{1.394649in}{0.967980in}}%
\pgfpathlineto{\pgfqpoint{1.386849in}{0.960301in}}%
\pgfpathclose%
\pgfusepath{fill}%
\end{pgfscope}%
\begin{pgfscope}%
\pgfpathrectangle{\pgfqpoint{0.407592in}{0.222904in}}{\pgfqpoint{1.411163in}{1.411163in}}%
\pgfusepath{clip}%
\pgfsetbuttcap%
\pgfsetroundjoin%
\definecolor{currentfill}{rgb}{0.909804,0.698039,0.423529}%
\pgfsetfillcolor{currentfill}%
\pgfsetfillopacity{0.120000}%
\pgfsetlinewidth{0.000000pt}%
\definecolor{currentstroke}{rgb}{0.000000,0.000000,0.000000}%
\pgfsetstrokecolor{currentstroke}%
\pgfsetdash{}{0pt}%
\pgfpathmoveto{\pgfqpoint{1.733342in}{1.240411in}}%
\pgfpathlineto{\pgfqpoint{1.751185in}{1.250055in}}%
\pgfpathlineto{\pgfqpoint{1.769107in}{1.259938in}}%
\pgfpathlineto{\pgfqpoint{1.776517in}{1.270911in}}%
\pgfpathlineto{\pgfqpoint{1.783915in}{1.282058in}}%
\pgfpathlineto{\pgfqpoint{1.766017in}{1.272415in}}%
\pgfpathlineto{\pgfqpoint{1.748197in}{1.263018in}}%
\pgfpathlineto{\pgfqpoint{1.740775in}{1.251626in}}%
\pgfpathlineto{\pgfqpoint{1.733342in}{1.240411in}}%
\pgfpathclose%
\pgfusepath{fill}%
\end{pgfscope}%
\begin{pgfscope}%
\pgfpathrectangle{\pgfqpoint{0.407592in}{0.222904in}}{\pgfqpoint{1.411163in}{1.411163in}}%
\pgfusepath{clip}%
\pgfsetbuttcap%
\pgfsetroundjoin%
\definecolor{currentfill}{rgb}{0.909804,0.698039,0.423529}%
\pgfsetfillcolor{currentfill}%
\pgfsetfillopacity{0.120000}%
\pgfsetlinewidth{0.000000pt}%
\definecolor{currentstroke}{rgb}{0.000000,0.000000,0.000000}%
\pgfsetstrokecolor{currentstroke}%
\pgfsetdash{}{0pt}%
\pgfpathmoveto{\pgfqpoint{0.657740in}{1.046500in}}%
\pgfpathlineto{\pgfqpoint{0.674089in}{1.031982in}}%
\pgfpathlineto{\pgfqpoint{0.690450in}{1.017870in}}%
\pgfpathlineto{\pgfqpoint{0.699507in}{1.015984in}}%
\pgfpathlineto{\pgfqpoint{0.708521in}{1.014388in}}%
\pgfpathlineto{\pgfqpoint{0.692205in}{1.028748in}}%
\pgfpathlineto{\pgfqpoint{0.675900in}{1.043512in}}%
\pgfpathlineto{\pgfqpoint{0.666841in}{1.044863in}}%
\pgfpathlineto{\pgfqpoint{0.657740in}{1.046500in}}%
\pgfpathclose%
\pgfusepath{fill}%
\end{pgfscope}%
\begin{pgfscope}%
\pgfpathrectangle{\pgfqpoint{0.407592in}{0.222904in}}{\pgfqpoint{1.411163in}{1.411163in}}%
\pgfusepath{clip}%
\pgfsetbuttcap%
\pgfsetroundjoin%
\definecolor{currentfill}{rgb}{0.909804,0.698039,0.423529}%
\pgfsetfillcolor{currentfill}%
\pgfsetfillopacity{0.120000}%
\pgfsetlinewidth{0.000000pt}%
\definecolor{currentstroke}{rgb}{0.000000,0.000000,0.000000}%
\pgfsetstrokecolor{currentstroke}%
\pgfsetdash{}{0pt}%
\pgfpathmoveto{\pgfqpoint{0.541278in}{1.143383in}}%
\pgfpathlineto{\pgfqpoint{0.557608in}{1.127170in}}%
\pgfpathlineto{\pgfqpoint{0.573948in}{1.111229in}}%
\pgfpathlineto{\pgfqpoint{0.583222in}{1.109542in}}%
\pgfpathlineto{\pgfqpoint{0.592453in}{1.108112in}}%
\pgfpathlineto{\pgfqpoint{0.576157in}{1.124282in}}%
\pgfpathlineto{\pgfqpoint{0.559870in}{1.140715in}}%
\pgfpathlineto{\pgfqpoint{0.550596in}{1.141923in}}%
\pgfpathlineto{\pgfqpoint{0.541278in}{1.143383in}}%
\pgfpathclose%
\pgfusepath{fill}%
\end{pgfscope}%
\begin{pgfscope}%
\pgfpathrectangle{\pgfqpoint{0.407592in}{0.222904in}}{\pgfqpoint{1.411163in}{1.411163in}}%
\pgfusepath{clip}%
\pgfsetbuttcap%
\pgfsetroundjoin%
\definecolor{currentfill}{rgb}{0.909804,0.698039,0.423529}%
\pgfsetfillcolor{currentfill}%
\pgfsetfillopacity{0.120000}%
\pgfsetlinewidth{0.000000pt}%
\definecolor{currentstroke}{rgb}{0.000000,0.000000,0.000000}%
\pgfsetstrokecolor{currentstroke}%
\pgfsetdash{}{0pt}%
\pgfpathmoveto{\pgfqpoint{0.857088in}{0.913884in}}%
\pgfpathlineto{\pgfqpoint{0.873495in}{0.904484in}}%
\pgfpathlineto{\pgfqpoint{0.889924in}{0.895847in}}%
\pgfpathlineto{\pgfqpoint{0.898594in}{0.894167in}}%
\pgfpathlineto{\pgfqpoint{0.907226in}{0.892859in}}%
\pgfpathlineto{\pgfqpoint{0.890842in}{0.901641in}}%
\pgfpathlineto{\pgfqpoint{0.874480in}{0.911207in}}%
\pgfpathlineto{\pgfqpoint{0.865804in}{0.912361in}}%
\pgfpathlineto{\pgfqpoint{0.857088in}{0.913884in}}%
\pgfpathclose%
\pgfusepath{fill}%
\end{pgfscope}%
\begin{pgfscope}%
\pgfpathrectangle{\pgfqpoint{0.407592in}{0.222904in}}{\pgfqpoint{1.411163in}{1.411163in}}%
\pgfusepath{clip}%
\pgfsetbuttcap%
\pgfsetroundjoin%
\definecolor{currentfill}{rgb}{0.909804,0.698039,0.423529}%
\pgfsetfillcolor{currentfill}%
\pgfsetfillopacity{0.120000}%
\pgfsetlinewidth{0.000000pt}%
\definecolor{currentstroke}{rgb}{0.000000,0.000000,0.000000}%
\pgfsetstrokecolor{currentstroke}%
\pgfsetdash{}{0pt}%
\pgfpathmoveto{\pgfqpoint{1.552242in}{1.080028in}}%
\pgfpathlineto{\pgfqpoint{1.569738in}{1.089249in}}%
\pgfpathlineto{\pgfqpoint{1.587310in}{1.098764in}}%
\pgfpathlineto{\pgfqpoint{1.594908in}{1.108264in}}%
\pgfpathlineto{\pgfqpoint{1.602493in}{1.118023in}}%
\pgfpathlineto{\pgfqpoint{1.584943in}{1.108737in}}%
\pgfpathlineto{\pgfqpoint{1.567470in}{1.099749in}}%
\pgfpathlineto{\pgfqpoint{1.559862in}{1.089754in}}%
\pgfpathlineto{\pgfqpoint{1.552242in}{1.080028in}}%
\pgfpathclose%
\pgfusepath{fill}%
\end{pgfscope}%
\begin{pgfscope}%
\pgfpathrectangle{\pgfqpoint{0.407592in}{0.222904in}}{\pgfqpoint{1.411163in}{1.411163in}}%
\pgfusepath{clip}%
\pgfsetbuttcap%
\pgfsetroundjoin%
\definecolor{currentfill}{rgb}{0.909804,0.698039,0.423529}%
\pgfsetfillcolor{currentfill}%
\pgfsetfillopacity{0.120000}%
\pgfsetlinewidth{0.000000pt}%
\definecolor{currentstroke}{rgb}{0.000000,0.000000,0.000000}%
\pgfsetstrokecolor{currentstroke}%
\pgfsetdash{}{0pt}%
\pgfpathmoveto{\pgfqpoint{1.172180in}{0.862877in}}%
\pgfpathlineto{\pgfqpoint{1.188918in}{0.865978in}}%
\pgfpathlineto{\pgfqpoint{1.205714in}{0.869841in}}%
\pgfpathlineto{\pgfqpoint{1.213790in}{0.873020in}}%
\pgfpathlineto{\pgfqpoint{1.221842in}{0.876697in}}%
\pgfpathlineto{\pgfqpoint{1.205074in}{0.872745in}}%
\pgfpathlineto{\pgfqpoint{1.188366in}{0.869556in}}%
\pgfpathlineto{\pgfqpoint{1.180286in}{0.865957in}}%
\pgfpathlineto{\pgfqpoint{1.172180in}{0.862877in}}%
\pgfpathclose%
\pgfusepath{fill}%
\end{pgfscope}%
\begin{pgfscope}%
\pgfpathrectangle{\pgfqpoint{0.407592in}{0.222904in}}{\pgfqpoint{1.411163in}{1.411163in}}%
\pgfusepath{clip}%
\pgfsetbuttcap%
\pgfsetroundjoin%
\definecolor{currentfill}{rgb}{0.909804,0.698039,0.423529}%
\pgfsetfillcolor{currentfill}%
\pgfsetfillopacity{0.120000}%
\pgfsetlinewidth{0.000000pt}%
\definecolor{currentstroke}{rgb}{0.000000,0.000000,0.000000}%
\pgfsetstrokecolor{currentstroke}%
\pgfsetdash{}{0pt}%
\pgfpathmoveto{\pgfqpoint{0.773911in}{0.961616in}}%
\pgfpathlineto{\pgfqpoint{0.790293in}{0.949759in}}%
\pgfpathlineto{\pgfqpoint{0.806691in}{0.938508in}}%
\pgfpathlineto{\pgfqpoint{0.815531in}{0.936504in}}%
\pgfpathlineto{\pgfqpoint{0.824331in}{0.934830in}}%
\pgfpathlineto{\pgfqpoint{0.807978in}{0.946302in}}%
\pgfpathlineto{\pgfqpoint{0.791640in}{0.958388in}}%
\pgfpathlineto{\pgfqpoint{0.782796in}{0.959839in}}%
\pgfpathlineto{\pgfqpoint{0.773911in}{0.961616in}}%
\pgfpathclose%
\pgfusepath{fill}%
\end{pgfscope}%
\begin{pgfscope}%
\pgfpathrectangle{\pgfqpoint{0.407592in}{0.222904in}}{\pgfqpoint{1.411163in}{1.411163in}}%
\pgfusepath{clip}%
\pgfsetbuttcap%
\pgfsetroundjoin%
\definecolor{currentfill}{rgb}{0.909804,0.698039,0.423529}%
\pgfsetfillcolor{currentfill}%
\pgfsetfillopacity{0.120000}%
\pgfsetlinewidth{0.000000pt}%
\definecolor{currentstroke}{rgb}{0.000000,0.000000,0.000000}%
\pgfsetstrokecolor{currentstroke}%
\pgfsetdash{}{0pt}%
\pgfpathmoveto{\pgfqpoint{1.221842in}{0.876697in}}%
\pgfpathlineto{\pgfqpoint{1.238671in}{0.881363in}}%
\pgfpathlineto{\pgfqpoint{1.255564in}{0.886694in}}%
\pgfpathlineto{\pgfqpoint{1.263565in}{0.890928in}}%
\pgfpathlineto{\pgfqpoint{1.271543in}{0.895643in}}%
\pgfpathlineto{\pgfqpoint{1.254677in}{0.890272in}}%
\pgfpathlineto{\pgfqpoint{1.237874in}{0.885564in}}%
\pgfpathlineto{\pgfqpoint{1.229870in}{0.880878in}}%
\pgfpathlineto{\pgfqpoint{1.221842in}{0.876697in}}%
\pgfpathclose%
\pgfusepath{fill}%
\end{pgfscope}%
\begin{pgfscope}%
\pgfpathrectangle{\pgfqpoint{0.407592in}{0.222904in}}{\pgfqpoint{1.411163in}{1.411163in}}%
\pgfusepath{clip}%
\pgfsetbuttcap%
\pgfsetroundjoin%
\definecolor{currentfill}{rgb}{0.909804,0.698039,0.423529}%
\pgfsetfillcolor{currentfill}%
\pgfsetfillopacity{0.120000}%
\pgfsetlinewidth{0.000000pt}%
\definecolor{currentstroke}{rgb}{0.000000,0.000000,0.000000}%
\pgfsetstrokecolor{currentstroke}%
\pgfsetdash{}{0pt}%
\pgfpathmoveto{\pgfqpoint{1.122515in}{0.854750in}}%
\pgfpathlineto{\pgfqpoint{1.139175in}{0.856073in}}%
\pgfpathlineto{\pgfqpoint{1.155887in}{0.858243in}}%
\pgfpathlineto{\pgfqpoint{1.164048in}{0.860309in}}%
\pgfpathlineto{\pgfqpoint{1.172180in}{0.862877in}}%
\pgfpathlineto{\pgfqpoint{1.155500in}{0.860586in}}%
\pgfpathlineto{\pgfqpoint{1.138873in}{0.859153in}}%
\pgfpathlineto{\pgfqpoint{1.130709in}{0.856693in}}%
\pgfpathlineto{\pgfqpoint{1.122515in}{0.854750in}}%
\pgfpathclose%
\pgfusepath{fill}%
\end{pgfscope}%
\begin{pgfscope}%
\pgfpathrectangle{\pgfqpoint{0.407592in}{0.222904in}}{\pgfqpoint{1.411163in}{1.411163in}}%
\pgfusepath{clip}%
\pgfsetbuttcap%
\pgfsetroundjoin%
\definecolor{currentfill}{rgb}{0.909804,0.698039,0.423529}%
\pgfsetfillcolor{currentfill}%
\pgfsetfillopacity{0.120000}%
\pgfsetlinewidth{0.000000pt}%
\definecolor{currentstroke}{rgb}{0.000000,0.000000,0.000000}%
\pgfsetstrokecolor{currentstroke}%
\pgfsetdash{}{0pt}%
\pgfpathmoveto{\pgfqpoint{1.436741in}{0.991806in}}%
\pgfpathlineto{\pgfqpoint{1.454004in}{1.000286in}}%
\pgfpathlineto{\pgfqpoint{1.471340in}{1.009129in}}%
\pgfpathlineto{\pgfqpoint{1.479065in}{1.017223in}}%
\pgfpathlineto{\pgfqpoint{1.486776in}{1.025652in}}%
\pgfpathlineto{\pgfqpoint{1.469461in}{1.016988in}}%
\pgfpathlineto{\pgfqpoint{1.452220in}{1.008691in}}%
\pgfpathlineto{\pgfqpoint{1.444488in}{1.000073in}}%
\pgfpathlineto{\pgfqpoint{1.436741in}{0.991806in}}%
\pgfpathclose%
\pgfusepath{fill}%
\end{pgfscope}%
\begin{pgfscope}%
\pgfpathrectangle{\pgfqpoint{0.407592in}{0.222904in}}{\pgfqpoint{1.411163in}{1.411163in}}%
\pgfusepath{clip}%
\pgfsetbuttcap%
\pgfsetroundjoin%
\definecolor{currentfill}{rgb}{0.909804,0.698039,0.423529}%
\pgfsetfillcolor{currentfill}%
\pgfsetfillopacity{0.120000}%
\pgfsetlinewidth{0.000000pt}%
\definecolor{currentstroke}{rgb}{0.000000,0.000000,0.000000}%
\pgfsetstrokecolor{currentstroke}%
\pgfsetdash{}{0pt}%
\pgfpathmoveto{\pgfqpoint{1.667949in}{1.177943in}}%
\pgfpathlineto{\pgfqpoint{1.685683in}{1.187570in}}%
\pgfpathlineto{\pgfqpoint{1.703495in}{1.197448in}}%
\pgfpathlineto{\pgfqpoint{1.710974in}{1.207891in}}%
\pgfpathlineto{\pgfqpoint{1.718441in}{1.218537in}}%
\pgfpathlineto{\pgfqpoint{1.700652in}{1.208899in}}%
\pgfpathlineto{\pgfqpoint{1.682941in}{1.199517in}}%
\pgfpathlineto{\pgfqpoint{1.675450in}{1.188627in}}%
\pgfpathlineto{\pgfqpoint{1.667949in}{1.177943in}}%
\pgfpathclose%
\pgfusepath{fill}%
\end{pgfscope}%
\begin{pgfscope}%
\pgfpathrectangle{\pgfqpoint{0.407592in}{0.222904in}}{\pgfqpoint{1.411163in}{1.411163in}}%
\pgfusepath{clip}%
\pgfsetbuttcap%
\pgfsetroundjoin%
\definecolor{currentfill}{rgb}{0.909804,0.698039,0.423529}%
\pgfsetfillcolor{currentfill}%
\pgfsetfillopacity{0.120000}%
\pgfsetlinewidth{0.000000pt}%
\definecolor{currentstroke}{rgb}{0.000000,0.000000,0.000000}%
\pgfsetstrokecolor{currentstroke}%
\pgfsetdash{}{0pt}%
\pgfpathmoveto{\pgfqpoint{0.940064in}{0.877782in}}%
\pgfpathlineto{\pgfqpoint{0.956523in}{0.871542in}}%
\pgfpathlineto{\pgfqpoint{0.973011in}{0.866197in}}%
\pgfpathlineto{\pgfqpoint{0.981518in}{0.865135in}}%
\pgfpathlineto{\pgfqpoint{0.989989in}{0.864492in}}%
\pgfpathlineto{\pgfqpoint{0.973544in}{0.869864in}}%
\pgfpathlineto{\pgfqpoint{0.957128in}{0.876160in}}%
\pgfpathlineto{\pgfqpoint{0.948615in}{0.876763in}}%
\pgfpathlineto{\pgfqpoint{0.940064in}{0.877782in}}%
\pgfpathclose%
\pgfusepath{fill}%
\end{pgfscope}%
\begin{pgfscope}%
\pgfpathrectangle{\pgfqpoint{0.407592in}{0.222904in}}{\pgfqpoint{1.411163in}{1.411163in}}%
\pgfusepath{clip}%
\pgfsetbuttcap%
\pgfsetroundjoin%
\definecolor{currentfill}{rgb}{0.909804,0.698039,0.423529}%
\pgfsetfillcolor{currentfill}%
\pgfsetfillopacity{0.120000}%
\pgfsetlinewidth{0.000000pt}%
\definecolor{currentstroke}{rgb}{0.000000,0.000000,0.000000}%
\pgfsetstrokecolor{currentstroke}%
\pgfsetdash{}{0pt}%
\pgfpathmoveto{\pgfqpoint{1.271543in}{0.895643in}}%
\pgfpathlineto{\pgfqpoint{1.288476in}{0.901629in}}%
\pgfpathlineto{\pgfqpoint{1.305475in}{0.908184in}}%
\pgfpathlineto{\pgfqpoint{1.313407in}{0.913378in}}%
\pgfpathlineto{\pgfqpoint{1.321320in}{0.919025in}}%
\pgfpathlineto{\pgfqpoint{1.304345in}{0.912484in}}%
\pgfpathlineto{\pgfqpoint{1.287437in}{0.906512in}}%
\pgfpathlineto{\pgfqpoint{1.279501in}{0.900838in}}%
\pgfpathlineto{\pgfqpoint{1.271543in}{0.895643in}}%
\pgfpathclose%
\pgfusepath{fill}%
\end{pgfscope}%
\begin{pgfscope}%
\pgfpathrectangle{\pgfqpoint{0.407592in}{0.222904in}}{\pgfqpoint{1.411163in}{1.411163in}}%
\pgfusepath{clip}%
\pgfsetbuttcap%
\pgfsetroundjoin%
\definecolor{currentfill}{rgb}{0.909804,0.698039,0.423529}%
\pgfsetfillcolor{currentfill}%
\pgfsetfillopacity{0.120000}%
\pgfsetlinewidth{0.000000pt}%
\definecolor{currentstroke}{rgb}{0.000000,0.000000,0.000000}%
\pgfsetstrokecolor{currentstroke}%
\pgfsetdash{}{0pt}%
\pgfpathmoveto{\pgfqpoint{1.072799in}{0.852644in}}%
\pgfpathlineto{\pgfqpoint{1.089396in}{0.852048in}}%
\pgfpathlineto{\pgfqpoint{1.106040in}{0.852358in}}%
\pgfpathlineto{\pgfqpoint{1.114293in}{0.853311in}}%
\pgfpathlineto{\pgfqpoint{1.122515in}{0.854750in}}%
\pgfpathlineto{\pgfqpoint{1.105907in}{0.854317in}}%
\pgfpathlineto{\pgfqpoint{1.089346in}{0.854809in}}%
\pgfpathlineto{\pgfqpoint{1.081088in}{0.853479in}}%
\pgfpathlineto{\pgfqpoint{1.072799in}{0.852644in}}%
\pgfpathclose%
\pgfusepath{fill}%
\end{pgfscope}%
\begin{pgfscope}%
\pgfpathrectangle{\pgfqpoint{0.407592in}{0.222904in}}{\pgfqpoint{1.411163in}{1.411163in}}%
\pgfusepath{clip}%
\pgfsetbuttcap%
\pgfsetroundjoin%
\definecolor{currentfill}{rgb}{0.909804,0.698039,0.423529}%
\pgfsetfillcolor{currentfill}%
\pgfsetfillopacity{0.120000}%
\pgfsetlinewidth{0.000000pt}%
\definecolor{currentstroke}{rgb}{0.000000,0.000000,0.000000}%
\pgfsetstrokecolor{currentstroke}%
\pgfsetdash{}{0pt}%
\pgfpathmoveto{\pgfqpoint{0.573948in}{1.111229in}}%
\pgfpathlineto{\pgfqpoint{0.590297in}{1.095579in}}%
\pgfpathlineto{\pgfqpoint{0.606658in}{1.080241in}}%
\pgfpathlineto{\pgfqpoint{0.615888in}{1.078313in}}%
\pgfpathlineto{\pgfqpoint{0.625076in}{1.076648in}}%
\pgfpathlineto{\pgfqpoint{0.608759in}{1.092227in}}%
\pgfpathlineto{\pgfqpoint{0.592453in}{1.108112in}}%
\pgfpathlineto{\pgfqpoint{0.583222in}{1.109542in}}%
\pgfpathlineto{\pgfqpoint{0.573948in}{1.111229in}}%
\pgfpathclose%
\pgfusepath{fill}%
\end{pgfscope}%
\begin{pgfscope}%
\pgfpathrectangle{\pgfqpoint{0.407592in}{0.222904in}}{\pgfqpoint{1.411163in}{1.411163in}}%
\pgfusepath{clip}%
\pgfsetbuttcap%
\pgfsetroundjoin%
\definecolor{currentfill}{rgb}{0.909804,0.698039,0.423529}%
\pgfsetfillcolor{currentfill}%
\pgfsetfillopacity{0.120000}%
\pgfsetlinewidth{0.000000pt}%
\definecolor{currentstroke}{rgb}{0.000000,0.000000,0.000000}%
\pgfsetstrokecolor{currentstroke}%
\pgfsetdash{}{0pt}%
\pgfpathmoveto{\pgfqpoint{0.690450in}{1.017870in}}%
\pgfpathlineto{\pgfqpoint{0.706822in}{1.004195in}}%
\pgfpathlineto{\pgfqpoint{0.723207in}{0.990986in}}%
\pgfpathlineto{\pgfqpoint{0.732219in}{0.988846in}}%
\pgfpathlineto{\pgfqpoint{0.741190in}{0.987002in}}%
\pgfpathlineto{\pgfqpoint{0.724850in}{1.000462in}}%
\pgfpathlineto{\pgfqpoint{0.708521in}{1.014388in}}%
\pgfpathlineto{\pgfqpoint{0.699507in}{1.015984in}}%
\pgfpathlineto{\pgfqpoint{0.690450in}{1.017870in}}%
\pgfpathclose%
\pgfusepath{fill}%
\end{pgfscope}%
\begin{pgfscope}%
\pgfpathrectangle{\pgfqpoint{0.407592in}{0.222904in}}{\pgfqpoint{1.411163in}{1.411163in}}%
\pgfusepath{clip}%
\pgfsetbuttcap%
\pgfsetroundjoin%
\definecolor{currentfill}{rgb}{0.909804,0.698039,0.423529}%
\pgfsetfillcolor{currentfill}%
\pgfsetfillopacity{0.120000}%
\pgfsetlinewidth{0.000000pt}%
\definecolor{currentstroke}{rgb}{0.000000,0.000000,0.000000}%
\pgfsetstrokecolor{currentstroke}%
\pgfsetdash{}{0pt}%
\pgfpathmoveto{\pgfqpoint{1.321320in}{0.919025in}}%
\pgfpathlineto{\pgfqpoint{1.338364in}{0.926091in}}%
\pgfpathlineto{\pgfqpoint{1.355477in}{0.933643in}}%
\pgfpathlineto{\pgfqpoint{1.363347in}{0.939685in}}%
\pgfpathlineto{\pgfqpoint{1.371199in}{0.946146in}}%
\pgfpathlineto{\pgfqpoint{1.354108in}{0.938661in}}%
\pgfpathlineto{\pgfqpoint{1.337087in}{0.931661in}}%
\pgfpathlineto{\pgfqpoint{1.329213in}{0.925121in}}%
\pgfpathlineto{\pgfqpoint{1.321320in}{0.919025in}}%
\pgfpathclose%
\pgfusepath{fill}%
\end{pgfscope}%
\begin{pgfscope}%
\pgfpathrectangle{\pgfqpoint{0.407592in}{0.222904in}}{\pgfqpoint{1.411163in}{1.411163in}}%
\pgfusepath{clip}%
\pgfsetbuttcap%
\pgfsetroundjoin%
\definecolor{currentfill}{rgb}{0.909804,0.698039,0.423529}%
\pgfsetfillcolor{currentfill}%
\pgfsetfillopacity{0.120000}%
\pgfsetlinewidth{0.000000pt}%
\definecolor{currentstroke}{rgb}{0.000000,0.000000,0.000000}%
\pgfsetstrokecolor{currentstroke}%
\pgfsetdash{}{0pt}%
\pgfpathmoveto{\pgfqpoint{1.602493in}{1.118023in}}%
\pgfpathlineto{\pgfqpoint{1.620118in}{1.127592in}}%
\pgfpathlineto{\pgfqpoint{1.637821in}{1.137428in}}%
\pgfpathlineto{\pgfqpoint{1.645371in}{1.147209in}}%
\pgfpathlineto{\pgfqpoint{1.652909in}{1.157226in}}%
\pgfpathlineto{\pgfqpoint{1.635229in}{1.147622in}}%
\pgfpathlineto{\pgfqpoint{1.617626in}{1.138290in}}%
\pgfpathlineto{\pgfqpoint{1.610065in}{1.128035in}}%
\pgfpathlineto{\pgfqpoint{1.602493in}{1.118023in}}%
\pgfpathclose%
\pgfusepath{fill}%
\end{pgfscope}%
\begin{pgfscope}%
\pgfpathrectangle{\pgfqpoint{0.407592in}{0.222904in}}{\pgfqpoint{1.411163in}{1.411163in}}%
\pgfusepath{clip}%
\pgfsetbuttcap%
\pgfsetroundjoin%
\definecolor{currentfill}{rgb}{0.909804,0.698039,0.423529}%
\pgfsetfillcolor{currentfill}%
\pgfsetfillopacity{0.120000}%
\pgfsetlinewidth{0.000000pt}%
\definecolor{currentstroke}{rgb}{0.000000,0.000000,0.000000}%
\pgfsetstrokecolor{currentstroke}%
\pgfsetdash{}{0pt}%
\pgfpathmoveto{\pgfqpoint{1.486776in}{1.025652in}}%
\pgfpathlineto{\pgfqpoint{1.504165in}{1.034660in}}%
\pgfpathlineto{\pgfqpoint{1.521629in}{1.043989in}}%
\pgfpathlineto{\pgfqpoint{1.529303in}{1.052551in}}%
\pgfpathlineto{\pgfqpoint{1.536963in}{1.061417in}}%
\pgfpathlineto{\pgfqpoint{1.519521in}{1.052286in}}%
\pgfpathlineto{\pgfqpoint{1.502153in}{1.043481in}}%
\pgfpathlineto{\pgfqpoint{1.494472in}{1.034408in}}%
\pgfpathlineto{\pgfqpoint{1.486776in}{1.025652in}}%
\pgfpathclose%
\pgfusepath{fill}%
\end{pgfscope}%
\begin{pgfscope}%
\pgfpathrectangle{\pgfqpoint{0.407592in}{0.222904in}}{\pgfqpoint{1.411163in}{1.411163in}}%
\pgfusepath{clip}%
\pgfsetbuttcap%
\pgfsetroundjoin%
\definecolor{currentfill}{rgb}{0.909804,0.698039,0.423529}%
\pgfsetfillcolor{currentfill}%
\pgfsetfillopacity{0.120000}%
\pgfsetlinewidth{0.000000pt}%
\definecolor{currentstroke}{rgb}{0.000000,0.000000,0.000000}%
\pgfsetstrokecolor{currentstroke}%
\pgfsetdash{}{0pt}%
\pgfpathmoveto{\pgfqpoint{0.457113in}{1.213391in}}%
\pgfpathlineto{\pgfqpoint{0.473448in}{1.196468in}}%
\pgfpathlineto{\pgfqpoint{0.489793in}{1.179757in}}%
\pgfpathlineto{\pgfqpoint{0.499242in}{1.178035in}}%
\pgfpathlineto{\pgfqpoint{0.508647in}{1.176548in}}%
\pgfpathlineto{\pgfqpoint{0.492346in}{1.193469in}}%
\pgfpathlineto{\pgfqpoint{0.476054in}{1.210594in}}%
\pgfpathlineto{\pgfqpoint{0.466606in}{1.211878in}}%
\pgfpathlineto{\pgfqpoint{0.457113in}{1.213391in}}%
\pgfpathclose%
\pgfusepath{fill}%
\end{pgfscope}%
\begin{pgfscope}%
\pgfpathrectangle{\pgfqpoint{0.407592in}{0.222904in}}{\pgfqpoint{1.411163in}{1.411163in}}%
\pgfusepath{clip}%
\pgfsetbuttcap%
\pgfsetroundjoin%
\definecolor{currentfill}{rgb}{0.909804,0.698039,0.423529}%
\pgfsetfillcolor{currentfill}%
\pgfsetfillopacity{0.120000}%
\pgfsetlinewidth{0.000000pt}%
\definecolor{currentstroke}{rgb}{0.000000,0.000000,0.000000}%
\pgfsetstrokecolor{currentstroke}%
\pgfsetdash{}{0pt}%
\pgfpathmoveto{\pgfqpoint{1.783915in}{1.282058in}}%
\pgfpathlineto{\pgfqpoint{1.801892in}{1.291935in}}%
\pgfpathlineto{\pgfqpoint{1.819947in}{1.302038in}}%
\pgfpathlineto{\pgfqpoint{1.827310in}{1.313119in}}%
\pgfpathlineto{\pgfqpoint{1.809267in}{1.303133in}}%
\pgfpathlineto{\pgfqpoint{1.791303in}{1.293374in}}%
\pgfusepath{fill}%
\end{pgfscope}%
\begin{pgfscope}%
\pgfpathrectangle{\pgfqpoint{0.407592in}{0.222904in}}{\pgfqpoint{1.411163in}{1.411163in}}%
\pgfusepath{clip}%
\pgfsetbuttcap%
\pgfsetroundjoin%
\definecolor{currentfill}{rgb}{0.909804,0.698039,0.423529}%
\pgfsetfillcolor{currentfill}%
\pgfsetfillopacity{0.120000}%
\pgfsetlinewidth{0.000000pt}%
\definecolor{currentstroke}{rgb}{0.000000,0.000000,0.000000}%
\pgfsetstrokecolor{currentstroke}%
\pgfsetdash{}{0pt}%
\pgfpathmoveto{\pgfqpoint{1.022980in}{0.856565in}}%
\pgfpathlineto{\pgfqpoint{1.039531in}{0.854013in}}%
\pgfpathlineto{\pgfqpoint{1.056122in}{0.852389in}}%
\pgfpathlineto{\pgfqpoint{1.064477in}{0.852288in}}%
\pgfpathlineto{\pgfqpoint{1.072799in}{0.852644in}}%
\pgfpathlineto{\pgfqpoint{1.056246in}{0.854176in}}%
\pgfpathlineto{\pgfqpoint{1.039735in}{0.856663in}}%
\pgfpathlineto{\pgfqpoint{1.031374in}{0.856383in}}%
\pgfpathlineto{\pgfqpoint{1.022980in}{0.856565in}}%
\pgfpathclose%
\pgfusepath{fill}%
\end{pgfscope}%
\begin{pgfscope}%
\pgfpathrectangle{\pgfqpoint{0.407592in}{0.222904in}}{\pgfqpoint{1.411163in}{1.411163in}}%
\pgfusepath{clip}%
\pgfsetbuttcap%
\pgfsetroundjoin%
\definecolor{currentfill}{rgb}{0.909804,0.698039,0.423529}%
\pgfsetfillcolor{currentfill}%
\pgfsetfillopacity{0.120000}%
\pgfsetlinewidth{0.000000pt}%
\definecolor{currentstroke}{rgb}{0.000000,0.000000,0.000000}%
\pgfsetstrokecolor{currentstroke}%
\pgfsetdash{}{0pt}%
\pgfpathmoveto{\pgfqpoint{1.718441in}{1.218537in}}%
\pgfpathlineto{\pgfqpoint{1.736308in}{1.228421in}}%
\pgfpathlineto{\pgfqpoint{1.754253in}{1.238540in}}%
\pgfpathlineto{\pgfqpoint{1.761685in}{1.249145in}}%
\pgfpathlineto{\pgfqpoint{1.769107in}{1.259938in}}%
\pgfpathlineto{\pgfqpoint{1.751185in}{1.250055in}}%
\pgfpathlineto{\pgfqpoint{1.733342in}{1.240411in}}%
\pgfpathlineto{\pgfqpoint{1.725897in}{1.229379in}}%
\pgfpathlineto{\pgfqpoint{1.718441in}{1.218537in}}%
\pgfpathclose%
\pgfusepath{fill}%
\end{pgfscope}%
\begin{pgfscope}%
\pgfpathrectangle{\pgfqpoint{0.407592in}{0.222904in}}{\pgfqpoint{1.411163in}{1.411163in}}%
\pgfusepath{clip}%
\pgfsetbuttcap%
\pgfsetroundjoin%
\definecolor{currentfill}{rgb}{0.909804,0.698039,0.423529}%
\pgfsetfillcolor{currentfill}%
\pgfsetfillopacity{0.120000}%
\pgfsetlinewidth{0.000000pt}%
\definecolor{currentstroke}{rgb}{0.000000,0.000000,0.000000}%
\pgfsetstrokecolor{currentstroke}%
\pgfsetdash{}{0pt}%
\pgfpathmoveto{\pgfqpoint{0.806691in}{0.938508in}}%
\pgfpathlineto{\pgfqpoint{0.823106in}{0.927897in}}%
\pgfpathlineto{\pgfqpoint{0.839539in}{0.917962in}}%
\pgfpathlineto{\pgfqpoint{0.848333in}{0.915756in}}%
\pgfpathlineto{\pgfqpoint{0.857088in}{0.913884in}}%
\pgfpathlineto{\pgfqpoint{0.840701in}{0.924011in}}%
\pgfpathlineto{\pgfqpoint{0.824331in}{0.934830in}}%
\pgfpathlineto{\pgfqpoint{0.815531in}{0.936504in}}%
\pgfpathlineto{\pgfqpoint{0.806691in}{0.938508in}}%
\pgfpathclose%
\pgfusepath{fill}%
\end{pgfscope}%
\begin{pgfscope}%
\pgfpathrectangle{\pgfqpoint{0.407592in}{0.222904in}}{\pgfqpoint{1.411163in}{1.411163in}}%
\pgfusepath{clip}%
\pgfsetbuttcap%
\pgfsetroundjoin%
\definecolor{currentfill}{rgb}{0.909804,0.698039,0.423529}%
\pgfsetfillcolor{currentfill}%
\pgfsetfillopacity{0.120000}%
\pgfsetlinewidth{0.000000pt}%
\definecolor{currentstroke}{rgb}{0.000000,0.000000,0.000000}%
\pgfsetstrokecolor{currentstroke}%
\pgfsetdash{}{0pt}%
\pgfpathmoveto{\pgfqpoint{0.889924in}{0.895847in}}%
\pgfpathlineto{\pgfqpoint{0.906375in}{0.888005in}}%
\pgfpathlineto{\pgfqpoint{0.922850in}{0.880983in}}%
\pgfpathlineto{\pgfqpoint{0.931476in}{0.879195in}}%
\pgfpathlineto{\pgfqpoint{0.940064in}{0.877782in}}%
\pgfpathlineto{\pgfqpoint{0.923633in}{0.884895in}}%
\pgfpathlineto{\pgfqpoint{0.907226in}{0.892859in}}%
\pgfpathlineto{\pgfqpoint{0.898594in}{0.894167in}}%
\pgfpathlineto{\pgfqpoint{0.889924in}{0.895847in}}%
\pgfpathclose%
\pgfusepath{fill}%
\end{pgfscope}%
\begin{pgfscope}%
\pgfpathrectangle{\pgfqpoint{0.407592in}{0.222904in}}{\pgfqpoint{1.411163in}{1.411163in}}%
\pgfusepath{clip}%
\pgfsetbuttcap%
\pgfsetroundjoin%
\definecolor{currentfill}{rgb}{0.909804,0.698039,0.423529}%
\pgfsetfillcolor{currentfill}%
\pgfsetfillopacity{0.120000}%
\pgfsetlinewidth{0.000000pt}%
\definecolor{currentstroke}{rgb}{0.000000,0.000000,0.000000}%
\pgfsetstrokecolor{currentstroke}%
\pgfsetdash{}{0pt}%
\pgfpathmoveto{\pgfqpoint{1.371199in}{0.946146in}}%
\pgfpathlineto{\pgfqpoint{1.388361in}{0.954080in}}%
\pgfpathlineto{\pgfqpoint{1.405595in}{0.962430in}}%
\pgfpathlineto{\pgfqpoint{1.413406in}{0.969204in}}%
\pgfpathlineto{\pgfqpoint{1.421200in}{0.976363in}}%
\pgfpathlineto{\pgfqpoint{1.403989in}{0.968123in}}%
\pgfpathlineto{\pgfqpoint{1.386849in}{0.960301in}}%
\pgfpathlineto{\pgfqpoint{1.379032in}{0.953020in}}%
\pgfpathlineto{\pgfqpoint{1.371199in}{0.946146in}}%
\pgfpathclose%
\pgfusepath{fill}%
\end{pgfscope}%
\begin{pgfscope}%
\pgfpathrectangle{\pgfqpoint{0.407592in}{0.222904in}}{\pgfqpoint{1.411163in}{1.411163in}}%
\pgfusepath{clip}%
\pgfsetbuttcap%
\pgfsetroundjoin%
\definecolor{currentfill}{rgb}{0.909804,0.698039,0.423529}%
\pgfsetfillcolor{currentfill}%
\pgfsetfillopacity{0.120000}%
\pgfsetlinewidth{0.000000pt}%
\definecolor{currentstroke}{rgb}{0.000000,0.000000,0.000000}%
\pgfsetstrokecolor{currentstroke}%
\pgfsetdash{}{0pt}%
\pgfpathmoveto{\pgfqpoint{0.606658in}{1.080241in}}%
\pgfpathlineto{\pgfqpoint{0.623028in}{1.065238in}}%
\pgfpathlineto{\pgfqpoint{0.639410in}{1.050592in}}%
\pgfpathlineto{\pgfqpoint{0.648597in}{1.048412in}}%
\pgfpathlineto{\pgfqpoint{0.657740in}{1.046500in}}%
\pgfpathlineto{\pgfqpoint{0.641403in}{1.061397in}}%
\pgfpathlineto{\pgfqpoint{0.625076in}{1.076648in}}%
\pgfpathlineto{\pgfqpoint{0.615888in}{1.078313in}}%
\pgfpathlineto{\pgfqpoint{0.606658in}{1.080241in}}%
\pgfpathclose%
\pgfusepath{fill}%
\end{pgfscope}%
\begin{pgfscope}%
\pgfpathrectangle{\pgfqpoint{0.407592in}{0.222904in}}{\pgfqpoint{1.411163in}{1.411163in}}%
\pgfusepath{clip}%
\pgfsetbuttcap%
\pgfsetroundjoin%
\definecolor{currentfill}{rgb}{0.909804,0.698039,0.423529}%
\pgfsetfillcolor{currentfill}%
\pgfsetfillopacity{0.120000}%
\pgfsetlinewidth{0.000000pt}%
\definecolor{currentstroke}{rgb}{0.000000,0.000000,0.000000}%
\pgfsetstrokecolor{currentstroke}%
\pgfsetdash{}{0pt}%
\pgfpathmoveto{\pgfqpoint{0.489793in}{1.179757in}}%
\pgfpathlineto{\pgfqpoint{0.506147in}{1.163273in}}%
\pgfpathlineto{\pgfqpoint{0.522510in}{1.147032in}}%
\pgfpathlineto{\pgfqpoint{0.531916in}{1.145088in}}%
\pgfpathlineto{\pgfqpoint{0.541278in}{1.143383in}}%
\pgfpathlineto{\pgfqpoint{0.524958in}{1.159848in}}%
\pgfpathlineto{\pgfqpoint{0.508647in}{1.176548in}}%
\pgfpathlineto{\pgfqpoint{0.499242in}{1.178035in}}%
\pgfpathlineto{\pgfqpoint{0.489793in}{1.179757in}}%
\pgfpathclose%
\pgfusepath{fill}%
\end{pgfscope}%
\begin{pgfscope}%
\pgfpathrectangle{\pgfqpoint{0.407592in}{0.222904in}}{\pgfqpoint{1.411163in}{1.411163in}}%
\pgfusepath{clip}%
\pgfsetbuttcap%
\pgfsetroundjoin%
\definecolor{currentfill}{rgb}{0.909804,0.698039,0.423529}%
\pgfsetfillcolor{currentfill}%
\pgfsetfillopacity{0.120000}%
\pgfsetlinewidth{0.000000pt}%
\definecolor{currentstroke}{rgb}{0.000000,0.000000,0.000000}%
\pgfsetstrokecolor{currentstroke}%
\pgfsetdash{}{0pt}%
\pgfpathmoveto{\pgfqpoint{0.723207in}{0.990986in}}%
\pgfpathlineto{\pgfqpoint{0.739605in}{0.978274in}}%
\pgfpathlineto{\pgfqpoint{0.756017in}{0.966091in}}%
\pgfpathlineto{\pgfqpoint{0.764985in}{0.963704in}}%
\pgfpathlineto{\pgfqpoint{0.773911in}{0.961616in}}%
\pgfpathlineto{\pgfqpoint{0.757544in}{0.974042in}}%
\pgfpathlineto{\pgfqpoint{0.741190in}{0.987002in}}%
\pgfpathlineto{\pgfqpoint{0.732219in}{0.988846in}}%
\pgfpathlineto{\pgfqpoint{0.723207in}{0.990986in}}%
\pgfpathclose%
\pgfusepath{fill}%
\end{pgfscope}%
\begin{pgfscope}%
\pgfpathrectangle{\pgfqpoint{0.407592in}{0.222904in}}{\pgfqpoint{1.411163in}{1.411163in}}%
\pgfusepath{clip}%
\pgfsetbuttcap%
\pgfsetroundjoin%
\definecolor{currentfill}{rgb}{0.909804,0.698039,0.423529}%
\pgfsetfillcolor{currentfill}%
\pgfsetfillopacity{0.120000}%
\pgfsetlinewidth{0.000000pt}%
\definecolor{currentstroke}{rgb}{0.000000,0.000000,0.000000}%
\pgfsetstrokecolor{currentstroke}%
\pgfsetdash{}{0pt}%
\pgfpathmoveto{\pgfqpoint{1.536963in}{1.061417in}}%
\pgfpathlineto{\pgfqpoint{1.554480in}{1.070854in}}%
\pgfpathlineto{\pgfqpoint{1.572074in}{1.080578in}}%
\pgfpathlineto{\pgfqpoint{1.579698in}{1.089533in}}%
\pgfpathlineto{\pgfqpoint{1.587310in}{1.098764in}}%
\pgfpathlineto{\pgfqpoint{1.569738in}{1.089249in}}%
\pgfpathlineto{\pgfqpoint{1.552242in}{1.080028in}}%
\pgfpathlineto{\pgfqpoint{1.544609in}{1.070579in}}%
\pgfpathlineto{\pgfqpoint{1.536963in}{1.061417in}}%
\pgfpathclose%
\pgfusepath{fill}%
\end{pgfscope}%
\begin{pgfscope}%
\pgfpathrectangle{\pgfqpoint{0.407592in}{0.222904in}}{\pgfqpoint{1.411163in}{1.411163in}}%
\pgfusepath{clip}%
\pgfsetbuttcap%
\pgfsetroundjoin%
\definecolor{currentfill}{rgb}{0.909804,0.698039,0.423529}%
\pgfsetfillcolor{currentfill}%
\pgfsetfillopacity{0.120000}%
\pgfsetlinewidth{0.000000pt}%
\definecolor{currentstroke}{rgb}{0.000000,0.000000,0.000000}%
\pgfsetstrokecolor{currentstroke}%
\pgfsetdash{}{0pt}%
\pgfpathmoveto{\pgfqpoint{0.973011in}{0.866197in}}%
\pgfpathlineto{\pgfqpoint{0.989531in}{0.861757in}}%
\pgfpathlineto{\pgfqpoint{1.006086in}{0.858228in}}%
\pgfpathlineto{\pgfqpoint{1.014551in}{0.857187in}}%
\pgfpathlineto{\pgfqpoint{1.022980in}{0.856565in}}%
\pgfpathlineto{\pgfqpoint{1.006467in}{0.860057in}}%
\pgfpathlineto{\pgfqpoint{0.989989in}{0.864492in}}%
\pgfpathlineto{\pgfqpoint{0.981518in}{0.865135in}}%
\pgfpathlineto{\pgfqpoint{0.973011in}{0.866197in}}%
\pgfpathclose%
\pgfusepath{fill}%
\end{pgfscope}%
\begin{pgfscope}%
\pgfpathrectangle{\pgfqpoint{0.407592in}{0.222904in}}{\pgfqpoint{1.411163in}{1.411163in}}%
\pgfusepath{clip}%
\pgfsetbuttcap%
\pgfsetroundjoin%
\definecolor{currentfill}{rgb}{0.909804,0.698039,0.423529}%
\pgfsetfillcolor{currentfill}%
\pgfsetfillopacity{0.120000}%
\pgfsetlinewidth{0.000000pt}%
\definecolor{currentstroke}{rgb}{0.000000,0.000000,0.000000}%
\pgfsetstrokecolor{currentstroke}%
\pgfsetdash{}{0pt}%
\pgfpathmoveto{\pgfqpoint{1.652909in}{1.157226in}}%
\pgfpathlineto{\pgfqpoint{1.670666in}{1.167090in}}%
\pgfpathlineto{\pgfqpoint{1.688501in}{1.177199in}}%
\pgfpathlineto{\pgfqpoint{1.696004in}{1.187215in}}%
\pgfpathlineto{\pgfqpoint{1.703495in}{1.197448in}}%
\pgfpathlineto{\pgfqpoint{1.685683in}{1.187570in}}%
\pgfpathlineto{\pgfqpoint{1.667949in}{1.177943in}}%
\pgfpathlineto{\pgfqpoint{1.660435in}{1.167474in}}%
\pgfpathlineto{\pgfqpoint{1.652909in}{1.157226in}}%
\pgfpathclose%
\pgfusepath{fill}%
\end{pgfscope}%
\begin{pgfscope}%
\pgfpathrectangle{\pgfqpoint{0.407592in}{0.222904in}}{\pgfqpoint{1.411163in}{1.411163in}}%
\pgfusepath{clip}%
\pgfsetbuttcap%
\pgfsetroundjoin%
\definecolor{currentfill}{rgb}{0.909804,0.698039,0.423529}%
\pgfsetfillcolor{currentfill}%
\pgfsetfillopacity{0.120000}%
\pgfsetlinewidth{0.000000pt}%
\definecolor{currentstroke}{rgb}{0.000000,0.000000,0.000000}%
\pgfsetstrokecolor{currentstroke}%
\pgfsetdash{}{0pt}%
\pgfpathmoveto{\pgfqpoint{1.421200in}{0.976363in}}%
\pgfpathlineto{\pgfqpoint{1.438485in}{0.984991in}}%
\pgfpathlineto{\pgfqpoint{1.455844in}{0.993978in}}%
\pgfpathlineto{\pgfqpoint{1.463600in}{1.001378in}}%
\pgfpathlineto{\pgfqpoint{1.471340in}{1.009129in}}%
\pgfpathlineto{\pgfqpoint{1.454004in}{1.000286in}}%
\pgfpathlineto{\pgfqpoint{1.436741in}{0.991806in}}%
\pgfpathlineto{\pgfqpoint{1.428979in}{0.983900in}}%
\pgfpathlineto{\pgfqpoint{1.421200in}{0.976363in}}%
\pgfpathclose%
\pgfusepath{fill}%
\end{pgfscope}%
\begin{pgfscope}%
\pgfpathrectangle{\pgfqpoint{0.407592in}{0.222904in}}{\pgfqpoint{1.411163in}{1.411163in}}%
\pgfusepath{clip}%
\pgfsetbuttcap%
\pgfsetroundjoin%
\definecolor{currentfill}{rgb}{0.909804,0.698039,0.423529}%
\pgfsetfillcolor{currentfill}%
\pgfsetfillopacity{0.120000}%
\pgfsetlinewidth{0.000000pt}%
\definecolor{currentstroke}{rgb}{0.000000,0.000000,0.000000}%
\pgfsetstrokecolor{currentstroke}%
\pgfsetdash{}{0pt}%
\pgfpathmoveto{\pgfqpoint{1.769107in}{1.259938in}}%
\pgfpathlineto{\pgfqpoint{1.787107in}{1.270050in}}%
\pgfpathlineto{\pgfqpoint{1.805186in}{1.280383in}}%
\pgfpathlineto{\pgfqpoint{1.812572in}{1.291124in}}%
\pgfpathlineto{\pgfqpoint{1.819947in}{1.302038in}}%
\pgfpathlineto{\pgfqpoint{1.801892in}{1.291935in}}%
\pgfpathlineto{\pgfqpoint{1.783915in}{1.282058in}}%
\pgfpathlineto{\pgfqpoint{1.776517in}{1.270911in}}%
\pgfpathlineto{\pgfqpoint{1.769107in}{1.259938in}}%
\pgfpathclose%
\pgfusepath{fill}%
\end{pgfscope}%
\begin{pgfscope}%
\pgfpathrectangle{\pgfqpoint{0.407592in}{0.222904in}}{\pgfqpoint{1.411163in}{1.411163in}}%
\pgfusepath{clip}%
\pgfsetbuttcap%
\pgfsetroundjoin%
\definecolor{currentfill}{rgb}{0.909804,0.698039,0.423529}%
\pgfsetfillcolor{currentfill}%
\pgfsetfillopacity{0.120000}%
\pgfsetlinewidth{0.000000pt}%
\definecolor{currentstroke}{rgb}{0.000000,0.000000,0.000000}%
\pgfsetstrokecolor{currentstroke}%
\pgfsetdash{}{0pt}%
\pgfpathmoveto{\pgfqpoint{1.155887in}{0.858243in}}%
\pgfpathlineto{\pgfqpoint{1.172656in}{0.861216in}}%
\pgfpathlineto{\pgfqpoint{1.189483in}{0.864945in}}%
\pgfpathlineto{\pgfqpoint{1.197611in}{0.867152in}}%
\pgfpathlineto{\pgfqpoint{1.205714in}{0.869841in}}%
\pgfpathlineto{\pgfqpoint{1.188918in}{0.865978in}}%
\pgfpathlineto{\pgfqpoint{1.172180in}{0.862877in}}%
\pgfpathlineto{\pgfqpoint{1.164048in}{0.860309in}}%
\pgfpathlineto{\pgfqpoint{1.155887in}{0.858243in}}%
\pgfpathclose%
\pgfusepath{fill}%
\end{pgfscope}%
\begin{pgfscope}%
\pgfpathrectangle{\pgfqpoint{0.407592in}{0.222904in}}{\pgfqpoint{1.411163in}{1.411163in}}%
\pgfusepath{clip}%
\pgfsetbuttcap%
\pgfsetroundjoin%
\definecolor{currentfill}{rgb}{0.909804,0.698039,0.423529}%
\pgfsetfillcolor{currentfill}%
\pgfsetfillopacity{0.120000}%
\pgfsetlinewidth{0.000000pt}%
\definecolor{currentstroke}{rgb}{0.000000,0.000000,0.000000}%
\pgfsetstrokecolor{currentstroke}%
\pgfsetdash{}{0pt}%
\pgfpathmoveto{\pgfqpoint{1.205714in}{0.869841in}}%
\pgfpathlineto{\pgfqpoint{1.222571in}{0.874416in}}%
\pgfpathlineto{\pgfqpoint{1.239491in}{0.879655in}}%
\pgfpathlineto{\pgfqpoint{1.247539in}{0.882937in}}%
\pgfpathlineto{\pgfqpoint{1.255564in}{0.886694in}}%
\pgfpathlineto{\pgfqpoint{1.238671in}{0.881363in}}%
\pgfpathlineto{\pgfqpoint{1.221842in}{0.876697in}}%
\pgfpathlineto{\pgfqpoint{1.213790in}{0.873020in}}%
\pgfpathlineto{\pgfqpoint{1.205714in}{0.869841in}}%
\pgfpathclose%
\pgfusepath{fill}%
\end{pgfscope}%
\begin{pgfscope}%
\pgfpathrectangle{\pgfqpoint{0.407592in}{0.222904in}}{\pgfqpoint{1.411163in}{1.411163in}}%
\pgfusepath{clip}%
\pgfsetbuttcap%
\pgfsetroundjoin%
\definecolor{currentfill}{rgb}{0.909804,0.698039,0.423529}%
\pgfsetfillcolor{currentfill}%
\pgfsetfillopacity{0.120000}%
\pgfsetlinewidth{0.000000pt}%
\definecolor{currentstroke}{rgb}{0.000000,0.000000,0.000000}%
\pgfsetstrokecolor{currentstroke}%
\pgfsetdash{}{0pt}%
\pgfpathmoveto{\pgfqpoint{0.639410in}{1.050592in}}%
\pgfpathlineto{\pgfqpoint{0.655803in}{1.036328in}}%
\pgfpathlineto{\pgfqpoint{0.672208in}{1.022474in}}%
\pgfpathlineto{\pgfqpoint{0.681350in}{1.020037in}}%
\pgfpathlineto{\pgfqpoint{0.690450in}{1.017870in}}%
\pgfpathlineto{\pgfqpoint{0.674089in}{1.031982in}}%
\pgfpathlineto{\pgfqpoint{0.657740in}{1.046500in}}%
\pgfpathlineto{\pgfqpoint{0.648597in}{1.048412in}}%
\pgfpathlineto{\pgfqpoint{0.639410in}{1.050592in}}%
\pgfpathclose%
\pgfusepath{fill}%
\end{pgfscope}%
\begin{pgfscope}%
\pgfpathrectangle{\pgfqpoint{0.407592in}{0.222904in}}{\pgfqpoint{1.411163in}{1.411163in}}%
\pgfusepath{clip}%
\pgfsetbuttcap%
\pgfsetroundjoin%
\definecolor{currentfill}{rgb}{0.909804,0.698039,0.423529}%
\pgfsetfillcolor{currentfill}%
\pgfsetfillopacity{0.120000}%
\pgfsetlinewidth{0.000000pt}%
\definecolor{currentstroke}{rgb}{0.000000,0.000000,0.000000}%
\pgfsetstrokecolor{currentstroke}%
\pgfsetdash{}{0pt}%
\pgfpathmoveto{\pgfqpoint{0.522510in}{1.147032in}}%
\pgfpathlineto{\pgfqpoint{0.538884in}{1.131051in}}%
\pgfpathlineto{\pgfqpoint{0.555267in}{1.115347in}}%
\pgfpathlineto{\pgfqpoint{0.564629in}{1.113167in}}%
\pgfpathlineto{\pgfqpoint{0.573948in}{1.111229in}}%
\pgfpathlineto{\pgfqpoint{0.557608in}{1.127170in}}%
\pgfpathlineto{\pgfqpoint{0.541278in}{1.143383in}}%
\pgfpathlineto{\pgfqpoint{0.531916in}{1.145088in}}%
\pgfpathlineto{\pgfqpoint{0.522510in}{1.147032in}}%
\pgfpathclose%
\pgfusepath{fill}%
\end{pgfscope}%
\begin{pgfscope}%
\pgfpathrectangle{\pgfqpoint{0.407592in}{0.222904in}}{\pgfqpoint{1.411163in}{1.411163in}}%
\pgfusepath{clip}%
\pgfsetbuttcap%
\pgfsetroundjoin%
\definecolor{currentfill}{rgb}{0.909804,0.698039,0.423529}%
\pgfsetfillcolor{currentfill}%
\pgfsetfillopacity{0.120000}%
\pgfsetlinewidth{0.000000pt}%
\definecolor{currentstroke}{rgb}{0.000000,0.000000,0.000000}%
\pgfsetstrokecolor{currentstroke}%
\pgfsetdash{}{0pt}%
\pgfpathmoveto{\pgfqpoint{1.106040in}{0.852358in}}%
\pgfpathlineto{\pgfqpoint{1.122734in}{0.853542in}}%
\pgfpathlineto{\pgfqpoint{1.139481in}{0.855560in}}%
\pgfpathlineto{\pgfqpoint{1.147699in}{0.856664in}}%
\pgfpathlineto{\pgfqpoint{1.155887in}{0.858243in}}%
\pgfpathlineto{\pgfqpoint{1.139175in}{0.856073in}}%
\pgfpathlineto{\pgfqpoint{1.122515in}{0.854750in}}%
\pgfpathlineto{\pgfqpoint{1.114293in}{0.853311in}}%
\pgfpathlineto{\pgfqpoint{1.106040in}{0.852358in}}%
\pgfpathclose%
\pgfusepath{fill}%
\end{pgfscope}%
\begin{pgfscope}%
\pgfpathrectangle{\pgfqpoint{0.407592in}{0.222904in}}{\pgfqpoint{1.411163in}{1.411163in}}%
\pgfusepath{clip}%
\pgfsetbuttcap%
\pgfsetroundjoin%
\definecolor{currentfill}{rgb}{0.909804,0.698039,0.423529}%
\pgfsetfillcolor{currentfill}%
\pgfsetfillopacity{0.120000}%
\pgfsetlinewidth{0.000000pt}%
\definecolor{currentstroke}{rgb}{0.000000,0.000000,0.000000}%
\pgfsetstrokecolor{currentstroke}%
\pgfsetdash{}{0pt}%
\pgfpathmoveto{\pgfqpoint{0.839539in}{0.917962in}}%
\pgfpathlineto{\pgfqpoint{0.855991in}{0.908735in}}%
\pgfpathlineto{\pgfqpoint{0.872465in}{0.900249in}}%
\pgfpathlineto{\pgfqpoint{0.881214in}{0.897881in}}%
\pgfpathlineto{\pgfqpoint{0.889924in}{0.895847in}}%
\pgfpathlineto{\pgfqpoint{0.873495in}{0.904484in}}%
\pgfpathlineto{\pgfqpoint{0.857088in}{0.913884in}}%
\pgfpathlineto{\pgfqpoint{0.848333in}{0.915756in}}%
\pgfpathlineto{\pgfqpoint{0.839539in}{0.917962in}}%
\pgfpathclose%
\pgfusepath{fill}%
\end{pgfscope}%
\begin{pgfscope}%
\pgfpathrectangle{\pgfqpoint{0.407592in}{0.222904in}}{\pgfqpoint{1.411163in}{1.411163in}}%
\pgfusepath{clip}%
\pgfsetbuttcap%
\pgfsetroundjoin%
\definecolor{currentfill}{rgb}{0.909804,0.698039,0.423529}%
\pgfsetfillcolor{currentfill}%
\pgfsetfillopacity{0.120000}%
\pgfsetlinewidth{0.000000pt}%
\definecolor{currentstroke}{rgb}{0.000000,0.000000,0.000000}%
\pgfsetstrokecolor{currentstroke}%
\pgfsetdash{}{0pt}%
\pgfpathmoveto{\pgfqpoint{1.255564in}{0.886694in}}%
\pgfpathlineto{\pgfqpoint{1.272521in}{0.892640in}}%
\pgfpathlineto{\pgfqpoint{1.289546in}{0.899158in}}%
\pgfpathlineto{\pgfqpoint{1.297521in}{0.903444in}}%
\pgfpathlineto{\pgfqpoint{1.305475in}{0.908184in}}%
\pgfpathlineto{\pgfqpoint{1.288476in}{0.901629in}}%
\pgfpathlineto{\pgfqpoint{1.271543in}{0.895643in}}%
\pgfpathlineto{\pgfqpoint{1.263565in}{0.890928in}}%
\pgfpathlineto{\pgfqpoint{1.255564in}{0.886694in}}%
\pgfpathclose%
\pgfusepath{fill}%
\end{pgfscope}%
\begin{pgfscope}%
\pgfpathrectangle{\pgfqpoint{0.407592in}{0.222904in}}{\pgfqpoint{1.411163in}{1.411163in}}%
\pgfusepath{clip}%
\pgfsetbuttcap%
\pgfsetroundjoin%
\definecolor{currentfill}{rgb}{0.909804,0.698039,0.423529}%
\pgfsetfillcolor{currentfill}%
\pgfsetfillopacity{0.120000}%
\pgfsetlinewidth{0.000000pt}%
\definecolor{currentstroke}{rgb}{0.000000,0.000000,0.000000}%
\pgfsetstrokecolor{currentstroke}%
\pgfsetdash{}{0pt}%
\pgfpathmoveto{\pgfqpoint{1.587310in}{1.098764in}}%
\pgfpathlineto{\pgfqpoint{1.604958in}{1.108555in}}%
\pgfpathlineto{\pgfqpoint{1.622683in}{1.118607in}}%
\pgfpathlineto{\pgfqpoint{1.630258in}{1.127892in}}%
\pgfpathlineto{\pgfqpoint{1.637821in}{1.137428in}}%
\pgfpathlineto{\pgfqpoint{1.620118in}{1.127592in}}%
\pgfpathlineto{\pgfqpoint{1.602493in}{1.118023in}}%
\pgfpathlineto{\pgfqpoint{1.594908in}{1.108264in}}%
\pgfpathlineto{\pgfqpoint{1.587310in}{1.098764in}}%
\pgfpathclose%
\pgfusepath{fill}%
\end{pgfscope}%
\begin{pgfscope}%
\pgfpathrectangle{\pgfqpoint{0.407592in}{0.222904in}}{\pgfqpoint{1.411163in}{1.411163in}}%
\pgfusepath{clip}%
\pgfsetbuttcap%
\pgfsetroundjoin%
\definecolor{currentfill}{rgb}{0.909804,0.698039,0.423529}%
\pgfsetfillcolor{currentfill}%
\pgfsetfillopacity{0.120000}%
\pgfsetlinewidth{0.000000pt}%
\definecolor{currentstroke}{rgb}{0.000000,0.000000,0.000000}%
\pgfsetstrokecolor{currentstroke}%
\pgfsetdash{}{0pt}%
\pgfpathmoveto{\pgfqpoint{0.756017in}{0.966091in}}%
\pgfpathlineto{\pgfqpoint{0.772445in}{0.954470in}}%
\pgfpathlineto{\pgfqpoint{0.788888in}{0.943444in}}%
\pgfpathlineto{\pgfqpoint{0.797810in}{0.940826in}}%
\pgfpathlineto{\pgfqpoint{0.806691in}{0.938508in}}%
\pgfpathlineto{\pgfqpoint{0.790293in}{0.949759in}}%
\pgfpathlineto{\pgfqpoint{0.773911in}{0.961616in}}%
\pgfpathlineto{\pgfqpoint{0.764985in}{0.963704in}}%
\pgfpathlineto{\pgfqpoint{0.756017in}{0.966091in}}%
\pgfpathclose%
\pgfusepath{fill}%
\end{pgfscope}%
\begin{pgfscope}%
\pgfpathrectangle{\pgfqpoint{0.407592in}{0.222904in}}{\pgfqpoint{1.411163in}{1.411163in}}%
\pgfusepath{clip}%
\pgfsetbuttcap%
\pgfsetroundjoin%
\definecolor{currentfill}{rgb}{0.909804,0.698039,0.423529}%
\pgfsetfillcolor{currentfill}%
\pgfsetfillopacity{0.120000}%
\pgfsetlinewidth{0.000000pt}%
\definecolor{currentstroke}{rgb}{0.000000,0.000000,0.000000}%
\pgfsetstrokecolor{currentstroke}%
\pgfsetdash{}{0pt}%
\pgfpathmoveto{\pgfqpoint{0.922850in}{0.880983in}}%
\pgfpathlineto{\pgfqpoint{0.939353in}{0.874807in}}%
\pgfpathlineto{\pgfqpoint{0.955886in}{0.869491in}}%
\pgfpathlineto{\pgfqpoint{0.964467in}{0.867656in}}%
\pgfpathlineto{\pgfqpoint{0.973011in}{0.866197in}}%
\pgfpathlineto{\pgfqpoint{0.956523in}{0.871542in}}%
\pgfpathlineto{\pgfqpoint{0.940064in}{0.877782in}}%
\pgfpathlineto{\pgfqpoint{0.931476in}{0.879195in}}%
\pgfpathlineto{\pgfqpoint{0.922850in}{0.880983in}}%
\pgfpathclose%
\pgfusepath{fill}%
\end{pgfscope}%
\begin{pgfscope}%
\pgfpathrectangle{\pgfqpoint{0.407592in}{0.222904in}}{\pgfqpoint{1.411163in}{1.411163in}}%
\pgfusepath{clip}%
\pgfsetbuttcap%
\pgfsetroundjoin%
\definecolor{currentfill}{rgb}{0.909804,0.698039,0.423529}%
\pgfsetfillcolor{currentfill}%
\pgfsetfillopacity{0.120000}%
\pgfsetlinewidth{0.000000pt}%
\definecolor{currentstroke}{rgb}{0.000000,0.000000,0.000000}%
\pgfsetstrokecolor{currentstroke}%
\pgfsetdash{}{0pt}%
\pgfpathmoveto{\pgfqpoint{1.471340in}{1.009129in}}%
\pgfpathlineto{\pgfqpoint{1.488751in}{1.018310in}}%
\pgfpathlineto{\pgfqpoint{1.506237in}{1.027808in}}%
\pgfpathlineto{\pgfqpoint{1.513940in}{1.035738in}}%
\pgfpathlineto{\pgfqpoint{1.521629in}{1.043989in}}%
\pgfpathlineto{\pgfqpoint{1.504165in}{1.034660in}}%
\pgfpathlineto{\pgfqpoint{1.486776in}{1.025652in}}%
\pgfpathlineto{\pgfqpoint{1.479065in}{1.017223in}}%
\pgfpathlineto{\pgfqpoint{1.471340in}{1.009129in}}%
\pgfpathclose%
\pgfusepath{fill}%
\end{pgfscope}%
\begin{pgfscope}%
\pgfpathrectangle{\pgfqpoint{0.407592in}{0.222904in}}{\pgfqpoint{1.411163in}{1.411163in}}%
\pgfusepath{clip}%
\pgfsetbuttcap%
\pgfsetroundjoin%
\definecolor{currentfill}{rgb}{0.909804,0.698039,0.423529}%
\pgfsetfillcolor{currentfill}%
\pgfsetfillopacity{0.120000}%
\pgfsetlinewidth{0.000000pt}%
\definecolor{currentstroke}{rgb}{0.000000,0.000000,0.000000}%
\pgfsetstrokecolor{currentstroke}%
\pgfsetdash{}{0pt}%
\pgfpathmoveto{\pgfqpoint{1.703495in}{1.197448in}}%
\pgfpathlineto{\pgfqpoint{1.721385in}{1.207567in}}%
\pgfpathlineto{\pgfqpoint{1.739354in}{1.217914in}}%
\pgfpathlineto{\pgfqpoint{1.746809in}{1.228127in}}%
\pgfpathlineto{\pgfqpoint{1.754253in}{1.238540in}}%
\pgfpathlineto{\pgfqpoint{1.736308in}{1.228421in}}%
\pgfpathlineto{\pgfqpoint{1.718441in}{1.218537in}}%
\pgfpathlineto{\pgfqpoint{1.710974in}{1.207891in}}%
\pgfpathlineto{\pgfqpoint{1.703495in}{1.197448in}}%
\pgfpathclose%
\pgfusepath{fill}%
\end{pgfscope}%
\begin{pgfscope}%
\pgfpathrectangle{\pgfqpoint{0.407592in}{0.222904in}}{\pgfqpoint{1.411163in}{1.411163in}}%
\pgfusepath{clip}%
\pgfsetbuttcap%
\pgfsetroundjoin%
\definecolor{currentfill}{rgb}{0.909804,0.698039,0.423529}%
\pgfsetfillcolor{currentfill}%
\pgfsetfillopacity{0.120000}%
\pgfsetlinewidth{0.000000pt}%
\definecolor{currentstroke}{rgb}{0.000000,0.000000,0.000000}%
\pgfsetstrokecolor{currentstroke}%
\pgfsetdash{}{0pt}%
\pgfpathmoveto{\pgfqpoint{1.305475in}{0.908184in}}%
\pgfpathlineto{\pgfqpoint{1.322543in}{0.915266in}}%
\pgfpathlineto{\pgfqpoint{1.339680in}{0.922834in}}%
\pgfpathlineto{\pgfqpoint{1.347589in}{0.928025in}}%
\pgfpathlineto{\pgfqpoint{1.355477in}{0.933643in}}%
\pgfpathlineto{\pgfqpoint{1.338364in}{0.926091in}}%
\pgfpathlineto{\pgfqpoint{1.321320in}{0.919025in}}%
\pgfpathlineto{\pgfqpoint{1.313407in}{0.913378in}}%
\pgfpathlineto{\pgfqpoint{1.305475in}{0.908184in}}%
\pgfpathclose%
\pgfusepath{fill}%
\end{pgfscope}%
\begin{pgfscope}%
\pgfpathrectangle{\pgfqpoint{0.407592in}{0.222904in}}{\pgfqpoint{1.411163in}{1.411163in}}%
\pgfusepath{clip}%
\pgfsetbuttcap%
\pgfsetroundjoin%
\definecolor{currentfill}{rgb}{0.909804,0.698039,0.423529}%
\pgfsetfillcolor{currentfill}%
\pgfsetfillopacity{0.120000}%
\pgfsetlinewidth{0.000000pt}%
\definecolor{currentstroke}{rgb}{0.000000,0.000000,0.000000}%
\pgfsetstrokecolor{currentstroke}%
\pgfsetdash{}{0pt}%
\pgfpathmoveto{\pgfqpoint{1.819947in}{1.302038in}}%
\pgfpathlineto{\pgfqpoint{1.838081in}{1.312357in}}%
\pgfpathlineto{\pgfqpoint{1.845433in}{1.323323in}}%
\pgfpathlineto{\pgfqpoint{1.827310in}{1.313119in}}%
\pgfusepath{fill}%
\end{pgfscope}%
\begin{pgfscope}%
\pgfpathrectangle{\pgfqpoint{0.407592in}{0.222904in}}{\pgfqpoint{1.411163in}{1.411163in}}%
\pgfusepath{clip}%
\pgfsetbuttcap%
\pgfsetroundjoin%
\definecolor{currentfill}{rgb}{0.909804,0.698039,0.423529}%
\pgfsetfillcolor{currentfill}%
\pgfsetfillopacity{0.120000}%
\pgfsetlinewidth{0.000000pt}%
\definecolor{currentstroke}{rgb}{0.000000,0.000000,0.000000}%
\pgfsetstrokecolor{currentstroke}%
\pgfsetdash{}{0pt}%
\pgfpathmoveto{\pgfqpoint{1.056122in}{0.852389in}}%
\pgfpathlineto{\pgfqpoint{1.072758in}{0.851675in}}%
\pgfpathlineto{\pgfqpoint{1.089440in}{0.851845in}}%
\pgfpathlineto{\pgfqpoint{1.097756in}{0.851876in}}%
\pgfpathlineto{\pgfqpoint{1.106040in}{0.852358in}}%
\pgfpathlineto{\pgfqpoint{1.089396in}{0.852048in}}%
\pgfpathlineto{\pgfqpoint{1.072799in}{0.852644in}}%
\pgfpathlineto{\pgfqpoint{1.064477in}{0.852288in}}%
\pgfpathlineto{\pgfqpoint{1.056122in}{0.852389in}}%
\pgfpathclose%
\pgfusepath{fill}%
\end{pgfscope}%
\begin{pgfscope}%
\pgfpathrectangle{\pgfqpoint{0.407592in}{0.222904in}}{\pgfqpoint{1.411163in}{1.411163in}}%
\pgfusepath{clip}%
\pgfsetbuttcap%
\pgfsetroundjoin%
\definecolor{currentfill}{rgb}{0.909804,0.698039,0.423529}%
\pgfsetfillcolor{currentfill}%
\pgfsetfillopacity{0.120000}%
\pgfsetlinewidth{0.000000pt}%
\definecolor{currentstroke}{rgb}{0.000000,0.000000,0.000000}%
\pgfsetstrokecolor{currentstroke}%
\pgfsetdash{}{0pt}%
\pgfpathmoveto{\pgfqpoint{0.555267in}{1.115347in}}%
\pgfpathlineto{\pgfqpoint{0.571661in}{1.099941in}}%
\pgfpathlineto{\pgfqpoint{0.588065in}{1.084851in}}%
\pgfpathlineto{\pgfqpoint{0.597383in}{1.082424in}}%
\pgfpathlineto{\pgfqpoint{0.606658in}{1.080241in}}%
\pgfpathlineto{\pgfqpoint{0.590297in}{1.095579in}}%
\pgfpathlineto{\pgfqpoint{0.573948in}{1.111229in}}%
\pgfpathlineto{\pgfqpoint{0.564629in}{1.113167in}}%
\pgfpathlineto{\pgfqpoint{0.555267in}{1.115347in}}%
\pgfpathclose%
\pgfusepath{fill}%
\end{pgfscope}%
\begin{pgfscope}%
\pgfpathrectangle{\pgfqpoint{0.407592in}{0.222904in}}{\pgfqpoint{1.411163in}{1.411163in}}%
\pgfusepath{clip}%
\pgfsetbuttcap%
\pgfsetroundjoin%
\definecolor{currentfill}{rgb}{0.909804,0.698039,0.423529}%
\pgfsetfillcolor{currentfill}%
\pgfsetfillopacity{0.120000}%
\pgfsetlinewidth{0.000000pt}%
\definecolor{currentstroke}{rgb}{0.000000,0.000000,0.000000}%
\pgfsetstrokecolor{currentstroke}%
\pgfsetdash{}{0pt}%
\pgfpathmoveto{\pgfqpoint{0.672208in}{1.022474in}}%
\pgfpathlineto{\pgfqpoint{0.688625in}{1.009055in}}%
\pgfpathlineto{\pgfqpoint{0.705055in}{0.996102in}}%
\pgfpathlineto{\pgfqpoint{0.714152in}{0.993408in}}%
\pgfpathlineto{\pgfqpoint{0.723207in}{0.990986in}}%
\pgfpathlineto{\pgfqpoint{0.706822in}{1.004195in}}%
\pgfpathlineto{\pgfqpoint{0.690450in}{1.017870in}}%
\pgfpathlineto{\pgfqpoint{0.681350in}{1.020037in}}%
\pgfpathlineto{\pgfqpoint{0.672208in}{1.022474in}}%
\pgfpathclose%
\pgfusepath{fill}%
\end{pgfscope}%
\begin{pgfscope}%
\pgfpathrectangle{\pgfqpoint{0.407592in}{0.222904in}}{\pgfqpoint{1.411163in}{1.411163in}}%
\pgfusepath{clip}%
\pgfsetbuttcap%
\pgfsetroundjoin%
\definecolor{currentfill}{rgb}{0.909804,0.698039,0.423529}%
\pgfsetfillcolor{currentfill}%
\pgfsetfillopacity{0.120000}%
\pgfsetlinewidth{0.000000pt}%
\definecolor{currentstroke}{rgb}{0.000000,0.000000,0.000000}%
\pgfsetstrokecolor{currentstroke}%
\pgfsetdash{}{0pt}%
\pgfpathmoveto{\pgfqpoint{1.355477in}{0.933643in}}%
\pgfpathlineto{\pgfqpoint{1.372662in}{0.941644in}}%
\pgfpathlineto{\pgfqpoint{1.389919in}{0.950059in}}%
\pgfpathlineto{\pgfqpoint{1.397766in}{0.956047in}}%
\pgfpathlineto{\pgfqpoint{1.405595in}{0.962430in}}%
\pgfpathlineto{\pgfqpoint{1.388361in}{0.954080in}}%
\pgfpathlineto{\pgfqpoint{1.371199in}{0.946146in}}%
\pgfpathlineto{\pgfqpoint{1.363347in}{0.939685in}}%
\pgfpathlineto{\pgfqpoint{1.355477in}{0.933643in}}%
\pgfpathclose%
\pgfusepath{fill}%
\end{pgfscope}%
\begin{pgfscope}%
\pgfpathrectangle{\pgfqpoint{0.407592in}{0.222904in}}{\pgfqpoint{1.411163in}{1.411163in}}%
\pgfusepath{clip}%
\pgfsetbuttcap%
\pgfsetroundjoin%
\definecolor{currentfill}{rgb}{0.909804,0.698039,0.423529}%
\pgfsetfillcolor{currentfill}%
\pgfsetfillopacity{0.120000}%
\pgfsetlinewidth{0.000000pt}%
\definecolor{currentstroke}{rgb}{0.000000,0.000000,0.000000}%
\pgfsetstrokecolor{currentstroke}%
\pgfsetdash{}{0pt}%
\pgfpathmoveto{\pgfqpoint{0.437993in}{1.217086in}}%
\pgfpathlineto{\pgfqpoint{0.454371in}{1.200373in}}%
\pgfpathlineto{\pgfqpoint{0.470759in}{1.183880in}}%
\pgfpathlineto{\pgfqpoint{0.480298in}{1.181707in}}%
\pgfpathlineto{\pgfqpoint{0.489793in}{1.179757in}}%
\pgfpathlineto{\pgfqpoint{0.473448in}{1.196468in}}%
\pgfpathlineto{\pgfqpoint{0.457113in}{1.213391in}}%
\pgfpathlineto{\pgfqpoint{0.447576in}{1.215129in}}%
\pgfpathlineto{\pgfqpoint{0.437993in}{1.217086in}}%
\pgfpathclose%
\pgfusepath{fill}%
\end{pgfscope}%
\begin{pgfscope}%
\pgfpathrectangle{\pgfqpoint{0.407592in}{0.222904in}}{\pgfqpoint{1.411163in}{1.411163in}}%
\pgfusepath{clip}%
\pgfsetbuttcap%
\pgfsetroundjoin%
\definecolor{currentfill}{rgb}{0.909804,0.698039,0.423529}%
\pgfsetfillcolor{currentfill}%
\pgfsetfillopacity{0.120000}%
\pgfsetlinewidth{0.000000pt}%
\definecolor{currentstroke}{rgb}{0.000000,0.000000,0.000000}%
\pgfsetstrokecolor{currentstroke}%
\pgfsetdash{}{0pt}%
\pgfpathmoveto{\pgfqpoint{1.006086in}{0.858228in}}%
\pgfpathlineto{\pgfqpoint{1.022679in}{0.855607in}}%
\pgfpathlineto{\pgfqpoint{1.039311in}{0.853885in}}%
\pgfpathlineto{\pgfqpoint{1.047734in}{0.852928in}}%
\pgfpathlineto{\pgfqpoint{1.056122in}{0.852389in}}%
\pgfpathlineto{\pgfqpoint{1.039531in}{0.854013in}}%
\pgfpathlineto{\pgfqpoint{1.022980in}{0.856565in}}%
\pgfpathlineto{\pgfqpoint{1.014551in}{0.857187in}}%
\pgfpathlineto{\pgfqpoint{1.006086in}{0.858228in}}%
\pgfpathclose%
\pgfusepath{fill}%
\end{pgfscope}%
\begin{pgfscope}%
\pgfpathrectangle{\pgfqpoint{0.407592in}{0.222904in}}{\pgfqpoint{1.411163in}{1.411163in}}%
\pgfusepath{clip}%
\pgfsetbuttcap%
\pgfsetroundjoin%
\definecolor{currentfill}{rgb}{0.909804,0.698039,0.423529}%
\pgfsetfillcolor{currentfill}%
\pgfsetfillopacity{0.120000}%
\pgfsetlinewidth{0.000000pt}%
\definecolor{currentstroke}{rgb}{0.000000,0.000000,0.000000}%
\pgfsetstrokecolor{currentstroke}%
\pgfsetdash{}{0pt}%
\pgfpathmoveto{\pgfqpoint{1.637821in}{1.137428in}}%
\pgfpathlineto{\pgfqpoint{1.655601in}{1.147516in}}%
\pgfpathlineto{\pgfqpoint{1.673460in}{1.157845in}}%
\pgfpathlineto{\pgfqpoint{1.680987in}{1.167407in}}%
\pgfpathlineto{\pgfqpoint{1.688501in}{1.177199in}}%
\pgfpathlineto{\pgfqpoint{1.670666in}{1.167090in}}%
\pgfpathlineto{\pgfqpoint{1.652909in}{1.157226in}}%
\pgfpathlineto{\pgfqpoint{1.645371in}{1.147209in}}%
\pgfpathlineto{\pgfqpoint{1.637821in}{1.137428in}}%
\pgfpathclose%
\pgfusepath{fill}%
\end{pgfscope}%
\begin{pgfscope}%
\pgfpathrectangle{\pgfqpoint{0.407592in}{0.222904in}}{\pgfqpoint{1.411163in}{1.411163in}}%
\pgfusepath{clip}%
\pgfsetbuttcap%
\pgfsetroundjoin%
\definecolor{currentfill}{rgb}{0.909804,0.698039,0.423529}%
\pgfsetfillcolor{currentfill}%
\pgfsetfillopacity{0.120000}%
\pgfsetlinewidth{0.000000pt}%
\definecolor{currentstroke}{rgb}{0.000000,0.000000,0.000000}%
\pgfsetstrokecolor{currentstroke}%
\pgfsetdash{}{0pt}%
\pgfpathmoveto{\pgfqpoint{1.521629in}{1.043989in}}%
\pgfpathlineto{\pgfqpoint{1.539168in}{1.053617in}}%
\pgfpathlineto{\pgfqpoint{1.556785in}{1.063528in}}%
\pgfpathlineto{\pgfqpoint{1.564436in}{1.071907in}}%
\pgfpathlineto{\pgfqpoint{1.572074in}{1.080578in}}%
\pgfpathlineto{\pgfqpoint{1.554480in}{1.070854in}}%
\pgfpathlineto{\pgfqpoint{1.536963in}{1.061417in}}%
\pgfpathlineto{\pgfqpoint{1.529303in}{1.052551in}}%
\pgfpathlineto{\pgfqpoint{1.521629in}{1.043989in}}%
\pgfpathclose%
\pgfusepath{fill}%
\end{pgfscope}%
\begin{pgfscope}%
\pgfpathrectangle{\pgfqpoint{0.407592in}{0.222904in}}{\pgfqpoint{1.411163in}{1.411163in}}%
\pgfusepath{clip}%
\pgfsetbuttcap%
\pgfsetroundjoin%
\definecolor{currentfill}{rgb}{0.909804,0.698039,0.423529}%
\pgfsetfillcolor{currentfill}%
\pgfsetfillopacity{0.120000}%
\pgfsetlinewidth{0.000000pt}%
\definecolor{currentstroke}{rgb}{0.000000,0.000000,0.000000}%
\pgfsetstrokecolor{currentstroke}%
\pgfsetdash{}{0pt}%
\pgfpathmoveto{\pgfqpoint{1.754253in}{1.238540in}}%
\pgfpathlineto{\pgfqpoint{1.772277in}{1.248883in}}%
\pgfpathlineto{\pgfqpoint{1.790381in}{1.259440in}}%
\pgfpathlineto{\pgfqpoint{1.797789in}{1.269820in}}%
\pgfpathlineto{\pgfqpoint{1.805186in}{1.280383in}}%
\pgfpathlineto{\pgfqpoint{1.787107in}{1.270050in}}%
\pgfpathlineto{\pgfqpoint{1.769107in}{1.259938in}}%
\pgfpathlineto{\pgfqpoint{1.761685in}{1.249145in}}%
\pgfpathlineto{\pgfqpoint{1.754253in}{1.238540in}}%
\pgfpathclose%
\pgfusepath{fill}%
\end{pgfscope}%
\begin{pgfscope}%
\pgfpathrectangle{\pgfqpoint{0.407592in}{0.222904in}}{\pgfqpoint{1.411163in}{1.411163in}}%
\pgfusepath{clip}%
\pgfsetbuttcap%
\pgfsetroundjoin%
\definecolor{currentfill}{rgb}{0.909804,0.698039,0.423529}%
\pgfsetfillcolor{currentfill}%
\pgfsetfillopacity{0.120000}%
\pgfsetlinewidth{0.000000pt}%
\definecolor{currentstroke}{rgb}{0.000000,0.000000,0.000000}%
\pgfsetstrokecolor{currentstroke}%
\pgfsetdash{}{0pt}%
\pgfpathmoveto{\pgfqpoint{0.872465in}{0.900249in}}%
\pgfpathlineto{\pgfqpoint{0.888961in}{0.892531in}}%
\pgfpathlineto{\pgfqpoint{0.905483in}{0.885606in}}%
\pgfpathlineto{\pgfqpoint{0.914186in}{0.883127in}}%
\pgfpathlineto{\pgfqpoint{0.922850in}{0.880983in}}%
\pgfpathlineto{\pgfqpoint{0.906375in}{0.888005in}}%
\pgfpathlineto{\pgfqpoint{0.889924in}{0.895847in}}%
\pgfpathlineto{\pgfqpoint{0.881214in}{0.897881in}}%
\pgfpathlineto{\pgfqpoint{0.872465in}{0.900249in}}%
\pgfpathclose%
\pgfusepath{fill}%
\end{pgfscope}%
\begin{pgfscope}%
\pgfpathrectangle{\pgfqpoint{0.407592in}{0.222904in}}{\pgfqpoint{1.411163in}{1.411163in}}%
\pgfusepath{clip}%
\pgfsetbuttcap%
\pgfsetroundjoin%
\definecolor{currentfill}{rgb}{0.909804,0.698039,0.423529}%
\pgfsetfillcolor{currentfill}%
\pgfsetfillopacity{0.120000}%
\pgfsetlinewidth{0.000000pt}%
\definecolor{currentstroke}{rgb}{0.000000,0.000000,0.000000}%
\pgfsetstrokecolor{currentstroke}%
\pgfsetdash{}{0pt}%
\pgfpathmoveto{\pgfqpoint{0.788888in}{0.943444in}}%
\pgfpathlineto{\pgfqpoint{0.805348in}{0.933045in}}%
\pgfpathlineto{\pgfqpoint{0.821827in}{0.923305in}}%
\pgfpathlineto{\pgfqpoint{0.830703in}{0.920483in}}%
\pgfpathlineto{\pgfqpoint{0.839539in}{0.917962in}}%
\pgfpathlineto{\pgfqpoint{0.823106in}{0.927897in}}%
\pgfpathlineto{\pgfqpoint{0.806691in}{0.938508in}}%
\pgfpathlineto{\pgfqpoint{0.797810in}{0.940826in}}%
\pgfpathlineto{\pgfqpoint{0.788888in}{0.943444in}}%
\pgfpathclose%
\pgfusepath{fill}%
\end{pgfscope}%
\begin{pgfscope}%
\pgfpathrectangle{\pgfqpoint{0.407592in}{0.222904in}}{\pgfqpoint{1.411163in}{1.411163in}}%
\pgfusepath{clip}%
\pgfsetbuttcap%
\pgfsetroundjoin%
\definecolor{currentfill}{rgb}{0.909804,0.698039,0.423529}%
\pgfsetfillcolor{currentfill}%
\pgfsetfillopacity{0.120000}%
\pgfsetlinewidth{0.000000pt}%
\definecolor{currentstroke}{rgb}{0.000000,0.000000,0.000000}%
\pgfsetstrokecolor{currentstroke}%
\pgfsetdash{}{0pt}%
\pgfpathmoveto{\pgfqpoint{1.405595in}{0.962430in}}%
\pgfpathlineto{\pgfqpoint{1.422902in}{0.971166in}}%
\pgfpathlineto{\pgfqpoint{1.440283in}{0.980259in}}%
\pgfpathlineto{\pgfqpoint{1.448072in}{0.986936in}}%
\pgfpathlineto{\pgfqpoint{1.455844in}{0.993978in}}%
\pgfpathlineto{\pgfqpoint{1.438485in}{0.984991in}}%
\pgfpathlineto{\pgfqpoint{1.421200in}{0.976363in}}%
\pgfpathlineto{\pgfqpoint{1.413406in}{0.969204in}}%
\pgfpathlineto{\pgfqpoint{1.405595in}{0.962430in}}%
\pgfpathclose%
\pgfusepath{fill}%
\end{pgfscope}%
\begin{pgfscope}%
\pgfpathrectangle{\pgfqpoint{0.407592in}{0.222904in}}{\pgfqpoint{1.411163in}{1.411163in}}%
\pgfusepath{clip}%
\pgfsetbuttcap%
\pgfsetroundjoin%
\definecolor{currentfill}{rgb}{0.909804,0.698039,0.423529}%
\pgfsetfillcolor{currentfill}%
\pgfsetfillopacity{0.120000}%
\pgfsetlinewidth{0.000000pt}%
\definecolor{currentstroke}{rgb}{0.000000,0.000000,0.000000}%
\pgfsetstrokecolor{currentstroke}%
\pgfsetdash{}{0pt}%
\pgfpathmoveto{\pgfqpoint{0.588065in}{1.084851in}}%
\pgfpathlineto{\pgfqpoint{0.604480in}{1.070101in}}%
\pgfpathlineto{\pgfqpoint{0.620906in}{1.055712in}}%
\pgfpathlineto{\pgfqpoint{0.630180in}{1.053028in}}%
\pgfpathlineto{\pgfqpoint{0.639410in}{1.050592in}}%
\pgfpathlineto{\pgfqpoint{0.623028in}{1.065238in}}%
\pgfpathlineto{\pgfqpoint{0.606658in}{1.080241in}}%
\pgfpathlineto{\pgfqpoint{0.597383in}{1.082424in}}%
\pgfpathlineto{\pgfqpoint{0.588065in}{1.084851in}}%
\pgfpathclose%
\pgfusepath{fill}%
\end{pgfscope}%
\begin{pgfscope}%
\pgfpathrectangle{\pgfqpoint{0.407592in}{0.222904in}}{\pgfqpoint{1.411163in}{1.411163in}}%
\pgfusepath{clip}%
\pgfsetbuttcap%
\pgfsetroundjoin%
\definecolor{currentfill}{rgb}{0.909804,0.698039,0.423529}%
\pgfsetfillcolor{currentfill}%
\pgfsetfillopacity{0.120000}%
\pgfsetlinewidth{0.000000pt}%
\definecolor{currentstroke}{rgb}{0.000000,0.000000,0.000000}%
\pgfsetstrokecolor{currentstroke}%
\pgfsetdash{}{0pt}%
\pgfpathmoveto{\pgfqpoint{1.805186in}{1.280383in}}%
\pgfpathlineto{\pgfqpoint{1.823345in}{1.290927in}}%
\pgfpathlineto{\pgfqpoint{1.830719in}{1.301556in}}%
\pgfpathlineto{\pgfqpoint{1.838081in}{1.312357in}}%
\pgfpathlineto{\pgfqpoint{1.819947in}{1.302038in}}%
\pgfpathlineto{\pgfqpoint{1.812572in}{1.291124in}}%
\pgfusepath{fill}%
\end{pgfscope}%
\begin{pgfscope}%
\pgfpathrectangle{\pgfqpoint{0.407592in}{0.222904in}}{\pgfqpoint{1.411163in}{1.411163in}}%
\pgfusepath{clip}%
\pgfsetbuttcap%
\pgfsetroundjoin%
\definecolor{currentfill}{rgb}{0.909804,0.698039,0.423529}%
\pgfsetfillcolor{currentfill}%
\pgfsetfillopacity{0.120000}%
\pgfsetlinewidth{0.000000pt}%
\definecolor{currentstroke}{rgb}{0.000000,0.000000,0.000000}%
\pgfsetstrokecolor{currentstroke}%
\pgfsetdash{}{0pt}%
\pgfpathmoveto{\pgfqpoint{0.470759in}{1.183880in}}%
\pgfpathlineto{\pgfqpoint{0.487157in}{1.167621in}}%
\pgfpathlineto{\pgfqpoint{0.503564in}{1.151611in}}%
\pgfpathlineto{\pgfqpoint{0.513059in}{1.149208in}}%
\pgfpathlineto{\pgfqpoint{0.522510in}{1.147032in}}%
\pgfpathlineto{\pgfqpoint{0.506147in}{1.163273in}}%
\pgfpathlineto{\pgfqpoint{0.489793in}{1.179757in}}%
\pgfpathlineto{\pgfqpoint{0.480298in}{1.181707in}}%
\pgfpathlineto{\pgfqpoint{0.470759in}{1.183880in}}%
\pgfpathclose%
\pgfusepath{fill}%
\end{pgfscope}%
\begin{pgfscope}%
\pgfpathrectangle{\pgfqpoint{0.407592in}{0.222904in}}{\pgfqpoint{1.411163in}{1.411163in}}%
\pgfusepath{clip}%
\pgfsetbuttcap%
\pgfsetroundjoin%
\definecolor{currentfill}{rgb}{0.909804,0.698039,0.423529}%
\pgfsetfillcolor{currentfill}%
\pgfsetfillopacity{0.120000}%
\pgfsetlinewidth{0.000000pt}%
\definecolor{currentstroke}{rgb}{0.000000,0.000000,0.000000}%
\pgfsetstrokecolor{currentstroke}%
\pgfsetdash{}{0pt}%
\pgfpathmoveto{\pgfqpoint{0.705055in}{0.996102in}}%
\pgfpathlineto{\pgfqpoint{0.721498in}{0.983641in}}%
\pgfpathlineto{\pgfqpoint{0.737956in}{0.971704in}}%
\pgfpathlineto{\pgfqpoint{0.747008in}{0.968762in}}%
\pgfpathlineto{\pgfqpoint{0.756017in}{0.966091in}}%
\pgfpathlineto{\pgfqpoint{0.739605in}{0.978274in}}%
\pgfpathlineto{\pgfqpoint{0.723207in}{0.990986in}}%
\pgfpathlineto{\pgfqpoint{0.714152in}{0.993408in}}%
\pgfpathlineto{\pgfqpoint{0.705055in}{0.996102in}}%
\pgfpathclose%
\pgfusepath{fill}%
\end{pgfscope}%
\begin{pgfscope}%
\pgfpathrectangle{\pgfqpoint{0.407592in}{0.222904in}}{\pgfqpoint{1.411163in}{1.411163in}}%
\pgfusepath{clip}%
\pgfsetbuttcap%
\pgfsetroundjoin%
\definecolor{currentfill}{rgb}{0.909804,0.698039,0.423529}%
\pgfsetfillcolor{currentfill}%
\pgfsetfillopacity{0.120000}%
\pgfsetlinewidth{0.000000pt}%
\definecolor{currentstroke}{rgb}{0.000000,0.000000,0.000000}%
\pgfsetstrokecolor{currentstroke}%
\pgfsetdash{}{0pt}%
\pgfpathmoveto{\pgfqpoint{0.955886in}{0.869491in}}%
\pgfpathlineto{\pgfqpoint{0.972450in}{0.865049in}}%
\pgfpathlineto{\pgfqpoint{0.989049in}{0.861483in}}%
\pgfpathlineto{\pgfqpoint{0.997586in}{0.859667in}}%
\pgfpathlineto{\pgfqpoint{1.006086in}{0.858228in}}%
\pgfpathlineto{\pgfqpoint{0.989531in}{0.861757in}}%
\pgfpathlineto{\pgfqpoint{0.973011in}{0.866197in}}%
\pgfpathlineto{\pgfqpoint{0.964467in}{0.867656in}}%
\pgfpathlineto{\pgfqpoint{0.955886in}{0.869491in}}%
\pgfpathclose%
\pgfusepath{fill}%
\end{pgfscope}%
\begin{pgfscope}%
\pgfpathrectangle{\pgfqpoint{0.407592in}{0.222904in}}{\pgfqpoint{1.411163in}{1.411163in}}%
\pgfusepath{clip}%
\pgfsetbuttcap%
\pgfsetroundjoin%
\definecolor{currentfill}{rgb}{0.909804,0.698039,0.423529}%
\pgfsetfillcolor{currentfill}%
\pgfsetfillopacity{0.120000}%
\pgfsetlinewidth{0.000000pt}%
\definecolor{currentstroke}{rgb}{0.000000,0.000000,0.000000}%
\pgfsetstrokecolor{currentstroke}%
\pgfsetdash{}{0pt}%
\pgfpathmoveto{\pgfqpoint{1.572074in}{1.080578in}}%
\pgfpathlineto{\pgfqpoint{1.589745in}{1.090573in}}%
\pgfpathlineto{\pgfqpoint{1.607493in}{1.100822in}}%
\pgfpathlineto{\pgfqpoint{1.615095in}{1.109582in}}%
\pgfpathlineto{\pgfqpoint{1.622683in}{1.118607in}}%
\pgfpathlineto{\pgfqpoint{1.604958in}{1.108555in}}%
\pgfpathlineto{\pgfqpoint{1.587310in}{1.098764in}}%
\pgfpathlineto{\pgfqpoint{1.579698in}{1.089533in}}%
\pgfpathlineto{\pgfqpoint{1.572074in}{1.080578in}}%
\pgfpathclose%
\pgfusepath{fill}%
\end{pgfscope}%
\begin{pgfscope}%
\pgfpathrectangle{\pgfqpoint{0.407592in}{0.222904in}}{\pgfqpoint{1.411163in}{1.411163in}}%
\pgfusepath{clip}%
\pgfsetbuttcap%
\pgfsetroundjoin%
\definecolor{currentfill}{rgb}{0.909804,0.698039,0.423529}%
\pgfsetfillcolor{currentfill}%
\pgfsetfillopacity{0.120000}%
\pgfsetlinewidth{0.000000pt}%
\definecolor{currentstroke}{rgb}{0.000000,0.000000,0.000000}%
\pgfsetstrokecolor{currentstroke}%
\pgfsetdash{}{0pt}%
\pgfpathmoveto{\pgfqpoint{1.688501in}{1.177199in}}%
\pgfpathlineto{\pgfqpoint{1.706415in}{1.187543in}}%
\pgfpathlineto{\pgfqpoint{1.724408in}{1.198109in}}%
\pgfpathlineto{\pgfqpoint{1.731887in}{1.207906in}}%
\pgfpathlineto{\pgfqpoint{1.739354in}{1.217914in}}%
\pgfpathlineto{\pgfqpoint{1.721385in}{1.207567in}}%
\pgfpathlineto{\pgfqpoint{1.703495in}{1.197448in}}%
\pgfpathlineto{\pgfqpoint{1.696004in}{1.187215in}}%
\pgfpathlineto{\pgfqpoint{1.688501in}{1.177199in}}%
\pgfpathclose%
\pgfusepath{fill}%
\end{pgfscope}%
\begin{pgfscope}%
\pgfpathrectangle{\pgfqpoint{0.407592in}{0.222904in}}{\pgfqpoint{1.411163in}{1.411163in}}%
\pgfusepath{clip}%
\pgfsetbuttcap%
\pgfsetroundjoin%
\definecolor{currentfill}{rgb}{0.909804,0.698039,0.423529}%
\pgfsetfillcolor{currentfill}%
\pgfsetfillopacity{0.120000}%
\pgfsetlinewidth{0.000000pt}%
\definecolor{currentstroke}{rgb}{0.000000,0.000000,0.000000}%
\pgfsetstrokecolor{currentstroke}%
\pgfsetdash{}{0pt}%
\pgfpathmoveto{\pgfqpoint{1.189483in}{0.864945in}}%
\pgfpathlineto{\pgfqpoint{1.206370in}{0.869384in}}%
\pgfpathlineto{\pgfqpoint{1.223319in}{0.874486in}}%
\pgfpathlineto{\pgfqpoint{1.231418in}{0.876840in}}%
\pgfpathlineto{\pgfqpoint{1.239491in}{0.879655in}}%
\pgfpathlineto{\pgfqpoint{1.222571in}{0.874416in}}%
\pgfpathlineto{\pgfqpoint{1.205714in}{0.869841in}}%
\pgfpathlineto{\pgfqpoint{1.197611in}{0.867152in}}%
\pgfpathlineto{\pgfqpoint{1.189483in}{0.864945in}}%
\pgfpathclose%
\pgfusepath{fill}%
\end{pgfscope}%
\begin{pgfscope}%
\pgfpathrectangle{\pgfqpoint{0.407592in}{0.222904in}}{\pgfqpoint{1.411163in}{1.411163in}}%
\pgfusepath{clip}%
\pgfsetbuttcap%
\pgfsetroundjoin%
\definecolor{currentfill}{rgb}{0.909804,0.698039,0.423529}%
\pgfsetfillcolor{currentfill}%
\pgfsetfillopacity{0.120000}%
\pgfsetlinewidth{0.000000pt}%
\definecolor{currentstroke}{rgb}{0.000000,0.000000,0.000000}%
\pgfsetstrokecolor{currentstroke}%
\pgfsetdash{}{0pt}%
\pgfpathmoveto{\pgfqpoint{1.455844in}{0.993978in}}%
\pgfpathlineto{\pgfqpoint{1.473277in}{1.003300in}}%
\pgfpathlineto{\pgfqpoint{1.490786in}{1.012935in}}%
\pgfpathlineto{\pgfqpoint{1.498519in}{1.020205in}}%
\pgfpathlineto{\pgfqpoint{1.506237in}{1.027808in}}%
\pgfpathlineto{\pgfqpoint{1.488751in}{1.018310in}}%
\pgfpathlineto{\pgfqpoint{1.471340in}{1.009129in}}%
\pgfpathlineto{\pgfqpoint{1.463600in}{1.001378in}}%
\pgfpathlineto{\pgfqpoint{1.455844in}{0.993978in}}%
\pgfpathclose%
\pgfusepath{fill}%
\end{pgfscope}%
\begin{pgfscope}%
\pgfpathrectangle{\pgfqpoint{0.407592in}{0.222904in}}{\pgfqpoint{1.411163in}{1.411163in}}%
\pgfusepath{clip}%
\pgfsetbuttcap%
\pgfsetroundjoin%
\definecolor{currentfill}{rgb}{0.909804,0.698039,0.423529}%
\pgfsetfillcolor{currentfill}%
\pgfsetfillopacity{0.120000}%
\pgfsetlinewidth{0.000000pt}%
\definecolor{currentstroke}{rgb}{0.000000,0.000000,0.000000}%
\pgfsetstrokecolor{currentstroke}%
\pgfsetdash{}{0pt}%
\pgfpathmoveto{\pgfqpoint{1.139481in}{0.855560in}}%
\pgfpathlineto{\pgfqpoint{1.156282in}{0.858370in}}%
\pgfpathlineto{\pgfqpoint{1.173142in}{0.861930in}}%
\pgfpathlineto{\pgfqpoint{1.181326in}{0.863208in}}%
\pgfpathlineto{\pgfqpoint{1.189483in}{0.864945in}}%
\pgfpathlineto{\pgfqpoint{1.172656in}{0.861216in}}%
\pgfpathlineto{\pgfqpoint{1.155887in}{0.858243in}}%
\pgfpathlineto{\pgfqpoint{1.147699in}{0.856664in}}%
\pgfpathlineto{\pgfqpoint{1.139481in}{0.855560in}}%
\pgfpathclose%
\pgfusepath{fill}%
\end{pgfscope}%
\begin{pgfscope}%
\pgfpathrectangle{\pgfqpoint{0.407592in}{0.222904in}}{\pgfqpoint{1.411163in}{1.411163in}}%
\pgfusepath{clip}%
\pgfsetbuttcap%
\pgfsetroundjoin%
\definecolor{currentfill}{rgb}{0.909804,0.698039,0.423529}%
\pgfsetfillcolor{currentfill}%
\pgfsetfillopacity{0.120000}%
\pgfsetlinewidth{0.000000pt}%
\definecolor{currentstroke}{rgb}{0.000000,0.000000,0.000000}%
\pgfsetstrokecolor{currentstroke}%
\pgfsetdash{}{0pt}%
\pgfpathmoveto{\pgfqpoint{1.239491in}{0.879655in}}%
\pgfpathlineto{\pgfqpoint{1.256476in}{0.885511in}}%
\pgfpathlineto{\pgfqpoint{1.273528in}{0.891941in}}%
\pgfpathlineto{\pgfqpoint{1.281548in}{0.895325in}}%
\pgfpathlineto{\pgfqpoint{1.289546in}{0.899158in}}%
\pgfpathlineto{\pgfqpoint{1.272521in}{0.892640in}}%
\pgfpathlineto{\pgfqpoint{1.255564in}{0.886694in}}%
\pgfpathlineto{\pgfqpoint{1.247539in}{0.882937in}}%
\pgfpathlineto{\pgfqpoint{1.239491in}{0.879655in}}%
\pgfpathclose%
\pgfusepath{fill}%
\end{pgfscope}%
\begin{pgfscope}%
\pgfpathrectangle{\pgfqpoint{0.407592in}{0.222904in}}{\pgfqpoint{1.411163in}{1.411163in}}%
\pgfusepath{clip}%
\pgfsetbuttcap%
\pgfsetroundjoin%
\definecolor{currentfill}{rgb}{0.909804,0.698039,0.423529}%
\pgfsetfillcolor{currentfill}%
\pgfsetfillopacity{0.120000}%
\pgfsetlinewidth{0.000000pt}%
\definecolor{currentstroke}{rgb}{0.000000,0.000000,0.000000}%
\pgfsetstrokecolor{currentstroke}%
\pgfsetdash{}{0pt}%
\pgfpathmoveto{\pgfqpoint{1.089440in}{0.851845in}}%
\pgfpathlineto{\pgfqpoint{1.106171in}{0.852869in}}%
\pgfpathlineto{\pgfqpoint{1.122953in}{0.854710in}}%
\pgfpathlineto{\pgfqpoint{1.131232in}{0.854914in}}%
\pgfpathlineto{\pgfqpoint{1.139481in}{0.855560in}}%
\pgfpathlineto{\pgfqpoint{1.122734in}{0.853542in}}%
\pgfpathlineto{\pgfqpoint{1.106040in}{0.852358in}}%
\pgfpathlineto{\pgfqpoint{1.097756in}{0.851876in}}%
\pgfpathlineto{\pgfqpoint{1.089440in}{0.851845in}}%
\pgfpathclose%
\pgfusepath{fill}%
\end{pgfscope}%
\begin{pgfscope}%
\pgfpathrectangle{\pgfqpoint{0.407592in}{0.222904in}}{\pgfqpoint{1.411163in}{1.411163in}}%
\pgfusepath{clip}%
\pgfsetbuttcap%
\pgfsetroundjoin%
\definecolor{currentfill}{rgb}{0.909804,0.698039,0.423529}%
\pgfsetfillcolor{currentfill}%
\pgfsetfillopacity{0.120000}%
\pgfsetlinewidth{0.000000pt}%
\definecolor{currentstroke}{rgb}{0.000000,0.000000,0.000000}%
\pgfsetstrokecolor{currentstroke}%
\pgfsetdash{}{0pt}%
\pgfpathmoveto{\pgfqpoint{0.620906in}{1.055712in}}%
\pgfpathlineto{\pgfqpoint{0.637343in}{1.041707in}}%
\pgfpathlineto{\pgfqpoint{0.653793in}{1.028113in}}%
\pgfpathlineto{\pgfqpoint{0.663022in}{1.025169in}}%
\pgfpathlineto{\pgfqpoint{0.672208in}{1.022474in}}%
\pgfpathlineto{\pgfqpoint{0.655803in}{1.036328in}}%
\pgfpathlineto{\pgfqpoint{0.639410in}{1.050592in}}%
\pgfpathlineto{\pgfqpoint{0.630180in}{1.053028in}}%
\pgfpathlineto{\pgfqpoint{0.620906in}{1.055712in}}%
\pgfpathclose%
\pgfusepath{fill}%
\end{pgfscope}%
\begin{pgfscope}%
\pgfpathrectangle{\pgfqpoint{0.407592in}{0.222904in}}{\pgfqpoint{1.411163in}{1.411163in}}%
\pgfusepath{clip}%
\pgfsetbuttcap%
\pgfsetroundjoin%
\definecolor{currentfill}{rgb}{0.909804,0.698039,0.423529}%
\pgfsetfillcolor{currentfill}%
\pgfsetfillopacity{0.120000}%
\pgfsetlinewidth{0.000000pt}%
\definecolor{currentstroke}{rgb}{0.000000,0.000000,0.000000}%
\pgfsetstrokecolor{currentstroke}%
\pgfsetdash{}{0pt}%
\pgfpathmoveto{\pgfqpoint{1.289546in}{0.899158in}}%
\pgfpathlineto{\pgfqpoint{1.306639in}{0.906204in}}%
\pgfpathlineto{\pgfqpoint{1.323802in}{0.913739in}}%
\pgfpathlineto{\pgfqpoint{1.331751in}{0.918072in}}%
\pgfpathlineto{\pgfqpoint{1.339680in}{0.922834in}}%
\pgfpathlineto{\pgfqpoint{1.322543in}{0.915266in}}%
\pgfpathlineto{\pgfqpoint{1.305475in}{0.908184in}}%
\pgfpathlineto{\pgfqpoint{1.297521in}{0.903444in}}%
\pgfpathlineto{\pgfqpoint{1.289546in}{0.899158in}}%
\pgfpathclose%
\pgfusepath{fill}%
\end{pgfscope}%
\begin{pgfscope}%
\pgfpathrectangle{\pgfqpoint{0.407592in}{0.222904in}}{\pgfqpoint{1.411163in}{1.411163in}}%
\pgfusepath{clip}%
\pgfsetbuttcap%
\pgfsetroundjoin%
\definecolor{currentfill}{rgb}{0.909804,0.698039,0.423529}%
\pgfsetfillcolor{currentfill}%
\pgfsetfillopacity{0.120000}%
\pgfsetlinewidth{0.000000pt}%
\definecolor{currentstroke}{rgb}{0.000000,0.000000,0.000000}%
\pgfsetstrokecolor{currentstroke}%
\pgfsetdash{}{0pt}%
\pgfpathmoveto{\pgfqpoint{0.503564in}{1.151611in}}%
\pgfpathlineto{\pgfqpoint{0.519981in}{1.135867in}}%
\pgfpathlineto{\pgfqpoint{0.536408in}{1.120407in}}%
\pgfpathlineto{\pgfqpoint{0.545860in}{1.117763in}}%
\pgfpathlineto{\pgfqpoint{0.555267in}{1.115347in}}%
\pgfpathlineto{\pgfqpoint{0.538884in}{1.131051in}}%
\pgfpathlineto{\pgfqpoint{0.522510in}{1.147032in}}%
\pgfpathlineto{\pgfqpoint{0.513059in}{1.149208in}}%
\pgfpathlineto{\pgfqpoint{0.503564in}{1.151611in}}%
\pgfpathclose%
\pgfusepath{fill}%
\end{pgfscope}%
\begin{pgfscope}%
\pgfpathrectangle{\pgfqpoint{0.407592in}{0.222904in}}{\pgfqpoint{1.411163in}{1.411163in}}%
\pgfusepath{clip}%
\pgfsetbuttcap%
\pgfsetroundjoin%
\definecolor{currentfill}{rgb}{0.909804,0.698039,0.423529}%
\pgfsetfillcolor{currentfill}%
\pgfsetfillopacity{0.120000}%
\pgfsetlinewidth{0.000000pt}%
\definecolor{currentstroke}{rgb}{0.000000,0.000000,0.000000}%
\pgfsetstrokecolor{currentstroke}%
\pgfsetdash{}{0pt}%
\pgfpathmoveto{\pgfqpoint{0.821827in}{0.923305in}}%
\pgfpathlineto{\pgfqpoint{0.838326in}{0.914253in}}%
\pgfpathlineto{\pgfqpoint{0.854846in}{0.905919in}}%
\pgfpathlineto{\pgfqpoint{0.863675in}{0.902934in}}%
\pgfpathlineto{\pgfqpoint{0.872465in}{0.900249in}}%
\pgfpathlineto{\pgfqpoint{0.855991in}{0.908735in}}%
\pgfpathlineto{\pgfqpoint{0.839539in}{0.917962in}}%
\pgfpathlineto{\pgfqpoint{0.830703in}{0.920483in}}%
\pgfpathlineto{\pgfqpoint{0.821827in}{0.923305in}}%
\pgfpathclose%
\pgfusepath{fill}%
\end{pgfscope}%
\begin{pgfscope}%
\pgfpathrectangle{\pgfqpoint{0.407592in}{0.222904in}}{\pgfqpoint{1.411163in}{1.411163in}}%
\pgfusepath{clip}%
\pgfsetbuttcap%
\pgfsetroundjoin%
\definecolor{currentfill}{rgb}{0.909804,0.698039,0.423529}%
\pgfsetfillcolor{currentfill}%
\pgfsetfillopacity{0.120000}%
\pgfsetlinewidth{0.000000pt}%
\definecolor{currentstroke}{rgb}{0.000000,0.000000,0.000000}%
\pgfsetstrokecolor{currentstroke}%
\pgfsetdash{}{0pt}%
\pgfpathmoveto{\pgfqpoint{1.622683in}{1.118607in}}%
\pgfpathlineto{\pgfqpoint{1.640486in}{1.128906in}}%
\pgfpathlineto{\pgfqpoint{1.658368in}{1.139438in}}%
\pgfpathlineto{\pgfqpoint{1.665920in}{1.148520in}}%
\pgfpathlineto{\pgfqpoint{1.673460in}{1.157845in}}%
\pgfpathlineto{\pgfqpoint{1.655601in}{1.147516in}}%
\pgfpathlineto{\pgfqpoint{1.637821in}{1.137428in}}%
\pgfpathlineto{\pgfqpoint{1.630258in}{1.127892in}}%
\pgfpathlineto{\pgfqpoint{1.622683in}{1.118607in}}%
\pgfpathclose%
\pgfusepath{fill}%
\end{pgfscope}%
\begin{pgfscope}%
\pgfpathrectangle{\pgfqpoint{0.407592in}{0.222904in}}{\pgfqpoint{1.411163in}{1.411163in}}%
\pgfusepath{clip}%
\pgfsetbuttcap%
\pgfsetroundjoin%
\definecolor{currentfill}{rgb}{0.909804,0.698039,0.423529}%
\pgfsetfillcolor{currentfill}%
\pgfsetfillopacity{0.120000}%
\pgfsetlinewidth{0.000000pt}%
\definecolor{currentstroke}{rgb}{0.000000,0.000000,0.000000}%
\pgfsetstrokecolor{currentstroke}%
\pgfsetdash{}{0pt}%
\pgfpathmoveto{\pgfqpoint{0.905483in}{0.885606in}}%
\pgfpathlineto{\pgfqpoint{0.922031in}{0.879494in}}%
\pgfpathlineto{\pgfqpoint{0.938609in}{0.874211in}}%
\pgfpathlineto{\pgfqpoint{0.947267in}{0.871683in}}%
\pgfpathlineto{\pgfqpoint{0.955886in}{0.869491in}}%
\pgfpathlineto{\pgfqpoint{0.939353in}{0.874807in}}%
\pgfpathlineto{\pgfqpoint{0.922850in}{0.880983in}}%
\pgfpathlineto{\pgfqpoint{0.914186in}{0.883127in}}%
\pgfpathlineto{\pgfqpoint{0.905483in}{0.885606in}}%
\pgfpathclose%
\pgfusepath{fill}%
\end{pgfscope}%
\begin{pgfscope}%
\pgfpathrectangle{\pgfqpoint{0.407592in}{0.222904in}}{\pgfqpoint{1.411163in}{1.411163in}}%
\pgfusepath{clip}%
\pgfsetbuttcap%
\pgfsetroundjoin%
\definecolor{currentfill}{rgb}{0.909804,0.698039,0.423529}%
\pgfsetfillcolor{currentfill}%
\pgfsetfillopacity{0.120000}%
\pgfsetlinewidth{0.000000pt}%
\definecolor{currentstroke}{rgb}{0.000000,0.000000,0.000000}%
\pgfsetstrokecolor{currentstroke}%
\pgfsetdash{}{0pt}%
\pgfpathmoveto{\pgfqpoint{0.737956in}{0.971704in}}%
\pgfpathlineto{\pgfqpoint{0.754429in}{0.960321in}}%
\pgfpathlineto{\pgfqpoint{0.770919in}{0.949521in}}%
\pgfpathlineto{\pgfqpoint{0.779924in}{0.946347in}}%
\pgfpathlineto{\pgfqpoint{0.788888in}{0.943444in}}%
\pgfpathlineto{\pgfqpoint{0.772445in}{0.954470in}}%
\pgfpathlineto{\pgfqpoint{0.756017in}{0.966091in}}%
\pgfpathlineto{\pgfqpoint{0.747008in}{0.968762in}}%
\pgfpathlineto{\pgfqpoint{0.737956in}{0.971704in}}%
\pgfpathclose%
\pgfusepath{fill}%
\end{pgfscope}%
\begin{pgfscope}%
\pgfpathrectangle{\pgfqpoint{0.407592in}{0.222904in}}{\pgfqpoint{1.411163in}{1.411163in}}%
\pgfusepath{clip}%
\pgfsetbuttcap%
\pgfsetroundjoin%
\definecolor{currentfill}{rgb}{0.909804,0.698039,0.423529}%
\pgfsetfillcolor{currentfill}%
\pgfsetfillopacity{0.120000}%
\pgfsetlinewidth{0.000000pt}%
\definecolor{currentstroke}{rgb}{0.000000,0.000000,0.000000}%
\pgfsetstrokecolor{currentstroke}%
\pgfsetdash{}{0pt}%
\pgfpathmoveto{\pgfqpoint{1.506237in}{1.027808in}}%
\pgfpathlineto{\pgfqpoint{1.523800in}{1.037601in}}%
\pgfpathlineto{\pgfqpoint{1.541439in}{1.047671in}}%
\pgfpathlineto{\pgfqpoint{1.549119in}{1.055447in}}%
\pgfpathlineto{\pgfqpoint{1.556785in}{1.063528in}}%
\pgfpathlineto{\pgfqpoint{1.539168in}{1.053617in}}%
\pgfpathlineto{\pgfqpoint{1.521629in}{1.043989in}}%
\pgfpathlineto{\pgfqpoint{1.513940in}{1.035738in}}%
\pgfpathlineto{\pgfqpoint{1.506237in}{1.027808in}}%
\pgfpathclose%
\pgfusepath{fill}%
\end{pgfscope}%
\begin{pgfscope}%
\pgfpathrectangle{\pgfqpoint{0.407592in}{0.222904in}}{\pgfqpoint{1.411163in}{1.411163in}}%
\pgfusepath{clip}%
\pgfsetbuttcap%
\pgfsetroundjoin%
\definecolor{currentfill}{rgb}{0.909804,0.698039,0.423529}%
\pgfsetfillcolor{currentfill}%
\pgfsetfillopacity{0.120000}%
\pgfsetlinewidth{0.000000pt}%
\definecolor{currentstroke}{rgb}{0.000000,0.000000,0.000000}%
\pgfsetstrokecolor{currentstroke}%
\pgfsetdash{}{0pt}%
\pgfpathmoveto{\pgfqpoint{1.039311in}{0.853885in}}%
\pgfpathlineto{\pgfqpoint{1.055987in}{0.853044in}}%
\pgfpathlineto{\pgfqpoint{1.072708in}{0.853062in}}%
\pgfpathlineto{\pgfqpoint{1.081091in}{0.852247in}}%
\pgfpathlineto{\pgfqpoint{1.089440in}{0.851845in}}%
\pgfpathlineto{\pgfqpoint{1.072758in}{0.851675in}}%
\pgfpathlineto{\pgfqpoint{1.056122in}{0.852389in}}%
\pgfpathlineto{\pgfqpoint{1.047734in}{0.852928in}}%
\pgfpathlineto{\pgfqpoint{1.039311in}{0.853885in}}%
\pgfpathclose%
\pgfusepath{fill}%
\end{pgfscope}%
\begin{pgfscope}%
\pgfpathrectangle{\pgfqpoint{0.407592in}{0.222904in}}{\pgfqpoint{1.411163in}{1.411163in}}%
\pgfusepath{clip}%
\pgfsetbuttcap%
\pgfsetroundjoin%
\definecolor{currentfill}{rgb}{0.909804,0.698039,0.423529}%
\pgfsetfillcolor{currentfill}%
\pgfsetfillopacity{0.120000}%
\pgfsetlinewidth{0.000000pt}%
\definecolor{currentstroke}{rgb}{0.000000,0.000000,0.000000}%
\pgfsetstrokecolor{currentstroke}%
\pgfsetdash{}{0pt}%
\pgfpathmoveto{\pgfqpoint{1.339680in}{0.922834in}}%
\pgfpathlineto{\pgfqpoint{1.356889in}{0.930853in}}%
\pgfpathlineto{\pgfqpoint{1.374170in}{0.939287in}}%
\pgfpathlineto{\pgfqpoint{1.382054in}{0.944472in}}%
\pgfpathlineto{\pgfqpoint{1.389919in}{0.950059in}}%
\pgfpathlineto{\pgfqpoint{1.372662in}{0.941644in}}%
\pgfpathlineto{\pgfqpoint{1.355477in}{0.933643in}}%
\pgfpathlineto{\pgfqpoint{1.347589in}{0.928025in}}%
\pgfpathlineto{\pgfqpoint{1.339680in}{0.922834in}}%
\pgfpathclose%
\pgfusepath{fill}%
\end{pgfscope}%
\begin{pgfscope}%
\pgfpathrectangle{\pgfqpoint{0.407592in}{0.222904in}}{\pgfqpoint{1.411163in}{1.411163in}}%
\pgfusepath{clip}%
\pgfsetbuttcap%
\pgfsetroundjoin%
\definecolor{currentfill}{rgb}{0.909804,0.698039,0.423529}%
\pgfsetfillcolor{currentfill}%
\pgfsetfillopacity{0.120000}%
\pgfsetlinewidth{0.000000pt}%
\definecolor{currentstroke}{rgb}{0.000000,0.000000,0.000000}%
\pgfsetstrokecolor{currentstroke}%
\pgfsetdash{}{0pt}%
\pgfpathmoveto{\pgfqpoint{1.739354in}{1.217914in}}%
\pgfpathlineto{\pgfqpoint{1.757402in}{1.228479in}}%
\pgfpathlineto{\pgfqpoint{1.775531in}{1.239254in}}%
\pgfpathlineto{\pgfqpoint{1.782961in}{1.249249in}}%
\pgfpathlineto{\pgfqpoint{1.790381in}{1.259440in}}%
\pgfpathlineto{\pgfqpoint{1.772277in}{1.248883in}}%
\pgfpathlineto{\pgfqpoint{1.754253in}{1.238540in}}%
\pgfpathlineto{\pgfqpoint{1.746809in}{1.228127in}}%
\pgfpathlineto{\pgfqpoint{1.739354in}{1.217914in}}%
\pgfpathclose%
\pgfusepath{fill}%
\end{pgfscope}%
\begin{pgfscope}%
\pgfpathrectangle{\pgfqpoint{0.407592in}{0.222904in}}{\pgfqpoint{1.411163in}{1.411163in}}%
\pgfusepath{clip}%
\pgfsetbuttcap%
\pgfsetroundjoin%
\definecolor{currentfill}{rgb}{0.909804,0.698039,0.423529}%
\pgfsetfillcolor{currentfill}%
\pgfsetfillopacity{0.120000}%
\pgfsetlinewidth{0.000000pt}%
\definecolor{currentstroke}{rgb}{0.000000,0.000000,0.000000}%
\pgfsetstrokecolor{currentstroke}%
\pgfsetdash{}{0pt}%
\pgfpathmoveto{\pgfqpoint{1.790381in}{1.259440in}}%
\pgfpathlineto{\pgfqpoint{1.808565in}{1.270203in}}%
\pgfpathlineto{\pgfqpoint{1.815961in}{1.280474in}}%
\pgfpathlineto{\pgfqpoint{1.823345in}{1.290927in}}%
\pgfpathlineto{\pgfqpoint{1.805186in}{1.280383in}}%
\pgfpathlineto{\pgfqpoint{1.797789in}{1.269820in}}%
\pgfusepath{fill}%
\end{pgfscope}%
\begin{pgfscope}%
\pgfpathrectangle{\pgfqpoint{0.407592in}{0.222904in}}{\pgfqpoint{1.411163in}{1.411163in}}%
\pgfusepath{clip}%
\pgfsetbuttcap%
\pgfsetroundjoin%
\definecolor{currentfill}{rgb}{0.909804,0.698039,0.423529}%
\pgfsetfillcolor{currentfill}%
\pgfsetfillopacity{0.120000}%
\pgfsetlinewidth{0.000000pt}%
\definecolor{currentstroke}{rgb}{0.000000,0.000000,0.000000}%
\pgfsetstrokecolor{currentstroke}%
\pgfsetdash{}{0pt}%
\pgfpathmoveto{\pgfqpoint{0.536408in}{1.120407in}}%
\pgfpathlineto{\pgfqpoint{0.552846in}{1.105248in}}%
\pgfpathlineto{\pgfqpoint{0.569295in}{1.090412in}}%
\pgfpathlineto{\pgfqpoint{0.578702in}{1.087517in}}%
\pgfpathlineto{\pgfqpoint{0.588065in}{1.084851in}}%
\pgfpathlineto{\pgfqpoint{0.571661in}{1.099941in}}%
\pgfpathlineto{\pgfqpoint{0.555267in}{1.115347in}}%
\pgfpathlineto{\pgfqpoint{0.545860in}{1.117763in}}%
\pgfpathlineto{\pgfqpoint{0.536408in}{1.120407in}}%
\pgfpathclose%
\pgfusepath{fill}%
\end{pgfscope}%
\begin{pgfscope}%
\pgfpathrectangle{\pgfqpoint{0.407592in}{0.222904in}}{\pgfqpoint{1.411163in}{1.411163in}}%
\pgfusepath{clip}%
\pgfsetbuttcap%
\pgfsetroundjoin%
\definecolor{currentfill}{rgb}{0.909804,0.698039,0.423529}%
\pgfsetfillcolor{currentfill}%
\pgfsetfillopacity{0.120000}%
\pgfsetlinewidth{0.000000pt}%
\definecolor{currentstroke}{rgb}{0.000000,0.000000,0.000000}%
\pgfsetstrokecolor{currentstroke}%
\pgfsetdash{}{0pt}%
\pgfpathmoveto{\pgfqpoint{0.653793in}{1.028113in}}%
\pgfpathlineto{\pgfqpoint{0.670255in}{1.014954in}}%
\pgfpathlineto{\pgfqpoint{0.686731in}{1.002256in}}%
\pgfpathlineto{\pgfqpoint{0.695914in}{0.999055in}}%
\pgfpathlineto{\pgfqpoint{0.705055in}{0.996102in}}%
\pgfpathlineto{\pgfqpoint{0.688625in}{1.009055in}}%
\pgfpathlineto{\pgfqpoint{0.672208in}{1.022474in}}%
\pgfpathlineto{\pgfqpoint{0.663022in}{1.025169in}}%
\pgfpathlineto{\pgfqpoint{0.653793in}{1.028113in}}%
\pgfpathclose%
\pgfusepath{fill}%
\end{pgfscope}%
\begin{pgfscope}%
\pgfpathrectangle{\pgfqpoint{0.407592in}{0.222904in}}{\pgfqpoint{1.411163in}{1.411163in}}%
\pgfusepath{clip}%
\pgfsetbuttcap%
\pgfsetroundjoin%
\definecolor{currentfill}{rgb}{0.909804,0.698039,0.423529}%
\pgfsetfillcolor{currentfill}%
\pgfsetfillopacity{0.120000}%
\pgfsetlinewidth{0.000000pt}%
\definecolor{currentstroke}{rgb}{0.000000,0.000000,0.000000}%
\pgfsetstrokecolor{currentstroke}%
\pgfsetdash{}{0pt}%
\pgfpathmoveto{\pgfqpoint{1.389919in}{0.950059in}}%
\pgfpathlineto{\pgfqpoint{1.407250in}{0.958860in}}%
\pgfpathlineto{\pgfqpoint{1.424654in}{0.968017in}}%
\pgfpathlineto{\pgfqpoint{1.432477in}{0.973951in}}%
\pgfpathlineto{\pgfqpoint{1.440283in}{0.980259in}}%
\pgfpathlineto{\pgfqpoint{1.422902in}{0.971166in}}%
\pgfpathlineto{\pgfqpoint{1.405595in}{0.962430in}}%
\pgfpathlineto{\pgfqpoint{1.397766in}{0.956047in}}%
\pgfpathlineto{\pgfqpoint{1.389919in}{0.950059in}}%
\pgfpathclose%
\pgfusepath{fill}%
\end{pgfscope}%
\begin{pgfscope}%
\pgfpathrectangle{\pgfqpoint{0.407592in}{0.222904in}}{\pgfqpoint{1.411163in}{1.411163in}}%
\pgfusepath{clip}%
\pgfsetbuttcap%
\pgfsetroundjoin%
\definecolor{currentfill}{rgb}{0.909804,0.698039,0.423529}%
\pgfsetfillcolor{currentfill}%
\pgfsetfillopacity{0.120000}%
\pgfsetlinewidth{0.000000pt}%
\definecolor{currentstroke}{rgb}{0.000000,0.000000,0.000000}%
\pgfsetstrokecolor{currentstroke}%
\pgfsetdash{}{0pt}%
\pgfpathmoveto{\pgfqpoint{0.418690in}{1.221637in}}%
\pgfpathlineto{\pgfqpoint{0.435112in}{1.205141in}}%
\pgfpathlineto{\pgfqpoint{0.451544in}{1.188871in}}%
\pgfpathlineto{\pgfqpoint{0.461174in}{1.186270in}}%
\pgfpathlineto{\pgfqpoint{0.470759in}{1.183880in}}%
\pgfpathlineto{\pgfqpoint{0.454371in}{1.200373in}}%
\pgfpathlineto{\pgfqpoint{0.437993in}{1.217086in}}%
\pgfpathlineto{\pgfqpoint{0.428365in}{1.219257in}}%
\pgfpathlineto{\pgfqpoint{0.418690in}{1.221637in}}%
\pgfpathclose%
\pgfusepath{fill}%
\end{pgfscope}%
\begin{pgfscope}%
\pgfpathrectangle{\pgfqpoint{0.407592in}{0.222904in}}{\pgfqpoint{1.411163in}{1.411163in}}%
\pgfusepath{clip}%
\pgfsetbuttcap%
\pgfsetroundjoin%
\definecolor{currentfill}{rgb}{0.909804,0.698039,0.423529}%
\pgfsetfillcolor{currentfill}%
\pgfsetfillopacity{0.120000}%
\pgfsetlinewidth{0.000000pt}%
\definecolor{currentstroke}{rgb}{0.000000,0.000000,0.000000}%
\pgfsetstrokecolor{currentstroke}%
\pgfsetdash{}{0pt}%
\pgfpathmoveto{\pgfqpoint{0.989049in}{0.861483in}}%
\pgfpathlineto{\pgfqpoint{1.005684in}{0.858792in}}%
\pgfpathlineto{\pgfqpoint{1.022360in}{0.856968in}}%
\pgfpathlineto{\pgfqpoint{1.030853in}{0.855238in}}%
\pgfpathlineto{\pgfqpoint{1.039311in}{0.853885in}}%
\pgfpathlineto{\pgfqpoint{1.022679in}{0.855607in}}%
\pgfpathlineto{\pgfqpoint{1.006086in}{0.858228in}}%
\pgfpathlineto{\pgfqpoint{0.997586in}{0.859667in}}%
\pgfpathlineto{\pgfqpoint{0.989049in}{0.861483in}}%
\pgfpathclose%
\pgfusepath{fill}%
\end{pgfscope}%
\begin{pgfscope}%
\pgfpathrectangle{\pgfqpoint{0.407592in}{0.222904in}}{\pgfqpoint{1.411163in}{1.411163in}}%
\pgfusepath{clip}%
\pgfsetbuttcap%
\pgfsetroundjoin%
\definecolor{currentfill}{rgb}{0.909804,0.698039,0.423529}%
\pgfsetfillcolor{currentfill}%
\pgfsetfillopacity{0.120000}%
\pgfsetlinewidth{0.000000pt}%
\definecolor{currentstroke}{rgb}{0.000000,0.000000,0.000000}%
\pgfsetstrokecolor{currentstroke}%
\pgfsetdash{}{0pt}%
\pgfpathmoveto{\pgfqpoint{1.556785in}{1.063528in}}%
\pgfpathlineto{\pgfqpoint{1.574478in}{1.073703in}}%
\pgfpathlineto{\pgfqpoint{1.592250in}{1.084128in}}%
\pgfpathlineto{\pgfqpoint{1.599879in}{1.092336in}}%
\pgfpathlineto{\pgfqpoint{1.607493in}{1.100822in}}%
\pgfpathlineto{\pgfqpoint{1.589745in}{1.090573in}}%
\pgfpathlineto{\pgfqpoint{1.572074in}{1.080578in}}%
\pgfpathlineto{\pgfqpoint{1.564436in}{1.071907in}}%
\pgfpathlineto{\pgfqpoint{1.556785in}{1.063528in}}%
\pgfpathclose%
\pgfusepath{fill}%
\end{pgfscope}%
\begin{pgfscope}%
\pgfpathrectangle{\pgfqpoint{0.407592in}{0.222904in}}{\pgfqpoint{1.411163in}{1.411163in}}%
\pgfusepath{clip}%
\pgfsetbuttcap%
\pgfsetroundjoin%
\definecolor{currentfill}{rgb}{0.909804,0.698039,0.423529}%
\pgfsetfillcolor{currentfill}%
\pgfsetfillopacity{0.120000}%
\pgfsetlinewidth{0.000000pt}%
\definecolor{currentstroke}{rgb}{0.000000,0.000000,0.000000}%
\pgfsetstrokecolor{currentstroke}%
\pgfsetdash{}{0pt}%
\pgfpathmoveto{\pgfqpoint{1.673460in}{1.157845in}}%
\pgfpathlineto{\pgfqpoint{1.691397in}{1.168401in}}%
\pgfpathlineto{\pgfqpoint{1.709414in}{1.179174in}}%
\pgfpathlineto{\pgfqpoint{1.716917in}{1.188529in}}%
\pgfpathlineto{\pgfqpoint{1.724408in}{1.198109in}}%
\pgfpathlineto{\pgfqpoint{1.706415in}{1.187543in}}%
\pgfpathlineto{\pgfqpoint{1.688501in}{1.177199in}}%
\pgfpathlineto{\pgfqpoint{1.680987in}{1.167407in}}%
\pgfpathlineto{\pgfqpoint{1.673460in}{1.157845in}}%
\pgfpathclose%
\pgfusepath{fill}%
\end{pgfscope}%
\begin{pgfscope}%
\pgfpathrectangle{\pgfqpoint{0.407592in}{0.222904in}}{\pgfqpoint{1.411163in}{1.411163in}}%
\pgfusepath{clip}%
\pgfsetbuttcap%
\pgfsetroundjoin%
\definecolor{currentfill}{rgb}{0.909804,0.698039,0.423529}%
\pgfsetfillcolor{currentfill}%
\pgfsetfillopacity{0.120000}%
\pgfsetlinewidth{0.000000pt}%
\definecolor{currentstroke}{rgb}{0.000000,0.000000,0.000000}%
\pgfsetstrokecolor{currentstroke}%
\pgfsetdash{}{0pt}%
\pgfpathmoveto{\pgfqpoint{0.854846in}{0.905919in}}%
\pgfpathlineto{\pgfqpoint{0.871389in}{0.898327in}}%
\pgfpathlineto{\pgfqpoint{0.887957in}{0.891498in}}%
\pgfpathlineto{\pgfqpoint{0.896740in}{0.888402in}}%
\pgfpathlineto{\pgfqpoint{0.905483in}{0.885606in}}%
\pgfpathlineto{\pgfqpoint{0.888961in}{0.892531in}}%
\pgfpathlineto{\pgfqpoint{0.872465in}{0.900249in}}%
\pgfpathlineto{\pgfqpoint{0.863675in}{0.902934in}}%
\pgfpathlineto{\pgfqpoint{0.854846in}{0.905919in}}%
\pgfpathclose%
\pgfusepath{fill}%
\end{pgfscope}%
\begin{pgfscope}%
\pgfpathrectangle{\pgfqpoint{0.407592in}{0.222904in}}{\pgfqpoint{1.411163in}{1.411163in}}%
\pgfusepath{clip}%
\pgfsetbuttcap%
\pgfsetroundjoin%
\definecolor{currentfill}{rgb}{0.909804,0.698039,0.423529}%
\pgfsetfillcolor{currentfill}%
\pgfsetfillopacity{0.120000}%
\pgfsetlinewidth{0.000000pt}%
\definecolor{currentstroke}{rgb}{0.000000,0.000000,0.000000}%
\pgfsetstrokecolor{currentstroke}%
\pgfsetdash{}{0pt}%
\pgfpathmoveto{\pgfqpoint{1.440283in}{0.980259in}}%
\pgfpathlineto{\pgfqpoint{1.457739in}{0.989684in}}%
\pgfpathlineto{\pgfqpoint{1.475271in}{0.999420in}}%
\pgfpathlineto{\pgfqpoint{1.483037in}{1.006005in}}%
\pgfpathlineto{\pgfqpoint{1.490786in}{1.012935in}}%
\pgfpathlineto{\pgfqpoint{1.473277in}{1.003300in}}%
\pgfpathlineto{\pgfqpoint{1.455844in}{0.993978in}}%
\pgfpathlineto{\pgfqpoint{1.448072in}{0.986936in}}%
\pgfpathlineto{\pgfqpoint{1.440283in}{0.980259in}}%
\pgfpathclose%
\pgfusepath{fill}%
\end{pgfscope}%
\begin{pgfscope}%
\pgfpathrectangle{\pgfqpoint{0.407592in}{0.222904in}}{\pgfqpoint{1.411163in}{1.411163in}}%
\pgfusepath{clip}%
\pgfsetbuttcap%
\pgfsetroundjoin%
\definecolor{currentfill}{rgb}{0.909804,0.698039,0.423529}%
\pgfsetfillcolor{currentfill}%
\pgfsetfillopacity{0.120000}%
\pgfsetlinewidth{0.000000pt}%
\definecolor{currentstroke}{rgb}{0.000000,0.000000,0.000000}%
\pgfsetstrokecolor{currentstroke}%
\pgfsetdash{}{0pt}%
\pgfpathmoveto{\pgfqpoint{0.770919in}{0.949521in}}%
\pgfpathlineto{\pgfqpoint{0.787425in}{0.939333in}}%
\pgfpathlineto{\pgfqpoint{0.803951in}{0.929787in}}%
\pgfpathlineto{\pgfqpoint{0.812910in}{0.926411in}}%
\pgfpathlineto{\pgfqpoint{0.821827in}{0.923305in}}%
\pgfpathlineto{\pgfqpoint{0.805348in}{0.933045in}}%
\pgfpathlineto{\pgfqpoint{0.788888in}{0.943444in}}%
\pgfpathlineto{\pgfqpoint{0.779924in}{0.946347in}}%
\pgfpathlineto{\pgfqpoint{0.770919in}{0.949521in}}%
\pgfpathclose%
\pgfusepath{fill}%
\end{pgfscope}%
\begin{pgfscope}%
\pgfpathrectangle{\pgfqpoint{0.407592in}{0.222904in}}{\pgfqpoint{1.411163in}{1.411163in}}%
\pgfusepath{clip}%
\pgfsetbuttcap%
\pgfsetroundjoin%
\definecolor{currentfill}{rgb}{0.909804,0.698039,0.423529}%
\pgfsetfillcolor{currentfill}%
\pgfsetfillopacity{0.120000}%
\pgfsetlinewidth{0.000000pt}%
\definecolor{currentstroke}{rgb}{0.000000,0.000000,0.000000}%
\pgfsetstrokecolor{currentstroke}%
\pgfsetdash{}{0pt}%
\pgfpathmoveto{\pgfqpoint{0.569295in}{1.090412in}}%
\pgfpathlineto{\pgfqpoint{0.585754in}{1.075918in}}%
\pgfpathlineto{\pgfqpoint{0.602225in}{1.061787in}}%
\pgfpathlineto{\pgfqpoint{0.611588in}{1.058634in}}%
\pgfpathlineto{\pgfqpoint{0.620906in}{1.055712in}}%
\pgfpathlineto{\pgfqpoint{0.604480in}{1.070101in}}%
\pgfpathlineto{\pgfqpoint{0.588065in}{1.084851in}}%
\pgfpathlineto{\pgfqpoint{0.578702in}{1.087517in}}%
\pgfpathlineto{\pgfqpoint{0.569295in}{1.090412in}}%
\pgfpathclose%
\pgfusepath{fill}%
\end{pgfscope}%
\begin{pgfscope}%
\pgfpathrectangle{\pgfqpoint{0.407592in}{0.222904in}}{\pgfqpoint{1.411163in}{1.411163in}}%
\pgfusepath{clip}%
\pgfsetbuttcap%
\pgfsetroundjoin%
\definecolor{currentfill}{rgb}{0.909804,0.698039,0.423529}%
\pgfsetfillcolor{currentfill}%
\pgfsetfillopacity{0.120000}%
\pgfsetlinewidth{0.000000pt}%
\definecolor{currentstroke}{rgb}{0.000000,0.000000,0.000000}%
\pgfsetstrokecolor{currentstroke}%
\pgfsetdash{}{0pt}%
\pgfpathmoveto{\pgfqpoint{0.451544in}{1.188871in}}%
\pgfpathlineto{\pgfqpoint{0.467985in}{1.172841in}}%
\pgfpathlineto{\pgfqpoint{0.484436in}{1.157067in}}%
\pgfpathlineto{\pgfqpoint{0.494023in}{1.154232in}}%
\pgfpathlineto{\pgfqpoint{0.503564in}{1.151611in}}%
\pgfpathlineto{\pgfqpoint{0.487157in}{1.167621in}}%
\pgfpathlineto{\pgfqpoint{0.470759in}{1.183880in}}%
\pgfpathlineto{\pgfqpoint{0.461174in}{1.186270in}}%
\pgfpathlineto{\pgfqpoint{0.451544in}{1.188871in}}%
\pgfpathclose%
\pgfusepath{fill}%
\end{pgfscope}%
\begin{pgfscope}%
\pgfpathrectangle{\pgfqpoint{0.407592in}{0.222904in}}{\pgfqpoint{1.411163in}{1.411163in}}%
\pgfusepath{clip}%
\pgfsetbuttcap%
\pgfsetroundjoin%
\definecolor{currentfill}{rgb}{0.909804,0.698039,0.423529}%
\pgfsetfillcolor{currentfill}%
\pgfsetfillopacity{0.120000}%
\pgfsetlinewidth{0.000000pt}%
\definecolor{currentstroke}{rgb}{0.000000,0.000000,0.000000}%
\pgfsetstrokecolor{currentstroke}%
\pgfsetdash{}{0pt}%
\pgfpathmoveto{\pgfqpoint{1.173142in}{0.861930in}}%
\pgfpathlineto{\pgfqpoint{1.190061in}{0.866194in}}%
\pgfpathlineto{\pgfqpoint{1.207042in}{0.871119in}}%
\pgfpathlineto{\pgfqpoint{1.215194in}{0.872582in}}%
\pgfpathlineto{\pgfqpoint{1.223319in}{0.874486in}}%
\pgfpathlineto{\pgfqpoint{1.206370in}{0.869384in}}%
\pgfpathlineto{\pgfqpoint{1.189483in}{0.864945in}}%
\pgfpathlineto{\pgfqpoint{1.181326in}{0.863208in}}%
\pgfpathlineto{\pgfqpoint{1.173142in}{0.861930in}}%
\pgfpathclose%
\pgfusepath{fill}%
\end{pgfscope}%
\begin{pgfscope}%
\pgfpathrectangle{\pgfqpoint{0.407592in}{0.222904in}}{\pgfqpoint{1.411163in}{1.411163in}}%
\pgfusepath{clip}%
\pgfsetbuttcap%
\pgfsetroundjoin%
\definecolor{currentfill}{rgb}{0.909804,0.698039,0.423529}%
\pgfsetfillcolor{currentfill}%
\pgfsetfillopacity{0.120000}%
\pgfsetlinewidth{0.000000pt}%
\definecolor{currentstroke}{rgb}{0.000000,0.000000,0.000000}%
\pgfsetstrokecolor{currentstroke}%
\pgfsetdash{}{0pt}%
\pgfpathmoveto{\pgfqpoint{1.223319in}{0.874486in}}%
\pgfpathlineto{\pgfqpoint{1.240333in}{0.880205in}}%
\pgfpathlineto{\pgfqpoint{1.257414in}{0.886499in}}%
\pgfpathlineto{\pgfqpoint{1.265483in}{0.889001in}}%
\pgfpathlineto{\pgfqpoint{1.273528in}{0.891941in}}%
\pgfpathlineto{\pgfqpoint{1.256476in}{0.885511in}}%
\pgfpathlineto{\pgfqpoint{1.239491in}{0.879655in}}%
\pgfpathlineto{\pgfqpoint{1.231418in}{0.876840in}}%
\pgfpathlineto{\pgfqpoint{1.223319in}{0.874486in}}%
\pgfpathclose%
\pgfusepath{fill}%
\end{pgfscope}%
\begin{pgfscope}%
\pgfpathrectangle{\pgfqpoint{0.407592in}{0.222904in}}{\pgfqpoint{1.411163in}{1.411163in}}%
\pgfusepath{clip}%
\pgfsetbuttcap%
\pgfsetroundjoin%
\definecolor{currentfill}{rgb}{0.909804,0.698039,0.423529}%
\pgfsetfillcolor{currentfill}%
\pgfsetfillopacity{0.120000}%
\pgfsetlinewidth{0.000000pt}%
\definecolor{currentstroke}{rgb}{0.000000,0.000000,0.000000}%
\pgfsetstrokecolor{currentstroke}%
\pgfsetdash{}{0pt}%
\pgfpathmoveto{\pgfqpoint{0.686731in}{1.002256in}}%
\pgfpathlineto{\pgfqpoint{0.703220in}{0.990048in}}%
\pgfpathlineto{\pgfqpoint{0.719724in}{0.978355in}}%
\pgfpathlineto{\pgfqpoint{0.728861in}{0.974906in}}%
\pgfpathlineto{\pgfqpoint{0.737956in}{0.971704in}}%
\pgfpathlineto{\pgfqpoint{0.721498in}{0.983641in}}%
\pgfpathlineto{\pgfqpoint{0.705055in}{0.996102in}}%
\pgfpathlineto{\pgfqpoint{0.695914in}{0.999055in}}%
\pgfpathlineto{\pgfqpoint{0.686731in}{1.002256in}}%
\pgfpathclose%
\pgfusepath{fill}%
\end{pgfscope}%
\begin{pgfscope}%
\pgfpathrectangle{\pgfqpoint{0.407592in}{0.222904in}}{\pgfqpoint{1.411163in}{1.411163in}}%
\pgfusepath{clip}%
\pgfsetbuttcap%
\pgfsetroundjoin%
\definecolor{currentfill}{rgb}{0.909804,0.698039,0.423529}%
\pgfsetfillcolor{currentfill}%
\pgfsetfillopacity{0.120000}%
\pgfsetlinewidth{0.000000pt}%
\definecolor{currentstroke}{rgb}{0.000000,0.000000,0.000000}%
\pgfsetstrokecolor{currentstroke}%
\pgfsetdash{}{0pt}%
\pgfpathmoveto{\pgfqpoint{0.938609in}{0.874211in}}%
\pgfpathlineto{\pgfqpoint{0.955219in}{0.869767in}}%
\pgfpathlineto{\pgfqpoint{0.971863in}{0.866165in}}%
\pgfpathlineto{\pgfqpoint{0.980475in}{0.863655in}}%
\pgfpathlineto{\pgfqpoint{0.989049in}{0.861483in}}%
\pgfpathlineto{\pgfqpoint{0.972450in}{0.865049in}}%
\pgfpathlineto{\pgfqpoint{0.955886in}{0.869491in}}%
\pgfpathlineto{\pgfqpoint{0.947267in}{0.871683in}}%
\pgfpathlineto{\pgfqpoint{0.938609in}{0.874211in}}%
\pgfpathclose%
\pgfusepath{fill}%
\end{pgfscope}%
\begin{pgfscope}%
\pgfpathrectangle{\pgfqpoint{0.407592in}{0.222904in}}{\pgfqpoint{1.411163in}{1.411163in}}%
\pgfusepath{clip}%
\pgfsetbuttcap%
\pgfsetroundjoin%
\definecolor{currentfill}{rgb}{0.909804,0.698039,0.423529}%
\pgfsetfillcolor{currentfill}%
\pgfsetfillopacity{0.120000}%
\pgfsetlinewidth{0.000000pt}%
\definecolor{currentstroke}{rgb}{0.000000,0.000000,0.000000}%
\pgfsetstrokecolor{currentstroke}%
\pgfsetdash{}{0pt}%
\pgfpathmoveto{\pgfqpoint{1.122953in}{0.854710in}}%
\pgfpathlineto{\pgfqpoint{1.139791in}{0.857332in}}%
\pgfpathlineto{\pgfqpoint{1.156685in}{0.860693in}}%
\pgfpathlineto{\pgfqpoint{1.164928in}{0.861096in}}%
\pgfpathlineto{\pgfqpoint{1.173142in}{0.861930in}}%
\pgfpathlineto{\pgfqpoint{1.156282in}{0.858370in}}%
\pgfpathlineto{\pgfqpoint{1.139481in}{0.855560in}}%
\pgfpathlineto{\pgfqpoint{1.131232in}{0.854914in}}%
\pgfpathlineto{\pgfqpoint{1.122953in}{0.854710in}}%
\pgfpathclose%
\pgfusepath{fill}%
\end{pgfscope}%
\begin{pgfscope}%
\pgfpathrectangle{\pgfqpoint{0.407592in}{0.222904in}}{\pgfqpoint{1.411163in}{1.411163in}}%
\pgfusepath{clip}%
\pgfsetbuttcap%
\pgfsetroundjoin%
\definecolor{currentfill}{rgb}{0.909804,0.698039,0.423529}%
\pgfsetfillcolor{currentfill}%
\pgfsetfillopacity{0.120000}%
\pgfsetlinewidth{0.000000pt}%
\definecolor{currentstroke}{rgb}{0.000000,0.000000,0.000000}%
\pgfsetstrokecolor{currentstroke}%
\pgfsetdash{}{0pt}%
\pgfpathmoveto{\pgfqpoint{1.607493in}{1.100822in}}%
\pgfpathlineto{\pgfqpoint{1.625320in}{1.111312in}}%
\pgfpathlineto{\pgfqpoint{1.643226in}{1.122030in}}%
\pgfpathlineto{\pgfqpoint{1.650804in}{1.130606in}}%
\pgfpathlineto{\pgfqpoint{1.658368in}{1.139438in}}%
\pgfpathlineto{\pgfqpoint{1.640486in}{1.128906in}}%
\pgfpathlineto{\pgfqpoint{1.622683in}{1.118607in}}%
\pgfpathlineto{\pgfqpoint{1.615095in}{1.109582in}}%
\pgfpathlineto{\pgfqpoint{1.607493in}{1.100822in}}%
\pgfpathclose%
\pgfusepath{fill}%
\end{pgfscope}%
\begin{pgfscope}%
\pgfpathrectangle{\pgfqpoint{0.407592in}{0.222904in}}{\pgfqpoint{1.411163in}{1.411163in}}%
\pgfusepath{clip}%
\pgfsetbuttcap%
\pgfsetroundjoin%
\definecolor{currentfill}{rgb}{0.909804,0.698039,0.423529}%
\pgfsetfillcolor{currentfill}%
\pgfsetfillopacity{0.120000}%
\pgfsetlinewidth{0.000000pt}%
\definecolor{currentstroke}{rgb}{0.000000,0.000000,0.000000}%
\pgfsetstrokecolor{currentstroke}%
\pgfsetdash{}{0pt}%
\pgfpathmoveto{\pgfqpoint{1.273528in}{0.891941in}}%
\pgfpathlineto{\pgfqpoint{1.290647in}{0.898901in}}%
\pgfpathlineto{\pgfqpoint{1.307836in}{0.906352in}}%
\pgfpathlineto{\pgfqpoint{1.315830in}{0.909833in}}%
\pgfpathlineto{\pgfqpoint{1.323802in}{0.913739in}}%
\pgfpathlineto{\pgfqpoint{1.306639in}{0.906204in}}%
\pgfpathlineto{\pgfqpoint{1.289546in}{0.899158in}}%
\pgfpathlineto{\pgfqpoint{1.281548in}{0.895325in}}%
\pgfpathlineto{\pgfqpoint{1.273528in}{0.891941in}}%
\pgfpathclose%
\pgfusepath{fill}%
\end{pgfscope}%
\begin{pgfscope}%
\pgfpathrectangle{\pgfqpoint{0.407592in}{0.222904in}}{\pgfqpoint{1.411163in}{1.411163in}}%
\pgfusepath{clip}%
\pgfsetbuttcap%
\pgfsetroundjoin%
\definecolor{currentfill}{rgb}{0.909804,0.698039,0.423529}%
\pgfsetfillcolor{currentfill}%
\pgfsetfillopacity{0.120000}%
\pgfsetlinewidth{0.000000pt}%
\definecolor{currentstroke}{rgb}{0.000000,0.000000,0.000000}%
\pgfsetstrokecolor{currentstroke}%
\pgfsetdash{}{0pt}%
\pgfpathmoveto{\pgfqpoint{1.724408in}{1.198109in}}%
\pgfpathlineto{\pgfqpoint{1.742481in}{1.208887in}}%
\pgfpathlineto{\pgfqpoint{1.760634in}{1.219869in}}%
\pgfpathlineto{\pgfqpoint{1.768088in}{1.229458in}}%
\pgfpathlineto{\pgfqpoint{1.775531in}{1.239254in}}%
\pgfpathlineto{\pgfqpoint{1.757402in}{1.228479in}}%
\pgfpathlineto{\pgfqpoint{1.739354in}{1.217914in}}%
\pgfpathlineto{\pgfqpoint{1.731887in}{1.207906in}}%
\pgfpathlineto{\pgfqpoint{1.724408in}{1.198109in}}%
\pgfpathclose%
\pgfusepath{fill}%
\end{pgfscope}%
\begin{pgfscope}%
\pgfpathrectangle{\pgfqpoint{0.407592in}{0.222904in}}{\pgfqpoint{1.411163in}{1.411163in}}%
\pgfusepath{clip}%
\pgfsetbuttcap%
\pgfsetroundjoin%
\definecolor{currentfill}{rgb}{0.909804,0.698039,0.423529}%
\pgfsetfillcolor{currentfill}%
\pgfsetfillopacity{0.120000}%
\pgfsetlinewidth{0.000000pt}%
\definecolor{currentstroke}{rgb}{0.000000,0.000000,0.000000}%
\pgfsetstrokecolor{currentstroke}%
\pgfsetdash{}{0pt}%
\pgfpathmoveto{\pgfqpoint{1.490786in}{1.012935in}}%
\pgfpathlineto{\pgfqpoint{1.508371in}{1.022861in}}%
\pgfpathlineto{\pgfqpoint{1.526034in}{1.033060in}}%
\pgfpathlineto{\pgfqpoint{1.533744in}{1.040207in}}%
\pgfpathlineto{\pgfqpoint{1.541439in}{1.047671in}}%
\pgfpathlineto{\pgfqpoint{1.523800in}{1.037601in}}%
\pgfpathlineto{\pgfqpoint{1.506237in}{1.027808in}}%
\pgfpathlineto{\pgfqpoint{1.498519in}{1.020205in}}%
\pgfpathlineto{\pgfqpoint{1.490786in}{1.012935in}}%
\pgfpathclose%
\pgfusepath{fill}%
\end{pgfscope}%
\begin{pgfscope}%
\pgfpathrectangle{\pgfqpoint{0.407592in}{0.222904in}}{\pgfqpoint{1.411163in}{1.411163in}}%
\pgfusepath{clip}%
\pgfsetbuttcap%
\pgfsetroundjoin%
\definecolor{currentfill}{rgb}{0.909804,0.698039,0.423529}%
\pgfsetfillcolor{currentfill}%
\pgfsetfillopacity{0.120000}%
\pgfsetlinewidth{0.000000pt}%
\definecolor{currentstroke}{rgb}{0.000000,0.000000,0.000000}%
\pgfsetstrokecolor{currentstroke}%
\pgfsetdash{}{0pt}%
\pgfpathmoveto{\pgfqpoint{1.072708in}{0.853062in}}%
\pgfpathlineto{\pgfqpoint{1.089478in}{0.853911in}}%
\pgfpathlineto{\pgfqpoint{1.106299in}{0.855560in}}%
\pgfpathlineto{\pgfqpoint{1.114643in}{0.854932in}}%
\pgfpathlineto{\pgfqpoint{1.122953in}{0.854710in}}%
\pgfpathlineto{\pgfqpoint{1.106171in}{0.852869in}}%
\pgfpathlineto{\pgfqpoint{1.089440in}{0.851845in}}%
\pgfpathlineto{\pgfqpoint{1.081091in}{0.852247in}}%
\pgfpathlineto{\pgfqpoint{1.072708in}{0.853062in}}%
\pgfpathclose%
\pgfusepath{fill}%
\end{pgfscope}%
\begin{pgfscope}%
\pgfpathrectangle{\pgfqpoint{0.407592in}{0.222904in}}{\pgfqpoint{1.411163in}{1.411163in}}%
\pgfusepath{clip}%
\pgfsetbuttcap%
\pgfsetroundjoin%
\definecolor{currentfill}{rgb}{0.909804,0.698039,0.423529}%
\pgfsetfillcolor{currentfill}%
\pgfsetfillopacity{0.120000}%
\pgfsetlinewidth{0.000000pt}%
\definecolor{currentstroke}{rgb}{0.000000,0.000000,0.000000}%
\pgfsetstrokecolor{currentstroke}%
\pgfsetdash{}{0pt}%
\pgfpathmoveto{\pgfqpoint{1.775531in}{1.239254in}}%
\pgfpathlineto{\pgfqpoint{1.793740in}{1.250228in}}%
\pgfpathlineto{\pgfqpoint{1.801158in}{1.260119in}}%
\pgfpathlineto{\pgfqpoint{1.808565in}{1.270203in}}%
\pgfpathlineto{\pgfqpoint{1.790381in}{1.259440in}}%
\pgfpathlineto{\pgfqpoint{1.782961in}{1.249249in}}%
\pgfusepath{fill}%
\end{pgfscope}%
\begin{pgfscope}%
\pgfpathrectangle{\pgfqpoint{0.407592in}{0.222904in}}{\pgfqpoint{1.411163in}{1.411163in}}%
\pgfusepath{clip}%
\pgfsetbuttcap%
\pgfsetroundjoin%
\definecolor{currentfill}{rgb}{0.909804,0.698039,0.423529}%
\pgfsetfillcolor{currentfill}%
\pgfsetfillopacity{0.120000}%
\pgfsetlinewidth{0.000000pt}%
\definecolor{currentstroke}{rgb}{0.000000,0.000000,0.000000}%
\pgfsetstrokecolor{currentstroke}%
\pgfsetdash{}{0pt}%
\pgfpathmoveto{\pgfqpoint{1.323802in}{0.913739in}}%
\pgfpathlineto{\pgfqpoint{1.341036in}{0.921725in}}%
\pgfpathlineto{\pgfqpoint{1.358341in}{0.930130in}}%
\pgfpathlineto{\pgfqpoint{1.366266in}{0.934506in}}%
\pgfpathlineto{\pgfqpoint{1.374170in}{0.939287in}}%
\pgfpathlineto{\pgfqpoint{1.356889in}{0.930853in}}%
\pgfpathlineto{\pgfqpoint{1.339680in}{0.922834in}}%
\pgfpathlineto{\pgfqpoint{1.331751in}{0.918072in}}%
\pgfpathlineto{\pgfqpoint{1.323802in}{0.913739in}}%
\pgfpathclose%
\pgfusepath{fill}%
\end{pgfscope}%
\begin{pgfscope}%
\pgfpathrectangle{\pgfqpoint{0.407592in}{0.222904in}}{\pgfqpoint{1.411163in}{1.411163in}}%
\pgfusepath{clip}%
\pgfsetbuttcap%
\pgfsetroundjoin%
\definecolor{currentfill}{rgb}{0.909804,0.698039,0.423529}%
\pgfsetfillcolor{currentfill}%
\pgfsetfillopacity{0.120000}%
\pgfsetlinewidth{0.000000pt}%
\definecolor{currentstroke}{rgb}{0.000000,0.000000,0.000000}%
\pgfsetstrokecolor{currentstroke}%
\pgfsetdash{}{0pt}%
\pgfpathmoveto{\pgfqpoint{0.602225in}{1.061787in}}%
\pgfpathlineto{\pgfqpoint{0.618708in}{1.048043in}}%
\pgfpathlineto{\pgfqpoint{0.635203in}{1.034708in}}%
\pgfpathlineto{\pgfqpoint{0.644520in}{1.031296in}}%
\pgfpathlineto{\pgfqpoint{0.653793in}{1.028113in}}%
\pgfpathlineto{\pgfqpoint{0.637343in}{1.041707in}}%
\pgfpathlineto{\pgfqpoint{0.620906in}{1.055712in}}%
\pgfpathlineto{\pgfqpoint{0.611588in}{1.058634in}}%
\pgfpathlineto{\pgfqpoint{0.602225in}{1.061787in}}%
\pgfpathclose%
\pgfusepath{fill}%
\end{pgfscope}%
\begin{pgfscope}%
\pgfpathrectangle{\pgfqpoint{0.407592in}{0.222904in}}{\pgfqpoint{1.411163in}{1.411163in}}%
\pgfusepath{clip}%
\pgfsetbuttcap%
\pgfsetroundjoin%
\definecolor{currentfill}{rgb}{0.909804,0.698039,0.423529}%
\pgfsetfillcolor{currentfill}%
\pgfsetfillopacity{0.120000}%
\pgfsetlinewidth{0.000000pt}%
\definecolor{currentstroke}{rgb}{0.000000,0.000000,0.000000}%
\pgfsetstrokecolor{currentstroke}%
\pgfsetdash{}{0pt}%
\pgfpathmoveto{\pgfqpoint{0.484436in}{1.157067in}}%
\pgfpathlineto{\pgfqpoint{0.500898in}{1.141564in}}%
\pgfpathlineto{\pgfqpoint{0.517369in}{1.126350in}}%
\pgfpathlineto{\pgfqpoint{0.526912in}{1.123271in}}%
\pgfpathlineto{\pgfqpoint{0.536408in}{1.120407in}}%
\pgfpathlineto{\pgfqpoint{0.519981in}{1.135867in}}%
\pgfpathlineto{\pgfqpoint{0.503564in}{1.151611in}}%
\pgfpathlineto{\pgfqpoint{0.494023in}{1.154232in}}%
\pgfpathlineto{\pgfqpoint{0.484436in}{1.157067in}}%
\pgfpathclose%
\pgfusepath{fill}%
\end{pgfscope}%
\begin{pgfscope}%
\pgfpathrectangle{\pgfqpoint{0.407592in}{0.222904in}}{\pgfqpoint{1.411163in}{1.411163in}}%
\pgfusepath{clip}%
\pgfsetbuttcap%
\pgfsetroundjoin%
\definecolor{currentfill}{rgb}{0.909804,0.698039,0.423529}%
\pgfsetfillcolor{currentfill}%
\pgfsetfillopacity{0.120000}%
\pgfsetlinewidth{0.000000pt}%
\definecolor{currentstroke}{rgb}{0.000000,0.000000,0.000000}%
\pgfsetstrokecolor{currentstroke}%
\pgfsetdash{}{0pt}%
\pgfpathmoveto{\pgfqpoint{0.803951in}{0.929787in}}%
\pgfpathlineto{\pgfqpoint{0.820497in}{0.920909in}}%
\pgfpathlineto{\pgfqpoint{0.837064in}{0.912725in}}%
\pgfpathlineto{\pgfqpoint{0.845975in}{0.909188in}}%
\pgfpathlineto{\pgfqpoint{0.854846in}{0.905919in}}%
\pgfpathlineto{\pgfqpoint{0.838326in}{0.914253in}}%
\pgfpathlineto{\pgfqpoint{0.821827in}{0.923305in}}%
\pgfpathlineto{\pgfqpoint{0.812910in}{0.926411in}}%
\pgfpathlineto{\pgfqpoint{0.803951in}{0.929787in}}%
\pgfpathclose%
\pgfusepath{fill}%
\end{pgfscope}%
\begin{pgfscope}%
\pgfpathrectangle{\pgfqpoint{0.407592in}{0.222904in}}{\pgfqpoint{1.411163in}{1.411163in}}%
\pgfusepath{clip}%
\pgfsetbuttcap%
\pgfsetroundjoin%
\definecolor{currentfill}{rgb}{0.909804,0.698039,0.423529}%
\pgfsetfillcolor{currentfill}%
\pgfsetfillopacity{0.120000}%
\pgfsetlinewidth{0.000000pt}%
\definecolor{currentstroke}{rgb}{0.000000,0.000000,0.000000}%
\pgfsetstrokecolor{currentstroke}%
\pgfsetdash{}{0pt}%
\pgfpathmoveto{\pgfqpoint{1.022360in}{0.856968in}}%
\pgfpathlineto{\pgfqpoint{1.039078in}{0.855995in}}%
\pgfpathlineto{\pgfqpoint{1.055841in}{0.855855in}}%
\pgfpathlineto{\pgfqpoint{1.064292in}{0.854271in}}%
\pgfpathlineto{\pgfqpoint{1.072708in}{0.853062in}}%
\pgfpathlineto{\pgfqpoint{1.055987in}{0.853044in}}%
\pgfpathlineto{\pgfqpoint{1.039311in}{0.853885in}}%
\pgfpathlineto{\pgfqpoint{1.030853in}{0.855238in}}%
\pgfpathlineto{\pgfqpoint{1.022360in}{0.856968in}}%
\pgfpathclose%
\pgfusepath{fill}%
\end{pgfscope}%
\begin{pgfscope}%
\pgfpathrectangle{\pgfqpoint{0.407592in}{0.222904in}}{\pgfqpoint{1.411163in}{1.411163in}}%
\pgfusepath{clip}%
\pgfsetbuttcap%
\pgfsetroundjoin%
\definecolor{currentfill}{rgb}{0.909804,0.698039,0.423529}%
\pgfsetfillcolor{currentfill}%
\pgfsetfillopacity{0.120000}%
\pgfsetlinewidth{0.000000pt}%
\definecolor{currentstroke}{rgb}{0.000000,0.000000,0.000000}%
\pgfsetstrokecolor{currentstroke}%
\pgfsetdash{}{0pt}%
\pgfpathmoveto{\pgfqpoint{1.658368in}{1.139438in}}%
\pgfpathlineto{\pgfqpoint{1.676330in}{1.150192in}}%
\pgfpathlineto{\pgfqpoint{1.694371in}{1.161157in}}%
\pgfpathlineto{\pgfqpoint{1.701899in}{1.170048in}}%
\pgfpathlineto{\pgfqpoint{1.709414in}{1.179174in}}%
\pgfpathlineto{\pgfqpoint{1.691397in}{1.168401in}}%
\pgfpathlineto{\pgfqpoint{1.673460in}{1.157845in}}%
\pgfpathlineto{\pgfqpoint{1.665920in}{1.148520in}}%
\pgfpathlineto{\pgfqpoint{1.658368in}{1.139438in}}%
\pgfpathclose%
\pgfusepath{fill}%
\end{pgfscope}%
\begin{pgfscope}%
\pgfpathrectangle{\pgfqpoint{0.407592in}{0.222904in}}{\pgfqpoint{1.411163in}{1.411163in}}%
\pgfusepath{clip}%
\pgfsetbuttcap%
\pgfsetroundjoin%
\definecolor{currentfill}{rgb}{0.909804,0.698039,0.423529}%
\pgfsetfillcolor{currentfill}%
\pgfsetfillopacity{0.120000}%
\pgfsetlinewidth{0.000000pt}%
\definecolor{currentstroke}{rgb}{0.000000,0.000000,0.000000}%
\pgfsetstrokecolor{currentstroke}%
\pgfsetdash{}{0pt}%
\pgfpathmoveto{\pgfqpoint{1.374170in}{0.939287in}}%
\pgfpathlineto{\pgfqpoint{1.391524in}{0.948107in}}%
\pgfpathlineto{\pgfqpoint{1.408953in}{0.957283in}}%
\pgfpathlineto{\pgfqpoint{1.416813in}{0.962460in}}%
\pgfpathlineto{\pgfqpoint{1.424654in}{0.968017in}}%
\pgfpathlineto{\pgfqpoint{1.407250in}{0.958860in}}%
\pgfpathlineto{\pgfqpoint{1.389919in}{0.950059in}}%
\pgfpathlineto{\pgfqpoint{1.382054in}{0.944472in}}%
\pgfpathlineto{\pgfqpoint{1.374170in}{0.939287in}}%
\pgfpathclose%
\pgfusepath{fill}%
\end{pgfscope}%
\begin{pgfscope}%
\pgfpathrectangle{\pgfqpoint{0.407592in}{0.222904in}}{\pgfqpoint{1.411163in}{1.411163in}}%
\pgfusepath{clip}%
\pgfsetbuttcap%
\pgfsetroundjoin%
\definecolor{currentfill}{rgb}{0.909804,0.698039,0.423529}%
\pgfsetfillcolor{currentfill}%
\pgfsetfillopacity{0.120000}%
\pgfsetlinewidth{0.000000pt}%
\definecolor{currentstroke}{rgb}{0.000000,0.000000,0.000000}%
\pgfsetstrokecolor{currentstroke}%
\pgfsetdash{}{0pt}%
\pgfpathmoveto{\pgfqpoint{0.887957in}{0.891498in}}%
\pgfpathlineto{\pgfqpoint{0.904553in}{0.885453in}}%
\pgfpathlineto{\pgfqpoint{0.921178in}{0.880203in}}%
\pgfpathlineto{\pgfqpoint{0.929913in}{0.877057in}}%
\pgfpathlineto{\pgfqpoint{0.938609in}{0.874211in}}%
\pgfpathlineto{\pgfqpoint{0.922031in}{0.879494in}}%
\pgfpathlineto{\pgfqpoint{0.905483in}{0.885606in}}%
\pgfpathlineto{\pgfqpoint{0.896740in}{0.888402in}}%
\pgfpathlineto{\pgfqpoint{0.887957in}{0.891498in}}%
\pgfpathclose%
\pgfusepath{fill}%
\end{pgfscope}%
\begin{pgfscope}%
\pgfpathrectangle{\pgfqpoint{0.407592in}{0.222904in}}{\pgfqpoint{1.411163in}{1.411163in}}%
\pgfusepath{clip}%
\pgfsetbuttcap%
\pgfsetroundjoin%
\definecolor{currentfill}{rgb}{0.909804,0.698039,0.423529}%
\pgfsetfillcolor{currentfill}%
\pgfsetfillopacity{0.120000}%
\pgfsetlinewidth{0.000000pt}%
\definecolor{currentstroke}{rgb}{0.000000,0.000000,0.000000}%
\pgfsetstrokecolor{currentstroke}%
\pgfsetdash{}{0pt}%
\pgfpathmoveto{\pgfqpoint{1.541439in}{1.047671in}}%
\pgfpathlineto{\pgfqpoint{1.559156in}{1.058001in}}%
\pgfpathlineto{\pgfqpoint{1.576951in}{1.068576in}}%
\pgfpathlineto{\pgfqpoint{1.584608in}{1.076206in}}%
\pgfpathlineto{\pgfqpoint{1.592250in}{1.084128in}}%
\pgfpathlineto{\pgfqpoint{1.574478in}{1.073703in}}%
\pgfpathlineto{\pgfqpoint{1.556785in}{1.063528in}}%
\pgfpathlineto{\pgfqpoint{1.549119in}{1.055447in}}%
\pgfpathlineto{\pgfqpoint{1.541439in}{1.047671in}}%
\pgfpathclose%
\pgfusepath{fill}%
\end{pgfscope}%
\begin{pgfscope}%
\pgfpathrectangle{\pgfqpoint{0.407592in}{0.222904in}}{\pgfqpoint{1.411163in}{1.411163in}}%
\pgfusepath{clip}%
\pgfsetbuttcap%
\pgfsetroundjoin%
\definecolor{currentfill}{rgb}{0.909804,0.698039,0.423529}%
\pgfsetfillcolor{currentfill}%
\pgfsetfillopacity{0.120000}%
\pgfsetlinewidth{0.000000pt}%
\definecolor{currentstroke}{rgb}{0.000000,0.000000,0.000000}%
\pgfsetstrokecolor{currentstroke}%
\pgfsetdash{}{0pt}%
\pgfpathmoveto{\pgfqpoint{0.719724in}{0.978355in}}%
\pgfpathlineto{\pgfqpoint{0.736244in}{0.967207in}}%
\pgfpathlineto{\pgfqpoint{0.752780in}{0.956630in}}%
\pgfpathlineto{\pgfqpoint{0.761871in}{0.952952in}}%
\pgfpathlineto{\pgfqpoint{0.770919in}{0.949521in}}%
\pgfpathlineto{\pgfqpoint{0.754429in}{0.960321in}}%
\pgfpathlineto{\pgfqpoint{0.737956in}{0.971704in}}%
\pgfpathlineto{\pgfqpoint{0.728861in}{0.974906in}}%
\pgfpathlineto{\pgfqpoint{0.719724in}{0.978355in}}%
\pgfpathclose%
\pgfusepath{fill}%
\end{pgfscope}%
\begin{pgfscope}%
\pgfpathrectangle{\pgfqpoint{0.407592in}{0.222904in}}{\pgfqpoint{1.411163in}{1.411163in}}%
\pgfusepath{clip}%
\pgfsetbuttcap%
\pgfsetroundjoin%
\definecolor{currentfill}{rgb}{0.909804,0.698039,0.423529}%
\pgfsetfillcolor{currentfill}%
\pgfsetfillopacity{0.120000}%
\pgfsetlinewidth{0.000000pt}%
\definecolor{currentstroke}{rgb}{0.000000,0.000000,0.000000}%
\pgfsetstrokecolor{currentstroke}%
\pgfsetdash{}{0pt}%
\pgfpathmoveto{\pgfqpoint{0.517369in}{1.126350in}}%
\pgfpathlineto{\pgfqpoint{0.533852in}{1.111442in}}%
\pgfpathlineto{\pgfqpoint{0.550345in}{1.096860in}}%
\pgfpathlineto{\pgfqpoint{0.559842in}{1.093529in}}%
\pgfpathlineto{\pgfqpoint{0.569295in}{1.090412in}}%
\pgfpathlineto{\pgfqpoint{0.552846in}{1.105248in}}%
\pgfpathlineto{\pgfqpoint{0.536408in}{1.120407in}}%
\pgfpathlineto{\pgfqpoint{0.526912in}{1.123271in}}%
\pgfpathlineto{\pgfqpoint{0.517369in}{1.126350in}}%
\pgfpathclose%
\pgfusepath{fill}%
\end{pgfscope}%
\begin{pgfscope}%
\pgfpathrectangle{\pgfqpoint{0.407592in}{0.222904in}}{\pgfqpoint{1.411163in}{1.411163in}}%
\pgfusepath{clip}%
\pgfsetbuttcap%
\pgfsetroundjoin%
\definecolor{currentfill}{rgb}{0.909804,0.698039,0.423529}%
\pgfsetfillcolor{currentfill}%
\pgfsetfillopacity{0.120000}%
\pgfsetlinewidth{0.000000pt}%
\definecolor{currentstroke}{rgb}{0.000000,0.000000,0.000000}%
\pgfsetstrokecolor{currentstroke}%
\pgfsetdash{}{0pt}%
\pgfpathmoveto{\pgfqpoint{1.424654in}{0.968017in}}%
\pgfpathlineto{\pgfqpoint{1.442134in}{0.977505in}}%
\pgfpathlineto{\pgfqpoint{1.459689in}{0.987302in}}%
\pgfpathlineto{\pgfqpoint{1.467489in}{0.993184in}}%
\pgfpathlineto{\pgfqpoint{1.475271in}{0.999420in}}%
\pgfpathlineto{\pgfqpoint{1.457739in}{0.989684in}}%
\pgfpathlineto{\pgfqpoint{1.440283in}{0.980259in}}%
\pgfpathlineto{\pgfqpoint{1.432477in}{0.973951in}}%
\pgfpathlineto{\pgfqpoint{1.424654in}{0.968017in}}%
\pgfpathclose%
\pgfusepath{fill}%
\end{pgfscope}%
\begin{pgfscope}%
\pgfpathrectangle{\pgfqpoint{0.407592in}{0.222904in}}{\pgfqpoint{1.411163in}{1.411163in}}%
\pgfusepath{clip}%
\pgfsetbuttcap%
\pgfsetroundjoin%
\definecolor{currentfill}{rgb}{0.909804,0.698039,0.423529}%
\pgfsetfillcolor{currentfill}%
\pgfsetfillopacity{0.120000}%
\pgfsetlinewidth{0.000000pt}%
\definecolor{currentstroke}{rgb}{0.000000,0.000000,0.000000}%
\pgfsetstrokecolor{currentstroke}%
\pgfsetdash{}{0pt}%
\pgfpathmoveto{\pgfqpoint{0.635203in}{1.034708in}}%
\pgfpathlineto{\pgfqpoint{0.651711in}{1.021807in}}%
\pgfpathlineto{\pgfqpoint{0.668232in}{1.009363in}}%
\pgfpathlineto{\pgfqpoint{0.677503in}{1.005696in}}%
\pgfpathlineto{\pgfqpoint{0.686731in}{1.002256in}}%
\pgfpathlineto{\pgfqpoint{0.670255in}{1.014954in}}%
\pgfpathlineto{\pgfqpoint{0.653793in}{1.028113in}}%
\pgfpathlineto{\pgfqpoint{0.644520in}{1.031296in}}%
\pgfpathlineto{\pgfqpoint{0.635203in}{1.034708in}}%
\pgfpathclose%
\pgfusepath{fill}%
\end{pgfscope}%
\begin{pgfscope}%
\pgfpathrectangle{\pgfqpoint{0.407592in}{0.222904in}}{\pgfqpoint{1.411163in}{1.411163in}}%
\pgfusepath{clip}%
\pgfsetbuttcap%
\pgfsetroundjoin%
\definecolor{currentfill}{rgb}{0.909804,0.698039,0.423529}%
\pgfsetfillcolor{currentfill}%
\pgfsetfillopacity{0.120000}%
\pgfsetlinewidth{0.000000pt}%
\definecolor{currentstroke}{rgb}{0.000000,0.000000,0.000000}%
\pgfsetstrokecolor{currentstroke}%
\pgfsetdash{}{0pt}%
\pgfpathmoveto{\pgfqpoint{0.971863in}{0.866165in}}%
\pgfpathlineto{\pgfqpoint{0.988543in}{0.863405in}}%
\pgfpathlineto{\pgfqpoint{1.005263in}{0.861478in}}%
\pgfpathlineto{\pgfqpoint{1.013830in}{0.859054in}}%
\pgfpathlineto{\pgfqpoint{1.022360in}{0.856968in}}%
\pgfpathlineto{\pgfqpoint{1.005684in}{0.858792in}}%
\pgfpathlineto{\pgfqpoint{0.989049in}{0.861483in}}%
\pgfpathlineto{\pgfqpoint{0.980475in}{0.863655in}}%
\pgfpathlineto{\pgfqpoint{0.971863in}{0.866165in}}%
\pgfpathclose%
\pgfusepath{fill}%
\end{pgfscope}%
\begin{pgfscope}%
\pgfpathrectangle{\pgfqpoint{0.407592in}{0.222904in}}{\pgfqpoint{1.411163in}{1.411163in}}%
\pgfusepath{clip}%
\pgfsetbuttcap%
\pgfsetroundjoin%
\definecolor{currentfill}{rgb}{0.909804,0.698039,0.423529}%
\pgfsetfillcolor{currentfill}%
\pgfsetfillopacity{0.120000}%
\pgfsetlinewidth{0.000000pt}%
\definecolor{currentstroke}{rgb}{0.000000,0.000000,0.000000}%
\pgfsetstrokecolor{currentstroke}%
\pgfsetdash{}{0pt}%
\pgfpathmoveto{\pgfqpoint{1.592250in}{1.084128in}}%
\pgfpathlineto{\pgfqpoint{1.610100in}{1.094788in}}%
\pgfpathlineto{\pgfqpoint{1.628030in}{1.105671in}}%
\pgfpathlineto{\pgfqpoint{1.635635in}{1.113717in}}%
\pgfpathlineto{\pgfqpoint{1.643226in}{1.122030in}}%
\pgfpathlineto{\pgfqpoint{1.625320in}{1.111312in}}%
\pgfpathlineto{\pgfqpoint{1.607493in}{1.100822in}}%
\pgfpathlineto{\pgfqpoint{1.599879in}{1.092336in}}%
\pgfpathlineto{\pgfqpoint{1.592250in}{1.084128in}}%
\pgfpathclose%
\pgfusepath{fill}%
\end{pgfscope}%
\begin{pgfscope}%
\pgfpathrectangle{\pgfqpoint{0.407592in}{0.222904in}}{\pgfqpoint{1.411163in}{1.411163in}}%
\pgfusepath{clip}%
\pgfsetbuttcap%
\pgfsetroundjoin%
\definecolor{currentfill}{rgb}{0.909804,0.698039,0.423529}%
\pgfsetfillcolor{currentfill}%
\pgfsetfillopacity{0.120000}%
\pgfsetlinewidth{0.000000pt}%
\definecolor{currentstroke}{rgb}{0.000000,0.000000,0.000000}%
\pgfsetstrokecolor{currentstroke}%
\pgfsetdash{}{0pt}%
\pgfpathmoveto{\pgfqpoint{1.709414in}{1.179174in}}%
\pgfpathlineto{\pgfqpoint{1.727511in}{1.190153in}}%
\pgfpathlineto{\pgfqpoint{1.745689in}{1.201330in}}%
\pgfpathlineto{\pgfqpoint{1.753168in}{1.210491in}}%
\pgfpathlineto{\pgfqpoint{1.760634in}{1.219869in}}%
\pgfpathlineto{\pgfqpoint{1.742481in}{1.208887in}}%
\pgfpathlineto{\pgfqpoint{1.724408in}{1.198109in}}%
\pgfpathlineto{\pgfqpoint{1.716917in}{1.188529in}}%
\pgfpathlineto{\pgfqpoint{1.709414in}{1.179174in}}%
\pgfpathclose%
\pgfusepath{fill}%
\end{pgfscope}%
\begin{pgfscope}%
\pgfpathrectangle{\pgfqpoint{0.407592in}{0.222904in}}{\pgfqpoint{1.411163in}{1.411163in}}%
\pgfusepath{clip}%
\pgfsetbuttcap%
\pgfsetroundjoin%
\definecolor{currentfill}{rgb}{0.909804,0.698039,0.423529}%
\pgfsetfillcolor{currentfill}%
\pgfsetfillopacity{0.120000}%
\pgfsetlinewidth{0.000000pt}%
\definecolor{currentstroke}{rgb}{0.000000,0.000000,0.000000}%
\pgfsetstrokecolor{currentstroke}%
\pgfsetdash{}{0pt}%
\pgfpathmoveto{\pgfqpoint{1.760634in}{1.219869in}}%
\pgfpathlineto{\pgfqpoint{1.778868in}{1.231045in}}%
\pgfpathlineto{\pgfqpoint{1.786310in}{1.240535in}}%
\pgfpathlineto{\pgfqpoint{1.793740in}{1.250228in}}%
\pgfpathlineto{\pgfqpoint{1.775531in}{1.239254in}}%
\pgfpathlineto{\pgfqpoint{1.768088in}{1.229458in}}%
\pgfusepath{fill}%
\end{pgfscope}%
\begin{pgfscope}%
\pgfpathrectangle{\pgfqpoint{0.407592in}{0.222904in}}{\pgfqpoint{1.411163in}{1.411163in}}%
\pgfusepath{clip}%
\pgfsetbuttcap%
\pgfsetroundjoin%
\definecolor{currentfill}{rgb}{0.909804,0.698039,0.423529}%
\pgfsetfillcolor{currentfill}%
\pgfsetfillopacity{0.120000}%
\pgfsetlinewidth{0.000000pt}%
\definecolor{currentstroke}{rgb}{0.000000,0.000000,0.000000}%
\pgfsetstrokecolor{currentstroke}%
\pgfsetdash{}{0pt}%
\pgfpathmoveto{\pgfqpoint{1.475271in}{0.999420in}}%
\pgfpathlineto{\pgfqpoint{1.492880in}{1.009444in}}%
\pgfpathlineto{\pgfqpoint{1.510566in}{1.019738in}}%
\pgfpathlineto{\pgfqpoint{1.518308in}{1.026235in}}%
\pgfpathlineto{\pgfqpoint{1.526034in}{1.033060in}}%
\pgfpathlineto{\pgfqpoint{1.508371in}{1.022861in}}%
\pgfpathlineto{\pgfqpoint{1.490786in}{1.012935in}}%
\pgfpathlineto{\pgfqpoint{1.483037in}{1.006005in}}%
\pgfpathlineto{\pgfqpoint{1.475271in}{0.999420in}}%
\pgfpathclose%
\pgfusepath{fill}%
\end{pgfscope}%
\begin{pgfscope}%
\pgfpathrectangle{\pgfqpoint{0.407592in}{0.222904in}}{\pgfqpoint{1.411163in}{1.411163in}}%
\pgfusepath{clip}%
\pgfsetbuttcap%
\pgfsetroundjoin%
\definecolor{currentfill}{rgb}{0.909804,0.698039,0.423529}%
\pgfsetfillcolor{currentfill}%
\pgfsetfillopacity{0.120000}%
\pgfsetlinewidth{0.000000pt}%
\definecolor{currentstroke}{rgb}{0.000000,0.000000,0.000000}%
\pgfsetstrokecolor{currentstroke}%
\pgfsetdash{}{0pt}%
\pgfpathmoveto{\pgfqpoint{0.837064in}{0.912725in}}%
\pgfpathlineto{\pgfqpoint{0.853655in}{0.905257in}}%
\pgfpathlineto{\pgfqpoint{0.870271in}{0.898525in}}%
\pgfpathlineto{\pgfqpoint{0.879135in}{0.894878in}}%
\pgfpathlineto{\pgfqpoint{0.887957in}{0.891498in}}%
\pgfpathlineto{\pgfqpoint{0.871389in}{0.898327in}}%
\pgfpathlineto{\pgfqpoint{0.854846in}{0.905919in}}%
\pgfpathlineto{\pgfqpoint{0.845975in}{0.909188in}}%
\pgfpathlineto{\pgfqpoint{0.837064in}{0.912725in}}%
\pgfpathclose%
\pgfusepath{fill}%
\end{pgfscope}%
\begin{pgfscope}%
\pgfpathrectangle{\pgfqpoint{0.407592in}{0.222904in}}{\pgfqpoint{1.411163in}{1.411163in}}%
\pgfusepath{clip}%
\pgfsetbuttcap%
\pgfsetroundjoin%
\definecolor{currentfill}{rgb}{0.909804,0.698039,0.423529}%
\pgfsetfillcolor{currentfill}%
\pgfsetfillopacity{0.120000}%
\pgfsetlinewidth{0.000000pt}%
\definecolor{currentstroke}{rgb}{0.000000,0.000000,0.000000}%
\pgfsetstrokecolor{currentstroke}%
\pgfsetdash{}{0pt}%
\pgfpathmoveto{\pgfqpoint{1.207042in}{0.871119in}}%
\pgfpathlineto{\pgfqpoint{1.224087in}{0.876661in}}%
\pgfpathlineto{\pgfqpoint{1.241197in}{0.882778in}}%
\pgfpathlineto{\pgfqpoint{1.249318in}{0.884428in}}%
\pgfpathlineto{\pgfqpoint{1.257414in}{0.886499in}}%
\pgfpathlineto{\pgfqpoint{1.240333in}{0.880205in}}%
\pgfpathlineto{\pgfqpoint{1.223319in}{0.874486in}}%
\pgfpathlineto{\pgfqpoint{1.215194in}{0.872582in}}%
\pgfpathlineto{\pgfqpoint{1.207042in}{0.871119in}}%
\pgfpathclose%
\pgfusepath{fill}%
\end{pgfscope}%
\begin{pgfscope}%
\pgfpathrectangle{\pgfqpoint{0.407592in}{0.222904in}}{\pgfqpoint{1.411163in}{1.411163in}}%
\pgfusepath{clip}%
\pgfsetbuttcap%
\pgfsetroundjoin%
\definecolor{currentfill}{rgb}{0.909804,0.698039,0.423529}%
\pgfsetfillcolor{currentfill}%
\pgfsetfillopacity{0.120000}%
\pgfsetlinewidth{0.000000pt}%
\definecolor{currentstroke}{rgb}{0.000000,0.000000,0.000000}%
\pgfsetstrokecolor{currentstroke}%
\pgfsetdash{}{0pt}%
\pgfpathmoveto{\pgfqpoint{0.752780in}{0.956630in}}%
\pgfpathlineto{\pgfqpoint{0.769334in}{0.946650in}}%
\pgfpathlineto{\pgfqpoint{0.785907in}{0.937294in}}%
\pgfpathlineto{\pgfqpoint{0.794950in}{0.933419in}}%
\pgfpathlineto{\pgfqpoint{0.803951in}{0.929787in}}%
\pgfpathlineto{\pgfqpoint{0.787425in}{0.939333in}}%
\pgfpathlineto{\pgfqpoint{0.770919in}{0.949521in}}%
\pgfpathlineto{\pgfqpoint{0.761871in}{0.952952in}}%
\pgfpathlineto{\pgfqpoint{0.752780in}{0.956630in}}%
\pgfpathclose%
\pgfusepath{fill}%
\end{pgfscope}%
\begin{pgfscope}%
\pgfpathrectangle{\pgfqpoint{0.407592in}{0.222904in}}{\pgfqpoint{1.411163in}{1.411163in}}%
\pgfusepath{clip}%
\pgfsetbuttcap%
\pgfsetroundjoin%
\definecolor{currentfill}{rgb}{0.909804,0.698039,0.423529}%
\pgfsetfillcolor{currentfill}%
\pgfsetfillopacity{0.120000}%
\pgfsetlinewidth{0.000000pt}%
\definecolor{currentstroke}{rgb}{0.000000,0.000000,0.000000}%
\pgfsetstrokecolor{currentstroke}%
\pgfsetdash{}{0pt}%
\pgfpathmoveto{\pgfqpoint{1.156685in}{0.860693in}}%
\pgfpathlineto{\pgfqpoint{1.173638in}{0.864751in}}%
\pgfpathlineto{\pgfqpoint{1.190652in}{0.869466in}}%
\pgfpathlineto{\pgfqpoint{1.198861in}{0.870084in}}%
\pgfpathlineto{\pgfqpoint{1.207042in}{0.871119in}}%
\pgfpathlineto{\pgfqpoint{1.190061in}{0.866194in}}%
\pgfpathlineto{\pgfqpoint{1.173142in}{0.861930in}}%
\pgfpathlineto{\pgfqpoint{1.164928in}{0.861096in}}%
\pgfpathlineto{\pgfqpoint{1.156685in}{0.860693in}}%
\pgfpathclose%
\pgfusepath{fill}%
\end{pgfscope}%
\begin{pgfscope}%
\pgfpathrectangle{\pgfqpoint{0.407592in}{0.222904in}}{\pgfqpoint{1.411163in}{1.411163in}}%
\pgfusepath{clip}%
\pgfsetbuttcap%
\pgfsetroundjoin%
\definecolor{currentfill}{rgb}{0.909804,0.698039,0.423529}%
\pgfsetfillcolor{currentfill}%
\pgfsetfillopacity{0.120000}%
\pgfsetlinewidth{0.000000pt}%
\definecolor{currentstroke}{rgb}{0.000000,0.000000,0.000000}%
\pgfsetstrokecolor{currentstroke}%
\pgfsetdash{}{0pt}%
\pgfpathmoveto{\pgfqpoint{0.550345in}{1.096860in}}%
\pgfpathlineto{\pgfqpoint{0.566849in}{1.082622in}}%
\pgfpathlineto{\pgfqpoint{0.583365in}{1.068749in}}%
\pgfpathlineto{\pgfqpoint{0.592818in}{1.065161in}}%
\pgfpathlineto{\pgfqpoint{0.602225in}{1.061787in}}%
\pgfpathlineto{\pgfqpoint{0.585754in}{1.075918in}}%
\pgfpathlineto{\pgfqpoint{0.569295in}{1.090412in}}%
\pgfpathlineto{\pgfqpoint{0.559842in}{1.093529in}}%
\pgfpathlineto{\pgfqpoint{0.550345in}{1.096860in}}%
\pgfpathclose%
\pgfusepath{fill}%
\end{pgfscope}%
\begin{pgfscope}%
\pgfpathrectangle{\pgfqpoint{0.407592in}{0.222904in}}{\pgfqpoint{1.411163in}{1.411163in}}%
\pgfusepath{clip}%
\pgfsetbuttcap%
\pgfsetroundjoin%
\definecolor{currentfill}{rgb}{0.909804,0.698039,0.423529}%
\pgfsetfillcolor{currentfill}%
\pgfsetfillopacity{0.120000}%
\pgfsetlinewidth{0.000000pt}%
\definecolor{currentstroke}{rgb}{0.000000,0.000000,0.000000}%
\pgfsetstrokecolor{currentstroke}%
\pgfsetdash{}{0pt}%
\pgfpathmoveto{\pgfqpoint{1.257414in}{0.886499in}}%
\pgfpathlineto{\pgfqpoint{1.274561in}{0.893327in}}%
\pgfpathlineto{\pgfqpoint{1.291778in}{0.900649in}}%
\pgfpathlineto{\pgfqpoint{1.299819in}{0.903293in}}%
\pgfpathlineto{\pgfqpoint{1.307836in}{0.906352in}}%
\pgfpathlineto{\pgfqpoint{1.290647in}{0.898901in}}%
\pgfpathlineto{\pgfqpoint{1.273528in}{0.891941in}}%
\pgfpathlineto{\pgfqpoint{1.265483in}{0.889001in}}%
\pgfpathlineto{\pgfqpoint{1.257414in}{0.886499in}}%
\pgfpathclose%
\pgfusepath{fill}%
\end{pgfscope}%
\begin{pgfscope}%
\pgfpathrectangle{\pgfqpoint{0.407592in}{0.222904in}}{\pgfqpoint{1.411163in}{1.411163in}}%
\pgfusepath{clip}%
\pgfsetbuttcap%
\pgfsetroundjoin%
\definecolor{currentfill}{rgb}{0.909804,0.698039,0.423529}%
\pgfsetfillcolor{currentfill}%
\pgfsetfillopacity{0.120000}%
\pgfsetlinewidth{0.000000pt}%
\definecolor{currentstroke}{rgb}{0.000000,0.000000,0.000000}%
\pgfsetstrokecolor{currentstroke}%
\pgfsetdash{}{0pt}%
\pgfpathmoveto{\pgfqpoint{1.106299in}{0.855560in}}%
\pgfpathlineto{\pgfqpoint{1.123174in}{0.857973in}}%
\pgfpathlineto{\pgfqpoint{1.140105in}{0.861114in}}%
\pgfpathlineto{\pgfqpoint{1.148411in}{0.860704in}}%
\pgfpathlineto{\pgfqpoint{1.156685in}{0.860693in}}%
\pgfpathlineto{\pgfqpoint{1.139791in}{0.857332in}}%
\pgfpathlineto{\pgfqpoint{1.122953in}{0.854710in}}%
\pgfpathlineto{\pgfqpoint{1.114643in}{0.854932in}}%
\pgfpathlineto{\pgfqpoint{1.106299in}{0.855560in}}%
\pgfpathclose%
\pgfusepath{fill}%
\end{pgfscope}%
\begin{pgfscope}%
\pgfpathrectangle{\pgfqpoint{0.407592in}{0.222904in}}{\pgfqpoint{1.411163in}{1.411163in}}%
\pgfusepath{clip}%
\pgfsetbuttcap%
\pgfsetroundjoin%
\definecolor{currentfill}{rgb}{0.909804,0.698039,0.423529}%
\pgfsetfillcolor{currentfill}%
\pgfsetfillopacity{0.120000}%
\pgfsetlinewidth{0.000000pt}%
\definecolor{currentstroke}{rgb}{0.000000,0.000000,0.000000}%
\pgfsetstrokecolor{currentstroke}%
\pgfsetdash{}{0pt}%
\pgfpathmoveto{\pgfqpoint{0.921178in}{0.880203in}}%
\pgfpathlineto{\pgfqpoint{0.937834in}{0.875758in}}%
\pgfpathlineto{\pgfqpoint{0.954524in}{0.872122in}}%
\pgfpathlineto{\pgfqpoint{0.963213in}{0.868993in}}%
\pgfpathlineto{\pgfqpoint{0.971863in}{0.866165in}}%
\pgfpathlineto{\pgfqpoint{0.955219in}{0.869767in}}%
\pgfpathlineto{\pgfqpoint{0.938609in}{0.874211in}}%
\pgfpathlineto{\pgfqpoint{0.929913in}{0.877057in}}%
\pgfpathlineto{\pgfqpoint{0.921178in}{0.880203in}}%
\pgfpathclose%
\pgfusepath{fill}%
\end{pgfscope}%
\begin{pgfscope}%
\pgfpathrectangle{\pgfqpoint{0.407592in}{0.222904in}}{\pgfqpoint{1.411163in}{1.411163in}}%
\pgfusepath{clip}%
\pgfsetbuttcap%
\pgfsetroundjoin%
\definecolor{currentfill}{rgb}{0.909804,0.698039,0.423529}%
\pgfsetfillcolor{currentfill}%
\pgfsetfillopacity{0.120000}%
\pgfsetlinewidth{0.000000pt}%
\definecolor{currentstroke}{rgb}{0.000000,0.000000,0.000000}%
\pgfsetstrokecolor{currentstroke}%
\pgfsetdash{}{0pt}%
\pgfpathmoveto{\pgfqpoint{1.307836in}{0.906352in}}%
\pgfpathlineto{\pgfqpoint{1.325096in}{0.914259in}}%
\pgfpathlineto{\pgfqpoint{1.342428in}{0.922588in}}%
\pgfpathlineto{\pgfqpoint{1.350396in}{0.926158in}}%
\pgfpathlineto{\pgfqpoint{1.358341in}{0.930130in}}%
\pgfpathlineto{\pgfqpoint{1.341036in}{0.921725in}}%
\pgfpathlineto{\pgfqpoint{1.323802in}{0.913739in}}%
\pgfpathlineto{\pgfqpoint{1.315830in}{0.909833in}}%
\pgfpathlineto{\pgfqpoint{1.307836in}{0.906352in}}%
\pgfpathclose%
\pgfusepath{fill}%
\end{pgfscope}%
\begin{pgfscope}%
\pgfpathrectangle{\pgfqpoint{0.407592in}{0.222904in}}{\pgfqpoint{1.411163in}{1.411163in}}%
\pgfusepath{clip}%
\pgfsetbuttcap%
\pgfsetroundjoin%
\definecolor{currentfill}{rgb}{0.909804,0.698039,0.423529}%
\pgfsetfillcolor{currentfill}%
\pgfsetfillopacity{0.120000}%
\pgfsetlinewidth{0.000000pt}%
\definecolor{currentstroke}{rgb}{0.000000,0.000000,0.000000}%
\pgfsetstrokecolor{currentstroke}%
\pgfsetdash{}{0pt}%
\pgfpathmoveto{\pgfqpoint{0.668232in}{1.009363in}}%
\pgfpathlineto{\pgfqpoint{0.684768in}{0.997403in}}%
\pgfpathlineto{\pgfqpoint{0.701319in}{0.985951in}}%
\pgfpathlineto{\pgfqpoint{0.710543in}{0.982041in}}%
\pgfpathlineto{\pgfqpoint{0.719724in}{0.978355in}}%
\pgfpathlineto{\pgfqpoint{0.703220in}{0.990048in}}%
\pgfpathlineto{\pgfqpoint{0.686731in}{1.002256in}}%
\pgfpathlineto{\pgfqpoint{0.677503in}{1.005696in}}%
\pgfpathlineto{\pgfqpoint{0.668232in}{1.009363in}}%
\pgfpathclose%
\pgfusepath{fill}%
\end{pgfscope}%
\begin{pgfscope}%
\pgfpathrectangle{\pgfqpoint{0.407592in}{0.222904in}}{\pgfqpoint{1.411163in}{1.411163in}}%
\pgfusepath{clip}%
\pgfsetbuttcap%
\pgfsetroundjoin%
\definecolor{currentfill}{rgb}{0.909804,0.698039,0.423529}%
\pgfsetfillcolor{currentfill}%
\pgfsetfillopacity{0.120000}%
\pgfsetlinewidth{0.000000pt}%
\definecolor{currentstroke}{rgb}{0.000000,0.000000,0.000000}%
\pgfsetstrokecolor{currentstroke}%
\pgfsetdash{}{0pt}%
\pgfpathmoveto{\pgfqpoint{1.643226in}{1.122030in}}%
\pgfpathlineto{\pgfqpoint{1.661211in}{1.132964in}}%
\pgfpathlineto{\pgfqpoint{1.679277in}{1.144104in}}%
\pgfpathlineto{\pgfqpoint{1.686830in}{1.152507in}}%
\pgfpathlineto{\pgfqpoint{1.694371in}{1.161157in}}%
\pgfpathlineto{\pgfqpoint{1.676330in}{1.150192in}}%
\pgfpathlineto{\pgfqpoint{1.658368in}{1.139438in}}%
\pgfpathlineto{\pgfqpoint{1.650804in}{1.130606in}}%
\pgfpathlineto{\pgfqpoint{1.643226in}{1.122030in}}%
\pgfpathclose%
\pgfusepath{fill}%
\end{pgfscope}%
\begin{pgfscope}%
\pgfpathrectangle{\pgfqpoint{0.407592in}{0.222904in}}{\pgfqpoint{1.411163in}{1.411163in}}%
\pgfusepath{clip}%
\pgfsetbuttcap%
\pgfsetroundjoin%
\definecolor{currentfill}{rgb}{0.909804,0.698039,0.423529}%
\pgfsetfillcolor{currentfill}%
\pgfsetfillopacity{0.120000}%
\pgfsetlinewidth{0.000000pt}%
\definecolor{currentstroke}{rgb}{0.000000,0.000000,0.000000}%
\pgfsetstrokecolor{currentstroke}%
\pgfsetdash{}{0pt}%
\pgfpathmoveto{\pgfqpoint{1.055841in}{0.855855in}}%
\pgfpathlineto{\pgfqpoint{1.072652in}{0.856521in}}%
\pgfpathlineto{\pgfqpoint{1.089513in}{0.857967in}}%
\pgfpathlineto{\pgfqpoint{1.097923in}{0.856578in}}%
\pgfpathlineto{\pgfqpoint{1.106299in}{0.855560in}}%
\pgfpathlineto{\pgfqpoint{1.089478in}{0.853911in}}%
\pgfpathlineto{\pgfqpoint{1.072708in}{0.853062in}}%
\pgfpathlineto{\pgfqpoint{1.064292in}{0.854271in}}%
\pgfpathlineto{\pgfqpoint{1.055841in}{0.855855in}}%
\pgfpathclose%
\pgfusepath{fill}%
\end{pgfscope}%
\begin{pgfscope}%
\pgfpathrectangle{\pgfqpoint{0.407592in}{0.222904in}}{\pgfqpoint{1.411163in}{1.411163in}}%
\pgfusepath{clip}%
\pgfsetbuttcap%
\pgfsetroundjoin%
\definecolor{currentfill}{rgb}{0.909804,0.698039,0.423529}%
\pgfsetfillcolor{currentfill}%
\pgfsetfillopacity{0.120000}%
\pgfsetlinewidth{0.000000pt}%
\definecolor{currentstroke}{rgb}{0.000000,0.000000,0.000000}%
\pgfsetstrokecolor{currentstroke}%
\pgfsetdash{}{0pt}%
\pgfpathmoveto{\pgfqpoint{1.526034in}{1.033060in}}%
\pgfpathlineto{\pgfqpoint{1.543774in}{1.043515in}}%
\pgfpathlineto{\pgfqpoint{1.561593in}{1.054211in}}%
\pgfpathlineto{\pgfqpoint{1.569279in}{1.061242in}}%
\pgfpathlineto{\pgfqpoint{1.576951in}{1.068576in}}%
\pgfpathlineto{\pgfqpoint{1.559156in}{1.058001in}}%
\pgfpathlineto{\pgfqpoint{1.541439in}{1.047671in}}%
\pgfpathlineto{\pgfqpoint{1.533744in}{1.040207in}}%
\pgfpathlineto{\pgfqpoint{1.526034in}{1.033060in}}%
\pgfpathclose%
\pgfusepath{fill}%
\end{pgfscope}%
\begin{pgfscope}%
\pgfpathrectangle{\pgfqpoint{0.407592in}{0.222904in}}{\pgfqpoint{1.411163in}{1.411163in}}%
\pgfusepath{clip}%
\pgfsetbuttcap%
\pgfsetroundjoin%
\definecolor{currentfill}{rgb}{0.909804,0.698039,0.423529}%
\pgfsetfillcolor{currentfill}%
\pgfsetfillopacity{0.120000}%
\pgfsetlinewidth{0.000000pt}%
\definecolor{currentstroke}{rgb}{0.000000,0.000000,0.000000}%
\pgfsetstrokecolor{currentstroke}%
\pgfsetdash{}{0pt}%
\pgfpathmoveto{\pgfqpoint{1.358341in}{0.930130in}}%
\pgfpathlineto{\pgfqpoint{1.375721in}{0.938923in}}%
\pgfpathlineto{\pgfqpoint{1.393174in}{0.948074in}}%
\pgfpathlineto{\pgfqpoint{1.401073in}{0.952488in}}%
\pgfpathlineto{\pgfqpoint{1.408953in}{0.957283in}}%
\pgfpathlineto{\pgfqpoint{1.391524in}{0.948107in}}%
\pgfpathlineto{\pgfqpoint{1.374170in}{0.939287in}}%
\pgfpathlineto{\pgfqpoint{1.366266in}{0.934506in}}%
\pgfpathlineto{\pgfqpoint{1.358341in}{0.930130in}}%
\pgfpathclose%
\pgfusepath{fill}%
\end{pgfscope}%
\begin{pgfscope}%
\pgfpathrectangle{\pgfqpoint{0.407592in}{0.222904in}}{\pgfqpoint{1.411163in}{1.411163in}}%
\pgfusepath{clip}%
\pgfsetbuttcap%
\pgfsetroundjoin%
\definecolor{currentfill}{rgb}{0.909804,0.698039,0.423529}%
\pgfsetfillcolor{currentfill}%
\pgfsetfillopacity{0.120000}%
\pgfsetlinewidth{0.000000pt}%
\definecolor{currentstroke}{rgb}{0.000000,0.000000,0.000000}%
\pgfsetstrokecolor{currentstroke}%
\pgfsetdash{}{0pt}%
\pgfpathmoveto{\pgfqpoint{0.583365in}{1.068749in}}%
\pgfpathlineto{\pgfqpoint{0.599894in}{1.055263in}}%
\pgfpathlineto{\pgfqpoint{0.616435in}{1.042185in}}%
\pgfpathlineto{\pgfqpoint{0.625841in}{1.038340in}}%
\pgfpathlineto{\pgfqpoint{0.635203in}{1.034708in}}%
\pgfpathlineto{\pgfqpoint{0.618708in}{1.048043in}}%
\pgfpathlineto{\pgfqpoint{0.602225in}{1.061787in}}%
\pgfpathlineto{\pgfqpoint{0.592818in}{1.065161in}}%
\pgfpathlineto{\pgfqpoint{0.583365in}{1.068749in}}%
\pgfpathclose%
\pgfusepath{fill}%
\end{pgfscope}%
\begin{pgfscope}%
\pgfpathrectangle{\pgfqpoint{0.407592in}{0.222904in}}{\pgfqpoint{1.411163in}{1.411163in}}%
\pgfusepath{clip}%
\pgfsetbuttcap%
\pgfsetroundjoin%
\definecolor{currentfill}{rgb}{0.909804,0.698039,0.423529}%
\pgfsetfillcolor{currentfill}%
\pgfsetfillopacity{0.120000}%
\pgfsetlinewidth{0.000000pt}%
\definecolor{currentstroke}{rgb}{0.000000,0.000000,0.000000}%
\pgfsetstrokecolor{currentstroke}%
\pgfsetdash{}{0pt}%
\pgfpathmoveto{\pgfqpoint{1.005263in}{0.861478in}}%
\pgfpathlineto{\pgfqpoint{1.022025in}{0.860373in}}%
\pgfpathlineto{\pgfqpoint{1.038832in}{0.860072in}}%
\pgfpathlineto{\pgfqpoint{1.047354in}{0.857795in}}%
\pgfpathlineto{\pgfqpoint{1.055841in}{0.855855in}}%
\pgfpathlineto{\pgfqpoint{1.039078in}{0.855995in}}%
\pgfpathlineto{\pgfqpoint{1.022360in}{0.856968in}}%
\pgfpathlineto{\pgfqpoint{1.013830in}{0.859054in}}%
\pgfpathlineto{\pgfqpoint{1.005263in}{0.861478in}}%
\pgfpathclose%
\pgfusepath{fill}%
\end{pgfscope}%
\begin{pgfscope}%
\pgfpathrectangle{\pgfqpoint{0.407592in}{0.222904in}}{\pgfqpoint{1.411163in}{1.411163in}}%
\pgfusepath{clip}%
\pgfsetbuttcap%
\pgfsetroundjoin%
\definecolor{currentfill}{rgb}{0.909804,0.698039,0.423529}%
\pgfsetfillcolor{currentfill}%
\pgfsetfillopacity{0.120000}%
\pgfsetlinewidth{0.000000pt}%
\definecolor{currentstroke}{rgb}{0.000000,0.000000,0.000000}%
\pgfsetstrokecolor{currentstroke}%
\pgfsetdash{}{0pt}%
\pgfpathmoveto{\pgfqpoint{0.785907in}{0.937294in}}%
\pgfpathlineto{\pgfqpoint{0.802501in}{0.928586in}}%
\pgfpathlineto{\pgfqpoint{0.819116in}{0.920549in}}%
\pgfpathlineto{\pgfqpoint{0.828111in}{0.916517in}}%
\pgfpathlineto{\pgfqpoint{0.837064in}{0.912725in}}%
\pgfpathlineto{\pgfqpoint{0.820497in}{0.920909in}}%
\pgfpathlineto{\pgfqpoint{0.803951in}{0.929787in}}%
\pgfpathlineto{\pgfqpoint{0.794950in}{0.933419in}}%
\pgfpathlineto{\pgfqpoint{0.785907in}{0.937294in}}%
\pgfpathclose%
\pgfusepath{fill}%
\end{pgfscope}%
\begin{pgfscope}%
\pgfpathrectangle{\pgfqpoint{0.407592in}{0.222904in}}{\pgfqpoint{1.411163in}{1.411163in}}%
\pgfusepath{clip}%
\pgfsetbuttcap%
\pgfsetroundjoin%
\definecolor{currentfill}{rgb}{0.909804,0.698039,0.423529}%
\pgfsetfillcolor{currentfill}%
\pgfsetfillopacity{0.120000}%
\pgfsetlinewidth{0.000000pt}%
\definecolor{currentstroke}{rgb}{0.000000,0.000000,0.000000}%
\pgfsetstrokecolor{currentstroke}%
\pgfsetdash{}{0pt}%
\pgfpathmoveto{\pgfqpoint{1.408953in}{0.957283in}}%
\pgfpathlineto{\pgfqpoint{1.426457in}{0.966792in}}%
\pgfpathlineto{\pgfqpoint{1.444036in}{0.976610in}}%
\pgfpathlineto{\pgfqpoint{1.451872in}{0.981776in}}%
\pgfpathlineto{\pgfqpoint{1.459689in}{0.987302in}}%
\pgfpathlineto{\pgfqpoint{1.442134in}{0.977505in}}%
\pgfpathlineto{\pgfqpoint{1.424654in}{0.968017in}}%
\pgfpathlineto{\pgfqpoint{1.416813in}{0.962460in}}%
\pgfpathlineto{\pgfqpoint{1.408953in}{0.957283in}}%
\pgfpathclose%
\pgfusepath{fill}%
\end{pgfscope}%
\begin{pgfscope}%
\pgfpathrectangle{\pgfqpoint{0.407592in}{0.222904in}}{\pgfqpoint{1.411163in}{1.411163in}}%
\pgfusepath{clip}%
\pgfsetbuttcap%
\pgfsetroundjoin%
\definecolor{currentfill}{rgb}{0.909804,0.698039,0.423529}%
\pgfsetfillcolor{currentfill}%
\pgfsetfillopacity{0.120000}%
\pgfsetlinewidth{0.000000pt}%
\definecolor{currentstroke}{rgb}{0.000000,0.000000,0.000000}%
\pgfsetstrokecolor{currentstroke}%
\pgfsetdash{}{0pt}%
\pgfpathmoveto{\pgfqpoint{1.694371in}{1.161157in}}%
\pgfpathlineto{\pgfqpoint{1.712493in}{1.172323in}}%
\pgfpathlineto{\pgfqpoint{1.730696in}{1.183680in}}%
\pgfpathlineto{\pgfqpoint{1.738199in}{1.192391in}}%
\pgfpathlineto{\pgfqpoint{1.745689in}{1.201330in}}%
\pgfpathlineto{\pgfqpoint{1.727511in}{1.190153in}}%
\pgfpathlineto{\pgfqpoint{1.709414in}{1.179174in}}%
\pgfpathlineto{\pgfqpoint{1.701899in}{1.170048in}}%
\pgfpathlineto{\pgfqpoint{1.694371in}{1.161157in}}%
\pgfpathclose%
\pgfusepath{fill}%
\end{pgfscope}%
\begin{pgfscope}%
\pgfpathrectangle{\pgfqpoint{0.407592in}{0.222904in}}{\pgfqpoint{1.411163in}{1.411163in}}%
\pgfusepath{clip}%
\pgfsetbuttcap%
\pgfsetroundjoin%
\definecolor{currentfill}{rgb}{0.909804,0.698039,0.423529}%
\pgfsetfillcolor{currentfill}%
\pgfsetfillopacity{0.120000}%
\pgfsetlinewidth{0.000000pt}%
\definecolor{currentstroke}{rgb}{0.000000,0.000000,0.000000}%
\pgfsetstrokecolor{currentstroke}%
\pgfsetdash{}{0pt}%
\pgfpathmoveto{\pgfqpoint{0.870271in}{0.898525in}}%
\pgfpathlineto{\pgfqpoint{0.886914in}{0.892544in}}%
\pgfpathlineto{\pgfqpoint{0.903587in}{0.887328in}}%
\pgfpathlineto{\pgfqpoint{0.912402in}{0.883632in}}%
\pgfpathlineto{\pgfqpoint{0.921178in}{0.880203in}}%
\pgfpathlineto{\pgfqpoint{0.904553in}{0.885453in}}%
\pgfpathlineto{\pgfqpoint{0.887957in}{0.891498in}}%
\pgfpathlineto{\pgfqpoint{0.879135in}{0.894878in}}%
\pgfpathlineto{\pgfqpoint{0.870271in}{0.898525in}}%
\pgfpathclose%
\pgfusepath{fill}%
\end{pgfscope}%
\begin{pgfscope}%
\pgfpathrectangle{\pgfqpoint{0.407592in}{0.222904in}}{\pgfqpoint{1.411163in}{1.411163in}}%
\pgfusepath{clip}%
\pgfsetbuttcap%
\pgfsetroundjoin%
\definecolor{currentfill}{rgb}{0.909804,0.698039,0.423529}%
\pgfsetfillcolor{currentfill}%
\pgfsetfillopacity{0.120000}%
\pgfsetlinewidth{0.000000pt}%
\definecolor{currentstroke}{rgb}{0.000000,0.000000,0.000000}%
\pgfsetstrokecolor{currentstroke}%
\pgfsetdash{}{0pt}%
\pgfpathmoveto{\pgfqpoint{1.576951in}{1.068576in}}%
\pgfpathlineto{\pgfqpoint{1.594825in}{1.079381in}}%
\pgfpathlineto{\pgfqpoint{1.612779in}{1.090404in}}%
\pgfpathlineto{\pgfqpoint{1.620411in}{1.097898in}}%
\pgfpathlineto{\pgfqpoint{1.628030in}{1.105671in}}%
\pgfpathlineto{\pgfqpoint{1.610100in}{1.094788in}}%
\pgfpathlineto{\pgfqpoint{1.592250in}{1.084128in}}%
\pgfpathlineto{\pgfqpoint{1.584608in}{1.076206in}}%
\pgfpathlineto{\pgfqpoint{1.576951in}{1.068576in}}%
\pgfpathclose%
\pgfusepath{fill}%
\end{pgfscope}%
\begin{pgfscope}%
\pgfpathrectangle{\pgfqpoint{0.407592in}{0.222904in}}{\pgfqpoint{1.411163in}{1.411163in}}%
\pgfusepath{clip}%
\pgfsetbuttcap%
\pgfsetroundjoin%
\definecolor{currentfill}{rgb}{0.909804,0.698039,0.423529}%
\pgfsetfillcolor{currentfill}%
\pgfsetfillopacity{0.120000}%
\pgfsetlinewidth{0.000000pt}%
\definecolor{currentstroke}{rgb}{0.000000,0.000000,0.000000}%
\pgfsetstrokecolor{currentstroke}%
\pgfsetdash{}{0pt}%
\pgfpathmoveto{\pgfqpoint{0.701319in}{0.985951in}}%
\pgfpathlineto{\pgfqpoint{0.717886in}{0.975033in}}%
\pgfpathlineto{\pgfqpoint{0.734469in}{0.964673in}}%
\pgfpathlineto{\pgfqpoint{0.743646in}{0.960540in}}%
\pgfpathlineto{\pgfqpoint{0.752780in}{0.956630in}}%
\pgfpathlineto{\pgfqpoint{0.736244in}{0.967207in}}%
\pgfpathlineto{\pgfqpoint{0.719724in}{0.978355in}}%
\pgfpathlineto{\pgfqpoint{0.710543in}{0.982041in}}%
\pgfpathlineto{\pgfqpoint{0.701319in}{0.985951in}}%
\pgfpathclose%
\pgfusepath{fill}%
\end{pgfscope}%
\begin{pgfscope}%
\pgfpathrectangle{\pgfqpoint{0.407592in}{0.222904in}}{\pgfqpoint{1.411163in}{1.411163in}}%
\pgfusepath{clip}%
\pgfsetbuttcap%
\pgfsetroundjoin%
\definecolor{currentfill}{rgb}{0.909804,0.698039,0.423529}%
\pgfsetfillcolor{currentfill}%
\pgfsetfillopacity{0.120000}%
\pgfsetlinewidth{0.000000pt}%
\definecolor{currentstroke}{rgb}{0.000000,0.000000,0.000000}%
\pgfsetstrokecolor{currentstroke}%
\pgfsetdash{}{0pt}%
\pgfpathmoveto{\pgfqpoint{1.745689in}{1.201330in}}%
\pgfpathlineto{\pgfqpoint{1.763949in}{1.212696in}}%
\pgfpathlineto{\pgfqpoint{1.771414in}{1.221764in}}%
\pgfpathlineto{\pgfqpoint{1.778868in}{1.231045in}}%
\pgfpathlineto{\pgfqpoint{1.760634in}{1.219869in}}%
\pgfpathlineto{\pgfqpoint{1.753168in}{1.210491in}}%
\pgfusepath{fill}%
\end{pgfscope}%
\begin{pgfscope}%
\pgfpathrectangle{\pgfqpoint{0.407592in}{0.222904in}}{\pgfqpoint{1.411163in}{1.411163in}}%
\pgfusepath{clip}%
\pgfsetbuttcap%
\pgfsetroundjoin%
\definecolor{currentfill}{rgb}{0.909804,0.698039,0.423529}%
\pgfsetfillcolor{currentfill}%
\pgfsetfillopacity{0.120000}%
\pgfsetlinewidth{0.000000pt}%
\definecolor{currentstroke}{rgb}{0.000000,0.000000,0.000000}%
\pgfsetstrokecolor{currentstroke}%
\pgfsetdash{}{0pt}%
\pgfpathmoveto{\pgfqpoint{1.459689in}{0.987302in}}%
\pgfpathlineto{\pgfqpoint{1.477322in}{0.997386in}}%
\pgfpathlineto{\pgfqpoint{1.495032in}{1.007740in}}%
\pgfpathlineto{\pgfqpoint{1.502807in}{1.013572in}}%
\pgfpathlineto{\pgfqpoint{1.510566in}{1.019738in}}%
\pgfpathlineto{\pgfqpoint{1.492880in}{1.009444in}}%
\pgfpathlineto{\pgfqpoint{1.475271in}{0.999420in}}%
\pgfpathlineto{\pgfqpoint{1.467489in}{0.993184in}}%
\pgfpathlineto{\pgfqpoint{1.459689in}{0.987302in}}%
\pgfpathclose%
\pgfusepath{fill}%
\end{pgfscope}%
\begin{pgfscope}%
\pgfpathrectangle{\pgfqpoint{0.407592in}{0.222904in}}{\pgfqpoint{1.411163in}{1.411163in}}%
\pgfusepath{clip}%
\pgfsetbuttcap%
\pgfsetroundjoin%
\definecolor{currentfill}{rgb}{0.909804,0.698039,0.423529}%
\pgfsetfillcolor{currentfill}%
\pgfsetfillopacity{0.120000}%
\pgfsetlinewidth{0.000000pt}%
\definecolor{currentstroke}{rgb}{0.000000,0.000000,0.000000}%
\pgfsetstrokecolor{currentstroke}%
\pgfsetdash{}{0pt}%
\pgfpathmoveto{\pgfqpoint{0.616435in}{1.042185in}}%
\pgfpathlineto{\pgfqpoint{0.632989in}{1.029537in}}%
\pgfpathlineto{\pgfqpoint{0.649557in}{1.017343in}}%
\pgfpathlineto{\pgfqpoint{0.658917in}{1.013249in}}%
\pgfpathlineto{\pgfqpoint{0.668232in}{1.009363in}}%
\pgfpathlineto{\pgfqpoint{0.651711in}{1.021807in}}%
\pgfpathlineto{\pgfqpoint{0.635203in}{1.034708in}}%
\pgfpathlineto{\pgfqpoint{0.625841in}{1.038340in}}%
\pgfpathlineto{\pgfqpoint{0.616435in}{1.042185in}}%
\pgfpathclose%
\pgfusepath{fill}%
\end{pgfscope}%
\begin{pgfscope}%
\pgfpathrectangle{\pgfqpoint{0.407592in}{0.222904in}}{\pgfqpoint{1.411163in}{1.411163in}}%
\pgfusepath{clip}%
\pgfsetbuttcap%
\pgfsetroundjoin%
\definecolor{currentfill}{rgb}{0.909804,0.698039,0.423529}%
\pgfsetfillcolor{currentfill}%
\pgfsetfillopacity{0.120000}%
\pgfsetlinewidth{0.000000pt}%
\definecolor{currentstroke}{rgb}{0.000000,0.000000,0.000000}%
\pgfsetstrokecolor{currentstroke}%
\pgfsetdash{}{0pt}%
\pgfpathmoveto{\pgfqpoint{0.954524in}{0.872122in}}%
\pgfpathlineto{\pgfqpoint{0.971251in}{0.869294in}}%
\pgfpathlineto{\pgfqpoint{0.988017in}{0.867267in}}%
\pgfpathlineto{\pgfqpoint{0.996659in}{0.864222in}}%
\pgfpathlineto{\pgfqpoint{1.005263in}{0.861478in}}%
\pgfpathlineto{\pgfqpoint{0.988543in}{0.863405in}}%
\pgfpathlineto{\pgfqpoint{0.971863in}{0.866165in}}%
\pgfpathlineto{\pgfqpoint{0.963213in}{0.868993in}}%
\pgfpathlineto{\pgfqpoint{0.954524in}{0.872122in}}%
\pgfpathclose%
\pgfusepath{fill}%
\end{pgfscope}%
\begin{pgfscope}%
\pgfpathrectangle{\pgfqpoint{0.407592in}{0.222904in}}{\pgfqpoint{1.411163in}{1.411163in}}%
\pgfusepath{clip}%
\pgfsetbuttcap%
\pgfsetroundjoin%
\definecolor{currentfill}{rgb}{0.909804,0.698039,0.423529}%
\pgfsetfillcolor{currentfill}%
\pgfsetfillopacity{0.120000}%
\pgfsetlinewidth{0.000000pt}%
\definecolor{currentstroke}{rgb}{0.000000,0.000000,0.000000}%
\pgfsetstrokecolor{currentstroke}%
\pgfsetdash{}{0pt}%
\pgfpathmoveto{\pgfqpoint{1.628030in}{1.105671in}}%
\pgfpathlineto{\pgfqpoint{1.646040in}{1.116764in}}%
\pgfpathlineto{\pgfqpoint{1.664130in}{1.128058in}}%
\pgfpathlineto{\pgfqpoint{1.671710in}{1.135952in}}%
\pgfpathlineto{\pgfqpoint{1.679277in}{1.144104in}}%
\pgfpathlineto{\pgfqpoint{1.661211in}{1.132964in}}%
\pgfpathlineto{\pgfqpoint{1.643226in}{1.122030in}}%
\pgfpathlineto{\pgfqpoint{1.635635in}{1.113717in}}%
\pgfpathlineto{\pgfqpoint{1.628030in}{1.105671in}}%
\pgfpathclose%
\pgfusepath{fill}%
\end{pgfscope}%
\begin{pgfscope}%
\pgfpathrectangle{\pgfqpoint{0.407592in}{0.222904in}}{\pgfqpoint{1.411163in}{1.411163in}}%
\pgfusepath{clip}%
\pgfsetbuttcap%
\pgfsetroundjoin%
\definecolor{currentfill}{rgb}{0.909804,0.698039,0.423529}%
\pgfsetfillcolor{currentfill}%
\pgfsetfillopacity{0.120000}%
\pgfsetlinewidth{0.000000pt}%
\definecolor{currentstroke}{rgb}{0.000000,0.000000,0.000000}%
\pgfsetstrokecolor{currentstroke}%
\pgfsetdash{}{0pt}%
\pgfpathmoveto{\pgfqpoint{1.190652in}{0.869466in}}%
\pgfpathlineto{\pgfqpoint{1.207730in}{0.874796in}}%
\pgfpathlineto{\pgfqpoint{1.224872in}{0.880702in}}%
\pgfpathlineto{\pgfqpoint{1.233048in}{0.881540in}}%
\pgfpathlineto{\pgfqpoint{1.241197in}{0.882778in}}%
\pgfpathlineto{\pgfqpoint{1.224087in}{0.876661in}}%
\pgfpathlineto{\pgfqpoint{1.207042in}{0.871119in}}%
\pgfpathlineto{\pgfqpoint{1.198861in}{0.870084in}}%
\pgfpathlineto{\pgfqpoint{1.190652in}{0.869466in}}%
\pgfpathclose%
\pgfusepath{fill}%
\end{pgfscope}%
\begin{pgfscope}%
\pgfpathrectangle{\pgfqpoint{0.407592in}{0.222904in}}{\pgfqpoint{1.411163in}{1.411163in}}%
\pgfusepath{clip}%
\pgfsetbuttcap%
\pgfsetroundjoin%
\definecolor{currentfill}{rgb}{0.909804,0.698039,0.423529}%
\pgfsetfillcolor{currentfill}%
\pgfsetfillopacity{0.120000}%
\pgfsetlinewidth{0.000000pt}%
\definecolor{currentstroke}{rgb}{0.000000,0.000000,0.000000}%
\pgfsetstrokecolor{currentstroke}%
\pgfsetdash{}{0pt}%
\pgfpathmoveto{\pgfqpoint{1.241197in}{0.882778in}}%
\pgfpathlineto{\pgfqpoint{1.258375in}{0.889431in}}%
\pgfpathlineto{\pgfqpoint{1.275621in}{0.896583in}}%
\pgfpathlineto{\pgfqpoint{1.283712in}{0.898415in}}%
\pgfpathlineto{\pgfqpoint{1.291778in}{0.900649in}}%
\pgfpathlineto{\pgfqpoint{1.274561in}{0.893327in}}%
\pgfpathlineto{\pgfqpoint{1.257414in}{0.886499in}}%
\pgfpathlineto{\pgfqpoint{1.249318in}{0.884428in}}%
\pgfpathlineto{\pgfqpoint{1.241197in}{0.882778in}}%
\pgfpathclose%
\pgfusepath{fill}%
\end{pgfscope}%
\begin{pgfscope}%
\pgfpathrectangle{\pgfqpoint{0.407592in}{0.222904in}}{\pgfqpoint{1.411163in}{1.411163in}}%
\pgfusepath{clip}%
\pgfsetbuttcap%
\pgfsetroundjoin%
\definecolor{currentfill}{rgb}{0.909804,0.698039,0.423529}%
\pgfsetfillcolor{currentfill}%
\pgfsetfillopacity{0.120000}%
\pgfsetlinewidth{0.000000pt}%
\definecolor{currentstroke}{rgb}{0.000000,0.000000,0.000000}%
\pgfsetstrokecolor{currentstroke}%
\pgfsetdash{}{0pt}%
\pgfpathmoveto{\pgfqpoint{1.140105in}{0.861114in}}%
\pgfpathlineto{\pgfqpoint{1.157094in}{0.864943in}}%
\pgfpathlineto{\pgfqpoint{1.174144in}{0.869424in}}%
\pgfpathlineto{\pgfqpoint{1.182413in}{0.869250in}}%
\pgfpathlineto{\pgfqpoint{1.190652in}{0.869466in}}%
\pgfpathlineto{\pgfqpoint{1.173638in}{0.864751in}}%
\pgfpathlineto{\pgfqpoint{1.156685in}{0.860693in}}%
\pgfpathlineto{\pgfqpoint{1.148411in}{0.860704in}}%
\pgfpathlineto{\pgfqpoint{1.140105in}{0.861114in}}%
\pgfpathclose%
\pgfusepath{fill}%
\end{pgfscope}%
\begin{pgfscope}%
\pgfpathrectangle{\pgfqpoint{0.407592in}{0.222904in}}{\pgfqpoint{1.411163in}{1.411163in}}%
\pgfusepath{clip}%
\pgfsetbuttcap%
\pgfsetroundjoin%
\definecolor{currentfill}{rgb}{0.909804,0.698039,0.423529}%
\pgfsetfillcolor{currentfill}%
\pgfsetfillopacity{0.120000}%
\pgfsetlinewidth{0.000000pt}%
\definecolor{currentstroke}{rgb}{0.000000,0.000000,0.000000}%
\pgfsetstrokecolor{currentstroke}%
\pgfsetdash{}{0pt}%
\pgfpathmoveto{\pgfqpoint{1.291778in}{0.900649in}}%
\pgfpathlineto{\pgfqpoint{1.309066in}{0.908431in}}%
\pgfpathlineto{\pgfqpoint{1.326425in}{0.916638in}}%
\pgfpathlineto{\pgfqpoint{1.334438in}{0.919416in}}%
\pgfpathlineto{\pgfqpoint{1.342428in}{0.922588in}}%
\pgfpathlineto{\pgfqpoint{1.325096in}{0.914259in}}%
\pgfpathlineto{\pgfqpoint{1.307836in}{0.906352in}}%
\pgfpathlineto{\pgfqpoint{1.299819in}{0.903293in}}%
\pgfpathlineto{\pgfqpoint{1.291778in}{0.900649in}}%
\pgfpathclose%
\pgfusepath{fill}%
\end{pgfscope}%
\begin{pgfscope}%
\pgfpathrectangle{\pgfqpoint{0.407592in}{0.222904in}}{\pgfqpoint{1.411163in}{1.411163in}}%
\pgfusepath{clip}%
\pgfsetbuttcap%
\pgfsetroundjoin%
\definecolor{currentfill}{rgb}{0.909804,0.698039,0.423529}%
\pgfsetfillcolor{currentfill}%
\pgfsetfillopacity{0.120000}%
\pgfsetlinewidth{0.000000pt}%
\definecolor{currentstroke}{rgb}{0.000000,0.000000,0.000000}%
\pgfsetstrokecolor{currentstroke}%
\pgfsetdash{}{0pt}%
\pgfpathmoveto{\pgfqpoint{1.510566in}{1.019738in}}%
\pgfpathlineto{\pgfqpoint{1.528330in}{1.030286in}}%
\pgfpathlineto{\pgfqpoint{1.546173in}{1.041071in}}%
\pgfpathlineto{\pgfqpoint{1.553891in}{1.047486in}}%
\pgfpathlineto{\pgfqpoint{1.561593in}{1.054211in}}%
\pgfpathlineto{\pgfqpoint{1.543774in}{1.043515in}}%
\pgfpathlineto{\pgfqpoint{1.526034in}{1.033060in}}%
\pgfpathlineto{\pgfqpoint{1.518308in}{1.026235in}}%
\pgfpathlineto{\pgfqpoint{1.510566in}{1.019738in}}%
\pgfpathclose%
\pgfusepath{fill}%
\end{pgfscope}%
\begin{pgfscope}%
\pgfpathrectangle{\pgfqpoint{0.407592in}{0.222904in}}{\pgfqpoint{1.411163in}{1.411163in}}%
\pgfusepath{clip}%
\pgfsetbuttcap%
\pgfsetroundjoin%
\definecolor{currentfill}{rgb}{0.909804,0.698039,0.423529}%
\pgfsetfillcolor{currentfill}%
\pgfsetfillopacity{0.120000}%
\pgfsetlinewidth{0.000000pt}%
\definecolor{currentstroke}{rgb}{0.000000,0.000000,0.000000}%
\pgfsetstrokecolor{currentstroke}%
\pgfsetdash{}{0pt}%
\pgfpathmoveto{\pgfqpoint{0.819116in}{0.920549in}}%
\pgfpathlineto{\pgfqpoint{0.835756in}{0.913201in}}%
\pgfpathlineto{\pgfqpoint{0.852421in}{0.906562in}}%
\pgfpathlineto{\pgfqpoint{0.861367in}{0.902424in}}%
\pgfpathlineto{\pgfqpoint{0.870271in}{0.898525in}}%
\pgfpathlineto{\pgfqpoint{0.853655in}{0.905257in}}%
\pgfpathlineto{\pgfqpoint{0.837064in}{0.912725in}}%
\pgfpathlineto{\pgfqpoint{0.828111in}{0.916517in}}%
\pgfpathlineto{\pgfqpoint{0.819116in}{0.920549in}}%
\pgfpathclose%
\pgfusepath{fill}%
\end{pgfscope}%
\begin{pgfscope}%
\pgfpathrectangle{\pgfqpoint{0.407592in}{0.222904in}}{\pgfqpoint{1.411163in}{1.411163in}}%
\pgfusepath{clip}%
\pgfsetbuttcap%
\pgfsetroundjoin%
\definecolor{currentfill}{rgb}{0.909804,0.698039,0.423529}%
\pgfsetfillcolor{currentfill}%
\pgfsetfillopacity{0.120000}%
\pgfsetlinewidth{0.000000pt}%
\definecolor{currentstroke}{rgb}{0.000000,0.000000,0.000000}%
\pgfsetstrokecolor{currentstroke}%
\pgfsetdash{}{0pt}%
\pgfpathmoveto{\pgfqpoint{0.734469in}{0.964673in}}%
\pgfpathlineto{\pgfqpoint{0.751071in}{0.954896in}}%
\pgfpathlineto{\pgfqpoint{0.767693in}{0.945724in}}%
\pgfpathlineto{\pgfqpoint{0.776821in}{0.941400in}}%
\pgfpathlineto{\pgfqpoint{0.785907in}{0.937294in}}%
\pgfpathlineto{\pgfqpoint{0.769334in}{0.946650in}}%
\pgfpathlineto{\pgfqpoint{0.752780in}{0.956630in}}%
\pgfpathlineto{\pgfqpoint{0.743646in}{0.960540in}}%
\pgfpathlineto{\pgfqpoint{0.734469in}{0.964673in}}%
\pgfpathclose%
\pgfusepath{fill}%
\end{pgfscope}%
\begin{pgfscope}%
\pgfpathrectangle{\pgfqpoint{0.407592in}{0.222904in}}{\pgfqpoint{1.411163in}{1.411163in}}%
\pgfusepath{clip}%
\pgfsetbuttcap%
\pgfsetroundjoin%
\definecolor{currentfill}{rgb}{0.909804,0.698039,0.423529}%
\pgfsetfillcolor{currentfill}%
\pgfsetfillopacity{0.120000}%
\pgfsetlinewidth{0.000000pt}%
\definecolor{currentstroke}{rgb}{0.000000,0.000000,0.000000}%
\pgfsetstrokecolor{currentstroke}%
\pgfsetdash{}{0pt}%
\pgfpathmoveto{\pgfqpoint{1.089513in}{0.857967in}}%
\pgfpathlineto{\pgfqpoint{1.106427in}{0.860160in}}%
\pgfpathlineto{\pgfqpoint{1.123397in}{0.863066in}}%
\pgfpathlineto{\pgfqpoint{1.131767in}{0.861907in}}%
\pgfpathlineto{\pgfqpoint{1.140105in}{0.861114in}}%
\pgfpathlineto{\pgfqpoint{1.123174in}{0.857973in}}%
\pgfpathlineto{\pgfqpoint{1.106299in}{0.855560in}}%
\pgfpathlineto{\pgfqpoint{1.097923in}{0.856578in}}%
\pgfpathlineto{\pgfqpoint{1.089513in}{0.857967in}}%
\pgfpathclose%
\pgfusepath{fill}%
\end{pgfscope}%
\begin{pgfscope}%
\pgfpathrectangle{\pgfqpoint{0.407592in}{0.222904in}}{\pgfqpoint{1.411163in}{1.411163in}}%
\pgfusepath{clip}%
\pgfsetbuttcap%
\pgfsetroundjoin%
\definecolor{currentfill}{rgb}{0.909804,0.698039,0.423529}%
\pgfsetfillcolor{currentfill}%
\pgfsetfillopacity{0.120000}%
\pgfsetlinewidth{0.000000pt}%
\definecolor{currentstroke}{rgb}{0.000000,0.000000,0.000000}%
\pgfsetstrokecolor{currentstroke}%
\pgfsetdash{}{0pt}%
\pgfpathmoveto{\pgfqpoint{1.342428in}{0.922588in}}%
\pgfpathlineto{\pgfqpoint{1.359833in}{0.931307in}}%
\pgfpathlineto{\pgfqpoint{1.377313in}{0.940390in}}%
\pgfpathlineto{\pgfqpoint{1.385254in}{0.944042in}}%
\pgfpathlineto{\pgfqpoint{1.393174in}{0.948074in}}%
\pgfpathlineto{\pgfqpoint{1.375721in}{0.938923in}}%
\pgfpathlineto{\pgfqpoint{1.358341in}{0.930130in}}%
\pgfpathlineto{\pgfqpoint{1.350396in}{0.926158in}}%
\pgfpathlineto{\pgfqpoint{1.342428in}{0.922588in}}%
\pgfpathclose%
\pgfusepath{fill}%
\end{pgfscope}%
\begin{pgfscope}%
\pgfpathrectangle{\pgfqpoint{0.407592in}{0.222904in}}{\pgfqpoint{1.411163in}{1.411163in}}%
\pgfusepath{clip}%
\pgfsetbuttcap%
\pgfsetroundjoin%
\definecolor{currentfill}{rgb}{0.909804,0.698039,0.423529}%
\pgfsetfillcolor{currentfill}%
\pgfsetfillopacity{0.120000}%
\pgfsetlinewidth{0.000000pt}%
\definecolor{currentstroke}{rgb}{0.000000,0.000000,0.000000}%
\pgfsetstrokecolor{currentstroke}%
\pgfsetdash{}{0pt}%
\pgfpathmoveto{\pgfqpoint{0.903587in}{0.887328in}}%
\pgfpathlineto{\pgfqpoint{0.920291in}{0.882884in}}%
\pgfpathlineto{\pgfqpoint{0.937030in}{0.879215in}}%
\pgfpathlineto{\pgfqpoint{0.945797in}{0.875535in}}%
\pgfpathlineto{\pgfqpoint{0.954524in}{0.872122in}}%
\pgfpathlineto{\pgfqpoint{0.937834in}{0.875758in}}%
\pgfpathlineto{\pgfqpoint{0.921178in}{0.880203in}}%
\pgfpathlineto{\pgfqpoint{0.912402in}{0.883632in}}%
\pgfpathlineto{\pgfqpoint{0.903587in}{0.887328in}}%
\pgfpathclose%
\pgfusepath{fill}%
\end{pgfscope}%
\begin{pgfscope}%
\pgfpathrectangle{\pgfqpoint{0.407592in}{0.222904in}}{\pgfqpoint{1.411163in}{1.411163in}}%
\pgfusepath{clip}%
\pgfsetbuttcap%
\pgfsetroundjoin%
\definecolor{currentfill}{rgb}{0.909804,0.698039,0.423529}%
\pgfsetfillcolor{currentfill}%
\pgfsetfillopacity{0.120000}%
\pgfsetlinewidth{0.000000pt}%
\definecolor{currentstroke}{rgb}{0.000000,0.000000,0.000000}%
\pgfsetstrokecolor{currentstroke}%
\pgfsetdash{}{0pt}%
\pgfpathmoveto{\pgfqpoint{0.649557in}{1.017343in}}%
\pgfpathlineto{\pgfqpoint{0.666140in}{1.005626in}}%
\pgfpathlineto{\pgfqpoint{0.682738in}{0.994408in}}%
\pgfpathlineto{\pgfqpoint{0.692050in}{0.990077in}}%
\pgfpathlineto{\pgfqpoint{0.701319in}{0.985951in}}%
\pgfpathlineto{\pgfqpoint{0.684768in}{0.997403in}}%
\pgfpathlineto{\pgfqpoint{0.668232in}{1.009363in}}%
\pgfpathlineto{\pgfqpoint{0.658917in}{1.013249in}}%
\pgfpathlineto{\pgfqpoint{0.649557in}{1.017343in}}%
\pgfpathclose%
\pgfusepath{fill}%
\end{pgfscope}%
\begin{pgfscope}%
\pgfpathrectangle{\pgfqpoint{0.407592in}{0.222904in}}{\pgfqpoint{1.411163in}{1.411163in}}%
\pgfusepath{clip}%
\pgfsetbuttcap%
\pgfsetroundjoin%
\definecolor{currentfill}{rgb}{0.909804,0.698039,0.423529}%
\pgfsetfillcolor{currentfill}%
\pgfsetfillopacity{0.120000}%
\pgfsetlinewidth{0.000000pt}%
\definecolor{currentstroke}{rgb}{0.000000,0.000000,0.000000}%
\pgfsetstrokecolor{currentstroke}%
\pgfsetdash{}{0pt}%
\pgfpathmoveto{\pgfqpoint{1.679277in}{1.144104in}}%
\pgfpathlineto{\pgfqpoint{1.697423in}{1.155439in}}%
\pgfpathlineto{\pgfqpoint{1.715651in}{1.166960in}}%
\pgfpathlineto{\pgfqpoint{1.723180in}{1.175201in}}%
\pgfpathlineto{\pgfqpoint{1.730696in}{1.183680in}}%
\pgfpathlineto{\pgfqpoint{1.712493in}{1.172323in}}%
\pgfpathlineto{\pgfqpoint{1.694371in}{1.161157in}}%
\pgfpathlineto{\pgfqpoint{1.686830in}{1.152507in}}%
\pgfpathlineto{\pgfqpoint{1.679277in}{1.144104in}}%
\pgfpathclose%
\pgfusepath{fill}%
\end{pgfscope}%
\begin{pgfscope}%
\pgfpathrectangle{\pgfqpoint{0.407592in}{0.222904in}}{\pgfqpoint{1.411163in}{1.411163in}}%
\pgfusepath{clip}%
\pgfsetbuttcap%
\pgfsetroundjoin%
\definecolor{currentfill}{rgb}{0.909804,0.698039,0.423529}%
\pgfsetfillcolor{currentfill}%
\pgfsetfillopacity{0.120000}%
\pgfsetlinewidth{0.000000pt}%
\definecolor{currentstroke}{rgb}{0.000000,0.000000,0.000000}%
\pgfsetstrokecolor{currentstroke}%
\pgfsetdash{}{0pt}%
\pgfpathmoveto{\pgfqpoint{1.038832in}{0.860072in}}%
\pgfpathlineto{\pgfqpoint{1.055685in}{0.860553in}}%
\pgfpathlineto{\pgfqpoint{1.072588in}{0.861791in}}%
\pgfpathlineto{\pgfqpoint{1.081068in}{0.859710in}}%
\pgfpathlineto{\pgfqpoint{1.089513in}{0.857967in}}%
\pgfpathlineto{\pgfqpoint{1.072652in}{0.856521in}}%
\pgfpathlineto{\pgfqpoint{1.055841in}{0.855855in}}%
\pgfpathlineto{\pgfqpoint{1.047354in}{0.857795in}}%
\pgfpathlineto{\pgfqpoint{1.038832in}{0.860072in}}%
\pgfpathclose%
\pgfusepath{fill}%
\end{pgfscope}%
\begin{pgfscope}%
\pgfpathrectangle{\pgfqpoint{0.407592in}{0.222904in}}{\pgfqpoint{1.411163in}{1.411163in}}%
\pgfusepath{clip}%
\pgfsetbuttcap%
\pgfsetroundjoin%
\definecolor{currentfill}{rgb}{0.909804,0.698039,0.423529}%
\pgfsetfillcolor{currentfill}%
\pgfsetfillopacity{0.120000}%
\pgfsetlinewidth{0.000000pt}%
\definecolor{currentstroke}{rgb}{0.000000,0.000000,0.000000}%
\pgfsetstrokecolor{currentstroke}%
\pgfsetdash{}{0pt}%
\pgfpathmoveto{\pgfqpoint{1.561593in}{1.054211in}}%
\pgfpathlineto{\pgfqpoint{1.579491in}{1.065133in}}%
\pgfpathlineto{\pgfqpoint{1.597469in}{1.076270in}}%
\pgfpathlineto{\pgfqpoint{1.605131in}{1.083193in}}%
\pgfpathlineto{\pgfqpoint{1.612779in}{1.090404in}}%
\pgfpathlineto{\pgfqpoint{1.594825in}{1.079381in}}%
\pgfpathlineto{\pgfqpoint{1.576951in}{1.068576in}}%
\pgfpathlineto{\pgfqpoint{1.569279in}{1.061242in}}%
\pgfpathlineto{\pgfqpoint{1.561593in}{1.054211in}}%
\pgfpathclose%
\pgfusepath{fill}%
\end{pgfscope}%
\begin{pgfscope}%
\pgfpathrectangle{\pgfqpoint{0.407592in}{0.222904in}}{\pgfqpoint{1.411163in}{1.411163in}}%
\pgfusepath{clip}%
\pgfsetbuttcap%
\pgfsetroundjoin%
\definecolor{currentfill}{rgb}{0.909804,0.698039,0.423529}%
\pgfsetfillcolor{currentfill}%
\pgfsetfillopacity{0.120000}%
\pgfsetlinewidth{0.000000pt}%
\definecolor{currentstroke}{rgb}{0.000000,0.000000,0.000000}%
\pgfsetstrokecolor{currentstroke}%
\pgfsetdash{}{0pt}%
\pgfpathmoveto{\pgfqpoint{1.393174in}{0.948074in}}%
\pgfpathlineto{\pgfqpoint{1.410703in}{0.957560in}}%
\pgfpathlineto{\pgfqpoint{1.428308in}{0.967357in}}%
\pgfpathlineto{\pgfqpoint{1.436182in}{0.971803in}}%
\pgfpathlineto{\pgfqpoint{1.444036in}{0.976610in}}%
\pgfpathlineto{\pgfqpoint{1.426457in}{0.966792in}}%
\pgfpathlineto{\pgfqpoint{1.408953in}{0.957283in}}%
\pgfpathlineto{\pgfqpoint{1.401073in}{0.952488in}}%
\pgfpathlineto{\pgfqpoint{1.393174in}{0.948074in}}%
\pgfpathclose%
\pgfusepath{fill}%
\end{pgfscope}%
\begin{pgfscope}%
\pgfpathrectangle{\pgfqpoint{0.407592in}{0.222904in}}{\pgfqpoint{1.411163in}{1.411163in}}%
\pgfusepath{clip}%
\pgfsetbuttcap%
\pgfsetroundjoin%
\definecolor{currentfill}{rgb}{0.909804,0.698039,0.423529}%
\pgfsetfillcolor{currentfill}%
\pgfsetfillopacity{0.120000}%
\pgfsetlinewidth{0.000000pt}%
\definecolor{currentstroke}{rgb}{0.000000,0.000000,0.000000}%
\pgfsetstrokecolor{currentstroke}%
\pgfsetdash{}{0pt}%
\pgfpathmoveto{\pgfqpoint{1.730696in}{1.183680in}}%
\pgfpathlineto{\pgfqpoint{1.748980in}{1.195221in}}%
\pgfpathlineto{\pgfqpoint{1.756471in}{1.203847in}}%
\pgfpathlineto{\pgfqpoint{1.763949in}{1.212696in}}%
\pgfpathlineto{\pgfqpoint{1.745689in}{1.201330in}}%
\pgfpathlineto{\pgfqpoint{1.738199in}{1.192391in}}%
\pgfusepath{fill}%
\end{pgfscope}%
\begin{pgfscope}%
\pgfpathrectangle{\pgfqpoint{0.407592in}{0.222904in}}{\pgfqpoint{1.411163in}{1.411163in}}%
\pgfusepath{clip}%
\pgfsetbuttcap%
\pgfsetroundjoin%
\definecolor{currentfill}{rgb}{0.909804,0.698039,0.423529}%
\pgfsetfillcolor{currentfill}%
\pgfsetfillopacity{0.120000}%
\pgfsetlinewidth{0.000000pt}%
\definecolor{currentstroke}{rgb}{0.000000,0.000000,0.000000}%
\pgfsetstrokecolor{currentstroke}%
\pgfsetdash{}{0pt}%
\pgfpathmoveto{\pgfqpoint{0.988017in}{0.867267in}}%
\pgfpathlineto{\pgfqpoint{1.004824in}{0.866031in}}%
\pgfpathlineto{\pgfqpoint{1.021676in}{0.865571in}}%
\pgfpathlineto{\pgfqpoint{1.030272in}{0.862670in}}%
\pgfpathlineto{\pgfqpoint{1.038832in}{0.860072in}}%
\pgfpathlineto{\pgfqpoint{1.022025in}{0.860373in}}%
\pgfpathlineto{\pgfqpoint{1.005263in}{0.861478in}}%
\pgfpathlineto{\pgfqpoint{0.996659in}{0.864222in}}%
\pgfpathlineto{\pgfqpoint{0.988017in}{0.867267in}}%
\pgfpathclose%
\pgfusepath{fill}%
\end{pgfscope}%
\begin{pgfscope}%
\pgfpathrectangle{\pgfqpoint{0.407592in}{0.222904in}}{\pgfqpoint{1.411163in}{1.411163in}}%
\pgfusepath{clip}%
\pgfsetbuttcap%
\pgfsetroundjoin%
\definecolor{currentfill}{rgb}{0.909804,0.698039,0.423529}%
\pgfsetfillcolor{currentfill}%
\pgfsetfillopacity{0.120000}%
\pgfsetlinewidth{0.000000pt}%
\definecolor{currentstroke}{rgb}{0.000000,0.000000,0.000000}%
\pgfsetstrokecolor{currentstroke}%
\pgfsetdash{}{0pt}%
\pgfpathmoveto{\pgfqpoint{1.444036in}{0.976610in}}%
\pgfpathlineto{\pgfqpoint{1.461693in}{0.986715in}}%
\pgfpathlineto{\pgfqpoint{1.479428in}{0.997090in}}%
\pgfpathlineto{\pgfqpoint{1.487239in}{1.002245in}}%
\pgfpathlineto{\pgfqpoint{1.495032in}{1.007740in}}%
\pgfpathlineto{\pgfqpoint{1.477322in}{0.997386in}}%
\pgfpathlineto{\pgfqpoint{1.459689in}{0.987302in}}%
\pgfpathlineto{\pgfqpoint{1.451872in}{0.981776in}}%
\pgfpathlineto{\pgfqpoint{1.444036in}{0.976610in}}%
\pgfpathclose%
\pgfusepath{fill}%
\end{pgfscope}%
\begin{pgfscope}%
\pgfpathrectangle{\pgfqpoint{0.407592in}{0.222904in}}{\pgfqpoint{1.411163in}{1.411163in}}%
\pgfusepath{clip}%
\pgfsetbuttcap%
\pgfsetroundjoin%
\definecolor{currentfill}{rgb}{0.909804,0.698039,0.423529}%
\pgfsetfillcolor{currentfill}%
\pgfsetfillopacity{0.120000}%
\pgfsetlinewidth{0.000000pt}%
\definecolor{currentstroke}{rgb}{0.000000,0.000000,0.000000}%
\pgfsetstrokecolor{currentstroke}%
\pgfsetdash{}{0pt}%
\pgfpathmoveto{\pgfqpoint{0.767693in}{0.945724in}}%
\pgfpathlineto{\pgfqpoint{0.784335in}{0.937180in}}%
\pgfpathlineto{\pgfqpoint{0.801000in}{0.929284in}}%
\pgfpathlineto{\pgfqpoint{0.810079in}{0.924808in}}%
\pgfpathlineto{\pgfqpoint{0.819116in}{0.920549in}}%
\pgfpathlineto{\pgfqpoint{0.802501in}{0.928586in}}%
\pgfpathlineto{\pgfqpoint{0.785907in}{0.937294in}}%
\pgfpathlineto{\pgfqpoint{0.776821in}{0.941400in}}%
\pgfpathlineto{\pgfqpoint{0.767693in}{0.945724in}}%
\pgfpathclose%
\pgfusepath{fill}%
\end{pgfscope}%
\begin{pgfscope}%
\pgfpathrectangle{\pgfqpoint{0.407592in}{0.222904in}}{\pgfqpoint{1.411163in}{1.411163in}}%
\pgfusepath{clip}%
\pgfsetbuttcap%
\pgfsetroundjoin%
\definecolor{currentfill}{rgb}{0.909804,0.698039,0.423529}%
\pgfsetfillcolor{currentfill}%
\pgfsetfillopacity{0.120000}%
\pgfsetlinewidth{0.000000pt}%
\definecolor{currentstroke}{rgb}{0.000000,0.000000,0.000000}%
\pgfsetstrokecolor{currentstroke}%
\pgfsetdash{}{0pt}%
\pgfpathmoveto{\pgfqpoint{0.852421in}{0.906562in}}%
\pgfpathlineto{\pgfqpoint{0.869113in}{0.900646in}}%
\pgfpathlineto{\pgfqpoint{0.885835in}{0.895462in}}%
\pgfpathlineto{\pgfqpoint{0.894731in}{0.891276in}}%
\pgfpathlineto{\pgfqpoint{0.903587in}{0.887328in}}%
\pgfpathlineto{\pgfqpoint{0.886914in}{0.892544in}}%
\pgfpathlineto{\pgfqpoint{0.870271in}{0.898525in}}%
\pgfpathlineto{\pgfqpoint{0.861367in}{0.902424in}}%
\pgfpathlineto{\pgfqpoint{0.852421in}{0.906562in}}%
\pgfpathclose%
\pgfusepath{fill}%
\end{pgfscope}%
\begin{pgfscope}%
\pgfpathrectangle{\pgfqpoint{0.407592in}{0.222904in}}{\pgfqpoint{1.411163in}{1.411163in}}%
\pgfusepath{clip}%
\pgfsetbuttcap%
\pgfsetroundjoin%
\definecolor{currentfill}{rgb}{0.909804,0.698039,0.423529}%
\pgfsetfillcolor{currentfill}%
\pgfsetfillopacity{0.120000}%
\pgfsetlinewidth{0.000000pt}%
\definecolor{currentstroke}{rgb}{0.000000,0.000000,0.000000}%
\pgfsetstrokecolor{currentstroke}%
\pgfsetdash{}{0pt}%
\pgfpathmoveto{\pgfqpoint{1.612779in}{1.090404in}}%
\pgfpathlineto{\pgfqpoint{1.630813in}{1.101634in}}%
\pgfpathlineto{\pgfqpoint{1.648927in}{1.113059in}}%
\pgfpathlineto{\pgfqpoint{1.656536in}{1.120425in}}%
\pgfpathlineto{\pgfqpoint{1.664130in}{1.128058in}}%
\pgfpathlineto{\pgfqpoint{1.646040in}{1.116764in}}%
\pgfpathlineto{\pgfqpoint{1.628030in}{1.105671in}}%
\pgfpathlineto{\pgfqpoint{1.620411in}{1.097898in}}%
\pgfpathlineto{\pgfqpoint{1.612779in}{1.090404in}}%
\pgfpathclose%
\pgfusepath{fill}%
\end{pgfscope}%
\begin{pgfscope}%
\pgfpathrectangle{\pgfqpoint{0.407592in}{0.222904in}}{\pgfqpoint{1.411163in}{1.411163in}}%
\pgfusepath{clip}%
\pgfsetbuttcap%
\pgfsetroundjoin%
\definecolor{currentfill}{rgb}{0.909804,0.698039,0.423529}%
\pgfsetfillcolor{currentfill}%
\pgfsetfillopacity{0.120000}%
\pgfsetlinewidth{0.000000pt}%
\definecolor{currentstroke}{rgb}{0.000000,0.000000,0.000000}%
\pgfsetstrokecolor{currentstroke}%
\pgfsetdash{}{0pt}%
\pgfpathmoveto{\pgfqpoint{0.682738in}{0.994408in}}%
\pgfpathlineto{\pgfqpoint{0.699353in}{0.983713in}}%
\pgfpathlineto{\pgfqpoint{0.715985in}{0.973563in}}%
\pgfpathlineto{\pgfqpoint{0.725249in}{0.969017in}}%
\pgfpathlineto{\pgfqpoint{0.734469in}{0.964673in}}%
\pgfpathlineto{\pgfqpoint{0.717886in}{0.975033in}}%
\pgfpathlineto{\pgfqpoint{0.701319in}{0.985951in}}%
\pgfpathlineto{\pgfqpoint{0.692050in}{0.990077in}}%
\pgfpathlineto{\pgfqpoint{0.682738in}{0.994408in}}%
\pgfpathclose%
\pgfusepath{fill}%
\end{pgfscope}%
\begin{pgfscope}%
\pgfpathrectangle{\pgfqpoint{0.407592in}{0.222904in}}{\pgfqpoint{1.411163in}{1.411163in}}%
\pgfusepath{clip}%
\pgfsetbuttcap%
\pgfsetroundjoin%
\definecolor{currentfill}{rgb}{0.909804,0.698039,0.423529}%
\pgfsetfillcolor{currentfill}%
\pgfsetfillopacity{0.120000}%
\pgfsetlinewidth{0.000000pt}%
\definecolor{currentstroke}{rgb}{0.000000,0.000000,0.000000}%
\pgfsetstrokecolor{currentstroke}%
\pgfsetdash{}{0pt}%
\pgfpathmoveto{\pgfqpoint{1.495032in}{1.007740in}}%
\pgfpathlineto{\pgfqpoint{1.512820in}{1.018345in}}%
\pgfpathlineto{\pgfqpoint{1.530687in}{1.029188in}}%
\pgfpathlineto{\pgfqpoint{1.538438in}{1.034971in}}%
\pgfpathlineto{\pgfqpoint{1.546173in}{1.041071in}}%
\pgfpathlineto{\pgfqpoint{1.528330in}{1.030286in}}%
\pgfpathlineto{\pgfqpoint{1.510566in}{1.019738in}}%
\pgfpathlineto{\pgfqpoint{1.502807in}{1.013572in}}%
\pgfpathlineto{\pgfqpoint{1.495032in}{1.007740in}}%
\pgfpathclose%
\pgfusepath{fill}%
\end{pgfscope}%
\begin{pgfscope}%
\pgfpathrectangle{\pgfqpoint{0.407592in}{0.222904in}}{\pgfqpoint{1.411163in}{1.411163in}}%
\pgfusepath{clip}%
\pgfsetbuttcap%
\pgfsetroundjoin%
\definecolor{currentfill}{rgb}{0.909804,0.698039,0.423529}%
\pgfsetfillcolor{currentfill}%
\pgfsetfillopacity{0.120000}%
\pgfsetlinewidth{0.000000pt}%
\definecolor{currentstroke}{rgb}{0.000000,0.000000,0.000000}%
\pgfsetstrokecolor{currentstroke}%
\pgfsetdash{}{0pt}%
\pgfpathmoveto{\pgfqpoint{1.224872in}{0.880702in}}%
\pgfpathlineto{\pgfqpoint{1.242081in}{0.887146in}}%
\pgfpathlineto{\pgfqpoint{1.259358in}{0.894093in}}%
\pgfpathlineto{\pgfqpoint{1.267503in}{0.895145in}}%
\pgfpathlineto{\pgfqpoint{1.275621in}{0.896583in}}%
\pgfpathlineto{\pgfqpoint{1.258375in}{0.889431in}}%
\pgfpathlineto{\pgfqpoint{1.241197in}{0.882778in}}%
\pgfpathlineto{\pgfqpoint{1.233048in}{0.881540in}}%
\pgfpathlineto{\pgfqpoint{1.224872in}{0.880702in}}%
\pgfpathclose%
\pgfusepath{fill}%
\end{pgfscope}%
\begin{pgfscope}%
\pgfpathrectangle{\pgfqpoint{0.407592in}{0.222904in}}{\pgfqpoint{1.411163in}{1.411163in}}%
\pgfusepath{clip}%
\pgfsetbuttcap%
\pgfsetroundjoin%
\definecolor{currentfill}{rgb}{0.909804,0.698039,0.423529}%
\pgfsetfillcolor{currentfill}%
\pgfsetfillopacity{0.120000}%
\pgfsetlinewidth{0.000000pt}%
\definecolor{currentstroke}{rgb}{0.000000,0.000000,0.000000}%
\pgfsetstrokecolor{currentstroke}%
\pgfsetdash{}{0pt}%
\pgfpathmoveto{\pgfqpoint{1.174144in}{0.869424in}}%
\pgfpathlineto{\pgfqpoint{1.191256in}{0.874516in}}%
\pgfpathlineto{\pgfqpoint{1.208432in}{0.880184in}}%
\pgfpathlineto{\pgfqpoint{1.216667in}{0.880254in}}%
\pgfpathlineto{\pgfqpoint{1.224872in}{0.880702in}}%
\pgfpathlineto{\pgfqpoint{1.207730in}{0.874796in}}%
\pgfpathlineto{\pgfqpoint{1.190652in}{0.869466in}}%
\pgfpathlineto{\pgfqpoint{1.182413in}{0.869250in}}%
\pgfpathlineto{\pgfqpoint{1.174144in}{0.869424in}}%
\pgfpathclose%
\pgfusepath{fill}%
\end{pgfscope}%
\begin{pgfscope}%
\pgfpathrectangle{\pgfqpoint{0.407592in}{0.222904in}}{\pgfqpoint{1.411163in}{1.411163in}}%
\pgfusepath{clip}%
\pgfsetbuttcap%
\pgfsetroundjoin%
\definecolor{currentfill}{rgb}{0.909804,0.698039,0.423529}%
\pgfsetfillcolor{currentfill}%
\pgfsetfillopacity{0.120000}%
\pgfsetlinewidth{0.000000pt}%
\definecolor{currentstroke}{rgb}{0.000000,0.000000,0.000000}%
\pgfsetstrokecolor{currentstroke}%
\pgfsetdash{}{0pt}%
\pgfpathmoveto{\pgfqpoint{0.937030in}{0.879215in}}%
\pgfpathlineto{\pgfqpoint{0.953804in}{0.876321in}}%
\pgfpathlineto{\pgfqpoint{0.970617in}{0.874197in}}%
\pgfpathlineto{\pgfqpoint{0.979337in}{0.870598in}}%
\pgfpathlineto{\pgfqpoint{0.988017in}{0.867267in}}%
\pgfpathlineto{\pgfqpoint{0.971251in}{0.869294in}}%
\pgfpathlineto{\pgfqpoint{0.954524in}{0.872122in}}%
\pgfpathlineto{\pgfqpoint{0.945797in}{0.875535in}}%
\pgfpathlineto{\pgfqpoint{0.937030in}{0.879215in}}%
\pgfpathclose%
\pgfusepath{fill}%
\end{pgfscope}%
\begin{pgfscope}%
\pgfpathrectangle{\pgfqpoint{0.407592in}{0.222904in}}{\pgfqpoint{1.411163in}{1.411163in}}%
\pgfusepath{clip}%
\pgfsetbuttcap%
\pgfsetroundjoin%
\definecolor{currentfill}{rgb}{0.909804,0.698039,0.423529}%
\pgfsetfillcolor{currentfill}%
\pgfsetfillopacity{0.120000}%
\pgfsetlinewidth{0.000000pt}%
\definecolor{currentstroke}{rgb}{0.000000,0.000000,0.000000}%
\pgfsetstrokecolor{currentstroke}%
\pgfsetdash{}{0pt}%
\pgfpathmoveto{\pgfqpoint{1.275621in}{0.896583in}}%
\pgfpathlineto{\pgfqpoint{1.292937in}{0.904199in}}%
\pgfpathlineto{\pgfqpoint{1.310325in}{0.912246in}}%
\pgfpathlineto{\pgfqpoint{1.318388in}{0.914250in}}%
\pgfpathlineto{\pgfqpoint{1.326425in}{0.916638in}}%
\pgfpathlineto{\pgfqpoint{1.309066in}{0.908431in}}%
\pgfpathlineto{\pgfqpoint{1.291778in}{0.900649in}}%
\pgfpathlineto{\pgfqpoint{1.283712in}{0.898415in}}%
\pgfpathlineto{\pgfqpoint{1.275621in}{0.896583in}}%
\pgfpathclose%
\pgfusepath{fill}%
\end{pgfscope}%
\begin{pgfscope}%
\pgfpathrectangle{\pgfqpoint{0.407592in}{0.222904in}}{\pgfqpoint{1.411163in}{1.411163in}}%
\pgfusepath{clip}%
\pgfsetbuttcap%
\pgfsetroundjoin%
\definecolor{currentfill}{rgb}{0.909804,0.698039,0.423529}%
\pgfsetfillcolor{currentfill}%
\pgfsetfillopacity{0.120000}%
\pgfsetlinewidth{0.000000pt}%
\definecolor{currentstroke}{rgb}{0.000000,0.000000,0.000000}%
\pgfsetstrokecolor{currentstroke}%
\pgfsetdash{}{0pt}%
\pgfpathmoveto{\pgfqpoint{1.123397in}{0.863066in}}%
\pgfpathlineto{\pgfqpoint{1.140424in}{0.866651in}}%
\pgfpathlineto{\pgfqpoint{1.157510in}{0.870880in}}%
\pgfpathlineto{\pgfqpoint{1.165843in}{0.869971in}}%
\pgfpathlineto{\pgfqpoint{1.174144in}{0.869424in}}%
\pgfpathlineto{\pgfqpoint{1.157094in}{0.864943in}}%
\pgfpathlineto{\pgfqpoint{1.140105in}{0.861114in}}%
\pgfpathlineto{\pgfqpoint{1.131767in}{0.861907in}}%
\pgfpathlineto{\pgfqpoint{1.123397in}{0.863066in}}%
\pgfpathclose%
\pgfusepath{fill}%
\end{pgfscope}%
\begin{pgfscope}%
\pgfpathrectangle{\pgfqpoint{0.407592in}{0.222904in}}{\pgfqpoint{1.411163in}{1.411163in}}%
\pgfusepath{clip}%
\pgfsetbuttcap%
\pgfsetroundjoin%
\definecolor{currentfill}{rgb}{0.909804,0.698039,0.423529}%
\pgfsetfillcolor{currentfill}%
\pgfsetfillopacity{0.120000}%
\pgfsetlinewidth{0.000000pt}%
\definecolor{currentstroke}{rgb}{0.000000,0.000000,0.000000}%
\pgfsetstrokecolor{currentstroke}%
\pgfsetdash{}{0pt}%
\pgfpathmoveto{\pgfqpoint{1.664130in}{1.128058in}}%
\pgfpathlineto{\pgfqpoint{1.682301in}{1.139542in}}%
\pgfpathlineto{\pgfqpoint{1.700554in}{1.151209in}}%
\pgfpathlineto{\pgfqpoint{1.708110in}{1.158961in}}%
\pgfpathlineto{\pgfqpoint{1.715651in}{1.166960in}}%
\pgfpathlineto{\pgfqpoint{1.697423in}{1.155439in}}%
\pgfpathlineto{\pgfqpoint{1.679277in}{1.144104in}}%
\pgfpathlineto{\pgfqpoint{1.671710in}{1.135952in}}%
\pgfpathlineto{\pgfqpoint{1.664130in}{1.128058in}}%
\pgfpathclose%
\pgfusepath{fill}%
\end{pgfscope}%
\begin{pgfscope}%
\pgfpathrectangle{\pgfqpoint{0.407592in}{0.222904in}}{\pgfqpoint{1.411163in}{1.411163in}}%
\pgfusepath{clip}%
\pgfsetbuttcap%
\pgfsetroundjoin%
\definecolor{currentfill}{rgb}{0.909804,0.698039,0.423529}%
\pgfsetfillcolor{currentfill}%
\pgfsetfillopacity{0.120000}%
\pgfsetlinewidth{0.000000pt}%
\definecolor{currentstroke}{rgb}{0.000000,0.000000,0.000000}%
\pgfsetstrokecolor{currentstroke}%
\pgfsetdash{}{0pt}%
\pgfpathmoveto{\pgfqpoint{1.326425in}{0.916638in}}%
\pgfpathlineto{\pgfqpoint{1.343857in}{0.925242in}}%
\pgfpathlineto{\pgfqpoint{1.361364in}{0.934214in}}%
\pgfpathlineto{\pgfqpoint{1.369350in}{0.937115in}}%
\pgfpathlineto{\pgfqpoint{1.377313in}{0.940390in}}%
\pgfpathlineto{\pgfqpoint{1.359833in}{0.931307in}}%
\pgfpathlineto{\pgfqpoint{1.342428in}{0.922588in}}%
\pgfpathlineto{\pgfqpoint{1.334438in}{0.919416in}}%
\pgfpathlineto{\pgfqpoint{1.326425in}{0.916638in}}%
\pgfpathclose%
\pgfusepath{fill}%
\end{pgfscope}%
\begin{pgfscope}%
\pgfpathrectangle{\pgfqpoint{0.407592in}{0.222904in}}{\pgfqpoint{1.411163in}{1.411163in}}%
\pgfusepath{clip}%
\pgfsetbuttcap%
\pgfsetroundjoin%
\definecolor{currentfill}{rgb}{0.909804,0.698039,0.423529}%
\pgfsetfillcolor{currentfill}%
\pgfsetfillopacity{0.120000}%
\pgfsetlinewidth{0.000000pt}%
\definecolor{currentstroke}{rgb}{0.000000,0.000000,0.000000}%
\pgfsetstrokecolor{currentstroke}%
\pgfsetdash{}{0pt}%
\pgfpathmoveto{\pgfqpoint{1.715651in}{1.166960in}}%
\pgfpathlineto{\pgfqpoint{1.733962in}{1.178660in}}%
\pgfpathlineto{\pgfqpoint{1.741478in}{1.186824in}}%
\pgfpathlineto{\pgfqpoint{1.748980in}{1.195221in}}%
\pgfpathlineto{\pgfqpoint{1.730696in}{1.183680in}}%
\pgfpathlineto{\pgfqpoint{1.723180in}{1.175201in}}%
\pgfusepath{fill}%
\end{pgfscope}%
\begin{pgfscope}%
\pgfpathrectangle{\pgfqpoint{0.407592in}{0.222904in}}{\pgfqpoint{1.411163in}{1.411163in}}%
\pgfusepath{clip}%
\pgfsetbuttcap%
\pgfsetroundjoin%
\definecolor{currentfill}{rgb}{0.909804,0.698039,0.423529}%
\pgfsetfillcolor{currentfill}%
\pgfsetfillopacity{0.120000}%
\pgfsetlinewidth{0.000000pt}%
\definecolor{currentstroke}{rgb}{0.000000,0.000000,0.000000}%
\pgfsetstrokecolor{currentstroke}%
\pgfsetdash{}{0pt}%
\pgfpathmoveto{\pgfqpoint{1.072588in}{0.861791in}}%
\pgfpathlineto{\pgfqpoint{1.089544in}{0.863757in}}%
\pgfpathlineto{\pgfqpoint{1.106554in}{0.866422in}}%
\pgfpathlineto{\pgfqpoint{1.114992in}{0.864577in}}%
\pgfpathlineto{\pgfqpoint{1.123397in}{0.863066in}}%
\pgfpathlineto{\pgfqpoint{1.106427in}{0.860160in}}%
\pgfpathlineto{\pgfqpoint{1.089513in}{0.857967in}}%
\pgfpathlineto{\pgfqpoint{1.081068in}{0.859710in}}%
\pgfpathlineto{\pgfqpoint{1.072588in}{0.861791in}}%
\pgfpathclose%
\pgfusepath{fill}%
\end{pgfscope}%
\begin{pgfscope}%
\pgfpathrectangle{\pgfqpoint{0.407592in}{0.222904in}}{\pgfqpoint{1.411163in}{1.411163in}}%
\pgfusepath{clip}%
\pgfsetbuttcap%
\pgfsetroundjoin%
\definecolor{currentfill}{rgb}{0.909804,0.698039,0.423529}%
\pgfsetfillcolor{currentfill}%
\pgfsetfillopacity{0.120000}%
\pgfsetlinewidth{0.000000pt}%
\definecolor{currentstroke}{rgb}{0.000000,0.000000,0.000000}%
\pgfsetstrokecolor{currentstroke}%
\pgfsetdash{}{0pt}%
\pgfpathmoveto{\pgfqpoint{1.546173in}{1.041071in}}%
\pgfpathlineto{\pgfqpoint{1.564095in}{1.052081in}}%
\pgfpathlineto{\pgfqpoint{1.582098in}{1.063303in}}%
\pgfpathlineto{\pgfqpoint{1.589791in}{1.069639in}}%
\pgfpathlineto{\pgfqpoint{1.597469in}{1.076270in}}%
\pgfpathlineto{\pgfqpoint{1.579491in}{1.065133in}}%
\pgfpathlineto{\pgfqpoint{1.561593in}{1.054211in}}%
\pgfpathlineto{\pgfqpoint{1.553891in}{1.047486in}}%
\pgfpathlineto{\pgfqpoint{1.546173in}{1.041071in}}%
\pgfpathclose%
\pgfusepath{fill}%
\end{pgfscope}%
\begin{pgfscope}%
\pgfpathrectangle{\pgfqpoint{0.407592in}{0.222904in}}{\pgfqpoint{1.411163in}{1.411163in}}%
\pgfusepath{clip}%
\pgfsetbuttcap%
\pgfsetroundjoin%
\definecolor{currentfill}{rgb}{0.909804,0.698039,0.423529}%
\pgfsetfillcolor{currentfill}%
\pgfsetfillopacity{0.120000}%
\pgfsetlinewidth{0.000000pt}%
\definecolor{currentstroke}{rgb}{0.000000,0.000000,0.000000}%
\pgfsetstrokecolor{currentstroke}%
\pgfsetdash{}{0pt}%
\pgfpathmoveto{\pgfqpoint{0.801000in}{0.929284in}}%
\pgfpathlineto{\pgfqpoint{0.817689in}{0.922053in}}%
\pgfpathlineto{\pgfqpoint{0.834403in}{0.915504in}}%
\pgfpathlineto{\pgfqpoint{0.843433in}{0.910926in}}%
\pgfpathlineto{\pgfqpoint{0.852421in}{0.906562in}}%
\pgfpathlineto{\pgfqpoint{0.835756in}{0.913201in}}%
\pgfpathlineto{\pgfqpoint{0.819116in}{0.920549in}}%
\pgfpathlineto{\pgfqpoint{0.810079in}{0.924808in}}%
\pgfpathlineto{\pgfqpoint{0.801000in}{0.929284in}}%
\pgfpathclose%
\pgfusepath{fill}%
\end{pgfscope}%
\begin{pgfscope}%
\pgfpathrectangle{\pgfqpoint{0.407592in}{0.222904in}}{\pgfqpoint{1.411163in}{1.411163in}}%
\pgfusepath{clip}%
\pgfsetbuttcap%
\pgfsetroundjoin%
\definecolor{currentfill}{rgb}{0.909804,0.698039,0.423529}%
\pgfsetfillcolor{currentfill}%
\pgfsetfillopacity{0.120000}%
\pgfsetlinewidth{0.000000pt}%
\definecolor{currentstroke}{rgb}{0.000000,0.000000,0.000000}%
\pgfsetstrokecolor{currentstroke}%
\pgfsetdash{}{0pt}%
\pgfpathmoveto{\pgfqpoint{0.715985in}{0.973563in}}%
\pgfpathlineto{\pgfqpoint{0.732635in}{0.963981in}}%
\pgfpathlineto{\pgfqpoint{0.749306in}{0.954986in}}%
\pgfpathlineto{\pgfqpoint{0.758521in}{0.950256in}}%
\pgfpathlineto{\pgfqpoint{0.767693in}{0.945724in}}%
\pgfpathlineto{\pgfqpoint{0.751071in}{0.954896in}}%
\pgfpathlineto{\pgfqpoint{0.734469in}{0.964673in}}%
\pgfpathlineto{\pgfqpoint{0.725249in}{0.969017in}}%
\pgfpathlineto{\pgfqpoint{0.715985in}{0.973563in}}%
\pgfpathclose%
\pgfusepath{fill}%
\end{pgfscope}%
\begin{pgfscope}%
\pgfpathrectangle{\pgfqpoint{0.407592in}{0.222904in}}{\pgfqpoint{1.411163in}{1.411163in}}%
\pgfusepath{clip}%
\pgfsetbuttcap%
\pgfsetroundjoin%
\definecolor{currentfill}{rgb}{0.909804,0.698039,0.423529}%
\pgfsetfillcolor{currentfill}%
\pgfsetfillopacity{0.120000}%
\pgfsetlinewidth{0.000000pt}%
\definecolor{currentstroke}{rgb}{0.000000,0.000000,0.000000}%
\pgfsetstrokecolor{currentstroke}%
\pgfsetdash{}{0pt}%
\pgfpathmoveto{\pgfqpoint{1.377313in}{0.940390in}}%
\pgfpathlineto{\pgfqpoint{1.394867in}{0.949810in}}%
\pgfpathlineto{\pgfqpoint{1.412498in}{0.959546in}}%
\pgfpathlineto{\pgfqpoint{1.420413in}{0.963272in}}%
\pgfpathlineto{\pgfqpoint{1.428308in}{0.967357in}}%
\pgfpathlineto{\pgfqpoint{1.410703in}{0.957560in}}%
\pgfpathlineto{\pgfqpoint{1.393174in}{0.948074in}}%
\pgfpathlineto{\pgfqpoint{1.385254in}{0.944042in}}%
\pgfpathlineto{\pgfqpoint{1.377313in}{0.940390in}}%
\pgfpathclose%
\pgfusepath{fill}%
\end{pgfscope}%
\begin{pgfscope}%
\pgfpathrectangle{\pgfqpoint{0.407592in}{0.222904in}}{\pgfqpoint{1.411163in}{1.411163in}}%
\pgfusepath{clip}%
\pgfsetbuttcap%
\pgfsetroundjoin%
\definecolor{currentfill}{rgb}{0.909804,0.698039,0.423529}%
\pgfsetfillcolor{currentfill}%
\pgfsetfillopacity{0.120000}%
\pgfsetlinewidth{0.000000pt}%
\definecolor{currentstroke}{rgb}{0.000000,0.000000,0.000000}%
\pgfsetstrokecolor{currentstroke}%
\pgfsetdash{}{0pt}%
\pgfpathmoveto{\pgfqpoint{0.885835in}{0.895462in}}%
\pgfpathlineto{\pgfqpoint{0.902588in}{0.891019in}}%
\pgfpathlineto{\pgfqpoint{0.919375in}{0.887319in}}%
\pgfpathlineto{\pgfqpoint{0.928222in}{0.883148in}}%
\pgfpathlineto{\pgfqpoint{0.937030in}{0.879215in}}%
\pgfpathlineto{\pgfqpoint{0.920291in}{0.882884in}}%
\pgfpathlineto{\pgfqpoint{0.903587in}{0.887328in}}%
\pgfpathlineto{\pgfqpoint{0.894731in}{0.891276in}}%
\pgfpathlineto{\pgfqpoint{0.885835in}{0.895462in}}%
\pgfpathclose%
\pgfusepath{fill}%
\end{pgfscope}%
\begin{pgfscope}%
\pgfpathrectangle{\pgfqpoint{0.407592in}{0.222904in}}{\pgfqpoint{1.411163in}{1.411163in}}%
\pgfusepath{clip}%
\pgfsetbuttcap%
\pgfsetroundjoin%
\definecolor{currentfill}{rgb}{0.909804,0.698039,0.423529}%
\pgfsetfillcolor{currentfill}%
\pgfsetfillopacity{0.120000}%
\pgfsetlinewidth{0.000000pt}%
\definecolor{currentstroke}{rgb}{0.000000,0.000000,0.000000}%
\pgfsetstrokecolor{currentstroke}%
\pgfsetdash{}{0pt}%
\pgfpathmoveto{\pgfqpoint{1.021676in}{0.865571in}}%
\pgfpathlineto{\pgfqpoint{1.038574in}{0.865867in}}%
\pgfpathlineto{\pgfqpoint{1.055521in}{0.866897in}}%
\pgfpathlineto{\pgfqpoint{1.064073in}{0.864192in}}%
\pgfpathlineto{\pgfqpoint{1.072588in}{0.861791in}}%
\pgfpathlineto{\pgfqpoint{1.055685in}{0.860553in}}%
\pgfpathlineto{\pgfqpoint{1.038832in}{0.860072in}}%
\pgfpathlineto{\pgfqpoint{1.030272in}{0.862670in}}%
\pgfpathlineto{\pgfqpoint{1.021676in}{0.865571in}}%
\pgfpathclose%
\pgfusepath{fill}%
\end{pgfscope}%
\begin{pgfscope}%
\pgfpathrectangle{\pgfqpoint{0.407592in}{0.222904in}}{\pgfqpoint{1.411163in}{1.411163in}}%
\pgfusepath{clip}%
\pgfsetbuttcap%
\pgfsetroundjoin%
\definecolor{currentfill}{rgb}{0.909804,0.698039,0.423529}%
\pgfsetfillcolor{currentfill}%
\pgfsetfillopacity{0.120000}%
\pgfsetlinewidth{0.000000pt}%
\definecolor{currentstroke}{rgb}{0.000000,0.000000,0.000000}%
\pgfsetstrokecolor{currentstroke}%
\pgfsetdash{}{0pt}%
\pgfpathmoveto{\pgfqpoint{1.428308in}{0.967357in}}%
\pgfpathlineto{\pgfqpoint{1.445989in}{0.977444in}}%
\pgfpathlineto{\pgfqpoint{1.463749in}{0.987802in}}%
\pgfpathlineto{\pgfqpoint{1.471598in}{0.992276in}}%
\pgfpathlineto{\pgfqpoint{1.479428in}{0.997090in}}%
\pgfpathlineto{\pgfqpoint{1.461693in}{0.986715in}}%
\pgfpathlineto{\pgfqpoint{1.444036in}{0.976610in}}%
\pgfpathlineto{\pgfqpoint{1.436182in}{0.971803in}}%
\pgfpathlineto{\pgfqpoint{1.428308in}{0.967357in}}%
\pgfpathclose%
\pgfusepath{fill}%
\end{pgfscope}%
\begin{pgfscope}%
\pgfpathrectangle{\pgfqpoint{0.407592in}{0.222904in}}{\pgfqpoint{1.411163in}{1.411163in}}%
\pgfusepath{clip}%
\pgfsetbuttcap%
\pgfsetroundjoin%
\definecolor{currentfill}{rgb}{0.909804,0.698039,0.423529}%
\pgfsetfillcolor{currentfill}%
\pgfsetfillopacity{0.120000}%
\pgfsetlinewidth{0.000000pt}%
\definecolor{currentstroke}{rgb}{0.000000,0.000000,0.000000}%
\pgfsetstrokecolor{currentstroke}%
\pgfsetdash{}{0pt}%
\pgfpathmoveto{\pgfqpoint{1.597469in}{1.076270in}}%
\pgfpathlineto{\pgfqpoint{1.615527in}{1.087610in}}%
\pgfpathlineto{\pgfqpoint{1.633667in}{1.099142in}}%
\pgfpathlineto{\pgfqpoint{1.641305in}{1.105963in}}%
\pgfpathlineto{\pgfqpoint{1.648927in}{1.113059in}}%
\pgfpathlineto{\pgfqpoint{1.630813in}{1.101634in}}%
\pgfpathlineto{\pgfqpoint{1.612779in}{1.090404in}}%
\pgfpathlineto{\pgfqpoint{1.605131in}{1.083193in}}%
\pgfpathlineto{\pgfqpoint{1.597469in}{1.076270in}}%
\pgfpathclose%
\pgfusepath{fill}%
\end{pgfscope}%
\begin{pgfscope}%
\pgfpathrectangle{\pgfqpoint{0.407592in}{0.222904in}}{\pgfqpoint{1.411163in}{1.411163in}}%
\pgfusepath{clip}%
\pgfsetbuttcap%
\pgfsetroundjoin%
\definecolor{currentfill}{rgb}{0.909804,0.698039,0.423529}%
\pgfsetfillcolor{currentfill}%
\pgfsetfillopacity{0.120000}%
\pgfsetlinewidth{0.000000pt}%
\definecolor{currentstroke}{rgb}{0.000000,0.000000,0.000000}%
\pgfsetstrokecolor{currentstroke}%
\pgfsetdash{}{0pt}%
\pgfpathmoveto{\pgfqpoint{0.970617in}{0.874197in}}%
\pgfpathlineto{\pgfqpoint{0.987471in}{0.872834in}}%
\pgfpathlineto{\pgfqpoint{1.004369in}{0.872218in}}%
\pgfpathlineto{\pgfqpoint{1.013042in}{0.868758in}}%
\pgfpathlineto{\pgfqpoint{1.021676in}{0.865571in}}%
\pgfpathlineto{\pgfqpoint{1.004824in}{0.866031in}}%
\pgfpathlineto{\pgfqpoint{0.988017in}{0.867267in}}%
\pgfpathlineto{\pgfqpoint{0.979337in}{0.870598in}}%
\pgfpathlineto{\pgfqpoint{0.970617in}{0.874197in}}%
\pgfpathclose%
\pgfusepath{fill}%
\end{pgfscope}%
\begin{pgfscope}%
\pgfpathrectangle{\pgfqpoint{0.407592in}{0.222904in}}{\pgfqpoint{1.411163in}{1.411163in}}%
\pgfusepath{clip}%
\pgfsetbuttcap%
\pgfsetroundjoin%
\definecolor{currentfill}{rgb}{0.909804,0.698039,0.423529}%
\pgfsetfillcolor{currentfill}%
\pgfsetfillopacity{0.120000}%
\pgfsetlinewidth{0.000000pt}%
\definecolor{currentstroke}{rgb}{0.000000,0.000000,0.000000}%
\pgfsetstrokecolor{currentstroke}%
\pgfsetdash{}{0pt}%
\pgfpathmoveto{\pgfqpoint{1.479428in}{0.997090in}}%
\pgfpathlineto{\pgfqpoint{1.497241in}{1.007717in}}%
\pgfpathlineto{\pgfqpoint{1.515133in}{1.018582in}}%
\pgfpathlineto{\pgfqpoint{1.522919in}{1.023724in}}%
\pgfpathlineto{\pgfqpoint{1.530687in}{1.029188in}}%
\pgfpathlineto{\pgfqpoint{1.512820in}{1.018345in}}%
\pgfpathlineto{\pgfqpoint{1.495032in}{1.007740in}}%
\pgfpathlineto{\pgfqpoint{1.487239in}{1.002245in}}%
\pgfpathlineto{\pgfqpoint{1.479428in}{0.997090in}}%
\pgfpathclose%
\pgfusepath{fill}%
\end{pgfscope}%
\begin{pgfscope}%
\pgfpathrectangle{\pgfqpoint{0.407592in}{0.222904in}}{\pgfqpoint{1.411163in}{1.411163in}}%
\pgfusepath{clip}%
\pgfsetbuttcap%
\pgfsetroundjoin%
\definecolor{currentfill}{rgb}{0.909804,0.698039,0.423529}%
\pgfsetfillcolor{currentfill}%
\pgfsetfillopacity{0.120000}%
\pgfsetlinewidth{0.000000pt}%
\definecolor{currentstroke}{rgb}{0.000000,0.000000,0.000000}%
\pgfsetstrokecolor{currentstroke}%
\pgfsetdash{}{0pt}%
\pgfpathmoveto{\pgfqpoint{0.749306in}{0.954986in}}%
\pgfpathlineto{\pgfqpoint{0.765998in}{0.946600in}}%
\pgfpathlineto{\pgfqpoint{0.782712in}{0.938839in}}%
\pgfpathlineto{\pgfqpoint{0.791878in}{0.933964in}}%
\pgfpathlineto{\pgfqpoint{0.801000in}{0.929284in}}%
\pgfpathlineto{\pgfqpoint{0.784335in}{0.937180in}}%
\pgfpathlineto{\pgfqpoint{0.767693in}{0.945724in}}%
\pgfpathlineto{\pgfqpoint{0.758521in}{0.950256in}}%
\pgfpathlineto{\pgfqpoint{0.749306in}{0.954986in}}%
\pgfpathclose%
\pgfusepath{fill}%
\end{pgfscope}%
\begin{pgfscope}%
\pgfpathrectangle{\pgfqpoint{0.407592in}{0.222904in}}{\pgfqpoint{1.411163in}{1.411163in}}%
\pgfusepath{clip}%
\pgfsetbuttcap%
\pgfsetroundjoin%
\definecolor{currentfill}{rgb}{0.909804,0.698039,0.423529}%
\pgfsetfillcolor{currentfill}%
\pgfsetfillopacity{0.120000}%
\pgfsetlinewidth{0.000000pt}%
\definecolor{currentstroke}{rgb}{0.000000,0.000000,0.000000}%
\pgfsetstrokecolor{currentstroke}%
\pgfsetdash{}{0pt}%
\pgfpathmoveto{\pgfqpoint{1.648927in}{1.113059in}}%
\pgfpathlineto{\pgfqpoint{1.667124in}{1.124671in}}%
\pgfpathlineto{\pgfqpoint{1.685403in}{1.136461in}}%
\pgfpathlineto{\pgfqpoint{1.692986in}{1.143707in}}%
\pgfpathlineto{\pgfqpoint{1.700554in}{1.151209in}}%
\pgfpathlineto{\pgfqpoint{1.682301in}{1.139542in}}%
\pgfpathlineto{\pgfqpoint{1.664130in}{1.128058in}}%
\pgfpathlineto{\pgfqpoint{1.656536in}{1.120425in}}%
\pgfpathlineto{\pgfqpoint{1.648927in}{1.113059in}}%
\pgfpathclose%
\pgfusepath{fill}%
\end{pgfscope}%
\begin{pgfscope}%
\pgfpathrectangle{\pgfqpoint{0.407592in}{0.222904in}}{\pgfqpoint{1.411163in}{1.411163in}}%
\pgfusepath{clip}%
\pgfsetbuttcap%
\pgfsetroundjoin%
\definecolor{currentfill}{rgb}{0.909804,0.698039,0.423529}%
\pgfsetfillcolor{currentfill}%
\pgfsetfillopacity{0.120000}%
\pgfsetlinewidth{0.000000pt}%
\definecolor{currentstroke}{rgb}{0.000000,0.000000,0.000000}%
\pgfsetstrokecolor{currentstroke}%
\pgfsetdash{}{0pt}%
\pgfpathmoveto{\pgfqpoint{0.834403in}{0.915504in}}%
\pgfpathlineto{\pgfqpoint{0.851145in}{0.909648in}}%
\pgfpathlineto{\pgfqpoint{0.867917in}{0.904496in}}%
\pgfpathlineto{\pgfqpoint{0.876896in}{0.899873in}}%
\pgfpathlineto{\pgfqpoint{0.885835in}{0.895462in}}%
\pgfpathlineto{\pgfqpoint{0.869113in}{0.900646in}}%
\pgfpathlineto{\pgfqpoint{0.852421in}{0.906562in}}%
\pgfpathlineto{\pgfqpoint{0.843433in}{0.910926in}}%
\pgfpathlineto{\pgfqpoint{0.834403in}{0.915504in}}%
\pgfpathclose%
\pgfusepath{fill}%
\end{pgfscope}%
\begin{pgfscope}%
\pgfpathrectangle{\pgfqpoint{0.407592in}{0.222904in}}{\pgfqpoint{1.411163in}{1.411163in}}%
\pgfusepath{clip}%
\pgfsetbuttcap%
\pgfsetroundjoin%
\definecolor{currentfill}{rgb}{0.909804,0.698039,0.423529}%
\pgfsetfillcolor{currentfill}%
\pgfsetfillopacity{0.120000}%
\pgfsetlinewidth{0.000000pt}%
\definecolor{currentstroke}{rgb}{0.000000,0.000000,0.000000}%
\pgfsetstrokecolor{currentstroke}%
\pgfsetdash{}{0pt}%
\pgfpathmoveto{\pgfqpoint{1.700554in}{1.151209in}}%
\pgfpathlineto{\pgfqpoint{1.718890in}{1.163049in}}%
\pgfpathlineto{\pgfqpoint{1.726433in}{1.170734in}}%
\pgfpathlineto{\pgfqpoint{1.733962in}{1.178660in}}%
\pgfpathlineto{\pgfqpoint{1.715651in}{1.166960in}}%
\pgfpathlineto{\pgfqpoint{1.708110in}{1.158961in}}%
\pgfusepath{fill}%
\end{pgfscope}%
\begin{pgfscope}%
\pgfpathrectangle{\pgfqpoint{0.407592in}{0.222904in}}{\pgfqpoint{1.411163in}{1.411163in}}%
\pgfusepath{clip}%
\pgfsetbuttcap%
\pgfsetroundjoin%
\definecolor{currentfill}{rgb}{0.909804,0.698039,0.423529}%
\pgfsetfillcolor{currentfill}%
\pgfsetfillopacity{0.120000}%
\pgfsetlinewidth{0.000000pt}%
\definecolor{currentstroke}{rgb}{0.000000,0.000000,0.000000}%
\pgfsetstrokecolor{currentstroke}%
\pgfsetdash{}{0pt}%
\pgfpathmoveto{\pgfqpoint{1.208432in}{0.880184in}}%
\pgfpathlineto{\pgfqpoint{1.225674in}{0.886391in}}%
\pgfpathlineto{\pgfqpoint{1.242984in}{0.893102in}}%
\pgfpathlineto{\pgfqpoint{1.251186in}{0.893415in}}%
\pgfpathlineto{\pgfqpoint{1.259358in}{0.894093in}}%
\pgfpathlineto{\pgfqpoint{1.242081in}{0.887146in}}%
\pgfpathlineto{\pgfqpoint{1.224872in}{0.880702in}}%
\pgfpathlineto{\pgfqpoint{1.216667in}{0.880254in}}%
\pgfpathlineto{\pgfqpoint{1.208432in}{0.880184in}}%
\pgfpathclose%
\pgfusepath{fill}%
\end{pgfscope}%
\begin{pgfscope}%
\pgfpathrectangle{\pgfqpoint{0.407592in}{0.222904in}}{\pgfqpoint{1.411163in}{1.411163in}}%
\pgfusepath{clip}%
\pgfsetbuttcap%
\pgfsetroundjoin%
\definecolor{currentfill}{rgb}{0.909804,0.698039,0.423529}%
\pgfsetfillcolor{currentfill}%
\pgfsetfillopacity{0.120000}%
\pgfsetlinewidth{0.000000pt}%
\definecolor{currentstroke}{rgb}{0.000000,0.000000,0.000000}%
\pgfsetstrokecolor{currentstroke}%
\pgfsetdash{}{0pt}%
\pgfpathmoveto{\pgfqpoint{1.259358in}{0.894093in}}%
\pgfpathlineto{\pgfqpoint{1.276706in}{0.901507in}}%
\pgfpathlineto{\pgfqpoint{1.294124in}{0.909358in}}%
\pgfpathlineto{\pgfqpoint{1.302238in}{0.910618in}}%
\pgfpathlineto{\pgfqpoint{1.310325in}{0.912246in}}%
\pgfpathlineto{\pgfqpoint{1.292937in}{0.904199in}}%
\pgfpathlineto{\pgfqpoint{1.275621in}{0.896583in}}%
\pgfpathlineto{\pgfqpoint{1.267503in}{0.895145in}}%
\pgfpathlineto{\pgfqpoint{1.259358in}{0.894093in}}%
\pgfpathclose%
\pgfusepath{fill}%
\end{pgfscope}%
\begin{pgfscope}%
\pgfpathrectangle{\pgfqpoint{0.407592in}{0.222904in}}{\pgfqpoint{1.411163in}{1.411163in}}%
\pgfusepath{clip}%
\pgfsetbuttcap%
\pgfsetroundjoin%
\definecolor{currentfill}{rgb}{0.909804,0.698039,0.423529}%
\pgfsetfillcolor{currentfill}%
\pgfsetfillopacity{0.120000}%
\pgfsetlinewidth{0.000000pt}%
\definecolor{currentstroke}{rgb}{0.000000,0.000000,0.000000}%
\pgfsetstrokecolor{currentstroke}%
\pgfsetdash{}{0pt}%
\pgfpathmoveto{\pgfqpoint{1.157510in}{0.870880in}}%
\pgfpathlineto{\pgfqpoint{1.174659in}{0.875716in}}%
\pgfpathlineto{\pgfqpoint{1.191871in}{0.881125in}}%
\pgfpathlineto{\pgfqpoint{1.200167in}{0.880478in}}%
\pgfpathlineto{\pgfqpoint{1.208432in}{0.880184in}}%
\pgfpathlineto{\pgfqpoint{1.191256in}{0.874516in}}%
\pgfpathlineto{\pgfqpoint{1.174144in}{0.869424in}}%
\pgfpathlineto{\pgfqpoint{1.165843in}{0.869971in}}%
\pgfpathlineto{\pgfqpoint{1.157510in}{0.870880in}}%
\pgfpathclose%
\pgfusepath{fill}%
\end{pgfscope}%
\begin{pgfscope}%
\pgfpathrectangle{\pgfqpoint{0.407592in}{0.222904in}}{\pgfqpoint{1.411163in}{1.411163in}}%
\pgfusepath{clip}%
\pgfsetbuttcap%
\pgfsetroundjoin%
\definecolor{currentfill}{rgb}{0.909804,0.698039,0.423529}%
\pgfsetfillcolor{currentfill}%
\pgfsetfillopacity{0.120000}%
\pgfsetlinewidth{0.000000pt}%
\definecolor{currentstroke}{rgb}{0.000000,0.000000,0.000000}%
\pgfsetstrokecolor{currentstroke}%
\pgfsetdash{}{0pt}%
\pgfpathmoveto{\pgfqpoint{1.530687in}{1.029188in}}%
\pgfpathlineto{\pgfqpoint{1.548635in}{1.040253in}}%
\pgfpathlineto{\pgfqpoint{1.566662in}{1.051529in}}%
\pgfpathlineto{\pgfqpoint{1.574388in}{1.057265in}}%
\pgfpathlineto{\pgfqpoint{1.582098in}{1.063303in}}%
\pgfpathlineto{\pgfqpoint{1.564095in}{1.052081in}}%
\pgfpathlineto{\pgfqpoint{1.546173in}{1.041071in}}%
\pgfpathlineto{\pgfqpoint{1.538438in}{1.034971in}}%
\pgfpathlineto{\pgfqpoint{1.530687in}{1.029188in}}%
\pgfpathclose%
\pgfusepath{fill}%
\end{pgfscope}%
\begin{pgfscope}%
\pgfpathrectangle{\pgfqpoint{0.407592in}{0.222904in}}{\pgfqpoint{1.411163in}{1.411163in}}%
\pgfusepath{clip}%
\pgfsetbuttcap%
\pgfsetroundjoin%
\definecolor{currentfill}{rgb}{0.909804,0.698039,0.423529}%
\pgfsetfillcolor{currentfill}%
\pgfsetfillopacity{0.120000}%
\pgfsetlinewidth{0.000000pt}%
\definecolor{currentstroke}{rgb}{0.000000,0.000000,0.000000}%
\pgfsetstrokecolor{currentstroke}%
\pgfsetdash{}{0pt}%
\pgfpathmoveto{\pgfqpoint{1.310325in}{0.912246in}}%
\pgfpathlineto{\pgfqpoint{1.327786in}{0.920695in}}%
\pgfpathlineto{\pgfqpoint{1.345321in}{0.929517in}}%
\pgfpathlineto{\pgfqpoint{1.353354in}{0.931683in}}%
\pgfpathlineto{\pgfqpoint{1.361364in}{0.934214in}}%
\pgfpathlineto{\pgfqpoint{1.343857in}{0.925242in}}%
\pgfpathlineto{\pgfqpoint{1.326425in}{0.916638in}}%
\pgfpathlineto{\pgfqpoint{1.318388in}{0.914250in}}%
\pgfpathlineto{\pgfqpoint{1.310325in}{0.912246in}}%
\pgfpathclose%
\pgfusepath{fill}%
\end{pgfscope}%
\begin{pgfscope}%
\pgfpathrectangle{\pgfqpoint{0.407592in}{0.222904in}}{\pgfqpoint{1.411163in}{1.411163in}}%
\pgfusepath{clip}%
\pgfsetbuttcap%
\pgfsetroundjoin%
\definecolor{currentfill}{rgb}{0.909804,0.698039,0.423529}%
\pgfsetfillcolor{currentfill}%
\pgfsetfillopacity{0.120000}%
\pgfsetlinewidth{0.000000pt}%
\definecolor{currentstroke}{rgb}{0.000000,0.000000,0.000000}%
\pgfsetstrokecolor{currentstroke}%
\pgfsetdash{}{0pt}%
\pgfpathmoveto{\pgfqpoint{1.106554in}{0.866422in}}%
\pgfpathlineto{\pgfqpoint{1.123621in}{0.869754in}}%
\pgfpathlineto{\pgfqpoint{1.140747in}{0.873720in}}%
\pgfpathlineto{\pgfqpoint{1.149145in}{0.872134in}}%
\pgfpathlineto{\pgfqpoint{1.157510in}{0.870880in}}%
\pgfpathlineto{\pgfqpoint{1.140424in}{0.866651in}}%
\pgfpathlineto{\pgfqpoint{1.123397in}{0.863066in}}%
\pgfpathlineto{\pgfqpoint{1.114992in}{0.864577in}}%
\pgfpathlineto{\pgfqpoint{1.106554in}{0.866422in}}%
\pgfpathclose%
\pgfusepath{fill}%
\end{pgfscope}%
\begin{pgfscope}%
\pgfpathrectangle{\pgfqpoint{0.407592in}{0.222904in}}{\pgfqpoint{1.411163in}{1.411163in}}%
\pgfusepath{clip}%
\pgfsetbuttcap%
\pgfsetroundjoin%
\definecolor{currentfill}{rgb}{0.909804,0.698039,0.423529}%
\pgfsetfillcolor{currentfill}%
\pgfsetfillopacity{0.120000}%
\pgfsetlinewidth{0.000000pt}%
\definecolor{currentstroke}{rgb}{0.000000,0.000000,0.000000}%
\pgfsetstrokecolor{currentstroke}%
\pgfsetdash{}{0pt}%
\pgfpathmoveto{\pgfqpoint{0.919375in}{0.887319in}}%
\pgfpathlineto{\pgfqpoint{0.936198in}{0.884363in}}%
\pgfpathlineto{\pgfqpoint{0.953060in}{0.882146in}}%
\pgfpathlineto{\pgfqpoint{0.961858in}{0.878052in}}%
\pgfpathlineto{\pgfqpoint{0.970617in}{0.874197in}}%
\pgfpathlineto{\pgfqpoint{0.953804in}{0.876321in}}%
\pgfpathlineto{\pgfqpoint{0.937030in}{0.879215in}}%
\pgfpathlineto{\pgfqpoint{0.928222in}{0.883148in}}%
\pgfpathlineto{\pgfqpoint{0.919375in}{0.887319in}}%
\pgfpathclose%
\pgfusepath{fill}%
\end{pgfscope}%
\begin{pgfscope}%
\pgfpathrectangle{\pgfqpoint{0.407592in}{0.222904in}}{\pgfqpoint{1.411163in}{1.411163in}}%
\pgfusepath{clip}%
\pgfsetbuttcap%
\pgfsetroundjoin%
\definecolor{currentfill}{rgb}{0.909804,0.698039,0.423529}%
\pgfsetfillcolor{currentfill}%
\pgfsetfillopacity{0.120000}%
\pgfsetlinewidth{0.000000pt}%
\definecolor{currentstroke}{rgb}{0.000000,0.000000,0.000000}%
\pgfsetstrokecolor{currentstroke}%
\pgfsetdash{}{0pt}%
\pgfpathmoveto{\pgfqpoint{1.361364in}{0.934214in}}%
\pgfpathlineto{\pgfqpoint{1.378945in}{0.943529in}}%
\pgfpathlineto{\pgfqpoint{1.396603in}{0.953164in}}%
\pgfpathlineto{\pgfqpoint{1.404561in}{0.956177in}}%
\pgfpathlineto{\pgfqpoint{1.412498in}{0.959546in}}%
\pgfpathlineto{\pgfqpoint{1.394867in}{0.949810in}}%
\pgfpathlineto{\pgfqpoint{1.377313in}{0.940390in}}%
\pgfpathlineto{\pgfqpoint{1.369350in}{0.937115in}}%
\pgfpathlineto{\pgfqpoint{1.361364in}{0.934214in}}%
\pgfpathclose%
\pgfusepath{fill}%
\end{pgfscope}%
\begin{pgfscope}%
\pgfpathrectangle{\pgfqpoint{0.407592in}{0.222904in}}{\pgfqpoint{1.411163in}{1.411163in}}%
\pgfusepath{clip}%
\pgfsetbuttcap%
\pgfsetroundjoin%
\definecolor{currentfill}{rgb}{0.909804,0.698039,0.423529}%
\pgfsetfillcolor{currentfill}%
\pgfsetfillopacity{0.120000}%
\pgfsetlinewidth{0.000000pt}%
\definecolor{currentstroke}{rgb}{0.000000,0.000000,0.000000}%
\pgfsetstrokecolor{currentstroke}%
\pgfsetdash{}{0pt}%
\pgfpathmoveto{\pgfqpoint{1.055521in}{0.866897in}}%
\pgfpathlineto{\pgfqpoint{1.072519in}{0.868636in}}%
\pgfpathlineto{\pgfqpoint{1.089572in}{0.871057in}}%
\pgfpathlineto{\pgfqpoint{1.098081in}{0.868587in}}%
\pgfpathlineto{\pgfqpoint{1.106554in}{0.866422in}}%
\pgfpathlineto{\pgfqpoint{1.089544in}{0.863757in}}%
\pgfpathlineto{\pgfqpoint{1.072588in}{0.861791in}}%
\pgfpathlineto{\pgfqpoint{1.064073in}{0.864192in}}%
\pgfpathlineto{\pgfqpoint{1.055521in}{0.866897in}}%
\pgfpathclose%
\pgfusepath{fill}%
\end{pgfscope}%
\begin{pgfscope}%
\pgfpathrectangle{\pgfqpoint{0.407592in}{0.222904in}}{\pgfqpoint{1.411163in}{1.411163in}}%
\pgfusepath{clip}%
\pgfsetbuttcap%
\pgfsetroundjoin%
\definecolor{currentfill}{rgb}{0.909804,0.698039,0.423529}%
\pgfsetfillcolor{currentfill}%
\pgfsetfillopacity{0.120000}%
\pgfsetlinewidth{0.000000pt}%
\definecolor{currentstroke}{rgb}{0.000000,0.000000,0.000000}%
\pgfsetstrokecolor{currentstroke}%
\pgfsetdash{}{0pt}%
\pgfpathmoveto{\pgfqpoint{1.582098in}{1.063303in}}%
\pgfpathlineto{\pgfqpoint{1.600181in}{1.074725in}}%
\pgfpathlineto{\pgfqpoint{1.618346in}{1.086338in}}%
\pgfpathlineto{\pgfqpoint{1.626015in}{1.092599in}}%
\pgfpathlineto{\pgfqpoint{1.633667in}{1.099142in}}%
\pgfpathlineto{\pgfqpoint{1.615527in}{1.087610in}}%
\pgfpathlineto{\pgfqpoint{1.597469in}{1.076270in}}%
\pgfpathlineto{\pgfqpoint{1.589791in}{1.069639in}}%
\pgfpathlineto{\pgfqpoint{1.582098in}{1.063303in}}%
\pgfpathclose%
\pgfusepath{fill}%
\end{pgfscope}%
\begin{pgfscope}%
\pgfpathrectangle{\pgfqpoint{0.407592in}{0.222904in}}{\pgfqpoint{1.411163in}{1.411163in}}%
\pgfusepath{clip}%
\pgfsetbuttcap%
\pgfsetroundjoin%
\definecolor{currentfill}{rgb}{0.909804,0.698039,0.423529}%
\pgfsetfillcolor{currentfill}%
\pgfsetfillopacity{0.120000}%
\pgfsetlinewidth{0.000000pt}%
\definecolor{currentstroke}{rgb}{0.000000,0.000000,0.000000}%
\pgfsetstrokecolor{currentstroke}%
\pgfsetdash{}{0pt}%
\pgfpathmoveto{\pgfqpoint{1.412498in}{0.959546in}}%
\pgfpathlineto{\pgfqpoint{1.430206in}{0.969575in}}%
\pgfpathlineto{\pgfqpoint{1.447992in}{0.979879in}}%
\pgfpathlineto{\pgfqpoint{1.455881in}{0.983670in}}%
\pgfpathlineto{\pgfqpoint{1.463749in}{0.987802in}}%
\pgfpathlineto{\pgfqpoint{1.445989in}{0.977444in}}%
\pgfpathlineto{\pgfqpoint{1.428308in}{0.967357in}}%
\pgfpathlineto{\pgfqpoint{1.420413in}{0.963272in}}%
\pgfpathlineto{\pgfqpoint{1.412498in}{0.959546in}}%
\pgfpathclose%
\pgfusepath{fill}%
\end{pgfscope}%
\begin{pgfscope}%
\pgfpathrectangle{\pgfqpoint{0.407592in}{0.222904in}}{\pgfqpoint{1.411163in}{1.411163in}}%
\pgfusepath{clip}%
\pgfsetbuttcap%
\pgfsetroundjoin%
\definecolor{currentfill}{rgb}{0.909804,0.698039,0.423529}%
\pgfsetfillcolor{currentfill}%
\pgfsetfillopacity{0.120000}%
\pgfsetlinewidth{0.000000pt}%
\definecolor{currentstroke}{rgb}{0.000000,0.000000,0.000000}%
\pgfsetstrokecolor{currentstroke}%
\pgfsetdash{}{0pt}%
\pgfpathmoveto{\pgfqpoint{0.782712in}{0.938839in}}%
\pgfpathlineto{\pgfqpoint{0.799451in}{0.931719in}}%
\pgfpathlineto{\pgfqpoint{0.816216in}{0.925255in}}%
\pgfpathlineto{\pgfqpoint{0.825331in}{0.920284in}}%
\pgfpathlineto{\pgfqpoint{0.834403in}{0.915504in}}%
\pgfpathlineto{\pgfqpoint{0.817689in}{0.922053in}}%
\pgfpathlineto{\pgfqpoint{0.801000in}{0.929284in}}%
\pgfpathlineto{\pgfqpoint{0.791878in}{0.933964in}}%
\pgfpathlineto{\pgfqpoint{0.782712in}{0.938839in}}%
\pgfpathclose%
\pgfusepath{fill}%
\end{pgfscope}%
\begin{pgfscope}%
\pgfpathrectangle{\pgfqpoint{0.407592in}{0.222904in}}{\pgfqpoint{1.411163in}{1.411163in}}%
\pgfusepath{clip}%
\pgfsetbuttcap%
\pgfsetroundjoin%
\definecolor{currentfill}{rgb}{0.909804,0.698039,0.423529}%
\pgfsetfillcolor{currentfill}%
\pgfsetfillopacity{0.120000}%
\pgfsetlinewidth{0.000000pt}%
\definecolor{currentstroke}{rgb}{0.000000,0.000000,0.000000}%
\pgfsetstrokecolor{currentstroke}%
\pgfsetdash{}{0pt}%
\pgfpathmoveto{\pgfqpoint{1.004369in}{0.872218in}}%
\pgfpathlineto{\pgfqpoint{1.021313in}{0.872333in}}%
\pgfpathlineto{\pgfqpoint{1.038306in}{0.873158in}}%
\pgfpathlineto{\pgfqpoint{1.046932in}{0.869890in}}%
\pgfpathlineto{\pgfqpoint{1.055521in}{0.866897in}}%
\pgfpathlineto{\pgfqpoint{1.038574in}{0.865867in}}%
\pgfpathlineto{\pgfqpoint{1.021676in}{0.865571in}}%
\pgfpathlineto{\pgfqpoint{1.013042in}{0.868758in}}%
\pgfpathlineto{\pgfqpoint{1.004369in}{0.872218in}}%
\pgfpathclose%
\pgfusepath{fill}%
\end{pgfscope}%
\begin{pgfscope}%
\pgfpathrectangle{\pgfqpoint{0.407592in}{0.222904in}}{\pgfqpoint{1.411163in}{1.411163in}}%
\pgfusepath{clip}%
\pgfsetbuttcap%
\pgfsetroundjoin%
\definecolor{currentfill}{rgb}{0.909804,0.698039,0.423529}%
\pgfsetfillcolor{currentfill}%
\pgfsetfillopacity{0.120000}%
\pgfsetlinewidth{0.000000pt}%
\definecolor{currentstroke}{rgb}{0.000000,0.000000,0.000000}%
\pgfsetstrokecolor{currentstroke}%
\pgfsetdash{}{0pt}%
\pgfpathmoveto{\pgfqpoint{0.867917in}{0.904496in}}%
\pgfpathlineto{\pgfqpoint{0.884720in}{0.900055in}}%
\pgfpathlineto{\pgfqpoint{0.901557in}{0.896327in}}%
\pgfpathlineto{\pgfqpoint{0.910487in}{0.891716in}}%
\pgfpathlineto{\pgfqpoint{0.919375in}{0.887319in}}%
\pgfpathlineto{\pgfqpoint{0.902588in}{0.891019in}}%
\pgfpathlineto{\pgfqpoint{0.885835in}{0.895462in}}%
\pgfpathlineto{\pgfqpoint{0.876896in}{0.899873in}}%
\pgfpathlineto{\pgfqpoint{0.867917in}{0.904496in}}%
\pgfpathclose%
\pgfusepath{fill}%
\end{pgfscope}%
\begin{pgfscope}%
\pgfpathrectangle{\pgfqpoint{0.407592in}{0.222904in}}{\pgfqpoint{1.411163in}{1.411163in}}%
\pgfusepath{clip}%
\pgfsetbuttcap%
\pgfsetroundjoin%
\definecolor{currentfill}{rgb}{0.909804,0.698039,0.423529}%
\pgfsetfillcolor{currentfill}%
\pgfsetfillopacity{0.120000}%
\pgfsetlinewidth{0.000000pt}%
\definecolor{currentstroke}{rgb}{0.000000,0.000000,0.000000}%
\pgfsetstrokecolor{currentstroke}%
\pgfsetdash{}{0pt}%
\pgfpathmoveto{\pgfqpoint{1.463749in}{0.987802in}}%
\pgfpathlineto{\pgfqpoint{1.481588in}{0.998415in}}%
\pgfpathlineto{\pgfqpoint{1.499505in}{1.009267in}}%
\pgfpathlineto{\pgfqpoint{1.507329in}{1.013763in}}%
\pgfpathlineto{\pgfqpoint{1.515133in}{1.018582in}}%
\pgfpathlineto{\pgfqpoint{1.497241in}{1.007717in}}%
\pgfpathlineto{\pgfqpoint{1.479428in}{0.997090in}}%
\pgfpathlineto{\pgfqpoint{1.471598in}{0.992276in}}%
\pgfpathlineto{\pgfqpoint{1.463749in}{0.987802in}}%
\pgfpathclose%
\pgfusepath{fill}%
\end{pgfscope}%
\begin{pgfscope}%
\pgfpathrectangle{\pgfqpoint{0.407592in}{0.222904in}}{\pgfqpoint{1.411163in}{1.411163in}}%
\pgfusepath{clip}%
\pgfsetbuttcap%
\pgfsetroundjoin%
\definecolor{currentfill}{rgb}{0.909804,0.698039,0.423529}%
\pgfsetfillcolor{currentfill}%
\pgfsetfillopacity{0.120000}%
\pgfsetlinewidth{0.000000pt}%
\definecolor{currentstroke}{rgb}{0.000000,0.000000,0.000000}%
\pgfsetstrokecolor{currentstroke}%
\pgfsetdash{}{0pt}%
\pgfpathmoveto{\pgfqpoint{1.633667in}{1.099142in}}%
\pgfpathlineto{\pgfqpoint{1.651889in}{1.110858in}}%
\pgfpathlineto{\pgfqpoint{1.670193in}{1.122749in}}%
\pgfpathlineto{\pgfqpoint{1.677805in}{1.129474in}}%
\pgfpathlineto{\pgfqpoint{1.685403in}{1.136461in}}%
\pgfpathlineto{\pgfqpoint{1.667124in}{1.124671in}}%
\pgfpathlineto{\pgfqpoint{1.648927in}{1.113059in}}%
\pgfpathlineto{\pgfqpoint{1.641305in}{1.105963in}}%
\pgfpathlineto{\pgfqpoint{1.633667in}{1.099142in}}%
\pgfpathclose%
\pgfusepath{fill}%
\end{pgfscope}%
\begin{pgfscope}%
\pgfpathrectangle{\pgfqpoint{0.407592in}{0.222904in}}{\pgfqpoint{1.411163in}{1.411163in}}%
\pgfusepath{clip}%
\pgfsetbuttcap%
\pgfsetroundjoin%
\definecolor{currentfill}{rgb}{0.909804,0.698039,0.423529}%
\pgfsetfillcolor{currentfill}%
\pgfsetfillopacity{0.120000}%
\pgfsetlinewidth{0.000000pt}%
\definecolor{currentstroke}{rgb}{0.000000,0.000000,0.000000}%
\pgfsetstrokecolor{currentstroke}%
\pgfsetdash{}{0pt}%
\pgfpathmoveto{\pgfqpoint{1.685403in}{1.136461in}}%
\pgfpathlineto{\pgfqpoint{1.703764in}{1.148422in}}%
\pgfpathlineto{\pgfqpoint{1.711334in}{1.155610in}}%
\pgfpathlineto{\pgfqpoint{1.718890in}{1.163049in}}%
\pgfpathlineto{\pgfqpoint{1.700554in}{1.151209in}}%
\pgfpathlineto{\pgfqpoint{1.692986in}{1.143707in}}%
\pgfusepath{fill}%
\end{pgfscope}%
\begin{pgfscope}%
\pgfpathrectangle{\pgfqpoint{0.407592in}{0.222904in}}{\pgfqpoint{1.411163in}{1.411163in}}%
\pgfusepath{clip}%
\pgfsetbuttcap%
\pgfsetroundjoin%
\definecolor{currentfill}{rgb}{0.909804,0.698039,0.423529}%
\pgfsetfillcolor{currentfill}%
\pgfsetfillopacity{0.120000}%
\pgfsetlinewidth{0.000000pt}%
\definecolor{currentstroke}{rgb}{0.000000,0.000000,0.000000}%
\pgfsetstrokecolor{currentstroke}%
\pgfsetdash{}{0pt}%
\pgfpathmoveto{\pgfqpoint{0.953060in}{0.882146in}}%
\pgfpathlineto{\pgfqpoint{0.969962in}{0.880661in}}%
\pgfpathlineto{\pgfqpoint{0.986908in}{0.879895in}}%
\pgfpathlineto{\pgfqpoint{0.995658in}{0.875934in}}%
\pgfpathlineto{\pgfqpoint{1.004369in}{0.872218in}}%
\pgfpathlineto{\pgfqpoint{0.987471in}{0.872834in}}%
\pgfpathlineto{\pgfqpoint{0.970617in}{0.874197in}}%
\pgfpathlineto{\pgfqpoint{0.961858in}{0.878052in}}%
\pgfpathlineto{\pgfqpoint{0.953060in}{0.882146in}}%
\pgfpathclose%
\pgfusepath{fill}%
\end{pgfscope}%
\begin{pgfscope}%
\pgfpathrectangle{\pgfqpoint{0.407592in}{0.222904in}}{\pgfqpoint{1.411163in}{1.411163in}}%
\pgfusepath{clip}%
\pgfsetbuttcap%
\pgfsetroundjoin%
\definecolor{currentfill}{rgb}{0.909804,0.698039,0.423529}%
\pgfsetfillcolor{currentfill}%
\pgfsetfillopacity{0.120000}%
\pgfsetlinewidth{0.000000pt}%
\definecolor{currentstroke}{rgb}{0.000000,0.000000,0.000000}%
\pgfsetstrokecolor{currentstroke}%
\pgfsetdash{}{0pt}%
\pgfpathmoveto{\pgfqpoint{1.515133in}{1.018582in}}%
\pgfpathlineto{\pgfqpoint{1.533105in}{1.029670in}}%
\pgfpathlineto{\pgfqpoint{1.551158in}{1.040969in}}%
\pgfpathlineto{\pgfqpoint{1.558919in}{1.046096in}}%
\pgfpathlineto{\pgfqpoint{1.566662in}{1.051529in}}%
\pgfpathlineto{\pgfqpoint{1.548635in}{1.040253in}}%
\pgfpathlineto{\pgfqpoint{1.530687in}{1.029188in}}%
\pgfpathlineto{\pgfqpoint{1.522919in}{1.023724in}}%
\pgfpathlineto{\pgfqpoint{1.515133in}{1.018582in}}%
\pgfpathclose%
\pgfusepath{fill}%
\end{pgfscope}%
\begin{pgfscope}%
\pgfpathrectangle{\pgfqpoint{0.407592in}{0.222904in}}{\pgfqpoint{1.411163in}{1.411163in}}%
\pgfusepath{clip}%
\pgfsetbuttcap%
\pgfsetroundjoin%
\definecolor{currentfill}{rgb}{0.909804,0.698039,0.423529}%
\pgfsetfillcolor{currentfill}%
\pgfsetfillopacity{0.120000}%
\pgfsetlinewidth{0.000000pt}%
\definecolor{currentstroke}{rgb}{0.000000,0.000000,0.000000}%
\pgfsetstrokecolor{currentstroke}%
\pgfsetdash{}{0pt}%
\pgfpathmoveto{\pgfqpoint{1.242984in}{0.893102in}}%
\pgfpathlineto{\pgfqpoint{1.260364in}{0.900287in}}%
\pgfpathlineto{\pgfqpoint{1.277814in}{0.907913in}}%
\pgfpathlineto{\pgfqpoint{1.285983in}{0.908460in}}%
\pgfpathlineto{\pgfqpoint{1.294124in}{0.909358in}}%
\pgfpathlineto{\pgfqpoint{1.276706in}{0.901507in}}%
\pgfpathlineto{\pgfqpoint{1.259358in}{0.894093in}}%
\pgfpathlineto{\pgfqpoint{1.251186in}{0.893415in}}%
\pgfpathlineto{\pgfqpoint{1.242984in}{0.893102in}}%
\pgfpathclose%
\pgfusepath{fill}%
\end{pgfscope}%
\begin{pgfscope}%
\pgfpathrectangle{\pgfqpoint{0.407592in}{0.222904in}}{\pgfqpoint{1.411163in}{1.411163in}}%
\pgfusepath{clip}%
\pgfsetbuttcap%
\pgfsetroundjoin%
\definecolor{currentfill}{rgb}{0.909804,0.698039,0.423529}%
\pgfsetfillcolor{currentfill}%
\pgfsetfillopacity{0.120000}%
\pgfsetlinewidth{0.000000pt}%
\definecolor{currentstroke}{rgb}{0.000000,0.000000,0.000000}%
\pgfsetstrokecolor{currentstroke}%
\pgfsetdash{}{0pt}%
\pgfpathmoveto{\pgfqpoint{1.191871in}{0.881125in}}%
\pgfpathlineto{\pgfqpoint{1.209148in}{0.887074in}}%
\pgfpathlineto{\pgfqpoint{1.226493in}{0.893529in}}%
\pgfpathlineto{\pgfqpoint{1.234754in}{0.893144in}}%
\pgfpathlineto{\pgfqpoint{1.242984in}{0.893102in}}%
\pgfpathlineto{\pgfqpoint{1.225674in}{0.886391in}}%
\pgfpathlineto{\pgfqpoint{1.208432in}{0.880184in}}%
\pgfpathlineto{\pgfqpoint{1.200167in}{0.880478in}}%
\pgfpathlineto{\pgfqpoint{1.191871in}{0.881125in}}%
\pgfpathclose%
\pgfusepath{fill}%
\end{pgfscope}%
\begin{pgfscope}%
\pgfpathrectangle{\pgfqpoint{0.407592in}{0.222904in}}{\pgfqpoint{1.411163in}{1.411163in}}%
\pgfusepath{clip}%
\pgfsetbuttcap%
\pgfsetroundjoin%
\definecolor{currentfill}{rgb}{0.909804,0.698039,0.423529}%
\pgfsetfillcolor{currentfill}%
\pgfsetfillopacity{0.120000}%
\pgfsetlinewidth{0.000000pt}%
\definecolor{currentstroke}{rgb}{0.000000,0.000000,0.000000}%
\pgfsetstrokecolor{currentstroke}%
\pgfsetdash{}{0pt}%
\pgfpathmoveto{\pgfqpoint{0.816216in}{0.925255in}}%
\pgfpathlineto{\pgfqpoint{0.833008in}{0.919459in}}%
\pgfpathlineto{\pgfqpoint{0.849831in}{0.914337in}}%
\pgfpathlineto{\pgfqpoint{0.858895in}{0.909321in}}%
\pgfpathlineto{\pgfqpoint{0.867917in}{0.904496in}}%
\pgfpathlineto{\pgfqpoint{0.851145in}{0.909648in}}%
\pgfpathlineto{\pgfqpoint{0.834403in}{0.915504in}}%
\pgfpathlineto{\pgfqpoint{0.825331in}{0.920284in}}%
\pgfpathlineto{\pgfqpoint{0.816216in}{0.925255in}}%
\pgfpathclose%
\pgfusepath{fill}%
\end{pgfscope}%
\begin{pgfscope}%
\pgfpathrectangle{\pgfqpoint{0.407592in}{0.222904in}}{\pgfqpoint{1.411163in}{1.411163in}}%
\pgfusepath{clip}%
\pgfsetbuttcap%
\pgfsetroundjoin%
\definecolor{currentfill}{rgb}{0.909804,0.698039,0.423529}%
\pgfsetfillcolor{currentfill}%
\pgfsetfillopacity{0.120000}%
\pgfsetlinewidth{0.000000pt}%
\definecolor{currentstroke}{rgb}{0.000000,0.000000,0.000000}%
\pgfsetstrokecolor{currentstroke}%
\pgfsetdash{}{0pt}%
\pgfpathmoveto{\pgfqpoint{1.294124in}{0.909358in}}%
\pgfpathlineto{\pgfqpoint{1.311614in}{0.917618in}}%
\pgfpathlineto{\pgfqpoint{1.329179in}{0.926257in}}%
\pgfpathlineto{\pgfqpoint{1.337263in}{0.927711in}}%
\pgfpathlineto{\pgfqpoint{1.345321in}{0.929517in}}%
\pgfpathlineto{\pgfqpoint{1.327786in}{0.920695in}}%
\pgfpathlineto{\pgfqpoint{1.310325in}{0.912246in}}%
\pgfpathlineto{\pgfqpoint{1.302238in}{0.910618in}}%
\pgfpathlineto{\pgfqpoint{1.294124in}{0.909358in}}%
\pgfpathclose%
\pgfusepath{fill}%
\end{pgfscope}%
\begin{pgfscope}%
\pgfpathrectangle{\pgfqpoint{0.407592in}{0.222904in}}{\pgfqpoint{1.411163in}{1.411163in}}%
\pgfusepath{clip}%
\pgfsetbuttcap%
\pgfsetroundjoin%
\definecolor{currentfill}{rgb}{0.909804,0.698039,0.423529}%
\pgfsetfillcolor{currentfill}%
\pgfsetfillopacity{0.120000}%
\pgfsetlinewidth{0.000000pt}%
\definecolor{currentstroke}{rgb}{0.000000,0.000000,0.000000}%
\pgfsetstrokecolor{currentstroke}%
\pgfsetdash{}{0pt}%
\pgfpathmoveto{\pgfqpoint{1.140747in}{0.873720in}}%
\pgfpathlineto{\pgfqpoint{1.157933in}{0.878288in}}%
\pgfpathlineto{\pgfqpoint{1.175183in}{0.883426in}}%
\pgfpathlineto{\pgfqpoint{1.183543in}{0.882112in}}%
\pgfpathlineto{\pgfqpoint{1.191871in}{0.881125in}}%
\pgfpathlineto{\pgfqpoint{1.174659in}{0.875716in}}%
\pgfpathlineto{\pgfqpoint{1.157510in}{0.870880in}}%
\pgfpathlineto{\pgfqpoint{1.149145in}{0.872134in}}%
\pgfpathlineto{\pgfqpoint{1.140747in}{0.873720in}}%
\pgfpathclose%
\pgfusepath{fill}%
\end{pgfscope}%
\begin{pgfscope}%
\pgfpathrectangle{\pgfqpoint{0.407592in}{0.222904in}}{\pgfqpoint{1.411163in}{1.411163in}}%
\pgfusepath{clip}%
\pgfsetbuttcap%
\pgfsetroundjoin%
\definecolor{currentfill}{rgb}{0.909804,0.698039,0.423529}%
\pgfsetfillcolor{currentfill}%
\pgfsetfillopacity{0.120000}%
\pgfsetlinewidth{0.000000pt}%
\definecolor{currentstroke}{rgb}{0.000000,0.000000,0.000000}%
\pgfsetstrokecolor{currentstroke}%
\pgfsetdash{}{0pt}%
\pgfpathmoveto{\pgfqpoint{1.345321in}{0.929517in}}%
\pgfpathlineto{\pgfqpoint{1.362930in}{0.938689in}}%
\pgfpathlineto{\pgfqpoint{1.380616in}{0.948187in}}%
\pgfpathlineto{\pgfqpoint{1.388621in}{0.950502in}}%
\pgfpathlineto{\pgfqpoint{1.396603in}{0.953164in}}%
\pgfpathlineto{\pgfqpoint{1.378945in}{0.943529in}}%
\pgfpathlineto{\pgfqpoint{1.361364in}{0.934214in}}%
\pgfpathlineto{\pgfqpoint{1.353354in}{0.931683in}}%
\pgfpathlineto{\pgfqpoint{1.345321in}{0.929517in}}%
\pgfpathclose%
\pgfusepath{fill}%
\end{pgfscope}%
\begin{pgfscope}%
\pgfpathrectangle{\pgfqpoint{0.407592in}{0.222904in}}{\pgfqpoint{1.411163in}{1.411163in}}%
\pgfusepath{clip}%
\pgfsetbuttcap%
\pgfsetroundjoin%
\definecolor{currentfill}{rgb}{0.909804,0.698039,0.423529}%
\pgfsetfillcolor{currentfill}%
\pgfsetfillopacity{0.120000}%
\pgfsetlinewidth{0.000000pt}%
\definecolor{currentstroke}{rgb}{0.000000,0.000000,0.000000}%
\pgfsetstrokecolor{currentstroke}%
\pgfsetdash{}{0pt}%
\pgfpathmoveto{\pgfqpoint{1.566662in}{1.051529in}}%
\pgfpathlineto{\pgfqpoint{1.584771in}{1.063005in}}%
\pgfpathlineto{\pgfqpoint{1.602961in}{1.074670in}}%
\pgfpathlineto{\pgfqpoint{1.610662in}{1.080360in}}%
\pgfpathlineto{\pgfqpoint{1.618346in}{1.086338in}}%
\pgfpathlineto{\pgfqpoint{1.600181in}{1.074725in}}%
\pgfpathlineto{\pgfqpoint{1.582098in}{1.063303in}}%
\pgfpathlineto{\pgfqpoint{1.574388in}{1.057265in}}%
\pgfpathlineto{\pgfqpoint{1.566662in}{1.051529in}}%
\pgfpathclose%
\pgfusepath{fill}%
\end{pgfscope}%
\begin{pgfscope}%
\pgfpathrectangle{\pgfqpoint{0.407592in}{0.222904in}}{\pgfqpoint{1.411163in}{1.411163in}}%
\pgfusepath{clip}%
\pgfsetbuttcap%
\pgfsetroundjoin%
\definecolor{currentfill}{rgb}{0.909804,0.698039,0.423529}%
\pgfsetfillcolor{currentfill}%
\pgfsetfillopacity{0.120000}%
\pgfsetlinewidth{0.000000pt}%
\definecolor{currentstroke}{rgb}{0.000000,0.000000,0.000000}%
\pgfsetstrokecolor{currentstroke}%
\pgfsetdash{}{0pt}%
\pgfpathmoveto{\pgfqpoint{1.089572in}{0.871057in}}%
\pgfpathlineto{\pgfqpoint{1.106680in}{0.874132in}}%
\pgfpathlineto{\pgfqpoint{1.123847in}{0.877831in}}%
\pgfpathlineto{\pgfqpoint{1.132314in}{0.875624in}}%
\pgfpathlineto{\pgfqpoint{1.140747in}{0.873720in}}%
\pgfpathlineto{\pgfqpoint{1.123621in}{0.869754in}}%
\pgfpathlineto{\pgfqpoint{1.106554in}{0.866422in}}%
\pgfpathlineto{\pgfqpoint{1.098081in}{0.868587in}}%
\pgfpathlineto{\pgfqpoint{1.089572in}{0.871057in}}%
\pgfpathclose%
\pgfusepath{fill}%
\end{pgfscope}%
\begin{pgfscope}%
\pgfpathrectangle{\pgfqpoint{0.407592in}{0.222904in}}{\pgfqpoint{1.411163in}{1.411163in}}%
\pgfusepath{clip}%
\pgfsetbuttcap%
\pgfsetroundjoin%
\definecolor{currentfill}{rgb}{0.909804,0.698039,0.423529}%
\pgfsetfillcolor{currentfill}%
\pgfsetfillopacity{0.120000}%
\pgfsetlinewidth{0.000000pt}%
\definecolor{currentstroke}{rgb}{0.000000,0.000000,0.000000}%
\pgfsetstrokecolor{currentstroke}%
\pgfsetdash{}{0pt}%
\pgfpathmoveto{\pgfqpoint{0.901557in}{0.896327in}}%
\pgfpathlineto{\pgfqpoint{0.918430in}{0.893312in}}%
\pgfpathlineto{\pgfqpoint{0.935341in}{0.891006in}}%
\pgfpathlineto{\pgfqpoint{0.944221in}{0.886469in}}%
\pgfpathlineto{\pgfqpoint{0.953060in}{0.882146in}}%
\pgfpathlineto{\pgfqpoint{0.936198in}{0.884363in}}%
\pgfpathlineto{\pgfqpoint{0.919375in}{0.887319in}}%
\pgfpathlineto{\pgfqpoint{0.910487in}{0.891716in}}%
\pgfpathlineto{\pgfqpoint{0.901557in}{0.896327in}}%
\pgfpathclose%
\pgfusepath{fill}%
\end{pgfscope}%
\begin{pgfscope}%
\pgfpathrectangle{\pgfqpoint{0.407592in}{0.222904in}}{\pgfqpoint{1.411163in}{1.411163in}}%
\pgfusepath{clip}%
\pgfsetbuttcap%
\pgfsetroundjoin%
\definecolor{currentfill}{rgb}{0.909804,0.698039,0.423529}%
\pgfsetfillcolor{currentfill}%
\pgfsetfillopacity{0.120000}%
\pgfsetlinewidth{0.000000pt}%
\definecolor{currentstroke}{rgb}{0.000000,0.000000,0.000000}%
\pgfsetstrokecolor{currentstroke}%
\pgfsetdash{}{0pt}%
\pgfpathmoveto{\pgfqpoint{1.396603in}{0.953164in}}%
\pgfpathlineto{\pgfqpoint{1.414337in}{0.963098in}}%
\pgfpathlineto{\pgfqpoint{1.432150in}{0.973311in}}%
\pgfpathlineto{\pgfqpoint{1.440082in}{0.976427in}}%
\pgfpathlineto{\pgfqpoint{1.447992in}{0.979879in}}%
\pgfpathlineto{\pgfqpoint{1.430206in}{0.969575in}}%
\pgfpathlineto{\pgfqpoint{1.412498in}{0.959546in}}%
\pgfpathlineto{\pgfqpoint{1.404561in}{0.956177in}}%
\pgfpathlineto{\pgfqpoint{1.396603in}{0.953164in}}%
\pgfpathclose%
\pgfusepath{fill}%
\end{pgfscope}%
\begin{pgfscope}%
\pgfpathrectangle{\pgfqpoint{0.407592in}{0.222904in}}{\pgfqpoint{1.411163in}{1.411163in}}%
\pgfusepath{clip}%
\pgfsetbuttcap%
\pgfsetroundjoin%
\definecolor{currentfill}{rgb}{0.909804,0.698039,0.423529}%
\pgfsetfillcolor{currentfill}%
\pgfsetfillopacity{0.120000}%
\pgfsetlinewidth{0.000000pt}%
\definecolor{currentstroke}{rgb}{0.000000,0.000000,0.000000}%
\pgfsetstrokecolor{currentstroke}%
\pgfsetdash{}{0pt}%
\pgfpathmoveto{\pgfqpoint{1.038306in}{0.873158in}}%
\pgfpathlineto{\pgfqpoint{1.055349in}{0.874673in}}%
\pgfpathlineto{\pgfqpoint{1.072445in}{0.876853in}}%
\pgfpathlineto{\pgfqpoint{1.081027in}{0.873817in}}%
\pgfpathlineto{\pgfqpoint{1.089572in}{0.871057in}}%
\pgfpathlineto{\pgfqpoint{1.072519in}{0.868636in}}%
\pgfpathlineto{\pgfqpoint{1.055521in}{0.866897in}}%
\pgfpathlineto{\pgfqpoint{1.046932in}{0.869890in}}%
\pgfpathlineto{\pgfqpoint{1.038306in}{0.873158in}}%
\pgfpathclose%
\pgfusepath{fill}%
\end{pgfscope}%
\begin{pgfscope}%
\pgfpathrectangle{\pgfqpoint{0.407592in}{0.222904in}}{\pgfqpoint{1.411163in}{1.411163in}}%
\pgfusepath{clip}%
\pgfsetbuttcap%
\pgfsetroundjoin%
\definecolor{currentfill}{rgb}{0.909804,0.698039,0.423529}%
\pgfsetfillcolor{currentfill}%
\pgfsetfillopacity{0.120000}%
\pgfsetlinewidth{0.000000pt}%
\definecolor{currentstroke}{rgb}{0.000000,0.000000,0.000000}%
\pgfsetstrokecolor{currentstroke}%
\pgfsetdash{}{0pt}%
\pgfpathmoveto{\pgfqpoint{1.618346in}{1.086338in}}%
\pgfpathlineto{\pgfqpoint{1.636594in}{1.098132in}}%
\pgfpathlineto{\pgfqpoint{1.654924in}{1.110099in}}%
\pgfpathlineto{\pgfqpoint{1.662566in}{1.116290in}}%
\pgfpathlineto{\pgfqpoint{1.670193in}{1.122749in}}%
\pgfpathlineto{\pgfqpoint{1.651889in}{1.110858in}}%
\pgfpathlineto{\pgfqpoint{1.633667in}{1.099142in}}%
\pgfpathlineto{\pgfqpoint{1.626015in}{1.092599in}}%
\pgfpathlineto{\pgfqpoint{1.618346in}{1.086338in}}%
\pgfpathclose%
\pgfusepath{fill}%
\end{pgfscope}%
\begin{pgfscope}%
\pgfpathrectangle{\pgfqpoint{0.407592in}{0.222904in}}{\pgfqpoint{1.411163in}{1.411163in}}%
\pgfusepath{clip}%
\pgfsetbuttcap%
\pgfsetroundjoin%
\definecolor{currentfill}{rgb}{0.909804,0.698039,0.423529}%
\pgfsetfillcolor{currentfill}%
\pgfsetfillopacity{0.120000}%
\pgfsetlinewidth{0.000000pt}%
\definecolor{currentstroke}{rgb}{0.000000,0.000000,0.000000}%
\pgfsetstrokecolor{currentstroke}%
\pgfsetdash{}{0pt}%
\pgfpathmoveto{\pgfqpoint{1.447992in}{0.979879in}}%
\pgfpathlineto{\pgfqpoint{1.465856in}{0.990441in}}%
\pgfpathlineto{\pgfqpoint{1.483800in}{1.001245in}}%
\pgfpathlineto{\pgfqpoint{1.491663in}{1.005095in}}%
\pgfpathlineto{\pgfqpoint{1.499505in}{1.009267in}}%
\pgfpathlineto{\pgfqpoint{1.481588in}{0.998415in}}%
\pgfpathlineto{\pgfqpoint{1.463749in}{0.987802in}}%
\pgfpathlineto{\pgfqpoint{1.455881in}{0.983670in}}%
\pgfpathlineto{\pgfqpoint{1.447992in}{0.979879in}}%
\pgfpathclose%
\pgfusepath{fill}%
\end{pgfscope}%
\begin{pgfscope}%
\pgfpathrectangle{\pgfqpoint{0.407592in}{0.222904in}}{\pgfqpoint{1.411163in}{1.411163in}}%
\pgfusepath{clip}%
\pgfsetbuttcap%
\pgfsetroundjoin%
\definecolor{currentfill}{rgb}{0.909804,0.698039,0.423529}%
\pgfsetfillcolor{currentfill}%
\pgfsetfillopacity{0.120000}%
\pgfsetlinewidth{0.000000pt}%
\definecolor{currentstroke}{rgb}{0.000000,0.000000,0.000000}%
\pgfsetstrokecolor{currentstroke}%
\pgfsetdash{}{0pt}%
\pgfpathmoveto{\pgfqpoint{1.670193in}{1.122749in}}%
\pgfpathlineto{\pgfqpoint{1.688581in}{1.134807in}}%
\pgfpathlineto{\pgfqpoint{1.696180in}{1.141486in}}%
\pgfpathlineto{\pgfqpoint{1.703764in}{1.148422in}}%
\pgfpathlineto{\pgfqpoint{1.685403in}{1.136461in}}%
\pgfpathlineto{\pgfqpoint{1.677805in}{1.129474in}}%
\pgfusepath{fill}%
\end{pgfscope}%
\begin{pgfscope}%
\pgfpathrectangle{\pgfqpoint{0.407592in}{0.222904in}}{\pgfqpoint{1.411163in}{1.411163in}}%
\pgfusepath{clip}%
\pgfsetbuttcap%
\pgfsetroundjoin%
\definecolor{currentfill}{rgb}{0.909804,0.698039,0.423529}%
\pgfsetfillcolor{currentfill}%
\pgfsetfillopacity{0.120000}%
\pgfsetlinewidth{0.000000pt}%
\definecolor{currentstroke}{rgb}{0.000000,0.000000,0.000000}%
\pgfsetstrokecolor{currentstroke}%
\pgfsetdash{}{0pt}%
\pgfpathmoveto{\pgfqpoint{0.986908in}{0.879895in}}%
\pgfpathlineto{\pgfqpoint{1.003899in}{0.879834in}}%
\pgfpathlineto{\pgfqpoint{1.020938in}{0.880462in}}%
\pgfpathlineto{\pgfqpoint{1.029641in}{0.876687in}}%
\pgfpathlineto{\pgfqpoint{1.038306in}{0.873158in}}%
\pgfpathlineto{\pgfqpoint{1.021313in}{0.872333in}}%
\pgfpathlineto{\pgfqpoint{1.004369in}{0.872218in}}%
\pgfpathlineto{\pgfqpoint{0.995658in}{0.875934in}}%
\pgfpathlineto{\pgfqpoint{0.986908in}{0.879895in}}%
\pgfpathclose%
\pgfusepath{fill}%
\end{pgfscope}%
\begin{pgfscope}%
\pgfpathrectangle{\pgfqpoint{0.407592in}{0.222904in}}{\pgfqpoint{1.411163in}{1.411163in}}%
\pgfusepath{clip}%
\pgfsetbuttcap%
\pgfsetroundjoin%
\definecolor{currentfill}{rgb}{0.909804,0.698039,0.423529}%
\pgfsetfillcolor{currentfill}%
\pgfsetfillopacity{0.120000}%
\pgfsetlinewidth{0.000000pt}%
\definecolor{currentstroke}{rgb}{0.000000,0.000000,0.000000}%
\pgfsetstrokecolor{currentstroke}%
\pgfsetdash{}{0pt}%
\pgfpathmoveto{\pgfqpoint{0.849831in}{0.914337in}}%
\pgfpathlineto{\pgfqpoint{0.866685in}{0.909898in}}%
\pgfpathlineto{\pgfqpoint{0.883573in}{0.906143in}}%
\pgfpathlineto{\pgfqpoint{0.892586in}{0.901139in}}%
\pgfpathlineto{\pgfqpoint{0.901557in}{0.896327in}}%
\pgfpathlineto{\pgfqpoint{0.884720in}{0.900055in}}%
\pgfpathlineto{\pgfqpoint{0.867917in}{0.904496in}}%
\pgfpathlineto{\pgfqpoint{0.858895in}{0.909321in}}%
\pgfpathlineto{\pgfqpoint{0.849831in}{0.914337in}}%
\pgfpathclose%
\pgfusepath{fill}%
\end{pgfscope}%
\begin{pgfscope}%
\pgfpathrectangle{\pgfqpoint{0.407592in}{0.222904in}}{\pgfqpoint{1.411163in}{1.411163in}}%
\pgfusepath{clip}%
\pgfsetbuttcap%
\pgfsetroundjoin%
\definecolor{currentfill}{rgb}{0.909804,0.698039,0.423529}%
\pgfsetfillcolor{currentfill}%
\pgfsetfillopacity{0.120000}%
\pgfsetlinewidth{0.000000pt}%
\definecolor{currentstroke}{rgb}{0.000000,0.000000,0.000000}%
\pgfsetstrokecolor{currentstroke}%
\pgfsetdash{}{0pt}%
\pgfpathmoveto{\pgfqpoint{1.499505in}{1.009267in}}%
\pgfpathlineto{\pgfqpoint{1.517503in}{1.020344in}}%
\pgfpathlineto{\pgfqpoint{1.535582in}{1.031634in}}%
\pgfpathlineto{\pgfqpoint{1.543380in}{1.036148in}}%
\pgfpathlineto{\pgfqpoint{1.551158in}{1.040969in}}%
\pgfpathlineto{\pgfqpoint{1.533105in}{1.029670in}}%
\pgfpathlineto{\pgfqpoint{1.515133in}{1.018582in}}%
\pgfpathlineto{\pgfqpoint{1.507329in}{1.013763in}}%
\pgfpathlineto{\pgfqpoint{1.499505in}{1.009267in}}%
\pgfpathclose%
\pgfusepath{fill}%
\end{pgfscope}%
\begin{pgfscope}%
\pgfpathrectangle{\pgfqpoint{0.407592in}{0.222904in}}{\pgfqpoint{1.411163in}{1.411163in}}%
\pgfusepath{clip}%
\pgfsetbuttcap%
\pgfsetroundjoin%
\definecolor{currentfill}{rgb}{0.909804,0.698039,0.423529}%
\pgfsetfillcolor{currentfill}%
\pgfsetfillopacity{0.120000}%
\pgfsetlinewidth{0.000000pt}%
\definecolor{currentstroke}{rgb}{0.000000,0.000000,0.000000}%
\pgfsetstrokecolor{currentstroke}%
\pgfsetdash{}{0pt}%
\pgfpathmoveto{\pgfqpoint{1.226493in}{0.893529in}}%
\pgfpathlineto{\pgfqpoint{1.243907in}{0.900461in}}%
\pgfpathlineto{\pgfqpoint{1.261390in}{0.907840in}}%
\pgfpathlineto{\pgfqpoint{1.269617in}{0.907710in}}%
\pgfpathlineto{\pgfqpoint{1.277814in}{0.907913in}}%
\pgfpathlineto{\pgfqpoint{1.260364in}{0.900287in}}%
\pgfpathlineto{\pgfqpoint{1.242984in}{0.893102in}}%
\pgfpathlineto{\pgfqpoint{1.234754in}{0.893144in}}%
\pgfpathlineto{\pgfqpoint{1.226493in}{0.893529in}}%
\pgfpathclose%
\pgfusepath{fill}%
\end{pgfscope}%
\begin{pgfscope}%
\pgfpathrectangle{\pgfqpoint{0.407592in}{0.222904in}}{\pgfqpoint{1.411163in}{1.411163in}}%
\pgfusepath{clip}%
\pgfsetbuttcap%
\pgfsetroundjoin%
\definecolor{currentfill}{rgb}{0.909804,0.698039,0.423529}%
\pgfsetfillcolor{currentfill}%
\pgfsetfillopacity{0.120000}%
\pgfsetlinewidth{0.000000pt}%
\definecolor{currentstroke}{rgb}{0.000000,0.000000,0.000000}%
\pgfsetstrokecolor{currentstroke}%
\pgfsetdash{}{0pt}%
\pgfpathmoveto{\pgfqpoint{0.935341in}{0.891006in}}%
\pgfpathlineto{\pgfqpoint{0.952293in}{0.889404in}}%
\pgfpathlineto{\pgfqpoint{0.969288in}{0.888495in}}%
\pgfpathlineto{\pgfqpoint{0.978118in}{0.884085in}}%
\pgfpathlineto{\pgfqpoint{0.986908in}{0.879895in}}%
\pgfpathlineto{\pgfqpoint{0.969962in}{0.880661in}}%
\pgfpathlineto{\pgfqpoint{0.953060in}{0.882146in}}%
\pgfpathlineto{\pgfqpoint{0.944221in}{0.886469in}}%
\pgfpathlineto{\pgfqpoint{0.935341in}{0.891006in}}%
\pgfpathclose%
\pgfusepath{fill}%
\end{pgfscope}%
\begin{pgfscope}%
\pgfpathrectangle{\pgfqpoint{0.407592in}{0.222904in}}{\pgfqpoint{1.411163in}{1.411163in}}%
\pgfusepath{clip}%
\pgfsetbuttcap%
\pgfsetroundjoin%
\definecolor{currentfill}{rgb}{0.909804,0.698039,0.423529}%
\pgfsetfillcolor{currentfill}%
\pgfsetfillopacity{0.120000}%
\pgfsetlinewidth{0.000000pt}%
\definecolor{currentstroke}{rgb}{0.000000,0.000000,0.000000}%
\pgfsetstrokecolor{currentstroke}%
\pgfsetdash{}{0pt}%
\pgfpathmoveto{\pgfqpoint{1.175183in}{0.883426in}}%
\pgfpathlineto{\pgfqpoint{1.192498in}{0.889102in}}%
\pgfpathlineto{\pgfqpoint{1.209879in}{0.895285in}}%
\pgfpathlineto{\pgfqpoint{1.218202in}{0.894247in}}%
\pgfpathlineto{\pgfqpoint{1.226493in}{0.893529in}}%
\pgfpathlineto{\pgfqpoint{1.209148in}{0.887074in}}%
\pgfpathlineto{\pgfqpoint{1.191871in}{0.881125in}}%
\pgfpathlineto{\pgfqpoint{1.183543in}{0.882112in}}%
\pgfpathlineto{\pgfqpoint{1.175183in}{0.883426in}}%
\pgfpathclose%
\pgfusepath{fill}%
\end{pgfscope}%
\begin{pgfscope}%
\pgfpathrectangle{\pgfqpoint{0.407592in}{0.222904in}}{\pgfqpoint{1.411163in}{1.411163in}}%
\pgfusepath{clip}%
\pgfsetbuttcap%
\pgfsetroundjoin%
\definecolor{currentfill}{rgb}{0.909804,0.698039,0.423529}%
\pgfsetfillcolor{currentfill}%
\pgfsetfillopacity{0.120000}%
\pgfsetlinewidth{0.000000pt}%
\definecolor{currentstroke}{rgb}{0.000000,0.000000,0.000000}%
\pgfsetstrokecolor{currentstroke}%
\pgfsetdash{}{0pt}%
\pgfpathmoveto{\pgfqpoint{1.277814in}{0.907913in}}%
\pgfpathlineto{\pgfqpoint{1.295336in}{0.915954in}}%
\pgfpathlineto{\pgfqpoint{1.312932in}{0.924381in}}%
\pgfpathlineto{\pgfqpoint{1.321069in}{0.925150in}}%
\pgfpathlineto{\pgfqpoint{1.329179in}{0.926257in}}%
\pgfpathlineto{\pgfqpoint{1.311614in}{0.917618in}}%
\pgfpathlineto{\pgfqpoint{1.294124in}{0.909358in}}%
\pgfpathlineto{\pgfqpoint{1.285983in}{0.908460in}}%
\pgfpathlineto{\pgfqpoint{1.277814in}{0.907913in}}%
\pgfpathclose%
\pgfusepath{fill}%
\end{pgfscope}%
\begin{pgfscope}%
\pgfpathrectangle{\pgfqpoint{0.407592in}{0.222904in}}{\pgfqpoint{1.411163in}{1.411163in}}%
\pgfusepath{clip}%
\pgfsetbuttcap%
\pgfsetroundjoin%
\definecolor{currentfill}{rgb}{0.909804,0.698039,0.423529}%
\pgfsetfillcolor{currentfill}%
\pgfsetfillopacity{0.120000}%
\pgfsetlinewidth{0.000000pt}%
\definecolor{currentstroke}{rgb}{0.000000,0.000000,0.000000}%
\pgfsetstrokecolor{currentstroke}%
\pgfsetdash{}{0pt}%
\pgfpathmoveto{\pgfqpoint{1.551158in}{1.040969in}}%
\pgfpathlineto{\pgfqpoint{1.569293in}{1.052467in}}%
\pgfpathlineto{\pgfqpoint{1.587509in}{1.064156in}}%
\pgfpathlineto{\pgfqpoint{1.595244in}{1.069268in}}%
\pgfpathlineto{\pgfqpoint{1.602961in}{1.074670in}}%
\pgfpathlineto{\pgfqpoint{1.584771in}{1.063005in}}%
\pgfpathlineto{\pgfqpoint{1.566662in}{1.051529in}}%
\pgfpathlineto{\pgfqpoint{1.558919in}{1.046096in}}%
\pgfpathlineto{\pgfqpoint{1.551158in}{1.040969in}}%
\pgfpathclose%
\pgfusepath{fill}%
\end{pgfscope}%
\begin{pgfscope}%
\pgfpathrectangle{\pgfqpoint{0.407592in}{0.222904in}}{\pgfqpoint{1.411163in}{1.411163in}}%
\pgfusepath{clip}%
\pgfsetbuttcap%
\pgfsetroundjoin%
\definecolor{currentfill}{rgb}{0.909804,0.698039,0.423529}%
\pgfsetfillcolor{currentfill}%
\pgfsetfillopacity{0.120000}%
\pgfsetlinewidth{0.000000pt}%
\definecolor{currentstroke}{rgb}{0.000000,0.000000,0.000000}%
\pgfsetstrokecolor{currentstroke}%
\pgfsetdash{}{0pt}%
\pgfpathmoveto{\pgfqpoint{1.329179in}{0.926257in}}%
\pgfpathlineto{\pgfqpoint{1.346818in}{0.935253in}}%
\pgfpathlineto{\pgfqpoint{1.364533in}{0.944582in}}%
\pgfpathlineto{\pgfqpoint{1.372587in}{0.946216in}}%
\pgfpathlineto{\pgfqpoint{1.380616in}{0.948187in}}%
\pgfpathlineto{\pgfqpoint{1.362930in}{0.938689in}}%
\pgfpathlineto{\pgfqpoint{1.345321in}{0.929517in}}%
\pgfpathlineto{\pgfqpoint{1.337263in}{0.927711in}}%
\pgfpathlineto{\pgfqpoint{1.329179in}{0.926257in}}%
\pgfpathclose%
\pgfusepath{fill}%
\end{pgfscope}%
\begin{pgfscope}%
\pgfpathrectangle{\pgfqpoint{0.407592in}{0.222904in}}{\pgfqpoint{1.411163in}{1.411163in}}%
\pgfusepath{clip}%
\pgfsetbuttcap%
\pgfsetroundjoin%
\definecolor{currentfill}{rgb}{0.909804,0.698039,0.423529}%
\pgfsetfillcolor{currentfill}%
\pgfsetfillopacity{0.120000}%
\pgfsetlinewidth{0.000000pt}%
\definecolor{currentstroke}{rgb}{0.000000,0.000000,0.000000}%
\pgfsetstrokecolor{currentstroke}%
\pgfsetdash{}{0pt}%
\pgfpathmoveto{\pgfqpoint{1.123847in}{0.877831in}}%
\pgfpathlineto{\pgfqpoint{1.141074in}{0.882125in}}%
\pgfpathlineto{\pgfqpoint{1.158363in}{0.886984in}}%
\pgfpathlineto{\pgfqpoint{1.166790in}{0.885054in}}%
\pgfpathlineto{\pgfqpoint{1.175183in}{0.883426in}}%
\pgfpathlineto{\pgfqpoint{1.157933in}{0.878288in}}%
\pgfpathlineto{\pgfqpoint{1.140747in}{0.873720in}}%
\pgfpathlineto{\pgfqpoint{1.132314in}{0.875624in}}%
\pgfpathlineto{\pgfqpoint{1.123847in}{0.877831in}}%
\pgfpathclose%
\pgfusepath{fill}%
\end{pgfscope}%
\begin{pgfscope}%
\pgfpathrectangle{\pgfqpoint{0.407592in}{0.222904in}}{\pgfqpoint{1.411163in}{1.411163in}}%
\pgfusepath{clip}%
\pgfsetbuttcap%
\pgfsetroundjoin%
\definecolor{currentfill}{rgb}{0.909804,0.698039,0.423529}%
\pgfsetfillcolor{currentfill}%
\pgfsetfillopacity{0.120000}%
\pgfsetlinewidth{0.000000pt}%
\definecolor{currentstroke}{rgb}{0.000000,0.000000,0.000000}%
\pgfsetstrokecolor{currentstroke}%
\pgfsetdash{}{0pt}%
\pgfpathmoveto{\pgfqpoint{1.380616in}{0.948187in}}%
\pgfpathlineto{\pgfqpoint{1.398379in}{0.957991in}}%
\pgfpathlineto{\pgfqpoint{1.416219in}{0.968080in}}%
\pgfpathlineto{\pgfqpoint{1.424196in}{0.970530in}}%
\pgfpathlineto{\pgfqpoint{1.432150in}{0.973311in}}%
\pgfpathlineto{\pgfqpoint{1.414337in}{0.963098in}}%
\pgfpathlineto{\pgfqpoint{1.396603in}{0.953164in}}%
\pgfpathlineto{\pgfqpoint{1.388621in}{0.950502in}}%
\pgfpathlineto{\pgfqpoint{1.380616in}{0.948187in}}%
\pgfpathclose%
\pgfusepath{fill}%
\end{pgfscope}%
\begin{pgfscope}%
\pgfpathrectangle{\pgfqpoint{0.407592in}{0.222904in}}{\pgfqpoint{1.411163in}{1.411163in}}%
\pgfusepath{clip}%
\pgfsetbuttcap%
\pgfsetroundjoin%
\definecolor{currentfill}{rgb}{0.909804,0.698039,0.423529}%
\pgfsetfillcolor{currentfill}%
\pgfsetfillopacity{0.120000}%
\pgfsetlinewidth{0.000000pt}%
\definecolor{currentstroke}{rgb}{0.000000,0.000000,0.000000}%
\pgfsetstrokecolor{currentstroke}%
\pgfsetdash{}{0pt}%
\pgfpathmoveto{\pgfqpoint{1.072445in}{0.876853in}}%
\pgfpathlineto{\pgfqpoint{1.089597in}{0.879673in}}%
\pgfpathlineto{\pgfqpoint{1.106806in}{0.883105in}}%
\pgfpathlineto{\pgfqpoint{1.115345in}{0.880329in}}%
\pgfpathlineto{\pgfqpoint{1.123847in}{0.877831in}}%
\pgfpathlineto{\pgfqpoint{1.106680in}{0.874132in}}%
\pgfpathlineto{\pgfqpoint{1.089572in}{0.871057in}}%
\pgfpathlineto{\pgfqpoint{1.081027in}{0.873817in}}%
\pgfpathlineto{\pgfqpoint{1.072445in}{0.876853in}}%
\pgfpathclose%
\pgfusepath{fill}%
\end{pgfscope}%
\begin{pgfscope}%
\pgfpathrectangle{\pgfqpoint{0.407592in}{0.222904in}}{\pgfqpoint{1.411163in}{1.411163in}}%
\pgfusepath{clip}%
\pgfsetbuttcap%
\pgfsetroundjoin%
\definecolor{currentfill}{rgb}{0.909804,0.698039,0.423529}%
\pgfsetfillcolor{currentfill}%
\pgfsetfillopacity{0.120000}%
\pgfsetlinewidth{0.000000pt}%
\definecolor{currentstroke}{rgb}{0.000000,0.000000,0.000000}%
\pgfsetstrokecolor{currentstroke}%
\pgfsetdash{}{0pt}%
\pgfpathmoveto{\pgfqpoint{1.602961in}{1.074670in}}%
\pgfpathlineto{\pgfqpoint{1.621235in}{1.086515in}}%
\pgfpathlineto{\pgfqpoint{1.639591in}{1.098532in}}%
\pgfpathlineto{\pgfqpoint{1.647265in}{1.104179in}}%
\pgfpathlineto{\pgfqpoint{1.654924in}{1.110099in}}%
\pgfpathlineto{\pgfqpoint{1.636594in}{1.098132in}}%
\pgfpathlineto{\pgfqpoint{1.618346in}{1.086338in}}%
\pgfpathlineto{\pgfqpoint{1.610662in}{1.080360in}}%
\pgfpathlineto{\pgfqpoint{1.602961in}{1.074670in}}%
\pgfpathclose%
\pgfusepath{fill}%
\end{pgfscope}%
\begin{pgfscope}%
\pgfpathrectangle{\pgfqpoint{0.407592in}{0.222904in}}{\pgfqpoint{1.411163in}{1.411163in}}%
\pgfusepath{clip}%
\pgfsetbuttcap%
\pgfsetroundjoin%
\definecolor{currentfill}{rgb}{0.909804,0.698039,0.423529}%
\pgfsetfillcolor{currentfill}%
\pgfsetfillopacity{0.120000}%
\pgfsetlinewidth{0.000000pt}%
\definecolor{currentstroke}{rgb}{0.000000,0.000000,0.000000}%
\pgfsetstrokecolor{currentstroke}%
\pgfsetdash{}{0pt}%
\pgfpathmoveto{\pgfqpoint{1.654924in}{1.110099in}}%
\pgfpathlineto{\pgfqpoint{1.673337in}{1.122231in}}%
\pgfpathlineto{\pgfqpoint{1.680967in}{1.128388in}}%
\pgfpathlineto{\pgfqpoint{1.688581in}{1.134807in}}%
\pgfpathlineto{\pgfqpoint{1.670193in}{1.122749in}}%
\pgfpathlineto{\pgfqpoint{1.662566in}{1.116290in}}%
\pgfusepath{fill}%
\end{pgfscope}%
\begin{pgfscope}%
\pgfpathrectangle{\pgfqpoint{0.407592in}{0.222904in}}{\pgfqpoint{1.411163in}{1.411163in}}%
\pgfusepath{clip}%
\pgfsetbuttcap%
\pgfsetroundjoin%
\definecolor{currentfill}{rgb}{0.909804,0.698039,0.423529}%
\pgfsetfillcolor{currentfill}%
\pgfsetfillopacity{0.120000}%
\pgfsetlinewidth{0.000000pt}%
\definecolor{currentstroke}{rgb}{0.000000,0.000000,0.000000}%
\pgfsetstrokecolor{currentstroke}%
\pgfsetdash{}{0pt}%
\pgfpathmoveto{\pgfqpoint{0.883573in}{0.906143in}}%
\pgfpathlineto{\pgfqpoint{0.900497in}{0.903072in}}%
\pgfpathlineto{\pgfqpoint{0.917459in}{0.900684in}}%
\pgfpathlineto{\pgfqpoint{0.926421in}{0.895748in}}%
\pgfpathlineto{\pgfqpoint{0.935341in}{0.891006in}}%
\pgfpathlineto{\pgfqpoint{0.918430in}{0.893312in}}%
\pgfpathlineto{\pgfqpoint{0.901557in}{0.896327in}}%
\pgfpathlineto{\pgfqpoint{0.892586in}{0.901139in}}%
\pgfpathlineto{\pgfqpoint{0.883573in}{0.906143in}}%
\pgfpathclose%
\pgfusepath{fill}%
\end{pgfscope}%
\begin{pgfscope}%
\pgfpathrectangle{\pgfqpoint{0.407592in}{0.222904in}}{\pgfqpoint{1.411163in}{1.411163in}}%
\pgfusepath{clip}%
\pgfsetbuttcap%
\pgfsetroundjoin%
\definecolor{currentfill}{rgb}{0.909804,0.698039,0.423529}%
\pgfsetfillcolor{currentfill}%
\pgfsetfillopacity{0.120000}%
\pgfsetlinewidth{0.000000pt}%
\definecolor{currentstroke}{rgb}{0.000000,0.000000,0.000000}%
\pgfsetstrokecolor{currentstroke}%
\pgfsetdash{}{0pt}%
\pgfpathmoveto{\pgfqpoint{1.432150in}{0.973311in}}%
\pgfpathlineto{\pgfqpoint{1.450041in}{0.983788in}}%
\pgfpathlineto{\pgfqpoint{1.468012in}{0.994511in}}%
\pgfpathlineto{\pgfqpoint{1.475917in}{0.997718in}}%
\pgfpathlineto{\pgfqpoint{1.483800in}{1.001245in}}%
\pgfpathlineto{\pgfqpoint{1.465856in}{0.990441in}}%
\pgfpathlineto{\pgfqpoint{1.447992in}{0.979879in}}%
\pgfpathlineto{\pgfqpoint{1.440082in}{0.976427in}}%
\pgfpathlineto{\pgfqpoint{1.432150in}{0.973311in}}%
\pgfpathclose%
\pgfusepath{fill}%
\end{pgfscope}%
\begin{pgfscope}%
\pgfpathrectangle{\pgfqpoint{0.407592in}{0.222904in}}{\pgfqpoint{1.411163in}{1.411163in}}%
\pgfusepath{clip}%
\pgfsetbuttcap%
\pgfsetroundjoin%
\definecolor{currentfill}{rgb}{0.909804,0.698039,0.423529}%
\pgfsetfillcolor{currentfill}%
\pgfsetfillopacity{0.120000}%
\pgfsetlinewidth{0.000000pt}%
\definecolor{currentstroke}{rgb}{0.000000,0.000000,0.000000}%
\pgfsetstrokecolor{currentstroke}%
\pgfsetdash{}{0pt}%
\pgfpathmoveto{\pgfqpoint{1.020938in}{0.880462in}}%
\pgfpathlineto{\pgfqpoint{1.038028in}{0.881758in}}%
\pgfpathlineto{\pgfqpoint{1.055170in}{0.883702in}}%
\pgfpathlineto{\pgfqpoint{1.063826in}{0.880152in}}%
\pgfpathlineto{\pgfqpoint{1.072445in}{0.876853in}}%
\pgfpathlineto{\pgfqpoint{1.055349in}{0.874673in}}%
\pgfpathlineto{\pgfqpoint{1.038306in}{0.873158in}}%
\pgfpathlineto{\pgfqpoint{1.029641in}{0.876687in}}%
\pgfpathlineto{\pgfqpoint{1.020938in}{0.880462in}}%
\pgfpathclose%
\pgfusepath{fill}%
\end{pgfscope}%
\begin{pgfscope}%
\pgfpathrectangle{\pgfqpoint{0.407592in}{0.222904in}}{\pgfqpoint{1.411163in}{1.411163in}}%
\pgfusepath{clip}%
\pgfsetbuttcap%
\pgfsetroundjoin%
\definecolor{currentfill}{rgb}{0.909804,0.698039,0.423529}%
\pgfsetfillcolor{currentfill}%
\pgfsetfillopacity{0.120000}%
\pgfsetlinewidth{0.000000pt}%
\definecolor{currentstroke}{rgb}{0.000000,0.000000,0.000000}%
\pgfsetstrokecolor{currentstroke}%
\pgfsetdash{}{0pt}%
\pgfpathmoveto{\pgfqpoint{1.483800in}{1.001245in}}%
\pgfpathlineto{\pgfqpoint{1.501825in}{1.012278in}}%
\pgfpathlineto{\pgfqpoint{1.519930in}{1.023527in}}%
\pgfpathlineto{\pgfqpoint{1.527766in}{1.027427in}}%
\pgfpathlineto{\pgfqpoint{1.535582in}{1.031634in}}%
\pgfpathlineto{\pgfqpoint{1.517503in}{1.020344in}}%
\pgfpathlineto{\pgfqpoint{1.499505in}{1.009267in}}%
\pgfpathlineto{\pgfqpoint{1.491663in}{1.005095in}}%
\pgfpathlineto{\pgfqpoint{1.483800in}{1.001245in}}%
\pgfpathclose%
\pgfusepath{fill}%
\end{pgfscope}%
\begin{pgfscope}%
\pgfpathrectangle{\pgfqpoint{0.407592in}{0.222904in}}{\pgfqpoint{1.411163in}{1.411163in}}%
\pgfusepath{clip}%
\pgfsetbuttcap%
\pgfsetroundjoin%
\definecolor{currentfill}{rgb}{0.909804,0.698039,0.423529}%
\pgfsetfillcolor{currentfill}%
\pgfsetfillopacity{0.120000}%
\pgfsetlinewidth{0.000000pt}%
\definecolor{currentstroke}{rgb}{0.000000,0.000000,0.000000}%
\pgfsetstrokecolor{currentstroke}%
\pgfsetdash{}{0pt}%
\pgfpathmoveto{\pgfqpoint{0.969288in}{0.888495in}}%
\pgfpathlineto{\pgfqpoint{0.986328in}{0.888266in}}%
\pgfpathlineto{\pgfqpoint{1.003415in}{0.888704in}}%
\pgfpathlineto{\pgfqpoint{1.012197in}{0.884471in}}%
\pgfpathlineto{\pgfqpoint{1.020938in}{0.880462in}}%
\pgfpathlineto{\pgfqpoint{1.003899in}{0.879834in}}%
\pgfpathlineto{\pgfqpoint{0.986908in}{0.879895in}}%
\pgfpathlineto{\pgfqpoint{0.978118in}{0.884085in}}%
\pgfpathlineto{\pgfqpoint{0.969288in}{0.888495in}}%
\pgfpathclose%
\pgfusepath{fill}%
\end{pgfscope}%
\begin{pgfscope}%
\pgfpathrectangle{\pgfqpoint{0.407592in}{0.222904in}}{\pgfqpoint{1.411163in}{1.411163in}}%
\pgfusepath{clip}%
\pgfsetbuttcap%
\pgfsetroundjoin%
\definecolor{currentfill}{rgb}{0.909804,0.698039,0.423529}%
\pgfsetfillcolor{currentfill}%
\pgfsetfillopacity{0.120000}%
\pgfsetlinewidth{0.000000pt}%
\definecolor{currentstroke}{rgb}{0.000000,0.000000,0.000000}%
\pgfsetstrokecolor{currentstroke}%
\pgfsetdash{}{0pt}%
\pgfpathmoveto{\pgfqpoint{1.535582in}{1.031634in}}%
\pgfpathlineto{\pgfqpoint{1.553743in}{1.043125in}}%
\pgfpathlineto{\pgfqpoint{1.571986in}{1.054807in}}%
\pgfpathlineto{\pgfqpoint{1.579757in}{1.059335in}}%
\pgfpathlineto{\pgfqpoint{1.587509in}{1.064156in}}%
\pgfpathlineto{\pgfqpoint{1.569293in}{1.052467in}}%
\pgfpathlineto{\pgfqpoint{1.551158in}{1.040969in}}%
\pgfpathlineto{\pgfqpoint{1.543380in}{1.036148in}}%
\pgfpathlineto{\pgfqpoint{1.535582in}{1.031634in}}%
\pgfpathclose%
\pgfusepath{fill}%
\end{pgfscope}%
\begin{pgfscope}%
\pgfpathrectangle{\pgfqpoint{0.407592in}{0.222904in}}{\pgfqpoint{1.411163in}{1.411163in}}%
\pgfusepath{clip}%
\pgfsetbuttcap%
\pgfsetroundjoin%
\definecolor{currentfill}{rgb}{0.909804,0.698039,0.423529}%
\pgfsetfillcolor{currentfill}%
\pgfsetfillopacity{0.120000}%
\pgfsetlinewidth{0.000000pt}%
\definecolor{currentstroke}{rgb}{0.000000,0.000000,0.000000}%
\pgfsetstrokecolor{currentstroke}%
\pgfsetdash{}{0pt}%
\pgfpathmoveto{\pgfqpoint{1.209879in}{0.895285in}}%
\pgfpathlineto{\pgfqpoint{1.227328in}{0.901948in}}%
\pgfpathlineto{\pgfqpoint{1.244847in}{0.909061in}}%
\pgfpathlineto{\pgfqpoint{1.253134in}{0.908293in}}%
\pgfpathlineto{\pgfqpoint{1.261390in}{0.907840in}}%
\pgfpathlineto{\pgfqpoint{1.243907in}{0.900461in}}%
\pgfpathlineto{\pgfqpoint{1.226493in}{0.893529in}}%
\pgfpathlineto{\pgfqpoint{1.218202in}{0.894247in}}%
\pgfpathlineto{\pgfqpoint{1.209879in}{0.895285in}}%
\pgfpathclose%
\pgfusepath{fill}%
\end{pgfscope}%
\begin{pgfscope}%
\pgfpathrectangle{\pgfqpoint{0.407592in}{0.222904in}}{\pgfqpoint{1.411163in}{1.411163in}}%
\pgfusepath{clip}%
\pgfsetbuttcap%
\pgfsetroundjoin%
\definecolor{currentfill}{rgb}{0.909804,0.698039,0.423529}%
\pgfsetfillcolor{currentfill}%
\pgfsetfillopacity{0.120000}%
\pgfsetlinewidth{0.000000pt}%
\definecolor{currentstroke}{rgb}{0.000000,0.000000,0.000000}%
\pgfsetstrokecolor{currentstroke}%
\pgfsetdash{}{0pt}%
\pgfpathmoveto{\pgfqpoint{1.261390in}{0.907840in}}%
\pgfpathlineto{\pgfqpoint{1.278945in}{0.915638in}}%
\pgfpathlineto{\pgfqpoint{1.296574in}{0.923830in}}%
\pgfpathlineto{\pgfqpoint{1.304767in}{0.923944in}}%
\pgfpathlineto{\pgfqpoint{1.312932in}{0.924381in}}%
\pgfpathlineto{\pgfqpoint{1.295336in}{0.915954in}}%
\pgfpathlineto{\pgfqpoint{1.277814in}{0.907913in}}%
\pgfpathlineto{\pgfqpoint{1.269617in}{0.907710in}}%
\pgfpathlineto{\pgfqpoint{1.261390in}{0.907840in}}%
\pgfpathclose%
\pgfusepath{fill}%
\end{pgfscope}%
\begin{pgfscope}%
\pgfpathrectangle{\pgfqpoint{0.407592in}{0.222904in}}{\pgfqpoint{1.411163in}{1.411163in}}%
\pgfusepath{clip}%
\pgfsetbuttcap%
\pgfsetroundjoin%
\definecolor{currentfill}{rgb}{0.909804,0.698039,0.423529}%
\pgfsetfillcolor{currentfill}%
\pgfsetfillopacity{0.120000}%
\pgfsetlinewidth{0.000000pt}%
\definecolor{currentstroke}{rgb}{0.000000,0.000000,0.000000}%
\pgfsetstrokecolor{currentstroke}%
\pgfsetdash{}{0pt}%
\pgfpathmoveto{\pgfqpoint{1.158363in}{0.886984in}}%
\pgfpathlineto{\pgfqpoint{1.175716in}{0.892379in}}%
\pgfpathlineto{\pgfqpoint{1.193136in}{0.898281in}}%
\pgfpathlineto{\pgfqpoint{1.201524in}{0.896634in}}%
\pgfpathlineto{\pgfqpoint{1.209879in}{0.895285in}}%
\pgfpathlineto{\pgfqpoint{1.192498in}{0.889102in}}%
\pgfpathlineto{\pgfqpoint{1.175183in}{0.883426in}}%
\pgfpathlineto{\pgfqpoint{1.166790in}{0.885054in}}%
\pgfpathlineto{\pgfqpoint{1.158363in}{0.886984in}}%
\pgfpathclose%
\pgfusepath{fill}%
\end{pgfscope}%
\begin{pgfscope}%
\pgfpathrectangle{\pgfqpoint{0.407592in}{0.222904in}}{\pgfqpoint{1.411163in}{1.411163in}}%
\pgfusepath{clip}%
\pgfsetbuttcap%
\pgfsetroundjoin%
\definecolor{currentfill}{rgb}{0.909804,0.698039,0.423529}%
\pgfsetfillcolor{currentfill}%
\pgfsetfillopacity{0.120000}%
\pgfsetlinewidth{0.000000pt}%
\definecolor{currentstroke}{rgb}{0.000000,0.000000,0.000000}%
\pgfsetstrokecolor{currentstroke}%
\pgfsetdash{}{0pt}%
\pgfpathmoveto{\pgfqpoint{1.312932in}{0.924381in}}%
\pgfpathlineto{\pgfqpoint{1.330601in}{0.933172in}}%
\pgfpathlineto{\pgfqpoint{1.348347in}{0.942304in}}%
\pgfpathlineto{\pgfqpoint{1.356453in}{0.943280in}}%
\pgfpathlineto{\pgfqpoint{1.364533in}{0.944582in}}%
\pgfpathlineto{\pgfqpoint{1.346818in}{0.935253in}}%
\pgfpathlineto{\pgfqpoint{1.329179in}{0.926257in}}%
\pgfpathlineto{\pgfqpoint{1.321069in}{0.925150in}}%
\pgfpathlineto{\pgfqpoint{1.312932in}{0.924381in}}%
\pgfpathclose%
\pgfusepath{fill}%
\end{pgfscope}%
\begin{pgfscope}%
\pgfpathrectangle{\pgfqpoint{0.407592in}{0.222904in}}{\pgfqpoint{1.411163in}{1.411163in}}%
\pgfusepath{clip}%
\pgfsetbuttcap%
\pgfsetroundjoin%
\definecolor{currentfill}{rgb}{0.909804,0.698039,0.423529}%
\pgfsetfillcolor{currentfill}%
\pgfsetfillopacity{0.120000}%
\pgfsetlinewidth{0.000000pt}%
\definecolor{currentstroke}{rgb}{0.000000,0.000000,0.000000}%
\pgfsetstrokecolor{currentstroke}%
\pgfsetdash{}{0pt}%
\pgfpathmoveto{\pgfqpoint{0.917459in}{0.900684in}}%
\pgfpathlineto{\pgfqpoint{0.934461in}{0.898971in}}%
\pgfpathlineto{\pgfqpoint{0.951506in}{0.897926in}}%
\pgfpathlineto{\pgfqpoint{0.960417in}{0.893112in}}%
\pgfpathlineto{\pgfqpoint{0.969288in}{0.888495in}}%
\pgfpathlineto{\pgfqpoint{0.952293in}{0.889404in}}%
\pgfpathlineto{\pgfqpoint{0.935341in}{0.891006in}}%
\pgfpathlineto{\pgfqpoint{0.926421in}{0.895748in}}%
\pgfpathlineto{\pgfqpoint{0.917459in}{0.900684in}}%
\pgfpathclose%
\pgfusepath{fill}%
\end{pgfscope}%
\begin{pgfscope}%
\pgfpathrectangle{\pgfqpoint{0.407592in}{0.222904in}}{\pgfqpoint{1.411163in}{1.411163in}}%
\pgfusepath{clip}%
\pgfsetbuttcap%
\pgfsetroundjoin%
\definecolor{currentfill}{rgb}{0.909804,0.698039,0.423529}%
\pgfsetfillcolor{currentfill}%
\pgfsetfillopacity{0.120000}%
\pgfsetlinewidth{0.000000pt}%
\definecolor{currentstroke}{rgb}{0.000000,0.000000,0.000000}%
\pgfsetstrokecolor{currentstroke}%
\pgfsetdash{}{0pt}%
\pgfpathmoveto{\pgfqpoint{1.587509in}{1.064156in}}%
\pgfpathlineto{\pgfqpoint{1.605808in}{1.076025in}}%
\pgfpathlineto{\pgfqpoint{1.624191in}{1.088066in}}%
\pgfpathlineto{\pgfqpoint{1.631900in}{1.093161in}}%
\pgfpathlineto{\pgfqpoint{1.639591in}{1.098532in}}%
\pgfpathlineto{\pgfqpoint{1.621235in}{1.086515in}}%
\pgfpathlineto{\pgfqpoint{1.602961in}{1.074670in}}%
\pgfpathlineto{\pgfqpoint{1.595244in}{1.069268in}}%
\pgfpathlineto{\pgfqpoint{1.587509in}{1.064156in}}%
\pgfpathclose%
\pgfusepath{fill}%
\end{pgfscope}%
\begin{pgfscope}%
\pgfpathrectangle{\pgfqpoint{0.407592in}{0.222904in}}{\pgfqpoint{1.411163in}{1.411163in}}%
\pgfusepath{clip}%
\pgfsetbuttcap%
\pgfsetroundjoin%
\definecolor{currentfill}{rgb}{0.909804,0.698039,0.423529}%
\pgfsetfillcolor{currentfill}%
\pgfsetfillopacity{0.120000}%
\pgfsetlinewidth{0.000000pt}%
\definecolor{currentstroke}{rgb}{0.000000,0.000000,0.000000}%
\pgfsetstrokecolor{currentstroke}%
\pgfsetdash{}{0pt}%
\pgfpathmoveto{\pgfqpoint{1.639591in}{1.098532in}}%
\pgfpathlineto{\pgfqpoint{1.658031in}{1.110714in}}%
\pgfpathlineto{\pgfqpoint{1.665692in}{1.116339in}}%
\pgfpathlineto{\pgfqpoint{1.673337in}{1.122231in}}%
\pgfpathlineto{\pgfqpoint{1.654924in}{1.110099in}}%
\pgfpathlineto{\pgfqpoint{1.647265in}{1.104179in}}%
\pgfusepath{fill}%
\end{pgfscope}%
\begin{pgfscope}%
\pgfpathrectangle{\pgfqpoint{0.407592in}{0.222904in}}{\pgfqpoint{1.411163in}{1.411163in}}%
\pgfusepath{clip}%
\pgfsetbuttcap%
\pgfsetroundjoin%
\definecolor{currentfill}{rgb}{0.909804,0.698039,0.423529}%
\pgfsetfillcolor{currentfill}%
\pgfsetfillopacity{0.120000}%
\pgfsetlinewidth{0.000000pt}%
\definecolor{currentstroke}{rgb}{0.000000,0.000000,0.000000}%
\pgfsetstrokecolor{currentstroke}%
\pgfsetdash{}{0pt}%
\pgfpathmoveto{\pgfqpoint{1.364533in}{0.944582in}}%
\pgfpathlineto{\pgfqpoint{1.382324in}{0.954223in}}%
\pgfpathlineto{\pgfqpoint{1.400194in}{0.964157in}}%
\pgfpathlineto{\pgfqpoint{1.408219in}{0.965957in}}%
\pgfpathlineto{\pgfqpoint{1.416219in}{0.968080in}}%
\pgfpathlineto{\pgfqpoint{1.398379in}{0.957991in}}%
\pgfpathlineto{\pgfqpoint{1.380616in}{0.948187in}}%
\pgfpathlineto{\pgfqpoint{1.372587in}{0.946216in}}%
\pgfpathlineto{\pgfqpoint{1.364533in}{0.944582in}}%
\pgfpathclose%
\pgfusepath{fill}%
\end{pgfscope}%
\begin{pgfscope}%
\pgfpathrectangle{\pgfqpoint{0.407592in}{0.222904in}}{\pgfqpoint{1.411163in}{1.411163in}}%
\pgfusepath{clip}%
\pgfsetbuttcap%
\pgfsetroundjoin%
\definecolor{currentfill}{rgb}{0.909804,0.698039,0.423529}%
\pgfsetfillcolor{currentfill}%
\pgfsetfillopacity{0.120000}%
\pgfsetlinewidth{0.000000pt}%
\definecolor{currentstroke}{rgb}{0.000000,0.000000,0.000000}%
\pgfsetstrokecolor{currentstroke}%
\pgfsetdash{}{0pt}%
\pgfpathmoveto{\pgfqpoint{1.106806in}{0.883105in}}%
\pgfpathlineto{\pgfqpoint{1.124075in}{0.887124in}}%
\pgfpathlineto{\pgfqpoint{1.141406in}{0.891701in}}%
\pgfpathlineto{\pgfqpoint{1.149902in}{0.889204in}}%
\pgfpathlineto{\pgfqpoint{1.158363in}{0.886984in}}%
\pgfpathlineto{\pgfqpoint{1.141074in}{0.882125in}}%
\pgfpathlineto{\pgfqpoint{1.123847in}{0.877831in}}%
\pgfpathlineto{\pgfqpoint{1.115345in}{0.880329in}}%
\pgfpathlineto{\pgfqpoint{1.106806in}{0.883105in}}%
\pgfpathclose%
\pgfusepath{fill}%
\end{pgfscope}%
\begin{pgfscope}%
\pgfpathrectangle{\pgfqpoint{0.407592in}{0.222904in}}{\pgfqpoint{1.411163in}{1.411163in}}%
\pgfusepath{clip}%
\pgfsetbuttcap%
\pgfsetroundjoin%
\definecolor{currentfill}{rgb}{0.909804,0.698039,0.423529}%
\pgfsetfillcolor{currentfill}%
\pgfsetfillopacity{0.120000}%
\pgfsetlinewidth{0.000000pt}%
\definecolor{currentstroke}{rgb}{0.000000,0.000000,0.000000}%
\pgfsetstrokecolor{currentstroke}%
\pgfsetdash{}{0pt}%
\pgfpathmoveto{\pgfqpoint{1.416219in}{0.968080in}}%
\pgfpathlineto{\pgfqpoint{1.434139in}{0.978439in}}%
\pgfpathlineto{\pgfqpoint{1.452138in}{0.989050in}}%
\pgfpathlineto{\pgfqpoint{1.460086in}{0.991623in}}%
\pgfpathlineto{\pgfqpoint{1.468012in}{0.994511in}}%
\pgfpathlineto{\pgfqpoint{1.450041in}{0.983788in}}%
\pgfpathlineto{\pgfqpoint{1.432150in}{0.973311in}}%
\pgfpathlineto{\pgfqpoint{1.424196in}{0.970530in}}%
\pgfpathlineto{\pgfqpoint{1.416219in}{0.968080in}}%
\pgfpathclose%
\pgfusepath{fill}%
\end{pgfscope}%
\begin{pgfscope}%
\pgfpathrectangle{\pgfqpoint{0.407592in}{0.222904in}}{\pgfqpoint{1.411163in}{1.411163in}}%
\pgfusepath{clip}%
\pgfsetbuttcap%
\pgfsetroundjoin%
\definecolor{currentfill}{rgb}{0.909804,0.698039,0.423529}%
\pgfsetfillcolor{currentfill}%
\pgfsetfillopacity{0.120000}%
\pgfsetlinewidth{0.000000pt}%
\definecolor{currentstroke}{rgb}{0.000000,0.000000,0.000000}%
\pgfsetstrokecolor{currentstroke}%
\pgfsetdash{}{0pt}%
\pgfpathmoveto{\pgfqpoint{1.055170in}{0.883702in}}%
\pgfpathlineto{\pgfqpoint{1.072367in}{0.886271in}}%
\pgfpathlineto{\pgfqpoint{1.089620in}{0.889440in}}%
\pgfpathlineto{\pgfqpoint{1.098232in}{0.886146in}}%
\pgfpathlineto{\pgfqpoint{1.106806in}{0.883105in}}%
\pgfpathlineto{\pgfqpoint{1.089597in}{0.879673in}}%
\pgfpathlineto{\pgfqpoint{1.072445in}{0.876853in}}%
\pgfpathlineto{\pgfqpoint{1.063826in}{0.880152in}}%
\pgfpathlineto{\pgfqpoint{1.055170in}{0.883702in}}%
\pgfpathclose%
\pgfusepath{fill}%
\end{pgfscope}%
\begin{pgfscope}%
\pgfpathrectangle{\pgfqpoint{0.407592in}{0.222904in}}{\pgfqpoint{1.411163in}{1.411163in}}%
\pgfusepath{clip}%
\pgfsetbuttcap%
\pgfsetroundjoin%
\definecolor{currentfill}{rgb}{0.909804,0.698039,0.423529}%
\pgfsetfillcolor{currentfill}%
\pgfsetfillopacity{0.120000}%
\pgfsetlinewidth{0.000000pt}%
\definecolor{currentstroke}{rgb}{0.000000,0.000000,0.000000}%
\pgfsetstrokecolor{currentstroke}%
\pgfsetdash{}{0pt}%
\pgfpathmoveto{\pgfqpoint{1.468012in}{0.994511in}}%
\pgfpathlineto{\pgfqpoint{1.486064in}{1.005468in}}%
\pgfpathlineto{\pgfqpoint{1.504197in}{1.016645in}}%
\pgfpathlineto{\pgfqpoint{1.512074in}{1.019934in}}%
\pgfpathlineto{\pgfqpoint{1.519930in}{1.023527in}}%
\pgfpathlineto{\pgfqpoint{1.501825in}{1.012278in}}%
\pgfpathlineto{\pgfqpoint{1.483800in}{1.001245in}}%
\pgfpathlineto{\pgfqpoint{1.475917in}{0.997718in}}%
\pgfpathlineto{\pgfqpoint{1.468012in}{0.994511in}}%
\pgfpathclose%
\pgfusepath{fill}%
\end{pgfscope}%
\begin{pgfscope}%
\pgfpathrectangle{\pgfqpoint{0.407592in}{0.222904in}}{\pgfqpoint{1.411163in}{1.411163in}}%
\pgfusepath{clip}%
\pgfsetbuttcap%
\pgfsetroundjoin%
\definecolor{currentfill}{rgb}{0.909804,0.698039,0.423529}%
\pgfsetfillcolor{currentfill}%
\pgfsetfillopacity{0.120000}%
\pgfsetlinewidth{0.000000pt}%
\definecolor{currentstroke}{rgb}{0.000000,0.000000,0.000000}%
\pgfsetstrokecolor{currentstroke}%
\pgfsetdash{}{0pt}%
\pgfpathmoveto{\pgfqpoint{1.003415in}{0.888704in}}%
\pgfpathlineto{\pgfqpoint{1.020552in}{0.889789in}}%
\pgfpathlineto{\pgfqpoint{1.037741in}{0.891505in}}%
\pgfpathlineto{\pgfqpoint{1.046475in}{0.887490in}}%
\pgfpathlineto{\pgfqpoint{1.055170in}{0.883702in}}%
\pgfpathlineto{\pgfqpoint{1.038028in}{0.881758in}}%
\pgfpathlineto{\pgfqpoint{1.020938in}{0.880462in}}%
\pgfpathlineto{\pgfqpoint{1.012197in}{0.884471in}}%
\pgfpathlineto{\pgfqpoint{1.003415in}{0.888704in}}%
\pgfpathclose%
\pgfusepath{fill}%
\end{pgfscope}%
\begin{pgfscope}%
\pgfpathrectangle{\pgfqpoint{0.407592in}{0.222904in}}{\pgfqpoint{1.411163in}{1.411163in}}%
\pgfusepath{clip}%
\pgfsetbuttcap%
\pgfsetroundjoin%
\definecolor{currentfill}{rgb}{0.909804,0.698039,0.423529}%
\pgfsetfillcolor{currentfill}%
\pgfsetfillopacity{0.120000}%
\pgfsetlinewidth{0.000000pt}%
\definecolor{currentstroke}{rgb}{0.000000,0.000000,0.000000}%
\pgfsetstrokecolor{currentstroke}%
\pgfsetdash{}{0pt}%
\pgfpathmoveto{\pgfqpoint{1.519930in}{1.023527in}}%
\pgfpathlineto{\pgfqpoint{1.538117in}{1.034980in}}%
\pgfpathlineto{\pgfqpoint{1.556387in}{1.046628in}}%
\pgfpathlineto{\pgfqpoint{1.564196in}{1.050571in}}%
\pgfpathlineto{\pgfqpoint{1.571986in}{1.054807in}}%
\pgfpathlineto{\pgfqpoint{1.553743in}{1.043125in}}%
\pgfpathlineto{\pgfqpoint{1.535582in}{1.031634in}}%
\pgfpathlineto{\pgfqpoint{1.527766in}{1.027427in}}%
\pgfpathlineto{\pgfqpoint{1.519930in}{1.023527in}}%
\pgfpathclose%
\pgfusepath{fill}%
\end{pgfscope}%
\begin{pgfscope}%
\pgfpathrectangle{\pgfqpoint{0.407592in}{0.222904in}}{\pgfqpoint{1.411163in}{1.411163in}}%
\pgfusepath{clip}%
\pgfsetbuttcap%
\pgfsetroundjoin%
\definecolor{currentfill}{rgb}{0.909804,0.698039,0.423529}%
\pgfsetfillcolor{currentfill}%
\pgfsetfillopacity{0.120000}%
\pgfsetlinewidth{0.000000pt}%
\definecolor{currentstroke}{rgb}{0.000000,0.000000,0.000000}%
\pgfsetstrokecolor{currentstroke}%
\pgfsetdash{}{0pt}%
\pgfpathmoveto{\pgfqpoint{0.951506in}{0.897926in}}%
\pgfpathlineto{\pgfqpoint{0.968596in}{0.897538in}}%
\pgfpathlineto{\pgfqpoint{0.985732in}{0.897793in}}%
\pgfpathlineto{\pgfqpoint{0.994594in}{0.893148in}}%
\pgfpathlineto{\pgfqpoint{1.003415in}{0.888704in}}%
\pgfpathlineto{\pgfqpoint{0.986328in}{0.888266in}}%
\pgfpathlineto{\pgfqpoint{0.969288in}{0.888495in}}%
\pgfpathlineto{\pgfqpoint{0.960417in}{0.893112in}}%
\pgfpathlineto{\pgfqpoint{0.951506in}{0.897926in}}%
\pgfpathclose%
\pgfusepath{fill}%
\end{pgfscope}%
\begin{pgfscope}%
\pgfpathrectangle{\pgfqpoint{0.407592in}{0.222904in}}{\pgfqpoint{1.411163in}{1.411163in}}%
\pgfusepath{clip}%
\pgfsetbuttcap%
\pgfsetroundjoin%
\definecolor{currentfill}{rgb}{0.909804,0.698039,0.423529}%
\pgfsetfillcolor{currentfill}%
\pgfsetfillopacity{0.120000}%
\pgfsetlinewidth{0.000000pt}%
\definecolor{currentstroke}{rgb}{0.000000,0.000000,0.000000}%
\pgfsetstrokecolor{currentstroke}%
\pgfsetdash{}{0pt}%
\pgfpathmoveto{\pgfqpoint{1.624191in}{1.088066in}}%
\pgfpathlineto{\pgfqpoint{1.642658in}{1.100272in}}%
\pgfpathlineto{\pgfqpoint{1.650353in}{1.105358in}}%
\pgfpathlineto{\pgfqpoint{1.658031in}{1.110714in}}%
\pgfpathlineto{\pgfqpoint{1.639591in}{1.098532in}}%
\pgfpathlineto{\pgfqpoint{1.631900in}{1.093161in}}%
\pgfusepath{fill}%
\end{pgfscope}%
\begin{pgfscope}%
\pgfpathrectangle{\pgfqpoint{0.407592in}{0.222904in}}{\pgfqpoint{1.411163in}{1.411163in}}%
\pgfusepath{clip}%
\pgfsetbuttcap%
\pgfsetroundjoin%
\definecolor{currentfill}{rgb}{0.909804,0.698039,0.423529}%
\pgfsetfillcolor{currentfill}%
\pgfsetfillopacity{0.120000}%
\pgfsetlinewidth{0.000000pt}%
\definecolor{currentstroke}{rgb}{0.000000,0.000000,0.000000}%
\pgfsetstrokecolor{currentstroke}%
\pgfsetdash{}{0pt}%
\pgfpathmoveto{\pgfqpoint{1.571986in}{1.054807in}}%
\pgfpathlineto{\pgfqpoint{1.590312in}{1.066672in}}%
\pgfpathlineto{\pgfqpoint{1.608721in}{1.078710in}}%
\pgfpathlineto{\pgfqpoint{1.616465in}{1.083249in}}%
\pgfpathlineto{\pgfqpoint{1.624191in}{1.088066in}}%
\pgfpathlineto{\pgfqpoint{1.605808in}{1.076025in}}%
\pgfpathlineto{\pgfqpoint{1.587509in}{1.064156in}}%
\pgfpathlineto{\pgfqpoint{1.579757in}{1.059335in}}%
\pgfpathlineto{\pgfqpoint{1.571986in}{1.054807in}}%
\pgfpathclose%
\pgfusepath{fill}%
\end{pgfscope}%
\begin{pgfscope}%
\pgfpathrectangle{\pgfqpoint{0.407592in}{0.222904in}}{\pgfqpoint{1.411163in}{1.411163in}}%
\pgfusepath{clip}%
\pgfsetbuttcap%
\pgfsetroundjoin%
\definecolor{currentfill}{rgb}{0.909804,0.698039,0.423529}%
\pgfsetfillcolor{currentfill}%
\pgfsetfillopacity{0.120000}%
\pgfsetlinewidth{0.000000pt}%
\definecolor{currentstroke}{rgb}{0.000000,0.000000,0.000000}%
\pgfsetstrokecolor{currentstroke}%
\pgfsetdash{}{0pt}%
\pgfpathmoveto{\pgfqpoint{1.244847in}{0.909061in}}%
\pgfpathlineto{\pgfqpoint{1.262437in}{0.916599in}}%
\pgfpathlineto{\pgfqpoint{1.280099in}{0.924536in}}%
\pgfpathlineto{\pgfqpoint{1.288351in}{0.924030in}}%
\pgfpathlineto{\pgfqpoint{1.296574in}{0.923830in}}%
\pgfpathlineto{\pgfqpoint{1.278945in}{0.915638in}}%
\pgfpathlineto{\pgfqpoint{1.261390in}{0.907840in}}%
\pgfpathlineto{\pgfqpoint{1.253134in}{0.908293in}}%
\pgfpathlineto{\pgfqpoint{1.244847in}{0.909061in}}%
\pgfpathclose%
\pgfusepath{fill}%
\end{pgfscope}%
\begin{pgfscope}%
\pgfpathrectangle{\pgfqpoint{0.407592in}{0.222904in}}{\pgfqpoint{1.411163in}{1.411163in}}%
\pgfusepath{clip}%
\pgfsetbuttcap%
\pgfsetroundjoin%
\definecolor{currentfill}{rgb}{0.909804,0.698039,0.423529}%
\pgfsetfillcolor{currentfill}%
\pgfsetfillopacity{0.120000}%
\pgfsetlinewidth{0.000000pt}%
\definecolor{currentstroke}{rgb}{0.000000,0.000000,0.000000}%
\pgfsetstrokecolor{currentstroke}%
\pgfsetdash{}{0pt}%
\pgfpathmoveto{\pgfqpoint{1.193136in}{0.898281in}}%
\pgfpathlineto{\pgfqpoint{1.210622in}{0.904663in}}%
\pgfpathlineto{\pgfqpoint{1.228178in}{0.911499in}}%
\pgfpathlineto{\pgfqpoint{1.236528in}{0.910133in}}%
\pgfpathlineto{\pgfqpoint{1.244847in}{0.909061in}}%
\pgfpathlineto{\pgfqpoint{1.227328in}{0.901948in}}%
\pgfpathlineto{\pgfqpoint{1.209879in}{0.895285in}}%
\pgfpathlineto{\pgfqpoint{1.201524in}{0.896634in}}%
\pgfpathlineto{\pgfqpoint{1.193136in}{0.898281in}}%
\pgfpathclose%
\pgfusepath{fill}%
\end{pgfscope}%
\begin{pgfscope}%
\pgfpathrectangle{\pgfqpoint{0.407592in}{0.222904in}}{\pgfqpoint{1.411163in}{1.411163in}}%
\pgfusepath{clip}%
\pgfsetbuttcap%
\pgfsetroundjoin%
\definecolor{currentfill}{rgb}{0.909804,0.698039,0.423529}%
\pgfsetfillcolor{currentfill}%
\pgfsetfillopacity{0.120000}%
\pgfsetlinewidth{0.000000pt}%
\definecolor{currentstroke}{rgb}{0.000000,0.000000,0.000000}%
\pgfsetstrokecolor{currentstroke}%
\pgfsetdash{}{0pt}%
\pgfpathmoveto{\pgfqpoint{1.296574in}{0.923830in}}%
\pgfpathlineto{\pgfqpoint{1.314276in}{0.932391in}}%
\pgfpathlineto{\pgfqpoint{1.332053in}{0.941301in}}%
\pgfpathlineto{\pgfqpoint{1.340214in}{0.941646in}}%
\pgfpathlineto{\pgfqpoint{1.348347in}{0.942304in}}%
\pgfpathlineto{\pgfqpoint{1.330601in}{0.933172in}}%
\pgfpathlineto{\pgfqpoint{1.312932in}{0.924381in}}%
\pgfpathlineto{\pgfqpoint{1.304767in}{0.923944in}}%
\pgfpathlineto{\pgfqpoint{1.296574in}{0.923830in}}%
\pgfpathclose%
\pgfusepath{fill}%
\end{pgfscope}%
\begin{pgfscope}%
\pgfpathrectangle{\pgfqpoint{0.407592in}{0.222904in}}{\pgfqpoint{1.411163in}{1.411163in}}%
\pgfusepath{clip}%
\pgfsetbuttcap%
\pgfsetroundjoin%
\definecolor{currentfill}{rgb}{0.909804,0.698039,0.423529}%
\pgfsetfillcolor{currentfill}%
\pgfsetfillopacity{0.120000}%
\pgfsetlinewidth{0.000000pt}%
\definecolor{currentstroke}{rgb}{0.000000,0.000000,0.000000}%
\pgfsetstrokecolor{currentstroke}%
\pgfsetdash{}{0pt}%
\pgfpathmoveto{\pgfqpoint{1.141406in}{0.891701in}}%
\pgfpathlineto{\pgfqpoint{1.158799in}{0.896811in}}%
\pgfpathlineto{\pgfqpoint{1.176259in}{0.902427in}}%
\pgfpathlineto{\pgfqpoint{1.184714in}{0.900215in}}%
\pgfpathlineto{\pgfqpoint{1.193136in}{0.898281in}}%
\pgfpathlineto{\pgfqpoint{1.175716in}{0.892379in}}%
\pgfpathlineto{\pgfqpoint{1.158363in}{0.886984in}}%
\pgfpathlineto{\pgfqpoint{1.149902in}{0.889204in}}%
\pgfpathlineto{\pgfqpoint{1.141406in}{0.891701in}}%
\pgfpathclose%
\pgfusepath{fill}%
\end{pgfscope}%
\begin{pgfscope}%
\pgfpathrectangle{\pgfqpoint{0.407592in}{0.222904in}}{\pgfqpoint{1.411163in}{1.411163in}}%
\pgfusepath{clip}%
\pgfsetbuttcap%
\pgfsetroundjoin%
\definecolor{currentfill}{rgb}{0.909804,0.698039,0.423529}%
\pgfsetfillcolor{currentfill}%
\pgfsetfillopacity{0.120000}%
\pgfsetlinewidth{0.000000pt}%
\definecolor{currentstroke}{rgb}{0.000000,0.000000,0.000000}%
\pgfsetstrokecolor{currentstroke}%
\pgfsetdash{}{0pt}%
\pgfpathmoveto{\pgfqpoint{1.348347in}{0.942304in}}%
\pgfpathlineto{\pgfqpoint{1.366169in}{0.951755in}}%
\pgfpathlineto{\pgfqpoint{1.384069in}{0.961507in}}%
\pgfpathlineto{\pgfqpoint{1.392144in}{0.962676in}}%
\pgfpathlineto{\pgfqpoint{1.400194in}{0.964157in}}%
\pgfpathlineto{\pgfqpoint{1.382324in}{0.954223in}}%
\pgfpathlineto{\pgfqpoint{1.364533in}{0.944582in}}%
\pgfpathlineto{\pgfqpoint{1.356453in}{0.943280in}}%
\pgfpathlineto{\pgfqpoint{1.348347in}{0.942304in}}%
\pgfpathclose%
\pgfusepath{fill}%
\end{pgfscope}%
\begin{pgfscope}%
\pgfpathrectangle{\pgfqpoint{0.407592in}{0.222904in}}{\pgfqpoint{1.411163in}{1.411163in}}%
\pgfusepath{clip}%
\pgfsetbuttcap%
\pgfsetroundjoin%
\definecolor{currentfill}{rgb}{0.909804,0.698039,0.423529}%
\pgfsetfillcolor{currentfill}%
\pgfsetfillopacity{0.120000}%
\pgfsetlinewidth{0.000000pt}%
\definecolor{currentstroke}{rgb}{0.000000,0.000000,0.000000}%
\pgfsetstrokecolor{currentstroke}%
\pgfsetdash{}{0pt}%
\pgfpathmoveto{\pgfqpoint{1.400194in}{0.964157in}}%
\pgfpathlineto{\pgfqpoint{1.418142in}{0.974368in}}%
\pgfpathlineto{\pgfqpoint{1.436170in}{0.984838in}}%
\pgfpathlineto{\pgfqpoint{1.444166in}{0.986789in}}%
\pgfpathlineto{\pgfqpoint{1.452138in}{0.989050in}}%
\pgfpathlineto{\pgfqpoint{1.434139in}{0.978439in}}%
\pgfpathlineto{\pgfqpoint{1.416219in}{0.968080in}}%
\pgfpathlineto{\pgfqpoint{1.408219in}{0.965957in}}%
\pgfpathlineto{\pgfqpoint{1.400194in}{0.964157in}}%
\pgfpathclose%
\pgfusepath{fill}%
\end{pgfscope}%
\begin{pgfscope}%
\pgfpathrectangle{\pgfqpoint{0.407592in}{0.222904in}}{\pgfqpoint{1.411163in}{1.411163in}}%
\pgfusepath{clip}%
\pgfsetbuttcap%
\pgfsetroundjoin%
\definecolor{currentfill}{rgb}{0.909804,0.698039,0.423529}%
\pgfsetfillcolor{currentfill}%
\pgfsetfillopacity{0.120000}%
\pgfsetlinewidth{0.000000pt}%
\definecolor{currentstroke}{rgb}{0.000000,0.000000,0.000000}%
\pgfsetstrokecolor{currentstroke}%
\pgfsetdash{}{0pt}%
\pgfpathmoveto{\pgfqpoint{1.089620in}{0.889440in}}%
\pgfpathlineto{\pgfqpoint{1.106933in}{0.893187in}}%
\pgfpathlineto{\pgfqpoint{1.124306in}{0.897485in}}%
\pgfpathlineto{\pgfqpoint{1.132874in}{0.894466in}}%
\pgfpathlineto{\pgfqpoint{1.141406in}{0.891701in}}%
\pgfpathlineto{\pgfqpoint{1.124075in}{0.887124in}}%
\pgfpathlineto{\pgfqpoint{1.106806in}{0.883105in}}%
\pgfpathlineto{\pgfqpoint{1.098232in}{0.886146in}}%
\pgfpathlineto{\pgfqpoint{1.089620in}{0.889440in}}%
\pgfpathclose%
\pgfusepath{fill}%
\end{pgfscope}%
\begin{pgfscope}%
\pgfpathrectangle{\pgfqpoint{0.407592in}{0.222904in}}{\pgfqpoint{1.411163in}{1.411163in}}%
\pgfusepath{clip}%
\pgfsetbuttcap%
\pgfsetroundjoin%
\definecolor{currentfill}{rgb}{0.909804,0.698039,0.423529}%
\pgfsetfillcolor{currentfill}%
\pgfsetfillopacity{0.120000}%
\pgfsetlinewidth{0.000000pt}%
\definecolor{currentstroke}{rgb}{0.000000,0.000000,0.000000}%
\pgfsetstrokecolor{currentstroke}%
\pgfsetdash{}{0pt}%
\pgfpathmoveto{\pgfqpoint{1.452138in}{0.989050in}}%
\pgfpathlineto{\pgfqpoint{1.470217in}{0.999900in}}%
\pgfpathlineto{\pgfqpoint{1.488378in}{1.010977in}}%
\pgfpathlineto{\pgfqpoint{1.496298in}{1.013660in}}%
\pgfpathlineto{\pgfqpoint{1.504197in}{1.016645in}}%
\pgfpathlineto{\pgfqpoint{1.486064in}{1.005468in}}%
\pgfpathlineto{\pgfqpoint{1.468012in}{0.994511in}}%
\pgfpathlineto{\pgfqpoint{1.460086in}{0.991623in}}%
\pgfpathlineto{\pgfqpoint{1.452138in}{0.989050in}}%
\pgfpathclose%
\pgfusepath{fill}%
\end{pgfscope}%
\begin{pgfscope}%
\pgfpathrectangle{\pgfqpoint{0.407592in}{0.222904in}}{\pgfqpoint{1.411163in}{1.411163in}}%
\pgfusepath{clip}%
\pgfsetbuttcap%
\pgfsetroundjoin%
\definecolor{currentfill}{rgb}{0.909804,0.698039,0.423529}%
\pgfsetfillcolor{currentfill}%
\pgfsetfillopacity{0.120000}%
\pgfsetlinewidth{0.000000pt}%
\definecolor{currentstroke}{rgb}{0.000000,0.000000,0.000000}%
\pgfsetstrokecolor{currentstroke}%
\pgfsetdash{}{0pt}%
\pgfpathmoveto{\pgfqpoint{1.037741in}{0.891505in}}%
\pgfpathlineto{\pgfqpoint{1.054984in}{0.893830in}}%
\pgfpathlineto{\pgfqpoint{1.072284in}{0.896744in}}%
\pgfpathlineto{\pgfqpoint{1.080971in}{0.892976in}}%
\pgfpathlineto{\pgfqpoint{1.089620in}{0.889440in}}%
\pgfpathlineto{\pgfqpoint{1.072367in}{0.886271in}}%
\pgfpathlineto{\pgfqpoint{1.055170in}{0.883702in}}%
\pgfpathlineto{\pgfqpoint{1.046475in}{0.887490in}}%
\pgfpathlineto{\pgfqpoint{1.037741in}{0.891505in}}%
\pgfpathclose%
\pgfusepath{fill}%
\end{pgfscope}%
\begin{pgfscope}%
\pgfpathrectangle{\pgfqpoint{0.407592in}{0.222904in}}{\pgfqpoint{1.411163in}{1.411163in}}%
\pgfusepath{clip}%
\pgfsetbuttcap%
\pgfsetroundjoin%
\definecolor{currentfill}{rgb}{0.909804,0.698039,0.423529}%
\pgfsetfillcolor{currentfill}%
\pgfsetfillopacity{0.120000}%
\pgfsetlinewidth{0.000000pt}%
\definecolor{currentstroke}{rgb}{0.000000,0.000000,0.000000}%
\pgfsetstrokecolor{currentstroke}%
\pgfsetdash{}{0pt}%
\pgfpathmoveto{\pgfqpoint{1.504197in}{1.016645in}}%
\pgfpathlineto{\pgfqpoint{1.522412in}{1.028031in}}%
\pgfpathlineto{\pgfqpoint{1.540709in}{1.039616in}}%
\pgfpathlineto{\pgfqpoint{1.548558in}{1.042976in}}%
\pgfpathlineto{\pgfqpoint{1.556387in}{1.046628in}}%
\pgfpathlineto{\pgfqpoint{1.538117in}{1.034980in}}%
\pgfpathlineto{\pgfqpoint{1.519930in}{1.023527in}}%
\pgfpathlineto{\pgfqpoint{1.512074in}{1.019934in}}%
\pgfpathlineto{\pgfqpoint{1.504197in}{1.016645in}}%
\pgfpathclose%
\pgfusepath{fill}%
\end{pgfscope}%
\begin{pgfscope}%
\pgfpathrectangle{\pgfqpoint{0.407592in}{0.222904in}}{\pgfqpoint{1.411163in}{1.411163in}}%
\pgfusepath{clip}%
\pgfsetbuttcap%
\pgfsetroundjoin%
\definecolor{currentfill}{rgb}{0.909804,0.698039,0.423529}%
\pgfsetfillcolor{currentfill}%
\pgfsetfillopacity{0.120000}%
\pgfsetlinewidth{0.000000pt}%
\definecolor{currentstroke}{rgb}{0.000000,0.000000,0.000000}%
\pgfsetstrokecolor{currentstroke}%
\pgfsetdash{}{0pt}%
\pgfpathmoveto{\pgfqpoint{0.985732in}{0.897793in}}%
\pgfpathlineto{\pgfqpoint{1.002918in}{0.898678in}}%
\pgfpathlineto{\pgfqpoint{1.020156in}{0.900174in}}%
\pgfpathlineto{\pgfqpoint{1.028968in}{0.895736in}}%
\pgfpathlineto{\pgfqpoint{1.037741in}{0.891505in}}%
\pgfpathlineto{\pgfqpoint{1.020552in}{0.889789in}}%
\pgfpathlineto{\pgfqpoint{1.003415in}{0.888704in}}%
\pgfpathlineto{\pgfqpoint{0.994594in}{0.893148in}}%
\pgfpathlineto{\pgfqpoint{0.985732in}{0.897793in}}%
\pgfpathclose%
\pgfusepath{fill}%
\end{pgfscope}%
\begin{pgfscope}%
\pgfpathrectangle{\pgfqpoint{0.407592in}{0.222904in}}{\pgfqpoint{1.411163in}{1.411163in}}%
\pgfusepath{clip}%
\pgfsetbuttcap%
\pgfsetroundjoin%
\definecolor{currentfill}{rgb}{0.909804,0.698039,0.423529}%
\pgfsetfillcolor{currentfill}%
\pgfsetfillopacity{0.120000}%
\pgfsetlinewidth{0.000000pt}%
\definecolor{currentstroke}{rgb}{0.000000,0.000000,0.000000}%
\pgfsetstrokecolor{currentstroke}%
\pgfsetdash{}{0pt}%
\pgfpathmoveto{\pgfqpoint{1.608721in}{1.078710in}}%
\pgfpathlineto{\pgfqpoint{1.627215in}{1.090915in}}%
\pgfpathlineto{\pgfqpoint{1.634946in}{1.095457in}}%
\pgfpathlineto{\pgfqpoint{1.642658in}{1.100272in}}%
\pgfpathlineto{\pgfqpoint{1.624191in}{1.088066in}}%
\pgfpathlineto{\pgfqpoint{1.616465in}{1.083249in}}%
\pgfusepath{fill}%
\end{pgfscope}%
\begin{pgfscope}%
\pgfpathrectangle{\pgfqpoint{0.407592in}{0.222904in}}{\pgfqpoint{1.411163in}{1.411163in}}%
\pgfusepath{clip}%
\pgfsetbuttcap%
\pgfsetroundjoin%
\definecolor{currentfill}{rgb}{0.909804,0.698039,0.423529}%
\pgfsetfillcolor{currentfill}%
\pgfsetfillopacity{0.120000}%
\pgfsetlinewidth{0.000000pt}%
\definecolor{currentstroke}{rgb}{0.000000,0.000000,0.000000}%
\pgfsetstrokecolor{currentstroke}%
\pgfsetdash{}{0pt}%
\pgfpathmoveto{\pgfqpoint{1.556387in}{1.046628in}}%
\pgfpathlineto{\pgfqpoint{1.574740in}{1.058460in}}%
\pgfpathlineto{\pgfqpoint{1.593177in}{1.070469in}}%
\pgfpathlineto{\pgfqpoint{1.600959in}{1.074450in}}%
\pgfpathlineto{\pgfqpoint{1.608721in}{1.078710in}}%
\pgfpathlineto{\pgfqpoint{1.590312in}{1.066672in}}%
\pgfpathlineto{\pgfqpoint{1.571986in}{1.054807in}}%
\pgfpathlineto{\pgfqpoint{1.564196in}{1.050571in}}%
\pgfpathlineto{\pgfqpoint{1.556387in}{1.046628in}}%
\pgfpathclose%
\pgfusepath{fill}%
\end{pgfscope}%
\begin{pgfscope}%
\pgfpathrectangle{\pgfqpoint{0.407592in}{0.222904in}}{\pgfqpoint{1.411163in}{1.411163in}}%
\pgfusepath{clip}%
\pgfsetbuttcap%
\pgfsetroundjoin%
\definecolor{currentfill}{rgb}{0.909804,0.698039,0.423529}%
\pgfsetfillcolor{currentfill}%
\pgfsetfillopacity{0.120000}%
\pgfsetlinewidth{0.000000pt}%
\definecolor{currentstroke}{rgb}{0.000000,0.000000,0.000000}%
\pgfsetstrokecolor{currentstroke}%
\pgfsetdash{}{0pt}%
\pgfpathmoveto{\pgfqpoint{1.228178in}{0.911499in}}%
\pgfpathlineto{\pgfqpoint{1.245805in}{0.918763in}}%
\pgfpathlineto{\pgfqpoint{1.263503in}{0.926433in}}%
\pgfpathlineto{\pgfqpoint{1.271817in}{0.925340in}}%
\pgfpathlineto{\pgfqpoint{1.280099in}{0.924536in}}%
\pgfpathlineto{\pgfqpoint{1.262437in}{0.916599in}}%
\pgfpathlineto{\pgfqpoint{1.244847in}{0.909061in}}%
\pgfpathlineto{\pgfqpoint{1.236528in}{0.910133in}}%
\pgfpathlineto{\pgfqpoint{1.228178in}{0.911499in}}%
\pgfpathclose%
\pgfusepath{fill}%
\end{pgfscope}%
\begin{pgfscope}%
\pgfpathrectangle{\pgfqpoint{0.407592in}{0.222904in}}{\pgfqpoint{1.411163in}{1.411163in}}%
\pgfusepath{clip}%
\pgfsetbuttcap%
\pgfsetroundjoin%
\definecolor{currentfill}{rgb}{0.909804,0.698039,0.423529}%
\pgfsetfillcolor{currentfill}%
\pgfsetfillopacity{0.120000}%
\pgfsetlinewidth{0.000000pt}%
\definecolor{currentstroke}{rgb}{0.000000,0.000000,0.000000}%
\pgfsetstrokecolor{currentstroke}%
\pgfsetdash{}{0pt}%
\pgfpathmoveto{\pgfqpoint{1.280099in}{0.924536in}}%
\pgfpathlineto{\pgfqpoint{1.297835in}{0.932850in}}%
\pgfpathlineto{\pgfqpoint{1.315646in}{0.941519in}}%
\pgfpathlineto{\pgfqpoint{1.323864in}{0.941262in}}%
\pgfpathlineto{\pgfqpoint{1.332053in}{0.941301in}}%
\pgfpathlineto{\pgfqpoint{1.314276in}{0.932391in}}%
\pgfpathlineto{\pgfqpoint{1.296574in}{0.923830in}}%
\pgfpathlineto{\pgfqpoint{1.288351in}{0.924030in}}%
\pgfpathlineto{\pgfqpoint{1.280099in}{0.924536in}}%
\pgfpathclose%
\pgfusepath{fill}%
\end{pgfscope}%
\begin{pgfscope}%
\pgfpathrectangle{\pgfqpoint{0.407592in}{0.222904in}}{\pgfqpoint{1.411163in}{1.411163in}}%
\pgfusepath{clip}%
\pgfsetbuttcap%
\pgfsetroundjoin%
\definecolor{currentfill}{rgb}{0.909804,0.698039,0.423529}%
\pgfsetfillcolor{currentfill}%
\pgfsetfillopacity{0.120000}%
\pgfsetlinewidth{0.000000pt}%
\definecolor{currentstroke}{rgb}{0.000000,0.000000,0.000000}%
\pgfsetstrokecolor{currentstroke}%
\pgfsetdash{}{0pt}%
\pgfpathmoveto{\pgfqpoint{1.176259in}{0.902427in}}%
\pgfpathlineto{\pgfqpoint{1.193784in}{0.908523in}}%
\pgfpathlineto{\pgfqpoint{1.211379in}{0.915074in}}%
\pgfpathlineto{\pgfqpoint{1.219795in}{0.913149in}}%
\pgfpathlineto{\pgfqpoint{1.228178in}{0.911499in}}%
\pgfpathlineto{\pgfqpoint{1.210622in}{0.904663in}}%
\pgfpathlineto{\pgfqpoint{1.193136in}{0.898281in}}%
\pgfpathlineto{\pgfqpoint{1.184714in}{0.900215in}}%
\pgfpathlineto{\pgfqpoint{1.176259in}{0.902427in}}%
\pgfpathclose%
\pgfusepath{fill}%
\end{pgfscope}%
\begin{pgfscope}%
\pgfpathrectangle{\pgfqpoint{0.407592in}{0.222904in}}{\pgfqpoint{1.411163in}{1.411163in}}%
\pgfusepath{clip}%
\pgfsetbuttcap%
\pgfsetroundjoin%
\definecolor{currentfill}{rgb}{0.909804,0.698039,0.423529}%
\pgfsetfillcolor{currentfill}%
\pgfsetfillopacity{0.120000}%
\pgfsetlinewidth{0.000000pt}%
\definecolor{currentstroke}{rgb}{0.000000,0.000000,0.000000}%
\pgfsetstrokecolor{currentstroke}%
\pgfsetdash{}{0pt}%
\pgfpathmoveto{\pgfqpoint{1.332053in}{0.941301in}}%
\pgfpathlineto{\pgfqpoint{1.349907in}{0.950539in}}%
\pgfpathlineto{\pgfqpoint{1.367839in}{0.960085in}}%
\pgfpathlineto{\pgfqpoint{1.375967in}{0.960645in}}%
\pgfpathlineto{\pgfqpoint{1.384069in}{0.961507in}}%
\pgfpathlineto{\pgfqpoint{1.366169in}{0.951755in}}%
\pgfpathlineto{\pgfqpoint{1.348347in}{0.942304in}}%
\pgfpathlineto{\pgfqpoint{1.340214in}{0.941646in}}%
\pgfpathlineto{\pgfqpoint{1.332053in}{0.941301in}}%
\pgfpathclose%
\pgfusepath{fill}%
\end{pgfscope}%
\begin{pgfscope}%
\pgfpathrectangle{\pgfqpoint{0.407592in}{0.222904in}}{\pgfqpoint{1.411163in}{1.411163in}}%
\pgfusepath{clip}%
\pgfsetbuttcap%
\pgfsetroundjoin%
\definecolor{currentfill}{rgb}{0.909804,0.698039,0.423529}%
\pgfsetfillcolor{currentfill}%
\pgfsetfillopacity{0.120000}%
\pgfsetlinewidth{0.000000pt}%
\definecolor{currentstroke}{rgb}{0.000000,0.000000,0.000000}%
\pgfsetstrokecolor{currentstroke}%
\pgfsetdash{}{0pt}%
\pgfpathmoveto{\pgfqpoint{1.124306in}{0.897485in}}%
\pgfpathlineto{\pgfqpoint{1.141742in}{0.902311in}}%
\pgfpathlineto{\pgfqpoint{1.159242in}{0.907639in}}%
\pgfpathlineto{\pgfqpoint{1.167768in}{0.904905in}}%
\pgfpathlineto{\pgfqpoint{1.176259in}{0.902427in}}%
\pgfpathlineto{\pgfqpoint{1.158799in}{0.896811in}}%
\pgfpathlineto{\pgfqpoint{1.141406in}{0.891701in}}%
\pgfpathlineto{\pgfqpoint{1.132874in}{0.894466in}}%
\pgfpathlineto{\pgfqpoint{1.124306in}{0.897485in}}%
\pgfpathclose%
\pgfusepath{fill}%
\end{pgfscope}%
\begin{pgfscope}%
\pgfpathrectangle{\pgfqpoint{0.407592in}{0.222904in}}{\pgfqpoint{1.411163in}{1.411163in}}%
\pgfusepath{clip}%
\pgfsetbuttcap%
\pgfsetroundjoin%
\definecolor{currentfill}{rgb}{0.909804,0.698039,0.423529}%
\pgfsetfillcolor{currentfill}%
\pgfsetfillopacity{0.120000}%
\pgfsetlinewidth{0.000000pt}%
\definecolor{currentstroke}{rgb}{0.000000,0.000000,0.000000}%
\pgfsetstrokecolor{currentstroke}%
\pgfsetdash{}{0pt}%
\pgfpathmoveto{\pgfqpoint{1.384069in}{0.961507in}}%
\pgfpathlineto{\pgfqpoint{1.402048in}{0.971542in}}%
\pgfpathlineto{\pgfqpoint{1.420106in}{0.981845in}}%
\pgfpathlineto{\pgfqpoint{1.428151in}{0.983191in}}%
\pgfpathlineto{\pgfqpoint{1.436170in}{0.984838in}}%
\pgfpathlineto{\pgfqpoint{1.418142in}{0.974368in}}%
\pgfpathlineto{\pgfqpoint{1.400194in}{0.964157in}}%
\pgfpathlineto{\pgfqpoint{1.392144in}{0.962676in}}%
\pgfpathlineto{\pgfqpoint{1.384069in}{0.961507in}}%
\pgfpathclose%
\pgfusepath{fill}%
\end{pgfscope}%
\begin{pgfscope}%
\pgfpathrectangle{\pgfqpoint{0.407592in}{0.222904in}}{\pgfqpoint{1.411163in}{1.411163in}}%
\pgfusepath{clip}%
\pgfsetbuttcap%
\pgfsetroundjoin%
\definecolor{currentfill}{rgb}{0.909804,0.698039,0.423529}%
\pgfsetfillcolor{currentfill}%
\pgfsetfillopacity{0.120000}%
\pgfsetlinewidth{0.000000pt}%
\definecolor{currentstroke}{rgb}{0.000000,0.000000,0.000000}%
\pgfsetstrokecolor{currentstroke}%
\pgfsetdash{}{0pt}%
\pgfpathmoveto{\pgfqpoint{1.436170in}{0.984838in}}%
\pgfpathlineto{\pgfqpoint{1.454279in}{0.995554in}}%
\pgfpathlineto{\pgfqpoint{1.472469in}{1.006502in}}%
\pgfpathlineto{\pgfqpoint{1.480435in}{1.008591in}}%
\pgfpathlineto{\pgfqpoint{1.488378in}{1.010977in}}%
\pgfpathlineto{\pgfqpoint{1.470217in}{0.999900in}}%
\pgfpathlineto{\pgfqpoint{1.452138in}{0.989050in}}%
\pgfpathlineto{\pgfqpoint{1.444166in}{0.986789in}}%
\pgfpathlineto{\pgfqpoint{1.436170in}{0.984838in}}%
\pgfpathclose%
\pgfusepath{fill}%
\end{pgfscope}%
\begin{pgfscope}%
\pgfpathrectangle{\pgfqpoint{0.407592in}{0.222904in}}{\pgfqpoint{1.411163in}{1.411163in}}%
\pgfusepath{clip}%
\pgfsetbuttcap%
\pgfsetroundjoin%
\definecolor{currentfill}{rgb}{0.909804,0.698039,0.423529}%
\pgfsetfillcolor{currentfill}%
\pgfsetfillopacity{0.120000}%
\pgfsetlinewidth{0.000000pt}%
\definecolor{currentstroke}{rgb}{0.000000,0.000000,0.000000}%
\pgfsetstrokecolor{currentstroke}%
\pgfsetdash{}{0pt}%
\pgfpathmoveto{\pgfqpoint{1.072284in}{0.896744in}}%
\pgfpathlineto{\pgfqpoint{1.089641in}{0.900223in}}%
\pgfpathlineto{\pgfqpoint{1.107059in}{0.904247in}}%
\pgfpathlineto{\pgfqpoint{1.115701in}{0.900749in}}%
\pgfpathlineto{\pgfqpoint{1.124306in}{0.897485in}}%
\pgfpathlineto{\pgfqpoint{1.106933in}{0.893187in}}%
\pgfpathlineto{\pgfqpoint{1.089620in}{0.889440in}}%
\pgfpathlineto{\pgfqpoint{1.080971in}{0.892976in}}%
\pgfpathlineto{\pgfqpoint{1.072284in}{0.896744in}}%
\pgfpathclose%
\pgfusepath{fill}%
\end{pgfscope}%
\begin{pgfscope}%
\pgfpathrectangle{\pgfqpoint{0.407592in}{0.222904in}}{\pgfqpoint{1.411163in}{1.411163in}}%
\pgfusepath{clip}%
\pgfsetbuttcap%
\pgfsetroundjoin%
\definecolor{currentfill}{rgb}{0.909804,0.698039,0.423529}%
\pgfsetfillcolor{currentfill}%
\pgfsetfillopacity{0.120000}%
\pgfsetlinewidth{0.000000pt}%
\definecolor{currentstroke}{rgb}{0.000000,0.000000,0.000000}%
\pgfsetstrokecolor{currentstroke}%
\pgfsetdash{}{0pt}%
\pgfpathmoveto{\pgfqpoint{1.488378in}{1.010977in}}%
\pgfpathlineto{\pgfqpoint{1.506621in}{1.022267in}}%
\pgfpathlineto{\pgfqpoint{1.524947in}{1.033761in}}%
\pgfpathlineto{\pgfqpoint{1.532839in}{1.036544in}}%
\pgfpathlineto{\pgfqpoint{1.540709in}{1.039616in}}%
\pgfpathlineto{\pgfqpoint{1.522412in}{1.028031in}}%
\pgfpathlineto{\pgfqpoint{1.504197in}{1.016645in}}%
\pgfpathlineto{\pgfqpoint{1.496298in}{1.013660in}}%
\pgfpathlineto{\pgfqpoint{1.488378in}{1.010977in}}%
\pgfpathclose%
\pgfusepath{fill}%
\end{pgfscope}%
\begin{pgfscope}%
\pgfpathrectangle{\pgfqpoint{0.407592in}{0.222904in}}{\pgfqpoint{1.411163in}{1.411163in}}%
\pgfusepath{clip}%
\pgfsetbuttcap%
\pgfsetroundjoin%
\definecolor{currentfill}{rgb}{0.909804,0.698039,0.423529}%
\pgfsetfillcolor{currentfill}%
\pgfsetfillopacity{0.120000}%
\pgfsetlinewidth{0.000000pt}%
\definecolor{currentstroke}{rgb}{0.000000,0.000000,0.000000}%
\pgfsetstrokecolor{currentstroke}%
\pgfsetdash{}{0pt}%
\pgfpathmoveto{\pgfqpoint{1.020156in}{0.900174in}}%
\pgfpathlineto{\pgfqpoint{1.037447in}{0.902265in}}%
\pgfpathlineto{\pgfqpoint{1.054793in}{0.904931in}}%
\pgfpathlineto{\pgfqpoint{1.063558in}{0.900732in}}%
\pgfpathlineto{\pgfqpoint{1.072284in}{0.896744in}}%
\pgfpathlineto{\pgfqpoint{1.054984in}{0.893830in}}%
\pgfpathlineto{\pgfqpoint{1.037741in}{0.891505in}}%
\pgfpathlineto{\pgfqpoint{1.028968in}{0.895736in}}%
\pgfpathlineto{\pgfqpoint{1.020156in}{0.900174in}}%
\pgfpathclose%
\pgfusepath{fill}%
\end{pgfscope}%
\begin{pgfscope}%
\pgfpathrectangle{\pgfqpoint{0.407592in}{0.222904in}}{\pgfqpoint{1.411163in}{1.411163in}}%
\pgfusepath{clip}%
\pgfsetbuttcap%
\pgfsetroundjoin%
\definecolor{currentfill}{rgb}{0.909804,0.698039,0.423529}%
\pgfsetfillcolor{currentfill}%
\pgfsetfillopacity{0.120000}%
\pgfsetlinewidth{0.000000pt}%
\definecolor{currentstroke}{rgb}{0.000000,0.000000,0.000000}%
\pgfsetstrokecolor{currentstroke}%
\pgfsetdash{}{0pt}%
\pgfpathmoveto{\pgfqpoint{1.593177in}{1.070469in}}%
\pgfpathlineto{\pgfqpoint{1.611699in}{1.082647in}}%
\pgfpathlineto{\pgfqpoint{1.619467in}{1.086644in}}%
\pgfpathlineto{\pgfqpoint{1.627215in}{1.090915in}}%
\pgfpathlineto{\pgfqpoint{1.608721in}{1.078710in}}%
\pgfpathlineto{\pgfqpoint{1.600959in}{1.074450in}}%
\pgfusepath{fill}%
\end{pgfscope}%
\begin{pgfscope}%
\pgfpathrectangle{\pgfqpoint{0.407592in}{0.222904in}}{\pgfqpoint{1.411163in}{1.411163in}}%
\pgfusepath{clip}%
\pgfsetbuttcap%
\pgfsetroundjoin%
\definecolor{currentfill}{rgb}{0.909804,0.698039,0.423529}%
\pgfsetfillcolor{currentfill}%
\pgfsetfillopacity{0.120000}%
\pgfsetlinewidth{0.000000pt}%
\definecolor{currentstroke}{rgb}{0.000000,0.000000,0.000000}%
\pgfsetstrokecolor{currentstroke}%
\pgfsetdash{}{0pt}%
\pgfpathmoveto{\pgfqpoint{1.540709in}{1.039616in}}%
\pgfpathlineto{\pgfqpoint{1.559090in}{1.051389in}}%
\pgfpathlineto{\pgfqpoint{1.577555in}{1.063342in}}%
\pgfpathlineto{\pgfqpoint{1.585376in}{1.066767in}}%
\pgfpathlineto{\pgfqpoint{1.593177in}{1.070469in}}%
\pgfpathlineto{\pgfqpoint{1.574740in}{1.058460in}}%
\pgfpathlineto{\pgfqpoint{1.556387in}{1.046628in}}%
\pgfpathlineto{\pgfqpoint{1.548558in}{1.042976in}}%
\pgfpathlineto{\pgfqpoint{1.540709in}{1.039616in}}%
\pgfpathclose%
\pgfusepath{fill}%
\end{pgfscope}%
\begin{pgfscope}%
\pgfpathrectangle{\pgfqpoint{0.407592in}{0.222904in}}{\pgfqpoint{1.411163in}{1.411163in}}%
\pgfusepath{clip}%
\pgfsetbuttcap%
\pgfsetroundjoin%
\definecolor{currentfill}{rgb}{0.909804,0.698039,0.423529}%
\pgfsetfillcolor{currentfill}%
\pgfsetfillopacity{0.120000}%
\pgfsetlinewidth{0.000000pt}%
\definecolor{currentstroke}{rgb}{0.000000,0.000000,0.000000}%
\pgfsetstrokecolor{currentstroke}%
\pgfsetdash{}{0pt}%
\pgfpathmoveto{\pgfqpoint{1.263503in}{0.926433in}}%
\pgfpathlineto{\pgfqpoint{1.281275in}{0.934485in}}%
\pgfpathlineto{\pgfqpoint{1.299121in}{0.942899in}}%
\pgfpathlineto{\pgfqpoint{1.307399in}{0.942068in}}%
\pgfpathlineto{\pgfqpoint{1.315646in}{0.941519in}}%
\pgfpathlineto{\pgfqpoint{1.297835in}{0.932850in}}%
\pgfpathlineto{\pgfqpoint{1.280099in}{0.924536in}}%
\pgfpathlineto{\pgfqpoint{1.271817in}{0.925340in}}%
\pgfpathlineto{\pgfqpoint{1.263503in}{0.926433in}}%
\pgfpathclose%
\pgfusepath{fill}%
\end{pgfscope}%
\begin{pgfscope}%
\pgfpathrectangle{\pgfqpoint{0.407592in}{0.222904in}}{\pgfqpoint{1.411163in}{1.411163in}}%
\pgfusepath{clip}%
\pgfsetbuttcap%
\pgfsetroundjoin%
\definecolor{currentfill}{rgb}{0.909804,0.698039,0.423529}%
\pgfsetfillcolor{currentfill}%
\pgfsetfillopacity{0.120000}%
\pgfsetlinewidth{0.000000pt}%
\definecolor{currentstroke}{rgb}{0.000000,0.000000,0.000000}%
\pgfsetstrokecolor{currentstroke}%
\pgfsetdash{}{0pt}%
\pgfpathmoveto{\pgfqpoint{1.211379in}{0.915074in}}%
\pgfpathlineto{\pgfqpoint{1.229043in}{0.922057in}}%
\pgfpathlineto{\pgfqpoint{1.246779in}{0.929450in}}%
\pgfpathlineto{\pgfqpoint{1.255157in}{0.927805in}}%
\pgfpathlineto{\pgfqpoint{1.263503in}{0.926433in}}%
\pgfpathlineto{\pgfqpoint{1.245805in}{0.918763in}}%
\pgfpathlineto{\pgfqpoint{1.228178in}{0.911499in}}%
\pgfpathlineto{\pgfqpoint{1.219795in}{0.913149in}}%
\pgfpathlineto{\pgfqpoint{1.211379in}{0.915074in}}%
\pgfpathclose%
\pgfusepath{fill}%
\end{pgfscope}%
\begin{pgfscope}%
\pgfpathrectangle{\pgfqpoint{0.407592in}{0.222904in}}{\pgfqpoint{1.411163in}{1.411163in}}%
\pgfusepath{clip}%
\pgfsetbuttcap%
\pgfsetroundjoin%
\definecolor{currentfill}{rgb}{0.909804,0.698039,0.423529}%
\pgfsetfillcolor{currentfill}%
\pgfsetfillopacity{0.120000}%
\pgfsetlinewidth{0.000000pt}%
\definecolor{currentstroke}{rgb}{0.000000,0.000000,0.000000}%
\pgfsetstrokecolor{currentstroke}%
\pgfsetdash{}{0pt}%
\pgfpathmoveto{\pgfqpoint{1.315646in}{0.941519in}}%
\pgfpathlineto{\pgfqpoint{1.333534in}{0.950524in}}%
\pgfpathlineto{\pgfqpoint{1.351498in}{0.959845in}}%
\pgfpathlineto{\pgfqpoint{1.359683in}{0.959820in}}%
\pgfpathlineto{\pgfqpoint{1.367839in}{0.960085in}}%
\pgfpathlineto{\pgfqpoint{1.349907in}{0.950539in}}%
\pgfpathlineto{\pgfqpoint{1.332053in}{0.941301in}}%
\pgfpathlineto{\pgfqpoint{1.323864in}{0.941262in}}%
\pgfpathlineto{\pgfqpoint{1.315646in}{0.941519in}}%
\pgfpathclose%
\pgfusepath{fill}%
\end{pgfscope}%
\begin{pgfscope}%
\pgfpathrectangle{\pgfqpoint{0.407592in}{0.222904in}}{\pgfqpoint{1.411163in}{1.411163in}}%
\pgfusepath{clip}%
\pgfsetbuttcap%
\pgfsetroundjoin%
\definecolor{currentfill}{rgb}{0.909804,0.698039,0.423529}%
\pgfsetfillcolor{currentfill}%
\pgfsetfillopacity{0.120000}%
\pgfsetlinewidth{0.000000pt}%
\definecolor{currentstroke}{rgb}{0.000000,0.000000,0.000000}%
\pgfsetstrokecolor{currentstroke}%
\pgfsetdash{}{0pt}%
\pgfpathmoveto{\pgfqpoint{1.159242in}{0.907639in}}%
\pgfpathlineto{\pgfqpoint{1.176809in}{0.913447in}}%
\pgfpathlineto{\pgfqpoint{1.194444in}{0.919711in}}%
\pgfpathlineto{\pgfqpoint{1.202929in}{0.917265in}}%
\pgfpathlineto{\pgfqpoint{1.211379in}{0.915074in}}%
\pgfpathlineto{\pgfqpoint{1.193784in}{0.908523in}}%
\pgfpathlineto{\pgfqpoint{1.176259in}{0.902427in}}%
\pgfpathlineto{\pgfqpoint{1.167768in}{0.904905in}}%
\pgfpathlineto{\pgfqpoint{1.159242in}{0.907639in}}%
\pgfpathclose%
\pgfusepath{fill}%
\end{pgfscope}%
\begin{pgfscope}%
\pgfpathrectangle{\pgfqpoint{0.407592in}{0.222904in}}{\pgfqpoint{1.411163in}{1.411163in}}%
\pgfusepath{clip}%
\pgfsetbuttcap%
\pgfsetroundjoin%
\definecolor{currentfill}{rgb}{0.909804,0.698039,0.423529}%
\pgfsetfillcolor{currentfill}%
\pgfsetfillopacity{0.120000}%
\pgfsetlinewidth{0.000000pt}%
\definecolor{currentstroke}{rgb}{0.000000,0.000000,0.000000}%
\pgfsetstrokecolor{currentstroke}%
\pgfsetdash{}{0pt}%
\pgfpathmoveto{\pgfqpoint{1.367839in}{0.960085in}}%
\pgfpathlineto{\pgfqpoint{1.385849in}{0.969922in}}%
\pgfpathlineto{\pgfqpoint{1.403939in}{0.980035in}}%
\pgfpathlineto{\pgfqpoint{1.412035in}{0.980795in}}%
\pgfpathlineto{\pgfqpoint{1.420106in}{0.981845in}}%
\pgfpathlineto{\pgfqpoint{1.402048in}{0.971542in}}%
\pgfpathlineto{\pgfqpoint{1.384069in}{0.961507in}}%
\pgfpathlineto{\pgfqpoint{1.375967in}{0.960645in}}%
\pgfpathlineto{\pgfqpoint{1.367839in}{0.960085in}}%
\pgfpathclose%
\pgfusepath{fill}%
\end{pgfscope}%
\begin{pgfscope}%
\pgfpathrectangle{\pgfqpoint{0.407592in}{0.222904in}}{\pgfqpoint{1.411163in}{1.411163in}}%
\pgfusepath{clip}%
\pgfsetbuttcap%
\pgfsetroundjoin%
\definecolor{currentfill}{rgb}{0.909804,0.698039,0.423529}%
\pgfsetfillcolor{currentfill}%
\pgfsetfillopacity{0.120000}%
\pgfsetlinewidth{0.000000pt}%
\definecolor{currentstroke}{rgb}{0.000000,0.000000,0.000000}%
\pgfsetstrokecolor{currentstroke}%
\pgfsetdash{}{0pt}%
\pgfpathmoveto{\pgfqpoint{1.107059in}{0.904247in}}%
\pgfpathlineto{\pgfqpoint{1.124539in}{0.908792in}}%
\pgfpathlineto{\pgfqpoint{1.142083in}{0.913837in}}%
\pgfpathlineto{\pgfqpoint{1.150681in}{0.910620in}}%
\pgfpathlineto{\pgfqpoint{1.159242in}{0.907639in}}%
\pgfpathlineto{\pgfqpoint{1.141742in}{0.902311in}}%
\pgfpathlineto{\pgfqpoint{1.124306in}{0.897485in}}%
\pgfpathlineto{\pgfqpoint{1.115701in}{0.900749in}}%
\pgfpathlineto{\pgfqpoint{1.107059in}{0.904247in}}%
\pgfpathclose%
\pgfusepath{fill}%
\end{pgfscope}%
\begin{pgfscope}%
\pgfpathrectangle{\pgfqpoint{0.407592in}{0.222904in}}{\pgfqpoint{1.411163in}{1.411163in}}%
\pgfusepath{clip}%
\pgfsetbuttcap%
\pgfsetroundjoin%
\definecolor{currentfill}{rgb}{0.909804,0.698039,0.423529}%
\pgfsetfillcolor{currentfill}%
\pgfsetfillopacity{0.120000}%
\pgfsetlinewidth{0.000000pt}%
\definecolor{currentstroke}{rgb}{0.000000,0.000000,0.000000}%
\pgfsetstrokecolor{currentstroke}%
\pgfsetdash{}{0pt}%
\pgfpathmoveto{\pgfqpoint{1.420106in}{0.981845in}}%
\pgfpathlineto{\pgfqpoint{1.438244in}{0.992401in}}%
\pgfpathlineto{\pgfqpoint{1.456465in}{1.003196in}}%
\pgfpathlineto{\pgfqpoint{1.464479in}{1.004705in}}%
\pgfpathlineto{\pgfqpoint{1.472469in}{1.006502in}}%
\pgfpathlineto{\pgfqpoint{1.454279in}{0.995554in}}%
\pgfpathlineto{\pgfqpoint{1.436170in}{0.984838in}}%
\pgfpathlineto{\pgfqpoint{1.428151in}{0.983191in}}%
\pgfpathlineto{\pgfqpoint{1.420106in}{0.981845in}}%
\pgfpathclose%
\pgfusepath{fill}%
\end{pgfscope}%
\begin{pgfscope}%
\pgfpathrectangle{\pgfqpoint{0.407592in}{0.222904in}}{\pgfqpoint{1.411163in}{1.411163in}}%
\pgfusepath{clip}%
\pgfsetbuttcap%
\pgfsetroundjoin%
\definecolor{currentfill}{rgb}{0.909804,0.698039,0.423529}%
\pgfsetfillcolor{currentfill}%
\pgfsetfillopacity{0.120000}%
\pgfsetlinewidth{0.000000pt}%
\definecolor{currentstroke}{rgb}{0.000000,0.000000,0.000000}%
\pgfsetstrokecolor{currentstroke}%
\pgfsetdash{}{0pt}%
\pgfpathmoveto{\pgfqpoint{1.472469in}{1.006502in}}%
\pgfpathlineto{\pgfqpoint{1.490741in}{1.017671in}}%
\pgfpathlineto{\pgfqpoint{1.509096in}{1.029049in}}%
\pgfpathlineto{\pgfqpoint{1.517033in}{1.031263in}}%
\pgfpathlineto{\pgfqpoint{1.524947in}{1.033761in}}%
\pgfpathlineto{\pgfqpoint{1.506621in}{1.022267in}}%
\pgfpathlineto{\pgfqpoint{1.488378in}{1.010977in}}%
\pgfpathlineto{\pgfqpoint{1.480435in}{1.008591in}}%
\pgfpathlineto{\pgfqpoint{1.472469in}{1.006502in}}%
\pgfpathclose%
\pgfusepath{fill}%
\end{pgfscope}%
\begin{pgfscope}%
\pgfpathrectangle{\pgfqpoint{0.407592in}{0.222904in}}{\pgfqpoint{1.411163in}{1.411163in}}%
\pgfusepath{clip}%
\pgfsetbuttcap%
\pgfsetroundjoin%
\definecolor{currentfill}{rgb}{0.909804,0.698039,0.423529}%
\pgfsetfillcolor{currentfill}%
\pgfsetfillopacity{0.120000}%
\pgfsetlinewidth{0.000000pt}%
\definecolor{currentstroke}{rgb}{0.000000,0.000000,0.000000}%
\pgfsetstrokecolor{currentstroke}%
\pgfsetdash{}{0pt}%
\pgfpathmoveto{\pgfqpoint{1.054793in}{0.904931in}}%
\pgfpathlineto{\pgfqpoint{1.072197in}{0.908152in}}%
\pgfpathlineto{\pgfqpoint{1.089661in}{0.911908in}}%
\pgfpathlineto{\pgfqpoint{1.098379in}{0.907970in}}%
\pgfpathlineto{\pgfqpoint{1.107059in}{0.904247in}}%
\pgfpathlineto{\pgfqpoint{1.089641in}{0.900223in}}%
\pgfpathlineto{\pgfqpoint{1.072284in}{0.896744in}}%
\pgfpathlineto{\pgfqpoint{1.063558in}{0.900732in}}%
\pgfpathlineto{\pgfqpoint{1.054793in}{0.904931in}}%
\pgfpathclose%
\pgfusepath{fill}%
\end{pgfscope}%
\begin{pgfscope}%
\pgfpathrectangle{\pgfqpoint{0.407592in}{0.222904in}}{\pgfqpoint{1.411163in}{1.411163in}}%
\pgfusepath{clip}%
\pgfsetbuttcap%
\pgfsetroundjoin%
\definecolor{currentfill}{rgb}{0.909804,0.698039,0.423529}%
\pgfsetfillcolor{currentfill}%
\pgfsetfillopacity{0.120000}%
\pgfsetlinewidth{0.000000pt}%
\definecolor{currentstroke}{rgb}{0.000000,0.000000,0.000000}%
\pgfsetstrokecolor{currentstroke}%
\pgfsetdash{}{0pt}%
\pgfpathmoveto{\pgfqpoint{1.577555in}{1.063342in}}%
\pgfpathlineto{\pgfqpoint{1.596105in}{1.075467in}}%
\pgfpathlineto{\pgfqpoint{1.603912in}{1.078921in}}%
\pgfpathlineto{\pgfqpoint{1.611699in}{1.082647in}}%
\pgfpathlineto{\pgfqpoint{1.593177in}{1.070469in}}%
\pgfpathlineto{\pgfqpoint{1.585376in}{1.066767in}}%
\pgfusepath{fill}%
\end{pgfscope}%
\begin{pgfscope}%
\pgfpathrectangle{\pgfqpoint{0.407592in}{0.222904in}}{\pgfqpoint{1.411163in}{1.411163in}}%
\pgfusepath{clip}%
\pgfsetbuttcap%
\pgfsetroundjoin%
\definecolor{currentfill}{rgb}{0.909804,0.698039,0.423529}%
\pgfsetfillcolor{currentfill}%
\pgfsetfillopacity{0.120000}%
\pgfsetlinewidth{0.000000pt}%
\definecolor{currentstroke}{rgb}{0.000000,0.000000,0.000000}%
\pgfsetstrokecolor{currentstroke}%
\pgfsetdash{}{0pt}%
\pgfpathmoveto{\pgfqpoint{1.524947in}{1.033761in}}%
\pgfpathlineto{\pgfqpoint{1.543356in}{1.045449in}}%
\pgfpathlineto{\pgfqpoint{1.561850in}{1.057321in}}%
\pgfpathlineto{\pgfqpoint{1.569713in}{1.060194in}}%
\pgfpathlineto{\pgfqpoint{1.577555in}{1.063342in}}%
\pgfpathlineto{\pgfqpoint{1.559090in}{1.051389in}}%
\pgfpathlineto{\pgfqpoint{1.540709in}{1.039616in}}%
\pgfpathlineto{\pgfqpoint{1.532839in}{1.036544in}}%
\pgfpathlineto{\pgfqpoint{1.524947in}{1.033761in}}%
\pgfpathclose%
\pgfusepath{fill}%
\end{pgfscope}%
\begin{pgfscope}%
\pgfpathrectangle{\pgfqpoint{0.407592in}{0.222904in}}{\pgfqpoint{1.411163in}{1.411163in}}%
\pgfusepath{clip}%
\pgfsetbuttcap%
\pgfsetroundjoin%
\definecolor{currentfill}{rgb}{0.909804,0.698039,0.423529}%
\pgfsetfillcolor{currentfill}%
\pgfsetfillopacity{0.120000}%
\pgfsetlinewidth{0.000000pt}%
\definecolor{currentstroke}{rgb}{0.000000,0.000000,0.000000}%
\pgfsetstrokecolor{currentstroke}%
\pgfsetdash{}{0pt}%
\pgfpathmoveto{\pgfqpoint{1.246779in}{0.929450in}}%
\pgfpathlineto{\pgfqpoint{1.264588in}{0.937230in}}%
\pgfpathlineto{\pgfqpoint{1.282471in}{0.945378in}}%
\pgfpathlineto{\pgfqpoint{1.290812in}{0.944004in}}%
\pgfpathlineto{\pgfqpoint{1.299121in}{0.942899in}}%
\pgfpathlineto{\pgfqpoint{1.281275in}{0.934485in}}%
\pgfpathlineto{\pgfqpoint{1.263503in}{0.926433in}}%
\pgfpathlineto{\pgfqpoint{1.255157in}{0.927805in}}%
\pgfpathlineto{\pgfqpoint{1.246779in}{0.929450in}}%
\pgfpathclose%
\pgfusepath{fill}%
\end{pgfscope}%
\begin{pgfscope}%
\pgfpathrectangle{\pgfqpoint{0.407592in}{0.222904in}}{\pgfqpoint{1.411163in}{1.411163in}}%
\pgfusepath{clip}%
\pgfsetbuttcap%
\pgfsetroundjoin%
\definecolor{currentfill}{rgb}{0.909804,0.698039,0.423529}%
\pgfsetfillcolor{currentfill}%
\pgfsetfillopacity{0.120000}%
\pgfsetlinewidth{0.000000pt}%
\definecolor{currentstroke}{rgb}{0.000000,0.000000,0.000000}%
\pgfsetstrokecolor{currentstroke}%
\pgfsetdash{}{0pt}%
\pgfpathmoveto{\pgfqpoint{1.299121in}{0.942899in}}%
\pgfpathlineto{\pgfqpoint{1.317043in}{0.951655in}}%
\pgfpathlineto{\pgfqpoint{1.335042in}{0.960734in}}%
\pgfpathlineto{\pgfqpoint{1.343285in}{0.960152in}}%
\pgfpathlineto{\pgfqpoint{1.351498in}{0.959845in}}%
\pgfpathlineto{\pgfqpoint{1.333534in}{0.950524in}}%
\pgfpathlineto{\pgfqpoint{1.315646in}{0.941519in}}%
\pgfpathlineto{\pgfqpoint{1.307399in}{0.942068in}}%
\pgfpathlineto{\pgfqpoint{1.299121in}{0.942899in}}%
\pgfpathclose%
\pgfusepath{fill}%
\end{pgfscope}%
\begin{pgfscope}%
\pgfpathrectangle{\pgfqpoint{0.407592in}{0.222904in}}{\pgfqpoint{1.411163in}{1.411163in}}%
\pgfusepath{clip}%
\pgfsetbuttcap%
\pgfsetroundjoin%
\definecolor{currentfill}{rgb}{0.909804,0.698039,0.423529}%
\pgfsetfillcolor{currentfill}%
\pgfsetfillopacity{0.120000}%
\pgfsetlinewidth{0.000000pt}%
\definecolor{currentstroke}{rgb}{0.000000,0.000000,0.000000}%
\pgfsetstrokecolor{currentstroke}%
\pgfsetdash{}{0pt}%
\pgfpathmoveto{\pgfqpoint{1.194444in}{0.919711in}}%
\pgfpathlineto{\pgfqpoint{1.212148in}{0.926409in}}%
\pgfpathlineto{\pgfqpoint{1.229923in}{0.933519in}}%
\pgfpathlineto{\pgfqpoint{1.238368in}{0.931357in}}%
\pgfpathlineto{\pgfqpoint{1.246779in}{0.929450in}}%
\pgfpathlineto{\pgfqpoint{1.229043in}{0.922057in}}%
\pgfpathlineto{\pgfqpoint{1.211379in}{0.915074in}}%
\pgfpathlineto{\pgfqpoint{1.202929in}{0.917265in}}%
\pgfpathlineto{\pgfqpoint{1.194444in}{0.919711in}}%
\pgfpathclose%
\pgfusepath{fill}%
\end{pgfscope}%
\begin{pgfscope}%
\pgfpathrectangle{\pgfqpoint{0.407592in}{0.222904in}}{\pgfqpoint{1.411163in}{1.411163in}}%
\pgfusepath{clip}%
\pgfsetbuttcap%
\pgfsetroundjoin%
\definecolor{currentfill}{rgb}{0.909804,0.698039,0.423529}%
\pgfsetfillcolor{currentfill}%
\pgfsetfillopacity{0.120000}%
\pgfsetlinewidth{0.000000pt}%
\definecolor{currentstroke}{rgb}{0.000000,0.000000,0.000000}%
\pgfsetstrokecolor{currentstroke}%
\pgfsetdash{}{0pt}%
\pgfpathmoveto{\pgfqpoint{1.351498in}{0.959845in}}%
\pgfpathlineto{\pgfqpoint{1.369541in}{0.969465in}}%
\pgfpathlineto{\pgfqpoint{1.387663in}{0.979368in}}%
\pgfpathlineto{\pgfqpoint{1.395815in}{0.979561in}}%
\pgfpathlineto{\pgfqpoint{1.403939in}{0.980035in}}%
\pgfpathlineto{\pgfqpoint{1.385849in}{0.969922in}}%
\pgfpathlineto{\pgfqpoint{1.367839in}{0.960085in}}%
\pgfpathlineto{\pgfqpoint{1.359683in}{0.959820in}}%
\pgfpathlineto{\pgfqpoint{1.351498in}{0.959845in}}%
\pgfpathclose%
\pgfusepath{fill}%
\end{pgfscope}%
\begin{pgfscope}%
\pgfpathrectangle{\pgfqpoint{0.407592in}{0.222904in}}{\pgfqpoint{1.411163in}{1.411163in}}%
\pgfusepath{clip}%
\pgfsetbuttcap%
\pgfsetroundjoin%
\definecolor{currentfill}{rgb}{0.909804,0.698039,0.423529}%
\pgfsetfillcolor{currentfill}%
\pgfsetfillopacity{0.120000}%
\pgfsetlinewidth{0.000000pt}%
\definecolor{currentstroke}{rgb}{0.000000,0.000000,0.000000}%
\pgfsetstrokecolor{currentstroke}%
\pgfsetdash{}{0pt}%
\pgfpathmoveto{\pgfqpoint{1.142083in}{0.913837in}}%
\pgfpathlineto{\pgfqpoint{1.159692in}{0.919358in}}%
\pgfpathlineto{\pgfqpoint{1.177369in}{0.925336in}}%
\pgfpathlineto{\pgfqpoint{1.185924in}{0.922404in}}%
\pgfpathlineto{\pgfqpoint{1.194444in}{0.919711in}}%
\pgfpathlineto{\pgfqpoint{1.176809in}{0.913447in}}%
\pgfpathlineto{\pgfqpoint{1.159242in}{0.907639in}}%
\pgfpathlineto{\pgfqpoint{1.150681in}{0.910620in}}%
\pgfpathlineto{\pgfqpoint{1.142083in}{0.913837in}}%
\pgfpathclose%
\pgfusepath{fill}%
\end{pgfscope}%
\begin{pgfscope}%
\pgfpathrectangle{\pgfqpoint{0.407592in}{0.222904in}}{\pgfqpoint{1.411163in}{1.411163in}}%
\pgfusepath{clip}%
\pgfsetbuttcap%
\pgfsetroundjoin%
\definecolor{currentfill}{rgb}{0.909804,0.698039,0.423529}%
\pgfsetfillcolor{currentfill}%
\pgfsetfillopacity{0.120000}%
\pgfsetlinewidth{0.000000pt}%
\definecolor{currentstroke}{rgb}{0.000000,0.000000,0.000000}%
\pgfsetstrokecolor{currentstroke}%
\pgfsetdash{}{0pt}%
\pgfpathmoveto{\pgfqpoint{1.403939in}{0.980035in}}%
\pgfpathlineto{\pgfqpoint{1.422108in}{0.990409in}}%
\pgfpathlineto{\pgfqpoint{1.440360in}{1.001030in}}%
\pgfpathlineto{\pgfqpoint{1.448425in}{1.001973in}}%
\pgfpathlineto{\pgfqpoint{1.456465in}{1.003196in}}%
\pgfpathlineto{\pgfqpoint{1.438244in}{0.992401in}}%
\pgfpathlineto{\pgfqpoint{1.420106in}{0.981845in}}%
\pgfpathlineto{\pgfqpoint{1.412035in}{0.980795in}}%
\pgfpathlineto{\pgfqpoint{1.403939in}{0.980035in}}%
\pgfpathclose%
\pgfusepath{fill}%
\end{pgfscope}%
\begin{pgfscope}%
\pgfpathrectangle{\pgfqpoint{0.407592in}{0.222904in}}{\pgfqpoint{1.411163in}{1.411163in}}%
\pgfusepath{clip}%
\pgfsetbuttcap%
\pgfsetroundjoin%
\definecolor{currentfill}{rgb}{0.909804,0.698039,0.423529}%
\pgfsetfillcolor{currentfill}%
\pgfsetfillopacity{0.120000}%
\pgfsetlinewidth{0.000000pt}%
\definecolor{currentstroke}{rgb}{0.000000,0.000000,0.000000}%
\pgfsetstrokecolor{currentstroke}%
\pgfsetdash{}{0pt}%
\pgfpathmoveto{\pgfqpoint{1.089661in}{0.911908in}}%
\pgfpathlineto{\pgfqpoint{1.107186in}{0.916179in}}%
\pgfpathlineto{\pgfqpoint{1.124775in}{0.920945in}}%
\pgfpathlineto{\pgfqpoint{1.133447in}{0.917282in}}%
\pgfpathlineto{\pgfqpoint{1.142083in}{0.913837in}}%
\pgfpathlineto{\pgfqpoint{1.124539in}{0.908792in}}%
\pgfpathlineto{\pgfqpoint{1.107059in}{0.904247in}}%
\pgfpathlineto{\pgfqpoint{1.098379in}{0.907970in}}%
\pgfpathlineto{\pgfqpoint{1.089661in}{0.911908in}}%
\pgfpathclose%
\pgfusepath{fill}%
\end{pgfscope}%
\begin{pgfscope}%
\pgfpathrectangle{\pgfqpoint{0.407592in}{0.222904in}}{\pgfqpoint{1.411163in}{1.411163in}}%
\pgfusepath{clip}%
\pgfsetbuttcap%
\pgfsetroundjoin%
\definecolor{currentfill}{rgb}{0.909804,0.698039,0.423529}%
\pgfsetfillcolor{currentfill}%
\pgfsetfillopacity{0.120000}%
\pgfsetlinewidth{0.000000pt}%
\definecolor{currentstroke}{rgb}{0.000000,0.000000,0.000000}%
\pgfsetstrokecolor{currentstroke}%
\pgfsetdash{}{0pt}%
\pgfpathmoveto{\pgfqpoint{1.456465in}{1.003196in}}%
\pgfpathlineto{\pgfqpoint{1.474767in}{1.014219in}}%
\pgfpathlineto{\pgfqpoint{1.493152in}{1.025459in}}%
\pgfpathlineto{\pgfqpoint{1.501136in}{1.027115in}}%
\pgfpathlineto{\pgfqpoint{1.509096in}{1.029049in}}%
\pgfpathlineto{\pgfqpoint{1.490741in}{1.017671in}}%
\pgfpathlineto{\pgfqpoint{1.472469in}{1.006502in}}%
\pgfpathlineto{\pgfqpoint{1.464479in}{1.004705in}}%
\pgfpathlineto{\pgfqpoint{1.456465in}{1.003196in}}%
\pgfpathclose%
\pgfusepath{fill}%
\end{pgfscope}%
\begin{pgfscope}%
\pgfpathrectangle{\pgfqpoint{0.407592in}{0.222904in}}{\pgfqpoint{1.411163in}{1.411163in}}%
\pgfusepath{clip}%
\pgfsetbuttcap%
\pgfsetroundjoin%
\definecolor{currentfill}{rgb}{0.909804,0.698039,0.423529}%
\pgfsetfillcolor{currentfill}%
\pgfsetfillopacity{0.120000}%
\pgfsetlinewidth{0.000000pt}%
\definecolor{currentstroke}{rgb}{0.000000,0.000000,0.000000}%
\pgfsetstrokecolor{currentstroke}%
\pgfsetdash{}{0pt}%
\pgfpathmoveto{\pgfqpoint{1.561850in}{1.057321in}}%
\pgfpathlineto{\pgfqpoint{1.580430in}{1.069371in}}%
\pgfpathlineto{\pgfqpoint{1.588278in}{1.072284in}}%
\pgfpathlineto{\pgfqpoint{1.596105in}{1.075467in}}%
\pgfpathlineto{\pgfqpoint{1.577555in}{1.063342in}}%
\pgfpathlineto{\pgfqpoint{1.569713in}{1.060194in}}%
\pgfusepath{fill}%
\end{pgfscope}%
\begin{pgfscope}%
\pgfpathrectangle{\pgfqpoint{0.407592in}{0.222904in}}{\pgfqpoint{1.411163in}{1.411163in}}%
\pgfusepath{clip}%
\pgfsetbuttcap%
\pgfsetroundjoin%
\definecolor{currentfill}{rgb}{0.909804,0.698039,0.423529}%
\pgfsetfillcolor{currentfill}%
\pgfsetfillopacity{0.120000}%
\pgfsetlinewidth{0.000000pt}%
\definecolor{currentstroke}{rgb}{0.000000,0.000000,0.000000}%
\pgfsetstrokecolor{currentstroke}%
\pgfsetdash{}{0pt}%
\pgfpathmoveto{\pgfqpoint{1.509096in}{1.029049in}}%
\pgfpathlineto{\pgfqpoint{1.527535in}{1.040627in}}%
\pgfpathlineto{\pgfqpoint{1.546059in}{1.052395in}}%
\pgfpathlineto{\pgfqpoint{1.553966in}{1.054722in}}%
\pgfpathlineto{\pgfqpoint{1.561850in}{1.057321in}}%
\pgfpathlineto{\pgfqpoint{1.543356in}{1.045449in}}%
\pgfpathlineto{\pgfqpoint{1.524947in}{1.033761in}}%
\pgfpathlineto{\pgfqpoint{1.517033in}{1.031263in}}%
\pgfpathlineto{\pgfqpoint{1.509096in}{1.029049in}}%
\pgfpathclose%
\pgfusepath{fill}%
\end{pgfscope}%
\begin{pgfscope}%
\pgfpathrectangle{\pgfqpoint{0.407592in}{0.222904in}}{\pgfqpoint{1.411163in}{1.411163in}}%
\pgfusepath{clip}%
\pgfsetbuttcap%
\pgfsetroundjoin%
\definecolor{currentfill}{rgb}{0.909804,0.698039,0.423529}%
\pgfsetfillcolor{currentfill}%
\pgfsetfillopacity{0.120000}%
\pgfsetlinewidth{0.000000pt}%
\definecolor{currentstroke}{rgb}{0.000000,0.000000,0.000000}%
\pgfsetstrokecolor{currentstroke}%
\pgfsetdash{}{0pt}%
\pgfpathmoveto{\pgfqpoint{1.282471in}{0.945378in}}%
\pgfpathlineto{\pgfqpoint{1.300430in}{0.953874in}}%
\pgfpathlineto{\pgfqpoint{1.318465in}{0.962701in}}%
\pgfpathlineto{\pgfqpoint{1.326769in}{0.961587in}}%
\pgfpathlineto{\pgfqpoint{1.335042in}{0.960734in}}%
\pgfpathlineto{\pgfqpoint{1.317043in}{0.951655in}}%
\pgfpathlineto{\pgfqpoint{1.299121in}{0.942899in}}%
\pgfpathlineto{\pgfqpoint{1.290812in}{0.944004in}}%
\pgfpathlineto{\pgfqpoint{1.282471in}{0.945378in}}%
\pgfpathclose%
\pgfusepath{fill}%
\end{pgfscope}%
\begin{pgfscope}%
\pgfpathrectangle{\pgfqpoint{0.407592in}{0.222904in}}{\pgfqpoint{1.411163in}{1.411163in}}%
\pgfusepath{clip}%
\pgfsetbuttcap%
\pgfsetroundjoin%
\definecolor{currentfill}{rgb}{0.909804,0.698039,0.423529}%
\pgfsetfillcolor{currentfill}%
\pgfsetfillopacity{0.120000}%
\pgfsetlinewidth{0.000000pt}%
\definecolor{currentstroke}{rgb}{0.000000,0.000000,0.000000}%
\pgfsetstrokecolor{currentstroke}%
\pgfsetdash{}{0pt}%
\pgfpathmoveto{\pgfqpoint{1.229923in}{0.933519in}}%
\pgfpathlineto{\pgfqpoint{1.247771in}{0.941021in}}%
\pgfpathlineto{\pgfqpoint{1.265692in}{0.948896in}}%
\pgfpathlineto{\pgfqpoint{1.274098in}{0.947011in}}%
\pgfpathlineto{\pgfqpoint{1.282471in}{0.945378in}}%
\pgfpathlineto{\pgfqpoint{1.264588in}{0.937230in}}%
\pgfpathlineto{\pgfqpoint{1.246779in}{0.929450in}}%
\pgfpathlineto{\pgfqpoint{1.238368in}{0.931357in}}%
\pgfpathlineto{\pgfqpoint{1.229923in}{0.933519in}}%
\pgfpathclose%
\pgfusepath{fill}%
\end{pgfscope}%
\begin{pgfscope}%
\pgfpathrectangle{\pgfqpoint{0.407592in}{0.222904in}}{\pgfqpoint{1.411163in}{1.411163in}}%
\pgfusepath{clip}%
\pgfsetbuttcap%
\pgfsetroundjoin%
\definecolor{currentfill}{rgb}{0.909804,0.698039,0.423529}%
\pgfsetfillcolor{currentfill}%
\pgfsetfillopacity{0.120000}%
\pgfsetlinewidth{0.000000pt}%
\definecolor{currentstroke}{rgb}{0.000000,0.000000,0.000000}%
\pgfsetstrokecolor{currentstroke}%
\pgfsetdash{}{0pt}%
\pgfpathmoveto{\pgfqpoint{1.335042in}{0.960734in}}%
\pgfpathlineto{\pgfqpoint{1.353119in}{0.970121in}}%
\pgfpathlineto{\pgfqpoint{1.371275in}{0.979799in}}%
\pgfpathlineto{\pgfqpoint{1.379484in}{0.979449in}}%
\pgfpathlineto{\pgfqpoint{1.387663in}{0.979368in}}%
\pgfpathlineto{\pgfqpoint{1.369541in}{0.969465in}}%
\pgfpathlineto{\pgfqpoint{1.351498in}{0.959845in}}%
\pgfpathlineto{\pgfqpoint{1.343285in}{0.960152in}}%
\pgfpathlineto{\pgfqpoint{1.335042in}{0.960734in}}%
\pgfpathclose%
\pgfusepath{fill}%
\end{pgfscope}%
\begin{pgfscope}%
\pgfpathrectangle{\pgfqpoint{0.407592in}{0.222904in}}{\pgfqpoint{1.411163in}{1.411163in}}%
\pgfusepath{clip}%
\pgfsetbuttcap%
\pgfsetroundjoin%
\definecolor{currentfill}{rgb}{0.909804,0.698039,0.423529}%
\pgfsetfillcolor{currentfill}%
\pgfsetfillopacity{0.120000}%
\pgfsetlinewidth{0.000000pt}%
\definecolor{currentstroke}{rgb}{0.000000,0.000000,0.000000}%
\pgfsetstrokecolor{currentstroke}%
\pgfsetdash{}{0pt}%
\pgfpathmoveto{\pgfqpoint{1.177369in}{0.925336in}}%
\pgfpathlineto{\pgfqpoint{1.195114in}{0.931747in}}%
\pgfpathlineto{\pgfqpoint{1.212930in}{0.938574in}}%
\pgfpathlineto{\pgfqpoint{1.221444in}{0.935927in}}%
\pgfpathlineto{\pgfqpoint{1.229923in}{0.933519in}}%
\pgfpathlineto{\pgfqpoint{1.212148in}{0.926409in}}%
\pgfpathlineto{\pgfqpoint{1.194444in}{0.919711in}}%
\pgfpathlineto{\pgfqpoint{1.185924in}{0.922404in}}%
\pgfpathlineto{\pgfqpoint{1.177369in}{0.925336in}}%
\pgfpathclose%
\pgfusepath{fill}%
\end{pgfscope}%
\begin{pgfscope}%
\pgfpathrectangle{\pgfqpoint{0.407592in}{0.222904in}}{\pgfqpoint{1.411163in}{1.411163in}}%
\pgfusepath{clip}%
\pgfsetbuttcap%
\pgfsetroundjoin%
\definecolor{currentfill}{rgb}{0.909804,0.698039,0.423529}%
\pgfsetfillcolor{currentfill}%
\pgfsetfillopacity{0.120000}%
\pgfsetlinewidth{0.000000pt}%
\definecolor{currentstroke}{rgb}{0.000000,0.000000,0.000000}%
\pgfsetstrokecolor{currentstroke}%
\pgfsetdash{}{0pt}%
\pgfpathmoveto{\pgfqpoint{1.387663in}{0.979368in}}%
\pgfpathlineto{\pgfqpoint{1.405866in}{0.989540in}}%
\pgfpathlineto{\pgfqpoint{1.424150in}{0.999967in}}%
\pgfpathlineto{\pgfqpoint{1.432268in}{1.000363in}}%
\pgfpathlineto{\pgfqpoint{1.440360in}{1.001030in}}%
\pgfpathlineto{\pgfqpoint{1.422108in}{0.990409in}}%
\pgfpathlineto{\pgfqpoint{1.403939in}{0.980035in}}%
\pgfpathlineto{\pgfqpoint{1.395815in}{0.979561in}}%
\pgfpathlineto{\pgfqpoint{1.387663in}{0.979368in}}%
\pgfpathclose%
\pgfusepath{fill}%
\end{pgfscope}%
\begin{pgfscope}%
\pgfpathrectangle{\pgfqpoint{0.407592in}{0.222904in}}{\pgfqpoint{1.411163in}{1.411163in}}%
\pgfusepath{clip}%
\pgfsetbuttcap%
\pgfsetroundjoin%
\definecolor{currentfill}{rgb}{0.909804,0.698039,0.423529}%
\pgfsetfillcolor{currentfill}%
\pgfsetfillopacity{0.120000}%
\pgfsetlinewidth{0.000000pt}%
\definecolor{currentstroke}{rgb}{0.000000,0.000000,0.000000}%
\pgfsetstrokecolor{currentstroke}%
\pgfsetdash{}{0pt}%
\pgfpathmoveto{\pgfqpoint{1.124775in}{0.920945in}}%
\pgfpathlineto{\pgfqpoint{1.142428in}{0.926184in}}%
\pgfpathlineto{\pgfqpoint{1.160148in}{0.931878in}}%
\pgfpathlineto{\pgfqpoint{1.168777in}{0.928496in}}%
\pgfpathlineto{\pgfqpoint{1.177369in}{0.925336in}}%
\pgfpathlineto{\pgfqpoint{1.159692in}{0.919358in}}%
\pgfpathlineto{\pgfqpoint{1.142083in}{0.913837in}}%
\pgfpathlineto{\pgfqpoint{1.133447in}{0.917282in}}%
\pgfpathlineto{\pgfqpoint{1.124775in}{0.920945in}}%
\pgfpathclose%
\pgfusepath{fill}%
\end{pgfscope}%
\begin{pgfscope}%
\pgfpathrectangle{\pgfqpoint{0.407592in}{0.222904in}}{\pgfqpoint{1.411163in}{1.411163in}}%
\pgfusepath{clip}%
\pgfsetbuttcap%
\pgfsetroundjoin%
\definecolor{currentfill}{rgb}{0.909804,0.698039,0.423529}%
\pgfsetfillcolor{currentfill}%
\pgfsetfillopacity{0.120000}%
\pgfsetlinewidth{0.000000pt}%
\definecolor{currentstroke}{rgb}{0.000000,0.000000,0.000000}%
\pgfsetstrokecolor{currentstroke}%
\pgfsetdash{}{0pt}%
\pgfpathmoveto{\pgfqpoint{1.440360in}{1.001030in}}%
\pgfpathlineto{\pgfqpoint{1.458693in}{1.011885in}}%
\pgfpathlineto{\pgfqpoint{1.477110in}{1.022965in}}%
\pgfpathlineto{\pgfqpoint{1.485144in}{1.024076in}}%
\pgfpathlineto{\pgfqpoint{1.493152in}{1.025459in}}%
\pgfpathlineto{\pgfqpoint{1.474767in}{1.014219in}}%
\pgfpathlineto{\pgfqpoint{1.456465in}{1.003196in}}%
\pgfpathlineto{\pgfqpoint{1.448425in}{1.001973in}}%
\pgfpathlineto{\pgfqpoint{1.440360in}{1.001030in}}%
\pgfpathclose%
\pgfusepath{fill}%
\end{pgfscope}%
\begin{pgfscope}%
\pgfpathrectangle{\pgfqpoint{0.407592in}{0.222904in}}{\pgfqpoint{1.411163in}{1.411163in}}%
\pgfusepath{clip}%
\pgfsetbuttcap%
\pgfsetroundjoin%
\definecolor{currentfill}{rgb}{0.909804,0.698039,0.423529}%
\pgfsetfillcolor{currentfill}%
\pgfsetfillopacity{0.120000}%
\pgfsetlinewidth{0.000000pt}%
\definecolor{currentstroke}{rgb}{0.000000,0.000000,0.000000}%
\pgfsetstrokecolor{currentstroke}%
\pgfsetdash{}{0pt}%
\pgfpathmoveto{\pgfqpoint{1.546059in}{1.052395in}}%
\pgfpathlineto{\pgfqpoint{1.564668in}{1.064346in}}%
\pgfpathlineto{\pgfqpoint{1.572560in}{1.066725in}}%
\pgfpathlineto{\pgfqpoint{1.580430in}{1.069371in}}%
\pgfpathlineto{\pgfqpoint{1.561850in}{1.057321in}}%
\pgfpathlineto{\pgfqpoint{1.553966in}{1.054722in}}%
\pgfusepath{fill}%
\end{pgfscope}%
\begin{pgfscope}%
\pgfpathrectangle{\pgfqpoint{0.407592in}{0.222904in}}{\pgfqpoint{1.411163in}{1.411163in}}%
\pgfusepath{clip}%
\pgfsetbuttcap%
\pgfsetroundjoin%
\definecolor{currentfill}{rgb}{0.909804,0.698039,0.423529}%
\pgfsetfillcolor{currentfill}%
\pgfsetfillopacity{0.120000}%
\pgfsetlinewidth{0.000000pt}%
\definecolor{currentstroke}{rgb}{0.000000,0.000000,0.000000}%
\pgfsetstrokecolor{currentstroke}%
\pgfsetdash{}{0pt}%
\pgfpathmoveto{\pgfqpoint{1.493152in}{1.025459in}}%
\pgfpathlineto{\pgfqpoint{1.511622in}{1.036904in}}%
\pgfpathlineto{\pgfqpoint{1.530176in}{1.048547in}}%
\pgfpathlineto{\pgfqpoint{1.538129in}{1.050337in}}%
\pgfpathlineto{\pgfqpoint{1.546059in}{1.052395in}}%
\pgfpathlineto{\pgfqpoint{1.527535in}{1.040627in}}%
\pgfpathlineto{\pgfqpoint{1.509096in}{1.029049in}}%
\pgfpathlineto{\pgfqpoint{1.501136in}{1.027115in}}%
\pgfpathlineto{\pgfqpoint{1.493152in}{1.025459in}}%
\pgfpathclose%
\pgfusepath{fill}%
\end{pgfscope}%
\begin{pgfscope}%
\pgfpathrectangle{\pgfqpoint{0.407592in}{0.222904in}}{\pgfqpoint{1.411163in}{1.411163in}}%
\pgfusepath{clip}%
\pgfsetbuttcap%
\pgfsetroundjoin%
\definecolor{currentfill}{rgb}{0.909804,0.698039,0.423529}%
\pgfsetfillcolor{currentfill}%
\pgfsetfillopacity{0.120000}%
\pgfsetlinewidth{0.000000pt}%
\definecolor{currentstroke}{rgb}{0.000000,0.000000,0.000000}%
\pgfsetstrokecolor{currentstroke}%
\pgfsetdash{}{0pt}%
\pgfpathmoveto{\pgfqpoint{1.265692in}{0.948896in}}%
\pgfpathlineto{\pgfqpoint{1.283689in}{0.957125in}}%
\pgfpathlineto{\pgfqpoint{1.301761in}{0.965691in}}%
\pgfpathlineto{\pgfqpoint{1.310129in}{0.964072in}}%
\pgfpathlineto{\pgfqpoint{1.318465in}{0.962701in}}%
\pgfpathlineto{\pgfqpoint{1.300430in}{0.953874in}}%
\pgfpathlineto{\pgfqpoint{1.282471in}{0.945378in}}%
\pgfpathlineto{\pgfqpoint{1.274098in}{0.947011in}}%
\pgfpathlineto{\pgfqpoint{1.265692in}{0.948896in}}%
\pgfpathclose%
\pgfusepath{fill}%
\end{pgfscope}%
\begin{pgfscope}%
\pgfpathrectangle{\pgfqpoint{0.407592in}{0.222904in}}{\pgfqpoint{1.411163in}{1.411163in}}%
\pgfusepath{clip}%
\pgfsetbuttcap%
\pgfsetroundjoin%
\definecolor{currentfill}{rgb}{0.909804,0.698039,0.423529}%
\pgfsetfillcolor{currentfill}%
\pgfsetfillopacity{0.120000}%
\pgfsetlinewidth{0.000000pt}%
\definecolor{currentstroke}{rgb}{0.000000,0.000000,0.000000}%
\pgfsetstrokecolor{currentstroke}%
\pgfsetdash{}{0pt}%
\pgfpathmoveto{\pgfqpoint{1.318465in}{0.962701in}}%
\pgfpathlineto{\pgfqpoint{1.336577in}{0.971843in}}%
\pgfpathlineto{\pgfqpoint{1.354769in}{0.981283in}}%
\pgfpathlineto{\pgfqpoint{1.363037in}{0.980412in}}%
\pgfpathlineto{\pgfqpoint{1.371275in}{0.979799in}}%
\pgfpathlineto{\pgfqpoint{1.353119in}{0.970121in}}%
\pgfpathlineto{\pgfqpoint{1.335042in}{0.960734in}}%
\pgfpathlineto{\pgfqpoint{1.326769in}{0.961587in}}%
\pgfpathlineto{\pgfqpoint{1.318465in}{0.962701in}}%
\pgfpathclose%
\pgfusepath{fill}%
\end{pgfscope}%
\begin{pgfscope}%
\pgfpathrectangle{\pgfqpoint{0.407592in}{0.222904in}}{\pgfqpoint{1.411163in}{1.411163in}}%
\pgfusepath{clip}%
\pgfsetbuttcap%
\pgfsetroundjoin%
\definecolor{currentfill}{rgb}{0.909804,0.698039,0.423529}%
\pgfsetfillcolor{currentfill}%
\pgfsetfillopacity{0.120000}%
\pgfsetlinewidth{0.000000pt}%
\definecolor{currentstroke}{rgb}{0.000000,0.000000,0.000000}%
\pgfsetstrokecolor{currentstroke}%
\pgfsetdash{}{0pt}%
\pgfpathmoveto{\pgfqpoint{1.212930in}{0.938574in}}%
\pgfpathlineto{\pgfqpoint{1.230818in}{0.945795in}}%
\pgfpathlineto{\pgfqpoint{1.248779in}{0.953394in}}%
\pgfpathlineto{\pgfqpoint{1.257253in}{0.951026in}}%
\pgfpathlineto{\pgfqpoint{1.265692in}{0.948896in}}%
\pgfpathlineto{\pgfqpoint{1.247771in}{0.941021in}}%
\pgfpathlineto{\pgfqpoint{1.229923in}{0.933519in}}%
\pgfpathlineto{\pgfqpoint{1.221444in}{0.935927in}}%
\pgfpathlineto{\pgfqpoint{1.212930in}{0.938574in}}%
\pgfpathclose%
\pgfusepath{fill}%
\end{pgfscope}%
\begin{pgfscope}%
\pgfpathrectangle{\pgfqpoint{0.407592in}{0.222904in}}{\pgfqpoint{1.411163in}{1.411163in}}%
\pgfusepath{clip}%
\pgfsetbuttcap%
\pgfsetroundjoin%
\definecolor{currentfill}{rgb}{0.909804,0.698039,0.423529}%
\pgfsetfillcolor{currentfill}%
\pgfsetfillopacity{0.120000}%
\pgfsetlinewidth{0.000000pt}%
\definecolor{currentstroke}{rgb}{0.000000,0.000000,0.000000}%
\pgfsetstrokecolor{currentstroke}%
\pgfsetdash{}{0pt}%
\pgfpathmoveto{\pgfqpoint{1.371275in}{0.979799in}}%
\pgfpathlineto{\pgfqpoint{1.389512in}{0.989754in}}%
\pgfpathlineto{\pgfqpoint{1.407829in}{0.999971in}}%
\pgfpathlineto{\pgfqpoint{1.416004in}{0.999838in}}%
\pgfpathlineto{\pgfqpoint{1.424150in}{0.999967in}}%
\pgfpathlineto{\pgfqpoint{1.405866in}{0.989540in}}%
\pgfpathlineto{\pgfqpoint{1.387663in}{0.979368in}}%
\pgfpathlineto{\pgfqpoint{1.379484in}{0.979449in}}%
\pgfpathlineto{\pgfqpoint{1.371275in}{0.979799in}}%
\pgfpathclose%
\pgfusepath{fill}%
\end{pgfscope}%
\begin{pgfscope}%
\pgfpathrectangle{\pgfqpoint{0.407592in}{0.222904in}}{\pgfqpoint{1.411163in}{1.411163in}}%
\pgfusepath{clip}%
\pgfsetbuttcap%
\pgfsetroundjoin%
\definecolor{currentfill}{rgb}{0.909804,0.698039,0.423529}%
\pgfsetfillcolor{currentfill}%
\pgfsetfillopacity{0.120000}%
\pgfsetlinewidth{0.000000pt}%
\definecolor{currentstroke}{rgb}{0.000000,0.000000,0.000000}%
\pgfsetstrokecolor{currentstroke}%
\pgfsetdash{}{0pt}%
\pgfpathmoveto{\pgfqpoint{1.160148in}{0.931878in}}%
\pgfpathlineto{\pgfqpoint{1.177937in}{0.938006in}}%
\pgfpathlineto{\pgfqpoint{1.195795in}{0.944550in}}%
\pgfpathlineto{\pgfqpoint{1.204381in}{0.941450in}}%
\pgfpathlineto{\pgfqpoint{1.212930in}{0.938574in}}%
\pgfpathlineto{\pgfqpoint{1.195114in}{0.931747in}}%
\pgfpathlineto{\pgfqpoint{1.177369in}{0.925336in}}%
\pgfpathlineto{\pgfqpoint{1.168777in}{0.928496in}}%
\pgfpathlineto{\pgfqpoint{1.160148in}{0.931878in}}%
\pgfpathclose%
\pgfusepath{fill}%
\end{pgfscope}%
\begin{pgfscope}%
\pgfpathrectangle{\pgfqpoint{0.407592in}{0.222904in}}{\pgfqpoint{1.411163in}{1.411163in}}%
\pgfusepath{clip}%
\pgfsetbuttcap%
\pgfsetroundjoin%
\definecolor{currentfill}{rgb}{0.909804,0.698039,0.423529}%
\pgfsetfillcolor{currentfill}%
\pgfsetfillopacity{0.120000}%
\pgfsetlinewidth{0.000000pt}%
\definecolor{currentstroke}{rgb}{0.000000,0.000000,0.000000}%
\pgfsetstrokecolor{currentstroke}%
\pgfsetdash{}{0pt}%
\pgfpathmoveto{\pgfqpoint{1.424150in}{0.999967in}}%
\pgfpathlineto{\pgfqpoint{1.442516in}{1.010637in}}%
\pgfpathlineto{\pgfqpoint{1.460965in}{1.021538in}}%
\pgfpathlineto{\pgfqpoint{1.469051in}{1.022120in}}%
\pgfpathlineto{\pgfqpoint{1.477110in}{1.022965in}}%
\pgfpathlineto{\pgfqpoint{1.458693in}{1.011885in}}%
\pgfpathlineto{\pgfqpoint{1.440360in}{1.001030in}}%
\pgfpathlineto{\pgfqpoint{1.432268in}{1.000363in}}%
\pgfpathlineto{\pgfqpoint{1.424150in}{0.999967in}}%
\pgfpathclose%
\pgfusepath{fill}%
\end{pgfscope}%
\begin{pgfscope}%
\pgfpathrectangle{\pgfqpoint{0.407592in}{0.222904in}}{\pgfqpoint{1.411163in}{1.411163in}}%
\pgfusepath{clip}%
\pgfsetbuttcap%
\pgfsetroundjoin%
\definecolor{currentfill}{rgb}{0.909804,0.698039,0.423529}%
\pgfsetfillcolor{currentfill}%
\pgfsetfillopacity{0.120000}%
\pgfsetlinewidth{0.000000pt}%
\definecolor{currentstroke}{rgb}{0.000000,0.000000,0.000000}%
\pgfsetstrokecolor{currentstroke}%
\pgfsetdash{}{0pt}%
\pgfpathmoveto{\pgfqpoint{1.530176in}{1.048547in}}%
\pgfpathlineto{\pgfqpoint{1.548815in}{1.060377in}}%
\pgfpathlineto{\pgfqpoint{1.556753in}{1.062231in}}%
\pgfpathlineto{\pgfqpoint{1.564668in}{1.064346in}}%
\pgfpathlineto{\pgfqpoint{1.546059in}{1.052395in}}%
\pgfpathlineto{\pgfqpoint{1.538129in}{1.050337in}}%
\pgfusepath{fill}%
\end{pgfscope}%
\begin{pgfscope}%
\pgfpathrectangle{\pgfqpoint{0.407592in}{0.222904in}}{\pgfqpoint{1.411163in}{1.411163in}}%
\pgfusepath{clip}%
\pgfsetbuttcap%
\pgfsetroundjoin%
\definecolor{currentfill}{rgb}{0.909804,0.698039,0.423529}%
\pgfsetfillcolor{currentfill}%
\pgfsetfillopacity{0.120000}%
\pgfsetlinewidth{0.000000pt}%
\definecolor{currentstroke}{rgb}{0.000000,0.000000,0.000000}%
\pgfsetstrokecolor{currentstroke}%
\pgfsetdash{}{0pt}%
\pgfpathmoveto{\pgfqpoint{1.477110in}{1.022965in}}%
\pgfpathlineto{\pgfqpoint{1.495611in}{1.034258in}}%
\pgfpathlineto{\pgfqpoint{1.514196in}{1.045754in}}%
\pgfpathlineto{\pgfqpoint{1.522198in}{1.047020in}}%
\pgfpathlineto{\pgfqpoint{1.530176in}{1.048547in}}%
\pgfpathlineto{\pgfqpoint{1.511622in}{1.036904in}}%
\pgfpathlineto{\pgfqpoint{1.493152in}{1.025459in}}%
\pgfpathlineto{\pgfqpoint{1.485144in}{1.024076in}}%
\pgfpathlineto{\pgfqpoint{1.477110in}{1.022965in}}%
\pgfpathclose%
\pgfusepath{fill}%
\end{pgfscope}%
\begin{pgfscope}%
\pgfpathrectangle{\pgfqpoint{0.407592in}{0.222904in}}{\pgfqpoint{1.411163in}{1.411163in}}%
\pgfusepath{clip}%
\pgfsetbuttcap%
\pgfsetroundjoin%
\definecolor{currentfill}{rgb}{0.909804,0.698039,0.423529}%
\pgfsetfillcolor{currentfill}%
\pgfsetfillopacity{0.120000}%
\pgfsetlinewidth{0.000000pt}%
\definecolor{currentstroke}{rgb}{0.000000,0.000000,0.000000}%
\pgfsetstrokecolor{currentstroke}%
\pgfsetdash{}{0pt}%
\pgfpathmoveto{\pgfqpoint{1.301761in}{0.965691in}}%
\pgfpathlineto{\pgfqpoint{1.319911in}{0.974578in}}%
\pgfpathlineto{\pgfqpoint{1.338140in}{0.983771in}}%
\pgfpathlineto{\pgfqpoint{1.346470in}{0.982404in}}%
\pgfpathlineto{\pgfqpoint{1.354769in}{0.981283in}}%
\pgfpathlineto{\pgfqpoint{1.336577in}{0.971843in}}%
\pgfpathlineto{\pgfqpoint{1.318465in}{0.962701in}}%
\pgfpathlineto{\pgfqpoint{1.310129in}{0.964072in}}%
\pgfpathlineto{\pgfqpoint{1.301761in}{0.965691in}}%
\pgfpathclose%
\pgfusepath{fill}%
\end{pgfscope}%
\begin{pgfscope}%
\pgfpathrectangle{\pgfqpoint{0.407592in}{0.222904in}}{\pgfqpoint{1.411163in}{1.411163in}}%
\pgfusepath{clip}%
\pgfsetbuttcap%
\pgfsetroundjoin%
\definecolor{currentfill}{rgb}{0.909804,0.698039,0.423529}%
\pgfsetfillcolor{currentfill}%
\pgfsetfillopacity{0.120000}%
\pgfsetlinewidth{0.000000pt}%
\definecolor{currentstroke}{rgb}{0.000000,0.000000,0.000000}%
\pgfsetstrokecolor{currentstroke}%
\pgfsetdash{}{0pt}%
\pgfpathmoveto{\pgfqpoint{1.248779in}{0.953394in}}%
\pgfpathlineto{\pgfqpoint{1.266815in}{0.961351in}}%
\pgfpathlineto{\pgfqpoint{1.284926in}{0.969651in}}%
\pgfpathlineto{\pgfqpoint{1.293360in}{0.967553in}}%
\pgfpathlineto{\pgfqpoint{1.301761in}{0.965691in}}%
\pgfpathlineto{\pgfqpoint{1.283689in}{0.957125in}}%
\pgfpathlineto{\pgfqpoint{1.265692in}{0.948896in}}%
\pgfpathlineto{\pgfqpoint{1.257253in}{0.951026in}}%
\pgfpathlineto{\pgfqpoint{1.248779in}{0.953394in}}%
\pgfpathclose%
\pgfusepath{fill}%
\end{pgfscope}%
\begin{pgfscope}%
\pgfpathrectangle{\pgfqpoint{0.407592in}{0.222904in}}{\pgfqpoint{1.411163in}{1.411163in}}%
\pgfusepath{clip}%
\pgfsetbuttcap%
\pgfsetroundjoin%
\definecolor{currentfill}{rgb}{0.909804,0.698039,0.423529}%
\pgfsetfillcolor{currentfill}%
\pgfsetfillopacity{0.120000}%
\pgfsetlinewidth{0.000000pt}%
\definecolor{currentstroke}{rgb}{0.000000,0.000000,0.000000}%
\pgfsetstrokecolor{currentstroke}%
\pgfsetdash{}{0pt}%
\pgfpathmoveto{\pgfqpoint{1.354769in}{0.981283in}}%
\pgfpathlineto{\pgfqpoint{1.373041in}{0.991007in}}%
\pgfpathlineto{\pgfqpoint{1.391393in}{1.001002in}}%
\pgfpathlineto{\pgfqpoint{1.399626in}{1.000361in}}%
\pgfpathlineto{\pgfqpoint{1.407829in}{0.999971in}}%
\pgfpathlineto{\pgfqpoint{1.389512in}{0.989754in}}%
\pgfpathlineto{\pgfqpoint{1.371275in}{0.979799in}}%
\pgfpathlineto{\pgfqpoint{1.363037in}{0.980412in}}%
\pgfpathlineto{\pgfqpoint{1.354769in}{0.981283in}}%
\pgfpathclose%
\pgfusepath{fill}%
\end{pgfscope}%
\begin{pgfscope}%
\pgfpathrectangle{\pgfqpoint{0.407592in}{0.222904in}}{\pgfqpoint{1.411163in}{1.411163in}}%
\pgfusepath{clip}%
\pgfsetbuttcap%
\pgfsetroundjoin%
\definecolor{currentfill}{rgb}{0.909804,0.698039,0.423529}%
\pgfsetfillcolor{currentfill}%
\pgfsetfillopacity{0.120000}%
\pgfsetlinewidth{0.000000pt}%
\definecolor{currentstroke}{rgb}{0.000000,0.000000,0.000000}%
\pgfsetstrokecolor{currentstroke}%
\pgfsetdash{}{0pt}%
\pgfpathmoveto{\pgfqpoint{1.195795in}{0.944550in}}%
\pgfpathlineto{\pgfqpoint{1.213725in}{0.951491in}}%
\pgfpathlineto{\pgfqpoint{1.231727in}{0.958812in}}%
\pgfpathlineto{\pgfqpoint{1.240271in}{0.955991in}}%
\pgfpathlineto{\pgfqpoint{1.248779in}{0.953394in}}%
\pgfpathlineto{\pgfqpoint{1.230818in}{0.945795in}}%
\pgfpathlineto{\pgfqpoint{1.212930in}{0.938574in}}%
\pgfpathlineto{\pgfqpoint{1.204381in}{0.941450in}}%
\pgfpathlineto{\pgfqpoint{1.195795in}{0.944550in}}%
\pgfpathclose%
\pgfusepath{fill}%
\end{pgfscope}%
\begin{pgfscope}%
\pgfpathrectangle{\pgfqpoint{0.407592in}{0.222904in}}{\pgfqpoint{1.411163in}{1.411163in}}%
\pgfusepath{clip}%
\pgfsetbuttcap%
\pgfsetroundjoin%
\definecolor{currentfill}{rgb}{0.909804,0.698039,0.423529}%
\pgfsetfillcolor{currentfill}%
\pgfsetfillopacity{0.120000}%
\pgfsetlinewidth{0.000000pt}%
\definecolor{currentstroke}{rgb}{0.000000,0.000000,0.000000}%
\pgfsetstrokecolor{currentstroke}%
\pgfsetdash{}{0pt}%
\pgfpathmoveto{\pgfqpoint{1.514196in}{1.045754in}}%
\pgfpathlineto{\pgfqpoint{1.532867in}{1.057446in}}%
\pgfpathlineto{\pgfqpoint{1.540853in}{1.058783in}}%
\pgfpathlineto{\pgfqpoint{1.548815in}{1.060377in}}%
\pgfpathlineto{\pgfqpoint{1.530176in}{1.048547in}}%
\pgfpathlineto{\pgfqpoint{1.522198in}{1.047020in}}%
\pgfusepath{fill}%
\end{pgfscope}%
\begin{pgfscope}%
\pgfpathrectangle{\pgfqpoint{0.407592in}{0.222904in}}{\pgfqpoint{1.411163in}{1.411163in}}%
\pgfusepath{clip}%
\pgfsetbuttcap%
\pgfsetroundjoin%
\definecolor{currentfill}{rgb}{0.909804,0.698039,0.423529}%
\pgfsetfillcolor{currentfill}%
\pgfsetfillopacity{0.120000}%
\pgfsetlinewidth{0.000000pt}%
\definecolor{currentstroke}{rgb}{0.000000,0.000000,0.000000}%
\pgfsetstrokecolor{currentstroke}%
\pgfsetdash{}{0pt}%
\pgfpathmoveto{\pgfqpoint{1.407829in}{0.999971in}}%
\pgfpathlineto{\pgfqpoint{1.426229in}{1.010439in}}%
\pgfpathlineto{\pgfqpoint{1.444712in}{1.021146in}}%
\pgfpathlineto{\pgfqpoint{1.452852in}{1.021215in}}%
\pgfpathlineto{\pgfqpoint{1.460965in}{1.021538in}}%
\pgfpathlineto{\pgfqpoint{1.442516in}{1.010637in}}%
\pgfpathlineto{\pgfqpoint{1.424150in}{0.999967in}}%
\pgfpathlineto{\pgfqpoint{1.416004in}{0.999838in}}%
\pgfpathlineto{\pgfqpoint{1.407829in}{0.999971in}}%
\pgfpathclose%
\pgfusepath{fill}%
\end{pgfscope}%
\begin{pgfscope}%
\pgfpathrectangle{\pgfqpoint{0.407592in}{0.222904in}}{\pgfqpoint{1.411163in}{1.411163in}}%
\pgfusepath{clip}%
\pgfsetbuttcap%
\pgfsetroundjoin%
\definecolor{currentfill}{rgb}{0.909804,0.698039,0.423529}%
\pgfsetfillcolor{currentfill}%
\pgfsetfillopacity{0.120000}%
\pgfsetlinewidth{0.000000pt}%
\definecolor{currentstroke}{rgb}{0.000000,0.000000,0.000000}%
\pgfsetstrokecolor{currentstroke}%
\pgfsetdash{}{0pt}%
\pgfpathmoveto{\pgfqpoint{1.460965in}{1.021538in}}%
\pgfpathlineto{\pgfqpoint{1.479498in}{1.032660in}}%
\pgfpathlineto{\pgfqpoint{1.498116in}{1.043993in}}%
\pgfpathlineto{\pgfqpoint{1.506169in}{1.044746in}}%
\pgfpathlineto{\pgfqpoint{1.514196in}{1.045754in}}%
\pgfpathlineto{\pgfqpoint{1.495611in}{1.034258in}}%
\pgfpathlineto{\pgfqpoint{1.477110in}{1.022965in}}%
\pgfpathlineto{\pgfqpoint{1.469051in}{1.022120in}}%
\pgfpathlineto{\pgfqpoint{1.460965in}{1.021538in}}%
\pgfpathclose%
\pgfusepath{fill}%
\end{pgfscope}%
\begin{pgfscope}%
\pgfpathrectangle{\pgfqpoint{0.407592in}{0.222904in}}{\pgfqpoint{1.411163in}{1.411163in}}%
\pgfusepath{clip}%
\pgfsetbuttcap%
\pgfsetroundjoin%
\definecolor{currentfill}{rgb}{0.909804,0.698039,0.423529}%
\pgfsetfillcolor{currentfill}%
\pgfsetfillopacity{0.120000}%
\pgfsetlinewidth{0.000000pt}%
\definecolor{currentstroke}{rgb}{0.000000,0.000000,0.000000}%
\pgfsetstrokecolor{currentstroke}%
\pgfsetdash{}{0pt}%
\pgfpathmoveto{\pgfqpoint{1.284926in}{0.969651in}}%
\pgfpathlineto{\pgfqpoint{1.303115in}{0.978277in}}%
\pgfpathlineto{\pgfqpoint{1.321382in}{0.987216in}}%
\pgfpathlineto{\pgfqpoint{1.329777in}{0.985377in}}%
\pgfpathlineto{\pgfqpoint{1.338140in}{0.983771in}}%
\pgfpathlineto{\pgfqpoint{1.319911in}{0.974578in}}%
\pgfpathlineto{\pgfqpoint{1.301761in}{0.965691in}}%
\pgfpathlineto{\pgfqpoint{1.293360in}{0.967553in}}%
\pgfpathlineto{\pgfqpoint{1.284926in}{0.969651in}}%
\pgfpathclose%
\pgfusepath{fill}%
\end{pgfscope}%
\begin{pgfscope}%
\pgfpathrectangle{\pgfqpoint{0.407592in}{0.222904in}}{\pgfqpoint{1.411163in}{1.411163in}}%
\pgfusepath{clip}%
\pgfsetbuttcap%
\pgfsetroundjoin%
\definecolor{currentfill}{rgb}{0.909804,0.698039,0.423529}%
\pgfsetfillcolor{currentfill}%
\pgfsetfillopacity{0.120000}%
\pgfsetlinewidth{0.000000pt}%
\definecolor{currentstroke}{rgb}{0.000000,0.000000,0.000000}%
\pgfsetstrokecolor{currentstroke}%
\pgfsetdash{}{0pt}%
\pgfpathmoveto{\pgfqpoint{1.498116in}{1.043993in}}%
\pgfpathlineto{\pgfqpoint{1.516819in}{1.055528in}}%
\pgfpathlineto{\pgfqpoint{1.524856in}{1.056362in}}%
\pgfpathlineto{\pgfqpoint{1.532867in}{1.057446in}}%
\pgfpathlineto{\pgfqpoint{1.514196in}{1.045754in}}%
\pgfpathlineto{\pgfqpoint{1.506169in}{1.044746in}}%
\pgfusepath{fill}%
\end{pgfscope}%
\begin{pgfscope}%
\pgfpathrectangle{\pgfqpoint{0.407592in}{0.222904in}}{\pgfqpoint{1.411163in}{1.411163in}}%
\pgfusepath{clip}%
\pgfsetbuttcap%
\pgfsetroundjoin%
\definecolor{currentfill}{rgb}{0.909804,0.698039,0.423529}%
\pgfsetfillcolor{currentfill}%
\pgfsetfillopacity{0.120000}%
\pgfsetlinewidth{0.000000pt}%
\definecolor{currentstroke}{rgb}{0.000000,0.000000,0.000000}%
\pgfsetstrokecolor{currentstroke}%
\pgfsetdash{}{0pt}%
\pgfpathmoveto{\pgfqpoint{1.338140in}{0.983771in}}%
\pgfpathlineto{\pgfqpoint{1.356448in}{0.993255in}}%
\pgfpathlineto{\pgfqpoint{1.374837in}{1.003016in}}%
\pgfpathlineto{\pgfqpoint{1.383130in}{1.001888in}}%
\pgfpathlineto{\pgfqpoint{1.391393in}{1.001002in}}%
\pgfpathlineto{\pgfqpoint{1.373041in}{0.991007in}}%
\pgfpathlineto{\pgfqpoint{1.354769in}{0.981283in}}%
\pgfpathlineto{\pgfqpoint{1.346470in}{0.982404in}}%
\pgfpathlineto{\pgfqpoint{1.338140in}{0.983771in}}%
\pgfpathclose%
\pgfusepath{fill}%
\end{pgfscope}%
\begin{pgfscope}%
\pgfpathrectangle{\pgfqpoint{0.407592in}{0.222904in}}{\pgfqpoint{1.411163in}{1.411163in}}%
\pgfusepath{clip}%
\pgfsetbuttcap%
\pgfsetroundjoin%
\definecolor{currentfill}{rgb}{0.909804,0.698039,0.423529}%
\pgfsetfillcolor{currentfill}%
\pgfsetfillopacity{0.120000}%
\pgfsetlinewidth{0.000000pt}%
\definecolor{currentstroke}{rgb}{0.000000,0.000000,0.000000}%
\pgfsetstrokecolor{currentstroke}%
\pgfsetdash{}{0pt}%
\pgfpathmoveto{\pgfqpoint{1.231727in}{0.958812in}}%
\pgfpathlineto{\pgfqpoint{1.249804in}{0.966496in}}%
\pgfpathlineto{\pgfqpoint{1.267956in}{0.974527in}}%
\pgfpathlineto{\pgfqpoint{1.276459in}{0.971978in}}%
\pgfpathlineto{\pgfqpoint{1.284926in}{0.969651in}}%
\pgfpathlineto{\pgfqpoint{1.266815in}{0.961351in}}%
\pgfpathlineto{\pgfqpoint{1.248779in}{0.953394in}}%
\pgfpathlineto{\pgfqpoint{1.240271in}{0.955991in}}%
\pgfpathlineto{\pgfqpoint{1.231727in}{0.958812in}}%
\pgfpathclose%
\pgfusepath{fill}%
\end{pgfscope}%
\begin{pgfscope}%
\pgfpathrectangle{\pgfqpoint{0.407592in}{0.222904in}}{\pgfqpoint{1.411163in}{1.411163in}}%
\pgfusepath{clip}%
\pgfsetbuttcap%
\pgfsetroundjoin%
\definecolor{currentfill}{rgb}{0.909804,0.698039,0.423529}%
\pgfsetfillcolor{currentfill}%
\pgfsetfillopacity{0.120000}%
\pgfsetlinewidth{0.000000pt}%
\definecolor{currentstroke}{rgb}{0.000000,0.000000,0.000000}%
\pgfsetstrokecolor{currentstroke}%
\pgfsetdash{}{0pt}%
\pgfpathmoveto{\pgfqpoint{1.391393in}{1.001002in}}%
\pgfpathlineto{\pgfqpoint{1.409828in}{1.011254in}}%
\pgfpathlineto{\pgfqpoint{1.428345in}{1.021754in}}%
\pgfpathlineto{\pgfqpoint{1.436543in}{1.021327in}}%
\pgfpathlineto{\pgfqpoint{1.444712in}{1.021146in}}%
\pgfpathlineto{\pgfqpoint{1.426229in}{1.010439in}}%
\pgfpathlineto{\pgfqpoint{1.407829in}{0.999971in}}%
\pgfpathlineto{\pgfqpoint{1.399626in}{1.000361in}}%
\pgfpathlineto{\pgfqpoint{1.391393in}{1.001002in}}%
\pgfpathclose%
\pgfusepath{fill}%
\end{pgfscope}%
\begin{pgfscope}%
\pgfpathrectangle{\pgfqpoint{0.407592in}{0.222904in}}{\pgfqpoint{1.411163in}{1.411163in}}%
\pgfusepath{clip}%
\pgfsetbuttcap%
\pgfsetroundjoin%
\definecolor{currentfill}{rgb}{0.909804,0.698039,0.423529}%
\pgfsetfillcolor{currentfill}%
\pgfsetfillopacity{0.120000}%
\pgfsetlinewidth{0.000000pt}%
\definecolor{currentstroke}{rgb}{0.000000,0.000000,0.000000}%
\pgfsetstrokecolor{currentstroke}%
\pgfsetdash{}{0pt}%
\pgfpathmoveto{\pgfqpoint{1.444712in}{1.021146in}}%
\pgfpathlineto{\pgfqpoint{1.463278in}{1.032081in}}%
\pgfpathlineto{\pgfqpoint{1.481930in}{1.043235in}}%
\pgfpathlineto{\pgfqpoint{1.490036in}{1.043490in}}%
\pgfpathlineto{\pgfqpoint{1.498116in}{1.043993in}}%
\pgfpathlineto{\pgfqpoint{1.479498in}{1.032660in}}%
\pgfpathlineto{\pgfqpoint{1.460965in}{1.021538in}}%
\pgfpathlineto{\pgfqpoint{1.452852in}{1.021215in}}%
\pgfpathlineto{\pgfqpoint{1.444712in}{1.021146in}}%
\pgfpathclose%
\pgfusepath{fill}%
\end{pgfscope}%
\begin{pgfscope}%
\pgfpathrectangle{\pgfqpoint{0.407592in}{0.222904in}}{\pgfqpoint{1.411163in}{1.411163in}}%
\pgfusepath{clip}%
\pgfsetbuttcap%
\pgfsetroundjoin%
\definecolor{currentfill}{rgb}{0.909804,0.698039,0.423529}%
\pgfsetfillcolor{currentfill}%
\pgfsetfillopacity{0.120000}%
\pgfsetlinewidth{0.000000pt}%
\definecolor{currentstroke}{rgb}{0.000000,0.000000,0.000000}%
\pgfsetstrokecolor{currentstroke}%
\pgfsetdash{}{0pt}%
\pgfpathmoveto{\pgfqpoint{1.481930in}{1.043235in}}%
\pgfpathlineto{\pgfqpoint{1.500667in}{1.054598in}}%
\pgfpathlineto{\pgfqpoint{1.508756in}{1.054941in}}%
\pgfpathlineto{\pgfqpoint{1.516819in}{1.055528in}}%
\pgfpathlineto{\pgfqpoint{1.498116in}{1.043993in}}%
\pgfpathlineto{\pgfqpoint{1.490036in}{1.043490in}}%
\pgfusepath{fill}%
\end{pgfscope}%
\begin{pgfscope}%
\pgfpathrectangle{\pgfqpoint{0.407592in}{0.222904in}}{\pgfqpoint{1.411163in}{1.411163in}}%
\pgfusepath{clip}%
\pgfsetbuttcap%
\pgfsetroundjoin%
\definecolor{currentfill}{rgb}{0.909804,0.698039,0.423529}%
\pgfsetfillcolor{currentfill}%
\pgfsetfillopacity{0.120000}%
\pgfsetlinewidth{0.000000pt}%
\definecolor{currentstroke}{rgb}{0.000000,0.000000,0.000000}%
\pgfsetstrokecolor{currentstroke}%
\pgfsetdash{}{0pt}%
\pgfpathmoveto{\pgfqpoint{1.321382in}{0.987216in}}%
\pgfpathlineto{\pgfqpoint{1.339728in}{0.996452in}}%
\pgfpathlineto{\pgfqpoint{1.358155in}{1.005973in}}%
\pgfpathlineto{\pgfqpoint{1.366512in}{1.004380in}}%
\pgfpathlineto{\pgfqpoint{1.374837in}{1.003016in}}%
\pgfpathlineto{\pgfqpoint{1.356448in}{0.993255in}}%
\pgfpathlineto{\pgfqpoint{1.338140in}{0.983771in}}%
\pgfpathlineto{\pgfqpoint{1.329777in}{0.985377in}}%
\pgfpathlineto{\pgfqpoint{1.321382in}{0.987216in}}%
\pgfpathclose%
\pgfusepath{fill}%
\end{pgfscope}%
\begin{pgfscope}%
\pgfpathrectangle{\pgfqpoint{0.407592in}{0.222904in}}{\pgfqpoint{1.411163in}{1.411163in}}%
\pgfusepath{clip}%
\pgfsetbuttcap%
\pgfsetroundjoin%
\definecolor{currentfill}{rgb}{0.909804,0.698039,0.423529}%
\pgfsetfillcolor{currentfill}%
\pgfsetfillopacity{0.120000}%
\pgfsetlinewidth{0.000000pt}%
\definecolor{currentstroke}{rgb}{0.000000,0.000000,0.000000}%
\pgfsetstrokecolor{currentstroke}%
\pgfsetdash{}{0pt}%
\pgfpathmoveto{\pgfqpoint{1.267956in}{0.974527in}}%
\pgfpathlineto{\pgfqpoint{1.286185in}{0.982890in}}%
\pgfpathlineto{\pgfqpoint{1.304492in}{0.991571in}}%
\pgfpathlineto{\pgfqpoint{1.312954in}{0.989282in}}%
\pgfpathlineto{\pgfqpoint{1.321382in}{0.987216in}}%
\pgfpathlineto{\pgfqpoint{1.303115in}{0.978277in}}%
\pgfpathlineto{\pgfqpoint{1.284926in}{0.969651in}}%
\pgfpathlineto{\pgfqpoint{1.276459in}{0.971978in}}%
\pgfpathlineto{\pgfqpoint{1.267956in}{0.974527in}}%
\pgfpathclose%
\pgfusepath{fill}%
\end{pgfscope}%
\begin{pgfscope}%
\pgfpathrectangle{\pgfqpoint{0.407592in}{0.222904in}}{\pgfqpoint{1.411163in}{1.411163in}}%
\pgfusepath{clip}%
\pgfsetbuttcap%
\pgfsetroundjoin%
\definecolor{currentfill}{rgb}{0.909804,0.698039,0.423529}%
\pgfsetfillcolor{currentfill}%
\pgfsetfillopacity{0.120000}%
\pgfsetlinewidth{0.000000pt}%
\definecolor{currentstroke}{rgb}{0.000000,0.000000,0.000000}%
\pgfsetstrokecolor{currentstroke}%
\pgfsetdash{}{0pt}%
\pgfpathmoveto{\pgfqpoint{1.374837in}{1.003016in}}%
\pgfpathlineto{\pgfqpoint{1.393308in}{1.013044in}}%
\pgfpathlineto{\pgfqpoint{1.411862in}{1.023325in}}%
\pgfpathlineto{\pgfqpoint{1.420118in}{1.022421in}}%
\pgfpathlineto{\pgfqpoint{1.428345in}{1.021754in}}%
\pgfpathlineto{\pgfqpoint{1.409828in}{1.011254in}}%
\pgfpathlineto{\pgfqpoint{1.391393in}{1.001002in}}%
\pgfpathlineto{\pgfqpoint{1.383130in}{1.001888in}}%
\pgfpathlineto{\pgfqpoint{1.374837in}{1.003016in}}%
\pgfpathclose%
\pgfusepath{fill}%
\end{pgfscope}%
\begin{pgfscope}%
\pgfpathrectangle{\pgfqpoint{0.407592in}{0.222904in}}{\pgfqpoint{1.411163in}{1.411163in}}%
\pgfusepath{clip}%
\pgfsetbuttcap%
\pgfsetroundjoin%
\definecolor{currentfill}{rgb}{0.909804,0.698039,0.423529}%
\pgfsetfillcolor{currentfill}%
\pgfsetfillopacity{0.120000}%
\pgfsetlinewidth{0.000000pt}%
\definecolor{currentstroke}{rgb}{0.000000,0.000000,0.000000}%
\pgfsetstrokecolor{currentstroke}%
\pgfsetdash{}{0pt}%
\pgfpathmoveto{\pgfqpoint{1.428345in}{1.021754in}}%
\pgfpathlineto{\pgfqpoint{1.446947in}{1.032489in}}%
\pgfpathlineto{\pgfqpoint{1.465633in}{1.043450in}}%
\pgfpathlineto{\pgfqpoint{1.473795in}{1.043223in}}%
\pgfpathlineto{\pgfqpoint{1.481930in}{1.043235in}}%
\pgfpathlineto{\pgfqpoint{1.463278in}{1.032081in}}%
\pgfpathlineto{\pgfqpoint{1.444712in}{1.021146in}}%
\pgfpathlineto{\pgfqpoint{1.436543in}{1.021327in}}%
\pgfpathlineto{\pgfqpoint{1.428345in}{1.021754in}}%
\pgfpathclose%
\pgfusepath{fill}%
\end{pgfscope}%
\begin{pgfscope}%
\pgfpathrectangle{\pgfqpoint{0.407592in}{0.222904in}}{\pgfqpoint{1.411163in}{1.411163in}}%
\pgfusepath{clip}%
\pgfsetbuttcap%
\pgfsetroundjoin%
\definecolor{currentfill}{rgb}{0.909804,0.698039,0.423529}%
\pgfsetfillcolor{currentfill}%
\pgfsetfillopacity{0.120000}%
\pgfsetlinewidth{0.000000pt}%
\definecolor{currentstroke}{rgb}{0.000000,0.000000,0.000000}%
\pgfsetstrokecolor{currentstroke}%
\pgfsetdash{}{0pt}%
\pgfpathmoveto{\pgfqpoint{1.465633in}{1.043450in}}%
\pgfpathlineto{\pgfqpoint{1.484405in}{1.054628in}}%
\pgfpathlineto{\pgfqpoint{1.492550in}{1.054495in}}%
\pgfpathlineto{\pgfqpoint{1.500667in}{1.054598in}}%
\pgfpathlineto{\pgfqpoint{1.481930in}{1.043235in}}%
\pgfpathlineto{\pgfqpoint{1.473795in}{1.043223in}}%
\pgfusepath{fill}%
\end{pgfscope}%
\begin{pgfscope}%
\pgfpathrectangle{\pgfqpoint{0.407592in}{0.222904in}}{\pgfqpoint{1.411163in}{1.411163in}}%
\pgfusepath{clip}%
\pgfsetbuttcap%
\pgfsetroundjoin%
\definecolor{currentfill}{rgb}{0.909804,0.698039,0.423529}%
\pgfsetfillcolor{currentfill}%
\pgfsetfillopacity{0.120000}%
\pgfsetlinewidth{0.000000pt}%
\definecolor{currentstroke}{rgb}{0.000000,0.000000,0.000000}%
\pgfsetstrokecolor{currentstroke}%
\pgfsetdash{}{0pt}%
\pgfpathmoveto{\pgfqpoint{1.304492in}{0.991571in}}%
\pgfpathlineto{\pgfqpoint{1.322878in}{1.000555in}}%
\pgfpathlineto{\pgfqpoint{1.341344in}{1.009830in}}%
\pgfpathlineto{\pgfqpoint{1.349766in}{1.007792in}}%
\pgfpathlineto{\pgfqpoint{1.358155in}{1.005973in}}%
\pgfpathlineto{\pgfqpoint{1.339728in}{0.996452in}}%
\pgfpathlineto{\pgfqpoint{1.321382in}{0.987216in}}%
\pgfpathlineto{\pgfqpoint{1.312954in}{0.989282in}}%
\pgfpathlineto{\pgfqpoint{1.304492in}{0.991571in}}%
\pgfpathclose%
\pgfusepath{fill}%
\end{pgfscope}%
\begin{pgfscope}%
\pgfpathrectangle{\pgfqpoint{0.407592in}{0.222904in}}{\pgfqpoint{1.411163in}{1.411163in}}%
\pgfusepath{clip}%
\pgfsetbuttcap%
\pgfsetroundjoin%
\definecolor{currentfill}{rgb}{0.909804,0.698039,0.423529}%
\pgfsetfillcolor{currentfill}%
\pgfsetfillopacity{0.120000}%
\pgfsetlinewidth{0.000000pt}%
\definecolor{currentstroke}{rgb}{0.000000,0.000000,0.000000}%
\pgfsetstrokecolor{currentstroke}%
\pgfsetdash{}{0pt}%
\pgfpathmoveto{\pgfqpoint{1.358155in}{1.005973in}}%
\pgfpathlineto{\pgfqpoint{1.376664in}{1.015767in}}%
\pgfpathlineto{\pgfqpoint{1.395255in}{1.025821in}}%
\pgfpathlineto{\pgfqpoint{1.403574in}{1.024460in}}%
\pgfpathlineto{\pgfqpoint{1.411862in}{1.023325in}}%
\pgfpathlineto{\pgfqpoint{1.393308in}{1.013044in}}%
\pgfpathlineto{\pgfqpoint{1.374837in}{1.003016in}}%
\pgfpathlineto{\pgfqpoint{1.366512in}{1.004380in}}%
\pgfpathlineto{\pgfqpoint{1.358155in}{1.005973in}}%
\pgfpathclose%
\pgfusepath{fill}%
\end{pgfscope}%
\begin{pgfscope}%
\pgfpathrectangle{\pgfqpoint{0.407592in}{0.222904in}}{\pgfqpoint{1.411163in}{1.411163in}}%
\pgfusepath{clip}%
\pgfsetbuttcap%
\pgfsetroundjoin%
\definecolor{currentfill}{rgb}{0.909804,0.698039,0.423529}%
\pgfsetfillcolor{currentfill}%
\pgfsetfillopacity{0.120000}%
\pgfsetlinewidth{0.000000pt}%
\definecolor{currentstroke}{rgb}{0.000000,0.000000,0.000000}%
\pgfsetstrokecolor{currentstroke}%
\pgfsetdash{}{0pt}%
\pgfpathmoveto{\pgfqpoint{1.411862in}{1.023325in}}%
\pgfpathlineto{\pgfqpoint{1.430499in}{1.033849in}}%
\pgfpathlineto{\pgfqpoint{1.449221in}{1.044607in}}%
\pgfpathlineto{\pgfqpoint{1.457442in}{1.043913in}}%
\pgfpathlineto{\pgfqpoint{1.465633in}{1.043450in}}%
\pgfpathlineto{\pgfqpoint{1.446947in}{1.032489in}}%
\pgfpathlineto{\pgfqpoint{1.428345in}{1.021754in}}%
\pgfpathlineto{\pgfqpoint{1.420118in}{1.022421in}}%
\pgfpathlineto{\pgfqpoint{1.411862in}{1.023325in}}%
\pgfpathclose%
\pgfusepath{fill}%
\end{pgfscope}%
\begin{pgfscope}%
\pgfpathrectangle{\pgfqpoint{0.407592in}{0.222904in}}{\pgfqpoint{1.411163in}{1.411163in}}%
\pgfusepath{clip}%
\pgfsetbuttcap%
\pgfsetroundjoin%
\definecolor{currentfill}{rgb}{0.909804,0.698039,0.423529}%
\pgfsetfillcolor{currentfill}%
\pgfsetfillopacity{0.120000}%
\pgfsetlinewidth{0.000000pt}%
\definecolor{currentstroke}{rgb}{0.000000,0.000000,0.000000}%
\pgfsetstrokecolor{currentstroke}%
\pgfsetdash{}{0pt}%
\pgfpathmoveto{\pgfqpoint{1.449221in}{1.044607in}}%
\pgfpathlineto{\pgfqpoint{1.468029in}{1.055588in}}%
\pgfpathlineto{\pgfqpoint{1.476232in}{1.054994in}}%
\pgfpathlineto{\pgfqpoint{1.484405in}{1.054628in}}%
\pgfpathlineto{\pgfqpoint{1.465633in}{1.043450in}}%
\pgfpathlineto{\pgfqpoint{1.457442in}{1.043913in}}%
\pgfusepath{fill}%
\end{pgfscope}%
\begin{pgfscope}%
\pgfpathrectangle{\pgfqpoint{0.407592in}{0.222904in}}{\pgfqpoint{1.411163in}{1.411163in}}%
\pgfusepath{clip}%
\pgfsetbuttcap%
\pgfsetroundjoin%
\definecolor{currentfill}{rgb}{0.909804,0.698039,0.423529}%
\pgfsetfillcolor{currentfill}%
\pgfsetfillopacity{0.120000}%
\pgfsetlinewidth{0.000000pt}%
\definecolor{currentstroke}{rgb}{0.000000,0.000000,0.000000}%
\pgfsetstrokecolor{currentstroke}%
\pgfsetdash{}{0pt}%
\pgfpathmoveto{\pgfqpoint{1.341344in}{1.009830in}}%
\pgfpathlineto{\pgfqpoint{1.359891in}{1.019384in}}%
\pgfpathlineto{\pgfqpoint{1.378521in}{1.029206in}}%
\pgfpathlineto{\pgfqpoint{1.386904in}{1.027405in}}%
\pgfpathlineto{\pgfqpoint{1.395255in}{1.025821in}}%
\pgfpathlineto{\pgfqpoint{1.376664in}{1.015767in}}%
\pgfpathlineto{\pgfqpoint{1.358155in}{1.005973in}}%
\pgfpathlineto{\pgfqpoint{1.349766in}{1.007792in}}%
\pgfpathlineto{\pgfqpoint{1.341344in}{1.009830in}}%
\pgfpathclose%
\pgfusepath{fill}%
\end{pgfscope}%
\begin{pgfscope}%
\pgfpathrectangle{\pgfqpoint{0.407592in}{0.222904in}}{\pgfqpoint{1.411163in}{1.411163in}}%
\pgfusepath{clip}%
\pgfsetbuttcap%
\pgfsetroundjoin%
\definecolor{currentfill}{rgb}{0.909804,0.698039,0.423529}%
\pgfsetfillcolor{currentfill}%
\pgfsetfillopacity{0.120000}%
\pgfsetlinewidth{0.000000pt}%
\definecolor{currentstroke}{rgb}{0.000000,0.000000,0.000000}%
\pgfsetstrokecolor{currentstroke}%
\pgfsetdash{}{0pt}%
\pgfpathmoveto{\pgfqpoint{1.395255in}{1.025821in}}%
\pgfpathlineto{\pgfqpoint{1.413930in}{1.036126in}}%
\pgfpathlineto{\pgfqpoint{1.432689in}{1.046671in}}%
\pgfpathlineto{\pgfqpoint{1.440971in}{1.045528in}}%
\pgfpathlineto{\pgfqpoint{1.449221in}{1.044607in}}%
\pgfpathlineto{\pgfqpoint{1.430499in}{1.033849in}}%
\pgfpathlineto{\pgfqpoint{1.411862in}{1.023325in}}%
\pgfpathlineto{\pgfqpoint{1.403574in}{1.024460in}}%
\pgfpathlineto{\pgfqpoint{1.395255in}{1.025821in}}%
\pgfpathclose%
\pgfusepath{fill}%
\end{pgfscope}%
\begin{pgfscope}%
\pgfpathrectangle{\pgfqpoint{0.407592in}{0.222904in}}{\pgfqpoint{1.411163in}{1.411163in}}%
\pgfusepath{clip}%
\pgfsetbuttcap%
\pgfsetroundjoin%
\definecolor{currentfill}{rgb}{0.909804,0.698039,0.423529}%
\pgfsetfillcolor{currentfill}%
\pgfsetfillopacity{0.120000}%
\pgfsetlinewidth{0.000000pt}%
\definecolor{currentstroke}{rgb}{0.000000,0.000000,0.000000}%
\pgfsetstrokecolor{currentstroke}%
\pgfsetdash{}{0pt}%
\pgfpathmoveto{\pgfqpoint{1.432689in}{1.046671in}}%
\pgfpathlineto{\pgfqpoint{1.451535in}{1.057447in}}%
\pgfpathlineto{\pgfqpoint{1.459797in}{1.056407in}}%
\pgfpathlineto{\pgfqpoint{1.468029in}{1.055588in}}%
\pgfpathlineto{\pgfqpoint{1.449221in}{1.044607in}}%
\pgfpathlineto{\pgfqpoint{1.440971in}{1.045528in}}%
\pgfusepath{fill}%
\end{pgfscope}%
\begin{pgfscope}%
\pgfpathrectangle{\pgfqpoint{0.407592in}{0.222904in}}{\pgfqpoint{1.411163in}{1.411163in}}%
\pgfusepath{clip}%
\pgfsetbuttcap%
\pgfsetroundjoin%
\definecolor{currentfill}{rgb}{0.909804,0.698039,0.423529}%
\pgfsetfillcolor{currentfill}%
\pgfsetfillopacity{0.120000}%
\pgfsetlinewidth{0.000000pt}%
\definecolor{currentstroke}{rgb}{0.000000,0.000000,0.000000}%
\pgfsetstrokecolor{currentstroke}%
\pgfsetdash{}{0pt}%
\pgfpathmoveto{\pgfqpoint{1.378521in}{1.029206in}}%
\pgfpathlineto{\pgfqpoint{1.397235in}{1.039285in}}%
\pgfpathlineto{\pgfqpoint{1.416033in}{1.049611in}}%
\pgfpathlineto{\pgfqpoint{1.424377in}{1.048034in}}%
\pgfpathlineto{\pgfqpoint{1.432689in}{1.046671in}}%
\pgfpathlineto{\pgfqpoint{1.413930in}{1.036126in}}%
\pgfpathlineto{\pgfqpoint{1.395255in}{1.025821in}}%
\pgfpathlineto{\pgfqpoint{1.386904in}{1.027405in}}%
\pgfpathlineto{\pgfqpoint{1.378521in}{1.029206in}}%
\pgfpathclose%
\pgfusepath{fill}%
\end{pgfscope}%
\begin{pgfscope}%
\pgfpathrectangle{\pgfqpoint{0.407592in}{0.222904in}}{\pgfqpoint{1.411163in}{1.411163in}}%
\pgfusepath{clip}%
\pgfsetbuttcap%
\pgfsetroundjoin%
\definecolor{currentfill}{rgb}{0.909804,0.698039,0.423529}%
\pgfsetfillcolor{currentfill}%
\pgfsetfillopacity{0.120000}%
\pgfsetlinewidth{0.000000pt}%
\definecolor{currentstroke}{rgb}{0.000000,0.000000,0.000000}%
\pgfsetstrokecolor{currentstroke}%
\pgfsetdash{}{0pt}%
\pgfpathmoveto{\pgfqpoint{1.416033in}{1.049611in}}%
\pgfpathlineto{\pgfqpoint{1.434916in}{1.060174in}}%
\pgfpathlineto{\pgfqpoint{1.443241in}{1.058704in}}%
\pgfpathlineto{\pgfqpoint{1.451535in}{1.057447in}}%
\pgfpathlineto{\pgfqpoint{1.432689in}{1.046671in}}%
\pgfpathlineto{\pgfqpoint{1.424377in}{1.048034in}}%
\pgfusepath{fill}%
\end{pgfscope}%
\begin{pgfscope}%
\pgfpathrectangle{\pgfqpoint{0.407592in}{0.222904in}}{\pgfqpoint{1.411163in}{1.411163in}}%
\pgfusepath{clip}%
\pgfsetbuttcap%
\pgfsetroundjoin%
\definecolor{currentfill}{rgb}{0.749020,0.772549,0.792157}%
\pgfsetfillcolor{currentfill}%
\pgfsetfillopacity{0.160000}%
\pgfsetlinewidth{0.000000pt}%
\definecolor{currentstroke}{rgb}{0.000000,0.000000,0.000000}%
\pgfsetstrokecolor{currentstroke}%
\pgfsetdash{}{0pt}%
\pgfpathmoveto{\pgfqpoint{1.207276in}{1.516005in}}%
\pgfpathlineto{\pgfqpoint{1.209257in}{1.531330in}}%
\pgfpathlineto{\pgfqpoint{1.162616in}{1.532749in}}%
\pgfpathlineto{\pgfqpoint{1.115711in}{1.533155in}}%
\pgfpathlineto{\pgfqpoint{1.115287in}{1.517809in}}%
\pgfpathlineto{\pgfqpoint{1.161411in}{1.517408in}}%
\pgfusepath{fill}%
\end{pgfscope}%
\begin{pgfscope}%
\pgfpathrectangle{\pgfqpoint{0.407592in}{0.222904in}}{\pgfqpoint{1.411163in}{1.411163in}}%
\pgfusepath{clip}%
\pgfsetbuttcap%
\pgfsetroundjoin%
\definecolor{currentfill}{rgb}{0.749020,0.772549,0.792157}%
\pgfsetfillcolor{currentfill}%
\pgfsetfillopacity{0.160000}%
\pgfsetlinewidth{0.000000pt}%
\definecolor{currentstroke}{rgb}{0.000000,0.000000,0.000000}%
\pgfsetstrokecolor{currentstroke}%
\pgfsetdash{}{0pt}%
\pgfpathmoveto{\pgfqpoint{1.115287in}{1.517809in}}%
\pgfpathlineto{\pgfqpoint{1.115711in}{1.533155in}}%
\pgfpathlineto{\pgfqpoint{1.068695in}{1.532546in}}%
\pgfpathlineto{\pgfqpoint{1.021719in}{1.530924in}}%
\pgfpathlineto{\pgfqpoint{1.022860in}{1.515603in}}%
\pgfpathlineto{\pgfqpoint{1.069053in}{1.517207in}}%
\pgfusepath{fill}%
\end{pgfscope}%
\begin{pgfscope}%
\pgfpathrectangle{\pgfqpoint{0.407592in}{0.222904in}}{\pgfqpoint{1.411163in}{1.411163in}}%
\pgfusepath{clip}%
\pgfsetbuttcap%
\pgfsetroundjoin%
\definecolor{currentfill}{rgb}{0.749020,0.772549,0.792157}%
\pgfsetfillcolor{currentfill}%
\pgfsetfillopacity{0.160000}%
\pgfsetlinewidth{0.000000pt}%
\definecolor{currentstroke}{rgb}{0.000000,0.000000,0.000000}%
\pgfsetstrokecolor{currentstroke}%
\pgfsetdash{}{0pt}%
\pgfpathmoveto{\pgfqpoint{1.203303in}{1.485270in}}%
\pgfpathlineto{\pgfqpoint{1.205291in}{1.500652in}}%
\pgfpathlineto{\pgfqpoint{1.207276in}{1.516005in}}%
\pgfpathlineto{\pgfqpoint{1.161411in}{1.517408in}}%
\pgfpathlineto{\pgfqpoint{1.115287in}{1.517809in}}%
\pgfpathlineto{\pgfqpoint{1.114863in}{1.502435in}}%
\pgfpathlineto{\pgfqpoint{1.114437in}{1.487033in}}%
\pgfpathlineto{\pgfqpoint{1.158996in}{1.486641in}}%
\pgfpathlineto{\pgfqpoint{1.203303in}{1.485270in}}%
\pgfpathclose%
\pgfusepath{fill}%
\end{pgfscope}%
\begin{pgfscope}%
\pgfpathrectangle{\pgfqpoint{0.407592in}{0.222904in}}{\pgfqpoint{1.411163in}{1.411163in}}%
\pgfusepath{clip}%
\pgfsetbuttcap%
\pgfsetroundjoin%
\definecolor{currentfill}{rgb}{0.749020,0.772549,0.792157}%
\pgfsetfillcolor{currentfill}%
\pgfsetfillopacity{0.160000}%
\pgfsetlinewidth{0.000000pt}%
\definecolor{currentstroke}{rgb}{0.000000,0.000000,0.000000}%
\pgfsetstrokecolor{currentstroke}%
\pgfsetdash{}{0pt}%
\pgfpathmoveto{\pgfqpoint{1.297627in}{1.510206in}}%
\pgfpathlineto{\pgfqpoint{1.301145in}{1.525467in}}%
\pgfpathlineto{\pgfqpoint{1.255484in}{1.528901in}}%
\pgfpathlineto{\pgfqpoint{1.209257in}{1.531330in}}%
\pgfpathlineto{\pgfqpoint{1.207276in}{1.516005in}}%
\pgfpathlineto{\pgfqpoint{1.252731in}{1.513603in}}%
\pgfusepath{fill}%
\end{pgfscope}%
\begin{pgfscope}%
\pgfpathrectangle{\pgfqpoint{0.407592in}{0.222904in}}{\pgfqpoint{1.411163in}{1.411163in}}%
\pgfusepath{clip}%
\pgfsetbuttcap%
\pgfsetroundjoin%
\definecolor{currentfill}{rgb}{0.749020,0.772549,0.792157}%
\pgfsetfillcolor{currentfill}%
\pgfsetfillopacity{0.160000}%
\pgfsetlinewidth{0.000000pt}%
\definecolor{currentstroke}{rgb}{0.000000,0.000000,0.000000}%
\pgfsetstrokecolor{currentstroke}%
\pgfsetdash{}{0pt}%
\pgfpathmoveto{\pgfqpoint{1.114437in}{1.487033in}}%
\pgfpathlineto{\pgfqpoint{1.114863in}{1.502435in}}%
\pgfpathlineto{\pgfqpoint{1.115287in}{1.517809in}}%
\pgfpathlineto{\pgfqpoint{1.069053in}{1.517207in}}%
\pgfpathlineto{\pgfqpoint{1.022860in}{1.515603in}}%
\pgfpathlineto{\pgfqpoint{1.024002in}{1.500254in}}%
\pgfpathlineto{\pgfqpoint{1.025147in}{1.484877in}}%
\pgfpathlineto{\pgfqpoint{1.069772in}{1.486445in}}%
\pgfpathlineto{\pgfqpoint{1.114437in}{1.487033in}}%
\pgfpathclose%
\pgfusepath{fill}%
\end{pgfscope}%
\begin{pgfscope}%
\pgfpathrectangle{\pgfqpoint{0.407592in}{0.222904in}}{\pgfqpoint{1.411163in}{1.411163in}}%
\pgfusepath{clip}%
\pgfsetbuttcap%
\pgfsetroundjoin%
\definecolor{currentfill}{rgb}{0.749020,0.772549,0.792157}%
\pgfsetfillcolor{currentfill}%
\pgfsetfillopacity{0.160000}%
\pgfsetlinewidth{0.000000pt}%
\definecolor{currentstroke}{rgb}{0.000000,0.000000,0.000000}%
\pgfsetstrokecolor{currentstroke}%
\pgfsetdash{}{0pt}%
\pgfpathmoveto{\pgfqpoint{1.022860in}{1.515603in}}%
\pgfpathlineto{\pgfqpoint{1.021719in}{1.530924in}}%
\pgfpathlineto{\pgfqpoint{0.974935in}{1.528293in}}%
\pgfpathlineto{\pgfqpoint{0.928496in}{1.524657in}}%
\pgfpathlineto{\pgfqpoint{0.931195in}{1.509405in}}%
\pgfpathlineto{\pgfqpoint{0.976856in}{1.513001in}}%
\pgfusepath{fill}%
\end{pgfscope}%
\begin{pgfscope}%
\pgfpathrectangle{\pgfqpoint{0.407592in}{0.222904in}}{\pgfqpoint{1.411163in}{1.411163in}}%
\pgfusepath{clip}%
\pgfsetbuttcap%
\pgfsetroundjoin%
\definecolor{currentfill}{rgb}{0.749020,0.772549,0.792157}%
\pgfsetfillcolor{currentfill}%
\pgfsetfillopacity{0.160000}%
\pgfsetlinewidth{0.000000pt}%
\definecolor{currentstroke}{rgb}{0.000000,0.000000,0.000000}%
\pgfsetstrokecolor{currentstroke}%
\pgfsetdash{}{0pt}%
\pgfpathmoveto{\pgfqpoint{1.290572in}{1.479604in}}%
\pgfpathlineto{\pgfqpoint{1.294103in}{1.494918in}}%
\pgfpathlineto{\pgfqpoint{1.297627in}{1.510206in}}%
\pgfpathlineto{\pgfqpoint{1.252731in}{1.513603in}}%
\pgfpathlineto{\pgfqpoint{1.207276in}{1.516005in}}%
\pgfpathlineto{\pgfqpoint{1.205291in}{1.500652in}}%
\pgfpathlineto{\pgfqpoint{1.203303in}{1.485270in}}%
\pgfpathlineto{\pgfqpoint{1.247210in}{1.482922in}}%
\pgfpathlineto{\pgfqpoint{1.290572in}{1.479604in}}%
\pgfpathclose%
\pgfusepath{fill}%
\end{pgfscope}%
\begin{pgfscope}%
\pgfpathrectangle{\pgfqpoint{0.407592in}{0.222904in}}{\pgfqpoint{1.411163in}{1.411163in}}%
\pgfusepath{clip}%
\pgfsetbuttcap%
\pgfsetroundjoin%
\definecolor{currentfill}{rgb}{0.749020,0.772549,0.792157}%
\pgfsetfillcolor{currentfill}%
\pgfsetfillopacity{0.160000}%
\pgfsetlinewidth{0.000000pt}%
\definecolor{currentstroke}{rgb}{0.000000,0.000000,0.000000}%
\pgfsetstrokecolor{currentstroke}%
\pgfsetdash{}{0pt}%
\pgfpathmoveto{\pgfqpoint{1.025147in}{1.484877in}}%
\pgfpathlineto{\pgfqpoint{1.024002in}{1.500254in}}%
\pgfpathlineto{\pgfqpoint{1.022860in}{1.515603in}}%
\pgfpathlineto{\pgfqpoint{0.976856in}{1.513001in}}%
\pgfpathlineto{\pgfqpoint{0.931195in}{1.509405in}}%
\pgfpathlineto{\pgfqpoint{0.933898in}{1.494127in}}%
\pgfpathlineto{\pgfqpoint{0.936606in}{1.478821in}}%
\pgfpathlineto{\pgfqpoint{0.980709in}{1.482334in}}%
\pgfpathlineto{\pgfqpoint{1.025147in}{1.484877in}}%
\pgfpathclose%
\pgfusepath{fill}%
\end{pgfscope}%
\begin{pgfscope}%
\pgfpathrectangle{\pgfqpoint{0.407592in}{0.222904in}}{\pgfqpoint{1.411163in}{1.411163in}}%
\pgfusepath{clip}%
\pgfsetbuttcap%
\pgfsetroundjoin%
\definecolor{currentfill}{rgb}{0.749020,0.772549,0.792157}%
\pgfsetfillcolor{currentfill}%
\pgfsetfillopacity{0.160000}%
\pgfsetlinewidth{0.000000pt}%
\definecolor{currentstroke}{rgb}{0.000000,0.000000,0.000000}%
\pgfsetstrokecolor{currentstroke}%
\pgfsetdash{}{0pt}%
\pgfpathmoveto{\pgfqpoint{1.199315in}{1.454422in}}%
\pgfpathlineto{\pgfqpoint{1.201311in}{1.469860in}}%
\pgfpathlineto{\pgfqpoint{1.203303in}{1.485270in}}%
\pgfpathlineto{\pgfqpoint{1.158996in}{1.486641in}}%
\pgfpathlineto{\pgfqpoint{1.114437in}{1.487033in}}%
\pgfpathlineto{\pgfqpoint{1.114011in}{1.471602in}}%
\pgfpathlineto{\pgfqpoint{1.113584in}{1.456142in}}%
\pgfpathlineto{\pgfqpoint{1.156573in}{1.455760in}}%
\pgfpathlineto{\pgfqpoint{1.199315in}{1.454422in}}%
\pgfpathclose%
\pgfusepath{fill}%
\end{pgfscope}%
\begin{pgfscope}%
\pgfpathrectangle{\pgfqpoint{0.407592in}{0.222904in}}{\pgfqpoint{1.411163in}{1.411163in}}%
\pgfusepath{clip}%
\pgfsetbuttcap%
\pgfsetroundjoin%
\definecolor{currentfill}{rgb}{0.749020,0.772549,0.792157}%
\pgfsetfillcolor{currentfill}%
\pgfsetfillopacity{0.160000}%
\pgfsetlinewidth{0.000000pt}%
\definecolor{currentstroke}{rgb}{0.000000,0.000000,0.000000}%
\pgfsetstrokecolor{currentstroke}%
\pgfsetdash{}{0pt}%
\pgfpathmoveto{\pgfqpoint{1.113584in}{1.456142in}}%
\pgfpathlineto{\pgfqpoint{1.114011in}{1.471602in}}%
\pgfpathlineto{\pgfqpoint{1.114437in}{1.487033in}}%
\pgfpathlineto{\pgfqpoint{1.069772in}{1.486445in}}%
\pgfpathlineto{\pgfqpoint{1.025147in}{1.484877in}}%
\pgfpathlineto{\pgfqpoint{1.026293in}{1.469472in}}%
\pgfpathlineto{\pgfqpoint{1.027442in}{1.454039in}}%
\pgfpathlineto{\pgfqpoint{1.070493in}{1.455568in}}%
\pgfpathlineto{\pgfqpoint{1.113584in}{1.456142in}}%
\pgfpathclose%
\pgfusepath{fill}%
\end{pgfscope}%
\begin{pgfscope}%
\pgfpathrectangle{\pgfqpoint{0.407592in}{0.222904in}}{\pgfqpoint{1.411163in}{1.411163in}}%
\pgfusepath{clip}%
\pgfsetbuttcap%
\pgfsetroundjoin%
\definecolor{currentfill}{rgb}{0.749020,0.772549,0.792157}%
\pgfsetfillcolor{currentfill}%
\pgfsetfillopacity{0.160000}%
\pgfsetlinewidth{0.000000pt}%
\definecolor{currentstroke}{rgb}{0.000000,0.000000,0.000000}%
\pgfsetstrokecolor{currentstroke}%
\pgfsetdash{}{0pt}%
\pgfpathmoveto{\pgfqpoint{1.385145in}{1.500462in}}%
\pgfpathlineto{\pgfqpoint{1.390166in}{1.515613in}}%
\pgfpathlineto{\pgfqpoint{1.346090in}{1.521034in}}%
\pgfpathlineto{\pgfqpoint{1.301145in}{1.525467in}}%
\pgfpathlineto{\pgfqpoint{1.297627in}{1.510206in}}%
\pgfpathlineto{\pgfqpoint{1.341815in}{1.505822in}}%
\pgfusepath{fill}%
\end{pgfscope}%
\begin{pgfscope}%
\pgfpathrectangle{\pgfqpoint{0.407592in}{0.222904in}}{\pgfqpoint{1.411163in}{1.411163in}}%
\pgfusepath{clip}%
\pgfsetbuttcap%
\pgfsetroundjoin%
\definecolor{currentfill}{rgb}{0.749020,0.772549,0.792157}%
\pgfsetfillcolor{currentfill}%
\pgfsetfillopacity{0.160000}%
\pgfsetlinewidth{0.000000pt}%
\definecolor{currentstroke}{rgb}{0.000000,0.000000,0.000000}%
\pgfsetstrokecolor{currentstroke}%
\pgfsetdash{}{0pt}%
\pgfpathmoveto{\pgfqpoint{1.283493in}{1.448893in}}%
\pgfpathlineto{\pgfqpoint{1.287036in}{1.464262in}}%
\pgfpathlineto{\pgfqpoint{1.290572in}{1.479604in}}%
\pgfpathlineto{\pgfqpoint{1.247210in}{1.482922in}}%
\pgfpathlineto{\pgfqpoint{1.203303in}{1.485270in}}%
\pgfpathlineto{\pgfqpoint{1.201311in}{1.469860in}}%
\pgfpathlineto{\pgfqpoint{1.199315in}{1.454422in}}%
\pgfpathlineto{\pgfqpoint{1.241669in}{1.452131in}}%
\pgfpathlineto{\pgfqpoint{1.283493in}{1.448893in}}%
\pgfpathclose%
\pgfusepath{fill}%
\end{pgfscope}%
\begin{pgfscope}%
\pgfpathrectangle{\pgfqpoint{0.407592in}{0.222904in}}{\pgfqpoint{1.411163in}{1.411163in}}%
\pgfusepath{clip}%
\pgfsetbuttcap%
\pgfsetroundjoin%
\definecolor{currentfill}{rgb}{0.749020,0.772549,0.792157}%
\pgfsetfillcolor{currentfill}%
\pgfsetfillopacity{0.160000}%
\pgfsetlinewidth{0.000000pt}%
\definecolor{currentstroke}{rgb}{0.000000,0.000000,0.000000}%
\pgfsetstrokecolor{currentstroke}%
\pgfsetdash{}{0pt}%
\pgfpathmoveto{\pgfqpoint{0.931195in}{1.509405in}}%
\pgfpathlineto{\pgfqpoint{0.928496in}{1.524657in}}%
\pgfpathlineto{\pgfqpoint{0.882555in}{1.520026in}}%
\pgfpathlineto{\pgfqpoint{0.837265in}{1.514408in}}%
\pgfpathlineto{\pgfqpoint{0.841501in}{1.499270in}}%
\pgfpathlineto{\pgfqpoint{0.886026in}{1.504825in}}%
\pgfusepath{fill}%
\end{pgfscope}%
\begin{pgfscope}%
\pgfpathrectangle{\pgfqpoint{0.407592in}{0.222904in}}{\pgfqpoint{1.411163in}{1.411163in}}%
\pgfusepath{clip}%
\pgfsetbuttcap%
\pgfsetroundjoin%
\definecolor{currentfill}{rgb}{0.749020,0.772549,0.792157}%
\pgfsetfillcolor{currentfill}%
\pgfsetfillopacity{0.160000}%
\pgfsetlinewidth{0.000000pt}%
\definecolor{currentstroke}{rgb}{0.000000,0.000000,0.000000}%
\pgfsetstrokecolor{currentstroke}%
\pgfsetdash{}{0pt}%
\pgfpathmoveto{\pgfqpoint{1.027442in}{1.454039in}}%
\pgfpathlineto{\pgfqpoint{1.026293in}{1.469472in}}%
\pgfpathlineto{\pgfqpoint{1.025147in}{1.484877in}}%
\pgfpathlineto{\pgfqpoint{0.980709in}{1.482334in}}%
\pgfpathlineto{\pgfqpoint{0.936606in}{1.478821in}}%
\pgfpathlineto{\pgfqpoint{0.939319in}{1.463489in}}%
\pgfpathlineto{\pgfqpoint{0.942036in}{1.448130in}}%
\pgfpathlineto{\pgfqpoint{0.984575in}{1.451557in}}%
\pgfpathlineto{\pgfqpoint{1.027442in}{1.454039in}}%
\pgfpathclose%
\pgfusepath{fill}%
\end{pgfscope}%
\begin{pgfscope}%
\pgfpathrectangle{\pgfqpoint{0.407592in}{0.222904in}}{\pgfqpoint{1.411163in}{1.411163in}}%
\pgfusepath{clip}%
\pgfsetbuttcap%
\pgfsetroundjoin%
\definecolor{currentfill}{rgb}{0.749020,0.772549,0.792157}%
\pgfsetfillcolor{currentfill}%
\pgfsetfillopacity{0.160000}%
\pgfsetlinewidth{0.000000pt}%
\definecolor{currentstroke}{rgb}{0.000000,0.000000,0.000000}%
\pgfsetstrokecolor{currentstroke}%
\pgfsetdash{}{0pt}%
\pgfpathmoveto{\pgfqpoint{1.195313in}{1.423460in}}%
\pgfpathlineto{\pgfqpoint{1.197316in}{1.438955in}}%
\pgfpathlineto{\pgfqpoint{1.199315in}{1.454422in}}%
\pgfpathlineto{\pgfqpoint{1.156573in}{1.455760in}}%
\pgfpathlineto{\pgfqpoint{1.113584in}{1.456142in}}%
\pgfpathlineto{\pgfqpoint{1.113156in}{1.440654in}}%
\pgfpathlineto{\pgfqpoint{1.112728in}{1.425136in}}%
\pgfpathlineto{\pgfqpoint{1.154140in}{1.424763in}}%
\pgfpathlineto{\pgfqpoint{1.195313in}{1.423460in}}%
\pgfpathclose%
\pgfusepath{fill}%
\end{pgfscope}%
\begin{pgfscope}%
\pgfpathrectangle{\pgfqpoint{0.407592in}{0.222904in}}{\pgfqpoint{1.411163in}{1.411163in}}%
\pgfusepath{clip}%
\pgfsetbuttcap%
\pgfsetroundjoin%
\definecolor{currentfill}{rgb}{0.749020,0.772549,0.792157}%
\pgfsetfillcolor{currentfill}%
\pgfsetfillopacity{0.160000}%
\pgfsetlinewidth{0.000000pt}%
\definecolor{currentstroke}{rgb}{0.000000,0.000000,0.000000}%
\pgfsetstrokecolor{currentstroke}%
\pgfsetdash{}{0pt}%
\pgfpathmoveto{\pgfqpoint{1.112728in}{1.425136in}}%
\pgfpathlineto{\pgfqpoint{1.113156in}{1.440654in}}%
\pgfpathlineto{\pgfqpoint{1.113584in}{1.456142in}}%
\pgfpathlineto{\pgfqpoint{1.070493in}{1.455568in}}%
\pgfpathlineto{\pgfqpoint{1.027442in}{1.454039in}}%
\pgfpathlineto{\pgfqpoint{1.028593in}{1.438577in}}%
\pgfpathlineto{\pgfqpoint{1.029746in}{1.423087in}}%
\pgfpathlineto{\pgfqpoint{1.071217in}{1.424577in}}%
\pgfpathlineto{\pgfqpoint{1.112728in}{1.425136in}}%
\pgfpathclose%
\pgfusepath{fill}%
\end{pgfscope}%
\begin{pgfscope}%
\pgfpathrectangle{\pgfqpoint{0.407592in}{0.222904in}}{\pgfqpoint{1.411163in}{1.411163in}}%
\pgfusepath{clip}%
\pgfsetbuttcap%
\pgfsetroundjoin%
\definecolor{currentfill}{rgb}{0.749020,0.772549,0.792157}%
\pgfsetfillcolor{currentfill}%
\pgfsetfillopacity{0.160000}%
\pgfsetlinewidth{0.000000pt}%
\definecolor{currentstroke}{rgb}{0.000000,0.000000,0.000000}%
\pgfsetstrokecolor{currentstroke}%
\pgfsetdash{}{0pt}%
\pgfpathmoveto{\pgfqpoint{1.375079in}{1.470084in}}%
\pgfpathlineto{\pgfqpoint{1.380116in}{1.485285in}}%
\pgfpathlineto{\pgfqpoint{1.385145in}{1.500462in}}%
\pgfpathlineto{\pgfqpoint{1.341815in}{1.505822in}}%
\pgfpathlineto{\pgfqpoint{1.297627in}{1.510206in}}%
\pgfpathlineto{\pgfqpoint{1.294103in}{1.494918in}}%
\pgfpathlineto{\pgfqpoint{1.290572in}{1.479604in}}%
\pgfpathlineto{\pgfqpoint{1.333244in}{1.475321in}}%
\pgfpathlineto{\pgfqpoint{1.375079in}{1.470084in}}%
\pgfpathclose%
\pgfusepath{fill}%
\end{pgfscope}%
\begin{pgfscope}%
\pgfpathrectangle{\pgfqpoint{0.407592in}{0.222904in}}{\pgfqpoint{1.411163in}{1.411163in}}%
\pgfusepath{clip}%
\pgfsetbuttcap%
\pgfsetroundjoin%
\definecolor{currentfill}{rgb}{0.749020,0.772549,0.792157}%
\pgfsetfillcolor{currentfill}%
\pgfsetfillopacity{0.160000}%
\pgfsetlinewidth{0.000000pt}%
\definecolor{currentstroke}{rgb}{0.000000,0.000000,0.000000}%
\pgfsetstrokecolor{currentstroke}%
\pgfsetdash{}{0pt}%
\pgfpathmoveto{\pgfqpoint{0.936606in}{1.478821in}}%
\pgfpathlineto{\pgfqpoint{0.933898in}{1.494127in}}%
\pgfpathlineto{\pgfqpoint{0.931195in}{1.509405in}}%
\pgfpathlineto{\pgfqpoint{0.886026in}{1.504825in}}%
\pgfpathlineto{\pgfqpoint{0.841501in}{1.499270in}}%
\pgfpathlineto{\pgfqpoint{0.845745in}{1.484107in}}%
\pgfpathlineto{\pgfqpoint{0.849995in}{1.468920in}}%
\pgfpathlineto{\pgfqpoint{0.892985in}{1.474347in}}%
\pgfpathlineto{\pgfqpoint{0.936606in}{1.478821in}}%
\pgfpathclose%
\pgfusepath{fill}%
\end{pgfscope}%
\begin{pgfscope}%
\pgfpathrectangle{\pgfqpoint{0.407592in}{0.222904in}}{\pgfqpoint{1.411163in}{1.411163in}}%
\pgfusepath{clip}%
\pgfsetbuttcap%
\pgfsetroundjoin%
\definecolor{currentfill}{rgb}{0.749020,0.772549,0.792157}%
\pgfsetfillcolor{currentfill}%
\pgfsetfillopacity{0.160000}%
\pgfsetlinewidth{0.000000pt}%
\definecolor{currentstroke}{rgb}{0.000000,0.000000,0.000000}%
\pgfsetstrokecolor{currentstroke}%
\pgfsetdash{}{0pt}%
\pgfpathmoveto{\pgfqpoint{1.276388in}{1.418074in}}%
\pgfpathlineto{\pgfqpoint{1.279944in}{1.433497in}}%
\pgfpathlineto{\pgfqpoint{1.283493in}{1.448893in}}%
\pgfpathlineto{\pgfqpoint{1.241669in}{1.452131in}}%
\pgfpathlineto{\pgfqpoint{1.199315in}{1.454422in}}%
\pgfpathlineto{\pgfqpoint{1.197316in}{1.438955in}}%
\pgfpathlineto{\pgfqpoint{1.195313in}{1.423460in}}%
\pgfpathlineto{\pgfqpoint{1.236109in}{1.421228in}}%
\pgfpathlineto{\pgfqpoint{1.276388in}{1.418074in}}%
\pgfpathclose%
\pgfusepath{fill}%
\end{pgfscope}%
\begin{pgfscope}%
\pgfpathrectangle{\pgfqpoint{0.407592in}{0.222904in}}{\pgfqpoint{1.411163in}{1.411163in}}%
\pgfusepath{clip}%
\pgfsetbuttcap%
\pgfsetroundjoin%
\definecolor{currentfill}{rgb}{0.749020,0.772549,0.792157}%
\pgfsetfillcolor{currentfill}%
\pgfsetfillopacity{0.160000}%
\pgfsetlinewidth{0.000000pt}%
\definecolor{currentstroke}{rgb}{0.000000,0.000000,0.000000}%
\pgfsetstrokecolor{currentstroke}%
\pgfsetdash{}{0pt}%
\pgfpathmoveto{\pgfqpoint{1.029746in}{1.423087in}}%
\pgfpathlineto{\pgfqpoint{1.028593in}{1.438577in}}%
\pgfpathlineto{\pgfqpoint{1.027442in}{1.454039in}}%
\pgfpathlineto{\pgfqpoint{0.984575in}{1.451557in}}%
\pgfpathlineto{\pgfqpoint{0.942036in}{1.448130in}}%
\pgfpathlineto{\pgfqpoint{0.944758in}{1.432744in}}%
\pgfpathlineto{\pgfqpoint{0.947485in}{1.417331in}}%
\pgfpathlineto{\pgfqpoint{0.988456in}{1.420669in}}%
\pgfpathlineto{\pgfqpoint{1.029746in}{1.423087in}}%
\pgfpathclose%
\pgfusepath{fill}%
\end{pgfscope}%
\begin{pgfscope}%
\pgfpathrectangle{\pgfqpoint{0.407592in}{0.222904in}}{\pgfqpoint{1.411163in}{1.411163in}}%
\pgfusepath{clip}%
\pgfsetbuttcap%
\pgfsetroundjoin%
\definecolor{currentfill}{rgb}{0.749020,0.772549,0.792157}%
\pgfsetfillcolor{currentfill}%
\pgfsetfillopacity{0.160000}%
\pgfsetlinewidth{0.000000pt}%
\definecolor{currentstroke}{rgb}{0.000000,0.000000,0.000000}%
\pgfsetstrokecolor{currentstroke}%
\pgfsetdash{}{0pt}%
\pgfpathmoveto{\pgfqpoint{1.191296in}{1.392384in}}%
\pgfpathlineto{\pgfqpoint{1.193306in}{1.407936in}}%
\pgfpathlineto{\pgfqpoint{1.195313in}{1.423460in}}%
\pgfpathlineto{\pgfqpoint{1.154140in}{1.424763in}}%
\pgfpathlineto{\pgfqpoint{1.112728in}{1.425136in}}%
\pgfpathlineto{\pgfqpoint{1.112298in}{1.409589in}}%
\pgfpathlineto{\pgfqpoint{1.111868in}{1.394014in}}%
\pgfpathlineto{\pgfqpoint{1.151698in}{1.393652in}}%
\pgfpathlineto{\pgfqpoint{1.191296in}{1.392384in}}%
\pgfpathclose%
\pgfusepath{fill}%
\end{pgfscope}%
\begin{pgfscope}%
\pgfpathrectangle{\pgfqpoint{0.407592in}{0.222904in}}{\pgfqpoint{1.411163in}{1.411163in}}%
\pgfusepath{clip}%
\pgfsetbuttcap%
\pgfsetroundjoin%
\definecolor{currentfill}{rgb}{0.749020,0.772549,0.792157}%
\pgfsetfillcolor{currentfill}%
\pgfsetfillopacity{0.160000}%
\pgfsetlinewidth{0.000000pt}%
\definecolor{currentstroke}{rgb}{0.000000,0.000000,0.000000}%
\pgfsetstrokecolor{currentstroke}%
\pgfsetdash{}{0pt}%
\pgfpathmoveto{\pgfqpoint{1.111868in}{1.394014in}}%
\pgfpathlineto{\pgfqpoint{1.112298in}{1.409589in}}%
\pgfpathlineto{\pgfqpoint{1.112728in}{1.425136in}}%
\pgfpathlineto{\pgfqpoint{1.071217in}{1.424577in}}%
\pgfpathlineto{\pgfqpoint{1.029746in}{1.423087in}}%
\pgfpathlineto{\pgfqpoint{1.030901in}{1.407568in}}%
\pgfpathlineto{\pgfqpoint{1.032059in}{1.392021in}}%
\pgfpathlineto{\pgfqpoint{1.071943in}{1.393470in}}%
\pgfpathlineto{\pgfqpoint{1.111868in}{1.394014in}}%
\pgfpathclose%
\pgfusepath{fill}%
\end{pgfscope}%
\begin{pgfscope}%
\pgfpathrectangle{\pgfqpoint{0.407592in}{0.222904in}}{\pgfqpoint{1.411163in}{1.411163in}}%
\pgfusepath{clip}%
\pgfsetbuttcap%
\pgfsetroundjoin%
\definecolor{currentfill}{rgb}{0.749020,0.772549,0.792157}%
\pgfsetfillcolor{currentfill}%
\pgfsetfillopacity{0.160000}%
\pgfsetlinewidth{0.000000pt}%
\definecolor{currentstroke}{rgb}{0.000000,0.000000,0.000000}%
\pgfsetstrokecolor{currentstroke}%
\pgfsetdash{}{0pt}%
\pgfpathmoveto{\pgfqpoint{1.364979in}{1.439607in}}%
\pgfpathlineto{\pgfqpoint{1.370033in}{1.454858in}}%
\pgfpathlineto{\pgfqpoint{1.375079in}{1.470084in}}%
\pgfpathlineto{\pgfqpoint{1.333244in}{1.475321in}}%
\pgfpathlineto{\pgfqpoint{1.290572in}{1.479604in}}%
\pgfpathlineto{\pgfqpoint{1.287036in}{1.464262in}}%
\pgfpathlineto{\pgfqpoint{1.283493in}{1.448893in}}%
\pgfpathlineto{\pgfqpoint{1.324643in}{1.444715in}}%
\pgfpathlineto{\pgfqpoint{1.364979in}{1.439607in}}%
\pgfpathclose%
\pgfusepath{fill}%
\end{pgfscope}%
\begin{pgfscope}%
\pgfpathrectangle{\pgfqpoint{0.407592in}{0.222904in}}{\pgfqpoint{1.411163in}{1.411163in}}%
\pgfusepath{clip}%
\pgfsetbuttcap%
\pgfsetroundjoin%
\definecolor{currentfill}{rgb}{0.749020,0.772549,0.792157}%
\pgfsetfillcolor{currentfill}%
\pgfsetfillopacity{0.160000}%
\pgfsetlinewidth{0.000000pt}%
\definecolor{currentstroke}{rgb}{0.000000,0.000000,0.000000}%
\pgfsetstrokecolor{currentstroke}%
\pgfsetdash{}{0pt}%
\pgfpathmoveto{\pgfqpoint{0.942036in}{1.448130in}}%
\pgfpathlineto{\pgfqpoint{0.939319in}{1.463489in}}%
\pgfpathlineto{\pgfqpoint{0.936606in}{1.478821in}}%
\pgfpathlineto{\pgfqpoint{0.892985in}{1.474347in}}%
\pgfpathlineto{\pgfqpoint{0.849995in}{1.468920in}}%
\pgfpathlineto{\pgfqpoint{0.854253in}{1.453708in}}%
\pgfpathlineto{\pgfqpoint{0.858517in}{1.438471in}}%
\pgfpathlineto{\pgfqpoint{0.899968in}{1.443765in}}%
\pgfpathlineto{\pgfqpoint{0.942036in}{1.448130in}}%
\pgfpathclose%
\pgfusepath{fill}%
\end{pgfscope}%
\begin{pgfscope}%
\pgfpathrectangle{\pgfqpoint{0.407592in}{0.222904in}}{\pgfqpoint{1.411163in}{1.411163in}}%
\pgfusepath{clip}%
\pgfsetbuttcap%
\pgfsetroundjoin%
\definecolor{currentfill}{rgb}{0.749020,0.772549,0.792157}%
\pgfsetfillcolor{currentfill}%
\pgfsetfillopacity{0.160000}%
\pgfsetlinewidth{0.000000pt}%
\definecolor{currentstroke}{rgb}{0.000000,0.000000,0.000000}%
\pgfsetstrokecolor{currentstroke}%
\pgfsetdash{}{0pt}%
\pgfpathmoveto{\pgfqpoint{1.269258in}{1.387146in}}%
\pgfpathlineto{\pgfqpoint{1.272827in}{1.402624in}}%
\pgfpathlineto{\pgfqpoint{1.276388in}{1.418074in}}%
\pgfpathlineto{\pgfqpoint{1.236109in}{1.421228in}}%
\pgfpathlineto{\pgfqpoint{1.195313in}{1.423460in}}%
\pgfpathlineto{\pgfqpoint{1.193306in}{1.407936in}}%
\pgfpathlineto{\pgfqpoint{1.191296in}{1.392384in}}%
\pgfpathlineto{\pgfqpoint{1.230528in}{1.390213in}}%
\pgfpathlineto{\pgfqpoint{1.269258in}{1.387146in}}%
\pgfpathclose%
\pgfusepath{fill}%
\end{pgfscope}%
\begin{pgfscope}%
\pgfpathrectangle{\pgfqpoint{0.407592in}{0.222904in}}{\pgfqpoint{1.411163in}{1.411163in}}%
\pgfusepath{clip}%
\pgfsetbuttcap%
\pgfsetroundjoin%
\definecolor{currentfill}{rgb}{0.749020,0.772549,0.792157}%
\pgfsetfillcolor{currentfill}%
\pgfsetfillopacity{0.160000}%
\pgfsetlinewidth{0.000000pt}%
\definecolor{currentstroke}{rgb}{0.000000,0.000000,0.000000}%
\pgfsetstrokecolor{currentstroke}%
\pgfsetdash{}{0pt}%
\pgfpathmoveto{\pgfqpoint{1.032059in}{1.392021in}}%
\pgfpathlineto{\pgfqpoint{1.030901in}{1.407568in}}%
\pgfpathlineto{\pgfqpoint{1.029746in}{1.423087in}}%
\pgfpathlineto{\pgfqpoint{0.988456in}{1.420669in}}%
\pgfpathlineto{\pgfqpoint{0.947485in}{1.417331in}}%
\pgfpathlineto{\pgfqpoint{0.950217in}{1.401891in}}%
\pgfpathlineto{\pgfqpoint{0.952954in}{1.386423in}}%
\pgfpathlineto{\pgfqpoint{0.992350in}{1.389670in}}%
\pgfpathlineto{\pgfqpoint{1.032059in}{1.392021in}}%
\pgfpathclose%
\pgfusepath{fill}%
\end{pgfscope}%
\begin{pgfscope}%
\pgfpathrectangle{\pgfqpoint{0.407592in}{0.222904in}}{\pgfqpoint{1.411163in}{1.411163in}}%
\pgfusepath{clip}%
\pgfsetbuttcap%
\pgfsetroundjoin%
\definecolor{currentfill}{rgb}{0.749020,0.772549,0.792157}%
\pgfsetfillcolor{currentfill}%
\pgfsetfillopacity{0.160000}%
\pgfsetlinewidth{0.000000pt}%
\definecolor{currentstroke}{rgb}{0.000000,0.000000,0.000000}%
\pgfsetstrokecolor{currentstroke}%
\pgfsetdash{}{0pt}%
\pgfpathmoveto{\pgfqpoint{1.468638in}{1.486859in}}%
\pgfpathlineto{\pgfqpoint{1.475114in}{1.501856in}}%
\pgfpathlineto{\pgfqpoint{1.433225in}{1.509216in}}%
\pgfpathlineto{\pgfqpoint{1.390166in}{1.515613in}}%
\pgfpathlineto{\pgfqpoint{1.385145in}{1.500462in}}%
\pgfpathlineto{\pgfqpoint{1.427469in}{1.494136in}}%
\pgfusepath{fill}%
\end{pgfscope}%
\begin{pgfscope}%
\pgfpathrectangle{\pgfqpoint{0.407592in}{0.222904in}}{\pgfqpoint{1.411163in}{1.411163in}}%
\pgfusepath{clip}%
\pgfsetbuttcap%
\pgfsetroundjoin%
\definecolor{currentfill}{rgb}{0.749020,0.772549,0.792157}%
\pgfsetfillcolor{currentfill}%
\pgfsetfillopacity{0.160000}%
\pgfsetlinewidth{0.000000pt}%
\definecolor{currentstroke}{rgb}{0.000000,0.000000,0.000000}%
\pgfsetstrokecolor{currentstroke}%
\pgfsetdash{}{0pt}%
\pgfpathmoveto{\pgfqpoint{1.187264in}{1.361193in}}%
\pgfpathlineto{\pgfqpoint{1.189282in}{1.376803in}}%
\pgfpathlineto{\pgfqpoint{1.191296in}{1.392384in}}%
\pgfpathlineto{\pgfqpoint{1.151698in}{1.393652in}}%
\pgfpathlineto{\pgfqpoint{1.111868in}{1.394014in}}%
\pgfpathlineto{\pgfqpoint{1.111437in}{1.378409in}}%
\pgfpathlineto{\pgfqpoint{1.111005in}{1.362775in}}%
\pgfpathlineto{\pgfqpoint{1.149246in}{1.362423in}}%
\pgfpathlineto{\pgfqpoint{1.187264in}{1.361193in}}%
\pgfpathclose%
\pgfusepath{fill}%
\end{pgfscope}%
\begin{pgfscope}%
\pgfpathrectangle{\pgfqpoint{0.407592in}{0.222904in}}{\pgfqpoint{1.411163in}{1.411163in}}%
\pgfusepath{clip}%
\pgfsetbuttcap%
\pgfsetroundjoin%
\definecolor{currentfill}{rgb}{0.749020,0.772549,0.792157}%
\pgfsetfillcolor{currentfill}%
\pgfsetfillopacity{0.160000}%
\pgfsetlinewidth{0.000000pt}%
\definecolor{currentstroke}{rgb}{0.000000,0.000000,0.000000}%
\pgfsetstrokecolor{currentstroke}%
\pgfsetdash{}{0pt}%
\pgfpathmoveto{\pgfqpoint{1.111005in}{1.362775in}}%
\pgfpathlineto{\pgfqpoint{1.111437in}{1.378409in}}%
\pgfpathlineto{\pgfqpoint{1.111868in}{1.394014in}}%
\pgfpathlineto{\pgfqpoint{1.071943in}{1.393470in}}%
\pgfpathlineto{\pgfqpoint{1.032059in}{1.392021in}}%
\pgfpathlineto{\pgfqpoint{1.033218in}{1.376445in}}%
\pgfpathlineto{\pgfqpoint{1.034380in}{1.360840in}}%
\pgfpathlineto{\pgfqpoint{1.072673in}{1.362247in}}%
\pgfpathlineto{\pgfqpoint{1.111005in}{1.362775in}}%
\pgfpathclose%
\pgfusepath{fill}%
\end{pgfscope}%
\begin{pgfscope}%
\pgfpathrectangle{\pgfqpoint{0.407592in}{0.222904in}}{\pgfqpoint{1.411163in}{1.411163in}}%
\pgfusepath{clip}%
\pgfsetbuttcap%
\pgfsetroundjoin%
\definecolor{currentfill}{rgb}{0.749020,0.772549,0.792157}%
\pgfsetfillcolor{currentfill}%
\pgfsetfillopacity{0.160000}%
\pgfsetlinewidth{0.000000pt}%
\definecolor{currentstroke}{rgb}{0.000000,0.000000,0.000000}%
\pgfsetstrokecolor{currentstroke}%
\pgfsetdash{}{0pt}%
\pgfpathmoveto{\pgfqpoint{1.354846in}{1.409030in}}%
\pgfpathlineto{\pgfqpoint{1.359916in}{1.424331in}}%
\pgfpathlineto{\pgfqpoint{1.364979in}{1.439607in}}%
\pgfpathlineto{\pgfqpoint{1.324643in}{1.444715in}}%
\pgfpathlineto{\pgfqpoint{1.283493in}{1.448893in}}%
\pgfpathlineto{\pgfqpoint{1.279944in}{1.433497in}}%
\pgfpathlineto{\pgfqpoint{1.276388in}{1.418074in}}%
\pgfpathlineto{\pgfqpoint{1.316013in}{1.414005in}}%
\pgfpathlineto{\pgfqpoint{1.354846in}{1.409030in}}%
\pgfpathclose%
\pgfusepath{fill}%
\end{pgfscope}%
\begin{pgfscope}%
\pgfpathrectangle{\pgfqpoint{0.407592in}{0.222904in}}{\pgfqpoint{1.411163in}{1.411163in}}%
\pgfusepath{clip}%
\pgfsetbuttcap%
\pgfsetroundjoin%
\definecolor{currentfill}{rgb}{0.749020,0.772549,0.792157}%
\pgfsetfillcolor{currentfill}%
\pgfsetfillopacity{0.160000}%
\pgfsetlinewidth{0.000000pt}%
\definecolor{currentstroke}{rgb}{0.000000,0.000000,0.000000}%
\pgfsetstrokecolor{currentstroke}%
\pgfsetdash{}{0pt}%
\pgfpathmoveto{\pgfqpoint{0.841501in}{1.499270in}}%
\pgfpathlineto{\pgfqpoint{0.837265in}{1.514408in}}%
\pgfpathlineto{\pgfqpoint{0.792780in}{1.507816in}}%
\pgfpathlineto{\pgfqpoint{0.749254in}{1.500265in}}%
\pgfpathlineto{\pgfqpoint{0.754995in}{1.485286in}}%
\pgfpathlineto{\pgfqpoint{0.797774in}{1.492752in}}%
\pgfusepath{fill}%
\end{pgfscope}%
\begin{pgfscope}%
\pgfpathrectangle{\pgfqpoint{0.407592in}{0.222904in}}{\pgfqpoint{1.411163in}{1.411163in}}%
\pgfusepath{clip}%
\pgfsetbuttcap%
\pgfsetroundjoin%
\definecolor{currentfill}{rgb}{0.749020,0.772549,0.792157}%
\pgfsetfillcolor{currentfill}%
\pgfsetfillopacity{0.160000}%
\pgfsetlinewidth{0.000000pt}%
\definecolor{currentstroke}{rgb}{0.000000,0.000000,0.000000}%
\pgfsetstrokecolor{currentstroke}%
\pgfsetdash{}{0pt}%
\pgfpathmoveto{\pgfqpoint{0.947485in}{1.417331in}}%
\pgfpathlineto{\pgfqpoint{0.944758in}{1.432744in}}%
\pgfpathlineto{\pgfqpoint{0.942036in}{1.448130in}}%
\pgfpathlineto{\pgfqpoint{0.899968in}{1.443765in}}%
\pgfpathlineto{\pgfqpoint{0.858517in}{1.438471in}}%
\pgfpathlineto{\pgfqpoint{0.862788in}{1.423210in}}%
\pgfpathlineto{\pgfqpoint{0.867066in}{1.407924in}}%
\pgfpathlineto{\pgfqpoint{0.906975in}{1.413079in}}%
\pgfpathlineto{\pgfqpoint{0.947485in}{1.417331in}}%
\pgfpathclose%
\pgfusepath{fill}%
\end{pgfscope}%
\begin{pgfscope}%
\pgfpathrectangle{\pgfqpoint{0.407592in}{0.222904in}}{\pgfqpoint{1.411163in}{1.411163in}}%
\pgfusepath{clip}%
\pgfsetbuttcap%
\pgfsetroundjoin%
\definecolor{currentfill}{rgb}{0.749020,0.772549,0.792157}%
\pgfsetfillcolor{currentfill}%
\pgfsetfillopacity{0.160000}%
\pgfsetlinewidth{0.000000pt}%
\definecolor{currentstroke}{rgb}{0.000000,0.000000,0.000000}%
\pgfsetstrokecolor{currentstroke}%
\pgfsetdash{}{0pt}%
\pgfpathmoveto{\pgfqpoint{1.455658in}{1.456799in}}%
\pgfpathlineto{\pgfqpoint{1.462153in}{1.471840in}}%
\pgfpathlineto{\pgfqpoint{1.468638in}{1.486859in}}%
\pgfpathlineto{\pgfqpoint{1.427469in}{1.494136in}}%
\pgfpathlineto{\pgfqpoint{1.385145in}{1.500462in}}%
\pgfpathlineto{\pgfqpoint{1.380116in}{1.485285in}}%
\pgfpathlineto{\pgfqpoint{1.375079in}{1.470084in}}%
\pgfpathlineto{\pgfqpoint{1.415932in}{1.463906in}}%
\pgfpathlineto{\pgfqpoint{1.455658in}{1.456799in}}%
\pgfpathclose%
\pgfusepath{fill}%
\end{pgfscope}%
\begin{pgfscope}%
\pgfpathrectangle{\pgfqpoint{0.407592in}{0.222904in}}{\pgfqpoint{1.411163in}{1.411163in}}%
\pgfusepath{clip}%
\pgfsetbuttcap%
\pgfsetroundjoin%
\definecolor{currentfill}{rgb}{0.749020,0.772549,0.792157}%
\pgfsetfillcolor{currentfill}%
\pgfsetfillopacity{0.160000}%
\pgfsetlinewidth{0.000000pt}%
\definecolor{currentstroke}{rgb}{0.000000,0.000000,0.000000}%
\pgfsetstrokecolor{currentstroke}%
\pgfsetdash{}{0pt}%
\pgfpathmoveto{\pgfqpoint{1.262103in}{1.356108in}}%
\pgfpathlineto{\pgfqpoint{1.265684in}{1.371641in}}%
\pgfpathlineto{\pgfqpoint{1.269258in}{1.387146in}}%
\pgfpathlineto{\pgfqpoint{1.230528in}{1.390213in}}%
\pgfpathlineto{\pgfqpoint{1.191296in}{1.392384in}}%
\pgfpathlineto{\pgfqpoint{1.189282in}{1.376803in}}%
\pgfpathlineto{\pgfqpoint{1.187264in}{1.361193in}}%
\pgfpathlineto{\pgfqpoint{1.224927in}{1.359086in}}%
\pgfpathlineto{\pgfqpoint{1.262103in}{1.356108in}}%
\pgfpathclose%
\pgfusepath{fill}%
\end{pgfscope}%
\begin{pgfscope}%
\pgfpathrectangle{\pgfqpoint{0.407592in}{0.222904in}}{\pgfqpoint{1.411163in}{1.411163in}}%
\pgfusepath{clip}%
\pgfsetbuttcap%
\pgfsetroundjoin%
\definecolor{currentfill}{rgb}{0.749020,0.772549,0.792157}%
\pgfsetfillcolor{currentfill}%
\pgfsetfillopacity{0.160000}%
\pgfsetlinewidth{0.000000pt}%
\definecolor{currentstroke}{rgb}{0.000000,0.000000,0.000000}%
\pgfsetstrokecolor{currentstroke}%
\pgfsetdash{}{0pt}%
\pgfpathmoveto{\pgfqpoint{1.034380in}{1.360840in}}%
\pgfpathlineto{\pgfqpoint{1.033218in}{1.376445in}}%
\pgfpathlineto{\pgfqpoint{1.032059in}{1.392021in}}%
\pgfpathlineto{\pgfqpoint{0.992350in}{1.389670in}}%
\pgfpathlineto{\pgfqpoint{0.952954in}{1.386423in}}%
\pgfpathlineto{\pgfqpoint{0.955695in}{1.370928in}}%
\pgfpathlineto{\pgfqpoint{0.958441in}{1.355406in}}%
\pgfpathlineto{\pgfqpoint{0.996258in}{1.358558in}}%
\pgfpathlineto{\pgfqpoint{1.034380in}{1.360840in}}%
\pgfpathclose%
\pgfusepath{fill}%
\end{pgfscope}%
\begin{pgfscope}%
\pgfpathrectangle{\pgfqpoint{0.407592in}{0.222904in}}{\pgfqpoint{1.411163in}{1.411163in}}%
\pgfusepath{clip}%
\pgfsetbuttcap%
\pgfsetroundjoin%
\definecolor{currentfill}{rgb}{0.749020,0.772549,0.792157}%
\pgfsetfillcolor{currentfill}%
\pgfsetfillopacity{0.160000}%
\pgfsetlinewidth{0.000000pt}%
\definecolor{currentstroke}{rgb}{0.000000,0.000000,0.000000}%
\pgfsetstrokecolor{currentstroke}%
\pgfsetdash{}{0pt}%
\pgfpathmoveto{\pgfqpoint{0.849995in}{1.468920in}}%
\pgfpathlineto{\pgfqpoint{0.845745in}{1.484107in}}%
\pgfpathlineto{\pgfqpoint{0.841501in}{1.499270in}}%
\pgfpathlineto{\pgfqpoint{0.797774in}{1.492752in}}%
\pgfpathlineto{\pgfqpoint{0.754995in}{1.485286in}}%
\pgfpathlineto{\pgfqpoint{0.760743in}{1.470286in}}%
\pgfpathlineto{\pgfqpoint{0.766500in}{1.455264in}}%
\pgfpathlineto{\pgfqpoint{0.807784in}{1.462554in}}%
\pgfpathlineto{\pgfqpoint{0.849995in}{1.468920in}}%
\pgfpathclose%
\pgfusepath{fill}%
\end{pgfscope}%
\begin{pgfscope}%
\pgfpathrectangle{\pgfqpoint{0.407592in}{0.222904in}}{\pgfqpoint{1.411163in}{1.411163in}}%
\pgfusepath{clip}%
\pgfsetbuttcap%
\pgfsetroundjoin%
\definecolor{currentfill}{rgb}{0.749020,0.772549,0.792157}%
\pgfsetfillcolor{currentfill}%
\pgfsetfillopacity{0.160000}%
\pgfsetlinewidth{0.000000pt}%
\definecolor{currentstroke}{rgb}{0.000000,0.000000,0.000000}%
\pgfsetstrokecolor{currentstroke}%
\pgfsetdash{}{0pt}%
\pgfpathmoveto{\pgfqpoint{1.183217in}{1.329886in}}%
\pgfpathlineto{\pgfqpoint{1.185242in}{1.345554in}}%
\pgfpathlineto{\pgfqpoint{1.187264in}{1.361193in}}%
\pgfpathlineto{\pgfqpoint{1.149246in}{1.362423in}}%
\pgfpathlineto{\pgfqpoint{1.111005in}{1.362775in}}%
\pgfpathlineto{\pgfqpoint{1.110573in}{1.347112in}}%
\pgfpathlineto{\pgfqpoint{1.110139in}{1.331419in}}%
\pgfpathlineto{\pgfqpoint{1.146786in}{1.331078in}}%
\pgfpathlineto{\pgfqpoint{1.183217in}{1.329886in}}%
\pgfpathclose%
\pgfusepath{fill}%
\end{pgfscope}%
\begin{pgfscope}%
\pgfpathrectangle{\pgfqpoint{0.407592in}{0.222904in}}{\pgfqpoint{1.411163in}{1.411163in}}%
\pgfusepath{clip}%
\pgfsetbuttcap%
\pgfsetroundjoin%
\definecolor{currentfill}{rgb}{0.749020,0.772549,0.792157}%
\pgfsetfillcolor{currentfill}%
\pgfsetfillopacity{0.160000}%
\pgfsetlinewidth{0.000000pt}%
\definecolor{currentstroke}{rgb}{0.000000,0.000000,0.000000}%
\pgfsetstrokecolor{currentstroke}%
\pgfsetdash{}{0pt}%
\pgfpathmoveto{\pgfqpoint{1.110139in}{1.331419in}}%
\pgfpathlineto{\pgfqpoint{1.110573in}{1.347112in}}%
\pgfpathlineto{\pgfqpoint{1.111005in}{1.362775in}}%
\pgfpathlineto{\pgfqpoint{1.072673in}{1.362247in}}%
\pgfpathlineto{\pgfqpoint{1.034380in}{1.360840in}}%
\pgfpathlineto{\pgfqpoint{1.035543in}{1.345206in}}%
\pgfpathlineto{\pgfqpoint{1.036709in}{1.329544in}}%
\pgfpathlineto{\pgfqpoint{1.073405in}{1.330907in}}%
\pgfpathlineto{\pgfqpoint{1.110139in}{1.331419in}}%
\pgfpathclose%
\pgfusepath{fill}%
\end{pgfscope}%
\begin{pgfscope}%
\pgfpathrectangle{\pgfqpoint{0.407592in}{0.222904in}}{\pgfqpoint{1.411163in}{1.411163in}}%
\pgfusepath{clip}%
\pgfsetbuttcap%
\pgfsetroundjoin%
\definecolor{currentfill}{rgb}{0.749020,0.772549,0.792157}%
\pgfsetfillcolor{currentfill}%
\pgfsetfillopacity{0.160000}%
\pgfsetlinewidth{0.000000pt}%
\definecolor{currentstroke}{rgb}{0.000000,0.000000,0.000000}%
\pgfsetstrokecolor{currentstroke}%
\pgfsetdash{}{0pt}%
\pgfpathmoveto{\pgfqpoint{1.344679in}{1.378352in}}%
\pgfpathlineto{\pgfqpoint{1.349767in}{1.393704in}}%
\pgfpathlineto{\pgfqpoint{1.354846in}{1.409030in}}%
\pgfpathlineto{\pgfqpoint{1.316013in}{1.414005in}}%
\pgfpathlineto{\pgfqpoint{1.276388in}{1.418074in}}%
\pgfpathlineto{\pgfqpoint{1.272827in}{1.402624in}}%
\pgfpathlineto{\pgfqpoint{1.269258in}{1.387146in}}%
\pgfpathlineto{\pgfqpoint{1.307354in}{1.383189in}}%
\pgfpathlineto{\pgfqpoint{1.344679in}{1.378352in}}%
\pgfpathclose%
\pgfusepath{fill}%
\end{pgfscope}%
\begin{pgfscope}%
\pgfpathrectangle{\pgfqpoint{0.407592in}{0.222904in}}{\pgfqpoint{1.411163in}{1.411163in}}%
\pgfusepath{clip}%
\pgfsetbuttcap%
\pgfsetroundjoin%
\definecolor{currentfill}{rgb}{0.749020,0.772549,0.792157}%
\pgfsetfillcolor{currentfill}%
\pgfsetfillopacity{0.160000}%
\pgfsetlinewidth{0.000000pt}%
\definecolor{currentstroke}{rgb}{0.000000,0.000000,0.000000}%
\pgfsetstrokecolor{currentstroke}%
\pgfsetdash{}{0pt}%
\pgfpathmoveto{\pgfqpoint{0.952954in}{1.386423in}}%
\pgfpathlineto{\pgfqpoint{0.950217in}{1.401891in}}%
\pgfpathlineto{\pgfqpoint{0.947485in}{1.417331in}}%
\pgfpathlineto{\pgfqpoint{0.906975in}{1.413079in}}%
\pgfpathlineto{\pgfqpoint{0.867066in}{1.407924in}}%
\pgfpathlineto{\pgfqpoint{0.871351in}{1.392613in}}%
\pgfpathlineto{\pgfqpoint{0.875643in}{1.377277in}}%
\pgfpathlineto{\pgfqpoint{0.914006in}{1.382289in}}%
\pgfpathlineto{\pgfqpoint{0.952954in}{1.386423in}}%
\pgfpathclose%
\pgfusepath{fill}%
\end{pgfscope}%
\begin{pgfscope}%
\pgfpathrectangle{\pgfqpoint{0.407592in}{0.222904in}}{\pgfqpoint{1.411163in}{1.411163in}}%
\pgfusepath{clip}%
\pgfsetbuttcap%
\pgfsetroundjoin%
\definecolor{currentfill}{rgb}{0.749020,0.772549,0.792157}%
\pgfsetfillcolor{currentfill}%
\pgfsetfillopacity{0.160000}%
\pgfsetlinewidth{0.000000pt}%
\definecolor{currentstroke}{rgb}{0.000000,0.000000,0.000000}%
\pgfsetstrokecolor{currentstroke}%
\pgfsetdash{}{0pt}%
\pgfpathmoveto{\pgfqpoint{1.442640in}{1.426652in}}%
\pgfpathlineto{\pgfqpoint{1.449154in}{1.441737in}}%
\pgfpathlineto{\pgfqpoint{1.455658in}{1.456799in}}%
\pgfpathlineto{\pgfqpoint{1.415932in}{1.463906in}}%
\pgfpathlineto{\pgfqpoint{1.375079in}{1.470084in}}%
\pgfpathlineto{\pgfqpoint{1.370033in}{1.454858in}}%
\pgfpathlineto{\pgfqpoint{1.364979in}{1.439607in}}%
\pgfpathlineto{\pgfqpoint{1.404358in}{1.433581in}}%
\pgfpathlineto{\pgfqpoint{1.442640in}{1.426652in}}%
\pgfpathclose%
\pgfusepath{fill}%
\end{pgfscope}%
\begin{pgfscope}%
\pgfpathrectangle{\pgfqpoint{0.407592in}{0.222904in}}{\pgfqpoint{1.411163in}{1.411163in}}%
\pgfusepath{clip}%
\pgfsetbuttcap%
\pgfsetroundjoin%
\definecolor{currentfill}{rgb}{0.749020,0.772549,0.792157}%
\pgfsetfillcolor{currentfill}%
\pgfsetfillopacity{0.160000}%
\pgfsetlinewidth{0.000000pt}%
\definecolor{currentstroke}{rgb}{0.000000,0.000000,0.000000}%
\pgfsetstrokecolor{currentstroke}%
\pgfsetdash{}{0pt}%
\pgfpathmoveto{\pgfqpoint{1.254923in}{1.324960in}}%
\pgfpathlineto{\pgfqpoint{1.258516in}{1.340548in}}%
\pgfpathlineto{\pgfqpoint{1.262103in}{1.356108in}}%
\pgfpathlineto{\pgfqpoint{1.224927in}{1.359086in}}%
\pgfpathlineto{\pgfqpoint{1.187264in}{1.361193in}}%
\pgfpathlineto{\pgfqpoint{1.185242in}{1.345554in}}%
\pgfpathlineto{\pgfqpoint{1.183217in}{1.329886in}}%
\pgfpathlineto{\pgfqpoint{1.219305in}{1.327844in}}%
\pgfpathlineto{\pgfqpoint{1.254923in}{1.324960in}}%
\pgfpathclose%
\pgfusepath{fill}%
\end{pgfscope}%
\begin{pgfscope}%
\pgfpathrectangle{\pgfqpoint{0.407592in}{0.222904in}}{\pgfqpoint{1.411163in}{1.411163in}}%
\pgfusepath{clip}%
\pgfsetbuttcap%
\pgfsetroundjoin%
\definecolor{currentfill}{rgb}{0.749020,0.772549,0.792157}%
\pgfsetfillcolor{currentfill}%
\pgfsetfillopacity{0.160000}%
\pgfsetlinewidth{0.000000pt}%
\definecolor{currentstroke}{rgb}{0.000000,0.000000,0.000000}%
\pgfsetstrokecolor{currentstroke}%
\pgfsetdash{}{0pt}%
\pgfpathmoveto{\pgfqpoint{0.858517in}{1.438471in}}%
\pgfpathlineto{\pgfqpoint{0.854253in}{1.453708in}}%
\pgfpathlineto{\pgfqpoint{0.849995in}{1.468920in}}%
\pgfpathlineto{\pgfqpoint{0.807784in}{1.462554in}}%
\pgfpathlineto{\pgfqpoint{0.766500in}{1.455264in}}%
\pgfpathlineto{\pgfqpoint{0.772266in}{1.440220in}}%
\pgfpathlineto{\pgfqpoint{0.778039in}{1.425154in}}%
\pgfpathlineto{\pgfqpoint{0.817826in}{1.432263in}}%
\pgfpathlineto{\pgfqpoint{0.858517in}{1.438471in}}%
\pgfpathclose%
\pgfusepath{fill}%
\end{pgfscope}%
\begin{pgfscope}%
\pgfpathrectangle{\pgfqpoint{0.407592in}{0.222904in}}{\pgfqpoint{1.411163in}{1.411163in}}%
\pgfusepath{clip}%
\pgfsetbuttcap%
\pgfsetroundjoin%
\definecolor{currentfill}{rgb}{0.749020,0.772549,0.792157}%
\pgfsetfillcolor{currentfill}%
\pgfsetfillopacity{0.160000}%
\pgfsetlinewidth{0.000000pt}%
\definecolor{currentstroke}{rgb}{0.000000,0.000000,0.000000}%
\pgfsetstrokecolor{currentstroke}%
\pgfsetdash{}{0pt}%
\pgfpathmoveto{\pgfqpoint{1.036709in}{1.329544in}}%
\pgfpathlineto{\pgfqpoint{1.035543in}{1.345206in}}%
\pgfpathlineto{\pgfqpoint{1.034380in}{1.360840in}}%
\pgfpathlineto{\pgfqpoint{0.996258in}{1.358558in}}%
\pgfpathlineto{\pgfqpoint{0.958441in}{1.355406in}}%
\pgfpathlineto{\pgfqpoint{0.961192in}{1.339857in}}%
\pgfpathlineto{\pgfqpoint{0.963948in}{1.324280in}}%
\pgfpathlineto{\pgfqpoint{1.000181in}{1.327333in}}%
\pgfpathlineto{\pgfqpoint{1.036709in}{1.329544in}}%
\pgfpathclose%
\pgfusepath{fill}%
\end{pgfscope}%
\begin{pgfscope}%
\pgfpathrectangle{\pgfqpoint{0.407592in}{0.222904in}}{\pgfqpoint{1.411163in}{1.411163in}}%
\pgfusepath{clip}%
\pgfsetbuttcap%
\pgfsetroundjoin%
\definecolor{currentfill}{rgb}{0.749020,0.772549,0.792157}%
\pgfsetfillcolor{currentfill}%
\pgfsetfillopacity{0.160000}%
\pgfsetlinewidth{0.000000pt}%
\definecolor{currentstroke}{rgb}{0.000000,0.000000,0.000000}%
\pgfsetstrokecolor{currentstroke}%
\pgfsetdash{}{0pt}%
\pgfpathmoveto{\pgfqpoint{1.334479in}{1.347573in}}%
\pgfpathlineto{\pgfqpoint{1.339583in}{1.362975in}}%
\pgfpathlineto{\pgfqpoint{1.344679in}{1.378352in}}%
\pgfpathlineto{\pgfqpoint{1.307354in}{1.383189in}}%
\pgfpathlineto{\pgfqpoint{1.269258in}{1.387146in}}%
\pgfpathlineto{\pgfqpoint{1.265684in}{1.371641in}}%
\pgfpathlineto{\pgfqpoint{1.262103in}{1.356108in}}%
\pgfpathlineto{\pgfqpoint{1.298664in}{1.352267in}}%
\pgfpathlineto{\pgfqpoint{1.334479in}{1.347573in}}%
\pgfpathclose%
\pgfusepath{fill}%
\end{pgfscope}%
\begin{pgfscope}%
\pgfpathrectangle{\pgfqpoint{0.407592in}{0.222904in}}{\pgfqpoint{1.411163in}{1.411163in}}%
\pgfusepath{clip}%
\pgfsetbuttcap%
\pgfsetroundjoin%
\definecolor{currentfill}{rgb}{0.749020,0.772549,0.792157}%
\pgfsetfillcolor{currentfill}%
\pgfsetfillopacity{0.160000}%
\pgfsetlinewidth{0.000000pt}%
\definecolor{currentstroke}{rgb}{0.000000,0.000000,0.000000}%
\pgfsetstrokecolor{currentstroke}%
\pgfsetdash{}{0pt}%
\pgfpathmoveto{\pgfqpoint{1.179155in}{1.298463in}}%
\pgfpathlineto{\pgfqpoint{1.181188in}{1.314189in}}%
\pgfpathlineto{\pgfqpoint{1.183217in}{1.329886in}}%
\pgfpathlineto{\pgfqpoint{1.146786in}{1.331078in}}%
\pgfpathlineto{\pgfqpoint{1.110139in}{1.331419in}}%
\pgfpathlineto{\pgfqpoint{1.109705in}{1.315697in}}%
\pgfpathlineto{\pgfqpoint{1.109270in}{1.299945in}}%
\pgfpathlineto{\pgfqpoint{1.144317in}{1.299616in}}%
\pgfpathlineto{\pgfqpoint{1.179155in}{1.298463in}}%
\pgfpathclose%
\pgfusepath{fill}%
\end{pgfscope}%
\begin{pgfscope}%
\pgfpathrectangle{\pgfqpoint{0.407592in}{0.222904in}}{\pgfqpoint{1.411163in}{1.411163in}}%
\pgfusepath{clip}%
\pgfsetbuttcap%
\pgfsetroundjoin%
\definecolor{currentfill}{rgb}{0.749020,0.772549,0.792157}%
\pgfsetfillcolor{currentfill}%
\pgfsetfillopacity{0.160000}%
\pgfsetlinewidth{0.000000pt}%
\definecolor{currentstroke}{rgb}{0.000000,0.000000,0.000000}%
\pgfsetstrokecolor{currentstroke}%
\pgfsetdash{}{0pt}%
\pgfpathmoveto{\pgfqpoint{1.109270in}{1.299945in}}%
\pgfpathlineto{\pgfqpoint{1.109705in}{1.315697in}}%
\pgfpathlineto{\pgfqpoint{1.110139in}{1.331419in}}%
\pgfpathlineto{\pgfqpoint{1.073405in}{1.330907in}}%
\pgfpathlineto{\pgfqpoint{1.036709in}{1.329544in}}%
\pgfpathlineto{\pgfqpoint{1.037877in}{1.313853in}}%
\pgfpathlineto{\pgfqpoint{1.039047in}{1.298132in}}%
\pgfpathlineto{\pgfqpoint{1.074139in}{1.299451in}}%
\pgfpathlineto{\pgfqpoint{1.109270in}{1.299945in}}%
\pgfpathclose%
\pgfusepath{fill}%
\end{pgfscope}%
\begin{pgfscope}%
\pgfpathrectangle{\pgfqpoint{0.407592in}{0.222904in}}{\pgfqpoint{1.411163in}{1.411163in}}%
\pgfusepath{clip}%
\pgfsetbuttcap%
\pgfsetroundjoin%
\definecolor{currentfill}{rgb}{0.749020,0.772549,0.792157}%
\pgfsetfillcolor{currentfill}%
\pgfsetfillopacity{0.160000}%
\pgfsetlinewidth{0.000000pt}%
\definecolor{currentstroke}{rgb}{0.000000,0.000000,0.000000}%
\pgfsetstrokecolor{currentstroke}%
\pgfsetdash{}{0pt}%
\pgfpathmoveto{\pgfqpoint{0.958441in}{1.355406in}}%
\pgfpathlineto{\pgfqpoint{0.955695in}{1.370928in}}%
\pgfpathlineto{\pgfqpoint{0.952954in}{1.386423in}}%
\pgfpathlineto{\pgfqpoint{0.914006in}{1.382289in}}%
\pgfpathlineto{\pgfqpoint{0.875643in}{1.377277in}}%
\pgfpathlineto{\pgfqpoint{0.879942in}{1.361916in}}%
\pgfpathlineto{\pgfqpoint{0.884248in}{1.346529in}}%
\pgfpathlineto{\pgfqpoint{0.921061in}{1.351393in}}%
\pgfpathlineto{\pgfqpoint{0.958441in}{1.355406in}}%
\pgfpathclose%
\pgfusepath{fill}%
\end{pgfscope}%
\begin{pgfscope}%
\pgfpathrectangle{\pgfqpoint{0.407592in}{0.222904in}}{\pgfqpoint{1.411163in}{1.411163in}}%
\pgfusepath{clip}%
\pgfsetbuttcap%
\pgfsetroundjoin%
\definecolor{currentfill}{rgb}{0.749020,0.772549,0.792157}%
\pgfsetfillcolor{currentfill}%
\pgfsetfillopacity{0.160000}%
\pgfsetlinewidth{0.000000pt}%
\definecolor{currentstroke}{rgb}{0.000000,0.000000,0.000000}%
\pgfsetstrokecolor{currentstroke}%
\pgfsetdash{}{0pt}%
\pgfpathmoveto{\pgfqpoint{1.429584in}{1.396416in}}%
\pgfpathlineto{\pgfqpoint{1.436117in}{1.411545in}}%
\pgfpathlineto{\pgfqpoint{1.442640in}{1.426652in}}%
\pgfpathlineto{\pgfqpoint{1.404358in}{1.433581in}}%
\pgfpathlineto{\pgfqpoint{1.364979in}{1.439607in}}%
\pgfpathlineto{\pgfqpoint{1.359916in}{1.424331in}}%
\pgfpathlineto{\pgfqpoint{1.354846in}{1.409030in}}%
\pgfpathlineto{\pgfqpoint{1.392748in}{1.403162in}}%
\pgfpathlineto{\pgfqpoint{1.429584in}{1.396416in}}%
\pgfpathclose%
\pgfusepath{fill}%
\end{pgfscope}%
\begin{pgfscope}%
\pgfpathrectangle{\pgfqpoint{0.407592in}{0.222904in}}{\pgfqpoint{1.411163in}{1.411163in}}%
\pgfusepath{clip}%
\pgfsetbuttcap%
\pgfsetroundjoin%
\definecolor{currentfill}{rgb}{0.749020,0.772549,0.792157}%
\pgfsetfillcolor{currentfill}%
\pgfsetfillopacity{0.160000}%
\pgfsetlinewidth{0.000000pt}%
\definecolor{currentstroke}{rgb}{0.000000,0.000000,0.000000}%
\pgfsetstrokecolor{currentstroke}%
\pgfsetdash{}{0pt}%
\pgfpathmoveto{\pgfqpoint{0.867066in}{1.407924in}}%
\pgfpathlineto{\pgfqpoint{0.862788in}{1.423210in}}%
\pgfpathlineto{\pgfqpoint{0.858517in}{1.438471in}}%
\pgfpathlineto{\pgfqpoint{0.817826in}{1.432263in}}%
\pgfpathlineto{\pgfqpoint{0.778039in}{1.425154in}}%
\pgfpathlineto{\pgfqpoint{0.783821in}{1.410067in}}%
\pgfpathlineto{\pgfqpoint{0.789612in}{1.394958in}}%
\pgfpathlineto{\pgfqpoint{0.827898in}{1.401879in}}%
\pgfpathlineto{\pgfqpoint{0.867066in}{1.407924in}}%
\pgfpathclose%
\pgfusepath{fill}%
\end{pgfscope}%
\begin{pgfscope}%
\pgfpathrectangle{\pgfqpoint{0.407592in}{0.222904in}}{\pgfqpoint{1.411163in}{1.411163in}}%
\pgfusepath{clip}%
\pgfsetbuttcap%
\pgfsetroundjoin%
\definecolor{currentfill}{rgb}{0.749020,0.772549,0.792157}%
\pgfsetfillcolor{currentfill}%
\pgfsetfillopacity{0.160000}%
\pgfsetlinewidth{0.000000pt}%
\definecolor{currentstroke}{rgb}{0.000000,0.000000,0.000000}%
\pgfsetstrokecolor{currentstroke}%
\pgfsetdash{}{0pt}%
\pgfpathmoveto{\pgfqpoint{1.247717in}{1.293700in}}%
\pgfpathlineto{\pgfqpoint{1.251323in}{1.309344in}}%
\pgfpathlineto{\pgfqpoint{1.254923in}{1.324960in}}%
\pgfpathlineto{\pgfqpoint{1.219305in}{1.327844in}}%
\pgfpathlineto{\pgfqpoint{1.183217in}{1.329886in}}%
\pgfpathlineto{\pgfqpoint{1.181188in}{1.314189in}}%
\pgfpathlineto{\pgfqpoint{1.179155in}{1.298463in}}%
\pgfpathlineto{\pgfqpoint{1.213663in}{1.296489in}}%
\pgfpathlineto{\pgfqpoint{1.247717in}{1.293700in}}%
\pgfpathclose%
\pgfusepath{fill}%
\end{pgfscope}%
\begin{pgfscope}%
\pgfpathrectangle{\pgfqpoint{0.407592in}{0.222904in}}{\pgfqpoint{1.411163in}{1.411163in}}%
\pgfusepath{clip}%
\pgfsetbuttcap%
\pgfsetroundjoin%
\definecolor{currentfill}{rgb}{0.749020,0.772549,0.792157}%
\pgfsetfillcolor{currentfill}%
\pgfsetfillopacity{0.160000}%
\pgfsetlinewidth{0.000000pt}%
\definecolor{currentstroke}{rgb}{0.000000,0.000000,0.000000}%
\pgfsetstrokecolor{currentstroke}%
\pgfsetdash{}{0pt}%
\pgfpathmoveto{\pgfqpoint{1.039047in}{1.298132in}}%
\pgfpathlineto{\pgfqpoint{1.037877in}{1.313853in}}%
\pgfpathlineto{\pgfqpoint{1.036709in}{1.329544in}}%
\pgfpathlineto{\pgfqpoint{1.000181in}{1.327333in}}%
\pgfpathlineto{\pgfqpoint{0.963948in}{1.324280in}}%
\pgfpathlineto{\pgfqpoint{0.966709in}{1.308675in}}%
\pgfpathlineto{\pgfqpoint{0.969475in}{1.293043in}}%
\pgfpathlineto{\pgfqpoint{1.004118in}{1.295995in}}%
\pgfpathlineto{\pgfqpoint{1.039047in}{1.298132in}}%
\pgfpathclose%
\pgfusepath{fill}%
\end{pgfscope}%
\begin{pgfscope}%
\pgfpathrectangle{\pgfqpoint{0.407592in}{0.222904in}}{\pgfqpoint{1.411163in}{1.411163in}}%
\pgfusepath{clip}%
\pgfsetbuttcap%
\pgfsetroundjoin%
\definecolor{currentfill}{rgb}{0.749020,0.772549,0.792157}%
\pgfsetfillcolor{currentfill}%
\pgfsetfillopacity{0.160000}%
\pgfsetlinewidth{0.000000pt}%
\definecolor{currentstroke}{rgb}{0.000000,0.000000,0.000000}%
\pgfsetstrokecolor{currentstroke}%
\pgfsetdash{}{0pt}%
\pgfpathmoveto{\pgfqpoint{1.324246in}{1.316692in}}%
\pgfpathlineto{\pgfqpoint{1.329367in}{1.332145in}}%
\pgfpathlineto{\pgfqpoint{1.334479in}{1.347573in}}%
\pgfpathlineto{\pgfqpoint{1.298664in}{1.352267in}}%
\pgfpathlineto{\pgfqpoint{1.262103in}{1.356108in}}%
\pgfpathlineto{\pgfqpoint{1.258516in}{1.340548in}}%
\pgfpathlineto{\pgfqpoint{1.254923in}{1.324960in}}%
\pgfpathlineto{\pgfqpoint{1.289945in}{1.321239in}}%
\pgfpathlineto{\pgfqpoint{1.324246in}{1.316692in}}%
\pgfpathclose%
\pgfusepath{fill}%
\end{pgfscope}%
\begin{pgfscope}%
\pgfpathrectangle{\pgfqpoint{0.407592in}{0.222904in}}{\pgfqpoint{1.411163in}{1.411163in}}%
\pgfusepath{clip}%
\pgfsetbuttcap%
\pgfsetroundjoin%
\definecolor{currentfill}{rgb}{0.749020,0.772549,0.792157}%
\pgfsetfillcolor{currentfill}%
\pgfsetfillopacity{0.160000}%
\pgfsetlinewidth{0.000000pt}%
\definecolor{currentstroke}{rgb}{0.000000,0.000000,0.000000}%
\pgfsetstrokecolor{currentstroke}%
\pgfsetdash{}{0pt}%
\pgfpathmoveto{\pgfqpoint{1.175078in}{1.266923in}}%
\pgfpathlineto{\pgfqpoint{1.177118in}{1.282707in}}%
\pgfpathlineto{\pgfqpoint{1.179155in}{1.298463in}}%
\pgfpathlineto{\pgfqpoint{1.144317in}{1.299616in}}%
\pgfpathlineto{\pgfqpoint{1.109270in}{1.299945in}}%
\pgfpathlineto{\pgfqpoint{1.108834in}{1.284164in}}%
\pgfpathlineto{\pgfqpoint{1.108398in}{1.268353in}}%
\pgfpathlineto{\pgfqpoint{1.141838in}{1.268035in}}%
\pgfpathlineto{\pgfqpoint{1.175078in}{1.266923in}}%
\pgfpathclose%
\pgfusepath{fill}%
\end{pgfscope}%
\begin{pgfscope}%
\pgfpathrectangle{\pgfqpoint{0.407592in}{0.222904in}}{\pgfqpoint{1.411163in}{1.411163in}}%
\pgfusepath{clip}%
\pgfsetbuttcap%
\pgfsetroundjoin%
\definecolor{currentfill}{rgb}{0.749020,0.772549,0.792157}%
\pgfsetfillcolor{currentfill}%
\pgfsetfillopacity{0.160000}%
\pgfsetlinewidth{0.000000pt}%
\definecolor{currentstroke}{rgb}{0.000000,0.000000,0.000000}%
\pgfsetstrokecolor{currentstroke}%
\pgfsetdash{}{0pt}%
\pgfpathmoveto{\pgfqpoint{1.108398in}{1.268353in}}%
\pgfpathlineto{\pgfqpoint{1.108834in}{1.284164in}}%
\pgfpathlineto{\pgfqpoint{1.109270in}{1.299945in}}%
\pgfpathlineto{\pgfqpoint{1.074139in}{1.299451in}}%
\pgfpathlineto{\pgfqpoint{1.039047in}{1.298132in}}%
\pgfpathlineto{\pgfqpoint{1.040220in}{1.282383in}}%
\pgfpathlineto{\pgfqpoint{1.041394in}{1.266604in}}%
\pgfpathlineto{\pgfqpoint{1.074877in}{1.267876in}}%
\pgfpathlineto{\pgfqpoint{1.108398in}{1.268353in}}%
\pgfpathclose%
\pgfusepath{fill}%
\end{pgfscope}%
\begin{pgfscope}%
\pgfpathrectangle{\pgfqpoint{0.407592in}{0.222904in}}{\pgfqpoint{1.411163in}{1.411163in}}%
\pgfusepath{clip}%
\pgfsetbuttcap%
\pgfsetroundjoin%
\definecolor{currentfill}{rgb}{0.749020,0.772549,0.792157}%
\pgfsetfillcolor{currentfill}%
\pgfsetfillopacity{0.160000}%
\pgfsetlinewidth{0.000000pt}%
\definecolor{currentstroke}{rgb}{0.000000,0.000000,0.000000}%
\pgfsetstrokecolor{currentstroke}%
\pgfsetdash{}{0pt}%
\pgfpathmoveto{\pgfqpoint{0.963948in}{1.324280in}}%
\pgfpathlineto{\pgfqpoint{0.961192in}{1.339857in}}%
\pgfpathlineto{\pgfqpoint{0.958441in}{1.355406in}}%
\pgfpathlineto{\pgfqpoint{0.921061in}{1.351393in}}%
\pgfpathlineto{\pgfqpoint{0.884248in}{1.346529in}}%
\pgfpathlineto{\pgfqpoint{0.888561in}{1.331118in}}%
\pgfpathlineto{\pgfqpoint{0.892882in}{1.315681in}}%
\pgfpathlineto{\pgfqpoint{0.928139in}{1.320393in}}%
\pgfpathlineto{\pgfqpoint{0.963948in}{1.324280in}}%
\pgfpathclose%
\pgfusepath{fill}%
\end{pgfscope}%
\begin{pgfscope}%
\pgfpathrectangle{\pgfqpoint{0.407592in}{0.222904in}}{\pgfqpoint{1.411163in}{1.411163in}}%
\pgfusepath{clip}%
\pgfsetbuttcap%
\pgfsetroundjoin%
\definecolor{currentfill}{rgb}{0.749020,0.772549,0.792157}%
\pgfsetfillcolor{currentfill}%
\pgfsetfillopacity{0.160000}%
\pgfsetlinewidth{0.000000pt}%
\definecolor{currentstroke}{rgb}{0.000000,0.000000,0.000000}%
\pgfsetstrokecolor{currentstroke}%
\pgfsetdash{}{0pt}%
\pgfpathmoveto{\pgfqpoint{1.546915in}{1.469523in}}%
\pgfpathlineto{\pgfqpoint{1.554781in}{1.484319in}}%
\pgfpathlineto{\pgfqpoint{1.515683in}{1.493550in}}%
\pgfpathlineto{\pgfqpoint{1.475114in}{1.501856in}}%
\pgfpathlineto{\pgfqpoint{1.468638in}{1.486859in}}%
\pgfpathlineto{\pgfqpoint{1.508503in}{1.478648in}}%
\pgfusepath{fill}%
\end{pgfscope}%
\begin{pgfscope}%
\pgfpathrectangle{\pgfqpoint{0.407592in}{0.222904in}}{\pgfqpoint{1.411163in}{1.411163in}}%
\pgfusepath{clip}%
\pgfsetbuttcap%
\pgfsetroundjoin%
\definecolor{currentfill}{rgb}{0.749020,0.772549,0.792157}%
\pgfsetfillcolor{currentfill}%
\pgfsetfillopacity{0.160000}%
\pgfsetlinewidth{0.000000pt}%
\definecolor{currentstroke}{rgb}{0.000000,0.000000,0.000000}%
\pgfsetstrokecolor{currentstroke}%
\pgfsetdash{}{0pt}%
\pgfpathmoveto{\pgfqpoint{1.416489in}{1.366090in}}%
\pgfpathlineto{\pgfqpoint{1.423041in}{1.381264in}}%
\pgfpathlineto{\pgfqpoint{1.429584in}{1.396416in}}%
\pgfpathlineto{\pgfqpoint{1.392748in}{1.403162in}}%
\pgfpathlineto{\pgfqpoint{1.354846in}{1.409030in}}%
\pgfpathlineto{\pgfqpoint{1.349767in}{1.393704in}}%
\pgfpathlineto{\pgfqpoint{1.344679in}{1.378352in}}%
\pgfpathlineto{\pgfqpoint{1.381102in}{1.372648in}}%
\pgfpathlineto{\pgfqpoint{1.416489in}{1.366090in}}%
\pgfpathclose%
\pgfusepath{fill}%
\end{pgfscope}%
\begin{pgfscope}%
\pgfpathrectangle{\pgfqpoint{0.407592in}{0.222904in}}{\pgfqpoint{1.411163in}{1.411163in}}%
\pgfusepath{clip}%
\pgfsetbuttcap%
\pgfsetroundjoin%
\definecolor{currentfill}{rgb}{0.749020,0.772549,0.792157}%
\pgfsetfillcolor{currentfill}%
\pgfsetfillopacity{0.160000}%
\pgfsetlinewidth{0.000000pt}%
\definecolor{currentstroke}{rgb}{0.000000,0.000000,0.000000}%
\pgfsetstrokecolor{currentstroke}%
\pgfsetdash{}{0pt}%
\pgfpathmoveto{\pgfqpoint{0.875643in}{1.377277in}}%
\pgfpathlineto{\pgfqpoint{0.871351in}{1.392613in}}%
\pgfpathlineto{\pgfqpoint{0.867066in}{1.407924in}}%
\pgfpathlineto{\pgfqpoint{0.827898in}{1.401879in}}%
\pgfpathlineto{\pgfqpoint{0.789612in}{1.394958in}}%
\pgfpathlineto{\pgfqpoint{0.795410in}{1.379827in}}%
\pgfpathlineto{\pgfqpoint{0.801217in}{1.364674in}}%
\pgfpathlineto{\pgfqpoint{0.838002in}{1.371400in}}%
\pgfpathlineto{\pgfqpoint{0.875643in}{1.377277in}}%
\pgfpathclose%
\pgfusepath{fill}%
\end{pgfscope}%
\begin{pgfscope}%
\pgfpathrectangle{\pgfqpoint{0.407592in}{0.222904in}}{\pgfqpoint{1.411163in}{1.411163in}}%
\pgfusepath{clip}%
\pgfsetbuttcap%
\pgfsetroundjoin%
\definecolor{currentfill}{rgb}{0.749020,0.772549,0.792157}%
\pgfsetfillcolor{currentfill}%
\pgfsetfillopacity{0.160000}%
\pgfsetlinewidth{0.000000pt}%
\definecolor{currentstroke}{rgb}{0.000000,0.000000,0.000000}%
\pgfsetstrokecolor{currentstroke}%
\pgfsetdash{}{0pt}%
\pgfpathmoveto{\pgfqpoint{0.754995in}{1.485286in}}%
\pgfpathlineto{\pgfqpoint{0.749254in}{1.500265in}}%
\pgfpathlineto{\pgfqpoint{0.706844in}{1.491772in}}%
\pgfpathlineto{\pgfqpoint{0.665703in}{1.482357in}}%
\pgfpathlineto{\pgfqpoint{0.672897in}{1.467584in}}%
\pgfpathlineto{\pgfqpoint{0.713318in}{1.476890in}}%
\pgfusepath{fill}%
\end{pgfscope}%
\begin{pgfscope}%
\pgfpathrectangle{\pgfqpoint{0.407592in}{0.222904in}}{\pgfqpoint{1.411163in}{1.411163in}}%
\pgfusepath{clip}%
\pgfsetbuttcap%
\pgfsetroundjoin%
\definecolor{currentfill}{rgb}{0.749020,0.772549,0.792157}%
\pgfsetfillcolor{currentfill}%
\pgfsetfillopacity{0.160000}%
\pgfsetlinewidth{0.000000pt}%
\definecolor{currentstroke}{rgb}{0.000000,0.000000,0.000000}%
\pgfsetstrokecolor{currentstroke}%
\pgfsetdash{}{0pt}%
\pgfpathmoveto{\pgfqpoint{1.240485in}{1.262329in}}%
\pgfpathlineto{\pgfqpoint{1.244104in}{1.278029in}}%
\pgfpathlineto{\pgfqpoint{1.247717in}{1.293700in}}%
\pgfpathlineto{\pgfqpoint{1.213663in}{1.296489in}}%
\pgfpathlineto{\pgfqpoint{1.179155in}{1.298463in}}%
\pgfpathlineto{\pgfqpoint{1.177118in}{1.282707in}}%
\pgfpathlineto{\pgfqpoint{1.175078in}{1.266923in}}%
\pgfpathlineto{\pgfqpoint{1.208000in}{1.265019in}}%
\pgfpathlineto{\pgfqpoint{1.240485in}{1.262329in}}%
\pgfpathclose%
\pgfusepath{fill}%
\end{pgfscope}%
\begin{pgfscope}%
\pgfpathrectangle{\pgfqpoint{0.407592in}{0.222904in}}{\pgfqpoint{1.411163in}{1.411163in}}%
\pgfusepath{clip}%
\pgfsetbuttcap%
\pgfsetroundjoin%
\definecolor{currentfill}{rgb}{0.749020,0.772549,0.792157}%
\pgfsetfillcolor{currentfill}%
\pgfsetfillopacity{0.160000}%
\pgfsetlinewidth{0.000000pt}%
\definecolor{currentstroke}{rgb}{0.000000,0.000000,0.000000}%
\pgfsetstrokecolor{currentstroke}%
\pgfsetdash{}{0pt}%
\pgfpathmoveto{\pgfqpoint{1.531153in}{1.439875in}}%
\pgfpathlineto{\pgfqpoint{1.539039in}{1.454708in}}%
\pgfpathlineto{\pgfqpoint{1.546915in}{1.469523in}}%
\pgfpathlineto{\pgfqpoint{1.508503in}{1.478648in}}%
\pgfpathlineto{\pgfqpoint{1.468638in}{1.486859in}}%
\pgfpathlineto{\pgfqpoint{1.462153in}{1.471840in}}%
\pgfpathlineto{\pgfqpoint{1.455658in}{1.456799in}}%
\pgfpathlineto{\pgfqpoint{1.494113in}{1.448783in}}%
\pgfpathlineto{\pgfqpoint{1.531153in}{1.439875in}}%
\pgfpathclose%
\pgfusepath{fill}%
\end{pgfscope}%
\begin{pgfscope}%
\pgfpathrectangle{\pgfqpoint{0.407592in}{0.222904in}}{\pgfqpoint{1.411163in}{1.411163in}}%
\pgfusepath{clip}%
\pgfsetbuttcap%
\pgfsetroundjoin%
\definecolor{currentfill}{rgb}{0.749020,0.772549,0.792157}%
\pgfsetfillcolor{currentfill}%
\pgfsetfillopacity{0.160000}%
\pgfsetlinewidth{0.000000pt}%
\definecolor{currentstroke}{rgb}{0.000000,0.000000,0.000000}%
\pgfsetstrokecolor{currentstroke}%
\pgfsetdash{}{0pt}%
\pgfpathmoveto{\pgfqpoint{1.041394in}{1.266604in}}%
\pgfpathlineto{\pgfqpoint{1.040220in}{1.282383in}}%
\pgfpathlineto{\pgfqpoint{1.039047in}{1.298132in}}%
\pgfpathlineto{\pgfqpoint{1.004118in}{1.295995in}}%
\pgfpathlineto{\pgfqpoint{0.969475in}{1.293043in}}%
\pgfpathlineto{\pgfqpoint{0.972246in}{1.277383in}}%
\pgfpathlineto{\pgfqpoint{0.975021in}{1.261696in}}%
\pgfpathlineto{\pgfqpoint{1.008069in}{1.264542in}}%
\pgfpathlineto{\pgfqpoint{1.041394in}{1.266604in}}%
\pgfpathclose%
\pgfusepath{fill}%
\end{pgfscope}%
\begin{pgfscope}%
\pgfpathrectangle{\pgfqpoint{0.407592in}{0.222904in}}{\pgfqpoint{1.411163in}{1.411163in}}%
\pgfusepath{clip}%
\pgfsetbuttcap%
\pgfsetroundjoin%
\definecolor{currentfill}{rgb}{0.749020,0.772549,0.792157}%
\pgfsetfillcolor{currentfill}%
\pgfsetfillopacity{0.160000}%
\pgfsetlinewidth{0.000000pt}%
\definecolor{currentstroke}{rgb}{0.000000,0.000000,0.000000}%
\pgfsetstrokecolor{currentstroke}%
\pgfsetdash{}{0pt}%
\pgfpathmoveto{\pgfqpoint{1.313978in}{1.285709in}}%
\pgfpathlineto{\pgfqpoint{1.319116in}{1.301213in}}%
\pgfpathlineto{\pgfqpoint{1.324246in}{1.316692in}}%
\pgfpathlineto{\pgfqpoint{1.289945in}{1.321239in}}%
\pgfpathlineto{\pgfqpoint{1.254923in}{1.324960in}}%
\pgfpathlineto{\pgfqpoint{1.251323in}{1.309344in}}%
\pgfpathlineto{\pgfqpoint{1.247717in}{1.293700in}}%
\pgfpathlineto{\pgfqpoint{1.281196in}{1.290104in}}%
\pgfpathlineto{\pgfqpoint{1.313978in}{1.285709in}}%
\pgfpathclose%
\pgfusepath{fill}%
\end{pgfscope}%
\begin{pgfscope}%
\pgfpathrectangle{\pgfqpoint{0.407592in}{0.222904in}}{\pgfqpoint{1.411163in}{1.411163in}}%
\pgfusepath{clip}%
\pgfsetbuttcap%
\pgfsetroundjoin%
\definecolor{currentfill}{rgb}{0.749020,0.772549,0.792157}%
\pgfsetfillcolor{currentfill}%
\pgfsetfillopacity{0.160000}%
\pgfsetlinewidth{0.000000pt}%
\definecolor{currentstroke}{rgb}{0.000000,0.000000,0.000000}%
\pgfsetstrokecolor{currentstroke}%
\pgfsetdash{}{0pt}%
\pgfpathmoveto{\pgfqpoint{1.403355in}{1.335676in}}%
\pgfpathlineto{\pgfqpoint{1.409927in}{1.350894in}}%
\pgfpathlineto{\pgfqpoint{1.416489in}{1.366090in}}%
\pgfpathlineto{\pgfqpoint{1.381102in}{1.372648in}}%
\pgfpathlineto{\pgfqpoint{1.344679in}{1.378352in}}%
\pgfpathlineto{\pgfqpoint{1.339583in}{1.362975in}}%
\pgfpathlineto{\pgfqpoint{1.334479in}{1.347573in}}%
\pgfpathlineto{\pgfqpoint{1.369419in}{1.342038in}}%
\pgfpathlineto{\pgfqpoint{1.403355in}{1.335676in}}%
\pgfpathclose%
\pgfusepath{fill}%
\end{pgfscope}%
\begin{pgfscope}%
\pgfpathrectangle{\pgfqpoint{0.407592in}{0.222904in}}{\pgfqpoint{1.411163in}{1.411163in}}%
\pgfusepath{clip}%
\pgfsetbuttcap%
\pgfsetroundjoin%
\definecolor{currentfill}{rgb}{0.749020,0.772549,0.792157}%
\pgfsetfillcolor{currentfill}%
\pgfsetfillopacity{0.160000}%
\pgfsetlinewidth{0.000000pt}%
\definecolor{currentstroke}{rgb}{0.000000,0.000000,0.000000}%
\pgfsetstrokecolor{currentstroke}%
\pgfsetdash{}{0pt}%
\pgfpathmoveto{\pgfqpoint{0.969475in}{1.293043in}}%
\pgfpathlineto{\pgfqpoint{0.966709in}{1.308675in}}%
\pgfpathlineto{\pgfqpoint{0.963948in}{1.324280in}}%
\pgfpathlineto{\pgfqpoint{0.928139in}{1.320393in}}%
\pgfpathlineto{\pgfqpoint{0.892882in}{1.315681in}}%
\pgfpathlineto{\pgfqpoint{0.897209in}{1.300220in}}%
\pgfpathlineto{\pgfqpoint{0.901543in}{1.284732in}}%
\pgfpathlineto{\pgfqpoint{0.935242in}{1.289286in}}%
\pgfpathlineto{\pgfqpoint{0.969475in}{1.293043in}}%
\pgfpathclose%
\pgfusepath{fill}%
\end{pgfscope}%
\begin{pgfscope}%
\pgfpathrectangle{\pgfqpoint{0.407592in}{0.222904in}}{\pgfqpoint{1.411163in}{1.411163in}}%
\pgfusepath{clip}%
\pgfsetbuttcap%
\pgfsetroundjoin%
\definecolor{currentfill}{rgb}{0.749020,0.772549,0.792157}%
\pgfsetfillcolor{currentfill}%
\pgfsetfillopacity{0.160000}%
\pgfsetlinewidth{0.000000pt}%
\definecolor{currentstroke}{rgb}{0.000000,0.000000,0.000000}%
\pgfsetstrokecolor{currentstroke}%
\pgfsetdash{}{0pt}%
\pgfpathmoveto{\pgfqpoint{0.766500in}{1.455264in}}%
\pgfpathlineto{\pgfqpoint{0.760743in}{1.470286in}}%
\pgfpathlineto{\pgfqpoint{0.754995in}{1.485286in}}%
\pgfpathlineto{\pgfqpoint{0.713318in}{1.476890in}}%
\pgfpathlineto{\pgfqpoint{0.672897in}{1.467584in}}%
\pgfpathlineto{\pgfqpoint{0.680099in}{1.452792in}}%
\pgfpathlineto{\pgfqpoint{0.687311in}{1.437983in}}%
\pgfpathlineto{\pgfqpoint{0.726293in}{1.447066in}}%
\pgfpathlineto{\pgfqpoint{0.766500in}{1.455264in}}%
\pgfpathclose%
\pgfusepath{fill}%
\end{pgfscope}%
\begin{pgfscope}%
\pgfpathrectangle{\pgfqpoint{0.407592in}{0.222904in}}{\pgfqpoint{1.411163in}{1.411163in}}%
\pgfusepath{clip}%
\pgfsetbuttcap%
\pgfsetroundjoin%
\definecolor{currentfill}{rgb}{0.749020,0.772549,0.792157}%
\pgfsetfillcolor{currentfill}%
\pgfsetfillopacity{0.160000}%
\pgfsetlinewidth{0.000000pt}%
\definecolor{currentstroke}{rgb}{0.000000,0.000000,0.000000}%
\pgfsetstrokecolor{currentstroke}%
\pgfsetdash{}{0pt}%
\pgfpathmoveto{\pgfqpoint{1.170986in}{1.235265in}}%
\pgfpathlineto{\pgfqpoint{1.173034in}{1.251109in}}%
\pgfpathlineto{\pgfqpoint{1.175078in}{1.266923in}}%
\pgfpathlineto{\pgfqpoint{1.141838in}{1.268035in}}%
\pgfpathlineto{\pgfqpoint{1.108398in}{1.268353in}}%
\pgfpathlineto{\pgfqpoint{1.107960in}{1.252512in}}%
\pgfpathlineto{\pgfqpoint{1.107522in}{1.236641in}}%
\pgfpathlineto{\pgfqpoint{1.139350in}{1.236336in}}%
\pgfpathlineto{\pgfqpoint{1.170986in}{1.235265in}}%
\pgfpathclose%
\pgfusepath{fill}%
\end{pgfscope}%
\begin{pgfscope}%
\pgfpathrectangle{\pgfqpoint{0.407592in}{0.222904in}}{\pgfqpoint{1.411163in}{1.411163in}}%
\pgfusepath{clip}%
\pgfsetbuttcap%
\pgfsetroundjoin%
\definecolor{currentfill}{rgb}{0.749020,0.772549,0.792157}%
\pgfsetfillcolor{currentfill}%
\pgfsetfillopacity{0.160000}%
\pgfsetlinewidth{0.000000pt}%
\definecolor{currentstroke}{rgb}{0.000000,0.000000,0.000000}%
\pgfsetstrokecolor{currentstroke}%
\pgfsetdash{}{0pt}%
\pgfpathmoveto{\pgfqpoint{1.107522in}{1.236641in}}%
\pgfpathlineto{\pgfqpoint{1.107960in}{1.252512in}}%
\pgfpathlineto{\pgfqpoint{1.108398in}{1.268353in}}%
\pgfpathlineto{\pgfqpoint{1.074877in}{1.267876in}}%
\pgfpathlineto{\pgfqpoint{1.041394in}{1.266604in}}%
\pgfpathlineto{\pgfqpoint{1.042571in}{1.250796in}}%
\pgfpathlineto{\pgfqpoint{1.043750in}{1.234959in}}%
\pgfpathlineto{\pgfqpoint{1.075617in}{1.236182in}}%
\pgfpathlineto{\pgfqpoint{1.107522in}{1.236641in}}%
\pgfpathclose%
\pgfusepath{fill}%
\end{pgfscope}%
\begin{pgfscope}%
\pgfpathrectangle{\pgfqpoint{0.407592in}{0.222904in}}{\pgfqpoint{1.411163in}{1.411163in}}%
\pgfusepath{clip}%
\pgfsetbuttcap%
\pgfsetroundjoin%
\definecolor{currentfill}{rgb}{0.749020,0.772549,0.792157}%
\pgfsetfillcolor{currentfill}%
\pgfsetfillopacity{0.160000}%
\pgfsetlinewidth{0.000000pt}%
\definecolor{currentstroke}{rgb}{0.000000,0.000000,0.000000}%
\pgfsetstrokecolor{currentstroke}%
\pgfsetdash{}{0pt}%
\pgfpathmoveto{\pgfqpoint{0.884248in}{1.346529in}}%
\pgfpathlineto{\pgfqpoint{0.879942in}{1.361916in}}%
\pgfpathlineto{\pgfqpoint{0.875643in}{1.377277in}}%
\pgfpathlineto{\pgfqpoint{0.838002in}{1.371400in}}%
\pgfpathlineto{\pgfqpoint{0.801217in}{1.364674in}}%
\pgfpathlineto{\pgfqpoint{0.807033in}{1.349499in}}%
\pgfpathlineto{\pgfqpoint{0.812857in}{1.334301in}}%
\pgfpathlineto{\pgfqpoint{0.848137in}{1.340827in}}%
\pgfpathlineto{\pgfqpoint{0.884248in}{1.346529in}}%
\pgfpathclose%
\pgfusepath{fill}%
\end{pgfscope}%
\begin{pgfscope}%
\pgfpathrectangle{\pgfqpoint{0.407592in}{0.222904in}}{\pgfqpoint{1.411163in}{1.411163in}}%
\pgfusepath{clip}%
\pgfsetbuttcap%
\pgfsetroundjoin%
\definecolor{currentfill}{rgb}{0.749020,0.772549,0.792157}%
\pgfsetfillcolor{currentfill}%
\pgfsetfillopacity{0.160000}%
\pgfsetlinewidth{0.000000pt}%
\definecolor{currentstroke}{rgb}{0.000000,0.000000,0.000000}%
\pgfsetstrokecolor{currentstroke}%
\pgfsetdash{}{0pt}%
\pgfpathmoveto{\pgfqpoint{1.515351in}{1.410154in}}%
\pgfpathlineto{\pgfqpoint{1.523257in}{1.425024in}}%
\pgfpathlineto{\pgfqpoint{1.531153in}{1.439875in}}%
\pgfpathlineto{\pgfqpoint{1.494113in}{1.448783in}}%
\pgfpathlineto{\pgfqpoint{1.455658in}{1.456799in}}%
\pgfpathlineto{\pgfqpoint{1.449154in}{1.441737in}}%
\pgfpathlineto{\pgfqpoint{1.442640in}{1.426652in}}%
\pgfpathlineto{\pgfqpoint{1.479684in}{1.418836in}}%
\pgfpathlineto{\pgfqpoint{1.515351in}{1.410154in}}%
\pgfpathclose%
\pgfusepath{fill}%
\end{pgfscope}%
\begin{pgfscope}%
\pgfpathrectangle{\pgfqpoint{0.407592in}{0.222904in}}{\pgfqpoint{1.411163in}{1.411163in}}%
\pgfusepath{clip}%
\pgfsetbuttcap%
\pgfsetroundjoin%
\definecolor{currentfill}{rgb}{0.749020,0.772549,0.792157}%
\pgfsetfillcolor{currentfill}%
\pgfsetfillopacity{0.160000}%
\pgfsetlinewidth{0.000000pt}%
\definecolor{currentstroke}{rgb}{0.000000,0.000000,0.000000}%
\pgfsetstrokecolor{currentstroke}%
\pgfsetdash{}{0pt}%
\pgfpathmoveto{\pgfqpoint{1.233227in}{1.230846in}}%
\pgfpathlineto{\pgfqpoint{1.236860in}{1.246602in}}%
\pgfpathlineto{\pgfqpoint{1.240485in}{1.262329in}}%
\pgfpathlineto{\pgfqpoint{1.208000in}{1.265019in}}%
\pgfpathlineto{\pgfqpoint{1.175078in}{1.266923in}}%
\pgfpathlineto{\pgfqpoint{1.173034in}{1.251109in}}%
\pgfpathlineto{\pgfqpoint{1.170986in}{1.235265in}}%
\pgfpathlineto{\pgfqpoint{1.202316in}{1.233434in}}%
\pgfpathlineto{\pgfqpoint{1.233227in}{1.230846in}}%
\pgfpathclose%
\pgfusepath{fill}%
\end{pgfscope}%
\begin{pgfscope}%
\pgfpathrectangle{\pgfqpoint{0.407592in}{0.222904in}}{\pgfqpoint{1.411163in}{1.411163in}}%
\pgfusepath{clip}%
\pgfsetbuttcap%
\pgfsetroundjoin%
\definecolor{currentfill}{rgb}{0.749020,0.772549,0.792157}%
\pgfsetfillcolor{currentfill}%
\pgfsetfillopacity{0.160000}%
\pgfsetlinewidth{0.000000pt}%
\definecolor{currentstroke}{rgb}{0.000000,0.000000,0.000000}%
\pgfsetstrokecolor{currentstroke}%
\pgfsetdash{}{0pt}%
\pgfpathmoveto{\pgfqpoint{1.043750in}{1.234959in}}%
\pgfpathlineto{\pgfqpoint{1.042571in}{1.250796in}}%
\pgfpathlineto{\pgfqpoint{1.041394in}{1.266604in}}%
\pgfpathlineto{\pgfqpoint{1.008069in}{1.264542in}}%
\pgfpathlineto{\pgfqpoint{0.975021in}{1.261696in}}%
\pgfpathlineto{\pgfqpoint{0.977802in}{1.245980in}}%
\pgfpathlineto{\pgfqpoint{0.980587in}{1.230237in}}%
\pgfpathlineto{\pgfqpoint{1.012035in}{1.232975in}}%
\pgfpathlineto{\pgfqpoint{1.043750in}{1.234959in}}%
\pgfpathclose%
\pgfusepath{fill}%
\end{pgfscope}%
\begin{pgfscope}%
\pgfpathrectangle{\pgfqpoint{0.407592in}{0.222904in}}{\pgfqpoint{1.411163in}{1.411163in}}%
\pgfusepath{clip}%
\pgfsetbuttcap%
\pgfsetroundjoin%
\definecolor{currentfill}{rgb}{0.749020,0.772549,0.792157}%
\pgfsetfillcolor{currentfill}%
\pgfsetfillopacity{0.160000}%
\pgfsetlinewidth{0.000000pt}%
\definecolor{currentstroke}{rgb}{0.000000,0.000000,0.000000}%
\pgfsetstrokecolor{currentstroke}%
\pgfsetdash{}{0pt}%
\pgfpathmoveto{\pgfqpoint{0.778039in}{1.425154in}}%
\pgfpathlineto{\pgfqpoint{0.772266in}{1.440220in}}%
\pgfpathlineto{\pgfqpoint{0.766500in}{1.455264in}}%
\pgfpathlineto{\pgfqpoint{0.726293in}{1.447066in}}%
\pgfpathlineto{\pgfqpoint{0.687311in}{1.437983in}}%
\pgfpathlineto{\pgfqpoint{0.694531in}{1.423155in}}%
\pgfpathlineto{\pgfqpoint{0.701760in}{1.408309in}}%
\pgfpathlineto{\pgfqpoint{0.739303in}{1.417163in}}%
\pgfpathlineto{\pgfqpoint{0.778039in}{1.425154in}}%
\pgfpathclose%
\pgfusepath{fill}%
\end{pgfscope}%
\begin{pgfscope}%
\pgfpathrectangle{\pgfqpoint{0.407592in}{0.222904in}}{\pgfqpoint{1.411163in}{1.411163in}}%
\pgfusepath{clip}%
\pgfsetbuttcap%
\pgfsetroundjoin%
\definecolor{currentfill}{rgb}{0.749020,0.772549,0.792157}%
\pgfsetfillcolor{currentfill}%
\pgfsetfillopacity{0.160000}%
\pgfsetlinewidth{0.000000pt}%
\definecolor{currentstroke}{rgb}{0.000000,0.000000,0.000000}%
\pgfsetstrokecolor{currentstroke}%
\pgfsetdash{}{0pt}%
\pgfpathmoveto{\pgfqpoint{1.303676in}{1.254623in}}%
\pgfpathlineto{\pgfqpoint{1.308831in}{1.270179in}}%
\pgfpathlineto{\pgfqpoint{1.313978in}{1.285709in}}%
\pgfpathlineto{\pgfqpoint{1.281196in}{1.290104in}}%
\pgfpathlineto{\pgfqpoint{1.247717in}{1.293700in}}%
\pgfpathlineto{\pgfqpoint{1.244104in}{1.278029in}}%
\pgfpathlineto{\pgfqpoint{1.240485in}{1.262329in}}%
\pgfpathlineto{\pgfqpoint{1.272416in}{1.258861in}}%
\pgfpathlineto{\pgfqpoint{1.303676in}{1.254623in}}%
\pgfpathclose%
\pgfusepath{fill}%
\end{pgfscope}%
\begin{pgfscope}%
\pgfpathrectangle{\pgfqpoint{0.407592in}{0.222904in}}{\pgfqpoint{1.411163in}{1.411163in}}%
\pgfusepath{clip}%
\pgfsetbuttcap%
\pgfsetroundjoin%
\definecolor{currentfill}{rgb}{0.749020,0.772549,0.792157}%
\pgfsetfillcolor{currentfill}%
\pgfsetfillopacity{0.160000}%
\pgfsetlinewidth{0.000000pt}%
\definecolor{currentstroke}{rgb}{0.000000,0.000000,0.000000}%
\pgfsetstrokecolor{currentstroke}%
\pgfsetdash{}{0pt}%
\pgfpathmoveto{\pgfqpoint{1.390183in}{1.305172in}}%
\pgfpathlineto{\pgfqpoint{1.396774in}{1.320435in}}%
\pgfpathlineto{\pgfqpoint{1.403355in}{1.335676in}}%
\pgfpathlineto{\pgfqpoint{1.369419in}{1.342038in}}%
\pgfpathlineto{\pgfqpoint{1.334479in}{1.347573in}}%
\pgfpathlineto{\pgfqpoint{1.329367in}{1.332145in}}%
\pgfpathlineto{\pgfqpoint{1.324246in}{1.316692in}}%
\pgfpathlineto{\pgfqpoint{1.357700in}{1.311332in}}%
\pgfpathlineto{\pgfqpoint{1.390183in}{1.305172in}}%
\pgfpathclose%
\pgfusepath{fill}%
\end{pgfscope}%
\begin{pgfscope}%
\pgfpathrectangle{\pgfqpoint{0.407592in}{0.222904in}}{\pgfqpoint{1.411163in}{1.411163in}}%
\pgfusepath{clip}%
\pgfsetbuttcap%
\pgfsetroundjoin%
\definecolor{currentfill}{rgb}{0.749020,0.772549,0.792157}%
\pgfsetfillcolor{currentfill}%
\pgfsetfillopacity{0.160000}%
\pgfsetlinewidth{0.000000pt}%
\definecolor{currentstroke}{rgb}{0.000000,0.000000,0.000000}%
\pgfsetstrokecolor{currentstroke}%
\pgfsetdash{}{0pt}%
\pgfpathmoveto{\pgfqpoint{0.975021in}{1.261696in}}%
\pgfpathlineto{\pgfqpoint{0.972246in}{1.277383in}}%
\pgfpathlineto{\pgfqpoint{0.969475in}{1.293043in}}%
\pgfpathlineto{\pgfqpoint{0.935242in}{1.289286in}}%
\pgfpathlineto{\pgfqpoint{0.901543in}{1.284732in}}%
\pgfpathlineto{\pgfqpoint{0.905885in}{1.269220in}}%
\pgfpathlineto{\pgfqpoint{0.910233in}{1.253682in}}%
\pgfpathlineto{\pgfqpoint{0.942370in}{1.258072in}}%
\pgfpathlineto{\pgfqpoint{0.975021in}{1.261696in}}%
\pgfpathclose%
\pgfusepath{fill}%
\end{pgfscope}%
\begin{pgfscope}%
\pgfpathrectangle{\pgfqpoint{0.407592in}{0.222904in}}{\pgfqpoint{1.411163in}{1.411163in}}%
\pgfusepath{clip}%
\pgfsetbuttcap%
\pgfsetroundjoin%
\definecolor{currentfill}{rgb}{0.749020,0.772549,0.792157}%
\pgfsetfillcolor{currentfill}%
\pgfsetfillopacity{0.160000}%
\pgfsetlinewidth{0.000000pt}%
\definecolor{currentstroke}{rgb}{0.000000,0.000000,0.000000}%
\pgfsetstrokecolor{currentstroke}%
\pgfsetdash{}{0pt}%
\pgfpathmoveto{\pgfqpoint{0.892882in}{1.315681in}}%
\pgfpathlineto{\pgfqpoint{0.888561in}{1.331118in}}%
\pgfpathlineto{\pgfqpoint{0.884248in}{1.346529in}}%
\pgfpathlineto{\pgfqpoint{0.848137in}{1.340827in}}%
\pgfpathlineto{\pgfqpoint{0.812857in}{1.334301in}}%
\pgfpathlineto{\pgfqpoint{0.818690in}{1.319082in}}%
\pgfpathlineto{\pgfqpoint{0.824531in}{1.303841in}}%
\pgfpathlineto{\pgfqpoint{0.858303in}{1.310159in}}%
\pgfpathlineto{\pgfqpoint{0.892882in}{1.315681in}}%
\pgfpathclose%
\pgfusepath{fill}%
\end{pgfscope}%
\begin{pgfscope}%
\pgfpathrectangle{\pgfqpoint{0.407592in}{0.222904in}}{\pgfqpoint{1.411163in}{1.411163in}}%
\pgfusepath{clip}%
\pgfsetbuttcap%
\pgfsetroundjoin%
\definecolor{currentfill}{rgb}{0.749020,0.772549,0.792157}%
\pgfsetfillcolor{currentfill}%
\pgfsetfillopacity{0.160000}%
\pgfsetlinewidth{0.000000pt}%
\definecolor{currentstroke}{rgb}{0.000000,0.000000,0.000000}%
\pgfsetstrokecolor{currentstroke}%
\pgfsetdash{}{0pt}%
\pgfpathmoveto{\pgfqpoint{1.166878in}{1.203490in}}%
\pgfpathlineto{\pgfqpoint{1.168934in}{1.219392in}}%
\pgfpathlineto{\pgfqpoint{1.170986in}{1.235265in}}%
\pgfpathlineto{\pgfqpoint{1.139350in}{1.236336in}}%
\pgfpathlineto{\pgfqpoint{1.107522in}{1.236641in}}%
\pgfpathlineto{\pgfqpoint{1.107083in}{1.220741in}}%
\pgfpathlineto{\pgfqpoint{1.106643in}{1.204810in}}%
\pgfpathlineto{\pgfqpoint{1.136852in}{1.204516in}}%
\pgfpathlineto{\pgfqpoint{1.166878in}{1.203490in}}%
\pgfpathclose%
\pgfusepath{fill}%
\end{pgfscope}%
\begin{pgfscope}%
\pgfpathrectangle{\pgfqpoint{0.407592in}{0.222904in}}{\pgfqpoint{1.411163in}{1.411163in}}%
\pgfusepath{clip}%
\pgfsetbuttcap%
\pgfsetroundjoin%
\definecolor{currentfill}{rgb}{0.749020,0.772549,0.792157}%
\pgfsetfillcolor{currentfill}%
\pgfsetfillopacity{0.160000}%
\pgfsetlinewidth{0.000000pt}%
\definecolor{currentstroke}{rgb}{0.000000,0.000000,0.000000}%
\pgfsetstrokecolor{currentstroke}%
\pgfsetdash{}{0pt}%
\pgfpathmoveto{\pgfqpoint{1.106643in}{1.204810in}}%
\pgfpathlineto{\pgfqpoint{1.107083in}{1.220741in}}%
\pgfpathlineto{\pgfqpoint{1.107522in}{1.236641in}}%
\pgfpathlineto{\pgfqpoint{1.075617in}{1.236182in}}%
\pgfpathlineto{\pgfqpoint{1.043750in}{1.234959in}}%
\pgfpathlineto{\pgfqpoint{1.044931in}{1.219092in}}%
\pgfpathlineto{\pgfqpoint{1.046114in}{1.203196in}}%
\pgfpathlineto{\pgfqpoint{1.076360in}{1.204369in}}%
\pgfpathlineto{\pgfqpoint{1.106643in}{1.204810in}}%
\pgfpathclose%
\pgfusepath{fill}%
\end{pgfscope}%
\begin{pgfscope}%
\pgfpathrectangle{\pgfqpoint{0.407592in}{0.222904in}}{\pgfqpoint{1.411163in}{1.411163in}}%
\pgfusepath{clip}%
\pgfsetbuttcap%
\pgfsetroundjoin%
\definecolor{currentfill}{rgb}{0.749020,0.772549,0.792157}%
\pgfsetfillcolor{currentfill}%
\pgfsetfillopacity{0.160000}%
\pgfsetlinewidth{0.000000pt}%
\definecolor{currentstroke}{rgb}{0.000000,0.000000,0.000000}%
\pgfsetstrokecolor{currentstroke}%
\pgfsetdash{}{0pt}%
\pgfpathmoveto{\pgfqpoint{1.499510in}{1.380358in}}%
\pgfpathlineto{\pgfqpoint{1.507436in}{1.395265in}}%
\pgfpathlineto{\pgfqpoint{1.515351in}{1.410154in}}%
\pgfpathlineto{\pgfqpoint{1.479684in}{1.418836in}}%
\pgfpathlineto{\pgfqpoint{1.442640in}{1.426652in}}%
\pgfpathlineto{\pgfqpoint{1.436117in}{1.411545in}}%
\pgfpathlineto{\pgfqpoint{1.429584in}{1.396416in}}%
\pgfpathlineto{\pgfqpoint{1.465216in}{1.388808in}}%
\pgfpathlineto{\pgfqpoint{1.499510in}{1.380358in}}%
\pgfpathclose%
\pgfusepath{fill}%
\end{pgfscope}%
\begin{pgfscope}%
\pgfpathrectangle{\pgfqpoint{0.407592in}{0.222904in}}{\pgfqpoint{1.411163in}{1.411163in}}%
\pgfusepath{clip}%
\pgfsetbuttcap%
\pgfsetroundjoin%
\definecolor{currentfill}{rgb}{0.749020,0.772549,0.792157}%
\pgfsetfillcolor{currentfill}%
\pgfsetfillopacity{0.160000}%
\pgfsetlinewidth{0.000000pt}%
\definecolor{currentstroke}{rgb}{0.000000,0.000000,0.000000}%
\pgfsetstrokecolor{currentstroke}%
\pgfsetdash{}{0pt}%
\pgfpathmoveto{\pgfqpoint{0.789612in}{1.394958in}}%
\pgfpathlineto{\pgfqpoint{0.783821in}{1.410067in}}%
\pgfpathlineto{\pgfqpoint{0.778039in}{1.425154in}}%
\pgfpathlineto{\pgfqpoint{0.739303in}{1.417163in}}%
\pgfpathlineto{\pgfqpoint{0.701760in}{1.408309in}}%
\pgfpathlineto{\pgfqpoint{0.708998in}{1.393446in}}%
\pgfpathlineto{\pgfqpoint{0.716244in}{1.378564in}}%
\pgfpathlineto{\pgfqpoint{0.752347in}{1.387179in}}%
\pgfpathlineto{\pgfqpoint{0.789612in}{1.394958in}}%
\pgfpathclose%
\pgfusepath{fill}%
\end{pgfscope}%
\begin{pgfscope}%
\pgfpathrectangle{\pgfqpoint{0.407592in}{0.222904in}}{\pgfqpoint{1.411163in}{1.411163in}}%
\pgfusepath{clip}%
\pgfsetbuttcap%
\pgfsetroundjoin%
\definecolor{currentfill}{rgb}{0.749020,0.772549,0.792157}%
\pgfsetfillcolor{currentfill}%
\pgfsetfillopacity{0.160000}%
\pgfsetlinewidth{0.000000pt}%
\definecolor{currentstroke}{rgb}{0.000000,0.000000,0.000000}%
\pgfsetstrokecolor{currentstroke}%
\pgfsetdash{}{0pt}%
\pgfpathmoveto{\pgfqpoint{1.225944in}{1.199251in}}%
\pgfpathlineto{\pgfqpoint{1.229589in}{1.215063in}}%
\pgfpathlineto{\pgfqpoint{1.233227in}{1.230846in}}%
\pgfpathlineto{\pgfqpoint{1.202316in}{1.233434in}}%
\pgfpathlineto{\pgfqpoint{1.170986in}{1.235265in}}%
\pgfpathlineto{\pgfqpoint{1.168934in}{1.219392in}}%
\pgfpathlineto{\pgfqpoint{1.166878in}{1.203490in}}%
\pgfpathlineto{\pgfqpoint{1.196612in}{1.201733in}}%
\pgfpathlineto{\pgfqpoint{1.225944in}{1.199251in}}%
\pgfpathclose%
\pgfusepath{fill}%
\end{pgfscope}%
\begin{pgfscope}%
\pgfpathrectangle{\pgfqpoint{0.407592in}{0.222904in}}{\pgfqpoint{1.411163in}{1.411163in}}%
\pgfusepath{clip}%
\pgfsetbuttcap%
\pgfsetroundjoin%
\definecolor{currentfill}{rgb}{0.749020,0.772549,0.792157}%
\pgfsetfillcolor{currentfill}%
\pgfsetfillopacity{0.160000}%
\pgfsetlinewidth{0.000000pt}%
\definecolor{currentstroke}{rgb}{0.000000,0.000000,0.000000}%
\pgfsetstrokecolor{currentstroke}%
\pgfsetdash{}{0pt}%
\pgfpathmoveto{\pgfqpoint{1.046114in}{1.203196in}}%
\pgfpathlineto{\pgfqpoint{1.044931in}{1.219092in}}%
\pgfpathlineto{\pgfqpoint{1.043750in}{1.234959in}}%
\pgfpathlineto{\pgfqpoint{1.012035in}{1.232975in}}%
\pgfpathlineto{\pgfqpoint{0.980587in}{1.230237in}}%
\pgfpathlineto{\pgfqpoint{0.983378in}{1.214465in}}%
\pgfpathlineto{\pgfqpoint{0.986173in}{1.198666in}}%
\pgfpathlineto{\pgfqpoint{1.016015in}{1.201293in}}%
\pgfpathlineto{\pgfqpoint{1.046114in}{1.203196in}}%
\pgfpathclose%
\pgfusepath{fill}%
\end{pgfscope}%
\begin{pgfscope}%
\pgfpathrectangle{\pgfqpoint{0.407592in}{0.222904in}}{\pgfqpoint{1.411163in}{1.411163in}}%
\pgfusepath{clip}%
\pgfsetbuttcap%
\pgfsetroundjoin%
\definecolor{currentfill}{rgb}{0.749020,0.772549,0.792157}%
\pgfsetfillcolor{currentfill}%
\pgfsetfillopacity{0.160000}%
\pgfsetlinewidth{0.000000pt}%
\definecolor{currentstroke}{rgb}{0.000000,0.000000,0.000000}%
\pgfsetstrokecolor{currentstroke}%
\pgfsetdash{}{0pt}%
\pgfpathmoveto{\pgfqpoint{1.376972in}{1.274577in}}%
\pgfpathlineto{\pgfqpoint{1.383583in}{1.289886in}}%
\pgfpathlineto{\pgfqpoint{1.390183in}{1.305172in}}%
\pgfpathlineto{\pgfqpoint{1.357700in}{1.311332in}}%
\pgfpathlineto{\pgfqpoint{1.324246in}{1.316692in}}%
\pgfpathlineto{\pgfqpoint{1.319116in}{1.301213in}}%
\pgfpathlineto{\pgfqpoint{1.313978in}{1.285709in}}%
\pgfpathlineto{\pgfqpoint{1.345943in}{1.280529in}}%
\pgfpathlineto{\pgfqpoint{1.376972in}{1.274577in}}%
\pgfpathclose%
\pgfusepath{fill}%
\end{pgfscope}%
\begin{pgfscope}%
\pgfpathrectangle{\pgfqpoint{0.407592in}{0.222904in}}{\pgfqpoint{1.411163in}{1.411163in}}%
\pgfusepath{clip}%
\pgfsetbuttcap%
\pgfsetroundjoin%
\definecolor{currentfill}{rgb}{0.749020,0.772549,0.792157}%
\pgfsetfillcolor{currentfill}%
\pgfsetfillopacity{0.160000}%
\pgfsetlinewidth{0.000000pt}%
\definecolor{currentstroke}{rgb}{0.000000,0.000000,0.000000}%
\pgfsetstrokecolor{currentstroke}%
\pgfsetdash{}{0pt}%
\pgfpathmoveto{\pgfqpoint{1.293341in}{1.223435in}}%
\pgfpathlineto{\pgfqpoint{1.298513in}{1.239042in}}%
\pgfpathlineto{\pgfqpoint{1.303676in}{1.254623in}}%
\pgfpathlineto{\pgfqpoint{1.272416in}{1.258861in}}%
\pgfpathlineto{\pgfqpoint{1.240485in}{1.262329in}}%
\pgfpathlineto{\pgfqpoint{1.236860in}{1.246602in}}%
\pgfpathlineto{\pgfqpoint{1.233227in}{1.230846in}}%
\pgfpathlineto{\pgfqpoint{1.263606in}{1.227510in}}%
\pgfpathlineto{\pgfqpoint{1.293341in}{1.223435in}}%
\pgfpathclose%
\pgfusepath{fill}%
\end{pgfscope}%
\begin{pgfscope}%
\pgfpathrectangle{\pgfqpoint{0.407592in}{0.222904in}}{\pgfqpoint{1.411163in}{1.411163in}}%
\pgfusepath{clip}%
\pgfsetbuttcap%
\pgfsetroundjoin%
\definecolor{currentfill}{rgb}{0.749020,0.772549,0.792157}%
\pgfsetfillcolor{currentfill}%
\pgfsetfillopacity{0.160000}%
\pgfsetlinewidth{0.000000pt}%
\definecolor{currentstroke}{rgb}{0.000000,0.000000,0.000000}%
\pgfsetstrokecolor{currentstroke}%
\pgfsetdash{}{0pt}%
\pgfpathmoveto{\pgfqpoint{0.980587in}{1.230237in}}%
\pgfpathlineto{\pgfqpoint{0.977802in}{1.245980in}}%
\pgfpathlineto{\pgfqpoint{0.975021in}{1.261696in}}%
\pgfpathlineto{\pgfqpoint{0.942370in}{1.258072in}}%
\pgfpathlineto{\pgfqpoint{0.910233in}{1.253682in}}%
\pgfpathlineto{\pgfqpoint{0.914589in}{1.238118in}}%
\pgfpathlineto{\pgfqpoint{0.918952in}{1.222529in}}%
\pgfpathlineto{\pgfqpoint{0.949521in}{1.226751in}}%
\pgfpathlineto{\pgfqpoint{0.980587in}{1.230237in}}%
\pgfpathclose%
\pgfusepath{fill}%
\end{pgfscope}%
\begin{pgfscope}%
\pgfpathrectangle{\pgfqpoint{0.407592in}{0.222904in}}{\pgfqpoint{1.411163in}{1.411163in}}%
\pgfusepath{clip}%
\pgfsetbuttcap%
\pgfsetroundjoin%
\definecolor{currentfill}{rgb}{0.749020,0.772549,0.792157}%
\pgfsetfillcolor{currentfill}%
\pgfsetfillopacity{0.160000}%
\pgfsetlinewidth{0.000000pt}%
\definecolor{currentstroke}{rgb}{0.000000,0.000000,0.000000}%
\pgfsetstrokecolor{currentstroke}%
\pgfsetdash{}{0pt}%
\pgfpathmoveto{\pgfqpoint{1.483630in}{1.350488in}}%
\pgfpathlineto{\pgfqpoint{1.491575in}{1.365433in}}%
\pgfpathlineto{\pgfqpoint{1.499510in}{1.380358in}}%
\pgfpathlineto{\pgfqpoint{1.465216in}{1.388808in}}%
\pgfpathlineto{\pgfqpoint{1.429584in}{1.396416in}}%
\pgfpathlineto{\pgfqpoint{1.423041in}{1.381264in}}%
\pgfpathlineto{\pgfqpoint{1.416489in}{1.366090in}}%
\pgfpathlineto{\pgfqpoint{1.450708in}{1.358697in}}%
\pgfpathlineto{\pgfqpoint{1.483630in}{1.350488in}}%
\pgfpathclose%
\pgfusepath{fill}%
\end{pgfscope}%
\begin{pgfscope}%
\pgfpathrectangle{\pgfqpoint{0.407592in}{0.222904in}}{\pgfqpoint{1.411163in}{1.411163in}}%
\pgfusepath{clip}%
\pgfsetbuttcap%
\pgfsetroundjoin%
\definecolor{currentfill}{rgb}{0.749020,0.772549,0.792157}%
\pgfsetfillcolor{currentfill}%
\pgfsetfillopacity{0.160000}%
\pgfsetlinewidth{0.000000pt}%
\definecolor{currentstroke}{rgb}{0.000000,0.000000,0.000000}%
\pgfsetstrokecolor{currentstroke}%
\pgfsetdash{}{0pt}%
\pgfpathmoveto{\pgfqpoint{0.901543in}{1.284732in}}%
\pgfpathlineto{\pgfqpoint{0.897209in}{1.300220in}}%
\pgfpathlineto{\pgfqpoint{0.892882in}{1.315681in}}%
\pgfpathlineto{\pgfqpoint{0.858303in}{1.310159in}}%
\pgfpathlineto{\pgfqpoint{0.824531in}{1.303841in}}%
\pgfpathlineto{\pgfqpoint{0.830380in}{1.288577in}}%
\pgfpathlineto{\pgfqpoint{0.836238in}{1.273292in}}%
\pgfpathlineto{\pgfqpoint{0.868501in}{1.279396in}}%
\pgfpathlineto{\pgfqpoint{0.901543in}{1.284732in}}%
\pgfpathclose%
\pgfusepath{fill}%
\end{pgfscope}%
\begin{pgfscope}%
\pgfpathrectangle{\pgfqpoint{0.407592in}{0.222904in}}{\pgfqpoint{1.411163in}{1.411163in}}%
\pgfusepath{clip}%
\pgfsetbuttcap%
\pgfsetroundjoin%
\definecolor{currentfill}{rgb}{0.749020,0.772549,0.792157}%
\pgfsetfillcolor{currentfill}%
\pgfsetfillopacity{0.160000}%
\pgfsetlinewidth{0.000000pt}%
\definecolor{currentstroke}{rgb}{0.000000,0.000000,0.000000}%
\pgfsetstrokecolor{currentstroke}%
\pgfsetdash{}{0pt}%
\pgfpathmoveto{\pgfqpoint{1.162756in}{1.171595in}}%
\pgfpathlineto{\pgfqpoint{1.164819in}{1.187557in}}%
\pgfpathlineto{\pgfqpoint{1.166878in}{1.203490in}}%
\pgfpathlineto{\pgfqpoint{1.136852in}{1.204516in}}%
\pgfpathlineto{\pgfqpoint{1.106643in}{1.204810in}}%
\pgfpathlineto{\pgfqpoint{1.106202in}{1.188849in}}%
\pgfpathlineto{\pgfqpoint{1.105760in}{1.172858in}}%
\pgfpathlineto{\pgfqpoint{1.134345in}{1.172577in}}%
\pgfpathlineto{\pgfqpoint{1.162756in}{1.171595in}}%
\pgfpathclose%
\pgfusepath{fill}%
\end{pgfscope}%
\begin{pgfscope}%
\pgfpathrectangle{\pgfqpoint{0.407592in}{0.222904in}}{\pgfqpoint{1.411163in}{1.411163in}}%
\pgfusepath{clip}%
\pgfsetbuttcap%
\pgfsetroundjoin%
\definecolor{currentfill}{rgb}{0.749020,0.772549,0.792157}%
\pgfsetfillcolor{currentfill}%
\pgfsetfillopacity{0.160000}%
\pgfsetlinewidth{0.000000pt}%
\definecolor{currentstroke}{rgb}{0.000000,0.000000,0.000000}%
\pgfsetstrokecolor{currentstroke}%
\pgfsetdash{}{0pt}%
\pgfpathmoveto{\pgfqpoint{1.105760in}{1.172858in}}%
\pgfpathlineto{\pgfqpoint{1.106202in}{1.188849in}}%
\pgfpathlineto{\pgfqpoint{1.106643in}{1.204810in}}%
\pgfpathlineto{\pgfqpoint{1.076360in}{1.204369in}}%
\pgfpathlineto{\pgfqpoint{1.046114in}{1.203196in}}%
\pgfpathlineto{\pgfqpoint{1.047299in}{1.187270in}}%
\pgfpathlineto{\pgfqpoint{1.048487in}{1.171314in}}%
\pgfpathlineto{\pgfqpoint{1.077106in}{1.172436in}}%
\pgfpathlineto{\pgfqpoint{1.105760in}{1.172858in}}%
\pgfpathclose%
\pgfusepath{fill}%
\end{pgfscope}%
\begin{pgfscope}%
\pgfpathrectangle{\pgfqpoint{0.407592in}{0.222904in}}{\pgfqpoint{1.411163in}{1.411163in}}%
\pgfusepath{clip}%
\pgfsetbuttcap%
\pgfsetroundjoin%
\definecolor{currentfill}{rgb}{0.749020,0.772549,0.792157}%
\pgfsetfillcolor{currentfill}%
\pgfsetfillopacity{0.160000}%
\pgfsetlinewidth{0.000000pt}%
\definecolor{currentstroke}{rgb}{0.000000,0.000000,0.000000}%
\pgfsetstrokecolor{currentstroke}%
\pgfsetdash{}{0pt}%
\pgfpathmoveto{\pgfqpoint{0.801217in}{1.364674in}}%
\pgfpathlineto{\pgfqpoint{0.795410in}{1.379827in}}%
\pgfpathlineto{\pgfqpoint{0.789612in}{1.394958in}}%
\pgfpathlineto{\pgfqpoint{0.752347in}{1.387179in}}%
\pgfpathlineto{\pgfqpoint{0.716244in}{1.378564in}}%
\pgfpathlineto{\pgfqpoint{0.723500in}{1.363663in}}%
\pgfpathlineto{\pgfqpoint{0.730764in}{1.348745in}}%
\pgfpathlineto{\pgfqpoint{0.765427in}{1.357115in}}%
\pgfpathlineto{\pgfqpoint{0.801217in}{1.364674in}}%
\pgfpathclose%
\pgfusepath{fill}%
\end{pgfscope}%
\begin{pgfscope}%
\pgfpathrectangle{\pgfqpoint{0.407592in}{0.222904in}}{\pgfqpoint{1.411163in}{1.411163in}}%
\pgfusepath{clip}%
\pgfsetbuttcap%
\pgfsetroundjoin%
\definecolor{currentfill}{rgb}{0.749020,0.772549,0.792157}%
\pgfsetfillcolor{currentfill}%
\pgfsetfillopacity{0.160000}%
\pgfsetlinewidth{0.000000pt}%
\definecolor{currentstroke}{rgb}{0.000000,0.000000,0.000000}%
\pgfsetstrokecolor{currentstroke}%
\pgfsetdash{}{0pt}%
\pgfpathmoveto{\pgfqpoint{1.363722in}{1.243892in}}%
\pgfpathlineto{\pgfqpoint{1.370352in}{1.259246in}}%
\pgfpathlineto{\pgfqpoint{1.376972in}{1.274577in}}%
\pgfpathlineto{\pgfqpoint{1.345943in}{1.280529in}}%
\pgfpathlineto{\pgfqpoint{1.313978in}{1.285709in}}%
\pgfpathlineto{\pgfqpoint{1.308831in}{1.270179in}}%
\pgfpathlineto{\pgfqpoint{1.303676in}{1.254623in}}%
\pgfpathlineto{\pgfqpoint{1.334150in}{1.249629in}}%
\pgfpathlineto{\pgfqpoint{1.363722in}{1.243892in}}%
\pgfpathclose%
\pgfusepath{fill}%
\end{pgfscope}%
\begin{pgfscope}%
\pgfpathrectangle{\pgfqpoint{0.407592in}{0.222904in}}{\pgfqpoint{1.411163in}{1.411163in}}%
\pgfusepath{clip}%
\pgfsetbuttcap%
\pgfsetroundjoin%
\definecolor{currentfill}{rgb}{0.749020,0.772549,0.792157}%
\pgfsetfillcolor{currentfill}%
\pgfsetfillopacity{0.160000}%
\pgfsetlinewidth{0.000000pt}%
\definecolor{currentstroke}{rgb}{0.000000,0.000000,0.000000}%
\pgfsetstrokecolor{currentstroke}%
\pgfsetdash{}{0pt}%
\pgfpathmoveto{\pgfqpoint{1.218634in}{1.167541in}}%
\pgfpathlineto{\pgfqpoint{1.222292in}{1.183410in}}%
\pgfpathlineto{\pgfqpoint{1.225944in}{1.199251in}}%
\pgfpathlineto{\pgfqpoint{1.196612in}{1.201733in}}%
\pgfpathlineto{\pgfqpoint{1.166878in}{1.203490in}}%
\pgfpathlineto{\pgfqpoint{1.164819in}{1.187557in}}%
\pgfpathlineto{\pgfqpoint{1.162756in}{1.171595in}}%
\pgfpathlineto{\pgfqpoint{1.190887in}{1.169915in}}%
\pgfpathlineto{\pgfqpoint{1.218634in}{1.167541in}}%
\pgfpathclose%
\pgfusepath{fill}%
\end{pgfscope}%
\begin{pgfscope}%
\pgfpathrectangle{\pgfqpoint{0.407592in}{0.222904in}}{\pgfqpoint{1.411163in}{1.411163in}}%
\pgfusepath{clip}%
\pgfsetbuttcap%
\pgfsetroundjoin%
\definecolor{currentfill}{rgb}{0.749020,0.772549,0.792157}%
\pgfsetfillcolor{currentfill}%
\pgfsetfillopacity{0.160000}%
\pgfsetlinewidth{0.000000pt}%
\definecolor{currentstroke}{rgb}{0.000000,0.000000,0.000000}%
\pgfsetstrokecolor{currentstroke}%
\pgfsetdash{}{0pt}%
\pgfpathmoveto{\pgfqpoint{1.467710in}{1.320544in}}%
\pgfpathlineto{\pgfqpoint{1.475675in}{1.335525in}}%
\pgfpathlineto{\pgfqpoint{1.483630in}{1.350488in}}%
\pgfpathlineto{\pgfqpoint{1.450708in}{1.358697in}}%
\pgfpathlineto{\pgfqpoint{1.416489in}{1.366090in}}%
\pgfpathlineto{\pgfqpoint{1.409927in}{1.350894in}}%
\pgfpathlineto{\pgfqpoint{1.403355in}{1.335676in}}%
\pgfpathlineto{\pgfqpoint{1.436161in}{1.328505in}}%
\pgfpathlineto{\pgfqpoint{1.467710in}{1.320544in}}%
\pgfpathclose%
\pgfusepath{fill}%
\end{pgfscope}%
\begin{pgfscope}%
\pgfpathrectangle{\pgfqpoint{0.407592in}{0.222904in}}{\pgfqpoint{1.411163in}{1.411163in}}%
\pgfusepath{clip}%
\pgfsetbuttcap%
\pgfsetroundjoin%
\definecolor{currentfill}{rgb}{0.749020,0.772549,0.792157}%
\pgfsetfillcolor{currentfill}%
\pgfsetfillopacity{0.160000}%
\pgfsetlinewidth{0.000000pt}%
\definecolor{currentstroke}{rgb}{0.000000,0.000000,0.000000}%
\pgfsetstrokecolor{currentstroke}%
\pgfsetdash{}{0pt}%
\pgfpathmoveto{\pgfqpoint{1.282971in}{1.192142in}}%
\pgfpathlineto{\pgfqpoint{1.288160in}{1.207801in}}%
\pgfpathlineto{\pgfqpoint{1.293341in}{1.223435in}}%
\pgfpathlineto{\pgfqpoint{1.263606in}{1.227510in}}%
\pgfpathlineto{\pgfqpoint{1.233227in}{1.230846in}}%
\pgfpathlineto{\pgfqpoint{1.229589in}{1.215063in}}%
\pgfpathlineto{\pgfqpoint{1.225944in}{1.199251in}}%
\pgfpathlineto{\pgfqpoint{1.254766in}{1.196051in}}%
\pgfpathlineto{\pgfqpoint{1.282971in}{1.192142in}}%
\pgfpathclose%
\pgfusepath{fill}%
\end{pgfscope}%
\begin{pgfscope}%
\pgfpathrectangle{\pgfqpoint{0.407592in}{0.222904in}}{\pgfqpoint{1.411163in}{1.411163in}}%
\pgfusepath{clip}%
\pgfsetbuttcap%
\pgfsetroundjoin%
\definecolor{currentfill}{rgb}{0.749020,0.772549,0.792157}%
\pgfsetfillcolor{currentfill}%
\pgfsetfillopacity{0.160000}%
\pgfsetlinewidth{0.000000pt}%
\definecolor{currentstroke}{rgb}{0.000000,0.000000,0.000000}%
\pgfsetstrokecolor{currentstroke}%
\pgfsetdash{}{0pt}%
\pgfpathmoveto{\pgfqpoint{1.048487in}{1.171314in}}%
\pgfpathlineto{\pgfqpoint{1.047299in}{1.187270in}}%
\pgfpathlineto{\pgfqpoint{1.046114in}{1.203196in}}%
\pgfpathlineto{\pgfqpoint{1.016015in}{1.201293in}}%
\pgfpathlineto{\pgfqpoint{0.986173in}{1.198666in}}%
\pgfpathlineto{\pgfqpoint{0.988973in}{1.182838in}}%
\pgfpathlineto{\pgfqpoint{0.991779in}{1.166982in}}%
\pgfpathlineto{\pgfqpoint{1.020010in}{1.169494in}}%
\pgfpathlineto{\pgfqpoint{1.048487in}{1.171314in}}%
\pgfpathclose%
\pgfusepath{fill}%
\end{pgfscope}%
\begin{pgfscope}%
\pgfpathrectangle{\pgfqpoint{0.407592in}{0.222904in}}{\pgfqpoint{1.411163in}{1.411163in}}%
\pgfusepath{clip}%
\pgfsetbuttcap%
\pgfsetroundjoin%
\definecolor{currentfill}{rgb}{0.749020,0.772549,0.792157}%
\pgfsetfillcolor{currentfill}%
\pgfsetfillopacity{0.160000}%
\pgfsetlinewidth{0.000000pt}%
\definecolor{currentstroke}{rgb}{0.000000,0.000000,0.000000}%
\pgfsetstrokecolor{currentstroke}%
\pgfsetdash{}{0pt}%
\pgfpathmoveto{\pgfqpoint{0.910233in}{1.253682in}}%
\pgfpathlineto{\pgfqpoint{0.905885in}{1.269220in}}%
\pgfpathlineto{\pgfqpoint{0.901543in}{1.284732in}}%
\pgfpathlineto{\pgfqpoint{0.868501in}{1.279396in}}%
\pgfpathlineto{\pgfqpoint{0.836238in}{1.273292in}}%
\pgfpathlineto{\pgfqpoint{0.842105in}{1.257983in}}%
\pgfpathlineto{\pgfqpoint{0.847980in}{1.242653in}}%
\pgfpathlineto{\pgfqpoint{0.878731in}{1.248537in}}%
\pgfpathlineto{\pgfqpoint{0.910233in}{1.253682in}}%
\pgfpathclose%
\pgfusepath{fill}%
\end{pgfscope}%
\begin{pgfscope}%
\pgfpathrectangle{\pgfqpoint{0.407592in}{0.222904in}}{\pgfqpoint{1.411163in}{1.411163in}}%
\pgfusepath{clip}%
\pgfsetbuttcap%
\pgfsetroundjoin%
\definecolor{currentfill}{rgb}{0.749020,0.772549,0.792157}%
\pgfsetfillcolor{currentfill}%
\pgfsetfillopacity{0.160000}%
\pgfsetlinewidth{0.000000pt}%
\definecolor{currentstroke}{rgb}{0.000000,0.000000,0.000000}%
\pgfsetstrokecolor{currentstroke}%
\pgfsetdash{}{0pt}%
\pgfpathmoveto{\pgfqpoint{0.986173in}{1.198666in}}%
\pgfpathlineto{\pgfqpoint{0.983378in}{1.214465in}}%
\pgfpathlineto{\pgfqpoint{0.980587in}{1.230237in}}%
\pgfpathlineto{\pgfqpoint{0.949521in}{1.226751in}}%
\pgfpathlineto{\pgfqpoint{0.918952in}{1.222529in}}%
\pgfpathlineto{\pgfqpoint{0.923322in}{1.206914in}}%
\pgfpathlineto{\pgfqpoint{0.927699in}{1.191274in}}%
\pgfpathlineto{\pgfqpoint{0.956698in}{1.195323in}}%
\pgfpathlineto{\pgfqpoint{0.986173in}{1.198666in}}%
\pgfpathclose%
\pgfusepath{fill}%
\end{pgfscope}%
\begin{pgfscope}%
\pgfpathrectangle{\pgfqpoint{0.407592in}{0.222904in}}{\pgfqpoint{1.411163in}{1.411163in}}%
\pgfusepath{clip}%
\pgfsetbuttcap%
\pgfsetroundjoin%
\definecolor{currentfill}{rgb}{0.749020,0.772549,0.792157}%
\pgfsetfillcolor{currentfill}%
\pgfsetfillopacity{0.160000}%
\pgfsetlinewidth{0.000000pt}%
\definecolor{currentstroke}{rgb}{0.000000,0.000000,0.000000}%
\pgfsetstrokecolor{currentstroke}%
\pgfsetdash{}{0pt}%
\pgfpathmoveto{\pgfqpoint{0.812857in}{1.334301in}}%
\pgfpathlineto{\pgfqpoint{0.807033in}{1.349499in}}%
\pgfpathlineto{\pgfqpoint{0.801217in}{1.364674in}}%
\pgfpathlineto{\pgfqpoint{0.765427in}{1.357115in}}%
\pgfpathlineto{\pgfqpoint{0.730764in}{1.348745in}}%
\pgfpathlineto{\pgfqpoint{0.738037in}{1.333808in}}%
\pgfpathlineto{\pgfqpoint{0.745320in}{1.318853in}}%
\pgfpathlineto{\pgfqpoint{0.778541in}{1.326970in}}%
\pgfpathlineto{\pgfqpoint{0.812857in}{1.334301in}}%
\pgfpathclose%
\pgfusepath{fill}%
\end{pgfscope}%
\begin{pgfscope}%
\pgfpathrectangle{\pgfqpoint{0.407592in}{0.222904in}}{\pgfqpoint{1.411163in}{1.411163in}}%
\pgfusepath{clip}%
\pgfsetbuttcap%
\pgfsetroundjoin%
\definecolor{currentfill}{rgb}{0.749020,0.772549,0.792157}%
\pgfsetfillcolor{currentfill}%
\pgfsetfillopacity{0.160000}%
\pgfsetlinewidth{0.000000pt}%
\definecolor{currentstroke}{rgb}{0.000000,0.000000,0.000000}%
\pgfsetstrokecolor{currentstroke}%
\pgfsetdash{}{0pt}%
\pgfpathmoveto{\pgfqpoint{1.618793in}{1.448621in}}%
\pgfpathlineto{\pgfqpoint{1.627969in}{1.463170in}}%
\pgfpathlineto{\pgfqpoint{1.592260in}{1.474184in}}%
\pgfpathlineto{\pgfqpoint{1.554781in}{1.484319in}}%
\pgfpathlineto{\pgfqpoint{1.546915in}{1.469523in}}%
\pgfpathlineto{\pgfqpoint{1.583727in}{1.459505in}}%
\pgfusepath{fill}%
\end{pgfscope}%
\begin{pgfscope}%
\pgfpathrectangle{\pgfqpoint{0.407592in}{0.222904in}}{\pgfqpoint{1.411163in}{1.411163in}}%
\pgfusepath{clip}%
\pgfsetbuttcap%
\pgfsetroundjoin%
\definecolor{currentfill}{rgb}{0.749020,0.772549,0.792157}%
\pgfsetfillcolor{currentfill}%
\pgfsetfillopacity{0.160000}%
\pgfsetlinewidth{0.000000pt}%
\definecolor{currentstroke}{rgb}{0.000000,0.000000,0.000000}%
\pgfsetstrokecolor{currentstroke}%
\pgfsetdash{}{0pt}%
\pgfpathmoveto{\pgfqpoint{1.158617in}{1.139581in}}%
\pgfpathlineto{\pgfqpoint{1.160688in}{1.155603in}}%
\pgfpathlineto{\pgfqpoint{1.162756in}{1.171595in}}%
\pgfpathlineto{\pgfqpoint{1.134345in}{1.172577in}}%
\pgfpathlineto{\pgfqpoint{1.105760in}{1.172858in}}%
\pgfpathlineto{\pgfqpoint{1.105318in}{1.156836in}}%
\pgfpathlineto{\pgfqpoint{1.104874in}{1.140784in}}%
\pgfpathlineto{\pgfqpoint{1.131829in}{1.140517in}}%
\pgfpathlineto{\pgfqpoint{1.158617in}{1.139581in}}%
\pgfpathclose%
\pgfusepath{fill}%
\end{pgfscope}%
\begin{pgfscope}%
\pgfpathrectangle{\pgfqpoint{0.407592in}{0.222904in}}{\pgfqpoint{1.411163in}{1.411163in}}%
\pgfusepath{clip}%
\pgfsetbuttcap%
\pgfsetroundjoin%
\definecolor{currentfill}{rgb}{0.749020,0.772549,0.792157}%
\pgfsetfillcolor{currentfill}%
\pgfsetfillopacity{0.160000}%
\pgfsetlinewidth{0.000000pt}%
\definecolor{currentstroke}{rgb}{0.000000,0.000000,0.000000}%
\pgfsetstrokecolor{currentstroke}%
\pgfsetdash{}{0pt}%
\pgfpathmoveto{\pgfqpoint{1.104874in}{1.140784in}}%
\pgfpathlineto{\pgfqpoint{1.105318in}{1.156836in}}%
\pgfpathlineto{\pgfqpoint{1.105760in}{1.172858in}}%
\pgfpathlineto{\pgfqpoint{1.077106in}{1.172436in}}%
\pgfpathlineto{\pgfqpoint{1.048487in}{1.171314in}}%
\pgfpathlineto{\pgfqpoint{1.049677in}{1.155328in}}%
\pgfpathlineto{\pgfqpoint{1.050869in}{1.139313in}}%
\pgfpathlineto{\pgfqpoint{1.077854in}{1.140383in}}%
\pgfpathlineto{\pgfqpoint{1.104874in}{1.140784in}}%
\pgfpathclose%
\pgfusepath{fill}%
\end{pgfscope}%
\begin{pgfscope}%
\pgfpathrectangle{\pgfqpoint{0.407592in}{0.222904in}}{\pgfqpoint{1.411163in}{1.411163in}}%
\pgfusepath{clip}%
\pgfsetbuttcap%
\pgfsetroundjoin%
\definecolor{currentfill}{rgb}{0.749020,0.772549,0.792157}%
\pgfsetfillcolor{currentfill}%
\pgfsetfillopacity{0.160000}%
\pgfsetlinewidth{0.000000pt}%
\definecolor{currentstroke}{rgb}{0.000000,0.000000,0.000000}%
\pgfsetstrokecolor{currentstroke}%
\pgfsetdash{}{0pt}%
\pgfpathmoveto{\pgfqpoint{1.451750in}{1.290524in}}%
\pgfpathlineto{\pgfqpoint{1.459735in}{1.305543in}}%
\pgfpathlineto{\pgfqpoint{1.467710in}{1.320544in}}%
\pgfpathlineto{\pgfqpoint{1.436161in}{1.328505in}}%
\pgfpathlineto{\pgfqpoint{1.403355in}{1.335676in}}%
\pgfpathlineto{\pgfqpoint{1.396774in}{1.320435in}}%
\pgfpathlineto{\pgfqpoint{1.390183in}{1.305172in}}%
\pgfpathlineto{\pgfqpoint{1.421574in}{1.298229in}}%
\pgfpathlineto{\pgfqpoint{1.451750in}{1.290524in}}%
\pgfpathclose%
\pgfusepath{fill}%
\end{pgfscope}%
\begin{pgfscope}%
\pgfpathrectangle{\pgfqpoint{0.407592in}{0.222904in}}{\pgfqpoint{1.411163in}{1.411163in}}%
\pgfusepath{clip}%
\pgfsetbuttcap%
\pgfsetroundjoin%
\definecolor{currentfill}{rgb}{0.749020,0.772549,0.792157}%
\pgfsetfillcolor{currentfill}%
\pgfsetfillopacity{0.160000}%
\pgfsetlinewidth{0.000000pt}%
\definecolor{currentstroke}{rgb}{0.000000,0.000000,0.000000}%
\pgfsetstrokecolor{currentstroke}%
\pgfsetdash{}{0pt}%
\pgfpathmoveto{\pgfqpoint{0.672897in}{1.467584in}}%
\pgfpathlineto{\pgfqpoint{0.665703in}{1.482357in}}%
\pgfpathlineto{\pgfqpoint{0.625987in}{1.472043in}}%
\pgfpathlineto{\pgfqpoint{0.587852in}{1.460857in}}%
\pgfpathlineto{\pgfqpoint{0.596432in}{1.446335in}}%
\pgfpathlineto{\pgfqpoint{0.633884in}{1.457390in}}%
\pgfusepath{fill}%
\end{pgfscope}%
\begin{pgfscope}%
\pgfpathrectangle{\pgfqpoint{0.407592in}{0.222904in}}{\pgfqpoint{1.411163in}{1.411163in}}%
\pgfusepath{clip}%
\pgfsetbuttcap%
\pgfsetroundjoin%
\definecolor{currentfill}{rgb}{0.749020,0.772549,0.792157}%
\pgfsetfillcolor{currentfill}%
\pgfsetfillopacity{0.160000}%
\pgfsetlinewidth{0.000000pt}%
\definecolor{currentstroke}{rgb}{0.000000,0.000000,0.000000}%
\pgfsetstrokecolor{currentstroke}%
\pgfsetdash{}{0pt}%
\pgfpathmoveto{\pgfqpoint{1.350432in}{1.213116in}}%
\pgfpathlineto{\pgfqpoint{1.357082in}{1.228515in}}%
\pgfpathlineto{\pgfqpoint{1.363722in}{1.243892in}}%
\pgfpathlineto{\pgfqpoint{1.334150in}{1.249629in}}%
\pgfpathlineto{\pgfqpoint{1.303676in}{1.254623in}}%
\pgfpathlineto{\pgfqpoint{1.298513in}{1.239042in}}%
\pgfpathlineto{\pgfqpoint{1.293341in}{1.223435in}}%
\pgfpathlineto{\pgfqpoint{1.322319in}{1.218632in}}%
\pgfpathlineto{\pgfqpoint{1.350432in}{1.213116in}}%
\pgfpathclose%
\pgfusepath{fill}%
\end{pgfscope}%
\begin{pgfscope}%
\pgfpathrectangle{\pgfqpoint{0.407592in}{0.222904in}}{\pgfqpoint{1.411163in}{1.411163in}}%
\pgfusepath{clip}%
\pgfsetbuttcap%
\pgfsetroundjoin%
\definecolor{currentfill}{rgb}{0.749020,0.772549,0.792157}%
\pgfsetfillcolor{currentfill}%
\pgfsetfillopacity{0.160000}%
\pgfsetlinewidth{0.000000pt}%
\definecolor{currentstroke}{rgb}{0.000000,0.000000,0.000000}%
\pgfsetstrokecolor{currentstroke}%
\pgfsetdash{}{0pt}%
\pgfpathmoveto{\pgfqpoint{1.600415in}{1.419480in}}%
\pgfpathlineto{\pgfqpoint{1.609608in}{1.434057in}}%
\pgfpathlineto{\pgfqpoint{1.618793in}{1.448621in}}%
\pgfpathlineto{\pgfqpoint{1.583727in}{1.459505in}}%
\pgfpathlineto{\pgfqpoint{1.546915in}{1.469523in}}%
\pgfpathlineto{\pgfqpoint{1.539039in}{1.454708in}}%
\pgfpathlineto{\pgfqpoint{1.531153in}{1.439875in}}%
\pgfpathlineto{\pgfqpoint{1.566634in}{1.430099in}}%
\pgfpathlineto{\pgfqpoint{1.600415in}{1.419480in}}%
\pgfpathclose%
\pgfusepath{fill}%
\end{pgfscope}%
\begin{pgfscope}%
\pgfpathrectangle{\pgfqpoint{0.407592in}{0.222904in}}{\pgfqpoint{1.411163in}{1.411163in}}%
\pgfusepath{clip}%
\pgfsetbuttcap%
\pgfsetroundjoin%
\definecolor{currentfill}{rgb}{0.749020,0.772549,0.792157}%
\pgfsetfillcolor{currentfill}%
\pgfsetfillopacity{0.160000}%
\pgfsetlinewidth{0.000000pt}%
\definecolor{currentstroke}{rgb}{0.000000,0.000000,0.000000}%
\pgfsetstrokecolor{currentstroke}%
\pgfsetdash{}{0pt}%
\pgfpathmoveto{\pgfqpoint{0.918952in}{1.222529in}}%
\pgfpathlineto{\pgfqpoint{0.914589in}{1.238118in}}%
\pgfpathlineto{\pgfqpoint{0.910233in}{1.253682in}}%
\pgfpathlineto{\pgfqpoint{0.878731in}{1.248537in}}%
\pgfpathlineto{\pgfqpoint{0.847980in}{1.242653in}}%
\pgfpathlineto{\pgfqpoint{0.853864in}{1.227300in}}%
\pgfpathlineto{\pgfqpoint{0.859756in}{1.211925in}}%
\pgfpathlineto{\pgfqpoint{0.888992in}{1.217582in}}%
\pgfpathlineto{\pgfqpoint{0.918952in}{1.222529in}}%
\pgfpathclose%
\pgfusepath{fill}%
\end{pgfscope}%
\begin{pgfscope}%
\pgfpathrectangle{\pgfqpoint{0.407592in}{0.222904in}}{\pgfqpoint{1.411163in}{1.411163in}}%
\pgfusepath{clip}%
\pgfsetbuttcap%
\pgfsetroundjoin%
\definecolor{currentfill}{rgb}{0.749020,0.772549,0.792157}%
\pgfsetfillcolor{currentfill}%
\pgfsetfillopacity{0.160000}%
\pgfsetlinewidth{0.000000pt}%
\definecolor{currentstroke}{rgb}{0.000000,0.000000,0.000000}%
\pgfsetstrokecolor{currentstroke}%
\pgfsetdash{}{0pt}%
\pgfpathmoveto{\pgfqpoint{1.211298in}{1.135718in}}%
\pgfpathlineto{\pgfqpoint{1.214969in}{1.151644in}}%
\pgfpathlineto{\pgfqpoint{1.218634in}{1.167541in}}%
\pgfpathlineto{\pgfqpoint{1.190887in}{1.169915in}}%
\pgfpathlineto{\pgfqpoint{1.162756in}{1.171595in}}%
\pgfpathlineto{\pgfqpoint{1.160688in}{1.155603in}}%
\pgfpathlineto{\pgfqpoint{1.158617in}{1.139581in}}%
\pgfpathlineto{\pgfqpoint{1.185140in}{1.137980in}}%
\pgfpathlineto{\pgfqpoint{1.211298in}{1.135718in}}%
\pgfpathclose%
\pgfusepath{fill}%
\end{pgfscope}%
\begin{pgfscope}%
\pgfpathrectangle{\pgfqpoint{0.407592in}{0.222904in}}{\pgfqpoint{1.411163in}{1.411163in}}%
\pgfusepath{clip}%
\pgfsetbuttcap%
\pgfsetroundjoin%
\definecolor{currentfill}{rgb}{0.749020,0.772549,0.792157}%
\pgfsetfillcolor{currentfill}%
\pgfsetfillopacity{0.160000}%
\pgfsetlinewidth{0.000000pt}%
\definecolor{currentstroke}{rgb}{0.000000,0.000000,0.000000}%
\pgfsetstrokecolor{currentstroke}%
\pgfsetdash{}{0pt}%
\pgfpathmoveto{\pgfqpoint{1.272566in}{1.160745in}}%
\pgfpathlineto{\pgfqpoint{1.277773in}{1.176457in}}%
\pgfpathlineto{\pgfqpoint{1.282971in}{1.192142in}}%
\pgfpathlineto{\pgfqpoint{1.254766in}{1.196051in}}%
\pgfpathlineto{\pgfqpoint{1.225944in}{1.199251in}}%
\pgfpathlineto{\pgfqpoint{1.222292in}{1.183410in}}%
\pgfpathlineto{\pgfqpoint{1.218634in}{1.167541in}}%
\pgfpathlineto{\pgfqpoint{1.245895in}{1.164482in}}%
\pgfpathlineto{\pgfqpoint{1.272566in}{1.160745in}}%
\pgfpathclose%
\pgfusepath{fill}%
\end{pgfscope}%
\begin{pgfscope}%
\pgfpathrectangle{\pgfqpoint{0.407592in}{0.222904in}}{\pgfqpoint{1.411163in}{1.411163in}}%
\pgfusepath{clip}%
\pgfsetbuttcap%
\pgfsetroundjoin%
\definecolor{currentfill}{rgb}{0.749020,0.772549,0.792157}%
\pgfsetfillcolor{currentfill}%
\pgfsetfillopacity{0.160000}%
\pgfsetlinewidth{0.000000pt}%
\definecolor{currentstroke}{rgb}{0.000000,0.000000,0.000000}%
\pgfsetstrokecolor{currentstroke}%
\pgfsetdash{}{0pt}%
\pgfpathmoveto{\pgfqpoint{0.824531in}{1.303841in}}%
\pgfpathlineto{\pgfqpoint{0.818690in}{1.319082in}}%
\pgfpathlineto{\pgfqpoint{0.812857in}{1.334301in}}%
\pgfpathlineto{\pgfqpoint{0.778541in}{1.326970in}}%
\pgfpathlineto{\pgfqpoint{0.745320in}{1.318853in}}%
\pgfpathlineto{\pgfqpoint{0.752611in}{1.303880in}}%
\pgfpathlineto{\pgfqpoint{0.759911in}{1.288888in}}%
\pgfpathlineto{\pgfqpoint{0.791691in}{1.296744in}}%
\pgfpathlineto{\pgfqpoint{0.824531in}{1.303841in}}%
\pgfpathclose%
\pgfusepath{fill}%
\end{pgfscope}%
\begin{pgfscope}%
\pgfpathrectangle{\pgfqpoint{0.407592in}{0.222904in}}{\pgfqpoint{1.411163in}{1.411163in}}%
\pgfusepath{clip}%
\pgfsetbuttcap%
\pgfsetroundjoin%
\definecolor{currentfill}{rgb}{0.749020,0.772549,0.792157}%
\pgfsetfillcolor{currentfill}%
\pgfsetfillopacity{0.160000}%
\pgfsetlinewidth{0.000000pt}%
\definecolor{currentstroke}{rgb}{0.000000,0.000000,0.000000}%
\pgfsetstrokecolor{currentstroke}%
\pgfsetdash{}{0pt}%
\pgfpathmoveto{\pgfqpoint{1.050869in}{1.139313in}}%
\pgfpathlineto{\pgfqpoint{1.049677in}{1.155328in}}%
\pgfpathlineto{\pgfqpoint{1.048487in}{1.171314in}}%
\pgfpathlineto{\pgfqpoint{1.020010in}{1.169494in}}%
\pgfpathlineto{\pgfqpoint{0.991779in}{1.166982in}}%
\pgfpathlineto{\pgfqpoint{0.994589in}{1.151098in}}%
\pgfpathlineto{\pgfqpoint{0.997404in}{1.135186in}}%
\pgfpathlineto{\pgfqpoint{1.024019in}{1.137579in}}%
\pgfpathlineto{\pgfqpoint{1.050869in}{1.139313in}}%
\pgfpathclose%
\pgfusepath{fill}%
\end{pgfscope}%
\begin{pgfscope}%
\pgfpathrectangle{\pgfqpoint{0.407592in}{0.222904in}}{\pgfqpoint{1.411163in}{1.411163in}}%
\pgfusepath{clip}%
\pgfsetbuttcap%
\pgfsetroundjoin%
\definecolor{currentfill}{rgb}{0.749020,0.772549,0.792157}%
\pgfsetfillcolor{currentfill}%
\pgfsetfillopacity{0.160000}%
\pgfsetlinewidth{0.000000pt}%
\definecolor{currentstroke}{rgb}{0.000000,0.000000,0.000000}%
\pgfsetstrokecolor{currentstroke}%
\pgfsetdash{}{0pt}%
\pgfpathmoveto{\pgfqpoint{0.991779in}{1.166982in}}%
\pgfpathlineto{\pgfqpoint{0.988973in}{1.182838in}}%
\pgfpathlineto{\pgfqpoint{0.986173in}{1.198666in}}%
\pgfpathlineto{\pgfqpoint{0.956698in}{1.195323in}}%
\pgfpathlineto{\pgfqpoint{0.927699in}{1.191274in}}%
\pgfpathlineto{\pgfqpoint{0.932084in}{1.175607in}}%
\pgfpathlineto{\pgfqpoint{0.936476in}{1.159915in}}%
\pgfpathlineto{\pgfqpoint{0.963899in}{1.163786in}}%
\pgfpathlineto{\pgfqpoint{0.991779in}{1.166982in}}%
\pgfpathclose%
\pgfusepath{fill}%
\end{pgfscope}%
\begin{pgfscope}%
\pgfpathrectangle{\pgfqpoint{0.407592in}{0.222904in}}{\pgfqpoint{1.411163in}{1.411163in}}%
\pgfusepath{clip}%
\pgfsetbuttcap%
\pgfsetroundjoin%
\definecolor{currentfill}{rgb}{0.749020,0.772549,0.792157}%
\pgfsetfillcolor{currentfill}%
\pgfsetfillopacity{0.160000}%
\pgfsetlinewidth{0.000000pt}%
\definecolor{currentstroke}{rgb}{0.000000,0.000000,0.000000}%
\pgfsetstrokecolor{currentstroke}%
\pgfsetdash{}{0pt}%
\pgfpathmoveto{\pgfqpoint{0.687311in}{1.437983in}}%
\pgfpathlineto{\pgfqpoint{0.680099in}{1.452792in}}%
\pgfpathlineto{\pgfqpoint{0.672897in}{1.467584in}}%
\pgfpathlineto{\pgfqpoint{0.633884in}{1.457390in}}%
\pgfpathlineto{\pgfqpoint{0.596432in}{1.446335in}}%
\pgfpathlineto{\pgfqpoint{0.605020in}{1.431799in}}%
\pgfpathlineto{\pgfqpoint{0.613616in}{1.417250in}}%
\pgfpathlineto{\pgfqpoint{0.649702in}{1.428035in}}%
\pgfpathlineto{\pgfqpoint{0.687311in}{1.437983in}}%
\pgfpathclose%
\pgfusepath{fill}%
\end{pgfscope}%
\begin{pgfscope}%
\pgfpathrectangle{\pgfqpoint{0.407592in}{0.222904in}}{\pgfqpoint{1.411163in}{1.411163in}}%
\pgfusepath{clip}%
\pgfsetbuttcap%
\pgfsetroundjoin%
\definecolor{currentfill}{rgb}{0.749020,0.772549,0.792157}%
\pgfsetfillcolor{currentfill}%
\pgfsetfillopacity{0.160000}%
\pgfsetlinewidth{0.000000pt}%
\definecolor{currentstroke}{rgb}{0.000000,0.000000,0.000000}%
\pgfsetstrokecolor{currentstroke}%
\pgfsetdash{}{0pt}%
\pgfpathmoveto{\pgfqpoint{1.582000in}{1.390282in}}%
\pgfpathlineto{\pgfqpoint{1.591212in}{1.404888in}}%
\pgfpathlineto{\pgfqpoint{1.600415in}{1.419480in}}%
\pgfpathlineto{\pgfqpoint{1.566634in}{1.430099in}}%
\pgfpathlineto{\pgfqpoint{1.531153in}{1.439875in}}%
\pgfpathlineto{\pgfqpoint{1.523257in}{1.425024in}}%
\pgfpathlineto{\pgfqpoint{1.515351in}{1.410154in}}%
\pgfpathlineto{\pgfqpoint{1.549502in}{1.400627in}}%
\pgfpathlineto{\pgfqpoint{1.582000in}{1.390282in}}%
\pgfpathclose%
\pgfusepath{fill}%
\end{pgfscope}%
\begin{pgfscope}%
\pgfpathrectangle{\pgfqpoint{0.407592in}{0.222904in}}{\pgfqpoint{1.411163in}{1.411163in}}%
\pgfusepath{clip}%
\pgfsetbuttcap%
\pgfsetroundjoin%
\definecolor{currentfill}{rgb}{0.749020,0.772549,0.792157}%
\pgfsetfillcolor{currentfill}%
\pgfsetfillopacity{0.160000}%
\pgfsetlinewidth{0.000000pt}%
\definecolor{currentstroke}{rgb}{0.000000,0.000000,0.000000}%
\pgfsetstrokecolor{currentstroke}%
\pgfsetdash{}{0pt}%
\pgfpathmoveto{\pgfqpoint{1.435749in}{1.260429in}}%
\pgfpathlineto{\pgfqpoint{1.443755in}{1.275486in}}%
\pgfpathlineto{\pgfqpoint{1.451750in}{1.290524in}}%
\pgfpathlineto{\pgfqpoint{1.421574in}{1.298229in}}%
\pgfpathlineto{\pgfqpoint{1.390183in}{1.305172in}}%
\pgfpathlineto{\pgfqpoint{1.383583in}{1.289886in}}%
\pgfpathlineto{\pgfqpoint{1.376972in}{1.274577in}}%
\pgfpathlineto{\pgfqpoint{1.406946in}{1.267871in}}%
\pgfpathlineto{\pgfqpoint{1.435749in}{1.260429in}}%
\pgfpathclose%
\pgfusepath{fill}%
\end{pgfscope}%
\begin{pgfscope}%
\pgfpathrectangle{\pgfqpoint{0.407592in}{0.222904in}}{\pgfqpoint{1.411163in}{1.411163in}}%
\pgfusepath{clip}%
\pgfsetbuttcap%
\pgfsetroundjoin%
\definecolor{currentfill}{rgb}{0.749020,0.772549,0.792157}%
\pgfsetfillcolor{currentfill}%
\pgfsetfillopacity{0.160000}%
\pgfsetlinewidth{0.000000pt}%
\definecolor{currentstroke}{rgb}{0.000000,0.000000,0.000000}%
\pgfsetstrokecolor{currentstroke}%
\pgfsetdash{}{0pt}%
\pgfpathmoveto{\pgfqpoint{1.154463in}{1.107447in}}%
\pgfpathlineto{\pgfqpoint{1.156542in}{1.123529in}}%
\pgfpathlineto{\pgfqpoint{1.158617in}{1.139581in}}%
\pgfpathlineto{\pgfqpoint{1.131829in}{1.140517in}}%
\pgfpathlineto{\pgfqpoint{1.104874in}{1.140784in}}%
\pgfpathlineto{\pgfqpoint{1.104430in}{1.124702in}}%
\pgfpathlineto{\pgfqpoint{1.103985in}{1.108589in}}%
\pgfpathlineto{\pgfqpoint{1.129303in}{1.108335in}}%
\pgfpathlineto{\pgfqpoint{1.154463in}{1.107447in}}%
\pgfpathclose%
\pgfusepath{fill}%
\end{pgfscope}%
\begin{pgfscope}%
\pgfpathrectangle{\pgfqpoint{0.407592in}{0.222904in}}{\pgfqpoint{1.411163in}{1.411163in}}%
\pgfusepath{clip}%
\pgfsetbuttcap%
\pgfsetroundjoin%
\definecolor{currentfill}{rgb}{0.749020,0.772549,0.792157}%
\pgfsetfillcolor{currentfill}%
\pgfsetfillopacity{0.160000}%
\pgfsetlinewidth{0.000000pt}%
\definecolor{currentstroke}{rgb}{0.000000,0.000000,0.000000}%
\pgfsetstrokecolor{currentstroke}%
\pgfsetdash{}{0pt}%
\pgfpathmoveto{\pgfqpoint{1.103985in}{1.108589in}}%
\pgfpathlineto{\pgfqpoint{1.104430in}{1.124702in}}%
\pgfpathlineto{\pgfqpoint{1.104874in}{1.140784in}}%
\pgfpathlineto{\pgfqpoint{1.077854in}{1.140383in}}%
\pgfpathlineto{\pgfqpoint{1.050869in}{1.139313in}}%
\pgfpathlineto{\pgfqpoint{1.052063in}{1.123268in}}%
\pgfpathlineto{\pgfqpoint{1.053260in}{1.107192in}}%
\pgfpathlineto{\pgfqpoint{1.078606in}{1.108208in}}%
\pgfpathlineto{\pgfqpoint{1.103985in}{1.108589in}}%
\pgfpathclose%
\pgfusepath{fill}%
\end{pgfscope}%
\begin{pgfscope}%
\pgfpathrectangle{\pgfqpoint{0.407592in}{0.222904in}}{\pgfqpoint{1.411163in}{1.411163in}}%
\pgfusepath{clip}%
\pgfsetbuttcap%
\pgfsetroundjoin%
\definecolor{currentfill}{rgb}{0.749020,0.772549,0.792157}%
\pgfsetfillcolor{currentfill}%
\pgfsetfillopacity{0.160000}%
\pgfsetlinewidth{0.000000pt}%
\definecolor{currentstroke}{rgb}{0.000000,0.000000,0.000000}%
\pgfsetstrokecolor{currentstroke}%
\pgfsetdash{}{0pt}%
\pgfpathmoveto{\pgfqpoint{1.337103in}{1.182248in}}%
\pgfpathlineto{\pgfqpoint{1.343773in}{1.197694in}}%
\pgfpathlineto{\pgfqpoint{1.350432in}{1.213116in}}%
\pgfpathlineto{\pgfqpoint{1.322319in}{1.218632in}}%
\pgfpathlineto{\pgfqpoint{1.293341in}{1.223435in}}%
\pgfpathlineto{\pgfqpoint{1.288160in}{1.207801in}}%
\pgfpathlineto{\pgfqpoint{1.282971in}{1.192142in}}%
\pgfpathlineto{\pgfqpoint{1.310451in}{1.187537in}}%
\pgfpathlineto{\pgfqpoint{1.337103in}{1.182248in}}%
\pgfpathclose%
\pgfusepath{fill}%
\end{pgfscope}%
\begin{pgfscope}%
\pgfpathrectangle{\pgfqpoint{0.407592in}{0.222904in}}{\pgfqpoint{1.411163in}{1.411163in}}%
\pgfusepath{clip}%
\pgfsetbuttcap%
\pgfsetroundjoin%
\definecolor{currentfill}{rgb}{0.749020,0.772549,0.792157}%
\pgfsetfillcolor{currentfill}%
\pgfsetfillopacity{0.160000}%
\pgfsetlinewidth{0.000000pt}%
\definecolor{currentstroke}{rgb}{0.000000,0.000000,0.000000}%
\pgfsetstrokecolor{currentstroke}%
\pgfsetdash{}{0pt}%
\pgfpathmoveto{\pgfqpoint{0.836238in}{1.273292in}}%
\pgfpathlineto{\pgfqpoint{0.830380in}{1.288577in}}%
\pgfpathlineto{\pgfqpoint{0.824531in}{1.303841in}}%
\pgfpathlineto{\pgfqpoint{0.791691in}{1.296744in}}%
\pgfpathlineto{\pgfqpoint{0.759911in}{1.288888in}}%
\pgfpathlineto{\pgfqpoint{0.767220in}{1.273878in}}%
\pgfpathlineto{\pgfqpoint{0.774538in}{1.258849in}}%
\pgfpathlineto{\pgfqpoint{0.804877in}{1.266436in}}%
\pgfpathlineto{\pgfqpoint{0.836238in}{1.273292in}}%
\pgfpathclose%
\pgfusepath{fill}%
\end{pgfscope}%
\begin{pgfscope}%
\pgfpathrectangle{\pgfqpoint{0.407592in}{0.222904in}}{\pgfqpoint{1.411163in}{1.411163in}}%
\pgfusepath{clip}%
\pgfsetbuttcap%
\pgfsetroundjoin%
\definecolor{currentfill}{rgb}{0.749020,0.772549,0.792157}%
\pgfsetfillcolor{currentfill}%
\pgfsetfillopacity{0.160000}%
\pgfsetlinewidth{0.000000pt}%
\definecolor{currentstroke}{rgb}{0.000000,0.000000,0.000000}%
\pgfsetstrokecolor{currentstroke}%
\pgfsetdash{}{0pt}%
\pgfpathmoveto{\pgfqpoint{0.701760in}{1.408309in}}%
\pgfpathlineto{\pgfqpoint{0.694531in}{1.423155in}}%
\pgfpathlineto{\pgfqpoint{0.687311in}{1.437983in}}%
\pgfpathlineto{\pgfqpoint{0.649702in}{1.428035in}}%
\pgfpathlineto{\pgfqpoint{0.613616in}{1.417250in}}%
\pgfpathlineto{\pgfqpoint{0.622220in}{1.402687in}}%
\pgfpathlineto{\pgfqpoint{0.630832in}{1.388110in}}%
\pgfpathlineto{\pgfqpoint{0.665555in}{1.398616in}}%
\pgfpathlineto{\pgfqpoint{0.701760in}{1.408309in}}%
\pgfpathclose%
\pgfusepath{fill}%
\end{pgfscope}%
\begin{pgfscope}%
\pgfpathrectangle{\pgfqpoint{0.407592in}{0.222904in}}{\pgfqpoint{1.411163in}{1.411163in}}%
\pgfusepath{clip}%
\pgfsetbuttcap%
\pgfsetroundjoin%
\definecolor{currentfill}{rgb}{0.749020,0.772549,0.792157}%
\pgfsetfillcolor{currentfill}%
\pgfsetfillopacity{0.160000}%
\pgfsetlinewidth{0.000000pt}%
\definecolor{currentstroke}{rgb}{0.000000,0.000000,0.000000}%
\pgfsetstrokecolor{currentstroke}%
\pgfsetdash{}{0pt}%
\pgfpathmoveto{\pgfqpoint{0.927699in}{1.191274in}}%
\pgfpathlineto{\pgfqpoint{0.923322in}{1.206914in}}%
\pgfpathlineto{\pgfqpoint{0.918952in}{1.222529in}}%
\pgfpathlineto{\pgfqpoint{0.888992in}{1.217582in}}%
\pgfpathlineto{\pgfqpoint{0.859756in}{1.211925in}}%
\pgfpathlineto{\pgfqpoint{0.865657in}{1.196527in}}%
\pgfpathlineto{\pgfqpoint{0.871567in}{1.181107in}}%
\pgfpathlineto{\pgfqpoint{0.899286in}{1.186530in}}%
\pgfpathlineto{\pgfqpoint{0.927699in}{1.191274in}}%
\pgfpathclose%
\pgfusepath{fill}%
\end{pgfscope}%
\begin{pgfscope}%
\pgfpathrectangle{\pgfqpoint{0.407592in}{0.222904in}}{\pgfqpoint{1.411163in}{1.411163in}}%
\pgfusepath{clip}%
\pgfsetbuttcap%
\pgfsetroundjoin%
\definecolor{currentfill}{rgb}{0.749020,0.772549,0.792157}%
\pgfsetfillcolor{currentfill}%
\pgfsetfillopacity{0.160000}%
\pgfsetlinewidth{0.000000pt}%
\definecolor{currentstroke}{rgb}{0.000000,0.000000,0.000000}%
\pgfsetstrokecolor{currentstroke}%
\pgfsetdash{}{0pt}%
\pgfpathmoveto{\pgfqpoint{1.262126in}{1.129244in}}%
\pgfpathlineto{\pgfqpoint{1.267351in}{1.145008in}}%
\pgfpathlineto{\pgfqpoint{1.272566in}{1.160745in}}%
\pgfpathlineto{\pgfqpoint{1.245895in}{1.164482in}}%
\pgfpathlineto{\pgfqpoint{1.218634in}{1.167541in}}%
\pgfpathlineto{\pgfqpoint{1.214969in}{1.151644in}}%
\pgfpathlineto{\pgfqpoint{1.211298in}{1.135718in}}%
\pgfpathlineto{\pgfqpoint{1.236993in}{1.132803in}}%
\pgfpathlineto{\pgfqpoint{1.262126in}{1.129244in}}%
\pgfpathclose%
\pgfusepath{fill}%
\end{pgfscope}%
\begin{pgfscope}%
\pgfpathrectangle{\pgfqpoint{0.407592in}{0.222904in}}{\pgfqpoint{1.411163in}{1.411163in}}%
\pgfusepath{clip}%
\pgfsetbuttcap%
\pgfsetroundjoin%
\definecolor{currentfill}{rgb}{0.749020,0.772549,0.792157}%
\pgfsetfillcolor{currentfill}%
\pgfsetfillopacity{0.160000}%
\pgfsetlinewidth{0.000000pt}%
\definecolor{currentstroke}{rgb}{0.000000,0.000000,0.000000}%
\pgfsetstrokecolor{currentstroke}%
\pgfsetdash{}{0pt}%
\pgfpathmoveto{\pgfqpoint{1.203936in}{1.103781in}}%
\pgfpathlineto{\pgfqpoint{1.207620in}{1.119764in}}%
\pgfpathlineto{\pgfqpoint{1.211298in}{1.135718in}}%
\pgfpathlineto{\pgfqpoint{1.185140in}{1.137980in}}%
\pgfpathlineto{\pgfqpoint{1.158617in}{1.139581in}}%
\pgfpathlineto{\pgfqpoint{1.156542in}{1.123529in}}%
\pgfpathlineto{\pgfqpoint{1.154463in}{1.107447in}}%
\pgfpathlineto{\pgfqpoint{1.179372in}{1.105927in}}%
\pgfpathlineto{\pgfqpoint{1.203936in}{1.103781in}}%
\pgfpathclose%
\pgfusepath{fill}%
\end{pgfscope}%
\begin{pgfscope}%
\pgfpathrectangle{\pgfqpoint{0.407592in}{0.222904in}}{\pgfqpoint{1.411163in}{1.411163in}}%
\pgfusepath{clip}%
\pgfsetbuttcap%
\pgfsetroundjoin%
\definecolor{currentfill}{rgb}{0.749020,0.772549,0.792157}%
\pgfsetfillcolor{currentfill}%
\pgfsetfillopacity{0.160000}%
\pgfsetlinewidth{0.000000pt}%
\definecolor{currentstroke}{rgb}{0.000000,0.000000,0.000000}%
\pgfsetstrokecolor{currentstroke}%
\pgfsetdash{}{0pt}%
\pgfpathmoveto{\pgfqpoint{1.563549in}{1.361026in}}%
\pgfpathlineto{\pgfqpoint{1.572779in}{1.375661in}}%
\pgfpathlineto{\pgfqpoint{1.582000in}{1.390282in}}%
\pgfpathlineto{\pgfqpoint{1.549502in}{1.400627in}}%
\pgfpathlineto{\pgfqpoint{1.515351in}{1.410154in}}%
\pgfpathlineto{\pgfqpoint{1.507436in}{1.395265in}}%
\pgfpathlineto{\pgfqpoint{1.499510in}{1.380358in}}%
\pgfpathlineto{\pgfqpoint{1.532332in}{1.371089in}}%
\pgfpathlineto{\pgfqpoint{1.563549in}{1.361026in}}%
\pgfpathclose%
\pgfusepath{fill}%
\end{pgfscope}%
\begin{pgfscope}%
\pgfpathrectangle{\pgfqpoint{0.407592in}{0.222904in}}{\pgfqpoint{1.411163in}{1.411163in}}%
\pgfusepath{clip}%
\pgfsetbuttcap%
\pgfsetroundjoin%
\definecolor{currentfill}{rgb}{0.749020,0.772549,0.792157}%
\pgfsetfillcolor{currentfill}%
\pgfsetfillopacity{0.160000}%
\pgfsetlinewidth{0.000000pt}%
\definecolor{currentstroke}{rgb}{0.000000,0.000000,0.000000}%
\pgfsetstrokecolor{currentstroke}%
\pgfsetdash{}{0pt}%
\pgfpathmoveto{\pgfqpoint{1.053260in}{1.107192in}}%
\pgfpathlineto{\pgfqpoint{1.052063in}{1.123268in}}%
\pgfpathlineto{\pgfqpoint{1.050869in}{1.139313in}}%
\pgfpathlineto{\pgfqpoint{1.024019in}{1.137579in}}%
\pgfpathlineto{\pgfqpoint{0.997404in}{1.135186in}}%
\pgfpathlineto{\pgfqpoint{1.000225in}{1.119245in}}%
\pgfpathlineto{\pgfqpoint{1.003050in}{1.103275in}}%
\pgfpathlineto{\pgfqpoint{1.028043in}{1.105547in}}%
\pgfpathlineto{\pgfqpoint{1.053260in}{1.107192in}}%
\pgfpathclose%
\pgfusepath{fill}%
\end{pgfscope}%
\begin{pgfscope}%
\pgfpathrectangle{\pgfqpoint{0.407592in}{0.222904in}}{\pgfqpoint{1.411163in}{1.411163in}}%
\pgfusepath{clip}%
\pgfsetbuttcap%
\pgfsetroundjoin%
\definecolor{currentfill}{rgb}{0.749020,0.772549,0.792157}%
\pgfsetfillcolor{currentfill}%
\pgfsetfillopacity{0.160000}%
\pgfsetlinewidth{0.000000pt}%
\definecolor{currentstroke}{rgb}{0.000000,0.000000,0.000000}%
\pgfsetstrokecolor{currentstroke}%
\pgfsetdash{}{0pt}%
\pgfpathmoveto{\pgfqpoint{0.997404in}{1.135186in}}%
\pgfpathlineto{\pgfqpoint{0.994589in}{1.151098in}}%
\pgfpathlineto{\pgfqpoint{0.991779in}{1.166982in}}%
\pgfpathlineto{\pgfqpoint{0.963899in}{1.163786in}}%
\pgfpathlineto{\pgfqpoint{0.936476in}{1.159915in}}%
\pgfpathlineto{\pgfqpoint{0.940875in}{1.144197in}}%
\pgfpathlineto{\pgfqpoint{0.945281in}{1.128453in}}%
\pgfpathlineto{\pgfqpoint{0.971125in}{1.132140in}}%
\pgfpathlineto{\pgfqpoint{0.997404in}{1.135186in}}%
\pgfpathclose%
\pgfusepath{fill}%
\end{pgfscope}%
\begin{pgfscope}%
\pgfpathrectangle{\pgfqpoint{0.407592in}{0.222904in}}{\pgfqpoint{1.411163in}{1.411163in}}%
\pgfusepath{clip}%
\pgfsetbuttcap%
\pgfsetroundjoin%
\definecolor{currentfill}{rgb}{0.749020,0.772549,0.792157}%
\pgfsetfillcolor{currentfill}%
\pgfsetfillopacity{0.160000}%
\pgfsetlinewidth{0.000000pt}%
\definecolor{currentstroke}{rgb}{0.000000,0.000000,0.000000}%
\pgfsetstrokecolor{currentstroke}%
\pgfsetdash{}{0pt}%
\pgfpathmoveto{\pgfqpoint{1.419709in}{1.230258in}}%
\pgfpathlineto{\pgfqpoint{1.427734in}{1.245353in}}%
\pgfpathlineto{\pgfqpoint{1.435749in}{1.260429in}}%
\pgfpathlineto{\pgfqpoint{1.406946in}{1.267871in}}%
\pgfpathlineto{\pgfqpoint{1.376972in}{1.274577in}}%
\pgfpathlineto{\pgfqpoint{1.370352in}{1.259246in}}%
\pgfpathlineto{\pgfqpoint{1.363722in}{1.243892in}}%
\pgfpathlineto{\pgfqpoint{1.392279in}{1.237429in}}%
\pgfpathlineto{\pgfqpoint{1.419709in}{1.230258in}}%
\pgfpathclose%
\pgfusepath{fill}%
\end{pgfscope}%
\begin{pgfscope}%
\pgfpathrectangle{\pgfqpoint{0.407592in}{0.222904in}}{\pgfqpoint{1.411163in}{1.411163in}}%
\pgfusepath{clip}%
\pgfsetbuttcap%
\pgfsetroundjoin%
\definecolor{currentfill}{rgb}{0.749020,0.772549,0.792157}%
\pgfsetfillcolor{currentfill}%
\pgfsetfillopacity{0.160000}%
\pgfsetlinewidth{0.000000pt}%
\definecolor{currentstroke}{rgb}{0.000000,0.000000,0.000000}%
\pgfsetstrokecolor{currentstroke}%
\pgfsetdash{}{0pt}%
\pgfpathmoveto{\pgfqpoint{0.716244in}{1.378564in}}%
\pgfpathlineto{\pgfqpoint{0.708998in}{1.393446in}}%
\pgfpathlineto{\pgfqpoint{0.701760in}{1.408309in}}%
\pgfpathlineto{\pgfqpoint{0.665555in}{1.398616in}}%
\pgfpathlineto{\pgfqpoint{0.630832in}{1.388110in}}%
\pgfpathlineto{\pgfqpoint{0.639452in}{1.373519in}}%
\pgfpathlineto{\pgfqpoint{0.648081in}{1.358914in}}%
\pgfpathlineto{\pgfqpoint{0.681443in}{1.369133in}}%
\pgfpathlineto{\pgfqpoint{0.716244in}{1.378564in}}%
\pgfpathclose%
\pgfusepath{fill}%
\end{pgfscope}%
\begin{pgfscope}%
\pgfpathrectangle{\pgfqpoint{0.407592in}{0.222904in}}{\pgfqpoint{1.411163in}{1.411163in}}%
\pgfusepath{clip}%
\pgfsetbuttcap%
\pgfsetroundjoin%
\definecolor{currentfill}{rgb}{0.749020,0.772549,0.792157}%
\pgfsetfillcolor{currentfill}%
\pgfsetfillopacity{0.160000}%
\pgfsetlinewidth{0.000000pt}%
\definecolor{currentstroke}{rgb}{0.000000,0.000000,0.000000}%
\pgfsetstrokecolor{currentstroke}%
\pgfsetdash{}{0pt}%
\pgfpathmoveto{\pgfqpoint{1.323735in}{1.151289in}}%
\pgfpathlineto{\pgfqpoint{1.330424in}{1.166780in}}%
\pgfpathlineto{\pgfqpoint{1.337103in}{1.182248in}}%
\pgfpathlineto{\pgfqpoint{1.310451in}{1.187537in}}%
\pgfpathlineto{\pgfqpoint{1.282971in}{1.192142in}}%
\pgfpathlineto{\pgfqpoint{1.277773in}{1.176457in}}%
\pgfpathlineto{\pgfqpoint{1.272566in}{1.160745in}}%
\pgfpathlineto{\pgfqpoint{1.298546in}{1.156343in}}%
\pgfpathlineto{\pgfqpoint{1.323735in}{1.151289in}}%
\pgfpathclose%
\pgfusepath{fill}%
\end{pgfscope}%
\begin{pgfscope}%
\pgfpathrectangle{\pgfqpoint{0.407592in}{0.222904in}}{\pgfqpoint{1.411163in}{1.411163in}}%
\pgfusepath{clip}%
\pgfsetbuttcap%
\pgfsetroundjoin%
\definecolor{currentfill}{rgb}{0.749020,0.772549,0.792157}%
\pgfsetfillcolor{currentfill}%
\pgfsetfillopacity{0.160000}%
\pgfsetlinewidth{0.000000pt}%
\definecolor{currentstroke}{rgb}{0.000000,0.000000,0.000000}%
\pgfsetstrokecolor{currentstroke}%
\pgfsetdash{}{0pt}%
\pgfpathmoveto{\pgfqpoint{0.847980in}{1.242653in}}%
\pgfpathlineto{\pgfqpoint{0.842105in}{1.257983in}}%
\pgfpathlineto{\pgfqpoint{0.836238in}{1.273292in}}%
\pgfpathlineto{\pgfqpoint{0.804877in}{1.266436in}}%
\pgfpathlineto{\pgfqpoint{0.774538in}{1.258849in}}%
\pgfpathlineto{\pgfqpoint{0.781865in}{1.243802in}}%
\pgfpathlineto{\pgfqpoint{0.789201in}{1.228737in}}%
\pgfpathlineto{\pgfqpoint{0.818098in}{1.236046in}}%
\pgfpathlineto{\pgfqpoint{0.847980in}{1.242653in}}%
\pgfpathclose%
\pgfusepath{fill}%
\end{pgfscope}%
\begin{pgfscope}%
\pgfpathrectangle{\pgfqpoint{0.407592in}{0.222904in}}{\pgfqpoint{1.411163in}{1.411163in}}%
\pgfusepath{clip}%
\pgfsetbuttcap%
\pgfsetroundjoin%
\definecolor{currentfill}{rgb}{0.749020,0.772549,0.792157}%
\pgfsetfillcolor{currentfill}%
\pgfsetfillopacity{0.160000}%
\pgfsetlinewidth{0.000000pt}%
\definecolor{currentstroke}{rgb}{0.000000,0.000000,0.000000}%
\pgfsetstrokecolor{currentstroke}%
\pgfsetdash{}{0pt}%
\pgfpathmoveto{\pgfqpoint{1.150294in}{1.075192in}}%
\pgfpathlineto{\pgfqpoint{1.152381in}{1.091334in}}%
\pgfpathlineto{\pgfqpoint{1.154463in}{1.107447in}}%
\pgfpathlineto{\pgfqpoint{1.129303in}{1.108335in}}%
\pgfpathlineto{\pgfqpoint{1.103985in}{1.108589in}}%
\pgfpathlineto{\pgfqpoint{1.103539in}{1.092445in}}%
\pgfpathlineto{\pgfqpoint{1.103093in}{1.076271in}}%
\pgfpathlineto{\pgfqpoint{1.126767in}{1.076031in}}%
\pgfpathlineto{\pgfqpoint{1.150294in}{1.075192in}}%
\pgfpathclose%
\pgfusepath{fill}%
\end{pgfscope}%
\begin{pgfscope}%
\pgfpathrectangle{\pgfqpoint{0.407592in}{0.222904in}}{\pgfqpoint{1.411163in}{1.411163in}}%
\pgfusepath{clip}%
\pgfsetbuttcap%
\pgfsetroundjoin%
\definecolor{currentfill}{rgb}{0.749020,0.772549,0.792157}%
\pgfsetfillcolor{currentfill}%
\pgfsetfillopacity{0.160000}%
\pgfsetlinewidth{0.000000pt}%
\definecolor{currentstroke}{rgb}{0.000000,0.000000,0.000000}%
\pgfsetstrokecolor{currentstroke}%
\pgfsetdash{}{0pt}%
\pgfpathmoveto{\pgfqpoint{1.103093in}{1.076271in}}%
\pgfpathlineto{\pgfqpoint{1.103539in}{1.092445in}}%
\pgfpathlineto{\pgfqpoint{1.103985in}{1.108589in}}%
\pgfpathlineto{\pgfqpoint{1.078606in}{1.108208in}}%
\pgfpathlineto{\pgfqpoint{1.053260in}{1.107192in}}%
\pgfpathlineto{\pgfqpoint{1.054459in}{1.091087in}}%
\pgfpathlineto{\pgfqpoint{1.055660in}{1.074951in}}%
\pgfpathlineto{\pgfqpoint{1.079360in}{1.075910in}}%
\pgfpathlineto{\pgfqpoint{1.103093in}{1.076271in}}%
\pgfpathclose%
\pgfusepath{fill}%
\end{pgfscope}%
\begin{pgfscope}%
\pgfpathrectangle{\pgfqpoint{0.407592in}{0.222904in}}{\pgfqpoint{1.411163in}{1.411163in}}%
\pgfusepath{clip}%
\pgfsetbuttcap%
\pgfsetroundjoin%
\definecolor{currentfill}{rgb}{0.749020,0.772549,0.792157}%
\pgfsetfillcolor{currentfill}%
\pgfsetfillopacity{0.160000}%
\pgfsetlinewidth{0.000000pt}%
\definecolor{currentstroke}{rgb}{0.000000,0.000000,0.000000}%
\pgfsetstrokecolor{currentstroke}%
\pgfsetdash{}{0pt}%
\pgfpathmoveto{\pgfqpoint{1.545062in}{1.331713in}}%
\pgfpathlineto{\pgfqpoint{1.554310in}{1.346377in}}%
\pgfpathlineto{\pgfqpoint{1.563549in}{1.361026in}}%
\pgfpathlineto{\pgfqpoint{1.532332in}{1.371089in}}%
\pgfpathlineto{\pgfqpoint{1.499510in}{1.380358in}}%
\pgfpathlineto{\pgfqpoint{1.491575in}{1.365433in}}%
\pgfpathlineto{\pgfqpoint{1.483630in}{1.350488in}}%
\pgfpathlineto{\pgfqpoint{1.515124in}{1.341485in}}%
\pgfpathlineto{\pgfqpoint{1.545062in}{1.331713in}}%
\pgfpathclose%
\pgfusepath{fill}%
\end{pgfscope}%
\begin{pgfscope}%
\pgfpathrectangle{\pgfqpoint{0.407592in}{0.222904in}}{\pgfqpoint{1.411163in}{1.411163in}}%
\pgfusepath{clip}%
\pgfsetbuttcap%
\pgfsetroundjoin%
\definecolor{currentfill}{rgb}{0.749020,0.772549,0.792157}%
\pgfsetfillcolor{currentfill}%
\pgfsetfillopacity{0.160000}%
\pgfsetlinewidth{0.000000pt}%
\definecolor{currentstroke}{rgb}{0.000000,0.000000,0.000000}%
\pgfsetstrokecolor{currentstroke}%
\pgfsetdash{}{0pt}%
\pgfpathmoveto{\pgfqpoint{0.936476in}{1.159915in}}%
\pgfpathlineto{\pgfqpoint{0.932084in}{1.175607in}}%
\pgfpathlineto{\pgfqpoint{0.927699in}{1.191274in}}%
\pgfpathlineto{\pgfqpoint{0.899286in}{1.186530in}}%
\pgfpathlineto{\pgfqpoint{0.871567in}{1.181107in}}%
\pgfpathlineto{\pgfqpoint{0.877485in}{1.165664in}}%
\pgfpathlineto{\pgfqpoint{0.883412in}{1.150198in}}%
\pgfpathlineto{\pgfqpoint{0.909612in}{1.155381in}}%
\pgfpathlineto{\pgfqpoint{0.936476in}{1.159915in}}%
\pgfpathclose%
\pgfusepath{fill}%
\end{pgfscope}%
\begin{pgfscope}%
\pgfpathrectangle{\pgfqpoint{0.407592in}{0.222904in}}{\pgfqpoint{1.411163in}{1.411163in}}%
\pgfusepath{clip}%
\pgfsetbuttcap%
\pgfsetroundjoin%
\definecolor{currentfill}{rgb}{0.749020,0.772549,0.792157}%
\pgfsetfillcolor{currentfill}%
\pgfsetfillopacity{0.160000}%
\pgfsetlinewidth{0.000000pt}%
\definecolor{currentstroke}{rgb}{0.000000,0.000000,0.000000}%
\pgfsetstrokecolor{currentstroke}%
\pgfsetdash{}{0pt}%
\pgfpathmoveto{\pgfqpoint{1.251652in}{1.097637in}}%
\pgfpathlineto{\pgfqpoint{1.256894in}{1.113454in}}%
\pgfpathlineto{\pgfqpoint{1.262126in}{1.129244in}}%
\pgfpathlineto{\pgfqpoint{1.236993in}{1.132803in}}%
\pgfpathlineto{\pgfqpoint{1.211298in}{1.135718in}}%
\pgfpathlineto{\pgfqpoint{1.207620in}{1.119764in}}%
\pgfpathlineto{\pgfqpoint{1.203936in}{1.103781in}}%
\pgfpathlineto{\pgfqpoint{1.228060in}{1.101014in}}%
\pgfpathlineto{\pgfqpoint{1.251652in}{1.097637in}}%
\pgfpathclose%
\pgfusepath{fill}%
\end{pgfscope}%
\begin{pgfscope}%
\pgfpathrectangle{\pgfqpoint{0.407592in}{0.222904in}}{\pgfqpoint{1.411163in}{1.411163in}}%
\pgfusepath{clip}%
\pgfsetbuttcap%
\pgfsetroundjoin%
\definecolor{currentfill}{rgb}{0.749020,0.772549,0.792157}%
\pgfsetfillcolor{currentfill}%
\pgfsetfillopacity{0.160000}%
\pgfsetlinewidth{0.000000pt}%
\definecolor{currentstroke}{rgb}{0.000000,0.000000,0.000000}%
\pgfsetstrokecolor{currentstroke}%
\pgfsetdash{}{0pt}%
\pgfpathmoveto{\pgfqpoint{1.403629in}{1.200012in}}%
\pgfpathlineto{\pgfqpoint{1.411674in}{1.215145in}}%
\pgfpathlineto{\pgfqpoint{1.419709in}{1.230258in}}%
\pgfpathlineto{\pgfqpoint{1.392279in}{1.237429in}}%
\pgfpathlineto{\pgfqpoint{1.363722in}{1.243892in}}%
\pgfpathlineto{\pgfqpoint{1.357082in}{1.228515in}}%
\pgfpathlineto{\pgfqpoint{1.350432in}{1.213116in}}%
\pgfpathlineto{\pgfqpoint{1.377571in}{1.206903in}}%
\pgfpathlineto{\pgfqpoint{1.403629in}{1.200012in}}%
\pgfpathclose%
\pgfusepath{fill}%
\end{pgfscope}%
\begin{pgfscope}%
\pgfpathrectangle{\pgfqpoint{0.407592in}{0.222904in}}{\pgfqpoint{1.411163in}{1.411163in}}%
\pgfusepath{clip}%
\pgfsetbuttcap%
\pgfsetroundjoin%
\definecolor{currentfill}{rgb}{0.749020,0.772549,0.792157}%
\pgfsetfillcolor{currentfill}%
\pgfsetfillopacity{0.160000}%
\pgfsetlinewidth{0.000000pt}%
\definecolor{currentstroke}{rgb}{0.000000,0.000000,0.000000}%
\pgfsetstrokecolor{currentstroke}%
\pgfsetdash{}{0pt}%
\pgfpathmoveto{\pgfqpoint{1.196547in}{1.071728in}}%
\pgfpathlineto{\pgfqpoint{1.200244in}{1.087769in}}%
\pgfpathlineto{\pgfqpoint{1.203936in}{1.103781in}}%
\pgfpathlineto{\pgfqpoint{1.179372in}{1.105927in}}%
\pgfpathlineto{\pgfqpoint{1.154463in}{1.107447in}}%
\pgfpathlineto{\pgfqpoint{1.152381in}{1.091334in}}%
\pgfpathlineto{\pgfqpoint{1.150294in}{1.075192in}}%
\pgfpathlineto{\pgfqpoint{1.173583in}{1.073756in}}%
\pgfpathlineto{\pgfqpoint{1.196547in}{1.071728in}}%
\pgfpathclose%
\pgfusepath{fill}%
\end{pgfscope}%
\begin{pgfscope}%
\pgfpathrectangle{\pgfqpoint{0.407592in}{0.222904in}}{\pgfqpoint{1.411163in}{1.411163in}}%
\pgfusepath{clip}%
\pgfsetbuttcap%
\pgfsetroundjoin%
\definecolor{currentfill}{rgb}{0.749020,0.772549,0.792157}%
\pgfsetfillcolor{currentfill}%
\pgfsetfillopacity{0.160000}%
\pgfsetlinewidth{0.000000pt}%
\definecolor{currentstroke}{rgb}{0.000000,0.000000,0.000000}%
\pgfsetstrokecolor{currentstroke}%
\pgfsetdash{}{0pt}%
\pgfpathmoveto{\pgfqpoint{1.003050in}{1.103275in}}%
\pgfpathlineto{\pgfqpoint{1.000225in}{1.119245in}}%
\pgfpathlineto{\pgfqpoint{0.997404in}{1.135186in}}%
\pgfpathlineto{\pgfqpoint{0.971125in}{1.132140in}}%
\pgfpathlineto{\pgfqpoint{0.945281in}{1.128453in}}%
\pgfpathlineto{\pgfqpoint{0.949694in}{1.112683in}}%
\pgfpathlineto{\pgfqpoint{0.954115in}{1.096887in}}%
\pgfpathlineto{\pgfqpoint{0.978376in}{1.100385in}}%
\pgfpathlineto{\pgfqpoint{1.003050in}{1.103275in}}%
\pgfpathclose%
\pgfusepath{fill}%
\end{pgfscope}%
\begin{pgfscope}%
\pgfpathrectangle{\pgfqpoint{0.407592in}{0.222904in}}{\pgfqpoint{1.411163in}{1.411163in}}%
\pgfusepath{clip}%
\pgfsetbuttcap%
\pgfsetroundjoin%
\definecolor{currentfill}{rgb}{0.749020,0.772549,0.792157}%
\pgfsetfillcolor{currentfill}%
\pgfsetfillopacity{0.160000}%
\pgfsetlinewidth{0.000000pt}%
\definecolor{currentstroke}{rgb}{0.000000,0.000000,0.000000}%
\pgfsetstrokecolor{currentstroke}%
\pgfsetdash{}{0pt}%
\pgfpathmoveto{\pgfqpoint{0.730764in}{1.348745in}}%
\pgfpathlineto{\pgfqpoint{0.723500in}{1.363663in}}%
\pgfpathlineto{\pgfqpoint{0.716244in}{1.378564in}}%
\pgfpathlineto{\pgfqpoint{0.681443in}{1.369133in}}%
\pgfpathlineto{\pgfqpoint{0.648081in}{1.358914in}}%
\pgfpathlineto{\pgfqpoint{0.656718in}{1.344295in}}%
\pgfpathlineto{\pgfqpoint{0.665363in}{1.329663in}}%
\pgfpathlineto{\pgfqpoint{0.697365in}{1.339586in}}%
\pgfpathlineto{\pgfqpoint{0.730764in}{1.348745in}}%
\pgfpathclose%
\pgfusepath{fill}%
\end{pgfscope}%
\begin{pgfscope}%
\pgfpathrectangle{\pgfqpoint{0.407592in}{0.222904in}}{\pgfqpoint{1.411163in}{1.411163in}}%
\pgfusepath{clip}%
\pgfsetbuttcap%
\pgfsetroundjoin%
\definecolor{currentfill}{rgb}{0.749020,0.772549,0.792157}%
\pgfsetfillcolor{currentfill}%
\pgfsetfillopacity{0.160000}%
\pgfsetlinewidth{0.000000pt}%
\definecolor{currentstroke}{rgb}{0.000000,0.000000,0.000000}%
\pgfsetstrokecolor{currentstroke}%
\pgfsetdash{}{0pt}%
\pgfpathmoveto{\pgfqpoint{1.055660in}{1.074951in}}%
\pgfpathlineto{\pgfqpoint{1.054459in}{1.091087in}}%
\pgfpathlineto{\pgfqpoint{1.053260in}{1.107192in}}%
\pgfpathlineto{\pgfqpoint{1.028043in}{1.105547in}}%
\pgfpathlineto{\pgfqpoint{1.003050in}{1.103275in}}%
\pgfpathlineto{\pgfqpoint{1.005881in}{1.087277in}}%
\pgfpathlineto{\pgfqpoint{1.008716in}{1.071250in}}%
\pgfpathlineto{\pgfqpoint{1.032082in}{1.073396in}}%
\pgfpathlineto{\pgfqpoint{1.055660in}{1.074951in}}%
\pgfpathclose%
\pgfusepath{fill}%
\end{pgfscope}%
\begin{pgfscope}%
\pgfpathrectangle{\pgfqpoint{0.407592in}{0.222904in}}{\pgfqpoint{1.411163in}{1.411163in}}%
\pgfusepath{clip}%
\pgfsetbuttcap%
\pgfsetroundjoin%
\definecolor{currentfill}{rgb}{0.749020,0.772549,0.792157}%
\pgfsetfillcolor{currentfill}%
\pgfsetfillopacity{0.160000}%
\pgfsetlinewidth{0.000000pt}%
\definecolor{currentstroke}{rgb}{0.000000,0.000000,0.000000}%
\pgfsetstrokecolor{currentstroke}%
\pgfsetdash{}{0pt}%
\pgfpathmoveto{\pgfqpoint{0.859756in}{1.211925in}}%
\pgfpathlineto{\pgfqpoint{0.853864in}{1.227300in}}%
\pgfpathlineto{\pgfqpoint{0.847980in}{1.242653in}}%
\pgfpathlineto{\pgfqpoint{0.818098in}{1.236046in}}%
\pgfpathlineto{\pgfqpoint{0.789201in}{1.228737in}}%
\pgfpathlineto{\pgfqpoint{0.796546in}{1.213653in}}%
\pgfpathlineto{\pgfqpoint{0.803900in}{1.198550in}}%
\pgfpathlineto{\pgfqpoint{0.831355in}{1.205574in}}%
\pgfpathlineto{\pgfqpoint{0.859756in}{1.211925in}}%
\pgfpathclose%
\pgfusepath{fill}%
\end{pgfscope}%
\begin{pgfscope}%
\pgfpathrectangle{\pgfqpoint{0.407592in}{0.222904in}}{\pgfqpoint{1.411163in}{1.411163in}}%
\pgfusepath{clip}%
\pgfsetbuttcap%
\pgfsetroundjoin%
\definecolor{currentfill}{rgb}{0.749020,0.772549,0.792157}%
\pgfsetfillcolor{currentfill}%
\pgfsetfillopacity{0.160000}%
\pgfsetlinewidth{0.000000pt}%
\definecolor{currentstroke}{rgb}{0.000000,0.000000,0.000000}%
\pgfsetstrokecolor{currentstroke}%
\pgfsetdash{}{0pt}%
\pgfpathmoveto{\pgfqpoint{1.526539in}{1.302343in}}%
\pgfpathlineto{\pgfqpoint{1.535805in}{1.317036in}}%
\pgfpathlineto{\pgfqpoint{1.545062in}{1.331713in}}%
\pgfpathlineto{\pgfqpoint{1.515124in}{1.341485in}}%
\pgfpathlineto{\pgfqpoint{1.483630in}{1.350488in}}%
\pgfpathlineto{\pgfqpoint{1.475675in}{1.335525in}}%
\pgfpathlineto{\pgfqpoint{1.467710in}{1.320544in}}%
\pgfpathlineto{\pgfqpoint{1.497877in}{1.311815in}}%
\pgfpathlineto{\pgfqpoint{1.526539in}{1.302343in}}%
\pgfpathclose%
\pgfusepath{fill}%
\end{pgfscope}%
\begin{pgfscope}%
\pgfpathrectangle{\pgfqpoint{0.407592in}{0.222904in}}{\pgfqpoint{1.411163in}{1.411163in}}%
\pgfusepath{clip}%
\pgfsetbuttcap%
\pgfsetroundjoin%
\definecolor{currentfill}{rgb}{0.749020,0.772549,0.792157}%
\pgfsetfillcolor{currentfill}%
\pgfsetfillopacity{0.160000}%
\pgfsetlinewidth{0.000000pt}%
\definecolor{currentstroke}{rgb}{0.000000,0.000000,0.000000}%
\pgfsetstrokecolor{currentstroke}%
\pgfsetdash{}{0pt}%
\pgfpathmoveto{\pgfqpoint{1.310327in}{1.120238in}}%
\pgfpathlineto{\pgfqpoint{1.317036in}{1.135775in}}%
\pgfpathlineto{\pgfqpoint{1.323735in}{1.151289in}}%
\pgfpathlineto{\pgfqpoint{1.298546in}{1.156343in}}%
\pgfpathlineto{\pgfqpoint{1.272566in}{1.160745in}}%
\pgfpathlineto{\pgfqpoint{1.267351in}{1.145008in}}%
\pgfpathlineto{\pgfqpoint{1.262126in}{1.129244in}}%
\pgfpathlineto{\pgfqpoint{1.286603in}{1.125051in}}%
\pgfpathlineto{\pgfqpoint{1.310327in}{1.120238in}}%
\pgfpathclose%
\pgfusepath{fill}%
\end{pgfscope}%
\begin{pgfscope}%
\pgfpathrectangle{\pgfqpoint{0.407592in}{0.222904in}}{\pgfqpoint{1.411163in}{1.411163in}}%
\pgfusepath{clip}%
\pgfsetbuttcap%
\pgfsetroundjoin%
\definecolor{currentfill}{rgb}{0.749020,0.772549,0.792157}%
\pgfsetfillcolor{currentfill}%
\pgfsetfillopacity{0.160000}%
\pgfsetlinewidth{0.000000pt}%
\definecolor{currentstroke}{rgb}{0.000000,0.000000,0.000000}%
\pgfsetstrokecolor{currentstroke}%
\pgfsetdash{}{0pt}%
\pgfpathmoveto{\pgfqpoint{0.945281in}{1.128453in}}%
\pgfpathlineto{\pgfqpoint{0.940875in}{1.144197in}}%
\pgfpathlineto{\pgfqpoint{0.936476in}{1.159915in}}%
\pgfpathlineto{\pgfqpoint{0.909612in}{1.155381in}}%
\pgfpathlineto{\pgfqpoint{0.883412in}{1.150198in}}%
\pgfpathlineto{\pgfqpoint{0.889347in}{1.134710in}}%
\pgfpathlineto{\pgfqpoint{0.895292in}{1.119199in}}%
\pgfpathlineto{\pgfqpoint{0.919970in}{1.124135in}}%
\pgfpathlineto{\pgfqpoint{0.945281in}{1.128453in}}%
\pgfpathclose%
\pgfusepath{fill}%
\end{pgfscope}%
\begin{pgfscope}%
\pgfpathrectangle{\pgfqpoint{0.407592in}{0.222904in}}{\pgfqpoint{1.411163in}{1.411163in}}%
\pgfusepath{clip}%
\pgfsetbuttcap%
\pgfsetroundjoin%
\definecolor{currentfill}{rgb}{0.749020,0.772549,0.792157}%
\pgfsetfillcolor{currentfill}%
\pgfsetfillopacity{0.160000}%
\pgfsetlinewidth{0.000000pt}%
\definecolor{currentstroke}{rgb}{0.000000,0.000000,0.000000}%
\pgfsetstrokecolor{currentstroke}%
\pgfsetdash{}{0pt}%
\pgfpathmoveto{\pgfqpoint{0.745320in}{1.318853in}}%
\pgfpathlineto{\pgfqpoint{0.738037in}{1.333808in}}%
\pgfpathlineto{\pgfqpoint{0.730764in}{1.348745in}}%
\pgfpathlineto{\pgfqpoint{0.697365in}{1.339586in}}%
\pgfpathlineto{\pgfqpoint{0.665363in}{1.329663in}}%
\pgfpathlineto{\pgfqpoint{0.674016in}{1.315016in}}%
\pgfpathlineto{\pgfqpoint{0.682677in}{1.300356in}}%
\pgfpathlineto{\pgfqpoint{0.713322in}{1.309973in}}%
\pgfpathlineto{\pgfqpoint{0.745320in}{1.318853in}}%
\pgfpathclose%
\pgfusepath{fill}%
\end{pgfscope}%
\begin{pgfscope}%
\pgfpathrectangle{\pgfqpoint{0.407592in}{0.222904in}}{\pgfqpoint{1.411163in}{1.411163in}}%
\pgfusepath{clip}%
\pgfsetbuttcap%
\pgfsetroundjoin%
\definecolor{currentfill}{rgb}{0.749020,0.772549,0.792157}%
\pgfsetfillcolor{currentfill}%
\pgfsetfillopacity{0.160000}%
\pgfsetlinewidth{0.000000pt}%
\definecolor{currentstroke}{rgb}{0.000000,0.000000,0.000000}%
\pgfsetstrokecolor{currentstroke}%
\pgfsetdash{}{0pt}%
\pgfpathmoveto{\pgfqpoint{1.146109in}{1.042815in}}%
\pgfpathlineto{\pgfqpoint{1.148203in}{1.059018in}}%
\pgfpathlineto{\pgfqpoint{1.150294in}{1.075192in}}%
\pgfpathlineto{\pgfqpoint{1.126767in}{1.076031in}}%
\pgfpathlineto{\pgfqpoint{1.103093in}{1.076271in}}%
\pgfpathlineto{\pgfqpoint{1.102645in}{1.060065in}}%
\pgfpathlineto{\pgfqpoint{1.102197in}{1.043829in}}%
\pgfpathlineto{\pgfqpoint{1.124222in}{1.043604in}}%
\pgfpathlineto{\pgfqpoint{1.146109in}{1.042815in}}%
\pgfpathclose%
\pgfusepath{fill}%
\end{pgfscope}%
\begin{pgfscope}%
\pgfpathrectangle{\pgfqpoint{0.407592in}{0.222904in}}{\pgfqpoint{1.411163in}{1.411163in}}%
\pgfusepath{clip}%
\pgfsetbuttcap%
\pgfsetroundjoin%
\definecolor{currentfill}{rgb}{0.749020,0.772549,0.792157}%
\pgfsetfillcolor{currentfill}%
\pgfsetfillopacity{0.160000}%
\pgfsetlinewidth{0.000000pt}%
\definecolor{currentstroke}{rgb}{0.000000,0.000000,0.000000}%
\pgfsetstrokecolor{currentstroke}%
\pgfsetdash{}{0pt}%
\pgfpathmoveto{\pgfqpoint{1.102197in}{1.043829in}}%
\pgfpathlineto{\pgfqpoint{1.102645in}{1.060065in}}%
\pgfpathlineto{\pgfqpoint{1.103093in}{1.076271in}}%
\pgfpathlineto{\pgfqpoint{1.079360in}{1.075910in}}%
\pgfpathlineto{\pgfqpoint{1.055660in}{1.074951in}}%
\pgfpathlineto{\pgfqpoint{1.056863in}{1.058785in}}%
\pgfpathlineto{\pgfqpoint{1.058069in}{1.042589in}}%
\pgfpathlineto{\pgfqpoint{1.080117in}{1.043491in}}%
\pgfpathlineto{\pgfqpoint{1.102197in}{1.043829in}}%
\pgfpathclose%
\pgfusepath{fill}%
\end{pgfscope}%
\begin{pgfscope}%
\pgfpathrectangle{\pgfqpoint{0.407592in}{0.222904in}}{\pgfqpoint{1.411163in}{1.411163in}}%
\pgfusepath{clip}%
\pgfsetbuttcap%
\pgfsetroundjoin%
\definecolor{currentfill}{rgb}{0.749020,0.772549,0.792157}%
\pgfsetfillcolor{currentfill}%
\pgfsetfillopacity{0.160000}%
\pgfsetlinewidth{0.000000pt}%
\definecolor{currentstroke}{rgb}{0.000000,0.000000,0.000000}%
\pgfsetstrokecolor{currentstroke}%
\pgfsetdash{}{0pt}%
\pgfpathmoveto{\pgfqpoint{1.387508in}{1.169690in}}%
\pgfpathlineto{\pgfqpoint{1.395573in}{1.184860in}}%
\pgfpathlineto{\pgfqpoint{1.403629in}{1.200012in}}%
\pgfpathlineto{\pgfqpoint{1.377571in}{1.206903in}}%
\pgfpathlineto{\pgfqpoint{1.350432in}{1.213116in}}%
\pgfpathlineto{\pgfqpoint{1.343773in}{1.197694in}}%
\pgfpathlineto{\pgfqpoint{1.337103in}{1.182248in}}%
\pgfpathlineto{\pgfqpoint{1.362823in}{1.176293in}}%
\pgfpathlineto{\pgfqpoint{1.387508in}{1.169690in}}%
\pgfpathclose%
\pgfusepath{fill}%
\end{pgfscope}%
\begin{pgfscope}%
\pgfpathrectangle{\pgfqpoint{0.407592in}{0.222904in}}{\pgfqpoint{1.411163in}{1.411163in}}%
\pgfusepath{clip}%
\pgfsetbuttcap%
\pgfsetroundjoin%
\definecolor{currentfill}{rgb}{0.749020,0.772549,0.792157}%
\pgfsetfillcolor{currentfill}%
\pgfsetfillopacity{0.160000}%
\pgfsetlinewidth{0.000000pt}%
\definecolor{currentstroke}{rgb}{0.000000,0.000000,0.000000}%
\pgfsetstrokecolor{currentstroke}%
\pgfsetdash{}{0pt}%
\pgfpathmoveto{\pgfqpoint{1.241143in}{1.065924in}}%
\pgfpathlineto{\pgfqpoint{1.246402in}{1.081794in}}%
\pgfpathlineto{\pgfqpoint{1.251652in}{1.097637in}}%
\pgfpathlineto{\pgfqpoint{1.228060in}{1.101014in}}%
\pgfpathlineto{\pgfqpoint{1.203936in}{1.103781in}}%
\pgfpathlineto{\pgfqpoint{1.200244in}{1.087769in}}%
\pgfpathlineto{\pgfqpoint{1.196547in}{1.071728in}}%
\pgfpathlineto{\pgfqpoint{1.219096in}{1.069114in}}%
\pgfpathlineto{\pgfqpoint{1.241143in}{1.065924in}}%
\pgfpathclose%
\pgfusepath{fill}%
\end{pgfscope}%
\begin{pgfscope}%
\pgfpathrectangle{\pgfqpoint{0.407592in}{0.222904in}}{\pgfqpoint{1.411163in}{1.411163in}}%
\pgfusepath{clip}%
\pgfsetbuttcap%
\pgfsetroundjoin%
\definecolor{currentfill}{rgb}{0.749020,0.772549,0.792157}%
\pgfsetfillcolor{currentfill}%
\pgfsetfillopacity{0.160000}%
\pgfsetlinewidth{0.000000pt}%
\definecolor{currentstroke}{rgb}{0.000000,0.000000,0.000000}%
\pgfsetstrokecolor{currentstroke}%
\pgfsetdash{}{0pt}%
\pgfpathmoveto{\pgfqpoint{1.507980in}{1.272915in}}%
\pgfpathlineto{\pgfqpoint{1.517264in}{1.287636in}}%
\pgfpathlineto{\pgfqpoint{1.526539in}{1.302343in}}%
\pgfpathlineto{\pgfqpoint{1.497877in}{1.311815in}}%
\pgfpathlineto{\pgfqpoint{1.467710in}{1.320544in}}%
\pgfpathlineto{\pgfqpoint{1.459735in}{1.305543in}}%
\pgfpathlineto{\pgfqpoint{1.451750in}{1.290524in}}%
\pgfpathlineto{\pgfqpoint{1.480591in}{1.282078in}}%
\pgfpathlineto{\pgfqpoint{1.507980in}{1.272915in}}%
\pgfpathclose%
\pgfusepath{fill}%
\end{pgfscope}%
\begin{pgfscope}%
\pgfpathrectangle{\pgfqpoint{0.407592in}{0.222904in}}{\pgfqpoint{1.411163in}{1.411163in}}%
\pgfusepath{clip}%
\pgfsetbuttcap%
\pgfsetroundjoin%
\definecolor{currentfill}{rgb}{0.749020,0.772549,0.792157}%
\pgfsetfillcolor{currentfill}%
\pgfsetfillopacity{0.160000}%
\pgfsetlinewidth{0.000000pt}%
\definecolor{currentstroke}{rgb}{0.000000,0.000000,0.000000}%
\pgfsetstrokecolor{currentstroke}%
\pgfsetdash{}{0pt}%
\pgfpathmoveto{\pgfqpoint{1.008716in}{1.071250in}}%
\pgfpathlineto{\pgfqpoint{1.005881in}{1.087277in}}%
\pgfpathlineto{\pgfqpoint{1.003050in}{1.103275in}}%
\pgfpathlineto{\pgfqpoint{0.978376in}{1.100385in}}%
\pgfpathlineto{\pgfqpoint{0.954115in}{1.096887in}}%
\pgfpathlineto{\pgfqpoint{0.958543in}{1.081064in}}%
\pgfpathlineto{\pgfqpoint{0.962979in}{1.065216in}}%
\pgfpathlineto{\pgfqpoint{0.985652in}{1.068520in}}%
\pgfpathlineto{\pgfqpoint{1.008716in}{1.071250in}}%
\pgfpathclose%
\pgfusepath{fill}%
\end{pgfscope}%
\begin{pgfscope}%
\pgfpathrectangle{\pgfqpoint{0.407592in}{0.222904in}}{\pgfqpoint{1.411163in}{1.411163in}}%
\pgfusepath{clip}%
\pgfsetbuttcap%
\pgfsetroundjoin%
\definecolor{currentfill}{rgb}{0.749020,0.772549,0.792157}%
\pgfsetfillcolor{currentfill}%
\pgfsetfillopacity{0.160000}%
\pgfsetlinewidth{0.000000pt}%
\definecolor{currentstroke}{rgb}{0.000000,0.000000,0.000000}%
\pgfsetstrokecolor{currentstroke}%
\pgfsetdash{}{0pt}%
\pgfpathmoveto{\pgfqpoint{0.871567in}{1.181107in}}%
\pgfpathlineto{\pgfqpoint{0.865657in}{1.196527in}}%
\pgfpathlineto{\pgfqpoint{0.859756in}{1.211925in}}%
\pgfpathlineto{\pgfqpoint{0.831355in}{1.205574in}}%
\pgfpathlineto{\pgfqpoint{0.803900in}{1.198550in}}%
\pgfpathlineto{\pgfqpoint{0.811263in}{1.183428in}}%
\pgfpathlineto{\pgfqpoint{0.818636in}{1.168289in}}%
\pgfpathlineto{\pgfqpoint{0.844648in}{1.175020in}}%
\pgfpathlineto{\pgfqpoint{0.871567in}{1.181107in}}%
\pgfpathclose%
\pgfusepath{fill}%
\end{pgfscope}%
\begin{pgfscope}%
\pgfpathrectangle{\pgfqpoint{0.407592in}{0.222904in}}{\pgfqpoint{1.411163in}{1.411163in}}%
\pgfusepath{clip}%
\pgfsetbuttcap%
\pgfsetroundjoin%
\definecolor{currentfill}{rgb}{0.749020,0.772549,0.792157}%
\pgfsetfillcolor{currentfill}%
\pgfsetfillopacity{0.160000}%
\pgfsetlinewidth{0.000000pt}%
\definecolor{currentstroke}{rgb}{0.000000,0.000000,0.000000}%
\pgfsetstrokecolor{currentstroke}%
\pgfsetdash{}{0pt}%
\pgfpathmoveto{\pgfqpoint{1.189131in}{1.039559in}}%
\pgfpathlineto{\pgfqpoint{1.192842in}{1.055658in}}%
\pgfpathlineto{\pgfqpoint{1.196547in}{1.071728in}}%
\pgfpathlineto{\pgfqpoint{1.173583in}{1.073756in}}%
\pgfpathlineto{\pgfqpoint{1.150294in}{1.075192in}}%
\pgfpathlineto{\pgfqpoint{1.148203in}{1.059018in}}%
\pgfpathlineto{\pgfqpoint{1.146109in}{1.042815in}}%
\pgfpathlineto{\pgfqpoint{1.167773in}{1.041465in}}%
\pgfpathlineto{\pgfqpoint{1.189131in}{1.039559in}}%
\pgfpathclose%
\pgfusepath{fill}%
\end{pgfscope}%
\begin{pgfscope}%
\pgfpathrectangle{\pgfqpoint{0.407592in}{0.222904in}}{\pgfqpoint{1.411163in}{1.411163in}}%
\pgfusepath{clip}%
\pgfsetbuttcap%
\pgfsetroundjoin%
\definecolor{currentfill}{rgb}{0.749020,0.772549,0.792157}%
\pgfsetfillcolor{currentfill}%
\pgfsetfillopacity{0.160000}%
\pgfsetlinewidth{0.000000pt}%
\definecolor{currentstroke}{rgb}{0.000000,0.000000,0.000000}%
\pgfsetstrokecolor{currentstroke}%
\pgfsetdash{}{0pt}%
\pgfpathmoveto{\pgfqpoint{1.058069in}{1.042589in}}%
\pgfpathlineto{\pgfqpoint{1.056863in}{1.058785in}}%
\pgfpathlineto{\pgfqpoint{1.055660in}{1.074951in}}%
\pgfpathlineto{\pgfqpoint{1.032082in}{1.073396in}}%
\pgfpathlineto{\pgfqpoint{1.008716in}{1.071250in}}%
\pgfpathlineto{\pgfqpoint{1.011557in}{1.055195in}}%
\pgfpathlineto{\pgfqpoint{1.014403in}{1.039110in}}%
\pgfpathlineto{\pgfqpoint{1.036136in}{1.041127in}}%
\pgfpathlineto{\pgfqpoint{1.058069in}{1.042589in}}%
\pgfpathclose%
\pgfusepath{fill}%
\end{pgfscope}%
\begin{pgfscope}%
\pgfpathrectangle{\pgfqpoint{0.407592in}{0.222904in}}{\pgfqpoint{1.411163in}{1.411163in}}%
\pgfusepath{clip}%
\pgfsetbuttcap%
\pgfsetroundjoin%
\definecolor{currentfill}{rgb}{0.749020,0.772549,0.792157}%
\pgfsetfillcolor{currentfill}%
\pgfsetfillopacity{0.160000}%
\pgfsetlinewidth{0.000000pt}%
\definecolor{currentstroke}{rgb}{0.000000,0.000000,0.000000}%
\pgfsetstrokecolor{currentstroke}%
\pgfsetdash{}{0pt}%
\pgfpathmoveto{\pgfqpoint{0.759911in}{1.288888in}}%
\pgfpathlineto{\pgfqpoint{0.752611in}{1.303880in}}%
\pgfpathlineto{\pgfqpoint{0.745320in}{1.318853in}}%
\pgfpathlineto{\pgfqpoint{0.713322in}{1.309973in}}%
\pgfpathlineto{\pgfqpoint{0.682677in}{1.300356in}}%
\pgfpathlineto{\pgfqpoint{0.691347in}{1.285681in}}%
\pgfpathlineto{\pgfqpoint{0.700025in}{1.270993in}}%
\pgfpathlineto{\pgfqpoint{0.729314in}{1.280296in}}%
\pgfpathlineto{\pgfqpoint{0.759911in}{1.288888in}}%
\pgfpathclose%
\pgfusepath{fill}%
\end{pgfscope}%
\begin{pgfscope}%
\pgfpathrectangle{\pgfqpoint{0.407592in}{0.222904in}}{\pgfqpoint{1.411163in}{1.411163in}}%
\pgfusepath{clip}%
\pgfsetbuttcap%
\pgfsetroundjoin%
\definecolor{currentfill}{rgb}{0.749020,0.772549,0.792157}%
\pgfsetfillcolor{currentfill}%
\pgfsetfillopacity{0.160000}%
\pgfsetlinewidth{0.000000pt}%
\definecolor{currentstroke}{rgb}{0.000000,0.000000,0.000000}%
\pgfsetstrokecolor{currentstroke}%
\pgfsetdash{}{0pt}%
\pgfpathmoveto{\pgfqpoint{1.296878in}{1.089094in}}%
\pgfpathlineto{\pgfqpoint{1.303607in}{1.104677in}}%
\pgfpathlineto{\pgfqpoint{1.310327in}{1.120238in}}%
\pgfpathlineto{\pgfqpoint{1.286603in}{1.125051in}}%
\pgfpathlineto{\pgfqpoint{1.262126in}{1.129244in}}%
\pgfpathlineto{\pgfqpoint{1.256894in}{1.113454in}}%
\pgfpathlineto{\pgfqpoint{1.251652in}{1.097637in}}%
\pgfpathlineto{\pgfqpoint{1.274622in}{1.093659in}}%
\pgfpathlineto{\pgfqpoint{1.296878in}{1.089094in}}%
\pgfpathclose%
\pgfusepath{fill}%
\end{pgfscope}%
\begin{pgfscope}%
\pgfpathrectangle{\pgfqpoint{0.407592in}{0.222904in}}{\pgfqpoint{1.411163in}{1.411163in}}%
\pgfusepath{clip}%
\pgfsetbuttcap%
\pgfsetroundjoin%
\definecolor{currentfill}{rgb}{0.749020,0.772549,0.792157}%
\pgfsetfillcolor{currentfill}%
\pgfsetfillopacity{0.160000}%
\pgfsetlinewidth{0.000000pt}%
\definecolor{currentstroke}{rgb}{0.000000,0.000000,0.000000}%
\pgfsetstrokecolor{currentstroke}%
\pgfsetdash{}{0pt}%
\pgfpathmoveto{\pgfqpoint{0.954115in}{1.096887in}}%
\pgfpathlineto{\pgfqpoint{0.949694in}{1.112683in}}%
\pgfpathlineto{\pgfqpoint{0.945281in}{1.128453in}}%
\pgfpathlineto{\pgfqpoint{0.919970in}{1.124135in}}%
\pgfpathlineto{\pgfqpoint{0.895292in}{1.119199in}}%
\pgfpathlineto{\pgfqpoint{0.901245in}{1.103665in}}%
\pgfpathlineto{\pgfqpoint{0.907207in}{1.088109in}}%
\pgfpathlineto{\pgfqpoint{0.930361in}{1.092790in}}%
\pgfpathlineto{\pgfqpoint{0.954115in}{1.096887in}}%
\pgfpathclose%
\pgfusepath{fill}%
\end{pgfscope}%
\begin{pgfscope}%
\pgfpathrectangle{\pgfqpoint{0.407592in}{0.222904in}}{\pgfqpoint{1.411163in}{1.411163in}}%
\pgfusepath{clip}%
\pgfsetbuttcap%
\pgfsetroundjoin%
\definecolor{currentfill}{rgb}{0.749020,0.772549,0.792157}%
\pgfsetfillcolor{currentfill}%
\pgfsetfillopacity{0.160000}%
\pgfsetlinewidth{0.000000pt}%
\definecolor{currentstroke}{rgb}{0.000000,0.000000,0.000000}%
\pgfsetstrokecolor{currentstroke}%
\pgfsetdash{}{0pt}%
\pgfpathmoveto{\pgfqpoint{1.489383in}{1.243429in}}%
\pgfpathlineto{\pgfqpoint{1.498686in}{1.258179in}}%
\pgfpathlineto{\pgfqpoint{1.507980in}{1.272915in}}%
\pgfpathlineto{\pgfqpoint{1.480591in}{1.282078in}}%
\pgfpathlineto{\pgfqpoint{1.451750in}{1.290524in}}%
\pgfpathlineto{\pgfqpoint{1.443755in}{1.275486in}}%
\pgfpathlineto{\pgfqpoint{1.435749in}{1.260429in}}%
\pgfpathlineto{\pgfqpoint{1.463266in}{1.252274in}}%
\pgfpathlineto{\pgfqpoint{1.489383in}{1.243429in}}%
\pgfpathclose%
\pgfusepath{fill}%
\end{pgfscope}%
\begin{pgfscope}%
\pgfpathrectangle{\pgfqpoint{0.407592in}{0.222904in}}{\pgfqpoint{1.411163in}{1.411163in}}%
\pgfusepath{clip}%
\pgfsetbuttcap%
\pgfsetroundjoin%
\definecolor{currentfill}{rgb}{0.749020,0.772549,0.792157}%
\pgfsetfillcolor{currentfill}%
\pgfsetfillopacity{0.160000}%
\pgfsetlinewidth{0.000000pt}%
\definecolor{currentstroke}{rgb}{0.000000,0.000000,0.000000}%
\pgfsetstrokecolor{currentstroke}%
\pgfsetdash{}{0pt}%
\pgfpathmoveto{\pgfqpoint{1.371346in}{1.139291in}}%
\pgfpathlineto{\pgfqpoint{1.379432in}{1.154500in}}%
\pgfpathlineto{\pgfqpoint{1.387508in}{1.169690in}}%
\pgfpathlineto{\pgfqpoint{1.362823in}{1.176293in}}%
\pgfpathlineto{\pgfqpoint{1.337103in}{1.182248in}}%
\pgfpathlineto{\pgfqpoint{1.330424in}{1.166780in}}%
\pgfpathlineto{\pgfqpoint{1.323735in}{1.151289in}}%
\pgfpathlineto{\pgfqpoint{1.348034in}{1.145599in}}%
\pgfpathlineto{\pgfqpoint{1.371346in}{1.139291in}}%
\pgfpathclose%
\pgfusepath{fill}%
\end{pgfscope}%
\begin{pgfscope}%
\pgfpathrectangle{\pgfqpoint{0.407592in}{0.222904in}}{\pgfqpoint{1.411163in}{1.411163in}}%
\pgfusepath{clip}%
\pgfsetbuttcap%
\pgfsetroundjoin%
\definecolor{currentfill}{rgb}{0.749020,0.772549,0.792157}%
\pgfsetfillcolor{currentfill}%
\pgfsetfillopacity{0.160000}%
\pgfsetlinewidth{0.000000pt}%
\definecolor{currentstroke}{rgb}{0.000000,0.000000,0.000000}%
\pgfsetstrokecolor{currentstroke}%
\pgfsetdash{}{0pt}%
\pgfpathmoveto{\pgfqpoint{1.141908in}{1.010316in}}%
\pgfpathlineto{\pgfqpoint{1.144010in}{1.026581in}}%
\pgfpathlineto{\pgfqpoint{1.146109in}{1.042815in}}%
\pgfpathlineto{\pgfqpoint{1.124222in}{1.043604in}}%
\pgfpathlineto{\pgfqpoint{1.102197in}{1.043829in}}%
\pgfpathlineto{\pgfqpoint{1.101747in}{1.027562in}}%
\pgfpathlineto{\pgfqpoint{1.101297in}{1.011264in}}%
\pgfpathlineto{\pgfqpoint{1.121667in}{1.011053in}}%
\pgfpathlineto{\pgfqpoint{1.141908in}{1.010316in}}%
\pgfpathclose%
\pgfusepath{fill}%
\end{pgfscope}%
\begin{pgfscope}%
\pgfpathrectangle{\pgfqpoint{0.407592in}{0.222904in}}{\pgfqpoint{1.411163in}{1.411163in}}%
\pgfusepath{clip}%
\pgfsetbuttcap%
\pgfsetroundjoin%
\definecolor{currentfill}{rgb}{0.749020,0.772549,0.792157}%
\pgfsetfillcolor{currentfill}%
\pgfsetfillopacity{0.160000}%
\pgfsetlinewidth{0.000000pt}%
\definecolor{currentstroke}{rgb}{0.000000,0.000000,0.000000}%
\pgfsetstrokecolor{currentstroke}%
\pgfsetdash{}{0pt}%
\pgfpathmoveto{\pgfqpoint{1.101297in}{1.011264in}}%
\pgfpathlineto{\pgfqpoint{1.101747in}{1.027562in}}%
\pgfpathlineto{\pgfqpoint{1.102197in}{1.043829in}}%
\pgfpathlineto{\pgfqpoint{1.080117in}{1.043491in}}%
\pgfpathlineto{\pgfqpoint{1.058069in}{1.042589in}}%
\pgfpathlineto{\pgfqpoint{1.059277in}{1.026362in}}%
\pgfpathlineto{\pgfqpoint{1.060487in}{1.010105in}}%
\pgfpathlineto{\pgfqpoint{1.080877in}{1.010947in}}%
\pgfpathlineto{\pgfqpoint{1.101297in}{1.011264in}}%
\pgfpathclose%
\pgfusepath{fill}%
\end{pgfscope}%
\begin{pgfscope}%
\pgfpathrectangle{\pgfqpoint{0.407592in}{0.222904in}}{\pgfqpoint{1.411163in}{1.411163in}}%
\pgfusepath{clip}%
\pgfsetbuttcap%
\pgfsetroundjoin%
\definecolor{currentfill}{rgb}{0.749020,0.772549,0.792157}%
\pgfsetfillcolor{currentfill}%
\pgfsetfillopacity{0.160000}%
\pgfsetlinewidth{0.000000pt}%
\definecolor{currentstroke}{rgb}{0.000000,0.000000,0.000000}%
\pgfsetstrokecolor{currentstroke}%
\pgfsetdash{}{0pt}%
\pgfpathmoveto{\pgfqpoint{0.883412in}{1.150198in}}%
\pgfpathlineto{\pgfqpoint{0.877485in}{1.165664in}}%
\pgfpathlineto{\pgfqpoint{0.871567in}{1.181107in}}%
\pgfpathlineto{\pgfqpoint{0.844648in}{1.175020in}}%
\pgfpathlineto{\pgfqpoint{0.818636in}{1.168289in}}%
\pgfpathlineto{\pgfqpoint{0.826017in}{1.153130in}}%
\pgfpathlineto{\pgfqpoint{0.833408in}{1.137953in}}%
\pgfpathlineto{\pgfqpoint{0.857977in}{1.144383in}}%
\pgfpathlineto{\pgfqpoint{0.883412in}{1.150198in}}%
\pgfpathclose%
\pgfusepath{fill}%
\end{pgfscope}%
\begin{pgfscope}%
\pgfpathrectangle{\pgfqpoint{0.407592in}{0.222904in}}{\pgfqpoint{1.411163in}{1.411163in}}%
\pgfusepath{clip}%
\pgfsetbuttcap%
\pgfsetroundjoin%
\definecolor{currentfill}{rgb}{0.749020,0.772549,0.792157}%
\pgfsetfillcolor{currentfill}%
\pgfsetfillopacity{0.160000}%
\pgfsetlinewidth{0.000000pt}%
\definecolor{currentstroke}{rgb}{0.000000,0.000000,0.000000}%
\pgfsetstrokecolor{currentstroke}%
\pgfsetdash{}{0pt}%
\pgfpathmoveto{\pgfqpoint{1.230598in}{1.034105in}}%
\pgfpathlineto{\pgfqpoint{1.235875in}{1.050028in}}%
\pgfpathlineto{\pgfqpoint{1.241143in}{1.065924in}}%
\pgfpathlineto{\pgfqpoint{1.219096in}{1.069114in}}%
\pgfpathlineto{\pgfqpoint{1.196547in}{1.071728in}}%
\pgfpathlineto{\pgfqpoint{1.192842in}{1.055658in}}%
\pgfpathlineto{\pgfqpoint{1.189131in}{1.039559in}}%
\pgfpathlineto{\pgfqpoint{1.210100in}{1.037103in}}%
\pgfpathlineto{\pgfqpoint{1.230598in}{1.034105in}}%
\pgfpathclose%
\pgfusepath{fill}%
\end{pgfscope}%
\begin{pgfscope}%
\pgfpathrectangle{\pgfqpoint{0.407592in}{0.222904in}}{\pgfqpoint{1.411163in}{1.411163in}}%
\pgfusepath{clip}%
\pgfsetbuttcap%
\pgfsetroundjoin%
\definecolor{currentfill}{rgb}{0.749020,0.772549,0.792157}%
\pgfsetfillcolor{currentfill}%
\pgfsetfillopacity{0.160000}%
\pgfsetlinewidth{0.000000pt}%
\definecolor{currentstroke}{rgb}{0.000000,0.000000,0.000000}%
\pgfsetstrokecolor{currentstroke}%
\pgfsetdash{}{0pt}%
\pgfpathmoveto{\pgfqpoint{0.774538in}{1.258849in}}%
\pgfpathlineto{\pgfqpoint{0.767220in}{1.273878in}}%
\pgfpathlineto{\pgfqpoint{0.759911in}{1.288888in}}%
\pgfpathlineto{\pgfqpoint{0.729314in}{1.280296in}}%
\pgfpathlineto{\pgfqpoint{0.700025in}{1.270993in}}%
\pgfpathlineto{\pgfqpoint{0.708711in}{1.256291in}}%
\pgfpathlineto{\pgfqpoint{0.717406in}{1.241574in}}%
\pgfpathlineto{\pgfqpoint{0.745342in}{1.250554in}}%
\pgfpathlineto{\pgfqpoint{0.774538in}{1.258849in}}%
\pgfpathclose%
\pgfusepath{fill}%
\end{pgfscope}%
\begin{pgfscope}%
\pgfpathrectangle{\pgfqpoint{0.407592in}{0.222904in}}{\pgfqpoint{1.411163in}{1.411163in}}%
\pgfusepath{clip}%
\pgfsetbuttcap%
\pgfsetroundjoin%
\definecolor{currentfill}{rgb}{0.749020,0.772549,0.792157}%
\pgfsetfillcolor{currentfill}%
\pgfsetfillopacity{0.160000}%
\pgfsetlinewidth{0.000000pt}%
\definecolor{currentstroke}{rgb}{0.000000,0.000000,0.000000}%
\pgfsetstrokecolor{currentstroke}%
\pgfsetdash{}{0pt}%
\pgfpathmoveto{\pgfqpoint{1.014403in}{1.039110in}}%
\pgfpathlineto{\pgfqpoint{1.011557in}{1.055195in}}%
\pgfpathlineto{\pgfqpoint{1.008716in}{1.071250in}}%
\pgfpathlineto{\pgfqpoint{0.985652in}{1.068520in}}%
\pgfpathlineto{\pgfqpoint{0.962979in}{1.065216in}}%
\pgfpathlineto{\pgfqpoint{0.967422in}{1.049341in}}%
\pgfpathlineto{\pgfqpoint{0.971872in}{1.033439in}}%
\pgfpathlineto{\pgfqpoint{0.992953in}{1.036545in}}%
\pgfpathlineto{\pgfqpoint{1.014403in}{1.039110in}}%
\pgfpathclose%
\pgfusepath{fill}%
\end{pgfscope}%
\begin{pgfscope}%
\pgfpathrectangle{\pgfqpoint{0.407592in}{0.222904in}}{\pgfqpoint{1.411163in}{1.411163in}}%
\pgfusepath{clip}%
\pgfsetbuttcap%
\pgfsetroundjoin%
\definecolor{currentfill}{rgb}{0.749020,0.772549,0.792157}%
\pgfsetfillcolor{currentfill}%
\pgfsetfillopacity{0.160000}%
\pgfsetlinewidth{0.000000pt}%
\definecolor{currentstroke}{rgb}{0.000000,0.000000,0.000000}%
\pgfsetstrokecolor{currentstroke}%
\pgfsetdash{}{0pt}%
\pgfpathmoveto{\pgfqpoint{1.683104in}{1.424367in}}%
\pgfpathlineto{\pgfqpoint{1.693489in}{1.438622in}}%
\pgfpathlineto{\pgfqpoint{1.661761in}{1.451305in}}%
\pgfpathlineto{\pgfqpoint{1.627969in}{1.463170in}}%
\pgfpathlineto{\pgfqpoint{1.618793in}{1.448621in}}%
\pgfpathlineto{\pgfqpoint{1.651967in}{1.436897in}}%
\pgfusepath{fill}%
\end{pgfscope}%
\begin{pgfscope}%
\pgfpathrectangle{\pgfqpoint{0.407592in}{0.222904in}}{\pgfqpoint{1.411163in}{1.411163in}}%
\pgfusepath{clip}%
\pgfsetbuttcap%
\pgfsetroundjoin%
\definecolor{currentfill}{rgb}{0.749020,0.772549,0.792157}%
\pgfsetfillcolor{currentfill}%
\pgfsetfillopacity{0.160000}%
\pgfsetlinewidth{0.000000pt}%
\definecolor{currentstroke}{rgb}{0.000000,0.000000,0.000000}%
\pgfsetstrokecolor{currentstroke}%
\pgfsetdash{}{0pt}%
\pgfpathmoveto{\pgfqpoint{1.181688in}{1.007274in}}%
\pgfpathlineto{\pgfqpoint{1.185413in}{1.023431in}}%
\pgfpathlineto{\pgfqpoint{1.189131in}{1.039559in}}%
\pgfpathlineto{\pgfqpoint{1.167773in}{1.041465in}}%
\pgfpathlineto{\pgfqpoint{1.146109in}{1.042815in}}%
\pgfpathlineto{\pgfqpoint{1.144010in}{1.026581in}}%
\pgfpathlineto{\pgfqpoint{1.141908in}{1.010316in}}%
\pgfpathlineto{\pgfqpoint{1.161941in}{1.009055in}}%
\pgfpathlineto{\pgfqpoint{1.181688in}{1.007274in}}%
\pgfpathclose%
\pgfusepath{fill}%
\end{pgfscope}%
\begin{pgfscope}%
\pgfpathrectangle{\pgfqpoint{0.407592in}{0.222904in}}{\pgfqpoint{1.411163in}{1.411163in}}%
\pgfusepath{clip}%
\pgfsetbuttcap%
\pgfsetroundjoin%
\definecolor{currentfill}{rgb}{0.749020,0.772549,0.792157}%
\pgfsetfillcolor{currentfill}%
\pgfsetfillopacity{0.160000}%
\pgfsetlinewidth{0.000000pt}%
\definecolor{currentstroke}{rgb}{0.000000,0.000000,0.000000}%
\pgfsetstrokecolor{currentstroke}%
\pgfsetdash{}{0pt}%
\pgfpathmoveto{\pgfqpoint{1.060487in}{1.010105in}}%
\pgfpathlineto{\pgfqpoint{1.059277in}{1.026362in}}%
\pgfpathlineto{\pgfqpoint{1.058069in}{1.042589in}}%
\pgfpathlineto{\pgfqpoint{1.036136in}{1.041127in}}%
\pgfpathlineto{\pgfqpoint{1.014403in}{1.039110in}}%
\pgfpathlineto{\pgfqpoint{1.017254in}{1.022997in}}%
\pgfpathlineto{\pgfqpoint{1.020110in}{1.006855in}}%
\pgfpathlineto{\pgfqpoint{1.040205in}{1.008739in}}%
\pgfpathlineto{\pgfqpoint{1.060487in}{1.010105in}}%
\pgfpathclose%
\pgfusepath{fill}%
\end{pgfscope}%
\begin{pgfscope}%
\pgfpathrectangle{\pgfqpoint{0.407592in}{0.222904in}}{\pgfqpoint{1.411163in}{1.411163in}}%
\pgfusepath{clip}%
\pgfsetbuttcap%
\pgfsetroundjoin%
\definecolor{currentfill}{rgb}{0.749020,0.772549,0.792157}%
\pgfsetfillcolor{currentfill}%
\pgfsetfillopacity{0.160000}%
\pgfsetlinewidth{0.000000pt}%
\definecolor{currentstroke}{rgb}{0.000000,0.000000,0.000000}%
\pgfsetstrokecolor{currentstroke}%
\pgfsetdash{}{0pt}%
\pgfpathmoveto{\pgfqpoint{1.470751in}{1.213885in}}%
\pgfpathlineto{\pgfqpoint{1.480072in}{1.228664in}}%
\pgfpathlineto{\pgfqpoint{1.489383in}{1.243429in}}%
\pgfpathlineto{\pgfqpoint{1.463266in}{1.252274in}}%
\pgfpathlineto{\pgfqpoint{1.435749in}{1.260429in}}%
\pgfpathlineto{\pgfqpoint{1.427734in}{1.245353in}}%
\pgfpathlineto{\pgfqpoint{1.419709in}{1.230258in}}%
\pgfpathlineto{\pgfqpoint{1.445902in}{1.222403in}}%
\pgfpathlineto{\pgfqpoint{1.470751in}{1.213885in}}%
\pgfpathclose%
\pgfusepath{fill}%
\end{pgfscope}%
\begin{pgfscope}%
\pgfpathrectangle{\pgfqpoint{0.407592in}{0.222904in}}{\pgfqpoint{1.411163in}{1.411163in}}%
\pgfusepath{clip}%
\pgfsetbuttcap%
\pgfsetroundjoin%
\definecolor{currentfill}{rgb}{0.749020,0.772549,0.792157}%
\pgfsetfillcolor{currentfill}%
\pgfsetfillopacity{0.160000}%
\pgfsetlinewidth{0.000000pt}%
\definecolor{currentstroke}{rgb}{0.000000,0.000000,0.000000}%
\pgfsetstrokecolor{currentstroke}%
\pgfsetdash{}{0pt}%
\pgfpathmoveto{\pgfqpoint{1.283390in}{1.057857in}}%
\pgfpathlineto{\pgfqpoint{1.290139in}{1.073487in}}%
\pgfpathlineto{\pgfqpoint{1.296878in}{1.089094in}}%
\pgfpathlineto{\pgfqpoint{1.274622in}{1.093659in}}%
\pgfpathlineto{\pgfqpoint{1.251652in}{1.097637in}}%
\pgfpathlineto{\pgfqpoint{1.246402in}{1.081794in}}%
\pgfpathlineto{\pgfqpoint{1.241143in}{1.065924in}}%
\pgfpathlineto{\pgfqpoint{1.262602in}{1.062168in}}%
\pgfpathlineto{\pgfqpoint{1.283390in}{1.057857in}}%
\pgfpathclose%
\pgfusepath{fill}%
\end{pgfscope}%
\begin{pgfscope}%
\pgfpathrectangle{\pgfqpoint{0.407592in}{0.222904in}}{\pgfqpoint{1.411163in}{1.411163in}}%
\pgfusepath{clip}%
\pgfsetbuttcap%
\pgfsetroundjoin%
\definecolor{currentfill}{rgb}{0.749020,0.772549,0.792157}%
\pgfsetfillcolor{currentfill}%
\pgfsetfillopacity{0.160000}%
\pgfsetlinewidth{0.000000pt}%
\definecolor{currentstroke}{rgb}{0.000000,0.000000,0.000000}%
\pgfsetstrokecolor{currentstroke}%
\pgfsetdash{}{0pt}%
\pgfpathmoveto{\pgfqpoint{0.962979in}{1.065216in}}%
\pgfpathlineto{\pgfqpoint{0.958543in}{1.081064in}}%
\pgfpathlineto{\pgfqpoint{0.954115in}{1.096887in}}%
\pgfpathlineto{\pgfqpoint{0.930361in}{1.092790in}}%
\pgfpathlineto{\pgfqpoint{0.907207in}{1.088109in}}%
\pgfpathlineto{\pgfqpoint{0.913177in}{1.072529in}}%
\pgfpathlineto{\pgfqpoint{0.919157in}{1.056927in}}%
\pgfpathlineto{\pgfqpoint{0.940784in}{1.061347in}}%
\pgfpathlineto{\pgfqpoint{0.962979in}{1.065216in}}%
\pgfpathclose%
\pgfusepath{fill}%
\end{pgfscope}%
\begin{pgfscope}%
\pgfpathrectangle{\pgfqpoint{0.407592in}{0.222904in}}{\pgfqpoint{1.411163in}{1.411163in}}%
\pgfusepath{clip}%
\pgfsetbuttcap%
\pgfsetroundjoin%
\definecolor{currentfill}{rgb}{0.749020,0.772549,0.792157}%
\pgfsetfillcolor{currentfill}%
\pgfsetfillopacity{0.160000}%
\pgfsetlinewidth{0.000000pt}%
\definecolor{currentstroke}{rgb}{0.000000,0.000000,0.000000}%
\pgfsetstrokecolor{currentstroke}%
\pgfsetdash{}{0pt}%
\pgfpathmoveto{\pgfqpoint{1.355143in}{1.108815in}}%
\pgfpathlineto{\pgfqpoint{1.363250in}{1.124063in}}%
\pgfpathlineto{\pgfqpoint{1.371346in}{1.139291in}}%
\pgfpathlineto{\pgfqpoint{1.348034in}{1.145599in}}%
\pgfpathlineto{\pgfqpoint{1.323735in}{1.151289in}}%
\pgfpathlineto{\pgfqpoint{1.317036in}{1.135775in}}%
\pgfpathlineto{\pgfqpoint{1.310327in}{1.120238in}}%
\pgfpathlineto{\pgfqpoint{1.333204in}{1.114820in}}%
\pgfpathlineto{\pgfqpoint{1.355143in}{1.108815in}}%
\pgfpathclose%
\pgfusepath{fill}%
\end{pgfscope}%
\begin{pgfscope}%
\pgfpathrectangle{\pgfqpoint{0.407592in}{0.222904in}}{\pgfqpoint{1.411163in}{1.411163in}}%
\pgfusepath{clip}%
\pgfsetbuttcap%
\pgfsetroundjoin%
\definecolor{currentfill}{rgb}{0.749020,0.772549,0.792157}%
\pgfsetfillcolor{currentfill}%
\pgfsetfillopacity{0.160000}%
\pgfsetlinewidth{0.000000pt}%
\definecolor{currentstroke}{rgb}{0.000000,0.000000,0.000000}%
\pgfsetstrokecolor{currentstroke}%
\pgfsetdash{}{0pt}%
\pgfpathmoveto{\pgfqpoint{1.662313in}{1.395827in}}%
\pgfpathlineto{\pgfqpoint{1.672712in}{1.410102in}}%
\pgfpathlineto{\pgfqpoint{1.683104in}{1.424367in}}%
\pgfpathlineto{\pgfqpoint{1.651967in}{1.436897in}}%
\pgfpathlineto{\pgfqpoint{1.618793in}{1.448621in}}%
\pgfpathlineto{\pgfqpoint{1.609608in}{1.434057in}}%
\pgfpathlineto{\pgfqpoint{1.600415in}{1.419480in}}%
\pgfpathlineto{\pgfqpoint{1.632354in}{1.408045in}}%
\pgfpathlineto{\pgfqpoint{1.662313in}{1.395827in}}%
\pgfpathclose%
\pgfusepath{fill}%
\end{pgfscope}%
\begin{pgfscope}%
\pgfpathrectangle{\pgfqpoint{0.407592in}{0.222904in}}{\pgfqpoint{1.411163in}{1.411163in}}%
\pgfusepath{clip}%
\pgfsetbuttcap%
\pgfsetroundjoin%
\definecolor{currentfill}{rgb}{0.749020,0.772549,0.792157}%
\pgfsetfillcolor{currentfill}%
\pgfsetfillopacity{0.160000}%
\pgfsetlinewidth{0.000000pt}%
\definecolor{currentstroke}{rgb}{0.000000,0.000000,0.000000}%
\pgfsetstrokecolor{currentstroke}%
\pgfsetdash{}{0pt}%
\pgfpathmoveto{\pgfqpoint{0.596432in}{1.446335in}}%
\pgfpathlineto{\pgfqpoint{0.587852in}{1.460857in}}%
\pgfpathlineto{\pgfqpoint{0.551453in}{1.448826in}}%
\pgfpathlineto{\pgfqpoint{0.516943in}{1.435982in}}%
\pgfpathlineto{\pgfqpoint{0.526822in}{1.421759in}}%
\pgfpathlineto{\pgfqpoint{0.560694in}{1.434447in}}%
\pgfusepath{fill}%
\end{pgfscope}%
\begin{pgfscope}%
\pgfpathrectangle{\pgfqpoint{0.407592in}{0.222904in}}{\pgfqpoint{1.411163in}{1.411163in}}%
\pgfusepath{clip}%
\pgfsetbuttcap%
\pgfsetroundjoin%
\definecolor{currentfill}{rgb}{0.749020,0.772549,0.792157}%
\pgfsetfillcolor{currentfill}%
\pgfsetfillopacity{0.160000}%
\pgfsetlinewidth{0.000000pt}%
\definecolor{currentstroke}{rgb}{0.000000,0.000000,0.000000}%
\pgfsetstrokecolor{currentstroke}%
\pgfsetdash{}{0pt}%
\pgfpathmoveto{\pgfqpoint{0.789201in}{1.228737in}}%
\pgfpathlineto{\pgfqpoint{0.781865in}{1.243802in}}%
\pgfpathlineto{\pgfqpoint{0.774538in}{1.258849in}}%
\pgfpathlineto{\pgfqpoint{0.745342in}{1.250554in}}%
\pgfpathlineto{\pgfqpoint{0.717406in}{1.241574in}}%
\pgfpathlineto{\pgfqpoint{0.726109in}{1.226844in}}%
\pgfpathlineto{\pgfqpoint{0.734820in}{1.212099in}}%
\pgfpathlineto{\pgfqpoint{0.761404in}{1.220746in}}%
\pgfpathlineto{\pgfqpoint{0.789201in}{1.228737in}}%
\pgfpathclose%
\pgfusepath{fill}%
\end{pgfscope}%
\begin{pgfscope}%
\pgfpathrectangle{\pgfqpoint{0.407592in}{0.222904in}}{\pgfqpoint{1.411163in}{1.411163in}}%
\pgfusepath{clip}%
\pgfsetbuttcap%
\pgfsetroundjoin%
\definecolor{currentfill}{rgb}{0.749020,0.772549,0.792157}%
\pgfsetfillcolor{currentfill}%
\pgfsetfillopacity{0.160000}%
\pgfsetlinewidth{0.000000pt}%
\definecolor{currentstroke}{rgb}{0.000000,0.000000,0.000000}%
\pgfsetstrokecolor{currentstroke}%
\pgfsetdash{}{0pt}%
\pgfpathmoveto{\pgfqpoint{0.895292in}{1.119199in}}%
\pgfpathlineto{\pgfqpoint{0.889347in}{1.134710in}}%
\pgfpathlineto{\pgfqpoint{0.883412in}{1.150198in}}%
\pgfpathlineto{\pgfqpoint{0.857977in}{1.144383in}}%
\pgfpathlineto{\pgfqpoint{0.833408in}{1.137953in}}%
\pgfpathlineto{\pgfqpoint{0.840807in}{1.122756in}}%
\pgfpathlineto{\pgfqpoint{0.848216in}{1.107541in}}%
\pgfpathlineto{\pgfqpoint{0.871342in}{1.113662in}}%
\pgfpathlineto{\pgfqpoint{0.895292in}{1.119199in}}%
\pgfpathclose%
\pgfusepath{fill}%
\end{pgfscope}%
\begin{pgfscope}%
\pgfpathrectangle{\pgfqpoint{0.407592in}{0.222904in}}{\pgfqpoint{1.411163in}{1.411163in}}%
\pgfusepath{clip}%
\pgfsetbuttcap%
\pgfsetroundjoin%
\definecolor{currentfill}{rgb}{0.749020,0.772549,0.792157}%
\pgfsetfillcolor{currentfill}%
\pgfsetfillopacity{0.160000}%
\pgfsetlinewidth{0.000000pt}%
\definecolor{currentstroke}{rgb}{0.000000,0.000000,0.000000}%
\pgfsetstrokecolor{currentstroke}%
\pgfsetdash{}{0pt}%
\pgfpathmoveto{\pgfqpoint{1.137691in}{0.977694in}}%
\pgfpathlineto{\pgfqpoint{1.139801in}{0.994020in}}%
\pgfpathlineto{\pgfqpoint{1.141908in}{1.010316in}}%
\pgfpathlineto{\pgfqpoint{1.121667in}{1.011053in}}%
\pgfpathlineto{\pgfqpoint{1.101297in}{1.011264in}}%
\pgfpathlineto{\pgfqpoint{1.100846in}{0.994934in}}%
\pgfpathlineto{\pgfqpoint{1.100394in}{0.978573in}}%
\pgfpathlineto{\pgfqpoint{1.119102in}{0.978378in}}%
\pgfpathlineto{\pgfqpoint{1.137691in}{0.977694in}}%
\pgfpathclose%
\pgfusepath{fill}%
\end{pgfscope}%
\begin{pgfscope}%
\pgfpathrectangle{\pgfqpoint{0.407592in}{0.222904in}}{\pgfqpoint{1.411163in}{1.411163in}}%
\pgfusepath{clip}%
\pgfsetbuttcap%
\pgfsetroundjoin%
\definecolor{currentfill}{rgb}{0.749020,0.772549,0.792157}%
\pgfsetfillcolor{currentfill}%
\pgfsetfillopacity{0.160000}%
\pgfsetlinewidth{0.000000pt}%
\definecolor{currentstroke}{rgb}{0.000000,0.000000,0.000000}%
\pgfsetstrokecolor{currentstroke}%
\pgfsetdash{}{0pt}%
\pgfpathmoveto{\pgfqpoint{1.641494in}{1.367249in}}%
\pgfpathlineto{\pgfqpoint{1.651907in}{1.381543in}}%
\pgfpathlineto{\pgfqpoint{1.662313in}{1.395827in}}%
\pgfpathlineto{\pgfqpoint{1.632354in}{1.408045in}}%
\pgfpathlineto{\pgfqpoint{1.600415in}{1.419480in}}%
\pgfpathlineto{\pgfqpoint{1.591212in}{1.404888in}}%
\pgfpathlineto{\pgfqpoint{1.582000in}{1.390282in}}%
\pgfpathlineto{\pgfqpoint{1.612709in}{1.379145in}}%
\pgfpathlineto{\pgfqpoint{1.641494in}{1.367249in}}%
\pgfpathclose%
\pgfusepath{fill}%
\end{pgfscope}%
\begin{pgfscope}%
\pgfpathrectangle{\pgfqpoint{0.407592in}{0.222904in}}{\pgfqpoint{1.411163in}{1.411163in}}%
\pgfusepath{clip}%
\pgfsetbuttcap%
\pgfsetroundjoin%
\definecolor{currentfill}{rgb}{0.749020,0.772549,0.792157}%
\pgfsetfillcolor{currentfill}%
\pgfsetfillopacity{0.160000}%
\pgfsetlinewidth{0.000000pt}%
\definecolor{currentstroke}{rgb}{0.000000,0.000000,0.000000}%
\pgfsetstrokecolor{currentstroke}%
\pgfsetdash{}{0pt}%
\pgfpathmoveto{\pgfqpoint{1.452081in}{1.184283in}}%
\pgfpathlineto{\pgfqpoint{1.461420in}{1.199091in}}%
\pgfpathlineto{\pgfqpoint{1.470751in}{1.213885in}}%
\pgfpathlineto{\pgfqpoint{1.445902in}{1.222403in}}%
\pgfpathlineto{\pgfqpoint{1.419709in}{1.230258in}}%
\pgfpathlineto{\pgfqpoint{1.411674in}{1.215145in}}%
\pgfpathlineto{\pgfqpoint{1.403629in}{1.200012in}}%
\pgfpathlineto{\pgfqpoint{1.428500in}{1.192464in}}%
\pgfpathlineto{\pgfqpoint{1.452081in}{1.184283in}}%
\pgfpathclose%
\pgfusepath{fill}%
\end{pgfscope}%
\begin{pgfscope}%
\pgfpathrectangle{\pgfqpoint{0.407592in}{0.222904in}}{\pgfqpoint{1.411163in}{1.411163in}}%
\pgfusepath{clip}%
\pgfsetbuttcap%
\pgfsetroundjoin%
\definecolor{currentfill}{rgb}{0.749020,0.772549,0.792157}%
\pgfsetfillcolor{currentfill}%
\pgfsetfillopacity{0.160000}%
\pgfsetlinewidth{0.000000pt}%
\definecolor{currentstroke}{rgb}{0.000000,0.000000,0.000000}%
\pgfsetstrokecolor{currentstroke}%
\pgfsetdash{}{0pt}%
\pgfpathmoveto{\pgfqpoint{1.100394in}{0.978573in}}%
\pgfpathlineto{\pgfqpoint{1.100846in}{0.994934in}}%
\pgfpathlineto{\pgfqpoint{1.101297in}{1.011264in}}%
\pgfpathlineto{\pgfqpoint{1.080877in}{1.010947in}}%
\pgfpathlineto{\pgfqpoint{1.060487in}{1.010105in}}%
\pgfpathlineto{\pgfqpoint{1.061699in}{0.993817in}}%
\pgfpathlineto{\pgfqpoint{1.062914in}{0.977498in}}%
\pgfpathlineto{\pgfqpoint{1.081640in}{0.978280in}}%
\pgfpathlineto{\pgfqpoint{1.100394in}{0.978573in}}%
\pgfpathclose%
\pgfusepath{fill}%
\end{pgfscope}%
\begin{pgfscope}%
\pgfpathrectangle{\pgfqpoint{0.407592in}{0.222904in}}{\pgfqpoint{1.411163in}{1.411163in}}%
\pgfusepath{clip}%
\pgfsetbuttcap%
\pgfsetroundjoin%
\definecolor{currentfill}{rgb}{0.749020,0.772549,0.792157}%
\pgfsetfillcolor{currentfill}%
\pgfsetfillopacity{0.160000}%
\pgfsetlinewidth{0.000000pt}%
\definecolor{currentstroke}{rgb}{0.000000,0.000000,0.000000}%
\pgfsetstrokecolor{currentstroke}%
\pgfsetdash{}{0pt}%
\pgfpathmoveto{\pgfqpoint{0.613616in}{1.417250in}}%
\pgfpathlineto{\pgfqpoint{0.605020in}{1.431799in}}%
\pgfpathlineto{\pgfqpoint{0.596432in}{1.446335in}}%
\pgfpathlineto{\pgfqpoint{0.560694in}{1.434447in}}%
\pgfpathlineto{\pgfqpoint{0.526822in}{1.421759in}}%
\pgfpathlineto{\pgfqpoint{0.536706in}{1.407526in}}%
\pgfpathlineto{\pgfqpoint{0.546597in}{1.393285in}}%
\pgfpathlineto{\pgfqpoint{0.579199in}{1.405656in}}%
\pgfpathlineto{\pgfqpoint{0.613616in}{1.417250in}}%
\pgfpathclose%
\pgfusepath{fill}%
\end{pgfscope}%
\begin{pgfscope}%
\pgfpathrectangle{\pgfqpoint{0.407592in}{0.222904in}}{\pgfqpoint{1.411163in}{1.411163in}}%
\pgfusepath{clip}%
\pgfsetbuttcap%
\pgfsetroundjoin%
\definecolor{currentfill}{rgb}{0.749020,0.772549,0.792157}%
\pgfsetfillcolor{currentfill}%
\pgfsetfillopacity{0.160000}%
\pgfsetlinewidth{0.000000pt}%
\definecolor{currentstroke}{rgb}{0.000000,0.000000,0.000000}%
\pgfsetstrokecolor{currentstroke}%
\pgfsetdash{}{0pt}%
\pgfpathmoveto{\pgfqpoint{1.220018in}{1.002179in}}%
\pgfpathlineto{\pgfqpoint{1.225313in}{1.018156in}}%
\pgfpathlineto{\pgfqpoint{1.230598in}{1.034105in}}%
\pgfpathlineto{\pgfqpoint{1.210100in}{1.037103in}}%
\pgfpathlineto{\pgfqpoint{1.189131in}{1.039559in}}%
\pgfpathlineto{\pgfqpoint{1.185413in}{1.023431in}}%
\pgfpathlineto{\pgfqpoint{1.181688in}{1.007274in}}%
\pgfpathlineto{\pgfqpoint{1.201073in}{1.004980in}}%
\pgfpathlineto{\pgfqpoint{1.220018in}{1.002179in}}%
\pgfpathclose%
\pgfusepath{fill}%
\end{pgfscope}%
\begin{pgfscope}%
\pgfpathrectangle{\pgfqpoint{0.407592in}{0.222904in}}{\pgfqpoint{1.411163in}{1.411163in}}%
\pgfusepath{clip}%
\pgfsetbuttcap%
\pgfsetroundjoin%
\definecolor{currentfill}{rgb}{0.749020,0.772549,0.792157}%
\pgfsetfillcolor{currentfill}%
\pgfsetfillopacity{0.160000}%
\pgfsetlinewidth{0.000000pt}%
\definecolor{currentstroke}{rgb}{0.000000,0.000000,0.000000}%
\pgfsetstrokecolor{currentstroke}%
\pgfsetdash{}{0pt}%
\pgfpathmoveto{\pgfqpoint{1.020110in}{1.006855in}}%
\pgfpathlineto{\pgfqpoint{1.017254in}{1.022997in}}%
\pgfpathlineto{\pgfqpoint{1.014403in}{1.039110in}}%
\pgfpathlineto{\pgfqpoint{0.992953in}{1.036545in}}%
\pgfpathlineto{\pgfqpoint{0.971872in}{1.033439in}}%
\pgfpathlineto{\pgfqpoint{0.976330in}{1.017512in}}%
\pgfpathlineto{\pgfqpoint{0.980795in}{1.001558in}}%
\pgfpathlineto{\pgfqpoint{1.000280in}{1.004458in}}%
\pgfpathlineto{\pgfqpoint{1.020110in}{1.006855in}}%
\pgfpathclose%
\pgfusepath{fill}%
\end{pgfscope}%
\begin{pgfscope}%
\pgfpathrectangle{\pgfqpoint{0.407592in}{0.222904in}}{\pgfqpoint{1.411163in}{1.411163in}}%
\pgfusepath{clip}%
\pgfsetbuttcap%
\pgfsetroundjoin%
\definecolor{currentfill}{rgb}{0.749020,0.772549,0.792157}%
\pgfsetfillcolor{currentfill}%
\pgfsetfillopacity{0.160000}%
\pgfsetlinewidth{0.000000pt}%
\definecolor{currentstroke}{rgb}{0.000000,0.000000,0.000000}%
\pgfsetstrokecolor{currentstroke}%
\pgfsetdash{}{0pt}%
\pgfpathmoveto{\pgfqpoint{1.269861in}{1.026526in}}%
\pgfpathlineto{\pgfqpoint{1.276630in}{1.042203in}}%
\pgfpathlineto{\pgfqpoint{1.283390in}{1.057857in}}%
\pgfpathlineto{\pgfqpoint{1.262602in}{1.062168in}}%
\pgfpathlineto{\pgfqpoint{1.241143in}{1.065924in}}%
\pgfpathlineto{\pgfqpoint{1.235875in}{1.050028in}}%
\pgfpathlineto{\pgfqpoint{1.230598in}{1.034105in}}%
\pgfpathlineto{\pgfqpoint{1.250545in}{1.030576in}}%
\pgfpathlineto{\pgfqpoint{1.269861in}{1.026526in}}%
\pgfpathclose%
\pgfusepath{fill}%
\end{pgfscope}%
\begin{pgfscope}%
\pgfpathrectangle{\pgfqpoint{0.407592in}{0.222904in}}{\pgfqpoint{1.411163in}{1.411163in}}%
\pgfusepath{clip}%
\pgfsetbuttcap%
\pgfsetroundjoin%
\definecolor{currentfill}{rgb}{0.749020,0.772549,0.792157}%
\pgfsetfillcolor{currentfill}%
\pgfsetfillopacity{0.160000}%
\pgfsetlinewidth{0.000000pt}%
\definecolor{currentstroke}{rgb}{0.000000,0.000000,0.000000}%
\pgfsetstrokecolor{currentstroke}%
\pgfsetdash{}{0pt}%
\pgfpathmoveto{\pgfqpoint{1.338900in}{1.078263in}}%
\pgfpathlineto{\pgfqpoint{1.347027in}{1.093549in}}%
\pgfpathlineto{\pgfqpoint{1.355143in}{1.108815in}}%
\pgfpathlineto{\pgfqpoint{1.333204in}{1.114820in}}%
\pgfpathlineto{\pgfqpoint{1.310327in}{1.120238in}}%
\pgfpathlineto{\pgfqpoint{1.303607in}{1.104677in}}%
\pgfpathlineto{\pgfqpoint{1.296878in}{1.089094in}}%
\pgfpathlineto{\pgfqpoint{1.318333in}{1.083956in}}%
\pgfpathlineto{\pgfqpoint{1.338900in}{1.078263in}}%
\pgfpathclose%
\pgfusepath{fill}%
\end{pgfscope}%
\begin{pgfscope}%
\pgfpathrectangle{\pgfqpoint{0.407592in}{0.222904in}}{\pgfqpoint{1.411163in}{1.411163in}}%
\pgfusepath{clip}%
\pgfsetbuttcap%
\pgfsetroundjoin%
\definecolor{currentfill}{rgb}{0.749020,0.772549,0.792157}%
\pgfsetfillcolor{currentfill}%
\pgfsetfillopacity{0.160000}%
\pgfsetlinewidth{0.000000pt}%
\definecolor{currentstroke}{rgb}{0.000000,0.000000,0.000000}%
\pgfsetstrokecolor{currentstroke}%
\pgfsetdash{}{0pt}%
\pgfpathmoveto{\pgfqpoint{0.803900in}{1.198550in}}%
\pgfpathlineto{\pgfqpoint{0.796546in}{1.213653in}}%
\pgfpathlineto{\pgfqpoint{0.789201in}{1.228737in}}%
\pgfpathlineto{\pgfqpoint{0.761404in}{1.220746in}}%
\pgfpathlineto{\pgfqpoint{0.734820in}{1.212099in}}%
\pgfpathlineto{\pgfqpoint{0.743540in}{1.197340in}}%
\pgfpathlineto{\pgfqpoint{0.752268in}{1.182567in}}%
\pgfpathlineto{\pgfqpoint{0.777502in}{1.190873in}}%
\pgfpathlineto{\pgfqpoint{0.803900in}{1.198550in}}%
\pgfpathclose%
\pgfusepath{fill}%
\end{pgfscope}%
\begin{pgfscope}%
\pgfpathrectangle{\pgfqpoint{0.407592in}{0.222904in}}{\pgfqpoint{1.411163in}{1.411163in}}%
\pgfusepath{clip}%
\pgfsetbuttcap%
\pgfsetroundjoin%
\definecolor{currentfill}{rgb}{0.749020,0.772549,0.792157}%
\pgfsetfillcolor{currentfill}%
\pgfsetfillopacity{0.160000}%
\pgfsetlinewidth{0.000000pt}%
\definecolor{currentstroke}{rgb}{0.000000,0.000000,0.000000}%
\pgfsetstrokecolor{currentstroke}%
\pgfsetdash{}{0pt}%
\pgfpathmoveto{\pgfqpoint{1.174219in}{0.974872in}}%
\pgfpathlineto{\pgfqpoint{1.177957in}{0.991087in}}%
\pgfpathlineto{\pgfqpoint{1.181688in}{1.007274in}}%
\pgfpathlineto{\pgfqpoint{1.161941in}{1.009055in}}%
\pgfpathlineto{\pgfqpoint{1.141908in}{1.010316in}}%
\pgfpathlineto{\pgfqpoint{1.139801in}{0.994020in}}%
\pgfpathlineto{\pgfqpoint{1.137691in}{0.977694in}}%
\pgfpathlineto{\pgfqpoint{1.156087in}{0.976524in}}%
\pgfpathlineto{\pgfqpoint{1.174219in}{0.974872in}}%
\pgfpathclose%
\pgfusepath{fill}%
\end{pgfscope}%
\begin{pgfscope}%
\pgfpathrectangle{\pgfqpoint{0.407592in}{0.222904in}}{\pgfqpoint{1.411163in}{1.411163in}}%
\pgfusepath{clip}%
\pgfsetbuttcap%
\pgfsetroundjoin%
\definecolor{currentfill}{rgb}{0.749020,0.772549,0.792157}%
\pgfsetfillcolor{currentfill}%
\pgfsetfillopacity{0.160000}%
\pgfsetlinewidth{0.000000pt}%
\definecolor{currentstroke}{rgb}{0.000000,0.000000,0.000000}%
\pgfsetstrokecolor{currentstroke}%
\pgfsetdash{}{0pt}%
\pgfpathmoveto{\pgfqpoint{1.062914in}{0.977498in}}%
\pgfpathlineto{\pgfqpoint{1.061699in}{0.993817in}}%
\pgfpathlineto{\pgfqpoint{1.060487in}{1.010105in}}%
\pgfpathlineto{\pgfqpoint{1.040205in}{1.008739in}}%
\pgfpathlineto{\pgfqpoint{1.020110in}{1.006855in}}%
\pgfpathlineto{\pgfqpoint{1.022971in}{0.990683in}}%
\pgfpathlineto{\pgfqpoint{1.025837in}{0.974483in}}%
\pgfpathlineto{\pgfqpoint{1.044288in}{0.976231in}}%
\pgfpathlineto{\pgfqpoint{1.062914in}{0.977498in}}%
\pgfpathclose%
\pgfusepath{fill}%
\end{pgfscope}%
\begin{pgfscope}%
\pgfpathrectangle{\pgfqpoint{0.407592in}{0.222904in}}{\pgfqpoint{1.411163in}{1.411163in}}%
\pgfusepath{clip}%
\pgfsetbuttcap%
\pgfsetroundjoin%
\definecolor{currentfill}{rgb}{0.749020,0.772549,0.792157}%
\pgfsetfillcolor{currentfill}%
\pgfsetfillopacity{0.160000}%
\pgfsetlinewidth{0.000000pt}%
\definecolor{currentstroke}{rgb}{0.000000,0.000000,0.000000}%
\pgfsetstrokecolor{currentstroke}%
\pgfsetdash{}{0pt}%
\pgfpathmoveto{\pgfqpoint{0.971872in}{1.033439in}}%
\pgfpathlineto{\pgfqpoint{0.967422in}{1.049341in}}%
\pgfpathlineto{\pgfqpoint{0.962979in}{1.065216in}}%
\pgfpathlineto{\pgfqpoint{0.940784in}{1.061347in}}%
\pgfpathlineto{\pgfqpoint{0.919157in}{1.056927in}}%
\pgfpathlineto{\pgfqpoint{0.925145in}{1.041301in}}%
\pgfpathlineto{\pgfqpoint{0.931142in}{1.025653in}}%
\pgfpathlineto{\pgfqpoint{0.951241in}{1.029805in}}%
\pgfpathlineto{\pgfqpoint{0.971872in}{1.033439in}}%
\pgfpathclose%
\pgfusepath{fill}%
\end{pgfscope}%
\begin{pgfscope}%
\pgfpathrectangle{\pgfqpoint{0.407592in}{0.222904in}}{\pgfqpoint{1.411163in}{1.411163in}}%
\pgfusepath{clip}%
\pgfsetbuttcap%
\pgfsetroundjoin%
\definecolor{currentfill}{rgb}{0.749020,0.772549,0.792157}%
\pgfsetfillcolor{currentfill}%
\pgfsetfillopacity{0.160000}%
\pgfsetlinewidth{0.000000pt}%
\definecolor{currentstroke}{rgb}{0.000000,0.000000,0.000000}%
\pgfsetstrokecolor{currentstroke}%
\pgfsetdash{}{0pt}%
\pgfpathmoveto{\pgfqpoint{1.620647in}{1.338633in}}%
\pgfpathlineto{\pgfqpoint{1.631074in}{1.352946in}}%
\pgfpathlineto{\pgfqpoint{1.641494in}{1.367249in}}%
\pgfpathlineto{\pgfqpoint{1.612709in}{1.379145in}}%
\pgfpathlineto{\pgfqpoint{1.582000in}{1.390282in}}%
\pgfpathlineto{\pgfqpoint{1.572779in}{1.375661in}}%
\pgfpathlineto{\pgfqpoint{1.563549in}{1.361026in}}%
\pgfpathlineto{\pgfqpoint{1.593031in}{1.350197in}}%
\pgfpathlineto{\pgfqpoint{1.620647in}{1.338633in}}%
\pgfpathclose%
\pgfusepath{fill}%
\end{pgfscope}%
\begin{pgfscope}%
\pgfpathrectangle{\pgfqpoint{0.407592in}{0.222904in}}{\pgfqpoint{1.411163in}{1.411163in}}%
\pgfusepath{clip}%
\pgfsetbuttcap%
\pgfsetroundjoin%
\definecolor{currentfill}{rgb}{0.749020,0.772549,0.792157}%
\pgfsetfillcolor{currentfill}%
\pgfsetfillopacity{0.160000}%
\pgfsetlinewidth{0.000000pt}%
\definecolor{currentstroke}{rgb}{0.000000,0.000000,0.000000}%
\pgfsetstrokecolor{currentstroke}%
\pgfsetdash{}{0pt}%
\pgfpathmoveto{\pgfqpoint{0.630832in}{1.388110in}}%
\pgfpathlineto{\pgfqpoint{0.622220in}{1.402687in}}%
\pgfpathlineto{\pgfqpoint{0.613616in}{1.417250in}}%
\pgfpathlineto{\pgfqpoint{0.579199in}{1.405656in}}%
\pgfpathlineto{\pgfqpoint{0.546597in}{1.393285in}}%
\pgfpathlineto{\pgfqpoint{0.556495in}{1.379034in}}%
\pgfpathlineto{\pgfqpoint{0.566398in}{1.364774in}}%
\pgfpathlineto{\pgfqpoint{0.597732in}{1.376818in}}%
\pgfpathlineto{\pgfqpoint{0.630832in}{1.388110in}}%
\pgfpathclose%
\pgfusepath{fill}%
\end{pgfscope}%
\begin{pgfscope}%
\pgfpathrectangle{\pgfqpoint{0.407592in}{0.222904in}}{\pgfqpoint{1.411163in}{1.411163in}}%
\pgfusepath{clip}%
\pgfsetbuttcap%
\pgfsetroundjoin%
\definecolor{currentfill}{rgb}{0.749020,0.772549,0.792157}%
\pgfsetfillcolor{currentfill}%
\pgfsetfillopacity{0.160000}%
\pgfsetlinewidth{0.000000pt}%
\definecolor{currentstroke}{rgb}{0.000000,0.000000,0.000000}%
\pgfsetstrokecolor{currentstroke}%
\pgfsetdash{}{0pt}%
\pgfpathmoveto{\pgfqpoint{0.907207in}{1.088109in}}%
\pgfpathlineto{\pgfqpoint{0.901245in}{1.103665in}}%
\pgfpathlineto{\pgfqpoint{0.895292in}{1.119199in}}%
\pgfpathlineto{\pgfqpoint{0.871342in}{1.113662in}}%
\pgfpathlineto{\pgfqpoint{0.848216in}{1.107541in}}%
\pgfpathlineto{\pgfqpoint{0.855634in}{1.092308in}}%
\pgfpathlineto{\pgfqpoint{0.863061in}{1.077055in}}%
\pgfpathlineto{\pgfqpoint{0.884743in}{1.082858in}}%
\pgfpathlineto{\pgfqpoint{0.907207in}{1.088109in}}%
\pgfpathclose%
\pgfusepath{fill}%
\end{pgfscope}%
\begin{pgfscope}%
\pgfpathrectangle{\pgfqpoint{0.407592in}{0.222904in}}{\pgfqpoint{1.411163in}{1.411163in}}%
\pgfusepath{clip}%
\pgfsetbuttcap%
\pgfsetroundjoin%
\definecolor{currentfill}{rgb}{0.749020,0.772549,0.792157}%
\pgfsetfillcolor{currentfill}%
\pgfsetfillopacity{0.160000}%
\pgfsetlinewidth{0.000000pt}%
\definecolor{currentstroke}{rgb}{0.000000,0.000000,0.000000}%
\pgfsetstrokecolor{currentstroke}%
\pgfsetdash{}{0pt}%
\pgfpathmoveto{\pgfqpoint{1.433375in}{1.154622in}}%
\pgfpathlineto{\pgfqpoint{1.442732in}{1.169460in}}%
\pgfpathlineto{\pgfqpoint{1.452081in}{1.184283in}}%
\pgfpathlineto{\pgfqpoint{1.428500in}{1.192464in}}%
\pgfpathlineto{\pgfqpoint{1.403629in}{1.200012in}}%
\pgfpathlineto{\pgfqpoint{1.395573in}{1.184860in}}%
\pgfpathlineto{\pgfqpoint{1.387508in}{1.169690in}}%
\pgfpathlineto{\pgfqpoint{1.411057in}{1.162458in}}%
\pgfpathlineto{\pgfqpoint{1.433375in}{1.154622in}}%
\pgfpathclose%
\pgfusepath{fill}%
\end{pgfscope}%
\begin{pgfscope}%
\pgfpathrectangle{\pgfqpoint{0.407592in}{0.222904in}}{\pgfqpoint{1.411163in}{1.411163in}}%
\pgfusepath{clip}%
\pgfsetbuttcap%
\pgfsetroundjoin%
\definecolor{currentfill}{rgb}{0.749020,0.772549,0.792157}%
\pgfsetfillcolor{currentfill}%
\pgfsetfillopacity{0.160000}%
\pgfsetlinewidth{0.000000pt}%
\definecolor{currentstroke}{rgb}{0.000000,0.000000,0.000000}%
\pgfsetstrokecolor{currentstroke}%
\pgfsetdash{}{0pt}%
\pgfpathmoveto{\pgfqpoint{0.648081in}{1.358914in}}%
\pgfpathlineto{\pgfqpoint{0.639452in}{1.373519in}}%
\pgfpathlineto{\pgfqpoint{0.630832in}{1.388110in}}%
\pgfpathlineto{\pgfqpoint{0.597732in}{1.376818in}}%
\pgfpathlineto{\pgfqpoint{0.566398in}{1.364774in}}%
\pgfpathlineto{\pgfqpoint{0.576308in}{1.350505in}}%
\pgfpathlineto{\pgfqpoint{0.586224in}{1.336227in}}%
\pgfpathlineto{\pgfqpoint{0.616296in}{1.347935in}}%
\pgfpathlineto{\pgfqpoint{0.648081in}{1.358914in}}%
\pgfpathclose%
\pgfusepath{fill}%
\end{pgfscope}%
\begin{pgfscope}%
\pgfpathrectangle{\pgfqpoint{0.407592in}{0.222904in}}{\pgfqpoint{1.411163in}{1.411163in}}%
\pgfusepath{clip}%
\pgfsetbuttcap%
\pgfsetroundjoin%
\definecolor{currentfill}{rgb}{0.749020,0.772549,0.792157}%
\pgfsetfillcolor{currentfill}%
\pgfsetfillopacity{0.160000}%
\pgfsetlinewidth{0.000000pt}%
\definecolor{currentstroke}{rgb}{0.000000,0.000000,0.000000}%
\pgfsetstrokecolor{currentstroke}%
\pgfsetdash{}{0pt}%
\pgfpathmoveto{\pgfqpoint{0.818636in}{1.168289in}}%
\pgfpathlineto{\pgfqpoint{0.811263in}{1.183428in}}%
\pgfpathlineto{\pgfqpoint{0.803900in}{1.198550in}}%
\pgfpathlineto{\pgfqpoint{0.777502in}{1.190873in}}%
\pgfpathlineto{\pgfqpoint{0.752268in}{1.182567in}}%
\pgfpathlineto{\pgfqpoint{0.761004in}{1.167780in}}%
\pgfpathlineto{\pgfqpoint{0.769749in}{1.152979in}}%
\pgfpathlineto{\pgfqpoint{0.793635in}{1.160934in}}%
\pgfpathlineto{\pgfqpoint{0.818636in}{1.168289in}}%
\pgfpathclose%
\pgfusepath{fill}%
\end{pgfscope}%
\begin{pgfscope}%
\pgfpathrectangle{\pgfqpoint{0.407592in}{0.222904in}}{\pgfqpoint{1.411163in}{1.411163in}}%
\pgfusepath{clip}%
\pgfsetbuttcap%
\pgfsetroundjoin%
\definecolor{currentfill}{rgb}{0.749020,0.772549,0.792157}%
\pgfsetfillcolor{currentfill}%
\pgfsetfillopacity{0.160000}%
\pgfsetlinewidth{0.000000pt}%
\definecolor{currentstroke}{rgb}{0.000000,0.000000,0.000000}%
\pgfsetstrokecolor{currentstroke}%
\pgfsetdash{}{0pt}%
\pgfpathmoveto{\pgfqpoint{1.599772in}{1.309978in}}%
\pgfpathlineto{\pgfqpoint{1.610213in}{1.324310in}}%
\pgfpathlineto{\pgfqpoint{1.620647in}{1.338633in}}%
\pgfpathlineto{\pgfqpoint{1.593031in}{1.350197in}}%
\pgfpathlineto{\pgfqpoint{1.563549in}{1.361026in}}%
\pgfpathlineto{\pgfqpoint{1.554310in}{1.346377in}}%
\pgfpathlineto{\pgfqpoint{1.545062in}{1.331713in}}%
\pgfpathlineto{\pgfqpoint{1.573320in}{1.321201in}}%
\pgfpathlineto{\pgfqpoint{1.599772in}{1.309978in}}%
\pgfpathclose%
\pgfusepath{fill}%
\end{pgfscope}%
\begin{pgfscope}%
\pgfpathrectangle{\pgfqpoint{0.407592in}{0.222904in}}{\pgfqpoint{1.411163in}{1.411163in}}%
\pgfusepath{clip}%
\pgfsetbuttcap%
\pgfsetroundjoin%
\definecolor{currentfill}{rgb}{0.749020,0.772549,0.792157}%
\pgfsetfillcolor{currentfill}%
\pgfsetfillopacity{0.160000}%
\pgfsetlinewidth{0.000000pt}%
\definecolor{currentstroke}{rgb}{0.000000,0.000000,0.000000}%
\pgfsetstrokecolor{currentstroke}%
\pgfsetdash{}{0pt}%
\pgfpathmoveto{\pgfqpoint{1.209403in}{0.970146in}}%
\pgfpathlineto{\pgfqpoint{1.214715in}{0.986176in}}%
\pgfpathlineto{\pgfqpoint{1.220018in}{1.002179in}}%
\pgfpathlineto{\pgfqpoint{1.201073in}{1.004980in}}%
\pgfpathlineto{\pgfqpoint{1.181688in}{1.007274in}}%
\pgfpathlineto{\pgfqpoint{1.177957in}{0.991087in}}%
\pgfpathlineto{\pgfqpoint{1.174219in}{0.974872in}}%
\pgfpathlineto{\pgfqpoint{1.192014in}{0.972743in}}%
\pgfpathlineto{\pgfqpoint{1.209403in}{0.970146in}}%
\pgfpathclose%
\pgfusepath{fill}%
\end{pgfscope}%
\begin{pgfscope}%
\pgfpathrectangle{\pgfqpoint{0.407592in}{0.222904in}}{\pgfqpoint{1.411163in}{1.411163in}}%
\pgfusepath{clip}%
\pgfsetbuttcap%
\pgfsetroundjoin%
\definecolor{currentfill}{rgb}{0.749020,0.772549,0.792157}%
\pgfsetfillcolor{currentfill}%
\pgfsetfillopacity{0.160000}%
\pgfsetlinewidth{0.000000pt}%
\definecolor{currentstroke}{rgb}{0.000000,0.000000,0.000000}%
\pgfsetstrokecolor{currentstroke}%
\pgfsetdash{}{0pt}%
\pgfpathmoveto{\pgfqpoint{1.133458in}{0.944948in}}%
\pgfpathlineto{\pgfqpoint{1.135576in}{0.961336in}}%
\pgfpathlineto{\pgfqpoint{1.137691in}{0.977694in}}%
\pgfpathlineto{\pgfqpoint{1.119102in}{0.978378in}}%
\pgfpathlineto{\pgfqpoint{1.100394in}{0.978573in}}%
\pgfpathlineto{\pgfqpoint{1.099942in}{0.962181in}}%
\pgfpathlineto{\pgfqpoint{1.099488in}{0.945757in}}%
\pgfpathlineto{\pgfqpoint{1.116528in}{0.945577in}}%
\pgfpathlineto{\pgfqpoint{1.133458in}{0.944948in}}%
\pgfpathclose%
\pgfusepath{fill}%
\end{pgfscope}%
\begin{pgfscope}%
\pgfpathrectangle{\pgfqpoint{0.407592in}{0.222904in}}{\pgfqpoint{1.411163in}{1.411163in}}%
\pgfusepath{clip}%
\pgfsetbuttcap%
\pgfsetroundjoin%
\definecolor{currentfill}{rgb}{0.749020,0.772549,0.792157}%
\pgfsetfillcolor{currentfill}%
\pgfsetfillopacity{0.160000}%
\pgfsetlinewidth{0.000000pt}%
\definecolor{currentstroke}{rgb}{0.000000,0.000000,0.000000}%
\pgfsetstrokecolor{currentstroke}%
\pgfsetdash{}{0pt}%
\pgfpathmoveto{\pgfqpoint{1.322616in}{1.047633in}}%
\pgfpathlineto{\pgfqpoint{1.330763in}{1.062957in}}%
\pgfpathlineto{\pgfqpoint{1.338900in}{1.078263in}}%
\pgfpathlineto{\pgfqpoint{1.318333in}{1.083956in}}%
\pgfpathlineto{\pgfqpoint{1.296878in}{1.089094in}}%
\pgfpathlineto{\pgfqpoint{1.290139in}{1.073487in}}%
\pgfpathlineto{\pgfqpoint{1.283390in}{1.057857in}}%
\pgfpathlineto{\pgfqpoint{1.303421in}{1.053007in}}%
\pgfpathlineto{\pgfqpoint{1.322616in}{1.047633in}}%
\pgfpathclose%
\pgfusepath{fill}%
\end{pgfscope}%
\begin{pgfscope}%
\pgfpathrectangle{\pgfqpoint{0.407592in}{0.222904in}}{\pgfqpoint{1.411163in}{1.411163in}}%
\pgfusepath{clip}%
\pgfsetbuttcap%
\pgfsetroundjoin%
\definecolor{currentfill}{rgb}{0.749020,0.772549,0.792157}%
\pgfsetfillcolor{currentfill}%
\pgfsetfillopacity{0.160000}%
\pgfsetlinewidth{0.000000pt}%
\definecolor{currentstroke}{rgb}{0.000000,0.000000,0.000000}%
\pgfsetstrokecolor{currentstroke}%
\pgfsetdash{}{0pt}%
\pgfpathmoveto{\pgfqpoint{1.099488in}{0.945757in}}%
\pgfpathlineto{\pgfqpoint{1.099942in}{0.962181in}}%
\pgfpathlineto{\pgfqpoint{1.100394in}{0.978573in}}%
\pgfpathlineto{\pgfqpoint{1.081640in}{0.978280in}}%
\pgfpathlineto{\pgfqpoint{1.062914in}{0.977498in}}%
\pgfpathlineto{\pgfqpoint{1.064131in}{0.961148in}}%
\pgfpathlineto{\pgfqpoint{1.065350in}{0.944768in}}%
\pgfpathlineto{\pgfqpoint{1.082406in}{0.945487in}}%
\pgfpathlineto{\pgfqpoint{1.099488in}{0.945757in}}%
\pgfpathclose%
\pgfusepath{fill}%
\end{pgfscope}%
\begin{pgfscope}%
\pgfpathrectangle{\pgfqpoint{0.407592in}{0.222904in}}{\pgfqpoint{1.411163in}{1.411163in}}%
\pgfusepath{clip}%
\pgfsetbuttcap%
\pgfsetroundjoin%
\definecolor{currentfill}{rgb}{0.749020,0.772549,0.792157}%
\pgfsetfillcolor{currentfill}%
\pgfsetfillopacity{0.160000}%
\pgfsetlinewidth{0.000000pt}%
\definecolor{currentstroke}{rgb}{0.000000,0.000000,0.000000}%
\pgfsetstrokecolor{currentstroke}%
\pgfsetdash{}{0pt}%
\pgfpathmoveto{\pgfqpoint{1.025837in}{0.974483in}}%
\pgfpathlineto{\pgfqpoint{1.022971in}{0.990683in}}%
\pgfpathlineto{\pgfqpoint{1.020110in}{1.006855in}}%
\pgfpathlineto{\pgfqpoint{1.000280in}{1.004458in}}%
\pgfpathlineto{\pgfqpoint{0.980795in}{1.001558in}}%
\pgfpathlineto{\pgfqpoint{0.985267in}{0.985577in}}%
\pgfpathlineto{\pgfqpoint{0.989747in}{0.969569in}}%
\pgfpathlineto{\pgfqpoint{1.007632in}{0.972260in}}%
\pgfpathlineto{\pgfqpoint{1.025837in}{0.974483in}}%
\pgfpathclose%
\pgfusepath{fill}%
\end{pgfscope}%
\begin{pgfscope}%
\pgfpathrectangle{\pgfqpoint{0.407592in}{0.222904in}}{\pgfqpoint{1.411163in}{1.411163in}}%
\pgfusepath{clip}%
\pgfsetbuttcap%
\pgfsetroundjoin%
\definecolor{currentfill}{rgb}{0.749020,0.772549,0.792157}%
\pgfsetfillcolor{currentfill}%
\pgfsetfillopacity{0.160000}%
\pgfsetlinewidth{0.000000pt}%
\definecolor{currentstroke}{rgb}{0.000000,0.000000,0.000000}%
\pgfsetstrokecolor{currentstroke}%
\pgfsetdash{}{0pt}%
\pgfpathmoveto{\pgfqpoint{1.256291in}{0.995102in}}%
\pgfpathlineto{\pgfqpoint{1.263081in}{1.010826in}}%
\pgfpathlineto{\pgfqpoint{1.269861in}{1.026526in}}%
\pgfpathlineto{\pgfqpoint{1.250545in}{1.030576in}}%
\pgfpathlineto{\pgfqpoint{1.230598in}{1.034105in}}%
\pgfpathlineto{\pgfqpoint{1.225313in}{1.018156in}}%
\pgfpathlineto{\pgfqpoint{1.220018in}{1.002179in}}%
\pgfpathlineto{\pgfqpoint{1.238449in}{0.998883in}}%
\pgfpathlineto{\pgfqpoint{1.256291in}{0.995102in}}%
\pgfpathclose%
\pgfusepath{fill}%
\end{pgfscope}%
\begin{pgfscope}%
\pgfpathrectangle{\pgfqpoint{0.407592in}{0.222904in}}{\pgfqpoint{1.411163in}{1.411163in}}%
\pgfusepath{clip}%
\pgfsetbuttcap%
\pgfsetroundjoin%
\definecolor{currentfill}{rgb}{0.749020,0.772549,0.792157}%
\pgfsetfillcolor{currentfill}%
\pgfsetfillopacity{0.160000}%
\pgfsetlinewidth{0.000000pt}%
\definecolor{currentstroke}{rgb}{0.000000,0.000000,0.000000}%
\pgfsetstrokecolor{currentstroke}%
\pgfsetdash{}{0pt}%
\pgfpathmoveto{\pgfqpoint{1.414631in}{1.124902in}}%
\pgfpathlineto{\pgfqpoint{1.424008in}{1.139770in}}%
\pgfpathlineto{\pgfqpoint{1.433375in}{1.154622in}}%
\pgfpathlineto{\pgfqpoint{1.411057in}{1.162458in}}%
\pgfpathlineto{\pgfqpoint{1.387508in}{1.169690in}}%
\pgfpathlineto{\pgfqpoint{1.379432in}{1.154500in}}%
\pgfpathlineto{\pgfqpoint{1.371346in}{1.139291in}}%
\pgfpathlineto{\pgfqpoint{1.393576in}{1.132384in}}%
\pgfpathlineto{\pgfqpoint{1.414631in}{1.124902in}}%
\pgfpathclose%
\pgfusepath{fill}%
\end{pgfscope}%
\begin{pgfscope}%
\pgfpathrectangle{\pgfqpoint{0.407592in}{0.222904in}}{\pgfqpoint{1.411163in}{1.411163in}}%
\pgfusepath{clip}%
\pgfsetbuttcap%
\pgfsetroundjoin%
\definecolor{currentfill}{rgb}{0.749020,0.772549,0.792157}%
\pgfsetfillcolor{currentfill}%
\pgfsetfillopacity{0.160000}%
\pgfsetlinewidth{0.000000pt}%
\definecolor{currentstroke}{rgb}{0.000000,0.000000,0.000000}%
\pgfsetstrokecolor{currentstroke}%
\pgfsetdash{}{0pt}%
\pgfpathmoveto{\pgfqpoint{0.919157in}{1.056927in}}%
\pgfpathlineto{\pgfqpoint{0.913177in}{1.072529in}}%
\pgfpathlineto{\pgfqpoint{0.907207in}{1.088109in}}%
\pgfpathlineto{\pgfqpoint{0.884743in}{1.082858in}}%
\pgfpathlineto{\pgfqpoint{0.863061in}{1.077055in}}%
\pgfpathlineto{\pgfqpoint{0.870497in}{1.061784in}}%
\pgfpathlineto{\pgfqpoint{0.877943in}{1.046493in}}%
\pgfpathlineto{\pgfqpoint{0.898181in}{1.051970in}}%
\pgfpathlineto{\pgfqpoint{0.919157in}{1.056927in}}%
\pgfpathclose%
\pgfusepath{fill}%
\end{pgfscope}%
\begin{pgfscope}%
\pgfpathrectangle{\pgfqpoint{0.407592in}{0.222904in}}{\pgfqpoint{1.411163in}{1.411163in}}%
\pgfusepath{clip}%
\pgfsetbuttcap%
\pgfsetroundjoin%
\definecolor{currentfill}{rgb}{0.749020,0.772549,0.792157}%
\pgfsetfillcolor{currentfill}%
\pgfsetfillopacity{0.160000}%
\pgfsetlinewidth{0.000000pt}%
\definecolor{currentstroke}{rgb}{0.000000,0.000000,0.000000}%
\pgfsetstrokecolor{currentstroke}%
\pgfsetdash{}{0pt}%
\pgfpathmoveto{\pgfqpoint{0.980795in}{1.001558in}}%
\pgfpathlineto{\pgfqpoint{0.976330in}{1.017512in}}%
\pgfpathlineto{\pgfqpoint{0.971872in}{1.033439in}}%
\pgfpathlineto{\pgfqpoint{0.951241in}{1.029805in}}%
\pgfpathlineto{\pgfqpoint{0.931142in}{1.025653in}}%
\pgfpathlineto{\pgfqpoint{0.937148in}{1.009981in}}%
\pgfpathlineto{\pgfqpoint{0.943162in}{0.994286in}}%
\pgfpathlineto{\pgfqpoint{0.961730in}{0.998163in}}%
\pgfpathlineto{\pgfqpoint{0.980795in}{1.001558in}}%
\pgfpathclose%
\pgfusepath{fill}%
\end{pgfscope}%
\begin{pgfscope}%
\pgfpathrectangle{\pgfqpoint{0.407592in}{0.222904in}}{\pgfqpoint{1.411163in}{1.411163in}}%
\pgfusepath{clip}%
\pgfsetbuttcap%
\pgfsetroundjoin%
\definecolor{currentfill}{rgb}{0.749020,0.772549,0.792157}%
\pgfsetfillcolor{currentfill}%
\pgfsetfillopacity{0.160000}%
\pgfsetlinewidth{0.000000pt}%
\definecolor{currentstroke}{rgb}{0.000000,0.000000,0.000000}%
\pgfsetstrokecolor{currentstroke}%
\pgfsetdash{}{0pt}%
\pgfpathmoveto{\pgfqpoint{0.665363in}{1.329663in}}%
\pgfpathlineto{\pgfqpoint{0.656718in}{1.344295in}}%
\pgfpathlineto{\pgfqpoint{0.648081in}{1.358914in}}%
\pgfpathlineto{\pgfqpoint{0.616296in}{1.347935in}}%
\pgfpathlineto{\pgfqpoint{0.586224in}{1.336227in}}%
\pgfpathlineto{\pgfqpoint{0.596147in}{1.321940in}}%
\pgfpathlineto{\pgfqpoint{0.606076in}{1.307644in}}%
\pgfpathlineto{\pgfqpoint{0.634889in}{1.319005in}}%
\pgfpathlineto{\pgfqpoint{0.665363in}{1.329663in}}%
\pgfpathclose%
\pgfusepath{fill}%
\end{pgfscope}%
\begin{pgfscope}%
\pgfpathrectangle{\pgfqpoint{0.407592in}{0.222904in}}{\pgfqpoint{1.411163in}{1.411163in}}%
\pgfusepath{clip}%
\pgfsetbuttcap%
\pgfsetroundjoin%
\definecolor{currentfill}{rgb}{0.749020,0.772549,0.792157}%
\pgfsetfillcolor{currentfill}%
\pgfsetfillopacity{0.160000}%
\pgfsetlinewidth{0.000000pt}%
\definecolor{currentstroke}{rgb}{0.000000,0.000000,0.000000}%
\pgfsetstrokecolor{currentstroke}%
\pgfsetdash{}{0pt}%
\pgfpathmoveto{\pgfqpoint{1.166722in}{0.942352in}}%
\pgfpathlineto{\pgfqpoint{1.170474in}{0.958626in}}%
\pgfpathlineto{\pgfqpoint{1.174219in}{0.974872in}}%
\pgfpathlineto{\pgfqpoint{1.156087in}{0.976524in}}%
\pgfpathlineto{\pgfqpoint{1.137691in}{0.977694in}}%
\pgfpathlineto{\pgfqpoint{1.135576in}{0.961336in}}%
\pgfpathlineto{\pgfqpoint{1.133458in}{0.944948in}}%
\pgfpathlineto{\pgfqpoint{1.150212in}{0.943872in}}%
\pgfpathlineto{\pgfqpoint{1.166722in}{0.942352in}}%
\pgfpathclose%
\pgfusepath{fill}%
\end{pgfscope}%
\begin{pgfscope}%
\pgfpathrectangle{\pgfqpoint{0.407592in}{0.222904in}}{\pgfqpoint{1.411163in}{1.411163in}}%
\pgfusepath{clip}%
\pgfsetbuttcap%
\pgfsetroundjoin%
\definecolor{currentfill}{rgb}{0.749020,0.772549,0.792157}%
\pgfsetfillcolor{currentfill}%
\pgfsetfillopacity{0.160000}%
\pgfsetlinewidth{0.000000pt}%
\definecolor{currentstroke}{rgb}{0.000000,0.000000,0.000000}%
\pgfsetstrokecolor{currentstroke}%
\pgfsetdash{}{0pt}%
\pgfpathmoveto{\pgfqpoint{1.578870in}{1.281285in}}%
\pgfpathlineto{\pgfqpoint{1.589325in}{1.295637in}}%
\pgfpathlineto{\pgfqpoint{1.599772in}{1.309978in}}%
\pgfpathlineto{\pgfqpoint{1.573320in}{1.321201in}}%
\pgfpathlineto{\pgfqpoint{1.545062in}{1.331713in}}%
\pgfpathlineto{\pgfqpoint{1.535805in}{1.317036in}}%
\pgfpathlineto{\pgfqpoint{1.526539in}{1.302343in}}%
\pgfpathlineto{\pgfqpoint{1.553576in}{1.292156in}}%
\pgfpathlineto{\pgfqpoint{1.578870in}{1.281285in}}%
\pgfpathclose%
\pgfusepath{fill}%
\end{pgfscope}%
\begin{pgfscope}%
\pgfpathrectangle{\pgfqpoint{0.407592in}{0.222904in}}{\pgfqpoint{1.411163in}{1.411163in}}%
\pgfusepath{clip}%
\pgfsetbuttcap%
\pgfsetroundjoin%
\definecolor{currentfill}{rgb}{0.749020,0.772549,0.792157}%
\pgfsetfillcolor{currentfill}%
\pgfsetfillopacity{0.160000}%
\pgfsetlinewidth{0.000000pt}%
\definecolor{currentstroke}{rgb}{0.000000,0.000000,0.000000}%
\pgfsetstrokecolor{currentstroke}%
\pgfsetdash{}{0pt}%
\pgfpathmoveto{\pgfqpoint{1.065350in}{0.944768in}}%
\pgfpathlineto{\pgfqpoint{1.064131in}{0.961148in}}%
\pgfpathlineto{\pgfqpoint{1.062914in}{0.977498in}}%
\pgfpathlineto{\pgfqpoint{1.044288in}{0.976231in}}%
\pgfpathlineto{\pgfqpoint{1.025837in}{0.974483in}}%
\pgfpathlineto{\pgfqpoint{1.028708in}{0.958253in}}%
\pgfpathlineto{\pgfqpoint{1.031585in}{0.941994in}}%
\pgfpathlineto{\pgfqpoint{1.048387in}{0.943602in}}%
\pgfpathlineto{\pgfqpoint{1.065350in}{0.944768in}}%
\pgfpathclose%
\pgfusepath{fill}%
\end{pgfscope}%
\begin{pgfscope}%
\pgfpathrectangle{\pgfqpoint{0.407592in}{0.222904in}}{\pgfqpoint{1.411163in}{1.411163in}}%
\pgfusepath{clip}%
\pgfsetbuttcap%
\pgfsetroundjoin%
\definecolor{currentfill}{rgb}{0.749020,0.772549,0.792157}%
\pgfsetfillcolor{currentfill}%
\pgfsetfillopacity{0.160000}%
\pgfsetlinewidth{0.000000pt}%
\definecolor{currentstroke}{rgb}{0.000000,0.000000,0.000000}%
\pgfsetstrokecolor{currentstroke}%
\pgfsetdash{}{0pt}%
\pgfpathmoveto{\pgfqpoint{0.833408in}{1.137953in}}%
\pgfpathlineto{\pgfqpoint{0.826017in}{1.153130in}}%
\pgfpathlineto{\pgfqpoint{0.818636in}{1.168289in}}%
\pgfpathlineto{\pgfqpoint{0.793635in}{1.160934in}}%
\pgfpathlineto{\pgfqpoint{0.769749in}{1.152979in}}%
\pgfpathlineto{\pgfqpoint{0.778502in}{1.138164in}}%
\pgfpathlineto{\pgfqpoint{0.787263in}{1.123334in}}%
\pgfpathlineto{\pgfqpoint{0.809804in}{1.130929in}}%
\pgfpathlineto{\pgfqpoint{0.833408in}{1.137953in}}%
\pgfpathclose%
\pgfusepath{fill}%
\end{pgfscope}%
\begin{pgfscope}%
\pgfpathrectangle{\pgfqpoint{0.407592in}{0.222904in}}{\pgfqpoint{1.411163in}{1.411163in}}%
\pgfusepath{clip}%
\pgfsetbuttcap%
\pgfsetroundjoin%
\definecolor{currentfill}{rgb}{0.749020,0.772549,0.792157}%
\pgfsetfillcolor{currentfill}%
\pgfsetfillopacity{0.160000}%
\pgfsetlinewidth{0.000000pt}%
\definecolor{currentstroke}{rgb}{0.000000,0.000000,0.000000}%
\pgfsetstrokecolor{currentstroke}%
\pgfsetdash{}{0pt}%
\pgfpathmoveto{\pgfqpoint{0.682677in}{1.300356in}}%
\pgfpathlineto{\pgfqpoint{0.674016in}{1.315016in}}%
\pgfpathlineto{\pgfqpoint{0.665363in}{1.329663in}}%
\pgfpathlineto{\pgfqpoint{0.634889in}{1.319005in}}%
\pgfpathlineto{\pgfqpoint{0.606076in}{1.307644in}}%
\pgfpathlineto{\pgfqpoint{0.616011in}{1.293339in}}%
\pgfpathlineto{\pgfqpoint{0.625953in}{1.279025in}}%
\pgfpathlineto{\pgfqpoint{0.653512in}{1.290029in}}%
\pgfpathlineto{\pgfqpoint{0.682677in}{1.300356in}}%
\pgfpathclose%
\pgfusepath{fill}%
\end{pgfscope}%
\begin{pgfscope}%
\pgfpathrectangle{\pgfqpoint{0.407592in}{0.222904in}}{\pgfqpoint{1.411163in}{1.411163in}}%
\pgfusepath{clip}%
\pgfsetbuttcap%
\pgfsetroundjoin%
\definecolor{currentfill}{rgb}{0.749020,0.772549,0.792157}%
\pgfsetfillcolor{currentfill}%
\pgfsetfillopacity{0.160000}%
\pgfsetlinewidth{0.000000pt}%
\definecolor{currentstroke}{rgb}{0.000000,0.000000,0.000000}%
\pgfsetstrokecolor{currentstroke}%
\pgfsetdash{}{0pt}%
\pgfpathmoveto{\pgfqpoint{1.306290in}{1.016925in}}%
\pgfpathlineto{\pgfqpoint{1.314458in}{1.032289in}}%
\pgfpathlineto{\pgfqpoint{1.322616in}{1.047633in}}%
\pgfpathlineto{\pgfqpoint{1.303421in}{1.053007in}}%
\pgfpathlineto{\pgfqpoint{1.283390in}{1.057857in}}%
\pgfpathlineto{\pgfqpoint{1.276630in}{1.042203in}}%
\pgfpathlineto{\pgfqpoint{1.269861in}{1.026526in}}%
\pgfpathlineto{\pgfqpoint{1.288468in}{1.021971in}}%
\pgfpathlineto{\pgfqpoint{1.306290in}{1.016925in}}%
\pgfpathclose%
\pgfusepath{fill}%
\end{pgfscope}%
\begin{pgfscope}%
\pgfpathrectangle{\pgfqpoint{0.407592in}{0.222904in}}{\pgfqpoint{1.411163in}{1.411163in}}%
\pgfusepath{clip}%
\pgfsetbuttcap%
\pgfsetroundjoin%
\definecolor{currentfill}{rgb}{0.749020,0.772549,0.792157}%
\pgfsetfillcolor{currentfill}%
\pgfsetfillopacity{0.160000}%
\pgfsetlinewidth{0.000000pt}%
\definecolor{currentstroke}{rgb}{0.000000,0.000000,0.000000}%
\pgfsetstrokecolor{currentstroke}%
\pgfsetdash{}{0pt}%
\pgfpathmoveto{\pgfqpoint{1.557939in}{1.252553in}}%
\pgfpathlineto{\pgfqpoint{1.568408in}{1.266924in}}%
\pgfpathlineto{\pgfqpoint{1.578870in}{1.281285in}}%
\pgfpathlineto{\pgfqpoint{1.553576in}{1.292156in}}%
\pgfpathlineto{\pgfqpoint{1.526539in}{1.302343in}}%
\pgfpathlineto{\pgfqpoint{1.517264in}{1.287636in}}%
\pgfpathlineto{\pgfqpoint{1.507980in}{1.272915in}}%
\pgfpathlineto{\pgfqpoint{1.533800in}{1.263064in}}%
\pgfpathlineto{\pgfqpoint{1.557939in}{1.252553in}}%
\pgfpathclose%
\pgfusepath{fill}%
\end{pgfscope}%
\begin{pgfscope}%
\pgfpathrectangle{\pgfqpoint{0.407592in}{0.222904in}}{\pgfqpoint{1.411163in}{1.411163in}}%
\pgfusepath{clip}%
\pgfsetbuttcap%
\pgfsetroundjoin%
\definecolor{currentfill}{rgb}{0.749020,0.772549,0.792157}%
\pgfsetfillcolor{currentfill}%
\pgfsetfillopacity{0.160000}%
\pgfsetlinewidth{0.000000pt}%
\definecolor{currentstroke}{rgb}{0.000000,0.000000,0.000000}%
\pgfsetstrokecolor{currentstroke}%
\pgfsetdash{}{0pt}%
\pgfpathmoveto{\pgfqpoint{1.395851in}{1.095124in}}%
\pgfpathlineto{\pgfqpoint{1.405246in}{1.110021in}}%
\pgfpathlineto{\pgfqpoint{1.414631in}{1.124902in}}%
\pgfpathlineto{\pgfqpoint{1.393576in}{1.132384in}}%
\pgfpathlineto{\pgfqpoint{1.371346in}{1.139291in}}%
\pgfpathlineto{\pgfqpoint{1.363250in}{1.124063in}}%
\pgfpathlineto{\pgfqpoint{1.355143in}{1.108815in}}%
\pgfpathlineto{\pgfqpoint{1.376055in}{1.102243in}}%
\pgfpathlineto{\pgfqpoint{1.395851in}{1.095124in}}%
\pgfpathclose%
\pgfusepath{fill}%
\end{pgfscope}%
\begin{pgfscope}%
\pgfpathrectangle{\pgfqpoint{0.407592in}{0.222904in}}{\pgfqpoint{1.411163in}{1.411163in}}%
\pgfusepath{clip}%
\pgfsetbuttcap%
\pgfsetroundjoin%
\definecolor{currentfill}{rgb}{0.749020,0.772549,0.792157}%
\pgfsetfillcolor{currentfill}%
\pgfsetfillopacity{0.160000}%
\pgfsetlinewidth{0.000000pt}%
\definecolor{currentstroke}{rgb}{0.000000,0.000000,0.000000}%
\pgfsetstrokecolor{currentstroke}%
\pgfsetdash{}{0pt}%
\pgfpathmoveto{\pgfqpoint{1.198751in}{0.938005in}}%
\pgfpathlineto{\pgfqpoint{1.204082in}{0.954089in}}%
\pgfpathlineto{\pgfqpoint{1.209403in}{0.970146in}}%
\pgfpathlineto{\pgfqpoint{1.192014in}{0.972743in}}%
\pgfpathlineto{\pgfqpoint{1.174219in}{0.974872in}}%
\pgfpathlineto{\pgfqpoint{1.170474in}{0.958626in}}%
\pgfpathlineto{\pgfqpoint{1.166722in}{0.942352in}}%
\pgfpathlineto{\pgfqpoint{1.182924in}{0.940394in}}%
\pgfpathlineto{\pgfqpoint{1.198751in}{0.938005in}}%
\pgfpathclose%
\pgfusepath{fill}%
\end{pgfscope}%
\begin{pgfscope}%
\pgfpathrectangle{\pgfqpoint{0.407592in}{0.222904in}}{\pgfqpoint{1.411163in}{1.411163in}}%
\pgfusepath{clip}%
\pgfsetbuttcap%
\pgfsetroundjoin%
\definecolor{currentfill}{rgb}{0.749020,0.772549,0.792157}%
\pgfsetfillcolor{currentfill}%
\pgfsetfillopacity{0.160000}%
\pgfsetlinewidth{0.000000pt}%
\definecolor{currentstroke}{rgb}{0.000000,0.000000,0.000000}%
\pgfsetstrokecolor{currentstroke}%
\pgfsetdash{}{0pt}%
\pgfpathmoveto{\pgfqpoint{0.931142in}{1.025653in}}%
\pgfpathlineto{\pgfqpoint{0.925145in}{1.041301in}}%
\pgfpathlineto{\pgfqpoint{0.919157in}{1.056927in}}%
\pgfpathlineto{\pgfqpoint{0.898181in}{1.051970in}}%
\pgfpathlineto{\pgfqpoint{0.877943in}{1.046493in}}%
\pgfpathlineto{\pgfqpoint{0.885397in}{1.031184in}}%
\pgfpathlineto{\pgfqpoint{0.892861in}{1.015856in}}%
\pgfpathlineto{\pgfqpoint{0.911656in}{1.020998in}}%
\pgfpathlineto{\pgfqpoint{0.931142in}{1.025653in}}%
\pgfpathclose%
\pgfusepath{fill}%
\end{pgfscope}%
\begin{pgfscope}%
\pgfpathrectangle{\pgfqpoint{0.407592in}{0.222904in}}{\pgfqpoint{1.411163in}{1.411163in}}%
\pgfusepath{clip}%
\pgfsetbuttcap%
\pgfsetroundjoin%
\definecolor{currentfill}{rgb}{0.749020,0.772549,0.792157}%
\pgfsetfillcolor{currentfill}%
\pgfsetfillopacity{0.160000}%
\pgfsetlinewidth{0.000000pt}%
\definecolor{currentstroke}{rgb}{0.000000,0.000000,0.000000}%
\pgfsetstrokecolor{currentstroke}%
\pgfsetdash{}{0pt}%
\pgfpathmoveto{\pgfqpoint{1.242681in}{0.963583in}}%
\pgfpathlineto{\pgfqpoint{1.249491in}{0.979355in}}%
\pgfpathlineto{\pgfqpoint{1.256291in}{0.995102in}}%
\pgfpathlineto{\pgfqpoint{1.238449in}{0.998883in}}%
\pgfpathlineto{\pgfqpoint{1.220018in}{1.002179in}}%
\pgfpathlineto{\pgfqpoint{1.214715in}{0.986176in}}%
\pgfpathlineto{\pgfqpoint{1.209403in}{0.970146in}}%
\pgfpathlineto{\pgfqpoint{1.226314in}{0.967089in}}%
\pgfpathlineto{\pgfqpoint{1.242681in}{0.963583in}}%
\pgfpathclose%
\pgfusepath{fill}%
\end{pgfscope}%
\begin{pgfscope}%
\pgfpathrectangle{\pgfqpoint{0.407592in}{0.222904in}}{\pgfqpoint{1.411163in}{1.411163in}}%
\pgfusepath{clip}%
\pgfsetbuttcap%
\pgfsetroundjoin%
\definecolor{currentfill}{rgb}{0.749020,0.772549,0.792157}%
\pgfsetfillcolor{currentfill}%
\pgfsetfillopacity{0.160000}%
\pgfsetlinewidth{0.000000pt}%
\definecolor{currentstroke}{rgb}{0.000000,0.000000,0.000000}%
\pgfsetstrokecolor{currentstroke}%
\pgfsetdash{}{0pt}%
\pgfpathmoveto{\pgfqpoint{1.031585in}{0.941994in}}%
\pgfpathlineto{\pgfqpoint{1.028708in}{0.958253in}}%
\pgfpathlineto{\pgfqpoint{1.025837in}{0.974483in}}%
\pgfpathlineto{\pgfqpoint{1.007632in}{0.972260in}}%
\pgfpathlineto{\pgfqpoint{0.989747in}{0.969569in}}%
\pgfpathlineto{\pgfqpoint{0.994234in}{0.953535in}}%
\pgfpathlineto{\pgfqpoint{0.998729in}{0.937475in}}%
\pgfpathlineto{\pgfqpoint{1.015010in}{0.939949in}}%
\pgfpathlineto{\pgfqpoint{1.031585in}{0.941994in}}%
\pgfpathclose%
\pgfusepath{fill}%
\end{pgfscope}%
\begin{pgfscope}%
\pgfpathrectangle{\pgfqpoint{0.407592in}{0.222904in}}{\pgfqpoint{1.411163in}{1.411163in}}%
\pgfusepath{clip}%
\pgfsetbuttcap%
\pgfsetroundjoin%
\definecolor{currentfill}{rgb}{0.749020,0.772549,0.792157}%
\pgfsetfillcolor{currentfill}%
\pgfsetfillopacity{0.160000}%
\pgfsetlinewidth{0.000000pt}%
\definecolor{currentstroke}{rgb}{0.000000,0.000000,0.000000}%
\pgfsetstrokecolor{currentstroke}%
\pgfsetdash{}{0pt}%
\pgfpathmoveto{\pgfqpoint{1.129209in}{0.912078in}}%
\pgfpathlineto{\pgfqpoint{1.131335in}{0.928529in}}%
\pgfpathlineto{\pgfqpoint{1.133458in}{0.944948in}}%
\pgfpathlineto{\pgfqpoint{1.116528in}{0.945577in}}%
\pgfpathlineto{\pgfqpoint{1.099488in}{0.945757in}}%
\pgfpathlineto{\pgfqpoint{1.099034in}{0.929302in}}%
\pgfpathlineto{\pgfqpoint{1.098578in}{0.912815in}}%
\pgfpathlineto{\pgfqpoint{1.113943in}{0.912651in}}%
\pgfpathlineto{\pgfqpoint{1.129209in}{0.912078in}}%
\pgfpathclose%
\pgfusepath{fill}%
\end{pgfscope}%
\begin{pgfscope}%
\pgfpathrectangle{\pgfqpoint{0.407592in}{0.222904in}}{\pgfqpoint{1.411163in}{1.411163in}}%
\pgfusepath{clip}%
\pgfsetbuttcap%
\pgfsetroundjoin%
\definecolor{currentfill}{rgb}{0.749020,0.772549,0.792157}%
\pgfsetfillcolor{currentfill}%
\pgfsetfillopacity{0.160000}%
\pgfsetlinewidth{0.000000pt}%
\definecolor{currentstroke}{rgb}{0.000000,0.000000,0.000000}%
\pgfsetstrokecolor{currentstroke}%
\pgfsetdash{}{0pt}%
\pgfpathmoveto{\pgfqpoint{1.098578in}{0.912815in}}%
\pgfpathlineto{\pgfqpoint{1.099034in}{0.929302in}}%
\pgfpathlineto{\pgfqpoint{1.099488in}{0.945757in}}%
\pgfpathlineto{\pgfqpoint{1.082406in}{0.945487in}}%
\pgfpathlineto{\pgfqpoint{1.065350in}{0.944768in}}%
\pgfpathlineto{\pgfqpoint{1.066572in}{0.928356in}}%
\pgfpathlineto{\pgfqpoint{1.067796in}{0.911914in}}%
\pgfpathlineto{\pgfqpoint{1.083175in}{0.912569in}}%
\pgfpathlineto{\pgfqpoint{1.098578in}{0.912815in}}%
\pgfpathclose%
\pgfusepath{fill}%
\end{pgfscope}%
\begin{pgfscope}%
\pgfpathrectangle{\pgfqpoint{0.407592in}{0.222904in}}{\pgfqpoint{1.411163in}{1.411163in}}%
\pgfusepath{clip}%
\pgfsetbuttcap%
\pgfsetroundjoin%
\definecolor{currentfill}{rgb}{0.749020,0.772549,0.792157}%
\pgfsetfillcolor{currentfill}%
\pgfsetfillopacity{0.160000}%
\pgfsetlinewidth{0.000000pt}%
\definecolor{currentstroke}{rgb}{0.000000,0.000000,0.000000}%
\pgfsetstrokecolor{currentstroke}%
\pgfsetdash{}{0pt}%
\pgfpathmoveto{\pgfqpoint{0.848216in}{1.107541in}}%
\pgfpathlineto{\pgfqpoint{0.840807in}{1.122756in}}%
\pgfpathlineto{\pgfqpoint{0.833408in}{1.137953in}}%
\pgfpathlineto{\pgfqpoint{0.809804in}{1.130929in}}%
\pgfpathlineto{\pgfqpoint{0.787263in}{1.123334in}}%
\pgfpathlineto{\pgfqpoint{0.796033in}{1.108491in}}%
\pgfpathlineto{\pgfqpoint{0.804811in}{1.093633in}}%
\pgfpathlineto{\pgfqpoint{0.826008in}{1.100858in}}%
\pgfpathlineto{\pgfqpoint{0.848216in}{1.107541in}}%
\pgfpathclose%
\pgfusepath{fill}%
\end{pgfscope}%
\begin{pgfscope}%
\pgfpathrectangle{\pgfqpoint{0.407592in}{0.222904in}}{\pgfqpoint{1.411163in}{1.411163in}}%
\pgfusepath{clip}%
\pgfsetbuttcap%
\pgfsetroundjoin%
\definecolor{currentfill}{rgb}{0.749020,0.772549,0.792157}%
\pgfsetfillcolor{currentfill}%
\pgfsetfillopacity{0.160000}%
\pgfsetlinewidth{0.000000pt}%
\definecolor{currentstroke}{rgb}{0.000000,0.000000,0.000000}%
\pgfsetstrokecolor{currentstroke}%
\pgfsetdash{}{0pt}%
\pgfpathmoveto{\pgfqpoint{0.989747in}{0.969569in}}%
\pgfpathlineto{\pgfqpoint{0.985267in}{0.985577in}}%
\pgfpathlineto{\pgfqpoint{0.980795in}{1.001558in}}%
\pgfpathlineto{\pgfqpoint{0.961730in}{0.998163in}}%
\pgfpathlineto{\pgfqpoint{0.943162in}{0.994286in}}%
\pgfpathlineto{\pgfqpoint{0.949186in}{0.978568in}}%
\pgfpathlineto{\pgfqpoint{0.955219in}{0.962827in}}%
\pgfpathlineto{\pgfqpoint{0.972252in}{0.966422in}}%
\pgfpathlineto{\pgfqpoint{0.989747in}{0.969569in}}%
\pgfpathclose%
\pgfusepath{fill}%
\end{pgfscope}%
\begin{pgfscope}%
\pgfpathrectangle{\pgfqpoint{0.407592in}{0.222904in}}{\pgfqpoint{1.411163in}{1.411163in}}%
\pgfusepath{clip}%
\pgfsetbuttcap%
\pgfsetroundjoin%
\definecolor{currentfill}{rgb}{0.749020,0.772549,0.792157}%
\pgfsetfillcolor{currentfill}%
\pgfsetfillopacity{0.160000}%
\pgfsetlinewidth{0.000000pt}%
\definecolor{currentstroke}{rgb}{0.000000,0.000000,0.000000}%
\pgfsetstrokecolor{currentstroke}%
\pgfsetdash{}{0pt}%
\pgfpathmoveto{\pgfqpoint{0.700025in}{1.270993in}}%
\pgfpathlineto{\pgfqpoint{0.691347in}{1.285681in}}%
\pgfpathlineto{\pgfqpoint{0.682677in}{1.300356in}}%
\pgfpathlineto{\pgfqpoint{0.653512in}{1.290029in}}%
\pgfpathlineto{\pgfqpoint{0.625953in}{1.279025in}}%
\pgfpathlineto{\pgfqpoint{0.635901in}{1.264701in}}%
\pgfpathlineto{\pgfqpoint{0.645855in}{1.250368in}}%
\pgfpathlineto{\pgfqpoint{0.672165in}{1.261007in}}%
\pgfpathlineto{\pgfqpoint{0.700025in}{1.270993in}}%
\pgfpathclose%
\pgfusepath{fill}%
\end{pgfscope}%
\begin{pgfscope}%
\pgfpathrectangle{\pgfqpoint{0.407592in}{0.222904in}}{\pgfqpoint{1.411163in}{1.411163in}}%
\pgfusepath{clip}%
\pgfsetbuttcap%
\pgfsetroundjoin%
\definecolor{currentfill}{rgb}{0.749020,0.772549,0.792157}%
\pgfsetfillcolor{currentfill}%
\pgfsetfillopacity{0.160000}%
\pgfsetlinewidth{0.000000pt}%
\definecolor{currentstroke}{rgb}{0.000000,0.000000,0.000000}%
\pgfsetstrokecolor{currentstroke}%
\pgfsetdash{}{0pt}%
\pgfpathmoveto{\pgfqpoint{1.536980in}{1.223783in}}%
\pgfpathlineto{\pgfqpoint{1.547463in}{1.238173in}}%
\pgfpathlineto{\pgfqpoint{1.557939in}{1.252553in}}%
\pgfpathlineto{\pgfqpoint{1.533800in}{1.263064in}}%
\pgfpathlineto{\pgfqpoint{1.507980in}{1.272915in}}%
\pgfpathlineto{\pgfqpoint{1.498686in}{1.258179in}}%
\pgfpathlineto{\pgfqpoint{1.489383in}{1.243429in}}%
\pgfpathlineto{\pgfqpoint{1.513991in}{1.233922in}}%
\pgfpathlineto{\pgfqpoint{1.536980in}{1.223783in}}%
\pgfpathclose%
\pgfusepath{fill}%
\end{pgfscope}%
\begin{pgfscope}%
\pgfpathrectangle{\pgfqpoint{0.407592in}{0.222904in}}{\pgfqpoint{1.411163in}{1.411163in}}%
\pgfusepath{clip}%
\pgfsetbuttcap%
\pgfsetroundjoin%
\definecolor{currentfill}{rgb}{0.749020,0.772549,0.792157}%
\pgfsetfillcolor{currentfill}%
\pgfsetfillopacity{0.160000}%
\pgfsetlinewidth{0.000000pt}%
\definecolor{currentstroke}{rgb}{0.000000,0.000000,0.000000}%
\pgfsetstrokecolor{currentstroke}%
\pgfsetdash{}{0pt}%
\pgfpathmoveto{\pgfqpoint{1.159198in}{0.909713in}}%
\pgfpathlineto{\pgfqpoint{1.162964in}{0.926047in}}%
\pgfpathlineto{\pgfqpoint{1.166722in}{0.942352in}}%
\pgfpathlineto{\pgfqpoint{1.150212in}{0.943872in}}%
\pgfpathlineto{\pgfqpoint{1.133458in}{0.944948in}}%
\pgfpathlineto{\pgfqpoint{1.131335in}{0.928529in}}%
\pgfpathlineto{\pgfqpoint{1.129209in}{0.912078in}}%
\pgfpathlineto{\pgfqpoint{1.144314in}{0.911098in}}%
\pgfpathlineto{\pgfqpoint{1.159198in}{0.909713in}}%
\pgfpathclose%
\pgfusepath{fill}%
\end{pgfscope}%
\begin{pgfscope}%
\pgfpathrectangle{\pgfqpoint{0.407592in}{0.222904in}}{\pgfqpoint{1.411163in}{1.411163in}}%
\pgfusepath{clip}%
\pgfsetbuttcap%
\pgfsetroundjoin%
\definecolor{currentfill}{rgb}{0.749020,0.772549,0.792157}%
\pgfsetfillcolor{currentfill}%
\pgfsetfillopacity{0.160000}%
\pgfsetlinewidth{0.000000pt}%
\definecolor{currentstroke}{rgb}{0.000000,0.000000,0.000000}%
\pgfsetstrokecolor{currentstroke}%
\pgfsetdash{}{0pt}%
\pgfpathmoveto{\pgfqpoint{1.067796in}{0.911914in}}%
\pgfpathlineto{\pgfqpoint{1.066572in}{0.928356in}}%
\pgfpathlineto{\pgfqpoint{1.065350in}{0.944768in}}%
\pgfpathlineto{\pgfqpoint{1.048387in}{0.943602in}}%
\pgfpathlineto{\pgfqpoint{1.031585in}{0.941994in}}%
\pgfpathlineto{\pgfqpoint{1.034467in}{0.925705in}}%
\pgfpathlineto{\pgfqpoint{1.037354in}{0.909387in}}%
\pgfpathlineto{\pgfqpoint{1.052502in}{0.910852in}}%
\pgfpathlineto{\pgfqpoint{1.067796in}{0.911914in}}%
\pgfpathclose%
\pgfusepath{fill}%
\end{pgfscope}%
\begin{pgfscope}%
\pgfpathrectangle{\pgfqpoint{0.407592in}{0.222904in}}{\pgfqpoint{1.411163in}{1.411163in}}%
\pgfusepath{clip}%
\pgfsetbuttcap%
\pgfsetroundjoin%
\definecolor{currentfill}{rgb}{0.749020,0.772549,0.792157}%
\pgfsetfillcolor{currentfill}%
\pgfsetfillopacity{0.160000}%
\pgfsetlinewidth{0.000000pt}%
\definecolor{currentstroke}{rgb}{0.000000,0.000000,0.000000}%
\pgfsetstrokecolor{currentstroke}%
\pgfsetdash{}{0pt}%
\pgfpathmoveto{\pgfqpoint{1.377033in}{1.065287in}}%
\pgfpathlineto{\pgfqpoint{1.386447in}{1.080213in}}%
\pgfpathlineto{\pgfqpoint{1.395851in}{1.095124in}}%
\pgfpathlineto{\pgfqpoint{1.376055in}{1.102243in}}%
\pgfpathlineto{\pgfqpoint{1.355143in}{1.108815in}}%
\pgfpathlineto{\pgfqpoint{1.347027in}{1.093549in}}%
\pgfpathlineto{\pgfqpoint{1.338900in}{1.078263in}}%
\pgfpathlineto{\pgfqpoint{1.358494in}{1.072033in}}%
\pgfpathlineto{\pgfqpoint{1.377033in}{1.065287in}}%
\pgfpathclose%
\pgfusepath{fill}%
\end{pgfscope}%
\begin{pgfscope}%
\pgfpathrectangle{\pgfqpoint{0.407592in}{0.222904in}}{\pgfqpoint{1.411163in}{1.411163in}}%
\pgfusepath{clip}%
\pgfsetbuttcap%
\pgfsetroundjoin%
\definecolor{currentfill}{rgb}{0.749020,0.772549,0.792157}%
\pgfsetfillcolor{currentfill}%
\pgfsetfillopacity{0.160000}%
\pgfsetlinewidth{0.000000pt}%
\definecolor{currentstroke}{rgb}{0.000000,0.000000,0.000000}%
\pgfsetstrokecolor{currentstroke}%
\pgfsetdash{}{0pt}%
\pgfpathmoveto{\pgfqpoint{1.289923in}{0.986140in}}%
\pgfpathlineto{\pgfqpoint{1.298111in}{1.001543in}}%
\pgfpathlineto{\pgfqpoint{1.306290in}{1.016925in}}%
\pgfpathlineto{\pgfqpoint{1.288468in}{1.021971in}}%
\pgfpathlineto{\pgfqpoint{1.269861in}{1.026526in}}%
\pgfpathlineto{\pgfqpoint{1.263081in}{1.010826in}}%
\pgfpathlineto{\pgfqpoint{1.256291in}{0.995102in}}%
\pgfpathlineto{\pgfqpoint{1.273473in}{0.990849in}}%
\pgfpathlineto{\pgfqpoint{1.289923in}{0.986140in}}%
\pgfpathclose%
\pgfusepath{fill}%
\end{pgfscope}%
\begin{pgfscope}%
\pgfpathrectangle{\pgfqpoint{0.407592in}{0.222904in}}{\pgfqpoint{1.411163in}{1.411163in}}%
\pgfusepath{clip}%
\pgfsetbuttcap%
\pgfsetroundjoin%
\definecolor{currentfill}{rgb}{0.749020,0.772549,0.792157}%
\pgfsetfillcolor{currentfill}%
\pgfsetfillopacity{0.160000}%
\pgfsetlinewidth{0.000000pt}%
\definecolor{currentstroke}{rgb}{0.000000,0.000000,0.000000}%
\pgfsetstrokecolor{currentstroke}%
\pgfsetdash{}{0pt}%
\pgfpathmoveto{\pgfqpoint{0.717406in}{1.241574in}}%
\pgfpathlineto{\pgfqpoint{0.708711in}{1.256291in}}%
\pgfpathlineto{\pgfqpoint{0.700025in}{1.270993in}}%
\pgfpathlineto{\pgfqpoint{0.672165in}{1.261007in}}%
\pgfpathlineto{\pgfqpoint{0.645855in}{1.250368in}}%
\pgfpathlineto{\pgfqpoint{0.655815in}{1.236026in}}%
\pgfpathlineto{\pgfqpoint{0.665782in}{1.221675in}}%
\pgfpathlineto{\pgfqpoint{0.690848in}{1.231938in}}%
\pgfpathlineto{\pgfqpoint{0.717406in}{1.241574in}}%
\pgfpathclose%
\pgfusepath{fill}%
\end{pgfscope}%
\begin{pgfscope}%
\pgfpathrectangle{\pgfqpoint{0.407592in}{0.222904in}}{\pgfqpoint{1.411163in}{1.411163in}}%
\pgfusepath{clip}%
\pgfsetbuttcap%
\pgfsetroundjoin%
\definecolor{currentfill}{rgb}{0.749020,0.772549,0.792157}%
\pgfsetfillcolor{currentfill}%
\pgfsetfillopacity{0.160000}%
\pgfsetlinewidth{0.000000pt}%
\definecolor{currentstroke}{rgb}{0.000000,0.000000,0.000000}%
\pgfsetstrokecolor{currentstroke}%
\pgfsetdash{}{0pt}%
\pgfpathmoveto{\pgfqpoint{0.863061in}{1.077055in}}%
\pgfpathlineto{\pgfqpoint{0.855634in}{1.092308in}}%
\pgfpathlineto{\pgfqpoint{0.848216in}{1.107541in}}%
\pgfpathlineto{\pgfqpoint{0.826008in}{1.100858in}}%
\pgfpathlineto{\pgfqpoint{0.804811in}{1.093633in}}%
\pgfpathlineto{\pgfqpoint{0.813598in}{1.078760in}}%
\pgfpathlineto{\pgfqpoint{0.822393in}{1.063874in}}%
\pgfpathlineto{\pgfqpoint{0.842249in}{1.070720in}}%
\pgfpathlineto{\pgfqpoint{0.863061in}{1.077055in}}%
\pgfpathclose%
\pgfusepath{fill}%
\end{pgfscope}%
\begin{pgfscope}%
\pgfpathrectangle{\pgfqpoint{0.407592in}{0.222904in}}{\pgfqpoint{1.411163in}{1.411163in}}%
\pgfusepath{clip}%
\pgfsetbuttcap%
\pgfsetroundjoin%
\definecolor{currentfill}{rgb}{0.749020,0.772549,0.792157}%
\pgfsetfillcolor{currentfill}%
\pgfsetfillopacity{0.160000}%
\pgfsetlinewidth{0.000000pt}%
\definecolor{currentstroke}{rgb}{0.000000,0.000000,0.000000}%
\pgfsetstrokecolor{currentstroke}%
\pgfsetdash{}{0pt}%
\pgfpathmoveto{\pgfqpoint{1.515992in}{1.194974in}}%
\pgfpathlineto{\pgfqpoint{1.526489in}{1.209383in}}%
\pgfpathlineto{\pgfqpoint{1.536980in}{1.223783in}}%
\pgfpathlineto{\pgfqpoint{1.513991in}{1.233922in}}%
\pgfpathlineto{\pgfqpoint{1.489383in}{1.243429in}}%
\pgfpathlineto{\pgfqpoint{1.480072in}{1.228664in}}%
\pgfpathlineto{\pgfqpoint{1.470751in}{1.213885in}}%
\pgfpathlineto{\pgfqpoint{1.494148in}{1.204732in}}%
\pgfpathlineto{\pgfqpoint{1.515992in}{1.194974in}}%
\pgfpathclose%
\pgfusepath{fill}%
\end{pgfscope}%
\begin{pgfscope}%
\pgfpathrectangle{\pgfqpoint{0.407592in}{0.222904in}}{\pgfqpoint{1.411163in}{1.411163in}}%
\pgfusepath{clip}%
\pgfsetbuttcap%
\pgfsetroundjoin%
\definecolor{currentfill}{rgb}{0.749020,0.772549,0.792157}%
\pgfsetfillcolor{currentfill}%
\pgfsetfillopacity{0.160000}%
\pgfsetlinewidth{0.000000pt}%
\definecolor{currentstroke}{rgb}{0.000000,0.000000,0.000000}%
\pgfsetstrokecolor{currentstroke}%
\pgfsetdash{}{0pt}%
\pgfpathmoveto{\pgfqpoint{0.943162in}{0.994286in}}%
\pgfpathlineto{\pgfqpoint{0.937148in}{1.009981in}}%
\pgfpathlineto{\pgfqpoint{0.931142in}{1.025653in}}%
\pgfpathlineto{\pgfqpoint{0.911656in}{1.020998in}}%
\pgfpathlineto{\pgfqpoint{0.892861in}{1.015856in}}%
\pgfpathlineto{\pgfqpoint{0.900335in}{1.000508in}}%
\pgfpathlineto{\pgfqpoint{0.907817in}{0.985142in}}%
\pgfpathlineto{\pgfqpoint{0.925167in}{0.989941in}}%
\pgfpathlineto{\pgfqpoint{0.943162in}{0.994286in}}%
\pgfpathclose%
\pgfusepath{fill}%
\end{pgfscope}%
\begin{pgfscope}%
\pgfpathrectangle{\pgfqpoint{0.407592in}{0.222904in}}{\pgfqpoint{1.411163in}{1.411163in}}%
\pgfusepath{clip}%
\pgfsetbuttcap%
\pgfsetroundjoin%
\definecolor{currentfill}{rgb}{0.749020,0.772549,0.792157}%
\pgfsetfillcolor{currentfill}%
\pgfsetfillopacity{0.160000}%
\pgfsetlinewidth{0.000000pt}%
\definecolor{currentstroke}{rgb}{0.000000,0.000000,0.000000}%
\pgfsetstrokecolor{currentstroke}%
\pgfsetdash{}{0pt}%
\pgfpathmoveto{\pgfqpoint{1.229030in}{0.931970in}}%
\pgfpathlineto{\pgfqpoint{1.235861in}{0.947789in}}%
\pgfpathlineto{\pgfqpoint{1.242681in}{0.963583in}}%
\pgfpathlineto{\pgfqpoint{1.226314in}{0.967089in}}%
\pgfpathlineto{\pgfqpoint{1.209403in}{0.970146in}}%
\pgfpathlineto{\pgfqpoint{1.204082in}{0.954089in}}%
\pgfpathlineto{\pgfqpoint{1.198751in}{0.938005in}}%
\pgfpathlineto{\pgfqpoint{1.214141in}{0.935194in}}%
\pgfpathlineto{\pgfqpoint{1.229030in}{0.931970in}}%
\pgfpathclose%
\pgfusepath{fill}%
\end{pgfscope}%
\begin{pgfscope}%
\pgfpathrectangle{\pgfqpoint{0.407592in}{0.222904in}}{\pgfqpoint{1.411163in}{1.411163in}}%
\pgfusepath{clip}%
\pgfsetbuttcap%
\pgfsetroundjoin%
\definecolor{currentfill}{rgb}{0.749020,0.772549,0.792157}%
\pgfsetfillcolor{currentfill}%
\pgfsetfillopacity{0.160000}%
\pgfsetlinewidth{0.000000pt}%
\definecolor{currentstroke}{rgb}{0.000000,0.000000,0.000000}%
\pgfsetstrokecolor{currentstroke}%
\pgfsetdash{}{0pt}%
\pgfpathmoveto{\pgfqpoint{1.188064in}{0.905755in}}%
\pgfpathlineto{\pgfqpoint{1.193412in}{0.921894in}}%
\pgfpathlineto{\pgfqpoint{1.198751in}{0.938005in}}%
\pgfpathlineto{\pgfqpoint{1.182924in}{0.940394in}}%
\pgfpathlineto{\pgfqpoint{1.166722in}{0.942352in}}%
\pgfpathlineto{\pgfqpoint{1.162964in}{0.926047in}}%
\pgfpathlineto{\pgfqpoint{1.159198in}{0.909713in}}%
\pgfpathlineto{\pgfqpoint{1.173801in}{0.907930in}}%
\pgfpathlineto{\pgfqpoint{1.188064in}{0.905755in}}%
\pgfpathclose%
\pgfusepath{fill}%
\end{pgfscope}%
\begin{pgfscope}%
\pgfpathrectangle{\pgfqpoint{0.407592in}{0.222904in}}{\pgfqpoint{1.411163in}{1.411163in}}%
\pgfusepath{clip}%
\pgfsetbuttcap%
\pgfsetroundjoin%
\definecolor{currentfill}{rgb}{0.749020,0.772549,0.792157}%
\pgfsetfillcolor{currentfill}%
\pgfsetfillopacity{0.160000}%
\pgfsetlinewidth{0.000000pt}%
\definecolor{currentstroke}{rgb}{0.000000,0.000000,0.000000}%
\pgfsetstrokecolor{currentstroke}%
\pgfsetdash{}{0pt}%
\pgfpathmoveto{\pgfqpoint{0.998729in}{0.937475in}}%
\pgfpathlineto{\pgfqpoint{0.994234in}{0.953535in}}%
\pgfpathlineto{\pgfqpoint{0.989747in}{0.969569in}}%
\pgfpathlineto{\pgfqpoint{0.972252in}{0.966422in}}%
\pgfpathlineto{\pgfqpoint{0.955219in}{0.962827in}}%
\pgfpathlineto{\pgfqpoint{0.961260in}{0.947063in}}%
\pgfpathlineto{\pgfqpoint{0.967310in}{0.931275in}}%
\pgfpathlineto{\pgfqpoint{0.982808in}{0.934580in}}%
\pgfpathlineto{\pgfqpoint{0.998729in}{0.937475in}}%
\pgfpathclose%
\pgfusepath{fill}%
\end{pgfscope}%
\begin{pgfscope}%
\pgfpathrectangle{\pgfqpoint{0.407592in}{0.222904in}}{\pgfqpoint{1.411163in}{1.411163in}}%
\pgfusepath{clip}%
\pgfsetbuttcap%
\pgfsetroundjoin%
\definecolor{currentfill}{rgb}{0.749020,0.772549,0.792157}%
\pgfsetfillcolor{currentfill}%
\pgfsetfillopacity{0.160000}%
\pgfsetlinewidth{0.000000pt}%
\definecolor{currentstroke}{rgb}{0.000000,0.000000,0.000000}%
\pgfsetstrokecolor{currentstroke}%
\pgfsetdash{}{0pt}%
\pgfpathmoveto{\pgfqpoint{1.037354in}{0.909387in}}%
\pgfpathlineto{\pgfqpoint{1.034467in}{0.925705in}}%
\pgfpathlineto{\pgfqpoint{1.031585in}{0.941994in}}%
\pgfpathlineto{\pgfqpoint{1.015010in}{0.939949in}}%
\pgfpathlineto{\pgfqpoint{0.998729in}{0.937475in}}%
\pgfpathlineto{\pgfqpoint{1.003232in}{0.921387in}}%
\pgfpathlineto{\pgfqpoint{1.007742in}{0.905272in}}%
\pgfpathlineto{\pgfqpoint{1.022414in}{0.907525in}}%
\pgfpathlineto{\pgfqpoint{1.037354in}{0.909387in}}%
\pgfpathclose%
\pgfusepath{fill}%
\end{pgfscope}%
\begin{pgfscope}%
\pgfpathrectangle{\pgfqpoint{0.407592in}{0.222904in}}{\pgfqpoint{1.411163in}{1.411163in}}%
\pgfusepath{clip}%
\pgfsetbuttcap%
\pgfsetroundjoin%
\definecolor{currentfill}{rgb}{0.749020,0.772549,0.792157}%
\pgfsetfillcolor{currentfill}%
\pgfsetfillopacity{0.160000}%
\pgfsetlinewidth{0.000000pt}%
\definecolor{currentstroke}{rgb}{0.000000,0.000000,0.000000}%
\pgfsetstrokecolor{currentstroke}%
\pgfsetdash{}{0pt}%
\pgfpathmoveto{\pgfqpoint{0.734820in}{1.212099in}}%
\pgfpathlineto{\pgfqpoint{0.726109in}{1.226844in}}%
\pgfpathlineto{\pgfqpoint{0.717406in}{1.241574in}}%
\pgfpathlineto{\pgfqpoint{0.690848in}{1.231938in}}%
\pgfpathlineto{\pgfqpoint{0.665782in}{1.221675in}}%
\pgfpathlineto{\pgfqpoint{0.675756in}{1.207315in}}%
\pgfpathlineto{\pgfqpoint{0.685736in}{1.192946in}}%
\pgfpathlineto{\pgfqpoint{0.709561in}{1.202822in}}%
\pgfpathlineto{\pgfqpoint{0.734820in}{1.212099in}}%
\pgfpathclose%
\pgfusepath{fill}%
\end{pgfscope}%
\begin{pgfscope}%
\pgfpathrectangle{\pgfqpoint{0.407592in}{0.222904in}}{\pgfqpoint{1.411163in}{1.411163in}}%
\pgfusepath{clip}%
\pgfsetbuttcap%
\pgfsetroundjoin%
\definecolor{currentfill}{rgb}{0.749020,0.772549,0.792157}%
\pgfsetfillcolor{currentfill}%
\pgfsetfillopacity{0.160000}%
\pgfsetlinewidth{0.000000pt}%
\definecolor{currentstroke}{rgb}{0.000000,0.000000,0.000000}%
\pgfsetstrokecolor{currentstroke}%
\pgfsetdash{}{0pt}%
\pgfpathmoveto{\pgfqpoint{1.358178in}{1.035391in}}%
\pgfpathlineto{\pgfqpoint{1.367610in}{1.050346in}}%
\pgfpathlineto{\pgfqpoint{1.377033in}{1.065287in}}%
\pgfpathlineto{\pgfqpoint{1.358494in}{1.072033in}}%
\pgfpathlineto{\pgfqpoint{1.338900in}{1.078263in}}%
\pgfpathlineto{\pgfqpoint{1.330763in}{1.062957in}}%
\pgfpathlineto{\pgfqpoint{1.322616in}{1.047633in}}%
\pgfpathlineto{\pgfqpoint{1.340894in}{1.041754in}}%
\pgfpathlineto{\pgfqpoint{1.358178in}{1.035391in}}%
\pgfpathclose%
\pgfusepath{fill}%
\end{pgfscope}%
\begin{pgfscope}%
\pgfpathrectangle{\pgfqpoint{0.407592in}{0.222904in}}{\pgfqpoint{1.411163in}{1.411163in}}%
\pgfusepath{clip}%
\pgfsetbuttcap%
\pgfsetroundjoin%
\definecolor{currentfill}{rgb}{0.749020,0.772549,0.792157}%
\pgfsetfillcolor{currentfill}%
\pgfsetfillopacity{0.160000}%
\pgfsetlinewidth{0.000000pt}%
\definecolor{currentstroke}{rgb}{0.000000,0.000000,0.000000}%
\pgfsetstrokecolor{currentstroke}%
\pgfsetdash{}{0pt}%
\pgfpathmoveto{\pgfqpoint{1.124944in}{0.879083in}}%
\pgfpathlineto{\pgfqpoint{1.127078in}{0.895596in}}%
\pgfpathlineto{\pgfqpoint{1.129209in}{0.912078in}}%
\pgfpathlineto{\pgfqpoint{1.113943in}{0.912651in}}%
\pgfpathlineto{\pgfqpoint{1.098578in}{0.912815in}}%
\pgfpathlineto{\pgfqpoint{1.098122in}{0.896297in}}%
\pgfpathlineto{\pgfqpoint{1.097665in}{0.879746in}}%
\pgfpathlineto{\pgfqpoint{1.111349in}{0.879599in}}%
\pgfpathlineto{\pgfqpoint{1.124944in}{0.879083in}}%
\pgfpathclose%
\pgfusepath{fill}%
\end{pgfscope}%
\begin{pgfscope}%
\pgfpathrectangle{\pgfqpoint{0.407592in}{0.222904in}}{\pgfqpoint{1.411163in}{1.411163in}}%
\pgfusepath{clip}%
\pgfsetbuttcap%
\pgfsetroundjoin%
\definecolor{currentfill}{rgb}{0.749020,0.772549,0.792157}%
\pgfsetfillcolor{currentfill}%
\pgfsetfillopacity{0.160000}%
\pgfsetlinewidth{0.000000pt}%
\definecolor{currentstroke}{rgb}{0.000000,0.000000,0.000000}%
\pgfsetstrokecolor{currentstroke}%
\pgfsetdash{}{0pt}%
\pgfpathmoveto{\pgfqpoint{1.097665in}{0.879746in}}%
\pgfpathlineto{\pgfqpoint{1.098122in}{0.896297in}}%
\pgfpathlineto{\pgfqpoint{1.098578in}{0.912815in}}%
\pgfpathlineto{\pgfqpoint{1.083175in}{0.912569in}}%
\pgfpathlineto{\pgfqpoint{1.067796in}{0.911914in}}%
\pgfpathlineto{\pgfqpoint{1.069022in}{0.895440in}}%
\pgfpathlineto{\pgfqpoint{1.070251in}{0.878935in}}%
\pgfpathlineto{\pgfqpoint{1.083947in}{0.879525in}}%
\pgfpathlineto{\pgfqpoint{1.097665in}{0.879746in}}%
\pgfpathclose%
\pgfusepath{fill}%
\end{pgfscope}%
\begin{pgfscope}%
\pgfpathrectangle{\pgfqpoint{0.407592in}{0.222904in}}{\pgfqpoint{1.411163in}{1.411163in}}%
\pgfusepath{clip}%
\pgfsetbuttcap%
\pgfsetroundjoin%
\definecolor{currentfill}{rgb}{0.749020,0.772549,0.792157}%
\pgfsetfillcolor{currentfill}%
\pgfsetfillopacity{0.160000}%
\pgfsetlinewidth{0.000000pt}%
\definecolor{currentstroke}{rgb}{0.000000,0.000000,0.000000}%
\pgfsetstrokecolor{currentstroke}%
\pgfsetdash{}{0pt}%
\pgfpathmoveto{\pgfqpoint{1.494977in}{1.166126in}}%
\pgfpathlineto{\pgfqpoint{1.505488in}{1.180555in}}%
\pgfpathlineto{\pgfqpoint{1.515992in}{1.194974in}}%
\pgfpathlineto{\pgfqpoint{1.494148in}{1.204732in}}%
\pgfpathlineto{\pgfqpoint{1.470751in}{1.213885in}}%
\pgfpathlineto{\pgfqpoint{1.461420in}{1.199091in}}%
\pgfpathlineto{\pgfqpoint{1.452081in}{1.184283in}}%
\pgfpathlineto{\pgfqpoint{1.474273in}{1.175494in}}%
\pgfpathlineto{\pgfqpoint{1.494977in}{1.166126in}}%
\pgfpathclose%
\pgfusepath{fill}%
\end{pgfscope}%
\begin{pgfscope}%
\pgfpathrectangle{\pgfqpoint{0.407592in}{0.222904in}}{\pgfqpoint{1.411163in}{1.411163in}}%
\pgfusepath{clip}%
\pgfsetbuttcap%
\pgfsetroundjoin%
\definecolor{currentfill}{rgb}{0.749020,0.772549,0.792157}%
\pgfsetfillcolor{currentfill}%
\pgfsetfillopacity{0.160000}%
\pgfsetlinewidth{0.000000pt}%
\definecolor{currentstroke}{rgb}{0.000000,0.000000,0.000000}%
\pgfsetstrokecolor{currentstroke}%
\pgfsetdash{}{0pt}%
\pgfpathmoveto{\pgfqpoint{0.877943in}{1.046493in}}%
\pgfpathlineto{\pgfqpoint{0.870497in}{1.061784in}}%
\pgfpathlineto{\pgfqpoint{0.863061in}{1.077055in}}%
\pgfpathlineto{\pgfqpoint{0.842249in}{1.070720in}}%
\pgfpathlineto{\pgfqpoint{0.822393in}{1.063874in}}%
\pgfpathlineto{\pgfqpoint{0.831196in}{1.048973in}}%
\pgfpathlineto{\pgfqpoint{0.840008in}{1.034058in}}%
\pgfpathlineto{\pgfqpoint{0.858525in}{1.040516in}}%
\pgfpathlineto{\pgfqpoint{0.877943in}{1.046493in}}%
\pgfpathclose%
\pgfusepath{fill}%
\end{pgfscope}%
\begin{pgfscope}%
\pgfpathrectangle{\pgfqpoint{0.407592in}{0.222904in}}{\pgfqpoint{1.411163in}{1.411163in}}%
\pgfusepath{clip}%
\pgfsetbuttcap%
\pgfsetroundjoin%
\definecolor{currentfill}{rgb}{0.749020,0.772549,0.792157}%
\pgfsetfillcolor{currentfill}%
\pgfsetfillopacity{0.160000}%
\pgfsetlinewidth{0.000000pt}%
\definecolor{currentstroke}{rgb}{0.000000,0.000000,0.000000}%
\pgfsetstrokecolor{currentstroke}%
\pgfsetdash{}{0pt}%
\pgfpathmoveto{\pgfqpoint{1.273514in}{0.955277in}}%
\pgfpathlineto{\pgfqpoint{1.281724in}{0.970718in}}%
\pgfpathlineto{\pgfqpoint{1.289923in}{0.986140in}}%
\pgfpathlineto{\pgfqpoint{1.273473in}{0.990849in}}%
\pgfpathlineto{\pgfqpoint{1.256291in}{0.995102in}}%
\pgfpathlineto{\pgfqpoint{1.249491in}{0.979355in}}%
\pgfpathlineto{\pgfqpoint{1.242681in}{0.963583in}}%
\pgfpathlineto{\pgfqpoint{1.258436in}{0.959641in}}%
\pgfpathlineto{\pgfqpoint{1.273514in}{0.955277in}}%
\pgfpathclose%
\pgfusepath{fill}%
\end{pgfscope}%
\begin{pgfscope}%
\pgfpathrectangle{\pgfqpoint{0.407592in}{0.222904in}}{\pgfqpoint{1.411163in}{1.411163in}}%
\pgfusepath{clip}%
\pgfsetbuttcap%
\pgfsetroundjoin%
\definecolor{currentfill}{rgb}{0.749020,0.772549,0.792157}%
\pgfsetfillcolor{currentfill}%
\pgfsetfillopacity{0.160000}%
\pgfsetlinewidth{0.000000pt}%
\definecolor{currentstroke}{rgb}{0.000000,0.000000,0.000000}%
\pgfsetstrokecolor{currentstroke}%
\pgfsetdash{}{0pt}%
\pgfpathmoveto{\pgfqpoint{1.151647in}{0.876956in}}%
\pgfpathlineto{\pgfqpoint{1.155426in}{0.893350in}}%
\pgfpathlineto{\pgfqpoint{1.159198in}{0.909713in}}%
\pgfpathlineto{\pgfqpoint{1.144314in}{0.911098in}}%
\pgfpathlineto{\pgfqpoint{1.129209in}{0.912078in}}%
\pgfpathlineto{\pgfqpoint{1.127078in}{0.895596in}}%
\pgfpathlineto{\pgfqpoint{1.124944in}{0.879083in}}%
\pgfpathlineto{\pgfqpoint{1.138395in}{0.878201in}}%
\pgfpathlineto{\pgfqpoint{1.151647in}{0.876956in}}%
\pgfpathclose%
\pgfusepath{fill}%
\end{pgfscope}%
\begin{pgfscope}%
\pgfpathrectangle{\pgfqpoint{0.407592in}{0.222904in}}{\pgfqpoint{1.411163in}{1.411163in}}%
\pgfusepath{clip}%
\pgfsetbuttcap%
\pgfsetroundjoin%
\definecolor{currentfill}{rgb}{0.749020,0.772549,0.792157}%
\pgfsetfillcolor{currentfill}%
\pgfsetfillopacity{0.160000}%
\pgfsetlinewidth{0.000000pt}%
\definecolor{currentstroke}{rgb}{0.000000,0.000000,0.000000}%
\pgfsetstrokecolor{currentstroke}%
\pgfsetdash{}{0pt}%
\pgfpathmoveto{\pgfqpoint{1.070251in}{0.878935in}}%
\pgfpathlineto{\pgfqpoint{1.069022in}{0.895440in}}%
\pgfpathlineto{\pgfqpoint{1.067796in}{0.911914in}}%
\pgfpathlineto{\pgfqpoint{1.052502in}{0.910852in}}%
\pgfpathlineto{\pgfqpoint{1.037354in}{0.909387in}}%
\pgfpathlineto{\pgfqpoint{1.040246in}{0.893040in}}%
\pgfpathlineto{\pgfqpoint{1.043144in}{0.876663in}}%
\pgfpathlineto{\pgfqpoint{1.056631in}{0.877980in}}%
\pgfpathlineto{\pgfqpoint{1.070251in}{0.878935in}}%
\pgfpathclose%
\pgfusepath{fill}%
\end{pgfscope}%
\begin{pgfscope}%
\pgfpathrectangle{\pgfqpoint{0.407592in}{0.222904in}}{\pgfqpoint{1.411163in}{1.411163in}}%
\pgfusepath{clip}%
\pgfsetbuttcap%
\pgfsetroundjoin%
\definecolor{currentfill}{rgb}{0.749020,0.772549,0.792157}%
\pgfsetfillcolor{currentfill}%
\pgfsetfillopacity{0.160000}%
\pgfsetlinewidth{0.000000pt}%
\definecolor{currentstroke}{rgb}{0.000000,0.000000,0.000000}%
\pgfsetstrokecolor{currentstroke}%
\pgfsetdash{}{0pt}%
\pgfpathmoveto{\pgfqpoint{0.955219in}{0.962827in}}%
\pgfpathlineto{\pgfqpoint{0.949186in}{0.978568in}}%
\pgfpathlineto{\pgfqpoint{0.943162in}{0.994286in}}%
\pgfpathlineto{\pgfqpoint{0.925167in}{0.989941in}}%
\pgfpathlineto{\pgfqpoint{0.907817in}{0.985142in}}%
\pgfpathlineto{\pgfqpoint{0.915309in}{0.969756in}}%
\pgfpathlineto{\pgfqpoint{0.922810in}{0.954352in}}%
\pgfpathlineto{\pgfqpoint{0.938715in}{0.958799in}}%
\pgfpathlineto{\pgfqpoint{0.955219in}{0.962827in}}%
\pgfpathclose%
\pgfusepath{fill}%
\end{pgfscope}%
\begin{pgfscope}%
\pgfpathrectangle{\pgfqpoint{0.407592in}{0.222904in}}{\pgfqpoint{1.411163in}{1.411163in}}%
\pgfusepath{clip}%
\pgfsetbuttcap%
\pgfsetroundjoin%
\definecolor{currentfill}{rgb}{0.749020,0.772549,0.792157}%
\pgfsetfillcolor{currentfill}%
\pgfsetfillopacity{0.160000}%
\pgfsetlinewidth{0.000000pt}%
\definecolor{currentstroke}{rgb}{0.000000,0.000000,0.000000}%
\pgfsetstrokecolor{currentstroke}%
\pgfsetdash{}{0pt}%
\pgfpathmoveto{\pgfqpoint{0.752268in}{1.182567in}}%
\pgfpathlineto{\pgfqpoint{0.743540in}{1.197340in}}%
\pgfpathlineto{\pgfqpoint{0.734820in}{1.212099in}}%
\pgfpathlineto{\pgfqpoint{0.709561in}{1.202822in}}%
\pgfpathlineto{\pgfqpoint{0.685736in}{1.192946in}}%
\pgfpathlineto{\pgfqpoint{0.695722in}{1.178567in}}%
\pgfpathlineto{\pgfqpoint{0.705714in}{1.164179in}}%
\pgfpathlineto{\pgfqpoint{0.728304in}{1.173660in}}%
\pgfpathlineto{\pgfqpoint{0.752268in}{1.182567in}}%
\pgfpathclose%
\pgfusepath{fill}%
\end{pgfscope}%
\begin{pgfscope}%
\pgfpathrectangle{\pgfqpoint{0.407592in}{0.222904in}}{\pgfqpoint{1.411163in}{1.411163in}}%
\pgfusepath{clip}%
\pgfsetbuttcap%
\pgfsetroundjoin%
\definecolor{currentfill}{rgb}{0.749020,0.772549,0.792157}%
\pgfsetfillcolor{currentfill}%
\pgfsetfillopacity{0.160000}%
\pgfsetlinewidth{0.000000pt}%
\definecolor{currentstroke}{rgb}{0.000000,0.000000,0.000000}%
\pgfsetstrokecolor{currentstroke}%
\pgfsetdash{}{0pt}%
\pgfpathmoveto{\pgfqpoint{1.473933in}{1.137239in}}%
\pgfpathlineto{\pgfqpoint{1.484458in}{1.151687in}}%
\pgfpathlineto{\pgfqpoint{1.494977in}{1.166126in}}%
\pgfpathlineto{\pgfqpoint{1.474273in}{1.175494in}}%
\pgfpathlineto{\pgfqpoint{1.452081in}{1.184283in}}%
\pgfpathlineto{\pgfqpoint{1.442732in}{1.169460in}}%
\pgfpathlineto{\pgfqpoint{1.433375in}{1.154622in}}%
\pgfpathlineto{\pgfqpoint{1.454364in}{1.146206in}}%
\pgfpathlineto{\pgfqpoint{1.473933in}{1.137239in}}%
\pgfpathclose%
\pgfusepath{fill}%
\end{pgfscope}%
\begin{pgfscope}%
\pgfpathrectangle{\pgfqpoint{0.407592in}{0.222904in}}{\pgfqpoint{1.411163in}{1.411163in}}%
\pgfusepath{clip}%
\pgfsetbuttcap%
\pgfsetroundjoin%
\definecolor{currentfill}{rgb}{0.749020,0.772549,0.792157}%
\pgfsetfillcolor{currentfill}%
\pgfsetfillopacity{0.160000}%
\pgfsetlinewidth{0.000000pt}%
\definecolor{currentstroke}{rgb}{0.000000,0.000000,0.000000}%
\pgfsetstrokecolor{currentstroke}%
\pgfsetdash{}{0pt}%
\pgfpathmoveto{\pgfqpoint{1.339286in}{1.005435in}}%
\pgfpathlineto{\pgfqpoint{1.348737in}{1.020420in}}%
\pgfpathlineto{\pgfqpoint{1.358178in}{1.035391in}}%
\pgfpathlineto{\pgfqpoint{1.340894in}{1.041754in}}%
\pgfpathlineto{\pgfqpoint{1.322616in}{1.047633in}}%
\pgfpathlineto{\pgfqpoint{1.314458in}{1.032289in}}%
\pgfpathlineto{\pgfqpoint{1.306290in}{1.016925in}}%
\pgfpathlineto{\pgfqpoint{1.323253in}{1.011407in}}%
\pgfpathlineto{\pgfqpoint{1.339286in}{1.005435in}}%
\pgfpathclose%
\pgfusepath{fill}%
\end{pgfscope}%
\begin{pgfscope}%
\pgfpathrectangle{\pgfqpoint{0.407592in}{0.222904in}}{\pgfqpoint{1.411163in}{1.411163in}}%
\pgfusepath{clip}%
\pgfsetbuttcap%
\pgfsetroundjoin%
\definecolor{currentfill}{rgb}{0.749020,0.772549,0.792157}%
\pgfsetfillcolor{currentfill}%
\pgfsetfillopacity{0.160000}%
\pgfsetlinewidth{0.000000pt}%
\definecolor{currentstroke}{rgb}{0.000000,0.000000,0.000000}%
\pgfsetstrokecolor{currentstroke}%
\pgfsetdash{}{0pt}%
\pgfpathmoveto{\pgfqpoint{1.215338in}{0.900262in}}%
\pgfpathlineto{\pgfqpoint{1.222189in}{0.916128in}}%
\pgfpathlineto{\pgfqpoint{1.229030in}{0.931970in}}%
\pgfpathlineto{\pgfqpoint{1.214141in}{0.935194in}}%
\pgfpathlineto{\pgfqpoint{1.198751in}{0.938005in}}%
\pgfpathlineto{\pgfqpoint{1.193412in}{0.921894in}}%
\pgfpathlineto{\pgfqpoint{1.188064in}{0.905755in}}%
\pgfpathlineto{\pgfqpoint{1.201929in}{0.903196in}}%
\pgfpathlineto{\pgfqpoint{1.215338in}{0.900262in}}%
\pgfpathclose%
\pgfusepath{fill}%
\end{pgfscope}%
\begin{pgfscope}%
\pgfpathrectangle{\pgfqpoint{0.407592in}{0.222904in}}{\pgfqpoint{1.411163in}{1.411163in}}%
\pgfusepath{clip}%
\pgfsetbuttcap%
\pgfsetroundjoin%
\definecolor{currentfill}{rgb}{0.749020,0.772549,0.792157}%
\pgfsetfillcolor{currentfill}%
\pgfsetfillopacity{0.160000}%
\pgfsetlinewidth{0.000000pt}%
\definecolor{currentstroke}{rgb}{0.000000,0.000000,0.000000}%
\pgfsetstrokecolor{currentstroke}%
\pgfsetdash{}{0pt}%
\pgfpathmoveto{\pgfqpoint{1.007742in}{0.905272in}}%
\pgfpathlineto{\pgfqpoint{1.003232in}{0.921387in}}%
\pgfpathlineto{\pgfqpoint{0.998729in}{0.937475in}}%
\pgfpathlineto{\pgfqpoint{0.982808in}{0.934580in}}%
\pgfpathlineto{\pgfqpoint{0.967310in}{0.931275in}}%
\pgfpathlineto{\pgfqpoint{0.973370in}{0.915464in}}%
\pgfpathlineto{\pgfqpoint{0.979438in}{0.899629in}}%
\pgfpathlineto{\pgfqpoint{0.993397in}{0.902637in}}%
\pgfpathlineto{\pgfqpoint{1.007742in}{0.905272in}}%
\pgfpathclose%
\pgfusepath{fill}%
\end{pgfscope}%
\begin{pgfscope}%
\pgfpathrectangle{\pgfqpoint{0.407592in}{0.222904in}}{\pgfqpoint{1.411163in}{1.411163in}}%
\pgfusepath{clip}%
\pgfsetbuttcap%
\pgfsetroundjoin%
\definecolor{currentfill}{rgb}{0.749020,0.772549,0.792157}%
\pgfsetfillcolor{currentfill}%
\pgfsetfillopacity{0.160000}%
\pgfsetlinewidth{0.000000pt}%
\definecolor{currentstroke}{rgb}{0.000000,0.000000,0.000000}%
\pgfsetstrokecolor{currentstroke}%
\pgfsetdash{}{0pt}%
\pgfpathmoveto{\pgfqpoint{1.177341in}{0.873396in}}%
\pgfpathlineto{\pgfqpoint{1.182707in}{0.889589in}}%
\pgfpathlineto{\pgfqpoint{1.188064in}{0.905755in}}%
\pgfpathlineto{\pgfqpoint{1.173801in}{0.907930in}}%
\pgfpathlineto{\pgfqpoint{1.159198in}{0.909713in}}%
\pgfpathlineto{\pgfqpoint{1.155426in}{0.893350in}}%
\pgfpathlineto{\pgfqpoint{1.151647in}{0.876956in}}%
\pgfpathlineto{\pgfqpoint{1.164646in}{0.875352in}}%
\pgfpathlineto{\pgfqpoint{1.177341in}{0.873396in}}%
\pgfpathclose%
\pgfusepath{fill}%
\end{pgfscope}%
\begin{pgfscope}%
\pgfpathrectangle{\pgfqpoint{0.407592in}{0.222904in}}{\pgfqpoint{1.411163in}{1.411163in}}%
\pgfusepath{clip}%
\pgfsetbuttcap%
\pgfsetroundjoin%
\definecolor{currentfill}{rgb}{0.749020,0.772549,0.792157}%
\pgfsetfillcolor{currentfill}%
\pgfsetfillopacity{0.160000}%
\pgfsetlinewidth{0.000000pt}%
\definecolor{currentstroke}{rgb}{0.000000,0.000000,0.000000}%
\pgfsetstrokecolor{currentstroke}%
\pgfsetdash{}{0pt}%
\pgfpathmoveto{\pgfqpoint{0.892861in}{1.015856in}}%
\pgfpathlineto{\pgfqpoint{0.885397in}{1.031184in}}%
\pgfpathlineto{\pgfqpoint{0.877943in}{1.046493in}}%
\pgfpathlineto{\pgfqpoint{0.858525in}{1.040516in}}%
\pgfpathlineto{\pgfqpoint{0.840008in}{1.034058in}}%
\pgfpathlineto{\pgfqpoint{0.848829in}{1.019128in}}%
\pgfpathlineto{\pgfqpoint{0.857658in}{1.004184in}}%
\pgfpathlineto{\pgfqpoint{0.874837in}{1.010245in}}%
\pgfpathlineto{\pgfqpoint{0.892861in}{1.015856in}}%
\pgfpathclose%
\pgfusepath{fill}%
\end{pgfscope}%
\begin{pgfscope}%
\pgfpathrectangle{\pgfqpoint{0.407592in}{0.222904in}}{\pgfqpoint{1.411163in}{1.411163in}}%
\pgfusepath{clip}%
\pgfsetbuttcap%
\pgfsetroundjoin%
\definecolor{currentfill}{rgb}{0.749020,0.772549,0.792157}%
\pgfsetfillcolor{currentfill}%
\pgfsetfillopacity{0.160000}%
\pgfsetlinewidth{0.000000pt}%
\definecolor{currentstroke}{rgb}{0.000000,0.000000,0.000000}%
\pgfsetstrokecolor{currentstroke}%
\pgfsetdash{}{0pt}%
\pgfpathmoveto{\pgfqpoint{1.043144in}{0.876663in}}%
\pgfpathlineto{\pgfqpoint{1.040246in}{0.893040in}}%
\pgfpathlineto{\pgfqpoint{1.037354in}{0.909387in}}%
\pgfpathlineto{\pgfqpoint{1.022414in}{0.907525in}}%
\pgfpathlineto{\pgfqpoint{1.007742in}{0.905272in}}%
\pgfpathlineto{\pgfqpoint{1.012259in}{0.889131in}}%
\pgfpathlineto{\pgfqpoint{1.016784in}{0.872962in}}%
\pgfpathlineto{\pgfqpoint{1.029844in}{0.874988in}}%
\pgfpathlineto{\pgfqpoint{1.043144in}{0.876663in}}%
\pgfpathclose%
\pgfusepath{fill}%
\end{pgfscope}%
\begin{pgfscope}%
\pgfpathrectangle{\pgfqpoint{0.407592in}{0.222904in}}{\pgfqpoint{1.411163in}{1.411163in}}%
\pgfusepath{clip}%
\pgfsetbuttcap%
\pgfsetroundjoin%
\definecolor{currentfill}{rgb}{0.749020,0.772549,0.792157}%
\pgfsetfillcolor{currentfill}%
\pgfsetfillopacity{0.160000}%
\pgfsetlinewidth{0.000000pt}%
\definecolor{currentstroke}{rgb}{0.000000,0.000000,0.000000}%
\pgfsetstrokecolor{currentstroke}%
\pgfsetdash{}{0pt}%
\pgfpathmoveto{\pgfqpoint{0.769749in}{1.152979in}}%
\pgfpathlineto{\pgfqpoint{0.761004in}{1.167780in}}%
\pgfpathlineto{\pgfqpoint{0.752268in}{1.182567in}}%
\pgfpathlineto{\pgfqpoint{0.728304in}{1.173660in}}%
\pgfpathlineto{\pgfqpoint{0.705714in}{1.164179in}}%
\pgfpathlineto{\pgfqpoint{0.715713in}{1.149782in}}%
\pgfpathlineto{\pgfqpoint{0.725719in}{1.135376in}}%
\pgfpathlineto{\pgfqpoint{0.747077in}{1.144450in}}%
\pgfpathlineto{\pgfqpoint{0.769749in}{1.152979in}}%
\pgfpathclose%
\pgfusepath{fill}%
\end{pgfscope}%
\begin{pgfscope}%
\pgfpathrectangle{\pgfqpoint{0.407592in}{0.222904in}}{\pgfqpoint{1.411163in}{1.411163in}}%
\pgfusepath{clip}%
\pgfsetbuttcap%
\pgfsetroundjoin%
\definecolor{currentfill}{rgb}{0.749020,0.772549,0.792157}%
\pgfsetfillcolor{currentfill}%
\pgfsetfillopacity{0.160000}%
\pgfsetlinewidth{0.000000pt}%
\definecolor{currentstroke}{rgb}{0.000000,0.000000,0.000000}%
\pgfsetstrokecolor{currentstroke}%
\pgfsetdash{}{0pt}%
\pgfpathmoveto{\pgfqpoint{1.257064in}{0.924335in}}%
\pgfpathlineto{\pgfqpoint{1.265294in}{0.939816in}}%
\pgfpathlineto{\pgfqpoint{1.273514in}{0.955277in}}%
\pgfpathlineto{\pgfqpoint{1.258436in}{0.959641in}}%
\pgfpathlineto{\pgfqpoint{1.242681in}{0.963583in}}%
\pgfpathlineto{\pgfqpoint{1.235861in}{0.947789in}}%
\pgfpathlineto{\pgfqpoint{1.229030in}{0.931970in}}%
\pgfpathlineto{\pgfqpoint{1.243358in}{0.928346in}}%
\pgfpathlineto{\pgfqpoint{1.257064in}{0.924335in}}%
\pgfpathclose%
\pgfusepath{fill}%
\end{pgfscope}%
\begin{pgfscope}%
\pgfpathrectangle{\pgfqpoint{0.407592in}{0.222904in}}{\pgfqpoint{1.411163in}{1.411163in}}%
\pgfusepath{clip}%
\pgfsetbuttcap%
\pgfsetroundjoin%
\definecolor{currentfill}{rgb}{0.749020,0.772549,0.792157}%
\pgfsetfillcolor{currentfill}%
\pgfsetfillopacity{0.160000}%
\pgfsetlinewidth{0.000000pt}%
\definecolor{currentstroke}{rgb}{0.000000,0.000000,0.000000}%
\pgfsetstrokecolor{currentstroke}%
\pgfsetdash{}{0pt}%
\pgfpathmoveto{\pgfqpoint{1.452860in}{1.108313in}}%
\pgfpathlineto{\pgfqpoint{1.463400in}{1.122781in}}%
\pgfpathlineto{\pgfqpoint{1.473933in}{1.137239in}}%
\pgfpathlineto{\pgfqpoint{1.454364in}{1.146206in}}%
\pgfpathlineto{\pgfqpoint{1.433375in}{1.154622in}}%
\pgfpathlineto{\pgfqpoint{1.424008in}{1.139770in}}%
\pgfpathlineto{\pgfqpoint{1.414631in}{1.124902in}}%
\pgfpathlineto{\pgfqpoint{1.434422in}{1.116870in}}%
\pgfpathlineto{\pgfqpoint{1.452860in}{1.108313in}}%
\pgfpathclose%
\pgfusepath{fill}%
\end{pgfscope}%
\begin{pgfscope}%
\pgfpathrectangle{\pgfqpoint{0.407592in}{0.222904in}}{\pgfqpoint{1.411163in}{1.411163in}}%
\pgfusepath{clip}%
\pgfsetbuttcap%
\pgfsetroundjoin%
\definecolor{currentfill}{rgb}{0.749020,0.772549,0.792157}%
\pgfsetfillcolor{currentfill}%
\pgfsetfillopacity{0.160000}%
\pgfsetlinewidth{0.000000pt}%
\definecolor{currentstroke}{rgb}{0.000000,0.000000,0.000000}%
\pgfsetstrokecolor{currentstroke}%
\pgfsetdash{}{0pt}%
\pgfpathmoveto{\pgfqpoint{1.120662in}{0.845962in}}%
\pgfpathlineto{\pgfqpoint{1.122805in}{0.862538in}}%
\pgfpathlineto{\pgfqpoint{1.124944in}{0.879083in}}%
\pgfpathlineto{\pgfqpoint{1.111349in}{0.879599in}}%
\pgfpathlineto{\pgfqpoint{1.097665in}{0.879746in}}%
\pgfpathlineto{\pgfqpoint{1.097207in}{0.863164in}}%
\pgfpathlineto{\pgfqpoint{1.096748in}{0.846549in}}%
\pgfpathlineto{\pgfqpoint{1.108745in}{0.846419in}}%
\pgfpathlineto{\pgfqpoint{1.120662in}{0.845962in}}%
\pgfpathclose%
\pgfusepath{fill}%
\end{pgfscope}%
\begin{pgfscope}%
\pgfpathrectangle{\pgfqpoint{0.407592in}{0.222904in}}{\pgfqpoint{1.411163in}{1.411163in}}%
\pgfusepath{clip}%
\pgfsetbuttcap%
\pgfsetroundjoin%
\definecolor{currentfill}{rgb}{0.749020,0.772549,0.792157}%
\pgfsetfillcolor{currentfill}%
\pgfsetfillopacity{0.160000}%
\pgfsetlinewidth{0.000000pt}%
\definecolor{currentstroke}{rgb}{0.000000,0.000000,0.000000}%
\pgfsetstrokecolor{currentstroke}%
\pgfsetdash{}{0pt}%
\pgfpathmoveto{\pgfqpoint{1.096748in}{0.846549in}}%
\pgfpathlineto{\pgfqpoint{1.097207in}{0.863164in}}%
\pgfpathlineto{\pgfqpoint{1.097665in}{0.879746in}}%
\pgfpathlineto{\pgfqpoint{1.083947in}{0.879525in}}%
\pgfpathlineto{\pgfqpoint{1.070251in}{0.878935in}}%
\pgfpathlineto{\pgfqpoint{1.071482in}{0.862399in}}%
\pgfpathlineto{\pgfqpoint{1.072715in}{0.845831in}}%
\pgfpathlineto{\pgfqpoint{1.084722in}{0.846353in}}%
\pgfpathlineto{\pgfqpoint{1.096748in}{0.846549in}}%
\pgfpathclose%
\pgfusepath{fill}%
\end{pgfscope}%
\begin{pgfscope}%
\pgfpathrectangle{\pgfqpoint{0.407592in}{0.222904in}}{\pgfqpoint{1.411163in}{1.411163in}}%
\pgfusepath{clip}%
\pgfsetbuttcap%
\pgfsetroundjoin%
\definecolor{currentfill}{rgb}{0.749020,0.772549,0.792157}%
\pgfsetfillcolor{currentfill}%
\pgfsetfillopacity{0.160000}%
\pgfsetlinewidth{0.000000pt}%
\definecolor{currentstroke}{rgb}{0.000000,0.000000,0.000000}%
\pgfsetstrokecolor{currentstroke}%
\pgfsetdash{}{0pt}%
\pgfpathmoveto{\pgfqpoint{0.967310in}{0.931275in}}%
\pgfpathlineto{\pgfqpoint{0.961260in}{0.947063in}}%
\pgfpathlineto{\pgfqpoint{0.955219in}{0.962827in}}%
\pgfpathlineto{\pgfqpoint{0.938715in}{0.958799in}}%
\pgfpathlineto{\pgfqpoint{0.922810in}{0.954352in}}%
\pgfpathlineto{\pgfqpoint{0.930321in}{0.938928in}}%
\pgfpathlineto{\pgfqpoint{0.937841in}{0.923485in}}%
\pgfpathlineto{\pgfqpoint{0.952301in}{0.927572in}}%
\pgfpathlineto{\pgfqpoint{0.967310in}{0.931275in}}%
\pgfpathclose%
\pgfusepath{fill}%
\end{pgfscope}%
\begin{pgfscope}%
\pgfpathrectangle{\pgfqpoint{0.407592in}{0.222904in}}{\pgfqpoint{1.411163in}{1.411163in}}%
\pgfusepath{clip}%
\pgfsetbuttcap%
\pgfsetroundjoin%
\definecolor{currentfill}{rgb}{0.749020,0.772549,0.792157}%
\pgfsetfillcolor{currentfill}%
\pgfsetfillopacity{0.160000}%
\pgfsetlinewidth{0.000000pt}%
\definecolor{currentstroke}{rgb}{0.000000,0.000000,0.000000}%
\pgfsetstrokecolor{currentstroke}%
\pgfsetdash{}{0pt}%
\pgfpathmoveto{\pgfqpoint{1.320356in}{0.975420in}}%
\pgfpathlineto{\pgfqpoint{1.329826in}{0.990435in}}%
\pgfpathlineto{\pgfqpoint{1.339286in}{1.005435in}}%
\pgfpathlineto{\pgfqpoint{1.323253in}{1.011407in}}%
\pgfpathlineto{\pgfqpoint{1.306290in}{1.016925in}}%
\pgfpathlineto{\pgfqpoint{1.298111in}{1.001543in}}%
\pgfpathlineto{\pgfqpoint{1.289923in}{0.986140in}}%
\pgfpathlineto{\pgfqpoint{1.305573in}{0.980991in}}%
\pgfpathlineto{\pgfqpoint{1.320356in}{0.975420in}}%
\pgfpathclose%
\pgfusepath{fill}%
\end{pgfscope}%
\begin{pgfscope}%
\pgfpathrectangle{\pgfqpoint{0.407592in}{0.222904in}}{\pgfqpoint{1.411163in}{1.411163in}}%
\pgfusepath{clip}%
\pgfsetbuttcap%
\pgfsetroundjoin%
\definecolor{currentfill}{rgb}{0.749020,0.772549,0.792157}%
\pgfsetfillcolor{currentfill}%
\pgfsetfillopacity{0.160000}%
\pgfsetlinewidth{0.000000pt}%
\definecolor{currentstroke}{rgb}{0.000000,0.000000,0.000000}%
\pgfsetstrokecolor{currentstroke}%
\pgfsetdash{}{0pt}%
\pgfpathmoveto{\pgfqpoint{1.738704in}{1.397027in}}%
\pgfpathlineto{\pgfqpoint{1.750178in}{1.410942in}}%
\pgfpathlineto{\pgfqpoint{1.723009in}{1.425155in}}%
\pgfpathlineto{\pgfqpoint{1.693489in}{1.438622in}}%
\pgfpathlineto{\pgfqpoint{1.683104in}{1.424367in}}%
\pgfpathlineto{\pgfqpoint{1.712063in}{1.411064in}}%
\pgfusepath{fill}%
\end{pgfscope}%
\begin{pgfscope}%
\pgfpathrectangle{\pgfqpoint{0.407592in}{0.222904in}}{\pgfqpoint{1.411163in}{1.411163in}}%
\pgfusepath{clip}%
\pgfsetbuttcap%
\pgfsetroundjoin%
\definecolor{currentfill}{rgb}{0.749020,0.772549,0.792157}%
\pgfsetfillcolor{currentfill}%
\pgfsetfillopacity{0.160000}%
\pgfsetlinewidth{0.000000pt}%
\definecolor{currentstroke}{rgb}{0.000000,0.000000,0.000000}%
\pgfsetstrokecolor{currentstroke}%
\pgfsetdash{}{0pt}%
\pgfpathmoveto{\pgfqpoint{0.787263in}{1.123334in}}%
\pgfpathlineto{\pgfqpoint{0.778502in}{1.138164in}}%
\pgfpathlineto{\pgfqpoint{0.769749in}{1.152979in}}%
\pgfpathlineto{\pgfqpoint{0.747077in}{1.144450in}}%
\pgfpathlineto{\pgfqpoint{0.725719in}{1.135376in}}%
\pgfpathlineto{\pgfqpoint{0.735730in}{1.120961in}}%
\pgfpathlineto{\pgfqpoint{0.745749in}{1.106536in}}%
\pgfpathlineto{\pgfqpoint{0.765880in}{1.115194in}}%
\pgfpathlineto{\pgfqpoint{0.787263in}{1.123334in}}%
\pgfpathclose%
\pgfusepath{fill}%
\end{pgfscope}%
\begin{pgfscope}%
\pgfpathrectangle{\pgfqpoint{0.407592in}{0.222904in}}{\pgfqpoint{1.411163in}{1.411163in}}%
\pgfusepath{clip}%
\pgfsetbuttcap%
\pgfsetroundjoin%
\definecolor{currentfill}{rgb}{0.749020,0.772549,0.792157}%
\pgfsetfillcolor{currentfill}%
\pgfsetfillopacity{0.160000}%
\pgfsetlinewidth{0.000000pt}%
\definecolor{currentstroke}{rgb}{0.000000,0.000000,0.000000}%
\pgfsetstrokecolor{currentstroke}%
\pgfsetdash{}{0pt}%
\pgfpathmoveto{\pgfqpoint{1.144068in}{0.844079in}}%
\pgfpathlineto{\pgfqpoint{1.147861in}{0.860532in}}%
\pgfpathlineto{\pgfqpoint{1.151647in}{0.876956in}}%
\pgfpathlineto{\pgfqpoint{1.138395in}{0.878201in}}%
\pgfpathlineto{\pgfqpoint{1.124944in}{0.879083in}}%
\pgfpathlineto{\pgfqpoint{1.122805in}{0.862538in}}%
\pgfpathlineto{\pgfqpoint{1.120662in}{0.845962in}}%
\pgfpathlineto{\pgfqpoint{1.132453in}{0.845181in}}%
\pgfpathlineto{\pgfqpoint{1.144068in}{0.844079in}}%
\pgfpathclose%
\pgfusepath{fill}%
\end{pgfscope}%
\begin{pgfscope}%
\pgfpathrectangle{\pgfqpoint{0.407592in}{0.222904in}}{\pgfqpoint{1.411163in}{1.411163in}}%
\pgfusepath{clip}%
\pgfsetbuttcap%
\pgfsetroundjoin%
\definecolor{currentfill}{rgb}{0.749020,0.772549,0.792157}%
\pgfsetfillcolor{currentfill}%
\pgfsetfillopacity{0.160000}%
\pgfsetlinewidth{0.000000pt}%
\definecolor{currentstroke}{rgb}{0.000000,0.000000,0.000000}%
\pgfsetstrokecolor{currentstroke}%
\pgfsetdash{}{0pt}%
\pgfpathmoveto{\pgfqpoint{1.072715in}{0.845831in}}%
\pgfpathlineto{\pgfqpoint{1.071482in}{0.862399in}}%
\pgfpathlineto{\pgfqpoint{1.070251in}{0.878935in}}%
\pgfpathlineto{\pgfqpoint{1.056631in}{0.877980in}}%
\pgfpathlineto{\pgfqpoint{1.043144in}{0.876663in}}%
\pgfpathlineto{\pgfqpoint{1.046047in}{0.860256in}}%
\pgfpathlineto{\pgfqpoint{1.048955in}{0.843819in}}%
\pgfpathlineto{\pgfqpoint{1.060776in}{0.844986in}}%
\pgfpathlineto{\pgfqpoint{1.072715in}{0.845831in}}%
\pgfpathclose%
\pgfusepath{fill}%
\end{pgfscope}%
\begin{pgfscope}%
\pgfpathrectangle{\pgfqpoint{0.407592in}{0.222904in}}{\pgfqpoint{1.411163in}{1.411163in}}%
\pgfusepath{clip}%
\pgfsetbuttcap%
\pgfsetroundjoin%
\definecolor{currentfill}{rgb}{0.749020,0.772549,0.792157}%
\pgfsetfillcolor{currentfill}%
\pgfsetfillopacity{0.160000}%
\pgfsetlinewidth{0.000000pt}%
\definecolor{currentstroke}{rgb}{0.000000,0.000000,0.000000}%
\pgfsetstrokecolor{currentstroke}%
\pgfsetdash{}{0pt}%
\pgfpathmoveto{\pgfqpoint{0.907817in}{0.985142in}}%
\pgfpathlineto{\pgfqpoint{0.900335in}{1.000508in}}%
\pgfpathlineto{\pgfqpoint{0.892861in}{1.015856in}}%
\pgfpathlineto{\pgfqpoint{0.874837in}{1.010245in}}%
\pgfpathlineto{\pgfqpoint{0.857658in}{1.004184in}}%
\pgfpathlineto{\pgfqpoint{0.866495in}{0.989226in}}%
\pgfpathlineto{\pgfqpoint{0.875341in}{0.974253in}}%
\pgfpathlineto{\pgfqpoint{0.891185in}{0.979906in}}%
\pgfpathlineto{\pgfqpoint{0.907817in}{0.985142in}}%
\pgfpathclose%
\pgfusepath{fill}%
\end{pgfscope}%
\begin{pgfscope}%
\pgfpathrectangle{\pgfqpoint{0.407592in}{0.222904in}}{\pgfqpoint{1.411163in}{1.411163in}}%
\pgfusepath{clip}%
\pgfsetbuttcap%
\pgfsetroundjoin%
\definecolor{currentfill}{rgb}{0.749020,0.772549,0.792157}%
\pgfsetfillcolor{currentfill}%
\pgfsetfillopacity{0.160000}%
\pgfsetlinewidth{0.000000pt}%
\definecolor{currentstroke}{rgb}{0.000000,0.000000,0.000000}%
\pgfsetstrokecolor{currentstroke}%
\pgfsetdash{}{0pt}%
\pgfpathmoveto{\pgfqpoint{1.431759in}{1.079348in}}%
\pgfpathlineto{\pgfqpoint{1.442313in}{1.093835in}}%
\pgfpathlineto{\pgfqpoint{1.452860in}{1.108313in}}%
\pgfpathlineto{\pgfqpoint{1.434422in}{1.116870in}}%
\pgfpathlineto{\pgfqpoint{1.414631in}{1.124902in}}%
\pgfpathlineto{\pgfqpoint{1.405246in}{1.110021in}}%
\pgfpathlineto{\pgfqpoint{1.395851in}{1.095124in}}%
\pgfpathlineto{\pgfqpoint{1.414446in}{1.087484in}}%
\pgfpathlineto{\pgfqpoint{1.431759in}{1.079348in}}%
\pgfpathclose%
\pgfusepath{fill}%
\end{pgfscope}%
\begin{pgfscope}%
\pgfpathrectangle{\pgfqpoint{0.407592in}{0.222904in}}{\pgfqpoint{1.411163in}{1.411163in}}%
\pgfusepath{clip}%
\pgfsetbuttcap%
\pgfsetroundjoin%
\definecolor{currentfill}{rgb}{0.749020,0.772549,0.792157}%
\pgfsetfillcolor{currentfill}%
\pgfsetfillopacity{0.160000}%
\pgfsetlinewidth{0.000000pt}%
\definecolor{currentstroke}{rgb}{0.000000,0.000000,0.000000}%
\pgfsetstrokecolor{currentstroke}%
\pgfsetdash{}{0pt}%
\pgfpathmoveto{\pgfqpoint{1.715746in}{1.369183in}}%
\pgfpathlineto{\pgfqpoint{1.727227in}{1.383107in}}%
\pgfpathlineto{\pgfqpoint{1.738704in}{1.397027in}}%
\pgfpathlineto{\pgfqpoint{1.712063in}{1.411064in}}%
\pgfpathlineto{\pgfqpoint{1.683104in}{1.424367in}}%
\pgfpathlineto{\pgfqpoint{1.672712in}{1.410102in}}%
\pgfpathlineto{\pgfqpoint{1.662313in}{1.395827in}}%
\pgfpathlineto{\pgfqpoint{1.690155in}{1.382860in}}%
\pgfpathlineto{\pgfqpoint{1.715746in}{1.369183in}}%
\pgfpathclose%
\pgfusepath{fill}%
\end{pgfscope}%
\begin{pgfscope}%
\pgfpathrectangle{\pgfqpoint{0.407592in}{0.222904in}}{\pgfqpoint{1.411163in}{1.411163in}}%
\pgfusepath{clip}%
\pgfsetbuttcap%
\pgfsetroundjoin%
\definecolor{currentfill}{rgb}{0.749020,0.772549,0.792157}%
\pgfsetfillcolor{currentfill}%
\pgfsetfillopacity{0.160000}%
\pgfsetlinewidth{0.000000pt}%
\definecolor{currentstroke}{rgb}{0.000000,0.000000,0.000000}%
\pgfsetstrokecolor{currentstroke}%
\pgfsetdash{}{0pt}%
\pgfpathmoveto{\pgfqpoint{1.201605in}{0.868458in}}%
\pgfpathlineto{\pgfqpoint{1.208477in}{0.884372in}}%
\pgfpathlineto{\pgfqpoint{1.215338in}{0.900262in}}%
\pgfpathlineto{\pgfqpoint{1.201929in}{0.903196in}}%
\pgfpathlineto{\pgfqpoint{1.188064in}{0.905755in}}%
\pgfpathlineto{\pgfqpoint{1.182707in}{0.889589in}}%
\pgfpathlineto{\pgfqpoint{1.177341in}{0.873396in}}%
\pgfpathlineto{\pgfqpoint{1.189677in}{0.871095in}}%
\pgfpathlineto{\pgfqpoint{1.201605in}{0.868458in}}%
\pgfpathclose%
\pgfusepath{fill}%
\end{pgfscope}%
\begin{pgfscope}%
\pgfpathrectangle{\pgfqpoint{0.407592in}{0.222904in}}{\pgfqpoint{1.411163in}{1.411163in}}%
\pgfusepath{clip}%
\pgfsetbuttcap%
\pgfsetroundjoin%
\definecolor{currentfill}{rgb}{0.749020,0.772549,0.792157}%
\pgfsetfillcolor{currentfill}%
\pgfsetfillopacity{0.160000}%
\pgfsetlinewidth{0.000000pt}%
\definecolor{currentstroke}{rgb}{0.000000,0.000000,0.000000}%
\pgfsetstrokecolor{currentstroke}%
\pgfsetdash{}{0pt}%
\pgfpathmoveto{\pgfqpoint{1.016784in}{0.872962in}}%
\pgfpathlineto{\pgfqpoint{1.012259in}{0.889131in}}%
\pgfpathlineto{\pgfqpoint{1.007742in}{0.905272in}}%
\pgfpathlineto{\pgfqpoint{0.993397in}{0.902637in}}%
\pgfpathlineto{\pgfqpoint{0.979438in}{0.899629in}}%
\pgfpathlineto{\pgfqpoint{0.985516in}{0.883771in}}%
\pgfpathlineto{\pgfqpoint{0.991602in}{0.867889in}}%
\pgfpathlineto{\pgfqpoint{1.004020in}{0.870593in}}%
\pgfpathlineto{\pgfqpoint{1.016784in}{0.872962in}}%
\pgfpathclose%
\pgfusepath{fill}%
\end{pgfscope}%
\begin{pgfscope}%
\pgfpathrectangle{\pgfqpoint{0.407592in}{0.222904in}}{\pgfqpoint{1.411163in}{1.411163in}}%
\pgfusepath{clip}%
\pgfsetbuttcap%
\pgfsetroundjoin%
\definecolor{currentfill}{rgb}{0.749020,0.772549,0.792157}%
\pgfsetfillcolor{currentfill}%
\pgfsetfillopacity{0.160000}%
\pgfsetlinewidth{0.000000pt}%
\definecolor{currentstroke}{rgb}{0.000000,0.000000,0.000000}%
\pgfsetstrokecolor{currentstroke}%
\pgfsetdash{}{0pt}%
\pgfpathmoveto{\pgfqpoint{0.804811in}{1.093633in}}%
\pgfpathlineto{\pgfqpoint{0.796033in}{1.108491in}}%
\pgfpathlineto{\pgfqpoint{0.787263in}{1.123334in}}%
\pgfpathlineto{\pgfqpoint{0.765880in}{1.115194in}}%
\pgfpathlineto{\pgfqpoint{0.745749in}{1.106536in}}%
\pgfpathlineto{\pgfqpoint{0.755773in}{1.092102in}}%
\pgfpathlineto{\pgfqpoint{0.765804in}{1.077659in}}%
\pgfpathlineto{\pgfqpoint{0.784714in}{1.085891in}}%
\pgfpathlineto{\pgfqpoint{0.804811in}{1.093633in}}%
\pgfpathclose%
\pgfusepath{fill}%
\end{pgfscope}%
\begin{pgfscope}%
\pgfpathrectangle{\pgfqpoint{0.407592in}{0.222904in}}{\pgfqpoint{1.411163in}{1.411163in}}%
\pgfusepath{clip}%
\pgfsetbuttcap%
\pgfsetroundjoin%
\definecolor{currentfill}{rgb}{0.749020,0.772549,0.792157}%
\pgfsetfillcolor{currentfill}%
\pgfsetfillopacity{0.160000}%
\pgfsetlinewidth{0.000000pt}%
\definecolor{currentstroke}{rgb}{0.000000,0.000000,0.000000}%
\pgfsetstrokecolor{currentstroke}%
\pgfsetdash{}{0pt}%
\pgfpathmoveto{\pgfqpoint{1.166581in}{0.840928in}}%
\pgfpathlineto{\pgfqpoint{1.171965in}{0.857176in}}%
\pgfpathlineto{\pgfqpoint{1.177341in}{0.873396in}}%
\pgfpathlineto{\pgfqpoint{1.164646in}{0.875352in}}%
\pgfpathlineto{\pgfqpoint{1.151647in}{0.876956in}}%
\pgfpathlineto{\pgfqpoint{1.147861in}{0.860532in}}%
\pgfpathlineto{\pgfqpoint{1.144068in}{0.844079in}}%
\pgfpathlineto{\pgfqpoint{1.155459in}{0.842659in}}%
\pgfpathlineto{\pgfqpoint{1.166581in}{0.840928in}}%
\pgfpathclose%
\pgfusepath{fill}%
\end{pgfscope}%
\begin{pgfscope}%
\pgfpathrectangle{\pgfqpoint{0.407592in}{0.222904in}}{\pgfqpoint{1.411163in}{1.411163in}}%
\pgfusepath{clip}%
\pgfsetbuttcap%
\pgfsetroundjoin%
\definecolor{currentfill}{rgb}{0.749020,0.772549,0.792157}%
\pgfsetfillcolor{currentfill}%
\pgfsetfillopacity{0.160000}%
\pgfsetlinewidth{0.000000pt}%
\definecolor{currentstroke}{rgb}{0.000000,0.000000,0.000000}%
\pgfsetstrokecolor{currentstroke}%
\pgfsetdash{}{0pt}%
\pgfpathmoveto{\pgfqpoint{1.240571in}{0.893314in}}%
\pgfpathlineto{\pgfqpoint{1.248823in}{0.908834in}}%
\pgfpathlineto{\pgfqpoint{1.257064in}{0.924335in}}%
\pgfpathlineto{\pgfqpoint{1.243358in}{0.928346in}}%
\pgfpathlineto{\pgfqpoint{1.229030in}{0.931970in}}%
\pgfpathlineto{\pgfqpoint{1.222189in}{0.916128in}}%
\pgfpathlineto{\pgfqpoint{1.215338in}{0.900262in}}%
\pgfpathlineto{\pgfqpoint{1.228237in}{0.896964in}}%
\pgfpathlineto{\pgfqpoint{1.240571in}{0.893314in}}%
\pgfpathclose%
\pgfusepath{fill}%
\end{pgfscope}%
\begin{pgfscope}%
\pgfpathrectangle{\pgfqpoint{0.407592in}{0.222904in}}{\pgfqpoint{1.411163in}{1.411163in}}%
\pgfusepath{clip}%
\pgfsetbuttcap%
\pgfsetroundjoin%
\definecolor{currentfill}{rgb}{0.749020,0.772549,0.792157}%
\pgfsetfillcolor{currentfill}%
\pgfsetfillopacity{0.160000}%
\pgfsetlinewidth{0.000000pt}%
\definecolor{currentstroke}{rgb}{0.000000,0.000000,0.000000}%
\pgfsetstrokecolor{currentstroke}%
\pgfsetdash{}{0pt}%
\pgfpathmoveto{\pgfqpoint{1.692774in}{1.341320in}}%
\pgfpathlineto{\pgfqpoint{1.704262in}{1.355254in}}%
\pgfpathlineto{\pgfqpoint{1.715746in}{1.369183in}}%
\pgfpathlineto{\pgfqpoint{1.690155in}{1.382860in}}%
\pgfpathlineto{\pgfqpoint{1.662313in}{1.395827in}}%
\pgfpathlineto{\pgfqpoint{1.651907in}{1.381543in}}%
\pgfpathlineto{\pgfqpoint{1.641494in}{1.367249in}}%
\pgfpathlineto{\pgfqpoint{1.668225in}{1.354628in}}%
\pgfpathlineto{\pgfqpoint{1.692774in}{1.341320in}}%
\pgfpathclose%
\pgfusepath{fill}%
\end{pgfscope}%
\begin{pgfscope}%
\pgfpathrectangle{\pgfqpoint{0.407592in}{0.222904in}}{\pgfqpoint{1.411163in}{1.411163in}}%
\pgfusepath{clip}%
\pgfsetbuttcap%
\pgfsetroundjoin%
\definecolor{currentfill}{rgb}{0.749020,0.772549,0.792157}%
\pgfsetfillcolor{currentfill}%
\pgfsetfillopacity{0.160000}%
\pgfsetlinewidth{0.000000pt}%
\definecolor{currentstroke}{rgb}{0.000000,0.000000,0.000000}%
\pgfsetstrokecolor{currentstroke}%
\pgfsetdash{}{0pt}%
\pgfpathmoveto{\pgfqpoint{1.048955in}{0.843819in}}%
\pgfpathlineto{\pgfqpoint{1.046047in}{0.860256in}}%
\pgfpathlineto{\pgfqpoint{1.043144in}{0.876663in}}%
\pgfpathlineto{\pgfqpoint{1.029844in}{0.874988in}}%
\pgfpathlineto{\pgfqpoint{1.016784in}{0.872962in}}%
\pgfpathlineto{\pgfqpoint{1.021317in}{0.856767in}}%
\pgfpathlineto{\pgfqpoint{1.025857in}{0.840544in}}%
\pgfpathlineto{\pgfqpoint{1.037299in}{0.842337in}}%
\pgfpathlineto{\pgfqpoint{1.048955in}{0.843819in}}%
\pgfpathclose%
\pgfusepath{fill}%
\end{pgfscope}%
\begin{pgfscope}%
\pgfpathrectangle{\pgfqpoint{0.407592in}{0.222904in}}{\pgfqpoint{1.411163in}{1.411163in}}%
\pgfusepath{clip}%
\pgfsetbuttcap%
\pgfsetroundjoin%
\definecolor{currentfill}{rgb}{0.749020,0.772549,0.792157}%
\pgfsetfillcolor{currentfill}%
\pgfsetfillopacity{0.160000}%
\pgfsetlinewidth{0.000000pt}%
\definecolor{currentstroke}{rgb}{0.000000,0.000000,0.000000}%
\pgfsetstrokecolor{currentstroke}%
\pgfsetdash{}{0pt}%
\pgfpathmoveto{\pgfqpoint{0.526822in}{1.421759in}}%
\pgfpathlineto{\pgfqpoint{0.516943in}{1.435982in}}%
\pgfpathlineto{\pgfqpoint{0.484477in}{1.422362in}}%
\pgfpathlineto{\pgfqpoint{0.454204in}{1.408005in}}%
\pgfpathlineto{\pgfqpoint{0.465273in}{1.394126in}}%
\pgfpathlineto{\pgfqpoint{0.494965in}{1.408305in}}%
\pgfusepath{fill}%
\end{pgfscope}%
\begin{pgfscope}%
\pgfpathrectangle{\pgfqpoint{0.407592in}{0.222904in}}{\pgfqpoint{1.411163in}{1.411163in}}%
\pgfusepath{clip}%
\pgfsetbuttcap%
\pgfsetroundjoin%
\definecolor{currentfill}{rgb}{0.749020,0.772549,0.792157}%
\pgfsetfillcolor{currentfill}%
\pgfsetfillopacity{0.160000}%
\pgfsetlinewidth{0.000000pt}%
\definecolor{currentstroke}{rgb}{0.000000,0.000000,0.000000}%
\pgfsetstrokecolor{currentstroke}%
\pgfsetdash{}{0pt}%
\pgfpathmoveto{\pgfqpoint{1.301388in}{0.945345in}}%
\pgfpathlineto{\pgfqpoint{1.310877in}{0.960390in}}%
\pgfpathlineto{\pgfqpoint{1.320356in}{0.975420in}}%
\pgfpathlineto{\pgfqpoint{1.305573in}{0.980991in}}%
\pgfpathlineto{\pgfqpoint{1.289923in}{0.986140in}}%
\pgfpathlineto{\pgfqpoint{1.281724in}{0.970718in}}%
\pgfpathlineto{\pgfqpoint{1.273514in}{0.955277in}}%
\pgfpathlineto{\pgfqpoint{1.287852in}{0.950506in}}%
\pgfpathlineto{\pgfqpoint{1.301388in}{0.945345in}}%
\pgfpathclose%
\pgfusepath{fill}%
\end{pgfscope}%
\begin{pgfscope}%
\pgfpathrectangle{\pgfqpoint{0.407592in}{0.222904in}}{\pgfqpoint{1.411163in}{1.411163in}}%
\pgfusepath{clip}%
\pgfsetbuttcap%
\pgfsetroundjoin%
\definecolor{currentfill}{rgb}{0.749020,0.772549,0.792157}%
\pgfsetfillcolor{currentfill}%
\pgfsetfillopacity{0.160000}%
\pgfsetlinewidth{0.000000pt}%
\definecolor{currentstroke}{rgb}{0.000000,0.000000,0.000000}%
\pgfsetstrokecolor{currentstroke}%
\pgfsetdash{}{0pt}%
\pgfpathmoveto{\pgfqpoint{0.979438in}{0.899629in}}%
\pgfpathlineto{\pgfqpoint{0.973370in}{0.915464in}}%
\pgfpathlineto{\pgfqpoint{0.967310in}{0.931275in}}%
\pgfpathlineto{\pgfqpoint{0.952301in}{0.927572in}}%
\pgfpathlineto{\pgfqpoint{0.937841in}{0.923485in}}%
\pgfpathlineto{\pgfqpoint{0.945370in}{0.908022in}}%
\pgfpathlineto{\pgfqpoint{0.952909in}{0.892541in}}%
\pgfpathlineto{\pgfqpoint{0.965923in}{0.896259in}}%
\pgfpathlineto{\pgfqpoint{0.979438in}{0.899629in}}%
\pgfpathclose%
\pgfusepath{fill}%
\end{pgfscope}%
\begin{pgfscope}%
\pgfpathrectangle{\pgfqpoint{0.407592in}{0.222904in}}{\pgfqpoint{1.411163in}{1.411163in}}%
\pgfusepath{clip}%
\pgfsetbuttcap%
\pgfsetroundjoin%
\definecolor{currentfill}{rgb}{0.749020,0.772549,0.792157}%
\pgfsetfillcolor{currentfill}%
\pgfsetfillopacity{0.160000}%
\pgfsetlinewidth{0.000000pt}%
\definecolor{currentstroke}{rgb}{0.000000,0.000000,0.000000}%
\pgfsetstrokecolor{currentstroke}%
\pgfsetdash{}{0pt}%
\pgfpathmoveto{\pgfqpoint{1.410630in}{1.050344in}}%
\pgfpathlineto{\pgfqpoint{1.421198in}{1.064851in}}%
\pgfpathlineto{\pgfqpoint{1.431759in}{1.079348in}}%
\pgfpathlineto{\pgfqpoint{1.414446in}{1.087484in}}%
\pgfpathlineto{\pgfqpoint{1.395851in}{1.095124in}}%
\pgfpathlineto{\pgfqpoint{1.386447in}{1.080213in}}%
\pgfpathlineto{\pgfqpoint{1.377033in}{1.065287in}}%
\pgfpathlineto{\pgfqpoint{1.394437in}{1.058049in}}%
\pgfpathlineto{\pgfqpoint{1.410630in}{1.050344in}}%
\pgfpathclose%
\pgfusepath{fill}%
\end{pgfscope}%
\begin{pgfscope}%
\pgfpathrectangle{\pgfqpoint{0.407592in}{0.222904in}}{\pgfqpoint{1.411163in}{1.411163in}}%
\pgfusepath{clip}%
\pgfsetbuttcap%
\pgfsetroundjoin%
\definecolor{currentfill}{rgb}{0.749020,0.772549,0.792157}%
\pgfsetfillcolor{currentfill}%
\pgfsetfillopacity{0.160000}%
\pgfsetlinewidth{0.000000pt}%
\definecolor{currentstroke}{rgb}{0.000000,0.000000,0.000000}%
\pgfsetstrokecolor{currentstroke}%
\pgfsetdash{}{0pt}%
\pgfpathmoveto{\pgfqpoint{0.546597in}{1.393285in}}%
\pgfpathlineto{\pgfqpoint{0.536706in}{1.407526in}}%
\pgfpathlineto{\pgfqpoint{0.526822in}{1.421759in}}%
\pgfpathlineto{\pgfqpoint{0.494965in}{1.408305in}}%
\pgfpathlineto{\pgfqpoint{0.465273in}{1.394126in}}%
\pgfpathlineto{\pgfqpoint{0.476345in}{1.380243in}}%
\pgfpathlineto{\pgfqpoint{0.487421in}{1.366357in}}%
\pgfpathlineto{\pgfqpoint{0.515957in}{1.380172in}}%
\pgfpathlineto{\pgfqpoint{0.546597in}{1.393285in}}%
\pgfpathclose%
\pgfusepath{fill}%
\end{pgfscope}%
\begin{pgfscope}%
\pgfpathrectangle{\pgfqpoint{0.407592in}{0.222904in}}{\pgfqpoint{1.411163in}{1.411163in}}%
\pgfusepath{clip}%
\pgfsetbuttcap%
\pgfsetroundjoin%
\definecolor{currentfill}{rgb}{0.749020,0.772549,0.792157}%
\pgfsetfillcolor{currentfill}%
\pgfsetfillopacity{0.160000}%
\pgfsetlinewidth{0.000000pt}%
\definecolor{currentstroke}{rgb}{0.000000,0.000000,0.000000}%
\pgfsetstrokecolor{currentstroke}%
\pgfsetdash{}{0pt}%
\pgfpathmoveto{\pgfqpoint{0.922810in}{0.954352in}}%
\pgfpathlineto{\pgfqpoint{0.915309in}{0.969756in}}%
\pgfpathlineto{\pgfqpoint{0.907817in}{0.985142in}}%
\pgfpathlineto{\pgfqpoint{0.891185in}{0.979906in}}%
\pgfpathlineto{\pgfqpoint{0.875341in}{0.974253in}}%
\pgfpathlineto{\pgfqpoint{0.884195in}{0.959266in}}%
\pgfpathlineto{\pgfqpoint{0.893058in}{0.944265in}}%
\pgfpathlineto{\pgfqpoint{0.907570in}{0.949501in}}%
\pgfpathlineto{\pgfqpoint{0.922810in}{0.954352in}}%
\pgfpathclose%
\pgfusepath{fill}%
\end{pgfscope}%
\begin{pgfscope}%
\pgfpathrectangle{\pgfqpoint{0.407592in}{0.222904in}}{\pgfqpoint{1.411163in}{1.411163in}}%
\pgfusepath{clip}%
\pgfsetbuttcap%
\pgfsetroundjoin%
\definecolor{currentfill}{rgb}{0.749020,0.772549,0.792157}%
\pgfsetfillcolor{currentfill}%
\pgfsetfillopacity{0.160000}%
\pgfsetlinewidth{0.000000pt}%
\definecolor{currentstroke}{rgb}{0.000000,0.000000,0.000000}%
\pgfsetstrokecolor{currentstroke}%
\pgfsetdash{}{0pt}%
\pgfpathmoveto{\pgfqpoint{1.669786in}{1.313440in}}%
\pgfpathlineto{\pgfqpoint{1.681282in}{1.327382in}}%
\pgfpathlineto{\pgfqpoint{1.692774in}{1.341320in}}%
\pgfpathlineto{\pgfqpoint{1.668225in}{1.354628in}}%
\pgfpathlineto{\pgfqpoint{1.641494in}{1.367249in}}%
\pgfpathlineto{\pgfqpoint{1.631074in}{1.352946in}}%
\pgfpathlineto{\pgfqpoint{1.620647in}{1.338633in}}%
\pgfpathlineto{\pgfqpoint{1.646273in}{1.326368in}}%
\pgfpathlineto{\pgfqpoint{1.669786in}{1.313440in}}%
\pgfpathclose%
\pgfusepath{fill}%
\end{pgfscope}%
\begin{pgfscope}%
\pgfpathrectangle{\pgfqpoint{0.407592in}{0.222904in}}{\pgfqpoint{1.411163in}{1.411163in}}%
\pgfusepath{clip}%
\pgfsetbuttcap%
\pgfsetroundjoin%
\definecolor{currentfill}{rgb}{0.749020,0.772549,0.792157}%
\pgfsetfillcolor{currentfill}%
\pgfsetfillopacity{0.160000}%
\pgfsetlinewidth{0.000000pt}%
\definecolor{currentstroke}{rgb}{0.000000,0.000000,0.000000}%
\pgfsetstrokecolor{currentstroke}%
\pgfsetdash{}{0pt}%
\pgfpathmoveto{\pgfqpoint{1.116365in}{0.812715in}}%
\pgfpathlineto{\pgfqpoint{1.118516in}{0.829354in}}%
\pgfpathlineto{\pgfqpoint{1.120662in}{0.845962in}}%
\pgfpathlineto{\pgfqpoint{1.108745in}{0.846419in}}%
\pgfpathlineto{\pgfqpoint{1.096748in}{0.846549in}}%
\pgfpathlineto{\pgfqpoint{1.096288in}{0.829903in}}%
\pgfpathlineto{\pgfqpoint{1.095828in}{0.813224in}}%
\pgfpathlineto{\pgfqpoint{1.106130in}{0.813111in}}%
\pgfpathlineto{\pgfqpoint{1.116365in}{0.812715in}}%
\pgfpathclose%
\pgfusepath{fill}%
\end{pgfscope}%
\begin{pgfscope}%
\pgfpathrectangle{\pgfqpoint{0.407592in}{0.222904in}}{\pgfqpoint{1.411163in}{1.411163in}}%
\pgfusepath{clip}%
\pgfsetbuttcap%
\pgfsetroundjoin%
\definecolor{currentfill}{rgb}{0.749020,0.772549,0.792157}%
\pgfsetfillcolor{currentfill}%
\pgfsetfillopacity{0.160000}%
\pgfsetlinewidth{0.000000pt}%
\definecolor{currentstroke}{rgb}{0.000000,0.000000,0.000000}%
\pgfsetstrokecolor{currentstroke}%
\pgfsetdash{}{0pt}%
\pgfpathmoveto{\pgfqpoint{1.095828in}{0.813224in}}%
\pgfpathlineto{\pgfqpoint{1.096288in}{0.829903in}}%
\pgfpathlineto{\pgfqpoint{1.096748in}{0.846549in}}%
\pgfpathlineto{\pgfqpoint{1.084722in}{0.846353in}}%
\pgfpathlineto{\pgfqpoint{1.072715in}{0.845831in}}%
\pgfpathlineto{\pgfqpoint{1.073950in}{0.829232in}}%
\pgfpathlineto{\pgfqpoint{1.075188in}{0.812601in}}%
\pgfpathlineto{\pgfqpoint{1.085499in}{0.813054in}}%
\pgfpathlineto{\pgfqpoint{1.095828in}{0.813224in}}%
\pgfpathclose%
\pgfusepath{fill}%
\end{pgfscope}%
\begin{pgfscope}%
\pgfpathrectangle{\pgfqpoint{0.407592in}{0.222904in}}{\pgfqpoint{1.411163in}{1.411163in}}%
\pgfusepath{clip}%
\pgfsetbuttcap%
\pgfsetroundjoin%
\definecolor{currentfill}{rgb}{0.749020,0.772549,0.792157}%
\pgfsetfillcolor{currentfill}%
\pgfsetfillopacity{0.160000}%
\pgfsetlinewidth{0.000000pt}%
\definecolor{currentstroke}{rgb}{0.000000,0.000000,0.000000}%
\pgfsetstrokecolor{currentstroke}%
\pgfsetdash{}{0pt}%
\pgfpathmoveto{\pgfqpoint{0.822393in}{1.063874in}}%
\pgfpathlineto{\pgfqpoint{0.813598in}{1.078760in}}%
\pgfpathlineto{\pgfqpoint{0.804811in}{1.093633in}}%
\pgfpathlineto{\pgfqpoint{0.784714in}{1.085891in}}%
\pgfpathlineto{\pgfqpoint{0.765804in}{1.077659in}}%
\pgfpathlineto{\pgfqpoint{0.775842in}{1.063206in}}%
\pgfpathlineto{\pgfqpoint{0.785886in}{1.048744in}}%
\pgfpathlineto{\pgfqpoint{0.803578in}{1.056540in}}%
\pgfpathlineto{\pgfqpoint{0.822393in}{1.063874in}}%
\pgfpathclose%
\pgfusepath{fill}%
\end{pgfscope}%
\begin{pgfscope}%
\pgfpathrectangle{\pgfqpoint{0.407592in}{0.222904in}}{\pgfqpoint{1.411163in}{1.411163in}}%
\pgfusepath{clip}%
\pgfsetbuttcap%
\pgfsetroundjoin%
\definecolor{currentfill}{rgb}{0.749020,0.772549,0.792157}%
\pgfsetfillcolor{currentfill}%
\pgfsetfillopacity{0.160000}%
\pgfsetlinewidth{0.000000pt}%
\definecolor{currentstroke}{rgb}{0.000000,0.000000,0.000000}%
\pgfsetstrokecolor{currentstroke}%
\pgfsetdash{}{0pt}%
\pgfpathmoveto{\pgfqpoint{0.566398in}{1.364774in}}%
\pgfpathlineto{\pgfqpoint{0.556495in}{1.379034in}}%
\pgfpathlineto{\pgfqpoint{0.546597in}{1.393285in}}%
\pgfpathlineto{\pgfqpoint{0.515957in}{1.380172in}}%
\pgfpathlineto{\pgfqpoint{0.487421in}{1.366357in}}%
\pgfpathlineto{\pgfqpoint{0.498499in}{1.352466in}}%
\pgfpathlineto{\pgfqpoint{0.509581in}{1.338571in}}%
\pgfpathlineto{\pgfqpoint{0.536969in}{1.352012in}}%
\pgfpathlineto{\pgfqpoint{0.566398in}{1.364774in}}%
\pgfpathclose%
\pgfusepath{fill}%
\end{pgfscope}%
\begin{pgfscope}%
\pgfpathrectangle{\pgfqpoint{0.407592in}{0.222904in}}{\pgfqpoint{1.411163in}{1.411163in}}%
\pgfusepath{clip}%
\pgfsetbuttcap%
\pgfsetroundjoin%
\definecolor{currentfill}{rgb}{0.749020,0.772549,0.792157}%
\pgfsetfillcolor{currentfill}%
\pgfsetfillopacity{0.160000}%
\pgfsetlinewidth{0.000000pt}%
\definecolor{currentstroke}{rgb}{0.000000,0.000000,0.000000}%
\pgfsetstrokecolor{currentstroke}%
\pgfsetdash{}{0pt}%
\pgfpathmoveto{\pgfqpoint{1.646783in}{1.285542in}}%
\pgfpathlineto{\pgfqpoint{1.658286in}{1.299493in}}%
\pgfpathlineto{\pgfqpoint{1.669786in}{1.313440in}}%
\pgfpathlineto{\pgfqpoint{1.646273in}{1.326368in}}%
\pgfpathlineto{\pgfqpoint{1.620647in}{1.338633in}}%
\pgfpathlineto{\pgfqpoint{1.610213in}{1.324310in}}%
\pgfpathlineto{\pgfqpoint{1.599772in}{1.309978in}}%
\pgfpathlineto{\pgfqpoint{1.624299in}{1.298080in}}%
\pgfpathlineto{\pgfqpoint{1.646783in}{1.285542in}}%
\pgfpathclose%
\pgfusepath{fill}%
\end{pgfscope}%
\begin{pgfscope}%
\pgfpathrectangle{\pgfqpoint{0.407592in}{0.222904in}}{\pgfqpoint{1.411163in}{1.411163in}}%
\pgfusepath{clip}%
\pgfsetbuttcap%
\pgfsetroundjoin%
\definecolor{currentfill}{rgb}{0.749020,0.772549,0.792157}%
\pgfsetfillcolor{currentfill}%
\pgfsetfillopacity{0.160000}%
\pgfsetlinewidth{0.000000pt}%
\definecolor{currentstroke}{rgb}{0.000000,0.000000,0.000000}%
\pgfsetstrokecolor{currentstroke}%
\pgfsetdash{}{0pt}%
\pgfpathmoveto{\pgfqpoint{1.136461in}{0.811081in}}%
\pgfpathlineto{\pgfqpoint{1.140268in}{0.827595in}}%
\pgfpathlineto{\pgfqpoint{1.144068in}{0.844079in}}%
\pgfpathlineto{\pgfqpoint{1.132453in}{0.845181in}}%
\pgfpathlineto{\pgfqpoint{1.120662in}{0.845962in}}%
\pgfpathlineto{\pgfqpoint{1.118516in}{0.829354in}}%
\pgfpathlineto{\pgfqpoint{1.116365in}{0.812715in}}%
\pgfpathlineto{\pgfqpoint{1.126489in}{0.812037in}}%
\pgfpathlineto{\pgfqpoint{1.136461in}{0.811081in}}%
\pgfpathclose%
\pgfusepath{fill}%
\end{pgfscope}%
\begin{pgfscope}%
\pgfpathrectangle{\pgfqpoint{0.407592in}{0.222904in}}{\pgfqpoint{1.411163in}{1.411163in}}%
\pgfusepath{clip}%
\pgfsetbuttcap%
\pgfsetroundjoin%
\definecolor{currentfill}{rgb}{0.749020,0.772549,0.792157}%
\pgfsetfillcolor{currentfill}%
\pgfsetfillopacity{0.160000}%
\pgfsetlinewidth{0.000000pt}%
\definecolor{currentstroke}{rgb}{0.000000,0.000000,0.000000}%
\pgfsetstrokecolor{currentstroke}%
\pgfsetdash{}{0pt}%
\pgfpathmoveto{\pgfqpoint{1.075188in}{0.812601in}}%
\pgfpathlineto{\pgfqpoint{1.073950in}{0.829232in}}%
\pgfpathlineto{\pgfqpoint{1.072715in}{0.845831in}}%
\pgfpathlineto{\pgfqpoint{1.060776in}{0.844986in}}%
\pgfpathlineto{\pgfqpoint{1.048955in}{0.843819in}}%
\pgfpathlineto{\pgfqpoint{1.051868in}{0.827353in}}%
\pgfpathlineto{\pgfqpoint{1.054787in}{0.810856in}}%
\pgfpathlineto{\pgfqpoint{1.064937in}{0.811868in}}%
\pgfpathlineto{\pgfqpoint{1.075188in}{0.812601in}}%
\pgfpathclose%
\pgfusepath{fill}%
\end{pgfscope}%
\begin{pgfscope}%
\pgfpathrectangle{\pgfqpoint{0.407592in}{0.222904in}}{\pgfqpoint{1.411163in}{1.411163in}}%
\pgfusepath{clip}%
\pgfsetbuttcap%
\pgfsetroundjoin%
\definecolor{currentfill}{rgb}{0.749020,0.772549,0.792157}%
\pgfsetfillcolor{currentfill}%
\pgfsetfillopacity{0.160000}%
\pgfsetlinewidth{0.000000pt}%
\definecolor{currentstroke}{rgb}{0.000000,0.000000,0.000000}%
\pgfsetstrokecolor{currentstroke}%
\pgfsetdash{}{0pt}%
\pgfpathmoveto{\pgfqpoint{1.389472in}{1.021301in}}%
\pgfpathlineto{\pgfqpoint{1.400055in}{1.035827in}}%
\pgfpathlineto{\pgfqpoint{1.410630in}{1.050344in}}%
\pgfpathlineto{\pgfqpoint{1.394437in}{1.058049in}}%
\pgfpathlineto{\pgfqpoint{1.377033in}{1.065287in}}%
\pgfpathlineto{\pgfqpoint{1.367610in}{1.050346in}}%
\pgfpathlineto{\pgfqpoint{1.358178in}{1.035391in}}%
\pgfpathlineto{\pgfqpoint{1.374395in}{1.028565in}}%
\pgfpathlineto{\pgfqpoint{1.389472in}{1.021301in}}%
\pgfpathclose%
\pgfusepath{fill}%
\end{pgfscope}%
\begin{pgfscope}%
\pgfpathrectangle{\pgfqpoint{0.407592in}{0.222904in}}{\pgfqpoint{1.411163in}{1.411163in}}%
\pgfusepath{clip}%
\pgfsetbuttcap%
\pgfsetroundjoin%
\definecolor{currentfill}{rgb}{0.749020,0.772549,0.792157}%
\pgfsetfillcolor{currentfill}%
\pgfsetfillopacity{0.160000}%
\pgfsetlinewidth{0.000000pt}%
\definecolor{currentstroke}{rgb}{0.000000,0.000000,0.000000}%
\pgfsetstrokecolor{currentstroke}%
\pgfsetdash{}{0pt}%
\pgfpathmoveto{\pgfqpoint{1.187830in}{0.836558in}}%
\pgfpathlineto{\pgfqpoint{1.194722in}{0.852520in}}%
\pgfpathlineto{\pgfqpoint{1.201605in}{0.868458in}}%
\pgfpathlineto{\pgfqpoint{1.189677in}{0.871095in}}%
\pgfpathlineto{\pgfqpoint{1.177341in}{0.873396in}}%
\pgfpathlineto{\pgfqpoint{1.171965in}{0.857176in}}%
\pgfpathlineto{\pgfqpoint{1.166581in}{0.840928in}}%
\pgfpathlineto{\pgfqpoint{1.177386in}{0.838892in}}%
\pgfpathlineto{\pgfqpoint{1.187830in}{0.836558in}}%
\pgfpathclose%
\pgfusepath{fill}%
\end{pgfscope}%
\begin{pgfscope}%
\pgfpathrectangle{\pgfqpoint{0.407592in}{0.222904in}}{\pgfqpoint{1.411163in}{1.411163in}}%
\pgfusepath{clip}%
\pgfsetbuttcap%
\pgfsetroundjoin%
\definecolor{currentfill}{rgb}{0.749020,0.772549,0.792157}%
\pgfsetfillcolor{currentfill}%
\pgfsetfillopacity{0.160000}%
\pgfsetlinewidth{0.000000pt}%
\definecolor{currentstroke}{rgb}{0.000000,0.000000,0.000000}%
\pgfsetstrokecolor{currentstroke}%
\pgfsetdash{}{0pt}%
\pgfpathmoveto{\pgfqpoint{1.282383in}{0.915211in}}%
\pgfpathlineto{\pgfqpoint{1.291891in}{0.930285in}}%
\pgfpathlineto{\pgfqpoint{1.301388in}{0.945345in}}%
\pgfpathlineto{\pgfqpoint{1.287852in}{0.950506in}}%
\pgfpathlineto{\pgfqpoint{1.273514in}{0.955277in}}%
\pgfpathlineto{\pgfqpoint{1.265294in}{0.939816in}}%
\pgfpathlineto{\pgfqpoint{1.257064in}{0.924335in}}%
\pgfpathlineto{\pgfqpoint{1.270091in}{0.919951in}}%
\pgfpathlineto{\pgfqpoint{1.282383in}{0.915211in}}%
\pgfpathclose%
\pgfusepath{fill}%
\end{pgfscope}%
\begin{pgfscope}%
\pgfpathrectangle{\pgfqpoint{0.407592in}{0.222904in}}{\pgfqpoint{1.411163in}{1.411163in}}%
\pgfusepath{clip}%
\pgfsetbuttcap%
\pgfsetroundjoin%
\definecolor{currentfill}{rgb}{0.749020,0.772549,0.792157}%
\pgfsetfillcolor{currentfill}%
\pgfsetfillopacity{0.160000}%
\pgfsetlinewidth{0.000000pt}%
\definecolor{currentstroke}{rgb}{0.000000,0.000000,0.000000}%
\pgfsetstrokecolor{currentstroke}%
\pgfsetdash{}{0pt}%
\pgfpathmoveto{\pgfqpoint{1.025857in}{0.840544in}}%
\pgfpathlineto{\pgfqpoint{1.021317in}{0.856767in}}%
\pgfpathlineto{\pgfqpoint{1.016784in}{0.872962in}}%
\pgfpathlineto{\pgfqpoint{1.004020in}{0.870593in}}%
\pgfpathlineto{\pgfqpoint{0.991602in}{0.867889in}}%
\pgfpathlineto{\pgfqpoint{0.997698in}{0.851984in}}%
\pgfpathlineto{\pgfqpoint{1.003802in}{0.836054in}}%
\pgfpathlineto{\pgfqpoint{1.014676in}{0.838447in}}%
\pgfpathlineto{\pgfqpoint{1.025857in}{0.840544in}}%
\pgfpathclose%
\pgfusepath{fill}%
\end{pgfscope}%
\begin{pgfscope}%
\pgfpathrectangle{\pgfqpoint{0.407592in}{0.222904in}}{\pgfqpoint{1.411163in}{1.411163in}}%
\pgfusepath{clip}%
\pgfsetbuttcap%
\pgfsetroundjoin%
\definecolor{currentfill}{rgb}{0.749020,0.772549,0.792157}%
\pgfsetfillcolor{currentfill}%
\pgfsetfillopacity{0.160000}%
\pgfsetlinewidth{0.000000pt}%
\definecolor{currentstroke}{rgb}{0.000000,0.000000,0.000000}%
\pgfsetstrokecolor{currentstroke}%
\pgfsetdash{}{0pt}%
\pgfpathmoveto{\pgfqpoint{1.224037in}{0.862214in}}%
\pgfpathlineto{\pgfqpoint{1.232309in}{0.877774in}}%
\pgfpathlineto{\pgfqpoint{1.240571in}{0.893314in}}%
\pgfpathlineto{\pgfqpoint{1.228237in}{0.896964in}}%
\pgfpathlineto{\pgfqpoint{1.215338in}{0.900262in}}%
\pgfpathlineto{\pgfqpoint{1.208477in}{0.884372in}}%
\pgfpathlineto{\pgfqpoint{1.201605in}{0.868458in}}%
\pgfpathlineto{\pgfqpoint{1.213074in}{0.865494in}}%
\pgfpathlineto{\pgfqpoint{1.224037in}{0.862214in}}%
\pgfpathclose%
\pgfusepath{fill}%
\end{pgfscope}%
\begin{pgfscope}%
\pgfpathrectangle{\pgfqpoint{0.407592in}{0.222904in}}{\pgfqpoint{1.411163in}{1.411163in}}%
\pgfusepath{clip}%
\pgfsetbuttcap%
\pgfsetroundjoin%
\definecolor{currentfill}{rgb}{0.749020,0.772549,0.792157}%
\pgfsetfillcolor{currentfill}%
\pgfsetfillopacity{0.160000}%
\pgfsetlinewidth{0.000000pt}%
\definecolor{currentstroke}{rgb}{0.000000,0.000000,0.000000}%
\pgfsetstrokecolor{currentstroke}%
\pgfsetdash{}{0pt}%
\pgfpathmoveto{\pgfqpoint{0.586224in}{1.336227in}}%
\pgfpathlineto{\pgfqpoint{0.576308in}{1.350505in}}%
\pgfpathlineto{\pgfqpoint{0.566398in}{1.364774in}}%
\pgfpathlineto{\pgfqpoint{0.536969in}{1.352012in}}%
\pgfpathlineto{\pgfqpoint{0.509581in}{1.338571in}}%
\pgfpathlineto{\pgfqpoint{0.520666in}{1.324673in}}%
\pgfpathlineto{\pgfqpoint{0.531754in}{1.310770in}}%
\pgfpathlineto{\pgfqpoint{0.558000in}{1.323826in}}%
\pgfpathlineto{\pgfqpoint{0.586224in}{1.336227in}}%
\pgfpathclose%
\pgfusepath{fill}%
\end{pgfscope}%
\begin{pgfscope}%
\pgfpathrectangle{\pgfqpoint{0.407592in}{0.222904in}}{\pgfqpoint{1.411163in}{1.411163in}}%
\pgfusepath{clip}%
\pgfsetbuttcap%
\pgfsetroundjoin%
\definecolor{currentfill}{rgb}{0.749020,0.772549,0.792157}%
\pgfsetfillcolor{currentfill}%
\pgfsetfillopacity{0.160000}%
\pgfsetlinewidth{0.000000pt}%
\definecolor{currentstroke}{rgb}{0.000000,0.000000,0.000000}%
\pgfsetstrokecolor{currentstroke}%
\pgfsetdash{}{0pt}%
\pgfpathmoveto{\pgfqpoint{1.623766in}{1.257625in}}%
\pgfpathlineto{\pgfqpoint{1.635277in}{1.271586in}}%
\pgfpathlineto{\pgfqpoint{1.646783in}{1.285542in}}%
\pgfpathlineto{\pgfqpoint{1.624299in}{1.298080in}}%
\pgfpathlineto{\pgfqpoint{1.599772in}{1.309978in}}%
\pgfpathlineto{\pgfqpoint{1.589325in}{1.295637in}}%
\pgfpathlineto{\pgfqpoint{1.578870in}{1.281285in}}%
\pgfpathlineto{\pgfqpoint{1.602304in}{1.269763in}}%
\pgfpathlineto{\pgfqpoint{1.623766in}{1.257625in}}%
\pgfpathclose%
\pgfusepath{fill}%
\end{pgfscope}%
\begin{pgfscope}%
\pgfpathrectangle{\pgfqpoint{0.407592in}{0.222904in}}{\pgfqpoint{1.411163in}{1.411163in}}%
\pgfusepath{clip}%
\pgfsetbuttcap%
\pgfsetroundjoin%
\definecolor{currentfill}{rgb}{0.749020,0.772549,0.792157}%
\pgfsetfillcolor{currentfill}%
\pgfsetfillopacity{0.160000}%
\pgfsetlinewidth{0.000000pt}%
\definecolor{currentstroke}{rgb}{0.000000,0.000000,0.000000}%
\pgfsetstrokecolor{currentstroke}%
\pgfsetdash{}{0pt}%
\pgfpathmoveto{\pgfqpoint{0.840008in}{1.034058in}}%
\pgfpathlineto{\pgfqpoint{0.831196in}{1.048973in}}%
\pgfpathlineto{\pgfqpoint{0.822393in}{1.063874in}}%
\pgfpathlineto{\pgfqpoint{0.803578in}{1.056540in}}%
\pgfpathlineto{\pgfqpoint{0.785886in}{1.048744in}}%
\pgfpathlineto{\pgfqpoint{0.795936in}{1.034273in}}%
\pgfpathlineto{\pgfqpoint{0.805993in}{1.019793in}}%
\pgfpathlineto{\pgfqpoint{0.822472in}{1.027142in}}%
\pgfpathlineto{\pgfqpoint{0.840008in}{1.034058in}}%
\pgfpathclose%
\pgfusepath{fill}%
\end{pgfscope}%
\begin{pgfscope}%
\pgfpathrectangle{\pgfqpoint{0.407592in}{0.222904in}}{\pgfqpoint{1.411163in}{1.411163in}}%
\pgfusepath{clip}%
\pgfsetbuttcap%
\pgfsetroundjoin%
\definecolor{currentfill}{rgb}{0.749020,0.772549,0.792157}%
\pgfsetfillcolor{currentfill}%
\pgfsetfillopacity{0.160000}%
\pgfsetlinewidth{0.000000pt}%
\definecolor{currentstroke}{rgb}{0.000000,0.000000,0.000000}%
\pgfsetstrokecolor{currentstroke}%
\pgfsetdash{}{0pt}%
\pgfpathmoveto{\pgfqpoint{0.991602in}{0.867889in}}%
\pgfpathlineto{\pgfqpoint{0.985516in}{0.883771in}}%
\pgfpathlineto{\pgfqpoint{0.979438in}{0.899629in}}%
\pgfpathlineto{\pgfqpoint{0.965923in}{0.896259in}}%
\pgfpathlineto{\pgfqpoint{0.952909in}{0.892541in}}%
\pgfpathlineto{\pgfqpoint{0.960457in}{0.877040in}}%
\pgfpathlineto{\pgfqpoint{0.968014in}{0.861519in}}%
\pgfpathlineto{\pgfqpoint{0.979583in}{0.864861in}}%
\pgfpathlineto{\pgfqpoint{0.991602in}{0.867889in}}%
\pgfpathclose%
\pgfusepath{fill}%
\end{pgfscope}%
\begin{pgfscope}%
\pgfpathrectangle{\pgfqpoint{0.407592in}{0.222904in}}{\pgfqpoint{1.411163in}{1.411163in}}%
\pgfusepath{clip}%
\pgfsetbuttcap%
\pgfsetroundjoin%
\definecolor{currentfill}{rgb}{0.749020,0.772549,0.792157}%
\pgfsetfillcolor{currentfill}%
\pgfsetfillopacity{0.160000}%
\pgfsetlinewidth{0.000000pt}%
\definecolor{currentstroke}{rgb}{0.000000,0.000000,0.000000}%
\pgfsetstrokecolor{currentstroke}%
\pgfsetdash{}{0pt}%
\pgfpathmoveto{\pgfqpoint{0.937841in}{0.923485in}}%
\pgfpathlineto{\pgfqpoint{0.930321in}{0.938928in}}%
\pgfpathlineto{\pgfqpoint{0.922810in}{0.954352in}}%
\pgfpathlineto{\pgfqpoint{0.907570in}{0.949501in}}%
\pgfpathlineto{\pgfqpoint{0.893058in}{0.944265in}}%
\pgfpathlineto{\pgfqpoint{0.901930in}{0.929249in}}%
\pgfpathlineto{\pgfqpoint{0.910810in}{0.914218in}}%
\pgfpathlineto{\pgfqpoint{0.923991in}{0.919028in}}%
\pgfpathlineto{\pgfqpoint{0.937841in}{0.923485in}}%
\pgfpathclose%
\pgfusepath{fill}%
\end{pgfscope}%
\begin{pgfscope}%
\pgfpathrectangle{\pgfqpoint{0.407592in}{0.222904in}}{\pgfqpoint{1.411163in}{1.411163in}}%
\pgfusepath{clip}%
\pgfsetbuttcap%
\pgfsetroundjoin%
\definecolor{currentfill}{rgb}{0.749020,0.772549,0.792157}%
\pgfsetfillcolor{currentfill}%
\pgfsetfillopacity{0.160000}%
\pgfsetlinewidth{0.000000pt}%
\definecolor{currentstroke}{rgb}{0.000000,0.000000,0.000000}%
\pgfsetstrokecolor{currentstroke}%
\pgfsetdash{}{0pt}%
\pgfpathmoveto{\pgfqpoint{1.155785in}{0.808349in}}%
\pgfpathlineto{\pgfqpoint{1.161187in}{0.824653in}}%
\pgfpathlineto{\pgfqpoint{1.166581in}{0.840928in}}%
\pgfpathlineto{\pgfqpoint{1.155459in}{0.842659in}}%
\pgfpathlineto{\pgfqpoint{1.144068in}{0.844079in}}%
\pgfpathlineto{\pgfqpoint{1.140268in}{0.827595in}}%
\pgfpathlineto{\pgfqpoint{1.136461in}{0.811081in}}%
\pgfpathlineto{\pgfqpoint{1.146240in}{0.809851in}}%
\pgfpathlineto{\pgfqpoint{1.155785in}{0.808349in}}%
\pgfpathclose%
\pgfusepath{fill}%
\end{pgfscope}%
\begin{pgfscope}%
\pgfpathrectangle{\pgfqpoint{0.407592in}{0.222904in}}{\pgfqpoint{1.411163in}{1.411163in}}%
\pgfusepath{clip}%
\pgfsetbuttcap%
\pgfsetroundjoin%
\definecolor{currentfill}{rgb}{0.749020,0.772549,0.792157}%
\pgfsetfillcolor{currentfill}%
\pgfsetfillopacity{0.160000}%
\pgfsetlinewidth{0.000000pt}%
\definecolor{currentstroke}{rgb}{0.000000,0.000000,0.000000}%
\pgfsetstrokecolor{currentstroke}%
\pgfsetdash{}{0pt}%
\pgfpathmoveto{\pgfqpoint{1.054787in}{0.810856in}}%
\pgfpathlineto{\pgfqpoint{1.051868in}{0.827353in}}%
\pgfpathlineto{\pgfqpoint{1.048955in}{0.843819in}}%
\pgfpathlineto{\pgfqpoint{1.037299in}{0.842337in}}%
\pgfpathlineto{\pgfqpoint{1.025857in}{0.840544in}}%
\pgfpathlineto{\pgfqpoint{1.030405in}{0.824294in}}%
\pgfpathlineto{\pgfqpoint{1.034960in}{0.808016in}}%
\pgfpathlineto{\pgfqpoint{1.044781in}{0.809571in}}%
\pgfpathlineto{\pgfqpoint{1.054787in}{0.810856in}}%
\pgfpathclose%
\pgfusepath{fill}%
\end{pgfscope}%
\begin{pgfscope}%
\pgfpathrectangle{\pgfqpoint{0.407592in}{0.222904in}}{\pgfqpoint{1.411163in}{1.411163in}}%
\pgfusepath{clip}%
\pgfsetbuttcap%
\pgfsetroundjoin%
\definecolor{currentfill}{rgb}{0.749020,0.772549,0.792157}%
\pgfsetfillcolor{currentfill}%
\pgfsetfillopacity{0.160000}%
\pgfsetlinewidth{0.000000pt}%
\definecolor{currentstroke}{rgb}{0.000000,0.000000,0.000000}%
\pgfsetstrokecolor{currentstroke}%
\pgfsetdash{}{0pt}%
\pgfpathmoveto{\pgfqpoint{0.606076in}{1.307644in}}%
\pgfpathlineto{\pgfqpoint{0.596147in}{1.321940in}}%
\pgfpathlineto{\pgfqpoint{0.586224in}{1.336227in}}%
\pgfpathlineto{\pgfqpoint{0.558000in}{1.323826in}}%
\pgfpathlineto{\pgfqpoint{0.531754in}{1.310770in}}%
\pgfpathlineto{\pgfqpoint{0.542845in}{1.296863in}}%
\pgfpathlineto{\pgfqpoint{0.553940in}{1.282953in}}%
\pgfpathlineto{\pgfqpoint{0.579051in}{1.295614in}}%
\pgfpathlineto{\pgfqpoint{0.606076in}{1.307644in}}%
\pgfpathclose%
\pgfusepath{fill}%
\end{pgfscope}%
\begin{pgfscope}%
\pgfpathrectangle{\pgfqpoint{0.407592in}{0.222904in}}{\pgfqpoint{1.411163in}{1.411163in}}%
\pgfusepath{clip}%
\pgfsetbuttcap%
\pgfsetroundjoin%
\definecolor{currentfill}{rgb}{0.749020,0.772549,0.792157}%
\pgfsetfillcolor{currentfill}%
\pgfsetfillopacity{0.160000}%
\pgfsetlinewidth{0.000000pt}%
\definecolor{currentstroke}{rgb}{0.000000,0.000000,0.000000}%
\pgfsetstrokecolor{currentstroke}%
\pgfsetdash{}{0pt}%
\pgfpathmoveto{\pgfqpoint{1.368286in}{0.992218in}}%
\pgfpathlineto{\pgfqpoint{1.378882in}{1.006764in}}%
\pgfpathlineto{\pgfqpoint{1.389472in}{1.021301in}}%
\pgfpathlineto{\pgfqpoint{1.374395in}{1.028565in}}%
\pgfpathlineto{\pgfqpoint{1.358178in}{1.035391in}}%
\pgfpathlineto{\pgfqpoint{1.348737in}{1.020420in}}%
\pgfpathlineto{\pgfqpoint{1.339286in}{1.005435in}}%
\pgfpathlineto{\pgfqpoint{1.354319in}{0.999031in}}%
\pgfpathlineto{\pgfqpoint{1.368286in}{0.992218in}}%
\pgfpathclose%
\pgfusepath{fill}%
\end{pgfscope}%
\begin{pgfscope}%
\pgfpathrectangle{\pgfqpoint{0.407592in}{0.222904in}}{\pgfqpoint{1.411163in}{1.411163in}}%
\pgfusepath{clip}%
\pgfsetbuttcap%
\pgfsetroundjoin%
\definecolor{currentfill}{rgb}{0.749020,0.772549,0.792157}%
\pgfsetfillcolor{currentfill}%
\pgfsetfillopacity{0.160000}%
\pgfsetlinewidth{0.000000pt}%
\definecolor{currentstroke}{rgb}{0.000000,0.000000,0.000000}%
\pgfsetstrokecolor{currentstroke}%
\pgfsetdash{}{0pt}%
\pgfpathmoveto{\pgfqpoint{1.600734in}{1.229691in}}%
\pgfpathlineto{\pgfqpoint{1.612252in}{1.243660in}}%
\pgfpathlineto{\pgfqpoint{1.623766in}{1.257625in}}%
\pgfpathlineto{\pgfqpoint{1.602304in}{1.269763in}}%
\pgfpathlineto{\pgfqpoint{1.578870in}{1.281285in}}%
\pgfpathlineto{\pgfqpoint{1.568408in}{1.266924in}}%
\pgfpathlineto{\pgfqpoint{1.557939in}{1.252553in}}%
\pgfpathlineto{\pgfqpoint{1.580285in}{1.241417in}}%
\pgfpathlineto{\pgfqpoint{1.600734in}{1.229691in}}%
\pgfpathclose%
\pgfusepath{fill}%
\end{pgfscope}%
\begin{pgfscope}%
\pgfpathrectangle{\pgfqpoint{0.407592in}{0.222904in}}{\pgfqpoint{1.411163in}{1.411163in}}%
\pgfusepath{clip}%
\pgfsetbuttcap%
\pgfsetroundjoin%
\definecolor{currentfill}{rgb}{0.749020,0.772549,0.792157}%
\pgfsetfillcolor{currentfill}%
\pgfsetfillopacity{0.160000}%
\pgfsetlinewidth{0.000000pt}%
\definecolor{currentstroke}{rgb}{0.000000,0.000000,0.000000}%
\pgfsetstrokecolor{currentstroke}%
\pgfsetdash{}{0pt}%
\pgfpathmoveto{\pgfqpoint{0.625953in}{1.279025in}}%
\pgfpathlineto{\pgfqpoint{0.616011in}{1.293339in}}%
\pgfpathlineto{\pgfqpoint{0.606076in}{1.307644in}}%
\pgfpathlineto{\pgfqpoint{0.579051in}{1.295614in}}%
\pgfpathlineto{\pgfqpoint{0.553940in}{1.282953in}}%
\pgfpathlineto{\pgfqpoint{0.565037in}{1.269038in}}%
\pgfpathlineto{\pgfqpoint{0.576138in}{1.255120in}}%
\pgfpathlineto{\pgfqpoint{0.600121in}{1.267376in}}%
\pgfpathlineto{\pgfqpoint{0.625953in}{1.279025in}}%
\pgfpathclose%
\pgfusepath{fill}%
\end{pgfscope}%
\begin{pgfscope}%
\pgfpathrectangle{\pgfqpoint{0.407592in}{0.222904in}}{\pgfqpoint{1.411163in}{1.411163in}}%
\pgfusepath{clip}%
\pgfsetbuttcap%
\pgfsetroundjoin%
\definecolor{currentfill}{rgb}{0.749020,0.772549,0.792157}%
\pgfsetfillcolor{currentfill}%
\pgfsetfillopacity{0.160000}%
\pgfsetlinewidth{0.000000pt}%
\definecolor{currentstroke}{rgb}{0.000000,0.000000,0.000000}%
\pgfsetstrokecolor{currentstroke}%
\pgfsetdash{}{0pt}%
\pgfpathmoveto{\pgfqpoint{0.857658in}{1.004184in}}%
\pgfpathlineto{\pgfqpoint{0.848829in}{1.019128in}}%
\pgfpathlineto{\pgfqpoint{0.840008in}{1.034058in}}%
\pgfpathlineto{\pgfqpoint{0.822472in}{1.027142in}}%
\pgfpathlineto{\pgfqpoint{0.805993in}{1.019793in}}%
\pgfpathlineto{\pgfqpoint{0.816056in}{1.005303in}}%
\pgfpathlineto{\pgfqpoint{0.826126in}{0.990804in}}%
\pgfpathlineto{\pgfqpoint{0.841397in}{0.997696in}}%
\pgfpathlineto{\pgfqpoint{0.857658in}{1.004184in}}%
\pgfpathclose%
\pgfusepath{fill}%
\end{pgfscope}%
\begin{pgfscope}%
\pgfpathrectangle{\pgfqpoint{0.407592in}{0.222904in}}{\pgfqpoint{1.411163in}{1.411163in}}%
\pgfusepath{clip}%
\pgfsetbuttcap%
\pgfsetroundjoin%
\definecolor{currentfill}{rgb}{0.749020,0.772549,0.792157}%
\pgfsetfillcolor{currentfill}%
\pgfsetfillopacity{0.160000}%
\pgfsetlinewidth{0.000000pt}%
\definecolor{currentstroke}{rgb}{0.000000,0.000000,0.000000}%
\pgfsetstrokecolor{currentstroke}%
\pgfsetdash{}{0pt}%
\pgfpathmoveto{\pgfqpoint{1.263340in}{0.885016in}}%
\pgfpathlineto{\pgfqpoint{1.272866in}{0.900121in}}%
\pgfpathlineto{\pgfqpoint{1.282383in}{0.915211in}}%
\pgfpathlineto{\pgfqpoint{1.270091in}{0.919951in}}%
\pgfpathlineto{\pgfqpoint{1.257064in}{0.924335in}}%
\pgfpathlineto{\pgfqpoint{1.248823in}{0.908834in}}%
\pgfpathlineto{\pgfqpoint{1.240571in}{0.893314in}}%
\pgfpathlineto{\pgfqpoint{1.252289in}{0.889326in}}%
\pgfpathlineto{\pgfqpoint{1.263340in}{0.885016in}}%
\pgfpathclose%
\pgfusepath{fill}%
\end{pgfscope}%
\begin{pgfscope}%
\pgfpathrectangle{\pgfqpoint{0.407592in}{0.222904in}}{\pgfqpoint{1.411163in}{1.411163in}}%
\pgfusepath{clip}%
\pgfsetbuttcap%
\pgfsetroundjoin%
\definecolor{currentfill}{rgb}{0.749020,0.772549,0.792157}%
\pgfsetfillcolor{currentfill}%
\pgfsetfillopacity{0.160000}%
\pgfsetlinewidth{0.000000pt}%
\definecolor{currentstroke}{rgb}{0.000000,0.000000,0.000000}%
\pgfsetstrokecolor{currentstroke}%
\pgfsetdash{}{0pt}%
\pgfpathmoveto{\pgfqpoint{1.577687in}{1.201739in}}%
\pgfpathlineto{\pgfqpoint{1.589212in}{1.215717in}}%
\pgfpathlineto{\pgfqpoint{1.600734in}{1.229691in}}%
\pgfpathlineto{\pgfqpoint{1.580285in}{1.241417in}}%
\pgfpathlineto{\pgfqpoint{1.557939in}{1.252553in}}%
\pgfpathlineto{\pgfqpoint{1.547463in}{1.238173in}}%
\pgfpathlineto{\pgfqpoint{1.536980in}{1.223783in}}%
\pgfpathlineto{\pgfqpoint{1.558245in}{1.213043in}}%
\pgfpathlineto{\pgfqpoint{1.577687in}{1.201739in}}%
\pgfpathclose%
\pgfusepath{fill}%
\end{pgfscope}%
\begin{pgfscope}%
\pgfpathrectangle{\pgfqpoint{0.407592in}{0.222904in}}{\pgfqpoint{1.411163in}{1.411163in}}%
\pgfusepath{clip}%
\pgfsetbuttcap%
\pgfsetroundjoin%
\definecolor{currentfill}{rgb}{0.749020,0.772549,0.792157}%
\pgfsetfillcolor{currentfill}%
\pgfsetfillopacity{0.160000}%
\pgfsetlinewidth{0.000000pt}%
\definecolor{currentstroke}{rgb}{0.000000,0.000000,0.000000}%
\pgfsetstrokecolor{currentstroke}%
\pgfsetdash{}{0pt}%
\pgfpathmoveto{\pgfqpoint{1.112050in}{0.779340in}}%
\pgfpathlineto{\pgfqpoint{1.114210in}{0.796043in}}%
\pgfpathlineto{\pgfqpoint{1.116365in}{0.812715in}}%
\pgfpathlineto{\pgfqpoint{1.106130in}{0.813111in}}%
\pgfpathlineto{\pgfqpoint{1.095828in}{0.813224in}}%
\pgfpathlineto{\pgfqpoint{1.095366in}{0.796513in}}%
\pgfpathlineto{\pgfqpoint{1.094904in}{0.779770in}}%
\pgfpathlineto{\pgfqpoint{1.103506in}{0.779674in}}%
\pgfpathlineto{\pgfqpoint{1.112050in}{0.779340in}}%
\pgfpathclose%
\pgfusepath{fill}%
\end{pgfscope}%
\begin{pgfscope}%
\pgfpathrectangle{\pgfqpoint{0.407592in}{0.222904in}}{\pgfqpoint{1.411163in}{1.411163in}}%
\pgfusepath{clip}%
\pgfsetbuttcap%
\pgfsetroundjoin%
\definecolor{currentfill}{rgb}{0.749020,0.772549,0.792157}%
\pgfsetfillcolor{currentfill}%
\pgfsetfillopacity{0.160000}%
\pgfsetlinewidth{0.000000pt}%
\definecolor{currentstroke}{rgb}{0.000000,0.000000,0.000000}%
\pgfsetstrokecolor{currentstroke}%
\pgfsetdash{}{0pt}%
\pgfpathmoveto{\pgfqpoint{1.094904in}{0.779770in}}%
\pgfpathlineto{\pgfqpoint{1.095366in}{0.796513in}}%
\pgfpathlineto{\pgfqpoint{1.095828in}{0.813224in}}%
\pgfpathlineto{\pgfqpoint{1.085499in}{0.813054in}}%
\pgfpathlineto{\pgfqpoint{1.075188in}{0.812601in}}%
\pgfpathlineto{\pgfqpoint{1.076429in}{0.795939in}}%
\pgfpathlineto{\pgfqpoint{1.077671in}{0.779245in}}%
\pgfpathlineto{\pgfqpoint{1.086280in}{0.779626in}}%
\pgfpathlineto{\pgfqpoint{1.094904in}{0.779770in}}%
\pgfpathclose%
\pgfusepath{fill}%
\end{pgfscope}%
\begin{pgfscope}%
\pgfpathrectangle{\pgfqpoint{0.407592in}{0.222904in}}{\pgfqpoint{1.411163in}{1.411163in}}%
\pgfusepath{clip}%
\pgfsetbuttcap%
\pgfsetroundjoin%
\definecolor{currentfill}{rgb}{0.749020,0.772549,0.792157}%
\pgfsetfillcolor{currentfill}%
\pgfsetfillopacity{0.160000}%
\pgfsetlinewidth{0.000000pt}%
\definecolor{currentstroke}{rgb}{0.000000,0.000000,0.000000}%
\pgfsetstrokecolor{currentstroke}%
\pgfsetdash{}{0pt}%
\pgfpathmoveto{\pgfqpoint{0.952909in}{0.892541in}}%
\pgfpathlineto{\pgfqpoint{0.945370in}{0.908022in}}%
\pgfpathlineto{\pgfqpoint{0.937841in}{0.923485in}}%
\pgfpathlineto{\pgfqpoint{0.923991in}{0.919028in}}%
\pgfpathlineto{\pgfqpoint{0.910810in}{0.914218in}}%
\pgfpathlineto{\pgfqpoint{0.919699in}{0.899173in}}%
\pgfpathlineto{\pgfqpoint{0.928596in}{0.884114in}}%
\pgfpathlineto{\pgfqpoint{0.940448in}{0.888487in}}%
\pgfpathlineto{\pgfqpoint{0.952909in}{0.892541in}}%
\pgfpathclose%
\pgfusepath{fill}%
\end{pgfscope}%
\begin{pgfscope}%
\pgfpathrectangle{\pgfqpoint{0.407592in}{0.222904in}}{\pgfqpoint{1.411163in}{1.411163in}}%
\pgfusepath{clip}%
\pgfsetbuttcap%
\pgfsetroundjoin%
\definecolor{currentfill}{rgb}{0.749020,0.772549,0.792157}%
\pgfsetfillcolor{currentfill}%
\pgfsetfillopacity{0.160000}%
\pgfsetlinewidth{0.000000pt}%
\definecolor{currentstroke}{rgb}{0.000000,0.000000,0.000000}%
\pgfsetstrokecolor{currentstroke}%
\pgfsetdash{}{0pt}%
\pgfpathmoveto{\pgfqpoint{1.207460in}{0.831035in}}%
\pgfpathlineto{\pgfqpoint{1.215754in}{0.846635in}}%
\pgfpathlineto{\pgfqpoint{1.224037in}{0.862214in}}%
\pgfpathlineto{\pgfqpoint{1.213074in}{0.865494in}}%
\pgfpathlineto{\pgfqpoint{1.201605in}{0.868458in}}%
\pgfpathlineto{\pgfqpoint{1.194722in}{0.852520in}}%
\pgfpathlineto{\pgfqpoint{1.187830in}{0.836558in}}%
\pgfpathlineto{\pgfqpoint{1.197869in}{0.833936in}}%
\pgfpathlineto{\pgfqpoint{1.207460in}{0.831035in}}%
\pgfpathclose%
\pgfusepath{fill}%
\end{pgfscope}%
\begin{pgfscope}%
\pgfpathrectangle{\pgfqpoint{0.407592in}{0.222904in}}{\pgfqpoint{1.411163in}{1.411163in}}%
\pgfusepath{clip}%
\pgfsetbuttcap%
\pgfsetroundjoin%
\definecolor{currentfill}{rgb}{0.749020,0.772549,0.792157}%
\pgfsetfillcolor{currentfill}%
\pgfsetfillopacity{0.160000}%
\pgfsetlinewidth{0.000000pt}%
\definecolor{currentstroke}{rgb}{0.000000,0.000000,0.000000}%
\pgfsetstrokecolor{currentstroke}%
\pgfsetdash{}{0pt}%
\pgfpathmoveto{\pgfqpoint{1.347070in}{0.963096in}}%
\pgfpathlineto{\pgfqpoint{1.357682in}{0.977662in}}%
\pgfpathlineto{\pgfqpoint{1.368286in}{0.992218in}}%
\pgfpathlineto{\pgfqpoint{1.354319in}{0.999031in}}%
\pgfpathlineto{\pgfqpoint{1.339286in}{1.005435in}}%
\pgfpathlineto{\pgfqpoint{1.329826in}{0.990435in}}%
\pgfpathlineto{\pgfqpoint{1.320356in}{0.975420in}}%
\pgfpathlineto{\pgfqpoint{1.334209in}{0.969448in}}%
\pgfpathlineto{\pgfqpoint{1.347070in}{0.963096in}}%
\pgfpathclose%
\pgfusepath{fill}%
\end{pgfscope}%
\begin{pgfscope}%
\pgfpathrectangle{\pgfqpoint{0.407592in}{0.222904in}}{\pgfqpoint{1.411163in}{1.411163in}}%
\pgfusepath{clip}%
\pgfsetbuttcap%
\pgfsetroundjoin%
\definecolor{currentfill}{rgb}{0.749020,0.772549,0.792157}%
\pgfsetfillcolor{currentfill}%
\pgfsetfillopacity{0.160000}%
\pgfsetlinewidth{0.000000pt}%
\definecolor{currentstroke}{rgb}{0.000000,0.000000,0.000000}%
\pgfsetstrokecolor{currentstroke}%
\pgfsetdash{}{0pt}%
\pgfpathmoveto{\pgfqpoint{0.645855in}{1.250368in}}%
\pgfpathlineto{\pgfqpoint{0.635901in}{1.264701in}}%
\pgfpathlineto{\pgfqpoint{0.625953in}{1.279025in}}%
\pgfpathlineto{\pgfqpoint{0.600121in}{1.267376in}}%
\pgfpathlineto{\pgfqpoint{0.576138in}{1.255120in}}%
\pgfpathlineto{\pgfqpoint{0.587242in}{1.241197in}}%
\pgfpathlineto{\pgfqpoint{0.598349in}{1.227271in}}%
\pgfpathlineto{\pgfqpoint{0.621211in}{1.239111in}}%
\pgfpathlineto{\pgfqpoint{0.645855in}{1.250368in}}%
\pgfpathclose%
\pgfusepath{fill}%
\end{pgfscope}%
\begin{pgfscope}%
\pgfpathrectangle{\pgfqpoint{0.407592in}{0.222904in}}{\pgfqpoint{1.411163in}{1.411163in}}%
\pgfusepath{clip}%
\pgfsetbuttcap%
\pgfsetroundjoin%
\definecolor{currentfill}{rgb}{0.749020,0.772549,0.792157}%
\pgfsetfillcolor{currentfill}%
\pgfsetfillopacity{0.160000}%
\pgfsetlinewidth{0.000000pt}%
\definecolor{currentstroke}{rgb}{0.000000,0.000000,0.000000}%
\pgfsetstrokecolor{currentstroke}%
\pgfsetdash{}{0pt}%
\pgfpathmoveto{\pgfqpoint{1.174013in}{0.804561in}}%
\pgfpathlineto{\pgfqpoint{1.180927in}{0.820572in}}%
\pgfpathlineto{\pgfqpoint{1.187830in}{0.836558in}}%
\pgfpathlineto{\pgfqpoint{1.177386in}{0.838892in}}%
\pgfpathlineto{\pgfqpoint{1.166581in}{0.840928in}}%
\pgfpathlineto{\pgfqpoint{1.161187in}{0.824653in}}%
\pgfpathlineto{\pgfqpoint{1.155785in}{0.808349in}}%
\pgfpathlineto{\pgfqpoint{1.165055in}{0.806584in}}%
\pgfpathlineto{\pgfqpoint{1.174013in}{0.804561in}}%
\pgfpathclose%
\pgfusepath{fill}%
\end{pgfscope}%
\begin{pgfscope}%
\pgfpathrectangle{\pgfqpoint{0.407592in}{0.222904in}}{\pgfqpoint{1.411163in}{1.411163in}}%
\pgfusepath{clip}%
\pgfsetbuttcap%
\pgfsetroundjoin%
\definecolor{currentfill}{rgb}{0.749020,0.772549,0.792157}%
\pgfsetfillcolor{currentfill}%
\pgfsetfillopacity{0.160000}%
\pgfsetlinewidth{0.000000pt}%
\definecolor{currentstroke}{rgb}{0.000000,0.000000,0.000000}%
\pgfsetstrokecolor{currentstroke}%
\pgfsetdash{}{0pt}%
\pgfpathmoveto{\pgfqpoint{1.034960in}{0.808016in}}%
\pgfpathlineto{\pgfqpoint{1.030405in}{0.824294in}}%
\pgfpathlineto{\pgfqpoint{1.025857in}{0.840544in}}%
\pgfpathlineto{\pgfqpoint{1.014676in}{0.838447in}}%
\pgfpathlineto{\pgfqpoint{1.003802in}{0.836054in}}%
\pgfpathlineto{\pgfqpoint{1.009916in}{0.820102in}}%
\pgfpathlineto{\pgfqpoint{1.016039in}{0.804125in}}%
\pgfpathlineto{\pgfqpoint{1.025366in}{0.806199in}}%
\pgfpathlineto{\pgfqpoint{1.034960in}{0.808016in}}%
\pgfpathclose%
\pgfusepath{fill}%
\end{pgfscope}%
\begin{pgfscope}%
\pgfpathrectangle{\pgfqpoint{0.407592in}{0.222904in}}{\pgfqpoint{1.411163in}{1.411163in}}%
\pgfusepath{clip}%
\pgfsetbuttcap%
\pgfsetroundjoin%
\definecolor{currentfill}{rgb}{0.749020,0.772549,0.792157}%
\pgfsetfillcolor{currentfill}%
\pgfsetfillopacity{0.160000}%
\pgfsetlinewidth{0.000000pt}%
\definecolor{currentstroke}{rgb}{0.000000,0.000000,0.000000}%
\pgfsetstrokecolor{currentstroke}%
\pgfsetdash{}{0pt}%
\pgfpathmoveto{\pgfqpoint{1.128826in}{0.777963in}}%
\pgfpathlineto{\pgfqpoint{1.132647in}{0.794537in}}%
\pgfpathlineto{\pgfqpoint{1.136461in}{0.811081in}}%
\pgfpathlineto{\pgfqpoint{1.126489in}{0.812037in}}%
\pgfpathlineto{\pgfqpoint{1.116365in}{0.812715in}}%
\pgfpathlineto{\pgfqpoint{1.114210in}{0.796043in}}%
\pgfpathlineto{\pgfqpoint{1.112050in}{0.779340in}}%
\pgfpathlineto{\pgfqpoint{1.120503in}{0.778769in}}%
\pgfpathlineto{\pgfqpoint{1.128826in}{0.777963in}}%
\pgfpathclose%
\pgfusepath{fill}%
\end{pgfscope}%
\begin{pgfscope}%
\pgfpathrectangle{\pgfqpoint{0.407592in}{0.222904in}}{\pgfqpoint{1.411163in}{1.411163in}}%
\pgfusepath{clip}%
\pgfsetbuttcap%
\pgfsetroundjoin%
\definecolor{currentfill}{rgb}{0.749020,0.772549,0.792157}%
\pgfsetfillcolor{currentfill}%
\pgfsetfillopacity{0.160000}%
\pgfsetlinewidth{0.000000pt}%
\definecolor{currentstroke}{rgb}{0.000000,0.000000,0.000000}%
\pgfsetstrokecolor{currentstroke}%
\pgfsetdash{}{0pt}%
\pgfpathmoveto{\pgfqpoint{1.003802in}{0.836054in}}%
\pgfpathlineto{\pgfqpoint{0.997698in}{0.851984in}}%
\pgfpathlineto{\pgfqpoint{0.991602in}{0.867889in}}%
\pgfpathlineto{\pgfqpoint{0.979583in}{0.864861in}}%
\pgfpathlineto{\pgfqpoint{0.968014in}{0.861519in}}%
\pgfpathlineto{\pgfqpoint{0.975581in}{0.845980in}}%
\pgfpathlineto{\pgfqpoint{0.983157in}{0.830421in}}%
\pgfpathlineto{\pgfqpoint{0.993281in}{0.833376in}}%
\pgfpathlineto{\pgfqpoint{1.003802in}{0.836054in}}%
\pgfpathclose%
\pgfusepath{fill}%
\end{pgfscope}%
\begin{pgfscope}%
\pgfpathrectangle{\pgfqpoint{0.407592in}{0.222904in}}{\pgfqpoint{1.411163in}{1.411163in}}%
\pgfusepath{clip}%
\pgfsetbuttcap%
\pgfsetroundjoin%
\definecolor{currentfill}{rgb}{0.749020,0.772549,0.792157}%
\pgfsetfillcolor{currentfill}%
\pgfsetfillopacity{0.160000}%
\pgfsetlinewidth{0.000000pt}%
\definecolor{currentstroke}{rgb}{0.000000,0.000000,0.000000}%
\pgfsetstrokecolor{currentstroke}%
\pgfsetdash{}{0pt}%
\pgfpathmoveto{\pgfqpoint{1.077671in}{0.779245in}}%
\pgfpathlineto{\pgfqpoint{1.076429in}{0.795939in}}%
\pgfpathlineto{\pgfqpoint{1.075188in}{0.812601in}}%
\pgfpathlineto{\pgfqpoint{1.064937in}{0.811868in}}%
\pgfpathlineto{\pgfqpoint{1.054787in}{0.810856in}}%
\pgfpathlineto{\pgfqpoint{1.057711in}{0.794330in}}%
\pgfpathlineto{\pgfqpoint{1.060640in}{0.777773in}}%
\pgfpathlineto{\pgfqpoint{1.069113in}{0.778626in}}%
\pgfpathlineto{\pgfqpoint{1.077671in}{0.779245in}}%
\pgfpathclose%
\pgfusepath{fill}%
\end{pgfscope}%
\begin{pgfscope}%
\pgfpathrectangle{\pgfqpoint{0.407592in}{0.222904in}}{\pgfqpoint{1.411163in}{1.411163in}}%
\pgfusepath{clip}%
\pgfsetbuttcap%
\pgfsetroundjoin%
\definecolor{currentfill}{rgb}{0.749020,0.772549,0.792157}%
\pgfsetfillcolor{currentfill}%
\pgfsetfillopacity{0.160000}%
\pgfsetlinewidth{0.000000pt}%
\definecolor{currentstroke}{rgb}{0.000000,0.000000,0.000000}%
\pgfsetstrokecolor{currentstroke}%
\pgfsetdash{}{0pt}%
\pgfpathmoveto{\pgfqpoint{1.554624in}{1.173768in}}%
\pgfpathlineto{\pgfqpoint{1.566157in}{1.187756in}}%
\pgfpathlineto{\pgfqpoint{1.577687in}{1.201739in}}%
\pgfpathlineto{\pgfqpoint{1.558245in}{1.213043in}}%
\pgfpathlineto{\pgfqpoint{1.536980in}{1.223783in}}%
\pgfpathlineto{\pgfqpoint{1.526489in}{1.209383in}}%
\pgfpathlineto{\pgfqpoint{1.515992in}{1.194974in}}%
\pgfpathlineto{\pgfqpoint{1.536183in}{1.184641in}}%
\pgfpathlineto{\pgfqpoint{1.554624in}{1.173768in}}%
\pgfpathclose%
\pgfusepath{fill}%
\end{pgfscope}%
\begin{pgfscope}%
\pgfpathrectangle{\pgfqpoint{0.407592in}{0.222904in}}{\pgfqpoint{1.411163in}{1.411163in}}%
\pgfusepath{clip}%
\pgfsetbuttcap%
\pgfsetroundjoin%
\definecolor{currentfill}{rgb}{0.749020,0.772549,0.792157}%
\pgfsetfillcolor{currentfill}%
\pgfsetfillopacity{0.160000}%
\pgfsetlinewidth{0.000000pt}%
\definecolor{currentstroke}{rgb}{0.000000,0.000000,0.000000}%
\pgfsetstrokecolor{currentstroke}%
\pgfsetdash{}{0pt}%
\pgfpathmoveto{\pgfqpoint{0.875341in}{0.974253in}}%
\pgfpathlineto{\pgfqpoint{0.866495in}{0.989226in}}%
\pgfpathlineto{\pgfqpoint{0.857658in}{1.004184in}}%
\pgfpathlineto{\pgfqpoint{0.841397in}{0.997696in}}%
\pgfpathlineto{\pgfqpoint{0.826126in}{0.990804in}}%
\pgfpathlineto{\pgfqpoint{0.836202in}{0.976296in}}%
\pgfpathlineto{\pgfqpoint{0.846285in}{0.961778in}}%
\pgfpathlineto{\pgfqpoint{0.860352in}{0.968203in}}%
\pgfpathlineto{\pgfqpoint{0.875341in}{0.974253in}}%
\pgfpathclose%
\pgfusepath{fill}%
\end{pgfscope}%
\begin{pgfscope}%
\pgfpathrectangle{\pgfqpoint{0.407592in}{0.222904in}}{\pgfqpoint{1.411163in}{1.411163in}}%
\pgfusepath{clip}%
\pgfsetbuttcap%
\pgfsetroundjoin%
\definecolor{currentfill}{rgb}{0.749020,0.772549,0.792157}%
\pgfsetfillcolor{currentfill}%
\pgfsetfillopacity{0.160000}%
\pgfsetlinewidth{0.000000pt}%
\definecolor{currentstroke}{rgb}{0.000000,0.000000,0.000000}%
\pgfsetstrokecolor{currentstroke}%
\pgfsetdash{}{0pt}%
\pgfpathmoveto{\pgfqpoint{0.665782in}{1.221675in}}%
\pgfpathlineto{\pgfqpoint{0.655815in}{1.236026in}}%
\pgfpathlineto{\pgfqpoint{0.645855in}{1.250368in}}%
\pgfpathlineto{\pgfqpoint{0.621211in}{1.239111in}}%
\pgfpathlineto{\pgfqpoint{0.598349in}{1.227271in}}%
\pgfpathlineto{\pgfqpoint{0.609459in}{1.213340in}}%
\pgfpathlineto{\pgfqpoint{0.620573in}{1.199406in}}%
\pgfpathlineto{\pgfqpoint{0.642321in}{1.210819in}}%
\pgfpathlineto{\pgfqpoint{0.665782in}{1.221675in}}%
\pgfpathclose%
\pgfusepath{fill}%
\end{pgfscope}%
\begin{pgfscope}%
\pgfpathrectangle{\pgfqpoint{0.407592in}{0.222904in}}{\pgfqpoint{1.411163in}{1.411163in}}%
\pgfusepath{clip}%
\pgfsetbuttcap%
\pgfsetroundjoin%
\definecolor{currentfill}{rgb}{0.749020,0.772549,0.792157}%
\pgfsetfillcolor{currentfill}%
\pgfsetfillopacity{0.160000}%
\pgfsetlinewidth{0.000000pt}%
\definecolor{currentstroke}{rgb}{0.000000,0.000000,0.000000}%
\pgfsetstrokecolor{currentstroke}%
\pgfsetdash{}{0pt}%
\pgfpathmoveto{\pgfqpoint{1.531547in}{1.145779in}}%
\pgfpathlineto{\pgfqpoint{1.543088in}{1.159776in}}%
\pgfpathlineto{\pgfqpoint{1.554624in}{1.173768in}}%
\pgfpathlineto{\pgfqpoint{1.536183in}{1.184641in}}%
\pgfpathlineto{\pgfqpoint{1.515992in}{1.194974in}}%
\pgfpathlineto{\pgfqpoint{1.505488in}{1.180555in}}%
\pgfpathlineto{\pgfqpoint{1.494977in}{1.166126in}}%
\pgfpathlineto{\pgfqpoint{1.514098in}{1.156210in}}%
\pgfpathlineto{\pgfqpoint{1.531547in}{1.145779in}}%
\pgfpathclose%
\pgfusepath{fill}%
\end{pgfscope}%
\begin{pgfscope}%
\pgfpathrectangle{\pgfqpoint{0.407592in}{0.222904in}}{\pgfqpoint{1.411163in}{1.411163in}}%
\pgfusepath{clip}%
\pgfsetbuttcap%
\pgfsetroundjoin%
\definecolor{currentfill}{rgb}{0.749020,0.772549,0.792157}%
\pgfsetfillcolor{currentfill}%
\pgfsetfillopacity{0.160000}%
\pgfsetlinewidth{0.000000pt}%
\definecolor{currentstroke}{rgb}{0.000000,0.000000,0.000000}%
\pgfsetstrokecolor{currentstroke}%
\pgfsetdash{}{0pt}%
\pgfpathmoveto{\pgfqpoint{1.144952in}{0.775660in}}%
\pgfpathlineto{\pgfqpoint{1.150373in}{0.792019in}}%
\pgfpathlineto{\pgfqpoint{1.155785in}{0.808349in}}%
\pgfpathlineto{\pgfqpoint{1.146240in}{0.809851in}}%
\pgfpathlineto{\pgfqpoint{1.136461in}{0.811081in}}%
\pgfpathlineto{\pgfqpoint{1.132647in}{0.794537in}}%
\pgfpathlineto{\pgfqpoint{1.128826in}{0.777963in}}%
\pgfpathlineto{\pgfqpoint{1.136988in}{0.776925in}}%
\pgfpathlineto{\pgfqpoint{1.144952in}{0.775660in}}%
\pgfpathclose%
\pgfusepath{fill}%
\end{pgfscope}%
\begin{pgfscope}%
\pgfpathrectangle{\pgfqpoint{0.407592in}{0.222904in}}{\pgfqpoint{1.411163in}{1.411163in}}%
\pgfusepath{clip}%
\pgfsetbuttcap%
\pgfsetroundjoin%
\definecolor{currentfill}{rgb}{0.749020,0.772549,0.792157}%
\pgfsetfillcolor{currentfill}%
\pgfsetfillopacity{0.160000}%
\pgfsetlinewidth{0.000000pt}%
\definecolor{currentstroke}{rgb}{0.000000,0.000000,0.000000}%
\pgfsetstrokecolor{currentstroke}%
\pgfsetdash{}{0pt}%
\pgfpathmoveto{\pgfqpoint{1.244259in}{0.854761in}}%
\pgfpathlineto{\pgfqpoint{1.253804in}{0.869896in}}%
\pgfpathlineto{\pgfqpoint{1.263340in}{0.885016in}}%
\pgfpathlineto{\pgfqpoint{1.252289in}{0.889326in}}%
\pgfpathlineto{\pgfqpoint{1.240571in}{0.893314in}}%
\pgfpathlineto{\pgfqpoint{1.232309in}{0.877774in}}%
\pgfpathlineto{\pgfqpoint{1.224037in}{0.862214in}}%
\pgfpathlineto{\pgfqpoint{1.234447in}{0.858632in}}%
\pgfpathlineto{\pgfqpoint{1.244259in}{0.854761in}}%
\pgfpathclose%
\pgfusepath{fill}%
\end{pgfscope}%
\begin{pgfscope}%
\pgfpathrectangle{\pgfqpoint{0.407592in}{0.222904in}}{\pgfqpoint{1.411163in}{1.411163in}}%
\pgfusepath{clip}%
\pgfsetbuttcap%
\pgfsetroundjoin%
\definecolor{currentfill}{rgb}{0.749020,0.772549,0.792157}%
\pgfsetfillcolor{currentfill}%
\pgfsetfillopacity{0.160000}%
\pgfsetlinewidth{0.000000pt}%
\definecolor{currentstroke}{rgb}{0.000000,0.000000,0.000000}%
\pgfsetstrokecolor{currentstroke}%
\pgfsetdash{}{0pt}%
\pgfpathmoveto{\pgfqpoint{1.060640in}{0.777773in}}%
\pgfpathlineto{\pgfqpoint{1.057711in}{0.794330in}}%
\pgfpathlineto{\pgfqpoint{1.054787in}{0.810856in}}%
\pgfpathlineto{\pgfqpoint{1.044781in}{0.809571in}}%
\pgfpathlineto{\pgfqpoint{1.034960in}{0.808016in}}%
\pgfpathlineto{\pgfqpoint{1.039524in}{0.791712in}}%
\pgfpathlineto{\pgfqpoint{1.044094in}{0.775379in}}%
\pgfpathlineto{\pgfqpoint{1.052289in}{0.776690in}}%
\pgfpathlineto{\pgfqpoint{1.060640in}{0.777773in}}%
\pgfpathclose%
\pgfusepath{fill}%
\end{pgfscope}%
\begin{pgfscope}%
\pgfpathrectangle{\pgfqpoint{0.407592in}{0.222904in}}{\pgfqpoint{1.411163in}{1.411163in}}%
\pgfusepath{clip}%
\pgfsetbuttcap%
\pgfsetroundjoin%
\definecolor{currentfill}{rgb}{0.749020,0.772549,0.792157}%
\pgfsetfillcolor{currentfill}%
\pgfsetfillopacity{0.160000}%
\pgfsetlinewidth{0.000000pt}%
\definecolor{currentstroke}{rgb}{0.000000,0.000000,0.000000}%
\pgfsetstrokecolor{currentstroke}%
\pgfsetdash{}{0pt}%
\pgfpathmoveto{\pgfqpoint{1.325827in}{0.933935in}}%
\pgfpathlineto{\pgfqpoint{1.336452in}{0.948521in}}%
\pgfpathlineto{\pgfqpoint{1.347070in}{0.963096in}}%
\pgfpathlineto{\pgfqpoint{1.334209in}{0.969448in}}%
\pgfpathlineto{\pgfqpoint{1.320356in}{0.975420in}}%
\pgfpathlineto{\pgfqpoint{1.310877in}{0.960390in}}%
\pgfpathlineto{\pgfqpoint{1.301388in}{0.945345in}}%
\pgfpathlineto{\pgfqpoint{1.314065in}{0.939815in}}%
\pgfpathlineto{\pgfqpoint{1.325827in}{0.933935in}}%
\pgfpathclose%
\pgfusepath{fill}%
\end{pgfscope}%
\begin{pgfscope}%
\pgfpathrectangle{\pgfqpoint{0.407592in}{0.222904in}}{\pgfqpoint{1.411163in}{1.411163in}}%
\pgfusepath{clip}%
\pgfsetbuttcap%
\pgfsetroundjoin%
\definecolor{currentfill}{rgb}{0.749020,0.772549,0.792157}%
\pgfsetfillcolor{currentfill}%
\pgfsetfillopacity{0.160000}%
\pgfsetlinewidth{0.000000pt}%
\definecolor{currentstroke}{rgb}{0.000000,0.000000,0.000000}%
\pgfsetstrokecolor{currentstroke}%
\pgfsetdash{}{0pt}%
\pgfpathmoveto{\pgfqpoint{0.685736in}{1.192946in}}%
\pgfpathlineto{\pgfqpoint{0.675756in}{1.207315in}}%
\pgfpathlineto{\pgfqpoint{0.665782in}{1.221675in}}%
\pgfpathlineto{\pgfqpoint{0.642321in}{1.210819in}}%
\pgfpathlineto{\pgfqpoint{0.620573in}{1.199406in}}%
\pgfpathlineto{\pgfqpoint{0.631689in}{1.185467in}}%
\pgfpathlineto{\pgfqpoint{0.642809in}{1.171525in}}%
\pgfpathlineto{\pgfqpoint{0.663451in}{1.182502in}}%
\pgfpathlineto{\pgfqpoint{0.685736in}{1.192946in}}%
\pgfpathclose%
\pgfusepath{fill}%
\end{pgfscope}%
\begin{pgfscope}%
\pgfpathrectangle{\pgfqpoint{0.407592in}{0.222904in}}{\pgfqpoint{1.411163in}{1.411163in}}%
\pgfusepath{clip}%
\pgfsetbuttcap%
\pgfsetroundjoin%
\definecolor{currentfill}{rgb}{0.749020,0.772549,0.792157}%
\pgfsetfillcolor{currentfill}%
\pgfsetfillopacity{0.160000}%
\pgfsetlinewidth{0.000000pt}%
\definecolor{currentstroke}{rgb}{0.000000,0.000000,0.000000}%
\pgfsetstrokecolor{currentstroke}%
\pgfsetdash{}{0pt}%
\pgfpathmoveto{\pgfqpoint{0.968014in}{0.861519in}}%
\pgfpathlineto{\pgfqpoint{0.960457in}{0.877040in}}%
\pgfpathlineto{\pgfqpoint{0.952909in}{0.892541in}}%
\pgfpathlineto{\pgfqpoint{0.940448in}{0.888487in}}%
\pgfpathlineto{\pgfqpoint{0.928596in}{0.884114in}}%
\pgfpathlineto{\pgfqpoint{0.937502in}{0.869040in}}%
\pgfpathlineto{\pgfqpoint{0.946416in}{0.853951in}}%
\pgfpathlineto{\pgfqpoint{0.956943in}{0.857878in}}%
\pgfpathlineto{\pgfqpoint{0.968014in}{0.861519in}}%
\pgfpathclose%
\pgfusepath{fill}%
\end{pgfscope}%
\begin{pgfscope}%
\pgfpathrectangle{\pgfqpoint{0.407592in}{0.222904in}}{\pgfqpoint{1.411163in}{1.411163in}}%
\pgfusepath{clip}%
\pgfsetbuttcap%
\pgfsetroundjoin%
\definecolor{currentfill}{rgb}{0.749020,0.772549,0.792157}%
\pgfsetfillcolor{currentfill}%
\pgfsetfillopacity{0.160000}%
\pgfsetlinewidth{0.000000pt}%
\definecolor{currentstroke}{rgb}{0.000000,0.000000,0.000000}%
\pgfsetstrokecolor{currentstroke}%
\pgfsetdash{}{0pt}%
\pgfpathmoveto{\pgfqpoint{0.893058in}{0.944265in}}%
\pgfpathlineto{\pgfqpoint{0.884195in}{0.959266in}}%
\pgfpathlineto{\pgfqpoint{0.875341in}{0.974253in}}%
\pgfpathlineto{\pgfqpoint{0.860352in}{0.968203in}}%
\pgfpathlineto{\pgfqpoint{0.846285in}{0.961778in}}%
\pgfpathlineto{\pgfqpoint{0.856374in}{0.947251in}}%
\pgfpathlineto{\pgfqpoint{0.866469in}{0.932715in}}%
\pgfpathlineto{\pgfqpoint{0.879338in}{0.938663in}}%
\pgfpathlineto{\pgfqpoint{0.893058in}{0.944265in}}%
\pgfpathclose%
\pgfusepath{fill}%
\end{pgfscope}%
\begin{pgfscope}%
\pgfpathrectangle{\pgfqpoint{0.407592in}{0.222904in}}{\pgfqpoint{1.411163in}{1.411163in}}%
\pgfusepath{clip}%
\pgfsetbuttcap%
\pgfsetroundjoin%
\definecolor{currentfill}{rgb}{0.749020,0.772549,0.792157}%
\pgfsetfillcolor{currentfill}%
\pgfsetfillopacity{0.160000}%
\pgfsetlinewidth{0.000000pt}%
\definecolor{currentstroke}{rgb}{0.000000,0.000000,0.000000}%
\pgfsetstrokecolor{currentstroke}%
\pgfsetdash{}{0pt}%
\pgfpathmoveto{\pgfqpoint{1.508456in}{1.117773in}}%
\pgfpathlineto{\pgfqpoint{1.520003in}{1.131778in}}%
\pgfpathlineto{\pgfqpoint{1.531547in}{1.145779in}}%
\pgfpathlineto{\pgfqpoint{1.514098in}{1.156210in}}%
\pgfpathlineto{\pgfqpoint{1.494977in}{1.166126in}}%
\pgfpathlineto{\pgfqpoint{1.484458in}{1.151687in}}%
\pgfpathlineto{\pgfqpoint{1.473933in}{1.137239in}}%
\pgfpathlineto{\pgfqpoint{1.491992in}{1.127750in}}%
\pgfpathlineto{\pgfqpoint{1.508456in}{1.117773in}}%
\pgfpathclose%
\pgfusepath{fill}%
\end{pgfscope}%
\begin{pgfscope}%
\pgfpathrectangle{\pgfqpoint{0.407592in}{0.222904in}}{\pgfqpoint{1.411163in}{1.411163in}}%
\pgfusepath{clip}%
\pgfsetbuttcap%
\pgfsetroundjoin%
\definecolor{currentfill}{rgb}{0.749020,0.772549,0.792157}%
\pgfsetfillcolor{currentfill}%
\pgfsetfillopacity{0.160000}%
\pgfsetlinewidth{0.000000pt}%
\definecolor{currentstroke}{rgb}{0.000000,0.000000,0.000000}%
\pgfsetstrokecolor{currentstroke}%
\pgfsetdash{}{0pt}%
\pgfpathmoveto{\pgfqpoint{1.190842in}{0.799776in}}%
\pgfpathlineto{\pgfqpoint{1.199156in}{0.815416in}}%
\pgfpathlineto{\pgfqpoint{1.207460in}{0.831035in}}%
\pgfpathlineto{\pgfqpoint{1.197869in}{0.833936in}}%
\pgfpathlineto{\pgfqpoint{1.187830in}{0.836558in}}%
\pgfpathlineto{\pgfqpoint{1.180927in}{0.820572in}}%
\pgfpathlineto{\pgfqpoint{1.174013in}{0.804561in}}%
\pgfpathlineto{\pgfqpoint{1.182621in}{0.802289in}}%
\pgfpathlineto{\pgfqpoint{1.190842in}{0.799776in}}%
\pgfpathclose%
\pgfusepath{fill}%
\end{pgfscope}%
\begin{pgfscope}%
\pgfpathrectangle{\pgfqpoint{0.407592in}{0.222904in}}{\pgfqpoint{1.411163in}{1.411163in}}%
\pgfusepath{clip}%
\pgfsetbuttcap%
\pgfsetroundjoin%
\definecolor{currentfill}{rgb}{0.749020,0.772549,0.792157}%
\pgfsetfillcolor{currentfill}%
\pgfsetfillopacity{0.160000}%
\pgfsetlinewidth{0.000000pt}%
\definecolor{currentstroke}{rgb}{0.000000,0.000000,0.000000}%
\pgfsetstrokecolor{currentstroke}%
\pgfsetdash{}{0pt}%
\pgfpathmoveto{\pgfqpoint{0.705714in}{1.164179in}}%
\pgfpathlineto{\pgfqpoint{0.695722in}{1.178567in}}%
\pgfpathlineto{\pgfqpoint{0.685736in}{1.192946in}}%
\pgfpathlineto{\pgfqpoint{0.663451in}{1.182502in}}%
\pgfpathlineto{\pgfqpoint{0.642809in}{1.171525in}}%
\pgfpathlineto{\pgfqpoint{0.653932in}{1.157578in}}%
\pgfpathlineto{\pgfqpoint{0.665058in}{1.143628in}}%
\pgfpathlineto{\pgfqpoint{0.684600in}{1.154157in}}%
\pgfpathlineto{\pgfqpoint{0.705714in}{1.164179in}}%
\pgfpathclose%
\pgfusepath{fill}%
\end{pgfscope}%
\begin{pgfscope}%
\pgfpathrectangle{\pgfqpoint{0.407592in}{0.222904in}}{\pgfqpoint{1.411163in}{1.411163in}}%
\pgfusepath{clip}%
\pgfsetbuttcap%
\pgfsetroundjoin%
\definecolor{currentfill}{rgb}{0.749020,0.772549,0.792157}%
\pgfsetfillcolor{currentfill}%
\pgfsetfillopacity{0.160000}%
\pgfsetlinewidth{0.000000pt}%
\definecolor{currentstroke}{rgb}{0.000000,0.000000,0.000000}%
\pgfsetstrokecolor{currentstroke}%
\pgfsetdash{}{0pt}%
\pgfpathmoveto{\pgfqpoint{1.016039in}{0.804125in}}%
\pgfpathlineto{\pgfqpoint{1.009916in}{0.820102in}}%
\pgfpathlineto{\pgfqpoint{1.003802in}{0.836054in}}%
\pgfpathlineto{\pgfqpoint{0.993281in}{0.833376in}}%
\pgfpathlineto{\pgfqpoint{0.983157in}{0.830421in}}%
\pgfpathlineto{\pgfqpoint{0.990743in}{0.814842in}}%
\pgfpathlineto{\pgfqpoint{0.998338in}{0.799244in}}%
\pgfpathlineto{\pgfqpoint{1.007017in}{0.801804in}}%
\pgfpathlineto{\pgfqpoint{1.016039in}{0.804125in}}%
\pgfpathclose%
\pgfusepath{fill}%
\end{pgfscope}%
\begin{pgfscope}%
\pgfpathrectangle{\pgfqpoint{0.407592in}{0.222904in}}{\pgfqpoint{1.411163in}{1.411163in}}%
\pgfusepath{clip}%
\pgfsetbuttcap%
\pgfsetroundjoin%
\definecolor{currentfill}{rgb}{0.749020,0.772549,0.792157}%
\pgfsetfillcolor{currentfill}%
\pgfsetfillopacity{0.160000}%
\pgfsetlinewidth{0.000000pt}%
\definecolor{currentstroke}{rgb}{0.000000,0.000000,0.000000}%
\pgfsetstrokecolor{currentstroke}%
\pgfsetdash{}{0pt}%
\pgfpathmoveto{\pgfqpoint{1.160155in}{0.772468in}}%
\pgfpathlineto{\pgfqpoint{1.167090in}{0.788527in}}%
\pgfpathlineto{\pgfqpoint{1.174013in}{0.804561in}}%
\pgfpathlineto{\pgfqpoint{1.165055in}{0.806584in}}%
\pgfpathlineto{\pgfqpoint{1.155785in}{0.808349in}}%
\pgfpathlineto{\pgfqpoint{1.150373in}{0.792019in}}%
\pgfpathlineto{\pgfqpoint{1.144952in}{0.775660in}}%
\pgfpathlineto{\pgfqpoint{1.152685in}{0.774172in}}%
\pgfpathlineto{\pgfqpoint{1.160155in}{0.772468in}}%
\pgfpathclose%
\pgfusepath{fill}%
\end{pgfscope}%
\begin{pgfscope}%
\pgfpathrectangle{\pgfqpoint{0.407592in}{0.222904in}}{\pgfqpoint{1.411163in}{1.411163in}}%
\pgfusepath{clip}%
\pgfsetbuttcap%
\pgfsetroundjoin%
\definecolor{currentfill}{rgb}{0.749020,0.772549,0.792157}%
\pgfsetfillcolor{currentfill}%
\pgfsetfillopacity{0.160000}%
\pgfsetlinewidth{0.000000pt}%
\definecolor{currentstroke}{rgb}{0.000000,0.000000,0.000000}%
\pgfsetstrokecolor{currentstroke}%
\pgfsetdash{}{0pt}%
\pgfpathmoveto{\pgfqpoint{1.304554in}{0.904734in}}%
\pgfpathlineto{\pgfqpoint{1.315194in}{0.919339in}}%
\pgfpathlineto{\pgfqpoint{1.325827in}{0.933935in}}%
\pgfpathlineto{\pgfqpoint{1.314065in}{0.939815in}}%
\pgfpathlineto{\pgfqpoint{1.301388in}{0.945345in}}%
\pgfpathlineto{\pgfqpoint{1.291891in}{0.930285in}}%
\pgfpathlineto{\pgfqpoint{1.282383in}{0.915211in}}%
\pgfpathlineto{\pgfqpoint{1.293888in}{0.910132in}}%
\pgfpathlineto{\pgfqpoint{1.304554in}{0.904734in}}%
\pgfpathclose%
\pgfusepath{fill}%
\end{pgfscope}%
\begin{pgfscope}%
\pgfpathrectangle{\pgfqpoint{0.407592in}{0.222904in}}{\pgfqpoint{1.411163in}{1.411163in}}%
\pgfusepath{clip}%
\pgfsetbuttcap%
\pgfsetroundjoin%
\definecolor{currentfill}{rgb}{0.749020,0.772549,0.792157}%
\pgfsetfillcolor{currentfill}%
\pgfsetfillopacity{0.160000}%
\pgfsetlinewidth{0.000000pt}%
\definecolor{currentstroke}{rgb}{0.000000,0.000000,0.000000}%
\pgfsetstrokecolor{currentstroke}%
\pgfsetdash{}{0pt}%
\pgfpathmoveto{\pgfqpoint{1.485349in}{1.089748in}}%
\pgfpathlineto{\pgfqpoint{1.496904in}{1.103763in}}%
\pgfpathlineto{\pgfqpoint{1.508456in}{1.117773in}}%
\pgfpathlineto{\pgfqpoint{1.491992in}{1.127750in}}%
\pgfpathlineto{\pgfqpoint{1.473933in}{1.137239in}}%
\pgfpathlineto{\pgfqpoint{1.463400in}{1.122781in}}%
\pgfpathlineto{\pgfqpoint{1.452860in}{1.108313in}}%
\pgfpathlineto{\pgfqpoint{1.469863in}{1.099262in}}%
\pgfpathlineto{\pgfqpoint{1.485349in}{1.089748in}}%
\pgfpathclose%
\pgfusepath{fill}%
\end{pgfscope}%
\begin{pgfscope}%
\pgfpathrectangle{\pgfqpoint{0.407592in}{0.222904in}}{\pgfqpoint{1.411163in}{1.411163in}}%
\pgfusepath{clip}%
\pgfsetbuttcap%
\pgfsetroundjoin%
\definecolor{currentfill}{rgb}{0.749020,0.772549,0.792157}%
\pgfsetfillcolor{currentfill}%
\pgfsetfillopacity{0.160000}%
\pgfsetlinewidth{0.000000pt}%
\definecolor{currentstroke}{rgb}{0.000000,0.000000,0.000000}%
\pgfsetstrokecolor{currentstroke}%
\pgfsetdash{}{0pt}%
\pgfpathmoveto{\pgfqpoint{1.044094in}{0.775379in}}%
\pgfpathlineto{\pgfqpoint{1.039524in}{0.791712in}}%
\pgfpathlineto{\pgfqpoint{1.034960in}{0.808016in}}%
\pgfpathlineto{\pgfqpoint{1.025366in}{0.806199in}}%
\pgfpathlineto{\pgfqpoint{1.016039in}{0.804125in}}%
\pgfpathlineto{\pgfqpoint{1.022170in}{0.788125in}}%
\pgfpathlineto{\pgfqpoint{1.028311in}{0.772101in}}%
\pgfpathlineto{\pgfqpoint{1.036091in}{0.773848in}}%
\pgfpathlineto{\pgfqpoint{1.044094in}{0.775379in}}%
\pgfpathclose%
\pgfusepath{fill}%
\end{pgfscope}%
\begin{pgfscope}%
\pgfpathrectangle{\pgfqpoint{0.407592in}{0.222904in}}{\pgfqpoint{1.411163in}{1.411163in}}%
\pgfusepath{clip}%
\pgfsetbuttcap%
\pgfsetroundjoin%
\definecolor{currentfill}{rgb}{0.749020,0.772549,0.792157}%
\pgfsetfillcolor{currentfill}%
\pgfsetfillopacity{0.160000}%
\pgfsetlinewidth{0.000000pt}%
\definecolor{currentstroke}{rgb}{0.000000,0.000000,0.000000}%
\pgfsetstrokecolor{currentstroke}%
\pgfsetdash{}{0pt}%
\pgfpathmoveto{\pgfqpoint{1.107720in}{0.745838in}}%
\pgfpathlineto{\pgfqpoint{1.109887in}{0.762605in}}%
\pgfpathlineto{\pgfqpoint{1.112050in}{0.779340in}}%
\pgfpathlineto{\pgfqpoint{1.103506in}{0.779674in}}%
\pgfpathlineto{\pgfqpoint{1.094904in}{0.779770in}}%
\pgfpathlineto{\pgfqpoint{1.094440in}{0.762994in}}%
\pgfpathlineto{\pgfqpoint{1.093976in}{0.746185in}}%
\pgfpathlineto{\pgfqpoint{1.100871in}{0.746108in}}%
\pgfpathlineto{\pgfqpoint{1.107720in}{0.745838in}}%
\pgfpathclose%
\pgfusepath{fill}%
\end{pgfscope}%
\begin{pgfscope}%
\pgfpathrectangle{\pgfqpoint{0.407592in}{0.222904in}}{\pgfqpoint{1.411163in}{1.411163in}}%
\pgfusepath{clip}%
\pgfsetbuttcap%
\pgfsetroundjoin%
\definecolor{currentfill}{rgb}{0.749020,0.772549,0.792157}%
\pgfsetfillcolor{currentfill}%
\pgfsetfillopacity{0.160000}%
\pgfsetlinewidth{0.000000pt}%
\definecolor{currentstroke}{rgb}{0.000000,0.000000,0.000000}%
\pgfsetstrokecolor{currentstroke}%
\pgfsetdash{}{0pt}%
\pgfpathmoveto{\pgfqpoint{1.093976in}{0.746185in}}%
\pgfpathlineto{\pgfqpoint{1.094440in}{0.762994in}}%
\pgfpathlineto{\pgfqpoint{1.094904in}{0.779770in}}%
\pgfpathlineto{\pgfqpoint{1.086280in}{0.779626in}}%
\pgfpathlineto{\pgfqpoint{1.077671in}{0.779245in}}%
\pgfpathlineto{\pgfqpoint{1.078916in}{0.762518in}}%
\pgfpathlineto{\pgfqpoint{1.080164in}{0.745760in}}%
\pgfpathlineto{\pgfqpoint{1.087064in}{0.746069in}}%
\pgfpathlineto{\pgfqpoint{1.093976in}{0.746185in}}%
\pgfpathclose%
\pgfusepath{fill}%
\end{pgfscope}%
\begin{pgfscope}%
\pgfpathrectangle{\pgfqpoint{0.407592in}{0.222904in}}{\pgfqpoint{1.411163in}{1.411163in}}%
\pgfusepath{clip}%
\pgfsetbuttcap%
\pgfsetroundjoin%
\definecolor{currentfill}{rgb}{0.749020,0.772549,0.792157}%
\pgfsetfillcolor{currentfill}%
\pgfsetfillopacity{0.160000}%
\pgfsetlinewidth{0.000000pt}%
\definecolor{currentstroke}{rgb}{0.000000,0.000000,0.000000}%
\pgfsetstrokecolor{currentstroke}%
\pgfsetdash{}{0pt}%
\pgfpathmoveto{\pgfqpoint{0.725719in}{1.135376in}}%
\pgfpathlineto{\pgfqpoint{0.715713in}{1.149782in}}%
\pgfpathlineto{\pgfqpoint{0.705714in}{1.164179in}}%
\pgfpathlineto{\pgfqpoint{0.684600in}{1.154157in}}%
\pgfpathlineto{\pgfqpoint{0.665058in}{1.143628in}}%
\pgfpathlineto{\pgfqpoint{0.676188in}{1.129674in}}%
\pgfpathlineto{\pgfqpoint{0.687320in}{1.115715in}}%
\pgfpathlineto{\pgfqpoint{0.705769in}{1.125787in}}%
\pgfpathlineto{\pgfqpoint{0.725719in}{1.135376in}}%
\pgfpathclose%
\pgfusepath{fill}%
\end{pgfscope}%
\begin{pgfscope}%
\pgfpathrectangle{\pgfqpoint{0.407592in}{0.222904in}}{\pgfqpoint{1.411163in}{1.411163in}}%
\pgfusepath{clip}%
\pgfsetbuttcap%
\pgfsetroundjoin%
\definecolor{currentfill}{rgb}{0.749020,0.772549,0.792157}%
\pgfsetfillcolor{currentfill}%
\pgfsetfillopacity{0.160000}%
\pgfsetlinewidth{0.000000pt}%
\definecolor{currentstroke}{rgb}{0.000000,0.000000,0.000000}%
\pgfsetstrokecolor{currentstroke}%
\pgfsetdash{}{0pt}%
\pgfpathmoveto{\pgfqpoint{1.225140in}{0.824446in}}%
\pgfpathlineto{\pgfqpoint{1.234704in}{0.839611in}}%
\pgfpathlineto{\pgfqpoint{1.244259in}{0.854761in}}%
\pgfpathlineto{\pgfqpoint{1.234447in}{0.858632in}}%
\pgfpathlineto{\pgfqpoint{1.224037in}{0.862214in}}%
\pgfpathlineto{\pgfqpoint{1.215754in}{0.846635in}}%
\pgfpathlineto{\pgfqpoint{1.207460in}{0.831035in}}%
\pgfpathlineto{\pgfqpoint{1.216564in}{0.827868in}}%
\pgfpathlineto{\pgfqpoint{1.225140in}{0.824446in}}%
\pgfpathclose%
\pgfusepath{fill}%
\end{pgfscope}%
\begin{pgfscope}%
\pgfpathrectangle{\pgfqpoint{0.407592in}{0.222904in}}{\pgfqpoint{1.411163in}{1.411163in}}%
\pgfusepath{clip}%
\pgfsetbuttcap%
\pgfsetroundjoin%
\definecolor{currentfill}{rgb}{0.749020,0.772549,0.792157}%
\pgfsetfillcolor{currentfill}%
\pgfsetfillopacity{0.160000}%
\pgfsetlinewidth{0.000000pt}%
\definecolor{currentstroke}{rgb}{0.000000,0.000000,0.000000}%
\pgfsetstrokecolor{currentstroke}%
\pgfsetdash{}{0pt}%
\pgfpathmoveto{\pgfqpoint{0.910810in}{0.914218in}}%
\pgfpathlineto{\pgfqpoint{0.901930in}{0.929249in}}%
\pgfpathlineto{\pgfqpoint{0.893058in}{0.944265in}}%
\pgfpathlineto{\pgfqpoint{0.879338in}{0.938663in}}%
\pgfpathlineto{\pgfqpoint{0.866469in}{0.932715in}}%
\pgfpathlineto{\pgfqpoint{0.876572in}{0.918169in}}%
\pgfpathlineto{\pgfqpoint{0.886680in}{0.903614in}}%
\pgfpathlineto{\pgfqpoint{0.898355in}{0.909074in}}%
\pgfpathlineto{\pgfqpoint{0.910810in}{0.914218in}}%
\pgfpathclose%
\pgfusepath{fill}%
\end{pgfscope}%
\begin{pgfscope}%
\pgfpathrectangle{\pgfqpoint{0.407592in}{0.222904in}}{\pgfqpoint{1.411163in}{1.411163in}}%
\pgfusepath{clip}%
\pgfsetbuttcap%
\pgfsetroundjoin%
\definecolor{currentfill}{rgb}{0.749020,0.772549,0.792157}%
\pgfsetfillcolor{currentfill}%
\pgfsetfillopacity{0.160000}%
\pgfsetlinewidth{0.000000pt}%
\definecolor{currentstroke}{rgb}{0.000000,0.000000,0.000000}%
\pgfsetstrokecolor{currentstroke}%
\pgfsetdash{}{0pt}%
\pgfpathmoveto{\pgfqpoint{1.121164in}{0.744723in}}%
\pgfpathlineto{\pgfqpoint{1.124999in}{0.761358in}}%
\pgfpathlineto{\pgfqpoint{1.128826in}{0.777963in}}%
\pgfpathlineto{\pgfqpoint{1.120503in}{0.778769in}}%
\pgfpathlineto{\pgfqpoint{1.112050in}{0.779340in}}%
\pgfpathlineto{\pgfqpoint{1.109887in}{0.762605in}}%
\pgfpathlineto{\pgfqpoint{1.107720in}{0.745838in}}%
\pgfpathlineto{\pgfqpoint{1.114494in}{0.745375in}}%
\pgfpathlineto{\pgfqpoint{1.121164in}{0.744723in}}%
\pgfpathclose%
\pgfusepath{fill}%
\end{pgfscope}%
\begin{pgfscope}%
\pgfpathrectangle{\pgfqpoint{0.407592in}{0.222904in}}{\pgfqpoint{1.411163in}{1.411163in}}%
\pgfusepath{clip}%
\pgfsetbuttcap%
\pgfsetroundjoin%
\definecolor{currentfill}{rgb}{0.749020,0.772549,0.792157}%
\pgfsetfillcolor{currentfill}%
\pgfsetfillopacity{0.160000}%
\pgfsetlinewidth{0.000000pt}%
\definecolor{currentstroke}{rgb}{0.000000,0.000000,0.000000}%
\pgfsetstrokecolor{currentstroke}%
\pgfsetdash{}{0pt}%
\pgfpathmoveto{\pgfqpoint{1.080164in}{0.745760in}}%
\pgfpathlineto{\pgfqpoint{1.078916in}{0.762518in}}%
\pgfpathlineto{\pgfqpoint{1.077671in}{0.779245in}}%
\pgfpathlineto{\pgfqpoint{1.069113in}{0.778626in}}%
\pgfpathlineto{\pgfqpoint{1.060640in}{0.777773in}}%
\pgfpathlineto{\pgfqpoint{1.063575in}{0.761187in}}%
\pgfpathlineto{\pgfqpoint{1.066515in}{0.744570in}}%
\pgfpathlineto{\pgfqpoint{1.073304in}{0.745260in}}%
\pgfpathlineto{\pgfqpoint{1.080164in}{0.745760in}}%
\pgfpathclose%
\pgfusepath{fill}%
\end{pgfscope}%
\begin{pgfscope}%
\pgfpathrectangle{\pgfqpoint{0.407592in}{0.222904in}}{\pgfqpoint{1.411163in}{1.411163in}}%
\pgfusepath{clip}%
\pgfsetbuttcap%
\pgfsetroundjoin%
\definecolor{currentfill}{rgb}{0.749020,0.772549,0.792157}%
\pgfsetfillcolor{currentfill}%
\pgfsetfillopacity{0.160000}%
\pgfsetlinewidth{0.000000pt}%
\definecolor{currentstroke}{rgb}{0.000000,0.000000,0.000000}%
\pgfsetstrokecolor{currentstroke}%
\pgfsetdash{}{0pt}%
\pgfpathmoveto{\pgfqpoint{0.983157in}{0.830421in}}%
\pgfpathlineto{\pgfqpoint{0.975581in}{0.845980in}}%
\pgfpathlineto{\pgfqpoint{0.968014in}{0.861519in}}%
\pgfpathlineto{\pgfqpoint{0.956943in}{0.857878in}}%
\pgfpathlineto{\pgfqpoint{0.946416in}{0.853951in}}%
\pgfpathlineto{\pgfqpoint{0.955340in}{0.838848in}}%
\pgfpathlineto{\pgfqpoint{0.964271in}{0.823730in}}%
\pgfpathlineto{\pgfqpoint{0.973474in}{0.827201in}}%
\pgfpathlineto{\pgfqpoint{0.983157in}{0.830421in}}%
\pgfpathclose%
\pgfusepath{fill}%
\end{pgfscope}%
\begin{pgfscope}%
\pgfpathrectangle{\pgfqpoint{0.407592in}{0.222904in}}{\pgfqpoint{1.411163in}{1.411163in}}%
\pgfusepath{clip}%
\pgfsetbuttcap%
\pgfsetroundjoin%
\definecolor{currentfill}{rgb}{0.749020,0.772549,0.792157}%
\pgfsetfillcolor{currentfill}%
\pgfsetfillopacity{0.160000}%
\pgfsetlinewidth{0.000000pt}%
\definecolor{currentstroke}{rgb}{0.000000,0.000000,0.000000}%
\pgfsetstrokecolor{currentstroke}%
\pgfsetdash{}{0pt}%
\pgfpathmoveto{\pgfqpoint{1.462227in}{1.061705in}}%
\pgfpathlineto{\pgfqpoint{1.473790in}{1.075729in}}%
\pgfpathlineto{\pgfqpoint{1.485349in}{1.089748in}}%
\pgfpathlineto{\pgfqpoint{1.469863in}{1.099262in}}%
\pgfpathlineto{\pgfqpoint{1.452860in}{1.108313in}}%
\pgfpathlineto{\pgfqpoint{1.442313in}{1.093835in}}%
\pgfpathlineto{\pgfqpoint{1.431759in}{1.079348in}}%
\pgfpathlineto{\pgfqpoint{1.447711in}{1.070745in}}%
\pgfpathlineto{\pgfqpoint{1.462227in}{1.061705in}}%
\pgfpathclose%
\pgfusepath{fill}%
\end{pgfscope}%
\begin{pgfscope}%
\pgfpathrectangle{\pgfqpoint{0.407592in}{0.222904in}}{\pgfqpoint{1.411163in}{1.411163in}}%
\pgfusepath{clip}%
\pgfsetbuttcap%
\pgfsetroundjoin%
\definecolor{currentfill}{rgb}{0.749020,0.772549,0.792157}%
\pgfsetfillcolor{currentfill}%
\pgfsetfillopacity{0.160000}%
\pgfsetlinewidth{0.000000pt}%
\definecolor{currentstroke}{rgb}{0.000000,0.000000,0.000000}%
\pgfsetstrokecolor{currentstroke}%
\pgfsetdash{}{0pt}%
\pgfpathmoveto{\pgfqpoint{0.745749in}{1.106536in}}%
\pgfpathlineto{\pgfqpoint{0.735730in}{1.120961in}}%
\pgfpathlineto{\pgfqpoint{0.725719in}{1.135376in}}%
\pgfpathlineto{\pgfqpoint{0.705769in}{1.125787in}}%
\pgfpathlineto{\pgfqpoint{0.687320in}{1.115715in}}%
\pgfpathlineto{\pgfqpoint{0.698456in}{1.101753in}}%
\pgfpathlineto{\pgfqpoint{0.709595in}{1.087786in}}%
\pgfpathlineto{\pgfqpoint{0.726958in}{1.097389in}}%
\pgfpathlineto{\pgfqpoint{0.745749in}{1.106536in}}%
\pgfpathclose%
\pgfusepath{fill}%
\end{pgfscope}%
\begin{pgfscope}%
\pgfpathrectangle{\pgfqpoint{0.407592in}{0.222904in}}{\pgfqpoint{1.411163in}{1.411163in}}%
\pgfusepath{clip}%
\pgfsetbuttcap%
\pgfsetroundjoin%
\definecolor{currentfill}{rgb}{0.749020,0.772549,0.792157}%
\pgfsetfillcolor{currentfill}%
\pgfsetfillopacity{0.160000}%
\pgfsetlinewidth{0.000000pt}%
\definecolor{currentstroke}{rgb}{0.000000,0.000000,0.000000}%
\pgfsetstrokecolor{currentstroke}%
\pgfsetdash{}{0pt}%
\pgfpathmoveto{\pgfqpoint{1.134082in}{0.742860in}}%
\pgfpathlineto{\pgfqpoint{1.139521in}{0.759274in}}%
\pgfpathlineto{\pgfqpoint{1.144952in}{0.775660in}}%
\pgfpathlineto{\pgfqpoint{1.136988in}{0.776925in}}%
\pgfpathlineto{\pgfqpoint{1.128826in}{0.777963in}}%
\pgfpathlineto{\pgfqpoint{1.124999in}{0.761358in}}%
\pgfpathlineto{\pgfqpoint{1.121164in}{0.744723in}}%
\pgfpathlineto{\pgfqpoint{1.127702in}{0.743883in}}%
\pgfpathlineto{\pgfqpoint{1.134082in}{0.742860in}}%
\pgfpathclose%
\pgfusepath{fill}%
\end{pgfscope}%
\begin{pgfscope}%
\pgfpathrectangle{\pgfqpoint{0.407592in}{0.222904in}}{\pgfqpoint{1.411163in}{1.411163in}}%
\pgfusepath{clip}%
\pgfsetbuttcap%
\pgfsetroundjoin%
\definecolor{currentfill}{rgb}{0.749020,0.772549,0.792157}%
\pgfsetfillcolor{currentfill}%
\pgfsetfillopacity{0.160000}%
\pgfsetlinewidth{0.000000pt}%
\definecolor{currentstroke}{rgb}{0.000000,0.000000,0.000000}%
\pgfsetstrokecolor{currentstroke}%
\pgfsetdash{}{0pt}%
\pgfpathmoveto{\pgfqpoint{1.283252in}{0.875494in}}%
\pgfpathlineto{\pgfqpoint{1.293907in}{0.890119in}}%
\pgfpathlineto{\pgfqpoint{1.304554in}{0.904734in}}%
\pgfpathlineto{\pgfqpoint{1.293888in}{0.910132in}}%
\pgfpathlineto{\pgfqpoint{1.282383in}{0.915211in}}%
\pgfpathlineto{\pgfqpoint{1.272866in}{0.900121in}}%
\pgfpathlineto{\pgfqpoint{1.263340in}{0.885016in}}%
\pgfpathlineto{\pgfqpoint{1.273676in}{0.880399in}}%
\pgfpathlineto{\pgfqpoint{1.283252in}{0.875494in}}%
\pgfpathclose%
\pgfusepath{fill}%
\end{pgfscope}%
\begin{pgfscope}%
\pgfpathrectangle{\pgfqpoint{0.407592in}{0.222904in}}{\pgfqpoint{1.411163in}{1.411163in}}%
\pgfusepath{clip}%
\pgfsetbuttcap%
\pgfsetroundjoin%
\definecolor{currentfill}{rgb}{0.749020,0.772549,0.792157}%
\pgfsetfillcolor{currentfill}%
\pgfsetfillopacity{0.160000}%
\pgfsetlinewidth{0.000000pt}%
\definecolor{currentstroke}{rgb}{0.000000,0.000000,0.000000}%
\pgfsetstrokecolor{currentstroke}%
\pgfsetdash{}{0pt}%
\pgfpathmoveto{\pgfqpoint{1.066515in}{0.744570in}}%
\pgfpathlineto{\pgfqpoint{1.063575in}{0.761187in}}%
\pgfpathlineto{\pgfqpoint{1.060640in}{0.777773in}}%
\pgfpathlineto{\pgfqpoint{1.052289in}{0.776690in}}%
\pgfpathlineto{\pgfqpoint{1.044094in}{0.775379in}}%
\pgfpathlineto{\pgfqpoint{1.048673in}{0.759020in}}%
\pgfpathlineto{\pgfqpoint{1.053259in}{0.742632in}}%
\pgfpathlineto{\pgfqpoint{1.059824in}{0.743693in}}%
\pgfpathlineto{\pgfqpoint{1.066515in}{0.744570in}}%
\pgfpathclose%
\pgfusepath{fill}%
\end{pgfscope}%
\begin{pgfscope}%
\pgfpathrectangle{\pgfqpoint{0.407592in}{0.222904in}}{\pgfqpoint{1.411163in}{1.411163in}}%
\pgfusepath{clip}%
\pgfsetbuttcap%
\pgfsetroundjoin%
\definecolor{currentfill}{rgb}{0.749020,0.772549,0.792157}%
\pgfsetfillcolor{currentfill}%
\pgfsetfillopacity{0.160000}%
\pgfsetlinewidth{0.000000pt}%
\definecolor{currentstroke}{rgb}{0.000000,0.000000,0.000000}%
\pgfsetstrokecolor{currentstroke}%
\pgfsetdash{}{0pt}%
\pgfpathmoveto{\pgfqpoint{1.439090in}{1.033644in}}%
\pgfpathlineto{\pgfqpoint{1.450660in}{1.047676in}}%
\pgfpathlineto{\pgfqpoint{1.462227in}{1.061705in}}%
\pgfpathlineto{\pgfqpoint{1.447711in}{1.070745in}}%
\pgfpathlineto{\pgfqpoint{1.431759in}{1.079348in}}%
\pgfpathlineto{\pgfqpoint{1.421198in}{1.064851in}}%
\pgfpathlineto{\pgfqpoint{1.410630in}{1.050344in}}%
\pgfpathlineto{\pgfqpoint{1.425538in}{1.042199in}}%
\pgfpathlineto{\pgfqpoint{1.439090in}{1.033644in}}%
\pgfpathclose%
\pgfusepath{fill}%
\end{pgfscope}%
\begin{pgfscope}%
\pgfpathrectangle{\pgfqpoint{0.407592in}{0.222904in}}{\pgfqpoint{1.411163in}{1.411163in}}%
\pgfusepath{clip}%
\pgfsetbuttcap%
\pgfsetroundjoin%
\definecolor{currentfill}{rgb}{0.749020,0.772549,0.792157}%
\pgfsetfillcolor{currentfill}%
\pgfsetfillopacity{0.160000}%
\pgfsetlinewidth{0.000000pt}%
\definecolor{currentstroke}{rgb}{0.000000,0.000000,0.000000}%
\pgfsetstrokecolor{currentstroke}%
\pgfsetdash{}{0pt}%
\pgfpathmoveto{\pgfqpoint{1.174180in}{0.768437in}}%
\pgfpathlineto{\pgfqpoint{1.182516in}{0.784117in}}%
\pgfpathlineto{\pgfqpoint{1.190842in}{0.799776in}}%
\pgfpathlineto{\pgfqpoint{1.182621in}{0.802289in}}%
\pgfpathlineto{\pgfqpoint{1.174013in}{0.804561in}}%
\pgfpathlineto{\pgfqpoint{1.167090in}{0.788527in}}%
\pgfpathlineto{\pgfqpoint{1.160155in}{0.772468in}}%
\pgfpathlineto{\pgfqpoint{1.167330in}{0.770554in}}%
\pgfpathlineto{\pgfqpoint{1.174180in}{0.768437in}}%
\pgfpathclose%
\pgfusepath{fill}%
\end{pgfscope}%
\begin{pgfscope}%
\pgfpathrectangle{\pgfqpoint{0.407592in}{0.222904in}}{\pgfqpoint{1.411163in}{1.411163in}}%
\pgfusepath{clip}%
\pgfsetbuttcap%
\pgfsetroundjoin%
\definecolor{currentfill}{rgb}{0.749020,0.772549,0.792157}%
\pgfsetfillcolor{currentfill}%
\pgfsetfillopacity{0.160000}%
\pgfsetlinewidth{0.000000pt}%
\definecolor{currentstroke}{rgb}{0.000000,0.000000,0.000000}%
\pgfsetstrokecolor{currentstroke}%
\pgfsetdash{}{0pt}%
\pgfpathmoveto{\pgfqpoint{0.765804in}{1.077659in}}%
\pgfpathlineto{\pgfqpoint{0.755773in}{1.092102in}}%
\pgfpathlineto{\pgfqpoint{0.745749in}{1.106536in}}%
\pgfpathlineto{\pgfqpoint{0.726958in}{1.097389in}}%
\pgfpathlineto{\pgfqpoint{0.709595in}{1.087786in}}%
\pgfpathlineto{\pgfqpoint{0.720737in}{1.073816in}}%
\pgfpathlineto{\pgfqpoint{0.731883in}{1.059841in}}%
\pgfpathlineto{\pgfqpoint{0.748167in}{1.068965in}}%
\pgfpathlineto{\pgfqpoint{0.765804in}{1.077659in}}%
\pgfpathclose%
\pgfusepath{fill}%
\end{pgfscope}%
\begin{pgfscope}%
\pgfpathrectangle{\pgfqpoint{0.407592in}{0.222904in}}{\pgfqpoint{1.411163in}{1.411163in}}%
\pgfusepath{clip}%
\pgfsetbuttcap%
\pgfsetroundjoin%
\definecolor{currentfill}{rgb}{0.749020,0.772549,0.792157}%
\pgfsetfillcolor{currentfill}%
\pgfsetfillopacity{0.160000}%
\pgfsetlinewidth{0.000000pt}%
\definecolor{currentstroke}{rgb}{0.000000,0.000000,0.000000}%
\pgfsetstrokecolor{currentstroke}%
\pgfsetdash{}{0pt}%
\pgfpathmoveto{\pgfqpoint{0.928596in}{0.884114in}}%
\pgfpathlineto{\pgfqpoint{0.919699in}{0.899173in}}%
\pgfpathlineto{\pgfqpoint{0.910810in}{0.914218in}}%
\pgfpathlineto{\pgfqpoint{0.898355in}{0.909074in}}%
\pgfpathlineto{\pgfqpoint{0.886680in}{0.903614in}}%
\pgfpathlineto{\pgfqpoint{0.896795in}{0.889050in}}%
\pgfpathlineto{\pgfqpoint{0.906917in}{0.874476in}}%
\pgfpathlineto{\pgfqpoint{0.917403in}{0.879437in}}%
\pgfpathlineto{\pgfqpoint{0.928596in}{0.884114in}}%
\pgfpathclose%
\pgfusepath{fill}%
\end{pgfscope}%
\begin{pgfscope}%
\pgfpathrectangle{\pgfqpoint{0.407592in}{0.222904in}}{\pgfqpoint{1.411163in}{1.411163in}}%
\pgfusepath{clip}%
\pgfsetbuttcap%
\pgfsetroundjoin%
\definecolor{currentfill}{rgb}{0.749020,0.772549,0.792157}%
\pgfsetfillcolor{currentfill}%
\pgfsetfillopacity{0.160000}%
\pgfsetlinewidth{0.000000pt}%
\definecolor{currentstroke}{rgb}{0.000000,0.000000,0.000000}%
\pgfsetstrokecolor{currentstroke}%
\pgfsetdash{}{0pt}%
\pgfpathmoveto{\pgfqpoint{1.028311in}{0.772101in}}%
\pgfpathlineto{\pgfqpoint{1.022170in}{0.788125in}}%
\pgfpathlineto{\pgfqpoint{1.016039in}{0.804125in}}%
\pgfpathlineto{\pgfqpoint{1.007017in}{0.801804in}}%
\pgfpathlineto{\pgfqpoint{0.998338in}{0.799244in}}%
\pgfpathlineto{\pgfqpoint{1.005943in}{0.783626in}}%
\pgfpathlineto{\pgfqpoint{1.013558in}{0.767989in}}%
\pgfpathlineto{\pgfqpoint{1.020790in}{0.770145in}}%
\pgfpathlineto{\pgfqpoint{1.028311in}{0.772101in}}%
\pgfpathclose%
\pgfusepath{fill}%
\end{pgfscope}%
\begin{pgfscope}%
\pgfpathrectangle{\pgfqpoint{0.407592in}{0.222904in}}{\pgfqpoint{1.411163in}{1.411163in}}%
\pgfusepath{clip}%
\pgfsetbuttcap%
\pgfsetroundjoin%
\definecolor{currentfill}{rgb}{0.749020,0.772549,0.792157}%
\pgfsetfillcolor{currentfill}%
\pgfsetfillopacity{0.160000}%
\pgfsetlinewidth{0.000000pt}%
\definecolor{currentstroke}{rgb}{0.000000,0.000000,0.000000}%
\pgfsetstrokecolor{currentstroke}%
\pgfsetdash{}{0pt}%
\pgfpathmoveto{\pgfqpoint{1.205982in}{0.794070in}}%
\pgfpathlineto{\pgfqpoint{1.215566in}{0.809266in}}%
\pgfpathlineto{\pgfqpoint{1.225140in}{0.824446in}}%
\pgfpathlineto{\pgfqpoint{1.216564in}{0.827868in}}%
\pgfpathlineto{\pgfqpoint{1.207460in}{0.831035in}}%
\pgfpathlineto{\pgfqpoint{1.199156in}{0.815416in}}%
\pgfpathlineto{\pgfqpoint{1.190842in}{0.799776in}}%
\pgfpathlineto{\pgfqpoint{1.198640in}{0.797033in}}%
\pgfpathlineto{\pgfqpoint{1.205982in}{0.794070in}}%
\pgfpathclose%
\pgfusepath{fill}%
\end{pgfscope}%
\begin{pgfscope}%
\pgfpathrectangle{\pgfqpoint{0.407592in}{0.222904in}}{\pgfqpoint{1.411163in}{1.411163in}}%
\pgfusepath{clip}%
\pgfsetbuttcap%
\pgfsetroundjoin%
\definecolor{currentfill}{rgb}{0.749020,0.772549,0.792157}%
\pgfsetfillcolor{currentfill}%
\pgfsetfillopacity{0.160000}%
\pgfsetlinewidth{0.000000pt}%
\definecolor{currentstroke}{rgb}{0.000000,0.000000,0.000000}%
\pgfsetstrokecolor{currentstroke}%
\pgfsetdash{}{0pt}%
\pgfpathmoveto{\pgfqpoint{1.415938in}{1.005564in}}%
\pgfpathlineto{\pgfqpoint{1.427516in}{1.019606in}}%
\pgfpathlineto{\pgfqpoint{1.439090in}{1.033644in}}%
\pgfpathlineto{\pgfqpoint{1.425538in}{1.042199in}}%
\pgfpathlineto{\pgfqpoint{1.410630in}{1.050344in}}%
\pgfpathlineto{\pgfqpoint{1.400055in}{1.035827in}}%
\pgfpathlineto{\pgfqpoint{1.389472in}{1.021301in}}%
\pgfpathlineto{\pgfqpoint{1.403342in}{1.013625in}}%
\pgfpathlineto{\pgfqpoint{1.415938in}{1.005564in}}%
\pgfpathclose%
\pgfusepath{fill}%
\end{pgfscope}%
\begin{pgfscope}%
\pgfpathrectangle{\pgfqpoint{0.407592in}{0.222904in}}{\pgfqpoint{1.411163in}{1.411163in}}%
\pgfusepath{clip}%
\pgfsetbuttcap%
\pgfsetroundjoin%
\definecolor{currentfill}{rgb}{0.749020,0.772549,0.792157}%
\pgfsetfillcolor{currentfill}%
\pgfsetfillopacity{0.160000}%
\pgfsetlinewidth{0.000000pt}%
\definecolor{currentstroke}{rgb}{0.000000,0.000000,0.000000}%
\pgfsetstrokecolor{currentstroke}%
\pgfsetdash{}{0pt}%
\pgfpathmoveto{\pgfqpoint{1.146255in}{0.740277in}}%
\pgfpathlineto{\pgfqpoint{1.153210in}{0.756385in}}%
\pgfpathlineto{\pgfqpoint{1.160155in}{0.772468in}}%
\pgfpathlineto{\pgfqpoint{1.152685in}{0.774172in}}%
\pgfpathlineto{\pgfqpoint{1.144952in}{0.775660in}}%
\pgfpathlineto{\pgfqpoint{1.139521in}{0.759274in}}%
\pgfpathlineto{\pgfqpoint{1.134082in}{0.742860in}}%
\pgfpathlineto{\pgfqpoint{1.140275in}{0.741656in}}%
\pgfpathlineto{\pgfqpoint{1.146255in}{0.740277in}}%
\pgfpathclose%
\pgfusepath{fill}%
\end{pgfscope}%
\begin{pgfscope}%
\pgfpathrectangle{\pgfqpoint{0.407592in}{0.222904in}}{\pgfqpoint{1.411163in}{1.411163in}}%
\pgfusepath{clip}%
\pgfsetbuttcap%
\pgfsetroundjoin%
\definecolor{currentfill}{rgb}{0.749020,0.772549,0.792157}%
\pgfsetfillcolor{currentfill}%
\pgfsetfillopacity{0.160000}%
\pgfsetlinewidth{0.000000pt}%
\definecolor{currentstroke}{rgb}{0.000000,0.000000,0.000000}%
\pgfsetstrokecolor{currentstroke}%
\pgfsetdash{}{0pt}%
\pgfpathmoveto{\pgfqpoint{0.998338in}{0.799244in}}%
\pgfpathlineto{\pgfqpoint{0.990743in}{0.814842in}}%
\pgfpathlineto{\pgfqpoint{0.983157in}{0.830421in}}%
\pgfpathlineto{\pgfqpoint{0.973474in}{0.827201in}}%
\pgfpathlineto{\pgfqpoint{0.964271in}{0.823730in}}%
\pgfpathlineto{\pgfqpoint{0.973212in}{0.808597in}}%
\pgfpathlineto{\pgfqpoint{0.982161in}{0.793450in}}%
\pgfpathlineto{\pgfqpoint{0.990041in}{0.796456in}}%
\pgfpathlineto{\pgfqpoint{0.998338in}{0.799244in}}%
\pgfpathclose%
\pgfusepath{fill}%
\end{pgfscope}%
\begin{pgfscope}%
\pgfpathrectangle{\pgfqpoint{0.407592in}{0.222904in}}{\pgfqpoint{1.411163in}{1.411163in}}%
\pgfusepath{clip}%
\pgfsetbuttcap%
\pgfsetroundjoin%
\definecolor{currentfill}{rgb}{0.749020,0.772549,0.792157}%
\pgfsetfillcolor{currentfill}%
\pgfsetfillopacity{0.160000}%
\pgfsetlinewidth{0.000000pt}%
\definecolor{currentstroke}{rgb}{0.000000,0.000000,0.000000}%
\pgfsetstrokecolor{currentstroke}%
\pgfsetdash{}{0pt}%
\pgfpathmoveto{\pgfqpoint{0.785886in}{1.048744in}}%
\pgfpathlineto{\pgfqpoint{0.775842in}{1.063206in}}%
\pgfpathlineto{\pgfqpoint{0.765804in}{1.077659in}}%
\pgfpathlineto{\pgfqpoint{0.748167in}{1.068965in}}%
\pgfpathlineto{\pgfqpoint{0.731883in}{1.059841in}}%
\pgfpathlineto{\pgfqpoint{0.743031in}{1.045863in}}%
\pgfpathlineto{\pgfqpoint{0.754183in}{1.031880in}}%
\pgfpathlineto{\pgfqpoint{0.769395in}{1.040515in}}%
\pgfpathlineto{\pgfqpoint{0.785886in}{1.048744in}}%
\pgfpathclose%
\pgfusepath{fill}%
\end{pgfscope}%
\begin{pgfscope}%
\pgfpathrectangle{\pgfqpoint{0.407592in}{0.222904in}}{\pgfqpoint{1.411163in}{1.411163in}}%
\pgfusepath{clip}%
\pgfsetbuttcap%
\pgfsetroundjoin%
\definecolor{currentfill}{rgb}{0.749020,0.772549,0.792157}%
\pgfsetfillcolor{currentfill}%
\pgfsetfillopacity{0.160000}%
\pgfsetlinewidth{0.000000pt}%
\definecolor{currentstroke}{rgb}{0.000000,0.000000,0.000000}%
\pgfsetstrokecolor{currentstroke}%
\pgfsetdash{}{0pt}%
\pgfpathmoveto{\pgfqpoint{1.053259in}{0.742632in}}%
\pgfpathlineto{\pgfqpoint{1.048673in}{0.759020in}}%
\pgfpathlineto{\pgfqpoint{1.044094in}{0.775379in}}%
\pgfpathlineto{\pgfqpoint{1.036091in}{0.773848in}}%
\pgfpathlineto{\pgfqpoint{1.028311in}{0.772101in}}%
\pgfpathlineto{\pgfqpoint{1.034462in}{0.756052in}}%
\pgfpathlineto{\pgfqpoint{1.040621in}{0.739980in}}%
\pgfpathlineto{\pgfqpoint{1.046849in}{0.741393in}}%
\pgfpathlineto{\pgfqpoint{1.053259in}{0.742632in}}%
\pgfpathclose%
\pgfusepath{fill}%
\end{pgfscope}%
\begin{pgfscope}%
\pgfpathrectangle{\pgfqpoint{0.407592in}{0.222904in}}{\pgfqpoint{1.411163in}{1.411163in}}%
\pgfusepath{clip}%
\pgfsetbuttcap%
\pgfsetroundjoin%
\definecolor{currentfill}{rgb}{0.749020,0.772549,0.792157}%
\pgfsetfillcolor{currentfill}%
\pgfsetfillopacity{0.160000}%
\pgfsetlinewidth{0.000000pt}%
\definecolor{currentstroke}{rgb}{0.000000,0.000000,0.000000}%
\pgfsetstrokecolor{currentstroke}%
\pgfsetdash{}{0pt}%
\pgfpathmoveto{\pgfqpoint{1.261922in}{0.846213in}}%
\pgfpathlineto{\pgfqpoint{1.272591in}{0.860858in}}%
\pgfpathlineto{\pgfqpoint{1.283252in}{0.875494in}}%
\pgfpathlineto{\pgfqpoint{1.273676in}{0.880399in}}%
\pgfpathlineto{\pgfqpoint{1.263340in}{0.885016in}}%
\pgfpathlineto{\pgfqpoint{1.253804in}{0.869896in}}%
\pgfpathlineto{\pgfqpoint{1.244259in}{0.854761in}}%
\pgfpathlineto{\pgfqpoint{1.253431in}{0.850616in}}%
\pgfpathlineto{\pgfqpoint{1.261922in}{0.846213in}}%
\pgfpathclose%
\pgfusepath{fill}%
\end{pgfscope}%
\begin{pgfscope}%
\pgfpathrectangle{\pgfqpoint{0.407592in}{0.222904in}}{\pgfqpoint{1.411163in}{1.411163in}}%
\pgfusepath{clip}%
\pgfsetbuttcap%
\pgfsetroundjoin%
\definecolor{currentfill}{rgb}{0.749020,0.772549,0.792157}%
\pgfsetfillcolor{currentfill}%
\pgfsetfillopacity{0.160000}%
\pgfsetlinewidth{0.000000pt}%
\definecolor{currentstroke}{rgb}{0.000000,0.000000,0.000000}%
\pgfsetstrokecolor{currentstroke}%
\pgfsetdash{}{0pt}%
\pgfpathmoveto{\pgfqpoint{0.946416in}{0.853951in}}%
\pgfpathlineto{\pgfqpoint{0.937502in}{0.869040in}}%
\pgfpathlineto{\pgfqpoint{0.928596in}{0.884114in}}%
\pgfpathlineto{\pgfqpoint{0.917403in}{0.879437in}}%
\pgfpathlineto{\pgfqpoint{0.906917in}{0.874476in}}%
\pgfpathlineto{\pgfqpoint{0.917045in}{0.859893in}}%
\pgfpathlineto{\pgfqpoint{0.927180in}{0.845300in}}%
\pgfpathlineto{\pgfqpoint{0.936481in}{0.849753in}}%
\pgfpathlineto{\pgfqpoint{0.946416in}{0.853951in}}%
\pgfpathclose%
\pgfusepath{fill}%
\end{pgfscope}%
\begin{pgfscope}%
\pgfpathrectangle{\pgfqpoint{0.407592in}{0.222904in}}{\pgfqpoint{1.411163in}{1.411163in}}%
\pgfusepath{clip}%
\pgfsetbuttcap%
\pgfsetroundjoin%
\definecolor{currentfill}{rgb}{0.749020,0.772549,0.792157}%
\pgfsetfillcolor{currentfill}%
\pgfsetfillopacity{0.160000}%
\pgfsetlinewidth{0.000000pt}%
\definecolor{currentstroke}{rgb}{0.000000,0.000000,0.000000}%
\pgfsetstrokecolor{currentstroke}%
\pgfsetdash{}{0pt}%
\pgfpathmoveto{\pgfqpoint{1.392771in}{0.977467in}}%
\pgfpathlineto{\pgfqpoint{1.404357in}{0.991518in}}%
\pgfpathlineto{\pgfqpoint{1.415938in}{1.005564in}}%
\pgfpathlineto{\pgfqpoint{1.403342in}{1.013625in}}%
\pgfpathlineto{\pgfqpoint{1.389472in}{1.021301in}}%
\pgfpathlineto{\pgfqpoint{1.378882in}{1.006764in}}%
\pgfpathlineto{\pgfqpoint{1.368286in}{0.992218in}}%
\pgfpathlineto{\pgfqpoint{1.381123in}{0.985021in}}%
\pgfpathlineto{\pgfqpoint{1.392771in}{0.977467in}}%
\pgfpathclose%
\pgfusepath{fill}%
\end{pgfscope}%
\begin{pgfscope}%
\pgfpathrectangle{\pgfqpoint{0.407592in}{0.222904in}}{\pgfqpoint{1.411163in}{1.411163in}}%
\pgfusepath{clip}%
\pgfsetbuttcap%
\pgfsetroundjoin%
\definecolor{currentfill}{rgb}{0.749020,0.772549,0.792157}%
\pgfsetfillcolor{currentfill}%
\pgfsetfillopacity{0.160000}%
\pgfsetlinewidth{0.000000pt}%
\definecolor{currentstroke}{rgb}{0.000000,0.000000,0.000000}%
\pgfsetstrokecolor{currentstroke}%
\pgfsetdash{}{0pt}%
\pgfpathmoveto{\pgfqpoint{1.103372in}{0.712206in}}%
\pgfpathlineto{\pgfqpoint{1.105548in}{0.729038in}}%
\pgfpathlineto{\pgfqpoint{1.107720in}{0.745838in}}%
\pgfpathlineto{\pgfqpoint{1.100871in}{0.746108in}}%
\pgfpathlineto{\pgfqpoint{1.093976in}{0.746185in}}%
\pgfpathlineto{\pgfqpoint{1.093511in}{0.729344in}}%
\pgfpathlineto{\pgfqpoint{1.093045in}{0.712470in}}%
\pgfpathlineto{\pgfqpoint{1.098226in}{0.712411in}}%
\pgfpathlineto{\pgfqpoint{1.103372in}{0.712206in}}%
\pgfpathclose%
\pgfusepath{fill}%
\end{pgfscope}%
\begin{pgfscope}%
\pgfpathrectangle{\pgfqpoint{0.407592in}{0.222904in}}{\pgfqpoint{1.411163in}{1.411163in}}%
\pgfusepath{clip}%
\pgfsetbuttcap%
\pgfsetroundjoin%
\definecolor{currentfill}{rgb}{0.749020,0.772549,0.792157}%
\pgfsetfillcolor{currentfill}%
\pgfsetfillopacity{0.160000}%
\pgfsetlinewidth{0.000000pt}%
\definecolor{currentstroke}{rgb}{0.000000,0.000000,0.000000}%
\pgfsetstrokecolor{currentstroke}%
\pgfsetdash{}{0pt}%
\pgfpathmoveto{\pgfqpoint{1.093045in}{0.712470in}}%
\pgfpathlineto{\pgfqpoint{1.093511in}{0.729344in}}%
\pgfpathlineto{\pgfqpoint{1.093976in}{0.746185in}}%
\pgfpathlineto{\pgfqpoint{1.087064in}{0.746069in}}%
\pgfpathlineto{\pgfqpoint{1.080164in}{0.745760in}}%
\pgfpathlineto{\pgfqpoint{1.081414in}{0.728970in}}%
\pgfpathlineto{\pgfqpoint{1.082666in}{0.712148in}}%
\pgfpathlineto{\pgfqpoint{1.087851in}{0.712382in}}%
\pgfpathlineto{\pgfqpoint{1.093045in}{0.712470in}}%
\pgfpathclose%
\pgfusepath{fill}%
\end{pgfscope}%
\begin{pgfscope}%
\pgfpathrectangle{\pgfqpoint{0.407592in}{0.222904in}}{\pgfqpoint{1.411163in}{1.411163in}}%
\pgfusepath{clip}%
\pgfsetbuttcap%
\pgfsetroundjoin%
\definecolor{currentfill}{rgb}{0.749020,0.772549,0.792157}%
\pgfsetfillcolor{currentfill}%
\pgfsetfillopacity{0.160000}%
\pgfsetlinewidth{0.000000pt}%
\definecolor{currentstroke}{rgb}{0.000000,0.000000,0.000000}%
\pgfsetstrokecolor{currentstroke}%
\pgfsetdash{}{0pt}%
\pgfpathmoveto{\pgfqpoint{0.805993in}{1.019793in}}%
\pgfpathlineto{\pgfqpoint{0.795936in}{1.034273in}}%
\pgfpathlineto{\pgfqpoint{0.785886in}{1.048744in}}%
\pgfpathlineto{\pgfqpoint{0.769395in}{1.040515in}}%
\pgfpathlineto{\pgfqpoint{0.754183in}{1.031880in}}%
\pgfpathlineto{\pgfqpoint{0.765338in}{1.017894in}}%
\pgfpathlineto{\pgfqpoint{0.776496in}{1.003903in}}%
\pgfpathlineto{\pgfqpoint{0.790644in}{1.012037in}}%
\pgfpathlineto{\pgfqpoint{0.805993in}{1.019793in}}%
\pgfpathclose%
\pgfusepath{fill}%
\end{pgfscope}%
\begin{pgfscope}%
\pgfpathrectangle{\pgfqpoint{0.407592in}{0.222904in}}{\pgfqpoint{1.411163in}{1.411163in}}%
\pgfusepath{clip}%
\pgfsetbuttcap%
\pgfsetroundjoin%
\definecolor{currentfill}{rgb}{0.749020,0.772549,0.792157}%
\pgfsetfillcolor{currentfill}%
\pgfsetfillopacity{0.160000}%
\pgfsetlinewidth{0.000000pt}%
\definecolor{currentstroke}{rgb}{0.000000,0.000000,0.000000}%
\pgfsetstrokecolor{currentstroke}%
\pgfsetdash{}{0pt}%
\pgfpathmoveto{\pgfqpoint{1.113473in}{0.711361in}}%
\pgfpathlineto{\pgfqpoint{1.117322in}{0.728057in}}%
\pgfpathlineto{\pgfqpoint{1.121164in}{0.744723in}}%
\pgfpathlineto{\pgfqpoint{1.114494in}{0.745375in}}%
\pgfpathlineto{\pgfqpoint{1.107720in}{0.745838in}}%
\pgfpathlineto{\pgfqpoint{1.105548in}{0.729038in}}%
\pgfpathlineto{\pgfqpoint{1.103372in}{0.712206in}}%
\pgfpathlineto{\pgfqpoint{1.108462in}{0.711856in}}%
\pgfpathlineto{\pgfqpoint{1.113473in}{0.711361in}}%
\pgfpathclose%
\pgfusepath{fill}%
\end{pgfscope}%
\begin{pgfscope}%
\pgfpathrectangle{\pgfqpoint{0.407592in}{0.222904in}}{\pgfqpoint{1.411163in}{1.411163in}}%
\pgfusepath{clip}%
\pgfsetbuttcap%
\pgfsetroundjoin%
\definecolor{currentfill}{rgb}{0.749020,0.772549,0.792157}%
\pgfsetfillcolor{currentfill}%
\pgfsetfillopacity{0.160000}%
\pgfsetlinewidth{0.000000pt}%
\definecolor{currentstroke}{rgb}{0.000000,0.000000,0.000000}%
\pgfsetstrokecolor{currentstroke}%
\pgfsetdash{}{0pt}%
\pgfpathmoveto{\pgfqpoint{1.082666in}{0.712148in}}%
\pgfpathlineto{\pgfqpoint{1.081414in}{0.728970in}}%
\pgfpathlineto{\pgfqpoint{1.080164in}{0.745760in}}%
\pgfpathlineto{\pgfqpoint{1.073304in}{0.745260in}}%
\pgfpathlineto{\pgfqpoint{1.066515in}{0.744570in}}%
\pgfpathlineto{\pgfqpoint{1.069460in}{0.727922in}}%
\pgfpathlineto{\pgfqpoint{1.072411in}{0.711244in}}%
\pgfpathlineto{\pgfqpoint{1.077512in}{0.711768in}}%
\pgfpathlineto{\pgfqpoint{1.082666in}{0.712148in}}%
\pgfpathclose%
\pgfusepath{fill}%
\end{pgfscope}%
\begin{pgfscope}%
\pgfpathrectangle{\pgfqpoint{0.407592in}{0.222904in}}{\pgfqpoint{1.411163in}{1.411163in}}%
\pgfusepath{clip}%
\pgfsetbuttcap%
\pgfsetroundjoin%
\definecolor{currentfill}{rgb}{0.749020,0.772549,0.792157}%
\pgfsetfillcolor{currentfill}%
\pgfsetfillopacity{0.160000}%
\pgfsetlinewidth{0.000000pt}%
\definecolor{currentstroke}{rgb}{0.000000,0.000000,0.000000}%
\pgfsetstrokecolor{currentstroke}%
\pgfsetdash{}{0pt}%
\pgfpathmoveto{\pgfqpoint{1.240562in}{0.816893in}}%
\pgfpathlineto{\pgfqpoint{1.251246in}{0.831558in}}%
\pgfpathlineto{\pgfqpoint{1.261922in}{0.846213in}}%
\pgfpathlineto{\pgfqpoint{1.253431in}{0.850616in}}%
\pgfpathlineto{\pgfqpoint{1.244259in}{0.854761in}}%
\pgfpathlineto{\pgfqpoint{1.234704in}{0.839611in}}%
\pgfpathlineto{\pgfqpoint{1.225140in}{0.824446in}}%
\pgfpathlineto{\pgfqpoint{1.233151in}{0.820783in}}%
\pgfpathlineto{\pgfqpoint{1.240562in}{0.816893in}}%
\pgfpathclose%
\pgfusepath{fill}%
\end{pgfscope}%
\begin{pgfscope}%
\pgfpathrectangle{\pgfqpoint{0.407592in}{0.222904in}}{\pgfqpoint{1.411163in}{1.411163in}}%
\pgfusepath{clip}%
\pgfsetbuttcap%
\pgfsetroundjoin%
\definecolor{currentfill}{rgb}{0.749020,0.772549,0.792157}%
\pgfsetfillcolor{currentfill}%
\pgfsetfillopacity{0.160000}%
\pgfsetlinewidth{0.000000pt}%
\definecolor{currentstroke}{rgb}{0.000000,0.000000,0.000000}%
\pgfsetstrokecolor{currentstroke}%
\pgfsetdash{}{0pt}%
\pgfpathmoveto{\pgfqpoint{1.157476in}{0.737018in}}%
\pgfpathlineto{\pgfqpoint{1.165833in}{0.752738in}}%
\pgfpathlineto{\pgfqpoint{1.174180in}{0.768437in}}%
\pgfpathlineto{\pgfqpoint{1.167330in}{0.770554in}}%
\pgfpathlineto{\pgfqpoint{1.160155in}{0.772468in}}%
\pgfpathlineto{\pgfqpoint{1.153210in}{0.756385in}}%
\pgfpathlineto{\pgfqpoint{1.146255in}{0.740277in}}%
\pgfpathlineto{\pgfqpoint{1.151997in}{0.738729in}}%
\pgfpathlineto{\pgfqpoint{1.157476in}{0.737018in}}%
\pgfpathclose%
\pgfusepath{fill}%
\end{pgfscope}%
\begin{pgfscope}%
\pgfpathrectangle{\pgfqpoint{0.407592in}{0.222904in}}{\pgfqpoint{1.411163in}{1.411163in}}%
\pgfusepath{clip}%
\pgfsetbuttcap%
\pgfsetroundjoin%
\definecolor{currentfill}{rgb}{0.749020,0.772549,0.792157}%
\pgfsetfillcolor{currentfill}%
\pgfsetfillopacity{0.160000}%
\pgfsetlinewidth{0.000000pt}%
\definecolor{currentstroke}{rgb}{0.000000,0.000000,0.000000}%
\pgfsetstrokecolor{currentstroke}%
\pgfsetdash{}{0pt}%
\pgfpathmoveto{\pgfqpoint{0.826126in}{0.990804in}}%
\pgfpathlineto{\pgfqpoint{0.816056in}{1.005303in}}%
\pgfpathlineto{\pgfqpoint{0.805993in}{1.019793in}}%
\pgfpathlineto{\pgfqpoint{0.790644in}{1.012037in}}%
\pgfpathlineto{\pgfqpoint{0.776496in}{1.003903in}}%
\pgfpathlineto{\pgfqpoint{0.787657in}{0.989909in}}%
\pgfpathlineto{\pgfqpoint{0.798822in}{0.975910in}}%
\pgfpathlineto{\pgfqpoint{0.811912in}{0.983533in}}%
\pgfpathlineto{\pgfqpoint{0.826126in}{0.990804in}}%
\pgfpathclose%
\pgfusepath{fill}%
\end{pgfscope}%
\begin{pgfscope}%
\pgfpathrectangle{\pgfqpoint{0.407592in}{0.222904in}}{\pgfqpoint{1.411163in}{1.411163in}}%
\pgfusepath{clip}%
\pgfsetbuttcap%
\pgfsetroundjoin%
\definecolor{currentfill}{rgb}{0.749020,0.772549,0.792157}%
\pgfsetfillcolor{currentfill}%
\pgfsetfillopacity{0.160000}%
\pgfsetlinewidth{0.000000pt}%
\definecolor{currentstroke}{rgb}{0.000000,0.000000,0.000000}%
\pgfsetstrokecolor{currentstroke}%
\pgfsetdash{}{0pt}%
\pgfpathmoveto{\pgfqpoint{1.186787in}{0.763633in}}%
\pgfpathlineto{\pgfqpoint{1.196389in}{0.778859in}}%
\pgfpathlineto{\pgfqpoint{1.205982in}{0.794070in}}%
\pgfpathlineto{\pgfqpoint{1.198640in}{0.797033in}}%
\pgfpathlineto{\pgfqpoint{1.190842in}{0.799776in}}%
\pgfpathlineto{\pgfqpoint{1.182516in}{0.784117in}}%
\pgfpathlineto{\pgfqpoint{1.174180in}{0.768437in}}%
\pgfpathlineto{\pgfqpoint{1.180675in}{0.766128in}}%
\pgfpathlineto{\pgfqpoint{1.186787in}{0.763633in}}%
\pgfpathclose%
\pgfusepath{fill}%
\end{pgfscope}%
\begin{pgfscope}%
\pgfpathrectangle{\pgfqpoint{0.407592in}{0.222904in}}{\pgfqpoint{1.411163in}{1.411163in}}%
\pgfusepath{clip}%
\pgfsetbuttcap%
\pgfsetroundjoin%
\definecolor{currentfill}{rgb}{0.749020,0.772549,0.792157}%
\pgfsetfillcolor{currentfill}%
\pgfsetfillopacity{0.160000}%
\pgfsetlinewidth{0.000000pt}%
\definecolor{currentstroke}{rgb}{0.000000,0.000000,0.000000}%
\pgfsetstrokecolor{currentstroke}%
\pgfsetdash{}{0pt}%
\pgfpathmoveto{\pgfqpoint{1.369589in}{0.949351in}}%
\pgfpathlineto{\pgfqpoint{1.381182in}{0.963411in}}%
\pgfpathlineto{\pgfqpoint{1.392771in}{0.977467in}}%
\pgfpathlineto{\pgfqpoint{1.381123in}{0.985021in}}%
\pgfpathlineto{\pgfqpoint{1.368286in}{0.992218in}}%
\pgfpathlineto{\pgfqpoint{1.357682in}{0.977662in}}%
\pgfpathlineto{\pgfqpoint{1.347070in}{0.963096in}}%
\pgfpathlineto{\pgfqpoint{1.358882in}{0.956389in}}%
\pgfpathlineto{\pgfqpoint{1.369589in}{0.949351in}}%
\pgfpathclose%
\pgfusepath{fill}%
\end{pgfscope}%
\begin{pgfscope}%
\pgfpathrectangle{\pgfqpoint{0.407592in}{0.222904in}}{\pgfqpoint{1.411163in}{1.411163in}}%
\pgfusepath{clip}%
\pgfsetbuttcap%
\pgfsetroundjoin%
\definecolor{currentfill}{rgb}{0.749020,0.772549,0.792157}%
\pgfsetfillcolor{currentfill}%
\pgfsetfillopacity{0.160000}%
\pgfsetlinewidth{0.000000pt}%
\definecolor{currentstroke}{rgb}{0.000000,0.000000,0.000000}%
\pgfsetstrokecolor{currentstroke}%
\pgfsetdash{}{0pt}%
\pgfpathmoveto{\pgfqpoint{1.123175in}{0.709947in}}%
\pgfpathlineto{\pgfqpoint{1.128633in}{0.726417in}}%
\pgfpathlineto{\pgfqpoint{1.134082in}{0.742860in}}%
\pgfpathlineto{\pgfqpoint{1.127702in}{0.743883in}}%
\pgfpathlineto{\pgfqpoint{1.121164in}{0.744723in}}%
\pgfpathlineto{\pgfqpoint{1.117322in}{0.728057in}}%
\pgfpathlineto{\pgfqpoint{1.113473in}{0.711361in}}%
\pgfpathlineto{\pgfqpoint{1.118384in}{0.710724in}}%
\pgfpathlineto{\pgfqpoint{1.123175in}{0.709947in}}%
\pgfpathclose%
\pgfusepath{fill}%
\end{pgfscope}%
\begin{pgfscope}%
\pgfpathrectangle{\pgfqpoint{0.407592in}{0.222904in}}{\pgfqpoint{1.411163in}{1.411163in}}%
\pgfusepath{clip}%
\pgfsetbuttcap%
\pgfsetroundjoin%
\definecolor{currentfill}{rgb}{0.749020,0.772549,0.792157}%
\pgfsetfillcolor{currentfill}%
\pgfsetfillopacity{0.160000}%
\pgfsetlinewidth{0.000000pt}%
\definecolor{currentstroke}{rgb}{0.000000,0.000000,0.000000}%
\pgfsetstrokecolor{currentstroke}%
\pgfsetdash{}{0pt}%
\pgfpathmoveto{\pgfqpoint{1.040621in}{0.739980in}}%
\pgfpathlineto{\pgfqpoint{1.034462in}{0.756052in}}%
\pgfpathlineto{\pgfqpoint{1.028311in}{0.772101in}}%
\pgfpathlineto{\pgfqpoint{1.020790in}{0.770145in}}%
\pgfpathlineto{\pgfqpoint{1.013558in}{0.767989in}}%
\pgfpathlineto{\pgfqpoint{1.021182in}{0.752333in}}%
\pgfpathlineto{\pgfqpoint{1.028815in}{0.736656in}}%
\pgfpathlineto{\pgfqpoint{1.034601in}{0.738399in}}%
\pgfpathlineto{\pgfqpoint{1.040621in}{0.739980in}}%
\pgfpathclose%
\pgfusepath{fill}%
\end{pgfscope}%
\begin{pgfscope}%
\pgfpathrectangle{\pgfqpoint{0.407592in}{0.222904in}}{\pgfqpoint{1.411163in}{1.411163in}}%
\pgfusepath{clip}%
\pgfsetbuttcap%
\pgfsetroundjoin%
\definecolor{currentfill}{rgb}{0.749020,0.772549,0.792157}%
\pgfsetfillcolor{currentfill}%
\pgfsetfillopacity{0.160000}%
\pgfsetlinewidth{0.000000pt}%
\definecolor{currentstroke}{rgb}{0.000000,0.000000,0.000000}%
\pgfsetstrokecolor{currentstroke}%
\pgfsetdash{}{0pt}%
\pgfpathmoveto{\pgfqpoint{1.072411in}{0.711244in}}%
\pgfpathlineto{\pgfqpoint{1.069460in}{0.727922in}}%
\pgfpathlineto{\pgfqpoint{1.066515in}{0.744570in}}%
\pgfpathlineto{\pgfqpoint{1.059824in}{0.743693in}}%
\pgfpathlineto{\pgfqpoint{1.053259in}{0.742632in}}%
\pgfpathlineto{\pgfqpoint{1.057853in}{0.726217in}}%
\pgfpathlineto{\pgfqpoint{1.062455in}{0.709775in}}%
\pgfpathlineto{\pgfqpoint{1.067385in}{0.710579in}}%
\pgfpathlineto{\pgfqpoint{1.072411in}{0.711244in}}%
\pgfpathclose%
\pgfusepath{fill}%
\end{pgfscope}%
\begin{pgfscope}%
\pgfpathrectangle{\pgfqpoint{0.407592in}{0.222904in}}{\pgfqpoint{1.411163in}{1.411163in}}%
\pgfusepath{clip}%
\pgfsetbuttcap%
\pgfsetroundjoin%
\definecolor{currentfill}{rgb}{0.749020,0.772549,0.792157}%
\pgfsetfillcolor{currentfill}%
\pgfsetfillopacity{0.160000}%
\pgfsetlinewidth{0.000000pt}%
\definecolor{currentstroke}{rgb}{0.000000,0.000000,0.000000}%
\pgfsetstrokecolor{currentstroke}%
\pgfsetdash{}{0pt}%
\pgfpathmoveto{\pgfqpoint{1.013558in}{0.767989in}}%
\pgfpathlineto{\pgfqpoint{1.005943in}{0.783626in}}%
\pgfpathlineto{\pgfqpoint{0.998338in}{0.799244in}}%
\pgfpathlineto{\pgfqpoint{0.990041in}{0.796456in}}%
\pgfpathlineto{\pgfqpoint{0.982161in}{0.793450in}}%
\pgfpathlineto{\pgfqpoint{0.991119in}{0.778288in}}%
\pgfpathlineto{\pgfqpoint{1.000085in}{0.763112in}}%
\pgfpathlineto{\pgfqpoint{1.006646in}{0.765642in}}%
\pgfpathlineto{\pgfqpoint{1.013558in}{0.767989in}}%
\pgfpathclose%
\pgfusepath{fill}%
\end{pgfscope}%
\begin{pgfscope}%
\pgfpathrectangle{\pgfqpoint{0.407592in}{0.222904in}}{\pgfqpoint{1.411163in}{1.411163in}}%
\pgfusepath{clip}%
\pgfsetbuttcap%
\pgfsetroundjoin%
\definecolor{currentfill}{rgb}{0.749020,0.772549,0.792157}%
\pgfsetfillcolor{currentfill}%
\pgfsetfillopacity{0.160000}%
\pgfsetlinewidth{0.000000pt}%
\definecolor{currentstroke}{rgb}{0.000000,0.000000,0.000000}%
\pgfsetstrokecolor{currentstroke}%
\pgfsetdash{}{0pt}%
\pgfpathmoveto{\pgfqpoint{0.964271in}{0.823730in}}%
\pgfpathlineto{\pgfqpoint{0.955340in}{0.838848in}}%
\pgfpathlineto{\pgfqpoint{0.946416in}{0.853951in}}%
\pgfpathlineto{\pgfqpoint{0.936481in}{0.849753in}}%
\pgfpathlineto{\pgfqpoint{0.927180in}{0.845300in}}%
\pgfpathlineto{\pgfqpoint{0.937321in}{0.830698in}}%
\pgfpathlineto{\pgfqpoint{0.947469in}{0.816087in}}%
\pgfpathlineto{\pgfqpoint{0.955591in}{0.820020in}}%
\pgfpathlineto{\pgfqpoint{0.964271in}{0.823730in}}%
\pgfpathclose%
\pgfusepath{fill}%
\end{pgfscope}%
\begin{pgfscope}%
\pgfpathrectangle{\pgfqpoint{0.407592in}{0.222904in}}{\pgfqpoint{1.411163in}{1.411163in}}%
\pgfusepath{clip}%
\pgfsetbuttcap%
\pgfsetroundjoin%
\definecolor{currentfill}{rgb}{0.749020,0.772549,0.792157}%
\pgfsetfillcolor{currentfill}%
\pgfsetfillopacity{0.160000}%
\pgfsetlinewidth{0.000000pt}%
\definecolor{currentstroke}{rgb}{0.000000,0.000000,0.000000}%
\pgfsetstrokecolor{currentstroke}%
\pgfsetdash{}{0pt}%
\pgfpathmoveto{\pgfqpoint{0.846285in}{0.961778in}}%
\pgfpathlineto{\pgfqpoint{0.836202in}{0.976296in}}%
\pgfpathlineto{\pgfqpoint{0.826126in}{0.990804in}}%
\pgfpathlineto{\pgfqpoint{0.811912in}{0.983533in}}%
\pgfpathlineto{\pgfqpoint{0.798822in}{0.975910in}}%
\pgfpathlineto{\pgfqpoint{0.809990in}{0.961908in}}%
\pgfpathlineto{\pgfqpoint{0.821161in}{0.947901in}}%
\pgfpathlineto{\pgfqpoint{0.833201in}{0.955003in}}%
\pgfpathlineto{\pgfqpoint{0.846285in}{0.961778in}}%
\pgfpathclose%
\pgfusepath{fill}%
\end{pgfscope}%
\begin{pgfscope}%
\pgfpathrectangle{\pgfqpoint{0.407592in}{0.222904in}}{\pgfqpoint{1.411163in}{1.411163in}}%
\pgfusepath{clip}%
\pgfsetbuttcap%
\pgfsetroundjoin%
\definecolor{currentfill}{rgb}{0.749020,0.772549,0.792157}%
\pgfsetfillcolor{currentfill}%
\pgfsetfillopacity{0.160000}%
\pgfsetlinewidth{0.000000pt}%
\definecolor{currentstroke}{rgb}{0.000000,0.000000,0.000000}%
\pgfsetstrokecolor{currentstroke}%
\pgfsetdash{}{0pt}%
\pgfpathmoveto{\pgfqpoint{1.346392in}{0.921217in}}%
\pgfpathlineto{\pgfqpoint{1.357993in}{0.935286in}}%
\pgfpathlineto{\pgfqpoint{1.369589in}{0.949351in}}%
\pgfpathlineto{\pgfqpoint{1.358882in}{0.956389in}}%
\pgfpathlineto{\pgfqpoint{1.347070in}{0.963096in}}%
\pgfpathlineto{\pgfqpoint{1.336452in}{0.948521in}}%
\pgfpathlineto{\pgfqpoint{1.325827in}{0.933935in}}%
\pgfpathlineto{\pgfqpoint{1.336619in}{0.927728in}}%
\pgfpathlineto{\pgfqpoint{1.346392in}{0.921217in}}%
\pgfpathclose%
\pgfusepath{fill}%
\end{pgfscope}%
\begin{pgfscope}%
\pgfpathrectangle{\pgfqpoint{0.407592in}{0.222904in}}{\pgfqpoint{1.411163in}{1.411163in}}%
\pgfusepath{clip}%
\pgfsetbuttcap%
\pgfsetroundjoin%
\definecolor{currentfill}{rgb}{0.749020,0.772549,0.792157}%
\pgfsetfillcolor{currentfill}%
\pgfsetfillopacity{0.160000}%
\pgfsetlinewidth{0.000000pt}%
\definecolor{currentstroke}{rgb}{0.000000,0.000000,0.000000}%
\pgfsetstrokecolor{currentstroke}%
\pgfsetdash{}{0pt}%
\pgfpathmoveto{\pgfqpoint{1.132312in}{0.707989in}}%
\pgfpathlineto{\pgfqpoint{1.139289in}{0.724145in}}%
\pgfpathlineto{\pgfqpoint{1.146255in}{0.740277in}}%
\pgfpathlineto{\pgfqpoint{1.140275in}{0.741656in}}%
\pgfpathlineto{\pgfqpoint{1.134082in}{0.742860in}}%
\pgfpathlineto{\pgfqpoint{1.128633in}{0.726417in}}%
\pgfpathlineto{\pgfqpoint{1.123175in}{0.709947in}}%
\pgfpathlineto{\pgfqpoint{1.127824in}{0.709034in}}%
\pgfpathlineto{\pgfqpoint{1.132312in}{0.707989in}}%
\pgfpathclose%
\pgfusepath{fill}%
\end{pgfscope}%
\begin{pgfscope}%
\pgfpathrectangle{\pgfqpoint{0.407592in}{0.222904in}}{\pgfqpoint{1.411163in}{1.411163in}}%
\pgfusepath{clip}%
\pgfsetbuttcap%
\pgfsetroundjoin%
\definecolor{currentfill}{rgb}{0.749020,0.772549,0.792157}%
\pgfsetfillcolor{currentfill}%
\pgfsetfillopacity{0.160000}%
\pgfsetlinewidth{0.000000pt}%
\definecolor{currentstroke}{rgb}{0.000000,0.000000,0.000000}%
\pgfsetstrokecolor{currentstroke}%
\pgfsetdash{}{0pt}%
\pgfpathmoveto{\pgfqpoint{1.062455in}{0.709775in}}%
\pgfpathlineto{\pgfqpoint{1.057853in}{0.726217in}}%
\pgfpathlineto{\pgfqpoint{1.053259in}{0.742632in}}%
\pgfpathlineto{\pgfqpoint{1.046849in}{0.741393in}}%
\pgfpathlineto{\pgfqpoint{1.040621in}{0.739980in}}%
\pgfpathlineto{\pgfqpoint{1.046790in}{0.723884in}}%
\pgfpathlineto{\pgfqpoint{1.052967in}{0.707764in}}%
\pgfpathlineto{\pgfqpoint{1.057642in}{0.708835in}}%
\pgfpathlineto{\pgfqpoint{1.062455in}{0.709775in}}%
\pgfpathclose%
\pgfusepath{fill}%
\end{pgfscope}%
\begin{pgfscope}%
\pgfpathrectangle{\pgfqpoint{0.407592in}{0.222904in}}{\pgfqpoint{1.411163in}{1.411163in}}%
\pgfusepath{clip}%
\pgfsetbuttcap%
\pgfsetroundjoin%
\definecolor{currentfill}{rgb}{0.749020,0.772549,0.792157}%
\pgfsetfillcolor{currentfill}%
\pgfsetfillopacity{0.160000}%
\pgfsetlinewidth{0.000000pt}%
\definecolor{currentstroke}{rgb}{0.000000,0.000000,0.000000}%
\pgfsetstrokecolor{currentstroke}%
\pgfsetdash{}{0pt}%
\pgfpathmoveto{\pgfqpoint{1.219174in}{0.787533in}}%
\pgfpathlineto{\pgfqpoint{1.229872in}{0.802218in}}%
\pgfpathlineto{\pgfqpoint{1.240562in}{0.816893in}}%
\pgfpathlineto{\pgfqpoint{1.233151in}{0.820783in}}%
\pgfpathlineto{\pgfqpoint{1.225140in}{0.824446in}}%
\pgfpathlineto{\pgfqpoint{1.215566in}{0.809266in}}%
\pgfpathlineto{\pgfqpoint{1.205982in}{0.794070in}}%
\pgfpathlineto{\pgfqpoint{1.212837in}{0.790899in}}%
\pgfpathlineto{\pgfqpoint{1.219174in}{0.787533in}}%
\pgfpathclose%
\pgfusepath{fill}%
\end{pgfscope}%
\begin{pgfscope}%
\pgfpathrectangle{\pgfqpoint{0.407592in}{0.222904in}}{\pgfqpoint{1.411163in}{1.411163in}}%
\pgfusepath{clip}%
\pgfsetbuttcap%
\pgfsetroundjoin%
\definecolor{currentfill}{rgb}{0.749020,0.772549,0.792157}%
\pgfsetfillcolor{currentfill}%
\pgfsetfillopacity{0.160000}%
\pgfsetlinewidth{0.000000pt}%
\definecolor{currentstroke}{rgb}{0.000000,0.000000,0.000000}%
\pgfsetstrokecolor{currentstroke}%
\pgfsetdash{}{0pt}%
\pgfpathmoveto{\pgfqpoint{0.866469in}{0.932715in}}%
\pgfpathlineto{\pgfqpoint{0.856374in}{0.947251in}}%
\pgfpathlineto{\pgfqpoint{0.846285in}{0.961778in}}%
\pgfpathlineto{\pgfqpoint{0.833201in}{0.955003in}}%
\pgfpathlineto{\pgfqpoint{0.821161in}{0.947901in}}%
\pgfpathlineto{\pgfqpoint{0.832335in}{0.933891in}}%
\pgfpathlineto{\pgfqpoint{0.843512in}{0.919876in}}%
\pgfpathlineto{\pgfqpoint{0.854509in}{0.926445in}}%
\pgfpathlineto{\pgfqpoint{0.866469in}{0.932715in}}%
\pgfpathclose%
\pgfusepath{fill}%
\end{pgfscope}%
\begin{pgfscope}%
\pgfpathrectangle{\pgfqpoint{0.407592in}{0.222904in}}{\pgfqpoint{1.411163in}{1.411163in}}%
\pgfusepath{clip}%
\pgfsetbuttcap%
\pgfsetroundjoin%
\definecolor{currentfill}{rgb}{0.749020,0.772549,0.792157}%
\pgfsetfillcolor{currentfill}%
\pgfsetfillopacity{0.160000}%
\pgfsetlinewidth{0.000000pt}%
\definecolor{currentstroke}{rgb}{0.000000,0.000000,0.000000}%
\pgfsetstrokecolor{currentstroke}%
\pgfsetdash{}{0pt}%
\pgfpathmoveto{\pgfqpoint{1.323181in}{0.893064in}}%
\pgfpathlineto{\pgfqpoint{1.334788in}{0.907143in}}%
\pgfpathlineto{\pgfqpoint{1.346392in}{0.921217in}}%
\pgfpathlineto{\pgfqpoint{1.336619in}{0.927728in}}%
\pgfpathlineto{\pgfqpoint{1.325827in}{0.933935in}}%
\pgfpathlineto{\pgfqpoint{1.315194in}{0.919339in}}%
\pgfpathlineto{\pgfqpoint{1.304554in}{0.904734in}}%
\pgfpathlineto{\pgfqpoint{1.314333in}{0.899038in}}%
\pgfpathlineto{\pgfqpoint{1.323181in}{0.893064in}}%
\pgfpathclose%
\pgfusepath{fill}%
\end{pgfscope}%
\begin{pgfscope}%
\pgfpathrectangle{\pgfqpoint{0.407592in}{0.222904in}}{\pgfqpoint{1.411163in}{1.411163in}}%
\pgfusepath{clip}%
\pgfsetbuttcap%
\pgfsetroundjoin%
\definecolor{currentfill}{rgb}{0.749020,0.772549,0.792157}%
\pgfsetfillcolor{currentfill}%
\pgfsetfillopacity{0.160000}%
\pgfsetlinewidth{0.000000pt}%
\definecolor{currentstroke}{rgb}{0.000000,0.000000,0.000000}%
\pgfsetstrokecolor{currentstroke}%
\pgfsetdash{}{0pt}%
\pgfpathmoveto{\pgfqpoint{0.982161in}{0.793450in}}%
\pgfpathlineto{\pgfqpoint{0.973212in}{0.808597in}}%
\pgfpathlineto{\pgfqpoint{0.964271in}{0.823730in}}%
\pgfpathlineto{\pgfqpoint{0.955591in}{0.820020in}}%
\pgfpathlineto{\pgfqpoint{0.947469in}{0.816087in}}%
\pgfpathlineto{\pgfqpoint{0.957624in}{0.801466in}}%
\pgfpathlineto{\pgfqpoint{0.967785in}{0.786836in}}%
\pgfpathlineto{\pgfqpoint{0.974731in}{0.790239in}}%
\pgfpathlineto{\pgfqpoint{0.982161in}{0.793450in}}%
\pgfpathclose%
\pgfusepath{fill}%
\end{pgfscope}%
\begin{pgfscope}%
\pgfpathrectangle{\pgfqpoint{0.407592in}{0.222904in}}{\pgfqpoint{1.411163in}{1.411163in}}%
\pgfusepath{clip}%
\pgfsetbuttcap%
\pgfsetroundjoin%
\definecolor{currentfill}{rgb}{0.749020,0.772549,0.792157}%
\pgfsetfillcolor{currentfill}%
\pgfsetfillopacity{0.160000}%
\pgfsetlinewidth{0.000000pt}%
\definecolor{currentstroke}{rgb}{0.000000,0.000000,0.000000}%
\pgfsetstrokecolor{currentstroke}%
\pgfsetdash{}{0pt}%
\pgfpathmoveto{\pgfqpoint{1.167553in}{0.733136in}}%
\pgfpathlineto{\pgfqpoint{1.177174in}{0.748392in}}%
\pgfpathlineto{\pgfqpoint{1.186787in}{0.763633in}}%
\pgfpathlineto{\pgfqpoint{1.180675in}{0.766128in}}%
\pgfpathlineto{\pgfqpoint{1.174180in}{0.768437in}}%
\pgfpathlineto{\pgfqpoint{1.165833in}{0.752738in}}%
\pgfpathlineto{\pgfqpoint{1.157476in}{0.737018in}}%
\pgfpathlineto{\pgfqpoint{1.162669in}{0.735151in}}%
\pgfpathlineto{\pgfqpoint{1.167553in}{0.733136in}}%
\pgfpathclose%
\pgfusepath{fill}%
\end{pgfscope}%
\begin{pgfscope}%
\pgfpathrectangle{\pgfqpoint{0.407592in}{0.222904in}}{\pgfqpoint{1.411163in}{1.411163in}}%
\pgfusepath{clip}%
\pgfsetbuttcap%
\pgfsetroundjoin%
\definecolor{currentfill}{rgb}{0.749020,0.772549,0.792157}%
\pgfsetfillcolor{currentfill}%
\pgfsetfillopacity{0.160000}%
\pgfsetlinewidth{0.000000pt}%
\definecolor{currentstroke}{rgb}{0.000000,0.000000,0.000000}%
\pgfsetstrokecolor{currentstroke}%
\pgfsetdash{}{0pt}%
\pgfpathmoveto{\pgfqpoint{1.028815in}{0.736656in}}%
\pgfpathlineto{\pgfqpoint{1.021182in}{0.752333in}}%
\pgfpathlineto{\pgfqpoint{1.013558in}{0.767989in}}%
\pgfpathlineto{\pgfqpoint{1.006646in}{0.765642in}}%
\pgfpathlineto{\pgfqpoint{1.000085in}{0.763112in}}%
\pgfpathlineto{\pgfqpoint{1.009060in}{0.747920in}}%
\pgfpathlineto{\pgfqpoint{1.018044in}{0.732714in}}%
\pgfpathlineto{\pgfqpoint{1.023288in}{0.734759in}}%
\pgfpathlineto{\pgfqpoint{1.028815in}{0.736656in}}%
\pgfpathclose%
\pgfusepath{fill}%
\end{pgfscope}%
\begin{pgfscope}%
\pgfpathrectangle{\pgfqpoint{0.407592in}{0.222904in}}{\pgfqpoint{1.411163in}{1.411163in}}%
\pgfusepath{clip}%
\pgfsetbuttcap%
\pgfsetroundjoin%
\definecolor{currentfill}{rgb}{0.749020,0.772549,0.792157}%
\pgfsetfillcolor{currentfill}%
\pgfsetfillopacity{0.160000}%
\pgfsetlinewidth{0.000000pt}%
\definecolor{currentstroke}{rgb}{0.000000,0.000000,0.000000}%
\pgfsetstrokecolor{currentstroke}%
\pgfsetdash{}{0pt}%
\pgfpathmoveto{\pgfqpoint{1.140729in}{0.705519in}}%
\pgfpathlineto{\pgfqpoint{1.149108in}{0.721279in}}%
\pgfpathlineto{\pgfqpoint{1.157476in}{0.737018in}}%
\pgfpathlineto{\pgfqpoint{1.151997in}{0.738729in}}%
\pgfpathlineto{\pgfqpoint{1.146255in}{0.740277in}}%
\pgfpathlineto{\pgfqpoint{1.139289in}{0.724145in}}%
\pgfpathlineto{\pgfqpoint{1.132312in}{0.707989in}}%
\pgfpathlineto{\pgfqpoint{1.136620in}{0.706815in}}%
\pgfpathlineto{\pgfqpoint{1.140729in}{0.705519in}}%
\pgfpathclose%
\pgfusepath{fill}%
\end{pgfscope}%
\begin{pgfscope}%
\pgfpathrectangle{\pgfqpoint{0.407592in}{0.222904in}}{\pgfqpoint{1.411163in}{1.411163in}}%
\pgfusepath{clip}%
\pgfsetbuttcap%
\pgfsetroundjoin%
\definecolor{currentfill}{rgb}{0.749020,0.772549,0.792157}%
\pgfsetfillcolor{currentfill}%
\pgfsetfillopacity{0.160000}%
\pgfsetlinewidth{0.000000pt}%
\definecolor{currentstroke}{rgb}{0.000000,0.000000,0.000000}%
\pgfsetstrokecolor{currentstroke}%
\pgfsetdash{}{0pt}%
\pgfpathmoveto{\pgfqpoint{0.886680in}{0.903614in}}%
\pgfpathlineto{\pgfqpoint{0.876572in}{0.918169in}}%
\pgfpathlineto{\pgfqpoint{0.866469in}{0.932715in}}%
\pgfpathlineto{\pgfqpoint{0.854509in}{0.926445in}}%
\pgfpathlineto{\pgfqpoint{0.843512in}{0.919876in}}%
\pgfpathlineto{\pgfqpoint{0.854693in}{0.905857in}}%
\pgfpathlineto{\pgfqpoint{0.865877in}{0.891835in}}%
\pgfpathlineto{\pgfqpoint{0.875838in}{0.897861in}}%
\pgfpathlineto{\pgfqpoint{0.886680in}{0.903614in}}%
\pgfpathclose%
\pgfusepath{fill}%
\end{pgfscope}%
\begin{pgfscope}%
\pgfpathrectangle{\pgfqpoint{0.407592in}{0.222904in}}{\pgfqpoint{1.411163in}{1.411163in}}%
\pgfusepath{clip}%
\pgfsetbuttcap%
\pgfsetroundjoin%
\definecolor{currentfill}{rgb}{0.749020,0.772549,0.792157}%
\pgfsetfillcolor{currentfill}%
\pgfsetfillopacity{0.160000}%
\pgfsetlinewidth{0.000000pt}%
\definecolor{currentstroke}{rgb}{0.000000,0.000000,0.000000}%
\pgfsetstrokecolor{currentstroke}%
\pgfsetdash{}{0pt}%
\pgfpathmoveto{\pgfqpoint{1.099008in}{0.678446in}}%
\pgfpathlineto{\pgfqpoint{1.101192in}{0.695342in}}%
\pgfpathlineto{\pgfqpoint{1.103372in}{0.712206in}}%
\pgfpathlineto{\pgfqpoint{1.098226in}{0.712411in}}%
\pgfpathlineto{\pgfqpoint{1.093045in}{0.712470in}}%
\pgfpathlineto{\pgfqpoint{1.092578in}{0.695563in}}%
\pgfpathlineto{\pgfqpoint{1.092110in}{0.678624in}}%
\pgfpathlineto{\pgfqpoint{1.095571in}{0.678584in}}%
\pgfpathlineto{\pgfqpoint{1.099008in}{0.678446in}}%
\pgfpathclose%
\pgfusepath{fill}%
\end{pgfscope}%
\begin{pgfscope}%
\pgfpathrectangle{\pgfqpoint{0.407592in}{0.222904in}}{\pgfqpoint{1.411163in}{1.411163in}}%
\pgfusepath{clip}%
\pgfsetbuttcap%
\pgfsetroundjoin%
\definecolor{currentfill}{rgb}{0.749020,0.772549,0.792157}%
\pgfsetfillcolor{currentfill}%
\pgfsetfillopacity{0.160000}%
\pgfsetlinewidth{0.000000pt}%
\definecolor{currentstroke}{rgb}{0.000000,0.000000,0.000000}%
\pgfsetstrokecolor{currentstroke}%
\pgfsetdash{}{0pt}%
\pgfpathmoveto{\pgfqpoint{1.092110in}{0.678624in}}%
\pgfpathlineto{\pgfqpoint{1.092578in}{0.695563in}}%
\pgfpathlineto{\pgfqpoint{1.093045in}{0.712470in}}%
\pgfpathlineto{\pgfqpoint{1.087851in}{0.712382in}}%
\pgfpathlineto{\pgfqpoint{1.082666in}{0.712148in}}%
\pgfpathlineto{\pgfqpoint{1.083920in}{0.695293in}}%
\pgfpathlineto{\pgfqpoint{1.085177in}{0.678406in}}%
\pgfpathlineto{\pgfqpoint{1.088641in}{0.678564in}}%
\pgfpathlineto{\pgfqpoint{1.092110in}{0.678624in}}%
\pgfpathclose%
\pgfusepath{fill}%
\end{pgfscope}%
\begin{pgfscope}%
\pgfpathrectangle{\pgfqpoint{0.407592in}{0.222904in}}{\pgfqpoint{1.411163in}{1.411163in}}%
\pgfusepath{clip}%
\pgfsetbuttcap%
\pgfsetroundjoin%
\definecolor{currentfill}{rgb}{0.749020,0.772549,0.792157}%
\pgfsetfillcolor{currentfill}%
\pgfsetfillopacity{0.160000}%
\pgfsetlinewidth{0.000000pt}%
\definecolor{currentstroke}{rgb}{0.000000,0.000000,0.000000}%
\pgfsetstrokecolor{currentstroke}%
\pgfsetdash{}{0pt}%
\pgfpathmoveto{\pgfqpoint{1.052967in}{0.707764in}}%
\pgfpathlineto{\pgfqpoint{1.046790in}{0.723884in}}%
\pgfpathlineto{\pgfqpoint{1.040621in}{0.739980in}}%
\pgfpathlineto{\pgfqpoint{1.034601in}{0.738399in}}%
\pgfpathlineto{\pgfqpoint{1.028815in}{0.736656in}}%
\pgfpathlineto{\pgfqpoint{1.036458in}{0.720960in}}%
\pgfpathlineto{\pgfqpoint{1.044111in}{0.705244in}}%
\pgfpathlineto{\pgfqpoint{1.048450in}{0.706565in}}%
\pgfpathlineto{\pgfqpoint{1.052967in}{0.707764in}}%
\pgfpathclose%
\pgfusepath{fill}%
\end{pgfscope}%
\begin{pgfscope}%
\pgfpathrectangle{\pgfqpoint{0.407592in}{0.222904in}}{\pgfqpoint{1.411163in}{1.411163in}}%
\pgfusepath{clip}%
\pgfsetbuttcap%
\pgfsetroundjoin%
\definecolor{currentfill}{rgb}{0.749020,0.772549,0.792157}%
\pgfsetfillcolor{currentfill}%
\pgfsetfillopacity{0.160000}%
\pgfsetlinewidth{0.000000pt}%
\definecolor{currentstroke}{rgb}{0.000000,0.000000,0.000000}%
\pgfsetstrokecolor{currentstroke}%
\pgfsetdash{}{0pt}%
\pgfpathmoveto{\pgfqpoint{1.299953in}{0.864894in}}%
\pgfpathlineto{\pgfqpoint{1.311569in}{0.878981in}}%
\pgfpathlineto{\pgfqpoint{1.323181in}{0.893064in}}%
\pgfpathlineto{\pgfqpoint{1.314333in}{0.899038in}}%
\pgfpathlineto{\pgfqpoint{1.304554in}{0.904734in}}%
\pgfpathlineto{\pgfqpoint{1.293907in}{0.890119in}}%
\pgfpathlineto{\pgfqpoint{1.283252in}{0.875494in}}%
\pgfpathlineto{\pgfqpoint{1.292025in}{0.870318in}}%
\pgfpathlineto{\pgfqpoint{1.299953in}{0.864894in}}%
\pgfpathclose%
\pgfusepath{fill}%
\end{pgfscope}%
\begin{pgfscope}%
\pgfpathrectangle{\pgfqpoint{0.407592in}{0.222904in}}{\pgfqpoint{1.411163in}{1.411163in}}%
\pgfusepath{clip}%
\pgfsetbuttcap%
\pgfsetroundjoin%
\definecolor{currentfill}{rgb}{0.749020,0.772549,0.792157}%
\pgfsetfillcolor{currentfill}%
\pgfsetfillopacity{0.160000}%
\pgfsetlinewidth{0.000000pt}%
\definecolor{currentstroke}{rgb}{0.000000,0.000000,0.000000}%
\pgfsetstrokecolor{currentstroke}%
\pgfsetdash{}{0pt}%
\pgfpathmoveto{\pgfqpoint{1.105754in}{0.677875in}}%
\pgfpathlineto{\pgfqpoint{1.109617in}{0.694633in}}%
\pgfpathlineto{\pgfqpoint{1.113473in}{0.711361in}}%
\pgfpathlineto{\pgfqpoint{1.108462in}{0.711856in}}%
\pgfpathlineto{\pgfqpoint{1.103372in}{0.712206in}}%
\pgfpathlineto{\pgfqpoint{1.101192in}{0.695342in}}%
\pgfpathlineto{\pgfqpoint{1.099008in}{0.678446in}}%
\pgfpathlineto{\pgfqpoint{1.102407in}{0.678209in}}%
\pgfpathlineto{\pgfqpoint{1.105754in}{0.677875in}}%
\pgfpathclose%
\pgfusepath{fill}%
\end{pgfscope}%
\begin{pgfscope}%
\pgfpathrectangle{\pgfqpoint{0.407592in}{0.222904in}}{\pgfqpoint{1.411163in}{1.411163in}}%
\pgfusepath{clip}%
\pgfsetbuttcap%
\pgfsetroundjoin%
\definecolor{currentfill}{rgb}{0.749020,0.772549,0.792157}%
\pgfsetfillcolor{currentfill}%
\pgfsetfillopacity{0.160000}%
\pgfsetlinewidth{0.000000pt}%
\definecolor{currentstroke}{rgb}{0.000000,0.000000,0.000000}%
\pgfsetstrokecolor{currentstroke}%
\pgfsetdash{}{0pt}%
\pgfpathmoveto{\pgfqpoint{1.085177in}{0.678406in}}%
\pgfpathlineto{\pgfqpoint{1.083920in}{0.695293in}}%
\pgfpathlineto{\pgfqpoint{1.082666in}{0.712148in}}%
\pgfpathlineto{\pgfqpoint{1.077512in}{0.711768in}}%
\pgfpathlineto{\pgfqpoint{1.072411in}{0.711244in}}%
\pgfpathlineto{\pgfqpoint{1.075367in}{0.694536in}}%
\pgfpathlineto{\pgfqpoint{1.078329in}{0.677797in}}%
\pgfpathlineto{\pgfqpoint{1.081735in}{0.678150in}}%
\pgfpathlineto{\pgfqpoint{1.085177in}{0.678406in}}%
\pgfpathclose%
\pgfusepath{fill}%
\end{pgfscope}%
\begin{pgfscope}%
\pgfpathrectangle{\pgfqpoint{0.407592in}{0.222904in}}{\pgfqpoint{1.411163in}{1.411163in}}%
\pgfusepath{clip}%
\pgfsetbuttcap%
\pgfsetroundjoin%
\definecolor{currentfill}{rgb}{0.749020,0.772549,0.792157}%
\pgfsetfillcolor{currentfill}%
\pgfsetfillopacity{0.160000}%
\pgfsetlinewidth{0.000000pt}%
\definecolor{currentstroke}{rgb}{0.000000,0.000000,0.000000}%
\pgfsetstrokecolor{currentstroke}%
\pgfsetdash{}{0pt}%
\pgfpathmoveto{\pgfqpoint{1.197756in}{0.758134in}}%
\pgfpathlineto{\pgfqpoint{1.208469in}{0.772839in}}%
\pgfpathlineto{\pgfqpoint{1.219174in}{0.787533in}}%
\pgfpathlineto{\pgfqpoint{1.212837in}{0.790899in}}%
\pgfpathlineto{\pgfqpoint{1.205982in}{0.794070in}}%
\pgfpathlineto{\pgfqpoint{1.196389in}{0.778859in}}%
\pgfpathlineto{\pgfqpoint{1.186787in}{0.763633in}}%
\pgfpathlineto{\pgfqpoint{1.192489in}{0.760965in}}%
\pgfpathlineto{\pgfqpoint{1.197756in}{0.758134in}}%
\pgfpathclose%
\pgfusepath{fill}%
\end{pgfscope}%
\begin{pgfscope}%
\pgfpathrectangle{\pgfqpoint{0.407592in}{0.222904in}}{\pgfqpoint{1.411163in}{1.411163in}}%
\pgfusepath{clip}%
\pgfsetbuttcap%
\pgfsetroundjoin%
\definecolor{currentfill}{rgb}{0.749020,0.772549,0.792157}%
\pgfsetfillcolor{currentfill}%
\pgfsetfillopacity{0.160000}%
\pgfsetlinewidth{0.000000pt}%
\definecolor{currentstroke}{rgb}{0.000000,0.000000,0.000000}%
\pgfsetstrokecolor{currentstroke}%
\pgfsetdash{}{0pt}%
\pgfpathmoveto{\pgfqpoint{1.000085in}{0.763112in}}%
\pgfpathlineto{\pgfqpoint{0.991119in}{0.778288in}}%
\pgfpathlineto{\pgfqpoint{0.982161in}{0.793450in}}%
\pgfpathlineto{\pgfqpoint{0.974731in}{0.790239in}}%
\pgfpathlineto{\pgfqpoint{0.967785in}{0.786836in}}%
\pgfpathlineto{\pgfqpoint{0.977952in}{0.772196in}}%
\pgfpathlineto{\pgfqpoint{0.988126in}{0.757547in}}%
\pgfpathlineto{\pgfqpoint{0.993903in}{0.760410in}}%
\pgfpathlineto{\pgfqpoint{1.000085in}{0.763112in}}%
\pgfpathclose%
\pgfusepath{fill}%
\end{pgfscope}%
\begin{pgfscope}%
\pgfpathrectangle{\pgfqpoint{0.407592in}{0.222904in}}{\pgfqpoint{1.411163in}{1.411163in}}%
\pgfusepath{clip}%
\pgfsetbuttcap%
\pgfsetroundjoin%
\definecolor{currentfill}{rgb}{0.749020,0.772549,0.792157}%
\pgfsetfillcolor{currentfill}%
\pgfsetfillopacity{0.160000}%
\pgfsetlinewidth{0.000000pt}%
\definecolor{currentstroke}{rgb}{0.000000,0.000000,0.000000}%
\pgfsetstrokecolor{currentstroke}%
\pgfsetdash{}{0pt}%
\pgfpathmoveto{\pgfqpoint{1.112231in}{0.676922in}}%
\pgfpathlineto{\pgfqpoint{1.117707in}{0.693449in}}%
\pgfpathlineto{\pgfqpoint{1.123175in}{0.709947in}}%
\pgfpathlineto{\pgfqpoint{1.118384in}{0.710724in}}%
\pgfpathlineto{\pgfqpoint{1.113473in}{0.711361in}}%
\pgfpathlineto{\pgfqpoint{1.109617in}{0.694633in}}%
\pgfpathlineto{\pgfqpoint{1.105754in}{0.677875in}}%
\pgfpathlineto{\pgfqpoint{1.109033in}{0.677446in}}%
\pgfpathlineto{\pgfqpoint{1.112231in}{0.676922in}}%
\pgfpathclose%
\pgfusepath{fill}%
\end{pgfscope}%
\begin{pgfscope}%
\pgfpathrectangle{\pgfqpoint{0.407592in}{0.222904in}}{\pgfqpoint{1.411163in}{1.411163in}}%
\pgfusepath{clip}%
\pgfsetbuttcap%
\pgfsetroundjoin%
\definecolor{currentfill}{rgb}{0.749020,0.772549,0.792157}%
\pgfsetfillcolor{currentfill}%
\pgfsetfillopacity{0.160000}%
\pgfsetlinewidth{0.000000pt}%
\definecolor{currentstroke}{rgb}{0.000000,0.000000,0.000000}%
\pgfsetstrokecolor{currentstroke}%
\pgfsetdash{}{0pt}%
\pgfpathmoveto{\pgfqpoint{0.906917in}{0.874476in}}%
\pgfpathlineto{\pgfqpoint{0.896795in}{0.889050in}}%
\pgfpathlineto{\pgfqpoint{0.886680in}{0.903614in}}%
\pgfpathlineto{\pgfqpoint{0.875838in}{0.897861in}}%
\pgfpathlineto{\pgfqpoint{0.865877in}{0.891835in}}%
\pgfpathlineto{\pgfqpoint{0.877064in}{0.877808in}}%
\pgfpathlineto{\pgfqpoint{0.888254in}{0.863777in}}%
\pgfpathlineto{\pgfqpoint{0.897186in}{0.869249in}}%
\pgfpathlineto{\pgfqpoint{0.906917in}{0.874476in}}%
\pgfpathclose%
\pgfusepath{fill}%
\end{pgfscope}%
\begin{pgfscope}%
\pgfpathrectangle{\pgfqpoint{0.407592in}{0.222904in}}{\pgfqpoint{1.411163in}{1.411163in}}%
\pgfusepath{clip}%
\pgfsetbuttcap%
\pgfsetroundjoin%
\definecolor{currentfill}{rgb}{0.749020,0.772549,0.792157}%
\pgfsetfillcolor{currentfill}%
\pgfsetfillopacity{0.160000}%
\pgfsetlinewidth{0.000000pt}%
\definecolor{currentstroke}{rgb}{0.000000,0.000000,0.000000}%
\pgfsetstrokecolor{currentstroke}%
\pgfsetdash{}{0pt}%
\pgfpathmoveto{\pgfqpoint{1.078329in}{0.677797in}}%
\pgfpathlineto{\pgfqpoint{1.075367in}{0.694536in}}%
\pgfpathlineto{\pgfqpoint{1.072411in}{0.711244in}}%
\pgfpathlineto{\pgfqpoint{1.067385in}{0.710579in}}%
\pgfpathlineto{\pgfqpoint{1.062455in}{0.709775in}}%
\pgfpathlineto{\pgfqpoint{1.067065in}{0.693304in}}%
\pgfpathlineto{\pgfqpoint{1.071682in}{0.676806in}}%
\pgfpathlineto{\pgfqpoint{1.074973in}{0.677348in}}%
\pgfpathlineto{\pgfqpoint{1.078329in}{0.677797in}}%
\pgfpathclose%
\pgfusepath{fill}%
\end{pgfscope}%
\begin{pgfscope}%
\pgfpathrectangle{\pgfqpoint{0.407592in}{0.222904in}}{\pgfqpoint{1.411163in}{1.411163in}}%
\pgfusepath{clip}%
\pgfsetbuttcap%
\pgfsetroundjoin%
\definecolor{currentfill}{rgb}{0.749020,0.772549,0.792157}%
\pgfsetfillcolor{currentfill}%
\pgfsetfillopacity{0.160000}%
\pgfsetlinewidth{0.000000pt}%
\definecolor{currentstroke}{rgb}{0.000000,0.000000,0.000000}%
\pgfsetstrokecolor{currentstroke}%
\pgfsetdash{}{0pt}%
\pgfpathmoveto{\pgfqpoint{1.276711in}{0.836705in}}%
\pgfpathlineto{\pgfqpoint{1.288334in}{0.850802in}}%
\pgfpathlineto{\pgfqpoint{1.299953in}{0.864894in}}%
\pgfpathlineto{\pgfqpoint{1.292025in}{0.870318in}}%
\pgfpathlineto{\pgfqpoint{1.283252in}{0.875494in}}%
\pgfpathlineto{\pgfqpoint{1.272591in}{0.860858in}}%
\pgfpathlineto{\pgfqpoint{1.261922in}{0.846213in}}%
\pgfpathlineto{\pgfqpoint{1.269694in}{0.841570in}}%
\pgfpathlineto{\pgfqpoint{1.276711in}{0.836705in}}%
\pgfpathclose%
\pgfusepath{fill}%
\end{pgfscope}%
\begin{pgfscope}%
\pgfpathrectangle{\pgfqpoint{0.407592in}{0.222904in}}{\pgfqpoint{1.411163in}{1.411163in}}%
\pgfusepath{clip}%
\pgfsetbuttcap%
\pgfsetroundjoin%
\definecolor{currentfill}{rgb}{0.749020,0.772549,0.792157}%
\pgfsetfillcolor{currentfill}%
\pgfsetfillopacity{0.160000}%
\pgfsetlinewidth{0.000000pt}%
\definecolor{currentstroke}{rgb}{0.000000,0.000000,0.000000}%
\pgfsetstrokecolor{currentstroke}%
\pgfsetdash{}{0pt}%
\pgfpathmoveto{\pgfqpoint{1.118327in}{0.675602in}}%
\pgfpathlineto{\pgfqpoint{1.125325in}{0.691808in}}%
\pgfpathlineto{\pgfqpoint{1.132312in}{0.707989in}}%
\pgfpathlineto{\pgfqpoint{1.127824in}{0.709034in}}%
\pgfpathlineto{\pgfqpoint{1.123175in}{0.709947in}}%
\pgfpathlineto{\pgfqpoint{1.117707in}{0.693449in}}%
\pgfpathlineto{\pgfqpoint{1.112231in}{0.676922in}}%
\pgfpathlineto{\pgfqpoint{1.115333in}{0.676307in}}%
\pgfpathlineto{\pgfqpoint{1.118327in}{0.675602in}}%
\pgfpathclose%
\pgfusepath{fill}%
\end{pgfscope}%
\begin{pgfscope}%
\pgfpathrectangle{\pgfqpoint{0.407592in}{0.222904in}}{\pgfqpoint{1.411163in}{1.411163in}}%
\pgfusepath{clip}%
\pgfsetbuttcap%
\pgfsetroundjoin%
\definecolor{currentfill}{rgb}{0.749020,0.772549,0.792157}%
\pgfsetfillcolor{currentfill}%
\pgfsetfillopacity{0.160000}%
\pgfsetlinewidth{0.000000pt}%
\definecolor{currentstroke}{rgb}{0.000000,0.000000,0.000000}%
\pgfsetstrokecolor{currentstroke}%
\pgfsetdash{}{0pt}%
\pgfpathmoveto{\pgfqpoint{1.148280in}{0.702577in}}%
\pgfpathlineto{\pgfqpoint{1.157921in}{0.717864in}}%
\pgfpathlineto{\pgfqpoint{1.167553in}{0.733136in}}%
\pgfpathlineto{\pgfqpoint{1.162669in}{0.735151in}}%
\pgfpathlineto{\pgfqpoint{1.157476in}{0.737018in}}%
\pgfpathlineto{\pgfqpoint{1.149108in}{0.721279in}}%
\pgfpathlineto{\pgfqpoint{1.140729in}{0.705519in}}%
\pgfpathlineto{\pgfqpoint{1.144621in}{0.704104in}}%
\pgfpathlineto{\pgfqpoint{1.148280in}{0.702577in}}%
\pgfpathclose%
\pgfusepath{fill}%
\end{pgfscope}%
\begin{pgfscope}%
\pgfpathrectangle{\pgfqpoint{0.407592in}{0.222904in}}{\pgfqpoint{1.411163in}{1.411163in}}%
\pgfusepath{clip}%
\pgfsetbuttcap%
\pgfsetroundjoin%
\definecolor{currentfill}{rgb}{0.749020,0.772549,0.792157}%
\pgfsetfillcolor{currentfill}%
\pgfsetfillopacity{0.160000}%
\pgfsetlinewidth{0.000000pt}%
\definecolor{currentstroke}{rgb}{0.000000,0.000000,0.000000}%
\pgfsetstrokecolor{currentstroke}%
\pgfsetdash{}{0pt}%
\pgfpathmoveto{\pgfqpoint{1.071682in}{0.676806in}}%
\pgfpathlineto{\pgfqpoint{1.067065in}{0.693304in}}%
\pgfpathlineto{\pgfqpoint{1.062455in}{0.709775in}}%
\pgfpathlineto{\pgfqpoint{1.057642in}{0.708835in}}%
\pgfpathlineto{\pgfqpoint{1.052967in}{0.707764in}}%
\pgfpathlineto{\pgfqpoint{1.059155in}{0.691619in}}%
\pgfpathlineto{\pgfqpoint{1.065351in}{0.675451in}}%
\pgfpathlineto{\pgfqpoint{1.068470in}{0.676173in}}%
\pgfpathlineto{\pgfqpoint{1.071682in}{0.676806in}}%
\pgfpathclose%
\pgfusepath{fill}%
\end{pgfscope}%
\begin{pgfscope}%
\pgfpathrectangle{\pgfqpoint{0.407592in}{0.222904in}}{\pgfqpoint{1.411163in}{1.411163in}}%
\pgfusepath{clip}%
\pgfsetbuttcap%
\pgfsetroundjoin%
\definecolor{currentfill}{rgb}{0.749020,0.772549,0.792157}%
\pgfsetfillcolor{currentfill}%
\pgfsetfillopacity{0.160000}%
\pgfsetlinewidth{0.000000pt}%
\definecolor{currentstroke}{rgb}{0.000000,0.000000,0.000000}%
\pgfsetstrokecolor{currentstroke}%
\pgfsetdash{}{0pt}%
\pgfpathmoveto{\pgfqpoint{0.927180in}{0.845300in}}%
\pgfpathlineto{\pgfqpoint{0.917045in}{0.859893in}}%
\pgfpathlineto{\pgfqpoint{0.906917in}{0.874476in}}%
\pgfpathlineto{\pgfqpoint{0.897186in}{0.869249in}}%
\pgfpathlineto{\pgfqpoint{0.888254in}{0.863777in}}%
\pgfpathlineto{\pgfqpoint{0.899447in}{0.849743in}}%
\pgfpathlineto{\pgfqpoint{0.910644in}{0.835704in}}%
\pgfpathlineto{\pgfqpoint{0.918555in}{0.840611in}}%
\pgfpathlineto{\pgfqpoint{0.927180in}{0.845300in}}%
\pgfpathclose%
\pgfusepath{fill}%
\end{pgfscope}%
\begin{pgfscope}%
\pgfpathrectangle{\pgfqpoint{0.407592in}{0.222904in}}{\pgfqpoint{1.411163in}{1.411163in}}%
\pgfusepath{clip}%
\pgfsetbuttcap%
\pgfsetroundjoin%
\definecolor{currentfill}{rgb}{0.749020,0.772549,0.792157}%
\pgfsetfillcolor{currentfill}%
\pgfsetfillopacity{0.160000}%
\pgfsetlinewidth{0.000000pt}%
\definecolor{currentstroke}{rgb}{0.000000,0.000000,0.000000}%
\pgfsetstrokecolor{currentstroke}%
\pgfsetdash{}{0pt}%
\pgfpathmoveto{\pgfqpoint{1.044111in}{0.705244in}}%
\pgfpathlineto{\pgfqpoint{1.036458in}{0.720960in}}%
\pgfpathlineto{\pgfqpoint{1.028815in}{0.736656in}}%
\pgfpathlineto{\pgfqpoint{1.023288in}{0.734759in}}%
\pgfpathlineto{\pgfqpoint{1.018044in}{0.732714in}}%
\pgfpathlineto{\pgfqpoint{1.027036in}{0.717494in}}%
\pgfpathlineto{\pgfqpoint{1.036038in}{0.702258in}}%
\pgfpathlineto{\pgfqpoint{1.039967in}{0.703806in}}%
\pgfpathlineto{\pgfqpoint{1.044111in}{0.705244in}}%
\pgfpathclose%
\pgfusepath{fill}%
\end{pgfscope}%
\begin{pgfscope}%
\pgfpathrectangle{\pgfqpoint{0.407592in}{0.222904in}}{\pgfqpoint{1.411163in}{1.411163in}}%
\pgfusepath{clip}%
\pgfsetbuttcap%
\pgfsetroundjoin%
\definecolor{currentfill}{rgb}{0.749020,0.772549,0.792157}%
\pgfsetfillcolor{currentfill}%
\pgfsetfillopacity{0.160000}%
\pgfsetlinewidth{0.000000pt}%
\definecolor{currentstroke}{rgb}{0.000000,0.000000,0.000000}%
\pgfsetstrokecolor{currentstroke}%
\pgfsetdash{}{0pt}%
\pgfpathmoveto{\pgfqpoint{1.176309in}{0.728694in}}%
\pgfpathlineto{\pgfqpoint{1.187036in}{0.743419in}}%
\pgfpathlineto{\pgfqpoint{1.197756in}{0.758134in}}%
\pgfpathlineto{\pgfqpoint{1.192489in}{0.760965in}}%
\pgfpathlineto{\pgfqpoint{1.186787in}{0.763633in}}%
\pgfpathlineto{\pgfqpoint{1.177174in}{0.748392in}}%
\pgfpathlineto{\pgfqpoint{1.167553in}{0.733136in}}%
\pgfpathlineto{\pgfqpoint{1.172106in}{0.730980in}}%
\pgfpathlineto{\pgfqpoint{1.176309in}{0.728694in}}%
\pgfpathclose%
\pgfusepath{fill}%
\end{pgfscope}%
\begin{pgfscope}%
\pgfpathrectangle{\pgfqpoint{0.407592in}{0.222904in}}{\pgfqpoint{1.411163in}{1.411163in}}%
\pgfusepath{clip}%
\pgfsetbuttcap%
\pgfsetroundjoin%
\definecolor{currentfill}{rgb}{0.749020,0.772549,0.792157}%
\pgfsetfillcolor{currentfill}%
\pgfsetfillopacity{0.160000}%
\pgfsetlinewidth{0.000000pt}%
\definecolor{currentstroke}{rgb}{0.000000,0.000000,0.000000}%
\pgfsetstrokecolor{currentstroke}%
\pgfsetdash{}{0pt}%
\pgfpathmoveto{\pgfqpoint{1.253454in}{0.808498in}}%
\pgfpathlineto{\pgfqpoint{1.265085in}{0.822604in}}%
\pgfpathlineto{\pgfqpoint{1.276711in}{0.836705in}}%
\pgfpathlineto{\pgfqpoint{1.269694in}{0.841570in}}%
\pgfpathlineto{\pgfqpoint{1.261922in}{0.846213in}}%
\pgfpathlineto{\pgfqpoint{1.251246in}{0.831558in}}%
\pgfpathlineto{\pgfqpoint{1.240562in}{0.816893in}}%
\pgfpathlineto{\pgfqpoint{1.247340in}{0.812793in}}%
\pgfpathlineto{\pgfqpoint{1.253454in}{0.808498in}}%
\pgfpathclose%
\pgfusepath{fill}%
\end{pgfscope}%
\begin{pgfscope}%
\pgfpathrectangle{\pgfqpoint{0.407592in}{0.222904in}}{\pgfqpoint{1.411163in}{1.411163in}}%
\pgfusepath{clip}%
\pgfsetbuttcap%
\pgfsetroundjoin%
\definecolor{currentfill}{rgb}{0.749020,0.772549,0.792157}%
\pgfsetfillcolor{currentfill}%
\pgfsetfillopacity{0.160000}%
\pgfsetlinewidth{0.000000pt}%
\definecolor{currentstroke}{rgb}{0.000000,0.000000,0.000000}%
\pgfsetstrokecolor{currentstroke}%
\pgfsetdash{}{0pt}%
\pgfpathmoveto{\pgfqpoint{1.784494in}{1.366922in}}%
\pgfpathlineto{\pgfqpoint{1.796912in}{1.380452in}}%
\pgfpathlineto{\pgfqpoint{1.774857in}{1.396026in}}%
\pgfpathlineto{\pgfqpoint{1.750178in}{1.410942in}}%
\pgfpathlineto{\pgfqpoint{1.738704in}{1.397027in}}%
\pgfpathlineto{\pgfqpoint{1.762892in}{1.382298in}}%
\pgfusepath{fill}%
\end{pgfscope}%
\begin{pgfscope}%
\pgfpathrectangle{\pgfqpoint{0.407592in}{0.222904in}}{\pgfqpoint{1.411163in}{1.411163in}}%
\pgfusepath{clip}%
\pgfsetbuttcap%
\pgfsetroundjoin%
\definecolor{currentfill}{rgb}{0.749020,0.772549,0.792157}%
\pgfsetfillcolor{currentfill}%
\pgfsetfillopacity{0.160000}%
\pgfsetlinewidth{0.000000pt}%
\definecolor{currentstroke}{rgb}{0.000000,0.000000,0.000000}%
\pgfsetstrokecolor{currentstroke}%
\pgfsetdash{}{0pt}%
\pgfpathmoveto{\pgfqpoint{1.759659in}{1.339865in}}%
\pgfpathlineto{\pgfqpoint{1.772076in}{1.353393in}}%
\pgfpathlineto{\pgfqpoint{1.784494in}{1.366922in}}%
\pgfpathlineto{\pgfqpoint{1.762892in}{1.382298in}}%
\pgfpathlineto{\pgfqpoint{1.738704in}{1.397027in}}%
\pgfpathlineto{\pgfqpoint{1.727227in}{1.383107in}}%
\pgfpathlineto{\pgfqpoint{1.715746in}{1.369183in}}%
\pgfpathlineto{\pgfqpoint{1.738956in}{1.354836in}}%
\pgfpathlineto{\pgfqpoint{1.759659in}{1.339865in}}%
\pgfpathclose%
\pgfusepath{fill}%
\end{pgfscope}%
\begin{pgfscope}%
\pgfpathrectangle{\pgfqpoint{0.407592in}{0.222904in}}{\pgfqpoint{1.411163in}{1.411163in}}%
\pgfusepath{clip}%
\pgfsetbuttcap%
\pgfsetroundjoin%
\definecolor{currentfill}{rgb}{0.749020,0.772549,0.792157}%
\pgfsetfillcolor{currentfill}%
\pgfsetfillopacity{0.160000}%
\pgfsetlinewidth{0.000000pt}%
\definecolor{currentstroke}{rgb}{0.000000,0.000000,0.000000}%
\pgfsetstrokecolor{currentstroke}%
\pgfsetdash{}{0pt}%
\pgfpathmoveto{\pgfqpoint{1.018044in}{0.732714in}}%
\pgfpathlineto{\pgfqpoint{1.009060in}{0.747920in}}%
\pgfpathlineto{\pgfqpoint{1.000085in}{0.763112in}}%
\pgfpathlineto{\pgfqpoint{0.993903in}{0.760410in}}%
\pgfpathlineto{\pgfqpoint{0.988126in}{0.757547in}}%
\pgfpathlineto{\pgfqpoint{0.998307in}{0.742888in}}%
\pgfpathlineto{\pgfqpoint{1.008494in}{0.728220in}}%
\pgfpathlineto{\pgfqpoint{1.013105in}{0.730532in}}%
\pgfpathlineto{\pgfqpoint{1.018044in}{0.732714in}}%
\pgfpathclose%
\pgfusepath{fill}%
\end{pgfscope}%
\begin{pgfscope}%
\pgfpathrectangle{\pgfqpoint{0.407592in}{0.222904in}}{\pgfqpoint{1.411163in}{1.411163in}}%
\pgfusepath{clip}%
\pgfsetbuttcap%
\pgfsetroundjoin%
\definecolor{currentfill}{rgb}{0.749020,0.772549,0.792157}%
\pgfsetfillcolor{currentfill}%
\pgfsetfillopacity{0.160000}%
\pgfsetlinewidth{0.000000pt}%
\definecolor{currentstroke}{rgb}{0.000000,0.000000,0.000000}%
\pgfsetstrokecolor{currentstroke}%
\pgfsetdash{}{0pt}%
\pgfpathmoveto{\pgfqpoint{1.734827in}{1.312811in}}%
\pgfpathlineto{\pgfqpoint{1.747242in}{1.326337in}}%
\pgfpathlineto{\pgfqpoint{1.759659in}{1.339865in}}%
\pgfpathlineto{\pgfqpoint{1.738956in}{1.354836in}}%
\pgfpathlineto{\pgfqpoint{1.715746in}{1.369183in}}%
\pgfpathlineto{\pgfqpoint{1.704262in}{1.355254in}}%
\pgfpathlineto{\pgfqpoint{1.692774in}{1.341320in}}%
\pgfpathlineto{\pgfqpoint{1.715014in}{1.327366in}}%
\pgfpathlineto{\pgfqpoint{1.734827in}{1.312811in}}%
\pgfpathclose%
\pgfusepath{fill}%
\end{pgfscope}%
\begin{pgfscope}%
\pgfpathrectangle{\pgfqpoint{0.407592in}{0.222904in}}{\pgfqpoint{1.411163in}{1.411163in}}%
\pgfusepath{clip}%
\pgfsetbuttcap%
\pgfsetroundjoin%
\definecolor{currentfill}{rgb}{0.749020,0.772549,0.792157}%
\pgfsetfillcolor{currentfill}%
\pgfsetfillopacity{0.160000}%
\pgfsetlinewidth{0.000000pt}%
\definecolor{currentstroke}{rgb}{0.000000,0.000000,0.000000}%
\pgfsetstrokecolor{currentstroke}%
\pgfsetdash{}{0pt}%
\pgfpathmoveto{\pgfqpoint{1.123939in}{0.673938in}}%
\pgfpathlineto{\pgfqpoint{1.132340in}{0.689739in}}%
\pgfpathlineto{\pgfqpoint{1.140729in}{0.705519in}}%
\pgfpathlineto{\pgfqpoint{1.136620in}{0.706815in}}%
\pgfpathlineto{\pgfqpoint{1.132312in}{0.707989in}}%
\pgfpathlineto{\pgfqpoint{1.125325in}{0.691808in}}%
\pgfpathlineto{\pgfqpoint{1.118327in}{0.675602in}}%
\pgfpathlineto{\pgfqpoint{1.121200in}{0.674812in}}%
\pgfpathlineto{\pgfqpoint{1.123939in}{0.673938in}}%
\pgfpathclose%
\pgfusepath{fill}%
\end{pgfscope}%
\begin{pgfscope}%
\pgfpathrectangle{\pgfqpoint{0.407592in}{0.222904in}}{\pgfqpoint{1.411163in}{1.411163in}}%
\pgfusepath{clip}%
\pgfsetbuttcap%
\pgfsetroundjoin%
\definecolor{currentfill}{rgb}{0.749020,0.772549,0.792157}%
\pgfsetfillcolor{currentfill}%
\pgfsetfillopacity{0.160000}%
\pgfsetlinewidth{0.000000pt}%
\definecolor{currentstroke}{rgb}{0.000000,0.000000,0.000000}%
\pgfsetstrokecolor{currentstroke}%
\pgfsetdash{}{0pt}%
\pgfpathmoveto{\pgfqpoint{0.947469in}{0.816087in}}%
\pgfpathlineto{\pgfqpoint{0.937321in}{0.830698in}}%
\pgfpathlineto{\pgfqpoint{0.927180in}{0.845300in}}%
\pgfpathlineto{\pgfqpoint{0.918555in}{0.840611in}}%
\pgfpathlineto{\pgfqpoint{0.910644in}{0.835704in}}%
\pgfpathlineto{\pgfqpoint{0.921844in}{0.821661in}}%
\pgfpathlineto{\pgfqpoint{0.933047in}{0.807614in}}%
\pgfpathlineto{\pgfqpoint{0.939943in}{0.811946in}}%
\pgfpathlineto{\pgfqpoint{0.947469in}{0.816087in}}%
\pgfpathclose%
\pgfusepath{fill}%
\end{pgfscope}%
\begin{pgfscope}%
\pgfpathrectangle{\pgfqpoint{0.407592in}{0.222904in}}{\pgfqpoint{1.411163in}{1.411163in}}%
\pgfusepath{clip}%
\pgfsetbuttcap%
\pgfsetroundjoin%
\definecolor{currentfill}{rgb}{0.749020,0.772549,0.792157}%
\pgfsetfillcolor{currentfill}%
\pgfsetfillopacity{0.160000}%
\pgfsetlinewidth{0.000000pt}%
\definecolor{currentstroke}{rgb}{0.000000,0.000000,0.000000}%
\pgfsetstrokecolor{currentstroke}%
\pgfsetdash{}{0pt}%
\pgfpathmoveto{\pgfqpoint{1.065351in}{0.675451in}}%
\pgfpathlineto{\pgfqpoint{1.059155in}{0.691619in}}%
\pgfpathlineto{\pgfqpoint{1.052967in}{0.707764in}}%
\pgfpathlineto{\pgfqpoint{1.048450in}{0.706565in}}%
\pgfpathlineto{\pgfqpoint{1.044111in}{0.705244in}}%
\pgfpathlineto{\pgfqpoint{1.051773in}{0.689509in}}%
\pgfpathlineto{\pgfqpoint{1.059445in}{0.673753in}}%
\pgfpathlineto{\pgfqpoint{1.062338in}{0.674643in}}%
\pgfpathlineto{\pgfqpoint{1.065351in}{0.675451in}}%
\pgfpathclose%
\pgfusepath{fill}%
\end{pgfscope}%
\begin{pgfscope}%
\pgfpathrectangle{\pgfqpoint{0.407592in}{0.222904in}}{\pgfqpoint{1.411163in}{1.411163in}}%
\pgfusepath{clip}%
\pgfsetbuttcap%
\pgfsetroundjoin%
\definecolor{currentfill}{rgb}{0.749020,0.772549,0.792157}%
\pgfsetfillcolor{currentfill}%
\pgfsetfillopacity{0.160000}%
\pgfsetlinewidth{0.000000pt}%
\definecolor{currentstroke}{rgb}{0.000000,0.000000,0.000000}%
\pgfsetstrokecolor{currentstroke}%
\pgfsetdash{}{0pt}%
\pgfpathmoveto{\pgfqpoint{1.709998in}{1.285760in}}%
\pgfpathlineto{\pgfqpoint{1.722412in}{1.299285in}}%
\pgfpathlineto{\pgfqpoint{1.734827in}{1.312811in}}%
\pgfpathlineto{\pgfqpoint{1.715014in}{1.327366in}}%
\pgfpathlineto{\pgfqpoint{1.692774in}{1.341320in}}%
\pgfpathlineto{\pgfqpoint{1.681282in}{1.327382in}}%
\pgfpathlineto{\pgfqpoint{1.669786in}{1.313440in}}%
\pgfpathlineto{\pgfqpoint{1.691066in}{1.299889in}}%
\pgfpathlineto{\pgfqpoint{1.709998in}{1.285760in}}%
\pgfpathclose%
\pgfusepath{fill}%
\end{pgfscope}%
\begin{pgfscope}%
\pgfpathrectangle{\pgfqpoint{0.407592in}{0.222904in}}{\pgfqpoint{1.411163in}{1.411163in}}%
\pgfusepath{clip}%
\pgfsetbuttcap%
\pgfsetroundjoin%
\definecolor{currentfill}{rgb}{0.749020,0.772549,0.792157}%
\pgfsetfillcolor{currentfill}%
\pgfsetfillopacity{0.160000}%
\pgfsetlinewidth{0.000000pt}%
\definecolor{currentstroke}{rgb}{0.000000,0.000000,0.000000}%
\pgfsetstrokecolor{currentstroke}%
\pgfsetdash{}{0pt}%
\pgfpathmoveto{\pgfqpoint{1.230182in}{0.780272in}}%
\pgfpathlineto{\pgfqpoint{1.241820in}{0.794387in}}%
\pgfpathlineto{\pgfqpoint{1.253454in}{0.808498in}}%
\pgfpathlineto{\pgfqpoint{1.247340in}{0.812793in}}%
\pgfpathlineto{\pgfqpoint{1.240562in}{0.816893in}}%
\pgfpathlineto{\pgfqpoint{1.229872in}{0.802218in}}%
\pgfpathlineto{\pgfqpoint{1.219174in}{0.787533in}}%
\pgfpathlineto{\pgfqpoint{1.224964in}{0.783986in}}%
\pgfpathlineto{\pgfqpoint{1.230182in}{0.780272in}}%
\pgfpathclose%
\pgfusepath{fill}%
\end{pgfscope}%
\begin{pgfscope}%
\pgfpathrectangle{\pgfqpoint{0.407592in}{0.222904in}}{\pgfqpoint{1.411163in}{1.411163in}}%
\pgfusepath{clip}%
\pgfsetbuttcap%
\pgfsetroundjoin%
\definecolor{currentfill}{rgb}{0.749020,0.772549,0.792157}%
\pgfsetfillcolor{currentfill}%
\pgfsetfillopacity{0.160000}%
\pgfsetlinewidth{0.000000pt}%
\definecolor{currentstroke}{rgb}{0.000000,0.000000,0.000000}%
\pgfsetstrokecolor{currentstroke}%
\pgfsetdash{}{0pt}%
\pgfpathmoveto{\pgfqpoint{1.685171in}{1.258712in}}%
\pgfpathlineto{\pgfqpoint{1.697584in}{1.272236in}}%
\pgfpathlineto{\pgfqpoint{1.709998in}{1.285760in}}%
\pgfpathlineto{\pgfqpoint{1.691066in}{1.299889in}}%
\pgfpathlineto{\pgfqpoint{1.669786in}{1.313440in}}%
\pgfpathlineto{\pgfqpoint{1.658286in}{1.299493in}}%
\pgfpathlineto{\pgfqpoint{1.646783in}{1.285542in}}%
\pgfpathlineto{\pgfqpoint{1.667110in}{1.272405in}}%
\pgfpathlineto{\pgfqpoint{1.685171in}{1.258712in}}%
\pgfpathclose%
\pgfusepath{fill}%
\end{pgfscope}%
\begin{pgfscope}%
\pgfpathrectangle{\pgfqpoint{0.407592in}{0.222904in}}{\pgfqpoint{1.411163in}{1.411163in}}%
\pgfusepath{clip}%
\pgfsetbuttcap%
\pgfsetroundjoin%
\definecolor{currentfill}{rgb}{0.749020,0.772549,0.792157}%
\pgfsetfillcolor{currentfill}%
\pgfsetfillopacity{0.160000}%
\pgfsetlinewidth{0.000000pt}%
\definecolor{currentstroke}{rgb}{0.000000,0.000000,0.000000}%
\pgfsetstrokecolor{currentstroke}%
\pgfsetdash{}{0pt}%
\pgfpathmoveto{\pgfqpoint{1.660348in}{1.231668in}}%
\pgfpathlineto{\pgfqpoint{1.672759in}{1.245190in}}%
\pgfpathlineto{\pgfqpoint{1.685171in}{1.258712in}}%
\pgfpathlineto{\pgfqpoint{1.667110in}{1.272405in}}%
\pgfpathlineto{\pgfqpoint{1.646783in}{1.285542in}}%
\pgfpathlineto{\pgfqpoint{1.635277in}{1.271586in}}%
\pgfpathlineto{\pgfqpoint{1.623766in}{1.257625in}}%
\pgfpathlineto{\pgfqpoint{1.643149in}{1.244913in}}%
\pgfpathlineto{\pgfqpoint{1.660348in}{1.231668in}}%
\pgfpathclose%
\pgfusepath{fill}%
\end{pgfscope}%
\begin{pgfscope}%
\pgfpathrectangle{\pgfqpoint{0.407592in}{0.222904in}}{\pgfqpoint{1.411163in}{1.411163in}}%
\pgfusepath{clip}%
\pgfsetbuttcap%
\pgfsetroundjoin%
\definecolor{currentfill}{rgb}{0.749020,0.772549,0.792157}%
\pgfsetfillcolor{currentfill}%
\pgfsetfillopacity{0.160000}%
\pgfsetlinewidth{0.000000pt}%
\definecolor{currentstroke}{rgb}{0.000000,0.000000,0.000000}%
\pgfsetstrokecolor{currentstroke}%
\pgfsetdash{}{0pt}%
\pgfpathmoveto{\pgfqpoint{0.967785in}{0.786836in}}%
\pgfpathlineto{\pgfqpoint{0.957624in}{0.801466in}}%
\pgfpathlineto{\pgfqpoint{0.947469in}{0.816087in}}%
\pgfpathlineto{\pgfqpoint{0.939943in}{0.811946in}}%
\pgfpathlineto{\pgfqpoint{0.933047in}{0.807614in}}%
\pgfpathlineto{\pgfqpoint{0.944253in}{0.793563in}}%
\pgfpathlineto{\pgfqpoint{0.955463in}{0.779508in}}%
\pgfpathlineto{\pgfqpoint{0.961352in}{0.783254in}}%
\pgfpathlineto{\pgfqpoint{0.967785in}{0.786836in}}%
\pgfpathclose%
\pgfusepath{fill}%
\end{pgfscope}%
\begin{pgfscope}%
\pgfpathrectangle{\pgfqpoint{0.407592in}{0.222904in}}{\pgfqpoint{1.411163in}{1.411163in}}%
\pgfusepath{clip}%
\pgfsetbuttcap%
\pgfsetroundjoin%
\definecolor{currentfill}{rgb}{0.749020,0.772549,0.792157}%
\pgfsetfillcolor{currentfill}%
\pgfsetfillopacity{0.160000}%
\pgfsetlinewidth{0.000000pt}%
\definecolor{currentstroke}{rgb}{0.000000,0.000000,0.000000}%
\pgfsetstrokecolor{currentstroke}%
\pgfsetdash{}{0pt}%
\pgfpathmoveto{\pgfqpoint{1.635528in}{1.204627in}}%
\pgfpathlineto{\pgfqpoint{1.647938in}{1.218147in}}%
\pgfpathlineto{\pgfqpoint{1.660348in}{1.231668in}}%
\pgfpathlineto{\pgfqpoint{1.643149in}{1.244913in}}%
\pgfpathlineto{\pgfqpoint{1.623766in}{1.257625in}}%
\pgfpathlineto{\pgfqpoint{1.612252in}{1.243660in}}%
\pgfpathlineto{\pgfqpoint{1.600734in}{1.229691in}}%
\pgfpathlineto{\pgfqpoint{1.619181in}{1.217414in}}%
\pgfpathlineto{\pgfqpoint{1.635528in}{1.204627in}}%
\pgfpathclose%
\pgfusepath{fill}%
\end{pgfscope}%
\begin{pgfscope}%
\pgfpathrectangle{\pgfqpoint{0.407592in}{0.222904in}}{\pgfqpoint{1.411163in}{1.411163in}}%
\pgfusepath{clip}%
\pgfsetbuttcap%
\pgfsetroundjoin%
\definecolor{currentfill}{rgb}{0.749020,0.772549,0.792157}%
\pgfsetfillcolor{currentfill}%
\pgfsetfillopacity{0.160000}%
\pgfsetlinewidth{0.000000pt}%
\definecolor{currentstroke}{rgb}{0.000000,0.000000,0.000000}%
\pgfsetstrokecolor{currentstroke}%
\pgfsetdash{}{0pt}%
\pgfpathmoveto{\pgfqpoint{1.154833in}{0.699214in}}%
\pgfpathlineto{\pgfqpoint{1.165575in}{0.713959in}}%
\pgfpathlineto{\pgfqpoint{1.176309in}{0.728694in}}%
\pgfpathlineto{\pgfqpoint{1.172106in}{0.730980in}}%
\pgfpathlineto{\pgfqpoint{1.167553in}{0.733136in}}%
\pgfpathlineto{\pgfqpoint{1.157921in}{0.717864in}}%
\pgfpathlineto{\pgfqpoint{1.148280in}{0.702577in}}%
\pgfpathlineto{\pgfqpoint{1.151689in}{0.700945in}}%
\pgfpathlineto{\pgfqpoint{1.154833in}{0.699214in}}%
\pgfpathclose%
\pgfusepath{fill}%
\end{pgfscope}%
\begin{pgfscope}%
\pgfpathrectangle{\pgfqpoint{0.407592in}{0.222904in}}{\pgfqpoint{1.411163in}{1.411163in}}%
\pgfusepath{clip}%
\pgfsetbuttcap%
\pgfsetroundjoin%
\definecolor{currentfill}{rgb}{0.749020,0.772549,0.792157}%
\pgfsetfillcolor{currentfill}%
\pgfsetfillopacity{0.160000}%
\pgfsetlinewidth{0.000000pt}%
\definecolor{currentstroke}{rgb}{0.000000,0.000000,0.000000}%
\pgfsetstrokecolor{currentstroke}%
\pgfsetdash{}{0pt}%
\pgfpathmoveto{\pgfqpoint{1.094627in}{0.644555in}}%
\pgfpathlineto{\pgfqpoint{1.096820in}{0.661517in}}%
\pgfpathlineto{\pgfqpoint{1.099008in}{0.678446in}}%
\pgfpathlineto{\pgfqpoint{1.095571in}{0.678584in}}%
\pgfpathlineto{\pgfqpoint{1.092110in}{0.678624in}}%
\pgfpathlineto{\pgfqpoint{1.091641in}{0.661651in}}%
\pgfpathlineto{\pgfqpoint{1.091172in}{0.644645in}}%
\pgfpathlineto{\pgfqpoint{1.092905in}{0.644625in}}%
\pgfpathlineto{\pgfqpoint{1.094627in}{0.644555in}}%
\pgfpathclose%
\pgfusepath{fill}%
\end{pgfscope}%
\begin{pgfscope}%
\pgfpathrectangle{\pgfqpoint{0.407592in}{0.222904in}}{\pgfqpoint{1.411163in}{1.411163in}}%
\pgfusepath{clip}%
\pgfsetbuttcap%
\pgfsetroundjoin%
\definecolor{currentfill}{rgb}{0.749020,0.772549,0.792157}%
\pgfsetfillcolor{currentfill}%
\pgfsetfillopacity{0.160000}%
\pgfsetlinewidth{0.000000pt}%
\definecolor{currentstroke}{rgb}{0.000000,0.000000,0.000000}%
\pgfsetstrokecolor{currentstroke}%
\pgfsetdash{}{0pt}%
\pgfpathmoveto{\pgfqpoint{1.091172in}{0.644645in}}%
\pgfpathlineto{\pgfqpoint{1.091641in}{0.661651in}}%
\pgfpathlineto{\pgfqpoint{1.092110in}{0.678624in}}%
\pgfpathlineto{\pgfqpoint{1.088641in}{0.678564in}}%
\pgfpathlineto{\pgfqpoint{1.085177in}{0.678406in}}%
\pgfpathlineto{\pgfqpoint{1.086437in}{0.661487in}}%
\pgfpathlineto{\pgfqpoint{1.087699in}{0.644535in}}%
\pgfpathlineto{\pgfqpoint{1.089434in}{0.644615in}}%
\pgfpathlineto{\pgfqpoint{1.091172in}{0.644645in}}%
\pgfpathclose%
\pgfusepath{fill}%
\end{pgfscope}%
\begin{pgfscope}%
\pgfpathrectangle{\pgfqpoint{0.407592in}{0.222904in}}{\pgfqpoint{1.411163in}{1.411163in}}%
\pgfusepath{clip}%
\pgfsetbuttcap%
\pgfsetroundjoin%
\definecolor{currentfill}{rgb}{0.749020,0.772549,0.792157}%
\pgfsetfillcolor{currentfill}%
\pgfsetfillopacity{0.160000}%
\pgfsetlinewidth{0.000000pt}%
\definecolor{currentstroke}{rgb}{0.000000,0.000000,0.000000}%
\pgfsetstrokecolor{currentstroke}%
\pgfsetdash{}{0pt}%
\pgfpathmoveto{\pgfqpoint{1.128969in}{0.671957in}}%
\pgfpathlineto{\pgfqpoint{1.138629in}{0.687275in}}%
\pgfpathlineto{\pgfqpoint{1.148280in}{0.702577in}}%
\pgfpathlineto{\pgfqpoint{1.144621in}{0.704104in}}%
\pgfpathlineto{\pgfqpoint{1.140729in}{0.705519in}}%
\pgfpathlineto{\pgfqpoint{1.132340in}{0.689739in}}%
\pgfpathlineto{\pgfqpoint{1.123939in}{0.673938in}}%
\pgfpathlineto{\pgfqpoint{1.126532in}{0.672985in}}%
\pgfpathlineto{\pgfqpoint{1.128969in}{0.671957in}}%
\pgfpathclose%
\pgfusepath{fill}%
\end{pgfscope}%
\begin{pgfscope}%
\pgfpathrectangle{\pgfqpoint{0.407592in}{0.222904in}}{\pgfqpoint{1.411163in}{1.411163in}}%
\pgfusepath{clip}%
\pgfsetbuttcap%
\pgfsetroundjoin%
\definecolor{currentfill}{rgb}{0.749020,0.772549,0.792157}%
\pgfsetfillcolor{currentfill}%
\pgfsetfillopacity{0.160000}%
\pgfsetlinewidth{0.000000pt}%
\definecolor{currentstroke}{rgb}{0.000000,0.000000,0.000000}%
\pgfsetstrokecolor{currentstroke}%
\pgfsetdash{}{0pt}%
\pgfpathmoveto{\pgfqpoint{1.610711in}{1.177589in}}%
\pgfpathlineto{\pgfqpoint{1.623119in}{1.191107in}}%
\pgfpathlineto{\pgfqpoint{1.635528in}{1.204627in}}%
\pgfpathlineto{\pgfqpoint{1.619181in}{1.217414in}}%
\pgfpathlineto{\pgfqpoint{1.600734in}{1.229691in}}%
\pgfpathlineto{\pgfqpoint{1.589212in}{1.215717in}}%
\pgfpathlineto{\pgfqpoint{1.577687in}{1.201739in}}%
\pgfpathlineto{\pgfqpoint{1.595206in}{1.189907in}}%
\pgfpathlineto{\pgfqpoint{1.610711in}{1.177589in}}%
\pgfpathclose%
\pgfusepath{fill}%
\end{pgfscope}%
\begin{pgfscope}%
\pgfpathrectangle{\pgfqpoint{0.407592in}{0.222904in}}{\pgfqpoint{1.411163in}{1.411163in}}%
\pgfusepath{clip}%
\pgfsetbuttcap%
\pgfsetroundjoin%
\definecolor{currentfill}{rgb}{0.749020,0.772549,0.792157}%
\pgfsetfillcolor{currentfill}%
\pgfsetfillopacity{0.160000}%
\pgfsetlinewidth{0.000000pt}%
\definecolor{currentstroke}{rgb}{0.000000,0.000000,0.000000}%
\pgfsetstrokecolor{currentstroke}%
\pgfsetdash{}{0pt}%
\pgfpathmoveto{\pgfqpoint{1.036038in}{0.702258in}}%
\pgfpathlineto{\pgfqpoint{1.027036in}{0.717494in}}%
\pgfpathlineto{\pgfqpoint{1.018044in}{0.732714in}}%
\pgfpathlineto{\pgfqpoint{1.013105in}{0.730532in}}%
\pgfpathlineto{\pgfqpoint{1.008494in}{0.728220in}}%
\pgfpathlineto{\pgfqpoint{1.018688in}{0.713543in}}%
\pgfpathlineto{\pgfqpoint{1.028888in}{0.698855in}}%
\pgfpathlineto{\pgfqpoint{1.032339in}{0.700605in}}%
\pgfpathlineto{\pgfqpoint{1.036038in}{0.702258in}}%
\pgfpathclose%
\pgfusepath{fill}%
\end{pgfscope}%
\begin{pgfscope}%
\pgfpathrectangle{\pgfqpoint{0.407592in}{0.222904in}}{\pgfqpoint{1.411163in}{1.411163in}}%
\pgfusepath{clip}%
\pgfsetbuttcap%
\pgfsetroundjoin%
\definecolor{currentfill}{rgb}{0.749020,0.772549,0.792157}%
\pgfsetfillcolor{currentfill}%
\pgfsetfillopacity{0.160000}%
\pgfsetlinewidth{0.000000pt}%
\definecolor{currentstroke}{rgb}{0.000000,0.000000,0.000000}%
\pgfsetstrokecolor{currentstroke}%
\pgfsetdash{}{0pt}%
\pgfpathmoveto{\pgfqpoint{1.098006in}{0.644266in}}%
\pgfpathlineto{\pgfqpoint{1.101883in}{0.661086in}}%
\pgfpathlineto{\pgfqpoint{1.105754in}{0.677875in}}%
\pgfpathlineto{\pgfqpoint{1.102407in}{0.678209in}}%
\pgfpathlineto{\pgfqpoint{1.099008in}{0.678446in}}%
\pgfpathlineto{\pgfqpoint{1.096820in}{0.661517in}}%
\pgfpathlineto{\pgfqpoint{1.094627in}{0.644555in}}%
\pgfpathlineto{\pgfqpoint{1.096330in}{0.644435in}}%
\pgfpathlineto{\pgfqpoint{1.098006in}{0.644266in}}%
\pgfpathclose%
\pgfusepath{fill}%
\end{pgfscope}%
\begin{pgfscope}%
\pgfpathrectangle{\pgfqpoint{0.407592in}{0.222904in}}{\pgfqpoint{1.411163in}{1.411163in}}%
\pgfusepath{clip}%
\pgfsetbuttcap%
\pgfsetroundjoin%
\definecolor{currentfill}{rgb}{0.749020,0.772549,0.792157}%
\pgfsetfillcolor{currentfill}%
\pgfsetfillopacity{0.160000}%
\pgfsetlinewidth{0.000000pt}%
\definecolor{currentstroke}{rgb}{0.000000,0.000000,0.000000}%
\pgfsetstrokecolor{currentstroke}%
\pgfsetdash{}{0pt}%
\pgfpathmoveto{\pgfqpoint{1.087699in}{0.644535in}}%
\pgfpathlineto{\pgfqpoint{1.086437in}{0.661487in}}%
\pgfpathlineto{\pgfqpoint{1.085177in}{0.678406in}}%
\pgfpathlineto{\pgfqpoint{1.081735in}{0.678150in}}%
\pgfpathlineto{\pgfqpoint{1.078329in}{0.677797in}}%
\pgfpathlineto{\pgfqpoint{1.081296in}{0.661027in}}%
\pgfpathlineto{\pgfqpoint{1.084268in}{0.644227in}}%
\pgfpathlineto{\pgfqpoint{1.085974in}{0.644405in}}%
\pgfpathlineto{\pgfqpoint{1.087699in}{0.644535in}}%
\pgfpathclose%
\pgfusepath{fill}%
\end{pgfscope}%
\begin{pgfscope}%
\pgfpathrectangle{\pgfqpoint{0.407592in}{0.222904in}}{\pgfqpoint{1.411163in}{1.411163in}}%
\pgfusepath{clip}%
\pgfsetbuttcap%
\pgfsetroundjoin%
\definecolor{currentfill}{rgb}{0.749020,0.772549,0.792157}%
\pgfsetfillcolor{currentfill}%
\pgfsetfillopacity{0.160000}%
\pgfsetlinewidth{0.000000pt}%
\definecolor{currentstroke}{rgb}{0.000000,0.000000,0.000000}%
\pgfsetstrokecolor{currentstroke}%
\pgfsetdash{}{0pt}%
\pgfpathmoveto{\pgfqpoint{1.059445in}{0.673753in}}%
\pgfpathlineto{\pgfqpoint{1.051773in}{0.689509in}}%
\pgfpathlineto{\pgfqpoint{1.044111in}{0.705244in}}%
\pgfpathlineto{\pgfqpoint{1.039967in}{0.703806in}}%
\pgfpathlineto{\pgfqpoint{1.036038in}{0.702258in}}%
\pgfpathlineto{\pgfqpoint{1.045048in}{0.687008in}}%
\pgfpathlineto{\pgfqpoint{1.054067in}{0.671742in}}%
\pgfpathlineto{\pgfqpoint{1.056684in}{0.672785in}}%
\pgfpathlineto{\pgfqpoint{1.059445in}{0.673753in}}%
\pgfpathclose%
\pgfusepath{fill}%
\end{pgfscope}%
\begin{pgfscope}%
\pgfpathrectangle{\pgfqpoint{0.407592in}{0.222904in}}{\pgfqpoint{1.411163in}{1.411163in}}%
\pgfusepath{clip}%
\pgfsetbuttcap%
\pgfsetroundjoin%
\definecolor{currentfill}{rgb}{0.749020,0.772549,0.792157}%
\pgfsetfillcolor{currentfill}%
\pgfsetfillopacity{0.160000}%
\pgfsetlinewidth{0.000000pt}%
\definecolor{currentstroke}{rgb}{0.000000,0.000000,0.000000}%
\pgfsetstrokecolor{currentstroke}%
\pgfsetdash{}{0pt}%
\pgfpathmoveto{\pgfqpoint{1.206894in}{0.752028in}}%
\pgfpathlineto{\pgfqpoint{1.218540in}{0.766152in}}%
\pgfpathlineto{\pgfqpoint{1.230182in}{0.780272in}}%
\pgfpathlineto{\pgfqpoint{1.224964in}{0.783986in}}%
\pgfpathlineto{\pgfqpoint{1.219174in}{0.787533in}}%
\pgfpathlineto{\pgfqpoint{1.208469in}{0.772839in}}%
\pgfpathlineto{\pgfqpoint{1.197756in}{0.758134in}}%
\pgfpathlineto{\pgfqpoint{1.202565in}{0.755151in}}%
\pgfpathlineto{\pgfqpoint{1.206894in}{0.752028in}}%
\pgfpathclose%
\pgfusepath{fill}%
\end{pgfscope}%
\begin{pgfscope}%
\pgfpathrectangle{\pgfqpoint{0.407592in}{0.222904in}}{\pgfqpoint{1.411163in}{1.411163in}}%
\pgfusepath{clip}%
\pgfsetbuttcap%
\pgfsetroundjoin%
\definecolor{currentfill}{rgb}{0.749020,0.772549,0.792157}%
\pgfsetfillcolor{currentfill}%
\pgfsetfillopacity{0.160000}%
\pgfsetlinewidth{0.000000pt}%
\definecolor{currentstroke}{rgb}{0.000000,0.000000,0.000000}%
\pgfsetstrokecolor{currentstroke}%
\pgfsetdash{}{0pt}%
\pgfpathmoveto{\pgfqpoint{1.585897in}{1.150554in}}%
\pgfpathlineto{\pgfqpoint{1.598303in}{1.164071in}}%
\pgfpathlineto{\pgfqpoint{1.610711in}{1.177589in}}%
\pgfpathlineto{\pgfqpoint{1.595206in}{1.189907in}}%
\pgfpathlineto{\pgfqpoint{1.577687in}{1.201739in}}%
\pgfpathlineto{\pgfqpoint{1.566157in}{1.187756in}}%
\pgfpathlineto{\pgfqpoint{1.554624in}{1.173768in}}%
\pgfpathlineto{\pgfqpoint{1.571225in}{1.162392in}}%
\pgfpathlineto{\pgfqpoint{1.585897in}{1.150554in}}%
\pgfpathclose%
\pgfusepath{fill}%
\end{pgfscope}%
\begin{pgfscope}%
\pgfpathrectangle{\pgfqpoint{0.407592in}{0.222904in}}{\pgfqpoint{1.411163in}{1.411163in}}%
\pgfusepath{clip}%
\pgfsetbuttcap%
\pgfsetroundjoin%
\definecolor{currentfill}{rgb}{0.749020,0.772549,0.792157}%
\pgfsetfillcolor{currentfill}%
\pgfsetfillopacity{0.160000}%
\pgfsetlinewidth{0.000000pt}%
\definecolor{currentstroke}{rgb}{0.000000,0.000000,0.000000}%
\pgfsetstrokecolor{currentstroke}%
\pgfsetdash{}{0pt}%
\pgfpathmoveto{\pgfqpoint{0.465273in}{1.394126in}}%
\pgfpathlineto{\pgfqpoint{0.454204in}{1.408005in}}%
\pgfpathlineto{\pgfqpoint{0.426275in}{1.392953in}}%
\pgfpathlineto{\pgfqpoint{0.400835in}{1.377253in}}%
\pgfpathlineto{\pgfqpoint{0.412963in}{1.363764in}}%
\pgfpathlineto{\pgfqpoint{0.437892in}{1.379263in}}%
\pgfusepath{fill}%
\end{pgfscope}%
\begin{pgfscope}%
\pgfpathrectangle{\pgfqpoint{0.407592in}{0.222904in}}{\pgfqpoint{1.411163in}{1.411163in}}%
\pgfusepath{clip}%
\pgfsetbuttcap%
\pgfsetroundjoin%
\definecolor{currentfill}{rgb}{0.749020,0.772549,0.792157}%
\pgfsetfillcolor{currentfill}%
\pgfsetfillopacity{0.160000}%
\pgfsetlinewidth{0.000000pt}%
\definecolor{currentstroke}{rgb}{0.000000,0.000000,0.000000}%
\pgfsetstrokecolor{currentstroke}%
\pgfsetdash{}{0pt}%
\pgfpathmoveto{\pgfqpoint{1.101249in}{0.643784in}}%
\pgfpathlineto{\pgfqpoint{1.106744in}{0.660368in}}%
\pgfpathlineto{\pgfqpoint{1.112231in}{0.676922in}}%
\pgfpathlineto{\pgfqpoint{1.109033in}{0.677446in}}%
\pgfpathlineto{\pgfqpoint{1.105754in}{0.677875in}}%
\pgfpathlineto{\pgfqpoint{1.101883in}{0.661086in}}%
\pgfpathlineto{\pgfqpoint{1.098006in}{0.644266in}}%
\pgfpathlineto{\pgfqpoint{1.099648in}{0.644049in}}%
\pgfpathlineto{\pgfqpoint{1.101249in}{0.643784in}}%
\pgfpathclose%
\pgfusepath{fill}%
\end{pgfscope}%
\begin{pgfscope}%
\pgfpathrectangle{\pgfqpoint{0.407592in}{0.222904in}}{\pgfqpoint{1.411163in}{1.411163in}}%
\pgfusepath{clip}%
\pgfsetbuttcap%
\pgfsetroundjoin%
\definecolor{currentfill}{rgb}{0.749020,0.772549,0.792157}%
\pgfsetfillcolor{currentfill}%
\pgfsetfillopacity{0.160000}%
\pgfsetlinewidth{0.000000pt}%
\definecolor{currentstroke}{rgb}{0.000000,0.000000,0.000000}%
\pgfsetstrokecolor{currentstroke}%
\pgfsetdash{}{0pt}%
\pgfpathmoveto{\pgfqpoint{1.084268in}{0.644227in}}%
\pgfpathlineto{\pgfqpoint{1.081296in}{0.661027in}}%
\pgfpathlineto{\pgfqpoint{1.078329in}{0.677797in}}%
\pgfpathlineto{\pgfqpoint{1.074973in}{0.677348in}}%
\pgfpathlineto{\pgfqpoint{1.071682in}{0.676806in}}%
\pgfpathlineto{\pgfqpoint{1.076307in}{0.660280in}}%
\pgfpathlineto{\pgfqpoint{1.080940in}{0.643726in}}%
\pgfpathlineto{\pgfqpoint{1.082588in}{0.644000in}}%
\pgfpathlineto{\pgfqpoint{1.084268in}{0.644227in}}%
\pgfpathclose%
\pgfusepath{fill}%
\end{pgfscope}%
\begin{pgfscope}%
\pgfpathrectangle{\pgfqpoint{0.407592in}{0.222904in}}{\pgfqpoint{1.411163in}{1.411163in}}%
\pgfusepath{clip}%
\pgfsetbuttcap%
\pgfsetroundjoin%
\definecolor{currentfill}{rgb}{0.749020,0.772549,0.792157}%
\pgfsetfillcolor{currentfill}%
\pgfsetfillopacity{0.160000}%
\pgfsetlinewidth{0.000000pt}%
\definecolor{currentstroke}{rgb}{0.000000,0.000000,0.000000}%
\pgfsetstrokecolor{currentstroke}%
\pgfsetdash{}{0pt}%
\pgfpathmoveto{\pgfqpoint{1.561085in}{1.123522in}}%
\pgfpathlineto{\pgfqpoint{1.573491in}{1.137038in}}%
\pgfpathlineto{\pgfqpoint{1.585897in}{1.150554in}}%
\pgfpathlineto{\pgfqpoint{1.571225in}{1.162392in}}%
\pgfpathlineto{\pgfqpoint{1.554624in}{1.173768in}}%
\pgfpathlineto{\pgfqpoint{1.543088in}{1.159776in}}%
\pgfpathlineto{\pgfqpoint{1.531547in}{1.145779in}}%
\pgfpathlineto{\pgfqpoint{1.547237in}{1.134871in}}%
\pgfpathlineto{\pgfqpoint{1.561085in}{1.123522in}}%
\pgfpathclose%
\pgfusepath{fill}%
\end{pgfscope}%
\begin{pgfscope}%
\pgfpathrectangle{\pgfqpoint{0.407592in}{0.222904in}}{\pgfqpoint{1.411163in}{1.411163in}}%
\pgfusepath{clip}%
\pgfsetbuttcap%
\pgfsetroundjoin%
\definecolor{currentfill}{rgb}{0.749020,0.772549,0.792157}%
\pgfsetfillcolor{currentfill}%
\pgfsetfillopacity{0.160000}%
\pgfsetlinewidth{0.000000pt}%
\definecolor{currentstroke}{rgb}{0.000000,0.000000,0.000000}%
\pgfsetstrokecolor{currentstroke}%
\pgfsetdash{}{0pt}%
\pgfpathmoveto{\pgfqpoint{0.487421in}{1.366357in}}%
\pgfpathlineto{\pgfqpoint{0.476345in}{1.380243in}}%
\pgfpathlineto{\pgfqpoint{0.465273in}{1.394126in}}%
\pgfpathlineto{\pgfqpoint{0.437892in}{1.379263in}}%
\pgfpathlineto{\pgfqpoint{0.412963in}{1.363764in}}%
\pgfpathlineto{\pgfqpoint{0.425090in}{1.350276in}}%
\pgfpathlineto{\pgfqpoint{0.437216in}{1.336790in}}%
\pgfpathlineto{\pgfqpoint{0.461128in}{1.351881in}}%
\pgfpathlineto{\pgfqpoint{0.487421in}{1.366357in}}%
\pgfpathclose%
\pgfusepath{fill}%
\end{pgfscope}%
\begin{pgfscope}%
\pgfpathrectangle{\pgfqpoint{0.407592in}{0.222904in}}{\pgfqpoint{1.411163in}{1.411163in}}%
\pgfusepath{clip}%
\pgfsetbuttcap%
\pgfsetroundjoin%
\definecolor{currentfill}{rgb}{0.749020,0.772549,0.792157}%
\pgfsetfillcolor{currentfill}%
\pgfsetfillopacity{0.160000}%
\pgfsetlinewidth{0.000000pt}%
\definecolor{currentstroke}{rgb}{0.000000,0.000000,0.000000}%
\pgfsetstrokecolor{currentstroke}%
\pgfsetdash{}{0pt}%
\pgfpathmoveto{\pgfqpoint{0.988126in}{0.757547in}}%
\pgfpathlineto{\pgfqpoint{0.977952in}{0.772196in}}%
\pgfpathlineto{\pgfqpoint{0.967785in}{0.786836in}}%
\pgfpathlineto{\pgfqpoint{0.961352in}{0.783254in}}%
\pgfpathlineto{\pgfqpoint{0.955463in}{0.779508in}}%
\pgfpathlineto{\pgfqpoint{0.966676in}{0.765449in}}%
\pgfpathlineto{\pgfqpoint{0.977892in}{0.751386in}}%
\pgfpathlineto{\pgfqpoint{0.982781in}{0.754535in}}%
\pgfpathlineto{\pgfqpoint{0.988126in}{0.757547in}}%
\pgfpathclose%
\pgfusepath{fill}%
\end{pgfscope}%
\begin{pgfscope}%
\pgfpathrectangle{\pgfqpoint{0.407592in}{0.222904in}}{\pgfqpoint{1.411163in}{1.411163in}}%
\pgfusepath{clip}%
\pgfsetbuttcap%
\pgfsetroundjoin%
\definecolor{currentfill}{rgb}{0.749020,0.772549,0.792157}%
\pgfsetfillcolor{currentfill}%
\pgfsetfillopacity{0.160000}%
\pgfsetlinewidth{0.000000pt}%
\definecolor{currentstroke}{rgb}{0.000000,0.000000,0.000000}%
\pgfsetstrokecolor{currentstroke}%
\pgfsetdash{}{0pt}%
\pgfpathmoveto{\pgfqpoint{0.509581in}{1.338571in}}%
\pgfpathlineto{\pgfqpoint{0.498499in}{1.352466in}}%
\pgfpathlineto{\pgfqpoint{0.487421in}{1.366357in}}%
\pgfpathlineto{\pgfqpoint{0.461128in}{1.351881in}}%
\pgfpathlineto{\pgfqpoint{0.437216in}{1.336790in}}%
\pgfpathlineto{\pgfqpoint{0.449341in}{1.323306in}}%
\pgfpathlineto{\pgfqpoint{0.461465in}{1.309822in}}%
\pgfpathlineto{\pgfqpoint{0.484369in}{1.324493in}}%
\pgfpathlineto{\pgfqpoint{0.509581in}{1.338571in}}%
\pgfpathclose%
\pgfusepath{fill}%
\end{pgfscope}%
\begin{pgfscope}%
\pgfpathrectangle{\pgfqpoint{0.407592in}{0.222904in}}{\pgfqpoint{1.411163in}{1.411163in}}%
\pgfusepath{clip}%
\pgfsetbuttcap%
\pgfsetroundjoin%
\definecolor{currentfill}{rgb}{0.749020,0.772549,0.792157}%
\pgfsetfillcolor{currentfill}%
\pgfsetfillopacity{0.160000}%
\pgfsetlinewidth{0.000000pt}%
\definecolor{currentstroke}{rgb}{0.000000,0.000000,0.000000}%
\pgfsetstrokecolor{currentstroke}%
\pgfsetdash{}{0pt}%
\pgfpathmoveto{\pgfqpoint{1.536277in}{1.096494in}}%
\pgfpathlineto{\pgfqpoint{1.548681in}{1.110008in}}%
\pgfpathlineto{\pgfqpoint{1.561085in}{1.123522in}}%
\pgfpathlineto{\pgfqpoint{1.547237in}{1.134871in}}%
\pgfpathlineto{\pgfqpoint{1.531547in}{1.145779in}}%
\pgfpathlineto{\pgfqpoint{1.520003in}{1.131778in}}%
\pgfpathlineto{\pgfqpoint{1.508456in}{1.117773in}}%
\pgfpathlineto{\pgfqpoint{1.523243in}{1.107341in}}%
\pgfpathlineto{\pgfqpoint{1.536277in}{1.096494in}}%
\pgfpathclose%
\pgfusepath{fill}%
\end{pgfscope}%
\begin{pgfscope}%
\pgfpathrectangle{\pgfqpoint{0.407592in}{0.222904in}}{\pgfqpoint{1.411163in}{1.411163in}}%
\pgfusepath{clip}%
\pgfsetbuttcap%
\pgfsetroundjoin%
\definecolor{currentfill}{rgb}{0.749020,0.772549,0.792157}%
\pgfsetfillcolor{currentfill}%
\pgfsetfillopacity{0.160000}%
\pgfsetlinewidth{0.000000pt}%
\definecolor{currentstroke}{rgb}{0.000000,0.000000,0.000000}%
\pgfsetstrokecolor{currentstroke}%
\pgfsetdash{}{0pt}%
\pgfpathmoveto{\pgfqpoint{1.104300in}{0.643117in}}%
\pgfpathlineto{\pgfqpoint{1.111319in}{0.659372in}}%
\pgfpathlineto{\pgfqpoint{1.118327in}{0.675602in}}%
\pgfpathlineto{\pgfqpoint{1.115333in}{0.676307in}}%
\pgfpathlineto{\pgfqpoint{1.112231in}{0.676922in}}%
\pgfpathlineto{\pgfqpoint{1.106744in}{0.660368in}}%
\pgfpathlineto{\pgfqpoint{1.101249in}{0.643784in}}%
\pgfpathlineto{\pgfqpoint{1.102802in}{0.643473in}}%
\pgfpathlineto{\pgfqpoint{1.104300in}{0.643117in}}%
\pgfpathclose%
\pgfusepath{fill}%
\end{pgfscope}%
\begin{pgfscope}%
\pgfpathrectangle{\pgfqpoint{0.407592in}{0.222904in}}{\pgfqpoint{1.411163in}{1.411163in}}%
\pgfusepath{clip}%
\pgfsetbuttcap%
\pgfsetroundjoin%
\definecolor{currentfill}{rgb}{0.749020,0.772549,0.792157}%
\pgfsetfillcolor{currentfill}%
\pgfsetfillopacity{0.160000}%
\pgfsetlinewidth{0.000000pt}%
\definecolor{currentstroke}{rgb}{0.000000,0.000000,0.000000}%
\pgfsetstrokecolor{currentstroke}%
\pgfsetdash{}{0pt}%
\pgfpathmoveto{\pgfqpoint{0.531754in}{1.310770in}}%
\pgfpathlineto{\pgfqpoint{0.520666in}{1.324673in}}%
\pgfpathlineto{\pgfqpoint{0.509581in}{1.338571in}}%
\pgfpathlineto{\pgfqpoint{0.484369in}{1.324493in}}%
\pgfpathlineto{\pgfqpoint{0.461465in}{1.309822in}}%
\pgfpathlineto{\pgfqpoint{0.473587in}{1.296340in}}%
\pgfpathlineto{\pgfqpoint{0.485708in}{1.282860in}}%
\pgfpathlineto{\pgfqpoint{0.507615in}{1.297099in}}%
\pgfpathlineto{\pgfqpoint{0.531754in}{1.310770in}}%
\pgfpathclose%
\pgfusepath{fill}%
\end{pgfscope}%
\begin{pgfscope}%
\pgfpathrectangle{\pgfqpoint{0.407592in}{0.222904in}}{\pgfqpoint{1.411163in}{1.411163in}}%
\pgfusepath{clip}%
\pgfsetbuttcap%
\pgfsetroundjoin%
\definecolor{currentfill}{rgb}{0.749020,0.772549,0.792157}%
\pgfsetfillcolor{currentfill}%
\pgfsetfillopacity{0.160000}%
\pgfsetlinewidth{0.000000pt}%
\definecolor{currentstroke}{rgb}{0.000000,0.000000,0.000000}%
\pgfsetstrokecolor{currentstroke}%
\pgfsetdash{}{0pt}%
\pgfpathmoveto{\pgfqpoint{1.080940in}{0.643726in}}%
\pgfpathlineto{\pgfqpoint{1.076307in}{0.660280in}}%
\pgfpathlineto{\pgfqpoint{1.071682in}{0.676806in}}%
\pgfpathlineto{\pgfqpoint{1.068470in}{0.676173in}}%
\pgfpathlineto{\pgfqpoint{1.065351in}{0.675451in}}%
\pgfpathlineto{\pgfqpoint{1.071557in}{0.659258in}}%
\pgfpathlineto{\pgfqpoint{1.077772in}{0.643040in}}%
\pgfpathlineto{\pgfqpoint{1.079332in}{0.643405in}}%
\pgfpathlineto{\pgfqpoint{1.080940in}{0.643726in}}%
\pgfpathclose%
\pgfusepath{fill}%
\end{pgfscope}%
\begin{pgfscope}%
\pgfpathrectangle{\pgfqpoint{0.407592in}{0.222904in}}{\pgfqpoint{1.411163in}{1.411163in}}%
\pgfusepath{clip}%
\pgfsetbuttcap%
\pgfsetroundjoin%
\definecolor{currentfill}{rgb}{0.749020,0.772549,0.792157}%
\pgfsetfillcolor{currentfill}%
\pgfsetfillopacity{0.160000}%
\pgfsetlinewidth{0.000000pt}%
\definecolor{currentstroke}{rgb}{0.000000,0.000000,0.000000}%
\pgfsetstrokecolor{currentstroke}%
\pgfsetdash{}{0pt}%
\pgfpathmoveto{\pgfqpoint{1.511471in}{1.069469in}}%
\pgfpathlineto{\pgfqpoint{1.523874in}{1.082981in}}%
\pgfpathlineto{\pgfqpoint{1.536277in}{1.096494in}}%
\pgfpathlineto{\pgfqpoint{1.523243in}{1.107341in}}%
\pgfpathlineto{\pgfqpoint{1.508456in}{1.117773in}}%
\pgfpathlineto{\pgfqpoint{1.496904in}{1.103763in}}%
\pgfpathlineto{\pgfqpoint{1.485349in}{1.089748in}}%
\pgfpathlineto{\pgfqpoint{1.499242in}{1.079805in}}%
\pgfpathlineto{\pgfqpoint{1.511471in}{1.069469in}}%
\pgfpathclose%
\pgfusepath{fill}%
\end{pgfscope}%
\begin{pgfscope}%
\pgfpathrectangle{\pgfqpoint{0.407592in}{0.222904in}}{\pgfqpoint{1.411163in}{1.411163in}}%
\pgfusepath{clip}%
\pgfsetbuttcap%
\pgfsetroundjoin%
\definecolor{currentfill}{rgb}{0.749020,0.772549,0.792157}%
\pgfsetfillcolor{currentfill}%
\pgfsetfillopacity{0.160000}%
\pgfsetlinewidth{0.000000pt}%
\definecolor{currentstroke}{rgb}{0.000000,0.000000,0.000000}%
\pgfsetstrokecolor{currentstroke}%
\pgfsetdash{}{0pt}%
\pgfpathmoveto{\pgfqpoint{0.553940in}{1.282953in}}%
\pgfpathlineto{\pgfqpoint{0.542845in}{1.296863in}}%
\pgfpathlineto{\pgfqpoint{0.531754in}{1.310770in}}%
\pgfpathlineto{\pgfqpoint{0.507615in}{1.297099in}}%
\pgfpathlineto{\pgfqpoint{0.485708in}{1.282860in}}%
\pgfpathlineto{\pgfqpoint{0.497828in}{1.269380in}}%
\pgfpathlineto{\pgfqpoint{0.509947in}{1.255902in}}%
\pgfpathlineto{\pgfqpoint{0.530865in}{1.269701in}}%
\pgfpathlineto{\pgfqpoint{0.553940in}{1.282953in}}%
\pgfpathclose%
\pgfusepath{fill}%
\end{pgfscope}%
\begin{pgfscope}%
\pgfpathrectangle{\pgfqpoint{0.407592in}{0.222904in}}{\pgfqpoint{1.411163in}{1.411163in}}%
\pgfusepath{clip}%
\pgfsetbuttcap%
\pgfsetroundjoin%
\definecolor{currentfill}{rgb}{0.749020,0.772549,0.792157}%
\pgfsetfillcolor{currentfill}%
\pgfsetfillopacity{0.160000}%
\pgfsetlinewidth{0.000000pt}%
\definecolor{currentstroke}{rgb}{0.000000,0.000000,0.000000}%
\pgfsetstrokecolor{currentstroke}%
\pgfsetdash{}{0pt}%
\pgfpathmoveto{\pgfqpoint{1.183592in}{0.723766in}}%
\pgfpathlineto{\pgfqpoint{1.195245in}{0.737899in}}%
\pgfpathlineto{\pgfqpoint{1.206894in}{0.752028in}}%
\pgfpathlineto{\pgfqpoint{1.202565in}{0.755151in}}%
\pgfpathlineto{\pgfqpoint{1.197756in}{0.758134in}}%
\pgfpathlineto{\pgfqpoint{1.187036in}{0.743419in}}%
\pgfpathlineto{\pgfqpoint{1.176309in}{0.728694in}}%
\pgfpathlineto{\pgfqpoint{1.180144in}{0.726286in}}%
\pgfpathlineto{\pgfqpoint{1.183592in}{0.723766in}}%
\pgfpathclose%
\pgfusepath{fill}%
\end{pgfscope}%
\begin{pgfscope}%
\pgfpathrectangle{\pgfqpoint{0.407592in}{0.222904in}}{\pgfqpoint{1.411163in}{1.411163in}}%
\pgfusepath{clip}%
\pgfsetbuttcap%
\pgfsetroundjoin%
\definecolor{currentfill}{rgb}{0.749020,0.772549,0.792157}%
\pgfsetfillcolor{currentfill}%
\pgfsetfillopacity{0.160000}%
\pgfsetlinewidth{0.000000pt}%
\definecolor{currentstroke}{rgb}{0.000000,0.000000,0.000000}%
\pgfsetstrokecolor{currentstroke}%
\pgfsetdash{}{0pt}%
\pgfpathmoveto{\pgfqpoint{0.576138in}{1.255120in}}%
\pgfpathlineto{\pgfqpoint{0.565037in}{1.269038in}}%
\pgfpathlineto{\pgfqpoint{0.553940in}{1.282953in}}%
\pgfpathlineto{\pgfqpoint{0.530865in}{1.269701in}}%
\pgfpathlineto{\pgfqpoint{0.509947in}{1.255902in}}%
\pgfpathlineto{\pgfqpoint{0.522064in}{1.242426in}}%
\pgfpathlineto{\pgfqpoint{0.534180in}{1.228950in}}%
\pgfpathlineto{\pgfqpoint{0.554120in}{1.242297in}}%
\pgfpathlineto{\pgfqpoint{0.576138in}{1.255120in}}%
\pgfpathclose%
\pgfusepath{fill}%
\end{pgfscope}%
\begin{pgfscope}%
\pgfpathrectangle{\pgfqpoint{0.407592in}{0.222904in}}{\pgfqpoint{1.411163in}{1.411163in}}%
\pgfusepath{clip}%
\pgfsetbuttcap%
\pgfsetroundjoin%
\definecolor{currentfill}{rgb}{0.749020,0.772549,0.792157}%
\pgfsetfillcolor{currentfill}%
\pgfsetfillopacity{0.160000}%
\pgfsetlinewidth{0.000000pt}%
\definecolor{currentstroke}{rgb}{0.000000,0.000000,0.000000}%
\pgfsetstrokecolor{currentstroke}%
\pgfsetdash{}{0pt}%
\pgfpathmoveto{\pgfqpoint{1.486669in}{1.042447in}}%
\pgfpathlineto{\pgfqpoint{1.499070in}{1.055958in}}%
\pgfpathlineto{\pgfqpoint{1.511471in}{1.069469in}}%
\pgfpathlineto{\pgfqpoint{1.499242in}{1.079805in}}%
\pgfpathlineto{\pgfqpoint{1.485349in}{1.089748in}}%
\pgfpathlineto{\pgfqpoint{1.473790in}{1.075729in}}%
\pgfpathlineto{\pgfqpoint{1.462227in}{1.061705in}}%
\pgfpathlineto{\pgfqpoint{1.475235in}{1.052261in}}%
\pgfpathlineto{\pgfqpoint{1.486669in}{1.042447in}}%
\pgfpathclose%
\pgfusepath{fill}%
\end{pgfscope}%
\begin{pgfscope}%
\pgfpathrectangle{\pgfqpoint{0.407592in}{0.222904in}}{\pgfqpoint{1.411163in}{1.411163in}}%
\pgfusepath{clip}%
\pgfsetbuttcap%
\pgfsetroundjoin%
\definecolor{currentfill}{rgb}{0.749020,0.772549,0.792157}%
\pgfsetfillcolor{currentfill}%
\pgfsetfillopacity{0.160000}%
\pgfsetlinewidth{0.000000pt}%
\definecolor{currentstroke}{rgb}{0.000000,0.000000,0.000000}%
\pgfsetstrokecolor{currentstroke}%
\pgfsetdash{}{0pt}%
\pgfpathmoveto{\pgfqpoint{0.598349in}{1.227271in}}%
\pgfpathlineto{\pgfqpoint{0.587242in}{1.241197in}}%
\pgfpathlineto{\pgfqpoint{0.576138in}{1.255120in}}%
\pgfpathlineto{\pgfqpoint{0.554120in}{1.242297in}}%
\pgfpathlineto{\pgfqpoint{0.534180in}{1.228950in}}%
\pgfpathlineto{\pgfqpoint{0.546296in}{1.215477in}}%
\pgfpathlineto{\pgfqpoint{0.558409in}{1.202004in}}%
\pgfpathlineto{\pgfqpoint{0.577379in}{1.214888in}}%
\pgfpathlineto{\pgfqpoint{0.598349in}{1.227271in}}%
\pgfpathclose%
\pgfusepath{fill}%
\end{pgfscope}%
\begin{pgfscope}%
\pgfpathrectangle{\pgfqpoint{0.407592in}{0.222904in}}{\pgfqpoint{1.411163in}{1.411163in}}%
\pgfusepath{clip}%
\pgfsetbuttcap%
\pgfsetroundjoin%
\definecolor{currentfill}{rgb}{0.749020,0.772549,0.792157}%
\pgfsetfillcolor{currentfill}%
\pgfsetfillopacity{0.160000}%
\pgfsetlinewidth{0.000000pt}%
\definecolor{currentstroke}{rgb}{0.000000,0.000000,0.000000}%
\pgfsetstrokecolor{currentstroke}%
\pgfsetdash{}{0pt}%
\pgfpathmoveto{\pgfqpoint{1.133328in}{0.669694in}}%
\pgfpathlineto{\pgfqpoint{1.144084in}{0.684459in}}%
\pgfpathlineto{\pgfqpoint{1.154833in}{0.699214in}}%
\pgfpathlineto{\pgfqpoint{1.151689in}{0.700945in}}%
\pgfpathlineto{\pgfqpoint{1.148280in}{0.702577in}}%
\pgfpathlineto{\pgfqpoint{1.138629in}{0.687275in}}%
\pgfpathlineto{\pgfqpoint{1.128969in}{0.671957in}}%
\pgfpathlineto{\pgfqpoint{1.131237in}{0.670859in}}%
\pgfpathlineto{\pgfqpoint{1.133328in}{0.669694in}}%
\pgfpathclose%
\pgfusepath{fill}%
\end{pgfscope}%
\begin{pgfscope}%
\pgfpathrectangle{\pgfqpoint{0.407592in}{0.222904in}}{\pgfqpoint{1.411163in}{1.411163in}}%
\pgfusepath{clip}%
\pgfsetbuttcap%
\pgfsetroundjoin%
\definecolor{currentfill}{rgb}{0.749020,0.772549,0.792157}%
\pgfsetfillcolor{currentfill}%
\pgfsetfillopacity{0.160000}%
\pgfsetlinewidth{0.000000pt}%
\definecolor{currentstroke}{rgb}{0.000000,0.000000,0.000000}%
\pgfsetstrokecolor{currentstroke}%
\pgfsetdash{}{0pt}%
\pgfpathmoveto{\pgfqpoint{1.107106in}{0.642276in}}%
\pgfpathlineto{\pgfqpoint{1.115528in}{0.658117in}}%
\pgfpathlineto{\pgfqpoint{1.123939in}{0.673938in}}%
\pgfpathlineto{\pgfqpoint{1.121200in}{0.674812in}}%
\pgfpathlineto{\pgfqpoint{1.118327in}{0.675602in}}%
\pgfpathlineto{\pgfqpoint{1.111319in}{0.659372in}}%
\pgfpathlineto{\pgfqpoint{1.104300in}{0.643117in}}%
\pgfpathlineto{\pgfqpoint{1.105737in}{0.642718in}}%
\pgfpathlineto{\pgfqpoint{1.107106in}{0.642276in}}%
\pgfpathclose%
\pgfusepath{fill}%
\end{pgfscope}%
\begin{pgfscope}%
\pgfpathrectangle{\pgfqpoint{0.407592in}{0.222904in}}{\pgfqpoint{1.411163in}{1.411163in}}%
\pgfusepath{clip}%
\pgfsetbuttcap%
\pgfsetroundjoin%
\definecolor{currentfill}{rgb}{0.749020,0.772549,0.792157}%
\pgfsetfillcolor{currentfill}%
\pgfsetfillopacity{0.160000}%
\pgfsetlinewidth{0.000000pt}%
\definecolor{currentstroke}{rgb}{0.000000,0.000000,0.000000}%
\pgfsetstrokecolor{currentstroke}%
\pgfsetdash{}{0pt}%
\pgfpathmoveto{\pgfqpoint{1.008494in}{0.728220in}}%
\pgfpathlineto{\pgfqpoint{0.998307in}{0.742888in}}%
\pgfpathlineto{\pgfqpoint{0.988126in}{0.757547in}}%
\pgfpathlineto{\pgfqpoint{0.982781in}{0.754535in}}%
\pgfpathlineto{\pgfqpoint{0.977892in}{0.751386in}}%
\pgfpathlineto{\pgfqpoint{0.989111in}{0.737319in}}%
\pgfpathlineto{\pgfqpoint{1.000333in}{0.723248in}}%
\pgfpathlineto{\pgfqpoint{1.004230in}{0.725789in}}%
\pgfpathlineto{\pgfqpoint{1.008494in}{0.728220in}}%
\pgfpathclose%
\pgfusepath{fill}%
\end{pgfscope}%
\begin{pgfscope}%
\pgfpathrectangle{\pgfqpoint{0.407592in}{0.222904in}}{\pgfqpoint{1.411163in}{1.411163in}}%
\pgfusepath{clip}%
\pgfsetbuttcap%
\pgfsetroundjoin%
\definecolor{currentfill}{rgb}{0.749020,0.772549,0.792157}%
\pgfsetfillcolor{currentfill}%
\pgfsetfillopacity{0.160000}%
\pgfsetlinewidth{0.000000pt}%
\definecolor{currentstroke}{rgb}{0.000000,0.000000,0.000000}%
\pgfsetstrokecolor{currentstroke}%
\pgfsetdash{}{0pt}%
\pgfpathmoveto{\pgfqpoint{1.077772in}{0.643040in}}%
\pgfpathlineto{\pgfqpoint{1.071557in}{0.659258in}}%
\pgfpathlineto{\pgfqpoint{1.065351in}{0.675451in}}%
\pgfpathlineto{\pgfqpoint{1.062338in}{0.674643in}}%
\pgfpathlineto{\pgfqpoint{1.059445in}{0.673753in}}%
\pgfpathlineto{\pgfqpoint{1.067127in}{0.657978in}}%
\pgfpathlineto{\pgfqpoint{1.074818in}{0.642183in}}%
\pgfpathlineto{\pgfqpoint{1.076265in}{0.642632in}}%
\pgfpathlineto{\pgfqpoint{1.077772in}{0.643040in}}%
\pgfpathclose%
\pgfusepath{fill}%
\end{pgfscope}%
\begin{pgfscope}%
\pgfpathrectangle{\pgfqpoint{0.407592in}{0.222904in}}{\pgfqpoint{1.411163in}{1.411163in}}%
\pgfusepath{clip}%
\pgfsetbuttcap%
\pgfsetroundjoin%
\definecolor{currentfill}{rgb}{0.749020,0.772549,0.792157}%
\pgfsetfillcolor{currentfill}%
\pgfsetfillopacity{0.160000}%
\pgfsetlinewidth{0.000000pt}%
\definecolor{currentstroke}{rgb}{0.000000,0.000000,0.000000}%
\pgfsetstrokecolor{currentstroke}%
\pgfsetdash{}{0pt}%
\pgfpathmoveto{\pgfqpoint{1.461869in}{1.015428in}}%
\pgfpathlineto{\pgfqpoint{1.474269in}{1.028937in}}%
\pgfpathlineto{\pgfqpoint{1.486669in}{1.042447in}}%
\pgfpathlineto{\pgfqpoint{1.475235in}{1.052261in}}%
\pgfpathlineto{\pgfqpoint{1.462227in}{1.061705in}}%
\pgfpathlineto{\pgfqpoint{1.450660in}{1.047676in}}%
\pgfpathlineto{\pgfqpoint{1.439090in}{1.033644in}}%
\pgfpathlineto{\pgfqpoint{1.451221in}{1.024709in}}%
\pgfpathlineto{\pgfqpoint{1.461869in}{1.015428in}}%
\pgfpathclose%
\pgfusepath{fill}%
\end{pgfscope}%
\begin{pgfscope}%
\pgfpathrectangle{\pgfqpoint{0.407592in}{0.222904in}}{\pgfqpoint{1.411163in}{1.411163in}}%
\pgfusepath{clip}%
\pgfsetbuttcap%
\pgfsetroundjoin%
\definecolor{currentfill}{rgb}{0.749020,0.772549,0.792157}%
\pgfsetfillcolor{currentfill}%
\pgfsetfillopacity{0.160000}%
\pgfsetlinewidth{0.000000pt}%
\definecolor{currentstroke}{rgb}{0.000000,0.000000,0.000000}%
\pgfsetstrokecolor{currentstroke}%
\pgfsetdash{}{0pt}%
\pgfpathmoveto{\pgfqpoint{0.620573in}{1.199406in}}%
\pgfpathlineto{\pgfqpoint{0.609459in}{1.213340in}}%
\pgfpathlineto{\pgfqpoint{0.598349in}{1.227271in}}%
\pgfpathlineto{\pgfqpoint{0.577379in}{1.214888in}}%
\pgfpathlineto{\pgfqpoint{0.558409in}{1.202004in}}%
\pgfpathlineto{\pgfqpoint{0.570522in}{1.188533in}}%
\pgfpathlineto{\pgfqpoint{0.582634in}{1.175063in}}%
\pgfpathlineto{\pgfqpoint{0.600643in}{1.187473in}}%
\pgfpathlineto{\pgfqpoint{0.620573in}{1.199406in}}%
\pgfpathclose%
\pgfusepath{fill}%
\end{pgfscope}%
\begin{pgfscope}%
\pgfpathrectangle{\pgfqpoint{0.407592in}{0.222904in}}{\pgfqpoint{1.411163in}{1.411163in}}%
\pgfusepath{clip}%
\pgfsetbuttcap%
\pgfsetroundjoin%
\definecolor{currentfill}{rgb}{0.749020,0.772549,0.792157}%
\pgfsetfillcolor{currentfill}%
\pgfsetfillopacity{0.160000}%
\pgfsetlinewidth{0.000000pt}%
\definecolor{currentstroke}{rgb}{0.000000,0.000000,0.000000}%
\pgfsetstrokecolor{currentstroke}%
\pgfsetdash{}{0pt}%
\pgfpathmoveto{\pgfqpoint{1.054067in}{0.671742in}}%
\pgfpathlineto{\pgfqpoint{1.045048in}{0.687008in}}%
\pgfpathlineto{\pgfqpoint{1.036038in}{0.702258in}}%
\pgfpathlineto{\pgfqpoint{1.032339in}{0.700605in}}%
\pgfpathlineto{\pgfqpoint{1.028888in}{0.698855in}}%
\pgfpathlineto{\pgfqpoint{1.039095in}{0.684159in}}%
\pgfpathlineto{\pgfqpoint{1.049309in}{0.669453in}}%
\pgfpathlineto{\pgfqpoint{1.051605in}{0.670630in}}%
\pgfpathlineto{\pgfqpoint{1.054067in}{0.671742in}}%
\pgfpathclose%
\pgfusepath{fill}%
\end{pgfscope}%
\begin{pgfscope}%
\pgfpathrectangle{\pgfqpoint{0.407592in}{0.222904in}}{\pgfqpoint{1.411163in}{1.411163in}}%
\pgfusepath{clip}%
\pgfsetbuttcap%
\pgfsetroundjoin%
\definecolor{currentfill}{rgb}{0.749020,0.772549,0.792157}%
\pgfsetfillcolor{currentfill}%
\pgfsetfillopacity{0.160000}%
\pgfsetlinewidth{0.000000pt}%
\definecolor{currentstroke}{rgb}{0.000000,0.000000,0.000000}%
\pgfsetstrokecolor{currentstroke}%
\pgfsetdash{}{0pt}%
\pgfpathmoveto{\pgfqpoint{0.642809in}{1.171525in}}%
\pgfpathlineto{\pgfqpoint{0.631689in}{1.185467in}}%
\pgfpathlineto{\pgfqpoint{0.620573in}{1.199406in}}%
\pgfpathlineto{\pgfqpoint{0.600643in}{1.187473in}}%
\pgfpathlineto{\pgfqpoint{0.582634in}{1.175063in}}%
\pgfpathlineto{\pgfqpoint{0.594744in}{1.161594in}}%
\pgfpathlineto{\pgfqpoint{0.606853in}{1.148127in}}%
\pgfpathlineto{\pgfqpoint{0.623911in}{1.160053in}}%
\pgfpathlineto{\pgfqpoint{0.642809in}{1.171525in}}%
\pgfpathclose%
\pgfusepath{fill}%
\end{pgfscope}%
\begin{pgfscope}%
\pgfpathrectangle{\pgfqpoint{0.407592in}{0.222904in}}{\pgfqpoint{1.411163in}{1.411163in}}%
\pgfusepath{clip}%
\pgfsetbuttcap%
\pgfsetroundjoin%
\definecolor{currentfill}{rgb}{0.749020,0.772549,0.792157}%
\pgfsetfillcolor{currentfill}%
\pgfsetfillopacity{0.160000}%
\pgfsetlinewidth{0.000000pt}%
\definecolor{currentstroke}{rgb}{0.000000,0.000000,0.000000}%
\pgfsetstrokecolor{currentstroke}%
\pgfsetdash{}{0pt}%
\pgfpathmoveto{\pgfqpoint{1.437073in}{0.988413in}}%
\pgfpathlineto{\pgfqpoint{1.449471in}{1.001920in}}%
\pgfpathlineto{\pgfqpoint{1.461869in}{1.015428in}}%
\pgfpathlineto{\pgfqpoint{1.451221in}{1.024709in}}%
\pgfpathlineto{\pgfqpoint{1.439090in}{1.033644in}}%
\pgfpathlineto{\pgfqpoint{1.427516in}{1.019606in}}%
\pgfpathlineto{\pgfqpoint{1.415938in}{1.005564in}}%
\pgfpathlineto{\pgfqpoint{1.427201in}{0.997150in}}%
\pgfpathlineto{\pgfqpoint{1.437073in}{0.988413in}}%
\pgfpathclose%
\pgfusepath{fill}%
\end{pgfscope}%
\begin{pgfscope}%
\pgfpathrectangle{\pgfqpoint{0.407592in}{0.222904in}}{\pgfqpoint{1.411163in}{1.411163in}}%
\pgfusepath{clip}%
\pgfsetbuttcap%
\pgfsetroundjoin%
\definecolor{currentfill}{rgb}{0.749020,0.772549,0.792157}%
\pgfsetfillcolor{currentfill}%
\pgfsetfillopacity{0.160000}%
\pgfsetlinewidth{0.000000pt}%
\definecolor{currentstroke}{rgb}{0.000000,0.000000,0.000000}%
\pgfsetstrokecolor{currentstroke}%
\pgfsetdash{}{0pt}%
\pgfpathmoveto{\pgfqpoint{0.665058in}{1.143628in}}%
\pgfpathlineto{\pgfqpoint{0.653932in}{1.157578in}}%
\pgfpathlineto{\pgfqpoint{0.642809in}{1.171525in}}%
\pgfpathlineto{\pgfqpoint{0.623911in}{1.160053in}}%
\pgfpathlineto{\pgfqpoint{0.606853in}{1.148127in}}%
\pgfpathlineto{\pgfqpoint{0.618961in}{1.134661in}}%
\pgfpathlineto{\pgfqpoint{0.631067in}{1.121196in}}%
\pgfpathlineto{\pgfqpoint{0.647184in}{1.132628in}}%
\pgfpathlineto{\pgfqpoint{0.665058in}{1.143628in}}%
\pgfpathclose%
\pgfusepath{fill}%
\end{pgfscope}%
\begin{pgfscope}%
\pgfpathrectangle{\pgfqpoint{0.407592in}{0.222904in}}{\pgfqpoint{1.411163in}{1.411163in}}%
\pgfusepath{clip}%
\pgfsetbuttcap%
\pgfsetroundjoin%
\definecolor{currentfill}{rgb}{0.749020,0.772549,0.792157}%
\pgfsetfillcolor{currentfill}%
\pgfsetfillopacity{0.160000}%
\pgfsetlinewidth{0.000000pt}%
\definecolor{currentstroke}{rgb}{0.000000,0.000000,0.000000}%
\pgfsetstrokecolor{currentstroke}%
\pgfsetdash{}{0pt}%
\pgfpathmoveto{\pgfqpoint{1.412279in}{0.961401in}}%
\pgfpathlineto{\pgfqpoint{1.424676in}{0.974906in}}%
\pgfpathlineto{\pgfqpoint{1.437073in}{0.988413in}}%
\pgfpathlineto{\pgfqpoint{1.427201in}{0.997150in}}%
\pgfpathlineto{\pgfqpoint{1.415938in}{1.005564in}}%
\pgfpathlineto{\pgfqpoint{1.404357in}{0.991518in}}%
\pgfpathlineto{\pgfqpoint{1.392771in}{0.977467in}}%
\pgfpathlineto{\pgfqpoint{1.403174in}{0.969583in}}%
\pgfpathlineto{\pgfqpoint{1.412279in}{0.961401in}}%
\pgfpathclose%
\pgfusepath{fill}%
\end{pgfscope}%
\begin{pgfscope}%
\pgfpathrectangle{\pgfqpoint{0.407592in}{0.222904in}}{\pgfqpoint{1.411163in}{1.411163in}}%
\pgfusepath{clip}%
\pgfsetbuttcap%
\pgfsetroundjoin%
\definecolor{currentfill}{rgb}{0.749020,0.772549,0.792157}%
\pgfsetfillcolor{currentfill}%
\pgfsetfillopacity{0.160000}%
\pgfsetlinewidth{0.000000pt}%
\definecolor{currentstroke}{rgb}{0.000000,0.000000,0.000000}%
\pgfsetstrokecolor{currentstroke}%
\pgfsetdash{}{0pt}%
\pgfpathmoveto{\pgfqpoint{1.160274in}{0.695485in}}%
\pgfpathlineto{\pgfqpoint{1.171935in}{0.709628in}}%
\pgfpathlineto{\pgfqpoint{1.183592in}{0.723766in}}%
\pgfpathlineto{\pgfqpoint{1.180144in}{0.726286in}}%
\pgfpathlineto{\pgfqpoint{1.176309in}{0.728694in}}%
\pgfpathlineto{\pgfqpoint{1.165575in}{0.713959in}}%
\pgfpathlineto{\pgfqpoint{1.154833in}{0.699214in}}%
\pgfpathlineto{\pgfqpoint{1.157699in}{0.697392in}}%
\pgfpathlineto{\pgfqpoint{1.160274in}{0.695485in}}%
\pgfpathclose%
\pgfusepath{fill}%
\end{pgfscope}%
\begin{pgfscope}%
\pgfpathrectangle{\pgfqpoint{0.407592in}{0.222904in}}{\pgfqpoint{1.411163in}{1.411163in}}%
\pgfusepath{clip}%
\pgfsetbuttcap%
\pgfsetroundjoin%
\definecolor{currentfill}{rgb}{0.749020,0.772549,0.792157}%
\pgfsetfillcolor{currentfill}%
\pgfsetfillopacity{0.160000}%
\pgfsetlinewidth{0.000000pt}%
\definecolor{currentstroke}{rgb}{0.000000,0.000000,0.000000}%
\pgfsetstrokecolor{currentstroke}%
\pgfsetdash{}{0pt}%
\pgfpathmoveto{\pgfqpoint{0.687320in}{1.115715in}}%
\pgfpathlineto{\pgfqpoint{0.676188in}{1.129674in}}%
\pgfpathlineto{\pgfqpoint{0.665058in}{1.143628in}}%
\pgfpathlineto{\pgfqpoint{0.647184in}{1.132628in}}%
\pgfpathlineto{\pgfqpoint{0.631067in}{1.121196in}}%
\pgfpathlineto{\pgfqpoint{0.643173in}{1.107733in}}%
\pgfpathlineto{\pgfqpoint{0.655277in}{1.094271in}}%
\pgfpathlineto{\pgfqpoint{0.670462in}{1.105197in}}%
\pgfpathlineto{\pgfqpoint{0.687320in}{1.115715in}}%
\pgfpathclose%
\pgfusepath{fill}%
\end{pgfscope}%
\begin{pgfscope}%
\pgfpathrectangle{\pgfqpoint{0.407592in}{0.222904in}}{\pgfqpoint{1.411163in}{1.411163in}}%
\pgfusepath{clip}%
\pgfsetbuttcap%
\pgfsetroundjoin%
\definecolor{currentfill}{rgb}{0.749020,0.772549,0.792157}%
\pgfsetfillcolor{currentfill}%
\pgfsetfillopacity{0.160000}%
\pgfsetlinewidth{0.000000pt}%
\definecolor{currentstroke}{rgb}{0.000000,0.000000,0.000000}%
\pgfsetstrokecolor{currentstroke}%
\pgfsetdash{}{0pt}%
\pgfpathmoveto{\pgfqpoint{1.109619in}{0.641276in}}%
\pgfpathlineto{\pgfqpoint{1.119298in}{0.656624in}}%
\pgfpathlineto{\pgfqpoint{1.128969in}{0.671957in}}%
\pgfpathlineto{\pgfqpoint{1.126532in}{0.672985in}}%
\pgfpathlineto{\pgfqpoint{1.123939in}{0.673938in}}%
\pgfpathlineto{\pgfqpoint{1.115528in}{0.658117in}}%
\pgfpathlineto{\pgfqpoint{1.107106in}{0.642276in}}%
\pgfpathlineto{\pgfqpoint{1.108402in}{0.641795in}}%
\pgfpathlineto{\pgfqpoint{1.109619in}{0.641276in}}%
\pgfpathclose%
\pgfusepath{fill}%
\end{pgfscope}%
\begin{pgfscope}%
\pgfpathrectangle{\pgfqpoint{0.407592in}{0.222904in}}{\pgfqpoint{1.411163in}{1.411163in}}%
\pgfusepath{clip}%
\pgfsetbuttcap%
\pgfsetroundjoin%
\definecolor{currentfill}{rgb}{0.749020,0.772549,0.792157}%
\pgfsetfillcolor{currentfill}%
\pgfsetfillopacity{0.160000}%
\pgfsetlinewidth{0.000000pt}%
\definecolor{currentstroke}{rgb}{0.000000,0.000000,0.000000}%
\pgfsetstrokecolor{currentstroke}%
\pgfsetdash{}{0pt}%
\pgfpathmoveto{\pgfqpoint{1.074818in}{0.642183in}}%
\pgfpathlineto{\pgfqpoint{1.067127in}{0.657978in}}%
\pgfpathlineto{\pgfqpoint{1.059445in}{0.673753in}}%
\pgfpathlineto{\pgfqpoint{1.056684in}{0.672785in}}%
\pgfpathlineto{\pgfqpoint{1.054067in}{0.671742in}}%
\pgfpathlineto{\pgfqpoint{1.063094in}{0.656462in}}%
\pgfpathlineto{\pgfqpoint{1.072130in}{0.641167in}}%
\pgfpathlineto{\pgfqpoint{1.073438in}{0.641694in}}%
\pgfpathlineto{\pgfqpoint{1.074818in}{0.642183in}}%
\pgfpathclose%
\pgfusepath{fill}%
\end{pgfscope}%
\begin{pgfscope}%
\pgfpathrectangle{\pgfqpoint{0.407592in}{0.222904in}}{\pgfqpoint{1.411163in}{1.411163in}}%
\pgfusepath{clip}%
\pgfsetbuttcap%
\pgfsetroundjoin%
\definecolor{currentfill}{rgb}{0.749020,0.772549,0.792157}%
\pgfsetfillcolor{currentfill}%
\pgfsetfillopacity{0.160000}%
\pgfsetlinewidth{0.000000pt}%
\definecolor{currentstroke}{rgb}{0.000000,0.000000,0.000000}%
\pgfsetstrokecolor{currentstroke}%
\pgfsetdash{}{0pt}%
\pgfpathmoveto{\pgfqpoint{1.387488in}{0.934391in}}%
\pgfpathlineto{\pgfqpoint{1.399883in}{0.947896in}}%
\pgfpathlineto{\pgfqpoint{1.412279in}{0.961401in}}%
\pgfpathlineto{\pgfqpoint{1.403174in}{0.969583in}}%
\pgfpathlineto{\pgfqpoint{1.392771in}{0.977467in}}%
\pgfpathlineto{\pgfqpoint{1.381182in}{0.963411in}}%
\pgfpathlineto{\pgfqpoint{1.369589in}{0.949351in}}%
\pgfpathlineto{\pgfqpoint{1.379141in}{0.942009in}}%
\pgfpathlineto{\pgfqpoint{1.387488in}{0.934391in}}%
\pgfpathclose%
\pgfusepath{fill}%
\end{pgfscope}%
\begin{pgfscope}%
\pgfpathrectangle{\pgfqpoint{0.407592in}{0.222904in}}{\pgfqpoint{1.411163in}{1.411163in}}%
\pgfusepath{clip}%
\pgfsetbuttcap%
\pgfsetroundjoin%
\definecolor{currentfill}{rgb}{0.749020,0.772549,0.792157}%
\pgfsetfillcolor{currentfill}%
\pgfsetfillopacity{0.160000}%
\pgfsetlinewidth{0.000000pt}%
\definecolor{currentstroke}{rgb}{0.000000,0.000000,0.000000}%
\pgfsetstrokecolor{currentstroke}%
\pgfsetdash{}{0pt}%
\pgfpathmoveto{\pgfqpoint{0.709595in}{1.087786in}}%
\pgfpathlineto{\pgfqpoint{0.698456in}{1.101753in}}%
\pgfpathlineto{\pgfqpoint{0.687320in}{1.115715in}}%
\pgfpathlineto{\pgfqpoint{0.670462in}{1.105197in}}%
\pgfpathlineto{\pgfqpoint{0.655277in}{1.094271in}}%
\pgfpathlineto{\pgfqpoint{0.667380in}{1.080811in}}%
\pgfpathlineto{\pgfqpoint{0.679482in}{1.067351in}}%
\pgfpathlineto{\pgfqpoint{0.693743in}{1.077761in}}%
\pgfpathlineto{\pgfqpoint{0.709595in}{1.087786in}}%
\pgfpathclose%
\pgfusepath{fill}%
\end{pgfscope}%
\begin{pgfscope}%
\pgfpathrectangle{\pgfqpoint{0.407592in}{0.222904in}}{\pgfqpoint{1.411163in}{1.411163in}}%
\pgfusepath{clip}%
\pgfsetbuttcap%
\pgfsetroundjoin%
\definecolor{currentfill}{rgb}{0.749020,0.772549,0.792157}%
\pgfsetfillcolor{currentfill}%
\pgfsetfillopacity{0.160000}%
\pgfsetlinewidth{0.000000pt}%
\definecolor{currentstroke}{rgb}{0.000000,0.000000,0.000000}%
\pgfsetstrokecolor{currentstroke}%
\pgfsetdash{}{0pt}%
\pgfpathmoveto{\pgfqpoint{1.028888in}{0.698855in}}%
\pgfpathlineto{\pgfqpoint{1.018688in}{0.713543in}}%
\pgfpathlineto{\pgfqpoint{1.008494in}{0.728220in}}%
\pgfpathlineto{\pgfqpoint{1.004230in}{0.725789in}}%
\pgfpathlineto{\pgfqpoint{1.000333in}{0.723248in}}%
\pgfpathlineto{\pgfqpoint{1.011559in}{0.709173in}}%
\pgfpathlineto{\pgfqpoint{1.022788in}{0.695094in}}%
\pgfpathlineto{\pgfqpoint{1.025700in}{0.697016in}}%
\pgfpathlineto{\pgfqpoint{1.028888in}{0.698855in}}%
\pgfpathclose%
\pgfusepath{fill}%
\end{pgfscope}%
\begin{pgfscope}%
\pgfpathrectangle{\pgfqpoint{0.407592in}{0.222904in}}{\pgfqpoint{1.411163in}{1.411163in}}%
\pgfusepath{clip}%
\pgfsetbuttcap%
\pgfsetroundjoin%
\definecolor{currentfill}{rgb}{0.749020,0.772549,0.792157}%
\pgfsetfillcolor{currentfill}%
\pgfsetfillopacity{0.160000}%
\pgfsetlinewidth{0.000000pt}%
\definecolor{currentstroke}{rgb}{0.000000,0.000000,0.000000}%
\pgfsetstrokecolor{currentstroke}%
\pgfsetdash{}{0pt}%
\pgfpathmoveto{\pgfqpoint{0.731883in}{1.059841in}}%
\pgfpathlineto{\pgfqpoint{0.720737in}{1.073816in}}%
\pgfpathlineto{\pgfqpoint{0.709595in}{1.087786in}}%
\pgfpathlineto{\pgfqpoint{0.693743in}{1.077761in}}%
\pgfpathlineto{\pgfqpoint{0.679482in}{1.067351in}}%
\pgfpathlineto{\pgfqpoint{0.691583in}{1.053893in}}%
\pgfpathlineto{\pgfqpoint{0.703682in}{1.040437in}}%
\pgfpathlineto{\pgfqpoint{0.717030in}{1.050320in}}%
\pgfpathlineto{\pgfqpoint{0.731883in}{1.059841in}}%
\pgfpathclose%
\pgfusepath{fill}%
\end{pgfscope}%
\begin{pgfscope}%
\pgfpathrectangle{\pgfqpoint{0.407592in}{0.222904in}}{\pgfqpoint{1.411163in}{1.411163in}}%
\pgfusepath{clip}%
\pgfsetbuttcap%
\pgfsetroundjoin%
\definecolor{currentfill}{rgb}{0.749020,0.772549,0.792157}%
\pgfsetfillcolor{currentfill}%
\pgfsetfillopacity{0.160000}%
\pgfsetlinewidth{0.000000pt}%
\definecolor{currentstroke}{rgb}{0.000000,0.000000,0.000000}%
\pgfsetstrokecolor{currentstroke}%
\pgfsetdash{}{0pt}%
\pgfpathmoveto{\pgfqpoint{1.362701in}{0.907386in}}%
\pgfpathlineto{\pgfqpoint{1.375094in}{0.920888in}}%
\pgfpathlineto{\pgfqpoint{1.387488in}{0.934391in}}%
\pgfpathlineto{\pgfqpoint{1.379141in}{0.942009in}}%
\pgfpathlineto{\pgfqpoint{1.369589in}{0.949351in}}%
\pgfpathlineto{\pgfqpoint{1.357993in}{0.935286in}}%
\pgfpathlineto{\pgfqpoint{1.346392in}{0.921217in}}%
\pgfpathlineto{\pgfqpoint{1.355101in}{0.914427in}}%
\pgfpathlineto{\pgfqpoint{1.362701in}{0.907386in}}%
\pgfpathclose%
\pgfusepath{fill}%
\end{pgfscope}%
\begin{pgfscope}%
\pgfpathrectangle{\pgfqpoint{0.407592in}{0.222904in}}{\pgfqpoint{1.411163in}{1.411163in}}%
\pgfusepath{clip}%
\pgfsetbuttcap%
\pgfsetroundjoin%
\definecolor{currentfill}{rgb}{0.749020,0.772549,0.792157}%
\pgfsetfillcolor{currentfill}%
\pgfsetfillopacity{0.160000}%
\pgfsetlinewidth{0.000000pt}%
\definecolor{currentstroke}{rgb}{0.000000,0.000000,0.000000}%
\pgfsetstrokecolor{currentstroke}%
\pgfsetdash{}{0pt}%
\pgfpathmoveto{\pgfqpoint{0.754183in}{1.031880in}}%
\pgfpathlineto{\pgfqpoint{0.743031in}{1.045863in}}%
\pgfpathlineto{\pgfqpoint{0.731883in}{1.059841in}}%
\pgfpathlineto{\pgfqpoint{0.717030in}{1.050320in}}%
\pgfpathlineto{\pgfqpoint{0.703682in}{1.040437in}}%
\pgfpathlineto{\pgfqpoint{0.715780in}{1.026982in}}%
\pgfpathlineto{\pgfqpoint{0.727877in}{1.013528in}}%
\pgfpathlineto{\pgfqpoint{0.740321in}{1.022874in}}%
\pgfpathlineto{\pgfqpoint{0.754183in}{1.031880in}}%
\pgfpathclose%
\pgfusepath{fill}%
\end{pgfscope}%
\begin{pgfscope}%
\pgfpathrectangle{\pgfqpoint{0.407592in}{0.222904in}}{\pgfqpoint{1.411163in}{1.411163in}}%
\pgfusepath{clip}%
\pgfsetbuttcap%
\pgfsetroundjoin%
\definecolor{currentfill}{rgb}{0.749020,0.772549,0.792157}%
\pgfsetfillcolor{currentfill}%
\pgfsetfillopacity{0.160000}%
\pgfsetlinewidth{0.000000pt}%
\definecolor{currentstroke}{rgb}{0.000000,0.000000,0.000000}%
\pgfsetstrokecolor{currentstroke}%
\pgfsetdash{}{0pt}%
\pgfpathmoveto{\pgfqpoint{1.337916in}{0.880383in}}%
\pgfpathlineto{\pgfqpoint{1.350308in}{0.893884in}}%
\pgfpathlineto{\pgfqpoint{1.362701in}{0.907386in}}%
\pgfpathlineto{\pgfqpoint{1.355101in}{0.914427in}}%
\pgfpathlineto{\pgfqpoint{1.346392in}{0.921217in}}%
\pgfpathlineto{\pgfqpoint{1.334788in}{0.907143in}}%
\pgfpathlineto{\pgfqpoint{1.323181in}{0.893064in}}%
\pgfpathlineto{\pgfqpoint{1.331054in}{0.886838in}}%
\pgfpathlineto{\pgfqpoint{1.337916in}{0.880383in}}%
\pgfpathclose%
\pgfusepath{fill}%
\end{pgfscope}%
\begin{pgfscope}%
\pgfpathrectangle{\pgfqpoint{0.407592in}{0.222904in}}{\pgfqpoint{1.411163in}{1.411163in}}%
\pgfusepath{clip}%
\pgfsetbuttcap%
\pgfsetroundjoin%
\definecolor{currentfill}{rgb}{0.749020,0.772549,0.792157}%
\pgfsetfillcolor{currentfill}%
\pgfsetfillopacity{0.160000}%
\pgfsetlinewidth{0.000000pt}%
\definecolor{currentstroke}{rgb}{0.000000,0.000000,0.000000}%
\pgfsetstrokecolor{currentstroke}%
\pgfsetdash{}{0pt}%
\pgfpathmoveto{\pgfqpoint{0.776496in}{1.003903in}}%
\pgfpathlineto{\pgfqpoint{0.765338in}{1.017894in}}%
\pgfpathlineto{\pgfqpoint{0.754183in}{1.031880in}}%
\pgfpathlineto{\pgfqpoint{0.740321in}{1.022874in}}%
\pgfpathlineto{\pgfqpoint{0.727877in}{1.013528in}}%
\pgfpathlineto{\pgfqpoint{0.739973in}{1.000075in}}%
\pgfpathlineto{\pgfqpoint{0.752068in}{0.986624in}}%
\pgfpathlineto{\pgfqpoint{0.763616in}{0.995422in}}%
\pgfpathlineto{\pgfqpoint{0.776496in}{1.003903in}}%
\pgfpathclose%
\pgfusepath{fill}%
\end{pgfscope}%
\begin{pgfscope}%
\pgfpathrectangle{\pgfqpoint{0.407592in}{0.222904in}}{\pgfqpoint{1.411163in}{1.411163in}}%
\pgfusepath{clip}%
\pgfsetbuttcap%
\pgfsetroundjoin%
\definecolor{currentfill}{rgb}{0.749020,0.772549,0.792157}%
\pgfsetfillcolor{currentfill}%
\pgfsetfillopacity{0.160000}%
\pgfsetlinewidth{0.000000pt}%
\definecolor{currentstroke}{rgb}{0.000000,0.000000,0.000000}%
\pgfsetstrokecolor{currentstroke}%
\pgfsetdash{}{0pt}%
\pgfpathmoveto{\pgfqpoint{0.798822in}{0.975910in}}%
\pgfpathlineto{\pgfqpoint{0.787657in}{0.989909in}}%
\pgfpathlineto{\pgfqpoint{0.776496in}{1.003903in}}%
\pgfpathlineto{\pgfqpoint{0.763616in}{0.995422in}}%
\pgfpathlineto{\pgfqpoint{0.752068in}{0.986624in}}%
\pgfpathlineto{\pgfqpoint{0.764161in}{0.973174in}}%
\pgfpathlineto{\pgfqpoint{0.776253in}{0.959726in}}%
\pgfpathlineto{\pgfqpoint{0.786916in}{0.967964in}}%
\pgfpathlineto{\pgfqpoint{0.798822in}{0.975910in}}%
\pgfpathclose%
\pgfusepath{fill}%
\end{pgfscope}%
\begin{pgfscope}%
\pgfpathrectangle{\pgfqpoint{0.407592in}{0.222904in}}{\pgfqpoint{1.411163in}{1.411163in}}%
\pgfusepath{clip}%
\pgfsetbuttcap%
\pgfsetroundjoin%
\definecolor{currentfill}{rgb}{0.749020,0.772549,0.792157}%
\pgfsetfillcolor{currentfill}%
\pgfsetfillopacity{0.160000}%
\pgfsetlinewidth{0.000000pt}%
\definecolor{currentstroke}{rgb}{0.000000,0.000000,0.000000}%
\pgfsetstrokecolor{currentstroke}%
\pgfsetdash{}{0pt}%
\pgfpathmoveto{\pgfqpoint{1.313134in}{0.853384in}}%
\pgfpathlineto{\pgfqpoint{1.325524in}{0.866883in}}%
\pgfpathlineto{\pgfqpoint{1.337916in}{0.880383in}}%
\pgfpathlineto{\pgfqpoint{1.331054in}{0.886838in}}%
\pgfpathlineto{\pgfqpoint{1.323181in}{0.893064in}}%
\pgfpathlineto{\pgfqpoint{1.311569in}{0.878981in}}%
\pgfpathlineto{\pgfqpoint{1.299953in}{0.864894in}}%
\pgfpathlineto{\pgfqpoint{1.307001in}{0.859241in}}%
\pgfpathlineto{\pgfqpoint{1.313134in}{0.853384in}}%
\pgfpathclose%
\pgfusepath{fill}%
\end{pgfscope}%
\begin{pgfscope}%
\pgfpathrectangle{\pgfqpoint{0.407592in}{0.222904in}}{\pgfqpoint{1.411163in}{1.411163in}}%
\pgfusepath{clip}%
\pgfsetbuttcap%
\pgfsetroundjoin%
\definecolor{currentfill}{rgb}{0.749020,0.772549,0.792157}%
\pgfsetfillcolor{currentfill}%
\pgfsetfillopacity{0.160000}%
\pgfsetlinewidth{0.000000pt}%
\definecolor{currentstroke}{rgb}{0.000000,0.000000,0.000000}%
\pgfsetstrokecolor{currentstroke}%
\pgfsetdash{}{0pt}%
\pgfpathmoveto{\pgfqpoint{1.111794in}{0.640134in}}%
\pgfpathlineto{\pgfqpoint{1.122564in}{0.654919in}}%
\pgfpathlineto{\pgfqpoint{1.133328in}{0.669694in}}%
\pgfpathlineto{\pgfqpoint{1.131237in}{0.670859in}}%
\pgfpathlineto{\pgfqpoint{1.128969in}{0.671957in}}%
\pgfpathlineto{\pgfqpoint{1.119298in}{0.656624in}}%
\pgfpathlineto{\pgfqpoint{1.109619in}{0.641276in}}%
\pgfpathlineto{\pgfqpoint{1.110751in}{0.640721in}}%
\pgfpathlineto{\pgfqpoint{1.111794in}{0.640134in}}%
\pgfpathclose%
\pgfusepath{fill}%
\end{pgfscope}%
\begin{pgfscope}%
\pgfpathrectangle{\pgfqpoint{0.407592in}{0.222904in}}{\pgfqpoint{1.411163in}{1.411163in}}%
\pgfusepath{clip}%
\pgfsetbuttcap%
\pgfsetroundjoin%
\definecolor{currentfill}{rgb}{0.749020,0.772549,0.792157}%
\pgfsetfillcolor{currentfill}%
\pgfsetfillopacity{0.160000}%
\pgfsetlinewidth{0.000000pt}%
\definecolor{currentstroke}{rgb}{0.000000,0.000000,0.000000}%
\pgfsetstrokecolor{currentstroke}%
\pgfsetdash{}{0pt}%
\pgfpathmoveto{\pgfqpoint{1.136941in}{0.667186in}}%
\pgfpathlineto{\pgfqpoint{1.148609in}{0.681338in}}%
\pgfpathlineto{\pgfqpoint{1.160274in}{0.695485in}}%
\pgfpathlineto{\pgfqpoint{1.157699in}{0.697392in}}%
\pgfpathlineto{\pgfqpoint{1.154833in}{0.699214in}}%
\pgfpathlineto{\pgfqpoint{1.144084in}{0.684459in}}%
\pgfpathlineto{\pgfqpoint{1.133328in}{0.669694in}}%
\pgfpathlineto{\pgfqpoint{1.135232in}{0.668468in}}%
\pgfpathlineto{\pgfqpoint{1.136941in}{0.667186in}}%
\pgfpathclose%
\pgfusepath{fill}%
\end{pgfscope}%
\begin{pgfscope}%
\pgfpathrectangle{\pgfqpoint{0.407592in}{0.222904in}}{\pgfqpoint{1.411163in}{1.411163in}}%
\pgfusepath{clip}%
\pgfsetbuttcap%
\pgfsetroundjoin%
\definecolor{currentfill}{rgb}{0.749020,0.772549,0.792157}%
\pgfsetfillcolor{currentfill}%
\pgfsetfillopacity{0.160000}%
\pgfsetlinewidth{0.000000pt}%
\definecolor{currentstroke}{rgb}{0.000000,0.000000,0.000000}%
\pgfsetstrokecolor{currentstroke}%
\pgfsetdash{}{0pt}%
\pgfpathmoveto{\pgfqpoint{1.072130in}{0.641167in}}%
\pgfpathlineto{\pgfqpoint{1.063094in}{0.656462in}}%
\pgfpathlineto{\pgfqpoint{1.054067in}{0.671742in}}%
\pgfpathlineto{\pgfqpoint{1.051605in}{0.670630in}}%
\pgfpathlineto{\pgfqpoint{1.049309in}{0.669453in}}%
\pgfpathlineto{\pgfqpoint{1.059529in}{0.654737in}}%
\pgfpathlineto{\pgfqpoint{1.069756in}{0.640012in}}%
\pgfpathlineto{\pgfqpoint{1.070901in}{0.640606in}}%
\pgfpathlineto{\pgfqpoint{1.072130in}{0.641167in}}%
\pgfpathclose%
\pgfusepath{fill}%
\end{pgfscope}%
\begin{pgfscope}%
\pgfpathrectangle{\pgfqpoint{0.407592in}{0.222904in}}{\pgfqpoint{1.411163in}{1.411163in}}%
\pgfusepath{clip}%
\pgfsetbuttcap%
\pgfsetroundjoin%
\definecolor{currentfill}{rgb}{0.749020,0.772549,0.792157}%
\pgfsetfillcolor{currentfill}%
\pgfsetfillopacity{0.160000}%
\pgfsetlinewidth{0.000000pt}%
\definecolor{currentstroke}{rgb}{0.000000,0.000000,0.000000}%
\pgfsetstrokecolor{currentstroke}%
\pgfsetdash{}{0pt}%
\pgfpathmoveto{\pgfqpoint{0.821161in}{0.947901in}}%
\pgfpathlineto{\pgfqpoint{0.809990in}{0.961908in}}%
\pgfpathlineto{\pgfqpoint{0.798822in}{0.975910in}}%
\pgfpathlineto{\pgfqpoint{0.786916in}{0.967964in}}%
\pgfpathlineto{\pgfqpoint{0.776253in}{0.959726in}}%
\pgfpathlineto{\pgfqpoint{0.788344in}{0.946279in}}%
\pgfpathlineto{\pgfqpoint{0.800434in}{0.932833in}}%
\pgfpathlineto{\pgfqpoint{0.810221in}{0.940502in}}%
\pgfpathlineto{\pgfqpoint{0.821161in}{0.947901in}}%
\pgfpathclose%
\pgfusepath{fill}%
\end{pgfscope}%
\begin{pgfscope}%
\pgfpathrectangle{\pgfqpoint{0.407592in}{0.222904in}}{\pgfqpoint{1.411163in}{1.411163in}}%
\pgfusepath{clip}%
\pgfsetbuttcap%
\pgfsetroundjoin%
\definecolor{currentfill}{rgb}{0.749020,0.772549,0.792157}%
\pgfsetfillcolor{currentfill}%
\pgfsetfillopacity{0.160000}%
\pgfsetlinewidth{0.000000pt}%
\definecolor{currentstroke}{rgb}{0.000000,0.000000,0.000000}%
\pgfsetstrokecolor{currentstroke}%
\pgfsetdash{}{0pt}%
\pgfpathmoveto{\pgfqpoint{1.049309in}{0.669453in}}%
\pgfpathlineto{\pgfqpoint{1.039095in}{0.684159in}}%
\pgfpathlineto{\pgfqpoint{1.028888in}{0.698855in}}%
\pgfpathlineto{\pgfqpoint{1.025700in}{0.697016in}}%
\pgfpathlineto{\pgfqpoint{1.022788in}{0.695094in}}%
\pgfpathlineto{\pgfqpoint{1.034020in}{0.681010in}}%
\pgfpathlineto{\pgfqpoint{1.045256in}{0.666923in}}%
\pgfpathlineto{\pgfqpoint{1.047189in}{0.668215in}}%
\pgfpathlineto{\pgfqpoint{1.049309in}{0.669453in}}%
\pgfpathclose%
\pgfusepath{fill}%
\end{pgfscope}%
\begin{pgfscope}%
\pgfpathrectangle{\pgfqpoint{0.407592in}{0.222904in}}{\pgfqpoint{1.411163in}{1.411163in}}%
\pgfusepath{clip}%
\pgfsetbuttcap%
\pgfsetroundjoin%
\definecolor{currentfill}{rgb}{0.749020,0.772549,0.792157}%
\pgfsetfillcolor{currentfill}%
\pgfsetfillopacity{0.160000}%
\pgfsetlinewidth{0.000000pt}%
\definecolor{currentstroke}{rgb}{0.000000,0.000000,0.000000}%
\pgfsetstrokecolor{currentstroke}%
\pgfsetdash{}{0pt}%
\pgfpathmoveto{\pgfqpoint{1.288355in}{0.826387in}}%
\pgfpathlineto{\pgfqpoint{1.300744in}{0.839885in}}%
\pgfpathlineto{\pgfqpoint{1.313134in}{0.853384in}}%
\pgfpathlineto{\pgfqpoint{1.307001in}{0.859241in}}%
\pgfpathlineto{\pgfqpoint{1.299953in}{0.864894in}}%
\pgfpathlineto{\pgfqpoint{1.288334in}{0.850802in}}%
\pgfpathlineto{\pgfqpoint{1.276711in}{0.836705in}}%
\pgfpathlineto{\pgfqpoint{1.282942in}{0.831637in}}%
\pgfpathlineto{\pgfqpoint{1.288355in}{0.826387in}}%
\pgfpathclose%
\pgfusepath{fill}%
\end{pgfscope}%
\begin{pgfscope}%
\pgfpathrectangle{\pgfqpoint{0.407592in}{0.222904in}}{\pgfqpoint{1.411163in}{1.411163in}}%
\pgfusepath{clip}%
\pgfsetbuttcap%
\pgfsetroundjoin%
\definecolor{currentfill}{rgb}{0.749020,0.772549,0.792157}%
\pgfsetfillcolor{currentfill}%
\pgfsetfillopacity{0.160000}%
\pgfsetlinewidth{0.000000pt}%
\definecolor{currentstroke}{rgb}{0.000000,0.000000,0.000000}%
\pgfsetstrokecolor{currentstroke}%
\pgfsetdash{}{0pt}%
\pgfpathmoveto{\pgfqpoint{0.843512in}{0.919876in}}%
\pgfpathlineto{\pgfqpoint{0.832335in}{0.933891in}}%
\pgfpathlineto{\pgfqpoint{0.821161in}{0.947901in}}%
\pgfpathlineto{\pgfqpoint{0.810221in}{0.940502in}}%
\pgfpathlineto{\pgfqpoint{0.800434in}{0.932833in}}%
\pgfpathlineto{\pgfqpoint{0.812523in}{0.919388in}}%
\pgfpathlineto{\pgfqpoint{0.824610in}{0.905945in}}%
\pgfpathlineto{\pgfqpoint{0.833530in}{0.913034in}}%
\pgfpathlineto{\pgfqpoint{0.843512in}{0.919876in}}%
\pgfpathclose%
\pgfusepath{fill}%
\end{pgfscope}%
\begin{pgfscope}%
\pgfpathrectangle{\pgfqpoint{0.407592in}{0.222904in}}{\pgfqpoint{1.411163in}{1.411163in}}%
\pgfusepath{clip}%
\pgfsetbuttcap%
\pgfsetroundjoin%
\definecolor{currentfill}{rgb}{0.749020,0.772549,0.792157}%
\pgfsetfillcolor{currentfill}%
\pgfsetfillopacity{0.160000}%
\pgfsetlinewidth{0.000000pt}%
\definecolor{currentstroke}{rgb}{0.000000,0.000000,0.000000}%
\pgfsetstrokecolor{currentstroke}%
\pgfsetdash{}{0pt}%
\pgfpathmoveto{\pgfqpoint{1.263579in}{0.799394in}}%
\pgfpathlineto{\pgfqpoint{1.275967in}{0.812890in}}%
\pgfpathlineto{\pgfqpoint{1.288355in}{0.826387in}}%
\pgfpathlineto{\pgfqpoint{1.282942in}{0.831637in}}%
\pgfpathlineto{\pgfqpoint{1.276711in}{0.836705in}}%
\pgfpathlineto{\pgfqpoint{1.265085in}{0.822604in}}%
\pgfpathlineto{\pgfqpoint{1.253454in}{0.808498in}}%
\pgfpathlineto{\pgfqpoint{1.258876in}{0.804025in}}%
\pgfpathlineto{\pgfqpoint{1.263579in}{0.799394in}}%
\pgfpathclose%
\pgfusepath{fill}%
\end{pgfscope}%
\begin{pgfscope}%
\pgfpathrectangle{\pgfqpoint{0.407592in}{0.222904in}}{\pgfqpoint{1.411163in}{1.411163in}}%
\pgfusepath{clip}%
\pgfsetbuttcap%
\pgfsetroundjoin%
\definecolor{currentfill}{rgb}{0.749020,0.772549,0.792157}%
\pgfsetfillcolor{currentfill}%
\pgfsetfillopacity{0.160000}%
\pgfsetlinewidth{0.000000pt}%
\definecolor{currentstroke}{rgb}{0.000000,0.000000,0.000000}%
\pgfsetstrokecolor{currentstroke}%
\pgfsetdash{}{0pt}%
\pgfpathmoveto{\pgfqpoint{0.865877in}{0.891835in}}%
\pgfpathlineto{\pgfqpoint{0.854693in}{0.905857in}}%
\pgfpathlineto{\pgfqpoint{0.843512in}{0.919876in}}%
\pgfpathlineto{\pgfqpoint{0.833530in}{0.913034in}}%
\pgfpathlineto{\pgfqpoint{0.824610in}{0.905945in}}%
\pgfpathlineto{\pgfqpoint{0.836696in}{0.892503in}}%
\pgfpathlineto{\pgfqpoint{0.848781in}{0.879063in}}%
\pgfpathlineto{\pgfqpoint{0.856844in}{0.885560in}}%
\pgfpathlineto{\pgfqpoint{0.865877in}{0.891835in}}%
\pgfpathclose%
\pgfusepath{fill}%
\end{pgfscope}%
\begin{pgfscope}%
\pgfpathrectangle{\pgfqpoint{0.407592in}{0.222904in}}{\pgfqpoint{1.411163in}{1.411163in}}%
\pgfusepath{clip}%
\pgfsetbuttcap%
\pgfsetroundjoin%
\definecolor{currentfill}{rgb}{0.749020,0.772549,0.792157}%
\pgfsetfillcolor{currentfill}%
\pgfsetfillopacity{0.160000}%
\pgfsetlinewidth{0.000000pt}%
\definecolor{currentstroke}{rgb}{0.000000,0.000000,0.000000}%
\pgfsetstrokecolor{currentstroke}%
\pgfsetdash{}{0pt}%
\pgfpathmoveto{\pgfqpoint{1.090230in}{0.610533in}}%
\pgfpathlineto{\pgfqpoint{1.092431in}{0.627560in}}%
\pgfpathlineto{\pgfqpoint{1.094627in}{0.644555in}}%
\pgfpathlineto{\pgfqpoint{1.092905in}{0.644625in}}%
\pgfpathlineto{\pgfqpoint{1.091172in}{0.644645in}}%
\pgfpathlineto{\pgfqpoint{1.090701in}{0.627606in}}%
\pgfpathlineto{\pgfqpoint{1.090230in}{0.610533in}}%
\pgfpathlineto{\pgfqpoint{1.090230in}{0.610533in}}%
\pgfpathlineto{\pgfqpoint{1.090230in}{0.610533in}}%
\pgfpathclose%
\pgfusepath{fill}%
\end{pgfscope}%
\begin{pgfscope}%
\pgfpathrectangle{\pgfqpoint{0.407592in}{0.222904in}}{\pgfqpoint{1.411163in}{1.411163in}}%
\pgfusepath{clip}%
\pgfsetbuttcap%
\pgfsetroundjoin%
\definecolor{currentfill}{rgb}{0.749020,0.772549,0.792157}%
\pgfsetfillcolor{currentfill}%
\pgfsetfillopacity{0.160000}%
\pgfsetlinewidth{0.000000pt}%
\definecolor{currentstroke}{rgb}{0.000000,0.000000,0.000000}%
\pgfsetstrokecolor{currentstroke}%
\pgfsetdash{}{0pt}%
\pgfpathmoveto{\pgfqpoint{1.090230in}{0.610533in}}%
\pgfpathlineto{\pgfqpoint{1.090701in}{0.627606in}}%
\pgfpathlineto{\pgfqpoint{1.091172in}{0.644645in}}%
\pgfpathlineto{\pgfqpoint{1.089434in}{0.644615in}}%
\pgfpathlineto{\pgfqpoint{1.087699in}{0.644535in}}%
\pgfpathlineto{\pgfqpoint{1.088963in}{0.627550in}}%
\pgfpathlineto{\pgfqpoint{1.090230in}{0.610533in}}%
\pgfpathlineto{\pgfqpoint{1.090230in}{0.610533in}}%
\pgfpathlineto{\pgfqpoint{1.090230in}{0.610533in}}%
\pgfpathclose%
\pgfusepath{fill}%
\end{pgfscope}%
\begin{pgfscope}%
\pgfpathrectangle{\pgfqpoint{0.407592in}{0.222904in}}{\pgfqpoint{1.411163in}{1.411163in}}%
\pgfusepath{clip}%
\pgfsetbuttcap%
\pgfsetroundjoin%
\definecolor{currentfill}{rgb}{0.749020,0.772549,0.792157}%
\pgfsetfillcolor{currentfill}%
\pgfsetfillopacity{0.160000}%
\pgfsetlinewidth{0.000000pt}%
\definecolor{currentstroke}{rgb}{0.000000,0.000000,0.000000}%
\pgfsetstrokecolor{currentstroke}%
\pgfsetdash{}{0pt}%
\pgfpathmoveto{\pgfqpoint{1.090230in}{0.610533in}}%
\pgfpathlineto{\pgfqpoint{1.094121in}{0.627415in}}%
\pgfpathlineto{\pgfqpoint{1.098006in}{0.644266in}}%
\pgfpathlineto{\pgfqpoint{1.096330in}{0.644435in}}%
\pgfpathlineto{\pgfqpoint{1.094627in}{0.644555in}}%
\pgfpathlineto{\pgfqpoint{1.092431in}{0.627560in}}%
\pgfpathlineto{\pgfqpoint{1.090230in}{0.610533in}}%
\pgfpathlineto{\pgfqpoint{1.090230in}{0.610533in}}%
\pgfpathlineto{\pgfqpoint{1.090230in}{0.610533in}}%
\pgfpathclose%
\pgfusepath{fill}%
\end{pgfscope}%
\begin{pgfscope}%
\pgfpathrectangle{\pgfqpoint{0.407592in}{0.222904in}}{\pgfqpoint{1.411163in}{1.411163in}}%
\pgfusepath{clip}%
\pgfsetbuttcap%
\pgfsetroundjoin%
\definecolor{currentfill}{rgb}{0.749020,0.772549,0.792157}%
\pgfsetfillcolor{currentfill}%
\pgfsetfillopacity{0.160000}%
\pgfsetlinewidth{0.000000pt}%
\definecolor{currentstroke}{rgb}{0.000000,0.000000,0.000000}%
\pgfsetstrokecolor{currentstroke}%
\pgfsetdash{}{0pt}%
\pgfpathmoveto{\pgfqpoint{1.090230in}{0.610533in}}%
\pgfpathlineto{\pgfqpoint{1.088963in}{0.627550in}}%
\pgfpathlineto{\pgfqpoint{1.087699in}{0.644535in}}%
\pgfpathlineto{\pgfqpoint{1.085974in}{0.644405in}}%
\pgfpathlineto{\pgfqpoint{1.084268in}{0.644227in}}%
\pgfpathlineto{\pgfqpoint{1.087246in}{0.627395in}}%
\pgfpathlineto{\pgfqpoint{1.090230in}{0.610533in}}%
\pgfpathlineto{\pgfqpoint{1.090230in}{0.610533in}}%
\pgfpathlineto{\pgfqpoint{1.090230in}{0.610533in}}%
\pgfpathclose%
\pgfusepath{fill}%
\end{pgfscope}%
\begin{pgfscope}%
\pgfpathrectangle{\pgfqpoint{0.407592in}{0.222904in}}{\pgfqpoint{1.411163in}{1.411163in}}%
\pgfusepath{clip}%
\pgfsetbuttcap%
\pgfsetroundjoin%
\definecolor{currentfill}{rgb}{0.749020,0.772549,0.792157}%
\pgfsetfillcolor{currentfill}%
\pgfsetfillopacity{0.160000}%
\pgfsetlinewidth{0.000000pt}%
\definecolor{currentstroke}{rgb}{0.000000,0.000000,0.000000}%
\pgfsetstrokecolor{currentstroke}%
\pgfsetdash{}{0pt}%
\pgfpathmoveto{\pgfqpoint{0.888254in}{0.863777in}}%
\pgfpathlineto{\pgfqpoint{0.877064in}{0.877808in}}%
\pgfpathlineto{\pgfqpoint{0.865877in}{0.891835in}}%
\pgfpathlineto{\pgfqpoint{0.856844in}{0.885560in}}%
\pgfpathlineto{\pgfqpoint{0.848781in}{0.879063in}}%
\pgfpathlineto{\pgfqpoint{0.860865in}{0.865624in}}%
\pgfpathlineto{\pgfqpoint{0.872948in}{0.852186in}}%
\pgfpathlineto{\pgfqpoint{0.880162in}{0.858082in}}%
\pgfpathlineto{\pgfqpoint{0.888254in}{0.863777in}}%
\pgfpathclose%
\pgfusepath{fill}%
\end{pgfscope}%
\begin{pgfscope}%
\pgfpathrectangle{\pgfqpoint{0.407592in}{0.222904in}}{\pgfqpoint{1.411163in}{1.411163in}}%
\pgfusepath{clip}%
\pgfsetbuttcap%
\pgfsetroundjoin%
\definecolor{currentfill}{rgb}{0.749020,0.772549,0.792157}%
\pgfsetfillcolor{currentfill}%
\pgfsetfillopacity{0.160000}%
\pgfsetlinewidth{0.000000pt}%
\definecolor{currentstroke}{rgb}{0.000000,0.000000,0.000000}%
\pgfsetstrokecolor{currentstroke}%
\pgfsetdash{}{0pt}%
\pgfpathmoveto{\pgfqpoint{1.090230in}{0.610533in}}%
\pgfpathlineto{\pgfqpoint{1.095744in}{0.627173in}}%
\pgfpathlineto{\pgfqpoint{1.101249in}{0.643784in}}%
\pgfpathlineto{\pgfqpoint{1.099648in}{0.644049in}}%
\pgfpathlineto{\pgfqpoint{1.098006in}{0.644266in}}%
\pgfpathlineto{\pgfqpoint{1.094121in}{0.627415in}}%
\pgfpathlineto{\pgfqpoint{1.090230in}{0.610533in}}%
\pgfpathlineto{\pgfqpoint{1.090230in}{0.610533in}}%
\pgfpathlineto{\pgfqpoint{1.090230in}{0.610533in}}%
\pgfpathclose%
\pgfusepath{fill}%
\end{pgfscope}%
\begin{pgfscope}%
\pgfpathrectangle{\pgfqpoint{0.407592in}{0.222904in}}{\pgfqpoint{1.411163in}{1.411163in}}%
\pgfusepath{clip}%
\pgfsetbuttcap%
\pgfsetroundjoin%
\definecolor{currentfill}{rgb}{0.749020,0.772549,0.792157}%
\pgfsetfillcolor{currentfill}%
\pgfsetfillopacity{0.160000}%
\pgfsetlinewidth{0.000000pt}%
\definecolor{currentstroke}{rgb}{0.000000,0.000000,0.000000}%
\pgfsetstrokecolor{currentstroke}%
\pgfsetdash{}{0pt}%
\pgfpathmoveto{\pgfqpoint{1.090230in}{0.610533in}}%
\pgfpathlineto{\pgfqpoint{1.087246in}{0.627395in}}%
\pgfpathlineto{\pgfqpoint{1.084268in}{0.644227in}}%
\pgfpathlineto{\pgfqpoint{1.082588in}{0.644000in}}%
\pgfpathlineto{\pgfqpoint{1.080940in}{0.643726in}}%
\pgfpathlineto{\pgfqpoint{1.085581in}{0.627143in}}%
\pgfpathlineto{\pgfqpoint{1.090230in}{0.610533in}}%
\pgfpathlineto{\pgfqpoint{1.090230in}{0.610533in}}%
\pgfpathlineto{\pgfqpoint{1.090230in}{0.610533in}}%
\pgfpathclose%
\pgfusepath{fill}%
\end{pgfscope}%
\begin{pgfscope}%
\pgfpathrectangle{\pgfqpoint{0.407592in}{0.222904in}}{\pgfqpoint{1.411163in}{1.411163in}}%
\pgfusepath{clip}%
\pgfsetbuttcap%
\pgfsetroundjoin%
\definecolor{currentfill}{rgb}{0.749020,0.772549,0.792157}%
\pgfsetfillcolor{currentfill}%
\pgfsetfillopacity{0.160000}%
\pgfsetlinewidth{0.000000pt}%
\definecolor{currentstroke}{rgb}{0.000000,0.000000,0.000000}%
\pgfsetstrokecolor{currentstroke}%
\pgfsetdash{}{0pt}%
\pgfpathmoveto{\pgfqpoint{1.238806in}{0.772404in}}%
\pgfpathlineto{\pgfqpoint{1.251192in}{0.785899in}}%
\pgfpathlineto{\pgfqpoint{1.263579in}{0.799394in}}%
\pgfpathlineto{\pgfqpoint{1.258876in}{0.804025in}}%
\pgfpathlineto{\pgfqpoint{1.253454in}{0.808498in}}%
\pgfpathlineto{\pgfqpoint{1.241820in}{0.794387in}}%
\pgfpathlineto{\pgfqpoint{1.230182in}{0.780272in}}%
\pgfpathlineto{\pgfqpoint{1.234803in}{0.776406in}}%
\pgfpathlineto{\pgfqpoint{1.238806in}{0.772404in}}%
\pgfpathclose%
\pgfusepath{fill}%
\end{pgfscope}%
\begin{pgfscope}%
\pgfpathrectangle{\pgfqpoint{0.407592in}{0.222904in}}{\pgfqpoint{1.411163in}{1.411163in}}%
\pgfusepath{clip}%
\pgfsetbuttcap%
\pgfsetroundjoin%
\definecolor{currentfill}{rgb}{0.749020,0.772549,0.792157}%
\pgfsetfillcolor{currentfill}%
\pgfsetfillopacity{0.160000}%
\pgfsetlinewidth{0.000000pt}%
\definecolor{currentstroke}{rgb}{0.000000,0.000000,0.000000}%
\pgfsetstrokecolor{currentstroke}%
\pgfsetdash{}{0pt}%
\pgfpathmoveto{\pgfqpoint{1.090230in}{0.610533in}}%
\pgfpathlineto{\pgfqpoint{1.097270in}{0.626837in}}%
\pgfpathlineto{\pgfqpoint{1.104300in}{0.643117in}}%
\pgfpathlineto{\pgfqpoint{1.102802in}{0.643473in}}%
\pgfpathlineto{\pgfqpoint{1.101249in}{0.643784in}}%
\pgfpathlineto{\pgfqpoint{1.095744in}{0.627173in}}%
\pgfpathlineto{\pgfqpoint{1.090230in}{0.610533in}}%
\pgfpathlineto{\pgfqpoint{1.090230in}{0.610533in}}%
\pgfpathlineto{\pgfqpoint{1.090230in}{0.610533in}}%
\pgfpathclose%
\pgfusepath{fill}%
\end{pgfscope}%
\begin{pgfscope}%
\pgfpathrectangle{\pgfqpoint{0.407592in}{0.222904in}}{\pgfqpoint{1.411163in}{1.411163in}}%
\pgfusepath{clip}%
\pgfsetbuttcap%
\pgfsetroundjoin%
\definecolor{currentfill}{rgb}{0.749020,0.772549,0.792157}%
\pgfsetfillcolor{currentfill}%
\pgfsetfillopacity{0.160000}%
\pgfsetlinewidth{0.000000pt}%
\definecolor{currentstroke}{rgb}{0.000000,0.000000,0.000000}%
\pgfsetstrokecolor{currentstroke}%
\pgfsetdash{}{0pt}%
\pgfpathmoveto{\pgfqpoint{0.910644in}{0.835704in}}%
\pgfpathlineto{\pgfqpoint{0.899447in}{0.849743in}}%
\pgfpathlineto{\pgfqpoint{0.888254in}{0.863777in}}%
\pgfpathlineto{\pgfqpoint{0.880162in}{0.858082in}}%
\pgfpathlineto{\pgfqpoint{0.872948in}{0.852186in}}%
\pgfpathlineto{\pgfqpoint{0.885029in}{0.838749in}}%
\pgfpathlineto{\pgfqpoint{0.897109in}{0.825314in}}%
\pgfpathlineto{\pgfqpoint{0.903484in}{0.830598in}}%
\pgfpathlineto{\pgfqpoint{0.910644in}{0.835704in}}%
\pgfpathclose%
\pgfusepath{fill}%
\end{pgfscope}%
\begin{pgfscope}%
\pgfpathrectangle{\pgfqpoint{0.407592in}{0.222904in}}{\pgfqpoint{1.411163in}{1.411163in}}%
\pgfusepath{clip}%
\pgfsetbuttcap%
\pgfsetroundjoin%
\definecolor{currentfill}{rgb}{0.749020,0.772549,0.792157}%
\pgfsetfillcolor{currentfill}%
\pgfsetfillopacity{0.160000}%
\pgfsetlinewidth{0.000000pt}%
\definecolor{currentstroke}{rgb}{0.000000,0.000000,0.000000}%
\pgfsetstrokecolor{currentstroke}%
\pgfsetdash{}{0pt}%
\pgfpathmoveto{\pgfqpoint{1.090230in}{0.610533in}}%
\pgfpathlineto{\pgfqpoint{1.085581in}{0.627143in}}%
\pgfpathlineto{\pgfqpoint{1.080940in}{0.643726in}}%
\pgfpathlineto{\pgfqpoint{1.079332in}{0.643405in}}%
\pgfpathlineto{\pgfqpoint{1.077772in}{0.643040in}}%
\pgfpathlineto{\pgfqpoint{1.083996in}{0.626799in}}%
\pgfpathlineto{\pgfqpoint{1.090230in}{0.610533in}}%
\pgfpathlineto{\pgfqpoint{1.090230in}{0.610533in}}%
\pgfpathlineto{\pgfqpoint{1.090230in}{0.610533in}}%
\pgfpathclose%
\pgfusepath{fill}%
\end{pgfscope}%
\begin{pgfscope}%
\pgfpathrectangle{\pgfqpoint{0.407592in}{0.222904in}}{\pgfqpoint{1.411163in}{1.411163in}}%
\pgfusepath{clip}%
\pgfsetbuttcap%
\pgfsetroundjoin%
\definecolor{currentfill}{rgb}{0.749020,0.772549,0.792157}%
\pgfsetfillcolor{currentfill}%
\pgfsetfillopacity{0.160000}%
\pgfsetlinewidth{0.000000pt}%
\definecolor{currentstroke}{rgb}{0.000000,0.000000,0.000000}%
\pgfsetstrokecolor{currentstroke}%
\pgfsetdash{}{0pt}%
\pgfpathmoveto{\pgfqpoint{1.113593in}{0.638869in}}%
\pgfpathlineto{\pgfqpoint{1.125269in}{0.653030in}}%
\pgfpathlineto{\pgfqpoint{1.136941in}{0.667186in}}%
\pgfpathlineto{\pgfqpoint{1.135232in}{0.668468in}}%
\pgfpathlineto{\pgfqpoint{1.133328in}{0.669694in}}%
\pgfpathlineto{\pgfqpoint{1.122564in}{0.654919in}}%
\pgfpathlineto{\pgfqpoint{1.111794in}{0.640134in}}%
\pgfpathlineto{\pgfqpoint{1.112742in}{0.639515in}}%
\pgfpathlineto{\pgfqpoint{1.113593in}{0.638869in}}%
\pgfpathclose%
\pgfusepath{fill}%
\end{pgfscope}%
\begin{pgfscope}%
\pgfpathrectangle{\pgfqpoint{0.407592in}{0.222904in}}{\pgfqpoint{1.411163in}{1.411163in}}%
\pgfusepath{clip}%
\pgfsetbuttcap%
\pgfsetroundjoin%
\definecolor{currentfill}{rgb}{0.749020,0.772549,0.792157}%
\pgfsetfillcolor{currentfill}%
\pgfsetfillopacity{0.160000}%
\pgfsetlinewidth{0.000000pt}%
\definecolor{currentstroke}{rgb}{0.000000,0.000000,0.000000}%
\pgfsetstrokecolor{currentstroke}%
\pgfsetdash{}{0pt}%
\pgfpathmoveto{\pgfqpoint{1.214036in}{0.745418in}}%
\pgfpathlineto{\pgfqpoint{1.226421in}{0.758911in}}%
\pgfpathlineto{\pgfqpoint{1.238806in}{0.772404in}}%
\pgfpathlineto{\pgfqpoint{1.234803in}{0.776406in}}%
\pgfpathlineto{\pgfqpoint{1.230182in}{0.780272in}}%
\pgfpathlineto{\pgfqpoint{1.218540in}{0.766152in}}%
\pgfpathlineto{\pgfqpoint{1.206894in}{0.752028in}}%
\pgfpathlineto{\pgfqpoint{1.210724in}{0.748779in}}%
\pgfpathlineto{\pgfqpoint{1.214036in}{0.745418in}}%
\pgfpathclose%
\pgfusepath{fill}%
\end{pgfscope}%
\begin{pgfscope}%
\pgfpathrectangle{\pgfqpoint{0.407592in}{0.222904in}}{\pgfqpoint{1.411163in}{1.411163in}}%
\pgfusepath{clip}%
\pgfsetbuttcap%
\pgfsetroundjoin%
\definecolor{currentfill}{rgb}{0.749020,0.772549,0.792157}%
\pgfsetfillcolor{currentfill}%
\pgfsetfillopacity{0.160000}%
\pgfsetlinewidth{0.000000pt}%
\definecolor{currentstroke}{rgb}{0.000000,0.000000,0.000000}%
\pgfsetstrokecolor{currentstroke}%
\pgfsetdash{}{0pt}%
\pgfpathmoveto{\pgfqpoint{0.933047in}{0.807614in}}%
\pgfpathlineto{\pgfqpoint{0.921844in}{0.821661in}}%
\pgfpathlineto{\pgfqpoint{0.910644in}{0.835704in}}%
\pgfpathlineto{\pgfqpoint{0.903484in}{0.830598in}}%
\pgfpathlineto{\pgfqpoint{0.897109in}{0.825314in}}%
\pgfpathlineto{\pgfqpoint{0.909188in}{0.811880in}}%
\pgfpathlineto{\pgfqpoint{0.921266in}{0.798448in}}%
\pgfpathlineto{\pgfqpoint{0.926812in}{0.803108in}}%
\pgfpathlineto{\pgfqpoint{0.933047in}{0.807614in}}%
\pgfpathclose%
\pgfusepath{fill}%
\end{pgfscope}%
\begin{pgfscope}%
\pgfpathrectangle{\pgfqpoint{0.407592in}{0.222904in}}{\pgfqpoint{1.411163in}{1.411163in}}%
\pgfusepath{clip}%
\pgfsetbuttcap%
\pgfsetroundjoin%
\definecolor{currentfill}{rgb}{0.749020,0.772549,0.792157}%
\pgfsetfillcolor{currentfill}%
\pgfsetfillopacity{0.160000}%
\pgfsetlinewidth{0.000000pt}%
\definecolor{currentstroke}{rgb}{0.000000,0.000000,0.000000}%
\pgfsetstrokecolor{currentstroke}%
\pgfsetdash{}{0pt}%
\pgfpathmoveto{\pgfqpoint{1.090230in}{0.610533in}}%
\pgfpathlineto{\pgfqpoint{1.098673in}{0.626415in}}%
\pgfpathlineto{\pgfqpoint{1.107106in}{0.642276in}}%
\pgfpathlineto{\pgfqpoint{1.105737in}{0.642718in}}%
\pgfpathlineto{\pgfqpoint{1.104300in}{0.643117in}}%
\pgfpathlineto{\pgfqpoint{1.097270in}{0.626837in}}%
\pgfpathlineto{\pgfqpoint{1.090230in}{0.610533in}}%
\pgfpathlineto{\pgfqpoint{1.090230in}{0.610533in}}%
\pgfpathlineto{\pgfqpoint{1.090230in}{0.610533in}}%
\pgfpathclose%
\pgfusepath{fill}%
\end{pgfscope}%
\begin{pgfscope}%
\pgfpathrectangle{\pgfqpoint{0.407592in}{0.222904in}}{\pgfqpoint{1.411163in}{1.411163in}}%
\pgfusepath{clip}%
\pgfsetbuttcap%
\pgfsetroundjoin%
\definecolor{currentfill}{rgb}{0.749020,0.772549,0.792157}%
\pgfsetfillcolor{currentfill}%
\pgfsetfillopacity{0.160000}%
\pgfsetlinewidth{0.000000pt}%
\definecolor{currentstroke}{rgb}{0.000000,0.000000,0.000000}%
\pgfsetstrokecolor{currentstroke}%
\pgfsetdash{}{0pt}%
\pgfpathmoveto{\pgfqpoint{1.069756in}{0.640012in}}%
\pgfpathlineto{\pgfqpoint{1.059529in}{0.654737in}}%
\pgfpathlineto{\pgfqpoint{1.049309in}{0.669453in}}%
\pgfpathlineto{\pgfqpoint{1.047189in}{0.668215in}}%
\pgfpathlineto{\pgfqpoint{1.045256in}{0.666923in}}%
\pgfpathlineto{\pgfqpoint{1.056494in}{0.652831in}}%
\pgfpathlineto{\pgfqpoint{1.067736in}{0.638736in}}%
\pgfpathlineto{\pgfqpoint{1.068699in}{0.639388in}}%
\pgfpathlineto{\pgfqpoint{1.069756in}{0.640012in}}%
\pgfpathclose%
\pgfusepath{fill}%
\end{pgfscope}%
\begin{pgfscope}%
\pgfpathrectangle{\pgfqpoint{0.407592in}{0.222904in}}{\pgfqpoint{1.411163in}{1.411163in}}%
\pgfusepath{clip}%
\pgfsetbuttcap%
\pgfsetroundjoin%
\definecolor{currentfill}{rgb}{0.749020,0.772549,0.792157}%
\pgfsetfillcolor{currentfill}%
\pgfsetfillopacity{0.160000}%
\pgfsetlinewidth{0.000000pt}%
\definecolor{currentstroke}{rgb}{0.000000,0.000000,0.000000}%
\pgfsetstrokecolor{currentstroke}%
\pgfsetdash{}{0pt}%
\pgfpathmoveto{\pgfqpoint{1.090230in}{0.610533in}}%
\pgfpathlineto{\pgfqpoint{1.083996in}{0.626799in}}%
\pgfpathlineto{\pgfqpoint{1.077772in}{0.643040in}}%
\pgfpathlineto{\pgfqpoint{1.076265in}{0.642632in}}%
\pgfpathlineto{\pgfqpoint{1.074818in}{0.642183in}}%
\pgfpathlineto{\pgfqpoint{1.082519in}{0.626368in}}%
\pgfpathlineto{\pgfqpoint{1.090230in}{0.610533in}}%
\pgfpathlineto{\pgfqpoint{1.090230in}{0.610533in}}%
\pgfpathlineto{\pgfqpoint{1.090230in}{0.610533in}}%
\pgfpathclose%
\pgfusepath{fill}%
\end{pgfscope}%
\begin{pgfscope}%
\pgfpathrectangle{\pgfqpoint{0.407592in}{0.222904in}}{\pgfqpoint{1.411163in}{1.411163in}}%
\pgfusepath{clip}%
\pgfsetbuttcap%
\pgfsetroundjoin%
\definecolor{currentfill}{rgb}{0.749020,0.772549,0.792157}%
\pgfsetfillcolor{currentfill}%
\pgfsetfillopacity{0.160000}%
\pgfsetlinewidth{0.000000pt}%
\definecolor{currentstroke}{rgb}{0.000000,0.000000,0.000000}%
\pgfsetstrokecolor{currentstroke}%
\pgfsetdash{}{0pt}%
\pgfpathmoveto{\pgfqpoint{0.955463in}{0.779508in}}%
\pgfpathlineto{\pgfqpoint{0.944253in}{0.793563in}}%
\pgfpathlineto{\pgfqpoint{0.933047in}{0.807614in}}%
\pgfpathlineto{\pgfqpoint{0.926812in}{0.803108in}}%
\pgfpathlineto{\pgfqpoint{0.921266in}{0.798448in}}%
\pgfpathlineto{\pgfqpoint{0.933343in}{0.785017in}}%
\pgfpathlineto{\pgfqpoint{0.945418in}{0.771587in}}%
\pgfpathlineto{\pgfqpoint{0.950143in}{0.775614in}}%
\pgfpathlineto{\pgfqpoint{0.955463in}{0.779508in}}%
\pgfpathclose%
\pgfusepath{fill}%
\end{pgfscope}%
\begin{pgfscope}%
\pgfpathrectangle{\pgfqpoint{0.407592in}{0.222904in}}{\pgfqpoint{1.411163in}{1.411163in}}%
\pgfusepath{clip}%
\pgfsetbuttcap%
\pgfsetroundjoin%
\definecolor{currentfill}{rgb}{0.749020,0.772549,0.792157}%
\pgfsetfillcolor{currentfill}%
\pgfsetfillopacity{0.160000}%
\pgfsetlinewidth{0.000000pt}%
\definecolor{currentstroke}{rgb}{0.000000,0.000000,0.000000}%
\pgfsetstrokecolor{currentstroke}%
\pgfsetdash{}{0pt}%
\pgfpathmoveto{\pgfqpoint{1.189269in}{0.718434in}}%
\pgfpathlineto{\pgfqpoint{1.201652in}{0.731926in}}%
\pgfpathlineto{\pgfqpoint{1.214036in}{0.745418in}}%
\pgfpathlineto{\pgfqpoint{1.210724in}{0.748779in}}%
\pgfpathlineto{\pgfqpoint{1.206894in}{0.752028in}}%
\pgfpathlineto{\pgfqpoint{1.195245in}{0.737899in}}%
\pgfpathlineto{\pgfqpoint{1.183592in}{0.723766in}}%
\pgfpathlineto{\pgfqpoint{1.186638in}{0.721145in}}%
\pgfpathlineto{\pgfqpoint{1.189269in}{0.718434in}}%
\pgfpathclose%
\pgfusepath{fill}%
\end{pgfscope}%
\begin{pgfscope}%
\pgfpathrectangle{\pgfqpoint{0.407592in}{0.222904in}}{\pgfqpoint{1.411163in}{1.411163in}}%
\pgfusepath{clip}%
\pgfsetbuttcap%
\pgfsetroundjoin%
\definecolor{currentfill}{rgb}{0.749020,0.772549,0.792157}%
\pgfsetfillcolor{currentfill}%
\pgfsetfillopacity{0.160000}%
\pgfsetlinewidth{0.000000pt}%
\definecolor{currentstroke}{rgb}{0.000000,0.000000,0.000000}%
\pgfsetstrokecolor{currentstroke}%
\pgfsetdash{}{0pt}%
\pgfpathmoveto{\pgfqpoint{1.090230in}{0.610533in}}%
\pgfpathlineto{\pgfqpoint{1.099929in}{0.625912in}}%
\pgfpathlineto{\pgfqpoint{1.109619in}{0.641276in}}%
\pgfpathlineto{\pgfqpoint{1.108402in}{0.641795in}}%
\pgfpathlineto{\pgfqpoint{1.107106in}{0.642276in}}%
\pgfpathlineto{\pgfqpoint{1.098673in}{0.626415in}}%
\pgfpathlineto{\pgfqpoint{1.090230in}{0.610533in}}%
\pgfpathlineto{\pgfqpoint{1.090230in}{0.610533in}}%
\pgfpathlineto{\pgfqpoint{1.090230in}{0.610533in}}%
\pgfpathclose%
\pgfusepath{fill}%
\end{pgfscope}%
\begin{pgfscope}%
\pgfpathrectangle{\pgfqpoint{0.407592in}{0.222904in}}{\pgfqpoint{1.411163in}{1.411163in}}%
\pgfusepath{clip}%
\pgfsetbuttcap%
\pgfsetroundjoin%
\definecolor{currentfill}{rgb}{0.749020,0.772549,0.792157}%
\pgfsetfillcolor{currentfill}%
\pgfsetfillopacity{0.160000}%
\pgfsetlinewidth{0.000000pt}%
\definecolor{currentstroke}{rgb}{0.000000,0.000000,0.000000}%
\pgfsetstrokecolor{currentstroke}%
\pgfsetdash{}{0pt}%
\pgfpathmoveto{\pgfqpoint{1.090230in}{0.610533in}}%
\pgfpathlineto{\pgfqpoint{1.082519in}{0.626368in}}%
\pgfpathlineto{\pgfqpoint{1.074818in}{0.642183in}}%
\pgfpathlineto{\pgfqpoint{1.073438in}{0.641694in}}%
\pgfpathlineto{\pgfqpoint{1.072130in}{0.641167in}}%
\pgfpathlineto{\pgfqpoint{1.081176in}{0.625858in}}%
\pgfpathlineto{\pgfqpoint{1.090230in}{0.610533in}}%
\pgfpathlineto{\pgfqpoint{1.090230in}{0.610533in}}%
\pgfpathlineto{\pgfqpoint{1.090230in}{0.610533in}}%
\pgfpathclose%
\pgfusepath{fill}%
\end{pgfscope}%
\begin{pgfscope}%
\pgfpathrectangle{\pgfqpoint{0.407592in}{0.222904in}}{\pgfqpoint{1.411163in}{1.411163in}}%
\pgfusepath{clip}%
\pgfsetbuttcap%
\pgfsetroundjoin%
\definecolor{currentfill}{rgb}{0.749020,0.772549,0.792157}%
\pgfsetfillcolor{currentfill}%
\pgfsetfillopacity{0.160000}%
\pgfsetlinewidth{0.000000pt}%
\definecolor{currentstroke}{rgb}{0.000000,0.000000,0.000000}%
\pgfsetstrokecolor{currentstroke}%
\pgfsetdash{}{0pt}%
\pgfpathmoveto{\pgfqpoint{0.977892in}{0.751386in}}%
\pgfpathlineto{\pgfqpoint{0.966676in}{0.765449in}}%
\pgfpathlineto{\pgfqpoint{0.955463in}{0.779508in}}%
\pgfpathlineto{\pgfqpoint{0.950143in}{0.775614in}}%
\pgfpathlineto{\pgfqpoint{0.945418in}{0.771587in}}%
\pgfpathlineto{\pgfqpoint{0.957492in}{0.758158in}}%
\pgfpathlineto{\pgfqpoint{0.969565in}{0.744731in}}%
\pgfpathlineto{\pgfqpoint{0.973480in}{0.748114in}}%
\pgfpathlineto{\pgfqpoint{0.977892in}{0.751386in}}%
\pgfpathclose%
\pgfusepath{fill}%
\end{pgfscope}%
\begin{pgfscope}%
\pgfpathrectangle{\pgfqpoint{0.407592in}{0.222904in}}{\pgfqpoint{1.411163in}{1.411163in}}%
\pgfusepath{clip}%
\pgfsetbuttcap%
\pgfsetroundjoin%
\definecolor{currentfill}{rgb}{0.749020,0.772549,0.792157}%
\pgfsetfillcolor{currentfill}%
\pgfsetfillopacity{0.160000}%
\pgfsetlinewidth{0.000000pt}%
\definecolor{currentstroke}{rgb}{0.000000,0.000000,0.000000}%
\pgfsetstrokecolor{currentstroke}%
\pgfsetdash{}{0pt}%
\pgfpathmoveto{\pgfqpoint{1.164505in}{0.691454in}}%
\pgfpathlineto{\pgfqpoint{1.176886in}{0.704944in}}%
\pgfpathlineto{\pgfqpoint{1.189269in}{0.718434in}}%
\pgfpathlineto{\pgfqpoint{1.186638in}{0.721145in}}%
\pgfpathlineto{\pgfqpoint{1.183592in}{0.723766in}}%
\pgfpathlineto{\pgfqpoint{1.171935in}{0.709628in}}%
\pgfpathlineto{\pgfqpoint{1.160274in}{0.695485in}}%
\pgfpathlineto{\pgfqpoint{1.162546in}{0.693503in}}%
\pgfpathlineto{\pgfqpoint{1.164505in}{0.691454in}}%
\pgfpathclose%
\pgfusepath{fill}%
\end{pgfscope}%
\begin{pgfscope}%
\pgfpathrectangle{\pgfqpoint{0.407592in}{0.222904in}}{\pgfqpoint{1.411163in}{1.411163in}}%
\pgfusepath{clip}%
\pgfsetbuttcap%
\pgfsetroundjoin%
\definecolor{currentfill}{rgb}{0.749020,0.772549,0.792157}%
\pgfsetfillcolor{currentfill}%
\pgfsetfillopacity{0.160000}%
\pgfsetlinewidth{0.000000pt}%
\definecolor{currentstroke}{rgb}{0.000000,0.000000,0.000000}%
\pgfsetstrokecolor{currentstroke}%
\pgfsetdash{}{0pt}%
\pgfpathmoveto{\pgfqpoint{1.000333in}{0.723248in}}%
\pgfpathlineto{\pgfqpoint{0.989111in}{0.737319in}}%
\pgfpathlineto{\pgfqpoint{0.977892in}{0.751386in}}%
\pgfpathlineto{\pgfqpoint{0.973480in}{0.748114in}}%
\pgfpathlineto{\pgfqpoint{0.969565in}{0.744731in}}%
\pgfpathlineto{\pgfqpoint{0.981637in}{0.731305in}}%
\pgfpathlineto{\pgfqpoint{0.993708in}{0.717881in}}%
\pgfpathlineto{\pgfqpoint{0.996821in}{0.720608in}}%
\pgfpathlineto{\pgfqpoint{1.000333in}{0.723248in}}%
\pgfpathclose%
\pgfusepath{fill}%
\end{pgfscope}%
\begin{pgfscope}%
\pgfpathrectangle{\pgfqpoint{0.407592in}{0.222904in}}{\pgfqpoint{1.411163in}{1.411163in}}%
\pgfusepath{clip}%
\pgfsetbuttcap%
\pgfsetroundjoin%
\definecolor{currentfill}{rgb}{0.749020,0.772549,0.792157}%
\pgfsetfillcolor{currentfill}%
\pgfsetfillopacity{0.160000}%
\pgfsetlinewidth{0.000000pt}%
\definecolor{currentstroke}{rgb}{0.000000,0.000000,0.000000}%
\pgfsetstrokecolor{currentstroke}%
\pgfsetdash{}{0pt}%
\pgfpathmoveto{\pgfqpoint{1.090230in}{0.610533in}}%
\pgfpathlineto{\pgfqpoint{1.101015in}{0.625338in}}%
\pgfpathlineto{\pgfqpoint{1.111794in}{0.640134in}}%
\pgfpathlineto{\pgfqpoint{1.110751in}{0.640721in}}%
\pgfpathlineto{\pgfqpoint{1.109619in}{0.641276in}}%
\pgfpathlineto{\pgfqpoint{1.099929in}{0.625912in}}%
\pgfpathlineto{\pgfqpoint{1.090230in}{0.610533in}}%
\pgfpathlineto{\pgfqpoint{1.090230in}{0.610533in}}%
\pgfpathlineto{\pgfqpoint{1.090230in}{0.610533in}}%
\pgfpathclose%
\pgfusepath{fill}%
\end{pgfscope}%
\begin{pgfscope}%
\pgfpathrectangle{\pgfqpoint{0.407592in}{0.222904in}}{\pgfqpoint{1.411163in}{1.411163in}}%
\pgfusepath{clip}%
\pgfsetbuttcap%
\pgfsetroundjoin%
\definecolor{currentfill}{rgb}{0.749020,0.772549,0.792157}%
\pgfsetfillcolor{currentfill}%
\pgfsetfillopacity{0.160000}%
\pgfsetlinewidth{0.000000pt}%
\definecolor{currentstroke}{rgb}{0.000000,0.000000,0.000000}%
\pgfsetstrokecolor{currentstroke}%
\pgfsetdash{}{0pt}%
\pgfpathmoveto{\pgfqpoint{1.090230in}{0.610533in}}%
\pgfpathlineto{\pgfqpoint{1.081176in}{0.625858in}}%
\pgfpathlineto{\pgfqpoint{1.072130in}{0.641167in}}%
\pgfpathlineto{\pgfqpoint{1.070901in}{0.640606in}}%
\pgfpathlineto{\pgfqpoint{1.069756in}{0.640012in}}%
\pgfpathlineto{\pgfqpoint{1.079989in}{0.625277in}}%
\pgfpathlineto{\pgfqpoint{1.090230in}{0.610533in}}%
\pgfpathlineto{\pgfqpoint{1.090230in}{0.610533in}}%
\pgfpathlineto{\pgfqpoint{1.090230in}{0.610533in}}%
\pgfpathclose%
\pgfusepath{fill}%
\end{pgfscope}%
\begin{pgfscope}%
\pgfpathrectangle{\pgfqpoint{0.407592in}{0.222904in}}{\pgfqpoint{1.411163in}{1.411163in}}%
\pgfusepath{clip}%
\pgfsetbuttcap%
\pgfsetroundjoin%
\definecolor{currentfill}{rgb}{0.749020,0.772549,0.792157}%
\pgfsetfillcolor{currentfill}%
\pgfsetfillopacity{0.160000}%
\pgfsetlinewidth{0.000000pt}%
\definecolor{currentstroke}{rgb}{0.000000,0.000000,0.000000}%
\pgfsetstrokecolor{currentstroke}%
\pgfsetdash{}{0pt}%
\pgfpathmoveto{\pgfqpoint{1.139743in}{0.664477in}}%
\pgfpathlineto{\pgfqpoint{1.152124in}{0.677965in}}%
\pgfpathlineto{\pgfqpoint{1.164505in}{0.691454in}}%
\pgfpathlineto{\pgfqpoint{1.162546in}{0.693503in}}%
\pgfpathlineto{\pgfqpoint{1.160274in}{0.695485in}}%
\pgfpathlineto{\pgfqpoint{1.148609in}{0.681338in}}%
\pgfpathlineto{\pgfqpoint{1.136941in}{0.667186in}}%
\pgfpathlineto{\pgfqpoint{1.138447in}{0.665854in}}%
\pgfpathlineto{\pgfqpoint{1.139743in}{0.664477in}}%
\pgfpathclose%
\pgfusepath{fill}%
\end{pgfscope}%
\begin{pgfscope}%
\pgfpathrectangle{\pgfqpoint{0.407592in}{0.222904in}}{\pgfqpoint{1.411163in}{1.411163in}}%
\pgfusepath{clip}%
\pgfsetbuttcap%
\pgfsetroundjoin%
\definecolor{currentfill}{rgb}{0.749020,0.772549,0.792157}%
\pgfsetfillcolor{currentfill}%
\pgfsetfillopacity{0.160000}%
\pgfsetlinewidth{0.000000pt}%
\definecolor{currentstroke}{rgb}{0.000000,0.000000,0.000000}%
\pgfsetstrokecolor{currentstroke}%
\pgfsetdash{}{0pt}%
\pgfpathmoveto{\pgfqpoint{1.022788in}{0.695094in}}%
\pgfpathlineto{\pgfqpoint{1.011559in}{0.709173in}}%
\pgfpathlineto{\pgfqpoint{1.000333in}{0.723248in}}%
\pgfpathlineto{\pgfqpoint{0.996821in}{0.720608in}}%
\pgfpathlineto{\pgfqpoint{0.993708in}{0.717881in}}%
\pgfpathlineto{\pgfqpoint{1.005777in}{0.704458in}}%
\pgfpathlineto{\pgfqpoint{1.017845in}{0.691036in}}%
\pgfpathlineto{\pgfqpoint{1.020166in}{0.693097in}}%
\pgfpathlineto{\pgfqpoint{1.022788in}{0.695094in}}%
\pgfpathclose%
\pgfusepath{fill}%
\end{pgfscope}%
\begin{pgfscope}%
\pgfpathrectangle{\pgfqpoint{0.407592in}{0.222904in}}{\pgfqpoint{1.411163in}{1.411163in}}%
\pgfusepath{clip}%
\pgfsetbuttcap%
\pgfsetroundjoin%
\definecolor{currentfill}{rgb}{0.749020,0.772549,0.792157}%
\pgfsetfillcolor{currentfill}%
\pgfsetfillopacity{0.160000}%
\pgfsetlinewidth{0.000000pt}%
\definecolor{currentstroke}{rgb}{0.000000,0.000000,0.000000}%
\pgfsetstrokecolor{currentstroke}%
\pgfsetdash{}{0pt}%
\pgfpathmoveto{\pgfqpoint{1.090230in}{0.610533in}}%
\pgfpathlineto{\pgfqpoint{1.101913in}{0.624703in}}%
\pgfpathlineto{\pgfqpoint{1.113593in}{0.638869in}}%
\pgfpathlineto{\pgfqpoint{1.112742in}{0.639515in}}%
\pgfpathlineto{\pgfqpoint{1.111794in}{0.640134in}}%
\pgfpathlineto{\pgfqpoint{1.101015in}{0.625338in}}%
\pgfpathlineto{\pgfqpoint{1.090230in}{0.610533in}}%
\pgfpathlineto{\pgfqpoint{1.090230in}{0.610533in}}%
\pgfpathlineto{\pgfqpoint{1.090230in}{0.610533in}}%
\pgfpathclose%
\pgfusepath{fill}%
\end{pgfscope}%
\begin{pgfscope}%
\pgfpathrectangle{\pgfqpoint{0.407592in}{0.222904in}}{\pgfqpoint{1.411163in}{1.411163in}}%
\pgfusepath{clip}%
\pgfsetbuttcap%
\pgfsetroundjoin%
\definecolor{currentfill}{rgb}{0.749020,0.772549,0.792157}%
\pgfsetfillcolor{currentfill}%
\pgfsetfillopacity{0.160000}%
\pgfsetlinewidth{0.000000pt}%
\definecolor{currentstroke}{rgb}{0.000000,0.000000,0.000000}%
\pgfsetstrokecolor{currentstroke}%
\pgfsetdash{}{0pt}%
\pgfpathmoveto{\pgfqpoint{1.045256in}{0.666923in}}%
\pgfpathlineto{\pgfqpoint{1.034020in}{0.681010in}}%
\pgfpathlineto{\pgfqpoint{1.022788in}{0.695094in}}%
\pgfpathlineto{\pgfqpoint{1.020166in}{0.693097in}}%
\pgfpathlineto{\pgfqpoint{1.017845in}{0.691036in}}%
\pgfpathlineto{\pgfqpoint{1.029912in}{0.677615in}}%
\pgfpathlineto{\pgfqpoint{1.041978in}{0.664196in}}%
\pgfpathlineto{\pgfqpoint{1.043516in}{0.665581in}}%
\pgfpathlineto{\pgfqpoint{1.045256in}{0.666923in}}%
\pgfpathclose%
\pgfusepath{fill}%
\end{pgfscope}%
\begin{pgfscope}%
\pgfpathrectangle{\pgfqpoint{0.407592in}{0.222904in}}{\pgfqpoint{1.411163in}{1.411163in}}%
\pgfusepath{clip}%
\pgfsetbuttcap%
\pgfsetroundjoin%
\definecolor{currentfill}{rgb}{0.749020,0.772549,0.792157}%
\pgfsetfillcolor{currentfill}%
\pgfsetfillopacity{0.160000}%
\pgfsetlinewidth{0.000000pt}%
\definecolor{currentstroke}{rgb}{0.000000,0.000000,0.000000}%
\pgfsetstrokecolor{currentstroke}%
\pgfsetdash{}{0pt}%
\pgfpathmoveto{\pgfqpoint{1.090230in}{0.610533in}}%
\pgfpathlineto{\pgfqpoint{1.079989in}{0.625277in}}%
\pgfpathlineto{\pgfqpoint{1.069756in}{0.640012in}}%
\pgfpathlineto{\pgfqpoint{1.068699in}{0.639388in}}%
\pgfpathlineto{\pgfqpoint{1.067736in}{0.638736in}}%
\pgfpathlineto{\pgfqpoint{1.078981in}{0.624636in}}%
\pgfpathlineto{\pgfqpoint{1.090230in}{0.610533in}}%
\pgfpathlineto{\pgfqpoint{1.090230in}{0.610533in}}%
\pgfpathlineto{\pgfqpoint{1.090230in}{0.610533in}}%
\pgfpathclose%
\pgfusepath{fill}%
\end{pgfscope}%
\begin{pgfscope}%
\pgfpathrectangle{\pgfqpoint{0.407592in}{0.222904in}}{\pgfqpoint{1.411163in}{1.411163in}}%
\pgfusepath{clip}%
\pgfsetbuttcap%
\pgfsetroundjoin%
\definecolor{currentfill}{rgb}{0.749020,0.772549,0.792157}%
\pgfsetfillcolor{currentfill}%
\pgfsetfillopacity{0.160000}%
\pgfsetlinewidth{0.000000pt}%
\definecolor{currentstroke}{rgb}{0.000000,0.000000,0.000000}%
\pgfsetstrokecolor{currentstroke}%
\pgfsetdash{}{0pt}%
\pgfpathmoveto{\pgfqpoint{1.114985in}{0.637503in}}%
\pgfpathlineto{\pgfqpoint{1.127364in}{0.650990in}}%
\pgfpathlineto{\pgfqpoint{1.139743in}{0.664477in}}%
\pgfpathlineto{\pgfqpoint{1.138447in}{0.665854in}}%
\pgfpathlineto{\pgfqpoint{1.136941in}{0.667186in}}%
\pgfpathlineto{\pgfqpoint{1.125269in}{0.653030in}}%
\pgfpathlineto{\pgfqpoint{1.113593in}{0.638869in}}%
\pgfpathlineto{\pgfqpoint{1.114342in}{0.638197in}}%
\pgfpathlineto{\pgfqpoint{1.114985in}{0.637503in}}%
\pgfpathclose%
\pgfusepath{fill}%
\end{pgfscope}%
\begin{pgfscope}%
\pgfpathrectangle{\pgfqpoint{0.407592in}{0.222904in}}{\pgfqpoint{1.411163in}{1.411163in}}%
\pgfusepath{clip}%
\pgfsetbuttcap%
\pgfsetroundjoin%
\definecolor{currentfill}{rgb}{0.749020,0.772549,0.792157}%
\pgfsetfillcolor{currentfill}%
\pgfsetfillopacity{0.160000}%
\pgfsetlinewidth{0.000000pt}%
\definecolor{currentstroke}{rgb}{0.000000,0.000000,0.000000}%
\pgfsetstrokecolor{currentstroke}%
\pgfsetdash{}{0pt}%
\pgfpathmoveto{\pgfqpoint{1.067736in}{0.638736in}}%
\pgfpathlineto{\pgfqpoint{1.056494in}{0.652831in}}%
\pgfpathlineto{\pgfqpoint{1.045256in}{0.666923in}}%
\pgfpathlineto{\pgfqpoint{1.043516in}{0.665581in}}%
\pgfpathlineto{\pgfqpoint{1.041978in}{0.664196in}}%
\pgfpathlineto{\pgfqpoint{1.054043in}{0.650778in}}%
\pgfpathlineto{\pgfqpoint{1.066106in}{0.637362in}}%
\pgfpathlineto{\pgfqpoint{1.066870in}{0.638060in}}%
\pgfpathlineto{\pgfqpoint{1.067736in}{0.638736in}}%
\pgfpathclose%
\pgfusepath{fill}%
\end{pgfscope}%
\begin{pgfscope}%
\pgfpathrectangle{\pgfqpoint{0.407592in}{0.222904in}}{\pgfqpoint{1.411163in}{1.411163in}}%
\pgfusepath{clip}%
\pgfsetbuttcap%
\pgfsetroundjoin%
\definecolor{currentfill}{rgb}{0.749020,0.772549,0.792157}%
\pgfsetfillcolor{currentfill}%
\pgfsetfillopacity{0.160000}%
\pgfsetlinewidth{0.000000pt}%
\definecolor{currentstroke}{rgb}{0.000000,0.000000,0.000000}%
\pgfsetstrokecolor{currentstroke}%
\pgfsetdash{}{0pt}%
\pgfpathmoveto{\pgfqpoint{1.090230in}{0.610533in}}%
\pgfpathlineto{\pgfqpoint{1.102607in}{0.624018in}}%
\pgfpathlineto{\pgfqpoint{1.114985in}{0.637503in}}%
\pgfpathlineto{\pgfqpoint{1.114342in}{0.638197in}}%
\pgfpathlineto{\pgfqpoint{1.113593in}{0.638869in}}%
\pgfpathlineto{\pgfqpoint{1.101913in}{0.624703in}}%
\pgfpathlineto{\pgfqpoint{1.090230in}{0.610533in}}%
\pgfpathlineto{\pgfqpoint{1.090230in}{0.610533in}}%
\pgfpathlineto{\pgfqpoint{1.090230in}{0.610533in}}%
\pgfpathclose%
\pgfusepath{fill}%
\end{pgfscope}%
\begin{pgfscope}%
\pgfpathrectangle{\pgfqpoint{0.407592in}{0.222904in}}{\pgfqpoint{1.411163in}{1.411163in}}%
\pgfusepath{clip}%
\pgfsetbuttcap%
\pgfsetroundjoin%
\definecolor{currentfill}{rgb}{0.749020,0.772549,0.792157}%
\pgfsetfillcolor{currentfill}%
\pgfsetfillopacity{0.160000}%
\pgfsetlinewidth{0.000000pt}%
\definecolor{currentstroke}{rgb}{0.000000,0.000000,0.000000}%
\pgfsetstrokecolor{currentstroke}%
\pgfsetdash{}{0pt}%
\pgfpathmoveto{\pgfqpoint{1.090230in}{0.610533in}}%
\pgfpathlineto{\pgfqpoint{1.078981in}{0.624636in}}%
\pgfpathlineto{\pgfqpoint{1.067736in}{0.638736in}}%
\pgfpathlineto{\pgfqpoint{1.066870in}{0.638060in}}%
\pgfpathlineto{\pgfqpoint{1.066106in}{0.637362in}}%
\pgfpathlineto{\pgfqpoint{1.078169in}{0.623947in}}%
\pgfpathlineto{\pgfqpoint{1.090230in}{0.610533in}}%
\pgfpathlineto{\pgfqpoint{1.090230in}{0.610533in}}%
\pgfpathlineto{\pgfqpoint{1.090230in}{0.610533in}}%
\pgfpathclose%
\pgfusepath{fill}%
\end{pgfscope}%
\begin{pgfscope}%
\pgfpathrectangle{\pgfqpoint{0.407592in}{0.222904in}}{\pgfqpoint{1.411163in}{1.411163in}}%
\pgfusepath{clip}%
\pgfsetbuttcap%
\pgfsetroundjoin%
\definecolor{currentfill}{rgb}{0.749020,0.772549,0.792157}%
\pgfsetfillcolor{currentfill}%
\pgfsetfillopacity{0.160000}%
\pgfsetlinewidth{0.000000pt}%
\definecolor{currentstroke}{rgb}{0.000000,0.000000,0.000000}%
\pgfsetstrokecolor{currentstroke}%
\pgfsetdash{}{0pt}%
\pgfpathmoveto{\pgfqpoint{1.090230in}{0.610533in}}%
\pgfpathlineto{\pgfqpoint{1.103084in}{0.623294in}}%
\pgfpathlineto{\pgfqpoint{1.115945in}{0.636061in}}%
\pgfpathlineto{\pgfqpoint{1.115520in}{0.636790in}}%
\pgfpathlineto{\pgfqpoint{1.114985in}{0.637503in}}%
\pgfpathlineto{\pgfqpoint{1.102607in}{0.624018in}}%
\pgfpathlineto{\pgfqpoint{1.090230in}{0.610533in}}%
\pgfpathlineto{\pgfqpoint{1.090230in}{0.610533in}}%
\pgfpathlineto{\pgfqpoint{1.090230in}{0.610533in}}%
\pgfpathclose%
\pgfusepath{fill}%
\end{pgfscope}%
\begin{pgfscope}%
\pgfpathrectangle{\pgfqpoint{0.407592in}{0.222904in}}{\pgfqpoint{1.411163in}{1.411163in}}%
\pgfusepath{clip}%
\pgfsetbuttcap%
\pgfsetroundjoin%
\definecolor{currentfill}{rgb}{0.749020,0.772549,0.792157}%
\pgfsetfillcolor{currentfill}%
\pgfsetfillopacity{0.160000}%
\pgfsetlinewidth{0.000000pt}%
\definecolor{currentstroke}{rgb}{0.000000,0.000000,0.000000}%
\pgfsetstrokecolor{currentstroke}%
\pgfsetdash{}{0pt}%
\pgfpathmoveto{\pgfqpoint{1.090230in}{0.610533in}}%
\pgfpathlineto{\pgfqpoint{1.078169in}{0.623947in}}%
\pgfpathlineto{\pgfqpoint{1.066106in}{0.637362in}}%
\pgfpathlineto{\pgfqpoint{1.065447in}{0.636645in}}%
\pgfpathlineto{\pgfqpoint{1.064895in}{0.635913in}}%
\pgfpathlineto{\pgfqpoint{1.077566in}{0.623220in}}%
\pgfpathlineto{\pgfqpoint{1.090230in}{0.610533in}}%
\pgfpathlineto{\pgfqpoint{1.090230in}{0.610533in}}%
\pgfpathlineto{\pgfqpoint{1.090230in}{0.610533in}}%
\pgfpathclose%
\pgfusepath{fill}%
\end{pgfscope}%
\begin{pgfscope}%
\pgfpathrectangle{\pgfqpoint{0.407592in}{0.222904in}}{\pgfqpoint{1.411163in}{1.411163in}}%
\pgfusepath{clip}%
\pgfsetbuttcap%
\pgfsetroundjoin%
\definecolor{currentfill}{rgb}{0.749020,0.772549,0.792157}%
\pgfsetfillcolor{currentfill}%
\pgfsetfillopacity{0.160000}%
\pgfsetlinewidth{0.000000pt}%
\definecolor{currentstroke}{rgb}{0.000000,0.000000,0.000000}%
\pgfsetstrokecolor{currentstroke}%
\pgfsetdash{}{0pt}%
\pgfpathmoveto{\pgfqpoint{1.090230in}{0.610533in}}%
\pgfpathlineto{\pgfqpoint{1.103337in}{0.622545in}}%
\pgfpathlineto{\pgfqpoint{1.116456in}{0.634567in}}%
\pgfpathlineto{\pgfqpoint{1.116257in}{0.635319in}}%
\pgfpathlineto{\pgfqpoint{1.115945in}{0.636061in}}%
\pgfpathlineto{\pgfqpoint{1.103084in}{0.623294in}}%
\pgfpathlineto{\pgfqpoint{1.090230in}{0.610533in}}%
\pgfpathlineto{\pgfqpoint{1.090230in}{0.610533in}}%
\pgfpathlineto{\pgfqpoint{1.090230in}{0.610533in}}%
\pgfpathclose%
\pgfusepath{fill}%
\end{pgfscope}%
\begin{pgfscope}%
\pgfpathrectangle{\pgfqpoint{0.407592in}{0.222904in}}{\pgfqpoint{1.411163in}{1.411163in}}%
\pgfusepath{clip}%
\pgfsetbuttcap%
\pgfsetroundjoin%
\definecolor{currentfill}{rgb}{0.749020,0.772549,0.792157}%
\pgfsetfillcolor{currentfill}%
\pgfsetfillopacity{0.160000}%
\pgfsetlinewidth{0.000000pt}%
\definecolor{currentstroke}{rgb}{0.000000,0.000000,0.000000}%
\pgfsetstrokecolor{currentstroke}%
\pgfsetdash{}{0pt}%
\pgfpathmoveto{\pgfqpoint{1.090230in}{0.610533in}}%
\pgfpathlineto{\pgfqpoint{1.077566in}{0.623220in}}%
\pgfpathlineto{\pgfqpoint{1.064895in}{0.635913in}}%
\pgfpathlineto{\pgfqpoint{1.064454in}{0.635169in}}%
\pgfpathlineto{\pgfqpoint{1.064125in}{0.634416in}}%
\pgfpathlineto{\pgfqpoint{1.077183in}{0.622469in}}%
\pgfpathlineto{\pgfqpoint{1.090230in}{0.610533in}}%
\pgfpathlineto{\pgfqpoint{1.090230in}{0.610533in}}%
\pgfpathlineto{\pgfqpoint{1.090230in}{0.610533in}}%
\pgfpathclose%
\pgfusepath{fill}%
\end{pgfscope}%
\begin{pgfscope}%
\pgfpathrectangle{\pgfqpoint{0.407592in}{0.222904in}}{\pgfqpoint{1.411163in}{1.411163in}}%
\pgfusepath{clip}%
\pgfsetbuttcap%
\pgfsetroundjoin%
\definecolor{currentfill}{rgb}{0.749020,0.772549,0.792157}%
\pgfsetfillcolor{currentfill}%
\pgfsetfillopacity{0.160000}%
\pgfsetlinewidth{0.000000pt}%
\definecolor{currentstroke}{rgb}{0.000000,0.000000,0.000000}%
\pgfsetstrokecolor{currentstroke}%
\pgfsetdash{}{0pt}%
\pgfpathmoveto{\pgfqpoint{1.115945in}{0.636061in}}%
\pgfpathlineto{\pgfqpoint{1.128812in}{0.648834in}}%
\pgfpathlineto{\pgfqpoint{1.141684in}{0.661613in}}%
\pgfpathlineto{\pgfqpoint{1.140824in}{0.663062in}}%
\pgfpathlineto{\pgfqpoint{1.139743in}{0.664477in}}%
\pgfpathlineto{\pgfqpoint{1.127364in}{0.650990in}}%
\pgfpathlineto{\pgfqpoint{1.114985in}{0.637503in}}%
\pgfpathlineto{\pgfqpoint{1.115520in}{0.636790in}}%
\pgfpathlineto{\pgfqpoint{1.115945in}{0.636061in}}%
\pgfpathclose%
\pgfusepath{fill}%
\end{pgfscope}%
\begin{pgfscope}%
\pgfpathrectangle{\pgfqpoint{0.407592in}{0.222904in}}{\pgfqpoint{1.411163in}{1.411163in}}%
\pgfusepath{clip}%
\pgfsetbuttcap%
\pgfsetroundjoin%
\definecolor{currentfill}{rgb}{0.749020,0.772549,0.792157}%
\pgfsetfillcolor{currentfill}%
\pgfsetfillopacity{0.160000}%
\pgfsetlinewidth{0.000000pt}%
\definecolor{currentstroke}{rgb}{0.000000,0.000000,0.000000}%
\pgfsetstrokecolor{currentstroke}%
\pgfsetdash{}{0pt}%
\pgfpathmoveto{\pgfqpoint{1.066106in}{0.637362in}}%
\pgfpathlineto{\pgfqpoint{1.054043in}{0.650778in}}%
\pgfpathlineto{\pgfqpoint{1.041978in}{0.664196in}}%
\pgfpathlineto{\pgfqpoint{1.040649in}{0.662774in}}%
\pgfpathlineto{\pgfqpoint{1.039536in}{0.661320in}}%
\pgfpathlineto{\pgfqpoint{1.052219in}{0.648613in}}%
\pgfpathlineto{\pgfqpoint{1.064895in}{0.635913in}}%
\pgfpathlineto{\pgfqpoint{1.065447in}{0.636645in}}%
\pgfpathlineto{\pgfqpoint{1.066106in}{0.637362in}}%
\pgfpathclose%
\pgfusepath{fill}%
\end{pgfscope}%
\begin{pgfscope}%
\pgfpathrectangle{\pgfqpoint{0.407592in}{0.222904in}}{\pgfqpoint{1.411163in}{1.411163in}}%
\pgfusepath{clip}%
\pgfsetbuttcap%
\pgfsetroundjoin%
\definecolor{currentfill}{rgb}{0.749020,0.772549,0.792157}%
\pgfsetfillcolor{currentfill}%
\pgfsetfillopacity{0.160000}%
\pgfsetlinewidth{0.000000pt}%
\definecolor{currentstroke}{rgb}{0.000000,0.000000,0.000000}%
\pgfsetstrokecolor{currentstroke}%
\pgfsetdash{}{0pt}%
\pgfpathmoveto{\pgfqpoint{1.090230in}{0.610533in}}%
\pgfpathlineto{\pgfqpoint{1.103360in}{0.621783in}}%
\pgfpathlineto{\pgfqpoint{1.116508in}{0.633047in}}%
\pgfpathlineto{\pgfqpoint{1.116539in}{0.633809in}}%
\pgfpathlineto{\pgfqpoint{1.116456in}{0.634567in}}%
\pgfpathlineto{\pgfqpoint{1.103337in}{0.622545in}}%
\pgfpathlineto{\pgfqpoint{1.090230in}{0.610533in}}%
\pgfpathlineto{\pgfqpoint{1.090230in}{0.610533in}}%
\pgfpathlineto{\pgfqpoint{1.090230in}{0.610533in}}%
\pgfpathclose%
\pgfusepath{fill}%
\end{pgfscope}%
\begin{pgfscope}%
\pgfpathrectangle{\pgfqpoint{0.407592in}{0.222904in}}{\pgfqpoint{1.411163in}{1.411163in}}%
\pgfusepath{clip}%
\pgfsetbuttcap%
\pgfsetroundjoin%
\definecolor{currentfill}{rgb}{0.749020,0.772549,0.792157}%
\pgfsetfillcolor{currentfill}%
\pgfsetfillopacity{0.160000}%
\pgfsetlinewidth{0.000000pt}%
\definecolor{currentstroke}{rgb}{0.000000,0.000000,0.000000}%
\pgfsetstrokecolor{currentstroke}%
\pgfsetdash{}{0pt}%
\pgfpathmoveto{\pgfqpoint{1.090230in}{0.610533in}}%
\pgfpathlineto{\pgfqpoint{1.077183in}{0.622469in}}%
\pgfpathlineto{\pgfqpoint{1.064125in}{0.634416in}}%
\pgfpathlineto{\pgfqpoint{1.063910in}{0.633656in}}%
\pgfpathlineto{\pgfqpoint{1.063809in}{0.632895in}}%
\pgfpathlineto{\pgfqpoint{1.077028in}{0.621706in}}%
\pgfpathlineto{\pgfqpoint{1.090230in}{0.610533in}}%
\pgfpathlineto{\pgfqpoint{1.090230in}{0.610533in}}%
\pgfpathlineto{\pgfqpoint{1.090230in}{0.610533in}}%
\pgfpathclose%
\pgfusepath{fill}%
\end{pgfscope}%
\begin{pgfscope}%
\pgfpathrectangle{\pgfqpoint{0.407592in}{0.222904in}}{\pgfqpoint{1.411163in}{1.411163in}}%
\pgfusepath{clip}%
\pgfsetbuttcap%
\pgfsetroundjoin%
\definecolor{currentfill}{rgb}{0.749020,0.772549,0.792157}%
\pgfsetfillcolor{currentfill}%
\pgfsetfillopacity{0.160000}%
\pgfsetlinewidth{0.000000pt}%
\definecolor{currentstroke}{rgb}{0.000000,0.000000,0.000000}%
\pgfsetstrokecolor{currentstroke}%
\pgfsetdash{}{0pt}%
\pgfpathmoveto{\pgfqpoint{1.141684in}{0.661613in}}%
\pgfpathlineto{\pgfqpoint{1.154563in}{0.674399in}}%
\pgfpathlineto{\pgfqpoint{1.167448in}{0.687189in}}%
\pgfpathlineto{\pgfqpoint{1.166141in}{0.689347in}}%
\pgfpathlineto{\pgfqpoint{1.164505in}{0.691454in}}%
\pgfpathlineto{\pgfqpoint{1.152124in}{0.677965in}}%
\pgfpathlineto{\pgfqpoint{1.139743in}{0.664477in}}%
\pgfpathlineto{\pgfqpoint{1.140824in}{0.663062in}}%
\pgfpathlineto{\pgfqpoint{1.141684in}{0.661613in}}%
\pgfpathclose%
\pgfusepath{fill}%
\end{pgfscope}%
\begin{pgfscope}%
\pgfpathrectangle{\pgfqpoint{0.407592in}{0.222904in}}{\pgfqpoint{1.411163in}{1.411163in}}%
\pgfusepath{clip}%
\pgfsetbuttcap%
\pgfsetroundjoin%
\definecolor{currentfill}{rgb}{0.749020,0.772549,0.792157}%
\pgfsetfillcolor{currentfill}%
\pgfsetfillopacity{0.160000}%
\pgfsetlinewidth{0.000000pt}%
\definecolor{currentstroke}{rgb}{0.000000,0.000000,0.000000}%
\pgfsetstrokecolor{currentstroke}%
\pgfsetdash{}{0pt}%
\pgfpathmoveto{\pgfqpoint{1.041978in}{0.664196in}}%
\pgfpathlineto{\pgfqpoint{1.029912in}{0.677615in}}%
\pgfpathlineto{\pgfqpoint{1.017845in}{0.691036in}}%
\pgfpathlineto{\pgfqpoint{1.015837in}{0.688918in}}%
\pgfpathlineto{\pgfqpoint{1.014150in}{0.686752in}}%
\pgfpathlineto{\pgfqpoint{1.026846in}{0.674032in}}%
\pgfpathlineto{\pgfqpoint{1.039536in}{0.661320in}}%
\pgfpathlineto{\pgfqpoint{1.040649in}{0.662774in}}%
\pgfpathlineto{\pgfqpoint{1.041978in}{0.664196in}}%
\pgfpathclose%
\pgfusepath{fill}%
\end{pgfscope}%
\begin{pgfscope}%
\pgfpathrectangle{\pgfqpoint{0.407592in}{0.222904in}}{\pgfqpoint{1.411163in}{1.411163in}}%
\pgfusepath{clip}%
\pgfsetbuttcap%
\pgfsetroundjoin%
\definecolor{currentfill}{rgb}{0.749020,0.772549,0.792157}%
\pgfsetfillcolor{currentfill}%
\pgfsetfillopacity{0.160000}%
\pgfsetlinewidth{0.000000pt}%
\definecolor{currentstroke}{rgb}{0.000000,0.000000,0.000000}%
\pgfsetstrokecolor{currentstroke}%
\pgfsetdash{}{0pt}%
\pgfpathmoveto{\pgfqpoint{1.090230in}{0.610533in}}%
\pgfpathlineto{\pgfqpoint{1.103153in}{0.621021in}}%
\pgfpathlineto{\pgfqpoint{1.116100in}{0.631528in}}%
\pgfpathlineto{\pgfqpoint{1.116361in}{0.632286in}}%
\pgfpathlineto{\pgfqpoint{1.116508in}{0.633047in}}%
\pgfpathlineto{\pgfqpoint{1.103360in}{0.621783in}}%
\pgfpathlineto{\pgfqpoint{1.090230in}{0.610533in}}%
\pgfpathlineto{\pgfqpoint{1.090230in}{0.610533in}}%
\pgfpathlineto{\pgfqpoint{1.090230in}{0.610533in}}%
\pgfpathclose%
\pgfusepath{fill}%
\end{pgfscope}%
\begin{pgfscope}%
\pgfpathrectangle{\pgfqpoint{0.407592in}{0.222904in}}{\pgfqpoint{1.411163in}{1.411163in}}%
\pgfusepath{clip}%
\pgfsetbuttcap%
\pgfsetroundjoin%
\definecolor{currentfill}{rgb}{0.749020,0.772549,0.792157}%
\pgfsetfillcolor{currentfill}%
\pgfsetfillopacity{0.160000}%
\pgfsetlinewidth{0.000000pt}%
\definecolor{currentstroke}{rgb}{0.000000,0.000000,0.000000}%
\pgfsetstrokecolor{currentstroke}%
\pgfsetdash{}{0pt}%
\pgfpathmoveto{\pgfqpoint{1.090230in}{0.610533in}}%
\pgfpathlineto{\pgfqpoint{1.077028in}{0.621706in}}%
\pgfpathlineto{\pgfqpoint{1.063809in}{0.632895in}}%
\pgfpathlineto{\pgfqpoint{1.063824in}{0.632134in}}%
\pgfpathlineto{\pgfqpoint{1.063954in}{0.631377in}}%
\pgfpathlineto{\pgfqpoint{1.077104in}{0.620946in}}%
\pgfpathlineto{\pgfqpoint{1.090230in}{0.610533in}}%
\pgfpathlineto{\pgfqpoint{1.090230in}{0.610533in}}%
\pgfpathlineto{\pgfqpoint{1.090230in}{0.610533in}}%
\pgfpathclose%
\pgfusepath{fill}%
\end{pgfscope}%
\begin{pgfscope}%
\pgfpathrectangle{\pgfqpoint{0.407592in}{0.222904in}}{\pgfqpoint{1.411163in}{1.411163in}}%
\pgfusepath{clip}%
\pgfsetbuttcap%
\pgfsetroundjoin%
\definecolor{currentfill}{rgb}{0.749020,0.772549,0.792157}%
\pgfsetfillcolor{currentfill}%
\pgfsetfillopacity{0.160000}%
\pgfsetlinewidth{0.000000pt}%
\definecolor{currentstroke}{rgb}{0.000000,0.000000,0.000000}%
\pgfsetstrokecolor{currentstroke}%
\pgfsetdash{}{0pt}%
\pgfpathmoveto{\pgfqpoint{1.167448in}{0.687189in}}%
\pgfpathlineto{\pgfqpoint{1.180339in}{0.699986in}}%
\pgfpathlineto{\pgfqpoint{1.193235in}{0.712789in}}%
\pgfpathlineto{\pgfqpoint{1.191471in}{0.715645in}}%
\pgfpathlineto{\pgfqpoint{1.189269in}{0.718434in}}%
\pgfpathlineto{\pgfqpoint{1.176886in}{0.704944in}}%
\pgfpathlineto{\pgfqpoint{1.164505in}{0.691454in}}%
\pgfpathlineto{\pgfqpoint{1.166141in}{0.689347in}}%
\pgfpathlineto{\pgfqpoint{1.167448in}{0.687189in}}%
\pgfpathclose%
\pgfusepath{fill}%
\end{pgfscope}%
\begin{pgfscope}%
\pgfpathrectangle{\pgfqpoint{0.407592in}{0.222904in}}{\pgfqpoint{1.411163in}{1.411163in}}%
\pgfusepath{clip}%
\pgfsetbuttcap%
\pgfsetroundjoin%
\definecolor{currentfill}{rgb}{0.749020,0.772549,0.792157}%
\pgfsetfillcolor{currentfill}%
\pgfsetfillopacity{0.160000}%
\pgfsetlinewidth{0.000000pt}%
\definecolor{currentstroke}{rgb}{0.000000,0.000000,0.000000}%
\pgfsetstrokecolor{currentstroke}%
\pgfsetdash{}{0pt}%
\pgfpathmoveto{\pgfqpoint{1.090230in}{0.610533in}}%
\pgfpathlineto{\pgfqpoint{1.102720in}{0.620274in}}%
\pgfpathlineto{\pgfqpoint{1.115238in}{0.630036in}}%
\pgfpathlineto{\pgfqpoint{1.115725in}{0.630777in}}%
\pgfpathlineto{\pgfqpoint{1.116100in}{0.631528in}}%
\pgfpathlineto{\pgfqpoint{1.103153in}{0.621021in}}%
\pgfpathlineto{\pgfqpoint{1.090230in}{0.610533in}}%
\pgfpathlineto{\pgfqpoint{1.090230in}{0.610533in}}%
\pgfpathlineto{\pgfqpoint{1.090230in}{0.610533in}}%
\pgfpathclose%
\pgfusepath{fill}%
\end{pgfscope}%
\begin{pgfscope}%
\pgfpathrectangle{\pgfqpoint{0.407592in}{0.222904in}}{\pgfqpoint{1.411163in}{1.411163in}}%
\pgfusepath{clip}%
\pgfsetbuttcap%
\pgfsetroundjoin%
\definecolor{currentfill}{rgb}{0.749020,0.772549,0.792157}%
\pgfsetfillcolor{currentfill}%
\pgfsetfillopacity{0.160000}%
\pgfsetlinewidth{0.000000pt}%
\definecolor{currentstroke}{rgb}{0.000000,0.000000,0.000000}%
\pgfsetstrokecolor{currentstroke}%
\pgfsetdash{}{0pt}%
\pgfpathmoveto{\pgfqpoint{1.090230in}{0.610533in}}%
\pgfpathlineto{\pgfqpoint{1.077104in}{0.620946in}}%
\pgfpathlineto{\pgfqpoint{1.063954in}{0.631377in}}%
\pgfpathlineto{\pgfqpoint{1.064199in}{0.630628in}}%
\pgfpathlineto{\pgfqpoint{1.064558in}{0.629889in}}%
\pgfpathlineto{\pgfqpoint{1.077408in}{0.620200in}}%
\pgfpathlineto{\pgfqpoint{1.090230in}{0.610533in}}%
\pgfpathlineto{\pgfqpoint{1.090230in}{0.610533in}}%
\pgfpathlineto{\pgfqpoint{1.090230in}{0.610533in}}%
\pgfpathclose%
\pgfusepath{fill}%
\end{pgfscope}%
\begin{pgfscope}%
\pgfpathrectangle{\pgfqpoint{0.407592in}{0.222904in}}{\pgfqpoint{1.411163in}{1.411163in}}%
\pgfusepath{clip}%
\pgfsetbuttcap%
\pgfsetroundjoin%
\definecolor{currentfill}{rgb}{0.749020,0.772549,0.792157}%
\pgfsetfillcolor{currentfill}%
\pgfsetfillopacity{0.160000}%
\pgfsetlinewidth{0.000000pt}%
\definecolor{currentstroke}{rgb}{0.000000,0.000000,0.000000}%
\pgfsetstrokecolor{currentstroke}%
\pgfsetdash{}{0pt}%
\pgfpathmoveto{\pgfqpoint{1.116456in}{0.634567in}}%
\pgfpathlineto{\pgfqpoint{1.129586in}{0.646600in}}%
\pgfpathlineto{\pgfqpoint{1.142729in}{0.658644in}}%
\pgfpathlineto{\pgfqpoint{1.142320in}{0.660139in}}%
\pgfpathlineto{\pgfqpoint{1.141684in}{0.661613in}}%
\pgfpathlineto{\pgfqpoint{1.128812in}{0.648834in}}%
\pgfpathlineto{\pgfqpoint{1.115945in}{0.636061in}}%
\pgfpathlineto{\pgfqpoint{1.116257in}{0.635319in}}%
\pgfpathlineto{\pgfqpoint{1.116456in}{0.634567in}}%
\pgfpathclose%
\pgfusepath{fill}%
\end{pgfscope}%
\begin{pgfscope}%
\pgfpathrectangle{\pgfqpoint{0.407592in}{0.222904in}}{\pgfqpoint{1.411163in}{1.411163in}}%
\pgfusepath{clip}%
\pgfsetbuttcap%
\pgfsetroundjoin%
\definecolor{currentfill}{rgb}{0.749020,0.772549,0.792157}%
\pgfsetfillcolor{currentfill}%
\pgfsetfillopacity{0.160000}%
\pgfsetlinewidth{0.000000pt}%
\definecolor{currentstroke}{rgb}{0.000000,0.000000,0.000000}%
\pgfsetstrokecolor{currentstroke}%
\pgfsetdash{}{0pt}%
\pgfpathmoveto{\pgfqpoint{1.017845in}{0.691036in}}%
\pgfpathlineto{\pgfqpoint{1.005777in}{0.704458in}}%
\pgfpathlineto{\pgfqpoint{0.993708in}{0.717881in}}%
\pgfpathlineto{\pgfqpoint{0.991010in}{0.715077in}}%
\pgfpathlineto{\pgfqpoint{0.988739in}{0.712209in}}%
\pgfpathlineto{\pgfqpoint{1.001448in}{0.699477in}}%
\pgfpathlineto{\pgfqpoint{1.014150in}{0.686752in}}%
\pgfpathlineto{\pgfqpoint{1.015837in}{0.688918in}}%
\pgfpathlineto{\pgfqpoint{1.017845in}{0.691036in}}%
\pgfpathclose%
\pgfusepath{fill}%
\end{pgfscope}%
\begin{pgfscope}%
\pgfpathrectangle{\pgfqpoint{0.407592in}{0.222904in}}{\pgfqpoint{1.411163in}{1.411163in}}%
\pgfusepath{clip}%
\pgfsetbuttcap%
\pgfsetroundjoin%
\definecolor{currentfill}{rgb}{0.749020,0.772549,0.792157}%
\pgfsetfillcolor{currentfill}%
\pgfsetfillopacity{0.160000}%
\pgfsetlinewidth{0.000000pt}%
\definecolor{currentstroke}{rgb}{0.000000,0.000000,0.000000}%
\pgfsetstrokecolor{currentstroke}%
\pgfsetdash{}{0pt}%
\pgfpathmoveto{\pgfqpoint{1.064895in}{0.635913in}}%
\pgfpathlineto{\pgfqpoint{1.052219in}{0.648613in}}%
\pgfpathlineto{\pgfqpoint{1.039536in}{0.661320in}}%
\pgfpathlineto{\pgfqpoint{1.038642in}{0.659840in}}%
\pgfpathlineto{\pgfqpoint{1.037972in}{0.658343in}}%
\pgfpathlineto{\pgfqpoint{1.051055in}{0.646373in}}%
\pgfpathlineto{\pgfqpoint{1.064125in}{0.634416in}}%
\pgfpathlineto{\pgfqpoint{1.064454in}{0.635169in}}%
\pgfpathlineto{\pgfqpoint{1.064895in}{0.635913in}}%
\pgfpathclose%
\pgfusepath{fill}%
\end{pgfscope}%
\begin{pgfscope}%
\pgfpathrectangle{\pgfqpoint{0.407592in}{0.222904in}}{\pgfqpoint{1.411163in}{1.411163in}}%
\pgfusepath{clip}%
\pgfsetbuttcap%
\pgfsetroundjoin%
\definecolor{currentfill}{rgb}{0.749020,0.772549,0.792157}%
\pgfsetfillcolor{currentfill}%
\pgfsetfillopacity{0.160000}%
\pgfsetlinewidth{0.000000pt}%
\definecolor{currentstroke}{rgb}{0.000000,0.000000,0.000000}%
\pgfsetstrokecolor{currentstroke}%
\pgfsetdash{}{0pt}%
\pgfpathmoveto{\pgfqpoint{1.193235in}{0.712789in}}%
\pgfpathlineto{\pgfqpoint{1.206138in}{0.725598in}}%
\pgfpathlineto{\pgfqpoint{1.219047in}{0.738413in}}%
\pgfpathlineto{\pgfqpoint{1.216815in}{0.741958in}}%
\pgfpathlineto{\pgfqpoint{1.214036in}{0.745418in}}%
\pgfpathlineto{\pgfqpoint{1.201652in}{0.731926in}}%
\pgfpathlineto{\pgfqpoint{1.189269in}{0.718434in}}%
\pgfpathlineto{\pgfqpoint{1.191471in}{0.715645in}}%
\pgfpathlineto{\pgfqpoint{1.193235in}{0.712789in}}%
\pgfpathclose%
\pgfusepath{fill}%
\end{pgfscope}%
\begin{pgfscope}%
\pgfpathrectangle{\pgfqpoint{0.407592in}{0.222904in}}{\pgfqpoint{1.411163in}{1.411163in}}%
\pgfusepath{clip}%
\pgfsetbuttcap%
\pgfsetroundjoin%
\definecolor{currentfill}{rgb}{0.749020,0.772549,0.792157}%
\pgfsetfillcolor{currentfill}%
\pgfsetfillopacity{0.160000}%
\pgfsetlinewidth{0.000000pt}%
\definecolor{currentstroke}{rgb}{0.000000,0.000000,0.000000}%
\pgfsetstrokecolor{currentstroke}%
\pgfsetdash{}{0pt}%
\pgfpathmoveto{\pgfqpoint{1.090230in}{0.610533in}}%
\pgfpathlineto{\pgfqpoint{1.102068in}{0.619554in}}%
\pgfpathlineto{\pgfqpoint{1.113937in}{0.628598in}}%
\pgfpathlineto{\pgfqpoint{1.114641in}{0.629309in}}%
\pgfpathlineto{\pgfqpoint{1.115238in}{0.630036in}}%
\pgfpathlineto{\pgfqpoint{1.102720in}{0.620274in}}%
\pgfpathlineto{\pgfqpoint{1.090230in}{0.610533in}}%
\pgfpathlineto{\pgfqpoint{1.090230in}{0.610533in}}%
\pgfpathlineto{\pgfqpoint{1.090230in}{0.610533in}}%
\pgfpathclose%
\pgfusepath{fill}%
\end{pgfscope}%
\begin{pgfscope}%
\pgfpathrectangle{\pgfqpoint{0.407592in}{0.222904in}}{\pgfqpoint{1.411163in}{1.411163in}}%
\pgfusepath{clip}%
\pgfsetbuttcap%
\pgfsetroundjoin%
\definecolor{currentfill}{rgb}{0.749020,0.772549,0.792157}%
\pgfsetfillcolor{currentfill}%
\pgfsetfillopacity{0.160000}%
\pgfsetlinewidth{0.000000pt}%
\definecolor{currentstroke}{rgb}{0.000000,0.000000,0.000000}%
\pgfsetstrokecolor{currentstroke}%
\pgfsetdash{}{0pt}%
\pgfpathmoveto{\pgfqpoint{1.090230in}{0.610533in}}%
\pgfpathlineto{\pgfqpoint{1.077408in}{0.620200in}}%
\pgfpathlineto{\pgfqpoint{1.064558in}{0.629889in}}%
\pgfpathlineto{\pgfqpoint{1.065029in}{0.629165in}}%
\pgfpathlineto{\pgfqpoint{1.065610in}{0.628457in}}%
\pgfpathlineto{\pgfqpoint{1.077936in}{0.619484in}}%
\pgfpathlineto{\pgfqpoint{1.090230in}{0.610533in}}%
\pgfpathlineto{\pgfqpoint{1.090230in}{0.610533in}}%
\pgfpathlineto{\pgfqpoint{1.090230in}{0.610533in}}%
\pgfpathclose%
\pgfusepath{fill}%
\end{pgfscope}%
\begin{pgfscope}%
\pgfpathrectangle{\pgfqpoint{0.407592in}{0.222904in}}{\pgfqpoint{1.411163in}{1.411163in}}%
\pgfusepath{clip}%
\pgfsetbuttcap%
\pgfsetroundjoin%
\definecolor{currentfill}{rgb}{0.749020,0.772549,0.792157}%
\pgfsetfillcolor{currentfill}%
\pgfsetfillopacity{0.160000}%
\pgfsetlinewidth{0.000000pt}%
\definecolor{currentstroke}{rgb}{0.000000,0.000000,0.000000}%
\pgfsetstrokecolor{currentstroke}%
\pgfsetdash{}{0pt}%
\pgfpathmoveto{\pgfqpoint{0.993708in}{0.717881in}}%
\pgfpathlineto{\pgfqpoint{0.981637in}{0.731305in}}%
\pgfpathlineto{\pgfqpoint{0.969565in}{0.744731in}}%
\pgfpathlineto{\pgfqpoint{0.966167in}{0.741253in}}%
\pgfpathlineto{\pgfqpoint{0.963301in}{0.737693in}}%
\pgfpathlineto{\pgfqpoint{0.976023in}{0.724948in}}%
\pgfpathlineto{\pgfqpoint{0.988739in}{0.712209in}}%
\pgfpathlineto{\pgfqpoint{0.991010in}{0.715077in}}%
\pgfpathlineto{\pgfqpoint{0.993708in}{0.717881in}}%
\pgfpathclose%
\pgfusepath{fill}%
\end{pgfscope}%
\begin{pgfscope}%
\pgfpathrectangle{\pgfqpoint{0.407592in}{0.222904in}}{\pgfqpoint{1.411163in}{1.411163in}}%
\pgfusepath{clip}%
\pgfsetbuttcap%
\pgfsetroundjoin%
\definecolor{currentfill}{rgb}{0.749020,0.772549,0.792157}%
\pgfsetfillcolor{currentfill}%
\pgfsetfillopacity{0.160000}%
\pgfsetlinewidth{0.000000pt}%
\definecolor{currentstroke}{rgb}{0.000000,0.000000,0.000000}%
\pgfsetstrokecolor{currentstroke}%
\pgfsetdash{}{0pt}%
\pgfpathmoveto{\pgfqpoint{1.090230in}{0.610533in}}%
\pgfpathlineto{\pgfqpoint{1.101208in}{0.618873in}}%
\pgfpathlineto{\pgfqpoint{1.112219in}{0.627238in}}%
\pgfpathlineto{\pgfqpoint{1.113128in}{0.627907in}}%
\pgfpathlineto{\pgfqpoint{1.113937in}{0.628598in}}%
\pgfpathlineto{\pgfqpoint{1.102068in}{0.619554in}}%
\pgfpathlineto{\pgfqpoint{1.090230in}{0.610533in}}%
\pgfpathlineto{\pgfqpoint{1.090230in}{0.610533in}}%
\pgfpathlineto{\pgfqpoint{1.090230in}{0.610533in}}%
\pgfpathclose%
\pgfusepath{fill}%
\end{pgfscope}%
\begin{pgfscope}%
\pgfpathrectangle{\pgfqpoint{0.407592in}{0.222904in}}{\pgfqpoint{1.411163in}{1.411163in}}%
\pgfusepath{clip}%
\pgfsetbuttcap%
\pgfsetroundjoin%
\definecolor{currentfill}{rgb}{0.749020,0.772549,0.792157}%
\pgfsetfillcolor{currentfill}%
\pgfsetfillopacity{0.160000}%
\pgfsetlinewidth{0.000000pt}%
\definecolor{currentstroke}{rgb}{0.000000,0.000000,0.000000}%
\pgfsetstrokecolor{currentstroke}%
\pgfsetdash{}{0pt}%
\pgfpathmoveto{\pgfqpoint{1.090230in}{0.610533in}}%
\pgfpathlineto{\pgfqpoint{1.077936in}{0.619484in}}%
\pgfpathlineto{\pgfqpoint{1.065610in}{0.628457in}}%
\pgfpathlineto{\pgfqpoint{1.066299in}{0.627771in}}%
\pgfpathlineto{\pgfqpoint{1.067093in}{0.627107in}}%
\pgfpathlineto{\pgfqpoint{1.078678in}{0.618808in}}%
\pgfpathlineto{\pgfqpoint{1.090230in}{0.610533in}}%
\pgfpathlineto{\pgfqpoint{1.090230in}{0.610533in}}%
\pgfpathlineto{\pgfqpoint{1.090230in}{0.610533in}}%
\pgfpathclose%
\pgfusepath{fill}%
\end{pgfscope}%
\begin{pgfscope}%
\pgfpathrectangle{\pgfqpoint{0.407592in}{0.222904in}}{\pgfqpoint{1.411163in}{1.411163in}}%
\pgfusepath{clip}%
\pgfsetbuttcap%
\pgfsetroundjoin%
\definecolor{currentfill}{rgb}{0.749020,0.772549,0.792157}%
\pgfsetfillcolor{currentfill}%
\pgfsetfillopacity{0.160000}%
\pgfsetlinewidth{0.000000pt}%
\definecolor{currentstroke}{rgb}{0.000000,0.000000,0.000000}%
\pgfsetstrokecolor{currentstroke}%
\pgfsetdash{}{0pt}%
\pgfpathmoveto{\pgfqpoint{1.219047in}{0.738413in}}%
\pgfpathlineto{\pgfqpoint{1.231962in}{0.751234in}}%
\pgfpathlineto{\pgfqpoint{1.244882in}{0.764061in}}%
\pgfpathlineto{\pgfqpoint{1.242171in}{0.768284in}}%
\pgfpathlineto{\pgfqpoint{1.238806in}{0.772404in}}%
\pgfpathlineto{\pgfqpoint{1.226421in}{0.758911in}}%
\pgfpathlineto{\pgfqpoint{1.214036in}{0.745418in}}%
\pgfpathlineto{\pgfqpoint{1.216815in}{0.741958in}}%
\pgfpathlineto{\pgfqpoint{1.219047in}{0.738413in}}%
\pgfpathclose%
\pgfusepath{fill}%
\end{pgfscope}%
\begin{pgfscope}%
\pgfpathrectangle{\pgfqpoint{0.407592in}{0.222904in}}{\pgfqpoint{1.411163in}{1.411163in}}%
\pgfusepath{clip}%
\pgfsetbuttcap%
\pgfsetroundjoin%
\definecolor{currentfill}{rgb}{0.749020,0.772549,0.792157}%
\pgfsetfillcolor{currentfill}%
\pgfsetfillopacity{0.160000}%
\pgfsetlinewidth{0.000000pt}%
\definecolor{currentstroke}{rgb}{0.000000,0.000000,0.000000}%
\pgfsetstrokecolor{currentstroke}%
\pgfsetdash{}{0pt}%
\pgfpathmoveto{\pgfqpoint{1.090230in}{0.610533in}}%
\pgfpathlineto{\pgfqpoint{1.100705in}{0.618552in}}%
\pgfpathlineto{\pgfqpoint{1.111212in}{0.626595in}}%
\pgfpathlineto{\pgfqpoint{1.112219in}{0.627238in}}%
\pgfpathlineto{\pgfqpoint{1.101208in}{0.618873in}}%
\pgfpathlineto{\pgfqpoint{1.090230in}{0.610533in}}%
\pgfusepath{fill}%
\end{pgfscope}%
\begin{pgfscope}%
\pgfpathrectangle{\pgfqpoint{0.407592in}{0.222904in}}{\pgfqpoint{1.411163in}{1.411163in}}%
\pgfusepath{clip}%
\pgfsetbuttcap%
\pgfsetroundjoin%
\definecolor{currentfill}{rgb}{0.749020,0.772549,0.792157}%
\pgfsetfillcolor{currentfill}%
\pgfsetfillopacity{0.160000}%
\pgfsetlinewidth{0.000000pt}%
\definecolor{currentstroke}{rgb}{0.000000,0.000000,0.000000}%
\pgfsetstrokecolor{currentstroke}%
\pgfsetdash{}{0pt}%
\pgfpathmoveto{\pgfqpoint{1.090230in}{0.610533in}}%
\pgfpathlineto{\pgfqpoint{1.078678in}{0.618808in}}%
\pgfpathlineto{\pgfqpoint{1.067093in}{0.627107in}}%
\pgfpathlineto{\pgfqpoint{1.067988in}{0.626470in}}%
\pgfpathlineto{\pgfqpoint{1.068981in}{0.625862in}}%
\pgfpathlineto{\pgfqpoint{1.079623in}{0.618184in}}%
\pgfpathlineto{\pgfqpoint{1.090230in}{0.610533in}}%
\pgfpathlineto{\pgfqpoint{1.090230in}{0.610533in}}%
\pgfpathlineto{\pgfqpoint{1.090230in}{0.610533in}}%
\pgfpathclose%
\pgfusepath{fill}%
\end{pgfscope}%
\begin{pgfscope}%
\pgfpathrectangle{\pgfqpoint{0.407592in}{0.222904in}}{\pgfqpoint{1.411163in}{1.411163in}}%
\pgfusepath{clip}%
\pgfsetbuttcap%
\pgfsetroundjoin%
\definecolor{currentfill}{rgb}{0.749020,0.772549,0.792157}%
\pgfsetfillcolor{currentfill}%
\pgfsetfillopacity{0.160000}%
\pgfsetlinewidth{0.000000pt}%
\definecolor{currentstroke}{rgb}{0.000000,0.000000,0.000000}%
\pgfsetstrokecolor{currentstroke}%
\pgfsetdash{}{0pt}%
\pgfpathmoveto{\pgfqpoint{1.090230in}{0.610533in}}%
\pgfpathlineto{\pgfqpoint{1.099562in}{0.617953in}}%
\pgfpathlineto{\pgfqpoint{1.108927in}{0.625398in}}%
\pgfpathlineto{\pgfqpoint{1.110114in}{0.625981in}}%
\pgfpathlineto{\pgfqpoint{1.111212in}{0.626595in}}%
\pgfpathlineto{\pgfqpoint{1.100705in}{0.618552in}}%
\pgfpathlineto{\pgfqpoint{1.090230in}{0.610533in}}%
\pgfpathlineto{\pgfqpoint{1.090230in}{0.610533in}}%
\pgfpathlineto{\pgfqpoint{1.090230in}{0.610533in}}%
\pgfpathclose%
\pgfusepath{fill}%
\end{pgfscope}%
\begin{pgfscope}%
\pgfpathrectangle{\pgfqpoint{0.407592in}{0.222904in}}{\pgfqpoint{1.411163in}{1.411163in}}%
\pgfusepath{clip}%
\pgfsetbuttcap%
\pgfsetroundjoin%
\definecolor{currentfill}{rgb}{0.749020,0.772549,0.792157}%
\pgfsetfillcolor{currentfill}%
\pgfsetfillopacity{0.160000}%
\pgfsetlinewidth{0.000000pt}%
\definecolor{currentstroke}{rgb}{0.000000,0.000000,0.000000}%
\pgfsetstrokecolor{currentstroke}%
\pgfsetdash{}{0pt}%
\pgfpathmoveto{\pgfqpoint{1.116508in}{0.633047in}}%
\pgfpathlineto{\pgfqpoint{1.129673in}{0.644327in}}%
\pgfpathlineto{\pgfqpoint{1.142855in}{0.655621in}}%
\pgfpathlineto{\pgfqpoint{1.142907in}{0.657136in}}%
\pgfpathlineto{\pgfqpoint{1.142729in}{0.658644in}}%
\pgfpathlineto{\pgfqpoint{1.129586in}{0.646600in}}%
\pgfpathlineto{\pgfqpoint{1.116456in}{0.634567in}}%
\pgfpathlineto{\pgfqpoint{1.116539in}{0.633809in}}%
\pgfpathlineto{\pgfqpoint{1.116508in}{0.633047in}}%
\pgfpathclose%
\pgfusepath{fill}%
\end{pgfscope}%
\begin{pgfscope}%
\pgfpathrectangle{\pgfqpoint{0.407592in}{0.222904in}}{\pgfqpoint{1.411163in}{1.411163in}}%
\pgfusepath{clip}%
\pgfsetbuttcap%
\pgfsetroundjoin%
\definecolor{currentfill}{rgb}{0.749020,0.772549,0.792157}%
\pgfsetfillcolor{currentfill}%
\pgfsetfillopacity{0.160000}%
\pgfsetlinewidth{0.000000pt}%
\definecolor{currentstroke}{rgb}{0.000000,0.000000,0.000000}%
\pgfsetstrokecolor{currentstroke}%
\pgfsetdash{}{0pt}%
\pgfpathmoveto{\pgfqpoint{0.969565in}{0.744731in}}%
\pgfpathlineto{\pgfqpoint{0.957492in}{0.758158in}}%
\pgfpathlineto{\pgfqpoint{0.945418in}{0.771587in}}%
\pgfpathlineto{\pgfqpoint{0.941310in}{0.767444in}}%
\pgfpathlineto{\pgfqpoint{0.937838in}{0.763203in}}%
\pgfpathlineto{\pgfqpoint{0.950573in}{0.750445in}}%
\pgfpathlineto{\pgfqpoint{0.963301in}{0.737693in}}%
\pgfpathlineto{\pgfqpoint{0.966167in}{0.741253in}}%
\pgfpathlineto{\pgfqpoint{0.969565in}{0.744731in}}%
\pgfpathclose%
\pgfusepath{fill}%
\end{pgfscope}%
\begin{pgfscope}%
\pgfpathrectangle{\pgfqpoint{0.407592in}{0.222904in}}{\pgfqpoint{1.411163in}{1.411163in}}%
\pgfusepath{clip}%
\pgfsetbuttcap%
\pgfsetroundjoin%
\definecolor{currentfill}{rgb}{0.749020,0.772549,0.792157}%
\pgfsetfillcolor{currentfill}%
\pgfsetfillopacity{0.160000}%
\pgfsetlinewidth{0.000000pt}%
\definecolor{currentstroke}{rgb}{0.000000,0.000000,0.000000}%
\pgfsetstrokecolor{currentstroke}%
\pgfsetdash{}{0pt}%
\pgfpathmoveto{\pgfqpoint{1.142729in}{0.658644in}}%
\pgfpathlineto{\pgfqpoint{1.155882in}{0.670699in}}%
\pgfpathlineto{\pgfqpoint{1.169048in}{0.682764in}}%
\pgfpathlineto{\pgfqpoint{1.168418in}{0.684992in}}%
\pgfpathlineto{\pgfqpoint{1.167448in}{0.687189in}}%
\pgfpathlineto{\pgfqpoint{1.154563in}{0.674399in}}%
\pgfpathlineto{\pgfqpoint{1.141684in}{0.661613in}}%
\pgfpathlineto{\pgfqpoint{1.142320in}{0.660139in}}%
\pgfpathlineto{\pgfqpoint{1.142729in}{0.658644in}}%
\pgfpathclose%
\pgfusepath{fill}%
\end{pgfscope}%
\begin{pgfscope}%
\pgfpathrectangle{\pgfqpoint{0.407592in}{0.222904in}}{\pgfqpoint{1.411163in}{1.411163in}}%
\pgfusepath{clip}%
\pgfsetbuttcap%
\pgfsetroundjoin%
\definecolor{currentfill}{rgb}{0.749020,0.772549,0.792157}%
\pgfsetfillcolor{currentfill}%
\pgfsetfillopacity{0.160000}%
\pgfsetlinewidth{0.000000pt}%
\definecolor{currentstroke}{rgb}{0.000000,0.000000,0.000000}%
\pgfsetstrokecolor{currentstroke}%
\pgfsetdash{}{0pt}%
\pgfpathmoveto{\pgfqpoint{1.244882in}{0.764061in}}%
\pgfpathlineto{\pgfqpoint{1.257809in}{0.776894in}}%
\pgfpathlineto{\pgfqpoint{1.270742in}{0.789733in}}%
\pgfpathlineto{\pgfqpoint{1.267541in}{0.794623in}}%
\pgfpathlineto{\pgfqpoint{1.263579in}{0.799394in}}%
\pgfpathlineto{\pgfqpoint{1.251192in}{0.785899in}}%
\pgfpathlineto{\pgfqpoint{1.238806in}{0.772404in}}%
\pgfpathlineto{\pgfqpoint{1.242171in}{0.768284in}}%
\pgfpathlineto{\pgfqpoint{1.244882in}{0.764061in}}%
\pgfpathclose%
\pgfusepath{fill}%
\end{pgfscope}%
\begin{pgfscope}%
\pgfpathrectangle{\pgfqpoint{0.407592in}{0.222904in}}{\pgfqpoint{1.411163in}{1.411163in}}%
\pgfusepath{clip}%
\pgfsetbuttcap%
\pgfsetroundjoin%
\definecolor{currentfill}{rgb}{0.749020,0.772549,0.792157}%
\pgfsetfillcolor{currentfill}%
\pgfsetfillopacity{0.160000}%
\pgfsetlinewidth{0.000000pt}%
\definecolor{currentstroke}{rgb}{0.000000,0.000000,0.000000}%
\pgfsetstrokecolor{currentstroke}%
\pgfsetdash{}{0pt}%
\pgfpathmoveto{\pgfqpoint{1.064125in}{0.634416in}}%
\pgfpathlineto{\pgfqpoint{1.051055in}{0.646373in}}%
\pgfpathlineto{\pgfqpoint{1.037972in}{0.658343in}}%
\pgfpathlineto{\pgfqpoint{1.037530in}{0.656833in}}%
\pgfpathlineto{\pgfqpoint{1.037317in}{0.655317in}}%
\pgfpathlineto{\pgfqpoint{1.050572in}{0.644098in}}%
\pgfpathlineto{\pgfqpoint{1.063809in}{0.632895in}}%
\pgfpathlineto{\pgfqpoint{1.063910in}{0.633656in}}%
\pgfpathlineto{\pgfqpoint{1.064125in}{0.634416in}}%
\pgfpathclose%
\pgfusepath{fill}%
\end{pgfscope}%
\begin{pgfscope}%
\pgfpathrectangle{\pgfqpoint{0.407592in}{0.222904in}}{\pgfqpoint{1.411163in}{1.411163in}}%
\pgfusepath{clip}%
\pgfsetbuttcap%
\pgfsetroundjoin%
\definecolor{currentfill}{rgb}{0.749020,0.772549,0.792157}%
\pgfsetfillcolor{currentfill}%
\pgfsetfillopacity{0.160000}%
\pgfsetlinewidth{0.000000pt}%
\definecolor{currentstroke}{rgb}{0.000000,0.000000,0.000000}%
\pgfsetstrokecolor{currentstroke}%
\pgfsetdash{}{0pt}%
\pgfpathmoveto{\pgfqpoint{1.090230in}{0.610533in}}%
\pgfpathlineto{\pgfqpoint{1.079623in}{0.618184in}}%
\pgfpathlineto{\pgfqpoint{1.068981in}{0.625862in}}%
\pgfpathlineto{\pgfqpoint{1.070067in}{0.625285in}}%
\pgfpathlineto{\pgfqpoint{1.071242in}{0.624743in}}%
\pgfpathlineto{\pgfqpoint{1.080753in}{0.617625in}}%
\pgfpathlineto{\pgfqpoint{1.090230in}{0.610533in}}%
\pgfpathlineto{\pgfqpoint{1.090230in}{0.610533in}}%
\pgfpathlineto{\pgfqpoint{1.090230in}{0.610533in}}%
\pgfpathclose%
\pgfusepath{fill}%
\end{pgfscope}%
\begin{pgfscope}%
\pgfpathrectangle{\pgfqpoint{0.407592in}{0.222904in}}{\pgfqpoint{1.411163in}{1.411163in}}%
\pgfusepath{clip}%
\pgfsetbuttcap%
\pgfsetroundjoin%
\definecolor{currentfill}{rgb}{0.749020,0.772549,0.792157}%
\pgfsetfillcolor{currentfill}%
\pgfsetfillopacity{0.160000}%
\pgfsetlinewidth{0.000000pt}%
\definecolor{currentstroke}{rgb}{0.000000,0.000000,0.000000}%
\pgfsetstrokecolor{currentstroke}%
\pgfsetdash{}{0pt}%
\pgfpathmoveto{\pgfqpoint{1.039536in}{0.661320in}}%
\pgfpathlineto{\pgfqpoint{1.026846in}{0.674032in}}%
\pgfpathlineto{\pgfqpoint{1.014150in}{0.686752in}}%
\pgfpathlineto{\pgfqpoint{1.012792in}{0.684547in}}%
\pgfpathlineto{\pgfqpoint{1.011770in}{0.682314in}}%
\pgfpathlineto{\pgfqpoint{1.024877in}{0.670323in}}%
\pgfpathlineto{\pgfqpoint{1.037972in}{0.658343in}}%
\pgfpathlineto{\pgfqpoint{1.038642in}{0.659840in}}%
\pgfpathlineto{\pgfqpoint{1.039536in}{0.661320in}}%
\pgfpathclose%
\pgfusepath{fill}%
\end{pgfscope}%
\begin{pgfscope}%
\pgfpathrectangle{\pgfqpoint{0.407592in}{0.222904in}}{\pgfqpoint{1.411163in}{1.411163in}}%
\pgfusepath{clip}%
\pgfsetbuttcap%
\pgfsetroundjoin%
\definecolor{currentfill}{rgb}{0.749020,0.772549,0.792157}%
\pgfsetfillcolor{currentfill}%
\pgfsetfillopacity{0.160000}%
\pgfsetlinewidth{0.000000pt}%
\definecolor{currentstroke}{rgb}{0.000000,0.000000,0.000000}%
\pgfsetstrokecolor{currentstroke}%
\pgfsetdash{}{0pt}%
\pgfpathmoveto{\pgfqpoint{1.090230in}{0.610533in}}%
\pgfpathlineto{\pgfqpoint{1.098256in}{0.617422in}}%
\pgfpathlineto{\pgfqpoint{1.106312in}{0.624336in}}%
\pgfpathlineto{\pgfqpoint{1.107658in}{0.624849in}}%
\pgfpathlineto{\pgfqpoint{1.108927in}{0.625398in}}%
\pgfpathlineto{\pgfqpoint{1.099562in}{0.617953in}}%
\pgfpathlineto{\pgfqpoint{1.090230in}{0.610533in}}%
\pgfpathlineto{\pgfqpoint{1.090230in}{0.610533in}}%
\pgfpathlineto{\pgfqpoint{1.090230in}{0.610533in}}%
\pgfpathclose%
\pgfusepath{fill}%
\end{pgfscope}%
\begin{pgfscope}%
\pgfpathrectangle{\pgfqpoint{0.407592in}{0.222904in}}{\pgfqpoint{1.411163in}{1.411163in}}%
\pgfusepath{clip}%
\pgfsetbuttcap%
\pgfsetroundjoin%
\definecolor{currentfill}{rgb}{0.749020,0.772549,0.792157}%
\pgfsetfillcolor{currentfill}%
\pgfsetfillopacity{0.160000}%
\pgfsetlinewidth{0.000000pt}%
\definecolor{currentstroke}{rgb}{0.000000,0.000000,0.000000}%
\pgfsetstrokecolor{currentstroke}%
\pgfsetdash{}{0pt}%
\pgfpathmoveto{\pgfqpoint{1.090230in}{0.610533in}}%
\pgfpathlineto{\pgfqpoint{1.080753in}{0.617625in}}%
\pgfpathlineto{\pgfqpoint{1.071242in}{0.624743in}}%
\pgfpathlineto{\pgfqpoint{1.072499in}{0.624238in}}%
\pgfpathlineto{\pgfqpoint{1.073835in}{0.623771in}}%
\pgfpathlineto{\pgfqpoint{1.082048in}{0.617139in}}%
\pgfpathlineto{\pgfqpoint{1.090230in}{0.610533in}}%
\pgfpathlineto{\pgfqpoint{1.090230in}{0.610533in}}%
\pgfpathlineto{\pgfqpoint{1.090230in}{0.610533in}}%
\pgfpathclose%
\pgfusepath{fill}%
\end{pgfscope}%
\begin{pgfscope}%
\pgfpathrectangle{\pgfqpoint{0.407592in}{0.222904in}}{\pgfqpoint{1.411163in}{1.411163in}}%
\pgfusepath{clip}%
\pgfsetbuttcap%
\pgfsetroundjoin%
\definecolor{currentfill}{rgb}{0.749020,0.772549,0.792157}%
\pgfsetfillcolor{currentfill}%
\pgfsetfillopacity{0.160000}%
\pgfsetlinewidth{0.000000pt}%
\definecolor{currentstroke}{rgb}{0.000000,0.000000,0.000000}%
\pgfsetstrokecolor{currentstroke}%
\pgfsetdash{}{0pt}%
\pgfpathmoveto{\pgfqpoint{1.090230in}{0.610533in}}%
\pgfpathlineto{\pgfqpoint{1.096808in}{0.616968in}}%
\pgfpathlineto{\pgfqpoint{1.103412in}{0.623428in}}%
\pgfpathlineto{\pgfqpoint{1.104895in}{0.623862in}}%
\pgfpathlineto{\pgfqpoint{1.106312in}{0.624336in}}%
\pgfpathlineto{\pgfqpoint{1.098256in}{0.617422in}}%
\pgfpathlineto{\pgfqpoint{1.090230in}{0.610533in}}%
\pgfpathlineto{\pgfqpoint{1.090230in}{0.610533in}}%
\pgfpathlineto{\pgfqpoint{1.090230in}{0.610533in}}%
\pgfpathclose%
\pgfusepath{fill}%
\end{pgfscope}%
\begin{pgfscope}%
\pgfpathrectangle{\pgfqpoint{0.407592in}{0.222904in}}{\pgfqpoint{1.411163in}{1.411163in}}%
\pgfusepath{clip}%
\pgfsetbuttcap%
\pgfsetroundjoin%
\definecolor{currentfill}{rgb}{0.749020,0.772549,0.792157}%
\pgfsetfillcolor{currentfill}%
\pgfsetfillopacity{0.160000}%
\pgfsetlinewidth{0.000000pt}%
\definecolor{currentstroke}{rgb}{0.000000,0.000000,0.000000}%
\pgfsetstrokecolor{currentstroke}%
\pgfsetdash{}{0pt}%
\pgfpathmoveto{\pgfqpoint{1.270742in}{0.789733in}}%
\pgfpathlineto{\pgfqpoint{1.283681in}{0.802577in}}%
\pgfpathlineto{\pgfqpoint{1.296626in}{0.815428in}}%
\pgfpathlineto{\pgfqpoint{1.292924in}{0.820977in}}%
\pgfpathlineto{\pgfqpoint{1.288355in}{0.826387in}}%
\pgfpathlineto{\pgfqpoint{1.275967in}{0.812890in}}%
\pgfpathlineto{\pgfqpoint{1.263579in}{0.799394in}}%
\pgfpathlineto{\pgfqpoint{1.267541in}{0.794623in}}%
\pgfpathlineto{\pgfqpoint{1.270742in}{0.789733in}}%
\pgfpathclose%
\pgfusepath{fill}%
\end{pgfscope}%
\begin{pgfscope}%
\pgfpathrectangle{\pgfqpoint{0.407592in}{0.222904in}}{\pgfqpoint{1.411163in}{1.411163in}}%
\pgfusepath{clip}%
\pgfsetbuttcap%
\pgfsetroundjoin%
\definecolor{currentfill}{rgb}{0.749020,0.772549,0.792157}%
\pgfsetfillcolor{currentfill}%
\pgfsetfillopacity{0.160000}%
\pgfsetlinewidth{0.000000pt}%
\definecolor{currentstroke}{rgb}{0.000000,0.000000,0.000000}%
\pgfsetstrokecolor{currentstroke}%
\pgfsetdash{}{0pt}%
\pgfpathmoveto{\pgfqpoint{0.945418in}{0.771587in}}%
\pgfpathlineto{\pgfqpoint{0.933343in}{0.785017in}}%
\pgfpathlineto{\pgfqpoint{0.921266in}{0.798448in}}%
\pgfpathlineto{\pgfqpoint{0.916437in}{0.793651in}}%
\pgfpathlineto{\pgfqpoint{0.912349in}{0.788739in}}%
\pgfpathlineto{\pgfqpoint{0.925097in}{0.775968in}}%
\pgfpathlineto{\pgfqpoint{0.937838in}{0.763203in}}%
\pgfpathlineto{\pgfqpoint{0.941310in}{0.767444in}}%
\pgfpathlineto{\pgfqpoint{0.945418in}{0.771587in}}%
\pgfpathclose%
\pgfusepath{fill}%
\end{pgfscope}%
\begin{pgfscope}%
\pgfpathrectangle{\pgfqpoint{0.407592in}{0.222904in}}{\pgfqpoint{1.411163in}{1.411163in}}%
\pgfusepath{clip}%
\pgfsetbuttcap%
\pgfsetroundjoin%
\definecolor{currentfill}{rgb}{0.749020,0.772549,0.792157}%
\pgfsetfillcolor{currentfill}%
\pgfsetfillopacity{0.160000}%
\pgfsetlinewidth{0.000000pt}%
\definecolor{currentstroke}{rgb}{0.000000,0.000000,0.000000}%
\pgfsetstrokecolor{currentstroke}%
\pgfsetdash{}{0pt}%
\pgfpathmoveto{\pgfqpoint{1.090230in}{0.610533in}}%
\pgfpathlineto{\pgfqpoint{1.082048in}{0.617139in}}%
\pgfpathlineto{\pgfqpoint{1.073835in}{0.623771in}}%
\pgfpathlineto{\pgfqpoint{1.075243in}{0.623346in}}%
\pgfpathlineto{\pgfqpoint{1.076716in}{0.622964in}}%
\pgfpathlineto{\pgfqpoint{1.083487in}{0.616735in}}%
\pgfpathlineto{\pgfqpoint{1.090230in}{0.610533in}}%
\pgfpathlineto{\pgfqpoint{1.090230in}{0.610533in}}%
\pgfpathlineto{\pgfqpoint{1.090230in}{0.610533in}}%
\pgfpathclose%
\pgfusepath{fill}%
\end{pgfscope}%
\begin{pgfscope}%
\pgfpathrectangle{\pgfqpoint{0.407592in}{0.222904in}}{\pgfqpoint{1.411163in}{1.411163in}}%
\pgfusepath{clip}%
\pgfsetbuttcap%
\pgfsetroundjoin%
\definecolor{currentfill}{rgb}{0.749020,0.772549,0.792157}%
\pgfsetfillcolor{currentfill}%
\pgfsetfillopacity{0.160000}%
\pgfsetlinewidth{0.000000pt}%
\definecolor{currentstroke}{rgb}{0.000000,0.000000,0.000000}%
\pgfsetstrokecolor{currentstroke}%
\pgfsetdash{}{0pt}%
\pgfpathmoveto{\pgfqpoint{1.090230in}{0.610533in}}%
\pgfpathlineto{\pgfqpoint{1.095244in}{0.616599in}}%
\pgfpathlineto{\pgfqpoint{1.100280in}{0.622690in}}%
\pgfpathlineto{\pgfqpoint{1.101872in}{0.623037in}}%
\pgfpathlineto{\pgfqpoint{1.103412in}{0.623428in}}%
\pgfpathlineto{\pgfqpoint{1.096808in}{0.616968in}}%
\pgfpathlineto{\pgfqpoint{1.090230in}{0.610533in}}%
\pgfpathlineto{\pgfqpoint{1.090230in}{0.610533in}}%
\pgfpathlineto{\pgfqpoint{1.090230in}{0.610533in}}%
\pgfpathclose%
\pgfusepath{fill}%
\end{pgfscope}%
\begin{pgfscope}%
\pgfpathrectangle{\pgfqpoint{0.407592in}{0.222904in}}{\pgfqpoint{1.411163in}{1.411163in}}%
\pgfusepath{clip}%
\pgfsetbuttcap%
\pgfsetroundjoin%
\definecolor{currentfill}{rgb}{0.749020,0.772549,0.792157}%
\pgfsetfillcolor{currentfill}%
\pgfsetfillopacity{0.160000}%
\pgfsetlinewidth{0.000000pt}%
\definecolor{currentstroke}{rgb}{0.000000,0.000000,0.000000}%
\pgfsetstrokecolor{currentstroke}%
\pgfsetdash{}{0pt}%
\pgfpathmoveto{\pgfqpoint{1.090230in}{0.610533in}}%
\pgfpathlineto{\pgfqpoint{1.083487in}{0.616735in}}%
\pgfpathlineto{\pgfqpoint{1.076716in}{0.622964in}}%
\pgfpathlineto{\pgfqpoint{1.078249in}{0.622626in}}%
\pgfpathlineto{\pgfqpoint{1.079835in}{0.622335in}}%
\pgfpathlineto{\pgfqpoint{1.085043in}{0.616421in}}%
\pgfpathlineto{\pgfqpoint{1.090230in}{0.610533in}}%
\pgfpathlineto{\pgfqpoint{1.090230in}{0.610533in}}%
\pgfpathlineto{\pgfqpoint{1.090230in}{0.610533in}}%
\pgfpathclose%
\pgfusepath{fill}%
\end{pgfscope}%
\begin{pgfscope}%
\pgfpathrectangle{\pgfqpoint{0.407592in}{0.222904in}}{\pgfqpoint{1.411163in}{1.411163in}}%
\pgfusepath{clip}%
\pgfsetbuttcap%
\pgfsetroundjoin%
\definecolor{currentfill}{rgb}{0.749020,0.772549,0.792157}%
\pgfsetfillcolor{currentfill}%
\pgfsetfillopacity{0.160000}%
\pgfsetlinewidth{0.000000pt}%
\definecolor{currentstroke}{rgb}{0.000000,0.000000,0.000000}%
\pgfsetstrokecolor{currentstroke}%
\pgfsetdash{}{0pt}%
\pgfpathmoveto{\pgfqpoint{1.090230in}{0.610533in}}%
\pgfpathlineto{\pgfqpoint{1.093592in}{0.616322in}}%
\pgfpathlineto{\pgfqpoint{1.096969in}{0.622136in}}%
\pgfpathlineto{\pgfqpoint{1.098643in}{0.622389in}}%
\pgfpathlineto{\pgfqpoint{1.100280in}{0.622690in}}%
\pgfpathlineto{\pgfqpoint{1.095244in}{0.616599in}}%
\pgfpathlineto{\pgfqpoint{1.090230in}{0.610533in}}%
\pgfpathlineto{\pgfqpoint{1.090230in}{0.610533in}}%
\pgfpathlineto{\pgfqpoint{1.090230in}{0.610533in}}%
\pgfpathclose%
\pgfusepath{fill}%
\end{pgfscope}%
\begin{pgfscope}%
\pgfpathrectangle{\pgfqpoint{0.407592in}{0.222904in}}{\pgfqpoint{1.411163in}{1.411163in}}%
\pgfusepath{clip}%
\pgfsetbuttcap%
\pgfsetroundjoin%
\definecolor{currentfill}{rgb}{0.749020,0.772549,0.792157}%
\pgfsetfillcolor{currentfill}%
\pgfsetfillopacity{0.160000}%
\pgfsetlinewidth{0.000000pt}%
\definecolor{currentstroke}{rgb}{0.000000,0.000000,0.000000}%
\pgfsetstrokecolor{currentstroke}%
\pgfsetdash{}{0pt}%
\pgfpathmoveto{\pgfqpoint{1.090230in}{0.610533in}}%
\pgfpathlineto{\pgfqpoint{1.085043in}{0.616421in}}%
\pgfpathlineto{\pgfqpoint{1.079835in}{0.622335in}}%
\pgfpathlineto{\pgfqpoint{1.081466in}{0.622091in}}%
\pgfpathlineto{\pgfqpoint{1.083135in}{0.621895in}}%
\pgfpathlineto{\pgfqpoint{1.086690in}{0.616201in}}%
\pgfpathlineto{\pgfqpoint{1.090230in}{0.610533in}}%
\pgfpathlineto{\pgfqpoint{1.090230in}{0.610533in}}%
\pgfpathlineto{\pgfqpoint{1.090230in}{0.610533in}}%
\pgfpathclose%
\pgfusepath{fill}%
\end{pgfscope}%
\begin{pgfscope}%
\pgfpathrectangle{\pgfqpoint{0.407592in}{0.222904in}}{\pgfqpoint{1.411163in}{1.411163in}}%
\pgfusepath{clip}%
\pgfsetbuttcap%
\pgfsetroundjoin%
\definecolor{currentfill}{rgb}{0.749020,0.772549,0.792157}%
\pgfsetfillcolor{currentfill}%
\pgfsetfillopacity{0.160000}%
\pgfsetlinewidth{0.000000pt}%
\definecolor{currentstroke}{rgb}{0.000000,0.000000,0.000000}%
\pgfsetstrokecolor{currentstroke}%
\pgfsetdash{}{0pt}%
\pgfpathmoveto{\pgfqpoint{1.090230in}{0.610533in}}%
\pgfpathlineto{\pgfqpoint{1.091880in}{0.616141in}}%
\pgfpathlineto{\pgfqpoint{1.093538in}{0.621774in}}%
\pgfpathlineto{\pgfqpoint{1.095265in}{0.621930in}}%
\pgfpathlineto{\pgfqpoint{1.096969in}{0.622136in}}%
\pgfpathlineto{\pgfqpoint{1.093592in}{0.616322in}}%
\pgfpathlineto{\pgfqpoint{1.090230in}{0.610533in}}%
\pgfpathlineto{\pgfqpoint{1.090230in}{0.610533in}}%
\pgfpathlineto{\pgfqpoint{1.090230in}{0.610533in}}%
\pgfpathclose%
\pgfusepath{fill}%
\end{pgfscope}%
\begin{pgfscope}%
\pgfpathrectangle{\pgfqpoint{0.407592in}{0.222904in}}{\pgfqpoint{1.411163in}{1.411163in}}%
\pgfusepath{clip}%
\pgfsetbuttcap%
\pgfsetroundjoin%
\definecolor{currentfill}{rgb}{0.749020,0.772549,0.792157}%
\pgfsetfillcolor{currentfill}%
\pgfsetfillopacity{0.160000}%
\pgfsetlinewidth{0.000000pt}%
\definecolor{currentstroke}{rgb}{0.000000,0.000000,0.000000}%
\pgfsetstrokecolor{currentstroke}%
\pgfsetdash{}{0pt}%
\pgfpathmoveto{\pgfqpoint{1.296626in}{0.815428in}}%
\pgfpathlineto{\pgfqpoint{1.309577in}{0.828285in}}%
\pgfpathlineto{\pgfqpoint{1.322534in}{0.841148in}}%
\pgfpathlineto{\pgfqpoint{1.318321in}{0.847344in}}%
\pgfpathlineto{\pgfqpoint{1.313134in}{0.853384in}}%
\pgfpathlineto{\pgfqpoint{1.300744in}{0.839885in}}%
\pgfpathlineto{\pgfqpoint{1.288355in}{0.826387in}}%
\pgfpathlineto{\pgfqpoint{1.292924in}{0.820977in}}%
\pgfpathlineto{\pgfqpoint{1.296626in}{0.815428in}}%
\pgfpathclose%
\pgfusepath{fill}%
\end{pgfscope}%
\begin{pgfscope}%
\pgfpathrectangle{\pgfqpoint{0.407592in}{0.222904in}}{\pgfqpoint{1.411163in}{1.411163in}}%
\pgfusepath{clip}%
\pgfsetbuttcap%
\pgfsetroundjoin%
\definecolor{currentfill}{rgb}{0.749020,0.772549,0.792157}%
\pgfsetfillcolor{currentfill}%
\pgfsetfillopacity{0.160000}%
\pgfsetlinewidth{0.000000pt}%
\definecolor{currentstroke}{rgb}{0.000000,0.000000,0.000000}%
\pgfsetstrokecolor{currentstroke}%
\pgfsetdash{}{0pt}%
\pgfpathmoveto{\pgfqpoint{1.090230in}{0.610533in}}%
\pgfpathlineto{\pgfqpoint{1.086690in}{0.616201in}}%
\pgfpathlineto{\pgfqpoint{1.083135in}{0.621895in}}%
\pgfpathlineto{\pgfqpoint{1.084836in}{0.621749in}}%
\pgfpathlineto{\pgfqpoint{1.086560in}{0.621653in}}%
\pgfpathlineto{\pgfqpoint{1.088399in}{0.616080in}}%
\pgfpathlineto{\pgfqpoint{1.090230in}{0.610533in}}%
\pgfpathlineto{\pgfqpoint{1.090230in}{0.610533in}}%
\pgfpathlineto{\pgfqpoint{1.090230in}{0.610533in}}%
\pgfpathclose%
\pgfusepath{fill}%
\end{pgfscope}%
\begin{pgfscope}%
\pgfpathrectangle{\pgfqpoint{0.407592in}{0.222904in}}{\pgfqpoint{1.411163in}{1.411163in}}%
\pgfusepath{clip}%
\pgfsetbuttcap%
\pgfsetroundjoin%
\definecolor{currentfill}{rgb}{0.749020,0.772549,0.792157}%
\pgfsetfillcolor{currentfill}%
\pgfsetfillopacity{0.160000}%
\pgfsetlinewidth{0.000000pt}%
\definecolor{currentstroke}{rgb}{0.000000,0.000000,0.000000}%
\pgfsetstrokecolor{currentstroke}%
\pgfsetdash{}{0pt}%
\pgfpathmoveto{\pgfqpoint{1.090230in}{0.610533in}}%
\pgfpathlineto{\pgfqpoint{1.090139in}{0.616060in}}%
\pgfpathlineto{\pgfqpoint{1.090048in}{0.621612in}}%
\pgfpathlineto{\pgfqpoint{1.091796in}{0.621668in}}%
\pgfpathlineto{\pgfqpoint{1.093538in}{0.621774in}}%
\pgfpathlineto{\pgfqpoint{1.091880in}{0.616141in}}%
\pgfpathlineto{\pgfqpoint{1.090230in}{0.610533in}}%
\pgfpathlineto{\pgfqpoint{1.090230in}{0.610533in}}%
\pgfpathlineto{\pgfqpoint{1.090230in}{0.610533in}}%
\pgfpathclose%
\pgfusepath{fill}%
\end{pgfscope}%
\begin{pgfscope}%
\pgfpathrectangle{\pgfqpoint{0.407592in}{0.222904in}}{\pgfqpoint{1.411163in}{1.411163in}}%
\pgfusepath{clip}%
\pgfsetbuttcap%
\pgfsetroundjoin%
\definecolor{currentfill}{rgb}{0.749020,0.772549,0.792157}%
\pgfsetfillcolor{currentfill}%
\pgfsetfillopacity{0.160000}%
\pgfsetlinewidth{0.000000pt}%
\definecolor{currentstroke}{rgb}{0.000000,0.000000,0.000000}%
\pgfsetstrokecolor{currentstroke}%
\pgfsetdash{}{0pt}%
\pgfpathmoveto{\pgfqpoint{1.090230in}{0.610533in}}%
\pgfpathlineto{\pgfqpoint{1.088399in}{0.616080in}}%
\pgfpathlineto{\pgfqpoint{1.086560in}{0.621653in}}%
\pgfpathlineto{\pgfqpoint{1.088300in}{0.621607in}}%
\pgfpathlineto{\pgfqpoint{1.090048in}{0.621612in}}%
\pgfpathlineto{\pgfqpoint{1.090139in}{0.616060in}}%
\pgfpathlineto{\pgfqpoint{1.090230in}{0.610533in}}%
\pgfpathlineto{\pgfqpoint{1.090230in}{0.610533in}}%
\pgfpathlineto{\pgfqpoint{1.090230in}{0.610533in}}%
\pgfpathclose%
\pgfusepath{fill}%
\end{pgfscope}%
\begin{pgfscope}%
\pgfpathrectangle{\pgfqpoint{0.407592in}{0.222904in}}{\pgfqpoint{1.411163in}{1.411163in}}%
\pgfusepath{clip}%
\pgfsetbuttcap%
\pgfsetroundjoin%
\definecolor{currentfill}{rgb}{0.749020,0.772549,0.792157}%
\pgfsetfillcolor{currentfill}%
\pgfsetfillopacity{0.160000}%
\pgfsetlinewidth{0.000000pt}%
\definecolor{currentstroke}{rgb}{0.000000,0.000000,0.000000}%
\pgfsetstrokecolor{currentstroke}%
\pgfsetdash{}{0pt}%
\pgfpathmoveto{\pgfqpoint{0.921266in}{0.798448in}}%
\pgfpathlineto{\pgfqpoint{0.909188in}{0.811880in}}%
\pgfpathlineto{\pgfqpoint{0.897109in}{0.825314in}}%
\pgfpathlineto{\pgfqpoint{0.891550in}{0.819874in}}%
\pgfpathlineto{\pgfqpoint{0.886834in}{0.814300in}}%
\pgfpathlineto{\pgfqpoint{0.899595in}{0.801516in}}%
\pgfpathlineto{\pgfqpoint{0.912349in}{0.788739in}}%
\pgfpathlineto{\pgfqpoint{0.916437in}{0.793651in}}%
\pgfpathlineto{\pgfqpoint{0.921266in}{0.798448in}}%
\pgfpathclose%
\pgfusepath{fill}%
\end{pgfscope}%
\begin{pgfscope}%
\pgfpathrectangle{\pgfqpoint{0.407592in}{0.222904in}}{\pgfqpoint{1.411163in}{1.411163in}}%
\pgfusepath{clip}%
\pgfsetbuttcap%
\pgfsetroundjoin%
\definecolor{currentfill}{rgb}{0.749020,0.772549,0.792157}%
\pgfsetfillcolor{currentfill}%
\pgfsetfillopacity{0.160000}%
\pgfsetlinewidth{0.000000pt}%
\definecolor{currentstroke}{rgb}{0.000000,0.000000,0.000000}%
\pgfsetstrokecolor{currentstroke}%
\pgfsetdash{}{0pt}%
\pgfpathmoveto{\pgfqpoint{1.116100in}{0.631528in}}%
\pgfpathlineto{\pgfqpoint{1.129068in}{0.642053in}}%
\pgfpathlineto{\pgfqpoint{1.142060in}{0.652597in}}%
\pgfpathlineto{\pgfqpoint{1.142573in}{0.654106in}}%
\pgfpathlineto{\pgfqpoint{1.142855in}{0.655621in}}%
\pgfpathlineto{\pgfqpoint{1.129673in}{0.644327in}}%
\pgfpathlineto{\pgfqpoint{1.116508in}{0.633047in}}%
\pgfpathlineto{\pgfqpoint{1.116361in}{0.632286in}}%
\pgfpathlineto{\pgfqpoint{1.116100in}{0.631528in}}%
\pgfpathclose%
\pgfusepath{fill}%
\end{pgfscope}%
\begin{pgfscope}%
\pgfpathrectangle{\pgfqpoint{0.407592in}{0.222904in}}{\pgfqpoint{1.411163in}{1.411163in}}%
\pgfusepath{clip}%
\pgfsetbuttcap%
\pgfsetroundjoin%
\definecolor{currentfill}{rgb}{0.749020,0.772549,0.792157}%
\pgfsetfillcolor{currentfill}%
\pgfsetfillopacity{0.160000}%
\pgfsetlinewidth{0.000000pt}%
\definecolor{currentstroke}{rgb}{0.000000,0.000000,0.000000}%
\pgfsetstrokecolor{currentstroke}%
\pgfsetdash{}{0pt}%
\pgfpathmoveto{\pgfqpoint{1.063809in}{0.632895in}}%
\pgfpathlineto{\pgfqpoint{1.050572in}{0.644098in}}%
\pgfpathlineto{\pgfqpoint{1.037317in}{0.655317in}}%
\pgfpathlineto{\pgfqpoint{1.037335in}{0.653802in}}%
\pgfpathlineto{\pgfqpoint{1.037585in}{0.652295in}}%
\pgfpathlineto{\pgfqpoint{1.050781in}{0.641827in}}%
\pgfpathlineto{\pgfqpoint{1.063954in}{0.631377in}}%
\pgfpathlineto{\pgfqpoint{1.063824in}{0.632134in}}%
\pgfpathlineto{\pgfqpoint{1.063809in}{0.632895in}}%
\pgfpathclose%
\pgfusepath{fill}%
\end{pgfscope}%
\begin{pgfscope}%
\pgfpathrectangle{\pgfqpoint{0.407592in}{0.222904in}}{\pgfqpoint{1.411163in}{1.411163in}}%
\pgfusepath{clip}%
\pgfsetbuttcap%
\pgfsetroundjoin%
\definecolor{currentfill}{rgb}{0.749020,0.772549,0.792157}%
\pgfsetfillcolor{currentfill}%
\pgfsetfillopacity{0.160000}%
\pgfsetlinewidth{0.000000pt}%
\definecolor{currentstroke}{rgb}{0.000000,0.000000,0.000000}%
\pgfsetstrokecolor{currentstroke}%
\pgfsetdash{}{0pt}%
\pgfpathmoveto{\pgfqpoint{1.169048in}{0.682764in}}%
\pgfpathlineto{\pgfqpoint{1.182225in}{0.694840in}}%
\pgfpathlineto{\pgfqpoint{1.195414in}{0.706927in}}%
\pgfpathlineto{\pgfqpoint{1.194552in}{0.709879in}}%
\pgfpathlineto{\pgfqpoint{1.193235in}{0.712789in}}%
\pgfpathlineto{\pgfqpoint{1.180339in}{0.699986in}}%
\pgfpathlineto{\pgfqpoint{1.167448in}{0.687189in}}%
\pgfpathlineto{\pgfqpoint{1.168418in}{0.684992in}}%
\pgfpathlineto{\pgfqpoint{1.169048in}{0.682764in}}%
\pgfpathclose%
\pgfusepath{fill}%
\end{pgfscope}%
\begin{pgfscope}%
\pgfpathrectangle{\pgfqpoint{0.407592in}{0.222904in}}{\pgfqpoint{1.411163in}{1.411163in}}%
\pgfusepath{clip}%
\pgfsetbuttcap%
\pgfsetroundjoin%
\definecolor{currentfill}{rgb}{0.749020,0.772549,0.792157}%
\pgfsetfillcolor{currentfill}%
\pgfsetfillopacity{0.160000}%
\pgfsetlinewidth{0.000000pt}%
\definecolor{currentstroke}{rgb}{0.000000,0.000000,0.000000}%
\pgfsetstrokecolor{currentstroke}%
\pgfsetdash{}{0pt}%
\pgfpathmoveto{\pgfqpoint{1.322534in}{0.841148in}}%
\pgfpathlineto{\pgfqpoint{1.335497in}{0.854017in}}%
\pgfpathlineto{\pgfqpoint{1.348467in}{0.866892in}}%
\pgfpathlineto{\pgfqpoint{1.343730in}{0.873725in}}%
\pgfpathlineto{\pgfqpoint{1.337916in}{0.880383in}}%
\pgfpathlineto{\pgfqpoint{1.325524in}{0.866883in}}%
\pgfpathlineto{\pgfqpoint{1.313134in}{0.853384in}}%
\pgfpathlineto{\pgfqpoint{1.318321in}{0.847344in}}%
\pgfpathlineto{\pgfqpoint{1.322534in}{0.841148in}}%
\pgfpathclose%
\pgfusepath{fill}%
\end{pgfscope}%
\begin{pgfscope}%
\pgfpathrectangle{\pgfqpoint{0.407592in}{0.222904in}}{\pgfqpoint{1.411163in}{1.411163in}}%
\pgfusepath{clip}%
\pgfsetbuttcap%
\pgfsetroundjoin%
\definecolor{currentfill}{rgb}{0.749020,0.772549,0.792157}%
\pgfsetfillcolor{currentfill}%
\pgfsetfillopacity{0.160000}%
\pgfsetlinewidth{0.000000pt}%
\definecolor{currentstroke}{rgb}{0.000000,0.000000,0.000000}%
\pgfsetstrokecolor{currentstroke}%
\pgfsetdash{}{0pt}%
\pgfpathmoveto{\pgfqpoint{1.014150in}{0.686752in}}%
\pgfpathlineto{\pgfqpoint{1.001448in}{0.699477in}}%
\pgfpathlineto{\pgfqpoint{0.988739in}{0.712209in}}%
\pgfpathlineto{\pgfqpoint{0.986906in}{0.709290in}}%
\pgfpathlineto{\pgfqpoint{0.985520in}{0.706331in}}%
\pgfpathlineto{\pgfqpoint{0.998651in}{0.694317in}}%
\pgfpathlineto{\pgfqpoint{1.011770in}{0.682314in}}%
\pgfpathlineto{\pgfqpoint{1.012792in}{0.684547in}}%
\pgfpathlineto{\pgfqpoint{1.014150in}{0.686752in}}%
\pgfpathclose%
\pgfusepath{fill}%
\end{pgfscope}%
\begin{pgfscope}%
\pgfpathrectangle{\pgfqpoint{0.407592in}{0.222904in}}{\pgfqpoint{1.411163in}{1.411163in}}%
\pgfusepath{clip}%
\pgfsetbuttcap%
\pgfsetroundjoin%
\definecolor{currentfill}{rgb}{0.749020,0.772549,0.792157}%
\pgfsetfillcolor{currentfill}%
\pgfsetfillopacity{0.160000}%
\pgfsetlinewidth{0.000000pt}%
\definecolor{currentstroke}{rgb}{0.000000,0.000000,0.000000}%
\pgfsetstrokecolor{currentstroke}%
\pgfsetdash{}{0pt}%
\pgfpathmoveto{\pgfqpoint{0.897109in}{0.825314in}}%
\pgfpathlineto{\pgfqpoint{0.885029in}{0.838749in}}%
\pgfpathlineto{\pgfqpoint{0.872948in}{0.852186in}}%
\pgfpathlineto{\pgfqpoint{0.866648in}{0.846113in}}%
\pgfpathlineto{\pgfqpoint{0.861293in}{0.839888in}}%
\pgfpathlineto{\pgfqpoint{0.874067in}{0.827091in}}%
\pgfpathlineto{\pgfqpoint{0.886834in}{0.814300in}}%
\pgfpathlineto{\pgfqpoint{0.891550in}{0.819874in}}%
\pgfpathlineto{\pgfqpoint{0.897109in}{0.825314in}}%
\pgfpathclose%
\pgfusepath{fill}%
\end{pgfscope}%
\begin{pgfscope}%
\pgfpathrectangle{\pgfqpoint{0.407592in}{0.222904in}}{\pgfqpoint{1.411163in}{1.411163in}}%
\pgfusepath{clip}%
\pgfsetbuttcap%
\pgfsetroundjoin%
\definecolor{currentfill}{rgb}{0.749020,0.772549,0.792157}%
\pgfsetfillcolor{currentfill}%
\pgfsetfillopacity{0.160000}%
\pgfsetlinewidth{0.000000pt}%
\definecolor{currentstroke}{rgb}{0.000000,0.000000,0.000000}%
\pgfsetstrokecolor{currentstroke}%
\pgfsetdash{}{0pt}%
\pgfpathmoveto{\pgfqpoint{1.142855in}{0.655621in}}%
\pgfpathlineto{\pgfqpoint{1.156055in}{0.666930in}}%
\pgfpathlineto{\pgfqpoint{1.169272in}{0.678255in}}%
\pgfpathlineto{\pgfqpoint{1.169333in}{0.680515in}}%
\pgfpathlineto{\pgfqpoint{1.169048in}{0.682764in}}%
\pgfpathlineto{\pgfqpoint{1.155882in}{0.670699in}}%
\pgfpathlineto{\pgfqpoint{1.142729in}{0.658644in}}%
\pgfpathlineto{\pgfqpoint{1.142907in}{0.657136in}}%
\pgfpathlineto{\pgfqpoint{1.142855in}{0.655621in}}%
\pgfpathclose%
\pgfusepath{fill}%
\end{pgfscope}%
\begin{pgfscope}%
\pgfpathrectangle{\pgfqpoint{0.407592in}{0.222904in}}{\pgfqpoint{1.411163in}{1.411163in}}%
\pgfusepath{clip}%
\pgfsetbuttcap%
\pgfsetroundjoin%
\definecolor{currentfill}{rgb}{0.749020,0.772549,0.792157}%
\pgfsetfillcolor{currentfill}%
\pgfsetfillopacity{0.160000}%
\pgfsetlinewidth{0.000000pt}%
\definecolor{currentstroke}{rgb}{0.000000,0.000000,0.000000}%
\pgfsetstrokecolor{currentstroke}%
\pgfsetdash{}{0pt}%
\pgfpathmoveto{\pgfqpoint{1.348467in}{0.866892in}}%
\pgfpathlineto{\pgfqpoint{1.361442in}{0.879773in}}%
\pgfpathlineto{\pgfqpoint{1.374423in}{0.892660in}}%
\pgfpathlineto{\pgfqpoint{1.369153in}{0.900120in}}%
\pgfpathlineto{\pgfqpoint{1.362701in}{0.907386in}}%
\pgfpathlineto{\pgfqpoint{1.350308in}{0.893884in}}%
\pgfpathlineto{\pgfqpoint{1.337916in}{0.880383in}}%
\pgfpathlineto{\pgfqpoint{1.343730in}{0.873725in}}%
\pgfpathlineto{\pgfqpoint{1.348467in}{0.866892in}}%
\pgfpathclose%
\pgfusepath{fill}%
\end{pgfscope}%
\begin{pgfscope}%
\pgfpathrectangle{\pgfqpoint{0.407592in}{0.222904in}}{\pgfqpoint{1.411163in}{1.411163in}}%
\pgfusepath{clip}%
\pgfsetbuttcap%
\pgfsetroundjoin%
\definecolor{currentfill}{rgb}{0.749020,0.772549,0.792157}%
\pgfsetfillcolor{currentfill}%
\pgfsetfillopacity{0.160000}%
\pgfsetlinewidth{0.000000pt}%
\definecolor{currentstroke}{rgb}{0.000000,0.000000,0.000000}%
\pgfsetstrokecolor{currentstroke}%
\pgfsetdash{}{0pt}%
\pgfpathmoveto{\pgfqpoint{1.037972in}{0.658343in}}%
\pgfpathlineto{\pgfqpoint{1.024877in}{0.670323in}}%
\pgfpathlineto{\pgfqpoint{1.011770in}{0.682314in}}%
\pgfpathlineto{\pgfqpoint{1.011089in}{0.680062in}}%
\pgfpathlineto{\pgfqpoint{1.010753in}{0.677801in}}%
\pgfpathlineto{\pgfqpoint{1.024044in}{0.666551in}}%
\pgfpathlineto{\pgfqpoint{1.037317in}{0.655317in}}%
\pgfpathlineto{\pgfqpoint{1.037530in}{0.656833in}}%
\pgfpathlineto{\pgfqpoint{1.037972in}{0.658343in}}%
\pgfpathclose%
\pgfusepath{fill}%
\end{pgfscope}%
\begin{pgfscope}%
\pgfpathrectangle{\pgfqpoint{0.407592in}{0.222904in}}{\pgfqpoint{1.411163in}{1.411163in}}%
\pgfusepath{clip}%
\pgfsetbuttcap%
\pgfsetroundjoin%
\definecolor{currentfill}{rgb}{0.749020,0.772549,0.792157}%
\pgfsetfillcolor{currentfill}%
\pgfsetfillopacity{0.160000}%
\pgfsetlinewidth{0.000000pt}%
\definecolor{currentstroke}{rgb}{0.000000,0.000000,0.000000}%
\pgfsetstrokecolor{currentstroke}%
\pgfsetdash{}{0pt}%
\pgfpathmoveto{\pgfqpoint{0.872948in}{0.852186in}}%
\pgfpathlineto{\pgfqpoint{0.860865in}{0.865624in}}%
\pgfpathlineto{\pgfqpoint{0.848781in}{0.879063in}}%
\pgfpathlineto{\pgfqpoint{0.841730in}{0.872367in}}%
\pgfpathlineto{\pgfqpoint{0.835726in}{0.865502in}}%
\pgfpathlineto{\pgfqpoint{0.848513in}{0.852692in}}%
\pgfpathlineto{\pgfqpoint{0.861293in}{0.839888in}}%
\pgfpathlineto{\pgfqpoint{0.866648in}{0.846113in}}%
\pgfpathlineto{\pgfqpoint{0.872948in}{0.852186in}}%
\pgfpathclose%
\pgfusepath{fill}%
\end{pgfscope}%
\begin{pgfscope}%
\pgfpathrectangle{\pgfqpoint{0.407592in}{0.222904in}}{\pgfqpoint{1.411163in}{1.411163in}}%
\pgfusepath{clip}%
\pgfsetbuttcap%
\pgfsetroundjoin%
\definecolor{currentfill}{rgb}{0.749020,0.772549,0.792157}%
\pgfsetfillcolor{currentfill}%
\pgfsetfillopacity{0.160000}%
\pgfsetlinewidth{0.000000pt}%
\definecolor{currentstroke}{rgb}{0.000000,0.000000,0.000000}%
\pgfsetstrokecolor{currentstroke}%
\pgfsetdash{}{0pt}%
\pgfpathmoveto{\pgfqpoint{1.115238in}{0.630036in}}%
\pgfpathlineto{\pgfqpoint{1.127782in}{0.639820in}}%
\pgfpathlineto{\pgfqpoint{1.140354in}{0.649624in}}%
\pgfpathlineto{\pgfqpoint{1.141320in}{0.651101in}}%
\pgfpathlineto{\pgfqpoint{1.142060in}{0.652597in}}%
\pgfpathlineto{\pgfqpoint{1.129068in}{0.642053in}}%
\pgfpathlineto{\pgfqpoint{1.116100in}{0.631528in}}%
\pgfpathlineto{\pgfqpoint{1.115725in}{0.630777in}}%
\pgfpathlineto{\pgfqpoint{1.115238in}{0.630036in}}%
\pgfpathclose%
\pgfusepath{fill}%
\end{pgfscope}%
\begin{pgfscope}%
\pgfpathrectangle{\pgfqpoint{0.407592in}{0.222904in}}{\pgfqpoint{1.411163in}{1.411163in}}%
\pgfusepath{clip}%
\pgfsetbuttcap%
\pgfsetroundjoin%
\definecolor{currentfill}{rgb}{0.749020,0.772549,0.792157}%
\pgfsetfillcolor{currentfill}%
\pgfsetfillopacity{0.160000}%
\pgfsetlinewidth{0.000000pt}%
\definecolor{currentstroke}{rgb}{0.000000,0.000000,0.000000}%
\pgfsetstrokecolor{currentstroke}%
\pgfsetdash{}{0pt}%
\pgfpathmoveto{\pgfqpoint{1.063954in}{0.631377in}}%
\pgfpathlineto{\pgfqpoint{1.050781in}{0.641827in}}%
\pgfpathlineto{\pgfqpoint{1.037585in}{0.652295in}}%
\pgfpathlineto{\pgfqpoint{1.038064in}{0.650803in}}%
\pgfpathlineto{\pgfqpoint{1.038772in}{0.649331in}}%
\pgfpathlineto{\pgfqpoint{1.051679in}{0.639599in}}%
\pgfpathlineto{\pgfqpoint{1.064558in}{0.629889in}}%
\pgfpathlineto{\pgfqpoint{1.064199in}{0.630628in}}%
\pgfpathlineto{\pgfqpoint{1.063954in}{0.631377in}}%
\pgfpathclose%
\pgfusepath{fill}%
\end{pgfscope}%
\begin{pgfscope}%
\pgfpathrectangle{\pgfqpoint{0.407592in}{0.222904in}}{\pgfqpoint{1.411163in}{1.411163in}}%
\pgfusepath{clip}%
\pgfsetbuttcap%
\pgfsetroundjoin%
\definecolor{currentfill}{rgb}{0.749020,0.772549,0.792157}%
\pgfsetfillcolor{currentfill}%
\pgfsetfillopacity{0.160000}%
\pgfsetlinewidth{0.000000pt}%
\definecolor{currentstroke}{rgb}{0.000000,0.000000,0.000000}%
\pgfsetstrokecolor{currentstroke}%
\pgfsetdash{}{0pt}%
\pgfpathmoveto{\pgfqpoint{1.374423in}{0.892660in}}%
\pgfpathlineto{\pgfqpoint{1.387411in}{0.905553in}}%
\pgfpathlineto{\pgfqpoint{1.400404in}{0.918452in}}%
\pgfpathlineto{\pgfqpoint{1.394589in}{0.926529in}}%
\pgfpathlineto{\pgfqpoint{1.387488in}{0.934391in}}%
\pgfpathlineto{\pgfqpoint{1.375094in}{0.920888in}}%
\pgfpathlineto{\pgfqpoint{1.362701in}{0.907386in}}%
\pgfpathlineto{\pgfqpoint{1.369153in}{0.900120in}}%
\pgfpathlineto{\pgfqpoint{1.374423in}{0.892660in}}%
\pgfpathclose%
\pgfusepath{fill}%
\end{pgfscope}%
\begin{pgfscope}%
\pgfpathrectangle{\pgfqpoint{0.407592in}{0.222904in}}{\pgfqpoint{1.411163in}{1.411163in}}%
\pgfusepath{clip}%
\pgfsetbuttcap%
\pgfsetroundjoin%
\definecolor{currentfill}{rgb}{0.749020,0.772549,0.792157}%
\pgfsetfillcolor{currentfill}%
\pgfsetfillopacity{0.160000}%
\pgfsetlinewidth{0.000000pt}%
\definecolor{currentstroke}{rgb}{0.000000,0.000000,0.000000}%
\pgfsetstrokecolor{currentstroke}%
\pgfsetdash{}{0pt}%
\pgfpathmoveto{\pgfqpoint{1.195414in}{0.706927in}}%
\pgfpathlineto{\pgfqpoint{1.208615in}{0.719024in}}%
\pgfpathlineto{\pgfqpoint{1.221828in}{0.731133in}}%
\pgfpathlineto{\pgfqpoint{1.220721in}{0.734800in}}%
\pgfpathlineto{\pgfqpoint{1.219047in}{0.738413in}}%
\pgfpathlineto{\pgfqpoint{1.206138in}{0.725598in}}%
\pgfpathlineto{\pgfqpoint{1.193235in}{0.712789in}}%
\pgfpathlineto{\pgfqpoint{1.194552in}{0.709879in}}%
\pgfpathlineto{\pgfqpoint{1.195414in}{0.706927in}}%
\pgfpathclose%
\pgfusepath{fill}%
\end{pgfscope}%
\begin{pgfscope}%
\pgfpathrectangle{\pgfqpoint{0.407592in}{0.222904in}}{\pgfqpoint{1.411163in}{1.411163in}}%
\pgfusepath{clip}%
\pgfsetbuttcap%
\pgfsetroundjoin%
\definecolor{currentfill}{rgb}{0.749020,0.772549,0.792157}%
\pgfsetfillcolor{currentfill}%
\pgfsetfillopacity{0.160000}%
\pgfsetlinewidth{0.000000pt}%
\definecolor{currentstroke}{rgb}{0.000000,0.000000,0.000000}%
\pgfsetstrokecolor{currentstroke}%
\pgfsetdash{}{0pt}%
\pgfpathmoveto{\pgfqpoint{0.848781in}{0.879063in}}%
\pgfpathlineto{\pgfqpoint{0.836696in}{0.892503in}}%
\pgfpathlineto{\pgfqpoint{0.824610in}{0.905945in}}%
\pgfpathlineto{\pgfqpoint{0.816798in}{0.898638in}}%
\pgfpathlineto{\pgfqpoint{0.810133in}{0.891142in}}%
\pgfpathlineto{\pgfqpoint{0.822932in}{0.878319in}}%
\pgfpathlineto{\pgfqpoint{0.835726in}{0.865502in}}%
\pgfpathlineto{\pgfqpoint{0.841730in}{0.872367in}}%
\pgfpathlineto{\pgfqpoint{0.848781in}{0.879063in}}%
\pgfpathclose%
\pgfusepath{fill}%
\end{pgfscope}%
\begin{pgfscope}%
\pgfpathrectangle{\pgfqpoint{0.407592in}{0.222904in}}{\pgfqpoint{1.411163in}{1.411163in}}%
\pgfusepath{clip}%
\pgfsetbuttcap%
\pgfsetroundjoin%
\definecolor{currentfill}{rgb}{0.749020,0.772549,0.792157}%
\pgfsetfillcolor{currentfill}%
\pgfsetfillopacity{0.160000}%
\pgfsetlinewidth{0.000000pt}%
\definecolor{currentstroke}{rgb}{0.000000,0.000000,0.000000}%
\pgfsetstrokecolor{currentstroke}%
\pgfsetdash{}{0pt}%
\pgfpathmoveto{\pgfqpoint{0.988739in}{0.712209in}}%
\pgfpathlineto{\pgfqpoint{0.976023in}{0.724948in}}%
\pgfpathlineto{\pgfqpoint{0.963301in}{0.737693in}}%
\pgfpathlineto{\pgfqpoint{0.960982in}{0.734068in}}%
\pgfpathlineto{\pgfqpoint{0.959220in}{0.730392in}}%
\pgfpathlineto{\pgfqpoint{0.972376in}{0.718356in}}%
\pgfpathlineto{\pgfqpoint{0.985520in}{0.706331in}}%
\pgfpathlineto{\pgfqpoint{0.986906in}{0.709290in}}%
\pgfpathlineto{\pgfqpoint{0.988739in}{0.712209in}}%
\pgfpathclose%
\pgfusepath{fill}%
\end{pgfscope}%
\begin{pgfscope}%
\pgfpathrectangle{\pgfqpoint{0.407592in}{0.222904in}}{\pgfqpoint{1.411163in}{1.411163in}}%
\pgfusepath{clip}%
\pgfsetbuttcap%
\pgfsetroundjoin%
\definecolor{currentfill}{rgb}{0.749020,0.772549,0.792157}%
\pgfsetfillcolor{currentfill}%
\pgfsetfillopacity{0.160000}%
\pgfsetlinewidth{0.000000pt}%
\definecolor{currentstroke}{rgb}{0.000000,0.000000,0.000000}%
\pgfsetstrokecolor{currentstroke}%
\pgfsetdash{}{0pt}%
\pgfpathmoveto{\pgfqpoint{1.400404in}{0.918452in}}%
\pgfpathlineto{\pgfqpoint{1.413404in}{0.931357in}}%
\pgfpathlineto{\pgfqpoint{1.426410in}{0.944268in}}%
\pgfpathlineto{\pgfqpoint{1.420039in}{0.952951in}}%
\pgfpathlineto{\pgfqpoint{1.412279in}{0.961401in}}%
\pgfpathlineto{\pgfqpoint{1.399883in}{0.947896in}}%
\pgfpathlineto{\pgfqpoint{1.387488in}{0.934391in}}%
\pgfpathlineto{\pgfqpoint{1.394589in}{0.926529in}}%
\pgfpathlineto{\pgfqpoint{1.400404in}{0.918452in}}%
\pgfpathclose%
\pgfusepath{fill}%
\end{pgfscope}%
\begin{pgfscope}%
\pgfpathrectangle{\pgfqpoint{0.407592in}{0.222904in}}{\pgfqpoint{1.411163in}{1.411163in}}%
\pgfusepath{clip}%
\pgfsetbuttcap%
\pgfsetroundjoin%
\definecolor{currentfill}{rgb}{0.749020,0.772549,0.792157}%
\pgfsetfillcolor{currentfill}%
\pgfsetfillopacity{0.160000}%
\pgfsetlinewidth{0.000000pt}%
\definecolor{currentstroke}{rgb}{0.000000,0.000000,0.000000}%
\pgfsetstrokecolor{currentstroke}%
\pgfsetdash{}{0pt}%
\pgfpathmoveto{\pgfqpoint{0.824610in}{0.905945in}}%
\pgfpathlineto{\pgfqpoint{0.812523in}{0.919388in}}%
\pgfpathlineto{\pgfqpoint{0.800434in}{0.932833in}}%
\pgfpathlineto{\pgfqpoint{0.791850in}{0.924924in}}%
\pgfpathlineto{\pgfqpoint{0.784513in}{0.916808in}}%
\pgfpathlineto{\pgfqpoint{0.797326in}{0.903972in}}%
\pgfpathlineto{\pgfqpoint{0.810133in}{0.891142in}}%
\pgfpathlineto{\pgfqpoint{0.816798in}{0.898638in}}%
\pgfpathlineto{\pgfqpoint{0.824610in}{0.905945in}}%
\pgfpathclose%
\pgfusepath{fill}%
\end{pgfscope}%
\begin{pgfscope}%
\pgfpathrectangle{\pgfqpoint{0.407592in}{0.222904in}}{\pgfqpoint{1.411163in}{1.411163in}}%
\pgfusepath{clip}%
\pgfsetbuttcap%
\pgfsetroundjoin%
\definecolor{currentfill}{rgb}{0.749020,0.772549,0.792157}%
\pgfsetfillcolor{currentfill}%
\pgfsetfillopacity{0.160000}%
\pgfsetlinewidth{0.000000pt}%
\definecolor{currentstroke}{rgb}{0.000000,0.000000,0.000000}%
\pgfsetstrokecolor{currentstroke}%
\pgfsetdash{}{0pt}%
\pgfpathmoveto{\pgfqpoint{1.426410in}{0.944268in}}%
\pgfpathlineto{\pgfqpoint{1.439422in}{0.957185in}}%
\pgfpathlineto{\pgfqpoint{1.452439in}{0.970108in}}%
\pgfpathlineto{\pgfqpoint{1.445501in}{0.979387in}}%
\pgfpathlineto{\pgfqpoint{1.437073in}{0.988413in}}%
\pgfpathlineto{\pgfqpoint{1.424676in}{0.974906in}}%
\pgfpathlineto{\pgfqpoint{1.412279in}{0.961401in}}%
\pgfpathlineto{\pgfqpoint{1.420039in}{0.952951in}}%
\pgfpathlineto{\pgfqpoint{1.426410in}{0.944268in}}%
\pgfpathclose%
\pgfusepath{fill}%
\end{pgfscope}%
\begin{pgfscope}%
\pgfpathrectangle{\pgfqpoint{0.407592in}{0.222904in}}{\pgfqpoint{1.411163in}{1.411163in}}%
\pgfusepath{clip}%
\pgfsetbuttcap%
\pgfsetroundjoin%
\definecolor{currentfill}{rgb}{0.749020,0.772549,0.792157}%
\pgfsetfillcolor{currentfill}%
\pgfsetfillopacity{0.160000}%
\pgfsetlinewidth{0.000000pt}%
\definecolor{currentstroke}{rgb}{0.000000,0.000000,0.000000}%
\pgfsetstrokecolor{currentstroke}%
\pgfsetdash{}{0pt}%
\pgfpathmoveto{\pgfqpoint{1.113937in}{0.628598in}}%
\pgfpathlineto{\pgfqpoint{1.125836in}{0.637665in}}%
\pgfpathlineto{\pgfqpoint{1.137766in}{0.646756in}}%
\pgfpathlineto{\pgfqpoint{1.139168in}{0.648174in}}%
\pgfpathlineto{\pgfqpoint{1.140354in}{0.649624in}}%
\pgfpathlineto{\pgfqpoint{1.127782in}{0.639820in}}%
\pgfpathlineto{\pgfqpoint{1.115238in}{0.630036in}}%
\pgfpathlineto{\pgfqpoint{1.114641in}{0.629309in}}%
\pgfpathlineto{\pgfqpoint{1.113937in}{0.628598in}}%
\pgfpathclose%
\pgfusepath{fill}%
\end{pgfscope}%
\begin{pgfscope}%
\pgfpathrectangle{\pgfqpoint{0.407592in}{0.222904in}}{\pgfqpoint{1.411163in}{1.411163in}}%
\pgfusepath{clip}%
\pgfsetbuttcap%
\pgfsetroundjoin%
\definecolor{currentfill}{rgb}{0.749020,0.772549,0.792157}%
\pgfsetfillcolor{currentfill}%
\pgfsetfillopacity{0.160000}%
\pgfsetlinewidth{0.000000pt}%
\definecolor{currentstroke}{rgb}{0.000000,0.000000,0.000000}%
\pgfsetstrokecolor{currentstroke}%
\pgfsetdash{}{0pt}%
\pgfpathmoveto{\pgfqpoint{1.064558in}{0.629889in}}%
\pgfpathlineto{\pgfqpoint{1.051679in}{0.639599in}}%
\pgfpathlineto{\pgfqpoint{1.038772in}{0.649331in}}%
\pgfpathlineto{\pgfqpoint{1.039706in}{0.647887in}}%
\pgfpathlineto{\pgfqpoint{1.040861in}{0.646476in}}%
\pgfpathlineto{\pgfqpoint{1.053252in}{0.637455in}}%
\pgfpathlineto{\pgfqpoint{1.065610in}{0.628457in}}%
\pgfpathlineto{\pgfqpoint{1.065029in}{0.629165in}}%
\pgfpathlineto{\pgfqpoint{1.064558in}{0.629889in}}%
\pgfpathclose%
\pgfusepath{fill}%
\end{pgfscope}%
\begin{pgfscope}%
\pgfpathrectangle{\pgfqpoint{0.407592in}{0.222904in}}{\pgfqpoint{1.411163in}{1.411163in}}%
\pgfusepath{clip}%
\pgfsetbuttcap%
\pgfsetroundjoin%
\definecolor{currentfill}{rgb}{0.749020,0.772549,0.792157}%
\pgfsetfillcolor{currentfill}%
\pgfsetfillopacity{0.160000}%
\pgfsetlinewidth{0.000000pt}%
\definecolor{currentstroke}{rgb}{0.000000,0.000000,0.000000}%
\pgfsetstrokecolor{currentstroke}%
\pgfsetdash{}{0pt}%
\pgfpathmoveto{\pgfqpoint{1.221828in}{0.731133in}}%
\pgfpathlineto{\pgfqpoint{1.235052in}{0.743252in}}%
\pgfpathlineto{\pgfqpoint{1.248288in}{0.755381in}}%
\pgfpathlineto{\pgfqpoint{1.246925in}{0.759754in}}%
\pgfpathlineto{\pgfqpoint{1.244882in}{0.764061in}}%
\pgfpathlineto{\pgfqpoint{1.231962in}{0.751234in}}%
\pgfpathlineto{\pgfqpoint{1.219047in}{0.738413in}}%
\pgfpathlineto{\pgfqpoint{1.220721in}{0.734800in}}%
\pgfpathlineto{\pgfqpoint{1.221828in}{0.731133in}}%
\pgfpathclose%
\pgfusepath{fill}%
\end{pgfscope}%
\begin{pgfscope}%
\pgfpathrectangle{\pgfqpoint{0.407592in}{0.222904in}}{\pgfqpoint{1.411163in}{1.411163in}}%
\pgfusepath{clip}%
\pgfsetbuttcap%
\pgfsetroundjoin%
\definecolor{currentfill}{rgb}{0.749020,0.772549,0.792157}%
\pgfsetfillcolor{currentfill}%
\pgfsetfillopacity{0.160000}%
\pgfsetlinewidth{0.000000pt}%
\definecolor{currentstroke}{rgb}{0.000000,0.000000,0.000000}%
\pgfsetstrokecolor{currentstroke}%
\pgfsetdash{}{0pt}%
\pgfpathmoveto{\pgfqpoint{0.800434in}{0.932833in}}%
\pgfpathlineto{\pgfqpoint{0.788344in}{0.946279in}}%
\pgfpathlineto{\pgfqpoint{0.776253in}{0.959726in}}%
\pgfpathlineto{\pgfqpoint{0.766887in}{0.951227in}}%
\pgfpathlineto{\pgfqpoint{0.758868in}{0.942501in}}%
\pgfpathlineto{\pgfqpoint{0.771694in}{0.929651in}}%
\pgfpathlineto{\pgfqpoint{0.784513in}{0.916808in}}%
\pgfpathlineto{\pgfqpoint{0.791850in}{0.924924in}}%
\pgfpathlineto{\pgfqpoint{0.800434in}{0.932833in}}%
\pgfpathclose%
\pgfusepath{fill}%
\end{pgfscope}%
\begin{pgfscope}%
\pgfpathrectangle{\pgfqpoint{0.407592in}{0.222904in}}{\pgfqpoint{1.411163in}{1.411163in}}%
\pgfusepath{clip}%
\pgfsetbuttcap%
\pgfsetroundjoin%
\definecolor{currentfill}{rgb}{0.749020,0.772549,0.792157}%
\pgfsetfillcolor{currentfill}%
\pgfsetfillopacity{0.160000}%
\pgfsetlinewidth{0.000000pt}%
\definecolor{currentstroke}{rgb}{0.000000,0.000000,0.000000}%
\pgfsetstrokecolor{currentstroke}%
\pgfsetdash{}{0pt}%
\pgfpathmoveto{\pgfqpoint{1.452439in}{0.970108in}}%
\pgfpathlineto{\pgfqpoint{1.465463in}{0.983037in}}%
\pgfpathlineto{\pgfqpoint{1.478494in}{0.995973in}}%
\pgfpathlineto{\pgfqpoint{1.470977in}{1.005837in}}%
\pgfpathlineto{\pgfqpoint{1.461869in}{1.015428in}}%
\pgfpathlineto{\pgfqpoint{1.449471in}{1.001920in}}%
\pgfpathlineto{\pgfqpoint{1.437073in}{0.988413in}}%
\pgfpathlineto{\pgfqpoint{1.445501in}{0.979387in}}%
\pgfpathlineto{\pgfqpoint{1.452439in}{0.970108in}}%
\pgfpathclose%
\pgfusepath{fill}%
\end{pgfscope}%
\begin{pgfscope}%
\pgfpathrectangle{\pgfqpoint{0.407592in}{0.222904in}}{\pgfqpoint{1.411163in}{1.411163in}}%
\pgfusepath{clip}%
\pgfsetbuttcap%
\pgfsetroundjoin%
\definecolor{currentfill}{rgb}{0.749020,0.772549,0.792157}%
\pgfsetfillcolor{currentfill}%
\pgfsetfillopacity{0.160000}%
\pgfsetlinewidth{0.000000pt}%
\definecolor{currentstroke}{rgb}{0.000000,0.000000,0.000000}%
\pgfsetstrokecolor{currentstroke}%
\pgfsetdash{}{0pt}%
\pgfpathmoveto{\pgfqpoint{1.142060in}{0.652597in}}%
\pgfpathlineto{\pgfqpoint{1.155074in}{0.663159in}}%
\pgfpathlineto{\pgfqpoint{1.168111in}{0.673739in}}%
\pgfpathlineto{\pgfqpoint{1.168865in}{0.675993in}}%
\pgfpathlineto{\pgfqpoint{1.169272in}{0.678255in}}%
\pgfpathlineto{\pgfqpoint{1.156055in}{0.666930in}}%
\pgfpathlineto{\pgfqpoint{1.142855in}{0.655621in}}%
\pgfpathlineto{\pgfqpoint{1.142573in}{0.654106in}}%
\pgfpathlineto{\pgfqpoint{1.142060in}{0.652597in}}%
\pgfpathclose%
\pgfusepath{fill}%
\end{pgfscope}%
\begin{pgfscope}%
\pgfpathrectangle{\pgfqpoint{0.407592in}{0.222904in}}{\pgfqpoint{1.411163in}{1.411163in}}%
\pgfusepath{clip}%
\pgfsetbuttcap%
\pgfsetroundjoin%
\definecolor{currentfill}{rgb}{0.749020,0.772549,0.792157}%
\pgfsetfillcolor{currentfill}%
\pgfsetfillopacity{0.160000}%
\pgfsetlinewidth{0.000000pt}%
\definecolor{currentstroke}{rgb}{0.000000,0.000000,0.000000}%
\pgfsetstrokecolor{currentstroke}%
\pgfsetdash{}{0pt}%
\pgfpathmoveto{\pgfqpoint{1.169272in}{0.678255in}}%
\pgfpathlineto{\pgfqpoint{1.182507in}{0.689594in}}%
\pgfpathlineto{\pgfqpoint{1.195759in}{0.700948in}}%
\pgfpathlineto{\pgfqpoint{1.195818in}{0.703945in}}%
\pgfpathlineto{\pgfqpoint{1.195414in}{0.706927in}}%
\pgfpathlineto{\pgfqpoint{1.182225in}{0.694840in}}%
\pgfpathlineto{\pgfqpoint{1.169048in}{0.682764in}}%
\pgfpathlineto{\pgfqpoint{1.169333in}{0.680515in}}%
\pgfpathlineto{\pgfqpoint{1.169272in}{0.678255in}}%
\pgfpathclose%
\pgfusepath{fill}%
\end{pgfscope}%
\begin{pgfscope}%
\pgfpathrectangle{\pgfqpoint{0.407592in}{0.222904in}}{\pgfqpoint{1.411163in}{1.411163in}}%
\pgfusepath{clip}%
\pgfsetbuttcap%
\pgfsetroundjoin%
\definecolor{currentfill}{rgb}{0.749020,0.772549,0.792157}%
\pgfsetfillcolor{currentfill}%
\pgfsetfillopacity{0.160000}%
\pgfsetlinewidth{0.000000pt}%
\definecolor{currentstroke}{rgb}{0.000000,0.000000,0.000000}%
\pgfsetstrokecolor{currentstroke}%
\pgfsetdash{}{0pt}%
\pgfpathmoveto{\pgfqpoint{1.037317in}{0.655317in}}%
\pgfpathlineto{\pgfqpoint{1.024044in}{0.666551in}}%
\pgfpathlineto{\pgfqpoint{1.010753in}{0.677801in}}%
\pgfpathlineto{\pgfqpoint{1.010763in}{0.675540in}}%
\pgfpathlineto{\pgfqpoint{1.011121in}{0.673289in}}%
\pgfpathlineto{\pgfqpoint{1.024364in}{0.662783in}}%
\pgfpathlineto{\pgfqpoint{1.037585in}{0.652295in}}%
\pgfpathlineto{\pgfqpoint{1.037335in}{0.653802in}}%
\pgfpathlineto{\pgfqpoint{1.037317in}{0.655317in}}%
\pgfpathclose%
\pgfusepath{fill}%
\end{pgfscope}%
\begin{pgfscope}%
\pgfpathrectangle{\pgfqpoint{0.407592in}{0.222904in}}{\pgfqpoint{1.411163in}{1.411163in}}%
\pgfusepath{clip}%
\pgfsetbuttcap%
\pgfsetroundjoin%
\definecolor{currentfill}{rgb}{0.749020,0.772549,0.792157}%
\pgfsetfillcolor{currentfill}%
\pgfsetfillopacity{0.160000}%
\pgfsetlinewidth{0.000000pt}%
\definecolor{currentstroke}{rgb}{0.000000,0.000000,0.000000}%
\pgfsetstrokecolor{currentstroke}%
\pgfsetdash{}{0pt}%
\pgfpathmoveto{\pgfqpoint{0.963301in}{0.737693in}}%
\pgfpathlineto{\pgfqpoint{0.950573in}{0.750445in}}%
\pgfpathlineto{\pgfqpoint{0.937838in}{0.763203in}}%
\pgfpathlineto{\pgfqpoint{0.935020in}{0.758881in}}%
\pgfpathlineto{\pgfqpoint{0.932871in}{0.754498in}}%
\pgfpathlineto{\pgfqpoint{0.946052in}{0.742439in}}%
\pgfpathlineto{\pgfqpoint{0.959220in}{0.730392in}}%
\pgfpathlineto{\pgfqpoint{0.960982in}{0.734068in}}%
\pgfpathlineto{\pgfqpoint{0.963301in}{0.737693in}}%
\pgfpathclose%
\pgfusepath{fill}%
\end{pgfscope}%
\begin{pgfscope}%
\pgfpathrectangle{\pgfqpoint{0.407592in}{0.222904in}}{\pgfqpoint{1.411163in}{1.411163in}}%
\pgfusepath{clip}%
\pgfsetbuttcap%
\pgfsetroundjoin%
\definecolor{currentfill}{rgb}{0.749020,0.772549,0.792157}%
\pgfsetfillcolor{currentfill}%
\pgfsetfillopacity{0.160000}%
\pgfsetlinewidth{0.000000pt}%
\definecolor{currentstroke}{rgb}{0.000000,0.000000,0.000000}%
\pgfsetstrokecolor{currentstroke}%
\pgfsetdash{}{0pt}%
\pgfpathmoveto{\pgfqpoint{1.011770in}{0.682314in}}%
\pgfpathlineto{\pgfqpoint{0.998651in}{0.694317in}}%
\pgfpathlineto{\pgfqpoint{0.985520in}{0.706331in}}%
\pgfpathlineto{\pgfqpoint{0.984588in}{0.703345in}}%
\pgfpathlineto{\pgfqpoint{0.984116in}{0.700345in}}%
\pgfpathlineto{\pgfqpoint{0.997444in}{0.689065in}}%
\pgfpathlineto{\pgfqpoint{1.010753in}{0.677801in}}%
\pgfpathlineto{\pgfqpoint{1.011089in}{0.680062in}}%
\pgfpathlineto{\pgfqpoint{1.011770in}{0.682314in}}%
\pgfpathclose%
\pgfusepath{fill}%
\end{pgfscope}%
\begin{pgfscope}%
\pgfpathrectangle{\pgfqpoint{0.407592in}{0.222904in}}{\pgfqpoint{1.411163in}{1.411163in}}%
\pgfusepath{clip}%
\pgfsetbuttcap%
\pgfsetroundjoin%
\definecolor{currentfill}{rgb}{0.749020,0.772549,0.792157}%
\pgfsetfillcolor{currentfill}%
\pgfsetfillopacity{0.160000}%
\pgfsetlinewidth{0.000000pt}%
\definecolor{currentstroke}{rgb}{0.000000,0.000000,0.000000}%
\pgfsetstrokecolor{currentstroke}%
\pgfsetdash{}{0pt}%
\pgfpathmoveto{\pgfqpoint{1.478494in}{0.995973in}}%
\pgfpathlineto{\pgfqpoint{1.491530in}{1.008914in}}%
\pgfpathlineto{\pgfqpoint{1.504572in}{1.021862in}}%
\pgfpathlineto{\pgfqpoint{1.496467in}{1.032301in}}%
\pgfpathlineto{\pgfqpoint{1.486669in}{1.042447in}}%
\pgfpathlineto{\pgfqpoint{1.474269in}{1.028937in}}%
\pgfpathlineto{\pgfqpoint{1.461869in}{1.015428in}}%
\pgfpathlineto{\pgfqpoint{1.470977in}{1.005837in}}%
\pgfpathlineto{\pgfqpoint{1.478494in}{0.995973in}}%
\pgfpathclose%
\pgfusepath{fill}%
\end{pgfscope}%
\begin{pgfscope}%
\pgfpathrectangle{\pgfqpoint{0.407592in}{0.222904in}}{\pgfqpoint{1.411163in}{1.411163in}}%
\pgfusepath{clip}%
\pgfsetbuttcap%
\pgfsetroundjoin%
\definecolor{currentfill}{rgb}{0.749020,0.772549,0.792157}%
\pgfsetfillcolor{currentfill}%
\pgfsetfillopacity{0.160000}%
\pgfsetlinewidth{0.000000pt}%
\definecolor{currentstroke}{rgb}{0.000000,0.000000,0.000000}%
\pgfsetstrokecolor{currentstroke}%
\pgfsetdash{}{0pt}%
\pgfpathmoveto{\pgfqpoint{0.776253in}{0.959726in}}%
\pgfpathlineto{\pgfqpoint{0.764161in}{0.973174in}}%
\pgfpathlineto{\pgfqpoint{0.752068in}{0.986624in}}%
\pgfpathlineto{\pgfqpoint{0.741910in}{0.977545in}}%
\pgfpathlineto{\pgfqpoint{0.733196in}{0.968219in}}%
\pgfpathlineto{\pgfqpoint{0.746035in}{0.955357in}}%
\pgfpathlineto{\pgfqpoint{0.758868in}{0.942501in}}%
\pgfpathlineto{\pgfqpoint{0.766887in}{0.951227in}}%
\pgfpathlineto{\pgfqpoint{0.776253in}{0.959726in}}%
\pgfpathclose%
\pgfusepath{fill}%
\end{pgfscope}%
\begin{pgfscope}%
\pgfpathrectangle{\pgfqpoint{0.407592in}{0.222904in}}{\pgfqpoint{1.411163in}{1.411163in}}%
\pgfusepath{clip}%
\pgfsetbuttcap%
\pgfsetroundjoin%
\definecolor{currentfill}{rgb}{0.749020,0.772549,0.792157}%
\pgfsetfillcolor{currentfill}%
\pgfsetfillopacity{0.160000}%
\pgfsetlinewidth{0.000000pt}%
\definecolor{currentstroke}{rgb}{0.000000,0.000000,0.000000}%
\pgfsetstrokecolor{currentstroke}%
\pgfsetdash{}{0pt}%
\pgfpathmoveto{\pgfqpoint{1.112219in}{0.627238in}}%
\pgfpathlineto{\pgfqpoint{1.123262in}{0.635628in}}%
\pgfpathlineto{\pgfqpoint{1.134338in}{0.644042in}}%
\pgfpathlineto{\pgfqpoint{1.136154in}{0.645377in}}%
\pgfpathlineto{\pgfqpoint{1.137766in}{0.646756in}}%
\pgfpathlineto{\pgfqpoint{1.125836in}{0.637665in}}%
\pgfpathlineto{\pgfqpoint{1.113937in}{0.628598in}}%
\pgfpathlineto{\pgfqpoint{1.113128in}{0.627907in}}%
\pgfpathlineto{\pgfqpoint{1.112219in}{0.627238in}}%
\pgfpathclose%
\pgfusepath{fill}%
\end{pgfscope}%
\begin{pgfscope}%
\pgfpathrectangle{\pgfqpoint{0.407592in}{0.222904in}}{\pgfqpoint{1.411163in}{1.411163in}}%
\pgfusepath{clip}%
\pgfsetbuttcap%
\pgfsetroundjoin%
\definecolor{currentfill}{rgb}{0.749020,0.772549,0.792157}%
\pgfsetfillcolor{currentfill}%
\pgfsetfillopacity{0.160000}%
\pgfsetlinewidth{0.000000pt}%
\definecolor{currentstroke}{rgb}{0.000000,0.000000,0.000000}%
\pgfsetstrokecolor{currentstroke}%
\pgfsetdash{}{0pt}%
\pgfpathmoveto{\pgfqpoint{1.065610in}{0.628457in}}%
\pgfpathlineto{\pgfqpoint{1.053252in}{0.637455in}}%
\pgfpathlineto{\pgfqpoint{1.040861in}{0.646476in}}%
\pgfpathlineto{\pgfqpoint{1.042234in}{0.645105in}}%
\pgfpathlineto{\pgfqpoint{1.043817in}{0.643781in}}%
\pgfpathlineto{\pgfqpoint{1.055472in}{0.635431in}}%
\pgfpathlineto{\pgfqpoint{1.067093in}{0.627107in}}%
\pgfpathlineto{\pgfqpoint{1.066299in}{0.627771in}}%
\pgfpathlineto{\pgfqpoint{1.065610in}{0.628457in}}%
\pgfpathclose%
\pgfusepath{fill}%
\end{pgfscope}%
\begin{pgfscope}%
\pgfpathrectangle{\pgfqpoint{0.407592in}{0.222904in}}{\pgfqpoint{1.411163in}{1.411163in}}%
\pgfusepath{clip}%
\pgfsetbuttcap%
\pgfsetroundjoin%
\definecolor{currentfill}{rgb}{0.749020,0.772549,0.792157}%
\pgfsetfillcolor{currentfill}%
\pgfsetfillopacity{0.160000}%
\pgfsetlinewidth{0.000000pt}%
\definecolor{currentstroke}{rgb}{0.000000,0.000000,0.000000}%
\pgfsetstrokecolor{currentstroke}%
\pgfsetdash{}{0pt}%
\pgfpathmoveto{\pgfqpoint{1.504572in}{1.021862in}}%
\pgfpathlineto{\pgfqpoint{1.517621in}{1.034815in}}%
\pgfpathlineto{\pgfqpoint{1.530675in}{1.047775in}}%
\pgfpathlineto{\pgfqpoint{1.521969in}{1.058779in}}%
\pgfpathlineto{\pgfqpoint{1.511471in}{1.069469in}}%
\pgfpathlineto{\pgfqpoint{1.499070in}{1.055958in}}%
\pgfpathlineto{\pgfqpoint{1.486669in}{1.042447in}}%
\pgfpathlineto{\pgfqpoint{1.496467in}{1.032301in}}%
\pgfpathlineto{\pgfqpoint{1.504572in}{1.021862in}}%
\pgfpathclose%
\pgfusepath{fill}%
\end{pgfscope}%
\begin{pgfscope}%
\pgfpathrectangle{\pgfqpoint{0.407592in}{0.222904in}}{\pgfqpoint{1.411163in}{1.411163in}}%
\pgfusepath{clip}%
\pgfsetbuttcap%
\pgfsetroundjoin%
\definecolor{currentfill}{rgb}{0.749020,0.772549,0.792157}%
\pgfsetfillcolor{currentfill}%
\pgfsetfillopacity{0.160000}%
\pgfsetlinewidth{0.000000pt}%
\definecolor{currentstroke}{rgb}{0.000000,0.000000,0.000000}%
\pgfsetstrokecolor{currentstroke}%
\pgfsetdash{}{0pt}%
\pgfpathmoveto{\pgfqpoint{1.248288in}{0.755381in}}%
\pgfpathlineto{\pgfqpoint{1.261536in}{0.767522in}}%
\pgfpathlineto{\pgfqpoint{1.274795in}{0.779674in}}%
\pgfpathlineto{\pgfqpoint{1.273165in}{0.784742in}}%
\pgfpathlineto{\pgfqpoint{1.270742in}{0.789733in}}%
\pgfpathlineto{\pgfqpoint{1.257809in}{0.776894in}}%
\pgfpathlineto{\pgfqpoint{1.244882in}{0.764061in}}%
\pgfpathlineto{\pgfqpoint{1.246925in}{0.759754in}}%
\pgfpathlineto{\pgfqpoint{1.248288in}{0.755381in}}%
\pgfpathclose%
\pgfusepath{fill}%
\end{pgfscope}%
\begin{pgfscope}%
\pgfpathrectangle{\pgfqpoint{0.407592in}{0.222904in}}{\pgfqpoint{1.411163in}{1.411163in}}%
\pgfusepath{clip}%
\pgfsetbuttcap%
\pgfsetroundjoin%
\definecolor{currentfill}{rgb}{0.749020,0.772549,0.792157}%
\pgfsetfillcolor{currentfill}%
\pgfsetfillopacity{0.160000}%
\pgfsetlinewidth{0.000000pt}%
\definecolor{currentstroke}{rgb}{0.000000,0.000000,0.000000}%
\pgfsetstrokecolor{currentstroke}%
\pgfsetdash{}{0pt}%
\pgfpathmoveto{\pgfqpoint{0.752068in}{0.986624in}}%
\pgfpathlineto{\pgfqpoint{0.739973in}{1.000075in}}%
\pgfpathlineto{\pgfqpoint{0.727877in}{1.013528in}}%
\pgfpathlineto{\pgfqpoint{0.716917in}{1.003879in}}%
\pgfpathlineto{\pgfqpoint{0.707498in}{0.993964in}}%
\pgfpathlineto{\pgfqpoint{0.720350in}{0.981088in}}%
\pgfpathlineto{\pgfqpoint{0.733196in}{0.968219in}}%
\pgfpathlineto{\pgfqpoint{0.741910in}{0.977545in}}%
\pgfpathlineto{\pgfqpoint{0.752068in}{0.986624in}}%
\pgfpathclose%
\pgfusepath{fill}%
\end{pgfscope}%
\begin{pgfscope}%
\pgfpathrectangle{\pgfqpoint{0.407592in}{0.222904in}}{\pgfqpoint{1.411163in}{1.411163in}}%
\pgfusepath{clip}%
\pgfsetbuttcap%
\pgfsetroundjoin%
\definecolor{currentfill}{rgb}{0.749020,0.772549,0.792157}%
\pgfsetfillcolor{currentfill}%
\pgfsetfillopacity{0.160000}%
\pgfsetlinewidth{0.000000pt}%
\definecolor{currentstroke}{rgb}{0.000000,0.000000,0.000000}%
\pgfsetstrokecolor{currentstroke}%
\pgfsetdash{}{0pt}%
\pgfpathmoveto{\pgfqpoint{1.111212in}{0.626595in}}%
\pgfpathlineto{\pgfqpoint{1.121753in}{0.634664in}}%
\pgfpathlineto{\pgfqpoint{1.132327in}{0.642759in}}%
\pgfpathlineto{\pgfqpoint{1.134338in}{0.644042in}}%
\pgfpathlineto{\pgfqpoint{1.123262in}{0.635628in}}%
\pgfpathlineto{\pgfqpoint{1.112219in}{0.627238in}}%
\pgfusepath{fill}%
\end{pgfscope}%
\begin{pgfscope}%
\pgfpathrectangle{\pgfqpoint{0.407592in}{0.222904in}}{\pgfqpoint{1.411163in}{1.411163in}}%
\pgfusepath{clip}%
\pgfsetbuttcap%
\pgfsetroundjoin%
\definecolor{currentfill}{rgb}{0.749020,0.772549,0.792157}%
\pgfsetfillcolor{currentfill}%
\pgfsetfillopacity{0.160000}%
\pgfsetlinewidth{0.000000pt}%
\definecolor{currentstroke}{rgb}{0.000000,0.000000,0.000000}%
\pgfsetstrokecolor{currentstroke}%
\pgfsetdash{}{0pt}%
\pgfpathmoveto{\pgfqpoint{1.530675in}{1.047775in}}%
\pgfpathlineto{\pgfqpoint{1.543736in}{1.060741in}}%
\pgfpathlineto{\pgfqpoint{1.556803in}{1.073712in}}%
\pgfpathlineto{\pgfqpoint{1.547486in}{1.085270in}}%
\pgfpathlineto{\pgfqpoint{1.536277in}{1.096494in}}%
\pgfpathlineto{\pgfqpoint{1.523874in}{1.082981in}}%
\pgfpathlineto{\pgfqpoint{1.511471in}{1.069469in}}%
\pgfpathlineto{\pgfqpoint{1.521969in}{1.058779in}}%
\pgfpathlineto{\pgfqpoint{1.530675in}{1.047775in}}%
\pgfpathclose%
\pgfusepath{fill}%
\end{pgfscope}%
\begin{pgfscope}%
\pgfpathrectangle{\pgfqpoint{0.407592in}{0.222904in}}{\pgfqpoint{1.411163in}{1.411163in}}%
\pgfusepath{clip}%
\pgfsetbuttcap%
\pgfsetroundjoin%
\definecolor{currentfill}{rgb}{0.749020,0.772549,0.792157}%
\pgfsetfillcolor{currentfill}%
\pgfsetfillopacity{0.160000}%
\pgfsetlinewidth{0.000000pt}%
\definecolor{currentstroke}{rgb}{0.000000,0.000000,0.000000}%
\pgfsetstrokecolor{currentstroke}%
\pgfsetdash{}{0pt}%
\pgfpathmoveto{\pgfqpoint{0.937838in}{0.763203in}}%
\pgfpathlineto{\pgfqpoint{0.925097in}{0.775968in}}%
\pgfpathlineto{\pgfqpoint{0.912349in}{0.788739in}}%
\pgfpathlineto{\pgfqpoint{0.909022in}{0.783731in}}%
\pgfpathlineto{\pgfqpoint{0.906473in}{0.778649in}}%
\pgfpathlineto{\pgfqpoint{0.919678in}{0.766568in}}%
\pgfpathlineto{\pgfqpoint{0.932871in}{0.754498in}}%
\pgfpathlineto{\pgfqpoint{0.935020in}{0.758881in}}%
\pgfpathlineto{\pgfqpoint{0.937838in}{0.763203in}}%
\pgfpathclose%
\pgfusepath{fill}%
\end{pgfscope}%
\begin{pgfscope}%
\pgfpathrectangle{\pgfqpoint{0.407592in}{0.222904in}}{\pgfqpoint{1.411163in}{1.411163in}}%
\pgfusepath{clip}%
\pgfsetbuttcap%
\pgfsetroundjoin%
\definecolor{currentfill}{rgb}{0.749020,0.772549,0.792157}%
\pgfsetfillcolor{currentfill}%
\pgfsetfillopacity{0.160000}%
\pgfsetlinewidth{0.000000pt}%
\definecolor{currentstroke}{rgb}{0.000000,0.000000,0.000000}%
\pgfsetstrokecolor{currentstroke}%
\pgfsetdash{}{0pt}%
\pgfpathmoveto{\pgfqpoint{0.727877in}{1.013528in}}%
\pgfpathlineto{\pgfqpoint{0.715780in}{1.026982in}}%
\pgfpathlineto{\pgfqpoint{0.703682in}{1.040437in}}%
\pgfpathlineto{\pgfqpoint{0.691909in}{1.030229in}}%
\pgfpathlineto{\pgfqpoint{0.681774in}{1.019736in}}%
\pgfpathlineto{\pgfqpoint{0.694639in}{1.006847in}}%
\pgfpathlineto{\pgfqpoint{0.707498in}{0.993964in}}%
\pgfpathlineto{\pgfqpoint{0.716917in}{1.003879in}}%
\pgfpathlineto{\pgfqpoint{0.727877in}{1.013528in}}%
\pgfpathclose%
\pgfusepath{fill}%
\end{pgfscope}%
\begin{pgfscope}%
\pgfpathrectangle{\pgfqpoint{0.407592in}{0.222904in}}{\pgfqpoint{1.411163in}{1.411163in}}%
\pgfusepath{clip}%
\pgfsetbuttcap%
\pgfsetroundjoin%
\definecolor{currentfill}{rgb}{0.749020,0.772549,0.792157}%
\pgfsetfillcolor{currentfill}%
\pgfsetfillopacity{0.160000}%
\pgfsetlinewidth{0.000000pt}%
\definecolor{currentstroke}{rgb}{0.000000,0.000000,0.000000}%
\pgfsetstrokecolor{currentstroke}%
\pgfsetdash{}{0pt}%
\pgfpathmoveto{\pgfqpoint{1.067093in}{0.627107in}}%
\pgfpathlineto{\pgfqpoint{1.055472in}{0.635431in}}%
\pgfpathlineto{\pgfqpoint{1.043817in}{0.643781in}}%
\pgfpathlineto{\pgfqpoint{1.045605in}{0.642508in}}%
\pgfpathlineto{\pgfqpoint{1.047590in}{0.641293in}}%
\pgfpathlineto{\pgfqpoint{1.058303in}{0.633564in}}%
\pgfpathlineto{\pgfqpoint{1.068981in}{0.625862in}}%
\pgfpathlineto{\pgfqpoint{1.067988in}{0.626470in}}%
\pgfpathlineto{\pgfqpoint{1.067093in}{0.627107in}}%
\pgfpathclose%
\pgfusepath{fill}%
\end{pgfscope}%
\begin{pgfscope}%
\pgfpathrectangle{\pgfqpoint{0.407592in}{0.222904in}}{\pgfqpoint{1.411163in}{1.411163in}}%
\pgfusepath{clip}%
\pgfsetbuttcap%
\pgfsetroundjoin%
\definecolor{currentfill}{rgb}{0.749020,0.772549,0.792157}%
\pgfsetfillcolor{currentfill}%
\pgfsetfillopacity{0.160000}%
\pgfsetlinewidth{0.000000pt}%
\definecolor{currentstroke}{rgb}{0.000000,0.000000,0.000000}%
\pgfsetstrokecolor{currentstroke}%
\pgfsetdash{}{0pt}%
\pgfpathmoveto{\pgfqpoint{1.556803in}{1.073712in}}%
\pgfpathlineto{\pgfqpoint{1.569876in}{1.086690in}}%
\pgfpathlineto{\pgfqpoint{1.582955in}{1.099674in}}%
\pgfpathlineto{\pgfqpoint{1.573015in}{1.111776in}}%
\pgfpathlineto{\pgfqpoint{1.561085in}{1.123522in}}%
\pgfpathlineto{\pgfqpoint{1.548681in}{1.110008in}}%
\pgfpathlineto{\pgfqpoint{1.536277in}{1.096494in}}%
\pgfpathlineto{\pgfqpoint{1.547486in}{1.085270in}}%
\pgfpathlineto{\pgfqpoint{1.556803in}{1.073712in}}%
\pgfpathclose%
\pgfusepath{fill}%
\end{pgfscope}%
\begin{pgfscope}%
\pgfpathrectangle{\pgfqpoint{0.407592in}{0.222904in}}{\pgfqpoint{1.411163in}{1.411163in}}%
\pgfusepath{clip}%
\pgfsetbuttcap%
\pgfsetroundjoin%
\definecolor{currentfill}{rgb}{0.749020,0.772549,0.792157}%
\pgfsetfillcolor{currentfill}%
\pgfsetfillopacity{0.160000}%
\pgfsetlinewidth{0.000000pt}%
\definecolor{currentstroke}{rgb}{0.000000,0.000000,0.000000}%
\pgfsetstrokecolor{currentstroke}%
\pgfsetdash{}{0pt}%
\pgfpathmoveto{\pgfqpoint{1.140354in}{0.649624in}}%
\pgfpathlineto{\pgfqpoint{1.152954in}{0.659450in}}%
\pgfpathlineto{\pgfqpoint{1.165580in}{0.669297in}}%
\pgfpathlineto{\pgfqpoint{1.167015in}{0.671504in}}%
\pgfpathlineto{\pgfqpoint{1.168111in}{0.673739in}}%
\pgfpathlineto{\pgfqpoint{1.155074in}{0.663159in}}%
\pgfpathlineto{\pgfqpoint{1.142060in}{0.652597in}}%
\pgfpathlineto{\pgfqpoint{1.141320in}{0.651101in}}%
\pgfpathlineto{\pgfqpoint{1.140354in}{0.649624in}}%
\pgfpathclose%
\pgfusepath{fill}%
\end{pgfscope}%
\begin{pgfscope}%
\pgfpathrectangle{\pgfqpoint{0.407592in}{0.222904in}}{\pgfqpoint{1.411163in}{1.411163in}}%
\pgfusepath{clip}%
\pgfsetbuttcap%
\pgfsetroundjoin%
\definecolor{currentfill}{rgb}{0.749020,0.772549,0.792157}%
\pgfsetfillcolor{currentfill}%
\pgfsetfillopacity{0.160000}%
\pgfsetlinewidth{0.000000pt}%
\definecolor{currentstroke}{rgb}{0.000000,0.000000,0.000000}%
\pgfsetstrokecolor{currentstroke}%
\pgfsetdash{}{0pt}%
\pgfpathmoveto{\pgfqpoint{1.108927in}{0.625398in}}%
\pgfpathlineto{\pgfqpoint{1.118325in}{0.632869in}}%
\pgfpathlineto{\pgfqpoint{1.127755in}{0.640367in}}%
\pgfpathlineto{\pgfqpoint{1.130130in}{0.641532in}}%
\pgfpathlineto{\pgfqpoint{1.132327in}{0.642759in}}%
\pgfpathlineto{\pgfqpoint{1.121753in}{0.634664in}}%
\pgfpathlineto{\pgfqpoint{1.111212in}{0.626595in}}%
\pgfpathlineto{\pgfqpoint{1.110114in}{0.625981in}}%
\pgfpathlineto{\pgfqpoint{1.108927in}{0.625398in}}%
\pgfpathclose%
\pgfusepath{fill}%
\end{pgfscope}%
\begin{pgfscope}%
\pgfpathrectangle{\pgfqpoint{0.407592in}{0.222904in}}{\pgfqpoint{1.411163in}{1.411163in}}%
\pgfusepath{clip}%
\pgfsetbuttcap%
\pgfsetroundjoin%
\definecolor{currentfill}{rgb}{0.749020,0.772549,0.792157}%
\pgfsetfillcolor{currentfill}%
\pgfsetfillopacity{0.160000}%
\pgfsetlinewidth{0.000000pt}%
\definecolor{currentstroke}{rgb}{0.000000,0.000000,0.000000}%
\pgfsetstrokecolor{currentstroke}%
\pgfsetdash{}{0pt}%
\pgfpathmoveto{\pgfqpoint{1.195759in}{0.700948in}}%
\pgfpathlineto{\pgfqpoint{1.209029in}{0.712317in}}%
\pgfpathlineto{\pgfqpoint{1.222316in}{0.723701in}}%
\pgfpathlineto{\pgfqpoint{1.222361in}{0.727428in}}%
\pgfpathlineto{\pgfqpoint{1.221828in}{0.731133in}}%
\pgfpathlineto{\pgfqpoint{1.208615in}{0.719024in}}%
\pgfpathlineto{\pgfqpoint{1.195414in}{0.706927in}}%
\pgfpathlineto{\pgfqpoint{1.195818in}{0.703945in}}%
\pgfpathlineto{\pgfqpoint{1.195759in}{0.700948in}}%
\pgfpathclose%
\pgfusepath{fill}%
\end{pgfscope}%
\begin{pgfscope}%
\pgfpathrectangle{\pgfqpoint{0.407592in}{0.222904in}}{\pgfqpoint{1.411163in}{1.411163in}}%
\pgfusepath{clip}%
\pgfsetbuttcap%
\pgfsetroundjoin%
\definecolor{currentfill}{rgb}{0.749020,0.772549,0.792157}%
\pgfsetfillcolor{currentfill}%
\pgfsetfillopacity{0.160000}%
\pgfsetlinewidth{0.000000pt}%
\definecolor{currentstroke}{rgb}{0.000000,0.000000,0.000000}%
\pgfsetstrokecolor{currentstroke}%
\pgfsetdash{}{0pt}%
\pgfpathmoveto{\pgfqpoint{1.037585in}{0.652295in}}%
\pgfpathlineto{\pgfqpoint{1.024364in}{0.662783in}}%
\pgfpathlineto{\pgfqpoint{1.011121in}{0.673289in}}%
\pgfpathlineto{\pgfqpoint{1.011825in}{0.671059in}}%
\pgfpathlineto{\pgfqpoint{1.012872in}{0.668859in}}%
\pgfpathlineto{\pgfqpoint{1.025837in}{0.659084in}}%
\pgfpathlineto{\pgfqpoint{1.038772in}{0.649331in}}%
\pgfpathlineto{\pgfqpoint{1.038064in}{0.650803in}}%
\pgfpathlineto{\pgfqpoint{1.037585in}{0.652295in}}%
\pgfpathclose%
\pgfusepath{fill}%
\end{pgfscope}%
\begin{pgfscope}%
\pgfpathrectangle{\pgfqpoint{0.407592in}{0.222904in}}{\pgfqpoint{1.411163in}{1.411163in}}%
\pgfusepath{clip}%
\pgfsetbuttcap%
\pgfsetroundjoin%
\definecolor{currentfill}{rgb}{0.749020,0.772549,0.792157}%
\pgfsetfillcolor{currentfill}%
\pgfsetfillopacity{0.160000}%
\pgfsetlinewidth{0.000000pt}%
\definecolor{currentstroke}{rgb}{0.000000,0.000000,0.000000}%
\pgfsetstrokecolor{currentstroke}%
\pgfsetdash{}{0pt}%
\pgfpathmoveto{\pgfqpoint{0.703682in}{1.040437in}}%
\pgfpathlineto{\pgfqpoint{0.691583in}{1.053893in}}%
\pgfpathlineto{\pgfqpoint{0.679482in}{1.067351in}}%
\pgfpathlineto{\pgfqpoint{0.666885in}{1.056595in}}%
\pgfpathlineto{\pgfqpoint{0.656023in}{1.045533in}}%
\pgfpathlineto{\pgfqpoint{0.668902in}{1.032631in}}%
\pgfpathlineto{\pgfqpoint{0.681774in}{1.019736in}}%
\pgfpathlineto{\pgfqpoint{0.691909in}{1.030229in}}%
\pgfpathlineto{\pgfqpoint{0.703682in}{1.040437in}}%
\pgfpathclose%
\pgfusepath{fill}%
\end{pgfscope}%
\begin{pgfscope}%
\pgfpathrectangle{\pgfqpoint{0.407592in}{0.222904in}}{\pgfqpoint{1.411163in}{1.411163in}}%
\pgfusepath{clip}%
\pgfsetbuttcap%
\pgfsetroundjoin%
\definecolor{currentfill}{rgb}{0.749020,0.772549,0.792157}%
\pgfsetfillcolor{currentfill}%
\pgfsetfillopacity{0.160000}%
\pgfsetlinewidth{0.000000pt}%
\definecolor{currentstroke}{rgb}{0.000000,0.000000,0.000000}%
\pgfsetstrokecolor{currentstroke}%
\pgfsetdash{}{0pt}%
\pgfpathmoveto{\pgfqpoint{1.274795in}{0.779674in}}%
\pgfpathlineto{\pgfqpoint{1.288067in}{0.791836in}}%
\pgfpathlineto{\pgfqpoint{1.301350in}{0.804009in}}%
\pgfpathlineto{\pgfqpoint{1.299440in}{0.809764in}}%
\pgfpathlineto{\pgfqpoint{1.296626in}{0.815428in}}%
\pgfpathlineto{\pgfqpoint{1.283681in}{0.802577in}}%
\pgfpathlineto{\pgfqpoint{1.270742in}{0.789733in}}%
\pgfpathlineto{\pgfqpoint{1.273165in}{0.784742in}}%
\pgfpathlineto{\pgfqpoint{1.274795in}{0.779674in}}%
\pgfpathclose%
\pgfusepath{fill}%
\end{pgfscope}%
\begin{pgfscope}%
\pgfpathrectangle{\pgfqpoint{0.407592in}{0.222904in}}{\pgfqpoint{1.411163in}{1.411163in}}%
\pgfusepath{clip}%
\pgfsetbuttcap%
\pgfsetroundjoin%
\definecolor{currentfill}{rgb}{0.749020,0.772549,0.792157}%
\pgfsetfillcolor{currentfill}%
\pgfsetfillopacity{0.160000}%
\pgfsetlinewidth{0.000000pt}%
\definecolor{currentstroke}{rgb}{0.000000,0.000000,0.000000}%
\pgfsetstrokecolor{currentstroke}%
\pgfsetdash{}{0pt}%
\pgfpathmoveto{\pgfqpoint{1.582955in}{1.099674in}}%
\pgfpathlineto{\pgfqpoint{1.596040in}{1.112664in}}%
\pgfpathlineto{\pgfqpoint{1.609132in}{1.125661in}}%
\pgfpathlineto{\pgfqpoint{1.598558in}{1.138295in}}%
\pgfpathlineto{\pgfqpoint{1.585897in}{1.150554in}}%
\pgfpathlineto{\pgfqpoint{1.573491in}{1.137038in}}%
\pgfpathlineto{\pgfqpoint{1.561085in}{1.123522in}}%
\pgfpathlineto{\pgfqpoint{1.573015in}{1.111776in}}%
\pgfpathlineto{\pgfqpoint{1.582955in}{1.099674in}}%
\pgfpathclose%
\pgfusepath{fill}%
\end{pgfscope}%
\begin{pgfscope}%
\pgfpathrectangle{\pgfqpoint{0.407592in}{0.222904in}}{\pgfqpoint{1.411163in}{1.411163in}}%
\pgfusepath{clip}%
\pgfsetbuttcap%
\pgfsetroundjoin%
\definecolor{currentfill}{rgb}{0.749020,0.772549,0.792157}%
\pgfsetfillcolor{currentfill}%
\pgfsetfillopacity{0.160000}%
\pgfsetlinewidth{0.000000pt}%
\definecolor{currentstroke}{rgb}{0.000000,0.000000,0.000000}%
\pgfsetstrokecolor{currentstroke}%
\pgfsetdash{}{0pt}%
\pgfpathmoveto{\pgfqpoint{0.985520in}{0.706331in}}%
\pgfpathlineto{\pgfqpoint{0.972376in}{0.718356in}}%
\pgfpathlineto{\pgfqpoint{0.959220in}{0.730392in}}%
\pgfpathlineto{\pgfqpoint{0.958026in}{0.726681in}}%
\pgfpathlineto{\pgfqpoint{0.957406in}{0.722952in}}%
\pgfpathlineto{\pgfqpoint{0.970770in}{0.711641in}}%
\pgfpathlineto{\pgfqpoint{0.984116in}{0.700345in}}%
\pgfpathlineto{\pgfqpoint{0.984588in}{0.703345in}}%
\pgfpathlineto{\pgfqpoint{0.985520in}{0.706331in}}%
\pgfpathclose%
\pgfusepath{fill}%
\end{pgfscope}%
\begin{pgfscope}%
\pgfpathrectangle{\pgfqpoint{0.407592in}{0.222904in}}{\pgfqpoint{1.411163in}{1.411163in}}%
\pgfusepath{clip}%
\pgfsetbuttcap%
\pgfsetroundjoin%
\definecolor{currentfill}{rgb}{0.749020,0.772549,0.792157}%
\pgfsetfillcolor{currentfill}%
\pgfsetfillopacity{0.160000}%
\pgfsetlinewidth{0.000000pt}%
\definecolor{currentstroke}{rgb}{0.000000,0.000000,0.000000}%
\pgfsetstrokecolor{currentstroke}%
\pgfsetdash{}{0pt}%
\pgfpathmoveto{\pgfqpoint{1.068981in}{0.625862in}}%
\pgfpathlineto{\pgfqpoint{1.058303in}{0.633564in}}%
\pgfpathlineto{\pgfqpoint{1.047590in}{0.641293in}}%
\pgfpathlineto{\pgfqpoint{1.049763in}{0.640141in}}%
\pgfpathlineto{\pgfqpoint{1.052114in}{0.639058in}}%
\pgfpathlineto{\pgfqpoint{1.061695in}{0.631887in}}%
\pgfpathlineto{\pgfqpoint{1.071242in}{0.624743in}}%
\pgfpathlineto{\pgfqpoint{1.070067in}{0.625285in}}%
\pgfpathlineto{\pgfqpoint{1.068981in}{0.625862in}}%
\pgfpathclose%
\pgfusepath{fill}%
\end{pgfscope}%
\begin{pgfscope}%
\pgfpathrectangle{\pgfqpoint{0.407592in}{0.222904in}}{\pgfqpoint{1.411163in}{1.411163in}}%
\pgfusepath{clip}%
\pgfsetbuttcap%
\pgfsetroundjoin%
\definecolor{currentfill}{rgb}{0.749020,0.772549,0.792157}%
\pgfsetfillcolor{currentfill}%
\pgfsetfillopacity{0.160000}%
\pgfsetlinewidth{0.000000pt}%
\definecolor{currentstroke}{rgb}{0.000000,0.000000,0.000000}%
\pgfsetstrokecolor{currentstroke}%
\pgfsetdash{}{0pt}%
\pgfpathmoveto{\pgfqpoint{0.679482in}{1.067351in}}%
\pgfpathlineto{\pgfqpoint{0.667380in}{1.080811in}}%
\pgfpathlineto{\pgfqpoint{0.655277in}{1.094271in}}%
\pgfpathlineto{\pgfqpoint{0.641847in}{1.082977in}}%
\pgfpathlineto{\pgfqpoint{0.630246in}{1.071357in}}%
\pgfpathlineto{\pgfqpoint{0.643138in}{1.058442in}}%
\pgfpathlineto{\pgfqpoint{0.656023in}{1.045533in}}%
\pgfpathlineto{\pgfqpoint{0.666885in}{1.056595in}}%
\pgfpathlineto{\pgfqpoint{0.679482in}{1.067351in}}%
\pgfpathclose%
\pgfusepath{fill}%
\end{pgfscope}%
\begin{pgfscope}%
\pgfpathrectangle{\pgfqpoint{0.407592in}{0.222904in}}{\pgfqpoint{1.411163in}{1.411163in}}%
\pgfusepath{clip}%
\pgfsetbuttcap%
\pgfsetroundjoin%
\definecolor{currentfill}{rgb}{0.749020,0.772549,0.792157}%
\pgfsetfillcolor{currentfill}%
\pgfsetfillopacity{0.160000}%
\pgfsetlinewidth{0.000000pt}%
\definecolor{currentstroke}{rgb}{0.000000,0.000000,0.000000}%
\pgfsetstrokecolor{currentstroke}%
\pgfsetdash{}{0pt}%
\pgfpathmoveto{\pgfqpoint{1.168111in}{0.673739in}}%
\pgfpathlineto{\pgfqpoint{1.181171in}{0.684338in}}%
\pgfpathlineto{\pgfqpoint{1.194254in}{0.694956in}}%
\pgfpathlineto{\pgfqpoint{1.195237in}{0.697947in}}%
\pgfpathlineto{\pgfqpoint{1.195759in}{0.700948in}}%
\pgfpathlineto{\pgfqpoint{1.182507in}{0.689594in}}%
\pgfpathlineto{\pgfqpoint{1.169272in}{0.678255in}}%
\pgfpathlineto{\pgfqpoint{1.168865in}{0.675993in}}%
\pgfpathlineto{\pgfqpoint{1.168111in}{0.673739in}}%
\pgfpathclose%
\pgfusepath{fill}%
\end{pgfscope}%
\begin{pgfscope}%
\pgfpathrectangle{\pgfqpoint{0.407592in}{0.222904in}}{\pgfqpoint{1.411163in}{1.411163in}}%
\pgfusepath{clip}%
\pgfsetbuttcap%
\pgfsetroundjoin%
\definecolor{currentfill}{rgb}{0.749020,0.772549,0.792157}%
\pgfsetfillcolor{currentfill}%
\pgfsetfillopacity{0.160000}%
\pgfsetlinewidth{0.000000pt}%
\definecolor{currentstroke}{rgb}{0.000000,0.000000,0.000000}%
\pgfsetstrokecolor{currentstroke}%
\pgfsetdash{}{0pt}%
\pgfpathmoveto{\pgfqpoint{1.609132in}{1.125661in}}%
\pgfpathlineto{\pgfqpoint{1.622229in}{1.138663in}}%
\pgfpathlineto{\pgfqpoint{1.635333in}{1.151671in}}%
\pgfpathlineto{\pgfqpoint{1.624114in}{1.164828in}}%
\pgfpathlineto{\pgfqpoint{1.610711in}{1.177589in}}%
\pgfpathlineto{\pgfqpoint{1.598303in}{1.164071in}}%
\pgfpathlineto{\pgfqpoint{1.585897in}{1.150554in}}%
\pgfpathlineto{\pgfqpoint{1.598558in}{1.138295in}}%
\pgfpathlineto{\pgfqpoint{1.609132in}{1.125661in}}%
\pgfpathclose%
\pgfusepath{fill}%
\end{pgfscope}%
\begin{pgfscope}%
\pgfpathrectangle{\pgfqpoint{0.407592in}{0.222904in}}{\pgfqpoint{1.411163in}{1.411163in}}%
\pgfusepath{clip}%
\pgfsetbuttcap%
\pgfsetroundjoin%
\definecolor{currentfill}{rgb}{0.749020,0.772549,0.792157}%
\pgfsetfillcolor{currentfill}%
\pgfsetfillopacity{0.160000}%
\pgfsetlinewidth{0.000000pt}%
\definecolor{currentstroke}{rgb}{0.000000,0.000000,0.000000}%
\pgfsetstrokecolor{currentstroke}%
\pgfsetdash{}{0pt}%
\pgfpathmoveto{\pgfqpoint{0.912349in}{0.788739in}}%
\pgfpathlineto{\pgfqpoint{0.899595in}{0.801516in}}%
\pgfpathlineto{\pgfqpoint{0.886834in}{0.814300in}}%
\pgfpathlineto{\pgfqpoint{0.882986in}{0.808616in}}%
\pgfpathlineto{\pgfqpoint{0.880026in}{0.802846in}}%
\pgfpathlineto{\pgfqpoint{0.893256in}{0.790742in}}%
\pgfpathlineto{\pgfqpoint{0.906473in}{0.778649in}}%
\pgfpathlineto{\pgfqpoint{0.909022in}{0.783731in}}%
\pgfpathlineto{\pgfqpoint{0.912349in}{0.788739in}}%
\pgfpathclose%
\pgfusepath{fill}%
\end{pgfscope}%
\begin{pgfscope}%
\pgfpathrectangle{\pgfqpoint{0.407592in}{0.222904in}}{\pgfqpoint{1.411163in}{1.411163in}}%
\pgfusepath{clip}%
\pgfsetbuttcap%
\pgfsetroundjoin%
\definecolor{currentfill}{rgb}{0.749020,0.772549,0.792157}%
\pgfsetfillcolor{currentfill}%
\pgfsetfillopacity{0.160000}%
\pgfsetlinewidth{0.000000pt}%
\definecolor{currentstroke}{rgb}{0.000000,0.000000,0.000000}%
\pgfsetstrokecolor{currentstroke}%
\pgfsetdash{}{0pt}%
\pgfpathmoveto{\pgfqpoint{1.106312in}{0.624336in}}%
\pgfpathlineto{\pgfqpoint{1.114399in}{0.631277in}}%
\pgfpathlineto{\pgfqpoint{1.122516in}{0.638244in}}%
\pgfpathlineto{\pgfqpoint{1.125213in}{0.639269in}}%
\pgfpathlineto{\pgfqpoint{1.127755in}{0.640367in}}%
\pgfpathlineto{\pgfqpoint{1.118325in}{0.632869in}}%
\pgfpathlineto{\pgfqpoint{1.108927in}{0.625398in}}%
\pgfpathlineto{\pgfqpoint{1.107658in}{0.624849in}}%
\pgfpathlineto{\pgfqpoint{1.106312in}{0.624336in}}%
\pgfpathclose%
\pgfusepath{fill}%
\end{pgfscope}%
\begin{pgfscope}%
\pgfpathrectangle{\pgfqpoint{0.407592in}{0.222904in}}{\pgfqpoint{1.411163in}{1.411163in}}%
\pgfusepath{clip}%
\pgfsetbuttcap%
\pgfsetroundjoin%
\definecolor{currentfill}{rgb}{0.749020,0.772549,0.792157}%
\pgfsetfillcolor{currentfill}%
\pgfsetfillopacity{0.160000}%
\pgfsetlinewidth{0.000000pt}%
\definecolor{currentstroke}{rgb}{0.000000,0.000000,0.000000}%
\pgfsetstrokecolor{currentstroke}%
\pgfsetdash{}{0pt}%
\pgfpathmoveto{\pgfqpoint{1.010753in}{0.677801in}}%
\pgfpathlineto{\pgfqpoint{0.997444in}{0.689065in}}%
\pgfpathlineto{\pgfqpoint{0.984116in}{0.700345in}}%
\pgfpathlineto{\pgfqpoint{0.984107in}{0.697345in}}%
\pgfpathlineto{\pgfqpoint{0.984562in}{0.694358in}}%
\pgfpathlineto{\pgfqpoint{0.997853in}{0.683814in}}%
\pgfpathlineto{\pgfqpoint{1.011121in}{0.673289in}}%
\pgfpathlineto{\pgfqpoint{1.010763in}{0.675540in}}%
\pgfpathlineto{\pgfqpoint{1.010753in}{0.677801in}}%
\pgfpathclose%
\pgfusepath{fill}%
\end{pgfscope}%
\begin{pgfscope}%
\pgfpathrectangle{\pgfqpoint{0.407592in}{0.222904in}}{\pgfqpoint{1.411163in}{1.411163in}}%
\pgfusepath{clip}%
\pgfsetbuttcap%
\pgfsetroundjoin%
\definecolor{currentfill}{rgb}{0.749020,0.772549,0.792157}%
\pgfsetfillcolor{currentfill}%
\pgfsetfillopacity{0.160000}%
\pgfsetlinewidth{0.000000pt}%
\definecolor{currentstroke}{rgb}{0.000000,0.000000,0.000000}%
\pgfsetstrokecolor{currentstroke}%
\pgfsetdash{}{0pt}%
\pgfpathmoveto{\pgfqpoint{0.655277in}{1.094271in}}%
\pgfpathlineto{\pgfqpoint{0.643173in}{1.107733in}}%
\pgfpathlineto{\pgfqpoint{0.631067in}{1.121196in}}%
\pgfpathlineto{\pgfqpoint{0.616794in}{1.109375in}}%
\pgfpathlineto{\pgfqpoint{0.604443in}{1.097208in}}%
\pgfpathlineto{\pgfqpoint{0.617348in}{1.084279in}}%
\pgfpathlineto{\pgfqpoint{0.630246in}{1.071357in}}%
\pgfpathlineto{\pgfqpoint{0.641847in}{1.082977in}}%
\pgfpathlineto{\pgfqpoint{0.655277in}{1.094271in}}%
\pgfpathclose%
\pgfusepath{fill}%
\end{pgfscope}%
\begin{pgfscope}%
\pgfpathrectangle{\pgfqpoint{0.407592in}{0.222904in}}{\pgfqpoint{1.411163in}{1.411163in}}%
\pgfusepath{clip}%
\pgfsetbuttcap%
\pgfsetroundjoin%
\definecolor{currentfill}{rgb}{0.749020,0.772549,0.792157}%
\pgfsetfillcolor{currentfill}%
\pgfsetfillopacity{0.160000}%
\pgfsetlinewidth{0.000000pt}%
\definecolor{currentstroke}{rgb}{0.000000,0.000000,0.000000}%
\pgfsetstrokecolor{currentstroke}%
\pgfsetdash{}{0pt}%
\pgfpathmoveto{\pgfqpoint{1.635333in}{1.151671in}}%
\pgfpathlineto{\pgfqpoint{1.648443in}{1.164686in}}%
\pgfpathlineto{\pgfqpoint{1.661559in}{1.177707in}}%
\pgfpathlineto{\pgfqpoint{1.649683in}{1.191375in}}%
\pgfpathlineto{\pgfqpoint{1.635528in}{1.204627in}}%
\pgfpathlineto{\pgfqpoint{1.623119in}{1.191107in}}%
\pgfpathlineto{\pgfqpoint{1.610711in}{1.177589in}}%
\pgfpathlineto{\pgfqpoint{1.624114in}{1.164828in}}%
\pgfpathlineto{\pgfqpoint{1.635333in}{1.151671in}}%
\pgfpathclose%
\pgfusepath{fill}%
\end{pgfscope}%
\begin{pgfscope}%
\pgfpathrectangle{\pgfqpoint{0.407592in}{0.222904in}}{\pgfqpoint{1.411163in}{1.411163in}}%
\pgfusepath{clip}%
\pgfsetbuttcap%
\pgfsetroundjoin%
\definecolor{currentfill}{rgb}{0.749020,0.772549,0.792157}%
\pgfsetfillcolor{currentfill}%
\pgfsetfillopacity{0.160000}%
\pgfsetlinewidth{0.000000pt}%
\definecolor{currentstroke}{rgb}{0.000000,0.000000,0.000000}%
\pgfsetstrokecolor{currentstroke}%
\pgfsetdash{}{0pt}%
\pgfpathmoveto{\pgfqpoint{1.071242in}{0.624743in}}%
\pgfpathlineto{\pgfqpoint{1.061695in}{0.631887in}}%
\pgfpathlineto{\pgfqpoint{1.052114in}{0.639058in}}%
\pgfpathlineto{\pgfqpoint{1.054634in}{0.638047in}}%
\pgfpathlineto{\pgfqpoint{1.057311in}{0.637114in}}%
\pgfpathlineto{\pgfqpoint{1.065590in}{0.630430in}}%
\pgfpathlineto{\pgfqpoint{1.073835in}{0.623771in}}%
\pgfpathlineto{\pgfqpoint{1.072499in}{0.624238in}}%
\pgfpathlineto{\pgfqpoint{1.071242in}{0.624743in}}%
\pgfpathclose%
\pgfusepath{fill}%
\end{pgfscope}%
\begin{pgfscope}%
\pgfpathrectangle{\pgfqpoint{0.407592in}{0.222904in}}{\pgfqpoint{1.411163in}{1.411163in}}%
\pgfusepath{clip}%
\pgfsetbuttcap%
\pgfsetroundjoin%
\definecolor{currentfill}{rgb}{0.749020,0.772549,0.792157}%
\pgfsetfillcolor{currentfill}%
\pgfsetfillopacity{0.160000}%
\pgfsetlinewidth{0.000000pt}%
\definecolor{currentstroke}{rgb}{0.000000,0.000000,0.000000}%
\pgfsetstrokecolor{currentstroke}%
\pgfsetdash{}{0pt}%
\pgfpathmoveto{\pgfqpoint{1.301350in}{0.804009in}}%
\pgfpathlineto{\pgfqpoint{1.314645in}{0.816193in}}%
\pgfpathlineto{\pgfqpoint{1.327952in}{0.828388in}}%
\pgfpathlineto{\pgfqpoint{1.325751in}{0.834820in}}%
\pgfpathlineto{\pgfqpoint{1.322534in}{0.841148in}}%
\pgfpathlineto{\pgfqpoint{1.309577in}{0.828285in}}%
\pgfpathlineto{\pgfqpoint{1.296626in}{0.815428in}}%
\pgfpathlineto{\pgfqpoint{1.299440in}{0.809764in}}%
\pgfpathlineto{\pgfqpoint{1.301350in}{0.804009in}}%
\pgfpathclose%
\pgfusepath{fill}%
\end{pgfscope}%
\begin{pgfscope}%
\pgfpathrectangle{\pgfqpoint{0.407592in}{0.222904in}}{\pgfqpoint{1.411163in}{1.411163in}}%
\pgfusepath{clip}%
\pgfsetbuttcap%
\pgfsetroundjoin%
\definecolor{currentfill}{rgb}{0.749020,0.772549,0.792157}%
\pgfsetfillcolor{currentfill}%
\pgfsetfillopacity{0.160000}%
\pgfsetlinewidth{0.000000pt}%
\definecolor{currentstroke}{rgb}{0.000000,0.000000,0.000000}%
\pgfsetstrokecolor{currentstroke}%
\pgfsetdash{}{0pt}%
\pgfpathmoveto{\pgfqpoint{1.103412in}{0.623428in}}%
\pgfpathlineto{\pgfqpoint{1.110044in}{0.629915in}}%
\pgfpathlineto{\pgfqpoint{1.116702in}{0.636428in}}%
\pgfpathlineto{\pgfqpoint{1.119675in}{0.637295in}}%
\pgfpathlineto{\pgfqpoint{1.122516in}{0.638244in}}%
\pgfpathlineto{\pgfqpoint{1.114399in}{0.631277in}}%
\pgfpathlineto{\pgfqpoint{1.106312in}{0.624336in}}%
\pgfpathlineto{\pgfqpoint{1.104895in}{0.623862in}}%
\pgfpathlineto{\pgfqpoint{1.103412in}{0.623428in}}%
\pgfpathclose%
\pgfusepath{fill}%
\end{pgfscope}%
\begin{pgfscope}%
\pgfpathrectangle{\pgfqpoint{0.407592in}{0.222904in}}{\pgfqpoint{1.411163in}{1.411163in}}%
\pgfusepath{clip}%
\pgfsetbuttcap%
\pgfsetroundjoin%
\definecolor{currentfill}{rgb}{0.749020,0.772549,0.792157}%
\pgfsetfillcolor{currentfill}%
\pgfsetfillopacity{0.160000}%
\pgfsetlinewidth{0.000000pt}%
\definecolor{currentstroke}{rgb}{0.000000,0.000000,0.000000}%
\pgfsetstrokecolor{currentstroke}%
\pgfsetdash{}{0pt}%
\pgfpathmoveto{\pgfqpoint{1.137766in}{0.646756in}}%
\pgfpathlineto{\pgfqpoint{1.149727in}{0.655870in}}%
\pgfpathlineto{\pgfqpoint{1.161718in}{0.665008in}}%
\pgfpathlineto{\pgfqpoint{1.163812in}{0.667129in}}%
\pgfpathlineto{\pgfqpoint{1.165580in}{0.669297in}}%
\pgfpathlineto{\pgfqpoint{1.152954in}{0.659450in}}%
\pgfpathlineto{\pgfqpoint{1.140354in}{0.649624in}}%
\pgfpathlineto{\pgfqpoint{1.139168in}{0.648174in}}%
\pgfpathlineto{\pgfqpoint{1.137766in}{0.646756in}}%
\pgfpathclose%
\pgfusepath{fill}%
\end{pgfscope}%
\begin{pgfscope}%
\pgfpathrectangle{\pgfqpoint{0.407592in}{0.222904in}}{\pgfqpoint{1.411163in}{1.411163in}}%
\pgfusepath{clip}%
\pgfsetbuttcap%
\pgfsetroundjoin%
\definecolor{currentfill}{rgb}{0.749020,0.772549,0.792157}%
\pgfsetfillcolor{currentfill}%
\pgfsetfillopacity{0.160000}%
\pgfsetlinewidth{0.000000pt}%
\definecolor{currentstroke}{rgb}{0.000000,0.000000,0.000000}%
\pgfsetstrokecolor{currentstroke}%
\pgfsetdash{}{0pt}%
\pgfpathmoveto{\pgfqpoint{1.661559in}{1.177707in}}%
\pgfpathlineto{\pgfqpoint{1.674681in}{1.190733in}}%
\pgfpathlineto{\pgfqpoint{1.687810in}{1.203766in}}%
\pgfpathlineto{\pgfqpoint{1.675266in}{1.217936in}}%
\pgfpathlineto{\pgfqpoint{1.660348in}{1.231668in}}%
\pgfpathlineto{\pgfqpoint{1.647938in}{1.218147in}}%
\pgfpathlineto{\pgfqpoint{1.635528in}{1.204627in}}%
\pgfpathlineto{\pgfqpoint{1.649683in}{1.191375in}}%
\pgfpathlineto{\pgfqpoint{1.661559in}{1.177707in}}%
\pgfpathclose%
\pgfusepath{fill}%
\end{pgfscope}%
\begin{pgfscope}%
\pgfpathrectangle{\pgfqpoint{0.407592in}{0.222904in}}{\pgfqpoint{1.411163in}{1.411163in}}%
\pgfusepath{clip}%
\pgfsetbuttcap%
\pgfsetroundjoin%
\definecolor{currentfill}{rgb}{0.749020,0.772549,0.792157}%
\pgfsetfillcolor{currentfill}%
\pgfsetfillopacity{0.160000}%
\pgfsetlinewidth{0.000000pt}%
\definecolor{currentstroke}{rgb}{0.000000,0.000000,0.000000}%
\pgfsetstrokecolor{currentstroke}%
\pgfsetdash{}{0pt}%
\pgfpathmoveto{\pgfqpoint{0.631067in}{1.121196in}}%
\pgfpathlineto{\pgfqpoint{0.618961in}{1.134661in}}%
\pgfpathlineto{\pgfqpoint{0.606853in}{1.148127in}}%
\pgfpathlineto{\pgfqpoint{0.591725in}{1.135789in}}%
\pgfpathlineto{\pgfqpoint{0.578613in}{1.123085in}}%
\pgfpathlineto{\pgfqpoint{0.591531in}{1.110143in}}%
\pgfpathlineto{\pgfqpoint{0.604443in}{1.097208in}}%
\pgfpathlineto{\pgfqpoint{0.616794in}{1.109375in}}%
\pgfpathlineto{\pgfqpoint{0.631067in}{1.121196in}}%
\pgfpathclose%
\pgfusepath{fill}%
\end{pgfscope}%
\begin{pgfscope}%
\pgfpathrectangle{\pgfqpoint{0.407592in}{0.222904in}}{\pgfqpoint{1.411163in}{1.411163in}}%
\pgfusepath{clip}%
\pgfsetbuttcap%
\pgfsetroundjoin%
\definecolor{currentfill}{rgb}{0.749020,0.772549,0.792157}%
\pgfsetfillcolor{currentfill}%
\pgfsetfillopacity{0.160000}%
\pgfsetlinewidth{0.000000pt}%
\definecolor{currentstroke}{rgb}{0.000000,0.000000,0.000000}%
\pgfsetstrokecolor{currentstroke}%
\pgfsetdash{}{0pt}%
\pgfpathmoveto{\pgfqpoint{1.038772in}{0.649331in}}%
\pgfpathlineto{\pgfqpoint{1.025837in}{0.659084in}}%
\pgfpathlineto{\pgfqpoint{1.012872in}{0.668859in}}%
\pgfpathlineto{\pgfqpoint{1.014260in}{0.666699in}}%
\pgfpathlineto{\pgfqpoint{1.015983in}{0.664589in}}%
\pgfpathlineto{\pgfqpoint{1.028438in}{0.655520in}}%
\pgfpathlineto{\pgfqpoint{1.040861in}{0.646476in}}%
\pgfpathlineto{\pgfqpoint{1.039706in}{0.647887in}}%
\pgfpathlineto{\pgfqpoint{1.038772in}{0.649331in}}%
\pgfpathclose%
\pgfusepath{fill}%
\end{pgfscope}%
\begin{pgfscope}%
\pgfpathrectangle{\pgfqpoint{0.407592in}{0.222904in}}{\pgfqpoint{1.411163in}{1.411163in}}%
\pgfusepath{clip}%
\pgfsetbuttcap%
\pgfsetroundjoin%
\definecolor{currentfill}{rgb}{0.749020,0.772549,0.792157}%
\pgfsetfillcolor{currentfill}%
\pgfsetfillopacity{0.160000}%
\pgfsetlinewidth{0.000000pt}%
\definecolor{currentstroke}{rgb}{0.000000,0.000000,0.000000}%
\pgfsetstrokecolor{currentstroke}%
\pgfsetdash{}{0pt}%
\pgfpathmoveto{\pgfqpoint{1.073835in}{0.623771in}}%
\pgfpathlineto{\pgfqpoint{1.065590in}{0.630430in}}%
\pgfpathlineto{\pgfqpoint{1.057311in}{0.637114in}}%
\pgfpathlineto{\pgfqpoint{1.060134in}{0.636264in}}%
\pgfpathlineto{\pgfqpoint{1.063090in}{0.635499in}}%
\pgfpathlineto{\pgfqpoint{1.069918in}{0.629218in}}%
\pgfpathlineto{\pgfqpoint{1.076716in}{0.622964in}}%
\pgfpathlineto{\pgfqpoint{1.075243in}{0.623346in}}%
\pgfpathlineto{\pgfqpoint{1.073835in}{0.623771in}}%
\pgfpathclose%
\pgfusepath{fill}%
\end{pgfscope}%
\begin{pgfscope}%
\pgfpathrectangle{\pgfqpoint{0.407592in}{0.222904in}}{\pgfqpoint{1.411163in}{1.411163in}}%
\pgfusepath{clip}%
\pgfsetbuttcap%
\pgfsetroundjoin%
\definecolor{currentfill}{rgb}{0.749020,0.772549,0.792157}%
\pgfsetfillcolor{currentfill}%
\pgfsetfillopacity{0.160000}%
\pgfsetlinewidth{0.000000pt}%
\definecolor{currentstroke}{rgb}{0.000000,0.000000,0.000000}%
\pgfsetstrokecolor{currentstroke}%
\pgfsetdash{}{0pt}%
\pgfpathmoveto{\pgfqpoint{1.222316in}{0.723701in}}%
\pgfpathlineto{\pgfqpoint{1.235621in}{0.735100in}}%
\pgfpathlineto{\pgfqpoint{1.248943in}{0.746514in}}%
\pgfpathlineto{\pgfqpoint{1.248962in}{0.750962in}}%
\pgfpathlineto{\pgfqpoint{1.248288in}{0.755381in}}%
\pgfpathlineto{\pgfqpoint{1.235052in}{0.743252in}}%
\pgfpathlineto{\pgfqpoint{1.221828in}{0.731133in}}%
\pgfpathlineto{\pgfqpoint{1.222361in}{0.727428in}}%
\pgfpathlineto{\pgfqpoint{1.222316in}{0.723701in}}%
\pgfpathclose%
\pgfusepath{fill}%
\end{pgfscope}%
\begin{pgfscope}%
\pgfpathrectangle{\pgfqpoint{0.407592in}{0.222904in}}{\pgfqpoint{1.411163in}{1.411163in}}%
\pgfusepath{clip}%
\pgfsetbuttcap%
\pgfsetroundjoin%
\definecolor{currentfill}{rgb}{0.749020,0.772549,0.792157}%
\pgfsetfillcolor{currentfill}%
\pgfsetfillopacity{0.160000}%
\pgfsetlinewidth{0.000000pt}%
\definecolor{currentstroke}{rgb}{0.000000,0.000000,0.000000}%
\pgfsetstrokecolor{currentstroke}%
\pgfsetdash{}{0pt}%
\pgfpathmoveto{\pgfqpoint{0.886834in}{0.814300in}}%
\pgfpathlineto{\pgfqpoint{0.874067in}{0.827091in}}%
\pgfpathlineto{\pgfqpoint{0.861293in}{0.839888in}}%
\pgfpathlineto{\pgfqpoint{0.856912in}{0.833537in}}%
\pgfpathlineto{\pgfqpoint{0.853529in}{0.827087in}}%
\pgfpathlineto{\pgfqpoint{0.866784in}{0.814961in}}%
\pgfpathlineto{\pgfqpoint{0.880026in}{0.802846in}}%
\pgfpathlineto{\pgfqpoint{0.882986in}{0.808616in}}%
\pgfpathlineto{\pgfqpoint{0.886834in}{0.814300in}}%
\pgfpathclose%
\pgfusepath{fill}%
\end{pgfscope}%
\begin{pgfscope}%
\pgfpathrectangle{\pgfqpoint{0.407592in}{0.222904in}}{\pgfqpoint{1.411163in}{1.411163in}}%
\pgfusepath{clip}%
\pgfsetbuttcap%
\pgfsetroundjoin%
\definecolor{currentfill}{rgb}{0.749020,0.772549,0.792157}%
\pgfsetfillcolor{currentfill}%
\pgfsetfillopacity{0.160000}%
\pgfsetlinewidth{0.000000pt}%
\definecolor{currentstroke}{rgb}{0.000000,0.000000,0.000000}%
\pgfsetstrokecolor{currentstroke}%
\pgfsetdash{}{0pt}%
\pgfpathmoveto{\pgfqpoint{1.687810in}{1.203766in}}%
\pgfpathlineto{\pgfqpoint{1.700944in}{1.216805in}}%
\pgfpathlineto{\pgfqpoint{1.714085in}{1.229851in}}%
\pgfpathlineto{\pgfqpoint{1.700862in}{1.244511in}}%
\pgfpathlineto{\pgfqpoint{1.685171in}{1.258712in}}%
\pgfpathlineto{\pgfqpoint{1.672759in}{1.245190in}}%
\pgfpathlineto{\pgfqpoint{1.660348in}{1.231668in}}%
\pgfpathlineto{\pgfqpoint{1.675266in}{1.217936in}}%
\pgfpathlineto{\pgfqpoint{1.687810in}{1.203766in}}%
\pgfpathclose%
\pgfusepath{fill}%
\end{pgfscope}%
\begin{pgfscope}%
\pgfpathrectangle{\pgfqpoint{0.407592in}{0.222904in}}{\pgfqpoint{1.411163in}{1.411163in}}%
\pgfusepath{clip}%
\pgfsetbuttcap%
\pgfsetroundjoin%
\definecolor{currentfill}{rgb}{0.749020,0.772549,0.792157}%
\pgfsetfillcolor{currentfill}%
\pgfsetfillopacity{0.160000}%
\pgfsetlinewidth{0.000000pt}%
\definecolor{currentstroke}{rgb}{0.000000,0.000000,0.000000}%
\pgfsetstrokecolor{currentstroke}%
\pgfsetdash{}{0pt}%
\pgfpathmoveto{\pgfqpoint{1.100280in}{0.622690in}}%
\pgfpathlineto{\pgfqpoint{1.105337in}{0.628808in}}%
\pgfpathlineto{\pgfqpoint{1.110416in}{0.634951in}}%
\pgfpathlineto{\pgfqpoint{1.113611in}{0.635645in}}%
\pgfpathlineto{\pgfqpoint{1.116702in}{0.636428in}}%
\pgfpathlineto{\pgfqpoint{1.110044in}{0.629915in}}%
\pgfpathlineto{\pgfqpoint{1.103412in}{0.623428in}}%
\pgfpathlineto{\pgfqpoint{1.101872in}{0.623037in}}%
\pgfpathlineto{\pgfqpoint{1.100280in}{0.622690in}}%
\pgfpathclose%
\pgfusepath{fill}%
\end{pgfscope}%
\begin{pgfscope}%
\pgfpathrectangle{\pgfqpoint{0.407592in}{0.222904in}}{\pgfqpoint{1.411163in}{1.411163in}}%
\pgfusepath{clip}%
\pgfsetbuttcap%
\pgfsetroundjoin%
\definecolor{currentfill}{rgb}{0.749020,0.772549,0.792157}%
\pgfsetfillcolor{currentfill}%
\pgfsetfillopacity{0.160000}%
\pgfsetlinewidth{0.000000pt}%
\definecolor{currentstroke}{rgb}{0.000000,0.000000,0.000000}%
\pgfsetstrokecolor{currentstroke}%
\pgfsetdash{}{0pt}%
\pgfpathmoveto{\pgfqpoint{0.606853in}{1.148127in}}%
\pgfpathlineto{\pgfqpoint{0.594744in}{1.161594in}}%
\pgfpathlineto{\pgfqpoint{0.582634in}{1.175063in}}%
\pgfpathlineto{\pgfqpoint{0.566641in}{1.162219in}}%
\pgfpathlineto{\pgfqpoint{0.552757in}{1.148988in}}%
\pgfpathlineto{\pgfqpoint{0.565688in}{1.136033in}}%
\pgfpathlineto{\pgfqpoint{0.578613in}{1.123085in}}%
\pgfpathlineto{\pgfqpoint{0.591725in}{1.135789in}}%
\pgfpathlineto{\pgfqpoint{0.606853in}{1.148127in}}%
\pgfpathclose%
\pgfusepath{fill}%
\end{pgfscope}%
\begin{pgfscope}%
\pgfpathrectangle{\pgfqpoint{0.407592in}{0.222904in}}{\pgfqpoint{1.411163in}{1.411163in}}%
\pgfusepath{clip}%
\pgfsetbuttcap%
\pgfsetroundjoin%
\definecolor{currentfill}{rgb}{0.749020,0.772549,0.792157}%
\pgfsetfillcolor{currentfill}%
\pgfsetfillopacity{0.160000}%
\pgfsetlinewidth{0.000000pt}%
\definecolor{currentstroke}{rgb}{0.000000,0.000000,0.000000}%
\pgfsetstrokecolor{currentstroke}%
\pgfsetdash{}{0pt}%
\pgfpathmoveto{\pgfqpoint{0.959220in}{0.730392in}}%
\pgfpathlineto{\pgfqpoint{0.946052in}{0.742439in}}%
\pgfpathlineto{\pgfqpoint{0.932871in}{0.754498in}}%
\pgfpathlineto{\pgfqpoint{0.931403in}{0.750071in}}%
\pgfpathlineto{\pgfqpoint{0.930624in}{0.745620in}}%
\pgfpathlineto{\pgfqpoint{0.944024in}{0.734278in}}%
\pgfpathlineto{\pgfqpoint{0.957406in}{0.722952in}}%
\pgfpathlineto{\pgfqpoint{0.958026in}{0.726681in}}%
\pgfpathlineto{\pgfqpoint{0.959220in}{0.730392in}}%
\pgfpathclose%
\pgfusepath{fill}%
\end{pgfscope}%
\begin{pgfscope}%
\pgfpathrectangle{\pgfqpoint{0.407592in}{0.222904in}}{\pgfqpoint{1.411163in}{1.411163in}}%
\pgfusepath{clip}%
\pgfsetbuttcap%
\pgfsetroundjoin%
\definecolor{currentfill}{rgb}{0.749020,0.772549,0.792157}%
\pgfsetfillcolor{currentfill}%
\pgfsetfillopacity{0.160000}%
\pgfsetlinewidth{0.000000pt}%
\definecolor{currentstroke}{rgb}{0.000000,0.000000,0.000000}%
\pgfsetstrokecolor{currentstroke}%
\pgfsetdash{}{0pt}%
\pgfpathmoveto{\pgfqpoint{1.076716in}{0.622964in}}%
\pgfpathlineto{\pgfqpoint{1.069918in}{0.629218in}}%
\pgfpathlineto{\pgfqpoint{1.063090in}{0.635499in}}%
\pgfpathlineto{\pgfqpoint{1.066167in}{0.634823in}}%
\pgfpathlineto{\pgfqpoint{1.069349in}{0.634239in}}%
\pgfpathlineto{\pgfqpoint{1.074603in}{0.628274in}}%
\pgfpathlineto{\pgfqpoint{1.079835in}{0.622335in}}%
\pgfpathlineto{\pgfqpoint{1.078249in}{0.622626in}}%
\pgfpathlineto{\pgfqpoint{1.076716in}{0.622964in}}%
\pgfpathclose%
\pgfusepath{fill}%
\end{pgfscope}%
\begin{pgfscope}%
\pgfpathrectangle{\pgfqpoint{0.407592in}{0.222904in}}{\pgfqpoint{1.411163in}{1.411163in}}%
\pgfusepath{clip}%
\pgfsetbuttcap%
\pgfsetroundjoin%
\definecolor{currentfill}{rgb}{0.749020,0.772549,0.792157}%
\pgfsetfillcolor{currentfill}%
\pgfsetfillopacity{0.160000}%
\pgfsetlinewidth{0.000000pt}%
\definecolor{currentstroke}{rgb}{0.000000,0.000000,0.000000}%
\pgfsetstrokecolor{currentstroke}%
\pgfsetdash{}{0pt}%
\pgfpathmoveto{\pgfqpoint{1.327952in}{0.828388in}}%
\pgfpathlineto{\pgfqpoint{1.341271in}{0.840594in}}%
\pgfpathlineto{\pgfqpoint{1.354602in}{0.852811in}}%
\pgfpathlineto{\pgfqpoint{1.352098in}{0.859910in}}%
\pgfpathlineto{\pgfqpoint{1.348467in}{0.866892in}}%
\pgfpathlineto{\pgfqpoint{1.335497in}{0.854017in}}%
\pgfpathlineto{\pgfqpoint{1.322534in}{0.841148in}}%
\pgfpathlineto{\pgfqpoint{1.325751in}{0.834820in}}%
\pgfpathlineto{\pgfqpoint{1.327952in}{0.828388in}}%
\pgfpathclose%
\pgfusepath{fill}%
\end{pgfscope}%
\begin{pgfscope}%
\pgfpathrectangle{\pgfqpoint{0.407592in}{0.222904in}}{\pgfqpoint{1.411163in}{1.411163in}}%
\pgfusepath{clip}%
\pgfsetbuttcap%
\pgfsetroundjoin%
\definecolor{currentfill}{rgb}{0.749020,0.772549,0.792157}%
\pgfsetfillcolor{currentfill}%
\pgfsetfillopacity{0.160000}%
\pgfsetlinewidth{0.000000pt}%
\definecolor{currentstroke}{rgb}{0.000000,0.000000,0.000000}%
\pgfsetstrokecolor{currentstroke}%
\pgfsetdash{}{0pt}%
\pgfpathmoveto{\pgfqpoint{1.714085in}{1.229851in}}%
\pgfpathlineto{\pgfqpoint{1.727232in}{1.242902in}}%
\pgfpathlineto{\pgfqpoint{1.740385in}{1.255959in}}%
\pgfpathlineto{\pgfqpoint{1.726472in}{1.271100in}}%
\pgfpathlineto{\pgfqpoint{1.709998in}{1.285760in}}%
\pgfpathlineto{\pgfqpoint{1.697584in}{1.272236in}}%
\pgfpathlineto{\pgfqpoint{1.685171in}{1.258712in}}%
\pgfpathlineto{\pgfqpoint{1.700862in}{1.244511in}}%
\pgfpathlineto{\pgfqpoint{1.714085in}{1.229851in}}%
\pgfpathclose%
\pgfusepath{fill}%
\end{pgfscope}%
\begin{pgfscope}%
\pgfpathrectangle{\pgfqpoint{0.407592in}{0.222904in}}{\pgfqpoint{1.411163in}{1.411163in}}%
\pgfusepath{clip}%
\pgfsetbuttcap%
\pgfsetroundjoin%
\definecolor{currentfill}{rgb}{0.749020,0.772549,0.792157}%
\pgfsetfillcolor{currentfill}%
\pgfsetfillopacity{0.160000}%
\pgfsetlinewidth{0.000000pt}%
\definecolor{currentstroke}{rgb}{0.000000,0.000000,0.000000}%
\pgfsetstrokecolor{currentstroke}%
\pgfsetdash{}{0pt}%
\pgfpathmoveto{\pgfqpoint{1.096969in}{0.622136in}}%
\pgfpathlineto{\pgfqpoint{1.100361in}{0.627975in}}%
\pgfpathlineto{\pgfqpoint{1.103768in}{0.633841in}}%
\pgfpathlineto{\pgfqpoint{1.107130in}{0.634349in}}%
\pgfpathlineto{\pgfqpoint{1.110416in}{0.634951in}}%
\pgfpathlineto{\pgfqpoint{1.105337in}{0.628808in}}%
\pgfpathlineto{\pgfqpoint{1.100280in}{0.622690in}}%
\pgfpathlineto{\pgfqpoint{1.098643in}{0.622389in}}%
\pgfpathlineto{\pgfqpoint{1.096969in}{0.622136in}}%
\pgfpathclose%
\pgfusepath{fill}%
\end{pgfscope}%
\begin{pgfscope}%
\pgfpathrectangle{\pgfqpoint{0.407592in}{0.222904in}}{\pgfqpoint{1.411163in}{1.411163in}}%
\pgfusepath{clip}%
\pgfsetbuttcap%
\pgfsetroundjoin%
\definecolor{currentfill}{rgb}{0.749020,0.772549,0.792157}%
\pgfsetfillcolor{currentfill}%
\pgfsetfillopacity{0.160000}%
\pgfsetlinewidth{0.000000pt}%
\definecolor{currentstroke}{rgb}{0.000000,0.000000,0.000000}%
\pgfsetstrokecolor{currentstroke}%
\pgfsetdash{}{0pt}%
\pgfpathmoveto{\pgfqpoint{0.582634in}{1.175063in}}%
\pgfpathlineto{\pgfqpoint{0.570522in}{1.188533in}}%
\pgfpathlineto{\pgfqpoint{0.558409in}{1.202004in}}%
\pgfpathlineto{\pgfqpoint{0.541542in}{1.188665in}}%
\pgfpathlineto{\pgfqpoint{0.526874in}{1.174919in}}%
\pgfpathlineto{\pgfqpoint{0.539819in}{1.161950in}}%
\pgfpathlineto{\pgfqpoint{0.552757in}{1.148988in}}%
\pgfpathlineto{\pgfqpoint{0.566641in}{1.162219in}}%
\pgfpathlineto{\pgfqpoint{0.582634in}{1.175063in}}%
\pgfpathclose%
\pgfusepath{fill}%
\end{pgfscope}%
\begin{pgfscope}%
\pgfpathrectangle{\pgfqpoint{0.407592in}{0.222904in}}{\pgfqpoint{1.411163in}{1.411163in}}%
\pgfusepath{clip}%
\pgfsetbuttcap%
\pgfsetroundjoin%
\definecolor{currentfill}{rgb}{0.749020,0.772549,0.792157}%
\pgfsetfillcolor{currentfill}%
\pgfsetfillopacity{0.160000}%
\pgfsetlinewidth{0.000000pt}%
\definecolor{currentstroke}{rgb}{0.000000,0.000000,0.000000}%
\pgfsetstrokecolor{currentstroke}%
\pgfsetdash{}{0pt}%
\pgfpathmoveto{\pgfqpoint{1.079835in}{0.622335in}}%
\pgfpathlineto{\pgfqpoint{1.074603in}{0.628274in}}%
\pgfpathlineto{\pgfqpoint{1.069349in}{0.634239in}}%
\pgfpathlineto{\pgfqpoint{1.072624in}{0.633751in}}%
\pgfpathlineto{\pgfqpoint{1.075977in}{0.633359in}}%
\pgfpathlineto{\pgfqpoint{1.079564in}{0.627614in}}%
\pgfpathlineto{\pgfqpoint{1.083135in}{0.621895in}}%
\pgfpathlineto{\pgfqpoint{1.081466in}{0.622091in}}%
\pgfpathlineto{\pgfqpoint{1.079835in}{0.622335in}}%
\pgfpathclose%
\pgfusepath{fill}%
\end{pgfscope}%
\begin{pgfscope}%
\pgfpathrectangle{\pgfqpoint{0.407592in}{0.222904in}}{\pgfqpoint{1.411163in}{1.411163in}}%
\pgfusepath{clip}%
\pgfsetbuttcap%
\pgfsetroundjoin%
\definecolor{currentfill}{rgb}{0.749020,0.772549,0.792157}%
\pgfsetfillcolor{currentfill}%
\pgfsetfillopacity{0.160000}%
\pgfsetlinewidth{0.000000pt}%
\definecolor{currentstroke}{rgb}{0.000000,0.000000,0.000000}%
\pgfsetstrokecolor{currentstroke}%
\pgfsetdash{}{0pt}%
\pgfpathmoveto{\pgfqpoint{1.093538in}{0.621774in}}%
\pgfpathlineto{\pgfqpoint{1.095203in}{0.627433in}}%
\pgfpathlineto{\pgfqpoint{1.096876in}{0.633117in}}%
\pgfpathlineto{\pgfqpoint{1.100345in}{0.633430in}}%
\pgfpathlineto{\pgfqpoint{1.103768in}{0.633841in}}%
\pgfpathlineto{\pgfqpoint{1.100361in}{0.627975in}}%
\pgfpathlineto{\pgfqpoint{1.096969in}{0.622136in}}%
\pgfpathlineto{\pgfqpoint{1.095265in}{0.621930in}}%
\pgfpathlineto{\pgfqpoint{1.093538in}{0.621774in}}%
\pgfpathclose%
\pgfusepath{fill}%
\end{pgfscope}%
\begin{pgfscope}%
\pgfpathrectangle{\pgfqpoint{0.407592in}{0.222904in}}{\pgfqpoint{1.411163in}{1.411163in}}%
\pgfusepath{clip}%
\pgfsetbuttcap%
\pgfsetroundjoin%
\definecolor{currentfill}{rgb}{0.749020,0.772549,0.792157}%
\pgfsetfillcolor{currentfill}%
\pgfsetfillopacity{0.160000}%
\pgfsetlinewidth{0.000000pt}%
\definecolor{currentstroke}{rgb}{0.000000,0.000000,0.000000}%
\pgfsetstrokecolor{currentstroke}%
\pgfsetdash{}{0pt}%
\pgfpathmoveto{\pgfqpoint{1.740385in}{1.255959in}}%
\pgfpathlineto{\pgfqpoint{1.753545in}{1.269023in}}%
\pgfpathlineto{\pgfqpoint{1.766710in}{1.282093in}}%
\pgfpathlineto{\pgfqpoint{1.752095in}{1.297702in}}%
\pgfpathlineto{\pgfqpoint{1.734827in}{1.312811in}}%
\pgfpathlineto{\pgfqpoint{1.722412in}{1.299285in}}%
\pgfpathlineto{\pgfqpoint{1.709998in}{1.285760in}}%
\pgfpathlineto{\pgfqpoint{1.726472in}{1.271100in}}%
\pgfpathlineto{\pgfqpoint{1.740385in}{1.255959in}}%
\pgfpathclose%
\pgfusepath{fill}%
\end{pgfscope}%
\begin{pgfscope}%
\pgfpathrectangle{\pgfqpoint{0.407592in}{0.222904in}}{\pgfqpoint{1.411163in}{1.411163in}}%
\pgfusepath{clip}%
\pgfsetbuttcap%
\pgfsetroundjoin%
\definecolor{currentfill}{rgb}{0.749020,0.772549,0.792157}%
\pgfsetfillcolor{currentfill}%
\pgfsetfillopacity{0.160000}%
\pgfsetlinewidth{0.000000pt}%
\definecolor{currentstroke}{rgb}{0.000000,0.000000,0.000000}%
\pgfsetstrokecolor{currentstroke}%
\pgfsetdash{}{0pt}%
\pgfpathmoveto{\pgfqpoint{1.083135in}{0.621895in}}%
\pgfpathlineto{\pgfqpoint{1.079564in}{0.627614in}}%
\pgfpathlineto{\pgfqpoint{1.075977in}{0.633359in}}%
\pgfpathlineto{\pgfqpoint{1.079393in}{0.633066in}}%
\pgfpathlineto{\pgfqpoint{1.082856in}{0.632874in}}%
\pgfpathlineto{\pgfqpoint{1.084712in}{0.627251in}}%
\pgfpathlineto{\pgfqpoint{1.086560in}{0.621653in}}%
\pgfpathlineto{\pgfqpoint{1.084836in}{0.621749in}}%
\pgfpathlineto{\pgfqpoint{1.083135in}{0.621895in}}%
\pgfpathclose%
\pgfusepath{fill}%
\end{pgfscope}%
\begin{pgfscope}%
\pgfpathrectangle{\pgfqpoint{0.407592in}{0.222904in}}{\pgfqpoint{1.411163in}{1.411163in}}%
\pgfusepath{clip}%
\pgfsetbuttcap%
\pgfsetroundjoin%
\definecolor{currentfill}{rgb}{0.749020,0.772549,0.792157}%
\pgfsetfillcolor{currentfill}%
\pgfsetfillopacity{0.160000}%
\pgfsetlinewidth{0.000000pt}%
\definecolor{currentstroke}{rgb}{0.000000,0.000000,0.000000}%
\pgfsetstrokecolor{currentstroke}%
\pgfsetdash{}{0pt}%
\pgfpathmoveto{\pgfqpoint{1.090048in}{0.621612in}}%
\pgfpathlineto{\pgfqpoint{1.089956in}{0.627190in}}%
\pgfpathlineto{\pgfqpoint{1.089864in}{0.632793in}}%
\pgfpathlineto{\pgfqpoint{1.093378in}{0.632904in}}%
\pgfpathlineto{\pgfqpoint{1.096876in}{0.633117in}}%
\pgfpathlineto{\pgfqpoint{1.095203in}{0.627433in}}%
\pgfpathlineto{\pgfqpoint{1.093538in}{0.621774in}}%
\pgfpathlineto{\pgfqpoint{1.091796in}{0.621668in}}%
\pgfpathlineto{\pgfqpoint{1.090048in}{0.621612in}}%
\pgfpathclose%
\pgfusepath{fill}%
\end{pgfscope}%
\begin{pgfscope}%
\pgfpathrectangle{\pgfqpoint{0.407592in}{0.222904in}}{\pgfqpoint{1.411163in}{1.411163in}}%
\pgfusepath{clip}%
\pgfsetbuttcap%
\pgfsetroundjoin%
\definecolor{currentfill}{rgb}{0.749020,0.772549,0.792157}%
\pgfsetfillcolor{currentfill}%
\pgfsetfillopacity{0.160000}%
\pgfsetlinewidth{0.000000pt}%
\definecolor{currentstroke}{rgb}{0.000000,0.000000,0.000000}%
\pgfsetstrokecolor{currentstroke}%
\pgfsetdash{}{0pt}%
\pgfpathmoveto{\pgfqpoint{1.086560in}{0.621653in}}%
\pgfpathlineto{\pgfqpoint{1.084712in}{0.627251in}}%
\pgfpathlineto{\pgfqpoint{1.082856in}{0.632874in}}%
\pgfpathlineto{\pgfqpoint{1.086352in}{0.632783in}}%
\pgfpathlineto{\pgfqpoint{1.089864in}{0.632793in}}%
\pgfpathlineto{\pgfqpoint{1.089956in}{0.627190in}}%
\pgfpathlineto{\pgfqpoint{1.090048in}{0.621612in}}%
\pgfpathlineto{\pgfqpoint{1.088300in}{0.621607in}}%
\pgfpathlineto{\pgfqpoint{1.086560in}{0.621653in}}%
\pgfpathclose%
\pgfusepath{fill}%
\end{pgfscope}%
\begin{pgfscope}%
\pgfpathrectangle{\pgfqpoint{0.407592in}{0.222904in}}{\pgfqpoint{1.411163in}{1.411163in}}%
\pgfusepath{clip}%
\pgfsetbuttcap%
\pgfsetroundjoin%
\definecolor{currentfill}{rgb}{0.749020,0.772549,0.792157}%
\pgfsetfillcolor{currentfill}%
\pgfsetfillopacity{0.160000}%
\pgfsetlinewidth{0.000000pt}%
\definecolor{currentstroke}{rgb}{0.000000,0.000000,0.000000}%
\pgfsetstrokecolor{currentstroke}%
\pgfsetdash{}{0pt}%
\pgfpathmoveto{\pgfqpoint{0.861293in}{0.839888in}}%
\pgfpathlineto{\pgfqpoint{0.848513in}{0.852692in}}%
\pgfpathlineto{\pgfqpoint{0.835726in}{0.865502in}}%
\pgfpathlineto{\pgfqpoint{0.830800in}{0.858495in}}%
\pgfpathlineto{\pgfqpoint{0.826982in}{0.851374in}}%
\pgfpathlineto{\pgfqpoint{0.840262in}{0.839225in}}%
\pgfpathlineto{\pgfqpoint{0.853529in}{0.827087in}}%
\pgfpathlineto{\pgfqpoint{0.856912in}{0.833537in}}%
\pgfpathlineto{\pgfqpoint{0.861293in}{0.839888in}}%
\pgfpathclose%
\pgfusepath{fill}%
\end{pgfscope}%
\begin{pgfscope}%
\pgfpathrectangle{\pgfqpoint{0.407592in}{0.222904in}}{\pgfqpoint{1.411163in}{1.411163in}}%
\pgfusepath{clip}%
\pgfsetbuttcap%
\pgfsetroundjoin%
\definecolor{currentfill}{rgb}{0.749020,0.772549,0.792157}%
\pgfsetfillcolor{currentfill}%
\pgfsetfillopacity{0.160000}%
\pgfsetlinewidth{0.000000pt}%
\definecolor{currentstroke}{rgb}{0.000000,0.000000,0.000000}%
\pgfsetstrokecolor{currentstroke}%
\pgfsetdash{}{0pt}%
\pgfpathmoveto{\pgfqpoint{1.165580in}{0.669297in}}%
\pgfpathlineto{\pgfqpoint{1.178234in}{0.679166in}}%
\pgfpathlineto{\pgfqpoint{1.190916in}{0.689056in}}%
\pgfpathlineto{\pgfqpoint{1.192812in}{0.691988in}}%
\pgfpathlineto{\pgfqpoint{1.194254in}{0.694956in}}%
\pgfpathlineto{\pgfqpoint{1.181171in}{0.684338in}}%
\pgfpathlineto{\pgfqpoint{1.168111in}{0.673739in}}%
\pgfpathlineto{\pgfqpoint{1.167015in}{0.671504in}}%
\pgfpathlineto{\pgfqpoint{1.165580in}{0.669297in}}%
\pgfpathclose%
\pgfusepath{fill}%
\end{pgfscope}%
\begin{pgfscope}%
\pgfpathrectangle{\pgfqpoint{0.407592in}{0.222904in}}{\pgfqpoint{1.411163in}{1.411163in}}%
\pgfusepath{clip}%
\pgfsetbuttcap%
\pgfsetroundjoin%
\definecolor{currentfill}{rgb}{0.749020,0.772549,0.792157}%
\pgfsetfillcolor{currentfill}%
\pgfsetfillopacity{0.160000}%
\pgfsetlinewidth{0.000000pt}%
\definecolor{currentstroke}{rgb}{0.000000,0.000000,0.000000}%
\pgfsetstrokecolor{currentstroke}%
\pgfsetdash{}{0pt}%
\pgfpathmoveto{\pgfqpoint{0.558409in}{1.202004in}}%
\pgfpathlineto{\pgfqpoint{0.546296in}{1.215477in}}%
\pgfpathlineto{\pgfqpoint{0.534180in}{1.228950in}}%
\pgfpathlineto{\pgfqpoint{0.516428in}{1.215127in}}%
\pgfpathlineto{\pgfqpoint{0.500965in}{1.200875in}}%
\pgfpathlineto{\pgfqpoint{0.513923in}{1.187894in}}%
\pgfpathlineto{\pgfqpoint{0.526874in}{1.174919in}}%
\pgfpathlineto{\pgfqpoint{0.541542in}{1.188665in}}%
\pgfpathlineto{\pgfqpoint{0.558409in}{1.202004in}}%
\pgfpathclose%
\pgfusepath{fill}%
\end{pgfscope}%
\begin{pgfscope}%
\pgfpathrectangle{\pgfqpoint{0.407592in}{0.222904in}}{\pgfqpoint{1.411163in}{1.411163in}}%
\pgfusepath{clip}%
\pgfsetbuttcap%
\pgfsetroundjoin%
\definecolor{currentfill}{rgb}{0.749020,0.772549,0.792157}%
\pgfsetfillcolor{currentfill}%
\pgfsetfillopacity{0.160000}%
\pgfsetlinewidth{0.000000pt}%
\definecolor{currentstroke}{rgb}{0.000000,0.000000,0.000000}%
\pgfsetstrokecolor{currentstroke}%
\pgfsetdash{}{0pt}%
\pgfpathmoveto{\pgfqpoint{1.194254in}{0.694956in}}%
\pgfpathlineto{\pgfqpoint{1.207360in}{0.705592in}}%
\pgfpathlineto{\pgfqpoint{1.220488in}{0.716247in}}%
\pgfpathlineto{\pgfqpoint{1.221691in}{0.719968in}}%
\pgfpathlineto{\pgfqpoint{1.222316in}{0.723701in}}%
\pgfpathlineto{\pgfqpoint{1.209029in}{0.712317in}}%
\pgfpathlineto{\pgfqpoint{1.195759in}{0.700948in}}%
\pgfpathlineto{\pgfqpoint{1.195237in}{0.697947in}}%
\pgfpathlineto{\pgfqpoint{1.194254in}{0.694956in}}%
\pgfpathclose%
\pgfusepath{fill}%
\end{pgfscope}%
\begin{pgfscope}%
\pgfpathrectangle{\pgfqpoint{0.407592in}{0.222904in}}{\pgfqpoint{1.411163in}{1.411163in}}%
\pgfusepath{clip}%
\pgfsetbuttcap%
\pgfsetroundjoin%
\definecolor{currentfill}{rgb}{0.749020,0.772549,0.792157}%
\pgfsetfillcolor{currentfill}%
\pgfsetfillopacity{0.160000}%
\pgfsetlinewidth{0.000000pt}%
\definecolor{currentstroke}{rgb}{0.000000,0.000000,0.000000}%
\pgfsetstrokecolor{currentstroke}%
\pgfsetdash{}{0pt}%
\pgfpathmoveto{\pgfqpoint{1.134338in}{0.644042in}}%
\pgfpathlineto{\pgfqpoint{1.145447in}{0.652482in}}%
\pgfpathlineto{\pgfqpoint{1.156589in}{0.660946in}}%
\pgfpathlineto{\pgfqpoint{1.159307in}{0.662944in}}%
\pgfpathlineto{\pgfqpoint{1.161718in}{0.665008in}}%
\pgfpathlineto{\pgfqpoint{1.149727in}{0.655870in}}%
\pgfpathlineto{\pgfqpoint{1.137766in}{0.646756in}}%
\pgfpathlineto{\pgfqpoint{1.136154in}{0.645377in}}%
\pgfpathlineto{\pgfqpoint{1.134338in}{0.644042in}}%
\pgfpathclose%
\pgfusepath{fill}%
\end{pgfscope}%
\begin{pgfscope}%
\pgfpathrectangle{\pgfqpoint{0.407592in}{0.222904in}}{\pgfqpoint{1.411163in}{1.411163in}}%
\pgfusepath{clip}%
\pgfsetbuttcap%
\pgfsetroundjoin%
\definecolor{currentfill}{rgb}{0.749020,0.772549,0.792157}%
\pgfsetfillcolor{currentfill}%
\pgfsetfillopacity{0.160000}%
\pgfsetlinewidth{0.000000pt}%
\definecolor{currentstroke}{rgb}{0.000000,0.000000,0.000000}%
\pgfsetstrokecolor{currentstroke}%
\pgfsetdash{}{0pt}%
\pgfpathmoveto{\pgfqpoint{1.766710in}{1.282093in}}%
\pgfpathlineto{\pgfqpoint{1.779882in}{1.295169in}}%
\pgfpathlineto{\pgfqpoint{1.793060in}{1.308251in}}%
\pgfpathlineto{\pgfqpoint{1.777732in}{1.324319in}}%
\pgfpathlineto{\pgfqpoint{1.759659in}{1.339865in}}%
\pgfpathlineto{\pgfqpoint{1.747242in}{1.326337in}}%
\pgfpathlineto{\pgfqpoint{1.734827in}{1.312811in}}%
\pgfpathlineto{\pgfqpoint{1.752095in}{1.297702in}}%
\pgfpathlineto{\pgfqpoint{1.766710in}{1.282093in}}%
\pgfpathclose%
\pgfusepath{fill}%
\end{pgfscope}%
\begin{pgfscope}%
\pgfpathrectangle{\pgfqpoint{0.407592in}{0.222904in}}{\pgfqpoint{1.411163in}{1.411163in}}%
\pgfusepath{clip}%
\pgfsetbuttcap%
\pgfsetroundjoin%
\definecolor{currentfill}{rgb}{0.749020,0.772549,0.792157}%
\pgfsetfillcolor{currentfill}%
\pgfsetfillopacity{0.160000}%
\pgfsetlinewidth{0.000000pt}%
\definecolor{currentstroke}{rgb}{0.000000,0.000000,0.000000}%
\pgfsetstrokecolor{currentstroke}%
\pgfsetdash{}{0pt}%
\pgfpathmoveto{\pgfqpoint{1.011121in}{0.673289in}}%
\pgfpathlineto{\pgfqpoint{0.997853in}{0.683814in}}%
\pgfpathlineto{\pgfqpoint{0.984562in}{0.694358in}}%
\pgfpathlineto{\pgfqpoint{0.985479in}{0.691396in}}%
\pgfpathlineto{\pgfqpoint{0.986858in}{0.688474in}}%
\pgfpathlineto{\pgfqpoint{0.999880in}{0.678655in}}%
\pgfpathlineto{\pgfqpoint{1.012872in}{0.668859in}}%
\pgfpathlineto{\pgfqpoint{1.011825in}{0.671059in}}%
\pgfpathlineto{\pgfqpoint{1.011121in}{0.673289in}}%
\pgfpathclose%
\pgfusepath{fill}%
\end{pgfscope}%
\begin{pgfscope}%
\pgfpathrectangle{\pgfqpoint{0.407592in}{0.222904in}}{\pgfqpoint{1.411163in}{1.411163in}}%
\pgfusepath{clip}%
\pgfsetbuttcap%
\pgfsetroundjoin%
\definecolor{currentfill}{rgb}{0.749020,0.772549,0.792157}%
\pgfsetfillcolor{currentfill}%
\pgfsetfillopacity{0.160000}%
\pgfsetlinewidth{0.000000pt}%
\definecolor{currentstroke}{rgb}{0.000000,0.000000,0.000000}%
\pgfsetstrokecolor{currentstroke}%
\pgfsetdash{}{0pt}%
\pgfpathmoveto{\pgfqpoint{1.040861in}{0.646476in}}%
\pgfpathlineto{\pgfqpoint{1.028438in}{0.655520in}}%
\pgfpathlineto{\pgfqpoint{1.015983in}{0.664589in}}%
\pgfpathlineto{\pgfqpoint{1.018033in}{0.662537in}}%
\pgfpathlineto{\pgfqpoint{1.020402in}{0.660554in}}%
\pgfpathlineto{\pgfqpoint{1.032127in}{0.652155in}}%
\pgfpathlineto{\pgfqpoint{1.043817in}{0.643781in}}%
\pgfpathlineto{\pgfqpoint{1.042234in}{0.645105in}}%
\pgfpathlineto{\pgfqpoint{1.040861in}{0.646476in}}%
\pgfpathclose%
\pgfusepath{fill}%
\end{pgfscope}%
\begin{pgfscope}%
\pgfpathrectangle{\pgfqpoint{0.407592in}{0.222904in}}{\pgfqpoint{1.411163in}{1.411163in}}%
\pgfusepath{clip}%
\pgfsetbuttcap%
\pgfsetroundjoin%
\definecolor{currentfill}{rgb}{0.749020,0.772549,0.792157}%
\pgfsetfillcolor{currentfill}%
\pgfsetfillopacity{0.160000}%
\pgfsetlinewidth{0.000000pt}%
\definecolor{currentstroke}{rgb}{0.000000,0.000000,0.000000}%
\pgfsetstrokecolor{currentstroke}%
\pgfsetdash{}{0pt}%
\pgfpathmoveto{\pgfqpoint{0.984116in}{0.700345in}}%
\pgfpathlineto{\pgfqpoint{0.970770in}{0.711641in}}%
\pgfpathlineto{\pgfqpoint{0.957406in}{0.722952in}}%
\pgfpathlineto{\pgfqpoint{0.957366in}{0.719220in}}%
\pgfpathlineto{\pgfqpoint{0.957907in}{0.715503in}}%
\pgfpathlineto{\pgfqpoint{0.971246in}{0.704921in}}%
\pgfpathlineto{\pgfqpoint{0.984562in}{0.694358in}}%
\pgfpathlineto{\pgfqpoint{0.984107in}{0.697345in}}%
\pgfpathlineto{\pgfqpoint{0.984116in}{0.700345in}}%
\pgfpathclose%
\pgfusepath{fill}%
\end{pgfscope}%
\begin{pgfscope}%
\pgfpathrectangle{\pgfqpoint{0.407592in}{0.222904in}}{\pgfqpoint{1.411163in}{1.411163in}}%
\pgfusepath{clip}%
\pgfsetbuttcap%
\pgfsetroundjoin%
\definecolor{currentfill}{rgb}{0.749020,0.772549,0.792157}%
\pgfsetfillcolor{currentfill}%
\pgfsetfillopacity{0.160000}%
\pgfsetlinewidth{0.000000pt}%
\definecolor{currentstroke}{rgb}{0.000000,0.000000,0.000000}%
\pgfsetstrokecolor{currentstroke}%
\pgfsetdash{}{0pt}%
\pgfpathmoveto{\pgfqpoint{1.354602in}{0.852811in}}%
\pgfpathlineto{\pgfqpoint{1.367945in}{0.865038in}}%
\pgfpathlineto{\pgfqpoint{1.381300in}{0.877277in}}%
\pgfpathlineto{\pgfqpoint{1.378481in}{0.885035in}}%
\pgfpathlineto{\pgfqpoint{1.374423in}{0.892660in}}%
\pgfpathlineto{\pgfqpoint{1.361442in}{0.879773in}}%
\pgfpathlineto{\pgfqpoint{1.348467in}{0.866892in}}%
\pgfpathlineto{\pgfqpoint{1.352098in}{0.859910in}}%
\pgfpathlineto{\pgfqpoint{1.354602in}{0.852811in}}%
\pgfpathclose%
\pgfusepath{fill}%
\end{pgfscope}%
\begin{pgfscope}%
\pgfpathrectangle{\pgfqpoint{0.407592in}{0.222904in}}{\pgfqpoint{1.411163in}{1.411163in}}%
\pgfusepath{clip}%
\pgfsetbuttcap%
\pgfsetroundjoin%
\definecolor{currentfill}{rgb}{0.749020,0.772549,0.792157}%
\pgfsetfillcolor{currentfill}%
\pgfsetfillopacity{0.160000}%
\pgfsetlinewidth{0.000000pt}%
\definecolor{currentstroke}{rgb}{0.000000,0.000000,0.000000}%
\pgfsetstrokecolor{currentstroke}%
\pgfsetdash{}{0pt}%
\pgfpathmoveto{\pgfqpoint{0.534180in}{1.228950in}}%
\pgfpathlineto{\pgfqpoint{0.522064in}{1.242426in}}%
\pgfpathlineto{\pgfqpoint{0.509947in}{1.255902in}}%
\pgfpathlineto{\pgfqpoint{0.491298in}{1.241605in}}%
\pgfpathlineto{\pgfqpoint{0.475029in}{1.226859in}}%
\pgfpathlineto{\pgfqpoint{0.488000in}{1.213864in}}%
\pgfpathlineto{\pgfqpoint{0.500965in}{1.200875in}}%
\pgfpathlineto{\pgfqpoint{0.516428in}{1.215127in}}%
\pgfpathlineto{\pgfqpoint{0.534180in}{1.228950in}}%
\pgfpathclose%
\pgfusepath{fill}%
\end{pgfscope}%
\begin{pgfscope}%
\pgfpathrectangle{\pgfqpoint{0.407592in}{0.222904in}}{\pgfqpoint{1.411163in}{1.411163in}}%
\pgfusepath{clip}%
\pgfsetbuttcap%
\pgfsetroundjoin%
\definecolor{currentfill}{rgb}{0.749020,0.772549,0.792157}%
\pgfsetfillcolor{currentfill}%
\pgfsetfillopacity{0.160000}%
\pgfsetlinewidth{0.000000pt}%
\definecolor{currentstroke}{rgb}{0.000000,0.000000,0.000000}%
\pgfsetstrokecolor{currentstroke}%
\pgfsetdash{}{0pt}%
\pgfpathmoveto{\pgfqpoint{1.248943in}{0.746514in}}%
\pgfpathlineto{\pgfqpoint{1.262283in}{0.757944in}}%
\pgfpathlineto{\pgfqpoint{1.275641in}{0.769388in}}%
\pgfpathlineto{\pgfqpoint{1.275623in}{0.774548in}}%
\pgfpathlineto{\pgfqpoint{1.274795in}{0.779674in}}%
\pgfpathlineto{\pgfqpoint{1.261536in}{0.767522in}}%
\pgfpathlineto{\pgfqpoint{1.248288in}{0.755381in}}%
\pgfpathlineto{\pgfqpoint{1.248962in}{0.750962in}}%
\pgfpathlineto{\pgfqpoint{1.248943in}{0.746514in}}%
\pgfpathclose%
\pgfusepath{fill}%
\end{pgfscope}%
\begin{pgfscope}%
\pgfpathrectangle{\pgfqpoint{0.407592in}{0.222904in}}{\pgfqpoint{1.411163in}{1.411163in}}%
\pgfusepath{clip}%
\pgfsetbuttcap%
\pgfsetroundjoin%
\definecolor{currentfill}{rgb}{0.749020,0.772549,0.792157}%
\pgfsetfillcolor{currentfill}%
\pgfsetfillopacity{0.160000}%
\pgfsetlinewidth{0.000000pt}%
\definecolor{currentstroke}{rgb}{0.000000,0.000000,0.000000}%
\pgfsetstrokecolor{currentstroke}%
\pgfsetdash{}{0pt}%
\pgfpathmoveto{\pgfqpoint{1.793060in}{1.308251in}}%
\pgfpathlineto{\pgfqpoint{1.806244in}{1.321339in}}%
\pgfpathlineto{\pgfqpoint{1.819435in}{1.334434in}}%
\pgfpathlineto{\pgfqpoint{1.803382in}{1.350950in}}%
\pgfpathlineto{\pgfqpoint{1.784494in}{1.366922in}}%
\pgfpathlineto{\pgfqpoint{1.772076in}{1.353393in}}%
\pgfpathlineto{\pgfqpoint{1.759659in}{1.339865in}}%
\pgfpathlineto{\pgfqpoint{1.777732in}{1.324319in}}%
\pgfpathlineto{\pgfqpoint{1.793060in}{1.308251in}}%
\pgfpathclose%
\pgfusepath{fill}%
\end{pgfscope}%
\begin{pgfscope}%
\pgfpathrectangle{\pgfqpoint{0.407592in}{0.222904in}}{\pgfqpoint{1.411163in}{1.411163in}}%
\pgfusepath{clip}%
\pgfsetbuttcap%
\pgfsetroundjoin%
\definecolor{currentfill}{rgb}{0.749020,0.772549,0.792157}%
\pgfsetfillcolor{currentfill}%
\pgfsetfillopacity{0.160000}%
\pgfsetlinewidth{0.000000pt}%
\definecolor{currentstroke}{rgb}{0.000000,0.000000,0.000000}%
\pgfsetstrokecolor{currentstroke}%
\pgfsetdash{}{0pt}%
\pgfpathmoveto{\pgfqpoint{1.819435in}{1.334434in}}%
\pgfpathlineto{\pgfqpoint{1.832631in}{1.347535in}}%
\pgfpathlineto{\pgfqpoint{1.816212in}{1.364270in}}%
\pgfpathlineto{\pgfqpoint{1.796912in}{1.380452in}}%
\pgfpathlineto{\pgfqpoint{1.784494in}{1.366922in}}%
\pgfpathlineto{\pgfqpoint{1.803382in}{1.350950in}}%
\pgfusepath{fill}%
\end{pgfscope}%
\begin{pgfscope}%
\pgfpathrectangle{\pgfqpoint{0.407592in}{0.222904in}}{\pgfqpoint{1.411163in}{1.411163in}}%
\pgfusepath{clip}%
\pgfsetbuttcap%
\pgfsetroundjoin%
\definecolor{currentfill}{rgb}{0.749020,0.772549,0.792157}%
\pgfsetfillcolor{currentfill}%
\pgfsetfillopacity{0.160000}%
\pgfsetlinewidth{0.000000pt}%
\definecolor{currentstroke}{rgb}{0.000000,0.000000,0.000000}%
\pgfsetstrokecolor{currentstroke}%
\pgfsetdash{}{0pt}%
\pgfpathmoveto{\pgfqpoint{0.932871in}{0.754498in}}%
\pgfpathlineto{\pgfqpoint{0.919678in}{0.766568in}}%
\pgfpathlineto{\pgfqpoint{0.906473in}{0.778649in}}%
\pgfpathlineto{\pgfqpoint{0.904718in}{0.773515in}}%
\pgfpathlineto{\pgfqpoint{0.903768in}{0.768351in}}%
\pgfpathlineto{\pgfqpoint{0.917205in}{0.756978in}}%
\pgfpathlineto{\pgfqpoint{0.930624in}{0.745620in}}%
\pgfpathlineto{\pgfqpoint{0.931403in}{0.750071in}}%
\pgfpathlineto{\pgfqpoint{0.932871in}{0.754498in}}%
\pgfpathclose%
\pgfusepath{fill}%
\end{pgfscope}%
\begin{pgfscope}%
\pgfpathrectangle{\pgfqpoint{0.407592in}{0.222904in}}{\pgfqpoint{1.411163in}{1.411163in}}%
\pgfusepath{clip}%
\pgfsetbuttcap%
\pgfsetroundjoin%
\definecolor{currentfill}{rgb}{0.749020,0.772549,0.792157}%
\pgfsetfillcolor{currentfill}%
\pgfsetfillopacity{0.160000}%
\pgfsetlinewidth{0.000000pt}%
\definecolor{currentstroke}{rgb}{0.000000,0.000000,0.000000}%
\pgfsetstrokecolor{currentstroke}%
\pgfsetdash{}{0pt}%
\pgfpathmoveto{\pgfqpoint{0.509947in}{1.255902in}}%
\pgfpathlineto{\pgfqpoint{0.497828in}{1.269380in}}%
\pgfpathlineto{\pgfqpoint{0.485708in}{1.282860in}}%
\pgfpathlineto{\pgfqpoint{0.466153in}{1.268099in}}%
\pgfpathlineto{\pgfqpoint{0.449066in}{1.252869in}}%
\pgfpathlineto{\pgfqpoint{0.462051in}{1.239861in}}%
\pgfpathlineto{\pgfqpoint{0.475029in}{1.226859in}}%
\pgfpathlineto{\pgfqpoint{0.491298in}{1.241605in}}%
\pgfpathlineto{\pgfqpoint{0.509947in}{1.255902in}}%
\pgfpathclose%
\pgfusepath{fill}%
\end{pgfscope}%
\begin{pgfscope}%
\pgfpathrectangle{\pgfqpoint{0.407592in}{0.222904in}}{\pgfqpoint{1.411163in}{1.411163in}}%
\pgfusepath{clip}%
\pgfsetbuttcap%
\pgfsetroundjoin%
\definecolor{currentfill}{rgb}{0.749020,0.772549,0.792157}%
\pgfsetfillcolor{currentfill}%
\pgfsetfillopacity{0.160000}%
\pgfsetlinewidth{0.000000pt}%
\definecolor{currentstroke}{rgb}{0.000000,0.000000,0.000000}%
\pgfsetstrokecolor{currentstroke}%
\pgfsetdash{}{0pt}%
\pgfpathmoveto{\pgfqpoint{0.835726in}{0.865502in}}%
\pgfpathlineto{\pgfqpoint{0.822932in}{0.878319in}}%
\pgfpathlineto{\pgfqpoint{0.810133in}{0.891142in}}%
\pgfpathlineto{\pgfqpoint{0.804651in}{0.883488in}}%
\pgfpathlineto{\pgfqpoint{0.800386in}{0.875707in}}%
\pgfpathlineto{\pgfqpoint{0.813690in}{0.863535in}}%
\pgfpathlineto{\pgfqpoint{0.826982in}{0.851374in}}%
\pgfpathlineto{\pgfqpoint{0.830800in}{0.858495in}}%
\pgfpathlineto{\pgfqpoint{0.835726in}{0.865502in}}%
\pgfpathclose%
\pgfusepath{fill}%
\end{pgfscope}%
\begin{pgfscope}%
\pgfpathrectangle{\pgfqpoint{0.407592in}{0.222904in}}{\pgfqpoint{1.411163in}{1.411163in}}%
\pgfusepath{clip}%
\pgfsetbuttcap%
\pgfsetroundjoin%
\definecolor{currentfill}{rgb}{0.749020,0.772549,0.792157}%
\pgfsetfillcolor{currentfill}%
\pgfsetfillopacity{0.160000}%
\pgfsetlinewidth{0.000000pt}%
\definecolor{currentstroke}{rgb}{0.000000,0.000000,0.000000}%
\pgfsetstrokecolor{currentstroke}%
\pgfsetdash{}{0pt}%
\pgfpathmoveto{\pgfqpoint{1.132327in}{0.642759in}}%
\pgfpathlineto{\pgfqpoint{1.142935in}{0.650879in}}%
\pgfpathlineto{\pgfqpoint{1.153575in}{0.659024in}}%
\pgfpathlineto{\pgfqpoint{1.156589in}{0.660946in}}%
\pgfpathlineto{\pgfqpoint{1.145447in}{0.652482in}}%
\pgfpathlineto{\pgfqpoint{1.134338in}{0.644042in}}%
\pgfusepath{fill}%
\end{pgfscope}%
\begin{pgfscope}%
\pgfpathrectangle{\pgfqpoint{0.407592in}{0.222904in}}{\pgfqpoint{1.411163in}{1.411163in}}%
\pgfusepath{clip}%
\pgfsetbuttcap%
\pgfsetroundjoin%
\definecolor{currentfill}{rgb}{0.749020,0.772549,0.792157}%
\pgfsetfillcolor{currentfill}%
\pgfsetfillopacity{0.160000}%
\pgfsetlinewidth{0.000000pt}%
\definecolor{currentstroke}{rgb}{0.000000,0.000000,0.000000}%
\pgfsetstrokecolor{currentstroke}%
\pgfsetdash{}{0pt}%
\pgfpathmoveto{\pgfqpoint{0.485708in}{1.282860in}}%
\pgfpathlineto{\pgfqpoint{0.473587in}{1.296340in}}%
\pgfpathlineto{\pgfqpoint{0.461465in}{1.309822in}}%
\pgfpathlineto{\pgfqpoint{0.440993in}{1.294609in}}%
\pgfpathlineto{\pgfqpoint{0.423077in}{1.278906in}}%
\pgfpathlineto{\pgfqpoint{0.436075in}{1.265884in}}%
\pgfpathlineto{\pgfqpoint{0.449066in}{1.252869in}}%
\pgfpathlineto{\pgfqpoint{0.466153in}{1.268099in}}%
\pgfpathlineto{\pgfqpoint{0.485708in}{1.282860in}}%
\pgfpathclose%
\pgfusepath{fill}%
\end{pgfscope}%
\begin{pgfscope}%
\pgfpathrectangle{\pgfqpoint{0.407592in}{0.222904in}}{\pgfqpoint{1.411163in}{1.411163in}}%
\pgfusepath{clip}%
\pgfsetbuttcap%
\pgfsetroundjoin%
\definecolor{currentfill}{rgb}{0.749020,0.772549,0.792157}%
\pgfsetfillcolor{currentfill}%
\pgfsetfillopacity{0.160000}%
\pgfsetlinewidth{0.000000pt}%
\definecolor{currentstroke}{rgb}{0.000000,0.000000,0.000000}%
\pgfsetstrokecolor{currentstroke}%
\pgfsetdash{}{0pt}%
\pgfpathmoveto{\pgfqpoint{1.381300in}{0.877277in}}%
\pgfpathlineto{\pgfqpoint{1.394666in}{0.889526in}}%
\pgfpathlineto{\pgfqpoint{1.408045in}{0.901787in}}%
\pgfpathlineto{\pgfqpoint{1.404899in}{0.910193in}}%
\pgfpathlineto{\pgfqpoint{1.400404in}{0.918452in}}%
\pgfpathlineto{\pgfqpoint{1.387411in}{0.905553in}}%
\pgfpathlineto{\pgfqpoint{1.374423in}{0.892660in}}%
\pgfpathlineto{\pgfqpoint{1.378481in}{0.885035in}}%
\pgfpathlineto{\pgfqpoint{1.381300in}{0.877277in}}%
\pgfpathclose%
\pgfusepath{fill}%
\end{pgfscope}%
\begin{pgfscope}%
\pgfpathrectangle{\pgfqpoint{0.407592in}{0.222904in}}{\pgfqpoint{1.411163in}{1.411163in}}%
\pgfusepath{clip}%
\pgfsetbuttcap%
\pgfsetroundjoin%
\definecolor{currentfill}{rgb}{0.749020,0.772549,0.792157}%
\pgfsetfillcolor{currentfill}%
\pgfsetfillopacity{0.160000}%
\pgfsetlinewidth{0.000000pt}%
\definecolor{currentstroke}{rgb}{0.000000,0.000000,0.000000}%
\pgfsetstrokecolor{currentstroke}%
\pgfsetdash{}{0pt}%
\pgfpathmoveto{\pgfqpoint{1.043817in}{0.643781in}}%
\pgfpathlineto{\pgfqpoint{1.032127in}{0.652155in}}%
\pgfpathlineto{\pgfqpoint{1.020402in}{0.660554in}}%
\pgfpathlineto{\pgfqpoint{1.023079in}{0.658648in}}%
\pgfpathlineto{\pgfqpoint{1.026055in}{0.656828in}}%
\pgfpathlineto{\pgfqpoint{1.036840in}{0.649047in}}%
\pgfpathlineto{\pgfqpoint{1.047590in}{0.641293in}}%
\pgfpathlineto{\pgfqpoint{1.045605in}{0.642508in}}%
\pgfpathlineto{\pgfqpoint{1.043817in}{0.643781in}}%
\pgfpathclose%
\pgfusepath{fill}%
\end{pgfscope}%
\begin{pgfscope}%
\pgfpathrectangle{\pgfqpoint{0.407592in}{0.222904in}}{\pgfqpoint{1.411163in}{1.411163in}}%
\pgfusepath{clip}%
\pgfsetbuttcap%
\pgfsetroundjoin%
\definecolor{currentfill}{rgb}{0.749020,0.772549,0.792157}%
\pgfsetfillcolor{currentfill}%
\pgfsetfillopacity{0.160000}%
\pgfsetlinewidth{0.000000pt}%
\definecolor{currentstroke}{rgb}{0.000000,0.000000,0.000000}%
\pgfsetstrokecolor{currentstroke}%
\pgfsetdash{}{0pt}%
\pgfpathmoveto{\pgfqpoint{0.461465in}{1.309822in}}%
\pgfpathlineto{\pgfqpoint{0.449341in}{1.323306in}}%
\pgfpathlineto{\pgfqpoint{0.437216in}{1.336790in}}%
\pgfpathlineto{\pgfqpoint{0.415818in}{1.321136in}}%
\pgfpathlineto{\pgfqpoint{0.397060in}{1.304970in}}%
\pgfpathlineto{\pgfqpoint{0.410072in}{1.291934in}}%
\pgfpathlineto{\pgfqpoint{0.423077in}{1.278906in}}%
\pgfpathlineto{\pgfqpoint{0.440993in}{1.294609in}}%
\pgfpathlineto{\pgfqpoint{0.461465in}{1.309822in}}%
\pgfpathclose%
\pgfusepath{fill}%
\end{pgfscope}%
\begin{pgfscope}%
\pgfpathrectangle{\pgfqpoint{0.407592in}{0.222904in}}{\pgfqpoint{1.411163in}{1.411163in}}%
\pgfusepath{clip}%
\pgfsetbuttcap%
\pgfsetroundjoin%
\definecolor{currentfill}{rgb}{0.749020,0.772549,0.792157}%
\pgfsetfillcolor{currentfill}%
\pgfsetfillopacity{0.160000}%
\pgfsetlinewidth{0.000000pt}%
\definecolor{currentstroke}{rgb}{0.000000,0.000000,0.000000}%
\pgfsetstrokecolor{currentstroke}%
\pgfsetdash{}{0pt}%
\pgfpathmoveto{\pgfqpoint{1.127755in}{0.640367in}}%
\pgfpathlineto{\pgfqpoint{1.137218in}{0.647890in}}%
\pgfpathlineto{\pgfqpoint{1.146714in}{0.655440in}}%
\pgfpathlineto{\pgfqpoint{1.150279in}{0.657186in}}%
\pgfpathlineto{\pgfqpoint{1.153575in}{0.659024in}}%
\pgfpathlineto{\pgfqpoint{1.142935in}{0.650879in}}%
\pgfpathlineto{\pgfqpoint{1.132327in}{0.642759in}}%
\pgfpathlineto{\pgfqpoint{1.130130in}{0.641532in}}%
\pgfpathlineto{\pgfqpoint{1.127755in}{0.640367in}}%
\pgfpathclose%
\pgfusepath{fill}%
\end{pgfscope}%
\begin{pgfscope}%
\pgfpathrectangle{\pgfqpoint{0.407592in}{0.222904in}}{\pgfqpoint{1.411163in}{1.411163in}}%
\pgfusepath{clip}%
\pgfsetbuttcap%
\pgfsetroundjoin%
\definecolor{currentfill}{rgb}{0.749020,0.772549,0.792157}%
\pgfsetfillcolor{currentfill}%
\pgfsetfillopacity{0.160000}%
\pgfsetlinewidth{0.000000pt}%
\definecolor{currentstroke}{rgb}{0.000000,0.000000,0.000000}%
\pgfsetstrokecolor{currentstroke}%
\pgfsetdash{}{0pt}%
\pgfpathmoveto{\pgfqpoint{1.275641in}{0.769388in}}%
\pgfpathlineto{\pgfqpoint{1.289016in}{0.780848in}}%
\pgfpathlineto{\pgfqpoint{1.302409in}{0.792323in}}%
\pgfpathlineto{\pgfqpoint{1.302343in}{0.798187in}}%
\pgfpathlineto{\pgfqpoint{1.301350in}{0.804009in}}%
\pgfpathlineto{\pgfqpoint{1.288067in}{0.791836in}}%
\pgfpathlineto{\pgfqpoint{1.274795in}{0.779674in}}%
\pgfpathlineto{\pgfqpoint{1.275623in}{0.774548in}}%
\pgfpathlineto{\pgfqpoint{1.275641in}{0.769388in}}%
\pgfpathclose%
\pgfusepath{fill}%
\end{pgfscope}%
\begin{pgfscope}%
\pgfpathrectangle{\pgfqpoint{0.407592in}{0.222904in}}{\pgfqpoint{1.411163in}{1.411163in}}%
\pgfusepath{clip}%
\pgfsetbuttcap%
\pgfsetroundjoin%
\definecolor{currentfill}{rgb}{0.749020,0.772549,0.792157}%
\pgfsetfillcolor{currentfill}%
\pgfsetfillopacity{0.160000}%
\pgfsetlinewidth{0.000000pt}%
\definecolor{currentstroke}{rgb}{0.000000,0.000000,0.000000}%
\pgfsetstrokecolor{currentstroke}%
\pgfsetdash{}{0pt}%
\pgfpathmoveto{\pgfqpoint{0.810133in}{0.891142in}}%
\pgfpathlineto{\pgfqpoint{0.797326in}{0.903972in}}%
\pgfpathlineto{\pgfqpoint{0.784513in}{0.916808in}}%
\pgfpathlineto{\pgfqpoint{0.778464in}{0.908517in}}%
\pgfpathlineto{\pgfqpoint{0.773739in}{0.900086in}}%
\pgfpathlineto{\pgfqpoint{0.787069in}{0.887891in}}%
\pgfpathlineto{\pgfqpoint{0.800386in}{0.875707in}}%
\pgfpathlineto{\pgfqpoint{0.804651in}{0.883488in}}%
\pgfpathlineto{\pgfqpoint{0.810133in}{0.891142in}}%
\pgfpathclose%
\pgfusepath{fill}%
\end{pgfscope}%
\begin{pgfscope}%
\pgfpathrectangle{\pgfqpoint{0.407592in}{0.222904in}}{\pgfqpoint{1.411163in}{1.411163in}}%
\pgfusepath{clip}%
\pgfsetbuttcap%
\pgfsetroundjoin%
\definecolor{currentfill}{rgb}{0.749020,0.772549,0.792157}%
\pgfsetfillcolor{currentfill}%
\pgfsetfillopacity{0.160000}%
\pgfsetlinewidth{0.000000pt}%
\definecolor{currentstroke}{rgb}{0.000000,0.000000,0.000000}%
\pgfsetstrokecolor{currentstroke}%
\pgfsetdash{}{0pt}%
\pgfpathmoveto{\pgfqpoint{1.161718in}{0.665008in}}%
\pgfpathlineto{\pgfqpoint{1.173741in}{0.674169in}}%
\pgfpathlineto{\pgfqpoint{1.185795in}{0.683354in}}%
\pgfpathlineto{\pgfqpoint{1.188574in}{0.686174in}}%
\pgfpathlineto{\pgfqpoint{1.190916in}{0.689056in}}%
\pgfpathlineto{\pgfqpoint{1.178234in}{0.679166in}}%
\pgfpathlineto{\pgfqpoint{1.165580in}{0.669297in}}%
\pgfpathlineto{\pgfqpoint{1.163812in}{0.667129in}}%
\pgfpathlineto{\pgfqpoint{1.161718in}{0.665008in}}%
\pgfpathclose%
\pgfusepath{fill}%
\end{pgfscope}%
\begin{pgfscope}%
\pgfpathrectangle{\pgfqpoint{0.407592in}{0.222904in}}{\pgfqpoint{1.411163in}{1.411163in}}%
\pgfusepath{clip}%
\pgfsetbuttcap%
\pgfsetroundjoin%
\definecolor{currentfill}{rgb}{0.749020,0.772549,0.792157}%
\pgfsetfillcolor{currentfill}%
\pgfsetfillopacity{0.160000}%
\pgfsetlinewidth{0.000000pt}%
\definecolor{currentstroke}{rgb}{0.000000,0.000000,0.000000}%
\pgfsetstrokecolor{currentstroke}%
\pgfsetdash{}{0pt}%
\pgfpathmoveto{\pgfqpoint{0.437216in}{1.336790in}}%
\pgfpathlineto{\pgfqpoint{0.425090in}{1.350276in}}%
\pgfpathlineto{\pgfqpoint{0.412963in}{1.363764in}}%
\pgfpathlineto{\pgfqpoint{0.390627in}{1.347678in}}%
\pgfpathlineto{\pgfqpoint{0.371018in}{1.331060in}}%
\pgfpathlineto{\pgfqpoint{0.384042in}{1.318011in}}%
\pgfpathlineto{\pgfqpoint{0.397060in}{1.304970in}}%
\pgfpathlineto{\pgfqpoint{0.415818in}{1.321136in}}%
\pgfpathlineto{\pgfqpoint{0.437216in}{1.336790in}}%
\pgfpathclose%
\pgfusepath{fill}%
\end{pgfscope}%
\begin{pgfscope}%
\pgfpathrectangle{\pgfqpoint{0.407592in}{0.222904in}}{\pgfqpoint{1.411163in}{1.411163in}}%
\pgfusepath{clip}%
\pgfsetbuttcap%
\pgfsetroundjoin%
\definecolor{currentfill}{rgb}{0.749020,0.772549,0.792157}%
\pgfsetfillcolor{currentfill}%
\pgfsetfillopacity{0.160000}%
\pgfsetlinewidth{0.000000pt}%
\definecolor{currentstroke}{rgb}{0.000000,0.000000,0.000000}%
\pgfsetstrokecolor{currentstroke}%
\pgfsetdash{}{0pt}%
\pgfpathmoveto{\pgfqpoint{1.220488in}{0.716247in}}%
\pgfpathlineto{\pgfqpoint{1.233640in}{0.726920in}}%
\pgfpathlineto{\pgfqpoint{1.246815in}{0.737613in}}%
\pgfpathlineto{\pgfqpoint{1.248227in}{0.742058in}}%
\pgfpathlineto{\pgfqpoint{1.248943in}{0.746514in}}%
\pgfpathlineto{\pgfqpoint{1.235621in}{0.735100in}}%
\pgfpathlineto{\pgfqpoint{1.222316in}{0.723701in}}%
\pgfpathlineto{\pgfqpoint{1.221691in}{0.719968in}}%
\pgfpathlineto{\pgfqpoint{1.220488in}{0.716247in}}%
\pgfpathclose%
\pgfusepath{fill}%
\end{pgfscope}%
\begin{pgfscope}%
\pgfpathrectangle{\pgfqpoint{0.407592in}{0.222904in}}{\pgfqpoint{1.411163in}{1.411163in}}%
\pgfusepath{clip}%
\pgfsetbuttcap%
\pgfsetroundjoin%
\definecolor{currentfill}{rgb}{0.749020,0.772549,0.792157}%
\pgfsetfillcolor{currentfill}%
\pgfsetfillopacity{0.160000}%
\pgfsetlinewidth{0.000000pt}%
\definecolor{currentstroke}{rgb}{0.000000,0.000000,0.000000}%
\pgfsetstrokecolor{currentstroke}%
\pgfsetdash{}{0pt}%
\pgfpathmoveto{\pgfqpoint{1.012872in}{0.668859in}}%
\pgfpathlineto{\pgfqpoint{0.999880in}{0.678655in}}%
\pgfpathlineto{\pgfqpoint{0.986858in}{0.688474in}}%
\pgfpathlineto{\pgfqpoint{0.988691in}{0.685603in}}%
\pgfpathlineto{\pgfqpoint{0.990974in}{0.682797in}}%
\pgfpathlineto{\pgfqpoint{1.003495in}{0.673681in}}%
\pgfpathlineto{\pgfqpoint{1.015983in}{0.664589in}}%
\pgfpathlineto{\pgfqpoint{1.014260in}{0.666699in}}%
\pgfpathlineto{\pgfqpoint{1.012872in}{0.668859in}}%
\pgfpathclose%
\pgfusepath{fill}%
\end{pgfscope}%
\begin{pgfscope}%
\pgfpathrectangle{\pgfqpoint{0.407592in}{0.222904in}}{\pgfqpoint{1.411163in}{1.411163in}}%
\pgfusepath{clip}%
\pgfsetbuttcap%
\pgfsetroundjoin%
\definecolor{currentfill}{rgb}{0.749020,0.772549,0.792157}%
\pgfsetfillcolor{currentfill}%
\pgfsetfillopacity{0.160000}%
\pgfsetlinewidth{0.000000pt}%
\definecolor{currentstroke}{rgb}{0.000000,0.000000,0.000000}%
\pgfsetstrokecolor{currentstroke}%
\pgfsetdash{}{0pt}%
\pgfpathmoveto{\pgfqpoint{1.408045in}{0.901787in}}%
\pgfpathlineto{\pgfqpoint{1.421436in}{0.914058in}}%
\pgfpathlineto{\pgfqpoint{1.434838in}{0.926341in}}%
\pgfpathlineto{\pgfqpoint{1.431354in}{0.935386in}}%
\pgfpathlineto{\pgfqpoint{1.426410in}{0.944268in}}%
\pgfpathlineto{\pgfqpoint{1.413404in}{0.931357in}}%
\pgfpathlineto{\pgfqpoint{1.400404in}{0.918452in}}%
\pgfpathlineto{\pgfqpoint{1.404899in}{0.910193in}}%
\pgfpathlineto{\pgfqpoint{1.408045in}{0.901787in}}%
\pgfpathclose%
\pgfusepath{fill}%
\end{pgfscope}%
\begin{pgfscope}%
\pgfpathrectangle{\pgfqpoint{0.407592in}{0.222904in}}{\pgfqpoint{1.411163in}{1.411163in}}%
\pgfusepath{clip}%
\pgfsetbuttcap%
\pgfsetroundjoin%
\definecolor{currentfill}{rgb}{0.749020,0.772549,0.792157}%
\pgfsetfillcolor{currentfill}%
\pgfsetfillopacity{0.160000}%
\pgfsetlinewidth{0.000000pt}%
\definecolor{currentstroke}{rgb}{0.000000,0.000000,0.000000}%
\pgfsetstrokecolor{currentstroke}%
\pgfsetdash{}{0pt}%
\pgfpathmoveto{\pgfqpoint{0.412963in}{1.363764in}}%
\pgfpathlineto{\pgfqpoint{0.400835in}{1.377253in}}%
\pgfpathlineto{\pgfqpoint{0.378026in}{1.360955in}}%
\pgfpathlineto{\pgfqpoint{0.357986in}{1.344115in}}%
\pgfpathlineto{\pgfqpoint{0.371018in}{1.331060in}}%
\pgfpathlineto{\pgfqpoint{0.390627in}{1.347678in}}%
\pgfusepath{fill}%
\end{pgfscope}%
\begin{pgfscope}%
\pgfpathrectangle{\pgfqpoint{0.407592in}{0.222904in}}{\pgfqpoint{1.411163in}{1.411163in}}%
\pgfusepath{clip}%
\pgfsetbuttcap%
\pgfsetroundjoin%
\definecolor{currentfill}{rgb}{0.749020,0.772549,0.792157}%
\pgfsetfillcolor{currentfill}%
\pgfsetfillopacity{0.160000}%
\pgfsetlinewidth{0.000000pt}%
\definecolor{currentstroke}{rgb}{0.000000,0.000000,0.000000}%
\pgfsetstrokecolor{currentstroke}%
\pgfsetdash{}{0pt}%
\pgfpathmoveto{\pgfqpoint{0.906473in}{0.778649in}}%
\pgfpathlineto{\pgfqpoint{0.893256in}{0.790742in}}%
\pgfpathlineto{\pgfqpoint{0.880026in}{0.802846in}}%
\pgfpathlineto{\pgfqpoint{0.877972in}{0.797013in}}%
\pgfpathlineto{\pgfqpoint{0.876838in}{0.791143in}}%
\pgfpathlineto{\pgfqpoint{0.890312in}{0.779739in}}%
\pgfpathlineto{\pgfqpoint{0.903768in}{0.768351in}}%
\pgfpathlineto{\pgfqpoint{0.904718in}{0.773515in}}%
\pgfpathlineto{\pgfqpoint{0.906473in}{0.778649in}}%
\pgfpathclose%
\pgfusepath{fill}%
\end{pgfscope}%
\begin{pgfscope}%
\pgfpathrectangle{\pgfqpoint{0.407592in}{0.222904in}}{\pgfqpoint{1.411163in}{1.411163in}}%
\pgfusepath{clip}%
\pgfsetbuttcap%
\pgfsetroundjoin%
\definecolor{currentfill}{rgb}{0.749020,0.772549,0.792157}%
\pgfsetfillcolor{currentfill}%
\pgfsetfillopacity{0.160000}%
\pgfsetlinewidth{0.000000pt}%
\definecolor{currentstroke}{rgb}{0.000000,0.000000,0.000000}%
\pgfsetstrokecolor{currentstroke}%
\pgfsetdash{}{0pt}%
\pgfpathmoveto{\pgfqpoint{0.957406in}{0.722952in}}%
\pgfpathlineto{\pgfqpoint{0.944024in}{0.734278in}}%
\pgfpathlineto{\pgfqpoint{0.930624in}{0.745620in}}%
\pgfpathlineto{\pgfqpoint{0.930541in}{0.741165in}}%
\pgfpathlineto{\pgfqpoint{0.931156in}{0.736724in}}%
\pgfpathlineto{\pgfqpoint{0.944544in}{0.726104in}}%
\pgfpathlineto{\pgfqpoint{0.957907in}{0.715503in}}%
\pgfpathlineto{\pgfqpoint{0.957366in}{0.719220in}}%
\pgfpathlineto{\pgfqpoint{0.957406in}{0.722952in}}%
\pgfpathclose%
\pgfusepath{fill}%
\end{pgfscope}%
\begin{pgfscope}%
\pgfpathrectangle{\pgfqpoint{0.407592in}{0.222904in}}{\pgfqpoint{1.411163in}{1.411163in}}%
\pgfusepath{clip}%
\pgfsetbuttcap%
\pgfsetroundjoin%
\definecolor{currentfill}{rgb}{0.749020,0.772549,0.792157}%
\pgfsetfillcolor{currentfill}%
\pgfsetfillopacity{0.160000}%
\pgfsetlinewidth{0.000000pt}%
\definecolor{currentstroke}{rgb}{0.000000,0.000000,0.000000}%
\pgfsetstrokecolor{currentstroke}%
\pgfsetdash{}{0pt}%
\pgfpathmoveto{\pgfqpoint{1.047590in}{0.641293in}}%
\pgfpathlineto{\pgfqpoint{1.036840in}{0.649047in}}%
\pgfpathlineto{\pgfqpoint{1.026055in}{0.656828in}}%
\pgfpathlineto{\pgfqpoint{1.029315in}{0.655102in}}%
\pgfpathlineto{\pgfqpoint{1.032846in}{0.653477in}}%
\pgfpathlineto{\pgfqpoint{1.042498in}{0.646254in}}%
\pgfpathlineto{\pgfqpoint{1.052114in}{0.639058in}}%
\pgfpathlineto{\pgfqpoint{1.049763in}{0.640141in}}%
\pgfpathlineto{\pgfqpoint{1.047590in}{0.641293in}}%
\pgfpathclose%
\pgfusepath{fill}%
\end{pgfscope}%
\begin{pgfscope}%
\pgfpathrectangle{\pgfqpoint{0.407592in}{0.222904in}}{\pgfqpoint{1.411163in}{1.411163in}}%
\pgfusepath{clip}%
\pgfsetbuttcap%
\pgfsetroundjoin%
\definecolor{currentfill}{rgb}{0.749020,0.772549,0.792157}%
\pgfsetfillcolor{currentfill}%
\pgfsetfillopacity{0.160000}%
\pgfsetlinewidth{0.000000pt}%
\definecolor{currentstroke}{rgb}{0.000000,0.000000,0.000000}%
\pgfsetstrokecolor{currentstroke}%
\pgfsetdash{}{0pt}%
\pgfpathmoveto{\pgfqpoint{1.190916in}{0.689056in}}%
\pgfpathlineto{\pgfqpoint{1.203625in}{0.698968in}}%
\pgfpathlineto{\pgfqpoint{1.216362in}{0.708901in}}%
\pgfpathlineto{\pgfqpoint{1.218710in}{0.712552in}}%
\pgfpathlineto{\pgfqpoint{1.220488in}{0.716247in}}%
\pgfpathlineto{\pgfqpoint{1.207360in}{0.705592in}}%
\pgfpathlineto{\pgfqpoint{1.194254in}{0.694956in}}%
\pgfpathlineto{\pgfqpoint{1.192812in}{0.691988in}}%
\pgfpathlineto{\pgfqpoint{1.190916in}{0.689056in}}%
\pgfpathclose%
\pgfusepath{fill}%
\end{pgfscope}%
\begin{pgfscope}%
\pgfpathrectangle{\pgfqpoint{0.407592in}{0.222904in}}{\pgfqpoint{1.411163in}{1.411163in}}%
\pgfusepath{clip}%
\pgfsetbuttcap%
\pgfsetroundjoin%
\definecolor{currentfill}{rgb}{0.749020,0.772549,0.792157}%
\pgfsetfillcolor{currentfill}%
\pgfsetfillopacity{0.160000}%
\pgfsetlinewidth{0.000000pt}%
\definecolor{currentstroke}{rgb}{0.000000,0.000000,0.000000}%
\pgfsetstrokecolor{currentstroke}%
\pgfsetdash{}{0pt}%
\pgfpathmoveto{\pgfqpoint{0.984562in}{0.694358in}}%
\pgfpathlineto{\pgfqpoint{0.971246in}{0.704921in}}%
\pgfpathlineto{\pgfqpoint{0.957907in}{0.715503in}}%
\pgfpathlineto{\pgfqpoint{0.959028in}{0.711816in}}%
\pgfpathlineto{\pgfqpoint{0.960727in}{0.708176in}}%
\pgfpathlineto{\pgfqpoint{0.973807in}{0.698314in}}%
\pgfpathlineto{\pgfqpoint{0.986858in}{0.688474in}}%
\pgfpathlineto{\pgfqpoint{0.985479in}{0.691396in}}%
\pgfpathlineto{\pgfqpoint{0.984562in}{0.694358in}}%
\pgfpathclose%
\pgfusepath{fill}%
\end{pgfscope}%
\begin{pgfscope}%
\pgfpathrectangle{\pgfqpoint{0.407592in}{0.222904in}}{\pgfqpoint{1.411163in}{1.411163in}}%
\pgfusepath{clip}%
\pgfsetbuttcap%
\pgfsetroundjoin%
\definecolor{currentfill}{rgb}{0.749020,0.772549,0.792157}%
\pgfsetfillcolor{currentfill}%
\pgfsetfillopacity{0.160000}%
\pgfsetlinewidth{0.000000pt}%
\definecolor{currentstroke}{rgb}{0.000000,0.000000,0.000000}%
\pgfsetstrokecolor{currentstroke}%
\pgfsetdash{}{0pt}%
\pgfpathmoveto{\pgfqpoint{0.784513in}{0.916808in}}%
\pgfpathlineto{\pgfqpoint{0.771694in}{0.929651in}}%
\pgfpathlineto{\pgfqpoint{0.758868in}{0.942501in}}%
\pgfpathlineto{\pgfqpoint{0.752240in}{0.933583in}}%
\pgfpathlineto{\pgfqpoint{0.747043in}{0.924510in}}%
\pgfpathlineto{\pgfqpoint{0.760397in}{0.912292in}}%
\pgfpathlineto{\pgfqpoint{0.773739in}{0.900086in}}%
\pgfpathlineto{\pgfqpoint{0.778464in}{0.908517in}}%
\pgfpathlineto{\pgfqpoint{0.784513in}{0.916808in}}%
\pgfpathclose%
\pgfusepath{fill}%
\end{pgfscope}%
\begin{pgfscope}%
\pgfpathrectangle{\pgfqpoint{0.407592in}{0.222904in}}{\pgfqpoint{1.411163in}{1.411163in}}%
\pgfusepath{clip}%
\pgfsetbuttcap%
\pgfsetroundjoin%
\definecolor{currentfill}{rgb}{0.749020,0.772549,0.792157}%
\pgfsetfillcolor{currentfill}%
\pgfsetfillopacity{0.160000}%
\pgfsetlinewidth{0.000000pt}%
\definecolor{currentstroke}{rgb}{0.000000,0.000000,0.000000}%
\pgfsetstrokecolor{currentstroke}%
\pgfsetdash{}{0pt}%
\pgfpathmoveto{\pgfqpoint{1.122516in}{0.638244in}}%
\pgfpathlineto{\pgfqpoint{1.130664in}{0.645237in}}%
\pgfpathlineto{\pgfqpoint{1.138843in}{0.652257in}}%
\pgfpathlineto{\pgfqpoint{1.142897in}{0.653795in}}%
\pgfpathlineto{\pgfqpoint{1.146714in}{0.655440in}}%
\pgfpathlineto{\pgfqpoint{1.137218in}{0.647890in}}%
\pgfpathlineto{\pgfqpoint{1.127755in}{0.640367in}}%
\pgfpathlineto{\pgfqpoint{1.125213in}{0.639269in}}%
\pgfpathlineto{\pgfqpoint{1.122516in}{0.638244in}}%
\pgfpathclose%
\pgfusepath{fill}%
\end{pgfscope}%
\begin{pgfscope}%
\pgfpathrectangle{\pgfqpoint{0.407592in}{0.222904in}}{\pgfqpoint{1.411163in}{1.411163in}}%
\pgfusepath{clip}%
\pgfsetbuttcap%
\pgfsetroundjoin%
\definecolor{currentfill}{rgb}{0.749020,0.772549,0.792157}%
\pgfsetfillcolor{currentfill}%
\pgfsetfillopacity{0.160000}%
\pgfsetlinewidth{0.000000pt}%
\definecolor{currentstroke}{rgb}{0.000000,0.000000,0.000000}%
\pgfsetstrokecolor{currentstroke}%
\pgfsetdash{}{0pt}%
\pgfpathmoveto{\pgfqpoint{1.434838in}{0.926341in}}%
\pgfpathlineto{\pgfqpoint{1.448253in}{0.938635in}}%
\pgfpathlineto{\pgfqpoint{1.461680in}{0.950939in}}%
\pgfpathlineto{\pgfqpoint{1.457844in}{0.960613in}}%
\pgfpathlineto{\pgfqpoint{1.452439in}{0.970108in}}%
\pgfpathlineto{\pgfqpoint{1.439422in}{0.957185in}}%
\pgfpathlineto{\pgfqpoint{1.426410in}{0.944268in}}%
\pgfpathlineto{\pgfqpoint{1.431354in}{0.935386in}}%
\pgfpathlineto{\pgfqpoint{1.434838in}{0.926341in}}%
\pgfpathclose%
\pgfusepath{fill}%
\end{pgfscope}%
\begin{pgfscope}%
\pgfpathrectangle{\pgfqpoint{0.407592in}{0.222904in}}{\pgfqpoint{1.411163in}{1.411163in}}%
\pgfusepath{clip}%
\pgfsetbuttcap%
\pgfsetroundjoin%
\definecolor{currentfill}{rgb}{0.749020,0.772549,0.792157}%
\pgfsetfillcolor{currentfill}%
\pgfsetfillopacity{0.160000}%
\pgfsetlinewidth{0.000000pt}%
\definecolor{currentstroke}{rgb}{0.000000,0.000000,0.000000}%
\pgfsetstrokecolor{currentstroke}%
\pgfsetdash{}{0pt}%
\pgfpathmoveto{\pgfqpoint{1.302409in}{0.792323in}}%
\pgfpathlineto{\pgfqpoint{1.315820in}{0.803813in}}%
\pgfpathlineto{\pgfqpoint{1.329249in}{0.815318in}}%
\pgfpathlineto{\pgfqpoint{1.329122in}{0.821878in}}%
\pgfpathlineto{\pgfqpoint{1.327952in}{0.828388in}}%
\pgfpathlineto{\pgfqpoint{1.314645in}{0.816193in}}%
\pgfpathlineto{\pgfqpoint{1.301350in}{0.804009in}}%
\pgfpathlineto{\pgfqpoint{1.302343in}{0.798187in}}%
\pgfpathlineto{\pgfqpoint{1.302409in}{0.792323in}}%
\pgfpathclose%
\pgfusepath{fill}%
\end{pgfscope}%
\begin{pgfscope}%
\pgfpathrectangle{\pgfqpoint{0.407592in}{0.222904in}}{\pgfqpoint{1.411163in}{1.411163in}}%
\pgfusepath{clip}%
\pgfsetbuttcap%
\pgfsetroundjoin%
\definecolor{currentfill}{rgb}{0.749020,0.772549,0.792157}%
\pgfsetfillcolor{currentfill}%
\pgfsetfillopacity{0.160000}%
\pgfsetlinewidth{0.000000pt}%
\definecolor{currentstroke}{rgb}{0.000000,0.000000,0.000000}%
\pgfsetstrokecolor{currentstroke}%
\pgfsetdash{}{0pt}%
\pgfpathmoveto{\pgfqpoint{1.052114in}{0.639058in}}%
\pgfpathlineto{\pgfqpoint{1.042498in}{0.646254in}}%
\pgfpathlineto{\pgfqpoint{1.032846in}{0.653477in}}%
\pgfpathlineto{\pgfqpoint{1.036632in}{0.651962in}}%
\pgfpathlineto{\pgfqpoint{1.040657in}{0.650563in}}%
\pgfpathlineto{\pgfqpoint{1.049000in}{0.643826in}}%
\pgfpathlineto{\pgfqpoint{1.057311in}{0.637114in}}%
\pgfpathlineto{\pgfqpoint{1.054634in}{0.638047in}}%
\pgfpathlineto{\pgfqpoint{1.052114in}{0.639058in}}%
\pgfpathclose%
\pgfusepath{fill}%
\end{pgfscope}%
\begin{pgfscope}%
\pgfpathrectangle{\pgfqpoint{0.407592in}{0.222904in}}{\pgfqpoint{1.411163in}{1.411163in}}%
\pgfusepath{clip}%
\pgfsetbuttcap%
\pgfsetroundjoin%
\definecolor{currentfill}{rgb}{0.749020,0.772549,0.792157}%
\pgfsetfillcolor{currentfill}%
\pgfsetfillopacity{0.160000}%
\pgfsetlinewidth{0.000000pt}%
\definecolor{currentstroke}{rgb}{0.000000,0.000000,0.000000}%
\pgfsetstrokecolor{currentstroke}%
\pgfsetdash{}{0pt}%
\pgfpathmoveto{\pgfqpoint{1.156589in}{0.660946in}}%
\pgfpathlineto{\pgfqpoint{1.167764in}{0.669436in}}%
\pgfpathlineto{\pgfqpoint{1.178973in}{0.677951in}}%
\pgfpathlineto{\pgfqpoint{1.182590in}{0.680609in}}%
\pgfpathlineto{\pgfqpoint{1.185795in}{0.683354in}}%
\pgfpathlineto{\pgfqpoint{1.173741in}{0.674169in}}%
\pgfpathlineto{\pgfqpoint{1.161718in}{0.665008in}}%
\pgfpathlineto{\pgfqpoint{1.159307in}{0.662944in}}%
\pgfpathlineto{\pgfqpoint{1.156589in}{0.660946in}}%
\pgfpathclose%
\pgfusepath{fill}%
\end{pgfscope}%
\begin{pgfscope}%
\pgfpathrectangle{\pgfqpoint{0.407592in}{0.222904in}}{\pgfqpoint{1.411163in}{1.411163in}}%
\pgfusepath{clip}%
\pgfsetbuttcap%
\pgfsetroundjoin%
\definecolor{currentfill}{rgb}{0.749020,0.772549,0.792157}%
\pgfsetfillcolor{currentfill}%
\pgfsetfillopacity{0.160000}%
\pgfsetlinewidth{0.000000pt}%
\definecolor{currentstroke}{rgb}{0.000000,0.000000,0.000000}%
\pgfsetstrokecolor{currentstroke}%
\pgfsetdash{}{0pt}%
\pgfpathmoveto{\pgfqpoint{0.758868in}{0.942501in}}%
\pgfpathlineto{\pgfqpoint{0.746035in}{0.955357in}}%
\pgfpathlineto{\pgfqpoint{0.733196in}{0.968219in}}%
\pgfpathlineto{\pgfqpoint{0.725977in}{0.958685in}}%
\pgfpathlineto{\pgfqpoint{0.720296in}{0.948980in}}%
\pgfpathlineto{\pgfqpoint{0.733675in}{0.936740in}}%
\pgfpathlineto{\pgfqpoint{0.747043in}{0.924510in}}%
\pgfpathlineto{\pgfqpoint{0.752240in}{0.933583in}}%
\pgfpathlineto{\pgfqpoint{0.758868in}{0.942501in}}%
\pgfpathclose%
\pgfusepath{fill}%
\end{pgfscope}%
\begin{pgfscope}%
\pgfpathrectangle{\pgfqpoint{0.407592in}{0.222904in}}{\pgfqpoint{1.411163in}{1.411163in}}%
\pgfusepath{clip}%
\pgfsetbuttcap%
\pgfsetroundjoin%
\definecolor{currentfill}{rgb}{0.749020,0.772549,0.792157}%
\pgfsetfillcolor{currentfill}%
\pgfsetfillopacity{0.160000}%
\pgfsetlinewidth{0.000000pt}%
\definecolor{currentstroke}{rgb}{0.000000,0.000000,0.000000}%
\pgfsetstrokecolor{currentstroke}%
\pgfsetdash{}{0pt}%
\pgfpathmoveto{\pgfqpoint{1.116702in}{0.636428in}}%
\pgfpathlineto{\pgfqpoint{1.123387in}{0.642967in}}%
\pgfpathlineto{\pgfqpoint{1.130099in}{0.649533in}}%
\pgfpathlineto{\pgfqpoint{1.134571in}{0.650834in}}%
\pgfpathlineto{\pgfqpoint{1.138843in}{0.652257in}}%
\pgfpathlineto{\pgfqpoint{1.130664in}{0.645237in}}%
\pgfpathlineto{\pgfqpoint{1.122516in}{0.638244in}}%
\pgfpathlineto{\pgfqpoint{1.119675in}{0.637295in}}%
\pgfpathlineto{\pgfqpoint{1.116702in}{0.636428in}}%
\pgfpathclose%
\pgfusepath{fill}%
\end{pgfscope}%
\begin{pgfscope}%
\pgfpathrectangle{\pgfqpoint{0.407592in}{0.222904in}}{\pgfqpoint{1.411163in}{1.411163in}}%
\pgfusepath{clip}%
\pgfsetbuttcap%
\pgfsetroundjoin%
\definecolor{currentfill}{rgb}{0.749020,0.772549,0.792157}%
\pgfsetfillcolor{currentfill}%
\pgfsetfillopacity{0.160000}%
\pgfsetlinewidth{0.000000pt}%
\definecolor{currentstroke}{rgb}{0.000000,0.000000,0.000000}%
\pgfsetstrokecolor{currentstroke}%
\pgfsetdash{}{0pt}%
\pgfpathmoveto{\pgfqpoint{0.880026in}{0.802846in}}%
\pgfpathlineto{\pgfqpoint{0.866784in}{0.814961in}}%
\pgfpathlineto{\pgfqpoint{0.853529in}{0.827087in}}%
\pgfpathlineto{\pgfqpoint{0.851164in}{0.820565in}}%
\pgfpathlineto{\pgfqpoint{0.849835in}{0.813999in}}%
\pgfpathlineto{\pgfqpoint{0.863346in}{0.802563in}}%
\pgfpathlineto{\pgfqpoint{0.876838in}{0.791143in}}%
\pgfpathlineto{\pgfqpoint{0.877972in}{0.797013in}}%
\pgfpathlineto{\pgfqpoint{0.880026in}{0.802846in}}%
\pgfpathclose%
\pgfusepath{fill}%
\end{pgfscope}%
\begin{pgfscope}%
\pgfpathrectangle{\pgfqpoint{0.407592in}{0.222904in}}{\pgfqpoint{1.411163in}{1.411163in}}%
\pgfusepath{clip}%
\pgfsetbuttcap%
\pgfsetroundjoin%
\definecolor{currentfill}{rgb}{0.749020,0.772549,0.792157}%
\pgfsetfillcolor{currentfill}%
\pgfsetfillopacity{0.160000}%
\pgfsetlinewidth{0.000000pt}%
\definecolor{currentstroke}{rgb}{0.000000,0.000000,0.000000}%
\pgfsetstrokecolor{currentstroke}%
\pgfsetdash{}{0pt}%
\pgfpathmoveto{\pgfqpoint{1.015983in}{0.664589in}}%
\pgfpathlineto{\pgfqpoint{1.003495in}{0.673681in}}%
\pgfpathlineto{\pgfqpoint{0.990974in}{0.682797in}}%
\pgfpathlineto{\pgfqpoint{0.993695in}{0.680068in}}%
\pgfpathlineto{\pgfqpoint{0.996845in}{0.677429in}}%
\pgfpathlineto{\pgfqpoint{1.008641in}{0.668979in}}%
\pgfpathlineto{\pgfqpoint{1.020402in}{0.660554in}}%
\pgfpathlineto{\pgfqpoint{1.018033in}{0.662537in}}%
\pgfpathlineto{\pgfqpoint{1.015983in}{0.664589in}}%
\pgfpathclose%
\pgfusepath{fill}%
\end{pgfscope}%
\begin{pgfscope}%
\pgfpathrectangle{\pgfqpoint{0.407592in}{0.222904in}}{\pgfqpoint{1.411163in}{1.411163in}}%
\pgfusepath{clip}%
\pgfsetbuttcap%
\pgfsetroundjoin%
\definecolor{currentfill}{rgb}{0.749020,0.772549,0.792157}%
\pgfsetfillcolor{currentfill}%
\pgfsetfillopacity{0.160000}%
\pgfsetlinewidth{0.000000pt}%
\definecolor{currentstroke}{rgb}{0.000000,0.000000,0.000000}%
\pgfsetstrokecolor{currentstroke}%
\pgfsetdash{}{0pt}%
\pgfpathmoveto{\pgfqpoint{1.461680in}{0.950939in}}%
\pgfpathlineto{\pgfqpoint{1.475119in}{0.963255in}}%
\pgfpathlineto{\pgfqpoint{1.488569in}{0.975581in}}%
\pgfpathlineto{\pgfqpoint{1.484371in}{0.985874in}}%
\pgfpathlineto{\pgfqpoint{1.478494in}{0.995973in}}%
\pgfpathlineto{\pgfqpoint{1.465463in}{0.983037in}}%
\pgfpathlineto{\pgfqpoint{1.452439in}{0.970108in}}%
\pgfpathlineto{\pgfqpoint{1.457844in}{0.960613in}}%
\pgfpathlineto{\pgfqpoint{1.461680in}{0.950939in}}%
\pgfpathclose%
\pgfusepath{fill}%
\end{pgfscope}%
\begin{pgfscope}%
\pgfpathrectangle{\pgfqpoint{0.407592in}{0.222904in}}{\pgfqpoint{1.411163in}{1.411163in}}%
\pgfusepath{clip}%
\pgfsetbuttcap%
\pgfsetroundjoin%
\definecolor{currentfill}{rgb}{0.749020,0.772549,0.792157}%
\pgfsetfillcolor{currentfill}%
\pgfsetfillopacity{0.160000}%
\pgfsetlinewidth{0.000000pt}%
\definecolor{currentstroke}{rgb}{0.000000,0.000000,0.000000}%
\pgfsetstrokecolor{currentstroke}%
\pgfsetdash{}{0pt}%
\pgfpathmoveto{\pgfqpoint{1.246815in}{0.737613in}}%
\pgfpathlineto{\pgfqpoint{1.260013in}{0.748324in}}%
\pgfpathlineto{\pgfqpoint{1.273234in}{0.759054in}}%
\pgfpathlineto{\pgfqpoint{1.274845in}{0.764216in}}%
\pgfpathlineto{\pgfqpoint{1.275641in}{0.769388in}}%
\pgfpathlineto{\pgfqpoint{1.262283in}{0.757944in}}%
\pgfpathlineto{\pgfqpoint{1.248943in}{0.746514in}}%
\pgfpathlineto{\pgfqpoint{1.248227in}{0.742058in}}%
\pgfpathlineto{\pgfqpoint{1.246815in}{0.737613in}}%
\pgfpathclose%
\pgfusepath{fill}%
\end{pgfscope}%
\begin{pgfscope}%
\pgfpathrectangle{\pgfqpoint{0.407592in}{0.222904in}}{\pgfqpoint{1.411163in}{1.411163in}}%
\pgfusepath{clip}%
\pgfsetbuttcap%
\pgfsetroundjoin%
\definecolor{currentfill}{rgb}{0.749020,0.772549,0.792157}%
\pgfsetfillcolor{currentfill}%
\pgfsetfillopacity{0.160000}%
\pgfsetlinewidth{0.000000pt}%
\definecolor{currentstroke}{rgb}{0.000000,0.000000,0.000000}%
\pgfsetstrokecolor{currentstroke}%
\pgfsetdash{}{0pt}%
\pgfpathmoveto{\pgfqpoint{1.057311in}{0.637114in}}%
\pgfpathlineto{\pgfqpoint{1.049000in}{0.643826in}}%
\pgfpathlineto{\pgfqpoint{1.040657in}{0.650563in}}%
\pgfpathlineto{\pgfqpoint{1.044902in}{0.649287in}}%
\pgfpathlineto{\pgfqpoint{1.049350in}{0.648138in}}%
\pgfpathlineto{\pgfqpoint{1.056235in}{0.641805in}}%
\pgfpathlineto{\pgfqpoint{1.063090in}{0.635499in}}%
\pgfpathlineto{\pgfqpoint{1.060134in}{0.636264in}}%
\pgfpathlineto{\pgfqpoint{1.057311in}{0.637114in}}%
\pgfpathclose%
\pgfusepath{fill}%
\end{pgfscope}%
\begin{pgfscope}%
\pgfpathrectangle{\pgfqpoint{0.407592in}{0.222904in}}{\pgfqpoint{1.411163in}{1.411163in}}%
\pgfusepath{clip}%
\pgfsetbuttcap%
\pgfsetroundjoin%
\definecolor{currentfill}{rgb}{0.749020,0.772549,0.792157}%
\pgfsetfillcolor{currentfill}%
\pgfsetfillopacity{0.160000}%
\pgfsetlinewidth{0.000000pt}%
\definecolor{currentstroke}{rgb}{0.000000,0.000000,0.000000}%
\pgfsetstrokecolor{currentstroke}%
\pgfsetdash{}{0pt}%
\pgfpathmoveto{\pgfqpoint{0.930624in}{0.745620in}}%
\pgfpathlineto{\pgfqpoint{0.917205in}{0.756978in}}%
\pgfpathlineto{\pgfqpoint{0.903768in}{0.768351in}}%
\pgfpathlineto{\pgfqpoint{0.903630in}{0.763179in}}%
\pgfpathlineto{\pgfqpoint{0.904309in}{0.758021in}}%
\pgfpathlineto{\pgfqpoint{0.917745in}{0.747363in}}%
\pgfpathlineto{\pgfqpoint{0.931156in}{0.736724in}}%
\pgfpathlineto{\pgfqpoint{0.930541in}{0.741165in}}%
\pgfpathlineto{\pgfqpoint{0.930624in}{0.745620in}}%
\pgfpathclose%
\pgfusepath{fill}%
\end{pgfscope}%
\begin{pgfscope}%
\pgfpathrectangle{\pgfqpoint{0.407592in}{0.222904in}}{\pgfqpoint{1.411163in}{1.411163in}}%
\pgfusepath{clip}%
\pgfsetbuttcap%
\pgfsetroundjoin%
\definecolor{currentfill}{rgb}{0.749020,0.772549,0.792157}%
\pgfsetfillcolor{currentfill}%
\pgfsetfillopacity{0.160000}%
\pgfsetlinewidth{0.000000pt}%
\definecolor{currentstroke}{rgb}{0.000000,0.000000,0.000000}%
\pgfsetstrokecolor{currentstroke}%
\pgfsetdash{}{0pt}%
\pgfpathmoveto{\pgfqpoint{1.110416in}{0.634951in}}%
\pgfpathlineto{\pgfqpoint{1.115516in}{0.641121in}}%
\pgfpathlineto{\pgfqpoint{1.120638in}{0.647317in}}%
\pgfpathlineto{\pgfqpoint{1.125448in}{0.648358in}}%
\pgfpathlineto{\pgfqpoint{1.130099in}{0.649533in}}%
\pgfpathlineto{\pgfqpoint{1.123387in}{0.642967in}}%
\pgfpathlineto{\pgfqpoint{1.116702in}{0.636428in}}%
\pgfpathlineto{\pgfqpoint{1.113611in}{0.635645in}}%
\pgfpathlineto{\pgfqpoint{1.110416in}{0.634951in}}%
\pgfpathclose%
\pgfusepath{fill}%
\end{pgfscope}%
\begin{pgfscope}%
\pgfpathrectangle{\pgfqpoint{0.407592in}{0.222904in}}{\pgfqpoint{1.411163in}{1.411163in}}%
\pgfusepath{clip}%
\pgfsetbuttcap%
\pgfsetroundjoin%
\definecolor{currentfill}{rgb}{0.749020,0.772549,0.792157}%
\pgfsetfillcolor{currentfill}%
\pgfsetfillopacity{0.160000}%
\pgfsetlinewidth{0.000000pt}%
\definecolor{currentstroke}{rgb}{0.000000,0.000000,0.000000}%
\pgfsetstrokecolor{currentstroke}%
\pgfsetdash{}{0pt}%
\pgfpathmoveto{\pgfqpoint{0.733196in}{0.968219in}}%
\pgfpathlineto{\pgfqpoint{0.720350in}{0.981088in}}%
\pgfpathlineto{\pgfqpoint{0.707498in}{0.993964in}}%
\pgfpathlineto{\pgfqpoint{0.699676in}{0.983823in}}%
\pgfpathlineto{\pgfqpoint{0.693499in}{0.973497in}}%
\pgfpathlineto{\pgfqpoint{0.706903in}{0.961233in}}%
\pgfpathlineto{\pgfqpoint{0.720296in}{0.948980in}}%
\pgfpathlineto{\pgfqpoint{0.725977in}{0.958685in}}%
\pgfpathlineto{\pgfqpoint{0.733196in}{0.968219in}}%
\pgfpathclose%
\pgfusepath{fill}%
\end{pgfscope}%
\begin{pgfscope}%
\pgfpathrectangle{\pgfqpoint{0.407592in}{0.222904in}}{\pgfqpoint{1.411163in}{1.411163in}}%
\pgfusepath{clip}%
\pgfsetbuttcap%
\pgfsetroundjoin%
\definecolor{currentfill}{rgb}{0.749020,0.772549,0.792157}%
\pgfsetfillcolor{currentfill}%
\pgfsetfillopacity{0.160000}%
\pgfsetlinewidth{0.000000pt}%
\definecolor{currentstroke}{rgb}{0.000000,0.000000,0.000000}%
\pgfsetstrokecolor{currentstroke}%
\pgfsetdash{}{0pt}%
\pgfpathmoveto{\pgfqpoint{1.329249in}{0.815318in}}%
\pgfpathlineto{\pgfqpoint{1.342696in}{0.826839in}}%
\pgfpathlineto{\pgfqpoint{1.356160in}{0.838375in}}%
\pgfpathlineto{\pgfqpoint{1.355961in}{0.845622in}}%
\pgfpathlineto{\pgfqpoint{1.354602in}{0.852811in}}%
\pgfpathlineto{\pgfqpoint{1.341271in}{0.840594in}}%
\pgfpathlineto{\pgfqpoint{1.327952in}{0.828388in}}%
\pgfpathlineto{\pgfqpoint{1.329122in}{0.821878in}}%
\pgfpathlineto{\pgfqpoint{1.329249in}{0.815318in}}%
\pgfpathclose%
\pgfusepath{fill}%
\end{pgfscope}%
\begin{pgfscope}%
\pgfpathrectangle{\pgfqpoint{0.407592in}{0.222904in}}{\pgfqpoint{1.411163in}{1.411163in}}%
\pgfusepath{clip}%
\pgfsetbuttcap%
\pgfsetroundjoin%
\definecolor{currentfill}{rgb}{0.749020,0.772549,0.792157}%
\pgfsetfillcolor{currentfill}%
\pgfsetfillopacity{0.160000}%
\pgfsetlinewidth{0.000000pt}%
\definecolor{currentstroke}{rgb}{0.000000,0.000000,0.000000}%
\pgfsetstrokecolor{currentstroke}%
\pgfsetdash{}{0pt}%
\pgfpathmoveto{\pgfqpoint{1.488569in}{0.975581in}}%
\pgfpathlineto{\pgfqpoint{1.502032in}{0.987919in}}%
\pgfpathlineto{\pgfqpoint{1.515507in}{1.000268in}}%
\pgfpathlineto{\pgfqpoint{1.510934in}{1.011170in}}%
\pgfpathlineto{\pgfqpoint{1.504572in}{1.021862in}}%
\pgfpathlineto{\pgfqpoint{1.491530in}{1.008914in}}%
\pgfpathlineto{\pgfqpoint{1.478494in}{0.995973in}}%
\pgfpathlineto{\pgfqpoint{1.484371in}{0.985874in}}%
\pgfpathlineto{\pgfqpoint{1.488569in}{0.975581in}}%
\pgfpathclose%
\pgfusepath{fill}%
\end{pgfscope}%
\begin{pgfscope}%
\pgfpathrectangle{\pgfqpoint{0.407592in}{0.222904in}}{\pgfqpoint{1.411163in}{1.411163in}}%
\pgfusepath{clip}%
\pgfsetbuttcap%
\pgfsetroundjoin%
\definecolor{currentfill}{rgb}{0.749020,0.772549,0.792157}%
\pgfsetfillcolor{currentfill}%
\pgfsetfillopacity{0.160000}%
\pgfsetlinewidth{0.000000pt}%
\definecolor{currentstroke}{rgb}{0.000000,0.000000,0.000000}%
\pgfsetstrokecolor{currentstroke}%
\pgfsetdash{}{0pt}%
\pgfpathmoveto{\pgfqpoint{1.063090in}{0.635499in}}%
\pgfpathlineto{\pgfqpoint{1.056235in}{0.641805in}}%
\pgfpathlineto{\pgfqpoint{1.049350in}{0.648138in}}%
\pgfpathlineto{\pgfqpoint{1.053980in}{0.647124in}}%
\pgfpathlineto{\pgfqpoint{1.058772in}{0.646248in}}%
\pgfpathlineto{\pgfqpoint{1.064072in}{0.640230in}}%
\pgfpathlineto{\pgfqpoint{1.069349in}{0.634239in}}%
\pgfpathlineto{\pgfqpoint{1.066167in}{0.634823in}}%
\pgfpathlineto{\pgfqpoint{1.063090in}{0.635499in}}%
\pgfpathclose%
\pgfusepath{fill}%
\end{pgfscope}%
\begin{pgfscope}%
\pgfpathrectangle{\pgfqpoint{0.407592in}{0.222904in}}{\pgfqpoint{1.411163in}{1.411163in}}%
\pgfusepath{clip}%
\pgfsetbuttcap%
\pgfsetroundjoin%
\definecolor{currentfill}{rgb}{0.749020,0.772549,0.792157}%
\pgfsetfillcolor{currentfill}%
\pgfsetfillopacity{0.160000}%
\pgfsetlinewidth{0.000000pt}%
\definecolor{currentstroke}{rgb}{0.000000,0.000000,0.000000}%
\pgfsetstrokecolor{currentstroke}%
\pgfsetdash{}{0pt}%
\pgfpathmoveto{\pgfqpoint{1.153575in}{0.659024in}}%
\pgfpathlineto{\pgfqpoint{1.164250in}{0.667195in}}%
\pgfpathlineto{\pgfqpoint{1.174958in}{0.675392in}}%
\pgfpathlineto{\pgfqpoint{1.178973in}{0.677951in}}%
\pgfpathlineto{\pgfqpoint{1.167764in}{0.669436in}}%
\pgfpathlineto{\pgfqpoint{1.156589in}{0.660946in}}%
\pgfusepath{fill}%
\end{pgfscope}%
\begin{pgfscope}%
\pgfpathrectangle{\pgfqpoint{0.407592in}{0.222904in}}{\pgfqpoint{1.411163in}{1.411163in}}%
\pgfusepath{clip}%
\pgfsetbuttcap%
\pgfsetroundjoin%
\definecolor{currentfill}{rgb}{0.749020,0.772549,0.792157}%
\pgfsetfillcolor{currentfill}%
\pgfsetfillopacity{0.160000}%
\pgfsetlinewidth{0.000000pt}%
\definecolor{currentstroke}{rgb}{0.000000,0.000000,0.000000}%
\pgfsetstrokecolor{currentstroke}%
\pgfsetdash{}{0pt}%
\pgfpathmoveto{\pgfqpoint{1.185795in}{0.683354in}}%
\pgfpathlineto{\pgfqpoint{1.197880in}{0.692563in}}%
\pgfpathlineto{\pgfqpoint{1.209996in}{0.701796in}}%
\pgfpathlineto{\pgfqpoint{1.213454in}{0.705310in}}%
\pgfpathlineto{\pgfqpoint{1.216362in}{0.708901in}}%
\pgfpathlineto{\pgfqpoint{1.203625in}{0.698968in}}%
\pgfpathlineto{\pgfqpoint{1.190916in}{0.689056in}}%
\pgfpathlineto{\pgfqpoint{1.188574in}{0.686174in}}%
\pgfpathlineto{\pgfqpoint{1.185795in}{0.683354in}}%
\pgfpathclose%
\pgfusepath{fill}%
\end{pgfscope}%
\begin{pgfscope}%
\pgfpathrectangle{\pgfqpoint{0.407592in}{0.222904in}}{\pgfqpoint{1.411163in}{1.411163in}}%
\pgfusepath{clip}%
\pgfsetbuttcap%
\pgfsetroundjoin%
\definecolor{currentfill}{rgb}{0.749020,0.772549,0.792157}%
\pgfsetfillcolor{currentfill}%
\pgfsetfillopacity{0.160000}%
\pgfsetlinewidth{0.000000pt}%
\definecolor{currentstroke}{rgb}{0.000000,0.000000,0.000000}%
\pgfsetstrokecolor{currentstroke}%
\pgfsetdash{}{0pt}%
\pgfpathmoveto{\pgfqpoint{1.216362in}{0.708901in}}%
\pgfpathlineto{\pgfqpoint{1.229127in}{0.718857in}}%
\pgfpathlineto{\pgfqpoint{1.241920in}{0.728833in}}%
\pgfpathlineto{\pgfqpoint{1.244710in}{0.733198in}}%
\pgfpathlineto{\pgfqpoint{1.246815in}{0.737613in}}%
\pgfpathlineto{\pgfqpoint{1.233640in}{0.726920in}}%
\pgfpathlineto{\pgfqpoint{1.220488in}{0.716247in}}%
\pgfpathlineto{\pgfqpoint{1.218710in}{0.712552in}}%
\pgfpathlineto{\pgfqpoint{1.216362in}{0.708901in}}%
\pgfpathclose%
\pgfusepath{fill}%
\end{pgfscope}%
\begin{pgfscope}%
\pgfpathrectangle{\pgfqpoint{0.407592in}{0.222904in}}{\pgfqpoint{1.411163in}{1.411163in}}%
\pgfusepath{clip}%
\pgfsetbuttcap%
\pgfsetroundjoin%
\definecolor{currentfill}{rgb}{0.749020,0.772549,0.792157}%
\pgfsetfillcolor{currentfill}%
\pgfsetfillopacity{0.160000}%
\pgfsetlinewidth{0.000000pt}%
\definecolor{currentstroke}{rgb}{0.000000,0.000000,0.000000}%
\pgfsetstrokecolor{currentstroke}%
\pgfsetdash{}{0pt}%
\pgfpathmoveto{\pgfqpoint{1.103768in}{0.633841in}}%
\pgfpathlineto{\pgfqpoint{1.107190in}{0.639732in}}%
\pgfpathlineto{\pgfqpoint{1.110627in}{0.645650in}}%
\pgfpathlineto{\pgfqpoint{1.115690in}{0.646412in}}%
\pgfpathlineto{\pgfqpoint{1.120638in}{0.647317in}}%
\pgfpathlineto{\pgfqpoint{1.115516in}{0.641121in}}%
\pgfpathlineto{\pgfqpoint{1.110416in}{0.634951in}}%
\pgfpathlineto{\pgfqpoint{1.107130in}{0.634349in}}%
\pgfpathlineto{\pgfqpoint{1.103768in}{0.633841in}}%
\pgfpathclose%
\pgfusepath{fill}%
\end{pgfscope}%
\begin{pgfscope}%
\pgfpathrectangle{\pgfqpoint{0.407592in}{0.222904in}}{\pgfqpoint{1.411163in}{1.411163in}}%
\pgfusepath{clip}%
\pgfsetbuttcap%
\pgfsetroundjoin%
\definecolor{currentfill}{rgb}{0.749020,0.772549,0.792157}%
\pgfsetfillcolor{currentfill}%
\pgfsetfillopacity{0.160000}%
\pgfsetlinewidth{0.000000pt}%
\definecolor{currentstroke}{rgb}{0.000000,0.000000,0.000000}%
\pgfsetstrokecolor{currentstroke}%
\pgfsetdash{}{0pt}%
\pgfpathmoveto{\pgfqpoint{0.986858in}{0.688474in}}%
\pgfpathlineto{\pgfqpoint{0.973807in}{0.698314in}}%
\pgfpathlineto{\pgfqpoint{0.960727in}{0.708176in}}%
\pgfpathlineto{\pgfqpoint{0.962998in}{0.704599in}}%
\pgfpathlineto{\pgfqpoint{0.965832in}{0.701101in}}%
\pgfpathlineto{\pgfqpoint{0.978420in}{0.691937in}}%
\pgfpathlineto{\pgfqpoint{0.990974in}{0.682797in}}%
\pgfpathlineto{\pgfqpoint{0.988691in}{0.685603in}}%
\pgfpathlineto{\pgfqpoint{0.986858in}{0.688474in}}%
\pgfpathclose%
\pgfusepath{fill}%
\end{pgfscope}%
\begin{pgfscope}%
\pgfpathrectangle{\pgfqpoint{0.407592in}{0.222904in}}{\pgfqpoint{1.411163in}{1.411163in}}%
\pgfusepath{clip}%
\pgfsetbuttcap%
\pgfsetroundjoin%
\definecolor{currentfill}{rgb}{0.749020,0.772549,0.792157}%
\pgfsetfillcolor{currentfill}%
\pgfsetfillopacity{0.160000}%
\pgfsetlinewidth{0.000000pt}%
\definecolor{currentstroke}{rgb}{0.000000,0.000000,0.000000}%
\pgfsetstrokecolor{currentstroke}%
\pgfsetdash{}{0pt}%
\pgfpathmoveto{\pgfqpoint{0.853529in}{0.827087in}}%
\pgfpathlineto{\pgfqpoint{0.840262in}{0.839225in}}%
\pgfpathlineto{\pgfqpoint{0.826982in}{0.851374in}}%
\pgfpathlineto{\pgfqpoint{0.824295in}{0.844172in}}%
\pgfpathlineto{\pgfqpoint{0.822757in}{0.836917in}}%
\pgfpathlineto{\pgfqpoint{0.836305in}{0.825450in}}%
\pgfpathlineto{\pgfqpoint{0.849835in}{0.813999in}}%
\pgfpathlineto{\pgfqpoint{0.851164in}{0.820565in}}%
\pgfpathlineto{\pgfqpoint{0.853529in}{0.827087in}}%
\pgfpathclose%
\pgfusepath{fill}%
\end{pgfscope}%
\begin{pgfscope}%
\pgfpathrectangle{\pgfqpoint{0.407592in}{0.222904in}}{\pgfqpoint{1.411163in}{1.411163in}}%
\pgfusepath{clip}%
\pgfsetbuttcap%
\pgfsetroundjoin%
\definecolor{currentfill}{rgb}{0.749020,0.772549,0.792157}%
\pgfsetfillcolor{currentfill}%
\pgfsetfillopacity{0.160000}%
\pgfsetlinewidth{0.000000pt}%
\definecolor{currentstroke}{rgb}{0.000000,0.000000,0.000000}%
\pgfsetstrokecolor{currentstroke}%
\pgfsetdash{}{0pt}%
\pgfpathmoveto{\pgfqpoint{1.069349in}{0.634239in}}%
\pgfpathlineto{\pgfqpoint{1.064072in}{0.640230in}}%
\pgfpathlineto{\pgfqpoint{1.058772in}{0.646248in}}%
\pgfpathlineto{\pgfqpoint{1.063704in}{0.645514in}}%
\pgfpathlineto{\pgfqpoint{1.068755in}{0.644926in}}%
\pgfpathlineto{\pgfqpoint{1.072374in}{0.639130in}}%
\pgfpathlineto{\pgfqpoint{1.075977in}{0.633359in}}%
\pgfpathlineto{\pgfqpoint{1.072624in}{0.633751in}}%
\pgfpathlineto{\pgfqpoint{1.069349in}{0.634239in}}%
\pgfpathclose%
\pgfusepath{fill}%
\end{pgfscope}%
\begin{pgfscope}%
\pgfpathrectangle{\pgfqpoint{0.407592in}{0.222904in}}{\pgfqpoint{1.411163in}{1.411163in}}%
\pgfusepath{clip}%
\pgfsetbuttcap%
\pgfsetroundjoin%
\definecolor{currentfill}{rgb}{0.749020,0.772549,0.792157}%
\pgfsetfillcolor{currentfill}%
\pgfsetfillopacity{0.160000}%
\pgfsetlinewidth{0.000000pt}%
\definecolor{currentstroke}{rgb}{0.000000,0.000000,0.000000}%
\pgfsetstrokecolor{currentstroke}%
\pgfsetdash{}{0pt}%
\pgfpathmoveto{\pgfqpoint{0.957907in}{0.715503in}}%
\pgfpathlineto{\pgfqpoint{0.944544in}{0.726104in}}%
\pgfpathlineto{\pgfqpoint{0.931156in}{0.736724in}}%
\pgfpathlineto{\pgfqpoint{0.932470in}{0.732318in}}%
\pgfpathlineto{\pgfqpoint{0.934480in}{0.727966in}}%
\pgfpathlineto{\pgfqpoint{0.947618in}{0.718060in}}%
\pgfpathlineto{\pgfqpoint{0.960727in}{0.708176in}}%
\pgfpathlineto{\pgfqpoint{0.959028in}{0.711816in}}%
\pgfpathlineto{\pgfqpoint{0.957907in}{0.715503in}}%
\pgfpathclose%
\pgfusepath{fill}%
\end{pgfscope}%
\begin{pgfscope}%
\pgfpathrectangle{\pgfqpoint{0.407592in}{0.222904in}}{\pgfqpoint{1.411163in}{1.411163in}}%
\pgfusepath{clip}%
\pgfsetbuttcap%
\pgfsetroundjoin%
\definecolor{currentfill}{rgb}{0.749020,0.772549,0.792157}%
\pgfsetfillcolor{currentfill}%
\pgfsetfillopacity{0.160000}%
\pgfsetlinewidth{0.000000pt}%
\definecolor{currentstroke}{rgb}{0.000000,0.000000,0.000000}%
\pgfsetstrokecolor{currentstroke}%
\pgfsetdash{}{0pt}%
\pgfpathmoveto{\pgfqpoint{1.020402in}{0.660554in}}%
\pgfpathlineto{\pgfqpoint{1.008641in}{0.668979in}}%
\pgfpathlineto{\pgfqpoint{0.996845in}{0.677429in}}%
\pgfpathlineto{\pgfqpoint{1.000410in}{0.674892in}}%
\pgfpathlineto{\pgfqpoint{1.004374in}{0.672468in}}%
\pgfpathlineto{\pgfqpoint{1.015233in}{0.664635in}}%
\pgfpathlineto{\pgfqpoint{1.026055in}{0.656828in}}%
\pgfpathlineto{\pgfqpoint{1.023079in}{0.658648in}}%
\pgfpathlineto{\pgfqpoint{1.020402in}{0.660554in}}%
\pgfpathclose%
\pgfusepath{fill}%
\end{pgfscope}%
\begin{pgfscope}%
\pgfpathrectangle{\pgfqpoint{0.407592in}{0.222904in}}{\pgfqpoint{1.411163in}{1.411163in}}%
\pgfusepath{clip}%
\pgfsetbuttcap%
\pgfsetroundjoin%
\definecolor{currentfill}{rgb}{0.749020,0.772549,0.792157}%
\pgfsetfillcolor{currentfill}%
\pgfsetfillopacity{0.160000}%
\pgfsetlinewidth{0.000000pt}%
\definecolor{currentstroke}{rgb}{0.000000,0.000000,0.000000}%
\pgfsetstrokecolor{currentstroke}%
\pgfsetdash{}{0pt}%
\pgfpathmoveto{\pgfqpoint{1.096876in}{0.633117in}}%
\pgfpathlineto{\pgfqpoint{1.098557in}{0.638827in}}%
\pgfpathlineto{\pgfqpoint{1.100245in}{0.644563in}}%
\pgfpathlineto{\pgfqpoint{1.105471in}{0.645032in}}%
\pgfpathlineto{\pgfqpoint{1.110627in}{0.645650in}}%
\pgfpathlineto{\pgfqpoint{1.107190in}{0.639732in}}%
\pgfpathlineto{\pgfqpoint{1.103768in}{0.633841in}}%
\pgfpathlineto{\pgfqpoint{1.100345in}{0.633430in}}%
\pgfpathlineto{\pgfqpoint{1.096876in}{0.633117in}}%
\pgfpathclose%
\pgfusepath{fill}%
\end{pgfscope}%
\begin{pgfscope}%
\pgfpathrectangle{\pgfqpoint{0.407592in}{0.222904in}}{\pgfqpoint{1.411163in}{1.411163in}}%
\pgfusepath{clip}%
\pgfsetbuttcap%
\pgfsetroundjoin%
\definecolor{currentfill}{rgb}{0.749020,0.772549,0.792157}%
\pgfsetfillcolor{currentfill}%
\pgfsetfillopacity{0.160000}%
\pgfsetlinewidth{0.000000pt}%
\definecolor{currentstroke}{rgb}{0.000000,0.000000,0.000000}%
\pgfsetstrokecolor{currentstroke}%
\pgfsetdash{}{0pt}%
\pgfpathmoveto{\pgfqpoint{0.707498in}{0.993964in}}%
\pgfpathlineto{\pgfqpoint{0.694639in}{1.006847in}}%
\pgfpathlineto{\pgfqpoint{0.681774in}{1.019736in}}%
\pgfpathlineto{\pgfqpoint{0.673337in}{1.008998in}}%
\pgfpathlineto{\pgfqpoint{0.666651in}{0.998060in}}%
\pgfpathlineto{\pgfqpoint{0.680081in}{0.985772in}}%
\pgfpathlineto{\pgfqpoint{0.693499in}{0.973497in}}%
\pgfpathlineto{\pgfqpoint{0.699676in}{0.983823in}}%
\pgfpathlineto{\pgfqpoint{0.707498in}{0.993964in}}%
\pgfpathclose%
\pgfusepath{fill}%
\end{pgfscope}%
\begin{pgfscope}%
\pgfpathrectangle{\pgfqpoint{0.407592in}{0.222904in}}{\pgfqpoint{1.411163in}{1.411163in}}%
\pgfusepath{clip}%
\pgfsetbuttcap%
\pgfsetroundjoin%
\definecolor{currentfill}{rgb}{0.749020,0.772549,0.792157}%
\pgfsetfillcolor{currentfill}%
\pgfsetfillopacity{0.160000}%
\pgfsetlinewidth{0.000000pt}%
\definecolor{currentstroke}{rgb}{0.000000,0.000000,0.000000}%
\pgfsetstrokecolor{currentstroke}%
\pgfsetdash{}{0pt}%
\pgfpathmoveto{\pgfqpoint{1.515507in}{1.000268in}}%
\pgfpathlineto{\pgfqpoint{1.528995in}{1.012628in}}%
\pgfpathlineto{\pgfqpoint{1.542494in}{1.024999in}}%
\pgfpathlineto{\pgfqpoint{1.537533in}{1.036500in}}%
\pgfpathlineto{\pgfqpoint{1.530675in}{1.047775in}}%
\pgfpathlineto{\pgfqpoint{1.517621in}{1.034815in}}%
\pgfpathlineto{\pgfqpoint{1.504572in}{1.021862in}}%
\pgfpathlineto{\pgfqpoint{1.510934in}{1.011170in}}%
\pgfpathlineto{\pgfqpoint{1.515507in}{1.000268in}}%
\pgfpathclose%
\pgfusepath{fill}%
\end{pgfscope}%
\begin{pgfscope}%
\pgfpathrectangle{\pgfqpoint{0.407592in}{0.222904in}}{\pgfqpoint{1.411163in}{1.411163in}}%
\pgfusepath{clip}%
\pgfsetbuttcap%
\pgfsetroundjoin%
\definecolor{currentfill}{rgb}{0.749020,0.772549,0.792157}%
\pgfsetfillcolor{currentfill}%
\pgfsetfillopacity{0.160000}%
\pgfsetlinewidth{0.000000pt}%
\definecolor{currentstroke}{rgb}{0.000000,0.000000,0.000000}%
\pgfsetstrokecolor{currentstroke}%
\pgfsetdash{}{0pt}%
\pgfpathmoveto{\pgfqpoint{1.075977in}{0.633359in}}%
\pgfpathlineto{\pgfqpoint{1.072374in}{0.639130in}}%
\pgfpathlineto{\pgfqpoint{1.068755in}{0.644926in}}%
\pgfpathlineto{\pgfqpoint{1.073901in}{0.644486in}}%
\pgfpathlineto{\pgfqpoint{1.079119in}{0.644197in}}%
\pgfpathlineto{\pgfqpoint{1.080992in}{0.638523in}}%
\pgfpathlineto{\pgfqpoint{1.082856in}{0.632874in}}%
\pgfpathlineto{\pgfqpoint{1.079393in}{0.633066in}}%
\pgfpathlineto{\pgfqpoint{1.075977in}{0.633359in}}%
\pgfpathclose%
\pgfusepath{fill}%
\end{pgfscope}%
\begin{pgfscope}%
\pgfpathrectangle{\pgfqpoint{0.407592in}{0.222904in}}{\pgfqpoint{1.411163in}{1.411163in}}%
\pgfusepath{clip}%
\pgfsetbuttcap%
\pgfsetroundjoin%
\definecolor{currentfill}{rgb}{0.749020,0.772549,0.792157}%
\pgfsetfillcolor{currentfill}%
\pgfsetfillopacity{0.160000}%
\pgfsetlinewidth{0.000000pt}%
\definecolor{currentstroke}{rgb}{0.000000,0.000000,0.000000}%
\pgfsetstrokecolor{currentstroke}%
\pgfsetdash{}{0pt}%
\pgfpathmoveto{\pgfqpoint{1.089864in}{0.632793in}}%
\pgfpathlineto{\pgfqpoint{1.089772in}{0.638421in}}%
\pgfpathlineto{\pgfqpoint{1.089679in}{0.644076in}}%
\pgfpathlineto{\pgfqpoint{1.094973in}{0.644243in}}%
\pgfpathlineto{\pgfqpoint{1.100245in}{0.644563in}}%
\pgfpathlineto{\pgfqpoint{1.098557in}{0.638827in}}%
\pgfpathlineto{\pgfqpoint{1.096876in}{0.633117in}}%
\pgfpathlineto{\pgfqpoint{1.093378in}{0.632904in}}%
\pgfpathlineto{\pgfqpoint{1.089864in}{0.632793in}}%
\pgfpathclose%
\pgfusepath{fill}%
\end{pgfscope}%
\begin{pgfscope}%
\pgfpathrectangle{\pgfqpoint{0.407592in}{0.222904in}}{\pgfqpoint{1.411163in}{1.411163in}}%
\pgfusepath{clip}%
\pgfsetbuttcap%
\pgfsetroundjoin%
\definecolor{currentfill}{rgb}{0.749020,0.772549,0.792157}%
\pgfsetfillcolor{currentfill}%
\pgfsetfillopacity{0.160000}%
\pgfsetlinewidth{0.000000pt}%
\definecolor{currentstroke}{rgb}{0.000000,0.000000,0.000000}%
\pgfsetstrokecolor{currentstroke}%
\pgfsetdash{}{0pt}%
\pgfpathmoveto{\pgfqpoint{1.082856in}{0.632874in}}%
\pgfpathlineto{\pgfqpoint{1.080992in}{0.638523in}}%
\pgfpathlineto{\pgfqpoint{1.079119in}{0.644197in}}%
\pgfpathlineto{\pgfqpoint{1.084387in}{0.644060in}}%
\pgfpathlineto{\pgfqpoint{1.089679in}{0.644076in}}%
\pgfpathlineto{\pgfqpoint{1.089772in}{0.638421in}}%
\pgfpathlineto{\pgfqpoint{1.089864in}{0.632793in}}%
\pgfpathlineto{\pgfqpoint{1.086352in}{0.632783in}}%
\pgfpathlineto{\pgfqpoint{1.082856in}{0.632874in}}%
\pgfpathclose%
\pgfusepath{fill}%
\end{pgfscope}%
\begin{pgfscope}%
\pgfpathrectangle{\pgfqpoint{0.407592in}{0.222904in}}{\pgfqpoint{1.411163in}{1.411163in}}%
\pgfusepath{clip}%
\pgfsetbuttcap%
\pgfsetroundjoin%
\definecolor{currentfill}{rgb}{0.749020,0.772549,0.792157}%
\pgfsetfillcolor{currentfill}%
\pgfsetfillopacity{0.160000}%
\pgfsetlinewidth{0.000000pt}%
\definecolor{currentstroke}{rgb}{0.000000,0.000000,0.000000}%
\pgfsetstrokecolor{currentstroke}%
\pgfsetdash{}{0pt}%
\pgfpathmoveto{\pgfqpoint{1.273234in}{0.759054in}}%
\pgfpathlineto{\pgfqpoint{1.286479in}{0.769803in}}%
\pgfpathlineto{\pgfqpoint{1.299747in}{0.780571in}}%
\pgfpathlineto{\pgfqpoint{1.301544in}{0.786442in}}%
\pgfpathlineto{\pgfqpoint{1.302409in}{0.792323in}}%
\pgfpathlineto{\pgfqpoint{1.289016in}{0.780848in}}%
\pgfpathlineto{\pgfqpoint{1.275641in}{0.769388in}}%
\pgfpathlineto{\pgfqpoint{1.274845in}{0.764216in}}%
\pgfpathlineto{\pgfqpoint{1.273234in}{0.759054in}}%
\pgfpathclose%
\pgfusepath{fill}%
\end{pgfscope}%
\begin{pgfscope}%
\pgfpathrectangle{\pgfqpoint{0.407592in}{0.222904in}}{\pgfqpoint{1.411163in}{1.411163in}}%
\pgfusepath{clip}%
\pgfsetbuttcap%
\pgfsetroundjoin%
\definecolor{currentfill}{rgb}{0.749020,0.772549,0.792157}%
\pgfsetfillcolor{currentfill}%
\pgfsetfillopacity{0.160000}%
\pgfsetlinewidth{0.000000pt}%
\definecolor{currentstroke}{rgb}{0.000000,0.000000,0.000000}%
\pgfsetstrokecolor{currentstroke}%
\pgfsetdash{}{0pt}%
\pgfpathmoveto{\pgfqpoint{1.146714in}{0.655440in}}%
\pgfpathlineto{\pgfqpoint{1.156244in}{0.663017in}}%
\pgfpathlineto{\pgfqpoint{1.165807in}{0.670619in}}%
\pgfpathlineto{\pgfqpoint{1.170563in}{0.672945in}}%
\pgfpathlineto{\pgfqpoint{1.174958in}{0.675392in}}%
\pgfpathlineto{\pgfqpoint{1.164250in}{0.667195in}}%
\pgfpathlineto{\pgfqpoint{1.153575in}{0.659024in}}%
\pgfpathlineto{\pgfqpoint{1.150279in}{0.657186in}}%
\pgfpathlineto{\pgfqpoint{1.146714in}{0.655440in}}%
\pgfpathclose%
\pgfusepath{fill}%
\end{pgfscope}%
\begin{pgfscope}%
\pgfpathrectangle{\pgfqpoint{0.407592in}{0.222904in}}{\pgfqpoint{1.411163in}{1.411163in}}%
\pgfusepath{clip}%
\pgfsetbuttcap%
\pgfsetroundjoin%
\definecolor{currentfill}{rgb}{0.749020,0.772549,0.792157}%
\pgfsetfillcolor{currentfill}%
\pgfsetfillopacity{0.160000}%
\pgfsetlinewidth{0.000000pt}%
\definecolor{currentstroke}{rgb}{0.000000,0.000000,0.000000}%
\pgfsetstrokecolor{currentstroke}%
\pgfsetdash{}{0pt}%
\pgfpathmoveto{\pgfqpoint{1.356160in}{0.838375in}}%
\pgfpathlineto{\pgfqpoint{1.369643in}{0.849927in}}%
\pgfpathlineto{\pgfqpoint{1.383143in}{0.861493in}}%
\pgfpathlineto{\pgfqpoint{1.382859in}{0.869419in}}%
\pgfpathlineto{\pgfqpoint{1.381300in}{0.877277in}}%
\pgfpathlineto{\pgfqpoint{1.367945in}{0.865038in}}%
\pgfpathlineto{\pgfqpoint{1.354602in}{0.852811in}}%
\pgfpathlineto{\pgfqpoint{1.355961in}{0.845622in}}%
\pgfpathlineto{\pgfqpoint{1.356160in}{0.838375in}}%
\pgfpathclose%
\pgfusepath{fill}%
\end{pgfscope}%
\begin{pgfscope}%
\pgfpathrectangle{\pgfqpoint{0.407592in}{0.222904in}}{\pgfqpoint{1.411163in}{1.411163in}}%
\pgfusepath{clip}%
\pgfsetbuttcap%
\pgfsetroundjoin%
\definecolor{currentfill}{rgb}{0.749020,0.772549,0.792157}%
\pgfsetfillcolor{currentfill}%
\pgfsetfillopacity{0.160000}%
\pgfsetlinewidth{0.000000pt}%
\definecolor{currentstroke}{rgb}{0.000000,0.000000,0.000000}%
\pgfsetstrokecolor{currentstroke}%
\pgfsetdash{}{0pt}%
\pgfpathmoveto{\pgfqpoint{0.903768in}{0.768351in}}%
\pgfpathlineto{\pgfqpoint{0.890312in}{0.779739in}}%
\pgfpathlineto{\pgfqpoint{0.876838in}{0.791143in}}%
\pgfpathlineto{\pgfqpoint{0.876634in}{0.785263in}}%
\pgfpathlineto{\pgfqpoint{0.877364in}{0.779396in}}%
\pgfpathlineto{\pgfqpoint{0.890849in}{0.768699in}}%
\pgfpathlineto{\pgfqpoint{0.904309in}{0.758021in}}%
\pgfpathlineto{\pgfqpoint{0.903630in}{0.763179in}}%
\pgfpathlineto{\pgfqpoint{0.903768in}{0.768351in}}%
\pgfpathclose%
\pgfusepath{fill}%
\end{pgfscope}%
\begin{pgfscope}%
\pgfpathrectangle{\pgfqpoint{0.407592in}{0.222904in}}{\pgfqpoint{1.411163in}{1.411163in}}%
\pgfusepath{clip}%
\pgfsetbuttcap%
\pgfsetroundjoin%
\definecolor{currentfill}{rgb}{0.749020,0.772549,0.792157}%
\pgfsetfillcolor{currentfill}%
\pgfsetfillopacity{0.160000}%
\pgfsetlinewidth{0.000000pt}%
\definecolor{currentstroke}{rgb}{0.000000,0.000000,0.000000}%
\pgfsetstrokecolor{currentstroke}%
\pgfsetdash{}{0pt}%
\pgfpathmoveto{\pgfqpoint{1.542494in}{1.024999in}}%
\pgfpathlineto{\pgfqpoint{1.556005in}{1.037381in}}%
\pgfpathlineto{\pgfqpoint{1.569529in}{1.049775in}}%
\pgfpathlineto{\pgfqpoint{1.564168in}{1.061865in}}%
\pgfpathlineto{\pgfqpoint{1.556803in}{1.073712in}}%
\pgfpathlineto{\pgfqpoint{1.543736in}{1.060741in}}%
\pgfpathlineto{\pgfqpoint{1.530675in}{1.047775in}}%
\pgfpathlineto{\pgfqpoint{1.537533in}{1.036500in}}%
\pgfpathlineto{\pgfqpoint{1.542494in}{1.024999in}}%
\pgfpathclose%
\pgfusepath{fill}%
\end{pgfscope}%
\begin{pgfscope}%
\pgfpathrectangle{\pgfqpoint{0.407592in}{0.222904in}}{\pgfqpoint{1.411163in}{1.411163in}}%
\pgfusepath{clip}%
\pgfsetbuttcap%
\pgfsetroundjoin%
\definecolor{currentfill}{rgb}{0.749020,0.772549,0.792157}%
\pgfsetfillcolor{currentfill}%
\pgfsetfillopacity{0.160000}%
\pgfsetlinewidth{0.000000pt}%
\definecolor{currentstroke}{rgb}{0.000000,0.000000,0.000000}%
\pgfsetstrokecolor{currentstroke}%
\pgfsetdash{}{0pt}%
\pgfpathmoveto{\pgfqpoint{0.681774in}{1.019736in}}%
\pgfpathlineto{\pgfqpoint{0.668902in}{1.032631in}}%
\pgfpathlineto{\pgfqpoint{0.656023in}{1.045533in}}%
\pgfpathlineto{\pgfqpoint{0.646960in}{1.034209in}}%
\pgfpathlineto{\pgfqpoint{0.639752in}{1.022669in}}%
\pgfpathlineto{\pgfqpoint{0.653208in}{1.010358in}}%
\pgfpathlineto{\pgfqpoint{0.666651in}{0.998060in}}%
\pgfpathlineto{\pgfqpoint{0.673337in}{1.008998in}}%
\pgfpathlineto{\pgfqpoint{0.681774in}{1.019736in}}%
\pgfpathclose%
\pgfusepath{fill}%
\end{pgfscope}%
\begin{pgfscope}%
\pgfpathrectangle{\pgfqpoint{0.407592in}{0.222904in}}{\pgfqpoint{1.411163in}{1.411163in}}%
\pgfusepath{clip}%
\pgfsetbuttcap%
\pgfsetroundjoin%
\definecolor{currentfill}{rgb}{0.749020,0.772549,0.792157}%
\pgfsetfillcolor{currentfill}%
\pgfsetfillopacity{0.160000}%
\pgfsetlinewidth{0.000000pt}%
\definecolor{currentstroke}{rgb}{0.000000,0.000000,0.000000}%
\pgfsetstrokecolor{currentstroke}%
\pgfsetdash{}{0pt}%
\pgfpathmoveto{\pgfqpoint{0.826982in}{0.851374in}}%
\pgfpathlineto{\pgfqpoint{0.813690in}{0.863535in}}%
\pgfpathlineto{\pgfqpoint{0.800386in}{0.875707in}}%
\pgfpathlineto{\pgfqpoint{0.797363in}{0.867833in}}%
\pgfpathlineto{\pgfqpoint{0.795604in}{0.859898in}}%
\pgfpathlineto{\pgfqpoint{0.809190in}{0.848400in}}%
\pgfpathlineto{\pgfqpoint{0.822757in}{0.836917in}}%
\pgfpathlineto{\pgfqpoint{0.824295in}{0.844172in}}%
\pgfpathlineto{\pgfqpoint{0.826982in}{0.851374in}}%
\pgfpathclose%
\pgfusepath{fill}%
\end{pgfscope}%
\begin{pgfscope}%
\pgfpathrectangle{\pgfqpoint{0.407592in}{0.222904in}}{\pgfqpoint{1.411163in}{1.411163in}}%
\pgfusepath{clip}%
\pgfsetbuttcap%
\pgfsetroundjoin%
\definecolor{currentfill}{rgb}{0.749020,0.772549,0.792157}%
\pgfsetfillcolor{currentfill}%
\pgfsetfillopacity{0.160000}%
\pgfsetlinewidth{0.000000pt}%
\definecolor{currentstroke}{rgb}{0.000000,0.000000,0.000000}%
\pgfsetstrokecolor{currentstroke}%
\pgfsetdash{}{0pt}%
\pgfpathmoveto{\pgfqpoint{1.026055in}{0.656828in}}%
\pgfpathlineto{\pgfqpoint{1.015233in}{0.664635in}}%
\pgfpathlineto{\pgfqpoint{1.004374in}{0.672468in}}%
\pgfpathlineto{\pgfqpoint{1.008723in}{0.670169in}}%
\pgfpathlineto{\pgfqpoint{1.013435in}{0.668004in}}%
\pgfpathlineto{\pgfqpoint{1.023158in}{0.660727in}}%
\pgfpathlineto{\pgfqpoint{1.032846in}{0.653477in}}%
\pgfpathlineto{\pgfqpoint{1.029315in}{0.655102in}}%
\pgfpathlineto{\pgfqpoint{1.026055in}{0.656828in}}%
\pgfpathclose%
\pgfusepath{fill}%
\end{pgfscope}%
\begin{pgfscope}%
\pgfpathrectangle{\pgfqpoint{0.407592in}{0.222904in}}{\pgfqpoint{1.411163in}{1.411163in}}%
\pgfusepath{clip}%
\pgfsetbuttcap%
\pgfsetroundjoin%
\definecolor{currentfill}{rgb}{0.749020,0.772549,0.792157}%
\pgfsetfillcolor{currentfill}%
\pgfsetfillopacity{0.160000}%
\pgfsetlinewidth{0.000000pt}%
\definecolor{currentstroke}{rgb}{0.000000,0.000000,0.000000}%
\pgfsetstrokecolor{currentstroke}%
\pgfsetdash{}{0pt}%
\pgfpathmoveto{\pgfqpoint{1.178973in}{0.677951in}}%
\pgfpathlineto{\pgfqpoint{1.190214in}{0.686491in}}%
\pgfpathlineto{\pgfqpoint{1.201490in}{0.695057in}}%
\pgfpathlineto{\pgfqpoint{1.206003in}{0.698373in}}%
\pgfpathlineto{\pgfqpoint{1.209996in}{0.701796in}}%
\pgfpathlineto{\pgfqpoint{1.197880in}{0.692563in}}%
\pgfpathlineto{\pgfqpoint{1.185795in}{0.683354in}}%
\pgfpathlineto{\pgfqpoint{1.182590in}{0.680609in}}%
\pgfpathlineto{\pgfqpoint{1.178973in}{0.677951in}}%
\pgfpathclose%
\pgfusepath{fill}%
\end{pgfscope}%
\begin{pgfscope}%
\pgfpathrectangle{\pgfqpoint{0.407592in}{0.222904in}}{\pgfqpoint{1.411163in}{1.411163in}}%
\pgfusepath{clip}%
\pgfsetbuttcap%
\pgfsetroundjoin%
\definecolor{currentfill}{rgb}{0.749020,0.772549,0.792157}%
\pgfsetfillcolor{currentfill}%
\pgfsetfillopacity{0.160000}%
\pgfsetlinewidth{0.000000pt}%
\definecolor{currentstroke}{rgb}{0.000000,0.000000,0.000000}%
\pgfsetstrokecolor{currentstroke}%
\pgfsetdash{}{0pt}%
\pgfpathmoveto{\pgfqpoint{1.569529in}{1.049775in}}%
\pgfpathlineto{\pgfqpoint{1.583065in}{1.062179in}}%
\pgfpathlineto{\pgfqpoint{1.596613in}{1.074595in}}%
\pgfpathlineto{\pgfqpoint{1.590840in}{1.087265in}}%
\pgfpathlineto{\pgfqpoint{1.582955in}{1.099674in}}%
\pgfpathlineto{\pgfqpoint{1.569876in}{1.086690in}}%
\pgfpathlineto{\pgfqpoint{1.556803in}{1.073712in}}%
\pgfpathlineto{\pgfqpoint{1.564168in}{1.061865in}}%
\pgfpathlineto{\pgfqpoint{1.569529in}{1.049775in}}%
\pgfpathclose%
\pgfusepath{fill}%
\end{pgfscope}%
\begin{pgfscope}%
\pgfpathrectangle{\pgfqpoint{0.407592in}{0.222904in}}{\pgfqpoint{1.411163in}{1.411163in}}%
\pgfusepath{clip}%
\pgfsetbuttcap%
\pgfsetroundjoin%
\definecolor{currentfill}{rgb}{0.749020,0.772549,0.792157}%
\pgfsetfillcolor{currentfill}%
\pgfsetfillopacity{0.160000}%
\pgfsetlinewidth{0.000000pt}%
\definecolor{currentstroke}{rgb}{0.000000,0.000000,0.000000}%
\pgfsetstrokecolor{currentstroke}%
\pgfsetdash{}{0pt}%
\pgfpathmoveto{\pgfqpoint{0.990974in}{0.682797in}}%
\pgfpathlineto{\pgfqpoint{0.978420in}{0.691937in}}%
\pgfpathlineto{\pgfqpoint{0.965832in}{0.701101in}}%
\pgfpathlineto{\pgfqpoint{0.969219in}{0.697698in}}%
\pgfpathlineto{\pgfqpoint{0.973145in}{0.694406in}}%
\pgfpathlineto{\pgfqpoint{0.985013in}{0.685905in}}%
\pgfpathlineto{\pgfqpoint{0.996845in}{0.677429in}}%
\pgfpathlineto{\pgfqpoint{0.993695in}{0.680068in}}%
\pgfpathlineto{\pgfqpoint{0.990974in}{0.682797in}}%
\pgfpathclose%
\pgfusepath{fill}%
\end{pgfscope}%
\begin{pgfscope}%
\pgfpathrectangle{\pgfqpoint{0.407592in}{0.222904in}}{\pgfqpoint{1.411163in}{1.411163in}}%
\pgfusepath{clip}%
\pgfsetbuttcap%
\pgfsetroundjoin%
\definecolor{currentfill}{rgb}{0.749020,0.772549,0.792157}%
\pgfsetfillcolor{currentfill}%
\pgfsetfillopacity{0.160000}%
\pgfsetlinewidth{0.000000pt}%
\definecolor{currentstroke}{rgb}{0.000000,0.000000,0.000000}%
\pgfsetstrokecolor{currentstroke}%
\pgfsetdash{}{0pt}%
\pgfpathmoveto{\pgfqpoint{0.656023in}{1.045533in}}%
\pgfpathlineto{\pgfqpoint{0.643138in}{1.058442in}}%
\pgfpathlineto{\pgfqpoint{0.630246in}{1.071357in}}%
\pgfpathlineto{\pgfqpoint{0.620544in}{1.059457in}}%
\pgfpathlineto{\pgfqpoint{0.612803in}{1.047324in}}%
\pgfpathlineto{\pgfqpoint{0.626284in}{1.034991in}}%
\pgfpathlineto{\pgfqpoint{0.639752in}{1.022669in}}%
\pgfpathlineto{\pgfqpoint{0.646960in}{1.034209in}}%
\pgfpathlineto{\pgfqpoint{0.656023in}{1.045533in}}%
\pgfpathclose%
\pgfusepath{fill}%
\end{pgfscope}%
\begin{pgfscope}%
\pgfpathrectangle{\pgfqpoint{0.407592in}{0.222904in}}{\pgfqpoint{1.411163in}{1.411163in}}%
\pgfusepath{clip}%
\pgfsetbuttcap%
\pgfsetroundjoin%
\definecolor{currentfill}{rgb}{0.749020,0.772549,0.792157}%
\pgfsetfillcolor{currentfill}%
\pgfsetfillopacity{0.160000}%
\pgfsetlinewidth{0.000000pt}%
\definecolor{currentstroke}{rgb}{0.000000,0.000000,0.000000}%
\pgfsetstrokecolor{currentstroke}%
\pgfsetdash{}{0pt}%
\pgfpathmoveto{\pgfqpoint{1.138843in}{0.652257in}}%
\pgfpathlineto{\pgfqpoint{1.147053in}{0.659304in}}%
\pgfpathlineto{\pgfqpoint{1.155295in}{0.666377in}}%
\pgfpathlineto{\pgfqpoint{1.160710in}{0.668427in}}%
\pgfpathlineto{\pgfqpoint{1.165807in}{0.670619in}}%
\pgfpathlineto{\pgfqpoint{1.156244in}{0.663017in}}%
\pgfpathlineto{\pgfqpoint{1.146714in}{0.655440in}}%
\pgfpathlineto{\pgfqpoint{1.142897in}{0.653795in}}%
\pgfpathlineto{\pgfqpoint{1.138843in}{0.652257in}}%
\pgfpathclose%
\pgfusepath{fill}%
\end{pgfscope}%
\begin{pgfscope}%
\pgfpathrectangle{\pgfqpoint{0.407592in}{0.222904in}}{\pgfqpoint{1.411163in}{1.411163in}}%
\pgfusepath{clip}%
\pgfsetbuttcap%
\pgfsetroundjoin%
\definecolor{currentfill}{rgb}{0.749020,0.772549,0.792157}%
\pgfsetfillcolor{currentfill}%
\pgfsetfillopacity{0.160000}%
\pgfsetlinewidth{0.000000pt}%
\definecolor{currentstroke}{rgb}{0.000000,0.000000,0.000000}%
\pgfsetstrokecolor{currentstroke}%
\pgfsetdash{}{0pt}%
\pgfpathmoveto{\pgfqpoint{1.383143in}{0.861493in}}%
\pgfpathlineto{\pgfqpoint{1.396662in}{0.873076in}}%
\pgfpathlineto{\pgfqpoint{1.410198in}{0.884673in}}%
\pgfpathlineto{\pgfqpoint{1.409818in}{0.893268in}}%
\pgfpathlineto{\pgfqpoint{1.408045in}{0.901787in}}%
\pgfpathlineto{\pgfqpoint{1.394666in}{0.889526in}}%
\pgfpathlineto{\pgfqpoint{1.381300in}{0.877277in}}%
\pgfpathlineto{\pgfqpoint{1.382859in}{0.869419in}}%
\pgfpathlineto{\pgfqpoint{1.383143in}{0.861493in}}%
\pgfpathclose%
\pgfusepath{fill}%
\end{pgfscope}%
\begin{pgfscope}%
\pgfpathrectangle{\pgfqpoint{0.407592in}{0.222904in}}{\pgfqpoint{1.411163in}{1.411163in}}%
\pgfusepath{clip}%
\pgfsetbuttcap%
\pgfsetroundjoin%
\definecolor{currentfill}{rgb}{0.749020,0.772549,0.792157}%
\pgfsetfillcolor{currentfill}%
\pgfsetfillopacity{0.160000}%
\pgfsetlinewidth{0.000000pt}%
\definecolor{currentstroke}{rgb}{0.000000,0.000000,0.000000}%
\pgfsetstrokecolor{currentstroke}%
\pgfsetdash{}{0pt}%
\pgfpathmoveto{\pgfqpoint{1.241920in}{0.728833in}}%
\pgfpathlineto{\pgfqpoint{1.254741in}{0.738832in}}%
\pgfpathlineto{\pgfqpoint{1.267590in}{0.748853in}}%
\pgfpathlineto{\pgfqpoint{1.270814in}{0.753925in}}%
\pgfpathlineto{\pgfqpoint{1.273234in}{0.759054in}}%
\pgfpathlineto{\pgfqpoint{1.260013in}{0.748324in}}%
\pgfpathlineto{\pgfqpoint{1.246815in}{0.737613in}}%
\pgfpathlineto{\pgfqpoint{1.244710in}{0.733198in}}%
\pgfpathlineto{\pgfqpoint{1.241920in}{0.728833in}}%
\pgfpathclose%
\pgfusepath{fill}%
\end{pgfscope}%
\begin{pgfscope}%
\pgfpathrectangle{\pgfqpoint{0.407592in}{0.222904in}}{\pgfqpoint{1.411163in}{1.411163in}}%
\pgfusepath{clip}%
\pgfsetbuttcap%
\pgfsetroundjoin%
\definecolor{currentfill}{rgb}{0.749020,0.772549,0.792157}%
\pgfsetfillcolor{currentfill}%
\pgfsetfillopacity{0.160000}%
\pgfsetlinewidth{0.000000pt}%
\definecolor{currentstroke}{rgb}{0.000000,0.000000,0.000000}%
\pgfsetstrokecolor{currentstroke}%
\pgfsetdash{}{0pt}%
\pgfpathmoveto{\pgfqpoint{0.931156in}{0.736724in}}%
\pgfpathlineto{\pgfqpoint{0.917745in}{0.747363in}}%
\pgfpathlineto{\pgfqpoint{0.904309in}{0.758021in}}%
\pgfpathlineto{\pgfqpoint{0.905805in}{0.752902in}}%
\pgfpathlineto{\pgfqpoint{0.908116in}{0.747844in}}%
\pgfpathlineto{\pgfqpoint{0.921313in}{0.737894in}}%
\pgfpathlineto{\pgfqpoint{0.934480in}{0.727966in}}%
\pgfpathlineto{\pgfqpoint{0.932470in}{0.732318in}}%
\pgfpathlineto{\pgfqpoint{0.931156in}{0.736724in}}%
\pgfpathclose%
\pgfusepath{fill}%
\end{pgfscope}%
\begin{pgfscope}%
\pgfpathrectangle{\pgfqpoint{0.407592in}{0.222904in}}{\pgfqpoint{1.411163in}{1.411163in}}%
\pgfusepath{clip}%
\pgfsetbuttcap%
\pgfsetroundjoin%
\definecolor{currentfill}{rgb}{0.749020,0.772549,0.792157}%
\pgfsetfillcolor{currentfill}%
\pgfsetfillopacity{0.160000}%
\pgfsetlinewidth{0.000000pt}%
\definecolor{currentstroke}{rgb}{0.000000,0.000000,0.000000}%
\pgfsetstrokecolor{currentstroke}%
\pgfsetdash{}{0pt}%
\pgfpathmoveto{\pgfqpoint{1.299747in}{0.780571in}}%
\pgfpathlineto{\pgfqpoint{1.313039in}{0.791358in}}%
\pgfpathlineto{\pgfqpoint{1.326354in}{0.802164in}}%
\pgfpathlineto{\pgfqpoint{1.328327in}{0.808737in}}%
\pgfpathlineto{\pgfqpoint{1.329249in}{0.815318in}}%
\pgfpathlineto{\pgfqpoint{1.315820in}{0.803813in}}%
\pgfpathlineto{\pgfqpoint{1.302409in}{0.792323in}}%
\pgfpathlineto{\pgfqpoint{1.301544in}{0.786442in}}%
\pgfpathlineto{\pgfqpoint{1.299747in}{0.780571in}}%
\pgfpathclose%
\pgfusepath{fill}%
\end{pgfscope}%
\begin{pgfscope}%
\pgfpathrectangle{\pgfqpoint{0.407592in}{0.222904in}}{\pgfqpoint{1.411163in}{1.411163in}}%
\pgfusepath{clip}%
\pgfsetbuttcap%
\pgfsetroundjoin%
\definecolor{currentfill}{rgb}{0.749020,0.772549,0.792157}%
\pgfsetfillcolor{currentfill}%
\pgfsetfillopacity{0.160000}%
\pgfsetlinewidth{0.000000pt}%
\definecolor{currentstroke}{rgb}{0.000000,0.000000,0.000000}%
\pgfsetstrokecolor{currentstroke}%
\pgfsetdash{}{0pt}%
\pgfpathmoveto{\pgfqpoint{1.209996in}{0.701796in}}%
\pgfpathlineto{\pgfqpoint{1.222144in}{0.711052in}}%
\pgfpathlineto{\pgfqpoint{1.234324in}{0.720333in}}%
\pgfpathlineto{\pgfqpoint{1.238453in}{0.724539in}}%
\pgfpathlineto{\pgfqpoint{1.241920in}{0.728833in}}%
\pgfpathlineto{\pgfqpoint{1.229127in}{0.718857in}}%
\pgfpathlineto{\pgfqpoint{1.216362in}{0.708901in}}%
\pgfpathlineto{\pgfqpoint{1.213454in}{0.705310in}}%
\pgfpathlineto{\pgfqpoint{1.209996in}{0.701796in}}%
\pgfpathclose%
\pgfusepath{fill}%
\end{pgfscope}%
\begin{pgfscope}%
\pgfpathrectangle{\pgfqpoint{0.407592in}{0.222904in}}{\pgfqpoint{1.411163in}{1.411163in}}%
\pgfusepath{clip}%
\pgfsetbuttcap%
\pgfsetroundjoin%
\definecolor{currentfill}{rgb}{0.749020,0.772549,0.792157}%
\pgfsetfillcolor{currentfill}%
\pgfsetfillopacity{0.160000}%
\pgfsetlinewidth{0.000000pt}%
\definecolor{currentstroke}{rgb}{0.000000,0.000000,0.000000}%
\pgfsetstrokecolor{currentstroke}%
\pgfsetdash{}{0pt}%
\pgfpathmoveto{\pgfqpoint{0.800386in}{0.875707in}}%
\pgfpathlineto{\pgfqpoint{0.787069in}{0.887891in}}%
\pgfpathlineto{\pgfqpoint{0.773739in}{0.900086in}}%
\pgfpathlineto{\pgfqpoint{0.770368in}{0.891549in}}%
\pgfpathlineto{\pgfqpoint{0.768377in}{0.882943in}}%
\pgfpathlineto{\pgfqpoint{0.782000in}{0.871413in}}%
\pgfpathlineto{\pgfqpoint{0.795604in}{0.859898in}}%
\pgfpathlineto{\pgfqpoint{0.797363in}{0.867833in}}%
\pgfpathlineto{\pgfqpoint{0.800386in}{0.875707in}}%
\pgfpathclose%
\pgfusepath{fill}%
\end{pgfscope}%
\begin{pgfscope}%
\pgfpathrectangle{\pgfqpoint{0.407592in}{0.222904in}}{\pgfqpoint{1.411163in}{1.411163in}}%
\pgfusepath{clip}%
\pgfsetbuttcap%
\pgfsetroundjoin%
\definecolor{currentfill}{rgb}{0.749020,0.772549,0.792157}%
\pgfsetfillcolor{currentfill}%
\pgfsetfillopacity{0.160000}%
\pgfsetlinewidth{0.000000pt}%
\definecolor{currentstroke}{rgb}{0.000000,0.000000,0.000000}%
\pgfsetstrokecolor{currentstroke}%
\pgfsetdash{}{0pt}%
\pgfpathmoveto{\pgfqpoint{1.596613in}{1.074595in}}%
\pgfpathlineto{\pgfqpoint{1.610173in}{1.087022in}}%
\pgfpathlineto{\pgfqpoint{1.623745in}{1.099460in}}%
\pgfpathlineto{\pgfqpoint{1.617549in}{1.112699in}}%
\pgfpathlineto{\pgfqpoint{1.609132in}{1.125661in}}%
\pgfpathlineto{\pgfqpoint{1.596040in}{1.112664in}}%
\pgfpathlineto{\pgfqpoint{1.582955in}{1.099674in}}%
\pgfpathlineto{\pgfqpoint{1.590840in}{1.087265in}}%
\pgfpathlineto{\pgfqpoint{1.596613in}{1.074595in}}%
\pgfpathclose%
\pgfusepath{fill}%
\end{pgfscope}%
\begin{pgfscope}%
\pgfpathrectangle{\pgfqpoint{0.407592in}{0.222904in}}{\pgfqpoint{1.411163in}{1.411163in}}%
\pgfusepath{clip}%
\pgfsetbuttcap%
\pgfsetroundjoin%
\definecolor{currentfill}{rgb}{0.749020,0.772549,0.792157}%
\pgfsetfillcolor{currentfill}%
\pgfsetfillopacity{0.160000}%
\pgfsetlinewidth{0.000000pt}%
\definecolor{currentstroke}{rgb}{0.000000,0.000000,0.000000}%
\pgfsetstrokecolor{currentstroke}%
\pgfsetdash{}{0pt}%
\pgfpathmoveto{\pgfqpoint{0.960727in}{0.708176in}}%
\pgfpathlineto{\pgfqpoint{0.947618in}{0.718060in}}%
\pgfpathlineto{\pgfqpoint{0.934480in}{0.727966in}}%
\pgfpathlineto{\pgfqpoint{0.937179in}{0.723687in}}%
\pgfpathlineto{\pgfqpoint{0.940558in}{0.719502in}}%
\pgfpathlineto{\pgfqpoint{0.953212in}{0.710289in}}%
\pgfpathlineto{\pgfqpoint{0.965832in}{0.701101in}}%
\pgfpathlineto{\pgfqpoint{0.962998in}{0.704599in}}%
\pgfpathlineto{\pgfqpoint{0.960727in}{0.708176in}}%
\pgfpathclose%
\pgfusepath{fill}%
\end{pgfscope}%
\begin{pgfscope}%
\pgfpathrectangle{\pgfqpoint{0.407592in}{0.222904in}}{\pgfqpoint{1.411163in}{1.411163in}}%
\pgfusepath{clip}%
\pgfsetbuttcap%
\pgfsetroundjoin%
\definecolor{currentfill}{rgb}{0.749020,0.772549,0.792157}%
\pgfsetfillcolor{currentfill}%
\pgfsetfillopacity{0.160000}%
\pgfsetlinewidth{0.000000pt}%
\definecolor{currentstroke}{rgb}{0.000000,0.000000,0.000000}%
\pgfsetstrokecolor{currentstroke}%
\pgfsetdash{}{0pt}%
\pgfpathmoveto{\pgfqpoint{1.032846in}{0.653477in}}%
\pgfpathlineto{\pgfqpoint{1.023158in}{0.660727in}}%
\pgfpathlineto{\pgfqpoint{1.013435in}{0.668004in}}%
\pgfpathlineto{\pgfqpoint{1.018491in}{0.665984in}}%
\pgfpathlineto{\pgfqpoint{1.023870in}{0.664119in}}%
\pgfpathlineto{\pgfqpoint{1.032280in}{0.657328in}}%
\pgfpathlineto{\pgfqpoint{1.040657in}{0.650563in}}%
\pgfpathlineto{\pgfqpoint{1.036632in}{0.651962in}}%
\pgfpathlineto{\pgfqpoint{1.032846in}{0.653477in}}%
\pgfpathclose%
\pgfusepath{fill}%
\end{pgfscope}%
\begin{pgfscope}%
\pgfpathrectangle{\pgfqpoint{0.407592in}{0.222904in}}{\pgfqpoint{1.411163in}{1.411163in}}%
\pgfusepath{clip}%
\pgfsetbuttcap%
\pgfsetroundjoin%
\definecolor{currentfill}{rgb}{0.749020,0.772549,0.792157}%
\pgfsetfillcolor{currentfill}%
\pgfsetfillopacity{0.160000}%
\pgfsetlinewidth{0.000000pt}%
\definecolor{currentstroke}{rgb}{0.000000,0.000000,0.000000}%
\pgfsetstrokecolor{currentstroke}%
\pgfsetdash{}{0pt}%
\pgfpathmoveto{\pgfqpoint{0.630246in}{1.071357in}}%
\pgfpathlineto{\pgfqpoint{0.617348in}{1.084279in}}%
\pgfpathlineto{\pgfqpoint{0.604443in}{1.097208in}}%
\pgfpathlineto{\pgfqpoint{0.594090in}{1.084742in}}%
\pgfpathlineto{\pgfqpoint{0.585803in}{1.072026in}}%
\pgfpathlineto{\pgfqpoint{0.599309in}{1.059669in}}%
\pgfpathlineto{\pgfqpoint{0.612803in}{1.047324in}}%
\pgfpathlineto{\pgfqpoint{0.620544in}{1.059457in}}%
\pgfpathlineto{\pgfqpoint{0.630246in}{1.071357in}}%
\pgfpathclose%
\pgfusepath{fill}%
\end{pgfscope}%
\begin{pgfscope}%
\pgfpathrectangle{\pgfqpoint{0.407592in}{0.222904in}}{\pgfqpoint{1.411163in}{1.411163in}}%
\pgfusepath{clip}%
\pgfsetbuttcap%
\pgfsetroundjoin%
\definecolor{currentfill}{rgb}{0.749020,0.772549,0.792157}%
\pgfsetfillcolor{currentfill}%
\pgfsetfillopacity{0.160000}%
\pgfsetlinewidth{0.000000pt}%
\definecolor{currentstroke}{rgb}{0.000000,0.000000,0.000000}%
\pgfsetstrokecolor{currentstroke}%
\pgfsetdash{}{0pt}%
\pgfpathmoveto{\pgfqpoint{0.876838in}{0.791143in}}%
\pgfpathlineto{\pgfqpoint{0.863346in}{0.802563in}}%
\pgfpathlineto{\pgfqpoint{0.849835in}{0.813999in}}%
\pgfpathlineto{\pgfqpoint{0.849551in}{0.807417in}}%
\pgfpathlineto{\pgfqpoint{0.850322in}{0.800848in}}%
\pgfpathlineto{\pgfqpoint{0.863855in}{0.790113in}}%
\pgfpathlineto{\pgfqpoint{0.877364in}{0.779396in}}%
\pgfpathlineto{\pgfqpoint{0.876634in}{0.785263in}}%
\pgfpathlineto{\pgfqpoint{0.876838in}{0.791143in}}%
\pgfpathclose%
\pgfusepath{fill}%
\end{pgfscope}%
\begin{pgfscope}%
\pgfpathrectangle{\pgfqpoint{0.407592in}{0.222904in}}{\pgfqpoint{1.411163in}{1.411163in}}%
\pgfusepath{clip}%
\pgfsetbuttcap%
\pgfsetroundjoin%
\definecolor{currentfill}{rgb}{0.749020,0.772549,0.792157}%
\pgfsetfillcolor{currentfill}%
\pgfsetfillopacity{0.160000}%
\pgfsetlinewidth{0.000000pt}%
\definecolor{currentstroke}{rgb}{0.000000,0.000000,0.000000}%
\pgfsetstrokecolor{currentstroke}%
\pgfsetdash{}{0pt}%
\pgfpathmoveto{\pgfqpoint{1.130099in}{0.649533in}}%
\pgfpathlineto{\pgfqpoint{1.136839in}{0.656125in}}%
\pgfpathlineto{\pgfqpoint{1.143606in}{0.662744in}}%
\pgfpathlineto{\pgfqpoint{1.149585in}{0.664480in}}%
\pgfpathlineto{\pgfqpoint{1.155295in}{0.666377in}}%
\pgfpathlineto{\pgfqpoint{1.147053in}{0.659304in}}%
\pgfpathlineto{\pgfqpoint{1.138843in}{0.652257in}}%
\pgfpathlineto{\pgfqpoint{1.134571in}{0.650834in}}%
\pgfpathlineto{\pgfqpoint{1.130099in}{0.649533in}}%
\pgfpathclose%
\pgfusepath{fill}%
\end{pgfscope}%
\begin{pgfscope}%
\pgfpathrectangle{\pgfqpoint{0.407592in}{0.222904in}}{\pgfqpoint{1.411163in}{1.411163in}}%
\pgfusepath{clip}%
\pgfsetbuttcap%
\pgfsetroundjoin%
\definecolor{currentfill}{rgb}{0.749020,0.772549,0.792157}%
\pgfsetfillcolor{currentfill}%
\pgfsetfillopacity{0.160000}%
\pgfsetlinewidth{0.000000pt}%
\definecolor{currentstroke}{rgb}{0.000000,0.000000,0.000000}%
\pgfsetstrokecolor{currentstroke}%
\pgfsetdash{}{0pt}%
\pgfpathmoveto{\pgfqpoint{1.174958in}{0.675392in}}%
\pgfpathlineto{\pgfqpoint{1.185700in}{0.683615in}}%
\pgfpathlineto{\pgfqpoint{1.196476in}{0.691865in}}%
\pgfpathlineto{\pgfqpoint{1.201490in}{0.695057in}}%
\pgfpathlineto{\pgfqpoint{1.190214in}{0.686491in}}%
\pgfpathlineto{\pgfqpoint{1.178973in}{0.677951in}}%
\pgfusepath{fill}%
\end{pgfscope}%
\begin{pgfscope}%
\pgfpathrectangle{\pgfqpoint{0.407592in}{0.222904in}}{\pgfqpoint{1.411163in}{1.411163in}}%
\pgfusepath{clip}%
\pgfsetbuttcap%
\pgfsetroundjoin%
\definecolor{currentfill}{rgb}{0.749020,0.772549,0.792157}%
\pgfsetfillcolor{currentfill}%
\pgfsetfillopacity{0.160000}%
\pgfsetlinewidth{0.000000pt}%
\definecolor{currentstroke}{rgb}{0.000000,0.000000,0.000000}%
\pgfsetstrokecolor{currentstroke}%
\pgfsetdash{}{0pt}%
\pgfpathmoveto{\pgfqpoint{1.410198in}{0.884673in}}%
\pgfpathlineto{\pgfqpoint{1.423753in}{0.896287in}}%
\pgfpathlineto{\pgfqpoint{1.437326in}{0.907915in}}%
\pgfpathlineto{\pgfqpoint{1.436836in}{0.917171in}}%
\pgfpathlineto{\pgfqpoint{1.434838in}{0.926341in}}%
\pgfpathlineto{\pgfqpoint{1.421436in}{0.914058in}}%
\pgfpathlineto{\pgfqpoint{1.408045in}{0.901787in}}%
\pgfpathlineto{\pgfqpoint{1.409818in}{0.893268in}}%
\pgfpathlineto{\pgfqpoint{1.410198in}{0.884673in}}%
\pgfpathclose%
\pgfusepath{fill}%
\end{pgfscope}%
\begin{pgfscope}%
\pgfpathrectangle{\pgfqpoint{0.407592in}{0.222904in}}{\pgfqpoint{1.411163in}{1.411163in}}%
\pgfusepath{clip}%
\pgfsetbuttcap%
\pgfsetroundjoin%
\definecolor{currentfill}{rgb}{0.749020,0.772549,0.792157}%
\pgfsetfillcolor{currentfill}%
\pgfsetfillopacity{0.160000}%
\pgfsetlinewidth{0.000000pt}%
\definecolor{currentstroke}{rgb}{0.000000,0.000000,0.000000}%
\pgfsetstrokecolor{currentstroke}%
\pgfsetdash{}{0pt}%
\pgfpathmoveto{\pgfqpoint{1.623745in}{1.099460in}}%
\pgfpathlineto{\pgfqpoint{1.637330in}{1.111909in}}%
\pgfpathlineto{\pgfqpoint{1.650927in}{1.124370in}}%
\pgfpathlineto{\pgfqpoint{1.644293in}{1.138168in}}%
\pgfpathlineto{\pgfqpoint{1.635333in}{1.151671in}}%
\pgfpathlineto{\pgfqpoint{1.622229in}{1.138663in}}%
\pgfpathlineto{\pgfqpoint{1.609132in}{1.125661in}}%
\pgfpathlineto{\pgfqpoint{1.617549in}{1.112699in}}%
\pgfpathlineto{\pgfqpoint{1.623745in}{1.099460in}}%
\pgfpathclose%
\pgfusepath{fill}%
\end{pgfscope}%
\begin{pgfscope}%
\pgfpathrectangle{\pgfqpoint{0.407592in}{0.222904in}}{\pgfqpoint{1.411163in}{1.411163in}}%
\pgfusepath{clip}%
\pgfsetbuttcap%
\pgfsetroundjoin%
\definecolor{currentfill}{rgb}{0.749020,0.772549,0.792157}%
\pgfsetfillcolor{currentfill}%
\pgfsetfillopacity{0.160000}%
\pgfsetlinewidth{0.000000pt}%
\definecolor{currentstroke}{rgb}{0.000000,0.000000,0.000000}%
\pgfsetstrokecolor{currentstroke}%
\pgfsetdash{}{0pt}%
\pgfpathmoveto{\pgfqpoint{0.604443in}{1.097208in}}%
\pgfpathlineto{\pgfqpoint{0.591531in}{1.110143in}}%
\pgfpathlineto{\pgfqpoint{0.578613in}{1.123085in}}%
\pgfpathlineto{\pgfqpoint{0.567598in}{1.110063in}}%
\pgfpathlineto{\pgfqpoint{0.558751in}{1.096775in}}%
\pgfpathlineto{\pgfqpoint{0.572283in}{1.084395in}}%
\pgfpathlineto{\pgfqpoint{0.585803in}{1.072026in}}%
\pgfpathlineto{\pgfqpoint{0.594090in}{1.084742in}}%
\pgfpathlineto{\pgfqpoint{0.604443in}{1.097208in}}%
\pgfpathclose%
\pgfusepath{fill}%
\end{pgfscope}%
\begin{pgfscope}%
\pgfpathrectangle{\pgfqpoint{0.407592in}{0.222904in}}{\pgfqpoint{1.411163in}{1.411163in}}%
\pgfusepath{clip}%
\pgfsetbuttcap%
\pgfsetroundjoin%
\definecolor{currentfill}{rgb}{0.749020,0.772549,0.792157}%
\pgfsetfillcolor{currentfill}%
\pgfsetfillopacity{0.160000}%
\pgfsetlinewidth{0.000000pt}%
\definecolor{currentstroke}{rgb}{0.000000,0.000000,0.000000}%
\pgfsetstrokecolor{currentstroke}%
\pgfsetdash{}{0pt}%
\pgfpathmoveto{\pgfqpoint{1.040657in}{0.650563in}}%
\pgfpathlineto{\pgfqpoint{1.032280in}{0.657328in}}%
\pgfpathlineto{\pgfqpoint{1.023870in}{0.664119in}}%
\pgfpathlineto{\pgfqpoint{1.029545in}{0.662416in}}%
\pgfpathlineto{\pgfqpoint{1.035494in}{0.660885in}}%
\pgfpathlineto{\pgfqpoint{1.042437in}{0.654498in}}%
\pgfpathlineto{\pgfqpoint{1.049350in}{0.648138in}}%
\pgfpathlineto{\pgfqpoint{1.044902in}{0.649287in}}%
\pgfpathlineto{\pgfqpoint{1.040657in}{0.650563in}}%
\pgfpathclose%
\pgfusepath{fill}%
\end{pgfscope}%
\begin{pgfscope}%
\pgfpathrectangle{\pgfqpoint{0.407592in}{0.222904in}}{\pgfqpoint{1.411163in}{1.411163in}}%
\pgfusepath{clip}%
\pgfsetbuttcap%
\pgfsetroundjoin%
\definecolor{currentfill}{rgb}{0.749020,0.772549,0.792157}%
\pgfsetfillcolor{currentfill}%
\pgfsetfillopacity{0.160000}%
\pgfsetlinewidth{0.000000pt}%
\definecolor{currentstroke}{rgb}{0.000000,0.000000,0.000000}%
\pgfsetstrokecolor{currentstroke}%
\pgfsetdash{}{0pt}%
\pgfpathmoveto{\pgfqpoint{0.996845in}{0.677429in}}%
\pgfpathlineto{\pgfqpoint{0.985013in}{0.685905in}}%
\pgfpathlineto{\pgfqpoint{0.973145in}{0.694406in}}%
\pgfpathlineto{\pgfqpoint{0.977594in}{0.691240in}}%
\pgfpathlineto{\pgfqpoint{0.982547in}{0.688214in}}%
\pgfpathlineto{\pgfqpoint{0.993479in}{0.680328in}}%
\pgfpathlineto{\pgfqpoint{1.004374in}{0.672468in}}%
\pgfpathlineto{\pgfqpoint{1.000410in}{0.674892in}}%
\pgfpathlineto{\pgfqpoint{0.996845in}{0.677429in}}%
\pgfpathclose%
\pgfusepath{fill}%
\end{pgfscope}%
\begin{pgfscope}%
\pgfpathrectangle{\pgfqpoint{0.407592in}{0.222904in}}{\pgfqpoint{1.411163in}{1.411163in}}%
\pgfusepath{clip}%
\pgfsetbuttcap%
\pgfsetroundjoin%
\definecolor{currentfill}{rgb}{0.749020,0.772549,0.792157}%
\pgfsetfillcolor{currentfill}%
\pgfsetfillopacity{0.160000}%
\pgfsetlinewidth{0.000000pt}%
\definecolor{currentstroke}{rgb}{0.000000,0.000000,0.000000}%
\pgfsetstrokecolor{currentstroke}%
\pgfsetdash{}{0pt}%
\pgfpathmoveto{\pgfqpoint{0.773739in}{0.900086in}}%
\pgfpathlineto{\pgfqpoint{0.760397in}{0.912292in}}%
\pgfpathlineto{\pgfqpoint{0.747043in}{0.924510in}}%
\pgfpathlineto{\pgfqpoint{0.743311in}{0.915320in}}%
\pgfpathlineto{\pgfqpoint{0.741074in}{0.906052in}}%
\pgfpathlineto{\pgfqpoint{0.754735in}{0.894489in}}%
\pgfpathlineto{\pgfqpoint{0.768377in}{0.882943in}}%
\pgfpathlineto{\pgfqpoint{0.770368in}{0.891549in}}%
\pgfpathlineto{\pgfqpoint{0.773739in}{0.900086in}}%
\pgfpathclose%
\pgfusepath{fill}%
\end{pgfscope}%
\begin{pgfscope}%
\pgfpathrectangle{\pgfqpoint{0.407592in}{0.222904in}}{\pgfqpoint{1.411163in}{1.411163in}}%
\pgfusepath{clip}%
\pgfsetbuttcap%
\pgfsetroundjoin%
\definecolor{currentfill}{rgb}{0.749020,0.772549,0.792157}%
\pgfsetfillcolor{currentfill}%
\pgfsetfillopacity{0.160000}%
\pgfsetlinewidth{0.000000pt}%
\definecolor{currentstroke}{rgb}{0.000000,0.000000,0.000000}%
\pgfsetstrokecolor{currentstroke}%
\pgfsetdash{}{0pt}%
\pgfpathmoveto{\pgfqpoint{1.120638in}{0.647317in}}%
\pgfpathlineto{\pgfqpoint{1.125782in}{0.653539in}}%
\pgfpathlineto{\pgfqpoint{1.130948in}{0.659788in}}%
\pgfpathlineto{\pgfqpoint{1.137384in}{0.661178in}}%
\pgfpathlineto{\pgfqpoint{1.143606in}{0.662744in}}%
\pgfpathlineto{\pgfqpoint{1.136839in}{0.656125in}}%
\pgfpathlineto{\pgfqpoint{1.130099in}{0.649533in}}%
\pgfpathlineto{\pgfqpoint{1.125448in}{0.648358in}}%
\pgfpathlineto{\pgfqpoint{1.120638in}{0.647317in}}%
\pgfpathclose%
\pgfusepath{fill}%
\end{pgfscope}%
\begin{pgfscope}%
\pgfpathrectangle{\pgfqpoint{0.407592in}{0.222904in}}{\pgfqpoint{1.411163in}{1.411163in}}%
\pgfusepath{clip}%
\pgfsetbuttcap%
\pgfsetroundjoin%
\definecolor{currentfill}{rgb}{0.749020,0.772549,0.792157}%
\pgfsetfillcolor{currentfill}%
\pgfsetfillopacity{0.160000}%
\pgfsetlinewidth{0.000000pt}%
\definecolor{currentstroke}{rgb}{0.000000,0.000000,0.000000}%
\pgfsetstrokecolor{currentstroke}%
\pgfsetdash{}{0pt}%
\pgfpathmoveto{\pgfqpoint{1.326354in}{0.802164in}}%
\pgfpathlineto{\pgfqpoint{1.339692in}{0.812989in}}%
\pgfpathlineto{\pgfqpoint{1.353055in}{0.823833in}}%
\pgfpathlineto{\pgfqpoint{1.355193in}{0.831102in}}%
\pgfpathlineto{\pgfqpoint{1.356160in}{0.838375in}}%
\pgfpathlineto{\pgfqpoint{1.342696in}{0.826839in}}%
\pgfpathlineto{\pgfqpoint{1.329249in}{0.815318in}}%
\pgfpathlineto{\pgfqpoint{1.328327in}{0.808737in}}%
\pgfpathlineto{\pgfqpoint{1.326354in}{0.802164in}}%
\pgfpathclose%
\pgfusepath{fill}%
\end{pgfscope}%
\begin{pgfscope}%
\pgfpathrectangle{\pgfqpoint{0.407592in}{0.222904in}}{\pgfqpoint{1.411163in}{1.411163in}}%
\pgfusepath{clip}%
\pgfsetbuttcap%
\pgfsetroundjoin%
\definecolor{currentfill}{rgb}{0.749020,0.772549,0.792157}%
\pgfsetfillcolor{currentfill}%
\pgfsetfillopacity{0.160000}%
\pgfsetlinewidth{0.000000pt}%
\definecolor{currentstroke}{rgb}{0.000000,0.000000,0.000000}%
\pgfsetstrokecolor{currentstroke}%
\pgfsetdash{}{0pt}%
\pgfpathmoveto{\pgfqpoint{1.650927in}{1.124370in}}%
\pgfpathlineto{\pgfqpoint{1.664536in}{1.136842in}}%
\pgfpathlineto{\pgfqpoint{1.678158in}{1.149325in}}%
\pgfpathlineto{\pgfqpoint{1.671075in}{1.163672in}}%
\pgfpathlineto{\pgfqpoint{1.661559in}{1.177707in}}%
\pgfpathlineto{\pgfqpoint{1.648443in}{1.164686in}}%
\pgfpathlineto{\pgfqpoint{1.635333in}{1.151671in}}%
\pgfpathlineto{\pgfqpoint{1.644293in}{1.138168in}}%
\pgfpathlineto{\pgfqpoint{1.650927in}{1.124370in}}%
\pgfpathclose%
\pgfusepath{fill}%
\end{pgfscope}%
\begin{pgfscope}%
\pgfpathrectangle{\pgfqpoint{0.407592in}{0.222904in}}{\pgfqpoint{1.411163in}{1.411163in}}%
\pgfusepath{clip}%
\pgfsetbuttcap%
\pgfsetroundjoin%
\definecolor{currentfill}{rgb}{0.749020,0.772549,0.792157}%
\pgfsetfillcolor{currentfill}%
\pgfsetfillopacity{0.160000}%
\pgfsetlinewidth{0.000000pt}%
\definecolor{currentstroke}{rgb}{0.000000,0.000000,0.000000}%
\pgfsetstrokecolor{currentstroke}%
\pgfsetdash{}{0pt}%
\pgfpathmoveto{\pgfqpoint{1.267590in}{0.748853in}}%
\pgfpathlineto{\pgfqpoint{1.280467in}{0.758895in}}%
\pgfpathlineto{\pgfqpoint{1.293372in}{0.768960in}}%
\pgfpathlineto{\pgfqpoint{1.297021in}{0.774735in}}%
\pgfpathlineto{\pgfqpoint{1.299747in}{0.780571in}}%
\pgfpathlineto{\pgfqpoint{1.286479in}{0.769803in}}%
\pgfpathlineto{\pgfqpoint{1.273234in}{0.759054in}}%
\pgfpathlineto{\pgfqpoint{1.270814in}{0.753925in}}%
\pgfpathlineto{\pgfqpoint{1.267590in}{0.748853in}}%
\pgfpathclose%
\pgfusepath{fill}%
\end{pgfscope}%
\begin{pgfscope}%
\pgfpathrectangle{\pgfqpoint{0.407592in}{0.222904in}}{\pgfqpoint{1.411163in}{1.411163in}}%
\pgfusepath{clip}%
\pgfsetbuttcap%
\pgfsetroundjoin%
\definecolor{currentfill}{rgb}{0.749020,0.772549,0.792157}%
\pgfsetfillcolor{currentfill}%
\pgfsetfillopacity{0.160000}%
\pgfsetlinewidth{0.000000pt}%
\definecolor{currentstroke}{rgb}{0.000000,0.000000,0.000000}%
\pgfsetstrokecolor{currentstroke}%
\pgfsetdash{}{0pt}%
\pgfpathmoveto{\pgfqpoint{0.578613in}{1.123085in}}%
\pgfpathlineto{\pgfqpoint{0.565688in}{1.136033in}}%
\pgfpathlineto{\pgfqpoint{0.552757in}{1.148988in}}%
\pgfpathlineto{\pgfqpoint{0.541067in}{1.135422in}}%
\pgfpathlineto{\pgfqpoint{0.531649in}{1.121571in}}%
\pgfpathlineto{\pgfqpoint{0.545207in}{1.109167in}}%
\pgfpathlineto{\pgfqpoint{0.558751in}{1.096775in}}%
\pgfpathlineto{\pgfqpoint{0.567598in}{1.110063in}}%
\pgfpathlineto{\pgfqpoint{0.578613in}{1.123085in}}%
\pgfpathclose%
\pgfusepath{fill}%
\end{pgfscope}%
\begin{pgfscope}%
\pgfpathrectangle{\pgfqpoint{0.407592in}{0.222904in}}{\pgfqpoint{1.411163in}{1.411163in}}%
\pgfusepath{clip}%
\pgfsetbuttcap%
\pgfsetroundjoin%
\definecolor{currentfill}{rgb}{0.749020,0.772549,0.792157}%
\pgfsetfillcolor{currentfill}%
\pgfsetfillopacity{0.160000}%
\pgfsetlinewidth{0.000000pt}%
\definecolor{currentstroke}{rgb}{0.000000,0.000000,0.000000}%
\pgfsetstrokecolor{currentstroke}%
\pgfsetdash{}{0pt}%
\pgfpathmoveto{\pgfqpoint{1.049350in}{0.648138in}}%
\pgfpathlineto{\pgfqpoint{1.042437in}{0.654498in}}%
\pgfpathlineto{\pgfqpoint{1.035494in}{0.660885in}}%
\pgfpathlineto{\pgfqpoint{1.041689in}{0.659531in}}%
\pgfpathlineto{\pgfqpoint{1.048102in}{0.658362in}}%
\pgfpathlineto{\pgfqpoint{1.053449in}{0.652292in}}%
\pgfpathlineto{\pgfqpoint{1.058772in}{0.646248in}}%
\pgfpathlineto{\pgfqpoint{1.053980in}{0.647124in}}%
\pgfpathlineto{\pgfqpoint{1.049350in}{0.648138in}}%
\pgfpathclose%
\pgfusepath{fill}%
\end{pgfscope}%
\begin{pgfscope}%
\pgfpathrectangle{\pgfqpoint{0.407592in}{0.222904in}}{\pgfqpoint{1.411163in}{1.411163in}}%
\pgfusepath{clip}%
\pgfsetbuttcap%
\pgfsetroundjoin%
\definecolor{currentfill}{rgb}{0.749020,0.772549,0.792157}%
\pgfsetfillcolor{currentfill}%
\pgfsetfillopacity{0.160000}%
\pgfsetlinewidth{0.000000pt}%
\definecolor{currentstroke}{rgb}{0.000000,0.000000,0.000000}%
\pgfsetstrokecolor{currentstroke}%
\pgfsetdash{}{0pt}%
\pgfpathmoveto{\pgfqpoint{1.165807in}{0.670619in}}%
\pgfpathlineto{\pgfqpoint{1.175404in}{0.678249in}}%
\pgfpathlineto{\pgfqpoint{1.185034in}{0.685906in}}%
\pgfpathlineto{\pgfqpoint{1.190983in}{0.688809in}}%
\pgfpathlineto{\pgfqpoint{1.196476in}{0.691865in}}%
\pgfpathlineto{\pgfqpoint{1.185700in}{0.683615in}}%
\pgfpathlineto{\pgfqpoint{1.174958in}{0.675392in}}%
\pgfpathlineto{\pgfqpoint{1.170563in}{0.672945in}}%
\pgfpathlineto{\pgfqpoint{1.165807in}{0.670619in}}%
\pgfpathclose%
\pgfusepath{fill}%
\end{pgfscope}%
\begin{pgfscope}%
\pgfpathrectangle{\pgfqpoint{0.407592in}{0.222904in}}{\pgfqpoint{1.411163in}{1.411163in}}%
\pgfusepath{clip}%
\pgfsetbuttcap%
\pgfsetroundjoin%
\definecolor{currentfill}{rgb}{0.749020,0.772549,0.792157}%
\pgfsetfillcolor{currentfill}%
\pgfsetfillopacity{0.160000}%
\pgfsetlinewidth{0.000000pt}%
\definecolor{currentstroke}{rgb}{0.000000,0.000000,0.000000}%
\pgfsetstrokecolor{currentstroke}%
\pgfsetdash{}{0pt}%
\pgfpathmoveto{\pgfqpoint{1.437326in}{0.907915in}}%
\pgfpathlineto{\pgfqpoint{1.450917in}{0.919560in}}%
\pgfpathlineto{\pgfqpoint{1.464526in}{0.931220in}}%
\pgfpathlineto{\pgfqpoint{1.463915in}{0.941128in}}%
\pgfpathlineto{\pgfqpoint{1.461680in}{0.950939in}}%
\pgfpathlineto{\pgfqpoint{1.448253in}{0.938635in}}%
\pgfpathlineto{\pgfqpoint{1.434838in}{0.926341in}}%
\pgfpathlineto{\pgfqpoint{1.436836in}{0.917171in}}%
\pgfpathlineto{\pgfqpoint{1.437326in}{0.907915in}}%
\pgfpathclose%
\pgfusepath{fill}%
\end{pgfscope}%
\begin{pgfscope}%
\pgfpathrectangle{\pgfqpoint{0.407592in}{0.222904in}}{\pgfqpoint{1.411163in}{1.411163in}}%
\pgfusepath{clip}%
\pgfsetbuttcap%
\pgfsetroundjoin%
\definecolor{currentfill}{rgb}{0.749020,0.772549,0.792157}%
\pgfsetfillcolor{currentfill}%
\pgfsetfillopacity{0.160000}%
\pgfsetlinewidth{0.000000pt}%
\definecolor{currentstroke}{rgb}{0.000000,0.000000,0.000000}%
\pgfsetstrokecolor{currentstroke}%
\pgfsetdash{}{0pt}%
\pgfpathmoveto{\pgfqpoint{0.849835in}{0.813999in}}%
\pgfpathlineto{\pgfqpoint{0.836305in}{0.825450in}}%
\pgfpathlineto{\pgfqpoint{0.822757in}{0.836917in}}%
\pgfpathlineto{\pgfqpoint{0.822383in}{0.829642in}}%
\pgfpathlineto{\pgfqpoint{0.823182in}{0.822378in}}%
\pgfpathlineto{\pgfqpoint{0.836764in}{0.811604in}}%
\pgfpathlineto{\pgfqpoint{0.850322in}{0.800848in}}%
\pgfpathlineto{\pgfqpoint{0.849551in}{0.807417in}}%
\pgfpathlineto{\pgfqpoint{0.849835in}{0.813999in}}%
\pgfpathclose%
\pgfusepath{fill}%
\end{pgfscope}%
\begin{pgfscope}%
\pgfpathrectangle{\pgfqpoint{0.407592in}{0.222904in}}{\pgfqpoint{1.411163in}{1.411163in}}%
\pgfusepath{clip}%
\pgfsetbuttcap%
\pgfsetroundjoin%
\definecolor{currentfill}{rgb}{0.749020,0.772549,0.792157}%
\pgfsetfillcolor{currentfill}%
\pgfsetfillopacity{0.160000}%
\pgfsetlinewidth{0.000000pt}%
\definecolor{currentstroke}{rgb}{0.000000,0.000000,0.000000}%
\pgfsetstrokecolor{currentstroke}%
\pgfsetdash{}{0pt}%
\pgfpathmoveto{\pgfqpoint{0.904309in}{0.758021in}}%
\pgfpathlineto{\pgfqpoint{0.890849in}{0.768699in}}%
\pgfpathlineto{\pgfqpoint{0.877364in}{0.779396in}}%
\pgfpathlineto{\pgfqpoint{0.879032in}{0.773570in}}%
\pgfpathlineto{\pgfqpoint{0.881633in}{0.767811in}}%
\pgfpathlineto{\pgfqpoint{0.894889in}{0.757816in}}%
\pgfpathlineto{\pgfqpoint{0.908116in}{0.747844in}}%
\pgfpathlineto{\pgfqpoint{0.905805in}{0.752902in}}%
\pgfpathlineto{\pgfqpoint{0.904309in}{0.758021in}}%
\pgfpathclose%
\pgfusepath{fill}%
\end{pgfscope}%
\begin{pgfscope}%
\pgfpathrectangle{\pgfqpoint{0.407592in}{0.222904in}}{\pgfqpoint{1.411163in}{1.411163in}}%
\pgfusepath{clip}%
\pgfsetbuttcap%
\pgfsetroundjoin%
\definecolor{currentfill}{rgb}{0.749020,0.772549,0.792157}%
\pgfsetfillcolor{currentfill}%
\pgfsetfillopacity{0.160000}%
\pgfsetlinewidth{0.000000pt}%
\definecolor{currentstroke}{rgb}{0.000000,0.000000,0.000000}%
\pgfsetstrokecolor{currentstroke}%
\pgfsetdash{}{0pt}%
\pgfpathmoveto{\pgfqpoint{1.110627in}{0.645650in}}%
\pgfpathlineto{\pgfqpoint{1.114079in}{0.651593in}}%
\pgfpathlineto{\pgfqpoint{1.117547in}{0.657564in}}%
\pgfpathlineto{\pgfqpoint{1.124325in}{0.658581in}}%
\pgfpathlineto{\pgfqpoint{1.130948in}{0.659788in}}%
\pgfpathlineto{\pgfqpoint{1.125782in}{0.653539in}}%
\pgfpathlineto{\pgfqpoint{1.120638in}{0.647317in}}%
\pgfpathlineto{\pgfqpoint{1.115690in}{0.646412in}}%
\pgfpathlineto{\pgfqpoint{1.110627in}{0.645650in}}%
\pgfpathclose%
\pgfusepath{fill}%
\end{pgfscope}%
\begin{pgfscope}%
\pgfpathrectangle{\pgfqpoint{0.407592in}{0.222904in}}{\pgfqpoint{1.411163in}{1.411163in}}%
\pgfusepath{clip}%
\pgfsetbuttcap%
\pgfsetroundjoin%
\definecolor{currentfill}{rgb}{0.749020,0.772549,0.792157}%
\pgfsetfillcolor{currentfill}%
\pgfsetfillopacity{0.160000}%
\pgfsetlinewidth{0.000000pt}%
\definecolor{currentstroke}{rgb}{0.000000,0.000000,0.000000}%
\pgfsetstrokecolor{currentstroke}%
\pgfsetdash{}{0pt}%
\pgfpathmoveto{\pgfqpoint{1.201490in}{0.695057in}}%
\pgfpathlineto{\pgfqpoint{1.212799in}{0.703649in}}%
\pgfpathlineto{\pgfqpoint{1.224142in}{0.712266in}}%
\pgfpathlineto{\pgfqpoint{1.229547in}{0.716236in}}%
\pgfpathlineto{\pgfqpoint{1.234324in}{0.720333in}}%
\pgfpathlineto{\pgfqpoint{1.222144in}{0.711052in}}%
\pgfpathlineto{\pgfqpoint{1.209996in}{0.701796in}}%
\pgfpathlineto{\pgfqpoint{1.206003in}{0.698373in}}%
\pgfpathlineto{\pgfqpoint{1.201490in}{0.695057in}}%
\pgfpathclose%
\pgfusepath{fill}%
\end{pgfscope}%
\begin{pgfscope}%
\pgfpathrectangle{\pgfqpoint{0.407592in}{0.222904in}}{\pgfqpoint{1.411163in}{1.411163in}}%
\pgfusepath{clip}%
\pgfsetbuttcap%
\pgfsetroundjoin%
\definecolor{currentfill}{rgb}{0.749020,0.772549,0.792157}%
\pgfsetfillcolor{currentfill}%
\pgfsetfillopacity{0.160000}%
\pgfsetlinewidth{0.000000pt}%
\definecolor{currentstroke}{rgb}{0.000000,0.000000,0.000000}%
\pgfsetstrokecolor{currentstroke}%
\pgfsetdash{}{0pt}%
\pgfpathmoveto{\pgfqpoint{1.678158in}{1.149325in}}%
\pgfpathlineto{\pgfqpoint{1.691792in}{1.161819in}}%
\pgfpathlineto{\pgfqpoint{1.705438in}{1.174325in}}%
\pgfpathlineto{\pgfqpoint{1.697893in}{1.189211in}}%
\pgfpathlineto{\pgfqpoint{1.687810in}{1.203766in}}%
\pgfpathlineto{\pgfqpoint{1.674681in}{1.190733in}}%
\pgfpathlineto{\pgfqpoint{1.661559in}{1.177707in}}%
\pgfpathlineto{\pgfqpoint{1.671075in}{1.163672in}}%
\pgfpathlineto{\pgfqpoint{1.678158in}{1.149325in}}%
\pgfpathclose%
\pgfusepath{fill}%
\end{pgfscope}%
\begin{pgfscope}%
\pgfpathrectangle{\pgfqpoint{0.407592in}{0.222904in}}{\pgfqpoint{1.411163in}{1.411163in}}%
\pgfusepath{clip}%
\pgfsetbuttcap%
\pgfsetroundjoin%
\definecolor{currentfill}{rgb}{0.749020,0.772549,0.792157}%
\pgfsetfillcolor{currentfill}%
\pgfsetfillopacity{0.160000}%
\pgfsetlinewidth{0.000000pt}%
\definecolor{currentstroke}{rgb}{0.000000,0.000000,0.000000}%
\pgfsetstrokecolor{currentstroke}%
\pgfsetdash{}{0pt}%
\pgfpathmoveto{\pgfqpoint{1.058772in}{0.646248in}}%
\pgfpathlineto{\pgfqpoint{1.053449in}{0.652292in}}%
\pgfpathlineto{\pgfqpoint{1.048102in}{0.658362in}}%
\pgfpathlineto{\pgfqpoint{1.054705in}{0.657383in}}%
\pgfpathlineto{\pgfqpoint{1.061467in}{0.656598in}}%
\pgfpathlineto{\pgfqpoint{1.065119in}{0.650749in}}%
\pgfpathlineto{\pgfqpoint{1.068755in}{0.644926in}}%
\pgfpathlineto{\pgfqpoint{1.063704in}{0.645514in}}%
\pgfpathlineto{\pgfqpoint{1.058772in}{0.646248in}}%
\pgfpathclose%
\pgfusepath{fill}%
\end{pgfscope}%
\begin{pgfscope}%
\pgfpathrectangle{\pgfqpoint{0.407592in}{0.222904in}}{\pgfqpoint{1.411163in}{1.411163in}}%
\pgfusepath{clip}%
\pgfsetbuttcap%
\pgfsetroundjoin%
\definecolor{currentfill}{rgb}{0.749020,0.772549,0.792157}%
\pgfsetfillcolor{currentfill}%
\pgfsetfillopacity{0.160000}%
\pgfsetlinewidth{0.000000pt}%
\definecolor{currentstroke}{rgb}{0.000000,0.000000,0.000000}%
\pgfsetstrokecolor{currentstroke}%
\pgfsetdash{}{0pt}%
\pgfpathmoveto{\pgfqpoint{1.234324in}{0.720333in}}%
\pgfpathlineto{\pgfqpoint{1.246535in}{0.729638in}}%
\pgfpathlineto{\pgfqpoint{1.258778in}{0.738968in}}%
\pgfpathlineto{\pgfqpoint{1.263573in}{0.743859in}}%
\pgfpathlineto{\pgfqpoint{1.267590in}{0.748853in}}%
\pgfpathlineto{\pgfqpoint{1.254741in}{0.738832in}}%
\pgfpathlineto{\pgfqpoint{1.241920in}{0.728833in}}%
\pgfpathlineto{\pgfqpoint{1.238453in}{0.724539in}}%
\pgfpathlineto{\pgfqpoint{1.234324in}{0.720333in}}%
\pgfpathclose%
\pgfusepath{fill}%
\end{pgfscope}%
\begin{pgfscope}%
\pgfpathrectangle{\pgfqpoint{0.407592in}{0.222904in}}{\pgfqpoint{1.411163in}{1.411163in}}%
\pgfusepath{clip}%
\pgfsetbuttcap%
\pgfsetroundjoin%
\definecolor{currentfill}{rgb}{0.749020,0.772549,0.792157}%
\pgfsetfillcolor{currentfill}%
\pgfsetfillopacity{0.160000}%
\pgfsetlinewidth{0.000000pt}%
\definecolor{currentstroke}{rgb}{0.000000,0.000000,0.000000}%
\pgfsetstrokecolor{currentstroke}%
\pgfsetdash{}{0pt}%
\pgfpathmoveto{\pgfqpoint{0.965832in}{0.701101in}}%
\pgfpathlineto{\pgfqpoint{0.953212in}{0.710289in}}%
\pgfpathlineto{\pgfqpoint{0.940558in}{0.719502in}}%
\pgfpathlineto{\pgfqpoint{0.944605in}{0.715429in}}%
\pgfpathlineto{\pgfqpoint{0.949302in}{0.711487in}}%
\pgfpathlineto{\pgfqpoint{0.961242in}{0.702934in}}%
\pgfpathlineto{\pgfqpoint{0.973145in}{0.694406in}}%
\pgfpathlineto{\pgfqpoint{0.969219in}{0.697698in}}%
\pgfpathlineto{\pgfqpoint{0.965832in}{0.701101in}}%
\pgfpathclose%
\pgfusepath{fill}%
\end{pgfscope}%
\begin{pgfscope}%
\pgfpathrectangle{\pgfqpoint{0.407592in}{0.222904in}}{\pgfqpoint{1.411163in}{1.411163in}}%
\pgfusepath{clip}%
\pgfsetbuttcap%
\pgfsetroundjoin%
\definecolor{currentfill}{rgb}{0.749020,0.772549,0.792157}%
\pgfsetfillcolor{currentfill}%
\pgfsetfillopacity{0.160000}%
\pgfsetlinewidth{0.000000pt}%
\definecolor{currentstroke}{rgb}{0.000000,0.000000,0.000000}%
\pgfsetstrokecolor{currentstroke}%
\pgfsetdash{}{0pt}%
\pgfpathmoveto{\pgfqpoint{0.747043in}{0.924510in}}%
\pgfpathlineto{\pgfqpoint{0.733675in}{0.936740in}}%
\pgfpathlineto{\pgfqpoint{0.720296in}{0.948980in}}%
\pgfpathlineto{\pgfqpoint{0.716191in}{0.939146in}}%
\pgfpathlineto{\pgfqpoint{0.713696in}{0.929224in}}%
\pgfpathlineto{\pgfqpoint{0.727395in}{0.917630in}}%
\pgfpathlineto{\pgfqpoint{0.741074in}{0.906052in}}%
\pgfpathlineto{\pgfqpoint{0.743311in}{0.915320in}}%
\pgfpathlineto{\pgfqpoint{0.747043in}{0.924510in}}%
\pgfpathclose%
\pgfusepath{fill}%
\end{pgfscope}%
\begin{pgfscope}%
\pgfpathrectangle{\pgfqpoint{0.407592in}{0.222904in}}{\pgfqpoint{1.411163in}{1.411163in}}%
\pgfusepath{clip}%
\pgfsetbuttcap%
\pgfsetroundjoin%
\definecolor{currentfill}{rgb}{0.749020,0.772549,0.792157}%
\pgfsetfillcolor{currentfill}%
\pgfsetfillopacity{0.160000}%
\pgfsetlinewidth{0.000000pt}%
\definecolor{currentstroke}{rgb}{0.000000,0.000000,0.000000}%
\pgfsetstrokecolor{currentstroke}%
\pgfsetdash{}{0pt}%
\pgfpathmoveto{\pgfqpoint{0.552757in}{1.148988in}}%
\pgfpathlineto{\pgfqpoint{0.539819in}{1.161950in}}%
\pgfpathlineto{\pgfqpoint{0.526874in}{1.174919in}}%
\pgfpathlineto{\pgfqpoint{0.514497in}{1.160817in}}%
\pgfpathlineto{\pgfqpoint{0.504495in}{1.146414in}}%
\pgfpathlineto{\pgfqpoint{0.518078in}{1.133987in}}%
\pgfpathlineto{\pgfqpoint{0.531649in}{1.121571in}}%
\pgfpathlineto{\pgfqpoint{0.541067in}{1.135422in}}%
\pgfpathlineto{\pgfqpoint{0.552757in}{1.148988in}}%
\pgfpathclose%
\pgfusepath{fill}%
\end{pgfscope}%
\begin{pgfscope}%
\pgfpathrectangle{\pgfqpoint{0.407592in}{0.222904in}}{\pgfqpoint{1.411163in}{1.411163in}}%
\pgfusepath{clip}%
\pgfsetbuttcap%
\pgfsetroundjoin%
\definecolor{currentfill}{rgb}{0.749020,0.772549,0.792157}%
\pgfsetfillcolor{currentfill}%
\pgfsetfillopacity{0.160000}%
\pgfsetlinewidth{0.000000pt}%
\definecolor{currentstroke}{rgb}{0.000000,0.000000,0.000000}%
\pgfsetstrokecolor{currentstroke}%
\pgfsetdash{}{0pt}%
\pgfpathmoveto{\pgfqpoint{1.100245in}{0.644563in}}%
\pgfpathlineto{\pgfqpoint{1.101941in}{0.650324in}}%
\pgfpathlineto{\pgfqpoint{1.103644in}{0.656113in}}%
\pgfpathlineto{\pgfqpoint{1.110643in}{0.656739in}}%
\pgfpathlineto{\pgfqpoint{1.117547in}{0.657564in}}%
\pgfpathlineto{\pgfqpoint{1.114079in}{0.651593in}}%
\pgfpathlineto{\pgfqpoint{1.110627in}{0.645650in}}%
\pgfpathlineto{\pgfqpoint{1.105471in}{0.645032in}}%
\pgfpathlineto{\pgfqpoint{1.100245in}{0.644563in}}%
\pgfpathclose%
\pgfusepath{fill}%
\end{pgfscope}%
\begin{pgfscope}%
\pgfpathrectangle{\pgfqpoint{0.407592in}{0.222904in}}{\pgfqpoint{1.411163in}{1.411163in}}%
\pgfusepath{clip}%
\pgfsetbuttcap%
\pgfsetroundjoin%
\definecolor{currentfill}{rgb}{0.749020,0.772549,0.792157}%
\pgfsetfillcolor{currentfill}%
\pgfsetfillopacity{0.160000}%
\pgfsetlinewidth{0.000000pt}%
\definecolor{currentstroke}{rgb}{0.000000,0.000000,0.000000}%
\pgfsetstrokecolor{currentstroke}%
\pgfsetdash{}{0pt}%
\pgfpathmoveto{\pgfqpoint{0.934480in}{0.727966in}}%
\pgfpathlineto{\pgfqpoint{0.921313in}{0.737894in}}%
\pgfpathlineto{\pgfqpoint{0.908116in}{0.747844in}}%
\pgfpathlineto{\pgfqpoint{0.911234in}{0.742869in}}%
\pgfpathlineto{\pgfqpoint{0.915150in}{0.738000in}}%
\pgfpathlineto{\pgfqpoint{0.927871in}{0.728739in}}%
\pgfpathlineto{\pgfqpoint{0.940558in}{0.719502in}}%
\pgfpathlineto{\pgfqpoint{0.937179in}{0.723687in}}%
\pgfpathlineto{\pgfqpoint{0.934480in}{0.727966in}}%
\pgfpathclose%
\pgfusepath{fill}%
\end{pgfscope}%
\begin{pgfscope}%
\pgfpathrectangle{\pgfqpoint{0.407592in}{0.222904in}}{\pgfqpoint{1.411163in}{1.411163in}}%
\pgfusepath{clip}%
\pgfsetbuttcap%
\pgfsetroundjoin%
\definecolor{currentfill}{rgb}{0.749020,0.772549,0.792157}%
\pgfsetfillcolor{currentfill}%
\pgfsetfillopacity{0.160000}%
\pgfsetlinewidth{0.000000pt}%
\definecolor{currentstroke}{rgb}{0.000000,0.000000,0.000000}%
\pgfsetstrokecolor{currentstroke}%
\pgfsetdash{}{0pt}%
\pgfpathmoveto{\pgfqpoint{1.068755in}{0.644926in}}%
\pgfpathlineto{\pgfqpoint{1.065119in}{0.650749in}}%
\pgfpathlineto{\pgfqpoint{1.061467in}{0.656598in}}%
\pgfpathlineto{\pgfqpoint{1.068358in}{0.656011in}}%
\pgfpathlineto{\pgfqpoint{1.075348in}{0.655625in}}%
\pgfpathlineto{\pgfqpoint{1.077238in}{0.649898in}}%
\pgfpathlineto{\pgfqpoint{1.079119in}{0.644197in}}%
\pgfpathlineto{\pgfqpoint{1.073901in}{0.644486in}}%
\pgfpathlineto{\pgfqpoint{1.068755in}{0.644926in}}%
\pgfpathclose%
\pgfusepath{fill}%
\end{pgfscope}%
\begin{pgfscope}%
\pgfpathrectangle{\pgfqpoint{0.407592in}{0.222904in}}{\pgfqpoint{1.411163in}{1.411163in}}%
\pgfusepath{clip}%
\pgfsetbuttcap%
\pgfsetroundjoin%
\definecolor{currentfill}{rgb}{0.749020,0.772549,0.792157}%
\pgfsetfillcolor{currentfill}%
\pgfsetfillopacity{0.160000}%
\pgfsetlinewidth{0.000000pt}%
\definecolor{currentstroke}{rgb}{0.000000,0.000000,0.000000}%
\pgfsetstrokecolor{currentstroke}%
\pgfsetdash{}{0pt}%
\pgfpathmoveto{\pgfqpoint{1.004374in}{0.672468in}}%
\pgfpathlineto{\pgfqpoint{0.993479in}{0.680328in}}%
\pgfpathlineto{\pgfqpoint{0.982547in}{0.688214in}}%
\pgfpathlineto{\pgfqpoint{0.987984in}{0.685342in}}%
\pgfpathlineto{\pgfqpoint{0.993880in}{0.682638in}}%
\pgfpathlineto{\pgfqpoint{1.003676in}{0.675308in}}%
\pgfpathlineto{\pgfqpoint{1.013435in}{0.668004in}}%
\pgfpathlineto{\pgfqpoint{1.008723in}{0.670169in}}%
\pgfpathlineto{\pgfqpoint{1.004374in}{0.672468in}}%
\pgfpathclose%
\pgfusepath{fill}%
\end{pgfscope}%
\begin{pgfscope}%
\pgfpathrectangle{\pgfqpoint{0.407592in}{0.222904in}}{\pgfqpoint{1.411163in}{1.411163in}}%
\pgfusepath{clip}%
\pgfsetbuttcap%
\pgfsetroundjoin%
\definecolor{currentfill}{rgb}{0.749020,0.772549,0.792157}%
\pgfsetfillcolor{currentfill}%
\pgfsetfillopacity{0.160000}%
\pgfsetlinewidth{0.000000pt}%
\definecolor{currentstroke}{rgb}{0.000000,0.000000,0.000000}%
\pgfsetstrokecolor{currentstroke}%
\pgfsetdash{}{0pt}%
\pgfpathmoveto{\pgfqpoint{1.089679in}{0.644076in}}%
\pgfpathlineto{\pgfqpoint{1.089586in}{0.649756in}}%
\pgfpathlineto{\pgfqpoint{1.089492in}{0.655462in}}%
\pgfpathlineto{\pgfqpoint{1.096584in}{0.655686in}}%
\pgfpathlineto{\pgfqpoint{1.103644in}{0.656113in}}%
\pgfpathlineto{\pgfqpoint{1.101941in}{0.650324in}}%
\pgfpathlineto{\pgfqpoint{1.100245in}{0.644563in}}%
\pgfpathlineto{\pgfqpoint{1.094973in}{0.644243in}}%
\pgfpathlineto{\pgfqpoint{1.089679in}{0.644076in}}%
\pgfpathclose%
\pgfusepath{fill}%
\end{pgfscope}%
\begin{pgfscope}%
\pgfpathrectangle{\pgfqpoint{0.407592in}{0.222904in}}{\pgfqpoint{1.411163in}{1.411163in}}%
\pgfusepath{clip}%
\pgfsetbuttcap%
\pgfsetroundjoin%
\definecolor{currentfill}{rgb}{0.749020,0.772549,0.792157}%
\pgfsetfillcolor{currentfill}%
\pgfsetfillopacity{0.160000}%
\pgfsetlinewidth{0.000000pt}%
\definecolor{currentstroke}{rgb}{0.000000,0.000000,0.000000}%
\pgfsetstrokecolor{currentstroke}%
\pgfsetdash{}{0pt}%
\pgfpathmoveto{\pgfqpoint{1.079119in}{0.644197in}}%
\pgfpathlineto{\pgfqpoint{1.077238in}{0.649898in}}%
\pgfpathlineto{\pgfqpoint{1.075348in}{0.655625in}}%
\pgfpathlineto{\pgfqpoint{1.082403in}{0.655442in}}%
\pgfpathlineto{\pgfqpoint{1.089492in}{0.655462in}}%
\pgfpathlineto{\pgfqpoint{1.089586in}{0.649756in}}%
\pgfpathlineto{\pgfqpoint{1.089679in}{0.644076in}}%
\pgfpathlineto{\pgfqpoint{1.084387in}{0.644060in}}%
\pgfpathlineto{\pgfqpoint{1.079119in}{0.644197in}}%
\pgfpathclose%
\pgfusepath{fill}%
\end{pgfscope}%
\begin{pgfscope}%
\pgfpathrectangle{\pgfqpoint{0.407592in}{0.222904in}}{\pgfqpoint{1.411163in}{1.411163in}}%
\pgfusepath{clip}%
\pgfsetbuttcap%
\pgfsetroundjoin%
\definecolor{currentfill}{rgb}{0.749020,0.772549,0.792157}%
\pgfsetfillcolor{currentfill}%
\pgfsetfillopacity{0.160000}%
\pgfsetlinewidth{0.000000pt}%
\definecolor{currentstroke}{rgb}{0.000000,0.000000,0.000000}%
\pgfsetstrokecolor{currentstroke}%
\pgfsetdash{}{0pt}%
\pgfpathmoveto{\pgfqpoint{1.705438in}{1.174325in}}%
\pgfpathlineto{\pgfqpoint{1.719097in}{1.186842in}}%
\pgfpathlineto{\pgfqpoint{1.732768in}{1.199371in}}%
\pgfpathlineto{\pgfqpoint{1.724748in}{1.214785in}}%
\pgfpathlineto{\pgfqpoint{1.714085in}{1.229851in}}%
\pgfpathlineto{\pgfqpoint{1.700944in}{1.216805in}}%
\pgfpathlineto{\pgfqpoint{1.687810in}{1.203766in}}%
\pgfpathlineto{\pgfqpoint{1.697893in}{1.189211in}}%
\pgfpathlineto{\pgfqpoint{1.705438in}{1.174325in}}%
\pgfpathclose%
\pgfusepath{fill}%
\end{pgfscope}%
\begin{pgfscope}%
\pgfpathrectangle{\pgfqpoint{0.407592in}{0.222904in}}{\pgfqpoint{1.411163in}{1.411163in}}%
\pgfusepath{clip}%
\pgfsetbuttcap%
\pgfsetroundjoin%
\definecolor{currentfill}{rgb}{0.749020,0.772549,0.792157}%
\pgfsetfillcolor{currentfill}%
\pgfsetfillopacity{0.160000}%
\pgfsetlinewidth{0.000000pt}%
\definecolor{currentstroke}{rgb}{0.000000,0.000000,0.000000}%
\pgfsetstrokecolor{currentstroke}%
\pgfsetdash{}{0pt}%
\pgfpathmoveto{\pgfqpoint{1.464526in}{0.931220in}}%
\pgfpathlineto{\pgfqpoint{1.478153in}{0.942895in}}%
\pgfpathlineto{\pgfqpoint{1.491799in}{0.954586in}}%
\pgfpathlineto{\pgfqpoint{1.491055in}{0.965138in}}%
\pgfpathlineto{\pgfqpoint{1.488569in}{0.975581in}}%
\pgfpathlineto{\pgfqpoint{1.475119in}{0.963255in}}%
\pgfpathlineto{\pgfqpoint{1.461680in}{0.950939in}}%
\pgfpathlineto{\pgfqpoint{1.463915in}{0.941128in}}%
\pgfpathlineto{\pgfqpoint{1.464526in}{0.931220in}}%
\pgfpathclose%
\pgfusepath{fill}%
\end{pgfscope}%
\begin{pgfscope}%
\pgfpathrectangle{\pgfqpoint{0.407592in}{0.222904in}}{\pgfqpoint{1.411163in}{1.411163in}}%
\pgfusepath{clip}%
\pgfsetbuttcap%
\pgfsetroundjoin%
\definecolor{currentfill}{rgb}{0.749020,0.772549,0.792157}%
\pgfsetfillcolor{currentfill}%
\pgfsetfillopacity{0.160000}%
\pgfsetlinewidth{0.000000pt}%
\definecolor{currentstroke}{rgb}{0.000000,0.000000,0.000000}%
\pgfsetstrokecolor{currentstroke}%
\pgfsetdash{}{0pt}%
\pgfpathmoveto{\pgfqpoint{1.353055in}{0.823833in}}%
\pgfpathlineto{\pgfqpoint{1.366440in}{0.834697in}}%
\pgfpathlineto{\pgfqpoint{1.379850in}{0.845580in}}%
\pgfpathlineto{\pgfqpoint{1.382142in}{0.853535in}}%
\pgfpathlineto{\pgfqpoint{1.383143in}{0.861493in}}%
\pgfpathlineto{\pgfqpoint{1.369643in}{0.849927in}}%
\pgfpathlineto{\pgfqpoint{1.356160in}{0.838375in}}%
\pgfpathlineto{\pgfqpoint{1.355193in}{0.831102in}}%
\pgfpathlineto{\pgfqpoint{1.353055in}{0.823833in}}%
\pgfpathclose%
\pgfusepath{fill}%
\end{pgfscope}%
\begin{pgfscope}%
\pgfpathrectangle{\pgfqpoint{0.407592in}{0.222904in}}{\pgfqpoint{1.411163in}{1.411163in}}%
\pgfusepath{clip}%
\pgfsetbuttcap%
\pgfsetroundjoin%
\definecolor{currentfill}{rgb}{0.749020,0.772549,0.792157}%
\pgfsetfillcolor{currentfill}%
\pgfsetfillopacity{0.160000}%
\pgfsetlinewidth{0.000000pt}%
\definecolor{currentstroke}{rgb}{0.000000,0.000000,0.000000}%
\pgfsetstrokecolor{currentstroke}%
\pgfsetdash{}{0pt}%
\pgfpathmoveto{\pgfqpoint{0.526874in}{1.174919in}}%
\pgfpathlineto{\pgfqpoint{0.513923in}{1.187894in}}%
\pgfpathlineto{\pgfqpoint{0.500965in}{1.200875in}}%
\pgfpathlineto{\pgfqpoint{0.487888in}{1.186249in}}%
\pgfpathlineto{\pgfqpoint{0.477289in}{1.171304in}}%
\pgfpathlineto{\pgfqpoint{0.490898in}{1.158853in}}%
\pgfpathlineto{\pgfqpoint{0.504495in}{1.146414in}}%
\pgfpathlineto{\pgfqpoint{0.514497in}{1.160817in}}%
\pgfpathlineto{\pgfqpoint{0.526874in}{1.174919in}}%
\pgfpathclose%
\pgfusepath{fill}%
\end{pgfscope}%
\begin{pgfscope}%
\pgfpathrectangle{\pgfqpoint{0.407592in}{0.222904in}}{\pgfqpoint{1.411163in}{1.411163in}}%
\pgfusepath{clip}%
\pgfsetbuttcap%
\pgfsetroundjoin%
\definecolor{currentfill}{rgb}{0.749020,0.772549,0.792157}%
\pgfsetfillcolor{currentfill}%
\pgfsetfillopacity{0.160000}%
\pgfsetlinewidth{0.000000pt}%
\definecolor{currentstroke}{rgb}{0.000000,0.000000,0.000000}%
\pgfsetstrokecolor{currentstroke}%
\pgfsetdash{}{0pt}%
\pgfpathmoveto{\pgfqpoint{1.155295in}{0.666377in}}%
\pgfpathlineto{\pgfqpoint{1.163568in}{0.673478in}}%
\pgfpathlineto{\pgfqpoint{1.171873in}{0.680606in}}%
\pgfpathlineto{\pgfqpoint{1.178654in}{0.683167in}}%
\pgfpathlineto{\pgfqpoint{1.185034in}{0.685906in}}%
\pgfpathlineto{\pgfqpoint{1.175404in}{0.678249in}}%
\pgfpathlineto{\pgfqpoint{1.165807in}{0.670619in}}%
\pgfpathlineto{\pgfqpoint{1.160710in}{0.668427in}}%
\pgfpathlineto{\pgfqpoint{1.155295in}{0.666377in}}%
\pgfpathclose%
\pgfusepath{fill}%
\end{pgfscope}%
\begin{pgfscope}%
\pgfpathrectangle{\pgfqpoint{0.407592in}{0.222904in}}{\pgfqpoint{1.411163in}{1.411163in}}%
\pgfusepath{clip}%
\pgfsetbuttcap%
\pgfsetroundjoin%
\definecolor{currentfill}{rgb}{0.749020,0.772549,0.792157}%
\pgfsetfillcolor{currentfill}%
\pgfsetfillopacity{0.160000}%
\pgfsetlinewidth{0.000000pt}%
\definecolor{currentstroke}{rgb}{0.000000,0.000000,0.000000}%
\pgfsetstrokecolor{currentstroke}%
\pgfsetdash{}{0pt}%
\pgfpathmoveto{\pgfqpoint{0.822757in}{0.836917in}}%
\pgfpathlineto{\pgfqpoint{0.809190in}{0.848400in}}%
\pgfpathlineto{\pgfqpoint{0.795604in}{0.859898in}}%
\pgfpathlineto{\pgfqpoint{0.795127in}{0.851938in}}%
\pgfpathlineto{\pgfqpoint{0.795943in}{0.843987in}}%
\pgfpathlineto{\pgfqpoint{0.809575in}{0.833173in}}%
\pgfpathlineto{\pgfqpoint{0.823182in}{0.822378in}}%
\pgfpathlineto{\pgfqpoint{0.822383in}{0.829642in}}%
\pgfpathlineto{\pgfqpoint{0.822757in}{0.836917in}}%
\pgfpathclose%
\pgfusepath{fill}%
\end{pgfscope}%
\begin{pgfscope}%
\pgfpathrectangle{\pgfqpoint{0.407592in}{0.222904in}}{\pgfqpoint{1.411163in}{1.411163in}}%
\pgfusepath{clip}%
\pgfsetbuttcap%
\pgfsetroundjoin%
\definecolor{currentfill}{rgb}{0.749020,0.772549,0.792157}%
\pgfsetfillcolor{currentfill}%
\pgfsetfillopacity{0.160000}%
\pgfsetlinewidth{0.000000pt}%
\definecolor{currentstroke}{rgb}{0.000000,0.000000,0.000000}%
\pgfsetstrokecolor{currentstroke}%
\pgfsetdash{}{0pt}%
\pgfpathmoveto{\pgfqpoint{1.732768in}{1.199371in}}%
\pgfpathlineto{\pgfqpoint{1.746451in}{1.211910in}}%
\pgfpathlineto{\pgfqpoint{1.760147in}{1.224461in}}%
\pgfpathlineto{\pgfqpoint{1.751640in}{1.240394in}}%
\pgfpathlineto{\pgfqpoint{1.740385in}{1.255959in}}%
\pgfpathlineto{\pgfqpoint{1.727232in}{1.242902in}}%
\pgfpathlineto{\pgfqpoint{1.714085in}{1.229851in}}%
\pgfpathlineto{\pgfqpoint{1.724748in}{1.214785in}}%
\pgfpathlineto{\pgfqpoint{1.732768in}{1.199371in}}%
\pgfpathclose%
\pgfusepath{fill}%
\end{pgfscope}%
\begin{pgfscope}%
\pgfpathrectangle{\pgfqpoint{0.407592in}{0.222904in}}{\pgfqpoint{1.411163in}{1.411163in}}%
\pgfusepath{clip}%
\pgfsetbuttcap%
\pgfsetroundjoin%
\definecolor{currentfill}{rgb}{0.749020,0.772549,0.792157}%
\pgfsetfillcolor{currentfill}%
\pgfsetfillopacity{0.160000}%
\pgfsetlinewidth{0.000000pt}%
\definecolor{currentstroke}{rgb}{0.000000,0.000000,0.000000}%
\pgfsetstrokecolor{currentstroke}%
\pgfsetdash{}{0pt}%
\pgfpathmoveto{\pgfqpoint{0.720296in}{0.948980in}}%
\pgfpathlineto{\pgfqpoint{0.706903in}{0.961233in}}%
\pgfpathlineto{\pgfqpoint{0.693499in}{0.973497in}}%
\pgfpathlineto{\pgfqpoint{0.689009in}{0.963028in}}%
\pgfpathlineto{\pgfqpoint{0.686242in}{0.952460in}}%
\pgfpathlineto{\pgfqpoint{0.699979in}{0.940834in}}%
\pgfpathlineto{\pgfqpoint{0.713696in}{0.929224in}}%
\pgfpathlineto{\pgfqpoint{0.716191in}{0.939146in}}%
\pgfpathlineto{\pgfqpoint{0.720296in}{0.948980in}}%
\pgfpathclose%
\pgfusepath{fill}%
\end{pgfscope}%
\begin{pgfscope}%
\pgfpathrectangle{\pgfqpoint{0.407592in}{0.222904in}}{\pgfqpoint{1.411163in}{1.411163in}}%
\pgfusepath{clip}%
\pgfsetbuttcap%
\pgfsetroundjoin%
\definecolor{currentfill}{rgb}{0.749020,0.772549,0.792157}%
\pgfsetfillcolor{currentfill}%
\pgfsetfillopacity{0.160000}%
\pgfsetlinewidth{0.000000pt}%
\definecolor{currentstroke}{rgb}{0.000000,0.000000,0.000000}%
\pgfsetstrokecolor{currentstroke}%
\pgfsetdash{}{0pt}%
\pgfpathmoveto{\pgfqpoint{1.293372in}{0.768960in}}%
\pgfpathlineto{\pgfqpoint{1.306306in}{0.779047in}}%
\pgfpathlineto{\pgfqpoint{1.319268in}{0.789156in}}%
\pgfpathlineto{\pgfqpoint{1.323332in}{0.795627in}}%
\pgfpathlineto{\pgfqpoint{1.326354in}{0.802164in}}%
\pgfpathlineto{\pgfqpoint{1.313039in}{0.791358in}}%
\pgfpathlineto{\pgfqpoint{1.299747in}{0.780571in}}%
\pgfpathlineto{\pgfqpoint{1.297021in}{0.774735in}}%
\pgfpathlineto{\pgfqpoint{1.293372in}{0.768960in}}%
\pgfpathclose%
\pgfusepath{fill}%
\end{pgfscope}%
\begin{pgfscope}%
\pgfpathrectangle{\pgfqpoint{0.407592in}{0.222904in}}{\pgfqpoint{1.411163in}{1.411163in}}%
\pgfusepath{clip}%
\pgfsetbuttcap%
\pgfsetroundjoin%
\definecolor{currentfill}{rgb}{0.749020,0.772549,0.792157}%
\pgfsetfillcolor{currentfill}%
\pgfsetfillopacity{0.160000}%
\pgfsetlinewidth{0.000000pt}%
\definecolor{currentstroke}{rgb}{0.000000,0.000000,0.000000}%
\pgfsetstrokecolor{currentstroke}%
\pgfsetdash{}{0pt}%
\pgfpathmoveto{\pgfqpoint{1.196476in}{0.691865in}}%
\pgfpathlineto{\pgfqpoint{1.207287in}{0.700140in}}%
\pgfpathlineto{\pgfqpoint{1.218131in}{0.708442in}}%
\pgfpathlineto{\pgfqpoint{1.224142in}{0.712266in}}%
\pgfpathlineto{\pgfqpoint{1.212799in}{0.703649in}}%
\pgfpathlineto{\pgfqpoint{1.201490in}{0.695057in}}%
\pgfusepath{fill}%
\end{pgfscope}%
\begin{pgfscope}%
\pgfpathrectangle{\pgfqpoint{0.407592in}{0.222904in}}{\pgfqpoint{1.411163in}{1.411163in}}%
\pgfusepath{clip}%
\pgfsetbuttcap%
\pgfsetroundjoin%
\definecolor{currentfill}{rgb}{0.749020,0.772549,0.792157}%
\pgfsetfillcolor{currentfill}%
\pgfsetfillopacity{0.160000}%
\pgfsetlinewidth{0.000000pt}%
\definecolor{currentstroke}{rgb}{0.000000,0.000000,0.000000}%
\pgfsetstrokecolor{currentstroke}%
\pgfsetdash{}{0pt}%
\pgfpathmoveto{\pgfqpoint{0.500965in}{1.200875in}}%
\pgfpathlineto{\pgfqpoint{0.488000in}{1.213864in}}%
\pgfpathlineto{\pgfqpoint{0.475029in}{1.226859in}}%
\pgfpathlineto{\pgfqpoint{0.461241in}{1.211719in}}%
\pgfpathlineto{\pgfqpoint{0.450032in}{1.196241in}}%
\pgfpathlineto{\pgfqpoint{0.463667in}{1.183767in}}%
\pgfpathlineto{\pgfqpoint{0.477289in}{1.171304in}}%
\pgfpathlineto{\pgfqpoint{0.487888in}{1.186249in}}%
\pgfpathlineto{\pgfqpoint{0.500965in}{1.200875in}}%
\pgfpathclose%
\pgfusepath{fill}%
\end{pgfscope}%
\begin{pgfscope}%
\pgfpathrectangle{\pgfqpoint{0.407592in}{0.222904in}}{\pgfqpoint{1.411163in}{1.411163in}}%
\pgfusepath{clip}%
\pgfsetbuttcap%
\pgfsetroundjoin%
\definecolor{currentfill}{rgb}{0.749020,0.772549,0.792157}%
\pgfsetfillcolor{currentfill}%
\pgfsetfillopacity{0.160000}%
\pgfsetlinewidth{0.000000pt}%
\definecolor{currentstroke}{rgb}{0.000000,0.000000,0.000000}%
\pgfsetstrokecolor{currentstroke}%
\pgfsetdash{}{0pt}%
\pgfpathmoveto{\pgfqpoint{0.877364in}{0.779396in}}%
\pgfpathlineto{\pgfqpoint{0.863855in}{0.790113in}}%
\pgfpathlineto{\pgfqpoint{0.850322in}{0.800848in}}%
\pgfpathlineto{\pgfqpoint{0.852150in}{0.794322in}}%
\pgfpathlineto{\pgfqpoint{0.855032in}{0.787868in}}%
\pgfpathlineto{\pgfqpoint{0.868348in}{0.777829in}}%
\pgfpathlineto{\pgfqpoint{0.881633in}{0.767811in}}%
\pgfpathlineto{\pgfqpoint{0.879032in}{0.773570in}}%
\pgfpathlineto{\pgfqpoint{0.877364in}{0.779396in}}%
\pgfpathclose%
\pgfusepath{fill}%
\end{pgfscope}%
\begin{pgfscope}%
\pgfpathrectangle{\pgfqpoint{0.407592in}{0.222904in}}{\pgfqpoint{1.411163in}{1.411163in}}%
\pgfusepath{clip}%
\pgfsetbuttcap%
\pgfsetroundjoin%
\definecolor{currentfill}{rgb}{0.749020,0.772549,0.792157}%
\pgfsetfillcolor{currentfill}%
\pgfsetfillopacity{0.160000}%
\pgfsetlinewidth{0.000000pt}%
\definecolor{currentstroke}{rgb}{0.000000,0.000000,0.000000}%
\pgfsetstrokecolor{currentstroke}%
\pgfsetdash{}{0pt}%
\pgfpathmoveto{\pgfqpoint{1.013435in}{0.668004in}}%
\pgfpathlineto{\pgfqpoint{1.003676in}{0.675308in}}%
\pgfpathlineto{\pgfqpoint{0.993880in}{0.682638in}}%
\pgfpathlineto{\pgfqpoint{1.000211in}{0.680115in}}%
\pgfpathlineto{\pgfqpoint{1.006948in}{0.677783in}}%
\pgfpathlineto{\pgfqpoint{1.015426in}{0.670937in}}%
\pgfpathlineto{\pgfqpoint{1.023870in}{0.664119in}}%
\pgfpathlineto{\pgfqpoint{1.018491in}{0.665984in}}%
\pgfpathlineto{\pgfqpoint{1.013435in}{0.668004in}}%
\pgfpathclose%
\pgfusepath{fill}%
\end{pgfscope}%
\begin{pgfscope}%
\pgfpathrectangle{\pgfqpoint{0.407592in}{0.222904in}}{\pgfqpoint{1.411163in}{1.411163in}}%
\pgfusepath{clip}%
\pgfsetbuttcap%
\pgfsetroundjoin%
\definecolor{currentfill}{rgb}{0.749020,0.772549,0.792157}%
\pgfsetfillcolor{currentfill}%
\pgfsetfillopacity{0.160000}%
\pgfsetlinewidth{0.000000pt}%
\definecolor{currentstroke}{rgb}{0.000000,0.000000,0.000000}%
\pgfsetstrokecolor{currentstroke}%
\pgfsetdash{}{0pt}%
\pgfpathmoveto{\pgfqpoint{1.491799in}{0.954586in}}%
\pgfpathlineto{\pgfqpoint{1.505463in}{0.966293in}}%
\pgfpathlineto{\pgfqpoint{1.519145in}{0.978016in}}%
\pgfpathlineto{\pgfqpoint{1.518255in}{0.989201in}}%
\pgfpathlineto{\pgfqpoint{1.515507in}{1.000268in}}%
\pgfpathlineto{\pgfqpoint{1.502032in}{0.987919in}}%
\pgfpathlineto{\pgfqpoint{1.488569in}{0.975581in}}%
\pgfpathlineto{\pgfqpoint{1.491055in}{0.965138in}}%
\pgfpathlineto{\pgfqpoint{1.491799in}{0.954586in}}%
\pgfpathclose%
\pgfusepath{fill}%
\end{pgfscope}%
\begin{pgfscope}%
\pgfpathrectangle{\pgfqpoint{0.407592in}{0.222904in}}{\pgfqpoint{1.411163in}{1.411163in}}%
\pgfusepath{clip}%
\pgfsetbuttcap%
\pgfsetroundjoin%
\definecolor{currentfill}{rgb}{0.749020,0.772549,0.792157}%
\pgfsetfillcolor{currentfill}%
\pgfsetfillopacity{0.160000}%
\pgfsetlinewidth{0.000000pt}%
\definecolor{currentstroke}{rgb}{0.000000,0.000000,0.000000}%
\pgfsetstrokecolor{currentstroke}%
\pgfsetdash{}{0pt}%
\pgfpathmoveto{\pgfqpoint{1.760147in}{1.224461in}}%
\pgfpathlineto{\pgfqpoint{1.773855in}{1.237024in}}%
\pgfpathlineto{\pgfqpoint{1.787575in}{1.249598in}}%
\pgfpathlineto{\pgfqpoint{1.778569in}{1.266038in}}%
\pgfpathlineto{\pgfqpoint{1.766710in}{1.282093in}}%
\pgfpathlineto{\pgfqpoint{1.753545in}{1.269023in}}%
\pgfpathlineto{\pgfqpoint{1.740385in}{1.255959in}}%
\pgfpathlineto{\pgfqpoint{1.751640in}{1.240394in}}%
\pgfpathlineto{\pgfqpoint{1.760147in}{1.224461in}}%
\pgfpathclose%
\pgfusepath{fill}%
\end{pgfscope}%
\begin{pgfscope}%
\pgfpathrectangle{\pgfqpoint{0.407592in}{0.222904in}}{\pgfqpoint{1.411163in}{1.411163in}}%
\pgfusepath{clip}%
\pgfsetbuttcap%
\pgfsetroundjoin%
\definecolor{currentfill}{rgb}{0.749020,0.772549,0.792157}%
\pgfsetfillcolor{currentfill}%
\pgfsetfillopacity{0.160000}%
\pgfsetlinewidth{0.000000pt}%
\definecolor{currentstroke}{rgb}{0.000000,0.000000,0.000000}%
\pgfsetstrokecolor{currentstroke}%
\pgfsetdash{}{0pt}%
\pgfpathmoveto{\pgfqpoint{0.973145in}{0.694406in}}%
\pgfpathlineto{\pgfqpoint{0.961242in}{0.702934in}}%
\pgfpathlineto{\pgfqpoint{0.949302in}{0.711487in}}%
\pgfpathlineto{\pgfqpoint{0.954632in}{0.707694in}}%
\pgfpathlineto{\pgfqpoint{0.960572in}{0.704067in}}%
\pgfpathlineto{\pgfqpoint{0.971578in}{0.696127in}}%
\pgfpathlineto{\pgfqpoint{0.982547in}{0.688214in}}%
\pgfpathlineto{\pgfqpoint{0.977594in}{0.691240in}}%
\pgfpathlineto{\pgfqpoint{0.973145in}{0.694406in}}%
\pgfpathclose%
\pgfusepath{fill}%
\end{pgfscope}%
\begin{pgfscope}%
\pgfpathrectangle{\pgfqpoint{0.407592in}{0.222904in}}{\pgfqpoint{1.411163in}{1.411163in}}%
\pgfusepath{clip}%
\pgfsetbuttcap%
\pgfsetroundjoin%
\definecolor{currentfill}{rgb}{0.749020,0.772549,0.792157}%
\pgfsetfillcolor{currentfill}%
\pgfsetfillopacity{0.160000}%
\pgfsetlinewidth{0.000000pt}%
\definecolor{currentstroke}{rgb}{0.000000,0.000000,0.000000}%
\pgfsetstrokecolor{currentstroke}%
\pgfsetdash{}{0pt}%
\pgfpathmoveto{\pgfqpoint{1.143606in}{0.662744in}}%
\pgfpathlineto{\pgfqpoint{1.150401in}{0.669391in}}%
\pgfpathlineto{\pgfqpoint{1.157223in}{0.676065in}}%
\pgfpathlineto{\pgfqpoint{1.164718in}{0.678235in}}%
\pgfpathlineto{\pgfqpoint{1.171873in}{0.680606in}}%
\pgfpathlineto{\pgfqpoint{1.163568in}{0.673478in}}%
\pgfpathlineto{\pgfqpoint{1.155295in}{0.666377in}}%
\pgfpathlineto{\pgfqpoint{1.149585in}{0.664480in}}%
\pgfpathlineto{\pgfqpoint{1.143606in}{0.662744in}}%
\pgfpathclose%
\pgfusepath{fill}%
\end{pgfscope}%
\begin{pgfscope}%
\pgfpathrectangle{\pgfqpoint{0.407592in}{0.222904in}}{\pgfqpoint{1.411163in}{1.411163in}}%
\pgfusepath{clip}%
\pgfsetbuttcap%
\pgfsetroundjoin%
\definecolor{currentfill}{rgb}{0.749020,0.772549,0.792157}%
\pgfsetfillcolor{currentfill}%
\pgfsetfillopacity{0.160000}%
\pgfsetlinewidth{0.000000pt}%
\definecolor{currentstroke}{rgb}{0.000000,0.000000,0.000000}%
\pgfsetstrokecolor{currentstroke}%
\pgfsetdash{}{0pt}%
\pgfpathmoveto{\pgfqpoint{1.379850in}{0.845580in}}%
\pgfpathlineto{\pgfqpoint{1.393284in}{0.856482in}}%
\pgfpathlineto{\pgfqpoint{1.406741in}{0.867403in}}%
\pgfpathlineto{\pgfqpoint{1.409175in}{0.876039in}}%
\pgfpathlineto{\pgfqpoint{1.410198in}{0.884673in}}%
\pgfpathlineto{\pgfqpoint{1.396662in}{0.873076in}}%
\pgfpathlineto{\pgfqpoint{1.383143in}{0.861493in}}%
\pgfpathlineto{\pgfqpoint{1.382142in}{0.853535in}}%
\pgfpathlineto{\pgfqpoint{1.379850in}{0.845580in}}%
\pgfpathclose%
\pgfusepath{fill}%
\end{pgfscope}%
\begin{pgfscope}%
\pgfpathrectangle{\pgfqpoint{0.407592in}{0.222904in}}{\pgfqpoint{1.411163in}{1.411163in}}%
\pgfusepath{clip}%
\pgfsetbuttcap%
\pgfsetroundjoin%
\definecolor{currentfill}{rgb}{0.749020,0.772549,0.792157}%
\pgfsetfillcolor{currentfill}%
\pgfsetfillopacity{0.160000}%
\pgfsetlinewidth{0.000000pt}%
\definecolor{currentstroke}{rgb}{0.000000,0.000000,0.000000}%
\pgfsetstrokecolor{currentstroke}%
\pgfsetdash{}{0pt}%
\pgfpathmoveto{\pgfqpoint{0.475029in}{1.226859in}}%
\pgfpathlineto{\pgfqpoint{0.462051in}{1.239861in}}%
\pgfpathlineto{\pgfqpoint{0.449066in}{1.252869in}}%
\pgfpathlineto{\pgfqpoint{0.434555in}{1.237225in}}%
\pgfpathlineto{\pgfqpoint{0.422723in}{1.221226in}}%
\pgfpathlineto{\pgfqpoint{0.436384in}{1.208728in}}%
\pgfpathlineto{\pgfqpoint{0.450032in}{1.196241in}}%
\pgfpathlineto{\pgfqpoint{0.461241in}{1.211719in}}%
\pgfpathlineto{\pgfqpoint{0.475029in}{1.226859in}}%
\pgfpathclose%
\pgfusepath{fill}%
\end{pgfscope}%
\begin{pgfscope}%
\pgfpathrectangle{\pgfqpoint{0.407592in}{0.222904in}}{\pgfqpoint{1.411163in}{1.411163in}}%
\pgfusepath{clip}%
\pgfsetbuttcap%
\pgfsetroundjoin%
\definecolor{currentfill}{rgb}{0.749020,0.772549,0.792157}%
\pgfsetfillcolor{currentfill}%
\pgfsetfillopacity{0.160000}%
\pgfsetlinewidth{0.000000pt}%
\definecolor{currentstroke}{rgb}{0.000000,0.000000,0.000000}%
\pgfsetstrokecolor{currentstroke}%
\pgfsetdash{}{0pt}%
\pgfpathmoveto{\pgfqpoint{1.258778in}{0.738968in}}%
\pgfpathlineto{\pgfqpoint{1.271053in}{0.748321in}}%
\pgfpathlineto{\pgfqpoint{1.283360in}{0.757700in}}%
\pgfpathlineto{\pgfqpoint{1.288814in}{0.763273in}}%
\pgfpathlineto{\pgfqpoint{1.293372in}{0.768960in}}%
\pgfpathlineto{\pgfqpoint{1.280467in}{0.758895in}}%
\pgfpathlineto{\pgfqpoint{1.267590in}{0.748853in}}%
\pgfpathlineto{\pgfqpoint{1.263573in}{0.743859in}}%
\pgfpathlineto{\pgfqpoint{1.258778in}{0.738968in}}%
\pgfpathclose%
\pgfusepath{fill}%
\end{pgfscope}%
\begin{pgfscope}%
\pgfpathrectangle{\pgfqpoint{0.407592in}{0.222904in}}{\pgfqpoint{1.411163in}{1.411163in}}%
\pgfusepath{clip}%
\pgfsetbuttcap%
\pgfsetroundjoin%
\definecolor{currentfill}{rgb}{0.749020,0.772549,0.792157}%
\pgfsetfillcolor{currentfill}%
\pgfsetfillopacity{0.160000}%
\pgfsetlinewidth{0.000000pt}%
\definecolor{currentstroke}{rgb}{0.000000,0.000000,0.000000}%
\pgfsetstrokecolor{currentstroke}%
\pgfsetdash{}{0pt}%
\pgfpathmoveto{\pgfqpoint{0.693499in}{0.973497in}}%
\pgfpathlineto{\pgfqpoint{0.680081in}{0.985772in}}%
\pgfpathlineto{\pgfqpoint{0.666651in}{0.998060in}}%
\pgfpathlineto{\pgfqpoint{0.661762in}{0.986965in}}%
\pgfpathlineto{\pgfqpoint{0.658712in}{0.975761in}}%
\pgfpathlineto{\pgfqpoint{0.672487in}{0.964103in}}%
\pgfpathlineto{\pgfqpoint{0.686242in}{0.952460in}}%
\pgfpathlineto{\pgfqpoint{0.689009in}{0.963028in}}%
\pgfpathlineto{\pgfqpoint{0.693499in}{0.973497in}}%
\pgfpathclose%
\pgfusepath{fill}%
\end{pgfscope}%
\begin{pgfscope}%
\pgfpathrectangle{\pgfqpoint{0.407592in}{0.222904in}}{\pgfqpoint{1.411163in}{1.411163in}}%
\pgfusepath{clip}%
\pgfsetbuttcap%
\pgfsetroundjoin%
\definecolor{currentfill}{rgb}{0.749020,0.772549,0.792157}%
\pgfsetfillcolor{currentfill}%
\pgfsetfillopacity{0.160000}%
\pgfsetlinewidth{0.000000pt}%
\definecolor{currentstroke}{rgb}{0.000000,0.000000,0.000000}%
\pgfsetstrokecolor{currentstroke}%
\pgfsetdash{}{0pt}%
\pgfpathmoveto{\pgfqpoint{1.224142in}{0.712266in}}%
\pgfpathlineto{\pgfqpoint{1.235519in}{0.720909in}}%
\pgfpathlineto{\pgfqpoint{1.246930in}{0.729579in}}%
\pgfpathlineto{\pgfqpoint{1.253223in}{0.734200in}}%
\pgfpathlineto{\pgfqpoint{1.258778in}{0.738968in}}%
\pgfpathlineto{\pgfqpoint{1.246535in}{0.729638in}}%
\pgfpathlineto{\pgfqpoint{1.234324in}{0.720333in}}%
\pgfpathlineto{\pgfqpoint{1.229547in}{0.716236in}}%
\pgfpathlineto{\pgfqpoint{1.224142in}{0.712266in}}%
\pgfpathclose%
\pgfusepath{fill}%
\end{pgfscope}%
\begin{pgfscope}%
\pgfpathrectangle{\pgfqpoint{0.407592in}{0.222904in}}{\pgfqpoint{1.411163in}{1.411163in}}%
\pgfusepath{clip}%
\pgfsetbuttcap%
\pgfsetroundjoin%
\definecolor{currentfill}{rgb}{0.749020,0.772549,0.792157}%
\pgfsetfillcolor{currentfill}%
\pgfsetfillopacity{0.160000}%
\pgfsetlinewidth{0.000000pt}%
\definecolor{currentstroke}{rgb}{0.000000,0.000000,0.000000}%
\pgfsetstrokecolor{currentstroke}%
\pgfsetdash{}{0pt}%
\pgfpathmoveto{\pgfqpoint{1.787575in}{1.249598in}}%
\pgfpathlineto{\pgfqpoint{1.801308in}{1.262183in}}%
\pgfpathlineto{\pgfqpoint{1.815054in}{1.274780in}}%
\pgfpathlineto{\pgfqpoint{1.805534in}{1.291718in}}%
\pgfpathlineto{\pgfqpoint{1.793060in}{1.308251in}}%
\pgfpathlineto{\pgfqpoint{1.779882in}{1.295169in}}%
\pgfpathlineto{\pgfqpoint{1.766710in}{1.282093in}}%
\pgfpathlineto{\pgfqpoint{1.778569in}{1.266038in}}%
\pgfpathlineto{\pgfqpoint{1.787575in}{1.249598in}}%
\pgfpathclose%
\pgfusepath{fill}%
\end{pgfscope}%
\begin{pgfscope}%
\pgfpathrectangle{\pgfqpoint{0.407592in}{0.222904in}}{\pgfqpoint{1.411163in}{1.411163in}}%
\pgfusepath{clip}%
\pgfsetbuttcap%
\pgfsetroundjoin%
\definecolor{currentfill}{rgb}{0.749020,0.772549,0.792157}%
\pgfsetfillcolor{currentfill}%
\pgfsetfillopacity{0.160000}%
\pgfsetlinewidth{0.000000pt}%
\definecolor{currentstroke}{rgb}{0.000000,0.000000,0.000000}%
\pgfsetstrokecolor{currentstroke}%
\pgfsetdash{}{0pt}%
\pgfpathmoveto{\pgfqpoint{0.908116in}{0.747844in}}%
\pgfpathlineto{\pgfqpoint{0.894889in}{0.757816in}}%
\pgfpathlineto{\pgfqpoint{0.881633in}{0.767811in}}%
\pgfpathlineto{\pgfqpoint{0.885162in}{0.762145in}}%
\pgfpathlineto{\pgfqpoint{0.889607in}{0.756597in}}%
\pgfpathlineto{\pgfqpoint{0.902396in}{0.747287in}}%
\pgfpathlineto{\pgfqpoint{0.915150in}{0.738000in}}%
\pgfpathlineto{\pgfqpoint{0.911234in}{0.742869in}}%
\pgfpathlineto{\pgfqpoint{0.908116in}{0.747844in}}%
\pgfpathclose%
\pgfusepath{fill}%
\end{pgfscope}%
\begin{pgfscope}%
\pgfpathrectangle{\pgfqpoint{0.407592in}{0.222904in}}{\pgfqpoint{1.411163in}{1.411163in}}%
\pgfusepath{clip}%
\pgfsetbuttcap%
\pgfsetroundjoin%
\definecolor{currentfill}{rgb}{0.749020,0.772549,0.792157}%
\pgfsetfillcolor{currentfill}%
\pgfsetfillopacity{0.160000}%
\pgfsetlinewidth{0.000000pt}%
\definecolor{currentstroke}{rgb}{0.000000,0.000000,0.000000}%
\pgfsetstrokecolor{currentstroke}%
\pgfsetdash{}{0pt}%
\pgfpathmoveto{\pgfqpoint{0.795604in}{0.859898in}}%
\pgfpathlineto{\pgfqpoint{0.782000in}{0.871413in}}%
\pgfpathlineto{\pgfqpoint{0.768377in}{0.882943in}}%
\pgfpathlineto{\pgfqpoint{0.767784in}{0.874305in}}%
\pgfpathlineto{\pgfqpoint{0.768605in}{0.865674in}}%
\pgfpathlineto{\pgfqpoint{0.782286in}{0.854821in}}%
\pgfpathlineto{\pgfqpoint{0.795943in}{0.843987in}}%
\pgfpathlineto{\pgfqpoint{0.795127in}{0.851938in}}%
\pgfpathlineto{\pgfqpoint{0.795604in}{0.859898in}}%
\pgfpathclose%
\pgfusepath{fill}%
\end{pgfscope}%
\begin{pgfscope}%
\pgfpathrectangle{\pgfqpoint{0.407592in}{0.222904in}}{\pgfqpoint{1.411163in}{1.411163in}}%
\pgfusepath{clip}%
\pgfsetbuttcap%
\pgfsetroundjoin%
\definecolor{currentfill}{rgb}{0.749020,0.772549,0.792157}%
\pgfsetfillcolor{currentfill}%
\pgfsetfillopacity{0.160000}%
\pgfsetlinewidth{0.000000pt}%
\definecolor{currentstroke}{rgb}{0.000000,0.000000,0.000000}%
\pgfsetstrokecolor{currentstroke}%
\pgfsetdash{}{0pt}%
\pgfpathmoveto{\pgfqpoint{1.185034in}{0.685906in}}%
\pgfpathlineto{\pgfqpoint{1.194698in}{0.693589in}}%
\pgfpathlineto{\pgfqpoint{1.204397in}{0.701300in}}%
\pgfpathlineto{\pgfqpoint{1.211540in}{0.704780in}}%
\pgfpathlineto{\pgfqpoint{1.218131in}{0.708442in}}%
\pgfpathlineto{\pgfqpoint{1.207287in}{0.700140in}}%
\pgfpathlineto{\pgfqpoint{1.196476in}{0.691865in}}%
\pgfpathlineto{\pgfqpoint{1.190983in}{0.688809in}}%
\pgfpathlineto{\pgfqpoint{1.185034in}{0.685906in}}%
\pgfpathclose%
\pgfusepath{fill}%
\end{pgfscope}%
\begin{pgfscope}%
\pgfpathrectangle{\pgfqpoint{0.407592in}{0.222904in}}{\pgfqpoint{1.411163in}{1.411163in}}%
\pgfusepath{clip}%
\pgfsetbuttcap%
\pgfsetroundjoin%
\definecolor{currentfill}{rgb}{0.749020,0.772549,0.792157}%
\pgfsetfillcolor{currentfill}%
\pgfsetfillopacity{0.160000}%
\pgfsetlinewidth{0.000000pt}%
\definecolor{currentstroke}{rgb}{0.000000,0.000000,0.000000}%
\pgfsetstrokecolor{currentstroke}%
\pgfsetdash{}{0pt}%
\pgfpathmoveto{\pgfqpoint{1.023870in}{0.664119in}}%
\pgfpathlineto{\pgfqpoint{1.015426in}{0.670937in}}%
\pgfpathlineto{\pgfqpoint{1.006948in}{0.677783in}}%
\pgfpathlineto{\pgfqpoint{1.014062in}{0.675654in}}%
\pgfpathlineto{\pgfqpoint{1.021521in}{0.673739in}}%
\pgfpathlineto{\pgfqpoint{1.028522in}{0.667298in}}%
\pgfpathlineto{\pgfqpoint{1.035494in}{0.660885in}}%
\pgfpathlineto{\pgfqpoint{1.029545in}{0.662416in}}%
\pgfpathlineto{\pgfqpoint{1.023870in}{0.664119in}}%
\pgfpathclose%
\pgfusepath{fill}%
\end{pgfscope}%
\begin{pgfscope}%
\pgfpathrectangle{\pgfqpoint{0.407592in}{0.222904in}}{\pgfqpoint{1.411163in}{1.411163in}}%
\pgfusepath{clip}%
\pgfsetbuttcap%
\pgfsetroundjoin%
\definecolor{currentfill}{rgb}{0.749020,0.772549,0.792157}%
\pgfsetfillcolor{currentfill}%
\pgfsetfillopacity{0.160000}%
\pgfsetlinewidth{0.000000pt}%
\definecolor{currentstroke}{rgb}{0.000000,0.000000,0.000000}%
\pgfsetstrokecolor{currentstroke}%
\pgfsetdash{}{0pt}%
\pgfpathmoveto{\pgfqpoint{0.940558in}{0.719502in}}%
\pgfpathlineto{\pgfqpoint{0.927871in}{0.728739in}}%
\pgfpathlineto{\pgfqpoint{0.915150in}{0.738000in}}%
\pgfpathlineto{\pgfqpoint{0.919849in}{0.733260in}}%
\pgfpathlineto{\pgfqpoint{0.925313in}{0.728671in}}%
\pgfpathlineto{\pgfqpoint{0.937326in}{0.720066in}}%
\pgfpathlineto{\pgfqpoint{0.949302in}{0.711487in}}%
\pgfpathlineto{\pgfqpoint{0.944605in}{0.715429in}}%
\pgfpathlineto{\pgfqpoint{0.940558in}{0.719502in}}%
\pgfpathclose%
\pgfusepath{fill}%
\end{pgfscope}%
\begin{pgfscope}%
\pgfpathrectangle{\pgfqpoint{0.407592in}{0.222904in}}{\pgfqpoint{1.411163in}{1.411163in}}%
\pgfusepath{clip}%
\pgfsetbuttcap%
\pgfsetroundjoin%
\definecolor{currentfill}{rgb}{0.749020,0.772549,0.792157}%
\pgfsetfillcolor{currentfill}%
\pgfsetfillopacity{0.160000}%
\pgfsetlinewidth{0.000000pt}%
\definecolor{currentstroke}{rgb}{0.000000,0.000000,0.000000}%
\pgfsetstrokecolor{currentstroke}%
\pgfsetdash{}{0pt}%
\pgfpathmoveto{\pgfqpoint{1.519145in}{0.978016in}}%
\pgfpathlineto{\pgfqpoint{1.532845in}{0.989754in}}%
\pgfpathlineto{\pgfqpoint{1.546564in}{1.001508in}}%
\pgfpathlineto{\pgfqpoint{1.545516in}{1.013319in}}%
\pgfpathlineto{\pgfqpoint{1.542494in}{1.024999in}}%
\pgfpathlineto{\pgfqpoint{1.528995in}{1.012628in}}%
\pgfpathlineto{\pgfqpoint{1.515507in}{1.000268in}}%
\pgfpathlineto{\pgfqpoint{1.518255in}{0.989201in}}%
\pgfpathlineto{\pgfqpoint{1.519145in}{0.978016in}}%
\pgfpathclose%
\pgfusepath{fill}%
\end{pgfscope}%
\begin{pgfscope}%
\pgfpathrectangle{\pgfqpoint{0.407592in}{0.222904in}}{\pgfqpoint{1.411163in}{1.411163in}}%
\pgfusepath{clip}%
\pgfsetbuttcap%
\pgfsetroundjoin%
\definecolor{currentfill}{rgb}{0.749020,0.772549,0.792157}%
\pgfsetfillcolor{currentfill}%
\pgfsetfillopacity{0.160000}%
\pgfsetlinewidth{0.000000pt}%
\definecolor{currentstroke}{rgb}{0.000000,0.000000,0.000000}%
\pgfsetstrokecolor{currentstroke}%
\pgfsetdash{}{0pt}%
\pgfpathmoveto{\pgfqpoint{0.449066in}{1.252869in}}%
\pgfpathlineto{\pgfqpoint{0.436075in}{1.265884in}}%
\pgfpathlineto{\pgfqpoint{0.423077in}{1.278906in}}%
\pgfpathlineto{\pgfqpoint{0.407830in}{1.262769in}}%
\pgfpathlineto{\pgfqpoint{0.395362in}{1.246258in}}%
\pgfpathlineto{\pgfqpoint{0.409049in}{1.233736in}}%
\pgfpathlineto{\pgfqpoint{0.422723in}{1.221226in}}%
\pgfpathlineto{\pgfqpoint{0.434555in}{1.237225in}}%
\pgfpathlineto{\pgfqpoint{0.449066in}{1.252869in}}%
\pgfpathclose%
\pgfusepath{fill}%
\end{pgfscope}%
\begin{pgfscope}%
\pgfpathrectangle{\pgfqpoint{0.407592in}{0.222904in}}{\pgfqpoint{1.411163in}{1.411163in}}%
\pgfusepath{clip}%
\pgfsetbuttcap%
\pgfsetroundjoin%
\definecolor{currentfill}{rgb}{0.749020,0.772549,0.792157}%
\pgfsetfillcolor{currentfill}%
\pgfsetfillopacity{0.160000}%
\pgfsetlinewidth{0.000000pt}%
\definecolor{currentstroke}{rgb}{0.000000,0.000000,0.000000}%
\pgfsetstrokecolor{currentstroke}%
\pgfsetdash{}{0pt}%
\pgfpathmoveto{\pgfqpoint{1.130948in}{0.659788in}}%
\pgfpathlineto{\pgfqpoint{1.136136in}{0.666064in}}%
\pgfpathlineto{\pgfqpoint{1.141347in}{0.672367in}}%
\pgfpathlineto{\pgfqpoint{1.149421in}{0.674106in}}%
\pgfpathlineto{\pgfqpoint{1.157223in}{0.676065in}}%
\pgfpathlineto{\pgfqpoint{1.150401in}{0.669391in}}%
\pgfpathlineto{\pgfqpoint{1.143606in}{0.662744in}}%
\pgfpathlineto{\pgfqpoint{1.137384in}{0.661178in}}%
\pgfpathlineto{\pgfqpoint{1.130948in}{0.659788in}}%
\pgfpathclose%
\pgfusepath{fill}%
\end{pgfscope}%
\begin{pgfscope}%
\pgfpathrectangle{\pgfqpoint{0.407592in}{0.222904in}}{\pgfqpoint{1.411163in}{1.411163in}}%
\pgfusepath{clip}%
\pgfsetbuttcap%
\pgfsetroundjoin%
\definecolor{currentfill}{rgb}{0.749020,0.772549,0.792157}%
\pgfsetfillcolor{currentfill}%
\pgfsetfillopacity{0.160000}%
\pgfsetlinewidth{0.000000pt}%
\definecolor{currentstroke}{rgb}{0.000000,0.000000,0.000000}%
\pgfsetstrokecolor{currentstroke}%
\pgfsetdash{}{0pt}%
\pgfpathmoveto{\pgfqpoint{1.319268in}{0.789156in}}%
\pgfpathlineto{\pgfqpoint{1.332259in}{0.799287in}}%
\pgfpathlineto{\pgfqpoint{1.345279in}{0.809441in}}%
\pgfpathlineto{\pgfqpoint{1.349747in}{0.816602in}}%
\pgfpathlineto{\pgfqpoint{1.353055in}{0.823833in}}%
\pgfpathlineto{\pgfqpoint{1.339692in}{0.812989in}}%
\pgfpathlineto{\pgfqpoint{1.326354in}{0.802164in}}%
\pgfpathlineto{\pgfqpoint{1.323332in}{0.795627in}}%
\pgfpathlineto{\pgfqpoint{1.319268in}{0.789156in}}%
\pgfpathclose%
\pgfusepath{fill}%
\end{pgfscope}%
\begin{pgfscope}%
\pgfpathrectangle{\pgfqpoint{0.407592in}{0.222904in}}{\pgfqpoint{1.411163in}{1.411163in}}%
\pgfusepath{clip}%
\pgfsetbuttcap%
\pgfsetroundjoin%
\definecolor{currentfill}{rgb}{0.749020,0.772549,0.792157}%
\pgfsetfillcolor{currentfill}%
\pgfsetfillopacity{0.160000}%
\pgfsetlinewidth{0.000000pt}%
\definecolor{currentstroke}{rgb}{0.000000,0.000000,0.000000}%
\pgfsetstrokecolor{currentstroke}%
\pgfsetdash{}{0pt}%
\pgfpathmoveto{\pgfqpoint{1.815054in}{1.274780in}}%
\pgfpathlineto{\pgfqpoint{1.828812in}{1.287388in}}%
\pgfpathlineto{\pgfqpoint{1.842582in}{1.300007in}}%
\pgfpathlineto{\pgfqpoint{1.832537in}{1.317433in}}%
\pgfpathlineto{\pgfqpoint{1.819435in}{1.334434in}}%
\pgfpathlineto{\pgfqpoint{1.806244in}{1.321339in}}%
\pgfpathlineto{\pgfqpoint{1.793060in}{1.308251in}}%
\pgfpathlineto{\pgfqpoint{1.805534in}{1.291718in}}%
\pgfpathlineto{\pgfqpoint{1.815054in}{1.274780in}}%
\pgfpathclose%
\pgfusepath{fill}%
\end{pgfscope}%
\begin{pgfscope}%
\pgfpathrectangle{\pgfqpoint{0.407592in}{0.222904in}}{\pgfqpoint{1.411163in}{1.411163in}}%
\pgfusepath{clip}%
\pgfsetbuttcap%
\pgfsetroundjoin%
\definecolor{currentfill}{rgb}{0.749020,0.772549,0.792157}%
\pgfsetfillcolor{currentfill}%
\pgfsetfillopacity{0.160000}%
\pgfsetlinewidth{0.000000pt}%
\definecolor{currentstroke}{rgb}{0.000000,0.000000,0.000000}%
\pgfsetstrokecolor{currentstroke}%
\pgfsetdash{}{0pt}%
\pgfpathmoveto{\pgfqpoint{0.666651in}{0.998060in}}%
\pgfpathlineto{\pgfqpoint{0.653208in}{1.010358in}}%
\pgfpathlineto{\pgfqpoint{0.639752in}{1.022669in}}%
\pgfpathlineto{\pgfqpoint{0.634453in}{1.010958in}}%
\pgfpathlineto{\pgfqpoint{0.631105in}{0.999127in}}%
\pgfpathlineto{\pgfqpoint{0.644918in}{0.987436in}}%
\pgfpathlineto{\pgfqpoint{0.658712in}{0.975761in}}%
\pgfpathlineto{\pgfqpoint{0.661762in}{0.986965in}}%
\pgfpathlineto{\pgfqpoint{0.666651in}{0.998060in}}%
\pgfpathclose%
\pgfusepath{fill}%
\end{pgfscope}%
\begin{pgfscope}%
\pgfpathrectangle{\pgfqpoint{0.407592in}{0.222904in}}{\pgfqpoint{1.411163in}{1.411163in}}%
\pgfusepath{clip}%
\pgfsetbuttcap%
\pgfsetroundjoin%
\definecolor{currentfill}{rgb}{0.749020,0.772549,0.792157}%
\pgfsetfillcolor{currentfill}%
\pgfsetfillopacity{0.160000}%
\pgfsetlinewidth{0.000000pt}%
\definecolor{currentstroke}{rgb}{0.000000,0.000000,0.000000}%
\pgfsetstrokecolor{currentstroke}%
\pgfsetdash{}{0pt}%
\pgfpathmoveto{\pgfqpoint{0.850322in}{0.800848in}}%
\pgfpathlineto{\pgfqpoint{0.836764in}{0.811604in}}%
\pgfpathlineto{\pgfqpoint{0.823182in}{0.822378in}}%
\pgfpathlineto{\pgfqpoint{0.825158in}{0.815159in}}%
\pgfpathlineto{\pgfqpoint{0.828311in}{0.808016in}}%
\pgfpathlineto{\pgfqpoint{0.841686in}{0.797931in}}%
\pgfpathlineto{\pgfqpoint{0.855032in}{0.787868in}}%
\pgfpathlineto{\pgfqpoint{0.852150in}{0.794322in}}%
\pgfpathlineto{\pgfqpoint{0.850322in}{0.800848in}}%
\pgfpathclose%
\pgfusepath{fill}%
\end{pgfscope}%
\begin{pgfscope}%
\pgfpathrectangle{\pgfqpoint{0.407592in}{0.222904in}}{\pgfqpoint{1.411163in}{1.411163in}}%
\pgfusepath{clip}%
\pgfsetbuttcap%
\pgfsetroundjoin%
\definecolor{currentfill}{rgb}{0.749020,0.772549,0.792157}%
\pgfsetfillcolor{currentfill}%
\pgfsetfillopacity{0.160000}%
\pgfsetlinewidth{0.000000pt}%
\definecolor{currentstroke}{rgb}{0.000000,0.000000,0.000000}%
\pgfsetstrokecolor{currentstroke}%
\pgfsetdash{}{0pt}%
\pgfpathmoveto{\pgfqpoint{1.035494in}{0.660885in}}%
\pgfpathlineto{\pgfqpoint{1.028522in}{0.667298in}}%
\pgfpathlineto{\pgfqpoint{1.021521in}{0.673739in}}%
\pgfpathlineto{\pgfqpoint{1.029291in}{0.672046in}}%
\pgfpathlineto{\pgfqpoint{1.037338in}{0.670583in}}%
\pgfpathlineto{\pgfqpoint{1.042732in}{0.664459in}}%
\pgfpathlineto{\pgfqpoint{1.048102in}{0.658362in}}%
\pgfpathlineto{\pgfqpoint{1.041689in}{0.659531in}}%
\pgfpathlineto{\pgfqpoint{1.035494in}{0.660885in}}%
\pgfpathclose%
\pgfusepath{fill}%
\end{pgfscope}%
\begin{pgfscope}%
\pgfpathrectangle{\pgfqpoint{0.407592in}{0.222904in}}{\pgfqpoint{1.411163in}{1.411163in}}%
\pgfusepath{clip}%
\pgfsetbuttcap%
\pgfsetroundjoin%
\definecolor{currentfill}{rgb}{0.749020,0.772549,0.792157}%
\pgfsetfillcolor{currentfill}%
\pgfsetfillopacity{0.160000}%
\pgfsetlinewidth{0.000000pt}%
\definecolor{currentstroke}{rgb}{0.000000,0.000000,0.000000}%
\pgfsetstrokecolor{currentstroke}%
\pgfsetdash{}{0pt}%
\pgfpathmoveto{\pgfqpoint{0.423077in}{1.278906in}}%
\pgfpathlineto{\pgfqpoint{0.410072in}{1.291934in}}%
\pgfpathlineto{\pgfqpoint{0.397060in}{1.304970in}}%
\pgfpathlineto{\pgfqpoint{0.381066in}{1.288350in}}%
\pgfpathlineto{\pgfqpoint{0.367948in}{1.271338in}}%
\pgfpathlineto{\pgfqpoint{0.381662in}{1.258792in}}%
\pgfpathlineto{\pgfqpoint{0.395362in}{1.246258in}}%
\pgfpathlineto{\pgfqpoint{0.407830in}{1.262769in}}%
\pgfpathlineto{\pgfqpoint{0.423077in}{1.278906in}}%
\pgfpathclose%
\pgfusepath{fill}%
\end{pgfscope}%
\begin{pgfscope}%
\pgfpathrectangle{\pgfqpoint{0.407592in}{0.222904in}}{\pgfqpoint{1.411163in}{1.411163in}}%
\pgfusepath{clip}%
\pgfsetbuttcap%
\pgfsetroundjoin%
\definecolor{currentfill}{rgb}{0.749020,0.772549,0.792157}%
\pgfsetfillcolor{currentfill}%
\pgfsetfillopacity{0.160000}%
\pgfsetlinewidth{0.000000pt}%
\definecolor{currentstroke}{rgb}{0.000000,0.000000,0.000000}%
\pgfsetstrokecolor{currentstroke}%
\pgfsetdash{}{0pt}%
\pgfpathmoveto{\pgfqpoint{1.406741in}{0.867403in}}%
\pgfpathlineto{\pgfqpoint{1.420222in}{0.878344in}}%
\pgfpathlineto{\pgfqpoint{1.433727in}{0.889305in}}%
\pgfpathlineto{\pgfqpoint{1.436292in}{0.898613in}}%
\pgfpathlineto{\pgfqpoint{1.437326in}{0.907915in}}%
\pgfpathlineto{\pgfqpoint{1.423753in}{0.896287in}}%
\pgfpathlineto{\pgfqpoint{1.410198in}{0.884673in}}%
\pgfpathlineto{\pgfqpoint{1.409175in}{0.876039in}}%
\pgfpathlineto{\pgfqpoint{1.406741in}{0.867403in}}%
\pgfpathclose%
\pgfusepath{fill}%
\end{pgfscope}%
\begin{pgfscope}%
\pgfpathrectangle{\pgfqpoint{0.407592in}{0.222904in}}{\pgfqpoint{1.411163in}{1.411163in}}%
\pgfusepath{clip}%
\pgfsetbuttcap%
\pgfsetroundjoin%
\definecolor{currentfill}{rgb}{0.749020,0.772549,0.792157}%
\pgfsetfillcolor{currentfill}%
\pgfsetfillopacity{0.160000}%
\pgfsetlinewidth{0.000000pt}%
\definecolor{currentstroke}{rgb}{0.000000,0.000000,0.000000}%
\pgfsetstrokecolor{currentstroke}%
\pgfsetdash{}{0pt}%
\pgfpathmoveto{\pgfqpoint{1.842582in}{1.300007in}}%
\pgfpathlineto{\pgfqpoint{1.856365in}{1.312638in}}%
\pgfpathlineto{\pgfqpoint{1.846053in}{1.330303in}}%
\pgfpathlineto{\pgfqpoint{1.832631in}{1.347535in}}%
\pgfpathlineto{\pgfqpoint{1.819435in}{1.334434in}}%
\pgfpathlineto{\pgfqpoint{1.832537in}{1.317433in}}%
\pgfusepath{fill}%
\end{pgfscope}%
\begin{pgfscope}%
\pgfpathrectangle{\pgfqpoint{0.407592in}{0.222904in}}{\pgfqpoint{1.411163in}{1.411163in}}%
\pgfusepath{clip}%
\pgfsetbuttcap%
\pgfsetroundjoin%
\definecolor{currentfill}{rgb}{0.749020,0.772549,0.792157}%
\pgfsetfillcolor{currentfill}%
\pgfsetfillopacity{0.160000}%
\pgfsetlinewidth{0.000000pt}%
\definecolor{currentstroke}{rgb}{0.000000,0.000000,0.000000}%
\pgfsetstrokecolor{currentstroke}%
\pgfsetdash{}{0pt}%
\pgfpathmoveto{\pgfqpoint{0.982547in}{0.688214in}}%
\pgfpathlineto{\pgfqpoint{0.971578in}{0.696127in}}%
\pgfpathlineto{\pgfqpoint{0.960572in}{0.704067in}}%
\pgfpathlineto{\pgfqpoint{0.967097in}{0.700624in}}%
\pgfpathlineto{\pgfqpoint{0.974180in}{0.697382in}}%
\pgfpathlineto{\pgfqpoint{0.984048in}{0.689996in}}%
\pgfpathlineto{\pgfqpoint{0.993880in}{0.682638in}}%
\pgfpathlineto{\pgfqpoint{0.987984in}{0.685342in}}%
\pgfpathlineto{\pgfqpoint{0.982547in}{0.688214in}}%
\pgfpathclose%
\pgfusepath{fill}%
\end{pgfscope}%
\begin{pgfscope}%
\pgfpathrectangle{\pgfqpoint{0.407592in}{0.222904in}}{\pgfqpoint{1.411163in}{1.411163in}}%
\pgfusepath{clip}%
\pgfsetbuttcap%
\pgfsetroundjoin%
\definecolor{currentfill}{rgb}{0.749020,0.772549,0.792157}%
\pgfsetfillcolor{currentfill}%
\pgfsetfillopacity{0.160000}%
\pgfsetlinewidth{0.000000pt}%
\definecolor{currentstroke}{rgb}{0.000000,0.000000,0.000000}%
\pgfsetstrokecolor{currentstroke}%
\pgfsetdash{}{0pt}%
\pgfpathmoveto{\pgfqpoint{1.546564in}{1.001508in}}%
\pgfpathlineto{\pgfqpoint{1.560302in}{1.013278in}}%
\pgfpathlineto{\pgfqpoint{1.574058in}{1.025063in}}%
\pgfpathlineto{\pgfqpoint{1.572838in}{1.037491in}}%
\pgfpathlineto{\pgfqpoint{1.569529in}{1.049775in}}%
\pgfpathlineto{\pgfqpoint{1.556005in}{1.037381in}}%
\pgfpathlineto{\pgfqpoint{1.542494in}{1.024999in}}%
\pgfpathlineto{\pgfqpoint{1.545516in}{1.013319in}}%
\pgfpathlineto{\pgfqpoint{1.546564in}{1.001508in}}%
\pgfpathclose%
\pgfusepath{fill}%
\end{pgfscope}%
\begin{pgfscope}%
\pgfpathrectangle{\pgfqpoint{0.407592in}{0.222904in}}{\pgfqpoint{1.411163in}{1.411163in}}%
\pgfusepath{clip}%
\pgfsetbuttcap%
\pgfsetroundjoin%
\definecolor{currentfill}{rgb}{0.749020,0.772549,0.792157}%
\pgfsetfillcolor{currentfill}%
\pgfsetfillopacity{0.160000}%
\pgfsetlinewidth{0.000000pt}%
\definecolor{currentstroke}{rgb}{0.000000,0.000000,0.000000}%
\pgfsetstrokecolor{currentstroke}%
\pgfsetdash{}{0pt}%
\pgfpathmoveto{\pgfqpoint{1.117547in}{0.657564in}}%
\pgfpathlineto{\pgfqpoint{1.121030in}{0.663560in}}%
\pgfpathlineto{\pgfqpoint{1.124529in}{0.669584in}}%
\pgfpathlineto{\pgfqpoint{1.133037in}{0.670858in}}%
\pgfpathlineto{\pgfqpoint{1.141347in}{0.672367in}}%
\pgfpathlineto{\pgfqpoint{1.136136in}{0.666064in}}%
\pgfpathlineto{\pgfqpoint{1.130948in}{0.659788in}}%
\pgfpathlineto{\pgfqpoint{1.124325in}{0.658581in}}%
\pgfpathlineto{\pgfqpoint{1.117547in}{0.657564in}}%
\pgfpathclose%
\pgfusepath{fill}%
\end{pgfscope}%
\begin{pgfscope}%
\pgfpathrectangle{\pgfqpoint{0.407592in}{0.222904in}}{\pgfqpoint{1.411163in}{1.411163in}}%
\pgfusepath{clip}%
\pgfsetbuttcap%
\pgfsetroundjoin%
\definecolor{currentfill}{rgb}{0.749020,0.772549,0.792157}%
\pgfsetfillcolor{currentfill}%
\pgfsetfillopacity{0.160000}%
\pgfsetlinewidth{0.000000pt}%
\definecolor{currentstroke}{rgb}{0.000000,0.000000,0.000000}%
\pgfsetstrokecolor{currentstroke}%
\pgfsetdash{}{0pt}%
\pgfpathmoveto{\pgfqpoint{0.768377in}{0.882943in}}%
\pgfpathlineto{\pgfqpoint{0.754735in}{0.894489in}}%
\pgfpathlineto{\pgfqpoint{0.741074in}{0.906052in}}%
\pgfpathlineto{\pgfqpoint{0.740354in}{0.896745in}}%
\pgfpathlineto{\pgfqpoint{0.741166in}{0.887440in}}%
\pgfpathlineto{\pgfqpoint{0.754898in}{0.876547in}}%
\pgfpathlineto{\pgfqpoint{0.768605in}{0.865674in}}%
\pgfpathlineto{\pgfqpoint{0.767784in}{0.874305in}}%
\pgfpathlineto{\pgfqpoint{0.768377in}{0.882943in}}%
\pgfpathclose%
\pgfusepath{fill}%
\end{pgfscope}%
\begin{pgfscope}%
\pgfpathrectangle{\pgfqpoint{0.407592in}{0.222904in}}{\pgfqpoint{1.411163in}{1.411163in}}%
\pgfusepath{clip}%
\pgfsetbuttcap%
\pgfsetroundjoin%
\definecolor{currentfill}{rgb}{0.749020,0.772549,0.792157}%
\pgfsetfillcolor{currentfill}%
\pgfsetfillopacity{0.160000}%
\pgfsetlinewidth{0.000000pt}%
\definecolor{currentstroke}{rgb}{0.000000,0.000000,0.000000}%
\pgfsetstrokecolor{currentstroke}%
\pgfsetdash{}{0pt}%
\pgfpathmoveto{\pgfqpoint{1.048102in}{0.658362in}}%
\pgfpathlineto{\pgfqpoint{1.042732in}{0.664459in}}%
\pgfpathlineto{\pgfqpoint{1.037338in}{0.670583in}}%
\pgfpathlineto{\pgfqpoint{1.045624in}{0.669358in}}%
\pgfpathlineto{\pgfqpoint{1.054113in}{0.668375in}}%
\pgfpathlineto{\pgfqpoint{1.057798in}{0.662473in}}%
\pgfpathlineto{\pgfqpoint{1.061467in}{0.656598in}}%
\pgfpathlineto{\pgfqpoint{1.054705in}{0.657383in}}%
\pgfpathlineto{\pgfqpoint{1.048102in}{0.658362in}}%
\pgfpathclose%
\pgfusepath{fill}%
\end{pgfscope}%
\begin{pgfscope}%
\pgfpathrectangle{\pgfqpoint{0.407592in}{0.222904in}}{\pgfqpoint{1.411163in}{1.411163in}}%
\pgfusepath{clip}%
\pgfsetbuttcap%
\pgfsetroundjoin%
\definecolor{currentfill}{rgb}{0.749020,0.772549,0.792157}%
\pgfsetfillcolor{currentfill}%
\pgfsetfillopacity{0.160000}%
\pgfsetlinewidth{0.000000pt}%
\definecolor{currentstroke}{rgb}{0.000000,0.000000,0.000000}%
\pgfsetstrokecolor{currentstroke}%
\pgfsetdash{}{0pt}%
\pgfpathmoveto{\pgfqpoint{1.218131in}{0.708442in}}%
\pgfpathlineto{\pgfqpoint{1.229011in}{0.716770in}}%
\pgfpathlineto{\pgfqpoint{1.239925in}{0.725125in}}%
\pgfpathlineto{\pgfqpoint{1.246930in}{0.729579in}}%
\pgfpathlineto{\pgfqpoint{1.235519in}{0.720909in}}%
\pgfpathlineto{\pgfqpoint{1.224142in}{0.712266in}}%
\pgfusepath{fill}%
\end{pgfscope}%
\begin{pgfscope}%
\pgfpathrectangle{\pgfqpoint{0.407592in}{0.222904in}}{\pgfqpoint{1.411163in}{1.411163in}}%
\pgfusepath{clip}%
\pgfsetbuttcap%
\pgfsetroundjoin%
\definecolor{currentfill}{rgb}{0.749020,0.772549,0.792157}%
\pgfsetfillcolor{currentfill}%
\pgfsetfillopacity{0.160000}%
\pgfsetlinewidth{0.000000pt}%
\definecolor{currentstroke}{rgb}{0.000000,0.000000,0.000000}%
\pgfsetstrokecolor{currentstroke}%
\pgfsetdash{}{0pt}%
\pgfpathmoveto{\pgfqpoint{0.397060in}{1.304970in}}%
\pgfpathlineto{\pgfqpoint{0.384042in}{1.318011in}}%
\pgfpathlineto{\pgfqpoint{0.371018in}{1.331060in}}%
\pgfpathlineto{\pgfqpoint{0.354262in}{1.313969in}}%
\pgfpathlineto{\pgfqpoint{0.340483in}{1.296466in}}%
\pgfpathlineto{\pgfqpoint{0.354222in}{1.283896in}}%
\pgfpathlineto{\pgfqpoint{0.367948in}{1.271338in}}%
\pgfpathlineto{\pgfqpoint{0.381066in}{1.288350in}}%
\pgfpathlineto{\pgfqpoint{0.397060in}{1.304970in}}%
\pgfpathclose%
\pgfusepath{fill}%
\end{pgfscope}%
\begin{pgfscope}%
\pgfpathrectangle{\pgfqpoint{0.407592in}{0.222904in}}{\pgfqpoint{1.411163in}{1.411163in}}%
\pgfusepath{clip}%
\pgfsetbuttcap%
\pgfsetroundjoin%
\definecolor{currentfill}{rgb}{0.749020,0.772549,0.792157}%
\pgfsetfillcolor{currentfill}%
\pgfsetfillopacity{0.160000}%
\pgfsetlinewidth{0.000000pt}%
\definecolor{currentstroke}{rgb}{0.000000,0.000000,0.000000}%
\pgfsetstrokecolor{currentstroke}%
\pgfsetdash{}{0pt}%
\pgfpathmoveto{\pgfqpoint{1.171873in}{0.680606in}}%
\pgfpathlineto{\pgfqpoint{1.180209in}{0.687761in}}%
\pgfpathlineto{\pgfqpoint{1.188578in}{0.694944in}}%
\pgfpathlineto{\pgfqpoint{1.196731in}{0.698015in}}%
\pgfpathlineto{\pgfqpoint{1.204397in}{0.701300in}}%
\pgfpathlineto{\pgfqpoint{1.194698in}{0.693589in}}%
\pgfpathlineto{\pgfqpoint{1.185034in}{0.685906in}}%
\pgfpathlineto{\pgfqpoint{1.178654in}{0.683167in}}%
\pgfpathlineto{\pgfqpoint{1.171873in}{0.680606in}}%
\pgfpathclose%
\pgfusepath{fill}%
\end{pgfscope}%
\begin{pgfscope}%
\pgfpathrectangle{\pgfqpoint{0.407592in}{0.222904in}}{\pgfqpoint{1.411163in}{1.411163in}}%
\pgfusepath{clip}%
\pgfsetbuttcap%
\pgfsetroundjoin%
\definecolor{currentfill}{rgb}{0.749020,0.772549,0.792157}%
\pgfsetfillcolor{currentfill}%
\pgfsetfillopacity{0.160000}%
\pgfsetlinewidth{0.000000pt}%
\definecolor{currentstroke}{rgb}{0.000000,0.000000,0.000000}%
\pgfsetstrokecolor{currentstroke}%
\pgfsetdash{}{0pt}%
\pgfpathmoveto{\pgfqpoint{1.103644in}{0.656113in}}%
\pgfpathlineto{\pgfqpoint{1.105356in}{0.661927in}}%
\pgfpathlineto{\pgfqpoint{1.107075in}{0.667768in}}%
\pgfpathlineto{\pgfqpoint{1.115861in}{0.668553in}}%
\pgfpathlineto{\pgfqpoint{1.124529in}{0.669584in}}%
\pgfpathlineto{\pgfqpoint{1.121030in}{0.663560in}}%
\pgfpathlineto{\pgfqpoint{1.117547in}{0.657564in}}%
\pgfpathlineto{\pgfqpoint{1.110643in}{0.656739in}}%
\pgfpathlineto{\pgfqpoint{1.103644in}{0.656113in}}%
\pgfpathclose%
\pgfusepath{fill}%
\end{pgfscope}%
\begin{pgfscope}%
\pgfpathrectangle{\pgfqpoint{0.407592in}{0.222904in}}{\pgfqpoint{1.411163in}{1.411163in}}%
\pgfusepath{clip}%
\pgfsetbuttcap%
\pgfsetroundjoin%
\definecolor{currentfill}{rgb}{0.749020,0.772549,0.792157}%
\pgfsetfillcolor{currentfill}%
\pgfsetfillopacity{0.160000}%
\pgfsetlinewidth{0.000000pt}%
\definecolor{currentstroke}{rgb}{0.000000,0.000000,0.000000}%
\pgfsetstrokecolor{currentstroke}%
\pgfsetdash{}{0pt}%
\pgfpathmoveto{\pgfqpoint{0.639752in}{1.022669in}}%
\pgfpathlineto{\pgfqpoint{0.626284in}{1.034991in}}%
\pgfpathlineto{\pgfqpoint{0.612803in}{1.047324in}}%
\pgfpathlineto{\pgfqpoint{0.607079in}{1.035007in}}%
\pgfpathlineto{\pgfqpoint{0.603422in}{1.022558in}}%
\pgfpathlineto{\pgfqpoint{0.617273in}{1.010834in}}%
\pgfpathlineto{\pgfqpoint{0.631105in}{0.999127in}}%
\pgfpathlineto{\pgfqpoint{0.634453in}{1.010958in}}%
\pgfpathlineto{\pgfqpoint{0.639752in}{1.022669in}}%
\pgfpathclose%
\pgfusepath{fill}%
\end{pgfscope}%
\begin{pgfscope}%
\pgfpathrectangle{\pgfqpoint{0.407592in}{0.222904in}}{\pgfqpoint{1.411163in}{1.411163in}}%
\pgfusepath{clip}%
\pgfsetbuttcap%
\pgfsetroundjoin%
\definecolor{currentfill}{rgb}{0.749020,0.772549,0.792157}%
\pgfsetfillcolor{currentfill}%
\pgfsetfillopacity{0.160000}%
\pgfsetlinewidth{0.000000pt}%
\definecolor{currentstroke}{rgb}{0.000000,0.000000,0.000000}%
\pgfsetstrokecolor{currentstroke}%
\pgfsetdash{}{0pt}%
\pgfpathmoveto{\pgfqpoint{1.283360in}{0.757700in}}%
\pgfpathlineto{\pgfqpoint{1.295700in}{0.767102in}}%
\pgfpathlineto{\pgfqpoint{1.308072in}{0.776530in}}%
\pgfpathlineto{\pgfqpoint{1.314176in}{0.782781in}}%
\pgfpathlineto{\pgfqpoint{1.319268in}{0.789156in}}%
\pgfpathlineto{\pgfqpoint{1.306306in}{0.779047in}}%
\pgfpathlineto{\pgfqpoint{1.293372in}{0.768960in}}%
\pgfpathlineto{\pgfqpoint{1.288814in}{0.763273in}}%
\pgfpathlineto{\pgfqpoint{1.283360in}{0.757700in}}%
\pgfpathclose%
\pgfusepath{fill}%
\end{pgfscope}%
\begin{pgfscope}%
\pgfpathrectangle{\pgfqpoint{0.407592in}{0.222904in}}{\pgfqpoint{1.411163in}{1.411163in}}%
\pgfusepath{clip}%
\pgfsetbuttcap%
\pgfsetroundjoin%
\definecolor{currentfill}{rgb}{0.749020,0.772549,0.792157}%
\pgfsetfillcolor{currentfill}%
\pgfsetfillopacity{0.160000}%
\pgfsetlinewidth{0.000000pt}%
\definecolor{currentstroke}{rgb}{0.000000,0.000000,0.000000}%
\pgfsetstrokecolor{currentstroke}%
\pgfsetdash{}{0pt}%
\pgfpathmoveto{\pgfqpoint{1.061467in}{0.656598in}}%
\pgfpathlineto{\pgfqpoint{1.057798in}{0.662473in}}%
\pgfpathlineto{\pgfqpoint{1.054113in}{0.668375in}}%
\pgfpathlineto{\pgfqpoint{1.062765in}{0.667641in}}%
\pgfpathlineto{\pgfqpoint{1.071541in}{0.667158in}}%
\pgfpathlineto{\pgfqpoint{1.073449in}{0.661378in}}%
\pgfpathlineto{\pgfqpoint{1.075348in}{0.655625in}}%
\pgfpathlineto{\pgfqpoint{1.068358in}{0.656011in}}%
\pgfpathlineto{\pgfqpoint{1.061467in}{0.656598in}}%
\pgfpathclose%
\pgfusepath{fill}%
\end{pgfscope}%
\begin{pgfscope}%
\pgfpathrectangle{\pgfqpoint{0.407592in}{0.222904in}}{\pgfqpoint{1.411163in}{1.411163in}}%
\pgfusepath{clip}%
\pgfsetbuttcap%
\pgfsetroundjoin%
\definecolor{currentfill}{rgb}{0.749020,0.772549,0.792157}%
\pgfsetfillcolor{currentfill}%
\pgfsetfillopacity{0.160000}%
\pgfsetlinewidth{0.000000pt}%
\definecolor{currentstroke}{rgb}{0.000000,0.000000,0.000000}%
\pgfsetstrokecolor{currentstroke}%
\pgfsetdash{}{0pt}%
\pgfpathmoveto{\pgfqpoint{1.089492in}{0.655462in}}%
\pgfpathlineto{\pgfqpoint{1.089398in}{0.661195in}}%
\pgfpathlineto{\pgfqpoint{1.089304in}{0.666954in}}%
\pgfpathlineto{\pgfqpoint{1.098209in}{0.667234in}}%
\pgfpathlineto{\pgfqpoint{1.107075in}{0.667768in}}%
\pgfpathlineto{\pgfqpoint{1.105356in}{0.661927in}}%
\pgfpathlineto{\pgfqpoint{1.103644in}{0.656113in}}%
\pgfpathlineto{\pgfqpoint{1.096584in}{0.655686in}}%
\pgfpathlineto{\pgfqpoint{1.089492in}{0.655462in}}%
\pgfpathclose%
\pgfusepath{fill}%
\end{pgfscope}%
\begin{pgfscope}%
\pgfpathrectangle{\pgfqpoint{0.407592in}{0.222904in}}{\pgfqpoint{1.411163in}{1.411163in}}%
\pgfusepath{clip}%
\pgfsetbuttcap%
\pgfsetroundjoin%
\definecolor{currentfill}{rgb}{0.749020,0.772549,0.792157}%
\pgfsetfillcolor{currentfill}%
\pgfsetfillopacity{0.160000}%
\pgfsetlinewidth{0.000000pt}%
\definecolor{currentstroke}{rgb}{0.000000,0.000000,0.000000}%
\pgfsetstrokecolor{currentstroke}%
\pgfsetdash{}{0pt}%
\pgfpathmoveto{\pgfqpoint{0.371018in}{1.331060in}}%
\pgfpathlineto{\pgfqpoint{0.357986in}{1.344115in}}%
\pgfpathlineto{\pgfqpoint{0.340846in}{1.326792in}}%
\pgfpathlineto{\pgfqpoint{0.326731in}{1.309048in}}%
\pgfpathlineto{\pgfqpoint{0.340483in}{1.296466in}}%
\pgfpathlineto{\pgfqpoint{0.354262in}{1.313969in}}%
\pgfusepath{fill}%
\end{pgfscope}%
\begin{pgfscope}%
\pgfpathrectangle{\pgfqpoint{0.407592in}{0.222904in}}{\pgfqpoint{1.411163in}{1.411163in}}%
\pgfusepath{clip}%
\pgfsetbuttcap%
\pgfsetroundjoin%
\definecolor{currentfill}{rgb}{0.749020,0.772549,0.792157}%
\pgfsetfillcolor{currentfill}%
\pgfsetfillopacity{0.160000}%
\pgfsetlinewidth{0.000000pt}%
\definecolor{currentstroke}{rgb}{0.000000,0.000000,0.000000}%
\pgfsetstrokecolor{currentstroke}%
\pgfsetdash{}{0pt}%
\pgfpathmoveto{\pgfqpoint{1.075348in}{0.655625in}}%
\pgfpathlineto{\pgfqpoint{1.073449in}{0.661378in}}%
\pgfpathlineto{\pgfqpoint{1.071541in}{0.667158in}}%
\pgfpathlineto{\pgfqpoint{1.080401in}{0.666929in}}%
\pgfpathlineto{\pgfqpoint{1.089304in}{0.666954in}}%
\pgfpathlineto{\pgfqpoint{1.089398in}{0.661195in}}%
\pgfpathlineto{\pgfqpoint{1.089492in}{0.655462in}}%
\pgfpathlineto{\pgfqpoint{1.082403in}{0.655442in}}%
\pgfpathlineto{\pgfqpoint{1.075348in}{0.655625in}}%
\pgfpathclose%
\pgfusepath{fill}%
\end{pgfscope}%
\begin{pgfscope}%
\pgfpathrectangle{\pgfqpoint{0.407592in}{0.222904in}}{\pgfqpoint{1.411163in}{1.411163in}}%
\pgfusepath{clip}%
\pgfsetbuttcap%
\pgfsetroundjoin%
\definecolor{currentfill}{rgb}{0.749020,0.772549,0.792157}%
\pgfsetfillcolor{currentfill}%
\pgfsetfillopacity{0.160000}%
\pgfsetlinewidth{0.000000pt}%
\definecolor{currentstroke}{rgb}{0.000000,0.000000,0.000000}%
\pgfsetstrokecolor{currentstroke}%
\pgfsetdash{}{0pt}%
\pgfpathmoveto{\pgfqpoint{1.574058in}{1.025063in}}%
\pgfpathlineto{\pgfqpoint{1.587832in}{1.036865in}}%
\pgfpathlineto{\pgfqpoint{1.601625in}{1.048682in}}%
\pgfpathlineto{\pgfqpoint{1.600222in}{1.061717in}}%
\pgfpathlineto{\pgfqpoint{1.596613in}{1.074595in}}%
\pgfpathlineto{\pgfqpoint{1.583065in}{1.062179in}}%
\pgfpathlineto{\pgfqpoint{1.569529in}{1.049775in}}%
\pgfpathlineto{\pgfqpoint{1.572838in}{1.037491in}}%
\pgfpathlineto{\pgfqpoint{1.574058in}{1.025063in}}%
\pgfpathclose%
\pgfusepath{fill}%
\end{pgfscope}%
\begin{pgfscope}%
\pgfpathrectangle{\pgfqpoint{0.407592in}{0.222904in}}{\pgfqpoint{1.411163in}{1.411163in}}%
\pgfusepath{clip}%
\pgfsetbuttcap%
\pgfsetroundjoin%
\definecolor{currentfill}{rgb}{0.749020,0.772549,0.792157}%
\pgfsetfillcolor{currentfill}%
\pgfsetfillopacity{0.160000}%
\pgfsetlinewidth{0.000000pt}%
\definecolor{currentstroke}{rgb}{0.000000,0.000000,0.000000}%
\pgfsetstrokecolor{currentstroke}%
\pgfsetdash{}{0pt}%
\pgfpathmoveto{\pgfqpoint{0.881633in}{0.767811in}}%
\pgfpathlineto{\pgfqpoint{0.868348in}{0.777829in}}%
\pgfpathlineto{\pgfqpoint{0.855032in}{0.787868in}}%
\pgfpathlineto{\pgfqpoint{0.858962in}{0.781516in}}%
\pgfpathlineto{\pgfqpoint{0.863928in}{0.775293in}}%
\pgfpathlineto{\pgfqpoint{0.876785in}{0.765933in}}%
\pgfpathlineto{\pgfqpoint{0.889607in}{0.756597in}}%
\pgfpathlineto{\pgfqpoint{0.885162in}{0.762145in}}%
\pgfpathlineto{\pgfqpoint{0.881633in}{0.767811in}}%
\pgfpathclose%
\pgfusepath{fill}%
\end{pgfscope}%
\begin{pgfscope}%
\pgfpathrectangle{\pgfqpoint{0.407592in}{0.222904in}}{\pgfqpoint{1.411163in}{1.411163in}}%
\pgfusepath{clip}%
\pgfsetbuttcap%
\pgfsetroundjoin%
\definecolor{currentfill}{rgb}{0.749020,0.772549,0.792157}%
\pgfsetfillcolor{currentfill}%
\pgfsetfillopacity{0.160000}%
\pgfsetlinewidth{0.000000pt}%
\definecolor{currentstroke}{rgb}{0.000000,0.000000,0.000000}%
\pgfsetstrokecolor{currentstroke}%
\pgfsetdash{}{0pt}%
\pgfpathmoveto{\pgfqpoint{1.345279in}{0.809441in}}%
\pgfpathlineto{\pgfqpoint{1.358327in}{0.819617in}}%
\pgfpathlineto{\pgfqpoint{1.371404in}{0.829816in}}%
\pgfpathlineto{\pgfqpoint{1.376269in}{0.837661in}}%
\pgfpathlineto{\pgfqpoint{1.379850in}{0.845580in}}%
\pgfpathlineto{\pgfqpoint{1.366440in}{0.834697in}}%
\pgfpathlineto{\pgfqpoint{1.353055in}{0.823833in}}%
\pgfpathlineto{\pgfqpoint{1.349747in}{0.816602in}}%
\pgfpathlineto{\pgfqpoint{1.345279in}{0.809441in}}%
\pgfpathclose%
\pgfusepath{fill}%
\end{pgfscope}%
\begin{pgfscope}%
\pgfpathrectangle{\pgfqpoint{0.407592in}{0.222904in}}{\pgfqpoint{1.411163in}{1.411163in}}%
\pgfusepath{clip}%
\pgfsetbuttcap%
\pgfsetroundjoin%
\definecolor{currentfill}{rgb}{0.749020,0.772549,0.792157}%
\pgfsetfillcolor{currentfill}%
\pgfsetfillopacity{0.160000}%
\pgfsetlinewidth{0.000000pt}%
\definecolor{currentstroke}{rgb}{0.000000,0.000000,0.000000}%
\pgfsetstrokecolor{currentstroke}%
\pgfsetdash{}{0pt}%
\pgfpathmoveto{\pgfqpoint{0.949302in}{0.711487in}}%
\pgfpathlineto{\pgfqpoint{0.937326in}{0.720066in}}%
\pgfpathlineto{\pgfqpoint{0.925313in}{0.728671in}}%
\pgfpathlineto{\pgfqpoint{0.931521in}{0.724253in}}%
\pgfpathlineto{\pgfqpoint{0.938446in}{0.720028in}}%
\pgfpathlineto{\pgfqpoint{0.949528in}{0.712034in}}%
\pgfpathlineto{\pgfqpoint{0.960572in}{0.704067in}}%
\pgfpathlineto{\pgfqpoint{0.954632in}{0.707694in}}%
\pgfpathlineto{\pgfqpoint{0.949302in}{0.711487in}}%
\pgfpathclose%
\pgfusepath{fill}%
\end{pgfscope}%
\begin{pgfscope}%
\pgfpathrectangle{\pgfqpoint{0.407592in}{0.222904in}}{\pgfqpoint{1.411163in}{1.411163in}}%
\pgfusepath{clip}%
\pgfsetbuttcap%
\pgfsetroundjoin%
\definecolor{currentfill}{rgb}{0.749020,0.772549,0.792157}%
\pgfsetfillcolor{currentfill}%
\pgfsetfillopacity{0.160000}%
\pgfsetlinewidth{0.000000pt}%
\definecolor{currentstroke}{rgb}{0.000000,0.000000,0.000000}%
\pgfsetstrokecolor{currentstroke}%
\pgfsetdash{}{0pt}%
\pgfpathmoveto{\pgfqpoint{1.433727in}{0.889305in}}%
\pgfpathlineto{\pgfqpoint{1.447257in}{0.900285in}}%
\pgfpathlineto{\pgfqpoint{1.460810in}{0.911284in}}%
\pgfpathlineto{\pgfqpoint{1.463494in}{0.921257in}}%
\pgfpathlineto{\pgfqpoint{1.464526in}{0.931220in}}%
\pgfpathlineto{\pgfqpoint{1.450917in}{0.919560in}}%
\pgfpathlineto{\pgfqpoint{1.437326in}{0.907915in}}%
\pgfpathlineto{\pgfqpoint{1.436292in}{0.898613in}}%
\pgfpathlineto{\pgfqpoint{1.433727in}{0.889305in}}%
\pgfpathclose%
\pgfusepath{fill}%
\end{pgfscope}%
\begin{pgfscope}%
\pgfpathrectangle{\pgfqpoint{0.407592in}{0.222904in}}{\pgfqpoint{1.411163in}{1.411163in}}%
\pgfusepath{clip}%
\pgfsetbuttcap%
\pgfsetroundjoin%
\definecolor{currentfill}{rgb}{0.749020,0.772549,0.792157}%
\pgfsetfillcolor{currentfill}%
\pgfsetfillopacity{0.160000}%
\pgfsetlinewidth{0.000000pt}%
\definecolor{currentstroke}{rgb}{0.000000,0.000000,0.000000}%
\pgfsetstrokecolor{currentstroke}%
\pgfsetdash{}{0pt}%
\pgfpathmoveto{\pgfqpoint{1.246930in}{0.729579in}}%
\pgfpathlineto{\pgfqpoint{1.258376in}{0.738274in}}%
\pgfpathlineto{\pgfqpoint{1.269856in}{0.746995in}}%
\pgfpathlineto{\pgfqpoint{1.277033in}{0.752265in}}%
\pgfpathlineto{\pgfqpoint{1.283360in}{0.757700in}}%
\pgfpathlineto{\pgfqpoint{1.271053in}{0.748321in}}%
\pgfpathlineto{\pgfqpoint{1.258778in}{0.738968in}}%
\pgfpathlineto{\pgfqpoint{1.253223in}{0.734200in}}%
\pgfpathlineto{\pgfqpoint{1.246930in}{0.729579in}}%
\pgfpathclose%
\pgfusepath{fill}%
\end{pgfscope}%
\begin{pgfscope}%
\pgfpathrectangle{\pgfqpoint{0.407592in}{0.222904in}}{\pgfqpoint{1.411163in}{1.411163in}}%
\pgfusepath{clip}%
\pgfsetbuttcap%
\pgfsetroundjoin%
\definecolor{currentfill}{rgb}{0.749020,0.772549,0.792157}%
\pgfsetfillcolor{currentfill}%
\pgfsetfillopacity{0.160000}%
\pgfsetlinewidth{0.000000pt}%
\definecolor{currentstroke}{rgb}{0.000000,0.000000,0.000000}%
\pgfsetstrokecolor{currentstroke}%
\pgfsetdash{}{0pt}%
\pgfpathmoveto{\pgfqpoint{0.993880in}{0.682638in}}%
\pgfpathlineto{\pgfqpoint{0.984048in}{0.689996in}}%
\pgfpathlineto{\pgfqpoint{0.974180in}{0.697382in}}%
\pgfpathlineto{\pgfqpoint{0.981789in}{0.694354in}}%
\pgfpathlineto{\pgfqpoint{0.989891in}{0.691556in}}%
\pgfpathlineto{\pgfqpoint{0.998437in}{0.684656in}}%
\pgfpathlineto{\pgfqpoint{1.006948in}{0.677783in}}%
\pgfpathlineto{\pgfqpoint{1.000211in}{0.680115in}}%
\pgfpathlineto{\pgfqpoint{0.993880in}{0.682638in}}%
\pgfpathclose%
\pgfusepath{fill}%
\end{pgfscope}%
\begin{pgfscope}%
\pgfpathrectangle{\pgfqpoint{0.407592in}{0.222904in}}{\pgfqpoint{1.411163in}{1.411163in}}%
\pgfusepath{clip}%
\pgfsetbuttcap%
\pgfsetroundjoin%
\definecolor{currentfill}{rgb}{0.749020,0.772549,0.792157}%
\pgfsetfillcolor{currentfill}%
\pgfsetfillopacity{0.160000}%
\pgfsetlinewidth{0.000000pt}%
\definecolor{currentstroke}{rgb}{0.000000,0.000000,0.000000}%
\pgfsetstrokecolor{currentstroke}%
\pgfsetdash{}{0pt}%
\pgfpathmoveto{\pgfqpoint{0.823182in}{0.822378in}}%
\pgfpathlineto{\pgfqpoint{0.809575in}{0.833173in}}%
\pgfpathlineto{\pgfqpoint{0.795943in}{0.843987in}}%
\pgfpathlineto{\pgfqpoint{0.798057in}{0.836080in}}%
\pgfpathlineto{\pgfqpoint{0.801469in}{0.828254in}}%
\pgfpathlineto{\pgfqpoint{0.814905in}{0.818123in}}%
\pgfpathlineto{\pgfqpoint{0.828311in}{0.808016in}}%
\pgfpathlineto{\pgfqpoint{0.825158in}{0.815159in}}%
\pgfpathlineto{\pgfqpoint{0.823182in}{0.822378in}}%
\pgfpathclose%
\pgfusepath{fill}%
\end{pgfscope}%
\begin{pgfscope}%
\pgfpathrectangle{\pgfqpoint{0.407592in}{0.222904in}}{\pgfqpoint{1.411163in}{1.411163in}}%
\pgfusepath{clip}%
\pgfsetbuttcap%
\pgfsetroundjoin%
\definecolor{currentfill}{rgb}{0.749020,0.772549,0.792157}%
\pgfsetfillcolor{currentfill}%
\pgfsetfillopacity{0.160000}%
\pgfsetlinewidth{0.000000pt}%
\definecolor{currentstroke}{rgb}{0.000000,0.000000,0.000000}%
\pgfsetstrokecolor{currentstroke}%
\pgfsetdash{}{0pt}%
\pgfpathmoveto{\pgfqpoint{0.915150in}{0.738000in}}%
\pgfpathlineto{\pgfqpoint{0.902396in}{0.747287in}}%
\pgfpathlineto{\pgfqpoint{0.889607in}{0.756597in}}%
\pgfpathlineto{\pgfqpoint{0.894953in}{0.751194in}}%
\pgfpathlineto{\pgfqpoint{0.901178in}{0.745960in}}%
\pgfpathlineto{\pgfqpoint{0.913264in}{0.737302in}}%
\pgfpathlineto{\pgfqpoint{0.925313in}{0.728671in}}%
\pgfpathlineto{\pgfqpoint{0.919849in}{0.733260in}}%
\pgfpathlineto{\pgfqpoint{0.915150in}{0.738000in}}%
\pgfpathclose%
\pgfusepath{fill}%
\end{pgfscope}%
\begin{pgfscope}%
\pgfpathrectangle{\pgfqpoint{0.407592in}{0.222904in}}{\pgfqpoint{1.411163in}{1.411163in}}%
\pgfusepath{clip}%
\pgfsetbuttcap%
\pgfsetroundjoin%
\definecolor{currentfill}{rgb}{0.749020,0.772549,0.792157}%
\pgfsetfillcolor{currentfill}%
\pgfsetfillopacity{0.160000}%
\pgfsetlinewidth{0.000000pt}%
\definecolor{currentstroke}{rgb}{0.000000,0.000000,0.000000}%
\pgfsetstrokecolor{currentstroke}%
\pgfsetdash{}{0pt}%
\pgfpathmoveto{\pgfqpoint{0.612803in}{1.047324in}}%
\pgfpathlineto{\pgfqpoint{0.599309in}{1.059669in}}%
\pgfpathlineto{\pgfqpoint{0.585803in}{1.072026in}}%
\pgfpathlineto{\pgfqpoint{0.579642in}{1.059112in}}%
\pgfpathlineto{\pgfqpoint{0.575662in}{1.046054in}}%
\pgfpathlineto{\pgfqpoint{0.589552in}{1.034297in}}%
\pgfpathlineto{\pgfqpoint{0.603422in}{1.022558in}}%
\pgfpathlineto{\pgfqpoint{0.607079in}{1.035007in}}%
\pgfpathlineto{\pgfqpoint{0.612803in}{1.047324in}}%
\pgfpathclose%
\pgfusepath{fill}%
\end{pgfscope}%
\begin{pgfscope}%
\pgfpathrectangle{\pgfqpoint{0.407592in}{0.222904in}}{\pgfqpoint{1.411163in}{1.411163in}}%
\pgfusepath{clip}%
\pgfsetbuttcap%
\pgfsetroundjoin%
\definecolor{currentfill}{rgb}{0.749020,0.772549,0.792157}%
\pgfsetfillcolor{currentfill}%
\pgfsetfillopacity{0.160000}%
\pgfsetlinewidth{0.000000pt}%
\definecolor{currentstroke}{rgb}{0.000000,0.000000,0.000000}%
\pgfsetstrokecolor{currentstroke}%
\pgfsetdash{}{0pt}%
\pgfpathmoveto{\pgfqpoint{0.741074in}{0.906052in}}%
\pgfpathlineto{\pgfqpoint{0.727395in}{0.917630in}}%
\pgfpathlineto{\pgfqpoint{0.713696in}{0.929224in}}%
\pgfpathlineto{\pgfqpoint{0.712835in}{0.919256in}}%
\pgfpathlineto{\pgfqpoint{0.713628in}{0.909286in}}%
\pgfpathlineto{\pgfqpoint{0.727410in}{0.898353in}}%
\pgfpathlineto{\pgfqpoint{0.741166in}{0.887440in}}%
\pgfpathlineto{\pgfqpoint{0.740354in}{0.896745in}}%
\pgfpathlineto{\pgfqpoint{0.741074in}{0.906052in}}%
\pgfpathclose%
\pgfusepath{fill}%
\end{pgfscope}%
\begin{pgfscope}%
\pgfpathrectangle{\pgfqpoint{0.407592in}{0.222904in}}{\pgfqpoint{1.411163in}{1.411163in}}%
\pgfusepath{clip}%
\pgfsetbuttcap%
\pgfsetroundjoin%
\definecolor{currentfill}{rgb}{0.749020,0.772549,0.792157}%
\pgfsetfillcolor{currentfill}%
\pgfsetfillopacity{0.160000}%
\pgfsetlinewidth{0.000000pt}%
\definecolor{currentstroke}{rgb}{0.000000,0.000000,0.000000}%
\pgfsetstrokecolor{currentstroke}%
\pgfsetdash{}{0pt}%
\pgfpathmoveto{\pgfqpoint{1.157223in}{0.676065in}}%
\pgfpathlineto{\pgfqpoint{1.164074in}{0.682765in}}%
\pgfpathlineto{\pgfqpoint{1.170952in}{0.689494in}}%
\pgfpathlineto{\pgfqpoint{1.179972in}{0.692099in}}%
\pgfpathlineto{\pgfqpoint{1.188578in}{0.694944in}}%
\pgfpathlineto{\pgfqpoint{1.180209in}{0.687761in}}%
\pgfpathlineto{\pgfqpoint{1.171873in}{0.680606in}}%
\pgfpathlineto{\pgfqpoint{1.164718in}{0.678235in}}%
\pgfpathlineto{\pgfqpoint{1.157223in}{0.676065in}}%
\pgfpathclose%
\pgfusepath{fill}%
\end{pgfscope}%
\begin{pgfscope}%
\pgfpathrectangle{\pgfqpoint{0.407592in}{0.222904in}}{\pgfqpoint{1.411163in}{1.411163in}}%
\pgfusepath{clip}%
\pgfsetbuttcap%
\pgfsetroundjoin%
\definecolor{currentfill}{rgb}{0.749020,0.772549,0.792157}%
\pgfsetfillcolor{currentfill}%
\pgfsetfillopacity{0.160000}%
\pgfsetlinewidth{0.000000pt}%
\definecolor{currentstroke}{rgb}{0.000000,0.000000,0.000000}%
\pgfsetstrokecolor{currentstroke}%
\pgfsetdash{}{0pt}%
\pgfpathmoveto{\pgfqpoint{1.204397in}{0.701300in}}%
\pgfpathlineto{\pgfqpoint{1.214129in}{0.709037in}}%
\pgfpathlineto{\pgfqpoint{1.223897in}{0.716803in}}%
\pgfpathlineto{\pgfqpoint{1.232236in}{0.720859in}}%
\pgfpathlineto{\pgfqpoint{1.239925in}{0.725125in}}%
\pgfpathlineto{\pgfqpoint{1.229011in}{0.716770in}}%
\pgfpathlineto{\pgfqpoint{1.218131in}{0.708442in}}%
\pgfpathlineto{\pgfqpoint{1.211540in}{0.704780in}}%
\pgfpathlineto{\pgfqpoint{1.204397in}{0.701300in}}%
\pgfpathclose%
\pgfusepath{fill}%
\end{pgfscope}%
\begin{pgfscope}%
\pgfpathrectangle{\pgfqpoint{0.407592in}{0.222904in}}{\pgfqpoint{1.411163in}{1.411163in}}%
\pgfusepath{clip}%
\pgfsetbuttcap%
\pgfsetroundjoin%
\definecolor{currentfill}{rgb}{0.749020,0.772549,0.792157}%
\pgfsetfillcolor{currentfill}%
\pgfsetfillopacity{0.160000}%
\pgfsetlinewidth{0.000000pt}%
\definecolor{currentstroke}{rgb}{0.000000,0.000000,0.000000}%
\pgfsetstrokecolor{currentstroke}%
\pgfsetdash{}{0pt}%
\pgfpathmoveto{\pgfqpoint{1.601625in}{1.048682in}}%
\pgfpathlineto{\pgfqpoint{1.615437in}{1.060516in}}%
\pgfpathlineto{\pgfqpoint{1.629267in}{1.072365in}}%
\pgfpathlineto{\pgfqpoint{1.627667in}{1.085997in}}%
\pgfpathlineto{\pgfqpoint{1.623745in}{1.099460in}}%
\pgfpathlineto{\pgfqpoint{1.610173in}{1.087022in}}%
\pgfpathlineto{\pgfqpoint{1.596613in}{1.074595in}}%
\pgfpathlineto{\pgfqpoint{1.600222in}{1.061717in}}%
\pgfpathlineto{\pgfqpoint{1.601625in}{1.048682in}}%
\pgfpathclose%
\pgfusepath{fill}%
\end{pgfscope}%
\begin{pgfscope}%
\pgfpathrectangle{\pgfqpoint{0.407592in}{0.222904in}}{\pgfqpoint{1.411163in}{1.411163in}}%
\pgfusepath{clip}%
\pgfsetbuttcap%
\pgfsetroundjoin%
\definecolor{currentfill}{rgb}{0.749020,0.772549,0.792157}%
\pgfsetfillcolor{currentfill}%
\pgfsetfillopacity{0.160000}%
\pgfsetlinewidth{0.000000pt}%
\definecolor{currentstroke}{rgb}{0.000000,0.000000,0.000000}%
\pgfsetstrokecolor{currentstroke}%
\pgfsetdash{}{0pt}%
\pgfpathmoveto{\pgfqpoint{1.460810in}{0.911284in}}%
\pgfpathlineto{\pgfqpoint{1.474388in}{0.922304in}}%
\pgfpathlineto{\pgfqpoint{1.487990in}{0.933343in}}%
\pgfpathlineto{\pgfqpoint{1.490781in}{0.943972in}}%
\pgfpathlineto{\pgfqpoint{1.491799in}{0.954586in}}%
\pgfpathlineto{\pgfqpoint{1.478153in}{0.942895in}}%
\pgfpathlineto{\pgfqpoint{1.464526in}{0.931220in}}%
\pgfpathlineto{\pgfqpoint{1.463494in}{0.921257in}}%
\pgfpathlineto{\pgfqpoint{1.460810in}{0.911284in}}%
\pgfpathclose%
\pgfusepath{fill}%
\end{pgfscope}%
\begin{pgfscope}%
\pgfpathrectangle{\pgfqpoint{0.407592in}{0.222904in}}{\pgfqpoint{1.411163in}{1.411163in}}%
\pgfusepath{clip}%
\pgfsetbuttcap%
\pgfsetroundjoin%
\definecolor{currentfill}{rgb}{0.749020,0.772549,0.792157}%
\pgfsetfillcolor{currentfill}%
\pgfsetfillopacity{0.160000}%
\pgfsetlinewidth{0.000000pt}%
\definecolor{currentstroke}{rgb}{0.000000,0.000000,0.000000}%
\pgfsetstrokecolor{currentstroke}%
\pgfsetdash{}{0pt}%
\pgfpathmoveto{\pgfqpoint{1.006948in}{0.677783in}}%
\pgfpathlineto{\pgfqpoint{0.998437in}{0.684656in}}%
\pgfpathlineto{\pgfqpoint{0.989891in}{0.691556in}}%
\pgfpathlineto{\pgfqpoint{0.998451in}{0.689002in}}%
\pgfpathlineto{\pgfqpoint{1.007429in}{0.686702in}}%
\pgfpathlineto{\pgfqpoint{1.014490in}{0.680207in}}%
\pgfpathlineto{\pgfqpoint{1.021521in}{0.673739in}}%
\pgfpathlineto{\pgfqpoint{1.014062in}{0.675654in}}%
\pgfpathlineto{\pgfqpoint{1.006948in}{0.677783in}}%
\pgfpathclose%
\pgfusepath{fill}%
\end{pgfscope}%
\begin{pgfscope}%
\pgfpathrectangle{\pgfqpoint{0.407592in}{0.222904in}}{\pgfqpoint{1.411163in}{1.411163in}}%
\pgfusepath{clip}%
\pgfsetbuttcap%
\pgfsetroundjoin%
\definecolor{currentfill}{rgb}{0.749020,0.772549,0.792157}%
\pgfsetfillcolor{currentfill}%
\pgfsetfillopacity{0.160000}%
\pgfsetlinewidth{0.000000pt}%
\definecolor{currentstroke}{rgb}{0.000000,0.000000,0.000000}%
\pgfsetstrokecolor{currentstroke}%
\pgfsetdash{}{0pt}%
\pgfpathmoveto{\pgfqpoint{0.585803in}{1.072026in}}%
\pgfpathlineto{\pgfqpoint{0.572283in}{1.084395in}}%
\pgfpathlineto{\pgfqpoint{0.558751in}{1.096775in}}%
\pgfpathlineto{\pgfqpoint{0.552141in}{1.083274in}}%
\pgfpathlineto{\pgfqpoint{0.547824in}{1.069615in}}%
\pgfpathlineto{\pgfqpoint{0.561752in}{1.057826in}}%
\pgfpathlineto{\pgfqpoint{0.575662in}{1.046054in}}%
\pgfpathlineto{\pgfqpoint{0.579642in}{1.059112in}}%
\pgfpathlineto{\pgfqpoint{0.585803in}{1.072026in}}%
\pgfpathclose%
\pgfusepath{fill}%
\end{pgfscope}%
\begin{pgfscope}%
\pgfpathrectangle{\pgfqpoint{0.407592in}{0.222904in}}{\pgfqpoint{1.411163in}{1.411163in}}%
\pgfusepath{clip}%
\pgfsetbuttcap%
\pgfsetroundjoin%
\definecolor{currentfill}{rgb}{0.749020,0.772549,0.792157}%
\pgfsetfillcolor{currentfill}%
\pgfsetfillopacity{0.160000}%
\pgfsetlinewidth{0.000000pt}%
\definecolor{currentstroke}{rgb}{0.000000,0.000000,0.000000}%
\pgfsetstrokecolor{currentstroke}%
\pgfsetdash{}{0pt}%
\pgfpathmoveto{\pgfqpoint{1.371404in}{0.829816in}}%
\pgfpathlineto{\pgfqpoint{1.384511in}{0.840037in}}%
\pgfpathlineto{\pgfqpoint{1.397646in}{0.850281in}}%
\pgfpathlineto{\pgfqpoint{1.402896in}{0.858804in}}%
\pgfpathlineto{\pgfqpoint{1.406741in}{0.867403in}}%
\pgfpathlineto{\pgfqpoint{1.393284in}{0.856482in}}%
\pgfpathlineto{\pgfqpoint{1.379850in}{0.845580in}}%
\pgfpathlineto{\pgfqpoint{1.376269in}{0.837661in}}%
\pgfpathlineto{\pgfqpoint{1.371404in}{0.829816in}}%
\pgfpathclose%
\pgfusepath{fill}%
\end{pgfscope}%
\begin{pgfscope}%
\pgfpathrectangle{\pgfqpoint{0.407592in}{0.222904in}}{\pgfqpoint{1.411163in}{1.411163in}}%
\pgfusepath{clip}%
\pgfsetbuttcap%
\pgfsetroundjoin%
\definecolor{currentfill}{rgb}{0.749020,0.772549,0.792157}%
\pgfsetfillcolor{currentfill}%
\pgfsetfillopacity{0.160000}%
\pgfsetlinewidth{0.000000pt}%
\definecolor{currentstroke}{rgb}{0.000000,0.000000,0.000000}%
\pgfsetstrokecolor{currentstroke}%
\pgfsetdash{}{0pt}%
\pgfpathmoveto{\pgfqpoint{1.308072in}{0.776530in}}%
\pgfpathlineto{\pgfqpoint{1.320476in}{0.785982in}}%
\pgfpathlineto{\pgfqpoint{1.332913in}{0.795459in}}%
\pgfpathlineto{\pgfqpoint{1.339661in}{0.802383in}}%
\pgfpathlineto{\pgfqpoint{1.345279in}{0.809441in}}%
\pgfpathlineto{\pgfqpoint{1.332259in}{0.799287in}}%
\pgfpathlineto{\pgfqpoint{1.319268in}{0.789156in}}%
\pgfpathlineto{\pgfqpoint{1.314176in}{0.782781in}}%
\pgfpathlineto{\pgfqpoint{1.308072in}{0.776530in}}%
\pgfpathclose%
\pgfusepath{fill}%
\end{pgfscope}%
\begin{pgfscope}%
\pgfpathrectangle{\pgfqpoint{0.407592in}{0.222904in}}{\pgfqpoint{1.411163in}{1.411163in}}%
\pgfusepath{clip}%
\pgfsetbuttcap%
\pgfsetroundjoin%
\definecolor{currentfill}{rgb}{0.749020,0.772549,0.792157}%
\pgfsetfillcolor{currentfill}%
\pgfsetfillopacity{0.160000}%
\pgfsetlinewidth{0.000000pt}%
\definecolor{currentstroke}{rgb}{0.000000,0.000000,0.000000}%
\pgfsetstrokecolor{currentstroke}%
\pgfsetdash{}{0pt}%
\pgfpathmoveto{\pgfqpoint{1.629267in}{1.072365in}}%
\pgfpathlineto{\pgfqpoint{1.643116in}{1.084231in}}%
\pgfpathlineto{\pgfqpoint{1.656983in}{1.096112in}}%
\pgfpathlineto{\pgfqpoint{1.655174in}{1.110332in}}%
\pgfpathlineto{\pgfqpoint{1.650927in}{1.124370in}}%
\pgfpathlineto{\pgfqpoint{1.637330in}{1.111909in}}%
\pgfpathlineto{\pgfqpoint{1.623745in}{1.099460in}}%
\pgfpathlineto{\pgfqpoint{1.627667in}{1.085997in}}%
\pgfpathlineto{\pgfqpoint{1.629267in}{1.072365in}}%
\pgfpathclose%
\pgfusepath{fill}%
\end{pgfscope}%
\begin{pgfscope}%
\pgfpathrectangle{\pgfqpoint{0.407592in}{0.222904in}}{\pgfqpoint{1.411163in}{1.411163in}}%
\pgfusepath{clip}%
\pgfsetbuttcap%
\pgfsetroundjoin%
\definecolor{currentfill}{rgb}{0.749020,0.772549,0.792157}%
\pgfsetfillcolor{currentfill}%
\pgfsetfillopacity{0.160000}%
\pgfsetlinewidth{0.000000pt}%
\definecolor{currentstroke}{rgb}{0.000000,0.000000,0.000000}%
\pgfsetstrokecolor{currentstroke}%
\pgfsetdash{}{0pt}%
\pgfpathmoveto{\pgfqpoint{1.141347in}{0.672367in}}%
\pgfpathlineto{\pgfqpoint{1.146580in}{0.678697in}}%
\pgfpathlineto{\pgfqpoint{1.151835in}{0.685055in}}%
\pgfpathlineto{\pgfqpoint{1.161559in}{0.687142in}}%
\pgfpathlineto{\pgfqpoint{1.170952in}{0.689494in}}%
\pgfpathlineto{\pgfqpoint{1.164074in}{0.682765in}}%
\pgfpathlineto{\pgfqpoint{1.157223in}{0.676065in}}%
\pgfpathlineto{\pgfqpoint{1.149421in}{0.674106in}}%
\pgfpathlineto{\pgfqpoint{1.141347in}{0.672367in}}%
\pgfpathclose%
\pgfusepath{fill}%
\end{pgfscope}%
\begin{pgfscope}%
\pgfpathrectangle{\pgfqpoint{0.407592in}{0.222904in}}{\pgfqpoint{1.411163in}{1.411163in}}%
\pgfusepath{clip}%
\pgfsetbuttcap%
\pgfsetroundjoin%
\definecolor{currentfill}{rgb}{0.749020,0.772549,0.792157}%
\pgfsetfillcolor{currentfill}%
\pgfsetfillopacity{0.160000}%
\pgfsetlinewidth{0.000000pt}%
\definecolor{currentstroke}{rgb}{0.000000,0.000000,0.000000}%
\pgfsetstrokecolor{currentstroke}%
\pgfsetdash{}{0pt}%
\pgfpathmoveto{\pgfqpoint{0.960572in}{0.704067in}}%
\pgfpathlineto{\pgfqpoint{0.949528in}{0.712034in}}%
\pgfpathlineto{\pgfqpoint{0.938446in}{0.720028in}}%
\pgfpathlineto{\pgfqpoint{0.946061in}{0.716016in}}%
\pgfpathlineto{\pgfqpoint{0.954332in}{0.712235in}}%
\pgfpathlineto{\pgfqpoint{0.964274in}{0.704795in}}%
\pgfpathlineto{\pgfqpoint{0.974180in}{0.697382in}}%
\pgfpathlineto{\pgfqpoint{0.967097in}{0.700624in}}%
\pgfpathlineto{\pgfqpoint{0.960572in}{0.704067in}}%
\pgfpathclose%
\pgfusepath{fill}%
\end{pgfscope}%
\begin{pgfscope}%
\pgfpathrectangle{\pgfqpoint{0.407592in}{0.222904in}}{\pgfqpoint{1.411163in}{1.411163in}}%
\pgfusepath{clip}%
\pgfsetbuttcap%
\pgfsetroundjoin%
\definecolor{currentfill}{rgb}{0.749020,0.772549,0.792157}%
\pgfsetfillcolor{currentfill}%
\pgfsetfillopacity{0.160000}%
\pgfsetlinewidth{0.000000pt}%
\definecolor{currentstroke}{rgb}{0.000000,0.000000,0.000000}%
\pgfsetstrokecolor{currentstroke}%
\pgfsetdash{}{0pt}%
\pgfpathmoveto{\pgfqpoint{1.239925in}{0.725125in}}%
\pgfpathlineto{\pgfqpoint{1.250874in}{0.733506in}}%
\pgfpathlineto{\pgfqpoint{1.261858in}{0.741915in}}%
\pgfpathlineto{\pgfqpoint{1.269856in}{0.746995in}}%
\pgfpathlineto{\pgfqpoint{1.258376in}{0.738274in}}%
\pgfpathlineto{\pgfqpoint{1.246930in}{0.729579in}}%
\pgfusepath{fill}%
\end{pgfscope}%
\begin{pgfscope}%
\pgfpathrectangle{\pgfqpoint{0.407592in}{0.222904in}}{\pgfqpoint{1.411163in}{1.411163in}}%
\pgfusepath{clip}%
\pgfsetbuttcap%
\pgfsetroundjoin%
\definecolor{currentfill}{rgb}{0.749020,0.772549,0.792157}%
\pgfsetfillcolor{currentfill}%
\pgfsetfillopacity{0.160000}%
\pgfsetlinewidth{0.000000pt}%
\definecolor{currentstroke}{rgb}{0.000000,0.000000,0.000000}%
\pgfsetstrokecolor{currentstroke}%
\pgfsetdash{}{0pt}%
\pgfpathmoveto{\pgfqpoint{0.713696in}{0.929224in}}%
\pgfpathlineto{\pgfqpoint{0.699979in}{0.940834in}}%
\pgfpathlineto{\pgfqpoint{0.686242in}{0.952460in}}%
\pgfpathlineto{\pgfqpoint{0.685228in}{0.941840in}}%
\pgfpathlineto{\pgfqpoint{0.685989in}{0.931212in}}%
\pgfpathlineto{\pgfqpoint{0.699821in}{0.920239in}}%
\pgfpathlineto{\pgfqpoint{0.713628in}{0.909286in}}%
\pgfpathlineto{\pgfqpoint{0.712835in}{0.919256in}}%
\pgfpathlineto{\pgfqpoint{0.713696in}{0.929224in}}%
\pgfpathclose%
\pgfusepath{fill}%
\end{pgfscope}%
\begin{pgfscope}%
\pgfpathrectangle{\pgfqpoint{0.407592in}{0.222904in}}{\pgfqpoint{1.411163in}{1.411163in}}%
\pgfusepath{clip}%
\pgfsetbuttcap%
\pgfsetroundjoin%
\definecolor{currentfill}{rgb}{0.749020,0.772549,0.792157}%
\pgfsetfillcolor{currentfill}%
\pgfsetfillopacity{0.160000}%
\pgfsetlinewidth{0.000000pt}%
\definecolor{currentstroke}{rgb}{0.000000,0.000000,0.000000}%
\pgfsetstrokecolor{currentstroke}%
\pgfsetdash{}{0pt}%
\pgfpathmoveto{\pgfqpoint{0.855032in}{0.787868in}}%
\pgfpathlineto{\pgfqpoint{0.841686in}{0.797931in}}%
\pgfpathlineto{\pgfqpoint{0.828311in}{0.808016in}}%
\pgfpathlineto{\pgfqpoint{0.832633in}{0.800982in}}%
\pgfpathlineto{\pgfqpoint{0.838112in}{0.794089in}}%
\pgfpathlineto{\pgfqpoint{0.851037in}{0.784679in}}%
\pgfpathlineto{\pgfqpoint{0.863928in}{0.775293in}}%
\pgfpathlineto{\pgfqpoint{0.858962in}{0.781516in}}%
\pgfpathlineto{\pgfqpoint{0.855032in}{0.787868in}}%
\pgfpathclose%
\pgfusepath{fill}%
\end{pgfscope}%
\begin{pgfscope}%
\pgfpathrectangle{\pgfqpoint{0.407592in}{0.222904in}}{\pgfqpoint{1.411163in}{1.411163in}}%
\pgfusepath{clip}%
\pgfsetbuttcap%
\pgfsetroundjoin%
\definecolor{currentfill}{rgb}{0.749020,0.772549,0.792157}%
\pgfsetfillcolor{currentfill}%
\pgfsetfillopacity{0.160000}%
\pgfsetlinewidth{0.000000pt}%
\definecolor{currentstroke}{rgb}{0.000000,0.000000,0.000000}%
\pgfsetstrokecolor{currentstroke}%
\pgfsetdash{}{0pt}%
\pgfpathmoveto{\pgfqpoint{0.795943in}{0.843987in}}%
\pgfpathlineto{\pgfqpoint{0.782286in}{0.854821in}}%
\pgfpathlineto{\pgfqpoint{0.768605in}{0.865674in}}%
\pgfpathlineto{\pgfqpoint{0.770845in}{0.857087in}}%
\pgfpathlineto{\pgfqpoint{0.774506in}{0.848584in}}%
\pgfpathlineto{\pgfqpoint{0.788003in}{0.838407in}}%
\pgfpathlineto{\pgfqpoint{0.801469in}{0.828254in}}%
\pgfpathlineto{\pgfqpoint{0.798057in}{0.836080in}}%
\pgfpathlineto{\pgfqpoint{0.795943in}{0.843987in}}%
\pgfpathclose%
\pgfusepath{fill}%
\end{pgfscope}%
\begin{pgfscope}%
\pgfpathrectangle{\pgfqpoint{0.407592in}{0.222904in}}{\pgfqpoint{1.411163in}{1.411163in}}%
\pgfusepath{clip}%
\pgfsetbuttcap%
\pgfsetroundjoin%
\definecolor{currentfill}{rgb}{0.749020,0.772549,0.792157}%
\pgfsetfillcolor{currentfill}%
\pgfsetfillopacity{0.160000}%
\pgfsetlinewidth{0.000000pt}%
\definecolor{currentstroke}{rgb}{0.000000,0.000000,0.000000}%
\pgfsetstrokecolor{currentstroke}%
\pgfsetdash{}{0pt}%
\pgfpathmoveto{\pgfqpoint{1.021521in}{0.673739in}}%
\pgfpathlineto{\pgfqpoint{1.014490in}{0.680207in}}%
\pgfpathlineto{\pgfqpoint{1.007429in}{0.686702in}}%
\pgfpathlineto{\pgfqpoint{1.016786in}{0.684669in}}%
\pgfpathlineto{\pgfqpoint{1.026478in}{0.682912in}}%
\pgfpathlineto{\pgfqpoint{1.031920in}{0.676734in}}%
\pgfpathlineto{\pgfqpoint{1.037338in}{0.670583in}}%
\pgfpathlineto{\pgfqpoint{1.029291in}{0.672046in}}%
\pgfpathlineto{\pgfqpoint{1.021521in}{0.673739in}}%
\pgfpathclose%
\pgfusepath{fill}%
\end{pgfscope}%
\begin{pgfscope}%
\pgfpathrectangle{\pgfqpoint{0.407592in}{0.222904in}}{\pgfqpoint{1.411163in}{1.411163in}}%
\pgfusepath{clip}%
\pgfsetbuttcap%
\pgfsetroundjoin%
\definecolor{currentfill}{rgb}{0.749020,0.772549,0.792157}%
\pgfsetfillcolor{currentfill}%
\pgfsetfillopacity{0.160000}%
\pgfsetlinewidth{0.000000pt}%
\definecolor{currentstroke}{rgb}{0.000000,0.000000,0.000000}%
\pgfsetstrokecolor{currentstroke}%
\pgfsetdash{}{0pt}%
\pgfpathmoveto{\pgfqpoint{1.269856in}{0.746995in}}%
\pgfpathlineto{\pgfqpoint{1.281371in}{0.755743in}}%
\pgfpathlineto{\pgfqpoint{1.292921in}{0.764517in}}%
\pgfpathlineto{\pgfqpoint{1.300978in}{0.770432in}}%
\pgfpathlineto{\pgfqpoint{1.308072in}{0.776530in}}%
\pgfpathlineto{\pgfqpoint{1.295700in}{0.767102in}}%
\pgfpathlineto{\pgfqpoint{1.283360in}{0.757700in}}%
\pgfpathlineto{\pgfqpoint{1.277033in}{0.752265in}}%
\pgfpathlineto{\pgfqpoint{1.269856in}{0.746995in}}%
\pgfpathclose%
\pgfusepath{fill}%
\end{pgfscope}%
\begin{pgfscope}%
\pgfpathrectangle{\pgfqpoint{0.407592in}{0.222904in}}{\pgfqpoint{1.411163in}{1.411163in}}%
\pgfusepath{clip}%
\pgfsetbuttcap%
\pgfsetroundjoin%
\definecolor{currentfill}{rgb}{0.749020,0.772549,0.792157}%
\pgfsetfillcolor{currentfill}%
\pgfsetfillopacity{0.160000}%
\pgfsetlinewidth{0.000000pt}%
\definecolor{currentstroke}{rgb}{0.000000,0.000000,0.000000}%
\pgfsetstrokecolor{currentstroke}%
\pgfsetdash{}{0pt}%
\pgfpathmoveto{\pgfqpoint{0.558751in}{1.096775in}}%
\pgfpathlineto{\pgfqpoint{0.545207in}{1.109167in}}%
\pgfpathlineto{\pgfqpoint{0.531649in}{1.121571in}}%
\pgfpathlineto{\pgfqpoint{0.524575in}{1.107492in}}%
\pgfpathlineto{\pgfqpoint{0.519908in}{1.093242in}}%
\pgfpathlineto{\pgfqpoint{0.533876in}{1.081420in}}%
\pgfpathlineto{\pgfqpoint{0.547824in}{1.069615in}}%
\pgfpathlineto{\pgfqpoint{0.552141in}{1.083274in}}%
\pgfpathlineto{\pgfqpoint{0.558751in}{1.096775in}}%
\pgfpathclose%
\pgfusepath{fill}%
\end{pgfscope}%
\begin{pgfscope}%
\pgfpathrectangle{\pgfqpoint{0.407592in}{0.222904in}}{\pgfqpoint{1.411163in}{1.411163in}}%
\pgfusepath{clip}%
\pgfsetbuttcap%
\pgfsetroundjoin%
\definecolor{currentfill}{rgb}{0.749020,0.772549,0.792157}%
\pgfsetfillcolor{currentfill}%
\pgfsetfillopacity{0.160000}%
\pgfsetlinewidth{0.000000pt}%
\definecolor{currentstroke}{rgb}{0.000000,0.000000,0.000000}%
\pgfsetstrokecolor{currentstroke}%
\pgfsetdash{}{0pt}%
\pgfpathmoveto{\pgfqpoint{1.188578in}{0.694944in}}%
\pgfpathlineto{\pgfqpoint{1.196979in}{0.702154in}}%
\pgfpathlineto{\pgfqpoint{1.205412in}{0.709392in}}%
\pgfpathlineto{\pgfqpoint{1.214942in}{0.712974in}}%
\pgfpathlineto{\pgfqpoint{1.223897in}{0.716803in}}%
\pgfpathlineto{\pgfqpoint{1.214129in}{0.709037in}}%
\pgfpathlineto{\pgfqpoint{1.204397in}{0.701300in}}%
\pgfpathlineto{\pgfqpoint{1.196731in}{0.698015in}}%
\pgfpathlineto{\pgfqpoint{1.188578in}{0.694944in}}%
\pgfpathclose%
\pgfusepath{fill}%
\end{pgfscope}%
\begin{pgfscope}%
\pgfpathrectangle{\pgfqpoint{0.407592in}{0.222904in}}{\pgfqpoint{1.411163in}{1.411163in}}%
\pgfusepath{clip}%
\pgfsetbuttcap%
\pgfsetroundjoin%
\definecolor{currentfill}{rgb}{0.749020,0.772549,0.792157}%
\pgfsetfillcolor{currentfill}%
\pgfsetfillopacity{0.160000}%
\pgfsetlinewidth{0.000000pt}%
\definecolor{currentstroke}{rgb}{0.000000,0.000000,0.000000}%
\pgfsetstrokecolor{currentstroke}%
\pgfsetdash{}{0pt}%
\pgfpathmoveto{\pgfqpoint{1.124529in}{0.669584in}}%
\pgfpathlineto{\pgfqpoint{1.128043in}{0.675635in}}%
\pgfpathlineto{\pgfqpoint{1.131573in}{0.681712in}}%
\pgfpathlineto{\pgfqpoint{1.141825in}{0.683242in}}%
\pgfpathlineto{\pgfqpoint{1.151835in}{0.685055in}}%
\pgfpathlineto{\pgfqpoint{1.146580in}{0.678697in}}%
\pgfpathlineto{\pgfqpoint{1.141347in}{0.672367in}}%
\pgfpathlineto{\pgfqpoint{1.133037in}{0.670858in}}%
\pgfpathlineto{\pgfqpoint{1.124529in}{0.669584in}}%
\pgfpathclose%
\pgfusepath{fill}%
\end{pgfscope}%
\begin{pgfscope}%
\pgfpathrectangle{\pgfqpoint{0.407592in}{0.222904in}}{\pgfqpoint{1.411163in}{1.411163in}}%
\pgfusepath{clip}%
\pgfsetbuttcap%
\pgfsetroundjoin%
\definecolor{currentfill}{rgb}{0.749020,0.772549,0.792157}%
\pgfsetfillcolor{currentfill}%
\pgfsetfillopacity{0.160000}%
\pgfsetlinewidth{0.000000pt}%
\definecolor{currentstroke}{rgb}{0.000000,0.000000,0.000000}%
\pgfsetstrokecolor{currentstroke}%
\pgfsetdash{}{0pt}%
\pgfpathmoveto{\pgfqpoint{1.656983in}{1.096112in}}%
\pgfpathlineto{\pgfqpoint{1.670870in}{1.108009in}}%
\pgfpathlineto{\pgfqpoint{1.684775in}{1.119923in}}%
\pgfpathlineto{\pgfqpoint{1.682743in}{1.134722in}}%
\pgfpathlineto{\pgfqpoint{1.678158in}{1.149325in}}%
\pgfpathlineto{\pgfqpoint{1.664536in}{1.136842in}}%
\pgfpathlineto{\pgfqpoint{1.650927in}{1.124370in}}%
\pgfpathlineto{\pgfqpoint{1.655174in}{1.110332in}}%
\pgfpathlineto{\pgfqpoint{1.656983in}{1.096112in}}%
\pgfpathclose%
\pgfusepath{fill}%
\end{pgfscope}%
\begin{pgfscope}%
\pgfpathrectangle{\pgfqpoint{0.407592in}{0.222904in}}{\pgfqpoint{1.411163in}{1.411163in}}%
\pgfusepath{clip}%
\pgfsetbuttcap%
\pgfsetroundjoin%
\definecolor{currentfill}{rgb}{0.749020,0.772549,0.792157}%
\pgfsetfillcolor{currentfill}%
\pgfsetfillopacity{0.160000}%
\pgfsetlinewidth{0.000000pt}%
\definecolor{currentstroke}{rgb}{0.000000,0.000000,0.000000}%
\pgfsetstrokecolor{currentstroke}%
\pgfsetdash{}{0pt}%
\pgfpathmoveto{\pgfqpoint{0.925313in}{0.728671in}}%
\pgfpathlineto{\pgfqpoint{0.913264in}{0.737302in}}%
\pgfpathlineto{\pgfqpoint{0.901178in}{0.745960in}}%
\pgfpathlineto{\pgfqpoint{0.908260in}{0.740920in}}%
\pgfpathlineto{\pgfqpoint{0.916170in}{0.736098in}}%
\pgfpathlineto{\pgfqpoint{0.927327in}{0.728050in}}%
\pgfpathlineto{\pgfqpoint{0.938446in}{0.720028in}}%
\pgfpathlineto{\pgfqpoint{0.931521in}{0.724253in}}%
\pgfpathlineto{\pgfqpoint{0.925313in}{0.728671in}}%
\pgfpathclose%
\pgfusepath{fill}%
\end{pgfscope}%
\begin{pgfscope}%
\pgfpathrectangle{\pgfqpoint{0.407592in}{0.222904in}}{\pgfqpoint{1.411163in}{1.411163in}}%
\pgfusepath{clip}%
\pgfsetbuttcap%
\pgfsetroundjoin%
\definecolor{currentfill}{rgb}{0.749020,0.772549,0.792157}%
\pgfsetfillcolor{currentfill}%
\pgfsetfillopacity{0.160000}%
\pgfsetlinewidth{0.000000pt}%
\definecolor{currentstroke}{rgb}{0.000000,0.000000,0.000000}%
\pgfsetstrokecolor{currentstroke}%
\pgfsetdash{}{0pt}%
\pgfpathmoveto{\pgfqpoint{1.487990in}{0.933343in}}%
\pgfpathlineto{\pgfqpoint{1.501616in}{0.944401in}}%
\pgfpathlineto{\pgfqpoint{1.515267in}{0.955480in}}%
\pgfpathlineto{\pgfqpoint{1.518154in}{0.966759in}}%
\pgfpathlineto{\pgfqpoint{1.519145in}{0.978016in}}%
\pgfpathlineto{\pgfqpoint{1.505463in}{0.966293in}}%
\pgfpathlineto{\pgfqpoint{1.491799in}{0.954586in}}%
\pgfpathlineto{\pgfqpoint{1.490781in}{0.943972in}}%
\pgfpathlineto{\pgfqpoint{1.487990in}{0.933343in}}%
\pgfpathclose%
\pgfusepath{fill}%
\end{pgfscope}%
\begin{pgfscope}%
\pgfpathrectangle{\pgfqpoint{0.407592in}{0.222904in}}{\pgfqpoint{1.411163in}{1.411163in}}%
\pgfusepath{clip}%
\pgfsetbuttcap%
\pgfsetroundjoin%
\definecolor{currentfill}{rgb}{0.749020,0.772549,0.792157}%
\pgfsetfillcolor{currentfill}%
\pgfsetfillopacity{0.160000}%
\pgfsetlinewidth{0.000000pt}%
\definecolor{currentstroke}{rgb}{0.000000,0.000000,0.000000}%
\pgfsetstrokecolor{currentstroke}%
\pgfsetdash{}{0pt}%
\pgfpathmoveto{\pgfqpoint{0.889607in}{0.756597in}}%
\pgfpathlineto{\pgfqpoint{0.876785in}{0.765933in}}%
\pgfpathlineto{\pgfqpoint{0.863928in}{0.775293in}}%
\pgfpathlineto{\pgfqpoint{0.869913in}{0.769230in}}%
\pgfpathlineto{\pgfqpoint{0.876896in}{0.763355in}}%
\pgfpathlineto{\pgfqpoint{0.889056in}{0.754644in}}%
\pgfpathlineto{\pgfqpoint{0.901178in}{0.745960in}}%
\pgfpathlineto{\pgfqpoint{0.894953in}{0.751194in}}%
\pgfpathlineto{\pgfqpoint{0.889607in}{0.756597in}}%
\pgfpathclose%
\pgfusepath{fill}%
\end{pgfscope}%
\begin{pgfscope}%
\pgfpathrectangle{\pgfqpoint{0.407592in}{0.222904in}}{\pgfqpoint{1.411163in}{1.411163in}}%
\pgfusepath{clip}%
\pgfsetbuttcap%
\pgfsetroundjoin%
\definecolor{currentfill}{rgb}{0.749020,0.772549,0.792157}%
\pgfsetfillcolor{currentfill}%
\pgfsetfillopacity{0.160000}%
\pgfsetlinewidth{0.000000pt}%
\definecolor{currentstroke}{rgb}{0.000000,0.000000,0.000000}%
\pgfsetstrokecolor{currentstroke}%
\pgfsetdash{}{0pt}%
\pgfpathmoveto{\pgfqpoint{1.037338in}{0.670583in}}%
\pgfpathlineto{\pgfqpoint{1.031920in}{0.676734in}}%
\pgfpathlineto{\pgfqpoint{1.026478in}{0.682912in}}%
\pgfpathlineto{\pgfqpoint{1.036462in}{0.681440in}}%
\pgfpathlineto{\pgfqpoint{1.046692in}{0.680260in}}%
\pgfpathlineto{\pgfqpoint{1.050411in}{0.674304in}}%
\pgfpathlineto{\pgfqpoint{1.054113in}{0.668375in}}%
\pgfpathlineto{\pgfqpoint{1.045624in}{0.669358in}}%
\pgfpathlineto{\pgfqpoint{1.037338in}{0.670583in}}%
\pgfpathclose%
\pgfusepath{fill}%
\end{pgfscope}%
\begin{pgfscope}%
\pgfpathrectangle{\pgfqpoint{0.407592in}{0.222904in}}{\pgfqpoint{1.411163in}{1.411163in}}%
\pgfusepath{clip}%
\pgfsetbuttcap%
\pgfsetroundjoin%
\definecolor{currentfill}{rgb}{0.749020,0.772549,0.792157}%
\pgfsetfillcolor{currentfill}%
\pgfsetfillopacity{0.160000}%
\pgfsetlinewidth{0.000000pt}%
\definecolor{currentstroke}{rgb}{0.000000,0.000000,0.000000}%
\pgfsetstrokecolor{currentstroke}%
\pgfsetdash{}{0pt}%
\pgfpathmoveto{\pgfqpoint{1.107075in}{0.667768in}}%
\pgfpathlineto{\pgfqpoint{1.108802in}{0.673636in}}%
\pgfpathlineto{\pgfqpoint{1.110537in}{0.679531in}}%
\pgfpathlineto{\pgfqpoint{1.121128in}{0.680473in}}%
\pgfpathlineto{\pgfqpoint{1.131573in}{0.681712in}}%
\pgfpathlineto{\pgfqpoint{1.128043in}{0.675635in}}%
\pgfpathlineto{\pgfqpoint{1.124529in}{0.669584in}}%
\pgfpathlineto{\pgfqpoint{1.115861in}{0.668553in}}%
\pgfpathlineto{\pgfqpoint{1.107075in}{0.667768in}}%
\pgfpathclose%
\pgfusepath{fill}%
\end{pgfscope}%
\begin{pgfscope}%
\pgfpathrectangle{\pgfqpoint{0.407592in}{0.222904in}}{\pgfqpoint{1.411163in}{1.411163in}}%
\pgfusepath{clip}%
\pgfsetbuttcap%
\pgfsetroundjoin%
\definecolor{currentfill}{rgb}{0.749020,0.772549,0.792157}%
\pgfsetfillcolor{currentfill}%
\pgfsetfillopacity{0.160000}%
\pgfsetlinewidth{0.000000pt}%
\definecolor{currentstroke}{rgb}{0.000000,0.000000,0.000000}%
\pgfsetstrokecolor{currentstroke}%
\pgfsetdash{}{0pt}%
\pgfpathmoveto{\pgfqpoint{0.686242in}{0.952460in}}%
\pgfpathlineto{\pgfqpoint{0.672487in}{0.964103in}}%
\pgfpathlineto{\pgfqpoint{0.658712in}{0.975761in}}%
\pgfpathlineto{\pgfqpoint{0.657532in}{0.964496in}}%
\pgfpathlineto{\pgfqpoint{0.658248in}{0.953219in}}%
\pgfpathlineto{\pgfqpoint{0.672131in}{0.942205in}}%
\pgfpathlineto{\pgfqpoint{0.685989in}{0.931212in}}%
\pgfpathlineto{\pgfqpoint{0.685228in}{0.941840in}}%
\pgfpathlineto{\pgfqpoint{0.686242in}{0.952460in}}%
\pgfpathclose%
\pgfusepath{fill}%
\end{pgfscope}%
\begin{pgfscope}%
\pgfpathrectangle{\pgfqpoint{0.407592in}{0.222904in}}{\pgfqpoint{1.411163in}{1.411163in}}%
\pgfusepath{clip}%
\pgfsetbuttcap%
\pgfsetroundjoin%
\definecolor{currentfill}{rgb}{0.749020,0.772549,0.792157}%
\pgfsetfillcolor{currentfill}%
\pgfsetfillopacity{0.160000}%
\pgfsetlinewidth{0.000000pt}%
\definecolor{currentstroke}{rgb}{0.000000,0.000000,0.000000}%
\pgfsetstrokecolor{currentstroke}%
\pgfsetdash{}{0pt}%
\pgfpathmoveto{\pgfqpoint{1.397646in}{0.850281in}}%
\pgfpathlineto{\pgfqpoint{1.410810in}{0.860548in}}%
\pgfpathlineto{\pgfqpoint{1.424004in}{0.870837in}}%
\pgfpathlineto{\pgfqpoint{1.429629in}{0.880032in}}%
\pgfpathlineto{\pgfqpoint{1.433727in}{0.889305in}}%
\pgfpathlineto{\pgfqpoint{1.420222in}{0.878344in}}%
\pgfpathlineto{\pgfqpoint{1.406741in}{0.867403in}}%
\pgfpathlineto{\pgfqpoint{1.402896in}{0.858804in}}%
\pgfpathlineto{\pgfqpoint{1.397646in}{0.850281in}}%
\pgfpathclose%
\pgfusepath{fill}%
\end{pgfscope}%
\begin{pgfscope}%
\pgfpathrectangle{\pgfqpoint{0.407592in}{0.222904in}}{\pgfqpoint{1.411163in}{1.411163in}}%
\pgfusepath{clip}%
\pgfsetbuttcap%
\pgfsetroundjoin%
\definecolor{currentfill}{rgb}{0.749020,0.772549,0.792157}%
\pgfsetfillcolor{currentfill}%
\pgfsetfillopacity{0.160000}%
\pgfsetlinewidth{0.000000pt}%
\definecolor{currentstroke}{rgb}{0.000000,0.000000,0.000000}%
\pgfsetstrokecolor{currentstroke}%
\pgfsetdash{}{0pt}%
\pgfpathmoveto{\pgfqpoint{0.974180in}{0.697382in}}%
\pgfpathlineto{\pgfqpoint{0.964274in}{0.704795in}}%
\pgfpathlineto{\pgfqpoint{0.954332in}{0.712235in}}%
\pgfpathlineto{\pgfqpoint{0.963224in}{0.708704in}}%
\pgfpathlineto{\pgfqpoint{0.972697in}{0.705441in}}%
\pgfpathlineto{\pgfqpoint{0.981312in}{0.698485in}}%
\pgfpathlineto{\pgfqpoint{0.989891in}{0.691556in}}%
\pgfpathlineto{\pgfqpoint{0.981789in}{0.694354in}}%
\pgfpathlineto{\pgfqpoint{0.974180in}{0.697382in}}%
\pgfpathclose%
\pgfusepath{fill}%
\end{pgfscope}%
\begin{pgfscope}%
\pgfpathrectangle{\pgfqpoint{0.407592in}{0.222904in}}{\pgfqpoint{1.411163in}{1.411163in}}%
\pgfusepath{clip}%
\pgfsetbuttcap%
\pgfsetroundjoin%
\definecolor{currentfill}{rgb}{0.749020,0.772549,0.792157}%
\pgfsetfillcolor{currentfill}%
\pgfsetfillopacity{0.160000}%
\pgfsetlinewidth{0.000000pt}%
\definecolor{currentstroke}{rgb}{0.000000,0.000000,0.000000}%
\pgfsetstrokecolor{currentstroke}%
\pgfsetdash{}{0pt}%
\pgfpathmoveto{\pgfqpoint{0.531649in}{1.121571in}}%
\pgfpathlineto{\pgfqpoint{0.518078in}{1.133987in}}%
\pgfpathlineto{\pgfqpoint{0.504495in}{1.146414in}}%
\pgfpathlineto{\pgfqpoint{0.496945in}{1.131767in}}%
\pgfpathlineto{\pgfqpoint{0.491914in}{1.116936in}}%
\pgfpathlineto{\pgfqpoint{0.505921in}{1.105081in}}%
\pgfpathlineto{\pgfqpoint{0.519908in}{1.093242in}}%
\pgfpathlineto{\pgfqpoint{0.524575in}{1.107492in}}%
\pgfpathlineto{\pgfqpoint{0.531649in}{1.121571in}}%
\pgfpathclose%
\pgfusepath{fill}%
\end{pgfscope}%
\begin{pgfscope}%
\pgfpathrectangle{\pgfqpoint{0.407592in}{0.222904in}}{\pgfqpoint{1.411163in}{1.411163in}}%
\pgfusepath{clip}%
\pgfsetbuttcap%
\pgfsetroundjoin%
\definecolor{currentfill}{rgb}{0.749020,0.772549,0.792157}%
\pgfsetfillcolor{currentfill}%
\pgfsetfillopacity{0.160000}%
\pgfsetlinewidth{0.000000pt}%
\definecolor{currentstroke}{rgb}{0.000000,0.000000,0.000000}%
\pgfsetstrokecolor{currentstroke}%
\pgfsetdash{}{0pt}%
\pgfpathmoveto{\pgfqpoint{1.054113in}{0.668375in}}%
\pgfpathlineto{\pgfqpoint{1.050411in}{0.674304in}}%
\pgfpathlineto{\pgfqpoint{1.046692in}{0.680260in}}%
\pgfpathlineto{\pgfqpoint{1.057121in}{0.679378in}}%
\pgfpathlineto{\pgfqpoint{1.067700in}{0.678798in}}%
\pgfpathlineto{\pgfqpoint{1.069625in}{0.672964in}}%
\pgfpathlineto{\pgfqpoint{1.071541in}{0.667158in}}%
\pgfpathlineto{\pgfqpoint{1.062765in}{0.667641in}}%
\pgfpathlineto{\pgfqpoint{1.054113in}{0.668375in}}%
\pgfpathclose%
\pgfusepath{fill}%
\end{pgfscope}%
\begin{pgfscope}%
\pgfpathrectangle{\pgfqpoint{0.407592in}{0.222904in}}{\pgfqpoint{1.411163in}{1.411163in}}%
\pgfusepath{clip}%
\pgfsetbuttcap%
\pgfsetroundjoin%
\definecolor{currentfill}{rgb}{0.749020,0.772549,0.792157}%
\pgfsetfillcolor{currentfill}%
\pgfsetfillopacity{0.160000}%
\pgfsetlinewidth{0.000000pt}%
\definecolor{currentstroke}{rgb}{0.000000,0.000000,0.000000}%
\pgfsetstrokecolor{currentstroke}%
\pgfsetdash{}{0pt}%
\pgfpathmoveto{\pgfqpoint{1.223897in}{0.716803in}}%
\pgfpathlineto{\pgfqpoint{1.233699in}{0.724596in}}%
\pgfpathlineto{\pgfqpoint{1.243536in}{0.732416in}}%
\pgfpathlineto{\pgfqpoint{1.253072in}{0.737047in}}%
\pgfpathlineto{\pgfqpoint{1.261858in}{0.741915in}}%
\pgfpathlineto{\pgfqpoint{1.250874in}{0.733506in}}%
\pgfpathlineto{\pgfqpoint{1.239925in}{0.725125in}}%
\pgfpathlineto{\pgfqpoint{1.232236in}{0.720859in}}%
\pgfpathlineto{\pgfqpoint{1.223897in}{0.716803in}}%
\pgfpathclose%
\pgfusepath{fill}%
\end{pgfscope}%
\begin{pgfscope}%
\pgfpathrectangle{\pgfqpoint{0.407592in}{0.222904in}}{\pgfqpoint{1.411163in}{1.411163in}}%
\pgfusepath{clip}%
\pgfsetbuttcap%
\pgfsetroundjoin%
\definecolor{currentfill}{rgb}{0.749020,0.772549,0.792157}%
\pgfsetfillcolor{currentfill}%
\pgfsetfillopacity{0.160000}%
\pgfsetlinewidth{0.000000pt}%
\definecolor{currentstroke}{rgb}{0.000000,0.000000,0.000000}%
\pgfsetstrokecolor{currentstroke}%
\pgfsetdash{}{0pt}%
\pgfpathmoveto{\pgfqpoint{1.089304in}{0.666954in}}%
\pgfpathlineto{\pgfqpoint{1.089209in}{0.672740in}}%
\pgfpathlineto{\pgfqpoint{1.089114in}{0.678553in}}%
\pgfpathlineto{\pgfqpoint{1.099849in}{0.678890in}}%
\pgfpathlineto{\pgfqpoint{1.110537in}{0.679531in}}%
\pgfpathlineto{\pgfqpoint{1.108802in}{0.673636in}}%
\pgfpathlineto{\pgfqpoint{1.107075in}{0.667768in}}%
\pgfpathlineto{\pgfqpoint{1.098209in}{0.667234in}}%
\pgfpathlineto{\pgfqpoint{1.089304in}{0.666954in}}%
\pgfpathclose%
\pgfusepath{fill}%
\end{pgfscope}%
\begin{pgfscope}%
\pgfpathrectangle{\pgfqpoint{0.407592in}{0.222904in}}{\pgfqpoint{1.411163in}{1.411163in}}%
\pgfusepath{clip}%
\pgfsetbuttcap%
\pgfsetroundjoin%
\definecolor{currentfill}{rgb}{0.749020,0.772549,0.792157}%
\pgfsetfillcolor{currentfill}%
\pgfsetfillopacity{0.160000}%
\pgfsetlinewidth{0.000000pt}%
\definecolor{currentstroke}{rgb}{0.000000,0.000000,0.000000}%
\pgfsetstrokecolor{currentstroke}%
\pgfsetdash{}{0pt}%
\pgfpathmoveto{\pgfqpoint{1.684775in}{1.119923in}}%
\pgfpathlineto{\pgfqpoint{1.698699in}{1.131853in}}%
\pgfpathlineto{\pgfqpoint{1.712642in}{1.143798in}}%
\pgfpathlineto{\pgfqpoint{1.710375in}{1.159167in}}%
\pgfpathlineto{\pgfqpoint{1.705438in}{1.174325in}}%
\pgfpathlineto{\pgfqpoint{1.691792in}{1.161819in}}%
\pgfpathlineto{\pgfqpoint{1.678158in}{1.149325in}}%
\pgfpathlineto{\pgfqpoint{1.682743in}{1.134722in}}%
\pgfpathlineto{\pgfqpoint{1.684775in}{1.119923in}}%
\pgfpathclose%
\pgfusepath{fill}%
\end{pgfscope}%
\begin{pgfscope}%
\pgfpathrectangle{\pgfqpoint{0.407592in}{0.222904in}}{\pgfqpoint{1.411163in}{1.411163in}}%
\pgfusepath{clip}%
\pgfsetbuttcap%
\pgfsetroundjoin%
\definecolor{currentfill}{rgb}{0.749020,0.772549,0.792157}%
\pgfsetfillcolor{currentfill}%
\pgfsetfillopacity{0.160000}%
\pgfsetlinewidth{0.000000pt}%
\definecolor{currentstroke}{rgb}{0.000000,0.000000,0.000000}%
\pgfsetstrokecolor{currentstroke}%
\pgfsetdash{}{0pt}%
\pgfpathmoveto{\pgfqpoint{1.332913in}{0.795459in}}%
\pgfpathlineto{\pgfqpoint{1.345383in}{0.804961in}}%
\pgfpathlineto{\pgfqpoint{1.357886in}{0.814488in}}%
\pgfpathlineto{\pgfqpoint{1.365270in}{0.822080in}}%
\pgfpathlineto{\pgfqpoint{1.371404in}{0.829816in}}%
\pgfpathlineto{\pgfqpoint{1.358327in}{0.819617in}}%
\pgfpathlineto{\pgfqpoint{1.345279in}{0.809441in}}%
\pgfpathlineto{\pgfqpoint{1.339661in}{0.802383in}}%
\pgfpathlineto{\pgfqpoint{1.332913in}{0.795459in}}%
\pgfpathclose%
\pgfusepath{fill}%
\end{pgfscope}%
\begin{pgfscope}%
\pgfpathrectangle{\pgfqpoint{0.407592in}{0.222904in}}{\pgfqpoint{1.411163in}{1.411163in}}%
\pgfusepath{clip}%
\pgfsetbuttcap%
\pgfsetroundjoin%
\definecolor{currentfill}{rgb}{0.749020,0.772549,0.792157}%
\pgfsetfillcolor{currentfill}%
\pgfsetfillopacity{0.160000}%
\pgfsetlinewidth{0.000000pt}%
\definecolor{currentstroke}{rgb}{0.000000,0.000000,0.000000}%
\pgfsetstrokecolor{currentstroke}%
\pgfsetdash{}{0pt}%
\pgfpathmoveto{\pgfqpoint{1.071541in}{0.667158in}}%
\pgfpathlineto{\pgfqpoint{1.069625in}{0.672964in}}%
\pgfpathlineto{\pgfqpoint{1.067700in}{0.678798in}}%
\pgfpathlineto{\pgfqpoint{1.078381in}{0.678522in}}%
\pgfpathlineto{\pgfqpoint{1.089114in}{0.678553in}}%
\pgfpathlineto{\pgfqpoint{1.089209in}{0.672740in}}%
\pgfpathlineto{\pgfqpoint{1.089304in}{0.666954in}}%
\pgfpathlineto{\pgfqpoint{1.080401in}{0.666929in}}%
\pgfpathlineto{\pgfqpoint{1.071541in}{0.667158in}}%
\pgfpathclose%
\pgfusepath{fill}%
\end{pgfscope}%
\begin{pgfscope}%
\pgfpathrectangle{\pgfqpoint{0.407592in}{0.222904in}}{\pgfqpoint{1.411163in}{1.411163in}}%
\pgfusepath{clip}%
\pgfsetbuttcap%
\pgfsetroundjoin%
\definecolor{currentfill}{rgb}{0.749020,0.772549,0.792157}%
\pgfsetfillcolor{currentfill}%
\pgfsetfillopacity{0.160000}%
\pgfsetlinewidth{0.000000pt}%
\definecolor{currentstroke}{rgb}{0.000000,0.000000,0.000000}%
\pgfsetstrokecolor{currentstroke}%
\pgfsetdash{}{0pt}%
\pgfpathmoveto{\pgfqpoint{0.768605in}{0.865674in}}%
\pgfpathlineto{\pgfqpoint{0.754898in}{0.876547in}}%
\pgfpathlineto{\pgfqpoint{0.741166in}{0.887440in}}%
\pgfpathlineto{\pgfqpoint{0.743521in}{0.878180in}}%
\pgfpathlineto{\pgfqpoint{0.747421in}{0.869006in}}%
\pgfpathlineto{\pgfqpoint{0.760979in}{0.858783in}}%
\pgfpathlineto{\pgfqpoint{0.774506in}{0.848584in}}%
\pgfpathlineto{\pgfqpoint{0.770845in}{0.857087in}}%
\pgfpathlineto{\pgfqpoint{0.768605in}{0.865674in}}%
\pgfpathclose%
\pgfusepath{fill}%
\end{pgfscope}%
\begin{pgfscope}%
\pgfpathrectangle{\pgfqpoint{0.407592in}{0.222904in}}{\pgfqpoint{1.411163in}{1.411163in}}%
\pgfusepath{clip}%
\pgfsetbuttcap%
\pgfsetroundjoin%
\definecolor{currentfill}{rgb}{0.749020,0.772549,0.792157}%
\pgfsetfillcolor{currentfill}%
\pgfsetfillopacity{0.160000}%
\pgfsetlinewidth{0.000000pt}%
\definecolor{currentstroke}{rgb}{0.000000,0.000000,0.000000}%
\pgfsetstrokecolor{currentstroke}%
\pgfsetdash{}{0pt}%
\pgfpathmoveto{\pgfqpoint{1.515267in}{0.955480in}}%
\pgfpathlineto{\pgfqpoint{1.528942in}{0.966578in}}%
\pgfpathlineto{\pgfqpoint{1.542642in}{0.977696in}}%
\pgfpathlineto{\pgfqpoint{1.545612in}{0.989616in}}%
\pgfpathlineto{\pgfqpoint{1.546564in}{1.001508in}}%
\pgfpathlineto{\pgfqpoint{1.532845in}{0.989754in}}%
\pgfpathlineto{\pgfqpoint{1.519145in}{0.978016in}}%
\pgfpathlineto{\pgfqpoint{1.518154in}{0.966759in}}%
\pgfpathlineto{\pgfqpoint{1.515267in}{0.955480in}}%
\pgfpathclose%
\pgfusepath{fill}%
\end{pgfscope}%
\begin{pgfscope}%
\pgfpathrectangle{\pgfqpoint{0.407592in}{0.222904in}}{\pgfqpoint{1.411163in}{1.411163in}}%
\pgfusepath{clip}%
\pgfsetbuttcap%
\pgfsetroundjoin%
\definecolor{currentfill}{rgb}{0.749020,0.772549,0.792157}%
\pgfsetfillcolor{currentfill}%
\pgfsetfillopacity{0.160000}%
\pgfsetlinewidth{0.000000pt}%
\definecolor{currentstroke}{rgb}{0.000000,0.000000,0.000000}%
\pgfsetstrokecolor{currentstroke}%
\pgfsetdash{}{0pt}%
\pgfpathmoveto{\pgfqpoint{0.828311in}{0.808016in}}%
\pgfpathlineto{\pgfqpoint{0.814905in}{0.818123in}}%
\pgfpathlineto{\pgfqpoint{0.801469in}{0.828254in}}%
\pgfpathlineto{\pgfqpoint{0.806173in}{0.820544in}}%
\pgfpathlineto{\pgfqpoint{0.812157in}{0.812986in}}%
\pgfpathlineto{\pgfqpoint{0.825152in}{0.803525in}}%
\pgfpathlineto{\pgfqpoint{0.838112in}{0.794089in}}%
\pgfpathlineto{\pgfqpoint{0.832633in}{0.800982in}}%
\pgfpathlineto{\pgfqpoint{0.828311in}{0.808016in}}%
\pgfpathclose%
\pgfusepath{fill}%
\end{pgfscope}%
\begin{pgfscope}%
\pgfpathrectangle{\pgfqpoint{0.407592in}{0.222904in}}{\pgfqpoint{1.411163in}{1.411163in}}%
\pgfusepath{clip}%
\pgfsetbuttcap%
\pgfsetroundjoin%
\definecolor{currentfill}{rgb}{0.749020,0.772549,0.792157}%
\pgfsetfillcolor{currentfill}%
\pgfsetfillopacity{0.160000}%
\pgfsetlinewidth{0.000000pt}%
\definecolor{currentstroke}{rgb}{0.000000,0.000000,0.000000}%
\pgfsetstrokecolor{currentstroke}%
\pgfsetdash{}{0pt}%
\pgfpathmoveto{\pgfqpoint{1.170952in}{0.689494in}}%
\pgfpathlineto{\pgfqpoint{1.177859in}{0.696250in}}%
\pgfpathlineto{\pgfqpoint{1.184795in}{0.703034in}}%
\pgfpathlineto{\pgfqpoint{1.195348in}{0.706073in}}%
\pgfpathlineto{\pgfqpoint{1.205412in}{0.709392in}}%
\pgfpathlineto{\pgfqpoint{1.196979in}{0.702154in}}%
\pgfpathlineto{\pgfqpoint{1.188578in}{0.694944in}}%
\pgfpathlineto{\pgfqpoint{1.179972in}{0.692099in}}%
\pgfpathlineto{\pgfqpoint{1.170952in}{0.689494in}}%
\pgfpathclose%
\pgfusepath{fill}%
\end{pgfscope}%
\begin{pgfscope}%
\pgfpathrectangle{\pgfqpoint{0.407592in}{0.222904in}}{\pgfqpoint{1.411163in}{1.411163in}}%
\pgfusepath{clip}%
\pgfsetbuttcap%
\pgfsetroundjoin%
\definecolor{currentfill}{rgb}{0.749020,0.772549,0.792157}%
\pgfsetfillcolor{currentfill}%
\pgfsetfillopacity{0.160000}%
\pgfsetlinewidth{0.000000pt}%
\definecolor{currentstroke}{rgb}{0.000000,0.000000,0.000000}%
\pgfsetstrokecolor{currentstroke}%
\pgfsetdash{}{0pt}%
\pgfpathmoveto{\pgfqpoint{0.504495in}{1.146414in}}%
\pgfpathlineto{\pgfqpoint{0.490898in}{1.158853in}}%
\pgfpathlineto{\pgfqpoint{0.477289in}{1.171304in}}%
\pgfpathlineto{\pgfqpoint{0.469249in}{1.156099in}}%
\pgfpathlineto{\pgfqpoint{0.463842in}{1.140695in}}%
\pgfpathlineto{\pgfqpoint{0.477888in}{1.128807in}}%
\pgfpathlineto{\pgfqpoint{0.491914in}{1.116936in}}%
\pgfpathlineto{\pgfqpoint{0.496945in}{1.131767in}}%
\pgfpathlineto{\pgfqpoint{0.504495in}{1.146414in}}%
\pgfpathclose%
\pgfusepath{fill}%
\end{pgfscope}%
\begin{pgfscope}%
\pgfpathrectangle{\pgfqpoint{0.407592in}{0.222904in}}{\pgfqpoint{1.411163in}{1.411163in}}%
\pgfusepath{clip}%
\pgfsetbuttcap%
\pgfsetroundjoin%
\definecolor{currentfill}{rgb}{0.749020,0.772549,0.792157}%
\pgfsetfillcolor{currentfill}%
\pgfsetfillopacity{0.160000}%
\pgfsetlinewidth{0.000000pt}%
\definecolor{currentstroke}{rgb}{0.000000,0.000000,0.000000}%
\pgfsetstrokecolor{currentstroke}%
\pgfsetdash{}{0pt}%
\pgfpathmoveto{\pgfqpoint{1.261858in}{0.741915in}}%
\pgfpathlineto{\pgfqpoint{1.272878in}{0.750350in}}%
\pgfpathlineto{\pgfqpoint{1.283933in}{0.758812in}}%
\pgfpathlineto{\pgfqpoint{1.292921in}{0.764517in}}%
\pgfpathlineto{\pgfqpoint{1.281371in}{0.755743in}}%
\pgfpathlineto{\pgfqpoint{1.269856in}{0.746995in}}%
\pgfusepath{fill}%
\end{pgfscope}%
\begin{pgfscope}%
\pgfpathrectangle{\pgfqpoint{0.407592in}{0.222904in}}{\pgfqpoint{1.411163in}{1.411163in}}%
\pgfusepath{clip}%
\pgfsetbuttcap%
\pgfsetroundjoin%
\definecolor{currentfill}{rgb}{0.749020,0.772549,0.792157}%
\pgfsetfillcolor{currentfill}%
\pgfsetfillopacity{0.160000}%
\pgfsetlinewidth{0.000000pt}%
\definecolor{currentstroke}{rgb}{0.000000,0.000000,0.000000}%
\pgfsetstrokecolor{currentstroke}%
\pgfsetdash{}{0pt}%
\pgfpathmoveto{\pgfqpoint{1.712642in}{1.143798in}}%
\pgfpathlineto{\pgfqpoint{1.726603in}{1.155760in}}%
\pgfpathlineto{\pgfqpoint{1.740584in}{1.167739in}}%
\pgfpathlineto{\pgfqpoint{1.738068in}{1.183667in}}%
\pgfpathlineto{\pgfqpoint{1.732768in}{1.199371in}}%
\pgfpathlineto{\pgfqpoint{1.719097in}{1.186842in}}%
\pgfpathlineto{\pgfqpoint{1.705438in}{1.174325in}}%
\pgfpathlineto{\pgfqpoint{1.710375in}{1.159167in}}%
\pgfpathlineto{\pgfqpoint{1.712642in}{1.143798in}}%
\pgfpathclose%
\pgfusepath{fill}%
\end{pgfscope}%
\begin{pgfscope}%
\pgfpathrectangle{\pgfqpoint{0.407592in}{0.222904in}}{\pgfqpoint{1.411163in}{1.411163in}}%
\pgfusepath{clip}%
\pgfsetbuttcap%
\pgfsetroundjoin%
\definecolor{currentfill}{rgb}{0.749020,0.772549,0.792157}%
\pgfsetfillcolor{currentfill}%
\pgfsetfillopacity{0.160000}%
\pgfsetlinewidth{0.000000pt}%
\definecolor{currentstroke}{rgb}{0.000000,0.000000,0.000000}%
\pgfsetstrokecolor{currentstroke}%
\pgfsetdash{}{0pt}%
\pgfpathmoveto{\pgfqpoint{0.658712in}{0.975761in}}%
\pgfpathlineto{\pgfqpoint{0.644918in}{0.987436in}}%
\pgfpathlineto{\pgfqpoint{0.631105in}{0.999127in}}%
\pgfpathlineto{\pgfqpoint{0.629747in}{0.987226in}}%
\pgfpathlineto{\pgfqpoint{0.630405in}{0.975306in}}%
\pgfpathlineto{\pgfqpoint{0.644339in}{0.964252in}}%
\pgfpathlineto{\pgfqpoint{0.658248in}{0.953219in}}%
\pgfpathlineto{\pgfqpoint{0.657532in}{0.964496in}}%
\pgfpathlineto{\pgfqpoint{0.658712in}{0.975761in}}%
\pgfpathclose%
\pgfusepath{fill}%
\end{pgfscope}%
\begin{pgfscope}%
\pgfpathrectangle{\pgfqpoint{0.407592in}{0.222904in}}{\pgfqpoint{1.411163in}{1.411163in}}%
\pgfusepath{clip}%
\pgfsetbuttcap%
\pgfsetroundjoin%
\definecolor{currentfill}{rgb}{0.749020,0.772549,0.792157}%
\pgfsetfillcolor{currentfill}%
\pgfsetfillopacity{0.160000}%
\pgfsetlinewidth{0.000000pt}%
\definecolor{currentstroke}{rgb}{0.000000,0.000000,0.000000}%
\pgfsetstrokecolor{currentstroke}%
\pgfsetdash{}{0pt}%
\pgfpathmoveto{\pgfqpoint{0.938446in}{0.720028in}}%
\pgfpathlineto{\pgfqpoint{0.927327in}{0.728050in}}%
\pgfpathlineto{\pgfqpoint{0.916170in}{0.736098in}}%
\pgfpathlineto{\pgfqpoint{0.924874in}{0.731518in}}%
\pgfpathlineto{\pgfqpoint{0.934335in}{0.727200in}}%
\pgfpathlineto{\pgfqpoint{0.944352in}{0.719704in}}%
\pgfpathlineto{\pgfqpoint{0.954332in}{0.712235in}}%
\pgfpathlineto{\pgfqpoint{0.946061in}{0.716016in}}%
\pgfpathlineto{\pgfqpoint{0.938446in}{0.720028in}}%
\pgfpathclose%
\pgfusepath{fill}%
\end{pgfscope}%
\begin{pgfscope}%
\pgfpathrectangle{\pgfqpoint{0.407592in}{0.222904in}}{\pgfqpoint{1.411163in}{1.411163in}}%
\pgfusepath{clip}%
\pgfsetbuttcap%
\pgfsetroundjoin%
\definecolor{currentfill}{rgb}{0.749020,0.772549,0.792157}%
\pgfsetfillcolor{currentfill}%
\pgfsetfillopacity{0.160000}%
\pgfsetlinewidth{0.000000pt}%
\definecolor{currentstroke}{rgb}{0.000000,0.000000,0.000000}%
\pgfsetstrokecolor{currentstroke}%
\pgfsetdash{}{0pt}%
\pgfpathmoveto{\pgfqpoint{1.292921in}{0.764517in}}%
\pgfpathlineto{\pgfqpoint{1.304505in}{0.773318in}}%
\pgfpathlineto{\pgfqpoint{1.316125in}{0.782146in}}%
\pgfpathlineto{\pgfqpoint{1.325058in}{0.788703in}}%
\pgfpathlineto{\pgfqpoint{1.332913in}{0.795459in}}%
\pgfpathlineto{\pgfqpoint{1.320476in}{0.785982in}}%
\pgfpathlineto{\pgfqpoint{1.308072in}{0.776530in}}%
\pgfpathlineto{\pgfqpoint{1.300978in}{0.770432in}}%
\pgfpathlineto{\pgfqpoint{1.292921in}{0.764517in}}%
\pgfpathclose%
\pgfusepath{fill}%
\end{pgfscope}%
\begin{pgfscope}%
\pgfpathrectangle{\pgfqpoint{0.407592in}{0.222904in}}{\pgfqpoint{1.411163in}{1.411163in}}%
\pgfusepath{clip}%
\pgfsetbuttcap%
\pgfsetroundjoin%
\definecolor{currentfill}{rgb}{0.749020,0.772549,0.792157}%
\pgfsetfillcolor{currentfill}%
\pgfsetfillopacity{0.160000}%
\pgfsetlinewidth{0.000000pt}%
\definecolor{currentstroke}{rgb}{0.000000,0.000000,0.000000}%
\pgfsetstrokecolor{currentstroke}%
\pgfsetdash{}{0pt}%
\pgfpathmoveto{\pgfqpoint{0.989891in}{0.691556in}}%
\pgfpathlineto{\pgfqpoint{0.981312in}{0.698485in}}%
\pgfpathlineto{\pgfqpoint{0.972697in}{0.705441in}}%
\pgfpathlineto{\pgfqpoint{0.982710in}{0.702460in}}%
\pgfpathlineto{\pgfqpoint{0.993217in}{0.699776in}}%
\pgfpathlineto{\pgfqpoint{1.000338in}{0.693225in}}%
\pgfpathlineto{\pgfqpoint{1.007429in}{0.686702in}}%
\pgfpathlineto{\pgfqpoint{0.998451in}{0.689002in}}%
\pgfpathlineto{\pgfqpoint{0.989891in}{0.691556in}}%
\pgfpathclose%
\pgfusepath{fill}%
\end{pgfscope}%
\begin{pgfscope}%
\pgfpathrectangle{\pgfqpoint{0.407592in}{0.222904in}}{\pgfqpoint{1.411163in}{1.411163in}}%
\pgfusepath{clip}%
\pgfsetbuttcap%
\pgfsetroundjoin%
\definecolor{currentfill}{rgb}{0.749020,0.772549,0.792157}%
\pgfsetfillcolor{currentfill}%
\pgfsetfillopacity{0.160000}%
\pgfsetlinewidth{0.000000pt}%
\definecolor{currentstroke}{rgb}{0.000000,0.000000,0.000000}%
\pgfsetstrokecolor{currentstroke}%
\pgfsetdash{}{0pt}%
\pgfpathmoveto{\pgfqpoint{1.424004in}{0.870837in}}%
\pgfpathlineto{\pgfqpoint{1.437227in}{0.881150in}}%
\pgfpathlineto{\pgfqpoint{1.450479in}{0.891485in}}%
\pgfpathlineto{\pgfqpoint{1.456470in}{0.901345in}}%
\pgfpathlineto{\pgfqpoint{1.460810in}{0.911284in}}%
\pgfpathlineto{\pgfqpoint{1.447257in}{0.900285in}}%
\pgfpathlineto{\pgfqpoint{1.433727in}{0.889305in}}%
\pgfpathlineto{\pgfqpoint{1.429629in}{0.880032in}}%
\pgfpathlineto{\pgfqpoint{1.424004in}{0.870837in}}%
\pgfpathclose%
\pgfusepath{fill}%
\end{pgfscope}%
\begin{pgfscope}%
\pgfpathrectangle{\pgfqpoint{0.407592in}{0.222904in}}{\pgfqpoint{1.411163in}{1.411163in}}%
\pgfusepath{clip}%
\pgfsetbuttcap%
\pgfsetroundjoin%
\definecolor{currentfill}{rgb}{0.749020,0.772549,0.792157}%
\pgfsetfillcolor{currentfill}%
\pgfsetfillopacity{0.160000}%
\pgfsetlinewidth{0.000000pt}%
\definecolor{currentstroke}{rgb}{0.000000,0.000000,0.000000}%
\pgfsetstrokecolor{currentstroke}%
\pgfsetdash{}{0pt}%
\pgfpathmoveto{\pgfqpoint{0.863928in}{0.775293in}}%
\pgfpathlineto{\pgfqpoint{0.851037in}{0.784679in}}%
\pgfpathlineto{\pgfqpoint{0.838112in}{0.794089in}}%
\pgfpathlineto{\pgfqpoint{0.844730in}{0.787370in}}%
\pgfpathlineto{\pgfqpoint{0.852464in}{0.780857in}}%
\pgfpathlineto{\pgfqpoint{0.864698in}{0.772093in}}%
\pgfpathlineto{\pgfqpoint{0.876896in}{0.763355in}}%
\pgfpathlineto{\pgfqpoint{0.869913in}{0.769230in}}%
\pgfpathlineto{\pgfqpoint{0.863928in}{0.775293in}}%
\pgfpathclose%
\pgfusepath{fill}%
\end{pgfscope}%
\begin{pgfscope}%
\pgfpathrectangle{\pgfqpoint{0.407592in}{0.222904in}}{\pgfqpoint{1.411163in}{1.411163in}}%
\pgfusepath{clip}%
\pgfsetbuttcap%
\pgfsetroundjoin%
\definecolor{currentfill}{rgb}{0.749020,0.772549,0.792157}%
\pgfsetfillcolor{currentfill}%
\pgfsetfillopacity{0.160000}%
\pgfsetlinewidth{0.000000pt}%
\definecolor{currentstroke}{rgb}{0.000000,0.000000,0.000000}%
\pgfsetstrokecolor{currentstroke}%
\pgfsetdash{}{0pt}%
\pgfpathmoveto{\pgfqpoint{1.542642in}{0.977696in}}%
\pgfpathlineto{\pgfqpoint{1.556366in}{0.988834in}}%
\pgfpathlineto{\pgfqpoint{1.570115in}{0.999993in}}%
\pgfpathlineto{\pgfqpoint{1.573157in}{1.012546in}}%
\pgfpathlineto{\pgfqpoint{1.574058in}{1.025063in}}%
\pgfpathlineto{\pgfqpoint{1.560302in}{1.013278in}}%
\pgfpathlineto{\pgfqpoint{1.546564in}{1.001508in}}%
\pgfpathlineto{\pgfqpoint{1.545612in}{0.989616in}}%
\pgfpathlineto{\pgfqpoint{1.542642in}{0.977696in}}%
\pgfpathclose%
\pgfusepath{fill}%
\end{pgfscope}%
\begin{pgfscope}%
\pgfpathrectangle{\pgfqpoint{0.407592in}{0.222904in}}{\pgfqpoint{1.411163in}{1.411163in}}%
\pgfusepath{clip}%
\pgfsetbuttcap%
\pgfsetroundjoin%
\definecolor{currentfill}{rgb}{0.749020,0.772549,0.792157}%
\pgfsetfillcolor{currentfill}%
\pgfsetfillopacity{0.160000}%
\pgfsetlinewidth{0.000000pt}%
\definecolor{currentstroke}{rgb}{0.000000,0.000000,0.000000}%
\pgfsetstrokecolor{currentstroke}%
\pgfsetdash{}{0pt}%
\pgfpathmoveto{\pgfqpoint{1.151835in}{0.685055in}}%
\pgfpathlineto{\pgfqpoint{1.157114in}{0.691440in}}%
\pgfpathlineto{\pgfqpoint{1.162415in}{0.697853in}}%
\pgfpathlineto{\pgfqpoint{1.173801in}{0.700290in}}%
\pgfpathlineto{\pgfqpoint{1.184795in}{0.703034in}}%
\pgfpathlineto{\pgfqpoint{1.177859in}{0.696250in}}%
\pgfpathlineto{\pgfqpoint{1.170952in}{0.689494in}}%
\pgfpathlineto{\pgfqpoint{1.161559in}{0.687142in}}%
\pgfpathlineto{\pgfqpoint{1.151835in}{0.685055in}}%
\pgfpathclose%
\pgfusepath{fill}%
\end{pgfscope}%
\begin{pgfscope}%
\pgfpathrectangle{\pgfqpoint{0.407592in}{0.222904in}}{\pgfqpoint{1.411163in}{1.411163in}}%
\pgfusepath{clip}%
\pgfsetbuttcap%
\pgfsetroundjoin%
\definecolor{currentfill}{rgb}{0.749020,0.772549,0.792157}%
\pgfsetfillcolor{currentfill}%
\pgfsetfillopacity{0.160000}%
\pgfsetlinewidth{0.000000pt}%
\definecolor{currentstroke}{rgb}{0.000000,0.000000,0.000000}%
\pgfsetstrokecolor{currentstroke}%
\pgfsetdash{}{0pt}%
\pgfpathmoveto{\pgfqpoint{0.901178in}{0.745960in}}%
\pgfpathlineto{\pgfqpoint{0.889056in}{0.754644in}}%
\pgfpathlineto{\pgfqpoint{0.876896in}{0.763355in}}%
\pgfpathlineto{\pgfqpoint{0.884849in}{0.757696in}}%
\pgfpathlineto{\pgfqpoint{0.893740in}{0.752279in}}%
\pgfpathlineto{\pgfqpoint{0.904974in}{0.744175in}}%
\pgfpathlineto{\pgfqpoint{0.916170in}{0.736098in}}%
\pgfpathlineto{\pgfqpoint{0.908260in}{0.740920in}}%
\pgfpathlineto{\pgfqpoint{0.901178in}{0.745960in}}%
\pgfpathclose%
\pgfusepath{fill}%
\end{pgfscope}%
\begin{pgfscope}%
\pgfpathrectangle{\pgfqpoint{0.407592in}{0.222904in}}{\pgfqpoint{1.411163in}{1.411163in}}%
\pgfusepath{clip}%
\pgfsetbuttcap%
\pgfsetroundjoin%
\definecolor{currentfill}{rgb}{0.749020,0.772549,0.792157}%
\pgfsetfillcolor{currentfill}%
\pgfsetfillopacity{0.160000}%
\pgfsetlinewidth{0.000000pt}%
\definecolor{currentstroke}{rgb}{0.000000,0.000000,0.000000}%
\pgfsetstrokecolor{currentstroke}%
\pgfsetdash{}{0pt}%
\pgfpathmoveto{\pgfqpoint{0.477289in}{1.171304in}}%
\pgfpathlineto{\pgfqpoint{0.463667in}{1.183767in}}%
\pgfpathlineto{\pgfqpoint{0.450032in}{1.196241in}}%
\pgfpathlineto{\pgfqpoint{0.441489in}{1.180488in}}%
\pgfpathlineto{\pgfqpoint{0.435692in}{1.164522in}}%
\pgfpathlineto{\pgfqpoint{0.449777in}{1.152600in}}%
\pgfpathlineto{\pgfqpoint{0.463842in}{1.140695in}}%
\pgfpathlineto{\pgfqpoint{0.469249in}{1.156099in}}%
\pgfpathlineto{\pgfqpoint{0.477289in}{1.171304in}}%
\pgfpathclose%
\pgfusepath{fill}%
\end{pgfscope}%
\begin{pgfscope}%
\pgfpathrectangle{\pgfqpoint{0.407592in}{0.222904in}}{\pgfqpoint{1.411163in}{1.411163in}}%
\pgfusepath{clip}%
\pgfsetbuttcap%
\pgfsetroundjoin%
\definecolor{currentfill}{rgb}{0.749020,0.772549,0.792157}%
\pgfsetfillcolor{currentfill}%
\pgfsetfillopacity{0.160000}%
\pgfsetlinewidth{0.000000pt}%
\definecolor{currentstroke}{rgb}{0.000000,0.000000,0.000000}%
\pgfsetstrokecolor{currentstroke}%
\pgfsetdash{}{0pt}%
\pgfpathmoveto{\pgfqpoint{0.741166in}{0.887440in}}%
\pgfpathlineto{\pgfqpoint{0.727410in}{0.898353in}}%
\pgfpathlineto{\pgfqpoint{0.713628in}{0.909286in}}%
\pgfpathlineto{\pgfqpoint{0.716086in}{0.899359in}}%
\pgfpathlineto{\pgfqpoint{0.720213in}{0.889520in}}%
\pgfpathlineto{\pgfqpoint{0.733832in}{0.879251in}}%
\pgfpathlineto{\pgfqpoint{0.747421in}{0.869006in}}%
\pgfpathlineto{\pgfqpoint{0.743521in}{0.878180in}}%
\pgfpathlineto{\pgfqpoint{0.741166in}{0.887440in}}%
\pgfpathclose%
\pgfusepath{fill}%
\end{pgfscope}%
\begin{pgfscope}%
\pgfpathrectangle{\pgfqpoint{0.407592in}{0.222904in}}{\pgfqpoint{1.411163in}{1.411163in}}%
\pgfusepath{clip}%
\pgfsetbuttcap%
\pgfsetroundjoin%
\definecolor{currentfill}{rgb}{0.749020,0.772549,0.792157}%
\pgfsetfillcolor{currentfill}%
\pgfsetfillopacity{0.160000}%
\pgfsetlinewidth{0.000000pt}%
\definecolor{currentstroke}{rgb}{0.000000,0.000000,0.000000}%
\pgfsetstrokecolor{currentstroke}%
\pgfsetdash{}{0pt}%
\pgfpathmoveto{\pgfqpoint{1.740584in}{1.167739in}}%
\pgfpathlineto{\pgfqpoint{1.754584in}{1.179733in}}%
\pgfpathlineto{\pgfqpoint{1.768602in}{1.191744in}}%
\pgfpathlineto{\pgfqpoint{1.765824in}{1.208223in}}%
\pgfpathlineto{\pgfqpoint{1.760147in}{1.224461in}}%
\pgfpathlineto{\pgfqpoint{1.746451in}{1.211910in}}%
\pgfpathlineto{\pgfqpoint{1.732768in}{1.199371in}}%
\pgfpathlineto{\pgfqpoint{1.738068in}{1.183667in}}%
\pgfpathlineto{\pgfqpoint{1.740584in}{1.167739in}}%
\pgfpathclose%
\pgfusepath{fill}%
\end{pgfscope}%
\begin{pgfscope}%
\pgfpathrectangle{\pgfqpoint{0.407592in}{0.222904in}}{\pgfqpoint{1.411163in}{1.411163in}}%
\pgfusepath{clip}%
\pgfsetbuttcap%
\pgfsetroundjoin%
\definecolor{currentfill}{rgb}{0.749020,0.772549,0.792157}%
\pgfsetfillcolor{currentfill}%
\pgfsetfillopacity{0.160000}%
\pgfsetlinewidth{0.000000pt}%
\definecolor{currentstroke}{rgb}{0.000000,0.000000,0.000000}%
\pgfsetstrokecolor{currentstroke}%
\pgfsetdash{}{0pt}%
\pgfpathmoveto{\pgfqpoint{1.205412in}{0.709392in}}%
\pgfpathlineto{\pgfqpoint{1.213878in}{0.716658in}}%
\pgfpathlineto{\pgfqpoint{1.222376in}{0.723952in}}%
\pgfpathlineto{\pgfqpoint{1.233289in}{0.728044in}}%
\pgfpathlineto{\pgfqpoint{1.243536in}{0.732416in}}%
\pgfpathlineto{\pgfqpoint{1.233699in}{0.724596in}}%
\pgfpathlineto{\pgfqpoint{1.223897in}{0.716803in}}%
\pgfpathlineto{\pgfqpoint{1.214942in}{0.712974in}}%
\pgfpathlineto{\pgfqpoint{1.205412in}{0.709392in}}%
\pgfpathclose%
\pgfusepath{fill}%
\end{pgfscope}%
\begin{pgfscope}%
\pgfpathrectangle{\pgfqpoint{0.407592in}{0.222904in}}{\pgfqpoint{1.411163in}{1.411163in}}%
\pgfusepath{clip}%
\pgfsetbuttcap%
\pgfsetroundjoin%
\definecolor{currentfill}{rgb}{0.749020,0.772549,0.792157}%
\pgfsetfillcolor{currentfill}%
\pgfsetfillopacity{0.160000}%
\pgfsetlinewidth{0.000000pt}%
\definecolor{currentstroke}{rgb}{0.000000,0.000000,0.000000}%
\pgfsetstrokecolor{currentstroke}%
\pgfsetdash{}{0pt}%
\pgfpathmoveto{\pgfqpoint{1.357886in}{0.814488in}}%
\pgfpathlineto{\pgfqpoint{1.370421in}{0.824040in}}%
\pgfpathlineto{\pgfqpoint{1.382990in}{0.833618in}}%
\pgfpathlineto{\pgfqpoint{1.391004in}{0.841873in}}%
\pgfpathlineto{\pgfqpoint{1.397646in}{0.850281in}}%
\pgfpathlineto{\pgfqpoint{1.384511in}{0.840037in}}%
\pgfpathlineto{\pgfqpoint{1.371404in}{0.829816in}}%
\pgfpathlineto{\pgfqpoint{1.365270in}{0.822080in}}%
\pgfpathlineto{\pgfqpoint{1.357886in}{0.814488in}}%
\pgfpathclose%
\pgfusepath{fill}%
\end{pgfscope}%
\begin{pgfscope}%
\pgfpathrectangle{\pgfqpoint{0.407592in}{0.222904in}}{\pgfqpoint{1.411163in}{1.411163in}}%
\pgfusepath{clip}%
\pgfsetbuttcap%
\pgfsetroundjoin%
\definecolor{currentfill}{rgb}{0.749020,0.772549,0.792157}%
\pgfsetfillcolor{currentfill}%
\pgfsetfillopacity{0.160000}%
\pgfsetlinewidth{0.000000pt}%
\definecolor{currentstroke}{rgb}{0.000000,0.000000,0.000000}%
\pgfsetstrokecolor{currentstroke}%
\pgfsetdash{}{0pt}%
\pgfpathmoveto{\pgfqpoint{0.631105in}{0.999127in}}%
\pgfpathlineto{\pgfqpoint{0.617273in}{1.010834in}}%
\pgfpathlineto{\pgfqpoint{0.603422in}{1.022558in}}%
\pgfpathlineto{\pgfqpoint{0.601871in}{1.010029in}}%
\pgfpathlineto{\pgfqpoint{0.602459in}{0.997475in}}%
\pgfpathlineto{\pgfqpoint{0.616445in}{0.986380in}}%
\pgfpathlineto{\pgfqpoint{0.630405in}{0.975306in}}%
\pgfpathlineto{\pgfqpoint{0.629747in}{0.987226in}}%
\pgfpathlineto{\pgfqpoint{0.631105in}{0.999127in}}%
\pgfpathclose%
\pgfusepath{fill}%
\end{pgfscope}%
\begin{pgfscope}%
\pgfpathrectangle{\pgfqpoint{0.407592in}{0.222904in}}{\pgfqpoint{1.411163in}{1.411163in}}%
\pgfusepath{clip}%
\pgfsetbuttcap%
\pgfsetroundjoin%
\definecolor{currentfill}{rgb}{0.749020,0.772549,0.792157}%
\pgfsetfillcolor{currentfill}%
\pgfsetfillopacity{0.160000}%
\pgfsetlinewidth{0.000000pt}%
\definecolor{currentstroke}{rgb}{0.000000,0.000000,0.000000}%
\pgfsetstrokecolor{currentstroke}%
\pgfsetdash{}{0pt}%
\pgfpathmoveto{\pgfqpoint{1.007429in}{0.686702in}}%
\pgfpathlineto{\pgfqpoint{1.000338in}{0.693225in}}%
\pgfpathlineto{\pgfqpoint{0.993217in}{0.699776in}}%
\pgfpathlineto{\pgfqpoint{1.004171in}{0.697403in}}%
\pgfpathlineto{\pgfqpoint{1.015522in}{0.695351in}}%
\pgfpathlineto{\pgfqpoint{1.021012in}{0.689118in}}%
\pgfpathlineto{\pgfqpoint{1.026478in}{0.682912in}}%
\pgfpathlineto{\pgfqpoint{1.016786in}{0.684669in}}%
\pgfpathlineto{\pgfqpoint{1.007429in}{0.686702in}}%
\pgfpathclose%
\pgfusepath{fill}%
\end{pgfscope}%
\begin{pgfscope}%
\pgfpathrectangle{\pgfqpoint{0.407592in}{0.222904in}}{\pgfqpoint{1.411163in}{1.411163in}}%
\pgfusepath{clip}%
\pgfsetbuttcap%
\pgfsetroundjoin%
\definecolor{currentfill}{rgb}{0.749020,0.772549,0.792157}%
\pgfsetfillcolor{currentfill}%
\pgfsetfillopacity{0.160000}%
\pgfsetlinewidth{0.000000pt}%
\definecolor{currentstroke}{rgb}{0.000000,0.000000,0.000000}%
\pgfsetstrokecolor{currentstroke}%
\pgfsetdash{}{0pt}%
\pgfpathmoveto{\pgfqpoint{0.801469in}{0.828254in}}%
\pgfpathlineto{\pgfqpoint{0.788003in}{0.838407in}}%
\pgfpathlineto{\pgfqpoint{0.774506in}{0.848584in}}%
\pgfpathlineto{\pgfqpoint{0.779583in}{0.840203in}}%
\pgfpathlineto{\pgfqpoint{0.786063in}{0.831984in}}%
\pgfpathlineto{\pgfqpoint{0.799127in}{0.822472in}}%
\pgfpathlineto{\pgfqpoint{0.812157in}{0.812986in}}%
\pgfpathlineto{\pgfqpoint{0.806173in}{0.820544in}}%
\pgfpathlineto{\pgfqpoint{0.801469in}{0.828254in}}%
\pgfpathclose%
\pgfusepath{fill}%
\end{pgfscope}%
\begin{pgfscope}%
\pgfpathrectangle{\pgfqpoint{0.407592in}{0.222904in}}{\pgfqpoint{1.411163in}{1.411163in}}%
\pgfusepath{clip}%
\pgfsetbuttcap%
\pgfsetroundjoin%
\definecolor{currentfill}{rgb}{0.749020,0.772549,0.792157}%
\pgfsetfillcolor{currentfill}%
\pgfsetfillopacity{0.160000}%
\pgfsetlinewidth{0.000000pt}%
\definecolor{currentstroke}{rgb}{0.000000,0.000000,0.000000}%
\pgfsetstrokecolor{currentstroke}%
\pgfsetdash{}{0pt}%
\pgfpathmoveto{\pgfqpoint{1.131573in}{0.681712in}}%
\pgfpathlineto{\pgfqpoint{1.135119in}{0.687818in}}%
\pgfpathlineto{\pgfqpoint{1.138681in}{0.693950in}}%
\pgfpathlineto{\pgfqpoint{1.150691in}{0.695737in}}%
\pgfpathlineto{\pgfqpoint{1.162415in}{0.697853in}}%
\pgfpathlineto{\pgfqpoint{1.157114in}{0.691440in}}%
\pgfpathlineto{\pgfqpoint{1.151835in}{0.685055in}}%
\pgfpathlineto{\pgfqpoint{1.141825in}{0.683242in}}%
\pgfpathlineto{\pgfqpoint{1.131573in}{0.681712in}}%
\pgfpathclose%
\pgfusepath{fill}%
\end{pgfscope}%
\begin{pgfscope}%
\pgfpathrectangle{\pgfqpoint{0.407592in}{0.222904in}}{\pgfqpoint{1.411163in}{1.411163in}}%
\pgfusepath{clip}%
\pgfsetbuttcap%
\pgfsetroundjoin%
\definecolor{currentfill}{rgb}{0.749020,0.772549,0.792157}%
\pgfsetfillcolor{currentfill}%
\pgfsetfillopacity{0.160000}%
\pgfsetlinewidth{0.000000pt}%
\definecolor{currentstroke}{rgb}{0.000000,0.000000,0.000000}%
\pgfsetstrokecolor{currentstroke}%
\pgfsetdash{}{0pt}%
\pgfpathmoveto{\pgfqpoint{1.243536in}{0.732416in}}%
\pgfpathlineto{\pgfqpoint{1.253408in}{0.740265in}}%
\pgfpathlineto{\pgfqpoint{1.263315in}{0.748141in}}%
\pgfpathlineto{\pgfqpoint{1.274050in}{0.753345in}}%
\pgfpathlineto{\pgfqpoint{1.283933in}{0.758812in}}%
\pgfpathlineto{\pgfqpoint{1.272878in}{0.750350in}}%
\pgfpathlineto{\pgfqpoint{1.261858in}{0.741915in}}%
\pgfpathlineto{\pgfqpoint{1.253072in}{0.737047in}}%
\pgfpathlineto{\pgfqpoint{1.243536in}{0.732416in}}%
\pgfpathclose%
\pgfusepath{fill}%
\end{pgfscope}%
\begin{pgfscope}%
\pgfpathrectangle{\pgfqpoint{0.407592in}{0.222904in}}{\pgfqpoint{1.411163in}{1.411163in}}%
\pgfusepath{clip}%
\pgfsetbuttcap%
\pgfsetroundjoin%
\definecolor{currentfill}{rgb}{0.749020,0.772549,0.792157}%
\pgfsetfillcolor{currentfill}%
\pgfsetfillopacity{0.160000}%
\pgfsetlinewidth{0.000000pt}%
\definecolor{currentstroke}{rgb}{0.000000,0.000000,0.000000}%
\pgfsetstrokecolor{currentstroke}%
\pgfsetdash{}{0pt}%
\pgfpathmoveto{\pgfqpoint{0.450032in}{1.196241in}}%
\pgfpathlineto{\pgfqpoint{0.436384in}{1.208728in}}%
\pgfpathlineto{\pgfqpoint{0.422723in}{1.221226in}}%
\pgfpathlineto{\pgfqpoint{0.413663in}{1.204934in}}%
\pgfpathlineto{\pgfqpoint{0.407462in}{1.188415in}}%
\pgfpathlineto{\pgfqpoint{0.421586in}{1.176460in}}%
\pgfpathlineto{\pgfqpoint{0.435692in}{1.164522in}}%
\pgfpathlineto{\pgfqpoint{0.441489in}{1.180488in}}%
\pgfpathlineto{\pgfqpoint{0.450032in}{1.196241in}}%
\pgfpathclose%
\pgfusepath{fill}%
\end{pgfscope}%
\begin{pgfscope}%
\pgfpathrectangle{\pgfqpoint{0.407592in}{0.222904in}}{\pgfqpoint{1.411163in}{1.411163in}}%
\pgfusepath{clip}%
\pgfsetbuttcap%
\pgfsetroundjoin%
\definecolor{currentfill}{rgb}{0.749020,0.772549,0.792157}%
\pgfsetfillcolor{currentfill}%
\pgfsetfillopacity{0.160000}%
\pgfsetlinewidth{0.000000pt}%
\definecolor{currentstroke}{rgb}{0.000000,0.000000,0.000000}%
\pgfsetstrokecolor{currentstroke}%
\pgfsetdash{}{0pt}%
\pgfpathmoveto{\pgfqpoint{1.768602in}{1.191744in}}%
\pgfpathlineto{\pgfqpoint{1.782640in}{1.203771in}}%
\pgfpathlineto{\pgfqpoint{1.796697in}{1.215815in}}%
\pgfpathlineto{\pgfqpoint{1.793644in}{1.232834in}}%
\pgfpathlineto{\pgfqpoint{1.787575in}{1.249598in}}%
\pgfpathlineto{\pgfqpoint{1.773855in}{1.237024in}}%
\pgfpathlineto{\pgfqpoint{1.760147in}{1.224461in}}%
\pgfpathlineto{\pgfqpoint{1.765824in}{1.208223in}}%
\pgfpathlineto{\pgfqpoint{1.768602in}{1.191744in}}%
\pgfpathclose%
\pgfusepath{fill}%
\end{pgfscope}%
\begin{pgfscope}%
\pgfpathrectangle{\pgfqpoint{0.407592in}{0.222904in}}{\pgfqpoint{1.411163in}{1.411163in}}%
\pgfusepath{clip}%
\pgfsetbuttcap%
\pgfsetroundjoin%
\definecolor{currentfill}{rgb}{0.749020,0.772549,0.792157}%
\pgfsetfillcolor{currentfill}%
\pgfsetfillopacity{0.160000}%
\pgfsetlinewidth{0.000000pt}%
\definecolor{currentstroke}{rgb}{0.000000,0.000000,0.000000}%
\pgfsetstrokecolor{currentstroke}%
\pgfsetdash{}{0pt}%
\pgfpathmoveto{\pgfqpoint{1.570115in}{0.999993in}}%
\pgfpathlineto{\pgfqpoint{1.583889in}{1.011171in}}%
\pgfpathlineto{\pgfqpoint{1.597687in}{1.022369in}}%
\pgfpathlineto{\pgfqpoint{1.600788in}{1.035548in}}%
\pgfpathlineto{\pgfqpoint{1.601625in}{1.048682in}}%
\pgfpathlineto{\pgfqpoint{1.587832in}{1.036865in}}%
\pgfpathlineto{\pgfqpoint{1.574058in}{1.025063in}}%
\pgfpathlineto{\pgfqpoint{1.573157in}{1.012546in}}%
\pgfpathlineto{\pgfqpoint{1.570115in}{0.999993in}}%
\pgfpathclose%
\pgfusepath{fill}%
\end{pgfscope}%
\begin{pgfscope}%
\pgfpathrectangle{\pgfqpoint{0.407592in}{0.222904in}}{\pgfqpoint{1.411163in}{1.411163in}}%
\pgfusepath{clip}%
\pgfsetbuttcap%
\pgfsetroundjoin%
\definecolor{currentfill}{rgb}{0.749020,0.772549,0.792157}%
\pgfsetfillcolor{currentfill}%
\pgfsetfillopacity{0.160000}%
\pgfsetlinewidth{0.000000pt}%
\definecolor{currentstroke}{rgb}{0.000000,0.000000,0.000000}%
\pgfsetstrokecolor{currentstroke}%
\pgfsetdash{}{0pt}%
\pgfpathmoveto{\pgfqpoint{1.450479in}{0.891485in}}%
\pgfpathlineto{\pgfqpoint{1.463761in}{0.901843in}}%
\pgfpathlineto{\pgfqpoint{1.477073in}{0.912225in}}%
\pgfpathlineto{\pgfqpoint{1.483419in}{0.922744in}}%
\pgfpathlineto{\pgfqpoint{1.487990in}{0.933343in}}%
\pgfpathlineto{\pgfqpoint{1.474388in}{0.922304in}}%
\pgfpathlineto{\pgfqpoint{1.460810in}{0.911284in}}%
\pgfpathlineto{\pgfqpoint{1.456470in}{0.901345in}}%
\pgfpathlineto{\pgfqpoint{1.450479in}{0.891485in}}%
\pgfpathclose%
\pgfusepath{fill}%
\end{pgfscope}%
\begin{pgfscope}%
\pgfpathrectangle{\pgfqpoint{0.407592in}{0.222904in}}{\pgfqpoint{1.411163in}{1.411163in}}%
\pgfusepath{clip}%
\pgfsetbuttcap%
\pgfsetroundjoin%
\definecolor{currentfill}{rgb}{0.749020,0.772549,0.792157}%
\pgfsetfillcolor{currentfill}%
\pgfsetfillopacity{0.160000}%
\pgfsetlinewidth{0.000000pt}%
\definecolor{currentstroke}{rgb}{0.000000,0.000000,0.000000}%
\pgfsetstrokecolor{currentstroke}%
\pgfsetdash{}{0pt}%
\pgfpathmoveto{\pgfqpoint{0.954332in}{0.712235in}}%
\pgfpathlineto{\pgfqpoint{0.944352in}{0.719704in}}%
\pgfpathlineto{\pgfqpoint{0.934335in}{0.727200in}}%
\pgfpathlineto{\pgfqpoint{0.944514in}{0.723167in}}%
\pgfpathlineto{\pgfqpoint{0.955364in}{0.719437in}}%
\pgfpathlineto{\pgfqpoint{0.964048in}{0.712425in}}%
\pgfpathlineto{\pgfqpoint{0.972697in}{0.705441in}}%
\pgfpathlineto{\pgfqpoint{0.963224in}{0.708704in}}%
\pgfpathlineto{\pgfqpoint{0.954332in}{0.712235in}}%
\pgfpathclose%
\pgfusepath{fill}%
\end{pgfscope}%
\begin{pgfscope}%
\pgfpathrectangle{\pgfqpoint{0.407592in}{0.222904in}}{\pgfqpoint{1.411163in}{1.411163in}}%
\pgfusepath{clip}%
\pgfsetbuttcap%
\pgfsetroundjoin%
\definecolor{currentfill}{rgb}{0.749020,0.772549,0.792157}%
\pgfsetfillcolor{currentfill}%
\pgfsetfillopacity{0.160000}%
\pgfsetlinewidth{0.000000pt}%
\definecolor{currentstroke}{rgb}{0.000000,0.000000,0.000000}%
\pgfsetstrokecolor{currentstroke}%
\pgfsetdash{}{0pt}%
\pgfpathmoveto{\pgfqpoint{1.026478in}{0.682912in}}%
\pgfpathlineto{\pgfqpoint{1.021012in}{0.689118in}}%
\pgfpathlineto{\pgfqpoint{1.015522in}{0.695351in}}%
\pgfpathlineto{\pgfqpoint{1.027217in}{0.693632in}}%
\pgfpathlineto{\pgfqpoint{1.039203in}{0.692254in}}%
\pgfpathlineto{\pgfqpoint{1.042956in}{0.686244in}}%
\pgfpathlineto{\pgfqpoint{1.046692in}{0.680260in}}%
\pgfpathlineto{\pgfqpoint{1.036462in}{0.681440in}}%
\pgfpathlineto{\pgfqpoint{1.026478in}{0.682912in}}%
\pgfpathclose%
\pgfusepath{fill}%
\end{pgfscope}%
\begin{pgfscope}%
\pgfpathrectangle{\pgfqpoint{0.407592in}{0.222904in}}{\pgfqpoint{1.411163in}{1.411163in}}%
\pgfusepath{clip}%
\pgfsetbuttcap%
\pgfsetroundjoin%
\definecolor{currentfill}{rgb}{0.749020,0.772549,0.792157}%
\pgfsetfillcolor{currentfill}%
\pgfsetfillopacity{0.160000}%
\pgfsetlinewidth{0.000000pt}%
\definecolor{currentstroke}{rgb}{0.000000,0.000000,0.000000}%
\pgfsetstrokecolor{currentstroke}%
\pgfsetdash{}{0pt}%
\pgfpathmoveto{\pgfqpoint{1.110537in}{0.679531in}}%
\pgfpathlineto{\pgfqpoint{1.112280in}{0.685453in}}%
\pgfpathlineto{\pgfqpoint{1.114031in}{0.691402in}}%
\pgfpathlineto{\pgfqpoint{1.126442in}{0.692503in}}%
\pgfpathlineto{\pgfqpoint{1.138681in}{0.693950in}}%
\pgfpathlineto{\pgfqpoint{1.135119in}{0.687818in}}%
\pgfpathlineto{\pgfqpoint{1.131573in}{0.681712in}}%
\pgfpathlineto{\pgfqpoint{1.121128in}{0.680473in}}%
\pgfpathlineto{\pgfqpoint{1.110537in}{0.679531in}}%
\pgfpathclose%
\pgfusepath{fill}%
\end{pgfscope}%
\begin{pgfscope}%
\pgfpathrectangle{\pgfqpoint{0.407592in}{0.222904in}}{\pgfqpoint{1.411163in}{1.411163in}}%
\pgfusepath{clip}%
\pgfsetbuttcap%
\pgfsetroundjoin%
\definecolor{currentfill}{rgb}{0.749020,0.772549,0.792157}%
\pgfsetfillcolor{currentfill}%
\pgfsetfillopacity{0.160000}%
\pgfsetlinewidth{0.000000pt}%
\definecolor{currentstroke}{rgb}{0.000000,0.000000,0.000000}%
\pgfsetstrokecolor{currentstroke}%
\pgfsetdash{}{0pt}%
\pgfpathmoveto{\pgfqpoint{1.283933in}{0.758812in}}%
\pgfpathlineto{\pgfqpoint{1.295023in}{0.767302in}}%
\pgfpathlineto{\pgfqpoint{1.306149in}{0.775819in}}%
\pgfpathlineto{\pgfqpoint{1.316125in}{0.782146in}}%
\pgfpathlineto{\pgfqpoint{1.304505in}{0.773318in}}%
\pgfpathlineto{\pgfqpoint{1.292921in}{0.764517in}}%
\pgfusepath{fill}%
\end{pgfscope}%
\begin{pgfscope}%
\pgfpathrectangle{\pgfqpoint{0.407592in}{0.222904in}}{\pgfqpoint{1.411163in}{1.411163in}}%
\pgfusepath{clip}%
\pgfsetbuttcap%
\pgfsetroundjoin%
\definecolor{currentfill}{rgb}{0.749020,0.772549,0.792157}%
\pgfsetfillcolor{currentfill}%
\pgfsetfillopacity{0.160000}%
\pgfsetlinewidth{0.000000pt}%
\definecolor{currentstroke}{rgb}{0.000000,0.000000,0.000000}%
\pgfsetstrokecolor{currentstroke}%
\pgfsetdash{}{0pt}%
\pgfpathmoveto{\pgfqpoint{0.713628in}{0.909286in}}%
\pgfpathlineto{\pgfqpoint{0.699821in}{0.920239in}}%
\pgfpathlineto{\pgfqpoint{0.685989in}{0.931212in}}%
\pgfpathlineto{\pgfqpoint{0.688537in}{0.920626in}}%
\pgfpathlineto{\pgfqpoint{0.692881in}{0.910128in}}%
\pgfpathlineto{\pgfqpoint{0.706562in}{0.899812in}}%
\pgfpathlineto{\pgfqpoint{0.720213in}{0.889520in}}%
\pgfpathlineto{\pgfqpoint{0.716086in}{0.899359in}}%
\pgfpathlineto{\pgfqpoint{0.713628in}{0.909286in}}%
\pgfpathclose%
\pgfusepath{fill}%
\end{pgfscope}%
\begin{pgfscope}%
\pgfpathrectangle{\pgfqpoint{0.407592in}{0.222904in}}{\pgfqpoint{1.411163in}{1.411163in}}%
\pgfusepath{clip}%
\pgfsetbuttcap%
\pgfsetroundjoin%
\definecolor{currentfill}{rgb}{0.749020,0.772549,0.792157}%
\pgfsetfillcolor{currentfill}%
\pgfsetfillopacity{0.160000}%
\pgfsetlinewidth{0.000000pt}%
\definecolor{currentstroke}{rgb}{0.000000,0.000000,0.000000}%
\pgfsetstrokecolor{currentstroke}%
\pgfsetdash{}{0pt}%
\pgfpathmoveto{\pgfqpoint{1.316125in}{0.782146in}}%
\pgfpathlineto{\pgfqpoint{1.327780in}{0.791000in}}%
\pgfpathlineto{\pgfqpoint{1.339471in}{0.799881in}}%
\pgfpathlineto{\pgfqpoint{1.349275in}{0.807077in}}%
\pgfpathlineto{\pgfqpoint{1.357886in}{0.814488in}}%
\pgfpathlineto{\pgfqpoint{1.345383in}{0.804961in}}%
\pgfpathlineto{\pgfqpoint{1.332913in}{0.795459in}}%
\pgfpathlineto{\pgfqpoint{1.325058in}{0.788703in}}%
\pgfpathlineto{\pgfqpoint{1.316125in}{0.782146in}}%
\pgfpathclose%
\pgfusepath{fill}%
\end{pgfscope}%
\begin{pgfscope}%
\pgfpathrectangle{\pgfqpoint{0.407592in}{0.222904in}}{\pgfqpoint{1.411163in}{1.411163in}}%
\pgfusepath{clip}%
\pgfsetbuttcap%
\pgfsetroundjoin%
\definecolor{currentfill}{rgb}{0.749020,0.772549,0.792157}%
\pgfsetfillcolor{currentfill}%
\pgfsetfillopacity{0.160000}%
\pgfsetlinewidth{0.000000pt}%
\definecolor{currentstroke}{rgb}{0.000000,0.000000,0.000000}%
\pgfsetstrokecolor{currentstroke}%
\pgfsetdash{}{0pt}%
\pgfpathmoveto{\pgfqpoint{0.603422in}{1.022558in}}%
\pgfpathlineto{\pgfqpoint{0.589552in}{1.034297in}}%
\pgfpathlineto{\pgfqpoint{0.575662in}{1.046054in}}%
\pgfpathlineto{\pgfqpoint{0.573906in}{1.032905in}}%
\pgfpathlineto{\pgfqpoint{0.574411in}{1.019725in}}%
\pgfpathlineto{\pgfqpoint{0.588448in}{1.008590in}}%
\pgfpathlineto{\pgfqpoint{0.602459in}{0.997475in}}%
\pgfpathlineto{\pgfqpoint{0.601871in}{1.010029in}}%
\pgfpathlineto{\pgfqpoint{0.603422in}{1.022558in}}%
\pgfpathclose%
\pgfusepath{fill}%
\end{pgfscope}%
\begin{pgfscope}%
\pgfpathrectangle{\pgfqpoint{0.407592in}{0.222904in}}{\pgfqpoint{1.411163in}{1.411163in}}%
\pgfusepath{clip}%
\pgfsetbuttcap%
\pgfsetroundjoin%
\definecolor{currentfill}{rgb}{0.749020,0.772549,0.792157}%
\pgfsetfillcolor{currentfill}%
\pgfsetfillopacity{0.160000}%
\pgfsetlinewidth{0.000000pt}%
\definecolor{currentstroke}{rgb}{0.000000,0.000000,0.000000}%
\pgfsetstrokecolor{currentstroke}%
\pgfsetdash{}{0pt}%
\pgfpathmoveto{\pgfqpoint{1.046692in}{0.680260in}}%
\pgfpathlineto{\pgfqpoint{1.042956in}{0.686244in}}%
\pgfpathlineto{\pgfqpoint{1.039203in}{0.692254in}}%
\pgfpathlineto{\pgfqpoint{1.051424in}{0.691224in}}%
\pgfpathlineto{\pgfqpoint{1.063823in}{0.690546in}}%
\pgfpathlineto{\pgfqpoint{1.065766in}{0.684658in}}%
\pgfpathlineto{\pgfqpoint{1.067700in}{0.678798in}}%
\pgfpathlineto{\pgfqpoint{1.057121in}{0.679378in}}%
\pgfpathlineto{\pgfqpoint{1.046692in}{0.680260in}}%
\pgfpathclose%
\pgfusepath{fill}%
\end{pgfscope}%
\begin{pgfscope}%
\pgfpathrectangle{\pgfqpoint{0.407592in}{0.222904in}}{\pgfqpoint{1.411163in}{1.411163in}}%
\pgfusepath{clip}%
\pgfsetbuttcap%
\pgfsetroundjoin%
\definecolor{currentfill}{rgb}{0.749020,0.772549,0.792157}%
\pgfsetfillcolor{currentfill}%
\pgfsetfillopacity{0.160000}%
\pgfsetlinewidth{0.000000pt}%
\definecolor{currentstroke}{rgb}{0.000000,0.000000,0.000000}%
\pgfsetstrokecolor{currentstroke}%
\pgfsetdash{}{0pt}%
\pgfpathmoveto{\pgfqpoint{1.184795in}{0.703034in}}%
\pgfpathlineto{\pgfqpoint{1.191759in}{0.709846in}}%
\pgfpathlineto{\pgfqpoint{1.198752in}{0.716687in}}%
\pgfpathlineto{\pgfqpoint{1.210847in}{0.720160in}}%
\pgfpathlineto{\pgfqpoint{1.222376in}{0.723952in}}%
\pgfpathlineto{\pgfqpoint{1.213878in}{0.716658in}}%
\pgfpathlineto{\pgfqpoint{1.205412in}{0.709392in}}%
\pgfpathlineto{\pgfqpoint{1.195348in}{0.706073in}}%
\pgfpathlineto{\pgfqpoint{1.184795in}{0.703034in}}%
\pgfpathclose%
\pgfusepath{fill}%
\end{pgfscope}%
\begin{pgfscope}%
\pgfpathrectangle{\pgfqpoint{0.407592in}{0.222904in}}{\pgfqpoint{1.411163in}{1.411163in}}%
\pgfusepath{clip}%
\pgfsetbuttcap%
\pgfsetroundjoin%
\definecolor{currentfill}{rgb}{0.749020,0.772549,0.792157}%
\pgfsetfillcolor{currentfill}%
\pgfsetfillopacity{0.160000}%
\pgfsetlinewidth{0.000000pt}%
\definecolor{currentstroke}{rgb}{0.000000,0.000000,0.000000}%
\pgfsetstrokecolor{currentstroke}%
\pgfsetdash{}{0pt}%
\pgfpathmoveto{\pgfqpoint{0.422723in}{1.221226in}}%
\pgfpathlineto{\pgfqpoint{0.409049in}{1.233736in}}%
\pgfpathlineto{\pgfqpoint{0.395362in}{1.246258in}}%
\pgfpathlineto{\pgfqpoint{0.385772in}{1.229438in}}%
\pgfpathlineto{\pgfqpoint{0.379152in}{1.212375in}}%
\pgfpathlineto{\pgfqpoint{0.393317in}{1.200387in}}%
\pgfpathlineto{\pgfqpoint{0.407462in}{1.188415in}}%
\pgfpathlineto{\pgfqpoint{0.413663in}{1.204934in}}%
\pgfpathlineto{\pgfqpoint{0.422723in}{1.221226in}}%
\pgfpathclose%
\pgfusepath{fill}%
\end{pgfscope}%
\begin{pgfscope}%
\pgfpathrectangle{\pgfqpoint{0.407592in}{0.222904in}}{\pgfqpoint{1.411163in}{1.411163in}}%
\pgfusepath{clip}%
\pgfsetbuttcap%
\pgfsetroundjoin%
\definecolor{currentfill}{rgb}{0.749020,0.772549,0.792157}%
\pgfsetfillcolor{currentfill}%
\pgfsetfillopacity{0.160000}%
\pgfsetlinewidth{0.000000pt}%
\definecolor{currentstroke}{rgb}{0.000000,0.000000,0.000000}%
\pgfsetstrokecolor{currentstroke}%
\pgfsetdash{}{0pt}%
\pgfpathmoveto{\pgfqpoint{1.796697in}{1.215815in}}%
\pgfpathlineto{\pgfqpoint{1.810773in}{1.227874in}}%
\pgfpathlineto{\pgfqpoint{1.824868in}{1.239951in}}%
\pgfpathlineto{\pgfqpoint{1.821526in}{1.257501in}}%
\pgfpathlineto{\pgfqpoint{1.815054in}{1.274780in}}%
\pgfpathlineto{\pgfqpoint{1.801308in}{1.262183in}}%
\pgfpathlineto{\pgfqpoint{1.787575in}{1.249598in}}%
\pgfpathlineto{\pgfqpoint{1.793644in}{1.232834in}}%
\pgfpathlineto{\pgfqpoint{1.796697in}{1.215815in}}%
\pgfpathclose%
\pgfusepath{fill}%
\end{pgfscope}%
\begin{pgfscope}%
\pgfpathrectangle{\pgfqpoint{0.407592in}{0.222904in}}{\pgfqpoint{1.411163in}{1.411163in}}%
\pgfusepath{clip}%
\pgfsetbuttcap%
\pgfsetroundjoin%
\definecolor{currentfill}{rgb}{0.749020,0.772549,0.792157}%
\pgfsetfillcolor{currentfill}%
\pgfsetfillopacity{0.160000}%
\pgfsetlinewidth{0.000000pt}%
\definecolor{currentstroke}{rgb}{0.000000,0.000000,0.000000}%
\pgfsetstrokecolor{currentstroke}%
\pgfsetdash{}{0pt}%
\pgfpathmoveto{\pgfqpoint{1.089114in}{0.678553in}}%
\pgfpathlineto{\pgfqpoint{1.089018in}{0.684393in}}%
\pgfpathlineto{\pgfqpoint{1.088921in}{0.690260in}}%
\pgfpathlineto{\pgfqpoint{1.101504in}{0.690653in}}%
\pgfpathlineto{\pgfqpoint{1.114031in}{0.691402in}}%
\pgfpathlineto{\pgfqpoint{1.112280in}{0.685453in}}%
\pgfpathlineto{\pgfqpoint{1.110537in}{0.679531in}}%
\pgfpathlineto{\pgfqpoint{1.099849in}{0.678890in}}%
\pgfpathlineto{\pgfqpoint{1.089114in}{0.678553in}}%
\pgfpathclose%
\pgfusepath{fill}%
\end{pgfscope}%
\begin{pgfscope}%
\pgfpathrectangle{\pgfqpoint{0.407592in}{0.222904in}}{\pgfqpoint{1.411163in}{1.411163in}}%
\pgfusepath{clip}%
\pgfsetbuttcap%
\pgfsetroundjoin%
\definecolor{currentfill}{rgb}{0.749020,0.772549,0.792157}%
\pgfsetfillcolor{currentfill}%
\pgfsetfillopacity{0.160000}%
\pgfsetlinewidth{0.000000pt}%
\definecolor{currentstroke}{rgb}{0.000000,0.000000,0.000000}%
\pgfsetstrokecolor{currentstroke}%
\pgfsetdash{}{0pt}%
\pgfpathmoveto{\pgfqpoint{0.838112in}{0.794089in}}%
\pgfpathlineto{\pgfqpoint{0.825152in}{0.803525in}}%
\pgfpathlineto{\pgfqpoint{0.812157in}{0.812986in}}%
\pgfpathlineto{\pgfqpoint{0.819401in}{0.805615in}}%
\pgfpathlineto{\pgfqpoint{0.827881in}{0.798467in}}%
\pgfpathlineto{\pgfqpoint{0.840191in}{0.789648in}}%
\pgfpathlineto{\pgfqpoint{0.852464in}{0.780857in}}%
\pgfpathlineto{\pgfqpoint{0.844730in}{0.787370in}}%
\pgfpathlineto{\pgfqpoint{0.838112in}{0.794089in}}%
\pgfpathclose%
\pgfusepath{fill}%
\end{pgfscope}%
\begin{pgfscope}%
\pgfpathrectangle{\pgfqpoint{0.407592in}{0.222904in}}{\pgfqpoint{1.411163in}{1.411163in}}%
\pgfusepath{clip}%
\pgfsetbuttcap%
\pgfsetroundjoin%
\definecolor{currentfill}{rgb}{0.749020,0.772549,0.792157}%
\pgfsetfillcolor{currentfill}%
\pgfsetfillopacity{0.160000}%
\pgfsetlinewidth{0.000000pt}%
\definecolor{currentstroke}{rgb}{0.000000,0.000000,0.000000}%
\pgfsetstrokecolor{currentstroke}%
\pgfsetdash{}{0pt}%
\pgfpathmoveto{\pgfqpoint{1.067700in}{0.678798in}}%
\pgfpathlineto{\pgfqpoint{1.065766in}{0.684658in}}%
\pgfpathlineto{\pgfqpoint{1.063823in}{0.690546in}}%
\pgfpathlineto{\pgfqpoint{1.076341in}{0.690224in}}%
\pgfpathlineto{\pgfqpoint{1.088921in}{0.690260in}}%
\pgfpathlineto{\pgfqpoint{1.089018in}{0.684393in}}%
\pgfpathlineto{\pgfqpoint{1.089114in}{0.678553in}}%
\pgfpathlineto{\pgfqpoint{1.078381in}{0.678522in}}%
\pgfpathlineto{\pgfqpoint{1.067700in}{0.678798in}}%
\pgfpathclose%
\pgfusepath{fill}%
\end{pgfscope}%
\begin{pgfscope}%
\pgfpathrectangle{\pgfqpoint{0.407592in}{0.222904in}}{\pgfqpoint{1.411163in}{1.411163in}}%
\pgfusepath{clip}%
\pgfsetbuttcap%
\pgfsetroundjoin%
\definecolor{currentfill}{rgb}{0.749020,0.772549,0.792157}%
\pgfsetfillcolor{currentfill}%
\pgfsetfillopacity{0.160000}%
\pgfsetlinewidth{0.000000pt}%
\definecolor{currentstroke}{rgb}{0.000000,0.000000,0.000000}%
\pgfsetstrokecolor{currentstroke}%
\pgfsetdash{}{0pt}%
\pgfpathmoveto{\pgfqpoint{1.382990in}{0.833618in}}%
\pgfpathlineto{\pgfqpoint{1.395592in}{0.843221in}}%
\pgfpathlineto{\pgfqpoint{1.408228in}{0.852849in}}%
\pgfpathlineto{\pgfqpoint{1.416863in}{0.861762in}}%
\pgfpathlineto{\pgfqpoint{1.424004in}{0.870837in}}%
\pgfpathlineto{\pgfqpoint{1.410810in}{0.860548in}}%
\pgfpathlineto{\pgfqpoint{1.397646in}{0.850281in}}%
\pgfpathlineto{\pgfqpoint{1.391004in}{0.841873in}}%
\pgfpathlineto{\pgfqpoint{1.382990in}{0.833618in}}%
\pgfpathclose%
\pgfusepath{fill}%
\end{pgfscope}%
\begin{pgfscope}%
\pgfpathrectangle{\pgfqpoint{0.407592in}{0.222904in}}{\pgfqpoint{1.411163in}{1.411163in}}%
\pgfusepath{clip}%
\pgfsetbuttcap%
\pgfsetroundjoin%
\definecolor{currentfill}{rgb}{0.749020,0.772549,0.792157}%
\pgfsetfillcolor{currentfill}%
\pgfsetfillopacity{0.160000}%
\pgfsetlinewidth{0.000000pt}%
\definecolor{currentstroke}{rgb}{0.000000,0.000000,0.000000}%
\pgfsetstrokecolor{currentstroke}%
\pgfsetdash{}{0pt}%
\pgfpathmoveto{\pgfqpoint{0.916170in}{0.736098in}}%
\pgfpathlineto{\pgfqpoint{0.904974in}{0.744175in}}%
\pgfpathlineto{\pgfqpoint{0.893740in}{0.752279in}}%
\pgfpathlineto{\pgfqpoint{0.903534in}{0.747131in}}%
\pgfpathlineto{\pgfqpoint{0.914188in}{0.742278in}}%
\pgfpathlineto{\pgfqpoint{0.924281in}{0.734725in}}%
\pgfpathlineto{\pgfqpoint{0.934335in}{0.727200in}}%
\pgfpathlineto{\pgfqpoint{0.924874in}{0.731518in}}%
\pgfpathlineto{\pgfqpoint{0.916170in}{0.736098in}}%
\pgfpathclose%
\pgfusepath{fill}%
\end{pgfscope}%
\begin{pgfscope}%
\pgfpathrectangle{\pgfqpoint{0.407592in}{0.222904in}}{\pgfqpoint{1.411163in}{1.411163in}}%
\pgfusepath{clip}%
\pgfsetbuttcap%
\pgfsetroundjoin%
\definecolor{currentfill}{rgb}{0.749020,0.772549,0.792157}%
\pgfsetfillcolor{currentfill}%
\pgfsetfillopacity{0.160000}%
\pgfsetlinewidth{0.000000pt}%
\definecolor{currentstroke}{rgb}{0.000000,0.000000,0.000000}%
\pgfsetstrokecolor{currentstroke}%
\pgfsetdash{}{0pt}%
\pgfpathmoveto{\pgfqpoint{1.597687in}{1.022369in}}%
\pgfpathlineto{\pgfqpoint{1.611511in}{1.033588in}}%
\pgfpathlineto{\pgfqpoint{1.625359in}{1.044827in}}%
\pgfpathlineto{\pgfqpoint{1.628507in}{1.058622in}}%
\pgfpathlineto{\pgfqpoint{1.629267in}{1.072365in}}%
\pgfpathlineto{\pgfqpoint{1.615437in}{1.060516in}}%
\pgfpathlineto{\pgfqpoint{1.601625in}{1.048682in}}%
\pgfpathlineto{\pgfqpoint{1.600788in}{1.035548in}}%
\pgfpathlineto{\pgfqpoint{1.597687in}{1.022369in}}%
\pgfpathclose%
\pgfusepath{fill}%
\end{pgfscope}%
\begin{pgfscope}%
\pgfpathrectangle{\pgfqpoint{0.407592in}{0.222904in}}{\pgfqpoint{1.411163in}{1.411163in}}%
\pgfusepath{clip}%
\pgfsetbuttcap%
\pgfsetroundjoin%
\definecolor{currentfill}{rgb}{0.749020,0.772549,0.792157}%
\pgfsetfillcolor{currentfill}%
\pgfsetfillopacity{0.160000}%
\pgfsetlinewidth{0.000000pt}%
\definecolor{currentstroke}{rgb}{0.000000,0.000000,0.000000}%
\pgfsetstrokecolor{currentstroke}%
\pgfsetdash{}{0pt}%
\pgfpathmoveto{\pgfqpoint{0.876896in}{0.763355in}}%
\pgfpathlineto{\pgfqpoint{0.864698in}{0.772093in}}%
\pgfpathlineto{\pgfqpoint{0.852464in}{0.780857in}}%
\pgfpathlineto{\pgfqpoint{0.861284in}{0.774581in}}%
\pgfpathlineto{\pgfqpoint{0.871156in}{0.768571in}}%
\pgfpathlineto{\pgfqpoint{0.882467in}{0.760411in}}%
\pgfpathlineto{\pgfqpoint{0.893740in}{0.752279in}}%
\pgfpathlineto{\pgfqpoint{0.884849in}{0.757696in}}%
\pgfpathlineto{\pgfqpoint{0.876896in}{0.763355in}}%
\pgfpathclose%
\pgfusepath{fill}%
\end{pgfscope}%
\begin{pgfscope}%
\pgfpathrectangle{\pgfqpoint{0.407592in}{0.222904in}}{\pgfqpoint{1.411163in}{1.411163in}}%
\pgfusepath{clip}%
\pgfsetbuttcap%
\pgfsetroundjoin%
\definecolor{currentfill}{rgb}{0.749020,0.772549,0.792157}%
\pgfsetfillcolor{currentfill}%
\pgfsetfillopacity{0.160000}%
\pgfsetlinewidth{0.000000pt}%
\definecolor{currentstroke}{rgb}{0.000000,0.000000,0.000000}%
\pgfsetstrokecolor{currentstroke}%
\pgfsetdash{}{0pt}%
\pgfpathmoveto{\pgfqpoint{0.774506in}{0.848584in}}%
\pgfpathlineto{\pgfqpoint{0.760979in}{0.858783in}}%
\pgfpathlineto{\pgfqpoint{0.747421in}{0.869006in}}%
\pgfpathlineto{\pgfqpoint{0.752860in}{0.859959in}}%
\pgfpathlineto{\pgfqpoint{0.759828in}{0.851084in}}%
\pgfpathlineto{\pgfqpoint{0.772963in}{0.841521in}}%
\pgfpathlineto{\pgfqpoint{0.786063in}{0.831984in}}%
\pgfpathlineto{\pgfqpoint{0.779583in}{0.840203in}}%
\pgfpathlineto{\pgfqpoint{0.774506in}{0.848584in}}%
\pgfpathclose%
\pgfusepath{fill}%
\end{pgfscope}%
\begin{pgfscope}%
\pgfpathrectangle{\pgfqpoint{0.407592in}{0.222904in}}{\pgfqpoint{1.411163in}{1.411163in}}%
\pgfusepath{clip}%
\pgfsetbuttcap%
\pgfsetroundjoin%
\definecolor{currentfill}{rgb}{0.749020,0.772549,0.792157}%
\pgfsetfillcolor{currentfill}%
\pgfsetfillopacity{0.160000}%
\pgfsetlinewidth{0.000000pt}%
\definecolor{currentstroke}{rgb}{0.000000,0.000000,0.000000}%
\pgfsetstrokecolor{currentstroke}%
\pgfsetdash{}{0pt}%
\pgfpathmoveto{\pgfqpoint{0.972697in}{0.705441in}}%
\pgfpathlineto{\pgfqpoint{0.964048in}{0.712425in}}%
\pgfpathlineto{\pgfqpoint{0.955364in}{0.719437in}}%
\pgfpathlineto{\pgfqpoint{0.966838in}{0.716030in}}%
\pgfpathlineto{\pgfqpoint{0.978884in}{0.712962in}}%
\pgfpathlineto{\pgfqpoint{0.986066in}{0.706355in}}%
\pgfpathlineto{\pgfqpoint{0.993217in}{0.699776in}}%
\pgfpathlineto{\pgfqpoint{0.982710in}{0.702460in}}%
\pgfpathlineto{\pgfqpoint{0.972697in}{0.705441in}}%
\pgfpathclose%
\pgfusepath{fill}%
\end{pgfscope}%
\begin{pgfscope}%
\pgfpathrectangle{\pgfqpoint{0.407592in}{0.222904in}}{\pgfqpoint{1.411163in}{1.411163in}}%
\pgfusepath{clip}%
\pgfsetbuttcap%
\pgfsetroundjoin%
\definecolor{currentfill}{rgb}{0.749020,0.772549,0.792157}%
\pgfsetfillcolor{currentfill}%
\pgfsetfillopacity{0.160000}%
\pgfsetlinewidth{0.000000pt}%
\definecolor{currentstroke}{rgb}{0.000000,0.000000,0.000000}%
\pgfsetstrokecolor{currentstroke}%
\pgfsetdash{}{0pt}%
\pgfpathmoveto{\pgfqpoint{1.477073in}{0.912225in}}%
\pgfpathlineto{\pgfqpoint{1.490414in}{0.922629in}}%
\pgfpathlineto{\pgfqpoint{1.503785in}{0.933057in}}%
\pgfpathlineto{\pgfqpoint{1.510477in}{0.944229in}}%
\pgfpathlineto{\pgfqpoint{1.515267in}{0.955480in}}%
\pgfpathlineto{\pgfqpoint{1.501616in}{0.944401in}}%
\pgfpathlineto{\pgfqpoint{1.487990in}{0.933343in}}%
\pgfpathlineto{\pgfqpoint{1.483419in}{0.922744in}}%
\pgfpathlineto{\pgfqpoint{1.477073in}{0.912225in}}%
\pgfpathclose%
\pgfusepath{fill}%
\end{pgfscope}%
\begin{pgfscope}%
\pgfpathrectangle{\pgfqpoint{0.407592in}{0.222904in}}{\pgfqpoint{1.411163in}{1.411163in}}%
\pgfusepath{clip}%
\pgfsetbuttcap%
\pgfsetroundjoin%
\definecolor{currentfill}{rgb}{0.749020,0.772549,0.792157}%
\pgfsetfillcolor{currentfill}%
\pgfsetfillopacity{0.160000}%
\pgfsetlinewidth{0.000000pt}%
\definecolor{currentstroke}{rgb}{0.000000,0.000000,0.000000}%
\pgfsetstrokecolor{currentstroke}%
\pgfsetdash{}{0pt}%
\pgfpathmoveto{\pgfqpoint{1.824868in}{1.239951in}}%
\pgfpathlineto{\pgfqpoint{1.838982in}{1.252044in}}%
\pgfpathlineto{\pgfqpoint{1.853116in}{1.264153in}}%
\pgfpathlineto{\pgfqpoint{1.849471in}{1.282224in}}%
\pgfpathlineto{\pgfqpoint{1.842582in}{1.300007in}}%
\pgfpathlineto{\pgfqpoint{1.828812in}{1.287388in}}%
\pgfpathlineto{\pgfqpoint{1.815054in}{1.274780in}}%
\pgfpathlineto{\pgfqpoint{1.821526in}{1.257501in}}%
\pgfpathlineto{\pgfqpoint{1.824868in}{1.239951in}}%
\pgfpathclose%
\pgfusepath{fill}%
\end{pgfscope}%
\begin{pgfscope}%
\pgfpathrectangle{\pgfqpoint{0.407592in}{0.222904in}}{\pgfqpoint{1.411163in}{1.411163in}}%
\pgfusepath{clip}%
\pgfsetbuttcap%
\pgfsetroundjoin%
\definecolor{currentfill}{rgb}{0.749020,0.772549,0.792157}%
\pgfsetfillcolor{currentfill}%
\pgfsetfillopacity{0.160000}%
\pgfsetlinewidth{0.000000pt}%
\definecolor{currentstroke}{rgb}{0.000000,0.000000,0.000000}%
\pgfsetstrokecolor{currentstroke}%
\pgfsetdash{}{0pt}%
\pgfpathmoveto{\pgfqpoint{0.395362in}{1.246258in}}%
\pgfpathlineto{\pgfqpoint{0.381662in}{1.258792in}}%
\pgfpathlineto{\pgfqpoint{0.367948in}{1.271338in}}%
\pgfpathlineto{\pgfqpoint{0.357815in}{1.254000in}}%
\pgfpathlineto{\pgfqpoint{0.350763in}{1.236403in}}%
\pgfpathlineto{\pgfqpoint{0.364968in}{1.224381in}}%
\pgfpathlineto{\pgfqpoint{0.379152in}{1.212375in}}%
\pgfpathlineto{\pgfqpoint{0.385772in}{1.229438in}}%
\pgfpathlineto{\pgfqpoint{0.395362in}{1.246258in}}%
\pgfpathclose%
\pgfusepath{fill}%
\end{pgfscope}%
\begin{pgfscope}%
\pgfpathrectangle{\pgfqpoint{0.407592in}{0.222904in}}{\pgfqpoint{1.411163in}{1.411163in}}%
\pgfusepath{clip}%
\pgfsetbuttcap%
\pgfsetroundjoin%
\definecolor{currentfill}{rgb}{0.749020,0.772549,0.792157}%
\pgfsetfillcolor{currentfill}%
\pgfsetfillopacity{0.160000}%
\pgfsetlinewidth{0.000000pt}%
\definecolor{currentstroke}{rgb}{0.000000,0.000000,0.000000}%
\pgfsetstrokecolor{currentstroke}%
\pgfsetdash{}{0pt}%
\pgfpathmoveto{\pgfqpoint{0.575662in}{1.046054in}}%
\pgfpathlineto{\pgfqpoint{0.561752in}{1.057826in}}%
\pgfpathlineto{\pgfqpoint{0.547824in}{1.069615in}}%
\pgfpathlineto{\pgfqpoint{0.545850in}{1.055857in}}%
\pgfpathlineto{\pgfqpoint{0.546258in}{1.042058in}}%
\pgfpathlineto{\pgfqpoint{0.560348in}{1.030882in}}%
\pgfpathlineto{\pgfqpoint{0.574411in}{1.019725in}}%
\pgfpathlineto{\pgfqpoint{0.573906in}{1.032905in}}%
\pgfpathlineto{\pgfqpoint{0.575662in}{1.046054in}}%
\pgfpathclose%
\pgfusepath{fill}%
\end{pgfscope}%
\begin{pgfscope}%
\pgfpathrectangle{\pgfqpoint{0.407592in}{0.222904in}}{\pgfqpoint{1.411163in}{1.411163in}}%
\pgfusepath{clip}%
\pgfsetbuttcap%
\pgfsetroundjoin%
\definecolor{currentfill}{rgb}{0.749020,0.772549,0.792157}%
\pgfsetfillcolor{currentfill}%
\pgfsetfillopacity{0.160000}%
\pgfsetlinewidth{0.000000pt}%
\definecolor{currentstroke}{rgb}{0.000000,0.000000,0.000000}%
\pgfsetstrokecolor{currentstroke}%
\pgfsetdash{}{0pt}%
\pgfpathmoveto{\pgfqpoint{1.222376in}{0.723952in}}%
\pgfpathlineto{\pgfqpoint{1.230908in}{0.731275in}}%
\pgfpathlineto{\pgfqpoint{1.239473in}{0.738626in}}%
\pgfpathlineto{\pgfqpoint{1.251772in}{0.743227in}}%
\pgfpathlineto{\pgfqpoint{1.263315in}{0.748141in}}%
\pgfpathlineto{\pgfqpoint{1.253408in}{0.740265in}}%
\pgfpathlineto{\pgfqpoint{1.243536in}{0.732416in}}%
\pgfpathlineto{\pgfqpoint{1.233289in}{0.728044in}}%
\pgfpathlineto{\pgfqpoint{1.222376in}{0.723952in}}%
\pgfpathclose%
\pgfusepath{fill}%
\end{pgfscope}%
\begin{pgfscope}%
\pgfpathrectangle{\pgfqpoint{0.407592in}{0.222904in}}{\pgfqpoint{1.411163in}{1.411163in}}%
\pgfusepath{clip}%
\pgfsetbuttcap%
\pgfsetroundjoin%
\definecolor{currentfill}{rgb}{0.749020,0.772549,0.792157}%
\pgfsetfillcolor{currentfill}%
\pgfsetfillopacity{0.160000}%
\pgfsetlinewidth{0.000000pt}%
\definecolor{currentstroke}{rgb}{0.000000,0.000000,0.000000}%
\pgfsetstrokecolor{currentstroke}%
\pgfsetdash{}{0pt}%
\pgfpathmoveto{\pgfqpoint{1.162415in}{0.697853in}}%
\pgfpathlineto{\pgfqpoint{1.167740in}{0.704294in}}%
\pgfpathlineto{\pgfqpoint{1.173088in}{0.710763in}}%
\pgfpathlineto{\pgfqpoint{1.186146in}{0.713549in}}%
\pgfpathlineto{\pgfqpoint{1.198752in}{0.716687in}}%
\pgfpathlineto{\pgfqpoint{1.191759in}{0.709846in}}%
\pgfpathlineto{\pgfqpoint{1.184795in}{0.703034in}}%
\pgfpathlineto{\pgfqpoint{1.173801in}{0.700290in}}%
\pgfpathlineto{\pgfqpoint{1.162415in}{0.697853in}}%
\pgfpathclose%
\pgfusepath{fill}%
\end{pgfscope}%
\begin{pgfscope}%
\pgfpathrectangle{\pgfqpoint{0.407592in}{0.222904in}}{\pgfqpoint{1.411163in}{1.411163in}}%
\pgfusepath{clip}%
\pgfsetbuttcap%
\pgfsetroundjoin%
\definecolor{currentfill}{rgb}{0.749020,0.772549,0.792157}%
\pgfsetfillcolor{currentfill}%
\pgfsetfillopacity{0.160000}%
\pgfsetlinewidth{0.000000pt}%
\definecolor{currentstroke}{rgb}{0.000000,0.000000,0.000000}%
\pgfsetstrokecolor{currentstroke}%
\pgfsetdash{}{0pt}%
\pgfpathmoveto{\pgfqpoint{0.685989in}{0.931212in}}%
\pgfpathlineto{\pgfqpoint{0.672131in}{0.942205in}}%
\pgfpathlineto{\pgfqpoint{0.658248in}{0.953219in}}%
\pgfpathlineto{\pgfqpoint{0.660876in}{0.941980in}}%
\pgfpathlineto{\pgfqpoint{0.665424in}{0.930830in}}%
\pgfpathlineto{\pgfqpoint{0.679168in}{0.920467in}}%
\pgfpathlineto{\pgfqpoint{0.692881in}{0.910128in}}%
\pgfpathlineto{\pgfqpoint{0.688537in}{0.920626in}}%
\pgfpathlineto{\pgfqpoint{0.685989in}{0.931212in}}%
\pgfpathclose%
\pgfusepath{fill}%
\end{pgfscope}%
\begin{pgfscope}%
\pgfpathrectangle{\pgfqpoint{0.407592in}{0.222904in}}{\pgfqpoint{1.411163in}{1.411163in}}%
\pgfusepath{clip}%
\pgfsetbuttcap%
\pgfsetroundjoin%
\definecolor{currentfill}{rgb}{0.749020,0.772549,0.792157}%
\pgfsetfillcolor{currentfill}%
\pgfsetfillopacity{0.160000}%
\pgfsetlinewidth{0.000000pt}%
\definecolor{currentstroke}{rgb}{0.000000,0.000000,0.000000}%
\pgfsetstrokecolor{currentstroke}%
\pgfsetdash{}{0pt}%
\pgfpathmoveto{\pgfqpoint{1.853116in}{1.264153in}}%
\pgfpathlineto{\pgfqpoint{1.867269in}{1.276279in}}%
\pgfpathlineto{\pgfqpoint{1.863468in}{1.294607in}}%
\pgfpathlineto{\pgfqpoint{1.856365in}{1.312638in}}%
\pgfpathlineto{\pgfqpoint{1.842582in}{1.300007in}}%
\pgfpathlineto{\pgfqpoint{1.849471in}{1.282224in}}%
\pgfusepath{fill}%
\end{pgfscope}%
\begin{pgfscope}%
\pgfpathrectangle{\pgfqpoint{0.407592in}{0.222904in}}{\pgfqpoint{1.411163in}{1.411163in}}%
\pgfusepath{clip}%
\pgfsetbuttcap%
\pgfsetroundjoin%
\definecolor{currentfill}{rgb}{0.749020,0.772549,0.792157}%
\pgfsetfillcolor{currentfill}%
\pgfsetfillopacity{0.160000}%
\pgfsetlinewidth{0.000000pt}%
\definecolor{currentstroke}{rgb}{0.000000,0.000000,0.000000}%
\pgfsetstrokecolor{currentstroke}%
\pgfsetdash{}{0pt}%
\pgfpathmoveto{\pgfqpoint{1.263315in}{0.748141in}}%
\pgfpathlineto{\pgfqpoint{1.273257in}{0.756046in}}%
\pgfpathlineto{\pgfqpoint{1.283236in}{0.763979in}}%
\pgfpathlineto{\pgfqpoint{1.295171in}{0.769754in}}%
\pgfpathlineto{\pgfqpoint{1.306149in}{0.775819in}}%
\pgfpathlineto{\pgfqpoint{1.295023in}{0.767302in}}%
\pgfpathlineto{\pgfqpoint{1.283933in}{0.758812in}}%
\pgfpathlineto{\pgfqpoint{1.274050in}{0.753345in}}%
\pgfpathlineto{\pgfqpoint{1.263315in}{0.748141in}}%
\pgfpathclose%
\pgfusepath{fill}%
\end{pgfscope}%
\begin{pgfscope}%
\pgfpathrectangle{\pgfqpoint{0.407592in}{0.222904in}}{\pgfqpoint{1.411163in}{1.411163in}}%
\pgfusepath{clip}%
\pgfsetbuttcap%
\pgfsetroundjoin%
\definecolor{currentfill}{rgb}{0.749020,0.772549,0.792157}%
\pgfsetfillcolor{currentfill}%
\pgfsetfillopacity{0.160000}%
\pgfsetlinewidth{0.000000pt}%
\definecolor{currentstroke}{rgb}{0.000000,0.000000,0.000000}%
\pgfsetstrokecolor{currentstroke}%
\pgfsetdash{}{0pt}%
\pgfpathmoveto{\pgfqpoint{1.625359in}{1.044827in}}%
\pgfpathlineto{\pgfqpoint{1.639232in}{1.056086in}}%
\pgfpathlineto{\pgfqpoint{1.653131in}{1.067365in}}%
\pgfpathlineto{\pgfqpoint{1.656313in}{1.081769in}}%
\pgfpathlineto{\pgfqpoint{1.656983in}{1.096112in}}%
\pgfpathlineto{\pgfqpoint{1.643116in}{1.084231in}}%
\pgfpathlineto{\pgfqpoint{1.629267in}{1.072365in}}%
\pgfpathlineto{\pgfqpoint{1.628507in}{1.058622in}}%
\pgfpathlineto{\pgfqpoint{1.625359in}{1.044827in}}%
\pgfpathclose%
\pgfusepath{fill}%
\end{pgfscope}%
\begin{pgfscope}%
\pgfpathrectangle{\pgfqpoint{0.407592in}{0.222904in}}{\pgfqpoint{1.411163in}{1.411163in}}%
\pgfusepath{clip}%
\pgfsetbuttcap%
\pgfsetroundjoin%
\definecolor{currentfill}{rgb}{0.749020,0.772549,0.792157}%
\pgfsetfillcolor{currentfill}%
\pgfsetfillopacity{0.160000}%
\pgfsetlinewidth{0.000000pt}%
\definecolor{currentstroke}{rgb}{0.000000,0.000000,0.000000}%
\pgfsetstrokecolor{currentstroke}%
\pgfsetdash{}{0pt}%
\pgfpathmoveto{\pgfqpoint{0.367948in}{1.271338in}}%
\pgfpathlineto{\pgfqpoint{0.354222in}{1.283896in}}%
\pgfpathlineto{\pgfqpoint{0.340483in}{1.296466in}}%
\pgfpathlineto{\pgfqpoint{0.329792in}{1.278619in}}%
\pgfpathlineto{\pgfqpoint{0.322294in}{1.260499in}}%
\pgfpathlineto{\pgfqpoint{0.336539in}{1.248443in}}%
\pgfpathlineto{\pgfqpoint{0.350763in}{1.236403in}}%
\pgfpathlineto{\pgfqpoint{0.357815in}{1.254000in}}%
\pgfpathlineto{\pgfqpoint{0.367948in}{1.271338in}}%
\pgfpathclose%
\pgfusepath{fill}%
\end{pgfscope}%
\begin{pgfscope}%
\pgfpathrectangle{\pgfqpoint{0.407592in}{0.222904in}}{\pgfqpoint{1.411163in}{1.411163in}}%
\pgfusepath{clip}%
\pgfsetbuttcap%
\pgfsetroundjoin%
\definecolor{currentfill}{rgb}{0.749020,0.772549,0.792157}%
\pgfsetfillcolor{currentfill}%
\pgfsetfillopacity{0.160000}%
\pgfsetlinewidth{0.000000pt}%
\definecolor{currentstroke}{rgb}{0.000000,0.000000,0.000000}%
\pgfsetstrokecolor{currentstroke}%
\pgfsetdash{}{0pt}%
\pgfpathmoveto{\pgfqpoint{1.339471in}{0.799881in}}%
\pgfpathlineto{\pgfqpoint{1.351197in}{0.808790in}}%
\pgfpathlineto{\pgfqpoint{1.362959in}{0.817725in}}%
\pgfpathlineto{\pgfqpoint{1.373631in}{0.825556in}}%
\pgfpathlineto{\pgfqpoint{1.382990in}{0.833618in}}%
\pgfpathlineto{\pgfqpoint{1.370421in}{0.824040in}}%
\pgfpathlineto{\pgfqpoint{1.357886in}{0.814488in}}%
\pgfpathlineto{\pgfqpoint{1.349275in}{0.807077in}}%
\pgfpathlineto{\pgfqpoint{1.339471in}{0.799881in}}%
\pgfpathclose%
\pgfusepath{fill}%
\end{pgfscope}%
\begin{pgfscope}%
\pgfpathrectangle{\pgfqpoint{0.407592in}{0.222904in}}{\pgfqpoint{1.411163in}{1.411163in}}%
\pgfusepath{clip}%
\pgfsetbuttcap%
\pgfsetroundjoin%
\definecolor{currentfill}{rgb}{0.749020,0.772549,0.792157}%
\pgfsetfillcolor{currentfill}%
\pgfsetfillopacity{0.160000}%
\pgfsetlinewidth{0.000000pt}%
\definecolor{currentstroke}{rgb}{0.000000,0.000000,0.000000}%
\pgfsetstrokecolor{currentstroke}%
\pgfsetdash{}{0pt}%
\pgfpathmoveto{\pgfqpoint{0.993217in}{0.699776in}}%
\pgfpathlineto{\pgfqpoint{0.986066in}{0.706355in}}%
\pgfpathlineto{\pgfqpoint{0.978884in}{0.712962in}}%
\pgfpathlineto{\pgfqpoint{0.991446in}{0.710248in}}%
\pgfpathlineto{\pgfqpoint{1.004468in}{0.707902in}}%
\pgfpathlineto{\pgfqpoint{1.010007in}{0.701613in}}%
\pgfpathlineto{\pgfqpoint{1.015522in}{0.695351in}}%
\pgfpathlineto{\pgfqpoint{1.004171in}{0.697403in}}%
\pgfpathlineto{\pgfqpoint{0.993217in}{0.699776in}}%
\pgfpathclose%
\pgfusepath{fill}%
\end{pgfscope}%
\begin{pgfscope}%
\pgfpathrectangle{\pgfqpoint{0.407592in}{0.222904in}}{\pgfqpoint{1.411163in}{1.411163in}}%
\pgfusepath{clip}%
\pgfsetbuttcap%
\pgfsetroundjoin%
\definecolor{currentfill}{rgb}{0.749020,0.772549,0.792157}%
\pgfsetfillcolor{currentfill}%
\pgfsetfillopacity{0.160000}%
\pgfsetlinewidth{0.000000pt}%
\definecolor{currentstroke}{rgb}{0.000000,0.000000,0.000000}%
\pgfsetstrokecolor{currentstroke}%
\pgfsetdash{}{0pt}%
\pgfpathmoveto{\pgfqpoint{1.306149in}{0.775819in}}%
\pgfpathlineto{\pgfqpoint{1.317311in}{0.784364in}}%
\pgfpathlineto{\pgfqpoint{1.328510in}{0.792936in}}%
\pgfpathlineto{\pgfqpoint{1.339471in}{0.799881in}}%
\pgfpathlineto{\pgfqpoint{1.327780in}{0.791000in}}%
\pgfpathlineto{\pgfqpoint{1.316125in}{0.782146in}}%
\pgfusepath{fill}%
\end{pgfscope}%
\begin{pgfscope}%
\pgfpathrectangle{\pgfqpoint{0.407592in}{0.222904in}}{\pgfqpoint{1.411163in}{1.411163in}}%
\pgfusepath{clip}%
\pgfsetbuttcap%
\pgfsetroundjoin%
\definecolor{currentfill}{rgb}{0.749020,0.772549,0.792157}%
\pgfsetfillcolor{currentfill}%
\pgfsetfillopacity{0.160000}%
\pgfsetlinewidth{0.000000pt}%
\definecolor{currentstroke}{rgb}{0.000000,0.000000,0.000000}%
\pgfsetstrokecolor{currentstroke}%
\pgfsetdash{}{0pt}%
\pgfpathmoveto{\pgfqpoint{1.408228in}{0.852849in}}%
\pgfpathlineto{\pgfqpoint{1.420897in}{0.862503in}}%
\pgfpathlineto{\pgfqpoint{1.433600in}{0.872182in}}%
\pgfpathlineto{\pgfqpoint{1.442849in}{0.881749in}}%
\pgfpathlineto{\pgfqpoint{1.450479in}{0.891485in}}%
\pgfpathlineto{\pgfqpoint{1.437227in}{0.881150in}}%
\pgfpathlineto{\pgfqpoint{1.424004in}{0.870837in}}%
\pgfpathlineto{\pgfqpoint{1.416863in}{0.861762in}}%
\pgfpathlineto{\pgfqpoint{1.408228in}{0.852849in}}%
\pgfpathclose%
\pgfusepath{fill}%
\end{pgfscope}%
\begin{pgfscope}%
\pgfpathrectangle{\pgfqpoint{0.407592in}{0.222904in}}{\pgfqpoint{1.411163in}{1.411163in}}%
\pgfusepath{clip}%
\pgfsetbuttcap%
\pgfsetroundjoin%
\definecolor{currentfill}{rgb}{0.749020,0.772549,0.792157}%
\pgfsetfillcolor{currentfill}%
\pgfsetfillopacity{0.160000}%
\pgfsetlinewidth{0.000000pt}%
\definecolor{currentstroke}{rgb}{0.000000,0.000000,0.000000}%
\pgfsetstrokecolor{currentstroke}%
\pgfsetdash{}{0pt}%
\pgfpathmoveto{\pgfqpoint{0.812157in}{0.812986in}}%
\pgfpathlineto{\pgfqpoint{0.799127in}{0.822472in}}%
\pgfpathlineto{\pgfqpoint{0.786063in}{0.831984in}}%
\pgfpathlineto{\pgfqpoint{0.793926in}{0.823965in}}%
\pgfpathlineto{\pgfqpoint{0.803146in}{0.816186in}}%
\pgfpathlineto{\pgfqpoint{0.815533in}{0.807313in}}%
\pgfpathlineto{\pgfqpoint{0.827881in}{0.798467in}}%
\pgfpathlineto{\pgfqpoint{0.819401in}{0.805615in}}%
\pgfpathlineto{\pgfqpoint{0.812157in}{0.812986in}}%
\pgfpathclose%
\pgfusepath{fill}%
\end{pgfscope}%
\begin{pgfscope}%
\pgfpathrectangle{\pgfqpoint{0.407592in}{0.222904in}}{\pgfqpoint{1.411163in}{1.411163in}}%
\pgfusepath{clip}%
\pgfsetbuttcap%
\pgfsetroundjoin%
\definecolor{currentfill}{rgb}{0.749020,0.772549,0.792157}%
\pgfsetfillcolor{currentfill}%
\pgfsetfillopacity{0.160000}%
\pgfsetlinewidth{0.000000pt}%
\definecolor{currentstroke}{rgb}{0.000000,0.000000,0.000000}%
\pgfsetstrokecolor{currentstroke}%
\pgfsetdash{}{0pt}%
\pgfpathmoveto{\pgfqpoint{0.934335in}{0.727200in}}%
\pgfpathlineto{\pgfqpoint{0.924281in}{0.734725in}}%
\pgfpathlineto{\pgfqpoint{0.914188in}{0.742278in}}%
\pgfpathlineto{\pgfqpoint{0.925657in}{0.737742in}}%
\pgfpathlineto{\pgfqpoint{0.937890in}{0.733547in}}%
\pgfpathlineto{\pgfqpoint{0.946645in}{0.726478in}}%
\pgfpathlineto{\pgfqpoint{0.955364in}{0.719437in}}%
\pgfpathlineto{\pgfqpoint{0.944514in}{0.723167in}}%
\pgfpathlineto{\pgfqpoint{0.934335in}{0.727200in}}%
\pgfpathclose%
\pgfusepath{fill}%
\end{pgfscope}%
\begin{pgfscope}%
\pgfpathrectangle{\pgfqpoint{0.407592in}{0.222904in}}{\pgfqpoint{1.411163in}{1.411163in}}%
\pgfusepath{clip}%
\pgfsetbuttcap%
\pgfsetroundjoin%
\definecolor{currentfill}{rgb}{0.749020,0.772549,0.792157}%
\pgfsetfillcolor{currentfill}%
\pgfsetfillopacity{0.160000}%
\pgfsetlinewidth{0.000000pt}%
\definecolor{currentstroke}{rgb}{0.000000,0.000000,0.000000}%
\pgfsetstrokecolor{currentstroke}%
\pgfsetdash{}{0pt}%
\pgfpathmoveto{\pgfqpoint{1.503785in}{0.933057in}}%
\pgfpathlineto{\pgfqpoint{1.517186in}{0.943508in}}%
\pgfpathlineto{\pgfqpoint{1.530618in}{0.953983in}}%
\pgfpathlineto{\pgfqpoint{1.537644in}{0.965800in}}%
\pgfpathlineto{\pgfqpoint{1.542642in}{0.977696in}}%
\pgfpathlineto{\pgfqpoint{1.528942in}{0.966578in}}%
\pgfpathlineto{\pgfqpoint{1.515267in}{0.955480in}}%
\pgfpathlineto{\pgfqpoint{1.510477in}{0.944229in}}%
\pgfpathlineto{\pgfqpoint{1.503785in}{0.933057in}}%
\pgfpathclose%
\pgfusepath{fill}%
\end{pgfscope}%
\begin{pgfscope}%
\pgfpathrectangle{\pgfqpoint{0.407592in}{0.222904in}}{\pgfqpoint{1.411163in}{1.411163in}}%
\pgfusepath{clip}%
\pgfsetbuttcap%
\pgfsetroundjoin%
\definecolor{currentfill}{rgb}{0.749020,0.772549,0.792157}%
\pgfsetfillcolor{currentfill}%
\pgfsetfillopacity{0.160000}%
\pgfsetlinewidth{0.000000pt}%
\definecolor{currentstroke}{rgb}{0.000000,0.000000,0.000000}%
\pgfsetstrokecolor{currentstroke}%
\pgfsetdash{}{0pt}%
\pgfpathmoveto{\pgfqpoint{0.547824in}{1.069615in}}%
\pgfpathlineto{\pgfqpoint{0.533876in}{1.081420in}}%
\pgfpathlineto{\pgfqpoint{0.519908in}{1.093242in}}%
\pgfpathlineto{\pgfqpoint{0.517702in}{1.078882in}}%
\pgfpathlineto{\pgfqpoint{0.518002in}{1.064474in}}%
\pgfpathlineto{\pgfqpoint{0.532143in}{1.053256in}}%
\pgfpathlineto{\pgfqpoint{0.546258in}{1.042058in}}%
\pgfpathlineto{\pgfqpoint{0.545850in}{1.055857in}}%
\pgfpathlineto{\pgfqpoint{0.547824in}{1.069615in}}%
\pgfpathclose%
\pgfusepath{fill}%
\end{pgfscope}%
\begin{pgfscope}%
\pgfpathrectangle{\pgfqpoint{0.407592in}{0.222904in}}{\pgfqpoint{1.411163in}{1.411163in}}%
\pgfusepath{clip}%
\pgfsetbuttcap%
\pgfsetroundjoin%
\definecolor{currentfill}{rgb}{0.749020,0.772549,0.792157}%
\pgfsetfillcolor{currentfill}%
\pgfsetfillopacity{0.160000}%
\pgfsetlinewidth{0.000000pt}%
\definecolor{currentstroke}{rgb}{0.000000,0.000000,0.000000}%
\pgfsetstrokecolor{currentstroke}%
\pgfsetdash{}{0pt}%
\pgfpathmoveto{\pgfqpoint{1.138681in}{0.693950in}}%
\pgfpathlineto{\pgfqpoint{1.142260in}{0.700111in}}%
\pgfpathlineto{\pgfqpoint{1.145854in}{0.706299in}}%
\pgfpathlineto{\pgfqpoint{1.159636in}{0.708342in}}%
\pgfpathlineto{\pgfqpoint{1.173088in}{0.710763in}}%
\pgfpathlineto{\pgfqpoint{1.167740in}{0.704294in}}%
\pgfpathlineto{\pgfqpoint{1.162415in}{0.697853in}}%
\pgfpathlineto{\pgfqpoint{1.150691in}{0.695737in}}%
\pgfpathlineto{\pgfqpoint{1.138681in}{0.693950in}}%
\pgfpathclose%
\pgfusepath{fill}%
\end{pgfscope}%
\begin{pgfscope}%
\pgfpathrectangle{\pgfqpoint{0.407592in}{0.222904in}}{\pgfqpoint{1.411163in}{1.411163in}}%
\pgfusepath{clip}%
\pgfsetbuttcap%
\pgfsetroundjoin%
\definecolor{currentfill}{rgb}{0.749020,0.772549,0.792157}%
\pgfsetfillcolor{currentfill}%
\pgfsetfillopacity{0.160000}%
\pgfsetlinewidth{0.000000pt}%
\definecolor{currentstroke}{rgb}{0.000000,0.000000,0.000000}%
\pgfsetstrokecolor{currentstroke}%
\pgfsetdash{}{0pt}%
\pgfpathmoveto{\pgfqpoint{0.340483in}{1.296466in}}%
\pgfpathlineto{\pgfqpoint{0.326731in}{1.309048in}}%
\pgfpathlineto{\pgfqpoint{0.315756in}{1.290951in}}%
\pgfpathlineto{\pgfqpoint{0.308030in}{1.272572in}}%
\pgfpathlineto{\pgfqpoint{0.322294in}{1.260499in}}%
\pgfpathlineto{\pgfqpoint{0.329792in}{1.278619in}}%
\pgfusepath{fill}%
\end{pgfscope}%
\begin{pgfscope}%
\pgfpathrectangle{\pgfqpoint{0.407592in}{0.222904in}}{\pgfqpoint{1.411163in}{1.411163in}}%
\pgfusepath{clip}%
\pgfsetbuttcap%
\pgfsetroundjoin%
\definecolor{currentfill}{rgb}{0.749020,0.772549,0.792157}%
\pgfsetfillcolor{currentfill}%
\pgfsetfillopacity{0.160000}%
\pgfsetlinewidth{0.000000pt}%
\definecolor{currentstroke}{rgb}{0.000000,0.000000,0.000000}%
\pgfsetstrokecolor{currentstroke}%
\pgfsetdash{}{0pt}%
\pgfpathmoveto{\pgfqpoint{0.747421in}{0.869006in}}%
\pgfpathlineto{\pgfqpoint{0.733832in}{0.879251in}}%
\pgfpathlineto{\pgfqpoint{0.720213in}{0.889520in}}%
\pgfpathlineto{\pgfqpoint{0.726005in}{0.879814in}}%
\pgfpathlineto{\pgfqpoint{0.733451in}{0.870288in}}%
\pgfpathlineto{\pgfqpoint{0.746657in}{0.860673in}}%
\pgfpathlineto{\pgfqpoint{0.759828in}{0.851084in}}%
\pgfpathlineto{\pgfqpoint{0.752860in}{0.859959in}}%
\pgfpathlineto{\pgfqpoint{0.747421in}{0.869006in}}%
\pgfpathclose%
\pgfusepath{fill}%
\end{pgfscope}%
\begin{pgfscope}%
\pgfpathrectangle{\pgfqpoint{0.407592in}{0.222904in}}{\pgfqpoint{1.411163in}{1.411163in}}%
\pgfusepath{clip}%
\pgfsetbuttcap%
\pgfsetroundjoin%
\definecolor{currentfill}{rgb}{0.749020,0.772549,0.792157}%
\pgfsetfillcolor{currentfill}%
\pgfsetfillopacity{0.160000}%
\pgfsetlinewidth{0.000000pt}%
\definecolor{currentstroke}{rgb}{0.000000,0.000000,0.000000}%
\pgfsetstrokecolor{currentstroke}%
\pgfsetdash{}{0pt}%
\pgfpathmoveto{\pgfqpoint{1.015522in}{0.695351in}}%
\pgfpathlineto{\pgfqpoint{1.010007in}{0.701613in}}%
\pgfpathlineto{\pgfqpoint{1.004468in}{0.707902in}}%
\pgfpathlineto{\pgfqpoint{1.017888in}{0.705935in}}%
\pgfpathlineto{\pgfqpoint{1.031645in}{0.704359in}}%
\pgfpathlineto{\pgfqpoint{1.035433in}{0.698292in}}%
\pgfpathlineto{\pgfqpoint{1.039203in}{0.692254in}}%
\pgfpathlineto{\pgfqpoint{1.027217in}{0.693632in}}%
\pgfpathlineto{\pgfqpoint{1.015522in}{0.695351in}}%
\pgfpathclose%
\pgfusepath{fill}%
\end{pgfscope}%
\begin{pgfscope}%
\pgfpathrectangle{\pgfqpoint{0.407592in}{0.222904in}}{\pgfqpoint{1.411163in}{1.411163in}}%
\pgfusepath{clip}%
\pgfsetbuttcap%
\pgfsetroundjoin%
\definecolor{currentfill}{rgb}{0.749020,0.772549,0.792157}%
\pgfsetfillcolor{currentfill}%
\pgfsetfillopacity{0.160000}%
\pgfsetlinewidth{0.000000pt}%
\definecolor{currentstroke}{rgb}{0.000000,0.000000,0.000000}%
\pgfsetstrokecolor{currentstroke}%
\pgfsetdash{}{0pt}%
\pgfpathmoveto{\pgfqpoint{1.653131in}{1.067365in}}%
\pgfpathlineto{\pgfqpoint{1.667054in}{1.078665in}}%
\pgfpathlineto{\pgfqpoint{1.681003in}{1.089985in}}%
\pgfpathlineto{\pgfqpoint{1.684207in}{1.104989in}}%
\pgfpathlineto{\pgfqpoint{1.684775in}{1.119923in}}%
\pgfpathlineto{\pgfqpoint{1.670870in}{1.108009in}}%
\pgfpathlineto{\pgfqpoint{1.656983in}{1.096112in}}%
\pgfpathlineto{\pgfqpoint{1.656313in}{1.081769in}}%
\pgfpathlineto{\pgfqpoint{1.653131in}{1.067365in}}%
\pgfpathclose%
\pgfusepath{fill}%
\end{pgfscope}%
\begin{pgfscope}%
\pgfpathrectangle{\pgfqpoint{0.407592in}{0.222904in}}{\pgfqpoint{1.411163in}{1.411163in}}%
\pgfusepath{clip}%
\pgfsetbuttcap%
\pgfsetroundjoin%
\definecolor{currentfill}{rgb}{0.749020,0.772549,0.792157}%
\pgfsetfillcolor{currentfill}%
\pgfsetfillopacity{0.160000}%
\pgfsetlinewidth{0.000000pt}%
\definecolor{currentstroke}{rgb}{0.000000,0.000000,0.000000}%
\pgfsetstrokecolor{currentstroke}%
\pgfsetdash{}{0pt}%
\pgfpathmoveto{\pgfqpoint{0.658248in}{0.953219in}}%
\pgfpathlineto{\pgfqpoint{0.644339in}{0.964252in}}%
\pgfpathlineto{\pgfqpoint{0.630405in}{0.975306in}}%
\pgfpathlineto{\pgfqpoint{0.633100in}{0.963422in}}%
\pgfpathlineto{\pgfqpoint{0.637842in}{0.951627in}}%
\pgfpathlineto{\pgfqpoint{0.651648in}{0.941217in}}%
\pgfpathlineto{\pgfqpoint{0.665424in}{0.930830in}}%
\pgfpathlineto{\pgfqpoint{0.660876in}{0.941980in}}%
\pgfpathlineto{\pgfqpoint{0.658248in}{0.953219in}}%
\pgfpathclose%
\pgfusepath{fill}%
\end{pgfscope}%
\begin{pgfscope}%
\pgfpathrectangle{\pgfqpoint{0.407592in}{0.222904in}}{\pgfqpoint{1.411163in}{1.411163in}}%
\pgfusepath{clip}%
\pgfsetbuttcap%
\pgfsetroundjoin%
\definecolor{currentfill}{rgb}{0.749020,0.772549,0.792157}%
\pgfsetfillcolor{currentfill}%
\pgfsetfillopacity{0.160000}%
\pgfsetlinewidth{0.000000pt}%
\definecolor{currentstroke}{rgb}{0.000000,0.000000,0.000000}%
\pgfsetstrokecolor{currentstroke}%
\pgfsetdash{}{0pt}%
\pgfpathmoveto{\pgfqpoint{1.198752in}{0.716687in}}%
\pgfpathlineto{\pgfqpoint{1.205774in}{0.723556in}}%
\pgfpathlineto{\pgfqpoint{1.212825in}{0.730453in}}%
\pgfpathlineto{\pgfqpoint{1.226471in}{0.734361in}}%
\pgfpathlineto{\pgfqpoint{1.239473in}{0.738626in}}%
\pgfpathlineto{\pgfqpoint{1.230908in}{0.731275in}}%
\pgfpathlineto{\pgfqpoint{1.222376in}{0.723952in}}%
\pgfpathlineto{\pgfqpoint{1.210847in}{0.720160in}}%
\pgfpathlineto{\pgfqpoint{1.198752in}{0.716687in}}%
\pgfpathclose%
\pgfusepath{fill}%
\end{pgfscope}%
\begin{pgfscope}%
\pgfpathrectangle{\pgfqpoint{0.407592in}{0.222904in}}{\pgfqpoint{1.411163in}{1.411163in}}%
\pgfusepath{clip}%
\pgfsetbuttcap%
\pgfsetroundjoin%
\definecolor{currentfill}{rgb}{0.749020,0.772549,0.792157}%
\pgfsetfillcolor{currentfill}%
\pgfsetfillopacity{0.160000}%
\pgfsetlinewidth{0.000000pt}%
\definecolor{currentstroke}{rgb}{0.000000,0.000000,0.000000}%
\pgfsetstrokecolor{currentstroke}%
\pgfsetdash{}{0pt}%
\pgfpathmoveto{\pgfqpoint{0.893740in}{0.752279in}}%
\pgfpathlineto{\pgfqpoint{0.882467in}{0.760411in}}%
\pgfpathlineto{\pgfqpoint{0.871156in}{0.768571in}}%
\pgfpathlineto{\pgfqpoint{0.882039in}{0.762858in}}%
\pgfpathlineto{\pgfqpoint{0.893888in}{0.757470in}}%
\pgfpathlineto{\pgfqpoint{0.904057in}{0.749860in}}%
\pgfpathlineto{\pgfqpoint{0.914188in}{0.742278in}}%
\pgfpathlineto{\pgfqpoint{0.903534in}{0.747131in}}%
\pgfpathlineto{\pgfqpoint{0.893740in}{0.752279in}}%
\pgfpathclose%
\pgfusepath{fill}%
\end{pgfscope}%
\begin{pgfscope}%
\pgfpathrectangle{\pgfqpoint{0.407592in}{0.222904in}}{\pgfqpoint{1.411163in}{1.411163in}}%
\pgfusepath{clip}%
\pgfsetbuttcap%
\pgfsetroundjoin%
\definecolor{currentfill}{rgb}{0.749020,0.772549,0.792157}%
\pgfsetfillcolor{currentfill}%
\pgfsetfillopacity{0.160000}%
\pgfsetlinewidth{0.000000pt}%
\definecolor{currentstroke}{rgb}{0.000000,0.000000,0.000000}%
\pgfsetstrokecolor{currentstroke}%
\pgfsetdash{}{0pt}%
\pgfpathmoveto{\pgfqpoint{0.852464in}{0.780857in}}%
\pgfpathlineto{\pgfqpoint{0.840191in}{0.789648in}}%
\pgfpathlineto{\pgfqpoint{0.827881in}{0.798467in}}%
\pgfpathlineto{\pgfqpoint{0.837565in}{0.791576in}}%
\pgfpathlineto{\pgfqpoint{0.848416in}{0.784976in}}%
\pgfpathlineto{\pgfqpoint{0.859805in}{0.776759in}}%
\pgfpathlineto{\pgfqpoint{0.871156in}{0.768571in}}%
\pgfpathlineto{\pgfqpoint{0.861284in}{0.774581in}}%
\pgfpathlineto{\pgfqpoint{0.852464in}{0.780857in}}%
\pgfpathclose%
\pgfusepath{fill}%
\end{pgfscope}%
\begin{pgfscope}%
\pgfpathrectangle{\pgfqpoint{0.407592in}{0.222904in}}{\pgfqpoint{1.411163in}{1.411163in}}%
\pgfusepath{clip}%
\pgfsetbuttcap%
\pgfsetroundjoin%
\definecolor{currentfill}{rgb}{0.749020,0.772549,0.792157}%
\pgfsetfillcolor{currentfill}%
\pgfsetfillopacity{0.160000}%
\pgfsetlinewidth{0.000000pt}%
\definecolor{currentstroke}{rgb}{0.000000,0.000000,0.000000}%
\pgfsetstrokecolor{currentstroke}%
\pgfsetdash{}{0pt}%
\pgfpathmoveto{\pgfqpoint{1.114031in}{0.691402in}}%
\pgfpathlineto{\pgfqpoint{1.115790in}{0.697379in}}%
\pgfpathlineto{\pgfqpoint{1.117557in}{0.703384in}}%
\pgfpathlineto{\pgfqpoint{1.131806in}{0.704643in}}%
\pgfpathlineto{\pgfqpoint{1.145854in}{0.706299in}}%
\pgfpathlineto{\pgfqpoint{1.142260in}{0.700111in}}%
\pgfpathlineto{\pgfqpoint{1.138681in}{0.693950in}}%
\pgfpathlineto{\pgfqpoint{1.126442in}{0.692503in}}%
\pgfpathlineto{\pgfqpoint{1.114031in}{0.691402in}}%
\pgfpathclose%
\pgfusepath{fill}%
\end{pgfscope}%
\begin{pgfscope}%
\pgfpathrectangle{\pgfqpoint{0.407592in}{0.222904in}}{\pgfqpoint{1.411163in}{1.411163in}}%
\pgfusepath{clip}%
\pgfsetbuttcap%
\pgfsetroundjoin%
\definecolor{currentfill}{rgb}{0.749020,0.772549,0.792157}%
\pgfsetfillcolor{currentfill}%
\pgfsetfillopacity{0.160000}%
\pgfsetlinewidth{0.000000pt}%
\definecolor{currentstroke}{rgb}{0.000000,0.000000,0.000000}%
\pgfsetstrokecolor{currentstroke}%
\pgfsetdash{}{0pt}%
\pgfpathmoveto{\pgfqpoint{1.039203in}{0.692254in}}%
\pgfpathlineto{\pgfqpoint{1.035433in}{0.698292in}}%
\pgfpathlineto{\pgfqpoint{1.031645in}{0.704359in}}%
\pgfpathlineto{\pgfqpoint{1.045674in}{0.703179in}}%
\pgfpathlineto{\pgfqpoint{1.059909in}{0.702404in}}%
\pgfpathlineto{\pgfqpoint{1.061870in}{0.696461in}}%
\pgfpathlineto{\pgfqpoint{1.063823in}{0.690546in}}%
\pgfpathlineto{\pgfqpoint{1.051424in}{0.691224in}}%
\pgfpathlineto{\pgfqpoint{1.039203in}{0.692254in}}%
\pgfpathclose%
\pgfusepath{fill}%
\end{pgfscope}%
\begin{pgfscope}%
\pgfpathrectangle{\pgfqpoint{0.407592in}{0.222904in}}{\pgfqpoint{1.411163in}{1.411163in}}%
\pgfusepath{clip}%
\pgfsetbuttcap%
\pgfsetroundjoin%
\definecolor{currentfill}{rgb}{0.749020,0.772549,0.792157}%
\pgfsetfillcolor{currentfill}%
\pgfsetfillopacity{0.160000}%
\pgfsetlinewidth{0.000000pt}%
\definecolor{currentstroke}{rgb}{0.000000,0.000000,0.000000}%
\pgfsetstrokecolor{currentstroke}%
\pgfsetdash{}{0pt}%
\pgfpathmoveto{\pgfqpoint{1.088921in}{0.690260in}}%
\pgfpathlineto{\pgfqpoint{1.088825in}{0.696154in}}%
\pgfpathlineto{\pgfqpoint{1.088728in}{0.702077in}}%
\pgfpathlineto{\pgfqpoint{1.103175in}{0.702527in}}%
\pgfpathlineto{\pgfqpoint{1.117557in}{0.703384in}}%
\pgfpathlineto{\pgfqpoint{1.115790in}{0.697379in}}%
\pgfpathlineto{\pgfqpoint{1.114031in}{0.691402in}}%
\pgfpathlineto{\pgfqpoint{1.101504in}{0.690653in}}%
\pgfpathlineto{\pgfqpoint{1.088921in}{0.690260in}}%
\pgfpathclose%
\pgfusepath{fill}%
\end{pgfscope}%
\begin{pgfscope}%
\pgfpathrectangle{\pgfqpoint{0.407592in}{0.222904in}}{\pgfqpoint{1.411163in}{1.411163in}}%
\pgfusepath{clip}%
\pgfsetbuttcap%
\pgfsetroundjoin%
\definecolor{currentfill}{rgb}{0.749020,0.772549,0.792157}%
\pgfsetfillcolor{currentfill}%
\pgfsetfillopacity{0.160000}%
\pgfsetlinewidth{0.000000pt}%
\definecolor{currentstroke}{rgb}{0.000000,0.000000,0.000000}%
\pgfsetstrokecolor{currentstroke}%
\pgfsetdash{}{0pt}%
\pgfpathmoveto{\pgfqpoint{0.519908in}{1.093242in}}%
\pgfpathlineto{\pgfqpoint{0.505921in}{1.105081in}}%
\pgfpathlineto{\pgfqpoint{0.491914in}{1.116936in}}%
\pgfpathlineto{\pgfqpoint{0.489463in}{1.101983in}}%
\pgfpathlineto{\pgfqpoint{0.489640in}{1.086973in}}%
\pgfpathlineto{\pgfqpoint{0.503834in}{1.075713in}}%
\pgfpathlineto{\pgfqpoint{0.518002in}{1.064474in}}%
\pgfpathlineto{\pgfqpoint{0.517702in}{1.078882in}}%
\pgfpathlineto{\pgfqpoint{0.519908in}{1.093242in}}%
\pgfpathclose%
\pgfusepath{fill}%
\end{pgfscope}%
\begin{pgfscope}%
\pgfpathrectangle{\pgfqpoint{0.407592in}{0.222904in}}{\pgfqpoint{1.411163in}{1.411163in}}%
\pgfusepath{clip}%
\pgfsetbuttcap%
\pgfsetroundjoin%
\definecolor{currentfill}{rgb}{0.749020,0.772549,0.792157}%
\pgfsetfillcolor{currentfill}%
\pgfsetfillopacity{0.160000}%
\pgfsetlinewidth{0.000000pt}%
\definecolor{currentstroke}{rgb}{0.000000,0.000000,0.000000}%
\pgfsetstrokecolor{currentstroke}%
\pgfsetdash{}{0pt}%
\pgfpathmoveto{\pgfqpoint{1.063823in}{0.690546in}}%
\pgfpathlineto{\pgfqpoint{1.061870in}{0.696461in}}%
\pgfpathlineto{\pgfqpoint{1.059909in}{0.702404in}}%
\pgfpathlineto{\pgfqpoint{1.074283in}{0.702035in}}%
\pgfpathlineto{\pgfqpoint{1.088728in}{0.702077in}}%
\pgfpathlineto{\pgfqpoint{1.088825in}{0.696154in}}%
\pgfpathlineto{\pgfqpoint{1.088921in}{0.690260in}}%
\pgfpathlineto{\pgfqpoint{1.076341in}{0.690224in}}%
\pgfpathlineto{\pgfqpoint{1.063823in}{0.690546in}}%
\pgfpathclose%
\pgfusepath{fill}%
\end{pgfscope}%
\begin{pgfscope}%
\pgfpathrectangle{\pgfqpoint{0.407592in}{0.222904in}}{\pgfqpoint{1.411163in}{1.411163in}}%
\pgfusepath{clip}%
\pgfsetbuttcap%
\pgfsetroundjoin%
\definecolor{currentfill}{rgb}{0.749020,0.772549,0.792157}%
\pgfsetfillcolor{currentfill}%
\pgfsetfillopacity{0.160000}%
\pgfsetlinewidth{0.000000pt}%
\definecolor{currentstroke}{rgb}{0.000000,0.000000,0.000000}%
\pgfsetstrokecolor{currentstroke}%
\pgfsetdash{}{0pt}%
\pgfpathmoveto{\pgfqpoint{1.530618in}{0.953983in}}%
\pgfpathlineto{\pgfqpoint{1.544079in}{0.964481in}}%
\pgfpathlineto{\pgfqpoint{1.557570in}{0.975003in}}%
\pgfpathlineto{\pgfqpoint{1.564920in}{0.987459in}}%
\pgfpathlineto{\pgfqpoint{1.570115in}{0.999993in}}%
\pgfpathlineto{\pgfqpoint{1.556366in}{0.988834in}}%
\pgfpathlineto{\pgfqpoint{1.542642in}{0.977696in}}%
\pgfpathlineto{\pgfqpoint{1.537644in}{0.965800in}}%
\pgfpathlineto{\pgfqpoint{1.530618in}{0.953983in}}%
\pgfpathclose%
\pgfusepath{fill}%
\end{pgfscope}%
\begin{pgfscope}%
\pgfpathrectangle{\pgfqpoint{0.407592in}{0.222904in}}{\pgfqpoint{1.411163in}{1.411163in}}%
\pgfusepath{clip}%
\pgfsetbuttcap%
\pgfsetroundjoin%
\definecolor{currentfill}{rgb}{0.749020,0.772549,0.792157}%
\pgfsetfillcolor{currentfill}%
\pgfsetfillopacity{0.160000}%
\pgfsetlinewidth{0.000000pt}%
\definecolor{currentstroke}{rgb}{0.000000,0.000000,0.000000}%
\pgfsetstrokecolor{currentstroke}%
\pgfsetdash{}{0pt}%
\pgfpathmoveto{\pgfqpoint{0.955364in}{0.719437in}}%
\pgfpathlineto{\pgfqpoint{0.946645in}{0.726478in}}%
\pgfpathlineto{\pgfqpoint{0.937890in}{0.733547in}}%
\pgfpathlineto{\pgfqpoint{0.950833in}{0.729714in}}%
\pgfpathlineto{\pgfqpoint{0.964427in}{0.726261in}}%
\pgfpathlineto{\pgfqpoint{0.971671in}{0.719597in}}%
\pgfpathlineto{\pgfqpoint{0.978884in}{0.712962in}}%
\pgfpathlineto{\pgfqpoint{0.966838in}{0.716030in}}%
\pgfpathlineto{\pgfqpoint{0.955364in}{0.719437in}}%
\pgfpathclose%
\pgfusepath{fill}%
\end{pgfscope}%
\begin{pgfscope}%
\pgfpathrectangle{\pgfqpoint{0.407592in}{0.222904in}}{\pgfqpoint{1.411163in}{1.411163in}}%
\pgfusepath{clip}%
\pgfsetbuttcap%
\pgfsetroundjoin%
\definecolor{currentfill}{rgb}{0.749020,0.772549,0.792157}%
\pgfsetfillcolor{currentfill}%
\pgfsetfillopacity{0.160000}%
\pgfsetlinewidth{0.000000pt}%
\definecolor{currentstroke}{rgb}{0.000000,0.000000,0.000000}%
\pgfsetstrokecolor{currentstroke}%
\pgfsetdash{}{0pt}%
\pgfpathmoveto{\pgfqpoint{1.433600in}{0.872182in}}%
\pgfpathlineto{\pgfqpoint{1.446336in}{0.881888in}}%
\pgfpathlineto{\pgfqpoint{1.459107in}{0.891619in}}%
\pgfpathlineto{\pgfqpoint{1.468962in}{0.901833in}}%
\pgfpathlineto{\pgfqpoint{1.477073in}{0.912225in}}%
\pgfpathlineto{\pgfqpoint{1.463761in}{0.901843in}}%
\pgfpathlineto{\pgfqpoint{1.450479in}{0.891485in}}%
\pgfpathlineto{\pgfqpoint{1.442849in}{0.881749in}}%
\pgfpathlineto{\pgfqpoint{1.433600in}{0.872182in}}%
\pgfpathclose%
\pgfusepath{fill}%
\end{pgfscope}%
\begin{pgfscope}%
\pgfpathrectangle{\pgfqpoint{0.407592in}{0.222904in}}{\pgfqpoint{1.411163in}{1.411163in}}%
\pgfusepath{clip}%
\pgfsetbuttcap%
\pgfsetroundjoin%
\definecolor{currentfill}{rgb}{0.749020,0.772549,0.792157}%
\pgfsetfillcolor{currentfill}%
\pgfsetfillopacity{0.160000}%
\pgfsetlinewidth{0.000000pt}%
\definecolor{currentstroke}{rgb}{0.000000,0.000000,0.000000}%
\pgfsetstrokecolor{currentstroke}%
\pgfsetdash{}{0pt}%
\pgfpathmoveto{\pgfqpoint{1.239473in}{0.738626in}}%
\pgfpathlineto{\pgfqpoint{1.248071in}{0.746005in}}%
\pgfpathlineto{\pgfqpoint{1.256703in}{0.753414in}}%
\pgfpathlineto{\pgfqpoint{1.270394in}{0.758523in}}%
\pgfpathlineto{\pgfqpoint{1.283236in}{0.763979in}}%
\pgfpathlineto{\pgfqpoint{1.273257in}{0.756046in}}%
\pgfpathlineto{\pgfqpoint{1.263315in}{0.748141in}}%
\pgfpathlineto{\pgfqpoint{1.251772in}{0.743227in}}%
\pgfpathlineto{\pgfqpoint{1.239473in}{0.738626in}}%
\pgfpathclose%
\pgfusepath{fill}%
\end{pgfscope}%
\begin{pgfscope}%
\pgfpathrectangle{\pgfqpoint{0.407592in}{0.222904in}}{\pgfqpoint{1.411163in}{1.411163in}}%
\pgfusepath{clip}%
\pgfsetbuttcap%
\pgfsetroundjoin%
\definecolor{currentfill}{rgb}{0.749020,0.772549,0.792157}%
\pgfsetfillcolor{currentfill}%
\pgfsetfillopacity{0.160000}%
\pgfsetlinewidth{0.000000pt}%
\definecolor{currentstroke}{rgb}{0.000000,0.000000,0.000000}%
\pgfsetstrokecolor{currentstroke}%
\pgfsetdash{}{0pt}%
\pgfpathmoveto{\pgfqpoint{1.681003in}{1.089985in}}%
\pgfpathlineto{\pgfqpoint{1.694977in}{1.101326in}}%
\pgfpathlineto{\pgfqpoint{1.708976in}{1.112688in}}%
\pgfpathlineto{\pgfqpoint{1.712189in}{1.128283in}}%
\pgfpathlineto{\pgfqpoint{1.712642in}{1.143798in}}%
\pgfpathlineto{\pgfqpoint{1.698699in}{1.131853in}}%
\pgfpathlineto{\pgfqpoint{1.684775in}{1.119923in}}%
\pgfpathlineto{\pgfqpoint{1.684207in}{1.104989in}}%
\pgfpathlineto{\pgfqpoint{1.681003in}{1.089985in}}%
\pgfpathclose%
\pgfusepath{fill}%
\end{pgfscope}%
\begin{pgfscope}%
\pgfpathrectangle{\pgfqpoint{0.407592in}{0.222904in}}{\pgfqpoint{1.411163in}{1.411163in}}%
\pgfusepath{clip}%
\pgfsetbuttcap%
\pgfsetroundjoin%
\definecolor{currentfill}{rgb}{0.749020,0.772549,0.792157}%
\pgfsetfillcolor{currentfill}%
\pgfsetfillopacity{0.160000}%
\pgfsetlinewidth{0.000000pt}%
\definecolor{currentstroke}{rgb}{0.000000,0.000000,0.000000}%
\pgfsetstrokecolor{currentstroke}%
\pgfsetdash{}{0pt}%
\pgfpathmoveto{\pgfqpoint{1.362959in}{0.817725in}}%
\pgfpathlineto{\pgfqpoint{1.374757in}{0.826688in}}%
\pgfpathlineto{\pgfqpoint{1.386591in}{0.835679in}}%
\pgfpathlineto{\pgfqpoint{1.398125in}{0.844141in}}%
\pgfpathlineto{\pgfqpoint{1.408228in}{0.852849in}}%
\pgfpathlineto{\pgfqpoint{1.395592in}{0.843221in}}%
\pgfpathlineto{\pgfqpoint{1.382990in}{0.833618in}}%
\pgfpathlineto{\pgfqpoint{1.373631in}{0.825556in}}%
\pgfpathlineto{\pgfqpoint{1.362959in}{0.817725in}}%
\pgfpathclose%
\pgfusepath{fill}%
\end{pgfscope}%
\begin{pgfscope}%
\pgfpathrectangle{\pgfqpoint{0.407592in}{0.222904in}}{\pgfqpoint{1.411163in}{1.411163in}}%
\pgfusepath{clip}%
\pgfsetbuttcap%
\pgfsetroundjoin%
\definecolor{currentfill}{rgb}{0.749020,0.772549,0.792157}%
\pgfsetfillcolor{currentfill}%
\pgfsetfillopacity{0.160000}%
\pgfsetlinewidth{0.000000pt}%
\definecolor{currentstroke}{rgb}{0.000000,0.000000,0.000000}%
\pgfsetstrokecolor{currentstroke}%
\pgfsetdash{}{0pt}%
\pgfpathmoveto{\pgfqpoint{1.283236in}{0.763979in}}%
\pgfpathlineto{\pgfqpoint{1.293250in}{0.771941in}}%
\pgfpathlineto{\pgfqpoint{1.303300in}{0.779931in}}%
\pgfpathlineto{\pgfqpoint{1.316436in}{0.786275in}}%
\pgfpathlineto{\pgfqpoint{1.328510in}{0.792936in}}%
\pgfpathlineto{\pgfqpoint{1.317311in}{0.784364in}}%
\pgfpathlineto{\pgfqpoint{1.306149in}{0.775819in}}%
\pgfpathlineto{\pgfqpoint{1.295171in}{0.769754in}}%
\pgfpathlineto{\pgfqpoint{1.283236in}{0.763979in}}%
\pgfpathclose%
\pgfusepath{fill}%
\end{pgfscope}%
\begin{pgfscope}%
\pgfpathrectangle{\pgfqpoint{0.407592in}{0.222904in}}{\pgfqpoint{1.411163in}{1.411163in}}%
\pgfusepath{clip}%
\pgfsetbuttcap%
\pgfsetroundjoin%
\definecolor{currentfill}{rgb}{0.749020,0.772549,0.792157}%
\pgfsetfillcolor{currentfill}%
\pgfsetfillopacity{0.160000}%
\pgfsetlinewidth{0.000000pt}%
\definecolor{currentstroke}{rgb}{0.000000,0.000000,0.000000}%
\pgfsetstrokecolor{currentstroke}%
\pgfsetdash{}{0pt}%
\pgfpathmoveto{\pgfqpoint{0.720213in}{0.889520in}}%
\pgfpathlineto{\pgfqpoint{0.706562in}{0.899812in}}%
\pgfpathlineto{\pgfqpoint{0.692881in}{0.910128in}}%
\pgfpathlineto{\pgfqpoint{0.699016in}{0.899768in}}%
\pgfpathlineto{\pgfqpoint{0.706932in}{0.889595in}}%
\pgfpathlineto{\pgfqpoint{0.720209in}{0.879928in}}%
\pgfpathlineto{\pgfqpoint{0.733451in}{0.870288in}}%
\pgfpathlineto{\pgfqpoint{0.726005in}{0.879814in}}%
\pgfpathlineto{\pgfqpoint{0.720213in}{0.889520in}}%
\pgfpathclose%
\pgfusepath{fill}%
\end{pgfscope}%
\begin{pgfscope}%
\pgfpathrectangle{\pgfqpoint{0.407592in}{0.222904in}}{\pgfqpoint{1.411163in}{1.411163in}}%
\pgfusepath{clip}%
\pgfsetbuttcap%
\pgfsetroundjoin%
\definecolor{currentfill}{rgb}{0.749020,0.772549,0.792157}%
\pgfsetfillcolor{currentfill}%
\pgfsetfillopacity{0.160000}%
\pgfsetlinewidth{0.000000pt}%
\definecolor{currentstroke}{rgb}{0.000000,0.000000,0.000000}%
\pgfsetstrokecolor{currentstroke}%
\pgfsetdash{}{0pt}%
\pgfpathmoveto{\pgfqpoint{0.630405in}{0.975306in}}%
\pgfpathlineto{\pgfqpoint{0.616445in}{0.986380in}}%
\pgfpathlineto{\pgfqpoint{0.602459in}{0.997475in}}%
\pgfpathlineto{\pgfqpoint{0.605209in}{0.984952in}}%
\pgfpathlineto{\pgfqpoint{0.610133in}{0.972519in}}%
\pgfpathlineto{\pgfqpoint{0.624003in}{0.962061in}}%
\pgfpathlineto{\pgfqpoint{0.637842in}{0.951627in}}%
\pgfpathlineto{\pgfqpoint{0.633100in}{0.963422in}}%
\pgfpathlineto{\pgfqpoint{0.630405in}{0.975306in}}%
\pgfpathclose%
\pgfusepath{fill}%
\end{pgfscope}%
\begin{pgfscope}%
\pgfpathrectangle{\pgfqpoint{0.407592in}{0.222904in}}{\pgfqpoint{1.411163in}{1.411163in}}%
\pgfusepath{clip}%
\pgfsetbuttcap%
\pgfsetroundjoin%
\definecolor{currentfill}{rgb}{0.749020,0.772549,0.792157}%
\pgfsetfillcolor{currentfill}%
\pgfsetfillopacity{0.160000}%
\pgfsetlinewidth{0.000000pt}%
\definecolor{currentstroke}{rgb}{0.000000,0.000000,0.000000}%
\pgfsetstrokecolor{currentstroke}%
\pgfsetdash{}{0pt}%
\pgfpathmoveto{\pgfqpoint{1.328510in}{0.792936in}}%
\pgfpathlineto{\pgfqpoint{1.339744in}{0.801536in}}%
\pgfpathlineto{\pgfqpoint{1.351015in}{0.810164in}}%
\pgfpathlineto{\pgfqpoint{1.362959in}{0.817725in}}%
\pgfpathlineto{\pgfqpoint{1.351197in}{0.808790in}}%
\pgfpathlineto{\pgfqpoint{1.339471in}{0.799881in}}%
\pgfusepath{fill}%
\end{pgfscope}%
\begin{pgfscope}%
\pgfpathrectangle{\pgfqpoint{0.407592in}{0.222904in}}{\pgfqpoint{1.411163in}{1.411163in}}%
\pgfusepath{clip}%
\pgfsetbuttcap%
\pgfsetroundjoin%
\definecolor{currentfill}{rgb}{0.749020,0.772549,0.792157}%
\pgfsetfillcolor{currentfill}%
\pgfsetfillopacity{0.160000}%
\pgfsetlinewidth{0.000000pt}%
\definecolor{currentstroke}{rgb}{0.000000,0.000000,0.000000}%
\pgfsetstrokecolor{currentstroke}%
\pgfsetdash{}{0pt}%
\pgfpathmoveto{\pgfqpoint{0.786063in}{0.831984in}}%
\pgfpathlineto{\pgfqpoint{0.772963in}{0.841521in}}%
\pgfpathlineto{\pgfqpoint{0.759828in}{0.851084in}}%
\pgfpathlineto{\pgfqpoint{0.768303in}{0.842422in}}%
\pgfpathlineto{\pgfqpoint{0.778258in}{0.834014in}}%
\pgfpathlineto{\pgfqpoint{0.790721in}{0.825086in}}%
\pgfpathlineto{\pgfqpoint{0.803146in}{0.816186in}}%
\pgfpathlineto{\pgfqpoint{0.793926in}{0.823965in}}%
\pgfpathlineto{\pgfqpoint{0.786063in}{0.831984in}}%
\pgfpathclose%
\pgfusepath{fill}%
\end{pgfscope}%
\begin{pgfscope}%
\pgfpathrectangle{\pgfqpoint{0.407592in}{0.222904in}}{\pgfqpoint{1.411163in}{1.411163in}}%
\pgfusepath{clip}%
\pgfsetbuttcap%
\pgfsetroundjoin%
\definecolor{currentfill}{rgb}{0.749020,0.772549,0.792157}%
\pgfsetfillcolor{currentfill}%
\pgfsetfillopacity{0.160000}%
\pgfsetlinewidth{0.000000pt}%
\definecolor{currentstroke}{rgb}{0.000000,0.000000,0.000000}%
\pgfsetstrokecolor{currentstroke}%
\pgfsetdash{}{0pt}%
\pgfpathmoveto{\pgfqpoint{1.173088in}{0.710763in}}%
\pgfpathlineto{\pgfqpoint{1.178459in}{0.717260in}}%
\pgfpathlineto{\pgfqpoint{1.183854in}{0.723786in}}%
\pgfpathlineto{\pgfqpoint{1.198598in}{0.726922in}}%
\pgfpathlineto{\pgfqpoint{1.212825in}{0.730453in}}%
\pgfpathlineto{\pgfqpoint{1.205774in}{0.723556in}}%
\pgfpathlineto{\pgfqpoint{1.198752in}{0.716687in}}%
\pgfpathlineto{\pgfqpoint{1.186146in}{0.713549in}}%
\pgfpathlineto{\pgfqpoint{1.173088in}{0.710763in}}%
\pgfpathclose%
\pgfusepath{fill}%
\end{pgfscope}%
\begin{pgfscope}%
\pgfpathrectangle{\pgfqpoint{0.407592in}{0.222904in}}{\pgfqpoint{1.411163in}{1.411163in}}%
\pgfusepath{clip}%
\pgfsetbuttcap%
\pgfsetroundjoin%
\definecolor{currentfill}{rgb}{0.749020,0.772549,0.792157}%
\pgfsetfillcolor{currentfill}%
\pgfsetfillopacity{0.160000}%
\pgfsetlinewidth{0.000000pt}%
\definecolor{currentstroke}{rgb}{0.000000,0.000000,0.000000}%
\pgfsetstrokecolor{currentstroke}%
\pgfsetdash{}{0pt}%
\pgfpathmoveto{\pgfqpoint{0.491914in}{1.116936in}}%
\pgfpathlineto{\pgfqpoint{0.477888in}{1.128807in}}%
\pgfpathlineto{\pgfqpoint{0.463842in}{1.140695in}}%
\pgfpathlineto{\pgfqpoint{0.461132in}{1.125158in}}%
\pgfpathlineto{\pgfqpoint{0.461173in}{1.109555in}}%
\pgfpathlineto{\pgfqpoint{0.475420in}{1.098254in}}%
\pgfpathlineto{\pgfqpoint{0.489640in}{1.086973in}}%
\pgfpathlineto{\pgfqpoint{0.489463in}{1.101983in}}%
\pgfpathlineto{\pgfqpoint{0.491914in}{1.116936in}}%
\pgfpathclose%
\pgfusepath{fill}%
\end{pgfscope}%
\begin{pgfscope}%
\pgfpathrectangle{\pgfqpoint{0.407592in}{0.222904in}}{\pgfqpoint{1.411163in}{1.411163in}}%
\pgfusepath{clip}%
\pgfsetbuttcap%
\pgfsetroundjoin%
\definecolor{currentfill}{rgb}{0.749020,0.772549,0.792157}%
\pgfsetfillcolor{currentfill}%
\pgfsetfillopacity{0.160000}%
\pgfsetlinewidth{0.000000pt}%
\definecolor{currentstroke}{rgb}{0.000000,0.000000,0.000000}%
\pgfsetstrokecolor{currentstroke}%
\pgfsetdash{}{0pt}%
\pgfpathmoveto{\pgfqpoint{1.708976in}{1.112688in}}%
\pgfpathlineto{\pgfqpoint{1.723001in}{1.124070in}}%
\pgfpathlineto{\pgfqpoint{1.737051in}{1.135472in}}%
\pgfpathlineto{\pgfqpoint{1.740260in}{1.151651in}}%
\pgfpathlineto{\pgfqpoint{1.740584in}{1.167739in}}%
\pgfpathlineto{\pgfqpoint{1.726603in}{1.155760in}}%
\pgfpathlineto{\pgfqpoint{1.712642in}{1.143798in}}%
\pgfpathlineto{\pgfqpoint{1.712189in}{1.128283in}}%
\pgfpathlineto{\pgfqpoint{1.708976in}{1.112688in}}%
\pgfpathclose%
\pgfusepath{fill}%
\end{pgfscope}%
\begin{pgfscope}%
\pgfpathrectangle{\pgfqpoint{0.407592in}{0.222904in}}{\pgfqpoint{1.411163in}{1.411163in}}%
\pgfusepath{clip}%
\pgfsetbuttcap%
\pgfsetroundjoin%
\definecolor{currentfill}{rgb}{0.749020,0.772549,0.792157}%
\pgfsetfillcolor{currentfill}%
\pgfsetfillopacity{0.160000}%
\pgfsetlinewidth{0.000000pt}%
\definecolor{currentstroke}{rgb}{0.000000,0.000000,0.000000}%
\pgfsetstrokecolor{currentstroke}%
\pgfsetdash{}{0pt}%
\pgfpathmoveto{\pgfqpoint{0.914188in}{0.742278in}}%
\pgfpathlineto{\pgfqpoint{0.904057in}{0.749860in}}%
\pgfpathlineto{\pgfqpoint{0.893888in}{0.757470in}}%
\pgfpathlineto{\pgfqpoint{0.906652in}{0.752433in}}%
\pgfpathlineto{\pgfqpoint{0.920274in}{0.747773in}}%
\pgfpathlineto{\pgfqpoint{0.929100in}{0.740646in}}%
\pgfpathlineto{\pgfqpoint{0.937890in}{0.733547in}}%
\pgfpathlineto{\pgfqpoint{0.925657in}{0.737742in}}%
\pgfpathlineto{\pgfqpoint{0.914188in}{0.742278in}}%
\pgfpathclose%
\pgfusepath{fill}%
\end{pgfscope}%
\begin{pgfscope}%
\pgfpathrectangle{\pgfqpoint{0.407592in}{0.222904in}}{\pgfqpoint{1.411163in}{1.411163in}}%
\pgfusepath{clip}%
\pgfsetbuttcap%
\pgfsetroundjoin%
\definecolor{currentfill}{rgb}{0.749020,0.772549,0.792157}%
\pgfsetfillcolor{currentfill}%
\pgfsetfillopacity{0.160000}%
\pgfsetlinewidth{0.000000pt}%
\definecolor{currentstroke}{rgb}{0.000000,0.000000,0.000000}%
\pgfsetstrokecolor{currentstroke}%
\pgfsetdash{}{0pt}%
\pgfpathmoveto{\pgfqpoint{0.978884in}{0.712962in}}%
\pgfpathlineto{\pgfqpoint{0.971671in}{0.719597in}}%
\pgfpathlineto{\pgfqpoint{0.964427in}{0.726261in}}%
\pgfpathlineto{\pgfqpoint{0.978609in}{0.723206in}}%
\pgfpathlineto{\pgfqpoint{0.993314in}{0.720565in}}%
\pgfpathlineto{\pgfqpoint{0.998904in}{0.714219in}}%
\pgfpathlineto{\pgfqpoint{1.004468in}{0.707902in}}%
\pgfpathlineto{\pgfqpoint{0.991446in}{0.710248in}}%
\pgfpathlineto{\pgfqpoint{0.978884in}{0.712962in}}%
\pgfpathclose%
\pgfusepath{fill}%
\end{pgfscope}%
\begin{pgfscope}%
\pgfpathrectangle{\pgfqpoint{0.407592in}{0.222904in}}{\pgfqpoint{1.411163in}{1.411163in}}%
\pgfusepath{clip}%
\pgfsetbuttcap%
\pgfsetroundjoin%
\definecolor{currentfill}{rgb}{0.749020,0.772549,0.792157}%
\pgfsetfillcolor{currentfill}%
\pgfsetfillopacity{0.160000}%
\pgfsetlinewidth{0.000000pt}%
\definecolor{currentstroke}{rgb}{0.000000,0.000000,0.000000}%
\pgfsetstrokecolor{currentstroke}%
\pgfsetdash{}{0pt}%
\pgfpathmoveto{\pgfqpoint{1.557570in}{0.975003in}}%
\pgfpathlineto{\pgfqpoint{1.571092in}{0.985549in}}%
\pgfpathlineto{\pgfqpoint{1.584645in}{0.996118in}}%
\pgfpathlineto{\pgfqpoint{1.592308in}{1.009206in}}%
\pgfpathlineto{\pgfqpoint{1.597687in}{1.022369in}}%
\pgfpathlineto{\pgfqpoint{1.583889in}{1.011171in}}%
\pgfpathlineto{\pgfqpoint{1.570115in}{0.999993in}}%
\pgfpathlineto{\pgfqpoint{1.564920in}{0.987459in}}%
\pgfpathlineto{\pgfqpoint{1.557570in}{0.975003in}}%
\pgfpathclose%
\pgfusepath{fill}%
\end{pgfscope}%
\begin{pgfscope}%
\pgfpathrectangle{\pgfqpoint{0.407592in}{0.222904in}}{\pgfqpoint{1.411163in}{1.411163in}}%
\pgfusepath{clip}%
\pgfsetbuttcap%
\pgfsetroundjoin%
\definecolor{currentfill}{rgb}{0.749020,0.772549,0.792157}%
\pgfsetfillcolor{currentfill}%
\pgfsetfillopacity{0.160000}%
\pgfsetlinewidth{0.000000pt}%
\definecolor{currentstroke}{rgb}{0.000000,0.000000,0.000000}%
\pgfsetstrokecolor{currentstroke}%
\pgfsetdash{}{0pt}%
\pgfpathmoveto{\pgfqpoint{0.827881in}{0.798467in}}%
\pgfpathlineto{\pgfqpoint{0.815533in}{0.807313in}}%
\pgfpathlineto{\pgfqpoint{0.803146in}{0.816186in}}%
\pgfpathlineto{\pgfqpoint{0.813690in}{0.808683in}}%
\pgfpathlineto{\pgfqpoint{0.825518in}{0.801494in}}%
\pgfpathlineto{\pgfqpoint{0.836986in}{0.793221in}}%
\pgfpathlineto{\pgfqpoint{0.848416in}{0.784976in}}%
\pgfpathlineto{\pgfqpoint{0.837565in}{0.791576in}}%
\pgfpathlineto{\pgfqpoint{0.827881in}{0.798467in}}%
\pgfpathclose%
\pgfusepath{fill}%
\end{pgfscope}%
\begin{pgfscope}%
\pgfpathrectangle{\pgfqpoint{0.407592in}{0.222904in}}{\pgfqpoint{1.411163in}{1.411163in}}%
\pgfusepath{clip}%
\pgfsetbuttcap%
\pgfsetroundjoin%
\definecolor{currentfill}{rgb}{0.749020,0.772549,0.792157}%
\pgfsetfillcolor{currentfill}%
\pgfsetfillopacity{0.160000}%
\pgfsetlinewidth{0.000000pt}%
\definecolor{currentstroke}{rgb}{0.000000,0.000000,0.000000}%
\pgfsetstrokecolor{currentstroke}%
\pgfsetdash{}{0pt}%
\pgfpathmoveto{\pgfqpoint{1.145854in}{0.706299in}}%
\pgfpathlineto{\pgfqpoint{1.149465in}{0.712515in}}%
\pgfpathlineto{\pgfqpoint{1.153092in}{0.718760in}}%
\pgfpathlineto{\pgfqpoint{1.168661in}{0.721061in}}%
\pgfpathlineto{\pgfqpoint{1.183854in}{0.723786in}}%
\pgfpathlineto{\pgfqpoint{1.178459in}{0.717260in}}%
\pgfpathlineto{\pgfqpoint{1.173088in}{0.710763in}}%
\pgfpathlineto{\pgfqpoint{1.159636in}{0.708342in}}%
\pgfpathlineto{\pgfqpoint{1.145854in}{0.706299in}}%
\pgfpathclose%
\pgfusepath{fill}%
\end{pgfscope}%
\begin{pgfscope}%
\pgfpathrectangle{\pgfqpoint{0.407592in}{0.222904in}}{\pgfqpoint{1.411163in}{1.411163in}}%
\pgfusepath{clip}%
\pgfsetbuttcap%
\pgfsetroundjoin%
\definecolor{currentfill}{rgb}{0.749020,0.772549,0.792157}%
\pgfsetfillcolor{currentfill}%
\pgfsetfillopacity{0.160000}%
\pgfsetlinewidth{0.000000pt}%
\definecolor{currentstroke}{rgb}{0.000000,0.000000,0.000000}%
\pgfsetstrokecolor{currentstroke}%
\pgfsetdash{}{0pt}%
\pgfpathmoveto{\pgfqpoint{1.459107in}{0.891619in}}%
\pgfpathlineto{\pgfqpoint{1.471911in}{0.901376in}}%
\pgfpathlineto{\pgfqpoint{1.484750in}{0.911159in}}%
\pgfpathlineto{\pgfqpoint{1.495203in}{0.922017in}}%
\pgfpathlineto{\pgfqpoint{1.503785in}{0.933057in}}%
\pgfpathlineto{\pgfqpoint{1.490414in}{0.922629in}}%
\pgfpathlineto{\pgfqpoint{1.477073in}{0.912225in}}%
\pgfpathlineto{\pgfqpoint{1.468962in}{0.901833in}}%
\pgfpathlineto{\pgfqpoint{1.459107in}{0.891619in}}%
\pgfpathclose%
\pgfusepath{fill}%
\end{pgfscope}%
\begin{pgfscope}%
\pgfpathrectangle{\pgfqpoint{0.407592in}{0.222904in}}{\pgfqpoint{1.411163in}{1.411163in}}%
\pgfusepath{clip}%
\pgfsetbuttcap%
\pgfsetroundjoin%
\definecolor{currentfill}{rgb}{0.749020,0.772549,0.792157}%
\pgfsetfillcolor{currentfill}%
\pgfsetfillopacity{0.160000}%
\pgfsetlinewidth{0.000000pt}%
\definecolor{currentstroke}{rgb}{0.000000,0.000000,0.000000}%
\pgfsetstrokecolor{currentstroke}%
\pgfsetdash{}{0pt}%
\pgfpathmoveto{\pgfqpoint{0.871156in}{0.768571in}}%
\pgfpathlineto{\pgfqpoint{0.859805in}{0.776759in}}%
\pgfpathlineto{\pgfqpoint{0.848416in}{0.784976in}}%
\pgfpathlineto{\pgfqpoint{0.860389in}{0.778699in}}%
\pgfpathlineto{\pgfqpoint{0.873435in}{0.772777in}}%
\pgfpathlineto{\pgfqpoint{0.883681in}{0.765109in}}%
\pgfpathlineto{\pgfqpoint{0.893888in}{0.757470in}}%
\pgfpathlineto{\pgfqpoint{0.882039in}{0.762858in}}%
\pgfpathlineto{\pgfqpoint{0.871156in}{0.768571in}}%
\pgfpathclose%
\pgfusepath{fill}%
\end{pgfscope}%
\begin{pgfscope}%
\pgfpathrectangle{\pgfqpoint{0.407592in}{0.222904in}}{\pgfqpoint{1.411163in}{1.411163in}}%
\pgfusepath{clip}%
\pgfsetbuttcap%
\pgfsetroundjoin%
\definecolor{currentfill}{rgb}{0.749020,0.772549,0.792157}%
\pgfsetfillcolor{currentfill}%
\pgfsetfillopacity{0.160000}%
\pgfsetlinewidth{0.000000pt}%
\definecolor{currentstroke}{rgb}{0.000000,0.000000,0.000000}%
\pgfsetstrokecolor{currentstroke}%
\pgfsetdash{}{0pt}%
\pgfpathmoveto{\pgfqpoint{0.602459in}{0.997475in}}%
\pgfpathlineto{\pgfqpoint{0.588448in}{1.008590in}}%
\pgfpathlineto{\pgfqpoint{0.574411in}{1.019725in}}%
\pgfpathlineto{\pgfqpoint{0.577202in}{1.006572in}}%
\pgfpathlineto{\pgfqpoint{0.582297in}{0.993507in}}%
\pgfpathlineto{\pgfqpoint{0.596231in}{0.983001in}}%
\pgfpathlineto{\pgfqpoint{0.610133in}{0.972519in}}%
\pgfpathlineto{\pgfqpoint{0.605209in}{0.984952in}}%
\pgfpathlineto{\pgfqpoint{0.602459in}{0.997475in}}%
\pgfpathclose%
\pgfusepath{fill}%
\end{pgfscope}%
\begin{pgfscope}%
\pgfpathrectangle{\pgfqpoint{0.407592in}{0.222904in}}{\pgfqpoint{1.411163in}{1.411163in}}%
\pgfusepath{clip}%
\pgfsetbuttcap%
\pgfsetroundjoin%
\definecolor{currentfill}{rgb}{0.749020,0.772549,0.792157}%
\pgfsetfillcolor{currentfill}%
\pgfsetfillopacity{0.160000}%
\pgfsetlinewidth{0.000000pt}%
\definecolor{currentstroke}{rgb}{0.000000,0.000000,0.000000}%
\pgfsetstrokecolor{currentstroke}%
\pgfsetdash{}{0pt}%
\pgfpathmoveto{\pgfqpoint{1.212825in}{0.730453in}}%
\pgfpathlineto{\pgfqpoint{1.219906in}{0.737379in}}%
\pgfpathlineto{\pgfqpoint{1.227016in}{0.744334in}}%
\pgfpathlineto{\pgfqpoint{1.242222in}{0.748676in}}%
\pgfpathlineto{\pgfqpoint{1.256703in}{0.753414in}}%
\pgfpathlineto{\pgfqpoint{1.248071in}{0.746005in}}%
\pgfpathlineto{\pgfqpoint{1.239473in}{0.738626in}}%
\pgfpathlineto{\pgfqpoint{1.226471in}{0.734361in}}%
\pgfpathlineto{\pgfqpoint{1.212825in}{0.730453in}}%
\pgfpathclose%
\pgfusepath{fill}%
\end{pgfscope}%
\begin{pgfscope}%
\pgfpathrectangle{\pgfqpoint{0.407592in}{0.222904in}}{\pgfqpoint{1.411163in}{1.411163in}}%
\pgfusepath{clip}%
\pgfsetbuttcap%
\pgfsetroundjoin%
\definecolor{currentfill}{rgb}{0.749020,0.772549,0.792157}%
\pgfsetfillcolor{currentfill}%
\pgfsetfillopacity{0.160000}%
\pgfsetlinewidth{0.000000pt}%
\definecolor{currentstroke}{rgb}{0.000000,0.000000,0.000000}%
\pgfsetstrokecolor{currentstroke}%
\pgfsetdash{}{0pt}%
\pgfpathmoveto{\pgfqpoint{0.463842in}{1.140695in}}%
\pgfpathlineto{\pgfqpoint{0.449777in}{1.152600in}}%
\pgfpathlineto{\pgfqpoint{0.435692in}{1.164522in}}%
\pgfpathlineto{\pgfqpoint{0.432709in}{1.148410in}}%
\pgfpathlineto{\pgfqpoint{0.432600in}{1.132222in}}%
\pgfpathlineto{\pgfqpoint{0.446900in}{1.120878in}}%
\pgfpathlineto{\pgfqpoint{0.461173in}{1.109555in}}%
\pgfpathlineto{\pgfqpoint{0.461132in}{1.125158in}}%
\pgfpathlineto{\pgfqpoint{0.463842in}{1.140695in}}%
\pgfpathclose%
\pgfusepath{fill}%
\end{pgfscope}%
\begin{pgfscope}%
\pgfpathrectangle{\pgfqpoint{0.407592in}{0.222904in}}{\pgfqpoint{1.411163in}{1.411163in}}%
\pgfusepath{clip}%
\pgfsetbuttcap%
\pgfsetroundjoin%
\definecolor{currentfill}{rgb}{0.749020,0.772549,0.792157}%
\pgfsetfillcolor{currentfill}%
\pgfsetfillopacity{0.160000}%
\pgfsetlinewidth{0.000000pt}%
\definecolor{currentstroke}{rgb}{0.000000,0.000000,0.000000}%
\pgfsetstrokecolor{currentstroke}%
\pgfsetdash{}{0pt}%
\pgfpathmoveto{\pgfqpoint{1.386591in}{0.835679in}}%
\pgfpathlineto{\pgfqpoint{1.398461in}{0.844696in}}%
\pgfpathlineto{\pgfqpoint{1.410368in}{0.853742in}}%
\pgfpathlineto{\pgfqpoint{1.422760in}{0.862832in}}%
\pgfpathlineto{\pgfqpoint{1.433600in}{0.872182in}}%
\pgfpathlineto{\pgfqpoint{1.420897in}{0.862503in}}%
\pgfpathlineto{\pgfqpoint{1.408228in}{0.852849in}}%
\pgfpathlineto{\pgfqpoint{1.398125in}{0.844141in}}%
\pgfpathlineto{\pgfqpoint{1.386591in}{0.835679in}}%
\pgfpathclose%
\pgfusepath{fill}%
\end{pgfscope}%
\begin{pgfscope}%
\pgfpathrectangle{\pgfqpoint{0.407592in}{0.222904in}}{\pgfqpoint{1.411163in}{1.411163in}}%
\pgfusepath{clip}%
\pgfsetbuttcap%
\pgfsetroundjoin%
\definecolor{currentfill}{rgb}{0.749020,0.772549,0.792157}%
\pgfsetfillcolor{currentfill}%
\pgfsetfillopacity{0.160000}%
\pgfsetlinewidth{0.000000pt}%
\definecolor{currentstroke}{rgb}{0.000000,0.000000,0.000000}%
\pgfsetstrokecolor{currentstroke}%
\pgfsetdash{}{0pt}%
\pgfpathmoveto{\pgfqpoint{1.004468in}{0.707902in}}%
\pgfpathlineto{\pgfqpoint{0.998904in}{0.714219in}}%
\pgfpathlineto{\pgfqpoint{0.993314in}{0.720565in}}%
\pgfpathlineto{\pgfqpoint{1.008474in}{0.718351in}}%
\pgfpathlineto{\pgfqpoint{1.024018in}{0.716575in}}%
\pgfpathlineto{\pgfqpoint{1.027840in}{0.710453in}}%
\pgfpathlineto{\pgfqpoint{1.031645in}{0.704359in}}%
\pgfpathlineto{\pgfqpoint{1.017888in}{0.705935in}}%
\pgfpathlineto{\pgfqpoint{1.004468in}{0.707902in}}%
\pgfpathclose%
\pgfusepath{fill}%
\end{pgfscope}%
\begin{pgfscope}%
\pgfpathrectangle{\pgfqpoint{0.407592in}{0.222904in}}{\pgfqpoint{1.411163in}{1.411163in}}%
\pgfusepath{clip}%
\pgfsetbuttcap%
\pgfsetroundjoin%
\definecolor{currentfill}{rgb}{0.749020,0.772549,0.792157}%
\pgfsetfillcolor{currentfill}%
\pgfsetfillopacity{0.160000}%
\pgfsetlinewidth{0.000000pt}%
\definecolor{currentstroke}{rgb}{0.000000,0.000000,0.000000}%
\pgfsetstrokecolor{currentstroke}%
\pgfsetdash{}{0pt}%
\pgfpathmoveto{\pgfqpoint{0.692881in}{0.910128in}}%
\pgfpathlineto{\pgfqpoint{0.679168in}{0.920467in}}%
\pgfpathlineto{\pgfqpoint{0.665424in}{0.930830in}}%
\pgfpathlineto{\pgfqpoint{0.671892in}{0.919822in}}%
\pgfpathlineto{\pgfqpoint{0.680268in}{0.909008in}}%
\pgfpathlineto{\pgfqpoint{0.693618in}{0.899288in}}%
\pgfpathlineto{\pgfqpoint{0.706932in}{0.889595in}}%
\pgfpathlineto{\pgfqpoint{0.699016in}{0.899768in}}%
\pgfpathlineto{\pgfqpoint{0.692881in}{0.910128in}}%
\pgfpathclose%
\pgfusepath{fill}%
\end{pgfscope}%
\begin{pgfscope}%
\pgfpathrectangle{\pgfqpoint{0.407592in}{0.222904in}}{\pgfqpoint{1.411163in}{1.411163in}}%
\pgfusepath{clip}%
\pgfsetbuttcap%
\pgfsetroundjoin%
\definecolor{currentfill}{rgb}{0.749020,0.772549,0.792157}%
\pgfsetfillcolor{currentfill}%
\pgfsetfillopacity{0.160000}%
\pgfsetlinewidth{0.000000pt}%
\definecolor{currentstroke}{rgb}{0.000000,0.000000,0.000000}%
\pgfsetstrokecolor{currentstroke}%
\pgfsetdash{}{0pt}%
\pgfpathmoveto{\pgfqpoint{1.737051in}{1.135472in}}%
\pgfpathlineto{\pgfqpoint{1.751127in}{1.146896in}}%
\pgfpathlineto{\pgfqpoint{1.765228in}{1.158340in}}%
\pgfpathlineto{\pgfqpoint{1.768420in}{1.175092in}}%
\pgfpathlineto{\pgfqpoint{1.768602in}{1.191744in}}%
\pgfpathlineto{\pgfqpoint{1.754584in}{1.179733in}}%
\pgfpathlineto{\pgfqpoint{1.740584in}{1.167739in}}%
\pgfpathlineto{\pgfqpoint{1.740260in}{1.151651in}}%
\pgfpathlineto{\pgfqpoint{1.737051in}{1.135472in}}%
\pgfpathclose%
\pgfusepath{fill}%
\end{pgfscope}%
\begin{pgfscope}%
\pgfpathrectangle{\pgfqpoint{0.407592in}{0.222904in}}{\pgfqpoint{1.411163in}{1.411163in}}%
\pgfusepath{clip}%
\pgfsetbuttcap%
\pgfsetroundjoin%
\definecolor{currentfill}{rgb}{0.749020,0.772549,0.792157}%
\pgfsetfillcolor{currentfill}%
\pgfsetfillopacity{0.160000}%
\pgfsetlinewidth{0.000000pt}%
\definecolor{currentstroke}{rgb}{0.000000,0.000000,0.000000}%
\pgfsetstrokecolor{currentstroke}%
\pgfsetdash{}{0pt}%
\pgfpathmoveto{\pgfqpoint{1.117557in}{0.703384in}}%
\pgfpathlineto{\pgfqpoint{1.119332in}{0.709416in}}%
\pgfpathlineto{\pgfqpoint{1.121116in}{0.715477in}}%
\pgfpathlineto{\pgfqpoint{1.137218in}{0.716895in}}%
\pgfpathlineto{\pgfqpoint{1.153092in}{0.718760in}}%
\pgfpathlineto{\pgfqpoint{1.149465in}{0.712515in}}%
\pgfpathlineto{\pgfqpoint{1.145854in}{0.706299in}}%
\pgfpathlineto{\pgfqpoint{1.131806in}{0.704643in}}%
\pgfpathlineto{\pgfqpoint{1.117557in}{0.703384in}}%
\pgfpathclose%
\pgfusepath{fill}%
\end{pgfscope}%
\begin{pgfscope}%
\pgfpathrectangle{\pgfqpoint{0.407592in}{0.222904in}}{\pgfqpoint{1.411163in}{1.411163in}}%
\pgfusepath{clip}%
\pgfsetbuttcap%
\pgfsetroundjoin%
\definecolor{currentfill}{rgb}{0.749020,0.772549,0.792157}%
\pgfsetfillcolor{currentfill}%
\pgfsetfillopacity{0.160000}%
\pgfsetlinewidth{0.000000pt}%
\definecolor{currentstroke}{rgb}{0.000000,0.000000,0.000000}%
\pgfsetstrokecolor{currentstroke}%
\pgfsetdash{}{0pt}%
\pgfpathmoveto{\pgfqpoint{0.759828in}{0.851084in}}%
\pgfpathlineto{\pgfqpoint{0.746657in}{0.860673in}}%
\pgfpathlineto{\pgfqpoint{0.733451in}{0.870288in}}%
\pgfpathlineto{\pgfqpoint{0.742531in}{0.860986in}}%
\pgfpathlineto{\pgfqpoint{0.753215in}{0.851954in}}%
\pgfpathlineto{\pgfqpoint{0.765756in}{0.842970in}}%
\pgfpathlineto{\pgfqpoint{0.778258in}{0.834014in}}%
\pgfpathlineto{\pgfqpoint{0.768303in}{0.842422in}}%
\pgfpathlineto{\pgfqpoint{0.759828in}{0.851084in}}%
\pgfpathclose%
\pgfusepath{fill}%
\end{pgfscope}%
\begin{pgfscope}%
\pgfpathrectangle{\pgfqpoint{0.407592in}{0.222904in}}{\pgfqpoint{1.411163in}{1.411163in}}%
\pgfusepath{clip}%
\pgfsetbuttcap%
\pgfsetroundjoin%
\definecolor{currentfill}{rgb}{0.749020,0.772549,0.792157}%
\pgfsetfillcolor{currentfill}%
\pgfsetfillopacity{0.160000}%
\pgfsetlinewidth{0.000000pt}%
\definecolor{currentstroke}{rgb}{0.000000,0.000000,0.000000}%
\pgfsetstrokecolor{currentstroke}%
\pgfsetdash{}{0pt}%
\pgfpathmoveto{\pgfqpoint{1.303300in}{0.779931in}}%
\pgfpathlineto{\pgfqpoint{1.313386in}{0.787950in}}%
\pgfpathlineto{\pgfqpoint{1.323509in}{0.795998in}}%
\pgfpathlineto{\pgfqpoint{1.337847in}{0.802910in}}%
\pgfpathlineto{\pgfqpoint{1.351015in}{0.810164in}}%
\pgfpathlineto{\pgfqpoint{1.339744in}{0.801536in}}%
\pgfpathlineto{\pgfqpoint{1.328510in}{0.792936in}}%
\pgfpathlineto{\pgfqpoint{1.316436in}{0.786275in}}%
\pgfpathlineto{\pgfqpoint{1.303300in}{0.779931in}}%
\pgfpathclose%
\pgfusepath{fill}%
\end{pgfscope}%
\begin{pgfscope}%
\pgfpathrectangle{\pgfqpoint{0.407592in}{0.222904in}}{\pgfqpoint{1.411163in}{1.411163in}}%
\pgfusepath{clip}%
\pgfsetbuttcap%
\pgfsetroundjoin%
\definecolor{currentfill}{rgb}{0.749020,0.772549,0.792157}%
\pgfsetfillcolor{currentfill}%
\pgfsetfillopacity{0.160000}%
\pgfsetlinewidth{0.000000pt}%
\definecolor{currentstroke}{rgb}{0.000000,0.000000,0.000000}%
\pgfsetstrokecolor{currentstroke}%
\pgfsetdash{}{0pt}%
\pgfpathmoveto{\pgfqpoint{1.256703in}{0.753414in}}%
\pgfpathlineto{\pgfqpoint{1.265368in}{0.760851in}}%
\pgfpathlineto{\pgfqpoint{1.274068in}{0.768318in}}%
\pgfpathlineto{\pgfqpoint{1.289156in}{0.773935in}}%
\pgfpathlineto{\pgfqpoint{1.303300in}{0.779931in}}%
\pgfpathlineto{\pgfqpoint{1.293250in}{0.771941in}}%
\pgfpathlineto{\pgfqpoint{1.283236in}{0.763979in}}%
\pgfpathlineto{\pgfqpoint{1.270394in}{0.758523in}}%
\pgfpathlineto{\pgfqpoint{1.256703in}{0.753414in}}%
\pgfpathclose%
\pgfusepath{fill}%
\end{pgfscope}%
\begin{pgfscope}%
\pgfpathrectangle{\pgfqpoint{0.407592in}{0.222904in}}{\pgfqpoint{1.411163in}{1.411163in}}%
\pgfusepath{clip}%
\pgfsetbuttcap%
\pgfsetroundjoin%
\definecolor{currentfill}{rgb}{0.749020,0.772549,0.792157}%
\pgfsetfillcolor{currentfill}%
\pgfsetfillopacity{0.160000}%
\pgfsetlinewidth{0.000000pt}%
\definecolor{currentstroke}{rgb}{0.000000,0.000000,0.000000}%
\pgfsetstrokecolor{currentstroke}%
\pgfsetdash{}{0pt}%
\pgfpathmoveto{\pgfqpoint{1.351015in}{0.810164in}}%
\pgfpathlineto{\pgfqpoint{1.362323in}{0.818820in}}%
\pgfpathlineto{\pgfqpoint{1.373668in}{0.827504in}}%
\pgfpathlineto{\pgfqpoint{1.386591in}{0.835679in}}%
\pgfpathlineto{\pgfqpoint{1.374757in}{0.826688in}}%
\pgfpathlineto{\pgfqpoint{1.362959in}{0.817725in}}%
\pgfusepath{fill}%
\end{pgfscope}%
\begin{pgfscope}%
\pgfpathrectangle{\pgfqpoint{0.407592in}{0.222904in}}{\pgfqpoint{1.411163in}{1.411163in}}%
\pgfusepath{clip}%
\pgfsetbuttcap%
\pgfsetroundjoin%
\definecolor{currentfill}{rgb}{0.749020,0.772549,0.792157}%
\pgfsetfillcolor{currentfill}%
\pgfsetfillopacity{0.160000}%
\pgfsetlinewidth{0.000000pt}%
\definecolor{currentstroke}{rgb}{0.000000,0.000000,0.000000}%
\pgfsetstrokecolor{currentstroke}%
\pgfsetdash{}{0pt}%
\pgfpathmoveto{\pgfqpoint{1.031645in}{0.704359in}}%
\pgfpathlineto{\pgfqpoint{1.027840in}{0.710453in}}%
\pgfpathlineto{\pgfqpoint{1.024018in}{0.716575in}}%
\pgfpathlineto{\pgfqpoint{1.039871in}{0.715247in}}%
\pgfpathlineto{\pgfqpoint{1.055959in}{0.714373in}}%
\pgfpathlineto{\pgfqpoint{1.057939in}{0.708374in}}%
\pgfpathlineto{\pgfqpoint{1.059909in}{0.702404in}}%
\pgfpathlineto{\pgfqpoint{1.045674in}{0.703179in}}%
\pgfpathlineto{\pgfqpoint{1.031645in}{0.704359in}}%
\pgfpathclose%
\pgfusepath{fill}%
\end{pgfscope}%
\begin{pgfscope}%
\pgfpathrectangle{\pgfqpoint{0.407592in}{0.222904in}}{\pgfqpoint{1.411163in}{1.411163in}}%
\pgfusepath{clip}%
\pgfsetbuttcap%
\pgfsetroundjoin%
\definecolor{currentfill}{rgb}{0.749020,0.772549,0.792157}%
\pgfsetfillcolor{currentfill}%
\pgfsetfillopacity{0.160000}%
\pgfsetlinewidth{0.000000pt}%
\definecolor{currentstroke}{rgb}{0.000000,0.000000,0.000000}%
\pgfsetstrokecolor{currentstroke}%
\pgfsetdash{}{0pt}%
\pgfpathmoveto{\pgfqpoint{1.584645in}{0.996118in}}%
\pgfpathlineto{\pgfqpoint{1.598228in}{1.006711in}}%
\pgfpathlineto{\pgfqpoint{1.611841in}{1.017328in}}%
\pgfpathlineto{\pgfqpoint{1.619806in}{1.031041in}}%
\pgfpathlineto{\pgfqpoint{1.625359in}{1.044827in}}%
\pgfpathlineto{\pgfqpoint{1.611511in}{1.033588in}}%
\pgfpathlineto{\pgfqpoint{1.597687in}{1.022369in}}%
\pgfpathlineto{\pgfqpoint{1.592308in}{1.009206in}}%
\pgfpathlineto{\pgfqpoint{1.584645in}{0.996118in}}%
\pgfpathclose%
\pgfusepath{fill}%
\end{pgfscope}%
\begin{pgfscope}%
\pgfpathrectangle{\pgfqpoint{0.407592in}{0.222904in}}{\pgfqpoint{1.411163in}{1.411163in}}%
\pgfusepath{clip}%
\pgfsetbuttcap%
\pgfsetroundjoin%
\definecolor{currentfill}{rgb}{0.749020,0.772549,0.792157}%
\pgfsetfillcolor{currentfill}%
\pgfsetfillopacity{0.160000}%
\pgfsetlinewidth{0.000000pt}%
\definecolor{currentstroke}{rgb}{0.000000,0.000000,0.000000}%
\pgfsetstrokecolor{currentstroke}%
\pgfsetdash{}{0pt}%
\pgfpathmoveto{\pgfqpoint{1.088728in}{0.702077in}}%
\pgfpathlineto{\pgfqpoint{1.088630in}{0.708026in}}%
\pgfpathlineto{\pgfqpoint{1.088532in}{0.714004in}}%
\pgfpathlineto{\pgfqpoint{1.104862in}{0.714512in}}%
\pgfpathlineto{\pgfqpoint{1.121116in}{0.715477in}}%
\pgfpathlineto{\pgfqpoint{1.119332in}{0.709416in}}%
\pgfpathlineto{\pgfqpoint{1.117557in}{0.703384in}}%
\pgfpathlineto{\pgfqpoint{1.103175in}{0.702527in}}%
\pgfpathlineto{\pgfqpoint{1.088728in}{0.702077in}}%
\pgfpathclose%
\pgfusepath{fill}%
\end{pgfscope}%
\begin{pgfscope}%
\pgfpathrectangle{\pgfqpoint{0.407592in}{0.222904in}}{\pgfqpoint{1.411163in}{1.411163in}}%
\pgfusepath{clip}%
\pgfsetbuttcap%
\pgfsetroundjoin%
\definecolor{currentfill}{rgb}{0.749020,0.772549,0.792157}%
\pgfsetfillcolor{currentfill}%
\pgfsetfillopacity{0.160000}%
\pgfsetlinewidth{0.000000pt}%
\definecolor{currentstroke}{rgb}{0.000000,0.000000,0.000000}%
\pgfsetstrokecolor{currentstroke}%
\pgfsetdash{}{0pt}%
\pgfpathmoveto{\pgfqpoint{0.937890in}{0.733547in}}%
\pgfpathlineto{\pgfqpoint{0.929100in}{0.740646in}}%
\pgfpathlineto{\pgfqpoint{0.920274in}{0.747773in}}%
\pgfpathlineto{\pgfqpoint{0.934693in}{0.743513in}}%
\pgfpathlineto{\pgfqpoint{0.949845in}{0.739675in}}%
\pgfpathlineto{\pgfqpoint{0.957151in}{0.732953in}}%
\pgfpathlineto{\pgfqpoint{0.964427in}{0.726261in}}%
\pgfpathlineto{\pgfqpoint{0.950833in}{0.729714in}}%
\pgfpathlineto{\pgfqpoint{0.937890in}{0.733547in}}%
\pgfpathclose%
\pgfusepath{fill}%
\end{pgfscope}%
\begin{pgfscope}%
\pgfpathrectangle{\pgfqpoint{0.407592in}{0.222904in}}{\pgfqpoint{1.411163in}{1.411163in}}%
\pgfusepath{clip}%
\pgfsetbuttcap%
\pgfsetroundjoin%
\definecolor{currentfill}{rgb}{0.749020,0.772549,0.792157}%
\pgfsetfillcolor{currentfill}%
\pgfsetfillopacity{0.160000}%
\pgfsetlinewidth{0.000000pt}%
\definecolor{currentstroke}{rgb}{0.000000,0.000000,0.000000}%
\pgfsetstrokecolor{currentstroke}%
\pgfsetdash{}{0pt}%
\pgfpathmoveto{\pgfqpoint{1.059909in}{0.702404in}}%
\pgfpathlineto{\pgfqpoint{1.057939in}{0.708374in}}%
\pgfpathlineto{\pgfqpoint{1.055959in}{0.714373in}}%
\pgfpathlineto{\pgfqpoint{1.072205in}{0.713958in}}%
\pgfpathlineto{\pgfqpoint{1.088532in}{0.714004in}}%
\pgfpathlineto{\pgfqpoint{1.088630in}{0.708026in}}%
\pgfpathlineto{\pgfqpoint{1.088728in}{0.702077in}}%
\pgfpathlineto{\pgfqpoint{1.074283in}{0.702035in}}%
\pgfpathlineto{\pgfqpoint{1.059909in}{0.702404in}}%
\pgfpathclose%
\pgfusepath{fill}%
\end{pgfscope}%
\begin{pgfscope}%
\pgfpathrectangle{\pgfqpoint{0.407592in}{0.222904in}}{\pgfqpoint{1.411163in}{1.411163in}}%
\pgfusepath{clip}%
\pgfsetbuttcap%
\pgfsetroundjoin%
\definecolor{currentfill}{rgb}{0.749020,0.772549,0.792157}%
\pgfsetfillcolor{currentfill}%
\pgfsetfillopacity{0.160000}%
\pgfsetlinewidth{0.000000pt}%
\definecolor{currentstroke}{rgb}{0.000000,0.000000,0.000000}%
\pgfsetstrokecolor{currentstroke}%
\pgfsetdash{}{0pt}%
\pgfpathmoveto{\pgfqpoint{0.435692in}{1.164522in}}%
\pgfpathlineto{\pgfqpoint{0.421586in}{1.176460in}}%
\pgfpathlineto{\pgfqpoint{0.407462in}{1.188415in}}%
\pgfpathlineto{\pgfqpoint{0.404193in}{1.171737in}}%
\pgfpathlineto{\pgfqpoint{0.403920in}{1.154973in}}%
\pgfpathlineto{\pgfqpoint{0.418274in}{1.143587in}}%
\pgfpathlineto{\pgfqpoint{0.432600in}{1.132222in}}%
\pgfpathlineto{\pgfqpoint{0.432709in}{1.148410in}}%
\pgfpathlineto{\pgfqpoint{0.435692in}{1.164522in}}%
\pgfpathclose%
\pgfusepath{fill}%
\end{pgfscope}%
\begin{pgfscope}%
\pgfpathrectangle{\pgfqpoint{0.407592in}{0.222904in}}{\pgfqpoint{1.411163in}{1.411163in}}%
\pgfusepath{clip}%
\pgfsetbuttcap%
\pgfsetroundjoin%
\definecolor{currentfill}{rgb}{0.749020,0.772549,0.792157}%
\pgfsetfillcolor{currentfill}%
\pgfsetfillopacity{0.160000}%
\pgfsetlinewidth{0.000000pt}%
\definecolor{currentstroke}{rgb}{0.000000,0.000000,0.000000}%
\pgfsetstrokecolor{currentstroke}%
\pgfsetdash{}{0pt}%
\pgfpathmoveto{\pgfqpoint{0.574411in}{1.019725in}}%
\pgfpathlineto{\pgfqpoint{0.560348in}{1.030882in}}%
\pgfpathlineto{\pgfqpoint{0.546258in}{1.042058in}}%
\pgfpathlineto{\pgfqpoint{0.549080in}{1.028282in}}%
\pgfpathlineto{\pgfqpoint{0.554333in}{1.014591in}}%
\pgfpathlineto{\pgfqpoint{0.568331in}{1.004037in}}%
\pgfpathlineto{\pgfqpoint{0.582297in}{0.993507in}}%
\pgfpathlineto{\pgfqpoint{0.577202in}{1.006572in}}%
\pgfpathlineto{\pgfqpoint{0.574411in}{1.019725in}}%
\pgfpathclose%
\pgfusepath{fill}%
\end{pgfscope}%
\begin{pgfscope}%
\pgfpathrectangle{\pgfqpoint{0.407592in}{0.222904in}}{\pgfqpoint{1.411163in}{1.411163in}}%
\pgfusepath{clip}%
\pgfsetbuttcap%
\pgfsetroundjoin%
\definecolor{currentfill}{rgb}{0.749020,0.772549,0.792157}%
\pgfsetfillcolor{currentfill}%
\pgfsetfillopacity{0.160000}%
\pgfsetlinewidth{0.000000pt}%
\definecolor{currentstroke}{rgb}{0.000000,0.000000,0.000000}%
\pgfsetstrokecolor{currentstroke}%
\pgfsetdash{}{0pt}%
\pgfpathmoveto{\pgfqpoint{1.765228in}{1.158340in}}%
\pgfpathlineto{\pgfqpoint{1.779355in}{1.169805in}}%
\pgfpathlineto{\pgfqpoint{1.793508in}{1.181291in}}%
\pgfpathlineto{\pgfqpoint{1.796670in}{1.198609in}}%
\pgfpathlineto{\pgfqpoint{1.796697in}{1.215815in}}%
\pgfpathlineto{\pgfqpoint{1.782640in}{1.203771in}}%
\pgfpathlineto{\pgfqpoint{1.768602in}{1.191744in}}%
\pgfpathlineto{\pgfqpoint{1.768420in}{1.175092in}}%
\pgfpathlineto{\pgfqpoint{1.765228in}{1.158340in}}%
\pgfpathclose%
\pgfusepath{fill}%
\end{pgfscope}%
\begin{pgfscope}%
\pgfpathrectangle{\pgfqpoint{0.407592in}{0.222904in}}{\pgfqpoint{1.411163in}{1.411163in}}%
\pgfusepath{clip}%
\pgfsetbuttcap%
\pgfsetroundjoin%
\definecolor{currentfill}{rgb}{0.749020,0.772549,0.792157}%
\pgfsetfillcolor{currentfill}%
\pgfsetfillopacity{0.160000}%
\pgfsetlinewidth{0.000000pt}%
\definecolor{currentstroke}{rgb}{0.000000,0.000000,0.000000}%
\pgfsetstrokecolor{currentstroke}%
\pgfsetdash{}{0pt}%
\pgfpathmoveto{\pgfqpoint{1.484750in}{0.911159in}}%
\pgfpathlineto{\pgfqpoint{1.497623in}{0.920969in}}%
\pgfpathlineto{\pgfqpoint{1.510531in}{0.930804in}}%
\pgfpathlineto{\pgfqpoint{1.521573in}{0.942299in}}%
\pgfpathlineto{\pgfqpoint{1.530618in}{0.953983in}}%
\pgfpathlineto{\pgfqpoint{1.517186in}{0.943508in}}%
\pgfpathlineto{\pgfqpoint{1.503785in}{0.933057in}}%
\pgfpathlineto{\pgfqpoint{1.495203in}{0.922017in}}%
\pgfpathlineto{\pgfqpoint{1.484750in}{0.911159in}}%
\pgfpathclose%
\pgfusepath{fill}%
\end{pgfscope}%
\begin{pgfscope}%
\pgfpathrectangle{\pgfqpoint{0.407592in}{0.222904in}}{\pgfqpoint{1.411163in}{1.411163in}}%
\pgfusepath{clip}%
\pgfsetbuttcap%
\pgfsetroundjoin%
\definecolor{currentfill}{rgb}{0.749020,0.772549,0.792157}%
\pgfsetfillcolor{currentfill}%
\pgfsetfillopacity{0.160000}%
\pgfsetlinewidth{0.000000pt}%
\definecolor{currentstroke}{rgb}{0.000000,0.000000,0.000000}%
\pgfsetstrokecolor{currentstroke}%
\pgfsetdash{}{0pt}%
\pgfpathmoveto{\pgfqpoint{1.183854in}{0.723786in}}%
\pgfpathlineto{\pgfqpoint{1.189272in}{0.730341in}}%
\pgfpathlineto{\pgfqpoint{1.194714in}{0.736924in}}%
\pgfpathlineto{\pgfqpoint{1.211156in}{0.740410in}}%
\pgfpathlineto{\pgfqpoint{1.227016in}{0.744334in}}%
\pgfpathlineto{\pgfqpoint{1.219906in}{0.737379in}}%
\pgfpathlineto{\pgfqpoint{1.212825in}{0.730453in}}%
\pgfpathlineto{\pgfqpoint{1.198598in}{0.726922in}}%
\pgfpathlineto{\pgfqpoint{1.183854in}{0.723786in}}%
\pgfpathclose%
\pgfusepath{fill}%
\end{pgfscope}%
\begin{pgfscope}%
\pgfpathrectangle{\pgfqpoint{0.407592in}{0.222904in}}{\pgfqpoint{1.411163in}{1.411163in}}%
\pgfusepath{clip}%
\pgfsetbuttcap%
\pgfsetroundjoin%
\definecolor{currentfill}{rgb}{0.749020,0.772549,0.792157}%
\pgfsetfillcolor{currentfill}%
\pgfsetfillopacity{0.160000}%
\pgfsetlinewidth{0.000000pt}%
\definecolor{currentstroke}{rgb}{0.000000,0.000000,0.000000}%
\pgfsetstrokecolor{currentstroke}%
\pgfsetdash{}{0pt}%
\pgfpathmoveto{\pgfqpoint{0.803146in}{0.816186in}}%
\pgfpathlineto{\pgfqpoint{0.790721in}{0.825086in}}%
\pgfpathlineto{\pgfqpoint{0.778258in}{0.834014in}}%
\pgfpathlineto{\pgfqpoint{0.789658in}{0.825903in}}%
\pgfpathlineto{\pgfqpoint{0.802460in}{0.818128in}}%
\pgfpathlineto{\pgfqpoint{0.814009in}{0.809797in}}%
\pgfpathlineto{\pgfqpoint{0.825518in}{0.801494in}}%
\pgfpathlineto{\pgfqpoint{0.813690in}{0.808683in}}%
\pgfpathlineto{\pgfqpoint{0.803146in}{0.816186in}}%
\pgfpathclose%
\pgfusepath{fill}%
\end{pgfscope}%
\begin{pgfscope}%
\pgfpathrectangle{\pgfqpoint{0.407592in}{0.222904in}}{\pgfqpoint{1.411163in}{1.411163in}}%
\pgfusepath{clip}%
\pgfsetbuttcap%
\pgfsetroundjoin%
\definecolor{currentfill}{rgb}{0.749020,0.772549,0.792157}%
\pgfsetfillcolor{currentfill}%
\pgfsetfillopacity{0.160000}%
\pgfsetlinewidth{0.000000pt}%
\definecolor{currentstroke}{rgb}{0.000000,0.000000,0.000000}%
\pgfsetstrokecolor{currentstroke}%
\pgfsetdash{}{0pt}%
\pgfpathmoveto{\pgfqpoint{0.665424in}{0.930830in}}%
\pgfpathlineto{\pgfqpoint{0.651648in}{0.941217in}}%
\pgfpathlineto{\pgfqpoint{0.637842in}{0.951627in}}%
\pgfpathlineto{\pgfqpoint{0.644631in}{0.939976in}}%
\pgfpathlineto{\pgfqpoint{0.653460in}{0.928526in}}%
\pgfpathlineto{\pgfqpoint{0.666882in}{0.918754in}}%
\pgfpathlineto{\pgfqpoint{0.680268in}{0.909008in}}%
\pgfpathlineto{\pgfqpoint{0.671892in}{0.919822in}}%
\pgfpathlineto{\pgfqpoint{0.665424in}{0.930830in}}%
\pgfpathclose%
\pgfusepath{fill}%
\end{pgfscope}%
\begin{pgfscope}%
\pgfpathrectangle{\pgfqpoint{0.407592in}{0.222904in}}{\pgfqpoint{1.411163in}{1.411163in}}%
\pgfusepath{clip}%
\pgfsetbuttcap%
\pgfsetroundjoin%
\definecolor{currentfill}{rgb}{0.749020,0.772549,0.792157}%
\pgfsetfillcolor{currentfill}%
\pgfsetfillopacity{0.160000}%
\pgfsetlinewidth{0.000000pt}%
\definecolor{currentstroke}{rgb}{0.000000,0.000000,0.000000}%
\pgfsetstrokecolor{currentstroke}%
\pgfsetdash{}{0pt}%
\pgfpathmoveto{\pgfqpoint{0.893888in}{0.757470in}}%
\pgfpathlineto{\pgfqpoint{0.883681in}{0.765109in}}%
\pgfpathlineto{\pgfqpoint{0.873435in}{0.772777in}}%
\pgfpathlineto{\pgfqpoint{0.887497in}{0.767239in}}%
\pgfpathlineto{\pgfqpoint{0.902514in}{0.762114in}}%
\pgfpathlineto{\pgfqpoint{0.911412in}{0.754929in}}%
\pgfpathlineto{\pgfqpoint{0.920274in}{0.747773in}}%
\pgfpathlineto{\pgfqpoint{0.906652in}{0.752433in}}%
\pgfpathlineto{\pgfqpoint{0.893888in}{0.757470in}}%
\pgfpathclose%
\pgfusepath{fill}%
\end{pgfscope}%
\begin{pgfscope}%
\pgfpathrectangle{\pgfqpoint{0.407592in}{0.222904in}}{\pgfqpoint{1.411163in}{1.411163in}}%
\pgfusepath{clip}%
\pgfsetbuttcap%
\pgfsetroundjoin%
\definecolor{currentfill}{rgb}{0.749020,0.772549,0.792157}%
\pgfsetfillcolor{currentfill}%
\pgfsetfillopacity{0.160000}%
\pgfsetlinewidth{0.000000pt}%
\definecolor{currentstroke}{rgb}{0.000000,0.000000,0.000000}%
\pgfsetstrokecolor{currentstroke}%
\pgfsetdash{}{0pt}%
\pgfpathmoveto{\pgfqpoint{1.410368in}{0.853742in}}%
\pgfpathlineto{\pgfqpoint{1.422311in}{0.862815in}}%
\pgfpathlineto{\pgfqpoint{1.434292in}{0.871917in}}%
\pgfpathlineto{\pgfqpoint{1.447537in}{0.881631in}}%
\pgfpathlineto{\pgfqpoint{1.459107in}{0.891619in}}%
\pgfpathlineto{\pgfqpoint{1.446336in}{0.881888in}}%
\pgfpathlineto{\pgfqpoint{1.433600in}{0.872182in}}%
\pgfpathlineto{\pgfqpoint{1.422760in}{0.862832in}}%
\pgfpathlineto{\pgfqpoint{1.410368in}{0.853742in}}%
\pgfpathclose%
\pgfusepath{fill}%
\end{pgfscope}%
\begin{pgfscope}%
\pgfpathrectangle{\pgfqpoint{0.407592in}{0.222904in}}{\pgfqpoint{1.411163in}{1.411163in}}%
\pgfusepath{clip}%
\pgfsetbuttcap%
\pgfsetroundjoin%
\definecolor{currentfill}{rgb}{0.749020,0.772549,0.792157}%
\pgfsetfillcolor{currentfill}%
\pgfsetfillopacity{0.160000}%
\pgfsetlinewidth{0.000000pt}%
\definecolor{currentstroke}{rgb}{0.000000,0.000000,0.000000}%
\pgfsetstrokecolor{currentstroke}%
\pgfsetdash{}{0pt}%
\pgfpathmoveto{\pgfqpoint{1.611841in}{1.017328in}}%
\pgfpathlineto{\pgfqpoint{1.625485in}{1.027969in}}%
\pgfpathlineto{\pgfqpoint{1.639161in}{1.038634in}}%
\pgfpathlineto{\pgfqpoint{1.647416in}{1.052965in}}%
\pgfpathlineto{\pgfqpoint{1.653131in}{1.067365in}}%
\pgfpathlineto{\pgfqpoint{1.639232in}{1.056086in}}%
\pgfpathlineto{\pgfqpoint{1.625359in}{1.044827in}}%
\pgfpathlineto{\pgfqpoint{1.619806in}{1.031041in}}%
\pgfpathlineto{\pgfqpoint{1.611841in}{1.017328in}}%
\pgfpathclose%
\pgfusepath{fill}%
\end{pgfscope}%
\begin{pgfscope}%
\pgfpathrectangle{\pgfqpoint{0.407592in}{0.222904in}}{\pgfqpoint{1.411163in}{1.411163in}}%
\pgfusepath{clip}%
\pgfsetbuttcap%
\pgfsetroundjoin%
\definecolor{currentfill}{rgb}{0.749020,0.772549,0.792157}%
\pgfsetfillcolor{currentfill}%
\pgfsetfillopacity{0.160000}%
\pgfsetlinewidth{0.000000pt}%
\definecolor{currentstroke}{rgb}{0.000000,0.000000,0.000000}%
\pgfsetstrokecolor{currentstroke}%
\pgfsetdash{}{0pt}%
\pgfpathmoveto{\pgfqpoint{0.407462in}{1.188415in}}%
\pgfpathlineto{\pgfqpoint{0.393317in}{1.200387in}}%
\pgfpathlineto{\pgfqpoint{0.379152in}{1.212375in}}%
\pgfpathlineto{\pgfqpoint{0.375583in}{1.195141in}}%
\pgfpathlineto{\pgfqpoint{0.375133in}{1.177810in}}%
\pgfpathlineto{\pgfqpoint{0.389540in}{1.166381in}}%
\pgfpathlineto{\pgfqpoint{0.403920in}{1.154973in}}%
\pgfpathlineto{\pgfqpoint{0.404193in}{1.171737in}}%
\pgfpathlineto{\pgfqpoint{0.407462in}{1.188415in}}%
\pgfpathclose%
\pgfusepath{fill}%
\end{pgfscope}%
\begin{pgfscope}%
\pgfpathrectangle{\pgfqpoint{0.407592in}{0.222904in}}{\pgfqpoint{1.411163in}{1.411163in}}%
\pgfusepath{clip}%
\pgfsetbuttcap%
\pgfsetroundjoin%
\definecolor{currentfill}{rgb}{0.749020,0.772549,0.792157}%
\pgfsetfillcolor{currentfill}%
\pgfsetfillopacity{0.160000}%
\pgfsetlinewidth{0.000000pt}%
\definecolor{currentstroke}{rgb}{0.000000,0.000000,0.000000}%
\pgfsetstrokecolor{currentstroke}%
\pgfsetdash{}{0pt}%
\pgfpathmoveto{\pgfqpoint{0.848416in}{0.784976in}}%
\pgfpathlineto{\pgfqpoint{0.836986in}{0.793221in}}%
\pgfpathlineto{\pgfqpoint{0.825518in}{0.801494in}}%
\pgfpathlineto{\pgfqpoint{0.838581in}{0.794655in}}%
\pgfpathlineto{\pgfqpoint{0.852825in}{0.788200in}}%
\pgfpathlineto{\pgfqpoint{0.863150in}{0.780474in}}%
\pgfpathlineto{\pgfqpoint{0.873435in}{0.772777in}}%
\pgfpathlineto{\pgfqpoint{0.860389in}{0.778699in}}%
\pgfpathlineto{\pgfqpoint{0.848416in}{0.784976in}}%
\pgfpathclose%
\pgfusepath{fill}%
\end{pgfscope}%
\begin{pgfscope}%
\pgfpathrectangle{\pgfqpoint{0.407592in}{0.222904in}}{\pgfqpoint{1.411163in}{1.411163in}}%
\pgfusepath{clip}%
\pgfsetbuttcap%
\pgfsetroundjoin%
\definecolor{currentfill}{rgb}{0.749020,0.772549,0.792157}%
\pgfsetfillcolor{currentfill}%
\pgfsetfillopacity{0.160000}%
\pgfsetlinewidth{0.000000pt}%
\definecolor{currentstroke}{rgb}{0.000000,0.000000,0.000000}%
\pgfsetstrokecolor{currentstroke}%
\pgfsetdash{}{0pt}%
\pgfpathmoveto{\pgfqpoint{0.964427in}{0.726261in}}%
\pgfpathlineto{\pgfqpoint{0.957151in}{0.732953in}}%
\pgfpathlineto{\pgfqpoint{0.949845in}{0.739675in}}%
\pgfpathlineto{\pgfqpoint{0.965658in}{0.736279in}}%
\pgfpathlineto{\pgfqpoint{0.982060in}{0.733342in}}%
\pgfpathlineto{\pgfqpoint{0.987700in}{0.726939in}}%
\pgfpathlineto{\pgfqpoint{0.993314in}{0.720565in}}%
\pgfpathlineto{\pgfqpoint{0.978609in}{0.723206in}}%
\pgfpathlineto{\pgfqpoint{0.964427in}{0.726261in}}%
\pgfpathclose%
\pgfusepath{fill}%
\end{pgfscope}%
\begin{pgfscope}%
\pgfpathrectangle{\pgfqpoint{0.407592in}{0.222904in}}{\pgfqpoint{1.411163in}{1.411163in}}%
\pgfusepath{clip}%
\pgfsetbuttcap%
\pgfsetroundjoin%
\definecolor{currentfill}{rgb}{0.749020,0.772549,0.792157}%
\pgfsetfillcolor{currentfill}%
\pgfsetfillopacity{0.160000}%
\pgfsetlinewidth{0.000000pt}%
\definecolor{currentstroke}{rgb}{0.000000,0.000000,0.000000}%
\pgfsetstrokecolor{currentstroke}%
\pgfsetdash{}{0pt}%
\pgfpathmoveto{\pgfqpoint{0.733451in}{0.870288in}}%
\pgfpathlineto{\pgfqpoint{0.720209in}{0.879928in}}%
\pgfpathlineto{\pgfqpoint{0.706932in}{0.889595in}}%
\pgfpathlineto{\pgfqpoint{0.716608in}{0.879658in}}%
\pgfpathlineto{\pgfqpoint{0.728016in}{0.870005in}}%
\pgfpathlineto{\pgfqpoint{0.740635in}{0.860966in}}%
\pgfpathlineto{\pgfqpoint{0.753215in}{0.851954in}}%
\pgfpathlineto{\pgfqpoint{0.742531in}{0.860986in}}%
\pgfpathlineto{\pgfqpoint{0.733451in}{0.870288in}}%
\pgfpathclose%
\pgfusepath{fill}%
\end{pgfscope}%
\begin{pgfscope}%
\pgfpathrectangle{\pgfqpoint{0.407592in}{0.222904in}}{\pgfqpoint{1.411163in}{1.411163in}}%
\pgfusepath{clip}%
\pgfsetbuttcap%
\pgfsetroundjoin%
\definecolor{currentfill}{rgb}{0.749020,0.772549,0.792157}%
\pgfsetfillcolor{currentfill}%
\pgfsetfillopacity{0.160000}%
\pgfsetlinewidth{0.000000pt}%
\definecolor{currentstroke}{rgb}{0.000000,0.000000,0.000000}%
\pgfsetstrokecolor{currentstroke}%
\pgfsetdash{}{0pt}%
\pgfpathmoveto{\pgfqpoint{1.793508in}{1.181291in}}%
\pgfpathlineto{\pgfqpoint{1.807686in}{1.192798in}}%
\pgfpathlineto{\pgfqpoint{1.821891in}{1.204326in}}%
\pgfpathlineto{\pgfqpoint{1.825009in}{1.222200in}}%
\pgfpathlineto{\pgfqpoint{1.824868in}{1.239951in}}%
\pgfpathlineto{\pgfqpoint{1.810773in}{1.227874in}}%
\pgfpathlineto{\pgfqpoint{1.796697in}{1.215815in}}%
\pgfpathlineto{\pgfqpoint{1.796670in}{1.198609in}}%
\pgfpathlineto{\pgfqpoint{1.793508in}{1.181291in}}%
\pgfpathclose%
\pgfusepath{fill}%
\end{pgfscope}%
\begin{pgfscope}%
\pgfpathrectangle{\pgfqpoint{0.407592in}{0.222904in}}{\pgfqpoint{1.411163in}{1.411163in}}%
\pgfusepath{clip}%
\pgfsetbuttcap%
\pgfsetroundjoin%
\definecolor{currentfill}{rgb}{0.749020,0.772549,0.792157}%
\pgfsetfillcolor{currentfill}%
\pgfsetfillopacity{0.160000}%
\pgfsetlinewidth{0.000000pt}%
\definecolor{currentstroke}{rgb}{0.000000,0.000000,0.000000}%
\pgfsetstrokecolor{currentstroke}%
\pgfsetdash{}{0pt}%
\pgfpathmoveto{\pgfqpoint{1.227016in}{0.744334in}}%
\pgfpathlineto{\pgfqpoint{1.234156in}{0.751318in}}%
\pgfpathlineto{\pgfqpoint{1.241326in}{0.758332in}}%
\pgfpathlineto{\pgfqpoint{1.258100in}{0.763108in}}%
\pgfpathlineto{\pgfqpoint{1.274068in}{0.768318in}}%
\pgfpathlineto{\pgfqpoint{1.265368in}{0.760851in}}%
\pgfpathlineto{\pgfqpoint{1.256703in}{0.753414in}}%
\pgfpathlineto{\pgfqpoint{1.242222in}{0.748676in}}%
\pgfpathlineto{\pgfqpoint{1.227016in}{0.744334in}}%
\pgfpathclose%
\pgfusepath{fill}%
\end{pgfscope}%
\begin{pgfscope}%
\pgfpathrectangle{\pgfqpoint{0.407592in}{0.222904in}}{\pgfqpoint{1.411163in}{1.411163in}}%
\pgfusepath{clip}%
\pgfsetbuttcap%
\pgfsetroundjoin%
\definecolor{currentfill}{rgb}{0.749020,0.772549,0.792157}%
\pgfsetfillcolor{currentfill}%
\pgfsetfillopacity{0.160000}%
\pgfsetlinewidth{0.000000pt}%
\definecolor{currentstroke}{rgb}{0.000000,0.000000,0.000000}%
\pgfsetstrokecolor{currentstroke}%
\pgfsetdash{}{0pt}%
\pgfpathmoveto{\pgfqpoint{1.153092in}{0.718760in}}%
\pgfpathlineto{\pgfqpoint{1.156735in}{0.725033in}}%
\pgfpathlineto{\pgfqpoint{1.160396in}{0.731335in}}%
\pgfpathlineto{\pgfqpoint{1.177768in}{0.733894in}}%
\pgfpathlineto{\pgfqpoint{1.194714in}{0.736924in}}%
\pgfpathlineto{\pgfqpoint{1.189272in}{0.730341in}}%
\pgfpathlineto{\pgfqpoint{1.183854in}{0.723786in}}%
\pgfpathlineto{\pgfqpoint{1.168661in}{0.721061in}}%
\pgfpathlineto{\pgfqpoint{1.153092in}{0.718760in}}%
\pgfpathclose%
\pgfusepath{fill}%
\end{pgfscope}%
\begin{pgfscope}%
\pgfpathrectangle{\pgfqpoint{0.407592in}{0.222904in}}{\pgfqpoint{1.411163in}{1.411163in}}%
\pgfusepath{clip}%
\pgfsetbuttcap%
\pgfsetroundjoin%
\definecolor{currentfill}{rgb}{0.749020,0.772549,0.792157}%
\pgfsetfillcolor{currentfill}%
\pgfsetfillopacity{0.160000}%
\pgfsetlinewidth{0.000000pt}%
\definecolor{currentstroke}{rgb}{0.000000,0.000000,0.000000}%
\pgfsetstrokecolor{currentstroke}%
\pgfsetdash{}{0pt}%
\pgfpathmoveto{\pgfqpoint{0.546258in}{1.042058in}}%
\pgfpathlineto{\pgfqpoint{0.532143in}{1.053256in}}%
\pgfpathlineto{\pgfqpoint{0.518002in}{1.064474in}}%
\pgfpathlineto{\pgfqpoint{0.520840in}{1.050082in}}%
\pgfpathlineto{\pgfqpoint{0.526239in}{1.035773in}}%
\pgfpathlineto{\pgfqpoint{0.540302in}{1.025170in}}%
\pgfpathlineto{\pgfqpoint{0.554333in}{1.014591in}}%
\pgfpathlineto{\pgfqpoint{0.549080in}{1.028282in}}%
\pgfpathlineto{\pgfqpoint{0.546258in}{1.042058in}}%
\pgfpathclose%
\pgfusepath{fill}%
\end{pgfscope}%
\begin{pgfscope}%
\pgfpathrectangle{\pgfqpoint{0.407592in}{0.222904in}}{\pgfqpoint{1.411163in}{1.411163in}}%
\pgfusepath{clip}%
\pgfsetbuttcap%
\pgfsetroundjoin%
\definecolor{currentfill}{rgb}{0.749020,0.772549,0.792157}%
\pgfsetfillcolor{currentfill}%
\pgfsetfillopacity{0.160000}%
\pgfsetlinewidth{0.000000pt}%
\definecolor{currentstroke}{rgb}{0.000000,0.000000,0.000000}%
\pgfsetstrokecolor{currentstroke}%
\pgfsetdash{}{0pt}%
\pgfpathmoveto{\pgfqpoint{1.373668in}{0.827504in}}%
\pgfpathlineto{\pgfqpoint{1.385049in}{0.836217in}}%
\pgfpathlineto{\pgfqpoint{1.396468in}{0.844958in}}%
\pgfpathlineto{\pgfqpoint{1.410368in}{0.853742in}}%
\pgfpathlineto{\pgfqpoint{1.398461in}{0.844696in}}%
\pgfpathlineto{\pgfqpoint{1.386591in}{0.835679in}}%
\pgfusepath{fill}%
\end{pgfscope}%
\begin{pgfscope}%
\pgfpathrectangle{\pgfqpoint{0.407592in}{0.222904in}}{\pgfqpoint{1.411163in}{1.411163in}}%
\pgfusepath{clip}%
\pgfsetbuttcap%
\pgfsetroundjoin%
\definecolor{currentfill}{rgb}{0.749020,0.772549,0.792157}%
\pgfsetfillcolor{currentfill}%
\pgfsetfillopacity{0.160000}%
\pgfsetlinewidth{0.000000pt}%
\definecolor{currentstroke}{rgb}{0.000000,0.000000,0.000000}%
\pgfsetstrokecolor{currentstroke}%
\pgfsetdash{}{0pt}%
\pgfpathmoveto{\pgfqpoint{1.323509in}{0.795998in}}%
\pgfpathlineto{\pgfqpoint{1.333668in}{0.804075in}}%
\pgfpathlineto{\pgfqpoint{1.343864in}{0.812181in}}%
\pgfpathlineto{\pgfqpoint{1.359406in}{0.819659in}}%
\pgfpathlineto{\pgfqpoint{1.373668in}{0.827504in}}%
\pgfpathlineto{\pgfqpoint{1.362323in}{0.818820in}}%
\pgfpathlineto{\pgfqpoint{1.351015in}{0.810164in}}%
\pgfpathlineto{\pgfqpoint{1.337847in}{0.802910in}}%
\pgfpathlineto{\pgfqpoint{1.323509in}{0.795998in}}%
\pgfpathclose%
\pgfusepath{fill}%
\end{pgfscope}%
\begin{pgfscope}%
\pgfpathrectangle{\pgfqpoint{0.407592in}{0.222904in}}{\pgfqpoint{1.411163in}{1.411163in}}%
\pgfusepath{clip}%
\pgfsetbuttcap%
\pgfsetroundjoin%
\definecolor{currentfill}{rgb}{0.749020,0.772549,0.792157}%
\pgfsetfillcolor{currentfill}%
\pgfsetfillopacity{0.160000}%
\pgfsetlinewidth{0.000000pt}%
\definecolor{currentstroke}{rgb}{0.000000,0.000000,0.000000}%
\pgfsetstrokecolor{currentstroke}%
\pgfsetdash{}{0pt}%
\pgfpathmoveto{\pgfqpoint{1.510531in}{0.930804in}}%
\pgfpathlineto{\pgfqpoint{1.523473in}{0.940667in}}%
\pgfpathlineto{\pgfqpoint{1.536450in}{0.950555in}}%
\pgfpathlineto{\pgfqpoint{1.548074in}{0.962682in}}%
\pgfpathlineto{\pgfqpoint{1.557570in}{0.975003in}}%
\pgfpathlineto{\pgfqpoint{1.544079in}{0.964481in}}%
\pgfpathlineto{\pgfqpoint{1.530618in}{0.953983in}}%
\pgfpathlineto{\pgfqpoint{1.521573in}{0.942299in}}%
\pgfpathlineto{\pgfqpoint{1.510531in}{0.930804in}}%
\pgfpathclose%
\pgfusepath{fill}%
\end{pgfscope}%
\begin{pgfscope}%
\pgfpathrectangle{\pgfqpoint{0.407592in}{0.222904in}}{\pgfqpoint{1.411163in}{1.411163in}}%
\pgfusepath{clip}%
\pgfsetbuttcap%
\pgfsetroundjoin%
\definecolor{currentfill}{rgb}{0.749020,0.772549,0.792157}%
\pgfsetfillcolor{currentfill}%
\pgfsetfillopacity{0.160000}%
\pgfsetlinewidth{0.000000pt}%
\definecolor{currentstroke}{rgb}{0.000000,0.000000,0.000000}%
\pgfsetstrokecolor{currentstroke}%
\pgfsetdash{}{0pt}%
\pgfpathmoveto{\pgfqpoint{1.274068in}{0.768318in}}%
\pgfpathlineto{\pgfqpoint{1.282801in}{0.775814in}}%
\pgfpathlineto{\pgfqpoint{1.291569in}{0.783339in}}%
\pgfpathlineto{\pgfqpoint{1.308060in}{0.789463in}}%
\pgfpathlineto{\pgfqpoint{1.323509in}{0.795998in}}%
\pgfpathlineto{\pgfqpoint{1.313386in}{0.787950in}}%
\pgfpathlineto{\pgfqpoint{1.303300in}{0.779931in}}%
\pgfpathlineto{\pgfqpoint{1.289156in}{0.773935in}}%
\pgfpathlineto{\pgfqpoint{1.274068in}{0.768318in}}%
\pgfpathclose%
\pgfusepath{fill}%
\end{pgfscope}%
\begin{pgfscope}%
\pgfpathrectangle{\pgfqpoint{0.407592in}{0.222904in}}{\pgfqpoint{1.411163in}{1.411163in}}%
\pgfusepath{clip}%
\pgfsetbuttcap%
\pgfsetroundjoin%
\definecolor{currentfill}{rgb}{0.749020,0.772549,0.792157}%
\pgfsetfillcolor{currentfill}%
\pgfsetfillopacity{0.160000}%
\pgfsetlinewidth{0.000000pt}%
\definecolor{currentstroke}{rgb}{0.000000,0.000000,0.000000}%
\pgfsetstrokecolor{currentstroke}%
\pgfsetdash{}{0pt}%
\pgfpathmoveto{\pgfqpoint{0.993314in}{0.720565in}}%
\pgfpathlineto{\pgfqpoint{0.987700in}{0.726939in}}%
\pgfpathlineto{\pgfqpoint{0.982060in}{0.733342in}}%
\pgfpathlineto{\pgfqpoint{0.998974in}{0.730880in}}%
\pgfpathlineto{\pgfqpoint{1.016319in}{0.728905in}}%
\pgfpathlineto{\pgfqpoint{1.020177in}{0.722725in}}%
\pgfpathlineto{\pgfqpoint{1.024018in}{0.716575in}}%
\pgfpathlineto{\pgfqpoint{1.008474in}{0.718351in}}%
\pgfpathlineto{\pgfqpoint{0.993314in}{0.720565in}}%
\pgfpathclose%
\pgfusepath{fill}%
\end{pgfscope}%
\begin{pgfscope}%
\pgfpathrectangle{\pgfqpoint{0.407592in}{0.222904in}}{\pgfqpoint{1.411163in}{1.411163in}}%
\pgfusepath{clip}%
\pgfsetbuttcap%
\pgfsetroundjoin%
\definecolor{currentfill}{rgb}{0.749020,0.772549,0.792157}%
\pgfsetfillcolor{currentfill}%
\pgfsetfillopacity{0.160000}%
\pgfsetlinewidth{0.000000pt}%
\definecolor{currentstroke}{rgb}{0.000000,0.000000,0.000000}%
\pgfsetstrokecolor{currentstroke}%
\pgfsetdash{}{0pt}%
\pgfpathmoveto{\pgfqpoint{0.379152in}{1.212375in}}%
\pgfpathlineto{\pgfqpoint{0.364968in}{1.224381in}}%
\pgfpathlineto{\pgfqpoint{0.350763in}{1.236403in}}%
\pgfpathlineto{\pgfqpoint{0.346879in}{1.218622in}}%
\pgfpathlineto{\pgfqpoint{0.346239in}{1.200731in}}%
\pgfpathlineto{\pgfqpoint{0.360700in}{1.189260in}}%
\pgfpathlineto{\pgfqpoint{0.375133in}{1.177810in}}%
\pgfpathlineto{\pgfqpoint{0.375583in}{1.195141in}}%
\pgfpathlineto{\pgfqpoint{0.379152in}{1.212375in}}%
\pgfpathclose%
\pgfusepath{fill}%
\end{pgfscope}%
\begin{pgfscope}%
\pgfpathrectangle{\pgfqpoint{0.407592in}{0.222904in}}{\pgfqpoint{1.411163in}{1.411163in}}%
\pgfusepath{clip}%
\pgfsetbuttcap%
\pgfsetroundjoin%
\definecolor{currentfill}{rgb}{0.749020,0.772549,0.792157}%
\pgfsetfillcolor{currentfill}%
\pgfsetfillopacity{0.160000}%
\pgfsetlinewidth{0.000000pt}%
\definecolor{currentstroke}{rgb}{0.000000,0.000000,0.000000}%
\pgfsetstrokecolor{currentstroke}%
\pgfsetdash{}{0pt}%
\pgfpathmoveto{\pgfqpoint{1.639161in}{1.038634in}}%
\pgfpathlineto{\pgfqpoint{1.652867in}{1.049323in}}%
\pgfpathlineto{\pgfqpoint{1.666604in}{1.060037in}}%
\pgfpathlineto{\pgfqpoint{1.675139in}{1.074978in}}%
\pgfpathlineto{\pgfqpoint{1.681003in}{1.089985in}}%
\pgfpathlineto{\pgfqpoint{1.667054in}{1.078665in}}%
\pgfpathlineto{\pgfqpoint{1.653131in}{1.067365in}}%
\pgfpathlineto{\pgfqpoint{1.647416in}{1.052965in}}%
\pgfpathlineto{\pgfqpoint{1.639161in}{1.038634in}}%
\pgfpathclose%
\pgfusepath{fill}%
\end{pgfscope}%
\begin{pgfscope}%
\pgfpathrectangle{\pgfqpoint{0.407592in}{0.222904in}}{\pgfqpoint{1.411163in}{1.411163in}}%
\pgfusepath{clip}%
\pgfsetbuttcap%
\pgfsetroundjoin%
\definecolor{currentfill}{rgb}{0.749020,0.772549,0.792157}%
\pgfsetfillcolor{currentfill}%
\pgfsetfillopacity{0.160000}%
\pgfsetlinewidth{0.000000pt}%
\definecolor{currentstroke}{rgb}{0.000000,0.000000,0.000000}%
\pgfsetstrokecolor{currentstroke}%
\pgfsetdash{}{0pt}%
\pgfpathmoveto{\pgfqpoint{1.821891in}{1.204326in}}%
\pgfpathlineto{\pgfqpoint{1.836121in}{1.215875in}}%
\pgfpathlineto{\pgfqpoint{1.850378in}{1.227445in}}%
\pgfpathlineto{\pgfqpoint{1.853439in}{1.245867in}}%
\pgfpathlineto{\pgfqpoint{1.853116in}{1.264153in}}%
\pgfpathlineto{\pgfqpoint{1.838982in}{1.252044in}}%
\pgfpathlineto{\pgfqpoint{1.824868in}{1.239951in}}%
\pgfpathlineto{\pgfqpoint{1.825009in}{1.222200in}}%
\pgfpathlineto{\pgfqpoint{1.821891in}{1.204326in}}%
\pgfpathclose%
\pgfusepath{fill}%
\end{pgfscope}%
\begin{pgfscope}%
\pgfpathrectangle{\pgfqpoint{0.407592in}{0.222904in}}{\pgfqpoint{1.411163in}{1.411163in}}%
\pgfusepath{clip}%
\pgfsetbuttcap%
\pgfsetroundjoin%
\definecolor{currentfill}{rgb}{0.749020,0.772549,0.792157}%
\pgfsetfillcolor{currentfill}%
\pgfsetfillopacity{0.160000}%
\pgfsetlinewidth{0.000000pt}%
\definecolor{currentstroke}{rgb}{0.000000,0.000000,0.000000}%
\pgfsetstrokecolor{currentstroke}%
\pgfsetdash{}{0pt}%
\pgfpathmoveto{\pgfqpoint{1.121116in}{0.715477in}}%
\pgfpathlineto{\pgfqpoint{1.122908in}{0.721566in}}%
\pgfpathlineto{\pgfqpoint{1.124709in}{0.727683in}}%
\pgfpathlineto{\pgfqpoint{1.142681in}{0.729261in}}%
\pgfpathlineto{\pgfqpoint{1.160396in}{0.731335in}}%
\pgfpathlineto{\pgfqpoint{1.156735in}{0.725033in}}%
\pgfpathlineto{\pgfqpoint{1.153092in}{0.718760in}}%
\pgfpathlineto{\pgfqpoint{1.137218in}{0.716895in}}%
\pgfpathlineto{\pgfqpoint{1.121116in}{0.715477in}}%
\pgfpathclose%
\pgfusepath{fill}%
\end{pgfscope}%
\begin{pgfscope}%
\pgfpathrectangle{\pgfqpoint{0.407592in}{0.222904in}}{\pgfqpoint{1.411163in}{1.411163in}}%
\pgfusepath{clip}%
\pgfsetbuttcap%
\pgfsetroundjoin%
\definecolor{currentfill}{rgb}{0.749020,0.772549,0.792157}%
\pgfsetfillcolor{currentfill}%
\pgfsetfillopacity{0.160000}%
\pgfsetlinewidth{0.000000pt}%
\definecolor{currentstroke}{rgb}{0.000000,0.000000,0.000000}%
\pgfsetstrokecolor{currentstroke}%
\pgfsetdash{}{0pt}%
\pgfpathmoveto{\pgfqpoint{0.637842in}{0.951627in}}%
\pgfpathlineto{\pgfqpoint{0.624003in}{0.962061in}}%
\pgfpathlineto{\pgfqpoint{0.610133in}{0.972519in}}%
\pgfpathlineto{\pgfqpoint{0.617234in}{0.960232in}}%
\pgfpathlineto{\pgfqpoint{0.626505in}{0.948151in}}%
\pgfpathlineto{\pgfqpoint{0.640001in}{0.938325in}}%
\pgfpathlineto{\pgfqpoint{0.653460in}{0.928526in}}%
\pgfpathlineto{\pgfqpoint{0.644631in}{0.939976in}}%
\pgfpathlineto{\pgfqpoint{0.637842in}{0.951627in}}%
\pgfpathclose%
\pgfusepath{fill}%
\end{pgfscope}%
\begin{pgfscope}%
\pgfpathrectangle{\pgfqpoint{0.407592in}{0.222904in}}{\pgfqpoint{1.411163in}{1.411163in}}%
\pgfusepath{clip}%
\pgfsetbuttcap%
\pgfsetroundjoin%
\definecolor{currentfill}{rgb}{0.749020,0.772549,0.792157}%
\pgfsetfillcolor{currentfill}%
\pgfsetfillopacity{0.160000}%
\pgfsetlinewidth{0.000000pt}%
\definecolor{currentstroke}{rgb}{0.000000,0.000000,0.000000}%
\pgfsetstrokecolor{currentstroke}%
\pgfsetdash{}{0pt}%
\pgfpathmoveto{\pgfqpoint{0.920274in}{0.747773in}}%
\pgfpathlineto{\pgfqpoint{0.911412in}{0.754929in}}%
\pgfpathlineto{\pgfqpoint{0.902514in}{0.762114in}}%
\pgfpathlineto{\pgfqpoint{0.918418in}{0.757428in}}%
\pgfpathlineto{\pgfqpoint{0.935136in}{0.753205in}}%
\pgfpathlineto{\pgfqpoint{0.942506in}{0.746425in}}%
\pgfpathlineto{\pgfqpoint{0.949845in}{0.739675in}}%
\pgfpathlineto{\pgfqpoint{0.934693in}{0.743513in}}%
\pgfpathlineto{\pgfqpoint{0.920274in}{0.747773in}}%
\pgfpathclose%
\pgfusepath{fill}%
\end{pgfscope}%
\begin{pgfscope}%
\pgfpathrectangle{\pgfqpoint{0.407592in}{0.222904in}}{\pgfqpoint{1.411163in}{1.411163in}}%
\pgfusepath{clip}%
\pgfsetbuttcap%
\pgfsetroundjoin%
\definecolor{currentfill}{rgb}{0.749020,0.772549,0.792157}%
\pgfsetfillcolor{currentfill}%
\pgfsetfillopacity{0.160000}%
\pgfsetlinewidth{0.000000pt}%
\definecolor{currentstroke}{rgb}{0.000000,0.000000,0.000000}%
\pgfsetstrokecolor{currentstroke}%
\pgfsetdash{}{0pt}%
\pgfpathmoveto{\pgfqpoint{1.024018in}{0.716575in}}%
\pgfpathlineto{\pgfqpoint{1.020177in}{0.722725in}}%
\pgfpathlineto{\pgfqpoint{1.016319in}{0.728905in}}%
\pgfpathlineto{\pgfqpoint{1.034013in}{0.727427in}}%
\pgfpathlineto{\pgfqpoint{1.051971in}{0.726455in}}%
\pgfpathlineto{\pgfqpoint{1.053970in}{0.720400in}}%
\pgfpathlineto{\pgfqpoint{1.055959in}{0.714373in}}%
\pgfpathlineto{\pgfqpoint{1.039871in}{0.715247in}}%
\pgfpathlineto{\pgfqpoint{1.024018in}{0.716575in}}%
\pgfpathclose%
\pgfusepath{fill}%
\end{pgfscope}%
\begin{pgfscope}%
\pgfpathrectangle{\pgfqpoint{0.407592in}{0.222904in}}{\pgfqpoint{1.411163in}{1.411163in}}%
\pgfusepath{clip}%
\pgfsetbuttcap%
\pgfsetroundjoin%
\definecolor{currentfill}{rgb}{0.749020,0.772549,0.792157}%
\pgfsetfillcolor{currentfill}%
\pgfsetfillopacity{0.160000}%
\pgfsetlinewidth{0.000000pt}%
\definecolor{currentstroke}{rgb}{0.000000,0.000000,0.000000}%
\pgfsetstrokecolor{currentstroke}%
\pgfsetdash{}{0pt}%
\pgfpathmoveto{\pgfqpoint{0.778258in}{0.834014in}}%
\pgfpathlineto{\pgfqpoint{0.765756in}{0.842970in}}%
\pgfpathlineto{\pgfqpoint{0.753215in}{0.851954in}}%
\pgfpathlineto{\pgfqpoint{0.765467in}{0.843236in}}%
\pgfpathlineto{\pgfqpoint{0.779242in}{0.834877in}}%
\pgfpathlineto{\pgfqpoint{0.790871in}{0.826488in}}%
\pgfpathlineto{\pgfqpoint{0.802460in}{0.818128in}}%
\pgfpathlineto{\pgfqpoint{0.789658in}{0.825903in}}%
\pgfpathlineto{\pgfqpoint{0.778258in}{0.834014in}}%
\pgfpathclose%
\pgfusepath{fill}%
\end{pgfscope}%
\begin{pgfscope}%
\pgfpathrectangle{\pgfqpoint{0.407592in}{0.222904in}}{\pgfqpoint{1.411163in}{1.411163in}}%
\pgfusepath{clip}%
\pgfsetbuttcap%
\pgfsetroundjoin%
\definecolor{currentfill}{rgb}{0.749020,0.772549,0.792157}%
\pgfsetfillcolor{currentfill}%
\pgfsetfillopacity{0.160000}%
\pgfsetlinewidth{0.000000pt}%
\definecolor{currentstroke}{rgb}{0.000000,0.000000,0.000000}%
\pgfsetstrokecolor{currentstroke}%
\pgfsetdash{}{0pt}%
\pgfpathmoveto{\pgfqpoint{1.434292in}{0.871917in}}%
\pgfpathlineto{\pgfqpoint{1.446309in}{0.881046in}}%
\pgfpathlineto{\pgfqpoint{1.458363in}{0.890204in}}%
\pgfpathlineto{\pgfqpoint{1.472457in}{0.900538in}}%
\pgfpathlineto{\pgfqpoint{1.484750in}{0.911159in}}%
\pgfpathlineto{\pgfqpoint{1.471911in}{0.901376in}}%
\pgfpathlineto{\pgfqpoint{1.459107in}{0.891619in}}%
\pgfpathlineto{\pgfqpoint{1.447537in}{0.881631in}}%
\pgfpathlineto{\pgfqpoint{1.434292in}{0.871917in}}%
\pgfpathclose%
\pgfusepath{fill}%
\end{pgfscope}%
\begin{pgfscope}%
\pgfpathrectangle{\pgfqpoint{0.407592in}{0.222904in}}{\pgfqpoint{1.411163in}{1.411163in}}%
\pgfusepath{clip}%
\pgfsetbuttcap%
\pgfsetroundjoin%
\definecolor{currentfill}{rgb}{0.749020,0.772549,0.792157}%
\pgfsetfillcolor{currentfill}%
\pgfsetfillopacity{0.160000}%
\pgfsetlinewidth{0.000000pt}%
\definecolor{currentstroke}{rgb}{0.000000,0.000000,0.000000}%
\pgfsetstrokecolor{currentstroke}%
\pgfsetdash{}{0pt}%
\pgfpathmoveto{\pgfqpoint{1.088532in}{0.714004in}}%
\pgfpathlineto{\pgfqpoint{1.088433in}{0.720011in}}%
\pgfpathlineto{\pgfqpoint{1.088334in}{0.726045in}}%
\pgfpathlineto{\pgfqpoint{1.106564in}{0.726610in}}%
\pgfpathlineto{\pgfqpoint{1.124709in}{0.727683in}}%
\pgfpathlineto{\pgfqpoint{1.122908in}{0.721566in}}%
\pgfpathlineto{\pgfqpoint{1.121116in}{0.715477in}}%
\pgfpathlineto{\pgfqpoint{1.104862in}{0.714512in}}%
\pgfpathlineto{\pgfqpoint{1.088532in}{0.714004in}}%
\pgfpathclose%
\pgfusepath{fill}%
\end{pgfscope}%
\begin{pgfscope}%
\pgfpathrectangle{\pgfqpoint{0.407592in}{0.222904in}}{\pgfqpoint{1.411163in}{1.411163in}}%
\pgfusepath{clip}%
\pgfsetbuttcap%
\pgfsetroundjoin%
\definecolor{currentfill}{rgb}{0.749020,0.772549,0.792157}%
\pgfsetfillcolor{currentfill}%
\pgfsetfillopacity{0.160000}%
\pgfsetlinewidth{0.000000pt}%
\definecolor{currentstroke}{rgb}{0.000000,0.000000,0.000000}%
\pgfsetstrokecolor{currentstroke}%
\pgfsetdash{}{0pt}%
\pgfpathmoveto{\pgfqpoint{0.518002in}{1.064474in}}%
\pgfpathlineto{\pgfqpoint{0.503834in}{1.075713in}}%
\pgfpathlineto{\pgfqpoint{0.489640in}{1.086973in}}%
\pgfpathlineto{\pgfqpoint{0.492483in}{1.071973in}}%
\pgfpathlineto{\pgfqpoint{0.498016in}{1.057053in}}%
\pgfpathlineto{\pgfqpoint{0.512144in}{1.046401in}}%
\pgfpathlineto{\pgfqpoint{0.526239in}{1.035773in}}%
\pgfpathlineto{\pgfqpoint{0.520840in}{1.050082in}}%
\pgfpathlineto{\pgfqpoint{0.518002in}{1.064474in}}%
\pgfpathclose%
\pgfusepath{fill}%
\end{pgfscope}%
\begin{pgfscope}%
\pgfpathrectangle{\pgfqpoint{0.407592in}{0.222904in}}{\pgfqpoint{1.411163in}{1.411163in}}%
\pgfusepath{clip}%
\pgfsetbuttcap%
\pgfsetroundjoin%
\definecolor{currentfill}{rgb}{0.749020,0.772549,0.792157}%
\pgfsetfillcolor{currentfill}%
\pgfsetfillopacity{0.160000}%
\pgfsetlinewidth{0.000000pt}%
\definecolor{currentstroke}{rgb}{0.000000,0.000000,0.000000}%
\pgfsetstrokecolor{currentstroke}%
\pgfsetdash{}{0pt}%
\pgfpathmoveto{\pgfqpoint{1.055959in}{0.714373in}}%
\pgfpathlineto{\pgfqpoint{1.053970in}{0.720400in}}%
\pgfpathlineto{\pgfqpoint{1.051971in}{0.726455in}}%
\pgfpathlineto{\pgfqpoint{1.070107in}{0.725994in}}%
\pgfpathlineto{\pgfqpoint{1.088334in}{0.726045in}}%
\pgfpathlineto{\pgfqpoint{1.088433in}{0.720011in}}%
\pgfpathlineto{\pgfqpoint{1.088532in}{0.714004in}}%
\pgfpathlineto{\pgfqpoint{1.072205in}{0.713958in}}%
\pgfpathlineto{\pgfqpoint{1.055959in}{0.714373in}}%
\pgfpathclose%
\pgfusepath{fill}%
\end{pgfscope}%
\begin{pgfscope}%
\pgfpathrectangle{\pgfqpoint{0.407592in}{0.222904in}}{\pgfqpoint{1.411163in}{1.411163in}}%
\pgfusepath{clip}%
\pgfsetbuttcap%
\pgfsetroundjoin%
\definecolor{currentfill}{rgb}{0.749020,0.772549,0.792157}%
\pgfsetfillcolor{currentfill}%
\pgfsetfillopacity{0.160000}%
\pgfsetlinewidth{0.000000pt}%
\definecolor{currentstroke}{rgb}{0.000000,0.000000,0.000000}%
\pgfsetstrokecolor{currentstroke}%
\pgfsetdash{}{0pt}%
\pgfpathmoveto{\pgfqpoint{1.850378in}{1.227445in}}%
\pgfpathlineto{\pgfqpoint{1.864660in}{1.239036in}}%
\pgfpathlineto{\pgfqpoint{1.867689in}{1.257729in}}%
\pgfpathlineto{\pgfqpoint{1.867269in}{1.276279in}}%
\pgfpathlineto{\pgfqpoint{1.853116in}{1.264153in}}%
\pgfpathlineto{\pgfqpoint{1.853439in}{1.245867in}}%
\pgfusepath{fill}%
\end{pgfscope}%
\begin{pgfscope}%
\pgfpathrectangle{\pgfqpoint{0.407592in}{0.222904in}}{\pgfqpoint{1.411163in}{1.411163in}}%
\pgfusepath{clip}%
\pgfsetbuttcap%
\pgfsetroundjoin%
\definecolor{currentfill}{rgb}{0.749020,0.772549,0.792157}%
\pgfsetfillcolor{currentfill}%
\pgfsetfillopacity{0.160000}%
\pgfsetlinewidth{0.000000pt}%
\definecolor{currentstroke}{rgb}{0.000000,0.000000,0.000000}%
\pgfsetstrokecolor{currentstroke}%
\pgfsetdash{}{0pt}%
\pgfpathmoveto{\pgfqpoint{0.350763in}{1.236403in}}%
\pgfpathlineto{\pgfqpoint{0.336539in}{1.248443in}}%
\pgfpathlineto{\pgfqpoint{0.322294in}{1.260499in}}%
\pgfpathlineto{\pgfqpoint{0.318082in}{1.242180in}}%
\pgfpathlineto{\pgfqpoint{0.317235in}{1.223739in}}%
\pgfpathlineto{\pgfqpoint{0.331751in}{1.212225in}}%
\pgfpathlineto{\pgfqpoint{0.346239in}{1.200731in}}%
\pgfpathlineto{\pgfqpoint{0.346879in}{1.218622in}}%
\pgfpathlineto{\pgfqpoint{0.350763in}{1.236403in}}%
\pgfpathclose%
\pgfusepath{fill}%
\end{pgfscope}%
\begin{pgfscope}%
\pgfpathrectangle{\pgfqpoint{0.407592in}{0.222904in}}{\pgfqpoint{1.411163in}{1.411163in}}%
\pgfusepath{clip}%
\pgfsetbuttcap%
\pgfsetroundjoin%
\definecolor{currentfill}{rgb}{0.749020,0.772549,0.792157}%
\pgfsetfillcolor{currentfill}%
\pgfsetfillopacity{0.160000}%
\pgfsetlinewidth{0.000000pt}%
\definecolor{currentstroke}{rgb}{0.000000,0.000000,0.000000}%
\pgfsetstrokecolor{currentstroke}%
\pgfsetdash{}{0pt}%
\pgfpathmoveto{\pgfqpoint{1.194714in}{0.736924in}}%
\pgfpathlineto{\pgfqpoint{1.200181in}{0.743536in}}%
\pgfpathlineto{\pgfqpoint{1.205671in}{0.750178in}}%
\pgfpathlineto{\pgfqpoint{1.223823in}{0.754014in}}%
\pgfpathlineto{\pgfqpoint{1.241326in}{0.758332in}}%
\pgfpathlineto{\pgfqpoint{1.234156in}{0.751318in}}%
\pgfpathlineto{\pgfqpoint{1.227016in}{0.744334in}}%
\pgfpathlineto{\pgfqpoint{1.211156in}{0.740410in}}%
\pgfpathlineto{\pgfqpoint{1.194714in}{0.736924in}}%
\pgfpathclose%
\pgfusepath{fill}%
\end{pgfscope}%
\begin{pgfscope}%
\pgfpathrectangle{\pgfqpoint{0.407592in}{0.222904in}}{\pgfqpoint{1.411163in}{1.411163in}}%
\pgfusepath{clip}%
\pgfsetbuttcap%
\pgfsetroundjoin%
\definecolor{currentfill}{rgb}{0.749020,0.772549,0.792157}%
\pgfsetfillcolor{currentfill}%
\pgfsetfillopacity{0.160000}%
\pgfsetlinewidth{0.000000pt}%
\definecolor{currentstroke}{rgb}{0.000000,0.000000,0.000000}%
\pgfsetstrokecolor{currentstroke}%
\pgfsetdash{}{0pt}%
\pgfpathmoveto{\pgfqpoint{0.706932in}{0.889595in}}%
\pgfpathlineto{\pgfqpoint{0.693618in}{0.899288in}}%
\pgfpathlineto{\pgfqpoint{0.680268in}{0.909008in}}%
\pgfpathlineto{\pgfqpoint{0.690534in}{0.898440in}}%
\pgfpathlineto{\pgfqpoint{0.702659in}{0.888170in}}%
\pgfpathlineto{\pgfqpoint{0.715357in}{0.879073in}}%
\pgfpathlineto{\pgfqpoint{0.728016in}{0.870005in}}%
\pgfpathlineto{\pgfqpoint{0.716608in}{0.879658in}}%
\pgfpathlineto{\pgfqpoint{0.706932in}{0.889595in}}%
\pgfpathclose%
\pgfusepath{fill}%
\end{pgfscope}%
\begin{pgfscope}%
\pgfpathrectangle{\pgfqpoint{0.407592in}{0.222904in}}{\pgfqpoint{1.411163in}{1.411163in}}%
\pgfusepath{clip}%
\pgfsetbuttcap%
\pgfsetroundjoin%
\definecolor{currentfill}{rgb}{0.749020,0.772549,0.792157}%
\pgfsetfillcolor{currentfill}%
\pgfsetfillopacity{0.160000}%
\pgfsetlinewidth{0.000000pt}%
\definecolor{currentstroke}{rgb}{0.000000,0.000000,0.000000}%
\pgfsetstrokecolor{currentstroke}%
\pgfsetdash{}{0pt}%
\pgfpathmoveto{\pgfqpoint{0.873435in}{0.772777in}}%
\pgfpathlineto{\pgfqpoint{0.863150in}{0.780474in}}%
\pgfpathlineto{\pgfqpoint{0.852825in}{0.788200in}}%
\pgfpathlineto{\pgfqpoint{0.868190in}{0.782163in}}%
\pgfpathlineto{\pgfqpoint{0.884607in}{0.776574in}}%
\pgfpathlineto{\pgfqpoint{0.893579in}{0.769329in}}%
\pgfpathlineto{\pgfqpoint{0.902514in}{0.762114in}}%
\pgfpathlineto{\pgfqpoint{0.887497in}{0.767239in}}%
\pgfpathlineto{\pgfqpoint{0.873435in}{0.772777in}}%
\pgfpathclose%
\pgfusepath{fill}%
\end{pgfscope}%
\begin{pgfscope}%
\pgfpathrectangle{\pgfqpoint{0.407592in}{0.222904in}}{\pgfqpoint{1.411163in}{1.411163in}}%
\pgfusepath{clip}%
\pgfsetbuttcap%
\pgfsetroundjoin%
\definecolor{currentfill}{rgb}{0.749020,0.772549,0.792157}%
\pgfsetfillcolor{currentfill}%
\pgfsetfillopacity{0.160000}%
\pgfsetlinewidth{0.000000pt}%
\definecolor{currentstroke}{rgb}{0.000000,0.000000,0.000000}%
\pgfsetstrokecolor{currentstroke}%
\pgfsetdash{}{0pt}%
\pgfpathmoveto{\pgfqpoint{0.825518in}{0.801494in}}%
\pgfpathlineto{\pgfqpoint{0.814009in}{0.809797in}}%
\pgfpathlineto{\pgfqpoint{0.802460in}{0.818128in}}%
\pgfpathlineto{\pgfqpoint{0.816613in}{0.810728in}}%
\pgfpathlineto{\pgfqpoint{0.832058in}{0.803742in}}%
\pgfpathlineto{\pgfqpoint{0.842461in}{0.795956in}}%
\pgfpathlineto{\pgfqpoint{0.852825in}{0.788200in}}%
\pgfpathlineto{\pgfqpoint{0.838581in}{0.794655in}}%
\pgfpathlineto{\pgfqpoint{0.825518in}{0.801494in}}%
\pgfpathclose%
\pgfusepath{fill}%
\end{pgfscope}%
\begin{pgfscope}%
\pgfpathrectangle{\pgfqpoint{0.407592in}{0.222904in}}{\pgfqpoint{1.411163in}{1.411163in}}%
\pgfusepath{clip}%
\pgfsetbuttcap%
\pgfsetroundjoin%
\definecolor{currentfill}{rgb}{0.749020,0.772549,0.792157}%
\pgfsetfillcolor{currentfill}%
\pgfsetfillopacity{0.160000}%
\pgfsetlinewidth{0.000000pt}%
\definecolor{currentstroke}{rgb}{0.000000,0.000000,0.000000}%
\pgfsetstrokecolor{currentstroke}%
\pgfsetdash{}{0pt}%
\pgfpathmoveto{\pgfqpoint{1.536450in}{0.950555in}}%
\pgfpathlineto{\pgfqpoint{1.549462in}{0.960470in}}%
\pgfpathlineto{\pgfqpoint{1.562509in}{0.970412in}}%
\pgfpathlineto{\pgfqpoint{1.574706in}{0.983166in}}%
\pgfpathlineto{\pgfqpoint{1.584645in}{0.996118in}}%
\pgfpathlineto{\pgfqpoint{1.571092in}{0.985549in}}%
\pgfpathlineto{\pgfqpoint{1.557570in}{0.975003in}}%
\pgfpathlineto{\pgfqpoint{1.548074in}{0.962682in}}%
\pgfpathlineto{\pgfqpoint{1.536450in}{0.950555in}}%
\pgfpathclose%
\pgfusepath{fill}%
\end{pgfscope}%
\begin{pgfscope}%
\pgfpathrectangle{\pgfqpoint{0.407592in}{0.222904in}}{\pgfqpoint{1.411163in}{1.411163in}}%
\pgfusepath{clip}%
\pgfsetbuttcap%
\pgfsetroundjoin%
\definecolor{currentfill}{rgb}{0.749020,0.772549,0.792157}%
\pgfsetfillcolor{currentfill}%
\pgfsetfillopacity{0.160000}%
\pgfsetlinewidth{0.000000pt}%
\definecolor{currentstroke}{rgb}{0.000000,0.000000,0.000000}%
\pgfsetstrokecolor{currentstroke}%
\pgfsetdash{}{0pt}%
\pgfpathmoveto{\pgfqpoint{1.666604in}{1.060037in}}%
\pgfpathlineto{\pgfqpoint{1.680372in}{1.070774in}}%
\pgfpathlineto{\pgfqpoint{1.694172in}{1.081536in}}%
\pgfpathlineto{\pgfqpoint{1.702976in}{1.097082in}}%
\pgfpathlineto{\pgfqpoint{1.708976in}{1.112688in}}%
\pgfpathlineto{\pgfqpoint{1.694977in}{1.101326in}}%
\pgfpathlineto{\pgfqpoint{1.681003in}{1.089985in}}%
\pgfpathlineto{\pgfqpoint{1.675139in}{1.074978in}}%
\pgfpathlineto{\pgfqpoint{1.666604in}{1.060037in}}%
\pgfpathclose%
\pgfusepath{fill}%
\end{pgfscope}%
\begin{pgfscope}%
\pgfpathrectangle{\pgfqpoint{0.407592in}{0.222904in}}{\pgfqpoint{1.411163in}{1.411163in}}%
\pgfusepath{clip}%
\pgfsetbuttcap%
\pgfsetroundjoin%
\definecolor{currentfill}{rgb}{0.749020,0.772549,0.792157}%
\pgfsetfillcolor{currentfill}%
\pgfsetfillopacity{0.160000}%
\pgfsetlinewidth{0.000000pt}%
\definecolor{currentstroke}{rgb}{0.000000,0.000000,0.000000}%
\pgfsetstrokecolor{currentstroke}%
\pgfsetdash{}{0pt}%
\pgfpathmoveto{\pgfqpoint{1.396468in}{0.844958in}}%
\pgfpathlineto{\pgfqpoint{1.407924in}{0.853728in}}%
\pgfpathlineto{\pgfqpoint{1.419417in}{0.862526in}}%
\pgfpathlineto{\pgfqpoint{1.434292in}{0.871917in}}%
\pgfpathlineto{\pgfqpoint{1.422311in}{0.862815in}}%
\pgfpathlineto{\pgfqpoint{1.410368in}{0.853742in}}%
\pgfusepath{fill}%
\end{pgfscope}%
\begin{pgfscope}%
\pgfpathrectangle{\pgfqpoint{0.407592in}{0.222904in}}{\pgfqpoint{1.411163in}{1.411163in}}%
\pgfusepath{clip}%
\pgfsetbuttcap%
\pgfsetroundjoin%
\definecolor{currentfill}{rgb}{0.749020,0.772549,0.792157}%
\pgfsetfillcolor{currentfill}%
\pgfsetfillopacity{0.160000}%
\pgfsetlinewidth{0.000000pt}%
\definecolor{currentstroke}{rgb}{0.000000,0.000000,0.000000}%
\pgfsetstrokecolor{currentstroke}%
\pgfsetdash{}{0pt}%
\pgfpathmoveto{\pgfqpoint{0.322294in}{1.260499in}}%
\pgfpathlineto{\pgfqpoint{0.308030in}{1.272572in}}%
\pgfpathlineto{\pgfqpoint{0.303647in}{1.253987in}}%
\pgfpathlineto{\pgfqpoint{0.302693in}{1.235276in}}%
\pgfpathlineto{\pgfqpoint{0.317235in}{1.223739in}}%
\pgfpathlineto{\pgfqpoint{0.318082in}{1.242180in}}%
\pgfusepath{fill}%
\end{pgfscope}%
\begin{pgfscope}%
\pgfpathrectangle{\pgfqpoint{0.407592in}{0.222904in}}{\pgfqpoint{1.411163in}{1.411163in}}%
\pgfusepath{clip}%
\pgfsetbuttcap%
\pgfsetroundjoin%
\definecolor{currentfill}{rgb}{0.749020,0.772549,0.792157}%
\pgfsetfillcolor{currentfill}%
\pgfsetfillopacity{0.160000}%
\pgfsetlinewidth{0.000000pt}%
\definecolor{currentstroke}{rgb}{0.000000,0.000000,0.000000}%
\pgfsetstrokecolor{currentstroke}%
\pgfsetdash{}{0pt}%
\pgfpathmoveto{\pgfqpoint{0.949845in}{0.739675in}}%
\pgfpathlineto{\pgfqpoint{0.942506in}{0.746425in}}%
\pgfpathlineto{\pgfqpoint{0.935136in}{0.753205in}}%
\pgfpathlineto{\pgfqpoint{0.952592in}{0.749468in}}%
\pgfpathlineto{\pgfqpoint{0.970704in}{0.746235in}}%
\pgfpathlineto{\pgfqpoint{0.976395in}{0.739774in}}%
\pgfpathlineto{\pgfqpoint{0.982060in}{0.733342in}}%
\pgfpathlineto{\pgfqpoint{0.965658in}{0.736279in}}%
\pgfpathlineto{\pgfqpoint{0.949845in}{0.739675in}}%
\pgfpathclose%
\pgfusepath{fill}%
\end{pgfscope}%
\begin{pgfscope}%
\pgfpathrectangle{\pgfqpoint{0.407592in}{0.222904in}}{\pgfqpoint{1.411163in}{1.411163in}}%
\pgfusepath{clip}%
\pgfsetbuttcap%
\pgfsetroundjoin%
\definecolor{currentfill}{rgb}{0.749020,0.772549,0.792157}%
\pgfsetfillcolor{currentfill}%
\pgfsetfillopacity{0.160000}%
\pgfsetlinewidth{0.000000pt}%
\definecolor{currentstroke}{rgb}{0.000000,0.000000,0.000000}%
\pgfsetstrokecolor{currentstroke}%
\pgfsetdash{}{0pt}%
\pgfpathmoveto{\pgfqpoint{0.610133in}{0.972519in}}%
\pgfpathlineto{\pgfqpoint{0.596231in}{0.983001in}}%
\pgfpathlineto{\pgfqpoint{0.582297in}{0.993507in}}%
\pgfpathlineto{\pgfqpoint{0.589699in}{0.980589in}}%
\pgfpathlineto{\pgfqpoint{0.599402in}{0.967883in}}%
\pgfpathlineto{\pgfqpoint{0.612972in}{0.958003in}}%
\pgfpathlineto{\pgfqpoint{0.626505in}{0.948151in}}%
\pgfpathlineto{\pgfqpoint{0.617234in}{0.960232in}}%
\pgfpathlineto{\pgfqpoint{0.610133in}{0.972519in}}%
\pgfpathclose%
\pgfusepath{fill}%
\end{pgfscope}%
\begin{pgfscope}%
\pgfpathrectangle{\pgfqpoint{0.407592in}{0.222904in}}{\pgfqpoint{1.411163in}{1.411163in}}%
\pgfusepath{clip}%
\pgfsetbuttcap%
\pgfsetroundjoin%
\definecolor{currentfill}{rgb}{0.749020,0.772549,0.792157}%
\pgfsetfillcolor{currentfill}%
\pgfsetfillopacity{0.160000}%
\pgfsetlinewidth{0.000000pt}%
\definecolor{currentstroke}{rgb}{0.000000,0.000000,0.000000}%
\pgfsetstrokecolor{currentstroke}%
\pgfsetdash{}{0pt}%
\pgfpathmoveto{\pgfqpoint{1.343864in}{0.812181in}}%
\pgfpathlineto{\pgfqpoint{1.354097in}{0.820316in}}%
\pgfpathlineto{\pgfqpoint{1.364368in}{0.828482in}}%
\pgfpathlineto{\pgfqpoint{1.381114in}{0.836524in}}%
\pgfpathlineto{\pgfqpoint{1.396468in}{0.844958in}}%
\pgfpathlineto{\pgfqpoint{1.385049in}{0.836217in}}%
\pgfpathlineto{\pgfqpoint{1.373668in}{0.827504in}}%
\pgfpathlineto{\pgfqpoint{1.359406in}{0.819659in}}%
\pgfpathlineto{\pgfqpoint{1.343864in}{0.812181in}}%
\pgfpathclose%
\pgfusepath{fill}%
\end{pgfscope}%
\begin{pgfscope}%
\pgfpathrectangle{\pgfqpoint{0.407592in}{0.222904in}}{\pgfqpoint{1.411163in}{1.411163in}}%
\pgfusepath{clip}%
\pgfsetbuttcap%
\pgfsetroundjoin%
\definecolor{currentfill}{rgb}{0.749020,0.772549,0.792157}%
\pgfsetfillcolor{currentfill}%
\pgfsetfillopacity{0.160000}%
\pgfsetlinewidth{0.000000pt}%
\definecolor{currentstroke}{rgb}{0.000000,0.000000,0.000000}%
\pgfsetstrokecolor{currentstroke}%
\pgfsetdash{}{0pt}%
\pgfpathmoveto{\pgfqpoint{1.241326in}{0.758332in}}%
\pgfpathlineto{\pgfqpoint{1.248526in}{0.765375in}}%
\pgfpathlineto{\pgfqpoint{1.255757in}{0.772447in}}%
\pgfpathlineto{\pgfqpoint{1.274108in}{0.777657in}}%
\pgfpathlineto{\pgfqpoint{1.291569in}{0.783339in}}%
\pgfpathlineto{\pgfqpoint{1.282801in}{0.775814in}}%
\pgfpathlineto{\pgfqpoint{1.274068in}{0.768318in}}%
\pgfpathlineto{\pgfqpoint{1.258100in}{0.763108in}}%
\pgfpathlineto{\pgfqpoint{1.241326in}{0.758332in}}%
\pgfpathclose%
\pgfusepath{fill}%
\end{pgfscope}%
\begin{pgfscope}%
\pgfpathrectangle{\pgfqpoint{0.407592in}{0.222904in}}{\pgfqpoint{1.411163in}{1.411163in}}%
\pgfusepath{clip}%
\pgfsetbuttcap%
\pgfsetroundjoin%
\definecolor{currentfill}{rgb}{0.749020,0.772549,0.792157}%
\pgfsetfillcolor{currentfill}%
\pgfsetfillopacity{0.160000}%
\pgfsetlinewidth{0.000000pt}%
\definecolor{currentstroke}{rgb}{0.000000,0.000000,0.000000}%
\pgfsetstrokecolor{currentstroke}%
\pgfsetdash{}{0pt}%
\pgfpathmoveto{\pgfqpoint{0.489640in}{1.086973in}}%
\pgfpathlineto{\pgfqpoint{0.475420in}{1.098254in}}%
\pgfpathlineto{\pgfqpoint{0.461173in}{1.109555in}}%
\pgfpathlineto{\pgfqpoint{0.464007in}{1.093956in}}%
\pgfpathlineto{\pgfqpoint{0.469662in}{1.078431in}}%
\pgfpathlineto{\pgfqpoint{0.483856in}{1.067730in}}%
\pgfpathlineto{\pgfqpoint{0.498016in}{1.057053in}}%
\pgfpathlineto{\pgfqpoint{0.492483in}{1.071973in}}%
\pgfpathlineto{\pgfqpoint{0.489640in}{1.086973in}}%
\pgfpathclose%
\pgfusepath{fill}%
\end{pgfscope}%
\begin{pgfscope}%
\pgfpathrectangle{\pgfqpoint{0.407592in}{0.222904in}}{\pgfqpoint{1.411163in}{1.411163in}}%
\pgfusepath{clip}%
\pgfsetbuttcap%
\pgfsetroundjoin%
\definecolor{currentfill}{rgb}{0.749020,0.772549,0.792157}%
\pgfsetfillcolor{currentfill}%
\pgfsetfillopacity{0.160000}%
\pgfsetlinewidth{0.000000pt}%
\definecolor{currentstroke}{rgb}{0.000000,0.000000,0.000000}%
\pgfsetstrokecolor{currentstroke}%
\pgfsetdash{}{0pt}%
\pgfpathmoveto{\pgfqpoint{1.160396in}{0.731335in}}%
\pgfpathlineto{\pgfqpoint{1.164073in}{0.737666in}}%
\pgfpathlineto{\pgfqpoint{1.167767in}{0.744025in}}%
\pgfpathlineto{\pgfqpoint{1.186956in}{0.746842in}}%
\pgfpathlineto{\pgfqpoint{1.205671in}{0.750178in}}%
\pgfpathlineto{\pgfqpoint{1.200181in}{0.743536in}}%
\pgfpathlineto{\pgfqpoint{1.194714in}{0.736924in}}%
\pgfpathlineto{\pgfqpoint{1.177768in}{0.733894in}}%
\pgfpathlineto{\pgfqpoint{1.160396in}{0.731335in}}%
\pgfpathclose%
\pgfusepath{fill}%
\end{pgfscope}%
\begin{pgfscope}%
\pgfpathrectangle{\pgfqpoint{0.407592in}{0.222904in}}{\pgfqpoint{1.411163in}{1.411163in}}%
\pgfusepath{clip}%
\pgfsetbuttcap%
\pgfsetroundjoin%
\definecolor{currentfill}{rgb}{0.749020,0.772549,0.792157}%
\pgfsetfillcolor{currentfill}%
\pgfsetfillopacity{0.160000}%
\pgfsetlinewidth{0.000000pt}%
\definecolor{currentstroke}{rgb}{0.000000,0.000000,0.000000}%
\pgfsetstrokecolor{currentstroke}%
\pgfsetdash{}{0pt}%
\pgfpathmoveto{\pgfqpoint{1.291569in}{0.783339in}}%
\pgfpathlineto{\pgfqpoint{1.300372in}{0.790894in}}%
\pgfpathlineto{\pgfqpoint{1.309209in}{0.798479in}}%
\pgfpathlineto{\pgfqpoint{1.327108in}{0.805109in}}%
\pgfpathlineto{\pgfqpoint{1.343864in}{0.812181in}}%
\pgfpathlineto{\pgfqpoint{1.333668in}{0.804075in}}%
\pgfpathlineto{\pgfqpoint{1.323509in}{0.795998in}}%
\pgfpathlineto{\pgfqpoint{1.308060in}{0.789463in}}%
\pgfpathlineto{\pgfqpoint{1.291569in}{0.783339in}}%
\pgfpathclose%
\pgfusepath{fill}%
\end{pgfscope}%
\begin{pgfscope}%
\pgfpathrectangle{\pgfqpoint{0.407592in}{0.222904in}}{\pgfqpoint{1.411163in}{1.411163in}}%
\pgfusepath{clip}%
\pgfsetbuttcap%
\pgfsetroundjoin%
\definecolor{currentfill}{rgb}{0.749020,0.772549,0.792157}%
\pgfsetfillcolor{currentfill}%
\pgfsetfillopacity{0.160000}%
\pgfsetlinewidth{0.000000pt}%
\definecolor{currentstroke}{rgb}{0.000000,0.000000,0.000000}%
\pgfsetstrokecolor{currentstroke}%
\pgfsetdash{}{0pt}%
\pgfpathmoveto{\pgfqpoint{1.458363in}{0.890204in}}%
\pgfpathlineto{\pgfqpoint{1.470455in}{0.899390in}}%
\pgfpathlineto{\pgfqpoint{1.482584in}{0.908604in}}%
\pgfpathlineto{\pgfqpoint{1.497521in}{0.919554in}}%
\pgfpathlineto{\pgfqpoint{1.510531in}{0.930804in}}%
\pgfpathlineto{\pgfqpoint{1.497623in}{0.920969in}}%
\pgfpathlineto{\pgfqpoint{1.484750in}{0.911159in}}%
\pgfpathlineto{\pgfqpoint{1.472457in}{0.900538in}}%
\pgfpathlineto{\pgfqpoint{1.458363in}{0.890204in}}%
\pgfpathclose%
\pgfusepath{fill}%
\end{pgfscope}%
\begin{pgfscope}%
\pgfpathrectangle{\pgfqpoint{0.407592in}{0.222904in}}{\pgfqpoint{1.411163in}{1.411163in}}%
\pgfusepath{clip}%
\pgfsetbuttcap%
\pgfsetroundjoin%
\definecolor{currentfill}{rgb}{0.749020,0.772549,0.792157}%
\pgfsetfillcolor{currentfill}%
\pgfsetfillopacity{0.160000}%
\pgfsetlinewidth{0.000000pt}%
\definecolor{currentstroke}{rgb}{0.000000,0.000000,0.000000}%
\pgfsetstrokecolor{currentstroke}%
\pgfsetdash{}{0pt}%
\pgfpathmoveto{\pgfqpoint{0.753215in}{0.851954in}}%
\pgfpathlineto{\pgfqpoint{0.740635in}{0.860966in}}%
\pgfpathlineto{\pgfqpoint{0.728016in}{0.870005in}}%
\pgfpathlineto{\pgfqpoint{0.741116in}{0.860685in}}%
\pgfpathlineto{\pgfqpoint{0.755860in}{0.851745in}}%
\pgfpathlineto{\pgfqpoint{0.767572in}{0.843296in}}%
\pgfpathlineto{\pgfqpoint{0.779242in}{0.834877in}}%
\pgfpathlineto{\pgfqpoint{0.765467in}{0.843236in}}%
\pgfpathlineto{\pgfqpoint{0.753215in}{0.851954in}}%
\pgfpathclose%
\pgfusepath{fill}%
\end{pgfscope}%
\begin{pgfscope}%
\pgfpathrectangle{\pgfqpoint{0.407592in}{0.222904in}}{\pgfqpoint{1.411163in}{1.411163in}}%
\pgfusepath{clip}%
\pgfsetbuttcap%
\pgfsetroundjoin%
\definecolor{currentfill}{rgb}{0.749020,0.772549,0.792157}%
\pgfsetfillcolor{currentfill}%
\pgfsetfillopacity{0.160000}%
\pgfsetlinewidth{0.000000pt}%
\definecolor{currentstroke}{rgb}{0.000000,0.000000,0.000000}%
\pgfsetstrokecolor{currentstroke}%
\pgfsetdash{}{0pt}%
\pgfpathmoveto{\pgfqpoint{1.694172in}{1.081536in}}%
\pgfpathlineto{\pgfqpoint{1.708003in}{1.092323in}}%
\pgfpathlineto{\pgfqpoint{1.721866in}{1.103134in}}%
\pgfpathlineto{\pgfqpoint{1.730926in}{1.119275in}}%
\pgfpathlineto{\pgfqpoint{1.737051in}{1.135472in}}%
\pgfpathlineto{\pgfqpoint{1.723001in}{1.124070in}}%
\pgfpathlineto{\pgfqpoint{1.708976in}{1.112688in}}%
\pgfpathlineto{\pgfqpoint{1.702976in}{1.097082in}}%
\pgfpathlineto{\pgfqpoint{1.694172in}{1.081536in}}%
\pgfpathclose%
\pgfusepath{fill}%
\end{pgfscope}%
\begin{pgfscope}%
\pgfpathrectangle{\pgfqpoint{0.407592in}{0.222904in}}{\pgfqpoint{1.411163in}{1.411163in}}%
\pgfusepath{clip}%
\pgfsetbuttcap%
\pgfsetroundjoin%
\definecolor{currentfill}{rgb}{0.749020,0.772549,0.792157}%
\pgfsetfillcolor{currentfill}%
\pgfsetfillopacity{0.160000}%
\pgfsetlinewidth{0.000000pt}%
\definecolor{currentstroke}{rgb}{0.000000,0.000000,0.000000}%
\pgfsetstrokecolor{currentstroke}%
\pgfsetdash{}{0pt}%
\pgfpathmoveto{\pgfqpoint{0.982060in}{0.733342in}}%
\pgfpathlineto{\pgfqpoint{0.976395in}{0.739774in}}%
\pgfpathlineto{\pgfqpoint{0.970704in}{0.746235in}}%
\pgfpathlineto{\pgfqpoint{0.989386in}{0.743524in}}%
\pgfpathlineto{\pgfqpoint{1.008549in}{0.741349in}}%
\pgfpathlineto{\pgfqpoint{1.012443in}{0.735113in}}%
\pgfpathlineto{\pgfqpoint{1.016319in}{0.728905in}}%
\pgfpathlineto{\pgfqpoint{0.998974in}{0.730880in}}%
\pgfpathlineto{\pgfqpoint{0.982060in}{0.733342in}}%
\pgfpathclose%
\pgfusepath{fill}%
\end{pgfscope}%
\begin{pgfscope}%
\pgfpathrectangle{\pgfqpoint{0.407592in}{0.222904in}}{\pgfqpoint{1.411163in}{1.411163in}}%
\pgfusepath{clip}%
\pgfsetbuttcap%
\pgfsetroundjoin%
\definecolor{currentfill}{rgb}{0.749020,0.772549,0.792157}%
\pgfsetfillcolor{currentfill}%
\pgfsetfillopacity{0.160000}%
\pgfsetlinewidth{0.000000pt}%
\definecolor{currentstroke}{rgb}{0.000000,0.000000,0.000000}%
\pgfsetstrokecolor{currentstroke}%
\pgfsetdash{}{0pt}%
\pgfpathmoveto{\pgfqpoint{1.562509in}{0.970412in}}%
\pgfpathlineto{\pgfqpoint{1.575592in}{0.980381in}}%
\pgfpathlineto{\pgfqpoint{1.588709in}{0.990377in}}%
\pgfpathlineto{\pgfqpoint{1.601470in}{1.003751in}}%
\pgfpathlineto{\pgfqpoint{1.611841in}{1.017328in}}%
\pgfpathlineto{\pgfqpoint{1.598228in}{1.006711in}}%
\pgfpathlineto{\pgfqpoint{1.584645in}{0.996118in}}%
\pgfpathlineto{\pgfqpoint{1.574706in}{0.983166in}}%
\pgfpathlineto{\pgfqpoint{1.562509in}{0.970412in}}%
\pgfpathclose%
\pgfusepath{fill}%
\end{pgfscope}%
\begin{pgfscope}%
\pgfpathrectangle{\pgfqpoint{0.407592in}{0.222904in}}{\pgfqpoint{1.411163in}{1.411163in}}%
\pgfusepath{clip}%
\pgfsetbuttcap%
\pgfsetroundjoin%
\definecolor{currentfill}{rgb}{0.749020,0.772549,0.792157}%
\pgfsetfillcolor{currentfill}%
\pgfsetfillopacity{0.160000}%
\pgfsetlinewidth{0.000000pt}%
\definecolor{currentstroke}{rgb}{0.000000,0.000000,0.000000}%
\pgfsetstrokecolor{currentstroke}%
\pgfsetdash{}{0pt}%
\pgfpathmoveto{\pgfqpoint{0.680268in}{0.909008in}}%
\pgfpathlineto{\pgfqpoint{0.666882in}{0.918754in}}%
\pgfpathlineto{\pgfqpoint{0.653460in}{0.928526in}}%
\pgfpathlineto{\pgfqpoint{0.664307in}{0.917332in}}%
\pgfpathlineto{\pgfqpoint{0.677143in}{0.906448in}}%
\pgfpathlineto{\pgfqpoint{0.689921in}{0.897295in}}%
\pgfpathlineto{\pgfqpoint{0.702659in}{0.888170in}}%
\pgfpathlineto{\pgfqpoint{0.690534in}{0.898440in}}%
\pgfpathlineto{\pgfqpoint{0.680268in}{0.909008in}}%
\pgfpathclose%
\pgfusepath{fill}%
\end{pgfscope}%
\begin{pgfscope}%
\pgfpathrectangle{\pgfqpoint{0.407592in}{0.222904in}}{\pgfqpoint{1.411163in}{1.411163in}}%
\pgfusepath{clip}%
\pgfsetbuttcap%
\pgfsetroundjoin%
\definecolor{currentfill}{rgb}{0.749020,0.772549,0.792157}%
\pgfsetfillcolor{currentfill}%
\pgfsetfillopacity{0.160000}%
\pgfsetlinewidth{0.000000pt}%
\definecolor{currentstroke}{rgb}{0.000000,0.000000,0.000000}%
\pgfsetstrokecolor{currentstroke}%
\pgfsetdash{}{0pt}%
\pgfpathmoveto{\pgfqpoint{0.902514in}{0.762114in}}%
\pgfpathlineto{\pgfqpoint{0.893579in}{0.769329in}}%
\pgfpathlineto{\pgfqpoint{0.884607in}{0.776574in}}%
\pgfpathlineto{\pgfqpoint{0.902004in}{0.771462in}}%
\pgfpathlineto{\pgfqpoint{0.920300in}{0.766854in}}%
\pgfpathlineto{\pgfqpoint{0.927734in}{0.760015in}}%
\pgfpathlineto{\pgfqpoint{0.935136in}{0.753205in}}%
\pgfpathlineto{\pgfqpoint{0.918418in}{0.757428in}}%
\pgfpathlineto{\pgfqpoint{0.902514in}{0.762114in}}%
\pgfpathclose%
\pgfusepath{fill}%
\end{pgfscope}%
\begin{pgfscope}%
\pgfpathrectangle{\pgfqpoint{0.407592in}{0.222904in}}{\pgfqpoint{1.411163in}{1.411163in}}%
\pgfusepath{clip}%
\pgfsetbuttcap%
\pgfsetroundjoin%
\definecolor{currentfill}{rgb}{0.749020,0.772549,0.792157}%
\pgfsetfillcolor{currentfill}%
\pgfsetfillopacity{0.160000}%
\pgfsetlinewidth{0.000000pt}%
\definecolor{currentstroke}{rgb}{0.000000,0.000000,0.000000}%
\pgfsetstrokecolor{currentstroke}%
\pgfsetdash{}{0pt}%
\pgfpathmoveto{\pgfqpoint{1.124709in}{0.727683in}}%
\pgfpathlineto{\pgfqpoint{1.126518in}{0.733830in}}%
\pgfpathlineto{\pgfqpoint{1.128335in}{0.740005in}}%
\pgfpathlineto{\pgfqpoint{1.148195in}{0.741742in}}%
\pgfpathlineto{\pgfqpoint{1.167767in}{0.744025in}}%
\pgfpathlineto{\pgfqpoint{1.164073in}{0.737666in}}%
\pgfpathlineto{\pgfqpoint{1.160396in}{0.731335in}}%
\pgfpathlineto{\pgfqpoint{1.142681in}{0.729261in}}%
\pgfpathlineto{\pgfqpoint{1.124709in}{0.727683in}}%
\pgfpathclose%
\pgfusepath{fill}%
\end{pgfscope}%
\begin{pgfscope}%
\pgfpathrectangle{\pgfqpoint{0.407592in}{0.222904in}}{\pgfqpoint{1.411163in}{1.411163in}}%
\pgfusepath{clip}%
\pgfsetbuttcap%
\pgfsetroundjoin%
\definecolor{currentfill}{rgb}{0.749020,0.772549,0.792157}%
\pgfsetfillcolor{currentfill}%
\pgfsetfillopacity{0.160000}%
\pgfsetlinewidth{0.000000pt}%
\definecolor{currentstroke}{rgb}{0.000000,0.000000,0.000000}%
\pgfsetstrokecolor{currentstroke}%
\pgfsetdash{}{0pt}%
\pgfpathmoveto{\pgfqpoint{0.461173in}{1.109555in}}%
\pgfpathlineto{\pgfqpoint{0.446900in}{1.120878in}}%
\pgfpathlineto{\pgfqpoint{0.432600in}{1.132222in}}%
\pgfpathlineto{\pgfqpoint{0.435411in}{1.116030in}}%
\pgfpathlineto{\pgfqpoint{0.441176in}{1.099909in}}%
\pgfpathlineto{\pgfqpoint{0.455436in}{1.089158in}}%
\pgfpathlineto{\pgfqpoint{0.469662in}{1.078431in}}%
\pgfpathlineto{\pgfqpoint{0.464007in}{1.093956in}}%
\pgfpathlineto{\pgfqpoint{0.461173in}{1.109555in}}%
\pgfpathclose%
\pgfusepath{fill}%
\end{pgfscope}%
\begin{pgfscope}%
\pgfpathrectangle{\pgfqpoint{0.407592in}{0.222904in}}{\pgfqpoint{1.411163in}{1.411163in}}%
\pgfusepath{clip}%
\pgfsetbuttcap%
\pgfsetroundjoin%
\definecolor{currentfill}{rgb}{0.749020,0.772549,0.792157}%
\pgfsetfillcolor{currentfill}%
\pgfsetfillopacity{0.160000}%
\pgfsetlinewidth{0.000000pt}%
\definecolor{currentstroke}{rgb}{0.000000,0.000000,0.000000}%
\pgfsetstrokecolor{currentstroke}%
\pgfsetdash{}{0pt}%
\pgfpathmoveto{\pgfqpoint{0.582297in}{0.993507in}}%
\pgfpathlineto{\pgfqpoint{0.568331in}{1.004037in}}%
\pgfpathlineto{\pgfqpoint{0.554333in}{1.014591in}}%
\pgfpathlineto{\pgfqpoint{0.562025in}{1.001050in}}%
\pgfpathlineto{\pgfqpoint{0.572151in}{0.987724in}}%
\pgfpathlineto{\pgfqpoint{0.585795in}{0.977790in}}%
\pgfpathlineto{\pgfqpoint{0.599402in}{0.967883in}}%
\pgfpathlineto{\pgfqpoint{0.589699in}{0.980589in}}%
\pgfpathlineto{\pgfqpoint{0.582297in}{0.993507in}}%
\pgfpathclose%
\pgfusepath{fill}%
\end{pgfscope}%
\begin{pgfscope}%
\pgfpathrectangle{\pgfqpoint{0.407592in}{0.222904in}}{\pgfqpoint{1.411163in}{1.411163in}}%
\pgfusepath{clip}%
\pgfsetbuttcap%
\pgfsetroundjoin%
\definecolor{currentfill}{rgb}{0.749020,0.772549,0.792157}%
\pgfsetfillcolor{currentfill}%
\pgfsetfillopacity{0.160000}%
\pgfsetlinewidth{0.000000pt}%
\definecolor{currentstroke}{rgb}{0.000000,0.000000,0.000000}%
\pgfsetstrokecolor{currentstroke}%
\pgfsetdash{}{0pt}%
\pgfpathmoveto{\pgfqpoint{1.016319in}{0.728905in}}%
\pgfpathlineto{\pgfqpoint{1.012443in}{0.735113in}}%
\pgfpathlineto{\pgfqpoint{1.008549in}{0.741349in}}%
\pgfpathlineto{\pgfqpoint{1.028100in}{0.739723in}}%
\pgfpathlineto{\pgfqpoint{1.047946in}{0.738652in}}%
\pgfpathlineto{\pgfqpoint{1.049963in}{0.732539in}}%
\pgfpathlineto{\pgfqpoint{1.051971in}{0.726455in}}%
\pgfpathlineto{\pgfqpoint{1.034013in}{0.727427in}}%
\pgfpathlineto{\pgfqpoint{1.016319in}{0.728905in}}%
\pgfpathclose%
\pgfusepath{fill}%
\end{pgfscope}%
\begin{pgfscope}%
\pgfpathrectangle{\pgfqpoint{0.407592in}{0.222904in}}{\pgfqpoint{1.411163in}{1.411163in}}%
\pgfusepath{clip}%
\pgfsetbuttcap%
\pgfsetroundjoin%
\definecolor{currentfill}{rgb}{0.749020,0.772549,0.792157}%
\pgfsetfillcolor{currentfill}%
\pgfsetfillopacity{0.160000}%
\pgfsetlinewidth{0.000000pt}%
\definecolor{currentstroke}{rgb}{0.000000,0.000000,0.000000}%
\pgfsetstrokecolor{currentstroke}%
\pgfsetdash{}{0pt}%
\pgfpathmoveto{\pgfqpoint{1.419417in}{0.862526in}}%
\pgfpathlineto{\pgfqpoint{1.430949in}{0.871353in}}%
\pgfpathlineto{\pgfqpoint{1.442518in}{0.880209in}}%
\pgfpathlineto{\pgfqpoint{1.458363in}{0.890204in}}%
\pgfpathlineto{\pgfqpoint{1.446309in}{0.881046in}}%
\pgfpathlineto{\pgfqpoint{1.434292in}{0.871917in}}%
\pgfusepath{fill}%
\end{pgfscope}%
\begin{pgfscope}%
\pgfpathrectangle{\pgfqpoint{0.407592in}{0.222904in}}{\pgfqpoint{1.411163in}{1.411163in}}%
\pgfusepath{clip}%
\pgfsetbuttcap%
\pgfsetroundjoin%
\definecolor{currentfill}{rgb}{0.749020,0.772549,0.792157}%
\pgfsetfillcolor{currentfill}%
\pgfsetfillopacity{0.160000}%
\pgfsetlinewidth{0.000000pt}%
\definecolor{currentstroke}{rgb}{0.000000,0.000000,0.000000}%
\pgfsetstrokecolor{currentstroke}%
\pgfsetdash{}{0pt}%
\pgfpathmoveto{\pgfqpoint{0.802460in}{0.818128in}}%
\pgfpathlineto{\pgfqpoint{0.790871in}{0.826488in}}%
\pgfpathlineto{\pgfqpoint{0.779242in}{0.834877in}}%
\pgfpathlineto{\pgfqpoint{0.794484in}{0.826919in}}%
\pgfpathlineto{\pgfqpoint{0.811131in}{0.819403in}}%
\pgfpathlineto{\pgfqpoint{0.821615in}{0.811558in}}%
\pgfpathlineto{\pgfqpoint{0.832058in}{0.803742in}}%
\pgfpathlineto{\pgfqpoint{0.816613in}{0.810728in}}%
\pgfpathlineto{\pgfqpoint{0.802460in}{0.818128in}}%
\pgfpathclose%
\pgfusepath{fill}%
\end{pgfscope}%
\begin{pgfscope}%
\pgfpathrectangle{\pgfqpoint{0.407592in}{0.222904in}}{\pgfqpoint{1.411163in}{1.411163in}}%
\pgfusepath{clip}%
\pgfsetbuttcap%
\pgfsetroundjoin%
\definecolor{currentfill}{rgb}{0.749020,0.772549,0.792157}%
\pgfsetfillcolor{currentfill}%
\pgfsetfillopacity{0.160000}%
\pgfsetlinewidth{0.000000pt}%
\definecolor{currentstroke}{rgb}{0.000000,0.000000,0.000000}%
\pgfsetstrokecolor{currentstroke}%
\pgfsetdash{}{0pt}%
\pgfpathmoveto{\pgfqpoint{1.088334in}{0.726045in}}%
\pgfpathlineto{\pgfqpoint{1.088235in}{0.732109in}}%
\pgfpathlineto{\pgfqpoint{1.088135in}{0.738201in}}%
\pgfpathlineto{\pgfqpoint{1.108283in}{0.738822in}}%
\pgfpathlineto{\pgfqpoint{1.128335in}{0.740005in}}%
\pgfpathlineto{\pgfqpoint{1.126518in}{0.733830in}}%
\pgfpathlineto{\pgfqpoint{1.124709in}{0.727683in}}%
\pgfpathlineto{\pgfqpoint{1.106564in}{0.726610in}}%
\pgfpathlineto{\pgfqpoint{1.088334in}{0.726045in}}%
\pgfpathclose%
\pgfusepath{fill}%
\end{pgfscope}%
\begin{pgfscope}%
\pgfpathrectangle{\pgfqpoint{0.407592in}{0.222904in}}{\pgfqpoint{1.411163in}{1.411163in}}%
\pgfusepath{clip}%
\pgfsetbuttcap%
\pgfsetroundjoin%
\definecolor{currentfill}{rgb}{0.749020,0.772549,0.792157}%
\pgfsetfillcolor{currentfill}%
\pgfsetfillopacity{0.160000}%
\pgfsetlinewidth{0.000000pt}%
\definecolor{currentstroke}{rgb}{0.000000,0.000000,0.000000}%
\pgfsetstrokecolor{currentstroke}%
\pgfsetdash{}{0pt}%
\pgfpathmoveto{\pgfqpoint{0.852825in}{0.788200in}}%
\pgfpathlineto{\pgfqpoint{0.842461in}{0.795956in}}%
\pgfpathlineto{\pgfqpoint{0.832058in}{0.803742in}}%
\pgfpathlineto{\pgfqpoint{0.848729in}{0.797205in}}%
\pgfpathlineto{\pgfqpoint{0.866553in}{0.791152in}}%
\pgfpathlineto{\pgfqpoint{0.875599in}{0.783848in}}%
\pgfpathlineto{\pgfqpoint{0.884607in}{0.776574in}}%
\pgfpathlineto{\pgfqpoint{0.868190in}{0.782163in}}%
\pgfpathlineto{\pgfqpoint{0.852825in}{0.788200in}}%
\pgfpathclose%
\pgfusepath{fill}%
\end{pgfscope}%
\begin{pgfscope}%
\pgfpathrectangle{\pgfqpoint{0.407592in}{0.222904in}}{\pgfqpoint{1.411163in}{1.411163in}}%
\pgfusepath{clip}%
\pgfsetbuttcap%
\pgfsetroundjoin%
\definecolor{currentfill}{rgb}{0.749020,0.772549,0.792157}%
\pgfsetfillcolor{currentfill}%
\pgfsetfillopacity{0.160000}%
\pgfsetlinewidth{0.000000pt}%
\definecolor{currentstroke}{rgb}{0.000000,0.000000,0.000000}%
\pgfsetstrokecolor{currentstroke}%
\pgfsetdash{}{0pt}%
\pgfpathmoveto{\pgfqpoint{1.205671in}{0.750178in}}%
\pgfpathlineto{\pgfqpoint{1.211186in}{0.756849in}}%
\pgfpathlineto{\pgfqpoint{1.216725in}{0.763549in}}%
\pgfpathlineto{\pgfqpoint{1.236599in}{0.767736in}}%
\pgfpathlineto{\pgfqpoint{1.255757in}{0.772447in}}%
\pgfpathlineto{\pgfqpoint{1.248526in}{0.765375in}}%
\pgfpathlineto{\pgfqpoint{1.241326in}{0.758332in}}%
\pgfpathlineto{\pgfqpoint{1.223823in}{0.754014in}}%
\pgfpathlineto{\pgfqpoint{1.205671in}{0.750178in}}%
\pgfpathclose%
\pgfusepath{fill}%
\end{pgfscope}%
\begin{pgfscope}%
\pgfpathrectangle{\pgfqpoint{0.407592in}{0.222904in}}{\pgfqpoint{1.411163in}{1.411163in}}%
\pgfusepath{clip}%
\pgfsetbuttcap%
\pgfsetroundjoin%
\definecolor{currentfill}{rgb}{0.749020,0.772549,0.792157}%
\pgfsetfillcolor{currentfill}%
\pgfsetfillopacity{0.160000}%
\pgfsetlinewidth{0.000000pt}%
\definecolor{currentstroke}{rgb}{0.000000,0.000000,0.000000}%
\pgfsetstrokecolor{currentstroke}%
\pgfsetdash{}{0pt}%
\pgfpathmoveto{\pgfqpoint{1.051971in}{0.726455in}}%
\pgfpathlineto{\pgfqpoint{1.049963in}{0.732539in}}%
\pgfpathlineto{\pgfqpoint{1.047946in}{0.738652in}}%
\pgfpathlineto{\pgfqpoint{1.067990in}{0.738144in}}%
\pgfpathlineto{\pgfqpoint{1.088135in}{0.738201in}}%
\pgfpathlineto{\pgfqpoint{1.088235in}{0.732109in}}%
\pgfpathlineto{\pgfqpoint{1.088334in}{0.726045in}}%
\pgfpathlineto{\pgfqpoint{1.070107in}{0.725994in}}%
\pgfpathlineto{\pgfqpoint{1.051971in}{0.726455in}}%
\pgfpathclose%
\pgfusepath{fill}%
\end{pgfscope}%
\begin{pgfscope}%
\pgfpathrectangle{\pgfqpoint{0.407592in}{0.222904in}}{\pgfqpoint{1.411163in}{1.411163in}}%
\pgfusepath{clip}%
\pgfsetbuttcap%
\pgfsetroundjoin%
\definecolor{currentfill}{rgb}{0.749020,0.772549,0.792157}%
\pgfsetfillcolor{currentfill}%
\pgfsetfillopacity{0.160000}%
\pgfsetlinewidth{0.000000pt}%
\definecolor{currentstroke}{rgb}{0.000000,0.000000,0.000000}%
\pgfsetstrokecolor{currentstroke}%
\pgfsetdash{}{0pt}%
\pgfpathmoveto{\pgfqpoint{1.364368in}{0.828482in}}%
\pgfpathlineto{\pgfqpoint{1.374675in}{0.836677in}}%
\pgfpathlineto{\pgfqpoint{1.385021in}{0.844902in}}%
\pgfpathlineto{\pgfqpoint{1.402972in}{0.853506in}}%
\pgfpathlineto{\pgfqpoint{1.419417in}{0.862526in}}%
\pgfpathlineto{\pgfqpoint{1.407924in}{0.853728in}}%
\pgfpathlineto{\pgfqpoint{1.396468in}{0.844958in}}%
\pgfpathlineto{\pgfqpoint{1.381114in}{0.836524in}}%
\pgfpathlineto{\pgfqpoint{1.364368in}{0.828482in}}%
\pgfpathclose%
\pgfusepath{fill}%
\end{pgfscope}%
\begin{pgfscope}%
\pgfpathrectangle{\pgfqpoint{0.407592in}{0.222904in}}{\pgfqpoint{1.411163in}{1.411163in}}%
\pgfusepath{clip}%
\pgfsetbuttcap%
\pgfsetroundjoin%
\definecolor{currentfill}{rgb}{0.749020,0.772549,0.792157}%
\pgfsetfillcolor{currentfill}%
\pgfsetfillopacity{0.160000}%
\pgfsetlinewidth{0.000000pt}%
\definecolor{currentstroke}{rgb}{0.000000,0.000000,0.000000}%
\pgfsetstrokecolor{currentstroke}%
\pgfsetdash{}{0pt}%
\pgfpathmoveto{\pgfqpoint{1.721866in}{1.103134in}}%
\pgfpathlineto{\pgfqpoint{1.735760in}{1.113970in}}%
\pgfpathlineto{\pgfqpoint{1.749686in}{1.124831in}}%
\pgfpathlineto{\pgfqpoint{1.758991in}{1.141560in}}%
\pgfpathlineto{\pgfqpoint{1.765228in}{1.158340in}}%
\pgfpathlineto{\pgfqpoint{1.751127in}{1.146896in}}%
\pgfpathlineto{\pgfqpoint{1.737051in}{1.135472in}}%
\pgfpathlineto{\pgfqpoint{1.730926in}{1.119275in}}%
\pgfpathlineto{\pgfqpoint{1.721866in}{1.103134in}}%
\pgfpathclose%
\pgfusepath{fill}%
\end{pgfscope}%
\begin{pgfscope}%
\pgfpathrectangle{\pgfqpoint{0.407592in}{0.222904in}}{\pgfqpoint{1.411163in}{1.411163in}}%
\pgfusepath{clip}%
\pgfsetbuttcap%
\pgfsetroundjoin%
\definecolor{currentfill}{rgb}{0.749020,0.772549,0.792157}%
\pgfsetfillcolor{currentfill}%
\pgfsetfillopacity{0.160000}%
\pgfsetlinewidth{0.000000pt}%
\definecolor{currentstroke}{rgb}{0.000000,0.000000,0.000000}%
\pgfsetstrokecolor{currentstroke}%
\pgfsetdash{}{0pt}%
\pgfpathmoveto{\pgfqpoint{1.482584in}{0.908604in}}%
\pgfpathlineto{\pgfqpoint{1.494750in}{0.917847in}}%
\pgfpathlineto{\pgfqpoint{1.506955in}{0.927119in}}%
\pgfpathlineto{\pgfqpoint{1.522730in}{0.938681in}}%
\pgfpathlineto{\pgfqpoint{1.536450in}{0.950555in}}%
\pgfpathlineto{\pgfqpoint{1.523473in}{0.940667in}}%
\pgfpathlineto{\pgfqpoint{1.510531in}{0.930804in}}%
\pgfpathlineto{\pgfqpoint{1.497521in}{0.919554in}}%
\pgfpathlineto{\pgfqpoint{1.482584in}{0.908604in}}%
\pgfpathclose%
\pgfusepath{fill}%
\end{pgfscope}%
\begin{pgfscope}%
\pgfpathrectangle{\pgfqpoint{0.407592in}{0.222904in}}{\pgfqpoint{1.411163in}{1.411163in}}%
\pgfusepath{clip}%
\pgfsetbuttcap%
\pgfsetroundjoin%
\definecolor{currentfill}{rgb}{0.749020,0.772549,0.792157}%
\pgfsetfillcolor{currentfill}%
\pgfsetfillopacity{0.160000}%
\pgfsetlinewidth{0.000000pt}%
\definecolor{currentstroke}{rgb}{0.000000,0.000000,0.000000}%
\pgfsetstrokecolor{currentstroke}%
\pgfsetdash{}{0pt}%
\pgfpathmoveto{\pgfqpoint{1.588709in}{0.990377in}}%
\pgfpathlineto{\pgfqpoint{1.601862in}{1.000399in}}%
\pgfpathlineto{\pgfqpoint{1.615051in}{1.010449in}}%
\pgfpathlineto{\pgfqpoint{1.628367in}{1.024439in}}%
\pgfpathlineto{\pgfqpoint{1.639161in}{1.038634in}}%
\pgfpathlineto{\pgfqpoint{1.625485in}{1.027969in}}%
\pgfpathlineto{\pgfqpoint{1.611841in}{1.017328in}}%
\pgfpathlineto{\pgfqpoint{1.601470in}{1.003751in}}%
\pgfpathlineto{\pgfqpoint{1.588709in}{0.990377in}}%
\pgfpathclose%
\pgfusepath{fill}%
\end{pgfscope}%
\begin{pgfscope}%
\pgfpathrectangle{\pgfqpoint{0.407592in}{0.222904in}}{\pgfqpoint{1.411163in}{1.411163in}}%
\pgfusepath{clip}%
\pgfsetbuttcap%
\pgfsetroundjoin%
\definecolor{currentfill}{rgb}{0.749020,0.772549,0.792157}%
\pgfsetfillcolor{currentfill}%
\pgfsetfillopacity{0.160000}%
\pgfsetlinewidth{0.000000pt}%
\definecolor{currentstroke}{rgb}{0.000000,0.000000,0.000000}%
\pgfsetstrokecolor{currentstroke}%
\pgfsetdash{}{0pt}%
\pgfpathmoveto{\pgfqpoint{1.255757in}{0.772447in}}%
\pgfpathlineto{\pgfqpoint{1.263018in}{0.779550in}}%
\pgfpathlineto{\pgfqpoint{1.270310in}{0.786683in}}%
\pgfpathlineto{\pgfqpoint{1.290248in}{0.792326in}}%
\pgfpathlineto{\pgfqpoint{1.309209in}{0.798479in}}%
\pgfpathlineto{\pgfqpoint{1.300372in}{0.790894in}}%
\pgfpathlineto{\pgfqpoint{1.291569in}{0.783339in}}%
\pgfpathlineto{\pgfqpoint{1.274108in}{0.777657in}}%
\pgfpathlineto{\pgfqpoint{1.255757in}{0.772447in}}%
\pgfpathclose%
\pgfusepath{fill}%
\end{pgfscope}%
\begin{pgfscope}%
\pgfpathrectangle{\pgfqpoint{0.407592in}{0.222904in}}{\pgfqpoint{1.411163in}{1.411163in}}%
\pgfusepath{clip}%
\pgfsetbuttcap%
\pgfsetroundjoin%
\definecolor{currentfill}{rgb}{0.749020,0.772549,0.792157}%
\pgfsetfillcolor{currentfill}%
\pgfsetfillopacity{0.160000}%
\pgfsetlinewidth{0.000000pt}%
\definecolor{currentstroke}{rgb}{0.000000,0.000000,0.000000}%
\pgfsetstrokecolor{currentstroke}%
\pgfsetdash{}{0pt}%
\pgfpathmoveto{\pgfqpoint{0.935136in}{0.753205in}}%
\pgfpathlineto{\pgfqpoint{0.927734in}{0.760015in}}%
\pgfpathlineto{\pgfqpoint{0.920300in}{0.766854in}}%
\pgfpathlineto{\pgfqpoint{0.939410in}{0.762774in}}%
\pgfpathlineto{\pgfqpoint{0.959245in}{0.759245in}}%
\pgfpathlineto{\pgfqpoint{0.964987in}{0.752725in}}%
\pgfpathlineto{\pgfqpoint{0.970704in}{0.746235in}}%
\pgfpathlineto{\pgfqpoint{0.952592in}{0.749468in}}%
\pgfpathlineto{\pgfqpoint{0.935136in}{0.753205in}}%
\pgfpathclose%
\pgfusepath{fill}%
\end{pgfscope}%
\begin{pgfscope}%
\pgfpathrectangle{\pgfqpoint{0.407592in}{0.222904in}}{\pgfqpoint{1.411163in}{1.411163in}}%
\pgfusepath{clip}%
\pgfsetbuttcap%
\pgfsetroundjoin%
\definecolor{currentfill}{rgb}{0.749020,0.772549,0.792157}%
\pgfsetfillcolor{currentfill}%
\pgfsetfillopacity{0.160000}%
\pgfsetlinewidth{0.000000pt}%
\definecolor{currentstroke}{rgb}{0.000000,0.000000,0.000000}%
\pgfsetstrokecolor{currentstroke}%
\pgfsetdash{}{0pt}%
\pgfpathmoveto{\pgfqpoint{1.309209in}{0.798479in}}%
\pgfpathlineto{\pgfqpoint{1.318081in}{0.806094in}}%
\pgfpathlineto{\pgfqpoint{1.326989in}{0.813739in}}%
\pgfpathlineto{\pgfqpoint{1.346300in}{0.820874in}}%
\pgfpathlineto{\pgfqpoint{1.364368in}{0.828482in}}%
\pgfpathlineto{\pgfqpoint{1.354097in}{0.820316in}}%
\pgfpathlineto{\pgfqpoint{1.343864in}{0.812181in}}%
\pgfpathlineto{\pgfqpoint{1.327108in}{0.805109in}}%
\pgfpathlineto{\pgfqpoint{1.309209in}{0.798479in}}%
\pgfpathclose%
\pgfusepath{fill}%
\end{pgfscope}%
\begin{pgfscope}%
\pgfpathrectangle{\pgfqpoint{0.407592in}{0.222904in}}{\pgfqpoint{1.411163in}{1.411163in}}%
\pgfusepath{clip}%
\pgfsetbuttcap%
\pgfsetroundjoin%
\definecolor{currentfill}{rgb}{0.749020,0.772549,0.792157}%
\pgfsetfillcolor{currentfill}%
\pgfsetfillopacity{0.160000}%
\pgfsetlinewidth{0.000000pt}%
\definecolor{currentstroke}{rgb}{0.000000,0.000000,0.000000}%
\pgfsetstrokecolor{currentstroke}%
\pgfsetdash{}{0pt}%
\pgfpathmoveto{\pgfqpoint{0.432600in}{1.132222in}}%
\pgfpathlineto{\pgfqpoint{0.418274in}{1.143587in}}%
\pgfpathlineto{\pgfqpoint{0.403920in}{1.154973in}}%
\pgfpathlineto{\pgfqpoint{0.406696in}{1.138197in}}%
\pgfpathlineto{\pgfqpoint{0.412558in}{1.121487in}}%
\pgfpathlineto{\pgfqpoint{0.426884in}{1.110686in}}%
\pgfpathlineto{\pgfqpoint{0.441176in}{1.099909in}}%
\pgfpathlineto{\pgfqpoint{0.435411in}{1.116030in}}%
\pgfpathlineto{\pgfqpoint{0.432600in}{1.132222in}}%
\pgfpathclose%
\pgfusepath{fill}%
\end{pgfscope}%
\begin{pgfscope}%
\pgfpathrectangle{\pgfqpoint{0.407592in}{0.222904in}}{\pgfqpoint{1.411163in}{1.411163in}}%
\pgfusepath{clip}%
\pgfsetbuttcap%
\pgfsetroundjoin%
\definecolor{currentfill}{rgb}{0.749020,0.772549,0.792157}%
\pgfsetfillcolor{currentfill}%
\pgfsetfillopacity{0.160000}%
\pgfsetlinewidth{0.000000pt}%
\definecolor{currentstroke}{rgb}{0.000000,0.000000,0.000000}%
\pgfsetstrokecolor{currentstroke}%
\pgfsetdash{}{0pt}%
\pgfpathmoveto{\pgfqpoint{0.653460in}{0.928526in}}%
\pgfpathlineto{\pgfqpoint{0.640001in}{0.938325in}}%
\pgfpathlineto{\pgfqpoint{0.626505in}{0.948151in}}%
\pgfpathlineto{\pgfqpoint{0.637925in}{0.936335in}}%
\pgfpathlineto{\pgfqpoint{0.651466in}{0.924842in}}%
\pgfpathlineto{\pgfqpoint{0.664324in}{0.915631in}}%
\pgfpathlineto{\pgfqpoint{0.677143in}{0.906448in}}%
\pgfpathlineto{\pgfqpoint{0.664307in}{0.917332in}}%
\pgfpathlineto{\pgfqpoint{0.653460in}{0.928526in}}%
\pgfpathclose%
\pgfusepath{fill}%
\end{pgfscope}%
\begin{pgfscope}%
\pgfpathrectangle{\pgfqpoint{0.407592in}{0.222904in}}{\pgfqpoint{1.411163in}{1.411163in}}%
\pgfusepath{clip}%
\pgfsetbuttcap%
\pgfsetroundjoin%
\definecolor{currentfill}{rgb}{0.749020,0.772549,0.792157}%
\pgfsetfillcolor{currentfill}%
\pgfsetfillopacity{0.160000}%
\pgfsetlinewidth{0.000000pt}%
\definecolor{currentstroke}{rgb}{0.000000,0.000000,0.000000}%
\pgfsetstrokecolor{currentstroke}%
\pgfsetdash{}{0pt}%
\pgfpathmoveto{\pgfqpoint{0.728016in}{0.870005in}}%
\pgfpathlineto{\pgfqpoint{0.715357in}{0.879073in}}%
\pgfpathlineto{\pgfqpoint{0.702659in}{0.888170in}}%
\pgfpathlineto{\pgfqpoint{0.716603in}{0.878250in}}%
\pgfpathlineto{\pgfqpoint{0.732315in}{0.868730in}}%
\pgfpathlineto{\pgfqpoint{0.744108in}{0.860222in}}%
\pgfpathlineto{\pgfqpoint{0.755860in}{0.851745in}}%
\pgfpathlineto{\pgfqpoint{0.741116in}{0.860685in}}%
\pgfpathlineto{\pgfqpoint{0.728016in}{0.870005in}}%
\pgfpathclose%
\pgfusepath{fill}%
\end{pgfscope}%
\begin{pgfscope}%
\pgfpathrectangle{\pgfqpoint{0.407592in}{0.222904in}}{\pgfqpoint{1.411163in}{1.411163in}}%
\pgfusepath{clip}%
\pgfsetbuttcap%
\pgfsetroundjoin%
\definecolor{currentfill}{rgb}{0.749020,0.772549,0.792157}%
\pgfsetfillcolor{currentfill}%
\pgfsetfillopacity{0.160000}%
\pgfsetlinewidth{0.000000pt}%
\definecolor{currentstroke}{rgb}{0.000000,0.000000,0.000000}%
\pgfsetstrokecolor{currentstroke}%
\pgfsetdash{}{0pt}%
\pgfpathmoveto{\pgfqpoint{0.554333in}{1.014591in}}%
\pgfpathlineto{\pgfqpoint{0.540302in}{1.025170in}}%
\pgfpathlineto{\pgfqpoint{0.526239in}{1.035773in}}%
\pgfpathlineto{\pgfqpoint{0.534210in}{1.021614in}}%
\pgfpathlineto{\pgfqpoint{0.544749in}{1.007673in}}%
\pgfpathlineto{\pgfqpoint{0.558469in}{0.997685in}}%
\pgfpathlineto{\pgfqpoint{0.572151in}{0.987724in}}%
\pgfpathlineto{\pgfqpoint{0.562025in}{1.001050in}}%
\pgfpathlineto{\pgfqpoint{0.554333in}{1.014591in}}%
\pgfpathclose%
\pgfusepath{fill}%
\end{pgfscope}%
\begin{pgfscope}%
\pgfpathrectangle{\pgfqpoint{0.407592in}{0.222904in}}{\pgfqpoint{1.411163in}{1.411163in}}%
\pgfusepath{clip}%
\pgfsetbuttcap%
\pgfsetroundjoin%
\definecolor{currentfill}{rgb}{0.749020,0.772549,0.792157}%
\pgfsetfillcolor{currentfill}%
\pgfsetfillopacity{0.160000}%
\pgfsetlinewidth{0.000000pt}%
\definecolor{currentstroke}{rgb}{0.000000,0.000000,0.000000}%
\pgfsetstrokecolor{currentstroke}%
\pgfsetdash{}{0pt}%
\pgfpathmoveto{\pgfqpoint{1.167767in}{0.744025in}}%
\pgfpathlineto{\pgfqpoint{1.171478in}{0.750415in}}%
\pgfpathlineto{\pgfqpoint{1.175206in}{0.756833in}}%
\pgfpathlineto{\pgfqpoint{1.196228in}{0.759908in}}%
\pgfpathlineto{\pgfqpoint{1.216725in}{0.763549in}}%
\pgfpathlineto{\pgfqpoint{1.211186in}{0.756849in}}%
\pgfpathlineto{\pgfqpoint{1.205671in}{0.750178in}}%
\pgfpathlineto{\pgfqpoint{1.186956in}{0.746842in}}%
\pgfpathlineto{\pgfqpoint{1.167767in}{0.744025in}}%
\pgfpathclose%
\pgfusepath{fill}%
\end{pgfscope}%
\begin{pgfscope}%
\pgfpathrectangle{\pgfqpoint{0.407592in}{0.222904in}}{\pgfqpoint{1.411163in}{1.411163in}}%
\pgfusepath{clip}%
\pgfsetbuttcap%
\pgfsetroundjoin%
\definecolor{currentfill}{rgb}{0.749020,0.772549,0.792157}%
\pgfsetfillcolor{currentfill}%
\pgfsetfillopacity{0.160000}%
\pgfsetlinewidth{0.000000pt}%
\definecolor{currentstroke}{rgb}{0.000000,0.000000,0.000000}%
\pgfsetstrokecolor{currentstroke}%
\pgfsetdash{}{0pt}%
\pgfpathmoveto{\pgfqpoint{1.749686in}{1.124831in}}%
\pgfpathlineto{\pgfqpoint{1.763644in}{1.135716in}}%
\pgfpathlineto{\pgfqpoint{1.777633in}{1.146626in}}%
\pgfpathlineto{\pgfqpoint{1.787171in}{1.163937in}}%
\pgfpathlineto{\pgfqpoint{1.793508in}{1.181291in}}%
\pgfpathlineto{\pgfqpoint{1.779355in}{1.169805in}}%
\pgfpathlineto{\pgfqpoint{1.765228in}{1.158340in}}%
\pgfpathlineto{\pgfqpoint{1.758991in}{1.141560in}}%
\pgfpathlineto{\pgfqpoint{1.749686in}{1.124831in}}%
\pgfpathclose%
\pgfusepath{fill}%
\end{pgfscope}%
\begin{pgfscope}%
\pgfpathrectangle{\pgfqpoint{0.407592in}{0.222904in}}{\pgfqpoint{1.411163in}{1.411163in}}%
\pgfusepath{clip}%
\pgfsetbuttcap%
\pgfsetroundjoin%
\definecolor{currentfill}{rgb}{0.749020,0.772549,0.792157}%
\pgfsetfillcolor{currentfill}%
\pgfsetfillopacity{0.160000}%
\pgfsetlinewidth{0.000000pt}%
\definecolor{currentstroke}{rgb}{0.000000,0.000000,0.000000}%
\pgfsetstrokecolor{currentstroke}%
\pgfsetdash{}{0pt}%
\pgfpathmoveto{\pgfqpoint{1.442518in}{0.880209in}}%
\pgfpathlineto{\pgfqpoint{1.454126in}{0.889095in}}%
\pgfpathlineto{\pgfqpoint{1.465771in}{0.898010in}}%
\pgfpathlineto{\pgfqpoint{1.482584in}{0.908604in}}%
\pgfpathlineto{\pgfqpoint{1.470455in}{0.899390in}}%
\pgfpathlineto{\pgfqpoint{1.458363in}{0.890204in}}%
\pgfusepath{fill}%
\end{pgfscope}%
\begin{pgfscope}%
\pgfpathrectangle{\pgfqpoint{0.407592in}{0.222904in}}{\pgfqpoint{1.411163in}{1.411163in}}%
\pgfusepath{clip}%
\pgfsetbuttcap%
\pgfsetroundjoin%
\definecolor{currentfill}{rgb}{0.749020,0.772549,0.792157}%
\pgfsetfillcolor{currentfill}%
\pgfsetfillopacity{0.160000}%
\pgfsetlinewidth{0.000000pt}%
\definecolor{currentstroke}{rgb}{0.000000,0.000000,0.000000}%
\pgfsetstrokecolor{currentstroke}%
\pgfsetdash{}{0pt}%
\pgfpathmoveto{\pgfqpoint{0.970704in}{0.746235in}}%
\pgfpathlineto{\pgfqpoint{0.964987in}{0.752725in}}%
\pgfpathlineto{\pgfqpoint{0.959245in}{0.759245in}}%
\pgfpathlineto{\pgfqpoint{0.979709in}{0.756286in}}%
\pgfpathlineto{\pgfqpoint{1.000705in}{0.753911in}}%
\pgfpathlineto{\pgfqpoint{1.004636in}{0.747616in}}%
\pgfpathlineto{\pgfqpoint{1.008549in}{0.741349in}}%
\pgfpathlineto{\pgfqpoint{0.989386in}{0.743524in}}%
\pgfpathlineto{\pgfqpoint{0.970704in}{0.746235in}}%
\pgfpathclose%
\pgfusepath{fill}%
\end{pgfscope}%
\begin{pgfscope}%
\pgfpathrectangle{\pgfqpoint{0.407592in}{0.222904in}}{\pgfqpoint{1.411163in}{1.411163in}}%
\pgfusepath{clip}%
\pgfsetbuttcap%
\pgfsetroundjoin%
\definecolor{currentfill}{rgb}{0.749020,0.772549,0.792157}%
\pgfsetfillcolor{currentfill}%
\pgfsetfillopacity{0.160000}%
\pgfsetlinewidth{0.000000pt}%
\definecolor{currentstroke}{rgb}{0.000000,0.000000,0.000000}%
\pgfsetstrokecolor{currentstroke}%
\pgfsetdash{}{0pt}%
\pgfpathmoveto{\pgfqpoint{0.884607in}{0.776574in}}%
\pgfpathlineto{\pgfqpoint{0.875599in}{0.783848in}}%
\pgfpathlineto{\pgfqpoint{0.866553in}{0.791152in}}%
\pgfpathlineto{\pgfqpoint{0.885450in}{0.785614in}}%
\pgfpathlineto{\pgfqpoint{0.905333in}{0.780622in}}%
\pgfpathlineto{\pgfqpoint{0.912833in}{0.773723in}}%
\pgfpathlineto{\pgfqpoint{0.920300in}{0.766854in}}%
\pgfpathlineto{\pgfqpoint{0.902004in}{0.771462in}}%
\pgfpathlineto{\pgfqpoint{0.884607in}{0.776574in}}%
\pgfpathclose%
\pgfusepath{fill}%
\end{pgfscope}%
\begin{pgfscope}%
\pgfpathrectangle{\pgfqpoint{0.407592in}{0.222904in}}{\pgfqpoint{1.411163in}{1.411163in}}%
\pgfusepath{clip}%
\pgfsetbuttcap%
\pgfsetroundjoin%
\definecolor{currentfill}{rgb}{0.749020,0.772549,0.792157}%
\pgfsetfillcolor{currentfill}%
\pgfsetfillopacity{0.160000}%
\pgfsetlinewidth{0.000000pt}%
\definecolor{currentstroke}{rgb}{0.000000,0.000000,0.000000}%
\pgfsetstrokecolor{currentstroke}%
\pgfsetdash{}{0pt}%
\pgfpathmoveto{\pgfqpoint{0.779242in}{0.834877in}}%
\pgfpathlineto{\pgfqpoint{0.767572in}{0.843296in}}%
\pgfpathlineto{\pgfqpoint{0.755860in}{0.851745in}}%
\pgfpathlineto{\pgfqpoint{0.772192in}{0.843230in}}%
\pgfpathlineto{\pgfqpoint{0.790043in}{0.835185in}}%
\pgfpathlineto{\pgfqpoint{0.800608in}{0.827279in}}%
\pgfpathlineto{\pgfqpoint{0.811131in}{0.819403in}}%
\pgfpathlineto{\pgfqpoint{0.794484in}{0.826919in}}%
\pgfpathlineto{\pgfqpoint{0.779242in}{0.834877in}}%
\pgfpathclose%
\pgfusepath{fill}%
\end{pgfscope}%
\begin{pgfscope}%
\pgfpathrectangle{\pgfqpoint{0.407592in}{0.222904in}}{\pgfqpoint{1.411163in}{1.411163in}}%
\pgfusepath{clip}%
\pgfsetbuttcap%
\pgfsetroundjoin%
\definecolor{currentfill}{rgb}{0.749020,0.772549,0.792157}%
\pgfsetfillcolor{currentfill}%
\pgfsetfillopacity{0.160000}%
\pgfsetlinewidth{0.000000pt}%
\definecolor{currentstroke}{rgb}{0.000000,0.000000,0.000000}%
\pgfsetstrokecolor{currentstroke}%
\pgfsetdash{}{0pt}%
\pgfpathmoveto{\pgfqpoint{1.615051in}{1.010449in}}%
\pgfpathlineto{\pgfqpoint{1.628276in}{1.020526in}}%
\pgfpathlineto{\pgfqpoint{1.641536in}{1.030631in}}%
\pgfpathlineto{\pgfqpoint{1.655398in}{1.045230in}}%
\pgfpathlineto{\pgfqpoint{1.666604in}{1.060037in}}%
\pgfpathlineto{\pgfqpoint{1.652867in}{1.049323in}}%
\pgfpathlineto{\pgfqpoint{1.639161in}{1.038634in}}%
\pgfpathlineto{\pgfqpoint{1.628367in}{1.024439in}}%
\pgfpathlineto{\pgfqpoint{1.615051in}{1.010449in}}%
\pgfpathclose%
\pgfusepath{fill}%
\end{pgfscope}%
\begin{pgfscope}%
\pgfpathrectangle{\pgfqpoint{0.407592in}{0.222904in}}{\pgfqpoint{1.411163in}{1.411163in}}%
\pgfusepath{clip}%
\pgfsetbuttcap%
\pgfsetroundjoin%
\definecolor{currentfill}{rgb}{0.749020,0.772549,0.792157}%
\pgfsetfillcolor{currentfill}%
\pgfsetfillopacity{0.160000}%
\pgfsetlinewidth{0.000000pt}%
\definecolor{currentstroke}{rgb}{0.000000,0.000000,0.000000}%
\pgfsetstrokecolor{currentstroke}%
\pgfsetdash{}{0pt}%
\pgfpathmoveto{\pgfqpoint{0.403920in}{1.154973in}}%
\pgfpathlineto{\pgfqpoint{0.389540in}{1.166381in}}%
\pgfpathlineto{\pgfqpoint{0.375133in}{1.177810in}}%
\pgfpathlineto{\pgfqpoint{0.377860in}{1.160458in}}%
\pgfpathlineto{\pgfqpoint{0.383806in}{1.143166in}}%
\pgfpathlineto{\pgfqpoint{0.398199in}{1.132314in}}%
\pgfpathlineto{\pgfqpoint{0.412558in}{1.121487in}}%
\pgfpathlineto{\pgfqpoint{0.406696in}{1.138197in}}%
\pgfpathlineto{\pgfqpoint{0.403920in}{1.154973in}}%
\pgfpathclose%
\pgfusepath{fill}%
\end{pgfscope}%
\begin{pgfscope}%
\pgfpathrectangle{\pgfqpoint{0.407592in}{0.222904in}}{\pgfqpoint{1.411163in}{1.411163in}}%
\pgfusepath{clip}%
\pgfsetbuttcap%
\pgfsetroundjoin%
\definecolor{currentfill}{rgb}{0.749020,0.772549,0.792157}%
\pgfsetfillcolor{currentfill}%
\pgfsetfillopacity{0.160000}%
\pgfsetlinewidth{0.000000pt}%
\definecolor{currentstroke}{rgb}{0.000000,0.000000,0.000000}%
\pgfsetstrokecolor{currentstroke}%
\pgfsetdash{}{0pt}%
\pgfpathmoveto{\pgfqpoint{1.128335in}{0.740005in}}%
\pgfpathlineto{\pgfqpoint{1.130161in}{0.746209in}}%
\pgfpathlineto{\pgfqpoint{1.131996in}{0.752442in}}%
\pgfpathlineto{\pgfqpoint{1.153761in}{0.754340in}}%
\pgfpathlineto{\pgfqpoint{1.175206in}{0.756833in}}%
\pgfpathlineto{\pgfqpoint{1.171478in}{0.750415in}}%
\pgfpathlineto{\pgfqpoint{1.167767in}{0.744025in}}%
\pgfpathlineto{\pgfqpoint{1.148195in}{0.741742in}}%
\pgfpathlineto{\pgfqpoint{1.128335in}{0.740005in}}%
\pgfpathclose%
\pgfusepath{fill}%
\end{pgfscope}%
\begin{pgfscope}%
\pgfpathrectangle{\pgfqpoint{0.407592in}{0.222904in}}{\pgfqpoint{1.411163in}{1.411163in}}%
\pgfusepath{clip}%
\pgfsetbuttcap%
\pgfsetroundjoin%
\definecolor{currentfill}{rgb}{0.749020,0.772549,0.792157}%
\pgfsetfillcolor{currentfill}%
\pgfsetfillopacity{0.160000}%
\pgfsetlinewidth{0.000000pt}%
\definecolor{currentstroke}{rgb}{0.000000,0.000000,0.000000}%
\pgfsetstrokecolor{currentstroke}%
\pgfsetdash{}{0pt}%
\pgfpathmoveto{\pgfqpoint{1.506955in}{0.927119in}}%
\pgfpathlineto{\pgfqpoint{1.519197in}{0.936420in}}%
\pgfpathlineto{\pgfqpoint{1.531478in}{0.945749in}}%
\pgfpathlineto{\pgfqpoint{1.548086in}{0.957919in}}%
\pgfpathlineto{\pgfqpoint{1.562509in}{0.970412in}}%
\pgfpathlineto{\pgfqpoint{1.549462in}{0.960470in}}%
\pgfpathlineto{\pgfqpoint{1.536450in}{0.950555in}}%
\pgfpathlineto{\pgfqpoint{1.522730in}{0.938681in}}%
\pgfpathlineto{\pgfqpoint{1.506955in}{0.927119in}}%
\pgfpathclose%
\pgfusepath{fill}%
\end{pgfscope}%
\begin{pgfscope}%
\pgfpathrectangle{\pgfqpoint{0.407592in}{0.222904in}}{\pgfqpoint{1.411163in}{1.411163in}}%
\pgfusepath{clip}%
\pgfsetbuttcap%
\pgfsetroundjoin%
\definecolor{currentfill}{rgb}{0.749020,0.772549,0.792157}%
\pgfsetfillcolor{currentfill}%
\pgfsetfillopacity{0.160000}%
\pgfsetlinewidth{0.000000pt}%
\definecolor{currentstroke}{rgb}{0.000000,0.000000,0.000000}%
\pgfsetstrokecolor{currentstroke}%
\pgfsetdash{}{0pt}%
\pgfpathmoveto{\pgfqpoint{0.832058in}{0.803742in}}%
\pgfpathlineto{\pgfqpoint{0.821615in}{0.811558in}}%
\pgfpathlineto{\pgfqpoint{0.811131in}{0.819403in}}%
\pgfpathlineto{\pgfqpoint{0.829113in}{0.812368in}}%
\pgfpathlineto{\pgfqpoint{0.848350in}{0.805852in}}%
\pgfpathlineto{\pgfqpoint{0.857470in}{0.798487in}}%
\pgfpathlineto{\pgfqpoint{0.866553in}{0.791152in}}%
\pgfpathlineto{\pgfqpoint{0.848729in}{0.797205in}}%
\pgfpathlineto{\pgfqpoint{0.832058in}{0.803742in}}%
\pgfpathclose%
\pgfusepath{fill}%
\end{pgfscope}%
\begin{pgfscope}%
\pgfpathrectangle{\pgfqpoint{0.407592in}{0.222904in}}{\pgfqpoint{1.411163in}{1.411163in}}%
\pgfusepath{clip}%
\pgfsetbuttcap%
\pgfsetroundjoin%
\definecolor{currentfill}{rgb}{0.749020,0.772549,0.792157}%
\pgfsetfillcolor{currentfill}%
\pgfsetfillopacity{0.160000}%
\pgfsetlinewidth{0.000000pt}%
\definecolor{currentstroke}{rgb}{0.000000,0.000000,0.000000}%
\pgfsetstrokecolor{currentstroke}%
\pgfsetdash{}{0pt}%
\pgfpathmoveto{\pgfqpoint{1.385021in}{0.844902in}}%
\pgfpathlineto{\pgfqpoint{1.395404in}{0.853157in}}%
\pgfpathlineto{\pgfqpoint{1.405825in}{0.861442in}}%
\pgfpathlineto{\pgfqpoint{1.424982in}{0.870606in}}%
\pgfpathlineto{\pgfqpoint{1.442518in}{0.880209in}}%
\pgfpathlineto{\pgfqpoint{1.430949in}{0.871353in}}%
\pgfpathlineto{\pgfqpoint{1.419417in}{0.862526in}}%
\pgfpathlineto{\pgfqpoint{1.402972in}{0.853506in}}%
\pgfpathlineto{\pgfqpoint{1.385021in}{0.844902in}}%
\pgfpathclose%
\pgfusepath{fill}%
\end{pgfscope}%
\begin{pgfscope}%
\pgfpathrectangle{\pgfqpoint{0.407592in}{0.222904in}}{\pgfqpoint{1.411163in}{1.411163in}}%
\pgfusepath{clip}%
\pgfsetbuttcap%
\pgfsetroundjoin%
\definecolor{currentfill}{rgb}{0.749020,0.772549,0.792157}%
\pgfsetfillcolor{currentfill}%
\pgfsetfillopacity{0.160000}%
\pgfsetlinewidth{0.000000pt}%
\definecolor{currentstroke}{rgb}{0.000000,0.000000,0.000000}%
\pgfsetstrokecolor{currentstroke}%
\pgfsetdash{}{0pt}%
\pgfpathmoveto{\pgfqpoint{1.216725in}{0.763549in}}%
\pgfpathlineto{\pgfqpoint{1.222289in}{0.770280in}}%
\pgfpathlineto{\pgfqpoint{1.227878in}{0.777040in}}%
\pgfpathlineto{\pgfqpoint{1.249487in}{0.781578in}}%
\pgfpathlineto{\pgfqpoint{1.270310in}{0.786683in}}%
\pgfpathlineto{\pgfqpoint{1.263018in}{0.779550in}}%
\pgfpathlineto{\pgfqpoint{1.255757in}{0.772447in}}%
\pgfpathlineto{\pgfqpoint{1.236599in}{0.767736in}}%
\pgfpathlineto{\pgfqpoint{1.216725in}{0.763549in}}%
\pgfpathclose%
\pgfusepath{fill}%
\end{pgfscope}%
\begin{pgfscope}%
\pgfpathrectangle{\pgfqpoint{0.407592in}{0.222904in}}{\pgfqpoint{1.411163in}{1.411163in}}%
\pgfusepath{clip}%
\pgfsetbuttcap%
\pgfsetroundjoin%
\definecolor{currentfill}{rgb}{0.749020,0.772549,0.792157}%
\pgfsetfillcolor{currentfill}%
\pgfsetfillopacity{0.160000}%
\pgfsetlinewidth{0.000000pt}%
\definecolor{currentstroke}{rgb}{0.000000,0.000000,0.000000}%
\pgfsetstrokecolor{currentstroke}%
\pgfsetdash{}{0pt}%
\pgfpathmoveto{\pgfqpoint{1.008549in}{0.741349in}}%
\pgfpathlineto{\pgfqpoint{1.004636in}{0.747616in}}%
\pgfpathlineto{\pgfqpoint{1.000705in}{0.753911in}}%
\pgfpathlineto{\pgfqpoint{1.022131in}{0.752134in}}%
\pgfpathlineto{\pgfqpoint{1.043882in}{0.750965in}}%
\pgfpathlineto{\pgfqpoint{1.045919in}{0.744794in}}%
\pgfpathlineto{\pgfqpoint{1.047946in}{0.738652in}}%
\pgfpathlineto{\pgfqpoint{1.028100in}{0.739723in}}%
\pgfpathlineto{\pgfqpoint{1.008549in}{0.741349in}}%
\pgfpathclose%
\pgfusepath{fill}%
\end{pgfscope}%
\begin{pgfscope}%
\pgfpathrectangle{\pgfqpoint{0.407592in}{0.222904in}}{\pgfqpoint{1.411163in}{1.411163in}}%
\pgfusepath{clip}%
\pgfsetbuttcap%
\pgfsetroundjoin%
\definecolor{currentfill}{rgb}{0.749020,0.772549,0.792157}%
\pgfsetfillcolor{currentfill}%
\pgfsetfillopacity{0.160000}%
\pgfsetlinewidth{0.000000pt}%
\definecolor{currentstroke}{rgb}{0.000000,0.000000,0.000000}%
\pgfsetstrokecolor{currentstroke}%
\pgfsetdash{}{0pt}%
\pgfpathmoveto{\pgfqpoint{0.526239in}{1.035773in}}%
\pgfpathlineto{\pgfqpoint{0.512144in}{1.046401in}}%
\pgfpathlineto{\pgfqpoint{0.498016in}{1.057053in}}%
\pgfpathlineto{\pgfqpoint{0.506254in}{1.042282in}}%
\pgfpathlineto{\pgfqpoint{0.517197in}{1.027733in}}%
\pgfpathlineto{\pgfqpoint{0.530992in}{1.017690in}}%
\pgfpathlineto{\pgfqpoint{0.544749in}{1.007673in}}%
\pgfpathlineto{\pgfqpoint{0.534210in}{1.021614in}}%
\pgfpathlineto{\pgfqpoint{0.526239in}{1.035773in}}%
\pgfpathclose%
\pgfusepath{fill}%
\end{pgfscope}%
\begin{pgfscope}%
\pgfpathrectangle{\pgfqpoint{0.407592in}{0.222904in}}{\pgfqpoint{1.411163in}{1.411163in}}%
\pgfusepath{clip}%
\pgfsetbuttcap%
\pgfsetroundjoin%
\definecolor{currentfill}{rgb}{0.749020,0.772549,0.792157}%
\pgfsetfillcolor{currentfill}%
\pgfsetfillopacity{0.160000}%
\pgfsetlinewidth{0.000000pt}%
\definecolor{currentstroke}{rgb}{0.000000,0.000000,0.000000}%
\pgfsetstrokecolor{currentstroke}%
\pgfsetdash{}{0pt}%
\pgfpathmoveto{\pgfqpoint{0.626505in}{0.948151in}}%
\pgfpathlineto{\pgfqpoint{0.612972in}{0.958003in}}%
\pgfpathlineto{\pgfqpoint{0.599402in}{0.967883in}}%
\pgfpathlineto{\pgfqpoint{0.611388in}{0.955450in}}%
\pgfpathlineto{\pgfqpoint{0.625626in}{0.943352in}}%
\pgfpathlineto{\pgfqpoint{0.638566in}{0.934083in}}%
\pgfpathlineto{\pgfqpoint{0.651466in}{0.924842in}}%
\pgfpathlineto{\pgfqpoint{0.637925in}{0.936335in}}%
\pgfpathlineto{\pgfqpoint{0.626505in}{0.948151in}}%
\pgfpathclose%
\pgfusepath{fill}%
\end{pgfscope}%
\begin{pgfscope}%
\pgfpathrectangle{\pgfqpoint{0.407592in}{0.222904in}}{\pgfqpoint{1.411163in}{1.411163in}}%
\pgfusepath{clip}%
\pgfsetbuttcap%
\pgfsetroundjoin%
\definecolor{currentfill}{rgb}{0.749020,0.772549,0.792157}%
\pgfsetfillcolor{currentfill}%
\pgfsetfillopacity{0.160000}%
\pgfsetlinewidth{0.000000pt}%
\definecolor{currentstroke}{rgb}{0.000000,0.000000,0.000000}%
\pgfsetstrokecolor{currentstroke}%
\pgfsetdash{}{0pt}%
\pgfpathmoveto{\pgfqpoint{1.777633in}{1.146626in}}%
\pgfpathlineto{\pgfqpoint{1.791655in}{1.157562in}}%
\pgfpathlineto{\pgfqpoint{1.805709in}{1.168522in}}%
\pgfpathlineto{\pgfqpoint{1.815467in}{1.186405in}}%
\pgfpathlineto{\pgfqpoint{1.821891in}{1.204326in}}%
\pgfpathlineto{\pgfqpoint{1.807686in}{1.192798in}}%
\pgfpathlineto{\pgfqpoint{1.793508in}{1.181291in}}%
\pgfpathlineto{\pgfqpoint{1.787171in}{1.163937in}}%
\pgfpathlineto{\pgfqpoint{1.777633in}{1.146626in}}%
\pgfpathclose%
\pgfusepath{fill}%
\end{pgfscope}%
\begin{pgfscope}%
\pgfpathrectangle{\pgfqpoint{0.407592in}{0.222904in}}{\pgfqpoint{1.411163in}{1.411163in}}%
\pgfusepath{clip}%
\pgfsetbuttcap%
\pgfsetroundjoin%
\definecolor{currentfill}{rgb}{0.749020,0.772549,0.792157}%
\pgfsetfillcolor{currentfill}%
\pgfsetfillopacity{0.160000}%
\pgfsetlinewidth{0.000000pt}%
\definecolor{currentstroke}{rgb}{0.000000,0.000000,0.000000}%
\pgfsetstrokecolor{currentstroke}%
\pgfsetdash{}{0pt}%
\pgfpathmoveto{\pgfqpoint{1.088135in}{0.738201in}}%
\pgfpathlineto{\pgfqpoint{1.088034in}{0.744322in}}%
\pgfpathlineto{\pgfqpoint{1.087933in}{0.750472in}}%
\pgfpathlineto{\pgfqpoint{1.110017in}{0.751151in}}%
\pgfpathlineto{\pgfqpoint{1.131996in}{0.752442in}}%
\pgfpathlineto{\pgfqpoint{1.130161in}{0.746209in}}%
\pgfpathlineto{\pgfqpoint{1.128335in}{0.740005in}}%
\pgfpathlineto{\pgfqpoint{1.108283in}{0.738822in}}%
\pgfpathlineto{\pgfqpoint{1.088135in}{0.738201in}}%
\pgfpathclose%
\pgfusepath{fill}%
\end{pgfscope}%
\begin{pgfscope}%
\pgfpathrectangle{\pgfqpoint{0.407592in}{0.222904in}}{\pgfqpoint{1.411163in}{1.411163in}}%
\pgfusepath{clip}%
\pgfsetbuttcap%
\pgfsetroundjoin%
\definecolor{currentfill}{rgb}{0.749020,0.772549,0.792157}%
\pgfsetfillcolor{currentfill}%
\pgfsetfillopacity{0.160000}%
\pgfsetlinewidth{0.000000pt}%
\definecolor{currentstroke}{rgb}{0.000000,0.000000,0.000000}%
\pgfsetstrokecolor{currentstroke}%
\pgfsetdash{}{0pt}%
\pgfpathmoveto{\pgfqpoint{1.047946in}{0.738652in}}%
\pgfpathlineto{\pgfqpoint{1.045919in}{0.744794in}}%
\pgfpathlineto{\pgfqpoint{1.043882in}{0.750965in}}%
\pgfpathlineto{\pgfqpoint{1.065852in}{0.750410in}}%
\pgfpathlineto{\pgfqpoint{1.087933in}{0.750472in}}%
\pgfpathlineto{\pgfqpoint{1.088034in}{0.744322in}}%
\pgfpathlineto{\pgfqpoint{1.088135in}{0.738201in}}%
\pgfpathlineto{\pgfqpoint{1.067990in}{0.738144in}}%
\pgfpathlineto{\pgfqpoint{1.047946in}{0.738652in}}%
\pgfpathclose%
\pgfusepath{fill}%
\end{pgfscope}%
\begin{pgfscope}%
\pgfpathrectangle{\pgfqpoint{0.407592in}{0.222904in}}{\pgfqpoint{1.411163in}{1.411163in}}%
\pgfusepath{clip}%
\pgfsetbuttcap%
\pgfsetroundjoin%
\definecolor{currentfill}{rgb}{0.749020,0.772549,0.792157}%
\pgfsetfillcolor{currentfill}%
\pgfsetfillopacity{0.160000}%
\pgfsetlinewidth{0.000000pt}%
\definecolor{currentstroke}{rgb}{0.000000,0.000000,0.000000}%
\pgfsetstrokecolor{currentstroke}%
\pgfsetdash{}{0pt}%
\pgfpathmoveto{\pgfqpoint{0.702659in}{0.888170in}}%
\pgfpathlineto{\pgfqpoint{0.689921in}{0.897295in}}%
\pgfpathlineto{\pgfqpoint{0.677143in}{0.906448in}}%
\pgfpathlineto{\pgfqpoint{0.691925in}{0.895932in}}%
\pgfpathlineto{\pgfqpoint{0.708603in}{0.885836in}}%
\pgfpathlineto{\pgfqpoint{0.720480in}{0.877268in}}%
\pgfpathlineto{\pgfqpoint{0.732315in}{0.868730in}}%
\pgfpathlineto{\pgfqpoint{0.716603in}{0.878250in}}%
\pgfpathlineto{\pgfqpoint{0.702659in}{0.888170in}}%
\pgfpathclose%
\pgfusepath{fill}%
\end{pgfscope}%
\begin{pgfscope}%
\pgfpathrectangle{\pgfqpoint{0.407592in}{0.222904in}}{\pgfqpoint{1.411163in}{1.411163in}}%
\pgfusepath{clip}%
\pgfsetbuttcap%
\pgfsetroundjoin%
\definecolor{currentfill}{rgb}{0.749020,0.772549,0.792157}%
\pgfsetfillcolor{currentfill}%
\pgfsetfillopacity{0.160000}%
\pgfsetlinewidth{0.000000pt}%
\definecolor{currentstroke}{rgb}{0.000000,0.000000,0.000000}%
\pgfsetstrokecolor{currentstroke}%
\pgfsetdash{}{0pt}%
\pgfpathmoveto{\pgfqpoint{1.326989in}{0.813739in}}%
\pgfpathlineto{\pgfqpoint{1.335931in}{0.821415in}}%
\pgfpathlineto{\pgfqpoint{1.344910in}{0.829121in}}%
\pgfpathlineto{\pgfqpoint{1.365639in}{0.836759in}}%
\pgfpathlineto{\pgfqpoint{1.385021in}{0.844902in}}%
\pgfpathlineto{\pgfqpoint{1.374675in}{0.836677in}}%
\pgfpathlineto{\pgfqpoint{1.364368in}{0.828482in}}%
\pgfpathlineto{\pgfqpoint{1.346300in}{0.820874in}}%
\pgfpathlineto{\pgfqpoint{1.326989in}{0.813739in}}%
\pgfpathclose%
\pgfusepath{fill}%
\end{pgfscope}%
\begin{pgfscope}%
\pgfpathrectangle{\pgfqpoint{0.407592in}{0.222904in}}{\pgfqpoint{1.411163in}{1.411163in}}%
\pgfusepath{clip}%
\pgfsetbuttcap%
\pgfsetroundjoin%
\definecolor{currentfill}{rgb}{0.749020,0.772549,0.792157}%
\pgfsetfillcolor{currentfill}%
\pgfsetfillopacity{0.160000}%
\pgfsetlinewidth{0.000000pt}%
\definecolor{currentstroke}{rgb}{0.000000,0.000000,0.000000}%
\pgfsetstrokecolor{currentstroke}%
\pgfsetdash{}{0pt}%
\pgfpathmoveto{\pgfqpoint{1.270310in}{0.786683in}}%
\pgfpathlineto{\pgfqpoint{1.277632in}{0.793845in}}%
\pgfpathlineto{\pgfqpoint{1.284986in}{0.801039in}}%
\pgfpathlineto{\pgfqpoint{1.306520in}{0.807116in}}%
\pgfpathlineto{\pgfqpoint{1.326989in}{0.813739in}}%
\pgfpathlineto{\pgfqpoint{1.318081in}{0.806094in}}%
\pgfpathlineto{\pgfqpoint{1.309209in}{0.798479in}}%
\pgfpathlineto{\pgfqpoint{1.290248in}{0.792326in}}%
\pgfpathlineto{\pgfqpoint{1.270310in}{0.786683in}}%
\pgfpathclose%
\pgfusepath{fill}%
\end{pgfscope}%
\begin{pgfscope}%
\pgfpathrectangle{\pgfqpoint{0.407592in}{0.222904in}}{\pgfqpoint{1.411163in}{1.411163in}}%
\pgfusepath{clip}%
\pgfsetbuttcap%
\pgfsetroundjoin%
\definecolor{currentfill}{rgb}{0.749020,0.772549,0.792157}%
\pgfsetfillcolor{currentfill}%
\pgfsetfillopacity{0.160000}%
\pgfsetlinewidth{0.000000pt}%
\definecolor{currentstroke}{rgb}{0.000000,0.000000,0.000000}%
\pgfsetstrokecolor{currentstroke}%
\pgfsetdash{}{0pt}%
\pgfpathmoveto{\pgfqpoint{0.375133in}{1.177810in}}%
\pgfpathlineto{\pgfqpoint{0.360700in}{1.189260in}}%
\pgfpathlineto{\pgfqpoint{0.346239in}{1.200731in}}%
\pgfpathlineto{\pgfqpoint{0.348903in}{1.182812in}}%
\pgfpathlineto{\pgfqpoint{0.354919in}{1.164946in}}%
\pgfpathlineto{\pgfqpoint{0.369379in}{1.154043in}}%
\pgfpathlineto{\pgfqpoint{0.383806in}{1.143166in}}%
\pgfpathlineto{\pgfqpoint{0.377860in}{1.160458in}}%
\pgfpathlineto{\pgfqpoint{0.375133in}{1.177810in}}%
\pgfpathclose%
\pgfusepath{fill}%
\end{pgfscope}%
\begin{pgfscope}%
\pgfpathrectangle{\pgfqpoint{0.407592in}{0.222904in}}{\pgfqpoint{1.411163in}{1.411163in}}%
\pgfusepath{clip}%
\pgfsetbuttcap%
\pgfsetroundjoin%
\definecolor{currentfill}{rgb}{0.749020,0.772549,0.792157}%
\pgfsetfillcolor{currentfill}%
\pgfsetfillopacity{0.160000}%
\pgfsetlinewidth{0.000000pt}%
\definecolor{currentstroke}{rgb}{0.000000,0.000000,0.000000}%
\pgfsetstrokecolor{currentstroke}%
\pgfsetdash{}{0pt}%
\pgfpathmoveto{\pgfqpoint{0.920300in}{0.766854in}}%
\pgfpathlineto{\pgfqpoint{0.912833in}{0.773723in}}%
\pgfpathlineto{\pgfqpoint{0.905333in}{0.780622in}}%
\pgfpathlineto{\pgfqpoint{0.926109in}{0.776200in}}%
\pgfpathlineto{\pgfqpoint{0.947680in}{0.772375in}}%
\pgfpathlineto{\pgfqpoint{0.953476in}{0.765795in}}%
\pgfpathlineto{\pgfqpoint{0.959245in}{0.759245in}}%
\pgfpathlineto{\pgfqpoint{0.939410in}{0.762774in}}%
\pgfpathlineto{\pgfqpoint{0.920300in}{0.766854in}}%
\pgfpathclose%
\pgfusepath{fill}%
\end{pgfscope}%
\begin{pgfscope}%
\pgfpathrectangle{\pgfqpoint{0.407592in}{0.222904in}}{\pgfqpoint{1.411163in}{1.411163in}}%
\pgfusepath{clip}%
\pgfsetbuttcap%
\pgfsetroundjoin%
\definecolor{currentfill}{rgb}{0.749020,0.772549,0.792157}%
\pgfsetfillcolor{currentfill}%
\pgfsetfillopacity{0.160000}%
\pgfsetlinewidth{0.000000pt}%
\definecolor{currentstroke}{rgb}{0.000000,0.000000,0.000000}%
\pgfsetstrokecolor{currentstroke}%
\pgfsetdash{}{0pt}%
\pgfpathmoveto{\pgfqpoint{1.465771in}{0.898010in}}%
\pgfpathlineto{\pgfqpoint{1.477455in}{0.906954in}}%
\pgfpathlineto{\pgfqpoint{1.489178in}{0.915928in}}%
\pgfpathlineto{\pgfqpoint{1.506955in}{0.927119in}}%
\pgfpathlineto{\pgfqpoint{1.494750in}{0.917847in}}%
\pgfpathlineto{\pgfqpoint{1.482584in}{0.908604in}}%
\pgfusepath{fill}%
\end{pgfscope}%
\begin{pgfscope}%
\pgfpathrectangle{\pgfqpoint{0.407592in}{0.222904in}}{\pgfqpoint{1.411163in}{1.411163in}}%
\pgfusepath{clip}%
\pgfsetbuttcap%
\pgfsetroundjoin%
\definecolor{currentfill}{rgb}{0.749020,0.772549,0.792157}%
\pgfsetfillcolor{currentfill}%
\pgfsetfillopacity{0.160000}%
\pgfsetlinewidth{0.000000pt}%
\definecolor{currentstroke}{rgb}{0.000000,0.000000,0.000000}%
\pgfsetstrokecolor{currentstroke}%
\pgfsetdash{}{0pt}%
\pgfpathmoveto{\pgfqpoint{1.641536in}{1.030631in}}%
\pgfpathlineto{\pgfqpoint{1.654833in}{1.040763in}}%
\pgfpathlineto{\pgfqpoint{1.668165in}{1.050923in}}%
\pgfpathlineto{\pgfqpoint{1.682564in}{1.066125in}}%
\pgfpathlineto{\pgfqpoint{1.694172in}{1.081536in}}%
\pgfpathlineto{\pgfqpoint{1.680372in}{1.070774in}}%
\pgfpathlineto{\pgfqpoint{1.666604in}{1.060037in}}%
\pgfpathlineto{\pgfqpoint{1.655398in}{1.045230in}}%
\pgfpathlineto{\pgfqpoint{1.641536in}{1.030631in}}%
\pgfpathclose%
\pgfusepath{fill}%
\end{pgfscope}%
\begin{pgfscope}%
\pgfpathrectangle{\pgfqpoint{0.407592in}{0.222904in}}{\pgfqpoint{1.411163in}{1.411163in}}%
\pgfusepath{clip}%
\pgfsetbuttcap%
\pgfsetroundjoin%
\definecolor{currentfill}{rgb}{0.749020,0.772549,0.792157}%
\pgfsetfillcolor{currentfill}%
\pgfsetfillopacity{0.160000}%
\pgfsetlinewidth{0.000000pt}%
\definecolor{currentstroke}{rgb}{0.000000,0.000000,0.000000}%
\pgfsetstrokecolor{currentstroke}%
\pgfsetdash{}{0pt}%
\pgfpathmoveto{\pgfqpoint{1.805709in}{1.168522in}}%
\pgfpathlineto{\pgfqpoint{1.819795in}{1.179507in}}%
\pgfpathlineto{\pgfqpoint{1.833914in}{1.190518in}}%
\pgfpathlineto{\pgfqpoint{1.843881in}{1.208967in}}%
\pgfpathlineto{\pgfqpoint{1.850378in}{1.227445in}}%
\pgfpathlineto{\pgfqpoint{1.836121in}{1.215875in}}%
\pgfpathlineto{\pgfqpoint{1.821891in}{1.204326in}}%
\pgfpathlineto{\pgfqpoint{1.815467in}{1.186405in}}%
\pgfpathlineto{\pgfqpoint{1.805709in}{1.168522in}}%
\pgfpathclose%
\pgfusepath{fill}%
\end{pgfscope}%
\begin{pgfscope}%
\pgfpathrectangle{\pgfqpoint{0.407592in}{0.222904in}}{\pgfqpoint{1.411163in}{1.411163in}}%
\pgfusepath{clip}%
\pgfsetbuttcap%
\pgfsetroundjoin%
\definecolor{currentfill}{rgb}{0.749020,0.772549,0.792157}%
\pgfsetfillcolor{currentfill}%
\pgfsetfillopacity{0.160000}%
\pgfsetlinewidth{0.000000pt}%
\definecolor{currentstroke}{rgb}{0.000000,0.000000,0.000000}%
\pgfsetstrokecolor{currentstroke}%
\pgfsetdash{}{0pt}%
\pgfpathmoveto{\pgfqpoint{1.175206in}{0.756833in}}%
\pgfpathlineto{\pgfqpoint{1.178951in}{0.763281in}}%
\pgfpathlineto{\pgfqpoint{1.182714in}{0.769759in}}%
\pgfpathlineto{\pgfqpoint{1.205584in}{0.773094in}}%
\pgfpathlineto{\pgfqpoint{1.227878in}{0.777040in}}%
\pgfpathlineto{\pgfqpoint{1.222289in}{0.770280in}}%
\pgfpathlineto{\pgfqpoint{1.216725in}{0.763549in}}%
\pgfpathlineto{\pgfqpoint{1.196228in}{0.759908in}}%
\pgfpathlineto{\pgfqpoint{1.175206in}{0.756833in}}%
\pgfpathclose%
\pgfusepath{fill}%
\end{pgfscope}%
\begin{pgfscope}%
\pgfpathrectangle{\pgfqpoint{0.407592in}{0.222904in}}{\pgfqpoint{1.411163in}{1.411163in}}%
\pgfusepath{clip}%
\pgfsetbuttcap%
\pgfsetroundjoin%
\definecolor{currentfill}{rgb}{0.749020,0.772549,0.792157}%
\pgfsetfillcolor{currentfill}%
\pgfsetfillopacity{0.160000}%
\pgfsetlinewidth{0.000000pt}%
\definecolor{currentstroke}{rgb}{0.000000,0.000000,0.000000}%
\pgfsetstrokecolor{currentstroke}%
\pgfsetdash{}{0pt}%
\pgfpathmoveto{\pgfqpoint{1.531478in}{0.945749in}}%
\pgfpathlineto{\pgfqpoint{1.543797in}{0.955108in}}%
\pgfpathlineto{\pgfqpoint{1.556155in}{0.964496in}}%
\pgfpathlineto{\pgfqpoint{1.573590in}{0.977270in}}%
\pgfpathlineto{\pgfqpoint{1.588709in}{0.990377in}}%
\pgfpathlineto{\pgfqpoint{1.575592in}{0.980381in}}%
\pgfpathlineto{\pgfqpoint{1.562509in}{0.970412in}}%
\pgfpathlineto{\pgfqpoint{1.548086in}{0.957919in}}%
\pgfpathlineto{\pgfqpoint{1.531478in}{0.945749in}}%
\pgfpathclose%
\pgfusepath{fill}%
\end{pgfscope}%
\begin{pgfscope}%
\pgfpathrectangle{\pgfqpoint{0.407592in}{0.222904in}}{\pgfqpoint{1.411163in}{1.411163in}}%
\pgfusepath{clip}%
\pgfsetbuttcap%
\pgfsetroundjoin%
\definecolor{currentfill}{rgb}{0.749020,0.772549,0.792157}%
\pgfsetfillcolor{currentfill}%
\pgfsetfillopacity{0.160000}%
\pgfsetlinewidth{0.000000pt}%
\definecolor{currentstroke}{rgb}{0.000000,0.000000,0.000000}%
\pgfsetstrokecolor{currentstroke}%
\pgfsetdash{}{0pt}%
\pgfpathmoveto{\pgfqpoint{0.755860in}{0.851745in}}%
\pgfpathlineto{\pgfqpoint{0.744108in}{0.860222in}}%
\pgfpathlineto{\pgfqpoint{0.732315in}{0.868730in}}%
\pgfpathlineto{\pgfqpoint{0.749735in}{0.859661in}}%
\pgfpathlineto{\pgfqpoint{0.768792in}{0.851089in}}%
\pgfpathlineto{\pgfqpoint{0.779438in}{0.843121in}}%
\pgfpathlineto{\pgfqpoint{0.790043in}{0.835185in}}%
\pgfpathlineto{\pgfqpoint{0.772192in}{0.843230in}}%
\pgfpathlineto{\pgfqpoint{0.755860in}{0.851745in}}%
\pgfpathclose%
\pgfusepath{fill}%
\end{pgfscope}%
\begin{pgfscope}%
\pgfpathrectangle{\pgfqpoint{0.407592in}{0.222904in}}{\pgfqpoint{1.411163in}{1.411163in}}%
\pgfusepath{clip}%
\pgfsetbuttcap%
\pgfsetroundjoin%
\definecolor{currentfill}{rgb}{0.749020,0.772549,0.792157}%
\pgfsetfillcolor{currentfill}%
\pgfsetfillopacity{0.160000}%
\pgfsetlinewidth{0.000000pt}%
\definecolor{currentstroke}{rgb}{0.000000,0.000000,0.000000}%
\pgfsetstrokecolor{currentstroke}%
\pgfsetdash{}{0pt}%
\pgfpathmoveto{\pgfqpoint{0.498016in}{1.057053in}}%
\pgfpathlineto{\pgfqpoint{0.483856in}{1.067730in}}%
\pgfpathlineto{\pgfqpoint{0.469662in}{1.078431in}}%
\pgfpathlineto{\pgfqpoint{0.478156in}{1.063056in}}%
\pgfpathlineto{\pgfqpoint{0.489491in}{1.047904in}}%
\pgfpathlineto{\pgfqpoint{0.503363in}{1.037805in}}%
\pgfpathlineto{\pgfqpoint{0.517197in}{1.027733in}}%
\pgfpathlineto{\pgfqpoint{0.506254in}{1.042282in}}%
\pgfpathlineto{\pgfqpoint{0.498016in}{1.057053in}}%
\pgfpathclose%
\pgfusepath{fill}%
\end{pgfscope}%
\begin{pgfscope}%
\pgfpathrectangle{\pgfqpoint{0.407592in}{0.222904in}}{\pgfqpoint{1.411163in}{1.411163in}}%
\pgfusepath{clip}%
\pgfsetbuttcap%
\pgfsetroundjoin%
\definecolor{currentfill}{rgb}{0.749020,0.772549,0.792157}%
\pgfsetfillcolor{currentfill}%
\pgfsetfillopacity{0.160000}%
\pgfsetlinewidth{0.000000pt}%
\definecolor{currentstroke}{rgb}{0.000000,0.000000,0.000000}%
\pgfsetstrokecolor{currentstroke}%
\pgfsetdash{}{0pt}%
\pgfpathmoveto{\pgfqpoint{1.405825in}{0.861442in}}%
\pgfpathlineto{\pgfqpoint{1.416284in}{0.869757in}}%
\pgfpathlineto{\pgfqpoint{1.426782in}{0.878103in}}%
\pgfpathlineto{\pgfqpoint{1.447146in}{0.887825in}}%
\pgfpathlineto{\pgfqpoint{1.465771in}{0.898010in}}%
\pgfpathlineto{\pgfqpoint{1.454126in}{0.889095in}}%
\pgfpathlineto{\pgfqpoint{1.442518in}{0.880209in}}%
\pgfpathlineto{\pgfqpoint{1.424982in}{0.870606in}}%
\pgfpathlineto{\pgfqpoint{1.405825in}{0.861442in}}%
\pgfpathclose%
\pgfusepath{fill}%
\end{pgfscope}%
\begin{pgfscope}%
\pgfpathrectangle{\pgfqpoint{0.407592in}{0.222904in}}{\pgfqpoint{1.411163in}{1.411163in}}%
\pgfusepath{clip}%
\pgfsetbuttcap%
\pgfsetroundjoin%
\definecolor{currentfill}{rgb}{0.749020,0.772549,0.792157}%
\pgfsetfillcolor{currentfill}%
\pgfsetfillopacity{0.160000}%
\pgfsetlinewidth{0.000000pt}%
\definecolor{currentstroke}{rgb}{0.000000,0.000000,0.000000}%
\pgfsetstrokecolor{currentstroke}%
\pgfsetdash{}{0pt}%
\pgfpathmoveto{\pgfqpoint{0.866553in}{0.791152in}}%
\pgfpathlineto{\pgfqpoint{0.857470in}{0.798487in}}%
\pgfpathlineto{\pgfqpoint{0.848350in}{0.805852in}}%
\pgfpathlineto{\pgfqpoint{0.868755in}{0.799888in}}%
\pgfpathlineto{\pgfqpoint{0.890235in}{0.794510in}}%
\pgfpathlineto{\pgfqpoint{0.897801in}{0.787551in}}%
\pgfpathlineto{\pgfqpoint{0.905333in}{0.780622in}}%
\pgfpathlineto{\pgfqpoint{0.885450in}{0.785614in}}%
\pgfpathlineto{\pgfqpoint{0.866553in}{0.791152in}}%
\pgfpathclose%
\pgfusepath{fill}%
\end{pgfscope}%
\begin{pgfscope}%
\pgfpathrectangle{\pgfqpoint{0.407592in}{0.222904in}}{\pgfqpoint{1.411163in}{1.411163in}}%
\pgfusepath{clip}%
\pgfsetbuttcap%
\pgfsetroundjoin%
\definecolor{currentfill}{rgb}{0.749020,0.772549,0.792157}%
\pgfsetfillcolor{currentfill}%
\pgfsetfillopacity{0.160000}%
\pgfsetlinewidth{0.000000pt}%
\definecolor{currentstroke}{rgb}{0.000000,0.000000,0.000000}%
\pgfsetstrokecolor{currentstroke}%
\pgfsetdash{}{0pt}%
\pgfpathmoveto{\pgfqpoint{0.599402in}{0.967883in}}%
\pgfpathlineto{\pgfqpoint{0.585795in}{0.977790in}}%
\pgfpathlineto{\pgfqpoint{0.572151in}{0.987724in}}%
\pgfpathlineto{\pgfqpoint{0.584694in}{0.974678in}}%
\pgfpathlineto{\pgfqpoint{0.599623in}{0.961980in}}%
\pgfpathlineto{\pgfqpoint{0.612645in}{0.952651in}}%
\pgfpathlineto{\pgfqpoint{0.625626in}{0.943352in}}%
\pgfpathlineto{\pgfqpoint{0.611388in}{0.955450in}}%
\pgfpathlineto{\pgfqpoint{0.599402in}{0.967883in}}%
\pgfpathclose%
\pgfusepath{fill}%
\end{pgfscope}%
\begin{pgfscope}%
\pgfpathrectangle{\pgfqpoint{0.407592in}{0.222904in}}{\pgfqpoint{1.411163in}{1.411163in}}%
\pgfusepath{clip}%
\pgfsetbuttcap%
\pgfsetroundjoin%
\definecolor{currentfill}{rgb}{0.749020,0.772549,0.792157}%
\pgfsetfillcolor{currentfill}%
\pgfsetfillopacity{0.160000}%
\pgfsetlinewidth{0.000000pt}%
\definecolor{currentstroke}{rgb}{0.000000,0.000000,0.000000}%
\pgfsetstrokecolor{currentstroke}%
\pgfsetdash{}{0pt}%
\pgfpathmoveto{\pgfqpoint{0.959245in}{0.759245in}}%
\pgfpathlineto{\pgfqpoint{0.953476in}{0.765795in}}%
\pgfpathlineto{\pgfqpoint{0.947680in}{0.772375in}}%
\pgfpathlineto{\pgfqpoint{0.969942in}{0.769166in}}%
\pgfpathlineto{\pgfqpoint{0.992788in}{0.766591in}}%
\pgfpathlineto{\pgfqpoint{0.996756in}{0.760236in}}%
\pgfpathlineto{\pgfqpoint{1.000705in}{0.753911in}}%
\pgfpathlineto{\pgfqpoint{0.979709in}{0.756286in}}%
\pgfpathlineto{\pgfqpoint{0.959245in}{0.759245in}}%
\pgfpathclose%
\pgfusepath{fill}%
\end{pgfscope}%
\begin{pgfscope}%
\pgfpathrectangle{\pgfqpoint{0.407592in}{0.222904in}}{\pgfqpoint{1.411163in}{1.411163in}}%
\pgfusepath{clip}%
\pgfsetbuttcap%
\pgfsetroundjoin%
\definecolor{currentfill}{rgb}{0.749020,0.772549,0.792157}%
\pgfsetfillcolor{currentfill}%
\pgfsetfillopacity{0.160000}%
\pgfsetlinewidth{0.000000pt}%
\definecolor{currentstroke}{rgb}{0.000000,0.000000,0.000000}%
\pgfsetstrokecolor{currentstroke}%
\pgfsetdash{}{0pt}%
\pgfpathmoveto{\pgfqpoint{0.346239in}{1.200731in}}%
\pgfpathlineto{\pgfqpoint{0.331751in}{1.212225in}}%
\pgfpathlineto{\pgfqpoint{0.317235in}{1.223739in}}%
\pgfpathlineto{\pgfqpoint{0.319823in}{1.205261in}}%
\pgfpathlineto{\pgfqpoint{0.325897in}{1.186828in}}%
\pgfpathlineto{\pgfqpoint{0.340425in}{1.175874in}}%
\pgfpathlineto{\pgfqpoint{0.354919in}{1.164946in}}%
\pgfpathlineto{\pgfqpoint{0.348903in}{1.182812in}}%
\pgfpathlineto{\pgfqpoint{0.346239in}{1.200731in}}%
\pgfpathclose%
\pgfusepath{fill}%
\end{pgfscope}%
\begin{pgfscope}%
\pgfpathrectangle{\pgfqpoint{0.407592in}{0.222904in}}{\pgfqpoint{1.411163in}{1.411163in}}%
\pgfusepath{clip}%
\pgfsetbuttcap%
\pgfsetroundjoin%
\definecolor{currentfill}{rgb}{0.749020,0.772549,0.792157}%
\pgfsetfillcolor{currentfill}%
\pgfsetfillopacity{0.160000}%
\pgfsetlinewidth{0.000000pt}%
\definecolor{currentstroke}{rgb}{0.000000,0.000000,0.000000}%
\pgfsetstrokecolor{currentstroke}%
\pgfsetdash{}{0pt}%
\pgfpathmoveto{\pgfqpoint{0.811131in}{0.819403in}}%
\pgfpathlineto{\pgfqpoint{0.800608in}{0.827279in}}%
\pgfpathlineto{\pgfqpoint{0.790043in}{0.835185in}}%
\pgfpathlineto{\pgfqpoint{0.809339in}{0.827653in}}%
\pgfpathlineto{\pgfqpoint{0.829995in}{0.820673in}}%
\pgfpathlineto{\pgfqpoint{0.839191in}{0.813247in}}%
\pgfpathlineto{\pgfqpoint{0.848350in}{0.805852in}}%
\pgfpathlineto{\pgfqpoint{0.829113in}{0.812368in}}%
\pgfpathlineto{\pgfqpoint{0.811131in}{0.819403in}}%
\pgfpathclose%
\pgfusepath{fill}%
\end{pgfscope}%
\begin{pgfscope}%
\pgfpathrectangle{\pgfqpoint{0.407592in}{0.222904in}}{\pgfqpoint{1.411163in}{1.411163in}}%
\pgfusepath{clip}%
\pgfsetbuttcap%
\pgfsetroundjoin%
\definecolor{currentfill}{rgb}{0.749020,0.772549,0.792157}%
\pgfsetfillcolor{currentfill}%
\pgfsetfillopacity{0.160000}%
\pgfsetlinewidth{0.000000pt}%
\definecolor{currentstroke}{rgb}{0.000000,0.000000,0.000000}%
\pgfsetstrokecolor{currentstroke}%
\pgfsetdash{}{0pt}%
\pgfpathmoveto{\pgfqpoint{1.833914in}{1.190518in}}%
\pgfpathlineto{\pgfqpoint{1.848065in}{1.201554in}}%
\pgfpathlineto{\pgfqpoint{1.858131in}{1.220283in}}%
\pgfpathlineto{\pgfqpoint{1.864660in}{1.239036in}}%
\pgfpathlineto{\pgfqpoint{1.850378in}{1.227445in}}%
\pgfpathlineto{\pgfqpoint{1.843881in}{1.208967in}}%
\pgfusepath{fill}%
\end{pgfscope}%
\begin{pgfscope}%
\pgfpathrectangle{\pgfqpoint{0.407592in}{0.222904in}}{\pgfqpoint{1.411163in}{1.411163in}}%
\pgfusepath{clip}%
\pgfsetbuttcap%
\pgfsetroundjoin%
\definecolor{currentfill}{rgb}{0.749020,0.772549,0.792157}%
\pgfsetfillcolor{currentfill}%
\pgfsetfillopacity{0.160000}%
\pgfsetlinewidth{0.000000pt}%
\definecolor{currentstroke}{rgb}{0.000000,0.000000,0.000000}%
\pgfsetstrokecolor{currentstroke}%
\pgfsetdash{}{0pt}%
\pgfpathmoveto{\pgfqpoint{0.677143in}{0.906448in}}%
\pgfpathlineto{\pgfqpoint{0.664324in}{0.915631in}}%
\pgfpathlineto{\pgfqpoint{0.651466in}{0.924842in}}%
\pgfpathlineto{\pgfqpoint{0.667083in}{0.913732in}}%
\pgfpathlineto{\pgfqpoint{0.684723in}{0.903063in}}%
\pgfpathlineto{\pgfqpoint{0.696684in}{0.894434in}}%
\pgfpathlineto{\pgfqpoint{0.708603in}{0.885836in}}%
\pgfpathlineto{\pgfqpoint{0.691925in}{0.895932in}}%
\pgfpathlineto{\pgfqpoint{0.677143in}{0.906448in}}%
\pgfpathclose%
\pgfusepath{fill}%
\end{pgfscope}%
\begin{pgfscope}%
\pgfpathrectangle{\pgfqpoint{0.407592in}{0.222904in}}{\pgfqpoint{1.411163in}{1.411163in}}%
\pgfusepath{clip}%
\pgfsetbuttcap%
\pgfsetroundjoin%
\definecolor{currentfill}{rgb}{0.749020,0.772549,0.792157}%
\pgfsetfillcolor{currentfill}%
\pgfsetfillopacity{0.160000}%
\pgfsetlinewidth{0.000000pt}%
\definecolor{currentstroke}{rgb}{0.000000,0.000000,0.000000}%
\pgfsetstrokecolor{currentstroke}%
\pgfsetdash{}{0pt}%
\pgfpathmoveto{\pgfqpoint{1.131996in}{0.752442in}}%
\pgfpathlineto{\pgfqpoint{1.133839in}{0.758705in}}%
\pgfpathlineto{\pgfqpoint{1.135691in}{0.764998in}}%
\pgfpathlineto{\pgfqpoint{1.159378in}{0.767056in}}%
\pgfpathlineto{\pgfqpoint{1.182714in}{0.769759in}}%
\pgfpathlineto{\pgfqpoint{1.178951in}{0.763281in}}%
\pgfpathlineto{\pgfqpoint{1.175206in}{0.756833in}}%
\pgfpathlineto{\pgfqpoint{1.153761in}{0.754340in}}%
\pgfpathlineto{\pgfqpoint{1.131996in}{0.752442in}}%
\pgfpathclose%
\pgfusepath{fill}%
\end{pgfscope}%
\begin{pgfscope}%
\pgfpathrectangle{\pgfqpoint{0.407592in}{0.222904in}}{\pgfqpoint{1.411163in}{1.411163in}}%
\pgfusepath{clip}%
\pgfsetbuttcap%
\pgfsetroundjoin%
\definecolor{currentfill}{rgb}{0.749020,0.772549,0.792157}%
\pgfsetfillcolor{currentfill}%
\pgfsetfillopacity{0.160000}%
\pgfsetlinewidth{0.000000pt}%
\definecolor{currentstroke}{rgb}{0.000000,0.000000,0.000000}%
\pgfsetstrokecolor{currentstroke}%
\pgfsetdash{}{0pt}%
\pgfpathmoveto{\pgfqpoint{1.668165in}{1.050923in}}%
\pgfpathlineto{\pgfqpoint{1.681534in}{1.061110in}}%
\pgfpathlineto{\pgfqpoint{1.694940in}{1.071325in}}%
\pgfpathlineto{\pgfqpoint{1.709867in}{1.087125in}}%
\pgfpathlineto{\pgfqpoint{1.721866in}{1.103134in}}%
\pgfpathlineto{\pgfqpoint{1.708003in}{1.092323in}}%
\pgfpathlineto{\pgfqpoint{1.694172in}{1.081536in}}%
\pgfpathlineto{\pgfqpoint{1.682564in}{1.066125in}}%
\pgfpathlineto{\pgfqpoint{1.668165in}{1.050923in}}%
\pgfpathclose%
\pgfusepath{fill}%
\end{pgfscope}%
\begin{pgfscope}%
\pgfpathrectangle{\pgfqpoint{0.407592in}{0.222904in}}{\pgfqpoint{1.411163in}{1.411163in}}%
\pgfusepath{clip}%
\pgfsetbuttcap%
\pgfsetroundjoin%
\definecolor{currentfill}{rgb}{0.749020,0.772549,0.792157}%
\pgfsetfillcolor{currentfill}%
\pgfsetfillopacity{0.160000}%
\pgfsetlinewidth{0.000000pt}%
\definecolor{currentstroke}{rgb}{0.000000,0.000000,0.000000}%
\pgfsetstrokecolor{currentstroke}%
\pgfsetdash{}{0pt}%
\pgfpathmoveto{\pgfqpoint{1.227878in}{0.777040in}}%
\pgfpathlineto{\pgfqpoint{1.233491in}{0.783831in}}%
\pgfpathlineto{\pgfqpoint{1.239130in}{0.790652in}}%
\pgfpathlineto{\pgfqpoint{1.262487in}{0.795541in}}%
\pgfpathlineto{\pgfqpoint{1.284986in}{0.801039in}}%
\pgfpathlineto{\pgfqpoint{1.277632in}{0.793845in}}%
\pgfpathlineto{\pgfqpoint{1.270310in}{0.786683in}}%
\pgfpathlineto{\pgfqpoint{1.249487in}{0.781578in}}%
\pgfpathlineto{\pgfqpoint{1.227878in}{0.777040in}}%
\pgfpathclose%
\pgfusepath{fill}%
\end{pgfscope}%
\begin{pgfscope}%
\pgfpathrectangle{\pgfqpoint{0.407592in}{0.222904in}}{\pgfqpoint{1.411163in}{1.411163in}}%
\pgfusepath{clip}%
\pgfsetbuttcap%
\pgfsetroundjoin%
\definecolor{currentfill}{rgb}{0.749020,0.772549,0.792157}%
\pgfsetfillcolor{currentfill}%
\pgfsetfillopacity{0.160000}%
\pgfsetlinewidth{0.000000pt}%
\definecolor{currentstroke}{rgb}{0.000000,0.000000,0.000000}%
\pgfsetstrokecolor{currentstroke}%
\pgfsetdash{}{0pt}%
\pgfpathmoveto{\pgfqpoint{1.000705in}{0.753911in}}%
\pgfpathlineto{\pgfqpoint{0.996756in}{0.760236in}}%
\pgfpathlineto{\pgfqpoint{0.992788in}{0.766591in}}%
\pgfpathlineto{\pgfqpoint{1.016105in}{0.764664in}}%
\pgfpathlineto{\pgfqpoint{1.039780in}{0.763396in}}%
\pgfpathlineto{\pgfqpoint{1.041836in}{0.757166in}}%
\pgfpathlineto{\pgfqpoint{1.043882in}{0.750965in}}%
\pgfpathlineto{\pgfqpoint{1.022131in}{0.752134in}}%
\pgfpathlineto{\pgfqpoint{1.000705in}{0.753911in}}%
\pgfpathclose%
\pgfusepath{fill}%
\end{pgfscope}%
\begin{pgfscope}%
\pgfpathrectangle{\pgfqpoint{0.407592in}{0.222904in}}{\pgfqpoint{1.411163in}{1.411163in}}%
\pgfusepath{clip}%
\pgfsetbuttcap%
\pgfsetroundjoin%
\definecolor{currentfill}{rgb}{0.749020,0.772549,0.792157}%
\pgfsetfillcolor{currentfill}%
\pgfsetfillopacity{0.160000}%
\pgfsetlinewidth{0.000000pt}%
\definecolor{currentstroke}{rgb}{0.000000,0.000000,0.000000}%
\pgfsetstrokecolor{currentstroke}%
\pgfsetdash{}{0pt}%
\pgfpathmoveto{\pgfqpoint{1.344910in}{0.829121in}}%
\pgfpathlineto{\pgfqpoint{1.353924in}{0.836857in}}%
\pgfpathlineto{\pgfqpoint{1.362974in}{0.844625in}}%
\pgfpathlineto{\pgfqpoint{1.385126in}{0.852766in}}%
\pgfpathlineto{\pgfqpoint{1.405825in}{0.861442in}}%
\pgfpathlineto{\pgfqpoint{1.395404in}{0.853157in}}%
\pgfpathlineto{\pgfqpoint{1.385021in}{0.844902in}}%
\pgfpathlineto{\pgfqpoint{1.365639in}{0.836759in}}%
\pgfpathlineto{\pgfqpoint{1.344910in}{0.829121in}}%
\pgfpathclose%
\pgfusepath{fill}%
\end{pgfscope}%
\begin{pgfscope}%
\pgfpathrectangle{\pgfqpoint{0.407592in}{0.222904in}}{\pgfqpoint{1.411163in}{1.411163in}}%
\pgfusepath{clip}%
\pgfsetbuttcap%
\pgfsetroundjoin%
\definecolor{currentfill}{rgb}{0.749020,0.772549,0.792157}%
\pgfsetfillcolor{currentfill}%
\pgfsetfillopacity{0.160000}%
\pgfsetlinewidth{0.000000pt}%
\definecolor{currentstroke}{rgb}{0.000000,0.000000,0.000000}%
\pgfsetstrokecolor{currentstroke}%
\pgfsetdash{}{0pt}%
\pgfpathmoveto{\pgfqpoint{1.489178in}{0.915928in}}%
\pgfpathlineto{\pgfqpoint{1.500940in}{0.924931in}}%
\pgfpathlineto{\pgfqpoint{1.512741in}{0.933965in}}%
\pgfpathlineto{\pgfqpoint{1.531478in}{0.945749in}}%
\pgfpathlineto{\pgfqpoint{1.519197in}{0.936420in}}%
\pgfpathlineto{\pgfqpoint{1.506955in}{0.927119in}}%
\pgfusepath{fill}%
\end{pgfscope}%
\begin{pgfscope}%
\pgfpathrectangle{\pgfqpoint{0.407592in}{0.222904in}}{\pgfqpoint{1.411163in}{1.411163in}}%
\pgfusepath{clip}%
\pgfsetbuttcap%
\pgfsetroundjoin%
\definecolor{currentfill}{rgb}{0.749020,0.772549,0.792157}%
\pgfsetfillcolor{currentfill}%
\pgfsetfillopacity{0.160000}%
\pgfsetlinewidth{0.000000pt}%
\definecolor{currentstroke}{rgb}{0.000000,0.000000,0.000000}%
\pgfsetstrokecolor{currentstroke}%
\pgfsetdash{}{0pt}%
\pgfpathmoveto{\pgfqpoint{1.087933in}{0.750472in}}%
\pgfpathlineto{\pgfqpoint{1.087832in}{0.756652in}}%
\pgfpathlineto{\pgfqpoint{1.087730in}{0.762862in}}%
\pgfpathlineto{\pgfqpoint{1.111769in}{0.763598in}}%
\pgfpathlineto{\pgfqpoint{1.135691in}{0.764998in}}%
\pgfpathlineto{\pgfqpoint{1.133839in}{0.758705in}}%
\pgfpathlineto{\pgfqpoint{1.131996in}{0.752442in}}%
\pgfpathlineto{\pgfqpoint{1.110017in}{0.751151in}}%
\pgfpathlineto{\pgfqpoint{1.087933in}{0.750472in}}%
\pgfpathclose%
\pgfusepath{fill}%
\end{pgfscope}%
\begin{pgfscope}%
\pgfpathrectangle{\pgfqpoint{0.407592in}{0.222904in}}{\pgfqpoint{1.411163in}{1.411163in}}%
\pgfusepath{clip}%
\pgfsetbuttcap%
\pgfsetroundjoin%
\definecolor{currentfill}{rgb}{0.749020,0.772549,0.792157}%
\pgfsetfillcolor{currentfill}%
\pgfsetfillopacity{0.160000}%
\pgfsetlinewidth{0.000000pt}%
\definecolor{currentstroke}{rgb}{0.000000,0.000000,0.000000}%
\pgfsetstrokecolor{currentstroke}%
\pgfsetdash{}{0pt}%
\pgfpathmoveto{\pgfqpoint{0.317235in}{1.223739in}}%
\pgfpathlineto{\pgfqpoint{0.302693in}{1.235276in}}%
\pgfpathlineto{\pgfqpoint{0.305237in}{1.216520in}}%
\pgfpathlineto{\pgfqpoint{0.311335in}{1.197808in}}%
\pgfpathlineto{\pgfqpoint{0.325897in}{1.186828in}}%
\pgfpathlineto{\pgfqpoint{0.319823in}{1.205261in}}%
\pgfusepath{fill}%
\end{pgfscope}%
\begin{pgfscope}%
\pgfpathrectangle{\pgfqpoint{0.407592in}{0.222904in}}{\pgfqpoint{1.411163in}{1.411163in}}%
\pgfusepath{clip}%
\pgfsetbuttcap%
\pgfsetroundjoin%
\definecolor{currentfill}{rgb}{0.749020,0.772549,0.792157}%
\pgfsetfillcolor{currentfill}%
\pgfsetfillopacity{0.160000}%
\pgfsetlinewidth{0.000000pt}%
\definecolor{currentstroke}{rgb}{0.000000,0.000000,0.000000}%
\pgfsetstrokecolor{currentstroke}%
\pgfsetdash{}{0pt}%
\pgfpathmoveto{\pgfqpoint{1.284986in}{0.801039in}}%
\pgfpathlineto{\pgfqpoint{1.292371in}{0.808263in}}%
\pgfpathlineto{\pgfqpoint{1.299788in}{0.815517in}}%
\pgfpathlineto{\pgfqpoint{1.322926in}{0.822027in}}%
\pgfpathlineto{\pgfqpoint{1.344910in}{0.829121in}}%
\pgfpathlineto{\pgfqpoint{1.335931in}{0.821415in}}%
\pgfpathlineto{\pgfqpoint{1.326989in}{0.813739in}}%
\pgfpathlineto{\pgfqpoint{1.306520in}{0.807116in}}%
\pgfpathlineto{\pgfqpoint{1.284986in}{0.801039in}}%
\pgfpathclose%
\pgfusepath{fill}%
\end{pgfscope}%
\begin{pgfscope}%
\pgfpathrectangle{\pgfqpoint{0.407592in}{0.222904in}}{\pgfqpoint{1.411163in}{1.411163in}}%
\pgfusepath{clip}%
\pgfsetbuttcap%
\pgfsetroundjoin%
\definecolor{currentfill}{rgb}{0.749020,0.772549,0.792157}%
\pgfsetfillcolor{currentfill}%
\pgfsetfillopacity{0.160000}%
\pgfsetlinewidth{0.000000pt}%
\definecolor{currentstroke}{rgb}{0.000000,0.000000,0.000000}%
\pgfsetstrokecolor{currentstroke}%
\pgfsetdash{}{0pt}%
\pgfpathmoveto{\pgfqpoint{1.043882in}{0.750965in}}%
\pgfpathlineto{\pgfqpoint{1.041836in}{0.757166in}}%
\pgfpathlineto{\pgfqpoint{1.039780in}{0.763396in}}%
\pgfpathlineto{\pgfqpoint{1.063694in}{0.762794in}}%
\pgfpathlineto{\pgfqpoint{1.087730in}{0.762862in}}%
\pgfpathlineto{\pgfqpoint{1.087832in}{0.756652in}}%
\pgfpathlineto{\pgfqpoint{1.087933in}{0.750472in}}%
\pgfpathlineto{\pgfqpoint{1.065852in}{0.750410in}}%
\pgfpathlineto{\pgfqpoint{1.043882in}{0.750965in}}%
\pgfpathclose%
\pgfusepath{fill}%
\end{pgfscope}%
\begin{pgfscope}%
\pgfpathrectangle{\pgfqpoint{0.407592in}{0.222904in}}{\pgfqpoint{1.411163in}{1.411163in}}%
\pgfusepath{clip}%
\pgfsetbuttcap%
\pgfsetroundjoin%
\definecolor{currentfill}{rgb}{0.749020,0.772549,0.792157}%
\pgfsetfillcolor{currentfill}%
\pgfsetfillopacity{0.160000}%
\pgfsetlinewidth{0.000000pt}%
\definecolor{currentstroke}{rgb}{0.000000,0.000000,0.000000}%
\pgfsetstrokecolor{currentstroke}%
\pgfsetdash{}{0pt}%
\pgfpathmoveto{\pgfqpoint{0.469662in}{1.078431in}}%
\pgfpathlineto{\pgfqpoint{0.455436in}{1.089158in}}%
\pgfpathlineto{\pgfqpoint{0.441176in}{1.099909in}}%
\pgfpathlineto{\pgfqpoint{0.449915in}{1.083936in}}%
\pgfpathlineto{\pgfqpoint{0.461632in}{1.068187in}}%
\pgfpathlineto{\pgfqpoint{0.475581in}{1.058032in}}%
\pgfpathlineto{\pgfqpoint{0.489491in}{1.047904in}}%
\pgfpathlineto{\pgfqpoint{0.478156in}{1.063056in}}%
\pgfpathlineto{\pgfqpoint{0.469662in}{1.078431in}}%
\pgfpathclose%
\pgfusepath{fill}%
\end{pgfscope}%
\begin{pgfscope}%
\pgfpathrectangle{\pgfqpoint{0.407592in}{0.222904in}}{\pgfqpoint{1.411163in}{1.411163in}}%
\pgfusepath{clip}%
\pgfsetbuttcap%
\pgfsetroundjoin%
\definecolor{currentfill}{rgb}{0.749020,0.772549,0.792157}%
\pgfsetfillcolor{currentfill}%
\pgfsetfillopacity{0.160000}%
\pgfsetlinewidth{0.000000pt}%
\definecolor{currentstroke}{rgb}{0.000000,0.000000,0.000000}%
\pgfsetstrokecolor{currentstroke}%
\pgfsetdash{}{0pt}%
\pgfpathmoveto{\pgfqpoint{1.556155in}{0.964496in}}%
\pgfpathlineto{\pgfqpoint{1.568551in}{0.973914in}}%
\pgfpathlineto{\pgfqpoint{1.580986in}{0.983361in}}%
\pgfpathlineto{\pgfqpoint{1.599243in}{0.996733in}}%
\pgfpathlineto{\pgfqpoint{1.615051in}{1.010449in}}%
\pgfpathlineto{\pgfqpoint{1.601862in}{1.000399in}}%
\pgfpathlineto{\pgfqpoint{1.588709in}{0.990377in}}%
\pgfpathlineto{\pgfqpoint{1.573590in}{0.977270in}}%
\pgfpathlineto{\pgfqpoint{1.556155in}{0.964496in}}%
\pgfpathclose%
\pgfusepath{fill}%
\end{pgfscope}%
\begin{pgfscope}%
\pgfpathrectangle{\pgfqpoint{0.407592in}{0.222904in}}{\pgfqpoint{1.411163in}{1.411163in}}%
\pgfusepath{clip}%
\pgfsetbuttcap%
\pgfsetroundjoin%
\definecolor{currentfill}{rgb}{0.749020,0.772549,0.792157}%
\pgfsetfillcolor{currentfill}%
\pgfsetfillopacity{0.160000}%
\pgfsetlinewidth{0.000000pt}%
\definecolor{currentstroke}{rgb}{0.000000,0.000000,0.000000}%
\pgfsetstrokecolor{currentstroke}%
\pgfsetdash{}{0pt}%
\pgfpathmoveto{\pgfqpoint{0.905333in}{0.780622in}}%
\pgfpathlineto{\pgfqpoint{0.897801in}{0.787551in}}%
\pgfpathlineto{\pgfqpoint{0.890235in}{0.794510in}}%
\pgfpathlineto{\pgfqpoint{0.912689in}{0.789747in}}%
\pgfpathlineto{\pgfqpoint{0.936010in}{0.785625in}}%
\pgfpathlineto{\pgfqpoint{0.941858in}{0.778985in}}%
\pgfpathlineto{\pgfqpoint{0.947680in}{0.772375in}}%
\pgfpathlineto{\pgfqpoint{0.926109in}{0.776200in}}%
\pgfpathlineto{\pgfqpoint{0.905333in}{0.780622in}}%
\pgfpathclose%
\pgfusepath{fill}%
\end{pgfscope}%
\begin{pgfscope}%
\pgfpathrectangle{\pgfqpoint{0.407592in}{0.222904in}}{\pgfqpoint{1.411163in}{1.411163in}}%
\pgfusepath{clip}%
\pgfsetbuttcap%
\pgfsetroundjoin%
\definecolor{currentfill}{rgb}{0.749020,0.772549,0.792157}%
\pgfsetfillcolor{currentfill}%
\pgfsetfillopacity{0.160000}%
\pgfsetlinewidth{0.000000pt}%
\definecolor{currentstroke}{rgb}{0.000000,0.000000,0.000000}%
\pgfsetstrokecolor{currentstroke}%
\pgfsetdash{}{0pt}%
\pgfpathmoveto{\pgfqpoint{0.572151in}{0.987724in}}%
\pgfpathlineto{\pgfqpoint{0.558469in}{0.997685in}}%
\pgfpathlineto{\pgfqpoint{0.544749in}{1.007673in}}%
\pgfpathlineto{\pgfqpoint{0.557841in}{0.994021in}}%
\pgfpathlineto{\pgfqpoint{0.573455in}{0.980725in}}%
\pgfpathlineto{\pgfqpoint{0.586560in}{0.971338in}}%
\pgfpathlineto{\pgfqpoint{0.599623in}{0.961980in}}%
\pgfpathlineto{\pgfqpoint{0.584694in}{0.974678in}}%
\pgfpathlineto{\pgfqpoint{0.572151in}{0.987724in}}%
\pgfpathclose%
\pgfusepath{fill}%
\end{pgfscope}%
\begin{pgfscope}%
\pgfpathrectangle{\pgfqpoint{0.407592in}{0.222904in}}{\pgfqpoint{1.411163in}{1.411163in}}%
\pgfusepath{clip}%
\pgfsetbuttcap%
\pgfsetroundjoin%
\definecolor{currentfill}{rgb}{0.749020,0.772549,0.792157}%
\pgfsetfillcolor{currentfill}%
\pgfsetfillopacity{0.160000}%
\pgfsetlinewidth{0.000000pt}%
\definecolor{currentstroke}{rgb}{0.000000,0.000000,0.000000}%
\pgfsetstrokecolor{currentstroke}%
\pgfsetdash{}{0pt}%
\pgfpathmoveto{\pgfqpoint{0.732315in}{0.868730in}}%
\pgfpathlineto{\pgfqpoint{0.720480in}{0.877268in}}%
\pgfpathlineto{\pgfqpoint{0.708603in}{0.885836in}}%
\pgfpathlineto{\pgfqpoint{0.727111in}{0.876213in}}%
\pgfpathlineto{\pgfqpoint{0.747376in}{0.867116in}}%
\pgfpathlineto{\pgfqpoint{0.758105in}{0.859087in}}%
\pgfpathlineto{\pgfqpoint{0.768792in}{0.851089in}}%
\pgfpathlineto{\pgfqpoint{0.749735in}{0.859661in}}%
\pgfpathlineto{\pgfqpoint{0.732315in}{0.868730in}}%
\pgfpathclose%
\pgfusepath{fill}%
\end{pgfscope}%
\begin{pgfscope}%
\pgfpathrectangle{\pgfqpoint{0.407592in}{0.222904in}}{\pgfqpoint{1.411163in}{1.411163in}}%
\pgfusepath{clip}%
\pgfsetbuttcap%
\pgfsetroundjoin%
\definecolor{currentfill}{rgb}{0.749020,0.772549,0.792157}%
\pgfsetfillcolor{currentfill}%
\pgfsetfillopacity{0.160000}%
\pgfsetlinewidth{0.000000pt}%
\definecolor{currentstroke}{rgb}{0.000000,0.000000,0.000000}%
\pgfsetstrokecolor{currentstroke}%
\pgfsetdash{}{0pt}%
\pgfpathmoveto{\pgfqpoint{1.426782in}{0.878103in}}%
\pgfpathlineto{\pgfqpoint{1.437319in}{0.886480in}}%
\pgfpathlineto{\pgfqpoint{1.447894in}{0.894888in}}%
\pgfpathlineto{\pgfqpoint{1.469466in}{0.905165in}}%
\pgfpathlineto{\pgfqpoint{1.489178in}{0.915928in}}%
\pgfpathlineto{\pgfqpoint{1.477455in}{0.906954in}}%
\pgfpathlineto{\pgfqpoint{1.465771in}{0.898010in}}%
\pgfpathlineto{\pgfqpoint{1.447146in}{0.887825in}}%
\pgfpathlineto{\pgfqpoint{1.426782in}{0.878103in}}%
\pgfpathclose%
\pgfusepath{fill}%
\end{pgfscope}%
\begin{pgfscope}%
\pgfpathrectangle{\pgfqpoint{0.407592in}{0.222904in}}{\pgfqpoint{1.411163in}{1.411163in}}%
\pgfusepath{clip}%
\pgfsetbuttcap%
\pgfsetroundjoin%
\definecolor{currentfill}{rgb}{0.749020,0.772549,0.792157}%
\pgfsetfillcolor{currentfill}%
\pgfsetfillopacity{0.160000}%
\pgfsetlinewidth{0.000000pt}%
\definecolor{currentstroke}{rgb}{0.000000,0.000000,0.000000}%
\pgfsetstrokecolor{currentstroke}%
\pgfsetdash{}{0pt}%
\pgfpathmoveto{\pgfqpoint{1.694940in}{1.071325in}}%
\pgfpathlineto{\pgfqpoint{1.708382in}{1.081568in}}%
\pgfpathlineto{\pgfqpoint{1.721861in}{1.091839in}}%
\pgfpathlineto{\pgfqpoint{1.737307in}{1.108230in}}%
\pgfpathlineto{\pgfqpoint{1.749686in}{1.124831in}}%
\pgfpathlineto{\pgfqpoint{1.735760in}{1.113970in}}%
\pgfpathlineto{\pgfqpoint{1.721866in}{1.103134in}}%
\pgfpathlineto{\pgfqpoint{1.709867in}{1.087125in}}%
\pgfpathlineto{\pgfqpoint{1.694940in}{1.071325in}}%
\pgfpathclose%
\pgfusepath{fill}%
\end{pgfscope}%
\begin{pgfscope}%
\pgfpathrectangle{\pgfqpoint{0.407592in}{0.222904in}}{\pgfqpoint{1.411163in}{1.411163in}}%
\pgfusepath{clip}%
\pgfsetbuttcap%
\pgfsetroundjoin%
\definecolor{currentfill}{rgb}{0.749020,0.772549,0.792157}%
\pgfsetfillcolor{currentfill}%
\pgfsetfillopacity{0.160000}%
\pgfsetlinewidth{0.000000pt}%
\definecolor{currentstroke}{rgb}{0.000000,0.000000,0.000000}%
\pgfsetstrokecolor{currentstroke}%
\pgfsetdash{}{0pt}%
\pgfpathmoveto{\pgfqpoint{1.182714in}{0.769759in}}%
\pgfpathlineto{\pgfqpoint{1.186494in}{0.776268in}}%
\pgfpathlineto{\pgfqpoint{1.190292in}{0.782806in}}%
\pgfpathlineto{\pgfqpoint{1.215026in}{0.786399in}}%
\pgfpathlineto{\pgfqpoint{1.239130in}{0.790652in}}%
\pgfpathlineto{\pgfqpoint{1.233491in}{0.783831in}}%
\pgfpathlineto{\pgfqpoint{1.227878in}{0.777040in}}%
\pgfpathlineto{\pgfqpoint{1.205584in}{0.773094in}}%
\pgfpathlineto{\pgfqpoint{1.182714in}{0.769759in}}%
\pgfpathclose%
\pgfusepath{fill}%
\end{pgfscope}%
\begin{pgfscope}%
\pgfpathrectangle{\pgfqpoint{0.407592in}{0.222904in}}{\pgfqpoint{1.411163in}{1.411163in}}%
\pgfusepath{clip}%
\pgfsetbuttcap%
\pgfsetroundjoin%
\definecolor{currentfill}{rgb}{0.749020,0.772549,0.792157}%
\pgfsetfillcolor{currentfill}%
\pgfsetfillopacity{0.160000}%
\pgfsetlinewidth{0.000000pt}%
\definecolor{currentstroke}{rgb}{0.000000,0.000000,0.000000}%
\pgfsetstrokecolor{currentstroke}%
\pgfsetdash{}{0pt}%
\pgfpathmoveto{\pgfqpoint{0.651466in}{0.924842in}}%
\pgfpathlineto{\pgfqpoint{0.638566in}{0.934083in}}%
\pgfpathlineto{\pgfqpoint{0.625626in}{0.943352in}}%
\pgfpathlineto{\pgfqpoint{0.642073in}{0.931653in}}%
\pgfpathlineto{\pgfqpoint{0.660673in}{0.920412in}}%
\pgfpathlineto{\pgfqpoint{0.672719in}{0.911722in}}%
\pgfpathlineto{\pgfqpoint{0.684723in}{0.903063in}}%
\pgfpathlineto{\pgfqpoint{0.667083in}{0.913732in}}%
\pgfpathlineto{\pgfqpoint{0.651466in}{0.924842in}}%
\pgfpathclose%
\pgfusepath{fill}%
\end{pgfscope}%
\begin{pgfscope}%
\pgfpathrectangle{\pgfqpoint{0.407592in}{0.222904in}}{\pgfqpoint{1.411163in}{1.411163in}}%
\pgfusepath{clip}%
\pgfsetbuttcap%
\pgfsetroundjoin%
\definecolor{currentfill}{rgb}{0.749020,0.772549,0.792157}%
\pgfsetfillcolor{currentfill}%
\pgfsetfillopacity{0.160000}%
\pgfsetlinewidth{0.000000pt}%
\definecolor{currentstroke}{rgb}{0.000000,0.000000,0.000000}%
\pgfsetstrokecolor{currentstroke}%
\pgfsetdash{}{0pt}%
\pgfpathmoveto{\pgfqpoint{0.848350in}{0.805852in}}%
\pgfpathlineto{\pgfqpoint{0.839191in}{0.813247in}}%
\pgfpathlineto{\pgfqpoint{0.829995in}{0.820673in}}%
\pgfpathlineto{\pgfqpoint{0.851917in}{0.814285in}}%
\pgfpathlineto{\pgfqpoint{0.875004in}{0.808522in}}%
\pgfpathlineto{\pgfqpoint{0.882636in}{0.801501in}}%
\pgfpathlineto{\pgfqpoint{0.890235in}{0.794510in}}%
\pgfpathlineto{\pgfqpoint{0.868755in}{0.799888in}}%
\pgfpathlineto{\pgfqpoint{0.848350in}{0.805852in}}%
\pgfpathclose%
\pgfusepath{fill}%
\end{pgfscope}%
\begin{pgfscope}%
\pgfpathrectangle{\pgfqpoint{0.407592in}{0.222904in}}{\pgfqpoint{1.411163in}{1.411163in}}%
\pgfusepath{clip}%
\pgfsetbuttcap%
\pgfsetroundjoin%
\definecolor{currentfill}{rgb}{0.749020,0.772549,0.792157}%
\pgfsetfillcolor{currentfill}%
\pgfsetfillopacity{0.160000}%
\pgfsetlinewidth{0.000000pt}%
\definecolor{currentstroke}{rgb}{0.000000,0.000000,0.000000}%
\pgfsetstrokecolor{currentstroke}%
\pgfsetdash{}{0pt}%
\pgfpathmoveto{\pgfqpoint{0.790043in}{0.835185in}}%
\pgfpathlineto{\pgfqpoint{0.779438in}{0.843121in}}%
\pgfpathlineto{\pgfqpoint{0.768792in}{0.851089in}}%
\pgfpathlineto{\pgfqpoint{0.789407in}{0.843060in}}%
\pgfpathlineto{\pgfqpoint{0.811487in}{0.835619in}}%
\pgfpathlineto{\pgfqpoint{0.820760in}{0.828130in}}%
\pgfpathlineto{\pgfqpoint{0.829995in}{0.820673in}}%
\pgfpathlineto{\pgfqpoint{0.809339in}{0.827653in}}%
\pgfpathlineto{\pgfqpoint{0.790043in}{0.835185in}}%
\pgfpathclose%
\pgfusepath{fill}%
\end{pgfscope}%
\begin{pgfscope}%
\pgfpathrectangle{\pgfqpoint{0.407592in}{0.222904in}}{\pgfqpoint{1.411163in}{1.411163in}}%
\pgfusepath{clip}%
\pgfsetbuttcap%
\pgfsetroundjoin%
\definecolor{currentfill}{rgb}{0.749020,0.772549,0.792157}%
\pgfsetfillcolor{currentfill}%
\pgfsetfillopacity{0.160000}%
\pgfsetlinewidth{0.000000pt}%
\definecolor{currentstroke}{rgb}{0.000000,0.000000,0.000000}%
\pgfsetstrokecolor{currentstroke}%
\pgfsetdash{}{0pt}%
\pgfpathmoveto{\pgfqpoint{0.441176in}{1.099909in}}%
\pgfpathlineto{\pgfqpoint{0.426884in}{1.110686in}}%
\pgfpathlineto{\pgfqpoint{0.412558in}{1.121487in}}%
\pgfpathlineto{\pgfqpoint{0.421529in}{1.104922in}}%
\pgfpathlineto{\pgfqpoint{0.433618in}{1.088583in}}%
\pgfpathlineto{\pgfqpoint{0.447645in}{1.078371in}}%
\pgfpathlineto{\pgfqpoint{0.461632in}{1.068187in}}%
\pgfpathlineto{\pgfqpoint{0.449915in}{1.083936in}}%
\pgfpathlineto{\pgfqpoint{0.441176in}{1.099909in}}%
\pgfpathclose%
\pgfusepath{fill}%
\end{pgfscope}%
\begin{pgfscope}%
\pgfpathrectangle{\pgfqpoint{0.407592in}{0.222904in}}{\pgfqpoint{1.411163in}{1.411163in}}%
\pgfusepath{clip}%
\pgfsetbuttcap%
\pgfsetroundjoin%
\definecolor{currentfill}{rgb}{0.749020,0.772549,0.792157}%
\pgfsetfillcolor{currentfill}%
\pgfsetfillopacity{0.160000}%
\pgfsetlinewidth{0.000000pt}%
\definecolor{currentstroke}{rgb}{0.000000,0.000000,0.000000}%
\pgfsetstrokecolor{currentstroke}%
\pgfsetdash{}{0pt}%
\pgfpathmoveto{\pgfqpoint{0.947680in}{0.772375in}}%
\pgfpathlineto{\pgfqpoint{0.941858in}{0.778985in}}%
\pgfpathlineto{\pgfqpoint{0.936010in}{0.785625in}}%
\pgfpathlineto{\pgfqpoint{0.960084in}{0.782167in}}%
\pgfpathlineto{\pgfqpoint{0.984796in}{0.779391in}}%
\pgfpathlineto{\pgfqpoint{0.988801in}{0.772976in}}%
\pgfpathlineto{\pgfqpoint{0.992788in}{0.766591in}}%
\pgfpathlineto{\pgfqpoint{0.969942in}{0.769166in}}%
\pgfpathlineto{\pgfqpoint{0.947680in}{0.772375in}}%
\pgfpathclose%
\pgfusepath{fill}%
\end{pgfscope}%
\begin{pgfscope}%
\pgfpathrectangle{\pgfqpoint{0.407592in}{0.222904in}}{\pgfqpoint{1.411163in}{1.411163in}}%
\pgfusepath{clip}%
\pgfsetbuttcap%
\pgfsetroundjoin%
\definecolor{currentfill}{rgb}{0.749020,0.772549,0.792157}%
\pgfsetfillcolor{currentfill}%
\pgfsetfillopacity{0.160000}%
\pgfsetlinewidth{0.000000pt}%
\definecolor{currentstroke}{rgb}{0.000000,0.000000,0.000000}%
\pgfsetstrokecolor{currentstroke}%
\pgfsetdash{}{0pt}%
\pgfpathmoveto{\pgfqpoint{1.512741in}{0.933965in}}%
\pgfpathlineto{\pgfqpoint{1.524581in}{0.943028in}}%
\pgfpathlineto{\pgfqpoint{1.536460in}{0.952122in}}%
\pgfpathlineto{\pgfqpoint{1.556155in}{0.964496in}}%
\pgfpathlineto{\pgfqpoint{1.543797in}{0.955108in}}%
\pgfpathlineto{\pgfqpoint{1.531478in}{0.945749in}}%
\pgfusepath{fill}%
\end{pgfscope}%
\begin{pgfscope}%
\pgfpathrectangle{\pgfqpoint{0.407592in}{0.222904in}}{\pgfqpoint{1.411163in}{1.411163in}}%
\pgfusepath{clip}%
\pgfsetbuttcap%
\pgfsetroundjoin%
\definecolor{currentfill}{rgb}{0.749020,0.772549,0.792157}%
\pgfsetfillcolor{currentfill}%
\pgfsetfillopacity{0.160000}%
\pgfsetlinewidth{0.000000pt}%
\definecolor{currentstroke}{rgb}{0.000000,0.000000,0.000000}%
\pgfsetstrokecolor{currentstroke}%
\pgfsetdash{}{0pt}%
\pgfpathmoveto{\pgfqpoint{1.580986in}{0.983361in}}%
\pgfpathlineto{\pgfqpoint{1.593461in}{0.992838in}}%
\pgfpathlineto{\pgfqpoint{1.605974in}{1.002344in}}%
\pgfpathlineto{\pgfqpoint{1.625047in}{1.016311in}}%
\pgfpathlineto{\pgfqpoint{1.641536in}{1.030631in}}%
\pgfpathlineto{\pgfqpoint{1.628276in}{1.020526in}}%
\pgfpathlineto{\pgfqpoint{1.615051in}{1.010449in}}%
\pgfpathlineto{\pgfqpoint{1.599243in}{0.996733in}}%
\pgfpathlineto{\pgfqpoint{1.580986in}{0.983361in}}%
\pgfpathclose%
\pgfusepath{fill}%
\end{pgfscope}%
\begin{pgfscope}%
\pgfpathrectangle{\pgfqpoint{0.407592in}{0.222904in}}{\pgfqpoint{1.411163in}{1.411163in}}%
\pgfusepath{clip}%
\pgfsetbuttcap%
\pgfsetroundjoin%
\definecolor{currentfill}{rgb}{0.749020,0.772549,0.792157}%
\pgfsetfillcolor{currentfill}%
\pgfsetfillopacity{0.160000}%
\pgfsetlinewidth{0.000000pt}%
\definecolor{currentstroke}{rgb}{0.000000,0.000000,0.000000}%
\pgfsetstrokecolor{currentstroke}%
\pgfsetdash{}{0pt}%
\pgfpathmoveto{\pgfqpoint{1.362974in}{0.844625in}}%
\pgfpathlineto{\pgfqpoint{1.372061in}{0.852424in}}%
\pgfpathlineto{\pgfqpoint{1.381184in}{0.860254in}}%
\pgfpathlineto{\pgfqpoint{1.404763in}{0.868896in}}%
\pgfpathlineto{\pgfqpoint{1.426782in}{0.878103in}}%
\pgfpathlineto{\pgfqpoint{1.416284in}{0.869757in}}%
\pgfpathlineto{\pgfqpoint{1.405825in}{0.861442in}}%
\pgfpathlineto{\pgfqpoint{1.385126in}{0.852766in}}%
\pgfpathlineto{\pgfqpoint{1.362974in}{0.844625in}}%
\pgfpathclose%
\pgfusepath{fill}%
\end{pgfscope}%
\begin{pgfscope}%
\pgfpathrectangle{\pgfqpoint{0.407592in}{0.222904in}}{\pgfqpoint{1.411163in}{1.411163in}}%
\pgfusepath{clip}%
\pgfsetbuttcap%
\pgfsetroundjoin%
\definecolor{currentfill}{rgb}{0.749020,0.772549,0.792157}%
\pgfsetfillcolor{currentfill}%
\pgfsetfillopacity{0.160000}%
\pgfsetlinewidth{0.000000pt}%
\definecolor{currentstroke}{rgb}{0.000000,0.000000,0.000000}%
\pgfsetstrokecolor{currentstroke}%
\pgfsetdash{}{0pt}%
\pgfpathmoveto{\pgfqpoint{1.239130in}{0.790652in}}%
\pgfpathlineto{\pgfqpoint{1.244795in}{0.797504in}}%
\pgfpathlineto{\pgfqpoint{1.250484in}{0.804387in}}%
\pgfpathlineto{\pgfqpoint{1.275602in}{0.809626in}}%
\pgfpathlineto{\pgfqpoint{1.299788in}{0.815517in}}%
\pgfpathlineto{\pgfqpoint{1.292371in}{0.808263in}}%
\pgfpathlineto{\pgfqpoint{1.284986in}{0.801039in}}%
\pgfpathlineto{\pgfqpoint{1.262487in}{0.795541in}}%
\pgfpathlineto{\pgfqpoint{1.239130in}{0.790652in}}%
\pgfpathclose%
\pgfusepath{fill}%
\end{pgfscope}%
\begin{pgfscope}%
\pgfpathrectangle{\pgfqpoint{0.407592in}{0.222904in}}{\pgfqpoint{1.411163in}{1.411163in}}%
\pgfusepath{clip}%
\pgfsetbuttcap%
\pgfsetroundjoin%
\definecolor{currentfill}{rgb}{0.749020,0.772549,0.792157}%
\pgfsetfillcolor{currentfill}%
\pgfsetfillopacity{0.160000}%
\pgfsetlinewidth{0.000000pt}%
\definecolor{currentstroke}{rgb}{0.000000,0.000000,0.000000}%
\pgfsetstrokecolor{currentstroke}%
\pgfsetdash{}{0pt}%
\pgfpathmoveto{\pgfqpoint{1.135691in}{0.764998in}}%
\pgfpathlineto{\pgfqpoint{1.137552in}{0.771321in}}%
\pgfpathlineto{\pgfqpoint{1.139422in}{0.777674in}}%
\pgfpathlineto{\pgfqpoint{1.165049in}{0.779892in}}%
\pgfpathlineto{\pgfqpoint{1.190292in}{0.782806in}}%
\pgfpathlineto{\pgfqpoint{1.186494in}{0.776268in}}%
\pgfpathlineto{\pgfqpoint{1.182714in}{0.769759in}}%
\pgfpathlineto{\pgfqpoint{1.159378in}{0.767056in}}%
\pgfpathlineto{\pgfqpoint{1.135691in}{0.764998in}}%
\pgfpathclose%
\pgfusepath{fill}%
\end{pgfscope}%
\begin{pgfscope}%
\pgfpathrectangle{\pgfqpoint{0.407592in}{0.222904in}}{\pgfqpoint{1.411163in}{1.411163in}}%
\pgfusepath{clip}%
\pgfsetbuttcap%
\pgfsetroundjoin%
\definecolor{currentfill}{rgb}{0.749020,0.772549,0.792157}%
\pgfsetfillcolor{currentfill}%
\pgfsetfillopacity{0.160000}%
\pgfsetlinewidth{0.000000pt}%
\definecolor{currentstroke}{rgb}{0.000000,0.000000,0.000000}%
\pgfsetstrokecolor{currentstroke}%
\pgfsetdash{}{0pt}%
\pgfpathmoveto{\pgfqpoint{0.544749in}{1.007673in}}%
\pgfpathlineto{\pgfqpoint{0.530992in}{1.017690in}}%
\pgfpathlineto{\pgfqpoint{0.517197in}{1.027733in}}%
\pgfpathlineto{\pgfqpoint{0.530828in}{1.013479in}}%
\pgfpathlineto{\pgfqpoint{0.547119in}{0.999591in}}%
\pgfpathlineto{\pgfqpoint{0.560308in}{0.990143in}}%
\pgfpathlineto{\pgfqpoint{0.573455in}{0.980725in}}%
\pgfpathlineto{\pgfqpoint{0.557841in}{0.994021in}}%
\pgfpathlineto{\pgfqpoint{0.544749in}{1.007673in}}%
\pgfpathclose%
\pgfusepath{fill}%
\end{pgfscope}%
\begin{pgfscope}%
\pgfpathrectangle{\pgfqpoint{0.407592in}{0.222904in}}{\pgfqpoint{1.411163in}{1.411163in}}%
\pgfusepath{clip}%
\pgfsetbuttcap%
\pgfsetroundjoin%
\definecolor{currentfill}{rgb}{0.749020,0.772549,0.792157}%
\pgfsetfillcolor{currentfill}%
\pgfsetfillopacity{0.160000}%
\pgfsetlinewidth{0.000000pt}%
\definecolor{currentstroke}{rgb}{0.000000,0.000000,0.000000}%
\pgfsetstrokecolor{currentstroke}%
\pgfsetdash{}{0pt}%
\pgfpathmoveto{\pgfqpoint{1.299788in}{0.815517in}}%
\pgfpathlineto{\pgfqpoint{1.307236in}{0.822803in}}%
\pgfpathlineto{\pgfqpoint{1.314716in}{0.830120in}}%
\pgfpathlineto{\pgfqpoint{1.339469in}{0.837062in}}%
\pgfpathlineto{\pgfqpoint{1.362974in}{0.844625in}}%
\pgfpathlineto{\pgfqpoint{1.353924in}{0.836857in}}%
\pgfpathlineto{\pgfqpoint{1.344910in}{0.829121in}}%
\pgfpathlineto{\pgfqpoint{1.322926in}{0.822027in}}%
\pgfpathlineto{\pgfqpoint{1.299788in}{0.815517in}}%
\pgfpathclose%
\pgfusepath{fill}%
\end{pgfscope}%
\begin{pgfscope}%
\pgfpathrectangle{\pgfqpoint{0.407592in}{0.222904in}}{\pgfqpoint{1.411163in}{1.411163in}}%
\pgfusepath{clip}%
\pgfsetbuttcap%
\pgfsetroundjoin%
\definecolor{currentfill}{rgb}{0.749020,0.772549,0.792157}%
\pgfsetfillcolor{currentfill}%
\pgfsetfillopacity{0.160000}%
\pgfsetlinewidth{0.000000pt}%
\definecolor{currentstroke}{rgb}{0.000000,0.000000,0.000000}%
\pgfsetstrokecolor{currentstroke}%
\pgfsetdash{}{0pt}%
\pgfpathmoveto{\pgfqpoint{1.721861in}{1.091839in}}%
\pgfpathlineto{\pgfqpoint{1.735377in}{1.102138in}}%
\pgfpathlineto{\pgfqpoint{1.748930in}{1.112466in}}%
\pgfpathlineto{\pgfqpoint{1.764885in}{1.129441in}}%
\pgfpathlineto{\pgfqpoint{1.777633in}{1.146626in}}%
\pgfpathlineto{\pgfqpoint{1.763644in}{1.135716in}}%
\pgfpathlineto{\pgfqpoint{1.749686in}{1.124831in}}%
\pgfpathlineto{\pgfqpoint{1.737307in}{1.108230in}}%
\pgfpathlineto{\pgfqpoint{1.721861in}{1.091839in}}%
\pgfpathclose%
\pgfusepath{fill}%
\end{pgfscope}%
\begin{pgfscope}%
\pgfpathrectangle{\pgfqpoint{0.407592in}{0.222904in}}{\pgfqpoint{1.411163in}{1.411163in}}%
\pgfusepath{clip}%
\pgfsetbuttcap%
\pgfsetroundjoin%
\definecolor{currentfill}{rgb}{0.749020,0.772549,0.792157}%
\pgfsetfillcolor{currentfill}%
\pgfsetfillopacity{0.160000}%
\pgfsetlinewidth{0.000000pt}%
\definecolor{currentstroke}{rgb}{0.000000,0.000000,0.000000}%
\pgfsetstrokecolor{currentstroke}%
\pgfsetdash{}{0pt}%
\pgfpathmoveto{\pgfqpoint{0.992788in}{0.766591in}}%
\pgfpathlineto{\pgfqpoint{0.988801in}{0.772976in}}%
\pgfpathlineto{\pgfqpoint{0.984796in}{0.779391in}}%
\pgfpathlineto{\pgfqpoint{1.010022in}{0.777314in}}%
\pgfpathlineto{\pgfqpoint{1.035637in}{0.775947in}}%
\pgfpathlineto{\pgfqpoint{1.037713in}{0.769657in}}%
\pgfpathlineto{\pgfqpoint{1.039780in}{0.763396in}}%
\pgfpathlineto{\pgfqpoint{1.016105in}{0.764664in}}%
\pgfpathlineto{\pgfqpoint{0.992788in}{0.766591in}}%
\pgfpathclose%
\pgfusepath{fill}%
\end{pgfscope}%
\begin{pgfscope}%
\pgfpathrectangle{\pgfqpoint{0.407592in}{0.222904in}}{\pgfqpoint{1.411163in}{1.411163in}}%
\pgfusepath{clip}%
\pgfsetbuttcap%
\pgfsetroundjoin%
\definecolor{currentfill}{rgb}{0.749020,0.772549,0.792157}%
\pgfsetfillcolor{currentfill}%
\pgfsetfillopacity{0.160000}%
\pgfsetlinewidth{0.000000pt}%
\definecolor{currentstroke}{rgb}{0.000000,0.000000,0.000000}%
\pgfsetstrokecolor{currentstroke}%
\pgfsetdash{}{0pt}%
\pgfpathmoveto{\pgfqpoint{1.087730in}{0.762862in}}%
\pgfpathlineto{\pgfqpoint{1.087628in}{0.769101in}}%
\pgfpathlineto{\pgfqpoint{1.087525in}{0.775370in}}%
\pgfpathlineto{\pgfqpoint{1.113537in}{0.776164in}}%
\pgfpathlineto{\pgfqpoint{1.139422in}{0.777674in}}%
\pgfpathlineto{\pgfqpoint{1.137552in}{0.771321in}}%
\pgfpathlineto{\pgfqpoint{1.135691in}{0.764998in}}%
\pgfpathlineto{\pgfqpoint{1.111769in}{0.763598in}}%
\pgfpathlineto{\pgfqpoint{1.087730in}{0.762862in}}%
\pgfpathclose%
\pgfusepath{fill}%
\end{pgfscope}%
\begin{pgfscope}%
\pgfpathrectangle{\pgfqpoint{0.407592in}{0.222904in}}{\pgfqpoint{1.411163in}{1.411163in}}%
\pgfusepath{clip}%
\pgfsetbuttcap%
\pgfsetroundjoin%
\definecolor{currentfill}{rgb}{0.749020,0.772549,0.792157}%
\pgfsetfillcolor{currentfill}%
\pgfsetfillopacity{0.160000}%
\pgfsetlinewidth{0.000000pt}%
\definecolor{currentstroke}{rgb}{0.000000,0.000000,0.000000}%
\pgfsetstrokecolor{currentstroke}%
\pgfsetdash{}{0pt}%
\pgfpathmoveto{\pgfqpoint{1.447894in}{0.894888in}}%
\pgfpathlineto{\pgfqpoint{1.458509in}{0.903327in}}%
\pgfpathlineto{\pgfqpoint{1.469162in}{0.911797in}}%
\pgfpathlineto{\pgfqpoint{1.491942in}{0.922627in}}%
\pgfpathlineto{\pgfqpoint{1.512741in}{0.933965in}}%
\pgfpathlineto{\pgfqpoint{1.500940in}{0.924931in}}%
\pgfpathlineto{\pgfqpoint{1.489178in}{0.915928in}}%
\pgfpathlineto{\pgfqpoint{1.469466in}{0.905165in}}%
\pgfpathlineto{\pgfqpoint{1.447894in}{0.894888in}}%
\pgfpathclose%
\pgfusepath{fill}%
\end{pgfscope}%
\begin{pgfscope}%
\pgfpathrectangle{\pgfqpoint{0.407592in}{0.222904in}}{\pgfqpoint{1.411163in}{1.411163in}}%
\pgfusepath{clip}%
\pgfsetbuttcap%
\pgfsetroundjoin%
\definecolor{currentfill}{rgb}{0.749020,0.772549,0.792157}%
\pgfsetfillcolor{currentfill}%
\pgfsetfillopacity{0.160000}%
\pgfsetlinewidth{0.000000pt}%
\definecolor{currentstroke}{rgb}{0.000000,0.000000,0.000000}%
\pgfsetstrokecolor{currentstroke}%
\pgfsetdash{}{0pt}%
\pgfpathmoveto{\pgfqpoint{1.039780in}{0.763396in}}%
\pgfpathlineto{\pgfqpoint{1.037713in}{0.769657in}}%
\pgfpathlineto{\pgfqpoint{1.035637in}{0.775947in}}%
\pgfpathlineto{\pgfqpoint{1.061515in}{0.775298in}}%
\pgfpathlineto{\pgfqpoint{1.087525in}{0.775370in}}%
\pgfpathlineto{\pgfqpoint{1.087628in}{0.769101in}}%
\pgfpathlineto{\pgfqpoint{1.087730in}{0.762862in}}%
\pgfpathlineto{\pgfqpoint{1.063694in}{0.762794in}}%
\pgfpathlineto{\pgfqpoint{1.039780in}{0.763396in}}%
\pgfpathclose%
\pgfusepath{fill}%
\end{pgfscope}%
\begin{pgfscope}%
\pgfpathrectangle{\pgfqpoint{0.407592in}{0.222904in}}{\pgfqpoint{1.411163in}{1.411163in}}%
\pgfusepath{clip}%
\pgfsetbuttcap%
\pgfsetroundjoin%
\definecolor{currentfill}{rgb}{0.749020,0.772549,0.792157}%
\pgfsetfillcolor{currentfill}%
\pgfsetfillopacity{0.160000}%
\pgfsetlinewidth{0.000000pt}%
\definecolor{currentstroke}{rgb}{0.000000,0.000000,0.000000}%
\pgfsetstrokecolor{currentstroke}%
\pgfsetdash{}{0pt}%
\pgfpathmoveto{\pgfqpoint{0.708603in}{0.885836in}}%
\pgfpathlineto{\pgfqpoint{0.696684in}{0.894434in}}%
\pgfpathlineto{\pgfqpoint{0.684723in}{0.903063in}}%
\pgfpathlineto{\pgfqpoint{0.704319in}{0.892890in}}%
\pgfpathlineto{\pgfqpoint{0.725793in}{0.883268in}}%
\pgfpathlineto{\pgfqpoint{0.736606in}{0.875176in}}%
\pgfpathlineto{\pgfqpoint{0.747376in}{0.867116in}}%
\pgfpathlineto{\pgfqpoint{0.727111in}{0.876213in}}%
\pgfpathlineto{\pgfqpoint{0.708603in}{0.885836in}}%
\pgfpathclose%
\pgfusepath{fill}%
\end{pgfscope}%
\begin{pgfscope}%
\pgfpathrectangle{\pgfqpoint{0.407592in}{0.222904in}}{\pgfqpoint{1.411163in}{1.411163in}}%
\pgfusepath{clip}%
\pgfsetbuttcap%
\pgfsetroundjoin%
\definecolor{currentfill}{rgb}{0.749020,0.772549,0.792157}%
\pgfsetfillcolor{currentfill}%
\pgfsetfillopacity{0.160000}%
\pgfsetlinewidth{0.000000pt}%
\definecolor{currentstroke}{rgb}{0.000000,0.000000,0.000000}%
\pgfsetstrokecolor{currentstroke}%
\pgfsetdash{}{0pt}%
\pgfpathmoveto{\pgfqpoint{0.412558in}{1.121487in}}%
\pgfpathlineto{\pgfqpoint{0.398199in}{1.132314in}}%
\pgfpathlineto{\pgfqpoint{0.383806in}{1.143166in}}%
\pgfpathlineto{\pgfqpoint{0.392998in}{1.126016in}}%
\pgfpathlineto{\pgfqpoint{0.405447in}{1.109093in}}%
\pgfpathlineto{\pgfqpoint{0.419552in}{1.098824in}}%
\pgfpathlineto{\pgfqpoint{0.433618in}{1.088583in}}%
\pgfpathlineto{\pgfqpoint{0.421529in}{1.104922in}}%
\pgfpathlineto{\pgfqpoint{0.412558in}{1.121487in}}%
\pgfpathclose%
\pgfusepath{fill}%
\end{pgfscope}%
\begin{pgfscope}%
\pgfpathrectangle{\pgfqpoint{0.407592in}{0.222904in}}{\pgfqpoint{1.411163in}{1.411163in}}%
\pgfusepath{clip}%
\pgfsetbuttcap%
\pgfsetroundjoin%
\definecolor{currentfill}{rgb}{0.749020,0.772549,0.792157}%
\pgfsetfillcolor{currentfill}%
\pgfsetfillopacity{0.160000}%
\pgfsetlinewidth{0.000000pt}%
\definecolor{currentstroke}{rgb}{0.000000,0.000000,0.000000}%
\pgfsetstrokecolor{currentstroke}%
\pgfsetdash{}{0pt}%
\pgfpathmoveto{\pgfqpoint{0.890235in}{0.794510in}}%
\pgfpathlineto{\pgfqpoint{0.882636in}{0.801501in}}%
\pgfpathlineto{\pgfqpoint{0.875004in}{0.808522in}}%
\pgfpathlineto{\pgfqpoint{0.899148in}{0.803417in}}%
\pgfpathlineto{\pgfqpoint{0.924231in}{0.798997in}}%
\pgfpathlineto{\pgfqpoint{0.930134in}{0.792296in}}%
\pgfpathlineto{\pgfqpoint{0.936010in}{0.785625in}}%
\pgfpathlineto{\pgfqpoint{0.912689in}{0.789747in}}%
\pgfpathlineto{\pgfqpoint{0.890235in}{0.794510in}}%
\pgfpathclose%
\pgfusepath{fill}%
\end{pgfscope}%
\begin{pgfscope}%
\pgfpathrectangle{\pgfqpoint{0.407592in}{0.222904in}}{\pgfqpoint{1.411163in}{1.411163in}}%
\pgfusepath{clip}%
\pgfsetbuttcap%
\pgfsetroundjoin%
\definecolor{currentfill}{rgb}{0.749020,0.772549,0.792157}%
\pgfsetfillcolor{currentfill}%
\pgfsetfillopacity{0.160000}%
\pgfsetlinewidth{0.000000pt}%
\definecolor{currentstroke}{rgb}{0.000000,0.000000,0.000000}%
\pgfsetstrokecolor{currentstroke}%
\pgfsetdash{}{0pt}%
\pgfpathmoveto{\pgfqpoint{0.625626in}{0.943352in}}%
\pgfpathlineto{\pgfqpoint{0.612645in}{0.952651in}}%
\pgfpathlineto{\pgfqpoint{0.599623in}{0.961980in}}%
\pgfpathlineto{\pgfqpoint{0.616895in}{0.949694in}}%
\pgfpathlineto{\pgfqpoint{0.636452in}{0.937885in}}%
\pgfpathlineto{\pgfqpoint{0.648584in}{0.929133in}}%
\pgfpathlineto{\pgfqpoint{0.660673in}{0.920412in}}%
\pgfpathlineto{\pgfqpoint{0.642073in}{0.931653in}}%
\pgfpathlineto{\pgfqpoint{0.625626in}{0.943352in}}%
\pgfpathclose%
\pgfusepath{fill}%
\end{pgfscope}%
\begin{pgfscope}%
\pgfpathrectangle{\pgfqpoint{0.407592in}{0.222904in}}{\pgfqpoint{1.411163in}{1.411163in}}%
\pgfusepath{clip}%
\pgfsetbuttcap%
\pgfsetroundjoin%
\definecolor{currentfill}{rgb}{0.749020,0.772549,0.792157}%
\pgfsetfillcolor{currentfill}%
\pgfsetfillopacity{0.160000}%
\pgfsetlinewidth{0.000000pt}%
\definecolor{currentstroke}{rgb}{0.000000,0.000000,0.000000}%
\pgfsetstrokecolor{currentstroke}%
\pgfsetdash{}{0pt}%
\pgfpathmoveto{\pgfqpoint{1.605974in}{1.002344in}}%
\pgfpathlineto{\pgfqpoint{1.618527in}{1.011881in}}%
\pgfpathlineto{\pgfqpoint{1.631120in}{1.021447in}}%
\pgfpathlineto{\pgfqpoint{1.651003in}{1.036004in}}%
\pgfpathlineto{\pgfqpoint{1.668165in}{1.050923in}}%
\pgfpathlineto{\pgfqpoint{1.654833in}{1.040763in}}%
\pgfpathlineto{\pgfqpoint{1.641536in}{1.030631in}}%
\pgfpathlineto{\pgfqpoint{1.625047in}{1.016311in}}%
\pgfpathlineto{\pgfqpoint{1.605974in}{1.002344in}}%
\pgfpathclose%
\pgfusepath{fill}%
\end{pgfscope}%
\begin{pgfscope}%
\pgfpathrectangle{\pgfqpoint{0.407592in}{0.222904in}}{\pgfqpoint{1.411163in}{1.411163in}}%
\pgfusepath{clip}%
\pgfsetbuttcap%
\pgfsetroundjoin%
\definecolor{currentfill}{rgb}{0.749020,0.772549,0.792157}%
\pgfsetfillcolor{currentfill}%
\pgfsetfillopacity{0.160000}%
\pgfsetlinewidth{0.000000pt}%
\definecolor{currentstroke}{rgb}{0.000000,0.000000,0.000000}%
\pgfsetstrokecolor{currentstroke}%
\pgfsetdash{}{0pt}%
\pgfpathmoveto{\pgfqpoint{1.536460in}{0.952122in}}%
\pgfpathlineto{\pgfqpoint{1.548380in}{0.961246in}}%
\pgfpathlineto{\pgfqpoint{1.560339in}{0.970401in}}%
\pgfpathlineto{\pgfqpoint{1.580986in}{0.983361in}}%
\pgfpathlineto{\pgfqpoint{1.568551in}{0.973914in}}%
\pgfpathlineto{\pgfqpoint{1.556155in}{0.964496in}}%
\pgfusepath{fill}%
\end{pgfscope}%
\begin{pgfscope}%
\pgfpathrectangle{\pgfqpoint{0.407592in}{0.222904in}}{\pgfqpoint{1.411163in}{1.411163in}}%
\pgfusepath{clip}%
\pgfsetbuttcap%
\pgfsetroundjoin%
\definecolor{currentfill}{rgb}{0.749020,0.772549,0.792157}%
\pgfsetfillcolor{currentfill}%
\pgfsetfillopacity{0.160000}%
\pgfsetlinewidth{0.000000pt}%
\definecolor{currentstroke}{rgb}{0.000000,0.000000,0.000000}%
\pgfsetstrokecolor{currentstroke}%
\pgfsetdash{}{0pt}%
\pgfpathmoveto{\pgfqpoint{1.190292in}{0.782806in}}%
\pgfpathlineto{\pgfqpoint{1.194108in}{0.789375in}}%
\pgfpathlineto{\pgfqpoint{1.197941in}{0.795975in}}%
\pgfpathlineto{\pgfqpoint{1.224555in}{0.799828in}}%
\pgfpathlineto{\pgfqpoint{1.250484in}{0.804387in}}%
\pgfpathlineto{\pgfqpoint{1.244795in}{0.797504in}}%
\pgfpathlineto{\pgfqpoint{1.239130in}{0.790652in}}%
\pgfpathlineto{\pgfqpoint{1.215026in}{0.786399in}}%
\pgfpathlineto{\pgfqpoint{1.190292in}{0.782806in}}%
\pgfpathclose%
\pgfusepath{fill}%
\end{pgfscope}%
\begin{pgfscope}%
\pgfpathrectangle{\pgfqpoint{0.407592in}{0.222904in}}{\pgfqpoint{1.411163in}{1.411163in}}%
\pgfusepath{clip}%
\pgfsetbuttcap%
\pgfsetroundjoin%
\definecolor{currentfill}{rgb}{0.749020,0.772549,0.792157}%
\pgfsetfillcolor{currentfill}%
\pgfsetfillopacity{0.160000}%
\pgfsetlinewidth{0.000000pt}%
\definecolor{currentstroke}{rgb}{0.000000,0.000000,0.000000}%
\pgfsetstrokecolor{currentstroke}%
\pgfsetdash{}{0pt}%
\pgfpathmoveto{\pgfqpoint{0.768792in}{0.851089in}}%
\pgfpathlineto{\pgfqpoint{0.758105in}{0.859087in}}%
\pgfpathlineto{\pgfqpoint{0.747376in}{0.867116in}}%
\pgfpathlineto{\pgfqpoint{0.769312in}{0.858592in}}%
\pgfpathlineto{\pgfqpoint{0.792823in}{0.850689in}}%
\pgfpathlineto{\pgfqpoint{0.802175in}{0.843138in}}%
\pgfpathlineto{\pgfqpoint{0.811487in}{0.835619in}}%
\pgfpathlineto{\pgfqpoint{0.789407in}{0.843060in}}%
\pgfpathlineto{\pgfqpoint{0.768792in}{0.851089in}}%
\pgfpathclose%
\pgfusepath{fill}%
\end{pgfscope}%
\begin{pgfscope}%
\pgfpathrectangle{\pgfqpoint{0.407592in}{0.222904in}}{\pgfqpoint{1.411163in}{1.411163in}}%
\pgfusepath{clip}%
\pgfsetbuttcap%
\pgfsetroundjoin%
\definecolor{currentfill}{rgb}{0.749020,0.772549,0.792157}%
\pgfsetfillcolor{currentfill}%
\pgfsetfillopacity{0.160000}%
\pgfsetlinewidth{0.000000pt}%
\definecolor{currentstroke}{rgb}{0.000000,0.000000,0.000000}%
\pgfsetstrokecolor{currentstroke}%
\pgfsetdash{}{0pt}%
\pgfpathmoveto{\pgfqpoint{0.829995in}{0.820673in}}%
\pgfpathlineto{\pgfqpoint{0.820760in}{0.828130in}}%
\pgfpathlineto{\pgfqpoint{0.811487in}{0.835619in}}%
\pgfpathlineto{\pgfqpoint{0.834933in}{0.828805in}}%
\pgfpathlineto{\pgfqpoint{0.859638in}{0.822658in}}%
\pgfpathlineto{\pgfqpoint{0.867338in}{0.815574in}}%
\pgfpathlineto{\pgfqpoint{0.875004in}{0.808522in}}%
\pgfpathlineto{\pgfqpoint{0.851917in}{0.814285in}}%
\pgfpathlineto{\pgfqpoint{0.829995in}{0.820673in}}%
\pgfpathclose%
\pgfusepath{fill}%
\end{pgfscope}%
\begin{pgfscope}%
\pgfpathrectangle{\pgfqpoint{0.407592in}{0.222904in}}{\pgfqpoint{1.411163in}{1.411163in}}%
\pgfusepath{clip}%
\pgfsetbuttcap%
\pgfsetroundjoin%
\definecolor{currentfill}{rgb}{0.749020,0.772549,0.792157}%
\pgfsetfillcolor{currentfill}%
\pgfsetfillopacity{0.160000}%
\pgfsetlinewidth{0.000000pt}%
\definecolor{currentstroke}{rgb}{0.000000,0.000000,0.000000}%
\pgfsetstrokecolor{currentstroke}%
\pgfsetdash{}{0pt}%
\pgfpathmoveto{\pgfqpoint{1.748930in}{1.112466in}}%
\pgfpathlineto{\pgfqpoint{1.762521in}{1.122822in}}%
\pgfpathlineto{\pgfqpoint{1.776149in}{1.133206in}}%
\pgfpathlineto{\pgfqpoint{1.792602in}{1.150760in}}%
\pgfpathlineto{\pgfqpoint{1.805709in}{1.168522in}}%
\pgfpathlineto{\pgfqpoint{1.791655in}{1.157562in}}%
\pgfpathlineto{\pgfqpoint{1.777633in}{1.146626in}}%
\pgfpathlineto{\pgfqpoint{1.764885in}{1.129441in}}%
\pgfpathlineto{\pgfqpoint{1.748930in}{1.112466in}}%
\pgfpathclose%
\pgfusepath{fill}%
\end{pgfscope}%
\begin{pgfscope}%
\pgfpathrectangle{\pgfqpoint{0.407592in}{0.222904in}}{\pgfqpoint{1.411163in}{1.411163in}}%
\pgfusepath{clip}%
\pgfsetbuttcap%
\pgfsetroundjoin%
\definecolor{currentfill}{rgb}{0.749020,0.772549,0.792157}%
\pgfsetfillcolor{currentfill}%
\pgfsetfillopacity{0.160000}%
\pgfsetlinewidth{0.000000pt}%
\definecolor{currentstroke}{rgb}{0.000000,0.000000,0.000000}%
\pgfsetstrokecolor{currentstroke}%
\pgfsetdash{}{0pt}%
\pgfpathmoveto{\pgfqpoint{0.517197in}{1.027733in}}%
\pgfpathlineto{\pgfqpoint{0.503363in}{1.037805in}}%
\pgfpathlineto{\pgfqpoint{0.489491in}{1.047904in}}%
\pgfpathlineto{\pgfqpoint{0.503654in}{1.033053in}}%
\pgfpathlineto{\pgfqpoint{0.520615in}{1.018577in}}%
\pgfpathlineto{\pgfqpoint{0.533889in}{1.009069in}}%
\pgfpathlineto{\pgfqpoint{0.547119in}{0.999591in}}%
\pgfpathlineto{\pgfqpoint{0.530828in}{1.013479in}}%
\pgfpathlineto{\pgfqpoint{0.517197in}{1.027733in}}%
\pgfpathclose%
\pgfusepath{fill}%
\end{pgfscope}%
\begin{pgfscope}%
\pgfpathrectangle{\pgfqpoint{0.407592in}{0.222904in}}{\pgfqpoint{1.411163in}{1.411163in}}%
\pgfusepath{clip}%
\pgfsetbuttcap%
\pgfsetroundjoin%
\definecolor{currentfill}{rgb}{0.749020,0.772549,0.792157}%
\pgfsetfillcolor{currentfill}%
\pgfsetfillopacity{0.160000}%
\pgfsetlinewidth{0.000000pt}%
\definecolor{currentstroke}{rgb}{0.000000,0.000000,0.000000}%
\pgfsetstrokecolor{currentstroke}%
\pgfsetdash{}{0pt}%
\pgfpathmoveto{\pgfqpoint{1.381184in}{0.860254in}}%
\pgfpathlineto{\pgfqpoint{1.390343in}{0.868115in}}%
\pgfpathlineto{\pgfqpoint{1.399540in}{0.876009in}}%
\pgfpathlineto{\pgfqpoint{1.424552in}{0.885152in}}%
\pgfpathlineto{\pgfqpoint{1.447894in}{0.894888in}}%
\pgfpathlineto{\pgfqpoint{1.437319in}{0.886480in}}%
\pgfpathlineto{\pgfqpoint{1.426782in}{0.878103in}}%
\pgfpathlineto{\pgfqpoint{1.404763in}{0.868896in}}%
\pgfpathlineto{\pgfqpoint{1.381184in}{0.860254in}}%
\pgfpathclose%
\pgfusepath{fill}%
\end{pgfscope}%
\begin{pgfscope}%
\pgfpathrectangle{\pgfqpoint{0.407592in}{0.222904in}}{\pgfqpoint{1.411163in}{1.411163in}}%
\pgfusepath{clip}%
\pgfsetbuttcap%
\pgfsetroundjoin%
\definecolor{currentfill}{rgb}{0.749020,0.772549,0.792157}%
\pgfsetfillcolor{currentfill}%
\pgfsetfillopacity{0.160000}%
\pgfsetlinewidth{0.000000pt}%
\definecolor{currentstroke}{rgb}{0.000000,0.000000,0.000000}%
\pgfsetstrokecolor{currentstroke}%
\pgfsetdash{}{0pt}%
\pgfpathmoveto{\pgfqpoint{0.936010in}{0.785625in}}%
\pgfpathlineto{\pgfqpoint{0.930134in}{0.792296in}}%
\pgfpathlineto{\pgfqpoint{0.924231in}{0.798997in}}%
\pgfpathlineto{\pgfqpoint{0.950134in}{0.795289in}}%
\pgfpathlineto{\pgfqpoint{0.976727in}{0.792313in}}%
\pgfpathlineto{\pgfqpoint{0.980771in}{0.785837in}}%
\pgfpathlineto{\pgfqpoint{0.984796in}{0.779391in}}%
\pgfpathlineto{\pgfqpoint{0.960084in}{0.782167in}}%
\pgfpathlineto{\pgfqpoint{0.936010in}{0.785625in}}%
\pgfpathclose%
\pgfusepath{fill}%
\end{pgfscope}%
\begin{pgfscope}%
\pgfpathrectangle{\pgfqpoint{0.407592in}{0.222904in}}{\pgfqpoint{1.411163in}{1.411163in}}%
\pgfusepath{clip}%
\pgfsetbuttcap%
\pgfsetroundjoin%
\definecolor{currentfill}{rgb}{0.749020,0.772549,0.792157}%
\pgfsetfillcolor{currentfill}%
\pgfsetfillopacity{0.160000}%
\pgfsetlinewidth{0.000000pt}%
\definecolor{currentstroke}{rgb}{0.000000,0.000000,0.000000}%
\pgfsetstrokecolor{currentstroke}%
\pgfsetdash{}{0pt}%
\pgfpathmoveto{\pgfqpoint{0.383806in}{1.143166in}}%
\pgfpathlineto{\pgfqpoint{0.369379in}{1.154043in}}%
\pgfpathlineto{\pgfqpoint{0.354919in}{1.164946in}}%
\pgfpathlineto{\pgfqpoint{0.364320in}{1.147219in}}%
\pgfpathlineto{\pgfqpoint{0.377119in}{1.129718in}}%
\pgfpathlineto{\pgfqpoint{0.391303in}{1.119391in}}%
\pgfpathlineto{\pgfqpoint{0.405447in}{1.109093in}}%
\pgfpathlineto{\pgfqpoint{0.392998in}{1.126016in}}%
\pgfpathlineto{\pgfqpoint{0.383806in}{1.143166in}}%
\pgfpathclose%
\pgfusepath{fill}%
\end{pgfscope}%
\begin{pgfscope}%
\pgfpathrectangle{\pgfqpoint{0.407592in}{0.222904in}}{\pgfqpoint{1.411163in}{1.411163in}}%
\pgfusepath{clip}%
\pgfsetbuttcap%
\pgfsetroundjoin%
\definecolor{currentfill}{rgb}{0.749020,0.772549,0.792157}%
\pgfsetfillcolor{currentfill}%
\pgfsetfillopacity{0.160000}%
\pgfsetlinewidth{0.000000pt}%
\definecolor{currentstroke}{rgb}{0.000000,0.000000,0.000000}%
\pgfsetstrokecolor{currentstroke}%
\pgfsetdash{}{0pt}%
\pgfpathmoveto{\pgfqpoint{1.250484in}{0.804387in}}%
\pgfpathlineto{\pgfqpoint{1.256200in}{0.811300in}}%
\pgfpathlineto{\pgfqpoint{1.261941in}{0.818245in}}%
\pgfpathlineto{\pgfqpoint{1.288832in}{0.823836in}}%
\pgfpathlineto{\pgfqpoint{1.314716in}{0.830120in}}%
\pgfpathlineto{\pgfqpoint{1.307236in}{0.822803in}}%
\pgfpathlineto{\pgfqpoint{1.299788in}{0.815517in}}%
\pgfpathlineto{\pgfqpoint{1.275602in}{0.809626in}}%
\pgfpathlineto{\pgfqpoint{1.250484in}{0.804387in}}%
\pgfpathclose%
\pgfusepath{fill}%
\end{pgfscope}%
\begin{pgfscope}%
\pgfpathrectangle{\pgfqpoint{0.407592in}{0.222904in}}{\pgfqpoint{1.411163in}{1.411163in}}%
\pgfusepath{clip}%
\pgfsetbuttcap%
\pgfsetroundjoin%
\definecolor{currentfill}{rgb}{0.749020,0.772549,0.792157}%
\pgfsetfillcolor{currentfill}%
\pgfsetfillopacity{0.160000}%
\pgfsetlinewidth{0.000000pt}%
\definecolor{currentstroke}{rgb}{0.000000,0.000000,0.000000}%
\pgfsetstrokecolor{currentstroke}%
\pgfsetdash{}{0pt}%
\pgfpathmoveto{\pgfqpoint{1.314716in}{0.830120in}}%
\pgfpathlineto{\pgfqpoint{1.322228in}{0.837468in}}%
\pgfpathlineto{\pgfqpoint{1.329773in}{0.844848in}}%
\pgfpathlineto{\pgfqpoint{1.356149in}{0.852223in}}%
\pgfpathlineto{\pgfqpoint{1.381184in}{0.860254in}}%
\pgfpathlineto{\pgfqpoint{1.372061in}{0.852424in}}%
\pgfpathlineto{\pgfqpoint{1.362974in}{0.844625in}}%
\pgfpathlineto{\pgfqpoint{1.339469in}{0.837062in}}%
\pgfpathlineto{\pgfqpoint{1.314716in}{0.830120in}}%
\pgfpathclose%
\pgfusepath{fill}%
\end{pgfscope}%
\begin{pgfscope}%
\pgfpathrectangle{\pgfqpoint{0.407592in}{0.222904in}}{\pgfqpoint{1.411163in}{1.411163in}}%
\pgfusepath{clip}%
\pgfsetbuttcap%
\pgfsetroundjoin%
\definecolor{currentfill}{rgb}{0.749020,0.772549,0.792157}%
\pgfsetfillcolor{currentfill}%
\pgfsetfillopacity{0.160000}%
\pgfsetlinewidth{0.000000pt}%
\definecolor{currentstroke}{rgb}{0.000000,0.000000,0.000000}%
\pgfsetstrokecolor{currentstroke}%
\pgfsetdash{}{0pt}%
\pgfpathmoveto{\pgfqpoint{1.139422in}{0.777674in}}%
\pgfpathlineto{\pgfqpoint{1.141300in}{0.784057in}}%
\pgfpathlineto{\pgfqpoint{1.143188in}{0.790471in}}%
\pgfpathlineto{\pgfqpoint{1.170774in}{0.792850in}}%
\pgfpathlineto{\pgfqpoint{1.197941in}{0.795975in}}%
\pgfpathlineto{\pgfqpoint{1.194108in}{0.789375in}}%
\pgfpathlineto{\pgfqpoint{1.190292in}{0.782806in}}%
\pgfpathlineto{\pgfqpoint{1.165049in}{0.779892in}}%
\pgfpathlineto{\pgfqpoint{1.139422in}{0.777674in}}%
\pgfpathclose%
\pgfusepath{fill}%
\end{pgfscope}%
\begin{pgfscope}%
\pgfpathrectangle{\pgfqpoint{0.407592in}{0.222904in}}{\pgfqpoint{1.411163in}{1.411163in}}%
\pgfusepath{clip}%
\pgfsetbuttcap%
\pgfsetroundjoin%
\definecolor{currentfill}{rgb}{0.749020,0.772549,0.792157}%
\pgfsetfillcolor{currentfill}%
\pgfsetfillopacity{0.160000}%
\pgfsetlinewidth{0.000000pt}%
\definecolor{currentstroke}{rgb}{0.000000,0.000000,0.000000}%
\pgfsetstrokecolor{currentstroke}%
\pgfsetdash{}{0pt}%
\pgfpathmoveto{\pgfqpoint{1.469162in}{0.911797in}}%
\pgfpathlineto{\pgfqpoint{1.479856in}{0.920299in}}%
\pgfpathlineto{\pgfqpoint{1.490589in}{0.928832in}}%
\pgfpathlineto{\pgfqpoint{1.514577in}{0.940213in}}%
\pgfpathlineto{\pgfqpoint{1.536460in}{0.952122in}}%
\pgfpathlineto{\pgfqpoint{1.524581in}{0.943028in}}%
\pgfpathlineto{\pgfqpoint{1.512741in}{0.933965in}}%
\pgfpathlineto{\pgfqpoint{1.491942in}{0.922627in}}%
\pgfpathlineto{\pgfqpoint{1.469162in}{0.911797in}}%
\pgfpathclose%
\pgfusepath{fill}%
\end{pgfscope}%
\begin{pgfscope}%
\pgfpathrectangle{\pgfqpoint{0.407592in}{0.222904in}}{\pgfqpoint{1.411163in}{1.411163in}}%
\pgfusepath{clip}%
\pgfsetbuttcap%
\pgfsetroundjoin%
\definecolor{currentfill}{rgb}{0.749020,0.772549,0.792157}%
\pgfsetfillcolor{currentfill}%
\pgfsetfillopacity{0.160000}%
\pgfsetlinewidth{0.000000pt}%
\definecolor{currentstroke}{rgb}{0.000000,0.000000,0.000000}%
\pgfsetstrokecolor{currentstroke}%
\pgfsetdash{}{0pt}%
\pgfpathmoveto{\pgfqpoint{0.684723in}{0.903063in}}%
\pgfpathlineto{\pgfqpoint{0.672719in}{0.911722in}}%
\pgfpathlineto{\pgfqpoint{0.660673in}{0.920412in}}%
\pgfpathlineto{\pgfqpoint{0.681356in}{0.909691in}}%
\pgfpathlineto{\pgfqpoint{0.704041in}{0.899547in}}%
\pgfpathlineto{\pgfqpoint{0.714938in}{0.891391in}}%
\pgfpathlineto{\pgfqpoint{0.725793in}{0.883268in}}%
\pgfpathlineto{\pgfqpoint{0.704319in}{0.892890in}}%
\pgfpathlineto{\pgfqpoint{0.684723in}{0.903063in}}%
\pgfpathclose%
\pgfusepath{fill}%
\end{pgfscope}%
\begin{pgfscope}%
\pgfpathrectangle{\pgfqpoint{0.407592in}{0.222904in}}{\pgfqpoint{1.411163in}{1.411163in}}%
\pgfusepath{clip}%
\pgfsetbuttcap%
\pgfsetroundjoin%
\definecolor{currentfill}{rgb}{0.749020,0.772549,0.792157}%
\pgfsetfillcolor{currentfill}%
\pgfsetfillopacity{0.160000}%
\pgfsetlinewidth{0.000000pt}%
\definecolor{currentstroke}{rgb}{0.000000,0.000000,0.000000}%
\pgfsetstrokecolor{currentstroke}%
\pgfsetdash{}{0pt}%
\pgfpathmoveto{\pgfqpoint{0.599623in}{0.961980in}}%
\pgfpathlineto{\pgfqpoint{0.586560in}{0.971338in}}%
\pgfpathlineto{\pgfqpoint{0.573455in}{0.980725in}}%
\pgfpathlineto{\pgfqpoint{0.591547in}{0.967857in}}%
\pgfpathlineto{\pgfqpoint{0.612057in}{0.955483in}}%
\pgfpathlineto{\pgfqpoint{0.624276in}{0.946668in}}%
\pgfpathlineto{\pgfqpoint{0.636452in}{0.937885in}}%
\pgfpathlineto{\pgfqpoint{0.616895in}{0.949694in}}%
\pgfpathlineto{\pgfqpoint{0.599623in}{0.961980in}}%
\pgfpathclose%
\pgfusepath{fill}%
\end{pgfscope}%
\begin{pgfscope}%
\pgfpathrectangle{\pgfqpoint{0.407592in}{0.222904in}}{\pgfqpoint{1.411163in}{1.411163in}}%
\pgfusepath{clip}%
\pgfsetbuttcap%
\pgfsetroundjoin%
\definecolor{currentfill}{rgb}{0.749020,0.772549,0.792157}%
\pgfsetfillcolor{currentfill}%
\pgfsetfillopacity{0.160000}%
\pgfsetlinewidth{0.000000pt}%
\definecolor{currentstroke}{rgb}{0.000000,0.000000,0.000000}%
\pgfsetstrokecolor{currentstroke}%
\pgfsetdash{}{0pt}%
\pgfpathmoveto{\pgfqpoint{1.631120in}{1.021447in}}%
\pgfpathlineto{\pgfqpoint{1.643752in}{1.031044in}}%
\pgfpathlineto{\pgfqpoint{1.656425in}{1.040672in}}%
\pgfpathlineto{\pgfqpoint{1.677111in}{1.055814in}}%
\pgfpathlineto{\pgfqpoint{1.694940in}{1.071325in}}%
\pgfpathlineto{\pgfqpoint{1.681534in}{1.061110in}}%
\pgfpathlineto{\pgfqpoint{1.668165in}{1.050923in}}%
\pgfpathlineto{\pgfqpoint{1.651003in}{1.036004in}}%
\pgfpathlineto{\pgfqpoint{1.631120in}{1.021447in}}%
\pgfpathclose%
\pgfusepath{fill}%
\end{pgfscope}%
\begin{pgfscope}%
\pgfpathrectangle{\pgfqpoint{0.407592in}{0.222904in}}{\pgfqpoint{1.411163in}{1.411163in}}%
\pgfusepath{clip}%
\pgfsetbuttcap%
\pgfsetroundjoin%
\definecolor{currentfill}{rgb}{0.749020,0.772549,0.792157}%
\pgfsetfillcolor{currentfill}%
\pgfsetfillopacity{0.160000}%
\pgfsetlinewidth{0.000000pt}%
\definecolor{currentstroke}{rgb}{0.000000,0.000000,0.000000}%
\pgfsetstrokecolor{currentstroke}%
\pgfsetdash{}{0pt}%
\pgfpathmoveto{\pgfqpoint{0.984796in}{0.779391in}}%
\pgfpathlineto{\pgfqpoint{0.980771in}{0.785837in}}%
\pgfpathlineto{\pgfqpoint{0.976727in}{0.792313in}}%
\pgfpathlineto{\pgfqpoint{1.003880in}{0.790085in}}%
\pgfpathlineto{\pgfqpoint{1.031455in}{0.788619in}}%
\pgfpathlineto{\pgfqpoint{1.033551in}{0.782268in}}%
\pgfpathlineto{\pgfqpoint{1.035637in}{0.775947in}}%
\pgfpathlineto{\pgfqpoint{1.010022in}{0.777314in}}%
\pgfpathlineto{\pgfqpoint{0.984796in}{0.779391in}}%
\pgfpathclose%
\pgfusepath{fill}%
\end{pgfscope}%
\begin{pgfscope}%
\pgfpathrectangle{\pgfqpoint{0.407592in}{0.222904in}}{\pgfqpoint{1.411163in}{1.411163in}}%
\pgfusepath{clip}%
\pgfsetbuttcap%
\pgfsetroundjoin%
\definecolor{currentfill}{rgb}{0.749020,0.772549,0.792157}%
\pgfsetfillcolor{currentfill}%
\pgfsetfillopacity{0.160000}%
\pgfsetlinewidth{0.000000pt}%
\definecolor{currentstroke}{rgb}{0.000000,0.000000,0.000000}%
\pgfsetstrokecolor{currentstroke}%
\pgfsetdash{}{0pt}%
\pgfpathmoveto{\pgfqpoint{1.776149in}{1.133206in}}%
\pgfpathlineto{\pgfqpoint{1.789814in}{1.143619in}}%
\pgfpathlineto{\pgfqpoint{1.803517in}{1.154061in}}%
\pgfpathlineto{\pgfqpoint{1.820459in}{1.172186in}}%
\pgfpathlineto{\pgfqpoint{1.833914in}{1.190518in}}%
\pgfpathlineto{\pgfqpoint{1.819795in}{1.179507in}}%
\pgfpathlineto{\pgfqpoint{1.805709in}{1.168522in}}%
\pgfpathlineto{\pgfqpoint{1.792602in}{1.150760in}}%
\pgfpathlineto{\pgfqpoint{1.776149in}{1.133206in}}%
\pgfpathclose%
\pgfusepath{fill}%
\end{pgfscope}%
\begin{pgfscope}%
\pgfpathrectangle{\pgfqpoint{0.407592in}{0.222904in}}{\pgfqpoint{1.411163in}{1.411163in}}%
\pgfusepath{clip}%
\pgfsetbuttcap%
\pgfsetroundjoin%
\definecolor{currentfill}{rgb}{0.749020,0.772549,0.792157}%
\pgfsetfillcolor{currentfill}%
\pgfsetfillopacity{0.160000}%
\pgfsetlinewidth{0.000000pt}%
\definecolor{currentstroke}{rgb}{0.000000,0.000000,0.000000}%
\pgfsetstrokecolor{currentstroke}%
\pgfsetdash{}{0pt}%
\pgfpathmoveto{\pgfqpoint{1.087525in}{0.775370in}}%
\pgfpathlineto{\pgfqpoint{1.087422in}{0.781670in}}%
\pgfpathlineto{\pgfqpoint{1.087318in}{0.788000in}}%
\pgfpathlineto{\pgfqpoint{1.115323in}{0.788852in}}%
\pgfpathlineto{\pgfqpoint{1.143188in}{0.790471in}}%
\pgfpathlineto{\pgfqpoint{1.141300in}{0.784057in}}%
\pgfpathlineto{\pgfqpoint{1.139422in}{0.777674in}}%
\pgfpathlineto{\pgfqpoint{1.113537in}{0.776164in}}%
\pgfpathlineto{\pgfqpoint{1.087525in}{0.775370in}}%
\pgfpathclose%
\pgfusepath{fill}%
\end{pgfscope}%
\begin{pgfscope}%
\pgfpathrectangle{\pgfqpoint{0.407592in}{0.222904in}}{\pgfqpoint{1.411163in}{1.411163in}}%
\pgfusepath{clip}%
\pgfsetbuttcap%
\pgfsetroundjoin%
\definecolor{currentfill}{rgb}{0.749020,0.772549,0.792157}%
\pgfsetfillcolor{currentfill}%
\pgfsetfillopacity{0.160000}%
\pgfsetlinewidth{0.000000pt}%
\definecolor{currentstroke}{rgb}{0.000000,0.000000,0.000000}%
\pgfsetstrokecolor{currentstroke}%
\pgfsetdash{}{0pt}%
\pgfpathmoveto{\pgfqpoint{1.560339in}{0.970401in}}%
\pgfpathlineto{\pgfqpoint{1.572338in}{0.979586in}}%
\pgfpathlineto{\pgfqpoint{1.584377in}{0.988802in}}%
\pgfpathlineto{\pgfqpoint{1.605974in}{1.002344in}}%
\pgfpathlineto{\pgfqpoint{1.593461in}{0.992838in}}%
\pgfpathlineto{\pgfqpoint{1.580986in}{0.983361in}}%
\pgfusepath{fill}%
\end{pgfscope}%
\begin{pgfscope}%
\pgfpathrectangle{\pgfqpoint{0.407592in}{0.222904in}}{\pgfqpoint{1.411163in}{1.411163in}}%
\pgfusepath{clip}%
\pgfsetbuttcap%
\pgfsetroundjoin%
\definecolor{currentfill}{rgb}{0.749020,0.772549,0.792157}%
\pgfsetfillcolor{currentfill}%
\pgfsetfillopacity{0.160000}%
\pgfsetlinewidth{0.000000pt}%
\definecolor{currentstroke}{rgb}{0.000000,0.000000,0.000000}%
\pgfsetstrokecolor{currentstroke}%
\pgfsetdash{}{0pt}%
\pgfpathmoveto{\pgfqpoint{1.035637in}{0.775947in}}%
\pgfpathlineto{\pgfqpoint{1.033551in}{0.782268in}}%
\pgfpathlineto{\pgfqpoint{1.031455in}{0.788619in}}%
\pgfpathlineto{\pgfqpoint{1.059315in}{0.787923in}}%
\pgfpathlineto{\pgfqpoint{1.087318in}{0.788000in}}%
\pgfpathlineto{\pgfqpoint{1.087422in}{0.781670in}}%
\pgfpathlineto{\pgfqpoint{1.087525in}{0.775370in}}%
\pgfpathlineto{\pgfqpoint{1.061515in}{0.775298in}}%
\pgfpathlineto{\pgfqpoint{1.035637in}{0.775947in}}%
\pgfpathclose%
\pgfusepath{fill}%
\end{pgfscope}%
\begin{pgfscope}%
\pgfpathrectangle{\pgfqpoint{0.407592in}{0.222904in}}{\pgfqpoint{1.411163in}{1.411163in}}%
\pgfusepath{clip}%
\pgfsetbuttcap%
\pgfsetroundjoin%
\definecolor{currentfill}{rgb}{0.749020,0.772549,0.792157}%
\pgfsetfillcolor{currentfill}%
\pgfsetfillopacity{0.160000}%
\pgfsetlinewidth{0.000000pt}%
\definecolor{currentstroke}{rgb}{0.000000,0.000000,0.000000}%
\pgfsetstrokecolor{currentstroke}%
\pgfsetdash{}{0pt}%
\pgfpathmoveto{\pgfqpoint{0.875004in}{0.808522in}}%
\pgfpathlineto{\pgfqpoint{0.867338in}{0.815574in}}%
\pgfpathlineto{\pgfqpoint{0.859638in}{0.822658in}}%
\pgfpathlineto{\pgfqpoint{0.885483in}{0.817210in}}%
\pgfpathlineto{\pgfqpoint{0.912344in}{0.812494in}}%
\pgfpathlineto{\pgfqpoint{0.918301in}{0.805730in}}%
\pgfpathlineto{\pgfqpoint{0.924231in}{0.798997in}}%
\pgfpathlineto{\pgfqpoint{0.899148in}{0.803417in}}%
\pgfpathlineto{\pgfqpoint{0.875004in}{0.808522in}}%
\pgfpathclose%
\pgfusepath{fill}%
\end{pgfscope}%
\begin{pgfscope}%
\pgfpathrectangle{\pgfqpoint{0.407592in}{0.222904in}}{\pgfqpoint{1.411163in}{1.411163in}}%
\pgfusepath{clip}%
\pgfsetbuttcap%
\pgfsetroundjoin%
\definecolor{currentfill}{rgb}{0.749020,0.772549,0.792157}%
\pgfsetfillcolor{currentfill}%
\pgfsetfillopacity{0.160000}%
\pgfsetlinewidth{0.000000pt}%
\definecolor{currentstroke}{rgb}{0.000000,0.000000,0.000000}%
\pgfsetstrokecolor{currentstroke}%
\pgfsetdash{}{0pt}%
\pgfpathmoveto{\pgfqpoint{0.489491in}{1.047904in}}%
\pgfpathlineto{\pgfqpoint{0.475581in}{1.058032in}}%
\pgfpathlineto{\pgfqpoint{0.461632in}{1.068187in}}%
\pgfpathlineto{\pgfqpoint{0.476317in}{1.052744in}}%
\pgfpathlineto{\pgfqpoint{0.493941in}{1.037685in}}%
\pgfpathlineto{\pgfqpoint{0.507300in}{1.028116in}}%
\pgfpathlineto{\pgfqpoint{0.520615in}{1.018577in}}%
\pgfpathlineto{\pgfqpoint{0.503654in}{1.033053in}}%
\pgfpathlineto{\pgfqpoint{0.489491in}{1.047904in}}%
\pgfpathclose%
\pgfusepath{fill}%
\end{pgfscope}%
\begin{pgfscope}%
\pgfpathrectangle{\pgfqpoint{0.407592in}{0.222904in}}{\pgfqpoint{1.411163in}{1.411163in}}%
\pgfusepath{clip}%
\pgfsetbuttcap%
\pgfsetroundjoin%
\definecolor{currentfill}{rgb}{0.749020,0.772549,0.792157}%
\pgfsetfillcolor{currentfill}%
\pgfsetfillopacity{0.160000}%
\pgfsetlinewidth{0.000000pt}%
\definecolor{currentstroke}{rgb}{0.000000,0.000000,0.000000}%
\pgfsetstrokecolor{currentstroke}%
\pgfsetdash{}{0pt}%
\pgfpathmoveto{\pgfqpoint{0.354919in}{1.164946in}}%
\pgfpathlineto{\pgfqpoint{0.340425in}{1.175874in}}%
\pgfpathlineto{\pgfqpoint{0.325897in}{1.186828in}}%
\pgfpathlineto{\pgfqpoint{0.335494in}{1.168530in}}%
\pgfpathlineto{\pgfqpoint{0.348632in}{1.150458in}}%
\pgfpathlineto{\pgfqpoint{0.362895in}{1.140073in}}%
\pgfpathlineto{\pgfqpoint{0.377119in}{1.129718in}}%
\pgfpathlineto{\pgfqpoint{0.364320in}{1.147219in}}%
\pgfpathlineto{\pgfqpoint{0.354919in}{1.164946in}}%
\pgfpathclose%
\pgfusepath{fill}%
\end{pgfscope}%
\begin{pgfscope}%
\pgfpathrectangle{\pgfqpoint{0.407592in}{0.222904in}}{\pgfqpoint{1.411163in}{1.411163in}}%
\pgfusepath{clip}%
\pgfsetbuttcap%
\pgfsetroundjoin%
\definecolor{currentfill}{rgb}{0.749020,0.772549,0.792157}%
\pgfsetfillcolor{currentfill}%
\pgfsetfillopacity{0.160000}%
\pgfsetlinewidth{0.000000pt}%
\definecolor{currentstroke}{rgb}{0.000000,0.000000,0.000000}%
\pgfsetstrokecolor{currentstroke}%
\pgfsetdash{}{0pt}%
\pgfpathmoveto{\pgfqpoint{0.747376in}{0.867116in}}%
\pgfpathlineto{\pgfqpoint{0.736606in}{0.875176in}}%
\pgfpathlineto{\pgfqpoint{0.725793in}{0.883268in}}%
\pgfpathlineto{\pgfqpoint{0.749055in}{0.874251in}}%
\pgfpathlineto{\pgfqpoint{0.774003in}{0.865887in}}%
\pgfpathlineto{\pgfqpoint{0.783433in}{0.858272in}}%
\pgfpathlineto{\pgfqpoint{0.792823in}{0.850689in}}%
\pgfpathlineto{\pgfqpoint{0.769312in}{0.858592in}}%
\pgfpathlineto{\pgfqpoint{0.747376in}{0.867116in}}%
\pgfpathclose%
\pgfusepath{fill}%
\end{pgfscope}%
\begin{pgfscope}%
\pgfpathrectangle{\pgfqpoint{0.407592in}{0.222904in}}{\pgfqpoint{1.411163in}{1.411163in}}%
\pgfusepath{clip}%
\pgfsetbuttcap%
\pgfsetroundjoin%
\definecolor{currentfill}{rgb}{0.749020,0.772549,0.792157}%
\pgfsetfillcolor{currentfill}%
\pgfsetfillopacity{0.160000}%
\pgfsetlinewidth{0.000000pt}%
\definecolor{currentstroke}{rgb}{0.000000,0.000000,0.000000}%
\pgfsetstrokecolor{currentstroke}%
\pgfsetdash{}{0pt}%
\pgfpathmoveto{\pgfqpoint{0.811487in}{0.835619in}}%
\pgfpathlineto{\pgfqpoint{0.802175in}{0.843138in}}%
\pgfpathlineto{\pgfqpoint{0.792823in}{0.850689in}}%
\pgfpathlineto{\pgfqpoint{0.817803in}{0.843452in}}%
\pgfpathlineto{\pgfqpoint{0.844134in}{0.836919in}}%
\pgfpathlineto{\pgfqpoint{0.851903in}{0.829773in}}%
\pgfpathlineto{\pgfqpoint{0.859638in}{0.822658in}}%
\pgfpathlineto{\pgfqpoint{0.834933in}{0.828805in}}%
\pgfpathlineto{\pgfqpoint{0.811487in}{0.835619in}}%
\pgfpathclose%
\pgfusepath{fill}%
\end{pgfscope}%
\begin{pgfscope}%
\pgfpathrectangle{\pgfqpoint{0.407592in}{0.222904in}}{\pgfqpoint{1.411163in}{1.411163in}}%
\pgfusepath{clip}%
\pgfsetbuttcap%
\pgfsetroundjoin%
\definecolor{currentfill}{rgb}{0.749020,0.772549,0.792157}%
\pgfsetfillcolor{currentfill}%
\pgfsetfillopacity{0.160000}%
\pgfsetlinewidth{0.000000pt}%
\definecolor{currentstroke}{rgb}{0.000000,0.000000,0.000000}%
\pgfsetstrokecolor{currentstroke}%
\pgfsetdash{}{0pt}%
\pgfpathmoveto{\pgfqpoint{1.197941in}{0.795975in}}%
\pgfpathlineto{\pgfqpoint{1.201793in}{0.802606in}}%
\pgfpathlineto{\pgfqpoint{1.205662in}{0.809268in}}%
\pgfpathlineto{\pgfqpoint{1.234172in}{0.813380in}}%
\pgfpathlineto{\pgfqpoint{1.261941in}{0.818245in}}%
\pgfpathlineto{\pgfqpoint{1.256200in}{0.811300in}}%
\pgfpathlineto{\pgfqpoint{1.250484in}{0.804387in}}%
\pgfpathlineto{\pgfqpoint{1.224555in}{0.799828in}}%
\pgfpathlineto{\pgfqpoint{1.197941in}{0.795975in}}%
\pgfpathclose%
\pgfusepath{fill}%
\end{pgfscope}%
\begin{pgfscope}%
\pgfpathrectangle{\pgfqpoint{0.407592in}{0.222904in}}{\pgfqpoint{1.411163in}{1.411163in}}%
\pgfusepath{clip}%
\pgfsetbuttcap%
\pgfsetroundjoin%
\definecolor{currentfill}{rgb}{0.749020,0.772549,0.792157}%
\pgfsetfillcolor{currentfill}%
\pgfsetfillopacity{0.160000}%
\pgfsetlinewidth{0.000000pt}%
\definecolor{currentstroke}{rgb}{0.000000,0.000000,0.000000}%
\pgfsetstrokecolor{currentstroke}%
\pgfsetdash{}{0pt}%
\pgfpathmoveto{\pgfqpoint{1.803517in}{1.154061in}}%
\pgfpathlineto{\pgfqpoint{1.817258in}{1.164532in}}%
\pgfpathlineto{\pgfqpoint{1.834441in}{1.182940in}}%
\pgfpathlineto{\pgfqpoint{1.848065in}{1.201554in}}%
\pgfpathlineto{\pgfqpoint{1.833914in}{1.190518in}}%
\pgfpathlineto{\pgfqpoint{1.820459in}{1.172186in}}%
\pgfusepath{fill}%
\end{pgfscope}%
\begin{pgfscope}%
\pgfpathrectangle{\pgfqpoint{0.407592in}{0.222904in}}{\pgfqpoint{1.411163in}{1.411163in}}%
\pgfusepath{clip}%
\pgfsetbuttcap%
\pgfsetroundjoin%
\definecolor{currentfill}{rgb}{0.749020,0.772549,0.792157}%
\pgfsetfillcolor{currentfill}%
\pgfsetfillopacity{0.160000}%
\pgfsetlinewidth{0.000000pt}%
\definecolor{currentstroke}{rgb}{0.000000,0.000000,0.000000}%
\pgfsetstrokecolor{currentstroke}%
\pgfsetdash{}{0pt}%
\pgfpathmoveto{\pgfqpoint{1.399540in}{0.876009in}}%
\pgfpathlineto{\pgfqpoint{1.408774in}{0.883934in}}%
\pgfpathlineto{\pgfqpoint{1.418045in}{0.891891in}}%
\pgfpathlineto{\pgfqpoint{1.444495in}{0.901533in}}%
\pgfpathlineto{\pgfqpoint{1.469162in}{0.911797in}}%
\pgfpathlineto{\pgfqpoint{1.458509in}{0.903327in}}%
\pgfpathlineto{\pgfqpoint{1.447894in}{0.894888in}}%
\pgfpathlineto{\pgfqpoint{1.424552in}{0.885152in}}%
\pgfpathlineto{\pgfqpoint{1.399540in}{0.876009in}}%
\pgfpathclose%
\pgfusepath{fill}%
\end{pgfscope}%
\begin{pgfscope}%
\pgfpathrectangle{\pgfqpoint{0.407592in}{0.222904in}}{\pgfqpoint{1.411163in}{1.411163in}}%
\pgfusepath{clip}%
\pgfsetbuttcap%
\pgfsetroundjoin%
\definecolor{currentfill}{rgb}{0.749020,0.772549,0.792157}%
\pgfsetfillcolor{currentfill}%
\pgfsetfillopacity{0.160000}%
\pgfsetlinewidth{0.000000pt}%
\definecolor{currentstroke}{rgb}{0.000000,0.000000,0.000000}%
\pgfsetstrokecolor{currentstroke}%
\pgfsetdash{}{0pt}%
\pgfpathmoveto{\pgfqpoint{1.656425in}{1.040672in}}%
\pgfpathlineto{\pgfqpoint{1.669138in}{1.050329in}}%
\pgfpathlineto{\pgfqpoint{1.681891in}{1.060018in}}%
\pgfpathlineto{\pgfqpoint{1.703375in}{1.075741in}}%
\pgfpathlineto{\pgfqpoint{1.721861in}{1.091839in}}%
\pgfpathlineto{\pgfqpoint{1.708382in}{1.081568in}}%
\pgfpathlineto{\pgfqpoint{1.694940in}{1.071325in}}%
\pgfpathlineto{\pgfqpoint{1.677111in}{1.055814in}}%
\pgfpathlineto{\pgfqpoint{1.656425in}{1.040672in}}%
\pgfpathclose%
\pgfusepath{fill}%
\end{pgfscope}%
\begin{pgfscope}%
\pgfpathrectangle{\pgfqpoint{0.407592in}{0.222904in}}{\pgfqpoint{1.411163in}{1.411163in}}%
\pgfusepath{clip}%
\pgfsetbuttcap%
\pgfsetroundjoin%
\definecolor{currentfill}{rgb}{0.749020,0.772549,0.792157}%
\pgfsetfillcolor{currentfill}%
\pgfsetfillopacity{0.160000}%
\pgfsetlinewidth{0.000000pt}%
\definecolor{currentstroke}{rgb}{0.000000,0.000000,0.000000}%
\pgfsetstrokecolor{currentstroke}%
\pgfsetdash{}{0pt}%
\pgfpathmoveto{\pgfqpoint{1.490589in}{0.928832in}}%
\pgfpathlineto{\pgfqpoint{1.501362in}{0.937397in}}%
\pgfpathlineto{\pgfqpoint{1.512175in}{0.945993in}}%
\pgfpathlineto{\pgfqpoint{1.537372in}{0.957923in}}%
\pgfpathlineto{\pgfqpoint{1.560339in}{0.970401in}}%
\pgfpathlineto{\pgfqpoint{1.548380in}{0.961246in}}%
\pgfpathlineto{\pgfqpoint{1.536460in}{0.952122in}}%
\pgfpathlineto{\pgfqpoint{1.514577in}{0.940213in}}%
\pgfpathlineto{\pgfqpoint{1.490589in}{0.928832in}}%
\pgfpathclose%
\pgfusepath{fill}%
\end{pgfscope}%
\begin{pgfscope}%
\pgfpathrectangle{\pgfqpoint{0.407592in}{0.222904in}}{\pgfqpoint{1.411163in}{1.411163in}}%
\pgfusepath{clip}%
\pgfsetbuttcap%
\pgfsetroundjoin%
\definecolor{currentfill}{rgb}{0.749020,0.772549,0.792157}%
\pgfsetfillcolor{currentfill}%
\pgfsetfillopacity{0.160000}%
\pgfsetlinewidth{0.000000pt}%
\definecolor{currentstroke}{rgb}{0.000000,0.000000,0.000000}%
\pgfsetstrokecolor{currentstroke}%
\pgfsetdash{}{0pt}%
\pgfpathmoveto{\pgfqpoint{0.573455in}{0.980725in}}%
\pgfpathlineto{\pgfqpoint{0.560308in}{0.990143in}}%
\pgfpathlineto{\pgfqpoint{0.547119in}{0.999591in}}%
\pgfpathlineto{\pgfqpoint{0.566026in}{0.986143in}}%
\pgfpathlineto{\pgfqpoint{0.587487in}{0.973208in}}%
\pgfpathlineto{\pgfqpoint{0.599794in}{0.964330in}}%
\pgfpathlineto{\pgfqpoint{0.612057in}{0.955483in}}%
\pgfpathlineto{\pgfqpoint{0.591547in}{0.967857in}}%
\pgfpathlineto{\pgfqpoint{0.573455in}{0.980725in}}%
\pgfpathclose%
\pgfusepath{fill}%
\end{pgfscope}%
\begin{pgfscope}%
\pgfpathrectangle{\pgfqpoint{0.407592in}{0.222904in}}{\pgfqpoint{1.411163in}{1.411163in}}%
\pgfusepath{clip}%
\pgfsetbuttcap%
\pgfsetroundjoin%
\definecolor{currentfill}{rgb}{0.749020,0.772549,0.792157}%
\pgfsetfillcolor{currentfill}%
\pgfsetfillopacity{0.160000}%
\pgfsetlinewidth{0.000000pt}%
\definecolor{currentstroke}{rgb}{0.000000,0.000000,0.000000}%
\pgfsetstrokecolor{currentstroke}%
\pgfsetdash{}{0pt}%
\pgfpathmoveto{\pgfqpoint{0.924231in}{0.798997in}}%
\pgfpathlineto{\pgfqpoint{0.918301in}{0.805730in}}%
\pgfpathlineto{\pgfqpoint{0.912344in}{0.812494in}}%
\pgfpathlineto{\pgfqpoint{0.940090in}{0.808535in}}%
\pgfpathlineto{\pgfqpoint{0.968582in}{0.805358in}}%
\pgfpathlineto{\pgfqpoint{0.972665in}{0.798820in}}%
\pgfpathlineto{\pgfqpoint{0.976727in}{0.792313in}}%
\pgfpathlineto{\pgfqpoint{0.950134in}{0.795289in}}%
\pgfpathlineto{\pgfqpoint{0.924231in}{0.798997in}}%
\pgfpathclose%
\pgfusepath{fill}%
\end{pgfscope}%
\begin{pgfscope}%
\pgfpathrectangle{\pgfqpoint{0.407592in}{0.222904in}}{\pgfqpoint{1.411163in}{1.411163in}}%
\pgfusepath{clip}%
\pgfsetbuttcap%
\pgfsetroundjoin%
\definecolor{currentfill}{rgb}{0.749020,0.772549,0.792157}%
\pgfsetfillcolor{currentfill}%
\pgfsetfillopacity{0.160000}%
\pgfsetlinewidth{0.000000pt}%
\definecolor{currentstroke}{rgb}{0.000000,0.000000,0.000000}%
\pgfsetstrokecolor{currentstroke}%
\pgfsetdash{}{0pt}%
\pgfpathmoveto{\pgfqpoint{0.325897in}{1.186828in}}%
\pgfpathlineto{\pgfqpoint{0.311335in}{1.197808in}}%
\pgfpathlineto{\pgfqpoint{0.321025in}{1.179228in}}%
\pgfpathlineto{\pgfqpoint{0.334328in}{1.160872in}}%
\pgfpathlineto{\pgfqpoint{0.348632in}{1.150458in}}%
\pgfpathlineto{\pgfqpoint{0.335494in}{1.168530in}}%
\pgfusepath{fill}%
\end{pgfscope}%
\begin{pgfscope}%
\pgfpathrectangle{\pgfqpoint{0.407592in}{0.222904in}}{\pgfqpoint{1.411163in}{1.411163in}}%
\pgfusepath{clip}%
\pgfsetbuttcap%
\pgfsetroundjoin%
\definecolor{currentfill}{rgb}{0.749020,0.772549,0.792157}%
\pgfsetfillcolor{currentfill}%
\pgfsetfillopacity{0.160000}%
\pgfsetlinewidth{0.000000pt}%
\definecolor{currentstroke}{rgb}{0.000000,0.000000,0.000000}%
\pgfsetstrokecolor{currentstroke}%
\pgfsetdash{}{0pt}%
\pgfpathmoveto{\pgfqpoint{1.261941in}{0.818245in}}%
\pgfpathlineto{\pgfqpoint{1.267708in}{0.825222in}}%
\pgfpathlineto{\pgfqpoint{1.273502in}{0.832230in}}%
\pgfpathlineto{\pgfqpoint{1.302180in}{0.838171in}}%
\pgfpathlineto{\pgfqpoint{1.329773in}{0.844848in}}%
\pgfpathlineto{\pgfqpoint{1.322228in}{0.837468in}}%
\pgfpathlineto{\pgfqpoint{1.314716in}{0.830120in}}%
\pgfpathlineto{\pgfqpoint{1.288832in}{0.823836in}}%
\pgfpathlineto{\pgfqpoint{1.261941in}{0.818245in}}%
\pgfpathclose%
\pgfusepath{fill}%
\end{pgfscope}%
\begin{pgfscope}%
\pgfpathrectangle{\pgfqpoint{0.407592in}{0.222904in}}{\pgfqpoint{1.411163in}{1.411163in}}%
\pgfusepath{clip}%
\pgfsetbuttcap%
\pgfsetroundjoin%
\definecolor{currentfill}{rgb}{0.749020,0.772549,0.792157}%
\pgfsetfillcolor{currentfill}%
\pgfsetfillopacity{0.160000}%
\pgfsetlinewidth{0.000000pt}%
\definecolor{currentstroke}{rgb}{0.000000,0.000000,0.000000}%
\pgfsetstrokecolor{currentstroke}%
\pgfsetdash{}{0pt}%
\pgfpathmoveto{\pgfqpoint{1.329773in}{0.844848in}}%
\pgfpathlineto{\pgfqpoint{1.337350in}{0.852260in}}%
\pgfpathlineto{\pgfqpoint{1.344960in}{0.859704in}}%
\pgfpathlineto{\pgfqpoint{1.372969in}{0.867510in}}%
\pgfpathlineto{\pgfqpoint{1.399540in}{0.876009in}}%
\pgfpathlineto{\pgfqpoint{1.390343in}{0.868115in}}%
\pgfpathlineto{\pgfqpoint{1.381184in}{0.860254in}}%
\pgfpathlineto{\pgfqpoint{1.356149in}{0.852223in}}%
\pgfpathlineto{\pgfqpoint{1.329773in}{0.844848in}}%
\pgfpathclose%
\pgfusepath{fill}%
\end{pgfscope}%
\begin{pgfscope}%
\pgfpathrectangle{\pgfqpoint{0.407592in}{0.222904in}}{\pgfqpoint{1.411163in}{1.411163in}}%
\pgfusepath{clip}%
\pgfsetbuttcap%
\pgfsetroundjoin%
\definecolor{currentfill}{rgb}{0.749020,0.772549,0.792157}%
\pgfsetfillcolor{currentfill}%
\pgfsetfillopacity{0.160000}%
\pgfsetlinewidth{0.000000pt}%
\definecolor{currentstroke}{rgb}{0.000000,0.000000,0.000000}%
\pgfsetstrokecolor{currentstroke}%
\pgfsetdash{}{0pt}%
\pgfpathmoveto{\pgfqpoint{0.660673in}{0.920412in}}%
\pgfpathlineto{\pgfqpoint{0.648584in}{0.929133in}}%
\pgfpathlineto{\pgfqpoint{0.636452in}{0.937885in}}%
\pgfpathlineto{\pgfqpoint{0.658221in}{0.926618in}}%
\pgfpathlineto{\pgfqpoint{0.682118in}{0.915953in}}%
\pgfpathlineto{\pgfqpoint{0.693101in}{0.907734in}}%
\pgfpathlineto{\pgfqpoint{0.704041in}{0.899547in}}%
\pgfpathlineto{\pgfqpoint{0.681356in}{0.909691in}}%
\pgfpathlineto{\pgfqpoint{0.660673in}{0.920412in}}%
\pgfpathclose%
\pgfusepath{fill}%
\end{pgfscope}%
\begin{pgfscope}%
\pgfpathrectangle{\pgfqpoint{0.407592in}{0.222904in}}{\pgfqpoint{1.411163in}{1.411163in}}%
\pgfusepath{clip}%
\pgfsetbuttcap%
\pgfsetroundjoin%
\definecolor{currentfill}{rgb}{0.749020,0.772549,0.792157}%
\pgfsetfillcolor{currentfill}%
\pgfsetfillopacity{0.160000}%
\pgfsetlinewidth{0.000000pt}%
\definecolor{currentstroke}{rgb}{0.000000,0.000000,0.000000}%
\pgfsetstrokecolor{currentstroke}%
\pgfsetdash{}{0pt}%
\pgfpathmoveto{\pgfqpoint{1.584377in}{0.988802in}}%
\pgfpathlineto{\pgfqpoint{1.596457in}{0.998049in}}%
\pgfpathlineto{\pgfqpoint{1.608577in}{1.007328in}}%
\pgfpathlineto{\pgfqpoint{1.631120in}{1.021447in}}%
\pgfpathlineto{\pgfqpoint{1.618527in}{1.011881in}}%
\pgfpathlineto{\pgfqpoint{1.605974in}{1.002344in}}%
\pgfusepath{fill}%
\end{pgfscope}%
\begin{pgfscope}%
\pgfpathrectangle{\pgfqpoint{0.407592in}{0.222904in}}{\pgfqpoint{1.411163in}{1.411163in}}%
\pgfusepath{clip}%
\pgfsetbuttcap%
\pgfsetroundjoin%
\definecolor{currentfill}{rgb}{0.749020,0.772549,0.792157}%
\pgfsetfillcolor{currentfill}%
\pgfsetfillopacity{0.160000}%
\pgfsetlinewidth{0.000000pt}%
\definecolor{currentstroke}{rgb}{0.000000,0.000000,0.000000}%
\pgfsetstrokecolor{currentstroke}%
\pgfsetdash{}{0pt}%
\pgfpathmoveto{\pgfqpoint{0.461632in}{1.068187in}}%
\pgfpathlineto{\pgfqpoint{0.447645in}{1.078371in}}%
\pgfpathlineto{\pgfqpoint{0.433618in}{1.088583in}}%
\pgfpathlineto{\pgfqpoint{0.448815in}{1.072554in}}%
\pgfpathlineto{\pgfqpoint{0.467095in}{1.056916in}}%
\pgfpathlineto{\pgfqpoint{0.480540in}{1.047285in}}%
\pgfpathlineto{\pgfqpoint{0.493941in}{1.037685in}}%
\pgfpathlineto{\pgfqpoint{0.476317in}{1.052744in}}%
\pgfpathlineto{\pgfqpoint{0.461632in}{1.068187in}}%
\pgfpathclose%
\pgfusepath{fill}%
\end{pgfscope}%
\begin{pgfscope}%
\pgfpathrectangle{\pgfqpoint{0.407592in}{0.222904in}}{\pgfqpoint{1.411163in}{1.411163in}}%
\pgfusepath{clip}%
\pgfsetbuttcap%
\pgfsetroundjoin%
\definecolor{currentfill}{rgb}{0.749020,0.772549,0.792157}%
\pgfsetfillcolor{currentfill}%
\pgfsetfillopacity{0.160000}%
\pgfsetlinewidth{0.000000pt}%
\definecolor{currentstroke}{rgb}{0.000000,0.000000,0.000000}%
\pgfsetstrokecolor{currentstroke}%
\pgfsetdash{}{0pt}%
\pgfpathmoveto{\pgfqpoint{1.143188in}{0.790471in}}%
\pgfpathlineto{\pgfqpoint{1.145085in}{0.796916in}}%
\pgfpathlineto{\pgfqpoint{1.146991in}{0.803392in}}%
\pgfpathlineto{\pgfqpoint{1.176553in}{0.805932in}}%
\pgfpathlineto{\pgfqpoint{1.205662in}{0.809268in}}%
\pgfpathlineto{\pgfqpoint{1.201793in}{0.802606in}}%
\pgfpathlineto{\pgfqpoint{1.197941in}{0.795975in}}%
\pgfpathlineto{\pgfqpoint{1.170774in}{0.792850in}}%
\pgfpathlineto{\pgfqpoint{1.143188in}{0.790471in}}%
\pgfpathclose%
\pgfusepath{fill}%
\end{pgfscope}%
\begin{pgfscope}%
\pgfpathrectangle{\pgfqpoint{0.407592in}{0.222904in}}{\pgfqpoint{1.411163in}{1.411163in}}%
\pgfusepath{clip}%
\pgfsetbuttcap%
\pgfsetroundjoin%
\definecolor{currentfill}{rgb}{0.749020,0.772549,0.792157}%
\pgfsetfillcolor{currentfill}%
\pgfsetfillopacity{0.160000}%
\pgfsetlinewidth{0.000000pt}%
\definecolor{currentstroke}{rgb}{0.000000,0.000000,0.000000}%
\pgfsetstrokecolor{currentstroke}%
\pgfsetdash{}{0pt}%
\pgfpathmoveto{\pgfqpoint{0.976727in}{0.792313in}}%
\pgfpathlineto{\pgfqpoint{0.972665in}{0.798820in}}%
\pgfpathlineto{\pgfqpoint{0.968582in}{0.805358in}}%
\pgfpathlineto{\pgfqpoint{0.997679in}{0.802979in}}%
\pgfpathlineto{\pgfqpoint{1.027232in}{0.801414in}}%
\pgfpathlineto{\pgfqpoint{1.029349in}{0.795001in}}%
\pgfpathlineto{\pgfqpoint{1.031455in}{0.788619in}}%
\pgfpathlineto{\pgfqpoint{1.003880in}{0.790085in}}%
\pgfpathlineto{\pgfqpoint{0.976727in}{0.792313in}}%
\pgfpathclose%
\pgfusepath{fill}%
\end{pgfscope}%
\begin{pgfscope}%
\pgfpathrectangle{\pgfqpoint{0.407592in}{0.222904in}}{\pgfqpoint{1.411163in}{1.411163in}}%
\pgfusepath{clip}%
\pgfsetbuttcap%
\pgfsetroundjoin%
\definecolor{currentfill}{rgb}{0.749020,0.772549,0.792157}%
\pgfsetfillcolor{currentfill}%
\pgfsetfillopacity{0.160000}%
\pgfsetlinewidth{0.000000pt}%
\definecolor{currentstroke}{rgb}{0.000000,0.000000,0.000000}%
\pgfsetstrokecolor{currentstroke}%
\pgfsetdash{}{0pt}%
\pgfpathmoveto{\pgfqpoint{1.087318in}{0.788000in}}%
\pgfpathlineto{\pgfqpoint{1.087213in}{0.794362in}}%
\pgfpathlineto{\pgfqpoint{1.087109in}{0.800754in}}%
\pgfpathlineto{\pgfqpoint{1.117125in}{0.801662in}}%
\pgfpathlineto{\pgfqpoint{1.146991in}{0.803392in}}%
\pgfpathlineto{\pgfqpoint{1.145085in}{0.796916in}}%
\pgfpathlineto{\pgfqpoint{1.143188in}{0.790471in}}%
\pgfpathlineto{\pgfqpoint{1.115323in}{0.788852in}}%
\pgfpathlineto{\pgfqpoint{1.087318in}{0.788000in}}%
\pgfpathclose%
\pgfusepath{fill}%
\end{pgfscope}%
\begin{pgfscope}%
\pgfpathrectangle{\pgfqpoint{0.407592in}{0.222904in}}{\pgfqpoint{1.411163in}{1.411163in}}%
\pgfusepath{clip}%
\pgfsetbuttcap%
\pgfsetroundjoin%
\definecolor{currentfill}{rgb}{0.749020,0.772549,0.792157}%
\pgfsetfillcolor{currentfill}%
\pgfsetfillopacity{0.160000}%
\pgfsetlinewidth{0.000000pt}%
\definecolor{currentstroke}{rgb}{0.000000,0.000000,0.000000}%
\pgfsetstrokecolor{currentstroke}%
\pgfsetdash{}{0pt}%
\pgfpathmoveto{\pgfqpoint{0.859638in}{0.822658in}}%
\pgfpathlineto{\pgfqpoint{0.851903in}{0.829773in}}%
\pgfpathlineto{\pgfqpoint{0.844134in}{0.836919in}}%
\pgfpathlineto{\pgfqpoint{0.871693in}{0.831130in}}%
\pgfpathlineto{\pgfqpoint{0.900346in}{0.826116in}}%
\pgfpathlineto{\pgfqpoint{0.906359in}{0.819289in}}%
\pgfpathlineto{\pgfqpoint{0.912344in}{0.812494in}}%
\pgfpathlineto{\pgfqpoint{0.885483in}{0.817210in}}%
\pgfpathlineto{\pgfqpoint{0.859638in}{0.822658in}}%
\pgfpathclose%
\pgfusepath{fill}%
\end{pgfscope}%
\begin{pgfscope}%
\pgfpathrectangle{\pgfqpoint{0.407592in}{0.222904in}}{\pgfqpoint{1.411163in}{1.411163in}}%
\pgfusepath{clip}%
\pgfsetbuttcap%
\pgfsetroundjoin%
\definecolor{currentfill}{rgb}{0.749020,0.772549,0.792157}%
\pgfsetfillcolor{currentfill}%
\pgfsetfillopacity{0.160000}%
\pgfsetlinewidth{0.000000pt}%
\definecolor{currentstroke}{rgb}{0.000000,0.000000,0.000000}%
\pgfsetstrokecolor{currentstroke}%
\pgfsetdash{}{0pt}%
\pgfpathmoveto{\pgfqpoint{0.725793in}{0.883268in}}%
\pgfpathlineto{\pgfqpoint{0.714938in}{0.891391in}}%
\pgfpathlineto{\pgfqpoint{0.704041in}{0.899547in}}%
\pgfpathlineto{\pgfqpoint{0.728632in}{0.890037in}}%
\pgfpathlineto{\pgfqpoint{0.755023in}{0.881214in}}%
\pgfpathlineto{\pgfqpoint{0.764533in}{0.873534in}}%
\pgfpathlineto{\pgfqpoint{0.774003in}{0.865887in}}%
\pgfpathlineto{\pgfqpoint{0.749055in}{0.874251in}}%
\pgfpathlineto{\pgfqpoint{0.725793in}{0.883268in}}%
\pgfpathclose%
\pgfusepath{fill}%
\end{pgfscope}%
\begin{pgfscope}%
\pgfpathrectangle{\pgfqpoint{0.407592in}{0.222904in}}{\pgfqpoint{1.411163in}{1.411163in}}%
\pgfusepath{clip}%
\pgfsetbuttcap%
\pgfsetroundjoin%
\definecolor{currentfill}{rgb}{0.749020,0.772549,0.792157}%
\pgfsetfillcolor{currentfill}%
\pgfsetfillopacity{0.160000}%
\pgfsetlinewidth{0.000000pt}%
\definecolor{currentstroke}{rgb}{0.000000,0.000000,0.000000}%
\pgfsetstrokecolor{currentstroke}%
\pgfsetdash{}{0pt}%
\pgfpathmoveto{\pgfqpoint{1.031455in}{0.788619in}}%
\pgfpathlineto{\pgfqpoint{1.029349in}{0.795001in}}%
\pgfpathlineto{\pgfqpoint{1.027232in}{0.801414in}}%
\pgfpathlineto{\pgfqpoint{1.057093in}{0.800671in}}%
\pgfpathlineto{\pgfqpoint{1.087109in}{0.800754in}}%
\pgfpathlineto{\pgfqpoint{1.087213in}{0.794362in}}%
\pgfpathlineto{\pgfqpoint{1.087318in}{0.788000in}}%
\pgfpathlineto{\pgfqpoint{1.059315in}{0.787923in}}%
\pgfpathlineto{\pgfqpoint{1.031455in}{0.788619in}}%
\pgfpathclose%
\pgfusepath{fill}%
\end{pgfscope}%
\begin{pgfscope}%
\pgfpathrectangle{\pgfqpoint{0.407592in}{0.222904in}}{\pgfqpoint{1.411163in}{1.411163in}}%
\pgfusepath{clip}%
\pgfsetbuttcap%
\pgfsetroundjoin%
\definecolor{currentfill}{rgb}{0.749020,0.772549,0.792157}%
\pgfsetfillcolor{currentfill}%
\pgfsetfillopacity{0.160000}%
\pgfsetlinewidth{0.000000pt}%
\definecolor{currentstroke}{rgb}{0.000000,0.000000,0.000000}%
\pgfsetstrokecolor{currentstroke}%
\pgfsetdash{}{0pt}%
\pgfpathmoveto{\pgfqpoint{1.681891in}{1.060018in}}%
\pgfpathlineto{\pgfqpoint{1.694685in}{1.069737in}}%
\pgfpathlineto{\pgfqpoint{1.707519in}{1.079488in}}%
\pgfpathlineto{\pgfqpoint{1.729795in}{1.095786in}}%
\pgfpathlineto{\pgfqpoint{1.748930in}{1.112466in}}%
\pgfpathlineto{\pgfqpoint{1.735377in}{1.102138in}}%
\pgfpathlineto{\pgfqpoint{1.721861in}{1.091839in}}%
\pgfpathlineto{\pgfqpoint{1.703375in}{1.075741in}}%
\pgfpathlineto{\pgfqpoint{1.681891in}{1.060018in}}%
\pgfpathclose%
\pgfusepath{fill}%
\end{pgfscope}%
\begin{pgfscope}%
\pgfpathrectangle{\pgfqpoint{0.407592in}{0.222904in}}{\pgfqpoint{1.411163in}{1.411163in}}%
\pgfusepath{clip}%
\pgfsetbuttcap%
\pgfsetroundjoin%
\definecolor{currentfill}{rgb}{0.749020,0.772549,0.792157}%
\pgfsetfillcolor{currentfill}%
\pgfsetfillopacity{0.160000}%
\pgfsetlinewidth{0.000000pt}%
\definecolor{currentstroke}{rgb}{0.000000,0.000000,0.000000}%
\pgfsetstrokecolor{currentstroke}%
\pgfsetdash{}{0pt}%
\pgfpathmoveto{\pgfqpoint{1.418045in}{0.891891in}}%
\pgfpathlineto{\pgfqpoint{1.427353in}{0.899881in}}%
\pgfpathlineto{\pgfqpoint{1.436700in}{0.907902in}}%
\pgfpathlineto{\pgfqpoint{1.464593in}{0.918042in}}%
\pgfpathlineto{\pgfqpoint{1.490589in}{0.928832in}}%
\pgfpathlineto{\pgfqpoint{1.479856in}{0.920299in}}%
\pgfpathlineto{\pgfqpoint{1.469162in}{0.911797in}}%
\pgfpathlineto{\pgfqpoint{1.444495in}{0.901533in}}%
\pgfpathlineto{\pgfqpoint{1.418045in}{0.891891in}}%
\pgfpathclose%
\pgfusepath{fill}%
\end{pgfscope}%
\begin{pgfscope}%
\pgfpathrectangle{\pgfqpoint{0.407592in}{0.222904in}}{\pgfqpoint{1.411163in}{1.411163in}}%
\pgfusepath{clip}%
\pgfsetbuttcap%
\pgfsetroundjoin%
\definecolor{currentfill}{rgb}{0.749020,0.772549,0.792157}%
\pgfsetfillcolor{currentfill}%
\pgfsetfillopacity{0.160000}%
\pgfsetlinewidth{0.000000pt}%
\definecolor{currentstroke}{rgb}{0.000000,0.000000,0.000000}%
\pgfsetstrokecolor{currentstroke}%
\pgfsetdash{}{0pt}%
\pgfpathmoveto{\pgfqpoint{0.792823in}{0.850689in}}%
\pgfpathlineto{\pgfqpoint{0.783433in}{0.858272in}}%
\pgfpathlineto{\pgfqpoint{0.774003in}{0.865887in}}%
\pgfpathlineto{\pgfqpoint{0.800523in}{0.858225in}}%
\pgfpathlineto{\pgfqpoint{0.828493in}{0.851308in}}%
\pgfpathlineto{\pgfqpoint{0.836331in}{0.844098in}}%
\pgfpathlineto{\pgfqpoint{0.844134in}{0.836919in}}%
\pgfpathlineto{\pgfqpoint{0.817803in}{0.843452in}}%
\pgfpathlineto{\pgfqpoint{0.792823in}{0.850689in}}%
\pgfpathclose%
\pgfusepath{fill}%
\end{pgfscope}%
\begin{pgfscope}%
\pgfpathrectangle{\pgfqpoint{0.407592in}{0.222904in}}{\pgfqpoint{1.411163in}{1.411163in}}%
\pgfusepath{clip}%
\pgfsetbuttcap%
\pgfsetroundjoin%
\definecolor{currentfill}{rgb}{0.749020,0.772549,0.792157}%
\pgfsetfillcolor{currentfill}%
\pgfsetfillopacity{0.160000}%
\pgfsetlinewidth{0.000000pt}%
\definecolor{currentstroke}{rgb}{0.000000,0.000000,0.000000}%
\pgfsetstrokecolor{currentstroke}%
\pgfsetdash{}{0pt}%
\pgfpathmoveto{\pgfqpoint{1.512175in}{0.945993in}}%
\pgfpathlineto{\pgfqpoint{1.523029in}{0.954622in}}%
\pgfpathlineto{\pgfqpoint{1.533923in}{0.963284in}}%
\pgfpathlineto{\pgfqpoint{1.560330in}{0.975758in}}%
\pgfpathlineto{\pgfqpoint{1.584377in}{0.988802in}}%
\pgfpathlineto{\pgfqpoint{1.572338in}{0.979586in}}%
\pgfpathlineto{\pgfqpoint{1.560339in}{0.970401in}}%
\pgfpathlineto{\pgfqpoint{1.537372in}{0.957923in}}%
\pgfpathlineto{\pgfqpoint{1.512175in}{0.945993in}}%
\pgfpathclose%
\pgfusepath{fill}%
\end{pgfscope}%
\begin{pgfscope}%
\pgfpathrectangle{\pgfqpoint{0.407592in}{0.222904in}}{\pgfqpoint{1.411163in}{1.411163in}}%
\pgfusepath{clip}%
\pgfsetbuttcap%
\pgfsetroundjoin%
\definecolor{currentfill}{rgb}{0.749020,0.772549,0.792157}%
\pgfsetfillcolor{currentfill}%
\pgfsetfillopacity{0.160000}%
\pgfsetlinewidth{0.000000pt}%
\definecolor{currentstroke}{rgb}{0.000000,0.000000,0.000000}%
\pgfsetstrokecolor{currentstroke}%
\pgfsetdash{}{0pt}%
\pgfpathmoveto{\pgfqpoint{1.205662in}{0.809268in}}%
\pgfpathlineto{\pgfqpoint{1.209550in}{0.815961in}}%
\pgfpathlineto{\pgfqpoint{1.213456in}{0.822686in}}%
\pgfpathlineto{\pgfqpoint{1.243878in}{0.827058in}}%
\pgfpathlineto{\pgfqpoint{1.273502in}{0.832230in}}%
\pgfpathlineto{\pgfqpoint{1.267708in}{0.825222in}}%
\pgfpathlineto{\pgfqpoint{1.261941in}{0.818245in}}%
\pgfpathlineto{\pgfqpoint{1.234172in}{0.813380in}}%
\pgfpathlineto{\pgfqpoint{1.205662in}{0.809268in}}%
\pgfpathclose%
\pgfusepath{fill}%
\end{pgfscope}%
\begin{pgfscope}%
\pgfpathrectangle{\pgfqpoint{0.407592in}{0.222904in}}{\pgfqpoint{1.411163in}{1.411163in}}%
\pgfusepath{clip}%
\pgfsetbuttcap%
\pgfsetroundjoin%
\definecolor{currentfill}{rgb}{0.749020,0.772549,0.792157}%
\pgfsetfillcolor{currentfill}%
\pgfsetfillopacity{0.160000}%
\pgfsetlinewidth{0.000000pt}%
\definecolor{currentstroke}{rgb}{0.000000,0.000000,0.000000}%
\pgfsetstrokecolor{currentstroke}%
\pgfsetdash{}{0pt}%
\pgfpathmoveto{\pgfqpoint{0.547119in}{0.999591in}}%
\pgfpathlineto{\pgfqpoint{0.533889in}{1.009069in}}%
\pgfpathlineto{\pgfqpoint{0.520615in}{1.018577in}}%
\pgfpathlineto{\pgfqpoint{0.540331in}{1.004554in}}%
\pgfpathlineto{\pgfqpoint{0.562740in}{0.991060in}}%
\pgfpathlineto{\pgfqpoint{0.575136in}{0.982118in}}%
\pgfpathlineto{\pgfqpoint{0.587487in}{0.973208in}}%
\pgfpathlineto{\pgfqpoint{0.566026in}{0.986143in}}%
\pgfpathlineto{\pgfqpoint{0.547119in}{0.999591in}}%
\pgfpathclose%
\pgfusepath{fill}%
\end{pgfscope}%
\begin{pgfscope}%
\pgfpathrectangle{\pgfqpoint{0.407592in}{0.222904in}}{\pgfqpoint{1.411163in}{1.411163in}}%
\pgfusepath{clip}%
\pgfsetbuttcap%
\pgfsetroundjoin%
\definecolor{currentfill}{rgb}{0.749020,0.772549,0.792157}%
\pgfsetfillcolor{currentfill}%
\pgfsetfillopacity{0.160000}%
\pgfsetlinewidth{0.000000pt}%
\definecolor{currentstroke}{rgb}{0.000000,0.000000,0.000000}%
\pgfsetstrokecolor{currentstroke}%
\pgfsetdash{}{0pt}%
\pgfpathmoveto{\pgfqpoint{0.433618in}{1.088583in}}%
\pgfpathlineto{\pgfqpoint{0.419552in}{1.098824in}}%
\pgfpathlineto{\pgfqpoint{0.405447in}{1.109093in}}%
\pgfpathlineto{\pgfqpoint{0.421148in}{1.092483in}}%
\pgfpathlineto{\pgfqpoint{0.440076in}{1.076272in}}%
\pgfpathlineto{\pgfqpoint{0.453607in}{1.066578in}}%
\pgfpathlineto{\pgfqpoint{0.467095in}{1.056916in}}%
\pgfpathlineto{\pgfqpoint{0.448815in}{1.072554in}}%
\pgfpathlineto{\pgfqpoint{0.433618in}{1.088583in}}%
\pgfpathclose%
\pgfusepath{fill}%
\end{pgfscope}%
\begin{pgfscope}%
\pgfpathrectangle{\pgfqpoint{0.407592in}{0.222904in}}{\pgfqpoint{1.411163in}{1.411163in}}%
\pgfusepath{clip}%
\pgfsetbuttcap%
\pgfsetroundjoin%
\definecolor{currentfill}{rgb}{0.749020,0.772549,0.792157}%
\pgfsetfillcolor{currentfill}%
\pgfsetfillopacity{0.160000}%
\pgfsetlinewidth{0.000000pt}%
\definecolor{currentstroke}{rgb}{0.000000,0.000000,0.000000}%
\pgfsetstrokecolor{currentstroke}%
\pgfsetdash{}{0pt}%
\pgfpathmoveto{\pgfqpoint{1.608577in}{1.007328in}}%
\pgfpathlineto{\pgfqpoint{1.620739in}{1.016637in}}%
\pgfpathlineto{\pgfqpoint{1.632941in}{1.025978in}}%
\pgfpathlineto{\pgfqpoint{1.656425in}{1.040672in}}%
\pgfpathlineto{\pgfqpoint{1.643752in}{1.031044in}}%
\pgfpathlineto{\pgfqpoint{1.631120in}{1.021447in}}%
\pgfusepath{fill}%
\end{pgfscope}%
\begin{pgfscope}%
\pgfpathrectangle{\pgfqpoint{0.407592in}{0.222904in}}{\pgfqpoint{1.411163in}{1.411163in}}%
\pgfusepath{clip}%
\pgfsetbuttcap%
\pgfsetroundjoin%
\definecolor{currentfill}{rgb}{0.749020,0.772549,0.792157}%
\pgfsetfillcolor{currentfill}%
\pgfsetfillopacity{0.160000}%
\pgfsetlinewidth{0.000000pt}%
\definecolor{currentstroke}{rgb}{0.000000,0.000000,0.000000}%
\pgfsetstrokecolor{currentstroke}%
\pgfsetdash{}{0pt}%
\pgfpathmoveto{\pgfqpoint{0.636452in}{0.937885in}}%
\pgfpathlineto{\pgfqpoint{0.624276in}{0.946668in}}%
\pgfpathlineto{\pgfqpoint{0.612057in}{0.955483in}}%
\pgfpathlineto{\pgfqpoint{0.634912in}{0.943672in}}%
\pgfpathlineto{\pgfqpoint{0.660022in}{0.932489in}}%
\pgfpathlineto{\pgfqpoint{0.671092in}{0.924205in}}%
\pgfpathlineto{\pgfqpoint{0.682118in}{0.915953in}}%
\pgfpathlineto{\pgfqpoint{0.658221in}{0.926618in}}%
\pgfpathlineto{\pgfqpoint{0.636452in}{0.937885in}}%
\pgfpathclose%
\pgfusepath{fill}%
\end{pgfscope}%
\begin{pgfscope}%
\pgfpathrectangle{\pgfqpoint{0.407592in}{0.222904in}}{\pgfqpoint{1.411163in}{1.411163in}}%
\pgfusepath{clip}%
\pgfsetbuttcap%
\pgfsetroundjoin%
\definecolor{currentfill}{rgb}{0.749020,0.772549,0.792157}%
\pgfsetfillcolor{currentfill}%
\pgfsetfillopacity{0.160000}%
\pgfsetlinewidth{0.000000pt}%
\definecolor{currentstroke}{rgb}{0.000000,0.000000,0.000000}%
\pgfsetstrokecolor{currentstroke}%
\pgfsetdash{}{0pt}%
\pgfpathmoveto{\pgfqpoint{1.344960in}{0.859704in}}%
\pgfpathlineto{\pgfqpoint{1.352603in}{0.867180in}}%
\pgfpathlineto{\pgfqpoint{1.360279in}{0.874689in}}%
\pgfpathlineto{\pgfqpoint{1.389930in}{0.882926in}}%
\pgfpathlineto{\pgfqpoint{1.418045in}{0.891891in}}%
\pgfpathlineto{\pgfqpoint{1.408774in}{0.883934in}}%
\pgfpathlineto{\pgfqpoint{1.399540in}{0.876009in}}%
\pgfpathlineto{\pgfqpoint{1.372969in}{0.867510in}}%
\pgfpathlineto{\pgfqpoint{1.344960in}{0.859704in}}%
\pgfpathclose%
\pgfusepath{fill}%
\end{pgfscope}%
\begin{pgfscope}%
\pgfpathrectangle{\pgfqpoint{0.407592in}{0.222904in}}{\pgfqpoint{1.411163in}{1.411163in}}%
\pgfusepath{clip}%
\pgfsetbuttcap%
\pgfsetroundjoin%
\definecolor{currentfill}{rgb}{0.749020,0.772549,0.792157}%
\pgfsetfillcolor{currentfill}%
\pgfsetfillopacity{0.160000}%
\pgfsetlinewidth{0.000000pt}%
\definecolor{currentstroke}{rgb}{0.000000,0.000000,0.000000}%
\pgfsetstrokecolor{currentstroke}%
\pgfsetdash{}{0pt}%
\pgfpathmoveto{\pgfqpoint{1.273502in}{0.832230in}}%
\pgfpathlineto{\pgfqpoint{1.279322in}{0.839270in}}%
\pgfpathlineto{\pgfqpoint{1.285168in}{0.846342in}}%
\pgfpathlineto{\pgfqpoint{1.315646in}{0.852634in}}%
\pgfpathlineto{\pgfqpoint{1.344960in}{0.859704in}}%
\pgfpathlineto{\pgfqpoint{1.337350in}{0.852260in}}%
\pgfpathlineto{\pgfqpoint{1.329773in}{0.844848in}}%
\pgfpathlineto{\pgfqpoint{1.302180in}{0.838171in}}%
\pgfpathlineto{\pgfqpoint{1.273502in}{0.832230in}}%
\pgfpathclose%
\pgfusepath{fill}%
\end{pgfscope}%
\begin{pgfscope}%
\pgfpathrectangle{\pgfqpoint{0.407592in}{0.222904in}}{\pgfqpoint{1.411163in}{1.411163in}}%
\pgfusepath{clip}%
\pgfsetbuttcap%
\pgfsetroundjoin%
\definecolor{currentfill}{rgb}{0.749020,0.772549,0.792157}%
\pgfsetfillcolor{currentfill}%
\pgfsetfillopacity{0.160000}%
\pgfsetlinewidth{0.000000pt}%
\definecolor{currentstroke}{rgb}{0.000000,0.000000,0.000000}%
\pgfsetstrokecolor{currentstroke}%
\pgfsetdash{}{0pt}%
\pgfpathmoveto{\pgfqpoint{0.912344in}{0.812494in}}%
\pgfpathlineto{\pgfqpoint{0.906359in}{0.819289in}}%
\pgfpathlineto{\pgfqpoint{0.900346in}{0.826116in}}%
\pgfpathlineto{\pgfqpoint{0.929950in}{0.821907in}}%
\pgfpathlineto{\pgfqpoint{0.960359in}{0.818528in}}%
\pgfpathlineto{\pgfqpoint{0.964480in}{0.811927in}}%
\pgfpathlineto{\pgfqpoint{0.968582in}{0.805358in}}%
\pgfpathlineto{\pgfqpoint{0.940090in}{0.808535in}}%
\pgfpathlineto{\pgfqpoint{0.912344in}{0.812494in}}%
\pgfpathclose%
\pgfusepath{fill}%
\end{pgfscope}%
\begin{pgfscope}%
\pgfpathrectangle{\pgfqpoint{0.407592in}{0.222904in}}{\pgfqpoint{1.411163in}{1.411163in}}%
\pgfusepath{clip}%
\pgfsetbuttcap%
\pgfsetroundjoin%
\definecolor{currentfill}{rgb}{0.749020,0.772549,0.792157}%
\pgfsetfillcolor{currentfill}%
\pgfsetfillopacity{0.160000}%
\pgfsetlinewidth{0.000000pt}%
\definecolor{currentstroke}{rgb}{0.000000,0.000000,0.000000}%
\pgfsetstrokecolor{currentstroke}%
\pgfsetdash{}{0pt}%
\pgfpathmoveto{\pgfqpoint{1.146991in}{0.803392in}}%
\pgfpathlineto{\pgfqpoint{1.148906in}{0.809899in}}%
\pgfpathlineto{\pgfqpoint{1.150830in}{0.816438in}}%
\pgfpathlineto{\pgfqpoint{1.182387in}{0.819139in}}%
\pgfpathlineto{\pgfqpoint{1.213456in}{0.822686in}}%
\pgfpathlineto{\pgfqpoint{1.209550in}{0.815961in}}%
\pgfpathlineto{\pgfqpoint{1.205662in}{0.809268in}}%
\pgfpathlineto{\pgfqpoint{1.176553in}{0.805932in}}%
\pgfpathlineto{\pgfqpoint{1.146991in}{0.803392in}}%
\pgfpathclose%
\pgfusepath{fill}%
\end{pgfscope}%
\begin{pgfscope}%
\pgfpathrectangle{\pgfqpoint{0.407592in}{0.222904in}}{\pgfqpoint{1.411163in}{1.411163in}}%
\pgfusepath{clip}%
\pgfsetbuttcap%
\pgfsetroundjoin%
\definecolor{currentfill}{rgb}{0.749020,0.772549,0.792157}%
\pgfsetfillcolor{currentfill}%
\pgfsetfillopacity{0.160000}%
\pgfsetlinewidth{0.000000pt}%
\definecolor{currentstroke}{rgb}{0.000000,0.000000,0.000000}%
\pgfsetstrokecolor{currentstroke}%
\pgfsetdash{}{0pt}%
\pgfpathmoveto{\pgfqpoint{1.707519in}{1.079488in}}%
\pgfpathlineto{\pgfqpoint{1.720395in}{1.089269in}}%
\pgfpathlineto{\pgfqpoint{1.733311in}{1.099082in}}%
\pgfpathlineto{\pgfqpoint{1.756371in}{1.115950in}}%
\pgfpathlineto{\pgfqpoint{1.776149in}{1.133206in}}%
\pgfpathlineto{\pgfqpoint{1.762521in}{1.122822in}}%
\pgfpathlineto{\pgfqpoint{1.748930in}{1.112466in}}%
\pgfpathlineto{\pgfqpoint{1.729795in}{1.095786in}}%
\pgfpathlineto{\pgfqpoint{1.707519in}{1.079488in}}%
\pgfpathclose%
\pgfusepath{fill}%
\end{pgfscope}%
\begin{pgfscope}%
\pgfpathrectangle{\pgfqpoint{0.407592in}{0.222904in}}{\pgfqpoint{1.411163in}{1.411163in}}%
\pgfusepath{clip}%
\pgfsetbuttcap%
\pgfsetroundjoin%
\definecolor{currentfill}{rgb}{0.749020,0.772549,0.792157}%
\pgfsetfillcolor{currentfill}%
\pgfsetfillopacity{0.160000}%
\pgfsetlinewidth{0.000000pt}%
\definecolor{currentstroke}{rgb}{0.000000,0.000000,0.000000}%
\pgfsetstrokecolor{currentstroke}%
\pgfsetdash{}{0pt}%
\pgfpathmoveto{\pgfqpoint{0.704041in}{0.899547in}}%
\pgfpathlineto{\pgfqpoint{0.693101in}{0.907734in}}%
\pgfpathlineto{\pgfqpoint{0.682118in}{0.915953in}}%
\pgfpathlineto{\pgfqpoint{0.708043in}{0.905952in}}%
\pgfpathlineto{\pgfqpoint{0.735882in}{0.896670in}}%
\pgfpathlineto{\pgfqpoint{0.745472in}{0.888925in}}%
\pgfpathlineto{\pgfqpoint{0.755023in}{0.881214in}}%
\pgfpathlineto{\pgfqpoint{0.728632in}{0.890037in}}%
\pgfpathlineto{\pgfqpoint{0.704041in}{0.899547in}}%
\pgfpathclose%
\pgfusepath{fill}%
\end{pgfscope}%
\begin{pgfscope}%
\pgfpathrectangle{\pgfqpoint{0.407592in}{0.222904in}}{\pgfqpoint{1.411163in}{1.411163in}}%
\pgfusepath{clip}%
\pgfsetbuttcap%
\pgfsetroundjoin%
\definecolor{currentfill}{rgb}{0.749020,0.772549,0.792157}%
\pgfsetfillcolor{currentfill}%
\pgfsetfillopacity{0.160000}%
\pgfsetlinewidth{0.000000pt}%
\definecolor{currentstroke}{rgb}{0.000000,0.000000,0.000000}%
\pgfsetstrokecolor{currentstroke}%
\pgfsetdash{}{0pt}%
\pgfpathmoveto{\pgfqpoint{0.968582in}{0.805358in}}%
\pgfpathlineto{\pgfqpoint{0.964480in}{0.811927in}}%
\pgfpathlineto{\pgfqpoint{0.960359in}{0.818528in}}%
\pgfpathlineto{\pgfqpoint{0.991417in}{0.815999in}}%
\pgfpathlineto{\pgfqpoint{1.022968in}{0.814334in}}%
\pgfpathlineto{\pgfqpoint{1.025106in}{0.807858in}}%
\pgfpathlineto{\pgfqpoint{1.027232in}{0.801414in}}%
\pgfpathlineto{\pgfqpoint{0.997679in}{0.802979in}}%
\pgfpathlineto{\pgfqpoint{0.968582in}{0.805358in}}%
\pgfpathclose%
\pgfusepath{fill}%
\end{pgfscope}%
\begin{pgfscope}%
\pgfpathrectangle{\pgfqpoint{0.407592in}{0.222904in}}{\pgfqpoint{1.411163in}{1.411163in}}%
\pgfusepath{clip}%
\pgfsetbuttcap%
\pgfsetroundjoin%
\definecolor{currentfill}{rgb}{0.749020,0.772549,0.792157}%
\pgfsetfillcolor{currentfill}%
\pgfsetfillopacity{0.160000}%
\pgfsetlinewidth{0.000000pt}%
\definecolor{currentstroke}{rgb}{0.000000,0.000000,0.000000}%
\pgfsetstrokecolor{currentstroke}%
\pgfsetdash{}{0pt}%
\pgfpathmoveto{\pgfqpoint{1.436700in}{0.907902in}}%
\pgfpathlineto{\pgfqpoint{1.446084in}{0.915957in}}%
\pgfpathlineto{\pgfqpoint{1.455507in}{0.924044in}}%
\pgfpathlineto{\pgfqpoint{1.484848in}{0.934679in}}%
\pgfpathlineto{\pgfqpoint{1.512175in}{0.945993in}}%
\pgfpathlineto{\pgfqpoint{1.501362in}{0.937397in}}%
\pgfpathlineto{\pgfqpoint{1.490589in}{0.928832in}}%
\pgfpathlineto{\pgfqpoint{1.464593in}{0.918042in}}%
\pgfpathlineto{\pgfqpoint{1.436700in}{0.907902in}}%
\pgfpathclose%
\pgfusepath{fill}%
\end{pgfscope}%
\begin{pgfscope}%
\pgfpathrectangle{\pgfqpoint{0.407592in}{0.222904in}}{\pgfqpoint{1.411163in}{1.411163in}}%
\pgfusepath{clip}%
\pgfsetbuttcap%
\pgfsetroundjoin%
\definecolor{currentfill}{rgb}{0.749020,0.772549,0.792157}%
\pgfsetfillcolor{currentfill}%
\pgfsetfillopacity{0.160000}%
\pgfsetlinewidth{0.000000pt}%
\definecolor{currentstroke}{rgb}{0.000000,0.000000,0.000000}%
\pgfsetstrokecolor{currentstroke}%
\pgfsetdash{}{0pt}%
\pgfpathmoveto{\pgfqpoint{0.405447in}{1.109093in}}%
\pgfpathlineto{\pgfqpoint{0.391303in}{1.119391in}}%
\pgfpathlineto{\pgfqpoint{0.377119in}{1.129718in}}%
\pgfpathlineto{\pgfqpoint{0.393314in}{1.112532in}}%
\pgfpathlineto{\pgfqpoint{0.412881in}{1.095753in}}%
\pgfpathlineto{\pgfqpoint{0.426500in}{1.085997in}}%
\pgfpathlineto{\pgfqpoint{0.440076in}{1.076272in}}%
\pgfpathlineto{\pgfqpoint{0.421148in}{1.092483in}}%
\pgfpathlineto{\pgfqpoint{0.405447in}{1.109093in}}%
\pgfpathclose%
\pgfusepath{fill}%
\end{pgfscope}%
\begin{pgfscope}%
\pgfpathrectangle{\pgfqpoint{0.407592in}{0.222904in}}{\pgfqpoint{1.411163in}{1.411163in}}%
\pgfusepath{clip}%
\pgfsetbuttcap%
\pgfsetroundjoin%
\definecolor{currentfill}{rgb}{0.749020,0.772549,0.792157}%
\pgfsetfillcolor{currentfill}%
\pgfsetfillopacity{0.160000}%
\pgfsetlinewidth{0.000000pt}%
\definecolor{currentstroke}{rgb}{0.000000,0.000000,0.000000}%
\pgfsetstrokecolor{currentstroke}%
\pgfsetdash{}{0pt}%
\pgfpathmoveto{\pgfqpoint{0.844134in}{0.836919in}}%
\pgfpathlineto{\pgfqpoint{0.836331in}{0.844098in}}%
\pgfpathlineto{\pgfqpoint{0.828493in}{0.851308in}}%
\pgfpathlineto{\pgfqpoint{0.857778in}{0.845177in}}%
\pgfpathlineto{\pgfqpoint{0.888235in}{0.839866in}}%
\pgfpathlineto{\pgfqpoint{0.894305in}{0.832975in}}%
\pgfpathlineto{\pgfqpoint{0.900346in}{0.826116in}}%
\pgfpathlineto{\pgfqpoint{0.871693in}{0.831130in}}%
\pgfpathlineto{\pgfqpoint{0.844134in}{0.836919in}}%
\pgfpathclose%
\pgfusepath{fill}%
\end{pgfscope}%
\begin{pgfscope}%
\pgfpathrectangle{\pgfqpoint{0.407592in}{0.222904in}}{\pgfqpoint{1.411163in}{1.411163in}}%
\pgfusepath{clip}%
\pgfsetbuttcap%
\pgfsetroundjoin%
\definecolor{currentfill}{rgb}{0.749020,0.772549,0.792157}%
\pgfsetfillcolor{currentfill}%
\pgfsetfillopacity{0.160000}%
\pgfsetlinewidth{0.000000pt}%
\definecolor{currentstroke}{rgb}{0.000000,0.000000,0.000000}%
\pgfsetstrokecolor{currentstroke}%
\pgfsetdash{}{0pt}%
\pgfpathmoveto{\pgfqpoint{1.533923in}{0.963284in}}%
\pgfpathlineto{\pgfqpoint{1.544858in}{0.971978in}}%
\pgfpathlineto{\pgfqpoint{1.555834in}{0.980704in}}%
\pgfpathlineto{\pgfqpoint{1.583451in}{0.993721in}}%
\pgfpathlineto{\pgfqpoint{1.608577in}{1.007328in}}%
\pgfpathlineto{\pgfqpoint{1.596457in}{0.998049in}}%
\pgfpathlineto{\pgfqpoint{1.584377in}{0.988802in}}%
\pgfpathlineto{\pgfqpoint{1.560330in}{0.975758in}}%
\pgfpathlineto{\pgfqpoint{1.533923in}{0.963284in}}%
\pgfpathclose%
\pgfusepath{fill}%
\end{pgfscope}%
\begin{pgfscope}%
\pgfpathrectangle{\pgfqpoint{0.407592in}{0.222904in}}{\pgfqpoint{1.411163in}{1.411163in}}%
\pgfusepath{clip}%
\pgfsetbuttcap%
\pgfsetroundjoin%
\definecolor{currentfill}{rgb}{0.749020,0.772549,0.792157}%
\pgfsetfillcolor{currentfill}%
\pgfsetfillopacity{0.160000}%
\pgfsetlinewidth{0.000000pt}%
\definecolor{currentstroke}{rgb}{0.000000,0.000000,0.000000}%
\pgfsetstrokecolor{currentstroke}%
\pgfsetdash{}{0pt}%
\pgfpathmoveto{\pgfqpoint{0.520615in}{1.018577in}}%
\pgfpathlineto{\pgfqpoint{0.507300in}{1.028116in}}%
\pgfpathlineto{\pgfqpoint{0.493941in}{1.037685in}}%
\pgfpathlineto{\pgfqpoint{0.514461in}{1.023091in}}%
\pgfpathlineto{\pgfqpoint{0.537814in}{1.009042in}}%
\pgfpathlineto{\pgfqpoint{0.550299in}{1.000035in}}%
\pgfpathlineto{\pgfqpoint{0.562740in}{0.991060in}}%
\pgfpathlineto{\pgfqpoint{0.540331in}{1.004554in}}%
\pgfpathlineto{\pgfqpoint{0.520615in}{1.018577in}}%
\pgfpathclose%
\pgfusepath{fill}%
\end{pgfscope}%
\begin{pgfscope}%
\pgfpathrectangle{\pgfqpoint{0.407592in}{0.222904in}}{\pgfqpoint{1.411163in}{1.411163in}}%
\pgfusepath{clip}%
\pgfsetbuttcap%
\pgfsetroundjoin%
\definecolor{currentfill}{rgb}{0.749020,0.772549,0.792157}%
\pgfsetfillcolor{currentfill}%
\pgfsetfillopacity{0.160000}%
\pgfsetlinewidth{0.000000pt}%
\definecolor{currentstroke}{rgb}{0.000000,0.000000,0.000000}%
\pgfsetstrokecolor{currentstroke}%
\pgfsetdash{}{0pt}%
\pgfpathmoveto{\pgfqpoint{1.087109in}{0.800754in}}%
\pgfpathlineto{\pgfqpoint{1.087003in}{0.807177in}}%
\pgfpathlineto{\pgfqpoint{1.086897in}{0.813631in}}%
\pgfpathlineto{\pgfqpoint{1.118946in}{0.814598in}}%
\pgfpathlineto{\pgfqpoint{1.150830in}{0.816438in}}%
\pgfpathlineto{\pgfqpoint{1.148906in}{0.809899in}}%
\pgfpathlineto{\pgfqpoint{1.146991in}{0.803392in}}%
\pgfpathlineto{\pgfqpoint{1.117125in}{0.801662in}}%
\pgfpathlineto{\pgfqpoint{1.087109in}{0.800754in}}%
\pgfpathclose%
\pgfusepath{fill}%
\end{pgfscope}%
\begin{pgfscope}%
\pgfpathrectangle{\pgfqpoint{0.407592in}{0.222904in}}{\pgfqpoint{1.411163in}{1.411163in}}%
\pgfusepath{clip}%
\pgfsetbuttcap%
\pgfsetroundjoin%
\definecolor{currentfill}{rgb}{0.749020,0.772549,0.792157}%
\pgfsetfillcolor{currentfill}%
\pgfsetfillopacity{0.160000}%
\pgfsetlinewidth{0.000000pt}%
\definecolor{currentstroke}{rgb}{0.000000,0.000000,0.000000}%
\pgfsetstrokecolor{currentstroke}%
\pgfsetdash{}{0pt}%
\pgfpathmoveto{\pgfqpoint{1.632941in}{1.025978in}}%
\pgfpathlineto{\pgfqpoint{1.645185in}{1.035351in}}%
\pgfpathlineto{\pgfqpoint{1.657470in}{1.044755in}}%
\pgfpathlineto{\pgfqpoint{1.681891in}{1.060018in}}%
\pgfpathlineto{\pgfqpoint{1.669138in}{1.050329in}}%
\pgfpathlineto{\pgfqpoint{1.656425in}{1.040672in}}%
\pgfusepath{fill}%
\end{pgfscope}%
\begin{pgfscope}%
\pgfpathrectangle{\pgfqpoint{0.407592in}{0.222904in}}{\pgfqpoint{1.411163in}{1.411163in}}%
\pgfusepath{clip}%
\pgfsetbuttcap%
\pgfsetroundjoin%
\definecolor{currentfill}{rgb}{0.749020,0.772549,0.792157}%
\pgfsetfillcolor{currentfill}%
\pgfsetfillopacity{0.160000}%
\pgfsetlinewidth{0.000000pt}%
\definecolor{currentstroke}{rgb}{0.000000,0.000000,0.000000}%
\pgfsetstrokecolor{currentstroke}%
\pgfsetdash{}{0pt}%
\pgfpathmoveto{\pgfqpoint{1.027232in}{0.801414in}}%
\pgfpathlineto{\pgfqpoint{1.025106in}{0.807858in}}%
\pgfpathlineto{\pgfqpoint{1.022968in}{0.814334in}}%
\pgfpathlineto{\pgfqpoint{1.054850in}{0.813543in}}%
\pgfpathlineto{\pgfqpoint{1.086897in}{0.813631in}}%
\pgfpathlineto{\pgfqpoint{1.087003in}{0.807177in}}%
\pgfpathlineto{\pgfqpoint{1.087109in}{0.800754in}}%
\pgfpathlineto{\pgfqpoint{1.057093in}{0.800671in}}%
\pgfpathlineto{\pgfqpoint{1.027232in}{0.801414in}}%
\pgfpathclose%
\pgfusepath{fill}%
\end{pgfscope}%
\begin{pgfscope}%
\pgfpathrectangle{\pgfqpoint{0.407592in}{0.222904in}}{\pgfqpoint{1.411163in}{1.411163in}}%
\pgfusepath{clip}%
\pgfsetbuttcap%
\pgfsetroundjoin%
\definecolor{currentfill}{rgb}{0.749020,0.772549,0.792157}%
\pgfsetfillcolor{currentfill}%
\pgfsetfillopacity{0.160000}%
\pgfsetlinewidth{0.000000pt}%
\definecolor{currentstroke}{rgb}{0.000000,0.000000,0.000000}%
\pgfsetstrokecolor{currentstroke}%
\pgfsetdash{}{0pt}%
\pgfpathmoveto{\pgfqpoint{0.774003in}{0.865887in}}%
\pgfpathlineto{\pgfqpoint{0.764533in}{0.873534in}}%
\pgfpathlineto{\pgfqpoint{0.755023in}{0.881214in}}%
\pgfpathlineto{\pgfqpoint{0.783092in}{0.873128in}}%
\pgfpathlineto{\pgfqpoint{0.812710in}{0.865827in}}%
\pgfpathlineto{\pgfqpoint{0.820619in}{0.858551in}}%
\pgfpathlineto{\pgfqpoint{0.828493in}{0.851308in}}%
\pgfpathlineto{\pgfqpoint{0.800523in}{0.858225in}}%
\pgfpathlineto{\pgfqpoint{0.774003in}{0.865887in}}%
\pgfpathclose%
\pgfusepath{fill}%
\end{pgfscope}%
\begin{pgfscope}%
\pgfpathrectangle{\pgfqpoint{0.407592in}{0.222904in}}{\pgfqpoint{1.411163in}{1.411163in}}%
\pgfusepath{clip}%
\pgfsetbuttcap%
\pgfsetroundjoin%
\definecolor{currentfill}{rgb}{0.749020,0.772549,0.792157}%
\pgfsetfillcolor{currentfill}%
\pgfsetfillopacity{0.160000}%
\pgfsetlinewidth{0.000000pt}%
\definecolor{currentstroke}{rgb}{0.000000,0.000000,0.000000}%
\pgfsetstrokecolor{currentstroke}%
\pgfsetdash{}{0pt}%
\pgfpathmoveto{\pgfqpoint{0.612057in}{0.955483in}}%
\pgfpathlineto{\pgfqpoint{0.599794in}{0.964330in}}%
\pgfpathlineto{\pgfqpoint{0.587487in}{0.973208in}}%
\pgfpathlineto{\pgfqpoint{0.611427in}{0.960856in}}%
\pgfpathlineto{\pgfqpoint{0.637751in}{0.949156in}}%
\pgfpathlineto{\pgfqpoint{0.648909in}{0.940806in}}%
\pgfpathlineto{\pgfqpoint{0.660022in}{0.932489in}}%
\pgfpathlineto{\pgfqpoint{0.634912in}{0.943672in}}%
\pgfpathlineto{\pgfqpoint{0.612057in}{0.955483in}}%
\pgfpathclose%
\pgfusepath{fill}%
\end{pgfscope}%
\begin{pgfscope}%
\pgfpathrectangle{\pgfqpoint{0.407592in}{0.222904in}}{\pgfqpoint{1.411163in}{1.411163in}}%
\pgfusepath{clip}%
\pgfsetbuttcap%
\pgfsetroundjoin%
\definecolor{currentfill}{rgb}{0.749020,0.772549,0.792157}%
\pgfsetfillcolor{currentfill}%
\pgfsetfillopacity{0.160000}%
\pgfsetlinewidth{0.000000pt}%
\definecolor{currentstroke}{rgb}{0.000000,0.000000,0.000000}%
\pgfsetstrokecolor{currentstroke}%
\pgfsetdash{}{0pt}%
\pgfpathmoveto{\pgfqpoint{1.213456in}{0.822686in}}%
\pgfpathlineto{\pgfqpoint{1.217380in}{0.829443in}}%
\pgfpathlineto{\pgfqpoint{1.221324in}{0.836231in}}%
\pgfpathlineto{\pgfqpoint{1.253675in}{0.840864in}}%
\pgfpathlineto{\pgfqpoint{1.285168in}{0.846342in}}%
\pgfpathlineto{\pgfqpoint{1.279322in}{0.839270in}}%
\pgfpathlineto{\pgfqpoint{1.273502in}{0.832230in}}%
\pgfpathlineto{\pgfqpoint{1.243878in}{0.827058in}}%
\pgfpathlineto{\pgfqpoint{1.213456in}{0.822686in}}%
\pgfpathclose%
\pgfusepath{fill}%
\end{pgfscope}%
\begin{pgfscope}%
\pgfpathrectangle{\pgfqpoint{0.407592in}{0.222904in}}{\pgfqpoint{1.411163in}{1.411163in}}%
\pgfusepath{clip}%
\pgfsetbuttcap%
\pgfsetroundjoin%
\definecolor{currentfill}{rgb}{0.749020,0.772549,0.792157}%
\pgfsetfillcolor{currentfill}%
\pgfsetfillopacity{0.160000}%
\pgfsetlinewidth{0.000000pt}%
\definecolor{currentstroke}{rgb}{0.000000,0.000000,0.000000}%
\pgfsetstrokecolor{currentstroke}%
\pgfsetdash{}{0pt}%
\pgfpathmoveto{\pgfqpoint{1.360279in}{0.874689in}}%
\pgfpathlineto{\pgfqpoint{1.367989in}{0.882230in}}%
\pgfpathlineto{\pgfqpoint{1.375732in}{0.889804in}}%
\pgfpathlineto{\pgfqpoint{1.407034in}{0.898472in}}%
\pgfpathlineto{\pgfqpoint{1.436700in}{0.907902in}}%
\pgfpathlineto{\pgfqpoint{1.427353in}{0.899881in}}%
\pgfpathlineto{\pgfqpoint{1.418045in}{0.891891in}}%
\pgfpathlineto{\pgfqpoint{1.389930in}{0.882926in}}%
\pgfpathlineto{\pgfqpoint{1.360279in}{0.874689in}}%
\pgfpathclose%
\pgfusepath{fill}%
\end{pgfscope}%
\begin{pgfscope}%
\pgfpathrectangle{\pgfqpoint{0.407592in}{0.222904in}}{\pgfqpoint{1.411163in}{1.411163in}}%
\pgfusepath{clip}%
\pgfsetbuttcap%
\pgfsetroundjoin%
\definecolor{currentfill}{rgb}{0.749020,0.772549,0.792157}%
\pgfsetfillcolor{currentfill}%
\pgfsetfillopacity{0.160000}%
\pgfsetlinewidth{0.000000pt}%
\definecolor{currentstroke}{rgb}{0.000000,0.000000,0.000000}%
\pgfsetstrokecolor{currentstroke}%
\pgfsetdash{}{0pt}%
\pgfpathmoveto{\pgfqpoint{1.733311in}{1.099082in}}%
\pgfpathlineto{\pgfqpoint{1.746270in}{1.108926in}}%
\pgfpathlineto{\pgfqpoint{1.759269in}{1.118802in}}%
\pgfpathlineto{\pgfqpoint{1.783107in}{1.136235in}}%
\pgfpathlineto{\pgfqpoint{1.803517in}{1.154061in}}%
\pgfpathlineto{\pgfqpoint{1.789814in}{1.143619in}}%
\pgfpathlineto{\pgfqpoint{1.776149in}{1.133206in}}%
\pgfpathlineto{\pgfqpoint{1.756371in}{1.115950in}}%
\pgfpathlineto{\pgfqpoint{1.733311in}{1.099082in}}%
\pgfpathclose%
\pgfusepath{fill}%
\end{pgfscope}%
\begin{pgfscope}%
\pgfpathrectangle{\pgfqpoint{0.407592in}{0.222904in}}{\pgfqpoint{1.411163in}{1.411163in}}%
\pgfusepath{clip}%
\pgfsetbuttcap%
\pgfsetroundjoin%
\definecolor{currentfill}{rgb}{0.749020,0.772549,0.792157}%
\pgfsetfillcolor{currentfill}%
\pgfsetfillopacity{0.160000}%
\pgfsetlinewidth{0.000000pt}%
\definecolor{currentstroke}{rgb}{0.000000,0.000000,0.000000}%
\pgfsetstrokecolor{currentstroke}%
\pgfsetdash{}{0pt}%
\pgfpathmoveto{\pgfqpoint{1.285168in}{0.846342in}}%
\pgfpathlineto{\pgfqpoint{1.291041in}{0.853447in}}%
\pgfpathlineto{\pgfqpoint{1.296941in}{0.860584in}}%
\pgfpathlineto{\pgfqpoint{1.329233in}{0.867227in}}%
\pgfpathlineto{\pgfqpoint{1.360279in}{0.874689in}}%
\pgfpathlineto{\pgfqpoint{1.352603in}{0.867180in}}%
\pgfpathlineto{\pgfqpoint{1.344960in}{0.859704in}}%
\pgfpathlineto{\pgfqpoint{1.315646in}{0.852634in}}%
\pgfpathlineto{\pgfqpoint{1.285168in}{0.846342in}}%
\pgfpathclose%
\pgfusepath{fill}%
\end{pgfscope}%
\begin{pgfscope}%
\pgfpathrectangle{\pgfqpoint{0.407592in}{0.222904in}}{\pgfqpoint{1.411163in}{1.411163in}}%
\pgfusepath{clip}%
\pgfsetbuttcap%
\pgfsetroundjoin%
\definecolor{currentfill}{rgb}{0.749020,0.772549,0.792157}%
\pgfsetfillcolor{currentfill}%
\pgfsetfillopacity{0.160000}%
\pgfsetlinewidth{0.000000pt}%
\definecolor{currentstroke}{rgb}{0.000000,0.000000,0.000000}%
\pgfsetstrokecolor{currentstroke}%
\pgfsetdash{}{0pt}%
\pgfpathmoveto{\pgfqpoint{0.900346in}{0.826116in}}%
\pgfpathlineto{\pgfqpoint{0.894305in}{0.832975in}}%
\pgfpathlineto{\pgfqpoint{0.888235in}{0.839866in}}%
\pgfpathlineto{\pgfqpoint{0.919714in}{0.835406in}}%
\pgfpathlineto{\pgfqpoint{0.952056in}{0.831826in}}%
\pgfpathlineto{\pgfqpoint{0.956217in}{0.825161in}}%
\pgfpathlineto{\pgfqpoint{0.960359in}{0.818528in}}%
\pgfpathlineto{\pgfqpoint{0.929950in}{0.821907in}}%
\pgfpathlineto{\pgfqpoint{0.900346in}{0.826116in}}%
\pgfpathclose%
\pgfusepath{fill}%
\end{pgfscope}%
\begin{pgfscope}%
\pgfpathrectangle{\pgfqpoint{0.407592in}{0.222904in}}{\pgfqpoint{1.411163in}{1.411163in}}%
\pgfusepath{clip}%
\pgfsetbuttcap%
\pgfsetroundjoin%
\definecolor{currentfill}{rgb}{0.749020,0.772549,0.792157}%
\pgfsetfillcolor{currentfill}%
\pgfsetfillopacity{0.160000}%
\pgfsetlinewidth{0.000000pt}%
\definecolor{currentstroke}{rgb}{0.000000,0.000000,0.000000}%
\pgfsetstrokecolor{currentstroke}%
\pgfsetdash{}{0pt}%
\pgfpathmoveto{\pgfqpoint{0.377119in}{1.129718in}}%
\pgfpathlineto{\pgfqpoint{0.362895in}{1.140073in}}%
\pgfpathlineto{\pgfqpoint{0.348632in}{1.150458in}}%
\pgfpathlineto{\pgfqpoint{0.365311in}{1.132703in}}%
\pgfpathlineto{\pgfqpoint{0.385509in}{1.115361in}}%
\pgfpathlineto{\pgfqpoint{0.399217in}{1.105541in}}%
\pgfpathlineto{\pgfqpoint{0.412881in}{1.095753in}}%
\pgfpathlineto{\pgfqpoint{0.393314in}{1.112532in}}%
\pgfpathlineto{\pgfqpoint{0.377119in}{1.129718in}}%
\pgfpathclose%
\pgfusepath{fill}%
\end{pgfscope}%
\begin{pgfscope}%
\pgfpathrectangle{\pgfqpoint{0.407592in}{0.222904in}}{\pgfqpoint{1.411163in}{1.411163in}}%
\pgfusepath{clip}%
\pgfsetbuttcap%
\pgfsetroundjoin%
\definecolor{currentfill}{rgb}{0.749020,0.772549,0.792157}%
\pgfsetfillcolor{currentfill}%
\pgfsetfillopacity{0.160000}%
\pgfsetlinewidth{0.000000pt}%
\definecolor{currentstroke}{rgb}{0.000000,0.000000,0.000000}%
\pgfsetstrokecolor{currentstroke}%
\pgfsetdash{}{0pt}%
\pgfpathmoveto{\pgfqpoint{0.493941in}{1.037685in}}%
\pgfpathlineto{\pgfqpoint{0.480540in}{1.047285in}}%
\pgfpathlineto{\pgfqpoint{0.467095in}{1.056916in}}%
\pgfpathlineto{\pgfqpoint{0.488413in}{1.041755in}}%
\pgfpathlineto{\pgfqpoint{0.512706in}{1.027154in}}%
\pgfpathlineto{\pgfqpoint{0.525283in}{1.018082in}}%
\pgfpathlineto{\pgfqpoint{0.537814in}{1.009042in}}%
\pgfpathlineto{\pgfqpoint{0.514461in}{1.023091in}}%
\pgfpathlineto{\pgfqpoint{0.493941in}{1.037685in}}%
\pgfpathclose%
\pgfusepath{fill}%
\end{pgfscope}%
\begin{pgfscope}%
\pgfpathrectangle{\pgfqpoint{0.407592in}{0.222904in}}{\pgfqpoint{1.411163in}{1.411163in}}%
\pgfusepath{clip}%
\pgfsetbuttcap%
\pgfsetroundjoin%
\definecolor{currentfill}{rgb}{0.749020,0.772549,0.792157}%
\pgfsetfillcolor{currentfill}%
\pgfsetfillopacity{0.160000}%
\pgfsetlinewidth{0.000000pt}%
\definecolor{currentstroke}{rgb}{0.000000,0.000000,0.000000}%
\pgfsetstrokecolor{currentstroke}%
\pgfsetdash{}{0pt}%
\pgfpathmoveto{\pgfqpoint{1.657470in}{1.044755in}}%
\pgfpathlineto{\pgfqpoint{1.669797in}{1.054191in}}%
\pgfpathlineto{\pgfqpoint{1.682166in}{1.063659in}}%
\pgfpathlineto{\pgfqpoint{1.707519in}{1.079488in}}%
\pgfpathlineto{\pgfqpoint{1.694685in}{1.069737in}}%
\pgfpathlineto{\pgfqpoint{1.681891in}{1.060018in}}%
\pgfusepath{fill}%
\end{pgfscope}%
\begin{pgfscope}%
\pgfpathrectangle{\pgfqpoint{0.407592in}{0.222904in}}{\pgfqpoint{1.411163in}{1.411163in}}%
\pgfusepath{clip}%
\pgfsetbuttcap%
\pgfsetroundjoin%
\definecolor{currentfill}{rgb}{0.749020,0.772549,0.792157}%
\pgfsetfillcolor{currentfill}%
\pgfsetfillopacity{0.160000}%
\pgfsetlinewidth{0.000000pt}%
\definecolor{currentstroke}{rgb}{0.000000,0.000000,0.000000}%
\pgfsetstrokecolor{currentstroke}%
\pgfsetdash{}{0pt}%
\pgfpathmoveto{\pgfqpoint{1.555834in}{0.980704in}}%
\pgfpathlineto{\pgfqpoint{1.566852in}{0.989463in}}%
\pgfpathlineto{\pgfqpoint{1.577911in}{0.998256in}}%
\pgfpathlineto{\pgfqpoint{1.606738in}{1.011813in}}%
\pgfpathlineto{\pgfqpoint{1.632941in}{1.025978in}}%
\pgfpathlineto{\pgfqpoint{1.620739in}{1.016637in}}%
\pgfpathlineto{\pgfqpoint{1.608577in}{1.007328in}}%
\pgfpathlineto{\pgfqpoint{1.583451in}{0.993721in}}%
\pgfpathlineto{\pgfqpoint{1.555834in}{0.980704in}}%
\pgfpathclose%
\pgfusepath{fill}%
\end{pgfscope}%
\begin{pgfscope}%
\pgfpathrectangle{\pgfqpoint{0.407592in}{0.222904in}}{\pgfqpoint{1.411163in}{1.411163in}}%
\pgfusepath{clip}%
\pgfsetbuttcap%
\pgfsetroundjoin%
\definecolor{currentfill}{rgb}{0.749020,0.772549,0.792157}%
\pgfsetfillcolor{currentfill}%
\pgfsetfillopacity{0.160000}%
\pgfsetlinewidth{0.000000pt}%
\definecolor{currentstroke}{rgb}{0.000000,0.000000,0.000000}%
\pgfsetstrokecolor{currentstroke}%
\pgfsetdash{}{0pt}%
\pgfpathmoveto{\pgfqpoint{1.150830in}{0.816438in}}%
\pgfpathlineto{\pgfqpoint{1.152764in}{0.823008in}}%
\pgfpathlineto{\pgfqpoint{1.154707in}{0.829610in}}%
\pgfpathlineto{\pgfqpoint{1.188278in}{0.832473in}}%
\pgfpathlineto{\pgfqpoint{1.221324in}{0.836231in}}%
\pgfpathlineto{\pgfqpoint{1.217380in}{0.829443in}}%
\pgfpathlineto{\pgfqpoint{1.213456in}{0.822686in}}%
\pgfpathlineto{\pgfqpoint{1.182387in}{0.819139in}}%
\pgfpathlineto{\pgfqpoint{1.150830in}{0.816438in}}%
\pgfpathclose%
\pgfusepath{fill}%
\end{pgfscope}%
\begin{pgfscope}%
\pgfpathrectangle{\pgfqpoint{0.407592in}{0.222904in}}{\pgfqpoint{1.411163in}{1.411163in}}%
\pgfusepath{clip}%
\pgfsetbuttcap%
\pgfsetroundjoin%
\definecolor{currentfill}{rgb}{0.749020,0.772549,0.792157}%
\pgfsetfillcolor{currentfill}%
\pgfsetfillopacity{0.160000}%
\pgfsetlinewidth{0.000000pt}%
\definecolor{currentstroke}{rgb}{0.000000,0.000000,0.000000}%
\pgfsetstrokecolor{currentstroke}%
\pgfsetdash{}{0pt}%
\pgfpathmoveto{\pgfqpoint{0.682118in}{0.915953in}}%
\pgfpathlineto{\pgfqpoint{0.671092in}{0.924205in}}%
\pgfpathlineto{\pgfqpoint{0.660022in}{0.932489in}}%
\pgfpathlineto{\pgfqpoint{0.687284in}{0.921998in}}%
\pgfpathlineto{\pgfqpoint{0.716577in}{0.912258in}}%
\pgfpathlineto{\pgfqpoint{0.726250in}{0.904448in}}%
\pgfpathlineto{\pgfqpoint{0.735882in}{0.896670in}}%
\pgfpathlineto{\pgfqpoint{0.708043in}{0.905952in}}%
\pgfpathlineto{\pgfqpoint{0.682118in}{0.915953in}}%
\pgfpathclose%
\pgfusepath{fill}%
\end{pgfscope}%
\begin{pgfscope}%
\pgfpathrectangle{\pgfqpoint{0.407592in}{0.222904in}}{\pgfqpoint{1.411163in}{1.411163in}}%
\pgfusepath{clip}%
\pgfsetbuttcap%
\pgfsetroundjoin%
\definecolor{currentfill}{rgb}{0.749020,0.772549,0.792157}%
\pgfsetfillcolor{currentfill}%
\pgfsetfillopacity{0.160000}%
\pgfsetlinewidth{0.000000pt}%
\definecolor{currentstroke}{rgb}{0.000000,0.000000,0.000000}%
\pgfsetstrokecolor{currentstroke}%
\pgfsetdash{}{0pt}%
\pgfpathmoveto{\pgfqpoint{1.455507in}{0.924044in}}%
\pgfpathlineto{\pgfqpoint{1.464968in}{0.932165in}}%
\pgfpathlineto{\pgfqpoint{1.474468in}{0.940318in}}%
\pgfpathlineto{\pgfqpoint{1.505262in}{0.951448in}}%
\pgfpathlineto{\pgfqpoint{1.533923in}{0.963284in}}%
\pgfpathlineto{\pgfqpoint{1.523029in}{0.954622in}}%
\pgfpathlineto{\pgfqpoint{1.512175in}{0.945993in}}%
\pgfpathlineto{\pgfqpoint{1.484848in}{0.934679in}}%
\pgfpathlineto{\pgfqpoint{1.455507in}{0.924044in}}%
\pgfpathclose%
\pgfusepath{fill}%
\end{pgfscope}%
\begin{pgfscope}%
\pgfpathrectangle{\pgfqpoint{0.407592in}{0.222904in}}{\pgfqpoint{1.411163in}{1.411163in}}%
\pgfusepath{clip}%
\pgfsetbuttcap%
\pgfsetroundjoin%
\definecolor{currentfill}{rgb}{0.749020,0.772549,0.792157}%
\pgfsetfillcolor{currentfill}%
\pgfsetfillopacity{0.160000}%
\pgfsetlinewidth{0.000000pt}%
\definecolor{currentstroke}{rgb}{0.000000,0.000000,0.000000}%
\pgfsetstrokecolor{currentstroke}%
\pgfsetdash{}{0pt}%
\pgfpathmoveto{\pgfqpoint{1.759269in}{1.118802in}}%
\pgfpathlineto{\pgfqpoint{1.772311in}{1.128710in}}%
\pgfpathlineto{\pgfqpoint{1.796535in}{1.146423in}}%
\pgfpathlineto{\pgfqpoint{1.817258in}{1.164532in}}%
\pgfpathlineto{\pgfqpoint{1.803517in}{1.154061in}}%
\pgfpathlineto{\pgfqpoint{1.783107in}{1.136235in}}%
\pgfusepath{fill}%
\end{pgfscope}%
\begin{pgfscope}%
\pgfpathrectangle{\pgfqpoint{0.407592in}{0.222904in}}{\pgfqpoint{1.411163in}{1.411163in}}%
\pgfusepath{clip}%
\pgfsetbuttcap%
\pgfsetroundjoin%
\definecolor{currentfill}{rgb}{0.749020,0.772549,0.792157}%
\pgfsetfillcolor{currentfill}%
\pgfsetfillopacity{0.160000}%
\pgfsetlinewidth{0.000000pt}%
\definecolor{currentstroke}{rgb}{0.000000,0.000000,0.000000}%
\pgfsetstrokecolor{currentstroke}%
\pgfsetdash{}{0pt}%
\pgfpathmoveto{\pgfqpoint{0.828493in}{0.851308in}}%
\pgfpathlineto{\pgfqpoint{0.820619in}{0.858551in}}%
\pgfpathlineto{\pgfqpoint{0.812710in}{0.865827in}}%
\pgfpathlineto{\pgfqpoint{0.843734in}{0.859353in}}%
\pgfpathlineto{\pgfqpoint{0.876011in}{0.853744in}}%
\pgfpathlineto{\pgfqpoint{0.882137in}{0.846789in}}%
\pgfpathlineto{\pgfqpoint{0.888235in}{0.839866in}}%
\pgfpathlineto{\pgfqpoint{0.857778in}{0.845177in}}%
\pgfpathlineto{\pgfqpoint{0.828493in}{0.851308in}}%
\pgfpathclose%
\pgfusepath{fill}%
\end{pgfscope}%
\begin{pgfscope}%
\pgfpathrectangle{\pgfqpoint{0.407592in}{0.222904in}}{\pgfqpoint{1.411163in}{1.411163in}}%
\pgfusepath{clip}%
\pgfsetbuttcap%
\pgfsetroundjoin%
\definecolor{currentfill}{rgb}{0.749020,0.772549,0.792157}%
\pgfsetfillcolor{currentfill}%
\pgfsetfillopacity{0.160000}%
\pgfsetlinewidth{0.000000pt}%
\definecolor{currentstroke}{rgb}{0.000000,0.000000,0.000000}%
\pgfsetstrokecolor{currentstroke}%
\pgfsetdash{}{0pt}%
\pgfpathmoveto{\pgfqpoint{0.960359in}{0.818528in}}%
\pgfpathlineto{\pgfqpoint{0.956217in}{0.825161in}}%
\pgfpathlineto{\pgfqpoint{0.952056in}{0.831826in}}%
\pgfpathlineto{\pgfqpoint{0.985095in}{0.829145in}}%
\pgfpathlineto{\pgfqpoint{1.018663in}{0.827381in}}%
\pgfpathlineto{\pgfqpoint{1.020821in}{0.820841in}}%
\pgfpathlineto{\pgfqpoint{1.022968in}{0.814334in}}%
\pgfpathlineto{\pgfqpoint{0.991417in}{0.815999in}}%
\pgfpathlineto{\pgfqpoint{0.960359in}{0.818528in}}%
\pgfpathclose%
\pgfusepath{fill}%
\end{pgfscope}%
\begin{pgfscope}%
\pgfpathrectangle{\pgfqpoint{0.407592in}{0.222904in}}{\pgfqpoint{1.411163in}{1.411163in}}%
\pgfusepath{clip}%
\pgfsetbuttcap%
\pgfsetroundjoin%
\definecolor{currentfill}{rgb}{0.749020,0.772549,0.792157}%
\pgfsetfillcolor{currentfill}%
\pgfsetfillopacity{0.160000}%
\pgfsetlinewidth{0.000000pt}%
\definecolor{currentstroke}{rgb}{0.000000,0.000000,0.000000}%
\pgfsetstrokecolor{currentstroke}%
\pgfsetdash{}{0pt}%
\pgfpathmoveto{\pgfqpoint{0.587487in}{0.973208in}}%
\pgfpathlineto{\pgfqpoint{0.575136in}{0.982118in}}%
\pgfpathlineto{\pgfqpoint{0.562740in}{0.991060in}}%
\pgfpathlineto{\pgfqpoint{0.587763in}{0.978170in}}%
\pgfpathlineto{\pgfqpoint{0.615303in}{0.965956in}}%
\pgfpathlineto{\pgfqpoint{0.626550in}{0.957539in}}%
\pgfpathlineto{\pgfqpoint{0.637751in}{0.949156in}}%
\pgfpathlineto{\pgfqpoint{0.611427in}{0.960856in}}%
\pgfpathlineto{\pgfqpoint{0.587487in}{0.973208in}}%
\pgfpathclose%
\pgfusepath{fill}%
\end{pgfscope}%
\begin{pgfscope}%
\pgfpathrectangle{\pgfqpoint{0.407592in}{0.222904in}}{\pgfqpoint{1.411163in}{1.411163in}}%
\pgfusepath{clip}%
\pgfsetbuttcap%
\pgfsetroundjoin%
\definecolor{currentfill}{rgb}{0.749020,0.772549,0.792157}%
\pgfsetfillcolor{currentfill}%
\pgfsetfillopacity{0.160000}%
\pgfsetlinewidth{0.000000pt}%
\definecolor{currentstroke}{rgb}{0.000000,0.000000,0.000000}%
\pgfsetstrokecolor{currentstroke}%
\pgfsetdash{}{0pt}%
\pgfpathmoveto{\pgfqpoint{0.755023in}{0.881214in}}%
\pgfpathlineto{\pgfqpoint{0.745472in}{0.888925in}}%
\pgfpathlineto{\pgfqpoint{0.735882in}{0.896670in}}%
\pgfpathlineto{\pgfqpoint{0.765509in}{0.888162in}}%
\pgfpathlineto{\pgfqpoint{0.796785in}{0.880477in}}%
\pgfpathlineto{\pgfqpoint{0.804765in}{0.873135in}}%
\pgfpathlineto{\pgfqpoint{0.812710in}{0.865827in}}%
\pgfpathlineto{\pgfqpoint{0.783092in}{0.873128in}}%
\pgfpathlineto{\pgfqpoint{0.755023in}{0.881214in}}%
\pgfpathclose%
\pgfusepath{fill}%
\end{pgfscope}%
\begin{pgfscope}%
\pgfpathrectangle{\pgfqpoint{0.407592in}{0.222904in}}{\pgfqpoint{1.411163in}{1.411163in}}%
\pgfusepath{clip}%
\pgfsetbuttcap%
\pgfsetroundjoin%
\definecolor{currentfill}{rgb}{0.749020,0.772549,0.792157}%
\pgfsetfillcolor{currentfill}%
\pgfsetfillopacity{0.160000}%
\pgfsetlinewidth{0.000000pt}%
\definecolor{currentstroke}{rgb}{0.000000,0.000000,0.000000}%
\pgfsetstrokecolor{currentstroke}%
\pgfsetdash{}{0pt}%
\pgfpathmoveto{\pgfqpoint{1.086897in}{0.813631in}}%
\pgfpathlineto{\pgfqpoint{1.086791in}{0.820118in}}%
\pgfpathlineto{\pgfqpoint{1.086684in}{0.826636in}}%
\pgfpathlineto{\pgfqpoint{1.120784in}{0.827661in}}%
\pgfpathlineto{\pgfqpoint{1.154707in}{0.829610in}}%
\pgfpathlineto{\pgfqpoint{1.152764in}{0.823008in}}%
\pgfpathlineto{\pgfqpoint{1.150830in}{0.816438in}}%
\pgfpathlineto{\pgfqpoint{1.118946in}{0.814598in}}%
\pgfpathlineto{\pgfqpoint{1.086897in}{0.813631in}}%
\pgfpathclose%
\pgfusepath{fill}%
\end{pgfscope}%
\begin{pgfscope}%
\pgfpathrectangle{\pgfqpoint{0.407592in}{0.222904in}}{\pgfqpoint{1.411163in}{1.411163in}}%
\pgfusepath{clip}%
\pgfsetbuttcap%
\pgfsetroundjoin%
\definecolor{currentfill}{rgb}{0.749020,0.772549,0.792157}%
\pgfsetfillcolor{currentfill}%
\pgfsetfillopacity{0.160000}%
\pgfsetlinewidth{0.000000pt}%
\definecolor{currentstroke}{rgb}{0.000000,0.000000,0.000000}%
\pgfsetstrokecolor{currentstroke}%
\pgfsetdash{}{0pt}%
\pgfpathmoveto{\pgfqpoint{0.348632in}{1.150458in}}%
\pgfpathlineto{\pgfqpoint{0.334328in}{1.160872in}}%
\pgfpathlineto{\pgfqpoint{0.351245in}{1.142835in}}%
\pgfpathlineto{\pgfqpoint{0.371757in}{1.125212in}}%
\pgfpathlineto{\pgfqpoint{0.385509in}{1.115361in}}%
\pgfpathlineto{\pgfqpoint{0.365311in}{1.132703in}}%
\pgfusepath{fill}%
\end{pgfscope}%
\begin{pgfscope}%
\pgfpathrectangle{\pgfqpoint{0.407592in}{0.222904in}}{\pgfqpoint{1.411163in}{1.411163in}}%
\pgfusepath{clip}%
\pgfsetbuttcap%
\pgfsetroundjoin%
\definecolor{currentfill}{rgb}{0.749020,0.772549,0.792157}%
\pgfsetfillcolor{currentfill}%
\pgfsetfillopacity{0.160000}%
\pgfsetlinewidth{0.000000pt}%
\definecolor{currentstroke}{rgb}{0.000000,0.000000,0.000000}%
\pgfsetstrokecolor{currentstroke}%
\pgfsetdash{}{0pt}%
\pgfpathmoveto{\pgfqpoint{1.022968in}{0.814334in}}%
\pgfpathlineto{\pgfqpoint{1.020821in}{0.820841in}}%
\pgfpathlineto{\pgfqpoint{1.018663in}{0.827381in}}%
\pgfpathlineto{\pgfqpoint{1.052584in}{0.826543in}}%
\pgfpathlineto{\pgfqpoint{1.086684in}{0.826636in}}%
\pgfpathlineto{\pgfqpoint{1.086791in}{0.820118in}}%
\pgfpathlineto{\pgfqpoint{1.086897in}{0.813631in}}%
\pgfpathlineto{\pgfqpoint{1.054850in}{0.813543in}}%
\pgfpathlineto{\pgfqpoint{1.022968in}{0.814334in}}%
\pgfpathclose%
\pgfusepath{fill}%
\end{pgfscope}%
\begin{pgfscope}%
\pgfpathrectangle{\pgfqpoint{0.407592in}{0.222904in}}{\pgfqpoint{1.411163in}{1.411163in}}%
\pgfusepath{clip}%
\pgfsetbuttcap%
\pgfsetroundjoin%
\definecolor{currentfill}{rgb}{0.749020,0.772549,0.792157}%
\pgfsetfillcolor{currentfill}%
\pgfsetfillopacity{0.160000}%
\pgfsetlinewidth{0.000000pt}%
\definecolor{currentstroke}{rgb}{0.000000,0.000000,0.000000}%
\pgfsetstrokecolor{currentstroke}%
\pgfsetdash{}{0pt}%
\pgfpathmoveto{\pgfqpoint{1.221324in}{0.836231in}}%
\pgfpathlineto{\pgfqpoint{1.225286in}{0.843052in}}%
\pgfpathlineto{\pgfqpoint{1.229266in}{0.849906in}}%
\pgfpathlineto{\pgfqpoint{1.263563in}{0.854799in}}%
\pgfpathlineto{\pgfqpoint{1.296941in}{0.860584in}}%
\pgfpathlineto{\pgfqpoint{1.291041in}{0.853447in}}%
\pgfpathlineto{\pgfqpoint{1.285168in}{0.846342in}}%
\pgfpathlineto{\pgfqpoint{1.253675in}{0.840864in}}%
\pgfpathlineto{\pgfqpoint{1.221324in}{0.836231in}}%
\pgfpathclose%
\pgfusepath{fill}%
\end{pgfscope}%
\begin{pgfscope}%
\pgfpathrectangle{\pgfqpoint{0.407592in}{0.222904in}}{\pgfqpoint{1.411163in}{1.411163in}}%
\pgfusepath{clip}%
\pgfsetbuttcap%
\pgfsetroundjoin%
\definecolor{currentfill}{rgb}{0.749020,0.772549,0.792157}%
\pgfsetfillcolor{currentfill}%
\pgfsetfillopacity{0.160000}%
\pgfsetlinewidth{0.000000pt}%
\definecolor{currentstroke}{rgb}{0.000000,0.000000,0.000000}%
\pgfsetstrokecolor{currentstroke}%
\pgfsetdash{}{0pt}%
\pgfpathmoveto{\pgfqpoint{1.375732in}{0.889804in}}%
\pgfpathlineto{\pgfqpoint{1.383509in}{0.897411in}}%
\pgfpathlineto{\pgfqpoint{1.391320in}{0.905052in}}%
\pgfpathlineto{\pgfqpoint{1.424283in}{0.914149in}}%
\pgfpathlineto{\pgfqpoint{1.455507in}{0.924044in}}%
\pgfpathlineto{\pgfqpoint{1.446084in}{0.915957in}}%
\pgfpathlineto{\pgfqpoint{1.436700in}{0.907902in}}%
\pgfpathlineto{\pgfqpoint{1.407034in}{0.898472in}}%
\pgfpathlineto{\pgfqpoint{1.375732in}{0.889804in}}%
\pgfpathclose%
\pgfusepath{fill}%
\end{pgfscope}%
\begin{pgfscope}%
\pgfpathrectangle{\pgfqpoint{0.407592in}{0.222904in}}{\pgfqpoint{1.411163in}{1.411163in}}%
\pgfusepath{clip}%
\pgfsetbuttcap%
\pgfsetroundjoin%
\definecolor{currentfill}{rgb}{0.749020,0.772549,0.792157}%
\pgfsetfillcolor{currentfill}%
\pgfsetfillopacity{0.160000}%
\pgfsetlinewidth{0.000000pt}%
\definecolor{currentstroke}{rgb}{0.000000,0.000000,0.000000}%
\pgfsetstrokecolor{currentstroke}%
\pgfsetdash{}{0pt}%
\pgfpathmoveto{\pgfqpoint{1.296941in}{0.860584in}}%
\pgfpathlineto{\pgfqpoint{1.302868in}{0.867753in}}%
\pgfpathlineto{\pgfqpoint{1.308823in}{0.874956in}}%
\pgfpathlineto{\pgfqpoint{1.342942in}{0.881950in}}%
\pgfpathlineto{\pgfqpoint{1.375732in}{0.889804in}}%
\pgfpathlineto{\pgfqpoint{1.367989in}{0.882230in}}%
\pgfpathlineto{\pgfqpoint{1.360279in}{0.874689in}}%
\pgfpathlineto{\pgfqpoint{1.329233in}{0.867227in}}%
\pgfpathlineto{\pgfqpoint{1.296941in}{0.860584in}}%
\pgfpathclose%
\pgfusepath{fill}%
\end{pgfscope}%
\begin{pgfscope}%
\pgfpathrectangle{\pgfqpoint{0.407592in}{0.222904in}}{\pgfqpoint{1.411163in}{1.411163in}}%
\pgfusepath{clip}%
\pgfsetbuttcap%
\pgfsetroundjoin%
\definecolor{currentfill}{rgb}{0.749020,0.772549,0.792157}%
\pgfsetfillcolor{currentfill}%
\pgfsetfillopacity{0.160000}%
\pgfsetlinewidth{0.000000pt}%
\definecolor{currentstroke}{rgb}{0.000000,0.000000,0.000000}%
\pgfsetstrokecolor{currentstroke}%
\pgfsetdash{}{0pt}%
\pgfpathmoveto{\pgfqpoint{1.682166in}{1.063659in}}%
\pgfpathlineto{\pgfqpoint{1.694577in}{1.073160in}}%
\pgfpathlineto{\pgfqpoint{1.707030in}{1.082693in}}%
\pgfpathlineto{\pgfqpoint{1.733311in}{1.099082in}}%
\pgfpathlineto{\pgfqpoint{1.720395in}{1.089269in}}%
\pgfpathlineto{\pgfqpoint{1.707519in}{1.079488in}}%
\pgfusepath{fill}%
\end{pgfscope}%
\begin{pgfscope}%
\pgfpathrectangle{\pgfqpoint{0.407592in}{0.222904in}}{\pgfqpoint{1.411163in}{1.411163in}}%
\pgfusepath{clip}%
\pgfsetbuttcap%
\pgfsetroundjoin%
\definecolor{currentfill}{rgb}{0.749020,0.772549,0.792157}%
\pgfsetfillcolor{currentfill}%
\pgfsetfillopacity{0.160000}%
\pgfsetlinewidth{0.000000pt}%
\definecolor{currentstroke}{rgb}{0.000000,0.000000,0.000000}%
\pgfsetstrokecolor{currentstroke}%
\pgfsetdash{}{0pt}%
\pgfpathmoveto{\pgfqpoint{0.467095in}{1.056916in}}%
\pgfpathlineto{\pgfqpoint{0.453607in}{1.066578in}}%
\pgfpathlineto{\pgfqpoint{0.440076in}{1.076272in}}%
\pgfpathlineto{\pgfqpoint{0.462186in}{1.060548in}}%
\pgfpathlineto{\pgfqpoint{0.487416in}{1.045399in}}%
\pgfpathlineto{\pgfqpoint{0.500084in}{1.036260in}}%
\pgfpathlineto{\pgfqpoint{0.512706in}{1.027154in}}%
\pgfpathlineto{\pgfqpoint{0.488413in}{1.041755in}}%
\pgfpathlineto{\pgfqpoint{0.467095in}{1.056916in}}%
\pgfpathclose%
\pgfusepath{fill}%
\end{pgfscope}%
\begin{pgfscope}%
\pgfpathrectangle{\pgfqpoint{0.407592in}{0.222904in}}{\pgfqpoint{1.411163in}{1.411163in}}%
\pgfusepath{clip}%
\pgfsetbuttcap%
\pgfsetroundjoin%
\definecolor{currentfill}{rgb}{0.749020,0.772549,0.792157}%
\pgfsetfillcolor{currentfill}%
\pgfsetfillopacity{0.160000}%
\pgfsetlinewidth{0.000000pt}%
\definecolor{currentstroke}{rgb}{0.000000,0.000000,0.000000}%
\pgfsetstrokecolor{currentstroke}%
\pgfsetdash{}{0pt}%
\pgfpathmoveto{\pgfqpoint{1.577911in}{0.998256in}}%
\pgfpathlineto{\pgfqpoint{1.589012in}{1.007081in}}%
\pgfpathlineto{\pgfqpoint{1.600155in}{1.015940in}}%
\pgfpathlineto{\pgfqpoint{1.630192in}{1.030035in}}%
\pgfpathlineto{\pgfqpoint{1.657470in}{1.044755in}}%
\pgfpathlineto{\pgfqpoint{1.645185in}{1.035351in}}%
\pgfpathlineto{\pgfqpoint{1.632941in}{1.025978in}}%
\pgfpathlineto{\pgfqpoint{1.606738in}{1.011813in}}%
\pgfpathlineto{\pgfqpoint{1.577911in}{0.998256in}}%
\pgfpathclose%
\pgfusepath{fill}%
\end{pgfscope}%
\begin{pgfscope}%
\pgfpathrectangle{\pgfqpoint{0.407592in}{0.222904in}}{\pgfqpoint{1.411163in}{1.411163in}}%
\pgfusepath{clip}%
\pgfsetbuttcap%
\pgfsetroundjoin%
\definecolor{currentfill}{rgb}{0.749020,0.772549,0.792157}%
\pgfsetfillcolor{currentfill}%
\pgfsetfillopacity{0.160000}%
\pgfsetlinewidth{0.000000pt}%
\definecolor{currentstroke}{rgb}{0.000000,0.000000,0.000000}%
\pgfsetstrokecolor{currentstroke}%
\pgfsetdash{}{0pt}%
\pgfpathmoveto{\pgfqpoint{0.888235in}{0.839866in}}%
\pgfpathlineto{\pgfqpoint{0.882137in}{0.846789in}}%
\pgfpathlineto{\pgfqpoint{0.876011in}{0.853744in}}%
\pgfpathlineto{\pgfqpoint{0.909380in}{0.849035in}}%
\pgfpathlineto{\pgfqpoint{0.943672in}{0.845252in}}%
\pgfpathlineto{\pgfqpoint{0.947874in}{0.838523in}}%
\pgfpathlineto{\pgfqpoint{0.952056in}{0.831826in}}%
\pgfpathlineto{\pgfqpoint{0.919714in}{0.835406in}}%
\pgfpathlineto{\pgfqpoint{0.888235in}{0.839866in}}%
\pgfpathclose%
\pgfusepath{fill}%
\end{pgfscope}%
\begin{pgfscope}%
\pgfpathrectangle{\pgfqpoint{0.407592in}{0.222904in}}{\pgfqpoint{1.411163in}{1.411163in}}%
\pgfusepath{clip}%
\pgfsetbuttcap%
\pgfsetroundjoin%
\definecolor{currentfill}{rgb}{0.749020,0.772549,0.792157}%
\pgfsetfillcolor{currentfill}%
\pgfsetfillopacity{0.160000}%
\pgfsetlinewidth{0.000000pt}%
\definecolor{currentstroke}{rgb}{0.000000,0.000000,0.000000}%
\pgfsetstrokecolor{currentstroke}%
\pgfsetdash{}{0pt}%
\pgfpathmoveto{\pgfqpoint{1.474468in}{0.940318in}}%
\pgfpathlineto{\pgfqpoint{1.484007in}{0.948506in}}%
\pgfpathlineto{\pgfqpoint{1.493586in}{0.956726in}}%
\pgfpathlineto{\pgfqpoint{1.525837in}{0.968349in}}%
\pgfpathlineto{\pgfqpoint{1.555834in}{0.980704in}}%
\pgfpathlineto{\pgfqpoint{1.544858in}{0.971978in}}%
\pgfpathlineto{\pgfqpoint{1.533923in}{0.963284in}}%
\pgfpathlineto{\pgfqpoint{1.505262in}{0.951448in}}%
\pgfpathlineto{\pgfqpoint{1.474468in}{0.940318in}}%
\pgfpathclose%
\pgfusepath{fill}%
\end{pgfscope}%
\begin{pgfscope}%
\pgfpathrectangle{\pgfqpoint{0.407592in}{0.222904in}}{\pgfqpoint{1.411163in}{1.411163in}}%
\pgfusepath{clip}%
\pgfsetbuttcap%
\pgfsetroundjoin%
\definecolor{currentfill}{rgb}{0.749020,0.772549,0.792157}%
\pgfsetfillcolor{currentfill}%
\pgfsetfillopacity{0.160000}%
\pgfsetlinewidth{0.000000pt}%
\definecolor{currentstroke}{rgb}{0.000000,0.000000,0.000000}%
\pgfsetstrokecolor{currentstroke}%
\pgfsetdash{}{0pt}%
\pgfpathmoveto{\pgfqpoint{0.660022in}{0.932489in}}%
\pgfpathlineto{\pgfqpoint{0.648909in}{0.940806in}}%
\pgfpathlineto{\pgfqpoint{0.637751in}{0.949156in}}%
\pgfpathlineto{\pgfqpoint{0.666353in}{0.938177in}}%
\pgfpathlineto{\pgfqpoint{0.697107in}{0.927980in}}%
\pgfpathlineto{\pgfqpoint{0.706863in}{0.920103in}}%
\pgfpathlineto{\pgfqpoint{0.716577in}{0.912258in}}%
\pgfpathlineto{\pgfqpoint{0.687284in}{0.921998in}}%
\pgfpathlineto{\pgfqpoint{0.660022in}{0.932489in}}%
\pgfpathclose%
\pgfusepath{fill}%
\end{pgfscope}%
\begin{pgfscope}%
\pgfpathrectangle{\pgfqpoint{0.407592in}{0.222904in}}{\pgfqpoint{1.411163in}{1.411163in}}%
\pgfusepath{clip}%
\pgfsetbuttcap%
\pgfsetroundjoin%
\definecolor{currentfill}{rgb}{0.749020,0.772549,0.792157}%
\pgfsetfillcolor{currentfill}%
\pgfsetfillopacity{0.160000}%
\pgfsetlinewidth{0.000000pt}%
\definecolor{currentstroke}{rgb}{0.000000,0.000000,0.000000}%
\pgfsetstrokecolor{currentstroke}%
\pgfsetdash{}{0pt}%
\pgfpathmoveto{\pgfqpoint{1.154707in}{0.829610in}}%
\pgfpathlineto{\pgfqpoint{1.156660in}{0.836245in}}%
\pgfpathlineto{\pgfqpoint{1.158622in}{0.842912in}}%
\pgfpathlineto{\pgfqpoint{1.194226in}{0.845936in}}%
\pgfpathlineto{\pgfqpoint{1.229266in}{0.849906in}}%
\pgfpathlineto{\pgfqpoint{1.225286in}{0.843052in}}%
\pgfpathlineto{\pgfqpoint{1.221324in}{0.836231in}}%
\pgfpathlineto{\pgfqpoint{1.188278in}{0.832473in}}%
\pgfpathlineto{\pgfqpoint{1.154707in}{0.829610in}}%
\pgfpathclose%
\pgfusepath{fill}%
\end{pgfscope}%
\begin{pgfscope}%
\pgfpathrectangle{\pgfqpoint{0.407592in}{0.222904in}}{\pgfqpoint{1.411163in}{1.411163in}}%
\pgfusepath{clip}%
\pgfsetbuttcap%
\pgfsetroundjoin%
\definecolor{currentfill}{rgb}{0.749020,0.772549,0.792157}%
\pgfsetfillcolor{currentfill}%
\pgfsetfillopacity{0.160000}%
\pgfsetlinewidth{0.000000pt}%
\definecolor{currentstroke}{rgb}{0.000000,0.000000,0.000000}%
\pgfsetstrokecolor{currentstroke}%
\pgfsetdash{}{0pt}%
\pgfpathmoveto{\pgfqpoint{0.562740in}{0.991060in}}%
\pgfpathlineto{\pgfqpoint{0.550299in}{1.000035in}}%
\pgfpathlineto{\pgfqpoint{0.537814in}{1.009042in}}%
\pgfpathlineto{\pgfqpoint{0.563918in}{0.995616in}}%
\pgfpathlineto{\pgfqpoint{0.592676in}{0.982890in}}%
\pgfpathlineto{\pgfqpoint{0.604012in}{0.974406in}}%
\pgfpathlineto{\pgfqpoint{0.615303in}{0.965956in}}%
\pgfpathlineto{\pgfqpoint{0.587763in}{0.978170in}}%
\pgfpathlineto{\pgfqpoint{0.562740in}{0.991060in}}%
\pgfpathclose%
\pgfusepath{fill}%
\end{pgfscope}%
\begin{pgfscope}%
\pgfpathrectangle{\pgfqpoint{0.407592in}{0.222904in}}{\pgfqpoint{1.411163in}{1.411163in}}%
\pgfusepath{clip}%
\pgfsetbuttcap%
\pgfsetroundjoin%
\definecolor{currentfill}{rgb}{0.749020,0.772549,0.792157}%
\pgfsetfillcolor{currentfill}%
\pgfsetfillopacity{0.160000}%
\pgfsetlinewidth{0.000000pt}%
\definecolor{currentstroke}{rgb}{0.000000,0.000000,0.000000}%
\pgfsetstrokecolor{currentstroke}%
\pgfsetdash{}{0pt}%
\pgfpathmoveto{\pgfqpoint{0.812710in}{0.865827in}}%
\pgfpathlineto{\pgfqpoint{0.804765in}{0.873135in}}%
\pgfpathlineto{\pgfqpoint{0.796785in}{0.880477in}}%
\pgfpathlineto{\pgfqpoint{0.829560in}{0.873661in}}%
\pgfpathlineto{\pgfqpoint{0.863671in}{0.867754in}}%
\pgfpathlineto{\pgfqpoint{0.869855in}{0.860733in}}%
\pgfpathlineto{\pgfqpoint{0.876011in}{0.853744in}}%
\pgfpathlineto{\pgfqpoint{0.843734in}{0.859353in}}%
\pgfpathlineto{\pgfqpoint{0.812710in}{0.865827in}}%
\pgfpathclose%
\pgfusepath{fill}%
\end{pgfscope}%
\begin{pgfscope}%
\pgfpathrectangle{\pgfqpoint{0.407592in}{0.222904in}}{\pgfqpoint{1.411163in}{1.411163in}}%
\pgfusepath{clip}%
\pgfsetbuttcap%
\pgfsetroundjoin%
\definecolor{currentfill}{rgb}{0.749020,0.772549,0.792157}%
\pgfsetfillcolor{currentfill}%
\pgfsetfillopacity{0.160000}%
\pgfsetlinewidth{0.000000pt}%
\definecolor{currentstroke}{rgb}{0.000000,0.000000,0.000000}%
\pgfsetstrokecolor{currentstroke}%
\pgfsetdash{}{0pt}%
\pgfpathmoveto{\pgfqpoint{0.735882in}{0.896670in}}%
\pgfpathlineto{\pgfqpoint{0.726250in}{0.904448in}}%
\pgfpathlineto{\pgfqpoint{0.716577in}{0.912258in}}%
\pgfpathlineto{\pgfqpoint{0.747770in}{0.903328in}}%
\pgfpathlineto{\pgfqpoint{0.780715in}{0.895259in}}%
\pgfpathlineto{\pgfqpoint{0.788768in}{0.887851in}}%
\pgfpathlineto{\pgfqpoint{0.796785in}{0.880477in}}%
\pgfpathlineto{\pgfqpoint{0.765509in}{0.888162in}}%
\pgfpathlineto{\pgfqpoint{0.735882in}{0.896670in}}%
\pgfpathclose%
\pgfusepath{fill}%
\end{pgfscope}%
\begin{pgfscope}%
\pgfpathrectangle{\pgfqpoint{0.407592in}{0.222904in}}{\pgfqpoint{1.411163in}{1.411163in}}%
\pgfusepath{clip}%
\pgfsetbuttcap%
\pgfsetroundjoin%
\definecolor{currentfill}{rgb}{0.749020,0.772549,0.792157}%
\pgfsetfillcolor{currentfill}%
\pgfsetfillopacity{0.160000}%
\pgfsetlinewidth{0.000000pt}%
\definecolor{currentstroke}{rgb}{0.000000,0.000000,0.000000}%
\pgfsetstrokecolor{currentstroke}%
\pgfsetdash{}{0pt}%
\pgfpathmoveto{\pgfqpoint{0.952056in}{0.831826in}}%
\pgfpathlineto{\pgfqpoint{0.947874in}{0.838523in}}%
\pgfpathlineto{\pgfqpoint{0.943672in}{0.845252in}}%
\pgfpathlineto{\pgfqpoint{0.978711in}{0.842420in}}%
\pgfpathlineto{\pgfqpoint{1.014314in}{0.840556in}}%
\pgfpathlineto{\pgfqpoint{1.016494in}{0.833952in}}%
\pgfpathlineto{\pgfqpoint{1.018663in}{0.827381in}}%
\pgfpathlineto{\pgfqpoint{0.985095in}{0.829145in}}%
\pgfpathlineto{\pgfqpoint{0.952056in}{0.831826in}}%
\pgfpathclose%
\pgfusepath{fill}%
\end{pgfscope}%
\begin{pgfscope}%
\pgfpathrectangle{\pgfqpoint{0.407592in}{0.222904in}}{\pgfqpoint{1.411163in}{1.411163in}}%
\pgfusepath{clip}%
\pgfsetbuttcap%
\pgfsetroundjoin%
\definecolor{currentfill}{rgb}{0.749020,0.772549,0.792157}%
\pgfsetfillcolor{currentfill}%
\pgfsetfillopacity{0.160000}%
\pgfsetlinewidth{0.000000pt}%
\definecolor{currentstroke}{rgb}{0.000000,0.000000,0.000000}%
\pgfsetstrokecolor{currentstroke}%
\pgfsetdash{}{0pt}%
\pgfpathmoveto{\pgfqpoint{1.086684in}{0.826636in}}%
\pgfpathlineto{\pgfqpoint{1.086576in}{0.833186in}}%
\pgfpathlineto{\pgfqpoint{1.086468in}{0.839769in}}%
\pgfpathlineto{\pgfqpoint{1.122640in}{0.840852in}}%
\pgfpathlineto{\pgfqpoint{1.158622in}{0.842912in}}%
\pgfpathlineto{\pgfqpoint{1.156660in}{0.836245in}}%
\pgfpathlineto{\pgfqpoint{1.154707in}{0.829610in}}%
\pgfpathlineto{\pgfqpoint{1.120784in}{0.827661in}}%
\pgfpathlineto{\pgfqpoint{1.086684in}{0.826636in}}%
\pgfpathclose%
\pgfusepath{fill}%
\end{pgfscope}%
\begin{pgfscope}%
\pgfpathrectangle{\pgfqpoint{0.407592in}{0.222904in}}{\pgfqpoint{1.411163in}{1.411163in}}%
\pgfusepath{clip}%
\pgfsetbuttcap%
\pgfsetroundjoin%
\definecolor{currentfill}{rgb}{0.749020,0.772549,0.792157}%
\pgfsetfillcolor{currentfill}%
\pgfsetfillopacity{0.160000}%
\pgfsetlinewidth{0.000000pt}%
\definecolor{currentstroke}{rgb}{0.000000,0.000000,0.000000}%
\pgfsetstrokecolor{currentstroke}%
\pgfsetdash{}{0pt}%
\pgfpathmoveto{\pgfqpoint{1.707030in}{1.082693in}}%
\pgfpathlineto{\pgfqpoint{1.719526in}{1.092258in}}%
\pgfpathlineto{\pgfqpoint{1.732064in}{1.101857in}}%
\pgfpathlineto{\pgfqpoint{1.759269in}{1.118802in}}%
\pgfpathlineto{\pgfqpoint{1.746270in}{1.108926in}}%
\pgfpathlineto{\pgfqpoint{1.733311in}{1.099082in}}%
\pgfusepath{fill}%
\end{pgfscope}%
\begin{pgfscope}%
\pgfpathrectangle{\pgfqpoint{0.407592in}{0.222904in}}{\pgfqpoint{1.411163in}{1.411163in}}%
\pgfusepath{clip}%
\pgfsetbuttcap%
\pgfsetroundjoin%
\definecolor{currentfill}{rgb}{0.749020,0.772549,0.792157}%
\pgfsetfillcolor{currentfill}%
\pgfsetfillopacity{0.160000}%
\pgfsetlinewidth{0.000000pt}%
\definecolor{currentstroke}{rgb}{0.000000,0.000000,0.000000}%
\pgfsetstrokecolor{currentstroke}%
\pgfsetdash{}{0pt}%
\pgfpathmoveto{\pgfqpoint{1.018663in}{0.827381in}}%
\pgfpathlineto{\pgfqpoint{1.016494in}{0.833952in}}%
\pgfpathlineto{\pgfqpoint{1.014314in}{0.840556in}}%
\pgfpathlineto{\pgfqpoint{1.050296in}{0.839670in}}%
\pgfpathlineto{\pgfqpoint{1.086468in}{0.839769in}}%
\pgfpathlineto{\pgfqpoint{1.086576in}{0.833186in}}%
\pgfpathlineto{\pgfqpoint{1.086684in}{0.826636in}}%
\pgfpathlineto{\pgfqpoint{1.052584in}{0.826543in}}%
\pgfpathlineto{\pgfqpoint{1.018663in}{0.827381in}}%
\pgfpathclose%
\pgfusepath{fill}%
\end{pgfscope}%
\begin{pgfscope}%
\pgfpathrectangle{\pgfqpoint{0.407592in}{0.222904in}}{\pgfqpoint{1.411163in}{1.411163in}}%
\pgfusepath{clip}%
\pgfsetbuttcap%
\pgfsetroundjoin%
\definecolor{currentfill}{rgb}{0.749020,0.772549,0.792157}%
\pgfsetfillcolor{currentfill}%
\pgfsetfillopacity{0.160000}%
\pgfsetlinewidth{0.000000pt}%
\definecolor{currentstroke}{rgb}{0.000000,0.000000,0.000000}%
\pgfsetstrokecolor{currentstroke}%
\pgfsetdash{}{0pt}%
\pgfpathmoveto{\pgfqpoint{1.391320in}{0.905052in}}%
\pgfpathlineto{\pgfqpoint{1.399165in}{0.912726in}}%
\pgfpathlineto{\pgfqpoint{1.407045in}{0.920434in}}%
\pgfpathlineto{\pgfqpoint{1.441679in}{0.929960in}}%
\pgfpathlineto{\pgfqpoint{1.474468in}{0.940318in}}%
\pgfpathlineto{\pgfqpoint{1.464968in}{0.932165in}}%
\pgfpathlineto{\pgfqpoint{1.455507in}{0.924044in}}%
\pgfpathlineto{\pgfqpoint{1.424283in}{0.914149in}}%
\pgfpathlineto{\pgfqpoint{1.391320in}{0.905052in}}%
\pgfpathclose%
\pgfusepath{fill}%
\end{pgfscope}%
\begin{pgfscope}%
\pgfpathrectangle{\pgfqpoint{0.407592in}{0.222904in}}{\pgfqpoint{1.411163in}{1.411163in}}%
\pgfusepath{clip}%
\pgfsetbuttcap%
\pgfsetroundjoin%
\definecolor{currentfill}{rgb}{0.749020,0.772549,0.792157}%
\pgfsetfillcolor{currentfill}%
\pgfsetfillopacity{0.160000}%
\pgfsetlinewidth{0.000000pt}%
\definecolor{currentstroke}{rgb}{0.000000,0.000000,0.000000}%
\pgfsetstrokecolor{currentstroke}%
\pgfsetdash{}{0pt}%
\pgfpathmoveto{\pgfqpoint{0.440076in}{1.076272in}}%
\pgfpathlineto{\pgfqpoint{0.426500in}{1.085997in}}%
\pgfpathlineto{\pgfqpoint{0.412881in}{1.095753in}}%
\pgfpathlineto{\pgfqpoint{0.435778in}{1.079471in}}%
\pgfpathlineto{\pgfqpoint{0.461940in}{1.063777in}}%
\pgfpathlineto{\pgfqpoint{0.474701in}{1.054571in}}%
\pgfpathlineto{\pgfqpoint{0.487416in}{1.045399in}}%
\pgfpathlineto{\pgfqpoint{0.462186in}{1.060548in}}%
\pgfpathlineto{\pgfqpoint{0.440076in}{1.076272in}}%
\pgfpathclose%
\pgfusepath{fill}%
\end{pgfscope}%
\begin{pgfscope}%
\pgfpathrectangle{\pgfqpoint{0.407592in}{0.222904in}}{\pgfqpoint{1.411163in}{1.411163in}}%
\pgfusepath{clip}%
\pgfsetbuttcap%
\pgfsetroundjoin%
\definecolor{currentfill}{rgb}{0.749020,0.772549,0.792157}%
\pgfsetfillcolor{currentfill}%
\pgfsetfillopacity{0.160000}%
\pgfsetlinewidth{0.000000pt}%
\definecolor{currentstroke}{rgb}{0.000000,0.000000,0.000000}%
\pgfsetstrokecolor{currentstroke}%
\pgfsetdash{}{0pt}%
\pgfpathmoveto{\pgfqpoint{1.229266in}{0.849906in}}%
\pgfpathlineto{\pgfqpoint{1.233266in}{0.856793in}}%
\pgfpathlineto{\pgfqpoint{1.237285in}{0.863712in}}%
\pgfpathlineto{\pgfqpoint{1.273544in}{0.868865in}}%
\pgfpathlineto{\pgfqpoint{1.308823in}{0.874956in}}%
\pgfpathlineto{\pgfqpoint{1.302868in}{0.867753in}}%
\pgfpathlineto{\pgfqpoint{1.296941in}{0.860584in}}%
\pgfpathlineto{\pgfqpoint{1.263563in}{0.854799in}}%
\pgfpathlineto{\pgfqpoint{1.229266in}{0.849906in}}%
\pgfpathclose%
\pgfusepath{fill}%
\end{pgfscope}%
\begin{pgfscope}%
\pgfpathrectangle{\pgfqpoint{0.407592in}{0.222904in}}{\pgfqpoint{1.411163in}{1.411163in}}%
\pgfusepath{clip}%
\pgfsetbuttcap%
\pgfsetroundjoin%
\definecolor{currentfill}{rgb}{0.749020,0.772549,0.792157}%
\pgfsetfillcolor{currentfill}%
\pgfsetfillopacity{0.160000}%
\pgfsetlinewidth{0.000000pt}%
\definecolor{currentstroke}{rgb}{0.000000,0.000000,0.000000}%
\pgfsetstrokecolor{currentstroke}%
\pgfsetdash{}{0pt}%
\pgfpathmoveto{\pgfqpoint{1.600155in}{1.015940in}}%
\pgfpathlineto{\pgfqpoint{1.611340in}{1.024833in}}%
\pgfpathlineto{\pgfqpoint{1.622568in}{1.033759in}}%
\pgfpathlineto{\pgfqpoint{1.653815in}{1.048388in}}%
\pgfpathlineto{\pgfqpoint{1.682166in}{1.063659in}}%
\pgfpathlineto{\pgfqpoint{1.669797in}{1.054191in}}%
\pgfpathlineto{\pgfqpoint{1.657470in}{1.044755in}}%
\pgfpathlineto{\pgfqpoint{1.630192in}{1.030035in}}%
\pgfpathlineto{\pgfqpoint{1.600155in}{1.015940in}}%
\pgfpathclose%
\pgfusepath{fill}%
\end{pgfscope}%
\begin{pgfscope}%
\pgfpathrectangle{\pgfqpoint{0.407592in}{0.222904in}}{\pgfqpoint{1.411163in}{1.411163in}}%
\pgfusepath{clip}%
\pgfsetbuttcap%
\pgfsetroundjoin%
\definecolor{currentfill}{rgb}{0.749020,0.772549,0.792157}%
\pgfsetfillcolor{currentfill}%
\pgfsetfillopacity{0.160000}%
\pgfsetlinewidth{0.000000pt}%
\definecolor{currentstroke}{rgb}{0.000000,0.000000,0.000000}%
\pgfsetstrokecolor{currentstroke}%
\pgfsetdash{}{0pt}%
\pgfpathmoveto{\pgfqpoint{1.308823in}{0.874956in}}%
\pgfpathlineto{\pgfqpoint{1.314805in}{0.882192in}}%
\pgfpathlineto{\pgfqpoint{1.320814in}{0.889462in}}%
\pgfpathlineto{\pgfqpoint{1.356774in}{0.896806in}}%
\pgfpathlineto{\pgfqpoint{1.391320in}{0.905052in}}%
\pgfpathlineto{\pgfqpoint{1.383509in}{0.897411in}}%
\pgfpathlineto{\pgfqpoint{1.375732in}{0.889804in}}%
\pgfpathlineto{\pgfqpoint{1.342942in}{0.881950in}}%
\pgfpathlineto{\pgfqpoint{1.308823in}{0.874956in}}%
\pgfpathclose%
\pgfusepath{fill}%
\end{pgfscope}%
\begin{pgfscope}%
\pgfpathrectangle{\pgfqpoint{0.407592in}{0.222904in}}{\pgfqpoint{1.411163in}{1.411163in}}%
\pgfusepath{clip}%
\pgfsetbuttcap%
\pgfsetroundjoin%
\definecolor{currentfill}{rgb}{0.749020,0.772549,0.792157}%
\pgfsetfillcolor{currentfill}%
\pgfsetfillopacity{0.160000}%
\pgfsetlinewidth{0.000000pt}%
\definecolor{currentstroke}{rgb}{0.000000,0.000000,0.000000}%
\pgfsetstrokecolor{currentstroke}%
\pgfsetdash{}{0pt}%
\pgfpathmoveto{\pgfqpoint{1.493586in}{0.956726in}}%
\pgfpathlineto{\pgfqpoint{1.503203in}{0.964981in}}%
\pgfpathlineto{\pgfqpoint{1.512861in}{0.973270in}}%
\pgfpathlineto{\pgfqpoint{1.546575in}{0.985383in}}%
\pgfpathlineto{\pgfqpoint{1.577911in}{0.998256in}}%
\pgfpathlineto{\pgfqpoint{1.566852in}{0.989463in}}%
\pgfpathlineto{\pgfqpoint{1.555834in}{0.980704in}}%
\pgfpathlineto{\pgfqpoint{1.525837in}{0.968349in}}%
\pgfpathlineto{\pgfqpoint{1.493586in}{0.956726in}}%
\pgfpathclose%
\pgfusepath{fill}%
\end{pgfscope}%
\begin{pgfscope}%
\pgfpathrectangle{\pgfqpoint{0.407592in}{0.222904in}}{\pgfqpoint{1.411163in}{1.411163in}}%
\pgfusepath{clip}%
\pgfsetbuttcap%
\pgfsetroundjoin%
\definecolor{currentfill}{rgb}{0.749020,0.772549,0.792157}%
\pgfsetfillcolor{currentfill}%
\pgfsetfillopacity{0.160000}%
\pgfsetlinewidth{0.000000pt}%
\definecolor{currentstroke}{rgb}{0.000000,0.000000,0.000000}%
\pgfsetstrokecolor{currentstroke}%
\pgfsetdash{}{0pt}%
\pgfpathmoveto{\pgfqpoint{0.637751in}{0.949156in}}%
\pgfpathlineto{\pgfqpoint{0.626550in}{0.957539in}}%
\pgfpathlineto{\pgfqpoint{0.615303in}{0.965956in}}%
\pgfpathlineto{\pgfqpoint{0.645249in}{0.954489in}}%
\pgfpathlineto{\pgfqpoint{0.677470in}{0.943837in}}%
\pgfpathlineto{\pgfqpoint{0.687310in}{0.935892in}}%
\pgfpathlineto{\pgfqpoint{0.697107in}{0.927980in}}%
\pgfpathlineto{\pgfqpoint{0.666353in}{0.938177in}}%
\pgfpathlineto{\pgfqpoint{0.637751in}{0.949156in}}%
\pgfpathclose%
\pgfusepath{fill}%
\end{pgfscope}%
\begin{pgfscope}%
\pgfpathrectangle{\pgfqpoint{0.407592in}{0.222904in}}{\pgfqpoint{1.411163in}{1.411163in}}%
\pgfusepath{clip}%
\pgfsetbuttcap%
\pgfsetroundjoin%
\definecolor{currentfill}{rgb}{0.749020,0.772549,0.792157}%
\pgfsetfillcolor{currentfill}%
\pgfsetfillopacity{0.160000}%
\pgfsetlinewidth{0.000000pt}%
\definecolor{currentstroke}{rgb}{0.000000,0.000000,0.000000}%
\pgfsetstrokecolor{currentstroke}%
\pgfsetdash{}{0pt}%
\pgfpathmoveto{\pgfqpoint{0.876011in}{0.853744in}}%
\pgfpathlineto{\pgfqpoint{0.869855in}{0.860733in}}%
\pgfpathlineto{\pgfqpoint{0.863671in}{0.867754in}}%
\pgfpathlineto{\pgfqpoint{0.898947in}{0.862794in}}%
\pgfpathlineto{\pgfqpoint{0.935207in}{0.858810in}}%
\pgfpathlineto{\pgfqpoint{0.939450in}{0.852015in}}%
\pgfpathlineto{\pgfqpoint{0.943672in}{0.845252in}}%
\pgfpathlineto{\pgfqpoint{0.909380in}{0.849035in}}%
\pgfpathlineto{\pgfqpoint{0.876011in}{0.853744in}}%
\pgfpathclose%
\pgfusepath{fill}%
\end{pgfscope}%
\begin{pgfscope}%
\pgfpathrectangle{\pgfqpoint{0.407592in}{0.222904in}}{\pgfqpoint{1.411163in}{1.411163in}}%
\pgfusepath{clip}%
\pgfsetbuttcap%
\pgfsetroundjoin%
\definecolor{currentfill}{rgb}{0.749020,0.772549,0.792157}%
\pgfsetfillcolor{currentfill}%
\pgfsetfillopacity{0.160000}%
\pgfsetlinewidth{0.000000pt}%
\definecolor{currentstroke}{rgb}{0.000000,0.000000,0.000000}%
\pgfsetstrokecolor{currentstroke}%
\pgfsetdash{}{0pt}%
\pgfpathmoveto{\pgfqpoint{0.537814in}{1.009042in}}%
\pgfpathlineto{\pgfqpoint{0.525283in}{1.018082in}}%
\pgfpathlineto{\pgfqpoint{0.512706in}{1.027154in}}%
\pgfpathlineto{\pgfqpoint{0.539892in}{1.013195in}}%
\pgfpathlineto{\pgfqpoint{0.569868in}{0.999959in}}%
\pgfpathlineto{\pgfqpoint{0.581295in}{0.991407in}}%
\pgfpathlineto{\pgfqpoint{0.592676in}{0.982890in}}%
\pgfpathlineto{\pgfqpoint{0.563918in}{0.995616in}}%
\pgfpathlineto{\pgfqpoint{0.537814in}{1.009042in}}%
\pgfpathclose%
\pgfusepath{fill}%
\end{pgfscope}%
\begin{pgfscope}%
\pgfpathrectangle{\pgfqpoint{0.407592in}{0.222904in}}{\pgfqpoint{1.411163in}{1.411163in}}%
\pgfusepath{clip}%
\pgfsetbuttcap%
\pgfsetroundjoin%
\definecolor{currentfill}{rgb}{0.749020,0.772549,0.792157}%
\pgfsetfillcolor{currentfill}%
\pgfsetfillopacity{0.160000}%
\pgfsetlinewidth{0.000000pt}%
\definecolor{currentstroke}{rgb}{0.000000,0.000000,0.000000}%
\pgfsetstrokecolor{currentstroke}%
\pgfsetdash{}{0pt}%
\pgfpathmoveto{\pgfqpoint{1.732064in}{1.101857in}}%
\pgfpathlineto{\pgfqpoint{1.744646in}{1.111488in}}%
\pgfpathlineto{\pgfqpoint{1.772311in}{1.128710in}}%
\pgfpathlineto{\pgfqpoint{1.759269in}{1.118802in}}%
\pgfusepath{fill}%
\end{pgfscope}%
\begin{pgfscope}%
\pgfpathrectangle{\pgfqpoint{0.407592in}{0.222904in}}{\pgfqpoint{1.411163in}{1.411163in}}%
\pgfusepath{clip}%
\pgfsetbuttcap%
\pgfsetroundjoin%
\definecolor{currentfill}{rgb}{0.749020,0.772549,0.792157}%
\pgfsetfillcolor{currentfill}%
\pgfsetfillopacity{0.160000}%
\pgfsetlinewidth{0.000000pt}%
\definecolor{currentstroke}{rgb}{0.000000,0.000000,0.000000}%
\pgfsetstrokecolor{currentstroke}%
\pgfsetdash{}{0pt}%
\pgfpathmoveto{\pgfqpoint{1.158622in}{0.842912in}}%
\pgfpathlineto{\pgfqpoint{1.160594in}{0.849611in}}%
\pgfpathlineto{\pgfqpoint{1.162575in}{0.856344in}}%
\pgfpathlineto{\pgfqpoint{1.200231in}{0.859530in}}%
\pgfpathlineto{\pgfqpoint{1.237285in}{0.863712in}}%
\pgfpathlineto{\pgfqpoint{1.233266in}{0.856793in}}%
\pgfpathlineto{\pgfqpoint{1.229266in}{0.849906in}}%
\pgfpathlineto{\pgfqpoint{1.194226in}{0.845936in}}%
\pgfpathlineto{\pgfqpoint{1.158622in}{0.842912in}}%
\pgfpathclose%
\pgfusepath{fill}%
\end{pgfscope}%
\begin{pgfscope}%
\pgfpathrectangle{\pgfqpoint{0.407592in}{0.222904in}}{\pgfqpoint{1.411163in}{1.411163in}}%
\pgfusepath{clip}%
\pgfsetbuttcap%
\pgfsetroundjoin%
\definecolor{currentfill}{rgb}{0.749020,0.772549,0.792157}%
\pgfsetfillcolor{currentfill}%
\pgfsetfillopacity{0.160000}%
\pgfsetlinewidth{0.000000pt}%
\definecolor{currentstroke}{rgb}{0.000000,0.000000,0.000000}%
\pgfsetstrokecolor{currentstroke}%
\pgfsetdash{}{0pt}%
\pgfpathmoveto{\pgfqpoint{0.716577in}{0.912258in}}%
\pgfpathlineto{\pgfqpoint{0.706863in}{0.920103in}}%
\pgfpathlineto{\pgfqpoint{0.697107in}{0.927980in}}%
\pgfpathlineto{\pgfqpoint{0.729874in}{0.918628in}}%
\pgfpathlineto{\pgfqpoint{0.764500in}{0.910176in}}%
\pgfpathlineto{\pgfqpoint{0.772626in}{0.902701in}}%
\pgfpathlineto{\pgfqpoint{0.780715in}{0.895259in}}%
\pgfpathlineto{\pgfqpoint{0.747770in}{0.903328in}}%
\pgfpathlineto{\pgfqpoint{0.716577in}{0.912258in}}%
\pgfpathclose%
\pgfusepath{fill}%
\end{pgfscope}%
\begin{pgfscope}%
\pgfpathrectangle{\pgfqpoint{0.407592in}{0.222904in}}{\pgfqpoint{1.411163in}{1.411163in}}%
\pgfusepath{clip}%
\pgfsetbuttcap%
\pgfsetroundjoin%
\definecolor{currentfill}{rgb}{0.749020,0.772549,0.792157}%
\pgfsetfillcolor{currentfill}%
\pgfsetfillopacity{0.160000}%
\pgfsetlinewidth{0.000000pt}%
\definecolor{currentstroke}{rgb}{0.000000,0.000000,0.000000}%
\pgfsetstrokecolor{currentstroke}%
\pgfsetdash{}{0pt}%
\pgfpathmoveto{\pgfqpoint{0.796785in}{0.880477in}}%
\pgfpathlineto{\pgfqpoint{0.788768in}{0.887851in}}%
\pgfpathlineto{\pgfqpoint{0.780715in}{0.895259in}}%
\pgfpathlineto{\pgfqpoint{0.815254in}{0.888101in}}%
\pgfpathlineto{\pgfqpoint{0.851214in}{0.881898in}}%
\pgfpathlineto{\pgfqpoint{0.857457in}{0.874809in}}%
\pgfpathlineto{\pgfqpoint{0.863671in}{0.867754in}}%
\pgfpathlineto{\pgfqpoint{0.829560in}{0.873661in}}%
\pgfpathlineto{\pgfqpoint{0.796785in}{0.880477in}}%
\pgfpathclose%
\pgfusepath{fill}%
\end{pgfscope}%
\begin{pgfscope}%
\pgfpathrectangle{\pgfqpoint{0.407592in}{0.222904in}}{\pgfqpoint{1.411163in}{1.411163in}}%
\pgfusepath{clip}%
\pgfsetbuttcap%
\pgfsetroundjoin%
\definecolor{currentfill}{rgb}{0.749020,0.772549,0.792157}%
\pgfsetfillcolor{currentfill}%
\pgfsetfillopacity{0.160000}%
\pgfsetlinewidth{0.000000pt}%
\definecolor{currentstroke}{rgb}{0.000000,0.000000,0.000000}%
\pgfsetstrokecolor{currentstroke}%
\pgfsetdash{}{0pt}%
\pgfpathmoveto{\pgfqpoint{0.412881in}{1.095753in}}%
\pgfpathlineto{\pgfqpoint{0.399217in}{1.105541in}}%
\pgfpathlineto{\pgfqpoint{0.385509in}{1.115361in}}%
\pgfpathlineto{\pgfqpoint{0.409186in}{1.098524in}}%
\pgfpathlineto{\pgfqpoint{0.436278in}{1.082289in}}%
\pgfpathlineto{\pgfqpoint{0.449132in}{1.073016in}}%
\pgfpathlineto{\pgfqpoint{0.461940in}{1.063777in}}%
\pgfpathlineto{\pgfqpoint{0.435778in}{1.079471in}}%
\pgfpathlineto{\pgfqpoint{0.412881in}{1.095753in}}%
\pgfpathclose%
\pgfusepath{fill}%
\end{pgfscope}%
\begin{pgfscope}%
\pgfpathrectangle{\pgfqpoint{0.407592in}{0.222904in}}{\pgfqpoint{1.411163in}{1.411163in}}%
\pgfusepath{clip}%
\pgfsetbuttcap%
\pgfsetroundjoin%
\definecolor{currentfill}{rgb}{0.749020,0.772549,0.792157}%
\pgfsetfillcolor{currentfill}%
\pgfsetfillopacity{0.160000}%
\pgfsetlinewidth{0.000000pt}%
\definecolor{currentstroke}{rgb}{0.000000,0.000000,0.000000}%
\pgfsetstrokecolor{currentstroke}%
\pgfsetdash{}{0pt}%
\pgfpathmoveto{\pgfqpoint{0.943672in}{0.845252in}}%
\pgfpathlineto{\pgfqpoint{0.939450in}{0.852015in}}%
\pgfpathlineto{\pgfqpoint{0.935207in}{0.858810in}}%
\pgfpathlineto{\pgfqpoint{0.972264in}{0.855826in}}%
\pgfpathlineto{\pgfqpoint{1.009923in}{0.853862in}}%
\pgfpathlineto{\pgfqpoint{1.012124in}{0.847192in}}%
\pgfpathlineto{\pgfqpoint{1.014314in}{0.840556in}}%
\pgfpathlineto{\pgfqpoint{0.978711in}{0.842420in}}%
\pgfpathlineto{\pgfqpoint{0.943672in}{0.845252in}}%
\pgfpathclose%
\pgfusepath{fill}%
\end{pgfscope}%
\begin{pgfscope}%
\pgfpathrectangle{\pgfqpoint{0.407592in}{0.222904in}}{\pgfqpoint{1.411163in}{1.411163in}}%
\pgfusepath{clip}%
\pgfsetbuttcap%
\pgfsetroundjoin%
\definecolor{currentfill}{rgb}{0.749020,0.772549,0.792157}%
\pgfsetfillcolor{currentfill}%
\pgfsetfillopacity{0.160000}%
\pgfsetlinewidth{0.000000pt}%
\definecolor{currentstroke}{rgb}{0.000000,0.000000,0.000000}%
\pgfsetstrokecolor{currentstroke}%
\pgfsetdash{}{0pt}%
\pgfpathmoveto{\pgfqpoint{1.622568in}{1.033759in}}%
\pgfpathlineto{\pgfqpoint{1.633838in}{1.042720in}}%
\pgfpathlineto{\pgfqpoint{1.645152in}{1.051715in}}%
\pgfpathlineto{\pgfqpoint{1.677609in}{1.066874in}}%
\pgfpathlineto{\pgfqpoint{1.707030in}{1.082693in}}%
\pgfpathlineto{\pgfqpoint{1.694577in}{1.073160in}}%
\pgfpathlineto{\pgfqpoint{1.682166in}{1.063659in}}%
\pgfpathlineto{\pgfqpoint{1.653815in}{1.048388in}}%
\pgfpathlineto{\pgfqpoint{1.622568in}{1.033759in}}%
\pgfpathclose%
\pgfusepath{fill}%
\end{pgfscope}%
\begin{pgfscope}%
\pgfpathrectangle{\pgfqpoint{0.407592in}{0.222904in}}{\pgfqpoint{1.411163in}{1.411163in}}%
\pgfusepath{clip}%
\pgfsetbuttcap%
\pgfsetroundjoin%
\definecolor{currentfill}{rgb}{0.749020,0.772549,0.792157}%
\pgfsetfillcolor{currentfill}%
\pgfsetfillopacity{0.160000}%
\pgfsetlinewidth{0.000000pt}%
\definecolor{currentstroke}{rgb}{0.000000,0.000000,0.000000}%
\pgfsetstrokecolor{currentstroke}%
\pgfsetdash{}{0pt}%
\pgfpathmoveto{\pgfqpoint{1.407045in}{0.920434in}}%
\pgfpathlineto{\pgfqpoint{1.414960in}{0.928176in}}%
\pgfpathlineto{\pgfqpoint{1.422909in}{0.935952in}}%
\pgfpathlineto{\pgfqpoint{1.459224in}{0.945906in}}%
\pgfpathlineto{\pgfqpoint{1.493586in}{0.956726in}}%
\pgfpathlineto{\pgfqpoint{1.484007in}{0.948506in}}%
\pgfpathlineto{\pgfqpoint{1.474468in}{0.940318in}}%
\pgfpathlineto{\pgfqpoint{1.441679in}{0.929960in}}%
\pgfpathlineto{\pgfqpoint{1.407045in}{0.920434in}}%
\pgfpathclose%
\pgfusepath{fill}%
\end{pgfscope}%
\begin{pgfscope}%
\pgfpathrectangle{\pgfqpoint{0.407592in}{0.222904in}}{\pgfqpoint{1.411163in}{1.411163in}}%
\pgfusepath{clip}%
\pgfsetbuttcap%
\pgfsetroundjoin%
\definecolor{currentfill}{rgb}{0.749020,0.772549,0.792157}%
\pgfsetfillcolor{currentfill}%
\pgfsetfillopacity{0.160000}%
\pgfsetlinewidth{0.000000pt}%
\definecolor{currentstroke}{rgb}{0.000000,0.000000,0.000000}%
\pgfsetstrokecolor{currentstroke}%
\pgfsetdash{}{0pt}%
\pgfpathmoveto{\pgfqpoint{1.086468in}{0.839769in}}%
\pgfpathlineto{\pgfqpoint{1.086360in}{0.846385in}}%
\pgfpathlineto{\pgfqpoint{1.086251in}{0.853033in}}%
\pgfpathlineto{\pgfqpoint{1.124515in}{0.854174in}}%
\pgfpathlineto{\pgfqpoint{1.162575in}{0.856344in}}%
\pgfpathlineto{\pgfqpoint{1.160594in}{0.849611in}}%
\pgfpathlineto{\pgfqpoint{1.158622in}{0.842912in}}%
\pgfpathlineto{\pgfqpoint{1.122640in}{0.840852in}}%
\pgfpathlineto{\pgfqpoint{1.086468in}{0.839769in}}%
\pgfpathclose%
\pgfusepath{fill}%
\end{pgfscope}%
\begin{pgfscope}%
\pgfpathrectangle{\pgfqpoint{0.407592in}{0.222904in}}{\pgfqpoint{1.411163in}{1.411163in}}%
\pgfusepath{clip}%
\pgfsetbuttcap%
\pgfsetroundjoin%
\definecolor{currentfill}{rgb}{0.749020,0.772549,0.792157}%
\pgfsetfillcolor{currentfill}%
\pgfsetfillopacity{0.160000}%
\pgfsetlinewidth{0.000000pt}%
\definecolor{currentstroke}{rgb}{0.000000,0.000000,0.000000}%
\pgfsetstrokecolor{currentstroke}%
\pgfsetdash{}{0pt}%
\pgfpathmoveto{\pgfqpoint{1.014314in}{0.840556in}}%
\pgfpathlineto{\pgfqpoint{1.012124in}{0.847192in}}%
\pgfpathlineto{\pgfqpoint{1.009923in}{0.853862in}}%
\pgfpathlineto{\pgfqpoint{1.047986in}{0.852929in}}%
\pgfpathlineto{\pgfqpoint{1.086251in}{0.853033in}}%
\pgfpathlineto{\pgfqpoint{1.086360in}{0.846385in}}%
\pgfpathlineto{\pgfqpoint{1.086468in}{0.839769in}}%
\pgfpathlineto{\pgfqpoint{1.050296in}{0.839670in}}%
\pgfpathlineto{\pgfqpoint{1.014314in}{0.840556in}}%
\pgfpathclose%
\pgfusepath{fill}%
\end{pgfscope}%
\begin{pgfscope}%
\pgfpathrectangle{\pgfqpoint{0.407592in}{0.222904in}}{\pgfqpoint{1.411163in}{1.411163in}}%
\pgfusepath{clip}%
\pgfsetbuttcap%
\pgfsetroundjoin%
\definecolor{currentfill}{rgb}{0.749020,0.772549,0.792157}%
\pgfsetfillcolor{currentfill}%
\pgfsetfillopacity{0.160000}%
\pgfsetlinewidth{0.000000pt}%
\definecolor{currentstroke}{rgb}{0.000000,0.000000,0.000000}%
\pgfsetstrokecolor{currentstroke}%
\pgfsetdash{}{0pt}%
\pgfpathmoveto{\pgfqpoint{1.237285in}{0.863712in}}%
\pgfpathlineto{\pgfqpoint{1.241324in}{0.870664in}}%
\pgfpathlineto{\pgfqpoint{1.245381in}{0.877651in}}%
\pgfpathlineto{\pgfqpoint{1.283620in}{0.883064in}}%
\pgfpathlineto{\pgfqpoint{1.320814in}{0.889462in}}%
\pgfpathlineto{\pgfqpoint{1.314805in}{0.882192in}}%
\pgfpathlineto{\pgfqpoint{1.308823in}{0.874956in}}%
\pgfpathlineto{\pgfqpoint{1.273544in}{0.868865in}}%
\pgfpathlineto{\pgfqpoint{1.237285in}{0.863712in}}%
\pgfpathclose%
\pgfusepath{fill}%
\end{pgfscope}%
\begin{pgfscope}%
\pgfpathrectangle{\pgfqpoint{0.407592in}{0.222904in}}{\pgfqpoint{1.411163in}{1.411163in}}%
\pgfusepath{clip}%
\pgfsetbuttcap%
\pgfsetroundjoin%
\definecolor{currentfill}{rgb}{0.749020,0.772549,0.792157}%
\pgfsetfillcolor{currentfill}%
\pgfsetfillopacity{0.160000}%
\pgfsetlinewidth{0.000000pt}%
\definecolor{currentstroke}{rgb}{0.000000,0.000000,0.000000}%
\pgfsetstrokecolor{currentstroke}%
\pgfsetdash{}{0pt}%
\pgfpathmoveto{\pgfqpoint{1.512861in}{0.973270in}}%
\pgfpathlineto{\pgfqpoint{1.522558in}{0.981593in}}%
\pgfpathlineto{\pgfqpoint{1.532296in}{0.989951in}}%
\pgfpathlineto{\pgfqpoint{1.567477in}{1.002553in}}%
\pgfpathlineto{\pgfqpoint{1.600155in}{1.015940in}}%
\pgfpathlineto{\pgfqpoint{1.589012in}{1.007081in}}%
\pgfpathlineto{\pgfqpoint{1.577911in}{0.998256in}}%
\pgfpathlineto{\pgfqpoint{1.546575in}{0.985383in}}%
\pgfpathlineto{\pgfqpoint{1.512861in}{0.973270in}}%
\pgfpathclose%
\pgfusepath{fill}%
\end{pgfscope}%
\begin{pgfscope}%
\pgfpathrectangle{\pgfqpoint{0.407592in}{0.222904in}}{\pgfqpoint{1.411163in}{1.411163in}}%
\pgfusepath{clip}%
\pgfsetbuttcap%
\pgfsetroundjoin%
\definecolor{currentfill}{rgb}{0.749020,0.772549,0.792157}%
\pgfsetfillcolor{currentfill}%
\pgfsetfillopacity{0.160000}%
\pgfsetlinewidth{0.000000pt}%
\definecolor{currentstroke}{rgb}{0.000000,0.000000,0.000000}%
\pgfsetstrokecolor{currentstroke}%
\pgfsetdash{}{0pt}%
\pgfpathmoveto{\pgfqpoint{1.320814in}{0.889462in}}%
\pgfpathlineto{\pgfqpoint{1.326852in}{0.896765in}}%
\pgfpathlineto{\pgfqpoint{1.332917in}{0.904103in}}%
\pgfpathlineto{\pgfqpoint{1.370732in}{0.911797in}}%
\pgfpathlineto{\pgfqpoint{1.407045in}{0.920434in}}%
\pgfpathlineto{\pgfqpoint{1.399165in}{0.912726in}}%
\pgfpathlineto{\pgfqpoint{1.391320in}{0.905052in}}%
\pgfpathlineto{\pgfqpoint{1.356774in}{0.896806in}}%
\pgfpathlineto{\pgfqpoint{1.320814in}{0.889462in}}%
\pgfpathclose%
\pgfusepath{fill}%
\end{pgfscope}%
\begin{pgfscope}%
\pgfpathrectangle{\pgfqpoint{0.407592in}{0.222904in}}{\pgfqpoint{1.411163in}{1.411163in}}%
\pgfusepath{clip}%
\pgfsetbuttcap%
\pgfsetroundjoin%
\definecolor{currentfill}{rgb}{0.749020,0.772549,0.792157}%
\pgfsetfillcolor{currentfill}%
\pgfsetfillopacity{0.160000}%
\pgfsetlinewidth{0.000000pt}%
\definecolor{currentstroke}{rgb}{0.000000,0.000000,0.000000}%
\pgfsetstrokecolor{currentstroke}%
\pgfsetdash{}{0pt}%
\pgfpathmoveto{\pgfqpoint{0.615303in}{0.965956in}}%
\pgfpathlineto{\pgfqpoint{0.604012in}{0.974406in}}%
\pgfpathlineto{\pgfqpoint{0.592676in}{0.982890in}}%
\pgfpathlineto{\pgfqpoint{0.623970in}{0.970938in}}%
\pgfpathlineto{\pgfqpoint{0.657663in}{0.959832in}}%
\pgfpathlineto{\pgfqpoint{0.667588in}{0.951817in}}%
\pgfpathlineto{\pgfqpoint{0.677470in}{0.943837in}}%
\pgfpathlineto{\pgfqpoint{0.645249in}{0.954489in}}%
\pgfpathlineto{\pgfqpoint{0.615303in}{0.965956in}}%
\pgfpathclose%
\pgfusepath{fill}%
\end{pgfscope}%
\begin{pgfscope}%
\pgfpathrectangle{\pgfqpoint{0.407592in}{0.222904in}}{\pgfqpoint{1.411163in}{1.411163in}}%
\pgfusepath{clip}%
\pgfsetbuttcap%
\pgfsetroundjoin%
\definecolor{currentfill}{rgb}{0.749020,0.772549,0.792157}%
\pgfsetfillcolor{currentfill}%
\pgfsetfillopacity{0.160000}%
\pgfsetlinewidth{0.000000pt}%
\definecolor{currentstroke}{rgb}{0.000000,0.000000,0.000000}%
\pgfsetstrokecolor{currentstroke}%
\pgfsetdash{}{0pt}%
\pgfpathmoveto{\pgfqpoint{0.512706in}{1.027154in}}%
\pgfpathlineto{\pgfqpoint{0.500084in}{1.036260in}}%
\pgfpathlineto{\pgfqpoint{0.487416in}{1.045399in}}%
\pgfpathlineto{\pgfqpoint{0.515681in}{1.030910in}}%
\pgfpathlineto{\pgfqpoint{0.546875in}{1.017166in}}%
\pgfpathlineto{\pgfqpoint{0.558395in}{1.008545in}}%
\pgfpathlineto{\pgfqpoint{0.569868in}{0.999959in}}%
\pgfpathlineto{\pgfqpoint{0.539892in}{1.013195in}}%
\pgfpathlineto{\pgfqpoint{0.512706in}{1.027154in}}%
\pgfpathclose%
\pgfusepath{fill}%
\end{pgfscope}%
\begin{pgfscope}%
\pgfpathrectangle{\pgfqpoint{0.407592in}{0.222904in}}{\pgfqpoint{1.411163in}{1.411163in}}%
\pgfusepath{clip}%
\pgfsetbuttcap%
\pgfsetroundjoin%
\definecolor{currentfill}{rgb}{0.749020,0.772549,0.792157}%
\pgfsetfillcolor{currentfill}%
\pgfsetfillopacity{0.160000}%
\pgfsetlinewidth{0.000000pt}%
\definecolor{currentstroke}{rgb}{0.000000,0.000000,0.000000}%
\pgfsetstrokecolor{currentstroke}%
\pgfsetdash{}{0pt}%
\pgfpathmoveto{\pgfqpoint{0.385509in}{1.115361in}}%
\pgfpathlineto{\pgfqpoint{0.371757in}{1.125212in}}%
\pgfpathlineto{\pgfqpoint{0.395821in}{1.108101in}}%
\pgfpathlineto{\pgfqpoint{0.423376in}{1.091597in}}%
\pgfpathlineto{\pgfqpoint{0.436278in}{1.082289in}}%
\pgfpathlineto{\pgfqpoint{0.409186in}{1.098524in}}%
\pgfusepath{fill}%
\end{pgfscope}%
\begin{pgfscope}%
\pgfpathrectangle{\pgfqpoint{0.407592in}{0.222904in}}{\pgfqpoint{1.411163in}{1.411163in}}%
\pgfusepath{clip}%
\pgfsetbuttcap%
\pgfsetroundjoin%
\definecolor{currentfill}{rgb}{0.749020,0.772549,0.792157}%
\pgfsetfillcolor{currentfill}%
\pgfsetfillopacity{0.160000}%
\pgfsetlinewidth{0.000000pt}%
\definecolor{currentstroke}{rgb}{0.000000,0.000000,0.000000}%
\pgfsetstrokecolor{currentstroke}%
\pgfsetdash{}{0pt}%
\pgfpathmoveto{\pgfqpoint{0.863671in}{0.867754in}}%
\pgfpathlineto{\pgfqpoint{0.857457in}{0.874809in}}%
\pgfpathlineto{\pgfqpoint{0.851214in}{0.881898in}}%
\pgfpathlineto{\pgfqpoint{0.888413in}{0.876686in}}%
\pgfpathlineto{\pgfqpoint{0.926659in}{0.872500in}}%
\pgfpathlineto{\pgfqpoint{0.930944in}{0.865638in}}%
\pgfpathlineto{\pgfqpoint{0.935207in}{0.858810in}}%
\pgfpathlineto{\pgfqpoint{0.898947in}{0.862794in}}%
\pgfpathlineto{\pgfqpoint{0.863671in}{0.867754in}}%
\pgfpathclose%
\pgfusepath{fill}%
\end{pgfscope}%
\begin{pgfscope}%
\pgfpathrectangle{\pgfqpoint{0.407592in}{0.222904in}}{\pgfqpoint{1.411163in}{1.411163in}}%
\pgfusepath{clip}%
\pgfsetbuttcap%
\pgfsetroundjoin%
\definecolor{currentfill}{rgb}{0.749020,0.772549,0.792157}%
\pgfsetfillcolor{currentfill}%
\pgfsetfillopacity{0.160000}%
\pgfsetlinewidth{0.000000pt}%
\definecolor{currentstroke}{rgb}{0.000000,0.000000,0.000000}%
\pgfsetstrokecolor{currentstroke}%
\pgfsetdash{}{0pt}%
\pgfpathmoveto{\pgfqpoint{1.645152in}{1.051715in}}%
\pgfpathlineto{\pgfqpoint{1.656509in}{1.060744in}}%
\pgfpathlineto{\pgfqpoint{1.667909in}{1.069807in}}%
\pgfpathlineto{\pgfqpoint{1.701576in}{1.085494in}}%
\pgfpathlineto{\pgfqpoint{1.732064in}{1.101857in}}%
\pgfpathlineto{\pgfqpoint{1.719526in}{1.092258in}}%
\pgfpathlineto{\pgfqpoint{1.707030in}{1.082693in}}%
\pgfpathlineto{\pgfqpoint{1.677609in}{1.066874in}}%
\pgfpathlineto{\pgfqpoint{1.645152in}{1.051715in}}%
\pgfpathclose%
\pgfusepath{fill}%
\end{pgfscope}%
\begin{pgfscope}%
\pgfpathrectangle{\pgfqpoint{0.407592in}{0.222904in}}{\pgfqpoint{1.411163in}{1.411163in}}%
\pgfusepath{clip}%
\pgfsetbuttcap%
\pgfsetroundjoin%
\definecolor{currentfill}{rgb}{0.749020,0.772549,0.792157}%
\pgfsetfillcolor{currentfill}%
\pgfsetfillopacity{0.160000}%
\pgfsetlinewidth{0.000000pt}%
\definecolor{currentstroke}{rgb}{0.000000,0.000000,0.000000}%
\pgfsetstrokecolor{currentstroke}%
\pgfsetdash{}{0pt}%
\pgfpathmoveto{\pgfqpoint{0.697107in}{0.927980in}}%
\pgfpathlineto{\pgfqpoint{0.687310in}{0.935892in}}%
\pgfpathlineto{\pgfqpoint{0.677470in}{0.943837in}}%
\pgfpathlineto{\pgfqpoint{0.711820in}{0.934065in}}%
\pgfpathlineto{\pgfqpoint{0.748135in}{0.925230in}}%
\pgfpathlineto{\pgfqpoint{0.756336in}{0.917686in}}%
\pgfpathlineto{\pgfqpoint{0.764500in}{0.910176in}}%
\pgfpathlineto{\pgfqpoint{0.729874in}{0.918628in}}%
\pgfpathlineto{\pgfqpoint{0.697107in}{0.927980in}}%
\pgfpathclose%
\pgfusepath{fill}%
\end{pgfscope}%
\begin{pgfscope}%
\pgfpathrectangle{\pgfqpoint{0.407592in}{0.222904in}}{\pgfqpoint{1.411163in}{1.411163in}}%
\pgfusepath{clip}%
\pgfsetbuttcap%
\pgfsetroundjoin%
\definecolor{currentfill}{rgb}{0.749020,0.772549,0.792157}%
\pgfsetfillcolor{currentfill}%
\pgfsetfillopacity{0.160000}%
\pgfsetlinewidth{0.000000pt}%
\definecolor{currentstroke}{rgb}{0.000000,0.000000,0.000000}%
\pgfsetstrokecolor{currentstroke}%
\pgfsetdash{}{0pt}%
\pgfpathmoveto{\pgfqpoint{1.162575in}{0.856344in}}%
\pgfpathlineto{\pgfqpoint{1.164567in}{0.863109in}}%
\pgfpathlineto{\pgfqpoint{1.166568in}{0.869908in}}%
\pgfpathlineto{\pgfqpoint{1.206296in}{0.873256in}}%
\pgfpathlineto{\pgfqpoint{1.245381in}{0.877651in}}%
\pgfpathlineto{\pgfqpoint{1.241324in}{0.870664in}}%
\pgfpathlineto{\pgfqpoint{1.237285in}{0.863712in}}%
\pgfpathlineto{\pgfqpoint{1.200231in}{0.859530in}}%
\pgfpathlineto{\pgfqpoint{1.162575in}{0.856344in}}%
\pgfpathclose%
\pgfusepath{fill}%
\end{pgfscope}%
\begin{pgfscope}%
\pgfpathrectangle{\pgfqpoint{0.407592in}{0.222904in}}{\pgfqpoint{1.411163in}{1.411163in}}%
\pgfusepath{clip}%
\pgfsetbuttcap%
\pgfsetroundjoin%
\definecolor{currentfill}{rgb}{0.749020,0.772549,0.792157}%
\pgfsetfillcolor{currentfill}%
\pgfsetfillopacity{0.160000}%
\pgfsetlinewidth{0.000000pt}%
\definecolor{currentstroke}{rgb}{0.000000,0.000000,0.000000}%
\pgfsetstrokecolor{currentstroke}%
\pgfsetdash{}{0pt}%
\pgfpathmoveto{\pgfqpoint{0.780715in}{0.895259in}}%
\pgfpathlineto{\pgfqpoint{0.772626in}{0.902701in}}%
\pgfpathlineto{\pgfqpoint{0.764500in}{0.910176in}}%
\pgfpathlineto{\pgfqpoint{0.800814in}{0.902677in}}%
\pgfpathlineto{\pgfqpoint{0.838638in}{0.896175in}}%
\pgfpathlineto{\pgfqpoint{0.844941in}{0.889020in}}%
\pgfpathlineto{\pgfqpoint{0.851214in}{0.881898in}}%
\pgfpathlineto{\pgfqpoint{0.815254in}{0.888101in}}%
\pgfpathlineto{\pgfqpoint{0.780715in}{0.895259in}}%
\pgfpathclose%
\pgfusepath{fill}%
\end{pgfscope}%
\begin{pgfscope}%
\pgfpathrectangle{\pgfqpoint{0.407592in}{0.222904in}}{\pgfqpoint{1.411163in}{1.411163in}}%
\pgfusepath{clip}%
\pgfsetbuttcap%
\pgfsetroundjoin%
\definecolor{currentfill}{rgb}{0.749020,0.772549,0.792157}%
\pgfsetfillcolor{currentfill}%
\pgfsetfillopacity{0.160000}%
\pgfsetlinewidth{0.000000pt}%
\definecolor{currentstroke}{rgb}{0.000000,0.000000,0.000000}%
\pgfsetstrokecolor{currentstroke}%
\pgfsetdash{}{0pt}%
\pgfpathmoveto{\pgfqpoint{1.422909in}{0.935952in}}%
\pgfpathlineto{\pgfqpoint{1.430894in}{0.943762in}}%
\pgfpathlineto{\pgfqpoint{1.438914in}{0.951607in}}%
\pgfpathlineto{\pgfqpoint{1.476919in}{0.961989in}}%
\pgfpathlineto{\pgfqpoint{1.512861in}{0.973270in}}%
\pgfpathlineto{\pgfqpoint{1.503203in}{0.964981in}}%
\pgfpathlineto{\pgfqpoint{1.493586in}{0.956726in}}%
\pgfpathlineto{\pgfqpoint{1.459224in}{0.945906in}}%
\pgfpathlineto{\pgfqpoint{1.422909in}{0.935952in}}%
\pgfpathclose%
\pgfusepath{fill}%
\end{pgfscope}%
\begin{pgfscope}%
\pgfpathrectangle{\pgfqpoint{0.407592in}{0.222904in}}{\pgfqpoint{1.411163in}{1.411163in}}%
\pgfusepath{clip}%
\pgfsetbuttcap%
\pgfsetroundjoin%
\definecolor{currentfill}{rgb}{0.749020,0.772549,0.792157}%
\pgfsetfillcolor{currentfill}%
\pgfsetfillopacity{0.160000}%
\pgfsetlinewidth{0.000000pt}%
\definecolor{currentstroke}{rgb}{0.000000,0.000000,0.000000}%
\pgfsetstrokecolor{currentstroke}%
\pgfsetdash{}{0pt}%
\pgfpathmoveto{\pgfqpoint{0.935207in}{0.858810in}}%
\pgfpathlineto{\pgfqpoint{0.930944in}{0.865638in}}%
\pgfpathlineto{\pgfqpoint{0.926659in}{0.872500in}}%
\pgfpathlineto{\pgfqpoint{0.965753in}{0.869365in}}%
\pgfpathlineto{\pgfqpoint{1.005488in}{0.867300in}}%
\pgfpathlineto{\pgfqpoint{1.007711in}{0.860564in}}%
\pgfpathlineto{\pgfqpoint{1.009923in}{0.853862in}}%
\pgfpathlineto{\pgfqpoint{0.972264in}{0.855826in}}%
\pgfpathlineto{\pgfqpoint{0.935207in}{0.858810in}}%
\pgfpathclose%
\pgfusepath{fill}%
\end{pgfscope}%
\begin{pgfscope}%
\pgfpathrectangle{\pgfqpoint{0.407592in}{0.222904in}}{\pgfqpoint{1.411163in}{1.411163in}}%
\pgfusepath{clip}%
\pgfsetbuttcap%
\pgfsetroundjoin%
\definecolor{currentfill}{rgb}{0.749020,0.772549,0.792157}%
\pgfsetfillcolor{currentfill}%
\pgfsetfillopacity{0.160000}%
\pgfsetlinewidth{0.000000pt}%
\definecolor{currentstroke}{rgb}{0.000000,0.000000,0.000000}%
\pgfsetstrokecolor{currentstroke}%
\pgfsetdash{}{0pt}%
\pgfpathmoveto{\pgfqpoint{1.532296in}{0.989951in}}%
\pgfpathlineto{\pgfqpoint{1.542074in}{0.998343in}}%
\pgfpathlineto{\pgfqpoint{1.551892in}{1.006770in}}%
\pgfpathlineto{\pgfqpoint{1.588546in}{1.019859in}}%
\pgfpathlineto{\pgfqpoint{1.622568in}{1.033759in}}%
\pgfpathlineto{\pgfqpoint{1.611340in}{1.024833in}}%
\pgfpathlineto{\pgfqpoint{1.600155in}{1.015940in}}%
\pgfpathlineto{\pgfqpoint{1.567477in}{1.002553in}}%
\pgfpathlineto{\pgfqpoint{1.532296in}{0.989951in}}%
\pgfpathclose%
\pgfusepath{fill}%
\end{pgfscope}%
\begin{pgfscope}%
\pgfpathrectangle{\pgfqpoint{0.407592in}{0.222904in}}{\pgfqpoint{1.411163in}{1.411163in}}%
\pgfusepath{clip}%
\pgfsetbuttcap%
\pgfsetroundjoin%
\definecolor{currentfill}{rgb}{0.749020,0.772549,0.792157}%
\pgfsetfillcolor{currentfill}%
\pgfsetfillopacity{0.160000}%
\pgfsetlinewidth{0.000000pt}%
\definecolor{currentstroke}{rgb}{0.000000,0.000000,0.000000}%
\pgfsetstrokecolor{currentstroke}%
\pgfsetdash{}{0pt}%
\pgfpathmoveto{\pgfqpoint{1.086251in}{0.853033in}}%
\pgfpathlineto{\pgfqpoint{1.086141in}{0.859714in}}%
\pgfpathlineto{\pgfqpoint{1.086031in}{0.866429in}}%
\pgfpathlineto{\pgfqpoint{1.126408in}{0.867628in}}%
\pgfpathlineto{\pgfqpoint{1.166568in}{0.869908in}}%
\pgfpathlineto{\pgfqpoint{1.164567in}{0.863109in}}%
\pgfpathlineto{\pgfqpoint{1.162575in}{0.856344in}}%
\pgfpathlineto{\pgfqpoint{1.124515in}{0.854174in}}%
\pgfpathlineto{\pgfqpoint{1.086251in}{0.853033in}}%
\pgfpathclose%
\pgfusepath{fill}%
\end{pgfscope}%
\begin{pgfscope}%
\pgfpathrectangle{\pgfqpoint{0.407592in}{0.222904in}}{\pgfqpoint{1.411163in}{1.411163in}}%
\pgfusepath{clip}%
\pgfsetbuttcap%
\pgfsetroundjoin%
\definecolor{currentfill}{rgb}{0.749020,0.772549,0.792157}%
\pgfsetfillcolor{currentfill}%
\pgfsetfillopacity{0.160000}%
\pgfsetlinewidth{0.000000pt}%
\definecolor{currentstroke}{rgb}{0.000000,0.000000,0.000000}%
\pgfsetstrokecolor{currentstroke}%
\pgfsetdash{}{0pt}%
\pgfpathmoveto{\pgfqpoint{1.245381in}{0.877651in}}%
\pgfpathlineto{\pgfqpoint{1.249459in}{0.884670in}}%
\pgfpathlineto{\pgfqpoint{1.253556in}{0.891724in}}%
\pgfpathlineto{\pgfqpoint{1.293792in}{0.897398in}}%
\pgfpathlineto{\pgfqpoint{1.332917in}{0.904103in}}%
\pgfpathlineto{\pgfqpoint{1.326852in}{0.896765in}}%
\pgfpathlineto{\pgfqpoint{1.320814in}{0.889462in}}%
\pgfpathlineto{\pgfqpoint{1.283620in}{0.883064in}}%
\pgfpathlineto{\pgfqpoint{1.245381in}{0.877651in}}%
\pgfpathclose%
\pgfusepath{fill}%
\end{pgfscope}%
\begin{pgfscope}%
\pgfpathrectangle{\pgfqpoint{0.407592in}{0.222904in}}{\pgfqpoint{1.411163in}{1.411163in}}%
\pgfusepath{clip}%
\pgfsetbuttcap%
\pgfsetroundjoin%
\definecolor{currentfill}{rgb}{0.749020,0.772549,0.792157}%
\pgfsetfillcolor{currentfill}%
\pgfsetfillopacity{0.160000}%
\pgfsetlinewidth{0.000000pt}%
\definecolor{currentstroke}{rgb}{0.000000,0.000000,0.000000}%
\pgfsetstrokecolor{currentstroke}%
\pgfsetdash{}{0pt}%
\pgfpathmoveto{\pgfqpoint{0.487416in}{1.045399in}}%
\pgfpathlineto{\pgfqpoint{0.474701in}{1.054571in}}%
\pgfpathlineto{\pgfqpoint{0.461940in}{1.063777in}}%
\pgfpathlineto{\pgfqpoint{0.491283in}{1.048761in}}%
\pgfpathlineto{\pgfqpoint{0.523697in}{1.034512in}}%
\pgfpathlineto{\pgfqpoint{0.535310in}{1.025821in}}%
\pgfpathlineto{\pgfqpoint{0.546875in}{1.017166in}}%
\pgfpathlineto{\pgfqpoint{0.515681in}{1.030910in}}%
\pgfpathlineto{\pgfqpoint{0.487416in}{1.045399in}}%
\pgfpathclose%
\pgfusepath{fill}%
\end{pgfscope}%
\begin{pgfscope}%
\pgfpathrectangle{\pgfqpoint{0.407592in}{0.222904in}}{\pgfqpoint{1.411163in}{1.411163in}}%
\pgfusepath{clip}%
\pgfsetbuttcap%
\pgfsetroundjoin%
\definecolor{currentfill}{rgb}{0.749020,0.772549,0.792157}%
\pgfsetfillcolor{currentfill}%
\pgfsetfillopacity{0.160000}%
\pgfsetlinewidth{0.000000pt}%
\definecolor{currentstroke}{rgb}{0.000000,0.000000,0.000000}%
\pgfsetstrokecolor{currentstroke}%
\pgfsetdash{}{0pt}%
\pgfpathmoveto{\pgfqpoint{1.009923in}{0.853862in}}%
\pgfpathlineto{\pgfqpoint{1.007711in}{0.860564in}}%
\pgfpathlineto{\pgfqpoint{1.005488in}{0.867300in}}%
\pgfpathlineto{\pgfqpoint{1.045652in}{0.866319in}}%
\pgfpathlineto{\pgfqpoint{1.086031in}{0.866429in}}%
\pgfpathlineto{\pgfqpoint{1.086141in}{0.859714in}}%
\pgfpathlineto{\pgfqpoint{1.086251in}{0.853033in}}%
\pgfpathlineto{\pgfqpoint{1.047986in}{0.852929in}}%
\pgfpathlineto{\pgfqpoint{1.009923in}{0.853862in}}%
\pgfpathclose%
\pgfusepath{fill}%
\end{pgfscope}%
\begin{pgfscope}%
\pgfpathrectangle{\pgfqpoint{0.407592in}{0.222904in}}{\pgfqpoint{1.411163in}{1.411163in}}%
\pgfusepath{clip}%
\pgfsetbuttcap%
\pgfsetroundjoin%
\definecolor{currentfill}{rgb}{0.749020,0.772549,0.792157}%
\pgfsetfillcolor{currentfill}%
\pgfsetfillopacity{0.160000}%
\pgfsetlinewidth{0.000000pt}%
\definecolor{currentstroke}{rgb}{0.000000,0.000000,0.000000}%
\pgfsetstrokecolor{currentstroke}%
\pgfsetdash{}{0pt}%
\pgfpathmoveto{\pgfqpoint{1.332917in}{0.904103in}}%
\pgfpathlineto{\pgfqpoint{1.339011in}{0.911474in}}%
\pgfpathlineto{\pgfqpoint{1.345134in}{0.918880in}}%
\pgfpathlineto{\pgfqpoint{1.384817in}{0.926925in}}%
\pgfpathlineto{\pgfqpoint{1.422909in}{0.935952in}}%
\pgfpathlineto{\pgfqpoint{1.414960in}{0.928176in}}%
\pgfpathlineto{\pgfqpoint{1.407045in}{0.920434in}}%
\pgfpathlineto{\pgfqpoint{1.370732in}{0.911797in}}%
\pgfpathlineto{\pgfqpoint{1.332917in}{0.904103in}}%
\pgfpathclose%
\pgfusepath{fill}%
\end{pgfscope}%
\begin{pgfscope}%
\pgfpathrectangle{\pgfqpoint{0.407592in}{0.222904in}}{\pgfqpoint{1.411163in}{1.411163in}}%
\pgfusepath{clip}%
\pgfsetbuttcap%
\pgfsetroundjoin%
\definecolor{currentfill}{rgb}{0.749020,0.772549,0.792157}%
\pgfsetfillcolor{currentfill}%
\pgfsetfillopacity{0.160000}%
\pgfsetlinewidth{0.000000pt}%
\definecolor{currentstroke}{rgb}{0.000000,0.000000,0.000000}%
\pgfsetstrokecolor{currentstroke}%
\pgfsetdash{}{0pt}%
\pgfpathmoveto{\pgfqpoint{0.592676in}{0.982890in}}%
\pgfpathlineto{\pgfqpoint{0.581295in}{0.991407in}}%
\pgfpathlineto{\pgfqpoint{0.569868in}{0.999959in}}%
\pgfpathlineto{\pgfqpoint{0.602512in}{0.987524in}}%
\pgfpathlineto{\pgfqpoint{0.637684in}{0.975965in}}%
\pgfpathlineto{\pgfqpoint{0.647695in}{0.967881in}}%
\pgfpathlineto{\pgfqpoint{0.657663in}{0.959832in}}%
\pgfpathlineto{\pgfqpoint{0.623970in}{0.970938in}}%
\pgfpathlineto{\pgfqpoint{0.592676in}{0.982890in}}%
\pgfpathclose%
\pgfusepath{fill}%
\end{pgfscope}%
\begin{pgfscope}%
\pgfpathrectangle{\pgfqpoint{0.407592in}{0.222904in}}{\pgfqpoint{1.411163in}{1.411163in}}%
\pgfusepath{clip}%
\pgfsetbuttcap%
\pgfsetroundjoin%
\definecolor{currentfill}{rgb}{0.749020,0.772549,0.792157}%
\pgfsetfillcolor{currentfill}%
\pgfsetfillopacity{0.160000}%
\pgfsetlinewidth{0.000000pt}%
\definecolor{currentstroke}{rgb}{0.000000,0.000000,0.000000}%
\pgfsetstrokecolor{currentstroke}%
\pgfsetdash{}{0pt}%
\pgfpathmoveto{\pgfqpoint{1.667909in}{1.069807in}}%
\pgfpathlineto{\pgfqpoint{1.679353in}{1.078906in}}%
\pgfpathlineto{\pgfqpoint{1.713624in}{1.094855in}}%
\pgfpathlineto{\pgfqpoint{1.744646in}{1.111488in}}%
\pgfpathlineto{\pgfqpoint{1.732064in}{1.101857in}}%
\pgfpathlineto{\pgfqpoint{1.701576in}{1.085494in}}%
\pgfusepath{fill}%
\end{pgfscope}%
\begin{pgfscope}%
\pgfpathrectangle{\pgfqpoint{0.407592in}{0.222904in}}{\pgfqpoint{1.411163in}{1.411163in}}%
\pgfusepath{clip}%
\pgfsetbuttcap%
\pgfsetroundjoin%
\definecolor{currentfill}{rgb}{0.749020,0.772549,0.792157}%
\pgfsetfillcolor{currentfill}%
\pgfsetfillopacity{0.160000}%
\pgfsetlinewidth{0.000000pt}%
\definecolor{currentstroke}{rgb}{0.000000,0.000000,0.000000}%
\pgfsetstrokecolor{currentstroke}%
\pgfsetdash{}{0pt}%
\pgfpathmoveto{\pgfqpoint{0.851214in}{0.881898in}}%
\pgfpathlineto{\pgfqpoint{0.844941in}{0.889020in}}%
\pgfpathlineto{\pgfqpoint{0.838638in}{0.896175in}}%
\pgfpathlineto{\pgfqpoint{0.877776in}{0.890713in}}%
\pgfpathlineto{\pgfqpoint{0.918027in}{0.886325in}}%
\pgfpathlineto{\pgfqpoint{0.922354in}{0.879395in}}%
\pgfpathlineto{\pgfqpoint{0.926659in}{0.872500in}}%
\pgfpathlineto{\pgfqpoint{0.888413in}{0.876686in}}%
\pgfpathlineto{\pgfqpoint{0.851214in}{0.881898in}}%
\pgfpathclose%
\pgfusepath{fill}%
\end{pgfscope}%
\begin{pgfscope}%
\pgfpathrectangle{\pgfqpoint{0.407592in}{0.222904in}}{\pgfqpoint{1.411163in}{1.411163in}}%
\pgfusepath{clip}%
\pgfsetbuttcap%
\pgfsetroundjoin%
\definecolor{currentfill}{rgb}{0.749020,0.772549,0.792157}%
\pgfsetfillcolor{currentfill}%
\pgfsetfillopacity{0.160000}%
\pgfsetlinewidth{0.000000pt}%
\definecolor{currentstroke}{rgb}{0.000000,0.000000,0.000000}%
\pgfsetstrokecolor{currentstroke}%
\pgfsetdash{}{0pt}%
\pgfpathmoveto{\pgfqpoint{0.677470in}{0.943837in}}%
\pgfpathlineto{\pgfqpoint{0.667588in}{0.951817in}}%
\pgfpathlineto{\pgfqpoint{0.657663in}{0.959832in}}%
\pgfpathlineto{\pgfqpoint{0.693604in}{0.949639in}}%
\pgfpathlineto{\pgfqpoint{0.731620in}{0.940422in}}%
\pgfpathlineto{\pgfqpoint{0.739897in}{0.932809in}}%
\pgfpathlineto{\pgfqpoint{0.748135in}{0.925230in}}%
\pgfpathlineto{\pgfqpoint{0.711820in}{0.934065in}}%
\pgfpathlineto{\pgfqpoint{0.677470in}{0.943837in}}%
\pgfpathclose%
\pgfusepath{fill}%
\end{pgfscope}%
\begin{pgfscope}%
\pgfpathrectangle{\pgfqpoint{0.407592in}{0.222904in}}{\pgfqpoint{1.411163in}{1.411163in}}%
\pgfusepath{clip}%
\pgfsetbuttcap%
\pgfsetroundjoin%
\definecolor{currentfill}{rgb}{0.749020,0.772549,0.792157}%
\pgfsetfillcolor{currentfill}%
\pgfsetfillopacity{0.160000}%
\pgfsetlinewidth{0.000000pt}%
\definecolor{currentstroke}{rgb}{0.000000,0.000000,0.000000}%
\pgfsetstrokecolor{currentstroke}%
\pgfsetdash{}{0pt}%
\pgfpathmoveto{\pgfqpoint{1.166568in}{0.869908in}}%
\pgfpathlineto{\pgfqpoint{1.168579in}{0.876741in}}%
\pgfpathlineto{\pgfqpoint{1.170600in}{0.883608in}}%
\pgfpathlineto{\pgfqpoint{1.212419in}{0.887118in}}%
\pgfpathlineto{\pgfqpoint{1.253556in}{0.891724in}}%
\pgfpathlineto{\pgfqpoint{1.249459in}{0.884670in}}%
\pgfpathlineto{\pgfqpoint{1.245381in}{0.877651in}}%
\pgfpathlineto{\pgfqpoint{1.206296in}{0.873256in}}%
\pgfpathlineto{\pgfqpoint{1.166568in}{0.869908in}}%
\pgfpathclose%
\pgfusepath{fill}%
\end{pgfscope}%
\begin{pgfscope}%
\pgfpathrectangle{\pgfqpoint{0.407592in}{0.222904in}}{\pgfqpoint{1.411163in}{1.411163in}}%
\pgfusepath{clip}%
\pgfsetbuttcap%
\pgfsetroundjoin%
\definecolor{currentfill}{rgb}{0.749020,0.772549,0.792157}%
\pgfsetfillcolor{currentfill}%
\pgfsetfillopacity{0.160000}%
\pgfsetlinewidth{0.000000pt}%
\definecolor{currentstroke}{rgb}{0.000000,0.000000,0.000000}%
\pgfsetstrokecolor{currentstroke}%
\pgfsetdash{}{0pt}%
\pgfpathmoveto{\pgfqpoint{0.764500in}{0.910176in}}%
\pgfpathlineto{\pgfqpoint{0.756336in}{0.917686in}}%
\pgfpathlineto{\pgfqpoint{0.748135in}{0.925230in}}%
\pgfpathlineto{\pgfqpoint{0.786240in}{0.917389in}}%
\pgfpathlineto{\pgfqpoint{0.825941in}{0.910590in}}%
\pgfpathlineto{\pgfqpoint{0.832305in}{0.903366in}}%
\pgfpathlineto{\pgfqpoint{0.838638in}{0.896175in}}%
\pgfpathlineto{\pgfqpoint{0.800814in}{0.902677in}}%
\pgfpathlineto{\pgfqpoint{0.764500in}{0.910176in}}%
\pgfpathclose%
\pgfusepath{fill}%
\end{pgfscope}%
\begin{pgfscope}%
\pgfpathrectangle{\pgfqpoint{0.407592in}{0.222904in}}{\pgfqpoint{1.411163in}{1.411163in}}%
\pgfusepath{clip}%
\pgfsetbuttcap%
\pgfsetroundjoin%
\definecolor{currentfill}{rgb}{0.749020,0.772549,0.792157}%
\pgfsetfillcolor{currentfill}%
\pgfsetfillopacity{0.160000}%
\pgfsetlinewidth{0.000000pt}%
\definecolor{currentstroke}{rgb}{0.000000,0.000000,0.000000}%
\pgfsetstrokecolor{currentstroke}%
\pgfsetdash{}{0pt}%
\pgfpathmoveto{\pgfqpoint{1.438914in}{0.951607in}}%
\pgfpathlineto{\pgfqpoint{1.446970in}{0.959487in}}%
\pgfpathlineto{\pgfqpoint{1.455062in}{0.967403in}}%
\pgfpathlineto{\pgfqpoint{1.494766in}{0.978210in}}%
\pgfpathlineto{\pgfqpoint{1.532296in}{0.989951in}}%
\pgfpathlineto{\pgfqpoint{1.522558in}{0.981593in}}%
\pgfpathlineto{\pgfqpoint{1.512861in}{0.973270in}}%
\pgfpathlineto{\pgfqpoint{1.476919in}{0.961989in}}%
\pgfpathlineto{\pgfqpoint{1.438914in}{0.951607in}}%
\pgfpathclose%
\pgfusepath{fill}%
\end{pgfscope}%
\begin{pgfscope}%
\pgfpathrectangle{\pgfqpoint{0.407592in}{0.222904in}}{\pgfqpoint{1.411163in}{1.411163in}}%
\pgfusepath{clip}%
\pgfsetbuttcap%
\pgfsetroundjoin%
\definecolor{currentfill}{rgb}{0.749020,0.772549,0.792157}%
\pgfsetfillcolor{currentfill}%
\pgfsetfillopacity{0.160000}%
\pgfsetlinewidth{0.000000pt}%
\definecolor{currentstroke}{rgb}{0.000000,0.000000,0.000000}%
\pgfsetstrokecolor{currentstroke}%
\pgfsetdash{}{0pt}%
\pgfpathmoveto{\pgfqpoint{1.551892in}{1.006770in}}%
\pgfpathlineto{\pgfqpoint{1.561752in}{1.015233in}}%
\pgfpathlineto{\pgfqpoint{1.571653in}{1.023730in}}%
\pgfpathlineto{\pgfqpoint{1.609784in}{1.037305in}}%
\pgfpathlineto{\pgfqpoint{1.645152in}{1.051715in}}%
\pgfpathlineto{\pgfqpoint{1.633838in}{1.042720in}}%
\pgfpathlineto{\pgfqpoint{1.622568in}{1.033759in}}%
\pgfpathlineto{\pgfqpoint{1.588546in}{1.019859in}}%
\pgfpathlineto{\pgfqpoint{1.551892in}{1.006770in}}%
\pgfpathclose%
\pgfusepath{fill}%
\end{pgfscope}%
\begin{pgfscope}%
\pgfpathrectangle{\pgfqpoint{0.407592in}{0.222904in}}{\pgfqpoint{1.411163in}{1.411163in}}%
\pgfusepath{clip}%
\pgfsetbuttcap%
\pgfsetroundjoin%
\definecolor{currentfill}{rgb}{0.749020,0.772549,0.792157}%
\pgfsetfillcolor{currentfill}%
\pgfsetfillopacity{0.160000}%
\pgfsetlinewidth{0.000000pt}%
\definecolor{currentstroke}{rgb}{0.000000,0.000000,0.000000}%
\pgfsetstrokecolor{currentstroke}%
\pgfsetdash{}{0pt}%
\pgfpathmoveto{\pgfqpoint{0.461940in}{1.063777in}}%
\pgfpathlineto{\pgfqpoint{0.449132in}{1.073016in}}%
\pgfpathlineto{\pgfqpoint{0.436278in}{1.082289in}}%
\pgfpathlineto{\pgfqpoint{0.466696in}{1.066750in}}%
\pgfpathlineto{\pgfqpoint{0.500331in}{1.051998in}}%
\pgfpathlineto{\pgfqpoint{0.512038in}{1.043237in}}%
\pgfpathlineto{\pgfqpoint{0.523697in}{1.034512in}}%
\pgfpathlineto{\pgfqpoint{0.491283in}{1.048761in}}%
\pgfpathlineto{\pgfqpoint{0.461940in}{1.063777in}}%
\pgfpathclose%
\pgfusepath{fill}%
\end{pgfscope}%
\begin{pgfscope}%
\pgfpathrectangle{\pgfqpoint{0.407592in}{0.222904in}}{\pgfqpoint{1.411163in}{1.411163in}}%
\pgfusepath{clip}%
\pgfsetbuttcap%
\pgfsetroundjoin%
\definecolor{currentfill}{rgb}{0.749020,0.772549,0.792157}%
\pgfsetfillcolor{currentfill}%
\pgfsetfillopacity{0.160000}%
\pgfsetlinewidth{0.000000pt}%
\definecolor{currentstroke}{rgb}{0.000000,0.000000,0.000000}%
\pgfsetstrokecolor{currentstroke}%
\pgfsetdash{}{0pt}%
\pgfpathmoveto{\pgfqpoint{0.926659in}{0.872500in}}%
\pgfpathlineto{\pgfqpoint{0.922354in}{0.879395in}}%
\pgfpathlineto{\pgfqpoint{0.918027in}{0.886325in}}%
\pgfpathlineto{\pgfqpoint{0.959177in}{0.883038in}}%
\pgfpathlineto{\pgfqpoint{1.001008in}{0.880873in}}%
\pgfpathlineto{\pgfqpoint{1.003254in}{0.874070in}}%
\pgfpathlineto{\pgfqpoint{1.005488in}{0.867300in}}%
\pgfpathlineto{\pgfqpoint{0.965753in}{0.869365in}}%
\pgfpathlineto{\pgfqpoint{0.926659in}{0.872500in}}%
\pgfpathclose%
\pgfusepath{fill}%
\end{pgfscope}%
\begin{pgfscope}%
\pgfpathrectangle{\pgfqpoint{0.407592in}{0.222904in}}{\pgfqpoint{1.411163in}{1.411163in}}%
\pgfusepath{clip}%
\pgfsetbuttcap%
\pgfsetroundjoin%
\definecolor{currentfill}{rgb}{0.749020,0.772549,0.792157}%
\pgfsetfillcolor{currentfill}%
\pgfsetfillopacity{0.160000}%
\pgfsetlinewidth{0.000000pt}%
\definecolor{currentstroke}{rgb}{0.000000,0.000000,0.000000}%
\pgfsetstrokecolor{currentstroke}%
\pgfsetdash{}{0pt}%
\pgfpathmoveto{\pgfqpoint{0.569868in}{0.999959in}}%
\pgfpathlineto{\pgfqpoint{0.558395in}{1.008545in}}%
\pgfpathlineto{\pgfqpoint{0.546875in}{1.017166in}}%
\pgfpathlineto{\pgfqpoint{0.580874in}{1.004249in}}%
\pgfpathlineto{\pgfqpoint{0.617531in}{0.992238in}}%
\pgfpathlineto{\pgfqpoint{0.627630in}{0.984084in}}%
\pgfpathlineto{\pgfqpoint{0.637684in}{0.975965in}}%
\pgfpathlineto{\pgfqpoint{0.602512in}{0.987524in}}%
\pgfpathlineto{\pgfqpoint{0.569868in}{0.999959in}}%
\pgfpathclose%
\pgfusepath{fill}%
\end{pgfscope}%
\begin{pgfscope}%
\pgfpathrectangle{\pgfqpoint{0.407592in}{0.222904in}}{\pgfqpoint{1.411163in}{1.411163in}}%
\pgfusepath{clip}%
\pgfsetbuttcap%
\pgfsetroundjoin%
\definecolor{currentfill}{rgb}{0.749020,0.772549,0.792157}%
\pgfsetfillcolor{currentfill}%
\pgfsetfillopacity{0.160000}%
\pgfsetlinewidth{0.000000pt}%
\definecolor{currentstroke}{rgb}{0.000000,0.000000,0.000000}%
\pgfsetstrokecolor{currentstroke}%
\pgfsetdash{}{0pt}%
\pgfpathmoveto{\pgfqpoint{1.253556in}{0.891724in}}%
\pgfpathlineto{\pgfqpoint{1.257673in}{0.898812in}}%
\pgfpathlineto{\pgfqpoint{1.261810in}{0.905934in}}%
\pgfpathlineto{\pgfqpoint{1.304060in}{0.911869in}}%
\pgfpathlineto{\pgfqpoint{1.345134in}{0.918880in}}%
\pgfpathlineto{\pgfqpoint{1.339011in}{0.911474in}}%
\pgfpathlineto{\pgfqpoint{1.332917in}{0.904103in}}%
\pgfpathlineto{\pgfqpoint{1.293792in}{0.897398in}}%
\pgfpathlineto{\pgfqpoint{1.253556in}{0.891724in}}%
\pgfpathclose%
\pgfusepath{fill}%
\end{pgfscope}%
\begin{pgfscope}%
\pgfpathrectangle{\pgfqpoint{0.407592in}{0.222904in}}{\pgfqpoint{1.411163in}{1.411163in}}%
\pgfusepath{clip}%
\pgfsetbuttcap%
\pgfsetroundjoin%
\definecolor{currentfill}{rgb}{0.749020,0.772549,0.792157}%
\pgfsetfillcolor{currentfill}%
\pgfsetfillopacity{0.160000}%
\pgfsetlinewidth{0.000000pt}%
\definecolor{currentstroke}{rgb}{0.000000,0.000000,0.000000}%
\pgfsetstrokecolor{currentstroke}%
\pgfsetdash{}{0pt}%
\pgfpathmoveto{\pgfqpoint{1.086031in}{0.866429in}}%
\pgfpathlineto{\pgfqpoint{1.085920in}{0.873177in}}%
\pgfpathlineto{\pgfqpoint{1.085809in}{0.879959in}}%
\pgfpathlineto{\pgfqpoint{1.128320in}{0.881216in}}%
\pgfpathlineto{\pgfqpoint{1.170600in}{0.883608in}}%
\pgfpathlineto{\pgfqpoint{1.168579in}{0.876741in}}%
\pgfpathlineto{\pgfqpoint{1.166568in}{0.869908in}}%
\pgfpathlineto{\pgfqpoint{1.126408in}{0.867628in}}%
\pgfpathlineto{\pgfqpoint{1.086031in}{0.866429in}}%
\pgfpathclose%
\pgfusepath{fill}%
\end{pgfscope}%
\begin{pgfscope}%
\pgfpathrectangle{\pgfqpoint{0.407592in}{0.222904in}}{\pgfqpoint{1.411163in}{1.411163in}}%
\pgfusepath{clip}%
\pgfsetbuttcap%
\pgfsetroundjoin%
\definecolor{currentfill}{rgb}{0.749020,0.772549,0.792157}%
\pgfsetfillcolor{currentfill}%
\pgfsetfillopacity{0.160000}%
\pgfsetlinewidth{0.000000pt}%
\definecolor{currentstroke}{rgb}{0.000000,0.000000,0.000000}%
\pgfsetstrokecolor{currentstroke}%
\pgfsetdash{}{0pt}%
\pgfpathmoveto{\pgfqpoint{1.345134in}{0.918880in}}%
\pgfpathlineto{\pgfqpoint{1.351285in}{0.926321in}}%
\pgfpathlineto{\pgfqpoint{1.357464in}{0.933796in}}%
\pgfpathlineto{\pgfqpoint{1.399030in}{0.942190in}}%
\pgfpathlineto{\pgfqpoint{1.438914in}{0.951607in}}%
\pgfpathlineto{\pgfqpoint{1.430894in}{0.943762in}}%
\pgfpathlineto{\pgfqpoint{1.422909in}{0.935952in}}%
\pgfpathlineto{\pgfqpoint{1.384817in}{0.926925in}}%
\pgfpathlineto{\pgfqpoint{1.345134in}{0.918880in}}%
\pgfpathclose%
\pgfusepath{fill}%
\end{pgfscope}%
\begin{pgfscope}%
\pgfpathrectangle{\pgfqpoint{0.407592in}{0.222904in}}{\pgfqpoint{1.411163in}{1.411163in}}%
\pgfusepath{clip}%
\pgfsetbuttcap%
\pgfsetroundjoin%
\definecolor{currentfill}{rgb}{0.749020,0.772549,0.792157}%
\pgfsetfillcolor{currentfill}%
\pgfsetfillopacity{0.160000}%
\pgfsetlinewidth{0.000000pt}%
\definecolor{currentstroke}{rgb}{0.000000,0.000000,0.000000}%
\pgfsetstrokecolor{currentstroke}%
\pgfsetdash{}{0pt}%
\pgfpathmoveto{\pgfqpoint{1.005488in}{0.867300in}}%
\pgfpathlineto{\pgfqpoint{1.003254in}{0.874070in}}%
\pgfpathlineto{\pgfqpoint{1.001008in}{0.880873in}}%
\pgfpathlineto{\pgfqpoint{1.043295in}{0.879845in}}%
\pgfpathlineto{\pgfqpoint{1.085809in}{0.879959in}}%
\pgfpathlineto{\pgfqpoint{1.085920in}{0.873177in}}%
\pgfpathlineto{\pgfqpoint{1.086031in}{0.866429in}}%
\pgfpathlineto{\pgfqpoint{1.045652in}{0.866319in}}%
\pgfpathlineto{\pgfqpoint{1.005488in}{0.867300in}}%
\pgfpathclose%
\pgfusepath{fill}%
\end{pgfscope}%
\begin{pgfscope}%
\pgfpathrectangle{\pgfqpoint{0.407592in}{0.222904in}}{\pgfqpoint{1.411163in}{1.411163in}}%
\pgfusepath{clip}%
\pgfsetbuttcap%
\pgfsetroundjoin%
\definecolor{currentfill}{rgb}{0.749020,0.772549,0.792157}%
\pgfsetfillcolor{currentfill}%
\pgfsetfillopacity{0.160000}%
\pgfsetlinewidth{0.000000pt}%
\definecolor{currentstroke}{rgb}{0.000000,0.000000,0.000000}%
\pgfsetstrokecolor{currentstroke}%
\pgfsetdash{}{0pt}%
\pgfpathmoveto{\pgfqpoint{0.838638in}{0.896175in}}%
\pgfpathlineto{\pgfqpoint{0.832305in}{0.903366in}}%
\pgfpathlineto{\pgfqpoint{0.825941in}{0.910590in}}%
\pgfpathlineto{\pgfqpoint{0.867036in}{0.904877in}}%
\pgfpathlineto{\pgfqpoint{0.909310in}{0.900286in}}%
\pgfpathlineto{\pgfqpoint{0.913679in}{0.893288in}}%
\pgfpathlineto{\pgfqpoint{0.918027in}{0.886325in}}%
\pgfpathlineto{\pgfqpoint{0.877776in}{0.890713in}}%
\pgfpathlineto{\pgfqpoint{0.838638in}{0.896175in}}%
\pgfpathclose%
\pgfusepath{fill}%
\end{pgfscope}%
\begin{pgfscope}%
\pgfpathrectangle{\pgfqpoint{0.407592in}{0.222904in}}{\pgfqpoint{1.411163in}{1.411163in}}%
\pgfusepath{clip}%
\pgfsetbuttcap%
\pgfsetroundjoin%
\definecolor{currentfill}{rgb}{0.749020,0.772549,0.792157}%
\pgfsetfillcolor{currentfill}%
\pgfsetfillopacity{0.160000}%
\pgfsetlinewidth{0.000000pt}%
\definecolor{currentstroke}{rgb}{0.000000,0.000000,0.000000}%
\pgfsetstrokecolor{currentstroke}%
\pgfsetdash{}{0pt}%
\pgfpathmoveto{\pgfqpoint{0.436278in}{1.082289in}}%
\pgfpathlineto{\pgfqpoint{0.423376in}{1.091597in}}%
\pgfpathlineto{\pgfqpoint{0.454331in}{1.075797in}}%
\pgfpathlineto{\pgfqpoint{0.488577in}{1.060795in}}%
\pgfpathlineto{\pgfqpoint{0.500331in}{1.051998in}}%
\pgfpathlineto{\pgfqpoint{0.466696in}{1.066750in}}%
\pgfusepath{fill}%
\end{pgfscope}%
\begin{pgfscope}%
\pgfpathrectangle{\pgfqpoint{0.407592in}{0.222904in}}{\pgfqpoint{1.411163in}{1.411163in}}%
\pgfusepath{clip}%
\pgfsetbuttcap%
\pgfsetroundjoin%
\definecolor{currentfill}{rgb}{0.749020,0.772549,0.792157}%
\pgfsetfillcolor{currentfill}%
\pgfsetfillopacity{0.160000}%
\pgfsetlinewidth{0.000000pt}%
\definecolor{currentstroke}{rgb}{0.000000,0.000000,0.000000}%
\pgfsetstrokecolor{currentstroke}%
\pgfsetdash{}{0pt}%
\pgfpathmoveto{\pgfqpoint{0.657663in}{0.959832in}}%
\pgfpathlineto{\pgfqpoint{0.647695in}{0.967881in}}%
\pgfpathlineto{\pgfqpoint{0.637684in}{0.975965in}}%
\pgfpathlineto{\pgfqpoint{0.675224in}{0.965353in}}%
\pgfpathlineto{\pgfqpoint{0.714953in}{0.955755in}}%
\pgfpathlineto{\pgfqpoint{0.723306in}{0.948071in}}%
\pgfpathlineto{\pgfqpoint{0.731620in}{0.940422in}}%
\pgfpathlineto{\pgfqpoint{0.693604in}{0.949639in}}%
\pgfpathlineto{\pgfqpoint{0.657663in}{0.959832in}}%
\pgfpathclose%
\pgfusepath{fill}%
\end{pgfscope}%
\begin{pgfscope}%
\pgfpathrectangle{\pgfqpoint{0.407592in}{0.222904in}}{\pgfqpoint{1.411163in}{1.411163in}}%
\pgfusepath{clip}%
\pgfsetbuttcap%
\pgfsetroundjoin%
\definecolor{currentfill}{rgb}{0.749020,0.772549,0.792157}%
\pgfsetfillcolor{currentfill}%
\pgfsetfillopacity{0.160000}%
\pgfsetlinewidth{0.000000pt}%
\definecolor{currentstroke}{rgb}{0.000000,0.000000,0.000000}%
\pgfsetstrokecolor{currentstroke}%
\pgfsetdash{}{0pt}%
\pgfpathmoveto{\pgfqpoint{1.455062in}{0.967403in}}%
\pgfpathlineto{\pgfqpoint{1.463190in}{0.975353in}}%
\pgfpathlineto{\pgfqpoint{1.471354in}{0.983339in}}%
\pgfpathlineto{\pgfqpoint{1.512768in}{0.994572in}}%
\pgfpathlineto{\pgfqpoint{1.551892in}{1.006770in}}%
\pgfpathlineto{\pgfqpoint{1.542074in}{0.998343in}}%
\pgfpathlineto{\pgfqpoint{1.532296in}{0.989951in}}%
\pgfpathlineto{\pgfqpoint{1.494766in}{0.978210in}}%
\pgfpathlineto{\pgfqpoint{1.455062in}{0.967403in}}%
\pgfpathclose%
\pgfusepath{fill}%
\end{pgfscope}%
\begin{pgfscope}%
\pgfpathrectangle{\pgfqpoint{0.407592in}{0.222904in}}{\pgfqpoint{1.411163in}{1.411163in}}%
\pgfusepath{clip}%
\pgfsetbuttcap%
\pgfsetroundjoin%
\definecolor{currentfill}{rgb}{0.749020,0.772549,0.792157}%
\pgfsetfillcolor{currentfill}%
\pgfsetfillopacity{0.160000}%
\pgfsetlinewidth{0.000000pt}%
\definecolor{currentstroke}{rgb}{0.000000,0.000000,0.000000}%
\pgfsetstrokecolor{currentstroke}%
\pgfsetdash{}{0pt}%
\pgfpathmoveto{\pgfqpoint{1.571653in}{1.023730in}}%
\pgfpathlineto{\pgfqpoint{1.581596in}{1.032264in}}%
\pgfpathlineto{\pgfqpoint{1.591580in}{1.040833in}}%
\pgfpathlineto{\pgfqpoint{1.631193in}{1.054890in}}%
\pgfpathlineto{\pgfqpoint{1.667909in}{1.069807in}}%
\pgfpathlineto{\pgfqpoint{1.656509in}{1.060744in}}%
\pgfpathlineto{\pgfqpoint{1.645152in}{1.051715in}}%
\pgfpathlineto{\pgfqpoint{1.609784in}{1.037305in}}%
\pgfpathlineto{\pgfqpoint{1.571653in}{1.023730in}}%
\pgfpathclose%
\pgfusepath{fill}%
\end{pgfscope}%
\begin{pgfscope}%
\pgfpathrectangle{\pgfqpoint{0.407592in}{0.222904in}}{\pgfqpoint{1.411163in}{1.411163in}}%
\pgfusepath{clip}%
\pgfsetbuttcap%
\pgfsetroundjoin%
\definecolor{currentfill}{rgb}{0.749020,0.772549,0.792157}%
\pgfsetfillcolor{currentfill}%
\pgfsetfillopacity{0.160000}%
\pgfsetlinewidth{0.000000pt}%
\definecolor{currentstroke}{rgb}{0.000000,0.000000,0.000000}%
\pgfsetstrokecolor{currentstroke}%
\pgfsetdash{}{0pt}%
\pgfpathmoveto{\pgfqpoint{0.748135in}{0.925230in}}%
\pgfpathlineto{\pgfqpoint{0.739897in}{0.932809in}}%
\pgfpathlineto{\pgfqpoint{0.731620in}{0.940422in}}%
\pgfpathlineto{\pgfqpoint{0.771527in}{0.932240in}}%
\pgfpathlineto{\pgfqpoint{0.813123in}{0.925144in}}%
\pgfpathlineto{\pgfqpoint{0.819547in}{0.917850in}}%
\pgfpathlineto{\pgfqpoint{0.825941in}{0.910590in}}%
\pgfpathlineto{\pgfqpoint{0.786240in}{0.917389in}}%
\pgfpathlineto{\pgfqpoint{0.748135in}{0.925230in}}%
\pgfpathclose%
\pgfusepath{fill}%
\end{pgfscope}%
\begin{pgfscope}%
\pgfpathrectangle{\pgfqpoint{0.407592in}{0.222904in}}{\pgfqpoint{1.411163in}{1.411163in}}%
\pgfusepath{clip}%
\pgfsetbuttcap%
\pgfsetroundjoin%
\definecolor{currentfill}{rgb}{0.749020,0.772549,0.792157}%
\pgfsetfillcolor{currentfill}%
\pgfsetfillopacity{0.160000}%
\pgfsetlinewidth{0.000000pt}%
\definecolor{currentstroke}{rgb}{0.000000,0.000000,0.000000}%
\pgfsetstrokecolor{currentstroke}%
\pgfsetdash{}{0pt}%
\pgfpathmoveto{\pgfqpoint{1.170600in}{0.883608in}}%
\pgfpathlineto{\pgfqpoint{1.172631in}{0.890508in}}%
\pgfpathlineto{\pgfqpoint{1.174672in}{0.897443in}}%
\pgfpathlineto{\pgfqpoint{1.218603in}{0.901116in}}%
\pgfpathlineto{\pgfqpoint{1.261810in}{0.905934in}}%
\pgfpathlineto{\pgfqpoint{1.257673in}{0.898812in}}%
\pgfpathlineto{\pgfqpoint{1.253556in}{0.891724in}}%
\pgfpathlineto{\pgfqpoint{1.212419in}{0.887118in}}%
\pgfpathlineto{\pgfqpoint{1.170600in}{0.883608in}}%
\pgfpathclose%
\pgfusepath{fill}%
\end{pgfscope}%
\begin{pgfscope}%
\pgfpathrectangle{\pgfqpoint{0.407592in}{0.222904in}}{\pgfqpoint{1.411163in}{1.411163in}}%
\pgfusepath{clip}%
\pgfsetbuttcap%
\pgfsetroundjoin%
\definecolor{currentfill}{rgb}{0.749020,0.772549,0.792157}%
\pgfsetfillcolor{currentfill}%
\pgfsetfillopacity{0.160000}%
\pgfsetlinewidth{0.000000pt}%
\definecolor{currentstroke}{rgb}{0.000000,0.000000,0.000000}%
\pgfsetstrokecolor{currentstroke}%
\pgfsetdash{}{0pt}%
\pgfpathmoveto{\pgfqpoint{0.546875in}{1.017166in}}%
\pgfpathlineto{\pgfqpoint{0.535310in}{1.025821in}}%
\pgfpathlineto{\pgfqpoint{0.523697in}{1.034512in}}%
\pgfpathlineto{\pgfqpoint{0.559054in}{1.021115in}}%
\pgfpathlineto{\pgfqpoint{0.597202in}{1.008654in}}%
\pgfpathlineto{\pgfqpoint{0.607389in}{1.000428in}}%
\pgfpathlineto{\pgfqpoint{0.617531in}{0.992238in}}%
\pgfpathlineto{\pgfqpoint{0.580874in}{1.004249in}}%
\pgfpathlineto{\pgfqpoint{0.546875in}{1.017166in}}%
\pgfpathclose%
\pgfusepath{fill}%
\end{pgfscope}%
\begin{pgfscope}%
\pgfpathrectangle{\pgfqpoint{0.407592in}{0.222904in}}{\pgfqpoint{1.411163in}{1.411163in}}%
\pgfusepath{clip}%
\pgfsetbuttcap%
\pgfsetroundjoin%
\definecolor{currentfill}{rgb}{0.749020,0.772549,0.792157}%
\pgfsetfillcolor{currentfill}%
\pgfsetfillopacity{0.160000}%
\pgfsetlinewidth{0.000000pt}%
\definecolor{currentstroke}{rgb}{0.000000,0.000000,0.000000}%
\pgfsetstrokecolor{currentstroke}%
\pgfsetdash{}{0pt}%
\pgfpathmoveto{\pgfqpoint{0.918027in}{0.886325in}}%
\pgfpathlineto{\pgfqpoint{0.913679in}{0.893288in}}%
\pgfpathlineto{\pgfqpoint{0.909310in}{0.900286in}}%
\pgfpathlineto{\pgfqpoint{0.952536in}{0.896847in}}%
\pgfpathlineto{\pgfqpoint{0.996484in}{0.894582in}}%
\pgfpathlineto{\pgfqpoint{0.998752in}{0.887710in}}%
\pgfpathlineto{\pgfqpoint{1.001008in}{0.880873in}}%
\pgfpathlineto{\pgfqpoint{0.959177in}{0.883038in}}%
\pgfpathlineto{\pgfqpoint{0.918027in}{0.886325in}}%
\pgfpathclose%
\pgfusepath{fill}%
\end{pgfscope}%
\begin{pgfscope}%
\pgfpathrectangle{\pgfqpoint{0.407592in}{0.222904in}}{\pgfqpoint{1.411163in}{1.411163in}}%
\pgfusepath{clip}%
\pgfsetbuttcap%
\pgfsetroundjoin%
\definecolor{currentfill}{rgb}{0.749020,0.772549,0.792157}%
\pgfsetfillcolor{currentfill}%
\pgfsetfillopacity{0.160000}%
\pgfsetlinewidth{0.000000pt}%
\definecolor{currentstroke}{rgb}{0.000000,0.000000,0.000000}%
\pgfsetstrokecolor{currentstroke}%
\pgfsetdash{}{0pt}%
\pgfpathmoveto{\pgfqpoint{1.357464in}{0.933796in}}%
\pgfpathlineto{\pgfqpoint{1.363673in}{0.941307in}}%
\pgfpathlineto{\pgfqpoint{1.369911in}{0.948853in}}%
\pgfpathlineto{\pgfqpoint{1.413375in}{0.957596in}}%
\pgfpathlineto{\pgfqpoint{1.455062in}{0.967403in}}%
\pgfpathlineto{\pgfqpoint{1.446970in}{0.959487in}}%
\pgfpathlineto{\pgfqpoint{1.438914in}{0.951607in}}%
\pgfpathlineto{\pgfqpoint{1.399030in}{0.942190in}}%
\pgfpathlineto{\pgfqpoint{1.357464in}{0.933796in}}%
\pgfpathclose%
\pgfusepath{fill}%
\end{pgfscope}%
\begin{pgfscope}%
\pgfpathrectangle{\pgfqpoint{0.407592in}{0.222904in}}{\pgfqpoint{1.411163in}{1.411163in}}%
\pgfusepath{clip}%
\pgfsetbuttcap%
\pgfsetroundjoin%
\definecolor{currentfill}{rgb}{0.749020,0.772549,0.792157}%
\pgfsetfillcolor{currentfill}%
\pgfsetfillopacity{0.160000}%
\pgfsetlinewidth{0.000000pt}%
\definecolor{currentstroke}{rgb}{0.000000,0.000000,0.000000}%
\pgfsetstrokecolor{currentstroke}%
\pgfsetdash{}{0pt}%
\pgfpathmoveto{\pgfqpoint{1.261810in}{0.905934in}}%
\pgfpathlineto{\pgfqpoint{1.265967in}{0.913092in}}%
\pgfpathlineto{\pgfqpoint{1.270144in}{0.920284in}}%
\pgfpathlineto{\pgfqpoint{1.314428in}{0.926478in}}%
\pgfpathlineto{\pgfqpoint{1.357464in}{0.933796in}}%
\pgfpathlineto{\pgfqpoint{1.351285in}{0.926321in}}%
\pgfpathlineto{\pgfqpoint{1.345134in}{0.918880in}}%
\pgfpathlineto{\pgfqpoint{1.304060in}{0.911869in}}%
\pgfpathlineto{\pgfqpoint{1.261810in}{0.905934in}}%
\pgfpathclose%
\pgfusepath{fill}%
\end{pgfscope}%
\begin{pgfscope}%
\pgfpathrectangle{\pgfqpoint{0.407592in}{0.222904in}}{\pgfqpoint{1.411163in}{1.411163in}}%
\pgfusepath{clip}%
\pgfsetbuttcap%
\pgfsetroundjoin%
\definecolor{currentfill}{rgb}{0.749020,0.772549,0.792157}%
\pgfsetfillcolor{currentfill}%
\pgfsetfillopacity{0.160000}%
\pgfsetlinewidth{0.000000pt}%
\definecolor{currentstroke}{rgb}{0.000000,0.000000,0.000000}%
\pgfsetstrokecolor{currentstroke}%
\pgfsetdash{}{0pt}%
\pgfpathmoveto{\pgfqpoint{1.085809in}{0.879959in}}%
\pgfpathlineto{\pgfqpoint{1.085697in}{0.886776in}}%
\pgfpathlineto{\pgfqpoint{1.085585in}{0.893626in}}%
\pgfpathlineto{\pgfqpoint{1.130252in}{0.894942in}}%
\pgfpathlineto{\pgfqpoint{1.174672in}{0.897443in}}%
\pgfpathlineto{\pgfqpoint{1.172631in}{0.890508in}}%
\pgfpathlineto{\pgfqpoint{1.170600in}{0.883608in}}%
\pgfpathlineto{\pgfqpoint{1.128320in}{0.881216in}}%
\pgfpathlineto{\pgfqpoint{1.085809in}{0.879959in}}%
\pgfpathclose%
\pgfusepath{fill}%
\end{pgfscope}%
\begin{pgfscope}%
\pgfpathrectangle{\pgfqpoint{0.407592in}{0.222904in}}{\pgfqpoint{1.411163in}{1.411163in}}%
\pgfusepath{clip}%
\pgfsetbuttcap%
\pgfsetroundjoin%
\definecolor{currentfill}{rgb}{0.749020,0.772549,0.792157}%
\pgfsetfillcolor{currentfill}%
\pgfsetfillopacity{0.160000}%
\pgfsetlinewidth{0.000000pt}%
\definecolor{currentstroke}{rgb}{0.000000,0.000000,0.000000}%
\pgfsetstrokecolor{currentstroke}%
\pgfsetdash{}{0pt}%
\pgfpathmoveto{\pgfqpoint{1.001008in}{0.880873in}}%
\pgfpathlineto{\pgfqpoint{0.998752in}{0.887710in}}%
\pgfpathlineto{\pgfqpoint{0.996484in}{0.894582in}}%
\pgfpathlineto{\pgfqpoint{1.040914in}{0.893506in}}%
\pgfpathlineto{\pgfqpoint{1.085585in}{0.893626in}}%
\pgfpathlineto{\pgfqpoint{1.085697in}{0.886776in}}%
\pgfpathlineto{\pgfqpoint{1.085809in}{0.879959in}}%
\pgfpathlineto{\pgfqpoint{1.043295in}{0.879845in}}%
\pgfpathlineto{\pgfqpoint{1.001008in}{0.880873in}}%
\pgfpathclose%
\pgfusepath{fill}%
\end{pgfscope}%
\begin{pgfscope}%
\pgfpathrectangle{\pgfqpoint{0.407592in}{0.222904in}}{\pgfqpoint{1.411163in}{1.411163in}}%
\pgfusepath{clip}%
\pgfsetbuttcap%
\pgfsetroundjoin%
\definecolor{currentfill}{rgb}{0.749020,0.772549,0.792157}%
\pgfsetfillcolor{currentfill}%
\pgfsetfillopacity{0.160000}%
\pgfsetlinewidth{0.000000pt}%
\definecolor{currentstroke}{rgb}{0.000000,0.000000,0.000000}%
\pgfsetstrokecolor{currentstroke}%
\pgfsetdash{}{0pt}%
\pgfpathmoveto{\pgfqpoint{1.591580in}{1.040833in}}%
\pgfpathlineto{\pgfqpoint{1.601606in}{1.049438in}}%
\pgfpathlineto{\pgfqpoint{1.641962in}{1.063736in}}%
\pgfpathlineto{\pgfqpoint{1.679353in}{1.078906in}}%
\pgfpathlineto{\pgfqpoint{1.667909in}{1.069807in}}%
\pgfpathlineto{\pgfqpoint{1.631193in}{1.054890in}}%
\pgfusepath{fill}%
\end{pgfscope}%
\begin{pgfscope}%
\pgfpathrectangle{\pgfqpoint{0.407592in}{0.222904in}}{\pgfqpoint{1.411163in}{1.411163in}}%
\pgfusepath{clip}%
\pgfsetbuttcap%
\pgfsetroundjoin%
\definecolor{currentfill}{rgb}{0.749020,0.772549,0.792157}%
\pgfsetfillcolor{currentfill}%
\pgfsetfillopacity{0.160000}%
\pgfsetlinewidth{0.000000pt}%
\definecolor{currentstroke}{rgb}{0.000000,0.000000,0.000000}%
\pgfsetstrokecolor{currentstroke}%
\pgfsetdash{}{0pt}%
\pgfpathmoveto{\pgfqpoint{0.637684in}{0.975965in}}%
\pgfpathlineto{\pgfqpoint{0.627630in}{0.984084in}}%
\pgfpathlineto{\pgfqpoint{0.617531in}{0.992238in}}%
\pgfpathlineto{\pgfqpoint{0.656679in}{0.981209in}}%
\pgfpathlineto{\pgfqpoint{0.698131in}{0.971230in}}%
\pgfpathlineto{\pgfqpoint{0.706562in}{0.963474in}}%
\pgfpathlineto{\pgfqpoint{0.714953in}{0.955755in}}%
\pgfpathlineto{\pgfqpoint{0.675224in}{0.965353in}}%
\pgfpathlineto{\pgfqpoint{0.637684in}{0.975965in}}%
\pgfpathclose%
\pgfusepath{fill}%
\end{pgfscope}%
\begin{pgfscope}%
\pgfpathrectangle{\pgfqpoint{0.407592in}{0.222904in}}{\pgfqpoint{1.411163in}{1.411163in}}%
\pgfusepath{clip}%
\pgfsetbuttcap%
\pgfsetroundjoin%
\definecolor{currentfill}{rgb}{0.749020,0.772549,0.792157}%
\pgfsetfillcolor{currentfill}%
\pgfsetfillopacity{0.160000}%
\pgfsetlinewidth{0.000000pt}%
\definecolor{currentstroke}{rgb}{0.000000,0.000000,0.000000}%
\pgfsetstrokecolor{currentstroke}%
\pgfsetdash{}{0pt}%
\pgfpathmoveto{\pgfqpoint{0.825941in}{0.910590in}}%
\pgfpathlineto{\pgfqpoint{0.819547in}{0.917850in}}%
\pgfpathlineto{\pgfqpoint{0.813123in}{0.925144in}}%
\pgfpathlineto{\pgfqpoint{0.856191in}{0.919180in}}%
\pgfpathlineto{\pgfqpoint{0.900505in}{0.914386in}}%
\pgfpathlineto{\pgfqpoint{0.904918in}{0.907319in}}%
\pgfpathlineto{\pgfqpoint{0.909310in}{0.900286in}}%
\pgfpathlineto{\pgfqpoint{0.867036in}{0.904877in}}%
\pgfpathlineto{\pgfqpoint{0.825941in}{0.910590in}}%
\pgfpathclose%
\pgfusepath{fill}%
\end{pgfscope}%
\begin{pgfscope}%
\pgfpathrectangle{\pgfqpoint{0.407592in}{0.222904in}}{\pgfqpoint{1.411163in}{1.411163in}}%
\pgfusepath{clip}%
\pgfsetbuttcap%
\pgfsetroundjoin%
\definecolor{currentfill}{rgb}{0.749020,0.772549,0.792157}%
\pgfsetfillcolor{currentfill}%
\pgfsetfillopacity{0.160000}%
\pgfsetlinewidth{0.000000pt}%
\definecolor{currentstroke}{rgb}{0.000000,0.000000,0.000000}%
\pgfsetstrokecolor{currentstroke}%
\pgfsetdash{}{0pt}%
\pgfpathmoveto{\pgfqpoint{1.471354in}{0.983339in}}%
\pgfpathlineto{\pgfqpoint{1.479555in}{0.991361in}}%
\pgfpathlineto{\pgfqpoint{1.487793in}{0.999419in}}%
\pgfpathlineto{\pgfqpoint{1.530926in}{1.011075in}}%
\pgfpathlineto{\pgfqpoint{1.571653in}{1.023730in}}%
\pgfpathlineto{\pgfqpoint{1.561752in}{1.015233in}}%
\pgfpathlineto{\pgfqpoint{1.551892in}{1.006770in}}%
\pgfpathlineto{\pgfqpoint{1.512768in}{0.994572in}}%
\pgfpathlineto{\pgfqpoint{1.471354in}{0.983339in}}%
\pgfpathclose%
\pgfusepath{fill}%
\end{pgfscope}%
\begin{pgfscope}%
\pgfpathrectangle{\pgfqpoint{0.407592in}{0.222904in}}{\pgfqpoint{1.411163in}{1.411163in}}%
\pgfusepath{clip}%
\pgfsetbuttcap%
\pgfsetroundjoin%
\definecolor{currentfill}{rgb}{0.749020,0.772549,0.792157}%
\pgfsetfillcolor{currentfill}%
\pgfsetfillopacity{0.160000}%
\pgfsetlinewidth{0.000000pt}%
\definecolor{currentstroke}{rgb}{0.000000,0.000000,0.000000}%
\pgfsetstrokecolor{currentstroke}%
\pgfsetdash{}{0pt}%
\pgfpathmoveto{\pgfqpoint{0.731620in}{0.940422in}}%
\pgfpathlineto{\pgfqpoint{0.723306in}{0.948071in}}%
\pgfpathlineto{\pgfqpoint{0.714953in}{0.955755in}}%
\pgfpathlineto{\pgfqpoint{0.756676in}{0.947232in}}%
\pgfpathlineto{\pgfqpoint{0.800180in}{0.939838in}}%
\pgfpathlineto{\pgfqpoint{0.806667in}{0.932473in}}%
\pgfpathlineto{\pgfqpoint{0.813123in}{0.925144in}}%
\pgfpathlineto{\pgfqpoint{0.771527in}{0.932240in}}%
\pgfpathlineto{\pgfqpoint{0.731620in}{0.940422in}}%
\pgfpathclose%
\pgfusepath{fill}%
\end{pgfscope}%
\begin{pgfscope}%
\pgfpathrectangle{\pgfqpoint{0.407592in}{0.222904in}}{\pgfqpoint{1.411163in}{1.411163in}}%
\pgfusepath{clip}%
\pgfsetbuttcap%
\pgfsetroundjoin%
\definecolor{currentfill}{rgb}{0.749020,0.772549,0.792157}%
\pgfsetfillcolor{currentfill}%
\pgfsetfillopacity{0.160000}%
\pgfsetlinewidth{0.000000pt}%
\definecolor{currentstroke}{rgb}{0.000000,0.000000,0.000000}%
\pgfsetstrokecolor{currentstroke}%
\pgfsetdash{}{0pt}%
\pgfpathmoveto{\pgfqpoint{0.523697in}{1.034512in}}%
\pgfpathlineto{\pgfqpoint{0.512038in}{1.043237in}}%
\pgfpathlineto{\pgfqpoint{0.500331in}{1.051998in}}%
\pgfpathlineto{\pgfqpoint{0.537050in}{1.038124in}}%
\pgfpathlineto{\pgfqpoint{0.576694in}{1.025215in}}%
\pgfpathlineto{\pgfqpoint{0.586970in}{1.016917in}}%
\pgfpathlineto{\pgfqpoint{0.597202in}{1.008654in}}%
\pgfpathlineto{\pgfqpoint{0.559054in}{1.021115in}}%
\pgfpathlineto{\pgfqpoint{0.523697in}{1.034512in}}%
\pgfpathclose%
\pgfusepath{fill}%
\end{pgfscope}%
\begin{pgfscope}%
\pgfpathrectangle{\pgfqpoint{0.407592in}{0.222904in}}{\pgfqpoint{1.411163in}{1.411163in}}%
\pgfusepath{clip}%
\pgfsetbuttcap%
\pgfsetroundjoin%
\definecolor{currentfill}{rgb}{0.749020,0.772549,0.792157}%
\pgfsetfillcolor{currentfill}%
\pgfsetfillopacity{0.160000}%
\pgfsetlinewidth{0.000000pt}%
\definecolor{currentstroke}{rgb}{0.000000,0.000000,0.000000}%
\pgfsetstrokecolor{currentstroke}%
\pgfsetdash{}{0pt}%
\pgfpathmoveto{\pgfqpoint{1.174672in}{0.897443in}}%
\pgfpathlineto{\pgfqpoint{1.176723in}{0.904413in}}%
\pgfpathlineto{\pgfqpoint{1.178785in}{0.911418in}}%
\pgfpathlineto{\pgfqpoint{1.224849in}{0.915253in}}%
\pgfpathlineto{\pgfqpoint{1.270144in}{0.920284in}}%
\pgfpathlineto{\pgfqpoint{1.265967in}{0.913092in}}%
\pgfpathlineto{\pgfqpoint{1.261810in}{0.905934in}}%
\pgfpathlineto{\pgfqpoint{1.218603in}{0.901116in}}%
\pgfpathlineto{\pgfqpoint{1.174672in}{0.897443in}}%
\pgfpathclose%
\pgfusepath{fill}%
\end{pgfscope}%
\begin{pgfscope}%
\pgfpathrectangle{\pgfqpoint{0.407592in}{0.222904in}}{\pgfqpoint{1.411163in}{1.411163in}}%
\pgfusepath{clip}%
\pgfsetbuttcap%
\pgfsetroundjoin%
\definecolor{currentfill}{rgb}{0.749020,0.772549,0.792157}%
\pgfsetfillcolor{currentfill}%
\pgfsetfillopacity{0.160000}%
\pgfsetlinewidth{0.000000pt}%
\definecolor{currentstroke}{rgb}{0.000000,0.000000,0.000000}%
\pgfsetstrokecolor{currentstroke}%
\pgfsetdash{}{0pt}%
\pgfpathmoveto{\pgfqpoint{1.369911in}{0.948853in}}%
\pgfpathlineto{\pgfqpoint{1.376179in}{0.956434in}}%
\pgfpathlineto{\pgfqpoint{1.382476in}{0.964052in}}%
\pgfpathlineto{\pgfqpoint{1.427851in}{0.973144in}}%
\pgfpathlineto{\pgfqpoint{1.471354in}{0.983339in}}%
\pgfpathlineto{\pgfqpoint{1.463190in}{0.975353in}}%
\pgfpathlineto{\pgfqpoint{1.455062in}{0.967403in}}%
\pgfpathlineto{\pgfqpoint{1.413375in}{0.957596in}}%
\pgfpathlineto{\pgfqpoint{1.369911in}{0.948853in}}%
\pgfpathclose%
\pgfusepath{fill}%
\end{pgfscope}%
\begin{pgfscope}%
\pgfpathrectangle{\pgfqpoint{0.407592in}{0.222904in}}{\pgfqpoint{1.411163in}{1.411163in}}%
\pgfusepath{clip}%
\pgfsetbuttcap%
\pgfsetroundjoin%
\definecolor{currentfill}{rgb}{0.749020,0.772549,0.792157}%
\pgfsetfillcolor{currentfill}%
\pgfsetfillopacity{0.160000}%
\pgfsetlinewidth{0.000000pt}%
\definecolor{currentstroke}{rgb}{0.000000,0.000000,0.000000}%
\pgfsetstrokecolor{currentstroke}%
\pgfsetdash{}{0pt}%
\pgfpathmoveto{\pgfqpoint{0.909310in}{0.900286in}}%
\pgfpathlineto{\pgfqpoint{0.904918in}{0.907319in}}%
\pgfpathlineto{\pgfqpoint{0.900505in}{0.914386in}}%
\pgfpathlineto{\pgfqpoint{0.945828in}{0.910795in}}%
\pgfpathlineto{\pgfqpoint{0.991914in}{0.908430in}}%
\pgfpathlineto{\pgfqpoint{0.994204in}{0.901489in}}%
\pgfpathlineto{\pgfqpoint{0.996484in}{0.894582in}}%
\pgfpathlineto{\pgfqpoint{0.952536in}{0.896847in}}%
\pgfpathlineto{\pgfqpoint{0.909310in}{0.900286in}}%
\pgfpathclose%
\pgfusepath{fill}%
\end{pgfscope}%
\begin{pgfscope}%
\pgfpathrectangle{\pgfqpoint{0.407592in}{0.222904in}}{\pgfqpoint{1.411163in}{1.411163in}}%
\pgfusepath{clip}%
\pgfsetbuttcap%
\pgfsetroundjoin%
\definecolor{currentfill}{rgb}{0.749020,0.772549,0.792157}%
\pgfsetfillcolor{currentfill}%
\pgfsetfillopacity{0.160000}%
\pgfsetlinewidth{0.000000pt}%
\definecolor{currentstroke}{rgb}{0.000000,0.000000,0.000000}%
\pgfsetstrokecolor{currentstroke}%
\pgfsetdash{}{0pt}%
\pgfpathmoveto{\pgfqpoint{1.270144in}{0.920284in}}%
\pgfpathlineto{\pgfqpoint{1.274342in}{0.927511in}}%
\pgfpathlineto{\pgfqpoint{1.278561in}{0.934774in}}%
\pgfpathlineto{\pgfqpoint{1.324895in}{0.941229in}}%
\pgfpathlineto{\pgfqpoint{1.369911in}{0.948853in}}%
\pgfpathlineto{\pgfqpoint{1.363673in}{0.941307in}}%
\pgfpathlineto{\pgfqpoint{1.357464in}{0.933796in}}%
\pgfpathlineto{\pgfqpoint{1.314428in}{0.926478in}}%
\pgfpathlineto{\pgfqpoint{1.270144in}{0.920284in}}%
\pgfpathclose%
\pgfusepath{fill}%
\end{pgfscope}%
\begin{pgfscope}%
\pgfpathrectangle{\pgfqpoint{0.407592in}{0.222904in}}{\pgfqpoint{1.411163in}{1.411163in}}%
\pgfusepath{clip}%
\pgfsetbuttcap%
\pgfsetroundjoin%
\definecolor{currentfill}{rgb}{0.749020,0.772549,0.792157}%
\pgfsetfillcolor{currentfill}%
\pgfsetfillopacity{0.160000}%
\pgfsetlinewidth{0.000000pt}%
\definecolor{currentstroke}{rgb}{0.000000,0.000000,0.000000}%
\pgfsetstrokecolor{currentstroke}%
\pgfsetdash{}{0pt}%
\pgfpathmoveto{\pgfqpoint{1.085585in}{0.893626in}}%
\pgfpathlineto{\pgfqpoint{1.085472in}{0.900511in}}%
\pgfpathlineto{\pgfqpoint{1.085358in}{0.907432in}}%
\pgfpathlineto{\pgfqpoint{1.132203in}{0.908805in}}%
\pgfpathlineto{\pgfqpoint{1.178785in}{0.911418in}}%
\pgfpathlineto{\pgfqpoint{1.176723in}{0.904413in}}%
\pgfpathlineto{\pgfqpoint{1.174672in}{0.897443in}}%
\pgfpathlineto{\pgfqpoint{1.130252in}{0.894942in}}%
\pgfpathlineto{\pgfqpoint{1.085585in}{0.893626in}}%
\pgfpathclose%
\pgfusepath{fill}%
\end{pgfscope}%
\begin{pgfscope}%
\pgfpathrectangle{\pgfqpoint{0.407592in}{0.222904in}}{\pgfqpoint{1.411163in}{1.411163in}}%
\pgfusepath{clip}%
\pgfsetbuttcap%
\pgfsetroundjoin%
\definecolor{currentfill}{rgb}{0.749020,0.772549,0.792157}%
\pgfsetfillcolor{currentfill}%
\pgfsetfillopacity{0.160000}%
\pgfsetlinewidth{0.000000pt}%
\definecolor{currentstroke}{rgb}{0.000000,0.000000,0.000000}%
\pgfsetstrokecolor{currentstroke}%
\pgfsetdash{}{0pt}%
\pgfpathmoveto{\pgfqpoint{0.996484in}{0.894582in}}%
\pgfpathlineto{\pgfqpoint{0.994204in}{0.901489in}}%
\pgfpathlineto{\pgfqpoint{0.991914in}{0.908430in}}%
\pgfpathlineto{\pgfqpoint{1.038509in}{0.907306in}}%
\pgfpathlineto{\pgfqpoint{1.085358in}{0.907432in}}%
\pgfpathlineto{\pgfqpoint{1.085472in}{0.900511in}}%
\pgfpathlineto{\pgfqpoint{1.085585in}{0.893626in}}%
\pgfpathlineto{\pgfqpoint{1.040914in}{0.893506in}}%
\pgfpathlineto{\pgfqpoint{0.996484in}{0.894582in}}%
\pgfpathclose%
\pgfusepath{fill}%
\end{pgfscope}%
\begin{pgfscope}%
\pgfpathrectangle{\pgfqpoint{0.407592in}{0.222904in}}{\pgfqpoint{1.411163in}{1.411163in}}%
\pgfusepath{clip}%
\pgfsetbuttcap%
\pgfsetroundjoin%
\definecolor{currentfill}{rgb}{0.749020,0.772549,0.792157}%
\pgfsetfillcolor{currentfill}%
\pgfsetfillopacity{0.160000}%
\pgfsetlinewidth{0.000000pt}%
\definecolor{currentstroke}{rgb}{0.000000,0.000000,0.000000}%
\pgfsetstrokecolor{currentstroke}%
\pgfsetdash{}{0pt}%
\pgfpathmoveto{\pgfqpoint{0.617531in}{0.992238in}}%
\pgfpathlineto{\pgfqpoint{0.607389in}{1.000428in}}%
\pgfpathlineto{\pgfqpoint{0.597202in}{1.008654in}}%
\pgfpathlineto{\pgfqpoint{0.637966in}{0.997208in}}%
\pgfpathlineto{\pgfqpoint{0.681152in}{0.986849in}}%
\pgfpathlineto{\pgfqpoint{0.689662in}{0.979021in}}%
\pgfpathlineto{\pgfqpoint{0.698131in}{0.971230in}}%
\pgfpathlineto{\pgfqpoint{0.656679in}{0.981209in}}%
\pgfpathlineto{\pgfqpoint{0.617531in}{0.992238in}}%
\pgfpathclose%
\pgfusepath{fill}%
\end{pgfscope}%
\begin{pgfscope}%
\pgfpathrectangle{\pgfqpoint{0.407592in}{0.222904in}}{\pgfqpoint{1.411163in}{1.411163in}}%
\pgfusepath{clip}%
\pgfsetbuttcap%
\pgfsetroundjoin%
\definecolor{currentfill}{rgb}{0.749020,0.772549,0.792157}%
\pgfsetfillcolor{currentfill}%
\pgfsetfillopacity{0.160000}%
\pgfsetlinewidth{0.000000pt}%
\definecolor{currentstroke}{rgb}{0.000000,0.000000,0.000000}%
\pgfsetstrokecolor{currentstroke}%
\pgfsetdash{}{0pt}%
\pgfpathmoveto{\pgfqpoint{0.500331in}{1.051998in}}%
\pgfpathlineto{\pgfqpoint{0.488577in}{1.060795in}}%
\pgfpathlineto{\pgfqpoint{0.525977in}{1.046683in}}%
\pgfpathlineto{\pgfqpoint{0.566372in}{1.033550in}}%
\pgfpathlineto{\pgfqpoint{0.576694in}{1.025215in}}%
\pgfpathlineto{\pgfqpoint{0.537050in}{1.038124in}}%
\pgfusepath{fill}%
\end{pgfscope}%
\begin{pgfscope}%
\pgfpathrectangle{\pgfqpoint{0.407592in}{0.222904in}}{\pgfqpoint{1.411163in}{1.411163in}}%
\pgfusepath{clip}%
\pgfsetbuttcap%
\pgfsetroundjoin%
\definecolor{currentfill}{rgb}{0.749020,0.772549,0.792157}%
\pgfsetfillcolor{currentfill}%
\pgfsetfillopacity{0.160000}%
\pgfsetlinewidth{0.000000pt}%
\definecolor{currentstroke}{rgb}{0.000000,0.000000,0.000000}%
\pgfsetstrokecolor{currentstroke}%
\pgfsetdash{}{0pt}%
\pgfpathmoveto{\pgfqpoint{1.487793in}{0.999419in}}%
\pgfpathlineto{\pgfqpoint{1.496068in}{1.007514in}}%
\pgfpathlineto{\pgfqpoint{1.504380in}{1.015645in}}%
\pgfpathlineto{\pgfqpoint{1.549243in}{1.027723in}}%
\pgfpathlineto{\pgfqpoint{1.591580in}{1.040833in}}%
\pgfpathlineto{\pgfqpoint{1.581596in}{1.032264in}}%
\pgfpathlineto{\pgfqpoint{1.571653in}{1.023730in}}%
\pgfpathlineto{\pgfqpoint{1.530926in}{1.011075in}}%
\pgfpathlineto{\pgfqpoint{1.487793in}{0.999419in}}%
\pgfpathclose%
\pgfusepath{fill}%
\end{pgfscope}%
\begin{pgfscope}%
\pgfpathrectangle{\pgfqpoint{0.407592in}{0.222904in}}{\pgfqpoint{1.411163in}{1.411163in}}%
\pgfusepath{clip}%
\pgfsetbuttcap%
\pgfsetroundjoin%
\definecolor{currentfill}{rgb}{0.749020,0.772549,0.792157}%
\pgfsetfillcolor{currentfill}%
\pgfsetfillopacity{0.160000}%
\pgfsetlinewidth{0.000000pt}%
\definecolor{currentstroke}{rgb}{0.000000,0.000000,0.000000}%
\pgfsetstrokecolor{currentstroke}%
\pgfsetdash{}{0pt}%
\pgfpathmoveto{\pgfqpoint{0.813123in}{0.925144in}}%
\pgfpathlineto{\pgfqpoint{0.806667in}{0.932473in}}%
\pgfpathlineto{\pgfqpoint{0.800180in}{0.939838in}}%
\pgfpathlineto{\pgfqpoint{0.845239in}{0.933623in}}%
\pgfpathlineto{\pgfqpoint{0.891613in}{0.928627in}}%
\pgfpathlineto{\pgfqpoint{0.896071in}{0.921489in}}%
\pgfpathlineto{\pgfqpoint{0.900505in}{0.914386in}}%
\pgfpathlineto{\pgfqpoint{0.856191in}{0.919180in}}%
\pgfpathlineto{\pgfqpoint{0.813123in}{0.925144in}}%
\pgfpathclose%
\pgfusepath{fill}%
\end{pgfscope}%
\begin{pgfscope}%
\pgfpathrectangle{\pgfqpoint{0.407592in}{0.222904in}}{\pgfqpoint{1.411163in}{1.411163in}}%
\pgfusepath{clip}%
\pgfsetbuttcap%
\pgfsetroundjoin%
\definecolor{currentfill}{rgb}{0.749020,0.772549,0.792157}%
\pgfsetfillcolor{currentfill}%
\pgfsetfillopacity{0.160000}%
\pgfsetlinewidth{0.000000pt}%
\definecolor{currentstroke}{rgb}{0.000000,0.000000,0.000000}%
\pgfsetstrokecolor{currentstroke}%
\pgfsetdash{}{0pt}%
\pgfpathmoveto{\pgfqpoint{0.714953in}{0.955755in}}%
\pgfpathlineto{\pgfqpoint{0.706562in}{0.963474in}}%
\pgfpathlineto{\pgfqpoint{0.698131in}{0.971230in}}%
\pgfpathlineto{\pgfqpoint{0.741683in}{0.962366in}}%
\pgfpathlineto{\pgfqpoint{0.787111in}{0.954676in}}%
\pgfpathlineto{\pgfqpoint{0.793661in}{0.947239in}}%
\pgfpathlineto{\pgfqpoint{0.800180in}{0.939838in}}%
\pgfpathlineto{\pgfqpoint{0.756676in}{0.947232in}}%
\pgfpathlineto{\pgfqpoint{0.714953in}{0.955755in}}%
\pgfpathclose%
\pgfusepath{fill}%
\end{pgfscope}%
\begin{pgfscope}%
\pgfpathrectangle{\pgfqpoint{0.407592in}{0.222904in}}{\pgfqpoint{1.411163in}{1.411163in}}%
\pgfusepath{clip}%
\pgfsetbuttcap%
\pgfsetroundjoin%
\definecolor{currentfill}{rgb}{0.749020,0.772549,0.792157}%
\pgfsetfillcolor{currentfill}%
\pgfsetfillopacity{0.160000}%
\pgfsetlinewidth{0.000000pt}%
\definecolor{currentstroke}{rgb}{0.000000,0.000000,0.000000}%
\pgfsetstrokecolor{currentstroke}%
\pgfsetdash{}{0pt}%
\pgfpathmoveto{\pgfqpoint{1.178785in}{0.911418in}}%
\pgfpathlineto{\pgfqpoint{1.180857in}{0.918458in}}%
\pgfpathlineto{\pgfqpoint{1.182939in}{0.925533in}}%
\pgfpathlineto{\pgfqpoint{1.231156in}{0.929530in}}%
\pgfpathlineto{\pgfqpoint{1.278561in}{0.934774in}}%
\pgfpathlineto{\pgfqpoint{1.274342in}{0.927511in}}%
\pgfpathlineto{\pgfqpoint{1.270144in}{0.920284in}}%
\pgfpathlineto{\pgfqpoint{1.224849in}{0.915253in}}%
\pgfpathlineto{\pgfqpoint{1.178785in}{0.911418in}}%
\pgfpathclose%
\pgfusepath{fill}%
\end{pgfscope}%
\begin{pgfscope}%
\pgfpathrectangle{\pgfqpoint{0.407592in}{0.222904in}}{\pgfqpoint{1.411163in}{1.411163in}}%
\pgfusepath{clip}%
\pgfsetbuttcap%
\pgfsetroundjoin%
\definecolor{currentfill}{rgb}{0.749020,0.772549,0.792157}%
\pgfsetfillcolor{currentfill}%
\pgfsetfillopacity{0.160000}%
\pgfsetlinewidth{0.000000pt}%
\definecolor{currentstroke}{rgb}{0.000000,0.000000,0.000000}%
\pgfsetstrokecolor{currentstroke}%
\pgfsetdash{}{0pt}%
\pgfpathmoveto{\pgfqpoint{1.382476in}{0.964052in}}%
\pgfpathlineto{\pgfqpoint{1.388803in}{0.971706in}}%
\pgfpathlineto{\pgfqpoint{1.395161in}{0.979396in}}%
\pgfpathlineto{\pgfqpoint{1.442462in}{0.988836in}}%
\pgfpathlineto{\pgfqpoint{1.487793in}{0.999419in}}%
\pgfpathlineto{\pgfqpoint{1.479555in}{0.991361in}}%
\pgfpathlineto{\pgfqpoint{1.471354in}{0.983339in}}%
\pgfpathlineto{\pgfqpoint{1.427851in}{0.973144in}}%
\pgfpathlineto{\pgfqpoint{1.382476in}{0.964052in}}%
\pgfpathclose%
\pgfusepath{fill}%
\end{pgfscope}%
\begin{pgfscope}%
\pgfpathrectangle{\pgfqpoint{0.407592in}{0.222904in}}{\pgfqpoint{1.411163in}{1.411163in}}%
\pgfusepath{clip}%
\pgfsetbuttcap%
\pgfsetroundjoin%
\definecolor{currentfill}{rgb}{0.749020,0.772549,0.792157}%
\pgfsetfillcolor{currentfill}%
\pgfsetfillopacity{0.160000}%
\pgfsetlinewidth{0.000000pt}%
\definecolor{currentstroke}{rgb}{0.000000,0.000000,0.000000}%
\pgfsetstrokecolor{currentstroke}%
\pgfsetdash{}{0pt}%
\pgfpathmoveto{\pgfqpoint{1.278561in}{0.934774in}}%
\pgfpathlineto{\pgfqpoint{1.282800in}{0.942072in}}%
\pgfpathlineto{\pgfqpoint{1.287060in}{0.949407in}}%
\pgfpathlineto{\pgfqpoint{1.335463in}{0.956122in}}%
\pgfpathlineto{\pgfqpoint{1.382476in}{0.964052in}}%
\pgfpathlineto{\pgfqpoint{1.376179in}{0.956434in}}%
\pgfpathlineto{\pgfqpoint{1.369911in}{0.948853in}}%
\pgfpathlineto{\pgfqpoint{1.324895in}{0.941229in}}%
\pgfpathlineto{\pgfqpoint{1.278561in}{0.934774in}}%
\pgfpathclose%
\pgfusepath{fill}%
\end{pgfscope}%
\begin{pgfscope}%
\pgfpathrectangle{\pgfqpoint{0.407592in}{0.222904in}}{\pgfqpoint{1.411163in}{1.411163in}}%
\pgfusepath{clip}%
\pgfsetbuttcap%
\pgfsetroundjoin%
\definecolor{currentfill}{rgb}{0.749020,0.772549,0.792157}%
\pgfsetfillcolor{currentfill}%
\pgfsetfillopacity{0.160000}%
\pgfsetlinewidth{0.000000pt}%
\definecolor{currentstroke}{rgb}{0.000000,0.000000,0.000000}%
\pgfsetstrokecolor{currentstroke}%
\pgfsetdash{}{0pt}%
\pgfpathmoveto{\pgfqpoint{0.900505in}{0.914386in}}%
\pgfpathlineto{\pgfqpoint{0.896071in}{0.921489in}}%
\pgfpathlineto{\pgfqpoint{0.891613in}{0.928627in}}%
\pgfpathlineto{\pgfqpoint{0.939053in}{0.924884in}}%
\pgfpathlineto{\pgfqpoint{0.987297in}{0.922418in}}%
\pgfpathlineto{\pgfqpoint{0.989611in}{0.915406in}}%
\pgfpathlineto{\pgfqpoint{0.991914in}{0.908430in}}%
\pgfpathlineto{\pgfqpoint{0.945828in}{0.910795in}}%
\pgfpathlineto{\pgfqpoint{0.900505in}{0.914386in}}%
\pgfpathclose%
\pgfusepath{fill}%
\end{pgfscope}%
\begin{pgfscope}%
\pgfpathrectangle{\pgfqpoint{0.407592in}{0.222904in}}{\pgfqpoint{1.411163in}{1.411163in}}%
\pgfusepath{clip}%
\pgfsetbuttcap%
\pgfsetroundjoin%
\definecolor{currentfill}{rgb}{0.749020,0.772549,0.792157}%
\pgfsetfillcolor{currentfill}%
\pgfsetfillopacity{0.160000}%
\pgfsetlinewidth{0.000000pt}%
\definecolor{currentstroke}{rgb}{0.000000,0.000000,0.000000}%
\pgfsetstrokecolor{currentstroke}%
\pgfsetdash{}{0pt}%
\pgfpathmoveto{\pgfqpoint{1.504380in}{1.015645in}}%
\pgfpathlineto{\pgfqpoint{1.512730in}{1.023812in}}%
\pgfpathlineto{\pgfqpoint{1.558462in}{1.036102in}}%
\pgfpathlineto{\pgfqpoint{1.601606in}{1.049438in}}%
\pgfpathlineto{\pgfqpoint{1.591580in}{1.040833in}}%
\pgfpathlineto{\pgfqpoint{1.549243in}{1.027723in}}%
\pgfusepath{fill}%
\end{pgfscope}%
\begin{pgfscope}%
\pgfpathrectangle{\pgfqpoint{0.407592in}{0.222904in}}{\pgfqpoint{1.411163in}{1.411163in}}%
\pgfusepath{clip}%
\pgfsetbuttcap%
\pgfsetroundjoin%
\definecolor{currentfill}{rgb}{0.749020,0.772549,0.792157}%
\pgfsetfillcolor{currentfill}%
\pgfsetfillopacity{0.160000}%
\pgfsetlinewidth{0.000000pt}%
\definecolor{currentstroke}{rgb}{0.000000,0.000000,0.000000}%
\pgfsetstrokecolor{currentstroke}%
\pgfsetdash{}{0pt}%
\pgfpathmoveto{\pgfqpoint{1.085358in}{0.907432in}}%
\pgfpathlineto{\pgfqpoint{1.085244in}{0.914387in}}%
\pgfpathlineto{\pgfqpoint{1.085129in}{0.921378in}}%
\pgfpathlineto{\pgfqpoint{1.134173in}{0.922810in}}%
\pgfpathlineto{\pgfqpoint{1.182939in}{0.925533in}}%
\pgfpathlineto{\pgfqpoint{1.180857in}{0.918458in}}%
\pgfpathlineto{\pgfqpoint{1.178785in}{0.911418in}}%
\pgfpathlineto{\pgfqpoint{1.132203in}{0.908805in}}%
\pgfpathlineto{\pgfqpoint{1.085358in}{0.907432in}}%
\pgfpathclose%
\pgfusepath{fill}%
\end{pgfscope}%
\begin{pgfscope}%
\pgfpathrectangle{\pgfqpoint{0.407592in}{0.222904in}}{\pgfqpoint{1.411163in}{1.411163in}}%
\pgfusepath{clip}%
\pgfsetbuttcap%
\pgfsetroundjoin%
\definecolor{currentfill}{rgb}{0.749020,0.772549,0.792157}%
\pgfsetfillcolor{currentfill}%
\pgfsetfillopacity{0.160000}%
\pgfsetlinewidth{0.000000pt}%
\definecolor{currentstroke}{rgb}{0.000000,0.000000,0.000000}%
\pgfsetstrokecolor{currentstroke}%
\pgfsetdash{}{0pt}%
\pgfpathmoveto{\pgfqpoint{0.991914in}{0.908430in}}%
\pgfpathlineto{\pgfqpoint{0.989611in}{0.915406in}}%
\pgfpathlineto{\pgfqpoint{0.987297in}{0.922418in}}%
\pgfpathlineto{\pgfqpoint{1.036079in}{0.921247in}}%
\pgfpathlineto{\pgfqpoint{1.085129in}{0.921378in}}%
\pgfpathlineto{\pgfqpoint{1.085244in}{0.914387in}}%
\pgfpathlineto{\pgfqpoint{1.085358in}{0.907432in}}%
\pgfpathlineto{\pgfqpoint{1.038509in}{0.907306in}}%
\pgfpathlineto{\pgfqpoint{0.991914in}{0.908430in}}%
\pgfpathclose%
\pgfusepath{fill}%
\end{pgfscope}%
\begin{pgfscope}%
\pgfpathrectangle{\pgfqpoint{0.407592in}{0.222904in}}{\pgfqpoint{1.411163in}{1.411163in}}%
\pgfusepath{clip}%
\pgfsetbuttcap%
\pgfsetroundjoin%
\definecolor{currentfill}{rgb}{0.749020,0.772549,0.792157}%
\pgfsetfillcolor{currentfill}%
\pgfsetfillopacity{0.160000}%
\pgfsetlinewidth{0.000000pt}%
\definecolor{currentstroke}{rgb}{0.000000,0.000000,0.000000}%
\pgfsetstrokecolor{currentstroke}%
\pgfsetdash{}{0pt}%
\pgfpathmoveto{\pgfqpoint{0.597202in}{1.008654in}}%
\pgfpathlineto{\pgfqpoint{0.586970in}{1.016917in}}%
\pgfpathlineto{\pgfqpoint{0.576694in}{1.025215in}}%
\pgfpathlineto{\pgfqpoint{0.619083in}{1.013352in}}%
\pgfpathlineto{\pgfqpoint{0.664014in}{1.002614in}}%
\pgfpathlineto{\pgfqpoint{0.672603in}{0.994713in}}%
\pgfpathlineto{\pgfqpoint{0.681152in}{0.986849in}}%
\pgfpathlineto{\pgfqpoint{0.637966in}{0.997208in}}%
\pgfpathlineto{\pgfqpoint{0.597202in}{1.008654in}}%
\pgfpathclose%
\pgfusepath{fill}%
\end{pgfscope}%
\begin{pgfscope}%
\pgfpathrectangle{\pgfqpoint{0.407592in}{0.222904in}}{\pgfqpoint{1.411163in}{1.411163in}}%
\pgfusepath{clip}%
\pgfsetbuttcap%
\pgfsetroundjoin%
\definecolor{currentfill}{rgb}{0.749020,0.772549,0.792157}%
\pgfsetfillcolor{currentfill}%
\pgfsetfillopacity{0.160000}%
\pgfsetlinewidth{0.000000pt}%
\definecolor{currentstroke}{rgb}{0.000000,0.000000,0.000000}%
\pgfsetstrokecolor{currentstroke}%
\pgfsetdash{}{0pt}%
\pgfpathmoveto{\pgfqpoint{0.800180in}{0.939838in}}%
\pgfpathlineto{\pgfqpoint{0.793661in}{0.947239in}}%
\pgfpathlineto{\pgfqpoint{0.787111in}{0.954676in}}%
\pgfpathlineto{\pgfqpoint{0.834178in}{0.948210in}}%
\pgfpathlineto{\pgfqpoint{0.882632in}{0.943011in}}%
\pgfpathlineto{\pgfqpoint{0.887134in}{0.935801in}}%
\pgfpathlineto{\pgfqpoint{0.891613in}{0.928627in}}%
\pgfpathlineto{\pgfqpoint{0.845239in}{0.933623in}}%
\pgfpathlineto{\pgfqpoint{0.800180in}{0.939838in}}%
\pgfpathclose%
\pgfusepath{fill}%
\end{pgfscope}%
\begin{pgfscope}%
\pgfpathrectangle{\pgfqpoint{0.407592in}{0.222904in}}{\pgfqpoint{1.411163in}{1.411163in}}%
\pgfusepath{clip}%
\pgfsetbuttcap%
\pgfsetroundjoin%
\definecolor{currentfill}{rgb}{0.749020,0.772549,0.792157}%
\pgfsetfillcolor{currentfill}%
\pgfsetfillopacity{0.160000}%
\pgfsetlinewidth{0.000000pt}%
\definecolor{currentstroke}{rgb}{0.000000,0.000000,0.000000}%
\pgfsetstrokecolor{currentstroke}%
\pgfsetdash{}{0pt}%
\pgfpathmoveto{\pgfqpoint{0.698131in}{0.971230in}}%
\pgfpathlineto{\pgfqpoint{0.689662in}{0.979021in}}%
\pgfpathlineto{\pgfqpoint{0.681152in}{0.986849in}}%
\pgfpathlineto{\pgfqpoint{0.726546in}{0.977645in}}%
\pgfpathlineto{\pgfqpoint{0.773915in}{0.969658in}}%
\pgfpathlineto{\pgfqpoint{0.780529in}{0.962149in}}%
\pgfpathlineto{\pgfqpoint{0.787111in}{0.954676in}}%
\pgfpathlineto{\pgfqpoint{0.741683in}{0.962366in}}%
\pgfpathlineto{\pgfqpoint{0.698131in}{0.971230in}}%
\pgfpathclose%
\pgfusepath{fill}%
\end{pgfscope}%
\begin{pgfscope}%
\pgfpathrectangle{\pgfqpoint{0.407592in}{0.222904in}}{\pgfqpoint{1.411163in}{1.411163in}}%
\pgfusepath{clip}%
\pgfsetbuttcap%
\pgfsetroundjoin%
\definecolor{currentfill}{rgb}{0.749020,0.772549,0.792157}%
\pgfsetfillcolor{currentfill}%
\pgfsetfillopacity{0.160000}%
\pgfsetlinewidth{0.000000pt}%
\definecolor{currentstroke}{rgb}{0.000000,0.000000,0.000000}%
\pgfsetstrokecolor{currentstroke}%
\pgfsetdash{}{0pt}%
\pgfpathmoveto{\pgfqpoint{1.182939in}{0.925533in}}%
\pgfpathlineto{\pgfqpoint{1.185032in}{0.932644in}}%
\pgfpathlineto{\pgfqpoint{1.187135in}{0.939791in}}%
\pgfpathlineto{\pgfqpoint{1.237527in}{0.943951in}}%
\pgfpathlineto{\pgfqpoint{1.287060in}{0.949407in}}%
\pgfpathlineto{\pgfqpoint{1.282800in}{0.942072in}}%
\pgfpathlineto{\pgfqpoint{1.278561in}{0.934774in}}%
\pgfpathlineto{\pgfqpoint{1.231156in}{0.929530in}}%
\pgfpathlineto{\pgfqpoint{1.182939in}{0.925533in}}%
\pgfpathclose%
\pgfusepath{fill}%
\end{pgfscope}%
\begin{pgfscope}%
\pgfpathrectangle{\pgfqpoint{0.407592in}{0.222904in}}{\pgfqpoint{1.411163in}{1.411163in}}%
\pgfusepath{clip}%
\pgfsetbuttcap%
\pgfsetroundjoin%
\definecolor{currentfill}{rgb}{0.749020,0.772549,0.792157}%
\pgfsetfillcolor{currentfill}%
\pgfsetfillopacity{0.160000}%
\pgfsetlinewidth{0.000000pt}%
\definecolor{currentstroke}{rgb}{0.000000,0.000000,0.000000}%
\pgfsetstrokecolor{currentstroke}%
\pgfsetdash{}{0pt}%
\pgfpathmoveto{\pgfqpoint{1.395161in}{0.979396in}}%
\pgfpathlineto{\pgfqpoint{1.401548in}{0.987123in}}%
\pgfpathlineto{\pgfqpoint{1.407967in}{0.994887in}}%
\pgfpathlineto{\pgfqpoint{1.457209in}{1.004675in}}%
\pgfpathlineto{\pgfqpoint{1.504380in}{1.015645in}}%
\pgfpathlineto{\pgfqpoint{1.496068in}{1.007514in}}%
\pgfpathlineto{\pgfqpoint{1.487793in}{0.999419in}}%
\pgfpathlineto{\pgfqpoint{1.442462in}{0.988836in}}%
\pgfpathlineto{\pgfqpoint{1.395161in}{0.979396in}}%
\pgfpathclose%
\pgfusepath{fill}%
\end{pgfscope}%
\begin{pgfscope}%
\pgfpathrectangle{\pgfqpoint{0.407592in}{0.222904in}}{\pgfqpoint{1.411163in}{1.411163in}}%
\pgfusepath{clip}%
\pgfsetbuttcap%
\pgfsetroundjoin%
\definecolor{currentfill}{rgb}{0.749020,0.772549,0.792157}%
\pgfsetfillcolor{currentfill}%
\pgfsetfillopacity{0.160000}%
\pgfsetlinewidth{0.000000pt}%
\definecolor{currentstroke}{rgb}{0.000000,0.000000,0.000000}%
\pgfsetstrokecolor{currentstroke}%
\pgfsetdash{}{0pt}%
\pgfpathmoveto{\pgfqpoint{0.576694in}{1.025215in}}%
\pgfpathlineto{\pgfqpoint{0.566372in}{1.033550in}}%
\pgfpathlineto{\pgfqpoint{0.609577in}{1.021480in}}%
\pgfpathlineto{\pgfqpoint{0.655385in}{1.010552in}}%
\pgfpathlineto{\pgfqpoint{0.664014in}{1.002614in}}%
\pgfpathlineto{\pgfqpoint{0.619083in}{1.013352in}}%
\pgfusepath{fill}%
\end{pgfscope}%
\begin{pgfscope}%
\pgfpathrectangle{\pgfqpoint{0.407592in}{0.222904in}}{\pgfqpoint{1.411163in}{1.411163in}}%
\pgfusepath{clip}%
\pgfsetbuttcap%
\pgfsetroundjoin%
\definecolor{currentfill}{rgb}{0.749020,0.772549,0.792157}%
\pgfsetfillcolor{currentfill}%
\pgfsetfillopacity{0.160000}%
\pgfsetlinewidth{0.000000pt}%
\definecolor{currentstroke}{rgb}{0.000000,0.000000,0.000000}%
\pgfsetstrokecolor{currentstroke}%
\pgfsetdash{}{0pt}%
\pgfpathmoveto{\pgfqpoint{1.287060in}{0.949407in}}%
\pgfpathlineto{\pgfqpoint{1.291341in}{0.956777in}}%
\pgfpathlineto{\pgfqpoint{1.295644in}{0.964185in}}%
\pgfpathlineto{\pgfqpoint{1.346134in}{0.971160in}}%
\pgfpathlineto{\pgfqpoint{1.395161in}{0.979396in}}%
\pgfpathlineto{\pgfqpoint{1.388803in}{0.971706in}}%
\pgfpathlineto{\pgfqpoint{1.382476in}{0.964052in}}%
\pgfpathlineto{\pgfqpoint{1.335463in}{0.956122in}}%
\pgfpathlineto{\pgfqpoint{1.287060in}{0.949407in}}%
\pgfpathclose%
\pgfusepath{fill}%
\end{pgfscope}%
\begin{pgfscope}%
\pgfpathrectangle{\pgfqpoint{0.407592in}{0.222904in}}{\pgfqpoint{1.411163in}{1.411163in}}%
\pgfusepath{clip}%
\pgfsetbuttcap%
\pgfsetroundjoin%
\definecolor{currentfill}{rgb}{0.749020,0.772549,0.792157}%
\pgfsetfillcolor{currentfill}%
\pgfsetfillopacity{0.160000}%
\pgfsetlinewidth{0.000000pt}%
\definecolor{currentstroke}{rgb}{0.000000,0.000000,0.000000}%
\pgfsetstrokecolor{currentstroke}%
\pgfsetdash{}{0pt}%
\pgfpathmoveto{\pgfqpoint{0.891613in}{0.928627in}}%
\pgfpathlineto{\pgfqpoint{0.887134in}{0.935801in}}%
\pgfpathlineto{\pgfqpoint{0.882632in}{0.943011in}}%
\pgfpathlineto{\pgfqpoint{0.932208in}{0.939116in}}%
\pgfpathlineto{\pgfqpoint{0.982633in}{0.936550in}}%
\pgfpathlineto{\pgfqpoint{0.984971in}{0.929466in}}%
\pgfpathlineto{\pgfqpoint{0.987297in}{0.922418in}}%
\pgfpathlineto{\pgfqpoint{0.939053in}{0.924884in}}%
\pgfpathlineto{\pgfqpoint{0.891613in}{0.928627in}}%
\pgfpathclose%
\pgfusepath{fill}%
\end{pgfscope}%
\begin{pgfscope}%
\pgfpathrectangle{\pgfqpoint{0.407592in}{0.222904in}}{\pgfqpoint{1.411163in}{1.411163in}}%
\pgfusepath{clip}%
\pgfsetbuttcap%
\pgfsetroundjoin%
\definecolor{currentfill}{rgb}{0.749020,0.772549,0.792157}%
\pgfsetfillcolor{currentfill}%
\pgfsetfillopacity{0.160000}%
\pgfsetlinewidth{0.000000pt}%
\definecolor{currentstroke}{rgb}{0.000000,0.000000,0.000000}%
\pgfsetstrokecolor{currentstroke}%
\pgfsetdash{}{0pt}%
\pgfpathmoveto{\pgfqpoint{1.085129in}{0.921378in}}%
\pgfpathlineto{\pgfqpoint{1.085014in}{0.928404in}}%
\pgfpathlineto{\pgfqpoint{1.084898in}{0.935467in}}%
\pgfpathlineto{\pgfqpoint{1.136164in}{0.936957in}}%
\pgfpathlineto{\pgfqpoint{1.187135in}{0.939791in}}%
\pgfpathlineto{\pgfqpoint{1.185032in}{0.932644in}}%
\pgfpathlineto{\pgfqpoint{1.182939in}{0.925533in}}%
\pgfpathlineto{\pgfqpoint{1.134173in}{0.922810in}}%
\pgfpathlineto{\pgfqpoint{1.085129in}{0.921378in}}%
\pgfpathclose%
\pgfusepath{fill}%
\end{pgfscope}%
\begin{pgfscope}%
\pgfpathrectangle{\pgfqpoint{0.407592in}{0.222904in}}{\pgfqpoint{1.411163in}{1.411163in}}%
\pgfusepath{clip}%
\pgfsetbuttcap%
\pgfsetroundjoin%
\definecolor{currentfill}{rgb}{0.749020,0.772549,0.792157}%
\pgfsetfillcolor{currentfill}%
\pgfsetfillopacity{0.160000}%
\pgfsetlinewidth{0.000000pt}%
\definecolor{currentstroke}{rgb}{0.000000,0.000000,0.000000}%
\pgfsetstrokecolor{currentstroke}%
\pgfsetdash{}{0pt}%
\pgfpathmoveto{\pgfqpoint{0.987297in}{0.922418in}}%
\pgfpathlineto{\pgfqpoint{0.984971in}{0.929466in}}%
\pgfpathlineto{\pgfqpoint{0.982633in}{0.936550in}}%
\pgfpathlineto{\pgfqpoint{1.033625in}{0.935330in}}%
\pgfpathlineto{\pgfqpoint{1.084898in}{0.935467in}}%
\pgfpathlineto{\pgfqpoint{1.085014in}{0.928404in}}%
\pgfpathlineto{\pgfqpoint{1.085129in}{0.921378in}}%
\pgfpathlineto{\pgfqpoint{1.036079in}{0.921247in}}%
\pgfpathlineto{\pgfqpoint{0.987297in}{0.922418in}}%
\pgfpathclose%
\pgfusepath{fill}%
\end{pgfscope}%
\begin{pgfscope}%
\pgfpathrectangle{\pgfqpoint{0.407592in}{0.222904in}}{\pgfqpoint{1.411163in}{1.411163in}}%
\pgfusepath{clip}%
\pgfsetbuttcap%
\pgfsetroundjoin%
\definecolor{currentfill}{rgb}{0.749020,0.772549,0.792157}%
\pgfsetfillcolor{currentfill}%
\pgfsetfillopacity{0.160000}%
\pgfsetlinewidth{0.000000pt}%
\definecolor{currentstroke}{rgb}{0.000000,0.000000,0.000000}%
\pgfsetstrokecolor{currentstroke}%
\pgfsetdash{}{0pt}%
\pgfpathmoveto{\pgfqpoint{0.787111in}{0.954676in}}%
\pgfpathlineto{\pgfqpoint{0.780529in}{0.962149in}}%
\pgfpathlineto{\pgfqpoint{0.773915in}{0.969658in}}%
\pgfpathlineto{\pgfqpoint{0.823008in}{0.962941in}}%
\pgfpathlineto{\pgfqpoint{0.873560in}{0.957540in}}%
\pgfpathlineto{\pgfqpoint{0.878108in}{0.950258in}}%
\pgfpathlineto{\pgfqpoint{0.882632in}{0.943011in}}%
\pgfpathlineto{\pgfqpoint{0.834178in}{0.948210in}}%
\pgfpathlineto{\pgfqpoint{0.787111in}{0.954676in}}%
\pgfpathclose%
\pgfusepath{fill}%
\end{pgfscope}%
\begin{pgfscope}%
\pgfpathrectangle{\pgfqpoint{0.407592in}{0.222904in}}{\pgfqpoint{1.411163in}{1.411163in}}%
\pgfusepath{clip}%
\pgfsetbuttcap%
\pgfsetroundjoin%
\definecolor{currentfill}{rgb}{0.749020,0.772549,0.792157}%
\pgfsetfillcolor{currentfill}%
\pgfsetfillopacity{0.160000}%
\pgfsetlinewidth{0.000000pt}%
\definecolor{currentstroke}{rgb}{0.000000,0.000000,0.000000}%
\pgfsetstrokecolor{currentstroke}%
\pgfsetdash{}{0pt}%
\pgfpathmoveto{\pgfqpoint{1.407967in}{0.994887in}}%
\pgfpathlineto{\pgfqpoint{1.414415in}{1.002688in}}%
\pgfpathlineto{\pgfqpoint{1.464634in}{1.012649in}}%
\pgfpathlineto{\pgfqpoint{1.512730in}{1.023812in}}%
\pgfpathlineto{\pgfqpoint{1.504380in}{1.015645in}}%
\pgfpathlineto{\pgfqpoint{1.457209in}{1.004675in}}%
\pgfusepath{fill}%
\end{pgfscope}%
\begin{pgfscope}%
\pgfpathrectangle{\pgfqpoint{0.407592in}{0.222904in}}{\pgfqpoint{1.411163in}{1.411163in}}%
\pgfusepath{clip}%
\pgfsetbuttcap%
\pgfsetroundjoin%
\definecolor{currentfill}{rgb}{0.749020,0.772549,0.792157}%
\pgfsetfillcolor{currentfill}%
\pgfsetfillopacity{0.160000}%
\pgfsetlinewidth{0.000000pt}%
\definecolor{currentstroke}{rgb}{0.000000,0.000000,0.000000}%
\pgfsetstrokecolor{currentstroke}%
\pgfsetdash{}{0pt}%
\pgfpathmoveto{\pgfqpoint{0.681152in}{0.986849in}}%
\pgfpathlineto{\pgfqpoint{0.672603in}{0.994713in}}%
\pgfpathlineto{\pgfqpoint{0.664014in}{1.002614in}}%
\pgfpathlineto{\pgfqpoint{0.711264in}{0.993071in}}%
\pgfpathlineto{\pgfqpoint{0.760589in}{0.984787in}}%
\pgfpathlineto{\pgfqpoint{0.767268in}{0.977204in}}%
\pgfpathlineto{\pgfqpoint{0.773915in}{0.969658in}}%
\pgfpathlineto{\pgfqpoint{0.726546in}{0.977645in}}%
\pgfpathlineto{\pgfqpoint{0.681152in}{0.986849in}}%
\pgfpathclose%
\pgfusepath{fill}%
\end{pgfscope}%
\begin{pgfscope}%
\pgfpathrectangle{\pgfqpoint{0.407592in}{0.222904in}}{\pgfqpoint{1.411163in}{1.411163in}}%
\pgfusepath{clip}%
\pgfsetbuttcap%
\pgfsetroundjoin%
\definecolor{currentfill}{rgb}{0.749020,0.772549,0.792157}%
\pgfsetfillcolor{currentfill}%
\pgfsetfillopacity{0.160000}%
\pgfsetlinewidth{0.000000pt}%
\definecolor{currentstroke}{rgb}{0.000000,0.000000,0.000000}%
\pgfsetstrokecolor{currentstroke}%
\pgfsetdash{}{0pt}%
\pgfpathmoveto{\pgfqpoint{1.187135in}{0.939791in}}%
\pgfpathlineto{\pgfqpoint{1.189250in}{0.946975in}}%
\pgfpathlineto{\pgfqpoint{1.191375in}{0.954195in}}%
\pgfpathlineto{\pgfqpoint{1.243962in}{0.958517in}}%
\pgfpathlineto{\pgfqpoint{1.295644in}{0.964185in}}%
\pgfpathlineto{\pgfqpoint{1.291341in}{0.956777in}}%
\pgfpathlineto{\pgfqpoint{1.287060in}{0.949407in}}%
\pgfpathlineto{\pgfqpoint{1.237527in}{0.943951in}}%
\pgfpathlineto{\pgfqpoint{1.187135in}{0.939791in}}%
\pgfpathclose%
\pgfusepath{fill}%
\end{pgfscope}%
\begin{pgfscope}%
\pgfpathrectangle{\pgfqpoint{0.407592in}{0.222904in}}{\pgfqpoint{1.411163in}{1.411163in}}%
\pgfusepath{clip}%
\pgfsetbuttcap%
\pgfsetroundjoin%
\definecolor{currentfill}{rgb}{0.749020,0.772549,0.792157}%
\pgfsetfillcolor{currentfill}%
\pgfsetfillopacity{0.160000}%
\pgfsetlinewidth{0.000000pt}%
\definecolor{currentstroke}{rgb}{0.000000,0.000000,0.000000}%
\pgfsetstrokecolor{currentstroke}%
\pgfsetdash{}{0pt}%
\pgfpathmoveto{\pgfqpoint{1.295644in}{0.964185in}}%
\pgfpathlineto{\pgfqpoint{1.299967in}{0.971629in}}%
\pgfpathlineto{\pgfqpoint{1.304313in}{0.979110in}}%
\pgfpathlineto{\pgfqpoint{1.356910in}{0.986345in}}%
\pgfpathlineto{\pgfqpoint{1.407967in}{0.994887in}}%
\pgfpathlineto{\pgfqpoint{1.401548in}{0.987123in}}%
\pgfpathlineto{\pgfqpoint{1.395161in}{0.979396in}}%
\pgfpathlineto{\pgfqpoint{1.346134in}{0.971160in}}%
\pgfpathlineto{\pgfqpoint{1.295644in}{0.964185in}}%
\pgfpathclose%
\pgfusepath{fill}%
\end{pgfscope}%
\begin{pgfscope}%
\pgfpathrectangle{\pgfqpoint{0.407592in}{0.222904in}}{\pgfqpoint{1.411163in}{1.411163in}}%
\pgfusepath{clip}%
\pgfsetbuttcap%
\pgfsetroundjoin%
\definecolor{currentfill}{rgb}{0.749020,0.772549,0.792157}%
\pgfsetfillcolor{currentfill}%
\pgfsetfillopacity{0.160000}%
\pgfsetlinewidth{0.000000pt}%
\definecolor{currentstroke}{rgb}{0.000000,0.000000,0.000000}%
\pgfsetstrokecolor{currentstroke}%
\pgfsetdash{}{0pt}%
\pgfpathmoveto{\pgfqpoint{0.882632in}{0.943011in}}%
\pgfpathlineto{\pgfqpoint{0.878108in}{0.950258in}}%
\pgfpathlineto{\pgfqpoint{0.873560in}{0.957540in}}%
\pgfpathlineto{\pgfqpoint{0.925294in}{0.953493in}}%
\pgfpathlineto{\pgfqpoint{0.977921in}{0.950826in}}%
\pgfpathlineto{\pgfqpoint{0.980283in}{0.943670in}}%
\pgfpathlineto{\pgfqpoint{0.982633in}{0.936550in}}%
\pgfpathlineto{\pgfqpoint{0.932208in}{0.939116in}}%
\pgfpathlineto{\pgfqpoint{0.882632in}{0.943011in}}%
\pgfpathclose%
\pgfusepath{fill}%
\end{pgfscope}%
\begin{pgfscope}%
\pgfpathrectangle{\pgfqpoint{0.407592in}{0.222904in}}{\pgfqpoint{1.411163in}{1.411163in}}%
\pgfusepath{clip}%
\pgfsetbuttcap%
\pgfsetroundjoin%
\definecolor{currentfill}{rgb}{0.749020,0.772549,0.792157}%
\pgfsetfillcolor{currentfill}%
\pgfsetfillopacity{0.160000}%
\pgfsetlinewidth{0.000000pt}%
\definecolor{currentstroke}{rgb}{0.000000,0.000000,0.000000}%
\pgfsetstrokecolor{currentstroke}%
\pgfsetdash{}{0pt}%
\pgfpathmoveto{\pgfqpoint{1.084898in}{0.935467in}}%
\pgfpathlineto{\pgfqpoint{1.084782in}{0.942565in}}%
\pgfpathlineto{\pgfqpoint{1.084665in}{0.949700in}}%
\pgfpathlineto{\pgfqpoint{1.138175in}{0.951249in}}%
\pgfpathlineto{\pgfqpoint{1.191375in}{0.954195in}}%
\pgfpathlineto{\pgfqpoint{1.189250in}{0.946975in}}%
\pgfpathlineto{\pgfqpoint{1.187135in}{0.939791in}}%
\pgfpathlineto{\pgfqpoint{1.136164in}{0.936957in}}%
\pgfpathlineto{\pgfqpoint{1.084898in}{0.935467in}}%
\pgfpathclose%
\pgfusepath{fill}%
\end{pgfscope}%
\begin{pgfscope}%
\pgfpathrectangle{\pgfqpoint{0.407592in}{0.222904in}}{\pgfqpoint{1.411163in}{1.411163in}}%
\pgfusepath{clip}%
\pgfsetbuttcap%
\pgfsetroundjoin%
\definecolor{currentfill}{rgb}{0.749020,0.772549,0.792157}%
\pgfsetfillcolor{currentfill}%
\pgfsetfillopacity{0.160000}%
\pgfsetlinewidth{0.000000pt}%
\definecolor{currentstroke}{rgb}{0.000000,0.000000,0.000000}%
\pgfsetstrokecolor{currentstroke}%
\pgfsetdash{}{0pt}%
\pgfpathmoveto{\pgfqpoint{0.664014in}{1.002614in}}%
\pgfpathlineto{\pgfqpoint{0.655385in}{1.010552in}}%
\pgfpathlineto{\pgfqpoint{0.703568in}{1.000840in}}%
\pgfpathlineto{\pgfqpoint{0.753877in}{0.992408in}}%
\pgfpathlineto{\pgfqpoint{0.760589in}{0.984787in}}%
\pgfpathlineto{\pgfqpoint{0.711264in}{0.993071in}}%
\pgfusepath{fill}%
\end{pgfscope}%
\begin{pgfscope}%
\pgfpathrectangle{\pgfqpoint{0.407592in}{0.222904in}}{\pgfqpoint{1.411163in}{1.411163in}}%
\pgfusepath{clip}%
\pgfsetbuttcap%
\pgfsetroundjoin%
\definecolor{currentfill}{rgb}{0.749020,0.772549,0.792157}%
\pgfsetfillcolor{currentfill}%
\pgfsetfillopacity{0.160000}%
\pgfsetlinewidth{0.000000pt}%
\definecolor{currentstroke}{rgb}{0.000000,0.000000,0.000000}%
\pgfsetstrokecolor{currentstroke}%
\pgfsetdash{}{0pt}%
\pgfpathmoveto{\pgfqpoint{0.982633in}{0.936550in}}%
\pgfpathlineto{\pgfqpoint{0.980283in}{0.943670in}}%
\pgfpathlineto{\pgfqpoint{0.977921in}{0.950826in}}%
\pgfpathlineto{\pgfqpoint{1.031145in}{0.949559in}}%
\pgfpathlineto{\pgfqpoint{1.084665in}{0.949700in}}%
\pgfpathlineto{\pgfqpoint{1.084782in}{0.942565in}}%
\pgfpathlineto{\pgfqpoint{1.084898in}{0.935467in}}%
\pgfpathlineto{\pgfqpoint{1.033625in}{0.935330in}}%
\pgfpathlineto{\pgfqpoint{0.982633in}{0.936550in}}%
\pgfpathclose%
\pgfusepath{fill}%
\end{pgfscope}%
\begin{pgfscope}%
\pgfpathrectangle{\pgfqpoint{0.407592in}{0.222904in}}{\pgfqpoint{1.411163in}{1.411163in}}%
\pgfusepath{clip}%
\pgfsetbuttcap%
\pgfsetroundjoin%
\definecolor{currentfill}{rgb}{0.749020,0.772549,0.792157}%
\pgfsetfillcolor{currentfill}%
\pgfsetfillopacity{0.160000}%
\pgfsetlinewidth{0.000000pt}%
\definecolor{currentstroke}{rgb}{0.000000,0.000000,0.000000}%
\pgfsetstrokecolor{currentstroke}%
\pgfsetdash{}{0pt}%
\pgfpathmoveto{\pgfqpoint{0.773915in}{0.969658in}}%
\pgfpathlineto{\pgfqpoint{0.767268in}{0.977204in}}%
\pgfpathlineto{\pgfqpoint{0.760589in}{0.984787in}}%
\pgfpathlineto{\pgfqpoint{0.811726in}{0.977820in}}%
\pgfpathlineto{\pgfqpoint{0.864397in}{0.972217in}}%
\pgfpathlineto{\pgfqpoint{0.868990in}{0.964860in}}%
\pgfpathlineto{\pgfqpoint{0.873560in}{0.957540in}}%
\pgfpathlineto{\pgfqpoint{0.823008in}{0.962941in}}%
\pgfpathlineto{\pgfqpoint{0.773915in}{0.969658in}}%
\pgfpathclose%
\pgfusepath{fill}%
\end{pgfscope}%
\begin{pgfscope}%
\pgfpathrectangle{\pgfqpoint{0.407592in}{0.222904in}}{\pgfqpoint{1.411163in}{1.411163in}}%
\pgfusepath{clip}%
\pgfsetbuttcap%
\pgfsetroundjoin%
\definecolor{currentfill}{rgb}{0.749020,0.772549,0.792157}%
\pgfsetfillcolor{currentfill}%
\pgfsetfillopacity{0.160000}%
\pgfsetlinewidth{0.000000pt}%
\definecolor{currentstroke}{rgb}{0.000000,0.000000,0.000000}%
\pgfsetstrokecolor{currentstroke}%
\pgfsetdash{}{0pt}%
\pgfpathmoveto{\pgfqpoint{1.304313in}{0.979110in}}%
\pgfpathlineto{\pgfqpoint{1.308680in}{0.986628in}}%
\pgfpathlineto{\pgfqpoint{1.362338in}{0.993994in}}%
\pgfpathlineto{\pgfqpoint{1.414415in}{1.002688in}}%
\pgfpathlineto{\pgfqpoint{1.407967in}{0.994887in}}%
\pgfpathlineto{\pgfqpoint{1.356910in}{0.986345in}}%
\pgfusepath{fill}%
\end{pgfscope}%
\begin{pgfscope}%
\pgfpathrectangle{\pgfqpoint{0.407592in}{0.222904in}}{\pgfqpoint{1.411163in}{1.411163in}}%
\pgfusepath{clip}%
\pgfsetbuttcap%
\pgfsetroundjoin%
\definecolor{currentfill}{rgb}{0.749020,0.772549,0.792157}%
\pgfsetfillcolor{currentfill}%
\pgfsetfillopacity{0.160000}%
\pgfsetlinewidth{0.000000pt}%
\definecolor{currentstroke}{rgb}{0.000000,0.000000,0.000000}%
\pgfsetstrokecolor{currentstroke}%
\pgfsetdash{}{0pt}%
\pgfpathmoveto{\pgfqpoint{1.191375in}{0.954195in}}%
\pgfpathlineto{\pgfqpoint{1.193510in}{0.961451in}}%
\pgfpathlineto{\pgfqpoint{1.195657in}{0.968745in}}%
\pgfpathlineto{\pgfqpoint{1.250462in}{0.973230in}}%
\pgfpathlineto{\pgfqpoint{1.304313in}{0.979110in}}%
\pgfpathlineto{\pgfqpoint{1.299967in}{0.971629in}}%
\pgfpathlineto{\pgfqpoint{1.295644in}{0.964185in}}%
\pgfpathlineto{\pgfqpoint{1.243962in}{0.958517in}}%
\pgfpathlineto{\pgfqpoint{1.191375in}{0.954195in}}%
\pgfpathclose%
\pgfusepath{fill}%
\end{pgfscope}%
\begin{pgfscope}%
\pgfpathrectangle{\pgfqpoint{0.407592in}{0.222904in}}{\pgfqpoint{1.411163in}{1.411163in}}%
\pgfusepath{clip}%
\pgfsetbuttcap%
\pgfsetroundjoin%
\definecolor{currentfill}{rgb}{0.749020,0.772549,0.792157}%
\pgfsetfillcolor{currentfill}%
\pgfsetfillopacity{0.160000}%
\pgfsetlinewidth{0.000000pt}%
\definecolor{currentstroke}{rgb}{0.000000,0.000000,0.000000}%
\pgfsetstrokecolor{currentstroke}%
\pgfsetdash{}{0pt}%
\pgfpathmoveto{\pgfqpoint{0.873560in}{0.957540in}}%
\pgfpathlineto{\pgfqpoint{0.868990in}{0.964860in}}%
\pgfpathlineto{\pgfqpoint{0.864397in}{0.972217in}}%
\pgfpathlineto{\pgfqpoint{0.918309in}{0.968017in}}%
\pgfpathlineto{\pgfqpoint{0.973161in}{0.965250in}}%
\pgfpathlineto{\pgfqpoint{0.975547in}{0.958019in}}%
\pgfpathlineto{\pgfqpoint{0.977921in}{0.950826in}}%
\pgfpathlineto{\pgfqpoint{0.925294in}{0.953493in}}%
\pgfpathlineto{\pgfqpoint{0.873560in}{0.957540in}}%
\pgfpathclose%
\pgfusepath{fill}%
\end{pgfscope}%
\begin{pgfscope}%
\pgfpathrectangle{\pgfqpoint{0.407592in}{0.222904in}}{\pgfqpoint{1.411163in}{1.411163in}}%
\pgfusepath{clip}%
\pgfsetbuttcap%
\pgfsetroundjoin%
\definecolor{currentfill}{rgb}{0.749020,0.772549,0.792157}%
\pgfsetfillcolor{currentfill}%
\pgfsetfillopacity{0.160000}%
\pgfsetlinewidth{0.000000pt}%
\definecolor{currentstroke}{rgb}{0.000000,0.000000,0.000000}%
\pgfsetstrokecolor{currentstroke}%
\pgfsetdash{}{0pt}%
\pgfpathmoveto{\pgfqpoint{1.084665in}{0.949700in}}%
\pgfpathlineto{\pgfqpoint{1.084547in}{0.956872in}}%
\pgfpathlineto{\pgfqpoint{1.084429in}{0.964081in}}%
\pgfpathlineto{\pgfqpoint{1.140207in}{0.965689in}}%
\pgfpathlineto{\pgfqpoint{1.195657in}{0.968745in}}%
\pgfpathlineto{\pgfqpoint{1.193510in}{0.961451in}}%
\pgfpathlineto{\pgfqpoint{1.191375in}{0.954195in}}%
\pgfpathlineto{\pgfqpoint{1.138175in}{0.951249in}}%
\pgfpathlineto{\pgfqpoint{1.084665in}{0.949700in}}%
\pgfpathclose%
\pgfusepath{fill}%
\end{pgfscope}%
\begin{pgfscope}%
\pgfpathrectangle{\pgfqpoint{0.407592in}{0.222904in}}{\pgfqpoint{1.411163in}{1.411163in}}%
\pgfusepath{clip}%
\pgfsetbuttcap%
\pgfsetroundjoin%
\definecolor{currentfill}{rgb}{0.749020,0.772549,0.792157}%
\pgfsetfillcolor{currentfill}%
\pgfsetfillopacity{0.160000}%
\pgfsetlinewidth{0.000000pt}%
\definecolor{currentstroke}{rgb}{0.000000,0.000000,0.000000}%
\pgfsetstrokecolor{currentstroke}%
\pgfsetdash{}{0pt}%
\pgfpathmoveto{\pgfqpoint{0.760589in}{0.984787in}}%
\pgfpathlineto{\pgfqpoint{0.753877in}{0.992408in}}%
\pgfpathlineto{\pgfqpoint{0.806042in}{0.985315in}}%
\pgfpathlineto{\pgfqpoint{0.859780in}{0.979611in}}%
\pgfpathlineto{\pgfqpoint{0.864397in}{0.972217in}}%
\pgfpathlineto{\pgfqpoint{0.811726in}{0.977820in}}%
\pgfusepath{fill}%
\end{pgfscope}%
\begin{pgfscope}%
\pgfpathrectangle{\pgfqpoint{0.407592in}{0.222904in}}{\pgfqpoint{1.411163in}{1.411163in}}%
\pgfusepath{clip}%
\pgfsetbuttcap%
\pgfsetroundjoin%
\definecolor{currentfill}{rgb}{0.749020,0.772549,0.792157}%
\pgfsetfillcolor{currentfill}%
\pgfsetfillopacity{0.160000}%
\pgfsetlinewidth{0.000000pt}%
\definecolor{currentstroke}{rgb}{0.000000,0.000000,0.000000}%
\pgfsetstrokecolor{currentstroke}%
\pgfsetdash{}{0pt}%
\pgfpathmoveto{\pgfqpoint{0.977921in}{0.950826in}}%
\pgfpathlineto{\pgfqpoint{0.975547in}{0.958019in}}%
\pgfpathlineto{\pgfqpoint{0.973161in}{0.965250in}}%
\pgfpathlineto{\pgfqpoint{1.028640in}{0.963935in}}%
\pgfpathlineto{\pgfqpoint{1.084429in}{0.964081in}}%
\pgfpathlineto{\pgfqpoint{1.084547in}{0.956872in}}%
\pgfpathlineto{\pgfqpoint{1.084665in}{0.949700in}}%
\pgfpathlineto{\pgfqpoint{1.031145in}{0.949559in}}%
\pgfpathlineto{\pgfqpoint{0.977921in}{0.950826in}}%
\pgfpathclose%
\pgfusepath{fill}%
\end{pgfscope}%
\begin{pgfscope}%
\pgfpathrectangle{\pgfqpoint{0.407592in}{0.222904in}}{\pgfqpoint{1.411163in}{1.411163in}}%
\pgfusepath{clip}%
\pgfsetbuttcap%
\pgfsetroundjoin%
\definecolor{currentfill}{rgb}{0.749020,0.772549,0.792157}%
\pgfsetfillcolor{currentfill}%
\pgfsetfillopacity{0.160000}%
\pgfsetlinewidth{0.000000pt}%
\definecolor{currentstroke}{rgb}{0.000000,0.000000,0.000000}%
\pgfsetstrokecolor{currentstroke}%
\pgfsetdash{}{0pt}%
\pgfpathmoveto{\pgfqpoint{1.195657in}{0.968745in}}%
\pgfpathlineto{\pgfqpoint{1.197815in}{0.976077in}}%
\pgfpathlineto{\pgfqpoint{1.253736in}{0.980642in}}%
\pgfpathlineto{\pgfqpoint{1.308680in}{0.986628in}}%
\pgfpathlineto{\pgfqpoint{1.304313in}{0.979110in}}%
\pgfpathlineto{\pgfqpoint{1.250462in}{0.973230in}}%
\pgfusepath{fill}%
\end{pgfscope}%
\begin{pgfscope}%
\pgfpathrectangle{\pgfqpoint{0.407592in}{0.222904in}}{\pgfqpoint{1.411163in}{1.411163in}}%
\pgfusepath{clip}%
\pgfsetbuttcap%
\pgfsetroundjoin%
\definecolor{currentfill}{rgb}{0.749020,0.772549,0.792157}%
\pgfsetfillcolor{currentfill}%
\pgfsetfillopacity{0.160000}%
\pgfsetlinewidth{0.000000pt}%
\definecolor{currentstroke}{rgb}{0.000000,0.000000,0.000000}%
\pgfsetstrokecolor{currentstroke}%
\pgfsetdash{}{0pt}%
\pgfpathmoveto{\pgfqpoint{0.864397in}{0.972217in}}%
\pgfpathlineto{\pgfqpoint{0.859780in}{0.979611in}}%
\pgfpathlineto{\pgfqpoint{0.914790in}{0.975335in}}%
\pgfpathlineto{\pgfqpoint{0.970763in}{0.972517in}}%
\pgfpathlineto{\pgfqpoint{0.973161in}{0.965250in}}%
\pgfpathlineto{\pgfqpoint{0.918309in}{0.968017in}}%
\pgfusepath{fill}%
\end{pgfscope}%
\begin{pgfscope}%
\pgfpathrectangle{\pgfqpoint{0.407592in}{0.222904in}}{\pgfqpoint{1.411163in}{1.411163in}}%
\pgfusepath{clip}%
\pgfsetbuttcap%
\pgfsetroundjoin%
\definecolor{currentfill}{rgb}{0.749020,0.772549,0.792157}%
\pgfsetfillcolor{currentfill}%
\pgfsetfillopacity{0.160000}%
\pgfsetlinewidth{0.000000pt}%
\definecolor{currentstroke}{rgb}{0.000000,0.000000,0.000000}%
\pgfsetstrokecolor{currentstroke}%
\pgfsetdash{}{0pt}%
\pgfpathmoveto{\pgfqpoint{1.084429in}{0.964081in}}%
\pgfpathlineto{\pgfqpoint{1.084310in}{0.971328in}}%
\pgfpathlineto{\pgfqpoint{1.141231in}{0.972965in}}%
\pgfpathlineto{\pgfqpoint{1.197815in}{0.976077in}}%
\pgfpathlineto{\pgfqpoint{1.195657in}{0.968745in}}%
\pgfpathlineto{\pgfqpoint{1.140207in}{0.965689in}}%
\pgfusepath{fill}%
\end{pgfscope}%
\begin{pgfscope}%
\pgfpathrectangle{\pgfqpoint{0.407592in}{0.222904in}}{\pgfqpoint{1.411163in}{1.411163in}}%
\pgfusepath{clip}%
\pgfsetbuttcap%
\pgfsetroundjoin%
\definecolor{currentfill}{rgb}{0.749020,0.772549,0.792157}%
\pgfsetfillcolor{currentfill}%
\pgfsetfillopacity{0.160000}%
\pgfsetlinewidth{0.000000pt}%
\definecolor{currentstroke}{rgb}{0.000000,0.000000,0.000000}%
\pgfsetstrokecolor{currentstroke}%
\pgfsetdash{}{0pt}%
\pgfpathmoveto{\pgfqpoint{0.973161in}{0.965250in}}%
\pgfpathlineto{\pgfqpoint{0.970763in}{0.972517in}}%
\pgfpathlineto{\pgfqpoint{1.027378in}{0.971179in}}%
\pgfpathlineto{\pgfqpoint{1.084310in}{0.971328in}}%
\pgfpathlineto{\pgfqpoint{1.084429in}{0.964081in}}%
\pgfpathlineto{\pgfqpoint{1.028640in}{0.963935in}}%
\pgfusepath{fill}%
\end{pgfscope}%
\end{pgfpicture}%
\makeatother%
\endgroup%
}
  \captionsetup{font=small}
  \caption{\textbf{Bures trajectories in the covariance cone.}
  (a) Equal diagonal blocks: invariant hyperbolas.
  (b) Unequal blocks: fixed-\(b\) sections of determinant leaves.
  The displayed trajectories converge to \(r=\beta\).}
  \label{fig:gaussian-w2-cone-geometry}
\end{wrapfigure}
Under special conditions on the conserved spectrum of $J\Sigma$, elimination of hidden blocks becomes possible: Appendix~\ref{app:intrinsic} describes these intrinsic dynamics.

\paragraph{Numerical illustration: the $2\times2$ covariance cone}
For scalar input and output, every covariance has the form
$\Sigma=\begin{psmallmatrix}p&r\\r&q\end{psmallmatrix}\succeq0$.
Take the loss $F(\Sigma):=(r-\beta)^2/2$, with scalar teacher $\beta$.
Writing $(a,b):=((p+q)/2,(p-q)/2)$ identifies $\Sigma\succeq0$ with
the Lorentz cone $a^2\ge r^2+b^2$, $a\ge0$. The Bures flow
\eqref{eq:block-flow-general} becomes
$(\dot a,\dot b,\dot r)=-(2r(r-\beta),0,2a(r-\beta))$,
preserving $b=b_0$ and $\alpha^2:=a^2-r^2-b^2=pq-r^2$.

When $p=q$, set $s:=p$. Then $b=0$ and trajectories follow hyperbolas
$s^2-r^2=\alpha^2$ in $s\ge|r|$ (Figure~\ref{fig:gaussian-w2-cone-geometry}(a)).
For $\alpha>0$, they converge to $(r,s)=(\beta,\sqrt{\beta^2+\alpha^2})$;
the spectral-margin bound is exact for the minimum $s$ on the leaf.
Unequal diagonal blocks give intersections of the planes $b=b_0$ with
determinant leaves in the three-dimensional cone (panel (b)).
If $b_0^2+\alpha^2>0$, the endpoint is
$(r_\infty,a_\infty)=(\beta,\sqrt{\beta^2+b_0^2+\alpha^2})$.

The invariants $2b$ and $-\alpha^2$ are the trace and determinant of
$J\Sigma$, whose eigenvalues are $b_0\pm\sqrt{b_0^2+\alpha^2}$.
Thus $\gap{\Sigma_0}=2\sqrt{b_0^2+\alpha^2}$: the convergence bound
covers both panels, including singular leaves with $b_0\ne0$.
Remark~\ref{rem:degenerate-boundary-leaves} in Appendix~\ref{app:gf}
treats the zero-support-gap boundary. Figure~\ref{fig:gaussian-w2-cone-geometry}
is obtained by integrating $\dot r=-2\sqrt{r^2+b_0^2+\alpha^2}(r-\beta)$.

%%% Local Variables:
%%% mode: LaTeX
%%% TeX-master: "../main"
%%% End:

\section{Linear Convergence of Training Dynamics}
\label{sec:convergence}\label{sec:balanced-convergence}
\label{sec:sampling}\label{sec:finite-width}

\paragraph{Convergence of gradient flow}
A positive conserved support gap $g(\Sigma_0) > 0$ guarantees global linear convergence for the Bures flow \eqref{eq:bures-G}, including some cases where the initial covariance $\Sigma_0$ is singular. Write $\Delta_t:=\ell(\Svu_t)-\ell_\star$ for the loss gap and $\cS_\star$ for the set of global predictor minimizers.

\begin{theorem}[Global convergence from a positive support gap]\label{thm:gf}\label{thm:balanced-convergence}
Under Assumption~\ref{ass:loss}, let $\Sigma_0\succeq0$ be a possibly
singular covariance with $\gap{\Sigma_0}>0$. 
Then the solution of \eqref{eq:bures-G} is global and satisfies
\begin{equation}\label{eq:gf-rate}
 \Delta_t\le\Delta_0e^{-2\kappa \gap{\Sigma_0}t},\qquad t\ge0.
\end{equation}
As $t \to \infty$, $\Sigma_t \to \Sigma_\infty\succeq0$ with
$\rankop\Sigma_\infty=\rankop\Sigma_0$ and $\Svu_\infty\in\cS_\star$.
The predictor $\Svu$ and any factor realization (as per \eqref{eq:mean-field-gradient-flow}) converge linearly with rate $\kappa\gap{\Sigma_0}$ in Frobenius norm.
\end{theorem}

The only nondegeneracy requirement is $\gap{\Sigma_0}>0$.
Positive definiteness is sufficient but not necessary: $\Sigma_0\succ0$ implies $\gap{\Sigma_0}\ge2\lambda_{\min}(\Sigma_0)>0$ and then
$\Sigma_\infty\succ0$. The exponent separates loss geometry (captured by the PL constant $\kappa$)
from initial covariance conditioning (the support gap $\gap{\Sigma_0}$), with no extra
dimension factor. The trajectory selects a global minimizer according to its initialization and dynamics. Proposition~\ref{prop:gf-lift-tail} supplies all explicit prefactors and bounds.

\begin{comment}
\emph{Proof sketch.}
\emph{Conserved conditioning.}
Proposition~\ref{lem:margins} gives
$\Suu_t\succeq \aplus{\Sigma_0}\Id$ and $\Svv_t\succeq \aminus{\Sigma_0}\Id$ throughout the
existence interval. The predictor equation
$\dot\Svu_t=-E_t\Suu_t-\Svv_tE_t$ therefore has a uniformly positive
preconditioner, with lower bound $\gap{\Sigma_0}$.
Only the sum of the two margins must be positive; either one may vanish.

\emph{Loss decay.}
Pairing this equation with $E_t=\nabla\ell(\Svu_t)$ gives
$-\dot\Delta_t\ge \gap{\Sigma_0}\norm{E_t}_{\F}^2$.
The PL inequality then yields
$-\dot\Delta_t\ge2\kappa \gap{\Sigma_0}\Delta_t$, proving \eqref{eq:gf-rate}
on the existence interval.

\emph{Global existence and convergence of the state.}
The same dissipation estimate applied to $\sqrt{\Delta_t}$ bounds
$\int_0^T\norm{G_t}_{\op}\,\dd t$ by
$\gap{\Sigma_0}^{-1}\sqrt{2\Delta_0/\kappa}$, uniformly over finite existence
intervals. Hence the fundamental matrix and its inverse remain bounded.
Covariance congruence prevents blow-up, so the flow is global. The
fundamental matrix converges to an invertible limit, preserving covariance
rank. Applying the length bound from time $t$ onward gives
exponent $\kappa \gap{\Sigma_0}$ for the remaining displacement and hence convergence
of the covariance and factors.
Appendix~\ref{app:gf} gives the full proof.\qed
\end{comment}

\begin{remark}[Singular covariance and predictor rank]\label{rem:singular-support}
The initialization $\Sigma_0=\diag(\Id_{d_u},0_{d_v})$ has
$\aplus{\Sigma_0}=1$, $\aminus{\Sigma_0}=0$, and $\gap{\Sigma_0}=1$.
Its rank is only $d_u$, yet the theorem applies. There is no rank
contradiction: a full-rank predictor needs rank $\min(d_u,d_v)$, not $m$.
More generally, $\aplus{\Sigma_0}>0$ implies $\rankop\Sigma_0\ge d_u$,
while $\aminus{\Sigma_0}>0$ implies $\rankop\Sigma_0\ge d_v$.
The condition cannot be dropped altogether: for $d_u=d_v=2$,
$\Sigma_0=\diag(1,0,1,0)$ and
$\ell(S):=\norm{S-e_2e_2^\top}_{\F}^2/2$, the flow is stationary and
nonoptimal. Here $g(\Sigma_0) = 0$, despite the nonzero signed eigenvalues $+1,-1$.
\end{remark}

\begin{remark}[Quadratic specialization]
For $\ell(S):=\frac12\ip{\quadH[S-\Svu_\star]}{S-\Svu_\star}$ with
$\quadH\succeq0$ a nonzero self-adjoint linear operator, the exponent is
$2\gap{\Sigma_0}\lambda_{\min}^{+}(\quadH)$. For a full-rank isotropic initialization
$\Sigma_0=\sigma^2\Id_m$, one has $\aplus{\Sigma_0}=\aminus{\Sigma_0}=\sigma^2$,
giving a rate of $4\sigma^2\lambda_{\min}^{+}(\quadH)$.
The operator $\quadH$ is allowed to have a nontrivial kernel: the predictor $\Svu$ converges to a point in
$\Svu_\star+\kerop\quadH$. If $\quadH$ is positive definite, the predictor converges to the teacher $\Svu_\star$.
\end{remark}

For a quadratic predictor loss, Figure~\ref{fig:bures-convergence-triptych} compares the initialization scale, loss of conditioning, and the finite-width sampling we discuss next. The dotted curves use a coarse determinant-based rate, derived for comparison in Appendix~\ref{app:gf}.
If the support gap $g(\Sigma_0)$ vanishes, trace and initial loss alone cannot guarantee a uniform positive rate (Example~\ref{ex:slow}).
Figure~\ref{fig:smooth-lasso-triptych} in Appendix~\ref{app:smooth-lasso} provides the counterpart for a non-quadratic smoothed-Lasso predictor loss.

\begin{figure}[t]
    \centering
    \resizebox{\linewidth}{!}{%% Creator: Matplotlib, PGF backend
%%
%% To include the figure in your LaTeX document, write
%%   \input{<filename>.pgf}
%%
%% Make sure the required packages are loaded in your preamble
%%   \usepackage{pgf}
%%
%% Also ensure that all the required font packages are loaded; for instance,
%% the lmodern package is sometimes necessary when using math font.
%%   \usepackage{lmodern}
%%
%% Figures using additional raster images can only be included by \input if
%% they are in the same directory as the main LaTeX file. For loading figures
%% from other directories you can use the `import` package
%%   \usepackage{import}
%%
%% and then include the figures with
%%   \import{<path to file>}{<filename>.pgf}
%%
%% Matplotlib used the following preamble
%%   \def\mathdefault#1{#1}
%%   \everymath=\expandafter{\the\everymath\displaystyle}
%%   \IfFileExists{scrextend.sty}{
%%     \usepackage[fontsize=6.000000pt]{scrextend}
%%   }{
%%     \renewcommand{\normalsize}{\fontsize{6.000000}{7.200000}\selectfont}
%%     \normalsize
%%   }
%%   \usepackage{amsfonts}
%%   \usepackage{amssymb}
%%   \ifdefined\pdftexversion\else  % non-pdftex case.
%%     \usepackage{fontspec}
%%   \fi
%%   \makeatletter\@ifpackageloaded{underscore}{}{\usepackage[strings]{underscore}}\makeatother
%%
\begingroup%
\makeatletter%
\begin{pgfpicture}%
\pgfpathrectangle{\pgfpointorigin}{\pgfqpoint{6.500000in}{1.687500in}}%
\pgfusepath{use as bounding box, clip}%
\begin{pgfscope}%
\pgfsetbuttcap%
\pgfsetmiterjoin%
\definecolor{currentfill}{rgb}{1.000000,1.000000,1.000000}%
\pgfsetfillcolor{currentfill}%
\pgfsetlinewidth{0.000000pt}%
\definecolor{currentstroke}{rgb}{1.000000,1.000000,1.000000}%
\pgfsetstrokecolor{currentstroke}%
\pgfsetdash{}{0pt}%
\pgfpathmoveto{\pgfqpoint{0.000000in}{0.000000in}}%
\pgfpathlineto{\pgfqpoint{6.500000in}{0.000000in}}%
\pgfpathlineto{\pgfqpoint{6.500000in}{1.687500in}}%
\pgfpathlineto{\pgfqpoint{0.000000in}{1.687500in}}%
\pgfpathlineto{\pgfqpoint{0.000000in}{0.000000in}}%
\pgfpathclose%
\pgfusepath{fill}%
\end{pgfscope}%
\begin{pgfscope}%
\pgfsetbuttcap%
\pgfsetmiterjoin%
\definecolor{currentfill}{rgb}{1.000000,1.000000,1.000000}%
\pgfsetfillcolor{currentfill}%
\pgfsetlinewidth{0.000000pt}%
\definecolor{currentstroke}{rgb}{0.000000,0.000000,0.000000}%
\pgfsetstrokecolor{currentstroke}%
\pgfsetstrokeopacity{0.000000}%
\pgfsetdash{}{0pt}%
\pgfpathmoveto{\pgfqpoint{0.487500in}{0.320625in}}%
\pgfpathlineto{\pgfqpoint{2.333179in}{0.320625in}}%
\pgfpathlineto{\pgfqpoint{2.333179in}{1.485000in}}%
\pgfpathlineto{\pgfqpoint{0.487500in}{1.485000in}}%
\pgfpathlineto{\pgfqpoint{0.487500in}{0.320625in}}%
\pgfpathclose%
\pgfusepath{fill}%
\end{pgfscope}%
\begin{pgfscope}%
\pgfpathrectangle{\pgfqpoint{0.487500in}{0.320625in}}{\pgfqpoint{1.845679in}{1.164375in}}%
\pgfusepath{clip}%
\pgfsetrectcap%
\pgfsetroundjoin%
\pgfsetlinewidth{0.803000pt}%
\definecolor{currentstroke}{rgb}{0.909804,0.909804,0.909804}%
\pgfsetstrokecolor{currentstroke}%
\pgfsetdash{}{0pt}%
\pgfpathmoveto{\pgfqpoint{0.571395in}{0.320625in}}%
\pgfpathlineto{\pgfqpoint{0.571395in}{1.485000in}}%
\pgfusepath{stroke}%
\end{pgfscope}%
\begin{pgfscope}%
\pgfsetbuttcap%
\pgfsetroundjoin%
\definecolor{currentfill}{rgb}{0.000000,0.000000,0.000000}%
\pgfsetfillcolor{currentfill}%
\pgfsetlinewidth{0.803000pt}%
\definecolor{currentstroke}{rgb}{0.000000,0.000000,0.000000}%
\pgfsetstrokecolor{currentstroke}%
\pgfsetdash{}{0pt}%
\pgfsys@defobject{currentmarker}{\pgfqpoint{0.000000in}{-0.048611in}}{\pgfqpoint{0.000000in}{0.000000in}}{%
\pgfpathmoveto{\pgfqpoint{0.000000in}{0.000000in}}%
\pgfpathlineto{\pgfqpoint{0.000000in}{-0.048611in}}%
\pgfusepath{stroke,fill}%
}%
\begin{pgfscope}%
\pgfsys@transformshift{0.571395in}{0.320625in}%
\pgfsys@useobject{currentmarker}{}%
\end{pgfscope}%
\end{pgfscope}%
\begin{pgfscope}%
\definecolor{textcolor}{rgb}{0.000000,0.000000,0.000000}%
\pgfsetstrokecolor{textcolor}%
\pgfsetfillcolor{textcolor}%
\pgftext[x=0.571395in,y=0.223403in,,top]{\color{textcolor}{\rmfamily\fontsize{6.000000}{7.200000}\selectfont\catcode`\^=\active\def^{\ifmmode\sp\else\^{}\fi}\catcode`\%=\active\def%{\%}0}}%
\end{pgfscope}%
\begin{pgfscope}%
\pgfpathrectangle{\pgfqpoint{0.487500in}{0.320625in}}{\pgfqpoint{1.845679in}{1.164375in}}%
\pgfusepath{clip}%
\pgfsetrectcap%
\pgfsetroundjoin%
\pgfsetlinewidth{0.803000pt}%
\definecolor{currentstroke}{rgb}{0.909804,0.909804,0.909804}%
\pgfsetstrokecolor{currentstroke}%
\pgfsetdash{}{0pt}%
\pgfpathmoveto{\pgfqpoint{0.990867in}{0.320625in}}%
\pgfpathlineto{\pgfqpoint{0.990867in}{1.485000in}}%
\pgfusepath{stroke}%
\end{pgfscope}%
\begin{pgfscope}%
\pgfsetbuttcap%
\pgfsetroundjoin%
\definecolor{currentfill}{rgb}{0.000000,0.000000,0.000000}%
\pgfsetfillcolor{currentfill}%
\pgfsetlinewidth{0.803000pt}%
\definecolor{currentstroke}{rgb}{0.000000,0.000000,0.000000}%
\pgfsetstrokecolor{currentstroke}%
\pgfsetdash{}{0pt}%
\pgfsys@defobject{currentmarker}{\pgfqpoint{0.000000in}{-0.048611in}}{\pgfqpoint{0.000000in}{0.000000in}}{%
\pgfpathmoveto{\pgfqpoint{0.000000in}{0.000000in}}%
\pgfpathlineto{\pgfqpoint{0.000000in}{-0.048611in}}%
\pgfusepath{stroke,fill}%
}%
\begin{pgfscope}%
\pgfsys@transformshift{0.990867in}{0.320625in}%
\pgfsys@useobject{currentmarker}{}%
\end{pgfscope}%
\end{pgfscope}%
\begin{pgfscope}%
\definecolor{textcolor}{rgb}{0.000000,0.000000,0.000000}%
\pgfsetstrokecolor{textcolor}%
\pgfsetfillcolor{textcolor}%
\pgftext[x=0.990867in,y=0.223403in,,top]{\color{textcolor}{\rmfamily\fontsize{6.000000}{7.200000}\selectfont\catcode`\^=\active\def^{\ifmmode\sp\else\^{}\fi}\catcode`\%=\active\def%{\%}10}}%
\end{pgfscope}%
\begin{pgfscope}%
\pgfpathrectangle{\pgfqpoint{0.487500in}{0.320625in}}{\pgfqpoint{1.845679in}{1.164375in}}%
\pgfusepath{clip}%
\pgfsetrectcap%
\pgfsetroundjoin%
\pgfsetlinewidth{0.803000pt}%
\definecolor{currentstroke}{rgb}{0.909804,0.909804,0.909804}%
\pgfsetstrokecolor{currentstroke}%
\pgfsetdash{}{0pt}%
\pgfpathmoveto{\pgfqpoint{1.410340in}{0.320625in}}%
\pgfpathlineto{\pgfqpoint{1.410340in}{1.485000in}}%
\pgfusepath{stroke}%
\end{pgfscope}%
\begin{pgfscope}%
\pgfsetbuttcap%
\pgfsetroundjoin%
\definecolor{currentfill}{rgb}{0.000000,0.000000,0.000000}%
\pgfsetfillcolor{currentfill}%
\pgfsetlinewidth{0.803000pt}%
\definecolor{currentstroke}{rgb}{0.000000,0.000000,0.000000}%
\pgfsetstrokecolor{currentstroke}%
\pgfsetdash{}{0pt}%
\pgfsys@defobject{currentmarker}{\pgfqpoint{0.000000in}{-0.048611in}}{\pgfqpoint{0.000000in}{0.000000in}}{%
\pgfpathmoveto{\pgfqpoint{0.000000in}{0.000000in}}%
\pgfpathlineto{\pgfqpoint{0.000000in}{-0.048611in}}%
\pgfusepath{stroke,fill}%
}%
\begin{pgfscope}%
\pgfsys@transformshift{1.410340in}{0.320625in}%
\pgfsys@useobject{currentmarker}{}%
\end{pgfscope}%
\end{pgfscope}%
\begin{pgfscope}%
\definecolor{textcolor}{rgb}{0.000000,0.000000,0.000000}%
\pgfsetstrokecolor{textcolor}%
\pgfsetfillcolor{textcolor}%
\pgftext[x=1.410340in,y=0.223403in,,top]{\color{textcolor}{\rmfamily\fontsize{6.000000}{7.200000}\selectfont\catcode`\^=\active\def^{\ifmmode\sp\else\^{}\fi}\catcode`\%=\active\def%{\%}20}}%
\end{pgfscope}%
\begin{pgfscope}%
\pgfpathrectangle{\pgfqpoint{0.487500in}{0.320625in}}{\pgfqpoint{1.845679in}{1.164375in}}%
\pgfusepath{clip}%
\pgfsetrectcap%
\pgfsetroundjoin%
\pgfsetlinewidth{0.803000pt}%
\definecolor{currentstroke}{rgb}{0.909804,0.909804,0.909804}%
\pgfsetstrokecolor{currentstroke}%
\pgfsetdash{}{0pt}%
\pgfpathmoveto{\pgfqpoint{1.829812in}{0.320625in}}%
\pgfpathlineto{\pgfqpoint{1.829812in}{1.485000in}}%
\pgfusepath{stroke}%
\end{pgfscope}%
\begin{pgfscope}%
\pgfsetbuttcap%
\pgfsetroundjoin%
\definecolor{currentfill}{rgb}{0.000000,0.000000,0.000000}%
\pgfsetfillcolor{currentfill}%
\pgfsetlinewidth{0.803000pt}%
\definecolor{currentstroke}{rgb}{0.000000,0.000000,0.000000}%
\pgfsetstrokecolor{currentstroke}%
\pgfsetdash{}{0pt}%
\pgfsys@defobject{currentmarker}{\pgfqpoint{0.000000in}{-0.048611in}}{\pgfqpoint{0.000000in}{0.000000in}}{%
\pgfpathmoveto{\pgfqpoint{0.000000in}{0.000000in}}%
\pgfpathlineto{\pgfqpoint{0.000000in}{-0.048611in}}%
\pgfusepath{stroke,fill}%
}%
\begin{pgfscope}%
\pgfsys@transformshift{1.829812in}{0.320625in}%
\pgfsys@useobject{currentmarker}{}%
\end{pgfscope}%
\end{pgfscope}%
\begin{pgfscope}%
\definecolor{textcolor}{rgb}{0.000000,0.000000,0.000000}%
\pgfsetstrokecolor{textcolor}%
\pgfsetfillcolor{textcolor}%
\pgftext[x=1.829812in,y=0.223403in,,top]{\color{textcolor}{\rmfamily\fontsize{6.000000}{7.200000}\selectfont\catcode`\^=\active\def^{\ifmmode\sp\else\^{}\fi}\catcode`\%=\active\def%{\%}30}}%
\end{pgfscope}%
\begin{pgfscope}%
\pgfpathrectangle{\pgfqpoint{0.487500in}{0.320625in}}{\pgfqpoint{1.845679in}{1.164375in}}%
\pgfusepath{clip}%
\pgfsetrectcap%
\pgfsetroundjoin%
\pgfsetlinewidth{0.803000pt}%
\definecolor{currentstroke}{rgb}{0.909804,0.909804,0.909804}%
\pgfsetstrokecolor{currentstroke}%
\pgfsetdash{}{0pt}%
\pgfpathmoveto{\pgfqpoint{2.249285in}{0.320625in}}%
\pgfpathlineto{\pgfqpoint{2.249285in}{1.485000in}}%
\pgfusepath{stroke}%
\end{pgfscope}%
\begin{pgfscope}%
\pgfsetbuttcap%
\pgfsetroundjoin%
\definecolor{currentfill}{rgb}{0.000000,0.000000,0.000000}%
\pgfsetfillcolor{currentfill}%
\pgfsetlinewidth{0.803000pt}%
\definecolor{currentstroke}{rgb}{0.000000,0.000000,0.000000}%
\pgfsetstrokecolor{currentstroke}%
\pgfsetdash{}{0pt}%
\pgfsys@defobject{currentmarker}{\pgfqpoint{0.000000in}{-0.048611in}}{\pgfqpoint{0.000000in}{0.000000in}}{%
\pgfpathmoveto{\pgfqpoint{0.000000in}{0.000000in}}%
\pgfpathlineto{\pgfqpoint{0.000000in}{-0.048611in}}%
\pgfusepath{stroke,fill}%
}%
\begin{pgfscope}%
\pgfsys@transformshift{2.249285in}{0.320625in}%
\pgfsys@useobject{currentmarker}{}%
\end{pgfscope}%
\end{pgfscope}%
\begin{pgfscope}%
\definecolor{textcolor}{rgb}{0.000000,0.000000,0.000000}%
\pgfsetstrokecolor{textcolor}%
\pgfsetfillcolor{textcolor}%
\pgftext[x=2.249285in,y=0.223403in,,top]{\color{textcolor}{\rmfamily\fontsize{6.000000}{7.200000}\selectfont\catcode`\^=\active\def^{\ifmmode\sp\else\^{}\fi}\catcode`\%=\active\def%{\%}40}}%
\end{pgfscope}%
\begin{pgfscope}%
\definecolor{textcolor}{rgb}{0.000000,0.000000,0.000000}%
\pgfsetstrokecolor{textcolor}%
\pgfsetfillcolor{textcolor}%
\pgftext[x=1.410340in,y=0.087200in,,top]{\color{textcolor}{\rmfamily\fontsize{6.000000}{7.200000}\selectfont\catcode`\^=\active\def^{\ifmmode\sp\else\^{}\fi}\catcode`\%=\active\def%{\%}time}}%
\end{pgfscope}%
\begin{pgfscope}%
\pgfpathrectangle{\pgfqpoint{0.487500in}{0.320625in}}{\pgfqpoint{1.845679in}{1.164375in}}%
\pgfusepath{clip}%
\pgfsetrectcap%
\pgfsetroundjoin%
\pgfsetlinewidth{0.803000pt}%
\definecolor{currentstroke}{rgb}{0.909804,0.909804,0.909804}%
\pgfsetstrokecolor{currentstroke}%
\pgfsetdash{}{0pt}%
\pgfpathmoveto{\pgfqpoint{0.487500in}{0.320625in}}%
\pgfpathlineto{\pgfqpoint{2.333179in}{0.320625in}}%
\pgfusepath{stroke}%
\end{pgfscope}%
\begin{pgfscope}%
\pgfsetbuttcap%
\pgfsetroundjoin%
\definecolor{currentfill}{rgb}{0.000000,0.000000,0.000000}%
\pgfsetfillcolor{currentfill}%
\pgfsetlinewidth{0.803000pt}%
\definecolor{currentstroke}{rgb}{0.000000,0.000000,0.000000}%
\pgfsetstrokecolor{currentstroke}%
\pgfsetdash{}{0pt}%
\pgfsys@defobject{currentmarker}{\pgfqpoint{-0.048611in}{0.000000in}}{\pgfqpoint{-0.000000in}{0.000000in}}{%
\pgfpathmoveto{\pgfqpoint{-0.000000in}{0.000000in}}%
\pgfpathlineto{\pgfqpoint{-0.048611in}{0.000000in}}%
\pgfusepath{stroke,fill}%
}%
\begin{pgfscope}%
\pgfsys@transformshift{0.487500in}{0.320625in}%
\pgfsys@useobject{currentmarker}{}%
\end{pgfscope}%
\end{pgfscope}%
\begin{pgfscope}%
\definecolor{textcolor}{rgb}{0.000000,0.000000,0.000000}%
\pgfsetstrokecolor{textcolor}%
\pgfsetfillcolor{textcolor}%
\pgftext[x=0.111728in, y=0.285187in, left, base]{\color{textcolor}{\rmfamily\fontsize{6.000000}{7.200000}\selectfont\catcode`\^=\active\def^{\ifmmode\sp\else\^{}\fi}\catcode`\%=\active\def%{\%}$\mathdefault{10^{-26}}$}}%
\end{pgfscope}%
\begin{pgfscope}%
\pgfpathrectangle{\pgfqpoint{0.487500in}{0.320625in}}{\pgfqpoint{1.845679in}{1.164375in}}%
\pgfusepath{clip}%
\pgfsetrectcap%
\pgfsetroundjoin%
\pgfsetlinewidth{0.803000pt}%
\definecolor{currentstroke}{rgb}{0.909804,0.909804,0.909804}%
\pgfsetstrokecolor{currentstroke}%
\pgfsetdash{}{0pt}%
\pgfpathmoveto{\pgfqpoint{0.487500in}{0.536250in}}%
\pgfpathlineto{\pgfqpoint{2.333179in}{0.536250in}}%
\pgfusepath{stroke}%
\end{pgfscope}%
\begin{pgfscope}%
\pgfsetbuttcap%
\pgfsetroundjoin%
\definecolor{currentfill}{rgb}{0.000000,0.000000,0.000000}%
\pgfsetfillcolor{currentfill}%
\pgfsetlinewidth{0.803000pt}%
\definecolor{currentstroke}{rgb}{0.000000,0.000000,0.000000}%
\pgfsetstrokecolor{currentstroke}%
\pgfsetdash{}{0pt}%
\pgfsys@defobject{currentmarker}{\pgfqpoint{-0.048611in}{0.000000in}}{\pgfqpoint{-0.000000in}{0.000000in}}{%
\pgfpathmoveto{\pgfqpoint{-0.000000in}{0.000000in}}%
\pgfpathlineto{\pgfqpoint{-0.048611in}{0.000000in}}%
\pgfusepath{stroke,fill}%
}%
\begin{pgfscope}%
\pgfsys@transformshift{0.487500in}{0.536250in}%
\pgfsys@useobject{currentmarker}{}%
\end{pgfscope}%
\end{pgfscope}%
\begin{pgfscope}%
\definecolor{textcolor}{rgb}{0.000000,0.000000,0.000000}%
\pgfsetstrokecolor{textcolor}%
\pgfsetfillcolor{textcolor}%
\pgftext[x=0.111728in, y=0.500812in, left, base]{\color{textcolor}{\rmfamily\fontsize{6.000000}{7.200000}\selectfont\catcode`\^=\active\def^{\ifmmode\sp\else\^{}\fi}\catcode`\%=\active\def%{\%}$\mathdefault{10^{-21}}$}}%
\end{pgfscope}%
\begin{pgfscope}%
\pgfpathrectangle{\pgfqpoint{0.487500in}{0.320625in}}{\pgfqpoint{1.845679in}{1.164375in}}%
\pgfusepath{clip}%
\pgfsetrectcap%
\pgfsetroundjoin%
\pgfsetlinewidth{0.803000pt}%
\definecolor{currentstroke}{rgb}{0.909804,0.909804,0.909804}%
\pgfsetstrokecolor{currentstroke}%
\pgfsetdash{}{0pt}%
\pgfpathmoveto{\pgfqpoint{0.487500in}{0.751875in}}%
\pgfpathlineto{\pgfqpoint{2.333179in}{0.751875in}}%
\pgfusepath{stroke}%
\end{pgfscope}%
\begin{pgfscope}%
\pgfsetbuttcap%
\pgfsetroundjoin%
\definecolor{currentfill}{rgb}{0.000000,0.000000,0.000000}%
\pgfsetfillcolor{currentfill}%
\pgfsetlinewidth{0.803000pt}%
\definecolor{currentstroke}{rgb}{0.000000,0.000000,0.000000}%
\pgfsetstrokecolor{currentstroke}%
\pgfsetdash{}{0pt}%
\pgfsys@defobject{currentmarker}{\pgfqpoint{-0.048611in}{0.000000in}}{\pgfqpoint{-0.000000in}{0.000000in}}{%
\pgfpathmoveto{\pgfqpoint{-0.000000in}{0.000000in}}%
\pgfpathlineto{\pgfqpoint{-0.048611in}{0.000000in}}%
\pgfusepath{stroke,fill}%
}%
\begin{pgfscope}%
\pgfsys@transformshift{0.487500in}{0.751875in}%
\pgfsys@useobject{currentmarker}{}%
\end{pgfscope}%
\end{pgfscope}%
\begin{pgfscope}%
\definecolor{textcolor}{rgb}{0.000000,0.000000,0.000000}%
\pgfsetstrokecolor{textcolor}%
\pgfsetfillcolor{textcolor}%
\pgftext[x=0.111728in, y=0.716437in, left, base]{\color{textcolor}{\rmfamily\fontsize{6.000000}{7.200000}\selectfont\catcode`\^=\active\def^{\ifmmode\sp\else\^{}\fi}\catcode`\%=\active\def%{\%}$\mathdefault{10^{-16}}$}}%
\end{pgfscope}%
\begin{pgfscope}%
\pgfpathrectangle{\pgfqpoint{0.487500in}{0.320625in}}{\pgfqpoint{1.845679in}{1.164375in}}%
\pgfusepath{clip}%
\pgfsetrectcap%
\pgfsetroundjoin%
\pgfsetlinewidth{0.803000pt}%
\definecolor{currentstroke}{rgb}{0.909804,0.909804,0.909804}%
\pgfsetstrokecolor{currentstroke}%
\pgfsetdash{}{0pt}%
\pgfpathmoveto{\pgfqpoint{0.487500in}{0.967500in}}%
\pgfpathlineto{\pgfqpoint{2.333179in}{0.967500in}}%
\pgfusepath{stroke}%
\end{pgfscope}%
\begin{pgfscope}%
\pgfsetbuttcap%
\pgfsetroundjoin%
\definecolor{currentfill}{rgb}{0.000000,0.000000,0.000000}%
\pgfsetfillcolor{currentfill}%
\pgfsetlinewidth{0.803000pt}%
\definecolor{currentstroke}{rgb}{0.000000,0.000000,0.000000}%
\pgfsetstrokecolor{currentstroke}%
\pgfsetdash{}{0pt}%
\pgfsys@defobject{currentmarker}{\pgfqpoint{-0.048611in}{0.000000in}}{\pgfqpoint{-0.000000in}{0.000000in}}{%
\pgfpathmoveto{\pgfqpoint{-0.000000in}{0.000000in}}%
\pgfpathlineto{\pgfqpoint{-0.048611in}{0.000000in}}%
\pgfusepath{stroke,fill}%
}%
\begin{pgfscope}%
\pgfsys@transformshift{0.487500in}{0.967500in}%
\pgfsys@useobject{currentmarker}{}%
\end{pgfscope}%
\end{pgfscope}%
\begin{pgfscope}%
\definecolor{textcolor}{rgb}{0.000000,0.000000,0.000000}%
\pgfsetstrokecolor{textcolor}%
\pgfsetfillcolor{textcolor}%
\pgftext[x=0.111728in, y=0.932062in, left, base]{\color{textcolor}{\rmfamily\fontsize{6.000000}{7.200000}\selectfont\catcode`\^=\active\def^{\ifmmode\sp\else\^{}\fi}\catcode`\%=\active\def%{\%}$\mathdefault{10^{-11}}$}}%
\end{pgfscope}%
\begin{pgfscope}%
\pgfpathrectangle{\pgfqpoint{0.487500in}{0.320625in}}{\pgfqpoint{1.845679in}{1.164375in}}%
\pgfusepath{clip}%
\pgfsetrectcap%
\pgfsetroundjoin%
\pgfsetlinewidth{0.803000pt}%
\definecolor{currentstroke}{rgb}{0.909804,0.909804,0.909804}%
\pgfsetstrokecolor{currentstroke}%
\pgfsetdash{}{0pt}%
\pgfpathmoveto{\pgfqpoint{0.487500in}{1.183125in}}%
\pgfpathlineto{\pgfqpoint{2.333179in}{1.183125in}}%
\pgfusepath{stroke}%
\end{pgfscope}%
\begin{pgfscope}%
\pgfsetbuttcap%
\pgfsetroundjoin%
\definecolor{currentfill}{rgb}{0.000000,0.000000,0.000000}%
\pgfsetfillcolor{currentfill}%
\pgfsetlinewidth{0.803000pt}%
\definecolor{currentstroke}{rgb}{0.000000,0.000000,0.000000}%
\pgfsetstrokecolor{currentstroke}%
\pgfsetdash{}{0pt}%
\pgfsys@defobject{currentmarker}{\pgfqpoint{-0.048611in}{0.000000in}}{\pgfqpoint{-0.000000in}{0.000000in}}{%
\pgfpathmoveto{\pgfqpoint{-0.000000in}{0.000000in}}%
\pgfpathlineto{\pgfqpoint{-0.048611in}{0.000000in}}%
\pgfusepath{stroke,fill}%
}%
\begin{pgfscope}%
\pgfsys@transformshift{0.487500in}{1.183125in}%
\pgfsys@useobject{currentmarker}{}%
\end{pgfscope}%
\end{pgfscope}%
\begin{pgfscope}%
\definecolor{textcolor}{rgb}{0.000000,0.000000,0.000000}%
\pgfsetstrokecolor{textcolor}%
\pgfsetfillcolor{textcolor}%
\pgftext[x=0.158989in, y=1.147687in, left, base]{\color{textcolor}{\rmfamily\fontsize{6.000000}{7.200000}\selectfont\catcode`\^=\active\def^{\ifmmode\sp\else\^{}\fi}\catcode`\%=\active\def%{\%}$\mathdefault{10^{-6}}$}}%
\end{pgfscope}%
\begin{pgfscope}%
\pgfpathrectangle{\pgfqpoint{0.487500in}{0.320625in}}{\pgfqpoint{1.845679in}{1.164375in}}%
\pgfusepath{clip}%
\pgfsetrectcap%
\pgfsetroundjoin%
\pgfsetlinewidth{0.803000pt}%
\definecolor{currentstroke}{rgb}{0.909804,0.909804,0.909804}%
\pgfsetstrokecolor{currentstroke}%
\pgfsetdash{}{0pt}%
\pgfpathmoveto{\pgfqpoint{0.487500in}{1.398750in}}%
\pgfpathlineto{\pgfqpoint{2.333179in}{1.398750in}}%
\pgfusepath{stroke}%
\end{pgfscope}%
\begin{pgfscope}%
\pgfsetbuttcap%
\pgfsetroundjoin%
\definecolor{currentfill}{rgb}{0.000000,0.000000,0.000000}%
\pgfsetfillcolor{currentfill}%
\pgfsetlinewidth{0.803000pt}%
\definecolor{currentstroke}{rgb}{0.000000,0.000000,0.000000}%
\pgfsetstrokecolor{currentstroke}%
\pgfsetdash{}{0pt}%
\pgfsys@defobject{currentmarker}{\pgfqpoint{-0.048611in}{0.000000in}}{\pgfqpoint{-0.000000in}{0.000000in}}{%
\pgfpathmoveto{\pgfqpoint{-0.000000in}{0.000000in}}%
\pgfpathlineto{\pgfqpoint{-0.048611in}{0.000000in}}%
\pgfusepath{stroke,fill}%
}%
\begin{pgfscope}%
\pgfsys@transformshift{0.487500in}{1.398750in}%
\pgfsys@useobject{currentmarker}{}%
\end{pgfscope}%
\end{pgfscope}%
\begin{pgfscope}%
\definecolor{textcolor}{rgb}{0.000000,0.000000,0.000000}%
\pgfsetstrokecolor{textcolor}%
\pgfsetfillcolor{textcolor}%
\pgftext[x=0.158989in, y=1.363312in, left, base]{\color{textcolor}{\rmfamily\fontsize{6.000000}{7.200000}\selectfont\catcode`\^=\active\def^{\ifmmode\sp\else\^{}\fi}\catcode`\%=\active\def%{\%}$\mathdefault{10^{-1}}$}}%
\end{pgfscope}%
\begin{pgfscope}%
\definecolor{textcolor}{rgb}{0.000000,0.000000,0.000000}%
\pgfsetstrokecolor{textcolor}%
\pgfsetfillcolor{textcolor}%
\pgftext[x=0.056172in,y=0.902813in,,bottom,rotate=90.000000]{\color{textcolor}{\rmfamily\fontsize{6.000000}{7.200000}\selectfont\catcode`\^=\active\def^{\ifmmode\sp\else\^{}\fi}\catcode`\%=\active\def%{\%}$\ell_t-\ell_\star$}}%
\end{pgfscope}%
\begin{pgfscope}%
\pgfpathrectangle{\pgfqpoint{0.487500in}{0.320625in}}{\pgfqpoint{1.845679in}{1.164375in}}%
\pgfusepath{clip}%
\pgfsetrectcap%
\pgfsetroundjoin%
\pgfsetlinewidth{1.254687pt}%
\definecolor{currentstroke}{rgb}{0.141176,0.345098,0.650980}%
\pgfsetstrokecolor{currentstroke}%
\pgfsetdash{}{0pt}%
\pgfpathmoveto{\pgfqpoint{0.571395in}{1.468525in}}%
\pgfpathlineto{\pgfqpoint{0.588173in}{1.465082in}}%
\pgfpathlineto{\pgfqpoint{0.598660in}{1.461710in}}%
\pgfpathlineto{\pgfqpoint{0.607050in}{1.457758in}}%
\pgfpathlineto{\pgfqpoint{0.615439in}{1.452402in}}%
\pgfpathlineto{\pgfqpoint{0.623829in}{1.445657in}}%
\pgfpathlineto{\pgfqpoint{0.636413in}{1.433755in}}%
\pgfpathlineto{\pgfqpoint{0.663678in}{1.407362in}}%
\pgfpathlineto{\pgfqpoint{0.680457in}{1.392974in}}%
\pgfpathlineto{\pgfqpoint{0.697236in}{1.380002in}}%
\pgfpathlineto{\pgfqpoint{0.718210in}{1.365204in}}%
\pgfpathlineto{\pgfqpoint{0.749670in}{1.344516in}}%
\pgfpathlineto{\pgfqpoint{1.087346in}{1.127817in}}%
\pgfpathlineto{\pgfqpoint{1.309666in}{0.987051in}}%
\pgfpathlineto{\pgfqpoint{1.731236in}{0.721516in}}%
\pgfpathlineto{\pgfqpoint{1.758502in}{0.704552in}}%
\pgfpathlineto{\pgfqpoint{1.810936in}{0.670820in}}%
\pgfpathlineto{\pgfqpoint{1.829812in}{0.660711in}}%
\pgfpathlineto{\pgfqpoint{1.846591in}{0.651559in}}%
\pgfpathlineto{\pgfqpoint{1.859175in}{0.644066in}}%
\pgfpathlineto{\pgfqpoint{1.867565in}{0.640854in}}%
\pgfpathlineto{\pgfqpoint{1.875954in}{0.637730in}}%
\pgfpathlineto{\pgfqpoint{1.880149in}{0.635243in}}%
\pgfpathlineto{\pgfqpoint{1.884343in}{0.631791in}}%
\pgfpathlineto{\pgfqpoint{1.896928in}{0.619924in}}%
\pgfpathlineto{\pgfqpoint{1.899025in}{0.619630in}}%
\pgfpathlineto{\pgfqpoint{1.901122in}{0.620253in}}%
\pgfpathlineto{\pgfqpoint{1.905317in}{0.623417in}}%
\pgfpathlineto{\pgfqpoint{1.909512in}{0.626985in}}%
\pgfpathlineto{\pgfqpoint{1.911609in}{0.628204in}}%
\pgfpathlineto{\pgfqpoint{1.913707in}{0.628753in}}%
\pgfpathlineto{\pgfqpoint{1.915804in}{0.628421in}}%
\pgfpathlineto{\pgfqpoint{1.917901in}{0.626958in}}%
\pgfpathlineto{\pgfqpoint{1.919999in}{0.624030in}}%
\pgfpathlineto{\pgfqpoint{1.922096in}{0.619308in}}%
\pgfpathlineto{\pgfqpoint{1.928388in}{0.600753in}}%
\pgfpathlineto{\pgfqpoint{1.930485in}{0.599600in}}%
\pgfpathlineto{\pgfqpoint{1.932583in}{0.602544in}}%
\pgfpathlineto{\pgfqpoint{1.936777in}{0.610531in}}%
\pgfpathlineto{\pgfqpoint{1.938875in}{0.612766in}}%
\pgfpathlineto{\pgfqpoint{1.940972in}{0.613312in}}%
\pgfpathlineto{\pgfqpoint{1.943070in}{0.612035in}}%
\pgfpathlineto{\pgfqpoint{1.945167in}{0.608767in}}%
\pgfpathlineto{\pgfqpoint{1.947264in}{0.603286in}}%
\pgfpathlineto{\pgfqpoint{1.951459in}{0.587383in}}%
\pgfpathlineto{\pgfqpoint{1.953556in}{0.585385in}}%
\pgfpathlineto{\pgfqpoint{1.955654in}{0.593400in}}%
\pgfpathlineto{\pgfqpoint{1.959848in}{0.613096in}}%
\pgfpathlineto{\pgfqpoint{1.961946in}{0.619497in}}%
\pgfpathlineto{\pgfqpoint{1.964043in}{0.623311in}}%
\pgfpathlineto{\pgfqpoint{1.966141in}{0.624363in}}%
\pgfpathlineto{\pgfqpoint{1.968238in}{0.621992in}}%
\pgfpathlineto{\pgfqpoint{1.970335in}{0.615667in}}%
\pgfpathlineto{\pgfqpoint{1.972433in}{0.605813in}}%
\pgfpathlineto{\pgfqpoint{1.976627in}{0.573871in}}%
\pgfpathlineto{\pgfqpoint{1.978725in}{0.571980in}}%
\pgfpathlineto{\pgfqpoint{1.982919in}{0.598860in}}%
\pgfpathlineto{\pgfqpoint{1.985017in}{0.606422in}}%
\pgfpathlineto{\pgfqpoint{1.987114in}{0.610625in}}%
\pgfpathlineto{\pgfqpoint{1.989212in}{0.612082in}}%
\pgfpathlineto{\pgfqpoint{1.991309in}{0.610972in}}%
\pgfpathlineto{\pgfqpoint{1.993406in}{0.607084in}}%
\pgfpathlineto{\pgfqpoint{1.995504in}{0.599678in}}%
\pgfpathlineto{\pgfqpoint{1.997601in}{0.586985in}}%
\pgfpathlineto{\pgfqpoint{1.999698in}{0.565811in}}%
\pgfpathlineto{\pgfqpoint{2.001796in}{0.558531in}}%
\pgfpathlineto{\pgfqpoint{2.003893in}{0.584413in}}%
\pgfpathlineto{\pgfqpoint{2.005990in}{0.602594in}}%
\pgfpathlineto{\pgfqpoint{2.008088in}{0.614372in}}%
\pgfpathlineto{\pgfqpoint{2.010185in}{0.622039in}}%
\pgfpathlineto{\pgfqpoint{2.012283in}{0.626508in}}%
\pgfpathlineto{\pgfqpoint{2.014380in}{0.627922in}}%
\pgfpathlineto{\pgfqpoint{2.016477in}{0.625720in}}%
\pgfpathlineto{\pgfqpoint{2.018575in}{0.618862in}}%
\pgfpathlineto{\pgfqpoint{2.020672in}{0.606524in}}%
\pgfpathlineto{\pgfqpoint{2.022769in}{0.584094in}}%
\pgfpathlineto{\pgfqpoint{2.024867in}{0.540075in}}%
\pgfpathlineto{\pgfqpoint{2.026964in}{0.579381in}}%
\pgfpathlineto{\pgfqpoint{2.029061in}{0.603488in}}%
\pgfpathlineto{\pgfqpoint{2.031159in}{0.616636in}}%
\pgfpathlineto{\pgfqpoint{2.033256in}{0.624596in}}%
\pgfpathlineto{\pgfqpoint{2.035354in}{0.629273in}}%
\pgfpathlineto{\pgfqpoint{2.037451in}{0.631441in}}%
\pgfpathlineto{\pgfqpoint{2.039548in}{0.631377in}}%
\pgfpathlineto{\pgfqpoint{2.041646in}{0.629038in}}%
\pgfpathlineto{\pgfqpoint{2.043743in}{0.624020in}}%
\pgfpathlineto{\pgfqpoint{2.045840in}{0.615284in}}%
\pgfpathlineto{\pgfqpoint{2.047938in}{0.600071in}}%
\pgfpathlineto{\pgfqpoint{2.050035in}{0.567555in}}%
\pgfpathlineto{\pgfqpoint{2.052132in}{0.552887in}}%
\pgfpathlineto{\pgfqpoint{2.054230in}{0.598028in}}%
\pgfpathlineto{\pgfqpoint{2.056327in}{0.618178in}}%
\pgfpathlineto{\pgfqpoint{2.058425in}{0.630437in}}%
\pgfpathlineto{\pgfqpoint{2.060522in}{0.638382in}}%
\pgfpathlineto{\pgfqpoint{2.062619in}{0.643190in}}%
\pgfpathlineto{\pgfqpoint{2.064717in}{0.645163in}}%
\pgfpathlineto{\pgfqpoint{2.066814in}{0.643977in}}%
\pgfpathlineto{\pgfqpoint{2.068911in}{0.638657in}}%
\pgfpathlineto{\pgfqpoint{2.071009in}{0.628969in}}%
\pgfpathlineto{\pgfqpoint{2.073106in}{0.612772in}}%
\pgfpathlineto{\pgfqpoint{2.077301in}{0.546799in}}%
\pgfpathlineto{\pgfqpoint{2.079398in}{0.599625in}}%
\pgfpathlineto{\pgfqpoint{2.081496in}{0.618811in}}%
\pgfpathlineto{\pgfqpoint{2.083593in}{0.629279in}}%
\pgfpathlineto{\pgfqpoint{2.085690in}{0.635230in}}%
\pgfpathlineto{\pgfqpoint{2.087788in}{0.638065in}}%
\pgfpathlineto{\pgfqpoint{2.089885in}{0.638284in}}%
\pgfpathlineto{\pgfqpoint{2.091982in}{0.635902in}}%
\pgfpathlineto{\pgfqpoint{2.094080in}{0.630459in}}%
\pgfpathlineto{\pgfqpoint{2.096177in}{0.620636in}}%
\pgfpathlineto{\pgfqpoint{2.098274in}{0.602508in}}%
\pgfpathlineto{\pgfqpoint{2.100372in}{0.553800in}}%
\pgfpathlineto{\pgfqpoint{2.102469in}{0.579340in}}%
\pgfpathlineto{\pgfqpoint{2.104566in}{0.614316in}}%
\pgfpathlineto{\pgfqpoint{2.106664in}{0.631743in}}%
\pgfpathlineto{\pgfqpoint{2.108761in}{0.642588in}}%
\pgfpathlineto{\pgfqpoint{2.110859in}{0.649471in}}%
\pgfpathlineto{\pgfqpoint{2.112956in}{0.653232in}}%
\pgfpathlineto{\pgfqpoint{2.115053in}{0.653916in}}%
\pgfpathlineto{\pgfqpoint{2.117151in}{0.650792in}}%
\pgfpathlineto{\pgfqpoint{2.119248in}{0.643921in}}%
\pgfpathlineto{\pgfqpoint{2.121345in}{0.633882in}}%
\pgfpathlineto{\pgfqpoint{2.123443in}{0.619116in}}%
\pgfpathlineto{\pgfqpoint{2.125540in}{0.594822in}}%
\pgfpathlineto{\pgfqpoint{2.127637in}{0.525487in}}%
\pgfpathlineto{\pgfqpoint{2.129735in}{0.574881in}}%
\pgfpathlineto{\pgfqpoint{2.131832in}{0.600946in}}%
\pgfpathlineto{\pgfqpoint{2.133930in}{0.612340in}}%
\pgfpathlineto{\pgfqpoint{2.136027in}{0.617373in}}%
\pgfpathlineto{\pgfqpoint{2.138124in}{0.618041in}}%
\pgfpathlineto{\pgfqpoint{2.140222in}{0.614624in}}%
\pgfpathlineto{\pgfqpoint{2.142319in}{0.606085in}}%
\pgfpathlineto{\pgfqpoint{2.144416in}{0.588368in}}%
\pgfpathlineto{\pgfqpoint{2.146514in}{0.534649in}}%
\pgfpathlineto{\pgfqpoint{2.150708in}{0.605276in}}%
\pgfpathlineto{\pgfqpoint{2.152806in}{0.622484in}}%
\pgfpathlineto{\pgfqpoint{2.154903in}{0.633180in}}%
\pgfpathlineto{\pgfqpoint{2.157001in}{0.639773in}}%
\pgfpathlineto{\pgfqpoint{2.159098in}{0.642954in}}%
\pgfpathlineto{\pgfqpoint{2.161195in}{0.642509in}}%
\pgfpathlineto{\pgfqpoint{2.163293in}{0.637584in}}%
\pgfpathlineto{\pgfqpoint{2.165390in}{0.629780in}}%
\pgfpathlineto{\pgfqpoint{2.167487in}{0.618958in}}%
\pgfpathlineto{\pgfqpoint{2.169585in}{0.603560in}}%
\pgfpathlineto{\pgfqpoint{2.171682in}{0.579000in}}%
\pgfpathlineto{\pgfqpoint{2.173779in}{0.514239in}}%
\pgfpathlineto{\pgfqpoint{2.175877in}{0.553707in}}%
\pgfpathlineto{\pgfqpoint{2.177974in}{0.579778in}}%
\pgfpathlineto{\pgfqpoint{2.180072in}{0.590042in}}%
\pgfpathlineto{\pgfqpoint{2.182169in}{0.593293in}}%
\pgfpathlineto{\pgfqpoint{2.184266in}{0.591235in}}%
\pgfpathlineto{\pgfqpoint{2.186364in}{0.583174in}}%
\pgfpathlineto{\pgfqpoint{2.188461in}{0.564510in}}%
\pgfpathlineto{\pgfqpoint{2.190558in}{0.493356in}}%
\pgfpathlineto{\pgfqpoint{2.192656in}{0.558714in}}%
\pgfpathlineto{\pgfqpoint{2.194753in}{0.589238in}}%
\pgfpathlineto{\pgfqpoint{2.196850in}{0.605700in}}%
\pgfpathlineto{\pgfqpoint{2.198948in}{0.616049in}}%
\pgfpathlineto{\pgfqpoint{2.201045in}{0.622367in}}%
\pgfpathlineto{\pgfqpoint{2.203143in}{0.625201in}}%
\pgfpathlineto{\pgfqpoint{2.205240in}{0.624180in}}%
\pgfpathlineto{\pgfqpoint{2.207337in}{0.618649in}}%
\pgfpathlineto{\pgfqpoint{2.209435in}{0.610514in}}%
\pgfpathlineto{\pgfqpoint{2.211532in}{0.599364in}}%
\pgfpathlineto{\pgfqpoint{2.213629in}{0.583602in}}%
\pgfpathlineto{\pgfqpoint{2.215727in}{0.558526in}}%
\pgfpathlineto{\pgfqpoint{2.217824in}{0.491163in}}%
\pgfpathlineto{\pgfqpoint{2.219921in}{0.533815in}}%
\pgfpathlineto{\pgfqpoint{2.222019in}{0.559016in}}%
\pgfpathlineto{\pgfqpoint{2.224116in}{0.568583in}}%
\pgfpathlineto{\pgfqpoint{2.226214in}{0.570966in}}%
\pgfpathlineto{\pgfqpoint{2.228311in}{0.567587in}}%
\pgfpathlineto{\pgfqpoint{2.230408in}{0.557019in}}%
\pgfpathlineto{\pgfqpoint{2.232506in}{0.530886in}}%
\pgfpathlineto{\pgfqpoint{2.234603in}{0.485308in}}%
\pgfpathlineto{\pgfqpoint{2.236700in}{0.552458in}}%
\pgfpathlineto{\pgfqpoint{2.238798in}{0.576067in}}%
\pgfpathlineto{\pgfqpoint{2.240895in}{0.589993in}}%
\pgfpathlineto{\pgfqpoint{2.242992in}{0.598832in}}%
\pgfpathlineto{\pgfqpoint{2.245090in}{0.603940in}}%
\pgfpathlineto{\pgfqpoint{2.247187in}{0.605520in}}%
\pgfpathlineto{\pgfqpoint{2.249285in}{0.602808in}}%
\pgfpathlineto{\pgfqpoint{2.249285in}{0.602808in}}%
\pgfusepath{stroke}%
\end{pgfscope}%
\begin{pgfscope}%
\pgfpathrectangle{\pgfqpoint{0.487500in}{0.320625in}}{\pgfqpoint{1.845679in}{1.164375in}}%
\pgfusepath{clip}%
\pgfsetbuttcap%
\pgfsetroundjoin%
\pgfsetlinewidth{0.752812pt}%
\definecolor{currentstroke}{rgb}{0.141176,0.345098,0.650980}%
\pgfsetstrokecolor{currentstroke}%
\pgfsetstrokeopacity{0.650000}%
\pgfsetdash{{2.775000pt}{1.200000pt}}{0.000000pt}%
\pgfpathmoveto{\pgfqpoint{0.571395in}{1.468525in}}%
\pgfpathlineto{\pgfqpoint{2.249285in}{1.168861in}}%
\pgfpathlineto{\pgfqpoint{2.249285in}{1.168861in}}%
\pgfusepath{stroke}%
\end{pgfscope}%
\begin{pgfscope}%
\pgfpathrectangle{\pgfqpoint{0.487500in}{0.320625in}}{\pgfqpoint{1.845679in}{1.164375in}}%
\pgfusepath{clip}%
\pgfsetbuttcap%
\pgfsetroundjoin%
\pgfsetlinewidth{0.803000pt}%
\definecolor{currentstroke}{rgb}{0.141176,0.345098,0.650980}%
\pgfsetstrokecolor{currentstroke}%
\pgfsetstrokeopacity{0.800000}%
\pgfsetdash{{0.800000pt}{1.320000pt}}{0.000000pt}%
\pgfpathmoveto{\pgfqpoint{0.571395in}{1.468525in}}%
\pgfpathlineto{\pgfqpoint{2.249285in}{1.468525in}}%
\pgfpathlineto{\pgfqpoint{2.249285in}{1.468525in}}%
\pgfusepath{stroke}%
\end{pgfscope}%
\begin{pgfscope}%
\pgfpathrectangle{\pgfqpoint{0.487500in}{0.320625in}}{\pgfqpoint{1.845679in}{1.164375in}}%
\pgfusepath{clip}%
\pgfsetrectcap%
\pgfsetroundjoin%
\pgfsetlinewidth{1.254687pt}%
\definecolor{currentstroke}{rgb}{0.819608,0.286275,0.356863}%
\pgfsetstrokecolor{currentstroke}%
\pgfsetdash{}{0pt}%
\pgfpathmoveto{\pgfqpoint{0.571395in}{1.468525in}}%
\pgfpathlineto{\pgfqpoint{0.581881in}{1.458695in}}%
\pgfpathlineto{\pgfqpoint{0.590271in}{1.449303in}}%
\pgfpathlineto{\pgfqpoint{0.598660in}{1.438055in}}%
\pgfpathlineto{\pgfqpoint{0.609147in}{1.421752in}}%
\pgfpathlineto{\pgfqpoint{0.659484in}{1.339277in}}%
\pgfpathlineto{\pgfqpoint{0.678360in}{1.312158in}}%
\pgfpathlineto{\pgfqpoint{0.699334in}{1.284091in}}%
\pgfpathlineto{\pgfqpoint{0.724502in}{1.252341in}}%
\pgfpathlineto{\pgfqpoint{0.755962in}{1.214487in}}%
\pgfpathlineto{\pgfqpoint{0.797910in}{1.165805in}}%
\pgfpathlineto{\pgfqpoint{0.858733in}{1.097049in}}%
\pgfpathlineto{\pgfqpoint{0.948920in}{0.996935in}}%
\pgfpathlineto{\pgfqpoint{1.083151in}{0.849760in}}%
\pgfpathlineto{\pgfqpoint{1.215285in}{0.706615in}}%
\pgfpathlineto{\pgfqpoint{1.227869in}{0.693571in}}%
\pgfpathlineto{\pgfqpoint{1.242551in}{0.676276in}}%
\pgfpathlineto{\pgfqpoint{1.257232in}{0.659096in}}%
\pgfpathlineto{\pgfqpoint{1.265621in}{0.650825in}}%
\pgfpathlineto{\pgfqpoint{1.274011in}{0.644232in}}%
\pgfpathlineto{\pgfqpoint{1.286595in}{0.635261in}}%
\pgfpathlineto{\pgfqpoint{1.290790in}{0.631044in}}%
\pgfpathlineto{\pgfqpoint{1.297082in}{0.623267in}}%
\pgfpathlineto{\pgfqpoint{1.305471in}{0.610723in}}%
\pgfpathlineto{\pgfqpoint{1.313861in}{0.595867in}}%
\pgfpathlineto{\pgfqpoint{1.328542in}{0.568981in}}%
\pgfpathlineto{\pgfqpoint{1.332737in}{0.563684in}}%
\pgfpathlineto{\pgfqpoint{1.336932in}{0.560797in}}%
\pgfpathlineto{\pgfqpoint{1.339029in}{0.560380in}}%
\pgfpathlineto{\pgfqpoint{1.341127in}{0.560634in}}%
\pgfpathlineto{\pgfqpoint{1.345321in}{0.562857in}}%
\pgfpathlineto{\pgfqpoint{1.351613in}{0.568758in}}%
\pgfpathlineto{\pgfqpoint{1.360003in}{0.576712in}}%
\pgfpathlineto{\pgfqpoint{1.364198in}{0.579265in}}%
\pgfpathlineto{\pgfqpoint{1.368392in}{0.580168in}}%
\pgfpathlineto{\pgfqpoint{1.370490in}{0.579797in}}%
\pgfpathlineto{\pgfqpoint{1.372587in}{0.578716in}}%
\pgfpathlineto{\pgfqpoint{1.376782in}{0.574038in}}%
\pgfpathlineto{\pgfqpoint{1.380976in}{0.566593in}}%
\pgfpathlineto{\pgfqpoint{1.385171in}{0.555850in}}%
\pgfpathlineto{\pgfqpoint{1.389366in}{0.540946in}}%
\pgfpathlineto{\pgfqpoint{1.393561in}{0.525670in}}%
\pgfpathlineto{\pgfqpoint{1.395658in}{0.524737in}}%
\pgfpathlineto{\pgfqpoint{1.397755in}{0.530092in}}%
\pgfpathlineto{\pgfqpoint{1.404047in}{0.551417in}}%
\pgfpathlineto{\pgfqpoint{1.408242in}{0.560987in}}%
\pgfpathlineto{\pgfqpoint{1.412437in}{0.567347in}}%
\pgfpathlineto{\pgfqpoint{1.416632in}{0.571280in}}%
\pgfpathlineto{\pgfqpoint{1.420826in}{0.573215in}}%
\pgfpathlineto{\pgfqpoint{1.425021in}{0.573331in}}%
\pgfpathlineto{\pgfqpoint{1.429216in}{0.571605in}}%
\pgfpathlineto{\pgfqpoint{1.433410in}{0.567795in}}%
\pgfpathlineto{\pgfqpoint{1.437605in}{0.561307in}}%
\pgfpathlineto{\pgfqpoint{1.441800in}{0.550787in}}%
\pgfpathlineto{\pgfqpoint{1.443897in}{0.543078in}}%
\pgfpathlineto{\pgfqpoint{1.445995in}{0.532592in}}%
\pgfpathlineto{\pgfqpoint{1.448092in}{0.517298in}}%
\pgfpathlineto{\pgfqpoint{1.450189in}{0.493295in}}%
\pgfpathlineto{\pgfqpoint{1.452287in}{0.485703in}}%
\pgfpathlineto{\pgfqpoint{1.454384in}{0.512914in}}%
\pgfpathlineto{\pgfqpoint{1.456481in}{0.531545in}}%
\pgfpathlineto{\pgfqpoint{1.460676in}{0.553973in}}%
\pgfpathlineto{\pgfqpoint{1.464871in}{0.567670in}}%
\pgfpathlineto{\pgfqpoint{1.469066in}{0.576945in}}%
\pgfpathlineto{\pgfqpoint{1.473260in}{0.583300in}}%
\pgfpathlineto{\pgfqpoint{1.477455in}{0.587333in}}%
\pgfpathlineto{\pgfqpoint{1.481650in}{0.589203in}}%
\pgfpathlineto{\pgfqpoint{1.483747in}{0.589282in}}%
\pgfpathlineto{\pgfqpoint{1.485845in}{0.588719in}}%
\pgfpathlineto{\pgfqpoint{1.487942in}{0.587415in}}%
\pgfpathlineto{\pgfqpoint{1.490039in}{0.585218in}}%
\pgfpathlineto{\pgfqpoint{1.496331in}{0.576813in}}%
\pgfpathlineto{\pgfqpoint{1.498429in}{0.577210in}}%
\pgfpathlineto{\pgfqpoint{1.500526in}{0.580317in}}%
\pgfpathlineto{\pgfqpoint{1.506818in}{0.593940in}}%
\pgfpathlineto{\pgfqpoint{1.511013in}{0.599555in}}%
\pgfpathlineto{\pgfqpoint{1.513110in}{0.600813in}}%
\pgfpathlineto{\pgfqpoint{1.515208in}{0.600991in}}%
\pgfpathlineto{\pgfqpoint{1.517305in}{0.600022in}}%
\pgfpathlineto{\pgfqpoint{1.519402in}{0.597776in}}%
\pgfpathlineto{\pgfqpoint{1.521500in}{0.594025in}}%
\pgfpathlineto{\pgfqpoint{1.523597in}{0.588368in}}%
\pgfpathlineto{\pgfqpoint{1.525694in}{0.580069in}}%
\pgfpathlineto{\pgfqpoint{1.527792in}{0.567732in}}%
\pgfpathlineto{\pgfqpoint{1.529889in}{0.549224in}}%
\pgfpathlineto{\pgfqpoint{1.531987in}{0.536199in}}%
\pgfpathlineto{\pgfqpoint{1.536181in}{0.574527in}}%
\pgfpathlineto{\pgfqpoint{1.538279in}{0.587318in}}%
\pgfpathlineto{\pgfqpoint{1.540376in}{0.596349in}}%
\pgfpathlineto{\pgfqpoint{1.542473in}{0.602757in}}%
\pgfpathlineto{\pgfqpoint{1.544571in}{0.607076in}}%
\pgfpathlineto{\pgfqpoint{1.546668in}{0.609503in}}%
\pgfpathlineto{\pgfqpoint{1.548765in}{0.609977in}}%
\pgfpathlineto{\pgfqpoint{1.550863in}{0.608132in}}%
\pgfpathlineto{\pgfqpoint{1.552960in}{0.603009in}}%
\pgfpathlineto{\pgfqpoint{1.555058in}{0.592321in}}%
\pgfpathlineto{\pgfqpoint{1.557155in}{0.572169in}}%
\pgfpathlineto{\pgfqpoint{1.559252in}{0.539881in}}%
\pgfpathlineto{\pgfqpoint{1.563447in}{0.592255in}}%
\pgfpathlineto{\pgfqpoint{1.565544in}{0.605736in}}%
\pgfpathlineto{\pgfqpoint{1.567642in}{0.614193in}}%
\pgfpathlineto{\pgfqpoint{1.569739in}{0.619435in}}%
\pgfpathlineto{\pgfqpoint{1.571836in}{0.622242in}}%
\pgfpathlineto{\pgfqpoint{1.573934in}{0.622932in}}%
\pgfpathlineto{\pgfqpoint{1.576031in}{0.621543in}}%
\pgfpathlineto{\pgfqpoint{1.578129in}{0.617831in}}%
\pgfpathlineto{\pgfqpoint{1.580226in}{0.611135in}}%
\pgfpathlineto{\pgfqpoint{1.582323in}{0.599841in}}%
\pgfpathlineto{\pgfqpoint{1.584421in}{0.579062in}}%
\pgfpathlineto{\pgfqpoint{1.586518in}{0.513327in}}%
\pgfpathlineto{\pgfqpoint{1.588615in}{0.568403in}}%
\pgfpathlineto{\pgfqpoint{1.590713in}{0.598283in}}%
\pgfpathlineto{\pgfqpoint{1.592810in}{0.614196in}}%
\pgfpathlineto{\pgfqpoint{1.594907in}{0.624282in}}%
\pgfpathlineto{\pgfqpoint{1.597005in}{0.630723in}}%
\pgfpathlineto{\pgfqpoint{1.599102in}{0.634240in}}%
\pgfpathlineto{\pgfqpoint{1.601199in}{0.634862in}}%
\pgfpathlineto{\pgfqpoint{1.603297in}{0.631928in}}%
\pgfpathlineto{\pgfqpoint{1.605394in}{0.624788in}}%
\pgfpathlineto{\pgfqpoint{1.607492in}{0.613399in}}%
\pgfpathlineto{\pgfqpoint{1.609589in}{0.595077in}}%
\pgfpathlineto{\pgfqpoint{1.611686in}{0.560621in}}%
\pgfpathlineto{\pgfqpoint{1.613784in}{0.556986in}}%
\pgfpathlineto{\pgfqpoint{1.615881in}{0.588707in}}%
\pgfpathlineto{\pgfqpoint{1.617978in}{0.603688in}}%
\pgfpathlineto{\pgfqpoint{1.620076in}{0.611530in}}%
\pgfpathlineto{\pgfqpoint{1.622173in}{0.615036in}}%
\pgfpathlineto{\pgfqpoint{1.624270in}{0.615099in}}%
\pgfpathlineto{\pgfqpoint{1.626368in}{0.611729in}}%
\pgfpathlineto{\pgfqpoint{1.628465in}{0.604050in}}%
\pgfpathlineto{\pgfqpoint{1.630563in}{0.589181in}}%
\pgfpathlineto{\pgfqpoint{1.632660in}{0.554073in}}%
\pgfpathlineto{\pgfqpoint{1.634757in}{0.545982in}}%
\pgfpathlineto{\pgfqpoint{1.636855in}{0.591182in}}%
\pgfpathlineto{\pgfqpoint{1.638952in}{0.610958in}}%
\pgfpathlineto{\pgfqpoint{1.641049in}{0.622788in}}%
\pgfpathlineto{\pgfqpoint{1.643147in}{0.630048in}}%
\pgfpathlineto{\pgfqpoint{1.645244in}{0.633713in}}%
\pgfpathlineto{\pgfqpoint{1.647341in}{0.633713in}}%
\pgfpathlineto{\pgfqpoint{1.649439in}{0.629026in}}%
\pgfpathlineto{\pgfqpoint{1.651536in}{0.620035in}}%
\pgfpathlineto{\pgfqpoint{1.653634in}{0.605666in}}%
\pgfpathlineto{\pgfqpoint{1.655731in}{0.580415in}}%
\pgfpathlineto{\pgfqpoint{1.657828in}{0.524118in}}%
\pgfpathlineto{\pgfqpoint{1.659926in}{0.575467in}}%
\pgfpathlineto{\pgfqpoint{1.662023in}{0.597955in}}%
\pgfpathlineto{\pgfqpoint{1.664120in}{0.609111in}}%
\pgfpathlineto{\pgfqpoint{1.666218in}{0.614782in}}%
\pgfpathlineto{\pgfqpoint{1.668315in}{0.616682in}}%
\pgfpathlineto{\pgfqpoint{1.670412in}{0.615263in}}%
\pgfpathlineto{\pgfqpoint{1.672510in}{0.610188in}}%
\pgfpathlineto{\pgfqpoint{1.674607in}{0.599978in}}%
\pgfpathlineto{\pgfqpoint{1.676705in}{0.579598in}}%
\pgfpathlineto{\pgfqpoint{1.678802in}{0.505737in}}%
\pgfpathlineto{\pgfqpoint{1.680899in}{0.572547in}}%
\pgfpathlineto{\pgfqpoint{1.682997in}{0.601951in}}%
\pgfpathlineto{\pgfqpoint{1.685094in}{0.617613in}}%
\pgfpathlineto{\pgfqpoint{1.687191in}{0.627270in}}%
\pgfpathlineto{\pgfqpoint{1.689289in}{0.632928in}}%
\pgfpathlineto{\pgfqpoint{1.691386in}{0.635063in}}%
\pgfpathlineto{\pgfqpoint{1.693483in}{0.633177in}}%
\pgfpathlineto{\pgfqpoint{1.695581in}{0.626381in}}%
\pgfpathlineto{\pgfqpoint{1.697678in}{0.614855in}}%
\pgfpathlineto{\pgfqpoint{1.699776in}{0.595316in}}%
\pgfpathlineto{\pgfqpoint{1.701873in}{0.550904in}}%
\pgfpathlineto{\pgfqpoint{1.703970in}{0.560907in}}%
\pgfpathlineto{\pgfqpoint{1.706068in}{0.594987in}}%
\pgfpathlineto{\pgfqpoint{1.708165in}{0.609996in}}%
\pgfpathlineto{\pgfqpoint{1.710262in}{0.618011in}}%
\pgfpathlineto{\pgfqpoint{1.712360in}{0.621862in}}%
\pgfpathlineto{\pgfqpoint{1.714457in}{0.622469in}}%
\pgfpathlineto{\pgfqpoint{1.716554in}{0.619949in}}%
\pgfpathlineto{\pgfqpoint{1.718652in}{0.613717in}}%
\pgfpathlineto{\pgfqpoint{1.720749in}{0.601866in}}%
\pgfpathlineto{\pgfqpoint{1.722847in}{0.577582in}}%
\pgfpathlineto{\pgfqpoint{1.724944in}{0.478217in}}%
\pgfpathlineto{\pgfqpoint{1.727041in}{0.584997in}}%
\pgfpathlineto{\pgfqpoint{1.729139in}{0.610172in}}%
\pgfpathlineto{\pgfqpoint{1.731236in}{0.624306in}}%
\pgfpathlineto{\pgfqpoint{1.733333in}{0.633077in}}%
\pgfpathlineto{\pgfqpoint{1.735431in}{0.638076in}}%
\pgfpathlineto{\pgfqpoint{1.737528in}{0.639608in}}%
\pgfpathlineto{\pgfqpoint{1.739625in}{0.637023in}}%
\pgfpathlineto{\pgfqpoint{1.741723in}{0.629694in}}%
\pgfpathlineto{\pgfqpoint{1.743820in}{0.617696in}}%
\pgfpathlineto{\pgfqpoint{1.745918in}{0.597483in}}%
\pgfpathlineto{\pgfqpoint{1.748015in}{0.549940in}}%
\pgfpathlineto{\pgfqpoint{1.750112in}{0.564571in}}%
\pgfpathlineto{\pgfqpoint{1.752210in}{0.597839in}}%
\pgfpathlineto{\pgfqpoint{1.754307in}{0.612349in}}%
\pgfpathlineto{\pgfqpoint{1.756404in}{0.619924in}}%
\pgfpathlineto{\pgfqpoint{1.758502in}{0.623297in}}%
\pgfpathlineto{\pgfqpoint{1.760599in}{0.623318in}}%
\pgfpathlineto{\pgfqpoint{1.762696in}{0.619989in}}%
\pgfpathlineto{\pgfqpoint{1.764794in}{0.612467in}}%
\pgfpathlineto{\pgfqpoint{1.766891in}{0.598026in}}%
\pgfpathlineto{\pgfqpoint{1.768988in}{0.564749in}}%
\pgfpathlineto{\pgfqpoint{1.771086in}{0.548000in}}%
\pgfpathlineto{\pgfqpoint{1.773183in}{0.597269in}}%
\pgfpathlineto{\pgfqpoint{1.775281in}{0.617691in}}%
\pgfpathlineto{\pgfqpoint{1.777378in}{0.629802in}}%
\pgfpathlineto{\pgfqpoint{1.779475in}{0.637242in}}%
\pgfpathlineto{\pgfqpoint{1.781573in}{0.641064in}}%
\pgfpathlineto{\pgfqpoint{1.783670in}{0.641245in}}%
\pgfpathlineto{\pgfqpoint{1.785767in}{0.636789in}}%
\pgfpathlineto{\pgfqpoint{1.787865in}{0.628198in}}%
\pgfpathlineto{\pgfqpoint{1.789962in}{0.614822in}}%
\pgfpathlineto{\pgfqpoint{1.792059in}{0.592393in}}%
\pgfpathlineto{\pgfqpoint{1.794157in}{0.533430in}}%
\pgfpathlineto{\pgfqpoint{1.798352in}{0.596993in}}%
\pgfpathlineto{\pgfqpoint{1.800449in}{0.609701in}}%
\pgfpathlineto{\pgfqpoint{1.802546in}{0.615871in}}%
\pgfpathlineto{\pgfqpoint{1.804644in}{0.617787in}}%
\pgfpathlineto{\pgfqpoint{1.806741in}{0.616005in}}%
\pgfpathlineto{\pgfqpoint{1.808838in}{0.610069in}}%
\pgfpathlineto{\pgfqpoint{1.810936in}{0.597937in}}%
\pgfpathlineto{\pgfqpoint{1.813033in}{0.571543in}}%
\pgfpathlineto{\pgfqpoint{1.815130in}{0.513578in}}%
\pgfpathlineto{\pgfqpoint{1.817228in}{0.587578in}}%
\pgfpathlineto{\pgfqpoint{1.819325in}{0.611068in}}%
\pgfpathlineto{\pgfqpoint{1.821423in}{0.624534in}}%
\pgfpathlineto{\pgfqpoint{1.823520in}{0.632818in}}%
\pgfpathlineto{\pgfqpoint{1.825617in}{0.637280in}}%
\pgfpathlineto{\pgfqpoint{1.827715in}{0.638037in}}%
\pgfpathlineto{\pgfqpoint{1.829812in}{0.634142in}}%
\pgfpathlineto{\pgfqpoint{1.831909in}{0.626031in}}%
\pgfpathlineto{\pgfqpoint{1.834007in}{0.613553in}}%
\pgfpathlineto{\pgfqpoint{1.836104in}{0.593354in}}%
\pgfpathlineto{\pgfqpoint{1.838201in}{0.548661in}}%
\pgfpathlineto{\pgfqpoint{1.840299in}{0.552355in}}%
\pgfpathlineto{\pgfqpoint{1.842396in}{0.588092in}}%
\pgfpathlineto{\pgfqpoint{1.844494in}{0.602460in}}%
\pgfpathlineto{\pgfqpoint{1.846591in}{0.609346in}}%
\pgfpathlineto{\pgfqpoint{1.848688in}{0.611634in}}%
\pgfpathlineto{\pgfqpoint{1.850786in}{0.610043in}}%
\pgfpathlineto{\pgfqpoint{1.852883in}{0.604160in}}%
\pgfpathlineto{\pgfqpoint{1.854980in}{0.591889in}}%
\pgfpathlineto{\pgfqpoint{1.857078in}{0.564675in}}%
\pgfpathlineto{\pgfqpoint{1.859175in}{0.515572in}}%
\pgfpathlineto{\pgfqpoint{1.861272in}{0.583568in}}%
\pgfpathlineto{\pgfqpoint{1.863370in}{0.606554in}}%
\pgfpathlineto{\pgfqpoint{1.865467in}{0.619814in}}%
\pgfpathlineto{\pgfqpoint{1.867565in}{0.627954in}}%
\pgfpathlineto{\pgfqpoint{1.869662in}{0.632263in}}%
\pgfpathlineto{\pgfqpoint{1.871759in}{0.632805in}}%
\pgfpathlineto{\pgfqpoint{1.873857in}{0.628580in}}%
\pgfpathlineto{\pgfqpoint{1.875954in}{0.619992in}}%
\pgfpathlineto{\pgfqpoint{1.878051in}{0.606380in}}%
\pgfpathlineto{\pgfqpoint{1.880149in}{0.582978in}}%
\pgfpathlineto{\pgfqpoint{1.882246in}{0.512183in}}%
\pgfpathlineto{\pgfqpoint{1.884343in}{0.565574in}}%
\pgfpathlineto{\pgfqpoint{1.886441in}{0.592372in}}%
\pgfpathlineto{\pgfqpoint{1.888538in}{0.604699in}}%
\pgfpathlineto{\pgfqpoint{1.890636in}{0.610854in}}%
\pgfpathlineto{\pgfqpoint{1.892733in}{0.612954in}}%
\pgfpathlineto{\pgfqpoint{1.894830in}{0.611557in}}%
\pgfpathlineto{\pgfqpoint{1.896928in}{0.606322in}}%
\pgfpathlineto{\pgfqpoint{1.899025in}{0.595638in}}%
\pgfpathlineto{\pgfqpoint{1.901122in}{0.573736in}}%
\pgfpathlineto{\pgfqpoint{1.903220in}{0.460888in}}%
\pgfpathlineto{\pgfqpoint{1.905317in}{0.573769in}}%
\pgfpathlineto{\pgfqpoint{1.907414in}{0.601121in}}%
\pgfpathlineto{\pgfqpoint{1.909512in}{0.616070in}}%
\pgfpathlineto{\pgfqpoint{1.911609in}{0.625285in}}%
\pgfpathlineto{\pgfqpoint{1.913707in}{0.630550in}}%
\pgfpathlineto{\pgfqpoint{1.915804in}{0.632221in}}%
\pgfpathlineto{\pgfqpoint{1.917901in}{0.629629in}}%
\pgfpathlineto{\pgfqpoint{1.919999in}{0.622130in}}%
\pgfpathlineto{\pgfqpoint{1.922096in}{0.609633in}}%
\pgfpathlineto{\pgfqpoint{1.924193in}{0.587907in}}%
\pgfpathlineto{\pgfqpoint{1.926291in}{0.527960in}}%
\pgfpathlineto{\pgfqpoint{1.930485in}{0.595230in}}%
\pgfpathlineto{\pgfqpoint{1.932583in}{0.608838in}}%
\pgfpathlineto{\pgfqpoint{1.934680in}{0.616095in}}%
\pgfpathlineto{\pgfqpoint{1.936777in}{0.619393in}}%
\pgfpathlineto{\pgfqpoint{1.938875in}{0.619500in}}%
\pgfpathlineto{\pgfqpoint{1.940972in}{0.616427in}}%
\pgfpathlineto{\pgfqpoint{1.943070in}{0.609438in}}%
\pgfpathlineto{\pgfqpoint{1.945167in}{0.596223in}}%
\pgfpathlineto{\pgfqpoint{1.947264in}{0.567738in}}%
\pgfpathlineto{\pgfqpoint{1.949362in}{0.524049in}}%
\pgfpathlineto{\pgfqpoint{1.951459in}{0.587987in}}%
\pgfpathlineto{\pgfqpoint{1.953556in}{0.610390in}}%
\pgfpathlineto{\pgfqpoint{1.955654in}{0.623391in}}%
\pgfpathlineto{\pgfqpoint{1.957751in}{0.631445in}}%
\pgfpathlineto{\pgfqpoint{1.959848in}{0.635840in}}%
\pgfpathlineto{\pgfqpoint{1.961946in}{0.636720in}}%
\pgfpathlineto{\pgfqpoint{1.964043in}{0.633229in}}%
\pgfpathlineto{\pgfqpoint{1.966141in}{0.625033in}}%
\pgfpathlineto{\pgfqpoint{1.968238in}{0.611454in}}%
\pgfpathlineto{\pgfqpoint{1.970335in}{0.587121in}}%
\pgfpathlineto{\pgfqpoint{1.972433in}{0.493684in}}%
\pgfpathlineto{\pgfqpoint{1.974530in}{0.577644in}}%
\pgfpathlineto{\pgfqpoint{1.976627in}{0.602823in}}%
\pgfpathlineto{\pgfqpoint{1.978725in}{0.615071in}}%
\pgfpathlineto{\pgfqpoint{1.980822in}{0.621583in}}%
\pgfpathlineto{\pgfqpoint{1.982919in}{0.624341in}}%
\pgfpathlineto{\pgfqpoint{1.985017in}{0.623962in}}%
\pgfpathlineto{\pgfqpoint{1.987114in}{0.620353in}}%
\pgfpathlineto{\pgfqpoint{1.989212in}{0.612625in}}%
\pgfpathlineto{\pgfqpoint{1.991309in}{0.598023in}}%
\pgfpathlineto{\pgfqpoint{1.993406in}{0.564578in}}%
\pgfpathlineto{\pgfqpoint{1.995504in}{0.547902in}}%
\pgfpathlineto{\pgfqpoint{1.997601in}{0.597070in}}%
\pgfpathlineto{\pgfqpoint{1.999698in}{0.617445in}}%
\pgfpathlineto{\pgfqpoint{2.001796in}{0.629540in}}%
\pgfpathlineto{\pgfqpoint{2.003893in}{0.636997in}}%
\pgfpathlineto{\pgfqpoint{2.005990in}{0.640877in}}%
\pgfpathlineto{\pgfqpoint{2.008088in}{0.641190in}}%
\pgfpathlineto{\pgfqpoint{2.010185in}{0.636965in}}%
\pgfpathlineto{\pgfqpoint{2.012283in}{0.628439in}}%
\pgfpathlineto{\pgfqpoint{2.014380in}{0.614923in}}%
\pgfpathlineto{\pgfqpoint{2.016477in}{0.591747in}}%
\pgfpathlineto{\pgfqpoint{2.018575in}{0.523035in}}%
\pgfpathlineto{\pgfqpoint{2.020672in}{0.573176in}}%
\pgfpathlineto{\pgfqpoint{2.022769in}{0.600395in}}%
\pgfpathlineto{\pgfqpoint{2.024867in}{0.612853in}}%
\pgfpathlineto{\pgfqpoint{2.026964in}{0.619081in}}%
\pgfpathlineto{\pgfqpoint{2.029061in}{0.621237in}}%
\pgfpathlineto{\pgfqpoint{2.031159in}{0.619894in}}%
\pgfpathlineto{\pgfqpoint{2.033256in}{0.614726in}}%
\pgfpathlineto{\pgfqpoint{2.035354in}{0.604150in}}%
\pgfpathlineto{\pgfqpoint{2.037451in}{0.582528in}}%
\pgfpathlineto{\pgfqpoint{2.039548in}{0.479244in}}%
\pgfpathlineto{\pgfqpoint{2.041646in}{0.581433in}}%
\pgfpathlineto{\pgfqpoint{2.043743in}{0.609114in}}%
\pgfpathlineto{\pgfqpoint{2.045840in}{0.624179in}}%
\pgfpathlineto{\pgfqpoint{2.047938in}{0.633464in}}%
\pgfpathlineto{\pgfqpoint{2.050035in}{0.638784in}}%
\pgfpathlineto{\pgfqpoint{2.052132in}{0.640511in}}%
\pgfpathlineto{\pgfqpoint{2.054230in}{0.637995in}}%
\pgfpathlineto{\pgfqpoint{2.056327in}{0.630867in}}%
\pgfpathlineto{\pgfqpoint{2.058425in}{0.619825in}}%
\pgfpathlineto{\pgfqpoint{2.060522in}{0.602518in}}%
\pgfpathlineto{\pgfqpoint{2.062619in}{0.569669in}}%
\pgfpathlineto{\pgfqpoint{2.064717in}{0.524704in}}%
\pgfpathlineto{\pgfqpoint{2.066814in}{0.582824in}}%
\pgfpathlineto{\pgfqpoint{2.068911in}{0.600842in}}%
\pgfpathlineto{\pgfqpoint{2.071009in}{0.609413in}}%
\pgfpathlineto{\pgfqpoint{2.073106in}{0.612836in}}%
\pgfpathlineto{\pgfqpoint{2.075203in}{0.612253in}}%
\pgfpathlineto{\pgfqpoint{2.077301in}{0.607540in}}%
\pgfpathlineto{\pgfqpoint{2.079398in}{0.597155in}}%
\pgfpathlineto{\pgfqpoint{2.081496in}{0.575223in}}%
\pgfpathlineto{\pgfqpoint{2.083593in}{0.444407in}}%
\pgfpathlineto{\pgfqpoint{2.085690in}{0.577056in}}%
\pgfpathlineto{\pgfqpoint{2.087788in}{0.604142in}}%
\pgfpathlineto{\pgfqpoint{2.089885in}{0.619024in}}%
\pgfpathlineto{\pgfqpoint{2.091982in}{0.628176in}}%
\pgfpathlineto{\pgfqpoint{2.094080in}{0.633326in}}%
\pgfpathlineto{\pgfqpoint{2.096177in}{0.634770in}}%
\pgfpathlineto{\pgfqpoint{2.098274in}{0.631710in}}%
\pgfpathlineto{\pgfqpoint{2.100372in}{0.624103in}}%
\pgfpathlineto{\pgfqpoint{2.102469in}{0.612242in}}%
\pgfpathlineto{\pgfqpoint{2.104566in}{0.593114in}}%
\pgfpathlineto{\pgfqpoint{2.106664in}{0.552744in}}%
\pgfpathlineto{\pgfqpoint{2.108761in}{0.543832in}}%
\pgfpathlineto{\pgfqpoint{2.110859in}{0.584252in}}%
\pgfpathlineto{\pgfqpoint{2.112956in}{0.599763in}}%
\pgfpathlineto{\pgfqpoint{2.115053in}{0.607349in}}%
\pgfpathlineto{\pgfqpoint{2.117151in}{0.610266in}}%
\pgfpathlineto{\pgfqpoint{2.119248in}{0.609400in}}%
\pgfpathlineto{\pgfqpoint{2.121345in}{0.604547in}}%
\pgfpathlineto{\pgfqpoint{2.123443in}{0.594169in}}%
\pgfpathlineto{\pgfqpoint{2.125540in}{0.572596in}}%
\pgfpathlineto{\pgfqpoint{2.127637in}{0.464956in}}%
\pgfpathlineto{\pgfqpoint{2.129735in}{0.572297in}}%
\pgfpathlineto{\pgfqpoint{2.131832in}{0.599888in}}%
\pgfpathlineto{\pgfqpoint{2.133930in}{0.614937in}}%
\pgfpathlineto{\pgfqpoint{2.136027in}{0.624197in}}%
\pgfpathlineto{\pgfqpoint{2.138124in}{0.629459in}}%
\pgfpathlineto{\pgfqpoint{2.140222in}{0.631061in}}%
\pgfpathlineto{\pgfqpoint{2.142319in}{0.628281in}}%
\pgfpathlineto{\pgfqpoint{2.144416in}{0.620655in}}%
\pgfpathlineto{\pgfqpoint{2.146514in}{0.608157in}}%
\pgfpathlineto{\pgfqpoint{2.148611in}{0.586722in}}%
\pgfpathlineto{\pgfqpoint{2.150708in}{0.530205in}}%
\pgfpathlineto{\pgfqpoint{2.154903in}{0.591856in}}%
\pgfpathlineto{\pgfqpoint{2.157001in}{0.605577in}}%
\pgfpathlineto{\pgfqpoint{2.159098in}{0.612771in}}%
\pgfpathlineto{\pgfqpoint{2.161195in}{0.615909in}}%
\pgfpathlineto{\pgfqpoint{2.163293in}{0.615754in}}%
\pgfpathlineto{\pgfqpoint{2.165390in}{0.612269in}}%
\pgfpathlineto{\pgfqpoint{2.167487in}{0.604572in}}%
\pgfpathlineto{\pgfqpoint{2.169585in}{0.589848in}}%
\pgfpathlineto{\pgfqpoint{2.171682in}{0.555555in}}%
\pgfpathlineto{\pgfqpoint{2.173779in}{0.542658in}}%
\pgfpathlineto{\pgfqpoint{2.175877in}{0.590214in}}%
\pgfpathlineto{\pgfqpoint{2.177974in}{0.610324in}}%
\pgfpathlineto{\pgfqpoint{2.180072in}{0.622288in}}%
\pgfpathlineto{\pgfqpoint{2.182169in}{0.629633in}}%
\pgfpathlineto{\pgfqpoint{2.184266in}{0.633378in}}%
\pgfpathlineto{\pgfqpoint{2.186364in}{0.633484in}}%
\pgfpathlineto{\pgfqpoint{2.188461in}{0.628919in}}%
\pgfpathlineto{\pgfqpoint{2.190558in}{0.619642in}}%
\pgfpathlineto{\pgfqpoint{2.192656in}{0.604084in}}%
\pgfpathlineto{\pgfqpoint{2.194753in}{0.573860in}}%
\pgfpathlineto{\pgfqpoint{2.196850in}{0.518785in}}%
\pgfpathlineto{\pgfqpoint{2.198948in}{0.585508in}}%
\pgfpathlineto{\pgfqpoint{2.201045in}{0.605632in}}%
\pgfpathlineto{\pgfqpoint{2.203143in}{0.616082in}}%
\pgfpathlineto{\pgfqpoint{2.205240in}{0.621621in}}%
\pgfpathlineto{\pgfqpoint{2.207337in}{0.623720in}}%
\pgfpathlineto{\pgfqpoint{2.209435in}{0.622809in}}%
\pgfpathlineto{\pgfqpoint{2.211532in}{0.618681in}}%
\pgfpathlineto{\pgfqpoint{2.213629in}{0.610308in}}%
\pgfpathlineto{\pgfqpoint{2.215727in}{0.594559in}}%
\pgfpathlineto{\pgfqpoint{2.217824in}{0.556556in}}%
\pgfpathlineto{\pgfqpoint{2.219921in}{0.555744in}}%
\pgfpathlineto{\pgfqpoint{2.222019in}{0.598432in}}%
\pgfpathlineto{\pgfqpoint{2.224116in}{0.617570in}}%
\pgfpathlineto{\pgfqpoint{2.226214in}{0.629082in}}%
\pgfpathlineto{\pgfqpoint{2.228311in}{0.636158in}}%
\pgfpathlineto{\pgfqpoint{2.230408in}{0.639730in}}%
\pgfpathlineto{\pgfqpoint{2.232506in}{0.639732in}}%
\pgfpathlineto{\pgfqpoint{2.234603in}{0.635324in}}%
\pgfpathlineto{\pgfqpoint{2.249285in}{0.594943in}}%
\pgfpathlineto{\pgfqpoint{2.249285in}{0.594943in}}%
\pgfusepath{stroke}%
\end{pgfscope}%
\begin{pgfscope}%
\pgfpathrectangle{\pgfqpoint{0.487500in}{0.320625in}}{\pgfqpoint{1.845679in}{1.164375in}}%
\pgfusepath{clip}%
\pgfsetbuttcap%
\pgfsetroundjoin%
\pgfsetlinewidth{0.752812pt}%
\definecolor{currentstroke}{rgb}{0.819608,0.286275,0.356863}%
\pgfsetstrokecolor{currentstroke}%
\pgfsetstrokeopacity{0.650000}%
\pgfsetdash{{2.775000pt}{1.200000pt}}{0.000000pt}%
\pgfpathmoveto{\pgfqpoint{0.571395in}{1.468525in}}%
\pgfpathlineto{\pgfqpoint{1.862470in}{0.315625in}}%
\pgfpathlineto{\pgfqpoint{1.862470in}{0.315625in}}%
\pgfusepath{stroke}%
\end{pgfscope}%
\begin{pgfscope}%
\pgfpathrectangle{\pgfqpoint{0.487500in}{0.320625in}}{\pgfqpoint{1.845679in}{1.164375in}}%
\pgfusepath{clip}%
\pgfsetbuttcap%
\pgfsetroundjoin%
\pgfsetlinewidth{0.803000pt}%
\definecolor{currentstroke}{rgb}{0.819608,0.286275,0.356863}%
\pgfsetstrokecolor{currentstroke}%
\pgfsetstrokeopacity{0.800000}%
\pgfsetdash{{0.800000pt}{1.320000pt}}{0.000000pt}%
\pgfpathmoveto{\pgfqpoint{0.571395in}{1.468525in}}%
\pgfpathlineto{\pgfqpoint{2.249285in}{1.468525in}}%
\pgfpathlineto{\pgfqpoint{2.249285in}{1.468525in}}%
\pgfusepath{stroke}%
\end{pgfscope}%
\begin{pgfscope}%
\pgfpathrectangle{\pgfqpoint{0.487500in}{0.320625in}}{\pgfqpoint{1.845679in}{1.164375in}}%
\pgfusepath{clip}%
\pgfsetrectcap%
\pgfsetroundjoin%
\pgfsetlinewidth{1.254687pt}%
\definecolor{currentstroke}{rgb}{0.000000,0.529412,0.423529}%
\pgfsetstrokecolor{currentstroke}%
\pgfsetdash{}{0pt}%
\pgfpathmoveto{\pgfqpoint{0.571395in}{1.468525in}}%
\pgfpathlineto{\pgfqpoint{0.581881in}{1.449020in}}%
\pgfpathlineto{\pgfqpoint{0.592368in}{1.426689in}}%
\pgfpathlineto{\pgfqpoint{0.604952in}{1.396826in}}%
\pgfpathlineto{\pgfqpoint{0.646900in}{1.295102in}}%
\pgfpathlineto{\pgfqpoint{0.669971in}{1.242905in}}%
\pgfpathlineto{\pgfqpoint{0.697236in}{1.184202in}}%
\pgfpathlineto{\pgfqpoint{0.730794in}{1.114812in}}%
\pgfpathlineto{\pgfqpoint{0.774839in}{1.026516in}}%
\pgfpathlineto{\pgfqpoint{0.835662in}{0.907336in}}%
\pgfpathlineto{\pgfqpoint{0.911167in}{0.761770in}}%
\pgfpathlineto{\pgfqpoint{0.923751in}{0.737306in}}%
\pgfpathlineto{\pgfqpoint{0.936336in}{0.713371in}}%
\pgfpathlineto{\pgfqpoint{0.948920in}{0.690581in}}%
\pgfpathlineto{\pgfqpoint{0.967796in}{0.651376in}}%
\pgfpathlineto{\pgfqpoint{0.971991in}{0.645504in}}%
\pgfpathlineto{\pgfqpoint{0.978283in}{0.639267in}}%
\pgfpathlineto{\pgfqpoint{0.982478in}{0.634500in}}%
\pgfpathlineto{\pgfqpoint{0.986672in}{0.626257in}}%
\pgfpathlineto{\pgfqpoint{0.990867in}{0.615244in}}%
\pgfpathlineto{\pgfqpoint{0.995062in}{0.601390in}}%
\pgfpathlineto{\pgfqpoint{1.001354in}{0.575320in}}%
\pgfpathlineto{\pgfqpoint{1.005549in}{0.557112in}}%
\pgfpathlineto{\pgfqpoint{1.007646in}{0.549758in}}%
\pgfpathlineto{\pgfqpoint{1.009743in}{0.544705in}}%
\pgfpathlineto{\pgfqpoint{1.011841in}{0.542704in}}%
\pgfpathlineto{\pgfqpoint{1.013938in}{0.544411in}}%
\pgfpathlineto{\pgfqpoint{1.016035in}{0.549625in}}%
\pgfpathlineto{\pgfqpoint{1.024425in}{0.577878in}}%
\pgfpathlineto{\pgfqpoint{1.026522in}{0.582692in}}%
\pgfpathlineto{\pgfqpoint{1.028620in}{0.586081in}}%
\pgfpathlineto{\pgfqpoint{1.030717in}{0.587857in}}%
\pgfpathlineto{\pgfqpoint{1.032814in}{0.587683in}}%
\pgfpathlineto{\pgfqpoint{1.034912in}{0.584901in}}%
\pgfpathlineto{\pgfqpoint{1.037009in}{0.579145in}}%
\pgfpathlineto{\pgfqpoint{1.039106in}{0.570189in}}%
\pgfpathlineto{\pgfqpoint{1.043301in}{0.544015in}}%
\pgfpathlineto{\pgfqpoint{1.045398in}{0.547323in}}%
\pgfpathlineto{\pgfqpoint{1.049593in}{0.575543in}}%
\pgfpathlineto{\pgfqpoint{1.053788in}{0.592092in}}%
\pgfpathlineto{\pgfqpoint{1.055885in}{0.597191in}}%
\pgfpathlineto{\pgfqpoint{1.057983in}{0.600781in}}%
\pgfpathlineto{\pgfqpoint{1.060080in}{0.603095in}}%
\pgfpathlineto{\pgfqpoint{1.062177in}{0.604258in}}%
\pgfpathlineto{\pgfqpoint{1.064275in}{0.604321in}}%
\pgfpathlineto{\pgfqpoint{1.066372in}{0.603260in}}%
\pgfpathlineto{\pgfqpoint{1.068469in}{0.600976in}}%
\pgfpathlineto{\pgfqpoint{1.070567in}{0.597259in}}%
\pgfpathlineto{\pgfqpoint{1.072664in}{0.591720in}}%
\pgfpathlineto{\pgfqpoint{1.074762in}{0.583604in}}%
\pgfpathlineto{\pgfqpoint{1.076859in}{0.571251in}}%
\pgfpathlineto{\pgfqpoint{1.078956in}{0.549873in}}%
\pgfpathlineto{\pgfqpoint{1.081054in}{0.493255in}}%
\pgfpathlineto{\pgfqpoint{1.083151in}{0.539372in}}%
\pgfpathlineto{\pgfqpoint{1.085248in}{0.568614in}}%
\pgfpathlineto{\pgfqpoint{1.087346in}{0.584842in}}%
\pgfpathlineto{\pgfqpoint{1.089443in}{0.595751in}}%
\pgfpathlineto{\pgfqpoint{1.093638in}{0.609282in}}%
\pgfpathlineto{\pgfqpoint{1.095735in}{0.613259in}}%
\pgfpathlineto{\pgfqpoint{1.097832in}{0.615689in}}%
\pgfpathlineto{\pgfqpoint{1.099930in}{0.616564in}}%
\pgfpathlineto{\pgfqpoint{1.102027in}{0.615690in}}%
\pgfpathlineto{\pgfqpoint{1.104125in}{0.612592in}}%
\pgfpathlineto{\pgfqpoint{1.106222in}{0.606569in}}%
\pgfpathlineto{\pgfqpoint{1.108319in}{0.598461in}}%
\pgfpathlineto{\pgfqpoint{1.110417in}{0.593906in}}%
\pgfpathlineto{\pgfqpoint{1.112514in}{0.600064in}}%
\pgfpathlineto{\pgfqpoint{1.116709in}{0.621158in}}%
\pgfpathlineto{\pgfqpoint{1.118806in}{0.628774in}}%
\pgfpathlineto{\pgfqpoint{1.120903in}{0.634258in}}%
\pgfpathlineto{\pgfqpoint{1.123001in}{0.637912in}}%
\pgfpathlineto{\pgfqpoint{1.125098in}{0.639936in}}%
\pgfpathlineto{\pgfqpoint{1.127196in}{0.640415in}}%
\pgfpathlineto{\pgfqpoint{1.129293in}{0.639315in}}%
\pgfpathlineto{\pgfqpoint{1.131390in}{0.636460in}}%
\pgfpathlineto{\pgfqpoint{1.133488in}{0.631451in}}%
\pgfpathlineto{\pgfqpoint{1.135585in}{0.623448in}}%
\pgfpathlineto{\pgfqpoint{1.137682in}{0.610481in}}%
\pgfpathlineto{\pgfqpoint{1.139780in}{0.586220in}}%
\pgfpathlineto{\pgfqpoint{1.141877in}{0.512108in}}%
\pgfpathlineto{\pgfqpoint{1.143974in}{0.589469in}}%
\pgfpathlineto{\pgfqpoint{1.146072in}{0.615099in}}%
\pgfpathlineto{\pgfqpoint{1.148169in}{0.629752in}}%
\pgfpathlineto{\pgfqpoint{1.150267in}{0.639403in}}%
\pgfpathlineto{\pgfqpoint{1.152364in}{0.645871in}}%
\pgfpathlineto{\pgfqpoint{1.154461in}{0.649839in}}%
\pgfpathlineto{\pgfqpoint{1.156559in}{0.651462in}}%
\pgfpathlineto{\pgfqpoint{1.158656in}{0.650472in}}%
\pgfpathlineto{\pgfqpoint{1.160753in}{0.645992in}}%
\pgfpathlineto{\pgfqpoint{1.162851in}{0.638158in}}%
\pgfpathlineto{\pgfqpoint{1.164948in}{0.627381in}}%
\pgfpathlineto{\pgfqpoint{1.171240in}{0.586277in}}%
\pgfpathlineto{\pgfqpoint{1.175435in}{0.596800in}}%
\pgfpathlineto{\pgfqpoint{1.177532in}{0.598099in}}%
\pgfpathlineto{\pgfqpoint{1.179630in}{0.593628in}}%
\pgfpathlineto{\pgfqpoint{1.181727in}{0.581265in}}%
\pgfpathlineto{\pgfqpoint{1.183824in}{0.557891in}}%
\pgfpathlineto{\pgfqpoint{1.185922in}{0.569312in}}%
\pgfpathlineto{\pgfqpoint{1.188019in}{0.596174in}}%
\pgfpathlineto{\pgfqpoint{1.190116in}{0.612442in}}%
\pgfpathlineto{\pgfqpoint{1.192214in}{0.622199in}}%
\pgfpathlineto{\pgfqpoint{1.194311in}{0.627024in}}%
\pgfpathlineto{\pgfqpoint{1.196409in}{0.626754in}}%
\pgfpathlineto{\pgfqpoint{1.198506in}{0.620323in}}%
\pgfpathlineto{\pgfqpoint{1.200603in}{0.609975in}}%
\pgfpathlineto{\pgfqpoint{1.202701in}{0.594794in}}%
\pgfpathlineto{\pgfqpoint{1.204798in}{0.571500in}}%
\pgfpathlineto{\pgfqpoint{1.206895in}{0.541785in}}%
\pgfpathlineto{\pgfqpoint{1.211090in}{0.574294in}}%
\pgfpathlineto{\pgfqpoint{1.213187in}{0.580408in}}%
\pgfpathlineto{\pgfqpoint{1.215285in}{0.579511in}}%
\pgfpathlineto{\pgfqpoint{1.217382in}{0.571151in}}%
\pgfpathlineto{\pgfqpoint{1.219480in}{0.549383in}}%
\pgfpathlineto{\pgfqpoint{1.221577in}{0.505861in}}%
\pgfpathlineto{\pgfqpoint{1.223674in}{0.562596in}}%
\pgfpathlineto{\pgfqpoint{1.225772in}{0.586830in}}%
\pgfpathlineto{\pgfqpoint{1.227869in}{0.600376in}}%
\pgfpathlineto{\pgfqpoint{1.229966in}{0.607948in}}%
\pgfpathlineto{\pgfqpoint{1.232064in}{0.610526in}}%
\pgfpathlineto{\pgfqpoint{1.234161in}{0.607128in}}%
\pgfpathlineto{\pgfqpoint{1.236258in}{0.596490in}}%
\pgfpathlineto{\pgfqpoint{1.238356in}{0.575190in}}%
\pgfpathlineto{\pgfqpoint{1.240453in}{0.518887in}}%
\pgfpathlineto{\pgfqpoint{1.242551in}{0.564813in}}%
\pgfpathlineto{\pgfqpoint{1.244648in}{0.590984in}}%
\pgfpathlineto{\pgfqpoint{1.246745in}{0.604019in}}%
\pgfpathlineto{\pgfqpoint{1.248843in}{0.611096in}}%
\pgfpathlineto{\pgfqpoint{1.250940in}{0.614257in}}%
\pgfpathlineto{\pgfqpoint{1.253037in}{0.614145in}}%
\pgfpathlineto{\pgfqpoint{1.255135in}{0.610671in}}%
\pgfpathlineto{\pgfqpoint{1.257232in}{0.602913in}}%
\pgfpathlineto{\pgfqpoint{1.259329in}{0.587975in}}%
\pgfpathlineto{\pgfqpoint{1.261427in}{0.552779in}}%
\pgfpathlineto{\pgfqpoint{1.263524in}{0.542958in}}%
\pgfpathlineto{\pgfqpoint{1.265621in}{0.588865in}}%
\pgfpathlineto{\pgfqpoint{1.267719in}{0.608337in}}%
\pgfpathlineto{\pgfqpoint{1.269816in}{0.619596in}}%
\pgfpathlineto{\pgfqpoint{1.271914in}{0.625960in}}%
\pgfpathlineto{\pgfqpoint{1.274011in}{0.628180in}}%
\pgfpathlineto{\pgfqpoint{1.276108in}{0.625618in}}%
\pgfpathlineto{\pgfqpoint{1.278206in}{0.617139in}}%
\pgfpathlineto{\pgfqpoint{1.280303in}{0.602126in}}%
\pgfpathlineto{\pgfqpoint{1.282400in}{0.573232in}}%
\pgfpathlineto{\pgfqpoint{1.284498in}{0.521101in}}%
\pgfpathlineto{\pgfqpoint{1.286595in}{0.580605in}}%
\pgfpathlineto{\pgfqpoint{1.288692in}{0.600055in}}%
\pgfpathlineto{\pgfqpoint{1.290790in}{0.609306in}}%
\pgfpathlineto{\pgfqpoint{1.292887in}{0.613055in}}%
\pgfpathlineto{\pgfqpoint{1.294985in}{0.612523in}}%
\pgfpathlineto{\pgfqpoint{1.297082in}{0.607538in}}%
\pgfpathlineto{\pgfqpoint{1.299179in}{0.596292in}}%
\pgfpathlineto{\pgfqpoint{1.301277in}{0.571453in}}%
\pgfpathlineto{\pgfqpoint{1.303374in}{0.497944in}}%
\pgfpathlineto{\pgfqpoint{1.305471in}{0.583484in}}%
\pgfpathlineto{\pgfqpoint{1.307569in}{0.607486in}}%
\pgfpathlineto{\pgfqpoint{1.309666in}{0.620660in}}%
\pgfpathlineto{\pgfqpoint{1.311763in}{0.628077in}}%
\pgfpathlineto{\pgfqpoint{1.313861in}{0.630873in}}%
\pgfpathlineto{\pgfqpoint{1.315958in}{0.628433in}}%
\pgfpathlineto{\pgfqpoint{1.318056in}{0.619749in}}%
\pgfpathlineto{\pgfqpoint{1.320153in}{0.604338in}}%
\pgfpathlineto{\pgfqpoint{1.322250in}{0.574027in}}%
\pgfpathlineto{\pgfqpoint{1.324348in}{0.525127in}}%
\pgfpathlineto{\pgfqpoint{1.326445in}{0.585025in}}%
\pgfpathlineto{\pgfqpoint{1.328542in}{0.603920in}}%
\pgfpathlineto{\pgfqpoint{1.330640in}{0.612917in}}%
\pgfpathlineto{\pgfqpoint{1.332737in}{0.616500in}}%
\pgfpathlineto{\pgfqpoint{1.334834in}{0.615827in}}%
\pgfpathlineto{\pgfqpoint{1.336932in}{0.610682in}}%
\pgfpathlineto{\pgfqpoint{1.339029in}{0.599186in}}%
\pgfpathlineto{\pgfqpoint{1.341127in}{0.573637in}}%
\pgfpathlineto{\pgfqpoint{1.343224in}{0.508956in}}%
\pgfpathlineto{\pgfqpoint{1.345321in}{0.587782in}}%
\pgfpathlineto{\pgfqpoint{1.347419in}{0.611299in}}%
\pgfpathlineto{\pgfqpoint{1.349516in}{0.624263in}}%
\pgfpathlineto{\pgfqpoint{1.351613in}{0.631534in}}%
\pgfpathlineto{\pgfqpoint{1.353711in}{0.634186in}}%
\pgfpathlineto{\pgfqpoint{1.355808in}{0.631549in}}%
\pgfpathlineto{\pgfqpoint{1.357905in}{0.622684in}}%
\pgfpathlineto{\pgfqpoint{1.360003in}{0.606991in}}%
\pgfpathlineto{\pgfqpoint{1.362100in}{0.575767in}}%
\pgfpathlineto{\pgfqpoint{1.364198in}{0.530180in}}%
\pgfpathlineto{\pgfqpoint{1.366295in}{0.588908in}}%
\pgfpathlineto{\pgfqpoint{1.368392in}{0.607424in}}%
\pgfpathlineto{\pgfqpoint{1.370490in}{0.616224in}}%
\pgfpathlineto{\pgfqpoint{1.372587in}{0.619654in}}%
\pgfpathlineto{\pgfqpoint{1.374684in}{0.618819in}}%
\pgfpathlineto{\pgfqpoint{1.376782in}{0.613460in}}%
\pgfpathlineto{\pgfqpoint{1.378879in}{0.601576in}}%
\pgfpathlineto{\pgfqpoint{1.380976in}{0.574822in}}%
\pgfpathlineto{\pgfqpoint{1.383074in}{0.522727in}}%
\pgfpathlineto{\pgfqpoint{1.385171in}{0.592517in}}%
\pgfpathlineto{\pgfqpoint{1.387269in}{0.615277in}}%
\pgfpathlineto{\pgfqpoint{1.389366in}{0.627910in}}%
\pgfpathlineto{\pgfqpoint{1.391463in}{0.634945in}}%
\pgfpathlineto{\pgfqpoint{1.393561in}{0.637359in}}%
\pgfpathlineto{\pgfqpoint{1.395658in}{0.634381in}}%
\pgfpathlineto{\pgfqpoint{1.397755in}{0.625426in}}%
\pgfpathlineto{\pgfqpoint{1.399853in}{0.610210in}}%
\pgfpathlineto{\pgfqpoint{1.401950in}{0.581595in}}%
\pgfpathlineto{\pgfqpoint{1.404047in}{0.503027in}}%
\pgfpathlineto{\pgfqpoint{1.406145in}{0.584416in}}%
\pgfpathlineto{\pgfqpoint{1.408242in}{0.604393in}}%
\pgfpathlineto{\pgfqpoint{1.410340in}{0.613378in}}%
\pgfpathlineto{\pgfqpoint{1.412437in}{0.616514in}}%
\pgfpathlineto{\pgfqpoint{1.414534in}{0.614970in}}%
\pgfpathlineto{\pgfqpoint{1.416632in}{0.608238in}}%
\pgfpathlineto{\pgfqpoint{1.418729in}{0.593271in}}%
\pgfpathlineto{\pgfqpoint{1.420826in}{0.553621in}}%
\pgfpathlineto{\pgfqpoint{1.422924in}{0.560557in}}%
\pgfpathlineto{\pgfqpoint{1.425021in}{0.600949in}}%
\pgfpathlineto{\pgfqpoint{1.427118in}{0.619286in}}%
\pgfpathlineto{\pgfqpoint{1.429216in}{0.629746in}}%
\pgfpathlineto{\pgfqpoint{1.431313in}{0.635075in}}%
\pgfpathlineto{\pgfqpoint{1.433410in}{0.635556in}}%
\pgfpathlineto{\pgfqpoint{1.435508in}{0.629916in}}%
\pgfpathlineto{\pgfqpoint{1.437605in}{0.618970in}}%
\pgfpathlineto{\pgfqpoint{1.439703in}{0.600756in}}%
\pgfpathlineto{\pgfqpoint{1.441800in}{0.563343in}}%
\pgfpathlineto{\pgfqpoint{1.443897in}{0.543116in}}%
\pgfpathlineto{\pgfqpoint{1.445995in}{0.587263in}}%
\pgfpathlineto{\pgfqpoint{1.448092in}{0.602571in}}%
\pgfpathlineto{\pgfqpoint{1.450189in}{0.609031in}}%
\pgfpathlineto{\pgfqpoint{1.452287in}{0.609956in}}%
\pgfpathlineto{\pgfqpoint{1.454384in}{0.605753in}}%
\pgfpathlineto{\pgfqpoint{1.456481in}{0.594716in}}%
\pgfpathlineto{\pgfqpoint{1.458579in}{0.568766in}}%
\pgfpathlineto{\pgfqpoint{1.460676in}{0.514514in}}%
\pgfpathlineto{\pgfqpoint{1.462774in}{0.586349in}}%
\pgfpathlineto{\pgfqpoint{1.464871in}{0.609429in}}%
\pgfpathlineto{\pgfqpoint{1.466968in}{0.622181in}}%
\pgfpathlineto{\pgfqpoint{1.469066in}{0.629176in}}%
\pgfpathlineto{\pgfqpoint{1.471163in}{0.631324in}}%
\pgfpathlineto{\pgfqpoint{1.473260in}{0.627643in}}%
\pgfpathlineto{\pgfqpoint{1.475358in}{0.618126in}}%
\pgfpathlineto{\pgfqpoint{1.477455in}{0.602157in}}%
\pgfpathlineto{\pgfqpoint{1.479552in}{0.571692in}}%
\pgfpathlineto{\pgfqpoint{1.481650in}{0.510415in}}%
\pgfpathlineto{\pgfqpoint{1.483747in}{0.579181in}}%
\pgfpathlineto{\pgfqpoint{1.485845in}{0.597863in}}%
\pgfpathlineto{\pgfqpoint{1.487942in}{0.606137in}}%
\pgfpathlineto{\pgfqpoint{1.490039in}{0.608631in}}%
\pgfpathlineto{\pgfqpoint{1.492137in}{0.606311in}}%
\pgfpathlineto{\pgfqpoint{1.494234in}{0.598363in}}%
\pgfpathlineto{\pgfqpoint{1.496331in}{0.580673in}}%
\pgfpathlineto{\pgfqpoint{1.498429in}{0.523463in}}%
\pgfpathlineto{\pgfqpoint{1.502623in}{0.598690in}}%
\pgfpathlineto{\pgfqpoint{1.504721in}{0.614862in}}%
\pgfpathlineto{\pgfqpoint{1.506818in}{0.624091in}}%
\pgfpathlineto{\pgfqpoint{1.508916in}{0.628369in}}%
\pgfpathlineto{\pgfqpoint{1.511013in}{0.627569in}}%
\pgfpathlineto{\pgfqpoint{1.513110in}{0.620375in}}%
\pgfpathlineto{\pgfqpoint{1.515208in}{0.607057in}}%
\pgfpathlineto{\pgfqpoint{1.517305in}{0.582492in}}%
\pgfpathlineto{\pgfqpoint{1.519402in}{0.483718in}}%
\pgfpathlineto{\pgfqpoint{1.521500in}{0.573187in}}%
\pgfpathlineto{\pgfqpoint{1.523597in}{0.597209in}}%
\pgfpathlineto{\pgfqpoint{1.525694in}{0.608141in}}%
\pgfpathlineto{\pgfqpoint{1.527792in}{0.612893in}}%
\pgfpathlineto{\pgfqpoint{1.529889in}{0.613168in}}%
\pgfpathlineto{\pgfqpoint{1.531987in}{0.609044in}}%
\pgfpathlineto{\pgfqpoint{1.534084in}{0.599085in}}%
\pgfpathlineto{\pgfqpoint{1.536181in}{0.577573in}}%
\pgfpathlineto{\pgfqpoint{1.538279in}{0.460729in}}%
\pgfpathlineto{\pgfqpoint{1.540376in}{0.578239in}}%
\pgfpathlineto{\pgfqpoint{1.542473in}{0.605243in}}%
\pgfpathlineto{\pgfqpoint{1.544571in}{0.619589in}}%
\pgfpathlineto{\pgfqpoint{1.546668in}{0.627756in}}%
\pgfpathlineto{\pgfqpoint{1.548765in}{0.631216in}}%
\pgfpathlineto{\pgfqpoint{1.550863in}{0.629597in}}%
\pgfpathlineto{\pgfqpoint{1.552960in}{0.621621in}}%
\pgfpathlineto{\pgfqpoint{1.555058in}{0.606992in}}%
\pgfpathlineto{\pgfqpoint{1.557155in}{0.578403in}}%
\pgfpathlineto{\pgfqpoint{1.559252in}{0.505664in}}%
\pgfpathlineto{\pgfqpoint{1.561350in}{0.584614in}}%
\pgfpathlineto{\pgfqpoint{1.563447in}{0.604932in}}%
\pgfpathlineto{\pgfqpoint{1.565544in}{0.614639in}}%
\pgfpathlineto{\pgfqpoint{1.567642in}{0.618805in}}%
\pgfpathlineto{\pgfqpoint{1.569739in}{0.618765in}}%
\pgfpathlineto{\pgfqpoint{1.571836in}{0.614490in}}%
\pgfpathlineto{\pgfqpoint{1.573934in}{0.604537in}}%
\pgfpathlineto{\pgfqpoint{1.576031in}{0.583390in}}%
\pgfpathlineto{\pgfqpoint{1.578129in}{0.483886in}}%
\pgfpathlineto{\pgfqpoint{1.580226in}{0.581892in}}%
\pgfpathlineto{\pgfqpoint{1.582323in}{0.609450in}}%
\pgfpathlineto{\pgfqpoint{1.584421in}{0.624006in}}%
\pgfpathlineto{\pgfqpoint{1.586518in}{0.632339in}}%
\pgfpathlineto{\pgfqpoint{1.588615in}{0.636010in}}%
\pgfpathlineto{\pgfqpoint{1.590713in}{0.634750in}}%
\pgfpathlineto{\pgfqpoint{1.592810in}{0.627298in}}%
\pgfpathlineto{\pgfqpoint{1.594907in}{0.613835in}}%
\pgfpathlineto{\pgfqpoint{1.597005in}{0.589186in}}%
\pgfpathlineto{\pgfqpoint{1.599102in}{0.491069in}}%
\pgfpathlineto{\pgfqpoint{1.601199in}{0.579340in}}%
\pgfpathlineto{\pgfqpoint{1.603297in}{0.603296in}}%
\pgfpathlineto{\pgfqpoint{1.605394in}{0.614084in}}%
\pgfpathlineto{\pgfqpoint{1.607492in}{0.618635in}}%
\pgfpathlineto{\pgfqpoint{1.609589in}{0.618627in}}%
\pgfpathlineto{\pgfqpoint{1.611686in}{0.614058in}}%
\pgfpathlineto{\pgfqpoint{1.613784in}{0.603255in}}%
\pgfpathlineto{\pgfqpoint{1.615881in}{0.579293in}}%
\pgfpathlineto{\pgfqpoint{1.617978in}{0.490331in}}%
\pgfpathlineto{\pgfqpoint{1.620076in}{0.589013in}}%
\pgfpathlineto{\pgfqpoint{1.622173in}{0.613638in}}%
\pgfpathlineto{\pgfqpoint{1.624270in}{0.627052in}}%
\pgfpathlineto{\pgfqpoint{1.626368in}{0.634588in}}%
\pgfpathlineto{\pgfqpoint{1.628465in}{0.637427in}}%
\pgfpathlineto{\pgfqpoint{1.630563in}{0.634942in}}%
\pgfpathlineto{\pgfqpoint{1.632660in}{0.626453in}}%
\pgfpathlineto{\pgfqpoint{1.634757in}{0.612285in}}%
\pgfpathlineto{\pgfqpoint{1.636855in}{0.587049in}}%
\pgfpathlineto{\pgfqpoint{1.638952in}{0.489553in}}%
\pgfpathlineto{\pgfqpoint{1.641049in}{0.575719in}}%
\pgfpathlineto{\pgfqpoint{1.643147in}{0.599045in}}%
\pgfpathlineto{\pgfqpoint{1.645244in}{0.609022in}}%
\pgfpathlineto{\pgfqpoint{1.647341in}{0.612514in}}%
\pgfpathlineto{\pgfqpoint{1.649439in}{0.610990in}}%
\pgfpathlineto{\pgfqpoint{1.651536in}{0.603929in}}%
\pgfpathlineto{\pgfqpoint{1.653634in}{0.587851in}}%
\pgfpathlineto{\pgfqpoint{1.655731in}{0.541136in}}%
\pgfpathlineto{\pgfqpoint{1.657828in}{0.565070in}}%
\pgfpathlineto{\pgfqpoint{1.659926in}{0.601003in}}%
\pgfpathlineto{\pgfqpoint{1.662023in}{0.618235in}}%
\pgfpathlineto{\pgfqpoint{1.664120in}{0.628067in}}%
\pgfpathlineto{\pgfqpoint{1.666218in}{0.632828in}}%
\pgfpathlineto{\pgfqpoint{1.668315in}{0.632559in}}%
\pgfpathlineto{\pgfqpoint{1.670412in}{0.626010in}}%
\pgfpathlineto{\pgfqpoint{1.672510in}{0.614202in}}%
\pgfpathlineto{\pgfqpoint{1.674607in}{0.594336in}}%
\pgfpathlineto{\pgfqpoint{1.676705in}{0.549811in}}%
\pgfpathlineto{\pgfqpoint{1.678802in}{0.552263in}}%
\pgfpathlineto{\pgfqpoint{1.680899in}{0.587838in}}%
\pgfpathlineto{\pgfqpoint{1.682997in}{0.601344in}}%
\pgfpathlineto{\pgfqpoint{1.685094in}{0.606847in}}%
\pgfpathlineto{\pgfqpoint{1.687191in}{0.606995in}}%
\pgfpathlineto{\pgfqpoint{1.689289in}{0.601903in}}%
\pgfpathlineto{\pgfqpoint{1.691386in}{0.589372in}}%
\pgfpathlineto{\pgfqpoint{1.695581in}{0.535919in}}%
\pgfpathlineto{\pgfqpoint{1.697678in}{0.588575in}}%
\pgfpathlineto{\pgfqpoint{1.699776in}{0.609337in}}%
\pgfpathlineto{\pgfqpoint{1.701873in}{0.621017in}}%
\pgfpathlineto{\pgfqpoint{1.703970in}{0.627228in}}%
\pgfpathlineto{\pgfqpoint{1.706068in}{0.628561in}}%
\pgfpathlineto{\pgfqpoint{1.708165in}{0.623812in}}%
\pgfpathlineto{\pgfqpoint{1.710262in}{0.613222in}}%
\pgfpathlineto{\pgfqpoint{1.712360in}{0.594919in}}%
\pgfpathlineto{\pgfqpoint{1.714457in}{0.555193in}}%
\pgfpathlineto{\pgfqpoint{1.716554in}{0.545737in}}%
\pgfpathlineto{\pgfqpoint{1.718652in}{0.586335in}}%
\pgfpathlineto{\pgfqpoint{1.720749in}{0.601539in}}%
\pgfpathlineto{\pgfqpoint{1.722847in}{0.608471in}}%
\pgfpathlineto{\pgfqpoint{1.724944in}{0.610283in}}%
\pgfpathlineto{\pgfqpoint{1.727041in}{0.607590in}}%
\pgfpathlineto{\pgfqpoint{1.729139in}{0.599482in}}%
\pgfpathlineto{\pgfqpoint{1.731236in}{0.581931in}}%
\pgfpathlineto{\pgfqpoint{1.733333in}{0.527494in}}%
\pgfpathlineto{\pgfqpoint{1.737528in}{0.598354in}}%
\pgfpathlineto{\pgfqpoint{1.739625in}{0.614728in}}%
\pgfpathlineto{\pgfqpoint{1.741723in}{0.624103in}}%
\pgfpathlineto{\pgfqpoint{1.743820in}{0.628571in}}%
\pgfpathlineto{\pgfqpoint{1.745918in}{0.628105in}}%
\pgfpathlineto{\pgfqpoint{1.748015in}{0.621314in}}%
\pgfpathlineto{\pgfqpoint{1.750112in}{0.608168in}}%
\pgfpathlineto{\pgfqpoint{1.752210in}{0.583449in}}%
\pgfpathlineto{\pgfqpoint{1.754307in}{0.470314in}}%
\pgfpathlineto{\pgfqpoint{1.756404in}{0.576517in}}%
\pgfpathlineto{\pgfqpoint{1.758502in}{0.600301in}}%
\pgfpathlineto{\pgfqpoint{1.760599in}{0.611382in}}%
\pgfpathlineto{\pgfqpoint{1.762696in}{0.616440in}}%
\pgfpathlineto{\pgfqpoint{1.764794in}{0.617188in}}%
\pgfpathlineto{\pgfqpoint{1.766891in}{0.613818in}}%
\pgfpathlineto{\pgfqpoint{1.768988in}{0.605252in}}%
\pgfpathlineto{\pgfqpoint{1.771086in}{0.587391in}}%
\pgfpathlineto{\pgfqpoint{1.773183in}{0.532883in}}%
\pgfpathlineto{\pgfqpoint{1.777378in}{0.603236in}}%
\pgfpathlineto{\pgfqpoint{1.779475in}{0.619584in}}%
\pgfpathlineto{\pgfqpoint{1.781573in}{0.629002in}}%
\pgfpathlineto{\pgfqpoint{1.783670in}{0.633618in}}%
\pgfpathlineto{\pgfqpoint{1.785767in}{0.633507in}}%
\pgfpathlineto{\pgfqpoint{1.787865in}{0.627303in}}%
\pgfpathlineto{\pgfqpoint{1.789962in}{0.614991in}}%
\pgfpathlineto{\pgfqpoint{1.792059in}{0.592476in}}%
\pgfpathlineto{\pgfqpoint{1.794157in}{0.522780in}}%
\pgfpathlineto{\pgfqpoint{1.796254in}{0.575188in}}%
\pgfpathlineto{\pgfqpoint{1.798352in}{0.602118in}}%
\pgfpathlineto{\pgfqpoint{1.800449in}{0.614181in}}%
\pgfpathlineto{\pgfqpoint{1.802546in}{0.619735in}}%
\pgfpathlineto{\pgfqpoint{1.804644in}{0.620801in}}%
\pgfpathlineto{\pgfqpoint{1.806741in}{0.617679in}}%
\pgfpathlineto{\pgfqpoint{1.808838in}{0.609361in}}%
\pgfpathlineto{\pgfqpoint{1.810936in}{0.591896in}}%
\pgfpathlineto{\pgfqpoint{1.813033in}{0.539807in}}%
\pgfpathlineto{\pgfqpoint{1.817228in}{0.606738in}}%
\pgfpathlineto{\pgfqpoint{1.819325in}{0.623313in}}%
\pgfpathlineto{\pgfqpoint{1.821423in}{0.632847in}}%
\pgfpathlineto{\pgfqpoint{1.823520in}{0.637534in}}%
\pgfpathlineto{\pgfqpoint{1.825617in}{0.637463in}}%
\pgfpathlineto{\pgfqpoint{1.827715in}{0.631334in}}%
\pgfpathlineto{\pgfqpoint{1.829812in}{0.619718in}}%
\pgfpathlineto{\pgfqpoint{1.831909in}{0.599799in}}%
\pgfpathlineto{\pgfqpoint{1.834007in}{0.553664in}}%
\pgfpathlineto{\pgfqpoint{1.836104in}{0.561362in}}%
\pgfpathlineto{\pgfqpoint{1.838201in}{0.595812in}}%
\pgfpathlineto{\pgfqpoint{1.840299in}{0.609373in}}%
\pgfpathlineto{\pgfqpoint{1.842396in}{0.615205in}}%
\pgfpathlineto{\pgfqpoint{1.844494in}{0.615926in}}%
\pgfpathlineto{\pgfqpoint{1.846591in}{0.611819in}}%
\pgfpathlineto{\pgfqpoint{1.848688in}{0.601317in}}%
\pgfpathlineto{\pgfqpoint{1.850786in}{0.577510in}}%
\pgfpathlineto{\pgfqpoint{1.852883in}{0.490395in}}%
\pgfpathlineto{\pgfqpoint{1.854980in}{0.587761in}}%
\pgfpathlineto{\pgfqpoint{1.857078in}{0.612395in}}%
\pgfpathlineto{\pgfqpoint{1.859175in}{0.625809in}}%
\pgfpathlineto{\pgfqpoint{1.861272in}{0.633296in}}%
\pgfpathlineto{\pgfqpoint{1.863370in}{0.635996in}}%
\pgfpathlineto{\pgfqpoint{1.865467in}{0.633177in}}%
\pgfpathlineto{\pgfqpoint{1.867565in}{0.624465in}}%
\pgfpathlineto{\pgfqpoint{1.869662in}{0.610195in}}%
\pgfpathlineto{\pgfqpoint{1.871759in}{0.585104in}}%
\pgfpathlineto{\pgfqpoint{1.873857in}{0.495051in}}%
\pgfpathlineto{\pgfqpoint{1.875954in}{0.571735in}}%
\pgfpathlineto{\pgfqpoint{1.878051in}{0.595285in}}%
\pgfpathlineto{\pgfqpoint{1.880149in}{0.605111in}}%
\pgfpathlineto{\pgfqpoint{1.882246in}{0.608279in}}%
\pgfpathlineto{\pgfqpoint{1.884343in}{0.606210in}}%
\pgfpathlineto{\pgfqpoint{1.886441in}{0.598145in}}%
\pgfpathlineto{\pgfqpoint{1.888538in}{0.579560in}}%
\pgfpathlineto{\pgfqpoint{1.890636in}{0.511016in}}%
\pgfpathlineto{\pgfqpoint{1.892733in}{0.571473in}}%
\pgfpathlineto{\pgfqpoint{1.894830in}{0.601756in}}%
\pgfpathlineto{\pgfqpoint{1.896928in}{0.617290in}}%
\pgfpathlineto{\pgfqpoint{1.899025in}{0.626107in}}%
\pgfpathlineto{\pgfqpoint{1.901122in}{0.629952in}}%
\pgfpathlineto{\pgfqpoint{1.903220in}{0.628495in}}%
\pgfpathlineto{\pgfqpoint{1.905317in}{0.620707in}}%
\pgfpathlineto{\pgfqpoint{1.907414in}{0.607309in}}%
\pgfpathlineto{\pgfqpoint{1.909512in}{0.583660in}}%
\pgfpathlineto{\pgfqpoint{1.911609in}{0.509584in}}%
\pgfpathlineto{\pgfqpoint{1.913707in}{0.566360in}}%
\pgfpathlineto{\pgfqpoint{1.915804in}{0.591920in}}%
\pgfpathlineto{\pgfqpoint{1.917901in}{0.602831in}}%
\pgfpathlineto{\pgfqpoint{1.919999in}{0.607053in}}%
\pgfpathlineto{\pgfqpoint{1.922096in}{0.606340in}}%
\pgfpathlineto{\pgfqpoint{1.924193in}{0.600476in}}%
\pgfpathlineto{\pgfqpoint{1.926291in}{0.586903in}}%
\pgfpathlineto{\pgfqpoint{1.928388in}{0.552808in}}%
\pgfpathlineto{\pgfqpoint{1.930485in}{0.542936in}}%
\pgfpathlineto{\pgfqpoint{1.932583in}{0.589568in}}%
\pgfpathlineto{\pgfqpoint{1.934680in}{0.609244in}}%
\pgfpathlineto{\pgfqpoint{1.936777in}{0.620391in}}%
\pgfpathlineto{\pgfqpoint{1.938875in}{0.626230in}}%
\pgfpathlineto{\pgfqpoint{1.940972in}{0.627223in}}%
\pgfpathlineto{\pgfqpoint{1.943070in}{0.622101in}}%
\pgfpathlineto{\pgfqpoint{1.945167in}{0.610874in}}%
\pgfpathlineto{\pgfqpoint{1.947264in}{0.590717in}}%
\pgfpathlineto{\pgfqpoint{1.949362in}{0.539871in}}%
\pgfpathlineto{\pgfqpoint{1.951459in}{0.560303in}}%
\pgfpathlineto{\pgfqpoint{1.953556in}{0.592392in}}%
\pgfpathlineto{\pgfqpoint{1.955654in}{0.605973in}}%
\pgfpathlineto{\pgfqpoint{1.957751in}{0.612375in}}%
\pgfpathlineto{\pgfqpoint{1.959848in}{0.614103in}}%
\pgfpathlineto{\pgfqpoint{1.961946in}{0.611662in}}%
\pgfpathlineto{\pgfqpoint{1.964043in}{0.604282in}}%
\pgfpathlineto{\pgfqpoint{1.966141in}{0.588745in}}%
\pgfpathlineto{\pgfqpoint{1.968238in}{0.547950in}}%
\pgfpathlineto{\pgfqpoint{1.970335in}{0.556936in}}%
\pgfpathlineto{\pgfqpoint{1.972433in}{0.596512in}}%
\pgfpathlineto{\pgfqpoint{1.974530in}{0.614621in}}%
\pgfpathlineto{\pgfqpoint{1.976627in}{0.625002in}}%
\pgfpathlineto{\pgfqpoint{1.978725in}{0.630373in}}%
\pgfpathlineto{\pgfqpoint{1.980822in}{0.631070in}}%
\pgfpathlineto{\pgfqpoint{1.982919in}{0.625807in}}%
\pgfpathlineto{\pgfqpoint{1.985017in}{0.614284in}}%
\pgfpathlineto{\pgfqpoint{1.987114in}{0.593090in}}%
\pgfpathlineto{\pgfqpoint{1.989212in}{0.533454in}}%
\pgfpathlineto{\pgfqpoint{1.993406in}{0.599989in}}%
\pgfpathlineto{\pgfqpoint{1.995504in}{0.612946in}}%
\pgfpathlineto{\pgfqpoint{1.997601in}{0.619182in}}%
\pgfpathlineto{\pgfqpoint{1.999698in}{0.620956in}}%
\pgfpathlineto{\pgfqpoint{2.001796in}{0.618744in}}%
\pgfpathlineto{\pgfqpoint{2.003893in}{0.611879in}}%
\pgfpathlineto{\pgfqpoint{2.005990in}{0.597600in}}%
\pgfpathlineto{\pgfqpoint{2.008088in}{0.562977in}}%
\pgfpathlineto{\pgfqpoint{2.010185in}{0.552862in}}%
\pgfpathlineto{\pgfqpoint{2.012283in}{0.599215in}}%
\pgfpathlineto{\pgfqpoint{2.014380in}{0.618761in}}%
\pgfpathlineto{\pgfqpoint{2.016477in}{0.629907in}}%
\pgfpathlineto{\pgfqpoint{2.018575in}{0.635915in}}%
\pgfpathlineto{\pgfqpoint{2.020672in}{0.637371in}}%
\pgfpathlineto{\pgfqpoint{2.022769in}{0.633202in}}%
\pgfpathlineto{\pgfqpoint{2.024867in}{0.623171in}}%
\pgfpathlineto{\pgfqpoint{2.026964in}{0.605782in}}%
\pgfpathlineto{\pgfqpoint{2.029061in}{0.569380in}}%
\pgfpathlineto{\pgfqpoint{2.031159in}{0.547953in}}%
\pgfpathlineto{\pgfqpoint{2.033256in}{0.593805in}}%
\pgfpathlineto{\pgfqpoint{2.035354in}{0.610075in}}%
\pgfpathlineto{\pgfqpoint{2.037451in}{0.617609in}}%
\pgfpathlineto{\pgfqpoint{2.039548in}{0.619939in}}%
\pgfpathlineto{\pgfqpoint{2.041646in}{0.617843in}}%
\pgfpathlineto{\pgfqpoint{2.043743in}{0.610639in}}%
\pgfpathlineto{\pgfqpoint{2.045840in}{0.595073in}}%
\pgfpathlineto{\pgfqpoint{2.047938in}{0.553240in}}%
\pgfpathlineto{\pgfqpoint{2.050035in}{0.565315in}}%
\pgfpathlineto{\pgfqpoint{2.052132in}{0.604081in}}%
\pgfpathlineto{\pgfqpoint{2.054230in}{0.622007in}}%
\pgfpathlineto{\pgfqpoint{2.056327in}{0.632263in}}%
\pgfpathlineto{\pgfqpoint{2.058425in}{0.637473in}}%
\pgfpathlineto{\pgfqpoint{2.060522in}{0.637894in}}%
\pgfpathlineto{\pgfqpoint{2.062619in}{0.632262in}}%
\pgfpathlineto{\pgfqpoint{2.064717in}{0.621404in}}%
\pgfpathlineto{\pgfqpoint{2.066814in}{0.603476in}}%
\pgfpathlineto{\pgfqpoint{2.071009in}{0.540382in}}%
\pgfpathlineto{\pgfqpoint{2.073106in}{0.587747in}}%
\pgfpathlineto{\pgfqpoint{2.075203in}{0.603460in}}%
\pgfpathlineto{\pgfqpoint{2.077301in}{0.610014in}}%
\pgfpathlineto{\pgfqpoint{2.079398in}{0.610902in}}%
\pgfpathlineto{\pgfqpoint{2.081496in}{0.606540in}}%
\pgfpathlineto{\pgfqpoint{2.083593in}{0.595098in}}%
\pgfpathlineto{\pgfqpoint{2.085690in}{0.567674in}}%
\pgfpathlineto{\pgfqpoint{2.087788in}{0.526111in}}%
\pgfpathlineto{\pgfqpoint{2.089885in}{0.589570in}}%
\pgfpathlineto{\pgfqpoint{2.091982in}{0.611814in}}%
\pgfpathlineto{\pgfqpoint{2.094080in}{0.624192in}}%
\pgfpathlineto{\pgfqpoint{2.096177in}{0.630906in}}%
\pgfpathlineto{\pgfqpoint{2.098274in}{0.632743in}}%
\pgfpathlineto{\pgfqpoint{2.100372in}{0.628621in}}%
\pgfpathlineto{\pgfqpoint{2.102469in}{0.618940in}}%
\pgfpathlineto{\pgfqpoint{2.104566in}{0.603080in}}%
\pgfpathlineto{\pgfqpoint{2.106664in}{0.573715in}}%
\pgfpathlineto{\pgfqpoint{2.108761in}{0.492587in}}%
\pgfpathlineto{\pgfqpoint{2.110859in}{0.576036in}}%
\pgfpathlineto{\pgfqpoint{2.112956in}{0.595296in}}%
\pgfpathlineto{\pgfqpoint{2.115053in}{0.603469in}}%
\pgfpathlineto{\pgfqpoint{2.117151in}{0.605547in}}%
\pgfpathlineto{\pgfqpoint{2.119248in}{0.602437in}}%
\pgfpathlineto{\pgfqpoint{2.121345in}{0.592929in}}%
\pgfpathlineto{\pgfqpoint{2.123443in}{0.570951in}}%
\pgfpathlineto{\pgfqpoint{2.125540in}{0.409298in}}%
\pgfpathlineto{\pgfqpoint{2.127637in}{0.576375in}}%
\pgfpathlineto{\pgfqpoint{2.129735in}{0.602448in}}%
\pgfpathlineto{\pgfqpoint{2.131832in}{0.616449in}}%
\pgfpathlineto{\pgfqpoint{2.133930in}{0.624288in}}%
\pgfpathlineto{\pgfqpoint{2.136027in}{0.627236in}}%
\pgfpathlineto{\pgfqpoint{2.138124in}{0.624608in}}%
\pgfpathlineto{\pgfqpoint{2.140222in}{0.615690in}}%
\pgfpathlineto{\pgfqpoint{2.142319in}{0.600202in}}%
\pgfpathlineto{\pgfqpoint{2.144416in}{0.570136in}}%
\pgfpathlineto{\pgfqpoint{2.146514in}{0.509046in}}%
\pgfpathlineto{\pgfqpoint{2.148611in}{0.578584in}}%
\pgfpathlineto{\pgfqpoint{2.150708in}{0.597761in}}%
\pgfpathlineto{\pgfqpoint{2.152806in}{0.606637in}}%
\pgfpathlineto{\pgfqpoint{2.154903in}{0.609921in}}%
\pgfpathlineto{\pgfqpoint{2.157001in}{0.608751in}}%
\pgfpathlineto{\pgfqpoint{2.159098in}{0.602756in}}%
\pgfpathlineto{\pgfqpoint{2.161195in}{0.589493in}}%
\pgfpathlineto{\pgfqpoint{2.163293in}{0.557631in}}%
\pgfpathlineto{\pgfqpoint{2.165390in}{0.537711in}}%
\pgfpathlineto{\pgfqpoint{2.167487in}{0.588689in}}%
\pgfpathlineto{\pgfqpoint{2.169585in}{0.609095in}}%
\pgfpathlineto{\pgfqpoint{2.171682in}{0.620624in}}%
\pgfpathlineto{\pgfqpoint{2.173779in}{0.626821in}}%
\pgfpathlineto{\pgfqpoint{2.175877in}{0.628316in}}%
\pgfpathlineto{\pgfqpoint{2.177974in}{0.623982in}}%
\pgfpathlineto{\pgfqpoint{2.180072in}{0.613312in}}%
\pgfpathlineto{\pgfqpoint{2.182169in}{0.593749in}}%
\pgfpathlineto{\pgfqpoint{2.184266in}{0.544610in}}%
\pgfpathlineto{\pgfqpoint{2.186364in}{0.562441in}}%
\pgfpathlineto{\pgfqpoint{2.188461in}{0.595527in}}%
\pgfpathlineto{\pgfqpoint{2.190558in}{0.609647in}}%
\pgfpathlineto{\pgfqpoint{2.192656in}{0.616578in}}%
\pgfpathlineto{\pgfqpoint{2.194753in}{0.618936in}}%
\pgfpathlineto{\pgfqpoint{2.196850in}{0.617364in}}%
\pgfpathlineto{\pgfqpoint{2.198948in}{0.611390in}}%
\pgfpathlineto{\pgfqpoint{2.201045in}{0.598814in}}%
\pgfpathlineto{\pgfqpoint{2.203143in}{0.570490in}}%
\pgfpathlineto{\pgfqpoint{2.205240in}{0.529213in}}%
\pgfpathlineto{\pgfqpoint{2.207337in}{0.591766in}}%
\pgfpathlineto{\pgfqpoint{2.209435in}{0.613731in}}%
\pgfpathlineto{\pgfqpoint{2.211532in}{0.626030in}}%
\pgfpathlineto{\pgfqpoint{2.213629in}{0.632892in}}%
\pgfpathlineto{\pgfqpoint{2.215727in}{0.635236in}}%
\pgfpathlineto{\pgfqpoint{2.217824in}{0.632313in}}%
\pgfpathlineto{\pgfqpoint{2.222019in}{0.615408in}}%
\pgfpathlineto{\pgfqpoint{2.226214in}{0.589430in}}%
\pgfpathlineto{\pgfqpoint{2.230408in}{0.551389in}}%
\pgfpathlineto{\pgfqpoint{2.232506in}{0.531114in}}%
\pgfpathlineto{\pgfqpoint{2.234603in}{0.523148in}}%
\pgfpathlineto{\pgfqpoint{2.236700in}{0.534327in}}%
\pgfpathlineto{\pgfqpoint{2.240895in}{0.567592in}}%
\pgfpathlineto{\pgfqpoint{2.242992in}{0.579213in}}%
\pgfpathlineto{\pgfqpoint{2.245090in}{0.586774in}}%
\pgfpathlineto{\pgfqpoint{2.247187in}{0.589959in}}%
\pgfpathlineto{\pgfqpoint{2.249285in}{0.587339in}}%
\pgfpathlineto{\pgfqpoint{2.249285in}{0.587339in}}%
\pgfusepath{stroke}%
\end{pgfscope}%
\begin{pgfscope}%
\pgfpathrectangle{\pgfqpoint{0.487500in}{0.320625in}}{\pgfqpoint{1.845679in}{1.164375in}}%
\pgfusepath{clip}%
\pgfsetbuttcap%
\pgfsetroundjoin%
\pgfsetlinewidth{0.752812pt}%
\definecolor{currentstroke}{rgb}{0.000000,0.529412,0.423529}%
\pgfsetstrokecolor{currentstroke}%
\pgfsetstrokeopacity{0.650000}%
\pgfsetdash{{2.775000pt}{1.200000pt}}{0.000000pt}%
\pgfpathmoveto{\pgfqpoint{0.571395in}{1.468525in}}%
\pgfpathlineto{\pgfqpoint{1.216932in}{0.315625in}}%
\pgfpathlineto{\pgfqpoint{1.216932in}{0.315625in}}%
\pgfusepath{stroke}%
\end{pgfscope}%
\begin{pgfscope}%
\pgfpathrectangle{\pgfqpoint{0.487500in}{0.320625in}}{\pgfqpoint{1.845679in}{1.164375in}}%
\pgfusepath{clip}%
\pgfsetbuttcap%
\pgfsetroundjoin%
\pgfsetlinewidth{0.803000pt}%
\definecolor{currentstroke}{rgb}{0.000000,0.529412,0.423529}%
\pgfsetstrokecolor{currentstroke}%
\pgfsetstrokeopacity{0.800000}%
\pgfsetdash{{0.800000pt}{1.320000pt}}{0.000000pt}%
\pgfpathmoveto{\pgfqpoint{0.571395in}{1.468525in}}%
\pgfpathlineto{\pgfqpoint{2.249285in}{1.468525in}}%
\pgfpathlineto{\pgfqpoint{2.249285in}{1.468525in}}%
\pgfusepath{stroke}%
\end{pgfscope}%
\begin{pgfscope}%
\pgfpathrectangle{\pgfqpoint{0.487500in}{0.320625in}}{\pgfqpoint{1.845679in}{1.164375in}}%
\pgfusepath{clip}%
\pgfsetrectcap%
\pgfsetroundjoin%
\pgfsetlinewidth{1.254687pt}%
\definecolor{currentstroke}{rgb}{0.478431,0.243137,0.615686}%
\pgfsetstrokecolor{currentstroke}%
\pgfsetdash{}{0pt}%
\pgfpathmoveto{\pgfqpoint{0.571395in}{1.468525in}}%
\pgfpathlineto{\pgfqpoint{0.598660in}{1.219935in}}%
\pgfpathlineto{\pgfqpoint{0.655289in}{0.703632in}}%
\pgfpathlineto{\pgfqpoint{0.657386in}{0.687880in}}%
\pgfpathlineto{\pgfqpoint{0.661581in}{0.635236in}}%
\pgfpathlineto{\pgfqpoint{0.663678in}{0.628303in}}%
\pgfpathlineto{\pgfqpoint{0.665776in}{0.635128in}}%
\pgfpathlineto{\pgfqpoint{0.667873in}{0.613144in}}%
\pgfpathlineto{\pgfqpoint{0.669971in}{0.545875in}}%
\pgfpathlineto{\pgfqpoint{0.672068in}{0.601897in}}%
\pgfpathlineto{\pgfqpoint{0.674165in}{0.583457in}}%
\pgfpathlineto{\pgfqpoint{0.676263in}{0.586511in}}%
\pgfpathlineto{\pgfqpoint{0.678360in}{0.624315in}}%
\pgfpathlineto{\pgfqpoint{0.680457in}{0.605347in}}%
\pgfpathlineto{\pgfqpoint{0.684652in}{0.666735in}}%
\pgfpathlineto{\pgfqpoint{0.686749in}{0.674127in}}%
\pgfpathlineto{\pgfqpoint{0.688847in}{0.663123in}}%
\pgfpathlineto{\pgfqpoint{0.690944in}{0.589738in}}%
\pgfpathlineto{\pgfqpoint{0.693042in}{0.660727in}}%
\pgfpathlineto{\pgfqpoint{0.695139in}{0.683119in}}%
\pgfpathlineto{\pgfqpoint{0.697236in}{0.676148in}}%
\pgfpathlineto{\pgfqpoint{0.699334in}{0.625452in}}%
\pgfpathlineto{\pgfqpoint{0.703528in}{0.633982in}}%
\pgfpathlineto{\pgfqpoint{0.705626in}{0.597412in}}%
\pgfpathlineto{\pgfqpoint{0.707723in}{0.662049in}}%
\pgfpathlineto{\pgfqpoint{0.709820in}{0.662056in}}%
\pgfpathlineto{\pgfqpoint{0.711918in}{0.612722in}}%
\pgfpathlineto{\pgfqpoint{0.714015in}{0.615012in}}%
\pgfpathlineto{\pgfqpoint{0.716113in}{0.620572in}}%
\pgfpathlineto{\pgfqpoint{0.718210in}{0.577875in}}%
\pgfpathlineto{\pgfqpoint{0.720307in}{0.647107in}}%
\pgfpathlineto{\pgfqpoint{0.722405in}{0.648095in}}%
\pgfpathlineto{\pgfqpoint{0.724502in}{0.532743in}}%
\pgfpathlineto{\pgfqpoint{0.726599in}{0.644753in}}%
\pgfpathlineto{\pgfqpoint{0.728697in}{0.653226in}}%
\pgfpathlineto{\pgfqpoint{0.730794in}{0.623917in}}%
\pgfpathlineto{\pgfqpoint{0.732891in}{0.632237in}}%
\pgfpathlineto{\pgfqpoint{0.734989in}{0.666421in}}%
\pgfpathlineto{\pgfqpoint{0.737086in}{0.655621in}}%
\pgfpathlineto{\pgfqpoint{0.739184in}{0.624539in}}%
\pgfpathlineto{\pgfqpoint{0.741281in}{0.674391in}}%
\pgfpathlineto{\pgfqpoint{0.743378in}{0.678589in}}%
\pgfpathlineto{\pgfqpoint{0.745476in}{0.648761in}}%
\pgfpathlineto{\pgfqpoint{0.747573in}{0.654151in}}%
\pgfpathlineto{\pgfqpoint{0.749670in}{0.689790in}}%
\pgfpathlineto{\pgfqpoint{0.751768in}{0.684393in}}%
\pgfpathlineto{\pgfqpoint{0.755962in}{0.606141in}}%
\pgfpathlineto{\pgfqpoint{0.758060in}{0.603575in}}%
\pgfpathlineto{\pgfqpoint{0.762254in}{0.664687in}}%
\pgfpathlineto{\pgfqpoint{0.764352in}{0.648469in}}%
\pgfpathlineto{\pgfqpoint{0.766449in}{0.582492in}}%
\pgfpathlineto{\pgfqpoint{0.768547in}{0.592695in}}%
\pgfpathlineto{\pgfqpoint{0.770644in}{0.519305in}}%
\pgfpathlineto{\pgfqpoint{0.772741in}{0.623417in}}%
\pgfpathlineto{\pgfqpoint{0.774839in}{0.638351in}}%
\pgfpathlineto{\pgfqpoint{0.776936in}{0.593974in}}%
\pgfpathlineto{\pgfqpoint{0.781131in}{0.629974in}}%
\pgfpathlineto{\pgfqpoint{0.783228in}{0.607548in}}%
\pgfpathlineto{\pgfqpoint{0.785325in}{0.604279in}}%
\pgfpathlineto{\pgfqpoint{0.787423in}{0.643810in}}%
\pgfpathlineto{\pgfqpoint{0.789520in}{0.634086in}}%
\pgfpathlineto{\pgfqpoint{0.791618in}{0.613281in}}%
\pgfpathlineto{\pgfqpoint{0.793715in}{0.659854in}}%
\pgfpathlineto{\pgfqpoint{0.795812in}{0.665337in}}%
\pgfpathlineto{\pgfqpoint{0.797910in}{0.642078in}}%
\pgfpathlineto{\pgfqpoint{0.800007in}{0.626189in}}%
\pgfpathlineto{\pgfqpoint{0.802104in}{0.672261in}}%
\pgfpathlineto{\pgfqpoint{0.804202in}{0.674007in}}%
\pgfpathlineto{\pgfqpoint{0.808396in}{0.595425in}}%
\pgfpathlineto{\pgfqpoint{0.810494in}{0.615360in}}%
\pgfpathlineto{\pgfqpoint{0.812591in}{0.590259in}}%
\pgfpathlineto{\pgfqpoint{0.814689in}{0.650247in}}%
\pgfpathlineto{\pgfqpoint{0.816786in}{0.648153in}}%
\pgfpathlineto{\pgfqpoint{0.818883in}{0.600571in}}%
\pgfpathlineto{\pgfqpoint{0.820981in}{0.588506in}}%
\pgfpathlineto{\pgfqpoint{0.823078in}{0.586934in}}%
\pgfpathlineto{\pgfqpoint{0.825175in}{0.592744in}}%
\pgfpathlineto{\pgfqpoint{0.827273in}{0.632724in}}%
\pgfpathlineto{\pgfqpoint{0.829370in}{0.619027in}}%
\pgfpathlineto{\pgfqpoint{0.831467in}{0.572616in}}%
\pgfpathlineto{\pgfqpoint{0.833565in}{0.624975in}}%
\pgfpathlineto{\pgfqpoint{0.835662in}{0.622212in}}%
\pgfpathlineto{\pgfqpoint{0.837760in}{0.510430in}}%
\pgfpathlineto{\pgfqpoint{0.839857in}{0.629774in}}%
\pgfpathlineto{\pgfqpoint{0.841954in}{0.642979in}}%
\pgfpathlineto{\pgfqpoint{0.844052in}{0.585247in}}%
\pgfpathlineto{\pgfqpoint{0.846149in}{0.640768in}}%
\pgfpathlineto{\pgfqpoint{0.848246in}{0.659921in}}%
\pgfpathlineto{\pgfqpoint{0.850344in}{0.651719in}}%
\pgfpathlineto{\pgfqpoint{0.852441in}{0.550710in}}%
\pgfpathlineto{\pgfqpoint{0.854538in}{0.655813in}}%
\pgfpathlineto{\pgfqpoint{0.856636in}{0.672234in}}%
\pgfpathlineto{\pgfqpoint{0.858733in}{0.650993in}}%
\pgfpathlineto{\pgfqpoint{0.860831in}{0.530303in}}%
\pgfpathlineto{\pgfqpoint{0.862928in}{0.616590in}}%
\pgfpathlineto{\pgfqpoint{0.865025in}{0.567172in}}%
\pgfpathlineto{\pgfqpoint{0.867123in}{0.632846in}}%
\pgfpathlineto{\pgfqpoint{0.869220in}{0.651187in}}%
\pgfpathlineto{\pgfqpoint{0.871317in}{0.622951in}}%
\pgfpathlineto{\pgfqpoint{0.873415in}{0.546568in}}%
\pgfpathlineto{\pgfqpoint{0.875512in}{0.596643in}}%
\pgfpathlineto{\pgfqpoint{0.877609in}{0.466427in}}%
\pgfpathlineto{\pgfqpoint{0.879707in}{0.619976in}}%
\pgfpathlineto{\pgfqpoint{0.881804in}{0.629257in}}%
\pgfpathlineto{\pgfqpoint{0.883902in}{0.569167in}}%
\pgfpathlineto{\pgfqpoint{0.885999in}{0.611208in}}%
\pgfpathlineto{\pgfqpoint{0.888096in}{0.624356in}}%
\pgfpathlineto{\pgfqpoint{0.890194in}{0.596172in}}%
\pgfpathlineto{\pgfqpoint{0.892291in}{0.605477in}}%
\pgfpathlineto{\pgfqpoint{0.894388in}{0.639049in}}%
\pgfpathlineto{\pgfqpoint{0.898583in}{0.610882in}}%
\pgfpathlineto{\pgfqpoint{0.900680in}{0.651769in}}%
\pgfpathlineto{\pgfqpoint{0.902778in}{0.654318in}}%
\pgfpathlineto{\pgfqpoint{0.904875in}{0.622249in}}%
\pgfpathlineto{\pgfqpoint{0.906973in}{0.632006in}}%
\pgfpathlineto{\pgfqpoint{0.909070in}{0.665874in}}%
\pgfpathlineto{\pgfqpoint{0.911167in}{0.659583in}}%
\pgfpathlineto{\pgfqpoint{0.913265in}{0.607846in}}%
\pgfpathlineto{\pgfqpoint{0.915362in}{0.608240in}}%
\pgfpathlineto{\pgfqpoint{0.917459in}{0.606161in}}%
\pgfpathlineto{\pgfqpoint{0.919557in}{0.602772in}}%
\pgfpathlineto{\pgfqpoint{0.921654in}{0.646633in}}%
\pgfpathlineto{\pgfqpoint{0.923751in}{0.637058in}}%
\pgfpathlineto{\pgfqpoint{0.925849in}{0.570476in}}%
\pgfpathlineto{\pgfqpoint{0.927946in}{0.594481in}}%
\pgfpathlineto{\pgfqpoint{0.930043in}{0.579993in}}%
\pgfpathlineto{\pgfqpoint{0.932141in}{0.596605in}}%
\pgfpathlineto{\pgfqpoint{0.934238in}{0.629996in}}%
\pgfpathlineto{\pgfqpoint{0.936336in}{0.610607in}}%
\pgfpathlineto{\pgfqpoint{0.938433in}{0.581135in}}%
\pgfpathlineto{\pgfqpoint{0.940530in}{0.622060in}}%
\pgfpathlineto{\pgfqpoint{0.942628in}{0.615453in}}%
\pgfpathlineto{\pgfqpoint{0.944725in}{0.529187in}}%
\pgfpathlineto{\pgfqpoint{0.946822in}{0.629317in}}%
\pgfpathlineto{\pgfqpoint{0.948920in}{0.637569in}}%
\pgfpathlineto{\pgfqpoint{0.951017in}{0.549414in}}%
\pgfpathlineto{\pgfqpoint{0.953114in}{0.638723in}}%
\pgfpathlineto{\pgfqpoint{0.955212in}{0.653755in}}%
\pgfpathlineto{\pgfqpoint{0.957309in}{0.640807in}}%
\pgfpathlineto{\pgfqpoint{0.959407in}{0.571518in}}%
\pgfpathlineto{\pgfqpoint{0.961504in}{0.655297in}}%
\pgfpathlineto{\pgfqpoint{0.963601in}{0.665002in}}%
\pgfpathlineto{\pgfqpoint{0.965699in}{0.634997in}}%
\pgfpathlineto{\pgfqpoint{0.967796in}{0.577192in}}%
\pgfpathlineto{\pgfqpoint{0.969893in}{0.616077in}}%
\pgfpathlineto{\pgfqpoint{0.971991in}{0.509204in}}%
\pgfpathlineto{\pgfqpoint{0.974088in}{0.634827in}}%
\pgfpathlineto{\pgfqpoint{0.976185in}{0.646238in}}%
\pgfpathlineto{\pgfqpoint{0.980380in}{0.574256in}}%
\pgfpathlineto{\pgfqpoint{0.982478in}{0.597489in}}%
\pgfpathlineto{\pgfqpoint{0.984575in}{0.524021in}}%
\pgfpathlineto{\pgfqpoint{0.986672in}{0.621426in}}%
\pgfpathlineto{\pgfqpoint{0.988770in}{0.626748in}}%
\pgfpathlineto{\pgfqpoint{0.990867in}{0.552636in}}%
\pgfpathlineto{\pgfqpoint{0.992964in}{0.612361in}}%
\pgfpathlineto{\pgfqpoint{0.995062in}{0.622440in}}%
\pgfpathlineto{\pgfqpoint{0.997159in}{0.588003in}}%
\pgfpathlineto{\pgfqpoint{0.999256in}{0.609851in}}%
\pgfpathlineto{\pgfqpoint{1.001354in}{0.638501in}}%
\pgfpathlineto{\pgfqpoint{1.003451in}{0.618254in}}%
\pgfpathlineto{\pgfqpoint{1.005549in}{0.613806in}}%
\pgfpathlineto{\pgfqpoint{1.007646in}{0.649054in}}%
\pgfpathlineto{\pgfqpoint{1.009743in}{0.648625in}}%
\pgfpathlineto{\pgfqpoint{1.011841in}{0.605026in}}%
\pgfpathlineto{\pgfqpoint{1.016035in}{0.663627in}}%
\pgfpathlineto{\pgfqpoint{1.018133in}{0.650575in}}%
\pgfpathlineto{\pgfqpoint{1.020230in}{0.573463in}}%
\pgfpathlineto{\pgfqpoint{1.022327in}{0.614131in}}%
\pgfpathlineto{\pgfqpoint{1.024425in}{0.599072in}}%
\pgfpathlineto{\pgfqpoint{1.026522in}{0.611916in}}%
\pgfpathlineto{\pgfqpoint{1.028620in}{0.646407in}}%
\pgfpathlineto{\pgfqpoint{1.030717in}{0.630927in}}%
\pgfpathlineto{\pgfqpoint{1.032814in}{0.526031in}}%
\pgfpathlineto{\pgfqpoint{1.034912in}{0.599419in}}%
\pgfpathlineto{\pgfqpoint{1.037009in}{0.575714in}}%
\pgfpathlineto{\pgfqpoint{1.041204in}{0.630770in}}%
\pgfpathlineto{\pgfqpoint{1.045398in}{0.587794in}}%
\pgfpathlineto{\pgfqpoint{1.047496in}{0.622311in}}%
\pgfpathlineto{\pgfqpoint{1.049593in}{0.612480in}}%
\pgfpathlineto{\pgfqpoint{1.051691in}{0.559576in}}%
\pgfpathlineto{\pgfqpoint{1.053788in}{0.631481in}}%
\pgfpathlineto{\pgfqpoint{1.055885in}{0.635558in}}%
\pgfpathlineto{\pgfqpoint{1.057983in}{0.463573in}}%
\pgfpathlineto{\pgfqpoint{1.060080in}{0.638825in}}%
\pgfpathlineto{\pgfqpoint{1.062177in}{0.650708in}}%
\pgfpathlineto{\pgfqpoint{1.064275in}{0.632934in}}%
\pgfpathlineto{\pgfqpoint{1.066372in}{0.600295in}}%
\pgfpathlineto{\pgfqpoint{1.068469in}{0.656499in}}%
\pgfpathlineto{\pgfqpoint{1.070567in}{0.660034in}}%
\pgfpathlineto{\pgfqpoint{1.072664in}{0.621006in}}%
\pgfpathlineto{\pgfqpoint{1.074762in}{0.598280in}}%
\pgfpathlineto{\pgfqpoint{1.076859in}{0.615170in}}%
\pgfpathlineto{\pgfqpoint{1.078956in}{0.543813in}}%
\pgfpathlineto{\pgfqpoint{1.081054in}{0.638137in}}%
\pgfpathlineto{\pgfqpoint{1.083151in}{0.643013in}}%
\pgfpathlineto{\pgfqpoint{1.085248in}{0.599662in}}%
\pgfpathlineto{\pgfqpoint{1.087346in}{0.587510in}}%
\pgfpathlineto{\pgfqpoint{1.089443in}{0.598169in}}%
\pgfpathlineto{\pgfqpoint{1.091540in}{0.551573in}}%
\pgfpathlineto{\pgfqpoint{1.093638in}{0.624656in}}%
\pgfpathlineto{\pgfqpoint{1.095735in}{0.625669in}}%
\pgfpathlineto{\pgfqpoint{1.097832in}{0.524856in}}%
\pgfpathlineto{\pgfqpoint{1.099930in}{0.614905in}}%
\pgfpathlineto{\pgfqpoint{1.102027in}{0.621754in}}%
\pgfpathlineto{\pgfqpoint{1.104125in}{0.578129in}}%
\pgfpathlineto{\pgfqpoint{1.106222in}{0.615563in}}%
\pgfpathlineto{\pgfqpoint{1.108319in}{0.639304in}}%
\pgfpathlineto{\pgfqpoint{1.110417in}{0.612093in}}%
\pgfpathlineto{\pgfqpoint{1.112514in}{0.618327in}}%
\pgfpathlineto{\pgfqpoint{1.114611in}{0.648048in}}%
\pgfpathlineto{\pgfqpoint{1.116709in}{0.644201in}}%
\pgfpathlineto{\pgfqpoint{1.118806in}{0.579370in}}%
\pgfpathlineto{\pgfqpoint{1.120903in}{0.642777in}}%
\pgfpathlineto{\pgfqpoint{1.123001in}{0.662380in}}%
\pgfpathlineto{\pgfqpoint{1.125098in}{0.642237in}}%
\pgfpathlineto{\pgfqpoint{1.127196in}{0.506817in}}%
\pgfpathlineto{\pgfqpoint{1.129293in}{0.617758in}}%
\pgfpathlineto{\pgfqpoint{1.131390in}{0.588983in}}%
\pgfpathlineto{\pgfqpoint{1.135585in}{0.646888in}}%
\pgfpathlineto{\pgfqpoint{1.137682in}{0.625101in}}%
\pgfpathlineto{\pgfqpoint{1.139780in}{0.529315in}}%
\pgfpathlineto{\pgfqpoint{1.141877in}{0.603360in}}%
\pgfpathlineto{\pgfqpoint{1.143974in}{0.568184in}}%
\pgfpathlineto{\pgfqpoint{1.146072in}{0.608767in}}%
\pgfpathlineto{\pgfqpoint{1.148169in}{0.632228in}}%
\pgfpathlineto{\pgfqpoint{1.150267in}{0.602176in}}%
\pgfpathlineto{\pgfqpoint{1.152364in}{0.595178in}}%
\pgfpathlineto{\pgfqpoint{1.154461in}{0.623308in}}%
\pgfpathlineto{\pgfqpoint{1.156559in}{0.609270in}}%
\pgfpathlineto{\pgfqpoint{1.158656in}{0.579341in}}%
\pgfpathlineto{\pgfqpoint{1.160753in}{0.634508in}}%
\pgfpathlineto{\pgfqpoint{1.162851in}{0.633446in}}%
\pgfpathlineto{\pgfqpoint{1.164948in}{0.551270in}}%
\pgfpathlineto{\pgfqpoint{1.167045in}{0.640038in}}%
\pgfpathlineto{\pgfqpoint{1.169143in}{0.648378in}}%
\pgfpathlineto{\pgfqpoint{1.171240in}{0.624052in}}%
\pgfpathlineto{\pgfqpoint{1.173338in}{0.616675in}}%
\pgfpathlineto{\pgfqpoint{1.175435in}{0.658081in}}%
\pgfpathlineto{\pgfqpoint{1.177532in}{0.654700in}}%
\pgfpathlineto{\pgfqpoint{1.179630in}{0.602739in}}%
\pgfpathlineto{\pgfqpoint{1.181727in}{0.609507in}}%
\pgfpathlineto{\pgfqpoint{1.183824in}{0.612765in}}%
\pgfpathlineto{\pgfqpoint{1.185922in}{0.580084in}}%
\pgfpathlineto{\pgfqpoint{1.188019in}{0.641641in}}%
\pgfpathlineto{\pgfqpoint{1.190116in}{0.639649in}}%
\pgfpathlineto{\pgfqpoint{1.192214in}{0.584173in}}%
\pgfpathlineto{\pgfqpoint{1.194311in}{0.596877in}}%
\pgfpathlineto{\pgfqpoint{1.196409in}{0.597517in}}%
\pgfpathlineto{\pgfqpoint{1.198506in}{0.572004in}}%
\pgfpathlineto{\pgfqpoint{1.200603in}{0.628374in}}%
\pgfpathlineto{\pgfqpoint{1.202701in}{0.624091in}}%
\pgfpathlineto{\pgfqpoint{1.204798in}{0.493040in}}%
\pgfpathlineto{\pgfqpoint{1.206895in}{0.617962in}}%
\pgfpathlineto{\pgfqpoint{1.208993in}{0.620917in}}%
\pgfpathlineto{\pgfqpoint{1.211090in}{0.560142in}}%
\pgfpathlineto{\pgfqpoint{1.213187in}{0.621702in}}%
\pgfpathlineto{\pgfqpoint{1.215285in}{0.640093in}}%
\pgfpathlineto{\pgfqpoint{1.217382in}{0.603107in}}%
\pgfpathlineto{\pgfqpoint{1.221577in}{0.647505in}}%
\pgfpathlineto{\pgfqpoint{1.223674in}{0.639404in}}%
\pgfpathlineto{\pgfqpoint{1.225772in}{0.418414in}}%
\pgfpathlineto{\pgfqpoint{1.227869in}{0.647710in}}%
\pgfpathlineto{\pgfqpoint{1.229966in}{0.660733in}}%
\pgfpathlineto{\pgfqpoint{1.232064in}{0.624098in}}%
\pgfpathlineto{\pgfqpoint{1.234161in}{0.631763in}}%
\pgfpathlineto{\pgfqpoint{1.236258in}{0.655077in}}%
\pgfpathlineto{\pgfqpoint{1.238356in}{0.639548in}}%
\pgfpathlineto{\pgfqpoint{1.240453in}{0.612249in}}%
\pgfpathlineto{\pgfqpoint{1.242551in}{0.665494in}}%
\pgfpathlineto{\pgfqpoint{1.244648in}{0.664539in}}%
\pgfpathlineto{\pgfqpoint{1.246745in}{0.614374in}}%
\pgfpathlineto{\pgfqpoint{1.250940in}{0.624120in}}%
\pgfpathlineto{\pgfqpoint{1.253037in}{0.581396in}}%
\pgfpathlineto{\pgfqpoint{1.255135in}{0.650370in}}%
\pgfpathlineto{\pgfqpoint{1.257232in}{0.650616in}}%
\pgfpathlineto{\pgfqpoint{1.259329in}{0.605894in}}%
\pgfpathlineto{\pgfqpoint{1.261427in}{0.587623in}}%
\pgfpathlineto{\pgfqpoint{1.263524in}{0.590819in}}%
\pgfpathlineto{\pgfqpoint{1.265621in}{0.589614in}}%
\pgfpathlineto{\pgfqpoint{1.267719in}{0.633428in}}%
\pgfpathlineto{\pgfqpoint{1.269816in}{0.622871in}}%
\pgfpathlineto{\pgfqpoint{1.271914in}{0.474768in}}%
\pgfpathlineto{\pgfqpoint{1.274011in}{0.608269in}}%
\pgfpathlineto{\pgfqpoint{1.276108in}{0.602454in}}%
\pgfpathlineto{\pgfqpoint{1.278206in}{0.561483in}}%
\pgfpathlineto{\pgfqpoint{1.280303in}{0.626091in}}%
\pgfpathlineto{\pgfqpoint{1.282400in}{0.625908in}}%
\pgfpathlineto{\pgfqpoint{1.284498in}{0.529754in}}%
\pgfpathlineto{\pgfqpoint{1.286595in}{0.629414in}}%
\pgfpathlineto{\pgfqpoint{1.288692in}{0.637304in}}%
\pgfpathlineto{\pgfqpoint{1.290790in}{0.609970in}}%
\pgfpathlineto{\pgfqpoint{1.292887in}{0.612038in}}%
\pgfpathlineto{\pgfqpoint{1.294985in}{0.648857in}}%
\pgfpathlineto{\pgfqpoint{1.297082in}{0.641036in}}%
\pgfpathlineto{\pgfqpoint{1.299179in}{0.589583in}}%
\pgfpathlineto{\pgfqpoint{1.301277in}{0.651765in}}%
\pgfpathlineto{\pgfqpoint{1.303374in}{0.656337in}}%
\pgfpathlineto{\pgfqpoint{1.305471in}{0.622666in}}%
\pgfpathlineto{\pgfqpoint{1.307569in}{0.638414in}}%
\pgfpathlineto{\pgfqpoint{1.309666in}{0.669632in}}%
\pgfpathlineto{\pgfqpoint{1.311763in}{0.658657in}}%
\pgfpathlineto{\pgfqpoint{1.313861in}{0.494093in}}%
\pgfpathlineto{\pgfqpoint{1.315958in}{0.642631in}}%
\pgfpathlineto{\pgfqpoint{1.318056in}{0.638696in}}%
\pgfpathlineto{\pgfqpoint{1.320153in}{0.587193in}}%
\pgfpathlineto{\pgfqpoint{1.322250in}{0.659911in}}%
\pgfpathlineto{\pgfqpoint{1.324348in}{0.662083in}}%
\pgfpathlineto{\pgfqpoint{1.328542in}{0.587578in}}%
\pgfpathlineto{\pgfqpoint{1.330640in}{0.597410in}}%
\pgfpathlineto{\pgfqpoint{1.332737in}{0.596846in}}%
\pgfpathlineto{\pgfqpoint{1.334834in}{0.641447in}}%
\pgfpathlineto{\pgfqpoint{1.336932in}{0.631963in}}%
\pgfpathlineto{\pgfqpoint{1.339029in}{0.557637in}}%
\pgfpathlineto{\pgfqpoint{1.341127in}{0.595607in}}%
\pgfpathlineto{\pgfqpoint{1.343224in}{0.583468in}}%
\pgfpathlineto{\pgfqpoint{1.345321in}{0.589660in}}%
\pgfpathlineto{\pgfqpoint{1.347419in}{0.627191in}}%
\pgfpathlineto{\pgfqpoint{1.349516in}{0.611349in}}%
\pgfpathlineto{\pgfqpoint{1.351613in}{0.576042in}}%
\pgfpathlineto{\pgfqpoint{1.353711in}{0.621878in}}%
\pgfpathlineto{\pgfqpoint{1.355808in}{0.618264in}}%
\pgfpathlineto{\pgfqpoint{1.357905in}{0.500081in}}%
\pgfpathlineto{\pgfqpoint{1.360003in}{0.625868in}}%
\pgfpathlineto{\pgfqpoint{1.362100in}{0.638887in}}%
\pgfpathlineto{\pgfqpoint{1.364198in}{0.584236in}}%
\pgfpathlineto{\pgfqpoint{1.366295in}{0.633585in}}%
\pgfpathlineto{\pgfqpoint{1.368392in}{0.652691in}}%
\pgfpathlineto{\pgfqpoint{1.370490in}{0.643089in}}%
\pgfpathlineto{\pgfqpoint{1.372587in}{0.494074in}}%
\pgfpathlineto{\pgfqpoint{1.374684in}{0.651714in}}%
\pgfpathlineto{\pgfqpoint{1.376782in}{0.664951in}}%
\pgfpathlineto{\pgfqpoint{1.378879in}{0.639016in}}%
\pgfpathlineto{\pgfqpoint{1.380976in}{0.551780in}}%
\pgfpathlineto{\pgfqpoint{1.383074in}{0.613882in}}%
\pgfpathlineto{\pgfqpoint{1.385171in}{0.543774in}}%
\pgfpathlineto{\pgfqpoint{1.387269in}{0.630900in}}%
\pgfpathlineto{\pgfqpoint{1.389366in}{0.645437in}}%
\pgfpathlineto{\pgfqpoint{1.391463in}{0.613104in}}%
\pgfpathlineto{\pgfqpoint{1.393561in}{0.563158in}}%
\pgfpathlineto{\pgfqpoint{1.395658in}{0.594678in}}%
\pgfpathlineto{\pgfqpoint{1.397755in}{0.503127in}}%
\pgfpathlineto{\pgfqpoint{1.399853in}{0.618004in}}%
\pgfpathlineto{\pgfqpoint{1.401950in}{0.624945in}}%
\pgfpathlineto{\pgfqpoint{1.404047in}{0.556737in}}%
\pgfpathlineto{\pgfqpoint{1.406145in}{0.609185in}}%
\pgfpathlineto{\pgfqpoint{1.408242in}{0.620348in}}%
\pgfpathlineto{\pgfqpoint{1.410340in}{0.588093in}}%
\pgfpathlineto{\pgfqpoint{1.412437in}{0.605866in}}%
\pgfpathlineto{\pgfqpoint{1.414534in}{0.635985in}}%
\pgfpathlineto{\pgfqpoint{1.416632in}{0.617261in}}%
\pgfpathlineto{\pgfqpoint{1.418729in}{0.611407in}}%
\pgfpathlineto{\pgfqpoint{1.420826in}{0.647857in}}%
\pgfpathlineto{\pgfqpoint{1.422924in}{0.648338in}}%
\pgfpathlineto{\pgfqpoint{1.425021in}{0.609098in}}%
\pgfpathlineto{\pgfqpoint{1.429216in}{0.662204in}}%
\pgfpathlineto{\pgfqpoint{1.431313in}{0.651430in}}%
\pgfpathlineto{\pgfqpoint{1.433410in}{0.584205in}}%
\pgfpathlineto{\pgfqpoint{1.435508in}{0.611438in}}%
\pgfpathlineto{\pgfqpoint{1.437605in}{0.600791in}}%
\pgfpathlineto{\pgfqpoint{1.439703in}{0.606754in}}%
\pgfpathlineto{\pgfqpoint{1.441800in}{0.644702in}}%
\pgfpathlineto{\pgfqpoint{1.443897in}{0.631477in}}%
\pgfpathlineto{\pgfqpoint{1.445995in}{0.543989in}}%
\pgfpathlineto{\pgfqpoint{1.448092in}{0.597099in}}%
\pgfpathlineto{\pgfqpoint{1.450189in}{0.577484in}}%
\pgfpathlineto{\pgfqpoint{1.452287in}{0.598104in}}%
\pgfpathlineto{\pgfqpoint{1.454384in}{0.629147in}}%
\pgfpathlineto{\pgfqpoint{1.458579in}{0.583670in}}%
\pgfpathlineto{\pgfqpoint{1.460676in}{0.621131in}}%
\pgfpathlineto{\pgfqpoint{1.462774in}{0.612916in}}%
\pgfpathlineto{\pgfqpoint{1.464871in}{0.546450in}}%
\pgfpathlineto{\pgfqpoint{1.466968in}{0.629351in}}%
\pgfpathlineto{\pgfqpoint{1.469066in}{0.635476in}}%
\pgfpathlineto{\pgfqpoint{1.471163in}{0.523972in}}%
\pgfpathlineto{\pgfqpoint{1.473260in}{0.637304in}}%
\pgfpathlineto{\pgfqpoint{1.475358in}{0.650621in}}%
\pgfpathlineto{\pgfqpoint{1.477455in}{0.634923in}}%
\pgfpathlineto{\pgfqpoint{1.479552in}{0.590549in}}%
\pgfpathlineto{\pgfqpoint{1.481650in}{0.654878in}}%
\pgfpathlineto{\pgfqpoint{1.483747in}{0.660820in}}%
\pgfpathlineto{\pgfqpoint{1.487942in}{0.592054in}}%
\pgfpathlineto{\pgfqpoint{1.490039in}{0.615033in}}%
\pgfpathlineto{\pgfqpoint{1.492137in}{0.511701in}}%
\pgfpathlineto{\pgfqpoint{1.494234in}{0.635979in}}%
\pgfpathlineto{\pgfqpoint{1.496331in}{0.643387in}}%
\pgfpathlineto{\pgfqpoint{1.498429in}{0.603116in}}%
\pgfpathlineto{\pgfqpoint{1.500526in}{0.582942in}}%
\pgfpathlineto{\pgfqpoint{1.502623in}{0.597662in}}%
\pgfpathlineto{\pgfqpoint{1.504721in}{0.540742in}}%
\pgfpathlineto{\pgfqpoint{1.506818in}{0.622669in}}%
\pgfpathlineto{\pgfqpoint{1.508916in}{0.625529in}}%
\pgfpathlineto{\pgfqpoint{1.511013in}{0.538370in}}%
\pgfpathlineto{\pgfqpoint{1.513110in}{0.613254in}}%
\pgfpathlineto{\pgfqpoint{1.515208in}{0.621428in}}%
\pgfpathlineto{\pgfqpoint{1.517305in}{0.581855in}}%
\pgfpathlineto{\pgfqpoint{1.521500in}{0.638387in}}%
\pgfpathlineto{\pgfqpoint{1.523597in}{0.613973in}}%
\pgfpathlineto{\pgfqpoint{1.525694in}{0.615938in}}%
\pgfpathlineto{\pgfqpoint{1.527792in}{0.647691in}}%
\pgfpathlineto{\pgfqpoint{1.529889in}{0.645111in}}%
\pgfpathlineto{\pgfqpoint{1.531987in}{0.589573in}}%
\pgfpathlineto{\pgfqpoint{1.534084in}{0.640414in}}%
\pgfpathlineto{\pgfqpoint{1.536181in}{0.662210in}}%
\pgfpathlineto{\pgfqpoint{1.538279in}{0.644489in}}%
\pgfpathlineto{\pgfqpoint{1.540376in}{0.510233in}}%
\pgfpathlineto{\pgfqpoint{1.542473in}{0.616470in}}%
\pgfpathlineto{\pgfqpoint{1.544571in}{0.592691in}}%
\pgfpathlineto{\pgfqpoint{1.548765in}{0.646298in}}%
\pgfpathlineto{\pgfqpoint{1.550863in}{0.626692in}}%
\pgfpathlineto{\pgfqpoint{1.552960in}{0.494367in}}%
\pgfpathlineto{\pgfqpoint{1.555058in}{0.601997in}}%
\pgfpathlineto{\pgfqpoint{1.557155in}{0.571251in}}%
\pgfpathlineto{\pgfqpoint{1.559252in}{0.606099in}}%
\pgfpathlineto{\pgfqpoint{1.561350in}{0.631437in}}%
\pgfpathlineto{\pgfqpoint{1.563447in}{0.603650in}}%
\pgfpathlineto{\pgfqpoint{1.565544in}{0.592394in}}%
\pgfpathlineto{\pgfqpoint{1.567642in}{0.622689in}}%
\pgfpathlineto{\pgfqpoint{1.569739in}{0.610195in}}%
\pgfpathlineto{\pgfqpoint{1.571836in}{0.573138in}}%
\pgfpathlineto{\pgfqpoint{1.573934in}{0.633204in}}%
\pgfpathlineto{\pgfqpoint{1.576031in}{0.633942in}}%
\pgfpathlineto{\pgfqpoint{1.578129in}{0.534303in}}%
\pgfpathlineto{\pgfqpoint{1.580226in}{0.639290in}}%
\pgfpathlineto{\pgfqpoint{1.582323in}{0.648806in}}%
\pgfpathlineto{\pgfqpoint{1.584421in}{0.626777in}}%
\pgfpathlineto{\pgfqpoint{1.586518in}{0.611972in}}%
\pgfpathlineto{\pgfqpoint{1.588615in}{0.657327in}}%
\pgfpathlineto{\pgfqpoint{1.590713in}{0.656106in}}%
\pgfpathlineto{\pgfqpoint{1.592810in}{0.608749in}}%
\pgfpathlineto{\pgfqpoint{1.594907in}{0.606540in}}%
\pgfpathlineto{\pgfqpoint{1.597005in}{0.613539in}}%
\pgfpathlineto{\pgfqpoint{1.599102in}{0.571508in}}%
\pgfpathlineto{\pgfqpoint{1.601199in}{0.640391in}}%
\pgfpathlineto{\pgfqpoint{1.603297in}{0.640506in}}%
\pgfpathlineto{\pgfqpoint{1.605394in}{0.589297in}}%
\pgfpathlineto{\pgfqpoint{1.607492in}{0.594268in}}%
\pgfpathlineto{\pgfqpoint{1.609589in}{0.597762in}}%
\pgfpathlineto{\pgfqpoint{1.611686in}{0.566252in}}%
\pgfpathlineto{\pgfqpoint{1.613784in}{0.627094in}}%
\pgfpathlineto{\pgfqpoint{1.615881in}{0.624475in}}%
\pgfpathlineto{\pgfqpoint{1.617978in}{0.395652in}}%
\pgfpathlineto{\pgfqpoint{1.620076in}{0.616923in}}%
\pgfpathlineto{\pgfqpoint{1.622173in}{0.621068in}}%
\pgfpathlineto{\pgfqpoint{1.624270in}{0.566610in}}%
\pgfpathlineto{\pgfqpoint{1.626368in}{0.619792in}}%
\pgfpathlineto{\pgfqpoint{1.628465in}{0.639758in}}%
\pgfpathlineto{\pgfqpoint{1.630563in}{0.605886in}}%
\pgfpathlineto{\pgfqpoint{1.632660in}{0.621932in}}%
\pgfpathlineto{\pgfqpoint{1.634757in}{0.647542in}}%
\pgfpathlineto{\pgfqpoint{1.636855in}{0.640732in}}%
\pgfpathlineto{\pgfqpoint{1.638952in}{0.531120in}}%
\pgfpathlineto{\pgfqpoint{1.641049in}{0.646281in}}%
\pgfpathlineto{\pgfqpoint{1.643147in}{0.661172in}}%
\pgfpathlineto{\pgfqpoint{1.645244in}{0.628382in}}%
\pgfpathlineto{\pgfqpoint{1.647341in}{0.629345in}}%
\pgfpathlineto{\pgfqpoint{1.649439in}{0.655438in}}%
\pgfpathlineto{\pgfqpoint{1.651536in}{0.642224in}}%
\pgfpathlineto{\pgfqpoint{1.653634in}{0.604362in}}%
\pgfpathlineto{\pgfqpoint{1.655731in}{0.664777in}}%
\pgfpathlineto{\pgfqpoint{1.657828in}{0.666323in}}%
\pgfpathlineto{\pgfqpoint{1.659926in}{0.620764in}}%
\pgfpathlineto{\pgfqpoint{1.662023in}{0.615880in}}%
\pgfpathlineto{\pgfqpoint{1.664120in}{0.625221in}}%
\pgfpathlineto{\pgfqpoint{1.666218in}{0.569752in}}%
\pgfpathlineto{\pgfqpoint{1.668315in}{0.649290in}}%
\pgfpathlineto{\pgfqpoint{1.670412in}{0.651846in}}%
\pgfpathlineto{\pgfqpoint{1.672510in}{0.610129in}}%
\pgfpathlineto{\pgfqpoint{1.674607in}{0.583266in}}%
\pgfpathlineto{\pgfqpoint{1.676705in}{0.591222in}}%
\pgfpathlineto{\pgfqpoint{1.678802in}{0.585878in}}%
\pgfpathlineto{\pgfqpoint{1.680899in}{0.632517in}}%
\pgfpathlineto{\pgfqpoint{1.682997in}{0.623604in}}%
\pgfpathlineto{\pgfqpoint{1.685094in}{0.461540in}}%
\pgfpathlineto{\pgfqpoint{1.687191in}{0.607140in}}%
\pgfpathlineto{\pgfqpoint{1.689289in}{0.602820in}}%
\pgfpathlineto{\pgfqpoint{1.691386in}{0.555124in}}%
\pgfpathlineto{\pgfqpoint{1.693483in}{0.624926in}}%
\pgfpathlineto{\pgfqpoint{1.695581in}{0.626163in}}%
\pgfpathlineto{\pgfqpoint{1.697678in}{0.512507in}}%
\pgfpathlineto{\pgfqpoint{1.699776in}{0.628777in}}%
\pgfpathlineto{\pgfqpoint{1.701873in}{0.637575in}}%
\pgfpathlineto{\pgfqpoint{1.703970in}{0.612236in}}%
\pgfpathlineto{\pgfqpoint{1.706068in}{0.608962in}}%
\pgfpathlineto{\pgfqpoint{1.708165in}{0.648296in}}%
\pgfpathlineto{\pgfqpoint{1.710262in}{0.642414in}}%
\pgfpathlineto{\pgfqpoint{1.712360in}{0.582061in}}%
\pgfpathlineto{\pgfqpoint{1.714457in}{0.651508in}}%
\pgfpathlineto{\pgfqpoint{1.716554in}{0.657349in}}%
\pgfpathlineto{\pgfqpoint{1.718652in}{0.627070in}}%
\pgfpathlineto{\pgfqpoint{1.720749in}{0.635432in}}%
\pgfpathlineto{\pgfqpoint{1.722847in}{0.669557in}}%
\pgfpathlineto{\pgfqpoint{1.724944in}{0.660968in}}%
\pgfpathlineto{\pgfqpoint{1.727041in}{0.556648in}}%
\pgfpathlineto{\pgfqpoint{1.729139in}{0.641895in}}%
\pgfpathlineto{\pgfqpoint{1.731236in}{0.640550in}}%
\pgfpathlineto{\pgfqpoint{1.733333in}{0.572315in}}%
\pgfpathlineto{\pgfqpoint{1.735431in}{0.658981in}}%
\pgfpathlineto{\pgfqpoint{1.737528in}{0.663640in}}%
\pgfpathlineto{\pgfqpoint{1.741723in}{0.581286in}}%
\pgfpathlineto{\pgfqpoint{1.743820in}{0.598024in}}%
\pgfpathlineto{\pgfqpoint{1.745918in}{0.592871in}}%
\pgfpathlineto{\pgfqpoint{1.748015in}{0.640719in}}%
\pgfpathlineto{\pgfqpoint{1.750112in}{0.633014in}}%
\pgfpathlineto{\pgfqpoint{1.752210in}{0.565593in}}%
\pgfpathlineto{\pgfqpoint{1.754307in}{0.593797in}}%
\pgfpathlineto{\pgfqpoint{1.756404in}{0.584166in}}%
\pgfpathlineto{\pgfqpoint{1.758502in}{0.586897in}}%
\pgfpathlineto{\pgfqpoint{1.760599in}{0.626370in}}%
\pgfpathlineto{\pgfqpoint{1.762696in}{0.611996in}}%
\pgfpathlineto{\pgfqpoint{1.764794in}{0.572680in}}%
\pgfpathlineto{\pgfqpoint{1.766891in}{0.621312in}}%
\pgfpathlineto{\pgfqpoint{1.768988in}{0.618633in}}%
\pgfpathlineto{\pgfqpoint{1.771086in}{0.523241in}}%
\pgfpathlineto{\pgfqpoint{1.773183in}{0.624536in}}%
\pgfpathlineto{\pgfqpoint{1.775281in}{0.638824in}}%
\pgfpathlineto{\pgfqpoint{1.777378in}{0.588556in}}%
\pgfpathlineto{\pgfqpoint{1.779475in}{0.632740in}}%
\pgfpathlineto{\pgfqpoint{1.781573in}{0.652977in}}%
\pgfpathlineto{\pgfqpoint{1.783670in}{0.644539in}}%
\pgfpathlineto{\pgfqpoint{1.785767in}{0.518390in}}%
\pgfpathlineto{\pgfqpoint{1.787865in}{0.650557in}}%
\pgfpathlineto{\pgfqpoint{1.789962in}{0.665495in}}%
\pgfpathlineto{\pgfqpoint{1.792059in}{0.641666in}}%
\pgfpathlineto{\pgfqpoint{1.794157in}{0.525138in}}%
\pgfpathlineto{\pgfqpoint{1.796254in}{0.613348in}}%
\pgfpathlineto{\pgfqpoint{1.798352in}{0.555281in}}%
\pgfpathlineto{\pgfqpoint{1.800449in}{0.629264in}}%
\pgfpathlineto{\pgfqpoint{1.802546in}{0.645723in}}%
\pgfpathlineto{\pgfqpoint{1.804644in}{0.615270in}}%
\pgfpathlineto{\pgfqpoint{1.806741in}{0.556055in}}%
\pgfpathlineto{\pgfqpoint{1.808838in}{0.594101in}}%
\pgfpathlineto{\pgfqpoint{1.810936in}{0.475645in}}%
\pgfpathlineto{\pgfqpoint{1.813033in}{0.616678in}}%
\pgfpathlineto{\pgfqpoint{1.815130in}{0.625016in}}%
\pgfpathlineto{\pgfqpoint{1.817228in}{0.561411in}}%
\pgfpathlineto{\pgfqpoint{1.819325in}{0.608070in}}%
\pgfpathlineto{\pgfqpoint{1.821423in}{0.620289in}}%
\pgfpathlineto{\pgfqpoint{1.823520in}{0.590186in}}%
\pgfpathlineto{\pgfqpoint{1.825617in}{0.603640in}}%
\pgfpathlineto{\pgfqpoint{1.827715in}{0.635443in}}%
\pgfpathlineto{\pgfqpoint{1.829812in}{0.618661in}}%
\pgfpathlineto{\pgfqpoint{1.831909in}{0.609404in}}%
\pgfpathlineto{\pgfqpoint{1.834007in}{0.647787in}}%
\pgfpathlineto{\pgfqpoint{1.836104in}{0.649187in}}%
\pgfpathlineto{\pgfqpoint{1.838201in}{0.613282in}}%
\pgfpathlineto{\pgfqpoint{1.840299in}{0.631346in}}%
\pgfpathlineto{\pgfqpoint{1.842396in}{0.662074in}}%
\pgfpathlineto{\pgfqpoint{1.844494in}{0.653168in}}%
\pgfpathlineto{\pgfqpoint{1.846591in}{0.592613in}}%
\pgfpathlineto{\pgfqpoint{1.848688in}{0.609658in}}%
\pgfpathlineto{\pgfqpoint{1.850786in}{0.602523in}}%
\pgfpathlineto{\pgfqpoint{1.852883in}{0.603348in}}%
\pgfpathlineto{\pgfqpoint{1.854980in}{0.644177in}}%
\pgfpathlineto{\pgfqpoint{1.857078in}{0.632725in}}%
\pgfpathlineto{\pgfqpoint{1.859175in}{0.555305in}}%
\pgfpathlineto{\pgfqpoint{1.861272in}{0.595604in}}%
\pgfpathlineto{\pgfqpoint{1.863370in}{0.578908in}}%
\pgfpathlineto{\pgfqpoint{1.865467in}{0.595722in}}%
\pgfpathlineto{\pgfqpoint{1.867565in}{0.628528in}}%
\pgfpathlineto{\pgfqpoint{1.869662in}{0.608542in}}%
\pgfpathlineto{\pgfqpoint{1.871759in}{0.580784in}}%
\pgfpathlineto{\pgfqpoint{1.873857in}{0.620671in}}%
\pgfpathlineto{\pgfqpoint{1.875954in}{0.613574in}}%
\pgfpathlineto{\pgfqpoint{1.878051in}{0.534656in}}%
\pgfpathlineto{\pgfqpoint{1.880149in}{0.628253in}}%
\pgfpathlineto{\pgfqpoint{1.882246in}{0.635777in}}%
\pgfpathlineto{\pgfqpoint{1.884343in}{0.542514in}}%
\pgfpathlineto{\pgfqpoint{1.886441in}{0.636645in}}%
\pgfpathlineto{\pgfqpoint{1.888538in}{0.651003in}}%
\pgfpathlineto{\pgfqpoint{1.890636in}{0.636732in}}%
\pgfpathlineto{\pgfqpoint{1.892733in}{0.581889in}}%
\pgfpathlineto{\pgfqpoint{1.894830in}{0.654066in}}%
\pgfpathlineto{\pgfqpoint{1.896928in}{0.661755in}}%
\pgfpathlineto{\pgfqpoint{1.901122in}{0.586640in}}%
\pgfpathlineto{\pgfqpoint{1.903220in}{0.615191in}}%
\pgfpathlineto{\pgfqpoint{1.905317in}{0.366262in}}%
\pgfpathlineto{\pgfqpoint{1.907414in}{0.634674in}}%
\pgfpathlineto{\pgfqpoint{1.909512in}{0.644074in}}%
\pgfpathlineto{\pgfqpoint{1.911609in}{0.606034in}}%
\pgfpathlineto{\pgfqpoint{1.913707in}{0.579281in}}%
\pgfpathlineto{\pgfqpoint{1.915804in}{0.597586in}}%
\pgfpathlineto{\pgfqpoint{1.917901in}{0.530341in}}%
\pgfpathlineto{\pgfqpoint{1.919999in}{0.621474in}}%
\pgfpathlineto{\pgfqpoint{1.922096in}{0.625843in}}%
\pgfpathlineto{\pgfqpoint{1.924193in}{0.547237in}}%
\pgfpathlineto{\pgfqpoint{1.926291in}{0.612227in}}%
\pgfpathlineto{\pgfqpoint{1.928388in}{0.621524in}}%
\pgfpathlineto{\pgfqpoint{1.930485in}{0.584925in}}%
\pgfpathlineto{\pgfqpoint{1.934680in}{0.637999in}}%
\pgfpathlineto{\pgfqpoint{1.936777in}{0.615867in}}%
\pgfpathlineto{\pgfqpoint{1.938875in}{0.614029in}}%
\pgfpathlineto{\pgfqpoint{1.940972in}{0.647646in}}%
\pgfpathlineto{\pgfqpoint{1.943070in}{0.646100in}}%
\pgfpathlineto{\pgfqpoint{1.945167in}{0.596476in}}%
\pgfpathlineto{\pgfqpoint{1.947264in}{0.638633in}}%
\pgfpathlineto{\pgfqpoint{1.949362in}{0.662297in}}%
\pgfpathlineto{\pgfqpoint{1.951459in}{0.646533in}}%
\pgfpathlineto{\pgfqpoint{1.953556in}{0.546001in}}%
\pgfpathlineto{\pgfqpoint{1.955654in}{0.615506in}}%
\pgfpathlineto{\pgfqpoint{1.957751in}{0.595585in}}%
\pgfpathlineto{\pgfqpoint{1.959848in}{0.614560in}}%
\pgfpathlineto{\pgfqpoint{1.961946in}{0.646021in}}%
\pgfpathlineto{\pgfqpoint{1.964043in}{0.628254in}}%
\pgfpathlineto{\pgfqpoint{1.966141in}{0.451045in}}%
\pgfpathlineto{\pgfqpoint{1.968238in}{0.601018in}}%
\pgfpathlineto{\pgfqpoint{1.970335in}{0.573826in}}%
\pgfpathlineto{\pgfqpoint{1.974530in}{0.631008in}}%
\pgfpathlineto{\pgfqpoint{1.976627in}{0.605207in}}%
\pgfpathlineto{\pgfqpoint{1.978725in}{0.589923in}}%
\pgfpathlineto{\pgfqpoint{1.980822in}{0.622343in}}%
\pgfpathlineto{\pgfqpoint{1.982919in}{0.611181in}}%
\pgfpathlineto{\pgfqpoint{1.985017in}{0.567051in}}%
\pgfpathlineto{\pgfqpoint{1.987114in}{0.632253in}}%
\pgfpathlineto{\pgfqpoint{1.989212in}{0.634575in}}%
\pgfpathlineto{\pgfqpoint{1.991309in}{0.506055in}}%
\pgfpathlineto{\pgfqpoint{1.993406in}{0.638698in}}%
\pgfpathlineto{\pgfqpoint{1.995504in}{0.649280in}}%
\pgfpathlineto{\pgfqpoint{1.999698in}{0.607449in}}%
\pgfpathlineto{\pgfqpoint{2.001796in}{0.656750in}}%
\pgfpathlineto{\pgfqpoint{2.003893in}{0.657409in}}%
\pgfpathlineto{\pgfqpoint{2.005990in}{0.613508in}}%
\pgfpathlineto{\pgfqpoint{2.008088in}{0.603758in}}%
\pgfpathlineto{\pgfqpoint{2.010185in}{0.614271in}}%
\pgfpathlineto{\pgfqpoint{2.012283in}{0.562122in}}%
\pgfpathlineto{\pgfqpoint{2.014380in}{0.639363in}}%
\pgfpathlineto{\pgfqpoint{2.016477in}{0.641414in}}%
\pgfpathlineto{\pgfqpoint{2.018575in}{0.593569in}}%
\pgfpathlineto{\pgfqpoint{2.020672in}{0.591903in}}%
\pgfpathlineto{\pgfqpoint{2.022769in}{0.598162in}}%
\pgfpathlineto{\pgfqpoint{2.024867in}{0.560242in}}%
\pgfpathlineto{\pgfqpoint{2.026964in}{0.626055in}}%
\pgfpathlineto{\pgfqpoint{2.029061in}{0.625060in}}%
\pgfpathlineto{\pgfqpoint{2.031159in}{0.496008in}}%
\pgfpathlineto{\pgfqpoint{2.033256in}{0.616021in}}%
\pgfpathlineto{\pgfqpoint{2.035354in}{0.621349in}}%
\pgfpathlineto{\pgfqpoint{2.037451in}{0.571990in}}%
\pgfpathlineto{\pgfqpoint{2.039548in}{0.618036in}}%
\pgfpathlineto{\pgfqpoint{2.041646in}{0.639568in}}%
\pgfpathlineto{\pgfqpoint{2.043743in}{0.608617in}}%
\pgfpathlineto{\pgfqpoint{2.045840in}{0.620291in}}%
\pgfpathlineto{\pgfqpoint{2.047938in}{0.647570in}}%
\pgfpathlineto{\pgfqpoint{2.050035in}{0.641948in}}%
\pgfpathlineto{\pgfqpoint{2.052132in}{0.556729in}}%
\pgfpathlineto{\pgfqpoint{2.054230in}{0.644905in}}%
\pgfpathlineto{\pgfqpoint{2.056327in}{0.661544in}}%
\pgfpathlineto{\pgfqpoint{2.058425in}{0.637782in}}%
\pgfpathlineto{\pgfqpoint{2.060522in}{0.557395in}}%
\pgfpathlineto{\pgfqpoint{2.062619in}{0.618778in}}%
\pgfpathlineto{\pgfqpoint{2.064717in}{0.581960in}}%
\pgfpathlineto{\pgfqpoint{2.066814in}{0.623972in}}%
\pgfpathlineto{\pgfqpoint{2.068911in}{0.646887in}}%
\pgfpathlineto{\pgfqpoint{2.071009in}{0.621869in}}%
\pgfpathlineto{\pgfqpoint{2.073106in}{0.553093in}}%
\pgfpathlineto{\pgfqpoint{2.075203in}{0.604802in}}%
\pgfpathlineto{\pgfqpoint{2.077301in}{0.562521in}}%
\pgfpathlineto{\pgfqpoint{2.079398in}{0.611893in}}%
\pgfpathlineto{\pgfqpoint{2.081496in}{0.632785in}}%
\pgfpathlineto{\pgfqpoint{2.083593in}{0.599256in}}%
\pgfpathlineto{\pgfqpoint{2.085690in}{0.598586in}}%
\pgfpathlineto{\pgfqpoint{2.087788in}{0.623740in}}%
\pgfpathlineto{\pgfqpoint{2.091982in}{0.586908in}}%
\pgfpathlineto{\pgfqpoint{2.094080in}{0.635932in}}%
\pgfpathlineto{\pgfqpoint{2.096177in}{0.632037in}}%
\pgfpathlineto{\pgfqpoint{2.098274in}{0.567758in}}%
\pgfpathlineto{\pgfqpoint{2.100372in}{0.640686in}}%
\pgfpathlineto{\pgfqpoint{2.102469in}{0.647163in}}%
\pgfpathlineto{\pgfqpoint{2.104566in}{0.618644in}}%
\pgfpathlineto{\pgfqpoint{2.106664in}{0.622725in}}%
\pgfpathlineto{\pgfqpoint{2.108761in}{0.658686in}}%
\pgfpathlineto{\pgfqpoint{2.110859in}{0.651177in}}%
\pgfpathlineto{\pgfqpoint{2.112956in}{0.558179in}}%
\pgfpathlineto{\pgfqpoint{2.115053in}{0.649410in}}%
\pgfpathlineto{\pgfqpoint{2.117151in}{0.652414in}}%
\pgfpathlineto{\pgfqpoint{2.119248in}{0.599186in}}%
\pgfpathlineto{\pgfqpoint{2.121345in}{0.649893in}}%
\pgfpathlineto{\pgfqpoint{2.123443in}{0.670554in}}%
\pgfpathlineto{\pgfqpoint{2.125540in}{0.648013in}}%
\pgfpathlineto{\pgfqpoint{2.127637in}{0.563414in}}%
\pgfpathlineto{\pgfqpoint{2.129735in}{0.629103in}}%
\pgfpathlineto{\pgfqpoint{2.131832in}{0.596628in}}%
\pgfpathlineto{\pgfqpoint{2.133930in}{0.631744in}}%
\pgfpathlineto{\pgfqpoint{2.136027in}{0.656712in}}%
\pgfpathlineto{\pgfqpoint{2.138124in}{0.634756in}}%
\pgfpathlineto{\pgfqpoint{2.140222in}{0.472020in}}%
\pgfpathlineto{\pgfqpoint{2.142319in}{0.599789in}}%
\pgfpathlineto{\pgfqpoint{2.144416in}{0.529219in}}%
\pgfpathlineto{\pgfqpoint{2.146514in}{0.620345in}}%
\pgfpathlineto{\pgfqpoint{2.148611in}{0.634912in}}%
\pgfpathlineto{\pgfqpoint{2.150708in}{0.596051in}}%
\pgfpathlineto{\pgfqpoint{2.152806in}{0.591195in}}%
\pgfpathlineto{\pgfqpoint{2.154903in}{0.611051in}}%
\pgfpathlineto{\pgfqpoint{2.157001in}{0.571442in}}%
\pgfpathlineto{\pgfqpoint{2.159098in}{0.608483in}}%
\pgfpathlineto{\pgfqpoint{2.161195in}{0.632119in}}%
\pgfpathlineto{\pgfqpoint{2.163293in}{0.601261in}}%
\pgfpathlineto{\pgfqpoint{2.165390in}{0.609911in}}%
\pgfpathlineto{\pgfqpoint{2.167487in}{0.637087in}}%
\pgfpathlineto{\pgfqpoint{2.169585in}{0.630023in}}%
\pgfpathlineto{\pgfqpoint{2.171682in}{0.478323in}}%
\pgfpathlineto{\pgfqpoint{2.173779in}{0.637604in}}%
\pgfpathlineto{\pgfqpoint{2.175877in}{0.650963in}}%
\pgfpathlineto{\pgfqpoint{2.177974in}{0.608272in}}%
\pgfpathlineto{\pgfqpoint{2.182169in}{0.657605in}}%
\pgfpathlineto{\pgfqpoint{2.184266in}{0.647169in}}%
\pgfpathlineto{\pgfqpoint{2.186364in}{0.564818in}}%
\pgfpathlineto{\pgfqpoint{2.188461in}{0.660222in}}%
\pgfpathlineto{\pgfqpoint{2.190558in}{0.669952in}}%
\pgfpathlineto{\pgfqpoint{2.192656in}{0.633120in}}%
\pgfpathlineto{\pgfqpoint{2.194753in}{0.624278in}}%
\pgfpathlineto{\pgfqpoint{2.196850in}{0.646258in}}%
\pgfpathlineto{\pgfqpoint{2.198948in}{0.610867in}}%
\pgfpathlineto{\pgfqpoint{2.203143in}{0.666897in}}%
\pgfpathlineto{\pgfqpoint{2.205240in}{0.647353in}}%
\pgfpathlineto{\pgfqpoint{2.207337in}{0.553574in}}%
\pgfpathlineto{\pgfqpoint{2.209435in}{0.604830in}}%
\pgfpathlineto{\pgfqpoint{2.211532in}{0.536160in}}%
\pgfpathlineto{\pgfqpoint{2.213629in}{0.628266in}}%
\pgfpathlineto{\pgfqpoint{2.215727in}{0.643206in}}%
\pgfpathlineto{\pgfqpoint{2.219921in}{0.570383in}}%
\pgfpathlineto{\pgfqpoint{2.222019in}{0.598150in}}%
\pgfpathlineto{\pgfqpoint{2.224116in}{0.460477in}}%
\pgfpathlineto{\pgfqpoint{2.226214in}{0.616367in}}%
\pgfpathlineto{\pgfqpoint{2.228311in}{0.626591in}}%
\pgfpathlineto{\pgfqpoint{2.230408in}{0.566721in}}%
\pgfpathlineto{\pgfqpoint{2.232506in}{0.610298in}}%
\pgfpathlineto{\pgfqpoint{2.234603in}{0.624217in}}%
\pgfpathlineto{\pgfqpoint{2.236700in}{0.599259in}}%
\pgfpathlineto{\pgfqpoint{2.238798in}{0.600039in}}%
\pgfpathlineto{\pgfqpoint{2.240895in}{0.637404in}}%
\pgfpathlineto{\pgfqpoint{2.242992in}{0.630163in}}%
\pgfpathlineto{\pgfqpoint{2.247187in}{0.590819in}}%
\pgfpathlineto{\pgfqpoint{2.249285in}{0.574712in}}%
\pgfpathlineto{\pgfqpoint{2.249285in}{0.574712in}}%
\pgfusepath{stroke}%
\end{pgfscope}%
\begin{pgfscope}%
\pgfpathrectangle{\pgfqpoint{0.487500in}{0.320625in}}{\pgfqpoint{1.845679in}{1.164375in}}%
\pgfusepath{clip}%
\pgfsetbuttcap%
\pgfsetroundjoin%
\pgfsetlinewidth{0.752812pt}%
\definecolor{currentstroke}{rgb}{0.478431,0.243137,0.615686}%
\pgfsetstrokecolor{currentstroke}%
\pgfsetstrokeopacity{0.650000}%
\pgfsetdash{{2.775000pt}{1.200000pt}}{0.000000pt}%
\pgfpathmoveto{\pgfqpoint{0.571395in}{1.468525in}}%
\pgfpathlineto{\pgfqpoint{0.700502in}{0.315625in}}%
\pgfpathlineto{\pgfqpoint{0.700502in}{0.315625in}}%
\pgfusepath{stroke}%
\end{pgfscope}%
\begin{pgfscope}%
\pgfpathrectangle{\pgfqpoint{0.487500in}{0.320625in}}{\pgfqpoint{1.845679in}{1.164375in}}%
\pgfusepath{clip}%
\pgfsetbuttcap%
\pgfsetroundjoin%
\pgfsetlinewidth{0.803000pt}%
\definecolor{currentstroke}{rgb}{0.478431,0.243137,0.615686}%
\pgfsetstrokecolor{currentstroke}%
\pgfsetstrokeopacity{0.800000}%
\pgfsetdash{{0.800000pt}{1.320000pt}}{0.000000pt}%
\pgfpathmoveto{\pgfqpoint{0.571395in}{1.468525in}}%
\pgfpathlineto{\pgfqpoint{2.249285in}{1.468525in}}%
\pgfpathlineto{\pgfqpoint{2.249285in}{1.468525in}}%
\pgfusepath{stroke}%
\end{pgfscope}%
\begin{pgfscope}%
\pgfsetrectcap%
\pgfsetmiterjoin%
\pgfsetlinewidth{0.803000pt}%
\definecolor{currentstroke}{rgb}{0.000000,0.000000,0.000000}%
\pgfsetstrokecolor{currentstroke}%
\pgfsetdash{}{0pt}%
\pgfpathmoveto{\pgfqpoint{0.487500in}{0.320625in}}%
\pgfpathlineto{\pgfqpoint{0.487500in}{1.485000in}}%
\pgfusepath{stroke}%
\end{pgfscope}%
\begin{pgfscope}%
\pgfsetrectcap%
\pgfsetmiterjoin%
\pgfsetlinewidth{0.803000pt}%
\definecolor{currentstroke}{rgb}{0.000000,0.000000,0.000000}%
\pgfsetstrokecolor{currentstroke}%
\pgfsetdash{}{0pt}%
\pgfpathmoveto{\pgfqpoint{2.333179in}{0.320625in}}%
\pgfpathlineto{\pgfqpoint{2.333179in}{1.485000in}}%
\pgfusepath{stroke}%
\end{pgfscope}%
\begin{pgfscope}%
\pgfsetrectcap%
\pgfsetmiterjoin%
\pgfsetlinewidth{0.803000pt}%
\definecolor{currentstroke}{rgb}{0.000000,0.000000,0.000000}%
\pgfsetstrokecolor{currentstroke}%
\pgfsetdash{}{0pt}%
\pgfpathmoveto{\pgfqpoint{0.487500in}{0.320625in}}%
\pgfpathlineto{\pgfqpoint{2.333179in}{0.320625in}}%
\pgfusepath{stroke}%
\end{pgfscope}%
\begin{pgfscope}%
\pgfsetrectcap%
\pgfsetmiterjoin%
\pgfsetlinewidth{0.803000pt}%
\definecolor{currentstroke}{rgb}{0.000000,0.000000,0.000000}%
\pgfsetstrokecolor{currentstroke}%
\pgfsetdash{}{0pt}%
\pgfpathmoveto{\pgfqpoint{0.487500in}{1.485000in}}%
\pgfpathlineto{\pgfqpoint{2.333179in}{1.485000in}}%
\pgfusepath{stroke}%
\end{pgfscope}%
\begin{pgfscope}%
\definecolor{textcolor}{rgb}{0.000000,0.000000,0.000000}%
\pgfsetstrokecolor{textcolor}%
\pgfsetfillcolor{textcolor}%
\pgftext[x=1.410340in,y=1.568333in,,base]{\color{textcolor}{\rmfamily\fontsize{8.000000}{9.600000}\selectfont\catcode`\^=\active\def^{\ifmmode\sp\else\^{}\fi}\catcode`\%=\active\def%{\%}(a) Initialization scale}}%
\end{pgfscope}%
\begin{pgfscope}%
\pgfsetbuttcap%
\pgfsetmiterjoin%
\definecolor{currentfill}{rgb}{1.000000,1.000000,1.000000}%
\pgfsetfillcolor{currentfill}%
\pgfsetfillopacity{0.800000}%
\pgfsetlinewidth{1.003750pt}%
\definecolor{currentstroke}{rgb}{0.800000,0.800000,0.800000}%
\pgfsetstrokecolor{currentstroke}%
\pgfsetstrokeopacity{0.800000}%
\pgfsetdash{}{0pt}%
\pgfpathmoveto{\pgfqpoint{1.409008in}{1.284674in}}%
\pgfpathlineto{\pgfqpoint{2.288457in}{1.284674in}}%
\pgfpathquadraticcurveto{\pgfqpoint{2.301235in}{1.284674in}}{\pgfqpoint{2.301235in}{1.297451in}}%
\pgfpathlineto{\pgfqpoint{2.301235in}{1.440278in}}%
\pgfpathquadraticcurveto{\pgfqpoint{2.301235in}{1.453056in}}{\pgfqpoint{2.288457in}{1.453056in}}%
\pgfpathlineto{\pgfqpoint{1.409008in}{1.453056in}}%
\pgfpathquadraticcurveto{\pgfqpoint{1.396230in}{1.453056in}}{\pgfqpoint{1.396230in}{1.440278in}}%
\pgfpathlineto{\pgfqpoint{1.396230in}{1.297451in}}%
\pgfpathquadraticcurveto{\pgfqpoint{1.396230in}{1.284674in}}{\pgfqpoint{1.409008in}{1.284674in}}%
\pgfpathlineto{\pgfqpoint{1.409008in}{1.284674in}}%
\pgfpathclose%
\pgfusepath{stroke,fill}%
\end{pgfscope}%
\begin{pgfscope}%
\pgfsetrectcap%
\pgfsetroundjoin%
\pgfsetlinewidth{1.254687pt}%
\definecolor{currentstroke}{rgb}{0.141176,0.345098,0.650980}%
\pgfsetstrokecolor{currentstroke}%
\pgfsetdash{}{0pt}%
\pgfpathmoveto{\pgfqpoint{1.412202in}{1.410904in}}%
\pgfpathlineto{\pgfqpoint{1.456924in}{1.410904in}}%
\pgfpathlineto{\pgfqpoint{1.501647in}{1.410904in}}%
\pgfusepath{stroke}%
\end{pgfscope}%
\begin{pgfscope}%
\definecolor{textcolor}{rgb}{0.000000,0.000000,0.000000}%
\pgfsetstrokecolor{textcolor}%
\pgfsetfillcolor{textcolor}%
\pgftext[x=1.552758in,y=1.388543in,left,base]{\color{textcolor}{\rmfamily\fontsize{4.600000}{5.520000}\selectfont\catcode`\^=\active\def^{\ifmmode\sp\else\^{}\fi}\catcode`\%=\active\def%{\%}$\alpha=0.1$}}%
\end{pgfscope}%
\begin{pgfscope}%
\pgfsetrectcap%
\pgfsetroundjoin%
\pgfsetlinewidth{1.254687pt}%
\definecolor{currentstroke}{rgb}{0.819608,0.286275,0.356863}%
\pgfsetstrokecolor{currentstroke}%
\pgfsetdash{}{0pt}%
\pgfpathmoveto{\pgfqpoint{1.412202in}{1.336296in}}%
\pgfpathlineto{\pgfqpoint{1.456924in}{1.336296in}}%
\pgfpathlineto{\pgfqpoint{1.501647in}{1.336296in}}%
\pgfusepath{stroke}%
\end{pgfscope}%
\begin{pgfscope}%
\definecolor{textcolor}{rgb}{0.000000,0.000000,0.000000}%
\pgfsetstrokecolor{textcolor}%
\pgfsetfillcolor{textcolor}%
\pgftext[x=1.552758in,y=1.313935in,left,base]{\color{textcolor}{\rmfamily\fontsize{4.600000}{5.520000}\selectfont\catcode`\^=\active\def^{\ifmmode\sp\else\^{}\fi}\catcode`\%=\active\def%{\%}$\alpha=0.5$}}%
\end{pgfscope}%
\begin{pgfscope}%
\pgfsetrectcap%
\pgfsetroundjoin%
\pgfsetlinewidth{1.254687pt}%
\definecolor{currentstroke}{rgb}{0.000000,0.529412,0.423529}%
\pgfsetstrokecolor{currentstroke}%
\pgfsetdash{}{0pt}%
\pgfpathmoveto{\pgfqpoint{1.905847in}{1.410904in}}%
\pgfpathlineto{\pgfqpoint{1.950569in}{1.410904in}}%
\pgfpathlineto{\pgfqpoint{1.995292in}{1.410904in}}%
\pgfusepath{stroke}%
\end{pgfscope}%
\begin{pgfscope}%
\definecolor{textcolor}{rgb}{0.000000,0.000000,0.000000}%
\pgfsetstrokecolor{textcolor}%
\pgfsetfillcolor{textcolor}%
\pgftext[x=2.046403in,y=1.388543in,left,base]{\color{textcolor}{\rmfamily\fontsize{4.600000}{5.520000}\selectfont\catcode`\^=\active\def^{\ifmmode\sp\else\^{}\fi}\catcode`\%=\active\def%{\%}$\alpha=1$}}%
\end{pgfscope}%
\begin{pgfscope}%
\pgfsetrectcap%
\pgfsetroundjoin%
\pgfsetlinewidth{1.254687pt}%
\definecolor{currentstroke}{rgb}{0.478431,0.243137,0.615686}%
\pgfsetstrokecolor{currentstroke}%
\pgfsetdash{}{0pt}%
\pgfpathmoveto{\pgfqpoint{1.905847in}{1.336296in}}%
\pgfpathlineto{\pgfqpoint{1.950569in}{1.336296in}}%
\pgfpathlineto{\pgfqpoint{1.995292in}{1.336296in}}%
\pgfusepath{stroke}%
\end{pgfscope}%
\begin{pgfscope}%
\definecolor{textcolor}{rgb}{0.000000,0.000000,0.000000}%
\pgfsetstrokecolor{textcolor}%
\pgfsetfillcolor{textcolor}%
\pgftext[x=2.046403in,y=1.313935in,left,base]{\color{textcolor}{\rmfamily\fontsize{4.600000}{5.520000}\selectfont\catcode`\^=\active\def^{\ifmmode\sp\else\^{}\fi}\catcode`\%=\active\def%{\%}$\alpha=5$}}%
\end{pgfscope}%
\begin{pgfscope}%
\pgfsetbuttcap%
\pgfsetmiterjoin%
\definecolor{currentfill}{rgb}{1.000000,1.000000,1.000000}%
\pgfsetfillcolor{currentfill}%
\pgfsetlinewidth{0.000000pt}%
\definecolor{currentstroke}{rgb}{0.000000,0.000000,0.000000}%
\pgfsetstrokecolor{currentstroke}%
\pgfsetstrokeopacity{0.000000}%
\pgfsetdash{}{0pt}%
\pgfpathmoveto{\pgfqpoint{2.554660in}{0.320625in}}%
\pgfpathlineto{\pgfqpoint{4.400340in}{0.320625in}}%
\pgfpathlineto{\pgfqpoint{4.400340in}{1.485000in}}%
\pgfpathlineto{\pgfqpoint{2.554660in}{1.485000in}}%
\pgfpathlineto{\pgfqpoint{2.554660in}{0.320625in}}%
\pgfpathclose%
\pgfusepath{fill}%
\end{pgfscope}%
\begin{pgfscope}%
\pgfpathrectangle{\pgfqpoint{2.554660in}{0.320625in}}{\pgfqpoint{1.845679in}{1.164375in}}%
\pgfusepath{clip}%
\pgfsetrectcap%
\pgfsetroundjoin%
\pgfsetlinewidth{0.803000pt}%
\definecolor{currentstroke}{rgb}{0.909804,0.909804,0.909804}%
\pgfsetstrokecolor{currentstroke}%
\pgfsetdash{}{0pt}%
\pgfpathmoveto{\pgfqpoint{2.638555in}{0.320625in}}%
\pgfpathlineto{\pgfqpoint{2.638555in}{1.485000in}}%
\pgfusepath{stroke}%
\end{pgfscope}%
\begin{pgfscope}%
\pgfsetbuttcap%
\pgfsetroundjoin%
\definecolor{currentfill}{rgb}{0.000000,0.000000,0.000000}%
\pgfsetfillcolor{currentfill}%
\pgfsetlinewidth{0.803000pt}%
\definecolor{currentstroke}{rgb}{0.000000,0.000000,0.000000}%
\pgfsetstrokecolor{currentstroke}%
\pgfsetdash{}{0pt}%
\pgfsys@defobject{currentmarker}{\pgfqpoint{0.000000in}{-0.048611in}}{\pgfqpoint{0.000000in}{0.000000in}}{%
\pgfpathmoveto{\pgfqpoint{0.000000in}{0.000000in}}%
\pgfpathlineto{\pgfqpoint{0.000000in}{-0.048611in}}%
\pgfusepath{stroke,fill}%
}%
\begin{pgfscope}%
\pgfsys@transformshift{2.638555in}{0.320625in}%
\pgfsys@useobject{currentmarker}{}%
\end{pgfscope}%
\end{pgfscope}%
\begin{pgfscope}%
\definecolor{textcolor}{rgb}{0.000000,0.000000,0.000000}%
\pgfsetstrokecolor{textcolor}%
\pgfsetfillcolor{textcolor}%
\pgftext[x=2.638555in,y=0.223403in,,top]{\color{textcolor}{\rmfamily\fontsize{6.000000}{7.200000}\selectfont\catcode`\^=\active\def^{\ifmmode\sp\else\^{}\fi}\catcode`\%=\active\def%{\%}0}}%
\end{pgfscope}%
\begin{pgfscope}%
\pgfpathrectangle{\pgfqpoint{2.554660in}{0.320625in}}{\pgfqpoint{1.845679in}{1.164375in}}%
\pgfusepath{clip}%
\pgfsetrectcap%
\pgfsetroundjoin%
\pgfsetlinewidth{0.803000pt}%
\definecolor{currentstroke}{rgb}{0.909804,0.909804,0.909804}%
\pgfsetstrokecolor{currentstroke}%
\pgfsetdash{}{0pt}%
\pgfpathmoveto{\pgfqpoint{3.058027in}{0.320625in}}%
\pgfpathlineto{\pgfqpoint{3.058027in}{1.485000in}}%
\pgfusepath{stroke}%
\end{pgfscope}%
\begin{pgfscope}%
\pgfsetbuttcap%
\pgfsetroundjoin%
\definecolor{currentfill}{rgb}{0.000000,0.000000,0.000000}%
\pgfsetfillcolor{currentfill}%
\pgfsetlinewidth{0.803000pt}%
\definecolor{currentstroke}{rgb}{0.000000,0.000000,0.000000}%
\pgfsetstrokecolor{currentstroke}%
\pgfsetdash{}{0pt}%
\pgfsys@defobject{currentmarker}{\pgfqpoint{0.000000in}{-0.048611in}}{\pgfqpoint{0.000000in}{0.000000in}}{%
\pgfpathmoveto{\pgfqpoint{0.000000in}{0.000000in}}%
\pgfpathlineto{\pgfqpoint{0.000000in}{-0.048611in}}%
\pgfusepath{stroke,fill}%
}%
\begin{pgfscope}%
\pgfsys@transformshift{3.058027in}{0.320625in}%
\pgfsys@useobject{currentmarker}{}%
\end{pgfscope}%
\end{pgfscope}%
\begin{pgfscope}%
\definecolor{textcolor}{rgb}{0.000000,0.000000,0.000000}%
\pgfsetstrokecolor{textcolor}%
\pgfsetfillcolor{textcolor}%
\pgftext[x=3.058027in,y=0.223403in,,top]{\color{textcolor}{\rmfamily\fontsize{6.000000}{7.200000}\selectfont\catcode`\^=\active\def^{\ifmmode\sp\else\^{}\fi}\catcode`\%=\active\def%{\%}10}}%
\end{pgfscope}%
\begin{pgfscope}%
\pgfpathrectangle{\pgfqpoint{2.554660in}{0.320625in}}{\pgfqpoint{1.845679in}{1.164375in}}%
\pgfusepath{clip}%
\pgfsetrectcap%
\pgfsetroundjoin%
\pgfsetlinewidth{0.803000pt}%
\definecolor{currentstroke}{rgb}{0.909804,0.909804,0.909804}%
\pgfsetstrokecolor{currentstroke}%
\pgfsetdash{}{0pt}%
\pgfpathmoveto{\pgfqpoint{3.477500in}{0.320625in}}%
\pgfpathlineto{\pgfqpoint{3.477500in}{1.485000in}}%
\pgfusepath{stroke}%
\end{pgfscope}%
\begin{pgfscope}%
\pgfsetbuttcap%
\pgfsetroundjoin%
\definecolor{currentfill}{rgb}{0.000000,0.000000,0.000000}%
\pgfsetfillcolor{currentfill}%
\pgfsetlinewidth{0.803000pt}%
\definecolor{currentstroke}{rgb}{0.000000,0.000000,0.000000}%
\pgfsetstrokecolor{currentstroke}%
\pgfsetdash{}{0pt}%
\pgfsys@defobject{currentmarker}{\pgfqpoint{0.000000in}{-0.048611in}}{\pgfqpoint{0.000000in}{0.000000in}}{%
\pgfpathmoveto{\pgfqpoint{0.000000in}{0.000000in}}%
\pgfpathlineto{\pgfqpoint{0.000000in}{-0.048611in}}%
\pgfusepath{stroke,fill}%
}%
\begin{pgfscope}%
\pgfsys@transformshift{3.477500in}{0.320625in}%
\pgfsys@useobject{currentmarker}{}%
\end{pgfscope}%
\end{pgfscope}%
\begin{pgfscope}%
\definecolor{textcolor}{rgb}{0.000000,0.000000,0.000000}%
\pgfsetstrokecolor{textcolor}%
\pgfsetfillcolor{textcolor}%
\pgftext[x=3.477500in,y=0.223403in,,top]{\color{textcolor}{\rmfamily\fontsize{6.000000}{7.200000}\selectfont\catcode`\^=\active\def^{\ifmmode\sp\else\^{}\fi}\catcode`\%=\active\def%{\%}20}}%
\end{pgfscope}%
\begin{pgfscope}%
\pgfpathrectangle{\pgfqpoint{2.554660in}{0.320625in}}{\pgfqpoint{1.845679in}{1.164375in}}%
\pgfusepath{clip}%
\pgfsetrectcap%
\pgfsetroundjoin%
\pgfsetlinewidth{0.803000pt}%
\definecolor{currentstroke}{rgb}{0.909804,0.909804,0.909804}%
\pgfsetstrokecolor{currentstroke}%
\pgfsetdash{}{0pt}%
\pgfpathmoveto{\pgfqpoint{3.896973in}{0.320625in}}%
\pgfpathlineto{\pgfqpoint{3.896973in}{1.485000in}}%
\pgfusepath{stroke}%
\end{pgfscope}%
\begin{pgfscope}%
\pgfsetbuttcap%
\pgfsetroundjoin%
\definecolor{currentfill}{rgb}{0.000000,0.000000,0.000000}%
\pgfsetfillcolor{currentfill}%
\pgfsetlinewidth{0.803000pt}%
\definecolor{currentstroke}{rgb}{0.000000,0.000000,0.000000}%
\pgfsetstrokecolor{currentstroke}%
\pgfsetdash{}{0pt}%
\pgfsys@defobject{currentmarker}{\pgfqpoint{0.000000in}{-0.048611in}}{\pgfqpoint{0.000000in}{0.000000in}}{%
\pgfpathmoveto{\pgfqpoint{0.000000in}{0.000000in}}%
\pgfpathlineto{\pgfqpoint{0.000000in}{-0.048611in}}%
\pgfusepath{stroke,fill}%
}%
\begin{pgfscope}%
\pgfsys@transformshift{3.896973in}{0.320625in}%
\pgfsys@useobject{currentmarker}{}%
\end{pgfscope}%
\end{pgfscope}%
\begin{pgfscope}%
\definecolor{textcolor}{rgb}{0.000000,0.000000,0.000000}%
\pgfsetstrokecolor{textcolor}%
\pgfsetfillcolor{textcolor}%
\pgftext[x=3.896973in,y=0.223403in,,top]{\color{textcolor}{\rmfamily\fontsize{6.000000}{7.200000}\selectfont\catcode`\^=\active\def^{\ifmmode\sp\else\^{}\fi}\catcode`\%=\active\def%{\%}30}}%
\end{pgfscope}%
\begin{pgfscope}%
\pgfpathrectangle{\pgfqpoint{2.554660in}{0.320625in}}{\pgfqpoint{1.845679in}{1.164375in}}%
\pgfusepath{clip}%
\pgfsetrectcap%
\pgfsetroundjoin%
\pgfsetlinewidth{0.803000pt}%
\definecolor{currentstroke}{rgb}{0.909804,0.909804,0.909804}%
\pgfsetstrokecolor{currentstroke}%
\pgfsetdash{}{0pt}%
\pgfpathmoveto{\pgfqpoint{4.316445in}{0.320625in}}%
\pgfpathlineto{\pgfqpoint{4.316445in}{1.485000in}}%
\pgfusepath{stroke}%
\end{pgfscope}%
\begin{pgfscope}%
\pgfsetbuttcap%
\pgfsetroundjoin%
\definecolor{currentfill}{rgb}{0.000000,0.000000,0.000000}%
\pgfsetfillcolor{currentfill}%
\pgfsetlinewidth{0.803000pt}%
\definecolor{currentstroke}{rgb}{0.000000,0.000000,0.000000}%
\pgfsetstrokecolor{currentstroke}%
\pgfsetdash{}{0pt}%
\pgfsys@defobject{currentmarker}{\pgfqpoint{0.000000in}{-0.048611in}}{\pgfqpoint{0.000000in}{0.000000in}}{%
\pgfpathmoveto{\pgfqpoint{0.000000in}{0.000000in}}%
\pgfpathlineto{\pgfqpoint{0.000000in}{-0.048611in}}%
\pgfusepath{stroke,fill}%
}%
\begin{pgfscope}%
\pgfsys@transformshift{4.316445in}{0.320625in}%
\pgfsys@useobject{currentmarker}{}%
\end{pgfscope}%
\end{pgfscope}%
\begin{pgfscope}%
\definecolor{textcolor}{rgb}{0.000000,0.000000,0.000000}%
\pgfsetstrokecolor{textcolor}%
\pgfsetfillcolor{textcolor}%
\pgftext[x=4.316445in,y=0.223403in,,top]{\color{textcolor}{\rmfamily\fontsize{6.000000}{7.200000}\selectfont\catcode`\^=\active\def^{\ifmmode\sp\else\^{}\fi}\catcode`\%=\active\def%{\%}40}}%
\end{pgfscope}%
\begin{pgfscope}%
\definecolor{textcolor}{rgb}{0.000000,0.000000,0.000000}%
\pgfsetstrokecolor{textcolor}%
\pgfsetfillcolor{textcolor}%
\pgftext[x=3.477500in,y=0.087200in,,top]{\color{textcolor}{\rmfamily\fontsize{6.000000}{7.200000}\selectfont\catcode`\^=\active\def^{\ifmmode\sp\else\^{}\fi}\catcode`\%=\active\def%{\%}time}}%
\end{pgfscope}%
\begin{pgfscope}%
\pgfpathrectangle{\pgfqpoint{2.554660in}{0.320625in}}{\pgfqpoint{1.845679in}{1.164375in}}%
\pgfusepath{clip}%
\pgfsetrectcap%
\pgfsetroundjoin%
\pgfsetlinewidth{0.803000pt}%
\definecolor{currentstroke}{rgb}{0.909804,0.909804,0.909804}%
\pgfsetstrokecolor{currentstroke}%
\pgfsetdash{}{0pt}%
\pgfpathmoveto{\pgfqpoint{2.554660in}{0.320625in}}%
\pgfpathlineto{\pgfqpoint{4.400340in}{0.320625in}}%
\pgfusepath{stroke}%
\end{pgfscope}%
\begin{pgfscope}%
\pgfsetbuttcap%
\pgfsetroundjoin%
\definecolor{currentfill}{rgb}{0.000000,0.000000,0.000000}%
\pgfsetfillcolor{currentfill}%
\pgfsetlinewidth{0.803000pt}%
\definecolor{currentstroke}{rgb}{0.000000,0.000000,0.000000}%
\pgfsetstrokecolor{currentstroke}%
\pgfsetdash{}{0pt}%
\pgfsys@defobject{currentmarker}{\pgfqpoint{-0.048611in}{0.000000in}}{\pgfqpoint{-0.000000in}{0.000000in}}{%
\pgfpathmoveto{\pgfqpoint{-0.000000in}{0.000000in}}%
\pgfpathlineto{\pgfqpoint{-0.048611in}{0.000000in}}%
\pgfusepath{stroke,fill}%
}%
\begin{pgfscope}%
\pgfsys@transformshift{2.554660in}{0.320625in}%
\pgfsys@useobject{currentmarker}{}%
\end{pgfscope}%
\end{pgfscope}%
\begin{pgfscope}%
\pgfpathrectangle{\pgfqpoint{2.554660in}{0.320625in}}{\pgfqpoint{1.845679in}{1.164375in}}%
\pgfusepath{clip}%
\pgfsetrectcap%
\pgfsetroundjoin%
\pgfsetlinewidth{0.803000pt}%
\definecolor{currentstroke}{rgb}{0.909804,0.909804,0.909804}%
\pgfsetstrokecolor{currentstroke}%
\pgfsetdash{}{0pt}%
\pgfpathmoveto{\pgfqpoint{2.554660in}{0.536250in}}%
\pgfpathlineto{\pgfqpoint{4.400340in}{0.536250in}}%
\pgfusepath{stroke}%
\end{pgfscope}%
\begin{pgfscope}%
\pgfsetbuttcap%
\pgfsetroundjoin%
\definecolor{currentfill}{rgb}{0.000000,0.000000,0.000000}%
\pgfsetfillcolor{currentfill}%
\pgfsetlinewidth{0.803000pt}%
\definecolor{currentstroke}{rgb}{0.000000,0.000000,0.000000}%
\pgfsetstrokecolor{currentstroke}%
\pgfsetdash{}{0pt}%
\pgfsys@defobject{currentmarker}{\pgfqpoint{-0.048611in}{0.000000in}}{\pgfqpoint{-0.000000in}{0.000000in}}{%
\pgfpathmoveto{\pgfqpoint{-0.000000in}{0.000000in}}%
\pgfpathlineto{\pgfqpoint{-0.048611in}{0.000000in}}%
\pgfusepath{stroke,fill}%
}%
\begin{pgfscope}%
\pgfsys@transformshift{2.554660in}{0.536250in}%
\pgfsys@useobject{currentmarker}{}%
\end{pgfscope}%
\end{pgfscope}%
\begin{pgfscope}%
\pgfpathrectangle{\pgfqpoint{2.554660in}{0.320625in}}{\pgfqpoint{1.845679in}{1.164375in}}%
\pgfusepath{clip}%
\pgfsetrectcap%
\pgfsetroundjoin%
\pgfsetlinewidth{0.803000pt}%
\definecolor{currentstroke}{rgb}{0.909804,0.909804,0.909804}%
\pgfsetstrokecolor{currentstroke}%
\pgfsetdash{}{0pt}%
\pgfpathmoveto{\pgfqpoint{2.554660in}{0.751875in}}%
\pgfpathlineto{\pgfqpoint{4.400340in}{0.751875in}}%
\pgfusepath{stroke}%
\end{pgfscope}%
\begin{pgfscope}%
\pgfsetbuttcap%
\pgfsetroundjoin%
\definecolor{currentfill}{rgb}{0.000000,0.000000,0.000000}%
\pgfsetfillcolor{currentfill}%
\pgfsetlinewidth{0.803000pt}%
\definecolor{currentstroke}{rgb}{0.000000,0.000000,0.000000}%
\pgfsetstrokecolor{currentstroke}%
\pgfsetdash{}{0pt}%
\pgfsys@defobject{currentmarker}{\pgfqpoint{-0.048611in}{0.000000in}}{\pgfqpoint{-0.000000in}{0.000000in}}{%
\pgfpathmoveto{\pgfqpoint{-0.000000in}{0.000000in}}%
\pgfpathlineto{\pgfqpoint{-0.048611in}{0.000000in}}%
\pgfusepath{stroke,fill}%
}%
\begin{pgfscope}%
\pgfsys@transformshift{2.554660in}{0.751875in}%
\pgfsys@useobject{currentmarker}{}%
\end{pgfscope}%
\end{pgfscope}%
\begin{pgfscope}%
\pgfpathrectangle{\pgfqpoint{2.554660in}{0.320625in}}{\pgfqpoint{1.845679in}{1.164375in}}%
\pgfusepath{clip}%
\pgfsetrectcap%
\pgfsetroundjoin%
\pgfsetlinewidth{0.803000pt}%
\definecolor{currentstroke}{rgb}{0.909804,0.909804,0.909804}%
\pgfsetstrokecolor{currentstroke}%
\pgfsetdash{}{0pt}%
\pgfpathmoveto{\pgfqpoint{2.554660in}{0.967500in}}%
\pgfpathlineto{\pgfqpoint{4.400340in}{0.967500in}}%
\pgfusepath{stroke}%
\end{pgfscope}%
\begin{pgfscope}%
\pgfsetbuttcap%
\pgfsetroundjoin%
\definecolor{currentfill}{rgb}{0.000000,0.000000,0.000000}%
\pgfsetfillcolor{currentfill}%
\pgfsetlinewidth{0.803000pt}%
\definecolor{currentstroke}{rgb}{0.000000,0.000000,0.000000}%
\pgfsetstrokecolor{currentstroke}%
\pgfsetdash{}{0pt}%
\pgfsys@defobject{currentmarker}{\pgfqpoint{-0.048611in}{0.000000in}}{\pgfqpoint{-0.000000in}{0.000000in}}{%
\pgfpathmoveto{\pgfqpoint{-0.000000in}{0.000000in}}%
\pgfpathlineto{\pgfqpoint{-0.048611in}{0.000000in}}%
\pgfusepath{stroke,fill}%
}%
\begin{pgfscope}%
\pgfsys@transformshift{2.554660in}{0.967500in}%
\pgfsys@useobject{currentmarker}{}%
\end{pgfscope}%
\end{pgfscope}%
\begin{pgfscope}%
\pgfpathrectangle{\pgfqpoint{2.554660in}{0.320625in}}{\pgfqpoint{1.845679in}{1.164375in}}%
\pgfusepath{clip}%
\pgfsetrectcap%
\pgfsetroundjoin%
\pgfsetlinewidth{0.803000pt}%
\definecolor{currentstroke}{rgb}{0.909804,0.909804,0.909804}%
\pgfsetstrokecolor{currentstroke}%
\pgfsetdash{}{0pt}%
\pgfpathmoveto{\pgfqpoint{2.554660in}{1.183125in}}%
\pgfpathlineto{\pgfqpoint{4.400340in}{1.183125in}}%
\pgfusepath{stroke}%
\end{pgfscope}%
\begin{pgfscope}%
\pgfsetbuttcap%
\pgfsetroundjoin%
\definecolor{currentfill}{rgb}{0.000000,0.000000,0.000000}%
\pgfsetfillcolor{currentfill}%
\pgfsetlinewidth{0.803000pt}%
\definecolor{currentstroke}{rgb}{0.000000,0.000000,0.000000}%
\pgfsetstrokecolor{currentstroke}%
\pgfsetdash{}{0pt}%
\pgfsys@defobject{currentmarker}{\pgfqpoint{-0.048611in}{0.000000in}}{\pgfqpoint{-0.000000in}{0.000000in}}{%
\pgfpathmoveto{\pgfqpoint{-0.000000in}{0.000000in}}%
\pgfpathlineto{\pgfqpoint{-0.048611in}{0.000000in}}%
\pgfusepath{stroke,fill}%
}%
\begin{pgfscope}%
\pgfsys@transformshift{2.554660in}{1.183125in}%
\pgfsys@useobject{currentmarker}{}%
\end{pgfscope}%
\end{pgfscope}%
\begin{pgfscope}%
\pgfpathrectangle{\pgfqpoint{2.554660in}{0.320625in}}{\pgfqpoint{1.845679in}{1.164375in}}%
\pgfusepath{clip}%
\pgfsetrectcap%
\pgfsetroundjoin%
\pgfsetlinewidth{0.803000pt}%
\definecolor{currentstroke}{rgb}{0.909804,0.909804,0.909804}%
\pgfsetstrokecolor{currentstroke}%
\pgfsetdash{}{0pt}%
\pgfpathmoveto{\pgfqpoint{2.554660in}{1.398750in}}%
\pgfpathlineto{\pgfqpoint{4.400340in}{1.398750in}}%
\pgfusepath{stroke}%
\end{pgfscope}%
\begin{pgfscope}%
\pgfsetbuttcap%
\pgfsetroundjoin%
\definecolor{currentfill}{rgb}{0.000000,0.000000,0.000000}%
\pgfsetfillcolor{currentfill}%
\pgfsetlinewidth{0.803000pt}%
\definecolor{currentstroke}{rgb}{0.000000,0.000000,0.000000}%
\pgfsetstrokecolor{currentstroke}%
\pgfsetdash{}{0pt}%
\pgfsys@defobject{currentmarker}{\pgfqpoint{-0.048611in}{0.000000in}}{\pgfqpoint{-0.000000in}{0.000000in}}{%
\pgfpathmoveto{\pgfqpoint{-0.000000in}{0.000000in}}%
\pgfpathlineto{\pgfqpoint{-0.048611in}{0.000000in}}%
\pgfusepath{stroke,fill}%
}%
\begin{pgfscope}%
\pgfsys@transformshift{2.554660in}{1.398750in}%
\pgfsys@useobject{currentmarker}{}%
\end{pgfscope}%
\end{pgfscope}%
\begin{pgfscope}%
\pgfpathrectangle{\pgfqpoint{2.554660in}{0.320625in}}{\pgfqpoint{1.845679in}{1.164375in}}%
\pgfusepath{clip}%
\pgfsetrectcap%
\pgfsetroundjoin%
\pgfsetlinewidth{1.154312pt}%
\definecolor{currentstroke}{rgb}{0.146785,0.068738,0.334011}%
\pgfsetstrokecolor{currentstroke}%
\pgfsetdash{}{0pt}%
\pgfpathmoveto{\pgfqpoint{2.638555in}{1.468525in}}%
\pgfpathlineto{\pgfqpoint{2.653237in}{1.443410in}}%
\pgfpathlineto{\pgfqpoint{2.665821in}{1.419142in}}%
\pgfpathlineto{\pgfqpoint{2.701476in}{1.348595in}}%
\pgfpathlineto{\pgfqpoint{2.716157in}{1.322658in}}%
\pgfpathlineto{\pgfqpoint{2.728742in}{1.302884in}}%
\pgfpathlineto{\pgfqpoint{2.739228in}{1.288502in}}%
\pgfpathlineto{\pgfqpoint{2.749715in}{1.276171in}}%
\pgfpathlineto{\pgfqpoint{2.760202in}{1.265775in}}%
\pgfpathlineto{\pgfqpoint{2.770689in}{1.256975in}}%
\pgfpathlineto{\pgfqpoint{2.783273in}{1.247925in}}%
\pgfpathlineto{\pgfqpoint{2.800052in}{1.237402in}}%
\pgfpathlineto{\pgfqpoint{2.825220in}{1.223142in}}%
\pgfpathlineto{\pgfqpoint{2.877654in}{1.195041in}}%
\pgfpathlineto{\pgfqpoint{3.280348in}{0.982203in}}%
\pgfpathlineto{\pgfqpoint{3.836149in}{0.689041in}}%
\pgfpathlineto{\pgfqpoint{3.842441in}{0.685643in}}%
\pgfpathlineto{\pgfqpoint{3.850831in}{0.680858in}}%
\pgfpathlineto{\pgfqpoint{3.863415in}{0.674351in}}%
\pgfpathlineto{\pgfqpoint{3.869707in}{0.670979in}}%
\pgfpathlineto{\pgfqpoint{3.880194in}{0.666874in}}%
\pgfpathlineto{\pgfqpoint{3.886486in}{0.662037in}}%
\pgfpathlineto{\pgfqpoint{3.890680in}{0.660034in}}%
\pgfpathlineto{\pgfqpoint{3.896973in}{0.657723in}}%
\pgfpathlineto{\pgfqpoint{3.901167in}{0.655092in}}%
\pgfpathlineto{\pgfqpoint{3.907459in}{0.651002in}}%
\pgfpathlineto{\pgfqpoint{3.909557in}{0.650634in}}%
\pgfpathlineto{\pgfqpoint{3.915849in}{0.651511in}}%
\pgfpathlineto{\pgfqpoint{3.917946in}{0.650640in}}%
\pgfpathlineto{\pgfqpoint{3.922141in}{0.645393in}}%
\pgfpathlineto{\pgfqpoint{3.924238in}{0.642673in}}%
\pgfpathlineto{\pgfqpoint{3.926336in}{0.640913in}}%
\pgfpathlineto{\pgfqpoint{3.928433in}{0.640399in}}%
\pgfpathlineto{\pgfqpoint{3.934725in}{0.641505in}}%
\pgfpathlineto{\pgfqpoint{3.936822in}{0.640828in}}%
\pgfpathlineto{\pgfqpoint{3.938920in}{0.639151in}}%
\pgfpathlineto{\pgfqpoint{3.945212in}{0.631011in}}%
\pgfpathlineto{\pgfqpoint{3.947309in}{0.630552in}}%
\pgfpathlineto{\pgfqpoint{3.949407in}{0.632665in}}%
\pgfpathlineto{\pgfqpoint{3.953601in}{0.639659in}}%
\pgfpathlineto{\pgfqpoint{3.955699in}{0.641515in}}%
\pgfpathlineto{\pgfqpoint{3.957796in}{0.640766in}}%
\pgfpathlineto{\pgfqpoint{3.959893in}{0.636562in}}%
\pgfpathlineto{\pgfqpoint{3.964088in}{0.623731in}}%
\pgfpathlineto{\pgfqpoint{3.966185in}{0.619870in}}%
\pgfpathlineto{\pgfqpoint{3.968283in}{0.620677in}}%
\pgfpathlineto{\pgfqpoint{3.972478in}{0.627544in}}%
\pgfpathlineto{\pgfqpoint{3.974575in}{0.629296in}}%
\pgfpathlineto{\pgfqpoint{3.976672in}{0.628941in}}%
\pgfpathlineto{\pgfqpoint{3.978770in}{0.626267in}}%
\pgfpathlineto{\pgfqpoint{3.980867in}{0.621231in}}%
\pgfpathlineto{\pgfqpoint{3.982964in}{0.614572in}}%
\pgfpathlineto{\pgfqpoint{3.985062in}{0.609883in}}%
\pgfpathlineto{\pgfqpoint{3.987159in}{0.612919in}}%
\pgfpathlineto{\pgfqpoint{3.991354in}{0.630104in}}%
\pgfpathlineto{\pgfqpoint{3.993451in}{0.635969in}}%
\pgfpathlineto{\pgfqpoint{3.995549in}{0.638454in}}%
\pgfpathlineto{\pgfqpoint{3.997646in}{0.636612in}}%
\pgfpathlineto{\pgfqpoint{3.999743in}{0.629818in}}%
\pgfpathlineto{\pgfqpoint{4.006035in}{0.598872in}}%
\pgfpathlineto{\pgfqpoint{4.008133in}{0.602209in}}%
\pgfpathlineto{\pgfqpoint{4.010230in}{0.609849in}}%
\pgfpathlineto{\pgfqpoint{4.012327in}{0.615239in}}%
\pgfpathlineto{\pgfqpoint{4.014425in}{0.617247in}}%
\pgfpathlineto{\pgfqpoint{4.016522in}{0.615772in}}%
\pgfpathlineto{\pgfqpoint{4.018620in}{0.610549in}}%
\pgfpathlineto{\pgfqpoint{4.022814in}{0.590831in}}%
\pgfpathlineto{\pgfqpoint{4.024912in}{0.593831in}}%
\pgfpathlineto{\pgfqpoint{4.029106in}{0.620684in}}%
\pgfpathlineto{\pgfqpoint{4.031204in}{0.628749in}}%
\pgfpathlineto{\pgfqpoint{4.033301in}{0.632580in}}%
\pgfpathlineto{\pgfqpoint{4.035398in}{0.631604in}}%
\pgfpathlineto{\pgfqpoint{4.037496in}{0.624811in}}%
\pgfpathlineto{\pgfqpoint{4.039593in}{0.613804in}}%
\pgfpathlineto{\pgfqpoint{4.043788in}{0.580929in}}%
\pgfpathlineto{\pgfqpoint{4.045885in}{0.582395in}}%
\pgfpathlineto{\pgfqpoint{4.047983in}{0.594632in}}%
\pgfpathlineto{\pgfqpoint{4.050080in}{0.602707in}}%
\pgfpathlineto{\pgfqpoint{4.052177in}{0.605898in}}%
\pgfpathlineto{\pgfqpoint{4.054275in}{0.604660in}}%
\pgfpathlineto{\pgfqpoint{4.056372in}{0.598677in}}%
\pgfpathlineto{\pgfqpoint{4.058469in}{0.586724in}}%
\pgfpathlineto{\pgfqpoint{4.060567in}{0.571082in}}%
\pgfpathlineto{\pgfqpoint{4.062664in}{0.579274in}}%
\pgfpathlineto{\pgfqpoint{4.064762in}{0.599248in}}%
\pgfpathlineto{\pgfqpoint{4.066859in}{0.613108in}}%
\pgfpathlineto{\pgfqpoint{4.068956in}{0.621499in}}%
\pgfpathlineto{\pgfqpoint{4.071054in}{0.625199in}}%
\pgfpathlineto{\pgfqpoint{4.073151in}{0.623658in}}%
\pgfpathlineto{\pgfqpoint{4.075248in}{0.615914in}}%
\pgfpathlineto{\pgfqpoint{4.077346in}{0.602908in}}%
\pgfpathlineto{\pgfqpoint{4.081540in}{0.558995in}}%
\pgfpathlineto{\pgfqpoint{4.085735in}{0.592544in}}%
\pgfpathlineto{\pgfqpoint{4.087832in}{0.601240in}}%
\pgfpathlineto{\pgfqpoint{4.089930in}{0.604498in}}%
\pgfpathlineto{\pgfqpoint{4.092027in}{0.603243in}}%
\pgfpathlineto{\pgfqpoint{4.094125in}{0.597098in}}%
\pgfpathlineto{\pgfqpoint{4.096222in}{0.583929in}}%
\pgfpathlineto{\pgfqpoint{4.098319in}{0.558730in}}%
\pgfpathlineto{\pgfqpoint{4.100417in}{0.561807in}}%
\pgfpathlineto{\pgfqpoint{4.102514in}{0.590666in}}%
\pgfpathlineto{\pgfqpoint{4.104611in}{0.607591in}}%
\pgfpathlineto{\pgfqpoint{4.106709in}{0.617559in}}%
\pgfpathlineto{\pgfqpoint{4.108806in}{0.622531in}}%
\pgfpathlineto{\pgfqpoint{4.110903in}{0.622551in}}%
\pgfpathlineto{\pgfqpoint{4.113001in}{0.616203in}}%
\pgfpathlineto{\pgfqpoint{4.115098in}{0.603537in}}%
\pgfpathlineto{\pgfqpoint{4.117196in}{0.580539in}}%
\pgfpathlineto{\pgfqpoint{4.119293in}{0.539317in}}%
\pgfpathlineto{\pgfqpoint{4.121390in}{0.574225in}}%
\pgfpathlineto{\pgfqpoint{4.123488in}{0.595942in}}%
\pgfpathlineto{\pgfqpoint{4.125585in}{0.606619in}}%
\pgfpathlineto{\pgfqpoint{4.127682in}{0.611502in}}%
\pgfpathlineto{\pgfqpoint{4.129780in}{0.612118in}}%
\pgfpathlineto{\pgfqpoint{4.131877in}{0.608604in}}%
\pgfpathlineto{\pgfqpoint{4.133974in}{0.599852in}}%
\pgfpathlineto{\pgfqpoint{4.136072in}{0.581834in}}%
\pgfpathlineto{\pgfqpoint{4.138169in}{0.537308in}}%
\pgfpathlineto{\pgfqpoint{4.142364in}{0.599917in}}%
\pgfpathlineto{\pgfqpoint{4.144461in}{0.615402in}}%
\pgfpathlineto{\pgfqpoint{4.146559in}{0.624312in}}%
\pgfpathlineto{\pgfqpoint{4.148656in}{0.628438in}}%
\pgfpathlineto{\pgfqpoint{4.150753in}{0.627646in}}%
\pgfpathlineto{\pgfqpoint{4.152851in}{0.620488in}}%
\pgfpathlineto{\pgfqpoint{4.154948in}{0.606603in}}%
\pgfpathlineto{\pgfqpoint{4.157045in}{0.579760in}}%
\pgfpathlineto{\pgfqpoint{4.159143in}{0.521751in}}%
\pgfpathlineto{\pgfqpoint{4.161240in}{0.583101in}}%
\pgfpathlineto{\pgfqpoint{4.163338in}{0.603877in}}%
\pgfpathlineto{\pgfqpoint{4.165435in}{0.613922in}}%
\pgfpathlineto{\pgfqpoint{4.167532in}{0.618381in}}%
\pgfpathlineto{\pgfqpoint{4.169630in}{0.618649in}}%
\pgfpathlineto{\pgfqpoint{4.171727in}{0.614770in}}%
\pgfpathlineto{\pgfqpoint{4.173824in}{0.605481in}}%
\pgfpathlineto{\pgfqpoint{4.175922in}{0.586084in}}%
\pgfpathlineto{\pgfqpoint{4.178019in}{0.524632in}}%
\pgfpathlineto{\pgfqpoint{4.180116in}{0.577773in}}%
\pgfpathlineto{\pgfqpoint{4.182214in}{0.607152in}}%
\pgfpathlineto{\pgfqpoint{4.184311in}{0.622371in}}%
\pgfpathlineto{\pgfqpoint{4.186409in}{0.631078in}}%
\pgfpathlineto{\pgfqpoint{4.188506in}{0.635026in}}%
\pgfpathlineto{\pgfqpoint{4.190603in}{0.634033in}}%
\pgfpathlineto{\pgfqpoint{4.192701in}{0.626801in}}%
\pgfpathlineto{\pgfqpoint{4.194798in}{0.613682in}}%
\pgfpathlineto{\pgfqpoint{4.196895in}{0.590146in}}%
\pgfpathlineto{\pgfqpoint{4.198993in}{0.519352in}}%
\pgfpathlineto{\pgfqpoint{4.201090in}{0.574781in}}%
\pgfpathlineto{\pgfqpoint{4.203187in}{0.599911in}}%
\pgfpathlineto{\pgfqpoint{4.205285in}{0.610880in}}%
\pgfpathlineto{\pgfqpoint{4.207382in}{0.615285in}}%
\pgfpathlineto{\pgfqpoint{4.209480in}{0.614868in}}%
\pgfpathlineto{\pgfqpoint{4.211577in}{0.609502in}}%
\pgfpathlineto{\pgfqpoint{4.213674in}{0.596990in}}%
\pgfpathlineto{\pgfqpoint{4.217869in}{0.539882in}}%
\pgfpathlineto{\pgfqpoint{4.219966in}{0.594136in}}%
\pgfpathlineto{\pgfqpoint{4.222064in}{0.615163in}}%
\pgfpathlineto{\pgfqpoint{4.224161in}{0.626927in}}%
\pgfpathlineto{\pgfqpoint{4.226258in}{0.633179in}}%
\pgfpathlineto{\pgfqpoint{4.228356in}{0.634551in}}%
\pgfpathlineto{\pgfqpoint{4.230453in}{0.629867in}}%
\pgfpathlineto{\pgfqpoint{4.232551in}{0.619761in}}%
\pgfpathlineto{\pgfqpoint{4.234648in}{0.603302in}}%
\pgfpathlineto{\pgfqpoint{4.236745in}{0.572560in}}%
\pgfpathlineto{\pgfqpoint{4.238843in}{0.510999in}}%
\pgfpathlineto{\pgfqpoint{4.240940in}{0.578158in}}%
\pgfpathlineto{\pgfqpoint{4.243037in}{0.596133in}}%
\pgfpathlineto{\pgfqpoint{4.245135in}{0.603430in}}%
\pgfpathlineto{\pgfqpoint{4.247232in}{0.604530in}}%
\pgfpathlineto{\pgfqpoint{4.249329in}{0.599998in}}%
\pgfpathlineto{\pgfqpoint{4.251427in}{0.587777in}}%
\pgfpathlineto{\pgfqpoint{4.253524in}{0.556950in}}%
\pgfpathlineto{\pgfqpoint{4.255621in}{0.536756in}}%
\pgfpathlineto{\pgfqpoint{4.257719in}{0.588112in}}%
\pgfpathlineto{\pgfqpoint{4.259816in}{0.608730in}}%
\pgfpathlineto{\pgfqpoint{4.261914in}{0.620274in}}%
\pgfpathlineto{\pgfqpoint{4.264011in}{0.626244in}}%
\pgfpathlineto{\pgfqpoint{4.266108in}{0.627101in}}%
\pgfpathlineto{\pgfqpoint{4.268206in}{0.621552in}}%
\pgfpathlineto{\pgfqpoint{4.270303in}{0.610568in}}%
\pgfpathlineto{\pgfqpoint{4.272400in}{0.592306in}}%
\pgfpathlineto{\pgfqpoint{4.274498in}{0.554923in}}%
\pgfpathlineto{\pgfqpoint{4.276595in}{0.534102in}}%
\pgfpathlineto{\pgfqpoint{4.278692in}{0.578232in}}%
\pgfpathlineto{\pgfqpoint{4.280790in}{0.593332in}}%
\pgfpathlineto{\pgfqpoint{4.282887in}{0.599472in}}%
\pgfpathlineto{\pgfqpoint{4.284985in}{0.599908in}}%
\pgfpathlineto{\pgfqpoint{4.287082in}{0.594887in}}%
\pgfpathlineto{\pgfqpoint{4.289179in}{0.582152in}}%
\pgfpathlineto{\pgfqpoint{4.291277in}{0.550040in}}%
\pgfpathlineto{\pgfqpoint{4.293374in}{0.534139in}}%
\pgfpathlineto{\pgfqpoint{4.295471in}{0.583488in}}%
\pgfpathlineto{\pgfqpoint{4.297569in}{0.603694in}}%
\pgfpathlineto{\pgfqpoint{4.299666in}{0.615032in}}%
\pgfpathlineto{\pgfqpoint{4.301763in}{0.620868in}}%
\pgfpathlineto{\pgfqpoint{4.303861in}{0.621630in}}%
\pgfpathlineto{\pgfqpoint{4.308056in}{0.609632in}}%
\pgfpathlineto{\pgfqpoint{4.316445in}{0.582819in}}%
\pgfpathlineto{\pgfqpoint{4.316445in}{0.582819in}}%
\pgfusepath{stroke}%
\end{pgfscope}%
\begin{pgfscope}%
\pgfpathrectangle{\pgfqpoint{2.554660in}{0.320625in}}{\pgfqpoint{1.845679in}{1.164375in}}%
\pgfusepath{clip}%
\pgfsetrectcap%
\pgfsetroundjoin%
\pgfsetlinewidth{1.154312pt}%
\definecolor{currentstroke}{rgb}{0.291366,0.064553,0.475462}%
\pgfsetstrokecolor{currentstroke}%
\pgfsetdash{}{0pt}%
\pgfpathmoveto{\pgfqpoint{2.638555in}{1.468525in}}%
\pgfpathlineto{\pgfqpoint{2.651139in}{1.446841in}}%
\pgfpathlineto{\pgfqpoint{2.661626in}{1.426459in}}%
\pgfpathlineto{\pgfqpoint{2.680502in}{1.386621in}}%
\pgfpathlineto{\pgfqpoint{2.701476in}{1.343293in}}%
\pgfpathlineto{\pgfqpoint{2.716157in}{1.315550in}}%
\pgfpathlineto{\pgfqpoint{2.728742in}{1.294055in}}%
\pgfpathlineto{\pgfqpoint{2.739228in}{1.277992in}}%
\pgfpathlineto{\pgfqpoint{2.749715in}{1.263706in}}%
\pgfpathlineto{\pgfqpoint{2.760202in}{1.251126in}}%
\pgfpathlineto{\pgfqpoint{2.772786in}{1.237951in}}%
\pgfpathlineto{\pgfqpoint{2.787468in}{1.224528in}}%
\pgfpathlineto{\pgfqpoint{2.804247in}{1.210775in}}%
\pgfpathlineto{\pgfqpoint{2.829415in}{1.191702in}}%
\pgfpathlineto{\pgfqpoint{2.877654in}{1.156844in}}%
\pgfpathlineto{\pgfqpoint{3.070612in}{1.019655in}}%
\pgfpathlineto{\pgfqpoint{3.515253in}{0.704447in}}%
\pgfpathlineto{\pgfqpoint{3.534129in}{0.691199in}}%
\pgfpathlineto{\pgfqpoint{3.550908in}{0.679476in}}%
\pgfpathlineto{\pgfqpoint{3.567687in}{0.668106in}}%
\pgfpathlineto{\pgfqpoint{3.578173in}{0.659320in}}%
\pgfpathlineto{\pgfqpoint{3.584465in}{0.654313in}}%
\pgfpathlineto{\pgfqpoint{3.588660in}{0.652198in}}%
\pgfpathlineto{\pgfqpoint{3.592855in}{0.651695in}}%
\pgfpathlineto{\pgfqpoint{3.597050in}{0.652691in}}%
\pgfpathlineto{\pgfqpoint{3.607536in}{0.656270in}}%
\pgfpathlineto{\pgfqpoint{3.611731in}{0.655726in}}%
\pgfpathlineto{\pgfqpoint{3.613829in}{0.654582in}}%
\pgfpathlineto{\pgfqpoint{3.615926in}{0.652666in}}%
\pgfpathlineto{\pgfqpoint{3.618023in}{0.649785in}}%
\pgfpathlineto{\pgfqpoint{3.628510in}{0.629809in}}%
\pgfpathlineto{\pgfqpoint{3.632705in}{0.625377in}}%
\pgfpathlineto{\pgfqpoint{3.638997in}{0.621046in}}%
\pgfpathlineto{\pgfqpoint{3.643192in}{0.619763in}}%
\pgfpathlineto{\pgfqpoint{3.647386in}{0.619378in}}%
\pgfpathlineto{\pgfqpoint{3.649484in}{0.618099in}}%
\pgfpathlineto{\pgfqpoint{3.659971in}{0.606128in}}%
\pgfpathlineto{\pgfqpoint{3.668360in}{0.599810in}}%
\pgfpathlineto{\pgfqpoint{3.674652in}{0.595992in}}%
\pgfpathlineto{\pgfqpoint{3.680944in}{0.592289in}}%
\pgfpathlineto{\pgfqpoint{3.685139in}{0.587889in}}%
\pgfpathlineto{\pgfqpoint{3.687236in}{0.586810in}}%
\pgfpathlineto{\pgfqpoint{3.689334in}{0.587818in}}%
\pgfpathlineto{\pgfqpoint{3.693528in}{0.594163in}}%
\pgfpathlineto{\pgfqpoint{3.697723in}{0.599739in}}%
\pgfpathlineto{\pgfqpoint{3.699820in}{0.600919in}}%
\pgfpathlineto{\pgfqpoint{3.701918in}{0.600762in}}%
\pgfpathlineto{\pgfqpoint{3.704015in}{0.599112in}}%
\pgfpathlineto{\pgfqpoint{3.706113in}{0.595764in}}%
\pgfpathlineto{\pgfqpoint{3.708210in}{0.590437in}}%
\pgfpathlineto{\pgfqpoint{3.714502in}{0.566967in}}%
\pgfpathlineto{\pgfqpoint{3.716599in}{0.571149in}}%
\pgfpathlineto{\pgfqpoint{3.720794in}{0.591932in}}%
\pgfpathlineto{\pgfqpoint{3.722891in}{0.599440in}}%
\pgfpathlineto{\pgfqpoint{3.724989in}{0.604349in}}%
\pgfpathlineto{\pgfqpoint{3.727086in}{0.606637in}}%
\pgfpathlineto{\pgfqpoint{3.729184in}{0.605939in}}%
\pgfpathlineto{\pgfqpoint{3.731281in}{0.601164in}}%
\pgfpathlineto{\pgfqpoint{3.733378in}{0.590643in}}%
\pgfpathlineto{\pgfqpoint{3.737573in}{0.550793in}}%
\pgfpathlineto{\pgfqpoint{3.741768in}{0.586900in}}%
\pgfpathlineto{\pgfqpoint{3.743865in}{0.597840in}}%
\pgfpathlineto{\pgfqpoint{3.745962in}{0.603795in}}%
\pgfpathlineto{\pgfqpoint{3.748060in}{0.606069in}}%
\pgfpathlineto{\pgfqpoint{3.750157in}{0.605046in}}%
\pgfpathlineto{\pgfqpoint{3.752254in}{0.600434in}}%
\pgfpathlineto{\pgfqpoint{3.754352in}{0.590975in}}%
\pgfpathlineto{\pgfqpoint{3.756449in}{0.572949in}}%
\pgfpathlineto{\pgfqpoint{3.758547in}{0.538428in}}%
\pgfpathlineto{\pgfqpoint{3.762741in}{0.590913in}}%
\pgfpathlineto{\pgfqpoint{3.764839in}{0.605813in}}%
\pgfpathlineto{\pgfqpoint{3.766936in}{0.614640in}}%
\pgfpathlineto{\pgfqpoint{3.769033in}{0.619034in}}%
\pgfpathlineto{\pgfqpoint{3.771131in}{0.619022in}}%
\pgfpathlineto{\pgfqpoint{3.773228in}{0.613172in}}%
\pgfpathlineto{\pgfqpoint{3.775325in}{0.600030in}}%
\pgfpathlineto{\pgfqpoint{3.777423in}{0.573167in}}%
\pgfpathlineto{\pgfqpoint{3.779520in}{0.528239in}}%
\pgfpathlineto{\pgfqpoint{3.781618in}{0.582298in}}%
\pgfpathlineto{\pgfqpoint{3.783715in}{0.602899in}}%
\pgfpathlineto{\pgfqpoint{3.785812in}{0.613571in}}%
\pgfpathlineto{\pgfqpoint{3.787910in}{0.619048in}}%
\pgfpathlineto{\pgfqpoint{3.790007in}{0.620779in}}%
\pgfpathlineto{\pgfqpoint{3.792104in}{0.619097in}}%
\pgfpathlineto{\pgfqpoint{3.794202in}{0.613542in}}%
\pgfpathlineto{\pgfqpoint{3.796299in}{0.602358in}}%
\pgfpathlineto{\pgfqpoint{3.798396in}{0.579332in}}%
\pgfpathlineto{\pgfqpoint{3.800494in}{0.506303in}}%
\pgfpathlineto{\pgfqpoint{3.802591in}{0.583656in}}%
\pgfpathlineto{\pgfqpoint{3.804689in}{0.609051in}}%
\pgfpathlineto{\pgfqpoint{3.806786in}{0.622872in}}%
\pgfpathlineto{\pgfqpoint{3.808883in}{0.630884in}}%
\pgfpathlineto{\pgfqpoint{3.810981in}{0.634501in}}%
\pgfpathlineto{\pgfqpoint{3.813078in}{0.633518in}}%
\pgfpathlineto{\pgfqpoint{3.815175in}{0.626391in}}%
\pgfpathlineto{\pgfqpoint{3.817273in}{0.611895in}}%
\pgfpathlineto{\pgfqpoint{3.819370in}{0.581682in}}%
\pgfpathlineto{\pgfqpoint{3.821467in}{0.536916in}}%
\pgfpathlineto{\pgfqpoint{3.823565in}{0.596544in}}%
\pgfpathlineto{\pgfqpoint{3.825662in}{0.615991in}}%
\pgfpathlineto{\pgfqpoint{3.827760in}{0.625837in}}%
\pgfpathlineto{\pgfqpoint{3.829857in}{0.630542in}}%
\pgfpathlineto{\pgfqpoint{3.831954in}{0.631407in}}%
\pgfpathlineto{\pgfqpoint{3.834052in}{0.628594in}}%
\pgfpathlineto{\pgfqpoint{3.836149in}{0.621297in}}%
\pgfpathlineto{\pgfqpoint{3.838246in}{0.606676in}}%
\pgfpathlineto{\pgfqpoint{3.840344in}{0.571682in}}%
\pgfpathlineto{\pgfqpoint{3.842441in}{0.562437in}}%
\pgfpathlineto{\pgfqpoint{3.844538in}{0.608121in}}%
\pgfpathlineto{\pgfqpoint{3.846636in}{0.627529in}}%
\pgfpathlineto{\pgfqpoint{3.848733in}{0.638635in}}%
\pgfpathlineto{\pgfqpoint{3.850831in}{0.644685in}}%
\pgfpathlineto{\pgfqpoint{3.852928in}{0.646303in}}%
\pgfpathlineto{\pgfqpoint{3.855025in}{0.642514in}}%
\pgfpathlineto{\pgfqpoint{3.857123in}{0.633223in}}%
\pgfpathlineto{\pgfqpoint{3.859220in}{0.618055in}}%
\pgfpathlineto{\pgfqpoint{3.861317in}{0.590770in}}%
\pgfpathlineto{\pgfqpoint{3.863415in}{0.475300in}}%
\pgfpathlineto{\pgfqpoint{3.865512in}{0.585603in}}%
\pgfpathlineto{\pgfqpoint{3.867609in}{0.606488in}}%
\pgfpathlineto{\pgfqpoint{3.869707in}{0.615168in}}%
\pgfpathlineto{\pgfqpoint{3.871804in}{0.617422in}}%
\pgfpathlineto{\pgfqpoint{3.873902in}{0.614290in}}%
\pgfpathlineto{\pgfqpoint{3.875999in}{0.604506in}}%
\pgfpathlineto{\pgfqpoint{3.878096in}{0.581446in}}%
\pgfpathlineto{\pgfqpoint{3.880194in}{0.489616in}}%
\pgfpathlineto{\pgfqpoint{3.882291in}{0.591137in}}%
\pgfpathlineto{\pgfqpoint{3.884388in}{0.616130in}}%
\pgfpathlineto{\pgfqpoint{3.886486in}{0.629670in}}%
\pgfpathlineto{\pgfqpoint{3.888583in}{0.637137in}}%
\pgfpathlineto{\pgfqpoint{3.890680in}{0.639632in}}%
\pgfpathlineto{\pgfqpoint{3.892778in}{0.636248in}}%
\pgfpathlineto{\pgfqpoint{3.894875in}{0.627555in}}%
\pgfpathlineto{\pgfqpoint{3.896973in}{0.614497in}}%
\pgfpathlineto{\pgfqpoint{3.899070in}{0.594133in}}%
\pgfpathlineto{\pgfqpoint{3.901167in}{0.553460in}}%
\pgfpathlineto{\pgfqpoint{3.903265in}{0.537039in}}%
\pgfpathlineto{\pgfqpoint{3.905362in}{0.576423in}}%
\pgfpathlineto{\pgfqpoint{3.907459in}{0.588427in}}%
\pgfpathlineto{\pgfqpoint{3.909557in}{0.590715in}}%
\pgfpathlineto{\pgfqpoint{3.911654in}{0.585207in}}%
\pgfpathlineto{\pgfqpoint{3.913751in}{0.568169in}}%
\pgfpathlineto{\pgfqpoint{3.915849in}{0.498807in}}%
\pgfpathlineto{\pgfqpoint{3.917946in}{0.563136in}}%
\pgfpathlineto{\pgfqpoint{3.920043in}{0.593695in}}%
\pgfpathlineto{\pgfqpoint{3.922141in}{0.609519in}}%
\pgfpathlineto{\pgfqpoint{3.924238in}{0.618453in}}%
\pgfpathlineto{\pgfqpoint{3.926336in}{0.622137in}}%
\pgfpathlineto{\pgfqpoint{3.928433in}{0.619970in}}%
\pgfpathlineto{\pgfqpoint{3.930530in}{0.611528in}}%
\pgfpathlineto{\pgfqpoint{3.932628in}{0.597853in}}%
\pgfpathlineto{\pgfqpoint{3.934725in}{0.574576in}}%
\pgfpathlineto{\pgfqpoint{3.936822in}{0.510744in}}%
\pgfpathlineto{\pgfqpoint{3.938920in}{0.550737in}}%
\pgfpathlineto{\pgfqpoint{3.941017in}{0.577293in}}%
\pgfpathlineto{\pgfqpoint{3.943114in}{0.587769in}}%
\pgfpathlineto{\pgfqpoint{3.945212in}{0.590969in}}%
\pgfpathlineto{\pgfqpoint{3.947309in}{0.588509in}}%
\pgfpathlineto{\pgfqpoint{3.949407in}{0.579395in}}%
\pgfpathlineto{\pgfqpoint{3.951504in}{0.557596in}}%
\pgfpathlineto{\pgfqpoint{3.953601in}{0.420463in}}%
\pgfpathlineto{\pgfqpoint{3.955699in}{0.563816in}}%
\pgfpathlineto{\pgfqpoint{3.957796in}{0.589841in}}%
\pgfpathlineto{\pgfqpoint{3.959893in}{0.603808in}}%
\pgfpathlineto{\pgfqpoint{3.961991in}{0.611546in}}%
\pgfpathlineto{\pgfqpoint{3.964088in}{0.614264in}}%
\pgfpathlineto{\pgfqpoint{3.966185in}{0.611108in}}%
\pgfpathlineto{\pgfqpoint{3.968283in}{0.601018in}}%
\pgfpathlineto{\pgfqpoint{3.970380in}{0.581655in}}%
\pgfpathlineto{\pgfqpoint{3.972478in}{0.530701in}}%
\pgfpathlineto{\pgfqpoint{3.974575in}{0.553857in}}%
\pgfpathlineto{\pgfqpoint{3.976672in}{0.586259in}}%
\pgfpathlineto{\pgfqpoint{3.978770in}{0.600552in}}%
\pgfpathlineto{\pgfqpoint{3.980867in}{0.607860in}}%
\pgfpathlineto{\pgfqpoint{3.982964in}{0.610763in}}%
\pgfpathlineto{\pgfqpoint{3.985062in}{0.609975in}}%
\pgfpathlineto{\pgfqpoint{3.987159in}{0.605251in}}%
\pgfpathlineto{\pgfqpoint{3.989256in}{0.595088in}}%
\pgfpathlineto{\pgfqpoint{3.991354in}{0.574134in}}%
\pgfpathlineto{\pgfqpoint{3.993451in}{0.485850in}}%
\pgfpathlineto{\pgfqpoint{3.995549in}{0.570268in}}%
\pgfpathlineto{\pgfqpoint{3.997646in}{0.598352in}}%
\pgfpathlineto{\pgfqpoint{3.999743in}{0.613096in}}%
\pgfpathlineto{\pgfqpoint{4.001841in}{0.621605in}}%
\pgfpathlineto{\pgfqpoint{4.003938in}{0.625542in}}%
\pgfpathlineto{\pgfqpoint{4.006035in}{0.624782in}}%
\pgfpathlineto{\pgfqpoint{4.008133in}{0.617663in}}%
\pgfpathlineto{\pgfqpoint{4.010230in}{0.602178in}}%
\pgfpathlineto{\pgfqpoint{4.012327in}{0.566306in}}%
\pgfpathlineto{\pgfqpoint{4.014425in}{0.552149in}}%
\pgfpathlineto{\pgfqpoint{4.016522in}{0.597229in}}%
\pgfpathlineto{\pgfqpoint{4.018620in}{0.614902in}}%
\pgfpathlineto{\pgfqpoint{4.020717in}{0.624334in}}%
\pgfpathlineto{\pgfqpoint{4.022814in}{0.629118in}}%
\pgfpathlineto{\pgfqpoint{4.024912in}{0.630395in}}%
\pgfpathlineto{\pgfqpoint{4.027009in}{0.628374in}}%
\pgfpathlineto{\pgfqpoint{4.029106in}{0.622524in}}%
\pgfpathlineto{\pgfqpoint{4.031204in}{0.610993in}}%
\pgfpathlineto{\pgfqpoint{4.033301in}{0.587148in}}%
\pgfpathlineto{\pgfqpoint{4.035398in}{0.470385in}}%
\pgfpathlineto{\pgfqpoint{4.037496in}{0.593084in}}%
\pgfpathlineto{\pgfqpoint{4.039593in}{0.618263in}}%
\pgfpathlineto{\pgfqpoint{4.041691in}{0.631979in}}%
\pgfpathlineto{\pgfqpoint{4.043788in}{0.639958in}}%
\pgfpathlineto{\pgfqpoint{4.045885in}{0.643611in}}%
\pgfpathlineto{\pgfqpoint{4.047983in}{0.642770in}}%
\pgfpathlineto{\pgfqpoint{4.050080in}{0.636054in}}%
\pgfpathlineto{\pgfqpoint{4.052177in}{0.623268in}}%
\pgfpathlineto{\pgfqpoint{4.054275in}{0.599847in}}%
\pgfpathlineto{\pgfqpoint{4.056372in}{0.520839in}}%
\pgfpathlineto{\pgfqpoint{4.058469in}{0.585856in}}%
\pgfpathlineto{\pgfqpoint{4.060567in}{0.611360in}}%
\pgfpathlineto{\pgfqpoint{4.062664in}{0.622783in}}%
\pgfpathlineto{\pgfqpoint{4.064762in}{0.627796in}}%
\pgfpathlineto{\pgfqpoint{4.066859in}{0.628247in}}%
\pgfpathlineto{\pgfqpoint{4.068956in}{0.624267in}}%
\pgfpathlineto{\pgfqpoint{4.071054in}{0.614464in}}%
\pgfpathlineto{\pgfqpoint{4.073151in}{0.593251in}}%
\pgfpathlineto{\pgfqpoint{4.075248in}{0.486720in}}%
\pgfpathlineto{\pgfqpoint{4.077346in}{0.592882in}}%
\pgfpathlineto{\pgfqpoint{4.079443in}{0.620181in}}%
\pgfpathlineto{\pgfqpoint{4.081540in}{0.634606in}}%
\pgfpathlineto{\pgfqpoint{4.083638in}{0.642778in}}%
\pgfpathlineto{\pgfqpoint{4.085735in}{0.646181in}}%
\pgfpathlineto{\pgfqpoint{4.087832in}{0.644402in}}%
\pgfpathlineto{\pgfqpoint{4.089930in}{0.636764in}}%
\pgfpathlineto{\pgfqpoint{4.092027in}{0.624922in}}%
\pgfpathlineto{\pgfqpoint{4.094125in}{0.606669in}}%
\pgfpathlineto{\pgfqpoint{4.096222in}{0.573087in}}%
\pgfpathlineto{\pgfqpoint{4.098319in}{0.521138in}}%
\pgfpathlineto{\pgfqpoint{4.100417in}{0.580104in}}%
\pgfpathlineto{\pgfqpoint{4.102514in}{0.595374in}}%
\pgfpathlineto{\pgfqpoint{4.104611in}{0.599762in}}%
\pgfpathlineto{\pgfqpoint{4.106709in}{0.596701in}}%
\pgfpathlineto{\pgfqpoint{4.108806in}{0.584536in}}%
\pgfpathlineto{\pgfqpoint{4.110903in}{0.549170in}}%
\pgfpathlineto{\pgfqpoint{4.113001in}{0.549289in}}%
\pgfpathlineto{\pgfqpoint{4.115098in}{0.593127in}}%
\pgfpathlineto{\pgfqpoint{4.117196in}{0.612565in}}%
\pgfpathlineto{\pgfqpoint{4.119293in}{0.623582in}}%
\pgfpathlineto{\pgfqpoint{4.121390in}{0.629078in}}%
\pgfpathlineto{\pgfqpoint{4.123488in}{0.629210in}}%
\pgfpathlineto{\pgfqpoint{4.125585in}{0.622822in}}%
\pgfpathlineto{\pgfqpoint{4.127682in}{0.612089in}}%
\pgfpathlineto{\pgfqpoint{4.129780in}{0.595773in}}%
\pgfpathlineto{\pgfqpoint{4.131877in}{0.567837in}}%
\pgfpathlineto{\pgfqpoint{4.133974in}{0.438143in}}%
\pgfpathlineto{\pgfqpoint{4.136072in}{0.556856in}}%
\pgfpathlineto{\pgfqpoint{4.138169in}{0.576602in}}%
\pgfpathlineto{\pgfqpoint{4.140267in}{0.583025in}}%
\pgfpathlineto{\pgfqpoint{4.142364in}{0.581689in}}%
\pgfpathlineto{\pgfqpoint{4.144461in}{0.572060in}}%
\pgfpathlineto{\pgfqpoint{4.146559in}{0.545444in}}%
\pgfpathlineto{\pgfqpoint{4.148656in}{0.508574in}}%
\pgfpathlineto{\pgfqpoint{4.150753in}{0.570405in}}%
\pgfpathlineto{\pgfqpoint{4.152851in}{0.592843in}}%
\pgfpathlineto{\pgfqpoint{4.154948in}{0.605364in}}%
\pgfpathlineto{\pgfqpoint{4.157045in}{0.612028in}}%
\pgfpathlineto{\pgfqpoint{4.159143in}{0.613464in}}%
\pgfpathlineto{\pgfqpoint{4.161240in}{0.608347in}}%
\pgfpathlineto{\pgfqpoint{4.163338in}{0.597007in}}%
\pgfpathlineto{\pgfqpoint{4.165435in}{0.576701in}}%
\pgfpathlineto{\pgfqpoint{4.167532in}{0.525421in}}%
\pgfpathlineto{\pgfqpoint{4.169630in}{0.546459in}}%
\pgfpathlineto{\pgfqpoint{4.171727in}{0.578234in}}%
\pgfpathlineto{\pgfqpoint{4.173824in}{0.591596in}}%
\pgfpathlineto{\pgfqpoint{4.175922in}{0.597747in}}%
\pgfpathlineto{\pgfqpoint{4.178019in}{0.599141in}}%
\pgfpathlineto{\pgfqpoint{4.180116in}{0.596205in}}%
\pgfpathlineto{\pgfqpoint{4.182214in}{0.587965in}}%
\pgfpathlineto{\pgfqpoint{4.184311in}{0.570416in}}%
\pgfpathlineto{\pgfqpoint{4.186409in}{0.516956in}}%
\pgfpathlineto{\pgfqpoint{4.190603in}{0.586085in}}%
\pgfpathlineto{\pgfqpoint{4.192701in}{0.602500in}}%
\pgfpathlineto{\pgfqpoint{4.194798in}{0.611887in}}%
\pgfpathlineto{\pgfqpoint{4.196895in}{0.616363in}}%
\pgfpathlineto{\pgfqpoint{4.198993in}{0.615920in}}%
\pgfpathlineto{\pgfqpoint{4.201090in}{0.608872in}}%
\pgfpathlineto{\pgfqpoint{4.203187in}{0.593144in}}%
\pgfpathlineto{\pgfqpoint{4.205285in}{0.555550in}}%
\pgfpathlineto{\pgfqpoint{4.207382in}{0.548265in}}%
\pgfpathlineto{\pgfqpoint{4.209480in}{0.590779in}}%
\pgfpathlineto{\pgfqpoint{4.211577in}{0.608096in}}%
\pgfpathlineto{\pgfqpoint{4.213674in}{0.617462in}}%
\pgfpathlineto{\pgfqpoint{4.215772in}{0.622293in}}%
\pgfpathlineto{\pgfqpoint{4.217869in}{0.623699in}}%
\pgfpathlineto{\pgfqpoint{4.219966in}{0.621902in}}%
\pgfpathlineto{\pgfqpoint{4.222064in}{0.616435in}}%
\pgfpathlineto{\pgfqpoint{4.224161in}{0.605658in}}%
\pgfpathlineto{\pgfqpoint{4.226258in}{0.583991in}}%
\pgfpathlineto{\pgfqpoint{4.228356in}{0.486761in}}%
\pgfpathlineto{\pgfqpoint{4.230453in}{0.581351in}}%
\pgfpathlineto{\pgfqpoint{4.232551in}{0.608892in}}%
\pgfpathlineto{\pgfqpoint{4.234648in}{0.623506in}}%
\pgfpathlineto{\pgfqpoint{4.236745in}{0.632080in}}%
\pgfpathlineto{\pgfqpoint{4.238843in}{0.636300in}}%
\pgfpathlineto{\pgfqpoint{4.240940in}{0.636198in}}%
\pgfpathlineto{\pgfqpoint{4.243037in}{0.630371in}}%
\pgfpathlineto{\pgfqpoint{4.245135in}{0.617462in}}%
\pgfpathlineto{\pgfqpoint{4.247232in}{0.591272in}}%
\pgfpathlineto{\pgfqpoint{4.249329in}{0.492719in}}%
\pgfpathlineto{\pgfqpoint{4.251427in}{0.594653in}}%
\pgfpathlineto{\pgfqpoint{4.253524in}{0.616901in}}%
\pgfpathlineto{\pgfqpoint{4.255621in}{0.627990in}}%
\pgfpathlineto{\pgfqpoint{4.257719in}{0.633565in}}%
\pgfpathlineto{\pgfqpoint{4.259816in}{0.635232in}}%
\pgfpathlineto{\pgfqpoint{4.261914in}{0.633339in}}%
\pgfpathlineto{\pgfqpoint{4.264011in}{0.627356in}}%
\pgfpathlineto{\pgfqpoint{4.266108in}{0.615247in}}%
\pgfpathlineto{\pgfqpoint{4.268206in}{0.589184in}}%
\pgfpathlineto{\pgfqpoint{4.270303in}{0.524946in}}%
\pgfpathlineto{\pgfqpoint{4.272400in}{0.603089in}}%
\pgfpathlineto{\pgfqpoint{4.274498in}{0.626441in}}%
\pgfpathlineto{\pgfqpoint{4.276595in}{0.639367in}}%
\pgfpathlineto{\pgfqpoint{4.278692in}{0.646734in}}%
\pgfpathlineto{\pgfqpoint{4.280790in}{0.649687in}}%
\pgfpathlineto{\pgfqpoint{4.282887in}{0.647765in}}%
\pgfpathlineto{\pgfqpoint{4.284985in}{0.640548in}}%
\pgfpathlineto{\pgfqpoint{4.287082in}{0.630589in}}%
\pgfpathlineto{\pgfqpoint{4.289179in}{0.617316in}}%
\pgfpathlineto{\pgfqpoint{4.291277in}{0.599158in}}%
\pgfpathlineto{\pgfqpoint{4.293374in}{0.571928in}}%
\pgfpathlineto{\pgfqpoint{4.295471in}{0.513825in}}%
\pgfpathlineto{\pgfqpoint{4.297569in}{0.521931in}}%
\pgfpathlineto{\pgfqpoint{4.299666in}{0.538867in}}%
\pgfpathlineto{\pgfqpoint{4.303861in}{0.505211in}}%
\pgfpathlineto{\pgfqpoint{4.305958in}{0.563338in}}%
\pgfpathlineto{\pgfqpoint{4.308056in}{0.587956in}}%
\pgfpathlineto{\pgfqpoint{4.310153in}{0.602642in}}%
\pgfpathlineto{\pgfqpoint{4.312250in}{0.611350in}}%
\pgfpathlineto{\pgfqpoint{4.314348in}{0.614934in}}%
\pgfpathlineto{\pgfqpoint{4.316445in}{0.612375in}}%
\pgfpathlineto{\pgfqpoint{4.316445in}{0.612375in}}%
\pgfusepath{stroke}%
\end{pgfscope}%
\begin{pgfscope}%
\pgfpathrectangle{\pgfqpoint{2.554660in}{0.320625in}}{\pgfqpoint{1.845679in}{1.164375in}}%
\pgfusepath{clip}%
\pgfsetrectcap%
\pgfsetroundjoin%
\pgfsetlinewidth{1.154312pt}%
\definecolor{currentstroke}{rgb}{0.432967,0.117855,0.506160}%
\pgfsetstrokecolor{currentstroke}%
\pgfsetdash{}{0pt}%
\pgfpathmoveto{\pgfqpoint{2.638555in}{1.468525in}}%
\pgfpathlineto{\pgfqpoint{2.651139in}{1.446491in}}%
\pgfpathlineto{\pgfqpoint{2.661626in}{1.425563in}}%
\pgfpathlineto{\pgfqpoint{2.678405in}{1.388854in}}%
\pgfpathlineto{\pgfqpoint{2.701476in}{1.338683in}}%
\pgfpathlineto{\pgfqpoint{2.716157in}{1.309203in}}%
\pgfpathlineto{\pgfqpoint{2.730839in}{1.282310in}}%
\pgfpathlineto{\pgfqpoint{2.743423in}{1.261542in}}%
\pgfpathlineto{\pgfqpoint{2.756007in}{1.242935in}}%
\pgfpathlineto{\pgfqpoint{2.768591in}{1.226335in}}%
\pgfpathlineto{\pgfqpoint{2.783273in}{1.209047in}}%
\pgfpathlineto{\pgfqpoint{2.800052in}{1.191228in}}%
\pgfpathlineto{\pgfqpoint{2.821026in}{1.170670in}}%
\pgfpathlineto{\pgfqpoint{2.854583in}{1.139567in}}%
\pgfpathlineto{\pgfqpoint{2.923796in}{1.077316in}}%
\pgfpathlineto{\pgfqpoint{3.265666in}{0.772378in}}%
\pgfpathlineto{\pgfqpoint{3.341171in}{0.705295in}}%
\pgfpathlineto{\pgfqpoint{3.372632in}{0.676453in}}%
\pgfpathlineto{\pgfqpoint{3.385216in}{0.666701in}}%
\pgfpathlineto{\pgfqpoint{3.399898in}{0.655629in}}%
\pgfpathlineto{\pgfqpoint{3.410384in}{0.645696in}}%
\pgfpathlineto{\pgfqpoint{3.422969in}{0.632113in}}%
\pgfpathlineto{\pgfqpoint{3.443942in}{0.609097in}}%
\pgfpathlineto{\pgfqpoint{3.450234in}{0.603572in}}%
\pgfpathlineto{\pgfqpoint{3.456526in}{0.599465in}}%
\pgfpathlineto{\pgfqpoint{3.462818in}{0.597105in}}%
\pgfpathlineto{\pgfqpoint{3.469111in}{0.596363in}}%
\pgfpathlineto{\pgfqpoint{3.483792in}{0.595798in}}%
\pgfpathlineto{\pgfqpoint{3.487987in}{0.594092in}}%
\pgfpathlineto{\pgfqpoint{3.492182in}{0.590682in}}%
\pgfpathlineto{\pgfqpoint{3.496376in}{0.585917in}}%
\pgfpathlineto{\pgfqpoint{3.498474in}{0.584380in}}%
\pgfpathlineto{\pgfqpoint{3.500571in}{0.584172in}}%
\pgfpathlineto{\pgfqpoint{3.502668in}{0.585210in}}%
\pgfpathlineto{\pgfqpoint{3.508960in}{0.589553in}}%
\pgfpathlineto{\pgfqpoint{3.511058in}{0.589648in}}%
\pgfpathlineto{\pgfqpoint{3.513155in}{0.588563in}}%
\pgfpathlineto{\pgfqpoint{3.515253in}{0.586096in}}%
\pgfpathlineto{\pgfqpoint{3.517350in}{0.582041in}}%
\pgfpathlineto{\pgfqpoint{3.521545in}{0.568960in}}%
\pgfpathlineto{\pgfqpoint{3.523642in}{0.562194in}}%
\pgfpathlineto{\pgfqpoint{3.525739in}{0.561092in}}%
\pgfpathlineto{\pgfqpoint{3.527837in}{0.567507in}}%
\pgfpathlineto{\pgfqpoint{3.532031in}{0.583831in}}%
\pgfpathlineto{\pgfqpoint{3.534129in}{0.589176in}}%
\pgfpathlineto{\pgfqpoint{3.536226in}{0.592055in}}%
\pgfpathlineto{\pgfqpoint{3.538324in}{0.592099in}}%
\pgfpathlineto{\pgfqpoint{3.540421in}{0.588375in}}%
\pgfpathlineto{\pgfqpoint{3.542518in}{0.576483in}}%
\pgfpathlineto{\pgfqpoint{3.544616in}{0.550056in}}%
\pgfpathlineto{\pgfqpoint{3.546713in}{0.561688in}}%
\pgfpathlineto{\pgfqpoint{3.548810in}{0.590125in}}%
\pgfpathlineto{\pgfqpoint{3.550908in}{0.606791in}}%
\pgfpathlineto{\pgfqpoint{3.553005in}{0.617388in}}%
\pgfpathlineto{\pgfqpoint{3.555102in}{0.624255in}}%
\pgfpathlineto{\pgfqpoint{3.557200in}{0.628411in}}%
\pgfpathlineto{\pgfqpoint{3.559297in}{0.630312in}}%
\pgfpathlineto{\pgfqpoint{3.561395in}{0.630104in}}%
\pgfpathlineto{\pgfqpoint{3.563492in}{0.627670in}}%
\pgfpathlineto{\pgfqpoint{3.565589in}{0.622568in}}%
\pgfpathlineto{\pgfqpoint{3.567687in}{0.613720in}}%
\pgfpathlineto{\pgfqpoint{3.569784in}{0.598286in}}%
\pgfpathlineto{\pgfqpoint{3.571881in}{0.564914in}}%
\pgfpathlineto{\pgfqpoint{3.573979in}{0.551913in}}%
\pgfpathlineto{\pgfqpoint{3.576076in}{0.596664in}}%
\pgfpathlineto{\pgfqpoint{3.578173in}{0.616453in}}%
\pgfpathlineto{\pgfqpoint{3.580271in}{0.628326in}}%
\pgfpathlineto{\pgfqpoint{3.582368in}{0.635768in}}%
\pgfpathlineto{\pgfqpoint{3.584465in}{0.639861in}}%
\pgfpathlineto{\pgfqpoint{3.586563in}{0.640730in}}%
\pgfpathlineto{\pgfqpoint{3.588660in}{0.637656in}}%
\pgfpathlineto{\pgfqpoint{3.590758in}{0.629958in}}%
\pgfpathlineto{\pgfqpoint{3.592855in}{0.618446in}}%
\pgfpathlineto{\pgfqpoint{3.594952in}{0.601290in}}%
\pgfpathlineto{\pgfqpoint{3.597050in}{0.572466in}}%
\pgfpathlineto{\pgfqpoint{3.599147in}{0.521842in}}%
\pgfpathlineto{\pgfqpoint{3.601244in}{0.566000in}}%
\pgfpathlineto{\pgfqpoint{3.603342in}{0.582455in}}%
\pgfpathlineto{\pgfqpoint{3.605439in}{0.586975in}}%
\pgfpathlineto{\pgfqpoint{3.607536in}{0.583322in}}%
\pgfpathlineto{\pgfqpoint{3.609634in}{0.569207in}}%
\pgfpathlineto{\pgfqpoint{3.611731in}{0.523929in}}%
\pgfpathlineto{\pgfqpoint{3.618023in}{0.605336in}}%
\pgfpathlineto{\pgfqpoint{3.620121in}{0.615267in}}%
\pgfpathlineto{\pgfqpoint{3.622218in}{0.619893in}}%
\pgfpathlineto{\pgfqpoint{3.624315in}{0.618983in}}%
\pgfpathlineto{\pgfqpoint{3.626413in}{0.611766in}}%
\pgfpathlineto{\pgfqpoint{3.628510in}{0.600628in}}%
\pgfpathlineto{\pgfqpoint{3.630607in}{0.584015in}}%
\pgfpathlineto{\pgfqpoint{3.632705in}{0.556214in}}%
\pgfpathlineto{\pgfqpoint{3.634802in}{0.489822in}}%
\pgfpathlineto{\pgfqpoint{3.636900in}{0.544186in}}%
\pgfpathlineto{\pgfqpoint{3.638997in}{0.562789in}}%
\pgfpathlineto{\pgfqpoint{3.641094in}{0.568239in}}%
\pgfpathlineto{\pgfqpoint{3.643192in}{0.565423in}}%
\pgfpathlineto{\pgfqpoint{3.645289in}{0.552705in}}%
\pgfpathlineto{\pgfqpoint{3.647386in}{0.513292in}}%
\pgfpathlineto{\pgfqpoint{3.649484in}{0.527652in}}%
\pgfpathlineto{\pgfqpoint{3.651581in}{0.566582in}}%
\pgfpathlineto{\pgfqpoint{3.653678in}{0.585070in}}%
\pgfpathlineto{\pgfqpoint{3.655776in}{0.595646in}}%
\pgfpathlineto{\pgfqpoint{3.657873in}{0.600811in}}%
\pgfpathlineto{\pgfqpoint{3.659971in}{0.600549in}}%
\pgfpathlineto{\pgfqpoint{3.662068in}{0.593227in}}%
\pgfpathlineto{\pgfqpoint{3.664165in}{0.577885in}}%
\pgfpathlineto{\pgfqpoint{3.666263in}{0.543927in}}%
\pgfpathlineto{\pgfqpoint{3.668360in}{0.521786in}}%
\pgfpathlineto{\pgfqpoint{3.670457in}{0.569602in}}%
\pgfpathlineto{\pgfqpoint{3.672555in}{0.587620in}}%
\pgfpathlineto{\pgfqpoint{3.674652in}{0.597059in}}%
\pgfpathlineto{\pgfqpoint{3.676749in}{0.601709in}}%
\pgfpathlineto{\pgfqpoint{3.678847in}{0.602741in}}%
\pgfpathlineto{\pgfqpoint{3.680944in}{0.600331in}}%
\pgfpathlineto{\pgfqpoint{3.683042in}{0.593825in}}%
\pgfpathlineto{\pgfqpoint{3.685139in}{0.580981in}}%
\pgfpathlineto{\pgfqpoint{3.687236in}{0.552938in}}%
\pgfpathlineto{\pgfqpoint{3.689334in}{0.507139in}}%
\pgfpathlineto{\pgfqpoint{3.691431in}{0.572222in}}%
\pgfpathlineto{\pgfqpoint{3.693528in}{0.594435in}}%
\pgfpathlineto{\pgfqpoint{3.695626in}{0.606907in}}%
\pgfpathlineto{\pgfqpoint{3.697723in}{0.614028in}}%
\pgfpathlineto{\pgfqpoint{3.699820in}{0.616840in}}%
\pgfpathlineto{\pgfqpoint{3.701918in}{0.614865in}}%
\pgfpathlineto{\pgfqpoint{3.704015in}{0.605589in}}%
\pgfpathlineto{\pgfqpoint{3.706113in}{0.582668in}}%
\pgfpathlineto{\pgfqpoint{3.708210in}{0.473059in}}%
\pgfpathlineto{\pgfqpoint{3.710307in}{0.587263in}}%
\pgfpathlineto{\pgfqpoint{3.712405in}{0.611978in}}%
\pgfpathlineto{\pgfqpoint{3.714502in}{0.625229in}}%
\pgfpathlineto{\pgfqpoint{3.716599in}{0.633103in}}%
\pgfpathlineto{\pgfqpoint{3.718697in}{0.637469in}}%
\pgfpathlineto{\pgfqpoint{3.720794in}{0.639039in}}%
\pgfpathlineto{\pgfqpoint{3.722891in}{0.637982in}}%
\pgfpathlineto{\pgfqpoint{3.724989in}{0.634033in}}%
\pgfpathlineto{\pgfqpoint{3.727086in}{0.626283in}}%
\pgfpathlineto{\pgfqpoint{3.729184in}{0.612214in}}%
\pgfpathlineto{\pgfqpoint{3.733378in}{0.547796in}}%
\pgfpathlineto{\pgfqpoint{3.735476in}{0.605449in}}%
\pgfpathlineto{\pgfqpoint{3.737573in}{0.626810in}}%
\pgfpathlineto{\pgfqpoint{3.739670in}{0.639122in}}%
\pgfpathlineto{\pgfqpoint{3.741768in}{0.646498in}}%
\pgfpathlineto{\pgfqpoint{3.743865in}{0.650064in}}%
\pgfpathlineto{\pgfqpoint{3.745962in}{0.649735in}}%
\pgfpathlineto{\pgfqpoint{3.748060in}{0.644359in}}%
\pgfpathlineto{\pgfqpoint{3.750157in}{0.635116in}}%
\pgfpathlineto{\pgfqpoint{3.752254in}{0.621925in}}%
\pgfpathlineto{\pgfqpoint{3.754352in}{0.602360in}}%
\pgfpathlineto{\pgfqpoint{3.756449in}{0.567479in}}%
\pgfpathlineto{\pgfqpoint{3.758547in}{0.514921in}}%
\pgfpathlineto{\pgfqpoint{3.760644in}{0.570706in}}%
\pgfpathlineto{\pgfqpoint{3.762741in}{0.583182in}}%
\pgfpathlineto{\pgfqpoint{3.764839in}{0.583203in}}%
\pgfpathlineto{\pgfqpoint{3.766936in}{0.571545in}}%
\pgfpathlineto{\pgfqpoint{3.769033in}{0.528639in}}%
\pgfpathlineto{\pgfqpoint{3.771131in}{0.552801in}}%
\pgfpathlineto{\pgfqpoint{3.773228in}{0.590603in}}%
\pgfpathlineto{\pgfqpoint{3.775325in}{0.609036in}}%
\pgfpathlineto{\pgfqpoint{3.777423in}{0.619687in}}%
\pgfpathlineto{\pgfqpoint{3.779520in}{0.624909in}}%
\pgfpathlineto{\pgfqpoint{3.781618in}{0.624608in}}%
\pgfpathlineto{\pgfqpoint{3.783715in}{0.617855in}}%
\pgfpathlineto{\pgfqpoint{3.785812in}{0.607479in}}%
\pgfpathlineto{\pgfqpoint{3.787910in}{0.592618in}}%
\pgfpathlineto{\pgfqpoint{3.790007in}{0.569783in}}%
\pgfpathlineto{\pgfqpoint{3.792104in}{0.522092in}}%
\pgfpathlineto{\pgfqpoint{3.794202in}{0.521789in}}%
\pgfpathlineto{\pgfqpoint{3.796299in}{0.552758in}}%
\pgfpathlineto{\pgfqpoint{3.798396in}{0.560469in}}%
\pgfpathlineto{\pgfqpoint{3.800494in}{0.557104in}}%
\pgfpathlineto{\pgfqpoint{3.802591in}{0.539870in}}%
\pgfpathlineto{\pgfqpoint{3.804689in}{0.432076in}}%
\pgfpathlineto{\pgfqpoint{3.806786in}{0.545276in}}%
\pgfpathlineto{\pgfqpoint{3.808883in}{0.573854in}}%
\pgfpathlineto{\pgfqpoint{3.810981in}{0.589272in}}%
\pgfpathlineto{\pgfqpoint{3.813078in}{0.598064in}}%
\pgfpathlineto{\pgfqpoint{3.815175in}{0.601609in}}%
\pgfpathlineto{\pgfqpoint{3.817273in}{0.599132in}}%
\pgfpathlineto{\pgfqpoint{3.819370in}{0.589560in}}%
\pgfpathlineto{\pgfqpoint{3.821467in}{0.571391in}}%
\pgfpathlineto{\pgfqpoint{3.823565in}{0.527509in}}%
\pgfpathlineto{\pgfqpoint{3.825662in}{0.534522in}}%
\pgfpathlineto{\pgfqpoint{3.827760in}{0.570949in}}%
\pgfpathlineto{\pgfqpoint{3.829857in}{0.586218in}}%
\pgfpathlineto{\pgfqpoint{3.831954in}{0.594004in}}%
\pgfpathlineto{\pgfqpoint{3.834052in}{0.597227in}}%
\pgfpathlineto{\pgfqpoint{3.836149in}{0.596713in}}%
\pgfpathlineto{\pgfqpoint{3.838246in}{0.592287in}}%
\pgfpathlineto{\pgfqpoint{3.840344in}{0.582564in}}%
\pgfpathlineto{\pgfqpoint{3.842441in}{0.562636in}}%
\pgfpathlineto{\pgfqpoint{3.844538in}{0.490386in}}%
\pgfpathlineto{\pgfqpoint{3.846636in}{0.554330in}}%
\pgfpathlineto{\pgfqpoint{3.848733in}{0.583881in}}%
\pgfpathlineto{\pgfqpoint{3.850831in}{0.599139in}}%
\pgfpathlineto{\pgfqpoint{3.852928in}{0.607968in}}%
\pgfpathlineto{\pgfqpoint{3.855025in}{0.612190in}}%
\pgfpathlineto{\pgfqpoint{3.857123in}{0.611776in}}%
\pgfpathlineto{\pgfqpoint{3.859220in}{0.604817in}}%
\pgfpathlineto{\pgfqpoint{3.861317in}{0.586439in}}%
\pgfpathlineto{\pgfqpoint{3.863415in}{0.522215in}}%
\pgfpathlineto{\pgfqpoint{3.865512in}{0.574254in}}%
\pgfpathlineto{\pgfqpoint{3.867609in}{0.604394in}}%
\pgfpathlineto{\pgfqpoint{3.869707in}{0.619567in}}%
\pgfpathlineto{\pgfqpoint{3.871804in}{0.628565in}}%
\pgfpathlineto{\pgfqpoint{3.873902in}{0.633777in}}%
\pgfpathlineto{\pgfqpoint{3.875999in}{0.636106in}}%
\pgfpathlineto{\pgfqpoint{3.878096in}{0.635843in}}%
\pgfpathlineto{\pgfqpoint{3.880194in}{0.632860in}}%
\pgfpathlineto{\pgfqpoint{3.882291in}{0.626507in}}%
\pgfpathlineto{\pgfqpoint{3.884388in}{0.615021in}}%
\pgfpathlineto{\pgfqpoint{3.886486in}{0.592631in}}%
\pgfpathlineto{\pgfqpoint{3.888583in}{0.486245in}}%
\pgfpathlineto{\pgfqpoint{3.890680in}{0.590817in}}%
\pgfpathlineto{\pgfqpoint{3.892778in}{0.617981in}}%
\pgfpathlineto{\pgfqpoint{3.894875in}{0.632617in}}%
\pgfpathlineto{\pgfqpoint{3.896973in}{0.641483in}}%
\pgfpathlineto{\pgfqpoint{3.899070in}{0.646347in}}%
\pgfpathlineto{\pgfqpoint{3.901167in}{0.647505in}}%
\pgfpathlineto{\pgfqpoint{3.903265in}{0.644128in}}%
\pgfpathlineto{\pgfqpoint{3.905362in}{0.636015in}}%
\pgfpathlineto{\pgfqpoint{3.907459in}{0.624148in}}%
\pgfpathlineto{\pgfqpoint{3.909557in}{0.606568in}}%
\pgfpathlineto{\pgfqpoint{3.911654in}{0.576623in}}%
\pgfpathlineto{\pgfqpoint{3.913751in}{0.448907in}}%
\pgfpathlineto{\pgfqpoint{3.915849in}{0.567736in}}%
\pgfpathlineto{\pgfqpoint{3.917946in}{0.584794in}}%
\pgfpathlineto{\pgfqpoint{3.920043in}{0.588503in}}%
\pgfpathlineto{\pgfqpoint{3.922141in}{0.582899in}}%
\pgfpathlineto{\pgfqpoint{3.924238in}{0.563288in}}%
\pgfpathlineto{\pgfqpoint{3.926336in}{0.431537in}}%
\pgfpathlineto{\pgfqpoint{3.928433in}{0.572715in}}%
\pgfpathlineto{\pgfqpoint{3.930530in}{0.599517in}}%
\pgfpathlineto{\pgfqpoint{3.932628in}{0.614105in}}%
\pgfpathlineto{\pgfqpoint{3.934725in}{0.622293in}}%
\pgfpathlineto{\pgfqpoint{3.936822in}{0.625263in}}%
\pgfpathlineto{\pgfqpoint{3.938920in}{0.622056in}}%
\pgfpathlineto{\pgfqpoint{3.941017in}{0.613552in}}%
\pgfpathlineto{\pgfqpoint{3.943114in}{0.601157in}}%
\pgfpathlineto{\pgfqpoint{3.945212in}{0.582651in}}%
\pgfpathlineto{\pgfqpoint{3.947309in}{0.550041in}}%
\pgfpathlineto{\pgfqpoint{3.949407in}{0.479337in}}%
\pgfpathlineto{\pgfqpoint{3.951504in}{0.549614in}}%
\pgfpathlineto{\pgfqpoint{3.953601in}{0.564459in}}%
\pgfpathlineto{\pgfqpoint{3.955699in}{0.567049in}}%
\pgfpathlineto{\pgfqpoint{3.957796in}{0.560256in}}%
\pgfpathlineto{\pgfqpoint{3.959893in}{0.537879in}}%
\pgfpathlineto{\pgfqpoint{3.961991in}{0.474636in}}%
\pgfpathlineto{\pgfqpoint{3.964088in}{0.555700in}}%
\pgfpathlineto{\pgfqpoint{3.966185in}{0.580366in}}%
\pgfpathlineto{\pgfqpoint{3.968283in}{0.594056in}}%
\pgfpathlineto{\pgfqpoint{3.970380in}{0.601618in}}%
\pgfpathlineto{\pgfqpoint{3.972478in}{0.603960in}}%
\pgfpathlineto{\pgfqpoint{3.974575in}{0.599895in}}%
\pgfpathlineto{\pgfqpoint{3.976672in}{0.588877in}}%
\pgfpathlineto{\pgfqpoint{3.978770in}{0.567817in}}%
\pgfpathlineto{\pgfqpoint{3.980867in}{0.505158in}}%
\pgfpathlineto{\pgfqpoint{3.982964in}{0.548686in}}%
\pgfpathlineto{\pgfqpoint{3.985062in}{0.577479in}}%
\pgfpathlineto{\pgfqpoint{3.987159in}{0.590675in}}%
\pgfpathlineto{\pgfqpoint{3.989256in}{0.597347in}}%
\pgfpathlineto{\pgfqpoint{3.991354in}{0.599743in}}%
\pgfpathlineto{\pgfqpoint{3.993451in}{0.598443in}}%
\pgfpathlineto{\pgfqpoint{3.995549in}{0.593065in}}%
\pgfpathlineto{\pgfqpoint{3.997646in}{0.581813in}}%
\pgfpathlineto{\pgfqpoint{3.999743in}{0.557953in}}%
\pgfpathlineto{\pgfqpoint{4.001841in}{0.452835in}}%
\pgfpathlineto{\pgfqpoint{4.003938in}{0.565201in}}%
\pgfpathlineto{\pgfqpoint{4.006035in}{0.590175in}}%
\pgfpathlineto{\pgfqpoint{4.008133in}{0.603760in}}%
\pgfpathlineto{\pgfqpoint{4.010230in}{0.611539in}}%
\pgfpathlineto{\pgfqpoint{4.012327in}{0.614830in}}%
\pgfpathlineto{\pgfqpoint{4.014425in}{0.613277in}}%
\pgfpathlineto{\pgfqpoint{4.016522in}{0.604580in}}%
\pgfpathlineto{\pgfqpoint{4.018620in}{0.583364in}}%
\pgfpathlineto{\pgfqpoint{4.020717in}{0.480708in}}%
\pgfpathlineto{\pgfqpoint{4.022814in}{0.580776in}}%
\pgfpathlineto{\pgfqpoint{4.024912in}{0.607319in}}%
\pgfpathlineto{\pgfqpoint{4.027009in}{0.621148in}}%
\pgfpathlineto{\pgfqpoint{4.029106in}{0.629286in}}%
\pgfpathlineto{\pgfqpoint{4.031204in}{0.633787in}}%
\pgfpathlineto{\pgfqpoint{4.033301in}{0.635416in}}%
\pgfpathlineto{\pgfqpoint{4.035398in}{0.634363in}}%
\pgfpathlineto{\pgfqpoint{4.037496in}{0.630360in}}%
\pgfpathlineto{\pgfqpoint{4.039593in}{0.622466in}}%
\pgfpathlineto{\pgfqpoint{4.041691in}{0.608039in}}%
\pgfpathlineto{\pgfqpoint{4.045885in}{0.550192in}}%
\pgfpathlineto{\pgfqpoint{4.047983in}{0.603510in}}%
\pgfpathlineto{\pgfqpoint{4.050080in}{0.624257in}}%
\pgfpathlineto{\pgfqpoint{4.052177in}{0.636270in}}%
\pgfpathlineto{\pgfqpoint{4.054275in}{0.643411in}}%
\pgfpathlineto{\pgfqpoint{4.056372in}{0.646719in}}%
\pgfpathlineto{\pgfqpoint{4.058469in}{0.646023in}}%
\pgfpathlineto{\pgfqpoint{4.060567in}{0.640160in}}%
\pgfpathlineto{\pgfqpoint{4.062664in}{0.630347in}}%
\pgfpathlineto{\pgfqpoint{4.064762in}{0.615973in}}%
\pgfpathlineto{\pgfqpoint{4.066859in}{0.593485in}}%
\pgfpathlineto{\pgfqpoint{4.068956in}{0.544883in}}%
\pgfpathlineto{\pgfqpoint{4.071054in}{0.549267in}}%
\pgfpathlineto{\pgfqpoint{4.073151in}{0.580381in}}%
\pgfpathlineto{\pgfqpoint{4.075248in}{0.589523in}}%
\pgfpathlineto{\pgfqpoint{4.077346in}{0.588964in}}%
\pgfpathlineto{\pgfqpoint{4.079443in}{0.578746in}}%
\pgfpathlineto{\pgfqpoint{4.081540in}{0.547112in}}%
\pgfpathlineto{\pgfqpoint{4.083638in}{0.538056in}}%
\pgfpathlineto{\pgfqpoint{4.085735in}{0.586175in}}%
\pgfpathlineto{\pgfqpoint{4.087832in}{0.606738in}}%
\pgfpathlineto{\pgfqpoint{4.089930in}{0.618445in}}%
\pgfpathlineto{\pgfqpoint{4.092027in}{0.624515in}}%
\pgfpathlineto{\pgfqpoint{4.094125in}{0.625261in}}%
\pgfpathlineto{\pgfqpoint{4.096222in}{0.619466in}}%
\pgfpathlineto{\pgfqpoint{4.098319in}{0.609494in}}%
\pgfpathlineto{\pgfqpoint{4.100417in}{0.594878in}}%
\pgfpathlineto{\pgfqpoint{4.102514in}{0.571878in}}%
\pgfpathlineto{\pgfqpoint{4.104611in}{0.520531in}}%
\pgfpathlineto{\pgfqpoint{4.106709in}{0.531343in}}%
\pgfpathlineto{\pgfqpoint{4.108806in}{0.560679in}}%
\pgfpathlineto{\pgfqpoint{4.110903in}{0.569320in}}%
\pgfpathlineto{\pgfqpoint{4.113001in}{0.568398in}}%
\pgfpathlineto{\pgfqpoint{4.115098in}{0.557671in}}%
\pgfpathlineto{\pgfqpoint{4.117196in}{0.524169in}}%
\pgfpathlineto{\pgfqpoint{4.119293in}{0.521431in}}%
\pgfpathlineto{\pgfqpoint{4.121390in}{0.566971in}}%
\pgfpathlineto{\pgfqpoint{4.123488in}{0.587014in}}%
\pgfpathlineto{\pgfqpoint{4.125585in}{0.598450in}}%
\pgfpathlineto{\pgfqpoint{4.127682in}{0.604303in}}%
\pgfpathlineto{\pgfqpoint{4.129780in}{0.604804in}}%
\pgfpathlineto{\pgfqpoint{4.131877in}{0.598430in}}%
\pgfpathlineto{\pgfqpoint{4.133974in}{0.585039in}}%
\pgfpathlineto{\pgfqpoint{4.136072in}{0.558504in}}%
\pgfpathlineto{\pgfqpoint{4.138169in}{0.457942in}}%
\pgfpathlineto{\pgfqpoint{4.140267in}{0.561053in}}%
\pgfpathlineto{\pgfqpoint{4.142364in}{0.583017in}}%
\pgfpathlineto{\pgfqpoint{4.144461in}{0.593770in}}%
\pgfpathlineto{\pgfqpoint{4.146559in}{0.598942in}}%
\pgfpathlineto{\pgfqpoint{4.148656in}{0.600093in}}%
\pgfpathlineto{\pgfqpoint{4.150753in}{0.597477in}}%
\pgfpathlineto{\pgfqpoint{4.152851in}{0.590332in}}%
\pgfpathlineto{\pgfqpoint{4.154948in}{0.575853in}}%
\pgfpathlineto{\pgfqpoint{4.157045in}{0.541083in}}%
\pgfpathlineto{\pgfqpoint{4.159143in}{0.530869in}}%
\pgfpathlineto{\pgfqpoint{4.161240in}{0.577152in}}%
\pgfpathlineto{\pgfqpoint{4.163338in}{0.596663in}}%
\pgfpathlineto{\pgfqpoint{4.165435in}{0.607810in}}%
\pgfpathlineto{\pgfqpoint{4.167532in}{0.613869in}}%
\pgfpathlineto{\pgfqpoint{4.169630in}{0.615462in}}%
\pgfpathlineto{\pgfqpoint{4.171727in}{0.611570in}}%
\pgfpathlineto{\pgfqpoint{4.173824in}{0.599290in}}%
\pgfpathlineto{\pgfqpoint{4.178019in}{0.534407in}}%
\pgfpathlineto{\pgfqpoint{4.180116in}{0.591779in}}%
\pgfpathlineto{\pgfqpoint{4.182214in}{0.612262in}}%
\pgfpathlineto{\pgfqpoint{4.184311in}{0.623634in}}%
\pgfpathlineto{\pgfqpoint{4.186409in}{0.630244in}}%
\pgfpathlineto{\pgfqpoint{4.188506in}{0.633541in}}%
\pgfpathlineto{\pgfqpoint{4.190603in}{0.634040in}}%
\pgfpathlineto{\pgfqpoint{4.192701in}{0.631742in}}%
\pgfpathlineto{\pgfqpoint{4.194798in}{0.626126in}}%
\pgfpathlineto{\pgfqpoint{4.196895in}{0.615674in}}%
\pgfpathlineto{\pgfqpoint{4.198993in}{0.595622in}}%
\pgfpathlineto{\pgfqpoint{4.201090in}{0.529479in}}%
\pgfpathlineto{\pgfqpoint{4.203187in}{0.583836in}}%
\pgfpathlineto{\pgfqpoint{4.205285in}{0.614226in}}%
\pgfpathlineto{\pgfqpoint{4.207382in}{0.629902in}}%
\pgfpathlineto{\pgfqpoint{4.209480in}{0.639298in}}%
\pgfpathlineto{\pgfqpoint{4.211577in}{0.644480in}}%
\pgfpathlineto{\pgfqpoint{4.213674in}{0.645825in}}%
\pgfpathlineto{\pgfqpoint{4.215772in}{0.642515in}}%
\pgfpathlineto{\pgfqpoint{4.217869in}{0.634327in}}%
\pgfpathlineto{\pgfqpoint{4.219966in}{0.622170in}}%
\pgfpathlineto{\pgfqpoint{4.222064in}{0.603795in}}%
\pgfpathlineto{\pgfqpoint{4.224161in}{0.570867in}}%
\pgfpathlineto{\pgfqpoint{4.226258in}{0.507966in}}%
\pgfpathlineto{\pgfqpoint{4.228356in}{0.573447in}}%
\pgfpathlineto{\pgfqpoint{4.230453in}{0.588567in}}%
\pgfpathlineto{\pgfqpoint{4.232551in}{0.592058in}}%
\pgfpathlineto{\pgfqpoint{4.234648in}{0.587164in}}%
\pgfpathlineto{\pgfqpoint{4.236745in}{0.570420in}}%
\pgfpathlineto{\pgfqpoint{4.238843in}{0.499428in}}%
\pgfpathlineto{\pgfqpoint{4.240940in}{0.566417in}}%
\pgfpathlineto{\pgfqpoint{4.243037in}{0.596977in}}%
\pgfpathlineto{\pgfqpoint{4.245135in}{0.612886in}}%
\pgfpathlineto{\pgfqpoint{4.247232in}{0.621914in}}%
\pgfpathlineto{\pgfqpoint{4.249329in}{0.625699in}}%
\pgfpathlineto{\pgfqpoint{4.251427in}{0.623660in}}%
\pgfpathlineto{\pgfqpoint{4.253524in}{0.615676in}}%
\pgfpathlineto{\pgfqpoint{4.255621in}{0.603765in}}%
\pgfpathlineto{\pgfqpoint{4.257719in}{0.585826in}}%
\pgfpathlineto{\pgfqpoint{4.259816in}{0.554214in}}%
\pgfpathlineto{\pgfqpoint{4.261914in}{0.475555in}}%
\pgfpathlineto{\pgfqpoint{4.264011in}{0.552916in}}%
\pgfpathlineto{\pgfqpoint{4.266108in}{0.568955in}}%
\pgfpathlineto{\pgfqpoint{4.268206in}{0.572913in}}%
\pgfpathlineto{\pgfqpoint{4.270303in}{0.568482in}}%
\pgfpathlineto{\pgfqpoint{4.272400in}{0.552630in}}%
\pgfpathlineto{\pgfqpoint{4.274498in}{0.491149in}}%
\pgfpathlineto{\pgfqpoint{4.276595in}{0.544720in}}%
\pgfpathlineto{\pgfqpoint{4.278692in}{0.576688in}}%
\pgfpathlineto{\pgfqpoint{4.280790in}{0.593074in}}%
\pgfpathlineto{\pgfqpoint{4.282887in}{0.602401in}}%
\pgfpathlineto{\pgfqpoint{4.284985in}{0.606459in}}%
\pgfpathlineto{\pgfqpoint{4.287082in}{0.604782in}}%
\pgfpathlineto{\pgfqpoint{4.291277in}{0.589258in}}%
\pgfpathlineto{\pgfqpoint{4.295471in}{0.568291in}}%
\pgfpathlineto{\pgfqpoint{4.301763in}{0.535032in}}%
\pgfpathlineto{\pgfqpoint{4.303861in}{0.530765in}}%
\pgfpathlineto{\pgfqpoint{4.305958in}{0.532543in}}%
\pgfpathlineto{\pgfqpoint{4.312250in}{0.549750in}}%
\pgfpathlineto{\pgfqpoint{4.314348in}{0.551823in}}%
\pgfpathlineto{\pgfqpoint{4.316445in}{0.548867in}}%
\pgfpathlineto{\pgfqpoint{4.316445in}{0.548867in}}%
\pgfusepath{stroke}%
\end{pgfscope}%
\begin{pgfscope}%
\pgfpathrectangle{\pgfqpoint{2.554660in}{0.320625in}}{\pgfqpoint{1.845679in}{1.164375in}}%
\pgfusepath{clip}%
\pgfsetrectcap%
\pgfsetroundjoin%
\pgfsetlinewidth{1.154312pt}%
\definecolor{currentstroke}{rgb}{0.569172,0.167454,0.504105}%
\pgfsetstrokecolor{currentstroke}%
\pgfsetdash{}{0pt}%
\pgfpathmoveto{\pgfqpoint{2.638555in}{1.468525in}}%
\pgfpathlineto{\pgfqpoint{2.651139in}{1.446163in}}%
\pgfpathlineto{\pgfqpoint{2.661626in}{1.424759in}}%
\pgfpathlineto{\pgfqpoint{2.676308in}{1.391774in}}%
\pgfpathlineto{\pgfqpoint{2.703573in}{1.330233in}}%
\pgfpathlineto{\pgfqpoint{2.720352in}{1.295284in}}%
\pgfpathlineto{\pgfqpoint{2.735034in}{1.267142in}}%
\pgfpathlineto{\pgfqpoint{2.749715in}{1.241403in}}%
\pgfpathlineto{\pgfqpoint{2.764397in}{1.218030in}}%
\pgfpathlineto{\pgfqpoint{2.779078in}{1.196792in}}%
\pgfpathlineto{\pgfqpoint{2.795857in}{1.174617in}}%
\pgfpathlineto{\pgfqpoint{2.814733in}{1.151543in}}%
\pgfpathlineto{\pgfqpoint{2.839902in}{1.122588in}}%
\pgfpathlineto{\pgfqpoint{2.879752in}{1.078641in}}%
\pgfpathlineto{\pgfqpoint{2.965744in}{0.985894in}}%
\pgfpathlineto{\pgfqpoint{3.227914in}{0.705023in}}%
\pgfpathlineto{\pgfqpoint{3.244693in}{0.687946in}}%
\pgfpathlineto{\pgfqpoint{3.259374in}{0.670785in}}%
\pgfpathlineto{\pgfqpoint{3.274056in}{0.653789in}}%
\pgfpathlineto{\pgfqpoint{3.282445in}{0.645631in}}%
\pgfpathlineto{\pgfqpoint{3.290835in}{0.639153in}}%
\pgfpathlineto{\pgfqpoint{3.303419in}{0.630452in}}%
\pgfpathlineto{\pgfqpoint{3.307614in}{0.626407in}}%
\pgfpathlineto{\pgfqpoint{3.313906in}{0.618816in}}%
\pgfpathlineto{\pgfqpoint{3.320198in}{0.609823in}}%
\pgfpathlineto{\pgfqpoint{3.328587in}{0.595676in}}%
\pgfpathlineto{\pgfqpoint{3.347464in}{0.560779in}}%
\pgfpathlineto{\pgfqpoint{3.351658in}{0.555886in}}%
\pgfpathlineto{\pgfqpoint{3.353756in}{0.554424in}}%
\pgfpathlineto{\pgfqpoint{3.355853in}{0.553695in}}%
\pgfpathlineto{\pgfqpoint{3.357950in}{0.553698in}}%
\pgfpathlineto{\pgfqpoint{3.362145in}{0.555651in}}%
\pgfpathlineto{\pgfqpoint{3.366340in}{0.559429in}}%
\pgfpathlineto{\pgfqpoint{3.378924in}{0.572177in}}%
\pgfpathlineto{\pgfqpoint{3.383119in}{0.574692in}}%
\pgfpathlineto{\pgfqpoint{3.387313in}{0.575513in}}%
\pgfpathlineto{\pgfqpoint{3.389411in}{0.575091in}}%
\pgfpathlineto{\pgfqpoint{3.391508in}{0.573957in}}%
\pgfpathlineto{\pgfqpoint{3.397800in}{0.568100in}}%
\pgfpathlineto{\pgfqpoint{3.399898in}{0.569539in}}%
\pgfpathlineto{\pgfqpoint{3.401995in}{0.574335in}}%
\pgfpathlineto{\pgfqpoint{3.408287in}{0.593587in}}%
\pgfpathlineto{\pgfqpoint{3.412482in}{0.602420in}}%
\pgfpathlineto{\pgfqpoint{3.414579in}{0.605299in}}%
\pgfpathlineto{\pgfqpoint{3.416676in}{0.607208in}}%
\pgfpathlineto{\pgfqpoint{3.418774in}{0.608168in}}%
\pgfpathlineto{\pgfqpoint{3.420871in}{0.608168in}}%
\pgfpathlineto{\pgfqpoint{3.422969in}{0.607159in}}%
\pgfpathlineto{\pgfqpoint{3.425066in}{0.605031in}}%
\pgfpathlineto{\pgfqpoint{3.427163in}{0.601593in}}%
\pgfpathlineto{\pgfqpoint{3.429261in}{0.596506in}}%
\pgfpathlineto{\pgfqpoint{3.431358in}{0.589153in}}%
\pgfpathlineto{\pgfqpoint{3.433455in}{0.578295in}}%
\pgfpathlineto{\pgfqpoint{3.435553in}{0.560991in}}%
\pgfpathlineto{\pgfqpoint{3.437650in}{0.530521in}}%
\pgfpathlineto{\pgfqpoint{3.439747in}{0.539227in}}%
\pgfpathlineto{\pgfqpoint{3.441845in}{0.568211in}}%
\pgfpathlineto{\pgfqpoint{3.443942in}{0.585393in}}%
\pgfpathlineto{\pgfqpoint{3.446040in}{0.596948in}}%
\pgfpathlineto{\pgfqpoint{3.450234in}{0.611315in}}%
\pgfpathlineto{\pgfqpoint{3.452332in}{0.615612in}}%
\pgfpathlineto{\pgfqpoint{3.454429in}{0.618345in}}%
\pgfpathlineto{\pgfqpoint{3.456526in}{0.619534in}}%
\pgfpathlineto{\pgfqpoint{3.458624in}{0.619017in}}%
\pgfpathlineto{\pgfqpoint{3.460721in}{0.616377in}}%
\pgfpathlineto{\pgfqpoint{3.462818in}{0.610700in}}%
\pgfpathlineto{\pgfqpoint{3.464916in}{0.600801in}}%
\pgfpathlineto{\pgfqpoint{3.467013in}{0.584853in}}%
\pgfpathlineto{\pgfqpoint{3.471208in}{0.526193in}}%
\pgfpathlineto{\pgfqpoint{3.473305in}{0.562165in}}%
\pgfpathlineto{\pgfqpoint{3.475403in}{0.578949in}}%
\pgfpathlineto{\pgfqpoint{3.477500in}{0.586350in}}%
\pgfpathlineto{\pgfqpoint{3.479597in}{0.588008in}}%
\pgfpathlineto{\pgfqpoint{3.481695in}{0.584571in}}%
\pgfpathlineto{\pgfqpoint{3.483792in}{0.574694in}}%
\pgfpathlineto{\pgfqpoint{3.485889in}{0.552081in}}%
\pgfpathlineto{\pgfqpoint{3.487987in}{0.489800in}}%
\pgfpathlineto{\pgfqpoint{3.490084in}{0.560751in}}%
\pgfpathlineto{\pgfqpoint{3.492182in}{0.585627in}}%
\pgfpathlineto{\pgfqpoint{3.494279in}{0.599219in}}%
\pgfpathlineto{\pgfqpoint{3.496376in}{0.606817in}}%
\pgfpathlineto{\pgfqpoint{3.498474in}{0.609570in}}%
\pgfpathlineto{\pgfqpoint{3.500571in}{0.606718in}}%
\pgfpathlineto{\pgfqpoint{3.502668in}{0.597937in}}%
\pgfpathlineto{\pgfqpoint{3.504766in}{0.583514in}}%
\pgfpathlineto{\pgfqpoint{3.506863in}{0.558131in}}%
\pgfpathlineto{\pgfqpoint{3.508960in}{0.495764in}}%
\pgfpathlineto{\pgfqpoint{3.511058in}{0.548912in}}%
\pgfpathlineto{\pgfqpoint{3.513155in}{0.570758in}}%
\pgfpathlineto{\pgfqpoint{3.515253in}{0.580195in}}%
\pgfpathlineto{\pgfqpoint{3.517350in}{0.583254in}}%
\pgfpathlineto{\pgfqpoint{3.519447in}{0.581221in}}%
\pgfpathlineto{\pgfqpoint{3.521545in}{0.573360in}}%
\pgfpathlineto{\pgfqpoint{3.523642in}{0.555448in}}%
\pgfpathlineto{\pgfqpoint{3.525739in}{0.495994in}}%
\pgfpathlineto{\pgfqpoint{3.527837in}{0.543857in}}%
\pgfpathlineto{\pgfqpoint{3.529934in}{0.575244in}}%
\pgfpathlineto{\pgfqpoint{3.532031in}{0.591153in}}%
\pgfpathlineto{\pgfqpoint{3.534129in}{0.600209in}}%
\pgfpathlineto{\pgfqpoint{3.536226in}{0.604287in}}%
\pgfpathlineto{\pgfqpoint{3.538324in}{0.603157in}}%
\pgfpathlineto{\pgfqpoint{3.540421in}{0.594728in}}%
\pgfpathlineto{\pgfqpoint{3.542518in}{0.574989in}}%
\pgfpathlineto{\pgfqpoint{3.544616in}{0.508210in}}%
\pgfpathlineto{\pgfqpoint{3.546713in}{0.564050in}}%
\pgfpathlineto{\pgfqpoint{3.548810in}{0.592693in}}%
\pgfpathlineto{\pgfqpoint{3.550908in}{0.607074in}}%
\pgfpathlineto{\pgfqpoint{3.553005in}{0.615392in}}%
\pgfpathlineto{\pgfqpoint{3.555102in}{0.619908in}}%
\pgfpathlineto{\pgfqpoint{3.557200in}{0.621445in}}%
\pgfpathlineto{\pgfqpoint{3.559297in}{0.620190in}}%
\pgfpathlineto{\pgfqpoint{3.561395in}{0.615833in}}%
\pgfpathlineto{\pgfqpoint{3.563492in}{0.607293in}}%
\pgfpathlineto{\pgfqpoint{3.565589in}{0.591426in}}%
\pgfpathlineto{\pgfqpoint{3.567687in}{0.553516in}}%
\pgfpathlineto{\pgfqpoint{3.569784in}{0.551646in}}%
\pgfpathlineto{\pgfqpoint{3.571881in}{0.594371in}}%
\pgfpathlineto{\pgfqpoint{3.573979in}{0.613232in}}%
\pgfpathlineto{\pgfqpoint{3.576076in}{0.624311in}}%
\pgfpathlineto{\pgfqpoint{3.578173in}{0.630742in}}%
\pgfpathlineto{\pgfqpoint{3.580271in}{0.633311in}}%
\pgfpathlineto{\pgfqpoint{3.582368in}{0.631597in}}%
\pgfpathlineto{\pgfqpoint{3.584465in}{0.624433in}}%
\pgfpathlineto{\pgfqpoint{3.586563in}{0.612273in}}%
\pgfpathlineto{\pgfqpoint{3.588660in}{0.592042in}}%
\pgfpathlineto{\pgfqpoint{3.590758in}{0.547077in}}%
\pgfpathlineto{\pgfqpoint{3.592855in}{0.550844in}}%
\pgfpathlineto{\pgfqpoint{3.594952in}{0.585256in}}%
\pgfpathlineto{\pgfqpoint{3.597050in}{0.598239in}}%
\pgfpathlineto{\pgfqpoint{3.599147in}{0.603189in}}%
\pgfpathlineto{\pgfqpoint{3.601244in}{0.602590in}}%
\pgfpathlineto{\pgfqpoint{3.603342in}{0.596282in}}%
\pgfpathlineto{\pgfqpoint{3.605439in}{0.581128in}}%
\pgfpathlineto{\pgfqpoint{3.607536in}{0.538171in}}%
\pgfpathlineto{\pgfqpoint{3.609634in}{0.554601in}}%
\pgfpathlineto{\pgfqpoint{3.611731in}{0.592494in}}%
\pgfpathlineto{\pgfqpoint{3.613829in}{0.610280in}}%
\pgfpathlineto{\pgfqpoint{3.615926in}{0.620400in}}%
\pgfpathlineto{\pgfqpoint{3.618023in}{0.625336in}}%
\pgfpathlineto{\pgfqpoint{3.620121in}{0.625162in}}%
\pgfpathlineto{\pgfqpoint{3.622218in}{0.618646in}}%
\pgfpathlineto{\pgfqpoint{3.624315in}{0.607040in}}%
\pgfpathlineto{\pgfqpoint{3.626413in}{0.587882in}}%
\pgfpathlineto{\pgfqpoint{3.628510in}{0.547538in}}%
\pgfpathlineto{\pgfqpoint{3.630607in}{0.537067in}}%
\pgfpathlineto{\pgfqpoint{3.632705in}{0.576621in}}%
\pgfpathlineto{\pgfqpoint{3.634802in}{0.590749in}}%
\pgfpathlineto{\pgfqpoint{3.636900in}{0.596321in}}%
\pgfpathlineto{\pgfqpoint{3.638997in}{0.596248in}}%
\pgfpathlineto{\pgfqpoint{3.641094in}{0.590580in}}%
\pgfpathlineto{\pgfqpoint{3.643192in}{0.576633in}}%
\pgfpathlineto{\pgfqpoint{3.645289in}{0.539475in}}%
\pgfpathlineto{\pgfqpoint{3.647386in}{0.540834in}}%
\pgfpathlineto{\pgfqpoint{3.649484in}{0.583367in}}%
\pgfpathlineto{\pgfqpoint{3.651581in}{0.602259in}}%
\pgfpathlineto{\pgfqpoint{3.653678in}{0.612978in}}%
\pgfpathlineto{\pgfqpoint{3.655776in}{0.618399in}}%
\pgfpathlineto{\pgfqpoint{3.657873in}{0.618785in}}%
\pgfpathlineto{\pgfqpoint{3.659971in}{0.612725in}}%
\pgfpathlineto{\pgfqpoint{3.662068in}{0.600077in}}%
\pgfpathlineto{\pgfqpoint{3.664165in}{0.576002in}}%
\pgfpathlineto{\pgfqpoint{3.666263in}{0.476636in}}%
\pgfpathlineto{\pgfqpoint{3.668360in}{0.568220in}}%
\pgfpathlineto{\pgfqpoint{3.670457in}{0.592659in}}%
\pgfpathlineto{\pgfqpoint{3.672555in}{0.604177in}}%
\pgfpathlineto{\pgfqpoint{3.674652in}{0.609675in}}%
\pgfpathlineto{\pgfqpoint{3.676749in}{0.610955in}}%
\pgfpathlineto{\pgfqpoint{3.678847in}{0.608331in}}%
\pgfpathlineto{\pgfqpoint{3.680944in}{0.601016in}}%
\pgfpathlineto{\pgfqpoint{3.683042in}{0.586007in}}%
\pgfpathlineto{\pgfqpoint{3.685139in}{0.548582in}}%
\pgfpathlineto{\pgfqpoint{3.687236in}{0.547554in}}%
\pgfpathlineto{\pgfqpoint{3.689334in}{0.590329in}}%
\pgfpathlineto{\pgfqpoint{3.691431in}{0.609143in}}%
\pgfpathlineto{\pgfqpoint{3.693528in}{0.619923in}}%
\pgfpathlineto{\pgfqpoint{3.695626in}{0.625672in}}%
\pgfpathlineto{\pgfqpoint{3.697723in}{0.626886in}}%
\pgfpathlineto{\pgfqpoint{3.699820in}{0.622410in}}%
\pgfpathlineto{\pgfqpoint{3.701918in}{0.611064in}}%
\pgfpathlineto{\pgfqpoint{3.704015in}{0.588694in}}%
\pgfpathlineto{\pgfqpoint{3.706113in}{0.507807in}}%
\pgfpathlineto{\pgfqpoint{3.708210in}{0.577895in}}%
\pgfpathlineto{\pgfqpoint{3.710307in}{0.604273in}}%
\pgfpathlineto{\pgfqpoint{3.712405in}{0.616921in}}%
\pgfpathlineto{\pgfqpoint{3.714502in}{0.623480in}}%
\pgfpathlineto{\pgfqpoint{3.716599in}{0.625980in}}%
\pgfpathlineto{\pgfqpoint{3.718697in}{0.624975in}}%
\pgfpathlineto{\pgfqpoint{3.720794in}{0.620171in}}%
\pgfpathlineto{\pgfqpoint{3.722891in}{0.610095in}}%
\pgfpathlineto{\pgfqpoint{3.724989in}{0.589655in}}%
\pgfpathlineto{\pgfqpoint{3.727086in}{0.512190in}}%
\pgfpathlineto{\pgfqpoint{3.729184in}{0.582891in}}%
\pgfpathlineto{\pgfqpoint{3.731281in}{0.611851in}}%
\pgfpathlineto{\pgfqpoint{3.733378in}{0.626932in}}%
\pgfpathlineto{\pgfqpoint{3.735476in}{0.635713in}}%
\pgfpathlineto{\pgfqpoint{3.737573in}{0.639985in}}%
\pgfpathlineto{\pgfqpoint{3.739670in}{0.639763in}}%
\pgfpathlineto{\pgfqpoint{3.741768in}{0.633699in}}%
\pgfpathlineto{\pgfqpoint{3.743865in}{0.621658in}}%
\pgfpathlineto{\pgfqpoint{3.745962in}{0.599893in}}%
\pgfpathlineto{\pgfqpoint{3.748060in}{0.538165in}}%
\pgfpathlineto{\pgfqpoint{3.750157in}{0.578173in}}%
\pgfpathlineto{\pgfqpoint{3.752254in}{0.606550in}}%
\pgfpathlineto{\pgfqpoint{3.754352in}{0.618980in}}%
\pgfpathlineto{\pgfqpoint{3.756449in}{0.624664in}}%
\pgfpathlineto{\pgfqpoint{3.758547in}{0.625749in}}%
\pgfpathlineto{\pgfqpoint{3.760644in}{0.622556in}}%
\pgfpathlineto{\pgfqpoint{3.762741in}{0.614032in}}%
\pgfpathlineto{\pgfqpoint{3.764839in}{0.595973in}}%
\pgfpathlineto{\pgfqpoint{3.766936in}{0.538923in}}%
\pgfpathlineto{\pgfqpoint{3.771131in}{0.613135in}}%
\pgfpathlineto{\pgfqpoint{3.773228in}{0.629271in}}%
\pgfpathlineto{\pgfqpoint{3.775325in}{0.638530in}}%
\pgfpathlineto{\pgfqpoint{3.777423in}{0.642950in}}%
\pgfpathlineto{\pgfqpoint{3.779520in}{0.642514in}}%
\pgfpathlineto{\pgfqpoint{3.781618in}{0.636046in}}%
\pgfpathlineto{\pgfqpoint{3.783715in}{0.624930in}}%
\pgfpathlineto{\pgfqpoint{3.785812in}{0.607328in}}%
\pgfpathlineto{\pgfqpoint{3.787910in}{0.574252in}}%
\pgfpathlineto{\pgfqpoint{3.790007in}{0.524672in}}%
\pgfpathlineto{\pgfqpoint{3.792104in}{0.583622in}}%
\pgfpathlineto{\pgfqpoint{3.794202in}{0.599981in}}%
\pgfpathlineto{\pgfqpoint{3.796299in}{0.605994in}}%
\pgfpathlineto{\pgfqpoint{3.798396in}{0.605555in}}%
\pgfpathlineto{\pgfqpoint{3.800494in}{0.598611in}}%
\pgfpathlineto{\pgfqpoint{3.802591in}{0.581047in}}%
\pgfpathlineto{\pgfqpoint{3.804689in}{0.517022in}}%
\pgfpathlineto{\pgfqpoint{3.806786in}{0.571987in}}%
\pgfpathlineto{\pgfqpoint{3.808883in}{0.603124in}}%
\pgfpathlineto{\pgfqpoint{3.810981in}{0.619008in}}%
\pgfpathlineto{\pgfqpoint{3.813078in}{0.628001in}}%
\pgfpathlineto{\pgfqpoint{3.815175in}{0.631882in}}%
\pgfpathlineto{\pgfqpoint{3.817273in}{0.630252in}}%
\pgfpathlineto{\pgfqpoint{3.819370in}{0.622467in}}%
\pgfpathlineto{\pgfqpoint{3.821467in}{0.610139in}}%
\pgfpathlineto{\pgfqpoint{3.823565in}{0.590467in}}%
\pgfpathlineto{\pgfqpoint{3.825662in}{0.550396in}}%
\pgfpathlineto{\pgfqpoint{3.827760in}{0.534936in}}%
\pgfpathlineto{\pgfqpoint{3.829857in}{0.575426in}}%
\pgfpathlineto{\pgfqpoint{3.831954in}{0.588758in}}%
\pgfpathlineto{\pgfqpoint{3.834052in}{0.592977in}}%
\pgfpathlineto{\pgfqpoint{3.836149in}{0.590742in}}%
\pgfpathlineto{\pgfqpoint{3.838246in}{0.581078in}}%
\pgfpathlineto{\pgfqpoint{3.840344in}{0.556439in}}%
\pgfpathlineto{\pgfqpoint{3.842441in}{0.497964in}}%
\pgfpathlineto{\pgfqpoint{3.844538in}{0.573700in}}%
\pgfpathlineto{\pgfqpoint{3.846636in}{0.597366in}}%
\pgfpathlineto{\pgfqpoint{3.848733in}{0.610416in}}%
\pgfpathlineto{\pgfqpoint{3.850831in}{0.617528in}}%
\pgfpathlineto{\pgfqpoint{3.852928in}{0.619577in}}%
\pgfpathlineto{\pgfqpoint{3.855025in}{0.615442in}}%
\pgfpathlineto{\pgfqpoint{3.857123in}{0.605215in}}%
\pgfpathlineto{\pgfqpoint{3.859220in}{0.587425in}}%
\pgfpathlineto{\pgfqpoint{3.861317in}{0.549389in}}%
\pgfpathlineto{\pgfqpoint{3.863415in}{0.534500in}}%
\pgfpathlineto{\pgfqpoint{3.865512in}{0.577448in}}%
\pgfpathlineto{\pgfqpoint{3.867609in}{0.593203in}}%
\pgfpathlineto{\pgfqpoint{3.869707in}{0.600487in}}%
\pgfpathlineto{\pgfqpoint{3.871804in}{0.602634in}}%
\pgfpathlineto{\pgfqpoint{3.873902in}{0.600355in}}%
\pgfpathlineto{\pgfqpoint{3.875999in}{0.592901in}}%
\pgfpathlineto{\pgfqpoint{3.878096in}{0.576831in}}%
\pgfpathlineto{\pgfqpoint{3.880194in}{0.532305in}}%
\pgfpathlineto{\pgfqpoint{3.882291in}{0.550672in}}%
\pgfpathlineto{\pgfqpoint{3.884388in}{0.587738in}}%
\pgfpathlineto{\pgfqpoint{3.886486in}{0.605226in}}%
\pgfpathlineto{\pgfqpoint{3.888583in}{0.615233in}}%
\pgfpathlineto{\pgfqpoint{3.890680in}{0.620231in}}%
\pgfpathlineto{\pgfqpoint{3.892778in}{0.620395in}}%
\pgfpathlineto{\pgfqpoint{3.894875in}{0.614250in}}%
\pgfpathlineto{\pgfqpoint{3.896973in}{0.600716in}}%
\pgfpathlineto{\pgfqpoint{3.899070in}{0.572608in}}%
\pgfpathlineto{\pgfqpoint{3.901167in}{0.508110in}}%
\pgfpathlineto{\pgfqpoint{3.903265in}{0.582511in}}%
\pgfpathlineto{\pgfqpoint{3.905362in}{0.603428in}}%
\pgfpathlineto{\pgfqpoint{3.907459in}{0.614087in}}%
\pgfpathlineto{\pgfqpoint{3.909557in}{0.619492in}}%
\pgfpathlineto{\pgfqpoint{3.911654in}{0.621121in}}%
\pgfpathlineto{\pgfqpoint{3.913751in}{0.619297in}}%
\pgfpathlineto{\pgfqpoint{3.915849in}{0.613522in}}%
\pgfpathlineto{\pgfqpoint{3.917946in}{0.601911in}}%
\pgfpathlineto{\pgfqpoint{3.920043in}{0.577522in}}%
\pgfpathlineto{\pgfqpoint{3.922141in}{0.483184in}}%
\pgfpathlineto{\pgfqpoint{3.924238in}{0.585835in}}%
\pgfpathlineto{\pgfqpoint{3.926336in}{0.610477in}}%
\pgfpathlineto{\pgfqpoint{3.928433in}{0.623969in}}%
\pgfpathlineto{\pgfqpoint{3.930530in}{0.631765in}}%
\pgfpathlineto{\pgfqpoint{3.932628in}{0.635194in}}%
\pgfpathlineto{\pgfqpoint{3.934725in}{0.633987in}}%
\pgfpathlineto{\pgfqpoint{3.936822in}{0.626630in}}%
\pgfpathlineto{\pgfqpoint{3.938920in}{0.612049in}}%
\pgfpathlineto{\pgfqpoint{3.941017in}{0.581952in}}%
\pgfpathlineto{\pgfqpoint{3.943114in}{0.531987in}}%
\pgfpathlineto{\pgfqpoint{3.945212in}{0.595565in}}%
\pgfpathlineto{\pgfqpoint{3.947309in}{0.615181in}}%
\pgfpathlineto{\pgfqpoint{3.949407in}{0.625031in}}%
\pgfpathlineto{\pgfqpoint{3.951504in}{0.629688in}}%
\pgfpathlineto{\pgfqpoint{3.953601in}{0.630468in}}%
\pgfpathlineto{\pgfqpoint{3.955699in}{0.627523in}}%
\pgfpathlineto{\pgfqpoint{3.957796in}{0.619998in}}%
\pgfpathlineto{\pgfqpoint{3.959893in}{0.604876in}}%
\pgfpathlineto{\pgfqpoint{3.961991in}{0.567574in}}%
\pgfpathlineto{\pgfqpoint{3.964088in}{0.565707in}}%
\pgfpathlineto{\pgfqpoint{3.966185in}{0.608708in}}%
\pgfpathlineto{\pgfqpoint{3.968283in}{0.627566in}}%
\pgfpathlineto{\pgfqpoint{3.970380in}{0.638403in}}%
\pgfpathlineto{\pgfqpoint{3.972478in}{0.644259in}}%
\pgfpathlineto{\pgfqpoint{3.974575in}{0.645685in}}%
\pgfpathlineto{\pgfqpoint{3.976672in}{0.641649in}}%
\pgfpathlineto{\pgfqpoint{3.978770in}{0.632200in}}%
\pgfpathlineto{\pgfqpoint{3.980867in}{0.616782in}}%
\pgfpathlineto{\pgfqpoint{3.982964in}{0.588782in}}%
\pgfpathlineto{\pgfqpoint{3.985062in}{0.471721in}}%
\pgfpathlineto{\pgfqpoint{3.987159in}{0.586236in}}%
\pgfpathlineto{\pgfqpoint{3.989256in}{0.606563in}}%
\pgfpathlineto{\pgfqpoint{3.991354in}{0.615093in}}%
\pgfpathlineto{\pgfqpoint{3.993451in}{0.617329in}}%
\pgfpathlineto{\pgfqpoint{3.995549in}{0.614279in}}%
\pgfpathlineto{\pgfqpoint{3.997646in}{0.604739in}}%
\pgfpathlineto{\pgfqpoint{3.999743in}{0.582498in}}%
\pgfpathlineto{\pgfqpoint{4.001841in}{0.452704in}}%
\pgfpathlineto{\pgfqpoint{4.003938in}{0.589146in}}%
\pgfpathlineto{\pgfqpoint{4.006035in}{0.614924in}}%
\pgfpathlineto{\pgfqpoint{4.008133in}{0.628804in}}%
\pgfpathlineto{\pgfqpoint{4.010230in}{0.636542in}}%
\pgfpathlineto{\pgfqpoint{4.012327in}{0.639361in}}%
\pgfpathlineto{\pgfqpoint{4.014425in}{0.636506in}}%
\pgfpathlineto{\pgfqpoint{4.016522in}{0.628096in}}%
\pgfpathlineto{\pgfqpoint{4.018620in}{0.615324in}}%
\pgfpathlineto{\pgfqpoint{4.020717in}{0.595297in}}%
\pgfpathlineto{\pgfqpoint{4.022814in}{0.555225in}}%
\pgfpathlineto{\pgfqpoint{4.024912in}{0.537424in}}%
\pgfpathlineto{\pgfqpoint{4.027009in}{0.578057in}}%
\pgfpathlineto{\pgfqpoint{4.029106in}{0.590728in}}%
\pgfpathlineto{\pgfqpoint{4.031204in}{0.593878in}}%
\pgfpathlineto{\pgfqpoint{4.033301in}{0.589831in}}%
\pgfpathlineto{\pgfqpoint{4.035398in}{0.576294in}}%
\pgfpathlineto{\pgfqpoint{4.037496in}{0.535371in}}%
\pgfpathlineto{\pgfqpoint{4.039593in}{0.549972in}}%
\pgfpathlineto{\pgfqpoint{4.041691in}{0.589096in}}%
\pgfpathlineto{\pgfqpoint{4.043788in}{0.607425in}}%
\pgfpathlineto{\pgfqpoint{4.045885in}{0.617866in}}%
\pgfpathlineto{\pgfqpoint{4.047983in}{0.622941in}}%
\pgfpathlineto{\pgfqpoint{4.050080in}{0.622657in}}%
\pgfpathlineto{\pgfqpoint{4.052177in}{0.615822in}}%
\pgfpathlineto{\pgfqpoint{4.054275in}{0.603903in}}%
\pgfpathlineto{\pgfqpoint{4.056372in}{0.584166in}}%
\pgfpathlineto{\pgfqpoint{4.058469in}{0.541326in}}%
\pgfpathlineto{\pgfqpoint{4.060567in}{0.537994in}}%
\pgfpathlineto{\pgfqpoint{4.062664in}{0.575049in}}%
\pgfpathlineto{\pgfqpoint{4.064762in}{0.588570in}}%
\pgfpathlineto{\pgfqpoint{4.066859in}{0.593776in}}%
\pgfpathlineto{\pgfqpoint{4.068956in}{0.593357in}}%
\pgfpathlineto{\pgfqpoint{4.071054in}{0.587228in}}%
\pgfpathlineto{\pgfqpoint{4.073151in}{0.572363in}}%
\pgfpathlineto{\pgfqpoint{4.075248in}{0.530745in}}%
\pgfpathlineto{\pgfqpoint{4.077346in}{0.544007in}}%
\pgfpathlineto{\pgfqpoint{4.079443in}{0.582805in}}%
\pgfpathlineto{\pgfqpoint{4.081540in}{0.600821in}}%
\pgfpathlineto{\pgfqpoint{4.083638in}{0.611065in}}%
\pgfpathlineto{\pgfqpoint{4.085735in}{0.616096in}}%
\pgfpathlineto{\pgfqpoint{4.087832in}{0.616022in}}%
\pgfpathlineto{\pgfqpoint{4.089930in}{0.609350in}}%
\pgfpathlineto{\pgfqpoint{4.092027in}{0.595558in}}%
\pgfpathlineto{\pgfqpoint{4.094125in}{0.567696in}}%
\pgfpathlineto{\pgfqpoint{4.096222in}{0.494773in}}%
\pgfpathlineto{\pgfqpoint{4.098319in}{0.574996in}}%
\pgfpathlineto{\pgfqpoint{4.100417in}{0.595991in}}%
\pgfpathlineto{\pgfqpoint{4.102514in}{0.606438in}}%
\pgfpathlineto{\pgfqpoint{4.104611in}{0.611496in}}%
\pgfpathlineto{\pgfqpoint{4.106709in}{0.612630in}}%
\pgfpathlineto{\pgfqpoint{4.108806in}{0.610079in}}%
\pgfpathlineto{\pgfqpoint{4.110903in}{0.603113in}}%
\pgfpathlineto{\pgfqpoint{4.113001in}{0.589091in}}%
\pgfpathlineto{\pgfqpoint{4.115098in}{0.556331in}}%
\pgfpathlineto{\pgfqpoint{4.117196in}{0.537995in}}%
\pgfpathlineto{\pgfqpoint{4.119293in}{0.587856in}}%
\pgfpathlineto{\pgfqpoint{4.121390in}{0.608012in}}%
\pgfpathlineto{\pgfqpoint{4.123488in}{0.619497in}}%
\pgfpathlineto{\pgfqpoint{4.125585in}{0.625841in}}%
\pgfpathlineto{\pgfqpoint{4.127682in}{0.627776in}}%
\pgfpathlineto{\pgfqpoint{4.129780in}{0.624444in}}%
\pgfpathlineto{\pgfqpoint{4.131877in}{0.614226in}}%
\pgfpathlineto{\pgfqpoint{4.133974in}{0.593649in}}%
\pgfpathlineto{\pgfqpoint{4.136072in}{0.529586in}}%
\pgfpathlineto{\pgfqpoint{4.138169in}{0.576798in}}%
\pgfpathlineto{\pgfqpoint{4.140267in}{0.605763in}}%
\pgfpathlineto{\pgfqpoint{4.142364in}{0.619460in}}%
\pgfpathlineto{\pgfqpoint{4.144461in}{0.626787in}}%
\pgfpathlineto{\pgfqpoint{4.146559in}{0.630032in}}%
\pgfpathlineto{\pgfqpoint{4.148656in}{0.629895in}}%
\pgfpathlineto{\pgfqpoint{4.150753in}{0.626285in}}%
\pgfpathlineto{\pgfqpoint{4.152851in}{0.618212in}}%
\pgfpathlineto{\pgfqpoint{4.154948in}{0.602491in}}%
\pgfpathlineto{\pgfqpoint{4.157045in}{0.563546in}}%
\pgfpathlineto{\pgfqpoint{4.159143in}{0.565870in}}%
\pgfpathlineto{\pgfqpoint{4.161240in}{0.607283in}}%
\pgfpathlineto{\pgfqpoint{4.163338in}{0.625806in}}%
\pgfpathlineto{\pgfqpoint{4.165435in}{0.636544in}}%
\pgfpathlineto{\pgfqpoint{4.167532in}{0.642456in}}%
\pgfpathlineto{\pgfqpoint{4.169630in}{0.644130in}}%
\pgfpathlineto{\pgfqpoint{4.171727in}{0.640689in}}%
\pgfpathlineto{\pgfqpoint{4.173824in}{0.631344in}}%
\pgfpathlineto{\pgfqpoint{4.175922in}{0.614980in}}%
\pgfpathlineto{\pgfqpoint{4.180116in}{0.545256in}}%
\pgfpathlineto{\pgfqpoint{4.182214in}{0.599136in}}%
\pgfpathlineto{\pgfqpoint{4.184311in}{0.616740in}}%
\pgfpathlineto{\pgfqpoint{4.186409in}{0.624996in}}%
\pgfpathlineto{\pgfqpoint{4.188506in}{0.627929in}}%
\pgfpathlineto{\pgfqpoint{4.190603in}{0.626506in}}%
\pgfpathlineto{\pgfqpoint{4.192701in}{0.620277in}}%
\pgfpathlineto{\pgfqpoint{4.194798in}{0.606694in}}%
\pgfpathlineto{\pgfqpoint{4.196895in}{0.573840in}}%
\pgfpathlineto{\pgfqpoint{4.198993in}{0.557882in}}%
\pgfpathlineto{\pgfqpoint{4.201090in}{0.606949in}}%
\pgfpathlineto{\pgfqpoint{4.203187in}{0.627013in}}%
\pgfpathlineto{\pgfqpoint{4.205285in}{0.638372in}}%
\pgfpathlineto{\pgfqpoint{4.207382in}{0.644447in}}%
\pgfpathlineto{\pgfqpoint{4.209480in}{0.645825in}}%
\pgfpathlineto{\pgfqpoint{4.211577in}{0.641396in}}%
\pgfpathlineto{\pgfqpoint{4.213674in}{0.632093in}}%
\pgfpathlineto{\pgfqpoint{4.215772in}{0.617934in}}%
\pgfpathlineto{\pgfqpoint{4.217869in}{0.594909in}}%
\pgfpathlineto{\pgfqpoint{4.219966in}{0.539049in}}%
\pgfpathlineto{\pgfqpoint{4.224161in}{0.590300in}}%
\pgfpathlineto{\pgfqpoint{4.226258in}{0.600075in}}%
\pgfpathlineto{\pgfqpoint{4.228356in}{0.601614in}}%
\pgfpathlineto{\pgfqpoint{4.230453in}{0.596077in}}%
\pgfpathlineto{\pgfqpoint{4.232551in}{0.580089in}}%
\pgfpathlineto{\pgfqpoint{4.234648in}{0.524959in}}%
\pgfpathlineto{\pgfqpoint{4.238843in}{0.600208in}}%
\pgfpathlineto{\pgfqpoint{4.240940in}{0.616830in}}%
\pgfpathlineto{\pgfqpoint{4.243037in}{0.626277in}}%
\pgfpathlineto{\pgfqpoint{4.245135in}{0.630505in}}%
\pgfpathlineto{\pgfqpoint{4.247232in}{0.629197in}}%
\pgfpathlineto{\pgfqpoint{4.249329in}{0.621647in}}%
\pgfpathlineto{\pgfqpoint{4.251427in}{0.609783in}}%
\pgfpathlineto{\pgfqpoint{4.253524in}{0.591303in}}%
\pgfpathlineto{\pgfqpoint{4.257719in}{0.514923in}}%
\pgfpathlineto{\pgfqpoint{4.259816in}{0.568014in}}%
\pgfpathlineto{\pgfqpoint{4.261914in}{0.582966in}}%
\pgfpathlineto{\pgfqpoint{4.264011in}{0.587616in}}%
\pgfpathlineto{\pgfqpoint{4.266108in}{0.585286in}}%
\pgfpathlineto{\pgfqpoint{4.268206in}{0.574882in}}%
\pgfpathlineto{\pgfqpoint{4.272400in}{0.515506in}}%
\pgfpathlineto{\pgfqpoint{4.274498in}{0.573918in}}%
\pgfpathlineto{\pgfqpoint{4.276595in}{0.595836in}}%
\pgfpathlineto{\pgfqpoint{4.278692in}{0.608113in}}%
\pgfpathlineto{\pgfqpoint{4.280790in}{0.614647in}}%
\pgfpathlineto{\pgfqpoint{4.282887in}{0.616043in}}%
\pgfpathlineto{\pgfqpoint{4.284985in}{0.611158in}}%
\pgfpathlineto{\pgfqpoint{4.287082in}{0.602797in}}%
\pgfpathlineto{\pgfqpoint{4.289179in}{0.592078in}}%
\pgfpathlineto{\pgfqpoint{4.291277in}{0.578426in}}%
\pgfpathlineto{\pgfqpoint{4.293374in}{0.560802in}}%
\pgfpathlineto{\pgfqpoint{4.295471in}{0.537282in}}%
\pgfpathlineto{\pgfqpoint{4.297569in}{0.504243in}}%
\pgfpathlineto{\pgfqpoint{4.299666in}{0.460347in}}%
\pgfpathlineto{\pgfqpoint{4.301763in}{0.461578in}}%
\pgfpathlineto{\pgfqpoint{4.303861in}{0.502350in}}%
\pgfpathlineto{\pgfqpoint{4.305958in}{0.531418in}}%
\pgfpathlineto{\pgfqpoint{4.308056in}{0.550885in}}%
\pgfpathlineto{\pgfqpoint{4.310153in}{0.563982in}}%
\pgfpathlineto{\pgfqpoint{4.312250in}{0.572130in}}%
\pgfpathlineto{\pgfqpoint{4.314348in}{0.575534in}}%
\pgfpathlineto{\pgfqpoint{4.316445in}{0.572962in}}%
\pgfpathlineto{\pgfqpoint{4.316445in}{0.572962in}}%
\pgfusepath{stroke}%
\end{pgfscope}%
\begin{pgfscope}%
\pgfpathrectangle{\pgfqpoint{2.554660in}{0.320625in}}{\pgfqpoint{1.845679in}{1.164375in}}%
\pgfusepath{clip}%
\pgfsetrectcap%
\pgfsetroundjoin%
\pgfsetlinewidth{1.154312pt}%
\definecolor{currentstroke}{rgb}{0.716387,0.214982,0.475290}%
\pgfsetstrokecolor{currentstroke}%
\pgfsetdash{}{0pt}%
\pgfpathmoveto{\pgfqpoint{2.638555in}{1.468525in}}%
\pgfpathlineto{\pgfqpoint{2.651139in}{1.445854in}}%
\pgfpathlineto{\pgfqpoint{2.661626in}{1.424042in}}%
\pgfpathlineto{\pgfqpoint{2.676308in}{1.390295in}}%
\pgfpathlineto{\pgfqpoint{2.705671in}{1.322202in}}%
\pgfpathlineto{\pgfqpoint{2.722449in}{1.286001in}}%
\pgfpathlineto{\pgfqpoint{2.739228in}{1.252283in}}%
\pgfpathlineto{\pgfqpoint{2.756007in}{1.221029in}}%
\pgfpathlineto{\pgfqpoint{2.772786in}{1.192109in}}%
\pgfpathlineto{\pgfqpoint{2.789565in}{1.165274in}}%
\pgfpathlineto{\pgfqpoint{2.808441in}{1.137126in}}%
\pgfpathlineto{\pgfqpoint{2.831512in}{1.104849in}}%
\pgfpathlineto{\pgfqpoint{2.860875in}{1.065831in}}%
\pgfpathlineto{\pgfqpoint{2.904920in}{1.009355in}}%
\pgfpathlineto{\pgfqpoint{2.993009in}{0.898608in}}%
\pgfpathlineto{\pgfqpoint{3.169188in}{0.678606in}}%
\pgfpathlineto{\pgfqpoint{3.188064in}{0.657756in}}%
\pgfpathlineto{\pgfqpoint{3.196453in}{0.645702in}}%
\pgfpathlineto{\pgfqpoint{3.217427in}{0.613787in}}%
\pgfpathlineto{\pgfqpoint{3.221622in}{0.609248in}}%
\pgfpathlineto{\pgfqpoint{3.225816in}{0.606039in}}%
\pgfpathlineto{\pgfqpoint{3.230011in}{0.604147in}}%
\pgfpathlineto{\pgfqpoint{3.236303in}{0.603011in}}%
\pgfpathlineto{\pgfqpoint{3.242595in}{0.601974in}}%
\pgfpathlineto{\pgfqpoint{3.246790in}{0.599863in}}%
\pgfpathlineto{\pgfqpoint{3.250985in}{0.595206in}}%
\pgfpathlineto{\pgfqpoint{3.255180in}{0.588502in}}%
\pgfpathlineto{\pgfqpoint{3.261472in}{0.575439in}}%
\pgfpathlineto{\pgfqpoint{3.267764in}{0.561636in}}%
\pgfpathlineto{\pgfqpoint{3.269861in}{0.558522in}}%
\pgfpathlineto{\pgfqpoint{3.271958in}{0.556855in}}%
\pgfpathlineto{\pgfqpoint{3.274056in}{0.556608in}}%
\pgfpathlineto{\pgfqpoint{3.278251in}{0.558575in}}%
\pgfpathlineto{\pgfqpoint{3.282445in}{0.560643in}}%
\pgfpathlineto{\pgfqpoint{3.284543in}{0.560984in}}%
\pgfpathlineto{\pgfqpoint{3.286640in}{0.560658in}}%
\pgfpathlineto{\pgfqpoint{3.288737in}{0.559551in}}%
\pgfpathlineto{\pgfqpoint{3.290835in}{0.557531in}}%
\pgfpathlineto{\pgfqpoint{3.292932in}{0.554419in}}%
\pgfpathlineto{\pgfqpoint{3.295029in}{0.549935in}}%
\pgfpathlineto{\pgfqpoint{3.297127in}{0.543608in}}%
\pgfpathlineto{\pgfqpoint{3.299224in}{0.534553in}}%
\pgfpathlineto{\pgfqpoint{3.301322in}{0.520871in}}%
\pgfpathlineto{\pgfqpoint{3.305516in}{0.477283in}}%
\pgfpathlineto{\pgfqpoint{3.307614in}{0.509183in}}%
\pgfpathlineto{\pgfqpoint{3.309711in}{0.530723in}}%
\pgfpathlineto{\pgfqpoint{3.311808in}{0.544885in}}%
\pgfpathlineto{\pgfqpoint{3.316003in}{0.563136in}}%
\pgfpathlineto{\pgfqpoint{3.320198in}{0.574378in}}%
\pgfpathlineto{\pgfqpoint{3.324393in}{0.581256in}}%
\pgfpathlineto{\pgfqpoint{3.326490in}{0.583332in}}%
\pgfpathlineto{\pgfqpoint{3.328587in}{0.584516in}}%
\pgfpathlineto{\pgfqpoint{3.330685in}{0.584751in}}%
\pgfpathlineto{\pgfqpoint{3.332782in}{0.583911in}}%
\pgfpathlineto{\pgfqpoint{3.334879in}{0.581769in}}%
\pgfpathlineto{\pgfqpoint{3.336977in}{0.577873in}}%
\pgfpathlineto{\pgfqpoint{3.341171in}{0.565252in}}%
\pgfpathlineto{\pgfqpoint{3.343269in}{0.564450in}}%
\pgfpathlineto{\pgfqpoint{3.345366in}{0.570921in}}%
\pgfpathlineto{\pgfqpoint{3.349561in}{0.586290in}}%
\pgfpathlineto{\pgfqpoint{3.351658in}{0.591166in}}%
\pgfpathlineto{\pgfqpoint{3.353756in}{0.593959in}}%
\pgfpathlineto{\pgfqpoint{3.355853in}{0.594732in}}%
\pgfpathlineto{\pgfqpoint{3.357950in}{0.593426in}}%
\pgfpathlineto{\pgfqpoint{3.360048in}{0.589761in}}%
\pgfpathlineto{\pgfqpoint{3.362145in}{0.583079in}}%
\pgfpathlineto{\pgfqpoint{3.364242in}{0.571903in}}%
\pgfpathlineto{\pgfqpoint{3.366340in}{0.552759in}}%
\pgfpathlineto{\pgfqpoint{3.368437in}{0.528500in}}%
\pgfpathlineto{\pgfqpoint{3.372632in}{0.574980in}}%
\pgfpathlineto{\pgfqpoint{3.374729in}{0.588673in}}%
\pgfpathlineto{\pgfqpoint{3.376827in}{0.597410in}}%
\pgfpathlineto{\pgfqpoint{3.378924in}{0.602573in}}%
\pgfpathlineto{\pgfqpoint{3.381021in}{0.604503in}}%
\pgfpathlineto{\pgfqpoint{3.383119in}{0.602740in}}%
\pgfpathlineto{\pgfqpoint{3.385216in}{0.595000in}}%
\pgfpathlineto{\pgfqpoint{3.387313in}{0.573549in}}%
\pgfpathlineto{\pgfqpoint{3.389411in}{0.530712in}}%
\pgfpathlineto{\pgfqpoint{3.391508in}{0.583476in}}%
\pgfpathlineto{\pgfqpoint{3.393605in}{0.607367in}}%
\pgfpathlineto{\pgfqpoint{3.395703in}{0.621018in}}%
\pgfpathlineto{\pgfqpoint{3.397800in}{0.629611in}}%
\pgfpathlineto{\pgfqpoint{3.399898in}{0.634917in}}%
\pgfpathlineto{\pgfqpoint{3.401995in}{0.637681in}}%
\pgfpathlineto{\pgfqpoint{3.404092in}{0.638192in}}%
\pgfpathlineto{\pgfqpoint{3.406190in}{0.636430in}}%
\pgfpathlineto{\pgfqpoint{3.408287in}{0.632051in}}%
\pgfpathlineto{\pgfqpoint{3.410384in}{0.624145in}}%
\pgfpathlineto{\pgfqpoint{3.412482in}{0.610353in}}%
\pgfpathlineto{\pgfqpoint{3.414579in}{0.581941in}}%
\pgfpathlineto{\pgfqpoint{3.416676in}{0.533428in}}%
\pgfpathlineto{\pgfqpoint{3.418774in}{0.599588in}}%
\pgfpathlineto{\pgfqpoint{3.420871in}{0.622058in}}%
\pgfpathlineto{\pgfqpoint{3.422969in}{0.635035in}}%
\pgfpathlineto{\pgfqpoint{3.425066in}{0.643113in}}%
\pgfpathlineto{\pgfqpoint{3.427163in}{0.647644in}}%
\pgfpathlineto{\pgfqpoint{3.429261in}{0.648859in}}%
\pgfpathlineto{\pgfqpoint{3.431358in}{0.646102in}}%
\pgfpathlineto{\pgfqpoint{3.433455in}{0.638666in}}%
\pgfpathlineto{\pgfqpoint{3.435553in}{0.627657in}}%
\pgfpathlineto{\pgfqpoint{3.437650in}{0.611616in}}%
\pgfpathlineto{\pgfqpoint{3.439747in}{0.586053in}}%
\pgfpathlineto{\pgfqpoint{3.441845in}{0.536431in}}%
\pgfpathlineto{\pgfqpoint{3.443942in}{0.565415in}}%
\pgfpathlineto{\pgfqpoint{3.446040in}{0.584798in}}%
\pgfpathlineto{\pgfqpoint{3.448137in}{0.590256in}}%
\pgfpathlineto{\pgfqpoint{3.450234in}{0.586773in}}%
\pgfpathlineto{\pgfqpoint{3.452332in}{0.571951in}}%
\pgfpathlineto{\pgfqpoint{3.454429in}{0.519918in}}%
\pgfpathlineto{\pgfqpoint{3.458624in}{0.595283in}}%
\pgfpathlineto{\pgfqpoint{3.460721in}{0.612035in}}%
\pgfpathlineto{\pgfqpoint{3.462818in}{0.621673in}}%
\pgfpathlineto{\pgfqpoint{3.464916in}{0.626072in}}%
\pgfpathlineto{\pgfqpoint{3.467013in}{0.624884in}}%
\pgfpathlineto{\pgfqpoint{3.469111in}{0.617523in}}%
\pgfpathlineto{\pgfqpoint{3.471208in}{0.606462in}}%
\pgfpathlineto{\pgfqpoint{3.473305in}{0.590238in}}%
\pgfpathlineto{\pgfqpoint{3.475403in}{0.563900in}}%
\pgfpathlineto{\pgfqpoint{3.477500in}{0.502650in}}%
\pgfpathlineto{\pgfqpoint{3.479597in}{0.544018in}}%
\pgfpathlineto{\pgfqpoint{3.481695in}{0.564147in}}%
\pgfpathlineto{\pgfqpoint{3.483792in}{0.569726in}}%
\pgfpathlineto{\pgfqpoint{3.485889in}{0.566341in}}%
\pgfpathlineto{\pgfqpoint{3.487987in}{0.551708in}}%
\pgfpathlineto{\pgfqpoint{3.490084in}{0.499049in}}%
\pgfpathlineto{\pgfqpoint{3.494279in}{0.574061in}}%
\pgfpathlineto{\pgfqpoint{3.496376in}{0.591020in}}%
\pgfpathlineto{\pgfqpoint{3.498474in}{0.600760in}}%
\pgfpathlineto{\pgfqpoint{3.500571in}{0.605232in}}%
\pgfpathlineto{\pgfqpoint{3.502668in}{0.604122in}}%
\pgfpathlineto{\pgfqpoint{3.504766in}{0.595901in}}%
\pgfpathlineto{\pgfqpoint{3.506863in}{0.579618in}}%
\pgfpathlineto{\pgfqpoint{3.508960in}{0.543060in}}%
\pgfpathlineto{\pgfqpoint{3.511058in}{0.530447in}}%
\pgfpathlineto{\pgfqpoint{3.513155in}{0.573706in}}%
\pgfpathlineto{\pgfqpoint{3.515253in}{0.590758in}}%
\pgfpathlineto{\pgfqpoint{3.517350in}{0.599636in}}%
\pgfpathlineto{\pgfqpoint{3.519447in}{0.603818in}}%
\pgfpathlineto{\pgfqpoint{3.521545in}{0.604360in}}%
\pgfpathlineto{\pgfqpoint{3.523642in}{0.601335in}}%
\pgfpathlineto{\pgfqpoint{3.525739in}{0.593901in}}%
\pgfpathlineto{\pgfqpoint{3.527837in}{0.579207in}}%
\pgfpathlineto{\pgfqpoint{3.529934in}{0.544244in}}%
\pgfpathlineto{\pgfqpoint{3.532031in}{0.534153in}}%
\pgfpathlineto{\pgfqpoint{3.534129in}{0.580293in}}%
\pgfpathlineto{\pgfqpoint{3.536226in}{0.599787in}}%
\pgfpathlineto{\pgfqpoint{3.538324in}{0.610984in}}%
\pgfpathlineto{\pgfqpoint{3.540421in}{0.617181in}}%
\pgfpathlineto{\pgfqpoint{3.542518in}{0.619061in}}%
\pgfpathlineto{\pgfqpoint{3.544616in}{0.615785in}}%
\pgfpathlineto{\pgfqpoint{3.546713in}{0.604373in}}%
\pgfpathlineto{\pgfqpoint{3.548810in}{0.576309in}}%
\pgfpathlineto{\pgfqpoint{3.550908in}{0.533635in}}%
\pgfpathlineto{\pgfqpoint{3.553005in}{0.595649in}}%
\pgfpathlineto{\pgfqpoint{3.555102in}{0.616958in}}%
\pgfpathlineto{\pgfqpoint{3.557200in}{0.628834in}}%
\pgfpathlineto{\pgfqpoint{3.559297in}{0.635879in}}%
\pgfpathlineto{\pgfqpoint{3.561395in}{0.639624in}}%
\pgfpathlineto{\pgfqpoint{3.563492in}{0.640644in}}%
\pgfpathlineto{\pgfqpoint{3.565589in}{0.639020in}}%
\pgfpathlineto{\pgfqpoint{3.567687in}{0.634383in}}%
\pgfpathlineto{\pgfqpoint{3.569784in}{0.625617in}}%
\pgfpathlineto{\pgfqpoint{3.571881in}{0.609529in}}%
\pgfpathlineto{\pgfqpoint{3.573979in}{0.571125in}}%
\pgfpathlineto{\pgfqpoint{3.576076in}{0.570431in}}%
\pgfpathlineto{\pgfqpoint{3.578173in}{0.612706in}}%
\pgfpathlineto{\pgfqpoint{3.580271in}{0.631501in}}%
\pgfpathlineto{\pgfqpoint{3.582368in}{0.642614in}}%
\pgfpathlineto{\pgfqpoint{3.584465in}{0.649168in}}%
\pgfpathlineto{\pgfqpoint{3.586563in}{0.651990in}}%
\pgfpathlineto{\pgfqpoint{3.588660in}{0.650776in}}%
\pgfpathlineto{\pgfqpoint{3.590758in}{0.644453in}}%
\pgfpathlineto{\pgfqpoint{3.592855in}{0.634521in}}%
\pgfpathlineto{\pgfqpoint{3.594952in}{0.620280in}}%
\pgfpathlineto{\pgfqpoint{3.597050in}{0.598613in}}%
\pgfpathlineto{\pgfqpoint{3.599147in}{0.555860in}}%
\pgfpathlineto{\pgfqpoint{3.601244in}{0.542468in}}%
\pgfpathlineto{\pgfqpoint{3.603342in}{0.578167in}}%
\pgfpathlineto{\pgfqpoint{3.605439in}{0.587502in}}%
\pgfpathlineto{\pgfqpoint{3.607536in}{0.585854in}}%
\pgfpathlineto{\pgfqpoint{3.609634in}{0.572311in}}%
\pgfpathlineto{\pgfqpoint{3.611731in}{0.520143in}}%
\pgfpathlineto{\pgfqpoint{3.615926in}{0.595393in}}%
\pgfpathlineto{\pgfqpoint{3.618023in}{0.612750in}}%
\pgfpathlineto{\pgfqpoint{3.620121in}{0.622771in}}%
\pgfpathlineto{\pgfqpoint{3.622218in}{0.627477in}}%
\pgfpathlineto{\pgfqpoint{3.624315in}{0.626594in}}%
\pgfpathlineto{\pgfqpoint{3.626413in}{0.619438in}}%
\pgfpathlineto{\pgfqpoint{3.628510in}{0.608725in}}%
\pgfpathlineto{\pgfqpoint{3.630607in}{0.593290in}}%
\pgfpathlineto{\pgfqpoint{3.632705in}{0.569106in}}%
\pgfpathlineto{\pgfqpoint{3.634802in}{0.513651in}}%
\pgfpathlineto{\pgfqpoint{3.636900in}{0.531944in}}%
\pgfpathlineto{\pgfqpoint{3.638997in}{0.558164in}}%
\pgfpathlineto{\pgfqpoint{3.641094in}{0.564814in}}%
\pgfpathlineto{\pgfqpoint{3.643192in}{0.561042in}}%
\pgfpathlineto{\pgfqpoint{3.645289in}{0.543768in}}%
\pgfpathlineto{\pgfqpoint{3.647386in}{0.444665in}}%
\pgfpathlineto{\pgfqpoint{3.649484in}{0.547914in}}%
\pgfpathlineto{\pgfqpoint{3.651581in}{0.576687in}}%
\pgfpathlineto{\pgfqpoint{3.653678in}{0.592156in}}%
\pgfpathlineto{\pgfqpoint{3.655776in}{0.600990in}}%
\pgfpathlineto{\pgfqpoint{3.657873in}{0.604612in}}%
\pgfpathlineto{\pgfqpoint{3.659971in}{0.602306in}}%
\pgfpathlineto{\pgfqpoint{3.662068in}{0.592961in}}%
\pgfpathlineto{\pgfqpoint{3.664165in}{0.575306in}}%
\pgfpathlineto{\pgfqpoint{3.666263in}{0.534077in}}%
\pgfpathlineto{\pgfqpoint{3.668360in}{0.533900in}}%
\pgfpathlineto{\pgfqpoint{3.670457in}{0.572546in}}%
\pgfpathlineto{\pgfqpoint{3.672555in}{0.588319in}}%
\pgfpathlineto{\pgfqpoint{3.674652in}{0.596359in}}%
\pgfpathlineto{\pgfqpoint{3.676749in}{0.599766in}}%
\pgfpathlineto{\pgfqpoint{3.678847in}{0.599427in}}%
\pgfpathlineto{\pgfqpoint{3.680944in}{0.595211in}}%
\pgfpathlineto{\pgfqpoint{3.683042in}{0.585818in}}%
\pgfpathlineto{\pgfqpoint{3.685139in}{0.566682in}}%
\pgfpathlineto{\pgfqpoint{3.687236in}{0.503327in}}%
\pgfpathlineto{\pgfqpoint{3.689334in}{0.554538in}}%
\pgfpathlineto{\pgfqpoint{3.691431in}{0.585440in}}%
\pgfpathlineto{\pgfqpoint{3.693528in}{0.601150in}}%
\pgfpathlineto{\pgfqpoint{3.695626in}{0.610263in}}%
\pgfpathlineto{\pgfqpoint{3.697723in}{0.614745in}}%
\pgfpathlineto{\pgfqpoint{3.699820in}{0.614658in}}%
\pgfpathlineto{\pgfqpoint{3.701918in}{0.608280in}}%
\pgfpathlineto{\pgfqpoint{3.704015in}{0.591337in}}%
\pgfpathlineto{\pgfqpoint{3.706113in}{0.539665in}}%
\pgfpathlineto{\pgfqpoint{3.710307in}{0.603966in}}%
\pgfpathlineto{\pgfqpoint{3.712405in}{0.619992in}}%
\pgfpathlineto{\pgfqpoint{3.714502in}{0.629405in}}%
\pgfpathlineto{\pgfqpoint{3.716599in}{0.634875in}}%
\pgfpathlineto{\pgfqpoint{3.718697in}{0.637394in}}%
\pgfpathlineto{\pgfqpoint{3.720794in}{0.637294in}}%
\pgfpathlineto{\pgfqpoint{3.722891in}{0.634472in}}%
\pgfpathlineto{\pgfqpoint{3.724989in}{0.628313in}}%
\pgfpathlineto{\pgfqpoint{3.727086in}{0.617132in}}%
\pgfpathlineto{\pgfqpoint{3.729184in}{0.595500in}}%
\pgfpathlineto{\pgfqpoint{3.731281in}{0.508000in}}%
\pgfpathlineto{\pgfqpoint{3.733378in}{0.590562in}}%
\pgfpathlineto{\pgfqpoint{3.735476in}{0.618671in}}%
\pgfpathlineto{\pgfqpoint{3.737573in}{0.633635in}}%
\pgfpathlineto{\pgfqpoint{3.739670in}{0.642690in}}%
\pgfpathlineto{\pgfqpoint{3.741768in}{0.647699in}}%
\pgfpathlineto{\pgfqpoint{3.743865in}{0.648997in}}%
\pgfpathlineto{\pgfqpoint{3.745962in}{0.645801in}}%
\pgfpathlineto{\pgfqpoint{3.748060in}{0.637808in}}%
\pgfpathlineto{\pgfqpoint{3.750157in}{0.626071in}}%
\pgfpathlineto{\pgfqpoint{3.752254in}{0.608673in}}%
\pgfpathlineto{\pgfqpoint{3.754352in}{0.579106in}}%
\pgfpathlineto{\pgfqpoint{3.756449in}{0.453982in}}%
\pgfpathlineto{\pgfqpoint{3.758547in}{0.569520in}}%
\pgfpathlineto{\pgfqpoint{3.760644in}{0.587045in}}%
\pgfpathlineto{\pgfqpoint{3.762741in}{0.591205in}}%
\pgfpathlineto{\pgfqpoint{3.764839in}{0.586340in}}%
\pgfpathlineto{\pgfqpoint{3.766936in}{0.568720in}}%
\pgfpathlineto{\pgfqpoint{3.769033in}{0.477363in}}%
\pgfpathlineto{\pgfqpoint{3.771131in}{0.570379in}}%
\pgfpathlineto{\pgfqpoint{3.773228in}{0.599385in}}%
\pgfpathlineto{\pgfqpoint{3.775325in}{0.614821in}}%
\pgfpathlineto{\pgfqpoint{3.777423in}{0.623591in}}%
\pgfpathlineto{\pgfqpoint{3.779520in}{0.627158in}}%
\pgfpathlineto{\pgfqpoint{3.781618in}{0.624830in}}%
\pgfpathlineto{\pgfqpoint{3.783715in}{0.616759in}}%
\pgfpathlineto{\pgfqpoint{3.785812in}{0.604921in}}%
\pgfpathlineto{\pgfqpoint{3.787910in}{0.587347in}}%
\pgfpathlineto{\pgfqpoint{3.790007in}{0.557295in}}%
\pgfpathlineto{\pgfqpoint{3.792104in}{0.442491in}}%
\pgfpathlineto{\pgfqpoint{3.794202in}{0.549382in}}%
\pgfpathlineto{\pgfqpoint{3.796299in}{0.566441in}}%
\pgfpathlineto{\pgfqpoint{3.798396in}{0.570382in}}%
\pgfpathlineto{\pgfqpoint{3.800494in}{0.565302in}}%
\pgfpathlineto{\pgfqpoint{3.802591in}{0.547249in}}%
\pgfpathlineto{\pgfqpoint{3.804689in}{0.444929in}}%
\pgfpathlineto{\pgfqpoint{3.806786in}{0.550513in}}%
\pgfpathlineto{\pgfqpoint{3.808883in}{0.579014in}}%
\pgfpathlineto{\pgfqpoint{3.810981in}{0.594261in}}%
\pgfpathlineto{\pgfqpoint{3.813078in}{0.602907in}}%
\pgfpathlineto{\pgfqpoint{3.815175in}{0.606357in}}%
\pgfpathlineto{\pgfqpoint{3.817273in}{0.603867in}}%
\pgfpathlineto{\pgfqpoint{3.819370in}{0.594381in}}%
\pgfpathlineto{\pgfqpoint{3.821467in}{0.576468in}}%
\pgfpathlineto{\pgfqpoint{3.823565in}{0.534032in}}%
\pgfpathlineto{\pgfqpoint{3.825662in}{0.537107in}}%
\pgfpathlineto{\pgfqpoint{3.827760in}{0.574722in}}%
\pgfpathlineto{\pgfqpoint{3.829857in}{0.590257in}}%
\pgfpathlineto{\pgfqpoint{3.831954in}{0.598174in}}%
\pgfpathlineto{\pgfqpoint{3.834052in}{0.601492in}}%
\pgfpathlineto{\pgfqpoint{3.836149in}{0.601069in}}%
\pgfpathlineto{\pgfqpoint{3.838246in}{0.596754in}}%
\pgfpathlineto{\pgfqpoint{3.840344in}{0.587207in}}%
\pgfpathlineto{\pgfqpoint{3.842441in}{0.567709in}}%
\pgfpathlineto{\pgfqpoint{3.844538in}{0.500561in}}%
\pgfpathlineto{\pgfqpoint{3.846636in}{0.557364in}}%
\pgfpathlineto{\pgfqpoint{3.848733in}{0.587626in}}%
\pgfpathlineto{\pgfqpoint{3.850831in}{0.603129in}}%
\pgfpathlineto{\pgfqpoint{3.852928in}{0.612118in}}%
\pgfpathlineto{\pgfqpoint{3.855025in}{0.616494in}}%
\pgfpathlineto{\pgfqpoint{3.857123in}{0.616283in}}%
\pgfpathlineto{\pgfqpoint{3.859220in}{0.609751in}}%
\pgfpathlineto{\pgfqpoint{3.861317in}{0.593026in}}%
\pgfpathlineto{\pgfqpoint{3.863415in}{0.544601in}}%
\pgfpathlineto{\pgfqpoint{3.865512in}{0.568630in}}%
\pgfpathlineto{\pgfqpoint{3.867609in}{0.603046in}}%
\pgfpathlineto{\pgfqpoint{3.869707in}{0.619254in}}%
\pgfpathlineto{\pgfqpoint{3.871804in}{0.628652in}}%
\pgfpathlineto{\pgfqpoint{3.873902in}{0.634020in}}%
\pgfpathlineto{\pgfqpoint{3.875999in}{0.636367in}}%
\pgfpathlineto{\pgfqpoint{3.878096in}{0.636014in}}%
\pgfpathlineto{\pgfqpoint{3.880194in}{0.632816in}}%
\pgfpathlineto{\pgfqpoint{3.882291in}{0.626049in}}%
\pgfpathlineto{\pgfqpoint{3.884388in}{0.613701in}}%
\pgfpathlineto{\pgfqpoint{3.886486in}{0.588713in}}%
\pgfpathlineto{\pgfqpoint{3.888583in}{0.496538in}}%
\pgfpathlineto{\pgfqpoint{3.890680in}{0.596848in}}%
\pgfpathlineto{\pgfqpoint{3.892778in}{0.621390in}}%
\pgfpathlineto{\pgfqpoint{3.894875in}{0.635014in}}%
\pgfpathlineto{\pgfqpoint{3.896973in}{0.643211in}}%
\pgfpathlineto{\pgfqpoint{3.899070in}{0.647448in}}%
\pgfpathlineto{\pgfqpoint{3.901167in}{0.647825in}}%
\pgfpathlineto{\pgfqpoint{3.903265in}{0.643258in}}%
\pgfpathlineto{\pgfqpoint{3.905362in}{0.634349in}}%
\pgfpathlineto{\pgfqpoint{3.907459in}{0.621269in}}%
\pgfpathlineto{\pgfqpoint{3.909557in}{0.601300in}}%
\pgfpathlineto{\pgfqpoint{3.913751in}{0.533312in}}%
\pgfpathlineto{\pgfqpoint{3.915849in}{0.578304in}}%
\pgfpathlineto{\pgfqpoint{3.917946in}{0.590673in}}%
\pgfpathlineto{\pgfqpoint{3.920043in}{0.592462in}}%
\pgfpathlineto{\pgfqpoint{3.922141in}{0.585522in}}%
\pgfpathlineto{\pgfqpoint{3.924238in}{0.564009in}}%
\pgfpathlineto{\pgfqpoint{3.926336in}{0.481364in}}%
\pgfpathlineto{\pgfqpoint{3.928433in}{0.577610in}}%
\pgfpathlineto{\pgfqpoint{3.930530in}{0.603042in}}%
\pgfpathlineto{\pgfqpoint{3.932628in}{0.617024in}}%
\pgfpathlineto{\pgfqpoint{3.934725in}{0.624804in}}%
\pgfpathlineto{\pgfqpoint{3.936822in}{0.627416in}}%
\pgfpathlineto{\pgfqpoint{3.938920in}{0.623821in}}%
\pgfpathlineto{\pgfqpoint{3.941017in}{0.615107in}}%
\pgfpathlineto{\pgfqpoint{3.943114in}{0.602305in}}%
\pgfpathlineto{\pgfqpoint{3.945212in}{0.582859in}}%
\pgfpathlineto{\pgfqpoint{3.949407in}{0.507096in}}%
\pgfpathlineto{\pgfqpoint{3.951504in}{0.557296in}}%
\pgfpathlineto{\pgfqpoint{3.953601in}{0.570486in}}%
\pgfpathlineto{\pgfqpoint{3.955699in}{0.572754in}}%
\pgfpathlineto{\pgfqpoint{3.957796in}{0.566356in}}%
\pgfpathlineto{\pgfqpoint{3.959893in}{0.546115in}}%
\pgfpathlineto{\pgfqpoint{3.961991in}{0.425225in}}%
\pgfpathlineto{\pgfqpoint{3.964088in}{0.555706in}}%
\pgfpathlineto{\pgfqpoint{3.966185in}{0.582184in}}%
\pgfpathlineto{\pgfqpoint{3.968283in}{0.596598in}}%
\pgfpathlineto{\pgfqpoint{3.970380in}{0.604676in}}%
\pgfpathlineto{\pgfqpoint{3.972478in}{0.607583in}}%
\pgfpathlineto{\pgfqpoint{3.974575in}{0.604371in}}%
\pgfpathlineto{\pgfqpoint{3.976672in}{0.594290in}}%
\pgfpathlineto{\pgfqpoint{3.978770in}{0.575367in}}%
\pgfpathlineto{\pgfqpoint{3.980867in}{0.528170in}}%
\pgfpathlineto{\pgfqpoint{3.982964in}{0.542708in}}%
\pgfpathlineto{\pgfqpoint{3.985062in}{0.576960in}}%
\pgfpathlineto{\pgfqpoint{3.987159in}{0.591626in}}%
\pgfpathlineto{\pgfqpoint{3.989256in}{0.599047in}}%
\pgfpathlineto{\pgfqpoint{3.991354in}{0.601963in}}%
\pgfpathlineto{\pgfqpoint{3.993451in}{0.601120in}}%
\pgfpathlineto{\pgfqpoint{3.995549in}{0.596259in}}%
\pgfpathlineto{\pgfqpoint{3.997646in}{0.585803in}}%
\pgfpathlineto{\pgfqpoint{3.999743in}{0.563992in}}%
\pgfpathlineto{\pgfqpoint{4.001841in}{0.450222in}}%
\pgfpathlineto{\pgfqpoint{4.003938in}{0.563895in}}%
\pgfpathlineto{\pgfqpoint{4.006035in}{0.590891in}}%
\pgfpathlineto{\pgfqpoint{4.008133in}{0.605258in}}%
\pgfpathlineto{\pgfqpoint{4.010230in}{0.613541in}}%
\pgfpathlineto{\pgfqpoint{4.012327in}{0.617287in}}%
\pgfpathlineto{\pgfqpoint{4.014425in}{0.616296in}}%
\pgfpathlineto{\pgfqpoint{4.016522in}{0.608577in}}%
\pgfpathlineto{\pgfqpoint{4.018620in}{0.589999in}}%
\pgfpathlineto{\pgfqpoint{4.020717in}{0.530200in}}%
\pgfpathlineto{\pgfqpoint{4.022814in}{0.574042in}}%
\pgfpathlineto{\pgfqpoint{4.024912in}{0.604760in}}%
\pgfpathlineto{\pgfqpoint{4.027009in}{0.619823in}}%
\pgfpathlineto{\pgfqpoint{4.029106in}{0.628551in}}%
\pgfpathlineto{\pgfqpoint{4.031204in}{0.633393in}}%
\pgfpathlineto{\pgfqpoint{4.033301in}{0.635244in}}%
\pgfpathlineto{\pgfqpoint{4.035398in}{0.634345in}}%
\pgfpathlineto{\pgfqpoint{4.037496in}{0.630452in}}%
\pgfpathlineto{\pgfqpoint{4.039593in}{0.622638in}}%
\pgfpathlineto{\pgfqpoint{4.041691in}{0.608271in}}%
\pgfpathlineto{\pgfqpoint{4.045885in}{0.550475in}}%
\pgfpathlineto{\pgfqpoint{4.047983in}{0.603816in}}%
\pgfpathlineto{\pgfqpoint{4.050080in}{0.624563in}}%
\pgfpathlineto{\pgfqpoint{4.052177in}{0.636556in}}%
\pgfpathlineto{\pgfqpoint{4.054275in}{0.643650in}}%
\pgfpathlineto{\pgfqpoint{4.056372in}{0.646872in}}%
\pgfpathlineto{\pgfqpoint{4.058469in}{0.646010in}}%
\pgfpathlineto{\pgfqpoint{4.060567in}{0.639916in}}%
\pgfpathlineto{\pgfqpoint{4.062664in}{0.629830in}}%
\pgfpathlineto{\pgfqpoint{4.064762in}{0.614879in}}%
\pgfpathlineto{\pgfqpoint{4.066859in}{0.590835in}}%
\pgfpathlineto{\pgfqpoint{4.068956in}{0.531403in}}%
\pgfpathlineto{\pgfqpoint{4.073151in}{0.585450in}}%
\pgfpathlineto{\pgfqpoint{4.075248in}{0.593667in}}%
\pgfpathlineto{\pgfqpoint{4.077346in}{0.593079in}}%
\pgfpathlineto{\pgfqpoint{4.079443in}{0.583664in}}%
\pgfpathlineto{\pgfqpoint{4.083638in}{0.526790in}}%
\pgfpathlineto{\pgfqpoint{4.085735in}{0.584697in}}%
\pgfpathlineto{\pgfqpoint{4.087832in}{0.606780in}}%
\pgfpathlineto{\pgfqpoint{4.089930in}{0.619220in}}%
\pgfpathlineto{\pgfqpoint{4.092027in}{0.625878in}}%
\pgfpathlineto{\pgfqpoint{4.094125in}{0.627340in}}%
\pgfpathlineto{\pgfqpoint{4.096222in}{0.622373in}}%
\pgfpathlineto{\pgfqpoint{4.098319in}{0.612850in}}%
\pgfpathlineto{\pgfqpoint{4.100417in}{0.598760in}}%
\pgfpathlineto{\pgfqpoint{4.102514in}{0.576583in}}%
\pgfpathlineto{\pgfqpoint{4.104611in}{0.528455in}}%
\pgfpathlineto{\pgfqpoint{4.106709in}{0.532810in}}%
\pgfpathlineto{\pgfqpoint{4.108806in}{0.564524in}}%
\pgfpathlineto{\pgfqpoint{4.110903in}{0.574382in}}%
\pgfpathlineto{\pgfqpoint{4.113001in}{0.574955in}}%
\pgfpathlineto{\pgfqpoint{4.115098in}{0.567045in}}%
\pgfpathlineto{\pgfqpoint{4.117196in}{0.543870in}}%
\pgfpathlineto{\pgfqpoint{4.119293in}{0.481680in}}%
\pgfpathlineto{\pgfqpoint{4.121390in}{0.561283in}}%
\pgfpathlineto{\pgfqpoint{4.123488in}{0.585626in}}%
\pgfpathlineto{\pgfqpoint{4.125585in}{0.599103in}}%
\pgfpathlineto{\pgfqpoint{4.127682in}{0.606519in}}%
\pgfpathlineto{\pgfqpoint{4.129780in}{0.608778in}}%
\pgfpathlineto{\pgfqpoint{4.131877in}{0.604720in}}%
\pgfpathlineto{\pgfqpoint{4.133974in}{0.593902in}}%
\pgfpathlineto{\pgfqpoint{4.136072in}{0.573588in}}%
\pgfpathlineto{\pgfqpoint{4.138169in}{0.518089in}}%
\pgfpathlineto{\pgfqpoint{4.142364in}{0.579750in}}%
\pgfpathlineto{\pgfqpoint{4.144461in}{0.593389in}}%
\pgfpathlineto{\pgfqpoint{4.146559in}{0.600226in}}%
\pgfpathlineto{\pgfqpoint{4.148656in}{0.602674in}}%
\pgfpathlineto{\pgfqpoint{4.150753in}{0.601351in}}%
\pgfpathlineto{\pgfqpoint{4.152851in}{0.595868in}}%
\pgfpathlineto{\pgfqpoint{4.154948in}{0.584350in}}%
\pgfpathlineto{\pgfqpoint{4.157045in}{0.559599in}}%
\pgfpathlineto{\pgfqpoint{4.159143in}{0.477951in}}%
\pgfpathlineto{\pgfqpoint{4.161240in}{0.570027in}}%
\pgfpathlineto{\pgfqpoint{4.163338in}{0.594256in}}%
\pgfpathlineto{\pgfqpoint{4.165435in}{0.607543in}}%
\pgfpathlineto{\pgfqpoint{4.167532in}{0.615119in}}%
\pgfpathlineto{\pgfqpoint{4.169630in}{0.618209in}}%
\pgfpathlineto{\pgfqpoint{4.171727in}{0.616374in}}%
\pgfpathlineto{\pgfqpoint{4.173824in}{0.607428in}}%
\pgfpathlineto{\pgfqpoint{4.175922in}{0.586810in}}%
\pgfpathlineto{\pgfqpoint{4.178019in}{0.506864in}}%
\pgfpathlineto{\pgfqpoint{4.180116in}{0.578843in}}%
\pgfpathlineto{\pgfqpoint{4.182214in}{0.606502in}}%
\pgfpathlineto{\pgfqpoint{4.184311in}{0.620501in}}%
\pgfpathlineto{\pgfqpoint{4.186409in}{0.628576in}}%
\pgfpathlineto{\pgfqpoint{4.188506in}{0.632893in}}%
\pgfpathlineto{\pgfqpoint{4.190603in}{0.634233in}}%
\pgfpathlineto{\pgfqpoint{4.192701in}{0.632754in}}%
\pgfpathlineto{\pgfqpoint{4.194798in}{0.628098in}}%
\pgfpathlineto{\pgfqpoint{4.196895in}{0.619080in}}%
\pgfpathlineto{\pgfqpoint{4.198993in}{0.602176in}}%
\pgfpathlineto{\pgfqpoint{4.201090in}{0.559106in}}%
\pgfpathlineto{\pgfqpoint{4.203187in}{0.571406in}}%
\pgfpathlineto{\pgfqpoint{4.205285in}{0.609620in}}%
\pgfpathlineto{\pgfqpoint{4.207382in}{0.627461in}}%
\pgfpathlineto{\pgfqpoint{4.209480in}{0.638014in}}%
\pgfpathlineto{\pgfqpoint{4.211577in}{0.644047in}}%
\pgfpathlineto{\pgfqpoint{4.213674in}{0.646211in}}%
\pgfpathlineto{\pgfqpoint{4.215772in}{0.643948in}}%
\pgfpathlineto{\pgfqpoint{4.217869in}{0.636480in}}%
\pgfpathlineto{\pgfqpoint{4.219966in}{0.625056in}}%
\pgfpathlineto{\pgfqpoint{4.222064in}{0.607797in}}%
\pgfpathlineto{\pgfqpoint{4.224161in}{0.577725in}}%
\pgfpathlineto{\pgfqpoint{4.226258in}{0.473722in}}%
\pgfpathlineto{\pgfqpoint{4.228356in}{0.572940in}}%
\pgfpathlineto{\pgfqpoint{4.230453in}{0.590489in}}%
\pgfpathlineto{\pgfqpoint{4.232551in}{0.595638in}}%
\pgfpathlineto{\pgfqpoint{4.234648in}{0.592874in}}%
\pgfpathlineto{\pgfqpoint{4.236745in}{0.580672in}}%
\pgfpathlineto{\pgfqpoint{4.238843in}{0.544030in}}%
\pgfpathlineto{\pgfqpoint{4.240940in}{0.548673in}}%
\pgfpathlineto{\pgfqpoint{4.243037in}{0.591122in}}%
\pgfpathlineto{\pgfqpoint{4.245135in}{0.610347in}}%
\pgfpathlineto{\pgfqpoint{4.247232in}{0.621314in}}%
\pgfpathlineto{\pgfqpoint{4.249329in}{0.626818in}}%
\pgfpathlineto{\pgfqpoint{4.251427in}{0.627005in}}%
\pgfpathlineto{\pgfqpoint{4.253524in}{0.620731in}}%
\pgfpathlineto{\pgfqpoint{4.255621in}{0.610272in}}%
\pgfpathlineto{\pgfqpoint{4.257719in}{0.594621in}}%
\pgfpathlineto{\pgfqpoint{4.259816in}{0.568742in}}%
\pgfpathlineto{\pgfqpoint{4.261914in}{0.490591in}}%
\pgfpathlineto{\pgfqpoint{4.264011in}{0.547751in}}%
\pgfpathlineto{\pgfqpoint{4.266108in}{0.570264in}}%
\pgfpathlineto{\pgfqpoint{4.268206in}{0.577479in}}%
\pgfpathlineto{\pgfqpoint{4.270303in}{0.576436in}}%
\pgfpathlineto{\pgfqpoint{4.272400in}{0.566793in}}%
\pgfpathlineto{\pgfqpoint{4.276595in}{0.508705in}}%
\pgfpathlineto{\pgfqpoint{4.278692in}{0.567128in}}%
\pgfpathlineto{\pgfqpoint{4.280790in}{0.589221in}}%
\pgfpathlineto{\pgfqpoint{4.282887in}{0.601644in}}%
\pgfpathlineto{\pgfqpoint{4.284985in}{0.608288in}}%
\pgfpathlineto{\pgfqpoint{4.287082in}{0.609751in}}%
\pgfpathlineto{\pgfqpoint{4.289179in}{0.604923in}}%
\pgfpathlineto{\pgfqpoint{4.293374in}{0.588748in}}%
\pgfpathlineto{\pgfqpoint{4.301763in}{0.549653in}}%
\pgfpathlineto{\pgfqpoint{4.303861in}{0.543607in}}%
\pgfpathlineto{\pgfqpoint{4.305958in}{0.541653in}}%
\pgfpathlineto{\pgfqpoint{4.308056in}{0.543323in}}%
\pgfpathlineto{\pgfqpoint{4.312250in}{0.549985in}}%
\pgfpathlineto{\pgfqpoint{4.314348in}{0.551059in}}%
\pgfpathlineto{\pgfqpoint{4.316445in}{0.547781in}}%
\pgfpathlineto{\pgfqpoint{4.316445in}{0.547781in}}%
\pgfusepath{stroke}%
\end{pgfscope}%
\begin{pgfscope}%
\pgfpathrectangle{\pgfqpoint{2.554660in}{0.320625in}}{\pgfqpoint{1.845679in}{1.164375in}}%
\pgfusepath{clip}%
\pgfsetrectcap%
\pgfsetroundjoin%
\pgfsetlinewidth{1.154312pt}%
\definecolor{currentstroke}{rgb}{0.852126,0.276106,0.418573}%
\pgfsetstrokecolor{currentstroke}%
\pgfsetdash{}{0pt}%
\pgfpathmoveto{\pgfqpoint{2.638555in}{1.468525in}}%
\pgfpathlineto{\pgfqpoint{2.649042in}{1.449644in}}%
\pgfpathlineto{\pgfqpoint{2.659529in}{1.428050in}}%
\pgfpathlineto{\pgfqpoint{2.672113in}{1.399087in}}%
\pgfpathlineto{\pgfqpoint{2.711963in}{1.305175in}}%
\pgfpathlineto{\pgfqpoint{2.730839in}{1.264095in}}%
\pgfpathlineto{\pgfqpoint{2.749715in}{1.225519in}}%
\pgfpathlineto{\pgfqpoint{2.770689in}{1.185267in}}%
\pgfpathlineto{\pgfqpoint{2.791662in}{1.147413in}}%
\pgfpathlineto{\pgfqpoint{2.814733in}{1.108125in}}%
\pgfpathlineto{\pgfqpoint{2.839902in}{1.067513in}}%
\pgfpathlineto{\pgfqpoint{2.869265in}{1.022269in}}%
\pgfpathlineto{\pgfqpoint{2.909115in}{0.963121in}}%
\pgfpathlineto{\pgfqpoint{2.967841in}{0.878191in}}%
\pgfpathlineto{\pgfqpoint{3.062222in}{0.743874in}}%
\pgfpathlineto{\pgfqpoint{3.076904in}{0.722611in}}%
\pgfpathlineto{\pgfqpoint{3.091585in}{0.701551in}}%
\pgfpathlineto{\pgfqpoint{3.114656in}{0.670328in}}%
\pgfpathlineto{\pgfqpoint{3.137727in}{0.634044in}}%
\pgfpathlineto{\pgfqpoint{3.141922in}{0.629569in}}%
\pgfpathlineto{\pgfqpoint{3.146117in}{0.626225in}}%
\pgfpathlineto{\pgfqpoint{3.158701in}{0.617911in}}%
\pgfpathlineto{\pgfqpoint{3.162896in}{0.612135in}}%
\pgfpathlineto{\pgfqpoint{3.167090in}{0.604532in}}%
\pgfpathlineto{\pgfqpoint{3.171285in}{0.595194in}}%
\pgfpathlineto{\pgfqpoint{3.177577in}{0.577621in}}%
\pgfpathlineto{\pgfqpoint{3.190161in}{0.538413in}}%
\pgfpathlineto{\pgfqpoint{3.196453in}{0.527277in}}%
\pgfpathlineto{\pgfqpoint{3.200648in}{0.520443in}}%
\pgfpathlineto{\pgfqpoint{3.202746in}{0.519628in}}%
\pgfpathlineto{\pgfqpoint{3.204843in}{0.523163in}}%
\pgfpathlineto{\pgfqpoint{3.209038in}{0.539036in}}%
\pgfpathlineto{\pgfqpoint{3.213232in}{0.554727in}}%
\pgfpathlineto{\pgfqpoint{3.217427in}{0.566353in}}%
\pgfpathlineto{\pgfqpoint{3.221622in}{0.573986in}}%
\pgfpathlineto{\pgfqpoint{3.223719in}{0.576346in}}%
\pgfpathlineto{\pgfqpoint{3.225816in}{0.577678in}}%
\pgfpathlineto{\pgfqpoint{3.227914in}{0.577863in}}%
\pgfpathlineto{\pgfqpoint{3.230011in}{0.576672in}}%
\pgfpathlineto{\pgfqpoint{3.232109in}{0.573717in}}%
\pgfpathlineto{\pgfqpoint{3.234206in}{0.569788in}}%
\pgfpathlineto{\pgfqpoint{3.236303in}{0.570973in}}%
\pgfpathlineto{\pgfqpoint{3.240498in}{0.594455in}}%
\pgfpathlineto{\pgfqpoint{3.244693in}{0.615555in}}%
\pgfpathlineto{\pgfqpoint{3.248887in}{0.629285in}}%
\pgfpathlineto{\pgfqpoint{3.253082in}{0.637885in}}%
\pgfpathlineto{\pgfqpoint{3.257277in}{0.642641in}}%
\pgfpathlineto{\pgfqpoint{3.259374in}{0.643767in}}%
\pgfpathlineto{\pgfqpoint{3.261472in}{0.644088in}}%
\pgfpathlineto{\pgfqpoint{3.263569in}{0.643586in}}%
\pgfpathlineto{\pgfqpoint{3.265666in}{0.642205in}}%
\pgfpathlineto{\pgfqpoint{3.267764in}{0.639841in}}%
\pgfpathlineto{\pgfqpoint{3.269861in}{0.636312in}}%
\pgfpathlineto{\pgfqpoint{3.271958in}{0.631313in}}%
\pgfpathlineto{\pgfqpoint{3.274056in}{0.624299in}}%
\pgfpathlineto{\pgfqpoint{3.276153in}{0.614186in}}%
\pgfpathlineto{\pgfqpoint{3.278251in}{0.598354in}}%
\pgfpathlineto{\pgfqpoint{3.282445in}{0.541344in}}%
\pgfpathlineto{\pgfqpoint{3.284543in}{0.591149in}}%
\pgfpathlineto{\pgfqpoint{3.286640in}{0.612533in}}%
\pgfpathlineto{\pgfqpoint{3.288737in}{0.625973in}}%
\pgfpathlineto{\pgfqpoint{3.292932in}{0.642647in}}%
\pgfpathlineto{\pgfqpoint{3.297127in}{0.651970in}}%
\pgfpathlineto{\pgfqpoint{3.299224in}{0.654664in}}%
\pgfpathlineto{\pgfqpoint{3.301322in}{0.656138in}}%
\pgfpathlineto{\pgfqpoint{3.303419in}{0.656317in}}%
\pgfpathlineto{\pgfqpoint{3.305516in}{0.654992in}}%
\pgfpathlineto{\pgfqpoint{3.307614in}{0.651730in}}%
\pgfpathlineto{\pgfqpoint{3.309711in}{0.645837in}}%
\pgfpathlineto{\pgfqpoint{3.313906in}{0.628108in}}%
\pgfpathlineto{\pgfqpoint{3.318100in}{0.602874in}}%
\pgfpathlineto{\pgfqpoint{3.322295in}{0.578604in}}%
\pgfpathlineto{\pgfqpoint{3.326490in}{0.567329in}}%
\pgfpathlineto{\pgfqpoint{3.328587in}{0.564701in}}%
\pgfpathlineto{\pgfqpoint{3.330685in}{0.571442in}}%
\pgfpathlineto{\pgfqpoint{3.334879in}{0.596928in}}%
\pgfpathlineto{\pgfqpoint{3.336977in}{0.605459in}}%
\pgfpathlineto{\pgfqpoint{3.339074in}{0.609897in}}%
\pgfpathlineto{\pgfqpoint{3.341171in}{0.609322in}}%
\pgfpathlineto{\pgfqpoint{3.343269in}{0.602827in}}%
\pgfpathlineto{\pgfqpoint{3.347464in}{0.582531in}}%
\pgfpathlineto{\pgfqpoint{3.351658in}{0.550875in}}%
\pgfpathlineto{\pgfqpoint{3.353756in}{0.532092in}}%
\pgfpathlineto{\pgfqpoint{3.355853in}{0.519726in}}%
\pgfpathlineto{\pgfqpoint{3.360048in}{0.509710in}}%
\pgfpathlineto{\pgfqpoint{3.362145in}{0.510666in}}%
\pgfpathlineto{\pgfqpoint{3.366340in}{0.547149in}}%
\pgfpathlineto{\pgfqpoint{3.368437in}{0.560154in}}%
\pgfpathlineto{\pgfqpoint{3.370535in}{0.568359in}}%
\pgfpathlineto{\pgfqpoint{3.372632in}{0.571920in}}%
\pgfpathlineto{\pgfqpoint{3.374729in}{0.569750in}}%
\pgfpathlineto{\pgfqpoint{3.376827in}{0.558364in}}%
\pgfpathlineto{\pgfqpoint{3.378924in}{0.527714in}}%
\pgfpathlineto{\pgfqpoint{3.381021in}{0.515297in}}%
\pgfpathlineto{\pgfqpoint{3.383119in}{0.559669in}}%
\pgfpathlineto{\pgfqpoint{3.385216in}{0.579374in}}%
\pgfpathlineto{\pgfqpoint{3.387313in}{0.591014in}}%
\pgfpathlineto{\pgfqpoint{3.389411in}{0.598272in}}%
\pgfpathlineto{\pgfqpoint{3.391508in}{0.602498in}}%
\pgfpathlineto{\pgfqpoint{3.393605in}{0.604253in}}%
\pgfpathlineto{\pgfqpoint{3.395703in}{0.603694in}}%
\pgfpathlineto{\pgfqpoint{3.397800in}{0.600654in}}%
\pgfpathlineto{\pgfqpoint{3.399898in}{0.594528in}}%
\pgfpathlineto{\pgfqpoint{3.401995in}{0.583780in}}%
\pgfpathlineto{\pgfqpoint{3.404092in}{0.563755in}}%
\pgfpathlineto{\pgfqpoint{3.406190in}{0.500996in}}%
\pgfpathlineto{\pgfqpoint{3.408287in}{0.549974in}}%
\pgfpathlineto{\pgfqpoint{3.410384in}{0.580971in}}%
\pgfpathlineto{\pgfqpoint{3.412482in}{0.596975in}}%
\pgfpathlineto{\pgfqpoint{3.414579in}{0.606787in}}%
\pgfpathlineto{\pgfqpoint{3.416676in}{0.612611in}}%
\pgfpathlineto{\pgfqpoint{3.418774in}{0.615032in}}%
\pgfpathlineto{\pgfqpoint{3.420871in}{0.613712in}}%
\pgfpathlineto{\pgfqpoint{3.422969in}{0.607214in}}%
\pgfpathlineto{\pgfqpoint{3.425066in}{0.593911in}}%
\pgfpathlineto{\pgfqpoint{3.427163in}{0.567391in}}%
\pgfpathlineto{\pgfqpoint{3.429261in}{0.502924in}}%
\pgfpathlineto{\pgfqpoint{3.431358in}{0.571934in}}%
\pgfpathlineto{\pgfqpoint{3.433455in}{0.593472in}}%
\pgfpathlineto{\pgfqpoint{3.435553in}{0.604273in}}%
\pgfpathlineto{\pgfqpoint{3.437650in}{0.609648in}}%
\pgfpathlineto{\pgfqpoint{3.439747in}{0.611134in}}%
\pgfpathlineto{\pgfqpoint{3.441845in}{0.609043in}}%
\pgfpathlineto{\pgfqpoint{3.443942in}{0.602791in}}%
\pgfpathlineto{\pgfqpoint{3.446040in}{0.590197in}}%
\pgfpathlineto{\pgfqpoint{3.448137in}{0.562611in}}%
\pgfpathlineto{\pgfqpoint{3.450234in}{0.513985in}}%
\pgfpathlineto{\pgfqpoint{3.452332in}{0.581026in}}%
\pgfpathlineto{\pgfqpoint{3.454429in}{0.603461in}}%
\pgfpathlineto{\pgfqpoint{3.456526in}{0.616021in}}%
\pgfpathlineto{\pgfqpoint{3.458624in}{0.623169in}}%
\pgfpathlineto{\pgfqpoint{3.460721in}{0.625954in}}%
\pgfpathlineto{\pgfqpoint{3.462818in}{0.623864in}}%
\pgfpathlineto{\pgfqpoint{3.464916in}{0.615381in}}%
\pgfpathlineto{\pgfqpoint{3.467013in}{0.598925in}}%
\pgfpathlineto{\pgfqpoint{3.469111in}{0.562162in}}%
\pgfpathlineto{\pgfqpoint{3.471208in}{0.549868in}}%
\pgfpathlineto{\pgfqpoint{3.473305in}{0.592803in}}%
\pgfpathlineto{\pgfqpoint{3.475403in}{0.609707in}}%
\pgfpathlineto{\pgfqpoint{3.477500in}{0.618443in}}%
\pgfpathlineto{\pgfqpoint{3.479597in}{0.622462in}}%
\pgfpathlineto{\pgfqpoint{3.481695in}{0.622797in}}%
\pgfpathlineto{\pgfqpoint{3.483792in}{0.619478in}}%
\pgfpathlineto{\pgfqpoint{3.485889in}{0.611556in}}%
\pgfpathlineto{\pgfqpoint{3.487987in}{0.595801in}}%
\pgfpathlineto{\pgfqpoint{3.490084in}{0.555944in}}%
\pgfpathlineto{\pgfqpoint{3.492182in}{0.561319in}}%
\pgfpathlineto{\pgfqpoint{3.494279in}{0.601803in}}%
\pgfpathlineto{\pgfqpoint{3.496376in}{0.620116in}}%
\pgfpathlineto{\pgfqpoint{3.498474in}{0.630702in}}%
\pgfpathlineto{\pgfqpoint{3.500571in}{0.636418in}}%
\pgfpathlineto{\pgfqpoint{3.502668in}{0.637768in}}%
\pgfpathlineto{\pgfqpoint{3.504766in}{0.633711in}}%
\pgfpathlineto{\pgfqpoint{3.506863in}{0.623363in}}%
\pgfpathlineto{\pgfqpoint{3.508960in}{0.604210in}}%
\pgfpathlineto{\pgfqpoint{3.511058in}{0.556888in}}%
\pgfpathlineto{\pgfqpoint{3.513155in}{0.571836in}}%
\pgfpathlineto{\pgfqpoint{3.515253in}{0.605582in}}%
\pgfpathlineto{\pgfqpoint{3.517350in}{0.620020in}}%
\pgfpathlineto{\pgfqpoint{3.519447in}{0.627225in}}%
\pgfpathlineto{\pgfqpoint{3.521545in}{0.629889in}}%
\pgfpathlineto{\pgfqpoint{3.523642in}{0.628713in}}%
\pgfpathlineto{\pgfqpoint{3.525739in}{0.623346in}}%
\pgfpathlineto{\pgfqpoint{3.527837in}{0.611944in}}%
\pgfpathlineto{\pgfqpoint{3.529934in}{0.587399in}}%
\pgfpathlineto{\pgfqpoint{3.532031in}{0.502763in}}%
\pgfpathlineto{\pgfqpoint{3.534129in}{0.597350in}}%
\pgfpathlineto{\pgfqpoint{3.536226in}{0.621716in}}%
\pgfpathlineto{\pgfqpoint{3.538324in}{0.635057in}}%
\pgfpathlineto{\pgfqpoint{3.540421in}{0.642658in}}%
\pgfpathlineto{\pgfqpoint{3.542518in}{0.645751in}}%
\pgfpathlineto{\pgfqpoint{3.544616in}{0.643890in}}%
\pgfpathlineto{\pgfqpoint{3.546713in}{0.636122in}}%
\pgfpathlineto{\pgfqpoint{3.548810in}{0.623174in}}%
\pgfpathlineto{\pgfqpoint{3.550908in}{0.601207in}}%
\pgfpathlineto{\pgfqpoint{3.553005in}{0.545924in}}%
\pgfpathlineto{\pgfqpoint{3.557200in}{0.601275in}}%
\pgfpathlineto{\pgfqpoint{3.559297in}{0.612908in}}%
\pgfpathlineto{\pgfqpoint{3.561395in}{0.617158in}}%
\pgfpathlineto{\pgfqpoint{3.563492in}{0.616040in}}%
\pgfpathlineto{\pgfqpoint{3.565589in}{0.609182in}}%
\pgfpathlineto{\pgfqpoint{3.567687in}{0.593084in}}%
\pgfpathlineto{\pgfqpoint{3.569784in}{0.545060in}}%
\pgfpathlineto{\pgfqpoint{3.571881in}{0.572103in}}%
\pgfpathlineto{\pgfqpoint{3.573979in}{0.607445in}}%
\pgfpathlineto{\pgfqpoint{3.576076in}{0.624540in}}%
\pgfpathlineto{\pgfqpoint{3.578173in}{0.634282in}}%
\pgfpathlineto{\pgfqpoint{3.580271in}{0.638925in}}%
\pgfpathlineto{\pgfqpoint{3.582368in}{0.638441in}}%
\pgfpathlineto{\pgfqpoint{3.584465in}{0.631747in}}%
\pgfpathlineto{\pgfqpoint{3.586563in}{0.620632in}}%
\pgfpathlineto{\pgfqpoint{3.588660in}{0.603365in}}%
\pgfpathlineto{\pgfqpoint{3.590758in}{0.572081in}}%
\pgfpathlineto{\pgfqpoint{3.592855in}{0.503805in}}%
\pgfpathlineto{\pgfqpoint{3.594952in}{0.574768in}}%
\pgfpathlineto{\pgfqpoint{3.597050in}{0.591901in}}%
\pgfpathlineto{\pgfqpoint{3.599147in}{0.597862in}}%
\pgfpathlineto{\pgfqpoint{3.601244in}{0.596927in}}%
\pgfpathlineto{\pgfqpoint{3.603342in}{0.588808in}}%
\pgfpathlineto{\pgfqpoint{3.605439in}{0.567802in}}%
\pgfpathlineto{\pgfqpoint{3.607536in}{0.419323in}}%
\pgfpathlineto{\pgfqpoint{3.609634in}{0.574295in}}%
\pgfpathlineto{\pgfqpoint{3.611731in}{0.600739in}}%
\pgfpathlineto{\pgfqpoint{3.613829in}{0.614975in}}%
\pgfpathlineto{\pgfqpoint{3.615926in}{0.622917in}}%
\pgfpathlineto{\pgfqpoint{3.618023in}{0.625791in}}%
\pgfpathlineto{\pgfqpoint{3.620121in}{0.622749in}}%
\pgfpathlineto{\pgfqpoint{3.622218in}{0.613729in}}%
\pgfpathlineto{\pgfqpoint{3.624315in}{0.598897in}}%
\pgfpathlineto{\pgfqpoint{3.626413in}{0.572244in}}%
\pgfpathlineto{\pgfqpoint{3.628510in}{0.440571in}}%
\pgfpathlineto{\pgfqpoint{3.630607in}{0.565623in}}%
\pgfpathlineto{\pgfqpoint{3.632705in}{0.587282in}}%
\pgfpathlineto{\pgfqpoint{3.634802in}{0.596485in}}%
\pgfpathlineto{\pgfqpoint{3.636900in}{0.599325in}}%
\pgfpathlineto{\pgfqpoint{3.638997in}{0.597030in}}%
\pgfpathlineto{\pgfqpoint{3.641094in}{0.588756in}}%
\pgfpathlineto{\pgfqpoint{3.643192in}{0.569843in}}%
\pgfpathlineto{\pgfqpoint{3.645289in}{0.498110in}}%
\pgfpathlineto{\pgfqpoint{3.647386in}{0.562881in}}%
\pgfpathlineto{\pgfqpoint{3.649484in}{0.592724in}}%
\pgfpathlineto{\pgfqpoint{3.651581in}{0.608103in}}%
\pgfpathlineto{\pgfqpoint{3.653678in}{0.616829in}}%
\pgfpathlineto{\pgfqpoint{3.655776in}{0.620607in}}%
\pgfpathlineto{\pgfqpoint{3.657873in}{0.619088in}}%
\pgfpathlineto{\pgfqpoint{3.659971in}{0.610865in}}%
\pgfpathlineto{\pgfqpoint{3.662068in}{0.595107in}}%
\pgfpathlineto{\pgfqpoint{3.666263in}{0.533151in}}%
\pgfpathlineto{\pgfqpoint{3.668360in}{0.583796in}}%
\pgfpathlineto{\pgfqpoint{3.670457in}{0.601744in}}%
\pgfpathlineto{\pgfqpoint{3.672555in}{0.610810in}}%
\pgfpathlineto{\pgfqpoint{3.674652in}{0.614913in}}%
\pgfpathlineto{\pgfqpoint{3.676749in}{0.615182in}}%
\pgfpathlineto{\pgfqpoint{3.678847in}{0.611642in}}%
\pgfpathlineto{\pgfqpoint{3.680944in}{0.603236in}}%
\pgfpathlineto{\pgfqpoint{3.683042in}{0.586264in}}%
\pgfpathlineto{\pgfqpoint{3.685139in}{0.539348in}}%
\pgfpathlineto{\pgfqpoint{3.687236in}{0.561759in}}%
\pgfpathlineto{\pgfqpoint{3.689334in}{0.597621in}}%
\pgfpathlineto{\pgfqpoint{3.691431in}{0.614783in}}%
\pgfpathlineto{\pgfqpoint{3.693528in}{0.624721in}}%
\pgfpathlineto{\pgfqpoint{3.695626in}{0.629874in}}%
\pgfpathlineto{\pgfqpoint{3.697723in}{0.630541in}}%
\pgfpathlineto{\pgfqpoint{3.699820in}{0.625437in}}%
\pgfpathlineto{\pgfqpoint{3.701918in}{0.613387in}}%
\pgfpathlineto{\pgfqpoint{3.704015in}{0.589356in}}%
\pgfpathlineto{\pgfqpoint{3.706113in}{0.466424in}}%
\pgfpathlineto{\pgfqpoint{3.708210in}{0.585353in}}%
\pgfpathlineto{\pgfqpoint{3.710307in}{0.609684in}}%
\pgfpathlineto{\pgfqpoint{3.712405in}{0.621634in}}%
\pgfpathlineto{\pgfqpoint{3.714502in}{0.627804in}}%
\pgfpathlineto{\pgfqpoint{3.716599in}{0.630022in}}%
\pgfpathlineto{\pgfqpoint{3.718697in}{0.628764in}}%
\pgfpathlineto{\pgfqpoint{3.720794in}{0.623676in}}%
\pgfpathlineto{\pgfqpoint{3.722891in}{0.613174in}}%
\pgfpathlineto{\pgfqpoint{3.724989in}{0.591713in}}%
\pgfpathlineto{\pgfqpoint{3.727086in}{0.494857in}}%
\pgfpathlineto{\pgfqpoint{3.729184in}{0.589266in}}%
\pgfpathlineto{\pgfqpoint{3.731281in}{0.616863in}}%
\pgfpathlineto{\pgfqpoint{3.733378in}{0.631476in}}%
\pgfpathlineto{\pgfqpoint{3.735476in}{0.639983in}}%
\pgfpathlineto{\pgfqpoint{3.737573in}{0.644038in}}%
\pgfpathlineto{\pgfqpoint{3.739670in}{0.643585in}}%
\pgfpathlineto{\pgfqpoint{3.741768in}{0.637294in}}%
\pgfpathlineto{\pgfqpoint{3.743865in}{0.625279in}}%
\pgfpathlineto{\pgfqpoint{3.745962in}{0.604046in}}%
\pgfpathlineto{\pgfqpoint{3.748060in}{0.548370in}}%
\pgfpathlineto{\pgfqpoint{3.752254in}{0.607107in}}%
\pgfpathlineto{\pgfqpoint{3.754352in}{0.619762in}}%
\pgfpathlineto{\pgfqpoint{3.756449in}{0.625369in}}%
\pgfpathlineto{\pgfqpoint{3.758547in}{0.626188in}}%
\pgfpathlineto{\pgfqpoint{3.760644in}{0.622491in}}%
\pgfpathlineto{\pgfqpoint{3.762741in}{0.612976in}}%
\pgfpathlineto{\pgfqpoint{3.764839in}{0.592264in}}%
\pgfpathlineto{\pgfqpoint{3.766936in}{0.497428in}}%
\pgfpathlineto{\pgfqpoint{3.769033in}{0.590376in}}%
\pgfpathlineto{\pgfqpoint{3.771131in}{0.618224in}}%
\pgfpathlineto{\pgfqpoint{3.773228in}{0.632878in}}%
\pgfpathlineto{\pgfqpoint{3.775325in}{0.641215in}}%
\pgfpathlineto{\pgfqpoint{3.777423in}{0.644782in}}%
\pgfpathlineto{\pgfqpoint{3.779520in}{0.643222in}}%
\pgfpathlineto{\pgfqpoint{3.781618in}{0.635733in}}%
\pgfpathlineto{\pgfqpoint{3.783715in}{0.623845in}}%
\pgfpathlineto{\pgfqpoint{3.785812in}{0.605156in}}%
\pgfpathlineto{\pgfqpoint{3.790007in}{0.535479in}}%
\pgfpathlineto{\pgfqpoint{3.792104in}{0.584474in}}%
\pgfpathlineto{\pgfqpoint{3.794202in}{0.599125in}}%
\pgfpathlineto{\pgfqpoint{3.796299in}{0.603897in}}%
\pgfpathlineto{\pgfqpoint{3.798396in}{0.601986in}}%
\pgfpathlineto{\pgfqpoint{3.800494in}{0.592579in}}%
\pgfpathlineto{\pgfqpoint{3.802591in}{0.568313in}}%
\pgfpathlineto{\pgfqpoint{3.804689in}{0.507049in}}%
\pgfpathlineto{\pgfqpoint{3.806786in}{0.584966in}}%
\pgfpathlineto{\pgfqpoint{3.808883in}{0.608851in}}%
\pgfpathlineto{\pgfqpoint{3.810981in}{0.622018in}}%
\pgfpathlineto{\pgfqpoint{3.813078in}{0.629221in}}%
\pgfpathlineto{\pgfqpoint{3.815175in}{0.631363in}}%
\pgfpathlineto{\pgfqpoint{3.817273in}{0.627358in}}%
\pgfpathlineto{\pgfqpoint{3.819370in}{0.618067in}}%
\pgfpathlineto{\pgfqpoint{3.821467in}{0.603604in}}%
\pgfpathlineto{\pgfqpoint{3.823565in}{0.579318in}}%
\pgfpathlineto{\pgfqpoint{3.825662in}{0.510235in}}%
\pgfpathlineto{\pgfqpoint{3.827760in}{0.556852in}}%
\pgfpathlineto{\pgfqpoint{3.829857in}{0.581902in}}%
\pgfpathlineto{\pgfqpoint{3.831954in}{0.591312in}}%
\pgfpathlineto{\pgfqpoint{3.834052in}{0.593258in}}%
\pgfpathlineto{\pgfqpoint{3.836149in}{0.588918in}}%
\pgfpathlineto{\pgfqpoint{3.838246in}{0.576113in}}%
\pgfpathlineto{\pgfqpoint{3.840344in}{0.541214in}}%
\pgfpathlineto{\pgfqpoint{3.842441in}{0.537519in}}%
\pgfpathlineto{\pgfqpoint{3.844538in}{0.582329in}}%
\pgfpathlineto{\pgfqpoint{3.846636in}{0.601841in}}%
\pgfpathlineto{\pgfqpoint{3.848733in}{0.612902in}}%
\pgfpathlineto{\pgfqpoint{3.850831in}{0.618535in}}%
\pgfpathlineto{\pgfqpoint{3.852928in}{0.619023in}}%
\pgfpathlineto{\pgfqpoint{3.855025in}{0.612984in}}%
\pgfpathlineto{\pgfqpoint{3.857123in}{0.600766in}}%
\pgfpathlineto{\pgfqpoint{3.859220in}{0.578494in}}%
\pgfpathlineto{\pgfqpoint{3.861317in}{0.511463in}}%
\pgfpathlineto{\pgfqpoint{3.863415in}{0.559917in}}%
\pgfpathlineto{\pgfqpoint{3.865512in}{0.587278in}}%
\pgfpathlineto{\pgfqpoint{3.867609in}{0.599465in}}%
\pgfpathlineto{\pgfqpoint{3.869707in}{0.605079in}}%
\pgfpathlineto{\pgfqpoint{3.871804in}{0.606180in}}%
\pgfpathlineto{\pgfqpoint{3.873902in}{0.603081in}}%
\pgfpathlineto{\pgfqpoint{3.875999in}{0.594780in}}%
\pgfpathlineto{\pgfqpoint{3.878096in}{0.577330in}}%
\pgfpathlineto{\pgfqpoint{3.880194in}{0.525279in}}%
\pgfpathlineto{\pgfqpoint{3.884388in}{0.592172in}}%
\pgfpathlineto{\pgfqpoint{3.886486in}{0.608750in}}%
\pgfpathlineto{\pgfqpoint{3.888583in}{0.618283in}}%
\pgfpathlineto{\pgfqpoint{3.890680in}{0.622965in}}%
\pgfpathlineto{\pgfqpoint{3.892778in}{0.622880in}}%
\pgfpathlineto{\pgfqpoint{3.894875in}{0.616511in}}%
\pgfpathlineto{\pgfqpoint{3.896973in}{0.602499in}}%
\pgfpathlineto{\pgfqpoint{3.899070in}{0.572511in}}%
\pgfpathlineto{\pgfqpoint{3.901167in}{0.526827in}}%
\pgfpathlineto{\pgfqpoint{3.903265in}{0.588499in}}%
\pgfpathlineto{\pgfqpoint{3.905362in}{0.608403in}}%
\pgfpathlineto{\pgfqpoint{3.907459in}{0.618751in}}%
\pgfpathlineto{\pgfqpoint{3.909557in}{0.624062in}}%
\pgfpathlineto{\pgfqpoint{3.911654in}{0.625716in}}%
\pgfpathlineto{\pgfqpoint{3.913751in}{0.624023in}}%
\pgfpathlineto{\pgfqpoint{3.915849in}{0.618528in}}%
\pgfpathlineto{\pgfqpoint{3.917946in}{0.607522in}}%
\pgfpathlineto{\pgfqpoint{3.920043in}{0.584988in}}%
\pgfpathlineto{\pgfqpoint{3.922141in}{0.445185in}}%
\pgfpathlineto{\pgfqpoint{3.924238in}{0.586346in}}%
\pgfpathlineto{\pgfqpoint{3.926336in}{0.612791in}}%
\pgfpathlineto{\pgfqpoint{3.928433in}{0.626999in}}%
\pgfpathlineto{\pgfqpoint{3.930530in}{0.635294in}}%
\pgfpathlineto{\pgfqpoint{3.932628in}{0.639224in}}%
\pgfpathlineto{\pgfqpoint{3.934725in}{0.638710in}}%
\pgfpathlineto{\pgfqpoint{3.936822in}{0.632323in}}%
\pgfpathlineto{\pgfqpoint{3.938920in}{0.619277in}}%
\pgfpathlineto{\pgfqpoint{3.941017in}{0.593894in}}%
\pgfpathlineto{\pgfqpoint{3.943114in}{0.412802in}}%
\pgfpathlineto{\pgfqpoint{3.945212in}{0.591952in}}%
\pgfpathlineto{\pgfqpoint{3.947309in}{0.614976in}}%
\pgfpathlineto{\pgfqpoint{3.949407in}{0.626096in}}%
\pgfpathlineto{\pgfqpoint{3.951504in}{0.631464in}}%
\pgfpathlineto{\pgfqpoint{3.953601in}{0.632749in}}%
\pgfpathlineto{\pgfqpoint{3.955699in}{0.630247in}}%
\pgfpathlineto{\pgfqpoint{3.957796in}{0.623229in}}%
\pgfpathlineto{\pgfqpoint{3.959893in}{0.608958in}}%
\pgfpathlineto{\pgfqpoint{3.961991in}{0.574952in}}%
\pgfpathlineto{\pgfqpoint{3.964088in}{0.561970in}}%
\pgfpathlineto{\pgfqpoint{3.966185in}{0.609455in}}%
\pgfpathlineto{\pgfqpoint{3.968283in}{0.629202in}}%
\pgfpathlineto{\pgfqpoint{3.970380in}{0.640476in}}%
\pgfpathlineto{\pgfqpoint{3.972478in}{0.646639in}}%
\pgfpathlineto{\pgfqpoint{3.974575in}{0.648348in}}%
\pgfpathlineto{\pgfqpoint{3.976672in}{0.644648in}}%
\pgfpathlineto{\pgfqpoint{3.978770in}{0.635606in}}%
\pgfpathlineto{\pgfqpoint{3.980867in}{0.621208in}}%
\pgfpathlineto{\pgfqpoint{3.982964in}{0.596495in}}%
\pgfpathlineto{\pgfqpoint{3.985062in}{0.518536in}}%
\pgfpathlineto{\pgfqpoint{3.987159in}{0.578704in}}%
\pgfpathlineto{\pgfqpoint{3.989256in}{0.602898in}}%
\pgfpathlineto{\pgfqpoint{3.991354in}{0.612484in}}%
\pgfpathlineto{\pgfqpoint{3.993451in}{0.615029in}}%
\pgfpathlineto{\pgfqpoint{3.995549in}{0.611846in}}%
\pgfpathlineto{\pgfqpoint{3.997646in}{0.601573in}}%
\pgfpathlineto{\pgfqpoint{3.999743in}{0.576502in}}%
\pgfpathlineto{\pgfqpoint{4.001841in}{0.517924in}}%
\pgfpathlineto{\pgfqpoint{4.003938in}{0.593312in}}%
\pgfpathlineto{\pgfqpoint{4.006035in}{0.616821in}}%
\pgfpathlineto{\pgfqpoint{4.008133in}{0.629778in}}%
\pgfpathlineto{\pgfqpoint{4.010230in}{0.636863in}}%
\pgfpathlineto{\pgfqpoint{4.012327in}{0.638981in}}%
\pgfpathlineto{\pgfqpoint{4.014425in}{0.635097in}}%
\pgfpathlineto{\pgfqpoint{4.016522in}{0.626149in}}%
\pgfpathlineto{\pgfqpoint{4.018620in}{0.612699in}}%
\pgfpathlineto{\pgfqpoint{4.020717in}{0.591479in}}%
\pgfpathlineto{\pgfqpoint{4.022814in}{0.546718in}}%
\pgfpathlineto{\pgfqpoint{4.024912in}{0.543493in}}%
\pgfpathlineto{\pgfqpoint{4.027009in}{0.578267in}}%
\pgfpathlineto{\pgfqpoint{4.029106in}{0.589469in}}%
\pgfpathlineto{\pgfqpoint{4.031204in}{0.591539in}}%
\pgfpathlineto{\pgfqpoint{4.033301in}{0.586129in}}%
\pgfpathlineto{\pgfqpoint{4.035398in}{0.569754in}}%
\pgfpathlineto{\pgfqpoint{4.037496in}{0.509004in}}%
\pgfpathlineto{\pgfqpoint{4.039593in}{0.560509in}}%
\pgfpathlineto{\pgfqpoint{4.041691in}{0.592474in}}%
\pgfpathlineto{\pgfqpoint{4.043788in}{0.608776in}}%
\pgfpathlineto{\pgfqpoint{4.045885in}{0.618051in}}%
\pgfpathlineto{\pgfqpoint{4.047983in}{0.622122in}}%
\pgfpathlineto{\pgfqpoint{4.050080in}{0.620586in}}%
\pgfpathlineto{\pgfqpoint{4.052177in}{0.612596in}}%
\pgfpathlineto{\pgfqpoint{4.054275in}{0.599172in}}%
\pgfpathlineto{\pgfqpoint{4.056372in}{0.575854in}}%
\pgfpathlineto{\pgfqpoint{4.058469in}{0.507817in}}%
\pgfpathlineto{\pgfqpoint{4.060567in}{0.555291in}}%
\pgfpathlineto{\pgfqpoint{4.062664in}{0.581560in}}%
\pgfpathlineto{\pgfqpoint{4.064762in}{0.592465in}}%
\pgfpathlineto{\pgfqpoint{4.066859in}{0.596455in}}%
\pgfpathlineto{\pgfqpoint{4.068956in}{0.595296in}}%
\pgfpathlineto{\pgfqpoint{4.071054in}{0.588604in}}%
\pgfpathlineto{\pgfqpoint{4.073151in}{0.573123in}}%
\pgfpathlineto{\pgfqpoint{4.075248in}{0.529443in}}%
\pgfpathlineto{\pgfqpoint{4.077346in}{0.546972in}}%
\pgfpathlineto{\pgfqpoint{4.079443in}{0.584494in}}%
\pgfpathlineto{\pgfqpoint{4.081540in}{0.602158in}}%
\pgfpathlineto{\pgfqpoint{4.083638in}{0.612227in}}%
\pgfpathlineto{\pgfqpoint{4.085735in}{0.617169in}}%
\pgfpathlineto{\pgfqpoint{4.087832in}{0.617086in}}%
\pgfpathlineto{\pgfqpoint{4.089930in}{0.610450in}}%
\pgfpathlineto{\pgfqpoint{4.092027in}{0.596344in}}%
\pgfpathlineto{\pgfqpoint{4.094125in}{0.566767in}}%
\pgfpathlineto{\pgfqpoint{4.096222in}{0.515694in}}%
\pgfpathlineto{\pgfqpoint{4.098319in}{0.580527in}}%
\pgfpathlineto{\pgfqpoint{4.100417in}{0.600564in}}%
\pgfpathlineto{\pgfqpoint{4.102514in}{0.610822in}}%
\pgfpathlineto{\pgfqpoint{4.104611in}{0.615948in}}%
\pgfpathlineto{\pgfqpoint{4.106709in}{0.617324in}}%
\pgfpathlineto{\pgfqpoint{4.108806in}{0.615217in}}%
\pgfpathlineto{\pgfqpoint{4.110903in}{0.609055in}}%
\pgfpathlineto{\pgfqpoint{4.113001in}{0.596747in}}%
\pgfpathlineto{\pgfqpoint{4.115098in}{0.570268in}}%
\pgfpathlineto{\pgfqpoint{4.117196in}{0.510195in}}%
\pgfpathlineto{\pgfqpoint{4.119293in}{0.585220in}}%
\pgfpathlineto{\pgfqpoint{4.121390in}{0.608348in}}%
\pgfpathlineto{\pgfqpoint{4.123488in}{0.621215in}}%
\pgfpathlineto{\pgfqpoint{4.125585in}{0.628598in}}%
\pgfpathlineto{\pgfqpoint{4.127682in}{0.631643in}}%
\pgfpathlineto{\pgfqpoint{4.129780in}{0.629948in}}%
\pgfpathlineto{\pgfqpoint{4.131877in}{0.621821in}}%
\pgfpathlineto{\pgfqpoint{4.133974in}{0.605157in}}%
\pgfpathlineto{\pgfqpoint{4.136072in}{0.565437in}}%
\pgfpathlineto{\pgfqpoint{4.138169in}{0.563209in}}%
\pgfpathlineto{\pgfqpoint{4.140267in}{0.603425in}}%
\pgfpathlineto{\pgfqpoint{4.142364in}{0.620031in}}%
\pgfpathlineto{\pgfqpoint{4.144461in}{0.628883in}}%
\pgfpathlineto{\pgfqpoint{4.146559in}{0.633223in}}%
\pgfpathlineto{\pgfqpoint{4.148656in}{0.634076in}}%
\pgfpathlineto{\pgfqpoint{4.150753in}{0.631569in}}%
\pgfpathlineto{\pgfqpoint{4.152851in}{0.625040in}}%
\pgfpathlineto{\pgfqpoint{4.154948in}{0.612276in}}%
\pgfpathlineto{\pgfqpoint{4.157045in}{0.584697in}}%
\pgfpathlineto{\pgfqpoint{4.159143in}{0.534340in}}%
\pgfpathlineto{\pgfqpoint{4.161240in}{0.602454in}}%
\pgfpathlineto{\pgfqpoint{4.163338in}{0.624967in}}%
\pgfpathlineto{\pgfqpoint{4.165435in}{0.637596in}}%
\pgfpathlineto{\pgfqpoint{4.167532in}{0.644869in}}%
\pgfpathlineto{\pgfqpoint{4.169630in}{0.647888in}}%
\pgfpathlineto{\pgfqpoint{4.171727in}{0.646269in}}%
\pgfpathlineto{\pgfqpoint{4.173824in}{0.638789in}}%
\pgfpathlineto{\pgfqpoint{4.175922in}{0.625600in}}%
\pgfpathlineto{\pgfqpoint{4.178019in}{0.602086in}}%
\pgfpathlineto{\pgfqpoint{4.180116in}{0.527289in}}%
\pgfpathlineto{\pgfqpoint{4.182214in}{0.585555in}}%
\pgfpathlineto{\pgfqpoint{4.184311in}{0.611219in}}%
\pgfpathlineto{\pgfqpoint{4.186409in}{0.622359in}}%
\pgfpathlineto{\pgfqpoint{4.188506in}{0.626902in}}%
\pgfpathlineto{\pgfqpoint{4.190603in}{0.626657in}}%
\pgfpathlineto{\pgfqpoint{4.192701in}{0.621562in}}%
\pgfpathlineto{\pgfqpoint{4.194798in}{0.609590in}}%
\pgfpathlineto{\pgfqpoint{4.196895in}{0.581660in}}%
\pgfpathlineto{\pgfqpoint{4.198993in}{0.541037in}}%
\pgfpathlineto{\pgfqpoint{4.201090in}{0.603586in}}%
\pgfpathlineto{\pgfqpoint{4.203187in}{0.625637in}}%
\pgfpathlineto{\pgfqpoint{4.205285in}{0.637931in}}%
\pgfpathlineto{\pgfqpoint{4.207382in}{0.644650in}}%
\pgfpathlineto{\pgfqpoint{4.209480in}{0.646618in}}%
\pgfpathlineto{\pgfqpoint{4.211577in}{0.642848in}}%
\pgfpathlineto{\pgfqpoint{4.213674in}{0.634060in}}%
\pgfpathlineto{\pgfqpoint{4.215772in}{0.620913in}}%
\pgfpathlineto{\pgfqpoint{4.217869in}{0.600414in}}%
\pgfpathlineto{\pgfqpoint{4.219966in}{0.559179in}}%
\pgfpathlineto{\pgfqpoint{4.222064in}{0.544369in}}%
\pgfpathlineto{\pgfqpoint{4.224161in}{0.583283in}}%
\pgfpathlineto{\pgfqpoint{4.226258in}{0.595229in}}%
\pgfpathlineto{\pgfqpoint{4.228356in}{0.597561in}}%
\pgfpathlineto{\pgfqpoint{4.230453in}{0.592204in}}%
\pgfpathlineto{\pgfqpoint{4.232551in}{0.575610in}}%
\pgfpathlineto{\pgfqpoint{4.234648in}{0.511390in}}%
\pgfpathlineto{\pgfqpoint{4.236745in}{0.568264in}}%
\pgfpathlineto{\pgfqpoint{4.238843in}{0.599653in}}%
\pgfpathlineto{\pgfqpoint{4.240940in}{0.615789in}}%
\pgfpathlineto{\pgfqpoint{4.243037in}{0.624966in}}%
\pgfpathlineto{\pgfqpoint{4.245135in}{0.628942in}}%
\pgfpathlineto{\pgfqpoint{4.247232in}{0.627262in}}%
\pgfpathlineto{\pgfqpoint{4.249329in}{0.619407in}}%
\pgfpathlineto{\pgfqpoint{4.251427in}{0.607122in}}%
\pgfpathlineto{\pgfqpoint{4.253524in}{0.587714in}}%
\pgfpathlineto{\pgfqpoint{4.255621in}{0.549150in}}%
\pgfpathlineto{\pgfqpoint{4.257719in}{0.527114in}}%
\pgfpathlineto{\pgfqpoint{4.259816in}{0.570246in}}%
\pgfpathlineto{\pgfqpoint{4.261914in}{0.583879in}}%
\pgfpathlineto{\pgfqpoint{4.264011in}{0.588067in}}%
\pgfpathlineto{\pgfqpoint{4.266108in}{0.585587in}}%
\pgfpathlineto{\pgfqpoint{4.268206in}{0.575296in}}%
\pgfpathlineto{\pgfqpoint{4.270303in}{0.548377in}}%
\pgfpathlineto{\pgfqpoint{4.272400in}{0.510019in}}%
\pgfpathlineto{\pgfqpoint{4.274498in}{0.572337in}}%
\pgfpathlineto{\pgfqpoint{4.276595in}{0.594714in}}%
\pgfpathlineto{\pgfqpoint{4.278692in}{0.607203in}}%
\pgfpathlineto{\pgfqpoint{4.280790in}{0.613914in}}%
\pgfpathlineto{\pgfqpoint{4.282887in}{0.615541in}}%
\pgfpathlineto{\pgfqpoint{4.284985in}{0.610956in}}%
\pgfpathlineto{\pgfqpoint{4.287082in}{0.602759in}}%
\pgfpathlineto{\pgfqpoint{4.291277in}{0.579072in}}%
\pgfpathlineto{\pgfqpoint{4.293374in}{0.562207in}}%
\pgfpathlineto{\pgfqpoint{4.295471in}{0.540236in}}%
\pgfpathlineto{\pgfqpoint{4.299666in}{0.477282in}}%
\pgfpathlineto{\pgfqpoint{4.301763in}{0.474537in}}%
\pgfpathlineto{\pgfqpoint{4.305958in}{0.531294in}}%
\pgfpathlineto{\pgfqpoint{4.308056in}{0.549927in}}%
\pgfpathlineto{\pgfqpoint{4.310153in}{0.562714in}}%
\pgfpathlineto{\pgfqpoint{4.312250in}{0.570735in}}%
\pgfpathlineto{\pgfqpoint{4.314348in}{0.574089in}}%
\pgfpathlineto{\pgfqpoint{4.316445in}{0.571509in}}%
\pgfpathlineto{\pgfqpoint{4.316445in}{0.571509in}}%
\pgfusepath{stroke}%
\end{pgfscope}%
\begin{pgfscope}%
\pgfpathrectangle{\pgfqpoint{2.554660in}{0.320625in}}{\pgfqpoint{1.845679in}{1.164375in}}%
\pgfusepath{clip}%
\pgfsetrectcap%
\pgfsetroundjoin%
\pgfsetlinewidth{1.154312pt}%
\definecolor{currentstroke}{rgb}{0.950210,0.390820,0.362468}%
\pgfsetstrokecolor{currentstroke}%
\pgfsetdash{}{0pt}%
\pgfpathmoveto{\pgfqpoint{2.638555in}{1.468525in}}%
\pgfpathlineto{\pgfqpoint{2.649042in}{1.449425in}}%
\pgfpathlineto{\pgfqpoint{2.659529in}{1.427545in}}%
\pgfpathlineto{\pgfqpoint{2.672113in}{1.398195in}}%
\pgfpathlineto{\pgfqpoint{2.714060in}{1.298024in}}%
\pgfpathlineto{\pgfqpoint{2.735034in}{1.251605in}}%
\pgfpathlineto{\pgfqpoint{2.758105in}{1.203355in}}%
\pgfpathlineto{\pgfqpoint{2.783273in}{1.153309in}}%
\pgfpathlineto{\pgfqpoint{2.812636in}{1.097504in}}%
\pgfpathlineto{\pgfqpoint{2.844097in}{1.040115in}}%
\pgfpathlineto{\pgfqpoint{2.879752in}{0.977452in}}%
\pgfpathlineto{\pgfqpoint{2.919602in}{0.909725in}}%
\pgfpathlineto{\pgfqpoint{2.965744in}{0.833485in}}%
\pgfpathlineto{\pgfqpoint{3.047541in}{0.701210in}}%
\pgfpathlineto{\pgfqpoint{3.066417in}{0.672950in}}%
\pgfpathlineto{\pgfqpoint{3.089488in}{0.632424in}}%
\pgfpathlineto{\pgfqpoint{3.093683in}{0.627903in}}%
\pgfpathlineto{\pgfqpoint{3.097877in}{0.624801in}}%
\pgfpathlineto{\pgfqpoint{3.106267in}{0.619444in}}%
\pgfpathlineto{\pgfqpoint{3.108364in}{0.617049in}}%
\pgfpathlineto{\pgfqpoint{3.112559in}{0.609696in}}%
\pgfpathlineto{\pgfqpoint{3.116754in}{0.599987in}}%
\pgfpathlineto{\pgfqpoint{3.120948in}{0.587840in}}%
\pgfpathlineto{\pgfqpoint{3.131435in}{0.554821in}}%
\pgfpathlineto{\pgfqpoint{3.133533in}{0.551543in}}%
\pgfpathlineto{\pgfqpoint{3.139825in}{0.545517in}}%
\pgfpathlineto{\pgfqpoint{3.141922in}{0.541986in}}%
\pgfpathlineto{\pgfqpoint{3.148214in}{0.525543in}}%
\pgfpathlineto{\pgfqpoint{3.150311in}{0.527760in}}%
\pgfpathlineto{\pgfqpoint{3.158701in}{0.565907in}}%
\pgfpathlineto{\pgfqpoint{3.162896in}{0.577179in}}%
\pgfpathlineto{\pgfqpoint{3.164993in}{0.580713in}}%
\pgfpathlineto{\pgfqpoint{3.167090in}{0.582880in}}%
\pgfpathlineto{\pgfqpoint{3.169188in}{0.583573in}}%
\pgfpathlineto{\pgfqpoint{3.171285in}{0.582524in}}%
\pgfpathlineto{\pgfqpoint{3.173382in}{0.579056in}}%
\pgfpathlineto{\pgfqpoint{3.175480in}{0.572621in}}%
\pgfpathlineto{\pgfqpoint{3.177577in}{0.574650in}}%
\pgfpathlineto{\pgfqpoint{3.183869in}{0.621376in}}%
\pgfpathlineto{\pgfqpoint{3.188064in}{0.640333in}}%
\pgfpathlineto{\pgfqpoint{3.192259in}{0.652584in}}%
\pgfpathlineto{\pgfqpoint{3.196453in}{0.660503in}}%
\pgfpathlineto{\pgfqpoint{3.200648in}{0.665184in}}%
\pgfpathlineto{\pgfqpoint{3.202746in}{0.666476in}}%
\pgfpathlineto{\pgfqpoint{3.204843in}{0.667106in}}%
\pgfpathlineto{\pgfqpoint{3.206940in}{0.667071in}}%
\pgfpathlineto{\pgfqpoint{3.209038in}{0.666350in}}%
\pgfpathlineto{\pgfqpoint{3.211135in}{0.664888in}}%
\pgfpathlineto{\pgfqpoint{3.213232in}{0.662598in}}%
\pgfpathlineto{\pgfqpoint{3.215330in}{0.659337in}}%
\pgfpathlineto{\pgfqpoint{3.217427in}{0.654874in}}%
\pgfpathlineto{\pgfqpoint{3.219524in}{0.648822in}}%
\pgfpathlineto{\pgfqpoint{3.221622in}{0.640478in}}%
\pgfpathlineto{\pgfqpoint{3.223719in}{0.628379in}}%
\pgfpathlineto{\pgfqpoint{3.225816in}{0.608707in}}%
\pgfpathlineto{\pgfqpoint{3.227914in}{0.568904in}}%
\pgfpathlineto{\pgfqpoint{3.230011in}{0.591948in}}%
\pgfpathlineto{\pgfqpoint{3.232109in}{0.621590in}}%
\pgfpathlineto{\pgfqpoint{3.234206in}{0.638452in}}%
\pgfpathlineto{\pgfqpoint{3.238401in}{0.658506in}}%
\pgfpathlineto{\pgfqpoint{3.242595in}{0.670157in}}%
\pgfpathlineto{\pgfqpoint{3.246790in}{0.676898in}}%
\pgfpathlineto{\pgfqpoint{3.248887in}{0.678749in}}%
\pgfpathlineto{\pgfqpoint{3.250985in}{0.679590in}}%
\pgfpathlineto{\pgfqpoint{3.253082in}{0.679330in}}%
\pgfpathlineto{\pgfqpoint{3.255180in}{0.677773in}}%
\pgfpathlineto{\pgfqpoint{3.257277in}{0.674550in}}%
\pgfpathlineto{\pgfqpoint{3.261472in}{0.662579in}}%
\pgfpathlineto{\pgfqpoint{3.274056in}{0.620213in}}%
\pgfpathlineto{\pgfqpoint{3.276153in}{0.616348in}}%
\pgfpathlineto{\pgfqpoint{3.278251in}{0.614367in}}%
\pgfpathlineto{\pgfqpoint{3.280348in}{0.614475in}}%
\pgfpathlineto{\pgfqpoint{3.284543in}{0.617909in}}%
\pgfpathlineto{\pgfqpoint{3.286640in}{0.618134in}}%
\pgfpathlineto{\pgfqpoint{3.288737in}{0.614689in}}%
\pgfpathlineto{\pgfqpoint{3.292932in}{0.600888in}}%
\pgfpathlineto{\pgfqpoint{3.301322in}{0.569147in}}%
\pgfpathlineto{\pgfqpoint{3.305516in}{0.558939in}}%
\pgfpathlineto{\pgfqpoint{3.307614in}{0.556681in}}%
\pgfpathlineto{\pgfqpoint{3.309711in}{0.556803in}}%
\pgfpathlineto{\pgfqpoint{3.316003in}{0.563188in}}%
\pgfpathlineto{\pgfqpoint{3.318100in}{0.561556in}}%
\pgfpathlineto{\pgfqpoint{3.320198in}{0.553606in}}%
\pgfpathlineto{\pgfqpoint{3.324393in}{0.521172in}}%
\pgfpathlineto{\pgfqpoint{3.326490in}{0.548929in}}%
\pgfpathlineto{\pgfqpoint{3.328587in}{0.568735in}}%
\pgfpathlineto{\pgfqpoint{3.330685in}{0.581172in}}%
\pgfpathlineto{\pgfqpoint{3.332782in}{0.589256in}}%
\pgfpathlineto{\pgfqpoint{3.334879in}{0.594333in}}%
\pgfpathlineto{\pgfqpoint{3.336977in}{0.597023in}}%
\pgfpathlineto{\pgfqpoint{3.339074in}{0.597574in}}%
\pgfpathlineto{\pgfqpoint{3.341171in}{0.595966in}}%
\pgfpathlineto{\pgfqpoint{3.343269in}{0.591890in}}%
\pgfpathlineto{\pgfqpoint{3.345366in}{0.584551in}}%
\pgfpathlineto{\pgfqpoint{3.347464in}{0.571961in}}%
\pgfpathlineto{\pgfqpoint{3.349561in}{0.547508in}}%
\pgfpathlineto{\pgfqpoint{3.351658in}{0.461959in}}%
\pgfpathlineto{\pgfqpoint{3.353756in}{0.552563in}}%
\pgfpathlineto{\pgfqpoint{3.355853in}{0.577759in}}%
\pgfpathlineto{\pgfqpoint{3.357950in}{0.591888in}}%
\pgfpathlineto{\pgfqpoint{3.360048in}{0.600755in}}%
\pgfpathlineto{\pgfqpoint{3.362145in}{0.606027in}}%
\pgfpathlineto{\pgfqpoint{3.364242in}{0.608129in}}%
\pgfpathlineto{\pgfqpoint{3.366340in}{0.606698in}}%
\pgfpathlineto{\pgfqpoint{3.368437in}{0.600266in}}%
\pgfpathlineto{\pgfqpoint{3.370535in}{0.586423in}}%
\pgfpathlineto{\pgfqpoint{3.374729in}{0.526809in}}%
\pgfpathlineto{\pgfqpoint{3.376827in}{0.574594in}}%
\pgfpathlineto{\pgfqpoint{3.378924in}{0.593880in}}%
\pgfpathlineto{\pgfqpoint{3.381021in}{0.604192in}}%
\pgfpathlineto{\pgfqpoint{3.383119in}{0.609642in}}%
\pgfpathlineto{\pgfqpoint{3.385216in}{0.611556in}}%
\pgfpathlineto{\pgfqpoint{3.387313in}{0.610272in}}%
\pgfpathlineto{\pgfqpoint{3.389411in}{0.605443in}}%
\pgfpathlineto{\pgfqpoint{3.391508in}{0.595691in}}%
\pgfpathlineto{\pgfqpoint{3.393605in}{0.576528in}}%
\pgfpathlineto{\pgfqpoint{3.395703in}{0.516264in}}%
\pgfpathlineto{\pgfqpoint{3.397800in}{0.562325in}}%
\pgfpathlineto{\pgfqpoint{3.399898in}{0.593755in}}%
\pgfpathlineto{\pgfqpoint{3.401995in}{0.609692in}}%
\pgfpathlineto{\pgfqpoint{3.404092in}{0.619084in}}%
\pgfpathlineto{\pgfqpoint{3.406190in}{0.624025in}}%
\pgfpathlineto{\pgfqpoint{3.408287in}{0.624784in}}%
\pgfpathlineto{\pgfqpoint{3.410384in}{0.620199in}}%
\pgfpathlineto{\pgfqpoint{3.412482in}{0.608454in}}%
\pgfpathlineto{\pgfqpoint{3.414579in}{0.584091in}}%
\pgfpathlineto{\pgfqpoint{3.416676in}{0.508786in}}%
\pgfpathlineto{\pgfqpoint{3.418774in}{0.585198in}}%
\pgfpathlineto{\pgfqpoint{3.420871in}{0.608639in}}%
\pgfpathlineto{\pgfqpoint{3.422969in}{0.620643in}}%
\pgfpathlineto{\pgfqpoint{3.425066in}{0.627148in}}%
\pgfpathlineto{\pgfqpoint{3.427163in}{0.629909in}}%
\pgfpathlineto{\pgfqpoint{3.429261in}{0.629461in}}%
\pgfpathlineto{\pgfqpoint{3.431358in}{0.625640in}}%
\pgfpathlineto{\pgfqpoint{3.433455in}{0.617427in}}%
\pgfpathlineto{\pgfqpoint{3.435553in}{0.601633in}}%
\pgfpathlineto{\pgfqpoint{3.437650in}{0.562804in}}%
\pgfpathlineto{\pgfqpoint{3.439747in}{0.564540in}}%
\pgfpathlineto{\pgfqpoint{3.441845in}{0.606108in}}%
\pgfpathlineto{\pgfqpoint{3.443942in}{0.624684in}}%
\pgfpathlineto{\pgfqpoint{3.446040in}{0.635498in}}%
\pgfpathlineto{\pgfqpoint{3.448137in}{0.641540in}}%
\pgfpathlineto{\pgfqpoint{3.450234in}{0.643452in}}%
\pgfpathlineto{\pgfqpoint{3.452332in}{0.640489in}}%
\pgfpathlineto{\pgfqpoint{3.454429in}{0.631323in}}%
\pgfpathlineto{\pgfqpoint{3.456526in}{0.614213in}}%
\pgfpathlineto{\pgfqpoint{3.458624in}{0.576037in}}%
\pgfpathlineto{\pgfqpoint{3.460721in}{0.567339in}}%
\pgfpathlineto{\pgfqpoint{3.462818in}{0.608433in}}%
\pgfpathlineto{\pgfqpoint{3.464916in}{0.624717in}}%
\pgfpathlineto{\pgfqpoint{3.467013in}{0.632952in}}%
\pgfpathlineto{\pgfqpoint{3.469111in}{0.636439in}}%
\pgfpathlineto{\pgfqpoint{3.471208in}{0.636113in}}%
\pgfpathlineto{\pgfqpoint{3.473305in}{0.631851in}}%
\pgfpathlineto{\pgfqpoint{3.475403in}{0.622315in}}%
\pgfpathlineto{\pgfqpoint{3.477500in}{0.602743in}}%
\pgfpathlineto{\pgfqpoint{3.479597in}{0.534386in}}%
\pgfpathlineto{\pgfqpoint{3.481695in}{0.593026in}}%
\pgfpathlineto{\pgfqpoint{3.483792in}{0.623100in}}%
\pgfpathlineto{\pgfqpoint{3.485889in}{0.638543in}}%
\pgfpathlineto{\pgfqpoint{3.487987in}{0.647490in}}%
\pgfpathlineto{\pgfqpoint{3.490084in}{0.651814in}}%
\pgfpathlineto{\pgfqpoint{3.492182in}{0.651520in}}%
\pgfpathlineto{\pgfqpoint{3.494279in}{0.645411in}}%
\pgfpathlineto{\pgfqpoint{3.496376in}{0.634766in}}%
\pgfpathlineto{\pgfqpoint{3.498474in}{0.618177in}}%
\pgfpathlineto{\pgfqpoint{3.500571in}{0.588597in}}%
\pgfpathlineto{\pgfqpoint{3.502668in}{0.507765in}}%
\pgfpathlineto{\pgfqpoint{3.504766in}{0.587066in}}%
\pgfpathlineto{\pgfqpoint{3.506863in}{0.605481in}}%
\pgfpathlineto{\pgfqpoint{3.508960in}{0.612257in}}%
\pgfpathlineto{\pgfqpoint{3.511058in}{0.612219in}}%
\pgfpathlineto{\pgfqpoint{3.513155in}{0.605527in}}%
\pgfpathlineto{\pgfqpoint{3.515253in}{0.588154in}}%
\pgfpathlineto{\pgfqpoint{3.517350in}{0.524772in}}%
\pgfpathlineto{\pgfqpoint{3.519447in}{0.579192in}}%
\pgfpathlineto{\pgfqpoint{3.521545in}{0.610428in}}%
\pgfpathlineto{\pgfqpoint{3.523642in}{0.626381in}}%
\pgfpathlineto{\pgfqpoint{3.525739in}{0.635432in}}%
\pgfpathlineto{\pgfqpoint{3.527837in}{0.639365in}}%
\pgfpathlineto{\pgfqpoint{3.529934in}{0.637791in}}%
\pgfpathlineto{\pgfqpoint{3.532031in}{0.630231in}}%
\pgfpathlineto{\pgfqpoint{3.534129in}{0.618864in}}%
\pgfpathlineto{\pgfqpoint{3.536226in}{0.601963in}}%
\pgfpathlineto{\pgfqpoint{3.538324in}{0.573440in}}%
\pgfpathlineto{\pgfqpoint{3.540421in}{0.467941in}}%
\pgfpathlineto{\pgfqpoint{3.542518in}{0.561447in}}%
\pgfpathlineto{\pgfqpoint{3.544616in}{0.580165in}}%
\pgfpathlineto{\pgfqpoint{3.546713in}{0.585279in}}%
\pgfpathlineto{\pgfqpoint{3.548810in}{0.581809in}}%
\pgfpathlineto{\pgfqpoint{3.550908in}{0.567514in}}%
\pgfpathlineto{\pgfqpoint{3.553005in}{0.517555in}}%
\pgfpathlineto{\pgfqpoint{3.557200in}{0.587456in}}%
\pgfpathlineto{\pgfqpoint{3.559297in}{0.604783in}}%
\pgfpathlineto{\pgfqpoint{3.561395in}{0.614715in}}%
\pgfpathlineto{\pgfqpoint{3.563492in}{0.619350in}}%
\pgfpathlineto{\pgfqpoint{3.565589in}{0.618456in}}%
\pgfpathlineto{\pgfqpoint{3.567687in}{0.610957in}}%
\pgfpathlineto{\pgfqpoint{3.569784in}{0.598062in}}%
\pgfpathlineto{\pgfqpoint{3.571881in}{0.575844in}}%
\pgfpathlineto{\pgfqpoint{3.573979in}{0.517059in}}%
\pgfpathlineto{\pgfqpoint{3.578173in}{0.578381in}}%
\pgfpathlineto{\pgfqpoint{3.580271in}{0.590020in}}%
\pgfpathlineto{\pgfqpoint{3.582368in}{0.594518in}}%
\pgfpathlineto{\pgfqpoint{3.584465in}{0.593876in}}%
\pgfpathlineto{\pgfqpoint{3.586563in}{0.587906in}}%
\pgfpathlineto{\pgfqpoint{3.588660in}{0.573895in}}%
\pgfpathlineto{\pgfqpoint{3.590758in}{0.537507in}}%
\pgfpathlineto{\pgfqpoint{3.592855in}{0.535755in}}%
\pgfpathlineto{\pgfqpoint{3.594952in}{0.579265in}}%
\pgfpathlineto{\pgfqpoint{3.597050in}{0.598333in}}%
\pgfpathlineto{\pgfqpoint{3.599147in}{0.609163in}}%
\pgfpathlineto{\pgfqpoint{3.601244in}{0.614730in}}%
\pgfpathlineto{\pgfqpoint{3.603342in}{0.615381in}}%
\pgfpathlineto{\pgfqpoint{3.605439in}{0.609668in}}%
\pgfpathlineto{\pgfqpoint{3.607536in}{0.596271in}}%
\pgfpathlineto{\pgfqpoint{3.609634in}{0.567667in}}%
\pgfpathlineto{\pgfqpoint{3.611731in}{0.512334in}}%
\pgfpathlineto{\pgfqpoint{3.613829in}{0.580618in}}%
\pgfpathlineto{\pgfqpoint{3.615926in}{0.601322in}}%
\pgfpathlineto{\pgfqpoint{3.618023in}{0.612105in}}%
\pgfpathlineto{\pgfqpoint{3.620121in}{0.617776in}}%
\pgfpathlineto{\pgfqpoint{3.622218in}{0.619802in}}%
\pgfpathlineto{\pgfqpoint{3.624315in}{0.618563in}}%
\pgfpathlineto{\pgfqpoint{3.626413in}{0.613719in}}%
\pgfpathlineto{\pgfqpoint{3.628510in}{0.603864in}}%
\pgfpathlineto{\pgfqpoint{3.630607in}{0.584350in}}%
\pgfpathlineto{\pgfqpoint{3.632705in}{0.520104in}}%
\pgfpathlineto{\pgfqpoint{3.634802in}{0.572189in}}%
\pgfpathlineto{\pgfqpoint{3.636900in}{0.602906in}}%
\pgfpathlineto{\pgfqpoint{3.638997in}{0.618598in}}%
\pgfpathlineto{\pgfqpoint{3.641094in}{0.627815in}}%
\pgfpathlineto{\pgfqpoint{3.643192in}{0.632567in}}%
\pgfpathlineto{\pgfqpoint{3.645289in}{0.633056in}}%
\pgfpathlineto{\pgfqpoint{3.647386in}{0.628005in}}%
\pgfpathlineto{\pgfqpoint{3.649484in}{0.615546in}}%
\pgfpathlineto{\pgfqpoint{3.651581in}{0.589132in}}%
\pgfpathlineto{\pgfqpoint{3.653678in}{0.509028in}}%
\pgfpathlineto{\pgfqpoint{3.655776in}{0.596344in}}%
\pgfpathlineto{\pgfqpoint{3.657873in}{0.618505in}}%
\pgfpathlineto{\pgfqpoint{3.659971in}{0.629971in}}%
\pgfpathlineto{\pgfqpoint{3.662068in}{0.636139in}}%
\pgfpathlineto{\pgfqpoint{3.664165in}{0.638632in}}%
\pgfpathlineto{\pgfqpoint{3.666263in}{0.637921in}}%
\pgfpathlineto{\pgfqpoint{3.668360in}{0.633786in}}%
\pgfpathlineto{\pgfqpoint{3.670457in}{0.625102in}}%
\pgfpathlineto{\pgfqpoint{3.672555in}{0.608321in}}%
\pgfpathlineto{\pgfqpoint{3.674652in}{0.564503in}}%
\pgfpathlineto{\pgfqpoint{3.676749in}{0.579225in}}%
\pgfpathlineto{\pgfqpoint{3.678847in}{0.616862in}}%
\pgfpathlineto{\pgfqpoint{3.680944in}{0.634513in}}%
\pgfpathlineto{\pgfqpoint{3.683042in}{0.644840in}}%
\pgfpathlineto{\pgfqpoint{3.685139in}{0.650506in}}%
\pgfpathlineto{\pgfqpoint{3.687236in}{0.652023in}}%
\pgfpathlineto{\pgfqpoint{3.689334in}{0.648495in}}%
\pgfpathlineto{\pgfqpoint{3.691431in}{0.639394in}}%
\pgfpathlineto{\pgfqpoint{3.693528in}{0.624151in}}%
\pgfpathlineto{\pgfqpoint{3.695626in}{0.595855in}}%
\pgfpathlineto{\pgfqpoint{3.697723in}{0.500952in}}%
\pgfpathlineto{\pgfqpoint{3.699820in}{0.596458in}}%
\pgfpathlineto{\pgfqpoint{3.701918in}{0.616686in}}%
\pgfpathlineto{\pgfqpoint{3.704015in}{0.625584in}}%
\pgfpathlineto{\pgfqpoint{3.706113in}{0.628481in}}%
\pgfpathlineto{\pgfqpoint{3.708210in}{0.626521in}}%
\pgfpathlineto{\pgfqpoint{3.710307in}{0.619027in}}%
\pgfpathlineto{\pgfqpoint{3.712405in}{0.602250in}}%
\pgfpathlineto{\pgfqpoint{3.714502in}{0.551666in}}%
\pgfpathlineto{\pgfqpoint{3.718697in}{0.617340in}}%
\pgfpathlineto{\pgfqpoint{3.720794in}{0.634111in}}%
\pgfpathlineto{\pgfqpoint{3.722891in}{0.643688in}}%
\pgfpathlineto{\pgfqpoint{3.724989in}{0.648258in}}%
\pgfpathlineto{\pgfqpoint{3.727086in}{0.647801in}}%
\pgfpathlineto{\pgfqpoint{3.729184in}{0.641301in}}%
\pgfpathlineto{\pgfqpoint{3.731281in}{0.630986in}}%
\pgfpathlineto{\pgfqpoint{3.733378in}{0.615953in}}%
\pgfpathlineto{\pgfqpoint{3.735476in}{0.592234in}}%
\pgfpathlineto{\pgfqpoint{3.737573in}{0.537366in}}%
\pgfpathlineto{\pgfqpoint{3.739670in}{0.555155in}}%
\pgfpathlineto{\pgfqpoint{3.741768in}{0.582537in}}%
\pgfpathlineto{\pgfqpoint{3.743865in}{0.590348in}}%
\pgfpathlineto{\pgfqpoint{3.745962in}{0.588520in}}%
\pgfpathlineto{\pgfqpoint{3.748060in}{0.576069in}}%
\pgfpathlineto{\pgfqpoint{3.750157in}{0.534186in}}%
\pgfpathlineto{\pgfqpoint{3.752254in}{0.553960in}}%
\pgfpathlineto{\pgfqpoint{3.754352in}{0.592413in}}%
\pgfpathlineto{\pgfqpoint{3.756449in}{0.610845in}}%
\pgfpathlineto{\pgfqpoint{3.758547in}{0.621445in}}%
\pgfpathlineto{\pgfqpoint{3.760644in}{0.626649in}}%
\pgfpathlineto{\pgfqpoint{3.762741in}{0.626421in}}%
\pgfpathlineto{\pgfqpoint{3.764839in}{0.619685in}}%
\pgfpathlineto{\pgfqpoint{3.766936in}{0.608632in}}%
\pgfpathlineto{\pgfqpoint{3.769033in}{0.591659in}}%
\pgfpathlineto{\pgfqpoint{3.771131in}{0.561631in}}%
\pgfpathlineto{\pgfqpoint{3.773228in}{0.467744in}}%
\pgfpathlineto{\pgfqpoint{3.775325in}{0.559166in}}%
\pgfpathlineto{\pgfqpoint{3.777423in}{0.577110in}}%
\pgfpathlineto{\pgfqpoint{3.779520in}{0.583133in}}%
\pgfpathlineto{\pgfqpoint{3.781618in}{0.581910in}}%
\pgfpathlineto{\pgfqpoint{3.783715in}{0.572993in}}%
\pgfpathlineto{\pgfqpoint{3.785812in}{0.549306in}}%
\pgfpathlineto{\pgfqpoint{3.787910in}{0.484270in}}%
\pgfpathlineto{\pgfqpoint{3.790007in}{0.565302in}}%
\pgfpathlineto{\pgfqpoint{3.792104in}{0.589503in}}%
\pgfpathlineto{\pgfqpoint{3.794202in}{0.602837in}}%
\pgfpathlineto{\pgfqpoint{3.796299in}{0.610154in}}%
\pgfpathlineto{\pgfqpoint{3.798396in}{0.612383in}}%
\pgfpathlineto{\pgfqpoint{3.800494in}{0.608414in}}%
\pgfpathlineto{\pgfqpoint{3.802591in}{0.597696in}}%
\pgfpathlineto{\pgfqpoint{3.804689in}{0.577520in}}%
\pgfpathlineto{\pgfqpoint{3.806786in}{0.522698in}}%
\pgfpathlineto{\pgfqpoint{3.810981in}{0.583504in}}%
\pgfpathlineto{\pgfqpoint{3.813078in}{0.597256in}}%
\pgfpathlineto{\pgfqpoint{3.815175in}{0.604190in}}%
\pgfpathlineto{\pgfqpoint{3.817273in}{0.606744in}}%
\pgfpathlineto{\pgfqpoint{3.819370in}{0.605560in}}%
\pgfpathlineto{\pgfqpoint{3.821467in}{0.600287in}}%
\pgfpathlineto{\pgfqpoint{3.823565in}{0.589164in}}%
\pgfpathlineto{\pgfqpoint{3.825662in}{0.565590in}}%
\pgfpathlineto{\pgfqpoint{3.827760in}{0.450445in}}%
\pgfpathlineto{\pgfqpoint{3.829857in}{0.571987in}}%
\pgfpathlineto{\pgfqpoint{3.831954in}{0.597205in}}%
\pgfpathlineto{\pgfqpoint{3.834052in}{0.610896in}}%
\pgfpathlineto{\pgfqpoint{3.836149in}{0.618752in}}%
\pgfpathlineto{\pgfqpoint{3.838246in}{0.622118in}}%
\pgfpathlineto{\pgfqpoint{3.840344in}{0.620664in}}%
\pgfpathlineto{\pgfqpoint{3.842441in}{0.612454in}}%
\pgfpathlineto{\pgfqpoint{3.844538in}{0.594140in}}%
\pgfpathlineto{\pgfqpoint{3.846636in}{0.541051in}}%
\pgfpathlineto{\pgfqpoint{3.850831in}{0.604284in}}%
\pgfpathlineto{\pgfqpoint{3.852928in}{0.619491in}}%
\pgfpathlineto{\pgfqpoint{3.855025in}{0.628051in}}%
\pgfpathlineto{\pgfqpoint{3.857123in}{0.632563in}}%
\pgfpathlineto{\pgfqpoint{3.859220in}{0.633929in}}%
\pgfpathlineto{\pgfqpoint{3.861317in}{0.632332in}}%
\pgfpathlineto{\pgfqpoint{3.863415in}{0.627377in}}%
\pgfpathlineto{\pgfqpoint{3.865512in}{0.617722in}}%
\pgfpathlineto{\pgfqpoint{3.867609in}{0.599182in}}%
\pgfpathlineto{\pgfqpoint{3.869707in}{0.545217in}}%
\pgfpathlineto{\pgfqpoint{3.873902in}{0.613369in}}%
\pgfpathlineto{\pgfqpoint{3.875999in}{0.629845in}}%
\pgfpathlineto{\pgfqpoint{3.878096in}{0.639613in}}%
\pgfpathlineto{\pgfqpoint{3.880194in}{0.644968in}}%
\pgfpathlineto{\pgfqpoint{3.882291in}{0.646332in}}%
\pgfpathlineto{\pgfqpoint{3.884388in}{0.642822in}}%
\pgfpathlineto{\pgfqpoint{3.886486in}{0.634397in}}%
\pgfpathlineto{\pgfqpoint{3.888583in}{0.621712in}}%
\pgfpathlineto{\pgfqpoint{3.890680in}{0.602020in}}%
\pgfpathlineto{\pgfqpoint{3.892778in}{0.563699in}}%
\pgfpathlineto{\pgfqpoint{3.894875in}{0.538049in}}%
\pgfpathlineto{\pgfqpoint{3.896973in}{0.581992in}}%
\pgfpathlineto{\pgfqpoint{3.899070in}{0.595058in}}%
\pgfpathlineto{\pgfqpoint{3.901167in}{0.598205in}}%
\pgfpathlineto{\pgfqpoint{3.903265in}{0.593888in}}%
\pgfpathlineto{\pgfqpoint{3.905362in}{0.579521in}}%
\pgfpathlineto{\pgfqpoint{3.907459in}{0.533210in}}%
\pgfpathlineto{\pgfqpoint{3.909557in}{0.559919in}}%
\pgfpathlineto{\pgfqpoint{3.911654in}{0.596126in}}%
\pgfpathlineto{\pgfqpoint{3.913751in}{0.613720in}}%
\pgfpathlineto{\pgfqpoint{3.915849in}{0.623768in}}%
\pgfpathlineto{\pgfqpoint{3.917946in}{0.628519in}}%
\pgfpathlineto{\pgfqpoint{3.920043in}{0.627841in}}%
\pgfpathlineto{\pgfqpoint{3.922141in}{0.620821in}}%
\pgfpathlineto{\pgfqpoint{3.924238in}{0.609682in}}%
\pgfpathlineto{\pgfqpoint{3.926336in}{0.592784in}}%
\pgfpathlineto{\pgfqpoint{3.928433in}{0.563435in}}%
\pgfpathlineto{\pgfqpoint{3.930530in}{0.439224in}}%
\pgfpathlineto{\pgfqpoint{3.932628in}{0.557186in}}%
\pgfpathlineto{\pgfqpoint{3.934725in}{0.575528in}}%
\pgfpathlineto{\pgfqpoint{3.936822in}{0.581320in}}%
\pgfpathlineto{\pgfqpoint{3.938920in}{0.579444in}}%
\pgfpathlineto{\pgfqpoint{3.941017in}{0.569027in}}%
\pgfpathlineto{\pgfqpoint{3.945212in}{0.516919in}}%
\pgfpathlineto{\pgfqpoint{3.947309in}{0.570992in}}%
\pgfpathlineto{\pgfqpoint{3.949407in}{0.592384in}}%
\pgfpathlineto{\pgfqpoint{3.951504in}{0.604452in}}%
\pgfpathlineto{\pgfqpoint{3.953601in}{0.610837in}}%
\pgfpathlineto{\pgfqpoint{3.955699in}{0.612049in}}%
\pgfpathlineto{\pgfqpoint{3.957796in}{0.606716in}}%
\pgfpathlineto{\pgfqpoint{3.959893in}{0.594795in}}%
\pgfpathlineto{\pgfqpoint{3.961991in}{0.572395in}}%
\pgfpathlineto{\pgfqpoint{3.964088in}{0.499382in}}%
\pgfpathlineto{\pgfqpoint{3.966185in}{0.557468in}}%
\pgfpathlineto{\pgfqpoint{3.968283in}{0.584267in}}%
\pgfpathlineto{\pgfqpoint{3.970380in}{0.596618in}}%
\pgfpathlineto{\pgfqpoint{3.972478in}{0.602633in}}%
\pgfpathlineto{\pgfqpoint{3.974575in}{0.604355in}}%
\pgfpathlineto{\pgfqpoint{3.976672in}{0.602216in}}%
\pgfpathlineto{\pgfqpoint{3.978770in}{0.595587in}}%
\pgfpathlineto{\pgfqpoint{3.980867in}{0.581925in}}%
\pgfpathlineto{\pgfqpoint{3.982964in}{0.550014in}}%
\pgfpathlineto{\pgfqpoint{3.985062in}{0.528546in}}%
\pgfpathlineto{\pgfqpoint{3.987159in}{0.580009in}}%
\pgfpathlineto{\pgfqpoint{3.989256in}{0.600446in}}%
\pgfpathlineto{\pgfqpoint{3.991354in}{0.612044in}}%
\pgfpathlineto{\pgfqpoint{3.993451in}{0.618422in}}%
\pgfpathlineto{\pgfqpoint{3.995549in}{0.620313in}}%
\pgfpathlineto{\pgfqpoint{3.997646in}{0.616804in}}%
\pgfpathlineto{\pgfqpoint{3.999743in}{0.605686in}}%
\pgfpathlineto{\pgfqpoint{4.001841in}{0.581101in}}%
\pgfpathlineto{\pgfqpoint{4.003938in}{0.473892in}}%
\pgfpathlineto{\pgfqpoint{4.006035in}{0.585645in}}%
\pgfpathlineto{\pgfqpoint{4.008133in}{0.609220in}}%
\pgfpathlineto{\pgfqpoint{4.010230in}{0.621572in}}%
\pgfpathlineto{\pgfqpoint{4.012327in}{0.628551in}}%
\pgfpathlineto{\pgfqpoint{4.014425in}{0.631926in}}%
\pgfpathlineto{\pgfqpoint{4.016522in}{0.632298in}}%
\pgfpathlineto{\pgfqpoint{4.018620in}{0.629661in}}%
\pgfpathlineto{\pgfqpoint{4.020717in}{0.623382in}}%
\pgfpathlineto{\pgfqpoint{4.022814in}{0.611566in}}%
\pgfpathlineto{\pgfqpoint{4.024912in}{0.587604in}}%
\pgfpathlineto{\pgfqpoint{4.027009in}{0.464981in}}%
\pgfpathlineto{\pgfqpoint{4.029106in}{0.592912in}}%
\pgfpathlineto{\pgfqpoint{4.031204in}{0.618194in}}%
\pgfpathlineto{\pgfqpoint{4.033301in}{0.632040in}}%
\pgfpathlineto{\pgfqpoint{4.035398in}{0.640262in}}%
\pgfpathlineto{\pgfqpoint{4.037496in}{0.644357in}}%
\pgfpathlineto{\pgfqpoint{4.039593in}{0.644353in}}%
\pgfpathlineto{\pgfqpoint{4.041691in}{0.639072in}}%
\pgfpathlineto{\pgfqpoint{4.043788in}{0.629376in}}%
\pgfpathlineto{\pgfqpoint{4.045885in}{0.614672in}}%
\pgfpathlineto{\pgfqpoint{4.047983in}{0.590526in}}%
\pgfpathlineto{\pgfqpoint{4.050080in}{0.526735in}}%
\pgfpathlineto{\pgfqpoint{4.052177in}{0.563653in}}%
\pgfpathlineto{\pgfqpoint{4.054275in}{0.589246in}}%
\pgfpathlineto{\pgfqpoint{4.056372in}{0.598140in}}%
\pgfpathlineto{\pgfqpoint{4.058469in}{0.598998in}}%
\pgfpathlineto{\pgfqpoint{4.060567in}{0.592590in}}%
\pgfpathlineto{\pgfqpoint{4.062664in}{0.574702in}}%
\pgfpathlineto{\pgfqpoint{4.064762in}{0.499505in}}%
\pgfpathlineto{\pgfqpoint{4.066859in}{0.570849in}}%
\pgfpathlineto{\pgfqpoint{4.068956in}{0.600703in}}%
\pgfpathlineto{\pgfqpoint{4.071054in}{0.616255in}}%
\pgfpathlineto{\pgfqpoint{4.073151in}{0.625057in}}%
\pgfpathlineto{\pgfqpoint{4.075248in}{0.628716in}}%
\pgfpathlineto{\pgfqpoint{4.077346in}{0.626681in}}%
\pgfpathlineto{\pgfqpoint{4.079443in}{0.618669in}}%
\pgfpathlineto{\pgfqpoint{4.081540in}{0.606394in}}%
\pgfpathlineto{\pgfqpoint{4.083638in}{0.587281in}}%
\pgfpathlineto{\pgfqpoint{4.087832in}{0.518683in}}%
\pgfpathlineto{\pgfqpoint{4.089930in}{0.566160in}}%
\pgfpathlineto{\pgfqpoint{4.092027in}{0.580132in}}%
\pgfpathlineto{\pgfqpoint{4.094125in}{0.584105in}}%
\pgfpathlineto{\pgfqpoint{4.096222in}{0.580970in}}%
\pgfpathlineto{\pgfqpoint{4.098319in}{0.569100in}}%
\pgfpathlineto{\pgfqpoint{4.100417in}{0.535729in}}%
\pgfpathlineto{\pgfqpoint{4.102514in}{0.528777in}}%
\pgfpathlineto{\pgfqpoint{4.104611in}{0.575254in}}%
\pgfpathlineto{\pgfqpoint{4.106709in}{0.595233in}}%
\pgfpathlineto{\pgfqpoint{4.108806in}{0.606570in}}%
\pgfpathlineto{\pgfqpoint{4.110903in}{0.612393in}}%
\pgfpathlineto{\pgfqpoint{4.113001in}{0.613010in}}%
\pgfpathlineto{\pgfqpoint{4.115098in}{0.606987in}}%
\pgfpathlineto{\pgfqpoint{4.117196in}{0.594342in}}%
\pgfpathlineto{\pgfqpoint{4.119293in}{0.570269in}}%
\pgfpathlineto{\pgfqpoint{4.121390in}{0.468236in}}%
\pgfpathlineto{\pgfqpoint{4.123488in}{0.562500in}}%
\pgfpathlineto{\pgfqpoint{4.125585in}{0.586977in}}%
\pgfpathlineto{\pgfqpoint{4.127682in}{0.598527in}}%
\pgfpathlineto{\pgfqpoint{4.129780in}{0.604064in}}%
\pgfpathlineto{\pgfqpoint{4.131877in}{0.605398in}}%
\pgfpathlineto{\pgfqpoint{4.133974in}{0.602854in}}%
\pgfpathlineto{\pgfqpoint{4.136072in}{0.595677in}}%
\pgfpathlineto{\pgfqpoint{4.138169in}{0.580972in}}%
\pgfpathlineto{\pgfqpoint{4.140267in}{0.544929in}}%
\pgfpathlineto{\pgfqpoint{4.142364in}{0.539440in}}%
\pgfpathlineto{\pgfqpoint{4.144461in}{0.583880in}}%
\pgfpathlineto{\pgfqpoint{4.146559in}{0.603045in}}%
\pgfpathlineto{\pgfqpoint{4.148656in}{0.614016in}}%
\pgfpathlineto{\pgfqpoint{4.150753in}{0.619926in}}%
\pgfpathlineto{\pgfqpoint{4.152851in}{0.621332in}}%
\pgfpathlineto{\pgfqpoint{4.154948in}{0.617137in}}%
\pgfpathlineto{\pgfqpoint{4.157045in}{0.605229in}}%
\pgfpathlineto{\pgfqpoint{4.159143in}{0.579064in}}%
\pgfpathlineto{\pgfqpoint{4.161240in}{0.502877in}}%
\pgfpathlineto{\pgfqpoint{4.163338in}{0.587785in}}%
\pgfpathlineto{\pgfqpoint{4.165435in}{0.610152in}}%
\pgfpathlineto{\pgfqpoint{4.167532in}{0.621946in}}%
\pgfpathlineto{\pgfqpoint{4.169630in}{0.628521in}}%
\pgfpathlineto{\pgfqpoint{4.171727in}{0.631522in}}%
\pgfpathlineto{\pgfqpoint{4.173824in}{0.631486in}}%
\pgfpathlineto{\pgfqpoint{4.175922in}{0.628330in}}%
\pgfpathlineto{\pgfqpoint{4.178019in}{0.621274in}}%
\pgfpathlineto{\pgfqpoint{4.180116in}{0.607987in}}%
\pgfpathlineto{\pgfqpoint{4.182214in}{0.579436in}}%
\pgfpathlineto{\pgfqpoint{4.184311in}{0.535167in}}%
\pgfpathlineto{\pgfqpoint{4.186409in}{0.599126in}}%
\pgfpathlineto{\pgfqpoint{4.188506in}{0.621213in}}%
\pgfpathlineto{\pgfqpoint{4.190603in}{0.633731in}}%
\pgfpathlineto{\pgfqpoint{4.192701in}{0.641060in}}%
\pgfpathlineto{\pgfqpoint{4.194798in}{0.644322in}}%
\pgfpathlineto{\pgfqpoint{4.196895in}{0.643275in}}%
\pgfpathlineto{\pgfqpoint{4.198993in}{0.636783in}}%
\pgfpathlineto{\pgfqpoint{4.201090in}{0.626045in}}%
\pgfpathlineto{\pgfqpoint{4.203187in}{0.609585in}}%
\pgfpathlineto{\pgfqpoint{4.205285in}{0.580928in}}%
\pgfpathlineto{\pgfqpoint{4.207382in}{0.430027in}}%
\pgfpathlineto{\pgfqpoint{4.209480in}{0.573980in}}%
\pgfpathlineto{\pgfqpoint{4.211577in}{0.593173in}}%
\pgfpathlineto{\pgfqpoint{4.213674in}{0.599786in}}%
\pgfpathlineto{\pgfqpoint{4.215772in}{0.599087in}}%
\pgfpathlineto{\pgfqpoint{4.217869in}{0.590925in}}%
\pgfpathlineto{\pgfqpoint{4.219966in}{0.569276in}}%
\pgfpathlineto{\pgfqpoint{4.222064in}{0.462782in}}%
\pgfpathlineto{\pgfqpoint{4.224161in}{0.578843in}}%
\pgfpathlineto{\pgfqpoint{4.226258in}{0.604605in}}%
\pgfpathlineto{\pgfqpoint{4.228356in}{0.618592in}}%
\pgfpathlineto{\pgfqpoint{4.230453in}{0.626364in}}%
\pgfpathlineto{\pgfqpoint{4.232551in}{0.629050in}}%
\pgfpathlineto{\pgfqpoint{4.234648in}{0.625714in}}%
\pgfpathlineto{\pgfqpoint{4.236745in}{0.616954in}}%
\pgfpathlineto{\pgfqpoint{4.238843in}{0.603599in}}%
\pgfpathlineto{\pgfqpoint{4.240940in}{0.582271in}}%
\pgfpathlineto{\pgfqpoint{4.243037in}{0.535977in}}%
\pgfpathlineto{\pgfqpoint{4.245135in}{0.538012in}}%
\pgfpathlineto{\pgfqpoint{4.247232in}{0.571572in}}%
\pgfpathlineto{\pgfqpoint{4.249329in}{0.582897in}}%
\pgfpathlineto{\pgfqpoint{4.251427in}{0.585484in}}%
\pgfpathlineto{\pgfqpoint{4.253524in}{0.581141in}}%
\pgfpathlineto{\pgfqpoint{4.255621in}{0.567456in}}%
\pgfpathlineto{\pgfqpoint{4.257719in}{0.526682in}}%
\pgfpathlineto{\pgfqpoint{4.259816in}{0.540494in}}%
\pgfpathlineto{\pgfqpoint{4.261914in}{0.579746in}}%
\pgfpathlineto{\pgfqpoint{4.264011in}{0.598080in}}%
\pgfpathlineto{\pgfqpoint{4.266108in}{0.608529in}}%
\pgfpathlineto{\pgfqpoint{4.268206in}{0.613636in}}%
\pgfpathlineto{\pgfqpoint{4.270303in}{0.613439in}}%
\pgfpathlineto{\pgfqpoint{4.272400in}{0.606493in}}%
\pgfpathlineto{\pgfqpoint{4.274498in}{0.592831in}}%
\pgfpathlineto{\pgfqpoint{4.276595in}{0.566278in}}%
\pgfpathlineto{\pgfqpoint{4.278692in}{0.456191in}}%
\pgfpathlineto{\pgfqpoint{4.280790in}{0.567355in}}%
\pgfpathlineto{\pgfqpoint{4.282887in}{0.589307in}}%
\pgfpathlineto{\pgfqpoint{4.284985in}{0.599877in}}%
\pgfpathlineto{\pgfqpoint{4.287082in}{0.604782in}}%
\pgfpathlineto{\pgfqpoint{4.289179in}{0.605563in}}%
\pgfpathlineto{\pgfqpoint{4.291277in}{0.602401in}}%
\pgfpathlineto{\pgfqpoint{4.293374in}{0.594329in}}%
\pgfpathlineto{\pgfqpoint{4.295471in}{0.577782in}}%
\pgfpathlineto{\pgfqpoint{4.297569in}{0.532625in}}%
\pgfpathlineto{\pgfqpoint{4.299666in}{0.551531in}}%
\pgfpathlineto{\pgfqpoint{4.301763in}{0.588288in}}%
\pgfpathlineto{\pgfqpoint{4.303861in}{0.605687in}}%
\pgfpathlineto{\pgfqpoint{4.305958in}{0.615722in}}%
\pgfpathlineto{\pgfqpoint{4.308056in}{0.620895in}}%
\pgfpathlineto{\pgfqpoint{4.310153in}{0.621494in}}%
\pgfpathlineto{\pgfqpoint{4.312250in}{0.616519in}}%
\pgfpathlineto{\pgfqpoint{4.316445in}{0.603256in}}%
\pgfpathlineto{\pgfqpoint{4.316445in}{0.603256in}}%
\pgfusepath{stroke}%
\end{pgfscope}%
\begin{pgfscope}%
\pgfpathrectangle{\pgfqpoint{2.554660in}{0.320625in}}{\pgfqpoint{1.845679in}{1.164375in}}%
\pgfusepath{clip}%
\pgfsetrectcap%
\pgfsetroundjoin%
\pgfsetlinewidth{1.154312pt}%
\definecolor{currentstroke}{rgb}{0.989363,0.557873,0.391671}%
\pgfsetstrokecolor{currentstroke}%
\pgfsetdash{}{0pt}%
\pgfpathmoveto{\pgfqpoint{2.638555in}{1.468525in}}%
\pgfpathlineto{\pgfqpoint{2.649042in}{1.449217in}}%
\pgfpathlineto{\pgfqpoint{2.659529in}{1.427093in}}%
\pgfpathlineto{\pgfqpoint{2.672113in}{1.397448in}}%
\pgfpathlineto{\pgfqpoint{2.716157in}{1.291455in}}%
\pgfpathlineto{\pgfqpoint{2.739228in}{1.239960in}}%
\pgfpathlineto{\pgfqpoint{2.764397in}{1.186560in}}%
\pgfpathlineto{\pgfqpoint{2.795857in}{1.122540in}}%
\pgfpathlineto{\pgfqpoint{2.835707in}{1.044156in}}%
\pgfpathlineto{\pgfqpoint{2.886044in}{0.947761in}}%
\pgfpathlineto{\pgfqpoint{2.948965in}{0.829875in}}%
\pgfpathlineto{\pgfqpoint{3.030762in}{0.679922in}}%
\pgfpathlineto{\pgfqpoint{3.043346in}{0.660579in}}%
\pgfpathlineto{\pgfqpoint{3.049638in}{0.647198in}}%
\pgfpathlineto{\pgfqpoint{3.064320in}{0.613034in}}%
\pgfpathlineto{\pgfqpoint{3.068514in}{0.606861in}}%
\pgfpathlineto{\pgfqpoint{3.070612in}{0.605072in}}%
\pgfpathlineto{\pgfqpoint{3.072709in}{0.604158in}}%
\pgfpathlineto{\pgfqpoint{3.076904in}{0.604339in}}%
\pgfpathlineto{\pgfqpoint{3.081098in}{0.605317in}}%
\pgfpathlineto{\pgfqpoint{3.083196in}{0.605215in}}%
\pgfpathlineto{\pgfqpoint{3.085293in}{0.604151in}}%
\pgfpathlineto{\pgfqpoint{3.087391in}{0.601536in}}%
\pgfpathlineto{\pgfqpoint{3.091585in}{0.591627in}}%
\pgfpathlineto{\pgfqpoint{3.095780in}{0.576075in}}%
\pgfpathlineto{\pgfqpoint{3.102072in}{0.544057in}}%
\pgfpathlineto{\pgfqpoint{3.104169in}{0.543118in}}%
\pgfpathlineto{\pgfqpoint{3.110462in}{0.561802in}}%
\pgfpathlineto{\pgfqpoint{3.112559in}{0.564976in}}%
\pgfpathlineto{\pgfqpoint{3.114656in}{0.566206in}}%
\pgfpathlineto{\pgfqpoint{3.116754in}{0.565530in}}%
\pgfpathlineto{\pgfqpoint{3.118851in}{0.562809in}}%
\pgfpathlineto{\pgfqpoint{3.120948in}{0.557620in}}%
\pgfpathlineto{\pgfqpoint{3.123046in}{0.548984in}}%
\pgfpathlineto{\pgfqpoint{3.125143in}{0.534502in}}%
\pgfpathlineto{\pgfqpoint{3.127240in}{0.507285in}}%
\pgfpathlineto{\pgfqpoint{3.129338in}{0.498422in}}%
\pgfpathlineto{\pgfqpoint{3.131435in}{0.533483in}}%
\pgfpathlineto{\pgfqpoint{3.133533in}{0.553377in}}%
\pgfpathlineto{\pgfqpoint{3.135630in}{0.566352in}}%
\pgfpathlineto{\pgfqpoint{3.137727in}{0.575531in}}%
\pgfpathlineto{\pgfqpoint{3.139825in}{0.582152in}}%
\pgfpathlineto{\pgfqpoint{3.141922in}{0.586768in}}%
\pgfpathlineto{\pgfqpoint{3.144019in}{0.589597in}}%
\pgfpathlineto{\pgfqpoint{3.146117in}{0.590625in}}%
\pgfpathlineto{\pgfqpoint{3.148214in}{0.589590in}}%
\pgfpathlineto{\pgfqpoint{3.150311in}{0.585576in}}%
\pgfpathlineto{\pgfqpoint{3.152409in}{0.576323in}}%
\pgfpathlineto{\pgfqpoint{3.154506in}{0.577521in}}%
\pgfpathlineto{\pgfqpoint{3.158701in}{0.616553in}}%
\pgfpathlineto{\pgfqpoint{3.162896in}{0.641299in}}%
\pgfpathlineto{\pgfqpoint{3.167090in}{0.656445in}}%
\pgfpathlineto{\pgfqpoint{3.171285in}{0.666082in}}%
\pgfpathlineto{\pgfqpoint{3.175480in}{0.671901in}}%
\pgfpathlineto{\pgfqpoint{3.177577in}{0.673631in}}%
\pgfpathlineto{\pgfqpoint{3.179675in}{0.674636in}}%
\pgfpathlineto{\pgfqpoint{3.181772in}{0.674930in}}%
\pgfpathlineto{\pgfqpoint{3.183869in}{0.674498in}}%
\pgfpathlineto{\pgfqpoint{3.185967in}{0.673297in}}%
\pgfpathlineto{\pgfqpoint{3.188064in}{0.671249in}}%
\pgfpathlineto{\pgfqpoint{3.190161in}{0.668218in}}%
\pgfpathlineto{\pgfqpoint{3.192259in}{0.663987in}}%
\pgfpathlineto{\pgfqpoint{3.194356in}{0.658188in}}%
\pgfpathlineto{\pgfqpoint{3.196453in}{0.650155in}}%
\pgfpathlineto{\pgfqpoint{3.198551in}{0.638525in}}%
\pgfpathlineto{\pgfqpoint{3.200648in}{0.619815in}}%
\pgfpathlineto{\pgfqpoint{3.202746in}{0.582224in}}%
\pgfpathlineto{\pgfqpoint{3.204843in}{0.597297in}}%
\pgfpathlineto{\pgfqpoint{3.206940in}{0.628780in}}%
\pgfpathlineto{\pgfqpoint{3.209038in}{0.646266in}}%
\pgfpathlineto{\pgfqpoint{3.213232in}{0.666717in}}%
\pgfpathlineto{\pgfqpoint{3.217427in}{0.678368in}}%
\pgfpathlineto{\pgfqpoint{3.221622in}{0.684846in}}%
\pgfpathlineto{\pgfqpoint{3.223719in}{0.686448in}}%
\pgfpathlineto{\pgfqpoint{3.225816in}{0.686931in}}%
\pgfpathlineto{\pgfqpoint{3.227914in}{0.686155in}}%
\pgfpathlineto{\pgfqpoint{3.230011in}{0.683833in}}%
\pgfpathlineto{\pgfqpoint{3.232109in}{0.679414in}}%
\pgfpathlineto{\pgfqpoint{3.236303in}{0.666606in}}%
\pgfpathlineto{\pgfqpoint{3.246790in}{0.632473in}}%
\pgfpathlineto{\pgfqpoint{3.248887in}{0.627951in}}%
\pgfpathlineto{\pgfqpoint{3.250985in}{0.624870in}}%
\pgfpathlineto{\pgfqpoint{3.253082in}{0.623549in}}%
\pgfpathlineto{\pgfqpoint{3.255180in}{0.623924in}}%
\pgfpathlineto{\pgfqpoint{3.259374in}{0.625714in}}%
\pgfpathlineto{\pgfqpoint{3.261472in}{0.623839in}}%
\pgfpathlineto{\pgfqpoint{3.265666in}{0.611386in}}%
\pgfpathlineto{\pgfqpoint{3.276153in}{0.574575in}}%
\pgfpathlineto{\pgfqpoint{3.278251in}{0.569588in}}%
\pgfpathlineto{\pgfqpoint{3.280348in}{0.566186in}}%
\pgfpathlineto{\pgfqpoint{3.282445in}{0.564780in}}%
\pgfpathlineto{\pgfqpoint{3.284543in}{0.565509in}}%
\pgfpathlineto{\pgfqpoint{3.288737in}{0.569401in}}%
\pgfpathlineto{\pgfqpoint{3.290835in}{0.569155in}}%
\pgfpathlineto{\pgfqpoint{3.292932in}{0.564477in}}%
\pgfpathlineto{\pgfqpoint{3.295029in}{0.552018in}}%
\pgfpathlineto{\pgfqpoint{3.297127in}{0.530933in}}%
\pgfpathlineto{\pgfqpoint{3.299224in}{0.538195in}}%
\pgfpathlineto{\pgfqpoint{3.301322in}{0.561815in}}%
\pgfpathlineto{\pgfqpoint{3.303419in}{0.577104in}}%
\pgfpathlineto{\pgfqpoint{3.305516in}{0.586912in}}%
\pgfpathlineto{\pgfqpoint{3.307614in}{0.593142in}}%
\pgfpathlineto{\pgfqpoint{3.309711in}{0.596685in}}%
\pgfpathlineto{\pgfqpoint{3.311808in}{0.597927in}}%
\pgfpathlineto{\pgfqpoint{3.313906in}{0.596939in}}%
\pgfpathlineto{\pgfqpoint{3.316003in}{0.593495in}}%
\pgfpathlineto{\pgfqpoint{3.318100in}{0.586921in}}%
\pgfpathlineto{\pgfqpoint{3.320198in}{0.575535in}}%
\pgfpathlineto{\pgfqpoint{3.322295in}{0.554118in}}%
\pgfpathlineto{\pgfqpoint{3.324393in}{0.481660in}}%
\pgfpathlineto{\pgfqpoint{3.326490in}{0.547337in}}%
\pgfpathlineto{\pgfqpoint{3.328587in}{0.575846in}}%
\pgfpathlineto{\pgfqpoint{3.330685in}{0.591123in}}%
\pgfpathlineto{\pgfqpoint{3.332782in}{0.600596in}}%
\pgfpathlineto{\pgfqpoint{3.334879in}{0.606251in}}%
\pgfpathlineto{\pgfqpoint{3.336977in}{0.608618in}}%
\pgfpathlineto{\pgfqpoint{3.339074in}{0.607372in}}%
\pgfpathlineto{\pgfqpoint{3.341171in}{0.601057in}}%
\pgfpathlineto{\pgfqpoint{3.343269in}{0.587332in}}%
\pgfpathlineto{\pgfqpoint{3.347464in}{0.529210in}}%
\pgfpathlineto{\pgfqpoint{3.349561in}{0.575625in}}%
\pgfpathlineto{\pgfqpoint{3.351658in}{0.594902in}}%
\pgfpathlineto{\pgfqpoint{3.353756in}{0.605264in}}%
\pgfpathlineto{\pgfqpoint{3.355853in}{0.610783in}}%
\pgfpathlineto{\pgfqpoint{3.357950in}{0.612788in}}%
\pgfpathlineto{\pgfqpoint{3.360048in}{0.611628in}}%
\pgfpathlineto{\pgfqpoint{3.362145in}{0.606986in}}%
\pgfpathlineto{\pgfqpoint{3.364242in}{0.597568in}}%
\pgfpathlineto{\pgfqpoint{3.366340in}{0.579226in}}%
\pgfpathlineto{\pgfqpoint{3.368437in}{0.525799in}}%
\pgfpathlineto{\pgfqpoint{3.372632in}{0.593527in}}%
\pgfpathlineto{\pgfqpoint{3.374729in}{0.610010in}}%
\pgfpathlineto{\pgfqpoint{3.376827in}{0.619724in}}%
\pgfpathlineto{\pgfqpoint{3.378924in}{0.624943in}}%
\pgfpathlineto{\pgfqpoint{3.381021in}{0.626031in}}%
\pgfpathlineto{\pgfqpoint{3.383119in}{0.621946in}}%
\pgfpathlineto{\pgfqpoint{3.385216in}{0.610816in}}%
\pgfpathlineto{\pgfqpoint{3.387313in}{0.587777in}}%
\pgfpathlineto{\pgfqpoint{3.389411in}{0.514528in}}%
\pgfpathlineto{\pgfqpoint{3.391508in}{0.584591in}}%
\pgfpathlineto{\pgfqpoint{3.393605in}{0.609261in}}%
\pgfpathlineto{\pgfqpoint{3.395703in}{0.621779in}}%
\pgfpathlineto{\pgfqpoint{3.397800in}{0.628627in}}%
\pgfpathlineto{\pgfqpoint{3.399898in}{0.631688in}}%
\pgfpathlineto{\pgfqpoint{3.401995in}{0.631557in}}%
\pgfpathlineto{\pgfqpoint{3.404092in}{0.628138in}}%
\pgfpathlineto{\pgfqpoint{3.406190in}{0.620543in}}%
\pgfpathlineto{\pgfqpoint{3.408287in}{0.606029in}}%
\pgfpathlineto{\pgfqpoint{3.410384in}{0.572727in}}%
\pgfpathlineto{\pgfqpoint{3.412482in}{0.555252in}}%
\pgfpathlineto{\pgfqpoint{3.414579in}{0.604502in}}%
\pgfpathlineto{\pgfqpoint{3.416676in}{0.624571in}}%
\pgfpathlineto{\pgfqpoint{3.418774in}{0.636131in}}%
\pgfpathlineto{\pgfqpoint{3.420871in}{0.642736in}}%
\pgfpathlineto{\pgfqpoint{3.422969in}{0.645228in}}%
\pgfpathlineto{\pgfqpoint{3.425066in}{0.643078in}}%
\pgfpathlineto{\pgfqpoint{3.427163in}{0.634885in}}%
\pgfpathlineto{\pgfqpoint{3.429261in}{0.619580in}}%
\pgfpathlineto{\pgfqpoint{3.433455in}{0.548829in}}%
\pgfpathlineto{\pgfqpoint{3.435553in}{0.604162in}}%
\pgfpathlineto{\pgfqpoint{3.437650in}{0.622821in}}%
\pgfpathlineto{\pgfqpoint{3.439747in}{0.632074in}}%
\pgfpathlineto{\pgfqpoint{3.441845in}{0.636168in}}%
\pgfpathlineto{\pgfqpoint{3.443942in}{0.636296in}}%
\pgfpathlineto{\pgfqpoint{3.446040in}{0.632454in}}%
\pgfpathlineto{\pgfqpoint{3.448137in}{0.623443in}}%
\pgfpathlineto{\pgfqpoint{3.450234in}{0.604927in}}%
\pgfpathlineto{\pgfqpoint{3.452332in}{0.546342in}}%
\pgfpathlineto{\pgfqpoint{3.454429in}{0.590327in}}%
\pgfpathlineto{\pgfqpoint{3.456526in}{0.622218in}}%
\pgfpathlineto{\pgfqpoint{3.458624in}{0.638238in}}%
\pgfpathlineto{\pgfqpoint{3.460721in}{0.647507in}}%
\pgfpathlineto{\pgfqpoint{3.462818in}{0.652071in}}%
\pgfpathlineto{\pgfqpoint{3.464916in}{0.652007in}}%
\pgfpathlineto{\pgfqpoint{3.467013in}{0.646121in}}%
\pgfpathlineto{\pgfqpoint{3.469111in}{0.635802in}}%
\pgfpathlineto{\pgfqpoint{3.471208in}{0.619988in}}%
\pgfpathlineto{\pgfqpoint{3.473305in}{0.592967in}}%
\pgfpathlineto{\pgfqpoint{3.475403in}{0.504467in}}%
\pgfpathlineto{\pgfqpoint{3.477500in}{0.580835in}}%
\pgfpathlineto{\pgfqpoint{3.479597in}{0.601577in}}%
\pgfpathlineto{\pgfqpoint{3.481695in}{0.608980in}}%
\pgfpathlineto{\pgfqpoint{3.483792in}{0.609032in}}%
\pgfpathlineto{\pgfqpoint{3.485889in}{0.601987in}}%
\pgfpathlineto{\pgfqpoint{3.487987in}{0.583197in}}%
\pgfpathlineto{\pgfqpoint{3.490084in}{0.497963in}}%
\pgfpathlineto{\pgfqpoint{3.492182in}{0.582033in}}%
\pgfpathlineto{\pgfqpoint{3.494279in}{0.610833in}}%
\pgfpathlineto{\pgfqpoint{3.496376in}{0.625984in}}%
\pgfpathlineto{\pgfqpoint{3.498474in}{0.634533in}}%
\pgfpathlineto{\pgfqpoint{3.500571in}{0.637977in}}%
\pgfpathlineto{\pgfqpoint{3.502668in}{0.635699in}}%
\pgfpathlineto{\pgfqpoint{3.504766in}{0.627696in}}%
\pgfpathlineto{\pgfqpoint{3.506863in}{0.615870in}}%
\pgfpathlineto{\pgfqpoint{3.508960in}{0.598214in}}%
\pgfpathlineto{\pgfqpoint{3.511058in}{0.567715in}}%
\pgfpathlineto{\pgfqpoint{3.513155in}{0.480158in}}%
\pgfpathlineto{\pgfqpoint{3.515253in}{0.562446in}}%
\pgfpathlineto{\pgfqpoint{3.517350in}{0.579153in}}%
\pgfpathlineto{\pgfqpoint{3.519447in}{0.583262in}}%
\pgfpathlineto{\pgfqpoint{3.521545in}{0.578756in}}%
\pgfpathlineto{\pgfqpoint{3.523642in}{0.562440in}}%
\pgfpathlineto{\pgfqpoint{3.525739in}{0.495149in}}%
\pgfpathlineto{\pgfqpoint{3.527837in}{0.557535in}}%
\pgfpathlineto{\pgfqpoint{3.529934in}{0.588548in}}%
\pgfpathlineto{\pgfqpoint{3.532031in}{0.604668in}}%
\pgfpathlineto{\pgfqpoint{3.534129in}{0.613869in}}%
\pgfpathlineto{\pgfqpoint{3.536226in}{0.617851in}}%
\pgfpathlineto{\pgfqpoint{3.538324in}{0.616116in}}%
\pgfpathlineto{\pgfqpoint{3.540421in}{0.607855in}}%
\pgfpathlineto{\pgfqpoint{3.542518in}{0.593813in}}%
\pgfpathlineto{\pgfqpoint{3.544616in}{0.568580in}}%
\pgfpathlineto{\pgfqpoint{3.546713in}{0.470473in}}%
\pgfpathlineto{\pgfqpoint{3.548810in}{0.558411in}}%
\pgfpathlineto{\pgfqpoint{3.550908in}{0.581681in}}%
\pgfpathlineto{\pgfqpoint{3.553005in}{0.591818in}}%
\pgfpathlineto{\pgfqpoint{3.555102in}{0.595583in}}%
\pgfpathlineto{\pgfqpoint{3.557200in}{0.594488in}}%
\pgfpathlineto{\pgfqpoint{3.559297in}{0.588179in}}%
\pgfpathlineto{\pgfqpoint{3.561395in}{0.573826in}}%
\pgfpathlineto{\pgfqpoint{3.563492in}{0.536561in}}%
\pgfpathlineto{\pgfqpoint{3.565589in}{0.537032in}}%
\pgfpathlineto{\pgfqpoint{3.567687in}{0.579648in}}%
\pgfpathlineto{\pgfqpoint{3.569784in}{0.598501in}}%
\pgfpathlineto{\pgfqpoint{3.571881in}{0.609227in}}%
\pgfpathlineto{\pgfqpoint{3.573979in}{0.614744in}}%
\pgfpathlineto{\pgfqpoint{3.576076in}{0.615394in}}%
\pgfpathlineto{\pgfqpoint{3.578173in}{0.609732in}}%
\pgfpathlineto{\pgfqpoint{3.580271in}{0.596188in}}%
\pgfpathlineto{\pgfqpoint{3.582368in}{0.566611in}}%
\pgfpathlineto{\pgfqpoint{3.584465in}{0.521022in}}%
\pgfpathlineto{\pgfqpoint{3.586563in}{0.583087in}}%
\pgfpathlineto{\pgfqpoint{3.588660in}{0.603312in}}%
\pgfpathlineto{\pgfqpoint{3.590758in}{0.613997in}}%
\pgfpathlineto{\pgfqpoint{3.592855in}{0.619697in}}%
\pgfpathlineto{\pgfqpoint{3.594952in}{0.621833in}}%
\pgfpathlineto{\pgfqpoint{3.597050in}{0.620789in}}%
\pgfpathlineto{\pgfqpoint{3.599147in}{0.616276in}}%
\pgfpathlineto{\pgfqpoint{3.601244in}{0.607039in}}%
\pgfpathlineto{\pgfqpoint{3.603342in}{0.589083in}}%
\pgfpathlineto{\pgfqpoint{3.605439in}{0.538196in}}%
\pgfpathlineto{\pgfqpoint{3.609634in}{0.601994in}}%
\pgfpathlineto{\pgfqpoint{3.611731in}{0.618735in}}%
\pgfpathlineto{\pgfqpoint{3.613829in}{0.628575in}}%
\pgfpathlineto{\pgfqpoint{3.615926in}{0.633870in}}%
\pgfpathlineto{\pgfqpoint{3.618023in}{0.635010in}}%
\pgfpathlineto{\pgfqpoint{3.620121in}{0.630963in}}%
\pgfpathlineto{\pgfqpoint{3.622218in}{0.619899in}}%
\pgfpathlineto{\pgfqpoint{3.624315in}{0.597050in}}%
\pgfpathlineto{\pgfqpoint{3.626413in}{0.495668in}}%
\pgfpathlineto{\pgfqpoint{3.628510in}{0.591576in}}%
\pgfpathlineto{\pgfqpoint{3.630607in}{0.617064in}}%
\pgfpathlineto{\pgfqpoint{3.632705in}{0.629752in}}%
\pgfpathlineto{\pgfqpoint{3.634802in}{0.636628in}}%
\pgfpathlineto{\pgfqpoint{3.636900in}{0.639647in}}%
\pgfpathlineto{\pgfqpoint{3.638997in}{0.639416in}}%
\pgfpathlineto{\pgfqpoint{3.641094in}{0.635818in}}%
\pgfpathlineto{\pgfqpoint{3.643192in}{0.627898in}}%
\pgfpathlineto{\pgfqpoint{3.645289in}{0.612657in}}%
\pgfpathlineto{\pgfqpoint{3.647386in}{0.576214in}}%
\pgfpathlineto{\pgfqpoint{3.649484in}{0.570549in}}%
\pgfpathlineto{\pgfqpoint{3.651581in}{0.614819in}}%
\pgfpathlineto{\pgfqpoint{3.653678in}{0.633954in}}%
\pgfpathlineto{\pgfqpoint{3.655776in}{0.645039in}}%
\pgfpathlineto{\pgfqpoint{3.657873in}{0.651265in}}%
\pgfpathlineto{\pgfqpoint{3.659971in}{0.653339in}}%
\pgfpathlineto{\pgfqpoint{3.662068in}{0.650565in}}%
\pgfpathlineto{\pgfqpoint{3.664165in}{0.642165in}}%
\pgfpathlineto{\pgfqpoint{3.666263in}{0.628257in}}%
\pgfpathlineto{\pgfqpoint{3.668360in}{0.603881in}}%
\pgfpathlineto{\pgfqpoint{3.670457in}{0.523962in}}%
\pgfpathlineto{\pgfqpoint{3.672555in}{0.588033in}}%
\pgfpathlineto{\pgfqpoint{3.674652in}{0.612578in}}%
\pgfpathlineto{\pgfqpoint{3.676749in}{0.622861in}}%
\pgfpathlineto{\pgfqpoint{3.678847in}{0.626423in}}%
\pgfpathlineto{\pgfqpoint{3.680944in}{0.624834in}}%
\pgfpathlineto{\pgfqpoint{3.683042in}{0.617538in}}%
\pgfpathlineto{\pgfqpoint{3.685139in}{0.600778in}}%
\pgfpathlineto{\pgfqpoint{3.687236in}{0.549090in}}%
\pgfpathlineto{\pgfqpoint{3.691431in}{0.616797in}}%
\pgfpathlineto{\pgfqpoint{3.693528in}{0.633481in}}%
\pgfpathlineto{\pgfqpoint{3.695626in}{0.642994in}}%
\pgfpathlineto{\pgfqpoint{3.697723in}{0.647465in}}%
\pgfpathlineto{\pgfqpoint{3.699820in}{0.646817in}}%
\pgfpathlineto{\pgfqpoint{3.701918in}{0.640121in}}%
\pgfpathlineto{\pgfqpoint{3.704015in}{0.629745in}}%
\pgfpathlineto{\pgfqpoint{3.706113in}{0.614747in}}%
\pgfpathlineto{\pgfqpoint{3.708210in}{0.591347in}}%
\pgfpathlineto{\pgfqpoint{3.710307in}{0.539416in}}%
\pgfpathlineto{\pgfqpoint{3.712405in}{0.550162in}}%
\pgfpathlineto{\pgfqpoint{3.714502in}{0.578780in}}%
\pgfpathlineto{\pgfqpoint{3.716599in}{0.586503in}}%
\pgfpathlineto{\pgfqpoint{3.718697in}{0.584041in}}%
\pgfpathlineto{\pgfqpoint{3.720794in}{0.569764in}}%
\pgfpathlineto{\pgfqpoint{3.722891in}{0.514130in}}%
\pgfpathlineto{\pgfqpoint{3.724989in}{0.560134in}}%
\pgfpathlineto{\pgfqpoint{3.727086in}{0.593547in}}%
\pgfpathlineto{\pgfqpoint{3.729184in}{0.610584in}}%
\pgfpathlineto{\pgfqpoint{3.731281in}{0.620395in}}%
\pgfpathlineto{\pgfqpoint{3.733378in}{0.624918in}}%
\pgfpathlineto{\pgfqpoint{3.735476in}{0.623825in}}%
\pgfpathlineto{\pgfqpoint{3.737573in}{0.616347in}}%
\pgfpathlineto{\pgfqpoint{3.739670in}{0.604510in}}%
\pgfpathlineto{\pgfqpoint{3.741768in}{0.586003in}}%
\pgfpathlineto{\pgfqpoint{3.745962in}{0.511838in}}%
\pgfpathlineto{\pgfqpoint{3.748060in}{0.563656in}}%
\pgfpathlineto{\pgfqpoint{3.750157in}{0.578591in}}%
\pgfpathlineto{\pgfqpoint{3.752254in}{0.583390in}}%
\pgfpathlineto{\pgfqpoint{3.754352in}{0.581371in}}%
\pgfpathlineto{\pgfqpoint{3.756449in}{0.571645in}}%
\pgfpathlineto{\pgfqpoint{3.758547in}{0.546214in}}%
\pgfpathlineto{\pgfqpoint{3.760644in}{0.497202in}}%
\pgfpathlineto{\pgfqpoint{3.762741in}{0.566471in}}%
\pgfpathlineto{\pgfqpoint{3.764839in}{0.589631in}}%
\pgfpathlineto{\pgfqpoint{3.766936in}{0.602487in}}%
\pgfpathlineto{\pgfqpoint{3.769033in}{0.609466in}}%
\pgfpathlineto{\pgfqpoint{3.771131in}{0.611369in}}%
\pgfpathlineto{\pgfqpoint{3.773228in}{0.606998in}}%
\pgfpathlineto{\pgfqpoint{3.775325in}{0.595695in}}%
\pgfpathlineto{\pgfqpoint{3.777423in}{0.573818in}}%
\pgfpathlineto{\pgfqpoint{3.779520in}{0.501728in}}%
\pgfpathlineto{\pgfqpoint{3.781618in}{0.559526in}}%
\pgfpathlineto{\pgfqpoint{3.783715in}{0.586809in}}%
\pgfpathlineto{\pgfqpoint{3.785812in}{0.599642in}}%
\pgfpathlineto{\pgfqpoint{3.787910in}{0.606213in}}%
\pgfpathlineto{\pgfqpoint{3.790007in}{0.608634in}}%
\pgfpathlineto{\pgfqpoint{3.792104in}{0.607465in}}%
\pgfpathlineto{\pgfqpoint{3.794202in}{0.602370in}}%
\pgfpathlineto{\pgfqpoint{3.796299in}{0.591720in}}%
\pgfpathlineto{\pgfqpoint{3.798396in}{0.569681in}}%
\pgfpathlineto{\pgfqpoint{3.800494in}{0.450437in}}%
\pgfpathlineto{\pgfqpoint{3.802591in}{0.569948in}}%
\pgfpathlineto{\pgfqpoint{3.804689in}{0.596785in}}%
\pgfpathlineto{\pgfqpoint{3.806786in}{0.611110in}}%
\pgfpathlineto{\pgfqpoint{3.808883in}{0.619409in}}%
\pgfpathlineto{\pgfqpoint{3.810981in}{0.623236in}}%
\pgfpathlineto{\pgfqpoint{3.813078in}{0.622441in}}%
\pgfpathlineto{\pgfqpoint{3.815175in}{0.615239in}}%
\pgfpathlineto{\pgfqpoint{3.817273in}{0.598561in}}%
\pgfpathlineto{\pgfqpoint{3.819370in}{0.554412in}}%
\pgfpathlineto{\pgfqpoint{3.821467in}{0.567088in}}%
\pgfpathlineto{\pgfqpoint{3.823565in}{0.603650in}}%
\pgfpathlineto{\pgfqpoint{3.825662in}{0.620025in}}%
\pgfpathlineto{\pgfqpoint{3.827760in}{0.629224in}}%
\pgfpathlineto{\pgfqpoint{3.829857in}{0.634213in}}%
\pgfpathlineto{\pgfqpoint{3.831954in}{0.636015in}}%
\pgfpathlineto{\pgfqpoint{3.834052in}{0.634894in}}%
\pgfpathlineto{\pgfqpoint{3.836149in}{0.630555in}}%
\pgfpathlineto{\pgfqpoint{3.838246in}{0.621885in}}%
\pgfpathlineto{\pgfqpoint{3.840344in}{0.605546in}}%
\pgfpathlineto{\pgfqpoint{3.842441in}{0.564976in}}%
\pgfpathlineto{\pgfqpoint{3.844538in}{0.571054in}}%
\pgfpathlineto{\pgfqpoint{3.846636in}{0.611136in}}%
\pgfpathlineto{\pgfqpoint{3.848733in}{0.629402in}}%
\pgfpathlineto{\pgfqpoint{3.850831in}{0.640135in}}%
\pgfpathlineto{\pgfqpoint{3.852928in}{0.646242in}}%
\pgfpathlineto{\pgfqpoint{3.855025in}{0.648397in}}%
\pgfpathlineto{\pgfqpoint{3.857123in}{0.646008in}}%
\pgfpathlineto{\pgfqpoint{3.859220in}{0.638422in}}%
\pgfpathlineto{\pgfqpoint{3.861317in}{0.626967in}}%
\pgfpathlineto{\pgfqpoint{3.863415in}{0.609791in}}%
\pgfpathlineto{\pgfqpoint{3.865512in}{0.580243in}}%
\pgfpathlineto{\pgfqpoint{3.867609in}{0.448333in}}%
\pgfpathlineto{\pgfqpoint{3.869707in}{0.572712in}}%
\pgfpathlineto{\pgfqpoint{3.871804in}{0.590609in}}%
\pgfpathlineto{\pgfqpoint{3.873902in}{0.595635in}}%
\pgfpathlineto{\pgfqpoint{3.875999in}{0.592432in}}%
\pgfpathlineto{\pgfqpoint{3.878096in}{0.579072in}}%
\pgfpathlineto{\pgfqpoint{3.880194in}{0.535881in}}%
\pgfpathlineto{\pgfqpoint{3.882291in}{0.557070in}}%
\pgfpathlineto{\pgfqpoint{3.884388in}{0.594814in}}%
\pgfpathlineto{\pgfqpoint{3.886486in}{0.612922in}}%
\pgfpathlineto{\pgfqpoint{3.888583in}{0.623291in}}%
\pgfpathlineto{\pgfqpoint{3.890680in}{0.628304in}}%
\pgfpathlineto{\pgfqpoint{3.892778in}{0.627897in}}%
\pgfpathlineto{\pgfqpoint{3.894875in}{0.621101in}}%
\pgfpathlineto{\pgfqpoint{3.896973in}{0.610327in}}%
\pgfpathlineto{\pgfqpoint{3.899070in}{0.594289in}}%
\pgfpathlineto{\pgfqpoint{3.901167in}{0.567705in}}%
\pgfpathlineto{\pgfqpoint{3.903265in}{0.481752in}}%
\pgfpathlineto{\pgfqpoint{3.905362in}{0.548725in}}%
\pgfpathlineto{\pgfqpoint{3.907459in}{0.570140in}}%
\pgfpathlineto{\pgfqpoint{3.909557in}{0.576655in}}%
\pgfpathlineto{\pgfqpoint{3.911654in}{0.574735in}}%
\pgfpathlineto{\pgfqpoint{3.913751in}{0.563413in}}%
\pgfpathlineto{\pgfqpoint{3.915849in}{0.529427in}}%
\pgfpathlineto{\pgfqpoint{3.917946in}{0.526818in}}%
\pgfpathlineto{\pgfqpoint{3.920043in}{0.572038in}}%
\pgfpathlineto{\pgfqpoint{3.922141in}{0.591928in}}%
\pgfpathlineto{\pgfqpoint{3.924238in}{0.603273in}}%
\pgfpathlineto{\pgfqpoint{3.926336in}{0.609089in}}%
\pgfpathlineto{\pgfqpoint{3.928433in}{0.609627in}}%
\pgfpathlineto{\pgfqpoint{3.930530in}{0.603416in}}%
\pgfpathlineto{\pgfqpoint{3.932628in}{0.590495in}}%
\pgfpathlineto{\pgfqpoint{3.934725in}{0.565644in}}%
\pgfpathlineto{\pgfqpoint{3.936822in}{0.429244in}}%
\pgfpathlineto{\pgfqpoint{3.938920in}{0.561180in}}%
\pgfpathlineto{\pgfqpoint{3.941017in}{0.584762in}}%
\pgfpathlineto{\pgfqpoint{3.943114in}{0.596023in}}%
\pgfpathlineto{\pgfqpoint{3.945212in}{0.601416in}}%
\pgfpathlineto{\pgfqpoint{3.947309in}{0.602664in}}%
\pgfpathlineto{\pgfqpoint{3.949407in}{0.600063in}}%
\pgfpathlineto{\pgfqpoint{3.951504in}{0.592844in}}%
\pgfpathlineto{\pgfqpoint{3.953601in}{0.578104in}}%
\pgfpathlineto{\pgfqpoint{3.955699in}{0.542008in}}%
\pgfpathlineto{\pgfqpoint{3.957796in}{0.536622in}}%
\pgfpathlineto{\pgfqpoint{3.959893in}{0.581009in}}%
\pgfpathlineto{\pgfqpoint{3.961991in}{0.600163in}}%
\pgfpathlineto{\pgfqpoint{3.964088in}{0.611133in}}%
\pgfpathlineto{\pgfqpoint{3.966185in}{0.617053in}}%
\pgfpathlineto{\pgfqpoint{3.968283in}{0.618484in}}%
\pgfpathlineto{\pgfqpoint{3.970380in}{0.614338in}}%
\pgfpathlineto{\pgfqpoint{3.972478in}{0.602068in}}%
\pgfpathlineto{\pgfqpoint{3.974575in}{0.573605in}}%
\pgfpathlineto{\pgfqpoint{3.976672in}{0.527055in}}%
\pgfpathlineto{\pgfqpoint{3.978770in}{0.590979in}}%
\pgfpathlineto{\pgfqpoint{3.980867in}{0.612011in}}%
\pgfpathlineto{\pgfqpoint{3.982964in}{0.623477in}}%
\pgfpathlineto{\pgfqpoint{3.985062in}{0.630034in}}%
\pgfpathlineto{\pgfqpoint{3.987159in}{0.633190in}}%
\pgfpathlineto{\pgfqpoint{3.989256in}{0.633464in}}%
\pgfpathlineto{\pgfqpoint{3.991354in}{0.630825in}}%
\pgfpathlineto{\pgfqpoint{3.993451in}{0.624659in}}%
\pgfpathlineto{\pgfqpoint{3.995549in}{0.613172in}}%
\pgfpathlineto{\pgfqpoint{3.997646in}{0.590314in}}%
\pgfpathlineto{\pgfqpoint{3.999743in}{0.452327in}}%
\pgfpathlineto{\pgfqpoint{4.001841in}{0.591205in}}%
\pgfpathlineto{\pgfqpoint{4.003938in}{0.617648in}}%
\pgfpathlineto{\pgfqpoint{4.006035in}{0.631958in}}%
\pgfpathlineto{\pgfqpoint{4.008133in}{0.640513in}}%
\pgfpathlineto{\pgfqpoint{4.010230in}{0.644959in}}%
\pgfpathlineto{\pgfqpoint{4.012327in}{0.645442in}}%
\pgfpathlineto{\pgfqpoint{4.014425in}{0.640866in}}%
\pgfpathlineto{\pgfqpoint{4.016522in}{0.631768in}}%
\pgfpathlineto{\pgfqpoint{4.018620in}{0.618105in}}%
\pgfpathlineto{\pgfqpoint{4.020717in}{0.596453in}}%
\pgfpathlineto{\pgfqpoint{4.022814in}{0.549646in}}%
\pgfpathlineto{\pgfqpoint{4.024912in}{0.551935in}}%
\pgfpathlineto{\pgfqpoint{4.027009in}{0.584920in}}%
\pgfpathlineto{\pgfqpoint{4.029106in}{0.595674in}}%
\pgfpathlineto{\pgfqpoint{4.031204in}{0.597458in}}%
\pgfpathlineto{\pgfqpoint{4.033301in}{0.591729in}}%
\pgfpathlineto{\pgfqpoint{4.035398in}{0.574711in}}%
\pgfpathlineto{\pgfqpoint{4.037496in}{0.507456in}}%
\pgfpathlineto{\pgfqpoint{4.039593in}{0.568432in}}%
\pgfpathlineto{\pgfqpoint{4.041691in}{0.599320in}}%
\pgfpathlineto{\pgfqpoint{4.043788in}{0.615266in}}%
\pgfpathlineto{\pgfqpoint{4.045885in}{0.624323in}}%
\pgfpathlineto{\pgfqpoint{4.047983in}{0.628201in}}%
\pgfpathlineto{\pgfqpoint{4.050080in}{0.626418in}}%
\pgfpathlineto{\pgfqpoint{4.052177in}{0.618589in}}%
\pgfpathlineto{\pgfqpoint{4.054275in}{0.606646in}}%
\pgfpathlineto{\pgfqpoint{4.056372in}{0.588324in}}%
\pgfpathlineto{\pgfqpoint{4.058469in}{0.554722in}}%
\pgfpathlineto{\pgfqpoint{4.060567in}{0.502017in}}%
\pgfpathlineto{\pgfqpoint{4.062664in}{0.561460in}}%
\pgfpathlineto{\pgfqpoint{4.064762in}{0.576717in}}%
\pgfpathlineto{\pgfqpoint{4.066859in}{0.581063in}}%
\pgfpathlineto{\pgfqpoint{4.068956in}{0.577930in}}%
\pgfpathlineto{\pgfqpoint{4.071054in}{0.565611in}}%
\pgfpathlineto{\pgfqpoint{4.073151in}{0.529505in}}%
\pgfpathlineto{\pgfqpoint{4.075248in}{0.531926in}}%
\pgfpathlineto{\pgfqpoint{4.077346in}{0.574996in}}%
\pgfpathlineto{\pgfqpoint{4.079443in}{0.594311in}}%
\pgfpathlineto{\pgfqpoint{4.081540in}{0.605315in}}%
\pgfpathlineto{\pgfqpoint{4.083638in}{0.610865in}}%
\pgfpathlineto{\pgfqpoint{4.085735in}{0.611142in}}%
\pgfpathlineto{\pgfqpoint{4.087832in}{0.604674in}}%
\pgfpathlineto{\pgfqpoint{4.089930in}{0.591431in}}%
\pgfpathlineto{\pgfqpoint{4.092027in}{0.565701in}}%
\pgfpathlineto{\pgfqpoint{4.094125in}{0.416105in}}%
\pgfpathlineto{\pgfqpoint{4.096222in}{0.564617in}}%
\pgfpathlineto{\pgfqpoint{4.098319in}{0.587314in}}%
\pgfpathlineto{\pgfqpoint{4.100417in}{0.598264in}}%
\pgfpathlineto{\pgfqpoint{4.102514in}{0.603486in}}%
\pgfpathlineto{\pgfqpoint{4.104611in}{0.604613in}}%
\pgfpathlineto{\pgfqpoint{4.106709in}{0.601904in}}%
\pgfpathlineto{\pgfqpoint{4.108806in}{0.594560in}}%
\pgfpathlineto{\pgfqpoint{4.110903in}{0.579604in}}%
\pgfpathlineto{\pgfqpoint{4.113001in}{0.542636in}}%
\pgfpathlineto{\pgfqpoint{4.115098in}{0.540020in}}%
\pgfpathlineto{\pgfqpoint{4.117196in}{0.583352in}}%
\pgfpathlineto{\pgfqpoint{4.119293in}{0.602288in}}%
\pgfpathlineto{\pgfqpoint{4.121390in}{0.613150in}}%
\pgfpathlineto{\pgfqpoint{4.123488in}{0.618994in}}%
\pgfpathlineto{\pgfqpoint{4.125585in}{0.620355in}}%
\pgfpathlineto{\pgfqpoint{4.127682in}{0.616125in}}%
\pgfpathlineto{\pgfqpoint{4.129780in}{0.603972in}}%
\pgfpathlineto{\pgfqpoint{4.131877in}{0.576540in}}%
\pgfpathlineto{\pgfqpoint{4.133974in}{0.518801in}}%
\pgfpathlineto{\pgfqpoint{4.136072in}{0.590027in}}%
\pgfpathlineto{\pgfqpoint{4.138169in}{0.611592in}}%
\pgfpathlineto{\pgfqpoint{4.140267in}{0.623170in}}%
\pgfpathlineto{\pgfqpoint{4.142364in}{0.629704in}}%
\pgfpathlineto{\pgfqpoint{4.144461in}{0.632758in}}%
\pgfpathlineto{\pgfqpoint{4.146559in}{0.632853in}}%
\pgfpathlineto{\pgfqpoint{4.148656in}{0.629930in}}%
\pgfpathlineto{\pgfqpoint{4.150753in}{0.623289in}}%
\pgfpathlineto{\pgfqpoint{4.152851in}{0.610865in}}%
\pgfpathlineto{\pgfqpoint{4.154948in}{0.585193in}}%
\pgfpathlineto{\pgfqpoint{4.157045in}{0.509487in}}%
\pgfpathlineto{\pgfqpoint{4.159143in}{0.596191in}}%
\pgfpathlineto{\pgfqpoint{4.161240in}{0.620089in}}%
\pgfpathlineto{\pgfqpoint{4.163338in}{0.633394in}}%
\pgfpathlineto{\pgfqpoint{4.165435in}{0.641279in}}%
\pgfpathlineto{\pgfqpoint{4.167532in}{0.645095in}}%
\pgfpathlineto{\pgfqpoint{4.169630in}{0.644787in}}%
\pgfpathlineto{\pgfqpoint{4.171727in}{0.639166in}}%
\pgfpathlineto{\pgfqpoint{4.173824in}{0.629260in}}%
\pgfpathlineto{\pgfqpoint{4.175922in}{0.614306in}}%
\pgfpathlineto{\pgfqpoint{4.178019in}{0.589726in}}%
\pgfpathlineto{\pgfqpoint{4.180116in}{0.522939in}}%
\pgfpathlineto{\pgfqpoint{4.182214in}{0.564325in}}%
\pgfpathlineto{\pgfqpoint{4.184311in}{0.589099in}}%
\pgfpathlineto{\pgfqpoint{4.186409in}{0.597556in}}%
\pgfpathlineto{\pgfqpoint{4.188506in}{0.597929in}}%
\pgfpathlineto{\pgfqpoint{4.190603in}{0.590731in}}%
\pgfpathlineto{\pgfqpoint{4.192701in}{0.570791in}}%
\pgfpathlineto{\pgfqpoint{4.194798in}{0.428031in}}%
\pgfpathlineto{\pgfqpoint{4.196895in}{0.575560in}}%
\pgfpathlineto{\pgfqpoint{4.198993in}{0.602807in}}%
\pgfpathlineto{\pgfqpoint{4.201090in}{0.617416in}}%
\pgfpathlineto{\pgfqpoint{4.203187in}{0.625609in}}%
\pgfpathlineto{\pgfqpoint{4.205285in}{0.628686in}}%
\pgfpathlineto{\pgfqpoint{4.207382in}{0.625836in}}%
\pgfpathlineto{\pgfqpoint{4.209480in}{0.617382in}}%
\pgfpathlineto{\pgfqpoint{4.211577in}{0.604597in}}%
\pgfpathlineto{\pgfqpoint{4.213674in}{0.584618in}}%
\pgfpathlineto{\pgfqpoint{4.215772in}{0.544921in}}%
\pgfpathlineto{\pgfqpoint{4.217869in}{0.525412in}}%
\pgfpathlineto{\pgfqpoint{4.219966in}{0.566702in}}%
\pgfpathlineto{\pgfqpoint{4.222064in}{0.579430in}}%
\pgfpathlineto{\pgfqpoint{4.224161in}{0.582544in}}%
\pgfpathlineto{\pgfqpoint{4.226258in}{0.578382in}}%
\pgfpathlineto{\pgfqpoint{4.228356in}{0.564540in}}%
\pgfpathlineto{\pgfqpoint{4.230453in}{0.521800in}}%
\pgfpathlineto{\pgfqpoint{4.232551in}{0.540715in}}%
\pgfpathlineto{\pgfqpoint{4.234648in}{0.578718in}}%
\pgfpathlineto{\pgfqpoint{4.236745in}{0.596779in}}%
\pgfpathlineto{\pgfqpoint{4.238843in}{0.607087in}}%
\pgfpathlineto{\pgfqpoint{4.240940in}{0.612065in}}%
\pgfpathlineto{\pgfqpoint{4.243037in}{0.611681in}}%
\pgfpathlineto{\pgfqpoint{4.245135in}{0.604485in}}%
\pgfpathlineto{\pgfqpoint{4.247232in}{0.590430in}}%
\pgfpathlineto{\pgfqpoint{4.249329in}{0.562658in}}%
\pgfpathlineto{\pgfqpoint{4.251427in}{0.482918in}}%
\pgfpathlineto{\pgfqpoint{4.253524in}{0.568035in}}%
\pgfpathlineto{\pgfqpoint{4.255621in}{0.589033in}}%
\pgfpathlineto{\pgfqpoint{4.257719in}{0.599269in}}%
\pgfpathlineto{\pgfqpoint{4.259816in}{0.604005in}}%
\pgfpathlineto{\pgfqpoint{4.261914in}{0.604686in}}%
\pgfpathlineto{\pgfqpoint{4.264011in}{0.601459in}}%
\pgfpathlineto{\pgfqpoint{4.266108in}{0.593345in}}%
\pgfpathlineto{\pgfqpoint{4.268206in}{0.576777in}}%
\pgfpathlineto{\pgfqpoint{4.270303in}{0.531670in}}%
\pgfpathlineto{\pgfqpoint{4.272400in}{0.550395in}}%
\pgfpathlineto{\pgfqpoint{4.274498in}{0.587184in}}%
\pgfpathlineto{\pgfqpoint{4.276595in}{0.604591in}}%
\pgfpathlineto{\pgfqpoint{4.278692in}{0.614639in}}%
\pgfpathlineto{\pgfqpoint{4.280790in}{0.619836in}}%
\pgfpathlineto{\pgfqpoint{4.282887in}{0.620482in}}%
\pgfpathlineto{\pgfqpoint{4.284985in}{0.615471in}}%
\pgfpathlineto{\pgfqpoint{4.287082in}{0.607174in}}%
\pgfpathlineto{\pgfqpoint{4.291277in}{0.583123in}}%
\pgfpathlineto{\pgfqpoint{4.293374in}{0.565884in}}%
\pgfpathlineto{\pgfqpoint{4.295471in}{0.543186in}}%
\pgfpathlineto{\pgfqpoint{4.299666in}{0.474866in}}%
\pgfpathlineto{\pgfqpoint{4.301763in}{0.474063in}}%
\pgfpathlineto{\pgfqpoint{4.305958in}{0.536032in}}%
\pgfpathlineto{\pgfqpoint{4.308056in}{0.554983in}}%
\pgfpathlineto{\pgfqpoint{4.310153in}{0.567874in}}%
\pgfpathlineto{\pgfqpoint{4.312250in}{0.575930in}}%
\pgfpathlineto{\pgfqpoint{4.314348in}{0.579295in}}%
\pgfpathlineto{\pgfqpoint{4.316445in}{0.576717in}}%
\pgfpathlineto{\pgfqpoint{4.316445in}{0.576717in}}%
\pgfusepath{stroke}%
\end{pgfscope}%
\begin{pgfscope}%
\pgfpathrectangle{\pgfqpoint{2.554660in}{0.320625in}}{\pgfqpoint{1.845679in}{1.164375in}}%
\pgfusepath{clip}%
\pgfsetrectcap%
\pgfsetroundjoin%
\pgfsetlinewidth{1.154312pt}%
\definecolor{currentstroke}{rgb}{0.997351,0.719089,0.493755}%
\pgfsetstrokecolor{currentstroke}%
\pgfsetdash{}{0pt}%
\pgfpathmoveto{\pgfqpoint{2.638555in}{1.468525in}}%
\pgfpathlineto{\pgfqpoint{2.649042in}{1.449020in}}%
\pgfpathlineto{\pgfqpoint{2.659529in}{1.426689in}}%
\pgfpathlineto{\pgfqpoint{2.672113in}{1.396826in}}%
\pgfpathlineto{\pgfqpoint{2.714060in}{1.295102in}}%
\pgfpathlineto{\pgfqpoint{2.737131in}{1.242905in}}%
\pgfpathlineto{\pgfqpoint{2.764397in}{1.184202in}}%
\pgfpathlineto{\pgfqpoint{2.797955in}{1.114812in}}%
\pgfpathlineto{\pgfqpoint{2.841999in}{1.026516in}}%
\pgfpathlineto{\pgfqpoint{2.902823in}{0.907336in}}%
\pgfpathlineto{\pgfqpoint{2.978328in}{0.761770in}}%
\pgfpathlineto{\pgfqpoint{2.990912in}{0.737306in}}%
\pgfpathlineto{\pgfqpoint{3.003496in}{0.713371in}}%
\pgfpathlineto{\pgfqpoint{3.016080in}{0.690581in}}%
\pgfpathlineto{\pgfqpoint{3.034957in}{0.651376in}}%
\pgfpathlineto{\pgfqpoint{3.039151in}{0.645504in}}%
\pgfpathlineto{\pgfqpoint{3.045443in}{0.639267in}}%
\pgfpathlineto{\pgfqpoint{3.049638in}{0.634500in}}%
\pgfpathlineto{\pgfqpoint{3.053833in}{0.626257in}}%
\pgfpathlineto{\pgfqpoint{3.058027in}{0.615244in}}%
\pgfpathlineto{\pgfqpoint{3.062222in}{0.601390in}}%
\pgfpathlineto{\pgfqpoint{3.068514in}{0.575320in}}%
\pgfpathlineto{\pgfqpoint{3.072709in}{0.557112in}}%
\pgfpathlineto{\pgfqpoint{3.074806in}{0.549758in}}%
\pgfpathlineto{\pgfqpoint{3.076904in}{0.544705in}}%
\pgfpathlineto{\pgfqpoint{3.079001in}{0.542704in}}%
\pgfpathlineto{\pgfqpoint{3.081098in}{0.544411in}}%
\pgfpathlineto{\pgfqpoint{3.083196in}{0.549625in}}%
\pgfpathlineto{\pgfqpoint{3.091585in}{0.577878in}}%
\pgfpathlineto{\pgfqpoint{3.093683in}{0.582692in}}%
\pgfpathlineto{\pgfqpoint{3.095780in}{0.586081in}}%
\pgfpathlineto{\pgfqpoint{3.097877in}{0.587857in}}%
\pgfpathlineto{\pgfqpoint{3.099975in}{0.587683in}}%
\pgfpathlineto{\pgfqpoint{3.102072in}{0.584901in}}%
\pgfpathlineto{\pgfqpoint{3.104169in}{0.579145in}}%
\pgfpathlineto{\pgfqpoint{3.106267in}{0.570189in}}%
\pgfpathlineto{\pgfqpoint{3.110462in}{0.544015in}}%
\pgfpathlineto{\pgfqpoint{3.112559in}{0.547323in}}%
\pgfpathlineto{\pgfqpoint{3.116754in}{0.575543in}}%
\pgfpathlineto{\pgfqpoint{3.120948in}{0.592092in}}%
\pgfpathlineto{\pgfqpoint{3.123046in}{0.597191in}}%
\pgfpathlineto{\pgfqpoint{3.125143in}{0.600781in}}%
\pgfpathlineto{\pgfqpoint{3.127240in}{0.603095in}}%
\pgfpathlineto{\pgfqpoint{3.129338in}{0.604258in}}%
\pgfpathlineto{\pgfqpoint{3.131435in}{0.604321in}}%
\pgfpathlineto{\pgfqpoint{3.133533in}{0.603260in}}%
\pgfpathlineto{\pgfqpoint{3.135630in}{0.600976in}}%
\pgfpathlineto{\pgfqpoint{3.137727in}{0.597259in}}%
\pgfpathlineto{\pgfqpoint{3.139825in}{0.591720in}}%
\pgfpathlineto{\pgfqpoint{3.141922in}{0.583604in}}%
\pgfpathlineto{\pgfqpoint{3.144019in}{0.571251in}}%
\pgfpathlineto{\pgfqpoint{3.146117in}{0.549873in}}%
\pgfpathlineto{\pgfqpoint{3.148214in}{0.493255in}}%
\pgfpathlineto{\pgfqpoint{3.150311in}{0.539372in}}%
\pgfpathlineto{\pgfqpoint{3.152409in}{0.568614in}}%
\pgfpathlineto{\pgfqpoint{3.154506in}{0.584842in}}%
\pgfpathlineto{\pgfqpoint{3.156604in}{0.595751in}}%
\pgfpathlineto{\pgfqpoint{3.160798in}{0.609282in}}%
\pgfpathlineto{\pgfqpoint{3.162896in}{0.613259in}}%
\pgfpathlineto{\pgfqpoint{3.164993in}{0.615689in}}%
\pgfpathlineto{\pgfqpoint{3.167090in}{0.616564in}}%
\pgfpathlineto{\pgfqpoint{3.169188in}{0.615690in}}%
\pgfpathlineto{\pgfqpoint{3.171285in}{0.612592in}}%
\pgfpathlineto{\pgfqpoint{3.173382in}{0.606569in}}%
\pgfpathlineto{\pgfqpoint{3.175480in}{0.598461in}}%
\pgfpathlineto{\pgfqpoint{3.177577in}{0.593906in}}%
\pgfpathlineto{\pgfqpoint{3.179675in}{0.600064in}}%
\pgfpathlineto{\pgfqpoint{3.183869in}{0.621158in}}%
\pgfpathlineto{\pgfqpoint{3.185967in}{0.628774in}}%
\pgfpathlineto{\pgfqpoint{3.188064in}{0.634258in}}%
\pgfpathlineto{\pgfqpoint{3.190161in}{0.637912in}}%
\pgfpathlineto{\pgfqpoint{3.192259in}{0.639936in}}%
\pgfpathlineto{\pgfqpoint{3.194356in}{0.640415in}}%
\pgfpathlineto{\pgfqpoint{3.196453in}{0.639315in}}%
\pgfpathlineto{\pgfqpoint{3.198551in}{0.636460in}}%
\pgfpathlineto{\pgfqpoint{3.200648in}{0.631451in}}%
\pgfpathlineto{\pgfqpoint{3.202746in}{0.623448in}}%
\pgfpathlineto{\pgfqpoint{3.204843in}{0.610481in}}%
\pgfpathlineto{\pgfqpoint{3.206940in}{0.586220in}}%
\pgfpathlineto{\pgfqpoint{3.209038in}{0.512108in}}%
\pgfpathlineto{\pgfqpoint{3.211135in}{0.589469in}}%
\pgfpathlineto{\pgfqpoint{3.213232in}{0.615099in}}%
\pgfpathlineto{\pgfqpoint{3.215330in}{0.629752in}}%
\pgfpathlineto{\pgfqpoint{3.217427in}{0.639403in}}%
\pgfpathlineto{\pgfqpoint{3.219524in}{0.645871in}}%
\pgfpathlineto{\pgfqpoint{3.221622in}{0.649839in}}%
\pgfpathlineto{\pgfqpoint{3.223719in}{0.651462in}}%
\pgfpathlineto{\pgfqpoint{3.225816in}{0.650472in}}%
\pgfpathlineto{\pgfqpoint{3.227914in}{0.645992in}}%
\pgfpathlineto{\pgfqpoint{3.230011in}{0.638158in}}%
\pgfpathlineto{\pgfqpoint{3.232109in}{0.627381in}}%
\pgfpathlineto{\pgfqpoint{3.238401in}{0.586277in}}%
\pgfpathlineto{\pgfqpoint{3.242595in}{0.596800in}}%
\pgfpathlineto{\pgfqpoint{3.244693in}{0.598099in}}%
\pgfpathlineto{\pgfqpoint{3.246790in}{0.593628in}}%
\pgfpathlineto{\pgfqpoint{3.248887in}{0.581265in}}%
\pgfpathlineto{\pgfqpoint{3.250985in}{0.557891in}}%
\pgfpathlineto{\pgfqpoint{3.253082in}{0.569312in}}%
\pgfpathlineto{\pgfqpoint{3.255180in}{0.596174in}}%
\pgfpathlineto{\pgfqpoint{3.257277in}{0.612442in}}%
\pgfpathlineto{\pgfqpoint{3.259374in}{0.622199in}}%
\pgfpathlineto{\pgfqpoint{3.261472in}{0.627024in}}%
\pgfpathlineto{\pgfqpoint{3.263569in}{0.626754in}}%
\pgfpathlineto{\pgfqpoint{3.265666in}{0.620323in}}%
\pgfpathlineto{\pgfqpoint{3.267764in}{0.609975in}}%
\pgfpathlineto{\pgfqpoint{3.269861in}{0.594794in}}%
\pgfpathlineto{\pgfqpoint{3.271958in}{0.571500in}}%
\pgfpathlineto{\pgfqpoint{3.274056in}{0.541785in}}%
\pgfpathlineto{\pgfqpoint{3.278251in}{0.574294in}}%
\pgfpathlineto{\pgfqpoint{3.280348in}{0.580408in}}%
\pgfpathlineto{\pgfqpoint{3.282445in}{0.579511in}}%
\pgfpathlineto{\pgfqpoint{3.284543in}{0.571151in}}%
\pgfpathlineto{\pgfqpoint{3.286640in}{0.549383in}}%
\pgfpathlineto{\pgfqpoint{3.288737in}{0.505861in}}%
\pgfpathlineto{\pgfqpoint{3.290835in}{0.562596in}}%
\pgfpathlineto{\pgfqpoint{3.292932in}{0.586830in}}%
\pgfpathlineto{\pgfqpoint{3.295029in}{0.600376in}}%
\pgfpathlineto{\pgfqpoint{3.297127in}{0.607948in}}%
\pgfpathlineto{\pgfqpoint{3.299224in}{0.610526in}}%
\pgfpathlineto{\pgfqpoint{3.301322in}{0.607128in}}%
\pgfpathlineto{\pgfqpoint{3.303419in}{0.596490in}}%
\pgfpathlineto{\pgfqpoint{3.305516in}{0.575190in}}%
\pgfpathlineto{\pgfqpoint{3.307614in}{0.518887in}}%
\pgfpathlineto{\pgfqpoint{3.309711in}{0.564813in}}%
\pgfpathlineto{\pgfqpoint{3.311808in}{0.590984in}}%
\pgfpathlineto{\pgfqpoint{3.313906in}{0.604019in}}%
\pgfpathlineto{\pgfqpoint{3.316003in}{0.611096in}}%
\pgfpathlineto{\pgfqpoint{3.318100in}{0.614257in}}%
\pgfpathlineto{\pgfqpoint{3.320198in}{0.614145in}}%
\pgfpathlineto{\pgfqpoint{3.322295in}{0.610671in}}%
\pgfpathlineto{\pgfqpoint{3.324393in}{0.602913in}}%
\pgfpathlineto{\pgfqpoint{3.326490in}{0.587975in}}%
\pgfpathlineto{\pgfqpoint{3.328587in}{0.552779in}}%
\pgfpathlineto{\pgfqpoint{3.330685in}{0.542958in}}%
\pgfpathlineto{\pgfqpoint{3.332782in}{0.588865in}}%
\pgfpathlineto{\pgfqpoint{3.334879in}{0.608337in}}%
\pgfpathlineto{\pgfqpoint{3.336977in}{0.619596in}}%
\pgfpathlineto{\pgfqpoint{3.339074in}{0.625960in}}%
\pgfpathlineto{\pgfqpoint{3.341171in}{0.628180in}}%
\pgfpathlineto{\pgfqpoint{3.343269in}{0.625618in}}%
\pgfpathlineto{\pgfqpoint{3.345366in}{0.617139in}}%
\pgfpathlineto{\pgfqpoint{3.347464in}{0.602126in}}%
\pgfpathlineto{\pgfqpoint{3.349561in}{0.573232in}}%
\pgfpathlineto{\pgfqpoint{3.351658in}{0.521101in}}%
\pgfpathlineto{\pgfqpoint{3.353756in}{0.580605in}}%
\pgfpathlineto{\pgfqpoint{3.355853in}{0.600055in}}%
\pgfpathlineto{\pgfqpoint{3.357950in}{0.609306in}}%
\pgfpathlineto{\pgfqpoint{3.360048in}{0.613055in}}%
\pgfpathlineto{\pgfqpoint{3.362145in}{0.612523in}}%
\pgfpathlineto{\pgfqpoint{3.364242in}{0.607538in}}%
\pgfpathlineto{\pgfqpoint{3.366340in}{0.596292in}}%
\pgfpathlineto{\pgfqpoint{3.368437in}{0.571453in}}%
\pgfpathlineto{\pgfqpoint{3.370535in}{0.497944in}}%
\pgfpathlineto{\pgfqpoint{3.372632in}{0.583484in}}%
\pgfpathlineto{\pgfqpoint{3.374729in}{0.607486in}}%
\pgfpathlineto{\pgfqpoint{3.376827in}{0.620660in}}%
\pgfpathlineto{\pgfqpoint{3.378924in}{0.628077in}}%
\pgfpathlineto{\pgfqpoint{3.381021in}{0.630873in}}%
\pgfpathlineto{\pgfqpoint{3.383119in}{0.628433in}}%
\pgfpathlineto{\pgfqpoint{3.385216in}{0.619749in}}%
\pgfpathlineto{\pgfqpoint{3.387313in}{0.604338in}}%
\pgfpathlineto{\pgfqpoint{3.389411in}{0.574027in}}%
\pgfpathlineto{\pgfqpoint{3.391508in}{0.525127in}}%
\pgfpathlineto{\pgfqpoint{3.393605in}{0.585025in}}%
\pgfpathlineto{\pgfqpoint{3.395703in}{0.603920in}}%
\pgfpathlineto{\pgfqpoint{3.397800in}{0.612917in}}%
\pgfpathlineto{\pgfqpoint{3.399898in}{0.616500in}}%
\pgfpathlineto{\pgfqpoint{3.401995in}{0.615827in}}%
\pgfpathlineto{\pgfqpoint{3.404092in}{0.610682in}}%
\pgfpathlineto{\pgfqpoint{3.406190in}{0.599186in}}%
\pgfpathlineto{\pgfqpoint{3.408287in}{0.573637in}}%
\pgfpathlineto{\pgfqpoint{3.410384in}{0.508956in}}%
\pgfpathlineto{\pgfqpoint{3.412482in}{0.587782in}}%
\pgfpathlineto{\pgfqpoint{3.414579in}{0.611299in}}%
\pgfpathlineto{\pgfqpoint{3.416676in}{0.624263in}}%
\pgfpathlineto{\pgfqpoint{3.418774in}{0.631534in}}%
\pgfpathlineto{\pgfqpoint{3.420871in}{0.634186in}}%
\pgfpathlineto{\pgfqpoint{3.422969in}{0.631549in}}%
\pgfpathlineto{\pgfqpoint{3.425066in}{0.622684in}}%
\pgfpathlineto{\pgfqpoint{3.427163in}{0.606991in}}%
\pgfpathlineto{\pgfqpoint{3.429261in}{0.575767in}}%
\pgfpathlineto{\pgfqpoint{3.431358in}{0.530180in}}%
\pgfpathlineto{\pgfqpoint{3.433455in}{0.588908in}}%
\pgfpathlineto{\pgfqpoint{3.435553in}{0.607424in}}%
\pgfpathlineto{\pgfqpoint{3.437650in}{0.616224in}}%
\pgfpathlineto{\pgfqpoint{3.439747in}{0.619654in}}%
\pgfpathlineto{\pgfqpoint{3.441845in}{0.618819in}}%
\pgfpathlineto{\pgfqpoint{3.443942in}{0.613460in}}%
\pgfpathlineto{\pgfqpoint{3.446040in}{0.601576in}}%
\pgfpathlineto{\pgfqpoint{3.448137in}{0.574822in}}%
\pgfpathlineto{\pgfqpoint{3.450234in}{0.522727in}}%
\pgfpathlineto{\pgfqpoint{3.452332in}{0.592517in}}%
\pgfpathlineto{\pgfqpoint{3.454429in}{0.615277in}}%
\pgfpathlineto{\pgfqpoint{3.456526in}{0.627910in}}%
\pgfpathlineto{\pgfqpoint{3.458624in}{0.634945in}}%
\pgfpathlineto{\pgfqpoint{3.460721in}{0.637359in}}%
\pgfpathlineto{\pgfqpoint{3.462818in}{0.634381in}}%
\pgfpathlineto{\pgfqpoint{3.464916in}{0.625426in}}%
\pgfpathlineto{\pgfqpoint{3.467013in}{0.610210in}}%
\pgfpathlineto{\pgfqpoint{3.469111in}{0.581595in}}%
\pgfpathlineto{\pgfqpoint{3.471208in}{0.503027in}}%
\pgfpathlineto{\pgfqpoint{3.473305in}{0.584416in}}%
\pgfpathlineto{\pgfqpoint{3.475403in}{0.604393in}}%
\pgfpathlineto{\pgfqpoint{3.477500in}{0.613378in}}%
\pgfpathlineto{\pgfqpoint{3.479597in}{0.616514in}}%
\pgfpathlineto{\pgfqpoint{3.481695in}{0.614970in}}%
\pgfpathlineto{\pgfqpoint{3.483792in}{0.608238in}}%
\pgfpathlineto{\pgfqpoint{3.485889in}{0.593271in}}%
\pgfpathlineto{\pgfqpoint{3.487987in}{0.553621in}}%
\pgfpathlineto{\pgfqpoint{3.490084in}{0.560557in}}%
\pgfpathlineto{\pgfqpoint{3.492182in}{0.600949in}}%
\pgfpathlineto{\pgfqpoint{3.494279in}{0.619286in}}%
\pgfpathlineto{\pgfqpoint{3.496376in}{0.629746in}}%
\pgfpathlineto{\pgfqpoint{3.498474in}{0.635075in}}%
\pgfpathlineto{\pgfqpoint{3.500571in}{0.635556in}}%
\pgfpathlineto{\pgfqpoint{3.502668in}{0.629916in}}%
\pgfpathlineto{\pgfqpoint{3.504766in}{0.618970in}}%
\pgfpathlineto{\pgfqpoint{3.506863in}{0.600756in}}%
\pgfpathlineto{\pgfqpoint{3.508960in}{0.563343in}}%
\pgfpathlineto{\pgfqpoint{3.511058in}{0.543116in}}%
\pgfpathlineto{\pgfqpoint{3.513155in}{0.587263in}}%
\pgfpathlineto{\pgfqpoint{3.515253in}{0.602571in}}%
\pgfpathlineto{\pgfqpoint{3.517350in}{0.609031in}}%
\pgfpathlineto{\pgfqpoint{3.519447in}{0.609956in}}%
\pgfpathlineto{\pgfqpoint{3.521545in}{0.605753in}}%
\pgfpathlineto{\pgfqpoint{3.523642in}{0.594716in}}%
\pgfpathlineto{\pgfqpoint{3.525739in}{0.568766in}}%
\pgfpathlineto{\pgfqpoint{3.527837in}{0.514514in}}%
\pgfpathlineto{\pgfqpoint{3.529934in}{0.586349in}}%
\pgfpathlineto{\pgfqpoint{3.532031in}{0.609429in}}%
\pgfpathlineto{\pgfqpoint{3.534129in}{0.622181in}}%
\pgfpathlineto{\pgfqpoint{3.536226in}{0.629176in}}%
\pgfpathlineto{\pgfqpoint{3.538324in}{0.631324in}}%
\pgfpathlineto{\pgfqpoint{3.540421in}{0.627643in}}%
\pgfpathlineto{\pgfqpoint{3.542518in}{0.618126in}}%
\pgfpathlineto{\pgfqpoint{3.544616in}{0.602157in}}%
\pgfpathlineto{\pgfqpoint{3.546713in}{0.571692in}}%
\pgfpathlineto{\pgfqpoint{3.548810in}{0.510415in}}%
\pgfpathlineto{\pgfqpoint{3.550908in}{0.579181in}}%
\pgfpathlineto{\pgfqpoint{3.553005in}{0.597863in}}%
\pgfpathlineto{\pgfqpoint{3.555102in}{0.606137in}}%
\pgfpathlineto{\pgfqpoint{3.557200in}{0.608631in}}%
\pgfpathlineto{\pgfqpoint{3.559297in}{0.606311in}}%
\pgfpathlineto{\pgfqpoint{3.561395in}{0.598363in}}%
\pgfpathlineto{\pgfqpoint{3.563492in}{0.580673in}}%
\pgfpathlineto{\pgfqpoint{3.565589in}{0.523463in}}%
\pgfpathlineto{\pgfqpoint{3.569784in}{0.598690in}}%
\pgfpathlineto{\pgfqpoint{3.571881in}{0.614862in}}%
\pgfpathlineto{\pgfqpoint{3.573979in}{0.624091in}}%
\pgfpathlineto{\pgfqpoint{3.576076in}{0.628369in}}%
\pgfpathlineto{\pgfqpoint{3.578173in}{0.627569in}}%
\pgfpathlineto{\pgfqpoint{3.580271in}{0.620375in}}%
\pgfpathlineto{\pgfqpoint{3.582368in}{0.607057in}}%
\pgfpathlineto{\pgfqpoint{3.584465in}{0.582492in}}%
\pgfpathlineto{\pgfqpoint{3.586563in}{0.483718in}}%
\pgfpathlineto{\pgfqpoint{3.588660in}{0.573187in}}%
\pgfpathlineto{\pgfqpoint{3.590758in}{0.597209in}}%
\pgfpathlineto{\pgfqpoint{3.592855in}{0.608141in}}%
\pgfpathlineto{\pgfqpoint{3.594952in}{0.612893in}}%
\pgfpathlineto{\pgfqpoint{3.597050in}{0.613168in}}%
\pgfpathlineto{\pgfqpoint{3.599147in}{0.609044in}}%
\pgfpathlineto{\pgfqpoint{3.601244in}{0.599085in}}%
\pgfpathlineto{\pgfqpoint{3.603342in}{0.577573in}}%
\pgfpathlineto{\pgfqpoint{3.605439in}{0.460729in}}%
\pgfpathlineto{\pgfqpoint{3.607536in}{0.578239in}}%
\pgfpathlineto{\pgfqpoint{3.609634in}{0.605243in}}%
\pgfpathlineto{\pgfqpoint{3.611731in}{0.619589in}}%
\pgfpathlineto{\pgfqpoint{3.613829in}{0.627756in}}%
\pgfpathlineto{\pgfqpoint{3.615926in}{0.631216in}}%
\pgfpathlineto{\pgfqpoint{3.618023in}{0.629597in}}%
\pgfpathlineto{\pgfqpoint{3.620121in}{0.621621in}}%
\pgfpathlineto{\pgfqpoint{3.622218in}{0.606992in}}%
\pgfpathlineto{\pgfqpoint{3.624315in}{0.578403in}}%
\pgfpathlineto{\pgfqpoint{3.626413in}{0.505664in}}%
\pgfpathlineto{\pgfqpoint{3.628510in}{0.584614in}}%
\pgfpathlineto{\pgfqpoint{3.630607in}{0.604932in}}%
\pgfpathlineto{\pgfqpoint{3.632705in}{0.614639in}}%
\pgfpathlineto{\pgfqpoint{3.634802in}{0.618805in}}%
\pgfpathlineto{\pgfqpoint{3.636900in}{0.618765in}}%
\pgfpathlineto{\pgfqpoint{3.638997in}{0.614490in}}%
\pgfpathlineto{\pgfqpoint{3.641094in}{0.604537in}}%
\pgfpathlineto{\pgfqpoint{3.643192in}{0.583390in}}%
\pgfpathlineto{\pgfqpoint{3.645289in}{0.483886in}}%
\pgfpathlineto{\pgfqpoint{3.647386in}{0.581892in}}%
\pgfpathlineto{\pgfqpoint{3.649484in}{0.609450in}}%
\pgfpathlineto{\pgfqpoint{3.651581in}{0.624006in}}%
\pgfpathlineto{\pgfqpoint{3.653678in}{0.632339in}}%
\pgfpathlineto{\pgfqpoint{3.655776in}{0.636010in}}%
\pgfpathlineto{\pgfqpoint{3.657873in}{0.634750in}}%
\pgfpathlineto{\pgfqpoint{3.659971in}{0.627298in}}%
\pgfpathlineto{\pgfqpoint{3.662068in}{0.613835in}}%
\pgfpathlineto{\pgfqpoint{3.664165in}{0.589186in}}%
\pgfpathlineto{\pgfqpoint{3.666263in}{0.491069in}}%
\pgfpathlineto{\pgfqpoint{3.668360in}{0.579340in}}%
\pgfpathlineto{\pgfqpoint{3.670457in}{0.603296in}}%
\pgfpathlineto{\pgfqpoint{3.672555in}{0.614084in}}%
\pgfpathlineto{\pgfqpoint{3.674652in}{0.618635in}}%
\pgfpathlineto{\pgfqpoint{3.676749in}{0.618627in}}%
\pgfpathlineto{\pgfqpoint{3.678847in}{0.614058in}}%
\pgfpathlineto{\pgfqpoint{3.680944in}{0.603255in}}%
\pgfpathlineto{\pgfqpoint{3.683042in}{0.579293in}}%
\pgfpathlineto{\pgfqpoint{3.685139in}{0.490331in}}%
\pgfpathlineto{\pgfqpoint{3.687236in}{0.589013in}}%
\pgfpathlineto{\pgfqpoint{3.689334in}{0.613638in}}%
\pgfpathlineto{\pgfqpoint{3.691431in}{0.627052in}}%
\pgfpathlineto{\pgfqpoint{3.693528in}{0.634588in}}%
\pgfpathlineto{\pgfqpoint{3.695626in}{0.637427in}}%
\pgfpathlineto{\pgfqpoint{3.697723in}{0.634942in}}%
\pgfpathlineto{\pgfqpoint{3.699820in}{0.626453in}}%
\pgfpathlineto{\pgfqpoint{3.701918in}{0.612285in}}%
\pgfpathlineto{\pgfqpoint{3.704015in}{0.587049in}}%
\pgfpathlineto{\pgfqpoint{3.706113in}{0.489553in}}%
\pgfpathlineto{\pgfqpoint{3.708210in}{0.575719in}}%
\pgfpathlineto{\pgfqpoint{3.710307in}{0.599045in}}%
\pgfpathlineto{\pgfqpoint{3.712405in}{0.609022in}}%
\pgfpathlineto{\pgfqpoint{3.714502in}{0.612514in}}%
\pgfpathlineto{\pgfqpoint{3.716599in}{0.610990in}}%
\pgfpathlineto{\pgfqpoint{3.718697in}{0.603929in}}%
\pgfpathlineto{\pgfqpoint{3.720794in}{0.587851in}}%
\pgfpathlineto{\pgfqpoint{3.722891in}{0.541136in}}%
\pgfpathlineto{\pgfqpoint{3.724989in}{0.565070in}}%
\pgfpathlineto{\pgfqpoint{3.727086in}{0.601003in}}%
\pgfpathlineto{\pgfqpoint{3.729184in}{0.618235in}}%
\pgfpathlineto{\pgfqpoint{3.731281in}{0.628067in}}%
\pgfpathlineto{\pgfqpoint{3.733378in}{0.632828in}}%
\pgfpathlineto{\pgfqpoint{3.735476in}{0.632559in}}%
\pgfpathlineto{\pgfqpoint{3.737573in}{0.626010in}}%
\pgfpathlineto{\pgfqpoint{3.739670in}{0.614202in}}%
\pgfpathlineto{\pgfqpoint{3.741768in}{0.594336in}}%
\pgfpathlineto{\pgfqpoint{3.743865in}{0.549811in}}%
\pgfpathlineto{\pgfqpoint{3.745962in}{0.552263in}}%
\pgfpathlineto{\pgfqpoint{3.748060in}{0.587838in}}%
\pgfpathlineto{\pgfqpoint{3.750157in}{0.601344in}}%
\pgfpathlineto{\pgfqpoint{3.752254in}{0.606847in}}%
\pgfpathlineto{\pgfqpoint{3.754352in}{0.606995in}}%
\pgfpathlineto{\pgfqpoint{3.756449in}{0.601903in}}%
\pgfpathlineto{\pgfqpoint{3.758547in}{0.589372in}}%
\pgfpathlineto{\pgfqpoint{3.762741in}{0.535919in}}%
\pgfpathlineto{\pgfqpoint{3.764839in}{0.588575in}}%
\pgfpathlineto{\pgfqpoint{3.766936in}{0.609337in}}%
\pgfpathlineto{\pgfqpoint{3.769033in}{0.621017in}}%
\pgfpathlineto{\pgfqpoint{3.771131in}{0.627228in}}%
\pgfpathlineto{\pgfqpoint{3.773228in}{0.628561in}}%
\pgfpathlineto{\pgfqpoint{3.775325in}{0.623812in}}%
\pgfpathlineto{\pgfqpoint{3.777423in}{0.613222in}}%
\pgfpathlineto{\pgfqpoint{3.779520in}{0.594919in}}%
\pgfpathlineto{\pgfqpoint{3.781618in}{0.555193in}}%
\pgfpathlineto{\pgfqpoint{3.783715in}{0.545737in}}%
\pgfpathlineto{\pgfqpoint{3.785812in}{0.586335in}}%
\pgfpathlineto{\pgfqpoint{3.787910in}{0.601539in}}%
\pgfpathlineto{\pgfqpoint{3.790007in}{0.608471in}}%
\pgfpathlineto{\pgfqpoint{3.792104in}{0.610283in}}%
\pgfpathlineto{\pgfqpoint{3.794202in}{0.607590in}}%
\pgfpathlineto{\pgfqpoint{3.796299in}{0.599482in}}%
\pgfpathlineto{\pgfqpoint{3.798396in}{0.581931in}}%
\pgfpathlineto{\pgfqpoint{3.800494in}{0.527494in}}%
\pgfpathlineto{\pgfqpoint{3.804689in}{0.598354in}}%
\pgfpathlineto{\pgfqpoint{3.806786in}{0.614728in}}%
\pgfpathlineto{\pgfqpoint{3.808883in}{0.624103in}}%
\pgfpathlineto{\pgfqpoint{3.810981in}{0.628571in}}%
\pgfpathlineto{\pgfqpoint{3.813078in}{0.628105in}}%
\pgfpathlineto{\pgfqpoint{3.815175in}{0.621314in}}%
\pgfpathlineto{\pgfqpoint{3.817273in}{0.608168in}}%
\pgfpathlineto{\pgfqpoint{3.819370in}{0.583449in}}%
\pgfpathlineto{\pgfqpoint{3.821467in}{0.470314in}}%
\pgfpathlineto{\pgfqpoint{3.823565in}{0.576517in}}%
\pgfpathlineto{\pgfqpoint{3.825662in}{0.600301in}}%
\pgfpathlineto{\pgfqpoint{3.827760in}{0.611382in}}%
\pgfpathlineto{\pgfqpoint{3.829857in}{0.616440in}}%
\pgfpathlineto{\pgfqpoint{3.831954in}{0.617188in}}%
\pgfpathlineto{\pgfqpoint{3.834052in}{0.613818in}}%
\pgfpathlineto{\pgfqpoint{3.836149in}{0.605252in}}%
\pgfpathlineto{\pgfqpoint{3.838246in}{0.587391in}}%
\pgfpathlineto{\pgfqpoint{3.840344in}{0.532883in}}%
\pgfpathlineto{\pgfqpoint{3.844538in}{0.603236in}}%
\pgfpathlineto{\pgfqpoint{3.846636in}{0.619584in}}%
\pgfpathlineto{\pgfqpoint{3.848733in}{0.629002in}}%
\pgfpathlineto{\pgfqpoint{3.850831in}{0.633618in}}%
\pgfpathlineto{\pgfqpoint{3.852928in}{0.633507in}}%
\pgfpathlineto{\pgfqpoint{3.855025in}{0.627303in}}%
\pgfpathlineto{\pgfqpoint{3.857123in}{0.614991in}}%
\pgfpathlineto{\pgfqpoint{3.859220in}{0.592476in}}%
\pgfpathlineto{\pgfqpoint{3.861317in}{0.522780in}}%
\pgfpathlineto{\pgfqpoint{3.863415in}{0.575188in}}%
\pgfpathlineto{\pgfqpoint{3.865512in}{0.602118in}}%
\pgfpathlineto{\pgfqpoint{3.867609in}{0.614181in}}%
\pgfpathlineto{\pgfqpoint{3.869707in}{0.619735in}}%
\pgfpathlineto{\pgfqpoint{3.871804in}{0.620801in}}%
\pgfpathlineto{\pgfqpoint{3.873902in}{0.617679in}}%
\pgfpathlineto{\pgfqpoint{3.875999in}{0.609361in}}%
\pgfpathlineto{\pgfqpoint{3.878096in}{0.591896in}}%
\pgfpathlineto{\pgfqpoint{3.880194in}{0.539807in}}%
\pgfpathlineto{\pgfqpoint{3.884388in}{0.606738in}}%
\pgfpathlineto{\pgfqpoint{3.886486in}{0.623313in}}%
\pgfpathlineto{\pgfqpoint{3.888583in}{0.632847in}}%
\pgfpathlineto{\pgfqpoint{3.890680in}{0.637534in}}%
\pgfpathlineto{\pgfqpoint{3.892778in}{0.637463in}}%
\pgfpathlineto{\pgfqpoint{3.894875in}{0.631334in}}%
\pgfpathlineto{\pgfqpoint{3.896973in}{0.619718in}}%
\pgfpathlineto{\pgfqpoint{3.899070in}{0.599799in}}%
\pgfpathlineto{\pgfqpoint{3.901167in}{0.553664in}}%
\pgfpathlineto{\pgfqpoint{3.903265in}{0.561362in}}%
\pgfpathlineto{\pgfqpoint{3.905362in}{0.595812in}}%
\pgfpathlineto{\pgfqpoint{3.907459in}{0.609373in}}%
\pgfpathlineto{\pgfqpoint{3.909557in}{0.615205in}}%
\pgfpathlineto{\pgfqpoint{3.911654in}{0.615926in}}%
\pgfpathlineto{\pgfqpoint{3.913751in}{0.611819in}}%
\pgfpathlineto{\pgfqpoint{3.915849in}{0.601317in}}%
\pgfpathlineto{\pgfqpoint{3.917946in}{0.577510in}}%
\pgfpathlineto{\pgfqpoint{3.920043in}{0.490395in}}%
\pgfpathlineto{\pgfqpoint{3.922141in}{0.587761in}}%
\pgfpathlineto{\pgfqpoint{3.924238in}{0.612395in}}%
\pgfpathlineto{\pgfqpoint{3.926336in}{0.625809in}}%
\pgfpathlineto{\pgfqpoint{3.928433in}{0.633296in}}%
\pgfpathlineto{\pgfqpoint{3.930530in}{0.635996in}}%
\pgfpathlineto{\pgfqpoint{3.932628in}{0.633177in}}%
\pgfpathlineto{\pgfqpoint{3.934725in}{0.624465in}}%
\pgfpathlineto{\pgfqpoint{3.936822in}{0.610195in}}%
\pgfpathlineto{\pgfqpoint{3.938920in}{0.585104in}}%
\pgfpathlineto{\pgfqpoint{3.941017in}{0.495051in}}%
\pgfpathlineto{\pgfqpoint{3.943114in}{0.571735in}}%
\pgfpathlineto{\pgfqpoint{3.945212in}{0.595285in}}%
\pgfpathlineto{\pgfqpoint{3.947309in}{0.605111in}}%
\pgfpathlineto{\pgfqpoint{3.949407in}{0.608279in}}%
\pgfpathlineto{\pgfqpoint{3.951504in}{0.606210in}}%
\pgfpathlineto{\pgfqpoint{3.953601in}{0.598145in}}%
\pgfpathlineto{\pgfqpoint{3.955699in}{0.579560in}}%
\pgfpathlineto{\pgfqpoint{3.957796in}{0.511016in}}%
\pgfpathlineto{\pgfqpoint{3.959893in}{0.571473in}}%
\pgfpathlineto{\pgfqpoint{3.961991in}{0.601756in}}%
\pgfpathlineto{\pgfqpoint{3.964088in}{0.617290in}}%
\pgfpathlineto{\pgfqpoint{3.966185in}{0.626107in}}%
\pgfpathlineto{\pgfqpoint{3.968283in}{0.629952in}}%
\pgfpathlineto{\pgfqpoint{3.970380in}{0.628495in}}%
\pgfpathlineto{\pgfqpoint{3.972478in}{0.620707in}}%
\pgfpathlineto{\pgfqpoint{3.974575in}{0.607309in}}%
\pgfpathlineto{\pgfqpoint{3.976672in}{0.583660in}}%
\pgfpathlineto{\pgfqpoint{3.978770in}{0.509584in}}%
\pgfpathlineto{\pgfqpoint{3.980867in}{0.566360in}}%
\pgfpathlineto{\pgfqpoint{3.982964in}{0.591920in}}%
\pgfpathlineto{\pgfqpoint{3.985062in}{0.602831in}}%
\pgfpathlineto{\pgfqpoint{3.987159in}{0.607053in}}%
\pgfpathlineto{\pgfqpoint{3.989256in}{0.606340in}}%
\pgfpathlineto{\pgfqpoint{3.991354in}{0.600476in}}%
\pgfpathlineto{\pgfqpoint{3.993451in}{0.586903in}}%
\pgfpathlineto{\pgfqpoint{3.995549in}{0.552808in}}%
\pgfpathlineto{\pgfqpoint{3.997646in}{0.542936in}}%
\pgfpathlineto{\pgfqpoint{3.999743in}{0.589568in}}%
\pgfpathlineto{\pgfqpoint{4.001841in}{0.609244in}}%
\pgfpathlineto{\pgfqpoint{4.003938in}{0.620391in}}%
\pgfpathlineto{\pgfqpoint{4.006035in}{0.626230in}}%
\pgfpathlineto{\pgfqpoint{4.008133in}{0.627223in}}%
\pgfpathlineto{\pgfqpoint{4.010230in}{0.622101in}}%
\pgfpathlineto{\pgfqpoint{4.012327in}{0.610874in}}%
\pgfpathlineto{\pgfqpoint{4.014425in}{0.590717in}}%
\pgfpathlineto{\pgfqpoint{4.016522in}{0.539871in}}%
\pgfpathlineto{\pgfqpoint{4.018620in}{0.560303in}}%
\pgfpathlineto{\pgfqpoint{4.020717in}{0.592392in}}%
\pgfpathlineto{\pgfqpoint{4.022814in}{0.605973in}}%
\pgfpathlineto{\pgfqpoint{4.024912in}{0.612375in}}%
\pgfpathlineto{\pgfqpoint{4.027009in}{0.614103in}}%
\pgfpathlineto{\pgfqpoint{4.029106in}{0.611662in}}%
\pgfpathlineto{\pgfqpoint{4.031204in}{0.604282in}}%
\pgfpathlineto{\pgfqpoint{4.033301in}{0.588745in}}%
\pgfpathlineto{\pgfqpoint{4.035398in}{0.547950in}}%
\pgfpathlineto{\pgfqpoint{4.037496in}{0.556936in}}%
\pgfpathlineto{\pgfqpoint{4.039593in}{0.596512in}}%
\pgfpathlineto{\pgfqpoint{4.041691in}{0.614621in}}%
\pgfpathlineto{\pgfqpoint{4.043788in}{0.625002in}}%
\pgfpathlineto{\pgfqpoint{4.045885in}{0.630373in}}%
\pgfpathlineto{\pgfqpoint{4.047983in}{0.631070in}}%
\pgfpathlineto{\pgfqpoint{4.050080in}{0.625807in}}%
\pgfpathlineto{\pgfqpoint{4.052177in}{0.614284in}}%
\pgfpathlineto{\pgfqpoint{4.054275in}{0.593090in}}%
\pgfpathlineto{\pgfqpoint{4.056372in}{0.533454in}}%
\pgfpathlineto{\pgfqpoint{4.060567in}{0.599989in}}%
\pgfpathlineto{\pgfqpoint{4.062664in}{0.612946in}}%
\pgfpathlineto{\pgfqpoint{4.064762in}{0.619182in}}%
\pgfpathlineto{\pgfqpoint{4.066859in}{0.620956in}}%
\pgfpathlineto{\pgfqpoint{4.068956in}{0.618744in}}%
\pgfpathlineto{\pgfqpoint{4.071054in}{0.611879in}}%
\pgfpathlineto{\pgfqpoint{4.073151in}{0.597600in}}%
\pgfpathlineto{\pgfqpoint{4.075248in}{0.562977in}}%
\pgfpathlineto{\pgfqpoint{4.077346in}{0.552862in}}%
\pgfpathlineto{\pgfqpoint{4.079443in}{0.599215in}}%
\pgfpathlineto{\pgfqpoint{4.081540in}{0.618761in}}%
\pgfpathlineto{\pgfqpoint{4.083638in}{0.629907in}}%
\pgfpathlineto{\pgfqpoint{4.085735in}{0.635915in}}%
\pgfpathlineto{\pgfqpoint{4.087832in}{0.637371in}}%
\pgfpathlineto{\pgfqpoint{4.089930in}{0.633202in}}%
\pgfpathlineto{\pgfqpoint{4.092027in}{0.623171in}}%
\pgfpathlineto{\pgfqpoint{4.094125in}{0.605782in}}%
\pgfpathlineto{\pgfqpoint{4.096222in}{0.569380in}}%
\pgfpathlineto{\pgfqpoint{4.098319in}{0.547953in}}%
\pgfpathlineto{\pgfqpoint{4.100417in}{0.593805in}}%
\pgfpathlineto{\pgfqpoint{4.102514in}{0.610075in}}%
\pgfpathlineto{\pgfqpoint{4.104611in}{0.617609in}}%
\pgfpathlineto{\pgfqpoint{4.106709in}{0.619939in}}%
\pgfpathlineto{\pgfqpoint{4.108806in}{0.617843in}}%
\pgfpathlineto{\pgfqpoint{4.110903in}{0.610639in}}%
\pgfpathlineto{\pgfqpoint{4.113001in}{0.595073in}}%
\pgfpathlineto{\pgfqpoint{4.115098in}{0.553240in}}%
\pgfpathlineto{\pgfqpoint{4.117196in}{0.565315in}}%
\pgfpathlineto{\pgfqpoint{4.119293in}{0.604081in}}%
\pgfpathlineto{\pgfqpoint{4.121390in}{0.622007in}}%
\pgfpathlineto{\pgfqpoint{4.123488in}{0.632263in}}%
\pgfpathlineto{\pgfqpoint{4.125585in}{0.637473in}}%
\pgfpathlineto{\pgfqpoint{4.127682in}{0.637894in}}%
\pgfpathlineto{\pgfqpoint{4.129780in}{0.632262in}}%
\pgfpathlineto{\pgfqpoint{4.131877in}{0.621404in}}%
\pgfpathlineto{\pgfqpoint{4.133974in}{0.603476in}}%
\pgfpathlineto{\pgfqpoint{4.138169in}{0.540382in}}%
\pgfpathlineto{\pgfqpoint{4.140267in}{0.587747in}}%
\pgfpathlineto{\pgfqpoint{4.142364in}{0.603460in}}%
\pgfpathlineto{\pgfqpoint{4.144461in}{0.610014in}}%
\pgfpathlineto{\pgfqpoint{4.146559in}{0.610902in}}%
\pgfpathlineto{\pgfqpoint{4.148656in}{0.606540in}}%
\pgfpathlineto{\pgfqpoint{4.150753in}{0.595098in}}%
\pgfpathlineto{\pgfqpoint{4.152851in}{0.567674in}}%
\pgfpathlineto{\pgfqpoint{4.154948in}{0.526111in}}%
\pgfpathlineto{\pgfqpoint{4.157045in}{0.589570in}}%
\pgfpathlineto{\pgfqpoint{4.159143in}{0.611814in}}%
\pgfpathlineto{\pgfqpoint{4.161240in}{0.624192in}}%
\pgfpathlineto{\pgfqpoint{4.163338in}{0.630906in}}%
\pgfpathlineto{\pgfqpoint{4.165435in}{0.632743in}}%
\pgfpathlineto{\pgfqpoint{4.167532in}{0.628621in}}%
\pgfpathlineto{\pgfqpoint{4.169630in}{0.618940in}}%
\pgfpathlineto{\pgfqpoint{4.171727in}{0.603080in}}%
\pgfpathlineto{\pgfqpoint{4.173824in}{0.573715in}}%
\pgfpathlineto{\pgfqpoint{4.175922in}{0.492587in}}%
\pgfpathlineto{\pgfqpoint{4.178019in}{0.576036in}}%
\pgfpathlineto{\pgfqpoint{4.180116in}{0.595296in}}%
\pgfpathlineto{\pgfqpoint{4.182214in}{0.603469in}}%
\pgfpathlineto{\pgfqpoint{4.184311in}{0.605547in}}%
\pgfpathlineto{\pgfqpoint{4.186409in}{0.602437in}}%
\pgfpathlineto{\pgfqpoint{4.188506in}{0.592929in}}%
\pgfpathlineto{\pgfqpoint{4.190603in}{0.570951in}}%
\pgfpathlineto{\pgfqpoint{4.192701in}{0.409298in}}%
\pgfpathlineto{\pgfqpoint{4.194798in}{0.576375in}}%
\pgfpathlineto{\pgfqpoint{4.196895in}{0.602448in}}%
\pgfpathlineto{\pgfqpoint{4.198993in}{0.616449in}}%
\pgfpathlineto{\pgfqpoint{4.201090in}{0.624288in}}%
\pgfpathlineto{\pgfqpoint{4.203187in}{0.627236in}}%
\pgfpathlineto{\pgfqpoint{4.205285in}{0.624608in}}%
\pgfpathlineto{\pgfqpoint{4.207382in}{0.615690in}}%
\pgfpathlineto{\pgfqpoint{4.209480in}{0.600202in}}%
\pgfpathlineto{\pgfqpoint{4.211577in}{0.570136in}}%
\pgfpathlineto{\pgfqpoint{4.213674in}{0.509046in}}%
\pgfpathlineto{\pgfqpoint{4.215772in}{0.578584in}}%
\pgfpathlineto{\pgfqpoint{4.217869in}{0.597761in}}%
\pgfpathlineto{\pgfqpoint{4.219966in}{0.606637in}}%
\pgfpathlineto{\pgfqpoint{4.222064in}{0.609921in}}%
\pgfpathlineto{\pgfqpoint{4.224161in}{0.608751in}}%
\pgfpathlineto{\pgfqpoint{4.226258in}{0.602756in}}%
\pgfpathlineto{\pgfqpoint{4.228356in}{0.589493in}}%
\pgfpathlineto{\pgfqpoint{4.230453in}{0.557631in}}%
\pgfpathlineto{\pgfqpoint{4.232551in}{0.537711in}}%
\pgfpathlineto{\pgfqpoint{4.234648in}{0.588689in}}%
\pgfpathlineto{\pgfqpoint{4.236745in}{0.609095in}}%
\pgfpathlineto{\pgfqpoint{4.238843in}{0.620624in}}%
\pgfpathlineto{\pgfqpoint{4.240940in}{0.626821in}}%
\pgfpathlineto{\pgfqpoint{4.243037in}{0.628316in}}%
\pgfpathlineto{\pgfqpoint{4.245135in}{0.623982in}}%
\pgfpathlineto{\pgfqpoint{4.247232in}{0.613312in}}%
\pgfpathlineto{\pgfqpoint{4.249329in}{0.593749in}}%
\pgfpathlineto{\pgfqpoint{4.251427in}{0.544610in}}%
\pgfpathlineto{\pgfqpoint{4.253524in}{0.562441in}}%
\pgfpathlineto{\pgfqpoint{4.255621in}{0.595527in}}%
\pgfpathlineto{\pgfqpoint{4.257719in}{0.609647in}}%
\pgfpathlineto{\pgfqpoint{4.259816in}{0.616578in}}%
\pgfpathlineto{\pgfqpoint{4.261914in}{0.618936in}}%
\pgfpathlineto{\pgfqpoint{4.264011in}{0.617364in}}%
\pgfpathlineto{\pgfqpoint{4.266108in}{0.611390in}}%
\pgfpathlineto{\pgfqpoint{4.268206in}{0.598814in}}%
\pgfpathlineto{\pgfqpoint{4.270303in}{0.570490in}}%
\pgfpathlineto{\pgfqpoint{4.272400in}{0.529213in}}%
\pgfpathlineto{\pgfqpoint{4.274498in}{0.591766in}}%
\pgfpathlineto{\pgfqpoint{4.276595in}{0.613731in}}%
\pgfpathlineto{\pgfqpoint{4.278692in}{0.626030in}}%
\pgfpathlineto{\pgfqpoint{4.280790in}{0.632892in}}%
\pgfpathlineto{\pgfqpoint{4.282887in}{0.635236in}}%
\pgfpathlineto{\pgfqpoint{4.284985in}{0.632313in}}%
\pgfpathlineto{\pgfqpoint{4.289179in}{0.615408in}}%
\pgfpathlineto{\pgfqpoint{4.293374in}{0.589430in}}%
\pgfpathlineto{\pgfqpoint{4.297569in}{0.551389in}}%
\pgfpathlineto{\pgfqpoint{4.299666in}{0.531114in}}%
\pgfpathlineto{\pgfqpoint{4.301763in}{0.523148in}}%
\pgfpathlineto{\pgfqpoint{4.303861in}{0.534327in}}%
\pgfpathlineto{\pgfqpoint{4.308056in}{0.567592in}}%
\pgfpathlineto{\pgfqpoint{4.310153in}{0.579213in}}%
\pgfpathlineto{\pgfqpoint{4.312250in}{0.586774in}}%
\pgfpathlineto{\pgfqpoint{4.314348in}{0.589959in}}%
\pgfpathlineto{\pgfqpoint{4.316445in}{0.587339in}}%
\pgfpathlineto{\pgfqpoint{4.316445in}{0.587339in}}%
\pgfusepath{stroke}%
\end{pgfscope}%
\begin{pgfscope}%
\pgfpathrectangle{\pgfqpoint{2.554660in}{0.320625in}}{\pgfqpoint{1.845679in}{1.164375in}}%
\pgfusepath{clip}%
\pgfsetbuttcap%
\pgfsetroundjoin%
\pgfsetlinewidth{0.803000pt}%
\definecolor{currentstroke}{rgb}{0.000000,0.000000,0.000000}%
\pgfsetstrokecolor{currentstroke}%
\pgfsetstrokeopacity{0.700000}%
\pgfsetdash{{2.960000pt}{1.280000pt}}{0.000000pt}%
\pgfpathmoveto{\pgfqpoint{2.638555in}{1.468525in}}%
\pgfpathlineto{\pgfqpoint{3.284093in}{0.315625in}}%
\pgfpathlineto{\pgfqpoint{3.284093in}{0.315625in}}%
\pgfusepath{stroke}%
\end{pgfscope}%
\begin{pgfscope}%
\pgfsetrectcap%
\pgfsetmiterjoin%
\pgfsetlinewidth{0.803000pt}%
\definecolor{currentstroke}{rgb}{0.000000,0.000000,0.000000}%
\pgfsetstrokecolor{currentstroke}%
\pgfsetdash{}{0pt}%
\pgfpathmoveto{\pgfqpoint{2.554660in}{0.320625in}}%
\pgfpathlineto{\pgfqpoint{2.554660in}{1.485000in}}%
\pgfusepath{stroke}%
\end{pgfscope}%
\begin{pgfscope}%
\pgfsetrectcap%
\pgfsetmiterjoin%
\pgfsetlinewidth{0.803000pt}%
\definecolor{currentstroke}{rgb}{0.000000,0.000000,0.000000}%
\pgfsetstrokecolor{currentstroke}%
\pgfsetdash{}{0pt}%
\pgfpathmoveto{\pgfqpoint{4.400340in}{0.320625in}}%
\pgfpathlineto{\pgfqpoint{4.400340in}{1.485000in}}%
\pgfusepath{stroke}%
\end{pgfscope}%
\begin{pgfscope}%
\pgfsetrectcap%
\pgfsetmiterjoin%
\pgfsetlinewidth{0.803000pt}%
\definecolor{currentstroke}{rgb}{0.000000,0.000000,0.000000}%
\pgfsetstrokecolor{currentstroke}%
\pgfsetdash{}{0pt}%
\pgfpathmoveto{\pgfqpoint{2.554660in}{0.320625in}}%
\pgfpathlineto{\pgfqpoint{4.400340in}{0.320625in}}%
\pgfusepath{stroke}%
\end{pgfscope}%
\begin{pgfscope}%
\pgfsetrectcap%
\pgfsetmiterjoin%
\pgfsetlinewidth{0.803000pt}%
\definecolor{currentstroke}{rgb}{0.000000,0.000000,0.000000}%
\pgfsetstrokecolor{currentstroke}%
\pgfsetdash{}{0pt}%
\pgfpathmoveto{\pgfqpoint{2.554660in}{1.485000in}}%
\pgfpathlineto{\pgfqpoint{4.400340in}{1.485000in}}%
\pgfusepath{stroke}%
\end{pgfscope}%
\begin{pgfscope}%
\definecolor{textcolor}{rgb}{0.000000,0.000000,0.000000}%
\pgfsetstrokecolor{textcolor}%
\pgfsetfillcolor{textcolor}%
\pgftext[x=3.477500in,y=1.568333in,,base]{\color{textcolor}{\rmfamily\fontsize{8.000000}{9.600000}\selectfont\catcode`\^=\active\def^{\ifmmode\sp\else\^{}\fi}\catcode`\%=\active\def%{\%}(b) Covariance conditioning}}%
\end{pgfscope}%
\begin{pgfscope}%
\pgfsetbuttcap%
\pgfsetmiterjoin%
\definecolor{currentfill}{rgb}{1.000000,1.000000,1.000000}%
\pgfsetfillcolor{currentfill}%
\pgfsetfillopacity{0.800000}%
\pgfsetlinewidth{1.003750pt}%
\definecolor{currentstroke}{rgb}{0.800000,0.800000,0.800000}%
\pgfsetstrokecolor{currentstroke}%
\pgfsetstrokeopacity{0.800000}%
\pgfsetdash{}{0pt}%
\pgfpathmoveto{\pgfqpoint{3.905118in}{1.210066in}}%
\pgfpathlineto{\pgfqpoint{4.355617in}{1.210066in}}%
\pgfpathquadraticcurveto{\pgfqpoint{4.368395in}{1.210066in}}{\pgfqpoint{4.368395in}{1.222844in}}%
\pgfpathlineto{\pgfqpoint{4.368395in}{1.440278in}}%
\pgfpathquadraticcurveto{\pgfqpoint{4.368395in}{1.453056in}}{\pgfqpoint{4.355617in}{1.453056in}}%
\pgfpathlineto{\pgfqpoint{3.905118in}{1.453056in}}%
\pgfpathquadraticcurveto{\pgfqpoint{3.892340in}{1.453056in}}{\pgfqpoint{3.892340in}{1.440278in}}%
\pgfpathlineto{\pgfqpoint{3.892340in}{1.222844in}}%
\pgfpathquadraticcurveto{\pgfqpoint{3.892340in}{1.210066in}}{\pgfqpoint{3.905118in}{1.210066in}}%
\pgfpathlineto{\pgfqpoint{3.905118in}{1.210066in}}%
\pgfpathclose%
\pgfusepath{stroke,fill}%
\end{pgfscope}%
\begin{pgfscope}%
\pgfsetrectcap%
\pgfsetroundjoin%
\pgfsetlinewidth{1.154312pt}%
\definecolor{currentstroke}{rgb}{0.146785,0.068738,0.334011}%
\pgfsetstrokecolor{currentstroke}%
\pgfsetdash{}{0pt}%
\pgfpathmoveto{\pgfqpoint{3.908312in}{1.410904in}}%
\pgfpathlineto{\pgfqpoint{3.953034in}{1.410904in}}%
\pgfpathlineto{\pgfqpoint{3.997757in}{1.410904in}}%
\pgfusepath{stroke}%
\end{pgfscope}%
\begin{pgfscope}%
\definecolor{textcolor}{rgb}{0.000000,0.000000,0.000000}%
\pgfsetstrokecolor{textcolor}%
\pgfsetfillcolor{textcolor}%
\pgftext[x=4.048868in,y=1.388543in,left,base]{\color{textcolor}{\rmfamily\fontsize{4.600000}{5.520000}\selectfont\catcode`\^=\active\def^{\ifmmode\sp\else\^{}\fi}\catcode`\%=\active\def%{\%}$r=0$}}%
\end{pgfscope}%
\begin{pgfscope}%
\pgfsetrectcap%
\pgfsetroundjoin%
\pgfsetlinewidth{1.154312pt}%
\definecolor{currentstroke}{rgb}{0.716387,0.214982,0.475290}%
\pgfsetstrokecolor{currentstroke}%
\pgfsetdash{}{0pt}%
\pgfpathmoveto{\pgfqpoint{3.908312in}{1.336296in}}%
\pgfpathlineto{\pgfqpoint{3.953034in}{1.336296in}}%
\pgfpathlineto{\pgfqpoint{3.997757in}{1.336296in}}%
\pgfusepath{stroke}%
\end{pgfscope}%
\begin{pgfscope}%
\definecolor{textcolor}{rgb}{0.000000,0.000000,0.000000}%
\pgfsetstrokecolor{textcolor}%
\pgfsetfillcolor{textcolor}%
\pgftext[x=4.048868in,y=1.313935in,left,base]{\color{textcolor}{\rmfamily\fontsize{4.600000}{5.520000}\selectfont\catcode`\^=\active\def^{\ifmmode\sp\else\^{}\fi}\catcode`\%=\active\def%{\%}$r=0.5$}}%
\end{pgfscope}%
\begin{pgfscope}%
\pgfsetrectcap%
\pgfsetroundjoin%
\pgfsetlinewidth{1.154312pt}%
\definecolor{currentstroke}{rgb}{0.997351,0.719089,0.493755}%
\pgfsetstrokecolor{currentstroke}%
\pgfsetdash{}{0pt}%
\pgfpathmoveto{\pgfqpoint{3.908312in}{1.261688in}}%
\pgfpathlineto{\pgfqpoint{3.953034in}{1.261688in}}%
\pgfpathlineto{\pgfqpoint{3.997757in}{1.261688in}}%
\pgfusepath{stroke}%
\end{pgfscope}%
\begin{pgfscope}%
\definecolor{textcolor}{rgb}{0.000000,0.000000,0.000000}%
\pgfsetstrokecolor{textcolor}%
\pgfsetfillcolor{textcolor}%
\pgftext[x=4.048868in,y=1.239327in,left,base]{\color{textcolor}{\rmfamily\fontsize{4.600000}{5.520000}\selectfont\catcode`\^=\active\def^{\ifmmode\sp\else\^{}\fi}\catcode`\%=\active\def%{\%}$r=1$}}%
\end{pgfscope}%
\begin{pgfscope}%
\pgfsetbuttcap%
\pgfsetmiterjoin%
\definecolor{currentfill}{rgb}{1.000000,1.000000,1.000000}%
\pgfsetfillcolor{currentfill}%
\pgfsetlinewidth{0.000000pt}%
\definecolor{currentstroke}{rgb}{0.000000,0.000000,0.000000}%
\pgfsetstrokecolor{currentstroke}%
\pgfsetstrokeopacity{0.000000}%
\pgfsetdash{}{0pt}%
\pgfpathmoveto{\pgfqpoint{4.621821in}{0.320625in}}%
\pgfpathlineto{\pgfqpoint{6.467500in}{0.320625in}}%
\pgfpathlineto{\pgfqpoint{6.467500in}{1.485000in}}%
\pgfpathlineto{\pgfqpoint{4.621821in}{1.485000in}}%
\pgfpathlineto{\pgfqpoint{4.621821in}{0.320625in}}%
\pgfpathclose%
\pgfusepath{fill}%
\end{pgfscope}%
\begin{pgfscope}%
\pgfpathrectangle{\pgfqpoint{4.621821in}{0.320625in}}{\pgfqpoint{1.845679in}{1.164375in}}%
\pgfusepath{clip}%
\pgfsetrectcap%
\pgfsetroundjoin%
\pgfsetlinewidth{0.803000pt}%
\definecolor{currentstroke}{rgb}{0.909804,0.909804,0.909804}%
\pgfsetstrokecolor{currentstroke}%
\pgfsetdash{}{0pt}%
\pgfpathmoveto{\pgfqpoint{4.705715in}{0.320625in}}%
\pgfpathlineto{\pgfqpoint{4.705715in}{1.485000in}}%
\pgfusepath{stroke}%
\end{pgfscope}%
\begin{pgfscope}%
\pgfsetbuttcap%
\pgfsetroundjoin%
\definecolor{currentfill}{rgb}{0.000000,0.000000,0.000000}%
\pgfsetfillcolor{currentfill}%
\pgfsetlinewidth{0.803000pt}%
\definecolor{currentstroke}{rgb}{0.000000,0.000000,0.000000}%
\pgfsetstrokecolor{currentstroke}%
\pgfsetdash{}{0pt}%
\pgfsys@defobject{currentmarker}{\pgfqpoint{0.000000in}{-0.048611in}}{\pgfqpoint{0.000000in}{0.000000in}}{%
\pgfpathmoveto{\pgfqpoint{0.000000in}{0.000000in}}%
\pgfpathlineto{\pgfqpoint{0.000000in}{-0.048611in}}%
\pgfusepath{stroke,fill}%
}%
\begin{pgfscope}%
\pgfsys@transformshift{4.705715in}{0.320625in}%
\pgfsys@useobject{currentmarker}{}%
\end{pgfscope}%
\end{pgfscope}%
\begin{pgfscope}%
\definecolor{textcolor}{rgb}{0.000000,0.000000,0.000000}%
\pgfsetstrokecolor{textcolor}%
\pgfsetfillcolor{textcolor}%
\pgftext[x=4.705715in,y=0.223403in,,top]{\color{textcolor}{\rmfamily\fontsize{6.000000}{7.200000}\selectfont\catcode`\^=\active\def^{\ifmmode\sp\else\^{}\fi}\catcode`\%=\active\def%{\%}0}}%
\end{pgfscope}%
\begin{pgfscope}%
\pgfpathrectangle{\pgfqpoint{4.621821in}{0.320625in}}{\pgfqpoint{1.845679in}{1.164375in}}%
\pgfusepath{clip}%
\pgfsetrectcap%
\pgfsetroundjoin%
\pgfsetlinewidth{0.803000pt}%
\definecolor{currentstroke}{rgb}{0.909804,0.909804,0.909804}%
\pgfsetstrokecolor{currentstroke}%
\pgfsetdash{}{0pt}%
\pgfpathmoveto{\pgfqpoint{5.125188in}{0.320625in}}%
\pgfpathlineto{\pgfqpoint{5.125188in}{1.485000in}}%
\pgfusepath{stroke}%
\end{pgfscope}%
\begin{pgfscope}%
\pgfsetbuttcap%
\pgfsetroundjoin%
\definecolor{currentfill}{rgb}{0.000000,0.000000,0.000000}%
\pgfsetfillcolor{currentfill}%
\pgfsetlinewidth{0.803000pt}%
\definecolor{currentstroke}{rgb}{0.000000,0.000000,0.000000}%
\pgfsetstrokecolor{currentstroke}%
\pgfsetdash{}{0pt}%
\pgfsys@defobject{currentmarker}{\pgfqpoint{0.000000in}{-0.048611in}}{\pgfqpoint{0.000000in}{0.000000in}}{%
\pgfpathmoveto{\pgfqpoint{0.000000in}{0.000000in}}%
\pgfpathlineto{\pgfqpoint{0.000000in}{-0.048611in}}%
\pgfusepath{stroke,fill}%
}%
\begin{pgfscope}%
\pgfsys@transformshift{5.125188in}{0.320625in}%
\pgfsys@useobject{currentmarker}{}%
\end{pgfscope}%
\end{pgfscope}%
\begin{pgfscope}%
\definecolor{textcolor}{rgb}{0.000000,0.000000,0.000000}%
\pgfsetstrokecolor{textcolor}%
\pgfsetfillcolor{textcolor}%
\pgftext[x=5.125188in,y=0.223403in,,top]{\color{textcolor}{\rmfamily\fontsize{6.000000}{7.200000}\selectfont\catcode`\^=\active\def^{\ifmmode\sp\else\^{}\fi}\catcode`\%=\active\def%{\%}10}}%
\end{pgfscope}%
\begin{pgfscope}%
\pgfpathrectangle{\pgfqpoint{4.621821in}{0.320625in}}{\pgfqpoint{1.845679in}{1.164375in}}%
\pgfusepath{clip}%
\pgfsetrectcap%
\pgfsetroundjoin%
\pgfsetlinewidth{0.803000pt}%
\definecolor{currentstroke}{rgb}{0.909804,0.909804,0.909804}%
\pgfsetstrokecolor{currentstroke}%
\pgfsetdash{}{0pt}%
\pgfpathmoveto{\pgfqpoint{5.544660in}{0.320625in}}%
\pgfpathlineto{\pgfqpoint{5.544660in}{1.485000in}}%
\pgfusepath{stroke}%
\end{pgfscope}%
\begin{pgfscope}%
\pgfsetbuttcap%
\pgfsetroundjoin%
\definecolor{currentfill}{rgb}{0.000000,0.000000,0.000000}%
\pgfsetfillcolor{currentfill}%
\pgfsetlinewidth{0.803000pt}%
\definecolor{currentstroke}{rgb}{0.000000,0.000000,0.000000}%
\pgfsetstrokecolor{currentstroke}%
\pgfsetdash{}{0pt}%
\pgfsys@defobject{currentmarker}{\pgfqpoint{0.000000in}{-0.048611in}}{\pgfqpoint{0.000000in}{0.000000in}}{%
\pgfpathmoveto{\pgfqpoint{0.000000in}{0.000000in}}%
\pgfpathlineto{\pgfqpoint{0.000000in}{-0.048611in}}%
\pgfusepath{stroke,fill}%
}%
\begin{pgfscope}%
\pgfsys@transformshift{5.544660in}{0.320625in}%
\pgfsys@useobject{currentmarker}{}%
\end{pgfscope}%
\end{pgfscope}%
\begin{pgfscope}%
\definecolor{textcolor}{rgb}{0.000000,0.000000,0.000000}%
\pgfsetstrokecolor{textcolor}%
\pgfsetfillcolor{textcolor}%
\pgftext[x=5.544660in,y=0.223403in,,top]{\color{textcolor}{\rmfamily\fontsize{6.000000}{7.200000}\selectfont\catcode`\^=\active\def^{\ifmmode\sp\else\^{}\fi}\catcode`\%=\active\def%{\%}20}}%
\end{pgfscope}%
\begin{pgfscope}%
\pgfpathrectangle{\pgfqpoint{4.621821in}{0.320625in}}{\pgfqpoint{1.845679in}{1.164375in}}%
\pgfusepath{clip}%
\pgfsetrectcap%
\pgfsetroundjoin%
\pgfsetlinewidth{0.803000pt}%
\definecolor{currentstroke}{rgb}{0.909804,0.909804,0.909804}%
\pgfsetstrokecolor{currentstroke}%
\pgfsetdash{}{0pt}%
\pgfpathmoveto{\pgfqpoint{5.964133in}{0.320625in}}%
\pgfpathlineto{\pgfqpoint{5.964133in}{1.485000in}}%
\pgfusepath{stroke}%
\end{pgfscope}%
\begin{pgfscope}%
\pgfsetbuttcap%
\pgfsetroundjoin%
\definecolor{currentfill}{rgb}{0.000000,0.000000,0.000000}%
\pgfsetfillcolor{currentfill}%
\pgfsetlinewidth{0.803000pt}%
\definecolor{currentstroke}{rgb}{0.000000,0.000000,0.000000}%
\pgfsetstrokecolor{currentstroke}%
\pgfsetdash{}{0pt}%
\pgfsys@defobject{currentmarker}{\pgfqpoint{0.000000in}{-0.048611in}}{\pgfqpoint{0.000000in}{0.000000in}}{%
\pgfpathmoveto{\pgfqpoint{0.000000in}{0.000000in}}%
\pgfpathlineto{\pgfqpoint{0.000000in}{-0.048611in}}%
\pgfusepath{stroke,fill}%
}%
\begin{pgfscope}%
\pgfsys@transformshift{5.964133in}{0.320625in}%
\pgfsys@useobject{currentmarker}{}%
\end{pgfscope}%
\end{pgfscope}%
\begin{pgfscope}%
\definecolor{textcolor}{rgb}{0.000000,0.000000,0.000000}%
\pgfsetstrokecolor{textcolor}%
\pgfsetfillcolor{textcolor}%
\pgftext[x=5.964133in,y=0.223403in,,top]{\color{textcolor}{\rmfamily\fontsize{6.000000}{7.200000}\selectfont\catcode`\^=\active\def^{\ifmmode\sp\else\^{}\fi}\catcode`\%=\active\def%{\%}30}}%
\end{pgfscope}%
\begin{pgfscope}%
\pgfpathrectangle{\pgfqpoint{4.621821in}{0.320625in}}{\pgfqpoint{1.845679in}{1.164375in}}%
\pgfusepath{clip}%
\pgfsetrectcap%
\pgfsetroundjoin%
\pgfsetlinewidth{0.803000pt}%
\definecolor{currentstroke}{rgb}{0.909804,0.909804,0.909804}%
\pgfsetstrokecolor{currentstroke}%
\pgfsetdash{}{0pt}%
\pgfpathmoveto{\pgfqpoint{6.383605in}{0.320625in}}%
\pgfpathlineto{\pgfqpoint{6.383605in}{1.485000in}}%
\pgfusepath{stroke}%
\end{pgfscope}%
\begin{pgfscope}%
\pgfsetbuttcap%
\pgfsetroundjoin%
\definecolor{currentfill}{rgb}{0.000000,0.000000,0.000000}%
\pgfsetfillcolor{currentfill}%
\pgfsetlinewidth{0.803000pt}%
\definecolor{currentstroke}{rgb}{0.000000,0.000000,0.000000}%
\pgfsetstrokecolor{currentstroke}%
\pgfsetdash{}{0pt}%
\pgfsys@defobject{currentmarker}{\pgfqpoint{0.000000in}{-0.048611in}}{\pgfqpoint{0.000000in}{0.000000in}}{%
\pgfpathmoveto{\pgfqpoint{0.000000in}{0.000000in}}%
\pgfpathlineto{\pgfqpoint{0.000000in}{-0.048611in}}%
\pgfusepath{stroke,fill}%
}%
\begin{pgfscope}%
\pgfsys@transformshift{6.383605in}{0.320625in}%
\pgfsys@useobject{currentmarker}{}%
\end{pgfscope}%
\end{pgfscope}%
\begin{pgfscope}%
\definecolor{textcolor}{rgb}{0.000000,0.000000,0.000000}%
\pgfsetstrokecolor{textcolor}%
\pgfsetfillcolor{textcolor}%
\pgftext[x=6.383605in,y=0.223403in,,top]{\color{textcolor}{\rmfamily\fontsize{6.000000}{7.200000}\selectfont\catcode`\^=\active\def^{\ifmmode\sp\else\^{}\fi}\catcode`\%=\active\def%{\%}40}}%
\end{pgfscope}%
\begin{pgfscope}%
\definecolor{textcolor}{rgb}{0.000000,0.000000,0.000000}%
\pgfsetstrokecolor{textcolor}%
\pgfsetfillcolor{textcolor}%
\pgftext[x=5.544660in,y=0.087200in,,top]{\color{textcolor}{\rmfamily\fontsize{6.000000}{7.200000}\selectfont\catcode`\^=\active\def^{\ifmmode\sp\else\^{}\fi}\catcode`\%=\active\def%{\%}time}}%
\end{pgfscope}%
\begin{pgfscope}%
\pgfpathrectangle{\pgfqpoint{4.621821in}{0.320625in}}{\pgfqpoint{1.845679in}{1.164375in}}%
\pgfusepath{clip}%
\pgfsetrectcap%
\pgfsetroundjoin%
\pgfsetlinewidth{0.803000pt}%
\definecolor{currentstroke}{rgb}{0.909804,0.909804,0.909804}%
\pgfsetstrokecolor{currentstroke}%
\pgfsetdash{}{0pt}%
\pgfpathmoveto{\pgfqpoint{4.621821in}{0.320625in}}%
\pgfpathlineto{\pgfqpoint{6.467500in}{0.320625in}}%
\pgfusepath{stroke}%
\end{pgfscope}%
\begin{pgfscope}%
\pgfsetbuttcap%
\pgfsetroundjoin%
\definecolor{currentfill}{rgb}{0.000000,0.000000,0.000000}%
\pgfsetfillcolor{currentfill}%
\pgfsetlinewidth{0.803000pt}%
\definecolor{currentstroke}{rgb}{0.000000,0.000000,0.000000}%
\pgfsetstrokecolor{currentstroke}%
\pgfsetdash{}{0pt}%
\pgfsys@defobject{currentmarker}{\pgfqpoint{-0.048611in}{0.000000in}}{\pgfqpoint{-0.000000in}{0.000000in}}{%
\pgfpathmoveto{\pgfqpoint{-0.000000in}{0.000000in}}%
\pgfpathlineto{\pgfqpoint{-0.048611in}{0.000000in}}%
\pgfusepath{stroke,fill}%
}%
\begin{pgfscope}%
\pgfsys@transformshift{4.621821in}{0.320625in}%
\pgfsys@useobject{currentmarker}{}%
\end{pgfscope}%
\end{pgfscope}%
\begin{pgfscope}%
\pgfpathrectangle{\pgfqpoint{4.621821in}{0.320625in}}{\pgfqpoint{1.845679in}{1.164375in}}%
\pgfusepath{clip}%
\pgfsetrectcap%
\pgfsetroundjoin%
\pgfsetlinewidth{0.803000pt}%
\definecolor{currentstroke}{rgb}{0.909804,0.909804,0.909804}%
\pgfsetstrokecolor{currentstroke}%
\pgfsetdash{}{0pt}%
\pgfpathmoveto{\pgfqpoint{4.621821in}{0.536250in}}%
\pgfpathlineto{\pgfqpoint{6.467500in}{0.536250in}}%
\pgfusepath{stroke}%
\end{pgfscope}%
\begin{pgfscope}%
\pgfsetbuttcap%
\pgfsetroundjoin%
\definecolor{currentfill}{rgb}{0.000000,0.000000,0.000000}%
\pgfsetfillcolor{currentfill}%
\pgfsetlinewidth{0.803000pt}%
\definecolor{currentstroke}{rgb}{0.000000,0.000000,0.000000}%
\pgfsetstrokecolor{currentstroke}%
\pgfsetdash{}{0pt}%
\pgfsys@defobject{currentmarker}{\pgfqpoint{-0.048611in}{0.000000in}}{\pgfqpoint{-0.000000in}{0.000000in}}{%
\pgfpathmoveto{\pgfqpoint{-0.000000in}{0.000000in}}%
\pgfpathlineto{\pgfqpoint{-0.048611in}{0.000000in}}%
\pgfusepath{stroke,fill}%
}%
\begin{pgfscope}%
\pgfsys@transformshift{4.621821in}{0.536250in}%
\pgfsys@useobject{currentmarker}{}%
\end{pgfscope}%
\end{pgfscope}%
\begin{pgfscope}%
\pgfpathrectangle{\pgfqpoint{4.621821in}{0.320625in}}{\pgfqpoint{1.845679in}{1.164375in}}%
\pgfusepath{clip}%
\pgfsetrectcap%
\pgfsetroundjoin%
\pgfsetlinewidth{0.803000pt}%
\definecolor{currentstroke}{rgb}{0.909804,0.909804,0.909804}%
\pgfsetstrokecolor{currentstroke}%
\pgfsetdash{}{0pt}%
\pgfpathmoveto{\pgfqpoint{4.621821in}{0.751875in}}%
\pgfpathlineto{\pgfqpoint{6.467500in}{0.751875in}}%
\pgfusepath{stroke}%
\end{pgfscope}%
\begin{pgfscope}%
\pgfsetbuttcap%
\pgfsetroundjoin%
\definecolor{currentfill}{rgb}{0.000000,0.000000,0.000000}%
\pgfsetfillcolor{currentfill}%
\pgfsetlinewidth{0.803000pt}%
\definecolor{currentstroke}{rgb}{0.000000,0.000000,0.000000}%
\pgfsetstrokecolor{currentstroke}%
\pgfsetdash{}{0pt}%
\pgfsys@defobject{currentmarker}{\pgfqpoint{-0.048611in}{0.000000in}}{\pgfqpoint{-0.000000in}{0.000000in}}{%
\pgfpathmoveto{\pgfqpoint{-0.000000in}{0.000000in}}%
\pgfpathlineto{\pgfqpoint{-0.048611in}{0.000000in}}%
\pgfusepath{stroke,fill}%
}%
\begin{pgfscope}%
\pgfsys@transformshift{4.621821in}{0.751875in}%
\pgfsys@useobject{currentmarker}{}%
\end{pgfscope}%
\end{pgfscope}%
\begin{pgfscope}%
\pgfpathrectangle{\pgfqpoint{4.621821in}{0.320625in}}{\pgfqpoint{1.845679in}{1.164375in}}%
\pgfusepath{clip}%
\pgfsetrectcap%
\pgfsetroundjoin%
\pgfsetlinewidth{0.803000pt}%
\definecolor{currentstroke}{rgb}{0.909804,0.909804,0.909804}%
\pgfsetstrokecolor{currentstroke}%
\pgfsetdash{}{0pt}%
\pgfpathmoveto{\pgfqpoint{4.621821in}{0.967500in}}%
\pgfpathlineto{\pgfqpoint{6.467500in}{0.967500in}}%
\pgfusepath{stroke}%
\end{pgfscope}%
\begin{pgfscope}%
\pgfsetbuttcap%
\pgfsetroundjoin%
\definecolor{currentfill}{rgb}{0.000000,0.000000,0.000000}%
\pgfsetfillcolor{currentfill}%
\pgfsetlinewidth{0.803000pt}%
\definecolor{currentstroke}{rgb}{0.000000,0.000000,0.000000}%
\pgfsetstrokecolor{currentstroke}%
\pgfsetdash{}{0pt}%
\pgfsys@defobject{currentmarker}{\pgfqpoint{-0.048611in}{0.000000in}}{\pgfqpoint{-0.000000in}{0.000000in}}{%
\pgfpathmoveto{\pgfqpoint{-0.000000in}{0.000000in}}%
\pgfpathlineto{\pgfqpoint{-0.048611in}{0.000000in}}%
\pgfusepath{stroke,fill}%
}%
\begin{pgfscope}%
\pgfsys@transformshift{4.621821in}{0.967500in}%
\pgfsys@useobject{currentmarker}{}%
\end{pgfscope}%
\end{pgfscope}%
\begin{pgfscope}%
\pgfpathrectangle{\pgfqpoint{4.621821in}{0.320625in}}{\pgfqpoint{1.845679in}{1.164375in}}%
\pgfusepath{clip}%
\pgfsetrectcap%
\pgfsetroundjoin%
\pgfsetlinewidth{0.803000pt}%
\definecolor{currentstroke}{rgb}{0.909804,0.909804,0.909804}%
\pgfsetstrokecolor{currentstroke}%
\pgfsetdash{}{0pt}%
\pgfpathmoveto{\pgfqpoint{4.621821in}{1.183125in}}%
\pgfpathlineto{\pgfqpoint{6.467500in}{1.183125in}}%
\pgfusepath{stroke}%
\end{pgfscope}%
\begin{pgfscope}%
\pgfsetbuttcap%
\pgfsetroundjoin%
\definecolor{currentfill}{rgb}{0.000000,0.000000,0.000000}%
\pgfsetfillcolor{currentfill}%
\pgfsetlinewidth{0.803000pt}%
\definecolor{currentstroke}{rgb}{0.000000,0.000000,0.000000}%
\pgfsetstrokecolor{currentstroke}%
\pgfsetdash{}{0pt}%
\pgfsys@defobject{currentmarker}{\pgfqpoint{-0.048611in}{0.000000in}}{\pgfqpoint{-0.000000in}{0.000000in}}{%
\pgfpathmoveto{\pgfqpoint{-0.000000in}{0.000000in}}%
\pgfpathlineto{\pgfqpoint{-0.048611in}{0.000000in}}%
\pgfusepath{stroke,fill}%
}%
\begin{pgfscope}%
\pgfsys@transformshift{4.621821in}{1.183125in}%
\pgfsys@useobject{currentmarker}{}%
\end{pgfscope}%
\end{pgfscope}%
\begin{pgfscope}%
\pgfpathrectangle{\pgfqpoint{4.621821in}{0.320625in}}{\pgfqpoint{1.845679in}{1.164375in}}%
\pgfusepath{clip}%
\pgfsetrectcap%
\pgfsetroundjoin%
\pgfsetlinewidth{0.803000pt}%
\definecolor{currentstroke}{rgb}{0.909804,0.909804,0.909804}%
\pgfsetstrokecolor{currentstroke}%
\pgfsetdash{}{0pt}%
\pgfpathmoveto{\pgfqpoint{4.621821in}{1.398750in}}%
\pgfpathlineto{\pgfqpoint{6.467500in}{1.398750in}}%
\pgfusepath{stroke}%
\end{pgfscope}%
\begin{pgfscope}%
\pgfsetbuttcap%
\pgfsetroundjoin%
\definecolor{currentfill}{rgb}{0.000000,0.000000,0.000000}%
\pgfsetfillcolor{currentfill}%
\pgfsetlinewidth{0.803000pt}%
\definecolor{currentstroke}{rgb}{0.000000,0.000000,0.000000}%
\pgfsetstrokecolor{currentstroke}%
\pgfsetdash{}{0pt}%
\pgfsys@defobject{currentmarker}{\pgfqpoint{-0.048611in}{0.000000in}}{\pgfqpoint{-0.000000in}{0.000000in}}{%
\pgfpathmoveto{\pgfqpoint{-0.000000in}{0.000000in}}%
\pgfpathlineto{\pgfqpoint{-0.048611in}{0.000000in}}%
\pgfusepath{stroke,fill}%
}%
\begin{pgfscope}%
\pgfsys@transformshift{4.621821in}{1.398750in}%
\pgfsys@useobject{currentmarker}{}%
\end{pgfscope}%
\end{pgfscope}%
\begin{pgfscope}%
\pgfpathrectangle{\pgfqpoint{4.621821in}{0.320625in}}{\pgfqpoint{1.845679in}{1.164375in}}%
\pgfusepath{clip}%
\pgfsetrectcap%
\pgfsetroundjoin%
\pgfsetlinewidth{1.304875pt}%
\definecolor{currentstroke}{rgb}{0.000000,0.000000,0.000000}%
\pgfsetstrokecolor{currentstroke}%
\pgfsetdash{}{0pt}%
\pgfpathmoveto{\pgfqpoint{4.705715in}{1.468525in}}%
\pgfpathlineto{\pgfqpoint{4.715783in}{1.449849in}}%
\pgfpathlineto{\pgfqpoint{4.725850in}{1.428582in}}%
\pgfpathlineto{\pgfqpoint{4.739273in}{1.396826in}}%
\pgfpathlineto{\pgfqpoint{4.779543in}{1.299011in}}%
\pgfpathlineto{\pgfqpoint{4.803033in}{1.245684in}}%
\pgfpathlineto{\pgfqpoint{4.829879in}{1.187742in}}%
\pgfpathlineto{\pgfqpoint{4.863437in}{1.118228in}}%
\pgfpathlineto{\pgfqpoint{4.907062in}{1.030672in}}%
\pgfpathlineto{\pgfqpoint{4.967466in}{0.912223in}}%
\pgfpathlineto{\pgfqpoint{5.044649in}{0.763330in}}%
\pgfpathlineto{\pgfqpoint{5.058072in}{0.737306in}}%
\pgfpathlineto{\pgfqpoint{5.068140in}{0.717867in}}%
\pgfpathlineto{\pgfqpoint{5.084919in}{0.687078in}}%
\pgfpathlineto{\pgfqpoint{5.101698in}{0.652063in}}%
\pgfpathlineto{\pgfqpoint{5.105053in}{0.647076in}}%
\pgfpathlineto{\pgfqpoint{5.111765in}{0.640040in}}%
\pgfpathlineto{\pgfqpoint{5.115121in}{0.636691in}}%
\pgfpathlineto{\pgfqpoint{5.118476in}{0.631615in}}%
\pgfpathlineto{\pgfqpoint{5.121832in}{0.624268in}}%
\pgfpathlineto{\pgfqpoint{5.125188in}{0.615244in}}%
\pgfpathlineto{\pgfqpoint{5.128544in}{0.604400in}}%
\pgfpathlineto{\pgfqpoint{5.135255in}{0.577200in}}%
\pgfpathlineto{\pgfqpoint{5.141967in}{0.549758in}}%
\pgfpathlineto{\pgfqpoint{5.145323in}{0.543089in}}%
\pgfpathlineto{\pgfqpoint{5.148678in}{0.545206in}}%
\pgfpathlineto{\pgfqpoint{5.152034in}{0.555371in}}%
\pgfpathlineto{\pgfqpoint{5.158746in}{0.577878in}}%
\pgfpathlineto{\pgfqpoint{5.162102in}{0.584907in}}%
\pgfpathlineto{\pgfqpoint{5.165457in}{0.587993in}}%
\pgfpathlineto{\pgfqpoint{5.168813in}{0.585714in}}%
\pgfpathlineto{\pgfqpoint{5.172169in}{0.575987in}}%
\pgfpathlineto{\pgfqpoint{5.175525in}{0.557412in}}%
\pgfpathlineto{\pgfqpoint{5.178880in}{0.543212in}}%
\pgfpathlineto{\pgfqpoint{5.182236in}{0.565484in}}%
\pgfpathlineto{\pgfqpoint{5.185592in}{0.583410in}}%
\pgfpathlineto{\pgfqpoint{5.188948in}{0.594334in}}%
\pgfpathlineto{\pgfqpoint{5.192304in}{0.600781in}}%
\pgfpathlineto{\pgfqpoint{5.195659in}{0.603926in}}%
\pgfpathlineto{\pgfqpoint{5.199015in}{0.604201in}}%
\pgfpathlineto{\pgfqpoint{5.202371in}{0.601539in}}%
\pgfpathlineto{\pgfqpoint{5.205727in}{0.595294in}}%
\pgfpathlineto{\pgfqpoint{5.209082in}{0.583604in}}%
\pgfpathlineto{\pgfqpoint{5.212438in}{0.560150in}}%
\pgfpathlineto{\pgfqpoint{5.215794in}{0.482906in}}%
\pgfpathlineto{\pgfqpoint{5.219150in}{0.564283in}}%
\pgfpathlineto{\pgfqpoint{5.222506in}{0.589670in}}%
\pgfpathlineto{\pgfqpoint{5.225861in}{0.603589in}}%
\pgfpathlineto{\pgfqpoint{5.229217in}{0.611859in}}%
\pgfpathlineto{\pgfqpoint{5.232573in}{0.615991in}}%
\pgfpathlineto{\pgfqpoint{5.235929in}{0.616022in}}%
\pgfpathlineto{\pgfqpoint{5.239285in}{0.610546in}}%
\pgfpathlineto{\pgfqpoint{5.242640in}{0.598461in}}%
\pgfpathlineto{\pgfqpoint{5.245996in}{0.596439in}}%
\pgfpathlineto{\pgfqpoint{5.249352in}{0.613372in}}%
\pgfpathlineto{\pgfqpoint{5.252708in}{0.627432in}}%
\pgfpathlineto{\pgfqpoint{5.256063in}{0.635926in}}%
\pgfpathlineto{\pgfqpoint{5.259419in}{0.639936in}}%
\pgfpathlineto{\pgfqpoint{5.262775in}{0.639952in}}%
\pgfpathlineto{\pgfqpoint{5.266131in}{0.635649in}}%
\pgfpathlineto{\pgfqpoint{5.269487in}{0.625356in}}%
\pgfpathlineto{\pgfqpoint{5.272842in}{0.602859in}}%
\pgfpathlineto{\pgfqpoint{5.276198in}{0.512108in}}%
\pgfpathlineto{\pgfqpoint{5.279554in}{0.606847in}}%
\pgfpathlineto{\pgfqpoint{5.282910in}{0.631998in}}%
\pgfpathlineto{\pgfqpoint{5.286265in}{0.644790in}}%
\pgfpathlineto{\pgfqpoint{5.289621in}{0.650774in}}%
\pgfpathlineto{\pgfqpoint{5.292977in}{0.650472in}}%
\pgfpathlineto{\pgfqpoint{5.296333in}{0.641618in}}%
\pgfpathlineto{\pgfqpoint{5.299689in}{0.624825in}}%
\pgfpathlineto{\pgfqpoint{5.303044in}{0.599643in}}%
\pgfpathlineto{\pgfqpoint{5.306400in}{0.587013in}}%
\pgfpathlineto{\pgfqpoint{5.309756in}{0.596800in}}%
\pgfpathlineto{\pgfqpoint{5.313112in}{0.596191in}}%
\pgfpathlineto{\pgfqpoint{5.316467in}{0.577447in}}%
\pgfpathlineto{\pgfqpoint{5.319823in}{0.562834in}}%
\pgfpathlineto{\pgfqpoint{5.323179in}{0.603664in}}%
\pgfpathlineto{\pgfqpoint{5.326535in}{0.622199in}}%
\pgfpathlineto{\pgfqpoint{5.329891in}{0.627556in}}%
\pgfpathlineto{\pgfqpoint{5.333246in}{0.618560in}}%
\pgfpathlineto{\pgfqpoint{5.336602in}{0.598332in}}%
\pgfpathlineto{\pgfqpoint{5.339958in}{0.558815in}}%
\pgfpathlineto{\pgfqpoint{5.343314in}{0.558340in}}%
\pgfpathlineto{\pgfqpoint{5.346669in}{0.578882in}}%
\pgfpathlineto{\pgfqpoint{5.350025in}{0.578510in}}%
\pgfpathlineto{\pgfqpoint{5.356737in}{0.535103in}}%
\pgfpathlineto{\pgfqpoint{5.360093in}{0.586830in}}%
\pgfpathlineto{\pgfqpoint{5.363448in}{0.605528in}}%
\pgfpathlineto{\pgfqpoint{5.366804in}{0.610396in}}%
\pgfpathlineto{\pgfqpoint{5.370160in}{0.599243in}}%
\pgfpathlineto{\pgfqpoint{5.373516in}{0.560179in}}%
\pgfpathlineto{\pgfqpoint{5.376871in}{0.564813in}}%
\pgfpathlineto{\pgfqpoint{5.380227in}{0.599722in}}%
\pgfpathlineto{\pgfqpoint{5.383583in}{0.612008in}}%
\pgfpathlineto{\pgfqpoint{5.386939in}{0.614424in}}%
\pgfpathlineto{\pgfqpoint{5.390295in}{0.608176in}}%
\pgfpathlineto{\pgfqpoint{5.393650in}{0.587975in}}%
\pgfpathlineto{\pgfqpoint{5.397006in}{0.447069in}}%
\pgfpathlineto{\pgfqpoint{5.400362in}{0.593784in}}%
\pgfpathlineto{\pgfqpoint{5.403718in}{0.617788in}}%
\pgfpathlineto{\pgfqpoint{5.407074in}{0.627349in}}%
\pgfpathlineto{\pgfqpoint{5.410429in}{0.625618in}}%
\pgfpathlineto{\pgfqpoint{5.413785in}{0.609150in}}%
\pgfpathlineto{\pgfqpoint{5.417141in}{0.563415in}}%
\pgfpathlineto{\pgfqpoint{5.420497in}{0.574230in}}%
\pgfpathlineto{\pgfqpoint{5.423852in}{0.604582in}}%
\pgfpathlineto{\pgfqpoint{5.427208in}{0.613055in}}%
\pgfpathlineto{\pgfqpoint{5.430564in}{0.610130in}}%
\pgfpathlineto{\pgfqpoint{5.433920in}{0.592877in}}%
\pgfpathlineto{\pgfqpoint{5.437276in}{0.488073in}}%
\pgfpathlineto{\pgfqpoint{5.440631in}{0.595125in}}%
\pgfpathlineto{\pgfqpoint{5.443987in}{0.620660in}}%
\pgfpathlineto{\pgfqpoint{5.447343in}{0.630310in}}%
\pgfpathlineto{\pgfqpoint{5.450699in}{0.627125in}}%
\pgfpathlineto{\pgfqpoint{5.454054in}{0.608164in}}%
\pgfpathlineto{\pgfqpoint{5.457410in}{0.548672in}}%
\pgfpathlineto{\pgfqpoint{5.460766in}{0.585025in}}%
\pgfpathlineto{\pgfqpoint{5.464122in}{0.610098in}}%
\pgfpathlineto{\pgfqpoint{5.467478in}{0.616696in}}%
\pgfpathlineto{\pgfqpoint{5.470833in}{0.612124in}}%
\pgfpathlineto{\pgfqpoint{5.474189in}{0.591603in}}%
\pgfpathlineto{\pgfqpoint{5.477545in}{0.508956in}}%
\pgfpathlineto{\pgfqpoint{5.480901in}{0.603721in}}%
\pgfpathlineto{\pgfqpoint{5.484256in}{0.626110in}}%
\pgfpathlineto{\pgfqpoint{5.487612in}{0.634032in}}%
\pgfpathlineto{\pgfqpoint{5.490968in}{0.628646in}}%
\pgfpathlineto{\pgfqpoint{5.494324in}{0.606991in}}%
\pgfpathlineto{\pgfqpoint{5.497680in}{0.522114in}}%
\pgfpathlineto{\pgfqpoint{5.501035in}{0.593890in}}%
\pgfpathlineto{\pgfqpoint{5.504391in}{0.614962in}}%
\pgfpathlineto{\pgfqpoint{5.507747in}{0.619819in}}%
\pgfpathlineto{\pgfqpoint{5.511103in}{0.613460in}}%
\pgfpathlineto{\pgfqpoint{5.514458in}{0.588647in}}%
\pgfpathlineto{\pgfqpoint{5.517814in}{0.549940in}}%
\pgfpathlineto{\pgfqpoint{5.521170in}{0.611809in}}%
\pgfpathlineto{\pgfqpoint{5.524526in}{0.631299in}}%
\pgfpathlineto{\pgfqpoint{5.527882in}{0.637359in}}%
\pgfpathlineto{\pgfqpoint{5.531237in}{0.629612in}}%
\pgfpathlineto{\pgfqpoint{5.534593in}{0.606011in}}%
\pgfpathlineto{\pgfqpoint{5.537949in}{0.502674in}}%
\pgfpathlineto{\pgfqpoint{5.541305in}{0.594478in}}%
\pgfpathlineto{\pgfqpoint{5.544660in}{0.613378in}}%
\pgfpathlineto{\pgfqpoint{5.548016in}{0.616146in}}%
\pgfpathlineto{\pgfqpoint{5.551372in}{0.606084in}}%
\pgfpathlineto{\pgfqpoint{5.554728in}{0.566419in}}%
\pgfpathlineto{\pgfqpoint{5.558084in}{0.582002in}}%
\pgfpathlineto{\pgfqpoint{5.561439in}{0.619286in}}%
\pgfpathlineto{\pgfqpoint{5.564795in}{0.633504in}}%
\pgfpathlineto{\pgfqpoint{5.568151in}{0.634970in}}%
\pgfpathlineto{\pgfqpoint{5.571507in}{0.621609in}}%
\pgfpathlineto{\pgfqpoint{5.574863in}{0.589703in}}%
\pgfpathlineto{\pgfqpoint{5.578218in}{0.543116in}}%
\pgfpathlineto{\pgfqpoint{5.581574in}{0.597888in}}%
\pgfpathlineto{\pgfqpoint{5.584930in}{0.609619in}}%
\pgfpathlineto{\pgfqpoint{5.588286in}{0.607046in}}%
\pgfpathlineto{\pgfqpoint{5.591641in}{0.587146in}}%
\pgfpathlineto{\pgfqpoint{5.594997in}{0.514514in}}%
\pgfpathlineto{\pgfqpoint{5.598353in}{0.601965in}}%
\pgfpathlineto{\pgfqpoint{5.601709in}{0.623982in}}%
\pgfpathlineto{\pgfqpoint{5.605065in}{0.631299in}}%
\pgfpathlineto{\pgfqpoint{5.608420in}{0.624434in}}%
\pgfpathlineto{\pgfqpoint{5.611776in}{0.602157in}}%
\pgfpathlineto{\pgfqpoint{5.615132in}{0.522924in}}%
\pgfpathlineto{\pgfqpoint{5.618488in}{0.584306in}}%
\pgfpathlineto{\pgfqpoint{5.621843in}{0.605012in}}%
\pgfpathlineto{\pgfqpoint{5.625199in}{0.608286in}}%
\pgfpathlineto{\pgfqpoint{5.628555in}{0.598363in}}%
\pgfpathlineto{\pgfqpoint{5.631911in}{0.558152in}}%
\pgfpathlineto{\pgfqpoint{5.635267in}{0.575430in}}%
\pgfpathlineto{\pgfqpoint{5.638622in}{0.612304in}}%
\pgfpathlineto{\pgfqpoint{5.641978in}{0.626366in}}%
\pgfpathlineto{\pgfqpoint{5.645334in}{0.627569in}}%
\pgfpathlineto{\pgfqpoint{5.648690in}{0.613272in}}%
\pgfpathlineto{\pgfqpoint{5.652045in}{0.574661in}}%
\pgfpathlineto{\pgfqpoint{5.655401in}{0.564593in}}%
\pgfpathlineto{\pgfqpoint{5.658757in}{0.602550in}}%
\pgfpathlineto{\pgfqpoint{5.662113in}{0.612893in}}%
\pgfpathlineto{\pgfqpoint{5.665469in}{0.611269in}}%
\pgfpathlineto{\pgfqpoint{5.668824in}{0.596058in}}%
\pgfpathlineto{\pgfqpoint{5.672180in}{0.525101in}}%
\pgfpathlineto{\pgfqpoint{5.675536in}{0.591550in}}%
\pgfpathlineto{\pgfqpoint{5.678892in}{0.619589in}}%
\pgfpathlineto{\pgfqpoint{5.682247in}{0.630384in}}%
\pgfpathlineto{\pgfqpoint{5.685603in}{0.628486in}}%
\pgfpathlineto{\pgfqpoint{5.688959in}{0.610631in}}%
\pgfpathlineto{\pgfqpoint{5.692315in}{0.555079in}}%
\pgfpathlineto{\pgfqpoint{5.699026in}{0.611567in}}%
\pgfpathlineto{\pgfqpoint{5.702382in}{0.619118in}}%
\pgfpathlineto{\pgfqpoint{5.705738in}{0.615725in}}%
\pgfpathlineto{\pgfqpoint{5.709094in}{0.598062in}}%
\pgfpathlineto{\pgfqpoint{5.712449in}{0.483886in}}%
\pgfpathlineto{\pgfqpoint{5.715805in}{0.600811in}}%
\pgfpathlineto{\pgfqpoint{5.719161in}{0.626086in}}%
\pgfpathlineto{\pgfqpoint{5.722517in}{0.635639in}}%
\pgfpathlineto{\pgfqpoint{5.725873in}{0.632444in}}%
\pgfpathlineto{\pgfqpoint{5.729228in}{0.613835in}}%
\pgfpathlineto{\pgfqpoint{5.732584in}{0.558212in}}%
\pgfpathlineto{\pgfqpoint{5.739296in}{0.612524in}}%
\pgfpathlineto{\pgfqpoint{5.742652in}{0.619150in}}%
\pgfpathlineto{\pgfqpoint{5.746007in}{0.614058in}}%
\pgfpathlineto{\pgfqpoint{5.749363in}{0.591524in}}%
\pgfpathlineto{\pgfqpoint{5.752719in}{0.537536in}}%
\pgfpathlineto{\pgfqpoint{5.756075in}{0.609938in}}%
\pgfpathlineto{\pgfqpoint{5.759430in}{0.630659in}}%
\pgfpathlineto{\pgfqpoint{5.762786in}{0.637427in}}%
\pgfpathlineto{\pgfqpoint{5.766142in}{0.630415in}}%
\pgfpathlineto{\pgfqpoint{5.769498in}{0.608451in}}%
\pgfpathlineto{\pgfqpoint{5.772854in}{0.535465in}}%
\pgfpathlineto{\pgfqpoint{5.776209in}{0.587722in}}%
\pgfpathlineto{\pgfqpoint{5.779565in}{0.609022in}}%
\pgfpathlineto{\pgfqpoint{5.782921in}{0.612194in}}%
\pgfpathlineto{\pgfqpoint{5.786277in}{0.601645in}}%
\pgfpathlineto{\pgfqpoint{5.789632in}{0.557223in}}%
\pgfpathlineto{\pgfqpoint{5.792988in}{0.583664in}}%
\pgfpathlineto{\pgfqpoint{5.796344in}{0.618235in}}%
\pgfpathlineto{\pgfqpoint{5.799700in}{0.631491in}}%
\pgfpathlineto{\pgfqpoint{5.803056in}{0.631770in}}%
\pgfpathlineto{\pgfqpoint{5.806411in}{0.617047in}}%
\pgfpathlineto{\pgfqpoint{5.813123in}{0.552263in}}%
\pgfpathlineto{\pgfqpoint{5.816479in}{0.597205in}}%
\pgfpathlineto{\pgfqpoint{5.819834in}{0.607275in}}%
\pgfpathlineto{\pgfqpoint{5.823190in}{0.603398in}}%
\pgfpathlineto{\pgfqpoint{5.826546in}{0.580692in}}%
\pgfpathlineto{\pgfqpoint{5.829902in}{0.535919in}}%
\pgfpathlineto{\pgfqpoint{5.833258in}{0.602533in}}%
\pgfpathlineto{\pgfqpoint{5.836613in}{0.622651in}}%
\pgfpathlineto{\pgfqpoint{5.839969in}{0.628716in}}%
\pgfpathlineto{\pgfqpoint{5.843325in}{0.620251in}}%
\pgfpathlineto{\pgfqpoint{5.846681in}{0.594919in}}%
\pgfpathlineto{\pgfqpoint{5.850036in}{0.464888in}}%
\pgfpathlineto{\pgfqpoint{5.853392in}{0.590390in}}%
\pgfpathlineto{\pgfqpoint{5.856748in}{0.607542in}}%
\pgfpathlineto{\pgfqpoint{5.860104in}{0.609760in}}%
\pgfpathlineto{\pgfqpoint{5.863460in}{0.599482in}}%
\pgfpathlineto{\pgfqpoint{5.866815in}{0.560016in}}%
\pgfpathlineto{\pgfqpoint{5.870171in}{0.574650in}}%
\pgfpathlineto{\pgfqpoint{5.873527in}{0.612138in}}%
\pgfpathlineto{\pgfqpoint{5.876883in}{0.626444in}}%
\pgfpathlineto{\pgfqpoint{5.880238in}{0.628105in}}%
\pgfpathlineto{\pgfqpoint{5.883594in}{0.614337in}}%
\pgfpathlineto{\pgfqpoint{5.886950in}{0.575464in}}%
\pgfpathlineto{\pgfqpoint{5.890306in}{0.568138in}}%
\pgfpathlineto{\pgfqpoint{5.893662in}{0.605679in}}%
\pgfpathlineto{\pgfqpoint{5.897017in}{0.616440in}}%
\pgfpathlineto{\pgfqpoint{5.900373in}{0.615694in}}%
\pgfpathlineto{\pgfqpoint{5.903729in}{0.602664in}}%
\pgfpathlineto{\pgfqpoint{5.907085in}{0.552754in}}%
\pgfpathlineto{\pgfqpoint{5.913796in}{0.619584in}}%
\pgfpathlineto{\pgfqpoint{5.917152in}{0.632309in}}%
\pgfpathlineto{\pgfqpoint{5.920508in}{0.632804in}}%
\pgfpathlineto{\pgfqpoint{5.923864in}{0.618034in}}%
\pgfpathlineto{\pgfqpoint{5.927219in}{0.576907in}}%
\pgfpathlineto{\pgfqpoint{5.930575in}{0.575188in}}%
\pgfpathlineto{\pgfqpoint{5.933931in}{0.610352in}}%
\pgfpathlineto{\pgfqpoint{5.937287in}{0.620281in}}%
\pgfpathlineto{\pgfqpoint{5.940643in}{0.618663in}}%
\pgfpathlineto{\pgfqpoint{5.943998in}{0.603912in}}%
\pgfpathlineto{\pgfqpoint{5.947354in}{0.539807in}}%
\pgfpathlineto{\pgfqpoint{5.950710in}{0.596627in}}%
\pgfpathlineto{\pgfqpoint{5.954066in}{0.625668in}}%
\pgfpathlineto{\pgfqpoint{5.957421in}{0.636958in}}%
\pgfpathlineto{\pgfqpoint{5.960777in}{0.635784in}}%
\pgfpathlineto{\pgfqpoint{5.964133in}{0.619718in}}%
\pgfpathlineto{\pgfqpoint{5.967489in}{0.578838in}}%
\pgfpathlineto{\pgfqpoint{5.970845in}{0.572048in}}%
\pgfpathlineto{\pgfqpoint{5.974200in}{0.607434in}}%
\pgfpathlineto{\pgfqpoint{5.977556in}{0.616055in}}%
\pgfpathlineto{\pgfqpoint{5.980912in}{0.611819in}}%
\pgfpathlineto{\pgfqpoint{5.984268in}{0.589701in}}%
\pgfpathlineto{\pgfqpoint{5.987623in}{0.536531in}}%
\pgfpathlineto{\pgfqpoint{5.990979in}{0.608692in}}%
\pgfpathlineto{\pgfqpoint{5.994335in}{0.629404in}}%
\pgfpathlineto{\pgfqpoint{5.997691in}{0.635996in}}%
\pgfpathlineto{\pgfqpoint{6.001047in}{0.628498in}}%
\pgfpathlineto{\pgfqpoint{6.004402in}{0.606360in}}%
\pgfpathlineto{\pgfqpoint{6.007758in}{0.535341in}}%
\pgfpathlineto{\pgfqpoint{6.011114in}{0.583918in}}%
\pgfpathlineto{\pgfqpoint{6.014470in}{0.605111in}}%
\pgfpathlineto{\pgfqpoint{6.017825in}{0.607668in}}%
\pgfpathlineto{\pgfqpoint{6.021181in}{0.595549in}}%
\pgfpathlineto{\pgfqpoint{6.024537in}{0.539262in}}%
\pgfpathlineto{\pgfqpoint{6.027893in}{0.586646in}}%
\pgfpathlineto{\pgfqpoint{6.031249in}{0.617290in}}%
\pgfpathlineto{\pgfqpoint{6.034604in}{0.628991in}}%
\pgfpathlineto{\pgfqpoint{6.037960in}{0.627382in}}%
\pgfpathlineto{\pgfqpoint{6.041316in}{0.610562in}}%
\pgfpathlineto{\pgfqpoint{6.044672in}{0.567444in}}%
\pgfpathlineto{\pgfqpoint{6.048027in}{0.566360in}}%
\pgfpathlineto{\pgfqpoint{6.051383in}{0.599471in}}%
\pgfpathlineto{\pgfqpoint{6.054739in}{0.607289in}}%
\pgfpathlineto{\pgfqpoint{6.058095in}{0.602132in}}%
\pgfpathlineto{\pgfqpoint{6.061451in}{0.577499in}}%
\pgfpathlineto{\pgfqpoint{6.064806in}{0.542936in}}%
\pgfpathlineto{\pgfqpoint{6.068162in}{0.602760in}}%
\pgfpathlineto{\pgfqpoint{6.071518in}{0.621943in}}%
\pgfpathlineto{\pgfqpoint{6.074874in}{0.627448in}}%
\pgfpathlineto{\pgfqpoint{6.078230in}{0.618358in}}%
\pgfpathlineto{\pgfqpoint{6.081585in}{0.590717in}}%
\pgfpathlineto{\pgfqpoint{6.084941in}{0.524879in}}%
\pgfpathlineto{\pgfqpoint{6.088297in}{0.595950in}}%
\pgfpathlineto{\pgfqpoint{6.091653in}{0.611509in}}%
\pgfpathlineto{\pgfqpoint{6.095008in}{0.613639in}}%
\pgfpathlineto{\pgfqpoint{6.098364in}{0.604282in}}%
\pgfpathlineto{\pgfqpoint{6.101720in}{0.570660in}}%
\pgfpathlineto{\pgfqpoint{6.105076in}{0.568819in}}%
\pgfpathlineto{\pgfqpoint{6.108432in}{0.611778in}}%
\pgfpathlineto{\pgfqpoint{6.111787in}{0.627697in}}%
\pgfpathlineto{\pgfqpoint{6.115143in}{0.631070in}}%
\pgfpathlineto{\pgfqpoint{6.118499in}{0.619709in}}%
\pgfpathlineto{\pgfqpoint{6.121855in}{0.586599in}}%
\pgfpathlineto{\pgfqpoint{6.125210in}{0.559430in}}%
\pgfpathlineto{\pgfqpoint{6.128566in}{0.606270in}}%
\pgfpathlineto{\pgfqpoint{6.131922in}{0.619182in}}%
\pgfpathlineto{\pgfqpoint{6.135278in}{0.620122in}}%
\pgfpathlineto{\pgfqpoint{6.138634in}{0.609773in}}%
\pgfpathlineto{\pgfqpoint{6.141989in}{0.573585in}}%
\pgfpathlineto{\pgfqpoint{6.145345in}{0.578333in}}%
\pgfpathlineto{\pgfqpoint{6.148701in}{0.618761in}}%
\pgfpathlineto{\pgfqpoint{6.152057in}{0.634055in}}%
\pgfpathlineto{\pgfqpoint{6.155412in}{0.637055in}}%
\pgfpathlineto{\pgfqpoint{6.158768in}{0.625640in}}%
\pgfpathlineto{\pgfqpoint{6.162124in}{0.595077in}}%
\pgfpathlineto{\pgfqpoint{6.165480in}{0.547953in}}%
\pgfpathlineto{\pgfqpoint{6.168836in}{0.605007in}}%
\pgfpathlineto{\pgfqpoint{6.172191in}{0.618443in}}%
\pgfpathlineto{\pgfqpoint{6.175547in}{0.618627in}}%
\pgfpathlineto{\pgfqpoint{6.178903in}{0.605774in}}%
\pgfpathlineto{\pgfqpoint{6.182259in}{0.553240in}}%
\pgfpathlineto{\pgfqpoint{6.185614in}{0.592905in}}%
\pgfpathlineto{\pgfqpoint{6.188970in}{0.624536in}}%
\pgfpathlineto{\pgfqpoint{6.192326in}{0.636799in}}%
\pgfpathlineto{\pgfqpoint{6.195682in}{0.636439in}}%
\pgfpathlineto{\pgfqpoint{6.199038in}{0.621404in}}%
\pgfpathlineto{\pgfqpoint{6.202393in}{0.585914in}}%
\pgfpathlineto{\pgfqpoint{6.205749in}{0.556946in}}%
\pgfpathlineto{\pgfqpoint{6.209105in}{0.601264in}}%
\pgfpathlineto{\pgfqpoint{6.212461in}{0.610983in}}%
\pgfpathlineto{\pgfqpoint{6.215816in}{0.606540in}}%
\pgfpathlineto{\pgfqpoint{6.219172in}{0.582037in}}%
\pgfpathlineto{\pgfqpoint{6.222528in}{0.549526in}}%
\pgfpathlineto{\pgfqpoint{6.225884in}{0.608409in}}%
\pgfpathlineto{\pgfqpoint{6.229240in}{0.627470in}}%
\pgfpathlineto{\pgfqpoint{6.232595in}{0.632743in}}%
\pgfpathlineto{\pgfqpoint{6.235951in}{0.623399in}}%
\pgfpathlineto{\pgfqpoint{6.239307in}{0.598751in}}%
\pgfpathlineto{\pgfqpoint{6.242663in}{0.486561in}}%
\pgfpathlineto{\pgfqpoint{6.246019in}{0.585813in}}%
\pgfpathlineto{\pgfqpoint{6.249374in}{0.603469in}}%
\pgfpathlineto{\pgfqpoint{6.252730in}{0.604328in}}%
\pgfpathlineto{\pgfqpoint{6.256086in}{0.589914in}}%
\pgfpathlineto{\pgfqpoint{6.259442in}{0.511172in}}%
\pgfpathlineto{\pgfqpoint{6.262797in}{0.589144in}}%
\pgfpathlineto{\pgfqpoint{6.266153in}{0.616449in}}%
\pgfpathlineto{\pgfqpoint{6.269509in}{0.626644in}}%
\pgfpathlineto{\pgfqpoint{6.272865in}{0.623231in}}%
\pgfpathlineto{\pgfqpoint{6.276221in}{0.604029in}}%
\pgfpathlineto{\pgfqpoint{6.279576in}{0.544999in}}%
\pgfpathlineto{\pgfqpoint{6.282932in}{0.578584in}}%
\pgfpathlineto{\pgfqpoint{6.286288in}{0.603891in}}%
\pgfpathlineto{\pgfqpoint{6.289644in}{0.610036in}}%
\pgfpathlineto{\pgfqpoint{6.292999in}{0.604412in}}%
\pgfpathlineto{\pgfqpoint{6.296355in}{0.580498in}}%
\pgfpathlineto{\pgfqpoint{6.299711in}{0.537711in}}%
\pgfpathlineto{\pgfqpoint{6.303067in}{0.602397in}}%
\pgfpathlineto{\pgfqpoint{6.306423in}{0.622244in}}%
\pgfpathlineto{\pgfqpoint{6.309778in}{0.628419in}}%
\pgfpathlineto{\pgfqpoint{6.313134in}{0.620470in}}%
\pgfpathlineto{\pgfqpoint{6.316490in}{0.593749in}}%
\pgfpathlineto{\pgfqpoint{6.319846in}{0.525069in}}%
\pgfpathlineto{\pgfqpoint{6.323201in}{0.599202in}}%
\pgfpathlineto{\pgfqpoint{6.326557in}{0.615600in}}%
\pgfpathlineto{\pgfqpoint{6.329913in}{0.618781in}}%
\pgfpathlineto{\pgfqpoint{6.333269in}{0.611390in}}%
\pgfpathlineto{\pgfqpoint{6.336625in}{0.585205in}}%
\pgfpathlineto{\pgfqpoint{6.339980in}{0.552242in}}%
\pgfpathlineto{\pgfqpoint{6.343336in}{0.610365in}}%
\pgfpathlineto{\pgfqpoint{6.346692in}{0.629336in}}%
\pgfpathlineto{\pgfqpoint{6.350048in}{0.635236in}}%
\pgfpathlineto{\pgfqpoint{6.353403in}{0.628049in}}%
\pgfpathlineto{\pgfqpoint{6.356759in}{0.613266in}}%
\pgfpathlineto{\pgfqpoint{6.360115in}{0.592531in}}%
\pgfpathlineto{\pgfqpoint{6.366827in}{0.531114in}}%
\pgfpathlineto{\pgfqpoint{6.370182in}{0.528216in}}%
\pgfpathlineto{\pgfqpoint{6.373538in}{0.555432in}}%
\pgfpathlineto{\pgfqpoint{6.376894in}{0.577212in}}%
\pgfpathlineto{\pgfqpoint{6.380250in}{0.588616in}}%
\pgfpathlineto{\pgfqpoint{6.383605in}{0.587339in}}%
\pgfpathlineto{\pgfqpoint{6.383605in}{0.587339in}}%
\pgfusepath{stroke}%
\end{pgfscope}%
\begin{pgfscope}%
\pgfpathrectangle{\pgfqpoint{4.621821in}{0.320625in}}{\pgfqpoint{1.845679in}{1.164375in}}%
\pgfusepath{clip}%
\pgfsetrectcap%
\pgfsetroundjoin%
\pgfsetlinewidth{0.602250pt}%
\definecolor{currentstroke}{rgb}{0.141176,0.345098,0.650980}%
\pgfsetstrokecolor{currentstroke}%
\pgfsetstrokeopacity{0.120000}%
\pgfsetdash{}{0pt}%
\pgfpathmoveto{\pgfqpoint{4.705715in}{1.479263in}}%
\pgfpathlineto{\pgfqpoint{4.712427in}{1.467492in}}%
\pgfpathlineto{\pgfqpoint{4.722494in}{1.453345in}}%
\pgfpathlineto{\pgfqpoint{4.742629in}{1.427571in}}%
\pgfpathlineto{\pgfqpoint{4.766120in}{1.399100in}}%
\pgfpathlineto{\pgfqpoint{4.782898in}{1.380408in}}%
\pgfpathlineto{\pgfqpoint{4.799677in}{1.363349in}}%
\pgfpathlineto{\pgfqpoint{4.819812in}{1.344689in}}%
\pgfpathlineto{\pgfqpoint{4.843302in}{1.324650in}}%
\pgfpathlineto{\pgfqpoint{4.873504in}{1.300485in}}%
\pgfpathlineto{\pgfqpoint{4.917130in}{1.267139in}}%
\pgfpathlineto{\pgfqpoint{4.990957in}{1.212326in}}%
\pgfpathlineto{\pgfqpoint{5.128544in}{1.111783in}}%
\pgfpathlineto{\pgfqpoint{5.410429in}{0.907360in}}%
\pgfpathlineto{\pgfqpoint{5.578218in}{0.786396in}}%
\pgfpathlineto{\pgfqpoint{5.584930in}{0.784073in}}%
\pgfpathlineto{\pgfqpoint{5.591641in}{0.776572in}}%
\pgfpathlineto{\pgfqpoint{5.594997in}{0.774523in}}%
\pgfpathlineto{\pgfqpoint{5.598353in}{0.773502in}}%
\pgfpathlineto{\pgfqpoint{5.601709in}{0.771288in}}%
\pgfpathlineto{\pgfqpoint{5.605065in}{0.767409in}}%
\pgfpathlineto{\pgfqpoint{5.608420in}{0.764612in}}%
\pgfpathlineto{\pgfqpoint{5.615132in}{0.768018in}}%
\pgfpathlineto{\pgfqpoint{5.618488in}{0.763483in}}%
\pgfpathlineto{\pgfqpoint{5.621843in}{0.756215in}}%
\pgfpathlineto{\pgfqpoint{5.625199in}{0.752248in}}%
\pgfpathlineto{\pgfqpoint{5.628555in}{0.751104in}}%
\pgfpathlineto{\pgfqpoint{5.631911in}{0.748776in}}%
\pgfpathlineto{\pgfqpoint{5.635267in}{0.744981in}}%
\pgfpathlineto{\pgfqpoint{5.638622in}{0.744668in}}%
\pgfpathlineto{\pgfqpoint{5.641978in}{0.748569in}}%
\pgfpathlineto{\pgfqpoint{5.645334in}{0.747722in}}%
\pgfpathlineto{\pgfqpoint{5.648690in}{0.739467in}}%
\pgfpathlineto{\pgfqpoint{5.652045in}{0.732964in}}%
\pgfpathlineto{\pgfqpoint{5.658757in}{0.730017in}}%
\pgfpathlineto{\pgfqpoint{5.662113in}{0.726177in}}%
\pgfpathlineto{\pgfqpoint{5.665469in}{0.723872in}}%
\pgfpathlineto{\pgfqpoint{5.668824in}{0.728076in}}%
\pgfpathlineto{\pgfqpoint{5.672180in}{0.731017in}}%
\pgfpathlineto{\pgfqpoint{5.675536in}{0.723491in}}%
\pgfpathlineto{\pgfqpoint{5.678892in}{0.713717in}}%
\pgfpathlineto{\pgfqpoint{5.682247in}{0.714554in}}%
\pgfpathlineto{\pgfqpoint{5.685603in}{0.718450in}}%
\pgfpathlineto{\pgfqpoint{5.688959in}{0.716062in}}%
\pgfpathlineto{\pgfqpoint{5.692315in}{0.706477in}}%
\pgfpathlineto{\pgfqpoint{5.695671in}{0.704159in}}%
\pgfpathlineto{\pgfqpoint{5.699026in}{0.718108in}}%
\pgfpathlineto{\pgfqpoint{5.702382in}{0.725287in}}%
\pgfpathlineto{\pgfqpoint{5.705738in}{0.715011in}}%
\pgfpathlineto{\pgfqpoint{5.709094in}{0.691667in}}%
\pgfpathlineto{\pgfqpoint{5.712449in}{0.711415in}}%
\pgfpathlineto{\pgfqpoint{5.715805in}{0.726449in}}%
\pgfpathlineto{\pgfqpoint{5.719161in}{0.728020in}}%
\pgfpathlineto{\pgfqpoint{5.722517in}{0.715791in}}%
\pgfpathlineto{\pgfqpoint{5.725873in}{0.681349in}}%
\pgfpathlineto{\pgfqpoint{5.729228in}{0.712631in}}%
\pgfpathlineto{\pgfqpoint{5.732584in}{0.735708in}}%
\pgfpathlineto{\pgfqpoint{5.735940in}{0.740075in}}%
\pgfpathlineto{\pgfqpoint{5.739296in}{0.721127in}}%
\pgfpathlineto{\pgfqpoint{5.742652in}{0.669872in}}%
\pgfpathlineto{\pgfqpoint{5.746007in}{0.723302in}}%
\pgfpathlineto{\pgfqpoint{5.749363in}{0.737739in}}%
\pgfpathlineto{\pgfqpoint{5.752719in}{0.735506in}}%
\pgfpathlineto{\pgfqpoint{5.756075in}{0.712559in}}%
\pgfpathlineto{\pgfqpoint{5.759430in}{0.683870in}}%
\pgfpathlineto{\pgfqpoint{5.762786in}{0.736857in}}%
\pgfpathlineto{\pgfqpoint{5.766142in}{0.752117in}}%
\pgfpathlineto{\pgfqpoint{5.769498in}{0.746943in}}%
\pgfpathlineto{\pgfqpoint{5.772854in}{0.718675in}}%
\pgfpathlineto{\pgfqpoint{5.776209in}{0.656417in}}%
\pgfpathlineto{\pgfqpoint{5.779565in}{0.713388in}}%
\pgfpathlineto{\pgfqpoint{5.782921in}{0.718149in}}%
\pgfpathlineto{\pgfqpoint{5.786277in}{0.692395in}}%
\pgfpathlineto{\pgfqpoint{5.789632in}{0.689421in}}%
\pgfpathlineto{\pgfqpoint{5.792988in}{0.732355in}}%
\pgfpathlineto{\pgfqpoint{5.796344in}{0.743258in}}%
\pgfpathlineto{\pgfqpoint{5.799700in}{0.730818in}}%
\pgfpathlineto{\pgfqpoint{5.803056in}{0.697432in}}%
\pgfpathlineto{\pgfqpoint{5.806411in}{0.640716in}}%
\pgfpathlineto{\pgfqpoint{5.809767in}{0.686139in}}%
\pgfpathlineto{\pgfqpoint{5.813123in}{0.676465in}}%
\pgfpathlineto{\pgfqpoint{5.816479in}{0.643518in}}%
\pgfpathlineto{\pgfqpoint{5.819834in}{0.705367in}}%
\pgfpathlineto{\pgfqpoint{5.823190in}{0.721942in}}%
\pgfpathlineto{\pgfqpoint{5.826546in}{0.714007in}}%
\pgfpathlineto{\pgfqpoint{5.833258in}{0.648828in}}%
\pgfpathlineto{\pgfqpoint{5.836613in}{0.688292in}}%
\pgfpathlineto{\pgfqpoint{5.839969in}{0.689059in}}%
\pgfpathlineto{\pgfqpoint{5.843325in}{0.655970in}}%
\pgfpathlineto{\pgfqpoint{5.846681in}{0.669456in}}%
\pgfpathlineto{\pgfqpoint{5.850036in}{0.706238in}}%
\pgfpathlineto{\pgfqpoint{5.853392in}{0.714580in}}%
\pgfpathlineto{\pgfqpoint{5.856748in}{0.696283in}}%
\pgfpathlineto{\pgfqpoint{5.860104in}{0.608287in}}%
\pgfpathlineto{\pgfqpoint{5.863460in}{0.699105in}}%
\pgfpathlineto{\pgfqpoint{5.866815in}{0.714803in}}%
\pgfpathlineto{\pgfqpoint{5.870171in}{0.713662in}}%
\pgfpathlineto{\pgfqpoint{5.873527in}{0.693234in}}%
\pgfpathlineto{\pgfqpoint{5.876883in}{0.637873in}}%
\pgfpathlineto{\pgfqpoint{5.880238in}{0.710638in}}%
\pgfpathlineto{\pgfqpoint{5.883594in}{0.728433in}}%
\pgfpathlineto{\pgfqpoint{5.886950in}{0.726128in}}%
\pgfpathlineto{\pgfqpoint{5.890306in}{0.688138in}}%
\pgfpathlineto{\pgfqpoint{5.893662in}{0.696340in}}%
\pgfpathlineto{\pgfqpoint{5.897017in}{0.729848in}}%
\pgfpathlineto{\pgfqpoint{5.900373in}{0.738219in}}%
\pgfpathlineto{\pgfqpoint{5.903729in}{0.732551in}}%
\pgfpathlineto{\pgfqpoint{5.907085in}{0.703741in}}%
\pgfpathlineto{\pgfqpoint{5.910441in}{0.689995in}}%
\pgfpathlineto{\pgfqpoint{5.913796in}{0.736969in}}%
\pgfpathlineto{\pgfqpoint{5.917152in}{0.750595in}}%
\pgfpathlineto{\pgfqpoint{5.920508in}{0.744657in}}%
\pgfpathlineto{\pgfqpoint{5.923864in}{0.710719in}}%
\pgfpathlineto{\pgfqpoint{5.927219in}{0.688261in}}%
\pgfpathlineto{\pgfqpoint{5.930575in}{0.728253in}}%
\pgfpathlineto{\pgfqpoint{5.933931in}{0.733241in}}%
\pgfpathlineto{\pgfqpoint{5.937287in}{0.716535in}}%
\pgfpathlineto{\pgfqpoint{5.940643in}{0.646523in}}%
\pgfpathlineto{\pgfqpoint{5.943998in}{0.732131in}}%
\pgfpathlineto{\pgfqpoint{5.947354in}{0.750613in}}%
\pgfpathlineto{\pgfqpoint{5.950710in}{0.746929in}}%
\pgfpathlineto{\pgfqpoint{5.954066in}{0.722763in}}%
\pgfpathlineto{\pgfqpoint{5.957421in}{0.649250in}}%
\pgfpathlineto{\pgfqpoint{5.960777in}{0.687829in}}%
\pgfpathlineto{\pgfqpoint{5.964133in}{0.694580in}}%
\pgfpathlineto{\pgfqpoint{5.967489in}{0.640173in}}%
\pgfpathlineto{\pgfqpoint{5.970845in}{0.698543in}}%
\pgfpathlineto{\pgfqpoint{5.974200in}{0.726909in}}%
\pgfpathlineto{\pgfqpoint{5.977556in}{0.729824in}}%
\pgfpathlineto{\pgfqpoint{5.980912in}{0.709312in}}%
\pgfpathlineto{\pgfqpoint{5.984268in}{0.654149in}}%
\pgfpathlineto{\pgfqpoint{5.987623in}{0.665216in}}%
\pgfpathlineto{\pgfqpoint{5.990979in}{0.680844in}}%
\pgfpathlineto{\pgfqpoint{5.994335in}{0.654554in}}%
\pgfpathlineto{\pgfqpoint{5.997691in}{0.662928in}}%
\pgfpathlineto{\pgfqpoint{6.001047in}{0.703690in}}%
\pgfpathlineto{\pgfqpoint{6.004402in}{0.713448in}}%
\pgfpathlineto{\pgfqpoint{6.007758in}{0.697041in}}%
\pgfpathlineto{\pgfqpoint{6.011114in}{0.609815in}}%
\pgfpathlineto{\pgfqpoint{6.014470in}{0.683839in}}%
\pgfpathlineto{\pgfqpoint{6.017825in}{0.701575in}}%
\pgfpathlineto{\pgfqpoint{6.021181in}{0.698919in}}%
\pgfpathlineto{\pgfqpoint{6.024537in}{0.670091in}}%
\pgfpathlineto{\pgfqpoint{6.027893in}{0.663150in}}%
\pgfpathlineto{\pgfqpoint{6.031249in}{0.707424in}}%
\pgfpathlineto{\pgfqpoint{6.034604in}{0.719249in}}%
\pgfpathlineto{\pgfqpoint{6.037960in}{0.707594in}}%
\pgfpathlineto{\pgfqpoint{6.041316in}{0.584525in}}%
\pgfpathlineto{\pgfqpoint{6.044672in}{0.707221in}}%
\pgfpathlineto{\pgfqpoint{6.048027in}{0.727454in}}%
\pgfpathlineto{\pgfqpoint{6.051383in}{0.730882in}}%
\pgfpathlineto{\pgfqpoint{6.054739in}{0.720361in}}%
\pgfpathlineto{\pgfqpoint{6.058095in}{0.674569in}}%
\pgfpathlineto{\pgfqpoint{6.064806in}{0.734852in}}%
\pgfpathlineto{\pgfqpoint{6.068162in}{0.743678in}}%
\pgfpathlineto{\pgfqpoint{6.071518in}{0.732675in}}%
\pgfpathlineto{\pgfqpoint{6.074874in}{0.698103in}}%
\pgfpathlineto{\pgfqpoint{6.078230in}{0.656371in}}%
\pgfpathlineto{\pgfqpoint{6.081585in}{0.699740in}}%
\pgfpathlineto{\pgfqpoint{6.084941in}{0.697026in}}%
\pgfpathlineto{\pgfqpoint{6.088297in}{0.633204in}}%
\pgfpathlineto{\pgfqpoint{6.091653in}{0.699085in}}%
\pgfpathlineto{\pgfqpoint{6.095008in}{0.725532in}}%
\pgfpathlineto{\pgfqpoint{6.098364in}{0.727271in}}%
\pgfpathlineto{\pgfqpoint{6.101720in}{0.705283in}}%
\pgfpathlineto{\pgfqpoint{6.105076in}{0.632467in}}%
\pgfpathlineto{\pgfqpoint{6.108432in}{0.676157in}}%
\pgfpathlineto{\pgfqpoint{6.111787in}{0.688450in}}%
\pgfpathlineto{\pgfqpoint{6.115143in}{0.667431in}}%
\pgfpathlineto{\pgfqpoint{6.118499in}{0.653246in}}%
\pgfpathlineto{\pgfqpoint{6.121855in}{0.703245in}}%
\pgfpathlineto{\pgfqpoint{6.125210in}{0.716048in}}%
\pgfpathlineto{\pgfqpoint{6.128566in}{0.703792in}}%
\pgfpathlineto{\pgfqpoint{6.131922in}{0.641678in}}%
\pgfpathlineto{\pgfqpoint{6.135278in}{0.683534in}}%
\pgfpathlineto{\pgfqpoint{6.138634in}{0.705428in}}%
\pgfpathlineto{\pgfqpoint{6.141989in}{0.705787in}}%
\pgfpathlineto{\pgfqpoint{6.145345in}{0.684787in}}%
\pgfpathlineto{\pgfqpoint{6.148701in}{0.641963in}}%
\pgfpathlineto{\pgfqpoint{6.152057in}{0.706481in}}%
\pgfpathlineto{\pgfqpoint{6.155412in}{0.722795in}}%
\pgfpathlineto{\pgfqpoint{6.158768in}{0.717516in}}%
\pgfpathlineto{\pgfqpoint{6.162124in}{0.668765in}}%
\pgfpathlineto{\pgfqpoint{6.168836in}{0.725199in}}%
\pgfpathlineto{\pgfqpoint{6.172191in}{0.731474in}}%
\pgfpathlineto{\pgfqpoint{6.175547in}{0.723532in}}%
\pgfpathlineto{\pgfqpoint{6.178903in}{0.687097in}}%
\pgfpathlineto{\pgfqpoint{6.182259in}{0.695950in}}%
\pgfpathlineto{\pgfqpoint{6.185614in}{0.733652in}}%
\pgfpathlineto{\pgfqpoint{6.188970in}{0.744401in}}%
\pgfpathlineto{\pgfqpoint{6.192326in}{0.734606in}}%
\pgfpathlineto{\pgfqpoint{6.195682in}{0.684066in}}%
\pgfpathlineto{\pgfqpoint{6.202393in}{0.734042in}}%
\pgfpathlineto{\pgfqpoint{6.205749in}{0.736761in}}%
\pgfpathlineto{\pgfqpoint{6.209105in}{0.720579in}}%
\pgfpathlineto{\pgfqpoint{6.212461in}{0.609371in}}%
\pgfpathlineto{\pgfqpoint{6.215816in}{0.730310in}}%
\pgfpathlineto{\pgfqpoint{6.219172in}{0.750770in}}%
\pgfpathlineto{\pgfqpoint{6.222528in}{0.750137in}}%
\pgfpathlineto{\pgfqpoint{6.225884in}{0.727451in}}%
\pgfpathlineto{\pgfqpoint{6.229240in}{0.649207in}}%
\pgfpathlineto{\pgfqpoint{6.232595in}{0.700083in}}%
\pgfpathlineto{\pgfqpoint{6.235951in}{0.711166in}}%
\pgfpathlineto{\pgfqpoint{6.239307in}{0.688164in}}%
\pgfpathlineto{\pgfqpoint{6.242663in}{0.680741in}}%
\pgfpathlineto{\pgfqpoint{6.246019in}{0.727371in}}%
\pgfpathlineto{\pgfqpoint{6.249374in}{0.739180in}}%
\pgfpathlineto{\pgfqpoint{6.252730in}{0.727167in}}%
\pgfpathlineto{\pgfqpoint{6.256086in}{0.693893in}}%
\pgfpathlineto{\pgfqpoint{6.259442in}{0.626727in}}%
\pgfpathlineto{\pgfqpoint{6.262797in}{0.682014in}}%
\pgfpathlineto{\pgfqpoint{6.266153in}{0.673032in}}%
\pgfpathlineto{\pgfqpoint{6.269509in}{0.630489in}}%
\pgfpathlineto{\pgfqpoint{6.272865in}{0.700716in}}%
\pgfpathlineto{\pgfqpoint{6.276221in}{0.718007in}}%
\pgfpathlineto{\pgfqpoint{6.279576in}{0.710646in}}%
\pgfpathlineto{\pgfqpoint{6.282932in}{0.674043in}}%
\pgfpathlineto{\pgfqpoint{6.286288in}{0.655304in}}%
\pgfpathlineto{\pgfqpoint{6.289644in}{0.691596in}}%
\pgfpathlineto{\pgfqpoint{6.292999in}{0.693625in}}%
\pgfpathlineto{\pgfqpoint{6.296355in}{0.668509in}}%
\pgfpathlineto{\pgfqpoint{6.299711in}{0.655404in}}%
\pgfpathlineto{\pgfqpoint{6.303067in}{0.703267in}}%
\pgfpathlineto{\pgfqpoint{6.306423in}{0.715676in}}%
\pgfpathlineto{\pgfqpoint{6.309778in}{0.703621in}}%
\pgfpathlineto{\pgfqpoint{6.313134in}{0.608362in}}%
\pgfpathlineto{\pgfqpoint{6.316490in}{0.698311in}}%
\pgfpathlineto{\pgfqpoint{6.319846in}{0.718593in}}%
\pgfpathlineto{\pgfqpoint{6.323201in}{0.721060in}}%
\pgfpathlineto{\pgfqpoint{6.326557in}{0.707935in}}%
\pgfpathlineto{\pgfqpoint{6.329913in}{0.638827in}}%
\pgfpathlineto{\pgfqpoint{6.333269in}{0.703395in}}%
\pgfpathlineto{\pgfqpoint{6.336625in}{0.729504in}}%
\pgfpathlineto{\pgfqpoint{6.339980in}{0.734541in}}%
\pgfpathlineto{\pgfqpoint{6.343336in}{0.714613in}}%
\pgfpathlineto{\pgfqpoint{6.346692in}{0.653387in}}%
\pgfpathlineto{\pgfqpoint{6.350048in}{0.724586in}}%
\pgfpathlineto{\pgfqpoint{6.353403in}{0.739478in}}%
\pgfpathlineto{\pgfqpoint{6.356759in}{0.738707in}}%
\pgfpathlineto{\pgfqpoint{6.360115in}{0.720381in}}%
\pgfpathlineto{\pgfqpoint{6.363471in}{0.629151in}}%
\pgfpathlineto{\pgfqpoint{6.366827in}{0.731152in}}%
\pgfpathlineto{\pgfqpoint{6.370182in}{0.751264in}}%
\pgfpathlineto{\pgfqpoint{6.373538in}{0.751845in}}%
\pgfpathlineto{\pgfqpoint{6.383605in}{0.710788in}}%
\pgfpathlineto{\pgfqpoint{6.383605in}{0.710788in}}%
\pgfusepath{stroke}%
\end{pgfscope}%
\begin{pgfscope}%
\pgfpathrectangle{\pgfqpoint{4.621821in}{0.320625in}}{\pgfqpoint{1.845679in}{1.164375in}}%
\pgfusepath{clip}%
\pgfsetrectcap%
\pgfsetroundjoin%
\pgfsetlinewidth{0.602250pt}%
\definecolor{currentstroke}{rgb}{0.141176,0.345098,0.650980}%
\pgfsetstrokecolor{currentstroke}%
\pgfsetstrokeopacity{0.120000}%
\pgfsetdash{}{0pt}%
\pgfpathmoveto{\pgfqpoint{4.705715in}{1.478401in}}%
\pgfpathlineto{\pgfqpoint{4.712427in}{1.466304in}}%
\pgfpathlineto{\pgfqpoint{4.722494in}{1.451315in}}%
\pgfpathlineto{\pgfqpoint{4.742629in}{1.423879in}}%
\pgfpathlineto{\pgfqpoint{4.759408in}{1.402982in}}%
\pgfpathlineto{\pgfqpoint{4.772831in}{1.388001in}}%
\pgfpathlineto{\pgfqpoint{4.786254in}{1.374521in}}%
\pgfpathlineto{\pgfqpoint{4.803033in}{1.359412in}}%
\pgfpathlineto{\pgfqpoint{4.823168in}{1.343101in}}%
\pgfpathlineto{\pgfqpoint{4.846658in}{1.325590in}}%
\pgfpathlineto{\pgfqpoint{4.880216in}{1.302021in}}%
\pgfpathlineto{\pgfqpoint{4.930553in}{1.268096in}}%
\pgfpathlineto{\pgfqpoint{5.001024in}{1.221994in}}%
\pgfpathlineto{\pgfqpoint{5.088274in}{1.166303in}}%
\pgfpathlineto{\pgfqpoint{5.195659in}{1.099148in}}%
\pgfpathlineto{\pgfqpoint{5.343314in}{1.008243in}}%
\pgfpathlineto{\pgfqpoint{5.581574in}{0.863012in}}%
\pgfpathlineto{\pgfqpoint{5.722517in}{0.778033in}}%
\pgfpathlineto{\pgfqpoint{5.725873in}{0.776765in}}%
\pgfpathlineto{\pgfqpoint{5.732584in}{0.771506in}}%
\pgfpathlineto{\pgfqpoint{5.739296in}{0.771404in}}%
\pgfpathlineto{\pgfqpoint{5.742652in}{0.768050in}}%
\pgfpathlineto{\pgfqpoint{5.746007in}{0.763269in}}%
\pgfpathlineto{\pgfqpoint{5.749363in}{0.763684in}}%
\pgfpathlineto{\pgfqpoint{5.752719in}{0.768343in}}%
\pgfpathlineto{\pgfqpoint{5.756075in}{0.766583in}}%
\pgfpathlineto{\pgfqpoint{5.759430in}{0.757651in}}%
\pgfpathlineto{\pgfqpoint{5.762786in}{0.752999in}}%
\pgfpathlineto{\pgfqpoint{5.766142in}{0.752438in}}%
\pgfpathlineto{\pgfqpoint{5.772854in}{0.747364in}}%
\pgfpathlineto{\pgfqpoint{5.776209in}{0.751384in}}%
\pgfpathlineto{\pgfqpoint{5.779565in}{0.752081in}}%
\pgfpathlineto{\pgfqpoint{5.782921in}{0.743769in}}%
\pgfpathlineto{\pgfqpoint{5.786277in}{0.738694in}}%
\pgfpathlineto{\pgfqpoint{5.789632in}{0.737924in}}%
\pgfpathlineto{\pgfqpoint{5.796344in}{0.732809in}}%
\pgfpathlineto{\pgfqpoint{5.799700in}{0.735820in}}%
\pgfpathlineto{\pgfqpoint{5.803056in}{0.737365in}}%
\pgfpathlineto{\pgfqpoint{5.806411in}{0.728971in}}%
\pgfpathlineto{\pgfqpoint{5.809767in}{0.725355in}}%
\pgfpathlineto{\pgfqpoint{5.813123in}{0.729603in}}%
\pgfpathlineto{\pgfqpoint{5.816479in}{0.729046in}}%
\pgfpathlineto{\pgfqpoint{5.819834in}{0.720824in}}%
\pgfpathlineto{\pgfqpoint{5.823190in}{0.719018in}}%
\pgfpathlineto{\pgfqpoint{5.826546in}{0.731374in}}%
\pgfpathlineto{\pgfqpoint{5.829902in}{0.733731in}}%
\pgfpathlineto{\pgfqpoint{5.833258in}{0.714465in}}%
\pgfpathlineto{\pgfqpoint{5.836613in}{0.721342in}}%
\pgfpathlineto{\pgfqpoint{5.839969in}{0.739839in}}%
\pgfpathlineto{\pgfqpoint{5.843325in}{0.742231in}}%
\pgfpathlineto{\pgfqpoint{5.846681in}{0.727448in}}%
\pgfpathlineto{\pgfqpoint{5.850036in}{0.701748in}}%
\pgfpathlineto{\pgfqpoint{5.853392in}{0.739099in}}%
\pgfpathlineto{\pgfqpoint{5.856748in}{0.753535in}}%
\pgfpathlineto{\pgfqpoint{5.860104in}{0.743179in}}%
\pgfpathlineto{\pgfqpoint{5.863460in}{0.707979in}}%
\pgfpathlineto{\pgfqpoint{5.866815in}{0.701128in}}%
\pgfpathlineto{\pgfqpoint{5.870171in}{0.712343in}}%
\pgfpathlineto{\pgfqpoint{5.873527in}{0.692095in}}%
\pgfpathlineto{\pgfqpoint{5.876883in}{0.709730in}}%
\pgfpathlineto{\pgfqpoint{5.880238in}{0.735214in}}%
\pgfpathlineto{\pgfqpoint{5.883594in}{0.730445in}}%
\pgfpathlineto{\pgfqpoint{5.886950in}{0.698955in}}%
\pgfpathlineto{\pgfqpoint{5.890306in}{0.682480in}}%
\pgfpathlineto{\pgfqpoint{5.893662in}{0.697239in}}%
\pgfpathlineto{\pgfqpoint{5.897017in}{0.681556in}}%
\pgfpathlineto{\pgfqpoint{5.900373in}{0.688524in}}%
\pgfpathlineto{\pgfqpoint{5.903729in}{0.717802in}}%
\pgfpathlineto{\pgfqpoint{5.907085in}{0.717271in}}%
\pgfpathlineto{\pgfqpoint{5.910441in}{0.678971in}}%
\pgfpathlineto{\pgfqpoint{5.913796in}{0.692407in}}%
\pgfpathlineto{\pgfqpoint{5.917152in}{0.712893in}}%
\pgfpathlineto{\pgfqpoint{5.920508in}{0.708289in}}%
\pgfpathlineto{\pgfqpoint{5.923864in}{0.667359in}}%
\pgfpathlineto{\pgfqpoint{5.927219in}{0.702141in}}%
\pgfpathlineto{\pgfqpoint{5.930575in}{0.727055in}}%
\pgfpathlineto{\pgfqpoint{5.933931in}{0.722266in}}%
\pgfpathlineto{\pgfqpoint{5.937287in}{0.650662in}}%
\pgfpathlineto{\pgfqpoint{5.940643in}{0.720938in}}%
\pgfpathlineto{\pgfqpoint{5.943998in}{0.738721in}}%
\pgfpathlineto{\pgfqpoint{5.947354in}{0.736180in}}%
\pgfpathlineto{\pgfqpoint{5.950710in}{0.706315in}}%
\pgfpathlineto{\pgfqpoint{5.954066in}{0.704187in}}%
\pgfpathlineto{\pgfqpoint{5.957421in}{0.744538in}}%
\pgfpathlineto{\pgfqpoint{5.960777in}{0.751832in}}%
\pgfpathlineto{\pgfqpoint{5.964133in}{0.731217in}}%
\pgfpathlineto{\pgfqpoint{5.967489in}{0.653821in}}%
\pgfpathlineto{\pgfqpoint{5.970845in}{0.701724in}}%
\pgfpathlineto{\pgfqpoint{5.974200in}{0.703555in}}%
\pgfpathlineto{\pgfqpoint{5.977556in}{0.622878in}}%
\pgfpathlineto{\pgfqpoint{5.980912in}{0.718981in}}%
\pgfpathlineto{\pgfqpoint{5.984268in}{0.736616in}}%
\pgfpathlineto{\pgfqpoint{5.987623in}{0.721328in}}%
\pgfpathlineto{\pgfqpoint{5.990979in}{0.663795in}}%
\pgfpathlineto{\pgfqpoint{5.994335in}{0.681878in}}%
\pgfpathlineto{\pgfqpoint{5.997691in}{0.691114in}}%
\pgfpathlineto{\pgfqpoint{6.001047in}{0.633989in}}%
\pgfpathlineto{\pgfqpoint{6.004402in}{0.697720in}}%
\pgfpathlineto{\pgfqpoint{6.007758in}{0.720877in}}%
\pgfpathlineto{\pgfqpoint{6.011114in}{0.709691in}}%
\pgfpathlineto{\pgfqpoint{6.014470in}{0.612253in}}%
\pgfpathlineto{\pgfqpoint{6.017825in}{0.700350in}}%
\pgfpathlineto{\pgfqpoint{6.021181in}{0.713320in}}%
\pgfpathlineto{\pgfqpoint{6.024537in}{0.700459in}}%
\pgfpathlineto{\pgfqpoint{6.027893in}{0.617172in}}%
\pgfpathlineto{\pgfqpoint{6.031249in}{0.713217in}}%
\pgfpathlineto{\pgfqpoint{6.034604in}{0.729279in}}%
\pgfpathlineto{\pgfqpoint{6.037960in}{0.712491in}}%
\pgfpathlineto{\pgfqpoint{6.041316in}{0.671383in}}%
\pgfpathlineto{\pgfqpoint{6.044672in}{0.728473in}}%
\pgfpathlineto{\pgfqpoint{6.048027in}{0.739057in}}%
\pgfpathlineto{\pgfqpoint{6.051383in}{0.730071in}}%
\pgfpathlineto{\pgfqpoint{6.054739in}{0.671006in}}%
\pgfpathlineto{\pgfqpoint{6.058095in}{0.723598in}}%
\pgfpathlineto{\pgfqpoint{6.061451in}{0.749147in}}%
\pgfpathlineto{\pgfqpoint{6.064806in}{0.747198in}}%
\pgfpathlineto{\pgfqpoint{6.068162in}{0.715753in}}%
\pgfpathlineto{\pgfqpoint{6.071518in}{0.659971in}}%
\pgfpathlineto{\pgfqpoint{6.074874in}{0.707779in}}%
\pgfpathlineto{\pgfqpoint{6.078230in}{0.693979in}}%
\pgfpathlineto{\pgfqpoint{6.081585in}{0.676865in}}%
\pgfpathlineto{\pgfqpoint{6.084941in}{0.728397in}}%
\pgfpathlineto{\pgfqpoint{6.088297in}{0.735536in}}%
\pgfpathlineto{\pgfqpoint{6.091653in}{0.709948in}}%
\pgfpathlineto{\pgfqpoint{6.095008in}{0.582671in}}%
\pgfpathlineto{\pgfqpoint{6.098364in}{0.691309in}}%
\pgfpathlineto{\pgfqpoint{6.101720in}{0.686090in}}%
\pgfpathlineto{\pgfqpoint{6.105076in}{0.632219in}}%
\pgfpathlineto{\pgfqpoint{6.108432in}{0.709362in}}%
\pgfpathlineto{\pgfqpoint{6.111787in}{0.722396in}}%
\pgfpathlineto{\pgfqpoint{6.115143in}{0.699009in}}%
\pgfpathlineto{\pgfqpoint{6.118499in}{0.655119in}}%
\pgfpathlineto{\pgfqpoint{6.121855in}{0.707766in}}%
\pgfpathlineto{\pgfqpoint{6.125210in}{0.712677in}}%
\pgfpathlineto{\pgfqpoint{6.128566in}{0.688587in}}%
\pgfpathlineto{\pgfqpoint{6.131922in}{0.675420in}}%
\pgfpathlineto{\pgfqpoint{6.135278in}{0.721795in}}%
\pgfpathlineto{\pgfqpoint{6.138634in}{0.729021in}}%
\pgfpathlineto{\pgfqpoint{6.141989in}{0.695432in}}%
\pgfpathlineto{\pgfqpoint{6.145345in}{0.700911in}}%
\pgfpathlineto{\pgfqpoint{6.148701in}{0.733500in}}%
\pgfpathlineto{\pgfqpoint{6.152057in}{0.737719in}}%
\pgfpathlineto{\pgfqpoint{6.155412in}{0.720398in}}%
\pgfpathlineto{\pgfqpoint{6.158768in}{0.655430in}}%
\pgfpathlineto{\pgfqpoint{6.162124in}{0.735659in}}%
\pgfpathlineto{\pgfqpoint{6.165480in}{0.751118in}}%
\pgfpathlineto{\pgfqpoint{6.168836in}{0.738046in}}%
\pgfpathlineto{\pgfqpoint{6.172191in}{0.516782in}}%
\pgfpathlineto{\pgfqpoint{6.175547in}{0.733122in}}%
\pgfpathlineto{\pgfqpoint{6.178903in}{0.745608in}}%
\pgfpathlineto{\pgfqpoint{6.182259in}{0.733559in}}%
\pgfpathlineto{\pgfqpoint{6.185614in}{0.598002in}}%
\pgfpathlineto{\pgfqpoint{6.188970in}{0.742918in}}%
\pgfpathlineto{\pgfqpoint{6.192326in}{0.760635in}}%
\pgfpathlineto{\pgfqpoint{6.195682in}{0.748975in}}%
\pgfpathlineto{\pgfqpoint{6.199038in}{0.699971in}}%
\pgfpathlineto{\pgfqpoint{6.202393in}{0.702311in}}%
\pgfpathlineto{\pgfqpoint{6.205749in}{0.718715in}}%
\pgfpathlineto{\pgfqpoint{6.209105in}{0.682113in}}%
\pgfpathlineto{\pgfqpoint{6.215816in}{0.744267in}}%
\pgfpathlineto{\pgfqpoint{6.219172in}{0.740111in}}%
\pgfpathlineto{\pgfqpoint{6.222528in}{0.706750in}}%
\pgfpathlineto{\pgfqpoint{6.225884in}{0.641957in}}%
\pgfpathlineto{\pgfqpoint{6.229240in}{0.688975in}}%
\pgfpathlineto{\pgfqpoint{6.232595in}{0.652613in}}%
\pgfpathlineto{\pgfqpoint{6.235951in}{0.693750in}}%
\pgfpathlineto{\pgfqpoint{6.239307in}{0.723299in}}%
\pgfpathlineto{\pgfqpoint{6.242663in}{0.717710in}}%
\pgfpathlineto{\pgfqpoint{6.246019in}{0.669202in}}%
\pgfpathlineto{\pgfqpoint{6.249374in}{0.680548in}}%
\pgfpathlineto{\pgfqpoint{6.252730in}{0.700925in}}%
\pgfpathlineto{\pgfqpoint{6.256086in}{0.686171in}}%
\pgfpathlineto{\pgfqpoint{6.259442in}{0.648299in}}%
\pgfpathlineto{\pgfqpoint{6.262797in}{0.710809in}}%
\pgfpathlineto{\pgfqpoint{6.266153in}{0.721485in}}%
\pgfpathlineto{\pgfqpoint{6.269509in}{0.693066in}}%
\pgfpathlineto{\pgfqpoint{6.272865in}{0.684796in}}%
\pgfpathlineto{\pgfqpoint{6.276221in}{0.722309in}}%
\pgfpathlineto{\pgfqpoint{6.279576in}{0.727643in}}%
\pgfpathlineto{\pgfqpoint{6.282932in}{0.711172in}}%
\pgfpathlineto{\pgfqpoint{6.286288in}{0.639972in}}%
\pgfpathlineto{\pgfqpoint{6.289644in}{0.725381in}}%
\pgfpathlineto{\pgfqpoint{6.292999in}{0.741166in}}%
\pgfpathlineto{\pgfqpoint{6.296355in}{0.727326in}}%
\pgfpathlineto{\pgfqpoint{6.299711in}{0.661340in}}%
\pgfpathlineto{\pgfqpoint{6.303067in}{0.736060in}}%
\pgfpathlineto{\pgfqpoint{6.306423in}{0.748285in}}%
\pgfpathlineto{\pgfqpoint{6.309778in}{0.740208in}}%
\pgfpathlineto{\pgfqpoint{6.313134in}{0.684707in}}%
\pgfpathlineto{\pgfqpoint{6.316490in}{0.732372in}}%
\pgfpathlineto{\pgfqpoint{6.319846in}{0.758667in}}%
\pgfpathlineto{\pgfqpoint{6.323201in}{0.756889in}}%
\pgfpathlineto{\pgfqpoint{6.326557in}{0.720361in}}%
\pgfpathlineto{\pgfqpoint{6.329913in}{0.704397in}}%
\pgfpathlineto{\pgfqpoint{6.333269in}{0.735722in}}%
\pgfpathlineto{\pgfqpoint{6.336625in}{0.728113in}}%
\pgfpathlineto{\pgfqpoint{6.339980in}{0.600994in}}%
\pgfpathlineto{\pgfqpoint{6.343336in}{0.740271in}}%
\pgfpathlineto{\pgfqpoint{6.346692in}{0.757243in}}%
\pgfpathlineto{\pgfqpoint{6.350048in}{0.743160in}}%
\pgfpathlineto{\pgfqpoint{6.353403in}{0.698460in}}%
\pgfpathlineto{\pgfqpoint{6.356759in}{0.676767in}}%
\pgfpathlineto{\pgfqpoint{6.360115in}{0.691605in}}%
\pgfpathlineto{\pgfqpoint{6.363471in}{0.622934in}}%
\pgfpathlineto{\pgfqpoint{6.366827in}{0.718060in}}%
\pgfpathlineto{\pgfqpoint{6.370182in}{0.733967in}}%
\pgfpathlineto{\pgfqpoint{6.373538in}{0.720717in}}%
\pgfpathlineto{\pgfqpoint{6.380250in}{0.688485in}}%
\pgfpathlineto{\pgfqpoint{6.383605in}{0.673850in}}%
\pgfpathlineto{\pgfqpoint{6.383605in}{0.673850in}}%
\pgfusepath{stroke}%
\end{pgfscope}%
\begin{pgfscope}%
\pgfpathrectangle{\pgfqpoint{4.621821in}{0.320625in}}{\pgfqpoint{1.845679in}{1.164375in}}%
\pgfusepath{clip}%
\pgfsetrectcap%
\pgfsetroundjoin%
\pgfsetlinewidth{0.602250pt}%
\definecolor{currentstroke}{rgb}{0.141176,0.345098,0.650980}%
\pgfsetstrokecolor{currentstroke}%
\pgfsetstrokeopacity{0.120000}%
\pgfsetdash{}{0pt}%
\pgfpathmoveto{\pgfqpoint{4.705715in}{1.477290in}}%
\pgfpathlineto{\pgfqpoint{4.712427in}{1.464886in}}%
\pgfpathlineto{\pgfqpoint{4.722494in}{1.450126in}}%
\pgfpathlineto{\pgfqpoint{4.735918in}{1.433216in}}%
\pgfpathlineto{\pgfqpoint{4.752696in}{1.414004in}}%
\pgfpathlineto{\pgfqpoint{4.772831in}{1.392624in}}%
\pgfpathlineto{\pgfqpoint{4.792966in}{1.372946in}}%
\pgfpathlineto{\pgfqpoint{4.813100in}{1.354784in}}%
\pgfpathlineto{\pgfqpoint{4.836591in}{1.335036in}}%
\pgfpathlineto{\pgfqpoint{4.866793in}{1.311192in}}%
\pgfpathlineto{\pgfqpoint{4.903707in}{1.283586in}}%
\pgfpathlineto{\pgfqpoint{4.947332in}{1.252409in}}%
\pgfpathlineto{\pgfqpoint{4.997668in}{1.217844in}}%
\pgfpathlineto{\pgfqpoint{5.051361in}{1.182319in}}%
\pgfpathlineto{\pgfqpoint{5.108409in}{1.145936in}}%
\pgfpathlineto{\pgfqpoint{5.165457in}{1.110870in}}%
\pgfpathlineto{\pgfqpoint{5.225861in}{1.075073in}}%
\pgfpathlineto{\pgfqpoint{5.289621in}{1.038593in}}%
\pgfpathlineto{\pgfqpoint{5.363448in}{0.997694in}}%
\pgfpathlineto{\pgfqpoint{5.450699in}{0.950675in}}%
\pgfpathlineto{\pgfqpoint{5.568151in}{0.888701in}}%
\pgfpathlineto{\pgfqpoint{5.722517in}{0.808354in}}%
\pgfpathlineto{\pgfqpoint{5.729228in}{0.805682in}}%
\pgfpathlineto{\pgfqpoint{5.735940in}{0.801376in}}%
\pgfpathlineto{\pgfqpoint{5.742652in}{0.798499in}}%
\pgfpathlineto{\pgfqpoint{5.749363in}{0.794500in}}%
\pgfpathlineto{\pgfqpoint{5.756075in}{0.793122in}}%
\pgfpathlineto{\pgfqpoint{5.759430in}{0.789893in}}%
\pgfpathlineto{\pgfqpoint{5.762786in}{0.787451in}}%
\pgfpathlineto{\pgfqpoint{5.769498in}{0.784908in}}%
\pgfpathlineto{\pgfqpoint{5.772854in}{0.782400in}}%
\pgfpathlineto{\pgfqpoint{5.776209in}{0.780937in}}%
\pgfpathlineto{\pgfqpoint{5.779565in}{0.781690in}}%
\pgfpathlineto{\pgfqpoint{5.782921in}{0.780440in}}%
\pgfpathlineto{\pgfqpoint{5.786277in}{0.776029in}}%
\pgfpathlineto{\pgfqpoint{5.789632in}{0.773633in}}%
\pgfpathlineto{\pgfqpoint{5.792988in}{0.772880in}}%
\pgfpathlineto{\pgfqpoint{5.796344in}{0.771009in}}%
\pgfpathlineto{\pgfqpoint{5.799700in}{0.768326in}}%
\pgfpathlineto{\pgfqpoint{5.803056in}{0.768304in}}%
\pgfpathlineto{\pgfqpoint{5.806411in}{0.769896in}}%
\pgfpathlineto{\pgfqpoint{5.809767in}{0.766330in}}%
\pgfpathlineto{\pgfqpoint{5.813123in}{0.761521in}}%
\pgfpathlineto{\pgfqpoint{5.816479in}{0.760567in}}%
\pgfpathlineto{\pgfqpoint{5.819834in}{0.760204in}}%
\pgfpathlineto{\pgfqpoint{5.826546in}{0.754612in}}%
\pgfpathlineto{\pgfqpoint{5.829902in}{0.757314in}}%
\pgfpathlineto{\pgfqpoint{5.833258in}{0.759273in}}%
\pgfpathlineto{\pgfqpoint{5.836613in}{0.752515in}}%
\pgfpathlineto{\pgfqpoint{5.839969in}{0.747604in}}%
\pgfpathlineto{\pgfqpoint{5.843325in}{0.750023in}}%
\pgfpathlineto{\pgfqpoint{5.846681in}{0.750352in}}%
\pgfpathlineto{\pgfqpoint{5.850036in}{0.744758in}}%
\pgfpathlineto{\pgfqpoint{5.853392in}{0.741367in}}%
\pgfpathlineto{\pgfqpoint{5.856748in}{0.749853in}}%
\pgfpathlineto{\pgfqpoint{5.860104in}{0.754081in}}%
\pgfpathlineto{\pgfqpoint{5.863460in}{0.742628in}}%
\pgfpathlineto{\pgfqpoint{5.866815in}{0.734264in}}%
\pgfpathlineto{\pgfqpoint{5.870171in}{0.743273in}}%
\pgfpathlineto{\pgfqpoint{5.873527in}{0.745715in}}%
\pgfpathlineto{\pgfqpoint{5.876883in}{0.735924in}}%
\pgfpathlineto{\pgfqpoint{5.880238in}{0.728347in}}%
\pgfpathlineto{\pgfqpoint{5.883594in}{0.747134in}}%
\pgfpathlineto{\pgfqpoint{5.886950in}{0.755393in}}%
\pgfpathlineto{\pgfqpoint{5.890306in}{0.740791in}}%
\pgfpathlineto{\pgfqpoint{5.893662in}{0.719790in}}%
\pgfpathlineto{\pgfqpoint{5.897017in}{0.735327in}}%
\pgfpathlineto{\pgfqpoint{5.900373in}{0.740233in}}%
\pgfpathlineto{\pgfqpoint{5.903729in}{0.726570in}}%
\pgfpathlineto{\pgfqpoint{5.907085in}{0.718503in}}%
\pgfpathlineto{\pgfqpoint{5.910441in}{0.746663in}}%
\pgfpathlineto{\pgfqpoint{5.913796in}{0.754206in}}%
\pgfpathlineto{\pgfqpoint{5.917152in}{0.735078in}}%
\pgfpathlineto{\pgfqpoint{5.920508in}{0.705905in}}%
\pgfpathlineto{\pgfqpoint{5.923864in}{0.725414in}}%
\pgfpathlineto{\pgfqpoint{5.927219in}{0.727913in}}%
\pgfpathlineto{\pgfqpoint{5.930575in}{0.706560in}}%
\pgfpathlineto{\pgfqpoint{5.933931in}{0.720347in}}%
\pgfpathlineto{\pgfqpoint{5.937287in}{0.744950in}}%
\pgfpathlineto{\pgfqpoint{5.940643in}{0.742809in}}%
\pgfpathlineto{\pgfqpoint{5.943998in}{0.712931in}}%
\pgfpathlineto{\pgfqpoint{5.947354in}{0.701646in}}%
\pgfpathlineto{\pgfqpoint{5.950710in}{0.720921in}}%
\pgfpathlineto{\pgfqpoint{5.954066in}{0.714160in}}%
\pgfpathlineto{\pgfqpoint{5.957421in}{0.687138in}}%
\pgfpathlineto{\pgfqpoint{5.960777in}{0.725650in}}%
\pgfpathlineto{\pgfqpoint{5.964133in}{0.740840in}}%
\pgfpathlineto{\pgfqpoint{5.967489in}{0.726555in}}%
\pgfpathlineto{\pgfqpoint{5.970845in}{0.680506in}}%
\pgfpathlineto{\pgfqpoint{5.974200in}{0.714311in}}%
\pgfpathlineto{\pgfqpoint{5.977556in}{0.724307in}}%
\pgfpathlineto{\pgfqpoint{5.980912in}{0.708709in}}%
\pgfpathlineto{\pgfqpoint{5.984268in}{0.683692in}}%
\pgfpathlineto{\pgfqpoint{5.987623in}{0.730350in}}%
\pgfpathlineto{\pgfqpoint{5.990979in}{0.741580in}}%
\pgfpathlineto{\pgfqpoint{5.994335in}{0.721373in}}%
\pgfpathlineto{\pgfqpoint{5.997691in}{0.676469in}}%
\pgfpathlineto{\pgfqpoint{6.001047in}{0.725804in}}%
\pgfpathlineto{\pgfqpoint{6.004402in}{0.733312in}}%
\pgfpathlineto{\pgfqpoint{6.007758in}{0.716474in}}%
\pgfpathlineto{\pgfqpoint{6.011114in}{0.676814in}}%
\pgfpathlineto{\pgfqpoint{6.014470in}{0.736010in}}%
\pgfpathlineto{\pgfqpoint{6.017825in}{0.749255in}}%
\pgfpathlineto{\pgfqpoint{6.021181in}{0.732591in}}%
\pgfpathlineto{\pgfqpoint{6.024537in}{0.653443in}}%
\pgfpathlineto{\pgfqpoint{6.027893in}{0.727364in}}%
\pgfpathlineto{\pgfqpoint{6.031249in}{0.737647in}}%
\pgfpathlineto{\pgfqpoint{6.034604in}{0.722232in}}%
\pgfpathlineto{\pgfqpoint{6.037960in}{0.670480in}}%
\pgfpathlineto{\pgfqpoint{6.041316in}{0.740189in}}%
\pgfpathlineto{\pgfqpoint{6.044672in}{0.754364in}}%
\pgfpathlineto{\pgfqpoint{6.048027in}{0.738970in}}%
\pgfpathlineto{\pgfqpoint{6.051383in}{0.663283in}}%
\pgfpathlineto{\pgfqpoint{6.054739in}{0.718311in}}%
\pgfpathlineto{\pgfqpoint{6.058095in}{0.729621in}}%
\pgfpathlineto{\pgfqpoint{6.061451in}{0.707812in}}%
\pgfpathlineto{\pgfqpoint{6.064806in}{0.696231in}}%
\pgfpathlineto{\pgfqpoint{6.068162in}{0.742379in}}%
\pgfpathlineto{\pgfqpoint{6.071518in}{0.749500in}}%
\pgfpathlineto{\pgfqpoint{6.074874in}{0.725192in}}%
\pgfpathlineto{\pgfqpoint{6.078230in}{0.635927in}}%
\pgfpathlineto{\pgfqpoint{6.081585in}{0.713630in}}%
\pgfpathlineto{\pgfqpoint{6.084941in}{0.714962in}}%
\pgfpathlineto{\pgfqpoint{6.088297in}{0.659077in}}%
\pgfpathlineto{\pgfqpoint{6.091653in}{0.715815in}}%
\pgfpathlineto{\pgfqpoint{6.095008in}{0.740400in}}%
\pgfpathlineto{\pgfqpoint{6.098364in}{0.733319in}}%
\pgfpathlineto{\pgfqpoint{6.101720in}{0.684857in}}%
\pgfpathlineto{\pgfqpoint{6.105076in}{0.697188in}}%
\pgfpathlineto{\pgfqpoint{6.108432in}{0.718154in}}%
\pgfpathlineto{\pgfqpoint{6.111787in}{0.706336in}}%
\pgfpathlineto{\pgfqpoint{6.115143in}{0.643028in}}%
\pgfpathlineto{\pgfqpoint{6.118499in}{0.723964in}}%
\pgfpathlineto{\pgfqpoint{6.121855in}{0.738684in}}%
\pgfpathlineto{\pgfqpoint{6.125210in}{0.721363in}}%
\pgfpathlineto{\pgfqpoint{6.128566in}{0.625485in}}%
\pgfpathlineto{\pgfqpoint{6.131922in}{0.718638in}}%
\pgfpathlineto{\pgfqpoint{6.135278in}{0.728575in}}%
\pgfpathlineto{\pgfqpoint{6.138634in}{0.713051in}}%
\pgfpathlineto{\pgfqpoint{6.141989in}{0.654015in}}%
\pgfpathlineto{\pgfqpoint{6.145345in}{0.730518in}}%
\pgfpathlineto{\pgfqpoint{6.148701in}{0.745062in}}%
\pgfpathlineto{\pgfqpoint{6.152057in}{0.729119in}}%
\pgfpathlineto{\pgfqpoint{6.155412in}{0.622274in}}%
\pgfpathlineto{\pgfqpoint{6.158768in}{0.726820in}}%
\pgfpathlineto{\pgfqpoint{6.162124in}{0.738038in}}%
\pgfpathlineto{\pgfqpoint{6.165480in}{0.725286in}}%
\pgfpathlineto{\pgfqpoint{6.168836in}{0.611827in}}%
\pgfpathlineto{\pgfqpoint{6.172191in}{0.735517in}}%
\pgfpathlineto{\pgfqpoint{6.175547in}{0.753327in}}%
\pgfpathlineto{\pgfqpoint{6.178903in}{0.742331in}}%
\pgfpathlineto{\pgfqpoint{6.182259in}{0.676690in}}%
\pgfpathlineto{\pgfqpoint{6.185614in}{0.719451in}}%
\pgfpathlineto{\pgfqpoint{6.188970in}{0.735503in}}%
\pgfpathlineto{\pgfqpoint{6.192326in}{0.722515in}}%
\pgfpathlineto{\pgfqpoint{6.195682in}{0.653932in}}%
\pgfpathlineto{\pgfqpoint{6.199038in}{0.739019in}}%
\pgfpathlineto{\pgfqpoint{6.202393in}{0.754095in}}%
\pgfpathlineto{\pgfqpoint{6.205749in}{0.739147in}}%
\pgfpathlineto{\pgfqpoint{6.209105in}{0.673481in}}%
\pgfpathlineto{\pgfqpoint{6.212461in}{0.709454in}}%
\pgfpathlineto{\pgfqpoint{6.215816in}{0.721590in}}%
\pgfpathlineto{\pgfqpoint{6.219172in}{0.693937in}}%
\pgfpathlineto{\pgfqpoint{6.222528in}{0.702072in}}%
\pgfpathlineto{\pgfqpoint{6.225884in}{0.740125in}}%
\pgfpathlineto{\pgfqpoint{6.229240in}{0.742793in}}%
\pgfpathlineto{\pgfqpoint{6.232595in}{0.712655in}}%
\pgfpathlineto{\pgfqpoint{6.235951in}{0.664729in}}%
\pgfpathlineto{\pgfqpoint{6.239307in}{0.712695in}}%
\pgfpathlineto{\pgfqpoint{6.242663in}{0.708927in}}%
\pgfpathlineto{\pgfqpoint{6.246019in}{0.607582in}}%
\pgfpathlineto{\pgfqpoint{6.249374in}{0.717339in}}%
\pgfpathlineto{\pgfqpoint{6.252730in}{0.737447in}}%
\pgfpathlineto{\pgfqpoint{6.256086in}{0.725890in}}%
\pgfpathlineto{\pgfqpoint{6.259442in}{0.649040in}}%
\pgfpathlineto{\pgfqpoint{6.262797in}{0.708078in}}%
\pgfpathlineto{\pgfqpoint{6.266153in}{0.722761in}}%
\pgfpathlineto{\pgfqpoint{6.269509in}{0.709845in}}%
\pgfpathlineto{\pgfqpoint{6.272865in}{0.632814in}}%
\pgfpathlineto{\pgfqpoint{6.276221in}{0.724725in}}%
\pgfpathlineto{\pgfqpoint{6.279576in}{0.740492in}}%
\pgfpathlineto{\pgfqpoint{6.282932in}{0.725200in}}%
\pgfpathlineto{\pgfqpoint{6.286288in}{0.618124in}}%
\pgfpathlineto{\pgfqpoint{6.289644in}{0.723825in}}%
\pgfpathlineto{\pgfqpoint{6.292999in}{0.735509in}}%
\pgfpathlineto{\pgfqpoint{6.296355in}{0.723974in}}%
\pgfpathlineto{\pgfqpoint{6.299711in}{0.586256in}}%
\pgfpathlineto{\pgfqpoint{6.303067in}{0.730466in}}%
\pgfpathlineto{\pgfqpoint{6.306423in}{0.750018in}}%
\pgfpathlineto{\pgfqpoint{6.309778in}{0.741040in}}%
\pgfpathlineto{\pgfqpoint{6.313134in}{0.674616in}}%
\pgfpathlineto{\pgfqpoint{6.316490in}{0.721981in}}%
\pgfpathlineto{\pgfqpoint{6.319846in}{0.739547in}}%
\pgfpathlineto{\pgfqpoint{6.323201in}{0.731040in}}%
\pgfpathlineto{\pgfqpoint{6.326557in}{0.646399in}}%
\pgfpathlineto{\pgfqpoint{6.329913in}{0.733776in}}%
\pgfpathlineto{\pgfqpoint{6.333269in}{0.755327in}}%
\pgfpathlineto{\pgfqpoint{6.336625in}{0.747800in}}%
\pgfpathlineto{\pgfqpoint{6.339980in}{0.702063in}}%
\pgfpathlineto{\pgfqpoint{6.343336in}{0.706018in}}%
\pgfpathlineto{\pgfqpoint{6.346692in}{0.728505in}}%
\pgfpathlineto{\pgfqpoint{6.350048in}{0.715339in}}%
\pgfpathlineto{\pgfqpoint{6.353403in}{0.668388in}}%
\pgfpathlineto{\pgfqpoint{6.356759in}{0.737689in}}%
\pgfpathlineto{\pgfqpoint{6.360115in}{0.750422in}}%
\pgfpathlineto{\pgfqpoint{6.363471in}{0.734366in}}%
\pgfpathlineto{\pgfqpoint{6.366827in}{0.697956in}}%
\pgfpathlineto{\pgfqpoint{6.370182in}{0.525856in}}%
\pgfpathlineto{\pgfqpoint{6.373538in}{0.621262in}}%
\pgfpathlineto{\pgfqpoint{6.376894in}{0.672253in}}%
\pgfpathlineto{\pgfqpoint{6.380250in}{0.709287in}}%
\pgfpathlineto{\pgfqpoint{6.383605in}{0.712891in}}%
\pgfpathlineto{\pgfqpoint{6.383605in}{0.712891in}}%
\pgfusepath{stroke}%
\end{pgfscope}%
\begin{pgfscope}%
\pgfpathrectangle{\pgfqpoint{4.621821in}{0.320625in}}{\pgfqpoint{1.845679in}{1.164375in}}%
\pgfusepath{clip}%
\pgfsetrectcap%
\pgfsetroundjoin%
\pgfsetlinewidth{0.602250pt}%
\definecolor{currentstroke}{rgb}{0.141176,0.345098,0.650980}%
\pgfsetstrokecolor{currentstroke}%
\pgfsetstrokeopacity{0.120000}%
\pgfsetdash{}{0pt}%
\pgfpathmoveto{\pgfqpoint{4.705715in}{1.473458in}}%
\pgfpathlineto{\pgfqpoint{4.712427in}{1.460817in}}%
\pgfpathlineto{\pgfqpoint{4.722494in}{1.444775in}}%
\pgfpathlineto{\pgfqpoint{4.739273in}{1.420598in}}%
\pgfpathlineto{\pgfqpoint{4.772831in}{1.374575in}}%
\pgfpathlineto{\pgfqpoint{4.792966in}{1.348740in}}%
\pgfpathlineto{\pgfqpoint{4.809745in}{1.328851in}}%
\pgfpathlineto{\pgfqpoint{4.826524in}{1.310490in}}%
\pgfpathlineto{\pgfqpoint{4.846658in}{1.290181in}}%
\pgfpathlineto{\pgfqpoint{4.870149in}{1.268251in}}%
\pgfpathlineto{\pgfqpoint{4.896995in}{1.244725in}}%
\pgfpathlineto{\pgfqpoint{4.933909in}{1.213999in}}%
\pgfpathlineto{\pgfqpoint{4.977534in}{1.179171in}}%
\pgfpathlineto{\pgfqpoint{5.027870in}{1.140408in}}%
\pgfpathlineto{\pgfqpoint{5.078207in}{1.103029in}}%
\pgfpathlineto{\pgfqpoint{5.125188in}{1.069460in}}%
\pgfpathlineto{\pgfqpoint{5.172169in}{1.037248in}}%
\pgfpathlineto{\pgfqpoint{5.219150in}{1.006391in}}%
\pgfpathlineto{\pgfqpoint{5.269487in}{0.974683in}}%
\pgfpathlineto{\pgfqpoint{5.326535in}{0.940104in}}%
\pgfpathlineto{\pgfqpoint{5.400362in}{0.896762in}}%
\pgfpathlineto{\pgfqpoint{5.487612in}{0.846869in}}%
\pgfpathlineto{\pgfqpoint{5.527882in}{0.823930in}}%
\pgfpathlineto{\pgfqpoint{5.578218in}{0.795651in}}%
\pgfpathlineto{\pgfqpoint{5.584930in}{0.792578in}}%
\pgfpathlineto{\pgfqpoint{5.591641in}{0.787847in}}%
\pgfpathlineto{\pgfqpoint{5.601709in}{0.782859in}}%
\pgfpathlineto{\pgfqpoint{5.608420in}{0.778500in}}%
\pgfpathlineto{\pgfqpoint{5.615132in}{0.777501in}}%
\pgfpathlineto{\pgfqpoint{5.621843in}{0.770862in}}%
\pgfpathlineto{\pgfqpoint{5.625199in}{0.769575in}}%
\pgfpathlineto{\pgfqpoint{5.628555in}{0.769278in}}%
\pgfpathlineto{\pgfqpoint{5.631911in}{0.767501in}}%
\pgfpathlineto{\pgfqpoint{5.635267in}{0.763920in}}%
\pgfpathlineto{\pgfqpoint{5.638622in}{0.761634in}}%
\pgfpathlineto{\pgfqpoint{5.645334in}{0.765961in}}%
\pgfpathlineto{\pgfqpoint{5.648690in}{0.761354in}}%
\pgfpathlineto{\pgfqpoint{5.652045in}{0.753999in}}%
\pgfpathlineto{\pgfqpoint{5.655401in}{0.754923in}}%
\pgfpathlineto{\pgfqpoint{5.658757in}{0.759277in}}%
\pgfpathlineto{\pgfqpoint{5.662113in}{0.758749in}}%
\pgfpathlineto{\pgfqpoint{5.668824in}{0.744402in}}%
\pgfpathlineto{\pgfqpoint{5.675536in}{0.763800in}}%
\pgfpathlineto{\pgfqpoint{5.678892in}{0.761507in}}%
\pgfpathlineto{\pgfqpoint{5.682247in}{0.744405in}}%
\pgfpathlineto{\pgfqpoint{5.685603in}{0.736870in}}%
\pgfpathlineto{\pgfqpoint{5.688959in}{0.749122in}}%
\pgfpathlineto{\pgfqpoint{5.692315in}{0.752562in}}%
\pgfpathlineto{\pgfqpoint{5.695671in}{0.743263in}}%
\pgfpathlineto{\pgfqpoint{5.699026in}{0.727284in}}%
\pgfpathlineto{\pgfqpoint{5.702382in}{0.747053in}}%
\pgfpathlineto{\pgfqpoint{5.705738in}{0.763168in}}%
\pgfpathlineto{\pgfqpoint{5.709094in}{0.761250in}}%
\pgfpathlineto{\pgfqpoint{5.715805in}{0.717947in}}%
\pgfpathlineto{\pgfqpoint{5.719161in}{0.732233in}}%
\pgfpathlineto{\pgfqpoint{5.722517in}{0.735639in}}%
\pgfpathlineto{\pgfqpoint{5.725873in}{0.721295in}}%
\pgfpathlineto{\pgfqpoint{5.729228in}{0.717957in}}%
\pgfpathlineto{\pgfqpoint{5.732584in}{0.745418in}}%
\pgfpathlineto{\pgfqpoint{5.735940in}{0.754945in}}%
\pgfpathlineto{\pgfqpoint{5.739296in}{0.742350in}}%
\pgfpathlineto{\pgfqpoint{5.742652in}{0.712577in}}%
\pgfpathlineto{\pgfqpoint{5.746007in}{0.708787in}}%
\pgfpathlineto{\pgfqpoint{5.749363in}{0.721875in}}%
\pgfpathlineto{\pgfqpoint{5.752719in}{0.715647in}}%
\pgfpathlineto{\pgfqpoint{5.756075in}{0.695161in}}%
\pgfpathlineto{\pgfqpoint{5.759430in}{0.723234in}}%
\pgfpathlineto{\pgfqpoint{5.762786in}{0.741415in}}%
\pgfpathlineto{\pgfqpoint{5.766142in}{0.737817in}}%
\pgfpathlineto{\pgfqpoint{5.769498in}{0.709928in}}%
\pgfpathlineto{\pgfqpoint{5.772854in}{0.694685in}}%
\pgfpathlineto{\pgfqpoint{5.776209in}{0.720679in}}%
\pgfpathlineto{\pgfqpoint{5.779565in}{0.724171in}}%
\pgfpathlineto{\pgfqpoint{5.782921in}{0.707259in}}%
\pgfpathlineto{\pgfqpoint{5.786277in}{0.685308in}}%
\pgfpathlineto{\pgfqpoint{5.789632in}{0.727391in}}%
\pgfpathlineto{\pgfqpoint{5.792988in}{0.741935in}}%
\pgfpathlineto{\pgfqpoint{5.796344in}{0.733612in}}%
\pgfpathlineto{\pgfqpoint{5.799700in}{0.687747in}}%
\pgfpathlineto{\pgfqpoint{5.806411in}{0.736535in}}%
\pgfpathlineto{\pgfqpoint{5.809767in}{0.738980in}}%
\pgfpathlineto{\pgfqpoint{5.813123in}{0.723397in}}%
\pgfpathlineto{\pgfqpoint{5.816479in}{0.661682in}}%
\pgfpathlineto{\pgfqpoint{5.819834in}{0.731540in}}%
\pgfpathlineto{\pgfqpoint{5.823190in}{0.751922in}}%
\pgfpathlineto{\pgfqpoint{5.826546in}{0.751183in}}%
\pgfpathlineto{\pgfqpoint{5.829902in}{0.720951in}}%
\pgfpathlineto{\pgfqpoint{5.833258in}{0.699614in}}%
\pgfpathlineto{\pgfqpoint{5.836613in}{0.741204in}}%
\pgfpathlineto{\pgfqpoint{5.839969in}{0.749267in}}%
\pgfpathlineto{\pgfqpoint{5.843325in}{0.739773in}}%
\pgfpathlineto{\pgfqpoint{5.846681in}{0.687718in}}%
\pgfpathlineto{\pgfqpoint{5.850036in}{0.730308in}}%
\pgfpathlineto{\pgfqpoint{5.853392in}{0.758771in}}%
\pgfpathlineto{\pgfqpoint{5.856748in}{0.763739in}}%
\pgfpathlineto{\pgfqpoint{5.860104in}{0.744871in}}%
\pgfpathlineto{\pgfqpoint{5.863460in}{0.670771in}}%
\pgfpathlineto{\pgfqpoint{5.866815in}{0.723560in}}%
\pgfpathlineto{\pgfqpoint{5.870171in}{0.738870in}}%
\pgfpathlineto{\pgfqpoint{5.873527in}{0.729092in}}%
\pgfpathlineto{\pgfqpoint{5.876883in}{0.641280in}}%
\pgfpathlineto{\pgfqpoint{5.880238in}{0.734576in}}%
\pgfpathlineto{\pgfqpoint{5.883594in}{0.757309in}}%
\pgfpathlineto{\pgfqpoint{5.886950in}{0.756049in}}%
\pgfpathlineto{\pgfqpoint{5.890306in}{0.731370in}}%
\pgfpathlineto{\pgfqpoint{5.893662in}{0.633286in}}%
\pgfpathlineto{\pgfqpoint{5.897017in}{0.709838in}}%
\pgfpathlineto{\pgfqpoint{5.900373in}{0.716702in}}%
\pgfpathlineto{\pgfqpoint{5.903729in}{0.684405in}}%
\pgfpathlineto{\pgfqpoint{5.907085in}{0.701756in}}%
\pgfpathlineto{\pgfqpoint{5.910441in}{0.738104in}}%
\pgfpathlineto{\pgfqpoint{5.913796in}{0.745330in}}%
\pgfpathlineto{\pgfqpoint{5.917152in}{0.726859in}}%
\pgfpathlineto{\pgfqpoint{5.920508in}{0.660680in}}%
\pgfpathlineto{\pgfqpoint{5.923864in}{0.699716in}}%
\pgfpathlineto{\pgfqpoint{5.927219in}{0.715794in}}%
\pgfpathlineto{\pgfqpoint{5.930575in}{0.704376in}}%
\pgfpathlineto{\pgfqpoint{5.933931in}{0.607218in}}%
\pgfpathlineto{\pgfqpoint{5.937287in}{0.716978in}}%
\pgfpathlineto{\pgfqpoint{5.940643in}{0.737020in}}%
\pgfpathlineto{\pgfqpoint{5.943998in}{0.732744in}}%
\pgfpathlineto{\pgfqpoint{5.947354in}{0.693694in}}%
\pgfpathlineto{\pgfqpoint{5.950710in}{0.694241in}}%
\pgfpathlineto{\pgfqpoint{5.954066in}{0.726570in}}%
\pgfpathlineto{\pgfqpoint{5.957421in}{0.731222in}}%
\pgfpathlineto{\pgfqpoint{5.960777in}{0.716837in}}%
\pgfpathlineto{\pgfqpoint{5.964133in}{0.583134in}}%
\pgfpathlineto{\pgfqpoint{5.967489in}{0.723189in}}%
\pgfpathlineto{\pgfqpoint{5.970845in}{0.744741in}}%
\pgfpathlineto{\pgfqpoint{5.974200in}{0.744256in}}%
\pgfpathlineto{\pgfqpoint{5.977556in}{0.711520in}}%
\pgfpathlineto{\pgfqpoint{5.980912in}{0.701449in}}%
\pgfpathlineto{\pgfqpoint{5.984268in}{0.740055in}}%
\pgfpathlineto{\pgfqpoint{5.987623in}{0.748272in}}%
\pgfpathlineto{\pgfqpoint{5.990979in}{0.740338in}}%
\pgfpathlineto{\pgfqpoint{5.994335in}{0.699111in}}%
\pgfpathlineto{\pgfqpoint{5.997691in}{0.720341in}}%
\pgfpathlineto{\pgfqpoint{6.001047in}{0.754110in}}%
\pgfpathlineto{\pgfqpoint{6.004402in}{0.762298in}}%
\pgfpathlineto{\pgfqpoint{6.007758in}{0.747481in}}%
\pgfpathlineto{\pgfqpoint{6.011114in}{0.679499in}}%
\pgfpathlineto{\pgfqpoint{6.014470in}{0.727432in}}%
\pgfpathlineto{\pgfqpoint{6.017825in}{0.746766in}}%
\pgfpathlineto{\pgfqpoint{6.021181in}{0.743644in}}%
\pgfpathlineto{\pgfqpoint{6.024537in}{0.710317in}}%
\pgfpathlineto{\pgfqpoint{6.027893in}{0.717733in}}%
\pgfpathlineto{\pgfqpoint{6.031249in}{0.756161in}}%
\pgfpathlineto{\pgfqpoint{6.034604in}{0.765298in}}%
\pgfpathlineto{\pgfqpoint{6.037960in}{0.751115in}}%
\pgfpathlineto{\pgfqpoint{6.041316in}{0.707776in}}%
\pgfpathlineto{\pgfqpoint{6.044672in}{0.698448in}}%
\pgfpathlineto{\pgfqpoint{6.048027in}{0.725196in}}%
\pgfpathlineto{\pgfqpoint{6.051383in}{0.715922in}}%
\pgfpathlineto{\pgfqpoint{6.054739in}{0.619168in}}%
\pgfpathlineto{\pgfqpoint{6.058095in}{0.731269in}}%
\pgfpathlineto{\pgfqpoint{6.061451in}{0.751096in}}%
\pgfpathlineto{\pgfqpoint{6.064806in}{0.746156in}}%
\pgfpathlineto{\pgfqpoint{6.068162in}{0.717375in}}%
\pgfpathlineto{\pgfqpoint{6.071518in}{0.628175in}}%
\pgfpathlineto{\pgfqpoint{6.074874in}{0.706248in}}%
\pgfpathlineto{\pgfqpoint{6.078230in}{0.707962in}}%
\pgfpathlineto{\pgfqpoint{6.081585in}{0.661162in}}%
\pgfpathlineto{\pgfqpoint{6.084941in}{0.703668in}}%
\pgfpathlineto{\pgfqpoint{6.088297in}{0.733596in}}%
\pgfpathlineto{\pgfqpoint{6.091653in}{0.737054in}}%
\pgfpathlineto{\pgfqpoint{6.095008in}{0.712677in}}%
\pgfpathlineto{\pgfqpoint{6.098364in}{0.634762in}}%
\pgfpathlineto{\pgfqpoint{6.101720in}{0.711983in}}%
\pgfpathlineto{\pgfqpoint{6.105076in}{0.722213in}}%
\pgfpathlineto{\pgfqpoint{6.108432in}{0.711094in}}%
\pgfpathlineto{\pgfqpoint{6.111787in}{0.621063in}}%
\pgfpathlineto{\pgfqpoint{6.115143in}{0.714049in}}%
\pgfpathlineto{\pgfqpoint{6.118499in}{0.737392in}}%
\pgfpathlineto{\pgfqpoint{6.121855in}{0.737425in}}%
\pgfpathlineto{\pgfqpoint{6.125210in}{0.704228in}}%
\pgfpathlineto{\pgfqpoint{6.128566in}{0.696524in}}%
\pgfpathlineto{\pgfqpoint{6.131922in}{0.734428in}}%
\pgfpathlineto{\pgfqpoint{6.135278in}{0.742686in}}%
\pgfpathlineto{\pgfqpoint{6.138634in}{0.735085in}}%
\pgfpathlineto{\pgfqpoint{6.141989in}{0.695923in}}%
\pgfpathlineto{\pgfqpoint{6.145345in}{0.712583in}}%
\pgfpathlineto{\pgfqpoint{6.148701in}{0.747676in}}%
\pgfpathlineto{\pgfqpoint{6.152057in}{0.756545in}}%
\pgfpathlineto{\pgfqpoint{6.155412in}{0.742230in}}%
\pgfpathlineto{\pgfqpoint{6.158768in}{0.651436in}}%
\pgfpathlineto{\pgfqpoint{6.162124in}{0.732739in}}%
\pgfpathlineto{\pgfqpoint{6.165480in}{0.751411in}}%
\pgfpathlineto{\pgfqpoint{6.168836in}{0.750712in}}%
\pgfpathlineto{\pgfqpoint{6.172191in}{0.728717in}}%
\pgfpathlineto{\pgfqpoint{6.175547in}{0.689914in}}%
\pgfpathlineto{\pgfqpoint{6.178903in}{0.751472in}}%
\pgfpathlineto{\pgfqpoint{6.182259in}{0.767176in}}%
\pgfpathlineto{\pgfqpoint{6.185614in}{0.761298in}}%
\pgfpathlineto{\pgfqpoint{6.188970in}{0.729807in}}%
\pgfpathlineto{\pgfqpoint{6.192326in}{0.684348in}}%
\pgfpathlineto{\pgfqpoint{6.195682in}{0.733468in}}%
\pgfpathlineto{\pgfqpoint{6.199038in}{0.736052in}}%
\pgfpathlineto{\pgfqpoint{6.202393in}{0.706203in}}%
\pgfpathlineto{\pgfqpoint{6.205749in}{0.711296in}}%
\pgfpathlineto{\pgfqpoint{6.209105in}{0.751354in}}%
\pgfpathlineto{\pgfqpoint{6.212461in}{0.760456in}}%
\pgfpathlineto{\pgfqpoint{6.215816in}{0.745449in}}%
\pgfpathlineto{\pgfqpoint{6.219172in}{0.704345in}}%
\pgfpathlineto{\pgfqpoint{6.222528in}{0.682153in}}%
\pgfpathlineto{\pgfqpoint{6.225884in}{0.711088in}}%
\pgfpathlineto{\pgfqpoint{6.229240in}{0.696257in}}%
\pgfpathlineto{\pgfqpoint{6.232595in}{0.671042in}}%
\pgfpathlineto{\pgfqpoint{6.235951in}{0.728489in}}%
\pgfpathlineto{\pgfqpoint{6.239307in}{0.742790in}}%
\pgfpathlineto{\pgfqpoint{6.242663in}{0.731661in}}%
\pgfpathlineto{\pgfqpoint{6.246019in}{0.689749in}}%
\pgfpathlineto{\pgfqpoint{6.249374in}{0.682199in}}%
\pgfpathlineto{\pgfqpoint{6.252730in}{0.711205in}}%
\pgfpathlineto{\pgfqpoint{6.256086in}{0.707437in}}%
\pgfpathlineto{\pgfqpoint{6.259442in}{0.655207in}}%
\pgfpathlineto{\pgfqpoint{6.262797in}{0.702402in}}%
\pgfpathlineto{\pgfqpoint{6.266153in}{0.730889in}}%
\pgfpathlineto{\pgfqpoint{6.269509in}{0.733839in}}%
\pgfpathlineto{\pgfqpoint{6.276221in}{0.674649in}}%
\pgfpathlineto{\pgfqpoint{6.279576in}{0.723410in}}%
\pgfpathlineto{\pgfqpoint{6.282932in}{0.733267in}}%
\pgfpathlineto{\pgfqpoint{6.286288in}{0.726114in}}%
\pgfpathlineto{\pgfqpoint{6.289644in}{0.686089in}}%
\pgfpathlineto{\pgfqpoint{6.292999in}{0.705508in}}%
\pgfpathlineto{\pgfqpoint{6.296355in}{0.739806in}}%
\pgfpathlineto{\pgfqpoint{6.299711in}{0.747965in}}%
\pgfpathlineto{\pgfqpoint{6.303067in}{0.731229in}}%
\pgfpathlineto{\pgfqpoint{6.306423in}{0.635649in}}%
\pgfpathlineto{\pgfqpoint{6.309778in}{0.734972in}}%
\pgfpathlineto{\pgfqpoint{6.313134in}{0.751235in}}%
\pgfpathlineto{\pgfqpoint{6.316490in}{0.750705in}}%
\pgfpathlineto{\pgfqpoint{6.319846in}{0.731693in}}%
\pgfpathlineto{\pgfqpoint{6.323201in}{0.659006in}}%
\pgfpathlineto{\pgfqpoint{6.326557in}{0.745330in}}%
\pgfpathlineto{\pgfqpoint{6.329913in}{0.764181in}}%
\pgfpathlineto{\pgfqpoint{6.333269in}{0.762668in}}%
\pgfpathlineto{\pgfqpoint{6.336625in}{0.734505in}}%
\pgfpathlineto{\pgfqpoint{6.339980in}{0.688017in}}%
\pgfpathlineto{\pgfqpoint{6.343336in}{0.742091in}}%
\pgfpathlineto{\pgfqpoint{6.346692in}{0.749795in}}%
\pgfpathlineto{\pgfqpoint{6.350048in}{0.736147in}}%
\pgfpathlineto{\pgfqpoint{6.353403in}{0.582363in}}%
\pgfpathlineto{\pgfqpoint{6.356759in}{0.745089in}}%
\pgfpathlineto{\pgfqpoint{6.360115in}{0.765842in}}%
\pgfpathlineto{\pgfqpoint{6.363471in}{0.763648in}}%
\pgfpathlineto{\pgfqpoint{6.366827in}{0.747273in}}%
\pgfpathlineto{\pgfqpoint{6.373538in}{0.710293in}}%
\pgfpathlineto{\pgfqpoint{6.376894in}{0.703472in}}%
\pgfpathlineto{\pgfqpoint{6.380250in}{0.706011in}}%
\pgfpathlineto{\pgfqpoint{6.383605in}{0.702602in}}%
\pgfpathlineto{\pgfqpoint{6.383605in}{0.702602in}}%
\pgfusepath{stroke}%
\end{pgfscope}%
\begin{pgfscope}%
\pgfpathrectangle{\pgfqpoint{4.621821in}{0.320625in}}{\pgfqpoint{1.845679in}{1.164375in}}%
\pgfusepath{clip}%
\pgfsetrectcap%
\pgfsetroundjoin%
\pgfsetlinewidth{0.602250pt}%
\definecolor{currentstroke}{rgb}{0.141176,0.345098,0.650980}%
\pgfsetstrokecolor{currentstroke}%
\pgfsetstrokeopacity{0.120000}%
\pgfsetdash{}{0pt}%
\pgfpathmoveto{\pgfqpoint{4.705715in}{1.473855in}}%
\pgfpathlineto{\pgfqpoint{4.712427in}{1.462409in}}%
\pgfpathlineto{\pgfqpoint{4.725850in}{1.443133in}}%
\pgfpathlineto{\pgfqpoint{4.749341in}{1.411474in}}%
\pgfpathlineto{\pgfqpoint{4.766120in}{1.390691in}}%
\pgfpathlineto{\pgfqpoint{4.779543in}{1.375643in}}%
\pgfpathlineto{\pgfqpoint{4.796322in}{1.358544in}}%
\pgfpathlineto{\pgfqpoint{4.816456in}{1.339850in}}%
\pgfpathlineto{\pgfqpoint{4.839947in}{1.319712in}}%
\pgfpathlineto{\pgfqpoint{4.870149in}{1.295469in}}%
\pgfpathlineto{\pgfqpoint{4.910418in}{1.264771in}}%
\pgfpathlineto{\pgfqpoint{4.967466in}{1.222837in}}%
\pgfpathlineto{\pgfqpoint{5.051361in}{1.162677in}}%
\pgfpathlineto{\pgfqpoint{5.158746in}{1.087097in}}%
\pgfpathlineto{\pgfqpoint{5.276198in}{1.005797in}}%
\pgfpathlineto{\pgfqpoint{5.397006in}{0.923533in}}%
\pgfpathlineto{\pgfqpoint{5.487612in}{0.862758in}}%
\pgfpathlineto{\pgfqpoint{5.564795in}{0.812011in}}%
\pgfpathlineto{\pgfqpoint{5.568151in}{0.810483in}}%
\pgfpathlineto{\pgfqpoint{5.578218in}{0.803097in}}%
\pgfpathlineto{\pgfqpoint{5.581574in}{0.802147in}}%
\pgfpathlineto{\pgfqpoint{5.584930in}{0.800330in}}%
\pgfpathlineto{\pgfqpoint{5.591641in}{0.793788in}}%
\pgfpathlineto{\pgfqpoint{5.594997in}{0.793816in}}%
\pgfpathlineto{\pgfqpoint{5.598353in}{0.795632in}}%
\pgfpathlineto{\pgfqpoint{5.601709in}{0.793211in}}%
\pgfpathlineto{\pgfqpoint{5.605065in}{0.786555in}}%
\pgfpathlineto{\pgfqpoint{5.608420in}{0.782842in}}%
\pgfpathlineto{\pgfqpoint{5.611776in}{0.782681in}}%
\pgfpathlineto{\pgfqpoint{5.615132in}{0.781214in}}%
\pgfpathlineto{\pgfqpoint{5.618488in}{0.776900in}}%
\pgfpathlineto{\pgfqpoint{5.621843in}{0.774789in}}%
\pgfpathlineto{\pgfqpoint{5.625199in}{0.779309in}}%
\pgfpathlineto{\pgfqpoint{5.628555in}{0.781649in}}%
\pgfpathlineto{\pgfqpoint{5.635267in}{0.766041in}}%
\pgfpathlineto{\pgfqpoint{5.638622in}{0.763270in}}%
\pgfpathlineto{\pgfqpoint{5.641978in}{0.762198in}}%
\pgfpathlineto{\pgfqpoint{5.645334in}{0.758970in}}%
\pgfpathlineto{\pgfqpoint{5.648690in}{0.757234in}}%
\pgfpathlineto{\pgfqpoint{5.652045in}{0.761309in}}%
\pgfpathlineto{\pgfqpoint{5.655401in}{0.762916in}}%
\pgfpathlineto{\pgfqpoint{5.662113in}{0.747753in}}%
\pgfpathlineto{\pgfqpoint{5.665469in}{0.746785in}}%
\pgfpathlineto{\pgfqpoint{5.668824in}{0.746402in}}%
\pgfpathlineto{\pgfqpoint{5.672180in}{0.742363in}}%
\pgfpathlineto{\pgfqpoint{5.675536in}{0.739292in}}%
\pgfpathlineto{\pgfqpoint{5.682247in}{0.748914in}}%
\pgfpathlineto{\pgfqpoint{5.685603in}{0.740851in}}%
\pgfpathlineto{\pgfqpoint{5.688959in}{0.730050in}}%
\pgfpathlineto{\pgfqpoint{5.695671in}{0.742797in}}%
\pgfpathlineto{\pgfqpoint{5.699026in}{0.740128in}}%
\pgfpathlineto{\pgfqpoint{5.702382in}{0.726800in}}%
\pgfpathlineto{\pgfqpoint{5.705738in}{0.723507in}}%
\pgfpathlineto{\pgfqpoint{5.709094in}{0.743291in}}%
\pgfpathlineto{\pgfqpoint{5.712449in}{0.751829in}}%
\pgfpathlineto{\pgfqpoint{5.715805in}{0.739800in}}%
\pgfpathlineto{\pgfqpoint{5.719161in}{0.710638in}}%
\pgfpathlineto{\pgfqpoint{5.722517in}{0.742094in}}%
\pgfpathlineto{\pgfqpoint{5.725873in}{0.756959in}}%
\pgfpathlineto{\pgfqpoint{5.729228in}{0.756750in}}%
\pgfpathlineto{\pgfqpoint{5.732584in}{0.739888in}}%
\pgfpathlineto{\pgfqpoint{5.735940in}{0.701178in}}%
\pgfpathlineto{\pgfqpoint{5.739296in}{0.750382in}}%
\pgfpathlineto{\pgfqpoint{5.742652in}{0.768742in}}%
\pgfpathlineto{\pgfqpoint{5.746007in}{0.766866in}}%
\pgfpathlineto{\pgfqpoint{5.749363in}{0.736821in}}%
\pgfpathlineto{\pgfqpoint{5.752719in}{0.715229in}}%
\pgfpathlineto{\pgfqpoint{5.756075in}{0.752901in}}%
\pgfpathlineto{\pgfqpoint{5.759430in}{0.758521in}}%
\pgfpathlineto{\pgfqpoint{5.762786in}{0.743252in}}%
\pgfpathlineto{\pgfqpoint{5.766142in}{0.684858in}}%
\pgfpathlineto{\pgfqpoint{5.769498in}{0.755727in}}%
\pgfpathlineto{\pgfqpoint{5.772854in}{0.774176in}}%
\pgfpathlineto{\pgfqpoint{5.776209in}{0.769223in}}%
\pgfpathlineto{\pgfqpoint{5.779565in}{0.741553in}}%
\pgfpathlineto{\pgfqpoint{5.782921in}{0.669775in}}%
\pgfpathlineto{\pgfqpoint{5.786277in}{0.723439in}}%
\pgfpathlineto{\pgfqpoint{5.789632in}{0.722303in}}%
\pgfpathlineto{\pgfqpoint{5.792988in}{0.663614in}}%
\pgfpathlineto{\pgfqpoint{5.796344in}{0.734068in}}%
\pgfpathlineto{\pgfqpoint{5.799700in}{0.756088in}}%
\pgfpathlineto{\pgfqpoint{5.803056in}{0.751823in}}%
\pgfpathlineto{\pgfqpoint{5.806411in}{0.724697in}}%
\pgfpathlineto{\pgfqpoint{5.809767in}{0.652269in}}%
\pgfpathlineto{\pgfqpoint{5.813123in}{0.701057in}}%
\pgfpathlineto{\pgfqpoint{5.816479in}{0.698822in}}%
\pgfpathlineto{\pgfqpoint{5.819834in}{0.648207in}}%
\pgfpathlineto{\pgfqpoint{5.823190in}{0.717898in}}%
\pgfpathlineto{\pgfqpoint{5.826546in}{0.737234in}}%
\pgfpathlineto{\pgfqpoint{5.829902in}{0.729950in}}%
\pgfpathlineto{\pgfqpoint{5.833258in}{0.687626in}}%
\pgfpathlineto{\pgfqpoint{5.836613in}{0.691780in}}%
\pgfpathlineto{\pgfqpoint{5.839969in}{0.718955in}}%
\pgfpathlineto{\pgfqpoint{5.843325in}{0.718387in}}%
\pgfpathlineto{\pgfqpoint{5.846681in}{0.688488in}}%
\pgfpathlineto{\pgfqpoint{5.850036in}{0.690731in}}%
\pgfpathlineto{\pgfqpoint{5.853392in}{0.730496in}}%
\pgfpathlineto{\pgfqpoint{5.856748in}{0.739377in}}%
\pgfpathlineto{\pgfqpoint{5.860104in}{0.719498in}}%
\pgfpathlineto{\pgfqpoint{5.863460in}{0.674755in}}%
\pgfpathlineto{\pgfqpoint{5.866815in}{0.734561in}}%
\pgfpathlineto{\pgfqpoint{5.870171in}{0.748073in}}%
\pgfpathlineto{\pgfqpoint{5.873527in}{0.745781in}}%
\pgfpathlineto{\pgfqpoint{5.876883in}{0.723349in}}%
\pgfpathlineto{\pgfqpoint{5.880238in}{0.682122in}}%
\pgfpathlineto{\pgfqpoint{5.883594in}{0.744678in}}%
\pgfpathlineto{\pgfqpoint{5.886950in}{0.760702in}}%
\pgfpathlineto{\pgfqpoint{5.890306in}{0.755742in}}%
\pgfpathlineto{\pgfqpoint{5.893662in}{0.725617in}}%
\pgfpathlineto{\pgfqpoint{5.897017in}{0.669103in}}%
\pgfpathlineto{\pgfqpoint{5.900373in}{0.724087in}}%
\pgfpathlineto{\pgfqpoint{5.903729in}{0.725614in}}%
\pgfpathlineto{\pgfqpoint{5.907085in}{0.686877in}}%
\pgfpathlineto{\pgfqpoint{5.910441in}{0.713422in}}%
\pgfpathlineto{\pgfqpoint{5.913796in}{0.746303in}}%
\pgfpathlineto{\pgfqpoint{5.917152in}{0.750848in}}%
\pgfpathlineto{\pgfqpoint{5.920508in}{0.728921in}}%
\pgfpathlineto{\pgfqpoint{5.923864in}{0.633452in}}%
\pgfpathlineto{\pgfqpoint{5.927219in}{0.708736in}}%
\pgfpathlineto{\pgfqpoint{5.930575in}{0.718985in}}%
\pgfpathlineto{\pgfqpoint{5.933931in}{0.697828in}}%
\pgfpathlineto{\pgfqpoint{5.937287in}{0.680972in}}%
\pgfpathlineto{\pgfqpoint{5.940643in}{0.731429in}}%
\pgfpathlineto{\pgfqpoint{5.943998in}{0.743173in}}%
\pgfpathlineto{\pgfqpoint{5.947354in}{0.727730in}}%
\pgfpathlineto{\pgfqpoint{5.950710in}{0.633866in}}%
\pgfpathlineto{\pgfqpoint{5.954066in}{0.716097in}}%
\pgfpathlineto{\pgfqpoint{5.957421in}{0.732641in}}%
\pgfpathlineto{\pgfqpoint{5.960777in}{0.727969in}}%
\pgfpathlineto{\pgfqpoint{5.964133in}{0.689950in}}%
\pgfpathlineto{\pgfqpoint{5.967489in}{0.707690in}}%
\pgfpathlineto{\pgfqpoint{5.970845in}{0.742358in}}%
\pgfpathlineto{\pgfqpoint{5.974200in}{0.749293in}}%
\pgfpathlineto{\pgfqpoint{5.977556in}{0.728207in}}%
\pgfpathlineto{\pgfqpoint{5.980912in}{0.671629in}}%
\pgfpathlineto{\pgfqpoint{5.984268in}{0.737547in}}%
\pgfpathlineto{\pgfqpoint{5.987623in}{0.750112in}}%
\pgfpathlineto{\pgfqpoint{5.990979in}{0.745264in}}%
\pgfpathlineto{\pgfqpoint{5.994335in}{0.712984in}}%
\pgfpathlineto{\pgfqpoint{5.997691in}{0.714125in}}%
\pgfpathlineto{\pgfqpoint{6.001047in}{0.754461in}}%
\pgfpathlineto{\pgfqpoint{6.004402in}{0.764578in}}%
\pgfpathlineto{\pgfqpoint{6.007758in}{0.750845in}}%
\pgfpathlineto{\pgfqpoint{6.011114in}{0.679186in}}%
\pgfpathlineto{\pgfqpoint{6.014470in}{0.733460in}}%
\pgfpathlineto{\pgfqpoint{6.017825in}{0.752247in}}%
\pgfpathlineto{\pgfqpoint{6.021181in}{0.748776in}}%
\pgfpathlineto{\pgfqpoint{6.024537in}{0.714095in}}%
\pgfpathlineto{\pgfqpoint{6.027893in}{0.724718in}}%
\pgfpathlineto{\pgfqpoint{6.031249in}{0.761737in}}%
\pgfpathlineto{\pgfqpoint{6.034604in}{0.769678in}}%
\pgfpathlineto{\pgfqpoint{6.037960in}{0.753126in}}%
\pgfpathlineto{\pgfqpoint{6.041316in}{0.698566in}}%
\pgfpathlineto{\pgfqpoint{6.048027in}{0.733763in}}%
\pgfpathlineto{\pgfqpoint{6.051383in}{0.719276in}}%
\pgfpathlineto{\pgfqpoint{6.054739in}{0.675187in}}%
\pgfpathlineto{\pgfqpoint{6.058095in}{0.743264in}}%
\pgfpathlineto{\pgfqpoint{6.061451in}{0.758700in}}%
\pgfpathlineto{\pgfqpoint{6.064806in}{0.748368in}}%
\pgfpathlineto{\pgfqpoint{6.071518in}{0.678687in}}%
\pgfpathlineto{\pgfqpoint{6.074874in}{0.715017in}}%
\pgfpathlineto{\pgfqpoint{6.078230in}{0.706416in}}%
\pgfpathlineto{\pgfqpoint{6.081585in}{0.633505in}}%
\pgfpathlineto{\pgfqpoint{6.084941in}{0.725492in}}%
\pgfpathlineto{\pgfqpoint{6.088297in}{0.743711in}}%
\pgfpathlineto{\pgfqpoint{6.091653in}{0.735463in}}%
\pgfpathlineto{\pgfqpoint{6.095008in}{0.692190in}}%
\pgfpathlineto{\pgfqpoint{6.098364in}{0.694331in}}%
\pgfpathlineto{\pgfqpoint{6.101720in}{0.722136in}}%
\pgfpathlineto{\pgfqpoint{6.105076in}{0.720494in}}%
\pgfpathlineto{\pgfqpoint{6.108432in}{0.684523in}}%
\pgfpathlineto{\pgfqpoint{6.111787in}{0.701024in}}%
\pgfpathlineto{\pgfqpoint{6.115143in}{0.736363in}}%
\pgfpathlineto{\pgfqpoint{6.118499in}{0.742899in}}%
\pgfpathlineto{\pgfqpoint{6.121855in}{0.719911in}}%
\pgfpathlineto{\pgfqpoint{6.125210in}{0.673999in}}%
\pgfpathlineto{\pgfqpoint{6.128566in}{0.732260in}}%
\pgfpathlineto{\pgfqpoint{6.131922in}{0.743743in}}%
\pgfpathlineto{\pgfqpoint{6.135278in}{0.737871in}}%
\pgfpathlineto{\pgfqpoint{6.138634in}{0.702011in}}%
\pgfpathlineto{\pgfqpoint{6.141989in}{0.712353in}}%
\pgfpathlineto{\pgfqpoint{6.145345in}{0.749298in}}%
\pgfpathlineto{\pgfqpoint{6.148701in}{0.758074in}}%
\pgfpathlineto{\pgfqpoint{6.152057in}{0.741702in}}%
\pgfpathlineto{\pgfqpoint{6.155412in}{0.598022in}}%
\pgfpathlineto{\pgfqpoint{6.158768in}{0.740131in}}%
\pgfpathlineto{\pgfqpoint{6.162124in}{0.755820in}}%
\pgfpathlineto{\pgfqpoint{6.165480in}{0.752866in}}%
\pgfpathlineto{\pgfqpoint{6.168836in}{0.724958in}}%
\pgfpathlineto{\pgfqpoint{6.172191in}{0.713393in}}%
\pgfpathlineto{\pgfqpoint{6.175547in}{0.759342in}}%
\pgfpathlineto{\pgfqpoint{6.178903in}{0.771131in}}%
\pgfpathlineto{\pgfqpoint{6.182259in}{0.759755in}}%
\pgfpathlineto{\pgfqpoint{6.185614in}{0.713707in}}%
\pgfpathlineto{\pgfqpoint{6.188970in}{0.719233in}}%
\pgfpathlineto{\pgfqpoint{6.192326in}{0.744483in}}%
\pgfpathlineto{\pgfqpoint{6.195682in}{0.739866in}}%
\pgfpathlineto{\pgfqpoint{6.199038in}{0.686287in}}%
\pgfpathlineto{\pgfqpoint{6.202393in}{0.734726in}}%
\pgfpathlineto{\pgfqpoint{6.205749in}{0.762513in}}%
\pgfpathlineto{\pgfqpoint{6.209105in}{0.764135in}}%
\pgfpathlineto{\pgfqpoint{6.212461in}{0.741342in}}%
\pgfpathlineto{\pgfqpoint{6.215816in}{0.666875in}}%
\pgfpathlineto{\pgfqpoint{6.219172in}{0.709297in}}%
\pgfpathlineto{\pgfqpoint{6.222528in}{0.717197in}}%
\pgfpathlineto{\pgfqpoint{6.225884in}{0.673963in}}%
\pgfpathlineto{\pgfqpoint{6.229240in}{0.714298in}}%
\pgfpathlineto{\pgfqpoint{6.232595in}{0.744824in}}%
\pgfpathlineto{\pgfqpoint{6.235951in}{0.747336in}}%
\pgfpathlineto{\pgfqpoint{6.239307in}{0.723029in}}%
\pgfpathlineto{\pgfqpoint{6.242663in}{0.554900in}}%
\pgfpathlineto{\pgfqpoint{6.246019in}{0.707314in}}%
\pgfpathlineto{\pgfqpoint{6.249374in}{0.714563in}}%
\pgfpathlineto{\pgfqpoint{6.252730in}{0.687588in}}%
\pgfpathlineto{\pgfqpoint{6.256086in}{0.689292in}}%
\pgfpathlineto{\pgfqpoint{6.259442in}{0.730926in}}%
\pgfpathlineto{\pgfqpoint{6.262797in}{0.739669in}}%
\pgfpathlineto{\pgfqpoint{6.266153in}{0.719689in}}%
\pgfpathlineto{\pgfqpoint{6.269509in}{0.634234in}}%
\pgfpathlineto{\pgfqpoint{6.272865in}{0.721060in}}%
\pgfpathlineto{\pgfqpoint{6.276221in}{0.733971in}}%
\pgfpathlineto{\pgfqpoint{6.279576in}{0.727695in}}%
\pgfpathlineto{\pgfqpoint{6.282932in}{0.686461in}}%
\pgfpathlineto{\pgfqpoint{6.286288in}{0.709718in}}%
\pgfpathlineto{\pgfqpoint{6.289644in}{0.742744in}}%
\pgfpathlineto{\pgfqpoint{6.292999in}{0.749200in}}%
\pgfpathlineto{\pgfqpoint{6.296355in}{0.727214in}}%
\pgfpathlineto{\pgfqpoint{6.299711in}{0.684145in}}%
\pgfpathlineto{\pgfqpoint{6.303067in}{0.741897in}}%
\pgfpathlineto{\pgfqpoint{6.306423in}{0.754235in}}%
\pgfpathlineto{\pgfqpoint{6.309778in}{0.750157in}}%
\pgfpathlineto{\pgfqpoint{6.313134in}{0.721931in}}%
\pgfpathlineto{\pgfqpoint{6.316490in}{0.709224in}}%
\pgfpathlineto{\pgfqpoint{6.319846in}{0.755597in}}%
\pgfpathlineto{\pgfqpoint{6.323201in}{0.767921in}}%
\pgfpathlineto{\pgfqpoint{6.326557in}{0.757758in}}%
\pgfpathlineto{\pgfqpoint{6.329913in}{0.705298in}}%
\pgfpathlineto{\pgfqpoint{6.336625in}{0.754012in}}%
\pgfpathlineto{\pgfqpoint{6.339980in}{0.753067in}}%
\pgfpathlineto{\pgfqpoint{6.343336in}{0.724961in}}%
\pgfpathlineto{\pgfqpoint{6.346692in}{0.718268in}}%
\pgfpathlineto{\pgfqpoint{6.350048in}{0.762260in}}%
\pgfpathlineto{\pgfqpoint{6.353403in}{0.772804in}}%
\pgfpathlineto{\pgfqpoint{6.356759in}{0.759314in}}%
\pgfpathlineto{\pgfqpoint{6.360115in}{0.717474in}}%
\pgfpathlineto{\pgfqpoint{6.363471in}{0.701531in}}%
\pgfpathlineto{\pgfqpoint{6.366827in}{0.729069in}}%
\pgfpathlineto{\pgfqpoint{6.370182in}{0.716030in}}%
\pgfpathlineto{\pgfqpoint{6.373538in}{0.676967in}}%
\pgfpathlineto{\pgfqpoint{6.376894in}{0.742609in}}%
\pgfpathlineto{\pgfqpoint{6.380250in}{0.757724in}}%
\pgfpathlineto{\pgfqpoint{6.383605in}{0.748501in}}%
\pgfpathlineto{\pgfqpoint{6.383605in}{0.748501in}}%
\pgfusepath{stroke}%
\end{pgfscope}%
\begin{pgfscope}%
\pgfpathrectangle{\pgfqpoint{4.621821in}{0.320625in}}{\pgfqpoint{1.845679in}{1.164375in}}%
\pgfusepath{clip}%
\pgfsetrectcap%
\pgfsetroundjoin%
\pgfsetlinewidth{1.154312pt}%
\definecolor{currentstroke}{rgb}{0.141176,0.345098,0.650980}%
\pgfsetstrokecolor{currentstroke}%
\pgfsetdash{}{0pt}%
\pgfpathmoveto{\pgfqpoint{4.705715in}{1.476601in}}%
\pgfpathlineto{\pgfqpoint{4.712427in}{1.464541in}}%
\pgfpathlineto{\pgfqpoint{4.722494in}{1.449699in}}%
\pgfpathlineto{\pgfqpoint{4.739273in}{1.427338in}}%
\pgfpathlineto{\pgfqpoint{4.759408in}{1.402349in}}%
\pgfpathlineto{\pgfqpoint{4.776187in}{1.383316in}}%
\pgfpathlineto{\pgfqpoint{4.792966in}{1.366051in}}%
\pgfpathlineto{\pgfqpoint{4.809745in}{1.350323in}}%
\pgfpathlineto{\pgfqpoint{4.829879in}{1.332966in}}%
\pgfpathlineto{\pgfqpoint{4.856726in}{1.311514in}}%
\pgfpathlineto{\pgfqpoint{4.890283in}{1.286286in}}%
\pgfpathlineto{\pgfqpoint{4.933909in}{1.254938in}}%
\pgfpathlineto{\pgfqpoint{4.990957in}{1.215389in}}%
\pgfpathlineto{\pgfqpoint{5.054717in}{1.172556in}}%
\pgfpathlineto{\pgfqpoint{5.121832in}{1.128781in}}%
\pgfpathlineto{\pgfqpoint{5.192304in}{1.084148in}}%
\pgfpathlineto{\pgfqpoint{5.266131in}{1.038756in}}%
\pgfpathlineto{\pgfqpoint{5.339958in}{0.994661in}}%
\pgfpathlineto{\pgfqpoint{5.420497in}{0.947892in}}%
\pgfpathlineto{\pgfqpoint{5.507747in}{0.898588in}}%
\pgfpathlineto{\pgfqpoint{5.594997in}{0.850483in}}%
\pgfpathlineto{\pgfqpoint{5.638622in}{0.826785in}}%
\pgfpathlineto{\pgfqpoint{5.702382in}{0.793017in}}%
\pgfpathlineto{\pgfqpoint{5.712449in}{0.787962in}}%
\pgfpathlineto{\pgfqpoint{5.719161in}{0.783819in}}%
\pgfpathlineto{\pgfqpoint{5.746007in}{0.772593in}}%
\pgfpathlineto{\pgfqpoint{5.749363in}{0.769385in}}%
\pgfpathlineto{\pgfqpoint{5.756075in}{0.768985in}}%
\pgfpathlineto{\pgfqpoint{5.762786in}{0.763691in}}%
\pgfpathlineto{\pgfqpoint{5.766142in}{0.762592in}}%
\pgfpathlineto{\pgfqpoint{5.769498in}{0.762522in}}%
\pgfpathlineto{\pgfqpoint{5.772854in}{0.763689in}}%
\pgfpathlineto{\pgfqpoint{5.776209in}{0.761622in}}%
\pgfpathlineto{\pgfqpoint{5.786277in}{0.749548in}}%
\pgfpathlineto{\pgfqpoint{5.789632in}{0.748532in}}%
\pgfpathlineto{\pgfqpoint{5.792988in}{0.749639in}}%
\pgfpathlineto{\pgfqpoint{5.796344in}{0.750042in}}%
\pgfpathlineto{\pgfqpoint{5.799700in}{0.749655in}}%
\pgfpathlineto{\pgfqpoint{5.806411in}{0.745668in}}%
\pgfpathlineto{\pgfqpoint{5.809767in}{0.741952in}}%
\pgfpathlineto{\pgfqpoint{5.813123in}{0.737175in}}%
\pgfpathlineto{\pgfqpoint{5.816479in}{0.734298in}}%
\pgfpathlineto{\pgfqpoint{5.819834in}{0.736290in}}%
\pgfpathlineto{\pgfqpoint{5.823190in}{0.741503in}}%
\pgfpathlineto{\pgfqpoint{5.826546in}{0.742591in}}%
\pgfpathlineto{\pgfqpoint{5.833258in}{0.731867in}}%
\pgfpathlineto{\pgfqpoint{5.836613in}{0.733458in}}%
\pgfpathlineto{\pgfqpoint{5.839969in}{0.738124in}}%
\pgfpathlineto{\pgfqpoint{5.843325in}{0.736505in}}%
\pgfpathlineto{\pgfqpoint{5.846681in}{0.726237in}}%
\pgfpathlineto{\pgfqpoint{5.850036in}{0.725060in}}%
\pgfpathlineto{\pgfqpoint{5.853392in}{0.742164in}}%
\pgfpathlineto{\pgfqpoint{5.856748in}{0.749646in}}%
\pgfpathlineto{\pgfqpoint{5.860104in}{0.739768in}}%
\pgfpathlineto{\pgfqpoint{5.863460in}{0.717447in}}%
\pgfpathlineto{\pgfqpoint{5.866815in}{0.725340in}}%
\pgfpathlineto{\pgfqpoint{5.870171in}{0.736486in}}%
\pgfpathlineto{\pgfqpoint{5.873527in}{0.732979in}}%
\pgfpathlineto{\pgfqpoint{5.876883in}{0.716471in}}%
\pgfpathlineto{\pgfqpoint{5.880238in}{0.725569in}}%
\pgfpathlineto{\pgfqpoint{5.883594in}{0.744643in}}%
\pgfpathlineto{\pgfqpoint{5.886950in}{0.749355in}}%
\pgfpathlineto{\pgfqpoint{5.890306in}{0.736286in}}%
\pgfpathlineto{\pgfqpoint{5.893662in}{0.709976in}}%
\pgfpathlineto{\pgfqpoint{5.900373in}{0.728471in}}%
\pgfpathlineto{\pgfqpoint{5.903729in}{0.722660in}}%
\pgfpathlineto{\pgfqpoint{5.907085in}{0.708840in}}%
\pgfpathlineto{\pgfqpoint{5.910441in}{0.728324in}}%
\pgfpathlineto{\pgfqpoint{5.913796in}{0.742760in}}%
\pgfpathlineto{\pgfqpoint{5.917152in}{0.740168in}}%
\pgfpathlineto{\pgfqpoint{5.920508in}{0.724591in}}%
\pgfpathlineto{\pgfqpoint{5.923864in}{0.705883in}}%
\pgfpathlineto{\pgfqpoint{5.930575in}{0.719547in}}%
\pgfpathlineto{\pgfqpoint{5.933931in}{0.717959in}}%
\pgfpathlineto{\pgfqpoint{5.937287in}{0.722191in}}%
\pgfpathlineto{\pgfqpoint{5.940643in}{0.730530in}}%
\pgfpathlineto{\pgfqpoint{5.943998in}{0.734310in}}%
\pgfpathlineto{\pgfqpoint{5.947354in}{0.732278in}}%
\pgfpathlineto{\pgfqpoint{5.950710in}{0.723434in}}%
\pgfpathlineto{\pgfqpoint{5.954066in}{0.718246in}}%
\pgfpathlineto{\pgfqpoint{5.957421in}{0.728056in}}%
\pgfpathlineto{\pgfqpoint{5.960777in}{0.731780in}}%
\pgfpathlineto{\pgfqpoint{5.967489in}{0.711443in}}%
\pgfpathlineto{\pgfqpoint{5.970845in}{0.728479in}}%
\pgfpathlineto{\pgfqpoint{5.974200in}{0.734817in}}%
\pgfpathlineto{\pgfqpoint{5.980912in}{0.707066in}}%
\pgfpathlineto{\pgfqpoint{5.984268in}{0.728970in}}%
\pgfpathlineto{\pgfqpoint{5.987623in}{0.736982in}}%
\pgfpathlineto{\pgfqpoint{5.990979in}{0.733268in}}%
\pgfpathlineto{\pgfqpoint{5.994335in}{0.705065in}}%
\pgfpathlineto{\pgfqpoint{5.997691in}{0.703827in}}%
\pgfpathlineto{\pgfqpoint{6.001047in}{0.739639in}}%
\pgfpathlineto{\pgfqpoint{6.004402in}{0.748920in}}%
\pgfpathlineto{\pgfqpoint{6.007758in}{0.735920in}}%
\pgfpathlineto{\pgfqpoint{6.011114in}{0.688035in}}%
\pgfpathlineto{\pgfqpoint{6.014470in}{0.723533in}}%
\pgfpathlineto{\pgfqpoint{6.017825in}{0.740900in}}%
\pgfpathlineto{\pgfqpoint{6.021181in}{0.735046in}}%
\pgfpathlineto{\pgfqpoint{6.024537in}{0.700618in}}%
\pgfpathlineto{\pgfqpoint{6.027893in}{0.714393in}}%
\pgfpathlineto{\pgfqpoint{6.031249in}{0.745920in}}%
\pgfpathlineto{\pgfqpoint{6.034604in}{0.753018in}}%
\pgfpathlineto{\pgfqpoint{6.037960in}{0.736990in}}%
\pgfpathlineto{\pgfqpoint{6.041316in}{0.715120in}}%
\pgfpathlineto{\pgfqpoint{6.044672in}{0.731899in}}%
\pgfpathlineto{\pgfqpoint{6.048027in}{0.733741in}}%
\pgfpathlineto{\pgfqpoint{6.051383in}{0.721143in}}%
\pgfpathlineto{\pgfqpoint{6.054739in}{0.703773in}}%
\pgfpathlineto{\pgfqpoint{6.058095in}{0.729401in}}%
\pgfpathlineto{\pgfqpoint{6.061451in}{0.744820in}}%
\pgfpathlineto{\pgfqpoint{6.064806in}{0.741011in}}%
\pgfpathlineto{\pgfqpoint{6.068162in}{0.731311in}}%
\pgfpathlineto{\pgfqpoint{6.071518in}{0.726186in}}%
\pgfpathlineto{\pgfqpoint{6.074874in}{0.712761in}}%
\pgfpathlineto{\pgfqpoint{6.078230in}{0.694828in}}%
\pgfpathlineto{\pgfqpoint{6.081585in}{0.693329in}}%
\pgfpathlineto{\pgfqpoint{6.084941in}{0.717529in}}%
\pgfpathlineto{\pgfqpoint{6.088297in}{0.728694in}}%
\pgfpathlineto{\pgfqpoint{6.095008in}{0.720803in}}%
\pgfpathlineto{\pgfqpoint{6.098364in}{0.715780in}}%
\pgfpathlineto{\pgfqpoint{6.101720in}{0.707352in}}%
\pgfpathlineto{\pgfqpoint{6.108432in}{0.705787in}}%
\pgfpathlineto{\pgfqpoint{6.111787in}{0.704386in}}%
\pgfpathlineto{\pgfqpoint{6.115143in}{0.713460in}}%
\pgfpathlineto{\pgfqpoint{6.118499in}{0.726888in}}%
\pgfpathlineto{\pgfqpoint{6.121855in}{0.726756in}}%
\pgfpathlineto{\pgfqpoint{6.125210in}{0.710915in}}%
\pgfpathlineto{\pgfqpoint{6.128566in}{0.709301in}}%
\pgfpathlineto{\pgfqpoint{6.131922in}{0.725623in}}%
\pgfpathlineto{\pgfqpoint{6.135278in}{0.730542in}}%
\pgfpathlineto{\pgfqpoint{6.138634in}{0.721401in}}%
\pgfpathlineto{\pgfqpoint{6.141989in}{0.699886in}}%
\pgfpathlineto{\pgfqpoint{6.145345in}{0.728134in}}%
\pgfpathlineto{\pgfqpoint{6.148701in}{0.743889in}}%
\pgfpathlineto{\pgfqpoint{6.152057in}{0.740466in}}%
\pgfpathlineto{\pgfqpoint{6.155412in}{0.721670in}}%
\pgfpathlineto{\pgfqpoint{6.158768in}{0.721099in}}%
\pgfpathlineto{\pgfqpoint{6.162124in}{0.738928in}}%
\pgfpathlineto{\pgfqpoint{6.165480in}{0.744015in}}%
\pgfpathlineto{\pgfqpoint{6.172191in}{0.724713in}}%
\pgfpathlineto{\pgfqpoint{6.175547in}{0.743485in}}%
\pgfpathlineto{\pgfqpoint{6.178903in}{0.752328in}}%
\pgfpathlineto{\pgfqpoint{6.182259in}{0.748749in}}%
\pgfpathlineto{\pgfqpoint{6.185614in}{0.737650in}}%
\pgfpathlineto{\pgfqpoint{6.188970in}{0.736421in}}%
\pgfpathlineto{\pgfqpoint{6.192326in}{0.741697in}}%
\pgfpathlineto{\pgfqpoint{6.199038in}{0.723904in}}%
\pgfpathlineto{\pgfqpoint{6.202393in}{0.735360in}}%
\pgfpathlineto{\pgfqpoint{6.205749in}{0.742324in}}%
\pgfpathlineto{\pgfqpoint{6.209105in}{0.743068in}}%
\pgfpathlineto{\pgfqpoint{6.215816in}{0.733764in}}%
\pgfpathlineto{\pgfqpoint{6.219172in}{0.731710in}}%
\pgfpathlineto{\pgfqpoint{6.225884in}{0.720516in}}%
\pgfpathlineto{\pgfqpoint{6.229240in}{0.718478in}}%
\pgfpathlineto{\pgfqpoint{6.232595in}{0.719565in}}%
\pgfpathlineto{\pgfqpoint{6.235951in}{0.725759in}}%
\pgfpathlineto{\pgfqpoint{6.239307in}{0.725212in}}%
\pgfpathlineto{\pgfqpoint{6.242663in}{0.712917in}}%
\pgfpathlineto{\pgfqpoint{6.246019in}{0.705110in}}%
\pgfpathlineto{\pgfqpoint{6.249374in}{0.718660in}}%
\pgfpathlineto{\pgfqpoint{6.252730in}{0.720620in}}%
\pgfpathlineto{\pgfqpoint{6.256086in}{0.706922in}}%
\pgfpathlineto{\pgfqpoint{6.259442in}{0.701627in}}%
\pgfpathlineto{\pgfqpoint{6.262797in}{0.718117in}}%
\pgfpathlineto{\pgfqpoint{6.266153in}{0.720352in}}%
\pgfpathlineto{\pgfqpoint{6.269509in}{0.709998in}}%
\pgfpathlineto{\pgfqpoint{6.272865in}{0.703397in}}%
\pgfpathlineto{\pgfqpoint{6.276221in}{0.721835in}}%
\pgfpathlineto{\pgfqpoint{6.279576in}{0.728341in}}%
\pgfpathlineto{\pgfqpoint{6.282932in}{0.716858in}}%
\pgfpathlineto{\pgfqpoint{6.286288in}{0.702964in}}%
\pgfpathlineto{\pgfqpoint{6.289644in}{0.724360in}}%
\pgfpathlineto{\pgfqpoint{6.292999in}{0.734502in}}%
\pgfpathlineto{\pgfqpoint{6.299711in}{0.718740in}}%
\pgfpathlineto{\pgfqpoint{6.303067in}{0.732131in}}%
\pgfpathlineto{\pgfqpoint{6.306423in}{0.742384in}}%
\pgfpathlineto{\pgfqpoint{6.309778in}{0.738826in}}%
\pgfpathlineto{\pgfqpoint{6.313134in}{0.725347in}}%
\pgfpathlineto{\pgfqpoint{6.316490in}{0.731167in}}%
\pgfpathlineto{\pgfqpoint{6.319846in}{0.746156in}}%
\pgfpathlineto{\pgfqpoint{6.323201in}{0.748570in}}%
\pgfpathlineto{\pgfqpoint{6.326557in}{0.737809in}}%
\pgfpathlineto{\pgfqpoint{6.329913in}{0.738700in}}%
\pgfpathlineto{\pgfqpoint{6.333269in}{0.746722in}}%
\pgfpathlineto{\pgfqpoint{6.336625in}{0.741705in}}%
\pgfpathlineto{\pgfqpoint{6.339980in}{0.730125in}}%
\pgfpathlineto{\pgfqpoint{6.343336in}{0.730441in}}%
\pgfpathlineto{\pgfqpoint{6.346692in}{0.740233in}}%
\pgfpathlineto{\pgfqpoint{6.353403in}{0.745942in}}%
\pgfpathlineto{\pgfqpoint{6.356759in}{0.743315in}}%
\pgfpathlineto{\pgfqpoint{6.360115in}{0.744756in}}%
\pgfpathlineto{\pgfqpoint{6.366827in}{0.730875in}}%
\pgfpathlineto{\pgfqpoint{6.370182in}{0.732504in}}%
\pgfpathlineto{\pgfqpoint{6.373538in}{0.726870in}}%
\pgfpathlineto{\pgfqpoint{6.376894in}{0.726035in}}%
\pgfpathlineto{\pgfqpoint{6.380250in}{0.732956in}}%
\pgfpathlineto{\pgfqpoint{6.383605in}{0.724494in}}%
\pgfpathlineto{\pgfqpoint{6.383605in}{0.724494in}}%
\pgfusepath{stroke}%
\end{pgfscope}%
\begin{pgfscope}%
\pgfpathrectangle{\pgfqpoint{4.621821in}{0.320625in}}{\pgfqpoint{1.845679in}{1.164375in}}%
\pgfusepath{clip}%
\pgfsetrectcap%
\pgfsetroundjoin%
\pgfsetlinewidth{0.602250pt}%
\definecolor{currentstroke}{rgb}{0.819608,0.286275,0.356863}%
\pgfsetstrokecolor{currentstroke}%
\pgfsetstrokeopacity{0.120000}%
\pgfsetdash{}{0pt}%
\pgfpathmoveto{\pgfqpoint{4.705715in}{1.471660in}}%
\pgfpathlineto{\pgfqpoint{4.729206in}{1.426942in}}%
\pgfpathlineto{\pgfqpoint{4.766120in}{1.351819in}}%
\pgfpathlineto{\pgfqpoint{4.782898in}{1.321199in}}%
\pgfpathlineto{\pgfqpoint{4.799677in}{1.292882in}}%
\pgfpathlineto{\pgfqpoint{4.819812in}{1.261068in}}%
\pgfpathlineto{\pgfqpoint{4.846658in}{1.221020in}}%
\pgfpathlineto{\pgfqpoint{4.876860in}{1.177984in}}%
\pgfpathlineto{\pgfqpoint{4.913774in}{1.127288in}}%
\pgfpathlineto{\pgfqpoint{4.954043in}{1.073792in}}%
\pgfpathlineto{\pgfqpoint{4.997668in}{1.017664in}}%
\pgfpathlineto{\pgfqpoint{5.044649in}{0.959093in}}%
\pgfpathlineto{\pgfqpoint{5.094986in}{0.898148in}}%
\pgfpathlineto{\pgfqpoint{5.145323in}{0.838542in}}%
\pgfpathlineto{\pgfqpoint{5.219150in}{0.752874in}}%
\pgfpathlineto{\pgfqpoint{5.235929in}{0.732196in}}%
\pgfpathlineto{\pgfqpoint{5.245996in}{0.722004in}}%
\pgfpathlineto{\pgfqpoint{5.266131in}{0.704358in}}%
\pgfpathlineto{\pgfqpoint{5.272842in}{0.695393in}}%
\pgfpathlineto{\pgfqpoint{5.282910in}{0.679176in}}%
\pgfpathlineto{\pgfqpoint{5.296333in}{0.655961in}}%
\pgfpathlineto{\pgfqpoint{5.303044in}{0.646529in}}%
\pgfpathlineto{\pgfqpoint{5.306400in}{0.642989in}}%
\pgfpathlineto{\pgfqpoint{5.309756in}{0.640592in}}%
\pgfpathlineto{\pgfqpoint{5.313112in}{0.639697in}}%
\pgfpathlineto{\pgfqpoint{5.316467in}{0.640500in}}%
\pgfpathlineto{\pgfqpoint{5.319823in}{0.642785in}}%
\pgfpathlineto{\pgfqpoint{5.329891in}{0.652200in}}%
\pgfpathlineto{\pgfqpoint{5.333246in}{0.654199in}}%
\pgfpathlineto{\pgfqpoint{5.336602in}{0.654857in}}%
\pgfpathlineto{\pgfqpoint{5.339958in}{0.653642in}}%
\pgfpathlineto{\pgfqpoint{5.343314in}{0.649668in}}%
\pgfpathlineto{\pgfqpoint{5.346669in}{0.651547in}}%
\pgfpathlineto{\pgfqpoint{5.353381in}{0.691051in}}%
\pgfpathlineto{\pgfqpoint{5.356737in}{0.703662in}}%
\pgfpathlineto{\pgfqpoint{5.360093in}{0.711768in}}%
\pgfpathlineto{\pgfqpoint{5.363448in}{0.716379in}}%
\pgfpathlineto{\pgfqpoint{5.366804in}{0.717947in}}%
\pgfpathlineto{\pgfqpoint{5.370160in}{0.716475in}}%
\pgfpathlineto{\pgfqpoint{5.373516in}{0.711434in}}%
\pgfpathlineto{\pgfqpoint{5.376871in}{0.701250in}}%
\pgfpathlineto{\pgfqpoint{5.380227in}{0.680858in}}%
\pgfpathlineto{\pgfqpoint{5.383583in}{0.614338in}}%
\pgfpathlineto{\pgfqpoint{5.386939in}{0.674250in}}%
\pgfpathlineto{\pgfqpoint{5.390295in}{0.702734in}}%
\pgfpathlineto{\pgfqpoint{5.393650in}{0.717543in}}%
\pgfpathlineto{\pgfqpoint{5.397006in}{0.725953in}}%
\pgfpathlineto{\pgfqpoint{5.400362in}{0.729614in}}%
\pgfpathlineto{\pgfqpoint{5.403718in}{0.728250in}}%
\pgfpathlineto{\pgfqpoint{5.407074in}{0.719173in}}%
\pgfpathlineto{\pgfqpoint{5.410429in}{0.696281in}}%
\pgfpathlineto{\pgfqpoint{5.413785in}{0.642767in}}%
\pgfpathlineto{\pgfqpoint{5.417141in}{0.679860in}}%
\pgfpathlineto{\pgfqpoint{5.420497in}{0.694773in}}%
\pgfpathlineto{\pgfqpoint{5.423852in}{0.691082in}}%
\pgfpathlineto{\pgfqpoint{5.427208in}{0.661663in}}%
\pgfpathlineto{\pgfqpoint{5.430564in}{0.657645in}}%
\pgfpathlineto{\pgfqpoint{5.433920in}{0.701604in}}%
\pgfpathlineto{\pgfqpoint{5.437276in}{0.716304in}}%
\pgfpathlineto{\pgfqpoint{5.440631in}{0.714428in}}%
\pgfpathlineto{\pgfqpoint{5.443987in}{0.694457in}}%
\pgfpathlineto{\pgfqpoint{5.447343in}{0.640327in}}%
\pgfpathlineto{\pgfqpoint{5.450699in}{0.669132in}}%
\pgfpathlineto{\pgfqpoint{5.454054in}{0.689096in}}%
\pgfpathlineto{\pgfqpoint{5.457410in}{0.688814in}}%
\pgfpathlineto{\pgfqpoint{5.460766in}{0.667840in}}%
\pgfpathlineto{\pgfqpoint{5.464122in}{0.624721in}}%
\pgfpathlineto{\pgfqpoint{5.467478in}{0.690210in}}%
\pgfpathlineto{\pgfqpoint{5.470833in}{0.709008in}}%
\pgfpathlineto{\pgfqpoint{5.474189in}{0.711075in}}%
\pgfpathlineto{\pgfqpoint{5.477545in}{0.688449in}}%
\pgfpathlineto{\pgfqpoint{5.480901in}{0.657259in}}%
\pgfpathlineto{\pgfqpoint{5.484256in}{0.711172in}}%
\pgfpathlineto{\pgfqpoint{5.487612in}{0.728311in}}%
\pgfpathlineto{\pgfqpoint{5.490968in}{0.733750in}}%
\pgfpathlineto{\pgfqpoint{5.494324in}{0.730440in}}%
\pgfpathlineto{\pgfqpoint{5.497680in}{0.715469in}}%
\pgfpathlineto{\pgfqpoint{5.501035in}{0.659463in}}%
\pgfpathlineto{\pgfqpoint{5.504391in}{0.702226in}}%
\pgfpathlineto{\pgfqpoint{5.507747in}{0.732617in}}%
\pgfpathlineto{\pgfqpoint{5.511103in}{0.744103in}}%
\pgfpathlineto{\pgfqpoint{5.514458in}{0.742827in}}%
\pgfpathlineto{\pgfqpoint{5.517814in}{0.725441in}}%
\pgfpathlineto{\pgfqpoint{5.521170in}{0.680751in}}%
\pgfpathlineto{\pgfqpoint{5.524526in}{0.683820in}}%
\pgfpathlineto{\pgfqpoint{5.527882in}{0.710337in}}%
\pgfpathlineto{\pgfqpoint{5.531237in}{0.711526in}}%
\pgfpathlineto{\pgfqpoint{5.534593in}{0.690442in}}%
\pgfpathlineto{\pgfqpoint{5.537949in}{0.653502in}}%
\pgfpathlineto{\pgfqpoint{5.541305in}{0.715508in}}%
\pgfpathlineto{\pgfqpoint{5.544660in}{0.733618in}}%
\pgfpathlineto{\pgfqpoint{5.548016in}{0.734414in}}%
\pgfpathlineto{\pgfqpoint{5.551372in}{0.716398in}}%
\pgfpathlineto{\pgfqpoint{5.554728in}{0.668015in}}%
\pgfpathlineto{\pgfqpoint{5.558084in}{0.679159in}}%
\pgfpathlineto{\pgfqpoint{5.561439in}{0.703911in}}%
\pgfpathlineto{\pgfqpoint{5.564795in}{0.704477in}}%
\pgfpathlineto{\pgfqpoint{5.568151in}{0.682606in}}%
\pgfpathlineto{\pgfqpoint{5.571507in}{0.649082in}}%
\pgfpathlineto{\pgfqpoint{5.574863in}{0.708878in}}%
\pgfpathlineto{\pgfqpoint{5.578218in}{0.726552in}}%
\pgfpathlineto{\pgfqpoint{5.581574in}{0.726863in}}%
\pgfpathlineto{\pgfqpoint{5.584930in}{0.705273in}}%
\pgfpathlineto{\pgfqpoint{5.588286in}{0.618063in}}%
\pgfpathlineto{\pgfqpoint{5.591641in}{0.703581in}}%
\pgfpathlineto{\pgfqpoint{5.594997in}{0.721055in}}%
\pgfpathlineto{\pgfqpoint{5.598353in}{0.723058in}}%
\pgfpathlineto{\pgfqpoint{5.601709in}{0.711421in}}%
\pgfpathlineto{\pgfqpoint{5.605065in}{0.662163in}}%
\pgfpathlineto{\pgfqpoint{5.611776in}{0.728222in}}%
\pgfpathlineto{\pgfqpoint{5.615132in}{0.738747in}}%
\pgfpathlineto{\pgfqpoint{5.618488in}{0.732822in}}%
\pgfpathlineto{\pgfqpoint{5.621843in}{0.699161in}}%
\pgfpathlineto{\pgfqpoint{5.625199in}{0.688656in}}%
\pgfpathlineto{\pgfqpoint{5.628555in}{0.728868in}}%
\pgfpathlineto{\pgfqpoint{5.631911in}{0.740510in}}%
\pgfpathlineto{\pgfqpoint{5.635267in}{0.739930in}}%
\pgfpathlineto{\pgfqpoint{5.638622in}{0.726079in}}%
\pgfpathlineto{\pgfqpoint{5.641978in}{0.666828in}}%
\pgfpathlineto{\pgfqpoint{5.645334in}{0.716331in}}%
\pgfpathlineto{\pgfqpoint{5.648690in}{0.745444in}}%
\pgfpathlineto{\pgfqpoint{5.652045in}{0.755154in}}%
\pgfpathlineto{\pgfqpoint{5.655401in}{0.749083in}}%
\pgfpathlineto{\pgfqpoint{5.658757in}{0.721413in}}%
\pgfpathlineto{\pgfqpoint{5.662113in}{0.666609in}}%
\pgfpathlineto{\pgfqpoint{5.665469in}{0.727797in}}%
\pgfpathlineto{\pgfqpoint{5.668824in}{0.740386in}}%
\pgfpathlineto{\pgfqpoint{5.672180in}{0.737480in}}%
\pgfpathlineto{\pgfqpoint{5.675536in}{0.714791in}}%
\pgfpathlineto{\pgfqpoint{5.678892in}{0.673448in}}%
\pgfpathlineto{\pgfqpoint{5.682247in}{0.737169in}}%
\pgfpathlineto{\pgfqpoint{5.685603in}{0.755403in}}%
\pgfpathlineto{\pgfqpoint{5.688959in}{0.757070in}}%
\pgfpathlineto{\pgfqpoint{5.692315in}{0.740935in}}%
\pgfpathlineto{\pgfqpoint{5.695671in}{0.702618in}}%
\pgfpathlineto{\pgfqpoint{5.699026in}{0.677159in}}%
\pgfpathlineto{\pgfqpoint{5.702382in}{0.713703in}}%
\pgfpathlineto{\pgfqpoint{5.705738in}{0.714601in}}%
\pgfpathlineto{\pgfqpoint{5.709094in}{0.685521in}}%
\pgfpathlineto{\pgfqpoint{5.712449in}{0.687029in}}%
\pgfpathlineto{\pgfqpoint{5.715805in}{0.729973in}}%
\pgfpathlineto{\pgfqpoint{5.719161in}{0.743743in}}%
\pgfpathlineto{\pgfqpoint{5.722517in}{0.739047in}}%
\pgfpathlineto{\pgfqpoint{5.725873in}{0.717103in}}%
\pgfpathlineto{\pgfqpoint{5.729228in}{0.656066in}}%
\pgfpathlineto{\pgfqpoint{5.732584in}{0.683157in}}%
\pgfpathlineto{\pgfqpoint{5.735940in}{0.701286in}}%
\pgfpathlineto{\pgfqpoint{5.739296in}{0.693021in}}%
\pgfpathlineto{\pgfqpoint{5.742652in}{0.615693in}}%
\pgfpathlineto{\pgfqpoint{5.746007in}{0.696754in}}%
\pgfpathlineto{\pgfqpoint{5.749363in}{0.723297in}}%
\pgfpathlineto{\pgfqpoint{5.752719in}{0.729851in}}%
\pgfpathlineto{\pgfqpoint{5.756075in}{0.716342in}}%
\pgfpathlineto{\pgfqpoint{5.759430in}{0.671702in}}%
\pgfpathlineto{\pgfqpoint{5.762786in}{0.678498in}}%
\pgfpathlineto{\pgfqpoint{5.766142in}{0.708704in}}%
\pgfpathlineto{\pgfqpoint{5.769498in}{0.714286in}}%
\pgfpathlineto{\pgfqpoint{5.772854in}{0.704444in}}%
\pgfpathlineto{\pgfqpoint{5.776209in}{0.657421in}}%
\pgfpathlineto{\pgfqpoint{5.782921in}{0.721918in}}%
\pgfpathlineto{\pgfqpoint{5.786277in}{0.732167in}}%
\pgfpathlineto{\pgfqpoint{5.789632in}{0.724599in}}%
\pgfpathlineto{\pgfqpoint{5.792988in}{0.684789in}}%
\pgfpathlineto{\pgfqpoint{5.796344in}{0.691099in}}%
\pgfpathlineto{\pgfqpoint{5.799700in}{0.725516in}}%
\pgfpathlineto{\pgfqpoint{5.803056in}{0.735705in}}%
\pgfpathlineto{\pgfqpoint{5.806411in}{0.734127in}}%
\pgfpathlineto{\pgfqpoint{5.809767in}{0.718786in}}%
\pgfpathlineto{\pgfqpoint{5.813123in}{0.645488in}}%
\pgfpathlineto{\pgfqpoint{5.816479in}{0.714418in}}%
\pgfpathlineto{\pgfqpoint{5.819834in}{0.741268in}}%
\pgfpathlineto{\pgfqpoint{5.823190in}{0.749886in}}%
\pgfpathlineto{\pgfqpoint{5.826546in}{0.742302in}}%
\pgfpathlineto{\pgfqpoint{5.829902in}{0.707135in}}%
\pgfpathlineto{\pgfqpoint{5.833258in}{0.696832in}}%
\pgfpathlineto{\pgfqpoint{5.836613in}{0.736189in}}%
\pgfpathlineto{\pgfqpoint{5.839969in}{0.746755in}}%
\pgfpathlineto{\pgfqpoint{5.843325in}{0.744471in}}%
\pgfpathlineto{\pgfqpoint{5.846681in}{0.726437in}}%
\pgfpathlineto{\pgfqpoint{5.850036in}{0.571158in}}%
\pgfpathlineto{\pgfqpoint{5.853392in}{0.733199in}}%
\pgfpathlineto{\pgfqpoint{5.856748in}{0.756044in}}%
\pgfpathlineto{\pgfqpoint{5.860104in}{0.761959in}}%
\pgfpathlineto{\pgfqpoint{5.863460in}{0.750731in}}%
\pgfpathlineto{\pgfqpoint{5.866815in}{0.719042in}}%
\pgfpathlineto{\pgfqpoint{5.870171in}{0.668352in}}%
\pgfpathlineto{\pgfqpoint{5.873527in}{0.722793in}}%
\pgfpathlineto{\pgfqpoint{5.876883in}{0.730450in}}%
\pgfpathlineto{\pgfqpoint{5.880238in}{0.717933in}}%
\pgfpathlineto{\pgfqpoint{5.883594in}{0.613424in}}%
\pgfpathlineto{\pgfqpoint{5.886950in}{0.723290in}}%
\pgfpathlineto{\pgfqpoint{5.890306in}{0.747901in}}%
\pgfpathlineto{\pgfqpoint{5.893662in}{0.753470in}}%
\pgfpathlineto{\pgfqpoint{5.897017in}{0.740587in}}%
\pgfpathlineto{\pgfqpoint{5.900373in}{0.711143in}}%
\pgfpathlineto{\pgfqpoint{5.903729in}{0.585186in}}%
\pgfpathlineto{\pgfqpoint{5.907085in}{0.695098in}}%
\pgfpathlineto{\pgfqpoint{5.910441in}{0.700281in}}%
\pgfpathlineto{\pgfqpoint{5.913796in}{0.670287in}}%
\pgfpathlineto{\pgfqpoint{5.917152in}{0.680424in}}%
\pgfpathlineto{\pgfqpoint{5.920508in}{0.721130in}}%
\pgfpathlineto{\pgfqpoint{5.923864in}{0.734278in}}%
\pgfpathlineto{\pgfqpoint{5.927219in}{0.728430in}}%
\pgfpathlineto{\pgfqpoint{5.930575in}{0.703275in}}%
\pgfpathlineto{\pgfqpoint{5.933931in}{0.555311in}}%
\pgfpathlineto{\pgfqpoint{5.937287in}{0.690698in}}%
\pgfpathlineto{\pgfqpoint{5.940643in}{0.704117in}}%
\pgfpathlineto{\pgfqpoint{5.943998in}{0.697210in}}%
\pgfpathlineto{\pgfqpoint{5.947354in}{0.651619in}}%
\pgfpathlineto{\pgfqpoint{5.954066in}{0.717279in}}%
\pgfpathlineto{\pgfqpoint{5.957421in}{0.727230in}}%
\pgfpathlineto{\pgfqpoint{5.960777in}{0.718127in}}%
\pgfpathlineto{\pgfqpoint{5.964133in}{0.676369in}}%
\pgfpathlineto{\pgfqpoint{5.967489in}{0.684319in}}%
\pgfpathlineto{\pgfqpoint{5.970845in}{0.717289in}}%
\pgfpathlineto{\pgfqpoint{5.974200in}{0.726326in}}%
\pgfpathlineto{\pgfqpoint{5.977556in}{0.723017in}}%
\pgfpathlineto{\pgfqpoint{5.980912in}{0.703335in}}%
\pgfpathlineto{\pgfqpoint{5.984268in}{0.614558in}}%
\pgfpathlineto{\pgfqpoint{5.987623in}{0.714313in}}%
\pgfpathlineto{\pgfqpoint{5.990979in}{0.735765in}}%
\pgfpathlineto{\pgfqpoint{5.994335in}{0.740848in}}%
\pgfpathlineto{\pgfqpoint{5.997691in}{0.727110in}}%
\pgfpathlineto{\pgfqpoint{6.001047in}{0.655641in}}%
\pgfpathlineto{\pgfqpoint{6.004402in}{0.716199in}}%
\pgfpathlineto{\pgfqpoint{6.007758in}{0.740330in}}%
\pgfpathlineto{\pgfqpoint{6.011114in}{0.747201in}}%
\pgfpathlineto{\pgfqpoint{6.014470in}{0.743063in}}%
\pgfpathlineto{\pgfqpoint{6.017825in}{0.723359in}}%
\pgfpathlineto{\pgfqpoint{6.021181in}{0.620039in}}%
\pgfpathlineto{\pgfqpoint{6.024537in}{0.732337in}}%
\pgfpathlineto{\pgfqpoint{6.027893in}{0.754420in}}%
\pgfpathlineto{\pgfqpoint{6.031249in}{0.760439in}}%
\pgfpathlineto{\pgfqpoint{6.034604in}{0.749801in}}%
\pgfpathlineto{\pgfqpoint{6.037960in}{0.713997in}}%
\pgfpathlineto{\pgfqpoint{6.041316in}{0.695705in}}%
\pgfpathlineto{\pgfqpoint{6.044672in}{0.735564in}}%
\pgfpathlineto{\pgfqpoint{6.048027in}{0.743796in}}%
\pgfpathlineto{\pgfqpoint{6.051383in}{0.736597in}}%
\pgfpathlineto{\pgfqpoint{6.054739in}{0.700726in}}%
\pgfpathlineto{\pgfqpoint{6.058095in}{0.709899in}}%
\pgfpathlineto{\pgfqpoint{6.061451in}{0.748557in}}%
\pgfpathlineto{\pgfqpoint{6.064806in}{0.761188in}}%
\pgfpathlineto{\pgfqpoint{6.068162in}{0.756751in}}%
\pgfpathlineto{\pgfqpoint{6.071518in}{0.736221in}}%
\pgfpathlineto{\pgfqpoint{6.074874in}{0.686624in}}%
\pgfpathlineto{\pgfqpoint{6.078230in}{0.687715in}}%
\pgfpathlineto{\pgfqpoint{6.081585in}{0.710612in}}%
\pgfpathlineto{\pgfqpoint{6.084941in}{0.701383in}}%
\pgfpathlineto{\pgfqpoint{6.088297in}{0.543954in}}%
\pgfpathlineto{\pgfqpoint{6.091653in}{0.714652in}}%
\pgfpathlineto{\pgfqpoint{6.095008in}{0.738326in}}%
\pgfpathlineto{\pgfqpoint{6.098364in}{0.742816in}}%
\pgfpathlineto{\pgfqpoint{6.101720in}{0.728233in}}%
\pgfpathlineto{\pgfqpoint{6.105076in}{0.695037in}}%
\pgfpathlineto{\pgfqpoint{6.108432in}{0.635240in}}%
\pgfpathlineto{\pgfqpoint{6.111787in}{0.690985in}}%
\pgfpathlineto{\pgfqpoint{6.115143in}{0.693359in}}%
\pgfpathlineto{\pgfqpoint{6.118499in}{0.661277in}}%
\pgfpathlineto{\pgfqpoint{6.121855in}{0.673971in}}%
\pgfpathlineto{\pgfqpoint{6.125210in}{0.713397in}}%
\pgfpathlineto{\pgfqpoint{6.128566in}{0.726030in}}%
\pgfpathlineto{\pgfqpoint{6.131922in}{0.719571in}}%
\pgfpathlineto{\pgfqpoint{6.135278in}{0.688886in}}%
\pgfpathlineto{\pgfqpoint{6.138634in}{0.652548in}}%
\pgfpathlineto{\pgfqpoint{6.141989in}{0.702838in}}%
\pgfpathlineto{\pgfqpoint{6.145345in}{0.713917in}}%
\pgfpathlineto{\pgfqpoint{6.148701in}{0.710110in}}%
\pgfpathlineto{\pgfqpoint{6.152057in}{0.685827in}}%
\pgfpathlineto{\pgfqpoint{6.155412in}{0.653101in}}%
\pgfpathlineto{\pgfqpoint{6.158768in}{0.711292in}}%
\pgfpathlineto{\pgfqpoint{6.162124in}{0.728667in}}%
\pgfpathlineto{\pgfqpoint{6.165480in}{0.729602in}}%
\pgfpathlineto{\pgfqpoint{6.168836in}{0.706979in}}%
\pgfpathlineto{\pgfqpoint{6.172191in}{0.652982in}}%
\pgfpathlineto{\pgfqpoint{6.175547in}{0.720170in}}%
\pgfpathlineto{\pgfqpoint{6.178903in}{0.736998in}}%
\pgfpathlineto{\pgfqpoint{6.182259in}{0.740638in}}%
\pgfpathlineto{\pgfqpoint{6.185614in}{0.733545in}}%
\pgfpathlineto{\pgfqpoint{6.188970in}{0.707484in}}%
\pgfpathlineto{\pgfqpoint{6.192326in}{0.674161in}}%
\pgfpathlineto{\pgfqpoint{6.195682in}{0.731906in}}%
\pgfpathlineto{\pgfqpoint{6.199038in}{0.749761in}}%
\pgfpathlineto{\pgfqpoint{6.202393in}{0.752954in}}%
\pgfpathlineto{\pgfqpoint{6.205749in}{0.739638in}}%
\pgfpathlineto{\pgfqpoint{6.209105in}{0.706363in}}%
\pgfpathlineto{\pgfqpoint{6.212461in}{0.658276in}}%
\pgfpathlineto{\pgfqpoint{6.215816in}{0.709136in}}%
\pgfpathlineto{\pgfqpoint{6.219172in}{0.714242in}}%
\pgfpathlineto{\pgfqpoint{6.222528in}{0.695132in}}%
\pgfpathlineto{\pgfqpoint{6.225884in}{0.654873in}}%
\pgfpathlineto{\pgfqpoint{6.229240in}{0.720027in}}%
\pgfpathlineto{\pgfqpoint{6.232595in}{0.738463in}}%
\pgfpathlineto{\pgfqpoint{6.235951in}{0.738807in}}%
\pgfpathlineto{\pgfqpoint{6.239307in}{0.720434in}}%
\pgfpathlineto{\pgfqpoint{6.242663in}{0.673456in}}%
\pgfpathlineto{\pgfqpoint{6.246019in}{0.674540in}}%
\pgfpathlineto{\pgfqpoint{6.249374in}{0.701281in}}%
\pgfpathlineto{\pgfqpoint{6.252730in}{0.699661in}}%
\pgfpathlineto{\pgfqpoint{6.256086in}{0.666669in}}%
\pgfpathlineto{\pgfqpoint{6.259442in}{0.676532in}}%
\pgfpathlineto{\pgfqpoint{6.262797in}{0.716097in}}%
\pgfpathlineto{\pgfqpoint{6.266153in}{0.728719in}}%
\pgfpathlineto{\pgfqpoint{6.269509in}{0.722656in}}%
\pgfpathlineto{\pgfqpoint{6.276221in}{0.659110in}}%
\pgfpathlineto{\pgfqpoint{6.279576in}{0.708015in}}%
\pgfpathlineto{\pgfqpoint{6.282932in}{0.719381in}}%
\pgfpathlineto{\pgfqpoint{6.286288in}{0.716405in}}%
\pgfpathlineto{\pgfqpoint{6.289644in}{0.694976in}}%
\pgfpathlineto{\pgfqpoint{6.292999in}{0.640427in}}%
\pgfpathlineto{\pgfqpoint{6.296355in}{0.713086in}}%
\pgfpathlineto{\pgfqpoint{6.299711in}{0.732443in}}%
\pgfpathlineto{\pgfqpoint{6.303067in}{0.735328in}}%
\pgfpathlineto{\pgfqpoint{6.306423in}{0.717162in}}%
\pgfpathlineto{\pgfqpoint{6.309778in}{0.475464in}}%
\pgfpathlineto{\pgfqpoint{6.313134in}{0.716463in}}%
\pgfpathlineto{\pgfqpoint{6.316490in}{0.736393in}}%
\pgfpathlineto{\pgfqpoint{6.319846in}{0.741130in}}%
\pgfpathlineto{\pgfqpoint{6.323201in}{0.734597in}}%
\pgfpathlineto{\pgfqpoint{6.326557in}{0.708925in}}%
\pgfpathlineto{\pgfqpoint{6.329913in}{0.675163in}}%
\pgfpathlineto{\pgfqpoint{6.333269in}{0.733283in}}%
\pgfpathlineto{\pgfqpoint{6.336625in}{0.751063in}}%
\pgfpathlineto{\pgfqpoint{6.339980in}{0.753819in}}%
\pgfpathlineto{\pgfqpoint{6.343336in}{0.738079in}}%
\pgfpathlineto{\pgfqpoint{6.346692in}{0.677286in}}%
\pgfpathlineto{\pgfqpoint{6.350048in}{0.717130in}}%
\pgfpathlineto{\pgfqpoint{6.353403in}{0.741414in}}%
\pgfpathlineto{\pgfqpoint{6.356759in}{0.746179in}}%
\pgfpathlineto{\pgfqpoint{6.360115in}{0.737266in}}%
\pgfpathlineto{\pgfqpoint{6.363471in}{0.699239in}}%
\pgfpathlineto{\pgfqpoint{6.366827in}{0.711562in}}%
\pgfpathlineto{\pgfqpoint{6.370182in}{0.748967in}}%
\pgfpathlineto{\pgfqpoint{6.373538in}{0.761441in}}%
\pgfpathlineto{\pgfqpoint{6.376894in}{0.757639in}}%
\pgfpathlineto{\pgfqpoint{6.383605in}{0.733396in}}%
\pgfpathlineto{\pgfqpoint{6.383605in}{0.733396in}}%
\pgfusepath{stroke}%
\end{pgfscope}%
\begin{pgfscope}%
\pgfpathrectangle{\pgfqpoint{4.621821in}{0.320625in}}{\pgfqpoint{1.845679in}{1.164375in}}%
\pgfusepath{clip}%
\pgfsetrectcap%
\pgfsetroundjoin%
\pgfsetlinewidth{0.602250pt}%
\definecolor{currentstroke}{rgb}{0.819608,0.286275,0.356863}%
\pgfsetstrokecolor{currentstroke}%
\pgfsetstrokeopacity{0.120000}%
\pgfsetdash{}{0pt}%
\pgfpathmoveto{\pgfqpoint{4.705715in}{1.471521in}}%
\pgfpathlineto{\pgfqpoint{4.729206in}{1.429543in}}%
\pgfpathlineto{\pgfqpoint{4.766120in}{1.360190in}}%
\pgfpathlineto{\pgfqpoint{4.786254in}{1.325552in}}%
\pgfpathlineto{\pgfqpoint{4.809745in}{1.287455in}}%
\pgfpathlineto{\pgfqpoint{4.839947in}{1.240860in}}%
\pgfpathlineto{\pgfqpoint{4.873504in}{1.191214in}}%
\pgfpathlineto{\pgfqpoint{4.913774in}{1.133723in}}%
\pgfpathlineto{\pgfqpoint{4.960755in}{1.068699in}}%
\pgfpathlineto{\pgfqpoint{5.017803in}{0.991748in}}%
\pgfpathlineto{\pgfqpoint{5.088274in}{0.898544in}}%
\pgfpathlineto{\pgfqpoint{5.145323in}{0.823984in}}%
\pgfpathlineto{\pgfqpoint{5.162102in}{0.802211in}}%
\pgfpathlineto{\pgfqpoint{5.182236in}{0.775408in}}%
\pgfpathlineto{\pgfqpoint{5.195659in}{0.760052in}}%
\pgfpathlineto{\pgfqpoint{5.202371in}{0.751874in}}%
\pgfpathlineto{\pgfqpoint{5.212438in}{0.737278in}}%
\pgfpathlineto{\pgfqpoint{5.229217in}{0.711447in}}%
\pgfpathlineto{\pgfqpoint{5.235929in}{0.703464in}}%
\pgfpathlineto{\pgfqpoint{5.242640in}{0.698332in}}%
\pgfpathlineto{\pgfqpoint{5.249352in}{0.695543in}}%
\pgfpathlineto{\pgfqpoint{5.256063in}{0.692717in}}%
\pgfpathlineto{\pgfqpoint{5.259419in}{0.690004in}}%
\pgfpathlineto{\pgfqpoint{5.262775in}{0.685608in}}%
\pgfpathlineto{\pgfqpoint{5.266131in}{0.680088in}}%
\pgfpathlineto{\pgfqpoint{5.272842in}{0.665602in}}%
\pgfpathlineto{\pgfqpoint{5.279554in}{0.645838in}}%
\pgfpathlineto{\pgfqpoint{5.286265in}{0.623894in}}%
\pgfpathlineto{\pgfqpoint{5.289621in}{0.617519in}}%
\pgfpathlineto{\pgfqpoint{5.292977in}{0.615884in}}%
\pgfpathlineto{\pgfqpoint{5.296333in}{0.615612in}}%
\pgfpathlineto{\pgfqpoint{5.299689in}{0.613628in}}%
\pgfpathlineto{\pgfqpoint{5.303044in}{0.608194in}}%
\pgfpathlineto{\pgfqpoint{5.309756in}{0.586513in}}%
\pgfpathlineto{\pgfqpoint{5.313112in}{0.594445in}}%
\pgfpathlineto{\pgfqpoint{5.316467in}{0.613691in}}%
\pgfpathlineto{\pgfqpoint{5.319823in}{0.629075in}}%
\pgfpathlineto{\pgfqpoint{5.323179in}{0.640355in}}%
\pgfpathlineto{\pgfqpoint{5.326535in}{0.648543in}}%
\pgfpathlineto{\pgfqpoint{5.329891in}{0.654225in}}%
\pgfpathlineto{\pgfqpoint{5.333246in}{0.657614in}}%
\pgfpathlineto{\pgfqpoint{5.336602in}{0.658603in}}%
\pgfpathlineto{\pgfqpoint{5.339958in}{0.656652in}}%
\pgfpathlineto{\pgfqpoint{5.343314in}{0.651641in}}%
\pgfpathlineto{\pgfqpoint{5.346669in}{0.656125in}}%
\pgfpathlineto{\pgfqpoint{5.353381in}{0.689628in}}%
\pgfpathlineto{\pgfqpoint{5.356737in}{0.700738in}}%
\pgfpathlineto{\pgfqpoint{5.360093in}{0.708125in}}%
\pgfpathlineto{\pgfqpoint{5.363448in}{0.712518in}}%
\pgfpathlineto{\pgfqpoint{5.366804in}{0.714276in}}%
\pgfpathlineto{\pgfqpoint{5.370160in}{0.713435in}}%
\pgfpathlineto{\pgfqpoint{5.373516in}{0.709667in}}%
\pgfpathlineto{\pgfqpoint{5.376871in}{0.702009in}}%
\pgfpathlineto{\pgfqpoint{5.380227in}{0.687835in}}%
\pgfpathlineto{\pgfqpoint{5.383583in}{0.656720in}}%
\pgfpathlineto{\pgfqpoint{5.386939in}{0.632437in}}%
\pgfpathlineto{\pgfqpoint{5.390295in}{0.683709in}}%
\pgfpathlineto{\pgfqpoint{5.393650in}{0.704383in}}%
\pgfpathlineto{\pgfqpoint{5.397006in}{0.716316in}}%
\pgfpathlineto{\pgfqpoint{5.400362in}{0.723276in}}%
\pgfpathlineto{\pgfqpoint{5.403718in}{0.726216in}}%
\pgfpathlineto{\pgfqpoint{5.407074in}{0.724775in}}%
\pgfpathlineto{\pgfqpoint{5.410429in}{0.716743in}}%
\pgfpathlineto{\pgfqpoint{5.413785in}{0.697855in}}%
\pgfpathlineto{\pgfqpoint{5.417141in}{0.659537in}}%
\pgfpathlineto{\pgfqpoint{5.420497in}{0.665686in}}%
\pgfpathlineto{\pgfqpoint{5.423852in}{0.685470in}}%
\pgfpathlineto{\pgfqpoint{5.427208in}{0.686624in}}%
\pgfpathlineto{\pgfqpoint{5.430564in}{0.669568in}}%
\pgfpathlineto{\pgfqpoint{5.433920in}{0.609086in}}%
\pgfpathlineto{\pgfqpoint{5.437276in}{0.682940in}}%
\pgfpathlineto{\pgfqpoint{5.440631in}{0.705134in}}%
\pgfpathlineto{\pgfqpoint{5.443987in}{0.711421in}}%
\pgfpathlineto{\pgfqpoint{5.447343in}{0.701588in}}%
\pgfpathlineto{\pgfqpoint{5.450699in}{0.675371in}}%
\pgfpathlineto{\pgfqpoint{5.454054in}{0.617403in}}%
\pgfpathlineto{\pgfqpoint{5.457410in}{0.667863in}}%
\pgfpathlineto{\pgfqpoint{5.460766in}{0.680974in}}%
\pgfpathlineto{\pgfqpoint{5.464122in}{0.677146in}}%
\pgfpathlineto{\pgfqpoint{5.467478in}{0.650072in}}%
\pgfpathlineto{\pgfqpoint{5.470833in}{0.635436in}}%
\pgfpathlineto{\pgfqpoint{5.474189in}{0.684449in}}%
\pgfpathlineto{\pgfqpoint{5.477545in}{0.701081in}}%
\pgfpathlineto{\pgfqpoint{5.480901in}{0.702770in}}%
\pgfpathlineto{\pgfqpoint{5.484256in}{0.681905in}}%
\pgfpathlineto{\pgfqpoint{5.487612in}{0.637574in}}%
\pgfpathlineto{\pgfqpoint{5.490968in}{0.700709in}}%
\pgfpathlineto{\pgfqpoint{5.494324in}{0.719857in}}%
\pgfpathlineto{\pgfqpoint{5.497680in}{0.727295in}}%
\pgfpathlineto{\pgfqpoint{5.501035in}{0.726977in}}%
\pgfpathlineto{\pgfqpoint{5.504391in}{0.718238in}}%
\pgfpathlineto{\pgfqpoint{5.507747in}{0.693227in}}%
\pgfpathlineto{\pgfqpoint{5.511103in}{0.642251in}}%
\pgfpathlineto{\pgfqpoint{5.514458in}{0.711104in}}%
\pgfpathlineto{\pgfqpoint{5.517814in}{0.731539in}}%
\pgfpathlineto{\pgfqpoint{5.521170in}{0.739098in}}%
\pgfpathlineto{\pgfqpoint{5.524526in}{0.734874in}}%
\pgfpathlineto{\pgfqpoint{5.527882in}{0.715927in}}%
\pgfpathlineto{\pgfqpoint{5.531237in}{0.665219in}}%
\pgfpathlineto{\pgfqpoint{5.534593in}{0.680194in}}%
\pgfpathlineto{\pgfqpoint{5.537949in}{0.705822in}}%
\pgfpathlineto{\pgfqpoint{5.541305in}{0.708597in}}%
\pgfpathlineto{\pgfqpoint{5.544660in}{0.693931in}}%
\pgfpathlineto{\pgfqpoint{5.548016in}{0.584144in}}%
\pgfpathlineto{\pgfqpoint{5.551372in}{0.699110in}}%
\pgfpathlineto{\pgfqpoint{5.554728in}{0.724055in}}%
\pgfpathlineto{\pgfqpoint{5.558084in}{0.731741in}}%
\pgfpathlineto{\pgfqpoint{5.561439in}{0.723192in}}%
\pgfpathlineto{\pgfqpoint{5.564795in}{0.697646in}}%
\pgfpathlineto{\pgfqpoint{5.568151in}{0.592670in}}%
\pgfpathlineto{\pgfqpoint{5.571507in}{0.687205in}}%
\pgfpathlineto{\pgfqpoint{5.574863in}{0.702173in}}%
\pgfpathlineto{\pgfqpoint{5.578218in}{0.699430in}}%
\pgfpathlineto{\pgfqpoint{5.584930in}{0.648821in}}%
\pgfpathlineto{\pgfqpoint{5.588286in}{0.704074in}}%
\pgfpathlineto{\pgfqpoint{5.591641in}{0.721809in}}%
\pgfpathlineto{\pgfqpoint{5.594997in}{0.724336in}}%
\pgfpathlineto{\pgfqpoint{5.598353in}{0.708574in}}%
\pgfpathlineto{\pgfqpoint{5.601709in}{0.653200in}}%
\pgfpathlineto{\pgfqpoint{5.608420in}{0.711797in}}%
\pgfpathlineto{\pgfqpoint{5.611776in}{0.719107in}}%
\pgfpathlineto{\pgfqpoint{5.615132in}{0.715043in}}%
\pgfpathlineto{\pgfqpoint{5.618488in}{0.695199in}}%
\pgfpathlineto{\pgfqpoint{5.621843in}{0.599677in}}%
\pgfpathlineto{\pgfqpoint{5.625199in}{0.705384in}}%
\pgfpathlineto{\pgfqpoint{5.628555in}{0.727729in}}%
\pgfpathlineto{\pgfqpoint{5.631911in}{0.734710in}}%
\pgfpathlineto{\pgfqpoint{5.635267in}{0.726409in}}%
\pgfpathlineto{\pgfqpoint{5.638622in}{0.692188in}}%
\pgfpathlineto{\pgfqpoint{5.641978in}{0.679389in}}%
\pgfpathlineto{\pgfqpoint{5.645334in}{0.720940in}}%
\pgfpathlineto{\pgfqpoint{5.648690in}{0.733495in}}%
\pgfpathlineto{\pgfqpoint{5.652045in}{0.734572in}}%
\pgfpathlineto{\pgfqpoint{5.655401in}{0.724675in}}%
\pgfpathlineto{\pgfqpoint{5.658757in}{0.690651in}}%
\pgfpathlineto{\pgfqpoint{5.662113in}{0.688642in}}%
\pgfpathlineto{\pgfqpoint{5.665469in}{0.731171in}}%
\pgfpathlineto{\pgfqpoint{5.668824in}{0.746505in}}%
\pgfpathlineto{\pgfqpoint{5.672180in}{0.748692in}}%
\pgfpathlineto{\pgfqpoint{5.675536in}{0.734885in}}%
\pgfpathlineto{\pgfqpoint{5.678892in}{0.694580in}}%
\pgfpathlineto{\pgfqpoint{5.682247in}{0.690166in}}%
\pgfpathlineto{\pgfqpoint{5.685603in}{0.725096in}}%
\pgfpathlineto{\pgfqpoint{5.688959in}{0.733567in}}%
\pgfpathlineto{\pgfqpoint{5.692315in}{0.728828in}}%
\pgfpathlineto{\pgfqpoint{5.695671in}{0.704741in}}%
\pgfpathlineto{\pgfqpoint{5.699026in}{0.667225in}}%
\pgfpathlineto{\pgfqpoint{5.702382in}{0.728495in}}%
\pgfpathlineto{\pgfqpoint{5.705738in}{0.747214in}}%
\pgfpathlineto{\pgfqpoint{5.709094in}{0.751173in}}%
\pgfpathlineto{\pgfqpoint{5.712449in}{0.739308in}}%
\pgfpathlineto{\pgfqpoint{5.715805in}{0.711658in}}%
\pgfpathlineto{\pgfqpoint{5.719161in}{0.593333in}}%
\pgfpathlineto{\pgfqpoint{5.722517in}{0.699210in}}%
\pgfpathlineto{\pgfqpoint{5.725873in}{0.710915in}}%
\pgfpathlineto{\pgfqpoint{5.729228in}{0.701934in}}%
\pgfpathlineto{\pgfqpoint{5.732584in}{0.641919in}}%
\pgfpathlineto{\pgfqpoint{5.735940in}{0.698590in}}%
\pgfpathlineto{\pgfqpoint{5.739296in}{0.728050in}}%
\pgfpathlineto{\pgfqpoint{5.742652in}{0.737934in}}%
\pgfpathlineto{\pgfqpoint{5.746007in}{0.731425in}}%
\pgfpathlineto{\pgfqpoint{5.749363in}{0.709913in}}%
\pgfpathlineto{\pgfqpoint{5.752719in}{0.652568in}}%
\pgfpathlineto{\pgfqpoint{5.759430in}{0.696099in}}%
\pgfpathlineto{\pgfqpoint{5.762786in}{0.693565in}}%
\pgfpathlineto{\pgfqpoint{5.766142in}{0.662101in}}%
\pgfpathlineto{\pgfqpoint{5.769498in}{0.666430in}}%
\pgfpathlineto{\pgfqpoint{5.772854in}{0.708399in}}%
\pgfpathlineto{\pgfqpoint{5.776209in}{0.723076in}}%
\pgfpathlineto{\pgfqpoint{5.779565in}{0.722017in}}%
\pgfpathlineto{\pgfqpoint{5.782921in}{0.702479in}}%
\pgfpathlineto{\pgfqpoint{5.786277in}{0.634054in}}%
\pgfpathlineto{\pgfqpoint{5.789632in}{0.682797in}}%
\pgfpathlineto{\pgfqpoint{5.792988in}{0.705250in}}%
\pgfpathlineto{\pgfqpoint{5.796344in}{0.709313in}}%
\pgfpathlineto{\pgfqpoint{5.799700in}{0.700042in}}%
\pgfpathlineto{\pgfqpoint{5.803056in}{0.661798in}}%
\pgfpathlineto{\pgfqpoint{5.806411in}{0.674508in}}%
\pgfpathlineto{\pgfqpoint{5.809767in}{0.712265in}}%
\pgfpathlineto{\pgfqpoint{5.813123in}{0.725793in}}%
\pgfpathlineto{\pgfqpoint{5.816479in}{0.724933in}}%
\pgfpathlineto{\pgfqpoint{5.819834in}{0.703857in}}%
\pgfpathlineto{\pgfqpoint{5.823190in}{0.590383in}}%
\pgfpathlineto{\pgfqpoint{5.826546in}{0.705102in}}%
\pgfpathlineto{\pgfqpoint{5.829902in}{0.724579in}}%
\pgfpathlineto{\pgfqpoint{5.833258in}{0.729896in}}%
\pgfpathlineto{\pgfqpoint{5.836613in}{0.725177in}}%
\pgfpathlineto{\pgfqpoint{5.839969in}{0.705827in}}%
\pgfpathlineto{\pgfqpoint{5.843325in}{0.553966in}}%
\pgfpathlineto{\pgfqpoint{5.846681in}{0.712041in}}%
\pgfpathlineto{\pgfqpoint{5.850036in}{0.735525in}}%
\pgfpathlineto{\pgfqpoint{5.853392in}{0.743587in}}%
\pgfpathlineto{\pgfqpoint{5.856748in}{0.737608in}}%
\pgfpathlineto{\pgfqpoint{5.860104in}{0.710485in}}%
\pgfpathlineto{\pgfqpoint{5.863460in}{0.659425in}}%
\pgfpathlineto{\pgfqpoint{5.866815in}{0.721467in}}%
\pgfpathlineto{\pgfqpoint{5.870171in}{0.736653in}}%
\pgfpathlineto{\pgfqpoint{5.873527in}{0.738710in}}%
\pgfpathlineto{\pgfqpoint{5.876883in}{0.729213in}}%
\pgfpathlineto{\pgfqpoint{5.880238in}{0.694885in}}%
\pgfpathlineto{\pgfqpoint{5.883594in}{0.694887in}}%
\pgfpathlineto{\pgfqpoint{5.886950in}{0.736735in}}%
\pgfpathlineto{\pgfqpoint{5.890306in}{0.751734in}}%
\pgfpathlineto{\pgfqpoint{5.893662in}{0.753202in}}%
\pgfpathlineto{\pgfqpoint{5.897017in}{0.738948in}}%
\pgfpathlineto{\pgfqpoint{5.903729in}{0.668916in}}%
\pgfpathlineto{\pgfqpoint{5.907085in}{0.714856in}}%
\pgfpathlineto{\pgfqpoint{5.910441in}{0.722696in}}%
\pgfpathlineto{\pgfqpoint{5.913796in}{0.713254in}}%
\pgfpathlineto{\pgfqpoint{5.917152in}{0.662348in}}%
\pgfpathlineto{\pgfqpoint{5.920508in}{0.703072in}}%
\pgfpathlineto{\pgfqpoint{5.923864in}{0.734598in}}%
\pgfpathlineto{\pgfqpoint{5.927219in}{0.745498in}}%
\pgfpathlineto{\pgfqpoint{5.930575in}{0.740555in}}%
\pgfpathlineto{\pgfqpoint{5.933931in}{0.721348in}}%
\pgfpathlineto{\pgfqpoint{5.937287in}{0.678437in}}%
\pgfpathlineto{\pgfqpoint{5.940643in}{0.663494in}}%
\pgfpathlineto{\pgfqpoint{5.943998in}{0.694525in}}%
\pgfpathlineto{\pgfqpoint{5.947354in}{0.693057in}}%
\pgfpathlineto{\pgfqpoint{5.950710in}{0.656921in}}%
\pgfpathlineto{\pgfqpoint{5.954066in}{0.676184in}}%
\pgfpathlineto{\pgfqpoint{5.957421in}{0.713847in}}%
\pgfpathlineto{\pgfqpoint{5.960777in}{0.727202in}}%
\pgfpathlineto{\pgfqpoint{5.964133in}{0.724320in}}%
\pgfpathlineto{\pgfqpoint{5.967489in}{0.704533in}}%
\pgfpathlineto{\pgfqpoint{5.970845in}{0.648395in}}%
\pgfpathlineto{\pgfqpoint{5.977556in}{0.696469in}}%
\pgfpathlineto{\pgfqpoint{5.980912in}{0.698353in}}%
\pgfpathlineto{\pgfqpoint{5.984268in}{0.682191in}}%
\pgfpathlineto{\pgfqpoint{5.987623in}{0.549664in}}%
\pgfpathlineto{\pgfqpoint{5.990979in}{0.691781in}}%
\pgfpathlineto{\pgfqpoint{5.994335in}{0.715152in}}%
\pgfpathlineto{\pgfqpoint{5.997691in}{0.721792in}}%
\pgfpathlineto{\pgfqpoint{6.001047in}{0.711220in}}%
\pgfpathlineto{\pgfqpoint{6.004402in}{0.671293in}}%
\pgfpathlineto{\pgfqpoint{6.007758in}{0.673422in}}%
\pgfpathlineto{\pgfqpoint{6.011114in}{0.708773in}}%
\pgfpathlineto{\pgfqpoint{6.014470in}{0.719449in}}%
\pgfpathlineto{\pgfqpoint{6.017825in}{0.718832in}}%
\pgfpathlineto{\pgfqpoint{6.021181in}{0.706056in}}%
\pgfpathlineto{\pgfqpoint{6.024537in}{0.657803in}}%
\pgfpathlineto{\pgfqpoint{6.031249in}{0.721543in}}%
\pgfpathlineto{\pgfqpoint{6.034604in}{0.733711in}}%
\pgfpathlineto{\pgfqpoint{6.037960in}{0.732467in}}%
\pgfpathlineto{\pgfqpoint{6.041316in}{0.711288in}}%
\pgfpathlineto{\pgfqpoint{6.044672in}{0.609380in}}%
\pgfpathlineto{\pgfqpoint{6.048027in}{0.714413in}}%
\pgfpathlineto{\pgfqpoint{6.051383in}{0.733742in}}%
\pgfpathlineto{\pgfqpoint{6.054739in}{0.739202in}}%
\pgfpathlineto{\pgfqpoint{6.058095in}{0.734864in}}%
\pgfpathlineto{\pgfqpoint{6.061451in}{0.716638in}}%
\pgfpathlineto{\pgfqpoint{6.064806in}{0.601113in}}%
\pgfpathlineto{\pgfqpoint{6.068162in}{0.718726in}}%
\pgfpathlineto{\pgfqpoint{6.071518in}{0.743554in}}%
\pgfpathlineto{\pgfqpoint{6.074874in}{0.752428in}}%
\pgfpathlineto{\pgfqpoint{6.078230in}{0.747720in}}%
\pgfpathlineto{\pgfqpoint{6.081585in}{0.725656in}}%
\pgfpathlineto{\pgfqpoint{6.084941in}{0.630125in}}%
\pgfpathlineto{\pgfqpoint{6.088297in}{0.712877in}}%
\pgfpathlineto{\pgfqpoint{6.091653in}{0.732004in}}%
\pgfpathlineto{\pgfqpoint{6.095008in}{0.734056in}}%
\pgfpathlineto{\pgfqpoint{6.098364in}{0.721671in}}%
\pgfpathlineto{\pgfqpoint{6.101720in}{0.664987in}}%
\pgfpathlineto{\pgfqpoint{6.105076in}{0.712276in}}%
\pgfpathlineto{\pgfqpoint{6.108432in}{0.742294in}}%
\pgfpathlineto{\pgfqpoint{6.111787in}{0.752802in}}%
\pgfpathlineto{\pgfqpoint{6.115143in}{0.748241in}}%
\pgfpathlineto{\pgfqpoint{6.118499in}{0.729495in}}%
\pgfpathlineto{\pgfqpoint{6.121855in}{0.687999in}}%
\pgfpathlineto{\pgfqpoint{6.125210in}{0.668634in}}%
\pgfpathlineto{\pgfqpoint{6.128566in}{0.701786in}}%
\pgfpathlineto{\pgfqpoint{6.131922in}{0.701212in}}%
\pgfpathlineto{\pgfqpoint{6.135278in}{0.668112in}}%
\pgfpathlineto{\pgfqpoint{6.138634in}{0.680444in}}%
\pgfpathlineto{\pgfqpoint{6.141989in}{0.720247in}}%
\pgfpathlineto{\pgfqpoint{6.145345in}{0.734324in}}%
\pgfpathlineto{\pgfqpoint{6.148701in}{0.732267in}}%
\pgfpathlineto{\pgfqpoint{6.152057in}{0.714598in}}%
\pgfpathlineto{\pgfqpoint{6.155412in}{0.675377in}}%
\pgfpathlineto{\pgfqpoint{6.158768in}{0.650191in}}%
\pgfpathlineto{\pgfqpoint{6.162124in}{0.687212in}}%
\pgfpathlineto{\pgfqpoint{6.165480in}{0.689206in}}%
\pgfpathlineto{\pgfqpoint{6.168836in}{0.665203in}}%
\pgfpathlineto{\pgfqpoint{6.172191in}{0.648493in}}%
\pgfpathlineto{\pgfqpoint{6.175547in}{0.700245in}}%
\pgfpathlineto{\pgfqpoint{6.178903in}{0.717441in}}%
\pgfpathlineto{\pgfqpoint{6.182259in}{0.718810in}}%
\pgfpathlineto{\pgfqpoint{6.185614in}{0.701411in}}%
\pgfpathlineto{\pgfqpoint{6.188970in}{0.642972in}}%
\pgfpathlineto{\pgfqpoint{6.192326in}{0.677494in}}%
\pgfpathlineto{\pgfqpoint{6.195682in}{0.702771in}}%
\pgfpathlineto{\pgfqpoint{6.199038in}{0.708583in}}%
\pgfpathlineto{\pgfqpoint{6.202393in}{0.701945in}}%
\pgfpathlineto{\pgfqpoint{6.205749in}{0.674027in}}%
\pgfpathlineto{\pgfqpoint{6.209105in}{0.654150in}}%
\pgfpathlineto{\pgfqpoint{6.212461in}{0.705130in}}%
\pgfpathlineto{\pgfqpoint{6.215816in}{0.722095in}}%
\pgfpathlineto{\pgfqpoint{6.219172in}{0.724771in}}%
\pgfpathlineto{\pgfqpoint{6.222528in}{0.708943in}}%
\pgfpathlineto{\pgfqpoint{6.225884in}{0.628095in}}%
\pgfpathlineto{\pgfqpoint{6.229240in}{0.700916in}}%
\pgfpathlineto{\pgfqpoint{6.232595in}{0.724837in}}%
\pgfpathlineto{\pgfqpoint{6.235951in}{0.732788in}}%
\pgfpathlineto{\pgfqpoint{6.239307in}{0.731177in}}%
\pgfpathlineto{\pgfqpoint{6.242663in}{0.718393in}}%
\pgfpathlineto{\pgfqpoint{6.246019in}{0.674138in}}%
\pgfpathlineto{\pgfqpoint{6.249374in}{0.696228in}}%
\pgfpathlineto{\pgfqpoint{6.252730in}{0.731008in}}%
\pgfpathlineto{\pgfqpoint{6.256086in}{0.744233in}}%
\pgfpathlineto{\pgfqpoint{6.259442in}{0.744923in}}%
\pgfpathlineto{\pgfqpoint{6.262797in}{0.728597in}}%
\pgfpathlineto{\pgfqpoint{6.266153in}{0.665949in}}%
\pgfpathlineto{\pgfqpoint{6.269509in}{0.710307in}}%
\pgfpathlineto{\pgfqpoint{6.272865in}{0.735543in}}%
\pgfpathlineto{\pgfqpoint{6.276221in}{0.742656in}}%
\pgfpathlineto{\pgfqpoint{6.279576in}{0.738818in}}%
\pgfpathlineto{\pgfqpoint{6.282932in}{0.720033in}}%
\pgfpathlineto{\pgfqpoint{6.286288in}{0.553858in}}%
\pgfpathlineto{\pgfqpoint{6.289644in}{0.726262in}}%
\pgfpathlineto{\pgfqpoint{6.292999in}{0.749794in}}%
\pgfpathlineto{\pgfqpoint{6.296355in}{0.757574in}}%
\pgfpathlineto{\pgfqpoint{6.299711in}{0.750847in}}%
\pgfpathlineto{\pgfqpoint{6.303067in}{0.729269in}}%
\pgfpathlineto{\pgfqpoint{6.306423in}{0.669045in}}%
\pgfpathlineto{\pgfqpoint{6.309778in}{0.696964in}}%
\pgfpathlineto{\pgfqpoint{6.313134in}{0.718041in}}%
\pgfpathlineto{\pgfqpoint{6.316490in}{0.716180in}}%
\pgfpathlineto{\pgfqpoint{6.319846in}{0.688420in}}%
\pgfpathlineto{\pgfqpoint{6.323201in}{0.680350in}}%
\pgfpathlineto{\pgfqpoint{6.326557in}{0.727229in}}%
\pgfpathlineto{\pgfqpoint{6.329913in}{0.743188in}}%
\pgfpathlineto{\pgfqpoint{6.333269in}{0.743554in}}%
\pgfpathlineto{\pgfqpoint{6.336625in}{0.728301in}}%
\pgfpathlineto{\pgfqpoint{6.339980in}{0.698397in}}%
\pgfpathlineto{\pgfqpoint{6.343336in}{0.551701in}}%
\pgfpathlineto{\pgfqpoint{6.346692in}{0.680284in}}%
\pgfpathlineto{\pgfqpoint{6.350048in}{0.686098in}}%
\pgfpathlineto{\pgfqpoint{6.353403in}{0.658121in}}%
\pgfpathlineto{\pgfqpoint{6.356759in}{0.663514in}}%
\pgfpathlineto{\pgfqpoint{6.360115in}{0.706775in}}%
\pgfpathlineto{\pgfqpoint{6.363471in}{0.722147in}}%
\pgfpathlineto{\pgfqpoint{6.366827in}{0.721352in}}%
\pgfpathlineto{\pgfqpoint{6.380250in}{0.676277in}}%
\pgfpathlineto{\pgfqpoint{6.383605in}{0.665727in}}%
\pgfpathlineto{\pgfqpoint{6.383605in}{0.665727in}}%
\pgfusepath{stroke}%
\end{pgfscope}%
\begin{pgfscope}%
\pgfpathrectangle{\pgfqpoint{4.621821in}{0.320625in}}{\pgfqpoint{1.845679in}{1.164375in}}%
\pgfusepath{clip}%
\pgfsetrectcap%
\pgfsetroundjoin%
\pgfsetlinewidth{0.602250pt}%
\definecolor{currentstroke}{rgb}{0.819608,0.286275,0.356863}%
\pgfsetstrokecolor{currentstroke}%
\pgfsetstrokeopacity{0.120000}%
\pgfsetdash{}{0pt}%
\pgfpathmoveto{\pgfqpoint{4.705715in}{1.473414in}}%
\pgfpathlineto{\pgfqpoint{4.729206in}{1.430786in}}%
\pgfpathlineto{\pgfqpoint{4.776187in}{1.339404in}}%
\pgfpathlineto{\pgfqpoint{4.792966in}{1.310245in}}%
\pgfpathlineto{\pgfqpoint{4.813100in}{1.277468in}}%
\pgfpathlineto{\pgfqpoint{4.839947in}{1.236225in}}%
\pgfpathlineto{\pgfqpoint{4.873504in}{1.186906in}}%
\pgfpathlineto{\pgfqpoint{4.920485in}{1.119923in}}%
\pgfpathlineto{\pgfqpoint{4.987601in}{1.026222in}}%
\pgfpathlineto{\pgfqpoint{5.068140in}{0.915610in}}%
\pgfpathlineto{\pgfqpoint{5.168813in}{0.779574in}}%
\pgfpathlineto{\pgfqpoint{5.185592in}{0.759155in}}%
\pgfpathlineto{\pgfqpoint{5.195659in}{0.744208in}}%
\pgfpathlineto{\pgfqpoint{5.209082in}{0.723408in}}%
\pgfpathlineto{\pgfqpoint{5.215794in}{0.714760in}}%
\pgfpathlineto{\pgfqpoint{5.222506in}{0.708578in}}%
\pgfpathlineto{\pgfqpoint{5.229217in}{0.704636in}}%
\pgfpathlineto{\pgfqpoint{5.235929in}{0.700840in}}%
\pgfpathlineto{\pgfqpoint{5.239285in}{0.697715in}}%
\pgfpathlineto{\pgfqpoint{5.242640in}{0.693053in}}%
\pgfpathlineto{\pgfqpoint{5.249352in}{0.680765in}}%
\pgfpathlineto{\pgfqpoint{5.256063in}{0.664175in}}%
\pgfpathlineto{\pgfqpoint{5.269487in}{0.623483in}}%
\pgfpathlineto{\pgfqpoint{5.272842in}{0.617351in}}%
\pgfpathlineto{\pgfqpoint{5.282910in}{0.603903in}}%
\pgfpathlineto{\pgfqpoint{5.286265in}{0.603748in}}%
\pgfpathlineto{\pgfqpoint{5.289621in}{0.612587in}}%
\pgfpathlineto{\pgfqpoint{5.296333in}{0.636615in}}%
\pgfpathlineto{\pgfqpoint{5.299689in}{0.645751in}}%
\pgfpathlineto{\pgfqpoint{5.303044in}{0.652591in}}%
\pgfpathlineto{\pgfqpoint{5.306400in}{0.657227in}}%
\pgfpathlineto{\pgfqpoint{5.309756in}{0.659590in}}%
\pgfpathlineto{\pgfqpoint{5.313112in}{0.659307in}}%
\pgfpathlineto{\pgfqpoint{5.319823in}{0.652443in}}%
\pgfpathlineto{\pgfqpoint{5.323179in}{0.664588in}}%
\pgfpathlineto{\pgfqpoint{5.326535in}{0.682446in}}%
\pgfpathlineto{\pgfqpoint{5.329891in}{0.696580in}}%
\pgfpathlineto{\pgfqpoint{5.333246in}{0.706813in}}%
\pgfpathlineto{\pgfqpoint{5.336602in}{0.714011in}}%
\pgfpathlineto{\pgfqpoint{5.339958in}{0.718782in}}%
\pgfpathlineto{\pgfqpoint{5.343314in}{0.721468in}}%
\pgfpathlineto{\pgfqpoint{5.346669in}{0.722207in}}%
\pgfpathlineto{\pgfqpoint{5.350025in}{0.720947in}}%
\pgfpathlineto{\pgfqpoint{5.353381in}{0.717407in}}%
\pgfpathlineto{\pgfqpoint{5.356737in}{0.710907in}}%
\pgfpathlineto{\pgfqpoint{5.360093in}{0.699839in}}%
\pgfpathlineto{\pgfqpoint{5.363448in}{0.679393in}}%
\pgfpathlineto{\pgfqpoint{5.366804in}{0.616024in}}%
\pgfpathlineto{\pgfqpoint{5.370160in}{0.668114in}}%
\pgfpathlineto{\pgfqpoint{5.373516in}{0.698349in}}%
\pgfpathlineto{\pgfqpoint{5.376871in}{0.714413in}}%
\pgfpathlineto{\pgfqpoint{5.380227in}{0.724575in}}%
\pgfpathlineto{\pgfqpoint{5.383583in}{0.731030in}}%
\pgfpathlineto{\pgfqpoint{5.386939in}{0.734484in}}%
\pgfpathlineto{\pgfqpoint{5.390295in}{0.734920in}}%
\pgfpathlineto{\pgfqpoint{5.393650in}{0.731564in}}%
\pgfpathlineto{\pgfqpoint{5.397006in}{0.722228in}}%
\pgfpathlineto{\pgfqpoint{5.400362in}{0.705525in}}%
\pgfpathlineto{\pgfqpoint{5.407074in}{0.653085in}}%
\pgfpathlineto{\pgfqpoint{5.410429in}{0.659247in}}%
\pgfpathlineto{\pgfqpoint{5.413785in}{0.654609in}}%
\pgfpathlineto{\pgfqpoint{5.417141in}{0.628700in}}%
\pgfpathlineto{\pgfqpoint{5.420497in}{0.665344in}}%
\pgfpathlineto{\pgfqpoint{5.423852in}{0.690502in}}%
\pgfpathlineto{\pgfqpoint{5.427208in}{0.700105in}}%
\pgfpathlineto{\pgfqpoint{5.430564in}{0.694432in}}%
\pgfpathlineto{\pgfqpoint{5.433920in}{0.677237in}}%
\pgfpathlineto{\pgfqpoint{5.437276in}{0.647262in}}%
\pgfpathlineto{\pgfqpoint{5.440631in}{0.610061in}}%
\pgfpathlineto{\pgfqpoint{5.443987in}{0.630918in}}%
\pgfpathlineto{\pgfqpoint{5.447343in}{0.629116in}}%
\pgfpathlineto{\pgfqpoint{5.450699in}{0.588173in}}%
\pgfpathlineto{\pgfqpoint{5.454054in}{0.633373in}}%
\pgfpathlineto{\pgfqpoint{5.457410in}{0.662774in}}%
\pgfpathlineto{\pgfqpoint{5.460766in}{0.673930in}}%
\pgfpathlineto{\pgfqpoint{5.464122in}{0.669144in}}%
\pgfpathlineto{\pgfqpoint{5.467478in}{0.622146in}}%
\pgfpathlineto{\pgfqpoint{5.470833in}{0.660845in}}%
\pgfpathlineto{\pgfqpoint{5.474189in}{0.692364in}}%
\pgfpathlineto{\pgfqpoint{5.477545in}{0.706280in}}%
\pgfpathlineto{\pgfqpoint{5.480901in}{0.712013in}}%
\pgfpathlineto{\pgfqpoint{5.484256in}{0.711514in}}%
\pgfpathlineto{\pgfqpoint{5.487612in}{0.704123in}}%
\pgfpathlineto{\pgfqpoint{5.490968in}{0.684827in}}%
\pgfpathlineto{\pgfqpoint{5.494324in}{0.595944in}}%
\pgfpathlineto{\pgfqpoint{5.497680in}{0.681997in}}%
\pgfpathlineto{\pgfqpoint{5.501035in}{0.708699in}}%
\pgfpathlineto{\pgfqpoint{5.504391in}{0.720422in}}%
\pgfpathlineto{\pgfqpoint{5.507747in}{0.722751in}}%
\pgfpathlineto{\pgfqpoint{5.511103in}{0.713111in}}%
\pgfpathlineto{\pgfqpoint{5.514458in}{0.685328in}}%
\pgfpathlineto{\pgfqpoint{5.517814in}{0.634593in}}%
\pgfpathlineto{\pgfqpoint{5.521170in}{0.690324in}}%
\pgfpathlineto{\pgfqpoint{5.524526in}{0.704189in}}%
\pgfpathlineto{\pgfqpoint{5.527882in}{0.704030in}}%
\pgfpathlineto{\pgfqpoint{5.531237in}{0.689577in}}%
\pgfpathlineto{\pgfqpoint{5.534593in}{0.619549in}}%
\pgfpathlineto{\pgfqpoint{5.537949in}{0.685507in}}%
\pgfpathlineto{\pgfqpoint{5.541305in}{0.713275in}}%
\pgfpathlineto{\pgfqpoint{5.544660in}{0.723206in}}%
\pgfpathlineto{\pgfqpoint{5.548016in}{0.718780in}}%
\pgfpathlineto{\pgfqpoint{5.551372in}{0.696856in}}%
\pgfpathlineto{\pgfqpoint{5.554728in}{0.622411in}}%
\pgfpathlineto{\pgfqpoint{5.558084in}{0.686013in}}%
\pgfpathlineto{\pgfqpoint{5.561439in}{0.705181in}}%
\pgfpathlineto{\pgfqpoint{5.564795in}{0.708224in}}%
\pgfpathlineto{\pgfqpoint{5.568151in}{0.698423in}}%
\pgfpathlineto{\pgfqpoint{5.571507in}{0.659611in}}%
\pgfpathlineto{\pgfqpoint{5.574863in}{0.673132in}}%
\pgfpathlineto{\pgfqpoint{5.578218in}{0.710660in}}%
\pgfpathlineto{\pgfqpoint{5.581574in}{0.724399in}}%
\pgfpathlineto{\pgfqpoint{5.584930in}{0.724383in}}%
\pgfpathlineto{\pgfqpoint{5.588286in}{0.706189in}}%
\pgfpathlineto{\pgfqpoint{5.591641in}{0.636732in}}%
\pgfpathlineto{\pgfqpoint{5.594997in}{0.693974in}}%
\pgfpathlineto{\pgfqpoint{5.598353in}{0.716815in}}%
\pgfpathlineto{\pgfqpoint{5.601709in}{0.723297in}}%
\pgfpathlineto{\pgfqpoint{5.605065in}{0.719289in}}%
\pgfpathlineto{\pgfqpoint{5.608420in}{0.700806in}}%
\pgfpathlineto{\pgfqpoint{5.611776in}{0.551570in}}%
\pgfpathlineto{\pgfqpoint{5.615132in}{0.705267in}}%
\pgfpathlineto{\pgfqpoint{5.618488in}{0.729474in}}%
\pgfpathlineto{\pgfqpoint{5.621843in}{0.737955in}}%
\pgfpathlineto{\pgfqpoint{5.625199in}{0.732550in}}%
\pgfpathlineto{\pgfqpoint{5.628555in}{0.708344in}}%
\pgfpathlineto{\pgfqpoint{5.631911in}{0.630998in}}%
\pgfpathlineto{\pgfqpoint{5.635267in}{0.708196in}}%
\pgfpathlineto{\pgfqpoint{5.638622in}{0.725195in}}%
\pgfpathlineto{\pgfqpoint{5.641978in}{0.727758in}}%
\pgfpathlineto{\pgfqpoint{5.645334in}{0.718237in}}%
\pgfpathlineto{\pgfqpoint{5.648690in}{0.682463in}}%
\pgfpathlineto{\pgfqpoint{5.652045in}{0.687386in}}%
\pgfpathlineto{\pgfqpoint{5.655401in}{0.727586in}}%
\pgfpathlineto{\pgfqpoint{5.658757in}{0.742162in}}%
\pgfpathlineto{\pgfqpoint{5.662113in}{0.743176in}}%
\pgfpathlineto{\pgfqpoint{5.665469in}{0.728488in}}%
\pgfpathlineto{\pgfqpoint{5.668824in}{0.693203in}}%
\pgfpathlineto{\pgfqpoint{5.672180in}{0.664879in}}%
\pgfpathlineto{\pgfqpoint{5.675536in}{0.706392in}}%
\pgfpathlineto{\pgfqpoint{5.678892in}{0.714079in}}%
\pgfpathlineto{\pgfqpoint{5.682247in}{0.705287in}}%
\pgfpathlineto{\pgfqpoint{5.685603in}{0.660349in}}%
\pgfpathlineto{\pgfqpoint{5.692315in}{0.723991in}}%
\pgfpathlineto{\pgfqpoint{5.695671in}{0.735997in}}%
\pgfpathlineto{\pgfqpoint{5.699026in}{0.732745in}}%
\pgfpathlineto{\pgfqpoint{5.702382in}{0.714440in}}%
\pgfpathlineto{\pgfqpoint{5.705738in}{0.672226in}}%
\pgfpathlineto{\pgfqpoint{5.709094in}{0.661954in}}%
\pgfpathlineto{\pgfqpoint{5.712449in}{0.692893in}}%
\pgfpathlineto{\pgfqpoint{5.715805in}{0.694917in}}%
\pgfpathlineto{\pgfqpoint{5.719161in}{0.674416in}}%
\pgfpathlineto{\pgfqpoint{5.722517in}{0.636260in}}%
\pgfpathlineto{\pgfqpoint{5.725873in}{0.699929in}}%
\pgfpathlineto{\pgfqpoint{5.729228in}{0.719290in}}%
\pgfpathlineto{\pgfqpoint{5.732584in}{0.722978in}}%
\pgfpathlineto{\pgfqpoint{5.735940in}{0.709499in}}%
\pgfpathlineto{\pgfqpoint{5.739296in}{0.674492in}}%
\pgfpathlineto{\pgfqpoint{5.742652in}{0.647641in}}%
\pgfpathlineto{\pgfqpoint{5.746007in}{0.690094in}}%
\pgfpathlineto{\pgfqpoint{5.749363in}{0.698951in}}%
\pgfpathlineto{\pgfqpoint{5.752719in}{0.692684in}}%
\pgfpathlineto{\pgfqpoint{5.756075in}{0.660938in}}%
\pgfpathlineto{\pgfqpoint{5.759430in}{0.659013in}}%
\pgfpathlineto{\pgfqpoint{5.762786in}{0.702355in}}%
\pgfpathlineto{\pgfqpoint{5.766142in}{0.717511in}}%
\pgfpathlineto{\pgfqpoint{5.769498in}{0.717980in}}%
\pgfpathlineto{\pgfqpoint{5.772854in}{0.699736in}}%
\pgfpathlineto{\pgfqpoint{5.776209in}{0.627389in}}%
\pgfpathlineto{\pgfqpoint{5.779565in}{0.685938in}}%
\pgfpathlineto{\pgfqpoint{5.782921in}{0.708919in}}%
\pgfpathlineto{\pgfqpoint{5.786277in}{0.715010in}}%
\pgfpathlineto{\pgfqpoint{5.789632in}{0.710112in}}%
\pgfpathlineto{\pgfqpoint{5.792988in}{0.688972in}}%
\pgfpathlineto{\pgfqpoint{5.796344in}{0.614837in}}%
\pgfpathlineto{\pgfqpoint{5.799700in}{0.702189in}}%
\pgfpathlineto{\pgfqpoint{5.803056in}{0.723649in}}%
\pgfpathlineto{\pgfqpoint{5.806411in}{0.730236in}}%
\pgfpathlineto{\pgfqpoint{5.809767in}{0.721719in}}%
\pgfpathlineto{\pgfqpoint{5.813123in}{0.688670in}}%
\pgfpathlineto{\pgfqpoint{5.816479in}{0.669855in}}%
\pgfpathlineto{\pgfqpoint{5.819834in}{0.713483in}}%
\pgfpathlineto{\pgfqpoint{5.823190in}{0.726172in}}%
\pgfpathlineto{\pgfqpoint{5.826546in}{0.727026in}}%
\pgfpathlineto{\pgfqpoint{5.829902in}{0.716443in}}%
\pgfpathlineto{\pgfqpoint{5.833258in}{0.678992in}}%
\pgfpathlineto{\pgfqpoint{5.836613in}{0.686907in}}%
\pgfpathlineto{\pgfqpoint{5.839969in}{0.725913in}}%
\pgfpathlineto{\pgfqpoint{5.843325in}{0.740236in}}%
\pgfpathlineto{\pgfqpoint{5.846681in}{0.741367in}}%
\pgfpathlineto{\pgfqpoint{5.850036in}{0.726437in}}%
\pgfpathlineto{\pgfqpoint{5.853392in}{0.685575in}}%
\pgfpathlineto{\pgfqpoint{5.856748in}{0.680451in}}%
\pgfpathlineto{\pgfqpoint{5.860104in}{0.714607in}}%
\pgfpathlineto{\pgfqpoint{5.863460in}{0.722110in}}%
\pgfpathlineto{\pgfqpoint{5.866815in}{0.715561in}}%
\pgfpathlineto{\pgfqpoint{5.870171in}{0.684919in}}%
\pgfpathlineto{\pgfqpoint{5.873527in}{0.677997in}}%
\pgfpathlineto{\pgfqpoint{5.876883in}{0.723214in}}%
\pgfpathlineto{\pgfqpoint{5.880238in}{0.738856in}}%
\pgfpathlineto{\pgfqpoint{5.883594in}{0.739982in}}%
\pgfpathlineto{\pgfqpoint{5.886950in}{0.725268in}}%
\pgfpathlineto{\pgfqpoint{5.890306in}{0.692944in}}%
\pgfpathlineto{\pgfqpoint{5.893662in}{0.635633in}}%
\pgfpathlineto{\pgfqpoint{5.897017in}{0.692390in}}%
\pgfpathlineto{\pgfqpoint{5.900373in}{0.699737in}}%
\pgfpathlineto{\pgfqpoint{5.903729in}{0.686116in}}%
\pgfpathlineto{\pgfqpoint{5.907085in}{0.545809in}}%
\pgfpathlineto{\pgfqpoint{5.910441in}{0.696408in}}%
\pgfpathlineto{\pgfqpoint{5.913796in}{0.720546in}}%
\pgfpathlineto{\pgfqpoint{5.917152in}{0.727531in}}%
\pgfpathlineto{\pgfqpoint{5.920508in}{0.717789in}}%
\pgfpathlineto{\pgfqpoint{5.923864in}{0.692026in}}%
\pgfpathlineto{\pgfqpoint{5.927219in}{0.585188in}}%
\pgfpathlineto{\pgfqpoint{5.930575in}{0.677094in}}%
\pgfpathlineto{\pgfqpoint{5.933931in}{0.691477in}}%
\pgfpathlineto{\pgfqpoint{5.937287in}{0.686370in}}%
\pgfpathlineto{\pgfqpoint{5.940643in}{0.651130in}}%
\pgfpathlineto{\pgfqpoint{5.943998in}{0.662685in}}%
\pgfpathlineto{\pgfqpoint{5.947354in}{0.701754in}}%
\pgfpathlineto{\pgfqpoint{5.950710in}{0.715553in}}%
\pgfpathlineto{\pgfqpoint{5.954066in}{0.713839in}}%
\pgfpathlineto{\pgfqpoint{5.957421in}{0.693277in}}%
\pgfpathlineto{\pgfqpoint{5.960777in}{0.607537in}}%
\pgfpathlineto{\pgfqpoint{5.964133in}{0.680665in}}%
\pgfpathlineto{\pgfqpoint{5.967489in}{0.701419in}}%
\pgfpathlineto{\pgfqpoint{5.970845in}{0.705419in}}%
\pgfpathlineto{\pgfqpoint{5.974200in}{0.696998in}}%
\pgfpathlineto{\pgfqpoint{5.977556in}{0.663906in}}%
\pgfpathlineto{\pgfqpoint{5.980912in}{0.661961in}}%
\pgfpathlineto{\pgfqpoint{5.984268in}{0.704771in}}%
\pgfpathlineto{\pgfqpoint{5.987623in}{0.719907in}}%
\pgfpathlineto{\pgfqpoint{5.990979in}{0.721073in}}%
\pgfpathlineto{\pgfqpoint{5.994335in}{0.703672in}}%
\pgfpathlineto{\pgfqpoint{5.997691in}{0.622906in}}%
\pgfpathlineto{\pgfqpoint{6.001047in}{0.693848in}}%
\pgfpathlineto{\pgfqpoint{6.004402in}{0.716977in}}%
\pgfpathlineto{\pgfqpoint{6.007758in}{0.723935in}}%
\pgfpathlineto{\pgfqpoint{6.011114in}{0.720754in}}%
\pgfpathlineto{\pgfqpoint{6.014470in}{0.704379in}}%
\pgfpathlineto{\pgfqpoint{6.017825in}{0.627409in}}%
\pgfpathlineto{\pgfqpoint{6.021181in}{0.700614in}}%
\pgfpathlineto{\pgfqpoint{6.024537in}{0.727599in}}%
\pgfpathlineto{\pgfqpoint{6.027893in}{0.737496in}}%
\pgfpathlineto{\pgfqpoint{6.031249in}{0.733948in}}%
\pgfpathlineto{\pgfqpoint{6.034604in}{0.711959in}}%
\pgfpathlineto{\pgfqpoint{6.037960in}{0.584584in}}%
\pgfpathlineto{\pgfqpoint{6.041316in}{0.706114in}}%
\pgfpathlineto{\pgfqpoint{6.044672in}{0.725158in}}%
\pgfpathlineto{\pgfqpoint{6.048027in}{0.728731in}}%
\pgfpathlineto{\pgfqpoint{6.051383in}{0.720330in}}%
\pgfpathlineto{\pgfqpoint{6.054739in}{0.688320in}}%
\pgfpathlineto{\pgfqpoint{6.058095in}{0.682281in}}%
\pgfpathlineto{\pgfqpoint{6.061451in}{0.726679in}}%
\pgfpathlineto{\pgfqpoint{6.064806in}{0.742272in}}%
\pgfpathlineto{\pgfqpoint{6.068162in}{0.744043in}}%
\pgfpathlineto{\pgfqpoint{6.071518in}{0.730117in}}%
\pgfpathlineto{\pgfqpoint{6.074874in}{0.697554in}}%
\pgfpathlineto{\pgfqpoint{6.078230in}{0.647294in}}%
\pgfpathlineto{\pgfqpoint{6.081585in}{0.701373in}}%
\pgfpathlineto{\pgfqpoint{6.084941in}{0.709804in}}%
\pgfpathlineto{\pgfqpoint{6.088297in}{0.699830in}}%
\pgfpathlineto{\pgfqpoint{6.091653in}{0.641270in}}%
\pgfpathlineto{\pgfqpoint{6.095008in}{0.694553in}}%
\pgfpathlineto{\pgfqpoint{6.098364in}{0.724245in}}%
\pgfpathlineto{\pgfqpoint{6.101720in}{0.734323in}}%
\pgfpathlineto{\pgfqpoint{6.105076in}{0.728352in}}%
\pgfpathlineto{\pgfqpoint{6.108432in}{0.707996in}}%
\pgfpathlineto{\pgfqpoint{6.111787in}{0.659116in}}%
\pgfpathlineto{\pgfqpoint{6.115143in}{0.661859in}}%
\pgfpathlineto{\pgfqpoint{6.118499in}{0.687406in}}%
\pgfpathlineto{\pgfqpoint{6.121855in}{0.685231in}}%
\pgfpathlineto{\pgfqpoint{6.125210in}{0.650978in}}%
\pgfpathlineto{\pgfqpoint{6.128566in}{0.664299in}}%
\pgfpathlineto{\pgfqpoint{6.131922in}{0.703506in}}%
\pgfpathlineto{\pgfqpoint{6.135278in}{0.717398in}}%
\pgfpathlineto{\pgfqpoint{6.138634in}{0.715355in}}%
\pgfpathlineto{\pgfqpoint{6.141989in}{0.695781in}}%
\pgfpathlineto{\pgfqpoint{6.145345in}{0.635926in}}%
\pgfpathlineto{\pgfqpoint{6.148701in}{0.668748in}}%
\pgfpathlineto{\pgfqpoint{6.152057in}{0.692235in}}%
\pgfpathlineto{\pgfqpoint{6.155412in}{0.695350in}}%
\pgfpathlineto{\pgfqpoint{6.158768in}{0.683029in}}%
\pgfpathlineto{\pgfqpoint{6.162124in}{0.621592in}}%
\pgfpathlineto{\pgfqpoint{6.165480in}{0.676757in}}%
\pgfpathlineto{\pgfqpoint{6.168836in}{0.705804in}}%
\pgfpathlineto{\pgfqpoint{6.172191in}{0.715784in}}%
\pgfpathlineto{\pgfqpoint{6.175547in}{0.710145in}}%
\pgfpathlineto{\pgfqpoint{6.178903in}{0.681264in}}%
\pgfpathlineto{\pgfqpoint{6.182259in}{0.646124in}}%
\pgfpathlineto{\pgfqpoint{6.185614in}{0.699132in}}%
\pgfpathlineto{\pgfqpoint{6.188970in}{0.713759in}}%
\pgfpathlineto{\pgfqpoint{6.192326in}{0.716139in}}%
\pgfpathlineto{\pgfqpoint{6.195682in}{0.707868in}}%
\pgfpathlineto{\pgfqpoint{6.199038in}{0.679151in}}%
\pgfpathlineto{\pgfqpoint{6.202393in}{0.658236in}}%
\pgfpathlineto{\pgfqpoint{6.205749in}{0.709334in}}%
\pgfpathlineto{\pgfqpoint{6.209105in}{0.726597in}}%
\pgfpathlineto{\pgfqpoint{6.212461in}{0.730409in}}%
\pgfpathlineto{\pgfqpoint{6.215816in}{0.718075in}}%
\pgfpathlineto{\pgfqpoint{6.219172in}{0.671479in}}%
\pgfpathlineto{\pgfqpoint{6.222528in}{0.689841in}}%
\pgfpathlineto{\pgfqpoint{6.225884in}{0.720850in}}%
\pgfpathlineto{\pgfqpoint{6.229240in}{0.730609in}}%
\pgfpathlineto{\pgfqpoint{6.232595in}{0.729682in}}%
\pgfpathlineto{\pgfqpoint{6.235951in}{0.716951in}}%
\pgfpathlineto{\pgfqpoint{6.239307in}{0.670371in}}%
\pgfpathlineto{\pgfqpoint{6.242663in}{0.697846in}}%
\pgfpathlineto{\pgfqpoint{6.246019in}{0.731383in}}%
\pgfpathlineto{\pgfqpoint{6.249374in}{0.743968in}}%
\pgfpathlineto{\pgfqpoint{6.252730in}{0.743494in}}%
\pgfpathlineto{\pgfqpoint{6.256086in}{0.726962in}}%
\pgfpathlineto{\pgfqpoint{6.259442in}{0.682048in}}%
\pgfpathlineto{\pgfqpoint{6.262797in}{0.685551in}}%
\pgfpathlineto{\pgfqpoint{6.266153in}{0.716270in}}%
\pgfpathlineto{\pgfqpoint{6.269509in}{0.722096in}}%
\pgfpathlineto{\pgfqpoint{6.272865in}{0.713037in}}%
\pgfpathlineto{\pgfqpoint{6.276221in}{0.671386in}}%
\pgfpathlineto{\pgfqpoint{6.279576in}{0.693286in}}%
\pgfpathlineto{\pgfqpoint{6.282932in}{0.728671in}}%
\pgfpathlineto{\pgfqpoint{6.286288in}{0.741274in}}%
\pgfpathlineto{\pgfqpoint{6.289644in}{0.738879in}}%
\pgfpathlineto{\pgfqpoint{6.292999in}{0.721718in}}%
\pgfpathlineto{\pgfqpoint{6.299711in}{0.644930in}}%
\pgfpathlineto{\pgfqpoint{6.303067in}{0.689119in}}%
\pgfpathlineto{\pgfqpoint{6.306423in}{0.691618in}}%
\pgfpathlineto{\pgfqpoint{6.309778in}{0.665907in}}%
\pgfpathlineto{\pgfqpoint{6.313134in}{0.658040in}}%
\pgfpathlineto{\pgfqpoint{6.316490in}{0.705889in}}%
\pgfpathlineto{\pgfqpoint{6.319846in}{0.722321in}}%
\pgfpathlineto{\pgfqpoint{6.323201in}{0.722948in}}%
\pgfpathlineto{\pgfqpoint{6.326557in}{0.706848in}}%
\pgfpathlineto{\pgfqpoint{6.329913in}{0.669986in}}%
\pgfpathlineto{\pgfqpoint{6.333269in}{0.639930in}}%
\pgfpathlineto{\pgfqpoint{6.336625in}{0.681406in}}%
\pgfpathlineto{\pgfqpoint{6.339980in}{0.686742in}}%
\pgfpathlineto{\pgfqpoint{6.343336in}{0.672093in}}%
\pgfpathlineto{\pgfqpoint{6.346692in}{0.535297in}}%
\pgfpathlineto{\pgfqpoint{6.350048in}{0.682737in}}%
\pgfpathlineto{\pgfqpoint{6.353403in}{0.706418in}}%
\pgfpathlineto{\pgfqpoint{6.356759in}{0.713118in}}%
\pgfpathlineto{\pgfqpoint{6.360115in}{0.704247in}}%
\pgfpathlineto{\pgfqpoint{6.366827in}{0.674755in}}%
\pgfpathlineto{\pgfqpoint{6.370182in}{0.659920in}}%
\pgfpathlineto{\pgfqpoint{6.373538in}{0.650719in}}%
\pgfpathlineto{\pgfqpoint{6.376894in}{0.649632in}}%
\pgfpathlineto{\pgfqpoint{6.380250in}{0.651686in}}%
\pgfpathlineto{\pgfqpoint{6.383605in}{0.647765in}}%
\pgfpathlineto{\pgfqpoint{6.383605in}{0.647765in}}%
\pgfusepath{stroke}%
\end{pgfscope}%
\begin{pgfscope}%
\pgfpathrectangle{\pgfqpoint{4.621821in}{0.320625in}}{\pgfqpoint{1.845679in}{1.164375in}}%
\pgfusepath{clip}%
\pgfsetrectcap%
\pgfsetroundjoin%
\pgfsetlinewidth{0.602250pt}%
\definecolor{currentstroke}{rgb}{0.819608,0.286275,0.356863}%
\pgfsetstrokecolor{currentstroke}%
\pgfsetstrokeopacity{0.120000}%
\pgfsetdash{}{0pt}%
\pgfpathmoveto{\pgfqpoint{4.705715in}{1.472285in}}%
\pgfpathlineto{\pgfqpoint{4.732562in}{1.424217in}}%
\pgfpathlineto{\pgfqpoint{4.769475in}{1.354185in}}%
\pgfpathlineto{\pgfqpoint{4.786254in}{1.325127in}}%
\pgfpathlineto{\pgfqpoint{4.803033in}{1.298083in}}%
\pgfpathlineto{\pgfqpoint{4.823168in}{1.267905in}}%
\pgfpathlineto{\pgfqpoint{4.843302in}{1.239770in}}%
\pgfpathlineto{\pgfqpoint{4.866793in}{1.209015in}}%
\pgfpathlineto{\pgfqpoint{4.890283in}{1.180023in}}%
\pgfpathlineto{\pgfqpoint{4.920485in}{1.144678in}}%
\pgfpathlineto{\pgfqpoint{4.954043in}{1.107162in}}%
\pgfpathlineto{\pgfqpoint{5.001024in}{1.056484in}}%
\pgfpathlineto{\pgfqpoint{5.068140in}{0.985916in}}%
\pgfpathlineto{\pgfqpoint{5.199015in}{0.850392in}}%
\pgfpathlineto{\pgfqpoint{5.269487in}{0.779158in}}%
\pgfpathlineto{\pgfqpoint{5.279554in}{0.769200in}}%
\pgfpathlineto{\pgfqpoint{5.292977in}{0.754016in}}%
\pgfpathlineto{\pgfqpoint{5.309756in}{0.734630in}}%
\pgfpathlineto{\pgfqpoint{5.316467in}{0.727959in}}%
\pgfpathlineto{\pgfqpoint{5.323179in}{0.722494in}}%
\pgfpathlineto{\pgfqpoint{5.329891in}{0.718274in}}%
\pgfpathlineto{\pgfqpoint{5.343314in}{0.710919in}}%
\pgfpathlineto{\pgfqpoint{5.346669in}{0.708226in}}%
\pgfpathlineto{\pgfqpoint{5.353381in}{0.701229in}}%
\pgfpathlineto{\pgfqpoint{5.356737in}{0.699461in}}%
\pgfpathlineto{\pgfqpoint{5.360093in}{0.701267in}}%
\pgfpathlineto{\pgfqpoint{5.363448in}{0.706223in}}%
\pgfpathlineto{\pgfqpoint{5.370160in}{0.717803in}}%
\pgfpathlineto{\pgfqpoint{5.373516in}{0.722257in}}%
\pgfpathlineto{\pgfqpoint{5.376871in}{0.725430in}}%
\pgfpathlineto{\pgfqpoint{5.380227in}{0.727282in}}%
\pgfpathlineto{\pgfqpoint{5.383583in}{0.727784in}}%
\pgfpathlineto{\pgfqpoint{5.386939in}{0.726854in}}%
\pgfpathlineto{\pgfqpoint{5.390295in}{0.724310in}}%
\pgfpathlineto{\pgfqpoint{5.393650in}{0.719789in}}%
\pgfpathlineto{\pgfqpoint{5.397006in}{0.712575in}}%
\pgfpathlineto{\pgfqpoint{5.400362in}{0.701087in}}%
\pgfpathlineto{\pgfqpoint{5.403718in}{0.680907in}}%
\pgfpathlineto{\pgfqpoint{5.407074in}{0.634232in}}%
\pgfpathlineto{\pgfqpoint{5.410429in}{0.670074in}}%
\pgfpathlineto{\pgfqpoint{5.413785in}{0.699052in}}%
\pgfpathlineto{\pgfqpoint{5.417141in}{0.715258in}}%
\pgfpathlineto{\pgfqpoint{5.420497in}{0.725975in}}%
\pgfpathlineto{\pgfqpoint{5.423852in}{0.733390in}}%
\pgfpathlineto{\pgfqpoint{5.427208in}{0.738359in}}%
\pgfpathlineto{\pgfqpoint{5.430564in}{0.741194in}}%
\pgfpathlineto{\pgfqpoint{5.433920in}{0.741859in}}%
\pgfpathlineto{\pgfqpoint{5.437276in}{0.739944in}}%
\pgfpathlineto{\pgfqpoint{5.440631in}{0.734360in}}%
\pgfpathlineto{\pgfqpoint{5.443987in}{0.723071in}}%
\pgfpathlineto{\pgfqpoint{5.447343in}{0.706114in}}%
\pgfpathlineto{\pgfqpoint{5.454054in}{0.659079in}}%
\pgfpathlineto{\pgfqpoint{5.457410in}{0.648569in}}%
\pgfpathlineto{\pgfqpoint{5.460766in}{0.649398in}}%
\pgfpathlineto{\pgfqpoint{5.467478in}{0.685206in}}%
\pgfpathlineto{\pgfqpoint{5.470833in}{0.692053in}}%
\pgfpathlineto{\pgfqpoint{5.474189in}{0.684024in}}%
\pgfpathlineto{\pgfqpoint{5.477545in}{0.668159in}}%
\pgfpathlineto{\pgfqpoint{5.484256in}{0.624212in}}%
\pgfpathlineto{\pgfqpoint{5.487612in}{0.614106in}}%
\pgfpathlineto{\pgfqpoint{5.490968in}{0.612250in}}%
\pgfpathlineto{\pgfqpoint{5.494324in}{0.625087in}}%
\pgfpathlineto{\pgfqpoint{5.497680in}{0.642481in}}%
\pgfpathlineto{\pgfqpoint{5.501035in}{0.651107in}}%
\pgfpathlineto{\pgfqpoint{5.504391in}{0.644094in}}%
\pgfpathlineto{\pgfqpoint{5.507747in}{0.593606in}}%
\pgfpathlineto{\pgfqpoint{5.511103in}{0.659692in}}%
\pgfpathlineto{\pgfqpoint{5.514458in}{0.685311in}}%
\pgfpathlineto{\pgfqpoint{5.517814in}{0.698060in}}%
\pgfpathlineto{\pgfqpoint{5.521170in}{0.703877in}}%
\pgfpathlineto{\pgfqpoint{5.524526in}{0.704295in}}%
\pgfpathlineto{\pgfqpoint{5.527882in}{0.699020in}}%
\pgfpathlineto{\pgfqpoint{5.531237in}{0.685245in}}%
\pgfpathlineto{\pgfqpoint{5.534593in}{0.648066in}}%
\pgfpathlineto{\pgfqpoint{5.537949in}{0.648984in}}%
\pgfpathlineto{\pgfqpoint{5.541305in}{0.690222in}}%
\pgfpathlineto{\pgfqpoint{5.544660in}{0.707126in}}%
\pgfpathlineto{\pgfqpoint{5.548016in}{0.714348in}}%
\pgfpathlineto{\pgfqpoint{5.551372in}{0.713008in}}%
\pgfpathlineto{\pgfqpoint{5.554728in}{0.698437in}}%
\pgfpathlineto{\pgfqpoint{5.558084in}{0.653357in}}%
\pgfpathlineto{\pgfqpoint{5.561439in}{0.665806in}}%
\pgfpathlineto{\pgfqpoint{5.564795in}{0.696333in}}%
\pgfpathlineto{\pgfqpoint{5.568151in}{0.705024in}}%
\pgfpathlineto{\pgfqpoint{5.571507in}{0.702247in}}%
\pgfpathlineto{\pgfqpoint{5.574863in}{0.685074in}}%
\pgfpathlineto{\pgfqpoint{5.578218in}{0.580664in}}%
\pgfpathlineto{\pgfqpoint{5.581574in}{0.687008in}}%
\pgfpathlineto{\pgfqpoint{5.584930in}{0.712207in}}%
\pgfpathlineto{\pgfqpoint{5.588286in}{0.721047in}}%
\pgfpathlineto{\pgfqpoint{5.591641in}{0.715728in}}%
\pgfpathlineto{\pgfqpoint{5.594997in}{0.691029in}}%
\pgfpathlineto{\pgfqpoint{5.598353in}{0.621503in}}%
\pgfpathlineto{\pgfqpoint{5.601709in}{0.693927in}}%
\pgfpathlineto{\pgfqpoint{5.605065in}{0.710662in}}%
\pgfpathlineto{\pgfqpoint{5.608420in}{0.713499in}}%
\pgfpathlineto{\pgfqpoint{5.611776in}{0.704854in}}%
\pgfpathlineto{\pgfqpoint{5.615132in}{0.673269in}}%
\pgfpathlineto{\pgfqpoint{5.618488in}{0.664972in}}%
\pgfpathlineto{\pgfqpoint{5.621843in}{0.710304in}}%
\pgfpathlineto{\pgfqpoint{5.625199in}{0.726318in}}%
\pgfpathlineto{\pgfqpoint{5.628555in}{0.728909in}}%
\pgfpathlineto{\pgfqpoint{5.631911in}{0.714594in}}%
\pgfpathlineto{\pgfqpoint{5.635267in}{0.661830in}}%
\pgfpathlineto{\pgfqpoint{5.641978in}{0.720768in}}%
\pgfpathlineto{\pgfqpoint{5.645334in}{0.729539in}}%
\pgfpathlineto{\pgfqpoint{5.648690in}{0.728000in}}%
\pgfpathlineto{\pgfqpoint{5.652045in}{0.714556in}}%
\pgfpathlineto{\pgfqpoint{5.655401in}{0.664517in}}%
\pgfpathlineto{\pgfqpoint{5.662113in}{0.730500in}}%
\pgfpathlineto{\pgfqpoint{5.665469in}{0.742751in}}%
\pgfpathlineto{\pgfqpoint{5.668824in}{0.742204in}}%
\pgfpathlineto{\pgfqpoint{5.672180in}{0.725266in}}%
\pgfpathlineto{\pgfqpoint{5.675536in}{0.674962in}}%
\pgfpathlineto{\pgfqpoint{5.678892in}{0.695834in}}%
\pgfpathlineto{\pgfqpoint{5.682247in}{0.722650in}}%
\pgfpathlineto{\pgfqpoint{5.685603in}{0.728847in}}%
\pgfpathlineto{\pgfqpoint{5.688959in}{0.722308in}}%
\pgfpathlineto{\pgfqpoint{5.692315in}{0.694147in}}%
\pgfpathlineto{\pgfqpoint{5.695671in}{0.676025in}}%
\pgfpathlineto{\pgfqpoint{5.699026in}{0.726232in}}%
\pgfpathlineto{\pgfqpoint{5.702382in}{0.743138in}}%
\pgfpathlineto{\pgfqpoint{5.705738in}{0.745871in}}%
\pgfpathlineto{\pgfqpoint{5.709094in}{0.733039in}}%
\pgfpathlineto{\pgfqpoint{5.712449in}{0.704479in}}%
\pgfpathlineto{\pgfqpoint{5.715805in}{0.632247in}}%
\pgfpathlineto{\pgfqpoint{5.719161in}{0.694995in}}%
\pgfpathlineto{\pgfqpoint{5.722517in}{0.705295in}}%
\pgfpathlineto{\pgfqpoint{5.725873in}{0.695225in}}%
\pgfpathlineto{\pgfqpoint{5.729228in}{0.627417in}}%
\pgfpathlineto{\pgfqpoint{5.732584in}{0.695145in}}%
\pgfpathlineto{\pgfqpoint{5.735940in}{0.723245in}}%
\pgfpathlineto{\pgfqpoint{5.739296in}{0.732578in}}%
\pgfpathlineto{\pgfqpoint{5.742652in}{0.725645in}}%
\pgfpathlineto{\pgfqpoint{5.746007in}{0.704280in}}%
\pgfpathlineto{\pgfqpoint{5.749363in}{0.651390in}}%
\pgfpathlineto{\pgfqpoint{5.752719in}{0.667045in}}%
\pgfpathlineto{\pgfqpoint{5.756075in}{0.688365in}}%
\pgfpathlineto{\pgfqpoint{5.759430in}{0.685319in}}%
\pgfpathlineto{\pgfqpoint{5.762786in}{0.651045in}}%
\pgfpathlineto{\pgfqpoint{5.766142in}{0.663517in}}%
\pgfpathlineto{\pgfqpoint{5.769498in}{0.702699in}}%
\pgfpathlineto{\pgfqpoint{5.772854in}{0.716606in}}%
\pgfpathlineto{\pgfqpoint{5.776209in}{0.714734in}}%
\pgfpathlineto{\pgfqpoint{5.779565in}{0.694729in}}%
\pgfpathlineto{\pgfqpoint{5.782921in}{0.628869in}}%
\pgfpathlineto{\pgfqpoint{5.786277in}{0.675319in}}%
\pgfpathlineto{\pgfqpoint{5.789632in}{0.696992in}}%
\pgfpathlineto{\pgfqpoint{5.792988in}{0.700466in}}%
\pgfpathlineto{\pgfqpoint{5.796344in}{0.690133in}}%
\pgfpathlineto{\pgfqpoint{5.799700in}{0.646293in}}%
\pgfpathlineto{\pgfqpoint{5.803056in}{0.671392in}}%
\pgfpathlineto{\pgfqpoint{5.806411in}{0.705763in}}%
\pgfpathlineto{\pgfqpoint{5.809767in}{0.718109in}}%
\pgfpathlineto{\pgfqpoint{5.813123in}{0.715823in}}%
\pgfpathlineto{\pgfqpoint{5.816479in}{0.692585in}}%
\pgfpathlineto{\pgfqpoint{5.819834in}{0.624287in}}%
\pgfpathlineto{\pgfqpoint{5.823190in}{0.699131in}}%
\pgfpathlineto{\pgfqpoint{5.826546in}{0.716879in}}%
\pgfpathlineto{\pgfqpoint{5.829902in}{0.721274in}}%
\pgfpathlineto{\pgfqpoint{5.833258in}{0.715551in}}%
\pgfpathlineto{\pgfqpoint{5.836613in}{0.693926in}}%
\pgfpathlineto{\pgfqpoint{5.839969in}{0.620019in}}%
\pgfpathlineto{\pgfqpoint{5.843325in}{0.706862in}}%
\pgfpathlineto{\pgfqpoint{5.846681in}{0.728321in}}%
\pgfpathlineto{\pgfqpoint{5.850036in}{0.735165in}}%
\pgfpathlineto{\pgfqpoint{5.853392in}{0.727479in}}%
\pgfpathlineto{\pgfqpoint{5.860104in}{0.668295in}}%
\pgfpathlineto{\pgfqpoint{5.863460in}{0.716268in}}%
\pgfpathlineto{\pgfqpoint{5.866815in}{0.729597in}}%
\pgfpathlineto{\pgfqpoint{5.870171in}{0.730587in}}%
\pgfpathlineto{\pgfqpoint{5.873527in}{0.719796in}}%
\pgfpathlineto{\pgfqpoint{5.876883in}{0.680585in}}%
\pgfpathlineto{\pgfqpoint{5.880238in}{0.693186in}}%
\pgfpathlineto{\pgfqpoint{5.883594in}{0.730696in}}%
\pgfpathlineto{\pgfqpoint{5.886950in}{0.744435in}}%
\pgfpathlineto{\pgfqpoint{5.890306in}{0.744700in}}%
\pgfpathlineto{\pgfqpoint{5.893662in}{0.729149in}}%
\pgfpathlineto{\pgfqpoint{5.897017in}{0.691159in}}%
\pgfpathlineto{\pgfqpoint{5.900373in}{0.671763in}}%
\pgfpathlineto{\pgfqpoint{5.903729in}{0.709159in}}%
\pgfpathlineto{\pgfqpoint{5.907085in}{0.715306in}}%
\pgfpathlineto{\pgfqpoint{5.910441in}{0.704203in}}%
\pgfpathlineto{\pgfqpoint{5.913796in}{0.642443in}}%
\pgfpathlineto{\pgfqpoint{5.917152in}{0.699563in}}%
\pgfpathlineto{\pgfqpoint{5.920508in}{0.728517in}}%
\pgfpathlineto{\pgfqpoint{5.923864in}{0.738197in}}%
\pgfpathlineto{\pgfqpoint{5.927219in}{0.731830in}}%
\pgfpathlineto{\pgfqpoint{5.930575in}{0.711555in}}%
\pgfpathlineto{\pgfqpoint{5.933931in}{0.665349in}}%
\pgfpathlineto{\pgfqpoint{5.937287in}{0.660209in}}%
\pgfpathlineto{\pgfqpoint{5.940643in}{0.686469in}}%
\pgfpathlineto{\pgfqpoint{5.943998in}{0.682507in}}%
\pgfpathlineto{\pgfqpoint{5.947354in}{0.635048in}}%
\pgfpathlineto{\pgfqpoint{5.950710in}{0.675863in}}%
\pgfpathlineto{\pgfqpoint{5.954066in}{0.708380in}}%
\pgfpathlineto{\pgfqpoint{5.957421in}{0.719643in}}%
\pgfpathlineto{\pgfqpoint{5.960777in}{0.714195in}}%
\pgfpathlineto{\pgfqpoint{5.964133in}{0.692116in}}%
\pgfpathlineto{\pgfqpoint{5.967489in}{0.622704in}}%
\pgfpathlineto{\pgfqpoint{5.970845in}{0.668196in}}%
\pgfpathlineto{\pgfqpoint{5.974200in}{0.687628in}}%
\pgfpathlineto{\pgfqpoint{5.977556in}{0.686796in}}%
\pgfpathlineto{\pgfqpoint{5.980912in}{0.664882in}}%
\pgfpathlineto{\pgfqpoint{5.984268in}{0.626615in}}%
\pgfpathlineto{\pgfqpoint{5.987623in}{0.689512in}}%
\pgfpathlineto{\pgfqpoint{5.990979in}{0.708429in}}%
\pgfpathlineto{\pgfqpoint{5.994335in}{0.711684in}}%
\pgfpathlineto{\pgfqpoint{5.997691in}{0.696229in}}%
\pgfpathlineto{\pgfqpoint{6.001047in}{0.637106in}}%
\pgfpathlineto{\pgfqpoint{6.004402in}{0.676645in}}%
\pgfpathlineto{\pgfqpoint{6.007758in}{0.702624in}}%
\pgfpathlineto{\pgfqpoint{6.011114in}{0.710107in}}%
\pgfpathlineto{\pgfqpoint{6.014470in}{0.706674in}}%
\pgfpathlineto{\pgfqpoint{6.017825in}{0.688774in}}%
\pgfpathlineto{\pgfqpoint{6.021181in}{0.562724in}}%
\pgfpathlineto{\pgfqpoint{6.024537in}{0.692073in}}%
\pgfpathlineto{\pgfqpoint{6.027893in}{0.716549in}}%
\pgfpathlineto{\pgfqpoint{6.031249in}{0.724838in}}%
\pgfpathlineto{\pgfqpoint{6.034604in}{0.718464in}}%
\pgfpathlineto{\pgfqpoint{6.037960in}{0.686624in}}%
\pgfpathlineto{\pgfqpoint{6.041316in}{0.669231in}}%
\pgfpathlineto{\pgfqpoint{6.044672in}{0.713745in}}%
\pgfpathlineto{\pgfqpoint{6.048027in}{0.727438in}}%
\pgfpathlineto{\pgfqpoint{6.051383in}{0.729730in}}%
\pgfpathlineto{\pgfqpoint{6.054739in}{0.721925in}}%
\pgfpathlineto{\pgfqpoint{6.058095in}{0.695525in}}%
\pgfpathlineto{\pgfqpoint{6.061451in}{0.661855in}}%
\pgfpathlineto{\pgfqpoint{6.064806in}{0.719913in}}%
\pgfpathlineto{\pgfqpoint{6.068162in}{0.738417in}}%
\pgfpathlineto{\pgfqpoint{6.071518in}{0.743342in}}%
\pgfpathlineto{\pgfqpoint{6.074874in}{0.733350in}}%
\pgfpathlineto{\pgfqpoint{6.078230in}{0.702281in}}%
\pgfpathlineto{\pgfqpoint{6.081585in}{0.662367in}}%
\pgfpathlineto{\pgfqpoint{6.084941in}{0.714446in}}%
\pgfpathlineto{\pgfqpoint{6.088297in}{0.726320in}}%
\pgfpathlineto{\pgfqpoint{6.091653in}{0.724073in}}%
\pgfpathlineto{\pgfqpoint{6.095008in}{0.704910in}}%
\pgfpathlineto{\pgfqpoint{6.098364in}{0.618191in}}%
\pgfpathlineto{\pgfqpoint{6.101720in}{0.716938in}}%
\pgfpathlineto{\pgfqpoint{6.105076in}{0.738881in}}%
\pgfpathlineto{\pgfqpoint{6.108432in}{0.744975in}}%
\pgfpathlineto{\pgfqpoint{6.111787in}{0.735329in}}%
\pgfpathlineto{\pgfqpoint{6.115143in}{0.711757in}}%
\pgfpathlineto{\pgfqpoint{6.118499in}{0.646573in}}%
\pgfpathlineto{\pgfqpoint{6.121855in}{0.678794in}}%
\pgfpathlineto{\pgfqpoint{6.125210in}{0.695947in}}%
\pgfpathlineto{\pgfqpoint{6.128566in}{0.687433in}}%
\pgfpathlineto{\pgfqpoint{6.131922in}{0.610394in}}%
\pgfpathlineto{\pgfqpoint{6.135278in}{0.691392in}}%
\pgfpathlineto{\pgfqpoint{6.138634in}{0.718796in}}%
\pgfpathlineto{\pgfqpoint{6.141989in}{0.727455in}}%
\pgfpathlineto{\pgfqpoint{6.145345in}{0.719147in}}%
\pgfpathlineto{\pgfqpoint{6.148701in}{0.696436in}}%
\pgfpathlineto{\pgfqpoint{6.152057in}{0.635727in}}%
\pgfpathlineto{\pgfqpoint{6.155412in}{0.661112in}}%
\pgfpathlineto{\pgfqpoint{6.158768in}{0.680124in}}%
\pgfpathlineto{\pgfqpoint{6.162124in}{0.673588in}}%
\pgfpathlineto{\pgfqpoint{6.165480in}{0.617585in}}%
\pgfpathlineto{\pgfqpoint{6.168836in}{0.670333in}}%
\pgfpathlineto{\pgfqpoint{6.172191in}{0.700659in}}%
\pgfpathlineto{\pgfqpoint{6.175547in}{0.710831in}}%
\pgfpathlineto{\pgfqpoint{6.178903in}{0.704000in}}%
\pgfpathlineto{\pgfqpoint{6.182259in}{0.676398in}}%
\pgfpathlineto{\pgfqpoint{6.185614in}{0.615536in}}%
\pgfpathlineto{\pgfqpoint{6.188970in}{0.682346in}}%
\pgfpathlineto{\pgfqpoint{6.192326in}{0.696532in}}%
\pgfpathlineto{\pgfqpoint{6.195682in}{0.696395in}}%
\pgfpathlineto{\pgfqpoint{6.199038in}{0.681679in}}%
\pgfpathlineto{\pgfqpoint{6.202393in}{0.606707in}}%
\pgfpathlineto{\pgfqpoint{6.205749in}{0.678875in}}%
\pgfpathlineto{\pgfqpoint{6.209105in}{0.705912in}}%
\pgfpathlineto{\pgfqpoint{6.212461in}{0.715053in}}%
\pgfpathlineto{\pgfqpoint{6.215816in}{0.708724in}}%
\pgfpathlineto{\pgfqpoint{6.219172in}{0.675009in}}%
\pgfpathlineto{\pgfqpoint{6.222528in}{0.665925in}}%
\pgfpathlineto{\pgfqpoint{6.225884in}{0.707331in}}%
\pgfpathlineto{\pgfqpoint{6.229240in}{0.720686in}}%
\pgfpathlineto{\pgfqpoint{6.232595in}{0.723078in}}%
\pgfpathlineto{\pgfqpoint{6.235951in}{0.715768in}}%
\pgfpathlineto{\pgfqpoint{6.239307in}{0.691309in}}%
\pgfpathlineto{\pgfqpoint{6.242663in}{0.643594in}}%
\pgfpathlineto{\pgfqpoint{6.246019in}{0.710564in}}%
\pgfpathlineto{\pgfqpoint{6.249374in}{0.730320in}}%
\pgfpathlineto{\pgfqpoint{6.252730in}{0.736292in}}%
\pgfpathlineto{\pgfqpoint{6.256086in}{0.727736in}}%
\pgfpathlineto{\pgfqpoint{6.259442in}{0.694067in}}%
\pgfpathlineto{\pgfqpoint{6.262797in}{0.677693in}}%
\pgfpathlineto{\pgfqpoint{6.266153in}{0.720439in}}%
\pgfpathlineto{\pgfqpoint{6.269509in}{0.732836in}}%
\pgfpathlineto{\pgfqpoint{6.272865in}{0.733366in}}%
\pgfpathlineto{\pgfqpoint{6.276221in}{0.722145in}}%
\pgfpathlineto{\pgfqpoint{6.279576in}{0.681514in}}%
\pgfpathlineto{\pgfqpoint{6.282932in}{0.697161in}}%
\pgfpathlineto{\pgfqpoint{6.286288in}{0.733727in}}%
\pgfpathlineto{\pgfqpoint{6.289644in}{0.747157in}}%
\pgfpathlineto{\pgfqpoint{6.292999in}{0.747131in}}%
\pgfpathlineto{\pgfqpoint{6.296355in}{0.731570in}}%
\pgfpathlineto{\pgfqpoint{6.303067in}{0.663157in}}%
\pgfpathlineto{\pgfqpoint{6.306423in}{0.706028in}}%
\pgfpathlineto{\pgfqpoint{6.309778in}{0.711811in}}%
\pgfpathlineto{\pgfqpoint{6.313134in}{0.697899in}}%
\pgfpathlineto{\pgfqpoint{6.316490in}{0.539589in}}%
\pgfpathlineto{\pgfqpoint{6.319846in}{0.706245in}}%
\pgfpathlineto{\pgfqpoint{6.323201in}{0.730581in}}%
\pgfpathlineto{\pgfqpoint{6.326557in}{0.737528in}}%
\pgfpathlineto{\pgfqpoint{6.329913in}{0.727905in}}%
\pgfpathlineto{\pgfqpoint{6.333269in}{0.705504in}}%
\pgfpathlineto{\pgfqpoint{6.336625in}{0.653372in}}%
\pgfpathlineto{\pgfqpoint{6.339980in}{0.656898in}}%
\pgfpathlineto{\pgfqpoint{6.343336in}{0.677642in}}%
\pgfpathlineto{\pgfqpoint{6.346692in}{0.665321in}}%
\pgfpathlineto{\pgfqpoint{6.350048in}{0.605474in}}%
\pgfpathlineto{\pgfqpoint{6.353403in}{0.688051in}}%
\pgfpathlineto{\pgfqpoint{6.356759in}{0.710044in}}%
\pgfpathlineto{\pgfqpoint{6.360115in}{0.715199in}}%
\pgfpathlineto{\pgfqpoint{6.363471in}{0.704943in}}%
\pgfpathlineto{\pgfqpoint{6.373538in}{0.666743in}}%
\pgfpathlineto{\pgfqpoint{6.376894in}{0.658751in}}%
\pgfpathlineto{\pgfqpoint{6.380250in}{0.653978in}}%
\pgfpathlineto{\pgfqpoint{6.383605in}{0.647083in}}%
\pgfpathlineto{\pgfqpoint{6.383605in}{0.647083in}}%
\pgfusepath{stroke}%
\end{pgfscope}%
\begin{pgfscope}%
\pgfpathrectangle{\pgfqpoint{4.621821in}{0.320625in}}{\pgfqpoint{1.845679in}{1.164375in}}%
\pgfusepath{clip}%
\pgfsetrectcap%
\pgfsetroundjoin%
\pgfsetlinewidth{0.602250pt}%
\definecolor{currentstroke}{rgb}{0.819608,0.286275,0.356863}%
\pgfsetstrokecolor{currentstroke}%
\pgfsetstrokeopacity{0.120000}%
\pgfsetdash{}{0pt}%
\pgfpathmoveto{\pgfqpoint{4.705715in}{1.469196in}}%
\pgfpathlineto{\pgfqpoint{4.725850in}{1.433697in}}%
\pgfpathlineto{\pgfqpoint{4.745985in}{1.394185in}}%
\pgfpathlineto{\pgfqpoint{4.766120in}{1.355343in}}%
\pgfpathlineto{\pgfqpoint{4.782898in}{1.325411in}}%
\pgfpathlineto{\pgfqpoint{4.803033in}{1.292118in}}%
\pgfpathlineto{\pgfqpoint{4.826524in}{1.255806in}}%
\pgfpathlineto{\pgfqpoint{4.853370in}{1.216465in}}%
\pgfpathlineto{\pgfqpoint{4.886928in}{1.169333in}}%
\pgfpathlineto{\pgfqpoint{4.930553in}{1.110113in}}%
\pgfpathlineto{\pgfqpoint{4.987601in}{1.034652in}}%
\pgfpathlineto{\pgfqpoint{5.064784in}{0.934470in}}%
\pgfpathlineto{\pgfqpoint{5.155390in}{0.818221in}}%
\pgfpathlineto{\pgfqpoint{5.175525in}{0.793136in}}%
\pgfpathlineto{\pgfqpoint{5.199015in}{0.761843in}}%
\pgfpathlineto{\pgfqpoint{5.209082in}{0.751148in}}%
\pgfpathlineto{\pgfqpoint{5.219150in}{0.741227in}}%
\pgfpathlineto{\pgfqpoint{5.225861in}{0.732366in}}%
\pgfpathlineto{\pgfqpoint{5.232573in}{0.721873in}}%
\pgfpathlineto{\pgfqpoint{5.256063in}{0.682552in}}%
\pgfpathlineto{\pgfqpoint{5.259419in}{0.679432in}}%
\pgfpathlineto{\pgfqpoint{5.262775in}{0.677664in}}%
\pgfpathlineto{\pgfqpoint{5.266131in}{0.677206in}}%
\pgfpathlineto{\pgfqpoint{5.269487in}{0.677800in}}%
\pgfpathlineto{\pgfqpoint{5.279554in}{0.681303in}}%
\pgfpathlineto{\pgfqpoint{5.282910in}{0.681317in}}%
\pgfpathlineto{\pgfqpoint{5.286265in}{0.679752in}}%
\pgfpathlineto{\pgfqpoint{5.289621in}{0.675802in}}%
\pgfpathlineto{\pgfqpoint{5.296333in}{0.664514in}}%
\pgfpathlineto{\pgfqpoint{5.299689in}{0.662426in}}%
\pgfpathlineto{\pgfqpoint{5.303044in}{0.666369in}}%
\pgfpathlineto{\pgfqpoint{5.309756in}{0.680627in}}%
\pgfpathlineto{\pgfqpoint{5.313112in}{0.686310in}}%
\pgfpathlineto{\pgfqpoint{5.316467in}{0.690371in}}%
\pgfpathlineto{\pgfqpoint{5.319823in}{0.692842in}}%
\pgfpathlineto{\pgfqpoint{5.323179in}{0.693756in}}%
\pgfpathlineto{\pgfqpoint{5.326535in}{0.693074in}}%
\pgfpathlineto{\pgfqpoint{5.329891in}{0.690629in}}%
\pgfpathlineto{\pgfqpoint{5.333246in}{0.686047in}}%
\pgfpathlineto{\pgfqpoint{5.336602in}{0.678546in}}%
\pgfpathlineto{\pgfqpoint{5.339958in}{0.666329in}}%
\pgfpathlineto{\pgfqpoint{5.343314in}{0.644014in}}%
\pgfpathlineto{\pgfqpoint{5.346669in}{0.592770in}}%
\pgfpathlineto{\pgfqpoint{5.350025in}{0.643297in}}%
\pgfpathlineto{\pgfqpoint{5.353381in}{0.669565in}}%
\pgfpathlineto{\pgfqpoint{5.356737in}{0.684709in}}%
\pgfpathlineto{\pgfqpoint{5.360093in}{0.694767in}}%
\pgfpathlineto{\pgfqpoint{5.363448in}{0.701626in}}%
\pgfpathlineto{\pgfqpoint{5.366804in}{0.706019in}}%
\pgfpathlineto{\pgfqpoint{5.370160in}{0.708156in}}%
\pgfpathlineto{\pgfqpoint{5.373516in}{0.707870in}}%
\pgfpathlineto{\pgfqpoint{5.376871in}{0.704493in}}%
\pgfpathlineto{\pgfqpoint{5.380227in}{0.696303in}}%
\pgfpathlineto{\pgfqpoint{5.386939in}{0.665916in}}%
\pgfpathlineto{\pgfqpoint{5.390295in}{0.667989in}}%
\pgfpathlineto{\pgfqpoint{5.393650in}{0.674756in}}%
\pgfpathlineto{\pgfqpoint{5.397006in}{0.673864in}}%
\pgfpathlineto{\pgfqpoint{5.400362in}{0.661944in}}%
\pgfpathlineto{\pgfqpoint{5.403718in}{0.638159in}}%
\pgfpathlineto{\pgfqpoint{5.410429in}{0.682559in}}%
\pgfpathlineto{\pgfqpoint{5.413785in}{0.691157in}}%
\pgfpathlineto{\pgfqpoint{5.417141in}{0.685561in}}%
\pgfpathlineto{\pgfqpoint{5.420497in}{0.663448in}}%
\pgfpathlineto{\pgfqpoint{5.423852in}{0.635231in}}%
\pgfpathlineto{\pgfqpoint{5.427208in}{0.663513in}}%
\pgfpathlineto{\pgfqpoint{5.430564in}{0.676816in}}%
\pgfpathlineto{\pgfqpoint{5.433920in}{0.677289in}}%
\pgfpathlineto{\pgfqpoint{5.437276in}{0.664479in}}%
\pgfpathlineto{\pgfqpoint{5.440631in}{0.614462in}}%
\pgfpathlineto{\pgfqpoint{5.443987in}{0.654085in}}%
\pgfpathlineto{\pgfqpoint{5.447343in}{0.683708in}}%
\pgfpathlineto{\pgfqpoint{5.450699in}{0.694249in}}%
\pgfpathlineto{\pgfqpoint{5.454054in}{0.689666in}}%
\pgfpathlineto{\pgfqpoint{5.457410in}{0.630078in}}%
\pgfpathlineto{\pgfqpoint{5.460766in}{0.703415in}}%
\pgfpathlineto{\pgfqpoint{5.464122in}{0.730577in}}%
\pgfpathlineto{\pgfqpoint{5.467478in}{0.743633in}}%
\pgfpathlineto{\pgfqpoint{5.470833in}{0.749242in}}%
\pgfpathlineto{\pgfqpoint{5.474189in}{0.748960in}}%
\pgfpathlineto{\pgfqpoint{5.477545in}{0.742168in}}%
\pgfpathlineto{\pgfqpoint{5.480901in}{0.724548in}}%
\pgfpathlineto{\pgfqpoint{5.484256in}{0.660414in}}%
\pgfpathlineto{\pgfqpoint{5.487612in}{0.714048in}}%
\pgfpathlineto{\pgfqpoint{5.490968in}{0.743511in}}%
\pgfpathlineto{\pgfqpoint{5.494324in}{0.756444in}}%
\pgfpathlineto{\pgfqpoint{5.497680in}{0.760026in}}%
\pgfpathlineto{\pgfqpoint{5.501035in}{0.752580in}}%
\pgfpathlineto{\pgfqpoint{5.504391in}{0.734685in}}%
\pgfpathlineto{\pgfqpoint{5.507747in}{0.704145in}}%
\pgfpathlineto{\pgfqpoint{5.511103in}{0.653071in}}%
\pgfpathlineto{\pgfqpoint{5.514458in}{0.665459in}}%
\pgfpathlineto{\pgfqpoint{5.517814in}{0.641278in}}%
\pgfpathlineto{\pgfqpoint{5.524526in}{0.711833in}}%
\pgfpathlineto{\pgfqpoint{5.527882in}{0.724893in}}%
\pgfpathlineto{\pgfqpoint{5.531237in}{0.719881in}}%
\pgfpathlineto{\pgfqpoint{5.534593in}{0.701582in}}%
\pgfpathlineto{\pgfqpoint{5.537949in}{0.669425in}}%
\pgfpathlineto{\pgfqpoint{5.541305in}{0.604689in}}%
\pgfpathlineto{\pgfqpoint{5.544660in}{0.632689in}}%
\pgfpathlineto{\pgfqpoint{5.548016in}{0.604035in}}%
\pgfpathlineto{\pgfqpoint{5.554728in}{0.680088in}}%
\pgfpathlineto{\pgfqpoint{5.558084in}{0.693116in}}%
\pgfpathlineto{\pgfqpoint{5.561439in}{0.687476in}}%
\pgfpathlineto{\pgfqpoint{5.564795in}{0.651755in}}%
\pgfpathlineto{\pgfqpoint{5.568151in}{0.648575in}}%
\pgfpathlineto{\pgfqpoint{5.571507in}{0.686537in}}%
\pgfpathlineto{\pgfqpoint{5.574863in}{0.698193in}}%
\pgfpathlineto{\pgfqpoint{5.578218in}{0.698310in}}%
\pgfpathlineto{\pgfqpoint{5.581574in}{0.686480in}}%
\pgfpathlineto{\pgfqpoint{5.584930in}{0.642845in}}%
\pgfpathlineto{\pgfqpoint{5.588286in}{0.664805in}}%
\pgfpathlineto{\pgfqpoint{5.591641in}{0.699407in}}%
\pgfpathlineto{\pgfqpoint{5.594997in}{0.711762in}}%
\pgfpathlineto{\pgfqpoint{5.598353in}{0.709804in}}%
\pgfpathlineto{\pgfqpoint{5.601709in}{0.684784in}}%
\pgfpathlineto{\pgfqpoint{5.605065in}{0.640923in}}%
\pgfpathlineto{\pgfqpoint{5.608420in}{0.700480in}}%
\pgfpathlineto{\pgfqpoint{5.611776in}{0.716350in}}%
\pgfpathlineto{\pgfqpoint{5.615132in}{0.719448in}}%
\pgfpathlineto{\pgfqpoint{5.618488in}{0.711809in}}%
\pgfpathlineto{\pgfqpoint{5.625199in}{0.657790in}}%
\pgfpathlineto{\pgfqpoint{5.628555in}{0.711742in}}%
\pgfpathlineto{\pgfqpoint{5.631911in}{0.729103in}}%
\pgfpathlineto{\pgfqpoint{5.635267in}{0.732181in}}%
\pgfpathlineto{\pgfqpoint{5.638622in}{0.716767in}}%
\pgfpathlineto{\pgfqpoint{5.641978in}{0.631487in}}%
\pgfpathlineto{\pgfqpoint{5.645334in}{0.709813in}}%
\pgfpathlineto{\pgfqpoint{5.648690in}{0.732723in}}%
\pgfpathlineto{\pgfqpoint{5.652045in}{0.739443in}}%
\pgfpathlineto{\pgfqpoint{5.655401in}{0.735661in}}%
\pgfpathlineto{\pgfqpoint{5.658757in}{0.717457in}}%
\pgfpathlineto{\pgfqpoint{5.662113in}{0.582257in}}%
\pgfpathlineto{\pgfqpoint{5.665469in}{0.720912in}}%
\pgfpathlineto{\pgfqpoint{5.668824in}{0.744953in}}%
\pgfpathlineto{\pgfqpoint{5.672180in}{0.752604in}}%
\pgfpathlineto{\pgfqpoint{5.675536in}{0.744897in}}%
\pgfpathlineto{\pgfqpoint{5.678892in}{0.716497in}}%
\pgfpathlineto{\pgfqpoint{5.682247in}{0.657811in}}%
\pgfpathlineto{\pgfqpoint{5.685603in}{0.721470in}}%
\pgfpathlineto{\pgfqpoint{5.688959in}{0.733678in}}%
\pgfpathlineto{\pgfqpoint{5.692315in}{0.730015in}}%
\pgfpathlineto{\pgfqpoint{5.699026in}{0.678450in}}%
\pgfpathlineto{\pgfqpoint{5.702382in}{0.733263in}}%
\pgfpathlineto{\pgfqpoint{5.705738in}{0.750221in}}%
\pgfpathlineto{\pgfqpoint{5.709094in}{0.750959in}}%
\pgfpathlineto{\pgfqpoint{5.712449in}{0.734959in}}%
\pgfpathlineto{\pgfqpoint{5.715805in}{0.701373in}}%
\pgfpathlineto{\pgfqpoint{5.719161in}{0.636891in}}%
\pgfpathlineto{\pgfqpoint{5.722517in}{0.694266in}}%
\pgfpathlineto{\pgfqpoint{5.725873in}{0.694906in}}%
\pgfpathlineto{\pgfqpoint{5.729228in}{0.651779in}}%
\pgfpathlineto{\pgfqpoint{5.735940in}{0.721739in}}%
\pgfpathlineto{\pgfqpoint{5.739296in}{0.732470in}}%
\pgfpathlineto{\pgfqpoint{5.742652in}{0.724697in}}%
\pgfpathlineto{\pgfqpoint{5.746007in}{0.701127in}}%
\pgfpathlineto{\pgfqpoint{5.749363in}{0.633736in}}%
\pgfpathlineto{\pgfqpoint{5.752719in}{0.668513in}}%
\pgfpathlineto{\pgfqpoint{5.756075in}{0.684114in}}%
\pgfpathlineto{\pgfqpoint{5.759430in}{0.671945in}}%
\pgfpathlineto{\pgfqpoint{5.762786in}{0.582769in}}%
\pgfpathlineto{\pgfqpoint{5.766142in}{0.688257in}}%
\pgfpathlineto{\pgfqpoint{5.769498in}{0.710665in}}%
\pgfpathlineto{\pgfqpoint{5.772854in}{0.714811in}}%
\pgfpathlineto{\pgfqpoint{5.776209in}{0.698530in}}%
\pgfpathlineto{\pgfqpoint{5.779565in}{0.639154in}}%
\pgfpathlineto{\pgfqpoint{5.782921in}{0.675763in}}%
\pgfpathlineto{\pgfqpoint{5.786277in}{0.700340in}}%
\pgfpathlineto{\pgfqpoint{5.789632in}{0.705159in}}%
\pgfpathlineto{\pgfqpoint{5.792988in}{0.696317in}}%
\pgfpathlineto{\pgfqpoint{5.796344in}{0.658646in}}%
\pgfpathlineto{\pgfqpoint{5.799700in}{0.670197in}}%
\pgfpathlineto{\pgfqpoint{5.803056in}{0.707983in}}%
\pgfpathlineto{\pgfqpoint{5.806411in}{0.720820in}}%
\pgfpathlineto{\pgfqpoint{5.809767in}{0.717758in}}%
\pgfpathlineto{\pgfqpoint{5.813123in}{0.687628in}}%
\pgfpathlineto{\pgfqpoint{5.816479in}{0.671080in}}%
\pgfpathlineto{\pgfqpoint{5.819834in}{0.716224in}}%
\pgfpathlineto{\pgfqpoint{5.823190in}{0.730223in}}%
\pgfpathlineto{\pgfqpoint{5.826546in}{0.732653in}}%
\pgfpathlineto{\pgfqpoint{5.829902in}{0.724795in}}%
\pgfpathlineto{\pgfqpoint{5.836613in}{0.667974in}}%
\pgfpathlineto{\pgfqpoint{5.839969in}{0.723626in}}%
\pgfpathlineto{\pgfqpoint{5.843325in}{0.741403in}}%
\pgfpathlineto{\pgfqpoint{5.846681in}{0.745142in}}%
\pgfpathlineto{\pgfqpoint{5.850036in}{0.733249in}}%
\pgfpathlineto{\pgfqpoint{5.853392in}{0.704145in}}%
\pgfpathlineto{\pgfqpoint{5.856748in}{0.605958in}}%
\pgfpathlineto{\pgfqpoint{5.860104in}{0.694787in}}%
\pgfpathlineto{\pgfqpoint{5.863460in}{0.703286in}}%
\pgfpathlineto{\pgfqpoint{5.866815in}{0.687959in}}%
\pgfpathlineto{\pgfqpoint{5.870171in}{0.618490in}}%
\pgfpathlineto{\pgfqpoint{5.873527in}{0.705905in}}%
\pgfpathlineto{\pgfqpoint{5.876883in}{0.727003in}}%
\pgfpathlineto{\pgfqpoint{5.880238in}{0.730250in}}%
\pgfpathlineto{\pgfqpoint{5.883594in}{0.715100in}}%
\pgfpathlineto{\pgfqpoint{5.890306in}{0.640524in}}%
\pgfpathlineto{\pgfqpoint{5.893662in}{0.685160in}}%
\pgfpathlineto{\pgfqpoint{5.897017in}{0.688531in}}%
\pgfpathlineto{\pgfqpoint{5.903729in}{0.645401in}}%
\pgfpathlineto{\pgfqpoint{5.907085in}{0.698758in}}%
\pgfpathlineto{\pgfqpoint{5.910441in}{0.715472in}}%
\pgfpathlineto{\pgfqpoint{5.913796in}{0.714635in}}%
\pgfpathlineto{\pgfqpoint{5.917152in}{0.692170in}}%
\pgfpathlineto{\pgfqpoint{5.920508in}{0.555581in}}%
\pgfpathlineto{\pgfqpoint{5.923864in}{0.689210in}}%
\pgfpathlineto{\pgfqpoint{5.927219in}{0.706491in}}%
\pgfpathlineto{\pgfqpoint{5.930575in}{0.708034in}}%
\pgfpathlineto{\pgfqpoint{5.933931in}{0.695379in}}%
\pgfpathlineto{\pgfqpoint{5.937287in}{0.638237in}}%
\pgfpathlineto{\pgfqpoint{5.940643in}{0.685676in}}%
\pgfpathlineto{\pgfqpoint{5.943998in}{0.715220in}}%
\pgfpathlineto{\pgfqpoint{5.947354in}{0.724862in}}%
\pgfpathlineto{\pgfqpoint{5.950710in}{0.717855in}}%
\pgfpathlineto{\pgfqpoint{5.954066in}{0.679131in}}%
\pgfpathlineto{\pgfqpoint{5.957421in}{0.684743in}}%
\pgfpathlineto{\pgfqpoint{5.960777in}{0.720301in}}%
\pgfpathlineto{\pgfqpoint{5.964133in}{0.731587in}}%
\pgfpathlineto{\pgfqpoint{5.967489in}{0.731739in}}%
\pgfpathlineto{\pgfqpoint{5.970845in}{0.720353in}}%
\pgfpathlineto{\pgfqpoint{5.974200in}{0.679764in}}%
\pgfpathlineto{\pgfqpoint{5.977556in}{0.694829in}}%
\pgfpathlineto{\pgfqpoint{5.980912in}{0.731305in}}%
\pgfpathlineto{\pgfqpoint{5.984268in}{0.744398in}}%
\pgfpathlineto{\pgfqpoint{5.987623in}{0.743506in}}%
\pgfpathlineto{\pgfqpoint{5.990979in}{0.723091in}}%
\pgfpathlineto{\pgfqpoint{5.994335in}{0.595707in}}%
\pgfpathlineto{\pgfqpoint{5.997691in}{0.716852in}}%
\pgfpathlineto{\pgfqpoint{6.001047in}{0.735981in}}%
\pgfpathlineto{\pgfqpoint{6.004402in}{0.738980in}}%
\pgfpathlineto{\pgfqpoint{6.007758in}{0.728872in}}%
\pgfpathlineto{\pgfqpoint{6.011114in}{0.688012in}}%
\pgfpathlineto{\pgfqpoint{6.014470in}{0.705927in}}%
\pgfpathlineto{\pgfqpoint{6.017825in}{0.741717in}}%
\pgfpathlineto{\pgfqpoint{6.021181in}{0.754012in}}%
\pgfpathlineto{\pgfqpoint{6.024537in}{0.750757in}}%
\pgfpathlineto{\pgfqpoint{6.027893in}{0.731191in}}%
\pgfpathlineto{\pgfqpoint{6.031249in}{0.681795in}}%
\pgfpathlineto{\pgfqpoint{6.034604in}{0.686968in}}%
\pgfpathlineto{\pgfqpoint{6.037960in}{0.711832in}}%
\pgfpathlineto{\pgfqpoint{6.041316in}{0.708830in}}%
\pgfpathlineto{\pgfqpoint{6.044672in}{0.670241in}}%
\pgfpathlineto{\pgfqpoint{6.048027in}{0.693036in}}%
\pgfpathlineto{\pgfqpoint{6.051383in}{0.728894in}}%
\pgfpathlineto{\pgfqpoint{6.054739in}{0.740521in}}%
\pgfpathlineto{\pgfqpoint{6.058095in}{0.734158in}}%
\pgfpathlineto{\pgfqpoint{6.061451in}{0.712796in}}%
\pgfpathlineto{\pgfqpoint{6.064806in}{0.662648in}}%
\pgfpathlineto{\pgfqpoint{6.068162in}{0.662602in}}%
\pgfpathlineto{\pgfqpoint{6.071518in}{0.684485in}}%
\pgfpathlineto{\pgfqpoint{6.074874in}{0.672303in}}%
\pgfpathlineto{\pgfqpoint{6.078230in}{0.613606in}}%
\pgfpathlineto{\pgfqpoint{6.081585in}{0.694982in}}%
\pgfpathlineto{\pgfqpoint{6.084941in}{0.716146in}}%
\pgfpathlineto{\pgfqpoint{6.088297in}{0.719216in}}%
\pgfpathlineto{\pgfqpoint{6.091653in}{0.702574in}}%
\pgfpathlineto{\pgfqpoint{6.095008in}{0.654356in}}%
\pgfpathlineto{\pgfqpoint{6.098364in}{0.664898in}}%
\pgfpathlineto{\pgfqpoint{6.101720in}{0.692515in}}%
\pgfpathlineto{\pgfqpoint{6.105076in}{0.695808in}}%
\pgfpathlineto{\pgfqpoint{6.108432in}{0.681268in}}%
\pgfpathlineto{\pgfqpoint{6.111787in}{0.562765in}}%
\pgfpathlineto{\pgfqpoint{6.115143in}{0.686539in}}%
\pgfpathlineto{\pgfqpoint{6.118499in}{0.710791in}}%
\pgfpathlineto{\pgfqpoint{6.121855in}{0.716897in}}%
\pgfpathlineto{\pgfqpoint{6.125210in}{0.703405in}}%
\pgfpathlineto{\pgfqpoint{6.128566in}{0.636901in}}%
\pgfpathlineto{\pgfqpoint{6.131922in}{0.690390in}}%
\pgfpathlineto{\pgfqpoint{6.135278in}{0.715408in}}%
\pgfpathlineto{\pgfqpoint{6.138634in}{0.722866in}}%
\pgfpathlineto{\pgfqpoint{6.141989in}{0.719599in}}%
\pgfpathlineto{\pgfqpoint{6.145345in}{0.702156in}}%
\pgfpathlineto{\pgfqpoint{6.148701in}{0.596701in}}%
\pgfpathlineto{\pgfqpoint{6.152057in}{0.703500in}}%
\pgfpathlineto{\pgfqpoint{6.155412in}{0.728283in}}%
\pgfpathlineto{\pgfqpoint{6.158768in}{0.736273in}}%
\pgfpathlineto{\pgfqpoint{6.162124in}{0.728603in}}%
\pgfpathlineto{\pgfqpoint{6.165480in}{0.689842in}}%
\pgfpathlineto{\pgfqpoint{6.168836in}{0.694596in}}%
\pgfpathlineto{\pgfqpoint{6.172191in}{0.730152in}}%
\pgfpathlineto{\pgfqpoint{6.175547in}{0.741236in}}%
\pgfpathlineto{\pgfqpoint{6.178903in}{0.741037in}}%
\pgfpathlineto{\pgfqpoint{6.182259in}{0.728876in}}%
\pgfpathlineto{\pgfqpoint{6.185614in}{0.683896in}}%
\pgfpathlineto{\pgfqpoint{6.188970in}{0.708454in}}%
\pgfpathlineto{\pgfqpoint{6.192326in}{0.742403in}}%
\pgfpathlineto{\pgfqpoint{6.195682in}{0.754547in}}%
\pgfpathlineto{\pgfqpoint{6.199038in}{0.752415in}}%
\pgfpathlineto{\pgfqpoint{6.202393in}{0.732706in}}%
\pgfpathlineto{\pgfqpoint{6.205749in}{0.670013in}}%
\pgfpathlineto{\pgfqpoint{6.209105in}{0.706376in}}%
\pgfpathlineto{\pgfqpoint{6.212461in}{0.728037in}}%
\pgfpathlineto{\pgfqpoint{6.215816in}{0.728538in}}%
\pgfpathlineto{\pgfqpoint{6.219172in}{0.709241in}}%
\pgfpathlineto{\pgfqpoint{6.222528in}{0.649908in}}%
\pgfpathlineto{\pgfqpoint{6.225884in}{0.727278in}}%
\pgfpathlineto{\pgfqpoint{6.229240in}{0.747258in}}%
\pgfpathlineto{\pgfqpoint{6.232595in}{0.750242in}}%
\pgfpathlineto{\pgfqpoint{6.235951in}{0.736051in}}%
\pgfpathlineto{\pgfqpoint{6.239307in}{0.707092in}}%
\pgfpathlineto{\pgfqpoint{6.242663in}{0.579189in}}%
\pgfpathlineto{\pgfqpoint{6.246019in}{0.683020in}}%
\pgfpathlineto{\pgfqpoint{6.249374in}{0.687134in}}%
\pgfpathlineto{\pgfqpoint{6.252730in}{0.642928in}}%
\pgfpathlineto{\pgfqpoint{6.256086in}{0.684316in}}%
\pgfpathlineto{\pgfqpoint{6.259442in}{0.717441in}}%
\pgfpathlineto{\pgfqpoint{6.262797in}{0.728108in}}%
\pgfpathlineto{\pgfqpoint{6.266153in}{0.720098in}}%
\pgfpathlineto{\pgfqpoint{6.269509in}{0.695607in}}%
\pgfpathlineto{\pgfqpoint{6.272865in}{0.614664in}}%
\pgfpathlineto{\pgfqpoint{6.276221in}{0.670122in}}%
\pgfpathlineto{\pgfqpoint{6.279576in}{0.684205in}}%
\pgfpathlineto{\pgfqpoint{6.282932in}{0.673106in}}%
\pgfpathlineto{\pgfqpoint{6.286288in}{0.508979in}}%
\pgfpathlineto{\pgfqpoint{6.289644in}{0.683386in}}%
\pgfpathlineto{\pgfqpoint{6.292999in}{0.707506in}}%
\pgfpathlineto{\pgfqpoint{6.296355in}{0.712959in}}%
\pgfpathlineto{\pgfqpoint{6.299711in}{0.698035in}}%
\pgfpathlineto{\pgfqpoint{6.303067in}{0.637689in}}%
\pgfpathlineto{\pgfqpoint{6.306423in}{0.678341in}}%
\pgfpathlineto{\pgfqpoint{6.309778in}{0.703449in}}%
\pgfpathlineto{\pgfqpoint{6.313134in}{0.709549in}}%
\pgfpathlineto{\pgfqpoint{6.316490in}{0.703392in}}%
\pgfpathlineto{\pgfqpoint{6.323201in}{0.649770in}}%
\pgfpathlineto{\pgfqpoint{6.326557in}{0.704358in}}%
\pgfpathlineto{\pgfqpoint{6.329913in}{0.721527in}}%
\pgfpathlineto{\pgfqpoint{6.333269in}{0.723535in}}%
\pgfpathlineto{\pgfqpoint{6.336625in}{0.704164in}}%
\pgfpathlineto{\pgfqpoint{6.339980in}{0.625761in}}%
\pgfpathlineto{\pgfqpoint{6.343336in}{0.712965in}}%
\pgfpathlineto{\pgfqpoint{6.346692in}{0.732129in}}%
\pgfpathlineto{\pgfqpoint{6.350048in}{0.737793in}}%
\pgfpathlineto{\pgfqpoint{6.353403in}{0.733794in}}%
\pgfpathlineto{\pgfqpoint{6.356759in}{0.716302in}}%
\pgfpathlineto{\pgfqpoint{6.360115in}{0.623215in}}%
\pgfpathlineto{\pgfqpoint{6.363471in}{0.715569in}}%
\pgfpathlineto{\pgfqpoint{6.366827in}{0.741069in}}%
\pgfpathlineto{\pgfqpoint{6.370182in}{0.749884in}}%
\pgfpathlineto{\pgfqpoint{6.373538in}{0.744591in}}%
\pgfpathlineto{\pgfqpoint{6.383605in}{0.708917in}}%
\pgfpathlineto{\pgfqpoint{6.383605in}{0.708917in}}%
\pgfusepath{stroke}%
\end{pgfscope}%
\begin{pgfscope}%
\pgfpathrectangle{\pgfqpoint{4.621821in}{0.320625in}}{\pgfqpoint{1.845679in}{1.164375in}}%
\pgfusepath{clip}%
\pgfsetrectcap%
\pgfsetroundjoin%
\pgfsetlinewidth{1.154312pt}%
\definecolor{currentstroke}{rgb}{0.819608,0.286275,0.356863}%
\pgfsetstrokecolor{currentstroke}%
\pgfsetdash{}{0pt}%
\pgfpathmoveto{\pgfqpoint{4.705715in}{1.471665in}}%
\pgfpathlineto{\pgfqpoint{4.729206in}{1.429105in}}%
\pgfpathlineto{\pgfqpoint{4.772831in}{1.345188in}}%
\pgfpathlineto{\pgfqpoint{4.789610in}{1.316203in}}%
\pgfpathlineto{\pgfqpoint{4.809745in}{1.283688in}}%
\pgfpathlineto{\pgfqpoint{4.833235in}{1.248033in}}%
\pgfpathlineto{\pgfqpoint{4.860081in}{1.209475in}}%
\pgfpathlineto{\pgfqpoint{4.886928in}{1.172754in}}%
\pgfpathlineto{\pgfqpoint{4.917130in}{1.133333in}}%
\pgfpathlineto{\pgfqpoint{4.947332in}{1.095721in}}%
\pgfpathlineto{\pgfqpoint{4.980889in}{1.055802in}}%
\pgfpathlineto{\pgfqpoint{5.017803in}{1.013749in}}%
\pgfpathlineto{\pgfqpoint{5.061428in}{0.965865in}}%
\pgfpathlineto{\pgfqpoint{5.118476in}{0.905069in}}%
\pgfpathlineto{\pgfqpoint{5.202371in}{0.817502in}}%
\pgfpathlineto{\pgfqpoint{5.235929in}{0.782739in}}%
\pgfpathlineto{\pgfqpoint{5.256063in}{0.761975in}}%
\pgfpathlineto{\pgfqpoint{5.282910in}{0.735825in}}%
\pgfpathlineto{\pgfqpoint{5.309756in}{0.705947in}}%
\pgfpathlineto{\pgfqpoint{5.316467in}{0.700707in}}%
\pgfpathlineto{\pgfqpoint{5.323179in}{0.697124in}}%
\pgfpathlineto{\pgfqpoint{5.326535in}{0.696559in}}%
\pgfpathlineto{\pgfqpoint{5.329891in}{0.696988in}}%
\pgfpathlineto{\pgfqpoint{5.339958in}{0.700290in}}%
\pgfpathlineto{\pgfqpoint{5.343314in}{0.700487in}}%
\pgfpathlineto{\pgfqpoint{5.350025in}{0.699523in}}%
\pgfpathlineto{\pgfqpoint{5.353381in}{0.699973in}}%
\pgfpathlineto{\pgfqpoint{5.356737in}{0.701648in}}%
\pgfpathlineto{\pgfqpoint{5.366804in}{0.708951in}}%
\pgfpathlineto{\pgfqpoint{5.370160in}{0.710533in}}%
\pgfpathlineto{\pgfqpoint{5.376871in}{0.712021in}}%
\pgfpathlineto{\pgfqpoint{5.383583in}{0.713227in}}%
\pgfpathlineto{\pgfqpoint{5.390295in}{0.716195in}}%
\pgfpathlineto{\pgfqpoint{5.393650in}{0.716928in}}%
\pgfpathlineto{\pgfqpoint{5.397006in}{0.716191in}}%
\pgfpathlineto{\pgfqpoint{5.400362in}{0.714527in}}%
\pgfpathlineto{\pgfqpoint{5.403718in}{0.711627in}}%
\pgfpathlineto{\pgfqpoint{5.407074in}{0.705673in}}%
\pgfpathlineto{\pgfqpoint{5.410429in}{0.695702in}}%
\pgfpathlineto{\pgfqpoint{5.413785in}{0.687765in}}%
\pgfpathlineto{\pgfqpoint{5.417141in}{0.691640in}}%
\pgfpathlineto{\pgfqpoint{5.420497in}{0.700800in}}%
\pgfpathlineto{\pgfqpoint{5.423852in}{0.707994in}}%
\pgfpathlineto{\pgfqpoint{5.427208in}{0.712061in}}%
\pgfpathlineto{\pgfqpoint{5.433920in}{0.714831in}}%
\pgfpathlineto{\pgfqpoint{5.437276in}{0.715499in}}%
\pgfpathlineto{\pgfqpoint{5.440631in}{0.712521in}}%
\pgfpathlineto{\pgfqpoint{5.443987in}{0.703793in}}%
\pgfpathlineto{\pgfqpoint{5.450699in}{0.678712in}}%
\pgfpathlineto{\pgfqpoint{5.454054in}{0.674593in}}%
\pgfpathlineto{\pgfqpoint{5.457410in}{0.668999in}}%
\pgfpathlineto{\pgfqpoint{5.460766in}{0.683369in}}%
\pgfpathlineto{\pgfqpoint{5.464122in}{0.702787in}}%
\pgfpathlineto{\pgfqpoint{5.467478in}{0.715448in}}%
\pgfpathlineto{\pgfqpoint{5.470833in}{0.722120in}}%
\pgfpathlineto{\pgfqpoint{5.474189in}{0.722908in}}%
\pgfpathlineto{\pgfqpoint{5.477545in}{0.717432in}}%
\pgfpathlineto{\pgfqpoint{5.484256in}{0.696662in}}%
\pgfpathlineto{\pgfqpoint{5.487612in}{0.708670in}}%
\pgfpathlineto{\pgfqpoint{5.490968in}{0.723742in}}%
\pgfpathlineto{\pgfqpoint{5.494324in}{0.732500in}}%
\pgfpathlineto{\pgfqpoint{5.497680in}{0.734568in}}%
\pgfpathlineto{\pgfqpoint{5.501035in}{0.728227in}}%
\pgfpathlineto{\pgfqpoint{5.504391in}{0.718142in}}%
\pgfpathlineto{\pgfqpoint{5.507747in}{0.714806in}}%
\pgfpathlineto{\pgfqpoint{5.511103in}{0.717596in}}%
\pgfpathlineto{\pgfqpoint{5.514458in}{0.717492in}}%
\pgfpathlineto{\pgfqpoint{5.517814in}{0.713449in}}%
\pgfpathlineto{\pgfqpoint{5.521170in}{0.713975in}}%
\pgfpathlineto{\pgfqpoint{5.524526in}{0.715184in}}%
\pgfpathlineto{\pgfqpoint{5.527882in}{0.713064in}}%
\pgfpathlineto{\pgfqpoint{5.531237in}{0.703180in}}%
\pgfpathlineto{\pgfqpoint{5.534593in}{0.683858in}}%
\pgfpathlineto{\pgfqpoint{5.537949in}{0.684370in}}%
\pgfpathlineto{\pgfqpoint{5.541305in}{0.704920in}}%
\pgfpathlineto{\pgfqpoint{5.544660in}{0.715899in}}%
\pgfpathlineto{\pgfqpoint{5.548016in}{0.715047in}}%
\pgfpathlineto{\pgfqpoint{5.551372in}{0.704191in}}%
\pgfpathlineto{\pgfqpoint{5.554728in}{0.700278in}}%
\pgfpathlineto{\pgfqpoint{5.558084in}{0.706363in}}%
\pgfpathlineto{\pgfqpoint{5.561439in}{0.705409in}}%
\pgfpathlineto{\pgfqpoint{5.564795in}{0.698445in}}%
\pgfpathlineto{\pgfqpoint{5.568151in}{0.688386in}}%
\pgfpathlineto{\pgfqpoint{5.571507in}{0.685469in}}%
\pgfpathlineto{\pgfqpoint{5.578218in}{0.708280in}}%
\pgfpathlineto{\pgfqpoint{5.581574in}{0.711261in}}%
\pgfpathlineto{\pgfqpoint{5.584930in}{0.706389in}}%
\pgfpathlineto{\pgfqpoint{5.588286in}{0.703024in}}%
\pgfpathlineto{\pgfqpoint{5.594997in}{0.712912in}}%
\pgfpathlineto{\pgfqpoint{5.598353in}{0.711345in}}%
\pgfpathlineto{\pgfqpoint{5.601709in}{0.705081in}}%
\pgfpathlineto{\pgfqpoint{5.605065in}{0.700714in}}%
\pgfpathlineto{\pgfqpoint{5.608420in}{0.705845in}}%
\pgfpathlineto{\pgfqpoint{5.611776in}{0.714727in}}%
\pgfpathlineto{\pgfqpoint{5.615132in}{0.719989in}}%
\pgfpathlineto{\pgfqpoint{5.618488in}{0.718464in}}%
\pgfpathlineto{\pgfqpoint{5.621843in}{0.714280in}}%
\pgfpathlineto{\pgfqpoint{5.625199in}{0.715995in}}%
\pgfpathlineto{\pgfqpoint{5.628555in}{0.723118in}}%
\pgfpathlineto{\pgfqpoint{5.631911in}{0.727759in}}%
\pgfpathlineto{\pgfqpoint{5.635267in}{0.725758in}}%
\pgfpathlineto{\pgfqpoint{5.638622in}{0.715868in}}%
\pgfpathlineto{\pgfqpoint{5.641978in}{0.708713in}}%
\pgfpathlineto{\pgfqpoint{5.648690in}{0.732190in}}%
\pgfpathlineto{\pgfqpoint{5.652045in}{0.737098in}}%
\pgfpathlineto{\pgfqpoint{5.655401in}{0.732731in}}%
\pgfpathlineto{\pgfqpoint{5.658757in}{0.722579in}}%
\pgfpathlineto{\pgfqpoint{5.662113in}{0.721595in}}%
\pgfpathlineto{\pgfqpoint{5.665469in}{0.731673in}}%
\pgfpathlineto{\pgfqpoint{5.668824in}{0.739796in}}%
\pgfpathlineto{\pgfqpoint{5.672180in}{0.739614in}}%
\pgfpathlineto{\pgfqpoint{5.675536in}{0.727148in}}%
\pgfpathlineto{\pgfqpoint{5.678892in}{0.704403in}}%
\pgfpathlineto{\pgfqpoint{5.682247in}{0.717385in}}%
\pgfpathlineto{\pgfqpoint{5.685603in}{0.734180in}}%
\pgfpathlineto{\pgfqpoint{5.688959in}{0.737477in}}%
\pgfpathlineto{\pgfqpoint{5.692315in}{0.728463in}}%
\pgfpathlineto{\pgfqpoint{5.695671in}{0.714506in}}%
\pgfpathlineto{\pgfqpoint{5.699026in}{0.714053in}}%
\pgfpathlineto{\pgfqpoint{5.702382in}{0.729941in}}%
\pgfpathlineto{\pgfqpoint{5.705738in}{0.739423in}}%
\pgfpathlineto{\pgfqpoint{5.709094in}{0.737481in}}%
\pgfpathlineto{\pgfqpoint{5.715805in}{0.710338in}}%
\pgfpathlineto{\pgfqpoint{5.719161in}{0.715426in}}%
\pgfpathlineto{\pgfqpoint{5.722517in}{0.714934in}}%
\pgfpathlineto{\pgfqpoint{5.729228in}{0.696294in}}%
\pgfpathlineto{\pgfqpoint{5.732584in}{0.700615in}}%
\pgfpathlineto{\pgfqpoint{5.735940in}{0.713555in}}%
\pgfpathlineto{\pgfqpoint{5.739296in}{0.722667in}}%
\pgfpathlineto{\pgfqpoint{5.742652in}{0.720974in}}%
\pgfpathlineto{\pgfqpoint{5.746007in}{0.711222in}}%
\pgfpathlineto{\pgfqpoint{5.749363in}{0.704077in}}%
\pgfpathlineto{\pgfqpoint{5.752719in}{0.703530in}}%
\pgfpathlineto{\pgfqpoint{5.756075in}{0.694565in}}%
\pgfpathlineto{\pgfqpoint{5.759430in}{0.681124in}}%
\pgfpathlineto{\pgfqpoint{5.762786in}{0.684925in}}%
\pgfpathlineto{\pgfqpoint{5.766142in}{0.699803in}}%
\pgfpathlineto{\pgfqpoint{5.769498in}{0.708436in}}%
\pgfpathlineto{\pgfqpoint{5.772854in}{0.709831in}}%
\pgfpathlineto{\pgfqpoint{5.779565in}{0.700294in}}%
\pgfpathlineto{\pgfqpoint{5.782921in}{0.704242in}}%
\pgfpathlineto{\pgfqpoint{5.786277in}{0.711246in}}%
\pgfpathlineto{\pgfqpoint{5.789632in}{0.708808in}}%
\pgfpathlineto{\pgfqpoint{5.792988in}{0.696596in}}%
\pgfpathlineto{\pgfqpoint{5.796344in}{0.690288in}}%
\pgfpathlineto{\pgfqpoint{5.803056in}{0.716617in}}%
\pgfpathlineto{\pgfqpoint{5.806411in}{0.721627in}}%
\pgfpathlineto{\pgfqpoint{5.809767in}{0.717969in}}%
\pgfpathlineto{\pgfqpoint{5.813123in}{0.707364in}}%
\pgfpathlineto{\pgfqpoint{5.816479in}{0.706386in}}%
\pgfpathlineto{\pgfqpoint{5.819834in}{0.720241in}}%
\pgfpathlineto{\pgfqpoint{5.823190in}{0.729666in}}%
\pgfpathlineto{\pgfqpoint{5.826546in}{0.728821in}}%
\pgfpathlineto{\pgfqpoint{5.833258in}{0.711596in}}%
\pgfpathlineto{\pgfqpoint{5.836613in}{0.716616in}}%
\pgfpathlineto{\pgfqpoint{5.839969in}{0.726911in}}%
\pgfpathlineto{\pgfqpoint{5.843325in}{0.733483in}}%
\pgfpathlineto{\pgfqpoint{5.846681in}{0.734040in}}%
\pgfpathlineto{\pgfqpoint{5.850036in}{0.728750in}}%
\pgfpathlineto{\pgfqpoint{5.853392in}{0.727905in}}%
\pgfpathlineto{\pgfqpoint{5.856748in}{0.732651in}}%
\pgfpathlineto{\pgfqpoint{5.860104in}{0.734888in}}%
\pgfpathlineto{\pgfqpoint{5.863460in}{0.727711in}}%
\pgfpathlineto{\pgfqpoint{5.866815in}{0.718724in}}%
\pgfpathlineto{\pgfqpoint{5.870171in}{0.717668in}}%
\pgfpathlineto{\pgfqpoint{5.873527in}{0.721589in}}%
\pgfpathlineto{\pgfqpoint{5.876883in}{0.723862in}}%
\pgfpathlineto{\pgfqpoint{5.883594in}{0.722492in}}%
\pgfpathlineto{\pgfqpoint{5.890306in}{0.739107in}}%
\pgfpathlineto{\pgfqpoint{5.893662in}{0.738817in}}%
\pgfpathlineto{\pgfqpoint{5.897017in}{0.724570in}}%
\pgfpathlineto{\pgfqpoint{5.900373in}{0.697791in}}%
\pgfpathlineto{\pgfqpoint{5.903729in}{0.685889in}}%
\pgfpathlineto{\pgfqpoint{5.907085in}{0.703970in}}%
\pgfpathlineto{\pgfqpoint{5.910441in}{0.710452in}}%
\pgfpathlineto{\pgfqpoint{5.917152in}{0.704827in}}%
\pgfpathlineto{\pgfqpoint{5.923864in}{0.727300in}}%
\pgfpathlineto{\pgfqpoint{5.927219in}{0.728419in}}%
\pgfpathlineto{\pgfqpoint{5.930575in}{0.718730in}}%
\pgfpathlineto{\pgfqpoint{5.937287in}{0.678236in}}%
\pgfpathlineto{\pgfqpoint{5.940643in}{0.686353in}}%
\pgfpathlineto{\pgfqpoint{5.943998in}{0.697569in}}%
\pgfpathlineto{\pgfqpoint{5.947354in}{0.702344in}}%
\pgfpathlineto{\pgfqpoint{5.950710in}{0.702471in}}%
\pgfpathlineto{\pgfqpoint{5.954066in}{0.705710in}}%
\pgfpathlineto{\pgfqpoint{5.957421in}{0.713662in}}%
\pgfpathlineto{\pgfqpoint{5.960777in}{0.716401in}}%
\pgfpathlineto{\pgfqpoint{5.964133in}{0.713645in}}%
\pgfpathlineto{\pgfqpoint{5.967489in}{0.709371in}}%
\pgfpathlineto{\pgfqpoint{5.974200in}{0.703469in}}%
\pgfpathlineto{\pgfqpoint{5.977556in}{0.702270in}}%
\pgfpathlineto{\pgfqpoint{5.980912in}{0.708124in}}%
\pgfpathlineto{\pgfqpoint{5.984268in}{0.717026in}}%
\pgfpathlineto{\pgfqpoint{5.987623in}{0.721571in}}%
\pgfpathlineto{\pgfqpoint{5.990979in}{0.721160in}}%
\pgfpathlineto{\pgfqpoint{5.994335in}{0.719533in}}%
\pgfpathlineto{\pgfqpoint{5.997691in}{0.714330in}}%
\pgfpathlineto{\pgfqpoint{6.001047in}{0.712016in}}%
\pgfpathlineto{\pgfqpoint{6.007758in}{0.724267in}}%
\pgfpathlineto{\pgfqpoint{6.011114in}{0.725297in}}%
\pgfpathlineto{\pgfqpoint{6.014470in}{0.722760in}}%
\pgfpathlineto{\pgfqpoint{6.017825in}{0.721854in}}%
\pgfpathlineto{\pgfqpoint{6.027893in}{0.735868in}}%
\pgfpathlineto{\pgfqpoint{6.031249in}{0.738296in}}%
\pgfpathlineto{\pgfqpoint{6.034604in}{0.730443in}}%
\pgfpathlineto{\pgfqpoint{6.037960in}{0.713251in}}%
\pgfpathlineto{\pgfqpoint{6.041316in}{0.702792in}}%
\pgfpathlineto{\pgfqpoint{6.044672in}{0.717613in}}%
\pgfpathlineto{\pgfqpoint{6.048027in}{0.727899in}}%
\pgfpathlineto{\pgfqpoint{6.051383in}{0.730621in}}%
\pgfpathlineto{\pgfqpoint{6.054739in}{0.727415in}}%
\pgfpathlineto{\pgfqpoint{6.058095in}{0.721191in}}%
\pgfpathlineto{\pgfqpoint{6.061451in}{0.727802in}}%
\pgfpathlineto{\pgfqpoint{6.064806in}{0.738387in}}%
\pgfpathlineto{\pgfqpoint{6.068162in}{0.739785in}}%
\pgfpathlineto{\pgfqpoint{6.071518in}{0.735169in}}%
\pgfpathlineto{\pgfqpoint{6.074874in}{0.729349in}}%
\pgfpathlineto{\pgfqpoint{6.078230in}{0.719939in}}%
\pgfpathlineto{\pgfqpoint{6.081585in}{0.708019in}}%
\pgfpathlineto{\pgfqpoint{6.084941in}{0.707137in}}%
\pgfpathlineto{\pgfqpoint{6.091653in}{0.717185in}}%
\pgfpathlineto{\pgfqpoint{6.095008in}{0.721827in}}%
\pgfpathlineto{\pgfqpoint{6.098364in}{0.722734in}}%
\pgfpathlineto{\pgfqpoint{6.101720in}{0.719365in}}%
\pgfpathlineto{\pgfqpoint{6.105076in}{0.721795in}}%
\pgfpathlineto{\pgfqpoint{6.108432in}{0.728200in}}%
\pgfpathlineto{\pgfqpoint{6.111787in}{0.729449in}}%
\pgfpathlineto{\pgfqpoint{6.115143in}{0.722171in}}%
\pgfpathlineto{\pgfqpoint{6.118499in}{0.707097in}}%
\pgfpathlineto{\pgfqpoint{6.121855in}{0.695904in}}%
\pgfpathlineto{\pgfqpoint{6.125210in}{0.697218in}}%
\pgfpathlineto{\pgfqpoint{6.128566in}{0.702807in}}%
\pgfpathlineto{\pgfqpoint{6.131922in}{0.702530in}}%
\pgfpathlineto{\pgfqpoint{6.135278in}{0.703952in}}%
\pgfpathlineto{\pgfqpoint{6.141989in}{0.716535in}}%
\pgfpathlineto{\pgfqpoint{6.145345in}{0.716839in}}%
\pgfpathlineto{\pgfqpoint{6.148701in}{0.709571in}}%
\pgfpathlineto{\pgfqpoint{6.152057in}{0.698225in}}%
\pgfpathlineto{\pgfqpoint{6.155412in}{0.702725in}}%
\pgfpathlineto{\pgfqpoint{6.158768in}{0.712185in}}%
\pgfpathlineto{\pgfqpoint{6.162124in}{0.712965in}}%
\pgfpathlineto{\pgfqpoint{6.168836in}{0.695468in}}%
\pgfpathlineto{\pgfqpoint{6.175547in}{0.722351in}}%
\pgfpathlineto{\pgfqpoint{6.178903in}{0.726241in}}%
\pgfpathlineto{\pgfqpoint{6.182259in}{0.722360in}}%
\pgfpathlineto{\pgfqpoint{6.185614in}{0.709856in}}%
\pgfpathlineto{\pgfqpoint{6.188970in}{0.702068in}}%
\pgfpathlineto{\pgfqpoint{6.192326in}{0.718424in}}%
\pgfpathlineto{\pgfqpoint{6.195682in}{0.731861in}}%
\pgfpathlineto{\pgfqpoint{6.199038in}{0.735320in}}%
\pgfpathlineto{\pgfqpoint{6.202393in}{0.729268in}}%
\pgfpathlineto{\pgfqpoint{6.205749in}{0.714289in}}%
\pgfpathlineto{\pgfqpoint{6.209105in}{0.709728in}}%
\pgfpathlineto{\pgfqpoint{6.212461in}{0.718177in}}%
\pgfpathlineto{\pgfqpoint{6.215816in}{0.718864in}}%
\pgfpathlineto{\pgfqpoint{6.219172in}{0.708830in}}%
\pgfpathlineto{\pgfqpoint{6.222528in}{0.691615in}}%
\pgfpathlineto{\pgfqpoint{6.229240in}{0.729823in}}%
\pgfpathlineto{\pgfqpoint{6.232595in}{0.736172in}}%
\pgfpathlineto{\pgfqpoint{6.235951in}{0.730401in}}%
\pgfpathlineto{\pgfqpoint{6.242663in}{0.695127in}}%
\pgfpathlineto{\pgfqpoint{6.249374in}{0.723885in}}%
\pgfpathlineto{\pgfqpoint{6.252730in}{0.728910in}}%
\pgfpathlineto{\pgfqpoint{6.256086in}{0.725801in}}%
\pgfpathlineto{\pgfqpoint{6.259442in}{0.720504in}}%
\pgfpathlineto{\pgfqpoint{6.262797in}{0.716732in}}%
\pgfpathlineto{\pgfqpoint{6.266153in}{0.718014in}}%
\pgfpathlineto{\pgfqpoint{6.269509in}{0.720455in}}%
\pgfpathlineto{\pgfqpoint{6.272865in}{0.720914in}}%
\pgfpathlineto{\pgfqpoint{6.276221in}{0.718668in}}%
\pgfpathlineto{\pgfqpoint{6.279576in}{0.714735in}}%
\pgfpathlineto{\pgfqpoint{6.282932in}{0.715515in}}%
\pgfpathlineto{\pgfqpoint{6.289644in}{0.730598in}}%
\pgfpathlineto{\pgfqpoint{6.292999in}{0.734403in}}%
\pgfpathlineto{\pgfqpoint{6.296355in}{0.734465in}}%
\pgfpathlineto{\pgfqpoint{6.299711in}{0.728164in}}%
\pgfpathlineto{\pgfqpoint{6.303067in}{0.716555in}}%
\pgfpathlineto{\pgfqpoint{6.306423in}{0.700094in}}%
\pgfpathlineto{\pgfqpoint{6.309778in}{0.696253in}}%
\pgfpathlineto{\pgfqpoint{6.313134in}{0.708073in}}%
\pgfpathlineto{\pgfqpoint{6.316490in}{0.716273in}}%
\pgfpathlineto{\pgfqpoint{6.319846in}{0.719953in}}%
\pgfpathlineto{\pgfqpoint{6.326557in}{0.721796in}}%
\pgfpathlineto{\pgfqpoint{6.333269in}{0.726913in}}%
\pgfpathlineto{\pgfqpoint{6.336625in}{0.727329in}}%
\pgfpathlineto{\pgfqpoint{6.339980in}{0.725229in}}%
\pgfpathlineto{\pgfqpoint{6.343336in}{0.713289in}}%
\pgfpathlineto{\pgfqpoint{6.346692in}{0.704514in}}%
\pgfpathlineto{\pgfqpoint{6.353403in}{0.723190in}}%
\pgfpathlineto{\pgfqpoint{6.356759in}{0.724017in}}%
\pgfpathlineto{\pgfqpoint{6.363471in}{0.709760in}}%
\pgfpathlineto{\pgfqpoint{6.366827in}{0.720387in}}%
\pgfpathlineto{\pgfqpoint{6.370182in}{0.733638in}}%
\pgfpathlineto{\pgfqpoint{6.373538in}{0.738254in}}%
\pgfpathlineto{\pgfqpoint{6.376894in}{0.732327in}}%
\pgfpathlineto{\pgfqpoint{6.383605in}{0.708423in}}%
\pgfpathlineto{\pgfqpoint{6.383605in}{0.708423in}}%
\pgfusepath{stroke}%
\end{pgfscope}%
\begin{pgfscope}%
\pgfpathrectangle{\pgfqpoint{4.621821in}{0.320625in}}{\pgfqpoint{1.845679in}{1.164375in}}%
\pgfusepath{clip}%
\pgfsetrectcap%
\pgfsetroundjoin%
\pgfsetlinewidth{0.602250pt}%
\definecolor{currentstroke}{rgb}{0.000000,0.529412,0.423529}%
\pgfsetstrokecolor{currentstroke}%
\pgfsetstrokeopacity{0.120000}%
\pgfsetdash{}{0pt}%
\pgfpathmoveto{\pgfqpoint{4.705715in}{1.468694in}}%
\pgfpathlineto{\pgfqpoint{4.719139in}{1.443565in}}%
\pgfpathlineto{\pgfqpoint{4.729206in}{1.422261in}}%
\pgfpathlineto{\pgfqpoint{4.749341in}{1.376161in}}%
\pgfpathlineto{\pgfqpoint{4.776187in}{1.315436in}}%
\pgfpathlineto{\pgfqpoint{4.799677in}{1.265399in}}%
\pgfpathlineto{\pgfqpoint{4.826524in}{1.211211in}}%
\pgfpathlineto{\pgfqpoint{4.856726in}{1.152950in}}%
\pgfpathlineto{\pgfqpoint{4.890283in}{1.090600in}}%
\pgfpathlineto{\pgfqpoint{4.927197in}{1.024254in}}%
\pgfpathlineto{\pgfqpoint{4.967466in}{0.954099in}}%
\pgfpathlineto{\pgfqpoint{5.011091in}{0.880300in}}%
\pgfpathlineto{\pgfqpoint{5.048005in}{0.819148in}}%
\pgfpathlineto{\pgfqpoint{5.074851in}{0.775440in}}%
\pgfpathlineto{\pgfqpoint{5.084919in}{0.758470in}}%
\pgfpathlineto{\pgfqpoint{5.091630in}{0.749057in}}%
\pgfpathlineto{\pgfqpoint{5.101698in}{0.737119in}}%
\pgfpathlineto{\pgfqpoint{5.108409in}{0.726010in}}%
\pgfpathlineto{\pgfqpoint{5.115121in}{0.711964in}}%
\pgfpathlineto{\pgfqpoint{5.128544in}{0.681915in}}%
\pgfpathlineto{\pgfqpoint{5.131900in}{0.676335in}}%
\pgfpathlineto{\pgfqpoint{5.135255in}{0.672485in}}%
\pgfpathlineto{\pgfqpoint{5.138611in}{0.670882in}}%
\pgfpathlineto{\pgfqpoint{5.141967in}{0.671709in}}%
\pgfpathlineto{\pgfqpoint{5.152034in}{0.679429in}}%
\pgfpathlineto{\pgfqpoint{5.155390in}{0.679232in}}%
\pgfpathlineto{\pgfqpoint{5.158746in}{0.675024in}}%
\pgfpathlineto{\pgfqpoint{5.165457in}{0.656851in}}%
\pgfpathlineto{\pgfqpoint{5.168813in}{0.657362in}}%
\pgfpathlineto{\pgfqpoint{5.175525in}{0.679925in}}%
\pgfpathlineto{\pgfqpoint{5.178880in}{0.687911in}}%
\pgfpathlineto{\pgfqpoint{5.182236in}{0.692947in}}%
\pgfpathlineto{\pgfqpoint{5.185592in}{0.695407in}}%
\pgfpathlineto{\pgfqpoint{5.188948in}{0.695436in}}%
\pgfpathlineto{\pgfqpoint{5.192304in}{0.692888in}}%
\pgfpathlineto{\pgfqpoint{5.195659in}{0.687191in}}%
\pgfpathlineto{\pgfqpoint{5.199015in}{0.676867in}}%
\pgfpathlineto{\pgfqpoint{5.202371in}{0.657560in}}%
\pgfpathlineto{\pgfqpoint{5.205727in}{0.607973in}}%
\pgfpathlineto{\pgfqpoint{5.209082in}{0.644755in}}%
\pgfpathlineto{\pgfqpoint{5.212438in}{0.674746in}}%
\pgfpathlineto{\pgfqpoint{5.215794in}{0.690639in}}%
\pgfpathlineto{\pgfqpoint{5.219150in}{0.700375in}}%
\pgfpathlineto{\pgfqpoint{5.222506in}{0.706055in}}%
\pgfpathlineto{\pgfqpoint{5.225861in}{0.708190in}}%
\pgfpathlineto{\pgfqpoint{5.229217in}{0.706303in}}%
\pgfpathlineto{\pgfqpoint{5.232573in}{0.698515in}}%
\pgfpathlineto{\pgfqpoint{5.235929in}{0.684674in}}%
\pgfpathlineto{\pgfqpoint{5.239285in}{0.678482in}}%
\pgfpathlineto{\pgfqpoint{5.245996in}{0.698279in}}%
\pgfpathlineto{\pgfqpoint{5.249352in}{0.700874in}}%
\pgfpathlineto{\pgfqpoint{5.252708in}{0.696420in}}%
\pgfpathlineto{\pgfqpoint{5.256063in}{0.681747in}}%
\pgfpathlineto{\pgfqpoint{5.259419in}{0.642148in}}%
\pgfpathlineto{\pgfqpoint{5.266131in}{0.697265in}}%
\pgfpathlineto{\pgfqpoint{5.269487in}{0.710040in}}%
\pgfpathlineto{\pgfqpoint{5.272842in}{0.712372in}}%
\pgfpathlineto{\pgfqpoint{5.276198in}{0.700267in}}%
\pgfpathlineto{\pgfqpoint{5.279554in}{0.653958in}}%
\pgfpathlineto{\pgfqpoint{5.282910in}{0.691144in}}%
\pgfpathlineto{\pgfqpoint{5.286265in}{0.717286in}}%
\pgfpathlineto{\pgfqpoint{5.289621in}{0.727987in}}%
\pgfpathlineto{\pgfqpoint{5.292977in}{0.730300in}}%
\pgfpathlineto{\pgfqpoint{5.296333in}{0.724949in}}%
\pgfpathlineto{\pgfqpoint{5.299689in}{0.707997in}}%
\pgfpathlineto{\pgfqpoint{5.303044in}{0.641852in}}%
\pgfpathlineto{\pgfqpoint{5.306400in}{0.699429in}}%
\pgfpathlineto{\pgfqpoint{5.309756in}{0.728153in}}%
\pgfpathlineto{\pgfqpoint{5.313112in}{0.740031in}}%
\pgfpathlineto{\pgfqpoint{5.316467in}{0.741217in}}%
\pgfpathlineto{\pgfqpoint{5.319823in}{0.729040in}}%
\pgfpathlineto{\pgfqpoint{5.323179in}{0.700368in}}%
\pgfpathlineto{\pgfqpoint{5.326535in}{0.643102in}}%
\pgfpathlineto{\pgfqpoint{5.329891in}{0.698305in}}%
\pgfpathlineto{\pgfqpoint{5.333246in}{0.710104in}}%
\pgfpathlineto{\pgfqpoint{5.336602in}{0.705462in}}%
\pgfpathlineto{\pgfqpoint{5.339958in}{0.676769in}}%
\pgfpathlineto{\pgfqpoint{5.343314in}{0.667153in}}%
\pgfpathlineto{\pgfqpoint{5.346669in}{0.714025in}}%
\pgfpathlineto{\pgfqpoint{5.350025in}{0.730385in}}%
\pgfpathlineto{\pgfqpoint{5.353381in}{0.732344in}}%
\pgfpathlineto{\pgfqpoint{5.356737in}{0.718226in}}%
\pgfpathlineto{\pgfqpoint{5.363448in}{0.647209in}}%
\pgfpathlineto{\pgfqpoint{5.366804in}{0.694069in}}%
\pgfpathlineto{\pgfqpoint{5.370160in}{0.703161in}}%
\pgfpathlineto{\pgfqpoint{5.373516in}{0.696243in}}%
\pgfpathlineto{\pgfqpoint{5.376871in}{0.660512in}}%
\pgfpathlineto{\pgfqpoint{5.380227in}{0.670634in}}%
\pgfpathlineto{\pgfqpoint{5.383583in}{0.709873in}}%
\pgfpathlineto{\pgfqpoint{5.386939in}{0.724087in}}%
\pgfpathlineto{\pgfqpoint{5.390295in}{0.723886in}}%
\pgfpathlineto{\pgfqpoint{5.393650in}{0.705566in}}%
\pgfpathlineto{\pgfqpoint{5.397006in}{0.638239in}}%
\pgfpathlineto{\pgfqpoint{5.400362in}{0.688494in}}%
\pgfpathlineto{\pgfqpoint{5.403718in}{0.712375in}}%
\pgfpathlineto{\pgfqpoint{5.407074in}{0.718797in}}%
\pgfpathlineto{\pgfqpoint{5.410429in}{0.714251in}}%
\pgfpathlineto{\pgfqpoint{5.413785in}{0.693976in}}%
\pgfpathlineto{\pgfqpoint{5.417141in}{0.604116in}}%
\pgfpathlineto{\pgfqpoint{5.420497in}{0.704813in}}%
\pgfpathlineto{\pgfqpoint{5.423852in}{0.727150in}}%
\pgfpathlineto{\pgfqpoint{5.427208in}{0.734542in}}%
\pgfpathlineto{\pgfqpoint{5.430564in}{0.727548in}}%
\pgfpathlineto{\pgfqpoint{5.437276in}{0.666774in}}%
\pgfpathlineto{\pgfqpoint{5.440631in}{0.717498in}}%
\pgfpathlineto{\pgfqpoint{5.443987in}{0.732088in}}%
\pgfpathlineto{\pgfqpoint{5.447343in}{0.734861in}}%
\pgfpathlineto{\pgfqpoint{5.450699in}{0.727629in}}%
\pgfpathlineto{\pgfqpoint{5.454054in}{0.702834in}}%
\pgfpathlineto{\pgfqpoint{5.457410in}{0.659010in}}%
\pgfpathlineto{\pgfqpoint{5.460766in}{0.723577in}}%
\pgfpathlineto{\pgfqpoint{5.464122in}{0.743188in}}%
\pgfpathlineto{\pgfqpoint{5.467478in}{0.749229in}}%
\pgfpathlineto{\pgfqpoint{5.470833in}{0.741254in}}%
\pgfpathlineto{\pgfqpoint{5.474189in}{0.714642in}}%
\pgfpathlineto{\pgfqpoint{5.477545in}{0.638526in}}%
\pgfpathlineto{\pgfqpoint{5.480901in}{0.716832in}}%
\pgfpathlineto{\pgfqpoint{5.484256in}{0.732245in}}%
\pgfpathlineto{\pgfqpoint{5.487612in}{0.733290in}}%
\pgfpathlineto{\pgfqpoint{5.490968in}{0.721039in}}%
\pgfpathlineto{\pgfqpoint{5.494324in}{0.670355in}}%
\pgfpathlineto{\pgfqpoint{5.501035in}{0.739004in}}%
\pgfpathlineto{\pgfqpoint{5.504391in}{0.750825in}}%
\pgfpathlineto{\pgfqpoint{5.507747in}{0.748890in}}%
\pgfpathlineto{\pgfqpoint{5.511103in}{0.732250in}}%
\pgfpathlineto{\pgfqpoint{5.514458in}{0.695210in}}%
\pgfpathlineto{\pgfqpoint{5.517814in}{0.664759in}}%
\pgfpathlineto{\pgfqpoint{5.521170in}{0.706313in}}%
\pgfpathlineto{\pgfqpoint{5.524526in}{0.711794in}}%
\pgfpathlineto{\pgfqpoint{5.527882in}{0.697657in}}%
\pgfpathlineto{\pgfqpoint{5.531237in}{0.527072in}}%
\pgfpathlineto{\pgfqpoint{5.534593in}{0.706691in}}%
\pgfpathlineto{\pgfqpoint{5.537949in}{0.731227in}}%
\pgfpathlineto{\pgfqpoint{5.541305in}{0.738954in}}%
\pgfpathlineto{\pgfqpoint{5.544660in}{0.730825in}}%
\pgfpathlineto{\pgfqpoint{5.548016in}{0.708570in}}%
\pgfpathlineto{\pgfqpoint{5.551372in}{0.647785in}}%
\pgfpathlineto{\pgfqpoint{5.554728in}{0.675651in}}%
\pgfpathlineto{\pgfqpoint{5.558084in}{0.696352in}}%
\pgfpathlineto{\pgfqpoint{5.561439in}{0.694036in}}%
\pgfpathlineto{\pgfqpoint{5.564795in}{0.664702in}}%
\pgfpathlineto{\pgfqpoint{5.568151in}{0.662169in}}%
\pgfpathlineto{\pgfqpoint{5.571507in}{0.706950in}}%
\pgfpathlineto{\pgfqpoint{5.574863in}{0.722787in}}%
\pgfpathlineto{\pgfqpoint{5.578218in}{0.723760in}}%
\pgfpathlineto{\pgfqpoint{5.581574in}{0.707349in}}%
\pgfpathlineto{\pgfqpoint{5.584930in}{0.659870in}}%
\pgfpathlineto{\pgfqpoint{5.588286in}{0.672093in}}%
\pgfpathlineto{\pgfqpoint{5.591641in}{0.701787in}}%
\pgfpathlineto{\pgfqpoint{5.594997in}{0.708893in}}%
\pgfpathlineto{\pgfqpoint{5.598353in}{0.703291in}}%
\pgfpathlineto{\pgfqpoint{5.601709in}{0.677802in}}%
\pgfpathlineto{\pgfqpoint{5.605065in}{0.647169in}}%
\pgfpathlineto{\pgfqpoint{5.608420in}{0.704297in}}%
\pgfpathlineto{\pgfqpoint{5.611776in}{0.722633in}}%
\pgfpathlineto{\pgfqpoint{5.615132in}{0.726827in}}%
\pgfpathlineto{\pgfqpoint{5.618488in}{0.714446in}}%
\pgfpathlineto{\pgfqpoint{5.621843in}{0.668205in}}%
\pgfpathlineto{\pgfqpoint{5.625199in}{0.686092in}}%
\pgfpathlineto{\pgfqpoint{5.628555in}{0.717596in}}%
\pgfpathlineto{\pgfqpoint{5.631911in}{0.727943in}}%
\pgfpathlineto{\pgfqpoint{5.635267in}{0.728009in}}%
\pgfpathlineto{\pgfqpoint{5.638622in}{0.717492in}}%
\pgfpathlineto{\pgfqpoint{5.641978in}{0.682741in}}%
\pgfpathlineto{\pgfqpoint{5.645334in}{0.682007in}}%
\pgfpathlineto{\pgfqpoint{5.648690in}{0.724076in}}%
\pgfpathlineto{\pgfqpoint{5.652045in}{0.739603in}}%
\pgfpathlineto{\pgfqpoint{5.655401in}{0.742620in}}%
\pgfpathlineto{\pgfqpoint{5.658757in}{0.730261in}}%
\pgfpathlineto{\pgfqpoint{5.662113in}{0.688743in}}%
\pgfpathlineto{\pgfqpoint{5.665469in}{0.693344in}}%
\pgfpathlineto{\pgfqpoint{5.668824in}{0.727718in}}%
\pgfpathlineto{\pgfqpoint{5.672180in}{0.738189in}}%
\pgfpathlineto{\pgfqpoint{5.675536in}{0.737606in}}%
\pgfpathlineto{\pgfqpoint{5.678892in}{0.725201in}}%
\pgfpathlineto{\pgfqpoint{5.682247in}{0.679896in}}%
\pgfpathlineto{\pgfqpoint{5.685603in}{0.705135in}}%
\pgfpathlineto{\pgfqpoint{5.688959in}{0.739329in}}%
\pgfpathlineto{\pgfqpoint{5.692315in}{0.752311in}}%
\pgfpathlineto{\pgfqpoint{5.695671in}{0.752519in}}%
\pgfpathlineto{\pgfqpoint{5.699026in}{0.737614in}}%
\pgfpathlineto{\pgfqpoint{5.705738in}{0.671921in}}%
\pgfpathlineto{\pgfqpoint{5.709094in}{0.715634in}}%
\pgfpathlineto{\pgfqpoint{5.712449in}{0.723844in}}%
\pgfpathlineto{\pgfqpoint{5.715805in}{0.716054in}}%
\pgfpathlineto{\pgfqpoint{5.719161in}{0.676903in}}%
\pgfpathlineto{\pgfqpoint{5.722517in}{0.694786in}}%
\pgfpathlineto{\pgfqpoint{5.725873in}{0.731665in}}%
\pgfpathlineto{\pgfqpoint{5.729228in}{0.745076in}}%
\pgfpathlineto{\pgfqpoint{5.732584in}{0.743952in}}%
\pgfpathlineto{\pgfqpoint{5.735940in}{0.728103in}}%
\pgfpathlineto{\pgfqpoint{5.739296in}{0.695699in}}%
\pgfpathlineto{\pgfqpoint{5.742652in}{0.626950in}}%
\pgfpathlineto{\pgfqpoint{5.746007in}{0.690701in}}%
\pgfpathlineto{\pgfqpoint{5.749363in}{0.697038in}}%
\pgfpathlineto{\pgfqpoint{5.752719in}{0.679264in}}%
\pgfpathlineto{\pgfqpoint{5.756075in}{0.632252in}}%
\pgfpathlineto{\pgfqpoint{5.759430in}{0.703131in}}%
\pgfpathlineto{\pgfqpoint{5.762786in}{0.723817in}}%
\pgfpathlineto{\pgfqpoint{5.766142in}{0.728843in}}%
\pgfpathlineto{\pgfqpoint{5.769498in}{0.717272in}}%
\pgfpathlineto{\pgfqpoint{5.772854in}{0.688255in}}%
\pgfpathlineto{\pgfqpoint{5.776209in}{0.610943in}}%
\pgfpathlineto{\pgfqpoint{5.779565in}{0.686463in}}%
\pgfpathlineto{\pgfqpoint{5.782921in}{0.698747in}}%
\pgfpathlineto{\pgfqpoint{5.786277in}{0.694394in}}%
\pgfpathlineto{\pgfqpoint{5.789632in}{0.666430in}}%
\pgfpathlineto{\pgfqpoint{5.792988in}{0.653583in}}%
\pgfpathlineto{\pgfqpoint{5.796344in}{0.702110in}}%
\pgfpathlineto{\pgfqpoint{5.799700in}{0.718765in}}%
\pgfpathlineto{\pgfqpoint{5.803056in}{0.720918in}}%
\pgfpathlineto{\pgfqpoint{5.806411in}{0.705130in}}%
\pgfpathlineto{\pgfqpoint{5.809767in}{0.647320in}}%
\pgfpathlineto{\pgfqpoint{5.813123in}{0.684264in}}%
\pgfpathlineto{\pgfqpoint{5.816479in}{0.711032in}}%
\pgfpathlineto{\pgfqpoint{5.819834in}{0.719204in}}%
\pgfpathlineto{\pgfqpoint{5.823190in}{0.716937in}}%
\pgfpathlineto{\pgfqpoint{5.826546in}{0.702024in}}%
\pgfpathlineto{\pgfqpoint{5.829902in}{0.640308in}}%
\pgfpathlineto{\pgfqpoint{5.833258in}{0.692965in}}%
\pgfpathlineto{\pgfqpoint{5.836613in}{0.722210in}}%
\pgfpathlineto{\pgfqpoint{5.839969in}{0.733300in}}%
\pgfpathlineto{\pgfqpoint{5.843325in}{0.731504in}}%
\pgfpathlineto{\pgfqpoint{5.846681in}{0.710404in}}%
\pgfpathlineto{\pgfqpoint{5.850036in}{0.594505in}}%
\pgfpathlineto{\pgfqpoint{5.853392in}{0.712502in}}%
\pgfpathlineto{\pgfqpoint{5.856748in}{0.732451in}}%
\pgfpathlineto{\pgfqpoint{5.860104in}{0.738597in}}%
\pgfpathlineto{\pgfqpoint{5.863460in}{0.735436in}}%
\pgfpathlineto{\pgfqpoint{5.866815in}{0.720196in}}%
\pgfpathlineto{\pgfqpoint{5.870171in}{0.660297in}}%
\pgfpathlineto{\pgfqpoint{5.873527in}{0.709602in}}%
\pgfpathlineto{\pgfqpoint{5.876883in}{0.739295in}}%
\pgfpathlineto{\pgfqpoint{5.880238in}{0.750838in}}%
\pgfpathlineto{\pgfqpoint{5.883594in}{0.750174in}}%
\pgfpathlineto{\pgfqpoint{5.886950in}{0.733775in}}%
\pgfpathlineto{\pgfqpoint{5.890306in}{0.686242in}}%
\pgfpathlineto{\pgfqpoint{5.893662in}{0.698677in}}%
\pgfpathlineto{\pgfqpoint{5.897017in}{0.728380in}}%
\pgfpathlineto{\pgfqpoint{5.900373in}{0.735532in}}%
\pgfpathlineto{\pgfqpoint{5.903729in}{0.730029in}}%
\pgfpathlineto{\pgfqpoint{5.907085in}{0.704884in}}%
\pgfpathlineto{\pgfqpoint{5.910441in}{0.672240in}}%
\pgfpathlineto{\pgfqpoint{5.913796in}{0.730543in}}%
\pgfpathlineto{\pgfqpoint{5.917152in}{0.749068in}}%
\pgfpathlineto{\pgfqpoint{5.920508in}{0.753418in}}%
\pgfpathlineto{\pgfqpoint{5.923864in}{0.742729in}}%
\pgfpathlineto{\pgfqpoint{5.927219in}{0.718699in}}%
\pgfpathlineto{\pgfqpoint{5.930575in}{0.649817in}}%
\pgfpathlineto{\pgfqpoint{5.933931in}{0.688286in}}%
\pgfpathlineto{\pgfqpoint{5.937287in}{0.705523in}}%
\pgfpathlineto{\pgfqpoint{5.940643in}{0.698889in}}%
\pgfpathlineto{\pgfqpoint{5.943998in}{0.646639in}}%
\pgfpathlineto{\pgfqpoint{5.947354in}{0.693240in}}%
\pgfpathlineto{\pgfqpoint{5.950710in}{0.724736in}}%
\pgfpathlineto{\pgfqpoint{5.954066in}{0.736049in}}%
\pgfpathlineto{\pgfqpoint{5.957421in}{0.731756in}}%
\pgfpathlineto{\pgfqpoint{5.960777in}{0.713264in}}%
\pgfpathlineto{\pgfqpoint{5.964133in}{0.672217in}}%
\pgfpathlineto{\pgfqpoint{5.967489in}{0.653181in}}%
\pgfpathlineto{\pgfqpoint{5.970845in}{0.687602in}}%
\pgfpathlineto{\pgfqpoint{5.974200in}{0.689404in}}%
\pgfpathlineto{\pgfqpoint{5.980912in}{0.644391in}}%
\pgfpathlineto{\pgfqpoint{5.984268in}{0.698999in}}%
\pgfpathlineto{\pgfqpoint{5.987623in}{0.717131in}}%
\pgfpathlineto{\pgfqpoint{5.990979in}{0.720127in}}%
\pgfpathlineto{\pgfqpoint{5.994335in}{0.705639in}}%
\pgfpathlineto{\pgfqpoint{5.997691in}{0.662803in}}%
\pgfpathlineto{\pgfqpoint{6.001047in}{0.665813in}}%
\pgfpathlineto{\pgfqpoint{6.004402in}{0.698939in}}%
\pgfpathlineto{\pgfqpoint{6.007758in}{0.707653in}}%
\pgfpathlineto{\pgfqpoint{6.011114in}{0.704138in}}%
\pgfpathlineto{\pgfqpoint{6.014470in}{0.684453in}}%
\pgfpathlineto{\pgfqpoint{6.017825in}{0.593781in}}%
\pgfpathlineto{\pgfqpoint{6.021181in}{0.695622in}}%
\pgfpathlineto{\pgfqpoint{6.024537in}{0.717978in}}%
\pgfpathlineto{\pgfqpoint{6.027893in}{0.725102in}}%
\pgfpathlineto{\pgfqpoint{6.031249in}{0.717083in}}%
\pgfpathlineto{\pgfqpoint{6.034604in}{0.681375in}}%
\pgfpathlineto{\pgfqpoint{6.037960in}{0.676841in}}%
\pgfpathlineto{\pgfqpoint{6.041316in}{0.716453in}}%
\pgfpathlineto{\pgfqpoint{6.044672in}{0.729647in}}%
\pgfpathlineto{\pgfqpoint{6.048027in}{0.732316in}}%
\pgfpathlineto{\pgfqpoint{6.051383in}{0.725862in}}%
\pgfpathlineto{\pgfqpoint{6.054739in}{0.704389in}}%
\pgfpathlineto{\pgfqpoint{6.058095in}{0.617248in}}%
\pgfpathlineto{\pgfqpoint{6.061451in}{0.714879in}}%
\pgfpathlineto{\pgfqpoint{6.064806in}{0.737214in}}%
\pgfpathlineto{\pgfqpoint{6.068162in}{0.745300in}}%
\pgfpathlineto{\pgfqpoint{6.071518in}{0.740624in}}%
\pgfpathlineto{\pgfqpoint{6.074874in}{0.716633in}}%
\pgfpathlineto{\pgfqpoint{6.078230in}{0.634482in}}%
\pgfpathlineto{\pgfqpoint{6.081585in}{0.720923in}}%
\pgfpathlineto{\pgfqpoint{6.084941in}{0.738707in}}%
\pgfpathlineto{\pgfqpoint{6.088297in}{0.742931in}}%
\pgfpathlineto{\pgfqpoint{6.091653in}{0.736905in}}%
\pgfpathlineto{\pgfqpoint{6.095008in}{0.714446in}}%
\pgfpathlineto{\pgfqpoint{6.098364in}{0.651839in}}%
\pgfpathlineto{\pgfqpoint{6.101720in}{0.729917in}}%
\pgfpathlineto{\pgfqpoint{6.105076in}{0.750876in}}%
\pgfpathlineto{\pgfqpoint{6.108432in}{0.757640in}}%
\pgfpathlineto{\pgfqpoint{6.111787in}{0.750514in}}%
\pgfpathlineto{\pgfqpoint{6.115143in}{0.729099in}}%
\pgfpathlineto{\pgfqpoint{6.118499in}{0.669663in}}%
\pgfpathlineto{\pgfqpoint{6.121855in}{0.697059in}}%
\pgfpathlineto{\pgfqpoint{6.125210in}{0.719149in}}%
\pgfpathlineto{\pgfqpoint{6.128566in}{0.719084in}}%
\pgfpathlineto{\pgfqpoint{6.131922in}{0.697819in}}%
\pgfpathlineto{\pgfqpoint{6.135278in}{0.657446in}}%
\pgfpathlineto{\pgfqpoint{6.138634in}{0.722180in}}%
\pgfpathlineto{\pgfqpoint{6.141989in}{0.741753in}}%
\pgfpathlineto{\pgfqpoint{6.145345in}{0.746149in}}%
\pgfpathlineto{\pgfqpoint{6.148701in}{0.734872in}}%
\pgfpathlineto{\pgfqpoint{6.152057in}{0.711439in}}%
\pgfpathlineto{\pgfqpoint{6.155412in}{0.653827in}}%
\pgfpathlineto{\pgfqpoint{6.158768in}{0.669628in}}%
\pgfpathlineto{\pgfqpoint{6.162124in}{0.688310in}}%
\pgfpathlineto{\pgfqpoint{6.165480in}{0.677512in}}%
\pgfpathlineto{\pgfqpoint{6.168836in}{0.569859in}}%
\pgfpathlineto{\pgfqpoint{6.172191in}{0.693277in}}%
\pgfpathlineto{\pgfqpoint{6.175547in}{0.717555in}}%
\pgfpathlineto{\pgfqpoint{6.178903in}{0.725075in}}%
\pgfpathlineto{\pgfqpoint{6.182259in}{0.716067in}}%
\pgfpathlineto{\pgfqpoint{6.185614in}{0.690624in}}%
\pgfpathlineto{\pgfqpoint{6.188970in}{0.547750in}}%
\pgfpathlineto{\pgfqpoint{6.192326in}{0.678896in}}%
\pgfpathlineto{\pgfqpoint{6.195682in}{0.694602in}}%
\pgfpathlineto{\pgfqpoint{6.199038in}{0.692724in}}%
\pgfpathlineto{\pgfqpoint{6.202393in}{0.670790in}}%
\pgfpathlineto{\pgfqpoint{6.205749in}{0.628328in}}%
\pgfpathlineto{\pgfqpoint{6.209105in}{0.693819in}}%
\pgfpathlineto{\pgfqpoint{6.212461in}{0.713403in}}%
\pgfpathlineto{\pgfqpoint{6.215816in}{0.718002in}}%
\pgfpathlineto{\pgfqpoint{6.219172in}{0.705298in}}%
\pgfpathlineto{\pgfqpoint{6.222528in}{0.656526in}}%
\pgfpathlineto{\pgfqpoint{6.225884in}{0.679835in}}%
\pgfpathlineto{\pgfqpoint{6.229240in}{0.710200in}}%
\pgfpathlineto{\pgfqpoint{6.232595in}{0.720304in}}%
\pgfpathlineto{\pgfqpoint{6.235951in}{0.720308in}}%
\pgfpathlineto{\pgfqpoint{6.239307in}{0.709852in}}%
\pgfpathlineto{\pgfqpoint{6.242663in}{0.675626in}}%
\pgfpathlineto{\pgfqpoint{6.246019in}{0.673128in}}%
\pgfpathlineto{\pgfqpoint{6.249374in}{0.715862in}}%
\pgfpathlineto{\pgfqpoint{6.252730in}{0.731584in}}%
\pgfpathlineto{\pgfqpoint{6.256086in}{0.734831in}}%
\pgfpathlineto{\pgfqpoint{6.259442in}{0.722297in}}%
\pgfpathlineto{\pgfqpoint{6.262797in}{0.669515in}}%
\pgfpathlineto{\pgfqpoint{6.269509in}{0.731657in}}%
\pgfpathlineto{\pgfqpoint{6.272865in}{0.742031in}}%
\pgfpathlineto{\pgfqpoint{6.276221in}{0.742875in}}%
\pgfpathlineto{\pgfqpoint{6.279576in}{0.734332in}}%
\pgfpathlineto{\pgfqpoint{6.282932in}{0.707713in}}%
\pgfpathlineto{\pgfqpoint{6.286288in}{0.672444in}}%
\pgfpathlineto{\pgfqpoint{6.289644in}{0.731370in}}%
\pgfpathlineto{\pgfqpoint{6.292999in}{0.750365in}}%
\pgfpathlineto{\pgfqpoint{6.296355in}{0.756436in}}%
\pgfpathlineto{\pgfqpoint{6.299711in}{0.749241in}}%
\pgfpathlineto{\pgfqpoint{6.303067in}{0.725079in}}%
\pgfpathlineto{\pgfqpoint{6.306423in}{0.551771in}}%
\pgfpathlineto{\pgfqpoint{6.309778in}{0.717251in}}%
\pgfpathlineto{\pgfqpoint{6.313134in}{0.734656in}}%
\pgfpathlineto{\pgfqpoint{6.316490in}{0.736065in}}%
\pgfpathlineto{\pgfqpoint{6.319846in}{0.723272in}}%
\pgfpathlineto{\pgfqpoint{6.323201in}{0.665299in}}%
\pgfpathlineto{\pgfqpoint{6.326557in}{0.714417in}}%
\pgfpathlineto{\pgfqpoint{6.329913in}{0.744342in}}%
\pgfpathlineto{\pgfqpoint{6.333269in}{0.755216in}}%
\pgfpathlineto{\pgfqpoint{6.336625in}{0.751821in}}%
\pgfpathlineto{\pgfqpoint{6.339980in}{0.735174in}}%
\pgfpathlineto{\pgfqpoint{6.343336in}{0.702698in}}%
\pgfpathlineto{\pgfqpoint{6.346692in}{0.620118in}}%
\pgfpathlineto{\pgfqpoint{6.350048in}{0.691847in}}%
\pgfpathlineto{\pgfqpoint{6.353403in}{0.695587in}}%
\pgfpathlineto{\pgfqpoint{6.356759in}{0.665796in}}%
\pgfpathlineto{\pgfqpoint{6.360115in}{0.674449in}}%
\pgfpathlineto{\pgfqpoint{6.363471in}{0.716540in}}%
\pgfpathlineto{\pgfqpoint{6.366827in}{0.731966in}}%
\pgfpathlineto{\pgfqpoint{6.370182in}{0.732155in}}%
\pgfpathlineto{\pgfqpoint{6.383605in}{0.688136in}}%
\pgfpathlineto{\pgfqpoint{6.383605in}{0.688136in}}%
\pgfusepath{stroke}%
\end{pgfscope}%
\begin{pgfscope}%
\pgfpathrectangle{\pgfqpoint{4.621821in}{0.320625in}}{\pgfqpoint{1.845679in}{1.164375in}}%
\pgfusepath{clip}%
\pgfsetrectcap%
\pgfsetroundjoin%
\pgfsetlinewidth{0.602250pt}%
\definecolor{currentstroke}{rgb}{0.000000,0.529412,0.423529}%
\pgfsetstrokecolor{currentstroke}%
\pgfsetstrokeopacity{0.120000}%
\pgfsetdash{}{0pt}%
\pgfpathmoveto{\pgfqpoint{4.705715in}{1.469525in}}%
\pgfpathlineto{\pgfqpoint{4.719139in}{1.445279in}}%
\pgfpathlineto{\pgfqpoint{4.729206in}{1.424522in}}%
\pgfpathlineto{\pgfqpoint{4.745985in}{1.386302in}}%
\pgfpathlineto{\pgfqpoint{4.772831in}{1.324674in}}%
\pgfpathlineto{\pgfqpoint{4.792966in}{1.281277in}}%
\pgfpathlineto{\pgfqpoint{4.816456in}{1.233644in}}%
\pgfpathlineto{\pgfqpoint{4.843302in}{1.182095in}}%
\pgfpathlineto{\pgfqpoint{4.873504in}{1.126593in}}%
\pgfpathlineto{\pgfqpoint{4.910418in}{1.061125in}}%
\pgfpathlineto{\pgfqpoint{4.954043in}{0.985985in}}%
\pgfpathlineto{\pgfqpoint{5.004380in}{0.901341in}}%
\pgfpathlineto{\pgfqpoint{5.064784in}{0.802356in}}%
\pgfpathlineto{\pgfqpoint{5.078207in}{0.779425in}}%
\pgfpathlineto{\pgfqpoint{5.088274in}{0.762505in}}%
\pgfpathlineto{\pgfqpoint{5.094986in}{0.752895in}}%
\pgfpathlineto{\pgfqpoint{5.105053in}{0.739785in}}%
\pgfpathlineto{\pgfqpoint{5.111765in}{0.727982in}}%
\pgfpathlineto{\pgfqpoint{5.121832in}{0.706505in}}%
\pgfpathlineto{\pgfqpoint{5.128544in}{0.692092in}}%
\pgfpathlineto{\pgfqpoint{5.131900in}{0.686091in}}%
\pgfpathlineto{\pgfqpoint{5.135255in}{0.681577in}}%
\pgfpathlineto{\pgfqpoint{5.138611in}{0.678916in}}%
\pgfpathlineto{\pgfqpoint{5.141967in}{0.678125in}}%
\pgfpathlineto{\pgfqpoint{5.152034in}{0.680368in}}%
\pgfpathlineto{\pgfqpoint{5.155390in}{0.679049in}}%
\pgfpathlineto{\pgfqpoint{5.158746in}{0.674274in}}%
\pgfpathlineto{\pgfqpoint{5.162102in}{0.666389in}}%
\pgfpathlineto{\pgfqpoint{5.168813in}{0.644085in}}%
\pgfpathlineto{\pgfqpoint{5.172169in}{0.638524in}}%
\pgfpathlineto{\pgfqpoint{5.175525in}{0.643772in}}%
\pgfpathlineto{\pgfqpoint{5.178880in}{0.651626in}}%
\pgfpathlineto{\pgfqpoint{5.182236in}{0.657236in}}%
\pgfpathlineto{\pgfqpoint{5.185592in}{0.659875in}}%
\pgfpathlineto{\pgfqpoint{5.188948in}{0.659492in}}%
\pgfpathlineto{\pgfqpoint{5.192304in}{0.655764in}}%
\pgfpathlineto{\pgfqpoint{5.195659in}{0.647614in}}%
\pgfpathlineto{\pgfqpoint{5.199015in}{0.631841in}}%
\pgfpathlineto{\pgfqpoint{5.202371in}{0.594713in}}%
\pgfpathlineto{\pgfqpoint{5.205727in}{0.602358in}}%
\pgfpathlineto{\pgfqpoint{5.209082in}{0.639255in}}%
\pgfpathlineto{\pgfqpoint{5.212438in}{0.657410in}}%
\pgfpathlineto{\pgfqpoint{5.215794in}{0.668344in}}%
\pgfpathlineto{\pgfqpoint{5.219150in}{0.674796in}}%
\pgfpathlineto{\pgfqpoint{5.222506in}{0.677470in}}%
\pgfpathlineto{\pgfqpoint{5.225861in}{0.675956in}}%
\pgfpathlineto{\pgfqpoint{5.229217in}{0.668278in}}%
\pgfpathlineto{\pgfqpoint{5.232573in}{0.659187in}}%
\pgfpathlineto{\pgfqpoint{5.235929in}{0.672012in}}%
\pgfpathlineto{\pgfqpoint{5.239285in}{0.689221in}}%
\pgfpathlineto{\pgfqpoint{5.242640in}{0.700482in}}%
\pgfpathlineto{\pgfqpoint{5.245996in}{0.706671in}}%
\pgfpathlineto{\pgfqpoint{5.249352in}{0.708652in}}%
\pgfpathlineto{\pgfqpoint{5.252708in}{0.706503in}}%
\pgfpathlineto{\pgfqpoint{5.256063in}{0.699295in}}%
\pgfpathlineto{\pgfqpoint{5.259419in}{0.683671in}}%
\pgfpathlineto{\pgfqpoint{5.262775in}{0.642406in}}%
\pgfpathlineto{\pgfqpoint{5.266131in}{0.654908in}}%
\pgfpathlineto{\pgfqpoint{5.269487in}{0.691951in}}%
\pgfpathlineto{\pgfqpoint{5.272842in}{0.708648in}}%
\pgfpathlineto{\pgfqpoint{5.276198in}{0.717049in}}%
\pgfpathlineto{\pgfqpoint{5.279554in}{0.719037in}}%
\pgfpathlineto{\pgfqpoint{5.282910in}{0.712888in}}%
\pgfpathlineto{\pgfqpoint{5.286265in}{0.694043in}}%
\pgfpathlineto{\pgfqpoint{5.289621in}{0.654025in}}%
\pgfpathlineto{\pgfqpoint{5.296333in}{0.698040in}}%
\pgfpathlineto{\pgfqpoint{5.299689in}{0.703673in}}%
\pgfpathlineto{\pgfqpoint{5.303044in}{0.698801in}}%
\pgfpathlineto{\pgfqpoint{5.306400in}{0.678470in}}%
\pgfpathlineto{\pgfqpoint{5.309756in}{0.605831in}}%
\pgfpathlineto{\pgfqpoint{5.313112in}{0.689560in}}%
\pgfpathlineto{\pgfqpoint{5.316467in}{0.712111in}}%
\pgfpathlineto{\pgfqpoint{5.319823in}{0.720297in}}%
\pgfpathlineto{\pgfqpoint{5.323179in}{0.715502in}}%
\pgfpathlineto{\pgfqpoint{5.326535in}{0.693961in}}%
\pgfpathlineto{\pgfqpoint{5.329891in}{0.633732in}}%
\pgfpathlineto{\pgfqpoint{5.333246in}{0.686803in}}%
\pgfpathlineto{\pgfqpoint{5.336602in}{0.706399in}}%
\pgfpathlineto{\pgfqpoint{5.339958in}{0.711417in}}%
\pgfpathlineto{\pgfqpoint{5.343314in}{0.706059in}}%
\pgfpathlineto{\pgfqpoint{5.346669in}{0.684921in}}%
\pgfpathlineto{\pgfqpoint{5.350025in}{0.608315in}}%
\pgfpathlineto{\pgfqpoint{5.353381in}{0.697394in}}%
\pgfpathlineto{\pgfqpoint{5.356737in}{0.719494in}}%
\pgfpathlineto{\pgfqpoint{5.360093in}{0.727310in}}%
\pgfpathlineto{\pgfqpoint{5.363448in}{0.721896in}}%
\pgfpathlineto{\pgfqpoint{5.366804in}{0.697412in}}%
\pgfpathlineto{\pgfqpoint{5.370160in}{0.637001in}}%
\pgfpathlineto{\pgfqpoint{5.373516in}{0.703238in}}%
\pgfpathlineto{\pgfqpoint{5.376871in}{0.720840in}}%
\pgfpathlineto{\pgfqpoint{5.380227in}{0.725533in}}%
\pgfpathlineto{\pgfqpoint{5.383583in}{0.720709in}}%
\pgfpathlineto{\pgfqpoint{5.386939in}{0.702009in}}%
\pgfpathlineto{\pgfqpoint{5.390295in}{0.576009in}}%
\pgfpathlineto{\pgfqpoint{5.393650in}{0.705172in}}%
\pgfpathlineto{\pgfqpoint{5.397006in}{0.729995in}}%
\pgfpathlineto{\pgfqpoint{5.400362in}{0.739489in}}%
\pgfpathlineto{\pgfqpoint{5.403718in}{0.736688in}}%
\pgfpathlineto{\pgfqpoint{5.407074in}{0.718439in}}%
\pgfpathlineto{\pgfqpoint{5.410429in}{0.665202in}}%
\pgfpathlineto{\pgfqpoint{5.417141in}{0.716792in}}%
\pgfpathlineto{\pgfqpoint{5.420497in}{0.723194in}}%
\pgfpathlineto{\pgfqpoint{5.423852in}{0.717604in}}%
\pgfpathlineto{\pgfqpoint{5.427208in}{0.693364in}}%
\pgfpathlineto{\pgfqpoint{5.430564in}{0.653677in}}%
\pgfpathlineto{\pgfqpoint{5.433920in}{0.716285in}}%
\pgfpathlineto{\pgfqpoint{5.437276in}{0.735727in}}%
\pgfpathlineto{\pgfqpoint{5.440631in}{0.741337in}}%
\pgfpathlineto{\pgfqpoint{5.443987in}{0.732647in}}%
\pgfpathlineto{\pgfqpoint{5.447343in}{0.710521in}}%
\pgfpathlineto{\pgfqpoint{5.450699in}{0.651065in}}%
\pgfpathlineto{\pgfqpoint{5.454054in}{0.679042in}}%
\pgfpathlineto{\pgfqpoint{5.457410in}{0.700270in}}%
\pgfpathlineto{\pgfqpoint{5.460766in}{0.699786in}}%
\pgfpathlineto{\pgfqpoint{5.464122in}{0.677730in}}%
\pgfpathlineto{\pgfqpoint{5.467478in}{0.642764in}}%
\pgfpathlineto{\pgfqpoint{5.470833in}{0.704170in}}%
\pgfpathlineto{\pgfqpoint{5.474189in}{0.723626in}}%
\pgfpathlineto{\pgfqpoint{5.477545in}{0.728503in}}%
\pgfpathlineto{\pgfqpoint{5.480901in}{0.717994in}}%
\pgfpathlineto{\pgfqpoint{5.484256in}{0.693402in}}%
\pgfpathlineto{\pgfqpoint{5.487612in}{0.612636in}}%
\pgfpathlineto{\pgfqpoint{5.490968in}{0.672040in}}%
\pgfpathlineto{\pgfqpoint{5.494324in}{0.689005in}}%
\pgfpathlineto{\pgfqpoint{5.497680in}{0.686200in}}%
\pgfpathlineto{\pgfqpoint{5.501035in}{0.659074in}}%
\pgfpathlineto{\pgfqpoint{5.504391in}{0.646606in}}%
\pgfpathlineto{\pgfqpoint{5.507747in}{0.695598in}}%
\pgfpathlineto{\pgfqpoint{5.511103in}{0.712751in}}%
\pgfpathlineto{\pgfqpoint{5.514458in}{0.715741in}}%
\pgfpathlineto{\pgfqpoint{5.517814in}{0.702403in}}%
\pgfpathlineto{\pgfqpoint{5.521170in}{0.665615in}}%
\pgfpathlineto{\pgfqpoint{5.524526in}{0.649874in}}%
\pgfpathlineto{\pgfqpoint{5.527882in}{0.689949in}}%
\pgfpathlineto{\pgfqpoint{5.531237in}{0.700519in}}%
\pgfpathlineto{\pgfqpoint{5.534593in}{0.698673in}}%
\pgfpathlineto{\pgfqpoint{5.537949in}{0.682472in}}%
\pgfpathlineto{\pgfqpoint{5.541305in}{0.593120in}}%
\pgfpathlineto{\pgfqpoint{5.544660in}{0.682600in}}%
\pgfpathlineto{\pgfqpoint{5.548016in}{0.708846in}}%
\pgfpathlineto{\pgfqpoint{5.551372in}{0.718504in}}%
\pgfpathlineto{\pgfqpoint{5.554728in}{0.714632in}}%
\pgfpathlineto{\pgfqpoint{5.558084in}{0.691238in}}%
\pgfpathlineto{\pgfqpoint{5.561439in}{0.606784in}}%
\pgfpathlineto{\pgfqpoint{5.564795in}{0.695860in}}%
\pgfpathlineto{\pgfqpoint{5.568151in}{0.714561in}}%
\pgfpathlineto{\pgfqpoint{5.571507in}{0.720031in}}%
\pgfpathlineto{\pgfqpoint{5.574863in}{0.716236in}}%
\pgfpathlineto{\pgfqpoint{5.578218in}{0.699914in}}%
\pgfpathlineto{\pgfqpoint{5.581574in}{0.629672in}}%
\pgfpathlineto{\pgfqpoint{5.584930in}{0.693759in}}%
\pgfpathlineto{\pgfqpoint{5.588286in}{0.721831in}}%
\pgfpathlineto{\pgfqpoint{5.591641in}{0.732893in}}%
\pgfpathlineto{\pgfqpoint{5.594997in}{0.732065in}}%
\pgfpathlineto{\pgfqpoint{5.598353in}{0.715044in}}%
\pgfpathlineto{\pgfqpoint{5.601709in}{0.656393in}}%
\pgfpathlineto{\pgfqpoint{5.605065in}{0.693541in}}%
\pgfpathlineto{\pgfqpoint{5.608420in}{0.719986in}}%
\pgfpathlineto{\pgfqpoint{5.611776in}{0.727980in}}%
\pgfpathlineto{\pgfqpoint{5.615132in}{0.725591in}}%
\pgfpathlineto{\pgfqpoint{5.618488in}{0.710575in}}%
\pgfpathlineto{\pgfqpoint{5.621843in}{0.648319in}}%
\pgfpathlineto{\pgfqpoint{5.625199in}{0.701825in}}%
\pgfpathlineto{\pgfqpoint{5.628555in}{0.731107in}}%
\pgfpathlineto{\pgfqpoint{5.631911in}{0.742464in}}%
\pgfpathlineto{\pgfqpoint{5.635267in}{0.741431in}}%
\pgfpathlineto{\pgfqpoint{5.638622in}{0.725911in}}%
\pgfpathlineto{\pgfqpoint{5.641978in}{0.689789in}}%
\pgfpathlineto{\pgfqpoint{5.645334in}{0.661318in}}%
\pgfpathlineto{\pgfqpoint{5.648690in}{0.704343in}}%
\pgfpathlineto{\pgfqpoint{5.652045in}{0.712774in}}%
\pgfpathlineto{\pgfqpoint{5.655401in}{0.705800in}}%
\pgfpathlineto{\pgfqpoint{5.658757in}{0.671186in}}%
\pgfpathlineto{\pgfqpoint{5.662113in}{0.677817in}}%
\pgfpathlineto{\pgfqpoint{5.665469in}{0.718374in}}%
\pgfpathlineto{\pgfqpoint{5.668824in}{0.733296in}}%
\pgfpathlineto{\pgfqpoint{5.672180in}{0.734483in}}%
\pgfpathlineto{\pgfqpoint{5.675536in}{0.720869in}}%
\pgfpathlineto{\pgfqpoint{5.678892in}{0.693095in}}%
\pgfpathlineto{\pgfqpoint{5.682247in}{0.552465in}}%
\pgfpathlineto{\pgfqpoint{5.685603in}{0.677095in}}%
\pgfpathlineto{\pgfqpoint{5.688959in}{0.689314in}}%
\pgfpathlineto{\pgfqpoint{5.692315in}{0.680419in}}%
\pgfpathlineto{\pgfqpoint{5.695671in}{0.619483in}}%
\pgfpathlineto{\pgfqpoint{5.699026in}{0.678034in}}%
\pgfpathlineto{\pgfqpoint{5.702382in}{0.707875in}}%
\pgfpathlineto{\pgfqpoint{5.705738in}{0.718861in}}%
\pgfpathlineto{\pgfqpoint{5.709094in}{0.715006in}}%
\pgfpathlineto{\pgfqpoint{5.712449in}{0.695793in}}%
\pgfpathlineto{\pgfqpoint{5.715805in}{0.644031in}}%
\pgfpathlineto{\pgfqpoint{5.719161in}{0.660599in}}%
\pgfpathlineto{\pgfqpoint{5.722517in}{0.687301in}}%
\pgfpathlineto{\pgfqpoint{5.725873in}{0.691772in}}%
\pgfpathlineto{\pgfqpoint{5.729228in}{0.681459in}}%
\pgfpathlineto{\pgfqpoint{5.732584in}{0.634511in}}%
\pgfpathlineto{\pgfqpoint{5.739296in}{0.699987in}}%
\pgfpathlineto{\pgfqpoint{5.742652in}{0.712495in}}%
\pgfpathlineto{\pgfqpoint{5.746007in}{0.711264in}}%
\pgfpathlineto{\pgfqpoint{5.749363in}{0.691902in}}%
\pgfpathlineto{\pgfqpoint{5.752719in}{0.605462in}}%
\pgfpathlineto{\pgfqpoint{5.756075in}{0.681866in}}%
\pgfpathlineto{\pgfqpoint{5.759430in}{0.704517in}}%
\pgfpathlineto{\pgfqpoint{5.762786in}{0.711403in}}%
\pgfpathlineto{\pgfqpoint{5.766142in}{0.708478in}}%
\pgfpathlineto{\pgfqpoint{5.769498in}{0.693114in}}%
\pgfpathlineto{\pgfqpoint{5.772854in}{0.630025in}}%
\pgfpathlineto{\pgfqpoint{5.776209in}{0.684432in}}%
\pgfpathlineto{\pgfqpoint{5.779565in}{0.713584in}}%
\pgfpathlineto{\pgfqpoint{5.782921in}{0.724988in}}%
\pgfpathlineto{\pgfqpoint{5.786277in}{0.724256in}}%
\pgfpathlineto{\pgfqpoint{5.789632in}{0.706187in}}%
\pgfpathlineto{\pgfqpoint{5.792988in}{0.623491in}}%
\pgfpathlineto{\pgfqpoint{5.796344in}{0.696814in}}%
\pgfpathlineto{\pgfqpoint{5.799700in}{0.720419in}}%
\pgfpathlineto{\pgfqpoint{5.803056in}{0.728310in}}%
\pgfpathlineto{\pgfqpoint{5.806411in}{0.726906in}}%
\pgfpathlineto{\pgfqpoint{5.809767in}{0.714895in}}%
\pgfpathlineto{\pgfqpoint{5.813123in}{0.675452in}}%
\pgfpathlineto{\pgfqpoint{5.816479in}{0.686589in}}%
\pgfpathlineto{\pgfqpoint{5.819834in}{0.724630in}}%
\pgfpathlineto{\pgfqpoint{5.823190in}{0.739332in}}%
\pgfpathlineto{\pgfqpoint{5.826546in}{0.742222in}}%
\pgfpathlineto{\pgfqpoint{5.829902in}{0.731026in}}%
\pgfpathlineto{\pgfqpoint{5.833258in}{0.700567in}}%
\pgfpathlineto{\pgfqpoint{5.836613in}{0.652622in}}%
\pgfpathlineto{\pgfqpoint{5.839969in}{0.709903in}}%
\pgfpathlineto{\pgfqpoint{5.843325in}{0.723002in}}%
\pgfpathlineto{\pgfqpoint{5.846681in}{0.722476in}}%
\pgfpathlineto{\pgfqpoint{5.850036in}{0.707802in}}%
\pgfpathlineto{\pgfqpoint{5.853392in}{0.636657in}}%
\pgfpathlineto{\pgfqpoint{5.856748in}{0.703878in}}%
\pgfpathlineto{\pgfqpoint{5.860104in}{0.731782in}}%
\pgfpathlineto{\pgfqpoint{5.863460in}{0.742162in}}%
\pgfpathlineto{\pgfqpoint{5.866815in}{0.739006in}}%
\pgfpathlineto{\pgfqpoint{5.870171in}{0.722445in}}%
\pgfpathlineto{\pgfqpoint{5.873527in}{0.688082in}}%
\pgfpathlineto{\pgfqpoint{5.876883in}{0.637863in}}%
\pgfpathlineto{\pgfqpoint{5.880238in}{0.689773in}}%
\pgfpathlineto{\pgfqpoint{5.883594in}{0.695974in}}%
\pgfpathlineto{\pgfqpoint{5.886950in}{0.680827in}}%
\pgfpathlineto{\pgfqpoint{5.890306in}{0.589486in}}%
\pgfpathlineto{\pgfqpoint{5.893662in}{0.695609in}}%
\pgfpathlineto{\pgfqpoint{5.897017in}{0.718974in}}%
\pgfpathlineto{\pgfqpoint{5.900373in}{0.726443in}}%
\pgfpathlineto{\pgfqpoint{5.903729in}{0.718539in}}%
\pgfpathlineto{\pgfqpoint{5.907085in}{0.697282in}}%
\pgfpathlineto{\pgfqpoint{5.910441in}{0.643589in}}%
\pgfpathlineto{\pgfqpoint{5.913796in}{0.658365in}}%
\pgfpathlineto{\pgfqpoint{5.917152in}{0.682325in}}%
\pgfpathlineto{\pgfqpoint{5.920508in}{0.682121in}}%
\pgfpathlineto{\pgfqpoint{5.927219in}{0.632813in}}%
\pgfpathlineto{\pgfqpoint{5.930575in}{0.689282in}}%
\pgfpathlineto{\pgfqpoint{5.933931in}{0.707963in}}%
\pgfpathlineto{\pgfqpoint{5.937287in}{0.712140in}}%
\pgfpathlineto{\pgfqpoint{5.940643in}{0.699926in}}%
\pgfpathlineto{\pgfqpoint{5.943998in}{0.664214in}}%
\pgfpathlineto{\pgfqpoint{5.947354in}{0.646643in}}%
\pgfpathlineto{\pgfqpoint{5.950710in}{0.688336in}}%
\pgfpathlineto{\pgfqpoint{5.954066in}{0.699774in}}%
\pgfpathlineto{\pgfqpoint{5.957421in}{0.699140in}}%
\pgfpathlineto{\pgfqpoint{5.960777in}{0.685729in}}%
\pgfpathlineto{\pgfqpoint{5.964133in}{0.631610in}}%
\pgfpathlineto{\pgfqpoint{5.967489in}{0.673335in}}%
\pgfpathlineto{\pgfqpoint{5.970845in}{0.704421in}}%
\pgfpathlineto{\pgfqpoint{5.974200in}{0.716340in}}%
\pgfpathlineto{\pgfqpoint{5.977556in}{0.715462in}}%
\pgfpathlineto{\pgfqpoint{5.980912in}{0.696068in}}%
\pgfpathlineto{\pgfqpoint{5.984268in}{0.561952in}}%
\pgfpathlineto{\pgfqpoint{5.987623in}{0.693883in}}%
\pgfpathlineto{\pgfqpoint{5.990979in}{0.715761in}}%
\pgfpathlineto{\pgfqpoint{5.994335in}{0.723291in}}%
\pgfpathlineto{\pgfqpoint{5.997691in}{0.721973in}}%
\pgfpathlineto{\pgfqpoint{6.001047in}{0.710586in}}%
\pgfpathlineto{\pgfqpoint{6.004402in}{0.674802in}}%
\pgfpathlineto{\pgfqpoint{6.007758in}{0.675766in}}%
\pgfpathlineto{\pgfqpoint{6.011114in}{0.717196in}}%
\pgfpathlineto{\pgfqpoint{6.014470in}{0.732927in}}%
\pgfpathlineto{\pgfqpoint{6.017825in}{0.736901in}}%
\pgfpathlineto{\pgfqpoint{6.021181in}{0.727101in}}%
\pgfpathlineto{\pgfqpoint{6.024537in}{0.694721in}}%
\pgfpathlineto{\pgfqpoint{6.027893in}{0.669429in}}%
\pgfpathlineto{\pgfqpoint{6.031249in}{0.716313in}}%
\pgfpathlineto{\pgfqpoint{6.034604in}{0.729808in}}%
\pgfpathlineto{\pgfqpoint{6.037960in}{0.731617in}}%
\pgfpathlineto{\pgfqpoint{6.041316in}{0.722953in}}%
\pgfpathlineto{\pgfqpoint{6.044672in}{0.693683in}}%
\pgfpathlineto{\pgfqpoint{6.048027in}{0.674660in}}%
\pgfpathlineto{\pgfqpoint{6.051383in}{0.724970in}}%
\pgfpathlineto{\pgfqpoint{6.054739in}{0.742475in}}%
\pgfpathlineto{\pgfqpoint{6.058095in}{0.747199in}}%
\pgfpathlineto{\pgfqpoint{6.061451in}{0.738121in}}%
\pgfpathlineto{\pgfqpoint{6.064806in}{0.715253in}}%
\pgfpathlineto{\pgfqpoint{6.068162in}{0.646834in}}%
\pgfpathlineto{\pgfqpoint{6.071518in}{0.688531in}}%
\pgfpathlineto{\pgfqpoint{6.074874in}{0.708462in}}%
\pgfpathlineto{\pgfqpoint{6.078230in}{0.707829in}}%
\pgfpathlineto{\pgfqpoint{6.081585in}{0.686469in}}%
\pgfpathlineto{\pgfqpoint{6.084941in}{0.644775in}}%
\pgfpathlineto{\pgfqpoint{6.088297in}{0.710381in}}%
\pgfpathlineto{\pgfqpoint{6.091653in}{0.730344in}}%
\pgfpathlineto{\pgfqpoint{6.095008in}{0.735616in}}%
\pgfpathlineto{\pgfqpoint{6.098364in}{0.725773in}}%
\pgfpathlineto{\pgfqpoint{6.101720in}{0.704097in}}%
\pgfpathlineto{\pgfqpoint{6.105076in}{0.654209in}}%
\pgfpathlineto{\pgfqpoint{6.108432in}{0.655834in}}%
\pgfpathlineto{\pgfqpoint{6.111787in}{0.680369in}}%
\pgfpathlineto{\pgfqpoint{6.115143in}{0.676075in}}%
\pgfpathlineto{\pgfqpoint{6.118499in}{0.629276in}}%
\pgfpathlineto{\pgfqpoint{6.125210in}{0.702104in}}%
\pgfpathlineto{\pgfqpoint{6.128566in}{0.714715in}}%
\pgfpathlineto{\pgfqpoint{6.131922in}{0.712593in}}%
\pgfpathlineto{\pgfqpoint{6.135278in}{0.694739in}}%
\pgfpathlineto{\pgfqpoint{6.138634in}{0.648798in}}%
\pgfpathlineto{\pgfqpoint{6.141989in}{0.653004in}}%
\pgfpathlineto{\pgfqpoint{6.145345in}{0.683357in}}%
\pgfpathlineto{\pgfqpoint{6.148701in}{0.689267in}}%
\pgfpathlineto{\pgfqpoint{6.152057in}{0.680804in}}%
\pgfpathlineto{\pgfqpoint{6.155412in}{0.642974in}}%
\pgfpathlineto{\pgfqpoint{6.158768in}{0.656322in}}%
\pgfpathlineto{\pgfqpoint{6.162124in}{0.694474in}}%
\pgfpathlineto{\pgfqpoint{6.165480in}{0.708674in}}%
\pgfpathlineto{\pgfqpoint{6.168836in}{0.709315in}}%
\pgfpathlineto{\pgfqpoint{6.172191in}{0.691910in}}%
\pgfpathlineto{\pgfqpoint{6.175547in}{0.616012in}}%
\pgfpathlineto{\pgfqpoint{6.178903in}{0.680762in}}%
\pgfpathlineto{\pgfqpoint{6.182259in}{0.705096in}}%
\pgfpathlineto{\pgfqpoint{6.185614in}{0.713315in}}%
\pgfpathlineto{\pgfqpoint{6.188970in}{0.712219in}}%
\pgfpathlineto{\pgfqpoint{6.192326in}{0.700715in}}%
\pgfpathlineto{\pgfqpoint{6.195682in}{0.663535in}}%
\pgfpathlineto{\pgfqpoint{6.199038in}{0.668875in}}%
\pgfpathlineto{\pgfqpoint{6.202393in}{0.708795in}}%
\pgfpathlineto{\pgfqpoint{6.205749in}{0.724039in}}%
\pgfpathlineto{\pgfqpoint{6.209105in}{0.727410in}}%
\pgfpathlineto{\pgfqpoint{6.212461in}{0.716057in}}%
\pgfpathlineto{\pgfqpoint{6.215816in}{0.672801in}}%
\pgfpathlineto{\pgfqpoint{6.219172in}{0.686071in}}%
\pgfpathlineto{\pgfqpoint{6.222528in}{0.719702in}}%
\pgfpathlineto{\pgfqpoint{6.225884in}{0.731473in}}%
\pgfpathlineto{\pgfqpoint{6.229240in}{0.733433in}}%
\pgfpathlineto{\pgfqpoint{6.232595in}{0.726484in}}%
\pgfpathlineto{\pgfqpoint{6.235951in}{0.704390in}}%
\pgfpathlineto{\pgfqpoint{6.239307in}{0.625564in}}%
\pgfpathlineto{\pgfqpoint{6.242663in}{0.716145in}}%
\pgfpathlineto{\pgfqpoint{6.246019in}{0.738279in}}%
\pgfpathlineto{\pgfqpoint{6.249374in}{0.746590in}}%
\pgfpathlineto{\pgfqpoint{6.252730in}{0.742847in}}%
\pgfpathlineto{\pgfqpoint{6.256086in}{0.724193in}}%
\pgfpathlineto{\pgfqpoint{6.259442in}{0.670751in}}%
\pgfpathlineto{\pgfqpoint{6.262797in}{0.693061in}}%
\pgfpathlineto{\pgfqpoint{6.266153in}{0.719694in}}%
\pgfpathlineto{\pgfqpoint{6.269509in}{0.725377in}}%
\pgfpathlineto{\pgfqpoint{6.272865in}{0.718033in}}%
\pgfpathlineto{\pgfqpoint{6.276221in}{0.687056in}}%
\pgfpathlineto{\pgfqpoint{6.279576in}{0.679810in}}%
\pgfpathlineto{\pgfqpoint{6.282932in}{0.725284in}}%
\pgfpathlineto{\pgfqpoint{6.286288in}{0.741543in}}%
\pgfpathlineto{\pgfqpoint{6.289644in}{0.744395in}}%
\pgfpathlineto{\pgfqpoint{6.292999in}{0.732788in}}%
\pgfpathlineto{\pgfqpoint{6.296355in}{0.709513in}}%
\pgfpathlineto{\pgfqpoint{6.299711in}{0.652578in}}%
\pgfpathlineto{\pgfqpoint{6.303067in}{0.668042in}}%
\pgfpathlineto{\pgfqpoint{6.306423in}{0.687902in}}%
\pgfpathlineto{\pgfqpoint{6.309778in}{0.679967in}}%
\pgfpathlineto{\pgfqpoint{6.313134in}{0.600488in}}%
\pgfpathlineto{\pgfqpoint{6.316490in}{0.685509in}}%
\pgfpathlineto{\pgfqpoint{6.319846in}{0.713499in}}%
\pgfpathlineto{\pgfqpoint{6.323201in}{0.723657in}}%
\pgfpathlineto{\pgfqpoint{6.326557in}{0.718607in}}%
\pgfpathlineto{\pgfqpoint{6.329913in}{0.700141in}}%
\pgfpathlineto{\pgfqpoint{6.333269in}{0.659773in}}%
\pgfpathlineto{\pgfqpoint{6.336625in}{0.638498in}}%
\pgfpathlineto{\pgfqpoint{6.339980in}{0.674387in}}%
\pgfpathlineto{\pgfqpoint{6.343336in}{0.677345in}}%
\pgfpathlineto{\pgfqpoint{6.346692in}{0.657849in}}%
\pgfpathlineto{\pgfqpoint{6.350048in}{0.614375in}}%
\pgfpathlineto{\pgfqpoint{6.353403in}{0.682223in}}%
\pgfpathlineto{\pgfqpoint{6.356759in}{0.702690in}}%
\pgfpathlineto{\pgfqpoint{6.360115in}{0.708189in}}%
\pgfpathlineto{\pgfqpoint{6.363471in}{0.699442in}}%
\pgfpathlineto{\pgfqpoint{6.373538in}{0.664528in}}%
\pgfpathlineto{\pgfqpoint{6.376894in}{0.655621in}}%
\pgfpathlineto{\pgfqpoint{6.383605in}{0.640787in}}%
\pgfpathlineto{\pgfqpoint{6.383605in}{0.640787in}}%
\pgfusepath{stroke}%
\end{pgfscope}%
\begin{pgfscope}%
\pgfpathrectangle{\pgfqpoint{4.621821in}{0.320625in}}{\pgfqpoint{1.845679in}{1.164375in}}%
\pgfusepath{clip}%
\pgfsetrectcap%
\pgfsetroundjoin%
\pgfsetlinewidth{0.602250pt}%
\definecolor{currentstroke}{rgb}{0.000000,0.529412,0.423529}%
\pgfsetstrokecolor{currentstroke}%
\pgfsetstrokeopacity{0.120000}%
\pgfsetdash{}{0pt}%
\pgfpathmoveto{\pgfqpoint{4.705715in}{1.469311in}}%
\pgfpathlineto{\pgfqpoint{4.719139in}{1.444863in}}%
\pgfpathlineto{\pgfqpoint{4.729206in}{1.423893in}}%
\pgfpathlineto{\pgfqpoint{4.745985in}{1.385546in}}%
\pgfpathlineto{\pgfqpoint{4.772831in}{1.324120in}}%
\pgfpathlineto{\pgfqpoint{4.792966in}{1.280817in}}%
\pgfpathlineto{\pgfqpoint{4.816456in}{1.233087in}}%
\pgfpathlineto{\pgfqpoint{4.843302in}{1.181222in}}%
\pgfpathlineto{\pgfqpoint{4.873504in}{1.125290in}}%
\pgfpathlineto{\pgfqpoint{4.910418in}{1.059402in}}%
\pgfpathlineto{\pgfqpoint{4.954043in}{0.984005in}}%
\pgfpathlineto{\pgfqpoint{5.004380in}{0.899187in}}%
\pgfpathlineto{\pgfqpoint{5.068140in}{0.794029in}}%
\pgfpathlineto{\pgfqpoint{5.081563in}{0.773601in}}%
\pgfpathlineto{\pgfqpoint{5.094986in}{0.749142in}}%
\pgfpathlineto{\pgfqpoint{5.101698in}{0.737350in}}%
\pgfpathlineto{\pgfqpoint{5.108409in}{0.727903in}}%
\pgfpathlineto{\pgfqpoint{5.115121in}{0.721381in}}%
\pgfpathlineto{\pgfqpoint{5.121832in}{0.715264in}}%
\pgfpathlineto{\pgfqpoint{5.125188in}{0.710278in}}%
\pgfpathlineto{\pgfqpoint{5.131900in}{0.696204in}}%
\pgfpathlineto{\pgfqpoint{5.138611in}{0.677484in}}%
\pgfpathlineto{\pgfqpoint{5.148678in}{0.647482in}}%
\pgfpathlineto{\pgfqpoint{5.152034in}{0.640357in}}%
\pgfpathlineto{\pgfqpoint{5.155390in}{0.634822in}}%
\pgfpathlineto{\pgfqpoint{5.158746in}{0.631770in}}%
\pgfpathlineto{\pgfqpoint{5.162102in}{0.634449in}}%
\pgfpathlineto{\pgfqpoint{5.175525in}{0.666989in}}%
\pgfpathlineto{\pgfqpoint{5.178880in}{0.669671in}}%
\pgfpathlineto{\pgfqpoint{5.182236in}{0.668319in}}%
\pgfpathlineto{\pgfqpoint{5.185592in}{0.660923in}}%
\pgfpathlineto{\pgfqpoint{5.188948in}{0.657969in}}%
\pgfpathlineto{\pgfqpoint{5.195659in}{0.696887in}}%
\pgfpathlineto{\pgfqpoint{5.199015in}{0.710435in}}%
\pgfpathlineto{\pgfqpoint{5.202371in}{0.719893in}}%
\pgfpathlineto{\pgfqpoint{5.205727in}{0.726374in}}%
\pgfpathlineto{\pgfqpoint{5.209082in}{0.730495in}}%
\pgfpathlineto{\pgfqpoint{5.212438in}{0.732567in}}%
\pgfpathlineto{\pgfqpoint{5.215794in}{0.732689in}}%
\pgfpathlineto{\pgfqpoint{5.219150in}{0.730756in}}%
\pgfpathlineto{\pgfqpoint{5.222506in}{0.726401in}}%
\pgfpathlineto{\pgfqpoint{5.225861in}{0.718766in}}%
\pgfpathlineto{\pgfqpoint{5.229217in}{0.705752in}}%
\pgfpathlineto{\pgfqpoint{5.232573in}{0.680406in}}%
\pgfpathlineto{\pgfqpoint{5.235929in}{0.631094in}}%
\pgfpathlineto{\pgfqpoint{5.239285in}{0.690724in}}%
\pgfpathlineto{\pgfqpoint{5.242640in}{0.714427in}}%
\pgfpathlineto{\pgfqpoint{5.245996in}{0.728197in}}%
\pgfpathlineto{\pgfqpoint{5.249352in}{0.737030in}}%
\pgfpathlineto{\pgfqpoint{5.252708in}{0.742466in}}%
\pgfpathlineto{\pgfqpoint{5.256063in}{0.744946in}}%
\pgfpathlineto{\pgfqpoint{5.259419in}{0.744233in}}%
\pgfpathlineto{\pgfqpoint{5.262775in}{0.739170in}}%
\pgfpathlineto{\pgfqpoint{5.266131in}{0.728230in}}%
\pgfpathlineto{\pgfqpoint{5.272842in}{0.694978in}}%
\pgfpathlineto{\pgfqpoint{5.276198in}{0.685876in}}%
\pgfpathlineto{\pgfqpoint{5.279554in}{0.683782in}}%
\pgfpathlineto{\pgfqpoint{5.282910in}{0.675829in}}%
\pgfpathlineto{\pgfqpoint{5.286265in}{0.664190in}}%
\pgfpathlineto{\pgfqpoint{5.292977in}{0.700521in}}%
\pgfpathlineto{\pgfqpoint{5.296333in}{0.708680in}}%
\pgfpathlineto{\pgfqpoint{5.299689in}{0.703397in}}%
\pgfpathlineto{\pgfqpoint{5.303044in}{0.687574in}}%
\pgfpathlineto{\pgfqpoint{5.309756in}{0.634037in}}%
\pgfpathlineto{\pgfqpoint{5.313112in}{0.645660in}}%
\pgfpathlineto{\pgfqpoint{5.316467in}{0.647536in}}%
\pgfpathlineto{\pgfqpoint{5.319823in}{0.627053in}}%
\pgfpathlineto{\pgfqpoint{5.323179in}{0.624065in}}%
\pgfpathlineto{\pgfqpoint{5.326535in}{0.663316in}}%
\pgfpathlineto{\pgfqpoint{5.329891in}{0.680359in}}%
\pgfpathlineto{\pgfqpoint{5.333246in}{0.683927in}}%
\pgfpathlineto{\pgfqpoint{5.336602in}{0.667586in}}%
\pgfpathlineto{\pgfqpoint{5.339958in}{0.613987in}}%
\pgfpathlineto{\pgfqpoint{5.343314in}{0.683651in}}%
\pgfpathlineto{\pgfqpoint{5.346669in}{0.705825in}}%
\pgfpathlineto{\pgfqpoint{5.350025in}{0.716188in}}%
\pgfpathlineto{\pgfqpoint{5.353381in}{0.719733in}}%
\pgfpathlineto{\pgfqpoint{5.356737in}{0.717401in}}%
\pgfpathlineto{\pgfqpoint{5.360093in}{0.707813in}}%
\pgfpathlineto{\pgfqpoint{5.363448in}{0.683691in}}%
\pgfpathlineto{\pgfqpoint{5.366804in}{0.613368in}}%
\pgfpathlineto{\pgfqpoint{5.370160in}{0.696302in}}%
\pgfpathlineto{\pgfqpoint{5.373516in}{0.718602in}}%
\pgfpathlineto{\pgfqpoint{5.376871in}{0.728576in}}%
\pgfpathlineto{\pgfqpoint{5.380227in}{0.729717in}}%
\pgfpathlineto{\pgfqpoint{5.383583in}{0.719150in}}%
\pgfpathlineto{\pgfqpoint{5.386939in}{0.691557in}}%
\pgfpathlineto{\pgfqpoint{5.390295in}{0.633422in}}%
\pgfpathlineto{\pgfqpoint{5.393650in}{0.693992in}}%
\pgfpathlineto{\pgfqpoint{5.397006in}{0.708675in}}%
\pgfpathlineto{\pgfqpoint{5.400362in}{0.709407in}}%
\pgfpathlineto{\pgfqpoint{5.403718in}{0.696958in}}%
\pgfpathlineto{\pgfqpoint{5.407074in}{0.645853in}}%
\pgfpathlineto{\pgfqpoint{5.413785in}{0.715464in}}%
\pgfpathlineto{\pgfqpoint{5.417141in}{0.727775in}}%
\pgfpathlineto{\pgfqpoint{5.420497in}{0.727261in}}%
\pgfpathlineto{\pgfqpoint{5.423852in}{0.711263in}}%
\pgfpathlineto{\pgfqpoint{5.427208in}{0.668943in}}%
\pgfpathlineto{\pgfqpoint{5.430564in}{0.670315in}}%
\pgfpathlineto{\pgfqpoint{5.433920in}{0.702670in}}%
\pgfpathlineto{\pgfqpoint{5.437276in}{0.711121in}}%
\pgfpathlineto{\pgfqpoint{5.440631in}{0.707371in}}%
\pgfpathlineto{\pgfqpoint{5.443987in}{0.687263in}}%
\pgfpathlineto{\pgfqpoint{5.447343in}{0.606293in}}%
\pgfpathlineto{\pgfqpoint{5.450699in}{0.699871in}}%
\pgfpathlineto{\pgfqpoint{5.454054in}{0.722171in}}%
\pgfpathlineto{\pgfqpoint{5.457410in}{0.729764in}}%
\pgfpathlineto{\pgfqpoint{5.460766in}{0.723368in}}%
\pgfpathlineto{\pgfqpoint{5.464122in}{0.697415in}}%
\pgfpathlineto{\pgfqpoint{5.467478in}{0.637738in}}%
\pgfpathlineto{\pgfqpoint{5.470833in}{0.705158in}}%
\pgfpathlineto{\pgfqpoint{5.474189in}{0.722006in}}%
\pgfpathlineto{\pgfqpoint{5.477545in}{0.726072in}}%
\pgfpathlineto{\pgfqpoint{5.480901in}{0.720371in}}%
\pgfpathlineto{\pgfqpoint{5.484256in}{0.699478in}}%
\pgfpathlineto{\pgfqpoint{5.487612in}{0.608954in}}%
\pgfpathlineto{\pgfqpoint{5.490968in}{0.709892in}}%
\pgfpathlineto{\pgfqpoint{5.494324in}{0.732518in}}%
\pgfpathlineto{\pgfqpoint{5.497680in}{0.740808in}}%
\pgfpathlineto{\pgfqpoint{5.501035in}{0.736399in}}%
\pgfpathlineto{\pgfqpoint{5.504391in}{0.714783in}}%
\pgfpathlineto{\pgfqpoint{5.507747in}{0.620561in}}%
\pgfpathlineto{\pgfqpoint{5.511103in}{0.707609in}}%
\pgfpathlineto{\pgfqpoint{5.514458in}{0.728084in}}%
\pgfpathlineto{\pgfqpoint{5.517814in}{0.733487in}}%
\pgfpathlineto{\pgfqpoint{5.521170in}{0.728597in}}%
\pgfpathlineto{\pgfqpoint{5.524526in}{0.708633in}}%
\pgfpathlineto{\pgfqpoint{5.527882in}{0.600887in}}%
\pgfpathlineto{\pgfqpoint{5.531237in}{0.717425in}}%
\pgfpathlineto{\pgfqpoint{5.534593in}{0.740582in}}%
\pgfpathlineto{\pgfqpoint{5.537949in}{0.749001in}}%
\pgfpathlineto{\pgfqpoint{5.541305in}{0.744445in}}%
\pgfpathlineto{\pgfqpoint{5.544660in}{0.725347in}}%
\pgfpathlineto{\pgfqpoint{5.548016in}{0.674078in}}%
\pgfpathlineto{\pgfqpoint{5.551372in}{0.691112in}}%
\pgfpathlineto{\pgfqpoint{5.554728in}{0.717747in}}%
\pgfpathlineto{\pgfqpoint{5.558084in}{0.722664in}}%
\pgfpathlineto{\pgfqpoint{5.561439in}{0.713482in}}%
\pgfpathlineto{\pgfqpoint{5.564795in}{0.673584in}}%
\pgfpathlineto{\pgfqpoint{5.568151in}{0.691324in}}%
\pgfpathlineto{\pgfqpoint{5.571507in}{0.728157in}}%
\pgfpathlineto{\pgfqpoint{5.574863in}{0.742028in}}%
\pgfpathlineto{\pgfqpoint{5.578218in}{0.742555in}}%
\pgfpathlineto{\pgfqpoint{5.581574in}{0.728405in}}%
\pgfpathlineto{\pgfqpoint{5.584930in}{0.698375in}}%
\pgfpathlineto{\pgfqpoint{5.588286in}{0.623180in}}%
\pgfpathlineto{\pgfqpoint{5.591641in}{0.693710in}}%
\pgfpathlineto{\pgfqpoint{5.594997in}{0.704553in}}%
\pgfpathlineto{\pgfqpoint{5.598353in}{0.697322in}}%
\pgfpathlineto{\pgfqpoint{5.601709in}{0.655805in}}%
\pgfpathlineto{\pgfqpoint{5.605065in}{0.680678in}}%
\pgfpathlineto{\pgfqpoint{5.608420in}{0.716263in}}%
\pgfpathlineto{\pgfqpoint{5.611776in}{0.729674in}}%
\pgfpathlineto{\pgfqpoint{5.615132in}{0.729112in}}%
\pgfpathlineto{\pgfqpoint{5.618488in}{0.713007in}}%
\pgfpathlineto{\pgfqpoint{5.621843in}{0.674128in}}%
\pgfpathlineto{\pgfqpoint{5.625199in}{0.658825in}}%
\pgfpathlineto{\pgfqpoint{5.628555in}{0.696170in}}%
\pgfpathlineto{\pgfqpoint{5.631911in}{0.704132in}}%
\pgfpathlineto{\pgfqpoint{5.635267in}{0.697764in}}%
\pgfpathlineto{\pgfqpoint{5.638622in}{0.667304in}}%
\pgfpathlineto{\pgfqpoint{5.641978in}{0.660478in}}%
\pgfpathlineto{\pgfqpoint{5.645334in}{0.706096in}}%
\pgfpathlineto{\pgfqpoint{5.648690in}{0.722439in}}%
\pgfpathlineto{\pgfqpoint{5.652045in}{0.725231in}}%
\pgfpathlineto{\pgfqpoint{5.655401in}{0.711809in}}%
\pgfpathlineto{\pgfqpoint{5.658757in}{0.670589in}}%
\pgfpathlineto{\pgfqpoint{5.662113in}{0.672667in}}%
\pgfpathlineto{\pgfqpoint{5.665469in}{0.707414in}}%
\pgfpathlineto{\pgfqpoint{5.668824in}{0.717945in}}%
\pgfpathlineto{\pgfqpoint{5.672180in}{0.717475in}}%
\pgfpathlineto{\pgfqpoint{5.675536in}{0.705411in}}%
\pgfpathlineto{\pgfqpoint{5.678892in}{0.662301in}}%
\pgfpathlineto{\pgfqpoint{5.682247in}{0.683069in}}%
\pgfpathlineto{\pgfqpoint{5.685603in}{0.718532in}}%
\pgfpathlineto{\pgfqpoint{5.688959in}{0.732239in}}%
\pgfpathlineto{\pgfqpoint{5.692315in}{0.733693in}}%
\pgfpathlineto{\pgfqpoint{5.695671in}{0.719217in}}%
\pgfpathlineto{\pgfqpoint{5.699026in}{0.669924in}}%
\pgfpathlineto{\pgfqpoint{5.702382in}{0.692363in}}%
\pgfpathlineto{\pgfqpoint{5.705738in}{0.722384in}}%
\pgfpathlineto{\pgfqpoint{5.709094in}{0.732231in}}%
\pgfpathlineto{\pgfqpoint{5.712449in}{0.731998in}}%
\pgfpathlineto{\pgfqpoint{5.715805in}{0.721233in}}%
\pgfpathlineto{\pgfqpoint{5.719161in}{0.685920in}}%
\pgfpathlineto{\pgfqpoint{5.722517in}{0.686716in}}%
\pgfpathlineto{\pgfqpoint{5.725873in}{0.728329in}}%
\pgfpathlineto{\pgfqpoint{5.729228in}{0.743977in}}%
\pgfpathlineto{\pgfqpoint{5.732584in}{0.747521in}}%
\pgfpathlineto{\pgfqpoint{5.735940in}{0.736982in}}%
\pgfpathlineto{\pgfqpoint{5.739296in}{0.707423in}}%
\pgfpathlineto{\pgfqpoint{5.742652in}{0.655621in}}%
\pgfpathlineto{\pgfqpoint{5.746007in}{0.715686in}}%
\pgfpathlineto{\pgfqpoint{5.749363in}{0.729526in}}%
\pgfpathlineto{\pgfqpoint{5.752719in}{0.729878in}}%
\pgfpathlineto{\pgfqpoint{5.756075in}{0.717103in}}%
\pgfpathlineto{\pgfqpoint{5.759430in}{0.664550in}}%
\pgfpathlineto{\pgfqpoint{5.762786in}{0.704289in}}%
\pgfpathlineto{\pgfqpoint{5.766142in}{0.735820in}}%
\pgfpathlineto{\pgfqpoint{5.769498in}{0.747874in}}%
\pgfpathlineto{\pgfqpoint{5.772854in}{0.746960in}}%
\pgfpathlineto{\pgfqpoint{5.776209in}{0.731750in}}%
\pgfpathlineto{\pgfqpoint{5.779565in}{0.698991in}}%
\pgfpathlineto{\pgfqpoint{5.782921in}{0.646392in}}%
\pgfpathlineto{\pgfqpoint{5.786277in}{0.701798in}}%
\pgfpathlineto{\pgfqpoint{5.789632in}{0.710805in}}%
\pgfpathlineto{\pgfqpoint{5.792988in}{0.702211in}}%
\pgfpathlineto{\pgfqpoint{5.796344in}{0.655116in}}%
\pgfpathlineto{\pgfqpoint{5.803056in}{0.723268in}}%
\pgfpathlineto{\pgfqpoint{5.806411in}{0.735760in}}%
\pgfpathlineto{\pgfqpoint{5.809767in}{0.734127in}}%
\pgfpathlineto{\pgfqpoint{5.813123in}{0.717873in}}%
\pgfpathlineto{\pgfqpoint{5.819834in}{0.646993in}}%
\pgfpathlineto{\pgfqpoint{5.823190in}{0.691915in}}%
\pgfpathlineto{\pgfqpoint{5.826546in}{0.699256in}}%
\pgfpathlineto{\pgfqpoint{5.829902in}{0.689155in}}%
\pgfpathlineto{\pgfqpoint{5.833258in}{0.633626in}}%
\pgfpathlineto{\pgfqpoint{5.836613in}{0.682154in}}%
\pgfpathlineto{\pgfqpoint{5.839969in}{0.712947in}}%
\pgfpathlineto{\pgfqpoint{5.843325in}{0.724421in}}%
\pgfpathlineto{\pgfqpoint{5.846681in}{0.721589in}}%
\pgfpathlineto{\pgfqpoint{5.850036in}{0.702030in}}%
\pgfpathlineto{\pgfqpoint{5.853392in}{0.637250in}}%
\pgfpathlineto{\pgfqpoint{5.856748in}{0.680811in}}%
\pgfpathlineto{\pgfqpoint{5.860104in}{0.704713in}}%
\pgfpathlineto{\pgfqpoint{5.863460in}{0.710423in}}%
\pgfpathlineto{\pgfqpoint{5.866815in}{0.704452in}}%
\pgfpathlineto{\pgfqpoint{5.870171in}{0.679783in}}%
\pgfpathlineto{\pgfqpoint{5.873527in}{0.641644in}}%
\pgfpathlineto{\pgfqpoint{5.876883in}{0.703222in}}%
\pgfpathlineto{\pgfqpoint{5.880238in}{0.722529in}}%
\pgfpathlineto{\pgfqpoint{5.883594in}{0.728138in}}%
\pgfpathlineto{\pgfqpoint{5.886950in}{0.718937in}}%
\pgfpathlineto{\pgfqpoint{5.890306in}{0.685250in}}%
\pgfpathlineto{\pgfqpoint{5.893662in}{0.669094in}}%
\pgfpathlineto{\pgfqpoint{5.897017in}{0.712553in}}%
\pgfpathlineto{\pgfqpoint{5.900373in}{0.726125in}}%
\pgfpathlineto{\pgfqpoint{5.903729in}{0.728710in}}%
\pgfpathlineto{\pgfqpoint{5.907085in}{0.721861in}}%
\pgfpathlineto{\pgfqpoint{5.910441in}{0.699039in}}%
\pgfpathlineto{\pgfqpoint{5.913796in}{0.634783in}}%
\pgfpathlineto{\pgfqpoint{5.917152in}{0.713926in}}%
\pgfpathlineto{\pgfqpoint{5.920508in}{0.735241in}}%
\pgfpathlineto{\pgfqpoint{5.923864in}{0.742874in}}%
\pgfpathlineto{\pgfqpoint{5.927219in}{0.737827in}}%
\pgfpathlineto{\pgfqpoint{5.930575in}{0.714939in}}%
\pgfpathlineto{\pgfqpoint{5.933931in}{0.579531in}}%
\pgfpathlineto{\pgfqpoint{5.937287in}{0.712625in}}%
\pgfpathlineto{\pgfqpoint{5.940643in}{0.731848in}}%
\pgfpathlineto{\pgfqpoint{5.943998in}{0.736745in}}%
\pgfpathlineto{\pgfqpoint{5.947354in}{0.731452in}}%
\pgfpathlineto{\pgfqpoint{5.950710in}{0.710772in}}%
\pgfpathlineto{\pgfqpoint{5.954066in}{0.620641in}}%
\pgfpathlineto{\pgfqpoint{5.957421in}{0.721418in}}%
\pgfpathlineto{\pgfqpoint{5.960777in}{0.743974in}}%
\pgfpathlineto{\pgfqpoint{5.964133in}{0.752036in}}%
\pgfpathlineto{\pgfqpoint{5.967489in}{0.747022in}}%
\pgfpathlineto{\pgfqpoint{5.970845in}{0.727869in}}%
\pgfpathlineto{\pgfqpoint{5.974200in}{0.678357in}}%
\pgfpathlineto{\pgfqpoint{5.977556in}{0.689170in}}%
\pgfpathlineto{\pgfqpoint{5.980912in}{0.716808in}}%
\pgfpathlineto{\pgfqpoint{5.984268in}{0.721096in}}%
\pgfpathlineto{\pgfqpoint{5.987623in}{0.709841in}}%
\pgfpathlineto{\pgfqpoint{5.990979in}{0.654851in}}%
\pgfpathlineto{\pgfqpoint{5.994335in}{0.701003in}}%
\pgfpathlineto{\pgfqpoint{5.997691in}{0.731865in}}%
\pgfpathlineto{\pgfqpoint{6.001047in}{0.743405in}}%
\pgfpathlineto{\pgfqpoint{6.004402in}{0.740988in}}%
\pgfpathlineto{\pgfqpoint{6.007758in}{0.724768in}}%
\pgfpathlineto{\pgfqpoint{6.011114in}{0.691205in}}%
\pgfpathlineto{\pgfqpoint{6.014470in}{0.636541in}}%
\pgfpathlineto{\pgfqpoint{6.017825in}{0.691617in}}%
\pgfpathlineto{\pgfqpoint{6.021181in}{0.698669in}}%
\pgfpathlineto{\pgfqpoint{6.024537in}{0.685161in}}%
\pgfpathlineto{\pgfqpoint{6.027893in}{0.515409in}}%
\pgfpathlineto{\pgfqpoint{6.031249in}{0.694922in}}%
\pgfpathlineto{\pgfqpoint{6.034604in}{0.719908in}}%
\pgfpathlineto{\pgfqpoint{6.037960in}{0.728434in}}%
\pgfpathlineto{\pgfqpoint{6.041316in}{0.721770in}}%
\pgfpathlineto{\pgfqpoint{6.044672in}{0.699737in}}%
\pgfpathlineto{\pgfqpoint{6.048027in}{0.628198in}}%
\pgfpathlineto{\pgfqpoint{6.051383in}{0.677683in}}%
\pgfpathlineto{\pgfqpoint{6.054739in}{0.698642in}}%
\pgfpathlineto{\pgfqpoint{6.058095in}{0.700898in}}%
\pgfpathlineto{\pgfqpoint{6.061451in}{0.688043in}}%
\pgfpathlineto{\pgfqpoint{6.064806in}{0.624518in}}%
\pgfpathlineto{\pgfqpoint{6.068162in}{0.682638in}}%
\pgfpathlineto{\pgfqpoint{6.071518in}{0.711734in}}%
\pgfpathlineto{\pgfqpoint{6.074874in}{0.722537in}}%
\pgfpathlineto{\pgfqpoint{6.078230in}{0.719409in}}%
\pgfpathlineto{\pgfqpoint{6.081585in}{0.697641in}}%
\pgfpathlineto{\pgfqpoint{6.084941in}{0.535517in}}%
\pgfpathlineto{\pgfqpoint{6.088297in}{0.694821in}}%
\pgfpathlineto{\pgfqpoint{6.091653in}{0.715001in}}%
\pgfpathlineto{\pgfqpoint{6.095008in}{0.720803in}}%
\pgfpathlineto{\pgfqpoint{6.098364in}{0.716910in}}%
\pgfpathlineto{\pgfqpoint{6.101720in}{0.699830in}}%
\pgfpathlineto{\pgfqpoint{6.105076in}{0.615833in}}%
\pgfpathlineto{\pgfqpoint{6.108432in}{0.697716in}}%
\pgfpathlineto{\pgfqpoint{6.111787in}{0.724375in}}%
\pgfpathlineto{\pgfqpoint{6.115143in}{0.734708in}}%
\pgfpathlineto{\pgfqpoint{6.118499in}{0.732818in}}%
\pgfpathlineto{\pgfqpoint{6.121855in}{0.713115in}}%
\pgfpathlineto{\pgfqpoint{6.125210in}{0.603112in}}%
\pgfpathlineto{\pgfqpoint{6.128566in}{0.708080in}}%
\pgfpathlineto{\pgfqpoint{6.131922in}{0.730041in}}%
\pgfpathlineto{\pgfqpoint{6.135278in}{0.737207in}}%
\pgfpathlineto{\pgfqpoint{6.138634in}{0.735165in}}%
\pgfpathlineto{\pgfqpoint{6.141989in}{0.722138in}}%
\pgfpathlineto{\pgfqpoint{6.145345in}{0.677772in}}%
\pgfpathlineto{\pgfqpoint{6.148701in}{0.699921in}}%
\pgfpathlineto{\pgfqpoint{6.152057in}{0.734917in}}%
\pgfpathlineto{\pgfqpoint{6.155412in}{0.748690in}}%
\pgfpathlineto{\pgfqpoint{6.158768in}{0.750781in}}%
\pgfpathlineto{\pgfqpoint{6.162124in}{0.738629in}}%
\pgfpathlineto{\pgfqpoint{6.165480in}{0.706398in}}%
\pgfpathlineto{\pgfqpoint{6.168836in}{0.668696in}}%
\pgfpathlineto{\pgfqpoint{6.172191in}{0.719576in}}%
\pgfpathlineto{\pgfqpoint{6.175547in}{0.731746in}}%
\pgfpathlineto{\pgfqpoint{6.178903in}{0.730592in}}%
\pgfpathlineto{\pgfqpoint{6.182259in}{0.714910in}}%
\pgfpathlineto{\pgfqpoint{6.185614in}{0.629800in}}%
\pgfpathlineto{\pgfqpoint{6.188970in}{0.714536in}}%
\pgfpathlineto{\pgfqpoint{6.192326in}{0.741118in}}%
\pgfpathlineto{\pgfqpoint{6.195682in}{0.750906in}}%
\pgfpathlineto{\pgfqpoint{6.199038in}{0.747115in}}%
\pgfpathlineto{\pgfqpoint{6.202393in}{0.730122in}}%
\pgfpathlineto{\pgfqpoint{6.205749in}{0.694458in}}%
\pgfpathlineto{\pgfqpoint{6.209105in}{0.652664in}}%
\pgfpathlineto{\pgfqpoint{6.212461in}{0.699514in}}%
\pgfpathlineto{\pgfqpoint{6.215816in}{0.705106in}}%
\pgfpathlineto{\pgfqpoint{6.219172in}{0.689653in}}%
\pgfpathlineto{\pgfqpoint{6.222528in}{0.600503in}}%
\pgfpathlineto{\pgfqpoint{6.225884in}{0.704555in}}%
\pgfpathlineto{\pgfqpoint{6.229240in}{0.727839in}}%
\pgfpathlineto{\pgfqpoint{6.232595in}{0.735330in}}%
\pgfpathlineto{\pgfqpoint{6.235951in}{0.727559in}}%
\pgfpathlineto{\pgfqpoint{6.239307in}{0.706410in}}%
\pgfpathlineto{\pgfqpoint{6.242663in}{0.652850in}}%
\pgfpathlineto{\pgfqpoint{6.246019in}{0.667798in}}%
\pgfpathlineto{\pgfqpoint{6.249374in}{0.692002in}}%
\pgfpathlineto{\pgfqpoint{6.252730in}{0.692294in}}%
\pgfpathlineto{\pgfqpoint{6.256086in}{0.670739in}}%
\pgfpathlineto{\pgfqpoint{6.259442in}{0.634489in}}%
\pgfpathlineto{\pgfqpoint{6.262797in}{0.697008in}}%
\pgfpathlineto{\pgfqpoint{6.266153in}{0.716672in}}%
\pgfpathlineto{\pgfqpoint{6.269509in}{0.721715in}}%
\pgfpathlineto{\pgfqpoint{6.272865in}{0.710704in}}%
\pgfpathlineto{\pgfqpoint{6.279576in}{0.648859in}}%
\pgfpathlineto{\pgfqpoint{6.282932in}{0.696285in}}%
\pgfpathlineto{\pgfqpoint{6.286288in}{0.708924in}}%
\pgfpathlineto{\pgfqpoint{6.289644in}{0.709221in}}%
\pgfpathlineto{\pgfqpoint{6.292999in}{0.697339in}}%
\pgfpathlineto{\pgfqpoint{6.296355in}{0.652767in}}%
\pgfpathlineto{\pgfqpoint{6.299711in}{0.677305in}}%
\pgfpathlineto{\pgfqpoint{6.303067in}{0.711905in}}%
\pgfpathlineto{\pgfqpoint{6.306423in}{0.725188in}}%
\pgfpathlineto{\pgfqpoint{6.309778in}{0.725868in}}%
\pgfpathlineto{\pgfqpoint{6.313134in}{0.709210in}}%
\pgfpathlineto{\pgfqpoint{6.316490in}{0.635756in}}%
\pgfpathlineto{\pgfqpoint{6.319846in}{0.698145in}}%
\pgfpathlineto{\pgfqpoint{6.323201in}{0.723100in}}%
\pgfpathlineto{\pgfqpoint{6.326557in}{0.731902in}}%
\pgfpathlineto{\pgfqpoint{6.329913in}{0.731636in}}%
\pgfpathlineto{\pgfqpoint{6.333269in}{0.721780in}}%
\pgfpathlineto{\pgfqpoint{6.336625in}{0.691512in}}%
\pgfpathlineto{\pgfqpoint{6.339980in}{0.673721in}}%
\pgfpathlineto{\pgfqpoint{6.343336in}{0.723169in}}%
\pgfpathlineto{\pgfqpoint{6.346692in}{0.740734in}}%
\pgfpathlineto{\pgfqpoint{6.350048in}{0.746142in}}%
\pgfpathlineto{\pgfqpoint{6.353403in}{0.738991in}}%
\pgfpathlineto{\pgfqpoint{6.356759in}{0.723852in}}%
\pgfpathlineto{\pgfqpoint{6.360115in}{0.702036in}}%
\pgfpathlineto{\pgfqpoint{6.363471in}{0.670511in}}%
\pgfpathlineto{\pgfqpoint{6.366827in}{0.629531in}}%
\pgfpathlineto{\pgfqpoint{6.370182in}{0.628842in}}%
\pgfpathlineto{\pgfqpoint{6.373538in}{0.664388in}}%
\pgfpathlineto{\pgfqpoint{6.376894in}{0.688733in}}%
\pgfpathlineto{\pgfqpoint{6.380250in}{0.700953in}}%
\pgfpathlineto{\pgfqpoint{6.383605in}{0.699836in}}%
\pgfpathlineto{\pgfqpoint{6.383605in}{0.699836in}}%
\pgfusepath{stroke}%
\end{pgfscope}%
\begin{pgfscope}%
\pgfpathrectangle{\pgfqpoint{4.621821in}{0.320625in}}{\pgfqpoint{1.845679in}{1.164375in}}%
\pgfusepath{clip}%
\pgfsetrectcap%
\pgfsetroundjoin%
\pgfsetlinewidth{0.602250pt}%
\definecolor{currentstroke}{rgb}{0.000000,0.529412,0.423529}%
\pgfsetstrokecolor{currentstroke}%
\pgfsetstrokeopacity{0.120000}%
\pgfsetdash{}{0pt}%
\pgfpathmoveto{\pgfqpoint{4.705715in}{1.468955in}}%
\pgfpathlineto{\pgfqpoint{4.719139in}{1.444562in}}%
\pgfpathlineto{\pgfqpoint{4.729206in}{1.423574in}}%
\pgfpathlineto{\pgfqpoint{4.745985in}{1.384957in}}%
\pgfpathlineto{\pgfqpoint{4.776187in}{1.314994in}}%
\pgfpathlineto{\pgfqpoint{4.799677in}{1.263905in}}%
\pgfpathlineto{\pgfqpoint{4.823168in}{1.215450in}}%
\pgfpathlineto{\pgfqpoint{4.853370in}{1.155847in}}%
\pgfpathlineto{\pgfqpoint{4.890283in}{1.085674in}}%
\pgfpathlineto{\pgfqpoint{4.937264in}{0.998965in}}%
\pgfpathlineto{\pgfqpoint{4.997668in}{0.889908in}}%
\pgfpathlineto{\pgfqpoint{5.061428in}{0.776667in}}%
\pgfpathlineto{\pgfqpoint{5.074851in}{0.755168in}}%
\pgfpathlineto{\pgfqpoint{5.084919in}{0.734444in}}%
\pgfpathlineto{\pgfqpoint{5.091630in}{0.720301in}}%
\pgfpathlineto{\pgfqpoint{5.098342in}{0.709049in}}%
\pgfpathlineto{\pgfqpoint{5.101698in}{0.705351in}}%
\pgfpathlineto{\pgfqpoint{5.105053in}{0.702949in}}%
\pgfpathlineto{\pgfqpoint{5.111765in}{0.699686in}}%
\pgfpathlineto{\pgfqpoint{5.115121in}{0.696622in}}%
\pgfpathlineto{\pgfqpoint{5.118476in}{0.690672in}}%
\pgfpathlineto{\pgfqpoint{5.121832in}{0.682573in}}%
\pgfpathlineto{\pgfqpoint{5.125188in}{0.672270in}}%
\pgfpathlineto{\pgfqpoint{5.128544in}{0.659483in}}%
\pgfpathlineto{\pgfqpoint{5.135255in}{0.630478in}}%
\pgfpathlineto{\pgfqpoint{5.138611in}{0.624482in}}%
\pgfpathlineto{\pgfqpoint{5.141967in}{0.624537in}}%
\pgfpathlineto{\pgfqpoint{5.145323in}{0.622558in}}%
\pgfpathlineto{\pgfqpoint{5.148678in}{0.614294in}}%
\pgfpathlineto{\pgfqpoint{5.152034in}{0.598509in}}%
\pgfpathlineto{\pgfqpoint{5.155390in}{0.603074in}}%
\pgfpathlineto{\pgfqpoint{5.158746in}{0.627651in}}%
\pgfpathlineto{\pgfqpoint{5.162102in}{0.645059in}}%
\pgfpathlineto{\pgfqpoint{5.165457in}{0.656350in}}%
\pgfpathlineto{\pgfqpoint{5.168813in}{0.663034in}}%
\pgfpathlineto{\pgfqpoint{5.172169in}{0.665389in}}%
\pgfpathlineto{\pgfqpoint{5.175525in}{0.662363in}}%
\pgfpathlineto{\pgfqpoint{5.178880in}{0.651730in}}%
\pgfpathlineto{\pgfqpoint{5.182236in}{0.665825in}}%
\pgfpathlineto{\pgfqpoint{5.185592in}{0.691682in}}%
\pgfpathlineto{\pgfqpoint{5.188948in}{0.709412in}}%
\pgfpathlineto{\pgfqpoint{5.192304in}{0.721491in}}%
\pgfpathlineto{\pgfqpoint{5.195659in}{0.729826in}}%
\pgfpathlineto{\pgfqpoint{5.199015in}{0.735405in}}%
\pgfpathlineto{\pgfqpoint{5.202371in}{0.738739in}}%
\pgfpathlineto{\pgfqpoint{5.205727in}{0.740057in}}%
\pgfpathlineto{\pgfqpoint{5.209082in}{0.739383in}}%
\pgfpathlineto{\pgfqpoint{5.212438in}{0.736516in}}%
\pgfpathlineto{\pgfqpoint{5.215794in}{0.730923in}}%
\pgfpathlineto{\pgfqpoint{5.219150in}{0.721359in}}%
\pgfpathlineto{\pgfqpoint{5.222506in}{0.704475in}}%
\pgfpathlineto{\pgfqpoint{5.225861in}{0.665621in}}%
\pgfpathlineto{\pgfqpoint{5.229217in}{0.672965in}}%
\pgfpathlineto{\pgfqpoint{5.232573in}{0.710378in}}%
\pgfpathlineto{\pgfqpoint{5.235929in}{0.728629in}}%
\pgfpathlineto{\pgfqpoint{5.239285in}{0.739902in}}%
\pgfpathlineto{\pgfqpoint{5.242640in}{0.747087in}}%
\pgfpathlineto{\pgfqpoint{5.245996in}{0.751122in}}%
\pgfpathlineto{\pgfqpoint{5.249352in}{0.752116in}}%
\pgfpathlineto{\pgfqpoint{5.252708in}{0.749439in}}%
\pgfpathlineto{\pgfqpoint{5.256063in}{0.741322in}}%
\pgfpathlineto{\pgfqpoint{5.259419in}{0.728568in}}%
\pgfpathlineto{\pgfqpoint{5.266131in}{0.697051in}}%
\pgfpathlineto{\pgfqpoint{5.269487in}{0.687722in}}%
\pgfpathlineto{\pgfqpoint{5.276198in}{0.674263in}}%
\pgfpathlineto{\pgfqpoint{5.279554in}{0.684036in}}%
\pgfpathlineto{\pgfqpoint{5.282910in}{0.699717in}}%
\pgfpathlineto{\pgfqpoint{5.286265in}{0.707512in}}%
\pgfpathlineto{\pgfqpoint{5.289621in}{0.702541in}}%
\pgfpathlineto{\pgfqpoint{5.292977in}{0.688156in}}%
\pgfpathlineto{\pgfqpoint{5.296333in}{0.667035in}}%
\pgfpathlineto{\pgfqpoint{5.299689in}{0.641262in}}%
\pgfpathlineto{\pgfqpoint{5.303044in}{0.631801in}}%
\pgfpathlineto{\pgfqpoint{5.306400in}{0.628050in}}%
\pgfpathlineto{\pgfqpoint{5.309756in}{0.614648in}}%
\pgfpathlineto{\pgfqpoint{5.316467in}{0.662371in}}%
\pgfpathlineto{\pgfqpoint{5.319823in}{0.673013in}}%
\pgfpathlineto{\pgfqpoint{5.323179in}{0.669635in}}%
\pgfpathlineto{\pgfqpoint{5.326535in}{0.636168in}}%
\pgfpathlineto{\pgfqpoint{5.329891in}{0.645157in}}%
\pgfpathlineto{\pgfqpoint{5.333246in}{0.682899in}}%
\pgfpathlineto{\pgfqpoint{5.336602in}{0.698718in}}%
\pgfpathlineto{\pgfqpoint{5.339958in}{0.705675in}}%
\pgfpathlineto{\pgfqpoint{5.343314in}{0.706407in}}%
\pgfpathlineto{\pgfqpoint{5.346669in}{0.700793in}}%
\pgfpathlineto{\pgfqpoint{5.350025in}{0.685497in}}%
\pgfpathlineto{\pgfqpoint{5.353381in}{0.638717in}}%
\pgfpathlineto{\pgfqpoint{5.356737in}{0.662950in}}%
\pgfpathlineto{\pgfqpoint{5.360093in}{0.697515in}}%
\pgfpathlineto{\pgfqpoint{5.363448in}{0.712290in}}%
\pgfpathlineto{\pgfqpoint{5.366804in}{0.717511in}}%
\pgfpathlineto{\pgfqpoint{5.370160in}{0.712639in}}%
\pgfpathlineto{\pgfqpoint{5.373516in}{0.692357in}}%
\pgfpathlineto{\pgfqpoint{5.376871in}{0.620610in}}%
\pgfpathlineto{\pgfqpoint{5.380227in}{0.680414in}}%
\pgfpathlineto{\pgfqpoint{5.383583in}{0.702368in}}%
\pgfpathlineto{\pgfqpoint{5.386939in}{0.708517in}}%
\pgfpathlineto{\pgfqpoint{5.390295in}{0.704430in}}%
\pgfpathlineto{\pgfqpoint{5.393650in}{0.686239in}}%
\pgfpathlineto{\pgfqpoint{5.397006in}{0.567030in}}%
\pgfpathlineto{\pgfqpoint{5.400362in}{0.689365in}}%
\pgfpathlineto{\pgfqpoint{5.403718in}{0.714375in}}%
\pgfpathlineto{\pgfqpoint{5.407074in}{0.723882in}}%
\pgfpathlineto{\pgfqpoint{5.410429in}{0.720855in}}%
\pgfpathlineto{\pgfqpoint{5.413785in}{0.701183in}}%
\pgfpathlineto{\pgfqpoint{5.417141in}{0.628936in}}%
\pgfpathlineto{\pgfqpoint{5.420497in}{0.687446in}}%
\pgfpathlineto{\pgfqpoint{5.423852in}{0.710301in}}%
\pgfpathlineto{\pgfqpoint{5.427208in}{0.716856in}}%
\pgfpathlineto{\pgfqpoint{5.430564in}{0.713187in}}%
\pgfpathlineto{\pgfqpoint{5.433920in}{0.695869in}}%
\pgfpathlineto{\pgfqpoint{5.437276in}{0.602284in}}%
\pgfpathlineto{\pgfqpoint{5.440631in}{0.695904in}}%
\pgfpathlineto{\pgfqpoint{5.443987in}{0.721953in}}%
\pgfpathlineto{\pgfqpoint{5.447343in}{0.731966in}}%
\pgfpathlineto{\pgfqpoint{5.450699in}{0.729563in}}%
\pgfpathlineto{\pgfqpoint{5.454054in}{0.710497in}}%
\pgfpathlineto{\pgfqpoint{5.457410in}{0.640105in}}%
\pgfpathlineto{\pgfqpoint{5.460766in}{0.695489in}}%
\pgfpathlineto{\pgfqpoint{5.464122in}{0.719128in}}%
\pgfpathlineto{\pgfqpoint{5.467478in}{0.726119in}}%
\pgfpathlineto{\pgfqpoint{5.470833in}{0.722976in}}%
\pgfpathlineto{\pgfqpoint{5.474189in}{0.706819in}}%
\pgfpathlineto{\pgfqpoint{5.477545in}{0.633383in}}%
\pgfpathlineto{\pgfqpoint{5.480901in}{0.702250in}}%
\pgfpathlineto{\pgfqpoint{5.484256in}{0.729946in}}%
\pgfpathlineto{\pgfqpoint{5.487612in}{0.740750in}}%
\pgfpathlineto{\pgfqpoint{5.490968in}{0.739351in}}%
\pgfpathlineto{\pgfqpoint{5.494324in}{0.723074in}}%
\pgfpathlineto{\pgfqpoint{5.497680in}{0.680971in}}%
\pgfpathlineto{\pgfqpoint{5.501035in}{0.678611in}}%
\pgfpathlineto{\pgfqpoint{5.504391in}{0.712207in}}%
\pgfpathlineto{\pgfqpoint{5.507747in}{0.720384in}}%
\pgfpathlineto{\pgfqpoint{5.511103in}{0.715744in}}%
\pgfpathlineto{\pgfqpoint{5.514458in}{0.692656in}}%
\pgfpathlineto{\pgfqpoint{5.517814in}{0.647447in}}%
\pgfpathlineto{\pgfqpoint{5.521170in}{0.713996in}}%
\pgfpathlineto{\pgfqpoint{5.524526in}{0.734020in}}%
\pgfpathlineto{\pgfqpoint{5.527882in}{0.740034in}}%
\pgfpathlineto{\pgfqpoint{5.531237in}{0.731770in}}%
\pgfpathlineto{\pgfqpoint{5.534593in}{0.709910in}}%
\pgfpathlineto{\pgfqpoint{5.537949in}{0.649697in}}%
\pgfpathlineto{\pgfqpoint{5.541305in}{0.679187in}}%
\pgfpathlineto{\pgfqpoint{5.544660in}{0.701218in}}%
\pgfpathlineto{\pgfqpoint{5.548016in}{0.702054in}}%
\pgfpathlineto{\pgfqpoint{5.551372in}{0.684111in}}%
\pgfpathlineto{\pgfqpoint{5.554728in}{0.606601in}}%
\pgfpathlineto{\pgfqpoint{5.558084in}{0.699202in}}%
\pgfpathlineto{\pgfqpoint{5.561439in}{0.721693in}}%
\pgfpathlineto{\pgfqpoint{5.564795in}{0.728933in}}%
\pgfpathlineto{\pgfqpoint{5.568151in}{0.721321in}}%
\pgfpathlineto{\pgfqpoint{5.571507in}{0.699158in}}%
\pgfpathlineto{\pgfqpoint{5.574863in}{0.631667in}}%
\pgfpathlineto{\pgfqpoint{5.578218in}{0.674430in}}%
\pgfpathlineto{\pgfqpoint{5.581574in}{0.695509in}}%
\pgfpathlineto{\pgfqpoint{5.584930in}{0.697188in}}%
\pgfpathlineto{\pgfqpoint{5.588286in}{0.682602in}}%
\pgfpathlineto{\pgfqpoint{5.591641in}{0.590584in}}%
\pgfpathlineto{\pgfqpoint{5.594997in}{0.685048in}}%
\pgfpathlineto{\pgfqpoint{5.598353in}{0.711378in}}%
\pgfpathlineto{\pgfqpoint{5.601709in}{0.720897in}}%
\pgfpathlineto{\pgfqpoint{5.605065in}{0.716201in}}%
\pgfpathlineto{\pgfqpoint{5.608420in}{0.694313in}}%
\pgfpathlineto{\pgfqpoint{5.611776in}{0.596974in}}%
\pgfpathlineto{\pgfqpoint{5.615132in}{0.683999in}}%
\pgfpathlineto{\pgfqpoint{5.618488in}{0.704532in}}%
\pgfpathlineto{\pgfqpoint{5.621843in}{0.709326in}}%
\pgfpathlineto{\pgfqpoint{5.625199in}{0.703095in}}%
\pgfpathlineto{\pgfqpoint{5.628555in}{0.678914in}}%
\pgfpathlineto{\pgfqpoint{5.631911in}{0.635222in}}%
\pgfpathlineto{\pgfqpoint{5.635267in}{0.700215in}}%
\pgfpathlineto{\pgfqpoint{5.638622in}{0.720047in}}%
\pgfpathlineto{\pgfqpoint{5.641978in}{0.726269in}}%
\pgfpathlineto{\pgfqpoint{5.645334in}{0.718411in}}%
\pgfpathlineto{\pgfqpoint{5.648690in}{0.689435in}}%
\pgfpathlineto{\pgfqpoint{5.652045in}{0.649117in}}%
\pgfpathlineto{\pgfqpoint{5.655401in}{0.704920in}}%
\pgfpathlineto{\pgfqpoint{5.658757in}{0.720251in}}%
\pgfpathlineto{\pgfqpoint{5.662113in}{0.723605in}}%
\pgfpathlineto{\pgfqpoint{5.665469in}{0.717303in}}%
\pgfpathlineto{\pgfqpoint{5.668824in}{0.695304in}}%
\pgfpathlineto{\pgfqpoint{5.672180in}{0.622588in}}%
\pgfpathlineto{\pgfqpoint{5.675536in}{0.708528in}}%
\pgfpathlineto{\pgfqpoint{5.678892in}{0.730352in}}%
\pgfpathlineto{\pgfqpoint{5.682247in}{0.738242in}}%
\pgfpathlineto{\pgfqpoint{5.685603in}{0.733442in}}%
\pgfpathlineto{\pgfqpoint{5.688959in}{0.712200in}}%
\pgfpathlineto{\pgfqpoint{5.692315in}{0.625078in}}%
\pgfpathlineto{\pgfqpoint{5.695671in}{0.699623in}}%
\pgfpathlineto{\pgfqpoint{5.699026in}{0.720952in}}%
\pgfpathlineto{\pgfqpoint{5.702382in}{0.726093in}}%
\pgfpathlineto{\pgfqpoint{5.705738in}{0.720201in}}%
\pgfpathlineto{\pgfqpoint{5.709094in}{0.696824in}}%
\pgfpathlineto{\pgfqpoint{5.712449in}{0.647073in}}%
\pgfpathlineto{\pgfqpoint{5.715805in}{0.716191in}}%
\pgfpathlineto{\pgfqpoint{5.719161in}{0.736503in}}%
\pgfpathlineto{\pgfqpoint{5.722517in}{0.743037in}}%
\pgfpathlineto{\pgfqpoint{5.725873in}{0.735866in}}%
\pgfpathlineto{\pgfqpoint{5.729228in}{0.714734in}}%
\pgfpathlineto{\pgfqpoint{5.732584in}{0.655768in}}%
\pgfpathlineto{\pgfqpoint{5.735940in}{0.684047in}}%
\pgfpathlineto{\pgfqpoint{5.739296in}{0.707310in}}%
\pgfpathlineto{\pgfqpoint{5.742652in}{0.709600in}}%
\pgfpathlineto{\pgfqpoint{5.746007in}{0.695340in}}%
\pgfpathlineto{\pgfqpoint{5.749363in}{0.604613in}}%
\pgfpathlineto{\pgfqpoint{5.752719in}{0.697792in}}%
\pgfpathlineto{\pgfqpoint{5.756075in}{0.724242in}}%
\pgfpathlineto{\pgfqpoint{5.759430in}{0.733790in}}%
\pgfpathlineto{\pgfqpoint{5.762786in}{0.729094in}}%
\pgfpathlineto{\pgfqpoint{5.766142in}{0.710749in}}%
\pgfpathlineto{\pgfqpoint{5.769498in}{0.668552in}}%
\pgfpathlineto{\pgfqpoint{5.772854in}{0.657293in}}%
\pgfpathlineto{\pgfqpoint{5.776209in}{0.690719in}}%
\pgfpathlineto{\pgfqpoint{5.779565in}{0.695200in}}%
\pgfpathlineto{\pgfqpoint{5.782921in}{0.681706in}}%
\pgfpathlineto{\pgfqpoint{5.786277in}{0.585588in}}%
\pgfpathlineto{\pgfqpoint{5.789632in}{0.685945in}}%
\pgfpathlineto{\pgfqpoint{5.792988in}{0.712259in}}%
\pgfpathlineto{\pgfqpoint{5.796344in}{0.721672in}}%
\pgfpathlineto{\pgfqpoint{5.799700in}{0.716491in}}%
\pgfpathlineto{\pgfqpoint{5.803056in}{0.695367in}}%
\pgfpathlineto{\pgfqpoint{5.806411in}{0.624593in}}%
\pgfpathlineto{\pgfqpoint{5.809767in}{0.675022in}}%
\pgfpathlineto{\pgfqpoint{5.813123in}{0.697143in}}%
\pgfpathlineto{\pgfqpoint{5.816479in}{0.701260in}}%
\pgfpathlineto{\pgfqpoint{5.819834in}{0.692566in}}%
\pgfpathlineto{\pgfqpoint{5.823190in}{0.657765in}}%
\pgfpathlineto{\pgfqpoint{5.826546in}{0.661518in}}%
\pgfpathlineto{\pgfqpoint{5.829902in}{0.702633in}}%
\pgfpathlineto{\pgfqpoint{5.833258in}{0.717883in}}%
\pgfpathlineto{\pgfqpoint{5.836613in}{0.720097in}}%
\pgfpathlineto{\pgfqpoint{5.839969in}{0.706127in}}%
\pgfpathlineto{\pgfqpoint{5.843325in}{0.660430in}}%
\pgfpathlineto{\pgfqpoint{5.846681in}{0.674595in}}%
\pgfpathlineto{\pgfqpoint{5.850036in}{0.706514in}}%
\pgfpathlineto{\pgfqpoint{5.853392in}{0.716795in}}%
\pgfpathlineto{\pgfqpoint{5.856748in}{0.716752in}}%
\pgfpathlineto{\pgfqpoint{5.860104in}{0.706100in}}%
\pgfpathlineto{\pgfqpoint{5.863460in}{0.670909in}}%
\pgfpathlineto{\pgfqpoint{5.866815in}{0.671565in}}%
\pgfpathlineto{\pgfqpoint{5.870171in}{0.713274in}}%
\pgfpathlineto{\pgfqpoint{5.873527in}{0.728960in}}%
\pgfpathlineto{\pgfqpoint{5.876883in}{0.732535in}}%
\pgfpathlineto{\pgfqpoint{5.880238in}{0.721687in}}%
\pgfpathlineto{\pgfqpoint{5.883594in}{0.686678in}}%
\pgfpathlineto{\pgfqpoint{5.886950in}{0.671740in}}%
\pgfpathlineto{\pgfqpoint{5.890306in}{0.713836in}}%
\pgfpathlineto{\pgfqpoint{5.893662in}{0.726592in}}%
\pgfpathlineto{\pgfqpoint{5.897017in}{0.728168in}}%
\pgfpathlineto{\pgfqpoint{5.900373in}{0.719538in}}%
\pgfpathlineto{\pgfqpoint{5.903729in}{0.690962in}}%
\pgfpathlineto{\pgfqpoint{5.907085in}{0.668303in}}%
\pgfpathlineto{\pgfqpoint{5.910441in}{0.720459in}}%
\pgfpathlineto{\pgfqpoint{5.913796in}{0.738419in}}%
\pgfpathlineto{\pgfqpoint{5.917152in}{0.743708in}}%
\pgfpathlineto{\pgfqpoint{5.920508in}{0.735482in}}%
\pgfpathlineto{\pgfqpoint{5.923864in}{0.711704in}}%
\pgfpathlineto{\pgfqpoint{5.927219in}{0.608775in}}%
\pgfpathlineto{\pgfqpoint{5.930575in}{0.698585in}}%
\pgfpathlineto{\pgfqpoint{5.933931in}{0.717133in}}%
\pgfpathlineto{\pgfqpoint{5.937287in}{0.718989in}}%
\pgfpathlineto{\pgfqpoint{5.940643in}{0.706723in}}%
\pgfpathlineto{\pgfqpoint{5.943998in}{0.652212in}}%
\pgfpathlineto{\pgfqpoint{5.947354in}{0.696240in}}%
\pgfpathlineto{\pgfqpoint{5.950710in}{0.727248in}}%
\pgfpathlineto{\pgfqpoint{5.954066in}{0.739028in}}%
\pgfpathlineto{\pgfqpoint{5.957421in}{0.737479in}}%
\pgfpathlineto{\pgfqpoint{5.960777in}{0.721923in}}%
\pgfpathlineto{\pgfqpoint{5.964133in}{0.689372in}}%
\pgfpathlineto{\pgfqpoint{5.967489in}{0.630318in}}%
\pgfpathlineto{\pgfqpoint{5.970845in}{0.689500in}}%
\pgfpathlineto{\pgfqpoint{5.974200in}{0.698164in}}%
\pgfpathlineto{\pgfqpoint{5.977556in}{0.688127in}}%
\pgfpathlineto{\pgfqpoint{5.980912in}{0.627800in}}%
\pgfpathlineto{\pgfqpoint{5.984268in}{0.684241in}}%
\pgfpathlineto{\pgfqpoint{5.987623in}{0.714164in}}%
\pgfpathlineto{\pgfqpoint{5.990979in}{0.725337in}}%
\pgfpathlineto{\pgfqpoint{5.994335in}{0.722164in}}%
\pgfpathlineto{\pgfqpoint{5.997691in}{0.704336in}}%
\pgfpathlineto{\pgfqpoint{6.001047in}{0.661331in}}%
\pgfpathlineto{\pgfqpoint{6.004402in}{0.655626in}}%
\pgfpathlineto{\pgfqpoint{6.007758in}{0.688416in}}%
\pgfpathlineto{\pgfqpoint{6.011114in}{0.694226in}}%
\pgfpathlineto{\pgfqpoint{6.014470in}{0.684544in}}%
\pgfpathlineto{\pgfqpoint{6.017825in}{0.638327in}}%
\pgfpathlineto{\pgfqpoint{6.024537in}{0.703589in}}%
\pgfpathlineto{\pgfqpoint{6.027893in}{0.716334in}}%
\pgfpathlineto{\pgfqpoint{6.031249in}{0.715447in}}%
\pgfpathlineto{\pgfqpoint{6.034604in}{0.697511in}}%
\pgfpathlineto{\pgfqpoint{6.037960in}{0.636932in}}%
\pgfpathlineto{\pgfqpoint{6.041316in}{0.676119in}}%
\pgfpathlineto{\pgfqpoint{6.044672in}{0.701896in}}%
\pgfpathlineto{\pgfqpoint{6.048027in}{0.709398in}}%
\pgfpathlineto{\pgfqpoint{6.051383in}{0.706355in}}%
\pgfpathlineto{\pgfqpoint{6.054739in}{0.689886in}}%
\pgfpathlineto{\pgfqpoint{6.058095in}{0.609261in}}%
\pgfpathlineto{\pgfqpoint{6.061451in}{0.687451in}}%
\pgfpathlineto{\pgfqpoint{6.064806in}{0.714409in}}%
\pgfpathlineto{\pgfqpoint{6.068162in}{0.724773in}}%
\pgfpathlineto{\pgfqpoint{6.071518in}{0.722616in}}%
\pgfpathlineto{\pgfqpoint{6.074874in}{0.702539in}}%
\pgfpathlineto{\pgfqpoint{6.078230in}{0.587455in}}%
\pgfpathlineto{\pgfqpoint{6.081585in}{0.697729in}}%
\pgfpathlineto{\pgfqpoint{6.084941in}{0.719504in}}%
\pgfpathlineto{\pgfqpoint{6.088297in}{0.726580in}}%
\pgfpathlineto{\pgfqpoint{6.091653in}{0.724473in}}%
\pgfpathlineto{\pgfqpoint{6.095008in}{0.711371in}}%
\pgfpathlineto{\pgfqpoint{6.098364in}{0.666692in}}%
\pgfpathlineto{\pgfqpoint{6.101720in}{0.689484in}}%
\pgfpathlineto{\pgfqpoint{6.105076in}{0.724348in}}%
\pgfpathlineto{\pgfqpoint{6.108432in}{0.738126in}}%
\pgfpathlineto{\pgfqpoint{6.111787in}{0.740307in}}%
\pgfpathlineto{\pgfqpoint{6.115143in}{0.728169in}}%
\pgfpathlineto{\pgfqpoint{6.118499in}{0.693902in}}%
\pgfpathlineto{\pgfqpoint{6.121855in}{0.670204in}}%
\pgfpathlineto{\pgfqpoint{6.125210in}{0.714776in}}%
\pgfpathlineto{\pgfqpoint{6.128566in}{0.726839in}}%
\pgfpathlineto{\pgfqpoint{6.131922in}{0.726727in}}%
\pgfpathlineto{\pgfqpoint{6.135278in}{0.714244in}}%
\pgfpathlineto{\pgfqpoint{6.138634in}{0.666401in}}%
\pgfpathlineto{\pgfqpoint{6.145345in}{0.730449in}}%
\pgfpathlineto{\pgfqpoint{6.148701in}{0.743297in}}%
\pgfpathlineto{\pgfqpoint{6.152057in}{0.743639in}}%
\pgfpathlineto{\pgfqpoint{6.155412in}{0.729775in}}%
\pgfpathlineto{\pgfqpoint{6.158768in}{0.700068in}}%
\pgfpathlineto{\pgfqpoint{6.162124in}{0.614791in}}%
\pgfpathlineto{\pgfqpoint{6.165480in}{0.695117in}}%
\pgfpathlineto{\pgfqpoint{6.168836in}{0.706628in}}%
\pgfpathlineto{\pgfqpoint{6.172191in}{0.700430in}}%
\pgfpathlineto{\pgfqpoint{6.175547in}{0.664198in}}%
\pgfpathlineto{\pgfqpoint{6.178903in}{0.677144in}}%
\pgfpathlineto{\pgfqpoint{6.182259in}{0.716077in}}%
\pgfpathlineto{\pgfqpoint{6.185614in}{0.730655in}}%
\pgfpathlineto{\pgfqpoint{6.188970in}{0.731477in}}%
\pgfpathlineto{\pgfqpoint{6.192326in}{0.717500in}}%
\pgfpathlineto{\pgfqpoint{6.195682in}{0.688635in}}%
\pgfpathlineto{\pgfqpoint{6.199038in}{0.568860in}}%
\pgfpathlineto{\pgfqpoint{6.202393in}{0.678224in}}%
\pgfpathlineto{\pgfqpoint{6.205749in}{0.689893in}}%
\pgfpathlineto{\pgfqpoint{6.209105in}{0.682201in}}%
\pgfpathlineto{\pgfqpoint{6.212461in}{0.635139in}}%
\pgfpathlineto{\pgfqpoint{6.219172in}{0.704490in}}%
\pgfpathlineto{\pgfqpoint{6.222528in}{0.717106in}}%
\pgfpathlineto{\pgfqpoint{6.225884in}{0.715609in}}%
\pgfpathlineto{\pgfqpoint{6.229240in}{0.697918in}}%
\pgfpathlineto{\pgfqpoint{6.232595in}{0.648095in}}%
\pgfpathlineto{\pgfqpoint{6.235951in}{0.664317in}}%
\pgfpathlineto{\pgfqpoint{6.239307in}{0.692978in}}%
\pgfpathlineto{\pgfqpoint{6.242663in}{0.699976in}}%
\pgfpathlineto{\pgfqpoint{6.246019in}{0.694711in}}%
\pgfpathlineto{\pgfqpoint{6.249374in}{0.670888in}}%
\pgfpathlineto{\pgfqpoint{6.252730in}{0.629087in}}%
\pgfpathlineto{\pgfqpoint{6.256086in}{0.693224in}}%
\pgfpathlineto{\pgfqpoint{6.259442in}{0.712921in}}%
\pgfpathlineto{\pgfqpoint{6.262797in}{0.718769in}}%
\pgfpathlineto{\pgfqpoint{6.266153in}{0.709812in}}%
\pgfpathlineto{\pgfqpoint{6.269509in}{0.676458in}}%
\pgfpathlineto{\pgfqpoint{6.272865in}{0.659584in}}%
\pgfpathlineto{\pgfqpoint{6.276221in}{0.703627in}}%
\pgfpathlineto{\pgfqpoint{6.279576in}{0.717496in}}%
\pgfpathlineto{\pgfqpoint{6.282932in}{0.720467in}}%
\pgfpathlineto{\pgfqpoint{6.286288in}{0.714305in}}%
\pgfpathlineto{\pgfqpoint{6.289644in}{0.693485in}}%
\pgfpathlineto{\pgfqpoint{6.292999in}{0.590525in}}%
\pgfpathlineto{\pgfqpoint{6.296355in}{0.702241in}}%
\pgfpathlineto{\pgfqpoint{6.299711in}{0.725304in}}%
\pgfpathlineto{\pgfqpoint{6.303067in}{0.734074in}}%
\pgfpathlineto{\pgfqpoint{6.306423in}{0.730775in}}%
\pgfpathlineto{\pgfqpoint{6.309778in}{0.710011in}}%
\pgfpathlineto{\pgfqpoint{6.313134in}{0.591858in}}%
\pgfpathlineto{\pgfqpoint{6.316490in}{0.704576in}}%
\pgfpathlineto{\pgfqpoint{6.319846in}{0.725845in}}%
\pgfpathlineto{\pgfqpoint{6.323201in}{0.732322in}}%
\pgfpathlineto{\pgfqpoint{6.326557in}{0.729274in}}%
\pgfpathlineto{\pgfqpoint{6.329913in}{0.714052in}}%
\pgfpathlineto{\pgfqpoint{6.333269in}{0.653596in}}%
\pgfpathlineto{\pgfqpoint{6.336625in}{0.703973in}}%
\pgfpathlineto{\pgfqpoint{6.339980in}{0.733696in}}%
\pgfpathlineto{\pgfqpoint{6.343336in}{0.745483in}}%
\pgfpathlineto{\pgfqpoint{6.346692in}{0.745463in}}%
\pgfpathlineto{\pgfqpoint{6.350048in}{0.731899in}}%
\pgfpathlineto{\pgfqpoint{6.353403in}{0.704123in}}%
\pgfpathlineto{\pgfqpoint{6.356759in}{0.531721in}}%
\pgfpathlineto{\pgfqpoint{6.360115in}{0.689573in}}%
\pgfpathlineto{\pgfqpoint{6.363471in}{0.702282in}}%
\pgfpathlineto{\pgfqpoint{6.366827in}{0.694976in}}%
\pgfpathlineto{\pgfqpoint{6.370182in}{0.648392in}}%
\pgfpathlineto{\pgfqpoint{6.376894in}{0.717458in}}%
\pgfpathlineto{\pgfqpoint{6.380250in}{0.730149in}}%
\pgfpathlineto{\pgfqpoint{6.383605in}{0.728699in}}%
\pgfpathlineto{\pgfqpoint{6.383605in}{0.728699in}}%
\pgfusepath{stroke}%
\end{pgfscope}%
\begin{pgfscope}%
\pgfpathrectangle{\pgfqpoint{4.621821in}{0.320625in}}{\pgfqpoint{1.845679in}{1.164375in}}%
\pgfusepath{clip}%
\pgfsetrectcap%
\pgfsetroundjoin%
\pgfsetlinewidth{0.602250pt}%
\definecolor{currentstroke}{rgb}{0.000000,0.529412,0.423529}%
\pgfsetstrokecolor{currentstroke}%
\pgfsetstrokeopacity{0.120000}%
\pgfsetdash{}{0pt}%
\pgfpathmoveto{\pgfqpoint{4.705715in}{1.469553in}}%
\pgfpathlineto{\pgfqpoint{4.719139in}{1.443911in}}%
\pgfpathlineto{\pgfqpoint{4.729206in}{1.422090in}}%
\pgfpathlineto{\pgfqpoint{4.752696in}{1.366863in}}%
\pgfpathlineto{\pgfqpoint{4.772831in}{1.320967in}}%
\pgfpathlineto{\pgfqpoint{4.789610in}{1.285221in}}%
\pgfpathlineto{\pgfqpoint{4.809745in}{1.244987in}}%
\pgfpathlineto{\pgfqpoint{4.833235in}{1.200748in}}%
\pgfpathlineto{\pgfqpoint{4.860081in}{1.152570in}}%
\pgfpathlineto{\pgfqpoint{4.893639in}{1.094621in}}%
\pgfpathlineto{\pgfqpoint{4.937264in}{1.021620in}}%
\pgfpathlineto{\pgfqpoint{4.994313in}{0.928554in}}%
\pgfpathlineto{\pgfqpoint{5.071496in}{0.805003in}}%
\pgfpathlineto{\pgfqpoint{5.084919in}{0.784559in}}%
\pgfpathlineto{\pgfqpoint{5.108409in}{0.745762in}}%
\pgfpathlineto{\pgfqpoint{5.115121in}{0.737492in}}%
\pgfpathlineto{\pgfqpoint{5.121832in}{0.729990in}}%
\pgfpathlineto{\pgfqpoint{5.125188in}{0.725289in}}%
\pgfpathlineto{\pgfqpoint{5.131900in}{0.713053in}}%
\pgfpathlineto{\pgfqpoint{5.138611in}{0.698106in}}%
\pgfpathlineto{\pgfqpoint{5.148678in}{0.673826in}}%
\pgfpathlineto{\pgfqpoint{5.152034in}{0.667158in}}%
\pgfpathlineto{\pgfqpoint{5.155390in}{0.662401in}}%
\pgfpathlineto{\pgfqpoint{5.158746in}{0.660161in}}%
\pgfpathlineto{\pgfqpoint{5.162102in}{0.660566in}}%
\pgfpathlineto{\pgfqpoint{5.165457in}{0.663010in}}%
\pgfpathlineto{\pgfqpoint{5.172169in}{0.669385in}}%
\pgfpathlineto{\pgfqpoint{5.175525in}{0.671051in}}%
\pgfpathlineto{\pgfqpoint{5.178880in}{0.670246in}}%
\pgfpathlineto{\pgfqpoint{5.182236in}{0.665474in}}%
\pgfpathlineto{\pgfqpoint{5.188948in}{0.648329in}}%
\pgfpathlineto{\pgfqpoint{5.192304in}{0.647820in}}%
\pgfpathlineto{\pgfqpoint{5.202371in}{0.675285in}}%
\pgfpathlineto{\pgfqpoint{5.205727in}{0.680332in}}%
\pgfpathlineto{\pgfqpoint{5.209082in}{0.682970in}}%
\pgfpathlineto{\pgfqpoint{5.212438in}{0.683333in}}%
\pgfpathlineto{\pgfqpoint{5.215794in}{0.681322in}}%
\pgfpathlineto{\pgfqpoint{5.219150in}{0.676496in}}%
\pgfpathlineto{\pgfqpoint{5.222506in}{0.667744in}}%
\pgfpathlineto{\pgfqpoint{5.225861in}{0.652082in}}%
\pgfpathlineto{\pgfqpoint{5.229217in}{0.617508in}}%
\pgfpathlineto{\pgfqpoint{5.232573in}{0.614031in}}%
\pgfpathlineto{\pgfqpoint{5.235929in}{0.654275in}}%
\pgfpathlineto{\pgfqpoint{5.239285in}{0.673325in}}%
\pgfpathlineto{\pgfqpoint{5.242640in}{0.684892in}}%
\pgfpathlineto{\pgfqpoint{5.245996in}{0.692113in}}%
\pgfpathlineto{\pgfqpoint{5.249352in}{0.695954in}}%
\pgfpathlineto{\pgfqpoint{5.252708in}{0.696451in}}%
\pgfpathlineto{\pgfqpoint{5.256063in}{0.692718in}}%
\pgfpathlineto{\pgfqpoint{5.259419in}{0.682318in}}%
\pgfpathlineto{\pgfqpoint{5.262775in}{0.666902in}}%
\pgfpathlineto{\pgfqpoint{5.266131in}{0.668114in}}%
\pgfpathlineto{\pgfqpoint{5.269487in}{0.681538in}}%
\pgfpathlineto{\pgfqpoint{5.272842in}{0.689351in}}%
\pgfpathlineto{\pgfqpoint{5.276198in}{0.690322in}}%
\pgfpathlineto{\pgfqpoint{5.279554in}{0.683819in}}%
\pgfpathlineto{\pgfqpoint{5.282910in}{0.665271in}}%
\pgfpathlineto{\pgfqpoint{5.286265in}{0.621202in}}%
\pgfpathlineto{\pgfqpoint{5.289621in}{0.667336in}}%
\pgfpathlineto{\pgfqpoint{5.292977in}{0.691216in}}%
\pgfpathlineto{\pgfqpoint{5.296333in}{0.701435in}}%
\pgfpathlineto{\pgfqpoint{5.299689in}{0.701300in}}%
\pgfpathlineto{\pgfqpoint{5.303044in}{0.685215in}}%
\pgfpathlineto{\pgfqpoint{5.306400in}{0.632505in}}%
\pgfpathlineto{\pgfqpoint{5.309756in}{0.682247in}}%
\pgfpathlineto{\pgfqpoint{5.313112in}{0.705429in}}%
\pgfpathlineto{\pgfqpoint{5.316467in}{0.714702in}}%
\pgfpathlineto{\pgfqpoint{5.319823in}{0.715869in}}%
\pgfpathlineto{\pgfqpoint{5.323179in}{0.709126in}}%
\pgfpathlineto{\pgfqpoint{5.326535in}{0.689223in}}%
\pgfpathlineto{\pgfqpoint{5.329891in}{0.569761in}}%
\pgfpathlineto{\pgfqpoint{5.333246in}{0.691990in}}%
\pgfpathlineto{\pgfqpoint{5.336602in}{0.716844in}}%
\pgfpathlineto{\pgfqpoint{5.339958in}{0.727120in}}%
\pgfpathlineto{\pgfqpoint{5.343314in}{0.726822in}}%
\pgfpathlineto{\pgfqpoint{5.346669in}{0.712939in}}%
\pgfpathlineto{\pgfqpoint{5.350025in}{0.679605in}}%
\pgfpathlineto{\pgfqpoint{5.353381in}{0.652749in}}%
\pgfpathlineto{\pgfqpoint{5.356737in}{0.693598in}}%
\pgfpathlineto{\pgfqpoint{5.360093in}{0.703734in}}%
\pgfpathlineto{\pgfqpoint{5.363448in}{0.700143in}}%
\pgfpathlineto{\pgfqpoint{5.366804in}{0.678323in}}%
\pgfpathlineto{\pgfqpoint{5.370160in}{0.627218in}}%
\pgfpathlineto{\pgfqpoint{5.373516in}{0.697758in}}%
\pgfpathlineto{\pgfqpoint{5.376871in}{0.718259in}}%
\pgfpathlineto{\pgfqpoint{5.380227in}{0.724351in}}%
\pgfpathlineto{\pgfqpoint{5.383583in}{0.715647in}}%
\pgfpathlineto{\pgfqpoint{5.386939in}{0.690467in}}%
\pgfpathlineto{\pgfqpoint{5.390295in}{0.611700in}}%
\pgfpathlineto{\pgfqpoint{5.393650in}{0.683524in}}%
\pgfpathlineto{\pgfqpoint{5.397006in}{0.699871in}}%
\pgfpathlineto{\pgfqpoint{5.400362in}{0.700622in}}%
\pgfpathlineto{\pgfqpoint{5.403718in}{0.686806in}}%
\pgfpathlineto{\pgfqpoint{5.407074in}{0.619801in}}%
\pgfpathlineto{\pgfqpoint{5.410429in}{0.682253in}}%
\pgfpathlineto{\pgfqpoint{5.413785in}{0.710647in}}%
\pgfpathlineto{\pgfqpoint{5.417141in}{0.721148in}}%
\pgfpathlineto{\pgfqpoint{5.420497in}{0.717810in}}%
\pgfpathlineto{\pgfqpoint{5.423852in}{0.696967in}}%
\pgfpathlineto{\pgfqpoint{5.427208in}{0.614894in}}%
\pgfpathlineto{\pgfqpoint{5.430564in}{0.686204in}}%
\pgfpathlineto{\pgfqpoint{5.433920in}{0.707370in}}%
\pgfpathlineto{\pgfqpoint{5.437276in}{0.712703in}}%
\pgfpathlineto{\pgfqpoint{5.440631in}{0.707252in}}%
\pgfpathlineto{\pgfqpoint{5.443987in}{0.685268in}}%
\pgfpathlineto{\pgfqpoint{5.447343in}{0.620894in}}%
\pgfpathlineto{\pgfqpoint{5.450699in}{0.700600in}}%
\pgfpathlineto{\pgfqpoint{5.454054in}{0.721831in}}%
\pgfpathlineto{\pgfqpoint{5.457410in}{0.728856in}}%
\pgfpathlineto{\pgfqpoint{5.460766in}{0.721976in}}%
\pgfpathlineto{\pgfqpoint{5.464122in}{0.695975in}}%
\pgfpathlineto{\pgfqpoint{5.467478in}{0.628196in}}%
\pgfpathlineto{\pgfqpoint{5.470833in}{0.701143in}}%
\pgfpathlineto{\pgfqpoint{5.474189in}{0.717653in}}%
\pgfpathlineto{\pgfqpoint{5.477545in}{0.720786in}}%
\pgfpathlineto{\pgfqpoint{5.480901in}{0.713131in}}%
\pgfpathlineto{\pgfqpoint{5.487612in}{0.659269in}}%
\pgfpathlineto{\pgfqpoint{5.490968in}{0.713641in}}%
\pgfpathlineto{\pgfqpoint{5.494324in}{0.731817in}}%
\pgfpathlineto{\pgfqpoint{5.497680in}{0.736802in}}%
\pgfpathlineto{\pgfqpoint{5.501035in}{0.727493in}}%
\pgfpathlineto{\pgfqpoint{5.504391in}{0.700776in}}%
\pgfpathlineto{\pgfqpoint{5.507747in}{0.614985in}}%
\pgfpathlineto{\pgfqpoint{5.511103in}{0.699347in}}%
\pgfpathlineto{\pgfqpoint{5.514458in}{0.714716in}}%
\pgfpathlineto{\pgfqpoint{5.517814in}{0.715173in}}%
\pgfpathlineto{\pgfqpoint{5.521170in}{0.701361in}}%
\pgfpathlineto{\pgfqpoint{5.524526in}{0.636149in}}%
\pgfpathlineto{\pgfqpoint{5.527882in}{0.695689in}}%
\pgfpathlineto{\pgfqpoint{5.531237in}{0.724400in}}%
\pgfpathlineto{\pgfqpoint{5.534593in}{0.735038in}}%
\pgfpathlineto{\pgfqpoint{5.537949in}{0.731921in}}%
\pgfpathlineto{\pgfqpoint{5.541305in}{0.714249in}}%
\pgfpathlineto{\pgfqpoint{5.544660in}{0.672009in}}%
\pgfpathlineto{\pgfqpoint{5.548016in}{0.663811in}}%
\pgfpathlineto{\pgfqpoint{5.551372in}{0.696963in}}%
\pgfpathlineto{\pgfqpoint{5.554728in}{0.702338in}}%
\pgfpathlineto{\pgfqpoint{5.558084in}{0.691302in}}%
\pgfpathlineto{\pgfqpoint{5.561439in}{0.633893in}}%
\pgfpathlineto{\pgfqpoint{5.564795in}{0.684398in}}%
\pgfpathlineto{\pgfqpoint{5.568151in}{0.714651in}}%
\pgfpathlineto{\pgfqpoint{5.571507in}{0.725774in}}%
\pgfpathlineto{\pgfqpoint{5.574863in}{0.722471in}}%
\pgfpathlineto{\pgfqpoint{5.578218in}{0.703977in}}%
\pgfpathlineto{\pgfqpoint{5.581574in}{0.656368in}}%
\pgfpathlineto{\pgfqpoint{5.584930in}{0.663669in}}%
\pgfpathlineto{\pgfqpoint{5.588286in}{0.692488in}}%
\pgfpathlineto{\pgfqpoint{5.591641in}{0.697353in}}%
\pgfpathlineto{\pgfqpoint{5.594997in}{0.686970in}}%
\pgfpathlineto{\pgfqpoint{5.598353in}{0.638037in}}%
\pgfpathlineto{\pgfqpoint{5.605065in}{0.706639in}}%
\pgfpathlineto{\pgfqpoint{5.608420in}{0.718764in}}%
\pgfpathlineto{\pgfqpoint{5.611776in}{0.716862in}}%
\pgfpathlineto{\pgfqpoint{5.615132in}{0.697981in}}%
\pgfpathlineto{\pgfqpoint{5.618488in}{0.635399in}}%
\pgfpathlineto{\pgfqpoint{5.621843in}{0.676162in}}%
\pgfpathlineto{\pgfqpoint{5.625199in}{0.700711in}}%
\pgfpathlineto{\pgfqpoint{5.628555in}{0.706828in}}%
\pgfpathlineto{\pgfqpoint{5.631911in}{0.701400in}}%
\pgfpathlineto{\pgfqpoint{5.635267in}{0.678275in}}%
\pgfpathlineto{\pgfqpoint{5.638622in}{0.628589in}}%
\pgfpathlineto{\pgfqpoint{5.641978in}{0.697750in}}%
\pgfpathlineto{\pgfqpoint{5.645334in}{0.717933in}}%
\pgfpathlineto{\pgfqpoint{5.648690in}{0.724055in}}%
\pgfpathlineto{\pgfqpoint{5.652045in}{0.715535in}}%
\pgfpathlineto{\pgfqpoint{5.658757in}{0.650457in}}%
\pgfpathlineto{\pgfqpoint{5.662113in}{0.702463in}}%
\pgfpathlineto{\pgfqpoint{5.665469in}{0.716771in}}%
\pgfpathlineto{\pgfqpoint{5.668824in}{0.719008in}}%
\pgfpathlineto{\pgfqpoint{5.672180in}{0.710762in}}%
\pgfpathlineto{\pgfqpoint{5.675536in}{0.682550in}}%
\pgfpathlineto{\pgfqpoint{5.678892in}{0.659113in}}%
\pgfpathlineto{\pgfqpoint{5.682247in}{0.711698in}}%
\pgfpathlineto{\pgfqpoint{5.685603in}{0.729587in}}%
\pgfpathlineto{\pgfqpoint{5.688959in}{0.734501in}}%
\pgfpathlineto{\pgfqpoint{5.692315in}{0.725151in}}%
\pgfpathlineto{\pgfqpoint{5.695671in}{0.696157in}}%
\pgfpathlineto{\pgfqpoint{5.699026in}{0.644732in}}%
\pgfpathlineto{\pgfqpoint{5.702382in}{0.705485in}}%
\pgfpathlineto{\pgfqpoint{5.705738in}{0.719839in}}%
\pgfpathlineto{\pgfqpoint{5.709094in}{0.720960in}}%
\pgfpathlineto{\pgfqpoint{5.712449in}{0.709836in}}%
\pgfpathlineto{\pgfqpoint{5.715805in}{0.667998in}}%
\pgfpathlineto{\pgfqpoint{5.719161in}{0.687186in}}%
\pgfpathlineto{\pgfqpoint{5.722517in}{0.723146in}}%
\pgfpathlineto{\pgfqpoint{5.725873in}{0.736766in}}%
\pgfpathlineto{\pgfqpoint{5.729228in}{0.737516in}}%
\pgfpathlineto{\pgfqpoint{5.732584in}{0.723297in}}%
\pgfpathlineto{\pgfqpoint{5.735940in}{0.689939in}}%
\pgfpathlineto{\pgfqpoint{5.739296in}{0.648684in}}%
\pgfpathlineto{\pgfqpoint{5.742652in}{0.699042in}}%
\pgfpathlineto{\pgfqpoint{5.746007in}{0.708927in}}%
\pgfpathlineto{\pgfqpoint{5.749363in}{0.703355in}}%
\pgfpathlineto{\pgfqpoint{5.752719in}{0.673424in}}%
\pgfpathlineto{\pgfqpoint{5.756075in}{0.666255in}}%
\pgfpathlineto{\pgfqpoint{5.759430in}{0.712169in}}%
\pgfpathlineto{\pgfqpoint{5.762786in}{0.728459in}}%
\pgfpathlineto{\pgfqpoint{5.766142in}{0.730837in}}%
\pgfpathlineto{\pgfqpoint{5.769498in}{0.717829in}}%
\pgfpathlineto{\pgfqpoint{5.772854in}{0.688500in}}%
\pgfpathlineto{\pgfqpoint{5.776209in}{0.604559in}}%
\pgfpathlineto{\pgfqpoint{5.779565in}{0.684302in}}%
\pgfpathlineto{\pgfqpoint{5.782921in}{0.696171in}}%
\pgfpathlineto{\pgfqpoint{5.786277in}{0.690674in}}%
\pgfpathlineto{\pgfqpoint{5.789632in}{0.657735in}}%
\pgfpathlineto{\pgfqpoint{5.792988in}{0.661807in}}%
\pgfpathlineto{\pgfqpoint{5.796344in}{0.703557in}}%
\pgfpathlineto{\pgfqpoint{5.799700in}{0.718769in}}%
\pgfpathlineto{\pgfqpoint{5.803056in}{0.719915in}}%
\pgfpathlineto{\pgfqpoint{5.806411in}{0.704777in}}%
\pgfpathlineto{\pgfqpoint{5.809767in}{0.665255in}}%
\pgfpathlineto{\pgfqpoint{5.813123in}{0.655143in}}%
\pgfpathlineto{\pgfqpoint{5.816479in}{0.691798in}}%
\pgfpathlineto{\pgfqpoint{5.819834in}{0.700562in}}%
\pgfpathlineto{\pgfqpoint{5.823190in}{0.696117in}}%
\pgfpathlineto{\pgfqpoint{5.826546in}{0.672918in}}%
\pgfpathlineto{\pgfqpoint{5.829902in}{0.629561in}}%
\pgfpathlineto{\pgfqpoint{5.833258in}{0.694889in}}%
\pgfpathlineto{\pgfqpoint{5.836613in}{0.714603in}}%
\pgfpathlineto{\pgfqpoint{5.839969in}{0.720081in}}%
\pgfpathlineto{\pgfqpoint{5.843325in}{0.710173in}}%
\pgfpathlineto{\pgfqpoint{5.846681in}{0.676927in}}%
\pgfpathlineto{\pgfqpoint{5.850036in}{0.654819in}}%
\pgfpathlineto{\pgfqpoint{5.853392in}{0.699922in}}%
\pgfpathlineto{\pgfqpoint{5.856748in}{0.712933in}}%
\pgfpathlineto{\pgfqpoint{5.860104in}{0.714257in}}%
\pgfpathlineto{\pgfqpoint{5.863460in}{0.704748in}}%
\pgfpathlineto{\pgfqpoint{5.866815in}{0.672291in}}%
\pgfpathlineto{\pgfqpoint{5.870171in}{0.665420in}}%
\pgfpathlineto{\pgfqpoint{5.873527in}{0.710140in}}%
\pgfpathlineto{\pgfqpoint{5.876883in}{0.726374in}}%
\pgfpathlineto{\pgfqpoint{5.880238in}{0.729952in}}%
\pgfpathlineto{\pgfqpoint{5.883594in}{0.718535in}}%
\pgfpathlineto{\pgfqpoint{5.886950in}{0.682035in}}%
\pgfpathlineto{\pgfqpoint{5.890306in}{0.671219in}}%
\pgfpathlineto{\pgfqpoint{5.893662in}{0.711102in}}%
\pgfpathlineto{\pgfqpoint{5.897017in}{0.722957in}}%
\pgfpathlineto{\pgfqpoint{5.900373in}{0.723448in}}%
\pgfpathlineto{\pgfqpoint{5.903729in}{0.712750in}}%
\pgfpathlineto{\pgfqpoint{5.907085in}{0.675716in}}%
\pgfpathlineto{\pgfqpoint{5.910441in}{0.682292in}}%
\pgfpathlineto{\pgfqpoint{5.913796in}{0.721923in}}%
\pgfpathlineto{\pgfqpoint{5.917152in}{0.736822in}}%
\pgfpathlineto{\pgfqpoint{5.920508in}{0.739230in}}%
\pgfpathlineto{\pgfqpoint{5.923864in}{0.726889in}}%
\pgfpathlineto{\pgfqpoint{5.927219in}{0.695201in}}%
\pgfpathlineto{\pgfqpoint{5.930575in}{0.650867in}}%
\pgfpathlineto{\pgfqpoint{5.933931in}{0.704824in}}%
\pgfpathlineto{\pgfqpoint{5.937287in}{0.716655in}}%
\pgfpathlineto{\pgfqpoint{5.940643in}{0.714309in}}%
\pgfpathlineto{\pgfqpoint{5.943998in}{0.695024in}}%
\pgfpathlineto{\pgfqpoint{5.947354in}{0.610952in}}%
\pgfpathlineto{\pgfqpoint{5.950710in}{0.707754in}}%
\pgfpathlineto{\pgfqpoint{5.954066in}{0.730081in}}%
\pgfpathlineto{\pgfqpoint{5.957421in}{0.737239in}}%
\pgfpathlineto{\pgfqpoint{5.960777in}{0.729799in}}%
\pgfpathlineto{\pgfqpoint{5.964133in}{0.708542in}}%
\pgfpathlineto{\pgfqpoint{5.967489in}{0.652079in}}%
\pgfpathlineto{\pgfqpoint{5.974200in}{0.696972in}}%
\pgfpathlineto{\pgfqpoint{5.977556in}{0.697513in}}%
\pgfpathlineto{\pgfqpoint{5.980912in}{0.677596in}}%
\pgfpathlineto{\pgfqpoint{5.984268in}{0.627216in}}%
\pgfpathlineto{\pgfqpoint{5.987623in}{0.698961in}}%
\pgfpathlineto{\pgfqpoint{5.990979in}{0.719629in}}%
\pgfpathlineto{\pgfqpoint{5.994335in}{0.725277in}}%
\pgfpathlineto{\pgfqpoint{5.997691in}{0.715306in}}%
\pgfpathlineto{\pgfqpoint{6.001047in}{0.689517in}}%
\pgfpathlineto{\pgfqpoint{6.004402in}{0.537378in}}%
\pgfpathlineto{\pgfqpoint{6.007758in}{0.678375in}}%
\pgfpathlineto{\pgfqpoint{6.011114in}{0.694118in}}%
\pgfpathlineto{\pgfqpoint{6.014470in}{0.692661in}}%
\pgfpathlineto{\pgfqpoint{6.017825in}{0.672314in}}%
\pgfpathlineto{\pgfqpoint{6.021181in}{0.615337in}}%
\pgfpathlineto{\pgfqpoint{6.024537in}{0.691089in}}%
\pgfpathlineto{\pgfqpoint{6.027893in}{0.712025in}}%
\pgfpathlineto{\pgfqpoint{6.031249in}{0.718040in}}%
\pgfpathlineto{\pgfqpoint{6.034604in}{0.708412in}}%
\pgfpathlineto{\pgfqpoint{6.037960in}{0.677543in}}%
\pgfpathlineto{\pgfqpoint{6.041316in}{0.640408in}}%
\pgfpathlineto{\pgfqpoint{6.044672in}{0.692669in}}%
\pgfpathlineto{\pgfqpoint{6.048027in}{0.706229in}}%
\pgfpathlineto{\pgfqpoint{6.051383in}{0.707217in}}%
\pgfpathlineto{\pgfqpoint{6.054739in}{0.696386in}}%
\pgfpathlineto{\pgfqpoint{6.058095in}{0.656990in}}%
\pgfpathlineto{\pgfqpoint{6.061451in}{0.670255in}}%
\pgfpathlineto{\pgfqpoint{6.064806in}{0.707849in}}%
\pgfpathlineto{\pgfqpoint{6.068162in}{0.722043in}}%
\pgfpathlineto{\pgfqpoint{6.071518in}{0.723507in}}%
\pgfpathlineto{\pgfqpoint{6.074874in}{0.708490in}}%
\pgfpathlineto{\pgfqpoint{6.078230in}{0.655609in}}%
\pgfpathlineto{\pgfqpoint{6.084941in}{0.713158in}}%
\pgfpathlineto{\pgfqpoint{6.088297in}{0.722400in}}%
\pgfpathlineto{\pgfqpoint{6.091653in}{0.721526in}}%
\pgfpathlineto{\pgfqpoint{6.095008in}{0.709580in}}%
\pgfpathlineto{\pgfqpoint{6.098364in}{0.668756in}}%
\pgfpathlineto{\pgfqpoint{6.101720in}{0.683828in}}%
\pgfpathlineto{\pgfqpoint{6.105076in}{0.720730in}}%
\pgfpathlineto{\pgfqpoint{6.108432in}{0.734922in}}%
\pgfpathlineto{\pgfqpoint{6.111787in}{0.736973in}}%
\pgfpathlineto{\pgfqpoint{6.115143in}{0.724067in}}%
\pgfpathlineto{\pgfqpoint{6.118499in}{0.687228in}}%
\pgfpathlineto{\pgfqpoint{6.121855in}{0.673016in}}%
\pgfpathlineto{\pgfqpoint{6.125210in}{0.712961in}}%
\pgfpathlineto{\pgfqpoint{6.128566in}{0.723789in}}%
\pgfpathlineto{\pgfqpoint{6.131922in}{0.722444in}}%
\pgfpathlineto{\pgfqpoint{6.135278in}{0.707516in}}%
\pgfpathlineto{\pgfqpoint{6.138634in}{0.638879in}}%
\pgfpathlineto{\pgfqpoint{6.141989in}{0.702180in}}%
\pgfpathlineto{\pgfqpoint{6.145345in}{0.730369in}}%
\pgfpathlineto{\pgfqpoint{6.148701in}{0.740944in}}%
\pgfpathlineto{\pgfqpoint{6.152057in}{0.738249in}}%
\pgfpathlineto{\pgfqpoint{6.155412in}{0.721197in}}%
\pgfpathlineto{\pgfqpoint{6.158768in}{0.681641in}}%
\pgfpathlineto{\pgfqpoint{6.162124in}{0.663830in}}%
\pgfpathlineto{\pgfqpoint{6.165480in}{0.700758in}}%
\pgfpathlineto{\pgfqpoint{6.168836in}{0.706796in}}%
\pgfpathlineto{\pgfqpoint{6.172191in}{0.696000in}}%
\pgfpathlineto{\pgfqpoint{6.175547in}{0.638307in}}%
\pgfpathlineto{\pgfqpoint{6.178903in}{0.689530in}}%
\pgfpathlineto{\pgfqpoint{6.182259in}{0.719746in}}%
\pgfpathlineto{\pgfqpoint{6.185614in}{0.730829in}}%
\pgfpathlineto{\pgfqpoint{6.188970in}{0.727389in}}%
\pgfpathlineto{\pgfqpoint{6.192326in}{0.709697in}}%
\pgfpathlineto{\pgfqpoint{6.195682in}{0.669667in}}%
\pgfpathlineto{\pgfqpoint{6.199038in}{0.650857in}}%
\pgfpathlineto{\pgfqpoint{6.202393in}{0.687138in}}%
\pgfpathlineto{\pgfqpoint{6.205749in}{0.691870in}}%
\pgfpathlineto{\pgfqpoint{6.209105in}{0.677772in}}%
\pgfpathlineto{\pgfqpoint{6.212461in}{0.544612in}}%
\pgfpathlineto{\pgfqpoint{6.215816in}{0.685242in}}%
\pgfpathlineto{\pgfqpoint{6.219172in}{0.710394in}}%
\pgfpathlineto{\pgfqpoint{6.222528in}{0.718882in}}%
\pgfpathlineto{\pgfqpoint{6.225884in}{0.712062in}}%
\pgfpathlineto{\pgfqpoint{6.229240in}{0.688114in}}%
\pgfpathlineto{\pgfqpoint{6.232595in}{0.550528in}}%
\pgfpathlineto{\pgfqpoint{6.235951in}{0.679253in}}%
\pgfpathlineto{\pgfqpoint{6.239307in}{0.697305in}}%
\pgfpathlineto{\pgfqpoint{6.242663in}{0.699453in}}%
\pgfpathlineto{\pgfqpoint{6.246019in}{0.688276in}}%
\pgfpathlineto{\pgfqpoint{6.249374in}{0.642464in}}%
\pgfpathlineto{\pgfqpoint{6.252730in}{0.670866in}}%
\pgfpathlineto{\pgfqpoint{6.256086in}{0.704707in}}%
\pgfpathlineto{\pgfqpoint{6.259442in}{0.717456in}}%
\pgfpathlineto{\pgfqpoint{6.262797in}{0.716856in}}%
\pgfpathlineto{\pgfqpoint{6.266153in}{0.698355in}}%
\pgfpathlineto{\pgfqpoint{6.269509in}{0.619014in}}%
\pgfpathlineto{\pgfqpoint{6.272865in}{0.686935in}}%
\pgfpathlineto{\pgfqpoint{6.276221in}{0.710353in}}%
\pgfpathlineto{\pgfqpoint{6.279576in}{0.717647in}}%
\pgfpathlineto{\pgfqpoint{6.282932in}{0.715161in}}%
\pgfpathlineto{\pgfqpoint{6.286288in}{0.700651in}}%
\pgfpathlineto{\pgfqpoint{6.289644in}{0.644942in}}%
\pgfpathlineto{\pgfqpoint{6.292999in}{0.688040in}}%
\pgfpathlineto{\pgfqpoint{6.296355in}{0.718771in}}%
\pgfpathlineto{\pgfqpoint{6.299711in}{0.730795in}}%
\pgfpathlineto{\pgfqpoint{6.303067in}{0.730704in}}%
\pgfpathlineto{\pgfqpoint{6.306423in}{0.714235in}}%
\pgfpathlineto{\pgfqpoint{6.309778in}{0.656234in}}%
\pgfpathlineto{\pgfqpoint{6.313134in}{0.692935in}}%
\pgfpathlineto{\pgfqpoint{6.316490in}{0.719719in}}%
\pgfpathlineto{\pgfqpoint{6.319846in}{0.727999in}}%
\pgfpathlineto{\pgfqpoint{6.323201in}{0.726009in}}%
\pgfpathlineto{\pgfqpoint{6.326557in}{0.711866in}}%
\pgfpathlineto{\pgfqpoint{6.329913in}{0.657299in}}%
\pgfpathlineto{\pgfqpoint{6.333269in}{0.698828in}}%
\pgfpathlineto{\pgfqpoint{6.336625in}{0.729833in}}%
\pgfpathlineto{\pgfqpoint{6.339980in}{0.741881in}}%
\pgfpathlineto{\pgfqpoint{6.343336in}{0.741612in}}%
\pgfpathlineto{\pgfqpoint{6.346692in}{0.726494in}}%
\pgfpathlineto{\pgfqpoint{6.350048in}{0.689398in}}%
\pgfpathlineto{\pgfqpoint{6.353403in}{0.668031in}}%
\pgfpathlineto{\pgfqpoint{6.356759in}{0.708576in}}%
\pgfpathlineto{\pgfqpoint{6.360115in}{0.717342in}}%
\pgfpathlineto{\pgfqpoint{6.363471in}{0.711769in}}%
\pgfpathlineto{\pgfqpoint{6.366827in}{0.683713in}}%
\pgfpathlineto{\pgfqpoint{6.370182in}{0.668115in}}%
\pgfpathlineto{\pgfqpoint{6.373538in}{0.717691in}}%
\pgfpathlineto{\pgfqpoint{6.376894in}{0.734804in}}%
\pgfpathlineto{\pgfqpoint{6.380250in}{0.738055in}}%
\pgfpathlineto{\pgfqpoint{6.383605in}{0.728425in}}%
\pgfpathlineto{\pgfqpoint{6.383605in}{0.728425in}}%
\pgfusepath{stroke}%
\end{pgfscope}%
\begin{pgfscope}%
\pgfpathrectangle{\pgfqpoint{4.621821in}{0.320625in}}{\pgfqpoint{1.845679in}{1.164375in}}%
\pgfusepath{clip}%
\pgfsetrectcap%
\pgfsetroundjoin%
\pgfsetlinewidth{1.154312pt}%
\definecolor{currentstroke}{rgb}{0.000000,0.529412,0.423529}%
\pgfsetstrokecolor{currentstroke}%
\pgfsetdash{}{0pt}%
\pgfpathmoveto{\pgfqpoint{4.705715in}{1.469210in}}%
\pgfpathlineto{\pgfqpoint{4.719139in}{1.444446in}}%
\pgfpathlineto{\pgfqpoint{4.729206in}{1.423292in}}%
\pgfpathlineto{\pgfqpoint{4.745985in}{1.384736in}}%
\pgfpathlineto{\pgfqpoint{4.772831in}{1.323071in}}%
\pgfpathlineto{\pgfqpoint{4.792966in}{1.279649in}}%
\pgfpathlineto{\pgfqpoint{4.816456in}{1.231891in}}%
\pgfpathlineto{\pgfqpoint{4.843302in}{1.180129in}}%
\pgfpathlineto{\pgfqpoint{4.873504in}{1.124407in}}%
\pgfpathlineto{\pgfqpoint{4.910418in}{1.058822in}}%
\pgfpathlineto{\pgfqpoint{4.954043in}{0.983823in}}%
\pgfpathlineto{\pgfqpoint{5.004380in}{0.899644in}}%
\pgfpathlineto{\pgfqpoint{5.064784in}{0.800999in}}%
\pgfpathlineto{\pgfqpoint{5.111765in}{0.726902in}}%
\pgfpathlineto{\pgfqpoint{5.121832in}{0.712699in}}%
\pgfpathlineto{\pgfqpoint{5.128544in}{0.700793in}}%
\pgfpathlineto{\pgfqpoint{5.141967in}{0.674829in}}%
\pgfpathlineto{\pgfqpoint{5.145323in}{0.671023in}}%
\pgfpathlineto{\pgfqpoint{5.148678in}{0.669038in}}%
\pgfpathlineto{\pgfqpoint{5.155390in}{0.666309in}}%
\pgfpathlineto{\pgfqpoint{5.158746in}{0.662719in}}%
\pgfpathlineto{\pgfqpoint{5.162102in}{0.658265in}}%
\pgfpathlineto{\pgfqpoint{5.165457in}{0.656023in}}%
\pgfpathlineto{\pgfqpoint{5.168813in}{0.658268in}}%
\pgfpathlineto{\pgfqpoint{5.175525in}{0.668109in}}%
\pgfpathlineto{\pgfqpoint{5.182236in}{0.674654in}}%
\pgfpathlineto{\pgfqpoint{5.185592in}{0.680521in}}%
\pgfpathlineto{\pgfqpoint{5.195659in}{0.704713in}}%
\pgfpathlineto{\pgfqpoint{5.199015in}{0.710719in}}%
\pgfpathlineto{\pgfqpoint{5.202371in}{0.715070in}}%
\pgfpathlineto{\pgfqpoint{5.205727in}{0.717810in}}%
\pgfpathlineto{\pgfqpoint{5.209082in}{0.718984in}}%
\pgfpathlineto{\pgfqpoint{5.212438in}{0.718590in}}%
\pgfpathlineto{\pgfqpoint{5.215794in}{0.716553in}}%
\pgfpathlineto{\pgfqpoint{5.219150in}{0.712718in}}%
\pgfpathlineto{\pgfqpoint{5.222506in}{0.706894in}}%
\pgfpathlineto{\pgfqpoint{5.229217in}{0.691593in}}%
\pgfpathlineto{\pgfqpoint{5.232573in}{0.691279in}}%
\pgfpathlineto{\pgfqpoint{5.242640in}{0.722420in}}%
\pgfpathlineto{\pgfqpoint{5.245996in}{0.728442in}}%
\pgfpathlineto{\pgfqpoint{5.249352in}{0.731463in}}%
\pgfpathlineto{\pgfqpoint{5.252708in}{0.731398in}}%
\pgfpathlineto{\pgfqpoint{5.256063in}{0.727852in}}%
\pgfpathlineto{\pgfqpoint{5.259419in}{0.721853in}}%
\pgfpathlineto{\pgfqpoint{5.266131in}{0.704928in}}%
\pgfpathlineto{\pgfqpoint{5.269487in}{0.700761in}}%
\pgfpathlineto{\pgfqpoint{5.272842in}{0.700776in}}%
\pgfpathlineto{\pgfqpoint{5.276198in}{0.699311in}}%
\pgfpathlineto{\pgfqpoint{5.279554in}{0.696363in}}%
\pgfpathlineto{\pgfqpoint{5.282910in}{0.695960in}}%
\pgfpathlineto{\pgfqpoint{5.286265in}{0.699623in}}%
\pgfpathlineto{\pgfqpoint{5.289621in}{0.704129in}}%
\pgfpathlineto{\pgfqpoint{5.292977in}{0.707673in}}%
\pgfpathlineto{\pgfqpoint{5.296333in}{0.707671in}}%
\pgfpathlineto{\pgfqpoint{5.299689in}{0.700236in}}%
\pgfpathlineto{\pgfqpoint{5.303044in}{0.682634in}}%
\pgfpathlineto{\pgfqpoint{5.306400in}{0.677001in}}%
\pgfpathlineto{\pgfqpoint{5.309756in}{0.699736in}}%
\pgfpathlineto{\pgfqpoint{5.313112in}{0.713858in}}%
\pgfpathlineto{\pgfqpoint{5.316467in}{0.718362in}}%
\pgfpathlineto{\pgfqpoint{5.319823in}{0.713463in}}%
\pgfpathlineto{\pgfqpoint{5.323179in}{0.700556in}}%
\pgfpathlineto{\pgfqpoint{5.326535in}{0.677562in}}%
\pgfpathlineto{\pgfqpoint{5.329891in}{0.675441in}}%
\pgfpathlineto{\pgfqpoint{5.333246in}{0.694293in}}%
\pgfpathlineto{\pgfqpoint{5.336602in}{0.704373in}}%
\pgfpathlineto{\pgfqpoint{5.339958in}{0.708202in}}%
\pgfpathlineto{\pgfqpoint{5.343314in}{0.707763in}}%
\pgfpathlineto{\pgfqpoint{5.346669in}{0.706208in}}%
\pgfpathlineto{\pgfqpoint{5.353381in}{0.711981in}}%
\pgfpathlineto{\pgfqpoint{5.360093in}{0.709615in}}%
\pgfpathlineto{\pgfqpoint{5.363448in}{0.705292in}}%
\pgfpathlineto{\pgfqpoint{5.366804in}{0.697905in}}%
\pgfpathlineto{\pgfqpoint{5.370160in}{0.695930in}}%
\pgfpathlineto{\pgfqpoint{5.376871in}{0.713767in}}%
\pgfpathlineto{\pgfqpoint{5.380227in}{0.717934in}}%
\pgfpathlineto{\pgfqpoint{5.383583in}{0.714658in}}%
\pgfpathlineto{\pgfqpoint{5.390295in}{0.699567in}}%
\pgfpathlineto{\pgfqpoint{5.393650in}{0.697107in}}%
\pgfpathlineto{\pgfqpoint{5.397006in}{0.707795in}}%
\pgfpathlineto{\pgfqpoint{5.400362in}{0.716443in}}%
\pgfpathlineto{\pgfqpoint{5.403718in}{0.717201in}}%
\pgfpathlineto{\pgfqpoint{5.407074in}{0.711119in}}%
\pgfpathlineto{\pgfqpoint{5.410429in}{0.703857in}}%
\pgfpathlineto{\pgfqpoint{5.413785in}{0.704755in}}%
\pgfpathlineto{\pgfqpoint{5.417141in}{0.712942in}}%
\pgfpathlineto{\pgfqpoint{5.420497in}{0.716598in}}%
\pgfpathlineto{\pgfqpoint{5.423852in}{0.715154in}}%
\pgfpathlineto{\pgfqpoint{5.427208in}{0.712385in}}%
\pgfpathlineto{\pgfqpoint{5.430564in}{0.706681in}}%
\pgfpathlineto{\pgfqpoint{5.433920in}{0.705502in}}%
\pgfpathlineto{\pgfqpoint{5.437276in}{0.714237in}}%
\pgfpathlineto{\pgfqpoint{5.440631in}{0.721061in}}%
\pgfpathlineto{\pgfqpoint{5.443987in}{0.721133in}}%
\pgfpathlineto{\pgfqpoint{5.447343in}{0.718903in}}%
\pgfpathlineto{\pgfqpoint{5.454054in}{0.712443in}}%
\pgfpathlineto{\pgfqpoint{5.460766in}{0.716414in}}%
\pgfpathlineto{\pgfqpoint{5.467478in}{0.723981in}}%
\pgfpathlineto{\pgfqpoint{5.470833in}{0.721888in}}%
\pgfpathlineto{\pgfqpoint{5.474189in}{0.717837in}}%
\pgfpathlineto{\pgfqpoint{5.477545in}{0.715929in}}%
\pgfpathlineto{\pgfqpoint{5.480901in}{0.715073in}}%
\pgfpathlineto{\pgfqpoint{5.487612in}{0.720394in}}%
\pgfpathlineto{\pgfqpoint{5.494324in}{0.720987in}}%
\pgfpathlineto{\pgfqpoint{5.497680in}{0.724272in}}%
\pgfpathlineto{\pgfqpoint{5.501035in}{0.725756in}}%
\pgfpathlineto{\pgfqpoint{5.504391in}{0.726249in}}%
\pgfpathlineto{\pgfqpoint{5.507747in}{0.723344in}}%
\pgfpathlineto{\pgfqpoint{5.511103in}{0.716944in}}%
\pgfpathlineto{\pgfqpoint{5.514458in}{0.713790in}}%
\pgfpathlineto{\pgfqpoint{5.517814in}{0.712169in}}%
\pgfpathlineto{\pgfqpoint{5.521170in}{0.711725in}}%
\pgfpathlineto{\pgfqpoint{5.524526in}{0.712439in}}%
\pgfpathlineto{\pgfqpoint{5.527882in}{0.714328in}}%
\pgfpathlineto{\pgfqpoint{5.531237in}{0.717454in}}%
\pgfpathlineto{\pgfqpoint{5.534593in}{0.725279in}}%
\pgfpathlineto{\pgfqpoint{5.537949in}{0.730099in}}%
\pgfpathlineto{\pgfqpoint{5.541305in}{0.727061in}}%
\pgfpathlineto{\pgfqpoint{5.544660in}{0.714323in}}%
\pgfpathlineto{\pgfqpoint{5.548016in}{0.698822in}}%
\pgfpathlineto{\pgfqpoint{5.551372in}{0.698633in}}%
\pgfpathlineto{\pgfqpoint{5.554728in}{0.703956in}}%
\pgfpathlineto{\pgfqpoint{5.558084in}{0.704594in}}%
\pgfpathlineto{\pgfqpoint{5.561439in}{0.703421in}}%
\pgfpathlineto{\pgfqpoint{5.564795in}{0.704387in}}%
\pgfpathlineto{\pgfqpoint{5.568151in}{0.709370in}}%
\pgfpathlineto{\pgfqpoint{5.571507in}{0.718982in}}%
\pgfpathlineto{\pgfqpoint{5.574863in}{0.724536in}}%
\pgfpathlineto{\pgfqpoint{5.578218in}{0.721482in}}%
\pgfpathlineto{\pgfqpoint{5.584930in}{0.688838in}}%
\pgfpathlineto{\pgfqpoint{5.588286in}{0.698086in}}%
\pgfpathlineto{\pgfqpoint{5.591641in}{0.709887in}}%
\pgfpathlineto{\pgfqpoint{5.594997in}{0.711767in}}%
\pgfpathlineto{\pgfqpoint{5.601709in}{0.694861in}}%
\pgfpathlineto{\pgfqpoint{5.605065in}{0.699716in}}%
\pgfpathlineto{\pgfqpoint{5.608420in}{0.712999in}}%
\pgfpathlineto{\pgfqpoint{5.611776in}{0.720754in}}%
\pgfpathlineto{\pgfqpoint{5.615132in}{0.719504in}}%
\pgfpathlineto{\pgfqpoint{5.618488in}{0.706906in}}%
\pgfpathlineto{\pgfqpoint{5.621843in}{0.686433in}}%
\pgfpathlineto{\pgfqpoint{5.625199in}{0.695383in}}%
\pgfpathlineto{\pgfqpoint{5.628555in}{0.713720in}}%
\pgfpathlineto{\pgfqpoint{5.631911in}{0.722310in}}%
\pgfpathlineto{\pgfqpoint{5.635267in}{0.721560in}}%
\pgfpathlineto{\pgfqpoint{5.641978in}{0.703575in}}%
\pgfpathlineto{\pgfqpoint{5.645334in}{0.706756in}}%
\pgfpathlineto{\pgfqpoint{5.648690in}{0.716969in}}%
\pgfpathlineto{\pgfqpoint{5.652045in}{0.722325in}}%
\pgfpathlineto{\pgfqpoint{5.655401in}{0.720237in}}%
\pgfpathlineto{\pgfqpoint{5.658757in}{0.709880in}}%
\pgfpathlineto{\pgfqpoint{5.662113in}{0.702630in}}%
\pgfpathlineto{\pgfqpoint{5.668824in}{0.722308in}}%
\pgfpathlineto{\pgfqpoint{5.672180in}{0.724322in}}%
\pgfpathlineto{\pgfqpoint{5.678892in}{0.712669in}}%
\pgfpathlineto{\pgfqpoint{5.682247in}{0.713571in}}%
\pgfpathlineto{\pgfqpoint{5.688959in}{0.728212in}}%
\pgfpathlineto{\pgfqpoint{5.692315in}{0.731287in}}%
\pgfpathlineto{\pgfqpoint{5.695671in}{0.726976in}}%
\pgfpathlineto{\pgfqpoint{5.699026in}{0.714893in}}%
\pgfpathlineto{\pgfqpoint{5.702382in}{0.710305in}}%
\pgfpathlineto{\pgfqpoint{5.705738in}{0.716536in}}%
\pgfpathlineto{\pgfqpoint{5.709094in}{0.719321in}}%
\pgfpathlineto{\pgfqpoint{5.712449in}{0.715828in}}%
\pgfpathlineto{\pgfqpoint{5.715805in}{0.708963in}}%
\pgfpathlineto{\pgfqpoint{5.719161in}{0.709741in}}%
\pgfpathlineto{\pgfqpoint{5.722517in}{0.720764in}}%
\pgfpathlineto{\pgfqpoint{5.725873in}{0.729776in}}%
\pgfpathlineto{\pgfqpoint{5.729228in}{0.734497in}}%
\pgfpathlineto{\pgfqpoint{5.732584in}{0.731369in}}%
\pgfpathlineto{\pgfqpoint{5.735940in}{0.717704in}}%
\pgfpathlineto{\pgfqpoint{5.739296in}{0.699323in}}%
\pgfpathlineto{\pgfqpoint{5.742652in}{0.698770in}}%
\pgfpathlineto{\pgfqpoint{5.746007in}{0.706658in}}%
\pgfpathlineto{\pgfqpoint{5.749363in}{0.707701in}}%
\pgfpathlineto{\pgfqpoint{5.752719in}{0.704616in}}%
\pgfpathlineto{\pgfqpoint{5.756075in}{0.705516in}}%
\pgfpathlineto{\pgfqpoint{5.762786in}{0.721738in}}%
\pgfpathlineto{\pgfqpoint{5.766142in}{0.725937in}}%
\pgfpathlineto{\pgfqpoint{5.769498in}{0.724875in}}%
\pgfpathlineto{\pgfqpoint{5.772854in}{0.718566in}}%
\pgfpathlineto{\pgfqpoint{5.776209in}{0.704935in}}%
\pgfpathlineto{\pgfqpoint{5.779565in}{0.698860in}}%
\pgfpathlineto{\pgfqpoint{5.782921in}{0.703355in}}%
\pgfpathlineto{\pgfqpoint{5.786277in}{0.703735in}}%
\pgfpathlineto{\pgfqpoint{5.789632in}{0.695420in}}%
\pgfpathlineto{\pgfqpoint{5.792988in}{0.692106in}}%
\pgfpathlineto{\pgfqpoint{5.799700in}{0.715472in}}%
\pgfpathlineto{\pgfqpoint{5.803056in}{0.720235in}}%
\pgfpathlineto{\pgfqpoint{5.806411in}{0.718711in}}%
\pgfpathlineto{\pgfqpoint{5.809767in}{0.710758in}}%
\pgfpathlineto{\pgfqpoint{5.813123in}{0.696946in}}%
\pgfpathlineto{\pgfqpoint{5.816479in}{0.697571in}}%
\pgfpathlineto{\pgfqpoint{5.819834in}{0.709436in}}%
\pgfpathlineto{\pgfqpoint{5.823190in}{0.716712in}}%
\pgfpathlineto{\pgfqpoint{5.826546in}{0.716346in}}%
\pgfpathlineto{\pgfqpoint{5.829902in}{0.706348in}}%
\pgfpathlineto{\pgfqpoint{5.833258in}{0.700392in}}%
\pgfpathlineto{\pgfqpoint{5.836613in}{0.710681in}}%
\pgfpathlineto{\pgfqpoint{5.839969in}{0.719176in}}%
\pgfpathlineto{\pgfqpoint{5.843325in}{0.719704in}}%
\pgfpathlineto{\pgfqpoint{5.846681in}{0.710543in}}%
\pgfpathlineto{\pgfqpoint{5.850036in}{0.696467in}}%
\pgfpathlineto{\pgfqpoint{5.853392in}{0.701668in}}%
\pgfpathlineto{\pgfqpoint{5.856748in}{0.715899in}}%
\pgfpathlineto{\pgfqpoint{5.860104in}{0.724117in}}%
\pgfpathlineto{\pgfqpoint{5.863460in}{0.725370in}}%
\pgfpathlineto{\pgfqpoint{5.866815in}{0.717430in}}%
\pgfpathlineto{\pgfqpoint{5.870171in}{0.703302in}}%
\pgfpathlineto{\pgfqpoint{5.873527in}{0.710278in}}%
\pgfpathlineto{\pgfqpoint{5.876883in}{0.725145in}}%
\pgfpathlineto{\pgfqpoint{5.880238in}{0.731679in}}%
\pgfpathlineto{\pgfqpoint{5.883594in}{0.728622in}}%
\pgfpathlineto{\pgfqpoint{5.886950in}{0.712570in}}%
\pgfpathlineto{\pgfqpoint{5.890306in}{0.691908in}}%
\pgfpathlineto{\pgfqpoint{5.897017in}{0.723102in}}%
\pgfpathlineto{\pgfqpoint{5.900373in}{0.727003in}}%
\pgfpathlineto{\pgfqpoint{5.903729in}{0.720431in}}%
\pgfpathlineto{\pgfqpoint{5.907085in}{0.702888in}}%
\pgfpathlineto{\pgfqpoint{5.910441in}{0.698423in}}%
\pgfpathlineto{\pgfqpoint{5.913796in}{0.722073in}}%
\pgfpathlineto{\pgfqpoint{5.917152in}{0.735728in}}%
\pgfpathlineto{\pgfqpoint{5.920508in}{0.738494in}}%
\pgfpathlineto{\pgfqpoint{5.923864in}{0.730762in}}%
\pgfpathlineto{\pgfqpoint{5.927219in}{0.714868in}}%
\pgfpathlineto{\pgfqpoint{5.930575in}{0.695119in}}%
\pgfpathlineto{\pgfqpoint{5.933931in}{0.702961in}}%
\pgfpathlineto{\pgfqpoint{5.937287in}{0.713729in}}%
\pgfpathlineto{\pgfqpoint{5.940643in}{0.714756in}}%
\pgfpathlineto{\pgfqpoint{5.943998in}{0.709181in}}%
\pgfpathlineto{\pgfqpoint{5.947354in}{0.706150in}}%
\pgfpathlineto{\pgfqpoint{5.954066in}{0.726766in}}%
\pgfpathlineto{\pgfqpoint{5.957421in}{0.729568in}}%
\pgfpathlineto{\pgfqpoint{5.964133in}{0.724494in}}%
\pgfpathlineto{\pgfqpoint{5.967489in}{0.717514in}}%
\pgfpathlineto{\pgfqpoint{5.970845in}{0.706375in}}%
\pgfpathlineto{\pgfqpoint{5.974200in}{0.700123in}}%
\pgfpathlineto{\pgfqpoint{5.977556in}{0.697674in}}%
\pgfpathlineto{\pgfqpoint{5.980912in}{0.694048in}}%
\pgfpathlineto{\pgfqpoint{5.984268in}{0.697963in}}%
\pgfpathlineto{\pgfqpoint{5.987623in}{0.708767in}}%
\pgfpathlineto{\pgfqpoint{5.990979in}{0.716489in}}%
\pgfpathlineto{\pgfqpoint{5.994335in}{0.717890in}}%
\pgfpathlineto{\pgfqpoint{5.997691in}{0.716974in}}%
\pgfpathlineto{\pgfqpoint{6.001047in}{0.717561in}}%
\pgfpathlineto{\pgfqpoint{6.004402in}{0.713391in}}%
\pgfpathlineto{\pgfqpoint{6.007758in}{0.704581in}}%
\pgfpathlineto{\pgfqpoint{6.011114in}{0.702916in}}%
\pgfpathlineto{\pgfqpoint{6.014470in}{0.707304in}}%
\pgfpathlineto{\pgfqpoint{6.017825in}{0.708990in}}%
\pgfpathlineto{\pgfqpoint{6.021181in}{0.703984in}}%
\pgfpathlineto{\pgfqpoint{6.024537in}{0.702293in}}%
\pgfpathlineto{\pgfqpoint{6.027893in}{0.709510in}}%
\pgfpathlineto{\pgfqpoint{6.031249in}{0.713972in}}%
\pgfpathlineto{\pgfqpoint{6.037960in}{0.714072in}}%
\pgfpathlineto{\pgfqpoint{6.044672in}{0.709608in}}%
\pgfpathlineto{\pgfqpoint{6.048027in}{0.710886in}}%
\pgfpathlineto{\pgfqpoint{6.054739in}{0.718267in}}%
\pgfpathlineto{\pgfqpoint{6.058095in}{0.718740in}}%
\pgfpathlineto{\pgfqpoint{6.061451in}{0.714958in}}%
\pgfpathlineto{\pgfqpoint{6.064806in}{0.718250in}}%
\pgfpathlineto{\pgfqpoint{6.068162in}{0.724689in}}%
\pgfpathlineto{\pgfqpoint{6.071518in}{0.724008in}}%
\pgfpathlineto{\pgfqpoint{6.074874in}{0.713078in}}%
\pgfpathlineto{\pgfqpoint{6.078230in}{0.697876in}}%
\pgfpathlineto{\pgfqpoint{6.081585in}{0.702608in}}%
\pgfpathlineto{\pgfqpoint{6.084941in}{0.717610in}}%
\pgfpathlineto{\pgfqpoint{6.088297in}{0.725810in}}%
\pgfpathlineto{\pgfqpoint{6.091653in}{0.727148in}}%
\pgfpathlineto{\pgfqpoint{6.095008in}{0.721068in}}%
\pgfpathlineto{\pgfqpoint{6.098364in}{0.705923in}}%
\pgfpathlineto{\pgfqpoint{6.101720in}{0.709191in}}%
\pgfpathlineto{\pgfqpoint{6.105076in}{0.727679in}}%
\pgfpathlineto{\pgfqpoint{6.108432in}{0.737383in}}%
\pgfpathlineto{\pgfqpoint{6.111787in}{0.736265in}}%
\pgfpathlineto{\pgfqpoint{6.115143in}{0.725492in}}%
\pgfpathlineto{\pgfqpoint{6.118499in}{0.706872in}}%
\pgfpathlineto{\pgfqpoint{6.121855in}{0.693309in}}%
\pgfpathlineto{\pgfqpoint{6.125210in}{0.709043in}}%
\pgfpathlineto{\pgfqpoint{6.128566in}{0.719624in}}%
\pgfpathlineto{\pgfqpoint{6.131922in}{0.720950in}}%
\pgfpathlineto{\pgfqpoint{6.135278in}{0.716053in}}%
\pgfpathlineto{\pgfqpoint{6.138634in}{0.713122in}}%
\pgfpathlineto{\pgfqpoint{6.145345in}{0.728263in}}%
\pgfpathlineto{\pgfqpoint{6.148701in}{0.731577in}}%
\pgfpathlineto{\pgfqpoint{6.152057in}{0.731332in}}%
\pgfpathlineto{\pgfqpoint{6.155412in}{0.727402in}}%
\pgfpathlineto{\pgfqpoint{6.158768in}{0.722613in}}%
\pgfpathlineto{\pgfqpoint{6.165480in}{0.700404in}}%
\pgfpathlineto{\pgfqpoint{6.168836in}{0.698823in}}%
\pgfpathlineto{\pgfqpoint{6.172191in}{0.703485in}}%
\pgfpathlineto{\pgfqpoint{6.175547in}{0.709256in}}%
\pgfpathlineto{\pgfqpoint{6.178903in}{0.713281in}}%
\pgfpathlineto{\pgfqpoint{6.182259in}{0.714965in}}%
\pgfpathlineto{\pgfqpoint{6.185614in}{0.717882in}}%
\pgfpathlineto{\pgfqpoint{6.192326in}{0.720032in}}%
\pgfpathlineto{\pgfqpoint{6.195682in}{0.722683in}}%
\pgfpathlineto{\pgfqpoint{6.199038in}{0.718345in}}%
\pgfpathlineto{\pgfqpoint{6.202393in}{0.707891in}}%
\pgfpathlineto{\pgfqpoint{6.205749in}{0.702142in}}%
\pgfpathlineto{\pgfqpoint{6.209105in}{0.702818in}}%
\pgfpathlineto{\pgfqpoint{6.212461in}{0.701470in}}%
\pgfpathlineto{\pgfqpoint{6.215816in}{0.699359in}}%
\pgfpathlineto{\pgfqpoint{6.219172in}{0.701409in}}%
\pgfpathlineto{\pgfqpoint{6.222528in}{0.709264in}}%
\pgfpathlineto{\pgfqpoint{6.225884in}{0.715087in}}%
\pgfpathlineto{\pgfqpoint{6.229240in}{0.718644in}}%
\pgfpathlineto{\pgfqpoint{6.232595in}{0.718916in}}%
\pgfpathlineto{\pgfqpoint{6.235951in}{0.711125in}}%
\pgfpathlineto{\pgfqpoint{6.239307in}{0.698731in}}%
\pgfpathlineto{\pgfqpoint{6.242663in}{0.698804in}}%
\pgfpathlineto{\pgfqpoint{6.246019in}{0.711875in}}%
\pgfpathlineto{\pgfqpoint{6.249374in}{0.720916in}}%
\pgfpathlineto{\pgfqpoint{6.252730in}{0.721958in}}%
\pgfpathlineto{\pgfqpoint{6.256086in}{0.716787in}}%
\pgfpathlineto{\pgfqpoint{6.259442in}{0.708947in}}%
\pgfpathlineto{\pgfqpoint{6.262797in}{0.706094in}}%
\pgfpathlineto{\pgfqpoint{6.266153in}{0.711091in}}%
\pgfpathlineto{\pgfqpoint{6.269509in}{0.717581in}}%
\pgfpathlineto{\pgfqpoint{6.272865in}{0.719858in}}%
\pgfpathlineto{\pgfqpoint{6.276221in}{0.718777in}}%
\pgfpathlineto{\pgfqpoint{6.279576in}{0.716034in}}%
\pgfpathlineto{\pgfqpoint{6.282932in}{0.715463in}}%
\pgfpathlineto{\pgfqpoint{6.292999in}{0.727643in}}%
\pgfpathlineto{\pgfqpoint{6.296355in}{0.730840in}}%
\pgfpathlineto{\pgfqpoint{6.299711in}{0.728806in}}%
\pgfpathlineto{\pgfqpoint{6.303067in}{0.723144in}}%
\pgfpathlineto{\pgfqpoint{6.306423in}{0.715878in}}%
\pgfpathlineto{\pgfqpoint{6.309778in}{0.710240in}}%
\pgfpathlineto{\pgfqpoint{6.313134in}{0.710355in}}%
\pgfpathlineto{\pgfqpoint{6.319846in}{0.720400in}}%
\pgfpathlineto{\pgfqpoint{6.323201in}{0.722645in}}%
\pgfpathlineto{\pgfqpoint{6.326557in}{0.722931in}}%
\pgfpathlineto{\pgfqpoint{6.329913in}{0.725310in}}%
\pgfpathlineto{\pgfqpoint{6.333269in}{0.728912in}}%
\pgfpathlineto{\pgfqpoint{6.339980in}{0.728122in}}%
\pgfpathlineto{\pgfqpoint{6.343336in}{0.730488in}}%
\pgfpathlineto{\pgfqpoint{6.346692in}{0.729662in}}%
\pgfpathlineto{\pgfqpoint{6.350048in}{0.724468in}}%
\pgfpathlineto{\pgfqpoint{6.356759in}{0.704827in}}%
\pgfpathlineto{\pgfqpoint{6.360115in}{0.703465in}}%
\pgfpathlineto{\pgfqpoint{6.366827in}{0.706921in}}%
\pgfpathlineto{\pgfqpoint{6.370182in}{0.703746in}}%
\pgfpathlineto{\pgfqpoint{6.373538in}{0.704573in}}%
\pgfpathlineto{\pgfqpoint{6.376894in}{0.715270in}}%
\pgfpathlineto{\pgfqpoint{6.380250in}{0.720194in}}%
\pgfpathlineto{\pgfqpoint{6.383605in}{0.714345in}}%
\pgfpathlineto{\pgfqpoint{6.383605in}{0.714345in}}%
\pgfusepath{stroke}%
\end{pgfscope}%
\begin{pgfscope}%
\pgfpathrectangle{\pgfqpoint{4.621821in}{0.320625in}}{\pgfqpoint{1.845679in}{1.164375in}}%
\pgfusepath{clip}%
\pgfsetrectcap%
\pgfsetroundjoin%
\pgfsetlinewidth{0.602250pt}%
\definecolor{currentstroke}{rgb}{0.478431,0.243137,0.615686}%
\pgfsetstrokecolor{currentstroke}%
\pgfsetstrokeopacity{0.120000}%
\pgfsetdash{}{0pt}%
\pgfpathmoveto{\pgfqpoint{4.705715in}{1.469034in}}%
\pgfpathlineto{\pgfqpoint{4.715783in}{1.450518in}}%
\pgfpathlineto{\pgfqpoint{4.725850in}{1.429517in}}%
\pgfpathlineto{\pgfqpoint{4.739273in}{1.398106in}}%
\pgfpathlineto{\pgfqpoint{4.779543in}{1.301177in}}%
\pgfpathlineto{\pgfqpoint{4.799677in}{1.255951in}}%
\pgfpathlineto{\pgfqpoint{4.823168in}{1.205836in}}%
\pgfpathlineto{\pgfqpoint{4.853370in}{1.144320in}}%
\pgfpathlineto{\pgfqpoint{4.890283in}{1.071929in}}%
\pgfpathlineto{\pgfqpoint{4.937264in}{0.982479in}}%
\pgfpathlineto{\pgfqpoint{4.994313in}{0.876411in}}%
\pgfpathlineto{\pgfqpoint{5.041293in}{0.790300in}}%
\pgfpathlineto{\pgfqpoint{5.051361in}{0.771953in}}%
\pgfpathlineto{\pgfqpoint{5.068140in}{0.744699in}}%
\pgfpathlineto{\pgfqpoint{5.078207in}{0.722462in}}%
\pgfpathlineto{\pgfqpoint{5.084919in}{0.707580in}}%
\pgfpathlineto{\pgfqpoint{5.088274in}{0.701455in}}%
\pgfpathlineto{\pgfqpoint{5.091630in}{0.696886in}}%
\pgfpathlineto{\pgfqpoint{5.094986in}{0.694117in}}%
\pgfpathlineto{\pgfqpoint{5.098342in}{0.692917in}}%
\pgfpathlineto{\pgfqpoint{5.105053in}{0.691648in}}%
\pgfpathlineto{\pgfqpoint{5.108409in}{0.688670in}}%
\pgfpathlineto{\pgfqpoint{5.111765in}{0.682029in}}%
\pgfpathlineto{\pgfqpoint{5.115121in}{0.672642in}}%
\pgfpathlineto{\pgfqpoint{5.118476in}{0.660278in}}%
\pgfpathlineto{\pgfqpoint{5.121832in}{0.645420in}}%
\pgfpathlineto{\pgfqpoint{5.125188in}{0.633731in}}%
\pgfpathlineto{\pgfqpoint{5.128544in}{0.634356in}}%
\pgfpathlineto{\pgfqpoint{5.131900in}{0.639967in}}%
\pgfpathlineto{\pgfqpoint{5.135255in}{0.642681in}}%
\pgfpathlineto{\pgfqpoint{5.138611in}{0.640612in}}%
\pgfpathlineto{\pgfqpoint{5.141967in}{0.632254in}}%
\pgfpathlineto{\pgfqpoint{5.145323in}{0.613028in}}%
\pgfpathlineto{\pgfqpoint{5.148678in}{0.585166in}}%
\pgfpathlineto{\pgfqpoint{5.152034in}{0.621932in}}%
\pgfpathlineto{\pgfqpoint{5.155390in}{0.645401in}}%
\pgfpathlineto{\pgfqpoint{5.158746in}{0.659026in}}%
\pgfpathlineto{\pgfqpoint{5.162102in}{0.666890in}}%
\pgfpathlineto{\pgfqpoint{5.165457in}{0.670008in}}%
\pgfpathlineto{\pgfqpoint{5.168813in}{0.667684in}}%
\pgfpathlineto{\pgfqpoint{5.172169in}{0.657265in}}%
\pgfpathlineto{\pgfqpoint{5.175525in}{0.661053in}}%
\pgfpathlineto{\pgfqpoint{5.178880in}{0.685393in}}%
\pgfpathlineto{\pgfqpoint{5.182236in}{0.703130in}}%
\pgfpathlineto{\pgfqpoint{5.185592in}{0.714717in}}%
\pgfpathlineto{\pgfqpoint{5.188948in}{0.722096in}}%
\pgfpathlineto{\pgfqpoint{5.192304in}{0.726265in}}%
\pgfpathlineto{\pgfqpoint{5.195659in}{0.727650in}}%
\pgfpathlineto{\pgfqpoint{5.199015in}{0.726268in}}%
\pgfpathlineto{\pgfqpoint{5.202371in}{0.721691in}}%
\pgfpathlineto{\pgfqpoint{5.205727in}{0.712674in}}%
\pgfpathlineto{\pgfqpoint{5.209082in}{0.695566in}}%
\pgfpathlineto{\pgfqpoint{5.212438in}{0.651537in}}%
\pgfpathlineto{\pgfqpoint{5.215794in}{0.668542in}}%
\pgfpathlineto{\pgfqpoint{5.219150in}{0.704791in}}%
\pgfpathlineto{\pgfqpoint{5.222506in}{0.722112in}}%
\pgfpathlineto{\pgfqpoint{5.225861in}{0.732225in}}%
\pgfpathlineto{\pgfqpoint{5.229217in}{0.737673in}}%
\pgfpathlineto{\pgfqpoint{5.232573in}{0.738886in}}%
\pgfpathlineto{\pgfqpoint{5.235929in}{0.734811in}}%
\pgfpathlineto{\pgfqpoint{5.239285in}{0.723225in}}%
\pgfpathlineto{\pgfqpoint{5.242640in}{0.704086in}}%
\pgfpathlineto{\pgfqpoint{5.245996in}{0.679417in}}%
\pgfpathlineto{\pgfqpoint{5.249352in}{0.677751in}}%
\pgfpathlineto{\pgfqpoint{5.252708in}{0.684854in}}%
\pgfpathlineto{\pgfqpoint{5.256063in}{0.679443in}}%
\pgfpathlineto{\pgfqpoint{5.259419in}{0.653421in}}%
\pgfpathlineto{\pgfqpoint{5.262775in}{0.665191in}}%
\pgfpathlineto{\pgfqpoint{5.266131in}{0.697093in}}%
\pgfpathlineto{\pgfqpoint{5.269487in}{0.711121in}}%
\pgfpathlineto{\pgfqpoint{5.272842in}{0.712372in}}%
\pgfpathlineto{\pgfqpoint{5.276198in}{0.699571in}}%
\pgfpathlineto{\pgfqpoint{5.279554in}{0.675072in}}%
\pgfpathlineto{\pgfqpoint{5.282910in}{0.630724in}}%
\pgfpathlineto{\pgfqpoint{5.286265in}{0.654449in}}%
\pgfpathlineto{\pgfqpoint{5.289621in}{0.667143in}}%
\pgfpathlineto{\pgfqpoint{5.292977in}{0.661396in}}%
\pgfpathlineto{\pgfqpoint{5.296333in}{0.625249in}}%
\pgfpathlineto{\pgfqpoint{5.299689in}{0.643981in}}%
\pgfpathlineto{\pgfqpoint{5.303044in}{0.680212in}}%
\pgfpathlineto{\pgfqpoint{5.306400in}{0.694447in}}%
\pgfpathlineto{\pgfqpoint{5.309756in}{0.695300in}}%
\pgfpathlineto{\pgfqpoint{5.313112in}{0.676149in}}%
\pgfpathlineto{\pgfqpoint{5.316467in}{0.610354in}}%
\pgfpathlineto{\pgfqpoint{5.319823in}{0.685616in}}%
\pgfpathlineto{\pgfqpoint{5.323179in}{0.706947in}}%
\pgfpathlineto{\pgfqpoint{5.326535in}{0.715890in}}%
\pgfpathlineto{\pgfqpoint{5.329891in}{0.717419in}}%
\pgfpathlineto{\pgfqpoint{5.333246in}{0.711856in}}%
\pgfpathlineto{\pgfqpoint{5.336602in}{0.695475in}}%
\pgfpathlineto{\pgfqpoint{5.339958in}{0.637885in}}%
\pgfpathlineto{\pgfqpoint{5.343314in}{0.682020in}}%
\pgfpathlineto{\pgfqpoint{5.346669in}{0.712750in}}%
\pgfpathlineto{\pgfqpoint{5.350025in}{0.725857in}}%
\pgfpathlineto{\pgfqpoint{5.353381in}{0.729025in}}%
\pgfpathlineto{\pgfqpoint{5.356737in}{0.720247in}}%
\pgfpathlineto{\pgfqpoint{5.360093in}{0.695550in}}%
\pgfpathlineto{\pgfqpoint{5.363448in}{0.611065in}}%
\pgfpathlineto{\pgfqpoint{5.366804in}{0.689528in}}%
\pgfpathlineto{\pgfqpoint{5.370160in}{0.707007in}}%
\pgfpathlineto{\pgfqpoint{5.373516in}{0.709490in}}%
\pgfpathlineto{\pgfqpoint{5.376871in}{0.699631in}}%
\pgfpathlineto{\pgfqpoint{5.380227in}{0.662099in}}%
\pgfpathlineto{\pgfqpoint{5.383583in}{0.672153in}}%
\pgfpathlineto{\pgfqpoint{5.386939in}{0.711029in}}%
\pgfpathlineto{\pgfqpoint{5.390295in}{0.725788in}}%
\pgfpathlineto{\pgfqpoint{5.393650in}{0.727926in}}%
\pgfpathlineto{\pgfqpoint{5.397006in}{0.715121in}}%
\pgfpathlineto{\pgfqpoint{5.403718in}{0.650954in}}%
\pgfpathlineto{\pgfqpoint{5.407074in}{0.697686in}}%
\pgfpathlineto{\pgfqpoint{5.410429in}{0.709570in}}%
\pgfpathlineto{\pgfqpoint{5.413785in}{0.708621in}}%
\pgfpathlineto{\pgfqpoint{5.417141in}{0.693885in}}%
\pgfpathlineto{\pgfqpoint{5.420497in}{0.625156in}}%
\pgfpathlineto{\pgfqpoint{5.423852in}{0.689013in}}%
\pgfpathlineto{\pgfqpoint{5.427208in}{0.717335in}}%
\pgfpathlineto{\pgfqpoint{5.430564in}{0.728201in}}%
\pgfpathlineto{\pgfqpoint{5.433920in}{0.726209in}}%
\pgfpathlineto{\pgfqpoint{5.437276in}{0.707745in}}%
\pgfpathlineto{\pgfqpoint{5.440631in}{0.645115in}}%
\pgfpathlineto{\pgfqpoint{5.443987in}{0.687966in}}%
\pgfpathlineto{\pgfqpoint{5.447343in}{0.713053in}}%
\pgfpathlineto{\pgfqpoint{5.450699in}{0.720407in}}%
\pgfpathlineto{\pgfqpoint{5.454054in}{0.717379in}}%
\pgfpathlineto{\pgfqpoint{5.457410in}{0.701160in}}%
\pgfpathlineto{\pgfqpoint{5.460766in}{0.625418in}}%
\pgfpathlineto{\pgfqpoint{5.464122in}{0.697422in}}%
\pgfpathlineto{\pgfqpoint{5.467478in}{0.724898in}}%
\pgfpathlineto{\pgfqpoint{5.470833in}{0.735621in}}%
\pgfpathlineto{\pgfqpoint{5.474189in}{0.734130in}}%
\pgfpathlineto{\pgfqpoint{5.477545in}{0.716929in}}%
\pgfpathlineto{\pgfqpoint{5.480901in}{0.664507in}}%
\pgfpathlineto{\pgfqpoint{5.487612in}{0.716945in}}%
\pgfpathlineto{\pgfqpoint{5.490968in}{0.724987in}}%
\pgfpathlineto{\pgfqpoint{5.494324in}{0.722132in}}%
\pgfpathlineto{\pgfqpoint{5.497680in}{0.705634in}}%
\pgfpathlineto{\pgfqpoint{5.501035in}{0.622046in}}%
\pgfpathlineto{\pgfqpoint{5.504391in}{0.704088in}}%
\pgfpathlineto{\pgfqpoint{5.507747in}{0.730849in}}%
\pgfpathlineto{\pgfqpoint{5.511103in}{0.741150in}}%
\pgfpathlineto{\pgfqpoint{5.514458in}{0.738930in}}%
\pgfpathlineto{\pgfqpoint{5.517814in}{0.722067in}}%
\pgfpathlineto{\pgfqpoint{5.521170in}{0.679390in}}%
\pgfpathlineto{\pgfqpoint{5.524526in}{0.676785in}}%
\pgfpathlineto{\pgfqpoint{5.527882in}{0.710124in}}%
\pgfpathlineto{\pgfqpoint{5.531237in}{0.717811in}}%
\pgfpathlineto{\pgfqpoint{5.534593in}{0.712276in}}%
\pgfpathlineto{\pgfqpoint{5.541305in}{0.659847in}}%
\pgfpathlineto{\pgfqpoint{5.544660in}{0.714958in}}%
\pgfpathlineto{\pgfqpoint{5.548016in}{0.733379in}}%
\pgfpathlineto{\pgfqpoint{5.551372in}{0.738267in}}%
\pgfpathlineto{\pgfqpoint{5.554728in}{0.728688in}}%
\pgfpathlineto{\pgfqpoint{5.558084in}{0.704697in}}%
\pgfpathlineto{\pgfqpoint{5.561439in}{0.619640in}}%
\pgfpathlineto{\pgfqpoint{5.564795in}{0.685867in}}%
\pgfpathlineto{\pgfqpoint{5.568151in}{0.704293in}}%
\pgfpathlineto{\pgfqpoint{5.571507in}{0.704487in}}%
\pgfpathlineto{\pgfqpoint{5.574863in}{0.687375in}}%
\pgfpathlineto{\pgfqpoint{5.578218in}{0.569971in}}%
\pgfpathlineto{\pgfqpoint{5.581574in}{0.697686in}}%
\pgfpathlineto{\pgfqpoint{5.584930in}{0.721430in}}%
\pgfpathlineto{\pgfqpoint{5.588286in}{0.729631in}}%
\pgfpathlineto{\pgfqpoint{5.591641in}{0.723507in}}%
\pgfpathlineto{\pgfqpoint{5.594997in}{0.702233in}}%
\pgfpathlineto{\pgfqpoint{5.598353in}{0.635039in}}%
\pgfpathlineto{\pgfqpoint{5.601709in}{0.679205in}}%
\pgfpathlineto{\pgfqpoint{5.605065in}{0.701570in}}%
\pgfpathlineto{\pgfqpoint{5.608420in}{0.705300in}}%
\pgfpathlineto{\pgfqpoint{5.611776in}{0.695536in}}%
\pgfpathlineto{\pgfqpoint{5.615132in}{0.655066in}}%
\pgfpathlineto{\pgfqpoint{5.618488in}{0.673291in}}%
\pgfpathlineto{\pgfqpoint{5.621843in}{0.709955in}}%
\pgfpathlineto{\pgfqpoint{5.625199in}{0.723922in}}%
\pgfpathlineto{\pgfqpoint{5.628555in}{0.724930in}}%
\pgfpathlineto{\pgfqpoint{5.631911in}{0.710007in}}%
\pgfpathlineto{\pgfqpoint{5.635267in}{0.666699in}}%
\pgfpathlineto{\pgfqpoint{5.638622in}{0.671863in}}%
\pgfpathlineto{\pgfqpoint{5.641978in}{0.705160in}}%
\pgfpathlineto{\pgfqpoint{5.645334in}{0.714929in}}%
\pgfpathlineto{\pgfqpoint{5.648690in}{0.713568in}}%
\pgfpathlineto{\pgfqpoint{5.652045in}{0.699746in}}%
\pgfpathlineto{\pgfqpoint{5.655401in}{0.645150in}}%
\pgfpathlineto{\pgfqpoint{5.658757in}{0.687272in}}%
\pgfpathlineto{\pgfqpoint{5.662113in}{0.718355in}}%
\pgfpathlineto{\pgfqpoint{5.665469in}{0.730540in}}%
\pgfpathlineto{\pgfqpoint{5.668824in}{0.730512in}}%
\pgfpathlineto{\pgfqpoint{5.672180in}{0.714614in}}%
\pgfpathlineto{\pgfqpoint{5.675536in}{0.664627in}}%
\pgfpathlineto{\pgfqpoint{5.678892in}{0.685846in}}%
\pgfpathlineto{\pgfqpoint{5.682247in}{0.715400in}}%
\pgfpathlineto{\pgfqpoint{5.685603in}{0.724527in}}%
\pgfpathlineto{\pgfqpoint{5.688959in}{0.723197in}}%
\pgfpathlineto{\pgfqpoint{5.692315in}{0.710137in}}%
\pgfpathlineto{\pgfqpoint{5.695671in}{0.662530in}}%
\pgfpathlineto{\pgfqpoint{5.702382in}{0.725447in}}%
\pgfpathlineto{\pgfqpoint{5.705738in}{0.738608in}}%
\pgfpathlineto{\pgfqpoint{5.709094in}{0.739826in}}%
\pgfpathlineto{\pgfqpoint{5.712449in}{0.726399in}}%
\pgfpathlineto{\pgfqpoint{5.715805in}{0.691275in}}%
\pgfpathlineto{\pgfqpoint{5.719161in}{0.667466in}}%
\pgfpathlineto{\pgfqpoint{5.722517in}{0.711157in}}%
\pgfpathlineto{\pgfqpoint{5.725873in}{0.722397in}}%
\pgfpathlineto{\pgfqpoint{5.729228in}{0.721035in}}%
\pgfpathlineto{\pgfqpoint{5.732584in}{0.705694in}}%
\pgfpathlineto{\pgfqpoint{5.735940in}{0.630070in}}%
\pgfpathlineto{\pgfqpoint{5.739296in}{0.702857in}}%
\pgfpathlineto{\pgfqpoint{5.742652in}{0.730363in}}%
\pgfpathlineto{\pgfqpoint{5.746007in}{0.740839in}}%
\pgfpathlineto{\pgfqpoint{5.749363in}{0.738363in}}%
\pgfpathlineto{\pgfqpoint{5.752719in}{0.721732in}}%
\pgfpathlineto{\pgfqpoint{5.756075in}{0.683092in}}%
\pgfpathlineto{\pgfqpoint{5.759430in}{0.663797in}}%
\pgfpathlineto{\pgfqpoint{5.762786in}{0.702377in}}%
\pgfpathlineto{\pgfqpoint{5.766142in}{0.709955in}}%
\pgfpathlineto{\pgfqpoint{5.769498in}{0.702444in}}%
\pgfpathlineto{\pgfqpoint{5.772854in}{0.666559in}}%
\pgfpathlineto{\pgfqpoint{5.776209in}{0.676186in}}%
\pgfpathlineto{\pgfqpoint{5.779565in}{0.715813in}}%
\pgfpathlineto{\pgfqpoint{5.782921in}{0.730681in}}%
\pgfpathlineto{\pgfqpoint{5.786277in}{0.732244in}}%
\pgfpathlineto{\pgfqpoint{5.789632in}{0.718801in}}%
\pgfpathlineto{\pgfqpoint{5.792988in}{0.688217in}}%
\pgfpathlineto{\pgfqpoint{5.796344in}{0.625189in}}%
\pgfpathlineto{\pgfqpoint{5.799700in}{0.690063in}}%
\pgfpathlineto{\pgfqpoint{5.803056in}{0.701983in}}%
\pgfpathlineto{\pgfqpoint{5.806411in}{0.698332in}}%
\pgfpathlineto{\pgfqpoint{5.809767in}{0.674059in}}%
\pgfpathlineto{\pgfqpoint{5.813123in}{0.643043in}}%
\pgfpathlineto{\pgfqpoint{5.816479in}{0.701337in}}%
\pgfpathlineto{\pgfqpoint{5.819834in}{0.720328in}}%
\pgfpathlineto{\pgfqpoint{5.823190in}{0.725414in}}%
\pgfpathlineto{\pgfqpoint{5.826546in}{0.715416in}}%
\pgfpathlineto{\pgfqpoint{5.829902in}{0.686815in}}%
\pgfpathlineto{\pgfqpoint{5.833258in}{0.626325in}}%
\pgfpathlineto{\pgfqpoint{5.836613in}{0.692846in}}%
\pgfpathlineto{\pgfqpoint{5.839969in}{0.707572in}}%
\pgfpathlineto{\pgfqpoint{5.843325in}{0.708784in}}%
\pgfpathlineto{\pgfqpoint{5.846681in}{0.697676in}}%
\pgfpathlineto{\pgfqpoint{5.850036in}{0.655662in}}%
\pgfpathlineto{\pgfqpoint{5.853392in}{0.675442in}}%
\pgfpathlineto{\pgfqpoint{5.856748in}{0.711406in}}%
\pgfpathlineto{\pgfqpoint{5.860104in}{0.725245in}}%
\pgfpathlineto{\pgfqpoint{5.863460in}{0.726568in}}%
\pgfpathlineto{\pgfqpoint{5.866815in}{0.712051in}}%
\pgfpathlineto{\pgfqpoint{5.870171in}{0.666578in}}%
\pgfpathlineto{\pgfqpoint{5.873527in}{0.679143in}}%
\pgfpathlineto{\pgfqpoint{5.876883in}{0.711179in}}%
\pgfpathlineto{\pgfqpoint{5.880238in}{0.721293in}}%
\pgfpathlineto{\pgfqpoint{5.883594in}{0.720949in}}%
\pgfpathlineto{\pgfqpoint{5.886950in}{0.709663in}}%
\pgfpathlineto{\pgfqpoint{5.890306in}{0.671494in}}%
\pgfpathlineto{\pgfqpoint{5.893662in}{0.680425in}}%
\pgfpathlineto{\pgfqpoint{5.897017in}{0.719329in}}%
\pgfpathlineto{\pgfqpoint{5.900373in}{0.734278in}}%
\pgfpathlineto{\pgfqpoint{5.903729in}{0.737250in}}%
\pgfpathlineto{\pgfqpoint{5.907085in}{0.725851in}}%
\pgfpathlineto{\pgfqpoint{5.910441in}{0.692710in}}%
\pgfpathlineto{\pgfqpoint{5.913796in}{0.665729in}}%
\pgfpathlineto{\pgfqpoint{5.917152in}{0.712469in}}%
\pgfpathlineto{\pgfqpoint{5.920508in}{0.725356in}}%
\pgfpathlineto{\pgfqpoint{5.923864in}{0.726267in}}%
\pgfpathlineto{\pgfqpoint{5.927219in}{0.715860in}}%
\pgfpathlineto{\pgfqpoint{5.930575in}{0.679405in}}%
\pgfpathlineto{\pgfqpoint{5.933931in}{0.684884in}}%
\pgfpathlineto{\pgfqpoint{5.937287in}{0.725055in}}%
\pgfpathlineto{\pgfqpoint{5.940643in}{0.740267in}}%
\pgfpathlineto{\pgfqpoint{5.943998in}{0.743176in}}%
\pgfpathlineto{\pgfqpoint{5.947354in}{0.731906in}}%
\pgfpathlineto{\pgfqpoint{5.950710in}{0.704351in}}%
\pgfpathlineto{\pgfqpoint{5.954066in}{0.607357in}}%
\pgfpathlineto{\pgfqpoint{5.957421in}{0.700626in}}%
\pgfpathlineto{\pgfqpoint{5.960777in}{0.715322in}}%
\pgfpathlineto{\pgfqpoint{5.964133in}{0.714501in}}%
\pgfpathlineto{\pgfqpoint{5.967489in}{0.697445in}}%
\pgfpathlineto{\pgfqpoint{5.970845in}{0.486993in}}%
\pgfpathlineto{\pgfqpoint{5.974200in}{0.704898in}}%
\pgfpathlineto{\pgfqpoint{5.977556in}{0.729250in}}%
\pgfpathlineto{\pgfqpoint{5.980912in}{0.737896in}}%
\pgfpathlineto{\pgfqpoint{5.984268in}{0.732615in}}%
\pgfpathlineto{\pgfqpoint{5.987623in}{0.713764in}}%
\pgfpathlineto{\pgfqpoint{5.990979in}{0.668649in}}%
\pgfpathlineto{\pgfqpoint{5.994335in}{0.666821in}}%
\pgfpathlineto{\pgfqpoint{5.997691in}{0.697485in}}%
\pgfpathlineto{\pgfqpoint{6.001047in}{0.701935in}}%
\pgfpathlineto{\pgfqpoint{6.004402in}{0.689690in}}%
\pgfpathlineto{\pgfqpoint{6.007758in}{0.621311in}}%
\pgfpathlineto{\pgfqpoint{6.011114in}{0.687421in}}%
\pgfpathlineto{\pgfqpoint{6.014470in}{0.716017in}}%
\pgfpathlineto{\pgfqpoint{6.017825in}{0.726717in}}%
\pgfpathlineto{\pgfqpoint{6.021181in}{0.723459in}}%
\pgfpathlineto{\pgfqpoint{6.024537in}{0.705037in}}%
\pgfpathlineto{\pgfqpoint{6.027893in}{0.655532in}}%
\pgfpathlineto{\pgfqpoint{6.031249in}{0.668966in}}%
\pgfpathlineto{\pgfqpoint{6.034604in}{0.697438in}}%
\pgfpathlineto{\pgfqpoint{6.037960in}{0.703588in}}%
\pgfpathlineto{\pgfqpoint{6.041316in}{0.696566in}}%
\pgfpathlineto{\pgfqpoint{6.044672in}{0.666311in}}%
\pgfpathlineto{\pgfqpoint{6.048027in}{0.656985in}}%
\pgfpathlineto{\pgfqpoint{6.051383in}{0.703533in}}%
\pgfpathlineto{\pgfqpoint{6.054739in}{0.720229in}}%
\pgfpathlineto{\pgfqpoint{6.058095in}{0.723742in}}%
\pgfpathlineto{\pgfqpoint{6.061451in}{0.711814in}}%
\pgfpathlineto{\pgfqpoint{6.064806in}{0.675744in}}%
\pgfpathlineto{\pgfqpoint{6.068162in}{0.662048in}}%
\pgfpathlineto{\pgfqpoint{6.071518in}{0.703014in}}%
\pgfpathlineto{\pgfqpoint{6.074874in}{0.715158in}}%
\pgfpathlineto{\pgfqpoint{6.078230in}{0.715993in}}%
\pgfpathlineto{\pgfqpoint{6.081585in}{0.706011in}}%
\pgfpathlineto{\pgfqpoint{6.084941in}{0.672250in}}%
\pgfpathlineto{\pgfqpoint{6.088297in}{0.669380in}}%
\pgfpathlineto{\pgfqpoint{6.091653in}{0.712558in}}%
\pgfpathlineto{\pgfqpoint{6.095008in}{0.728614in}}%
\pgfpathlineto{\pgfqpoint{6.098364in}{0.732402in}}%
\pgfpathlineto{\pgfqpoint{6.101720in}{0.721778in}}%
\pgfpathlineto{\pgfqpoint{6.105076in}{0.687915in}}%
\pgfpathlineto{\pgfqpoint{6.108432in}{0.668219in}}%
\pgfpathlineto{\pgfqpoint{6.111787in}{0.712461in}}%
\pgfpathlineto{\pgfqpoint{6.115143in}{0.725671in}}%
\pgfpathlineto{\pgfqpoint{6.118499in}{0.727587in}}%
\pgfpathlineto{\pgfqpoint{6.121855in}{0.719458in}}%
\pgfpathlineto{\pgfqpoint{6.128566in}{0.661748in}}%
\pgfpathlineto{\pgfqpoint{6.131922in}{0.718380in}}%
\pgfpathlineto{\pgfqpoint{6.135278in}{0.737165in}}%
\pgfpathlineto{\pgfqpoint{6.138634in}{0.743151in}}%
\pgfpathlineto{\pgfqpoint{6.141989in}{0.735942in}}%
\pgfpathlineto{\pgfqpoint{6.145345in}{0.712869in}}%
\pgfpathlineto{\pgfqpoint{6.148701in}{0.611311in}}%
\pgfpathlineto{\pgfqpoint{6.152057in}{0.700871in}}%
\pgfpathlineto{\pgfqpoint{6.155412in}{0.720389in}}%
\pgfpathlineto{\pgfqpoint{6.158768in}{0.723734in}}%
\pgfpathlineto{\pgfqpoint{6.162124in}{0.714728in}}%
\pgfpathlineto{\pgfqpoint{6.165480in}{0.680040in}}%
\pgfpathlineto{\pgfqpoint{6.168836in}{0.682778in}}%
\pgfpathlineto{\pgfqpoint{6.172191in}{0.724205in}}%
\pgfpathlineto{\pgfqpoint{6.175547in}{0.739646in}}%
\pgfpathlineto{\pgfqpoint{6.178903in}{0.742303in}}%
\pgfpathlineto{\pgfqpoint{6.182259in}{0.730590in}}%
\pgfpathlineto{\pgfqpoint{6.185614in}{0.704399in}}%
\pgfpathlineto{\pgfqpoint{6.188970in}{0.578343in}}%
\pgfpathlineto{\pgfqpoint{6.192326in}{0.690222in}}%
\pgfpathlineto{\pgfqpoint{6.195682in}{0.705666in}}%
\pgfpathlineto{\pgfqpoint{6.199038in}{0.703086in}}%
\pgfpathlineto{\pgfqpoint{6.205749in}{0.651586in}}%
\pgfpathlineto{\pgfqpoint{6.209105in}{0.708134in}}%
\pgfpathlineto{\pgfqpoint{6.212461in}{0.726906in}}%
\pgfpathlineto{\pgfqpoint{6.215816in}{0.731687in}}%
\pgfpathlineto{\pgfqpoint{6.219172in}{0.721494in}}%
\pgfpathlineto{\pgfqpoint{6.222528in}{0.696727in}}%
\pgfpathlineto{\pgfqpoint{6.225884in}{0.599601in}}%
\pgfpathlineto{\pgfqpoint{6.229240in}{0.680034in}}%
\pgfpathlineto{\pgfqpoint{6.232595in}{0.697404in}}%
\pgfpathlineto{\pgfqpoint{6.235951in}{0.696810in}}%
\pgfpathlineto{\pgfqpoint{6.239307in}{0.678002in}}%
\pgfpathlineto{\pgfqpoint{6.242663in}{0.603715in}}%
\pgfpathlineto{\pgfqpoint{6.246019in}{0.693281in}}%
\pgfpathlineto{\pgfqpoint{6.249374in}{0.715542in}}%
\pgfpathlineto{\pgfqpoint{6.252730in}{0.722825in}}%
\pgfpathlineto{\pgfqpoint{6.256086in}{0.715430in}}%
\pgfpathlineto{\pgfqpoint{6.259442in}{0.690326in}}%
\pgfpathlineto{\pgfqpoint{6.262797in}{0.582123in}}%
\pgfpathlineto{\pgfqpoint{6.266153in}{0.688685in}}%
\pgfpathlineto{\pgfqpoint{6.269509in}{0.706161in}}%
\pgfpathlineto{\pgfqpoint{6.272865in}{0.709337in}}%
\pgfpathlineto{\pgfqpoint{6.276221in}{0.701222in}}%
\pgfpathlineto{\pgfqpoint{6.279576in}{0.671664in}}%
\pgfpathlineto{\pgfqpoint{6.282932in}{0.655984in}}%
\pgfpathlineto{\pgfqpoint{6.286288in}{0.704979in}}%
\pgfpathlineto{\pgfqpoint{6.289644in}{0.722274in}}%
\pgfpathlineto{\pgfqpoint{6.292999in}{0.726738in}}%
\pgfpathlineto{\pgfqpoint{6.296355in}{0.716538in}}%
\pgfpathlineto{\pgfqpoint{6.299711in}{0.682299in}}%
\pgfpathlineto{\pgfqpoint{6.303067in}{0.665966in}}%
\pgfpathlineto{\pgfqpoint{6.306423in}{0.709164in}}%
\pgfpathlineto{\pgfqpoint{6.309778in}{0.722519in}}%
\pgfpathlineto{\pgfqpoint{6.313134in}{0.724876in}}%
\pgfpathlineto{\pgfqpoint{6.316490in}{0.717689in}}%
\pgfpathlineto{\pgfqpoint{6.319846in}{0.693937in}}%
\pgfpathlineto{\pgfqpoint{6.323201in}{0.639665in}}%
\pgfpathlineto{\pgfqpoint{6.326557in}{0.711462in}}%
\pgfpathlineto{\pgfqpoint{6.329913in}{0.732196in}}%
\pgfpathlineto{\pgfqpoint{6.333269in}{0.739585in}}%
\pgfpathlineto{\pgfqpoint{6.336625in}{0.734398in}}%
\pgfpathlineto{\pgfqpoint{6.339980in}{0.712464in}}%
\pgfpathlineto{\pgfqpoint{6.343336in}{0.603973in}}%
\pgfpathlineto{\pgfqpoint{6.346692in}{0.703894in}}%
\pgfpathlineto{\pgfqpoint{6.350048in}{0.724458in}}%
\pgfpathlineto{\pgfqpoint{6.353403in}{0.729691in}}%
\pgfpathlineto{\pgfqpoint{6.356759in}{0.724439in}}%
\pgfpathlineto{\pgfqpoint{6.360115in}{0.703418in}}%
\pgfpathlineto{\pgfqpoint{6.363471in}{0.623062in}}%
\pgfpathlineto{\pgfqpoint{6.366827in}{0.715671in}}%
\pgfpathlineto{\pgfqpoint{6.370182in}{0.737913in}}%
\pgfpathlineto{\pgfqpoint{6.373538in}{0.745868in}}%
\pgfpathlineto{\pgfqpoint{6.376894in}{0.740988in}}%
\pgfpathlineto{\pgfqpoint{6.383605in}{0.719746in}}%
\pgfpathlineto{\pgfqpoint{6.383605in}{0.719746in}}%
\pgfusepath{stroke}%
\end{pgfscope}%
\begin{pgfscope}%
\pgfpathrectangle{\pgfqpoint{4.621821in}{0.320625in}}{\pgfqpoint{1.845679in}{1.164375in}}%
\pgfusepath{clip}%
\pgfsetrectcap%
\pgfsetroundjoin%
\pgfsetlinewidth{0.602250pt}%
\definecolor{currentstroke}{rgb}{0.478431,0.243137,0.615686}%
\pgfsetstrokecolor{currentstroke}%
\pgfsetstrokeopacity{0.120000}%
\pgfsetdash{}{0pt}%
\pgfpathmoveto{\pgfqpoint{4.705715in}{1.468362in}}%
\pgfpathlineto{\pgfqpoint{4.715783in}{1.449591in}}%
\pgfpathlineto{\pgfqpoint{4.725850in}{1.428347in}}%
\pgfpathlineto{\pgfqpoint{4.739273in}{1.396827in}}%
\pgfpathlineto{\pgfqpoint{4.779543in}{1.300125in}}%
\pgfpathlineto{\pgfqpoint{4.803033in}{1.247570in}}%
\pgfpathlineto{\pgfqpoint{4.826524in}{1.197745in}}%
\pgfpathlineto{\pgfqpoint{4.853370in}{1.143258in}}%
\pgfpathlineto{\pgfqpoint{4.886928in}{1.077698in}}%
\pgfpathlineto{\pgfqpoint{4.927197in}{1.001516in}}%
\pgfpathlineto{\pgfqpoint{4.977534in}{0.908753in}}%
\pgfpathlineto{\pgfqpoint{5.017803in}{0.836030in}}%
\pgfpathlineto{\pgfqpoint{5.031226in}{0.811625in}}%
\pgfpathlineto{\pgfqpoint{5.044649in}{0.787776in}}%
\pgfpathlineto{\pgfqpoint{5.054717in}{0.771338in}}%
\pgfpathlineto{\pgfqpoint{5.064784in}{0.751660in}}%
\pgfpathlineto{\pgfqpoint{5.074851in}{0.731747in}}%
\pgfpathlineto{\pgfqpoint{5.081563in}{0.721660in}}%
\pgfpathlineto{\pgfqpoint{5.088274in}{0.715075in}}%
\pgfpathlineto{\pgfqpoint{5.091630in}{0.712042in}}%
\pgfpathlineto{\pgfqpoint{5.094986in}{0.707753in}}%
\pgfpathlineto{\pgfqpoint{5.098342in}{0.701273in}}%
\pgfpathlineto{\pgfqpoint{5.101698in}{0.693307in}}%
\pgfpathlineto{\pgfqpoint{5.108409in}{0.672874in}}%
\pgfpathlineto{\pgfqpoint{5.118476in}{0.638745in}}%
\pgfpathlineto{\pgfqpoint{5.121832in}{0.631145in}}%
\pgfpathlineto{\pgfqpoint{5.125188in}{0.625503in}}%
\pgfpathlineto{\pgfqpoint{5.128544in}{0.624086in}}%
\pgfpathlineto{\pgfqpoint{5.131900in}{0.631152in}}%
\pgfpathlineto{\pgfqpoint{5.138611in}{0.653426in}}%
\pgfpathlineto{\pgfqpoint{5.141967in}{0.661036in}}%
\pgfpathlineto{\pgfqpoint{5.145323in}{0.665049in}}%
\pgfpathlineto{\pgfqpoint{5.148678in}{0.664557in}}%
\pgfpathlineto{\pgfqpoint{5.152034in}{0.657210in}}%
\pgfpathlineto{\pgfqpoint{5.155390in}{0.646908in}}%
\pgfpathlineto{\pgfqpoint{5.158746in}{0.660521in}}%
\pgfpathlineto{\pgfqpoint{5.162102in}{0.680764in}}%
\pgfpathlineto{\pgfqpoint{5.165457in}{0.695082in}}%
\pgfpathlineto{\pgfqpoint{5.168813in}{0.704742in}}%
\pgfpathlineto{\pgfqpoint{5.172169in}{0.711047in}}%
\pgfpathlineto{\pgfqpoint{5.175525in}{0.714699in}}%
\pgfpathlineto{\pgfqpoint{5.178880in}{0.716011in}}%
\pgfpathlineto{\pgfqpoint{5.182236in}{0.715008in}}%
\pgfpathlineto{\pgfqpoint{5.185592in}{0.711411in}}%
\pgfpathlineto{\pgfqpoint{5.188948in}{0.704441in}}%
\pgfpathlineto{\pgfqpoint{5.192304in}{0.692129in}}%
\pgfpathlineto{\pgfqpoint{5.195659in}{0.668013in}}%
\pgfpathlineto{\pgfqpoint{5.199015in}{0.611711in}}%
\pgfpathlineto{\pgfqpoint{5.202371in}{0.674670in}}%
\pgfpathlineto{\pgfqpoint{5.205727in}{0.699160in}}%
\pgfpathlineto{\pgfqpoint{5.209082in}{0.713002in}}%
\pgfpathlineto{\pgfqpoint{5.212438in}{0.721552in}}%
\pgfpathlineto{\pgfqpoint{5.215794in}{0.726338in}}%
\pgfpathlineto{\pgfqpoint{5.219150in}{0.727629in}}%
\pgfpathlineto{\pgfqpoint{5.222506in}{0.724707in}}%
\pgfpathlineto{\pgfqpoint{5.225861in}{0.715594in}}%
\pgfpathlineto{\pgfqpoint{5.232573in}{0.684363in}}%
\pgfpathlineto{\pgfqpoint{5.235929in}{0.682761in}}%
\pgfpathlineto{\pgfqpoint{5.239285in}{0.687919in}}%
\pgfpathlineto{\pgfqpoint{5.242640in}{0.686438in}}%
\pgfpathlineto{\pgfqpoint{5.245996in}{0.673609in}}%
\pgfpathlineto{\pgfqpoint{5.249352in}{0.646728in}}%
\pgfpathlineto{\pgfqpoint{5.256063in}{0.696944in}}%
\pgfpathlineto{\pgfqpoint{5.259419in}{0.706563in}}%
\pgfpathlineto{\pgfqpoint{5.262775in}{0.703395in}}%
\pgfpathlineto{\pgfqpoint{5.266131in}{0.686197in}}%
\pgfpathlineto{\pgfqpoint{5.269487in}{0.650247in}}%
\pgfpathlineto{\pgfqpoint{5.272842in}{0.654411in}}%
\pgfpathlineto{\pgfqpoint{5.276198in}{0.677356in}}%
\pgfpathlineto{\pgfqpoint{5.279554in}{0.682976in}}%
\pgfpathlineto{\pgfqpoint{5.282910in}{0.675980in}}%
\pgfpathlineto{\pgfqpoint{5.286265in}{0.646858in}}%
\pgfpathlineto{\pgfqpoint{5.289621in}{0.636689in}}%
\pgfpathlineto{\pgfqpoint{5.292977in}{0.682309in}}%
\pgfpathlineto{\pgfqpoint{5.296333in}{0.699013in}}%
\pgfpathlineto{\pgfqpoint{5.299689in}{0.702519in}}%
\pgfpathlineto{\pgfqpoint{5.303044in}{0.688289in}}%
\pgfpathlineto{\pgfqpoint{5.306400in}{0.612358in}}%
\pgfpathlineto{\pgfqpoint{5.309756in}{0.692125in}}%
\pgfpathlineto{\pgfqpoint{5.313112in}{0.716327in}}%
\pgfpathlineto{\pgfqpoint{5.316467in}{0.726865in}}%
\pgfpathlineto{\pgfqpoint{5.319823in}{0.729854in}}%
\pgfpathlineto{\pgfqpoint{5.323179in}{0.726189in}}%
\pgfpathlineto{\pgfqpoint{5.326535in}{0.713570in}}%
\pgfpathlineto{\pgfqpoint{5.329891in}{0.677864in}}%
\pgfpathlineto{\pgfqpoint{5.333246in}{0.675667in}}%
\pgfpathlineto{\pgfqpoint{5.336602in}{0.718170in}}%
\pgfpathlineto{\pgfqpoint{5.339958in}{0.734807in}}%
\pgfpathlineto{\pgfqpoint{5.343314in}{0.740828in}}%
\pgfpathlineto{\pgfqpoint{5.346669in}{0.736162in}}%
\pgfpathlineto{\pgfqpoint{5.350025in}{0.719239in}}%
\pgfpathlineto{\pgfqpoint{5.356737in}{0.652138in}}%
\pgfpathlineto{\pgfqpoint{5.360093in}{0.692591in}}%
\pgfpathlineto{\pgfqpoint{5.363448in}{0.698428in}}%
\pgfpathlineto{\pgfqpoint{5.366804in}{0.685347in}}%
\pgfpathlineto{\pgfqpoint{5.370160in}{0.590672in}}%
\pgfpathlineto{\pgfqpoint{5.373516in}{0.690842in}}%
\pgfpathlineto{\pgfqpoint{5.376871in}{0.716898in}}%
\pgfpathlineto{\pgfqpoint{5.380227in}{0.726027in}}%
\pgfpathlineto{\pgfqpoint{5.383583in}{0.720251in}}%
\pgfpathlineto{\pgfqpoint{5.386939in}{0.700934in}}%
\pgfpathlineto{\pgfqpoint{5.390295in}{0.655992in}}%
\pgfpathlineto{\pgfqpoint{5.393650in}{0.652261in}}%
\pgfpathlineto{\pgfqpoint{5.397006in}{0.681793in}}%
\pgfpathlineto{\pgfqpoint{5.400362in}{0.684269in}}%
\pgfpathlineto{\pgfqpoint{5.403718in}{0.666556in}}%
\pgfpathlineto{\pgfqpoint{5.407074in}{0.599330in}}%
\pgfpathlineto{\pgfqpoint{5.410429in}{0.684071in}}%
\pgfpathlineto{\pgfqpoint{5.413785in}{0.706107in}}%
\pgfpathlineto{\pgfqpoint{5.417141in}{0.712797in}}%
\pgfpathlineto{\pgfqpoint{5.420497in}{0.703642in}}%
\pgfpathlineto{\pgfqpoint{5.427208in}{0.644964in}}%
\pgfpathlineto{\pgfqpoint{5.430564in}{0.692546in}}%
\pgfpathlineto{\pgfqpoint{5.433920in}{0.706233in}}%
\pgfpathlineto{\pgfqpoint{5.437276in}{0.708215in}}%
\pgfpathlineto{\pgfqpoint{5.440631in}{0.699830in}}%
\pgfpathlineto{\pgfqpoint{5.443987in}{0.671522in}}%
\pgfpathlineto{\pgfqpoint{5.447343in}{0.648176in}}%
\pgfpathlineto{\pgfqpoint{5.450699in}{0.700705in}}%
\pgfpathlineto{\pgfqpoint{5.454054in}{0.718647in}}%
\pgfpathlineto{\pgfqpoint{5.457410in}{0.723734in}}%
\pgfpathlineto{\pgfqpoint{5.460766in}{0.714288in}}%
\pgfpathlineto{\pgfqpoint{5.464122in}{0.674377in}}%
\pgfpathlineto{\pgfqpoint{5.467478in}{0.682340in}}%
\pgfpathlineto{\pgfqpoint{5.470833in}{0.718271in}}%
\pgfpathlineto{\pgfqpoint{5.474189in}{0.731184in}}%
\pgfpathlineto{\pgfqpoint{5.477545in}{0.734380in}}%
\pgfpathlineto{\pgfqpoint{5.480901in}{0.729384in}}%
\pgfpathlineto{\pgfqpoint{5.484256in}{0.712319in}}%
\pgfpathlineto{\pgfqpoint{5.487612in}{0.640884in}}%
\pgfpathlineto{\pgfqpoint{5.490968in}{0.705890in}}%
\pgfpathlineto{\pgfqpoint{5.494324in}{0.733915in}}%
\pgfpathlineto{\pgfqpoint{5.497680in}{0.745334in}}%
\pgfpathlineto{\pgfqpoint{5.501035in}{0.745703in}}%
\pgfpathlineto{\pgfqpoint{5.504391in}{0.733004in}}%
\pgfpathlineto{\pgfqpoint{5.507747in}{0.707351in}}%
\pgfpathlineto{\pgfqpoint{5.511103in}{0.625044in}}%
\pgfpathlineto{\pgfqpoint{5.514458in}{0.682188in}}%
\pgfpathlineto{\pgfqpoint{5.517814in}{0.697174in}}%
\pgfpathlineto{\pgfqpoint{5.521170in}{0.689419in}}%
\pgfpathlineto{\pgfqpoint{5.524526in}{0.632074in}}%
\pgfpathlineto{\pgfqpoint{5.527882in}{0.686200in}}%
\pgfpathlineto{\pgfqpoint{5.531237in}{0.716780in}}%
\pgfpathlineto{\pgfqpoint{5.534593in}{0.728142in}}%
\pgfpathlineto{\pgfqpoint{5.537949in}{0.724709in}}%
\pgfpathlineto{\pgfqpoint{5.541305in}{0.707746in}}%
\pgfpathlineto{\pgfqpoint{5.544660in}{0.672950in}}%
\pgfpathlineto{\pgfqpoint{5.548016in}{0.623165in}}%
\pgfpathlineto{\pgfqpoint{5.551372in}{0.674056in}}%
\pgfpathlineto{\pgfqpoint{5.554728in}{0.679317in}}%
\pgfpathlineto{\pgfqpoint{5.558084in}{0.661528in}}%
\pgfpathlineto{\pgfqpoint{5.561439in}{0.610503in}}%
\pgfpathlineto{\pgfqpoint{5.564795in}{0.684135in}}%
\pgfpathlineto{\pgfqpoint{5.568151in}{0.705389in}}%
\pgfpathlineto{\pgfqpoint{5.571507in}{0.711453in}}%
\pgfpathlineto{\pgfqpoint{5.574863in}{0.701294in}}%
\pgfpathlineto{\pgfqpoint{5.581574in}{0.641242in}}%
\pgfpathlineto{\pgfqpoint{5.584930in}{0.688517in}}%
\pgfpathlineto{\pgfqpoint{5.588286in}{0.701579in}}%
\pgfpathlineto{\pgfqpoint{5.591641in}{0.702609in}}%
\pgfpathlineto{\pgfqpoint{5.594997in}{0.692321in}}%
\pgfpathlineto{\pgfqpoint{5.598353in}{0.656224in}}%
\pgfpathlineto{\pgfqpoint{5.601709in}{0.660716in}}%
\pgfpathlineto{\pgfqpoint{5.605065in}{0.701142in}}%
\pgfpathlineto{\pgfqpoint{5.608420in}{0.716277in}}%
\pgfpathlineto{\pgfqpoint{5.611776in}{0.718867in}}%
\pgfpathlineto{\pgfqpoint{5.615132in}{0.704858in}}%
\pgfpathlineto{\pgfqpoint{5.618488in}{0.639517in}}%
\pgfpathlineto{\pgfqpoint{5.621843in}{0.693281in}}%
\pgfpathlineto{\pgfqpoint{5.625199in}{0.720167in}}%
\pgfpathlineto{\pgfqpoint{5.628555in}{0.730439in}}%
\pgfpathlineto{\pgfqpoint{5.631911in}{0.732013in}}%
\pgfpathlineto{\pgfqpoint{5.635267in}{0.725384in}}%
\pgfpathlineto{\pgfqpoint{5.638622in}{0.705189in}}%
\pgfpathlineto{\pgfqpoint{5.641978in}{0.398674in}}%
\pgfpathlineto{\pgfqpoint{5.645334in}{0.709631in}}%
\pgfpathlineto{\pgfqpoint{5.648690in}{0.733923in}}%
\pgfpathlineto{\pgfqpoint{5.652045in}{0.743704in}}%
\pgfpathlineto{\pgfqpoint{5.655401in}{0.742442in}}%
\pgfpathlineto{\pgfqpoint{5.658757in}{0.728167in}}%
\pgfpathlineto{\pgfqpoint{5.662113in}{0.699015in}}%
\pgfpathlineto{\pgfqpoint{5.665469in}{0.580977in}}%
\pgfpathlineto{\pgfqpoint{5.668824in}{0.688048in}}%
\pgfpathlineto{\pgfqpoint{5.672180in}{0.698921in}}%
\pgfpathlineto{\pgfqpoint{5.675536in}{0.689408in}}%
\pgfpathlineto{\pgfqpoint{5.678892in}{0.626108in}}%
\pgfpathlineto{\pgfqpoint{5.682247in}{0.687591in}}%
\pgfpathlineto{\pgfqpoint{5.685603in}{0.717019in}}%
\pgfpathlineto{\pgfqpoint{5.688959in}{0.727823in}}%
\pgfpathlineto{\pgfqpoint{5.692315in}{0.723844in}}%
\pgfpathlineto{\pgfqpoint{5.695671in}{0.706249in}}%
\pgfpathlineto{\pgfqpoint{5.702382in}{0.637082in}}%
\pgfpathlineto{\pgfqpoint{5.705738in}{0.678234in}}%
\pgfpathlineto{\pgfqpoint{5.709094in}{0.682739in}}%
\pgfpathlineto{\pgfqpoint{5.712449in}{0.665976in}}%
\pgfpathlineto{\pgfqpoint{5.715805in}{0.596313in}}%
\pgfpathlineto{\pgfqpoint{5.719161in}{0.683764in}}%
\pgfpathlineto{\pgfqpoint{5.722517in}{0.706076in}}%
\pgfpathlineto{\pgfqpoint{5.725873in}{0.712892in}}%
\pgfpathlineto{\pgfqpoint{5.729228in}{0.703809in}}%
\pgfpathlineto{\pgfqpoint{5.732584in}{0.673335in}}%
\pgfpathlineto{\pgfqpoint{5.735940in}{0.635852in}}%
\pgfpathlineto{\pgfqpoint{5.739296in}{0.688801in}}%
\pgfpathlineto{\pgfqpoint{5.742652in}{0.702841in}}%
\pgfpathlineto{\pgfqpoint{5.746007in}{0.704541in}}%
\pgfpathlineto{\pgfqpoint{5.749363in}{0.695187in}}%
\pgfpathlineto{\pgfqpoint{5.752719in}{0.662636in}}%
\pgfpathlineto{\pgfqpoint{5.756075in}{0.656583in}}%
\pgfpathlineto{\pgfqpoint{5.759430in}{0.701047in}}%
\pgfpathlineto{\pgfqpoint{5.762786in}{0.717233in}}%
\pgfpathlineto{\pgfqpoint{5.766142in}{0.720740in}}%
\pgfpathlineto{\pgfqpoint{5.769498in}{0.708436in}}%
\pgfpathlineto{\pgfqpoint{5.772854in}{0.656614in}}%
\pgfpathlineto{\pgfqpoint{5.779565in}{0.718466in}}%
\pgfpathlineto{\pgfqpoint{5.782921in}{0.729597in}}%
\pgfpathlineto{\pgfqpoint{5.786277in}{0.731589in}}%
\pgfpathlineto{\pgfqpoint{5.789632in}{0.725228in}}%
\pgfpathlineto{\pgfqpoint{5.792988in}{0.705305in}}%
\pgfpathlineto{\pgfqpoint{5.796344in}{0.522861in}}%
\pgfpathlineto{\pgfqpoint{5.799700in}{0.709327in}}%
\pgfpathlineto{\pgfqpoint{5.803056in}{0.733751in}}%
\pgfpathlineto{\pgfqpoint{5.806411in}{0.743505in}}%
\pgfpathlineto{\pgfqpoint{5.809767in}{0.742023in}}%
\pgfpathlineto{\pgfqpoint{5.813123in}{0.727426in}}%
\pgfpathlineto{\pgfqpoint{5.816479in}{0.697389in}}%
\pgfpathlineto{\pgfqpoint{5.819834in}{0.604836in}}%
\pgfpathlineto{\pgfqpoint{5.823190in}{0.689579in}}%
\pgfpathlineto{\pgfqpoint{5.826546in}{0.699735in}}%
\pgfpathlineto{\pgfqpoint{5.829902in}{0.689967in}}%
\pgfpathlineto{\pgfqpoint{5.833258in}{0.626476in}}%
\pgfpathlineto{\pgfqpoint{5.836613in}{0.687976in}}%
\pgfpathlineto{\pgfqpoint{5.839969in}{0.717349in}}%
\pgfpathlineto{\pgfqpoint{5.843325in}{0.728124in}}%
\pgfpathlineto{\pgfqpoint{5.846681in}{0.724144in}}%
\pgfpathlineto{\pgfqpoint{5.850036in}{0.706477in}}%
\pgfpathlineto{\pgfqpoint{5.853392in}{0.668259in}}%
\pgfpathlineto{\pgfqpoint{5.856748in}{0.639692in}}%
\pgfpathlineto{\pgfqpoint{5.860104in}{0.679586in}}%
\pgfpathlineto{\pgfqpoint{5.863460in}{0.684086in}}%
\pgfpathlineto{\pgfqpoint{5.866815in}{0.667818in}}%
\pgfpathlineto{\pgfqpoint{5.870171in}{0.588160in}}%
\pgfpathlineto{\pgfqpoint{5.873527in}{0.683773in}}%
\pgfpathlineto{\pgfqpoint{5.876883in}{0.706559in}}%
\pgfpathlineto{\pgfqpoint{5.880238in}{0.713693in}}%
\pgfpathlineto{\pgfqpoint{5.883594in}{0.705065in}}%
\pgfpathlineto{\pgfqpoint{5.886950in}{0.675679in}}%
\pgfpathlineto{\pgfqpoint{5.890306in}{0.631174in}}%
\pgfpathlineto{\pgfqpoint{5.893662in}{0.688313in}}%
\pgfpathlineto{\pgfqpoint{5.897017in}{0.702858in}}%
\pgfpathlineto{\pgfqpoint{5.900373in}{0.704783in}}%
\pgfpathlineto{\pgfqpoint{5.903729in}{0.695609in}}%
\pgfpathlineto{\pgfqpoint{5.907085in}{0.663468in}}%
\pgfpathlineto{\pgfqpoint{5.910441in}{0.656227in}}%
\pgfpathlineto{\pgfqpoint{5.913796in}{0.701210in}}%
\pgfpathlineto{\pgfqpoint{5.917152in}{0.717509in}}%
\pgfpathlineto{\pgfqpoint{5.920508in}{0.721075in}}%
\pgfpathlineto{\pgfqpoint{5.923864in}{0.708874in}}%
\pgfpathlineto{\pgfqpoint{5.927219in}{0.658704in}}%
\pgfpathlineto{\pgfqpoint{5.933931in}{0.717691in}}%
\pgfpathlineto{\pgfqpoint{5.937287in}{0.728878in}}%
\pgfpathlineto{\pgfqpoint{5.940643in}{0.730801in}}%
\pgfpathlineto{\pgfqpoint{5.943998in}{0.724237in}}%
\pgfpathlineto{\pgfqpoint{5.947354in}{0.703678in}}%
\pgfpathlineto{\pgfqpoint{5.950710in}{0.570303in}}%
\pgfpathlineto{\pgfqpoint{5.954066in}{0.710011in}}%
\pgfpathlineto{\pgfqpoint{5.957421in}{0.733723in}}%
\pgfpathlineto{\pgfqpoint{5.960777in}{0.743066in}}%
\pgfpathlineto{\pgfqpoint{5.964133in}{0.740990in}}%
\pgfpathlineto{\pgfqpoint{5.967489in}{0.725766in}}%
\pgfpathlineto{\pgfqpoint{5.970845in}{0.693993in}}%
\pgfpathlineto{\pgfqpoint{5.974200in}{0.627066in}}%
\pgfpathlineto{\pgfqpoint{5.977556in}{0.691662in}}%
\pgfpathlineto{\pgfqpoint{5.980912in}{0.700571in}}%
\pgfpathlineto{\pgfqpoint{5.984268in}{0.690163in}}%
\pgfpathlineto{\pgfqpoint{5.987623in}{0.623774in}}%
\pgfpathlineto{\pgfqpoint{5.990979in}{0.688866in}}%
\pgfpathlineto{\pgfqpoint{5.994335in}{0.717772in}}%
\pgfpathlineto{\pgfqpoint{5.997691in}{0.728323in}}%
\pgfpathlineto{\pgfqpoint{6.001047in}{0.724125in}}%
\pgfpathlineto{\pgfqpoint{6.004402in}{0.706192in}}%
\pgfpathlineto{\pgfqpoint{6.007758in}{0.666655in}}%
\pgfpathlineto{\pgfqpoint{6.011114in}{0.644138in}}%
\pgfpathlineto{\pgfqpoint{6.014470in}{0.681406in}}%
\pgfpathlineto{\pgfqpoint{6.017825in}{0.685651in}}%
\pgfpathlineto{\pgfqpoint{6.021181in}{0.669706in}}%
\pgfpathlineto{\pgfqpoint{6.024537in}{0.578713in}}%
\pgfpathlineto{\pgfqpoint{6.027893in}{0.683939in}}%
\pgfpathlineto{\pgfqpoint{6.031249in}{0.707141in}}%
\pgfpathlineto{\pgfqpoint{6.034604in}{0.714553in}}%
\pgfpathlineto{\pgfqpoint{6.037960in}{0.706337in}}%
\pgfpathlineto{\pgfqpoint{6.041316in}{0.677841in}}%
\pgfpathlineto{\pgfqpoint{6.044672in}{0.626805in}}%
\pgfpathlineto{\pgfqpoint{6.048027in}{0.688149in}}%
\pgfpathlineto{\pgfqpoint{6.051383in}{0.703161in}}%
\pgfpathlineto{\pgfqpoint{6.054739in}{0.705312in}}%
\pgfpathlineto{\pgfqpoint{6.058095in}{0.696355in}}%
\pgfpathlineto{\pgfqpoint{6.061451in}{0.664812in}}%
\pgfpathlineto{\pgfqpoint{6.064806in}{0.655653in}}%
\pgfpathlineto{\pgfqpoint{6.068162in}{0.701475in}}%
\pgfpathlineto{\pgfqpoint{6.071518in}{0.717964in}}%
\pgfpathlineto{\pgfqpoint{6.074874in}{0.721665in}}%
\pgfpathlineto{\pgfqpoint{6.078230in}{0.709710in}}%
\pgfpathlineto{\pgfqpoint{6.081585in}{0.661835in}}%
\pgfpathlineto{\pgfqpoint{6.084941in}{0.685569in}}%
\pgfpathlineto{\pgfqpoint{6.088297in}{0.716895in}}%
\pgfpathlineto{\pgfqpoint{6.091653in}{0.728228in}}%
\pgfpathlineto{\pgfqpoint{6.095008in}{0.730131in}}%
\pgfpathlineto{\pgfqpoint{6.098364in}{0.723404in}}%
\pgfpathlineto{\pgfqpoint{6.101720in}{0.702263in}}%
\pgfpathlineto{\pgfqpoint{6.105076in}{0.598595in}}%
\pgfpathlineto{\pgfqpoint{6.108432in}{0.710679in}}%
\pgfpathlineto{\pgfqpoint{6.111787in}{0.733771in}}%
\pgfpathlineto{\pgfqpoint{6.115143in}{0.742731in}}%
\pgfpathlineto{\pgfqpoint{6.118499in}{0.740065in}}%
\pgfpathlineto{\pgfqpoint{6.121855in}{0.720843in}}%
\pgfpathlineto{\pgfqpoint{6.125210in}{0.643411in}}%
\pgfpathlineto{\pgfqpoint{6.128566in}{0.707566in}}%
\pgfpathlineto{\pgfqpoint{6.131922in}{0.730731in}}%
\pgfpathlineto{\pgfqpoint{6.135278in}{0.737474in}}%
\pgfpathlineto{\pgfqpoint{6.138634in}{0.734027in}}%
\pgfpathlineto{\pgfqpoint{6.141989in}{0.717178in}}%
\pgfpathlineto{\pgfqpoint{6.145345in}{0.632875in}}%
\pgfpathlineto{\pgfqpoint{6.148701in}{0.715370in}}%
\pgfpathlineto{\pgfqpoint{6.152057in}{0.741978in}}%
\pgfpathlineto{\pgfqpoint{6.155412in}{0.752151in}}%
\pgfpathlineto{\pgfqpoint{6.158768in}{0.749758in}}%
\pgfpathlineto{\pgfqpoint{6.162124in}{0.733761in}}%
\pgfpathlineto{\pgfqpoint{6.165480in}{0.699715in}}%
\pgfpathlineto{\pgfqpoint{6.168836in}{0.651940in}}%
\pgfpathlineto{\pgfqpoint{6.172191in}{0.703585in}}%
\pgfpathlineto{\pgfqpoint{6.175547in}{0.711107in}}%
\pgfpathlineto{\pgfqpoint{6.178903in}{0.699725in}}%
\pgfpathlineto{\pgfqpoint{6.182259in}{0.625755in}}%
\pgfpathlineto{\pgfqpoint{6.185614in}{0.700422in}}%
\pgfpathlineto{\pgfqpoint{6.188970in}{0.728375in}}%
\pgfpathlineto{\pgfqpoint{6.192326in}{0.738459in}}%
\pgfpathlineto{\pgfqpoint{6.195682in}{0.733781in}}%
\pgfpathlineto{\pgfqpoint{6.199038in}{0.716594in}}%
\pgfpathlineto{\pgfqpoint{6.202393in}{0.683074in}}%
\pgfpathlineto{\pgfqpoint{6.205749in}{0.613671in}}%
\pgfpathlineto{\pgfqpoint{6.209105in}{0.675804in}}%
\pgfpathlineto{\pgfqpoint{6.212461in}{0.679675in}}%
\pgfpathlineto{\pgfqpoint{6.215816in}{0.653273in}}%
\pgfpathlineto{\pgfqpoint{6.219172in}{0.650936in}}%
\pgfpathlineto{\pgfqpoint{6.222528in}{0.697298in}}%
\pgfpathlineto{\pgfqpoint{6.225884in}{0.714252in}}%
\pgfpathlineto{\pgfqpoint{6.229240in}{0.716767in}}%
\pgfpathlineto{\pgfqpoint{6.232595in}{0.702958in}}%
\pgfpathlineto{\pgfqpoint{6.239307in}{0.635137in}}%
\pgfpathlineto{\pgfqpoint{6.242663in}{0.682362in}}%
\pgfpathlineto{\pgfqpoint{6.246019in}{0.692754in}}%
\pgfpathlineto{\pgfqpoint{6.249374in}{0.688928in}}%
\pgfpathlineto{\pgfqpoint{6.252730in}{0.665942in}}%
\pgfpathlineto{\pgfqpoint{6.256086in}{0.623854in}}%
\pgfpathlineto{\pgfqpoint{6.259442in}{0.688633in}}%
\pgfpathlineto{\pgfqpoint{6.262797in}{0.708349in}}%
\pgfpathlineto{\pgfqpoint{6.266153in}{0.713805in}}%
\pgfpathlineto{\pgfqpoint{6.269509in}{0.703448in}}%
\pgfpathlineto{\pgfqpoint{6.272865in}{0.662984in}}%
\pgfpathlineto{\pgfqpoint{6.276221in}{0.669746in}}%
\pgfpathlineto{\pgfqpoint{6.279576in}{0.705312in}}%
\pgfpathlineto{\pgfqpoint{6.282932in}{0.717620in}}%
\pgfpathlineto{\pgfqpoint{6.286288in}{0.719944in}}%
\pgfpathlineto{\pgfqpoint{6.289644in}{0.713431in}}%
\pgfpathlineto{\pgfqpoint{6.292999in}{0.692372in}}%
\pgfpathlineto{\pgfqpoint{6.296355in}{0.590605in}}%
\pgfpathlineto{\pgfqpoint{6.299711in}{0.701075in}}%
\pgfpathlineto{\pgfqpoint{6.303067in}{0.724085in}}%
\pgfpathlineto{\pgfqpoint{6.306423in}{0.732893in}}%
\pgfpathlineto{\pgfqpoint{6.309778in}{0.729807in}}%
\pgfpathlineto{\pgfqpoint{6.313134in}{0.707509in}}%
\pgfpathlineto{\pgfqpoint{6.316490in}{0.624597in}}%
\pgfpathlineto{\pgfqpoint{6.319846in}{0.714063in}}%
\pgfpathlineto{\pgfqpoint{6.323201in}{0.733544in}}%
\pgfpathlineto{\pgfqpoint{6.326557in}{0.740092in}}%
\pgfpathlineto{\pgfqpoint{6.329913in}{0.738054in}}%
\pgfpathlineto{\pgfqpoint{6.333269in}{0.725734in}}%
\pgfpathlineto{\pgfqpoint{6.336625in}{0.686410in}}%
\pgfpathlineto{\pgfqpoint{6.339980in}{0.696593in}}%
\pgfpathlineto{\pgfqpoint{6.343336in}{0.734889in}}%
\pgfpathlineto{\pgfqpoint{6.346692in}{0.749805in}}%
\pgfpathlineto{\pgfqpoint{6.350048in}{0.753203in}}%
\pgfpathlineto{\pgfqpoint{6.353403in}{0.743965in}}%
\pgfpathlineto{\pgfqpoint{6.356759in}{0.727162in}}%
\pgfpathlineto{\pgfqpoint{6.360115in}{0.701756in}}%
\pgfpathlineto{\pgfqpoint{6.363471in}{0.659769in}}%
\pgfpathlineto{\pgfqpoint{6.366827in}{0.559813in}}%
\pgfpathlineto{\pgfqpoint{6.370182in}{0.610316in}}%
\pgfpathlineto{\pgfqpoint{6.373538in}{0.670952in}}%
\pgfpathlineto{\pgfqpoint{6.376894in}{0.698737in}}%
\pgfpathlineto{\pgfqpoint{6.380250in}{0.711653in}}%
\pgfpathlineto{\pgfqpoint{6.383605in}{0.710640in}}%
\pgfpathlineto{\pgfqpoint{6.383605in}{0.710640in}}%
\pgfusepath{stroke}%
\end{pgfscope}%
\begin{pgfscope}%
\pgfpathrectangle{\pgfqpoint{4.621821in}{0.320625in}}{\pgfqpoint{1.845679in}{1.164375in}}%
\pgfusepath{clip}%
\pgfsetrectcap%
\pgfsetroundjoin%
\pgfsetlinewidth{0.602250pt}%
\definecolor{currentstroke}{rgb}{0.478431,0.243137,0.615686}%
\pgfsetstrokecolor{currentstroke}%
\pgfsetstrokeopacity{0.120000}%
\pgfsetdash{}{0pt}%
\pgfpathmoveto{\pgfqpoint{4.705715in}{1.468904in}}%
\pgfpathlineto{\pgfqpoint{4.715783in}{1.450600in}}%
\pgfpathlineto{\pgfqpoint{4.725850in}{1.429890in}}%
\pgfpathlineto{\pgfqpoint{4.739273in}{1.399049in}}%
\pgfpathlineto{\pgfqpoint{4.779543in}{1.304329in}}%
\pgfpathlineto{\pgfqpoint{4.803033in}{1.252902in}}%
\pgfpathlineto{\pgfqpoint{4.829879in}{1.197114in}}%
\pgfpathlineto{\pgfqpoint{4.863437in}{1.130244in}}%
\pgfpathlineto{\pgfqpoint{4.903707in}{1.052533in}}%
\pgfpathlineto{\pgfqpoint{4.954043in}{0.957770in}}%
\pgfpathlineto{\pgfqpoint{5.014447in}{0.846328in}}%
\pgfpathlineto{\pgfqpoint{5.031226in}{0.815563in}}%
\pgfpathlineto{\pgfqpoint{5.044649in}{0.791484in}}%
\pgfpathlineto{\pgfqpoint{5.054717in}{0.774394in}}%
\pgfpathlineto{\pgfqpoint{5.068140in}{0.747436in}}%
\pgfpathlineto{\pgfqpoint{5.074851in}{0.734927in}}%
\pgfpathlineto{\pgfqpoint{5.081563in}{0.725536in}}%
\pgfpathlineto{\pgfqpoint{5.091630in}{0.714424in}}%
\pgfpathlineto{\pgfqpoint{5.094986in}{0.708222in}}%
\pgfpathlineto{\pgfqpoint{5.101698in}{0.691663in}}%
\pgfpathlineto{\pgfqpoint{5.108409in}{0.670113in}}%
\pgfpathlineto{\pgfqpoint{5.115121in}{0.647781in}}%
\pgfpathlineto{\pgfqpoint{5.118476in}{0.639558in}}%
\pgfpathlineto{\pgfqpoint{5.121832in}{0.635074in}}%
\pgfpathlineto{\pgfqpoint{5.125188in}{0.635966in}}%
\pgfpathlineto{\pgfqpoint{5.128544in}{0.642373in}}%
\pgfpathlineto{\pgfqpoint{5.135255in}{0.658896in}}%
\pgfpathlineto{\pgfqpoint{5.138611in}{0.664239in}}%
\pgfpathlineto{\pgfqpoint{5.141967in}{0.666070in}}%
\pgfpathlineto{\pgfqpoint{5.145323in}{0.662713in}}%
\pgfpathlineto{\pgfqpoint{5.148678in}{0.652627in}}%
\pgfpathlineto{\pgfqpoint{5.152034in}{0.638927in}}%
\pgfpathlineto{\pgfqpoint{5.155390in}{0.643217in}}%
\pgfpathlineto{\pgfqpoint{5.158746in}{0.660709in}}%
\pgfpathlineto{\pgfqpoint{5.162102in}{0.674040in}}%
\pgfpathlineto{\pgfqpoint{5.165457in}{0.682802in}}%
\pgfpathlineto{\pgfqpoint{5.168813in}{0.688123in}}%
\pgfpathlineto{\pgfqpoint{5.172169in}{0.690642in}}%
\pgfpathlineto{\pgfqpoint{5.175525in}{0.690580in}}%
\pgfpathlineto{\pgfqpoint{5.178880in}{0.687790in}}%
\pgfpathlineto{\pgfqpoint{5.182236in}{0.681628in}}%
\pgfpathlineto{\pgfqpoint{5.185592in}{0.670368in}}%
\pgfpathlineto{\pgfqpoint{5.188948in}{0.648535in}}%
\pgfpathlineto{\pgfqpoint{5.192304in}{0.586399in}}%
\pgfpathlineto{\pgfqpoint{5.195659in}{0.646886in}}%
\pgfpathlineto{\pgfqpoint{5.199015in}{0.673583in}}%
\pgfpathlineto{\pgfqpoint{5.202371in}{0.688146in}}%
\pgfpathlineto{\pgfqpoint{5.205727in}{0.696956in}}%
\pgfpathlineto{\pgfqpoint{5.209082in}{0.701721in}}%
\pgfpathlineto{\pgfqpoint{5.212438in}{0.702685in}}%
\pgfpathlineto{\pgfqpoint{5.215794in}{0.698908in}}%
\pgfpathlineto{\pgfqpoint{5.219150in}{0.688047in}}%
\pgfpathlineto{\pgfqpoint{5.222506in}{0.673757in}}%
\pgfpathlineto{\pgfqpoint{5.225861in}{0.678659in}}%
\pgfpathlineto{\pgfqpoint{5.229217in}{0.692743in}}%
\pgfpathlineto{\pgfqpoint{5.232573in}{0.701076in}}%
\pgfpathlineto{\pgfqpoint{5.235929in}{0.703266in}}%
\pgfpathlineto{\pgfqpoint{5.239285in}{0.699242in}}%
\pgfpathlineto{\pgfqpoint{5.242640in}{0.686566in}}%
\pgfpathlineto{\pgfqpoint{5.245996in}{0.652882in}}%
\pgfpathlineto{\pgfqpoint{5.249352in}{0.652370in}}%
\pgfpathlineto{\pgfqpoint{5.252708in}{0.691631in}}%
\pgfpathlineto{\pgfqpoint{5.256063in}{0.708238in}}%
\pgfpathlineto{\pgfqpoint{5.259419in}{0.714799in}}%
\pgfpathlineto{\pgfqpoint{5.262775in}{0.711698in}}%
\pgfpathlineto{\pgfqpoint{5.269487in}{0.664897in}}%
\pgfpathlineto{\pgfqpoint{5.272842in}{0.711908in}}%
\pgfpathlineto{\pgfqpoint{5.276198in}{0.730453in}}%
\pgfpathlineto{\pgfqpoint{5.279554in}{0.738679in}}%
\pgfpathlineto{\pgfqpoint{5.282910in}{0.740229in}}%
\pgfpathlineto{\pgfqpoint{5.286265in}{0.735335in}}%
\pgfpathlineto{\pgfqpoint{5.289621in}{0.721066in}}%
\pgfpathlineto{\pgfqpoint{5.292977in}{0.679051in}}%
\pgfpathlineto{\pgfqpoint{5.296333in}{0.693239in}}%
\pgfpathlineto{\pgfqpoint{5.299689in}{0.730370in}}%
\pgfpathlineto{\pgfqpoint{5.303044in}{0.745836in}}%
\pgfpathlineto{\pgfqpoint{5.306400in}{0.751425in}}%
\pgfpathlineto{\pgfqpoint{5.309756in}{0.746846in}}%
\pgfpathlineto{\pgfqpoint{5.313112in}{0.731791in}}%
\pgfpathlineto{\pgfqpoint{5.316467in}{0.705115in}}%
\pgfpathlineto{\pgfqpoint{5.319823in}{0.651441in}}%
\pgfpathlineto{\pgfqpoint{5.323179in}{0.677642in}}%
\pgfpathlineto{\pgfqpoint{5.326535in}{0.686558in}}%
\pgfpathlineto{\pgfqpoint{5.333246in}{0.647664in}}%
\pgfpathlineto{\pgfqpoint{5.336602in}{0.701461in}}%
\pgfpathlineto{\pgfqpoint{5.339958in}{0.721022in}}%
\pgfpathlineto{\pgfqpoint{5.343314in}{0.726364in}}%
\pgfpathlineto{\pgfqpoint{5.346669in}{0.716708in}}%
\pgfpathlineto{\pgfqpoint{5.350025in}{0.696545in}}%
\pgfpathlineto{\pgfqpoint{5.353381in}{0.656655in}}%
\pgfpathlineto{\pgfqpoint{5.356737in}{0.633577in}}%
\pgfpathlineto{\pgfqpoint{5.360093in}{0.662916in}}%
\pgfpathlineto{\pgfqpoint{5.363448in}{0.659341in}}%
\pgfpathlineto{\pgfqpoint{5.366804in}{0.608660in}}%
\pgfpathlineto{\pgfqpoint{5.370160in}{0.658160in}}%
\pgfpathlineto{\pgfqpoint{5.373516in}{0.689811in}}%
\pgfpathlineto{\pgfqpoint{5.376871in}{0.702197in}}%
\pgfpathlineto{\pgfqpoint{5.380227in}{0.700322in}}%
\pgfpathlineto{\pgfqpoint{5.383583in}{0.680161in}}%
\pgfpathlineto{\pgfqpoint{5.386939in}{0.596467in}}%
\pgfpathlineto{\pgfqpoint{5.390295in}{0.671969in}}%
\pgfpathlineto{\pgfqpoint{5.393650in}{0.693936in}}%
\pgfpathlineto{\pgfqpoint{5.397006in}{0.700740in}}%
\pgfpathlineto{\pgfqpoint{5.400362in}{0.698012in}}%
\pgfpathlineto{\pgfqpoint{5.403718in}{0.683390in}}%
\pgfpathlineto{\pgfqpoint{5.407074in}{0.627857in}}%
\pgfpathlineto{\pgfqpoint{5.410429in}{0.670607in}}%
\pgfpathlineto{\pgfqpoint{5.413785in}{0.701499in}}%
\pgfpathlineto{\pgfqpoint{5.417141in}{0.713823in}}%
\pgfpathlineto{\pgfqpoint{5.420497in}{0.714485in}}%
\pgfpathlineto{\pgfqpoint{5.423852in}{0.696859in}}%
\pgfpathlineto{\pgfqpoint{5.427208in}{0.601933in}}%
\pgfpathlineto{\pgfqpoint{5.430564in}{0.704777in}}%
\pgfpathlineto{\pgfqpoint{5.433920in}{0.727175in}}%
\pgfpathlineto{\pgfqpoint{5.437276in}{0.736701in}}%
\pgfpathlineto{\pgfqpoint{5.440631in}{0.738777in}}%
\pgfpathlineto{\pgfqpoint{5.443987in}{0.733950in}}%
\pgfpathlineto{\pgfqpoint{5.447343in}{0.719102in}}%
\pgfpathlineto{\pgfqpoint{5.450699in}{0.672291in}}%
\pgfpathlineto{\pgfqpoint{5.454054in}{0.697221in}}%
\pgfpathlineto{\pgfqpoint{5.457410in}{0.731571in}}%
\pgfpathlineto{\pgfqpoint{5.460766in}{0.745933in}}%
\pgfpathlineto{\pgfqpoint{5.464122in}{0.750188in}}%
\pgfpathlineto{\pgfqpoint{5.467478in}{0.743142in}}%
\pgfpathlineto{\pgfqpoint{5.470833in}{0.725680in}}%
\pgfpathlineto{\pgfqpoint{5.474189in}{0.692520in}}%
\pgfpathlineto{\pgfqpoint{5.477545in}{0.618714in}}%
\pgfpathlineto{\pgfqpoint{5.480901in}{0.682198in}}%
\pgfpathlineto{\pgfqpoint{5.484256in}{0.685525in}}%
\pgfpathlineto{\pgfqpoint{5.487612in}{0.655297in}}%
\pgfpathlineto{\pgfqpoint{5.490968in}{0.665308in}}%
\pgfpathlineto{\pgfqpoint{5.494324in}{0.707267in}}%
\pgfpathlineto{\pgfqpoint{5.497680in}{0.723327in}}%
\pgfpathlineto{\pgfqpoint{5.501035in}{0.725386in}}%
\pgfpathlineto{\pgfqpoint{5.504391in}{0.712728in}}%
\pgfpathlineto{\pgfqpoint{5.507747in}{0.689082in}}%
\pgfpathlineto{\pgfqpoint{5.511103in}{0.633300in}}%
\pgfpathlineto{\pgfqpoint{5.514458in}{0.645035in}}%
\pgfpathlineto{\pgfqpoint{5.517814in}{0.664400in}}%
\pgfpathlineto{\pgfqpoint{5.521170in}{0.654056in}}%
\pgfpathlineto{\pgfqpoint{5.524526in}{0.540370in}}%
\pgfpathlineto{\pgfqpoint{5.527882in}{0.669408in}}%
\pgfpathlineto{\pgfqpoint{5.531237in}{0.694519in}}%
\pgfpathlineto{\pgfqpoint{5.534593in}{0.703440in}}%
\pgfpathlineto{\pgfqpoint{5.537949in}{0.697078in}}%
\pgfpathlineto{\pgfqpoint{5.541305in}{0.670767in}}%
\pgfpathlineto{\pgfqpoint{5.544660in}{0.607457in}}%
\pgfpathlineto{\pgfqpoint{5.548016in}{0.678545in}}%
\pgfpathlineto{\pgfqpoint{5.551372in}{0.695372in}}%
\pgfpathlineto{\pgfqpoint{5.554728in}{0.699399in}}%
\pgfpathlineto{\pgfqpoint{5.558084in}{0.693655in}}%
\pgfpathlineto{\pgfqpoint{5.561439in}{0.672684in}}%
\pgfpathlineto{\pgfqpoint{5.564795in}{0.583609in}}%
\pgfpathlineto{\pgfqpoint{5.568151in}{0.683340in}}%
\pgfpathlineto{\pgfqpoint{5.571507in}{0.705944in}}%
\pgfpathlineto{\pgfqpoint{5.574863in}{0.714287in}}%
\pgfpathlineto{\pgfqpoint{5.578218in}{0.710003in}}%
\pgfpathlineto{\pgfqpoint{5.581574in}{0.680294in}}%
\pgfpathlineto{\pgfqpoint{5.584930in}{0.666691in}}%
\pgfpathlineto{\pgfqpoint{5.588286in}{0.712999in}}%
\pgfpathlineto{\pgfqpoint{5.591641in}{0.729593in}}%
\pgfpathlineto{\pgfqpoint{5.594997in}{0.736154in}}%
\pgfpathlineto{\pgfqpoint{5.598353in}{0.735722in}}%
\pgfpathlineto{\pgfqpoint{5.601709in}{0.727691in}}%
\pgfpathlineto{\pgfqpoint{5.605065in}{0.705892in}}%
\pgfpathlineto{\pgfqpoint{5.608420in}{0.596563in}}%
\pgfpathlineto{\pgfqpoint{5.611776in}{0.713297in}}%
\pgfpathlineto{\pgfqpoint{5.615132in}{0.736816in}}%
\pgfpathlineto{\pgfqpoint{5.618488in}{0.746802in}}%
\pgfpathlineto{\pgfqpoint{5.621843in}{0.746884in}}%
\pgfpathlineto{\pgfqpoint{5.625199in}{0.734947in}}%
\pgfpathlineto{\pgfqpoint{5.628555in}{0.711921in}}%
\pgfpathlineto{\pgfqpoint{5.631911in}{0.657042in}}%
\pgfpathlineto{\pgfqpoint{5.635267in}{0.668474in}}%
\pgfpathlineto{\pgfqpoint{5.638622in}{0.689533in}}%
\pgfpathlineto{\pgfqpoint{5.641978in}{0.682700in}}%
\pgfpathlineto{\pgfqpoint{5.645334in}{0.617141in}}%
\pgfpathlineto{\pgfqpoint{5.648690in}{0.684663in}}%
\pgfpathlineto{\pgfqpoint{5.652045in}{0.714247in}}%
\pgfpathlineto{\pgfqpoint{5.655401in}{0.725483in}}%
\pgfpathlineto{\pgfqpoint{5.658757in}{0.722181in}}%
\pgfpathlineto{\pgfqpoint{5.662113in}{0.706061in}}%
\pgfpathlineto{\pgfqpoint{5.665469in}{0.675339in}}%
\pgfpathlineto{\pgfqpoint{5.668824in}{0.561035in}}%
\pgfpathlineto{\pgfqpoint{5.672180in}{0.661271in}}%
\pgfpathlineto{\pgfqpoint{5.675536in}{0.668631in}}%
\pgfpathlineto{\pgfqpoint{5.678892in}{0.649489in}}%
\pgfpathlineto{\pgfqpoint{5.682247in}{0.617711in}}%
\pgfpathlineto{\pgfqpoint{5.685603in}{0.679462in}}%
\pgfpathlineto{\pgfqpoint{5.688959in}{0.699641in}}%
\pgfpathlineto{\pgfqpoint{5.692315in}{0.705226in}}%
\pgfpathlineto{\pgfqpoint{5.695671in}{0.694618in}}%
\pgfpathlineto{\pgfqpoint{5.699026in}{0.660219in}}%
\pgfpathlineto{\pgfqpoint{5.702382in}{0.642257in}}%
\pgfpathlineto{\pgfqpoint{5.705738in}{0.685579in}}%
\pgfpathlineto{\pgfqpoint{5.709094in}{0.698520in}}%
\pgfpathlineto{\pgfqpoint{5.712449in}{0.700173in}}%
\pgfpathlineto{\pgfqpoint{5.715805in}{0.691601in}}%
\pgfpathlineto{\pgfqpoint{5.719161in}{0.663145in}}%
\pgfpathlineto{\pgfqpoint{5.722517in}{0.640047in}}%
\pgfpathlineto{\pgfqpoint{5.725873in}{0.692469in}}%
\pgfpathlineto{\pgfqpoint{5.729228in}{0.710514in}}%
\pgfpathlineto{\pgfqpoint{5.732584in}{0.715919in}}%
\pgfpathlineto{\pgfqpoint{5.735940in}{0.707166in}}%
\pgfpathlineto{\pgfqpoint{5.739296in}{0.664040in}}%
\pgfpathlineto{\pgfqpoint{5.742652in}{0.683494in}}%
\pgfpathlineto{\pgfqpoint{5.746007in}{0.717663in}}%
\pgfpathlineto{\pgfqpoint{5.749363in}{0.731112in}}%
\pgfpathlineto{\pgfqpoint{5.752719in}{0.735621in}}%
\pgfpathlineto{\pgfqpoint{5.756075in}{0.733074in}}%
\pgfpathlineto{\pgfqpoint{5.759430in}{0.721736in}}%
\pgfpathlineto{\pgfqpoint{5.762786in}{0.689856in}}%
\pgfpathlineto{\pgfqpoint{5.766142in}{0.675285in}}%
\pgfpathlineto{\pgfqpoint{5.769498in}{0.723029in}}%
\pgfpathlineto{\pgfqpoint{5.772854in}{0.740696in}}%
\pgfpathlineto{\pgfqpoint{5.776209in}{0.747169in}}%
\pgfpathlineto{\pgfqpoint{5.779565in}{0.742689in}}%
\pgfpathlineto{\pgfqpoint{5.782921in}{0.726774in}}%
\pgfpathlineto{\pgfqpoint{5.786277in}{0.696069in}}%
\pgfpathlineto{\pgfqpoint{5.789632in}{0.589568in}}%
\pgfpathlineto{\pgfqpoint{5.792988in}{0.683679in}}%
\pgfpathlineto{\pgfqpoint{5.796344in}{0.691651in}}%
\pgfpathlineto{\pgfqpoint{5.799700in}{0.675084in}}%
\pgfpathlineto{\pgfqpoint{5.803056in}{0.622118in}}%
\pgfpathlineto{\pgfqpoint{5.806411in}{0.697945in}}%
\pgfpathlineto{\pgfqpoint{5.809767in}{0.719771in}}%
\pgfpathlineto{\pgfqpoint{5.813123in}{0.726538in}}%
\pgfpathlineto{\pgfqpoint{5.816479in}{0.718207in}}%
\pgfpathlineto{\pgfqpoint{5.819834in}{0.698284in}}%
\pgfpathlineto{\pgfqpoint{5.823190in}{0.655879in}}%
\pgfpathlineto{\pgfqpoint{5.826546in}{0.636914in}}%
\pgfpathlineto{\pgfqpoint{5.829902in}{0.669647in}}%
\pgfpathlineto{\pgfqpoint{5.833258in}{0.669131in}}%
\pgfpathlineto{\pgfqpoint{5.836613in}{0.637068in}}%
\pgfpathlineto{\pgfqpoint{5.839969in}{0.647049in}}%
\pgfpathlineto{\pgfqpoint{5.843325in}{0.688294in}}%
\pgfpathlineto{\pgfqpoint{5.846681in}{0.703916in}}%
\pgfpathlineto{\pgfqpoint{5.850036in}{0.705585in}}%
\pgfpathlineto{\pgfqpoint{5.853392in}{0.690089in}}%
\pgfpathlineto{\pgfqpoint{5.856748in}{0.640520in}}%
\pgfpathlineto{\pgfqpoint{5.860104in}{0.661411in}}%
\pgfpathlineto{\pgfqpoint{5.863460in}{0.691236in}}%
\pgfpathlineto{\pgfqpoint{5.866815in}{0.700585in}}%
\pgfpathlineto{\pgfqpoint{5.870171in}{0.699565in}}%
\pgfpathlineto{\pgfqpoint{5.873527in}{0.687155in}}%
\pgfpathlineto{\pgfqpoint{5.876883in}{0.643508in}}%
\pgfpathlineto{\pgfqpoint{5.880238in}{0.665030in}}%
\pgfpathlineto{\pgfqpoint{5.883594in}{0.700314in}}%
\pgfpathlineto{\pgfqpoint{5.886950in}{0.714114in}}%
\pgfpathlineto{\pgfqpoint{5.890306in}{0.715979in}}%
\pgfpathlineto{\pgfqpoint{5.893662in}{0.701035in}}%
\pgfpathlineto{\pgfqpoint{5.897017in}{0.617114in}}%
\pgfpathlineto{\pgfqpoint{5.900373in}{0.696629in}}%
\pgfpathlineto{\pgfqpoint{5.903729in}{0.721979in}}%
\pgfpathlineto{\pgfqpoint{5.907085in}{0.732272in}}%
\pgfpathlineto{\pgfqpoint{5.910441in}{0.734483in}}%
\pgfpathlineto{\pgfqpoint{5.913796in}{0.729327in}}%
\pgfpathlineto{\pgfqpoint{5.917152in}{0.713232in}}%
\pgfpathlineto{\pgfqpoint{5.920508in}{0.656216in}}%
\pgfpathlineto{\pgfqpoint{5.923864in}{0.699666in}}%
\pgfpathlineto{\pgfqpoint{5.927219in}{0.730453in}}%
\pgfpathlineto{\pgfqpoint{5.930575in}{0.743399in}}%
\pgfpathlineto{\pgfqpoint{5.933931in}{0.746079in}}%
\pgfpathlineto{\pgfqpoint{5.937287in}{0.736651in}}%
\pgfpathlineto{\pgfqpoint{5.940643in}{0.716022in}}%
\pgfpathlineto{\pgfqpoint{5.943998in}{0.670081in}}%
\pgfpathlineto{\pgfqpoint{5.947354in}{0.663000in}}%
\pgfpathlineto{\pgfqpoint{5.950710in}{0.691546in}}%
\pgfpathlineto{\pgfqpoint{5.954066in}{0.690474in}}%
\pgfpathlineto{\pgfqpoint{5.957421in}{0.659291in}}%
\pgfpathlineto{\pgfqpoint{5.960777in}{0.665701in}}%
\pgfpathlineto{\pgfqpoint{5.964133in}{0.707971in}}%
\pgfpathlineto{\pgfqpoint{5.967489in}{0.723836in}}%
\pgfpathlineto{\pgfqpoint{5.970845in}{0.725772in}}%
\pgfpathlineto{\pgfqpoint{5.974200in}{0.712890in}}%
\pgfpathlineto{\pgfqpoint{5.977556in}{0.687609in}}%
\pgfpathlineto{\pgfqpoint{5.980912in}{0.612988in}}%
\pgfpathlineto{\pgfqpoint{5.984268in}{0.658829in}}%
\pgfpathlineto{\pgfqpoint{5.987623in}{0.674684in}}%
\pgfpathlineto{\pgfqpoint{5.990979in}{0.666894in}}%
\pgfpathlineto{\pgfqpoint{5.994335in}{0.606250in}}%
\pgfpathlineto{\pgfqpoint{5.997691in}{0.665602in}}%
\pgfpathlineto{\pgfqpoint{6.001047in}{0.695762in}}%
\pgfpathlineto{\pgfqpoint{6.004402in}{0.707199in}}%
\pgfpathlineto{\pgfqpoint{6.007758in}{0.704209in}}%
\pgfpathlineto{\pgfqpoint{6.011114in}{0.683228in}}%
\pgfpathlineto{\pgfqpoint{6.014470in}{0.583379in}}%
\pgfpathlineto{\pgfqpoint{6.017825in}{0.674488in}}%
\pgfpathlineto{\pgfqpoint{6.021181in}{0.695910in}}%
\pgfpathlineto{\pgfqpoint{6.024537in}{0.701972in}}%
\pgfpathlineto{\pgfqpoint{6.027893in}{0.698024in}}%
\pgfpathlineto{\pgfqpoint{6.031249in}{0.680434in}}%
\pgfpathlineto{\pgfqpoint{6.034604in}{0.582630in}}%
\pgfpathlineto{\pgfqpoint{6.037960in}{0.680940in}}%
\pgfpathlineto{\pgfqpoint{6.041316in}{0.706807in}}%
\pgfpathlineto{\pgfqpoint{6.044672in}{0.716799in}}%
\pgfpathlineto{\pgfqpoint{6.048027in}{0.714516in}}%
\pgfpathlineto{\pgfqpoint{6.051383in}{0.691549in}}%
\pgfpathlineto{\pgfqpoint{6.054739in}{0.636228in}}%
\pgfpathlineto{\pgfqpoint{6.058095in}{0.706156in}}%
\pgfpathlineto{\pgfqpoint{6.061451in}{0.725398in}}%
\pgfpathlineto{\pgfqpoint{6.064806in}{0.732837in}}%
\pgfpathlineto{\pgfqpoint{6.068162in}{0.732653in}}%
\pgfpathlineto{\pgfqpoint{6.071518in}{0.724362in}}%
\pgfpathlineto{\pgfqpoint{6.074874in}{0.701029in}}%
\pgfpathlineto{\pgfqpoint{6.078230in}{0.630755in}}%
\pgfpathlineto{\pgfqpoint{6.081585in}{0.714053in}}%
\pgfpathlineto{\pgfqpoint{6.084941in}{0.735986in}}%
\pgfpathlineto{\pgfqpoint{6.088297in}{0.744865in}}%
\pgfpathlineto{\pgfqpoint{6.091653in}{0.743162in}}%
\pgfpathlineto{\pgfqpoint{6.095008in}{0.729108in}}%
\pgfpathlineto{\pgfqpoint{6.098364in}{0.701628in}}%
\pgfpathlineto{\pgfqpoint{6.101720in}{0.586916in}}%
\pgfpathlineto{\pgfqpoint{6.105076in}{0.682421in}}%
\pgfpathlineto{\pgfqpoint{6.108432in}{0.695017in}}%
\pgfpathlineto{\pgfqpoint{6.111787in}{0.685625in}}%
\pgfpathlineto{\pgfqpoint{6.115143in}{0.616050in}}%
\pgfpathlineto{\pgfqpoint{6.118499in}{0.686796in}}%
\pgfpathlineto{\pgfqpoint{6.121855in}{0.715606in}}%
\pgfpathlineto{\pgfqpoint{6.125210in}{0.726384in}}%
\pgfpathlineto{\pgfqpoint{6.128566in}{0.722691in}}%
\pgfpathlineto{\pgfqpoint{6.131922in}{0.705945in}}%
\pgfpathlineto{\pgfqpoint{6.135278in}{0.672203in}}%
\pgfpathlineto{\pgfqpoint{6.138634in}{0.614035in}}%
\pgfpathlineto{\pgfqpoint{6.141989in}{0.670599in}}%
\pgfpathlineto{\pgfqpoint{6.145345in}{0.676875in}}%
\pgfpathlineto{\pgfqpoint{6.148701in}{0.660826in}}%
\pgfpathlineto{\pgfqpoint{6.152057in}{0.591301in}}%
\pgfpathlineto{\pgfqpoint{6.155412in}{0.679275in}}%
\pgfpathlineto{\pgfqpoint{6.158768in}{0.701924in}}%
\pgfpathlineto{\pgfqpoint{6.162124in}{0.709242in}}%
\pgfpathlineto{\pgfqpoint{6.165480in}{0.701102in}}%
\pgfpathlineto{\pgfqpoint{6.168836in}{0.672608in}}%
\pgfpathlineto{\pgfqpoint{6.172191in}{0.623571in}}%
\pgfpathlineto{\pgfqpoint{6.175547in}{0.684127in}}%
\pgfpathlineto{\pgfqpoint{6.178903in}{0.699521in}}%
\pgfpathlineto{\pgfqpoint{6.182259in}{0.702429in}}%
\pgfpathlineto{\pgfqpoint{6.185614in}{0.695070in}}%
\pgfpathlineto{\pgfqpoint{6.188970in}{0.669546in}}%
\pgfpathlineto{\pgfqpoint{6.192326in}{0.631639in}}%
\pgfpathlineto{\pgfqpoint{6.195682in}{0.692664in}}%
\pgfpathlineto{\pgfqpoint{6.199038in}{0.712057in}}%
\pgfpathlineto{\pgfqpoint{6.202393in}{0.718308in}}%
\pgfpathlineto{\pgfqpoint{6.205749in}{0.710814in}}%
\pgfpathlineto{\pgfqpoint{6.209105in}{0.675604in}}%
\pgfpathlineto{\pgfqpoint{6.212461in}{0.673246in}}%
\pgfpathlineto{\pgfqpoint{6.215816in}{0.713470in}}%
\pgfpathlineto{\pgfqpoint{6.219172in}{0.727983in}}%
\pgfpathlineto{\pgfqpoint{6.222528in}{0.732699in}}%
\pgfpathlineto{\pgfqpoint{6.225884in}{0.729882in}}%
\pgfpathlineto{\pgfqpoint{6.229240in}{0.717527in}}%
\pgfpathlineto{\pgfqpoint{6.232595in}{0.680633in}}%
\pgfpathlineto{\pgfqpoint{6.235951in}{0.683211in}}%
\pgfpathlineto{\pgfqpoint{6.239307in}{0.724019in}}%
\pgfpathlineto{\pgfqpoint{6.242663in}{0.739963in}}%
\pgfpathlineto{\pgfqpoint{6.246019in}{0.744900in}}%
\pgfpathlineto{\pgfqpoint{6.249374in}{0.737954in}}%
\pgfpathlineto{\pgfqpoint{6.252730in}{0.719393in}}%
\pgfpathlineto{\pgfqpoint{6.256086in}{0.679723in}}%
\pgfpathlineto{\pgfqpoint{6.259442in}{0.654671in}}%
\pgfpathlineto{\pgfqpoint{6.262797in}{0.692054in}}%
\pgfpathlineto{\pgfqpoint{6.266153in}{0.695183in}}%
\pgfpathlineto{\pgfqpoint{6.269509in}{0.675605in}}%
\pgfpathlineto{\pgfqpoint{6.272865in}{0.633746in}}%
\pgfpathlineto{\pgfqpoint{6.276221in}{0.700723in}}%
\pgfpathlineto{\pgfqpoint{6.279576in}{0.721300in}}%
\pgfpathlineto{\pgfqpoint{6.282932in}{0.727207in}}%
\pgfpathlineto{\pgfqpoint{6.286288in}{0.718047in}}%
\pgfpathlineto{\pgfqpoint{6.289644in}{0.696676in}}%
\pgfpathlineto{\pgfqpoint{6.292999in}{0.645776in}}%
\pgfpathlineto{\pgfqpoint{6.296355in}{0.652555in}}%
\pgfpathlineto{\pgfqpoint{6.299711in}{0.677551in}}%
\pgfpathlineto{\pgfqpoint{6.303067in}{0.676448in}}%
\pgfpathlineto{\pgfqpoint{6.306423in}{0.648483in}}%
\pgfpathlineto{\pgfqpoint{6.309778in}{0.643431in}}%
\pgfpathlineto{\pgfqpoint{6.313134in}{0.689896in}}%
\pgfpathlineto{\pgfqpoint{6.316490in}{0.706788in}}%
\pgfpathlineto{\pgfqpoint{6.319846in}{0.709738in}}%
\pgfpathlineto{\pgfqpoint{6.323201in}{0.696107in}}%
\pgfpathlineto{\pgfqpoint{6.326557in}{0.654864in}}%
\pgfpathlineto{\pgfqpoint{6.329913in}{0.656370in}}%
\pgfpathlineto{\pgfqpoint{6.333269in}{0.691427in}}%
\pgfpathlineto{\pgfqpoint{6.336625in}{0.702093in}}%
\pgfpathlineto{\pgfqpoint{6.339980in}{0.701835in}}%
\pgfpathlineto{\pgfqpoint{6.343336in}{0.690256in}}%
\pgfpathlineto{\pgfqpoint{6.346692in}{0.649868in}}%
\pgfpathlineto{\pgfqpoint{6.350048in}{0.664544in}}%
\pgfpathlineto{\pgfqpoint{6.353403in}{0.701723in}}%
\pgfpathlineto{\pgfqpoint{6.356759in}{0.716089in}}%
\pgfpathlineto{\pgfqpoint{6.360115in}{0.718370in}}%
\pgfpathlineto{\pgfqpoint{6.363471in}{0.708673in}}%
\pgfpathlineto{\pgfqpoint{6.373538in}{0.674296in}}%
\pgfpathlineto{\pgfqpoint{6.376894in}{0.665632in}}%
\pgfpathlineto{\pgfqpoint{6.383605in}{0.651079in}}%
\pgfpathlineto{\pgfqpoint{6.383605in}{0.651079in}}%
\pgfusepath{stroke}%
\end{pgfscope}%
\begin{pgfscope}%
\pgfpathrectangle{\pgfqpoint{4.621821in}{0.320625in}}{\pgfqpoint{1.845679in}{1.164375in}}%
\pgfusepath{clip}%
\pgfsetrectcap%
\pgfsetroundjoin%
\pgfsetlinewidth{0.602250pt}%
\definecolor{currentstroke}{rgb}{0.478431,0.243137,0.615686}%
\pgfsetstrokecolor{currentstroke}%
\pgfsetstrokeopacity{0.120000}%
\pgfsetdash{}{0pt}%
\pgfpathmoveto{\pgfqpoint{4.705715in}{1.468747in}}%
\pgfpathlineto{\pgfqpoint{4.715783in}{1.450186in}}%
\pgfpathlineto{\pgfqpoint{4.725850in}{1.429202in}}%
\pgfpathlineto{\pgfqpoint{4.739273in}{1.398027in}}%
\pgfpathlineto{\pgfqpoint{4.776187in}{1.310285in}}%
\pgfpathlineto{\pgfqpoint{4.796322in}{1.265658in}}%
\pgfpathlineto{\pgfqpoint{4.819812in}{1.216530in}}%
\pgfpathlineto{\pgfqpoint{4.846658in}{1.163239in}}%
\pgfpathlineto{\pgfqpoint{4.876860in}{1.105745in}}%
\pgfpathlineto{\pgfqpoint{4.917130in}{1.031652in}}%
\pgfpathlineto{\pgfqpoint{4.967466in}{0.941491in}}%
\pgfpathlineto{\pgfqpoint{5.041293in}{0.812060in}}%
\pgfpathlineto{\pgfqpoint{5.054717in}{0.789122in}}%
\pgfpathlineto{\pgfqpoint{5.071496in}{0.758771in}}%
\pgfpathlineto{\pgfqpoint{5.078207in}{0.748966in}}%
\pgfpathlineto{\pgfqpoint{5.084919in}{0.740051in}}%
\pgfpathlineto{\pgfqpoint{5.091630in}{0.728202in}}%
\pgfpathlineto{\pgfqpoint{5.098342in}{0.713428in}}%
\pgfpathlineto{\pgfqpoint{5.108409in}{0.689407in}}%
\pgfpathlineto{\pgfqpoint{5.111765in}{0.682773in}}%
\pgfpathlineto{\pgfqpoint{5.115121in}{0.677957in}}%
\pgfpathlineto{\pgfqpoint{5.118476in}{0.675445in}}%
\pgfpathlineto{\pgfqpoint{5.121832in}{0.675182in}}%
\pgfpathlineto{\pgfqpoint{5.131900in}{0.678568in}}%
\pgfpathlineto{\pgfqpoint{5.135255in}{0.676550in}}%
\pgfpathlineto{\pgfqpoint{5.138611in}{0.670331in}}%
\pgfpathlineto{\pgfqpoint{5.141967in}{0.660729in}}%
\pgfpathlineto{\pgfqpoint{5.148678in}{0.635955in}}%
\pgfpathlineto{\pgfqpoint{5.152034in}{0.636366in}}%
\pgfpathlineto{\pgfqpoint{5.158746in}{0.653251in}}%
\pgfpathlineto{\pgfqpoint{5.162102in}{0.657458in}}%
\pgfpathlineto{\pgfqpoint{5.165457in}{0.658215in}}%
\pgfpathlineto{\pgfqpoint{5.168813in}{0.655373in}}%
\pgfpathlineto{\pgfqpoint{5.172169in}{0.648007in}}%
\pgfpathlineto{\pgfqpoint{5.175525in}{0.633154in}}%
\pgfpathlineto{\pgfqpoint{5.178880in}{0.597748in}}%
\pgfpathlineto{\pgfqpoint{5.182236in}{0.597228in}}%
\pgfpathlineto{\pgfqpoint{5.185592in}{0.637804in}}%
\pgfpathlineto{\pgfqpoint{5.188948in}{0.656651in}}%
\pgfpathlineto{\pgfqpoint{5.192304in}{0.667716in}}%
\pgfpathlineto{\pgfqpoint{5.195659in}{0.673965in}}%
\pgfpathlineto{\pgfqpoint{5.199015in}{0.676044in}}%
\pgfpathlineto{\pgfqpoint{5.202371in}{0.673211in}}%
\pgfpathlineto{\pgfqpoint{5.205727in}{0.662860in}}%
\pgfpathlineto{\pgfqpoint{5.209082in}{0.661104in}}%
\pgfpathlineto{\pgfqpoint{5.212438in}{0.682026in}}%
\pgfpathlineto{\pgfqpoint{5.215794in}{0.698616in}}%
\pgfpathlineto{\pgfqpoint{5.219150in}{0.708881in}}%
\pgfpathlineto{\pgfqpoint{5.222506in}{0.714488in}}%
\pgfpathlineto{\pgfqpoint{5.225861in}{0.716247in}}%
\pgfpathlineto{\pgfqpoint{5.229217in}{0.714221in}}%
\pgfpathlineto{\pgfqpoint{5.232573in}{0.707610in}}%
\pgfpathlineto{\pgfqpoint{5.235929in}{0.693684in}}%
\pgfpathlineto{\pgfqpoint{5.239285in}{0.660384in}}%
\pgfpathlineto{\pgfqpoint{5.242640in}{0.646472in}}%
\pgfpathlineto{\pgfqpoint{5.245996in}{0.693448in}}%
\pgfpathlineto{\pgfqpoint{5.249352in}{0.712587in}}%
\pgfpathlineto{\pgfqpoint{5.252708in}{0.722721in}}%
\pgfpathlineto{\pgfqpoint{5.256063in}{0.726835in}}%
\pgfpathlineto{\pgfqpoint{5.259419in}{0.724583in}}%
\pgfpathlineto{\pgfqpoint{5.262775in}{0.712685in}}%
\pgfpathlineto{\pgfqpoint{5.269487in}{0.660894in}}%
\pgfpathlineto{\pgfqpoint{5.272842in}{0.690063in}}%
\pgfpathlineto{\pgfqpoint{5.276198in}{0.701762in}}%
\pgfpathlineto{\pgfqpoint{5.279554in}{0.701057in}}%
\pgfpathlineto{\pgfqpoint{5.282910in}{0.686745in}}%
\pgfpathlineto{\pgfqpoint{5.286265in}{0.627928in}}%
\pgfpathlineto{\pgfqpoint{5.289621in}{0.681660in}}%
\pgfpathlineto{\pgfqpoint{5.292977in}{0.709913in}}%
\pgfpathlineto{\pgfqpoint{5.296333in}{0.720869in}}%
\pgfpathlineto{\pgfqpoint{5.299689in}{0.719022in}}%
\pgfpathlineto{\pgfqpoint{5.303044in}{0.701825in}}%
\pgfpathlineto{\pgfqpoint{5.306400in}{0.657394in}}%
\pgfpathlineto{\pgfqpoint{5.309756in}{0.670288in}}%
\pgfpathlineto{\pgfqpoint{5.313112in}{0.697381in}}%
\pgfpathlineto{\pgfqpoint{5.316467in}{0.704558in}}%
\pgfpathlineto{\pgfqpoint{5.319823in}{0.699967in}}%
\pgfpathlineto{\pgfqpoint{5.323179in}{0.678463in}}%
\pgfpathlineto{\pgfqpoint{5.326535in}{0.616955in}}%
\pgfpathlineto{\pgfqpoint{5.329891in}{0.694458in}}%
\pgfpathlineto{\pgfqpoint{5.333246in}{0.715747in}}%
\pgfpathlineto{\pgfqpoint{5.336602in}{0.722782in}}%
\pgfpathlineto{\pgfqpoint{5.339958in}{0.715665in}}%
\pgfpathlineto{\pgfqpoint{5.346669in}{0.659572in}}%
\pgfpathlineto{\pgfqpoint{5.350025in}{0.707751in}}%
\pgfpathlineto{\pgfqpoint{5.353381in}{0.722907in}}%
\pgfpathlineto{\pgfqpoint{5.356737in}{0.726931in}}%
\pgfpathlineto{\pgfqpoint{5.360093in}{0.722212in}}%
\pgfpathlineto{\pgfqpoint{5.363448in}{0.704858in}}%
\pgfpathlineto{\pgfqpoint{5.366804in}{0.625132in}}%
\pgfpathlineto{\pgfqpoint{5.370160in}{0.701262in}}%
\pgfpathlineto{\pgfqpoint{5.373516in}{0.728346in}}%
\pgfpathlineto{\pgfqpoint{5.376871in}{0.739193in}}%
\pgfpathlineto{\pgfqpoint{5.380227in}{0.738586in}}%
\pgfpathlineto{\pgfqpoint{5.383583in}{0.722018in}}%
\pgfpathlineto{\pgfqpoint{5.386939in}{0.664005in}}%
\pgfpathlineto{\pgfqpoint{5.390295in}{0.702087in}}%
\pgfpathlineto{\pgfqpoint{5.393650in}{0.728637in}}%
\pgfpathlineto{\pgfqpoint{5.397006in}{0.737191in}}%
\pgfpathlineto{\pgfqpoint{5.400362in}{0.735830in}}%
\pgfpathlineto{\pgfqpoint{5.403718in}{0.723248in}}%
\pgfpathlineto{\pgfqpoint{5.407074in}{0.679605in}}%
\pgfpathlineto{\pgfqpoint{5.410429in}{0.700796in}}%
\pgfpathlineto{\pgfqpoint{5.413785in}{0.736085in}}%
\pgfpathlineto{\pgfqpoint{5.417141in}{0.749867in}}%
\pgfpathlineto{\pgfqpoint{5.420497in}{0.751733in}}%
\pgfpathlineto{\pgfqpoint{5.423852in}{0.739589in}}%
\pgfpathlineto{\pgfqpoint{5.427208in}{0.711602in}}%
\pgfpathlineto{\pgfqpoint{5.430564in}{0.620446in}}%
\pgfpathlineto{\pgfqpoint{5.433920in}{0.705753in}}%
\pgfpathlineto{\pgfqpoint{5.437276in}{0.719280in}}%
\pgfpathlineto{\pgfqpoint{5.440631in}{0.716167in}}%
\pgfpathlineto{\pgfqpoint{5.443987in}{0.691985in}}%
\pgfpathlineto{\pgfqpoint{5.447343in}{0.662353in}}%
\pgfpathlineto{\pgfqpoint{5.450699in}{0.719928in}}%
\pgfpathlineto{\pgfqpoint{5.454054in}{0.738706in}}%
\pgfpathlineto{\pgfqpoint{5.457410in}{0.743334in}}%
\pgfpathlineto{\pgfqpoint{5.460766in}{0.733139in}}%
\pgfpathlineto{\pgfqpoint{5.464122in}{0.711234in}}%
\pgfpathlineto{\pgfqpoint{5.467478in}{0.660406in}}%
\pgfpathlineto{\pgfqpoint{5.470833in}{0.664893in}}%
\pgfpathlineto{\pgfqpoint{5.474189in}{0.688597in}}%
\pgfpathlineto{\pgfqpoint{5.477545in}{0.684357in}}%
\pgfpathlineto{\pgfqpoint{5.480901in}{0.639350in}}%
\pgfpathlineto{\pgfqpoint{5.487612in}{0.709518in}}%
\pgfpathlineto{\pgfqpoint{5.490968in}{0.722463in}}%
\pgfpathlineto{\pgfqpoint{5.494324in}{0.720930in}}%
\pgfpathlineto{\pgfqpoint{5.497680in}{0.704096in}}%
\pgfpathlineto{\pgfqpoint{5.501035in}{0.664243in}}%
\pgfpathlineto{\pgfqpoint{5.504391in}{0.649595in}}%
\pgfpathlineto{\pgfqpoint{5.507747in}{0.686025in}}%
\pgfpathlineto{\pgfqpoint{5.511103in}{0.692955in}}%
\pgfpathlineto{\pgfqpoint{5.514458in}{0.684576in}}%
\pgfpathlineto{\pgfqpoint{5.517814in}{0.644739in}}%
\pgfpathlineto{\pgfqpoint{5.521170in}{0.663638in}}%
\pgfpathlineto{\pgfqpoint{5.524526in}{0.700375in}}%
\pgfpathlineto{\pgfqpoint{5.527882in}{0.714166in}}%
\pgfpathlineto{\pgfqpoint{5.531237in}{0.714336in}}%
\pgfpathlineto{\pgfqpoint{5.534593in}{0.696651in}}%
\pgfpathlineto{\pgfqpoint{5.537949in}{0.625431in}}%
\pgfpathlineto{\pgfqpoint{5.541305in}{0.682987in}}%
\pgfpathlineto{\pgfqpoint{5.544660in}{0.707585in}}%
\pgfpathlineto{\pgfqpoint{5.548016in}{0.715614in}}%
\pgfpathlineto{\pgfqpoint{5.551372in}{0.714068in}}%
\pgfpathlineto{\pgfqpoint{5.554728in}{0.701519in}}%
\pgfpathlineto{\pgfqpoint{5.558084in}{0.658826in}}%
\pgfpathlineto{\pgfqpoint{5.561439in}{0.677707in}}%
\pgfpathlineto{\pgfqpoint{5.564795in}{0.713574in}}%
\pgfpathlineto{\pgfqpoint{5.568151in}{0.727588in}}%
\pgfpathlineto{\pgfqpoint{5.571507in}{0.729787in}}%
\pgfpathlineto{\pgfqpoint{5.574863in}{0.716405in}}%
\pgfpathlineto{\pgfqpoint{5.578218in}{0.663989in}}%
\pgfpathlineto{\pgfqpoint{5.584930in}{0.724882in}}%
\pgfpathlineto{\pgfqpoint{5.588286in}{0.735517in}}%
\pgfpathlineto{\pgfqpoint{5.591641in}{0.736832in}}%
\pgfpathlineto{\pgfqpoint{5.594997in}{0.729291in}}%
\pgfpathlineto{\pgfqpoint{5.598353in}{0.706043in}}%
\pgfpathlineto{\pgfqpoint{5.601709in}{0.641923in}}%
\pgfpathlineto{\pgfqpoint{5.605065in}{0.720720in}}%
\pgfpathlineto{\pgfqpoint{5.608420in}{0.742103in}}%
\pgfpathlineto{\pgfqpoint{5.611776in}{0.750100in}}%
\pgfpathlineto{\pgfqpoint{5.615132in}{0.746166in}}%
\pgfpathlineto{\pgfqpoint{5.618488in}{0.726904in}}%
\pgfpathlineto{\pgfqpoint{5.621843in}{0.665950in}}%
\pgfpathlineto{\pgfqpoint{5.625199in}{0.702917in}}%
\pgfpathlineto{\pgfqpoint{5.628555in}{0.727618in}}%
\pgfpathlineto{\pgfqpoint{5.631911in}{0.733423in}}%
\pgfpathlineto{\pgfqpoint{5.635267in}{0.727285in}}%
\pgfpathlineto{\pgfqpoint{5.638622in}{0.701747in}}%
\pgfpathlineto{\pgfqpoint{5.641978in}{0.669486in}}%
\pgfpathlineto{\pgfqpoint{5.645334in}{0.727553in}}%
\pgfpathlineto{\pgfqpoint{5.648690in}{0.746345in}}%
\pgfpathlineto{\pgfqpoint{5.652045in}{0.751542in}}%
\pgfpathlineto{\pgfqpoint{5.655401in}{0.742587in}}%
\pgfpathlineto{\pgfqpoint{5.658757in}{0.721873in}}%
\pgfpathlineto{\pgfqpoint{5.662113in}{0.675536in}}%
\pgfpathlineto{\pgfqpoint{5.665469in}{0.669236in}}%
\pgfpathlineto{\pgfqpoint{5.668824in}{0.697161in}}%
\pgfpathlineto{\pgfqpoint{5.672180in}{0.695328in}}%
\pgfpathlineto{\pgfqpoint{5.675536in}{0.660641in}}%
\pgfpathlineto{\pgfqpoint{5.678892in}{0.676092in}}%
\pgfpathlineto{\pgfqpoint{5.682247in}{0.715324in}}%
\pgfpathlineto{\pgfqpoint{5.685603in}{0.730131in}}%
\pgfpathlineto{\pgfqpoint{5.688959in}{0.730704in}}%
\pgfpathlineto{\pgfqpoint{5.692315in}{0.716438in}}%
\pgfpathlineto{\pgfqpoint{5.695671in}{0.688133in}}%
\pgfpathlineto{\pgfqpoint{5.699026in}{0.526715in}}%
\pgfpathlineto{\pgfqpoint{5.702382in}{0.672000in}}%
\pgfpathlineto{\pgfqpoint{5.705738in}{0.683360in}}%
\pgfpathlineto{\pgfqpoint{5.709094in}{0.672743in}}%
\pgfpathlineto{\pgfqpoint{5.712449in}{0.586753in}}%
\pgfpathlineto{\pgfqpoint{5.715805in}{0.677586in}}%
\pgfpathlineto{\pgfqpoint{5.719161in}{0.704820in}}%
\pgfpathlineto{\pgfqpoint{5.722517in}{0.714619in}}%
\pgfpathlineto{\pgfqpoint{5.725873in}{0.709358in}}%
\pgfpathlineto{\pgfqpoint{5.729228in}{0.686222in}}%
\pgfpathlineto{\pgfqpoint{5.732584in}{0.532239in}}%
\pgfpathlineto{\pgfqpoint{5.735940in}{0.680007in}}%
\pgfpathlineto{\pgfqpoint{5.739296in}{0.699100in}}%
\pgfpathlineto{\pgfqpoint{5.742652in}{0.703118in}}%
\pgfpathlineto{\pgfqpoint{5.746007in}{0.695936in}}%
\pgfpathlineto{\pgfqpoint{5.752719in}{0.641676in}}%
\pgfpathlineto{\pgfqpoint{5.756075in}{0.696690in}}%
\pgfpathlineto{\pgfqpoint{5.759430in}{0.715013in}}%
\pgfpathlineto{\pgfqpoint{5.762786in}{0.720073in}}%
\pgfpathlineto{\pgfqpoint{5.766142in}{0.710068in}}%
\pgfpathlineto{\pgfqpoint{5.769498in}{0.668521in}}%
\pgfpathlineto{\pgfqpoint{5.772854in}{0.679938in}}%
\pgfpathlineto{\pgfqpoint{5.776209in}{0.714791in}}%
\pgfpathlineto{\pgfqpoint{5.779565in}{0.727402in}}%
\pgfpathlineto{\pgfqpoint{5.782921in}{0.730411in}}%
\pgfpathlineto{\pgfqpoint{5.786277in}{0.725233in}}%
\pgfpathlineto{\pgfqpoint{5.789632in}{0.707842in}}%
\pgfpathlineto{\pgfqpoint{5.792988in}{0.632454in}}%
\pgfpathlineto{\pgfqpoint{5.796344in}{0.702662in}}%
\pgfpathlineto{\pgfqpoint{5.799700in}{0.730239in}}%
\pgfpathlineto{\pgfqpoint{5.803056in}{0.741499in}}%
\pgfpathlineto{\pgfqpoint{5.806411in}{0.741760in}}%
\pgfpathlineto{\pgfqpoint{5.809767in}{0.728741in}}%
\pgfpathlineto{\pgfqpoint{5.813123in}{0.700789in}}%
\pgfpathlineto{\pgfqpoint{5.816479in}{0.561599in}}%
\pgfpathlineto{\pgfqpoint{5.819834in}{0.690206in}}%
\pgfpathlineto{\pgfqpoint{5.823190in}{0.703175in}}%
\pgfpathlineto{\pgfqpoint{5.826546in}{0.697683in}}%
\pgfpathlineto{\pgfqpoint{5.829902in}{0.662872in}}%
\pgfpathlineto{\pgfqpoint{5.833258in}{0.672894in}}%
\pgfpathlineto{\pgfqpoint{5.836613in}{0.712831in}}%
\pgfpathlineto{\pgfqpoint{5.839969in}{0.727667in}}%
\pgfpathlineto{\pgfqpoint{5.843325in}{0.728595in}}%
\pgfpathlineto{\pgfqpoint{5.846681in}{0.714417in}}%
\pgfpathlineto{\pgfqpoint{5.850036in}{0.684127in}}%
\pgfpathlineto{\pgfqpoint{5.853392in}{0.605633in}}%
\pgfpathlineto{\pgfqpoint{5.856748in}{0.680151in}}%
\pgfpathlineto{\pgfqpoint{5.860104in}{0.690937in}}%
\pgfpathlineto{\pgfqpoint{5.863460in}{0.683693in}}%
\pgfpathlineto{\pgfqpoint{5.866815in}{0.642233in}}%
\pgfpathlineto{\pgfqpoint{5.870171in}{0.666855in}}%
\pgfpathlineto{\pgfqpoint{5.873527in}{0.702510in}}%
\pgfpathlineto{\pgfqpoint{5.876883in}{0.715927in}}%
\pgfpathlineto{\pgfqpoint{5.880238in}{0.715351in}}%
\pgfpathlineto{\pgfqpoint{5.883594in}{0.697528in}}%
\pgfpathlineto{\pgfqpoint{5.886950in}{0.638297in}}%
\pgfpathlineto{\pgfqpoint{5.890306in}{0.675025in}}%
\pgfpathlineto{\pgfqpoint{5.893662in}{0.701021in}}%
\pgfpathlineto{\pgfqpoint{5.897017in}{0.708424in}}%
\pgfpathlineto{\pgfqpoint{5.900373in}{0.705076in}}%
\pgfpathlineto{\pgfqpoint{5.903729in}{0.687721in}}%
\pgfpathlineto{\pgfqpoint{5.907085in}{0.586617in}}%
\pgfpathlineto{\pgfqpoint{5.910441in}{0.688965in}}%
\pgfpathlineto{\pgfqpoint{5.913796in}{0.714613in}}%
\pgfpathlineto{\pgfqpoint{5.917152in}{0.724244in}}%
\pgfpathlineto{\pgfqpoint{5.920508in}{0.720968in}}%
\pgfpathlineto{\pgfqpoint{5.923864in}{0.697899in}}%
\pgfpathlineto{\pgfqpoint{5.927219in}{0.621657in}}%
\pgfpathlineto{\pgfqpoint{5.930575in}{0.705321in}}%
\pgfpathlineto{\pgfqpoint{5.933931in}{0.724348in}}%
\pgfpathlineto{\pgfqpoint{5.937287in}{0.730594in}}%
\pgfpathlineto{\pgfqpoint{5.940643in}{0.728211in}}%
\pgfpathlineto{\pgfqpoint{5.943998in}{0.715256in}}%
\pgfpathlineto{\pgfqpoint{5.947354in}{0.672774in}}%
\pgfpathlineto{\pgfqpoint{5.950710in}{0.690399in}}%
\pgfpathlineto{\pgfqpoint{5.954066in}{0.726549in}}%
\pgfpathlineto{\pgfqpoint{5.957421in}{0.740828in}}%
\pgfpathlineto{\pgfqpoint{5.960777in}{0.743680in}}%
\pgfpathlineto{\pgfqpoint{5.964133in}{0.732500in}}%
\pgfpathlineto{\pgfqpoint{5.967489in}{0.698229in}}%
\pgfpathlineto{\pgfqpoint{5.970845in}{0.678289in}}%
\pgfpathlineto{\pgfqpoint{5.974200in}{0.722006in}}%
\pgfpathlineto{\pgfqpoint{5.977556in}{0.734575in}}%
\pgfpathlineto{\pgfqpoint{5.980912in}{0.735485in}}%
\pgfpathlineto{\pgfqpoint{5.984268in}{0.725308in}}%
\pgfpathlineto{\pgfqpoint{5.987623in}{0.690208in}}%
\pgfpathlineto{\pgfqpoint{5.990979in}{0.691628in}}%
\pgfpathlineto{\pgfqpoint{5.994335in}{0.733145in}}%
\pgfpathlineto{\pgfqpoint{5.997691in}{0.748642in}}%
\pgfpathlineto{\pgfqpoint{6.001047in}{0.751699in}}%
\pgfpathlineto{\pgfqpoint{6.004402in}{0.740738in}}%
\pgfpathlineto{\pgfqpoint{6.007758in}{0.716042in}}%
\pgfpathlineto{\pgfqpoint{6.011114in}{0.632432in}}%
\pgfpathlineto{\pgfqpoint{6.014470in}{0.694134in}}%
\pgfpathlineto{\pgfqpoint{6.017825in}{0.711163in}}%
\pgfpathlineto{\pgfqpoint{6.021181in}{0.708209in}}%
\pgfpathlineto{\pgfqpoint{6.024537in}{0.680381in}}%
\pgfpathlineto{\pgfqpoint{6.027893in}{0.670645in}}%
\pgfpathlineto{\pgfqpoint{6.031249in}{0.718467in}}%
\pgfpathlineto{\pgfqpoint{6.034604in}{0.735414in}}%
\pgfpathlineto{\pgfqpoint{6.037960in}{0.738312in}}%
\pgfpathlineto{\pgfqpoint{6.041316in}{0.726232in}}%
\pgfpathlineto{\pgfqpoint{6.044672in}{0.701891in}}%
\pgfpathlineto{\pgfqpoint{6.048027in}{0.635522in}}%
\pgfpathlineto{\pgfqpoint{6.051383in}{0.668798in}}%
\pgfpathlineto{\pgfqpoint{6.054739in}{0.686205in}}%
\pgfpathlineto{\pgfqpoint{6.058095in}{0.678844in}}%
\pgfpathlineto{\pgfqpoint{6.061451in}{0.618305in}}%
\pgfpathlineto{\pgfqpoint{6.064806in}{0.677842in}}%
\pgfpathlineto{\pgfqpoint{6.068162in}{0.707958in}}%
\pgfpathlineto{\pgfqpoint{6.071518in}{0.719210in}}%
\pgfpathlineto{\pgfqpoint{6.074874in}{0.715668in}}%
\pgfpathlineto{\pgfqpoint{6.078230in}{0.696283in}}%
\pgfpathlineto{\pgfqpoint{6.081585in}{0.640507in}}%
\pgfpathlineto{\pgfqpoint{6.088297in}{0.691933in}}%
\pgfpathlineto{\pgfqpoint{6.091653in}{0.696948in}}%
\pgfpathlineto{\pgfqpoint{6.095008in}{0.688575in}}%
\pgfpathlineto{\pgfqpoint{6.098364in}{0.653477in}}%
\pgfpathlineto{\pgfqpoint{6.101720in}{0.658856in}}%
\pgfpathlineto{\pgfqpoint{6.105076in}{0.699486in}}%
\pgfpathlineto{\pgfqpoint{6.108432in}{0.714524in}}%
\pgfpathlineto{\pgfqpoint{6.111787in}{0.716305in}}%
\pgfpathlineto{\pgfqpoint{6.115143in}{0.700835in}}%
\pgfpathlineto{\pgfqpoint{6.118499in}{0.639164in}}%
\pgfpathlineto{\pgfqpoint{6.121855in}{0.684511in}}%
\pgfpathlineto{\pgfqpoint{6.125210in}{0.711301in}}%
\pgfpathlineto{\pgfqpoint{6.128566in}{0.720692in}}%
\pgfpathlineto{\pgfqpoint{6.131922in}{0.720785in}}%
\pgfpathlineto{\pgfqpoint{6.135278in}{0.711306in}}%
\pgfpathlineto{\pgfqpoint{6.138634in}{0.682036in}}%
\pgfpathlineto{\pgfqpoint{6.141989in}{0.660222in}}%
\pgfpathlineto{\pgfqpoint{6.145345in}{0.711714in}}%
\pgfpathlineto{\pgfqpoint{6.148701in}{0.729661in}}%
\pgfpathlineto{\pgfqpoint{6.152057in}{0.735307in}}%
\pgfpathlineto{\pgfqpoint{6.155412in}{0.727790in}}%
\pgfpathlineto{\pgfqpoint{6.162124in}{0.671601in}}%
\pgfpathlineto{\pgfqpoint{6.165480in}{0.720471in}}%
\pgfpathlineto{\pgfqpoint{6.168836in}{0.735526in}}%
\pgfpathlineto{\pgfqpoint{6.172191in}{0.739450in}}%
\pgfpathlineto{\pgfqpoint{6.175547in}{0.734587in}}%
\pgfpathlineto{\pgfqpoint{6.178903in}{0.716896in}}%
\pgfpathlineto{\pgfqpoint{6.182259in}{0.631209in}}%
\pgfpathlineto{\pgfqpoint{6.185614in}{0.714637in}}%
\pgfpathlineto{\pgfqpoint{6.188970in}{0.741228in}}%
\pgfpathlineto{\pgfqpoint{6.192326in}{0.751816in}}%
\pgfpathlineto{\pgfqpoint{6.195682in}{0.750856in}}%
\pgfpathlineto{\pgfqpoint{6.199038in}{0.735593in}}%
\pgfpathlineto{\pgfqpoint{6.202393in}{0.698154in}}%
\pgfpathlineto{\pgfqpoint{6.205749in}{0.678519in}}%
\pgfpathlineto{\pgfqpoint{6.209105in}{0.718554in}}%
\pgfpathlineto{\pgfqpoint{6.212461in}{0.727547in}}%
\pgfpathlineto{\pgfqpoint{6.215816in}{0.722646in}}%
\pgfpathlineto{\pgfqpoint{6.222528in}{0.668942in}}%
\pgfpathlineto{\pgfqpoint{6.225884in}{0.725218in}}%
\pgfpathlineto{\pgfqpoint{6.229240in}{0.743713in}}%
\pgfpathlineto{\pgfqpoint{6.232595in}{0.748351in}}%
\pgfpathlineto{\pgfqpoint{6.235951in}{0.738470in}}%
\pgfpathlineto{\pgfqpoint{6.239307in}{0.717311in}}%
\pgfpathlineto{\pgfqpoint{6.242663in}{0.670993in}}%
\pgfpathlineto{\pgfqpoint{6.246019in}{0.662162in}}%
\pgfpathlineto{\pgfqpoint{6.249374in}{0.689543in}}%
\pgfpathlineto{\pgfqpoint{6.252730in}{0.685290in}}%
\pgfpathlineto{\pgfqpoint{6.256086in}{0.633870in}}%
\pgfpathlineto{\pgfqpoint{6.259442in}{0.681755in}}%
\pgfpathlineto{\pgfqpoint{6.262797in}{0.713887in}}%
\pgfpathlineto{\pgfqpoint{6.266153in}{0.726122in}}%
\pgfpathlineto{\pgfqpoint{6.269509in}{0.723667in}}%
\pgfpathlineto{\pgfqpoint{6.272865in}{0.706904in}}%
\pgfpathlineto{\pgfqpoint{6.279576in}{0.634426in}}%
\pgfpathlineto{\pgfqpoint{6.282932in}{0.679242in}}%
\pgfpathlineto{\pgfqpoint{6.286288in}{0.685616in}}%
\pgfpathlineto{\pgfqpoint{6.289644in}{0.673071in}}%
\pgfpathlineto{\pgfqpoint{6.292999in}{0.582798in}}%
\pgfpathlineto{\pgfqpoint{6.296355in}{0.677045in}}%
\pgfpathlineto{\pgfqpoint{6.299711in}{0.703741in}}%
\pgfpathlineto{\pgfqpoint{6.303067in}{0.713254in}}%
\pgfpathlineto{\pgfqpoint{6.306423in}{0.707852in}}%
\pgfpathlineto{\pgfqpoint{6.309778in}{0.682874in}}%
\pgfpathlineto{\pgfqpoint{6.313134in}{0.604522in}}%
\pgfpathlineto{\pgfqpoint{6.316490in}{0.687403in}}%
\pgfpathlineto{\pgfqpoint{6.319846in}{0.704977in}}%
\pgfpathlineto{\pgfqpoint{6.323201in}{0.709250in}}%
\pgfpathlineto{\pgfqpoint{6.326557in}{0.703561in}}%
\pgfpathlineto{\pgfqpoint{6.329913in}{0.682380in}}%
\pgfpathlineto{\pgfqpoint{6.333269in}{0.599618in}}%
\pgfpathlineto{\pgfqpoint{6.336625in}{0.693982in}}%
\pgfpathlineto{\pgfqpoint{6.339980in}{0.716250in}}%
\pgfpathlineto{\pgfqpoint{6.343336in}{0.724233in}}%
\pgfpathlineto{\pgfqpoint{6.346692in}{0.719280in}}%
\pgfpathlineto{\pgfqpoint{6.350048in}{0.700184in}}%
\pgfpathlineto{\pgfqpoint{6.353403in}{0.650548in}}%
\pgfpathlineto{\pgfqpoint{6.356759in}{0.661834in}}%
\pgfpathlineto{\pgfqpoint{6.360115in}{0.689542in}}%
\pgfpathlineto{\pgfqpoint{6.363471in}{0.694032in}}%
\pgfpathlineto{\pgfqpoint{6.366827in}{0.683282in}}%
\pgfpathlineto{\pgfqpoint{6.370182in}{0.632438in}}%
\pgfpathlineto{\pgfqpoint{6.373538in}{0.671735in}}%
\pgfpathlineto{\pgfqpoint{6.376894in}{0.703749in}}%
\pgfpathlineto{\pgfqpoint{6.380250in}{0.715798in}}%
\pgfpathlineto{\pgfqpoint{6.383605in}{0.714071in}}%
\pgfpathlineto{\pgfqpoint{6.383605in}{0.714071in}}%
\pgfusepath{stroke}%
\end{pgfscope}%
\begin{pgfscope}%
\pgfpathrectangle{\pgfqpoint{4.621821in}{0.320625in}}{\pgfqpoint{1.845679in}{1.164375in}}%
\pgfusepath{clip}%
\pgfsetrectcap%
\pgfsetroundjoin%
\pgfsetlinewidth{0.602250pt}%
\definecolor{currentstroke}{rgb}{0.478431,0.243137,0.615686}%
\pgfsetstrokecolor{currentstroke}%
\pgfsetstrokeopacity{0.120000}%
\pgfsetdash{}{0pt}%
\pgfpathmoveto{\pgfqpoint{4.705715in}{1.468685in}}%
\pgfpathlineto{\pgfqpoint{4.715783in}{1.450011in}}%
\pgfpathlineto{\pgfqpoint{4.725850in}{1.428817in}}%
\pgfpathlineto{\pgfqpoint{4.739273in}{1.397199in}}%
\pgfpathlineto{\pgfqpoint{4.776187in}{1.308109in}}%
\pgfpathlineto{\pgfqpoint{4.796322in}{1.262917in}}%
\pgfpathlineto{\pgfqpoint{4.819812in}{1.213174in}}%
\pgfpathlineto{\pgfqpoint{4.846658in}{1.159034in}}%
\pgfpathlineto{\pgfqpoint{4.880216in}{1.093847in}}%
\pgfpathlineto{\pgfqpoint{4.923841in}{1.011522in}}%
\pgfpathlineto{\pgfqpoint{4.984245in}{0.900090in}}%
\pgfpathlineto{\pgfqpoint{5.044649in}{0.790481in}}%
\pgfpathlineto{\pgfqpoint{5.058072in}{0.767949in}}%
\pgfpathlineto{\pgfqpoint{5.078207in}{0.728137in}}%
\pgfpathlineto{\pgfqpoint{5.081563in}{0.723244in}}%
\pgfpathlineto{\pgfqpoint{5.088274in}{0.716381in}}%
\pgfpathlineto{\pgfqpoint{5.091630in}{0.713581in}}%
\pgfpathlineto{\pgfqpoint{5.094986in}{0.709939in}}%
\pgfpathlineto{\pgfqpoint{5.098342in}{0.704065in}}%
\pgfpathlineto{\pgfqpoint{5.101698in}{0.696468in}}%
\pgfpathlineto{\pgfqpoint{5.105053in}{0.687267in}}%
\pgfpathlineto{\pgfqpoint{5.111765in}{0.663927in}}%
\pgfpathlineto{\pgfqpoint{5.118476in}{0.638258in}}%
\pgfpathlineto{\pgfqpoint{5.121832in}{0.628839in}}%
\pgfpathlineto{\pgfqpoint{5.125188in}{0.622344in}}%
\pgfpathlineto{\pgfqpoint{5.128544in}{0.619982in}}%
\pgfpathlineto{\pgfqpoint{5.131900in}{0.626530in}}%
\pgfpathlineto{\pgfqpoint{5.138611in}{0.650875in}}%
\pgfpathlineto{\pgfqpoint{5.141967in}{0.659697in}}%
\pgfpathlineto{\pgfqpoint{5.145323in}{0.664980in}}%
\pgfpathlineto{\pgfqpoint{5.148678in}{0.666102in}}%
\pgfpathlineto{\pgfqpoint{5.152034in}{0.661270in}}%
\pgfpathlineto{\pgfqpoint{5.155390in}{0.649770in}}%
\pgfpathlineto{\pgfqpoint{5.158746in}{0.651698in}}%
\pgfpathlineto{\pgfqpoint{5.162102in}{0.671951in}}%
\pgfpathlineto{\pgfqpoint{5.165457in}{0.688209in}}%
\pgfpathlineto{\pgfqpoint{5.168813in}{0.699284in}}%
\pgfpathlineto{\pgfqpoint{5.172169in}{0.706630in}}%
\pgfpathlineto{\pgfqpoint{5.175525in}{0.711125in}}%
\pgfpathlineto{\pgfqpoint{5.178880in}{0.713198in}}%
\pgfpathlineto{\pgfqpoint{5.182236in}{0.712970in}}%
\pgfpathlineto{\pgfqpoint{5.185592in}{0.710271in}}%
\pgfpathlineto{\pgfqpoint{5.188948in}{0.704526in}}%
\pgfpathlineto{\pgfqpoint{5.192304in}{0.694293in}}%
\pgfpathlineto{\pgfqpoint{5.195659in}{0.675384in}}%
\pgfpathlineto{\pgfqpoint{5.199015in}{0.626458in}}%
\pgfpathlineto{\pgfqpoint{5.205727in}{0.690366in}}%
\pgfpathlineto{\pgfqpoint{5.209082in}{0.706673in}}%
\pgfpathlineto{\pgfqpoint{5.212438in}{0.716664in}}%
\pgfpathlineto{\pgfqpoint{5.215794in}{0.722585in}}%
\pgfpathlineto{\pgfqpoint{5.219150in}{0.725024in}}%
\pgfpathlineto{\pgfqpoint{5.222506in}{0.723610in}}%
\pgfpathlineto{\pgfqpoint{5.225861in}{0.716641in}}%
\pgfpathlineto{\pgfqpoint{5.229217in}{0.703241in}}%
\pgfpathlineto{\pgfqpoint{5.232573in}{0.685852in}}%
\pgfpathlineto{\pgfqpoint{5.235929in}{0.679347in}}%
\pgfpathlineto{\pgfqpoint{5.239285in}{0.685808in}}%
\pgfpathlineto{\pgfqpoint{5.242640in}{0.687947in}}%
\pgfpathlineto{\pgfqpoint{5.245996in}{0.680692in}}%
\pgfpathlineto{\pgfqpoint{5.249352in}{0.658170in}}%
\pgfpathlineto{\pgfqpoint{5.252708in}{0.654138in}}%
\pgfpathlineto{\pgfqpoint{5.256063in}{0.686815in}}%
\pgfpathlineto{\pgfqpoint{5.259419in}{0.702526in}}%
\pgfpathlineto{\pgfqpoint{5.262775in}{0.706444in}}%
\pgfpathlineto{\pgfqpoint{5.266131in}{0.696570in}}%
\pgfpathlineto{\pgfqpoint{5.269487in}{0.672379in}}%
\pgfpathlineto{\pgfqpoint{5.272842in}{0.634090in}}%
\pgfpathlineto{\pgfqpoint{5.276198in}{0.667352in}}%
\pgfpathlineto{\pgfqpoint{5.279554in}{0.681869in}}%
\pgfpathlineto{\pgfqpoint{5.282910in}{0.682826in}}%
\pgfpathlineto{\pgfqpoint{5.286265in}{0.670758in}}%
\pgfpathlineto{\pgfqpoint{5.289621in}{0.623507in}}%
\pgfpathlineto{\pgfqpoint{5.296333in}{0.688965in}}%
\pgfpathlineto{\pgfqpoint{5.299689in}{0.701595in}}%
\pgfpathlineto{\pgfqpoint{5.303044in}{0.701716in}}%
\pgfpathlineto{\pgfqpoint{5.306400in}{0.681329in}}%
\pgfpathlineto{\pgfqpoint{5.309756in}{0.634720in}}%
\pgfpathlineto{\pgfqpoint{5.313112in}{0.698646in}}%
\pgfpathlineto{\pgfqpoint{5.316467in}{0.718997in}}%
\pgfpathlineto{\pgfqpoint{5.319823in}{0.727924in}}%
\pgfpathlineto{\pgfqpoint{5.323179in}{0.729864in}}%
\pgfpathlineto{\pgfqpoint{5.326535in}{0.725239in}}%
\pgfpathlineto{\pgfqpoint{5.329891in}{0.711223in}}%
\pgfpathlineto{\pgfqpoint{5.333246in}{0.669915in}}%
\pgfpathlineto{\pgfqpoint{5.336602in}{0.682536in}}%
\pgfpathlineto{\pgfqpoint{5.339958in}{0.720112in}}%
\pgfpathlineto{\pgfqpoint{5.343314in}{0.735600in}}%
\pgfpathlineto{\pgfqpoint{5.346669in}{0.741023in}}%
\pgfpathlineto{\pgfqpoint{5.350025in}{0.735921in}}%
\pgfpathlineto{\pgfqpoint{5.353381in}{0.718968in}}%
\pgfpathlineto{\pgfqpoint{5.360093in}{0.652285in}}%
\pgfpathlineto{\pgfqpoint{5.363448in}{0.692456in}}%
\pgfpathlineto{\pgfqpoint{5.366804in}{0.698806in}}%
\pgfpathlineto{\pgfqpoint{5.370160in}{0.687098in}}%
\pgfpathlineto{\pgfqpoint{5.373516in}{0.616452in}}%
\pgfpathlineto{\pgfqpoint{5.376871in}{0.687156in}}%
\pgfpathlineto{\pgfqpoint{5.380227in}{0.715315in}}%
\pgfpathlineto{\pgfqpoint{5.383583in}{0.725816in}}%
\pgfpathlineto{\pgfqpoint{5.386939in}{0.722159in}}%
\pgfpathlineto{\pgfqpoint{5.390295in}{0.704952in}}%
\pgfpathlineto{\pgfqpoint{5.393650in}{0.667763in}}%
\pgfpathlineto{\pgfqpoint{5.397006in}{0.639784in}}%
\pgfpathlineto{\pgfqpoint{5.400362in}{0.679633in}}%
\pgfpathlineto{\pgfqpoint{5.403718in}{0.685707in}}%
\pgfpathlineto{\pgfqpoint{5.407074in}{0.673752in}}%
\pgfpathlineto{\pgfqpoint{5.410429in}{0.600063in}}%
\pgfpathlineto{\pgfqpoint{5.413785in}{0.674279in}}%
\pgfpathlineto{\pgfqpoint{5.417141in}{0.702217in}}%
\pgfpathlineto{\pgfqpoint{5.420497in}{0.712581in}}%
\pgfpathlineto{\pgfqpoint{5.423852in}{0.708684in}}%
\pgfpathlineto{\pgfqpoint{5.427208in}{0.686364in}}%
\pgfpathlineto{\pgfqpoint{5.430564in}{0.574386in}}%
\pgfpathlineto{\pgfqpoint{5.433920in}{0.683595in}}%
\pgfpathlineto{\pgfqpoint{5.437276in}{0.703500in}}%
\pgfpathlineto{\pgfqpoint{5.440631in}{0.709162in}}%
\pgfpathlineto{\pgfqpoint{5.443987in}{0.705160in}}%
\pgfpathlineto{\pgfqpoint{5.447343in}{0.687933in}}%
\pgfpathlineto{\pgfqpoint{5.450699in}{0.601768in}}%
\pgfpathlineto{\pgfqpoint{5.454054in}{0.686319in}}%
\pgfpathlineto{\pgfqpoint{5.457410in}{0.712884in}}%
\pgfpathlineto{\pgfqpoint{5.460766in}{0.723301in}}%
\pgfpathlineto{\pgfqpoint{5.464122in}{0.721739in}}%
\pgfpathlineto{\pgfqpoint{5.467478in}{0.701513in}}%
\pgfpathlineto{\pgfqpoint{5.470833in}{0.592483in}}%
\pgfpathlineto{\pgfqpoint{5.474189in}{0.705149in}}%
\pgfpathlineto{\pgfqpoint{5.477545in}{0.726334in}}%
\pgfpathlineto{\pgfqpoint{5.480901in}{0.734227in}}%
\pgfpathlineto{\pgfqpoint{5.484256in}{0.733951in}}%
\pgfpathlineto{\pgfqpoint{5.487612in}{0.724986in}}%
\pgfpathlineto{\pgfqpoint{5.490968in}{0.698966in}}%
\pgfpathlineto{\pgfqpoint{5.494324in}{0.656762in}}%
\pgfpathlineto{\pgfqpoint{5.497680in}{0.719938in}}%
\pgfpathlineto{\pgfqpoint{5.501035in}{0.739887in}}%
\pgfpathlineto{\pgfqpoint{5.504391in}{0.747375in}}%
\pgfpathlineto{\pgfqpoint{5.507747in}{0.743429in}}%
\pgfpathlineto{\pgfqpoint{5.511103in}{0.727445in}}%
\pgfpathlineto{\pgfqpoint{5.514458in}{0.695569in}}%
\pgfpathlineto{\pgfqpoint{5.517814in}{0.622123in}}%
\pgfpathlineto{\pgfqpoint{5.521170in}{0.689784in}}%
\pgfpathlineto{\pgfqpoint{5.524526in}{0.697714in}}%
\pgfpathlineto{\pgfqpoint{5.527882in}{0.684091in}}%
\pgfpathlineto{\pgfqpoint{5.531237in}{0.554433in}}%
\pgfpathlineto{\pgfqpoint{5.534593in}{0.696120in}}%
\pgfpathlineto{\pgfqpoint{5.537949in}{0.720763in}}%
\pgfpathlineto{\pgfqpoint{5.541305in}{0.729324in}}%
\pgfpathlineto{\pgfqpoint{5.544660in}{0.723081in}}%
\pgfpathlineto{\pgfqpoint{5.548016in}{0.704438in}}%
\pgfpathlineto{\pgfqpoint{5.551372in}{0.664878in}}%
\pgfpathlineto{\pgfqpoint{5.554728in}{0.638922in}}%
\pgfpathlineto{\pgfqpoint{5.558084in}{0.676375in}}%
\pgfpathlineto{\pgfqpoint{5.561439in}{0.679194in}}%
\pgfpathlineto{\pgfqpoint{5.564795in}{0.658446in}}%
\pgfpathlineto{\pgfqpoint{5.568151in}{0.625060in}}%
\pgfpathlineto{\pgfqpoint{5.571507in}{0.686672in}}%
\pgfpathlineto{\pgfqpoint{5.574863in}{0.706548in}}%
\pgfpathlineto{\pgfqpoint{5.578218in}{0.711940in}}%
\pgfpathlineto{\pgfqpoint{5.581574in}{0.701405in}}%
\pgfpathlineto{\pgfqpoint{5.584930in}{0.668720in}}%
\pgfpathlineto{\pgfqpoint{5.588286in}{0.641492in}}%
\pgfpathlineto{\pgfqpoint{5.591641in}{0.688750in}}%
\pgfpathlineto{\pgfqpoint{5.594997in}{0.702012in}}%
\pgfpathlineto{\pgfqpoint{5.598353in}{0.703441in}}%
\pgfpathlineto{\pgfqpoint{5.601709in}{0.694056in}}%
\pgfpathlineto{\pgfqpoint{5.605065in}{0.662034in}}%
\pgfpathlineto{\pgfqpoint{5.608420in}{0.653818in}}%
\pgfpathlineto{\pgfqpoint{5.611776in}{0.699224in}}%
\pgfpathlineto{\pgfqpoint{5.615132in}{0.715784in}}%
\pgfpathlineto{\pgfqpoint{5.618488in}{0.719899in}}%
\pgfpathlineto{\pgfqpoint{5.621843in}{0.708984in}}%
\pgfpathlineto{\pgfqpoint{5.625199in}{0.662670in}}%
\pgfpathlineto{\pgfqpoint{5.628555in}{0.685276in}}%
\pgfpathlineto{\pgfqpoint{5.631911in}{0.717650in}}%
\pgfpathlineto{\pgfqpoint{5.635267in}{0.729929in}}%
\pgfpathlineto{\pgfqpoint{5.638622in}{0.733112in}}%
\pgfpathlineto{\pgfqpoint{5.641978in}{0.728585in}}%
\pgfpathlineto{\pgfqpoint{5.645334in}{0.713087in}}%
\pgfpathlineto{\pgfqpoint{5.648690in}{0.658668in}}%
\pgfpathlineto{\pgfqpoint{5.652045in}{0.698184in}}%
\pgfpathlineto{\pgfqpoint{5.655401in}{0.729591in}}%
\pgfpathlineto{\pgfqpoint{5.658757in}{0.742602in}}%
\pgfpathlineto{\pgfqpoint{5.662113in}{0.745017in}}%
\pgfpathlineto{\pgfqpoint{5.665469in}{0.734899in}}%
\pgfpathlineto{\pgfqpoint{5.668824in}{0.712630in}}%
\pgfpathlineto{\pgfqpoint{5.672180in}{0.655553in}}%
\pgfpathlineto{\pgfqpoint{5.678892in}{0.697977in}}%
\pgfpathlineto{\pgfqpoint{5.682247in}{0.696452in}}%
\pgfpathlineto{\pgfqpoint{5.685603in}{0.669378in}}%
\pgfpathlineto{\pgfqpoint{5.688959in}{0.659736in}}%
\pgfpathlineto{\pgfqpoint{5.692315in}{0.708003in}}%
\pgfpathlineto{\pgfqpoint{5.695671in}{0.725273in}}%
\pgfpathlineto{\pgfqpoint{5.699026in}{0.728606in}}%
\pgfpathlineto{\pgfqpoint{5.702382in}{0.716905in}}%
\pgfpathlineto{\pgfqpoint{5.705738in}{0.692226in}}%
\pgfpathlineto{\pgfqpoint{5.709094in}{0.616535in}}%
\pgfpathlineto{\pgfqpoint{5.712449in}{0.666125in}}%
\pgfpathlineto{\pgfqpoint{5.715805in}{0.683304in}}%
\pgfpathlineto{\pgfqpoint{5.719161in}{0.679038in}}%
\pgfpathlineto{\pgfqpoint{5.722517in}{0.644345in}}%
\pgfpathlineto{\pgfqpoint{5.725873in}{0.656200in}}%
\pgfpathlineto{\pgfqpoint{5.729228in}{0.695999in}}%
\pgfpathlineto{\pgfqpoint{5.732584in}{0.711002in}}%
\pgfpathlineto{\pgfqpoint{5.735940in}{0.712226in}}%
\pgfpathlineto{\pgfqpoint{5.739296in}{0.696597in}}%
\pgfpathlineto{\pgfqpoint{5.742652in}{0.648402in}}%
\pgfpathlineto{\pgfqpoint{5.746007in}{0.665614in}}%
\pgfpathlineto{\pgfqpoint{5.749363in}{0.695956in}}%
\pgfpathlineto{\pgfqpoint{5.752719in}{0.705178in}}%
\pgfpathlineto{\pgfqpoint{5.756075in}{0.703746in}}%
\pgfpathlineto{\pgfqpoint{5.759430in}{0.690286in}}%
\pgfpathlineto{\pgfqpoint{5.762786in}{0.639602in}}%
\pgfpathlineto{\pgfqpoint{5.769498in}{0.707089in}}%
\pgfpathlineto{\pgfqpoint{5.772854in}{0.719777in}}%
\pgfpathlineto{\pgfqpoint{5.776209in}{0.720412in}}%
\pgfpathlineto{\pgfqpoint{5.779565in}{0.703799in}}%
\pgfpathlineto{\pgfqpoint{5.782921in}{0.619361in}}%
\pgfpathlineto{\pgfqpoint{5.786277in}{0.697341in}}%
\pgfpathlineto{\pgfqpoint{5.789632in}{0.721893in}}%
\pgfpathlineto{\pgfqpoint{5.792988in}{0.731228in}}%
\pgfpathlineto{\pgfqpoint{5.796344in}{0.732108in}}%
\pgfpathlineto{\pgfqpoint{5.799700in}{0.724664in}}%
\pgfpathlineto{\pgfqpoint{5.803056in}{0.702641in}}%
\pgfpathlineto{\pgfqpoint{5.806411in}{0.614683in}}%
\pgfpathlineto{\pgfqpoint{5.809767in}{0.712762in}}%
\pgfpathlineto{\pgfqpoint{5.813123in}{0.735451in}}%
\pgfpathlineto{\pgfqpoint{5.816479in}{0.744511in}}%
\pgfpathlineto{\pgfqpoint{5.819834in}{0.742582in}}%
\pgfpathlineto{\pgfqpoint{5.823190in}{0.727936in}}%
\pgfpathlineto{\pgfqpoint{5.826546in}{0.698184in}}%
\pgfpathlineto{\pgfqpoint{5.829902in}{0.599040in}}%
\pgfpathlineto{\pgfqpoint{5.833258in}{0.689611in}}%
\pgfpathlineto{\pgfqpoint{5.836613in}{0.700382in}}%
\pgfpathlineto{\pgfqpoint{5.839969in}{0.691702in}}%
\pgfpathlineto{\pgfqpoint{5.843325in}{0.638341in}}%
\pgfpathlineto{\pgfqpoint{5.846681in}{0.684894in}}%
\pgfpathlineto{\pgfqpoint{5.850036in}{0.716394in}}%
\pgfpathlineto{\pgfqpoint{5.853392in}{0.728365in}}%
\pgfpathlineto{\pgfqpoint{5.856748in}{0.726259in}}%
\pgfpathlineto{\pgfqpoint{5.860104in}{0.710290in}}%
\pgfpathlineto{\pgfqpoint{5.863460in}{0.677339in}}%
\pgfpathlineto{\pgfqpoint{5.866815in}{0.618888in}}%
\pgfpathlineto{\pgfqpoint{5.870171in}{0.677154in}}%
\pgfpathlineto{\pgfqpoint{5.873527in}{0.685245in}}%
\pgfpathlineto{\pgfqpoint{5.876883in}{0.673964in}}%
\pgfpathlineto{\pgfqpoint{5.880238in}{0.597912in}}%
\pgfpathlineto{\pgfqpoint{5.883594in}{0.675572in}}%
\pgfpathlineto{\pgfqpoint{5.886950in}{0.703503in}}%
\pgfpathlineto{\pgfqpoint{5.890306in}{0.713830in}}%
\pgfpathlineto{\pgfqpoint{5.893662in}{0.709709in}}%
\pgfpathlineto{\pgfqpoint{5.897017in}{0.687840in}}%
\pgfpathlineto{\pgfqpoint{5.900373in}{0.574891in}}%
\pgfpathlineto{\pgfqpoint{5.903729in}{0.680085in}}%
\pgfpathlineto{\pgfqpoint{5.907085in}{0.700637in}}%
\pgfpathlineto{\pgfqpoint{5.910441in}{0.706005in}}%
\pgfpathlineto{\pgfqpoint{5.913796in}{0.701044in}}%
\pgfpathlineto{\pgfqpoint{5.917152in}{0.680893in}}%
\pgfpathlineto{\pgfqpoint{5.920508in}{0.579137in}}%
\pgfpathlineto{\pgfqpoint{5.923864in}{0.690332in}}%
\pgfpathlineto{\pgfqpoint{5.927219in}{0.713364in}}%
\pgfpathlineto{\pgfqpoint{5.930575in}{0.721791in}}%
\pgfpathlineto{\pgfqpoint{5.933931in}{0.717303in}}%
\pgfpathlineto{\pgfqpoint{5.937287in}{0.691378in}}%
\pgfpathlineto{\pgfqpoint{5.940643in}{0.647165in}}%
\pgfpathlineto{\pgfqpoint{5.943998in}{0.708117in}}%
\pgfpathlineto{\pgfqpoint{5.947354in}{0.725685in}}%
\pgfpathlineto{\pgfqpoint{5.950710in}{0.731701in}}%
\pgfpathlineto{\pgfqpoint{5.954066in}{0.729680in}}%
\pgfpathlineto{\pgfqpoint{5.957421in}{0.718041in}}%
\pgfpathlineto{\pgfqpoint{5.960777in}{0.682930in}}%
\pgfpathlineto{\pgfqpoint{5.964133in}{0.680948in}}%
\pgfpathlineto{\pgfqpoint{5.967489in}{0.723471in}}%
\pgfpathlineto{\pgfqpoint{5.970845in}{0.739726in}}%
\pgfpathlineto{\pgfqpoint{5.974200in}{0.744632in}}%
\pgfpathlineto{\pgfqpoint{5.977556in}{0.737268in}}%
\pgfpathlineto{\pgfqpoint{5.980912in}{0.717736in}}%
\pgfpathlineto{\pgfqpoint{5.984268in}{0.672749in}}%
\pgfpathlineto{\pgfqpoint{5.987623in}{0.667367in}}%
\pgfpathlineto{\pgfqpoint{5.990979in}{0.697652in}}%
\pgfpathlineto{\pgfqpoint{5.994335in}{0.700159in}}%
\pgfpathlineto{\pgfqpoint{5.997691in}{0.682283in}}%
\pgfpathlineto{\pgfqpoint{6.001047in}{0.617234in}}%
\pgfpathlineto{\pgfqpoint{6.004402in}{0.700568in}}%
\pgfpathlineto{\pgfqpoint{6.007758in}{0.722630in}}%
\pgfpathlineto{\pgfqpoint{6.011114in}{0.729609in}}%
\pgfpathlineto{\pgfqpoint{6.014470in}{0.721712in}}%
\pgfpathlineto{\pgfqpoint{6.017825in}{0.701006in}}%
\pgfpathlineto{\pgfqpoint{6.021181in}{0.650868in}}%
\pgfpathlineto{\pgfqpoint{6.024537in}{0.658065in}}%
\pgfpathlineto{\pgfqpoint{6.027893in}{0.684218in}}%
\pgfpathlineto{\pgfqpoint{6.031249in}{0.685202in}}%
\pgfpathlineto{\pgfqpoint{6.034604in}{0.664757in}}%
\pgfpathlineto{\pgfqpoint{6.037960in}{0.622012in}}%
\pgfpathlineto{\pgfqpoint{6.041316in}{0.688939in}}%
\pgfpathlineto{\pgfqpoint{6.044672in}{0.709361in}}%
\pgfpathlineto{\pgfqpoint{6.048027in}{0.715216in}}%
\pgfpathlineto{\pgfqpoint{6.051383in}{0.705481in}}%
\pgfpathlineto{\pgfqpoint{6.054739in}{0.674803in}}%
\pgfpathlineto{\pgfqpoint{6.058095in}{0.637013in}}%
\pgfpathlineto{\pgfqpoint{6.061451in}{0.690046in}}%
\pgfpathlineto{\pgfqpoint{6.064806in}{0.704143in}}%
\pgfpathlineto{\pgfqpoint{6.068162in}{0.705993in}}%
\pgfpathlineto{\pgfqpoint{6.071518in}{0.697014in}}%
\pgfpathlineto{\pgfqpoint{6.074874in}{0.666095in}}%
\pgfpathlineto{\pgfqpoint{6.078230in}{0.654216in}}%
\pgfpathlineto{\pgfqpoint{6.081585in}{0.701267in}}%
\pgfpathlineto{\pgfqpoint{6.084941in}{0.718184in}}%
\pgfpathlineto{\pgfqpoint{6.088297in}{0.722540in}}%
\pgfpathlineto{\pgfqpoint{6.091653in}{0.712115in}}%
\pgfpathlineto{\pgfqpoint{6.095008in}{0.671203in}}%
\pgfpathlineto{\pgfqpoint{6.098364in}{0.680115in}}%
\pgfpathlineto{\pgfqpoint{6.101720in}{0.715511in}}%
\pgfpathlineto{\pgfqpoint{6.105076in}{0.728200in}}%
\pgfpathlineto{\pgfqpoint{6.108432in}{0.731232in}}%
\pgfpathlineto{\pgfqpoint{6.111787in}{0.726075in}}%
\pgfpathlineto{\pgfqpoint{6.115143in}{0.708735in}}%
\pgfpathlineto{\pgfqpoint{6.118499in}{0.634065in}}%
\pgfpathlineto{\pgfqpoint{6.121855in}{0.703368in}}%
\pgfpathlineto{\pgfqpoint{6.125210in}{0.731064in}}%
\pgfpathlineto{\pgfqpoint{6.128566in}{0.742441in}}%
\pgfpathlineto{\pgfqpoint{6.131922in}{0.742933in}}%
\pgfpathlineto{\pgfqpoint{6.135278in}{0.728762in}}%
\pgfpathlineto{\pgfqpoint{6.138634in}{0.684712in}}%
\pgfpathlineto{\pgfqpoint{6.141989in}{0.694114in}}%
\pgfpathlineto{\pgfqpoint{6.145345in}{0.727052in}}%
\pgfpathlineto{\pgfqpoint{6.148701in}{0.737418in}}%
\pgfpathlineto{\pgfqpoint{6.152057in}{0.737242in}}%
\pgfpathlineto{\pgfqpoint{6.155412in}{0.726166in}}%
\pgfpathlineto{\pgfqpoint{6.158768in}{0.688813in}}%
\pgfpathlineto{\pgfqpoint{6.162124in}{0.695686in}}%
\pgfpathlineto{\pgfqpoint{6.165480in}{0.735260in}}%
\pgfpathlineto{\pgfqpoint{6.168836in}{0.750383in}}%
\pgfpathlineto{\pgfqpoint{6.172191in}{0.753472in}}%
\pgfpathlineto{\pgfqpoint{6.175547in}{0.742955in}}%
\pgfpathlineto{\pgfqpoint{6.178903in}{0.719467in}}%
\pgfpathlineto{\pgfqpoint{6.182259in}{0.651047in}}%
\pgfpathlineto{\pgfqpoint{6.185614in}{0.691328in}}%
\pgfpathlineto{\pgfqpoint{6.188970in}{0.710656in}}%
\pgfpathlineto{\pgfqpoint{6.192326in}{0.708808in}}%
\pgfpathlineto{\pgfqpoint{6.195682in}{0.683266in}}%
\pgfpathlineto{\pgfqpoint{6.199038in}{0.665561in}}%
\pgfpathlineto{\pgfqpoint{6.202393in}{0.717422in}}%
\pgfpathlineto{\pgfqpoint{6.205749in}{0.735393in}}%
\pgfpathlineto{\pgfqpoint{6.209105in}{0.739386in}}%
\pgfpathlineto{\pgfqpoint{6.212461in}{0.728608in}}%
\pgfpathlineto{\pgfqpoint{6.215816in}{0.706759in}}%
\pgfpathlineto{\pgfqpoint{6.219172in}{0.657898in}}%
\pgfpathlineto{\pgfqpoint{6.222528in}{0.655602in}}%
\pgfpathlineto{\pgfqpoint{6.225884in}{0.680680in}}%
\pgfpathlineto{\pgfqpoint{6.229240in}{0.675599in}}%
\pgfpathlineto{\pgfqpoint{6.232595in}{0.620544in}}%
\pgfpathlineto{\pgfqpoint{6.235951in}{0.673775in}}%
\pgfpathlineto{\pgfqpoint{6.239307in}{0.705170in}}%
\pgfpathlineto{\pgfqpoint{6.242663in}{0.717258in}}%
\pgfpathlineto{\pgfqpoint{6.246019in}{0.714959in}}%
\pgfpathlineto{\pgfqpoint{6.249374in}{0.697176in}}%
\pgfpathlineto{\pgfqpoint{6.252730in}{0.651064in}}%
\pgfpathlineto{\pgfqpoint{6.256086in}{0.656548in}}%
\pgfpathlineto{\pgfqpoint{6.259442in}{0.687091in}}%
\pgfpathlineto{\pgfqpoint{6.262797in}{0.693729in}}%
\pgfpathlineto{\pgfqpoint{6.266153in}{0.686940in}}%
\pgfpathlineto{\pgfqpoint{6.269509in}{0.656903in}}%
\pgfpathlineto{\pgfqpoint{6.272865in}{0.647245in}}%
\pgfpathlineto{\pgfqpoint{6.276221in}{0.693990in}}%
\pgfpathlineto{\pgfqpoint{6.279576in}{0.710712in}}%
\pgfpathlineto{\pgfqpoint{6.282932in}{0.714170in}}%
\pgfpathlineto{\pgfqpoint{6.286288in}{0.701364in}}%
\pgfpathlineto{\pgfqpoint{6.289644in}{0.652852in}}%
\pgfpathlineto{\pgfqpoint{6.292999in}{0.676670in}}%
\pgfpathlineto{\pgfqpoint{6.296355in}{0.707726in}}%
\pgfpathlineto{\pgfqpoint{6.299711in}{0.718887in}}%
\pgfpathlineto{\pgfqpoint{6.303067in}{0.720641in}}%
\pgfpathlineto{\pgfqpoint{6.306423in}{0.713735in}}%
\pgfpathlineto{\pgfqpoint{6.309778in}{0.692198in}}%
\pgfpathlineto{\pgfqpoint{6.313134in}{0.599358in}}%
\pgfpathlineto{\pgfqpoint{6.316490in}{0.701850in}}%
\pgfpathlineto{\pgfqpoint{6.319846in}{0.724668in}}%
\pgfpathlineto{\pgfqpoint{6.323201in}{0.733589in}}%
\pgfpathlineto{\pgfqpoint{6.326557in}{0.731064in}}%
\pgfpathlineto{\pgfqpoint{6.329913in}{0.710378in}}%
\pgfpathlineto{\pgfqpoint{6.333269in}{0.578264in}}%
\pgfpathlineto{\pgfqpoint{6.336625in}{0.712287in}}%
\pgfpathlineto{\pgfqpoint{6.339980in}{0.733279in}}%
\pgfpathlineto{\pgfqpoint{6.343336in}{0.740729in}}%
\pgfpathlineto{\pgfqpoint{6.346692in}{0.739723in}}%
\pgfpathlineto{\pgfqpoint{6.350048in}{0.729298in}}%
\pgfpathlineto{\pgfqpoint{6.353403in}{0.698012in}}%
\pgfpathlineto{\pgfqpoint{6.356759in}{0.683457in}}%
\pgfpathlineto{\pgfqpoint{6.360115in}{0.731340in}}%
\pgfpathlineto{\pgfqpoint{6.363471in}{0.748740in}}%
\pgfpathlineto{\pgfqpoint{6.366827in}{0.754346in}}%
\pgfpathlineto{\pgfqpoint{6.370182in}{0.748002in}}%
\pgfpathlineto{\pgfqpoint{6.383605in}{0.705650in}}%
\pgfpathlineto{\pgfqpoint{6.383605in}{0.705650in}}%
\pgfusepath{stroke}%
\end{pgfscope}%
\begin{pgfscope}%
\pgfpathrectangle{\pgfqpoint{4.621821in}{0.320625in}}{\pgfqpoint{1.845679in}{1.164375in}}%
\pgfusepath{clip}%
\pgfsetrectcap%
\pgfsetroundjoin%
\pgfsetlinewidth{1.154312pt}%
\definecolor{currentstroke}{rgb}{0.478431,0.243137,0.615686}%
\pgfsetstrokecolor{currentstroke}%
\pgfsetdash{}{0pt}%
\pgfpathmoveto{\pgfqpoint{4.705715in}{1.468748in}}%
\pgfpathlineto{\pgfqpoint{4.715783in}{1.450185in}}%
\pgfpathlineto{\pgfqpoint{4.725850in}{1.429162in}}%
\pgfpathlineto{\pgfqpoint{4.739273in}{1.397858in}}%
\pgfpathlineto{\pgfqpoint{4.779543in}{1.301803in}}%
\pgfpathlineto{\pgfqpoint{4.799677in}{1.257107in}}%
\pgfpathlineto{\pgfqpoint{4.823168in}{1.207628in}}%
\pgfpathlineto{\pgfqpoint{4.850014in}{1.153585in}}%
\pgfpathlineto{\pgfqpoint{4.883572in}{1.088504in}}%
\pgfpathlineto{\pgfqpoint{4.927197in}{1.006508in}}%
\pgfpathlineto{\pgfqpoint{4.977534in}{0.914247in}}%
\pgfpathlineto{\pgfqpoint{5.024515in}{0.829849in}}%
\pgfpathlineto{\pgfqpoint{5.048005in}{0.788271in}}%
\pgfpathlineto{\pgfqpoint{5.058072in}{0.770693in}}%
\pgfpathlineto{\pgfqpoint{5.078207in}{0.733825in}}%
\pgfpathlineto{\pgfqpoint{5.094986in}{0.710118in}}%
\pgfpathlineto{\pgfqpoint{5.101698in}{0.696590in}}%
\pgfpathlineto{\pgfqpoint{5.108409in}{0.681229in}}%
\pgfpathlineto{\pgfqpoint{5.115121in}{0.664184in}}%
\pgfpathlineto{\pgfqpoint{5.118476in}{0.657028in}}%
\pgfpathlineto{\pgfqpoint{5.121832in}{0.652646in}}%
\pgfpathlineto{\pgfqpoint{5.125188in}{0.651760in}}%
\pgfpathlineto{\pgfqpoint{5.128544in}{0.653456in}}%
\pgfpathlineto{\pgfqpoint{5.135255in}{0.657966in}}%
\pgfpathlineto{\pgfqpoint{5.138611in}{0.658695in}}%
\pgfpathlineto{\pgfqpoint{5.141967in}{0.658804in}}%
\pgfpathlineto{\pgfqpoint{5.145323in}{0.657507in}}%
\pgfpathlineto{\pgfqpoint{5.148678in}{0.654006in}}%
\pgfpathlineto{\pgfqpoint{5.152034in}{0.648253in}}%
\pgfpathlineto{\pgfqpoint{5.155390in}{0.646296in}}%
\pgfpathlineto{\pgfqpoint{5.158746in}{0.657415in}}%
\pgfpathlineto{\pgfqpoint{5.162102in}{0.671737in}}%
\pgfpathlineto{\pgfqpoint{5.165457in}{0.682933in}}%
\pgfpathlineto{\pgfqpoint{5.168813in}{0.690751in}}%
\pgfpathlineto{\pgfqpoint{5.172169in}{0.695819in}}%
\pgfpathlineto{\pgfqpoint{5.175525in}{0.699083in}}%
\pgfpathlineto{\pgfqpoint{5.178880in}{0.701295in}}%
\pgfpathlineto{\pgfqpoint{5.182236in}{0.702742in}}%
\pgfpathlineto{\pgfqpoint{5.185592in}{0.703431in}}%
\pgfpathlineto{\pgfqpoint{5.188948in}{0.703281in}}%
\pgfpathlineto{\pgfqpoint{5.192304in}{0.702255in}}%
\pgfpathlineto{\pgfqpoint{5.202371in}{0.697283in}}%
\pgfpathlineto{\pgfqpoint{5.205727in}{0.698064in}}%
\pgfpathlineto{\pgfqpoint{5.209082in}{0.701583in}}%
\pgfpathlineto{\pgfqpoint{5.215794in}{0.711980in}}%
\pgfpathlineto{\pgfqpoint{5.219150in}{0.715766in}}%
\pgfpathlineto{\pgfqpoint{5.222506in}{0.717806in}}%
\pgfpathlineto{\pgfqpoint{5.225861in}{0.717924in}}%
\pgfpathlineto{\pgfqpoint{5.229217in}{0.717174in}}%
\pgfpathlineto{\pgfqpoint{5.232573in}{0.715499in}}%
\pgfpathlineto{\pgfqpoint{5.235929in}{0.711112in}}%
\pgfpathlineto{\pgfqpoint{5.239285in}{0.701888in}}%
\pgfpathlineto{\pgfqpoint{5.242640in}{0.689137in}}%
\pgfpathlineto{\pgfqpoint{5.245996in}{0.680012in}}%
\pgfpathlineto{\pgfqpoint{5.249352in}{0.687071in}}%
\pgfpathlineto{\pgfqpoint{5.252708in}{0.699156in}}%
\pgfpathlineto{\pgfqpoint{5.256063in}{0.707395in}}%
\pgfpathlineto{\pgfqpoint{5.259419in}{0.710084in}}%
\pgfpathlineto{\pgfqpoint{5.262775in}{0.705204in}}%
\pgfpathlineto{\pgfqpoint{5.266131in}{0.691756in}}%
\pgfpathlineto{\pgfqpoint{5.269487in}{0.686154in}}%
\pgfpathlineto{\pgfqpoint{5.272842in}{0.698151in}}%
\pgfpathlineto{\pgfqpoint{5.276198in}{0.707923in}}%
\pgfpathlineto{\pgfqpoint{5.279554in}{0.712967in}}%
\pgfpathlineto{\pgfqpoint{5.282910in}{0.712531in}}%
\pgfpathlineto{\pgfqpoint{5.286265in}{0.706235in}}%
\pgfpathlineto{\pgfqpoint{5.289621in}{0.694254in}}%
\pgfpathlineto{\pgfqpoint{5.292977in}{0.688043in}}%
\pgfpathlineto{\pgfqpoint{5.299689in}{0.713178in}}%
\pgfpathlineto{\pgfqpoint{5.303044in}{0.720119in}}%
\pgfpathlineto{\pgfqpoint{5.306400in}{0.722700in}}%
\pgfpathlineto{\pgfqpoint{5.309756in}{0.719106in}}%
\pgfpathlineto{\pgfqpoint{5.313112in}{0.712851in}}%
\pgfpathlineto{\pgfqpoint{5.316467in}{0.712125in}}%
\pgfpathlineto{\pgfqpoint{5.319823in}{0.714599in}}%
\pgfpathlineto{\pgfqpoint{5.323179in}{0.714840in}}%
\pgfpathlineto{\pgfqpoint{5.326535in}{0.710476in}}%
\pgfpathlineto{\pgfqpoint{5.329891in}{0.702061in}}%
\pgfpathlineto{\pgfqpoint{5.333246in}{0.698984in}}%
\pgfpathlineto{\pgfqpoint{5.339958in}{0.720270in}}%
\pgfpathlineto{\pgfqpoint{5.343314in}{0.726404in}}%
\pgfpathlineto{\pgfqpoint{5.346669in}{0.726302in}}%
\pgfpathlineto{\pgfqpoint{5.350025in}{0.721693in}}%
\pgfpathlineto{\pgfqpoint{5.356737in}{0.708012in}}%
\pgfpathlineto{\pgfqpoint{5.366804in}{0.682776in}}%
\pgfpathlineto{\pgfqpoint{5.370160in}{0.691260in}}%
\pgfpathlineto{\pgfqpoint{5.373516in}{0.707362in}}%
\pgfpathlineto{\pgfqpoint{5.376871in}{0.718154in}}%
\pgfpathlineto{\pgfqpoint{5.380227in}{0.720917in}}%
\pgfpathlineto{\pgfqpoint{5.383583in}{0.714301in}}%
\pgfpathlineto{\pgfqpoint{5.386939in}{0.704233in}}%
\pgfpathlineto{\pgfqpoint{5.390295in}{0.705488in}}%
\pgfpathlineto{\pgfqpoint{5.393650in}{0.713049in}}%
\pgfpathlineto{\pgfqpoint{5.397006in}{0.714742in}}%
\pgfpathlineto{\pgfqpoint{5.400362in}{0.710636in}}%
\pgfpathlineto{\pgfqpoint{5.403718in}{0.698352in}}%
\pgfpathlineto{\pgfqpoint{5.407074in}{0.677356in}}%
\pgfpathlineto{\pgfqpoint{5.410429in}{0.692502in}}%
\pgfpathlineto{\pgfqpoint{5.413785in}{0.715060in}}%
\pgfpathlineto{\pgfqpoint{5.417141in}{0.726199in}}%
\pgfpathlineto{\pgfqpoint{5.420497in}{0.727048in}}%
\pgfpathlineto{\pgfqpoint{5.423852in}{0.715587in}}%
\pgfpathlineto{\pgfqpoint{5.427208in}{0.699707in}}%
\pgfpathlineto{\pgfqpoint{5.430564in}{0.704871in}}%
\pgfpathlineto{\pgfqpoint{5.433920in}{0.715585in}}%
\pgfpathlineto{\pgfqpoint{5.437276in}{0.719502in}}%
\pgfpathlineto{\pgfqpoint{5.440631in}{0.717859in}}%
\pgfpathlineto{\pgfqpoint{5.443987in}{0.710674in}}%
\pgfpathlineto{\pgfqpoint{5.447343in}{0.701794in}}%
\pgfpathlineto{\pgfqpoint{5.450699in}{0.706679in}}%
\pgfpathlineto{\pgfqpoint{5.454054in}{0.719912in}}%
\pgfpathlineto{\pgfqpoint{5.457410in}{0.727845in}}%
\pgfpathlineto{\pgfqpoint{5.460766in}{0.728677in}}%
\pgfpathlineto{\pgfqpoint{5.464122in}{0.726627in}}%
\pgfpathlineto{\pgfqpoint{5.467478in}{0.721047in}}%
\pgfpathlineto{\pgfqpoint{5.470833in}{0.718530in}}%
\pgfpathlineto{\pgfqpoint{5.474189in}{0.719288in}}%
\pgfpathlineto{\pgfqpoint{5.477545in}{0.718274in}}%
\pgfpathlineto{\pgfqpoint{5.480901in}{0.715755in}}%
\pgfpathlineto{\pgfqpoint{5.484256in}{0.711711in}}%
\pgfpathlineto{\pgfqpoint{5.487612in}{0.708953in}}%
\pgfpathlineto{\pgfqpoint{5.490968in}{0.712193in}}%
\pgfpathlineto{\pgfqpoint{5.497680in}{0.726170in}}%
\pgfpathlineto{\pgfqpoint{5.501035in}{0.729324in}}%
\pgfpathlineto{\pgfqpoint{5.504391in}{0.727455in}}%
\pgfpathlineto{\pgfqpoint{5.507747in}{0.723853in}}%
\pgfpathlineto{\pgfqpoint{5.511103in}{0.719366in}}%
\pgfpathlineto{\pgfqpoint{5.514458in}{0.712336in}}%
\pgfpathlineto{\pgfqpoint{5.521170in}{0.679907in}}%
\pgfpathlineto{\pgfqpoint{5.524526in}{0.684807in}}%
\pgfpathlineto{\pgfqpoint{5.527882in}{0.699808in}}%
\pgfpathlineto{\pgfqpoint{5.531237in}{0.708649in}}%
\pgfpathlineto{\pgfqpoint{5.534593in}{0.711562in}}%
\pgfpathlineto{\pgfqpoint{5.541305in}{0.706378in}}%
\pgfpathlineto{\pgfqpoint{5.544660in}{0.707328in}}%
\pgfpathlineto{\pgfqpoint{5.548016in}{0.712692in}}%
\pgfpathlineto{\pgfqpoint{5.551372in}{0.714800in}}%
\pgfpathlineto{\pgfqpoint{5.554728in}{0.706432in}}%
\pgfpathlineto{\pgfqpoint{5.561439in}{0.667634in}}%
\pgfpathlineto{\pgfqpoint{5.564795in}{0.690888in}}%
\pgfpathlineto{\pgfqpoint{5.568151in}{0.707298in}}%
\pgfpathlineto{\pgfqpoint{5.571507in}{0.712764in}}%
\pgfpathlineto{\pgfqpoint{5.574863in}{0.707706in}}%
\pgfpathlineto{\pgfqpoint{5.578218in}{0.695505in}}%
\pgfpathlineto{\pgfqpoint{5.581574in}{0.691260in}}%
\pgfpathlineto{\pgfqpoint{5.584930in}{0.708363in}}%
\pgfpathlineto{\pgfqpoint{5.588286in}{0.720147in}}%
\pgfpathlineto{\pgfqpoint{5.591641in}{0.723147in}}%
\pgfpathlineto{\pgfqpoint{5.594997in}{0.720036in}}%
\pgfpathlineto{\pgfqpoint{5.598353in}{0.711913in}}%
\pgfpathlineto{\pgfqpoint{5.601709in}{0.702162in}}%
\pgfpathlineto{\pgfqpoint{5.605065in}{0.705411in}}%
\pgfpathlineto{\pgfqpoint{5.608420in}{0.718280in}}%
\pgfpathlineto{\pgfqpoint{5.611776in}{0.726906in}}%
\pgfpathlineto{\pgfqpoint{5.615132in}{0.728264in}}%
\pgfpathlineto{\pgfqpoint{5.625199in}{0.719666in}}%
\pgfpathlineto{\pgfqpoint{5.628555in}{0.721286in}}%
\pgfpathlineto{\pgfqpoint{5.631911in}{0.721615in}}%
\pgfpathlineto{\pgfqpoint{5.635267in}{0.718559in}}%
\pgfpathlineto{\pgfqpoint{5.641978in}{0.704946in}}%
\pgfpathlineto{\pgfqpoint{5.645334in}{0.713479in}}%
\pgfpathlineto{\pgfqpoint{5.648690in}{0.726525in}}%
\pgfpathlineto{\pgfqpoint{5.652045in}{0.733572in}}%
\pgfpathlineto{\pgfqpoint{5.655401in}{0.732376in}}%
\pgfpathlineto{\pgfqpoint{5.665469in}{0.716421in}}%
\pgfpathlineto{\pgfqpoint{5.668824in}{0.709829in}}%
\pgfpathlineto{\pgfqpoint{5.672180in}{0.696397in}}%
\pgfpathlineto{\pgfqpoint{5.675536in}{0.674827in}}%
\pgfpathlineto{\pgfqpoint{5.678892in}{0.680155in}}%
\pgfpathlineto{\pgfqpoint{5.682247in}{0.703091in}}%
\pgfpathlineto{\pgfqpoint{5.685603in}{0.715941in}}%
\pgfpathlineto{\pgfqpoint{5.688959in}{0.719440in}}%
\pgfpathlineto{\pgfqpoint{5.692315in}{0.713985in}}%
\pgfpathlineto{\pgfqpoint{5.695671in}{0.705390in}}%
\pgfpathlineto{\pgfqpoint{5.699026in}{0.702005in}}%
\pgfpathlineto{\pgfqpoint{5.702382in}{0.705375in}}%
\pgfpathlineto{\pgfqpoint{5.705738in}{0.712420in}}%
\pgfpathlineto{\pgfqpoint{5.709094in}{0.712891in}}%
\pgfpathlineto{\pgfqpoint{5.712449in}{0.701552in}}%
\pgfpathlineto{\pgfqpoint{5.715805in}{0.682673in}}%
\pgfpathlineto{\pgfqpoint{5.719161in}{0.685908in}}%
\pgfpathlineto{\pgfqpoint{5.722517in}{0.701690in}}%
\pgfpathlineto{\pgfqpoint{5.725873in}{0.708112in}}%
\pgfpathlineto{\pgfqpoint{5.729228in}{0.707190in}}%
\pgfpathlineto{\pgfqpoint{5.732584in}{0.702569in}}%
\pgfpathlineto{\pgfqpoint{5.735940in}{0.694797in}}%
\pgfpathlineto{\pgfqpoint{5.739296in}{0.694097in}}%
\pgfpathlineto{\pgfqpoint{5.742652in}{0.708524in}}%
\pgfpathlineto{\pgfqpoint{5.746007in}{0.718820in}}%
\pgfpathlineto{\pgfqpoint{5.749363in}{0.720313in}}%
\pgfpathlineto{\pgfqpoint{5.756075in}{0.709740in}}%
\pgfpathlineto{\pgfqpoint{5.759430in}{0.706877in}}%
\pgfpathlineto{\pgfqpoint{5.766142in}{0.706215in}}%
\pgfpathlineto{\pgfqpoint{5.769498in}{0.708267in}}%
\pgfpathlineto{\pgfqpoint{5.776209in}{0.724411in}}%
\pgfpathlineto{\pgfqpoint{5.779565in}{0.726262in}}%
\pgfpathlineto{\pgfqpoint{5.782921in}{0.725262in}}%
\pgfpathlineto{\pgfqpoint{5.786277in}{0.722384in}}%
\pgfpathlineto{\pgfqpoint{5.792988in}{0.707852in}}%
\pgfpathlineto{\pgfqpoint{5.796344in}{0.707253in}}%
\pgfpathlineto{\pgfqpoint{5.803056in}{0.723473in}}%
\pgfpathlineto{\pgfqpoint{5.806411in}{0.727164in}}%
\pgfpathlineto{\pgfqpoint{5.809767in}{0.725169in}}%
\pgfpathlineto{\pgfqpoint{5.813123in}{0.721992in}}%
\pgfpathlineto{\pgfqpoint{5.819834in}{0.719531in}}%
\pgfpathlineto{\pgfqpoint{5.823190in}{0.713318in}}%
\pgfpathlineto{\pgfqpoint{5.826546in}{0.700331in}}%
\pgfpathlineto{\pgfqpoint{5.829902in}{0.676423in}}%
\pgfpathlineto{\pgfqpoint{5.833258in}{0.670612in}}%
\pgfpathlineto{\pgfqpoint{5.836613in}{0.696948in}}%
\pgfpathlineto{\pgfqpoint{5.839969in}{0.711227in}}%
\pgfpathlineto{\pgfqpoint{5.843325in}{0.715210in}}%
\pgfpathlineto{\pgfqpoint{5.846681in}{0.709607in}}%
\pgfpathlineto{\pgfqpoint{5.850036in}{0.702396in}}%
\pgfpathlineto{\pgfqpoint{5.853392in}{0.702109in}}%
\pgfpathlineto{\pgfqpoint{5.856748in}{0.704419in}}%
\pgfpathlineto{\pgfqpoint{5.860104in}{0.705374in}}%
\pgfpathlineto{\pgfqpoint{5.863460in}{0.703106in}}%
\pgfpathlineto{\pgfqpoint{5.870171in}{0.678813in}}%
\pgfpathlineto{\pgfqpoint{5.873527in}{0.689477in}}%
\pgfpathlineto{\pgfqpoint{5.876883in}{0.703018in}}%
\pgfpathlineto{\pgfqpoint{5.880238in}{0.707899in}}%
\pgfpathlineto{\pgfqpoint{5.883594in}{0.705016in}}%
\pgfpathlineto{\pgfqpoint{5.886950in}{0.701146in}}%
\pgfpathlineto{\pgfqpoint{5.893662in}{0.698760in}}%
\pgfpathlineto{\pgfqpoint{5.897017in}{0.703644in}}%
\pgfpathlineto{\pgfqpoint{5.900373in}{0.712361in}}%
\pgfpathlineto{\pgfqpoint{5.903729in}{0.716700in}}%
\pgfpathlineto{\pgfqpoint{5.907085in}{0.714348in}}%
\pgfpathlineto{\pgfqpoint{5.910441in}{0.711031in}}%
\pgfpathlineto{\pgfqpoint{5.913796in}{0.711528in}}%
\pgfpathlineto{\pgfqpoint{5.917152in}{0.713948in}}%
\pgfpathlineto{\pgfqpoint{5.920508in}{0.713196in}}%
\pgfpathlineto{\pgfqpoint{5.923864in}{0.709137in}}%
\pgfpathlineto{\pgfqpoint{5.927219in}{0.712183in}}%
\pgfpathlineto{\pgfqpoint{5.930575in}{0.721209in}}%
\pgfpathlineto{\pgfqpoint{5.933931in}{0.726802in}}%
\pgfpathlineto{\pgfqpoint{5.937287in}{0.727156in}}%
\pgfpathlineto{\pgfqpoint{5.940643in}{0.726595in}}%
\pgfpathlineto{\pgfqpoint{5.943998in}{0.723653in}}%
\pgfpathlineto{\pgfqpoint{5.947354in}{0.714808in}}%
\pgfpathlineto{\pgfqpoint{5.950710in}{0.708642in}}%
\pgfpathlineto{\pgfqpoint{5.954066in}{0.715304in}}%
\pgfpathlineto{\pgfqpoint{5.957421in}{0.724672in}}%
\pgfpathlineto{\pgfqpoint{5.960777in}{0.728661in}}%
\pgfpathlineto{\pgfqpoint{5.964133in}{0.724659in}}%
\pgfpathlineto{\pgfqpoint{5.967489in}{0.717623in}}%
\pgfpathlineto{\pgfqpoint{5.970845in}{0.718378in}}%
\pgfpathlineto{\pgfqpoint{5.974200in}{0.723340in}}%
\pgfpathlineto{\pgfqpoint{5.977556in}{0.725560in}}%
\pgfpathlineto{\pgfqpoint{5.980912in}{0.723816in}}%
\pgfpathlineto{\pgfqpoint{5.984268in}{0.713902in}}%
\pgfpathlineto{\pgfqpoint{5.987623in}{0.691222in}}%
\pgfpathlineto{\pgfqpoint{5.990979in}{0.686500in}}%
\pgfpathlineto{\pgfqpoint{5.994335in}{0.712290in}}%
\pgfpathlineto{\pgfqpoint{5.997691in}{0.725379in}}%
\pgfpathlineto{\pgfqpoint{6.001047in}{0.727187in}}%
\pgfpathlineto{\pgfqpoint{6.004402in}{0.718283in}}%
\pgfpathlineto{\pgfqpoint{6.007758in}{0.706668in}}%
\pgfpathlineto{\pgfqpoint{6.011114in}{0.702961in}}%
\pgfpathlineto{\pgfqpoint{6.017825in}{0.708220in}}%
\pgfpathlineto{\pgfqpoint{6.021181in}{0.703798in}}%
\pgfpathlineto{\pgfqpoint{6.024537in}{0.689659in}}%
\pgfpathlineto{\pgfqpoint{6.027893in}{0.683359in}}%
\pgfpathlineto{\pgfqpoint{6.031249in}{0.700521in}}%
\pgfpathlineto{\pgfqpoint{6.034604in}{0.712654in}}%
\pgfpathlineto{\pgfqpoint{6.037960in}{0.714292in}}%
\pgfpathlineto{\pgfqpoint{6.041316in}{0.706800in}}%
\pgfpathlineto{\pgfqpoint{6.044672in}{0.701415in}}%
\pgfpathlineto{\pgfqpoint{6.048027in}{0.700326in}}%
\pgfpathlineto{\pgfqpoint{6.051383in}{0.698263in}}%
\pgfpathlineto{\pgfqpoint{6.054739in}{0.700169in}}%
\pgfpathlineto{\pgfqpoint{6.058095in}{0.703793in}}%
\pgfpathlineto{\pgfqpoint{6.061451in}{0.704952in}}%
\pgfpathlineto{\pgfqpoint{6.064806in}{0.708074in}}%
\pgfpathlineto{\pgfqpoint{6.068162in}{0.712675in}}%
\pgfpathlineto{\pgfqpoint{6.071518in}{0.714933in}}%
\pgfpathlineto{\pgfqpoint{6.074874in}{0.710914in}}%
\pgfpathlineto{\pgfqpoint{6.078230in}{0.699849in}}%
\pgfpathlineto{\pgfqpoint{6.081585in}{0.698993in}}%
\pgfpathlineto{\pgfqpoint{6.084941in}{0.713584in}}%
\pgfpathlineto{\pgfqpoint{6.088297in}{0.723586in}}%
\pgfpathlineto{\pgfqpoint{6.091653in}{0.725246in}}%
\pgfpathlineto{\pgfqpoint{6.095008in}{0.720693in}}%
\pgfpathlineto{\pgfqpoint{6.098364in}{0.714161in}}%
\pgfpathlineto{\pgfqpoint{6.101720in}{0.705565in}}%
\pgfpathlineto{\pgfqpoint{6.105076in}{0.704624in}}%
\pgfpathlineto{\pgfqpoint{6.108432in}{0.713333in}}%
\pgfpathlineto{\pgfqpoint{6.111787in}{0.720438in}}%
\pgfpathlineto{\pgfqpoint{6.115143in}{0.722225in}}%
\pgfpathlineto{\pgfqpoint{6.118499in}{0.718488in}}%
\pgfpathlineto{\pgfqpoint{6.121855in}{0.712612in}}%
\pgfpathlineto{\pgfqpoint{6.125210in}{0.716224in}}%
\pgfpathlineto{\pgfqpoint{6.128566in}{0.723618in}}%
\pgfpathlineto{\pgfqpoint{6.131922in}{0.727862in}}%
\pgfpathlineto{\pgfqpoint{6.135278in}{0.727205in}}%
\pgfpathlineto{\pgfqpoint{6.138634in}{0.722922in}}%
\pgfpathlineto{\pgfqpoint{6.141989in}{0.713671in}}%
\pgfpathlineto{\pgfqpoint{6.145345in}{0.709750in}}%
\pgfpathlineto{\pgfqpoint{6.152057in}{0.729648in}}%
\pgfpathlineto{\pgfqpoint{6.155412in}{0.732229in}}%
\pgfpathlineto{\pgfqpoint{6.158768in}{0.726264in}}%
\pgfpathlineto{\pgfqpoint{6.162124in}{0.714617in}}%
\pgfpathlineto{\pgfqpoint{6.165480in}{0.716309in}}%
\pgfpathlineto{\pgfqpoint{6.168836in}{0.727839in}}%
\pgfpathlineto{\pgfqpoint{6.172191in}{0.733845in}}%
\pgfpathlineto{\pgfqpoint{6.175547in}{0.731439in}}%
\pgfpathlineto{\pgfqpoint{6.178903in}{0.722724in}}%
\pgfpathlineto{\pgfqpoint{6.182259in}{0.704555in}}%
\pgfpathlineto{\pgfqpoint{6.185614in}{0.703005in}}%
\pgfpathlineto{\pgfqpoint{6.188970in}{0.721251in}}%
\pgfpathlineto{\pgfqpoint{6.192326in}{0.730820in}}%
\pgfpathlineto{\pgfqpoint{6.195682in}{0.729080in}}%
\pgfpathlineto{\pgfqpoint{6.202393in}{0.705833in}}%
\pgfpathlineto{\pgfqpoint{6.205749in}{0.710592in}}%
\pgfpathlineto{\pgfqpoint{6.209105in}{0.717858in}}%
\pgfpathlineto{\pgfqpoint{6.212461in}{0.718937in}}%
\pgfpathlineto{\pgfqpoint{6.215816in}{0.716935in}}%
\pgfpathlineto{\pgfqpoint{6.219172in}{0.710261in}}%
\pgfpathlineto{\pgfqpoint{6.222528in}{0.708136in}}%
\pgfpathlineto{\pgfqpoint{6.229240in}{0.721705in}}%
\pgfpathlineto{\pgfqpoint{6.232595in}{0.721347in}}%
\pgfpathlineto{\pgfqpoint{6.235951in}{0.711996in}}%
\pgfpathlineto{\pgfqpoint{6.239307in}{0.708289in}}%
\pgfpathlineto{\pgfqpoint{6.242663in}{0.715710in}}%
\pgfpathlineto{\pgfqpoint{6.246019in}{0.720232in}}%
\pgfpathlineto{\pgfqpoint{6.249374in}{0.716187in}}%
\pgfpathlineto{\pgfqpoint{6.252730in}{0.706007in}}%
\pgfpathlineto{\pgfqpoint{6.256086in}{0.688889in}}%
\pgfpathlineto{\pgfqpoint{6.259442in}{0.683841in}}%
\pgfpathlineto{\pgfqpoint{6.262797in}{0.700121in}}%
\pgfpathlineto{\pgfqpoint{6.266153in}{0.708665in}}%
\pgfpathlineto{\pgfqpoint{6.269509in}{0.704920in}}%
\pgfpathlineto{\pgfqpoint{6.272865in}{0.692328in}}%
\pgfpathlineto{\pgfqpoint{6.276221in}{0.691882in}}%
\pgfpathlineto{\pgfqpoint{6.279576in}{0.704822in}}%
\pgfpathlineto{\pgfqpoint{6.282932in}{0.711809in}}%
\pgfpathlineto{\pgfqpoint{6.286288in}{0.709640in}}%
\pgfpathlineto{\pgfqpoint{6.292999in}{0.700660in}}%
\pgfpathlineto{\pgfqpoint{6.296355in}{0.697196in}}%
\pgfpathlineto{\pgfqpoint{6.299711in}{0.702492in}}%
\pgfpathlineto{\pgfqpoint{6.303067in}{0.711225in}}%
\pgfpathlineto{\pgfqpoint{6.306423in}{0.714917in}}%
\pgfpathlineto{\pgfqpoint{6.309778in}{0.711719in}}%
\pgfpathlineto{\pgfqpoint{6.313134in}{0.702977in}}%
\pgfpathlineto{\pgfqpoint{6.316490in}{0.702257in}}%
\pgfpathlineto{\pgfqpoint{6.319846in}{0.712141in}}%
\pgfpathlineto{\pgfqpoint{6.323201in}{0.719937in}}%
\pgfpathlineto{\pgfqpoint{6.326557in}{0.722810in}}%
\pgfpathlineto{\pgfqpoint{6.329913in}{0.721112in}}%
\pgfpathlineto{\pgfqpoint{6.336625in}{0.713949in}}%
\pgfpathlineto{\pgfqpoint{6.339980in}{0.716667in}}%
\pgfpathlineto{\pgfqpoint{6.343336in}{0.725478in}}%
\pgfpathlineto{\pgfqpoint{6.346692in}{0.731392in}}%
\pgfpathlineto{\pgfqpoint{6.350048in}{0.731416in}}%
\pgfpathlineto{\pgfqpoint{6.356759in}{0.714521in}}%
\pgfpathlineto{\pgfqpoint{6.360115in}{0.714540in}}%
\pgfpathlineto{\pgfqpoint{6.363471in}{0.721721in}}%
\pgfpathlineto{\pgfqpoint{6.366827in}{0.727557in}}%
\pgfpathlineto{\pgfqpoint{6.370182in}{0.726907in}}%
\pgfpathlineto{\pgfqpoint{6.373538in}{0.725565in}}%
\pgfpathlineto{\pgfqpoint{6.380250in}{0.715956in}}%
\pgfpathlineto{\pgfqpoint{6.383605in}{0.709220in}}%
\pgfpathlineto{\pgfqpoint{6.383605in}{0.709220in}}%
\pgfusepath{stroke}%
\end{pgfscope}%
\begin{pgfscope}%
\pgfsetrectcap%
\pgfsetmiterjoin%
\pgfsetlinewidth{0.803000pt}%
\definecolor{currentstroke}{rgb}{0.000000,0.000000,0.000000}%
\pgfsetstrokecolor{currentstroke}%
\pgfsetdash{}{0pt}%
\pgfpathmoveto{\pgfqpoint{4.621821in}{0.320625in}}%
\pgfpathlineto{\pgfqpoint{4.621821in}{1.485000in}}%
\pgfusepath{stroke}%
\end{pgfscope}%
\begin{pgfscope}%
\pgfsetrectcap%
\pgfsetmiterjoin%
\pgfsetlinewidth{0.803000pt}%
\definecolor{currentstroke}{rgb}{0.000000,0.000000,0.000000}%
\pgfsetstrokecolor{currentstroke}%
\pgfsetdash{}{0pt}%
\pgfpathmoveto{\pgfqpoint{6.467500in}{0.320625in}}%
\pgfpathlineto{\pgfqpoint{6.467500in}{1.485000in}}%
\pgfusepath{stroke}%
\end{pgfscope}%
\begin{pgfscope}%
\pgfsetrectcap%
\pgfsetmiterjoin%
\pgfsetlinewidth{0.803000pt}%
\definecolor{currentstroke}{rgb}{0.000000,0.000000,0.000000}%
\pgfsetstrokecolor{currentstroke}%
\pgfsetdash{}{0pt}%
\pgfpathmoveto{\pgfqpoint{4.621821in}{0.320625in}}%
\pgfpathlineto{\pgfqpoint{6.467500in}{0.320625in}}%
\pgfusepath{stroke}%
\end{pgfscope}%
\begin{pgfscope}%
\pgfsetrectcap%
\pgfsetmiterjoin%
\pgfsetlinewidth{0.803000pt}%
\definecolor{currentstroke}{rgb}{0.000000,0.000000,0.000000}%
\pgfsetstrokecolor{currentstroke}%
\pgfsetdash{}{0pt}%
\pgfpathmoveto{\pgfqpoint{4.621821in}{1.485000in}}%
\pgfpathlineto{\pgfqpoint{6.467500in}{1.485000in}}%
\pgfusepath{stroke}%
\end{pgfscope}%
\begin{pgfscope}%
\definecolor{textcolor}{rgb}{0.000000,0.000000,0.000000}%
\pgfsetstrokecolor{textcolor}%
\pgfsetfillcolor{textcolor}%
\pgftext[x=5.544660in,y=1.568333in,,base]{\color{textcolor}{\rmfamily\fontsize{8.000000}{9.600000}\selectfont\catcode`\^=\active\def^{\ifmmode\sp\else\^{}\fi}\catcode`\%=\active\def%{\%}(c) Finite width}}%
\end{pgfscope}%
\begin{pgfscope}%
\pgfsetbuttcap%
\pgfsetmiterjoin%
\definecolor{currentfill}{rgb}{1.000000,1.000000,1.000000}%
\pgfsetfillcolor{currentfill}%
\pgfsetfillopacity{0.800000}%
\pgfsetlinewidth{1.003750pt}%
\definecolor{currentstroke}{rgb}{0.800000,0.800000,0.800000}%
\pgfsetstrokecolor{currentstroke}%
\pgfsetstrokeopacity{0.800000}%
\pgfsetdash{}{0pt}%
\pgfpathmoveto{\pgfqpoint{5.911892in}{1.084181in}}%
\pgfpathlineto{\pgfqpoint{6.424722in}{1.084181in}}%
\pgfpathquadraticcurveto{\pgfqpoint{6.436944in}{1.084181in}}{\pgfqpoint{6.436944in}{1.096403in}}%
\pgfpathlineto{\pgfqpoint{6.436944in}{1.442222in}}%
\pgfpathquadraticcurveto{\pgfqpoint{6.436944in}{1.454444in}}{\pgfqpoint{6.424722in}{1.454444in}}%
\pgfpathlineto{\pgfqpoint{5.911892in}{1.454444in}}%
\pgfpathquadraticcurveto{\pgfqpoint{5.899669in}{1.454444in}}{\pgfqpoint{5.899669in}{1.442222in}}%
\pgfpathlineto{\pgfqpoint{5.899669in}{1.096403in}}%
\pgfpathquadraticcurveto{\pgfqpoint{5.899669in}{1.084181in}}{\pgfqpoint{5.911892in}{1.084181in}}%
\pgfpathlineto{\pgfqpoint{5.911892in}{1.084181in}}%
\pgfpathclose%
\pgfusepath{stroke,fill}%
\end{pgfscope}%
\begin{pgfscope}%
\pgfsetrectcap%
\pgfsetroundjoin%
\pgfsetlinewidth{1.304875pt}%
\definecolor{currentstroke}{rgb}{0.000000,0.000000,0.000000}%
\pgfsetstrokecolor{currentstroke}%
\pgfsetdash{}{0pt}%
\pgfpathmoveto{\pgfqpoint{5.914947in}{1.414125in}}%
\pgfpathlineto{\pgfqpoint{5.954669in}{1.414125in}}%
\pgfpathlineto{\pgfqpoint{5.994392in}{1.414125in}}%
\pgfusepath{stroke}%
\end{pgfscope}%
\begin{pgfscope}%
\definecolor{textcolor}{rgb}{0.000000,0.000000,0.000000}%
\pgfsetstrokecolor{textcolor}%
\pgfsetfillcolor{textcolor}%
\pgftext[x=6.043281in,y=1.392737in,left,base]{\color{textcolor}{\rmfamily\fontsize{4.400000}{5.280000}\selectfont\catcode`\^=\active\def^{\ifmmode\sp\else\^{}\fi}\catcode`\%=\active\def%{\%}mean field}}%
\end{pgfscope}%
\begin{pgfscope}%
\pgfsetrectcap%
\pgfsetroundjoin%
\pgfsetlinewidth{1.154312pt}%
\definecolor{currentstroke}{rgb}{0.141176,0.345098,0.650980}%
\pgfsetstrokecolor{currentstroke}%
\pgfsetdash{}{0pt}%
\pgfpathmoveto{\pgfqpoint{5.914947in}{1.343984in}}%
\pgfpathlineto{\pgfqpoint{5.954669in}{1.343984in}}%
\pgfpathlineto{\pgfqpoint{5.994392in}{1.343984in}}%
\pgfusepath{stroke}%
\end{pgfscope}%
\begin{pgfscope}%
\definecolor{textcolor}{rgb}{0.000000,0.000000,0.000000}%
\pgfsetstrokecolor{textcolor}%
\pgfsetfillcolor{textcolor}%
\pgftext[x=6.043281in,y=1.322595in,left,base]{\color{textcolor}{\rmfamily\fontsize{4.400000}{5.280000}\selectfont\catcode`\^=\active\def^{\ifmmode\sp\else\^{}\fi}\catcode`\%=\active\def%{\%}$n=32$}}%
\end{pgfscope}%
\begin{pgfscope}%
\pgfsetrectcap%
\pgfsetroundjoin%
\pgfsetlinewidth{1.154312pt}%
\definecolor{currentstroke}{rgb}{0.819608,0.286275,0.356863}%
\pgfsetstrokecolor{currentstroke}%
\pgfsetdash{}{0pt}%
\pgfpathmoveto{\pgfqpoint{5.914947in}{1.273842in}}%
\pgfpathlineto{\pgfqpoint{5.954669in}{1.273842in}}%
\pgfpathlineto{\pgfqpoint{5.994392in}{1.273842in}}%
\pgfusepath{stroke}%
\end{pgfscope}%
\begin{pgfscope}%
\definecolor{textcolor}{rgb}{0.000000,0.000000,0.000000}%
\pgfsetstrokecolor{textcolor}%
\pgfsetfillcolor{textcolor}%
\pgftext[x=6.043281in,y=1.252453in,left,base]{\color{textcolor}{\rmfamily\fontsize{4.400000}{5.280000}\selectfont\catcode`\^=\active\def^{\ifmmode\sp\else\^{}\fi}\catcode`\%=\active\def%{\%}$n=128$}}%
\end{pgfscope}%
\begin{pgfscope}%
\pgfsetrectcap%
\pgfsetroundjoin%
\pgfsetlinewidth{1.154312pt}%
\definecolor{currentstroke}{rgb}{0.000000,0.529412,0.423529}%
\pgfsetstrokecolor{currentstroke}%
\pgfsetdash{}{0pt}%
\pgfpathmoveto{\pgfqpoint{5.914947in}{1.203701in}}%
\pgfpathlineto{\pgfqpoint{5.954669in}{1.203701in}}%
\pgfpathlineto{\pgfqpoint{5.994392in}{1.203701in}}%
\pgfusepath{stroke}%
\end{pgfscope}%
\begin{pgfscope}%
\definecolor{textcolor}{rgb}{0.000000,0.000000,0.000000}%
\pgfsetstrokecolor{textcolor}%
\pgfsetfillcolor{textcolor}%
\pgftext[x=6.043281in,y=1.182312in,left,base]{\color{textcolor}{\rmfamily\fontsize{4.400000}{5.280000}\selectfont\catcode`\^=\active\def^{\ifmmode\sp\else\^{}\fi}\catcode`\%=\active\def%{\%}$n=512$}}%
\end{pgfscope}%
\begin{pgfscope}%
\pgfsetrectcap%
\pgfsetroundjoin%
\pgfsetlinewidth{1.154312pt}%
\definecolor{currentstroke}{rgb}{0.478431,0.243137,0.615686}%
\pgfsetstrokecolor{currentstroke}%
\pgfsetdash{}{0pt}%
\pgfpathmoveto{\pgfqpoint{5.914947in}{1.133559in}}%
\pgfpathlineto{\pgfqpoint{5.954669in}{1.133559in}}%
\pgfpathlineto{\pgfqpoint{5.994392in}{1.133559in}}%
\pgfusepath{stroke}%
\end{pgfscope}%
\begin{pgfscope}%
\definecolor{textcolor}{rgb}{0.000000,0.000000,0.000000}%
\pgfsetstrokecolor{textcolor}%
\pgfsetfillcolor{textcolor}%
\pgftext[x=6.043281in,y=1.112170in,left,base]{\color{textcolor}{\rmfamily\fontsize{4.400000}{5.280000}\selectfont\catcode`\^=\active\def^{\ifmmode\sp\else\^{}\fi}\catcode`\%=\active\def%{\%}$n=2048$}}%
\end{pgfscope}%
\end{pgfpicture}%
\makeatother%
\endgroup%
}
    \captionsetup{font=footnotesize,skip=2pt}
    \caption{
      Numerical convergence for quadratic loss with $d_u=16$, $d_v=9$, a random teacher, and whitened input.
      \textbf{(a)} Bures flow from $\Sigma_0=\alpha\Id_m$ at several white scales. Dashed lines show the spectral-margin rate $4\alpha\kappa$ from Theorem~\ref{thm:balanced-convergence}; dotted lines show the coarse rate $4\kappa c_{\rm coarse}$ from the determinant comparison in Appendix~\ref{app:gf} (Remark~\ref{rem:coarse-rate}), using the exact initial sublevel radius $M=2\norm{\Svu_\star}_{\F}$.
      \textbf{(b)} Rate deterioration when one matched diagonal pair of $\Sigma_0$ is scaled by $r$: $r=1$ is the white reference and $r=0$ is rank-deficient.
      \textbf{(c)} Bures flow from the reference covariance $\tSigma_0=\Id_m$ compared with finite-width flows initialized by i.i.d. samples from $\mathcal N(0,\tSigma_0)$, for $n=32,128,512,2048$ and five independent trials.
    }
    \label{fig:bures-convergence-triptych}
    \label{fig:bures-scale-sweep}
    \label{fig:perturb-scale-sweep}
    \label{fig:bures-finite-width-comparison}
\end{figure}

\paragraph{Convergence under finite width sampling}
Theorem~\ref{thm:gf} applies generally to neuron law covariances with positive support gap.
We now focus on finite-width networks, whose initial empirical covariance is $\Sigma_0:=W_0W_0^\top/n$.
For random initialization, the goal is to make this
empirical covariance close to a prescribed population covariance
$\tSigma_0$, thereby retaining its support gap and convergence rate.
Draw the columns $w_i:=(u_i,v_i)$ of $W_0:=\binom{U_0}{V_0}$ independently
from the same \emph{subgaussian} distribution:
$w_i:=\tSigma_0^{1/2}\xi_i$, where $\mathbb E[\xi_i\xi_i^\top]=\Id_m$ and $\sup_{\norm z=1}\norm{\langle z,\xi_i\rangle}_{\psi_2}\le K$.
Here $\norm X_{\psi_2}:=\inf\{c>0:\mathbb E e^{X^2/c^2}\le2\}$ denotes the standard subgaussian norm \citep{Vershynin2018};
standard Gaussian sampling has dimension-independent $K$.
Thus $\mathbb E[w_iw_i^\top]=\tSigma_0$, whereas $\Sigma_0=W_0W_0^\top/n$ is the empirical covariance of finitely many sampled neurons.
The loss gap along the sampled trajectory remains $\Delta_t=\ell(\Svu_t)-\ell_\star$ and we write $\widetilde\Delta_0:=\ell(\tSvu_0)-\ell_\star$ for its population reference.

\begin{corollary}[Finite width preserves the population rate]
\label{cor:finite-width}\label{thm:finite-width}
Under Assumption~\ref{ass:loss}, let $\tSigma_0\succeq0$ with
$\gap{\tSigma_0}>0$, set $r_0:=\rankop\tSigma_0$, and fix
$\varepsilon,\delta\in(0,1)$.
There is a universal $C>0$ such that, if
\begin{equation}\label{eq:width}
 n\ge CK^4\varepsilon^{-2}\bigl(r_0+\log(2/\delta)\bigr),
\end{equation}
then with probability at least $1-\delta$, it holds that 
\((1-\varepsilon)\tSigma_0\preceq\Sigma_0 \preceq(1+\varepsilon)\tSigma_0\).
On this event,
$\gap{\Sigma_0}\ge(1-\varepsilon)\gap{\tSigma_0}$,
the empirical factor flow \eqref{eq:mean-field-gradient-flow} is global, and
\(\Delta_t\le\Delta_0e^{-2\kappa(1-\varepsilon)\gap{\tSigma_0}t}\)
for all $t\ge0$.
The factors $(U_t, V_t)$ converge to a limit with globally optimal predictor $\frac{1}{n} V_\infty U_\infty^\top = \Svu_\infty$.
\end{corollary}

\begin{comment}
\emph{Proof sketch.}
Concentration is needed only on the $r_0$-dimensional range of $\tSigma_0$.
Expressed in an orthonormal basis of this range, the projections of the
whitened samples $\xi_i$ have second-moment matrix $\Id_{r_0}$ and
subgaussian bound at most $K$.
The resulting relative covariance
bound gives $\Sigma_0\succeq(1-\varepsilon)\tSigma_0$, hence
$a_\pm(\Sigma_0)\ge(1-\varepsilon)a_\pm(\tSigma_0)$ and
$\gap{\Sigma_0}\ge(1-\varepsilon)\gap{\tSigma_0}>0$.
This positive empirical support gap supplies the hypothesis needed to apply Theorem~\ref{thm:gf}.
Appendix~\ref{app:sampling} gives the full argument.\qed
\end{comment}

\begin{remark}[Width and initialization geometry]\label{rem:rank-adaptive-width}
The condition $\gap{\tSigma_0}>0$ forces $r_0\ge\min(d_u,d_v)$,
and this minimum is attainable. If $d_u\le d_v$, take
$\tSigma_0=\diag(2\sigma^2\Id_{d_u},0_{d_v})$, $\sigma>0$:
then $r_0=d_u$ and $\gap{\tSigma_0}=2\sigma^2$.
With Gaussian sampling, Corollary~\ref{cor:finite-width} gives width proportional to
$\varepsilon^{-2}(d_u+\log(2/\delta))$ and rate $4(1-\varepsilon)\kappa\sigma^2$, the same as for full-rank
population covariance $\sigma^2\Id_m$. Here every $v_i$ starts at zero,
so $\Svu_0=0$ and $\Delta_0=\widetilde\Delta_0$ exactly.
Exchanging the blocks gives dependence on $\min(d_u,d_v)$ in general.
This improvement changes the initialization: for $\tSigma_0=\sigma^2\Id_m$,
$r_0=m$. Its empirical support gap can vanish at small width even though its
population support gap is positive (Example~\ref{ex:white-small-width}).
\end{remark}

For the sampling scheme in Corollary~\ref{cor:finite-width}, the prefactor is
the empirical initial loss gap $\Delta_0$, which need not equal the population
loss gap $\widetilde\Delta_0$. Controlling it is a separate
accuracy requirement: for isotropic Gaussian sampling and the standard quadratic loss,
the expected excess initial error is $\sigma^4d_ud_v/(2n)$
(Proposition~\ref{prop:gaussian}). Theorem~\ref{thm:sampling-budgets}
gives a separate accuracy budget and a deterministic loss envelope.
For population variance $\sigma^2 \Id_m$ and $\varepsilon=1/2$,
the certified exponent is $2\kappa\sigma^2$.

%%% Local Variables:
%%% mode: LaTeX
%%% TeX-master: "../main"
%%% End:

\paragraph{Convergence of gradient descent}
% \label{sec:gd}
Unlike finite-width sampling of neurons, discretized gradient descent in the factor matrices $(U, V)$ with a finite step size changes the covariance evolution itself.
We show that sufficiently small step sizes retain global linear convergence by controlling the accumulated conservation error.
% \paragraph{Gradient descent on the factors}
% \label{sec:discrete-factor-descent}
We consider discrete-time gradient descent (GD) on the factor matrices, using an Euler discretization of \eqref{eq:mean-field-gradient-flow} with step size $\tau$.
As in the gradient flow setting, define $\Sigma_k:=W_kW_k^\top/n$, $\Svu_k:=V_kU_k^\top/n$, and $\Delta_k:=\ell(\Svu_k)-\ell_\star$.
Setting $E_k:=\nabla\ell(\Svu_k)$ and $G_k:=G_{E_k}$, the iterates read
\begin{equation}\label{eq:gd-algorithm}
 U_{k+1}=U_k-\tau E_k^\top V_k,\qquad
 V_{k+1}=V_k-\tau E_kU_k.
\end{equation}
Equivalently, $W_{k+1}=W_k-n\tau\nabla_W\cL(U_k,V_k)$: this is GD for $n\cL$ with stepsize $\tau$, or for $\cL$ with stepsize $n\tau$.
Writing $W_{k+1}=M_kW_k$ with $M_k:=\Id-\tau G_k$ gives
\begin{equation}\label{eq:gd-covariance}
 \Sigma_{k+1}=M_k\Sigma_kM_k
 =\Sigma_k-\tau(G_k\Sigma_k+\Sigma_kG_k)+\tau^2G_k\Sigma_kG_k.
\end{equation}
The additional term $\tau^2G_k\Sigma_kG_k$ arises from taking the GD step in the individual factors $(U,V)$ and is absent from an Euler step taken directly on the Bures ODE.

A positive initial support gap $g(\Sigma_0)$ still guarantees global linear convergence for sufficiently small step sizes $\tau$. 
The condition on $\tau$ is given below and depends only on the initial covariance, initial loss gap, and loss parameters; it controls both 
descent and the accumulated conservation error.

\begin{theorem}[Small-step global convergence of factor descent]
\label{thm:gd}\label{thm:discrete-bootstrap}
Under Assumption~\ref{ass:loss}, let $\Sigma_0:=W_0W_0^\top/n$ satisfy
$\gap{\Sigma_0}>0$, and suppose $\Delta_0>0$. Set
$I_\star:=4\sqrt{\Delta_0/\kappa}/\gap{\Sigma_0}$. If
\begin{equation}\label{eq:gd-step}
 0<\tau\le\tau_{\rm GD}:=\gap{\Sigma_0}e^{-2I_\star}\min\left\{
 \tfrac{1}{4\norm{\Sigma_0}_{\op}\sqrt{L\Delta_0}
              +9\sqrt2 L\norm{\Sigma_0}_{\op}^2e^{2I_\star}},\quad
 \tfrac{\sqrt2-1}{4\Delta_0}
 \right\},
\end{equation}
then, with $q:=1-\kappa\tau\gap{\Sigma_0}/\sqrt2\in(0,1)$ it holds that \(\Delta_k\le q^k\Delta_0\).
The factors $W_k$ and covariance $\Sigma_k$ converge at rate $q^{k/2}$ to limits
$W_\infty$ and $\Sigma_\infty\succeq0$, with
$\rankop\Sigma_\infty=\rankop\Sigma_0$ and
$\Svu_\infty\in\cS_\star$; explicit prefactors are given in
\eqref{eq:gd-tails}.
\end{theorem}

The support-gap condition $g(\Sigma_0) > 0$ is implied by $\Sigma_0\succ0$, in which case $\Sigma_\infty$ is also positive definite.
The threshold $\tau_{\rm GD}$ is positive for every initialization with $\gap{\Sigma_0}>0$ and $\Delta_0 < \infty$; no initial loss gap assumption or dimension-dependent term is needed.
The first restriction ensures descent, and the second controls the accumulated conservation error.

\begin{comment}
\emph{Proof sketch.}
\emph{Approximate conservation.}
Transporting $J$ by the discrete factor updates \eqref{eq:gd-algorithm} transports the two inequalities defining the initial margins as well.
The transported matrix differs from $J$ by a defect that starts at zero for every initialization and gains a contribution of order $\tau^2$ at each step.
If the defect has operator norm at most $\eta<1$, then
$\Suu_k\succeq(1-\eta)\aplus{\Sigma_0}\Id$ and
$\Svv_k\succeq(1-\eta)\aminus{\Sigma_0}\Id$, so the sum of these spectral
lower bounds remains positive: $(1-\eta)\gap{\Sigma_0}>0$
(Lemma~\ref{lem:gd-support}).

\emph{Descent and length.}
Assuming an upper bound on $\norm{\Sigma_k}_{\op}$ and a sufficiently small
conservation defect, $L$-smoothness of the predictor loss and the step-size restriction \eqref{eq:gd-step} give
\(
\Delta_{k+1}\le\Delta_k-
\frac{\tau\gap{\Sigma_0}}{2\sqrt2}\norm{E_k}_{\F}^2
\). 
The PL inequality then yields geometric loss decay. Summing the decreases
in $\Delta_k$ and $\sqrt{\Delta_k}$ bounds
$\sum_k\tau\norm{E_k}_{\F}^2$ and $\sum_k\tau\norm{E_k}_{\F}$, respectively.

\emph{Closing the bounds.}
The bound on $\sum_k\tau\norm{E_k}_{\F}$ gives a uniform upper bound on
$\norm{\Sigma_k}_{\op}$; together with the bound on
$\sum_k\tau\norm{E_k}_{\F}^2$, it also controls the accumulated conservation defect.
The stepsize restrictions keep both within the same bounds at the next
iterate, closing an induction. Bounds on $\sum_{j=k}^{\infty}\tau\norm{E_j}_{\F}$ give convergence
of the factors and covariance; see Appendix~\ref{app:gd}.\qed
\end{comment}

Appendix~\ref{app:gd} gives the explicit expression for $\tau_{\rm GD}$ when
$\Sigma_0=\sigma^2\Id_m$ and a comparison with the step-size condition obtained
by linearizing near an equilibrium (Remark~\ref{rem:gd-local}).
Proposition~\ref{prop:gd-adjustable} shows that sufficiently small step sizes
yield a certified loss-decay exponent arbitrarily close to $2\kappa\gap{\Sigma_0}$, the gradient-flow exponent in
Theorem~\ref{thm:balanced-convergence}.
Corollary~\ref{cor:sampled-gd} gives a step size chosen before sampling the
factors, using the loss parameters and the bounds on the initial support gap,
covariance norm, and loss gap from Theorem~\ref{thm:sampling-budgets}.

%%% Local Variables:
%%% mode: LaTeX
%%% TeX-master: "../main"
%%% End:

\section{Deep Linear Residual Networks}
\label{sec:resnets}\label{sec:deep-resnets}
We now turn to deep linear residual networks (ResNets), including a continuous-depth formulation.
To stabilize optimization as depth increases, the networks considered here use skip connections between successive layers and scale each residual update by the inverse depth.
In this section, we demonstrate that linear ResNets also enjoy global exponential convergence rates that remain uniform across depths,
provided the average conserved layer support gap is bounded away from zero and the residual predictors satisfy an a priori operator-norm bound independent of depth throughout training.
The latter is an additional hypothesis: layerwise conservation alone does not control forward and backward propagation.
For simplicity, we state the results for gradient flow, as in Section~\ref{sec:balanced-convergence}.

\paragraph{Finite-depth setting}
We consider a depth-$S$, width-$n$ network with matching input and output dimensions $d_u=d_v=d$, and set $h:=1/S$.
For $s=0,\ldots,S-1$, layer $s$ has factors $U_s,V_s\in\R^{d\times n}$, covariance $\Sigma_s:=W_sW_s^\top/n$ with $W_s:=\binom{U_s}{V_s}$, and predictor $\Svu_s:=V_sU_s^\top/n$.
The network maps an input $x^0\in\R^d$ to its output $x^S$ through
\[
 x^{s+1}=x^s+\tfrac{1}{Sn}V_sU_s^\top x^s=A_sx^s, \qquad A_s:=\Id+h\Svu_s,\qquad s=0,\ldots,S-1.
\]
Thus $x^S=P_Sx^0$, where the endpoint predictor is the ordered product
\[
 P_S:=A_{S-1}\cdots A_0
 =\left(\Id+\tfrac{V_{S-1}U_{S-1}^\top}{Sn}\right)\cdots
  \left(\Id+\tfrac{V_0U_0^\top}{Sn}\right).
\]
Backpropagation in the factors couples the layers through these products, which can be singular even when every layer covariance is positive definite.
Let $P_{>s}:=A_{S-1}\cdots A_{s+1}$ and $P_{<s}:=A_{s-1}\cdots A_0$, with empty products equal to $\Id$.
For $E:=\nabla\ell(P_S)$, the backpropagated gradient is $\widetilde E_s:=P_{>s}^\top E P_{<s}^\top$.
In residual mean-field time, each layer obeys
\begin{equation}\label{eq:resnet-layer-bures}
 \dot\Sigma_s=-(\widetilde G_s\Sigma_s+\Sigma_s\widetilde G_s),
 \qquad \widetilde G_s:=G_{\widetilde E_s}.
\end{equation}
This is the covariance equation corresponding to factor gradient flow for $nS\ell(P_S)$. Since $\widetilde G_sJ+J\widetilde G_s=0$,
the signed margins are conserved separately at every layer.

\begin{theorem}[Finite-depth convergence from conserved layer margins]\label{thm:resnet-finite}
Under Assumption~\ref{ass:loss} for the endpoint loss, suppose
$\Sigma_s(0)\succ0$ for all $s$, and write
$\overline g_S:=S^{-1}\sum_{s=0}^{S-1}\gap{\Sigma_s(0)}$ for the average initial
support gap. Suppose
$\sup_s\norm{\Svu_s(t)}_{\op}\le B<S$ for some constant $0 \leq B < \infty$ on the maximal existence interval,
and define the propagation lower bound $\underline{\sigma}_S:=(1-B/S)^{S-1}$.
Then the flow is global and
% \begin{equation}\label{eq:resnet}
\(\Delta_t:=\ell(P_S(t))-\ell_\star
 \le\Delta_0e^{-2\kappa\overline g_S\underline{\sigma}_S^2t}\).
% \end{equation}
Every layer's covariance and factors converge, and the limiting endpoint
is globally optimal.
\end{theorem}

\begin{comment}
\paragraph{Proof sketch}
The path bound keeps every forward and backward product invertible,
so each backpropagated gradient has norm at least $\underline{\sigma}_S\norm E_{\F}$.
Conservation supplies spectral lower bounds on the hidden blocks, giving
$-\dot\Delta\ge\overline g_S\underline{\sigma}_S^2\norm E_{\F}^2$; PL yields the rate.
Upper bounds on the same products, combined with the PL estimate for
$\sqrt{\Delta_t}$, bound the time integral of each backpropagated gradient's
Frobenius norm. These bounds prevent finite-time blow-up and prove
convergence of every layer; see Appendix~\ref{app:resnets}.\qed
\end{comment}

The rate combines the average support gap $\overline g_S$ with
$\underline{\sigma}_S^2$. For fixed $B$, this product-conditioning factor tends to
$e^{-2B}$ and is at least $e^{-4B}$ when $S\ge2B$.
For fixed $B$, a positive lower bound on the average layer support gap suffices for a convergence rate uniform in depth.

\paragraph{Continuous-depth limit}
Formally in the limit $S \to \infty$, residual products become propagators of a linear ODE.
Theorem~\ref{thm:resnet-finite} can be extended to a continuum of residual layers when these propagators stay uniformly conditioned, but requires a well-defined global training trajectory.
Let $s\in[0,1]$ denote depth. At time $t$, the network maps $x$ to
$P(t)x$, where $P(t):=X_1(t)$ and $X_s(t)$ solves
\[
 \partial_s X_s(t)=\Svu_s(t)X_s(t),\qquad X_0(t)=\Id.
\]
Write $P_{b,a}(t)$ for the propagator from depth $a$ to depth $b$. The backpropagated gradient at depth $s$ is
$\widetilde E_s:=P_{1,s}^\top\nabla\ell(P)P_{s,0}^\top$.
We assume the covariance profile $s\mapsto\Sigma_s(t)$ is defined for all
training times $t\ge0$ and depends locally absolutely continuously on $t$
as an element of $L^1([0,1];\Ssym^{2d})$.
The layer covariances are positive semidefinite and satisfy
\eqref{eq:resnet-layer-bures} in $L^1$. Its precise time-integral form is
\eqref{eq:resnet-integral} in Appendix~\ref{app:resnets}.

\begin{theorem}[Continuous-depth convergence from averaged margins]
\label{thm:resnet-continuous}
Under Assumption~\ref{ass:loss} for the endpoint loss, consider a covariance profile satisfying the assumptions above, with $\Sigma_s(0)\succ0$ for almost every $s$.
Set $\overline g:=\int_0^1\gap{\Sigma_s(0)}\,\dd s\in(0,\infty)$.
If $\sup_{t\ge0}\int_0^1\norm{\Svu_s(t)}_{\op}\,\dd s\le B<\infty$, then
\begin{equation}\label{eq:continuous-resnet-rate}
 \ell(P(t))-\ell_\star
 \le(\ell(P(0))-\ell_\star)e^{-2\kappa\overline g e^{-2B}t}.
\end{equation}
As $t\to\infty$, $\Sigma_s(t) \to \Sigma_s^\infty$ for almost every
depth $s$, and profiles converge in $L^1([0,1];\Ssym^{2d})$.
The network defined by the limiting profile attains loss $\ell_\star$.
Moreover, $\Sigma_s^\infty\succ0$ for almost every $s$, although its smallest
eigenvalue need not be bounded away from zero uniformly over depth.
\end{theorem}

\begin{comment}
\paragraph{Proof sketch}
With $E=\nabla\ell(P(t))$, Gr\"onwall's inequality bounds the depth propagators $P_{b,a}(t)$ and their inverses, so
$e^{-B}\norm E_{\F}\le\lVert\widetilde E_s\rVert_{\F}\le e^B\norm E_{\F}$.
Integrating the layerwise loss-dissipation bounds over depth gives a lower
bound proportional to the integrated support gap $\overline g$; the PL
inequality then gives the stated loss-decay rate.
The resulting bound on $\int_0^\infty\lVert\widetilde E_s(t)\rVert_{\F}\,\dd t$
is uniform over depth. Together with
the integrable initial covariance, it yields almost-everywhere and
$L^1$ convergence. Continuity of the endpoint map gives optimality;
Appendix~\ref{app:resnets} provides the details.\qed
\end{comment}

The need for additional residual-path boundedness assumptions arises because the loss depends only on the endpoint $s = 1$: see Remark~\ref{rem:resnet-path-escape} in Appendix~\ref{app:resnets}.

%%% Local Variables:
%%% mode: LaTeX
%%% TeX-master: "../main"
%%% End:

\section{Heavy-ball Momentum Convergence}
\label{sec:momentum}
We now turn our attention to the analogous dynamics in an inertial setting that models heavy-ball momentum-based optimization.
Upon introducing momentum, the system no longer closes on the position covariance of the neuron weights \((u,v)\) alone: one must instead consider the phase space of positions and velocities tracked jointly.
We show that, in terms of the \emph{phase covariance}, dynamics are again closed for linear networks.
We derive conservation identities describing this phase geometry; energy dissipation and control of position deformation give the convergence certificate below.

\paragraph{Dynamics and phase closure}
Inertia adds velocity as a state variable, so the natural moment over which the dynamics close is the joint position-velocity covariance.
For $W:=\binom UV$, predictor $\Svu:=VU^\top/n$, and damping $\gamma>0$,
write $E:=\nabla\ell(\Svu)$. The factor dynamics are
\begin{equation}\label{eq:finite-momentum}
 \ddot W+\gamma\dot W+G(\Svu)W=0,\qquad
 G(\Svu):=\begin{psmallmatrix}0&E^\top\\E&0\end{psmallmatrix}.
\end{equation}
Write $(P,Q):=(\dot U,\dot V)$, so the columns of $W$ and $\dot W$ are
positions $x_i:=(u_i,v_i)$ and velocities $y_i:=(p_i,q_i)$ respectively.
In this section $\Phase:=n^{-1}\binom W{\dot W}\binom W{\dot W}^{\!\top}$
denotes the full $2m$-dimensional phase covariance. Its blocks are
$\Spos:=WW^\top/n$, $\Svel:=\dot W\dot W^\top/n$, and
$\Xi:=W\dot W^\top/n$; the predictor remains the $vu$ block of $\Spos$.
With $G_t:=G(\Svu_t)$, the closed equation is
\begin{equation}\label{eq:phase}
 \dot\Phase_t=\Mob_t\Phase_t+\Phase_t\Mob_t^\top,\qquad
 \Phase_t=\begin{psmallmatrix}\Spos_t&\Xi_t\\\Xi_t^\top&\Svel_t\end{psmallmatrix},\qquad
 \Mob_t:=\begin{psmallmatrix}0&\Id_m\\-G_t&-\gamma\Id_m\end{psmallmatrix}.
\end{equation}
Appendix~\ref{app:momentum-proof} gives the particle equations, measure
formulation, and blockwise expansion.

\paragraph{Conformal laws and mechanical energy}
Without damping, the fundamental matrix generated by $\Mob_t$ in \eqref{eq:phase} is symplectic;
damping rescales the associated bilinear form exponentially in time. 
\begin{proposition}[Conformal phase conservation laws]\label{prop:phase-laws}
Let $\Jcan:=\begin{psmallmatrix}0&\Id_m\\-\Id_m&0\end{psmallmatrix}$.
The phase fundamental matrix \(\dot\Psi_t=\Mob_t\Psi_t\),
\(\Psi_0=\Id_{2m}\) generated by $\Mob_t$ satisfies
\[
 (\Psi_t)^\top \Jcan\Psi_t=e^{-\gamma t}\Jcan,\qquad
 \Phase_t=\Psi_t\Phase_0(\Psi_t)^\top.
\]
Consequently, the spectrum of \(e^{\gamma t}\Jcan\Phase_t\)
and all its trace powers are constant. At the characteristic level,
these identities imply that, for any two trajectories of
the same flow, \(z_a:=(x_a,y_a)\), \(z_b:=(x_b,y_b)\),
\(e^{\gamma t}z_a^\top \Jcan z_b\) is conserved.
\end{proposition}

\begin{comment}
\emph{Proof sketch.}
Symmetry of $G_t$ gives $\Mob_t^\top\Jcan+\Jcan\Mob_t=-\gamma\Jcan$.
Differentiating $\Psi_t^\top\Jcan\Psi_t$ proves the conformal identity;
the identity $\Phase_t=\Psi_t\Phase_0\Psi_t^\top$ then shows that $e^{\gamma t}\Jcan\Phase_t$
is similar to $\Jcan\Phase_0$.\qed
\end{comment}

Special spectral values permit reduced phase dynamics, we describe these in Appendix~\ref{app:intrinsic}.
As before, write $\Delta_t:=\ell(\Svu_t)-\ell_\star$ for the predictor loss gap. Define $\mathcal{K}_t:=\tr(\Svel_t)/2$ and $\cE_t:=\Delta_t+\mathcal{K}_t$ as the kinetic and mechanical excess energies. Differentiating $\cE_t$ along \eqref{eq:phase} gives $\dot{\cE}_t=-2\gamma\mathcal{K}_t$.
This identity gives global existence of the solution to \eqref{eq:phase} from every $\Sigma_0 \succeq 0$ (Proposition~\ref{prop:global-phase}).
The remaining task is to derive spectral lower bounds for $\Spos_t$.

\paragraph{Global convergence}
We seek an initialization condition on $\Sigma_0$ that keeps $\Spos_t$ uniformly elliptic for the whole trajectory.
Let $a_0:=\lambda_{\min}(\Spos_0)>0$ and $b_0:=\lambda_{\max}(\Spos_0)$.
A deformation budget $s>0$ specifies a spectral interval $[a_s,b_s]:=[a_0e^{-2s},b_0e^{2s}]$ for $\Spos_t$.
For $\cE_0>0$, put $\rho_s:=4b_sL+2\sqrt{2L\cE_0}$ and define
\begin{equation}\label{eq:certificate}
 \lambda_s(\gamma):=\tfrac{a_s\kappa\gamma}{\rho_s+2\gamma^2},\qquad
 \mathfrak S_s(\gamma):=\tfrac4{\lambda_s(\gamma)}
             \sqrt{\tfrac{3\cE_0}{2a_s}}.
\end{equation}
Here, $\lambda_s$ is the exponential decay rate of $\cE_t$, and $\mathfrak S_s$ bounds the resulting total deformation in $\Spos_t$, conditional on the spectral interval $[a_s, b_s]$.

\begin{theorem}[Global linear rate from a phase-energy certificate]\label{thm:momentum}\label{thm:bootstrap-unconditional-rate}
Under Assumption~\ref{ass:loss}, let $\Phase_0\succeq0$ with
$\Spos_0\succ0$ and $\cE_0>0$. If
$\mathfrak S_s(\gamma)<s$ for some $s>0$, then for all $t\ge0$,
\[
 a_s\Id_m\preceq\Spos_t\preceq b_s\Id_m,\qquad
 \cE_t\le3\cE_0 e^{-\lambda_s(\gamma)t}.
\]
The covariance $\Spos_t$ and any associated factors
$(U_t,V_t)$ evolving under \eqref{eq:finite-momentum} converge exponentially with rate $\lambda_s/2$ to limits with optimal predictor and $\Spos_\infty\succ0$; $\Svel_t$ tends to zero.
\end{theorem}

\begin{comment}
\paragraph{Proof sketch}
\emph{Conditional energy decay.}
Global existence of $\Phase_t$ follows from the energy identity $\dot{\cE}_t=-2\gamma\mathcal{K}_t$.
When the eigenvalues of $\Spos_t$ lie in $[a_s,b_s]$, the PL inequality gives a lower bound on the mean squared particle force in terms
of $\Delta_t$. For sufficiently small $\epsilon>0$, we show that $\cE_t+\epsilon\dot\Delta_t$ is a Lyapunov function comparable to $\cE_t$, whose exponential decay yields the bound on $\cE_t$ (Lemma~\ref{lem:conditional}).

\emph{Controlled deformation.}
Whenever $\Spos_t\succ0$, define $D_t:=\Xi_t^\top(\Spos_t)^{-1}$.
Then $\dot\Spos_t=D_t\Spos_t+\Spos_tD_t^\top$ and
$\norm{D_t}_{\op}\le\sqrt{2\mathcal{K}_t/a_s}$ while the eigenvalues of
$\Spos_t$ remain in $[a_s,b_s]$. Using $\mathcal{K}_t\le\cE_t$ and the
energy bound gives $\int_0^t\norm{D_r}_{\op}\,\dd r\le\mathfrak S_s(\gamma)<s$.

\emph{Preserving the spectral bounds.}
Gr\"onwall's inequality and the integral bound above prevent the eigenvalues
of $\Spos_t$ from leaving $[a_s,b_s]$, so the energy bound and convergence
estimates hold globally. Appendix~\ref{app:momentum-proof} gives the full
proof and explicit convergence bounds; Appendix~\ref{app:momentum-details}
gives a sharper conditional energy bound.\qed
\end{comment}

\begin{corollary}[Two-scale white initialization]
\label{cor:two-scale-white-window}
Under Assumption~\ref{ass:loss}, let
\(\Spos_0=\sigma_{vu}^2\Id_m\), \(\Xi_0=0\), and
\(\Svel_0\succeq0\) have trace \(m\sigma_{qp}^2\), with
\(\sigma_{vu}>0\), \(\sigma_{qp}\ge0\).
Then \(\cE_0=\Delta+m\sigma_{qp}^2/2\), where
\(\Delta:=\ell(0)-\ell_\star\).
If \(\cE_0>0\), the certificate is
% \begin{equation}\label{eq:white-direct}
\[
 2\gamma^2-
 \tfrac{s\kappa\sigma_{vu}^3e^{-3s}}
 {4\sqrt{3\cE_0/2}}\gamma
 +4L\sigma_{vu}^2e^{2s}
 +2\sqrt{2L\cE_0}<0.
\]
% \end{equation}
For fixed \(s,\gamma>0,\Delta,\sigma_{qp},m,L,\kappa\), sufficiently
large \(\sigma_{vu}\) satisfies this condition.
\end{corollary}

Appendix~\ref{app:large-position-scale} gives explicit bounds on the admissible $\gamma$ interval and analyzes the large-$\sigma_{vu}$ asymptotics.

\paragraph{Numerical view of the bootstrap certificate}
\begin{wrapfigure}{r}{0.34\linewidth}
\centering
\resizebox{\linewidth}{!}{\subimport{figures/}{momentum_bootstrap_heatmap_notebook.pgf}}
\footnotesize
\captionsetup{font=footnotesize,skip=2pt}
\caption{Two-scale white momentum initialization. Curves show
$\mathfrak G_s(\gamma)=0$ for the color-bar scale $\sigma_{vu}$;
their interiors are certified damping windows.}
\label{fig:momentum-bootstrap-heatmap}\label{fig:momentum-window}
\end{wrapfigure}
Writing $\mathfrak G_s(\gamma):=\mathfrak S_s(\gamma)-s$, Figure~\ref{fig:momentum-bootstrap-heatmap} shows the boundaries
$\mathfrak G_s(\gamma)=0$ for $1\le\sigma_{vu}\le10^3$, with $(d_u,d_v)=(3,2)$,
$\kappa=L=1$, $\sigma_{qp}=1$, and $\ell(0)-\ell_\star=1$.
The regions where $\mathfrak G_s(\gamma)<0$ specify the pairs $(s,\gamma)$ for which Theorem~\ref{thm:momentum} guarantees spectral bounds and exponential energy decay.
For the two-scale initialization, sufficiently large $\sigma_{vu}$ yields a nonempty interval of $\gamma$ satisfying Corollary~\ref{cor:two-scale-white-window}.
For fixed initialization and damping, a smaller admissible $s$ gives a
tighter tube and faster rate; Figure~\ref{fig:momentum-bootstrap-lambda-multi-sigma}
illustrates the limiting rate at the feasibility boundary analyzed in
Appendix~\ref{app:large-position-scale}.

Figure~\ref{fig:momentum-empirical-vs-theory} compares empirical losses with the upper bounds on $\cE_t$ from Corollary~\ref{cor:explicit-tightened-conditional-momentum-rate}.
Each upper bound curve uses an explicit, approximately optimal rate certified by Theorem~\ref{thm:momentum} (see Proposition \ref{prop:optimal-tightened-conditional-momentum-rate}); no rate is fitted
to the empirical trajectory. All displayed losses stay below their upper bounds.
Notably in this inertial setting, the predictor loss may oscillate even while the mechanical energy decreases.

\begin{figure}[!htbp]
\centering
\resizebox{\linewidth}{!}{\subimport{figures/}{momentum_empirical_vs_tightened_theory_by_sigma.pgf}}
\captionsetup{font=footnotesize,skip=2pt}
\caption{
Empirical predictor-loss trajectories (solid) and the corresponding upper bounds
on mechanical excess energy (dotted) for two-scale white initializations with varying
\(\sigma_{vu}\) and damping \(\gamma\): \textbf{(a)} $\sigma_{vu}=200$,
\textbf{(b)} $\sigma_{vu}=500$, and \textbf{(c)} $\sigma_{vu}=1000$.
Each theoretical curve uses the explicit, approximately optimal rate on a tube certified by
Theorem~\ref{thm:bootstrap-unconditional-rate}.
}
\label{fig:momentum-empirical-vs-theory}\label{fig:momentum-envelopes}
\end{figure}
\WFclear

%%% Local Variables:
%%% mode: LaTeX
%%% TeX-master: "../main"
%%% End:

\FloatBarrier
\section{Conclusion}
The present work addresses two-layer linear networks, making use of the exact closure of mean-field training on the neuron law covariance as a finite-dimensional Bures flow.
This flow enjoys conservation laws that control the ellipticity of the predictor preconditioner through a conserved spectral support gap, yielding global linear convergence whenever this gap is positive.
In particular, this is guaranteed for every positive-definite initial covariance and for smooth PL predictor losses. 
We show that this guarantee survives finite-width sampling of neurons, and extends to discrete-time gradient descent in the factors through approximate versions of the same conservation laws.
The covariance viewpoint clarifies two further extensions: for deep linear ResNets, the same conservation law holds layerwise and gives linear convergence conditional on bounded residual paths.
Heavy-ball momentum dynamics close in the joint phase covariance, where conformal conservation laws and mechanical energy lead to an explicit damping certificate from general initial phase moments.
The certificate includes two-scale white initializations without imposing a phase-covariance identity. Covariance conditioning is the common quantitative mechanism throughout our analysis.

%%% Local Variables:
%%% mode: LaTeX
%%% TeX-master: "../main"
%%% End:


\subsection*{Acknowledgments}
Gabriel Peyr\'e acknowledges support from the European Research Council through
the WOLF project, and from France 2030 through grant ANR-23-IACL-0008
(PRAIRIE-PSAI), administered by the Agence Nationale de la Recherche.

\subsection*{AI use statement}
In this work, we used generative AI tools including GPT 5.6 and GPT 6 for aiding and polishing writing, research ideation and execution, drafting some sections of the paper, and deriving or proving mathematical claims as directed by the authors. Where outputs were produced by generative AI, these were refined through extensive iteration, verification, and manual editing of the results and exposition. 
We have reviewed all AI-assisted work, having verified the claims, proofs and other material presented in this paper.
We take responsibility for the final content of this work, including text, claims or artifacts produced with the aid of generative AI.

\newpage 
\bibliography{references}
\bibliographystyle{plainnat}
\appendix
\clearpage
\section{Notation table}\label{app:notation}
The table collects the symbols, their meanings, and their first definitions.
\begingroup
\small
\renewcommand{\arraystretch}{1.12}
\begin{longtable}{@{}>{\raggedright\arraybackslash}p{0.21\linewidth}
 >{\raggedright\arraybackslash}p{0.53\linewidth}
 >{\raggedright\arraybackslash}p{0.20\linewidth}@{}}
\toprule
Notation & Meaning & First reference\\
\midrule
\endfirsthead
\toprule
Notation & Meaning & First reference\\
\midrule
\endhead
\(d_u,d_v,n,m\)&Input/output dimensions, width, and \(m:=d_u+d_v\)&Section~\ref{sec:model}\\
\(U,V,W:=\binom UV,\cL\)&Factor matrices, stacked positions, and factor loss&Section~\ref{sec:model}\\
\(\Sigma,\Suu,\Svv,\Svu\)&Uncentered moment, hidden blocks, and predictor&Section~\ref{sec:laws}\\
\(\ell,\ell_\star,\cS_\star\)&Predictor loss, infimum, and minimizer set&Assumption~\ref{ass:loss}; Section~\ref{sec:convergence}\\
\(L,\kappa,\lambda\)&Smoothness, PL, and QG constants&Assumption~\ref{ass:loss}; Eq.~\eqref{eq:quadratic-growth}\\
\(\quadH,C_z,\Svu_\star\)&Quadratic Hessian, input second moment, teacher&Section~\ref{sec:model}\\
\(\mu,\widehat\mu_W\)&Neuron law and empirical law&Appendix~\ref{app:mean-field}\\
\(E,G_E\)&Predictor gradient and symmetric split generator&Section~\ref{sec:laws}\\
\(J,\Phi_t\)&Input/output signature and position fundamental matrix&Section~\ref{sec:laws}\\
\(\mathcal H,\mathcal H_{\varphi,r}\)&Mean-field law and Gram-test law&Def.~\ref{def:mean-field-law}; Eq.~\eqref{eq:gram-mean-field-law}\\
\(\mathcal H_k,\Gamma\)&Spectral laws and labeled Gram imbalance&Propositions~\ref{prop:covariance-closure-laws}, \ref{prop:mean-field-conservation-laws}\\
\(a_\pm(\Sigma),\gap{\Sigma}\)&Covariance support margins and gap&Def.~\ref{def:support-gap}\\
\(\Delta_t,I_0,B_\infty\)&Loss gap, force-length bound, covariance upper bound&Section~\ref{sec:convergence}; Proposition~\ref{prop:gf-lift-tail}\\
\(\varepsilon,\zeta,\beta\)&Rate tolerance, absolute accuracy, concentration tolerance&Corollary~\ref{cor:finite-width}; Eq.~\eqref{eq:sampling-budgets}\\
\(\tSigma_0,\widetilde\Delta_0,r_0\)&Population covariance, reference loss gap, and rank; empirical covariance is $\Sigma_0$&Corollary~\ref{cor:finite-width}\\
\(\underline g,\overline\Delta,\overline B_0\)&Empirical support-gap lower bound, loss-gap and covariance-norm upper bounds&Eqs.~\eqref{eq:sampling-budgets}, \eqref{eq:sampling-event}\\
\(\tau,I_\star\)&GD step and force-length bound&Eq.~\eqref{eq:gd-algorithm}; Theorem~\ref{thm:gd}\\
\(\alpha_\star,M_{\nabla},B_\star\)&Proof-local GD descent scale, initial gradient bound, and covariance bound&Eq.~\eqref{eq:gd-constants}\\
\(M_k,\mathcal J_k,\mathcal D_k\)&Factor update, transported signature, and its defect&Eq.~\eqref{eq:gd-algorithm}, \eqref{eq:gd-signature}\\
\(S,A_s,P_S,\widetilde E_s\)&Residual depth, layer map, endpoint, and backpropagated gradient&Section~\ref{sec:resnets}\\
\(\overline g_S,\overline g\)&Finite/continuous depth averages of $\gap{\Sigma_s(0)}$&Theorems~\ref{thm:resnet-finite}, \ref{thm:resnet-continuous}\\
\(s,B\)&Continuous depth and residual-path bound&Section~\ref{sec:resnets}\\
\(\underline{\sigma}_S,\overline{\sigma}_S\)&Lower and upper residual-propagation bounds&Section~\ref{sec:resnets}; Appendix~\ref{app:resnets}\\
\(x:=(u,v),y:=(p,q)\)&Phase position and velocity variables&Section~\ref{sec:momentum}\\
\(\Phase,\Spos,\Xi,\Svel\)&Full phase moment and its position, mixed, velocity blocks&Eq.~\eqref{eq:phase}\\
\(\gamma,\Jcan\)&Damping and canonical symplectic matrix&Eq.~\eqref{eq:finite-momentum}, Proposition~\ref{prop:phase-laws}\\
\(\Mob_t,\Psi_t\)&Phase generator and characteristic matrix&Eq.~\eqref{eq:phase}; Proposition~\ref{prop:phase-laws}\\
\(D,Q\)&Velocity regression and residual covariance (Schur complement)&Eq.~\eqref{eq:regression}\\
\(\mathcal{K},\cE\)&Kinetic and total mechanical excess energies; $\cE_t=\Delta_t+\mathcal{K}_t$&Section~\ref{sec:momentum}\\
\(s,a_s,b_s,\rho_s,\lambda_s,\mathfrak S_s\)&Deformation budget, position tube, force bound, rate, certificate; \(s\) here is not depth&Eq.~\eqref{eq:certificate}\\
\(\sigma_{vu},\sigma_{qp}\)&White-position scale and matched per-coordinate velocity scale&Corollary~\ref{cor:two-scale-white-window}\\
\(\varepsilon,\eta,q,\mathcal V\)&Momentum energy correction, Young parameter, equivalence constant, Lyapunov energy&Proposition~\ref{prop:tightened}\\
\(\cN,\mathfrak C_0\)&Mean squared particle force and tightened initial energy prefactor&Lemma~\ref{lem:conditional}, Corollary~\ref{cor:tightened-bootstrap}\\
\(\mathfrak o(J),\xi_A\)&Matrix Lie algebra and infinitesimal congruence field&Appendix~\ref{app:law-counting}\\
\bottomrule
\end{longtable}
\endgroup

% !TEX root = ../main.tex

\section{Measure Formulation and Wasserstein Conservation Laws}
\label{app:mean-field}
The covariance equations in the main text have the same meaning at both finite and infinite widths. In this section, we introduce the measure flow that captures training dynamics, define a law shared by all widths, and lifts the spectral invariants to functionals of a neuron distribution.

\subsection{Neuron distributions and Wasserstein flow}
We describe the configuration of any number of neurons by their law: for $w:=(u,v)\in\R^m$, let
$\mu\in\cP_2(\R^m)$ and write its uncentered second moment
\begin{equation}\label{eq:measure-covariance}
 \Sigma(\mu):=\int ww^\top\,\dd\mu(w).
\end{equation}
For $n$ neurons $W:=(w_1,\ldots,w_n)$, the neuron law is an empirical distribution 
$\widehat\mu_W:=n^{-1}\sum_{i=1}^n\delta_{w_i}$, for which 
$\Sigma(\widehat\mu_W)=WW^\top/n$. A general covariance
$\Sigma\succeq0$ is also realized by a centered Gaussian distribution.
The loss functional is $\cF(\mu):=\ell(\Svu(\mu))$, so
$\cF(\widehat\mu_W)=\cL(U,V)$.

For a functional with a sufficiently regular first variation, the
Wasserstein gradient is the spatial gradient
$\nabla_{W_2}\cF(\mu)(w):=\nabla_w(\delta\cF/\delta\mu)(\mu,w)$ and its negative is the velocity in the continuity equation
\begin{equation}\label{eq:wasserstein-pde}
 \partial_t\mu_t+\nabla_w\cdot(\mu_t b_t)=0,\qquad
 b_t:=-\nabla_{W_2}\cF(\mu_t).
\end{equation}
The PDE is understood distributionally, including for empirical
measures. In the present split model,
$\delta\cF/\delta\mu=\tfrac12 w^\top G_tw$ up to an additive constant.
Consequently, 
\begin{equation}\label{eq:split-wasserstein-pde}
 \partial_t\mu_t=\nabla_w\cdot(\mu_tG_tw),\qquad
 \mu_t=(\Phi_t)_\#\mu_0,
\end{equation}
where $\Phi_t$ is the fundamental matrix in \eqref{eq:orbit}.
For the empirical neuron law, the characteristics of \eqref{eq:wasserstein-pde} are exactly
$\dot u_i=-E_t^\top v_i$, $\dot v_i=-E_tu_i$.

Conversely, given a solution $\Sigma_t$ of \eqref{eq:bures-G}, transporting
any initial distribution with second moment $\Sigma_0$ by
$\dot\Phi_t=-G_t\Phi_t$ produces a Wasserstein gradient flow with second
moment $\Sigma_t$. Indeed, its second moment and $\Sigma_t$ solve the same
linear equation with the same initial data. Matching covariances alone
does not ensure that a measure curve follows this gradient flow.

\subsection{Covariance-dependent objectives}
Moment closure is not specific to the off-diagonal predictor:
the gradient flow of any sufficiently regular objective depending on second moments closes in terms of a Bures equation.
\begin{proposition}[Second-moment closure]
\label{prop:moment-closure}
Let $F$ be a differentiable scalar function on an open neighborhood of
$\Sp^m$ in the symmetric matrices, and define $\cF(\mu):=F(\Sigma(\mu))$.
Let $\mu_t$ be a sufficiently regular Wasserstein gradient flow of $\cF$, and set $\Sigma_t:=\Sigma(\mu_t)$ and $D_t:=\nabla_\Sigma F(\Sigma_t)$.
Then the characteristics satisfy \(\dot w_t=-2D_tw_t\), and the second moment evolves autonomously as
\(\dot\Sigma_t=-2(D_t\Sigma_t+\Sigma_tD_t)\).
Thus the Wasserstein flow of $\cF$ closes exactly on the Bures gradient flow of $F$ on second moments, with the normalization used throughout the paper.
\end{proposition}

\begin{proof}[Proof of Proposition~\ref{prop:moment-closure}]
The first variation of $\cF$ at $\mu_t$, up to an additive constant, is
$\frac{\delta\cF}{\delta\mu}(\mu_t,w)=\ip{D_t}{ww^\top}=w^\top D_tw$.
Taking its Euclidean gradient in $w$ gives $2D_tw$ and hence the claimed characteristic equation.
Finally, writing the flow map as $w_0 \mapsto w_t$ and differentiating $\Sigma_t=\int w_tw_t^\top\,\mathrm{d}\mu_0(w_0)$ gives
\[
\dot\Sigma_t=\int \left(\dot w_tw_t^\top+w_t\dot w_t^\top\right)\,\mathrm{d}\mu_0(w_0)=-2(D_t\Sigma_t+\Sigma_tD_t).
\]
\end{proof}


\subsection{Width consistency and covariance-level laws}
We seek conservation laws defined on neuron distributions that remain
conserved when evaluated on finite networks of any width.

\begin{definition}[Mean-field conservation law]\label{def:mean-field-law}
A scalar functional \(\mathcal H(\mu)\) is a mean-field conservation law
for a family of flows if it is constant along every admissible trajectory
of every flow in that family, for every admissible initialization.
It is covariance-level if \(\mathcal H(\mu)=h(\Sigma(\mu))\).
\end{definition}
For differentiable functionals, the chain rule expresses conservation as
orthogonality:
\[
 \frac{\dd}{\dd t}\mathcal H(\mu_t)
 =-\int\nabla_{W_2}\mathcal H(\mu_t)\cdot
          \nabla_{W_2}\cF(\mu_t)\,\dd\mu_t=0, 
\]
whenever taking gradients and the chain rule are well-defined.
The laws below also have a direct characteristic proof.
A single mean-field law \( \mathcal H\) induces laws
\(h_n(W):=\mathcal H(\widehat\mu_W)\) for every width, in the
finite-dimensional sense of \citet{MarcotteGribonvalPeyre2023}.
Unrelated laws found separately at each width need not arise this way.
Conversely, conservation at every finite width passes to the measure flow
whenever the empirical flows converge in \(W_2\) and \(\mathcal H\) is
continuous in that metric. For covariance-level laws no limiting argument is needed:
every positive semidefinite covariance has an empirical realization, and
the common closed moment equation gives the equivalence with Bures laws.

\subsection{Gram-function laws and their spectral restriction}
Preservation of pairwise indefinite products yields functionals
that need not depend only on covariance.
\begin{proposition}[Fundamental mean-field laws]
\label{prop:mean-field-conservation-laws}
Let $\mu_t=(\Phi_t)_\#\mu_0$ be the characteristic flow of
$\dot w_t=-G_tw_t$, where $G_t^\top=G_t$ and $G_tJ+JG_t=0$.
For every $r\ge1$ and Borel function $\varphi:\Ssym^r\to\R$,
the functional
\begin{equation}\label{eq:gram-mean-field-law}
 \mathcal H_{\varphi,r}(\mu)
 :=\int\varphi\!\left((w_a^\top Jw_b)_{1\le a,b\le r}\right)
       \prod_{a=1}^r\dd\mu(w_a)
\end{equation}
is conserved whenever this integral is absolutely integrable at
initialization. In particular, the spectral laws
$\mathcal H_k(\Sigma(\mu))$ of \eqref{eq:bures-G}
are conserved for every $\mu_0\in\cP_2(\R^m)$.
For particles, the Gram imbalance
$\Gamma_t:=V_t^\top V_t-U_t^\top U_t=-W_t^\top JW_t$ is constant and
\begin{equation}\label{eq:finite-spectrum-gamma}
 \mathcal H_k(\Sigma_t)=\tr\!\left((-\Gamma_t/n)^k\right),
 \qquad k\ge1.
\end{equation}
The nonzero spectra of $J\Sigma_t$ and $-\Gamma_t/n$ coincide.
\end{proposition}

\begin{proof}[Proof of Proposition~\ref{prop:mean-field-conservation-laws}]
Along characteristics,
\(
\frac{\mathrm{d}}{\mathrm{d}t}(w_a^\top Jw_b)
=-w_a^\top(G_tJ+JG_t)w_b=0,
\)
so \eqref{eq:gram-mean-field-law} is conserved.
Bounded Borel tests are integrable, and cyclic polynomial tests need
only second moments because each sampled variable appears quadratically.
For the cyclic polynomial test function, with \(w_{k+1}:=w_1\), repeated
integration gives
\[
\int\prod_{\ell=1}^k
\left(w_\ell^\top Jw_{\ell+1}\right)
\prod_{\ell=1}^k\mathrm{d}\mu(w_\ell)
=\tr\!\left((J\Sigma)^k\right).
\]
Thus \eqref{eq:bures-G} is a special case of
\eqref{eq:gram-mean-field-law}.
The first \(m\) trace powers determine the characteristic polynomial by Newton identities.
For empirical measures, $J\Sigma=(JW)(W^\top/n)$ while $(W^\top/n)(JW)=-\Gamma/n$.
The nonzero eigenvalues of $AB$ and $BA$ coincide, which gives the spectral statement and \eqref{eq:finite-spectrum-gamma}.
\end{proof}

%%% Local Variables:
%%% mode: LaTeX
%%% TeX-master: "../main"
%%% End:

\section{Technical details and proofs}\label{app:technical-details}\label{app:convergence-proofs}
\subsection{Loss estimates}\label{app:proofs}
\begin{lemma}[Smooth PL consequences]\label{lem:pl}
Under Assumption~\ref{ass:loss}, the minimizer set
$\cS_\star:=\{S:\ell(S)=\ell_\star\}$ is nonempty, and
\begin{equation}\label{eq:pl-consequences}
\norm{\nabla\ell(S)}_{\F}^2\le2L(\ell(S)-\ell_\star),\qquad
\dist(S,\cS_\star)^2\le2(\ell(S)-\ell_\star)/\kappa.
\end{equation}
For every perturbation $H$,
\begin{equation}\label{eq:perturb}
\sqrt{\ell(S+H)-\ell_\star}
\le\sqrt{\ell(S)-\ell_\star}+\sqrt{L/2}\norm H_{\F}.
\end{equation}
\end{lemma}

\begin{proof}[Proof of Lemma~\ref{lem:pl}]

Apply the descent lemma at $S-\nabla\ell(S)/L$ for the first inequality.
Since $\nabla\ell$ is globally Lipschitz, the Euclidean gradient flow
$\dot S=-\nabla\ell(S)$ exists for all $t\ge0$.
Writing $h:=\ell(S)-\ell_\star$ gives $\dot h=-\norm{\nabla\ell}^2$ and,
where $h>0$, $-\frac{\dd}{\dd t}\sqrt h\ge\sqrt{\kappa/2}\norm{\dot S}$.
Thus its total remaining length is at most $\sqrt{2h/\kappa}$.
The flow has a limit, with loss gap zero by PL, proving attainment and the distance
bound. Finally, the descent lemma and the first inequality bound the
perturbed loss gap by $(\sqrt h+\sqrt{L/2}\norm H_{\F})^2$.
\end{proof}
\begin{remark}[Quadratic growth and convexity]
Lemma~\ref{lem:pl} proves the quadratic-growth bound
\begin{equation}\label{eq:quadratic-growth}
 \ell(S)-\ell_\star\ge\frac\lambda2\dist(S,\cS_\star)^2
 \qquad\text{with }\lambda=\kappa.
\end{equation}
If a larger quadratic-growth constant $\lambda$ is known, it improves
the distance bounds while leaving the loss rates, which use $\kappa$, unchanged.
The PL-to-quadratic-growth implication does not need convexity
\citep{KarimiNutiniSchmidt2016}. Conversely, convex quadratic growth
with constant \(\lambda\) implies PL with constant \(\lambda/4\):
for a nearest minimizer \(S_\star\),
\(\Delta\le\ip{\nabla\ell(S)}{S-S_\star}
\le\norm{\nabla\ell(S)}_{\F}\sqrt{2\Delta/\lambda}\).
The best PL and quadratic-growth constants need not coincide.
For a nonzero quadratic Hessian $\quadH\succeq0$,
$\lambda=\kappa=\lambda_{\min}^{+}(\quadH)$ is valid and
$\cS_\star=\Svu_\star+\kerop\quadH$.
\end{remark}

\subsection{Covariance closure and conservation laws}
\label{app:covariance-laws}
Direct matrix differentiation proves closure and conservation at
every width. The measure interpretation is given separately in
Appendix~\ref{app:mean-field}.

\begin{proof}[Proof of Proposition~\ref{prop:covariance-closure-laws}]
Since $\dot W_t=-G_tW_t$ and $G_t^\top=G_t$,
differentiating $W_tW_t^\top/n$ gives \eqref{eq:bures-G}.
For any covariance solution, define
$\dot\Phi_t=-G_t\Phi_t$, $\Phi_0=\Id$. Uniqueness for the resulting
linear covariance equation yields
$\Sigma_t=\Phi_t\Sigma_0\Phi_t^\top$, so positive semidefiniteness
is preserved. The identity $G_tJ+JG_t=0$ gives
$\frac{\dd}{\dd t}(\Phi_t^\top J\Phi_t)=0$.
Invertibility of $\Phi_t$ then gives both identities in
\eqref{eq:orbit}. Equivalently,
$\frac{\dd}{\dd t}(J\Sigma_t)=[G_t,J\Sigma_t]$, so
$\frac{\dd}{\dd t}\tr((J\Sigma_t)^k)=0$ for every $k\ge1$.
The first $m$ trace powers determine the characteristic polynomial.
For later use, the block equations are
\begin{equation}\label{eq:block-flow-general}
 \begin{aligned}
 \dSuu_t&=-\Suv_tE_t-E_t^\top\Svu_t,\quad
 \dSvv_t=-\Svu_tE_t^\top-E_t\Suv_t,\quad
 \dSvu_t&=-E_t\Suu_t-\Svv_tE_t.
 \end{aligned}
\end{equation}
\end{proof}

\begin{proof}[Proof of Proposition~\ref{lem:margins}]
For any $c\in\R$, the characteristic congruence gives
\[
 \Sigma_t-cJ=\Phi_t(\Sigma_0-cJ)\Phi_t^\top.
\]
Invertibility of $\Phi_t$ preserves the entire feasible interval of $c$,
so both endpoints are conserved even when $\Sigma_0$ is singular.
Taking principal blocks proves the hidden bounds.
\end{proof}

\begin{proof}[Spectral formula \eqref{eq:gap-spectrum}]
The feasible sets in Definition~\ref{def:support-gap} contain zero, are
closed intervals, and are bounded above by the corresponding hidden-block
minimum eigenvalues. Thus the maxima exist.
For $\Sigma\succ0$, put $H:=\Sigma^{1/2}J\Sigma^{1/2}$.
Sylvester's inertia theorem gives $d_v$ negative and $d_u$ positive
eigenvalues. Since $H^{-1}=\Sigma^{-1/2}J\Sigma^{-1/2}$,
$\Sigma-cJ=\Sigma^{1/2}(\Id-cH^{-1})\Sigma^{1/2}$ gives
$\aplus{\Sigma}=\lambda_{d_v+1}(H)$ and
$\aminus{\Sigma}=-\lambda_{d_v}(H)$.
For singular $\Sigma$, apply this to $\Sigma+\epsilon\Id$.
As $\epsilon\downarrow0$, each margin decreases to a finite limit that
is feasible for $\Sigma$ by closedness, and hence equals its margin.
Continuity of the square root and of ordered symmetric eigenvalues
proves \eqref{eq:gap-spectrum} without deleting any zero eigenvalues.
Finally $H$ and $J\Sigma$ have the same characteristic polynomial by
the $AB,BA$ identity, even when they are not similar.
\end{proof}

\begin{remark}[Sharpness of the supporting constants]
For every $\Sigma\succeq0$, testing the supporting inequalities on the two
coordinate subspaces gives the upper bounds in
\[
 \lambda_{\min}(\Sigma)\le\aplus{\Sigma}\le\lambda_{\min}(\Suu),\qquad
 \lambda_{\min}(\Sigma)\le\aminus{\Sigma}\le\lambda_{\min}(\Svv).
\]
The lower bounds follow from $\Id\mp J\succeq0$.
When $\Svu=0$, the two upper bounds are attained.
Consequently $2\lambda_{\min}(\Sigma)\le\gap{\Sigma}\le2\norm{\Sigma}_{\op}$.
Both margins are positive exactly when $\Sigma\succ0$: if both are
positive and $x\in\kerop\Sigma$, the supporting inequalities force
$x^\top Jx=0$ and then $(\Sigma-\aplus{\Sigma}J)x=0$, so $x=0$.
The convergence theorem requires only their sum to be positive, so it
also covers some singular covariances. Each positive margin makes the
corresponding principal block positive definite, giving the rank bounds
in Remark~\ref{rem:singular-support}.
\end{remark}

\subsection{Linear Convergence}\label{app:gf}
We prove the spectral-margin theorem and its quantitative trajectory
bounds first. Conditional and determinant-based criteria follow as
supplementary comparisons.

\begin{proposition}[Quantitative trajectory bounds]\label{prop:gf-lift-tail}
Under Theorem~\ref{thm:gf}, set
$a_0:=\lambda_{\min}(\Sigma_0)$, $B_0:=\norm{\Sigma_0}_{\op}$,
$I_0:=\gap{\Sigma_0}^{-1}\sqrt{2\Delta_0/\kappa}$, and $B_\infty:=B_0e^{2I_0}$.
Then
\begin{align}
 a_0e^{-2I_0}\Id_m&\preceq\Sigma_t\preceq B_\infty\Id_m,\label{eq:gf-tube}\\
 \norm{\Sigma_t-\Sigma_\infty}_{\op}
 &\le2B_\infty I_0e^{-\kappa \gap{\Sigma_0}t},\label{eq:gf-tail}\\
 \norm{\Svu_t-\Svu_\infty}_{\F}
 &\le2B_\infty I_0e^{-\kappa \gap{\Sigma_0}t}.\label{eq:gf-predictor-tail}
\end{align}
For the characteristic lift of any initial law in $\cP_2$ with this
second moment, there is a limit $\mu_\infty$ and
\begin{equation}\label{eq:gf-w2}
 W_2(\mu_t,\mu_\infty)
 \le\min\!\left\{\sqrt{\tr\Sigma_0}\,e^{I_0},\sqrt{2B_\infty}\right\}
       I_0e^{-\kappa \gap{\Sigma_0} t}.
\end{equation}
The same bound applies to $\norm{W_t-W_\infty}_{\F}/\sqrt n$ for particles.
The smallest positive covariance eigenvalue also satisfies
$\lambda_{\min}^{+}(\Sigma_t)\ge
e^{-2I_0}\lambda_{\min}^{+}(\Sigma_0)$, including at $t=\infty$.
\end{proposition}

\begin{proof}[Proof of Theorem~\ref{thm:gf} and Proposition~\ref{prop:gf-lift-tail}]
The predictor equation \eqref{eq:block-flow-general} gives the exact
dissipation identity on every positive semidefinite trajectory:
\begin{equation}\label{eq:flow-dissipation}
\frac{\mathrm{d}}{\mathrm{d}t}\ell(\Svu_t)
=-\tr(E_t^\top E_t\Suu_t)-\tr(E_t^\top\Svv_tE_t)\le0.
\end{equation}
The transported hidden-block bounds from Proposition~\ref{lem:margins}
therefore give $-\dot\Delta\ge \gap{\Sigma_0}\norm E_{\F}^2\ge2\kappa \gap{\Sigma_0}\Delta$.
On any finite existence interval, the preceding dissipation and PL give,
where the loss gap is positive,
\[
 -\frac{\dd}{\dd t}\sqrt{\Delta_t}
 \ge\frac{\gap{\Sigma_0}\norm{E_t}_{\F}^2}{2\sqrt{\Delta_t}}
 \ge \gap{\Sigma_0}\sqrt{\kappa/2}\norm{E_t}_{\F}.
\]
At zero loss gap the flow is stationary. Hence, for every finite $T>t$,
\[
 \int_t^T\norm{E_r}_{\F}\,\dd r
 \le\frac1{\gap{\Sigma_0}}\sqrt{\frac2\kappa}
       (\sqrt{\Delta_t}-\sqrt{\Delta_T})
 \le I_0e^{-\kappa \gap{\Sigma_0}t}.
\]
Together with smoothness, this gives
\begin{equation}\label{eq:force-tail}
 \norm{G_t}_{\op}=\norm{E_t}_{\op}
 \le\sqrt{2L\Delta_0}e^{-\kappa \gap{\Sigma_0}t},\qquad
 \int_t^T\norm{G_r}_{\op}\,\dd r\le I_0e^{-\kappa \gap{\Sigma_0}t}.
\end{equation}
Both $\Phi_t$ and $\Phi_t^{-1}$ have norm at most $e^{I_0}$, even when
$\Sigma_0$ is singular. Their congruence representation proves
\eqref{eq:gf-tube} before the maximal existence time. The covariance vector
field is locally Lipschitz on all symmetric matrices, and the upper bound
prevents finite-time blow-up; hence the solution is global.
Integrating $\norm{\dot\Sigma_t}_{\op}\le2B_\infty\norm{G_t}_{\op}$
proves the Cauchy property and \eqref{eq:gf-tail}. Integrability of the
force also makes $\Phi_t$ Cauchy; its limit is invertible since
$\sigma_{\min}(\Phi_t)\ge e^{-I_0}$.
Thus $\Sigma_\infty=\Phi_\infty\Sigma_0\Phi_\infty^\top$ has the same
rank as $\Sigma_0$ and is optimal by \eqref{eq:gf-rate}.
The singular-value bound for congruence gives the claimed positive
eigenvalue floor, without requiring full rank.
Directly in predictor Frobenius norm,
$\norm{\dSvu_t}_{\F}\le2B_\infty\norm{E_t}_{\F}$;
its integral proves \eqref{eq:gf-predictor-tail} without a dimension factor.
Finally $\norm{\Phi_t-\Phi_\infty}_{\op}\le e^{I_0}I_0e^{-\kappa \gap{\Sigma_0}t}$.
Coupling the two characteristic images of the same initial particle gives
the first bound in \eqref{eq:gf-w2}. For the second, use
\[
 \left(\int\norm{G_t w}^2\,\dd\mu_t(w)\right)^{1/2}
 =\sqrt{\tr(G_t\Sigma_tG_t)}
 \le\sqrt{2B_\infty}\,\norm{E_t}_{\F},
\]
and integrate along the same characteristics. The identical estimate
holds for empirical factors after division by $\sqrt n$.
If $\Delta_0=0$, the first-order flow is stationary and
the same formulas hold with $I_0=0$.
\end{proof}
\paragraph{Alternative nondegeneracy criteria}
The same argument can use any independently established hidden-block
lower bounds. The common-bound formulation below uses $2c$ in place of
$\gap{\Sigma_0}$. For $\Sigma_0\succ0$, taking
$c:=\min(\aplus{\Sigma_0},\aminus{\Sigma_0})$ recovers a possibly weaker
rate than Theorem~\ref{thm:gf}.

\begin{lemma}[Conditional hidden-block criterion]
\label{thm:uniform-nondegeneracy-convergence}
Under Assumption~\ref{ass:loss}, start \eqref{eq:bures-G} from
$\Sigma_0\succ0$. If $\Suu_t\succeq c\Id$ and $\Svv_t\succeq c\Id$
for some $c>0$ throughout its maximal existence interval, the flow is
global and
\begin{equation}
\label{eq:uniform-nondegeneracy-energy-rate}
\Delta_t\le\Delta_0e^{-4\kappa ct}.
\end{equation}
Moreover $\Sigma_t$ converges to $\Sigma_\infty\succ0$ with
$\Svu_\infty\in\cS_\star$.
\end{lemma}

Quadratic growth further gives
$\dist(\Svu_t,\cS_\star)^2\le2\Delta_0e^{-4\kappa ct}/\lambda$,
where $\lambda=\kappa$ is always available.

\begin{proof}[Proof of Lemma~\ref{thm:uniform-nondegeneracy-convergence}]
The assumed hidden bounds give
$-\dot\Delta\ge2c\norm E_{\F}^2\ge4c\kappa\Delta$.
On every positive loss gap interval, differentiating $\sqrt{\Delta_t}$ gives
\[
 -\frac{\mathrm{d}}{\mathrm{d}t}\sqrt{\Delta_t}
 \ge c\sqrt{2\kappa}\,\norm{E_t}_{\F},\qquad
 \int_t^T\norm{G_r}_{\op}\,\mathrm{d}r
 \le\frac{\sqrt{\Delta_t}-\sqrt{\Delta_T}}{c\sqrt{2\kappa}}.
\]
The same bound holds through zero loss gap by stationarity and is valid
before the maximal existence time.
Both characteristic matrices $\Phi_t$ and $\Phi_t^{-1}$ are bounded by
the exponential of this integral. Their covariance congruence gives
\begin{equation}\label{eq:uniform-nondegeneracy-covariance-bounds}
 \norm{\Sigma_t}_{\op}\le\norm{\Sigma_0}_{\op}
  \exp\!\left(\frac1c\sqrt{\frac{2\Delta_0}{\kappa}}\right),\qquad
 \lambda_{\min}(\Sigma_t)\ge\lambda_{\min}(\Sigma_0)
  \exp\!\left(-\frac1c\sqrt{\frac{2\Delta_0}{\kappa}}\right).
\end{equation}
These bounds rule out finite-time blow-up. The tail of the integral proves convergence to a
positive-definite covariance with optimal predictor.
Quadratic growth gives the distance bound.
\end{proof}

For coercive losses, determinant conservation and the second spectral
law yield an alternative, generally much coarser, source of such bounds.
\begin{proposition}[Coercivity prevents covariance collapse]
\label{prop:coercive-nondegeneracy}
Under Assumption~\ref{ass:loss}, suppose $\ell$ is coercive and
$\Sigma_0\succ0$. Let $M:=\sup\{\norm S_{\F}:\ell(S)\le\ell(\Svu_0)\}$.
On its maximal existence interval the Bures flow satisfies
$c\Id_m\preceq\Sigma_t\preceq B\Id_m$, where
$B:=\sqrt{\tr((J\Sigma_0)^2)+2M^2}+M$ and
$c:=\det\Sigma_0/B^{m-1}>0$.
In particular both hidden blocks are bounded below by $c\Id$.
\end{proposition}

\begin{proof}[Proof of Proposition~\ref{prop:coercive-nondegeneracy}]
The characteristic congruence \eqref{eq:orbit} preserves positive
definiteness on the maximal existence interval. Loss dissipation gives
$\norm{\Svu_t}_{\F}\le M$. Since $\tr G_t=0$, Jacobi's formula yields
$\det\Sigma_t=\det\Sigma_0>0$. The second spectral law gives
\[
 \norm{\Suu_t}_{\F}^2+\norm{\Svv_t}_{\F}^2
 =\tr((J\Sigma_0)^2)+2\norm{\Svu_t}_{\F}^2
 \le\tr((J\Sigma_0)^2)+2M^2.
\]
Split $\Sigma_t$ into its block-diagonal and off-diagonal parts. Their
operator norms are bounded by the square root of the last expression
and by $\norm{\Svu_t}_{\op}\le M$, respectively. Thus
$\norm{\Sigma_t}_{\op}\le B$. Finally,
$\det\Sigma_0\le\lambda_{\min}(\Sigma_t)B^{m-1}$ gives
$\Sigma_t\succeq c\Id_m$, and the same lower bound holds on both
principal blocks.
\end{proof}

\begin{remark}[Coarse ellipticity versus the spectral-margin rate]
\label{rem:coarse-rate}
At $\Sigma_0=\alpha\Id_m$, the determinant argument gives
\[
c_{\rm coarse}
:=\alpha\left(\frac{\alpha}
{\sqrt{m\alpha^2+2M^2}+M}\right)^{m-1}.
\]
Its rate $4\kappa c_{\rm coarse}$ can be extremely small in dimension,
whereas the support gap $\gap{\Sigma_0}=2\alpha$ gives $4\alpha\kappa$ directly.
For general initial covariances, the margin argument likewise gives a
rate determined by the signed spectrum, with no additional dimension factor.
Trace and initial loss alone cannot provide a uniform positive rate as
the support gap tends to zero
(Example~\ref{ex:slow}).
\end{remark}

\begin{remark}[Degenerate boundary leaves]\label{rem:degenerate-boundary-leaves}
In the $2\times2$ covariance setting of Section~\ref{sec:laws}, write
$\Sigma=\begin{psmallmatrix}p&r\\r&q\end{psmallmatrix}$,
$b_0:=(p_0-q_0)/2$, and $\alpha^2:=p_0q_0-r_0^2$,
with predictor loss $\ell(r):=(r-\beta)^2/2$.
On the boundary \(b_0=\alpha=0\), the scalar equation is
\(\dot r=-2|r|(r-\beta)\).
If \(\beta\ne0\), the ray with \(r_0\beta>0\) converges to \(\beta\),
whereas \(r_0\beta<0\) converges to the nonoptimal apex \(r=0\);
the apex itself is stationary. For example, when \(\beta=1\) and
\(p_0=q_0=1\), \(r_0=-1\), one has
\(p_t=q_t=-r_t=(2e^{2t}-1)^{-1}\).
When \(\beta=0\), one has \(r_t=r_0/(1+2|r_0|t)\): convergence to the
optimal apex is algebraic for \(r_0\ne0\).
These boundary cases have $\gap{\Sigma_0}=0$ and are not covered by the
positive-support-gap convergence theorem.
\end{remark}

\begin{example}[Fixed trace and predictor, arbitrarily slow escape]\label{ex:slow}
Let $d_u=d_v=2$, $0<\delta<1$, and
\[
 \Sigma_0=\diag(\delta,2-\delta,\delta,2-\delta),\qquad
 \ell(S):=\tfrac12\norm{S-e_1e_1^\top}_{\F}^2.
\]
Here $L=\kappa=1$, $\tr\Sigma_0=4$, $\Svu_0=0$, and $\Delta_0=1/2$,
independently of $\delta$. The active scalar block evolves by
\[
 r'=2\sqrt{r^2+\delta^2}(1-r),\qquad r(0)=0.
\]
The instantaneous relative decay of the loss at zero is $4\delta$.
Until $r=1/2$, $r'\le2(r+\delta)$, so
$r(t)\le\delta(e^{2t}-1)$. The time to reach $1/2$ is at least
$\tfrac12\log(1+1/(2\delta))$, which diverges as $\delta\downarrow0$.
There can therefore be no uniform positive exponential rate with a fixed
finite prefactor determined only by the displayed trace, initial loss gap, $L$, and
$\kappa$. At $\delta=0$ the active block is stationary and nonoptimal.
Theorem~\ref{thm:gf} records exactly the missing scale: $\gap{\Sigma_0}=2\delta$.
The slow escape occurs even for a strongly convex loss with a unique
predictor minimizer.
\end{example}


\subsection{Finite Width Convergence}\label{app:sampling}
Write $r_0:=\rankop\tSigma_0$ for the population covariance rank.
\begin{theorem}[Relative initialization stability]\label{thm:relative}
Under Assumption~\ref{ass:loss}, let $\tSigma_0\succeq0$ with
$\gap{\tSigma_0}>0$, and put $\widetilde B_0:=\norm{\tSigma_0}_{\op}$ and
$\widetilde\Delta_0:=\ell(\tSvu_0)-\ell_\star$.
If an empirical covariance satisfies
\begin{equation}\label{eq:relative}
 (1-\varepsilon)\tSigma_0\preceq\Sigma_0
 \preceq(1+\varepsilon)\tSigma_0,\qquad 0<\varepsilon<1,
\end{equation}
then its first-order factor flow is global, converges to an optimal
predictor, and its loss gap $\Delta_t=\ell(\Svu_t)-\ell_\star$ satisfies
\begin{equation}\label{eq:sample-rate}
 \Delta_t\le\Delta_0
 e^{-2(1-\varepsilon)\kappa \gap{\tSigma_0}t}.
\end{equation}
Only the lower inequality in \eqref{eq:relative} is needed for this claim.
With both inequalities, $d_\wedge:=\min(d_u,d_v)$ gives the deterministic
prefactor bound
\begin{equation}\label{eq:sample-prefactor}
 \Delta_0\le
 \left(\sqrt{\widetilde\Delta_0}+\varepsilon \widetilde B_0\sqrt{Ld_\wedge/2}\right)^2.
\end{equation}
For the sampling law in Corollary~\ref{cor:finite-width}, the width
condition \eqref{eq:width} guarantees \eqref{eq:relative} with probability
at least $1-\delta$.
\end{theorem}


\begin{proof}[Proof of Theorem~\ref{thm:relative}]
From $\tSigma_0-\aplus{\tSigma_0}J\succeq0$ and the lower relative inequality,
$\Sigma_0-(1-\varepsilon)\aplus{\tSigma_0}J\succeq0$, and similarly for
$\aminus{\tSigma_0}$. Transporting these inequalities along the empirical
fundamental matrix allows us to apply the proof of Theorem~\ref{thm:gf} with margin
sum $(1-\varepsilon)\gap{\tSigma_0}$ and initial empirical loss gap. The two-sided relative
comparison gives
$\norm{\Sigma_0-\tSigma_0}_{\op}\le\varepsilon \widetilde B_0$ and hence
$\norm{\Svu_0-\tSvu_0}_{\F}\le\sqrt{d_\wedge}\varepsilon \widetilde B_0$.
Equation~\eqref{eq:sample-prefactor} follows from \eqref{eq:perturb}.
For the probability bound, choose an orthonormal basis
$T\in\R^{m\times r_0}$ of the range of $\tSigma_0$ and write
$\tSigma_0=T\Lambda T^\top$, where $\Lambda\succ0$.
The vectors $\eta_i:=T^\top\xi_i\in\R^{r_0}$ have identity second
moment and subgaussian norm at most $K$, and
$w_i=T\Lambda^{1/2}\eta_i$.
The covariance tail bound~\citep[Exercise 4.7.3]{Vershynin2018}
applied in this range gives
\[
 \norm{n^{-1}\textstyle\sum_i\eta_i\eta_i^\top-\Id_{r_0}}_{\op}
 \le CK^2\left(\sqrt{\frac{r_0+u}{n}}+\frac{r_0+u}{n}\right)
\]
with failure probability at most $2e^{-u}$. Set $u:=\log(2/\delta)$ and
adjust $C$ in \eqref{eq:width}. Multiplication on the left by
$T\Lambda^{1/2}$ and on the right by its transpose gives
\eqref{eq:relative}.
\end{proof}
\begin{proof}[Proof of Corollary~\ref{cor:finite-width}]
Theorem~\ref{thm:relative} gives the relative covariance event and the
loss rate. Its lower bound gives
$\gap{\Sigma_0}\ge(1-\varepsilon)\gap{\tSigma_0}>0$,
so Theorem~\ref{thm:gf} gives global existence and factor convergence.
The two-sided event preserves the kernel of $\tSigma_0$, so the empirical
initial covariance is positive definite exactly when $\tSigma_0$ is.
\end{proof}

\paragraph{Small width under full-rank isotropic initialization}
The rank improvement in Remark~\ref{rem:rank-adaptive-width} concerns
the chosen population law. A positive population support gap alone does not
guarantee a positive empirical support gap at small width when
$\tSigma_0=\sigma^2\Id_m$, $\sigma>0$.
\begin{example}[Isotropic Gaussian sampling can lose the support gap]
\label{ex:white-small-width}
Let $d_u=D$, $d_v=1$, and sample from $N(0,\tSigma_0)$ with
$\tSigma_0=\sigma^2\Id_{D+1}$, $\sigma>0$.
There is a universal $c>0$ such that, if
$1\le n\le D/16$, then
$\mathbb P\{\gap{\Sigma_0}=0\}\ge1-4e^{-cD}$,
although $\gap{\tSigma_0}=2\sigma^2$.
Thus a width bound depending only on $\min(d_u,d_v)$ cannot uniformly
preserve a positive support gap when sampling from $N(0,\sigma^2\Id_m)$. Convergence of
these trajectories is not decided by this example.
\end{example}
\begin{proof}
Write $U_0=\sigma X$ and $V_0=\sigma y^\top$, where
$X\in\R^{D\times n}$ and $y\in\R^n$ have independent standard
Gaussian entries. Gaussian singular-value and norm concentration
\citep[Section~7.3]{Vershynin2018}, with deviation $\sqrt D/8$, gives
\[
 \sigma_{\min}(X)\ge\sqrt D-\sqrt n-\sqrt D/8\ge5\sqrt D/8,
 \qquad \norm y\le\sqrt n+\sqrt D/8\le3\sqrt D/8
\]
with probability at least $1-4e^{-cD}$, for a universal $c>0$.
Hence $X^\top X-yy^\top\succeq(D/4)\Id_n$.
The nonzero spectrum of $J\Sigma_0$ is the spectrum of
$\sigma^2(X^\top X-yy^\top)/n$, so it has $n$ positive eigenvalues
and $D+1-n\ge2$ zero eigenvalues. In \eqref{eq:gap-spectrum},
$d_v=1$ therefore gives $\gap{\Sigma_0}=\lambda_2-\lambda_1=0$.
\end{proof}

\paragraph{A separate initial-accuracy budget}
A relative covariance bound preserves the rate; an absolute bound also
controls the initial loss. In the sampling setting of
Corollary~\ref{cor:finite-width}, write
$\widetilde B_0:=\norm{\tSigma_0}_{\op}$, $\widetilde\Delta_0:=\ell(\tSvu_0)-\ell_\star$, and
$d_\wedge:=\min(d_u,d_v)$.
Besides $\varepsilon,\delta\in(0,1)$, choose an absolute tolerance
$\zeta>0$ and set
\begin{equation}\label{eq:sampling-budgets}
 \beta:=\min\{\varepsilon,\zeta/\widetilde B_0\},\qquad
 \overline\Delta:=
 \left(\sqrt{\widetilde\Delta_0}+\zeta\sqrt{Ld_\wedge/2}\right)^2.
\end{equation}

\begin{theorem}[High-probability finite-width convergence]
\label{thm:sampling-budgets}
Under Assumption~\ref{ass:loss}, let $\tSigma_0\succeq0$ with
$\gap{\tSigma_0}>0$ and use the sampling law and tolerances above.
There is a universal $C>0$ such that, if
\begin{equation}\label{eq:budget-width}
 n\ge CK^4\beta^{-2}\bigl(r_0+\log(2/\delta)\bigr),
\end{equation}
then with probability at least $1-\delta$ the empirical covariance has
positive support gap, the factor flow is global, and
\[
 \Delta_t\le\Delta_0e^{-2\kappa(1-\varepsilon)\gap{\tSigma_0}t}
 \le\overline\Delta e^{-2\kappa(1-\varepsilon)\gap{\tSigma_0}t}.
\]
The factors converge, their limiting predictor belongs to $\cS_\star$, and
$\dist(\Svu_t,\cS_\star)^2\le2\Delta_t/\kappa$.
\end{theorem}

\begin{proof}[Proof of Theorem~\ref{thm:sampling-budgets}]
Apply the relative concentration event with tolerance $\beta$.
It implies the relative bounds with tolerance $\varepsilon$ and
$\norm{\Sigma_0-\tSigma_0}_{\op}\le\beta \widetilde B_0\le\zeta$.
The two initial supporting inequalities give the margin bounds, and
Lemma~\ref{lem:pl} gives the loss gap bound. More precisely, putting
$\underline g:=(1-\varepsilon)\gap{\tSigma_0}$ and $\overline B_0:=(1+\varepsilon)\widetilde B_0$,
the same event gives
\begin{equation}\label{eq:sampling-event}
 a_\pm(\Sigma_0)\ge(1-\varepsilon)a_\pm(\tSigma_0),\quad
 \norm{\Sigma_0}_{\op}\le\overline B_0,\quad
 \Delta_0\le\overline\Delta.
\end{equation}
Theorem~\ref{thm:gf} then applies
to the empirical covariance. The distance estimate follows from
Lemma~\ref{lem:pl}. More generally, any valid quadratic-growth constant
$\lambda$ gives $\dist(\Svu_t,\cS_\star)^2\le2\Delta_t/\lambda$.
\end{proof}
\begin{proposition}[The Gaussian initial-error cost]\label{prop:gaussian}
For independent $u_i\sim N(0,\sigma^2\Id_{d_u})$ and
$v_i\sim N(0,\sigma^2\Id_{d_v})$, and the quadratic loss
$\ell(S):=\norm{S-\Svu_\star}_{\F}^2/2$, one has exactly
\begin{equation}\label{eq:gaussian-error}
 \mathbb E\Delta_0
 =\tfrac12\norm{\Svu_\star}_{\F}^2+\frac{\sigma^4d_ud_v}{2n}.
\end{equation}
\end{proposition}

For balanced dimensions and fixed $\sigma$, width $n$ proportional to $m$
therefore gives an expected initial sampling error of order $m$, despite
the dimension-free exponent $4(1-\varepsilon)\sigma^2\kappa$ at any
fixed relative tolerance $\varepsilon\in(0,1)$. Keeping that error bounded
in expectation requires $n$ of order $d_ud_v$, when the teacher norm is
fixed. Thus width proportional to $m$ can preserve the convergence rate
while allowing the expected initial error to grow with dimension.

\begin{proof}[Proof of Proposition~\ref{prop:gaussian}]
Each entry of $n^{-1}\sum_i v_i u_i^\top$ has mean zero and variance
$\sigma^4/n$. Sum these variances; the cross term with the teacher has
zero expectation.
\end{proof}

\subsection{Discrete-Time Convergence}\label{app:gd}
The factor update transports the initial supporting inequalities even
though it no longer preserves $J$ exactly. To track this error, define
\begin{equation}\label{eq:gd-signature}
\Phi_0:=\Id,\quad\Phi_{k+1}:=M_k\Phi_k,\quad
\mathcal J_k:=\Phi_kJ\Phi_k^\top,\quad\mathcal D_k:=\mathcal J_k-J.
\end{equation}
The cross block of \eqref{eq:gd-covariance} also gives
\begin{equation}\label{eq:gd-predictor-update}
 \Svu_{k+1}-\Svu_k
 =-\tau(E_k\Suu_k+\Svv_kE_k)+\tau^2E_k(\Svu_k)^\top E_k.
\end{equation}
The next two lemmas control the accumulated defect $\mathcal D_k$ and
the quadratic correction in the predictor update.

\begin{lemma}[Transported support and accumulated defect]\label{lem:gd-support}
If $\Sigma_0\succeq0$, with initial margins $\aplus{\Sigma_0},\aminus{\Sigma_0}$ as in
\eqref{eq:margins}, then every finite iterate satisfies
\[
 \Sigma_k-\aplus{\Sigma_0}\mathcal J_k\succeq0,\qquad
 \Sigma_k+\aminus{\Sigma_0}\mathcal J_k\succeq0.
\]
If $\norm{\mathcal D_k}_{\op}\le\eta$ for some $\eta\in[0,1)$, then
$\Suu_k\succeq(1-\eta)\aplus{\Sigma_0}\Id$ and
$\Svv_k\succeq(1-\eta)\aminus{\Sigma_0}\Id$. With $e_k:=\norm{E_k}_{\F}$,
\begin{equation}\label{eq:gd-defect-recursion}
 \mathcal D_{k+1}=M_k\mathcal D_kM_k+\tau^2G_kJG_k,\qquad
 \norm{\mathcal D_N}_{\op}
 \le\tau\exp\left(2\sum_{j<N}\tau e_j\right)
             \sum_{j<N}\tau e_j^2.
\end{equation}
\end{lemma}

\begin{proof}[Proof of Lemma~\ref{lem:gd-support}]
The proof of Proposition~\ref{lem:margins} gives the initial supporting
inequalities. Congruence by $\Phi_k$ transports them, regardless of
invertibility. The $uu$ and $vv$ blocks of $\mathcal J_k$ are respectively
at least $(1-\eta)\Id$ and at most $-(1-\eta)\Id$, proving the hidden
lower bounds. Since $G_kJ+JG_k=0$, one has
$M_kJM_k=J+\tau^2G_kJG_k$, giving the recursion. Its norm is at most
$(1+\tau e_k)^2\norm{\mathcal D_k}_{\op}+\tau^2e_k^2$.
Unroll from zero and bound the amplification products by
$\exp(2\sum_{j<N}\tau e_j)$.
\end{proof}
\paragraph{Descent and length.}
The bootstrap needs two estimates: loss descent and a finite total
gradient length. The next lemma obtains both from the same step condition.
For the proof, abbreviate the descent scale, initial gradient bound, and
covariance bound as
\begin{equation}\label{eq:gd-constants}
 \alpha_\star:=\frac{\gap{\Sigma_0}}{2\sqrt2},\qquad
 M_{\nabla}:=\sqrt{2L\Delta_0},\qquad
 B_\star:=\norm{\Sigma_0}_{\op}e^{2I_\star}.
\end{equation}
Then $I_\star=\sqrt{2\Delta_0/\kappa}/\alpha_\star$ and
$q=1-2\alpha_\star\kappa\tau$. Condition~\eqref{eq:gd-step} is
equivalent to requiring both
$\tau\le\alpha_\star/(B_\star M_{\nabla}+\frac92LB_\star^2)$ and
$\tau\le(1-1/\sqrt2)\alpha_\star/(e^{2I_\star}\Delta_0)$.
\begin{lemma}[Descent with unequal hidden lower bounds]\label{lem:gd-descent}
Suppose $\norm{\Sigma_k}_{\op}\le B$, $\Delta_k\le\Delta_0$,
$\Suu_k\succeq b_+\Id$, and $\Svv_k\succeq b_-\Id$.
Write $h:=b_++b_->c>0$. If $\tau M_{\nabla}\le1$ and
$\tau(BM_{\nabla}+\frac92LB^2)\le h-c$, then
\begin{equation}\label{eq:gd-descent}
 \Delta_{k+1}\le\Delta_k-c\tau e_k^2
 \le(1-2c\kappa\tau)\Delta_k.
\end{equation}
If this descent holds for $j=r,\ldots,N-1$, then
\begin{equation}\label{eq:gd-length}
 \sum_{j=r}^{N-1}\tau e_j^2\le\Delta_r/c,\qquad
 \sum_{j=r}^{N-1}\tau e_j
 \le\frac1c\sqrt{\frac2\kappa}
            (\sqrt{\Delta_r}-\sqrt{\Delta_N}).
\end{equation}
\end{lemma}

\begin{proof}[Proof of Lemma~\ref{lem:gd-descent}]
Let $Z:=\Svu_{k+1}-\Svu_k$. By \eqref{eq:gd-covariance} and smoothness,
$\norm Z_{\F}\le3\tau Be_k$ and
$\ip{E_k}{Z}\le-\tau h e_k^2+\tau^2BM_{\nabla}e_k^2$.
The descent lemma therefore gives the first assertion. Telescope it for
the squared-gradient sum. At positive loss gap, PL gives
\[
 \sqrt{\Delta_j}-\sqrt{\Delta_{j+1}}
 \ge\frac{c\tau e_j^2}{2\sqrt{\Delta_j}}
 \ge c\tau\sqrt{\kappa/2}\,e_j.
\]
At zero loss gap $e_j=0$. Summation proves the length bound.
\end{proof}

The proof below also establishes the following explicit tails in
Theorem~\ref{thm:gd}:
\begin{equation}\label{eq:gd-tails}
 \frac{\norm{W_k-W_\infty}_{\F}}{\sqrt n}
 \le\sqrt{2B_\star}\,I_\star q^{k/2},\qquad
 \norm{\Sigma_k-\Sigma_\infty}_{\op}
 \le2B_\star I_\star q^{k/2}.
\end{equation}

\begin{proof}[Proof of Theorem~\ref{thm:gd}]
Let $\eta:=1-1/\sqrt2$. Since $\gap{\Sigma_0}\le2B_0$ for
$B_0:=\norm{\Sigma_0}_{\op}$, one has
$\alpha_\star\le B_0/\sqrt2\le B_\star/\sqrt2$.
The first restriction in \eqref{eq:gd-step} implies
\[
 \tau M_{\nabla}<1/\sqrt2,\qquad
 0<2\alpha_\star\kappa\tau
 \le\frac{4\kappa\alpha_\star^2}{9LB_\star^2}\le\frac29.
\]
In particular $q\in[7/9,1)$.
Inductively suppose, through iterate $N$, that
$\Delta_k\le q^k\Delta_0$,
$\norm{\Sigma_k}_{\op}\le B_\star$, and
$\norm{\mathcal D_k}_{\op}\le\eta$, with descent proved for $k<N$.
These assertions start at $N=0$ since $\mathcal D_0=0$.
At $N$, Lemma~\ref{lem:gd-support} supplies hidden lower bounds with sum
$(1-\eta)\gap{\Sigma_0}=2\alpha_\star$. Lemma~\ref{lem:gd-descent} applies with
$c=\alpha_\star$, $h=2\alpha_\star$, and $B=B_\star$, proving the next
loss-gap estimate. Its finite-horizon sums now give
\[
 \sum_{j=0}^{N}\tau e_j\le I_\star,\qquad
 \sum_{j=0}^{N}\tau e_j^2\le\Delta_0/\alpha_\star.
\]
Thus $\norm{\Phi_{N+1}}_{\op}\le e^{I_\star}$ implies the next covariance
upper bound, and \eqref{eq:gd-defect-recursion} gives
$\norm{\mathcal D_{N+1}}_{\op}
\le\tau e^{2I_\star}\Delta_0/\alpha_\star\le\eta$.
This closes the induction.

Taking tails in \eqref{eq:gd-length} yields
$\sum_{j\ge k}\tau e_j\le I_\star q^{k/2}$.
Factor increments have norm
\[
 \norm{W_{j+1}-W_j}_{\F}/\sqrt n
=\tau\sqrt{\tr(G_j\Sigma_jG_j)}
\le\tau\sqrt{2B_\star}\,e_j.
\]
This proves factor convergence and its tail estimate. Similarly
$\norm{\Phi_k-\Phi_\infty}_{\op}\le e^{I_\star}I_\star q^{k/2}$
gives the covariance tail bound by subtracting the two congruences.
The limit predictor is optimal by continuity and the loss-gap bound.
Finally, if $r:=\tau M_{\nabla}<1$, every factor update is invertible and
\[
 \sigma_{\min}(\Phi_k)\ge\exp\left(-\frac{I_\star}{1-r}\right),\qquad
 \lambda_{\min}^{+}(\Sigma_k)\ge\lambda_{\min}^{+}(\Sigma_0)
 \exp\left(-\frac{2I_\star}{1-r}\right).
\]
Indeed $\sigma_{\min}(M_j)\ge1-\tau e_j>0$ and
$\log(1-x)\ge-x/(1-r)$ for $0\le x\le r$. The bound makes
$\Phi_\infty$ invertible, proving rank preservation also at the limit.
\end{proof}
The factor-increment estimate also applies to population discrete
pushforward updates. Coupling each particle with its limiting position
gives a Wasserstein tail at most $\sqrt{2B_\star}I_\star q^{k/2}$.

\paragraph{Isotropic initialization and local stability}
For an exact isotropic empirical second moment
$\Sigma_0=W_0W_0^\top/n=\sigma^2\Id_m$ and $\Delta_0=\ell(0)-\ell_\star>0$,
$I_\star=2\sqrt{\Delta_0}/(\sigma^2\sqrt\kappa)$ and the explicit
condition is
\[
 \tau\le\min\left\{
 \frac1{2e^{2I_\star}\sqrt{L\Delta_0}
                 +\frac9{\sqrt2}L\sigma^2e^{4I_\star}},\quad
 \frac{(\sqrt2-1)\sigma^2}{2e^{2I_\star}\Delta_0}
 \right\}.
\]
The contraction factor is $1-\sqrt2\kappa\sigma^2\tau$.
Here the empirical covariance is exactly $\sigma^2\Id_m$. For iid Gaussian
initialization, its fluctuations are handled by Corollary~\ref{cor:sampled-gd}.
As $\Delta_0\downarrow0$ with $\sigma,L,\kappa$ fixed, the sufficient
threshold tends to $\sqrt2/(9L\sigma^2)$.


\begin{remark}[Local stability and the global stepsize]\label{rem:gd-local}
Suppose $\ell$ is $C^2$ near a minimizer at zero, and linearize at
$\Sigma=\sigma^2\Id$. The predictor perturbation evolves to first order by
$\delta\Svu\mapsto(\Id-2\tau\sigma^2\nabla^2\ell(0))[\delta\Svu]$.
Thus $0<\tau<1/(L\sigma^2)$ contracts its nonzero-curvature directions;
kernel directions are neutral at this order. Theorem~\ref{thm:gd} additionally controls
covariance growth and the accumulated signature defect along the entire
trajectory, explaining its more restrictive global threshold.
\end{remark}

\paragraph{Adjustable rate budgets.}
Smaller steps permit a smaller loss in the conserved conditioning,
bringing the discrete exponent closer to its continuous-time value.
\begin{proposition}[Adjustable rate budgets]\label{prop:gd-adjustable}
Under the loss and initialization assumptions of Theorem~\ref{thm:gd},
choose $0<c<h<\gap{\Sigma_0}$ and set
$I_c:=\sqrt{2\Delta_0/\kappa}/c$, $B_c:=B_0e^{2I_c}$.
The convergence and rank conclusions of Theorem~\ref{thm:gd} hold with
contraction factor $q_c:=1-2c\kappa\tau$, and the tail bounds
\eqref{eq:gd-tails} hold with $(I_\star,B_\star,q)$ replaced by
$(I_c,B_c,q_c)$, if
\[
 0<\tau\le\min\left\{\frac1{2M_{\nabla}},\
 \frac{h-c}{B_cM_{\nabla}+\frac92LB_c^2},\
 \frac{(1-h/\gap{\Sigma_0})c}{e^{2I_c}\Delta_0}\right\}.
\]
\end{proposition}

Taking $c=(1-\epsilon)\gap{\Sigma_0}$, $h=(1-\epsilon/2)\gap{\Sigma_0}$ for any
$\epsilon\in(0,1)$ certifies exponent $2\kappa(1-\epsilon)\gap{\Sigma_0}$ in time
$k\tau$, at an explicit sufficiently small step. This is arbitrarily
close to the flow exponent.

\begin{proof}[Proof of Proposition~\ref{prop:gd-adjustable}]
Use $\eta=1-h/\gap{\Sigma_0}$ in the preceding induction and
Lemma~\ref{lem:gd-descent} with $c,h$. The positivity guard follows from
$c(h-c)\le \gap{\Sigma_0}^2/4\le B_0^2$: $2c\kappa\tau\le4/9$.
The remaining estimates use $I_c,B_c$.
\end{proof}
\begin{corollary}[A stepsize fixed before sampling]\label{cor:sampled-gd}
Use the population parameters and width condition of
Theorem~\ref{thm:sampling-budgets}, with the bounds in
\eqref{eq:sampling-event}, and define before sampling
\[
 \overline\alpha_\star:=\frac{\underline g}{2\sqrt2},\quad
 \overline M_{\nabla}:=\sqrt{2L\overline\Delta},\quad
 \overline I_\star:=\frac1{\overline\alpha_\star}
                    \sqrt{\frac{2\overline\Delta}{\kappa}},\quad
 \overline B_\star:=\overline B_0e^{2\overline I_\star}.
\]
Choose a deterministic positive stepsize satisfying
\begin{equation}\label{eq:sampled-gd-step}
 \tau\le\min\left\{
 \frac{\overline\alpha_\star}
 {\overline B_\star\overline M_{\nabla}+\frac92L\overline B_\star^2},\quad
 \frac{(1-1/\sqrt2)\overline\alpha_\star}
 {e^{2\overline I_\star}\overline\Delta}\right\}.
\end{equation}
With probability at least $1-\delta$, factor GD converges to an optimal
predictor, and
\[
 \Delta_k
 \le(1-2\overline\alpha_\star\kappa\tau)^k\Delta_0
 \le(1-2\overline\alpha_\star\kappa\tau)^k\overline\Delta.
\]
Its factor/covariance tail estimates hold with the corresponding barred
constants and contraction factor.
\end{corollary}

At fixed relative tolerance, the sampled stepsize can deteriorate with
the dimension through $\overline\Delta$. The independent absolute
accuracy budget controls this initial-loss contribution.


\begin{proof}[Proof of Corollary~\ref{cor:sampled-gd}]
On \eqref{eq:sampling-event}, the empirical descent scale is at least
$\overline\alpha_\star$, whereas the empirical loss gap, force bound, length
bound, and covariance upper bound are at most their barred counterparts.
Both admissible stepsize terms increase with the descent scale and
decrease with the loss gap and initial covariance norm. Hence the chosen
step satisfies Theorem~\ref{thm:gd} for the empirical initialization.
Its contraction factor is no larger than the stated one. If the realized
initial loss gap vanishes, the empirical factors are stationary, and the
conclusions hold directly.
\end{proof}

\paragraph{Boundary cases.}
For the scalar loss $\ell(s):=(s+1)^2/2$, the singular initialization
$U_0=V_0\ne0$ preserves equality of the factors and keeps
$\Svu_k=\norm{U_k}^2/n\ge0$, so its loss gap is at least $1/2$.
Conversely, even full-rank isotropic initialization does not justify arbitrary steps:
for $\ell(s):=(s-1)^2/2$, $\Sigma_0=\Id_2$, and $\tau=1$, one gets
$\Sigma_1=\bigl(\begin{smallmatrix}2&2\\2&2\end{smallmatrix}\bigr)$
and $\Sigma_2=0$, after which training is trapped.

\subsection{Deep Linear Residual Networks}\label{app:resnets}
Write
$P_{b:a}:=A_b\cdots A_a$ when $a\le b$, and $P_{b:a}:=\Id$ otherwise.
Thus $P_{>s}=P_{S-1:s+1}$ and $P_{<s}=P_{s-1:0}$.
Differentiating $P_S=P_{>s}A_sP_{<s}$ gives
\begin{equation}
  \label{eq:resnet-layer-gradient}
\nabla_{\Svu_s}\ell(P_S)
=h\widetilde E_s,
\qquad
\widetilde E_s=P_{>s}^\top E P_{<s}^\top.
\end{equation}
In residual mean-field time the particle equations are
$(\dot u_{s,i},\dot v_{s,i})
=-(\widetilde E_s^\top v_{s,i},\widetilde E_su_{s,i})$.
The cross block of \eqref{eq:resnet-layer-bures} is therefore
$\dSvu_s=-\widetilde E_s\Suu_s-\Svv_s\widetilde E_s$.

In the finite-depth setting of Theorem~\ref{thm:resnet-finite}, put
$\overline{\sigma}_S:=(1+B/S)^{S-1}$ and
$I_{\rm res}:=\overline{\sigma}_S(\overline g_S\underline{\sigma}_S^2)^{-1}\sqrt{2\Delta_0/\kappa}$.
The proof gives the common layer-force tail
\begin{equation}\label{eq:resnet-length}
 \int_t^\infty\norm{\widetilde E_s(r)}_{\F}\,\dd r
 \le I_{\rm res}e^{-\kappa\overline g_S\underline{\sigma}_S^2t}.
\end{equation}

\begin{proof}[Proof of Theorem~\ref{thm:resnet-finite}]
Write $g_s:=\gap{\Sigma_s(0)}$ within this proof.
The endpoint derivative and layer covariance equations give
\[
 -\dot\Delta
 =\frac1S\sum_{s=0}^{S-1}
 \left[\tr(\widetilde E_s^\top \widetilde E_s\Suu_s)+\tr(\widetilde E_s\widetilde E_s^\top\Svv_s)\right]
 \ge\frac1S\sum_{s=0}^{S-1}g_s\norm{\widetilde E_s}_{\F}^2.
\]
Both the lower and upper singular-value bounds of the residual factors
multiply through the prefix and suffix, so
$\underline{\sigma}_S\norm E_{\F}\le\norm{\widetilde E_s}_{\F}\le \overline{\sigma}_S\norm E_{\F}$.
Therefore
\[
 -\dot\Delta\ge\overline g_S\underline{\sigma}_S^2\norm E_{\F}^2
 \ge2\kappa\overline g_S\underline{\sigma}_S^2\Delta.
\]
This proves the loss bound on each finite existence interval using the
average layer margin. The square-root PL
argument used in Theorem~\ref{thm:gf}, followed by the upper force bound,
gives \eqref{eq:resnet-length} first with any finite endpoint.
Thus all layer fundamental matrices and their inverses have norm at most
$e^{I_{\rm res}}$. Their covariances remain bounded and positive definite.
The finite-dimensional vector field is locally Lipschitz; these bounds
prevent a finite maximal existence time. The force-tail bound now holds
on $[t,\infty)$ and implies layer convergence by the argument of
Theorem~\ref{thm:gf}. Continuity of the finite residual product proves
optimality of its limit.
\end{proof}

\begin{proof}[Proof of Theorem~\ref{thm:resnet-continuous}]
The $L^1$ training equation in Section~\ref{sec:resnets} means
\begin{equation}\label{eq:resnet-integral}
 \Sigma_s(t)=\Sigma_s(0)-\int_0^t
 (\widetilde G_s(r)\Sigma_s(r)+\Sigma_s(r)\widetilde G_s(r))\,\dd r,
 \qquad\widetilde G_s=G_{\widetilde E_s}.
\end{equation}
This equality in $L^1$ gives an absolutely continuous representative for
almost every layer. The path bound makes the forces uniformly bounded
over depth on finite training intervals, justifying differentiation
of the endpoint and integration over depth.
An $L^1$ residual field has invertible propagators. Gronwall applied to
each propagator and its inverse gives
$e^{-B}\norm E_{\F}\le\norm{\widetilde E_s}_{\F}\le e^B\norm E_{\F}$.
The depth integrals of the conserved supporting inequalities yield
\[
 -\dot\Delta
 =\int_0^1\left[\tr(\widetilde E_s^\top \widetilde E_s\Suu_s)
                    +\tr(\widetilde E_s\widetilde E_s^\top\Svv_s)\right]\,\dd s
 \ge\overline g e^{-2B}\norm E_{\F}^2.
\]
These integrals are finite by the $L^1$ covariance assumption and the
uniform-in-depth force bounds. PL proves \eqref{eq:continuous-resnet-rate}.
As $\gap{\Sigma_s(0)}>0$ almost everywhere and
$\gap{\Sigma_s(0)}\le2\norm{\Sigma_s(0)}_{\op}$, its integral is finite and positive.
The square-root argument gives the common layer-force tail bound
\[
 \int_t^\infty\norm{\widetilde E_s(r)}_{\F}\,\dd r
 \le I_{\rm res}e^{-\kappa\overline g e^{-2B}t},\qquad
 I_{\rm res}:=\frac{e^{3B}}{\overline g}
                         \sqrt{\frac{2\Delta_0}{\kappa}}.
\]
The fundamental-matrix argument proves almost-everywhere convergence
and positive definiteness. The covariance tails are bounded by
$2e^{2I_{\rm res}}\norm{\Sigma_s(0)}_{\op}I_{\rm res}
e^{-\kappa\overline g e^{-2B}t}$, an integrable bound over depth.
This proves $L^1$ convergence. Continuity of the endpoint map on bounded
$L^1$ sets, obtained by variation of constants and Gronwall, proves
optimality of the limiting endpoint.
\end{proof}

\begin{remark}[Sharpness of the boundedness hypothesis]
\label{rem:resnet-path-escape}
Unlike the two-layer flow of Theorem~\ref{thm:gf}, positive conserved
layer margins do not rule out unbounded weight growth as $S\to\infty$
or along continuous-depth training. A linear neural ODE has positive
determinant, so it cannot fit an orientation-reversing target; training
may instead approach a singular endpoint through unbounded residuals.
For a concrete example, take the invertible scalar target $p_\star:=-1$ and loss
$\ell(p):=(p+1)^2/2$, with $L=\kappa=1$. Initialize every layer with
$\Sigma_s(0)=\Id_2$, realizable at width $n=2$; every conserved support
gap equals $2$. At each finite depth $S\ge2$, all residuals decrease
to $-S$, so $\sup_{t\ge0}\max_s|\Svu_s(t)|=S$ and no uniform $B<S$
exists. In continuous depth they decrease without bound, giving
$\sup_{t\ge0}\int_0^1|\Svu_s(t)|\,\dd s=\infty$.
Both flows are global, but their endpoints tend to zero and their losses
to $1/2>\ell_\star=0$. At finite depth the target is representable:
failure occurs along this symmetric trajectory, whereas the neural ODE
excludes the target altogether. In both cases, positive conserved margins
therefore fail to ensure the residual bounds required by the theorems.
\end{remark}

\begin{proof}[Verification of Remark~\ref{rem:resnet-path-escape}]
The scalar target $-1$ is invertible, and
$|\ell'(p)|^2=(p+1)^2=2\ell(p)$, so $L=\kappa=1$ and
$\ell_\star=0$. The width-two initialization
$U_s(0)=(\sqrt2,0)$, $V_s(0)=(0,\sqrt2)$ gives
$\Sigma_s(0)=\Id_2$ and initial endpoint $1$ at every depth.
For scalar layers write $a_s:=\Suu_s=\Svv_s$ and
$r_s:=\Svu_s$. Equality of the diagonal blocks is preserved, and the
layer equations reduce to
\[
 \dot a_s=-2\widetilde E_s r_s,\qquad
 \dot r_s=-2\widetilde E_s a_s.
\]
Thus $a_s^2-r_s^2=1$, and positivity gives
\[
 \Sigma_s(t)=\begin{pmatrix}
  \sqrt{1+r_s(t)^2}&r_s(t)\\
  r_s(t)&\sqrt{1+r_s(t)^2}
 \end{pmatrix}.
\]
The eigenvalues of $J\Sigma_s$ are $\pm1$, whereas those of
$\Sigma_s$ are $\sqrt{1+r_s^2}\pm|r_s|>0$.
Uniqueness and symmetry preserve identical layers. Denote their common
residual by $r_S$. At finite depth,
$\widetilde E_s=(p_S+1)(1+r_S/S)^{S-1}$, giving
\begin{equation}\label{eq:resnet-escape-finite}
 \begin{aligned}
 p_S&=(1+r_S/S)^S,\qquad r_S(0)=0,\\
 \dot r_S&=-2\sqrt{1+r_S^2}\big((1+r_S/S)^S+1\big)(1+r_S/S)^{S-1}.
 \end{aligned}
\end{equation}
For $S\ge2$, its vector field is locally Lipschitz and vanishes at $r_S=-S$.
Starting at zero, the solution decreases and cannot reach this equilibrium
in finite time. It stays in the compact interval $[-S,0]$, so it is global.
Any limit strictly above $-S$ would have strictly negative velocity;
hence its limit is $-S$, $p_S\to0$, and $\ell(p_S)\to1/2$.
In particular, every candidate bound $B<S$ is exceeded at a finite
training time, although $|r_S(t)|<S$ at every finite time.

The finite-depth architecture can represent $-1$ by taking one residual
factor equal to $-1$ and the others equal to $1$. The obstruction here
is the preserved symmetry of the initialized trajectory: all its
residual factors stay positive and approach zero. At $S=1$, this
obstruction disappears; $\dot r_1=-2\sqrt{1+r_1^2}(r_1+2)$ instead
gives $r_1\to-2$ and the optimal endpoint $-1$.

For a depth-constant field, the continuous propagators are
$P_{b,a}=e^{(b-a)r}$, so $\widetilde E_s=(e^r+1)e^r$ and
\begin{equation}\label{eq:resnet-escape-continuous}
 p=e^r,\qquad \dot r=-2\sqrt{1+r^2}(e^r+1)e^r,\qquad r(0)=0.
\end{equation}
For $r\le0$ its velocity is bounded, strictly negative, and has no finite
zero. The solution is therefore global and tends to $-\infty$.
Since $\int_0^1|\Svu_s(t)|\,\dd s=|r(t)|$, every finite integral
bound is eventually exceeded as well. The endpoint tends to zero.
The limiting loss $1/2$ is the infimum over positive scalar endpoints
and exceeds the ambient minimum $\ell_\star=0$.
The associated depth-constant covariance is a global locally absolutely
continuous $L^1$ trajectory, positive definite at every finite time.
Moreover, for $S\ge2$ and $-S<r_S\le0$,
$(1+r_S/S)^{S-1}\le e^{r_S/2}$ and $(1+r_S/S)^S+1\le2$.
Since $\sqrt{1+r^2}e^{r/2}\le1$ for $r\le0$, both $|\dot r_S|$ and
$|\dot r|$ are at most $4$. On $[0,T]$, all residuals therefore lie in
$[-4T,0]$. The finite-depth vector fields converge uniformly on this
compact interval, and continuous dependence gives $r_S\to r$ uniformly
on every bounded training interval $[0,T]$.
\end{proof}

\begin{remark}[Zero-target variant]
The same path bounds can fail while the loss tends to zero. With
$\ell(p):=p^2/2$ and the same initialization, the equations are
$\dot r_S=-2\sqrt{1+r_S^2}(1+r_S/S)^{2S-1}$ and
$\dot r=-2\sqrt{1+r^2}e^{2r}$. The preceding argument gives
$r_S\to-S$ for every $S\ge1$, $r\to-\infty$ in continuous depth,
and endpoints tending to the zero target. A uniform residual-path
bound would keep the endpoints bounded away from zero instead.
\end{remark}

\begin{remark}[Residual-path control and orientation]
Positive-definite layer covariances can coexist with singular residual
products. For a scalar depth-two network, take both layer predictors equal to $-2$,
so $A_0=A_1=0$, and the endpoint loss $\ell(p):=(p-1)^2/2$.
Both backpropagated gradients vanish although the endpoint is not optimal.
These predictors admit positive-definite layer covariances
$\bigl(\begin{smallmatrix}\sqrt5&-2\\-2&\sqrt5\end{smallmatrix}\bigr)$,
whose signed eigenvalues are $+1,-1$. Extending the earlier discretization and sampling arguments to this setting
requires control of the residual products along the corresponding training
trajectory; the two-layer GD theorem does not apply verbatim.

In continuous depth, $\det P(t)=
\exp(\int_0^1\tr\Svu_s(t)\,\dd s)>0$.
The uniform path bound also keeps $\sigma_{\min}(P(t))\ge e^{-B}$,
so any convergent endpoint remains in the orientation-preserving
invertible component. If every global minimizer has negative determinant,
these path bounds are incompatible with convergence to a global minimizer.
\end{remark}

\subsection{Heavy-ball Momentum}\label{app:momentum-proof}
We first expand the closed phase equation, then establish global
existence, position-deformation estimates, and the energy bootstrap.
For positions $x:=(u,v)$ and velocities $y:=(p,q)$, let
$\nu_t\in\cP_2(\R^{2m})$ be their joint distribution. Its covariance is
$\Phase_t:=\int\binom xy\binom xy^{\!\top}\,\dd\nu_t(x,y)$,
with mixed block $\Xi_t:=\int xy^\top\,\dd\nu_t(x,y)$.
Finite-width factors correspond to
$\nu_t:=n^{-1}\sum_i\delta_{(x_i(t),y_i(t))}$.
The characteristics satisfy
\begin{equation}
\label{eq:single-particle-momentum-main}
\dot x=y,\qquad
\dot y=-G(\Svu)x-\gamma y.
\end{equation}
Taking its second moments gives the block system
\begin{equation}
\label{eq:phase-block-ode}
\begin{aligned}
\dSpos_t&=\Xi_t+\Xi_t^\top,\qquad
\dot\Xi_t=\Svel_t-\Spos_tG_t-\gamma\Xi_t,\\
\dSvel_t&=-G_t\Xi_t-\Xi_t^\top G_t-2\gamma\Svel_t.
\end{aligned}
\end{equation}
In particular, the predictor velocity and acceleration are
\begin{equation}\label{eq:predictor-speed}
 \dSvu=\Sigma^{qu}+\Sigma^{vp},\qquad
 \ddSvu+\gamma\dSvu=2\Sqp-E\Suu-\Svv E.
\end{equation}

\begin{proof}[Proof of Proposition~\ref{prop:phase-laws}]
Since \(G_t\) is symmetric,
\(\Mob_t^\top \Jcan+\Jcan\Mob_t=-\gamma \Jcan\).
Differentiate the two-characteristic pairing and the fundamental-matrix
identity. Invertibility also gives
\(\Psi_t\Jcan(\Psi_t)^\top=e^{-\gamma t}\Jcan\), whence
\(\Jcan\Phase_t=e^{-\gamma t}
(\Psi_t)^{-\top}(\Jcan\Phase_0)(\Psi_t)^\top\).
\end{proof}

\begin{proposition}[Global full phase flow]\label{prop:global-phase}
If $\ell$ is $C^1$, has locally Lipschitz gradient, and is bounded below,
every $\Phase_0\succeq0$ has a global positive semidefinite solution of
\eqref{eq:phase-block-ode}.  This assertion alone gives neither uniform position ellipticity
nor convergence to a minimizer.
\end{proposition}

\begin{proof}[Proof of Proposition~\ref{prop:global-phase}]
Local well-posedness follows from local Lipschitz continuity. The linear
phase characteristic lift preserves positivity. The energy identity gives
$\tr\Svel_t\le2\cE_0$. Positivity and Cauchy--Schwarz yield
\[
 \sqrt{\tr\Spos_t}
 \le\sqrt{\tr\Spos_0}+t\sqrt{2\cE_0},\qquad
 \norm{\Xi_t}_{\F}^2\le\tr\Spos_t\tr\Svel_t.
\]
The first inequality also follows directly from the $L^2$ triangle
inequality along characteristics, including times when the trace is zero.
Every block is bounded on each finite interval; continuation proves
global existence.
\end{proof}
For the matrix calculations in the momentum sections, abbreviate
$A:=\Spos$ and $B:=\Svel$. Whenever $A\succ0$, define
\begin{equation}\label{eq:regression}
 D:=\Xi^\top A^{-1},\qquad Q:=B-DAD^\top\succeq0.
\end{equation}
Here $D$ is a regression matrix, not assumed symmetric, and $Q$ is its
positive semidefinite residual covariance (the Schur complement).
They encode all phase blocks by $\Xi=AD^\top$ and $B=DAD^\top+Q$.
\begin{proposition}[General covariance coordinates]\label{prop:regression}
On the chart $A\succ0$, \eqref{eq:phase-block-ode} is equivalent to
\begin{equation}\label{eq:regression-ode}
 \begin{aligned}
 \dot A&=DA+AD^\top,\\
 \dot D&=-G-\gamma D-D^2+QA^{-1},\\
 \dot Q&=-2\gamma Q-DQ-QD^\top.
 \end{aligned}
\end{equation}
Moreover, positive semidefiniteness controls the position deformation:
\begin{equation}\label{eq:regression-bound-main}
 \norm D_{\op}^2\le\frac{2\mathcal{K}}{\lambda_{\min}(A)}.
\end{equation}
\end{proposition}


\begin{proof}[Proof of Proposition~\ref{prop:regression}]
Differentiate $D=\Xi^\top A^{-1}$ and $Q=B-DAD^\top$, using
\eqref{eq:phase-block-ode}. This gives the three displayed identities.
Since $DAD^\top\preceq B$, its trace is at most $2\mathcal{K}$, whereas
$\lambda_{\min}(A)\norm D_{\op}^2\le\tr(DAD^\top)$.
\end{proof}
\begin{lemma}[Conditional phase-energy decay]\label{lem:conditional}
Under Assumption~\ref{ass:loss}, let a positive semidefinite solution of
\eqref{eq:phase-block-ode} satisfy $a\Id_m\preceq\Spos_t\preceq b\Id_m$ for
$0\le t\le T$, where $0<a\le b<\infty$. Put
$\mathfrak K:=4a\kappa$, $\mathfrak L:=4bL$, and
$\rho:=\mathfrak L+2\sqrt{2L\cE_0}$. Then
\begin{equation}\label{eq:conditional}
 \cE_t\le3\cE_0 e^{-\lambda t},\qquad
 \lambda:=\frac{a\kappa\gamma}{\rho+2\gamma^2},\qquad 0\le t\le T.
\end{equation}
\end{lemma}
\begin{proof}
Set $\cN:=\tr(G\Spos G)$, the mean squared particle force.
The hidden principal blocks are between $a\Id$ and $b\Id$, hence
\[
 2a\norm E_{\F}^2\le\cN\le2b\norm E_{\F}^2.
\]
Combining PL with the smooth-gradient upper bound in
Lemma~\ref{lem:pl}, and applying Cauchy--Schwarz to the phase law, gives
\begin{equation}\label{eq:conditional-force-bounds}
 \mathfrak K\Delta\le \cN\le\mathfrak L\Delta,\qquad
 |\dot\Delta|=|\tr(G\Xi)|\le\sqrt{2\mathcal{K} \cN}.
\end{equation}
All second derivatives below are understood almost everywhere. Lipschitz
continuity of $\nabla\ell$ implies that $E_t$ is locally absolutely
continuous and $\norm{\dot E_t}_{\F}\le L\norm{\dSvu_t}_{\F}$.
Differentiating \eqref{eq:predictor-speed} gives
\[
 \ddot\Delta=H-\cN-\gamma\dot\Delta,\qquad
 H:=\ip{\dot E}{\dSvu}+2\ip E{\Sqp}\le\rho\mathcal{K}.
\]
For the last bound, positive phase blocks give
$\norm{\Sigma^{qu}}_{\F}^2\le b\tr\Sigma^q$ and
$\norm{\Sigma^{vp}}_{\F}^2\le b\tr\Sigma^p$, hence
$\norm{\dSvu}_{\F}^2\le4b\mathcal{K}$. Also
$\norm{\Sqp}_{\F}\le\sqrt{\tr\Sigma^q\tr\Sigma^p}\le\mathcal{K}$ and
$\norm E_{\F}\le\sqrt{2L\cE_0}$.

Choose $\epsilon:=\gamma/[2(\rho+2\gamma^2)]$ and
$\mathcal V:=\cE+\epsilon\dot\Delta$. Since
\[
 \epsilon\le\frac1{4\sqrt{2\rho}}
 \le\frac1{\sqrt{2\mathfrak L}},\qquad
 |\dot\Delta|\le\sqrt{\mathfrak L/2}\,\cE,
\]
one has $\cE/2\le\mathcal V\le3\cE/2$.
The inequality $\gamma\sqrt{2\mathcal{K} \cN}\le \cN/4+2\gamma^2\mathcal{K}$ yields
\[
 \dot{\mathcal V}\le-\tfrac34\epsilon \cN
       -\bigl(2\gamma-\epsilon(\rho+2\gamma^2)\bigr)\mathcal{K}
 \le-\tfrac34\epsilon\mathfrak K\Delta-\gamma\mathcal{K}.
\]
Now $\tfrac34\epsilon\mathfrak K\le\gamma$, since $a\kappa\le bL$
and $\rho\ge4bL$. It follows that
$\dot{\mathcal V}\le-\tfrac23(\tfrac34\epsilon\mathfrak K)\mathcal V
=-\lambda\mathcal V$.
Integrating the a.e. inequality proves \eqref{eq:conditional}.

\end{proof}

Under Theorem~\ref{thm:momentum}, the position-covariance tail is
\begin{equation}\label{eq:momentum-tail}
 \norm{\Spos_t-\Spos_\infty}_{\F}
 \le\frac{4\sqrt{6b_s\cE_0}}{\lambda_s(\gamma)}
       e^{-\lambda_s(\gamma)t/2}.
\end{equation}
For the position factors $W$ or the position marginal of the characteristic
phase law, the normalized Frobenius or $W_2$ tail has prefactor
$2\sqrt{6\cE_0}/\lambda_s(\gamma)$ at the same exponent.

\begin{proof}[Proof of Theorem~\ref{thm:momentum}]
The full phase flow is global by Proposition~\ref{prop:global-phase}.
On the connected position-positive chart containing time zero, let
$h(t):=\int_0^t\norm{D_r}_{\op}\,\dd r$.
The equation $\dot A=DA+AD^\top$ yields
$A_t=Z_t A_0Z_t^\top$ with $\dot Z=DZ$, $Z_0=\Id_m$, and thus
\[
 a_0e^{-2h(t)}\Id_m\preceq A_t\preceq b_0e^{2h(t)}\Id_m.
\]
As long as $h\le s$, Lemma~\ref{lem:conditional} and
Proposition~\ref{prop:regression} imply
\[
 h(t)\le\int_0^t\sqrt{2\mathcal{K}_r/a_s}\,\dd r
 \le\frac{2}{\lambda_s}\sqrt{\frac{6\cE_0}{a_s}}
 =\mathfrak S_s(\gamma)<s.
\]
There is no first time at which $h=s$. Nor can the chart terminate earlier:
on each such interval $A\succeq a_s\Id$ and the globally defined full
phase blocks stay bounded on finite horizons. This proves the global tube
and energy estimate. Furthermore the positive-block factorization gives
$\norm\Xi_{\F}^2\le b_s\tr B\le6b_s\cE_0 e^{-\lambda_st}$.
Integrate $\dot A=\Xi+\Xi^\top$ to get \eqref{eq:momentum-tail}.
For particles or a characteristic phase law, the remaining $L^2$ position displacement is
at most $2\sqrt{6\cE_0}\,e^{-\lambda_st/2}/\lambda_s$.
The limit is therefore attained, with zero loss gap and full phase
covariance $\diag(A_\infty,0)$.
\end{proof}

If $\cE_0=0$, the nonnegative kinetic and potential energies both vanish,
the force vanishes, and the momentum trajectory is stationary. Positive
definiteness of the \emph{phase} covariance is never required; in particular
zero-speed initialization is covered.

\subsection{An exact decomposition of hidden kinetic energy}
Prescribing the position covariance and its derivative does not determine
the kinetic energy: velocity components that leave the covariance
unchanged can still contribute to it.

\begin{proposition}[Minimum-energy lift of a covariance velocity]\label{prop:kinetic}
Fix $A\succ0$ and a symmetric matrix $H$ as the desired value of $\dot A$.
Let $C=C^\top$ be the unique solution of $CA+AC=H$.
Every positive semidefinite phase covariance with this position moment and
derivative has regression coordinates
\[
 D=C+\Omega,\qquad \Omega A+A\Omega^\top=0,\qquad Q\succeq0,
\]
and its kinetic energy decomposes as
\begin{equation}\label{eq:kinetic-decomposition}
 \mathcal{K}=\tfrac12\tr(CAC)
    +\tfrac12\tr(\Omega A\Omega^\top)+\tfrac12\tr Q.
\end{equation}
Thus the graph lift $D=C$, $Q=0$ has minimum kinetic energy.
An additional positive semidefinite Schur complement contributes exactly
$\tr Q/2$.
\end{proposition}
\begin{proof}
In an eigenbasis of $A$, the Sylvester equation divides $H_{ij}$ by the
strictly positive sum of eigenvalues, proving existence and uniqueness.
Equation $DA+AD^\top=H$ is equivalent to the stated condition on $\Omega$.
Since $B=DAD^\top+Q$, expand its trace. The two cross terms vanish in sum:
$\Omega A$ is skew-symmetric and $C$ is symmetric. Each term in
\eqref{eq:kinetic-decomposition} is nonnegative.
\end{proof}

\begin{example}[Identical position covariance and covariance derivative, arbitrarily large energy]
Let $A=\sigma^2\Id_m$, $C=0$, choose any skew-symmetric $\Omega$, and set
$D=\Omega$, $Q=0$. The initial predictor is zero because
$A=\sigma^2\Id_m$, and $\dot A=0$ makes its derivative zero as well,
independently of $\Omega$. Nevertheless
$\mathcal{K}=\sigma^2\norm\Omega_{\F}^2/2$ can be arbitrarily large. The phase
covariance is positive semidefinite; nonzero $\Omega$ gives a nonsymmetric regression matrix.
For fixed $a_0,b_0,L,\kappa$, the best feasibility ratio
\eqref{eq:budget} tends to zero as $\cE_0\to\infty$. Thus the initial
position covariance and predictor speed do not determine whether the
energy-damping certificate holds.
\end{example}

%%% Local Variables:
%%% mode: LaTeX
%%% TeX-master: "../main"
%%% End:

% !TEX root = ../main.tex
% Section: Appendix: Tightened Conditional Momentum Rates

\section{Tightened Conditional Momentum Rates}
\label{app:tightened-conditional-momentum}\label{app:momentum-details}

This appendix sharpens Lemma~\ref{lem:conditional}. It gives the exact
optimizer within the same Lyapunov estimates, the simpler parameter
choice used in the numerical figures, and the corresponding global
bootstrap with its energy prefactor.

\begin{proposition}[Tightened conditional energy decay]
\label{prop:tightened-conditional-momentum-pl}\label{prop:tightened}
Under the loss and position bounds of Lemma~\ref{lem:conditional} on
$[0,T]$, fix $\eta\in(0,1)$ and $\varepsilon>0$. With its constants
$\mathfrak K,\mathfrak L,\rho$, set
\begin{align}
 q&:=\varepsilon\sqrt{\mathfrak L/2},\label{eq:tightened-momentum-q-definition}\\
 r_{\varepsilon,\eta}
 &:=\min\{\varepsilon(1-\eta)\mathfrak K,\,
          2\gamma-\varepsilon(\rho+\gamma^2/(2\eta))\},\label{eq:tightened-momentum-dissipation-definition}\\
 \lambda_{\varepsilon,\eta}
 &:=r_{\varepsilon,\eta}/(1+q).\label{eq:tightened-momentum-rate-definition}
\end{align}
If $q<1$ and $r_{\varepsilon,\eta}>0$, the corrected energy
$\mathcal V_t:=\cE_t+\varepsilon\dot\Delta_t$ satisfies, for $0\le t\le T$,
\begin{align}
 (1-q)\cE_t&\le\mathcal V_t\le(1+q)\cE_t,\label{eq:tightened-momentum-lyapunov-equivalence}\\
 \mathcal V_t&\le\mathcal V_0e^{-\lambda_{\varepsilon,\eta}t},\label{eq:tightened-conditional-momentum-rate}\\
 \ell(\Svu_t)-\ell_\star
 &\le\cE_t\le\frac{\mathcal V_0}{1-q}e^{-\lambda_{\varepsilon,\eta}t}
 \le\frac{1+q}{1-q}\cE_0e^{-\lambda_{\varepsilon,\eta}t}.\label{eq:tightened-conditional-momentum-loss-rate}
\end{align}
If $\dot\Delta_0=0$, the first energy prefactor is $\cE_0/(1-q)$.
\end{proposition}

\begin{proof}
The force bounds \eqref{eq:conditional-force-bounds} give
$|\dot\Delta|\le\sqrt{2\mathfrak L\Delta\mathcal{K}}
\le\sqrt{\mathfrak L/2}\,\cE$, proving the equivalence.
For the derivative, use $\dot{\cE}=-2\gamma\mathcal{K}$,
$\ddot\Delta\le\rho\mathcal{K}-\mathcal N-\gamma\dot\Delta$, and Young's inequality:
\begin{equation}\label{eq:tightened-momentum-parameterized-young}
 \gamma|\dot\Delta|\le\gamma\sqrt{2\mathcal{K}\mathcal N}
 \le\eta\mathcal N+\frac{\gamma^2}{2\eta}\mathcal{K}.
\end{equation}
Since $\mathcal N\ge\mathfrak K\Delta$, this yields
\[
 \dot{\mathcal V}
 \le-\varepsilon(1-\eta)\mathcal N
     -[2\gamma-\varepsilon(\rho+\gamma^2/(2\eta))]\mathcal{K}
 \le-r_{\varepsilon,\eta}\cE
 \le-\lambda_{\varepsilon,\eta}\mathcal V.
\]
Integrate on $[0,T]$ and use the equivalence once more.
\end{proof}

\begin{proposition}[Optimal rate within the Lyapunov bounds]
\label{prop:optimal-tightened-conditional-momentum-rate}
Under Proposition~\ref{prop:tightened}, fix $\eta\in(0,1)$ and write
\begin{equation}\label{eq:tightened-momentum-A-B-definition}
 A_\eta:=(1-\eta)\mathfrak K,\qquad
 B_\eta:=\rho+\gamma^2/(2\eta).
\end{equation}
The unique maximizing correction at fixed $\eta$ is admissible and gives
\begin{align}
 \varepsilon_\eta^\star&=\frac{2\gamma}{A_\eta+B_\eta},\label{eq:tightened-momentum-optimal-epsilon}\\
 q_\eta^\star&:=\varepsilon_\eta^\star\sqrt{\mathfrak L/2}<1,\label{eq:tightened-momentum-optimal-q}\\
 \lambda_\eta&:=\frac{2\gamma A_\eta}
 {A_\eta+B_\eta+\gamma\sqrt{2\mathfrak L}}.\label{eq:tightened-momentum-fixed-eta-rate}
\end{align}
The unique maximizer over $\eta$ and its rate are
\begin{align}
 \eta_\star&=\frac{\gamma}
 {\gamma+\sqrt{\gamma^2+2\rho+2\gamma\sqrt{2\mathfrak L}}}
 \in(0,1/2),\label{eq:tightened-momentum-optimal-eta}\\
 \lambda_\star&:=\max_{0<\eta<1}\lambda_\eta
 =\frac{2\gamma\mathfrak K}
 {\mathfrak K+\frac12
 (\gamma+\sqrt{\gamma^2+2\rho+2\gamma\sqrt{2\mathfrak L}})^2}.\label{eq:tightened-momentum-optimal-rate}
\end{align}
\end{proposition}

\begin{proof}
At fixed $\eta$, the rate is the minimum of
$\varepsilon A_\eta/(1+\varepsilon\sqrt{\mathfrak L/2})$ and
$(2\gamma-\varepsilon B_\eta)/(1+\varepsilon\sqrt{\mathfrak L/2})$.
The first is strictly increasing and the second strictly decreasing,
so their intersection gives the unique optimum.
It is admissible because
$B_\eta>\rho+\gamma^2/2\ge\gamma\sqrt{2\rho}
\ge\gamma\sqrt{2\mathfrak L}$, hence $q_\eta^\star<1$;
both dissipation terms there equal $\varepsilon_\eta^\star A_\eta>0$.

For the remaining optimization, set
$\widetilde\rho:=\rho+\gamma\sqrt{2\mathfrak L}$ and divide numerator
and denominator of \eqref{eq:tightened-momentum-fixed-eta-rate} by $1-\eta$.
Maximizing the rate is equivalent to minimizing
$(\widetilde\rho+\gamma^2/(2\eta))/(1-\eta)$.
Its derivative has the sign of
$\widetilde\rho\eta^2+\gamma^2\eta-\gamma^2/2$, which increases strictly
and vanishes only at \eqref{eq:tightened-momentum-optimal-eta}.
The quotient there is
$\frac12(\gamma+\sqrt{\gamma^2+2\widetilde\rho})^2$, proving the formula.
\end{proof}

The numerical figures use the simpler admissible choice
\begin{equation}\label{eq:tightened-momentum-eta-approx}
 \eta_{\mathrm{approx}}
 :=\frac{\gamma}{\gamma+\sqrt{\gamma^2+2\rho}}\in(0,1/2).
\end{equation}
It minimizes $B_\eta/A_\eta$: differentiation gives
$\rho\eta^2+\gamma^2\eta-\gamma^2/2=0$.
Thus it optimizes the dissipation rate before accounting for the
equivalence factor $1+q_\eta^\star$. That factor is retained in the
following rigorous bound.

\begin{corollary}[Explicit rate used in the figures]
\label{cor:explicit-tightened-conditional-momentum-rate}\label{cor:tightened-choices}
Under Proposition~\ref{prop:tightened}, use
$\eta_{\mathrm{approx}}$ from \eqref{eq:tightened-momentum-eta-approx} and set
\begin{align}
 A_{\mathrm{approx}}&:=(1-\eta_{\mathrm{approx}})\mathfrak K,&
 B_{\mathrm{approx}}&:=\rho+\frac{\gamma^2}{2\eta_{\mathrm{approx}}},\label{eq:tightened-momentum-explicit-A-B}\\
 \varepsilon_{\mathrm{approx}}&:=\frac{2\gamma}{A_{\mathrm{approx}}+B_{\mathrm{approx}}},&
 q_{\mathrm{approx}}&:=\varepsilon_{\mathrm{approx}}\sqrt{\mathfrak L/2}.\label{eq:tightened-momentum-explicit-epsilon-q}
\end{align}
These parameters are admissible. With
$\mathcal V_0=\cE_0+\varepsilon_{\mathrm{approx}}\dot\Delta_0$,
\begin{align}
 \lambda_{\mathrm{approx}}
 &:=\frac{\varepsilon_{\mathrm{approx}}A_{\mathrm{approx}}}
         {1+q_{\mathrm{approx}}},\label{eq:tightened-momentum-explicit-rate}\\
 \ell(\Svu_t)-\ell_\star
 &\le\cE_t\le\frac{\mathcal V_0}{1-q_{\mathrm{approx}}}
       e^{-\lambda_{\mathrm{approx}}t},\qquad 0\le t\le T.\label{eq:tightened-momentum-explicit-loss-rate}
\end{align}
If $\dot\Delta_0=0$, the prefactor is $\cE_0/(1-q_{\mathrm{approx}})$.
\end{corollary}

\begin{proof}
The fixed-$\eta$ optimum in
Proposition~\ref{prop:optimal-tightened-conditional-momentum-rate}
is admissible for every $\eta\in(0,1)$.
Apply Proposition~\ref{prop:tightened} with $\eta=\eta_{\mathrm{approx}}$.
\end{proof}

The approximate choice concerns only optimization: its bound is exact,
and $\lambda_{\mathrm{approx}}\le\lambda_\star$. Neither exponent is
claimed to be the optimal rate of the underlying training dynamics.

\begin{corollary}[Global bootstrap with the correct prefactor]
\label{cor:tightened-bootstrap}
Under the loss and initialization hypotheses of Theorem~\ref{thm:momentum},
choose a tube $a_s:=a_0e^{-2s}$, $b_s:=b_0e^{2s}$ and an admissible
$(\varepsilon,\eta)$ for its constants. Put
\[
 \mathfrak C_0:=\frac{\cE_0+\varepsilon\dot\Delta_0}{(1-q)\cE_0}.
\]
If $\cE_0>0$ and
\begin{equation}\label{eq:tightened-bootstrap}
 \frac2{\lambda_{\varepsilon,\eta}}
 \sqrt{\frac{2\mathfrak C_0\cE_0}{a_s}}<s,
\end{equation}
then the tube and \eqref{eq:tightened-conditional-momentum-loss-rate} hold globally.
A sufficient condition using only initial position bounds and total
energy replaces $\mathfrak C_0$ by $(1+q)/(1-q)$.
\end{corollary}
\begin{proof}
Within the initial connected positive-position chart,
$\int_0^t\norm D_{\op}\le\int_0^t\sqrt{2\cE_r/a_s}\,\dd r$.
Proposition~\ref{prop:tightened} bounds this by the left side of
\eqref{eq:tightened-bootstrap}. The strict first-exit proof of
Theorem~\ref{thm:momentum} then applies unchanged. Initial positivity and
the Lyapunov equivalence imply $1\le \mathfrak C_0\le(1+q)/(1-q)$.
\end{proof}

\paragraph{The role of the prefactor.}
A faster conditional exponent must be accompanied by its own prefactor
when testing the global deformation budget.
The sharper state-dependent prefactor uses the initial loss derivative
$\dot\Delta_0=\langle E_0,\Sigma^{qu}_0+\Sigma^{vp}_0\rangle$
as well as total energy. In particular, zero predictor speed gives
$\mathfrak C_0=1/(1-q)$. The parameter-only prefactor uses no extra initial derivative.
Maximizing the conditional exponent alone need not minimize the global
budget: the same choice of parameters also changes $q$ and hence the
prefactor. The closed-form optimizer should therefore be tested in
\eqref{eq:tightened-bootstrap}, not substituted there with an unchanged
prefactor.

\section{Large-Position-Scale Certificate Asymptotics}
\label{app:large-position-scale-certificate}\label{app:large-position-scale}
\label{subsec:large-position-scale-certificate}

Increasing position variance can open a damping window without reducing
the initial loss gap. We quantify that tradeoff and the best rates
certified by the same bootstrap.
For the two-scale white initialization of
Corollary~\ref{cor:two-scale-white-window}, the initial mechanical excess
energy is
\[
\cE_0
=
\ell(0)-\ell_\star+\frac{m\sigma_{qp}^2}{2}.
\]
This quantity is fixed when \(\sigma_{qp}\) is fixed and
\(\sigma_{vu}\to+\infty\).  It is useful to name the bootstrap gap
\[
\mathfrak G_s(\gamma):=\mathfrak S_s(\gamma)-s,
\qquad
s_{\mathrm{cert}}(\gamma)
:=
\inf\{s>0:\mathfrak G_s(\gamma)<0\}.
\]
Whenever this set is nonempty, define the corresponding boundary rate by
\[
\lambda_{\mathrm{cert}}(\gamma)
:=
\lambda_{s_{\mathrm{cert}}(\gamma)}(\gamma),
\]
using continuity in \(s\).  The strict bootstrap condition need not hold at
the boundary itself: the theorem applies at interior budgets
\(s>s_{\mathrm{cert}}\) with \(\mathfrak G_s(\gamma)<0\), whose rates approach
this boundary value.  Thus \(\lambda_{\mathrm{cert}}\) is a limiting supremal
certificate, rather than a rate asserted at a possibly nonadmissible boundary
point.
Define also the largest damping for which the bootstrap window is nonempty:
\[
\gamma^{\mathrm{cert}}_{\max}(\sigma_{vu})
:=
\sup\{\gamma>0:\text{ there exists }s>0
\text{ such that }\mathfrak G_s(\gamma)<0\}.
\]

The assumption \(\cE_0>0\) below excludes only the trivial zero-excess
case, in which both the predictor loss gap and kinetic energy vanish and the
trajectory is stationary.

\paragraph{Explicit damping interval and feasibility.}
At a fixed deformation budget, admissibility reduces to a quadratic
inequality, making both damping boundaries explicit.
For \(\cE_0>0\) and fixed deformation budget \(s>0\), write
\begin{equation}\label{eq:window}
 c_s:=\frac{s\kappa a_s^{3/2}}{4\sqrt{3\cE_0/2}}.
\end{equation}
The theorem's strict certificate is exactly
\(2\gamma^2-c_s\gamma+\rho_s<0\), or
\begin{equation}\label{eq:window-roots}
 c_s^2>8\rho_s,\qquad
 \frac{c_s-\sqrt{c_s^2-8\rho_s}}4<\gamma<
 \frac{c_s+\sqrt{c_s^2-8\rho_s}}4.
\end{equation}
These are sufficient boundaries, not instability thresholds.


\paragraph{Optimizing feasibility.}
To make this interval nonempty as easily as possible, maximize
\begin{equation}\label{eq:budget}
 \frac{c_s^2}{8\rho_s}
 =\frac{s^2\kappa^2a_0^3e^{-6s}}
 {192\cE_0(4b_0Le^{2s}+2\sqrt{2L\cE_0})}.
\end{equation}
Its unique maximizer $s_{\rm feas}$ belongs to $(1/4,1/3)$ and satisfies
\begin{equation}\label{eq:best-budget}
 \frac1{s_{\rm feas}}
 =3+\frac{4b_0Le^{2s_{\rm feas}}}
 {4b_0Le^{2s_{\rm feas}}+2\sqrt{2L\cE_0}}.
\end{equation}
For a proof, write the denominator term as $Ae^{2s}+B$, with $A,B>0$.
The logarithmic derivative is $2/s-6-2Ae^{2s}/(Ae^{2s}+B)$, and its
derivative is $-2/s^2-4ABe^{2s}/(Ae^{2s}+B)^2<0$.
It changes sign between $1/4$ and $1/3$. The resulting scalar equation
is explicit and uniquely solvable; no trajectory has to be computed.
This budget optimizes \emph{nonemptiness}, not the upper damping endpoint
or the decay exponent. Those are different optimization questions.


\paragraph{Asymptotic certificate bounds.}
The large-position limit distinguishes fixed-damping behavior from the
largest damping that the certificate can accommodate.
\begin{proposition}[Large-position-scale certificate asymptotics]
\label{prop:large-position-scale-certificate}\label{prop:large-position}
For the two-scale white initialization above, fix \(\sigma_{qp}\ge0\), assume
\(\cE_0>0\), and
let \(\sigma_{vu}\to+\infty\).  For each fixed \(\gamma>0\) and fixed
\(s>0\),
\[
\mathfrak S_s(\gamma)
=
\frac{16Le^{5s}}{\kappa\gamma\,\sigma_{vu}}
\sqrt{\frac{3\cE_0}{2}}
+O(\sigma_{vu}^{-3}).
\]
The expansion is uniform for \(s\) in bounded sets.  Consequently,
\[
s_{\mathrm{cert}}(\gamma)
=
\frac{16L}{\kappa\gamma\,\sigma_{vu}}
\sqrt{\frac{3\cE_0}{2}}
+o(\sigma_{vu}^{-1}),
\]
and
\[
\lambda_{\mathrm{cert}}(\gamma)
=
\frac{\kappa\gamma}{4L}
\;-\;
\frac{16}{\sigma_{vu}}\sqrt{\frac{3\cE_0}{2}}
+o(\sigma_{vu}^{-1}).
\]
Moreover, the upper damping boundary satisfies
\[
\gamma^{\mathrm{cert}}_{\max}(\sigma_{vu})
=
\frac{\kappa\sigma_{vu}^3}
{24e\sqrt{3\cE_0/2}}
\bigl(1+o(1)\bigr).
\]
\end{proposition}

\begin{proof}
For fixed $\gamma>0$, the exact expression \eqref{eq:certificate}
can be expanded directly. Since
\[
\rho_s=4L\sigma_{vu}^2e^{2s}
+2\sqrt{2L\cE_0},
\]
we have
\[
\lambda_s(\gamma)
=
\frac{\kappa\gamma}{4Le^{4s}}
+O(\sigma_{vu}^{-2}).
\]
Substituting this into the definition of \(\mathfrak S_s(\gamma)\), with
\(a_s=\sigma_{vu}^2e^{-2s}\), gives the first expansion.
Write
\[
c_0:=\frac{16L}{\kappa\gamma}
\sqrt{\frac{3\cE_0}{2}}.
\]
The expansion is uniform near \(s=0\).  Thus every
\(s=c/\sigma_{vu}\), with \(c>0\) fixed, satisfies
\[
\mathfrak G_s(\gamma)
=\frac{c_0-c}{\sigma_{vu}}+o(\sigma_{vu}^{-1}),
\]
uniformly when \(c\) ranges over a compact set.  Hence budgets whose leading
constant is below \(c_0\) are inadmissible, whereas those whose leading
constant is strictly above \(c_0\) are admissible for all sufficiently large
\(\sigma_{vu}\).  This proves the displayed expansion for
\(s_{\mathrm{cert}}(\gamma)\).  Evaluating
\[
\lambda_s(\gamma)
=
\frac{\kappa\gamma}{4L}e^{-4s}
+O(\sigma_{vu}^{-2})
\]
at \(s=s_{\mathrm{cert}}(\gamma)\) gives the expansion for
\(\lambda_{\mathrm{cert}}(\gamma)\).

It remains to identify the high-damping boundary.  First observe the exact
upper bound
\[
\lambda_s(\gamma)
=\frac{a_s\kappa\gamma}{\rho_s+2\gamma^2}
\le\frac{\kappa\sigma_{vu}^2e^{-2s}}{2\gamma}.
\]
Consequently every admissible pair satisfies
\[
\mathfrak S_s(\gamma)
\ge
\frac{8\gamma e^{3s}}{\kappa\sigma_{vu}^3}
\sqrt{\frac{3\cE_0}{2}},
\]
and hence
\[
\gamma
<
\frac{\kappa\sigma_{vu}^3}{8\sqrt{3\cE_0/2}}\,s e^{-3s}
\le
\frac{\kappa\sigma_{vu}^3}{24e\sqrt{3\cE_0/2}}.
\]
This proves the required asymptotic upper bound uniformly over all budgets,
not only over bounded \(s\).

For the matching lower bound, take \(\gamma\) of order
\(\sigma_{vu}^3\) and keep \(s\) in a compact subset of \((0,+\infty)\).
Then $2\gamma^2$ dominates $\rho_s$ in \eqref{eq:certificate}, and
\[
\lambda_s(\gamma)
=
\frac{\kappa\sigma_{vu}^2e^{-2s}}{2\gamma}
\bigl(1+o(1)\bigr).
\]
Consequently,
\[
\mathfrak S_s(\gamma)
=
\frac{8\gamma e^{3s}}{\kappa\sigma_{vu}^3}
\sqrt{\frac{3\cE_0}{2}}
\bigl(1+o(1)\bigr),
\]
so the condition \(\mathfrak S_s(\gamma)<s\) is asymptotically equivalent to
\[
\gamma
<
\frac{\kappa\sigma_{vu}^3}{8\sqrt{3\cE_0/2}}\,
s e^{-3s}.
\]
The right-hand side is maximized at \(s=1/3\).  Choosing this budget and any
fixed multiplicative factor below the maximum gives an admissible pair for all
sufficiently large \(\sigma_{vu}\), proving the matching lower bound and the
claimed asymptotic formula.
\end{proof}

Proposition~\ref{prop:large-position-scale-certificate} has two complementary
messages.  At fixed damping, increasing the initial position scale makes the
certified ellipticity deformation vanish and the certified rate converges to
\(\kappa\gamma/(4L)\).  At fixed \(\sigma_{vu}\), however, this
particular bootstrap theorem cannot certify arbitrarily large damping: the
admissible region in \((s,\gamma)\) eventually disappears, with the upper
edge scaling like \(\sigma_{vu}^3\).  This is a limitation of the certificate,
not a claim that the underlying damped dynamics fail in the large-damping
regime.

\begin{figure}[t]
  \centering
  \def\mathdefault#1{#1}
  \subimport{figures/}{momentum_bootstrap_lambda_multi_sigma.pgf}
\caption{
Limiting bootstrap rate at the admissibility boundary as a function of $\gamma$ for several initial position
scales $\sigma_{vu}$.
Solid curves are obtained by grid-searching the bootstrap certificate
for the boundary value \(\lambda_{s_{\mathrm{cert}}}(\gamma)\); rates at
strictly admissible nearby budgets approach these curves.  The dotted black curve is
the large-\(\sigma_{vu}\) limiting rate \(\kappa\gamma/(4L)\).  The
plot uses \((d_u, d_v) = (3, 2)\), \(\kappa=L=1\), \(\sigma_{qp}=1\), and
\(\Delta_0=\ell(0)-\ell_\star=1\).
For scales whose upper damping boundary lies within the scanned range, markers show numerical and asymptotic approximations to the maximal admissible damping and its corresponding boundary rate.
}
\label{fig:momentum-bootstrap-lambda-multi-sigma}\label{fig:momentum-boundary}
\end{figure}

\FloatBarrier

%%% Local Variables:
%%% mode: LaTeX
%%% TeX-master: "../main"
%%% End:

% !TEX root = ../main.tex
% Section: Appendix: Counting Covariance-Level Mean-Field Laws

\section{Counting Covariance-Level Mean-Field Laws}
\label{app:law-counting}

This appendix records the Lie-algebra computation behind the covariance-level
completeness statement used in Section~\ref{sec:laws}.
It is a continuum, covariance-level analogue of the local/generic finite-width
completeness theorem of \citet{MarcotteGribonvalPeyre2023}.

Width-consistent laws and their Wasserstein interpretation are defined
in Appendix~\ref{app:mean-field}. We work here directly with functions
of the covariance and the vector fields that must annihilate them.

\subsection{Split Bures vector fields}
\label{subsec:split-bures-fields}

To characterize the laws shared by the architecture, we allow every
smooth split loss:
\[
F(\Sigma):=\ell(\Svu),
\qquad
\Sigma=
\begin{pmatrix}
\Suu&\Suv\\
\Svu&\Svv
\end{pmatrix}.
\]
If \(E:=\nabla \ell(\Svu)\), set
\[
G_E:=
\begin{pmatrix}
0&E^\top\\
E&0
\end{pmatrix}\in\mathfrak p_J,
\qquad
\mathfrak p_J:=\{G\in\Ssym^m:GJ+JG=0\}.
\]
The corresponding Bures vector field is
\begin{equation}
\label{eq:appendix-split-vector-field}
X_\ell(\Sigma):=-(G_{\nabla \ell(\Svu)}\Sigma+\Sigma G_{\nabla \ell(\Svu)}).
\end{equation}
We denote the resulting space of split Bures vector fields by
\[
\mathcal X_{\rm split}
:=
\{X_\ell:\ell\in C^\infty(\R^{d_v\times d_u})\}.
\]
It is a real vector space.  Its constant-gradient subspace is obtained from
linear losses \(\ell_E(\Svu):=\ip{E}{\Svu}\):
\[
\mathcal X_{\rm lin}
:=
\{X_E:E\in\R^{d_v\times d_u}\},
\qquad
X_E(\Sigma):=-(G_E\Sigma+\Sigma G_E).
\]

\subsection{Generated Lie algebra}

Lie brackets reveal the directions generated by all split losses,
which determine how many independent common laws can remain.
Let
\[
\mathfrak o(J):=\{A\in\R^{m\times m}:A^\top J+JA=0\}
\]
be the Lie algebra of \(O(d_u,d_v)\).  For \(A\in\mathfrak o(J)\), define the
infinitesimal congruence field
\[
\xi_A(\Sigma):=-(A\Sigma+\Sigma A^\top).
\]
Its flow is \(\Sigma(t)=e^{-tA}\Sigma(0)e^{-tA^\top}\), so it is tangent to
the \(O(d_u,d_v)\)-congruence orbit of \(\Sigma\).  When \(A=G_E\in\mathfrak p_J\),
\(\xi_A=X_E\).

\begin{proposition}[Split fields generate \(\mathfrak o(d_u,d_v)\)]
\label{prop:appendix-lie-generated}
With the bracket convention \([X,Y]:=DY[X]-DX[Y]\), the map
\[
A\longmapsto \xi_A
\]
is a Lie-algebra homomorphism:
\[
[\xi_A,\xi_B]=\xi_{[A,B]}.
\]
Moreover,
\[
\operatorname{Lie}(\mathfrak p_J)=\mathfrak o(J).
\]
Consequently the Lie algebra generated by
\(\mathcal X_{\rm lin}\), the constant-gradient split Bures fields, is
\[
\{\xi_A:A\in\mathfrak o(J)\}.
\]
\end{proposition}

\begin{proof}
The homomorphism identity follows by differentiating the linear vector fields
\(\xi_A\).  It remains to compute the matrix Lie algebra generated by
\(\mathfrak p_J\).  For
\[
G_E=
\begin{pmatrix}
0&E^\top\\
E&0
\end{pmatrix},
\qquad
G_F=
\begin{pmatrix}
0&F^\top\\
F&0
\end{pmatrix},
\]
one has
\[
[G_E,G_F]
=
\begin{pmatrix}
E^\top F-F^\top E&0\\
0&EF^\top-FE^\top
\end{pmatrix}.
\]
These brackets span the block-diagonal skew-symmetric algebra
\[
\mathfrak k_J:=
\left\{
\begin{pmatrix}
A&0\\
0&D
\end{pmatrix}:A^\top=-A,\ D^\top=-D
\right\},
\]
by choosing \(E,F\) supported on one common row or one common column.  Since
\[
\mathfrak o(J)=\mathfrak k_J\oplus\mathfrak p_J,
\]
the generated Lie algebra is all of \(\mathfrak o(J)\).
\end{proof}

\paragraph{Fields versus their generated distribution.}
The full space \(\mathcal X_{\rm split}\) contains nonlinear fields and
is not identified with a finite-dimensional matrix Lie algebra. Each
of its fields is a smooth, state-dependent linear combination of a
basis of the \(X_E\)'s. Consequently all such fields and their brackets
remain tangent to the same congruence orbits, while the constant-gradient
fields already generate the whole orbit tangent space. It is this
distribution that determines the common conservation laws.

\subsection{Completeness of the spectral laws}

The generated directions coincide with the tangent directions of
covariance congruence orbits; their codimension gives the local law count.
For a constant-gradient split Bures field, \(B:=J\Sigma\) satisfies
\[
\dot B=[G_E,B].
\]
The trace powers $\mathcal H_k$ in \eqref{eq:bures-G} are
therefore conserved. The next corollary states that, locally and generically, these
are the only smooth covariance-level laws common to all split losses.

Let
\[
\Omega_{\rm simp}
:=
\{\Sigma\in\Spp^m:J\Sigma\text{ has simple spectrum}\}.
\]

\begin{corollary}[Marcotte-style covariance-level completeness]
\label{cor:appendix-covariance-completeness}
Let \(h\) be a smooth function on an open subset of \(\Omega_{\rm simp}\).
If \(h\) is conserved by every split Bures vector field \(X_\ell\) of
\eqref{eq:appendix-split-vector-field}, then \(h\) is locally a smooth
function of the eigenvalues of \(J\Sigma\), equivalently of
\[
\mathcal H_1(\Sigma),\ldots,\mathcal H_m(\Sigma).
\]
These \(m\) laws are functionally independent on \(\Omega_{\rm simp}\).
Thus the characteristic polynomial of \(J\Sigma\) gives a complete
local/generic list of covariance-level mean-field conservation laws for the
split two-layer linear architecture.
\end{corollary}

\begin{proof}
It is enough to use the linear split losses.  If \(X_Eh=0\) for every \(E\),
then \(h\) is also annihilated by all iterated Lie brackets of the \(X_E\)'s.
By Proposition~\ref{prop:appendix-lie-generated},
\[
\mathrm{d}h(\Sigma)[\xi_A(\Sigma)]=0
\qquad\text{for all }A\in\mathfrak o(J).
\]
Thus \(h\) is constant on the connected pieces of the
\(O(d_u,d_v)\)-congruence orbits.

The orbit invariants are exactly the eigenvalues of \(J\Sigma\) on
\(\Omega_{\rm simp}\).  Indeed,
\[
J(Q\Sigma Q^\top)=Q^{-\top}(J\Sigma)Q^\top
\qquad(Q^\top JQ=J),
\]
so the spectrum is preserved.  Conversely, for \(\Sigma\succ0\), the matrix
\(J\Sigma\) is similar to the symmetric matrix
\(\Sigma^{1/2}J\Sigma^{1/2}\), hence has real nonzero eigenvalues with
\(d_u\) positive and \(d_v\) negative signs.  On the simple-spectrum stratum
its \(J\)-orthogonal eigenbasis can be locally normalized to a
\(J\)-orthonormal basis, so equal locally ordered eigenvalues imply membership
in the same local \(O(d_u,d_v)\)-orbit.  Therefore any smooth orbit-constant
function factors locally through the eigenvalues.  Newton identities replace
the eigenvalues by \(\mathcal H_1,\ldots,\mathcal H_m\).
More explicitly, for the locally ordered eigenvalues \(\lambda_j\),
\[
 \frac{\partial\mathcal H_k}{\partial\lambda_j}
 =k\lambda_j^{k-1},\qquad 1\le k,j\le m.
\]
This scaled Vandermonde matrix is invertible at distinct eigenvalues.
The eigenvalues themselves vary independently by varying the diagonal
covariance in a fixed normalized eigenbasis. Hence the spectral map has
rank \(m\), and its local level sets have dimension
\(m(m+1)/2-m=m(m-1)/2\), exactly the orbit dimension.
\end{proof}

For empirical covariances \(\Sigma=WW^\top/n\) with
\(W=\binom{U}{V}\), the finite-width Gram imbalance is
\[
\Gamma:=V^\top V-U^\top U=-W^\top JW.
\]
The nonzero spectra of \(J\Sigma=JWW^\top/n\) and
\(-\Gamma/n=W^\top JW/n\) coincide.  Corollary~\ref{cor:appendix-covariance-completeness}
is therefore the covariance quotient of the Marcotte--Gribonval--Peyr\'e
completeness theorem: the full labeled finite-width conservation data include the matrix
\(\Gamma\), while the covariance-level law keeps only its spectrum.

\subsection{A ReLU contrast}

Homogeneity preserves a scalar balance beyond linear networks, but it
does not give the same covariance closure or spectral family.
For a two-layer ReLU or more generally one-homogeneous network, the characteristic
equations associated with a potential \(v^\top\psi(u)\), with \(\psi\) one-homogeneous
and possibly time-dependent,
have the form
\[
\dot u=-D\psi(u)^\top v,\qquad \dot v=-\psi(u).
\]
Euler's identity \(D\psi(u)u=\psi(u)\) gives
\[
\frac{\mathrm{d}}{\mathrm{d}t}\bigl(\norm u^2-\norm v^2\bigr)=0.
\]
For ReLU this calculation is interpreted almost everywhere along the
characteristics, with an activation-boundary convention satisfying Euler's
identity (for instance derivative zero at a zero preactivation).
Thus the full mean-field flow preserves the pushforward distribution of
\(\norm u^2-\norm v^2\), and its covariance-level shadow is the total imbalance
\[
\int\bigl(\norm u^2-\norm v^2\bigr)\,\mathrm{d}\mu(u,v).
\]
Unlike the linear split case above, this ReLU statement is not obtained from a
closed finite-dimensional Bures covariance equation; it is included only to
highlight how much smaller the covariance-level conservation-law problem can be
than the full particle-level one. No completeness claim for ReLU laws is
made here.

%%% Local Variables:
%%% mode: LaTeX
%%% TeX-master: "../main"
%%% End:

\clearpage
\section{Intrinsic dynamics}\label{app:intrinsic}
Special conserved spectral values express some covariance blocks in terms
of fewer variables. This appendix derives these exact reductions.

\subsection{A two-point signed spectrum: hyperbolic isotropy}
A special two-point spectrum determines the hidden covariance blocks
from the predictor and therefore closes a smaller equation.
\begin{definition}[Hyperbolic isotropy]\label{def:hyperbolic-isotropy}
For \(\alpha>0\), the hyperbolically isotropic (HI) leaf is
\begin{equation}\label{eq:intrinsic-hi}
 \mathcal M_\alpha:=\{\Sigma\succeq0:\Sigma J\Sigma=\alpha^2J\}.
\end{equation}
\end{definition}
The identity forces $\Sigma\succ0$ and is equivalent to
$(J\Sigma)^2=\alpha^2\Id_m$. The symmetric representative of $J\Sigma$
then has $d_u$ eigenvalues $+\alpha$ and $d_v$ eigenvalues $-\alpha$;
equivalently, $\mathcal H_k=\alpha^k(d_u+(-1)^kd_v)$ for $k=1,\ldots,m$.
Spectral conservation proves invariance of $\mathcal M_\alpha$ despite
the repeated eigenvalues. Directly, the derivative of $\Sigma J\Sigma$
is $-G(\Sigma J\Sigma)-(\Sigma J\Sigma)G$, with constant solution $\alpha^2J$.

\begin{proposition}[Predictor-only intrinsic equation]\label{prop:intrinsic-hi}
Every predictor \(\Svu\in\R^{d_v\times d_u}\) has a unique covariance
on \(\mathcal M_\alpha\), given by
\begin{equation}\label{eq:intrinsic-graph}
 \Sigma=
 \begin{pmatrix}
 ((\Svu)^\top\Svu+\alpha^2\Id_{d_u})^{1/2}&(\Svu)^\top\\
 \Svu&(\Svu(\Svu)^\top+\alpha^2\Id_{d_v})^{1/2}
 \end{pmatrix}.
\end{equation}
On this leaf, the Bures dynamics close on the predictor:
\begin{equation}\label{eq:intrinsic-hyperbolic-dynamics}
 \dSvu=-E((\Svu)^\top\Svu+\alpha^2\Id)^{1/2}
       -(\Svu(\Svu)^\top+\alpha^2\Id)^{1/2}E,
 \qquad E:=\nabla\ell(\Svu).
\end{equation}
For \(d_u=d_v=1\), this is
\(\dSvu=-2\sqrt{(\Svu)^2+\alpha^2}\,\ell'(\Svu)\).
\end{proposition}
\begin{proof}
Expanding \eqref{eq:intrinsic-hi} gives
\[
 (\Suu)^2-(\Svu)^\top\Svu=\alpha^2\Id,\quad
 (\Svv)^2-\Svu(\Svu)^\top=\alpha^2\Id,\quad
 \Svu\Suu=\Svv\Svu.
\]
Positivity selects the positive square roots. Conversely, an SVD proves
the intertwining identity and gives eigenvalues \(\alpha\) and
\(\sqrt{\sigma_i^2+\alpha^2}\pm\sigma_i>0\), proving membership in
\(\mathcal M_\alpha\). Substituting the diagonal blocks into
\eqref{eq:block-flow-general} gives the intrinsic equation.
\end{proof}
The same eigenvalue calculation gives the exact extremal envelopes
\[
 \lambda_{\max}(\Sigma)=\sqrt{\norm{\Svu}_{\op}^2+\alpha^2}
                         +\norm{\Svu}_{\op},\qquad
 \lambda_{\min}(\Sigma)=\frac{\alpha^2}{\lambda_{\max}(\Sigma)}.
\]
In the general convergence theorems this special spectral level has
\(\aplus{\Sigma_0}=\aminus{\Sigma_0}=\alpha\). Thus its loss exponent is \(4\alpha\kappa\),
and the GD contraction factor is \(1-\sqrt2\alpha\kappa\tau\).

\paragraph{Finite-width form.}
The covariance constraint has an exact finite-width form that depends
only on the product $WW^\top$.
For \(\Sigma=WW^\top/n\), condition \eqref{eq:intrinsic-hi} is
\(WW^\top JWW^\top=n^2\alpha^2J\). In terms of \(U,V\), it reads
\[
 \begin{gathered}
 (UU^\top)^2-UV^\top VU^\top=n^2\alpha^2\Id_{d_u},\qquad
 (VV^\top)^2-VU^\top UV^\top=n^2\alpha^2\Id_{d_v},\\
 VU^\top UU^\top=VV^\top VU^\top.
 \end{gathered}
\]
Any covariance in \eqref{eq:intrinsic-graph} admits empirical factors at
every width \(n\ge m\), for example \(W=\sqrt n\,\Sigma^{1/2}Q_0\)
with \(Q_0Q_0^\top=\Id_m\).

\paragraph{Comparison with relaxed balance.}
\citet[Theorem 3.8]{MarcottePeyreGribonval2026} derive intrinsic dynamics
under a different condition,
\(\Gamma_{\rm rb}:=(V^\top V-U^\top U)/n=\beta\Id_n\),
on a full-column-rank stratum. It implies
\[
 (\Suu)^2+\beta\Suu=(\Svu)^\top\Svu,\qquad
 (\Svv)^2-\beta\Svv=\Svu(\Svu)^\top.
\]
For \(\beta\ne0\), zero singular directions have ambiguous branches.
Where both factors have full column rank, the range projectors
\(\Pi_u,\Pi_v\) of \((\Svu)^\top,\Svu\) select
\[
 \begin{aligned}
 \Suu&=\Pi_u[-\tfrac\beta2\Id+
             ((\Svu)^\top\Svu+\tfrac{\beta^2}{4}\Id)^{1/2}],\\
 \Svv&=\Pi_v[\tfrac\beta2\Id+
             (\Svu(\Svu)^\top+\tfrac{\beta^2}{4}\Id)^{1/2}].
 \end{aligned}
\]
If \(n=d_u=d_v\) and both factors are invertible, the projectors are
identities and the predictor equation coincides with
\eqref{eq:intrinsic-hyperbolic-dynamics} for \(\alpha=|\beta|/2\).
The leaves remain distinct: for \(\beta\ne0\), relaxed balance is equivalent to
\(\Sigma J\Sigma=-\beta\Sigma\), \(\rankop\Sigma=n\).
Indeed a left inverse of \(W/\sqrt n\) proves the converse.
It forces \(n\le d_v\) if \(\beta>0\), and \(n\le d_u\) if \(\beta<0\);
positive HI instead has rank \(m\).
For \(\beta=0\), positivity selects the unshifted square roots without a
rank assumption, giving the boundary equation \(\alpha=0\), not a
positive-definite initialization.
Without the rank qualifications, relaxed balance does not ensure
predictor closure: \(n=1,d_u=d_v=2,U=0,V=e_1\) or \(e_2\) gives the
same predictor and \(\beta=1\), but different predictor velocities for
\(\ell(S):=\norm{S+\Id_2}_{\F}^2/2\).

\begin{remark}[Conserved values do not generally close the predictor]
For \(d_u=d_v=2\), take \(\Svu=0\), \(\Svv=\diag(3,4)\), and
\(\Suu=\diag(1,2)\) or \(\diag(2,1)\).
Both states have signed spectrum \(\{1,2,-3,-4\}\), hence the same
spectral laws. For \(\ell(S):=\norm{S+e_1e_1^\top}_{\F}^2/2\),
their predictor velocities are \(-4e_1e_1^\top\) and
\(-5e_1e_1^\top\).
The special two-point spectrum gives an algebraic graph unavailable
at these general spectral values.
\end{remark}

\subsection{A single phase spectral scale: conformal symplectic isotropy}
A special conformal spectral identity eliminates the velocity covariance
within a positive-position chart of the phase flow.
\begin{definition}[Conformal symplectic isotropy]\label{def:csi}
A phase covariance is conformally symplectically isotropic (CSI) at
scale \(\theta\ge0\) if
\begin{equation}\label{eq:intrinsic-csi}
 \Phase \Jcan\Phase=\theta^2\Jcan.
\end{equation}
\end{definition}
For $\theta>0$, this forces $\Phase\succ0$ and is equivalent to
$\Phase^{1/2}\Jcan\Phase^{1/2}$ having $m$ copies of each eigenvalue
$+i\theta,-i\theta$.
The renormalized trace powers in Proposition~\ref{prop:phase-laws}
are zero at odd orders and $2m(-1)^k\theta_0^{2k}$ at order $2k$.
Conformal transport preserves \eqref{eq:intrinsic-csi} with
$\theta_t=e^{-\gamma t}\theta_0$, including the singular boundary
$\theta_0=0$, which is specified algebraically, without diagonalizability.

\begin{proposition}[Reduced phase dynamics]\label{prop:intrinsic-csi}
On the chart \(A:=\Spos\succ0\), condition \eqref{eq:intrinsic-csi}
is equivalent to
\[
 D=C=C^\top,\qquad Q=\theta^2A^{-1},
 \qquad \Xi=AC,\quad \Svel=CAC+\theta^2A^{-1}.
\]
The reduced equations are
\begin{equation}\label{eq:csi-reduced}
 \dot A=CA+AC,\qquad
 \dot C=-G-\gamma C-C^2+\theta_t^2A^{-2},
 \qquad\dot\theta_t=-\gamma\theta_t.
\end{equation}
They use only the two symmetric matrices \(A,C\), with known scale
\(\theta_t\), instead of all blocks of the phase covariance.
\end{proposition}
\begin{proof}
Substitute \(\Xi=AD^\top\), \(\Svel=DAD^\top+Q\).
The upper-left block of \eqref{eq:intrinsic-csi} is
\(A(D-D^\top)A=0\), forcing symmetry.
The upper-right block then becomes \(AQ=\theta^2\Id\);
the other blocks follow.
Substitution in Proposition~\ref{prop:regression} gives
\eqref{eq:csi-reduced}. This calculation includes \(\theta=0\).
\end{proof}
White position and velocity blocks with \(\Xi_0=0\) give
\(\theta_0=\sigma_{vu}\sigma_{qp}\). Zero velocity gives
\(C_0=0,\theta_0=0\), allowing any \(A_0\succ0\).
Neither choice is needed for convergence: the general theorem controls
\(D\) by energy even when it is nonsymmetric and \(Q\) is independent.

\begin{remark}[What the reduced state still contains]
Even on CSI, the two cross blocks \(\Svu,\Sqp\) do not generally form a
closed state. With \(\theta=1,C=0\), choose \(A=\Id_2\) or \(2\Id_2\)
and \(\Svel=A^{-1}\). Both have
\(\Svu=\Sqp=\dSvu=0\), yet
\(\ddSvu=-2E\) or \(-4E\) for the same nonzero scalar loss gradient.
Thus the full position block remains part of the intrinsic variables.
\end{remark}

\clearpage
% !TEX root = ../main.tex
\section{A nonquadratic example: smoothed Lasso}
\label{app:smooth-lasso}

To illustrate the smooth PL loss class beyond quadratic objectives, we repeat the three experiments of
Figure~\ref{fig:bures-convergence-triptych} with a smoothed entrywise $\ell_1$ penalty on the predictor:
\begin{equation}\label{eq:smooth-lasso-loss}
  \ell(S)=\frac12\norm{S-S_{\rm teach}}_{\F}^2
  +\lambda\sum_{i,j}\Bigl(\sqrt{S_{ij}^2+\delta^2}-\delta\Bigr),
  \qquad \lambda=0.1,\quad \delta=0.01.
\end{equation}
The Hessian is diagonal, with entries
$1+\lambda\delta^2/(S_{ij}^2+\delta^2)^{3/2}\in[1,1+\lambda/\delta]$.
Thus $\ell$ is globally $1$-strongly convex and satisfies Assumption~\ref{ass:loss}
with $\kappa=1$ and $L=11$.
The penalty acts on $S$, not separately on the factors, so the covariance closure and spectral conservation laws apply without modification.
Its unique minimizer $S_{\rm reg}$ is generally different from the teacher; we compute it entrywise from
\[
  (S_{\rm reg})_{ij}-(S_{\rm teach})_{ij}
  +\lambda\frac{(S_{\rm reg})_{ij}}{\sqrt{(S_{\rm reg})_{ij}^2+\delta^2}}=0.
\]
All loss gaps below are relative to $\ell_\star=\ell(S_{\rm reg})$.
Positive smoothing does not generally yield exact sparsity; the nonsmooth limit $\delta=0$ is outside Assumption~\ref{ass:loss}.

\begin{figure}[htbp]
  \centering
  \resizebox{\linewidth}{!}{%% Creator: Matplotlib, PGF backend
%%
%% To include the figure in your LaTeX document, write
%%   \input{<filename>.pgf}
%%
%% Make sure the required packages are loaded in your preamble
%%   \usepackage{pgf}
%%
%% Also ensure that all the required font packages are loaded; for instance,
%% the lmodern package is sometimes necessary when using math font.
%%   \usepackage{lmodern}
%%
%% Figures using additional raster images can only be included by \input if
%% they are in the same directory as the main LaTeX file. For loading figures
%% from other directories you can use the `import` package
%%   \usepackage{import}
%%
%% and then include the figures with
%%   \import{<path to file>}{<filename>.pgf}
%%
%% Matplotlib used the following preamble
%%   \def\mathdefault#1{#1}
%%   \everymath=\expandafter{\the\everymath\displaystyle}
%%   \IfFileExists{scrextend.sty}{
%%     \usepackage[fontsize=6.000000pt]{scrextend}
%%   }{
%%     \renewcommand{\normalsize}{\fontsize{6.000000}{7.200000}\selectfont}
%%     \normalsize
%%   }
%%   \usepackage{amsfonts}
%%   \usepackage{amssymb}
%%   \makeatletter\@ifpackageloaded{underscore}{}{\usepackage[strings]{underscore}}\makeatother
%%
\begingroup%
\makeatletter%
\begin{pgfpicture}%
\pgfpathrectangle{\pgfpointorigin}{\pgfqpoint{6.500000in}{1.687500in}}%
\pgfusepath{use as bounding box, clip}%
\begin{pgfscope}%
\pgfsetbuttcap%
\pgfsetmiterjoin%
\definecolor{currentfill}{rgb}{1.000000,1.000000,1.000000}%
\pgfsetfillcolor{currentfill}%
\pgfsetlinewidth{0.000000pt}%
\definecolor{currentstroke}{rgb}{1.000000,1.000000,1.000000}%
\pgfsetstrokecolor{currentstroke}%
\pgfsetdash{}{0pt}%
\pgfpathmoveto{\pgfqpoint{0.000000in}{0.000000in}}%
\pgfpathlineto{\pgfqpoint{6.500000in}{0.000000in}}%
\pgfpathlineto{\pgfqpoint{6.500000in}{1.687500in}}%
\pgfpathlineto{\pgfqpoint{0.000000in}{1.687500in}}%
\pgfpathlineto{\pgfqpoint{0.000000in}{0.000000in}}%
\pgfpathclose%
\pgfusepath{fill}%
\end{pgfscope}%
\begin{pgfscope}%
\pgfsetbuttcap%
\pgfsetmiterjoin%
\definecolor{currentfill}{rgb}{1.000000,1.000000,1.000000}%
\pgfsetfillcolor{currentfill}%
\pgfsetlinewidth{0.000000pt}%
\definecolor{currentstroke}{rgb}{0.000000,0.000000,0.000000}%
\pgfsetstrokecolor{currentstroke}%
\pgfsetstrokeopacity{0.000000}%
\pgfsetdash{}{0pt}%
\pgfpathmoveto{\pgfqpoint{0.487500in}{0.320625in}}%
\pgfpathlineto{\pgfqpoint{2.303086in}{0.320625in}}%
\pgfpathlineto{\pgfqpoint{2.303086in}{1.485000in}}%
\pgfpathlineto{\pgfqpoint{0.487500in}{1.485000in}}%
\pgfpathlineto{\pgfqpoint{0.487500in}{0.320625in}}%
\pgfpathclose%
\pgfusepath{fill}%
\end{pgfscope}%
\begin{pgfscope}%
\pgfpathrectangle{\pgfqpoint{0.487500in}{0.320625in}}{\pgfqpoint{1.815586in}{1.164375in}}%
\pgfusepath{clip}%
\pgfsetrectcap%
\pgfsetroundjoin%
\pgfsetlinewidth{0.803000pt}%
\definecolor{currentstroke}{rgb}{0.909804,0.909804,0.909804}%
\pgfsetstrokecolor{currentstroke}%
\pgfsetdash{}{0pt}%
\pgfpathmoveto{\pgfqpoint{0.487500in}{0.320625in}}%
\pgfpathlineto{\pgfqpoint{0.487500in}{1.485000in}}%
\pgfusepath{stroke}%
\end{pgfscope}%
\begin{pgfscope}%
\pgfsetbuttcap%
\pgfsetroundjoin%
\definecolor{currentfill}{rgb}{0.000000,0.000000,0.000000}%
\pgfsetfillcolor{currentfill}%
\pgfsetlinewidth{0.803000pt}%
\definecolor{currentstroke}{rgb}{0.000000,0.000000,0.000000}%
\pgfsetstrokecolor{currentstroke}%
\pgfsetdash{}{0pt}%
\pgfsys@defobject{currentmarker}{\pgfqpoint{0.000000in}{-0.048611in}}{\pgfqpoint{0.000000in}{0.000000in}}{%
\pgfpathmoveto{\pgfqpoint{0.000000in}{0.000000in}}%
\pgfpathlineto{\pgfqpoint{0.000000in}{-0.048611in}}%
\pgfusepath{stroke,fill}%
}%
\begin{pgfscope}%
\pgfsys@transformshift{0.487500in}{0.320625in}%
\pgfsys@useobject{currentmarker}{}%
\end{pgfscope}%
\end{pgfscope}%
\begin{pgfscope}%
\definecolor{textcolor}{rgb}{0.000000,0.000000,0.000000}%
\pgfsetstrokecolor{textcolor}%
\pgfsetfillcolor{textcolor}%
\pgftext[x=0.487500in,y=0.223403in,,top]{\color{textcolor}{\rmfamily\fontsize{6.000000}{7.200000}\selectfont\catcode`\^=\active\def^{\ifmmode\sp\else\^{}\fi}\catcode`\%=\active\def%{\%}0}}%
\end{pgfscope}%
\begin{pgfscope}%
\pgfpathrectangle{\pgfqpoint{0.487500in}{0.320625in}}{\pgfqpoint{1.815586in}{1.164375in}}%
\pgfusepath{clip}%
\pgfsetrectcap%
\pgfsetroundjoin%
\pgfsetlinewidth{0.803000pt}%
\definecolor{currentstroke}{rgb}{0.909804,0.909804,0.909804}%
\pgfsetstrokecolor{currentstroke}%
\pgfsetdash{}{0pt}%
\pgfpathmoveto{\pgfqpoint{0.941397in}{0.320625in}}%
\pgfpathlineto{\pgfqpoint{0.941397in}{1.485000in}}%
\pgfusepath{stroke}%
\end{pgfscope}%
\begin{pgfscope}%
\pgfsetbuttcap%
\pgfsetroundjoin%
\definecolor{currentfill}{rgb}{0.000000,0.000000,0.000000}%
\pgfsetfillcolor{currentfill}%
\pgfsetlinewidth{0.803000pt}%
\definecolor{currentstroke}{rgb}{0.000000,0.000000,0.000000}%
\pgfsetstrokecolor{currentstroke}%
\pgfsetdash{}{0pt}%
\pgfsys@defobject{currentmarker}{\pgfqpoint{0.000000in}{-0.048611in}}{\pgfqpoint{0.000000in}{0.000000in}}{%
\pgfpathmoveto{\pgfqpoint{0.000000in}{0.000000in}}%
\pgfpathlineto{\pgfqpoint{0.000000in}{-0.048611in}}%
\pgfusepath{stroke,fill}%
}%
\begin{pgfscope}%
\pgfsys@transformshift{0.941397in}{0.320625in}%
\pgfsys@useobject{currentmarker}{}%
\end{pgfscope}%
\end{pgfscope}%
\begin{pgfscope}%
\definecolor{textcolor}{rgb}{0.000000,0.000000,0.000000}%
\pgfsetstrokecolor{textcolor}%
\pgfsetfillcolor{textcolor}%
\pgftext[x=0.941397in,y=0.223403in,,top]{\color{textcolor}{\rmfamily\fontsize{6.000000}{7.200000}\selectfont\catcode`\^=\active\def^{\ifmmode\sp\else\^{}\fi}\catcode`\%=\active\def%{\%}10}}%
\end{pgfscope}%
\begin{pgfscope}%
\pgfpathrectangle{\pgfqpoint{0.487500in}{0.320625in}}{\pgfqpoint{1.815586in}{1.164375in}}%
\pgfusepath{clip}%
\pgfsetrectcap%
\pgfsetroundjoin%
\pgfsetlinewidth{0.803000pt}%
\definecolor{currentstroke}{rgb}{0.909804,0.909804,0.909804}%
\pgfsetstrokecolor{currentstroke}%
\pgfsetdash{}{0pt}%
\pgfpathmoveto{\pgfqpoint{1.395293in}{0.320625in}}%
\pgfpathlineto{\pgfqpoint{1.395293in}{1.485000in}}%
\pgfusepath{stroke}%
\end{pgfscope}%
\begin{pgfscope}%
\pgfsetbuttcap%
\pgfsetroundjoin%
\definecolor{currentfill}{rgb}{0.000000,0.000000,0.000000}%
\pgfsetfillcolor{currentfill}%
\pgfsetlinewidth{0.803000pt}%
\definecolor{currentstroke}{rgb}{0.000000,0.000000,0.000000}%
\pgfsetstrokecolor{currentstroke}%
\pgfsetdash{}{0pt}%
\pgfsys@defobject{currentmarker}{\pgfqpoint{0.000000in}{-0.048611in}}{\pgfqpoint{0.000000in}{0.000000in}}{%
\pgfpathmoveto{\pgfqpoint{0.000000in}{0.000000in}}%
\pgfpathlineto{\pgfqpoint{0.000000in}{-0.048611in}}%
\pgfusepath{stroke,fill}%
}%
\begin{pgfscope}%
\pgfsys@transformshift{1.395293in}{0.320625in}%
\pgfsys@useobject{currentmarker}{}%
\end{pgfscope}%
\end{pgfscope}%
\begin{pgfscope}%
\definecolor{textcolor}{rgb}{0.000000,0.000000,0.000000}%
\pgfsetstrokecolor{textcolor}%
\pgfsetfillcolor{textcolor}%
\pgftext[x=1.395293in,y=0.223403in,,top]{\color{textcolor}{\rmfamily\fontsize{6.000000}{7.200000}\selectfont\catcode`\^=\active\def^{\ifmmode\sp\else\^{}\fi}\catcode`\%=\active\def%{\%}20}}%
\end{pgfscope}%
\begin{pgfscope}%
\pgfpathrectangle{\pgfqpoint{0.487500in}{0.320625in}}{\pgfqpoint{1.815586in}{1.164375in}}%
\pgfusepath{clip}%
\pgfsetrectcap%
\pgfsetroundjoin%
\pgfsetlinewidth{0.803000pt}%
\definecolor{currentstroke}{rgb}{0.909804,0.909804,0.909804}%
\pgfsetstrokecolor{currentstroke}%
\pgfsetdash{}{0pt}%
\pgfpathmoveto{\pgfqpoint{1.849190in}{0.320625in}}%
\pgfpathlineto{\pgfqpoint{1.849190in}{1.485000in}}%
\pgfusepath{stroke}%
\end{pgfscope}%
\begin{pgfscope}%
\pgfsetbuttcap%
\pgfsetroundjoin%
\definecolor{currentfill}{rgb}{0.000000,0.000000,0.000000}%
\pgfsetfillcolor{currentfill}%
\pgfsetlinewidth{0.803000pt}%
\definecolor{currentstroke}{rgb}{0.000000,0.000000,0.000000}%
\pgfsetstrokecolor{currentstroke}%
\pgfsetdash{}{0pt}%
\pgfsys@defobject{currentmarker}{\pgfqpoint{0.000000in}{-0.048611in}}{\pgfqpoint{0.000000in}{0.000000in}}{%
\pgfpathmoveto{\pgfqpoint{0.000000in}{0.000000in}}%
\pgfpathlineto{\pgfqpoint{0.000000in}{-0.048611in}}%
\pgfusepath{stroke,fill}%
}%
\begin{pgfscope}%
\pgfsys@transformshift{1.849190in}{0.320625in}%
\pgfsys@useobject{currentmarker}{}%
\end{pgfscope}%
\end{pgfscope}%
\begin{pgfscope}%
\definecolor{textcolor}{rgb}{0.000000,0.000000,0.000000}%
\pgfsetstrokecolor{textcolor}%
\pgfsetfillcolor{textcolor}%
\pgftext[x=1.849190in,y=0.223403in,,top]{\color{textcolor}{\rmfamily\fontsize{6.000000}{7.200000}\selectfont\catcode`\^=\active\def^{\ifmmode\sp\else\^{}\fi}\catcode`\%=\active\def%{\%}30}}%
\end{pgfscope}%
\begin{pgfscope}%
\pgfpathrectangle{\pgfqpoint{0.487500in}{0.320625in}}{\pgfqpoint{1.815586in}{1.164375in}}%
\pgfusepath{clip}%
\pgfsetrectcap%
\pgfsetroundjoin%
\pgfsetlinewidth{0.803000pt}%
\definecolor{currentstroke}{rgb}{0.909804,0.909804,0.909804}%
\pgfsetstrokecolor{currentstroke}%
\pgfsetdash{}{0pt}%
\pgfpathmoveto{\pgfqpoint{2.303086in}{0.320625in}}%
\pgfpathlineto{\pgfqpoint{2.303086in}{1.485000in}}%
\pgfusepath{stroke}%
\end{pgfscope}%
\begin{pgfscope}%
\pgfsetbuttcap%
\pgfsetroundjoin%
\definecolor{currentfill}{rgb}{0.000000,0.000000,0.000000}%
\pgfsetfillcolor{currentfill}%
\pgfsetlinewidth{0.803000pt}%
\definecolor{currentstroke}{rgb}{0.000000,0.000000,0.000000}%
\pgfsetstrokecolor{currentstroke}%
\pgfsetdash{}{0pt}%
\pgfsys@defobject{currentmarker}{\pgfqpoint{0.000000in}{-0.048611in}}{\pgfqpoint{0.000000in}{0.000000in}}{%
\pgfpathmoveto{\pgfqpoint{0.000000in}{0.000000in}}%
\pgfpathlineto{\pgfqpoint{0.000000in}{-0.048611in}}%
\pgfusepath{stroke,fill}%
}%
\begin{pgfscope}%
\pgfsys@transformshift{2.303086in}{0.320625in}%
\pgfsys@useobject{currentmarker}{}%
\end{pgfscope}%
\end{pgfscope}%
\begin{pgfscope}%
\definecolor{textcolor}{rgb}{0.000000,0.000000,0.000000}%
\pgfsetstrokecolor{textcolor}%
\pgfsetfillcolor{textcolor}%
\pgftext[x=2.303086in,y=0.223403in,,top]{\color{textcolor}{\rmfamily\fontsize{6.000000}{7.200000}\selectfont\catcode`\^=\active\def^{\ifmmode\sp\else\^{}\fi}\catcode`\%=\active\def%{\%}40}}%
\end{pgfscope}%
\begin{pgfscope}%
\definecolor{textcolor}{rgb}{0.000000,0.000000,0.000000}%
\pgfsetstrokecolor{textcolor}%
\pgfsetfillcolor{textcolor}%
\pgftext[x=1.395293in,y=0.084514in,,top]{\color{textcolor}{\rmfamily\fontsize{6.000000}{7.200000}\selectfont\catcode`\^=\active\def^{\ifmmode\sp\else\^{}\fi}\catcode`\%=\active\def%{\%}time}}%
\end{pgfscope}%
\begin{pgfscope}%
\pgfpathrectangle{\pgfqpoint{0.487500in}{0.320625in}}{\pgfqpoint{1.815586in}{1.164375in}}%
\pgfusepath{clip}%
\pgfsetrectcap%
\pgfsetroundjoin%
\pgfsetlinewidth{0.803000pt}%
\definecolor{currentstroke}{rgb}{0.909804,0.909804,0.909804}%
\pgfsetstrokecolor{currentstroke}%
\pgfsetdash{}{0pt}%
\pgfpathmoveto{\pgfqpoint{0.487500in}{0.466172in}}%
\pgfpathlineto{\pgfqpoint{2.303086in}{0.466172in}}%
\pgfusepath{stroke}%
\end{pgfscope}%
\begin{pgfscope}%
\pgfsetbuttcap%
\pgfsetroundjoin%
\definecolor{currentfill}{rgb}{0.000000,0.000000,0.000000}%
\pgfsetfillcolor{currentfill}%
\pgfsetlinewidth{0.803000pt}%
\definecolor{currentstroke}{rgb}{0.000000,0.000000,0.000000}%
\pgfsetstrokecolor{currentstroke}%
\pgfsetdash{}{0pt}%
\pgfsys@defobject{currentmarker}{\pgfqpoint{-0.048611in}{0.000000in}}{\pgfqpoint{-0.000000in}{0.000000in}}{%
\pgfpathmoveto{\pgfqpoint{-0.000000in}{0.000000in}}%
\pgfpathlineto{\pgfqpoint{-0.048611in}{0.000000in}}%
\pgfusepath{stroke,fill}%
}%
\begin{pgfscope}%
\pgfsys@transformshift{0.487500in}{0.466172in}%
\pgfsys@useobject{currentmarker}{}%
\end{pgfscope}%
\end{pgfscope}%
\begin{pgfscope}%
\definecolor{textcolor}{rgb}{0.000000,0.000000,0.000000}%
\pgfsetstrokecolor{textcolor}%
\pgfsetfillcolor{textcolor}%
\pgftext[x=0.111728in, y=0.422067in, left, base]{\color{textcolor}{\rmfamily\fontsize{6.000000}{7.200000}\selectfont\catcode`\^=\active\def^{\ifmmode\sp\else\^{}\fi}\catcode`\%=\active\def%{\%}$\mathdefault{10^{-20}}$}}%
\end{pgfscope}%
\begin{pgfscope}%
\pgfpathrectangle{\pgfqpoint{0.487500in}{0.320625in}}{\pgfqpoint{1.815586in}{1.164375in}}%
\pgfusepath{clip}%
\pgfsetrectcap%
\pgfsetroundjoin%
\pgfsetlinewidth{0.803000pt}%
\definecolor{currentstroke}{rgb}{0.909804,0.909804,0.909804}%
\pgfsetstrokecolor{currentstroke}%
\pgfsetdash{}{0pt}%
\pgfpathmoveto{\pgfqpoint{0.487500in}{0.660234in}}%
\pgfpathlineto{\pgfqpoint{2.303086in}{0.660234in}}%
\pgfusepath{stroke}%
\end{pgfscope}%
\begin{pgfscope}%
\pgfsetbuttcap%
\pgfsetroundjoin%
\definecolor{currentfill}{rgb}{0.000000,0.000000,0.000000}%
\pgfsetfillcolor{currentfill}%
\pgfsetlinewidth{0.803000pt}%
\definecolor{currentstroke}{rgb}{0.000000,0.000000,0.000000}%
\pgfsetstrokecolor{currentstroke}%
\pgfsetdash{}{0pt}%
\pgfsys@defobject{currentmarker}{\pgfqpoint{-0.048611in}{0.000000in}}{\pgfqpoint{-0.000000in}{0.000000in}}{%
\pgfpathmoveto{\pgfqpoint{-0.000000in}{0.000000in}}%
\pgfpathlineto{\pgfqpoint{-0.048611in}{0.000000in}}%
\pgfusepath{stroke,fill}%
}%
\begin{pgfscope}%
\pgfsys@transformshift{0.487500in}{0.660234in}%
\pgfsys@useobject{currentmarker}{}%
\end{pgfscope}%
\end{pgfscope}%
\begin{pgfscope}%
\definecolor{textcolor}{rgb}{0.000000,0.000000,0.000000}%
\pgfsetstrokecolor{textcolor}%
\pgfsetfillcolor{textcolor}%
\pgftext[x=0.111728in, y=0.616129in, left, base]{\color{textcolor}{\rmfamily\fontsize{6.000000}{7.200000}\selectfont\catcode`\^=\active\def^{\ifmmode\sp\else\^{}\fi}\catcode`\%=\active\def%{\%}$\mathdefault{10^{-16}}$}}%
\end{pgfscope}%
\begin{pgfscope}%
\pgfpathrectangle{\pgfqpoint{0.487500in}{0.320625in}}{\pgfqpoint{1.815586in}{1.164375in}}%
\pgfusepath{clip}%
\pgfsetrectcap%
\pgfsetroundjoin%
\pgfsetlinewidth{0.803000pt}%
\definecolor{currentstroke}{rgb}{0.909804,0.909804,0.909804}%
\pgfsetstrokecolor{currentstroke}%
\pgfsetdash{}{0pt}%
\pgfpathmoveto{\pgfqpoint{0.487500in}{0.854297in}}%
\pgfpathlineto{\pgfqpoint{2.303086in}{0.854297in}}%
\pgfusepath{stroke}%
\end{pgfscope}%
\begin{pgfscope}%
\pgfsetbuttcap%
\pgfsetroundjoin%
\definecolor{currentfill}{rgb}{0.000000,0.000000,0.000000}%
\pgfsetfillcolor{currentfill}%
\pgfsetlinewidth{0.803000pt}%
\definecolor{currentstroke}{rgb}{0.000000,0.000000,0.000000}%
\pgfsetstrokecolor{currentstroke}%
\pgfsetdash{}{0pt}%
\pgfsys@defobject{currentmarker}{\pgfqpoint{-0.048611in}{0.000000in}}{\pgfqpoint{-0.000000in}{0.000000in}}{%
\pgfpathmoveto{\pgfqpoint{-0.000000in}{0.000000in}}%
\pgfpathlineto{\pgfqpoint{-0.048611in}{0.000000in}}%
\pgfusepath{stroke,fill}%
}%
\begin{pgfscope}%
\pgfsys@transformshift{0.487500in}{0.854297in}%
\pgfsys@useobject{currentmarker}{}%
\end{pgfscope}%
\end{pgfscope}%
\begin{pgfscope}%
\definecolor{textcolor}{rgb}{0.000000,0.000000,0.000000}%
\pgfsetstrokecolor{textcolor}%
\pgfsetfillcolor{textcolor}%
\pgftext[x=0.111728in, y=0.810192in, left, base]{\color{textcolor}{\rmfamily\fontsize{6.000000}{7.200000}\selectfont\catcode`\^=\active\def^{\ifmmode\sp\else\^{}\fi}\catcode`\%=\active\def%{\%}$\mathdefault{10^{-12}}$}}%
\end{pgfscope}%
\begin{pgfscope}%
\pgfpathrectangle{\pgfqpoint{0.487500in}{0.320625in}}{\pgfqpoint{1.815586in}{1.164375in}}%
\pgfusepath{clip}%
\pgfsetrectcap%
\pgfsetroundjoin%
\pgfsetlinewidth{0.803000pt}%
\definecolor{currentstroke}{rgb}{0.909804,0.909804,0.909804}%
\pgfsetstrokecolor{currentstroke}%
\pgfsetdash{}{0pt}%
\pgfpathmoveto{\pgfqpoint{0.487500in}{1.048359in}}%
\pgfpathlineto{\pgfqpoint{2.303086in}{1.048359in}}%
\pgfusepath{stroke}%
\end{pgfscope}%
\begin{pgfscope}%
\pgfsetbuttcap%
\pgfsetroundjoin%
\definecolor{currentfill}{rgb}{0.000000,0.000000,0.000000}%
\pgfsetfillcolor{currentfill}%
\pgfsetlinewidth{0.803000pt}%
\definecolor{currentstroke}{rgb}{0.000000,0.000000,0.000000}%
\pgfsetstrokecolor{currentstroke}%
\pgfsetdash{}{0pt}%
\pgfsys@defobject{currentmarker}{\pgfqpoint{-0.048611in}{0.000000in}}{\pgfqpoint{-0.000000in}{0.000000in}}{%
\pgfpathmoveto{\pgfqpoint{-0.000000in}{0.000000in}}%
\pgfpathlineto{\pgfqpoint{-0.048611in}{0.000000in}}%
\pgfusepath{stroke,fill}%
}%
\begin{pgfscope}%
\pgfsys@transformshift{0.487500in}{1.048359in}%
\pgfsys@useobject{currentmarker}{}%
\end{pgfscope}%
\end{pgfscope}%
\begin{pgfscope}%
\definecolor{textcolor}{rgb}{0.000000,0.000000,0.000000}%
\pgfsetstrokecolor{textcolor}%
\pgfsetfillcolor{textcolor}%
\pgftext[x=0.158989in, y=1.004254in, left, base]{\color{textcolor}{\rmfamily\fontsize{6.000000}{7.200000}\selectfont\catcode`\^=\active\def^{\ifmmode\sp\else\^{}\fi}\catcode`\%=\active\def%{\%}$\mathdefault{10^{-8}}$}}%
\end{pgfscope}%
\begin{pgfscope}%
\pgfpathrectangle{\pgfqpoint{0.487500in}{0.320625in}}{\pgfqpoint{1.815586in}{1.164375in}}%
\pgfusepath{clip}%
\pgfsetrectcap%
\pgfsetroundjoin%
\pgfsetlinewidth{0.803000pt}%
\definecolor{currentstroke}{rgb}{0.909804,0.909804,0.909804}%
\pgfsetstrokecolor{currentstroke}%
\pgfsetdash{}{0pt}%
\pgfpathmoveto{\pgfqpoint{0.487500in}{1.242422in}}%
\pgfpathlineto{\pgfqpoint{2.303086in}{1.242422in}}%
\pgfusepath{stroke}%
\end{pgfscope}%
\begin{pgfscope}%
\pgfsetbuttcap%
\pgfsetroundjoin%
\definecolor{currentfill}{rgb}{0.000000,0.000000,0.000000}%
\pgfsetfillcolor{currentfill}%
\pgfsetlinewidth{0.803000pt}%
\definecolor{currentstroke}{rgb}{0.000000,0.000000,0.000000}%
\pgfsetstrokecolor{currentstroke}%
\pgfsetdash{}{0pt}%
\pgfsys@defobject{currentmarker}{\pgfqpoint{-0.048611in}{0.000000in}}{\pgfqpoint{-0.000000in}{0.000000in}}{%
\pgfpathmoveto{\pgfqpoint{-0.000000in}{0.000000in}}%
\pgfpathlineto{\pgfqpoint{-0.048611in}{0.000000in}}%
\pgfusepath{stroke,fill}%
}%
\begin{pgfscope}%
\pgfsys@transformshift{0.487500in}{1.242422in}%
\pgfsys@useobject{currentmarker}{}%
\end{pgfscope}%
\end{pgfscope}%
\begin{pgfscope}%
\definecolor{textcolor}{rgb}{0.000000,0.000000,0.000000}%
\pgfsetstrokecolor{textcolor}%
\pgfsetfillcolor{textcolor}%
\pgftext[x=0.158989in, y=1.198317in, left, base]{\color{textcolor}{\rmfamily\fontsize{6.000000}{7.200000}\selectfont\catcode`\^=\active\def^{\ifmmode\sp\else\^{}\fi}\catcode`\%=\active\def%{\%}$\mathdefault{10^{-4}}$}}%
\end{pgfscope}%
\begin{pgfscope}%
\pgfpathrectangle{\pgfqpoint{0.487500in}{0.320625in}}{\pgfqpoint{1.815586in}{1.164375in}}%
\pgfusepath{clip}%
\pgfsetrectcap%
\pgfsetroundjoin%
\pgfsetlinewidth{0.803000pt}%
\definecolor{currentstroke}{rgb}{0.909804,0.909804,0.909804}%
\pgfsetstrokecolor{currentstroke}%
\pgfsetdash{}{0pt}%
\pgfpathmoveto{\pgfqpoint{0.487500in}{1.436484in}}%
\pgfpathlineto{\pgfqpoint{2.303086in}{1.436484in}}%
\pgfusepath{stroke}%
\end{pgfscope}%
\begin{pgfscope}%
\pgfsetbuttcap%
\pgfsetroundjoin%
\definecolor{currentfill}{rgb}{0.000000,0.000000,0.000000}%
\pgfsetfillcolor{currentfill}%
\pgfsetlinewidth{0.803000pt}%
\definecolor{currentstroke}{rgb}{0.000000,0.000000,0.000000}%
\pgfsetstrokecolor{currentstroke}%
\pgfsetdash{}{0pt}%
\pgfsys@defobject{currentmarker}{\pgfqpoint{-0.048611in}{0.000000in}}{\pgfqpoint{-0.000000in}{0.000000in}}{%
\pgfpathmoveto{\pgfqpoint{-0.000000in}{0.000000in}}%
\pgfpathlineto{\pgfqpoint{-0.048611in}{0.000000in}}%
\pgfusepath{stroke,fill}%
}%
\begin{pgfscope}%
\pgfsys@transformshift{0.487500in}{1.436484in}%
\pgfsys@useobject{currentmarker}{}%
\end{pgfscope}%
\end{pgfscope}%
\begin{pgfscope}%
\definecolor{textcolor}{rgb}{0.000000,0.000000,0.000000}%
\pgfsetstrokecolor{textcolor}%
\pgfsetfillcolor{textcolor}%
\pgftext[x=0.234222in, y=1.393151in, left, base]{\color{textcolor}{\rmfamily\fontsize{6.000000}{7.200000}\selectfont\catcode`\^=\active\def^{\ifmmode\sp\else\^{}\fi}\catcode`\%=\active\def%{\%}$\mathdefault{10^{0}}$}}%
\end{pgfscope}%
\begin{pgfscope}%
\definecolor{textcolor}{rgb}{0.000000,0.000000,0.000000}%
\pgfsetstrokecolor{textcolor}%
\pgfsetfillcolor{textcolor}%
\pgftext[x=0.056172in,y=0.902813in,,bottom,rotate=90.000000]{\color{textcolor}{\rmfamily\fontsize{6.000000}{7.200000}\selectfont\catcode`\^=\active\def^{\ifmmode\sp\else\^{}\fi}\catcode`\%=\active\def%{\%}$\ell_t-\ell_\star$}}%
\end{pgfscope}%
\begin{pgfscope}%
\pgfpathrectangle{\pgfqpoint{0.487500in}{0.320625in}}{\pgfqpoint{1.815586in}{1.164375in}}%
\pgfusepath{clip}%
\pgfsetrectcap%
\pgfsetroundjoin%
\pgfsetlinewidth{1.254687pt}%
\definecolor{currentstroke}{rgb}{0.141176,0.345098,0.650980}%
\pgfsetstrokecolor{currentstroke}%
\pgfsetdash{}{0pt}%
\pgfpathmoveto{\pgfqpoint{0.487500in}{1.451682in}}%
\pgfpathlineto{\pgfqpoint{0.496578in}{1.449020in}}%
\pgfpathlineto{\pgfqpoint{0.517003in}{1.443375in}}%
\pgfpathlineto{\pgfqpoint{0.528351in}{1.439127in}}%
\pgfpathlineto{\pgfqpoint{0.539698in}{1.433676in}}%
\pgfpathlineto{\pgfqpoint{0.551046in}{1.426923in}}%
\pgfpathlineto{\pgfqpoint{0.564662in}{1.417276in}}%
\pgfpathlineto{\pgfqpoint{0.582818in}{1.402686in}}%
\pgfpathlineto{\pgfqpoint{0.632747in}{1.361662in}}%
\pgfpathlineto{\pgfqpoint{0.657711in}{1.343005in}}%
\pgfpathlineto{\pgfqpoint{0.687215in}{1.322384in}}%
\pgfpathlineto{\pgfqpoint{0.721257in}{1.300031in}}%
\pgfpathlineto{\pgfqpoint{0.757568in}{1.277597in}}%
\pgfpathlineto{\pgfqpoint{0.796150in}{1.255122in}}%
\pgfpathlineto{\pgfqpoint{0.837000in}{1.232665in}}%
\pgfpathlineto{\pgfqpoint{0.880121in}{1.210249in}}%
\pgfpathlineto{\pgfqpoint{0.927780in}{1.186751in}}%
\pgfpathlineto{\pgfqpoint{0.979978in}{1.162248in}}%
\pgfpathlineto{\pgfqpoint{1.038984in}{1.135766in}}%
\pgfpathlineto{\pgfqpoint{1.104799in}{1.107436in}}%
\pgfpathlineto{\pgfqpoint{1.177423in}{1.077392in}}%
\pgfpathlineto{\pgfqpoint{1.254585in}{1.046666in}}%
\pgfpathlineto{\pgfqpoint{1.338556in}{1.014445in}}%
\pgfpathlineto{\pgfqpoint{1.429335in}{0.980816in}}%
\pgfpathlineto{\pgfqpoint{1.533732in}{0.943371in}}%
\pgfpathlineto{\pgfqpoint{1.656284in}{0.900634in}}%
\pgfpathlineto{\pgfqpoint{1.812878in}{0.847258in}}%
\pgfpathlineto{\pgfqpoint{2.030748in}{0.774252in}}%
\pgfpathlineto{\pgfqpoint{2.303086in}{0.683958in}}%
\pgfpathlineto{\pgfqpoint{2.303086in}{0.683958in}}%
\pgfusepath{stroke}%
\end{pgfscope}%
\begin{pgfscope}%
\pgfpathrectangle{\pgfqpoint{0.487500in}{0.320625in}}{\pgfqpoint{1.815586in}{1.164375in}}%
\pgfusepath{clip}%
\pgfsetbuttcap%
\pgfsetroundjoin%
\pgfsetlinewidth{0.752812pt}%
\definecolor{currentstroke}{rgb}{0.141176,0.345098,0.650980}%
\pgfsetstrokecolor{currentstroke}%
\pgfsetstrokeopacity{0.650000}%
\pgfsetdash{{2.775000pt}{1.200000pt}}{0.000000pt}%
\pgfpathmoveto{\pgfqpoint{0.487500in}{1.451682in}}%
\pgfpathlineto{\pgfqpoint{2.303086in}{1.114561in}}%
\pgfpathlineto{\pgfqpoint{2.303086in}{1.114561in}}%
\pgfusepath{stroke}%
\end{pgfscope}%
\begin{pgfscope}%
\pgfpathrectangle{\pgfqpoint{0.487500in}{0.320625in}}{\pgfqpoint{1.815586in}{1.164375in}}%
\pgfusepath{clip}%
\pgfsetbuttcap%
\pgfsetroundjoin%
\pgfsetlinewidth{0.803000pt}%
\definecolor{currentstroke}{rgb}{0.141176,0.345098,0.650980}%
\pgfsetstrokecolor{currentstroke}%
\pgfsetstrokeopacity{0.800000}%
\pgfsetdash{{0.800000pt}{1.320000pt}}{0.000000pt}%
\pgfpathmoveto{\pgfqpoint{0.487500in}{1.451682in}}%
\pgfpathlineto{\pgfqpoint{2.303086in}{1.451682in}}%
\pgfpathlineto{\pgfqpoint{2.303086in}{1.451682in}}%
\pgfusepath{stroke}%
\end{pgfscope}%
\begin{pgfscope}%
\pgfpathrectangle{\pgfqpoint{0.487500in}{0.320625in}}{\pgfqpoint{1.815586in}{1.164375in}}%
\pgfusepath{clip}%
\pgfsetrectcap%
\pgfsetroundjoin%
\pgfsetlinewidth{1.254687pt}%
\definecolor{currentstroke}{rgb}{0.819608,0.286275,0.356863}%
\pgfsetstrokecolor{currentstroke}%
\pgfsetdash{}{0pt}%
\pgfpathmoveto{\pgfqpoint{0.487500in}{1.451682in}}%
\pgfpathlineto{\pgfqpoint{0.492039in}{1.446133in}}%
\pgfpathlineto{\pgfqpoint{0.507925in}{1.429504in}}%
\pgfpathlineto{\pgfqpoint{0.519273in}{1.415878in}}%
\pgfpathlineto{\pgfqpoint{0.532890in}{1.397493in}}%
\pgfpathlineto{\pgfqpoint{0.551046in}{1.370778in}}%
\pgfpathlineto{\pgfqpoint{0.607783in}{1.285872in}}%
\pgfpathlineto{\pgfqpoint{0.635016in}{1.247642in}}%
\pgfpathlineto{\pgfqpoint{0.664520in}{1.208201in}}%
\pgfpathlineto{\pgfqpoint{0.698562in}{1.164684in}}%
\pgfpathlineto{\pgfqpoint{0.737143in}{1.117217in}}%
\pgfpathlineto{\pgfqpoint{0.784802in}{1.060424in}}%
\pgfpathlineto{\pgfqpoint{0.843809in}{0.991967in}}%
\pgfpathlineto{\pgfqpoint{0.914163in}{0.912098in}}%
\pgfpathlineto{\pgfqpoint{1.000403in}{0.815918in}}%
\pgfpathlineto{\pgfqpoint{1.102530in}{0.703716in}}%
\pgfpathlineto{\pgfqpoint{1.202387in}{0.595445in}}%
\pgfpathlineto{\pgfqpoint{1.209196in}{0.588828in}}%
\pgfpathlineto{\pgfqpoint{1.213735in}{0.583075in}}%
\pgfpathlineto{\pgfqpoint{1.216004in}{0.581228in}}%
\pgfpathlineto{\pgfqpoint{1.220543in}{0.575773in}}%
\pgfpathlineto{\pgfqpoint{1.234160in}{0.561107in}}%
\pgfpathlineto{\pgfqpoint{1.236429in}{0.560744in}}%
\pgfpathlineto{\pgfqpoint{1.238699in}{0.556161in}}%
\pgfpathlineto{\pgfqpoint{1.240968in}{0.558660in}}%
\pgfpathlineto{\pgfqpoint{1.243238in}{0.551340in}}%
\pgfpathlineto{\pgfqpoint{1.245507in}{0.549293in}}%
\pgfpathlineto{\pgfqpoint{1.247777in}{0.548143in}}%
\pgfpathlineto{\pgfqpoint{1.252316in}{0.541882in}}%
\pgfpathlineto{\pgfqpoint{1.254585in}{0.539261in}}%
\pgfpathlineto{\pgfqpoint{1.256855in}{0.538232in}}%
\pgfpathlineto{\pgfqpoint{1.261394in}{0.532559in}}%
\pgfpathlineto{\pgfqpoint{1.263663in}{0.532171in}}%
\pgfpathlineto{\pgfqpoint{1.265933in}{0.527485in}}%
\pgfpathlineto{\pgfqpoint{1.268202in}{0.532444in}}%
\pgfpathlineto{\pgfqpoint{1.270472in}{0.523897in}}%
\pgfpathlineto{\pgfqpoint{1.272741in}{0.534109in}}%
\pgfpathlineto{\pgfqpoint{1.275011in}{0.521570in}}%
\pgfpathlineto{\pgfqpoint{1.279550in}{0.514498in}}%
\pgfpathlineto{\pgfqpoint{1.281819in}{0.519475in}}%
\pgfpathlineto{\pgfqpoint{1.284089in}{0.508300in}}%
\pgfpathlineto{\pgfqpoint{1.286358in}{0.505012in}}%
\pgfpathlineto{\pgfqpoint{1.288628in}{0.512398in}}%
\pgfpathlineto{\pgfqpoint{1.290897in}{0.500507in}}%
\pgfpathlineto{\pgfqpoint{1.293166in}{0.505005in}}%
\pgfpathlineto{\pgfqpoint{1.295436in}{0.499505in}}%
\pgfpathlineto{\pgfqpoint{1.297705in}{0.506227in}}%
\pgfpathlineto{\pgfqpoint{1.299975in}{0.507846in}}%
\pgfpathlineto{\pgfqpoint{1.302244in}{0.508319in}}%
\pgfpathlineto{\pgfqpoint{1.304514in}{0.510522in}}%
\pgfpathlineto{\pgfqpoint{1.306783in}{0.515241in}}%
\pgfpathlineto{\pgfqpoint{1.309053in}{0.489116in}}%
\pgfpathlineto{\pgfqpoint{1.311322in}{0.478212in}}%
\pgfpathlineto{\pgfqpoint{1.313592in}{0.512270in}}%
\pgfpathlineto{\pgfqpoint{1.315861in}{0.473355in}}%
\pgfpathlineto{\pgfqpoint{1.318131in}{0.477266in}}%
\pgfpathlineto{\pgfqpoint{1.320400in}{0.492566in}}%
\pgfpathlineto{\pgfqpoint{1.322670in}{0.470862in}}%
\pgfpathlineto{\pgfqpoint{1.324939in}{0.489400in}}%
\pgfpathlineto{\pgfqpoint{1.327209in}{0.461410in}}%
\pgfpathlineto{\pgfqpoint{1.329478in}{0.503355in}}%
\pgfpathlineto{\pgfqpoint{1.331748in}{0.475867in}}%
\pgfpathlineto{\pgfqpoint{1.334017in}{0.509244in}}%
\pgfpathlineto{\pgfqpoint{1.336287in}{0.474436in}}%
\pgfpathlineto{\pgfqpoint{1.338556in}{0.522885in}}%
\pgfpathlineto{\pgfqpoint{1.340826in}{0.487738in}}%
\pgfpathlineto{\pgfqpoint{1.343095in}{0.511084in}}%
\pgfpathlineto{\pgfqpoint{1.345365in}{0.510820in}}%
\pgfpathlineto{\pgfqpoint{1.347634in}{0.520139in}}%
\pgfpathlineto{\pgfqpoint{1.352173in}{0.434578in}}%
\pgfpathlineto{\pgfqpoint{1.354443in}{0.519425in}}%
\pgfpathlineto{\pgfqpoint{1.356712in}{0.434834in}}%
\pgfpathlineto{\pgfqpoint{1.358981in}{0.447404in}}%
\pgfpathlineto{\pgfqpoint{1.361251in}{0.492711in}}%
\pgfpathlineto{\pgfqpoint{1.363520in}{0.439546in}}%
\pgfpathlineto{\pgfqpoint{1.365790in}{0.470321in}}%
\pgfpathlineto{\pgfqpoint{1.368059in}{0.452341in}}%
\pgfpathlineto{\pgfqpoint{1.370329in}{0.484337in}}%
\pgfpathlineto{\pgfqpoint{1.372598in}{0.483507in}}%
\pgfpathlineto{\pgfqpoint{1.374868in}{0.486649in}}%
\pgfpathlineto{\pgfqpoint{1.377137in}{0.493444in}}%
\pgfpathlineto{\pgfqpoint{1.379407in}{0.493876in}}%
\pgfpathlineto{\pgfqpoint{1.381676in}{0.512305in}}%
\pgfpathlineto{\pgfqpoint{1.383946in}{0.498961in}}%
\pgfpathlineto{\pgfqpoint{1.386215in}{0.523349in}}%
\pgfpathlineto{\pgfqpoint{1.388485in}{0.507143in}}%
\pgfpathlineto{\pgfqpoint{1.390754in}{0.504787in}}%
\pgfpathlineto{\pgfqpoint{1.393024in}{0.393449in}}%
\pgfpathlineto{\pgfqpoint{1.395293in}{0.530519in}}%
\pgfpathlineto{\pgfqpoint{1.397563in}{0.429607in}}%
\pgfpathlineto{\pgfqpoint{1.399832in}{0.442790in}}%
\pgfpathlineto{\pgfqpoint{1.402102in}{0.501694in}}%
\pgfpathlineto{\pgfqpoint{1.404371in}{0.446939in}}%
\pgfpathlineto{\pgfqpoint{1.406641in}{0.454164in}}%
\pgfpathlineto{\pgfqpoint{1.408910in}{0.465387in}}%
\pgfpathlineto{\pgfqpoint{1.411180in}{0.464125in}}%
\pgfpathlineto{\pgfqpoint{1.413449in}{0.477980in}}%
\pgfpathlineto{\pgfqpoint{1.415719in}{0.443821in}}%
\pgfpathlineto{\pgfqpoint{1.417988in}{0.496444in}}%
\pgfpathlineto{\pgfqpoint{1.420258in}{0.362166in}}%
\pgfpathlineto{\pgfqpoint{1.422527in}{0.514489in}}%
\pgfpathlineto{\pgfqpoint{1.424796in}{0.446974in}}%
\pgfpathlineto{\pgfqpoint{1.427066in}{0.526599in}}%
\pgfpathlineto{\pgfqpoint{1.429335in}{0.476790in}}%
\pgfpathlineto{\pgfqpoint{1.431605in}{0.470776in}}%
\pgfpathlineto{\pgfqpoint{1.433874in}{0.485072in}}%
\pgfpathlineto{\pgfqpoint{1.436144in}{0.487509in}}%
\pgfpathlineto{\pgfqpoint{1.438413in}{0.446185in}}%
\pgfpathlineto{\pgfqpoint{1.440683in}{0.419473in}}%
\pgfpathlineto{\pgfqpoint{1.442952in}{0.483787in}}%
\pgfpathlineto{\pgfqpoint{1.445222in}{0.455909in}}%
\pgfpathlineto{\pgfqpoint{1.449761in}{0.480775in}}%
\pgfpathlineto{\pgfqpoint{1.452030in}{0.465472in}}%
\pgfpathlineto{\pgfqpoint{1.454300in}{0.507591in}}%
\pgfpathlineto{\pgfqpoint{1.456569in}{0.488962in}}%
\pgfpathlineto{\pgfqpoint{1.458839in}{0.514656in}}%
\pgfpathlineto{\pgfqpoint{1.461108in}{0.503031in}}%
\pgfpathlineto{\pgfqpoint{1.465647in}{0.440791in}}%
\pgfpathlineto{\pgfqpoint{1.467917in}{0.503462in}}%
\pgfpathlineto{\pgfqpoint{1.470186in}{0.413916in}}%
\pgfpathlineto{\pgfqpoint{1.472456in}{0.415050in}}%
\pgfpathlineto{\pgfqpoint{1.474725in}{0.487866in}}%
\pgfpathlineto{\pgfqpoint{1.476995in}{0.366548in}}%
\pgfpathlineto{\pgfqpoint{1.479264in}{0.479990in}}%
\pgfpathlineto{\pgfqpoint{1.481534in}{0.445790in}}%
\pgfpathlineto{\pgfqpoint{1.483803in}{0.494550in}}%
\pgfpathlineto{\pgfqpoint{1.486073in}{0.490277in}}%
\pgfpathlineto{\pgfqpoint{1.488342in}{0.506213in}}%
\pgfpathlineto{\pgfqpoint{1.490611in}{0.490146in}}%
\pgfpathlineto{\pgfqpoint{1.492881in}{0.519981in}}%
\pgfpathlineto{\pgfqpoint{1.495150in}{0.442266in}}%
\pgfpathlineto{\pgfqpoint{1.497420in}{0.427205in}}%
\pgfpathlineto{\pgfqpoint{1.499689in}{0.506982in}}%
\pgfpathlineto{\pgfqpoint{1.501959in}{0.424602in}}%
\pgfpathlineto{\pgfqpoint{1.506498in}{0.475977in}}%
\pgfpathlineto{\pgfqpoint{1.508767in}{0.452726in}}%
\pgfpathlineto{\pgfqpoint{1.511037in}{0.479219in}}%
\pgfpathlineto{\pgfqpoint{1.513306in}{0.400766in}}%
\pgfpathlineto{\pgfqpoint{1.515576in}{0.498632in}}%
\pgfpathlineto{\pgfqpoint{1.517845in}{0.444689in}}%
\pgfpathlineto{\pgfqpoint{1.520115in}{0.511155in}}%
\pgfpathlineto{\pgfqpoint{1.522384in}{0.413266in}}%
\pgfpathlineto{\pgfqpoint{1.524654in}{0.525507in}}%
\pgfpathlineto{\pgfqpoint{1.526923in}{0.420130in}}%
\pgfpathlineto{\pgfqpoint{1.529193in}{0.467872in}}%
\pgfpathlineto{\pgfqpoint{1.531462in}{0.496321in}}%
\pgfpathlineto{\pgfqpoint{1.533732in}{0.474423in}}%
\pgfpathlineto{\pgfqpoint{1.538271in}{0.444297in}}%
\pgfpathlineto{\pgfqpoint{1.540540in}{0.482654in}}%
\pgfpathlineto{\pgfqpoint{1.542810in}{0.466436in}}%
\pgfpathlineto{\pgfqpoint{1.545079in}{0.462239in}}%
\pgfpathlineto{\pgfqpoint{1.547349in}{0.489899in}}%
\pgfpathlineto{\pgfqpoint{1.549618in}{0.448732in}}%
\pgfpathlineto{\pgfqpoint{1.551888in}{0.508324in}}%
\pgfpathlineto{\pgfqpoint{1.554157in}{0.468358in}}%
\pgfpathlineto{\pgfqpoint{1.556427in}{0.520076in}}%
\pgfpathlineto{\pgfqpoint{1.558696in}{0.483210in}}%
\pgfpathlineto{\pgfqpoint{1.563235in}{0.469466in}}%
\pgfpathlineto{\pgfqpoint{1.565504in}{0.497247in}}%
\pgfpathlineto{\pgfqpoint{1.567774in}{0.444778in}}%
\pgfpathlineto{\pgfqpoint{1.570043in}{0.368565in}}%
\pgfpathlineto{\pgfqpoint{1.572313in}{0.492066in}}%
\pgfpathlineto{\pgfqpoint{1.574582in}{0.423481in}}%
\pgfpathlineto{\pgfqpoint{1.576852in}{0.478897in}}%
\pgfpathlineto{\pgfqpoint{1.579121in}{0.466423in}}%
\pgfpathlineto{\pgfqpoint{1.581391in}{0.489261in}}%
\pgfpathlineto{\pgfqpoint{1.583660in}{0.496662in}}%
\pgfpathlineto{\pgfqpoint{1.585930in}{0.496677in}}%
\pgfpathlineto{\pgfqpoint{1.588199in}{0.504545in}}%
\pgfpathlineto{\pgfqpoint{1.590469in}{0.506210in}}%
\pgfpathlineto{\pgfqpoint{1.592738in}{0.509993in}}%
\pgfpathlineto{\pgfqpoint{1.595008in}{0.487463in}}%
\pgfpathlineto{\pgfqpoint{1.597277in}{0.529370in}}%
\pgfpathlineto{\pgfqpoint{1.599547in}{0.492857in}}%
\pgfpathlineto{\pgfqpoint{1.601816in}{0.488253in}}%
\pgfpathlineto{\pgfqpoint{1.604086in}{0.480103in}}%
\pgfpathlineto{\pgfqpoint{1.606355in}{0.508978in}}%
\pgfpathlineto{\pgfqpoint{1.608625in}{0.433381in}}%
\pgfpathlineto{\pgfqpoint{1.610894in}{0.381727in}}%
\pgfpathlineto{\pgfqpoint{1.613164in}{0.494780in}}%
\pgfpathlineto{\pgfqpoint{1.615433in}{0.420897in}}%
\pgfpathlineto{\pgfqpoint{1.617703in}{0.456648in}}%
\pgfpathlineto{\pgfqpoint{1.619972in}{0.478877in}}%
\pgfpathlineto{\pgfqpoint{1.622242in}{0.446681in}}%
\pgfpathlineto{\pgfqpoint{1.624511in}{0.495871in}}%
\pgfpathlineto{\pgfqpoint{1.626780in}{0.453686in}}%
\pgfpathlineto{\pgfqpoint{1.629050in}{0.509218in}}%
\pgfpathlineto{\pgfqpoint{1.631319in}{0.419058in}}%
\pgfpathlineto{\pgfqpoint{1.633589in}{0.520539in}}%
\pgfpathlineto{\pgfqpoint{1.635858in}{0.442500in}}%
\pgfpathlineto{\pgfqpoint{1.638128in}{0.534540in}}%
\pgfpathlineto{\pgfqpoint{1.640397in}{0.451281in}}%
\pgfpathlineto{\pgfqpoint{1.642667in}{0.507119in}}%
\pgfpathlineto{\pgfqpoint{1.644936in}{0.493253in}}%
\pgfpathlineto{\pgfqpoint{1.647206in}{0.521304in}}%
\pgfpathlineto{\pgfqpoint{1.649475in}{0.434336in}}%
\pgfpathlineto{\pgfqpoint{1.651745in}{0.402581in}}%
\pgfpathlineto{\pgfqpoint{1.654014in}{0.503911in}}%
\pgfpathlineto{\pgfqpoint{1.656284in}{0.374702in}}%
\pgfpathlineto{\pgfqpoint{1.660823in}{0.483382in}}%
\pgfpathlineto{\pgfqpoint{1.663092in}{0.376953in}}%
\pgfpathlineto{\pgfqpoint{1.665362in}{0.482716in}}%
\pgfpathlineto{\pgfqpoint{1.667631in}{0.408347in}}%
\pgfpathlineto{\pgfqpoint{1.669901in}{0.499256in}}%
\pgfpathlineto{\pgfqpoint{1.672170in}{0.479238in}}%
\pgfpathlineto{\pgfqpoint{1.674440in}{0.513422in}}%
\pgfpathlineto{\pgfqpoint{1.676709in}{0.466094in}}%
\pgfpathlineto{\pgfqpoint{1.678979in}{0.528164in}}%
\pgfpathlineto{\pgfqpoint{1.681248in}{0.352607in}}%
\pgfpathlineto{\pgfqpoint{1.683518in}{0.450562in}}%
\pgfpathlineto{\pgfqpoint{1.685787in}{0.503485in}}%
\pgfpathlineto{\pgfqpoint{1.688057in}{0.454905in}}%
\pgfpathlineto{\pgfqpoint{1.690326in}{0.449006in}}%
\pgfpathlineto{\pgfqpoint{1.692595in}{0.465929in}}%
\pgfpathlineto{\pgfqpoint{1.694865in}{0.460534in}}%
\pgfpathlineto{\pgfqpoint{1.697134in}{0.474691in}}%
\pgfpathlineto{\pgfqpoint{1.699404in}{0.429723in}}%
\pgfpathlineto{\pgfqpoint{1.701673in}{0.494506in}}%
\pgfpathlineto{\pgfqpoint{1.703943in}{0.412898in}}%
\pgfpathlineto{\pgfqpoint{1.706212in}{0.513299in}}%
\pgfpathlineto{\pgfqpoint{1.708482in}{0.426667in}}%
\pgfpathlineto{\pgfqpoint{1.710751in}{0.525653in}}%
\pgfpathlineto{\pgfqpoint{1.713021in}{0.468229in}}%
\pgfpathlineto{\pgfqpoint{1.715290in}{0.471219in}}%
\pgfpathlineto{\pgfqpoint{1.717560in}{0.486291in}}%
\pgfpathlineto{\pgfqpoint{1.719829in}{0.485686in}}%
\pgfpathlineto{\pgfqpoint{1.724368in}{0.419639in}}%
\pgfpathlineto{\pgfqpoint{1.726638in}{0.484819in}}%
\pgfpathlineto{\pgfqpoint{1.728907in}{0.456906in}}%
\pgfpathlineto{\pgfqpoint{1.733446in}{0.481595in}}%
\pgfpathlineto{\pgfqpoint{1.735716in}{0.467884in}}%
\pgfpathlineto{\pgfqpoint{1.737985in}{0.507209in}}%
\pgfpathlineto{\pgfqpoint{1.740255in}{0.488227in}}%
\pgfpathlineto{\pgfqpoint{1.742524in}{0.515084in}}%
\pgfpathlineto{\pgfqpoint{1.744794in}{0.501779in}}%
\pgfpathlineto{\pgfqpoint{1.749333in}{0.443606in}}%
\pgfpathlineto{\pgfqpoint{1.751602in}{0.503253in}}%
\pgfpathlineto{\pgfqpoint{1.753872in}{0.417606in}}%
\pgfpathlineto{\pgfqpoint{1.756141in}{0.413995in}}%
\pgfpathlineto{\pgfqpoint{1.758410in}{0.488464in}}%
\pgfpathlineto{\pgfqpoint{1.760680in}{0.370264in}}%
\pgfpathlineto{\pgfqpoint{1.762949in}{0.479943in}}%
\pgfpathlineto{\pgfqpoint{1.765219in}{0.447734in}}%
\pgfpathlineto{\pgfqpoint{1.767488in}{0.494266in}}%
\pgfpathlineto{\pgfqpoint{1.769758in}{0.490610in}}%
\pgfpathlineto{\pgfqpoint{1.772027in}{0.505193in}}%
\pgfpathlineto{\pgfqpoint{1.774297in}{0.491748in}}%
\pgfpathlineto{\pgfqpoint{1.776566in}{0.518707in}}%
\pgfpathlineto{\pgfqpoint{1.778836in}{0.446904in}}%
\pgfpathlineto{\pgfqpoint{1.781105in}{0.423335in}}%
\pgfpathlineto{\pgfqpoint{1.783375in}{0.507242in}}%
\pgfpathlineto{\pgfqpoint{1.785644in}{0.419507in}}%
\pgfpathlineto{\pgfqpoint{1.790183in}{0.476975in}}%
\pgfpathlineto{\pgfqpoint{1.792453in}{0.452364in}}%
\pgfpathlineto{\pgfqpoint{1.794722in}{0.479648in}}%
\pgfpathlineto{\pgfqpoint{1.796992in}{0.396574in}}%
\pgfpathlineto{\pgfqpoint{1.799261in}{0.499074in}}%
\pgfpathlineto{\pgfqpoint{1.801531in}{0.446176in}}%
\pgfpathlineto{\pgfqpoint{1.803800in}{0.510656in}}%
\pgfpathlineto{\pgfqpoint{1.806070in}{0.423364in}}%
\pgfpathlineto{\pgfqpoint{1.808339in}{0.525162in}}%
\pgfpathlineto{\pgfqpoint{1.810609in}{0.407480in}}%
\pgfpathlineto{\pgfqpoint{1.812878in}{0.467312in}}%
\pgfpathlineto{\pgfqpoint{1.815148in}{0.497235in}}%
\pgfpathlineto{\pgfqpoint{1.817417in}{0.472932in}}%
\pgfpathlineto{\pgfqpoint{1.821956in}{0.446226in}}%
\pgfpathlineto{\pgfqpoint{1.824226in}{0.482669in}}%
\pgfpathlineto{\pgfqpoint{1.826495in}{0.467218in}}%
\pgfpathlineto{\pgfqpoint{1.828764in}{0.461594in}}%
\pgfpathlineto{\pgfqpoint{1.831034in}{0.490710in}}%
\pgfpathlineto{\pgfqpoint{1.833303in}{0.446655in}}%
\pgfpathlineto{\pgfqpoint{1.835573in}{0.508164in}}%
\pgfpathlineto{\pgfqpoint{1.837842in}{0.465095in}}%
\pgfpathlineto{\pgfqpoint{1.840112in}{0.520364in}}%
\pgfpathlineto{\pgfqpoint{1.842381in}{0.480163in}}%
\pgfpathlineto{\pgfqpoint{1.844651in}{0.476777in}}%
\pgfpathlineto{\pgfqpoint{1.846920in}{0.471613in}}%
\pgfpathlineto{\pgfqpoint{1.849190in}{0.496478in}}%
\pgfpathlineto{\pgfqpoint{1.851459in}{0.447121in}}%
\pgfpathlineto{\pgfqpoint{1.853729in}{0.359182in}}%
\pgfpathlineto{\pgfqpoint{1.855998in}{0.492525in}}%
\pgfpathlineto{\pgfqpoint{1.858268in}{0.426104in}}%
\pgfpathlineto{\pgfqpoint{1.860537in}{0.478811in}}%
\pgfpathlineto{\pgfqpoint{1.862807in}{0.468034in}}%
\pgfpathlineto{\pgfqpoint{1.865076in}{0.488712in}}%
\pgfpathlineto{\pgfqpoint{1.867346in}{0.496962in}}%
\pgfpathlineto{\pgfqpoint{1.869615in}{0.495321in}}%
\pgfpathlineto{\pgfqpoint{1.871885in}{0.505618in}}%
\pgfpathlineto{\pgfqpoint{1.874154in}{0.504176in}}%
\pgfpathlineto{\pgfqpoint{1.876424in}{0.510276in}}%
\pgfpathlineto{\pgfqpoint{1.878693in}{0.483792in}}%
\pgfpathlineto{\pgfqpoint{1.880963in}{0.529699in}}%
\pgfpathlineto{\pgfqpoint{1.883232in}{0.488902in}}%
\pgfpathlineto{\pgfqpoint{1.885502in}{0.488384in}}%
\pgfpathlineto{\pgfqpoint{1.887771in}{0.482641in}}%
\pgfpathlineto{\pgfqpoint{1.890041in}{0.507733in}}%
\pgfpathlineto{\pgfqpoint{1.892310in}{0.437672in}}%
\pgfpathlineto{\pgfqpoint{1.894579in}{0.389949in}}%
\pgfpathlineto{\pgfqpoint{1.896849in}{0.494909in}}%
\pgfpathlineto{\pgfqpoint{1.899118in}{0.424391in}}%
\pgfpathlineto{\pgfqpoint{1.901388in}{0.456429in}}%
\pgfpathlineto{\pgfqpoint{1.903657in}{0.479959in}}%
\pgfpathlineto{\pgfqpoint{1.905927in}{0.445277in}}%
\pgfpathlineto{\pgfqpoint{1.908196in}{0.496005in}}%
\pgfpathlineto{\pgfqpoint{1.910466in}{0.450807in}}%
\pgfpathlineto{\pgfqpoint{1.912735in}{0.509685in}}%
\pgfpathlineto{\pgfqpoint{1.915005in}{0.407034in}}%
\pgfpathlineto{\pgfqpoint{1.917274in}{0.520077in}}%
\pgfpathlineto{\pgfqpoint{1.919544in}{0.430403in}}%
\pgfpathlineto{\pgfqpoint{1.921813in}{0.534384in}}%
\pgfpathlineto{\pgfqpoint{1.924083in}{0.439710in}}%
\pgfpathlineto{\pgfqpoint{1.926352in}{0.506308in}}%
\pgfpathlineto{\pgfqpoint{1.928622in}{0.495963in}}%
\pgfpathlineto{\pgfqpoint{1.930891in}{0.519601in}}%
\pgfpathlineto{\pgfqpoint{1.933161in}{0.440282in}}%
\pgfpathlineto{\pgfqpoint{1.935430in}{0.410798in}}%
\pgfpathlineto{\pgfqpoint{1.937700in}{0.503662in}}%
\pgfpathlineto{\pgfqpoint{1.939969in}{0.387989in}}%
\pgfpathlineto{\pgfqpoint{1.942239in}{0.430489in}}%
\pgfpathlineto{\pgfqpoint{1.944508in}{0.484258in}}%
\pgfpathlineto{\pgfqpoint{1.946778in}{0.367989in}}%
\pgfpathlineto{\pgfqpoint{1.949047in}{0.482918in}}%
\pgfpathlineto{\pgfqpoint{1.951317in}{0.414344in}}%
\pgfpathlineto{\pgfqpoint{1.953586in}{0.499356in}}%
\pgfpathlineto{\pgfqpoint{1.955856in}{0.480397in}}%
\pgfpathlineto{\pgfqpoint{1.958125in}{0.512652in}}%
\pgfpathlineto{\pgfqpoint{1.960394in}{0.470305in}}%
\pgfpathlineto{\pgfqpoint{1.962664in}{0.527393in}}%
\pgfpathlineto{\pgfqpoint{1.964933in}{0.317702in}}%
\pgfpathlineto{\pgfqpoint{1.967203in}{0.448478in}}%
\pgfpathlineto{\pgfqpoint{1.969472in}{0.504170in}}%
\pgfpathlineto{\pgfqpoint{1.971742in}{0.451870in}}%
\pgfpathlineto{\pgfqpoint{1.974011in}{0.449594in}}%
\pgfpathlineto{\pgfqpoint{1.976281in}{0.467475in}}%
\pgfpathlineto{\pgfqpoint{1.978550in}{0.459665in}}%
\pgfpathlineto{\pgfqpoint{1.980820in}{0.475319in}}%
\pgfpathlineto{\pgfqpoint{1.983089in}{0.426675in}}%
\pgfpathlineto{\pgfqpoint{1.985359in}{0.495132in}}%
\pgfpathlineto{\pgfqpoint{1.987628in}{0.419588in}}%
\pgfpathlineto{\pgfqpoint{1.989898in}{0.512970in}}%
\pgfpathlineto{\pgfqpoint{1.992167in}{0.413702in}}%
\pgfpathlineto{\pgfqpoint{1.994437in}{0.525721in}}%
\pgfpathlineto{\pgfqpoint{1.996706in}{0.463222in}}%
\pgfpathlineto{\pgfqpoint{1.998976in}{0.471079in}}%
\pgfpathlineto{\pgfqpoint{2.001245in}{0.487984in}}%
\pgfpathlineto{\pgfqpoint{2.003515in}{0.484229in}}%
\pgfpathlineto{\pgfqpoint{2.008054in}{0.423863in}}%
\pgfpathlineto{\pgfqpoint{2.010323in}{0.484687in}}%
\pgfpathlineto{\pgfqpoint{2.012593in}{0.458439in}}%
\pgfpathlineto{\pgfqpoint{2.014862in}{0.468187in}}%
\pgfpathlineto{\pgfqpoint{2.017132in}{0.482942in}}%
\pgfpathlineto{\pgfqpoint{2.019401in}{0.465828in}}%
\pgfpathlineto{\pgfqpoint{2.021671in}{0.507393in}}%
\pgfpathlineto{\pgfqpoint{2.023940in}{0.485846in}}%
\pgfpathlineto{\pgfqpoint{2.026209in}{0.515977in}}%
\pgfpathlineto{\pgfqpoint{2.028479in}{0.499371in}}%
\pgfpathlineto{\pgfqpoint{2.033018in}{0.448491in}}%
\pgfpathlineto{\pgfqpoint{2.035287in}{0.502485in}}%
\pgfpathlineto{\pgfqpoint{2.037557in}{0.423301in}}%
\pgfpathlineto{\pgfqpoint{2.039826in}{0.410438in}}%
\pgfpathlineto{\pgfqpoint{2.042096in}{0.489135in}}%
\pgfpathlineto{\pgfqpoint{2.044365in}{0.383110in}}%
\pgfpathlineto{\pgfqpoint{2.046635in}{0.479895in}}%
\pgfpathlineto{\pgfqpoint{2.048904in}{0.450976in}}%
\pgfpathlineto{\pgfqpoint{2.051174in}{0.493732in}}%
\pgfpathlineto{\pgfqpoint{2.053443in}{0.491601in}}%
\pgfpathlineto{\pgfqpoint{2.055713in}{0.504060in}}%
\pgfpathlineto{\pgfqpoint{2.057982in}{0.494047in}}%
\pgfpathlineto{\pgfqpoint{2.060252in}{0.517275in}}%
\pgfpathlineto{\pgfqpoint{2.062521in}{0.451505in}}%
\pgfpathlineto{\pgfqpoint{2.064791in}{0.415890in}}%
\pgfpathlineto{\pgfqpoint{2.067060in}{0.507555in}}%
\pgfpathlineto{\pgfqpoint{2.069330in}{0.411032in}}%
\pgfpathlineto{\pgfqpoint{2.071599in}{0.451818in}}%
\pgfpathlineto{\pgfqpoint{2.073869in}{0.478427in}}%
\pgfpathlineto{\pgfqpoint{2.076138in}{0.450201in}}%
\pgfpathlineto{\pgfqpoint{2.078408in}{0.480147in}}%
\pgfpathlineto{\pgfqpoint{2.080677in}{0.384832in}}%
\pgfpathlineto{\pgfqpoint{2.082947in}{0.499465in}}%
\pgfpathlineto{\pgfqpoint{2.085216in}{0.450305in}}%
\pgfpathlineto{\pgfqpoint{2.087486in}{0.510174in}}%
\pgfpathlineto{\pgfqpoint{2.089755in}{0.434958in}}%
\pgfpathlineto{\pgfqpoint{2.092024in}{0.524812in}}%
\pgfpathlineto{\pgfqpoint{2.094294in}{0.384014in}}%
\pgfpathlineto{\pgfqpoint{2.096563in}{0.466039in}}%
\pgfpathlineto{\pgfqpoint{2.098833in}{0.498511in}}%
\pgfpathlineto{\pgfqpoint{2.101102in}{0.470458in}}%
\pgfpathlineto{\pgfqpoint{2.105641in}{0.449421in}}%
\pgfpathlineto{\pgfqpoint{2.107911in}{0.481984in}}%
\pgfpathlineto{\pgfqpoint{2.110180in}{0.468568in}}%
\pgfpathlineto{\pgfqpoint{2.112450in}{0.460000in}}%
\pgfpathlineto{\pgfqpoint{2.114719in}{0.491851in}}%
\pgfpathlineto{\pgfqpoint{2.116989in}{0.441734in}}%
\pgfpathlineto{\pgfqpoint{2.119258in}{0.508244in}}%
\pgfpathlineto{\pgfqpoint{2.121528in}{0.460206in}}%
\pgfpathlineto{\pgfqpoint{2.123797in}{0.520891in}}%
\pgfpathlineto{\pgfqpoint{2.126067in}{0.450094in}}%
\pgfpathlineto{\pgfqpoint{2.128336in}{0.517871in}}%
\pgfpathlineto{\pgfqpoint{2.130606in}{0.408791in}}%
\pgfpathlineto{\pgfqpoint{2.132875in}{0.535164in}}%
\pgfpathlineto{\pgfqpoint{2.135145in}{0.381477in}}%
\pgfpathlineto{\pgfqpoint{2.137414in}{0.478089in}}%
\pgfpathlineto{\pgfqpoint{2.139684in}{0.509537in}}%
\pgfpathlineto{\pgfqpoint{2.141953in}{0.484957in}}%
\pgfpathlineto{\pgfqpoint{2.144223in}{0.454074in}}%
\pgfpathlineto{\pgfqpoint{2.146492in}{0.465092in}}%
\pgfpathlineto{\pgfqpoint{2.148762in}{0.479187in}}%
\pgfpathlineto{\pgfqpoint{2.151031in}{0.459809in}}%
\pgfpathlineto{\pgfqpoint{2.153301in}{0.405992in}}%
\pgfpathlineto{\pgfqpoint{2.155570in}{0.490595in}}%
\pgfpathlineto{\pgfqpoint{2.157840in}{0.430260in}}%
\pgfpathlineto{\pgfqpoint{2.160109in}{0.497843in}}%
\pgfpathlineto{\pgfqpoint{2.162378in}{0.431250in}}%
\pgfpathlineto{\pgfqpoint{2.164648in}{0.512791in}}%
\pgfpathlineto{\pgfqpoint{2.166917in}{0.481170in}}%
\pgfpathlineto{\pgfqpoint{2.169187in}{0.518765in}}%
\pgfpathlineto{\pgfqpoint{2.171456in}{0.483920in}}%
\pgfpathlineto{\pgfqpoint{2.173726in}{0.533079in}}%
\pgfpathlineto{\pgfqpoint{2.175995in}{0.472675in}}%
\pgfpathlineto{\pgfqpoint{2.178265in}{0.492945in}}%
\pgfpathlineto{\pgfqpoint{2.180534in}{0.521573in}}%
\pgfpathlineto{\pgfqpoint{2.182804in}{0.498556in}}%
\pgfpathlineto{\pgfqpoint{2.185073in}{0.459923in}}%
\pgfpathlineto{\pgfqpoint{2.189612in}{0.488521in}}%
\pgfpathlineto{\pgfqpoint{2.191882in}{0.445153in}}%
\pgfpathlineto{\pgfqpoint{2.194151in}{0.383747in}}%
\pgfpathlineto{\pgfqpoint{2.196421in}{0.486069in}}%
\pgfpathlineto{\pgfqpoint{2.198690in}{0.449666in}}%
\pgfpathlineto{\pgfqpoint{2.200960in}{0.476227in}}%
\pgfpathlineto{\pgfqpoint{2.203229in}{0.473197in}}%
\pgfpathlineto{\pgfqpoint{2.205499in}{0.483692in}}%
\pgfpathlineto{\pgfqpoint{2.207768in}{0.504937in}}%
\pgfpathlineto{\pgfqpoint{2.210038in}{0.501299in}}%
\pgfpathlineto{\pgfqpoint{2.214577in}{0.514835in}}%
\pgfpathlineto{\pgfqpoint{2.219116in}{0.400542in}}%
\pgfpathlineto{\pgfqpoint{2.221385in}{0.506572in}}%
\pgfpathlineto{\pgfqpoint{2.223655in}{0.326626in}}%
\pgfpathlineto{\pgfqpoint{2.225924in}{0.424201in}}%
\pgfpathlineto{\pgfqpoint{2.228193in}{0.483942in}}%
\pgfpathlineto{\pgfqpoint{2.230463in}{0.351518in}}%
\pgfpathlineto{\pgfqpoint{2.232732in}{0.478665in}}%
\pgfpathlineto{\pgfqpoint{2.235002in}{0.428340in}}%
\pgfpathlineto{\pgfqpoint{2.237271in}{0.494764in}}%
\pgfpathlineto{\pgfqpoint{2.239541in}{0.485603in}}%
\pgfpathlineto{\pgfqpoint{2.241810in}{0.509958in}}%
\pgfpathlineto{\pgfqpoint{2.244080in}{0.476739in}}%
\pgfpathlineto{\pgfqpoint{2.246349in}{0.524777in}}%
\pgfpathlineto{\pgfqpoint{2.248619in}{0.397947in}}%
\pgfpathlineto{\pgfqpoint{2.253158in}{0.504368in}}%
\pgfpathlineto{\pgfqpoint{2.255427in}{0.447375in}}%
\pgfpathlineto{\pgfqpoint{2.257697in}{0.451913in}}%
\pgfpathlineto{\pgfqpoint{2.259966in}{0.468438in}}%
\pgfpathlineto{\pgfqpoint{2.262236in}{0.460906in}}%
\pgfpathlineto{\pgfqpoint{2.264505in}{0.476524in}}%
\pgfpathlineto{\pgfqpoint{2.266775in}{0.429084in}}%
\pgfpathlineto{\pgfqpoint{2.269044in}{0.496238in}}%
\pgfpathlineto{\pgfqpoint{2.271314in}{0.413072in}}%
\pgfpathlineto{\pgfqpoint{2.273583in}{0.512674in}}%
\pgfpathlineto{\pgfqpoint{2.275853in}{0.407534in}}%
\pgfpathlineto{\pgfqpoint{2.278122in}{0.525775in}}%
\pgfpathlineto{\pgfqpoint{2.280392in}{0.459878in}}%
\pgfpathlineto{\pgfqpoint{2.282661in}{0.471209in}}%
\pgfpathlineto{\pgfqpoint{2.284931in}{0.489065in}}%
\pgfpathlineto{\pgfqpoint{2.287200in}{0.483555in}}%
\pgfpathlineto{\pgfqpoint{2.291739in}{0.426008in}}%
\pgfpathlineto{\pgfqpoint{2.294008in}{0.484945in}}%
\pgfpathlineto{\pgfqpoint{2.296278in}{0.459051in}}%
\pgfpathlineto{\pgfqpoint{2.298547in}{0.467732in}}%
\pgfpathlineto{\pgfqpoint{2.300817in}{0.483793in}}%
\pgfpathlineto{\pgfqpoint{2.303086in}{0.481769in}}%
\pgfpathlineto{\pgfqpoint{2.303086in}{0.481769in}}%
\pgfusepath{stroke}%
\end{pgfscope}%
\begin{pgfscope}%
\pgfpathrectangle{\pgfqpoint{0.487500in}{0.320625in}}{\pgfqpoint{1.815586in}{1.164375in}}%
\pgfusepath{clip}%
\pgfsetbuttcap%
\pgfsetroundjoin%
\pgfsetlinewidth{0.752812pt}%
\definecolor{currentstroke}{rgb}{0.819608,0.286275,0.356863}%
\pgfsetstrokecolor{currentstroke}%
\pgfsetstrokeopacity{0.650000}%
\pgfsetdash{{2.775000pt}{1.200000pt}}{0.000000pt}%
\pgfpathmoveto{\pgfqpoint{0.487500in}{1.451682in}}%
\pgfpathlineto{\pgfqpoint{1.711161in}{0.315625in}}%
\pgfpathlineto{\pgfqpoint{1.711161in}{0.315625in}}%
\pgfusepath{stroke}%
\end{pgfscope}%
\begin{pgfscope}%
\pgfpathrectangle{\pgfqpoint{0.487500in}{0.320625in}}{\pgfqpoint{1.815586in}{1.164375in}}%
\pgfusepath{clip}%
\pgfsetbuttcap%
\pgfsetroundjoin%
\pgfsetlinewidth{0.803000pt}%
\definecolor{currentstroke}{rgb}{0.819608,0.286275,0.356863}%
\pgfsetstrokecolor{currentstroke}%
\pgfsetstrokeopacity{0.800000}%
\pgfsetdash{{0.800000pt}{1.320000pt}}{0.000000pt}%
\pgfpathmoveto{\pgfqpoint{0.487500in}{1.451682in}}%
\pgfpathlineto{\pgfqpoint{2.303086in}{1.451682in}}%
\pgfpathlineto{\pgfqpoint{2.303086in}{1.451682in}}%
\pgfusepath{stroke}%
\end{pgfscope}%
\begin{pgfscope}%
\pgfpathrectangle{\pgfqpoint{0.487500in}{0.320625in}}{\pgfqpoint{1.815586in}{1.164375in}}%
\pgfusepath{clip}%
\pgfsetrectcap%
\pgfsetroundjoin%
\pgfsetlinewidth{1.254687pt}%
\definecolor{currentstroke}{rgb}{0.000000,0.529412,0.423529}%
\pgfsetstrokecolor{currentstroke}%
\pgfsetdash{}{0pt}%
\pgfpathmoveto{\pgfqpoint{0.487500in}{1.451682in}}%
\pgfpathlineto{\pgfqpoint{0.492039in}{1.441674in}}%
\pgfpathlineto{\pgfqpoint{0.507925in}{1.409832in}}%
\pgfpathlineto{\pgfqpoint{0.523812in}{1.375075in}}%
\pgfpathlineto{\pgfqpoint{0.555584in}{1.301926in}}%
\pgfpathlineto{\pgfqpoint{0.600974in}{1.198168in}}%
\pgfpathlineto{\pgfqpoint{0.641825in}{1.107600in}}%
\pgfpathlineto{\pgfqpoint{0.687215in}{1.009674in}}%
\pgfpathlineto{\pgfqpoint{0.741682in}{0.894916in}}%
\pgfpathlineto{\pgfqpoint{0.807497in}{0.758978in}}%
\pgfpathlineto{\pgfqpoint{0.882390in}{0.606770in}}%
\pgfpathlineto{\pgfqpoint{0.886929in}{0.597610in}}%
\pgfpathlineto{\pgfqpoint{0.907354in}{0.557403in}}%
\pgfpathlineto{\pgfqpoint{0.909624in}{0.554293in}}%
\pgfpathlineto{\pgfqpoint{0.911893in}{0.549289in}}%
\pgfpathlineto{\pgfqpoint{0.914163in}{0.542665in}}%
\pgfpathlineto{\pgfqpoint{0.916432in}{0.540098in}}%
\pgfpathlineto{\pgfqpoint{0.918702in}{0.545433in}}%
\pgfpathlineto{\pgfqpoint{0.920971in}{0.530673in}}%
\pgfpathlineto{\pgfqpoint{0.923241in}{0.529633in}}%
\pgfpathlineto{\pgfqpoint{0.925510in}{0.519774in}}%
\pgfpathlineto{\pgfqpoint{0.927780in}{0.515249in}}%
\pgfpathlineto{\pgfqpoint{0.930049in}{0.513184in}}%
\pgfpathlineto{\pgfqpoint{0.932319in}{0.510191in}}%
\pgfpathlineto{\pgfqpoint{0.934588in}{0.513565in}}%
\pgfpathlineto{\pgfqpoint{0.936858in}{0.506455in}}%
\pgfpathlineto{\pgfqpoint{0.939127in}{0.492967in}}%
\pgfpathlineto{\pgfqpoint{0.941397in}{0.514557in}}%
\pgfpathlineto{\pgfqpoint{0.943666in}{0.518914in}}%
\pgfpathlineto{\pgfqpoint{0.945936in}{0.512938in}}%
\pgfpathlineto{\pgfqpoint{0.948205in}{0.480425in}}%
\pgfpathlineto{\pgfqpoint{0.950475in}{0.480844in}}%
\pgfpathlineto{\pgfqpoint{0.952744in}{0.517839in}}%
\pgfpathlineto{\pgfqpoint{0.955014in}{0.477202in}}%
\pgfpathlineto{\pgfqpoint{0.957283in}{0.502021in}}%
\pgfpathlineto{\pgfqpoint{0.959552in}{0.464112in}}%
\pgfpathlineto{\pgfqpoint{0.961822in}{0.461595in}}%
\pgfpathlineto{\pgfqpoint{0.964091in}{0.447965in}}%
\pgfpathlineto{\pgfqpoint{0.966361in}{0.485009in}}%
\pgfpathlineto{\pgfqpoint{0.968630in}{0.506115in}}%
\pgfpathlineto{\pgfqpoint{0.970900in}{0.497999in}}%
\pgfpathlineto{\pgfqpoint{0.973169in}{0.449318in}}%
\pgfpathlineto{\pgfqpoint{0.975439in}{0.521796in}}%
\pgfpathlineto{\pgfqpoint{0.977708in}{0.511324in}}%
\pgfpathlineto{\pgfqpoint{0.979978in}{0.515475in}}%
\pgfpathlineto{\pgfqpoint{0.982247in}{0.441602in}}%
\pgfpathlineto{\pgfqpoint{0.984517in}{0.415310in}}%
\pgfpathlineto{\pgfqpoint{0.986786in}{0.487226in}}%
\pgfpathlineto{\pgfqpoint{0.989056in}{0.448662in}}%
\pgfpathlineto{\pgfqpoint{0.991325in}{0.490994in}}%
\pgfpathlineto{\pgfqpoint{0.993595in}{0.474137in}}%
\pgfpathlineto{\pgfqpoint{0.995864in}{0.465675in}}%
\pgfpathlineto{\pgfqpoint{0.998134in}{0.383355in}}%
\pgfpathlineto{\pgfqpoint{1.000403in}{0.498139in}}%
\pgfpathlineto{\pgfqpoint{1.002673in}{0.517837in}}%
\pgfpathlineto{\pgfqpoint{1.004942in}{0.443897in}}%
\pgfpathlineto{\pgfqpoint{1.007212in}{0.449767in}}%
\pgfpathlineto{\pgfqpoint{1.009481in}{0.468412in}}%
\pgfpathlineto{\pgfqpoint{1.011751in}{0.451177in}}%
\pgfpathlineto{\pgfqpoint{1.014020in}{0.493635in}}%
\pgfpathlineto{\pgfqpoint{1.018559in}{0.471455in}}%
\pgfpathlineto{\pgfqpoint{1.020829in}{0.364063in}}%
\pgfpathlineto{\pgfqpoint{1.023098in}{0.498001in}}%
\pgfpathlineto{\pgfqpoint{1.025367in}{0.520610in}}%
\pgfpathlineto{\pgfqpoint{1.027637in}{0.493063in}}%
\pgfpathlineto{\pgfqpoint{1.029906in}{0.376493in}}%
\pgfpathlineto{\pgfqpoint{1.032176in}{0.512622in}}%
\pgfpathlineto{\pgfqpoint{1.034445in}{0.484924in}}%
\pgfpathlineto{\pgfqpoint{1.036715in}{0.516157in}}%
\pgfpathlineto{\pgfqpoint{1.038984in}{0.447659in}}%
\pgfpathlineto{\pgfqpoint{1.041254in}{0.475450in}}%
\pgfpathlineto{\pgfqpoint{1.043278in}{0.315625in}}%
\pgfpathmoveto{\pgfqpoint{1.043849in}{0.315625in}}%
\pgfpathlineto{\pgfqpoint{1.045793in}{0.431352in}}%
\pgfpathlineto{\pgfqpoint{1.048062in}{0.487786in}}%
\pgfpathlineto{\pgfqpoint{1.050332in}{0.497596in}}%
\pgfpathlineto{\pgfqpoint{1.052601in}{0.483086in}}%
\pgfpathlineto{\pgfqpoint{1.054871in}{0.412617in}}%
\pgfpathlineto{\pgfqpoint{1.057140in}{0.507383in}}%
\pgfpathlineto{\pgfqpoint{1.059410in}{0.529415in}}%
\pgfpathlineto{\pgfqpoint{1.061679in}{0.480906in}}%
\pgfpathlineto{\pgfqpoint{1.063949in}{0.458927in}}%
\pgfpathlineto{\pgfqpoint{1.066218in}{0.495801in}}%
\pgfpathlineto{\pgfqpoint{1.068488in}{0.449052in}}%
\pgfpathlineto{\pgfqpoint{1.070757in}{0.500495in}}%
\pgfpathlineto{\pgfqpoint{1.073027in}{0.461278in}}%
\pgfpathlineto{\pgfqpoint{1.075296in}{0.483924in}}%
\pgfpathlineto{\pgfqpoint{1.079835in}{0.421816in}}%
\pgfpathlineto{\pgfqpoint{1.082105in}{0.493734in}}%
\pgfpathlineto{\pgfqpoint{1.084374in}{0.511618in}}%
\pgfpathlineto{\pgfqpoint{1.088913in}{0.469847in}}%
\pgfpathlineto{\pgfqpoint{1.091182in}{0.507515in}}%
\pgfpathlineto{\pgfqpoint{1.093452in}{0.533291in}}%
\pgfpathlineto{\pgfqpoint{1.095721in}{0.463146in}}%
\pgfpathlineto{\pgfqpoint{1.097991in}{0.488397in}}%
\pgfpathlineto{\pgfqpoint{1.100260in}{0.449930in}}%
\pgfpathlineto{\pgfqpoint{1.102488in}{0.315625in}}%
\pgfpathmoveto{\pgfqpoint{1.102567in}{0.315625in}}%
\pgfpathlineto{\pgfqpoint{1.104799in}{0.470169in}}%
\pgfpathlineto{\pgfqpoint{1.107069in}{0.468522in}}%
\pgfpathlineto{\pgfqpoint{1.109338in}{0.489815in}}%
\pgfpathlineto{\pgfqpoint{1.111608in}{0.483469in}}%
\pgfpathlineto{\pgfqpoint{1.113877in}{0.442182in}}%
\pgfpathlineto{\pgfqpoint{1.116147in}{0.511362in}}%
\pgfpathlineto{\pgfqpoint{1.118416in}{0.481858in}}%
\pgfpathlineto{\pgfqpoint{1.120686in}{0.497013in}}%
\pgfpathlineto{\pgfqpoint{1.125225in}{0.329121in}}%
\pgfpathlineto{\pgfqpoint{1.127494in}{0.476750in}}%
\pgfpathlineto{\pgfqpoint{1.129764in}{0.486432in}}%
\pgfpathlineto{\pgfqpoint{1.132033in}{0.500367in}}%
\pgfpathlineto{\pgfqpoint{1.136572in}{0.450989in}}%
\pgfpathlineto{\pgfqpoint{1.138842in}{0.514513in}}%
\pgfpathlineto{\pgfqpoint{1.141111in}{0.505202in}}%
\pgfpathlineto{\pgfqpoint{1.143381in}{0.505825in}}%
\pgfpathlineto{\pgfqpoint{1.145650in}{0.422213in}}%
\pgfpathlineto{\pgfqpoint{1.147920in}{0.427219in}}%
\pgfpathlineto{\pgfqpoint{1.150189in}{0.501786in}}%
\pgfpathlineto{\pgfqpoint{1.152459in}{0.470280in}}%
\pgfpathlineto{\pgfqpoint{1.154728in}{0.498776in}}%
\pgfpathlineto{\pgfqpoint{1.159267in}{0.434018in}}%
\pgfpathlineto{\pgfqpoint{1.161536in}{0.458293in}}%
\pgfpathlineto{\pgfqpoint{1.163806in}{0.494716in}}%
\pgfpathlineto{\pgfqpoint{1.166075in}{0.509756in}}%
\pgfpathlineto{\pgfqpoint{1.170614in}{0.446371in}}%
\pgfpathlineto{\pgfqpoint{1.172884in}{0.518370in}}%
\pgfpathlineto{\pgfqpoint{1.175153in}{0.497883in}}%
\pgfpathlineto{\pgfqpoint{1.177423in}{0.512462in}}%
\pgfpathlineto{\pgfqpoint{1.179692in}{0.386121in}}%
\pgfpathlineto{\pgfqpoint{1.184231in}{0.470405in}}%
\pgfpathlineto{\pgfqpoint{1.186501in}{0.460776in}}%
\pgfpathlineto{\pgfqpoint{1.188770in}{0.496501in}}%
\pgfpathlineto{\pgfqpoint{1.191040in}{0.483722in}}%
\pgfpathlineto{\pgfqpoint{1.193309in}{0.457004in}}%
\pgfpathlineto{\pgfqpoint{1.195579in}{0.452742in}}%
\pgfpathlineto{\pgfqpoint{1.197848in}{0.504066in}}%
\pgfpathlineto{\pgfqpoint{1.200118in}{0.519224in}}%
\pgfpathlineto{\pgfqpoint{1.202387in}{0.476206in}}%
\pgfpathlineto{\pgfqpoint{1.204657in}{0.402592in}}%
\pgfpathlineto{\pgfqpoint{1.206926in}{0.513438in}}%
\pgfpathlineto{\pgfqpoint{1.209196in}{0.481144in}}%
\pgfpathlineto{\pgfqpoint{1.211465in}{0.511923in}}%
\pgfpathlineto{\pgfqpoint{1.213735in}{0.448428in}}%
\pgfpathlineto{\pgfqpoint{1.216004in}{0.468236in}}%
\pgfpathlineto{\pgfqpoint{1.218274in}{0.379116in}}%
\pgfpathlineto{\pgfqpoint{1.220543in}{0.452970in}}%
\pgfpathlineto{\pgfqpoint{1.222813in}{0.496522in}}%
\pgfpathlineto{\pgfqpoint{1.225082in}{0.499054in}}%
\pgfpathlineto{\pgfqpoint{1.227351in}{0.469093in}}%
\pgfpathlineto{\pgfqpoint{1.229621in}{0.470763in}}%
\pgfpathlineto{\pgfqpoint{1.231890in}{0.509760in}}%
\pgfpathlineto{\pgfqpoint{1.234160in}{0.526613in}}%
\pgfpathlineto{\pgfqpoint{1.236429in}{0.462654in}}%
\pgfpathlineto{\pgfqpoint{1.238699in}{0.445726in}}%
\pgfpathlineto{\pgfqpoint{1.240968in}{0.493284in}}%
\pgfpathlineto{\pgfqpoint{1.243238in}{0.451183in}}%
\pgfpathlineto{\pgfqpoint{1.245507in}{0.499125in}}%
\pgfpathlineto{\pgfqpoint{1.247777in}{0.465999in}}%
\pgfpathlineto{\pgfqpoint{1.250046in}{0.480607in}}%
\pgfpathlineto{\pgfqpoint{1.252316in}{0.446379in}}%
\pgfpathlineto{\pgfqpoint{1.254585in}{0.455437in}}%
\pgfpathlineto{\pgfqpoint{1.256855in}{0.504864in}}%
\pgfpathlineto{\pgfqpoint{1.259124in}{0.510699in}}%
\pgfpathlineto{\pgfqpoint{1.261394in}{0.475726in}}%
\pgfpathlineto{\pgfqpoint{1.263663in}{0.492780in}}%
\pgfpathlineto{\pgfqpoint{1.265933in}{0.505418in}}%
\pgfpathlineto{\pgfqpoint{1.268202in}{0.529277in}}%
\pgfpathlineto{\pgfqpoint{1.270472in}{0.452325in}}%
\pgfpathlineto{\pgfqpoint{1.272741in}{0.482378in}}%
\pgfpathlineto{\pgfqpoint{1.275011in}{0.443004in}}%
\pgfpathlineto{\pgfqpoint{1.277280in}{0.377016in}}%
\pgfpathlineto{\pgfqpoint{1.279550in}{0.474892in}}%
\pgfpathlineto{\pgfqpoint{1.281819in}{0.476639in}}%
\pgfpathlineto{\pgfqpoint{1.284089in}{0.489589in}}%
\pgfpathlineto{\pgfqpoint{1.286358in}{0.475194in}}%
\pgfpathlineto{\pgfqpoint{1.288628in}{0.472062in}}%
\pgfpathlineto{\pgfqpoint{1.290897in}{0.519585in}}%
\pgfpathlineto{\pgfqpoint{1.293166in}{0.514726in}}%
\pgfpathlineto{\pgfqpoint{1.295436in}{0.482252in}}%
\pgfpathlineto{\pgfqpoint{1.297705in}{0.501713in}}%
\pgfpathlineto{\pgfqpoint{1.299975in}{0.483215in}}%
\pgfpathlineto{\pgfqpoint{1.302244in}{0.522112in}}%
\pgfpathlineto{\pgfqpoint{1.304514in}{0.450945in}}%
\pgfpathlineto{\pgfqpoint{1.306783in}{0.490999in}}%
\pgfpathlineto{\pgfqpoint{1.309053in}{0.426167in}}%
\pgfpathlineto{\pgfqpoint{1.311322in}{0.411829in}}%
\pgfpathlineto{\pgfqpoint{1.315861in}{0.490661in}}%
\pgfpathlineto{\pgfqpoint{1.318131in}{0.500646in}}%
\pgfpathlineto{\pgfqpoint{1.320400in}{0.483072in}}%
\pgfpathlineto{\pgfqpoint{1.322670in}{0.492405in}}%
\pgfpathlineto{\pgfqpoint{1.324939in}{0.533387in}}%
\pgfpathlineto{\pgfqpoint{1.327209in}{0.506264in}}%
\pgfpathlineto{\pgfqpoint{1.329478in}{0.495159in}}%
\pgfpathlineto{\pgfqpoint{1.331748in}{0.489335in}}%
\pgfpathlineto{\pgfqpoint{1.334017in}{0.410257in}}%
\pgfpathlineto{\pgfqpoint{1.336287in}{0.495408in}}%
\pgfpathlineto{\pgfqpoint{1.338556in}{0.447773in}}%
\pgfpathlineto{\pgfqpoint{1.340826in}{0.487331in}}%
\pgfpathlineto{\pgfqpoint{1.343095in}{0.469388in}}%
\pgfpathlineto{\pgfqpoint{1.345365in}{0.439132in}}%
\pgfpathlineto{\pgfqpoint{1.347634in}{0.437790in}}%
\pgfpathlineto{\pgfqpoint{1.349904in}{0.509118in}}%
\pgfpathlineto{\pgfqpoint{1.352173in}{0.512177in}}%
\pgfpathlineto{\pgfqpoint{1.354443in}{0.429045in}}%
\pgfpathlineto{\pgfqpoint{1.356088in}{0.315625in}}%
\pgfpathmoveto{\pgfqpoint{1.357167in}{0.315625in}}%
\pgfpathlineto{\pgfqpoint{1.358981in}{0.486859in}}%
\pgfpathlineto{\pgfqpoint{1.361251in}{0.452162in}}%
\pgfpathlineto{\pgfqpoint{1.363520in}{0.489935in}}%
\pgfpathlineto{\pgfqpoint{1.365790in}{0.475359in}}%
\pgfpathlineto{\pgfqpoint{1.368059in}{0.447320in}}%
\pgfpathlineto{\pgfqpoint{1.370329in}{0.440169in}}%
\pgfpathlineto{\pgfqpoint{1.372598in}{0.508352in}}%
\pgfpathlineto{\pgfqpoint{1.374868in}{0.515675in}}%
\pgfpathlineto{\pgfqpoint{1.377137in}{0.388108in}}%
\pgfpathlineto{\pgfqpoint{1.379407in}{0.404440in}}%
\pgfpathlineto{\pgfqpoint{1.381676in}{0.482009in}}%
\pgfpathlineto{\pgfqpoint{1.383946in}{0.449136in}}%
\pgfpathlineto{\pgfqpoint{1.386215in}{0.490091in}}%
\pgfpathlineto{\pgfqpoint{1.388485in}{0.477366in}}%
\pgfpathlineto{\pgfqpoint{1.390754in}{0.459380in}}%
\pgfpathlineto{\pgfqpoint{1.393024in}{0.405455in}}%
\pgfpathlineto{\pgfqpoint{1.395293in}{0.505191in}}%
\pgfpathlineto{\pgfqpoint{1.397563in}{0.518535in}}%
\pgfpathlineto{\pgfqpoint{1.399832in}{0.394921in}}%
\pgfpathlineto{\pgfqpoint{1.404371in}{0.474630in}}%
\pgfpathlineto{\pgfqpoint{1.406641in}{0.445589in}}%
\pgfpathlineto{\pgfqpoint{1.408910in}{0.489748in}}%
\pgfpathlineto{\pgfqpoint{1.413449in}{0.468438in}}%
\pgfpathlineto{\pgfqpoint{1.415291in}{0.315625in}}%
\pgfpathmoveto{\pgfqpoint{1.416082in}{0.315625in}}%
\pgfpathlineto{\pgfqpoint{1.417988in}{0.501752in}}%
\pgfpathlineto{\pgfqpoint{1.420258in}{0.520700in}}%
\pgfpathlineto{\pgfqpoint{1.422527in}{0.492319in}}%
\pgfpathlineto{\pgfqpoint{1.424796in}{0.436582in}}%
\pgfpathlineto{\pgfqpoint{1.427066in}{0.524113in}}%
\pgfpathlineto{\pgfqpoint{1.429335in}{0.487435in}}%
\pgfpathlineto{\pgfqpoint{1.431605in}{0.516738in}}%
\pgfpathlineto{\pgfqpoint{1.433874in}{0.422404in}}%
\pgfpathlineto{\pgfqpoint{1.436144in}{0.470964in}}%
\pgfpathlineto{\pgfqpoint{1.438413in}{0.361132in}}%
\pgfpathlineto{\pgfqpoint{1.442952in}{0.482843in}}%
\pgfpathlineto{\pgfqpoint{1.445222in}{0.497370in}}%
\pgfpathlineto{\pgfqpoint{1.447491in}{0.479097in}}%
\pgfpathlineto{\pgfqpoint{1.449761in}{0.424052in}}%
\pgfpathlineto{\pgfqpoint{1.452030in}{0.515043in}}%
\pgfpathlineto{\pgfqpoint{1.454300in}{0.530552in}}%
\pgfpathlineto{\pgfqpoint{1.458839in}{0.405594in}}%
\pgfpathlineto{\pgfqpoint{1.461108in}{0.510108in}}%
\pgfpathlineto{\pgfqpoint{1.463378in}{0.450340in}}%
\pgfpathlineto{\pgfqpoint{1.465647in}{0.502336in}}%
\pgfpathlineto{\pgfqpoint{1.467917in}{0.452180in}}%
\pgfpathlineto{\pgfqpoint{1.470186in}{0.479420in}}%
\pgfpathlineto{\pgfqpoint{1.474725in}{0.420839in}}%
\pgfpathlineto{\pgfqpoint{1.476995in}{0.490968in}}%
\pgfpathlineto{\pgfqpoint{1.479264in}{0.514107in}}%
\pgfpathlineto{\pgfqpoint{1.481534in}{0.483445in}}%
\pgfpathlineto{\pgfqpoint{1.483803in}{0.498061in}}%
\pgfpathlineto{\pgfqpoint{1.486073in}{0.454437in}}%
\pgfpathlineto{\pgfqpoint{1.488342in}{0.503757in}}%
\pgfpathlineto{\pgfqpoint{1.490611in}{0.449152in}}%
\pgfpathlineto{\pgfqpoint{1.492881in}{0.482014in}}%
\pgfpathlineto{\pgfqpoint{1.495150in}{0.456639in}}%
\pgfpathlineto{\pgfqpoint{1.497420in}{0.352027in}}%
\pgfpathlineto{\pgfqpoint{1.499689in}{0.479110in}}%
\pgfpathlineto{\pgfqpoint{1.501959in}{0.512875in}}%
\pgfpathlineto{\pgfqpoint{1.504228in}{0.494537in}}%
\pgfpathlineto{\pgfqpoint{1.506498in}{0.484681in}}%
\pgfpathlineto{\pgfqpoint{1.508767in}{0.445469in}}%
\pgfpathlineto{\pgfqpoint{1.511037in}{0.500342in}}%
\pgfpathlineto{\pgfqpoint{1.513306in}{0.453505in}}%
\pgfpathlineto{\pgfqpoint{1.515576in}{0.485260in}}%
\pgfpathlineto{\pgfqpoint{1.517845in}{0.462417in}}%
\pgfpathlineto{\pgfqpoint{1.520115in}{0.345076in}}%
\pgfpathlineto{\pgfqpoint{1.522384in}{0.476001in}}%
\pgfpathlineto{\pgfqpoint{1.524654in}{0.512163in}}%
\pgfpathlineto{\pgfqpoint{1.526923in}{0.501012in}}%
\pgfpathlineto{\pgfqpoint{1.529193in}{0.473042in}}%
\pgfpathlineto{\pgfqpoint{1.531462in}{0.432505in}}%
\pgfpathlineto{\pgfqpoint{1.533732in}{0.497249in}}%
\pgfpathlineto{\pgfqpoint{1.536001in}{0.454896in}}%
\pgfpathlineto{\pgfqpoint{1.538271in}{0.487780in}}%
\pgfpathlineto{\pgfqpoint{1.540540in}{0.466962in}}%
\pgfpathlineto{\pgfqpoint{1.542810in}{0.404475in}}%
\pgfpathlineto{\pgfqpoint{1.545079in}{0.468359in}}%
\pgfpathlineto{\pgfqpoint{1.547349in}{0.510864in}}%
\pgfpathlineto{\pgfqpoint{1.549618in}{0.506695in}}%
\pgfpathlineto{\pgfqpoint{1.554157in}{0.407958in}}%
\pgfpathlineto{\pgfqpoint{1.556427in}{0.493193in}}%
\pgfpathlineto{\pgfqpoint{1.558696in}{0.455450in}}%
\pgfpathlineto{\pgfqpoint{1.560965in}{0.489820in}}%
\pgfpathlineto{\pgfqpoint{1.563235in}{0.471094in}}%
\pgfpathlineto{\pgfqpoint{1.565504in}{0.429166in}}%
\pgfpathlineto{\pgfqpoint{1.567774in}{0.458598in}}%
\pgfpathlineto{\pgfqpoint{1.570043in}{0.509210in}}%
\pgfpathlineto{\pgfqpoint{1.572313in}{0.511327in}}%
\pgfpathlineto{\pgfqpoint{1.574582in}{0.428516in}}%
\pgfpathlineto{\pgfqpoint{1.576762in}{0.315625in}}%
\pgfpathmoveto{\pgfqpoint{1.576911in}{0.315625in}}%
\pgfpathlineto{\pgfqpoint{1.579121in}{0.488297in}}%
\pgfpathlineto{\pgfqpoint{1.581391in}{0.454693in}}%
\pgfpathlineto{\pgfqpoint{1.583660in}{0.491197in}}%
\pgfpathlineto{\pgfqpoint{1.585930in}{0.474609in}}%
\pgfpathlineto{\pgfqpoint{1.588199in}{0.445698in}}%
\pgfpathlineto{\pgfqpoint{1.590469in}{0.444051in}}%
\pgfpathlineto{\pgfqpoint{1.592738in}{0.507027in}}%
\pgfpathlineto{\pgfqpoint{1.595008in}{0.515090in}}%
\pgfpathlineto{\pgfqpoint{1.597277in}{0.348355in}}%
\pgfpathlineto{\pgfqpoint{1.599547in}{0.403506in}}%
\pgfpathlineto{\pgfqpoint{1.601816in}{0.482283in}}%
\pgfpathlineto{\pgfqpoint{1.604086in}{0.452553in}}%
\pgfpathlineto{\pgfqpoint{1.606355in}{0.491810in}}%
\pgfpathlineto{\pgfqpoint{1.608625in}{0.477553in}}%
\pgfpathlineto{\pgfqpoint{1.610894in}{0.458105in}}%
\pgfpathlineto{\pgfqpoint{1.613164in}{0.419580in}}%
\pgfpathlineto{\pgfqpoint{1.615433in}{0.504226in}}%
\pgfpathlineto{\pgfqpoint{1.617703in}{0.518043in}}%
\pgfpathlineto{\pgfqpoint{1.619972in}{0.482798in}}%
\pgfpathlineto{\pgfqpoint{1.622242in}{0.461837in}}%
\pgfpathlineto{\pgfqpoint{1.624511in}{0.527162in}}%
\pgfpathlineto{\pgfqpoint{1.626780in}{0.487971in}}%
\pgfpathlineto{\pgfqpoint{1.629050in}{0.514313in}}%
\pgfpathlineto{\pgfqpoint{1.631319in}{0.401880in}}%
\pgfpathlineto{\pgfqpoint{1.633589in}{0.462643in}}%
\pgfpathlineto{\pgfqpoint{1.635858in}{0.412944in}}%
\pgfpathlineto{\pgfqpoint{1.638128in}{0.440152in}}%
\pgfpathlineto{\pgfqpoint{1.640397in}{0.487395in}}%
\pgfpathlineto{\pgfqpoint{1.642667in}{0.496148in}}%
\pgfpathlineto{\pgfqpoint{1.644936in}{0.468806in}}%
\pgfpathlineto{\pgfqpoint{1.647206in}{0.455982in}}%
\pgfpathlineto{\pgfqpoint{1.649475in}{0.516660in}}%
\pgfpathlineto{\pgfqpoint{1.651745in}{0.527401in}}%
\pgfpathlineto{\pgfqpoint{1.656284in}{0.390428in}}%
\pgfpathlineto{\pgfqpoint{1.658553in}{0.514378in}}%
\pgfpathlineto{\pgfqpoint{1.660823in}{0.456445in}}%
\pgfpathlineto{\pgfqpoint{1.663092in}{0.504165in}}%
\pgfpathlineto{\pgfqpoint{1.665362in}{0.450576in}}%
\pgfpathlineto{\pgfqpoint{1.667631in}{0.475493in}}%
\pgfpathlineto{\pgfqpoint{1.669901in}{0.441350in}}%
\pgfpathlineto{\pgfqpoint{1.672170in}{0.440268in}}%
\pgfpathlineto{\pgfqpoint{1.674440in}{0.495289in}}%
\pgfpathlineto{\pgfqpoint{1.676709in}{0.512525in}}%
\pgfpathlineto{\pgfqpoint{1.678979in}{0.470177in}}%
\pgfpathlineto{\pgfqpoint{1.681248in}{0.503783in}}%
\pgfpathlineto{\pgfqpoint{1.683518in}{0.461566in}}%
\pgfpathlineto{\pgfqpoint{1.685787in}{0.506124in}}%
\pgfpathlineto{\pgfqpoint{1.688057in}{0.448212in}}%
\pgfpathlineto{\pgfqpoint{1.690326in}{0.479035in}}%
\pgfpathlineto{\pgfqpoint{1.694865in}{0.410125in}}%
\pgfpathlineto{\pgfqpoint{1.697134in}{0.485990in}}%
\pgfpathlineto{\pgfqpoint{1.699404in}{0.512093in}}%
\pgfpathlineto{\pgfqpoint{1.701673in}{0.485491in}}%
\pgfpathlineto{\pgfqpoint{1.703943in}{0.492194in}}%
\pgfpathlineto{\pgfqpoint{1.706212in}{0.455642in}}%
\pgfpathlineto{\pgfqpoint{1.708482in}{0.503623in}}%
\pgfpathlineto{\pgfqpoint{1.710751in}{0.453731in}}%
\pgfpathlineto{\pgfqpoint{1.713021in}{0.482911in}}%
\pgfpathlineto{\pgfqpoint{1.715290in}{0.455165in}}%
\pgfpathlineto{\pgfqpoint{1.717560in}{0.386105in}}%
\pgfpathlineto{\pgfqpoint{1.719829in}{0.483885in}}%
\pgfpathlineto{\pgfqpoint{1.722099in}{0.511761in}}%
\pgfpathlineto{\pgfqpoint{1.724368in}{0.493717in}}%
\pgfpathlineto{\pgfqpoint{1.726638in}{0.482474in}}%
\pgfpathlineto{\pgfqpoint{1.728907in}{0.447254in}}%
\pgfpathlineto{\pgfqpoint{1.731177in}{0.501023in}}%
\pgfpathlineto{\pgfqpoint{1.733446in}{0.456109in}}%
\pgfpathlineto{\pgfqpoint{1.735716in}{0.486129in}}%
\pgfpathlineto{\pgfqpoint{1.737985in}{0.461149in}}%
\pgfpathlineto{\pgfqpoint{1.740069in}{0.315625in}}%
\pgfpathmoveto{\pgfqpoint{1.740422in}{0.315625in}}%
\pgfpathlineto{\pgfqpoint{1.742524in}{0.478241in}}%
\pgfpathlineto{\pgfqpoint{1.744794in}{0.510947in}}%
\pgfpathlineto{\pgfqpoint{1.747063in}{0.500836in}}%
\pgfpathlineto{\pgfqpoint{1.751602in}{0.433401in}}%
\pgfpathlineto{\pgfqpoint{1.753872in}{0.497558in}}%
\pgfpathlineto{\pgfqpoint{1.756141in}{0.457648in}}%
\pgfpathlineto{\pgfqpoint{1.758410in}{0.488871in}}%
\pgfpathlineto{\pgfqpoint{1.760680in}{0.466446in}}%
\pgfpathlineto{\pgfqpoint{1.762949in}{0.398575in}}%
\pgfpathlineto{\pgfqpoint{1.765219in}{0.471238in}}%
\pgfpathlineto{\pgfqpoint{1.767488in}{0.509787in}}%
\pgfpathlineto{\pgfqpoint{1.769758in}{0.506573in}}%
\pgfpathlineto{\pgfqpoint{1.772027in}{0.449631in}}%
\pgfpathlineto{\pgfqpoint{1.774297in}{0.407791in}}%
\pgfpathlineto{\pgfqpoint{1.776566in}{0.493358in}}%
\pgfpathlineto{\pgfqpoint{1.778836in}{0.457937in}}%
\pgfpathlineto{\pgfqpoint{1.781105in}{0.490996in}}%
\pgfpathlineto{\pgfqpoint{1.783375in}{0.470949in}}%
\pgfpathlineto{\pgfqpoint{1.785644in}{0.427432in}}%
\pgfpathlineto{\pgfqpoint{1.787914in}{0.461429in}}%
\pgfpathlineto{\pgfqpoint{1.790183in}{0.508143in}}%
\pgfpathlineto{\pgfqpoint{1.792453in}{0.511257in}}%
\pgfpathlineto{\pgfqpoint{1.794722in}{0.449884in}}%
\pgfpathlineto{\pgfqpoint{1.799261in}{0.533513in}}%
\pgfpathlineto{\pgfqpoint{1.801531in}{0.485065in}}%
\pgfpathlineto{\pgfqpoint{1.803800in}{0.507108in}}%
\pgfpathlineto{\pgfqpoint{1.806070in}{0.395035in}}%
\pgfpathlineto{\pgfqpoint{1.808339in}{0.439565in}}%
\pgfpathlineto{\pgfqpoint{1.810609in}{0.447937in}}%
\pgfpathlineto{\pgfqpoint{1.812878in}{0.454094in}}%
\pgfpathlineto{\pgfqpoint{1.815148in}{0.490435in}}%
\pgfpathlineto{\pgfqpoint{1.817417in}{0.492646in}}%
\pgfpathlineto{\pgfqpoint{1.819687in}{0.436439in}}%
\pgfpathlineto{\pgfqpoint{1.821956in}{0.484562in}}%
\pgfpathlineto{\pgfqpoint{1.824226in}{0.520118in}}%
\pgfpathlineto{\pgfqpoint{1.826495in}{0.519548in}}%
\pgfpathlineto{\pgfqpoint{1.828764in}{0.400210in}}%
\pgfpathlineto{\pgfqpoint{1.833303in}{0.523523in}}%
\pgfpathlineto{\pgfqpoint{1.835573in}{0.459761in}}%
\pgfpathlineto{\pgfqpoint{1.837842in}{0.503893in}}%
\pgfpathlineto{\pgfqpoint{1.840112in}{0.439627in}}%
\pgfpathlineto{\pgfqpoint{1.842381in}{0.464090in}}%
\pgfpathlineto{\pgfqpoint{1.844651in}{0.401202in}}%
\pgfpathlineto{\pgfqpoint{1.846920in}{0.459534in}}%
\pgfpathlineto{\pgfqpoint{1.849190in}{0.498779in}}%
\pgfpathlineto{\pgfqpoint{1.851459in}{0.509217in}}%
\pgfpathlineto{\pgfqpoint{1.853729in}{0.410252in}}%
\pgfpathlineto{\pgfqpoint{1.855998in}{0.515166in}}%
\pgfpathlineto{\pgfqpoint{1.858268in}{0.468538in}}%
\pgfpathlineto{\pgfqpoint{1.860537in}{0.506920in}}%
\pgfpathlineto{\pgfqpoint{1.862807in}{0.437979in}}%
\pgfpathlineto{\pgfqpoint{1.865076in}{0.469509in}}%
\pgfpathlineto{\pgfqpoint{1.867346in}{0.420002in}}%
\pgfpathlineto{\pgfqpoint{1.869615in}{0.444112in}}%
\pgfpathlineto{\pgfqpoint{1.871885in}{0.493230in}}%
\pgfpathlineto{\pgfqpoint{1.874154in}{0.510242in}}%
\pgfpathlineto{\pgfqpoint{1.876424in}{0.456069in}}%
\pgfpathlineto{\pgfqpoint{1.878693in}{0.506397in}}%
\pgfpathlineto{\pgfqpoint{1.880963in}{0.466948in}}%
\pgfpathlineto{\pgfqpoint{1.883232in}{0.506763in}}%
\pgfpathlineto{\pgfqpoint{1.885502in}{0.447406in}}%
\pgfpathlineto{\pgfqpoint{1.887771in}{0.474618in}}%
\pgfpathlineto{\pgfqpoint{1.890041in}{0.433141in}}%
\pgfpathlineto{\pgfqpoint{1.892310in}{0.437784in}}%
\pgfpathlineto{\pgfqpoint{1.894579in}{0.492953in}}%
\pgfpathlineto{\pgfqpoint{1.896849in}{0.510564in}}%
\pgfpathlineto{\pgfqpoint{1.899118in}{0.471539in}}%
\pgfpathlineto{\pgfqpoint{1.901388in}{0.499678in}}%
\pgfpathlineto{\pgfqpoint{1.903657in}{0.463590in}}%
\pgfpathlineto{\pgfqpoint{1.905927in}{0.506283in}}%
\pgfpathlineto{\pgfqpoint{1.908196in}{0.452400in}}%
\pgfpathlineto{\pgfqpoint{1.910466in}{0.479202in}}%
\pgfpathlineto{\pgfqpoint{1.912735in}{0.443603in}}%
\pgfpathlineto{\pgfqpoint{1.915005in}{0.423775in}}%
\pgfpathlineto{\pgfqpoint{1.917274in}{0.490081in}}%
\pgfpathlineto{\pgfqpoint{1.919544in}{0.510637in}}%
\pgfpathlineto{\pgfqpoint{1.921813in}{0.483749in}}%
\pgfpathlineto{\pgfqpoint{1.924083in}{0.490670in}}%
\pgfpathlineto{\pgfqpoint{1.926352in}{0.457947in}}%
\pgfpathlineto{\pgfqpoint{1.928622in}{0.504510in}}%
\pgfpathlineto{\pgfqpoint{1.930891in}{0.456350in}}%
\pgfpathlineto{\pgfqpoint{1.933161in}{0.483310in}}%
\pgfpathlineto{\pgfqpoint{1.935430in}{0.452262in}}%
\pgfpathlineto{\pgfqpoint{1.937700in}{0.400474in}}%
\pgfpathlineto{\pgfqpoint{1.939969in}{0.486422in}}%
\pgfpathlineto{\pgfqpoint{1.942239in}{0.510401in}}%
\pgfpathlineto{\pgfqpoint{1.946778in}{0.479655in}}%
\pgfpathlineto{\pgfqpoint{1.949047in}{0.449234in}}%
\pgfpathlineto{\pgfqpoint{1.951317in}{0.501682in}}%
\pgfpathlineto{\pgfqpoint{1.953586in}{0.458920in}}%
\pgfpathlineto{\pgfqpoint{1.955856in}{0.486857in}}%
\pgfpathlineto{\pgfqpoint{1.958125in}{0.459452in}}%
\pgfpathlineto{\pgfqpoint{1.960394in}{0.323994in}}%
\pgfpathlineto{\pgfqpoint{1.962664in}{0.481269in}}%
\pgfpathlineto{\pgfqpoint{1.964933in}{0.509776in}}%
\pgfpathlineto{\pgfqpoint{1.967203in}{0.500248in}}%
\pgfpathlineto{\pgfqpoint{1.971742in}{0.435315in}}%
\pgfpathlineto{\pgfqpoint{1.974011in}{0.498143in}}%
\pgfpathlineto{\pgfqpoint{1.976281in}{0.460305in}}%
\pgfpathlineto{\pgfqpoint{1.978550in}{0.489805in}}%
\pgfpathlineto{\pgfqpoint{1.980820in}{0.465471in}}%
\pgfpathlineto{\pgfqpoint{1.983089in}{0.388694in}}%
\pgfpathlineto{\pgfqpoint{1.985359in}{0.474340in}}%
\pgfpathlineto{\pgfqpoint{1.987628in}{0.508709in}}%
\pgfpathlineto{\pgfqpoint{1.989898in}{0.506154in}}%
\pgfpathlineto{\pgfqpoint{1.992167in}{0.405343in}}%
\pgfpathlineto{\pgfqpoint{1.994437in}{0.496857in}}%
\pgfpathlineto{\pgfqpoint{1.996706in}{0.534001in}}%
\pgfpathlineto{\pgfqpoint{1.998976in}{0.482265in}}%
\pgfpathlineto{\pgfqpoint{2.001245in}{0.501869in}}%
\pgfpathlineto{\pgfqpoint{2.003515in}{0.426275in}}%
\pgfpathlineto{\pgfqpoint{2.005784in}{0.418556in}}%
\pgfpathlineto{\pgfqpoint{2.008054in}{0.460291in}}%
\pgfpathlineto{\pgfqpoint{2.010323in}{0.459994in}}%
\pgfpathlineto{\pgfqpoint{2.012593in}{0.491213in}}%
\pgfpathlineto{\pgfqpoint{2.014862in}{0.489325in}}%
\pgfpathlineto{\pgfqpoint{2.017132in}{0.396377in}}%
\pgfpathlineto{\pgfqpoint{2.019401in}{0.493483in}}%
\pgfpathlineto{\pgfqpoint{2.021671in}{0.519976in}}%
\pgfpathlineto{\pgfqpoint{2.023940in}{0.513862in}}%
\pgfpathlineto{\pgfqpoint{2.026209in}{0.446892in}}%
\pgfpathlineto{\pgfqpoint{2.028479in}{0.478288in}}%
\pgfpathlineto{\pgfqpoint{2.030748in}{0.526138in}}%
\pgfpathlineto{\pgfqpoint{2.033018in}{0.460622in}}%
\pgfpathlineto{\pgfqpoint{2.035287in}{0.502494in}}%
\pgfpathlineto{\pgfqpoint{2.037557in}{0.431750in}}%
\pgfpathlineto{\pgfqpoint{2.039826in}{0.455987in}}%
\pgfpathlineto{\pgfqpoint{2.041146in}{0.315625in}}%
\pgfpathmoveto{\pgfqpoint{2.043005in}{0.315625in}}%
\pgfpathlineto{\pgfqpoint{2.044365in}{0.466823in}}%
\pgfpathlineto{\pgfqpoint{2.046635in}{0.499547in}}%
\pgfpathlineto{\pgfqpoint{2.048904in}{0.505762in}}%
\pgfpathlineto{\pgfqpoint{2.051174in}{0.380154in}}%
\pgfpathlineto{\pgfqpoint{2.053443in}{0.518737in}}%
\pgfpathlineto{\pgfqpoint{2.055713in}{0.470568in}}%
\pgfpathlineto{\pgfqpoint{2.057982in}{0.505972in}}%
\pgfpathlineto{\pgfqpoint{2.060252in}{0.430799in}}%
\pgfpathlineto{\pgfqpoint{2.062521in}{0.463057in}}%
\pgfpathlineto{\pgfqpoint{2.064791in}{0.383395in}}%
\pgfpathlineto{\pgfqpoint{2.067060in}{0.454616in}}%
\pgfpathlineto{\pgfqpoint{2.069330in}{0.495611in}}%
\pgfpathlineto{\pgfqpoint{2.071599in}{0.507834in}}%
\pgfpathlineto{\pgfqpoint{2.073869in}{0.427699in}}%
\pgfpathlineto{\pgfqpoint{2.076138in}{0.510919in}}%
\pgfpathlineto{\pgfqpoint{2.078408in}{0.470428in}}%
\pgfpathlineto{\pgfqpoint{2.080677in}{0.506760in}}%
\pgfpathlineto{\pgfqpoint{2.082947in}{0.443136in}}%
\pgfpathlineto{\pgfqpoint{2.085216in}{0.469154in}}%
\pgfpathlineto{\pgfqpoint{2.087486in}{0.410341in}}%
\pgfpathlineto{\pgfqpoint{2.092024in}{0.496090in}}%
\pgfpathlineto{\pgfqpoint{2.094294in}{0.508586in}}%
\pgfpathlineto{\pgfqpoint{2.096563in}{0.454206in}}%
\pgfpathlineto{\pgfqpoint{2.098833in}{0.505052in}}%
\pgfpathlineto{\pgfqpoint{2.101102in}{0.468468in}}%
\pgfpathlineto{\pgfqpoint{2.103372in}{0.507217in}}%
\pgfpathlineto{\pgfqpoint{2.105641in}{0.449788in}}%
\pgfpathlineto{\pgfqpoint{2.107911in}{0.474702in}}%
\pgfpathlineto{\pgfqpoint{2.110180in}{0.427975in}}%
\pgfpathlineto{\pgfqpoint{2.112450in}{0.441335in}}%
\pgfpathlineto{\pgfqpoint{2.114719in}{0.494303in}}%
\pgfpathlineto{\pgfqpoint{2.116989in}{0.509187in}}%
\pgfpathlineto{\pgfqpoint{2.119258in}{0.471816in}}%
\pgfpathlineto{\pgfqpoint{2.121528in}{0.497249in}}%
\pgfpathlineto{\pgfqpoint{2.123797in}{0.464698in}}%
\pgfpathlineto{\pgfqpoint{2.126067in}{0.506597in}}%
\pgfpathlineto{\pgfqpoint{2.128336in}{0.455073in}}%
\pgfpathlineto{\pgfqpoint{2.130606in}{0.479669in}}%
\pgfpathlineto{\pgfqpoint{2.132875in}{0.440802in}}%
\pgfpathlineto{\pgfqpoint{2.135145in}{0.428487in}}%
\pgfpathlineto{\pgfqpoint{2.137414in}{0.491751in}}%
\pgfpathlineto{\pgfqpoint{2.139684in}{0.509447in}}%
\pgfpathlineto{\pgfqpoint{2.141953in}{0.483957in}}%
\pgfpathlineto{\pgfqpoint{2.144223in}{0.487917in}}%
\pgfpathlineto{\pgfqpoint{2.146492in}{0.458694in}}%
\pgfpathlineto{\pgfqpoint{2.148762in}{0.504798in}}%
\pgfpathlineto{\pgfqpoint{2.151031in}{0.458763in}}%
\pgfpathlineto{\pgfqpoint{2.153301in}{0.484018in}}%
\pgfpathlineto{\pgfqpoint{2.155570in}{0.450759in}}%
\pgfpathlineto{\pgfqpoint{2.157840in}{0.406249in}}%
\pgfpathlineto{\pgfqpoint{2.160109in}{0.487936in}}%
\pgfpathlineto{\pgfqpoint{2.162378in}{0.509326in}}%
\pgfpathlineto{\pgfqpoint{2.164648in}{0.493225in}}%
\pgfpathlineto{\pgfqpoint{2.166917in}{0.450387in}}%
\pgfpathlineto{\pgfqpoint{2.169187in}{0.508360in}}%
\pgfpathlineto{\pgfqpoint{2.171456in}{0.534105in}}%
\pgfpathlineto{\pgfqpoint{2.173726in}{0.469431in}}%
\pgfpathlineto{\pgfqpoint{2.175995in}{0.487490in}}%
\pgfpathlineto{\pgfqpoint{2.178265in}{0.461859in}}%
\pgfpathlineto{\pgfqpoint{2.180451in}{0.315625in}}%
\pgfpathmoveto{\pgfqpoint{2.180611in}{0.315625in}}%
\pgfpathlineto{\pgfqpoint{2.182804in}{0.474830in}}%
\pgfpathlineto{\pgfqpoint{2.185073in}{0.465018in}}%
\pgfpathlineto{\pgfqpoint{2.187343in}{0.488615in}}%
\pgfpathlineto{\pgfqpoint{2.189612in}{0.481495in}}%
\pgfpathlineto{\pgfqpoint{2.191882in}{0.422455in}}%
\pgfpathlineto{\pgfqpoint{2.194151in}{0.504880in}}%
\pgfpathlineto{\pgfqpoint{2.196421in}{0.520059in}}%
\pgfpathlineto{\pgfqpoint{2.198690in}{0.498625in}}%
\pgfpathlineto{\pgfqpoint{2.200960in}{0.488403in}}%
\pgfpathlineto{\pgfqpoint{2.203229in}{0.494972in}}%
\pgfpathlineto{\pgfqpoint{2.205499in}{0.530559in}}%
\pgfpathlineto{\pgfqpoint{2.207768in}{0.453009in}}%
\pgfpathlineto{\pgfqpoint{2.210038in}{0.496783in}}%
\pgfpathlineto{\pgfqpoint{2.212307in}{0.392336in}}%
\pgfpathlineto{\pgfqpoint{2.214577in}{0.433532in}}%
\pgfpathlineto{\pgfqpoint{2.216846in}{0.424255in}}%
\pgfpathlineto{\pgfqpoint{2.219116in}{0.474924in}}%
\pgfpathlineto{\pgfqpoint{2.221385in}{0.497635in}}%
\pgfpathlineto{\pgfqpoint{2.223655in}{0.498392in}}%
\pgfpathlineto{\pgfqpoint{2.225924in}{0.462340in}}%
\pgfpathlineto{\pgfqpoint{2.228193in}{0.525870in}}%
\pgfpathlineto{\pgfqpoint{2.230463in}{0.468779in}}%
\pgfpathlineto{\pgfqpoint{2.232732in}{0.501155in}}%
\pgfpathlineto{\pgfqpoint{2.235002in}{0.392366in}}%
\pgfpathlineto{\pgfqpoint{2.237271in}{0.445150in}}%
\pgfpathlineto{\pgfqpoint{2.239541in}{0.400815in}}%
\pgfpathlineto{\pgfqpoint{2.241810in}{0.466320in}}%
\pgfpathlineto{\pgfqpoint{2.244080in}{0.496022in}}%
\pgfpathlineto{\pgfqpoint{2.246349in}{0.502573in}}%
\pgfpathlineto{\pgfqpoint{2.248619in}{0.425152in}}%
\pgfpathlineto{\pgfqpoint{2.250888in}{0.519814in}}%
\pgfpathlineto{\pgfqpoint{2.253158in}{0.472092in}}%
\pgfpathlineto{\pgfqpoint{2.255427in}{0.503662in}}%
\pgfpathlineto{\pgfqpoint{2.257697in}{0.421993in}}%
\pgfpathlineto{\pgfqpoint{2.259966in}{0.453554in}}%
\pgfpathlineto{\pgfqpoint{2.262236in}{0.368864in}}%
\pgfpathlineto{\pgfqpoint{2.264505in}{0.464748in}}%
\pgfpathlineto{\pgfqpoint{2.266775in}{0.497752in}}%
\pgfpathlineto{\pgfqpoint{2.269044in}{0.504045in}}%
\pgfpathlineto{\pgfqpoint{2.271314in}{0.332497in}}%
\pgfpathlineto{\pgfqpoint{2.273583in}{0.515556in}}%
\pgfpathlineto{\pgfqpoint{2.275853in}{0.472829in}}%
\pgfpathlineto{\pgfqpoint{2.278122in}{0.505768in}}%
\pgfpathlineto{\pgfqpoint{2.280392in}{0.435323in}}%
\pgfpathlineto{\pgfqpoint{2.282661in}{0.461371in}}%
\pgfpathlineto{\pgfqpoint{2.284931in}{0.328470in}}%
\pgfpathlineto{\pgfqpoint{2.287200in}{0.460073in}}%
\pgfpathlineto{\pgfqpoint{2.289470in}{0.497898in}}%
\pgfpathlineto{\pgfqpoint{2.291739in}{0.505637in}}%
\pgfpathlineto{\pgfqpoint{2.294008in}{0.421577in}}%
\pgfpathlineto{\pgfqpoint{2.296278in}{0.509865in}}%
\pgfpathlineto{\pgfqpoint{2.298547in}{0.472197in}}%
\pgfpathlineto{\pgfqpoint{2.300817in}{0.507068in}}%
\pgfpathlineto{\pgfqpoint{2.303086in}{0.457324in}}%
\pgfpathlineto{\pgfqpoint{2.303086in}{0.457324in}}%
\pgfusepath{stroke}%
\end{pgfscope}%
\begin{pgfscope}%
\pgfpathrectangle{\pgfqpoint{0.487500in}{0.320625in}}{\pgfqpoint{1.815586in}{1.164375in}}%
\pgfusepath{clip}%
\pgfsetbuttcap%
\pgfsetroundjoin%
\pgfsetlinewidth{0.752812pt}%
\definecolor{currentstroke}{rgb}{0.000000,0.529412,0.423529}%
\pgfsetstrokecolor{currentstroke}%
\pgfsetstrokeopacity{0.650000}%
\pgfsetdash{{2.775000pt}{1.200000pt}}{0.000000pt}%
\pgfpathmoveto{\pgfqpoint{0.487500in}{1.451682in}}%
\pgfpathlineto{\pgfqpoint{1.099331in}{0.315625in}}%
\pgfpathlineto{\pgfqpoint{1.099331in}{0.315625in}}%
\pgfusepath{stroke}%
\end{pgfscope}%
\begin{pgfscope}%
\pgfpathrectangle{\pgfqpoint{0.487500in}{0.320625in}}{\pgfqpoint{1.815586in}{1.164375in}}%
\pgfusepath{clip}%
\pgfsetbuttcap%
\pgfsetroundjoin%
\pgfsetlinewidth{0.803000pt}%
\definecolor{currentstroke}{rgb}{0.000000,0.529412,0.423529}%
\pgfsetstrokecolor{currentstroke}%
\pgfsetstrokeopacity{0.800000}%
\pgfsetdash{{0.800000pt}{1.320000pt}}{0.000000pt}%
\pgfpathmoveto{\pgfqpoint{0.487500in}{1.451682in}}%
\pgfpathlineto{\pgfqpoint{2.303086in}{1.451682in}}%
\pgfpathlineto{\pgfqpoint{2.303086in}{1.451682in}}%
\pgfusepath{stroke}%
\end{pgfscope}%
\begin{pgfscope}%
\pgfpathrectangle{\pgfqpoint{0.487500in}{0.320625in}}{\pgfqpoint{1.815586in}{1.164375in}}%
\pgfusepath{clip}%
\pgfsetrectcap%
\pgfsetroundjoin%
\pgfsetlinewidth{1.254687pt}%
\definecolor{currentstroke}{rgb}{0.478431,0.243137,0.615686}%
\pgfsetstrokecolor{currentstroke}%
\pgfsetdash{}{0pt}%
\pgfpathmoveto{\pgfqpoint{0.487500in}{1.451682in}}%
\pgfpathlineto{\pgfqpoint{0.496578in}{1.364963in}}%
\pgfpathlineto{\pgfqpoint{0.589627in}{0.488481in}}%
\pgfpathlineto{\pgfqpoint{0.591896in}{0.489143in}}%
\pgfpathlineto{\pgfqpoint{0.594166in}{0.478292in}}%
\pgfpathlineto{\pgfqpoint{0.596435in}{0.510544in}}%
\pgfpathlineto{\pgfqpoint{0.598705in}{0.482742in}}%
\pgfpathlineto{\pgfqpoint{0.600974in}{0.509429in}}%
\pgfpathlineto{\pgfqpoint{0.603244in}{0.467473in}}%
\pgfpathlineto{\pgfqpoint{0.605513in}{0.489698in}}%
\pgfpathlineto{\pgfqpoint{0.607783in}{0.475048in}}%
\pgfpathlineto{\pgfqpoint{0.610052in}{0.400980in}}%
\pgfpathlineto{\pgfqpoint{0.612322in}{0.495196in}}%
\pgfpathlineto{\pgfqpoint{0.614591in}{0.485052in}}%
\pgfpathlineto{\pgfqpoint{0.616861in}{0.320237in}}%
\pgfpathlineto{\pgfqpoint{0.619130in}{0.499125in}}%
\pgfpathlineto{\pgfqpoint{0.621399in}{0.505561in}}%
\pgfpathlineto{\pgfqpoint{0.623669in}{0.438635in}}%
\pgfpathlineto{\pgfqpoint{0.625938in}{0.494961in}}%
\pgfpathlineto{\pgfqpoint{0.630477in}{0.447492in}}%
\pgfpathlineto{\pgfqpoint{0.632747in}{0.507085in}}%
\pgfpathlineto{\pgfqpoint{0.635016in}{0.434260in}}%
\pgfpathlineto{\pgfqpoint{0.637286in}{0.472163in}}%
\pgfpathlineto{\pgfqpoint{0.639555in}{0.499377in}}%
\pgfpathlineto{\pgfqpoint{0.641825in}{0.452810in}}%
\pgfpathlineto{\pgfqpoint{0.644094in}{0.482918in}}%
\pgfpathlineto{\pgfqpoint{0.646364in}{0.489203in}}%
\pgfpathlineto{\pgfqpoint{0.648633in}{0.499103in}}%
\pgfpathlineto{\pgfqpoint{0.650903in}{0.384900in}}%
\pgfpathlineto{\pgfqpoint{0.653172in}{0.498536in}}%
\pgfpathlineto{\pgfqpoint{0.655442in}{0.465991in}}%
\pgfpathlineto{\pgfqpoint{0.657711in}{0.463951in}}%
\pgfpathlineto{\pgfqpoint{0.659981in}{0.517086in}}%
\pgfpathlineto{\pgfqpoint{0.664520in}{0.449421in}}%
\pgfpathlineto{\pgfqpoint{0.666789in}{0.511404in}}%
\pgfpathlineto{\pgfqpoint{0.669059in}{0.501686in}}%
\pgfpathlineto{\pgfqpoint{0.671328in}{0.418236in}}%
\pgfpathlineto{\pgfqpoint{0.673598in}{0.496394in}}%
\pgfpathlineto{\pgfqpoint{0.675867in}{0.440270in}}%
\pgfpathlineto{\pgfqpoint{0.678137in}{0.474335in}}%
\pgfpathlineto{\pgfqpoint{0.680406in}{0.524240in}}%
\pgfpathlineto{\pgfqpoint{0.682676in}{0.411642in}}%
\pgfpathlineto{\pgfqpoint{0.684945in}{0.503962in}}%
\pgfpathlineto{\pgfqpoint{0.687215in}{0.504445in}}%
\pgfpathlineto{\pgfqpoint{0.689484in}{0.422331in}}%
\pgfpathlineto{\pgfqpoint{0.691753in}{0.412807in}}%
\pgfpathlineto{\pgfqpoint{0.694023in}{0.469177in}}%
\pgfpathlineto{\pgfqpoint{0.696292in}{0.433517in}}%
\pgfpathlineto{\pgfqpoint{0.698562in}{0.451136in}}%
\pgfpathlineto{\pgfqpoint{0.700831in}{0.480890in}}%
\pgfpathlineto{\pgfqpoint{0.703101in}{0.478863in}}%
\pgfpathlineto{\pgfqpoint{0.705370in}{0.434955in}}%
\pgfpathlineto{\pgfqpoint{0.709909in}{0.450349in}}%
\pgfpathlineto{\pgfqpoint{0.712179in}{0.420093in}}%
\pgfpathlineto{\pgfqpoint{0.714448in}{0.411872in}}%
\pgfpathlineto{\pgfqpoint{0.716718in}{0.433683in}}%
\pgfpathlineto{\pgfqpoint{0.718987in}{0.447965in}}%
\pgfpathlineto{\pgfqpoint{0.721257in}{0.412294in}}%
\pgfpathlineto{\pgfqpoint{0.723526in}{0.466243in}}%
\pgfpathlineto{\pgfqpoint{0.725796in}{0.475418in}}%
\pgfpathlineto{\pgfqpoint{0.728065in}{0.514928in}}%
\pgfpathlineto{\pgfqpoint{0.730335in}{0.478260in}}%
\pgfpathlineto{\pgfqpoint{0.732604in}{0.509869in}}%
\pgfpathlineto{\pgfqpoint{0.734874in}{0.470879in}}%
\pgfpathlineto{\pgfqpoint{0.737143in}{0.487944in}}%
\pgfpathlineto{\pgfqpoint{0.739413in}{0.473917in}}%
\pgfpathlineto{\pgfqpoint{0.741682in}{0.405924in}}%
\pgfpathlineto{\pgfqpoint{0.743952in}{0.496467in}}%
\pgfpathlineto{\pgfqpoint{0.746221in}{0.452132in}}%
\pgfpathlineto{\pgfqpoint{0.748491in}{0.442184in}}%
\pgfpathlineto{\pgfqpoint{0.750760in}{0.493980in}}%
\pgfpathlineto{\pgfqpoint{0.753030in}{0.448021in}}%
\pgfpathlineto{\pgfqpoint{0.755299in}{0.466710in}}%
\pgfpathlineto{\pgfqpoint{0.757568in}{0.507399in}}%
\pgfpathlineto{\pgfqpoint{0.759838in}{0.457227in}}%
\pgfpathlineto{\pgfqpoint{0.762107in}{0.462820in}}%
\pgfpathlineto{\pgfqpoint{0.764377in}{0.514962in}}%
\pgfpathlineto{\pgfqpoint{0.766646in}{0.402949in}}%
\pgfpathlineto{\pgfqpoint{0.768916in}{0.490029in}}%
\pgfpathlineto{\pgfqpoint{0.771185in}{0.504248in}}%
\pgfpathlineto{\pgfqpoint{0.773455in}{0.421036in}}%
\pgfpathlineto{\pgfqpoint{0.777994in}{0.484650in}}%
\pgfpathlineto{\pgfqpoint{0.780263in}{0.485975in}}%
\pgfpathlineto{\pgfqpoint{0.782533in}{0.412194in}}%
\pgfpathlineto{\pgfqpoint{0.784802in}{0.495825in}}%
\pgfpathlineto{\pgfqpoint{0.787072in}{0.441105in}}%
\pgfpathlineto{\pgfqpoint{0.789341in}{0.475616in}}%
\pgfpathlineto{\pgfqpoint{0.791611in}{0.522431in}}%
\pgfpathlineto{\pgfqpoint{0.793880in}{0.464704in}}%
\pgfpathlineto{\pgfqpoint{0.796150in}{0.462553in}}%
\pgfpathlineto{\pgfqpoint{0.798419in}{0.509932in}}%
\pgfpathlineto{\pgfqpoint{0.800689in}{0.486819in}}%
\pgfpathlineto{\pgfqpoint{0.802958in}{0.445290in}}%
\pgfpathlineto{\pgfqpoint{0.805228in}{0.481147in}}%
\pgfpathlineto{\pgfqpoint{0.807497in}{0.388043in}}%
\pgfpathlineto{\pgfqpoint{0.809767in}{0.477209in}}%
\pgfpathlineto{\pgfqpoint{0.812036in}{0.524716in}}%
\pgfpathlineto{\pgfqpoint{0.814306in}{0.458966in}}%
\pgfpathlineto{\pgfqpoint{0.816575in}{0.510551in}}%
\pgfpathlineto{\pgfqpoint{0.821114in}{0.469635in}}%
\pgfpathlineto{\pgfqpoint{0.823383in}{0.462361in}}%
\pgfpathlineto{\pgfqpoint{0.825653in}{0.423171in}}%
\pgfpathlineto{\pgfqpoint{0.827922in}{0.459068in}}%
\pgfpathlineto{\pgfqpoint{0.830192in}{0.457571in}}%
\pgfpathlineto{\pgfqpoint{0.832461in}{0.425234in}}%
\pgfpathlineto{\pgfqpoint{0.834731in}{0.474139in}}%
\pgfpathlineto{\pgfqpoint{0.837000in}{0.476849in}}%
\pgfpathlineto{\pgfqpoint{0.839270in}{0.504134in}}%
\pgfpathlineto{\pgfqpoint{0.841539in}{0.475674in}}%
\pgfpathlineto{\pgfqpoint{0.843809in}{0.458298in}}%
\pgfpathlineto{\pgfqpoint{0.846078in}{0.449349in}}%
\pgfpathlineto{\pgfqpoint{0.848348in}{0.477665in}}%
\pgfpathlineto{\pgfqpoint{0.850617in}{0.442079in}}%
\pgfpathlineto{\pgfqpoint{0.852887in}{0.453881in}}%
\pgfpathlineto{\pgfqpoint{0.855156in}{0.494563in}}%
\pgfpathlineto{\pgfqpoint{0.857426in}{0.455603in}}%
\pgfpathlineto{\pgfqpoint{0.859695in}{0.492714in}}%
\pgfpathlineto{\pgfqpoint{0.861965in}{0.487100in}}%
\pgfpathlineto{\pgfqpoint{0.864234in}{0.502658in}}%
\pgfpathlineto{\pgfqpoint{0.866504in}{0.412020in}}%
\pgfpathlineto{\pgfqpoint{0.868773in}{0.497051in}}%
\pgfpathlineto{\pgfqpoint{0.871043in}{0.466897in}}%
\pgfpathlineto{\pgfqpoint{0.873312in}{0.462256in}}%
\pgfpathlineto{\pgfqpoint{0.875582in}{0.516393in}}%
\pgfpathlineto{\pgfqpoint{0.877851in}{0.407879in}}%
\pgfpathlineto{\pgfqpoint{0.880121in}{0.488385in}}%
\pgfpathlineto{\pgfqpoint{0.882390in}{0.521736in}}%
\pgfpathlineto{\pgfqpoint{0.884660in}{0.395814in}}%
\pgfpathlineto{\pgfqpoint{0.886929in}{0.488998in}}%
\pgfpathlineto{\pgfqpoint{0.889198in}{0.527402in}}%
\pgfpathlineto{\pgfqpoint{0.891468in}{0.421310in}}%
\pgfpathlineto{\pgfqpoint{0.896007in}{0.526285in}}%
\pgfpathlineto{\pgfqpoint{0.898276in}{0.423544in}}%
\pgfpathlineto{\pgfqpoint{0.900546in}{0.506606in}}%
\pgfpathlineto{\pgfqpoint{0.902815in}{0.502896in}}%
\pgfpathlineto{\pgfqpoint{0.905085in}{0.433509in}}%
\pgfpathlineto{\pgfqpoint{0.907354in}{0.413890in}}%
\pgfpathlineto{\pgfqpoint{0.909624in}{0.462232in}}%
\pgfpathlineto{\pgfqpoint{0.911893in}{0.419894in}}%
\pgfpathlineto{\pgfqpoint{0.916432in}{0.474795in}}%
\pgfpathlineto{\pgfqpoint{0.918702in}{0.423956in}}%
\pgfpathlineto{\pgfqpoint{0.920971in}{0.480240in}}%
\pgfpathlineto{\pgfqpoint{0.923241in}{0.519663in}}%
\pgfpathlineto{\pgfqpoint{0.925510in}{0.465442in}}%
\pgfpathlineto{\pgfqpoint{0.927780in}{0.507826in}}%
\pgfpathlineto{\pgfqpoint{0.930049in}{0.531832in}}%
\pgfpathlineto{\pgfqpoint{0.932319in}{0.471175in}}%
\pgfpathlineto{\pgfqpoint{0.934588in}{0.495686in}}%
\pgfpathlineto{\pgfqpoint{0.936858in}{0.529879in}}%
\pgfpathlineto{\pgfqpoint{0.939127in}{0.475289in}}%
\pgfpathlineto{\pgfqpoint{0.941397in}{0.475727in}}%
\pgfpathlineto{\pgfqpoint{0.943666in}{0.515057in}}%
\pgfpathlineto{\pgfqpoint{0.945936in}{0.480554in}}%
\pgfpathlineto{\pgfqpoint{0.948205in}{0.509946in}}%
\pgfpathlineto{\pgfqpoint{0.950475in}{0.462248in}}%
\pgfpathlineto{\pgfqpoint{0.952744in}{0.490768in}}%
\pgfpathlineto{\pgfqpoint{0.955014in}{0.472549in}}%
\pgfpathlineto{\pgfqpoint{0.957283in}{0.432138in}}%
\pgfpathlineto{\pgfqpoint{0.959552in}{0.502438in}}%
\pgfpathlineto{\pgfqpoint{0.961822in}{0.444186in}}%
\pgfpathlineto{\pgfqpoint{0.964091in}{0.461878in}}%
\pgfpathlineto{\pgfqpoint{0.966361in}{0.498859in}}%
\pgfpathlineto{\pgfqpoint{0.968630in}{0.456842in}}%
\pgfpathlineto{\pgfqpoint{0.970900in}{0.452257in}}%
\pgfpathlineto{\pgfqpoint{0.973169in}{0.506157in}}%
\pgfpathlineto{\pgfqpoint{0.975439in}{0.420083in}}%
\pgfpathlineto{\pgfqpoint{0.979978in}{0.513804in}}%
\pgfpathlineto{\pgfqpoint{0.982247in}{0.367673in}}%
\pgfpathlineto{\pgfqpoint{0.984517in}{0.494946in}}%
\pgfpathlineto{\pgfqpoint{0.986786in}{0.505533in}}%
\pgfpathlineto{\pgfqpoint{0.989056in}{0.366631in}}%
\pgfpathlineto{\pgfqpoint{0.991325in}{0.443648in}}%
\pgfpathlineto{\pgfqpoint{0.993595in}{0.478868in}}%
\pgfpathlineto{\pgfqpoint{0.995864in}{0.480318in}}%
\pgfpathlineto{\pgfqpoint{0.998134in}{0.423869in}}%
\pgfpathlineto{\pgfqpoint{1.000403in}{0.491822in}}%
\pgfpathlineto{\pgfqpoint{1.002673in}{0.424018in}}%
\pgfpathlineto{\pgfqpoint{1.007212in}{0.523673in}}%
\pgfpathlineto{\pgfqpoint{1.009481in}{0.429953in}}%
\pgfpathlineto{\pgfqpoint{1.011751in}{0.504365in}}%
\pgfpathlineto{\pgfqpoint{1.014020in}{0.501178in}}%
\pgfpathlineto{\pgfqpoint{1.016290in}{0.444332in}}%
\pgfpathlineto{\pgfqpoint{1.018559in}{0.499033in}}%
\pgfpathlineto{\pgfqpoint{1.020829in}{0.515111in}}%
\pgfpathlineto{\pgfqpoint{1.023098in}{0.460957in}}%
\pgfpathlineto{\pgfqpoint{1.025367in}{0.514169in}}%
\pgfpathlineto{\pgfqpoint{1.027637in}{0.528512in}}%
\pgfpathlineto{\pgfqpoint{1.029906in}{0.467077in}}%
\pgfpathlineto{\pgfqpoint{1.032176in}{0.513298in}}%
\pgfpathlineto{\pgfqpoint{1.034445in}{0.486269in}}%
\pgfpathlineto{\pgfqpoint{1.036715in}{0.477151in}}%
\pgfpathlineto{\pgfqpoint{1.038984in}{0.464997in}}%
\pgfpathlineto{\pgfqpoint{1.041254in}{0.363140in}}%
\pgfpathlineto{\pgfqpoint{1.043523in}{0.474360in}}%
\pgfpathlineto{\pgfqpoint{1.045793in}{0.453700in}}%
\pgfpathlineto{\pgfqpoint{1.047673in}{0.315625in}}%
\pgfpathmoveto{\pgfqpoint{1.048396in}{0.315625in}}%
\pgfpathlineto{\pgfqpoint{1.050332in}{0.481874in}}%
\pgfpathlineto{\pgfqpoint{1.052601in}{0.472731in}}%
\pgfpathlineto{\pgfqpoint{1.054871in}{0.502909in}}%
\pgfpathlineto{\pgfqpoint{1.057140in}{0.485691in}}%
\pgfpathlineto{\pgfqpoint{1.059410in}{0.505645in}}%
\pgfpathlineto{\pgfqpoint{1.061679in}{0.518181in}}%
\pgfpathlineto{\pgfqpoint{1.063949in}{0.490022in}}%
\pgfpathlineto{\pgfqpoint{1.066218in}{0.485494in}}%
\pgfpathlineto{\pgfqpoint{1.068488in}{0.511797in}}%
\pgfpathlineto{\pgfqpoint{1.070757in}{0.497381in}}%
\pgfpathlineto{\pgfqpoint{1.073027in}{0.453003in}}%
\pgfpathlineto{\pgfqpoint{1.075296in}{0.491467in}}%
\pgfpathlineto{\pgfqpoint{1.077566in}{0.484725in}}%
\pgfpathlineto{\pgfqpoint{1.079835in}{0.499632in}}%
\pgfpathlineto{\pgfqpoint{1.082105in}{0.373129in}}%
\pgfpathlineto{\pgfqpoint{1.084374in}{0.496585in}}%
\pgfpathlineto{\pgfqpoint{1.086644in}{0.460993in}}%
\pgfpathlineto{\pgfqpoint{1.088913in}{0.467815in}}%
\pgfpathlineto{\pgfqpoint{1.091182in}{0.519511in}}%
\pgfpathlineto{\pgfqpoint{1.093452in}{0.358831in}}%
\pgfpathlineto{\pgfqpoint{1.095721in}{0.494874in}}%
\pgfpathlineto{\pgfqpoint{1.097991in}{0.506446in}}%
\pgfpathlineto{\pgfqpoint{1.100260in}{0.347143in}}%
\pgfpathlineto{\pgfqpoint{1.102530in}{0.353484in}}%
\pgfpathlineto{\pgfqpoint{1.104799in}{0.481819in}}%
\pgfpathlineto{\pgfqpoint{1.107069in}{0.464628in}}%
\pgfpathlineto{\pgfqpoint{1.109338in}{0.421773in}}%
\pgfpathlineto{\pgfqpoint{1.111608in}{0.483491in}}%
\pgfpathlineto{\pgfqpoint{1.113877in}{0.332511in}}%
\pgfpathlineto{\pgfqpoint{1.116147in}{0.439057in}}%
\pgfpathlineto{\pgfqpoint{1.118416in}{0.495728in}}%
\pgfpathlineto{\pgfqpoint{1.120686in}{0.449289in}}%
\pgfpathlineto{\pgfqpoint{1.122955in}{0.424216in}}%
\pgfpathlineto{\pgfqpoint{1.125225in}{0.444842in}}%
\pgfpathlineto{\pgfqpoint{1.127494in}{0.338512in}}%
\pgfpathlineto{\pgfqpoint{1.129764in}{0.450980in}}%
\pgfpathlineto{\pgfqpoint{1.132033in}{0.460225in}}%
\pgfpathlineto{\pgfqpoint{1.134303in}{0.446695in}}%
\pgfpathlineto{\pgfqpoint{1.138842in}{0.518074in}}%
\pgfpathlineto{\pgfqpoint{1.141111in}{0.471987in}}%
\pgfpathlineto{\pgfqpoint{1.143381in}{0.509202in}}%
\pgfpathlineto{\pgfqpoint{1.145650in}{0.480803in}}%
\pgfpathlineto{\pgfqpoint{1.147920in}{0.482131in}}%
\pgfpathlineto{\pgfqpoint{1.150189in}{0.514189in}}%
\pgfpathlineto{\pgfqpoint{1.154728in}{0.483235in}}%
\pgfpathlineto{\pgfqpoint{1.156997in}{0.519731in}}%
\pgfpathlineto{\pgfqpoint{1.159267in}{0.516276in}}%
\pgfpathlineto{\pgfqpoint{1.161536in}{0.483116in}}%
\pgfpathlineto{\pgfqpoint{1.163806in}{0.511582in}}%
\pgfpathlineto{\pgfqpoint{1.166075in}{0.453289in}}%
\pgfpathlineto{\pgfqpoint{1.168345in}{0.492313in}}%
\pgfpathlineto{\pgfqpoint{1.172884in}{0.447659in}}%
\pgfpathlineto{\pgfqpoint{1.175153in}{0.508075in}}%
\pgfpathlineto{\pgfqpoint{1.177423in}{0.432571in}}%
\pgfpathlineto{\pgfqpoint{1.179692in}{0.475815in}}%
\pgfpathlineto{\pgfqpoint{1.181962in}{0.501570in}}%
\pgfpathlineto{\pgfqpoint{1.184231in}{0.445847in}}%
\pgfpathlineto{\pgfqpoint{1.188770in}{0.489457in}}%
\pgfpathlineto{\pgfqpoint{1.191040in}{0.492436in}}%
\pgfpathlineto{\pgfqpoint{1.193309in}{0.372095in}}%
\pgfpathlineto{\pgfqpoint{1.195579in}{0.498735in}}%
\pgfpathlineto{\pgfqpoint{1.197848in}{0.508983in}}%
\pgfpathlineto{\pgfqpoint{1.200118in}{0.412938in}}%
\pgfpathlineto{\pgfqpoint{1.202387in}{0.483899in}}%
\pgfpathlineto{\pgfqpoint{1.204657in}{0.502950in}}%
\pgfpathlineto{\pgfqpoint{1.206926in}{0.453456in}}%
\pgfpathlineto{\pgfqpoint{1.209196in}{0.471881in}}%
\pgfpathlineto{\pgfqpoint{1.211465in}{0.475267in}}%
\pgfpathlineto{\pgfqpoint{1.213735in}{0.430161in}}%
\pgfpathlineto{\pgfqpoint{1.216004in}{0.483463in}}%
\pgfpathlineto{\pgfqpoint{1.218274in}{0.378325in}}%
\pgfpathlineto{\pgfqpoint{1.220543in}{0.478500in}}%
\pgfpathlineto{\pgfqpoint{1.222813in}{0.524758in}}%
\pgfpathlineto{\pgfqpoint{1.225082in}{0.446427in}}%
\pgfpathlineto{\pgfqpoint{1.227351in}{0.507708in}}%
\pgfpathlineto{\pgfqpoint{1.229621in}{0.497125in}}%
\pgfpathlineto{\pgfqpoint{1.231890in}{0.458285in}}%
\pgfpathlineto{\pgfqpoint{1.234160in}{0.457252in}}%
\pgfpathlineto{\pgfqpoint{1.236429in}{0.448671in}}%
\pgfpathlineto{\pgfqpoint{1.238699in}{0.432060in}}%
\pgfpathlineto{\pgfqpoint{1.240968in}{0.461139in}}%
\pgfpathlineto{\pgfqpoint{1.243238in}{0.453122in}}%
\pgfpathlineto{\pgfqpoint{1.247777in}{0.493193in}}%
\pgfpathlineto{\pgfqpoint{1.250046in}{0.473445in}}%
\pgfpathlineto{\pgfqpoint{1.252316in}{0.483186in}}%
\pgfpathlineto{\pgfqpoint{1.254585in}{0.464530in}}%
\pgfpathlineto{\pgfqpoint{1.256855in}{0.391226in}}%
\pgfpathlineto{\pgfqpoint{1.259124in}{0.486329in}}%
\pgfpathlineto{\pgfqpoint{1.261394in}{0.448409in}}%
\pgfpathlineto{\pgfqpoint{1.263663in}{0.426876in}}%
\pgfpathlineto{\pgfqpoint{1.265933in}{0.487912in}}%
\pgfpathlineto{\pgfqpoint{1.268202in}{0.467259in}}%
\pgfpathlineto{\pgfqpoint{1.270472in}{0.499357in}}%
\pgfpathlineto{\pgfqpoint{1.272741in}{0.486950in}}%
\pgfpathlineto{\pgfqpoint{1.275011in}{0.504817in}}%
\pgfpathlineto{\pgfqpoint{1.277280in}{0.437975in}}%
\pgfpathlineto{\pgfqpoint{1.279550in}{0.496227in}}%
\pgfpathlineto{\pgfqpoint{1.284089in}{0.451607in}}%
\pgfpathlineto{\pgfqpoint{1.286358in}{0.511717in}}%
\pgfpathlineto{\pgfqpoint{1.288628in}{0.485692in}}%
\pgfpathlineto{\pgfqpoint{1.290897in}{0.437041in}}%
\pgfpathlineto{\pgfqpoint{1.293166in}{0.509547in}}%
\pgfpathlineto{\pgfqpoint{1.295436in}{0.504316in}}%
\pgfpathlineto{\pgfqpoint{1.297705in}{0.389567in}}%
\pgfpathlineto{\pgfqpoint{1.299975in}{0.497927in}}%
\pgfpathlineto{\pgfqpoint{1.302244in}{0.451840in}}%
\pgfpathlineto{\pgfqpoint{1.304514in}{0.471616in}}%
\pgfpathlineto{\pgfqpoint{1.306783in}{0.522557in}}%
\pgfpathlineto{\pgfqpoint{1.309053in}{0.379123in}}%
\pgfpathlineto{\pgfqpoint{1.311322in}{0.501016in}}%
\pgfpathlineto{\pgfqpoint{1.313592in}{0.505644in}}%
\pgfpathlineto{\pgfqpoint{1.315861in}{0.395686in}}%
\pgfpathlineto{\pgfqpoint{1.318131in}{0.370155in}}%
\pgfpathlineto{\pgfqpoint{1.320400in}{0.473892in}}%
\pgfpathlineto{\pgfqpoint{1.322670in}{0.450566in}}%
\pgfpathlineto{\pgfqpoint{1.324939in}{0.447734in}}%
\pgfpathlineto{\pgfqpoint{1.327209in}{0.485068in}}%
\pgfpathlineto{\pgfqpoint{1.329478in}{0.485374in}}%
\pgfpathlineto{\pgfqpoint{1.331748in}{0.428599in}}%
\pgfpathlineto{\pgfqpoint{1.334017in}{0.451425in}}%
\pgfpathlineto{\pgfqpoint{1.336287in}{0.462401in}}%
\pgfpathlineto{\pgfqpoint{1.338556in}{0.406607in}}%
\pgfpathlineto{\pgfqpoint{1.340826in}{0.424730in}}%
\pgfpathlineto{\pgfqpoint{1.343095in}{0.411090in}}%
\pgfpathlineto{\pgfqpoint{1.345365in}{0.447770in}}%
\pgfpathlineto{\pgfqpoint{1.347634in}{0.429163in}}%
\pgfpathlineto{\pgfqpoint{1.349904in}{0.461196in}}%
\pgfpathlineto{\pgfqpoint{1.352173in}{0.476295in}}%
\pgfpathlineto{\pgfqpoint{1.354443in}{0.516396in}}%
\pgfpathlineto{\pgfqpoint{1.356712in}{0.476550in}}%
\pgfpathlineto{\pgfqpoint{1.358981in}{0.510040in}}%
\pgfpathlineto{\pgfqpoint{1.361251in}{0.473876in}}%
\pgfpathlineto{\pgfqpoint{1.363520in}{0.486432in}}%
\pgfpathlineto{\pgfqpoint{1.365790in}{0.473622in}}%
\pgfpathlineto{\pgfqpoint{1.368059in}{0.391936in}}%
\pgfpathlineto{\pgfqpoint{1.370329in}{0.493931in}}%
\pgfpathlineto{\pgfqpoint{1.372598in}{0.453592in}}%
\pgfpathlineto{\pgfqpoint{1.374868in}{0.433987in}}%
\pgfpathlineto{\pgfqpoint{1.377137in}{0.490558in}}%
\pgfpathlineto{\pgfqpoint{1.379407in}{0.449405in}}%
\pgfpathlineto{\pgfqpoint{1.381676in}{0.463860in}}%
\pgfpathlineto{\pgfqpoint{1.383946in}{0.505181in}}%
\pgfpathlineto{\pgfqpoint{1.386215in}{0.458074in}}%
\pgfpathlineto{\pgfqpoint{1.388485in}{0.460903in}}%
\pgfpathlineto{\pgfqpoint{1.390754in}{0.513698in}}%
\pgfpathlineto{\pgfqpoint{1.393024in}{0.409141in}}%
\pgfpathlineto{\pgfqpoint{1.395293in}{0.488226in}}%
\pgfpathlineto{\pgfqpoint{1.397563in}{0.503764in}}%
\pgfpathlineto{\pgfqpoint{1.399832in}{0.426732in}}%
\pgfpathlineto{\pgfqpoint{1.404371in}{0.485489in}}%
\pgfpathlineto{\pgfqpoint{1.406641in}{0.487713in}}%
\pgfpathlineto{\pgfqpoint{1.408910in}{0.406224in}}%
\pgfpathlineto{\pgfqpoint{1.411180in}{0.496383in}}%
\pgfpathlineto{\pgfqpoint{1.413449in}{0.444590in}}%
\pgfpathlineto{\pgfqpoint{1.415719in}{0.474915in}}%
\pgfpathlineto{\pgfqpoint{1.417988in}{0.522165in}}%
\pgfpathlineto{\pgfqpoint{1.420258in}{0.466914in}}%
\pgfpathlineto{\pgfqpoint{1.422527in}{0.461940in}}%
\pgfpathlineto{\pgfqpoint{1.424796in}{0.510338in}}%
\pgfpathlineto{\pgfqpoint{1.427066in}{0.488471in}}%
\pgfpathlineto{\pgfqpoint{1.429335in}{0.444291in}}%
\pgfpathlineto{\pgfqpoint{1.431605in}{0.482860in}}%
\pgfpathlineto{\pgfqpoint{1.433874in}{0.369554in}}%
\pgfpathlineto{\pgfqpoint{1.436144in}{0.477253in}}%
\pgfpathlineto{\pgfqpoint{1.438413in}{0.524938in}}%
\pgfpathlineto{\pgfqpoint{1.440683in}{0.457047in}}%
\pgfpathlineto{\pgfqpoint{1.442952in}{0.510298in}}%
\pgfpathlineto{\pgfqpoint{1.445222in}{0.492651in}}%
\pgfpathlineto{\pgfqpoint{1.447491in}{0.467832in}}%
\pgfpathlineto{\pgfqpoint{1.449761in}{0.461018in}}%
\pgfpathlineto{\pgfqpoint{1.452030in}{0.428280in}}%
\pgfpathlineto{\pgfqpoint{1.454300in}{0.454684in}}%
\pgfpathlineto{\pgfqpoint{1.456569in}{0.457667in}}%
\pgfpathlineto{\pgfqpoint{1.458839in}{0.432278in}}%
\pgfpathlineto{\pgfqpoint{1.461108in}{0.472064in}}%
\pgfpathlineto{\pgfqpoint{1.463378in}{0.477384in}}%
\pgfpathlineto{\pgfqpoint{1.465647in}{0.505502in}}%
\pgfpathlineto{\pgfqpoint{1.467917in}{0.473072in}}%
\pgfpathlineto{\pgfqpoint{1.470186in}{0.459246in}}%
\pgfpathlineto{\pgfqpoint{1.472456in}{0.454465in}}%
\pgfpathlineto{\pgfqpoint{1.474725in}{0.474440in}}%
\pgfpathlineto{\pgfqpoint{1.476995in}{0.443388in}}%
\pgfpathlineto{\pgfqpoint{1.479264in}{0.449653in}}%
\pgfpathlineto{\pgfqpoint{1.481534in}{0.493561in}}%
\pgfpathlineto{\pgfqpoint{1.483803in}{0.457206in}}%
\pgfpathlineto{\pgfqpoint{1.486073in}{0.494251in}}%
\pgfpathlineto{\pgfqpoint{1.488342in}{0.487005in}}%
\pgfpathlineto{\pgfqpoint{1.490611in}{0.503309in}}%
\pgfpathlineto{\pgfqpoint{1.492881in}{0.417775in}}%
\pgfpathlineto{\pgfqpoint{1.495150in}{0.496890in}}%
\pgfpathlineto{\pgfqpoint{1.497420in}{0.467797in}}%
\pgfpathlineto{\pgfqpoint{1.499689in}{0.460886in}}%
\pgfpathlineto{\pgfqpoint{1.501959in}{0.515703in}}%
\pgfpathlineto{\pgfqpoint{1.504228in}{0.412751in}}%
\pgfpathlineto{\pgfqpoint{1.506498in}{0.487070in}}%
\pgfpathlineto{\pgfqpoint{1.508767in}{0.520785in}}%
\pgfpathlineto{\pgfqpoint{1.511037in}{0.402325in}}%
\pgfpathlineto{\pgfqpoint{1.513306in}{0.488295in}}%
\pgfpathlineto{\pgfqpoint{1.515576in}{0.526714in}}%
\pgfpathlineto{\pgfqpoint{1.517845in}{0.425296in}}%
\pgfpathlineto{\pgfqpoint{1.522384in}{0.525965in}}%
\pgfpathlineto{\pgfqpoint{1.524654in}{0.419557in}}%
\pgfpathlineto{\pgfqpoint{1.526923in}{0.506065in}}%
\pgfpathlineto{\pgfqpoint{1.529193in}{0.503437in}}%
\pgfpathlineto{\pgfqpoint{1.531462in}{0.429394in}}%
\pgfpathlineto{\pgfqpoint{1.533732in}{0.406573in}}%
\pgfpathlineto{\pgfqpoint{1.536001in}{0.463976in}}%
\pgfpathlineto{\pgfqpoint{1.538271in}{0.428132in}}%
\pgfpathlineto{\pgfqpoint{1.542810in}{0.476588in}}%
\pgfpathlineto{\pgfqpoint{1.545079in}{0.417499in}}%
\pgfpathlineto{\pgfqpoint{1.547349in}{0.480279in}}%
\pgfpathlineto{\pgfqpoint{1.549618in}{0.520178in}}%
\pgfpathlineto{\pgfqpoint{1.551888in}{0.463865in}}%
\pgfpathlineto{\pgfqpoint{1.554157in}{0.507699in}}%
\pgfpathlineto{\pgfqpoint{1.556427in}{0.532129in}}%
\pgfpathlineto{\pgfqpoint{1.558696in}{0.469716in}}%
\pgfpathlineto{\pgfqpoint{1.560965in}{0.495889in}}%
\pgfpathlineto{\pgfqpoint{1.563235in}{0.530422in}}%
\pgfpathlineto{\pgfqpoint{1.565504in}{0.473610in}}%
\pgfpathlineto{\pgfqpoint{1.567774in}{0.476310in}}%
\pgfpathlineto{\pgfqpoint{1.570043in}{0.515887in}}%
\pgfpathlineto{\pgfqpoint{1.572313in}{0.479943in}}%
\pgfpathlineto{\pgfqpoint{1.574582in}{0.510184in}}%
\pgfpathlineto{\pgfqpoint{1.576852in}{0.464207in}}%
\pgfpathlineto{\pgfqpoint{1.579121in}{0.490139in}}%
\pgfpathlineto{\pgfqpoint{1.581391in}{0.472557in}}%
\pgfpathlineto{\pgfqpoint{1.583660in}{0.428663in}}%
\pgfpathlineto{\pgfqpoint{1.585930in}{0.501183in}}%
\pgfpathlineto{\pgfqpoint{1.588199in}{0.445471in}}%
\pgfpathlineto{\pgfqpoint{1.590469in}{0.459008in}}%
\pgfpathlineto{\pgfqpoint{1.592738in}{0.498111in}}%
\pgfpathlineto{\pgfqpoint{1.595008in}{0.458422in}}%
\pgfpathlineto{\pgfqpoint{1.597277in}{0.457222in}}%
\pgfpathlineto{\pgfqpoint{1.599547in}{0.505217in}}%
\pgfpathlineto{\pgfqpoint{1.601816in}{0.423605in}}%
\pgfpathlineto{\pgfqpoint{1.604086in}{0.462370in}}%
\pgfpathlineto{\pgfqpoint{1.606355in}{0.512703in}}%
\pgfpathlineto{\pgfqpoint{1.608625in}{0.378236in}}%
\pgfpathlineto{\pgfqpoint{1.610894in}{0.493814in}}%
\pgfpathlineto{\pgfqpoint{1.613164in}{0.505340in}}%
\pgfpathlineto{\pgfqpoint{1.615433in}{0.380602in}}%
\pgfpathlineto{\pgfqpoint{1.617703in}{0.447678in}}%
\pgfpathlineto{\pgfqpoint{1.619972in}{0.479711in}}%
\pgfpathlineto{\pgfqpoint{1.622242in}{0.482042in}}%
\pgfpathlineto{\pgfqpoint{1.624511in}{0.420925in}}%
\pgfpathlineto{\pgfqpoint{1.626780in}{0.492522in}}%
\pgfpathlineto{\pgfqpoint{1.629050in}{0.428432in}}%
\pgfpathlineto{\pgfqpoint{1.633589in}{0.523582in}}%
\pgfpathlineto{\pgfqpoint{1.635858in}{0.426292in}}%
\pgfpathlineto{\pgfqpoint{1.638128in}{0.503899in}}%
\pgfpathlineto{\pgfqpoint{1.640397in}{0.501833in}}%
\pgfpathlineto{\pgfqpoint{1.642667in}{0.441153in}}%
\pgfpathlineto{\pgfqpoint{1.644936in}{0.497846in}}%
\pgfpathlineto{\pgfqpoint{1.647206in}{0.515588in}}%
\pgfpathlineto{\pgfqpoint{1.649475in}{0.459282in}}%
\pgfpathlineto{\pgfqpoint{1.651745in}{0.513667in}}%
\pgfpathlineto{\pgfqpoint{1.654014in}{0.528728in}}%
\pgfpathlineto{\pgfqpoint{1.656284in}{0.465840in}}%
\pgfpathlineto{\pgfqpoint{1.658553in}{0.513179in}}%
\pgfpathlineto{\pgfqpoint{1.660823in}{0.487351in}}%
\pgfpathlineto{\pgfqpoint{1.665362in}{0.464105in}}%
\pgfpathlineto{\pgfqpoint{1.667631in}{0.378404in}}%
\pgfpathlineto{\pgfqpoint{1.669901in}{0.471676in}}%
\pgfpathlineto{\pgfqpoint{1.672170in}{0.454003in}}%
\pgfpathlineto{\pgfqpoint{1.674440in}{0.336850in}}%
\pgfpathlineto{\pgfqpoint{1.676709in}{0.480409in}}%
\pgfpathlineto{\pgfqpoint{1.678979in}{0.473394in}}%
\pgfpathlineto{\pgfqpoint{1.681248in}{0.504004in}}%
\pgfpathlineto{\pgfqpoint{1.683518in}{0.485127in}}%
\pgfpathlineto{\pgfqpoint{1.685787in}{0.505977in}}%
\pgfpathlineto{\pgfqpoint{1.688057in}{0.519127in}}%
\pgfpathlineto{\pgfqpoint{1.690326in}{0.489506in}}%
\pgfpathlineto{\pgfqpoint{1.692595in}{0.486286in}}%
\pgfpathlineto{\pgfqpoint{1.694865in}{0.513062in}}%
\pgfpathlineto{\pgfqpoint{1.697134in}{0.496734in}}%
\pgfpathlineto{\pgfqpoint{1.699404in}{0.454736in}}%
\pgfpathlineto{\pgfqpoint{1.701673in}{0.492996in}}%
\pgfpathlineto{\pgfqpoint{1.703943in}{0.484895in}}%
\pgfpathlineto{\pgfqpoint{1.706212in}{0.500424in}}%
\pgfpathlineto{\pgfqpoint{1.708482in}{0.384399in}}%
\pgfpathlineto{\pgfqpoint{1.710751in}{0.496607in}}%
\pgfpathlineto{\pgfqpoint{1.713021in}{0.462134in}}%
\pgfpathlineto{\pgfqpoint{1.715290in}{0.466914in}}%
\pgfpathlineto{\pgfqpoint{1.717560in}{0.518997in}}%
\pgfpathlineto{\pgfqpoint{1.719829in}{0.372238in}}%
\pgfpathlineto{\pgfqpoint{1.722099in}{0.493947in}}%
\pgfpathlineto{\pgfqpoint{1.724368in}{0.506415in}}%
\pgfpathlineto{\pgfqpoint{1.726638in}{0.368556in}}%
\pgfpathlineto{\pgfqpoint{1.728907in}{0.377192in}}%
\pgfpathlineto{\pgfqpoint{1.731177in}{0.482674in}}%
\pgfpathlineto{\pgfqpoint{1.733446in}{0.467145in}}%
\pgfpathlineto{\pgfqpoint{1.735716in}{0.429579in}}%
\pgfpathlineto{\pgfqpoint{1.737985in}{0.484365in}}%
\pgfpathlineto{\pgfqpoint{1.739837in}{0.315625in}}%
\pgfpathmoveto{\pgfqpoint{1.740814in}{0.315625in}}%
\pgfpathlineto{\pgfqpoint{1.742524in}{0.432197in}}%
\pgfpathlineto{\pgfqpoint{1.744794in}{0.496448in}}%
\pgfpathlineto{\pgfqpoint{1.747063in}{0.446605in}}%
\pgfpathlineto{\pgfqpoint{1.749333in}{0.419217in}}%
\pgfpathlineto{\pgfqpoint{1.751602in}{0.447259in}}%
\pgfpathlineto{\pgfqpoint{1.753744in}{0.315625in}}%
\pgfpathmoveto{\pgfqpoint{1.753996in}{0.315625in}}%
\pgfpathlineto{\pgfqpoint{1.756141in}{0.450530in}}%
\pgfpathlineto{\pgfqpoint{1.758410in}{0.462779in}}%
\pgfpathlineto{\pgfqpoint{1.760680in}{0.443200in}}%
\pgfpathlineto{\pgfqpoint{1.765219in}{0.518698in}}%
\pgfpathlineto{\pgfqpoint{1.767488in}{0.470841in}}%
\pgfpathlineto{\pgfqpoint{1.769758in}{0.509199in}}%
\pgfpathlineto{\pgfqpoint{1.772027in}{0.482082in}}%
\pgfpathlineto{\pgfqpoint{1.774297in}{0.481074in}}%
\pgfpathlineto{\pgfqpoint{1.776566in}{0.513824in}}%
\pgfpathlineto{\pgfqpoint{1.778836in}{0.500095in}}%
\pgfpathlineto{\pgfqpoint{1.781105in}{0.482601in}}%
\pgfpathlineto{\pgfqpoint{1.783375in}{0.519702in}}%
\pgfpathlineto{\pgfqpoint{1.785644in}{0.517110in}}%
\pgfpathlineto{\pgfqpoint{1.787914in}{0.482713in}}%
\pgfpathlineto{\pgfqpoint{1.790183in}{0.511876in}}%
\pgfpathlineto{\pgfqpoint{1.792453in}{0.455479in}}%
\pgfpathlineto{\pgfqpoint{1.794722in}{0.491912in}}%
\pgfpathlineto{\pgfqpoint{1.799261in}{0.445663in}}%
\pgfpathlineto{\pgfqpoint{1.801531in}{0.507065in}}%
\pgfpathlineto{\pgfqpoint{1.803800in}{0.434406in}}%
\pgfpathlineto{\pgfqpoint{1.806070in}{0.473967in}}%
\pgfpathlineto{\pgfqpoint{1.808339in}{0.501009in}}%
\pgfpathlineto{\pgfqpoint{1.810609in}{0.447952in}}%
\pgfpathlineto{\pgfqpoint{1.815148in}{0.489631in}}%
\pgfpathlineto{\pgfqpoint{1.817417in}{0.493426in}}%
\pgfpathlineto{\pgfqpoint{1.819687in}{0.356547in}}%
\pgfpathlineto{\pgfqpoint{1.821956in}{0.498851in}}%
\pgfpathlineto{\pgfqpoint{1.824226in}{0.509589in}}%
\pgfpathlineto{\pgfqpoint{1.826495in}{0.418730in}}%
\pgfpathlineto{\pgfqpoint{1.828764in}{0.484627in}}%
\pgfpathlineto{\pgfqpoint{1.831034in}{0.503939in}}%
\pgfpathlineto{\pgfqpoint{1.833303in}{0.456596in}}%
\pgfpathlineto{\pgfqpoint{1.835573in}{0.472976in}}%
\pgfpathlineto{\pgfqpoint{1.837842in}{0.477356in}}%
\pgfpathlineto{\pgfqpoint{1.840112in}{0.428204in}}%
\pgfpathlineto{\pgfqpoint{1.842381in}{0.484538in}}%
\pgfpathlineto{\pgfqpoint{1.844651in}{0.390686in}}%
\pgfpathlineto{\pgfqpoint{1.846920in}{0.478317in}}%
\pgfpathlineto{\pgfqpoint{1.849190in}{0.524799in}}%
\pgfpathlineto{\pgfqpoint{1.851459in}{0.444147in}}%
\pgfpathlineto{\pgfqpoint{1.853729in}{0.507398in}}%
\pgfpathlineto{\pgfqpoint{1.855998in}{0.497880in}}%
\pgfpathlineto{\pgfqpoint{1.858268in}{0.456175in}}%
\pgfpathlineto{\pgfqpoint{1.860537in}{0.455587in}}%
\pgfpathlineto{\pgfqpoint{1.862807in}{0.451064in}}%
\pgfpathlineto{\pgfqpoint{1.865076in}{0.424861in}}%
\pgfpathlineto{\pgfqpoint{1.867346in}{0.458837in}}%
\pgfpathlineto{\pgfqpoint{1.869615in}{0.455172in}}%
\pgfpathlineto{\pgfqpoint{1.874154in}{0.492057in}}%
\pgfpathlineto{\pgfqpoint{1.876424in}{0.474829in}}%
\pgfpathlineto{\pgfqpoint{1.878693in}{0.482217in}}%
\pgfpathlineto{\pgfqpoint{1.880963in}{0.463792in}}%
\pgfpathlineto{\pgfqpoint{1.883232in}{0.381904in}}%
\pgfpathlineto{\pgfqpoint{1.885502in}{0.484380in}}%
\pgfpathlineto{\pgfqpoint{1.890041in}{0.420021in}}%
\pgfpathlineto{\pgfqpoint{1.892310in}{0.486746in}}%
\pgfpathlineto{\pgfqpoint{1.894579in}{0.468139in}}%
\pgfpathlineto{\pgfqpoint{1.896849in}{0.500560in}}%
\pgfpathlineto{\pgfqpoint{1.899118in}{0.486580in}}%
\pgfpathlineto{\pgfqpoint{1.901388in}{0.505255in}}%
\pgfpathlineto{\pgfqpoint{1.903657in}{0.441155in}}%
\pgfpathlineto{\pgfqpoint{1.905927in}{0.495863in}}%
\pgfpathlineto{\pgfqpoint{1.910466in}{0.449727in}}%
\pgfpathlineto{\pgfqpoint{1.912735in}{0.510939in}}%
\pgfpathlineto{\pgfqpoint{1.915005in}{0.485974in}}%
\pgfpathlineto{\pgfqpoint{1.917274in}{0.434807in}}%
\pgfpathlineto{\pgfqpoint{1.919544in}{0.509204in}}%
\pgfpathlineto{\pgfqpoint{1.921813in}{0.504632in}}%
\pgfpathlineto{\pgfqpoint{1.924083in}{0.381639in}}%
\pgfpathlineto{\pgfqpoint{1.926352in}{0.498087in}}%
\pgfpathlineto{\pgfqpoint{1.928622in}{0.453465in}}%
\pgfpathlineto{\pgfqpoint{1.930891in}{0.471030in}}%
\pgfpathlineto{\pgfqpoint{1.933161in}{0.522182in}}%
\pgfpathlineto{\pgfqpoint{1.935430in}{0.367862in}}%
\pgfpathlineto{\pgfqpoint{1.937700in}{0.500331in}}%
\pgfpathlineto{\pgfqpoint{1.939969in}{0.505824in}}%
\pgfpathlineto{\pgfqpoint{1.942239in}{0.385852in}}%
\pgfpathlineto{\pgfqpoint{1.944508in}{0.343377in}}%
\pgfpathlineto{\pgfqpoint{1.946778in}{0.475006in}}%
\pgfpathlineto{\pgfqpoint{1.949047in}{0.454236in}}%
\pgfpathlineto{\pgfqpoint{1.951317in}{0.446573in}}%
\pgfpathlineto{\pgfqpoint{1.953586in}{0.486190in}}%
\pgfpathlineto{\pgfqpoint{1.955856in}{0.487242in}}%
\pgfpathlineto{\pgfqpoint{1.958125in}{0.426015in}}%
\pgfpathlineto{\pgfqpoint{1.960394in}{0.453905in}}%
\pgfpathlineto{\pgfqpoint{1.962664in}{0.465995in}}%
\pgfpathlineto{\pgfqpoint{1.964933in}{0.399562in}}%
\pgfpathlineto{\pgfqpoint{1.967203in}{0.428652in}}%
\pgfpathlineto{\pgfqpoint{1.969472in}{0.397208in}}%
\pgfpathlineto{\pgfqpoint{1.971742in}{0.447566in}}%
\pgfpathlineto{\pgfqpoint{1.974011in}{0.434510in}}%
\pgfpathlineto{\pgfqpoint{1.976281in}{0.458786in}}%
\pgfpathlineto{\pgfqpoint{1.978550in}{0.476642in}}%
\pgfpathlineto{\pgfqpoint{1.980820in}{0.517100in}}%
\pgfpathlineto{\pgfqpoint{1.983089in}{0.475674in}}%
\pgfpathlineto{\pgfqpoint{1.985359in}{0.510136in}}%
\pgfpathlineto{\pgfqpoint{1.987628in}{0.475341in}}%
\pgfpathlineto{\pgfqpoint{1.989898in}{0.485600in}}%
\pgfpathlineto{\pgfqpoint{1.992167in}{0.473380in}}%
\pgfpathlineto{\pgfqpoint{1.994437in}{0.381868in}}%
\pgfpathlineto{\pgfqpoint{1.996706in}{0.492429in}}%
\pgfpathlineto{\pgfqpoint{1.998976in}{0.454300in}}%
\pgfpathlineto{\pgfqpoint{2.001245in}{0.428457in}}%
\pgfpathlineto{\pgfqpoint{2.003515in}{0.488318in}}%
\pgfpathlineto{\pgfqpoint{2.005784in}{0.450089in}}%
\pgfpathlineto{\pgfqpoint{2.008054in}{0.462031in}}%
\pgfpathlineto{\pgfqpoint{2.010323in}{0.503695in}}%
\pgfpathlineto{\pgfqpoint{2.012593in}{0.458477in}}%
\pgfpathlineto{\pgfqpoint{2.014862in}{0.459598in}}%
\pgfpathlineto{\pgfqpoint{2.017132in}{0.512831in}}%
\pgfpathlineto{\pgfqpoint{2.019401in}{0.412619in}}%
\pgfpathlineto{\pgfqpoint{2.021671in}{0.486986in}}%
\pgfpathlineto{\pgfqpoint{2.023940in}{0.503386in}}%
\pgfpathlineto{\pgfqpoint{2.026209in}{0.430197in}}%
\pgfpathlineto{\pgfqpoint{2.028479in}{0.462887in}}%
\pgfpathlineto{\pgfqpoint{2.030748in}{0.485983in}}%
\pgfpathlineto{\pgfqpoint{2.033018in}{0.488957in}}%
\pgfpathlineto{\pgfqpoint{2.035287in}{0.400877in}}%
\pgfpathlineto{\pgfqpoint{2.037557in}{0.496724in}}%
\pgfpathlineto{\pgfqpoint{2.039826in}{0.446972in}}%
\pgfpathlineto{\pgfqpoint{2.042096in}{0.474345in}}%
\pgfpathlineto{\pgfqpoint{2.044365in}{0.521954in}}%
\pgfpathlineto{\pgfqpoint{2.046635in}{0.468538in}}%
\pgfpathlineto{\pgfqpoint{2.048904in}{0.461389in}}%
\pgfpathlineto{\pgfqpoint{2.051174in}{0.510619in}}%
\pgfpathlineto{\pgfqpoint{2.053443in}{0.489787in}}%
\pgfpathlineto{\pgfqpoint{2.055713in}{0.443340in}}%
\pgfpathlineto{\pgfqpoint{2.057982in}{0.484101in}}%
\pgfpathlineto{\pgfqpoint{2.060252in}{0.345152in}}%
\pgfpathlineto{\pgfqpoint{2.062521in}{0.477244in}}%
\pgfpathlineto{\pgfqpoint{2.064791in}{0.525118in}}%
\pgfpathlineto{\pgfqpoint{2.067060in}{0.455396in}}%
\pgfpathlineto{\pgfqpoint{2.069330in}{0.510101in}}%
\pgfpathlineto{\pgfqpoint{2.071599in}{0.493522in}}%
\pgfpathlineto{\pgfqpoint{2.073869in}{0.466260in}}%
\pgfpathlineto{\pgfqpoint{2.076138in}{0.459721in}}%
\pgfpathlineto{\pgfqpoint{2.078408in}{0.432016in}}%
\pgfpathlineto{\pgfqpoint{2.080677in}{0.450549in}}%
\pgfpathlineto{\pgfqpoint{2.082947in}{0.457648in}}%
\pgfpathlineto{\pgfqpoint{2.085216in}{0.437322in}}%
\pgfpathlineto{\pgfqpoint{2.087486in}{0.470183in}}%
\pgfpathlineto{\pgfqpoint{2.089755in}{0.477790in}}%
\pgfpathlineto{\pgfqpoint{2.092024in}{0.506668in}}%
\pgfpathlineto{\pgfqpoint{2.094294in}{0.470650in}}%
\pgfpathlineto{\pgfqpoint{2.096563in}{0.459967in}}%
\pgfpathlineto{\pgfqpoint{2.098833in}{0.458421in}}%
\pgfpathlineto{\pgfqpoint{2.101102in}{0.471340in}}%
\pgfpathlineto{\pgfqpoint{2.103372in}{0.444377in}}%
\pgfpathlineto{\pgfqpoint{2.105641in}{0.445679in}}%
\pgfpathlineto{\pgfqpoint{2.107911in}{0.492633in}}%
\pgfpathlineto{\pgfqpoint{2.110180in}{0.458496in}}%
\pgfpathlineto{\pgfqpoint{2.112450in}{0.495621in}}%
\pgfpathlineto{\pgfqpoint{2.114719in}{0.486857in}}%
\pgfpathlineto{\pgfqpoint{2.116989in}{0.503882in}}%
\pgfpathlineto{\pgfqpoint{2.119258in}{0.422410in}}%
\pgfpathlineto{\pgfqpoint{2.121528in}{0.496702in}}%
\pgfpathlineto{\pgfqpoint{2.123797in}{0.468525in}}%
\pgfpathlineto{\pgfqpoint{2.126067in}{0.459606in}}%
\pgfpathlineto{\pgfqpoint{2.128336in}{0.515064in}}%
\pgfpathlineto{\pgfqpoint{2.130606in}{0.416530in}}%
\pgfpathlineto{\pgfqpoint{2.132875in}{0.485866in}}%
\pgfpathlineto{\pgfqpoint{2.135145in}{0.519893in}}%
\pgfpathlineto{\pgfqpoint{2.137414in}{0.407234in}}%
\pgfpathlineto{\pgfqpoint{2.139684in}{0.487646in}}%
\pgfpathlineto{\pgfqpoint{2.141953in}{0.526074in}}%
\pgfpathlineto{\pgfqpoint{2.144223in}{0.428475in}}%
\pgfpathlineto{\pgfqpoint{2.148762in}{0.525663in}}%
\pgfpathlineto{\pgfqpoint{2.151031in}{0.415635in}}%
\pgfpathlineto{\pgfqpoint{2.153301in}{0.505571in}}%
\pgfpathlineto{\pgfqpoint{2.155570in}{0.503871in}}%
\pgfpathlineto{\pgfqpoint{2.157840in}{0.425281in}}%
\pgfpathlineto{\pgfqpoint{2.160109in}{0.398013in}}%
\pgfpathlineto{\pgfqpoint{2.162378in}{0.465453in}}%
\pgfpathlineto{\pgfqpoint{2.164648in}{0.434521in}}%
\pgfpathlineto{\pgfqpoint{2.166917in}{0.448990in}}%
\pgfpathlineto{\pgfqpoint{2.169187in}{0.478118in}}%
\pgfpathlineto{\pgfqpoint{2.171456in}{0.410539in}}%
\pgfpathlineto{\pgfqpoint{2.173726in}{0.480281in}}%
\pgfpathlineto{\pgfqpoint{2.175995in}{0.520627in}}%
\pgfpathlineto{\pgfqpoint{2.178265in}{0.462337in}}%
\pgfpathlineto{\pgfqpoint{2.180534in}{0.507577in}}%
\pgfpathlineto{\pgfqpoint{2.182804in}{0.532392in}}%
\pgfpathlineto{\pgfqpoint{2.185073in}{0.468298in}}%
\pgfpathlineto{\pgfqpoint{2.187343in}{0.496059in}}%
\pgfpathlineto{\pgfqpoint{2.189612in}{0.530913in}}%
\pgfpathlineto{\pgfqpoint{2.191882in}{0.471970in}}%
\pgfpathlineto{\pgfqpoint{2.194151in}{0.476805in}}%
\pgfpathlineto{\pgfqpoint{2.196421in}{0.516665in}}%
\pgfpathlineto{\pgfqpoint{2.198690in}{0.479327in}}%
\pgfpathlineto{\pgfqpoint{2.200960in}{0.510401in}}%
\pgfpathlineto{\pgfqpoint{2.203229in}{0.465974in}}%
\pgfpathlineto{\pgfqpoint{2.205499in}{0.489514in}}%
\pgfpathlineto{\pgfqpoint{2.207768in}{0.472503in}}%
\pgfpathlineto{\pgfqpoint{2.210038in}{0.425114in}}%
\pgfpathlineto{\pgfqpoint{2.212307in}{0.499955in}}%
\pgfpathlineto{\pgfqpoint{2.214577in}{0.446565in}}%
\pgfpathlineto{\pgfqpoint{2.216846in}{0.456133in}}%
\pgfpathlineto{\pgfqpoint{2.219116in}{0.497366in}}%
\pgfpathlineto{\pgfqpoint{2.221385in}{0.459812in}}%
\pgfpathlineto{\pgfqpoint{2.223655in}{0.461469in}}%
\pgfpathlineto{\pgfqpoint{2.225924in}{0.504287in}}%
\pgfpathlineto{\pgfqpoint{2.228193in}{0.426591in}}%
\pgfpathlineto{\pgfqpoint{2.230463in}{0.458795in}}%
\pgfpathlineto{\pgfqpoint{2.232732in}{0.511612in}}%
\pgfpathlineto{\pgfqpoint{2.235002in}{0.386051in}}%
\pgfpathlineto{\pgfqpoint{2.237271in}{0.492714in}}%
\pgfpathlineto{\pgfqpoint{2.239541in}{0.505120in}}%
\pgfpathlineto{\pgfqpoint{2.241810in}{0.390455in}}%
\pgfpathlineto{\pgfqpoint{2.244080in}{0.451288in}}%
\pgfpathlineto{\pgfqpoint{2.246349in}{0.480453in}}%
\pgfpathlineto{\pgfqpoint{2.248619in}{0.483612in}}%
\pgfpathlineto{\pgfqpoint{2.250888in}{0.417776in}}%
\pgfpathlineto{\pgfqpoint{2.253158in}{0.493132in}}%
\pgfpathlineto{\pgfqpoint{2.255427in}{0.432182in}}%
\pgfpathlineto{\pgfqpoint{2.259966in}{0.523474in}}%
\pgfpathlineto{\pgfqpoint{2.262236in}{0.422498in}}%
\pgfpathlineto{\pgfqpoint{2.264505in}{0.503439in}}%
\pgfpathlineto{\pgfqpoint{2.266775in}{0.502436in}}%
\pgfpathlineto{\pgfqpoint{2.269044in}{0.437874in}}%
\pgfpathlineto{\pgfqpoint{2.271314in}{0.496654in}}%
\pgfpathlineto{\pgfqpoint{2.273583in}{0.516000in}}%
\pgfpathlineto{\pgfqpoint{2.275853in}{0.457601in}}%
\pgfpathlineto{\pgfqpoint{2.278122in}{0.513172in}}%
\pgfpathlineto{\pgfqpoint{2.280392in}{0.528908in}}%
\pgfpathlineto{\pgfqpoint{2.282661in}{0.464598in}}%
\pgfpathlineto{\pgfqpoint{2.284931in}{0.513056in}}%
\pgfpathlineto{\pgfqpoint{2.287200in}{0.488362in}}%
\pgfpathlineto{\pgfqpoint{2.291739in}{0.463154in}}%
\pgfpathlineto{\pgfqpoint{2.294008in}{0.389165in}}%
\pgfpathlineto{\pgfqpoint{2.296278in}{0.468904in}}%
\pgfpathlineto{\pgfqpoint{2.298547in}{0.454231in}}%
\pgfpathlineto{\pgfqpoint{2.300817in}{0.371014in}}%
\pgfpathlineto{\pgfqpoint{2.303086in}{0.454956in}}%
\pgfpathlineto{\pgfqpoint{2.303086in}{0.454956in}}%
\pgfusepath{stroke}%
\end{pgfscope}%
\begin{pgfscope}%
\pgfpathrectangle{\pgfqpoint{0.487500in}{0.320625in}}{\pgfqpoint{1.815586in}{1.164375in}}%
\pgfusepath{clip}%
\pgfsetbuttcap%
\pgfsetroundjoin%
\pgfsetlinewidth{0.752812pt}%
\definecolor{currentstroke}{rgb}{0.478431,0.243137,0.615686}%
\pgfsetstrokecolor{currentstroke}%
\pgfsetstrokeopacity{0.650000}%
\pgfsetdash{{2.775000pt}{1.200000pt}}{0.000000pt}%
\pgfpathmoveto{\pgfqpoint{0.487500in}{1.451682in}}%
\pgfpathlineto{\pgfqpoint{0.609866in}{0.315625in}}%
\pgfpathlineto{\pgfqpoint{0.609866in}{0.315625in}}%
\pgfusepath{stroke}%
\end{pgfscope}%
\begin{pgfscope}%
\pgfpathrectangle{\pgfqpoint{0.487500in}{0.320625in}}{\pgfqpoint{1.815586in}{1.164375in}}%
\pgfusepath{clip}%
\pgfsetbuttcap%
\pgfsetroundjoin%
\pgfsetlinewidth{0.803000pt}%
\definecolor{currentstroke}{rgb}{0.478431,0.243137,0.615686}%
\pgfsetstrokecolor{currentstroke}%
\pgfsetstrokeopacity{0.800000}%
\pgfsetdash{{0.800000pt}{1.320000pt}}{0.000000pt}%
\pgfpathmoveto{\pgfqpoint{0.487500in}{1.451682in}}%
\pgfpathlineto{\pgfqpoint{2.303086in}{1.451682in}}%
\pgfpathlineto{\pgfqpoint{2.303086in}{1.451682in}}%
\pgfusepath{stroke}%
\end{pgfscope}%
\begin{pgfscope}%
\pgfsetrectcap%
\pgfsetmiterjoin%
\pgfsetlinewidth{0.803000pt}%
\definecolor{currentstroke}{rgb}{0.000000,0.000000,0.000000}%
\pgfsetstrokecolor{currentstroke}%
\pgfsetdash{}{0pt}%
\pgfpathmoveto{\pgfqpoint{0.487500in}{0.320625in}}%
\pgfpathlineto{\pgfqpoint{0.487500in}{1.485000in}}%
\pgfusepath{stroke}%
\end{pgfscope}%
\begin{pgfscope}%
\pgfsetrectcap%
\pgfsetmiterjoin%
\pgfsetlinewidth{0.803000pt}%
\definecolor{currentstroke}{rgb}{0.000000,0.000000,0.000000}%
\pgfsetstrokecolor{currentstroke}%
\pgfsetdash{}{0pt}%
\pgfpathmoveto{\pgfqpoint{2.303086in}{0.320625in}}%
\pgfpathlineto{\pgfqpoint{2.303086in}{1.485000in}}%
\pgfusepath{stroke}%
\end{pgfscope}%
\begin{pgfscope}%
\pgfsetrectcap%
\pgfsetmiterjoin%
\pgfsetlinewidth{0.803000pt}%
\definecolor{currentstroke}{rgb}{0.000000,0.000000,0.000000}%
\pgfsetstrokecolor{currentstroke}%
\pgfsetdash{}{0pt}%
\pgfpathmoveto{\pgfqpoint{0.487500in}{0.320625in}}%
\pgfpathlineto{\pgfqpoint{2.303086in}{0.320625in}}%
\pgfusepath{stroke}%
\end{pgfscope}%
\begin{pgfscope}%
\pgfsetrectcap%
\pgfsetmiterjoin%
\pgfsetlinewidth{0.803000pt}%
\definecolor{currentstroke}{rgb}{0.000000,0.000000,0.000000}%
\pgfsetstrokecolor{currentstroke}%
\pgfsetdash{}{0pt}%
\pgfpathmoveto{\pgfqpoint{0.487500in}{1.485000in}}%
\pgfpathlineto{\pgfqpoint{2.303086in}{1.485000in}}%
\pgfusepath{stroke}%
\end{pgfscope}%
\begin{pgfscope}%
\definecolor{textcolor}{rgb}{0.000000,0.000000,0.000000}%
\pgfsetstrokecolor{textcolor}%
\pgfsetfillcolor{textcolor}%
\pgftext[x=1.395293in,y=1.568333in,,base]{\color{textcolor}{\rmfamily\fontsize{8.000000}{9.600000}\selectfont\catcode`\^=\active\def^{\ifmmode\sp\else\^{}\fi}\catcode`\%=\active\def%{\%}(a) Initialization scale}}%
\end{pgfscope}%
\begin{pgfscope}%
\pgfsetbuttcap%
\pgfsetmiterjoin%
\definecolor{currentfill}{rgb}{1.000000,1.000000,1.000000}%
\pgfsetfillcolor{currentfill}%
\pgfsetfillopacity{0.800000}%
\pgfsetlinewidth{1.003750pt}%
\definecolor{currentstroke}{rgb}{0.800000,0.800000,0.800000}%
\pgfsetstrokecolor{currentstroke}%
\pgfsetstrokeopacity{0.800000}%
\pgfsetdash{}{0pt}%
\pgfpathmoveto{\pgfqpoint{1.378915in}{1.280556in}}%
\pgfpathlineto{\pgfqpoint{2.258364in}{1.280556in}}%
\pgfpathquadraticcurveto{\pgfqpoint{2.271142in}{1.280556in}}{\pgfqpoint{2.271142in}{1.293333in}}%
\pgfpathlineto{\pgfqpoint{2.271142in}{1.440278in}}%
\pgfpathquadraticcurveto{\pgfqpoint{2.271142in}{1.453056in}}{\pgfqpoint{2.258364in}{1.453056in}}%
\pgfpathlineto{\pgfqpoint{1.378915in}{1.453056in}}%
\pgfpathquadraticcurveto{\pgfqpoint{1.366137in}{1.453056in}}{\pgfqpoint{1.366137in}{1.440278in}}%
\pgfpathlineto{\pgfqpoint{1.366137in}{1.293333in}}%
\pgfpathquadraticcurveto{\pgfqpoint{1.366137in}{1.280556in}}{\pgfqpoint{1.378915in}{1.280556in}}%
\pgfpathlineto{\pgfqpoint{1.378915in}{1.280556in}}%
\pgfpathclose%
\pgfusepath{stroke,fill}%
\end{pgfscope}%
\begin{pgfscope}%
\pgfsetrectcap%
\pgfsetroundjoin%
\pgfsetlinewidth{1.254687pt}%
\definecolor{currentstroke}{rgb}{0.141176,0.345098,0.650980}%
\pgfsetstrokecolor{currentstroke}%
\pgfsetdash{}{0pt}%
\pgfpathmoveto{\pgfqpoint{1.382110in}{1.410904in}}%
\pgfpathlineto{\pgfqpoint{1.426832in}{1.410904in}}%
\pgfpathlineto{\pgfqpoint{1.471554in}{1.410904in}}%
\pgfusepath{stroke}%
\end{pgfscope}%
\begin{pgfscope}%
\definecolor{textcolor}{rgb}{0.000000,0.000000,0.000000}%
\pgfsetstrokecolor{textcolor}%
\pgfsetfillcolor{textcolor}%
\pgftext[x=1.522665in,y=1.388543in,left,base]{\color{textcolor}{\rmfamily\fontsize{4.600000}{5.520000}\selectfont\catcode`\^=\active\def^{\ifmmode\sp\else\^{}\fi}\catcode`\%=\active\def%{\%}$\alpha=0.1$}}%
\end{pgfscope}%
\begin{pgfscope}%
\pgfsetrectcap%
\pgfsetroundjoin%
\pgfsetlinewidth{1.254687pt}%
\definecolor{currentstroke}{rgb}{0.819608,0.286275,0.356863}%
\pgfsetstrokecolor{currentstroke}%
\pgfsetdash{}{0pt}%
\pgfpathmoveto{\pgfqpoint{1.382110in}{1.334237in}}%
\pgfpathlineto{\pgfqpoint{1.426832in}{1.334237in}}%
\pgfpathlineto{\pgfqpoint{1.471554in}{1.334237in}}%
\pgfusepath{stroke}%
\end{pgfscope}%
\begin{pgfscope}%
\definecolor{textcolor}{rgb}{0.000000,0.000000,0.000000}%
\pgfsetstrokecolor{textcolor}%
\pgfsetfillcolor{textcolor}%
\pgftext[x=1.522665in,y=1.311876in,left,base]{\color{textcolor}{\rmfamily\fontsize{4.600000}{5.520000}\selectfont\catcode`\^=\active\def^{\ifmmode\sp\else\^{}\fi}\catcode`\%=\active\def%{\%}$\alpha=0.5$}}%
\end{pgfscope}%
\begin{pgfscope}%
\pgfsetrectcap%
\pgfsetroundjoin%
\pgfsetlinewidth{1.254687pt}%
\definecolor{currentstroke}{rgb}{0.000000,0.529412,0.423529}%
\pgfsetstrokecolor{currentstroke}%
\pgfsetdash{}{0pt}%
\pgfpathmoveto{\pgfqpoint{1.875754in}{1.410904in}}%
\pgfpathlineto{\pgfqpoint{1.920477in}{1.410904in}}%
\pgfpathlineto{\pgfqpoint{1.965199in}{1.410904in}}%
\pgfusepath{stroke}%
\end{pgfscope}%
\begin{pgfscope}%
\definecolor{textcolor}{rgb}{0.000000,0.000000,0.000000}%
\pgfsetstrokecolor{textcolor}%
\pgfsetfillcolor{textcolor}%
\pgftext[x=2.016310in,y=1.388543in,left,base]{\color{textcolor}{\rmfamily\fontsize{4.600000}{5.520000}\selectfont\catcode`\^=\active\def^{\ifmmode\sp\else\^{}\fi}\catcode`\%=\active\def%{\%}$\alpha=1$}}%
\end{pgfscope}%
\begin{pgfscope}%
\pgfsetrectcap%
\pgfsetroundjoin%
\pgfsetlinewidth{1.254687pt}%
\definecolor{currentstroke}{rgb}{0.478431,0.243137,0.615686}%
\pgfsetstrokecolor{currentstroke}%
\pgfsetdash{}{0pt}%
\pgfpathmoveto{\pgfqpoint{1.875754in}{1.334237in}}%
\pgfpathlineto{\pgfqpoint{1.920477in}{1.334237in}}%
\pgfpathlineto{\pgfqpoint{1.965199in}{1.334237in}}%
\pgfusepath{stroke}%
\end{pgfscope}%
\begin{pgfscope}%
\definecolor{textcolor}{rgb}{0.000000,0.000000,0.000000}%
\pgfsetstrokecolor{textcolor}%
\pgfsetfillcolor{textcolor}%
\pgftext[x=2.016310in,y=1.311876in,left,base]{\color{textcolor}{\rmfamily\fontsize{4.600000}{5.520000}\selectfont\catcode`\^=\active\def^{\ifmmode\sp\else\^{}\fi}\catcode`\%=\active\def%{\%}$\alpha=5$}}%
\end{pgfscope}%
\begin{pgfscope}%
\pgfsetbuttcap%
\pgfsetmiterjoin%
\definecolor{currentfill}{rgb}{1.000000,1.000000,1.000000}%
\pgfsetfillcolor{currentfill}%
\pgfsetlinewidth{0.000000pt}%
\definecolor{currentstroke}{rgb}{0.000000,0.000000,0.000000}%
\pgfsetstrokecolor{currentstroke}%
\pgfsetstrokeopacity{0.000000}%
\pgfsetdash{}{0pt}%
\pgfpathmoveto{\pgfqpoint{2.520957in}{0.320625in}}%
\pgfpathlineto{\pgfqpoint{4.336543in}{0.320625in}}%
\pgfpathlineto{\pgfqpoint{4.336543in}{1.485000in}}%
\pgfpathlineto{\pgfqpoint{2.520957in}{1.485000in}}%
\pgfpathlineto{\pgfqpoint{2.520957in}{0.320625in}}%
\pgfpathclose%
\pgfusepath{fill}%
\end{pgfscope}%
\begin{pgfscope}%
\pgfpathrectangle{\pgfqpoint{2.520957in}{0.320625in}}{\pgfqpoint{1.815586in}{1.164375in}}%
\pgfusepath{clip}%
\pgfsetrectcap%
\pgfsetroundjoin%
\pgfsetlinewidth{0.803000pt}%
\definecolor{currentstroke}{rgb}{0.909804,0.909804,0.909804}%
\pgfsetstrokecolor{currentstroke}%
\pgfsetdash{}{0pt}%
\pgfpathmoveto{\pgfqpoint{2.520957in}{0.320625in}}%
\pgfpathlineto{\pgfqpoint{2.520957in}{1.485000in}}%
\pgfusepath{stroke}%
\end{pgfscope}%
\begin{pgfscope}%
\pgfsetbuttcap%
\pgfsetroundjoin%
\definecolor{currentfill}{rgb}{0.000000,0.000000,0.000000}%
\pgfsetfillcolor{currentfill}%
\pgfsetlinewidth{0.803000pt}%
\definecolor{currentstroke}{rgb}{0.000000,0.000000,0.000000}%
\pgfsetstrokecolor{currentstroke}%
\pgfsetdash{}{0pt}%
\pgfsys@defobject{currentmarker}{\pgfqpoint{0.000000in}{-0.048611in}}{\pgfqpoint{0.000000in}{0.000000in}}{%
\pgfpathmoveto{\pgfqpoint{0.000000in}{0.000000in}}%
\pgfpathlineto{\pgfqpoint{0.000000in}{-0.048611in}}%
\pgfusepath{stroke,fill}%
}%
\begin{pgfscope}%
\pgfsys@transformshift{2.520957in}{0.320625in}%
\pgfsys@useobject{currentmarker}{}%
\end{pgfscope}%
\end{pgfscope}%
\begin{pgfscope}%
\definecolor{textcolor}{rgb}{0.000000,0.000000,0.000000}%
\pgfsetstrokecolor{textcolor}%
\pgfsetfillcolor{textcolor}%
\pgftext[x=2.520957in,y=0.223403in,,top]{\color{textcolor}{\rmfamily\fontsize{6.000000}{7.200000}\selectfont\catcode`\^=\active\def^{\ifmmode\sp\else\^{}\fi}\catcode`\%=\active\def%{\%}0}}%
\end{pgfscope}%
\begin{pgfscope}%
\pgfpathrectangle{\pgfqpoint{2.520957in}{0.320625in}}{\pgfqpoint{1.815586in}{1.164375in}}%
\pgfusepath{clip}%
\pgfsetrectcap%
\pgfsetroundjoin%
\pgfsetlinewidth{0.803000pt}%
\definecolor{currentstroke}{rgb}{0.909804,0.909804,0.909804}%
\pgfsetstrokecolor{currentstroke}%
\pgfsetdash{}{0pt}%
\pgfpathmoveto{\pgfqpoint{2.974853in}{0.320625in}}%
\pgfpathlineto{\pgfqpoint{2.974853in}{1.485000in}}%
\pgfusepath{stroke}%
\end{pgfscope}%
\begin{pgfscope}%
\pgfsetbuttcap%
\pgfsetroundjoin%
\definecolor{currentfill}{rgb}{0.000000,0.000000,0.000000}%
\pgfsetfillcolor{currentfill}%
\pgfsetlinewidth{0.803000pt}%
\definecolor{currentstroke}{rgb}{0.000000,0.000000,0.000000}%
\pgfsetstrokecolor{currentstroke}%
\pgfsetdash{}{0pt}%
\pgfsys@defobject{currentmarker}{\pgfqpoint{0.000000in}{-0.048611in}}{\pgfqpoint{0.000000in}{0.000000in}}{%
\pgfpathmoveto{\pgfqpoint{0.000000in}{0.000000in}}%
\pgfpathlineto{\pgfqpoint{0.000000in}{-0.048611in}}%
\pgfusepath{stroke,fill}%
}%
\begin{pgfscope}%
\pgfsys@transformshift{2.974853in}{0.320625in}%
\pgfsys@useobject{currentmarker}{}%
\end{pgfscope}%
\end{pgfscope}%
\begin{pgfscope}%
\definecolor{textcolor}{rgb}{0.000000,0.000000,0.000000}%
\pgfsetstrokecolor{textcolor}%
\pgfsetfillcolor{textcolor}%
\pgftext[x=2.974853in,y=0.223403in,,top]{\color{textcolor}{\rmfamily\fontsize{6.000000}{7.200000}\selectfont\catcode`\^=\active\def^{\ifmmode\sp\else\^{}\fi}\catcode`\%=\active\def%{\%}10}}%
\end{pgfscope}%
\begin{pgfscope}%
\pgfpathrectangle{\pgfqpoint{2.520957in}{0.320625in}}{\pgfqpoint{1.815586in}{1.164375in}}%
\pgfusepath{clip}%
\pgfsetrectcap%
\pgfsetroundjoin%
\pgfsetlinewidth{0.803000pt}%
\definecolor{currentstroke}{rgb}{0.909804,0.909804,0.909804}%
\pgfsetstrokecolor{currentstroke}%
\pgfsetdash{}{0pt}%
\pgfpathmoveto{\pgfqpoint{3.428750in}{0.320625in}}%
\pgfpathlineto{\pgfqpoint{3.428750in}{1.485000in}}%
\pgfusepath{stroke}%
\end{pgfscope}%
\begin{pgfscope}%
\pgfsetbuttcap%
\pgfsetroundjoin%
\definecolor{currentfill}{rgb}{0.000000,0.000000,0.000000}%
\pgfsetfillcolor{currentfill}%
\pgfsetlinewidth{0.803000pt}%
\definecolor{currentstroke}{rgb}{0.000000,0.000000,0.000000}%
\pgfsetstrokecolor{currentstroke}%
\pgfsetdash{}{0pt}%
\pgfsys@defobject{currentmarker}{\pgfqpoint{0.000000in}{-0.048611in}}{\pgfqpoint{0.000000in}{0.000000in}}{%
\pgfpathmoveto{\pgfqpoint{0.000000in}{0.000000in}}%
\pgfpathlineto{\pgfqpoint{0.000000in}{-0.048611in}}%
\pgfusepath{stroke,fill}%
}%
\begin{pgfscope}%
\pgfsys@transformshift{3.428750in}{0.320625in}%
\pgfsys@useobject{currentmarker}{}%
\end{pgfscope}%
\end{pgfscope}%
\begin{pgfscope}%
\definecolor{textcolor}{rgb}{0.000000,0.000000,0.000000}%
\pgfsetstrokecolor{textcolor}%
\pgfsetfillcolor{textcolor}%
\pgftext[x=3.428750in,y=0.223403in,,top]{\color{textcolor}{\rmfamily\fontsize{6.000000}{7.200000}\selectfont\catcode`\^=\active\def^{\ifmmode\sp\else\^{}\fi}\catcode`\%=\active\def%{\%}20}}%
\end{pgfscope}%
\begin{pgfscope}%
\pgfpathrectangle{\pgfqpoint{2.520957in}{0.320625in}}{\pgfqpoint{1.815586in}{1.164375in}}%
\pgfusepath{clip}%
\pgfsetrectcap%
\pgfsetroundjoin%
\pgfsetlinewidth{0.803000pt}%
\definecolor{currentstroke}{rgb}{0.909804,0.909804,0.909804}%
\pgfsetstrokecolor{currentstroke}%
\pgfsetdash{}{0pt}%
\pgfpathmoveto{\pgfqpoint{3.882647in}{0.320625in}}%
\pgfpathlineto{\pgfqpoint{3.882647in}{1.485000in}}%
\pgfusepath{stroke}%
\end{pgfscope}%
\begin{pgfscope}%
\pgfsetbuttcap%
\pgfsetroundjoin%
\definecolor{currentfill}{rgb}{0.000000,0.000000,0.000000}%
\pgfsetfillcolor{currentfill}%
\pgfsetlinewidth{0.803000pt}%
\definecolor{currentstroke}{rgb}{0.000000,0.000000,0.000000}%
\pgfsetstrokecolor{currentstroke}%
\pgfsetdash{}{0pt}%
\pgfsys@defobject{currentmarker}{\pgfqpoint{0.000000in}{-0.048611in}}{\pgfqpoint{0.000000in}{0.000000in}}{%
\pgfpathmoveto{\pgfqpoint{0.000000in}{0.000000in}}%
\pgfpathlineto{\pgfqpoint{0.000000in}{-0.048611in}}%
\pgfusepath{stroke,fill}%
}%
\begin{pgfscope}%
\pgfsys@transformshift{3.882647in}{0.320625in}%
\pgfsys@useobject{currentmarker}{}%
\end{pgfscope}%
\end{pgfscope}%
\begin{pgfscope}%
\definecolor{textcolor}{rgb}{0.000000,0.000000,0.000000}%
\pgfsetstrokecolor{textcolor}%
\pgfsetfillcolor{textcolor}%
\pgftext[x=3.882647in,y=0.223403in,,top]{\color{textcolor}{\rmfamily\fontsize{6.000000}{7.200000}\selectfont\catcode`\^=\active\def^{\ifmmode\sp\else\^{}\fi}\catcode`\%=\active\def%{\%}30}}%
\end{pgfscope}%
\begin{pgfscope}%
\pgfpathrectangle{\pgfqpoint{2.520957in}{0.320625in}}{\pgfqpoint{1.815586in}{1.164375in}}%
\pgfusepath{clip}%
\pgfsetrectcap%
\pgfsetroundjoin%
\pgfsetlinewidth{0.803000pt}%
\definecolor{currentstroke}{rgb}{0.909804,0.909804,0.909804}%
\pgfsetstrokecolor{currentstroke}%
\pgfsetdash{}{0pt}%
\pgfpathmoveto{\pgfqpoint{4.336543in}{0.320625in}}%
\pgfpathlineto{\pgfqpoint{4.336543in}{1.485000in}}%
\pgfusepath{stroke}%
\end{pgfscope}%
\begin{pgfscope}%
\pgfsetbuttcap%
\pgfsetroundjoin%
\definecolor{currentfill}{rgb}{0.000000,0.000000,0.000000}%
\pgfsetfillcolor{currentfill}%
\pgfsetlinewidth{0.803000pt}%
\definecolor{currentstroke}{rgb}{0.000000,0.000000,0.000000}%
\pgfsetstrokecolor{currentstroke}%
\pgfsetdash{}{0pt}%
\pgfsys@defobject{currentmarker}{\pgfqpoint{0.000000in}{-0.048611in}}{\pgfqpoint{0.000000in}{0.000000in}}{%
\pgfpathmoveto{\pgfqpoint{0.000000in}{0.000000in}}%
\pgfpathlineto{\pgfqpoint{0.000000in}{-0.048611in}}%
\pgfusepath{stroke,fill}%
}%
\begin{pgfscope}%
\pgfsys@transformshift{4.336543in}{0.320625in}%
\pgfsys@useobject{currentmarker}{}%
\end{pgfscope}%
\end{pgfscope}%
\begin{pgfscope}%
\definecolor{textcolor}{rgb}{0.000000,0.000000,0.000000}%
\pgfsetstrokecolor{textcolor}%
\pgfsetfillcolor{textcolor}%
\pgftext[x=4.336543in,y=0.223403in,,top]{\color{textcolor}{\rmfamily\fontsize{6.000000}{7.200000}\selectfont\catcode`\^=\active\def^{\ifmmode\sp\else\^{}\fi}\catcode`\%=\active\def%{\%}40}}%
\end{pgfscope}%
\begin{pgfscope}%
\definecolor{textcolor}{rgb}{0.000000,0.000000,0.000000}%
\pgfsetstrokecolor{textcolor}%
\pgfsetfillcolor{textcolor}%
\pgftext[x=3.428750in,y=0.084514in,,top]{\color{textcolor}{\rmfamily\fontsize{6.000000}{7.200000}\selectfont\catcode`\^=\active\def^{\ifmmode\sp\else\^{}\fi}\catcode`\%=\active\def%{\%}time}}%
\end{pgfscope}%
\begin{pgfscope}%
\pgfpathrectangle{\pgfqpoint{2.520957in}{0.320625in}}{\pgfqpoint{1.815586in}{1.164375in}}%
\pgfusepath{clip}%
\pgfsetrectcap%
\pgfsetroundjoin%
\pgfsetlinewidth{0.803000pt}%
\definecolor{currentstroke}{rgb}{0.909804,0.909804,0.909804}%
\pgfsetstrokecolor{currentstroke}%
\pgfsetdash{}{0pt}%
\pgfpathmoveto{\pgfqpoint{2.520957in}{0.466172in}}%
\pgfpathlineto{\pgfqpoint{4.336543in}{0.466172in}}%
\pgfusepath{stroke}%
\end{pgfscope}%
\begin{pgfscope}%
\pgfsetbuttcap%
\pgfsetroundjoin%
\definecolor{currentfill}{rgb}{0.000000,0.000000,0.000000}%
\pgfsetfillcolor{currentfill}%
\pgfsetlinewidth{0.803000pt}%
\definecolor{currentstroke}{rgb}{0.000000,0.000000,0.000000}%
\pgfsetstrokecolor{currentstroke}%
\pgfsetdash{}{0pt}%
\pgfsys@defobject{currentmarker}{\pgfqpoint{-0.048611in}{0.000000in}}{\pgfqpoint{-0.000000in}{0.000000in}}{%
\pgfpathmoveto{\pgfqpoint{-0.000000in}{0.000000in}}%
\pgfpathlineto{\pgfqpoint{-0.048611in}{0.000000in}}%
\pgfusepath{stroke,fill}%
}%
\begin{pgfscope}%
\pgfsys@transformshift{2.520957in}{0.466172in}%
\pgfsys@useobject{currentmarker}{}%
\end{pgfscope}%
\end{pgfscope}%
\begin{pgfscope}%
\pgfpathrectangle{\pgfqpoint{2.520957in}{0.320625in}}{\pgfqpoint{1.815586in}{1.164375in}}%
\pgfusepath{clip}%
\pgfsetrectcap%
\pgfsetroundjoin%
\pgfsetlinewidth{0.803000pt}%
\definecolor{currentstroke}{rgb}{0.909804,0.909804,0.909804}%
\pgfsetstrokecolor{currentstroke}%
\pgfsetdash{}{0pt}%
\pgfpathmoveto{\pgfqpoint{2.520957in}{0.660234in}}%
\pgfpathlineto{\pgfqpoint{4.336543in}{0.660234in}}%
\pgfusepath{stroke}%
\end{pgfscope}%
\begin{pgfscope}%
\pgfsetbuttcap%
\pgfsetroundjoin%
\definecolor{currentfill}{rgb}{0.000000,0.000000,0.000000}%
\pgfsetfillcolor{currentfill}%
\pgfsetlinewidth{0.803000pt}%
\definecolor{currentstroke}{rgb}{0.000000,0.000000,0.000000}%
\pgfsetstrokecolor{currentstroke}%
\pgfsetdash{}{0pt}%
\pgfsys@defobject{currentmarker}{\pgfqpoint{-0.048611in}{0.000000in}}{\pgfqpoint{-0.000000in}{0.000000in}}{%
\pgfpathmoveto{\pgfqpoint{-0.000000in}{0.000000in}}%
\pgfpathlineto{\pgfqpoint{-0.048611in}{0.000000in}}%
\pgfusepath{stroke,fill}%
}%
\begin{pgfscope}%
\pgfsys@transformshift{2.520957in}{0.660234in}%
\pgfsys@useobject{currentmarker}{}%
\end{pgfscope}%
\end{pgfscope}%
\begin{pgfscope}%
\pgfpathrectangle{\pgfqpoint{2.520957in}{0.320625in}}{\pgfqpoint{1.815586in}{1.164375in}}%
\pgfusepath{clip}%
\pgfsetrectcap%
\pgfsetroundjoin%
\pgfsetlinewidth{0.803000pt}%
\definecolor{currentstroke}{rgb}{0.909804,0.909804,0.909804}%
\pgfsetstrokecolor{currentstroke}%
\pgfsetdash{}{0pt}%
\pgfpathmoveto{\pgfqpoint{2.520957in}{0.854297in}}%
\pgfpathlineto{\pgfqpoint{4.336543in}{0.854297in}}%
\pgfusepath{stroke}%
\end{pgfscope}%
\begin{pgfscope}%
\pgfsetbuttcap%
\pgfsetroundjoin%
\definecolor{currentfill}{rgb}{0.000000,0.000000,0.000000}%
\pgfsetfillcolor{currentfill}%
\pgfsetlinewidth{0.803000pt}%
\definecolor{currentstroke}{rgb}{0.000000,0.000000,0.000000}%
\pgfsetstrokecolor{currentstroke}%
\pgfsetdash{}{0pt}%
\pgfsys@defobject{currentmarker}{\pgfqpoint{-0.048611in}{0.000000in}}{\pgfqpoint{-0.000000in}{0.000000in}}{%
\pgfpathmoveto{\pgfqpoint{-0.000000in}{0.000000in}}%
\pgfpathlineto{\pgfqpoint{-0.048611in}{0.000000in}}%
\pgfusepath{stroke,fill}%
}%
\begin{pgfscope}%
\pgfsys@transformshift{2.520957in}{0.854297in}%
\pgfsys@useobject{currentmarker}{}%
\end{pgfscope}%
\end{pgfscope}%
\begin{pgfscope}%
\pgfpathrectangle{\pgfqpoint{2.520957in}{0.320625in}}{\pgfqpoint{1.815586in}{1.164375in}}%
\pgfusepath{clip}%
\pgfsetrectcap%
\pgfsetroundjoin%
\pgfsetlinewidth{0.803000pt}%
\definecolor{currentstroke}{rgb}{0.909804,0.909804,0.909804}%
\pgfsetstrokecolor{currentstroke}%
\pgfsetdash{}{0pt}%
\pgfpathmoveto{\pgfqpoint{2.520957in}{1.048359in}}%
\pgfpathlineto{\pgfqpoint{4.336543in}{1.048359in}}%
\pgfusepath{stroke}%
\end{pgfscope}%
\begin{pgfscope}%
\pgfsetbuttcap%
\pgfsetroundjoin%
\definecolor{currentfill}{rgb}{0.000000,0.000000,0.000000}%
\pgfsetfillcolor{currentfill}%
\pgfsetlinewidth{0.803000pt}%
\definecolor{currentstroke}{rgb}{0.000000,0.000000,0.000000}%
\pgfsetstrokecolor{currentstroke}%
\pgfsetdash{}{0pt}%
\pgfsys@defobject{currentmarker}{\pgfqpoint{-0.048611in}{0.000000in}}{\pgfqpoint{-0.000000in}{0.000000in}}{%
\pgfpathmoveto{\pgfqpoint{-0.000000in}{0.000000in}}%
\pgfpathlineto{\pgfqpoint{-0.048611in}{0.000000in}}%
\pgfusepath{stroke,fill}%
}%
\begin{pgfscope}%
\pgfsys@transformshift{2.520957in}{1.048359in}%
\pgfsys@useobject{currentmarker}{}%
\end{pgfscope}%
\end{pgfscope}%
\begin{pgfscope}%
\pgfpathrectangle{\pgfqpoint{2.520957in}{0.320625in}}{\pgfqpoint{1.815586in}{1.164375in}}%
\pgfusepath{clip}%
\pgfsetrectcap%
\pgfsetroundjoin%
\pgfsetlinewidth{0.803000pt}%
\definecolor{currentstroke}{rgb}{0.909804,0.909804,0.909804}%
\pgfsetstrokecolor{currentstroke}%
\pgfsetdash{}{0pt}%
\pgfpathmoveto{\pgfqpoint{2.520957in}{1.242422in}}%
\pgfpathlineto{\pgfqpoint{4.336543in}{1.242422in}}%
\pgfusepath{stroke}%
\end{pgfscope}%
\begin{pgfscope}%
\pgfsetbuttcap%
\pgfsetroundjoin%
\definecolor{currentfill}{rgb}{0.000000,0.000000,0.000000}%
\pgfsetfillcolor{currentfill}%
\pgfsetlinewidth{0.803000pt}%
\definecolor{currentstroke}{rgb}{0.000000,0.000000,0.000000}%
\pgfsetstrokecolor{currentstroke}%
\pgfsetdash{}{0pt}%
\pgfsys@defobject{currentmarker}{\pgfqpoint{-0.048611in}{0.000000in}}{\pgfqpoint{-0.000000in}{0.000000in}}{%
\pgfpathmoveto{\pgfqpoint{-0.000000in}{0.000000in}}%
\pgfpathlineto{\pgfqpoint{-0.048611in}{0.000000in}}%
\pgfusepath{stroke,fill}%
}%
\begin{pgfscope}%
\pgfsys@transformshift{2.520957in}{1.242422in}%
\pgfsys@useobject{currentmarker}{}%
\end{pgfscope}%
\end{pgfscope}%
\begin{pgfscope}%
\pgfpathrectangle{\pgfqpoint{2.520957in}{0.320625in}}{\pgfqpoint{1.815586in}{1.164375in}}%
\pgfusepath{clip}%
\pgfsetrectcap%
\pgfsetroundjoin%
\pgfsetlinewidth{0.803000pt}%
\definecolor{currentstroke}{rgb}{0.909804,0.909804,0.909804}%
\pgfsetstrokecolor{currentstroke}%
\pgfsetdash{}{0pt}%
\pgfpathmoveto{\pgfqpoint{2.520957in}{1.436484in}}%
\pgfpathlineto{\pgfqpoint{4.336543in}{1.436484in}}%
\pgfusepath{stroke}%
\end{pgfscope}%
\begin{pgfscope}%
\pgfsetbuttcap%
\pgfsetroundjoin%
\definecolor{currentfill}{rgb}{0.000000,0.000000,0.000000}%
\pgfsetfillcolor{currentfill}%
\pgfsetlinewidth{0.803000pt}%
\definecolor{currentstroke}{rgb}{0.000000,0.000000,0.000000}%
\pgfsetstrokecolor{currentstroke}%
\pgfsetdash{}{0pt}%
\pgfsys@defobject{currentmarker}{\pgfqpoint{-0.048611in}{0.000000in}}{\pgfqpoint{-0.000000in}{0.000000in}}{%
\pgfpathmoveto{\pgfqpoint{-0.000000in}{0.000000in}}%
\pgfpathlineto{\pgfqpoint{-0.048611in}{0.000000in}}%
\pgfusepath{stroke,fill}%
}%
\begin{pgfscope}%
\pgfsys@transformshift{2.520957in}{1.436484in}%
\pgfsys@useobject{currentmarker}{}%
\end{pgfscope}%
\end{pgfscope}%
\begin{pgfscope}%
\pgfpathrectangle{\pgfqpoint{2.520957in}{0.320625in}}{\pgfqpoint{1.815586in}{1.164375in}}%
\pgfusepath{clip}%
\pgfsetrectcap%
\pgfsetroundjoin%
\pgfsetlinewidth{1.154312pt}%
\definecolor{currentstroke}{rgb}{0.146785,0.068738,0.334011}%
\pgfsetstrokecolor{currentstroke}%
\pgfsetdash{}{0pt}%
\pgfpathmoveto{\pgfqpoint{2.520957in}{1.451682in}}%
\pgfpathlineto{\pgfqpoint{2.525496in}{1.442427in}}%
\pgfpathlineto{\pgfqpoint{2.589041in}{1.331257in}}%
\pgfpathlineto{\pgfqpoint{2.611736in}{1.295326in}}%
\pgfpathlineto{\pgfqpoint{2.636700in}{1.258028in}}%
\pgfpathlineto{\pgfqpoint{2.661665in}{1.222794in}}%
\pgfpathlineto{\pgfqpoint{2.686629in}{1.189394in}}%
\pgfpathlineto{\pgfqpoint{2.716132in}{1.151870in}}%
\pgfpathlineto{\pgfqpoint{2.750175in}{1.110495in}}%
\pgfpathlineto{\pgfqpoint{2.793295in}{1.059936in}}%
\pgfpathlineto{\pgfqpoint{2.854571in}{0.989962in}}%
\pgfpathlineto{\pgfqpoint{2.945350in}{0.888116in}}%
\pgfpathlineto{\pgfqpoint{3.081519in}{0.737122in}}%
\pgfpathlineto{\pgfqpoint{3.215419in}{0.589913in}}%
\pgfpathlineto{\pgfqpoint{3.219958in}{0.584888in}}%
\pgfpathlineto{\pgfqpoint{3.226766in}{0.577622in}}%
\pgfpathlineto{\pgfqpoint{3.229035in}{0.575874in}}%
\pgfpathlineto{\pgfqpoint{3.233574in}{0.569904in}}%
\pgfpathlineto{\pgfqpoint{3.238113in}{0.565283in}}%
\pgfpathlineto{\pgfqpoint{3.240383in}{0.564172in}}%
\pgfpathlineto{\pgfqpoint{3.242652in}{0.560328in}}%
\pgfpathlineto{\pgfqpoint{3.249461in}{0.553037in}}%
\pgfpathlineto{\pgfqpoint{3.251730in}{0.553044in}}%
\pgfpathlineto{\pgfqpoint{3.254000in}{0.548593in}}%
\pgfpathlineto{\pgfqpoint{3.258539in}{0.542785in}}%
\pgfpathlineto{\pgfqpoint{3.260808in}{0.541155in}}%
\pgfpathlineto{\pgfqpoint{3.263078in}{0.543466in}}%
\pgfpathlineto{\pgfqpoint{3.265347in}{0.536921in}}%
\pgfpathlineto{\pgfqpoint{3.267617in}{0.534232in}}%
\pgfpathlineto{\pgfqpoint{3.269886in}{0.530300in}}%
\pgfpathlineto{\pgfqpoint{3.272156in}{0.528689in}}%
\pgfpathlineto{\pgfqpoint{3.274425in}{0.533012in}}%
\pgfpathlineto{\pgfqpoint{3.276695in}{0.526187in}}%
\pgfpathlineto{\pgfqpoint{3.278964in}{0.524269in}}%
\pgfpathlineto{\pgfqpoint{3.281234in}{0.517622in}}%
\pgfpathlineto{\pgfqpoint{3.283503in}{0.516825in}}%
\pgfpathlineto{\pgfqpoint{3.285773in}{0.525603in}}%
\pgfpathlineto{\pgfqpoint{3.288042in}{0.517780in}}%
\pgfpathlineto{\pgfqpoint{3.290312in}{0.516209in}}%
\pgfpathlineto{\pgfqpoint{3.292581in}{0.505218in}}%
\pgfpathlineto{\pgfqpoint{3.294851in}{0.505185in}}%
\pgfpathlineto{\pgfqpoint{3.297120in}{0.519563in}}%
\pgfpathlineto{\pgfqpoint{3.299389in}{0.508220in}}%
\pgfpathlineto{\pgfqpoint{3.301659in}{0.510777in}}%
\pgfpathlineto{\pgfqpoint{3.303928in}{0.493228in}}%
\pgfpathlineto{\pgfqpoint{3.306198in}{0.491439in}}%
\pgfpathlineto{\pgfqpoint{3.308467in}{0.508704in}}%
\pgfpathlineto{\pgfqpoint{3.310737in}{0.502688in}}%
\pgfpathlineto{\pgfqpoint{3.313006in}{0.508598in}}%
\pgfpathlineto{\pgfqpoint{3.315276in}{0.484187in}}%
\pgfpathlineto{\pgfqpoint{3.317545in}{0.479668in}}%
\pgfpathlineto{\pgfqpoint{3.319815in}{0.507585in}}%
\pgfpathlineto{\pgfqpoint{3.322084in}{0.500315in}}%
\pgfpathlineto{\pgfqpoint{3.324354in}{0.508972in}}%
\pgfpathlineto{\pgfqpoint{3.326623in}{0.473146in}}%
\pgfpathlineto{\pgfqpoint{3.328893in}{0.465915in}}%
\pgfpathlineto{\pgfqpoint{3.331162in}{0.499844in}}%
\pgfpathlineto{\pgfqpoint{3.333432in}{0.491880in}}%
\pgfpathlineto{\pgfqpoint{3.335701in}{0.507865in}}%
\pgfpathlineto{\pgfqpoint{3.337971in}{0.474434in}}%
\pgfpathlineto{\pgfqpoint{3.340240in}{0.455156in}}%
\pgfpathlineto{\pgfqpoint{3.342510in}{0.488026in}}%
\pgfpathlineto{\pgfqpoint{3.344779in}{0.492720in}}%
\pgfpathlineto{\pgfqpoint{3.347049in}{0.510603in}}%
\pgfpathlineto{\pgfqpoint{3.351588in}{0.445091in}}%
\pgfpathlineto{\pgfqpoint{3.353857in}{0.492022in}}%
\pgfpathlineto{\pgfqpoint{3.356127in}{0.490805in}}%
\pgfpathlineto{\pgfqpoint{3.358396in}{0.512269in}}%
\pgfpathlineto{\pgfqpoint{3.360666in}{0.472270in}}%
\pgfpathlineto{\pgfqpoint{3.362935in}{0.454234in}}%
\pgfpathlineto{\pgfqpoint{3.365204in}{0.474565in}}%
\pgfpathlineto{\pgfqpoint{3.367474in}{0.480839in}}%
\pgfpathlineto{\pgfqpoint{3.369743in}{0.509398in}}%
\pgfpathlineto{\pgfqpoint{3.372013in}{0.482895in}}%
\pgfpathlineto{\pgfqpoint{3.374282in}{0.464614in}}%
\pgfpathlineto{\pgfqpoint{3.376552in}{0.463191in}}%
\pgfpathlineto{\pgfqpoint{3.378821in}{0.485982in}}%
\pgfpathlineto{\pgfqpoint{3.381091in}{0.514455in}}%
\pgfpathlineto{\pgfqpoint{3.383360in}{0.484925in}}%
\pgfpathlineto{\pgfqpoint{3.385630in}{0.467850in}}%
\pgfpathlineto{\pgfqpoint{3.387899in}{0.468464in}}%
\pgfpathlineto{\pgfqpoint{3.390169in}{0.477860in}}%
\pgfpathlineto{\pgfqpoint{3.392438in}{0.512086in}}%
\pgfpathlineto{\pgfqpoint{3.394708in}{0.481661in}}%
\pgfpathlineto{\pgfqpoint{3.396977in}{0.479972in}}%
\pgfpathlineto{\pgfqpoint{3.399247in}{0.421575in}}%
\pgfpathlineto{\pgfqpoint{3.403786in}{0.507964in}}%
\pgfpathlineto{\pgfqpoint{3.406055in}{0.491470in}}%
\pgfpathlineto{\pgfqpoint{3.408325in}{0.486248in}}%
\pgfpathlineto{\pgfqpoint{3.410594in}{0.415194in}}%
\pgfpathlineto{\pgfqpoint{3.412864in}{0.474125in}}%
\pgfpathlineto{\pgfqpoint{3.415133in}{0.514486in}}%
\pgfpathlineto{\pgfqpoint{3.417403in}{0.489408in}}%
\pgfpathlineto{\pgfqpoint{3.419672in}{0.490470in}}%
\pgfpathlineto{\pgfqpoint{3.421942in}{0.407022in}}%
\pgfpathlineto{\pgfqpoint{3.426481in}{0.503794in}}%
\pgfpathlineto{\pgfqpoint{3.428750in}{0.486425in}}%
\pgfpathlineto{\pgfqpoint{3.431019in}{0.495611in}}%
\pgfpathlineto{\pgfqpoint{3.433289in}{0.432631in}}%
\pgfpathlineto{\pgfqpoint{3.435558in}{0.439868in}}%
\pgfpathlineto{\pgfqpoint{3.437828in}{0.502855in}}%
\pgfpathlineto{\pgfqpoint{3.440097in}{0.495805in}}%
\pgfpathlineto{\pgfqpoint{3.442367in}{0.501073in}}%
\pgfpathlineto{\pgfqpoint{3.444636in}{0.415491in}}%
\pgfpathlineto{\pgfqpoint{3.446906in}{0.445612in}}%
\pgfpathlineto{\pgfqpoint{3.449175in}{0.507177in}}%
\pgfpathlineto{\pgfqpoint{3.451445in}{0.488600in}}%
\pgfpathlineto{\pgfqpoint{3.453714in}{0.502886in}}%
\pgfpathlineto{\pgfqpoint{3.455984in}{0.438806in}}%
\pgfpathlineto{\pgfqpoint{3.458253in}{0.325850in}}%
\pgfpathlineto{\pgfqpoint{3.460523in}{0.488043in}}%
\pgfpathlineto{\pgfqpoint{3.462792in}{0.487140in}}%
\pgfpathlineto{\pgfqpoint{3.465062in}{0.505003in}}%
\pgfpathlineto{\pgfqpoint{3.467331in}{0.469319in}}%
\pgfpathlineto{\pgfqpoint{3.469590in}{0.315625in}}%
\pgfpathmoveto{\pgfqpoint{3.469610in}{0.315625in}}%
\pgfpathlineto{\pgfqpoint{3.471870in}{0.495321in}}%
\pgfpathlineto{\pgfqpoint{3.474140in}{0.495619in}}%
\pgfpathlineto{\pgfqpoint{3.476409in}{0.510889in}}%
\pgfpathlineto{\pgfqpoint{3.478679in}{0.456907in}}%
\pgfpathlineto{\pgfqpoint{3.480948in}{0.389575in}}%
\pgfpathlineto{\pgfqpoint{3.483218in}{0.490994in}}%
\pgfpathlineto{\pgfqpoint{3.485487in}{0.482008in}}%
\pgfpathlineto{\pgfqpoint{3.487757in}{0.507772in}}%
\pgfpathlineto{\pgfqpoint{3.490026in}{0.471224in}}%
\pgfpathlineto{\pgfqpoint{3.492296in}{0.450294in}}%
\pgfpathlineto{\pgfqpoint{3.494565in}{0.463155in}}%
\pgfpathlineto{\pgfqpoint{3.499104in}{0.511142in}}%
\pgfpathlineto{\pgfqpoint{3.501373in}{0.485082in}}%
\pgfpathlineto{\pgfqpoint{3.503643in}{0.448033in}}%
\pgfpathlineto{\pgfqpoint{3.505912in}{0.483187in}}%
\pgfpathlineto{\pgfqpoint{3.508182in}{0.489367in}}%
\pgfpathlineto{\pgfqpoint{3.510451in}{0.515524in}}%
\pgfpathlineto{\pgfqpoint{3.512721in}{0.473685in}}%
\pgfpathlineto{\pgfqpoint{3.517260in}{0.458265in}}%
\pgfpathlineto{\pgfqpoint{3.519529in}{0.468199in}}%
\pgfpathlineto{\pgfqpoint{3.521799in}{0.506186in}}%
\pgfpathlineto{\pgfqpoint{3.524068in}{0.485318in}}%
\pgfpathlineto{\pgfqpoint{3.526338in}{0.478823in}}%
\pgfpathlineto{\pgfqpoint{3.528607in}{0.410623in}}%
\pgfpathlineto{\pgfqpoint{3.530877in}{0.479400in}}%
\pgfpathlineto{\pgfqpoint{3.533146in}{0.514911in}}%
\pgfpathlineto{\pgfqpoint{3.535416in}{0.492612in}}%
\pgfpathlineto{\pgfqpoint{3.537685in}{0.480118in}}%
\pgfpathlineto{\pgfqpoint{3.539955in}{0.459313in}}%
\pgfpathlineto{\pgfqpoint{3.542224in}{0.472938in}}%
\pgfpathlineto{\pgfqpoint{3.544494in}{0.512658in}}%
\pgfpathlineto{\pgfqpoint{3.546763in}{0.480796in}}%
\pgfpathlineto{\pgfqpoint{3.549033in}{0.488794in}}%
\pgfpathlineto{\pgfqpoint{3.551302in}{0.361080in}}%
\pgfpathlineto{\pgfqpoint{3.553572in}{0.441894in}}%
\pgfpathlineto{\pgfqpoint{3.555841in}{0.499133in}}%
\pgfpathlineto{\pgfqpoint{3.558111in}{0.493212in}}%
\pgfpathlineto{\pgfqpoint{3.560380in}{0.495419in}}%
\pgfpathlineto{\pgfqpoint{3.562649in}{0.413478in}}%
\pgfpathlineto{\pgfqpoint{3.567188in}{0.514928in}}%
\pgfpathlineto{\pgfqpoint{3.569458in}{0.494068in}}%
\pgfpathlineto{\pgfqpoint{3.571727in}{0.498391in}}%
\pgfpathlineto{\pgfqpoint{3.573997in}{0.368309in}}%
\pgfpathlineto{\pgfqpoint{3.578536in}{0.498479in}}%
\pgfpathlineto{\pgfqpoint{3.580805in}{0.481698in}}%
\pgfpathlineto{\pgfqpoint{3.583075in}{0.499485in}}%
\pgfpathlineto{\pgfqpoint{3.585344in}{0.459260in}}%
\pgfpathlineto{\pgfqpoint{3.587614in}{0.341549in}}%
\pgfpathlineto{\pgfqpoint{3.589883in}{0.490801in}}%
\pgfpathlineto{\pgfqpoint{3.592153in}{0.497735in}}%
\pgfpathlineto{\pgfqpoint{3.594422in}{0.507471in}}%
\pgfpathlineto{\pgfqpoint{3.596692in}{0.452613in}}%
\pgfpathlineto{\pgfqpoint{3.598961in}{0.431996in}}%
\pgfpathlineto{\pgfqpoint{3.601231in}{0.508106in}}%
\pgfpathlineto{\pgfqpoint{3.603500in}{0.489257in}}%
\pgfpathlineto{\pgfqpoint{3.605770in}{0.507818in}}%
\pgfpathlineto{\pgfqpoint{3.608039in}{0.445361in}}%
\pgfpathlineto{\pgfqpoint{3.610309in}{0.425648in}}%
\pgfpathlineto{\pgfqpoint{3.612578in}{0.472857in}}%
\pgfpathlineto{\pgfqpoint{3.614848in}{0.478561in}}%
\pgfpathlineto{\pgfqpoint{3.617117in}{0.504872in}}%
\pgfpathlineto{\pgfqpoint{3.619387in}{0.482785in}}%
\pgfpathlineto{\pgfqpoint{3.621656in}{0.434939in}}%
\pgfpathlineto{\pgfqpoint{3.623926in}{0.482843in}}%
\pgfpathlineto{\pgfqpoint{3.628465in}{0.516289in}}%
\pgfpathlineto{\pgfqpoint{3.633003in}{0.425614in}}%
\pgfpathlineto{\pgfqpoint{3.635273in}{0.490856in}}%
\pgfpathlineto{\pgfqpoint{3.637542in}{0.476465in}}%
\pgfpathlineto{\pgfqpoint{3.639812in}{0.508432in}}%
\pgfpathlineto{\pgfqpoint{3.642081in}{0.471719in}}%
\pgfpathlineto{\pgfqpoint{3.644351in}{0.469650in}}%
\pgfpathlineto{\pgfqpoint{3.646620in}{0.414967in}}%
\pgfpathlineto{\pgfqpoint{3.648890in}{0.473423in}}%
\pgfpathlineto{\pgfqpoint{3.651159in}{0.508696in}}%
\pgfpathlineto{\pgfqpoint{3.653429in}{0.495387in}}%
\pgfpathlineto{\pgfqpoint{3.655698in}{0.468847in}}%
\pgfpathlineto{\pgfqpoint{3.657968in}{0.474019in}}%
\pgfpathlineto{\pgfqpoint{3.660237in}{0.492740in}}%
\pgfpathlineto{\pgfqpoint{3.662507in}{0.520538in}}%
\pgfpathlineto{\pgfqpoint{3.664776in}{0.475780in}}%
\pgfpathlineto{\pgfqpoint{3.667046in}{0.476690in}}%
\pgfpathlineto{\pgfqpoint{3.669315in}{0.450890in}}%
\pgfpathlineto{\pgfqpoint{3.671585in}{0.450282in}}%
\pgfpathlineto{\pgfqpoint{3.673854in}{0.498835in}}%
\pgfpathlineto{\pgfqpoint{3.676124in}{0.483016in}}%
\pgfpathlineto{\pgfqpoint{3.678393in}{0.488064in}}%
\pgfpathlineto{\pgfqpoint{3.680663in}{0.430277in}}%
\pgfpathlineto{\pgfqpoint{3.685202in}{0.512771in}}%
\pgfpathlineto{\pgfqpoint{3.689741in}{0.490091in}}%
\pgfpathlineto{\pgfqpoint{3.692010in}{0.457223in}}%
\pgfpathlineto{\pgfqpoint{3.694280in}{0.473927in}}%
\pgfpathlineto{\pgfqpoint{3.696549in}{0.516036in}}%
\pgfpathlineto{\pgfqpoint{3.698818in}{0.477132in}}%
\pgfpathlineto{\pgfqpoint{3.701088in}{0.494362in}}%
\pgfpathlineto{\pgfqpoint{3.703357in}{0.413469in}}%
\pgfpathlineto{\pgfqpoint{3.705627in}{0.362507in}}%
\pgfpathlineto{\pgfqpoint{3.707896in}{0.482580in}}%
\pgfpathlineto{\pgfqpoint{3.710166in}{0.490208in}}%
\pgfpathlineto{\pgfqpoint{3.712435in}{0.500616in}}%
\pgfpathlineto{\pgfqpoint{3.714705in}{0.465630in}}%
\pgfpathlineto{\pgfqpoint{3.716974in}{0.458035in}}%
\pgfpathlineto{\pgfqpoint{3.719244in}{0.515861in}}%
\pgfpathlineto{\pgfqpoint{3.721513in}{0.501713in}}%
\pgfpathlineto{\pgfqpoint{3.723783in}{0.505222in}}%
\pgfpathlineto{\pgfqpoint{3.726052in}{0.399293in}}%
\pgfpathlineto{\pgfqpoint{3.728322in}{0.407367in}}%
\pgfpathlineto{\pgfqpoint{3.730591in}{0.496634in}}%
\pgfpathlineto{\pgfqpoint{3.732861in}{0.472101in}}%
\pgfpathlineto{\pgfqpoint{3.735130in}{0.499403in}}%
\pgfpathlineto{\pgfqpoint{3.739669in}{0.433172in}}%
\pgfpathlineto{\pgfqpoint{3.744208in}{0.496846in}}%
\pgfpathlineto{\pgfqpoint{3.746478in}{0.511730in}}%
\pgfpathlineto{\pgfqpoint{3.748747in}{0.477990in}}%
\pgfpathlineto{\pgfqpoint{3.751017in}{0.428552in}}%
\pgfpathlineto{\pgfqpoint{3.753286in}{0.513960in}}%
\pgfpathlineto{\pgfqpoint{3.755556in}{0.493454in}}%
\pgfpathlineto{\pgfqpoint{3.757825in}{0.513169in}}%
\pgfpathlineto{\pgfqpoint{3.760095in}{0.433568in}}%
\pgfpathlineto{\pgfqpoint{3.762364in}{0.449685in}}%
\pgfpathlineto{\pgfqpoint{3.764633in}{0.457365in}}%
\pgfpathlineto{\pgfqpoint{3.766903in}{0.461391in}}%
\pgfpathlineto{\pgfqpoint{3.769172in}{0.498394in}}%
\pgfpathlineto{\pgfqpoint{3.771442in}{0.488389in}}%
\pgfpathlineto{\pgfqpoint{3.773711in}{0.464525in}}%
\pgfpathlineto{\pgfqpoint{3.775981in}{0.452189in}}%
\pgfpathlineto{\pgfqpoint{3.778250in}{0.502224in}}%
\pgfpathlineto{\pgfqpoint{3.780520in}{0.521597in}}%
\pgfpathlineto{\pgfqpoint{3.782789in}{0.480092in}}%
\pgfpathlineto{\pgfqpoint{3.785059in}{0.425770in}}%
\pgfpathlineto{\pgfqpoint{3.787328in}{0.500782in}}%
\pgfpathlineto{\pgfqpoint{3.789598in}{0.473408in}}%
\pgfpathlineto{\pgfqpoint{3.791867in}{0.510013in}}%
\pgfpathlineto{\pgfqpoint{3.794137in}{0.463163in}}%
\pgfpathlineto{\pgfqpoint{3.796406in}{0.478154in}}%
\pgfpathlineto{\pgfqpoint{3.798676in}{0.397244in}}%
\pgfpathlineto{\pgfqpoint{3.803215in}{0.498074in}}%
\pgfpathlineto{\pgfqpoint{3.805484in}{0.502334in}}%
\pgfpathlineto{\pgfqpoint{3.807754in}{0.482941in}}%
\pgfpathlineto{\pgfqpoint{3.810023in}{0.455868in}}%
\pgfpathlineto{\pgfqpoint{3.812293in}{0.501919in}}%
\pgfpathlineto{\pgfqpoint{3.814562in}{0.527962in}}%
\pgfpathlineto{\pgfqpoint{3.816832in}{0.476547in}}%
\pgfpathlineto{\pgfqpoint{3.819101in}{0.480620in}}%
\pgfpathlineto{\pgfqpoint{3.821371in}{0.465231in}}%
\pgfpathlineto{\pgfqpoint{3.823640in}{0.425736in}}%
\pgfpathlineto{\pgfqpoint{3.825910in}{0.490045in}}%
\pgfpathlineto{\pgfqpoint{3.828179in}{0.472463in}}%
\pgfpathlineto{\pgfqpoint{3.830448in}{0.489646in}}%
\pgfpathlineto{\pgfqpoint{3.834987in}{0.438686in}}%
\pgfpathlineto{\pgfqpoint{3.837257in}{0.504547in}}%
\pgfpathlineto{\pgfqpoint{3.839526in}{0.511995in}}%
\pgfpathlineto{\pgfqpoint{3.841796in}{0.497350in}}%
\pgfpathlineto{\pgfqpoint{3.844065in}{0.464855in}}%
\pgfpathlineto{\pgfqpoint{3.846335in}{0.487823in}}%
\pgfpathlineto{\pgfqpoint{3.848604in}{0.525898in}}%
\pgfpathlineto{\pgfqpoint{3.850874in}{0.469844in}}%
\pgfpathlineto{\pgfqpoint{3.853143in}{0.497469in}}%
\pgfpathlineto{\pgfqpoint{3.855413in}{0.389762in}}%
\pgfpathlineto{\pgfqpoint{3.857682in}{0.414862in}}%
\pgfpathlineto{\pgfqpoint{3.859952in}{0.455402in}}%
\pgfpathlineto{\pgfqpoint{3.864491in}{0.499083in}}%
\pgfpathlineto{\pgfqpoint{3.866760in}{0.489836in}}%
\pgfpathlineto{\pgfqpoint{3.869030in}{0.441181in}}%
\pgfpathlineto{\pgfqpoint{3.871299in}{0.516302in}}%
\pgfpathlineto{\pgfqpoint{3.873569in}{0.514340in}}%
\pgfpathlineto{\pgfqpoint{3.875838in}{0.509969in}}%
\pgfpathlineto{\pgfqpoint{3.878108in}{0.457022in}}%
\pgfpathlineto{\pgfqpoint{3.880377in}{0.431381in}}%
\pgfpathlineto{\pgfqpoint{3.882647in}{0.505545in}}%
\pgfpathlineto{\pgfqpoint{3.884916in}{0.457013in}}%
\pgfpathlineto{\pgfqpoint{3.887186in}{0.496288in}}%
\pgfpathlineto{\pgfqpoint{3.889455in}{0.462107in}}%
\pgfpathlineto{\pgfqpoint{3.891725in}{0.457070in}}%
\pgfpathlineto{\pgfqpoint{3.893994in}{0.351449in}}%
\pgfpathlineto{\pgfqpoint{3.896264in}{0.489137in}}%
\pgfpathlineto{\pgfqpoint{3.898533in}{0.511378in}}%
\pgfpathlineto{\pgfqpoint{3.900802in}{0.499847in}}%
\pgfpathlineto{\pgfqpoint{3.903072in}{0.445925in}}%
\pgfpathlineto{\pgfqpoint{3.905341in}{0.525551in}}%
\pgfpathlineto{\pgfqpoint{3.907611in}{0.505282in}}%
\pgfpathlineto{\pgfqpoint{3.909880in}{0.518874in}}%
\pgfpathlineto{\pgfqpoint{3.912150in}{0.386662in}}%
\pgfpathlineto{\pgfqpoint{3.914419in}{0.453442in}}%
\pgfpathlineto{\pgfqpoint{3.916689in}{0.460126in}}%
\pgfpathlineto{\pgfqpoint{3.918958in}{0.432171in}}%
\pgfpathlineto{\pgfqpoint{3.921228in}{0.483291in}}%
\pgfpathlineto{\pgfqpoint{3.923497in}{0.483292in}}%
\pgfpathlineto{\pgfqpoint{3.925767in}{0.475880in}}%
\pgfpathlineto{\pgfqpoint{3.928036in}{0.416987in}}%
\pgfpathlineto{\pgfqpoint{3.930306in}{0.504171in}}%
\pgfpathlineto{\pgfqpoint{3.932575in}{0.525091in}}%
\pgfpathlineto{\pgfqpoint{3.934845in}{0.428830in}}%
\pgfpathlineto{\pgfqpoint{3.937114in}{0.467587in}}%
\pgfpathlineto{\pgfqpoint{3.939384in}{0.455734in}}%
\pgfpathlineto{\pgfqpoint{3.941653in}{0.407056in}}%
\pgfpathlineto{\pgfqpoint{3.943923in}{0.475966in}}%
\pgfpathlineto{\pgfqpoint{3.948462in}{0.483612in}}%
\pgfpathlineto{\pgfqpoint{3.950731in}{0.464003in}}%
\pgfpathlineto{\pgfqpoint{3.955270in}{0.525180in}}%
\pgfpathlineto{\pgfqpoint{3.957540in}{0.456925in}}%
\pgfpathlineto{\pgfqpoint{3.959809in}{0.480801in}}%
\pgfpathlineto{\pgfqpoint{3.962079in}{0.432064in}}%
\pgfpathlineto{\pgfqpoint{3.964348in}{0.368458in}}%
\pgfpathlineto{\pgfqpoint{3.966617in}{0.470924in}}%
\pgfpathlineto{\pgfqpoint{3.971156in}{0.489801in}}%
\pgfpathlineto{\pgfqpoint{3.973426in}{0.475014in}}%
\pgfpathlineto{\pgfqpoint{3.975695in}{0.484601in}}%
\pgfpathlineto{\pgfqpoint{3.977965in}{0.525010in}}%
\pgfpathlineto{\pgfqpoint{3.980234in}{0.467419in}}%
\pgfpathlineto{\pgfqpoint{3.982504in}{0.490208in}}%
\pgfpathlineto{\pgfqpoint{3.984773in}{0.391022in}}%
\pgfpathlineto{\pgfqpoint{3.987043in}{0.369295in}}%
\pgfpathlineto{\pgfqpoint{3.989312in}{0.462319in}}%
\pgfpathlineto{\pgfqpoint{3.993851in}{0.494466in}}%
\pgfpathlineto{\pgfqpoint{3.996121in}{0.485005in}}%
\pgfpathlineto{\pgfqpoint{3.998390in}{0.469675in}}%
\pgfpathlineto{\pgfqpoint{4.000660in}{0.522653in}}%
\pgfpathlineto{\pgfqpoint{4.002929in}{0.473756in}}%
\pgfpathlineto{\pgfqpoint{4.005199in}{0.497409in}}%
\pgfpathlineto{\pgfqpoint{4.007468in}{0.371575in}}%
\pgfpathlineto{\pgfqpoint{4.009738in}{0.418957in}}%
\pgfpathlineto{\pgfqpoint{4.014277in}{0.476533in}}%
\pgfpathlineto{\pgfqpoint{4.016546in}{0.497988in}}%
\pgfpathlineto{\pgfqpoint{4.018816in}{0.492183in}}%
\pgfpathlineto{\pgfqpoint{4.021085in}{0.444431in}}%
\pgfpathlineto{\pgfqpoint{4.023355in}{0.517950in}}%
\pgfpathlineto{\pgfqpoint{4.025624in}{0.516753in}}%
\pgfpathlineto{\pgfqpoint{4.027894in}{0.509015in}}%
\pgfpathlineto{\pgfqpoint{4.032432in}{0.440156in}}%
\pgfpathlineto{\pgfqpoint{4.034702in}{0.508017in}}%
\pgfpathlineto{\pgfqpoint{4.036971in}{0.452921in}}%
\pgfpathlineto{\pgfqpoint{4.039241in}{0.495064in}}%
\pgfpathlineto{\pgfqpoint{4.041510in}{0.459899in}}%
\pgfpathlineto{\pgfqpoint{4.043780in}{0.457351in}}%
\pgfpathlineto{\pgfqpoint{4.046049in}{0.325364in}}%
\pgfpathlineto{\pgfqpoint{4.048319in}{0.488587in}}%
\pgfpathlineto{\pgfqpoint{4.050588in}{0.510964in}}%
\pgfpathlineto{\pgfqpoint{4.052858in}{0.502042in}}%
\pgfpathlineto{\pgfqpoint{4.055127in}{0.457574in}}%
\pgfpathlineto{\pgfqpoint{4.057397in}{0.529366in}}%
\pgfpathlineto{\pgfqpoint{4.059666in}{0.507159in}}%
\pgfpathlineto{\pgfqpoint{4.061936in}{0.518615in}}%
\pgfpathlineto{\pgfqpoint{4.064205in}{0.425975in}}%
\pgfpathlineto{\pgfqpoint{4.066475in}{0.451147in}}%
\pgfpathlineto{\pgfqpoint{4.068744in}{0.462382in}}%
\pgfpathlineto{\pgfqpoint{4.071014in}{0.426763in}}%
\pgfpathlineto{\pgfqpoint{4.073283in}{0.480256in}}%
\pgfpathlineto{\pgfqpoint{4.075553in}{0.481967in}}%
\pgfpathlineto{\pgfqpoint{4.077822in}{0.475325in}}%
\pgfpathlineto{\pgfqpoint{4.080092in}{0.428836in}}%
\pgfpathlineto{\pgfqpoint{4.082361in}{0.505935in}}%
\pgfpathlineto{\pgfqpoint{4.084631in}{0.525507in}}%
\pgfpathlineto{\pgfqpoint{4.086900in}{0.411571in}}%
\pgfpathlineto{\pgfqpoint{4.089170in}{0.465359in}}%
\pgfpathlineto{\pgfqpoint{4.091439in}{0.457827in}}%
\pgfpathlineto{\pgfqpoint{4.093709in}{0.402121in}}%
\pgfpathlineto{\pgfqpoint{4.095978in}{0.473581in}}%
\pgfpathlineto{\pgfqpoint{4.100517in}{0.482575in}}%
\pgfpathlineto{\pgfqpoint{4.102786in}{0.465810in}}%
\pgfpathlineto{\pgfqpoint{4.107325in}{0.525970in}}%
\pgfpathlineto{\pgfqpoint{4.109595in}{0.449987in}}%
\pgfpathlineto{\pgfqpoint{4.111864in}{0.478697in}}%
\pgfpathlineto{\pgfqpoint{4.114134in}{0.437338in}}%
\pgfpathlineto{\pgfqpoint{4.116403in}{0.363368in}}%
\pgfpathlineto{\pgfqpoint{4.118673in}{0.468835in}}%
\pgfpathlineto{\pgfqpoint{4.123212in}{0.488584in}}%
\pgfpathlineto{\pgfqpoint{4.125481in}{0.476043in}}%
\pgfpathlineto{\pgfqpoint{4.127751in}{0.487972in}}%
\pgfpathlineto{\pgfqpoint{4.130020in}{0.526103in}}%
\pgfpathlineto{\pgfqpoint{4.132290in}{0.463111in}}%
\pgfpathlineto{\pgfqpoint{4.134559in}{0.488336in}}%
\pgfpathlineto{\pgfqpoint{4.136829in}{0.405556in}}%
\pgfpathlineto{\pgfqpoint{4.139098in}{0.366090in}}%
\pgfpathlineto{\pgfqpoint{4.141368in}{0.460727in}}%
\pgfpathlineto{\pgfqpoint{4.145907in}{0.493232in}}%
\pgfpathlineto{\pgfqpoint{4.148176in}{0.485480in}}%
\pgfpathlineto{\pgfqpoint{4.150446in}{0.474715in}}%
\pgfpathlineto{\pgfqpoint{4.152715in}{0.524317in}}%
\pgfpathlineto{\pgfqpoint{4.154985in}{0.470993in}}%
\pgfpathlineto{\pgfqpoint{4.157254in}{0.495808in}}%
\pgfpathlineto{\pgfqpoint{4.159524in}{0.317654in}}%
\pgfpathlineto{\pgfqpoint{4.161793in}{0.415548in}}%
\pgfpathlineto{\pgfqpoint{4.166332in}{0.476111in}}%
\pgfpathlineto{\pgfqpoint{4.168601in}{0.496832in}}%
\pgfpathlineto{\pgfqpoint{4.170871in}{0.492437in}}%
\pgfpathlineto{\pgfqpoint{4.173140in}{0.453677in}}%
\pgfpathlineto{\pgfqpoint{4.175410in}{0.520428in}}%
\pgfpathlineto{\pgfqpoint{4.177679in}{0.475043in}}%
\pgfpathlineto{\pgfqpoint{4.179949in}{0.501426in}}%
\pgfpathlineto{\pgfqpoint{4.182218in}{0.411961in}}%
\pgfpathlineto{\pgfqpoint{4.184488in}{0.439766in}}%
\pgfpathlineto{\pgfqpoint{4.186757in}{0.430733in}}%
\pgfpathlineto{\pgfqpoint{4.189027in}{0.472403in}}%
\pgfpathlineto{\pgfqpoint{4.191296in}{0.499200in}}%
\pgfpathlineto{\pgfqpoint{4.193566in}{0.497794in}}%
\pgfpathlineto{\pgfqpoint{4.195835in}{0.403947in}}%
\pgfpathlineto{\pgfqpoint{4.198105in}{0.514556in}}%
\pgfpathlineto{\pgfqpoint{4.200374in}{0.476241in}}%
\pgfpathlineto{\pgfqpoint{4.202644in}{0.505360in}}%
\pgfpathlineto{\pgfqpoint{4.204913in}{0.436787in}}%
\pgfpathlineto{\pgfqpoint{4.207183in}{0.456048in}}%
\pgfpathlineto{\pgfqpoint{4.209452in}{0.393516in}}%
\pgfpathlineto{\pgfqpoint{4.211722in}{0.466718in}}%
\pgfpathlineto{\pgfqpoint{4.213991in}{0.500215in}}%
\pgfpathlineto{\pgfqpoint{4.216261in}{0.501782in}}%
\pgfpathlineto{\pgfqpoint{4.218530in}{0.416676in}}%
\pgfpathlineto{\pgfqpoint{4.220800in}{0.504942in}}%
\pgfpathlineto{\pgfqpoint{4.223069in}{0.516453in}}%
\pgfpathlineto{\pgfqpoint{4.225339in}{0.523027in}}%
\pgfpathlineto{\pgfqpoint{4.227608in}{0.350612in}}%
\pgfpathlineto{\pgfqpoint{4.229878in}{0.408282in}}%
\pgfpathlineto{\pgfqpoint{4.232147in}{0.495027in}}%
\pgfpathlineto{\pgfqpoint{4.234416in}{0.443415in}}%
\pgfpathlineto{\pgfqpoint{4.236686in}{0.492989in}}%
\pgfpathlineto{\pgfqpoint{4.238955in}{0.466871in}}%
\pgfpathlineto{\pgfqpoint{4.241225in}{0.475275in}}%
\pgfpathlineto{\pgfqpoint{4.243494in}{0.445854in}}%
\pgfpathlineto{\pgfqpoint{4.248033in}{0.512764in}}%
\pgfpathlineto{\pgfqpoint{4.250303in}{0.512241in}}%
\pgfpathlineto{\pgfqpoint{4.252572in}{0.427271in}}%
\pgfpathlineto{\pgfqpoint{4.254842in}{0.517069in}}%
\pgfpathlineto{\pgfqpoint{4.257111in}{0.505746in}}%
\pgfpathlineto{\pgfqpoint{4.259381in}{0.527039in}}%
\pgfpathlineto{\pgfqpoint{4.261650in}{0.409696in}}%
\pgfpathlineto{\pgfqpoint{4.263920in}{0.477999in}}%
\pgfpathlineto{\pgfqpoint{4.266189in}{0.427618in}}%
\pgfpathlineto{\pgfqpoint{4.268459in}{0.360442in}}%
\pgfpathlineto{\pgfqpoint{4.270728in}{0.467590in}}%
\pgfpathlineto{\pgfqpoint{4.272998in}{0.482034in}}%
\pgfpathlineto{\pgfqpoint{4.275267in}{0.488884in}}%
\pgfpathlineto{\pgfqpoint{4.277537in}{0.473922in}}%
\pgfpathlineto{\pgfqpoint{4.279806in}{0.493599in}}%
\pgfpathlineto{\pgfqpoint{4.282076in}{0.528269in}}%
\pgfpathlineto{\pgfqpoint{4.284345in}{0.456075in}}%
\pgfpathlineto{\pgfqpoint{4.286615in}{0.487187in}}%
\pgfpathlineto{\pgfqpoint{4.288884in}{0.416392in}}%
\pgfpathlineto{\pgfqpoint{4.291154in}{0.378687in}}%
\pgfpathlineto{\pgfqpoint{4.293423in}{0.456059in}}%
\pgfpathlineto{\pgfqpoint{4.297962in}{0.491973in}}%
\pgfpathlineto{\pgfqpoint{4.300231in}{0.488197in}}%
\pgfpathlineto{\pgfqpoint{4.302501in}{0.477199in}}%
\pgfpathlineto{\pgfqpoint{4.304770in}{0.525594in}}%
\pgfpathlineto{\pgfqpoint{4.307040in}{0.468004in}}%
\pgfpathlineto{\pgfqpoint{4.309309in}{0.494753in}}%
\pgfpathlineto{\pgfqpoint{4.311389in}{0.315625in}}%
\pgfpathmoveto{\pgfqpoint{4.311897in}{0.315625in}}%
\pgfpathlineto{\pgfqpoint{4.313848in}{0.416277in}}%
\pgfpathlineto{\pgfqpoint{4.318387in}{0.475231in}}%
\pgfpathlineto{\pgfqpoint{4.320657in}{0.495898in}}%
\pgfpathlineto{\pgfqpoint{4.322926in}{0.493704in}}%
\pgfpathlineto{\pgfqpoint{4.325196in}{0.459776in}}%
\pgfpathlineto{\pgfqpoint{4.327465in}{0.522634in}}%
\pgfpathlineto{\pgfqpoint{4.329735in}{0.472852in}}%
\pgfpathlineto{\pgfqpoint{4.332004in}{0.500463in}}%
\pgfpathlineto{\pgfqpoint{4.334274in}{0.456483in}}%
\pgfpathlineto{\pgfqpoint{4.336543in}{0.433044in}}%
\pgfpathlineto{\pgfqpoint{4.336543in}{0.433044in}}%
\pgfusepath{stroke}%
\end{pgfscope}%
\begin{pgfscope}%
\pgfpathrectangle{\pgfqpoint{2.520957in}{0.320625in}}{\pgfqpoint{1.815586in}{1.164375in}}%
\pgfusepath{clip}%
\pgfsetrectcap%
\pgfsetroundjoin%
\pgfsetlinewidth{1.154312pt}%
\definecolor{currentstroke}{rgb}{0.291366,0.064553,0.475462}%
\pgfsetstrokecolor{currentstroke}%
\pgfsetdash{}{0pt}%
\pgfpathmoveto{\pgfqpoint{2.520957in}{1.451682in}}%
\pgfpathlineto{\pgfqpoint{2.525496in}{1.442322in}}%
\pgfpathlineto{\pgfqpoint{2.548191in}{1.400764in}}%
\pgfpathlineto{\pgfqpoint{2.586772in}{1.329007in}}%
\pgfpathlineto{\pgfqpoint{2.609467in}{1.289602in}}%
\pgfpathlineto{\pgfqpoint{2.632161in}{1.252372in}}%
\pgfpathlineto{\pgfqpoint{2.657126in}{1.213506in}}%
\pgfpathlineto{\pgfqpoint{2.684360in}{1.173124in}}%
\pgfpathlineto{\pgfqpoint{2.716132in}{1.128073in}}%
\pgfpathlineto{\pgfqpoint{2.752444in}{1.078542in}}%
\pgfpathlineto{\pgfqpoint{2.800103in}{1.015520in}}%
\pgfpathlineto{\pgfqpoint{2.868188in}{0.927528in}}%
\pgfpathlineto{\pgfqpoint{2.977123in}{0.788779in}}%
\pgfpathlineto{\pgfqpoint{3.133717in}{0.591064in}}%
\pgfpathlineto{\pgfqpoint{3.138256in}{0.585322in}}%
\pgfpathlineto{\pgfqpoint{3.145065in}{0.576864in}}%
\pgfpathlineto{\pgfqpoint{3.147334in}{0.574808in}}%
\pgfpathlineto{\pgfqpoint{3.151873in}{0.568202in}}%
\pgfpathlineto{\pgfqpoint{3.156412in}{0.562917in}}%
\pgfpathlineto{\pgfqpoint{3.158682in}{0.561949in}}%
\pgfpathlineto{\pgfqpoint{3.160951in}{0.557185in}}%
\pgfpathlineto{\pgfqpoint{3.167759in}{0.548560in}}%
\pgfpathlineto{\pgfqpoint{3.170029in}{0.548535in}}%
\pgfpathlineto{\pgfqpoint{3.172298in}{0.543677in}}%
\pgfpathlineto{\pgfqpoint{3.174568in}{0.540830in}}%
\pgfpathlineto{\pgfqpoint{3.176837in}{0.536661in}}%
\pgfpathlineto{\pgfqpoint{3.179107in}{0.534872in}}%
\pgfpathlineto{\pgfqpoint{3.181376in}{0.538110in}}%
\pgfpathlineto{\pgfqpoint{3.183646in}{0.531748in}}%
\pgfpathlineto{\pgfqpoint{3.185915in}{0.528209in}}%
\pgfpathlineto{\pgfqpoint{3.188185in}{0.522756in}}%
\pgfpathlineto{\pgfqpoint{3.190454in}{0.521537in}}%
\pgfpathlineto{\pgfqpoint{3.192724in}{0.529461in}}%
\pgfpathlineto{\pgfqpoint{3.194993in}{0.518512in}}%
\pgfpathlineto{\pgfqpoint{3.197263in}{0.517909in}}%
\pgfpathlineto{\pgfqpoint{3.199532in}{0.508091in}}%
\pgfpathlineto{\pgfqpoint{3.201802in}{0.506219in}}%
\pgfpathlineto{\pgfqpoint{3.204071in}{0.515740in}}%
\pgfpathlineto{\pgfqpoint{3.206341in}{0.510185in}}%
\pgfpathlineto{\pgfqpoint{3.208610in}{0.511310in}}%
\pgfpathlineto{\pgfqpoint{3.210880in}{0.494828in}}%
\pgfpathlineto{\pgfqpoint{3.213149in}{0.494454in}}%
\pgfpathlineto{\pgfqpoint{3.215419in}{0.516251in}}%
\pgfpathlineto{\pgfqpoint{3.217688in}{0.505436in}}%
\pgfpathlineto{\pgfqpoint{3.219958in}{0.508007in}}%
\pgfpathlineto{\pgfqpoint{3.222227in}{0.479665in}}%
\pgfpathlineto{\pgfqpoint{3.224497in}{0.479076in}}%
\pgfpathlineto{\pgfqpoint{3.226766in}{0.506237in}}%
\pgfpathlineto{\pgfqpoint{3.229035in}{0.493247in}}%
\pgfpathlineto{\pgfqpoint{3.231305in}{0.505320in}}%
\pgfpathlineto{\pgfqpoint{3.233574in}{0.475245in}}%
\pgfpathlineto{\pgfqpoint{3.235844in}{0.462440in}}%
\pgfpathlineto{\pgfqpoint{3.238113in}{0.493368in}}%
\pgfpathlineto{\pgfqpoint{3.240383in}{0.497002in}}%
\pgfpathlineto{\pgfqpoint{3.242652in}{0.508719in}}%
\pgfpathlineto{\pgfqpoint{3.244922in}{0.472126in}}%
\pgfpathlineto{\pgfqpoint{3.247191in}{0.450666in}}%
\pgfpathlineto{\pgfqpoint{3.249461in}{0.504463in}}%
\pgfpathlineto{\pgfqpoint{3.251730in}{0.493853in}}%
\pgfpathlineto{\pgfqpoint{3.254000in}{0.510018in}}%
\pgfpathlineto{\pgfqpoint{3.256269in}{0.459438in}}%
\pgfpathlineto{\pgfqpoint{3.258539in}{0.442284in}}%
\pgfpathlineto{\pgfqpoint{3.260808in}{0.481756in}}%
\pgfpathlineto{\pgfqpoint{3.263078in}{0.481457in}}%
\pgfpathlineto{\pgfqpoint{3.265347in}{0.506431in}}%
\pgfpathlineto{\pgfqpoint{3.269886in}{0.448776in}}%
\pgfpathlineto{\pgfqpoint{3.272156in}{0.476595in}}%
\pgfpathlineto{\pgfqpoint{3.276695in}{0.514790in}}%
\pgfpathlineto{\pgfqpoint{3.281234in}{0.439575in}}%
\pgfpathlineto{\pgfqpoint{3.283503in}{0.488910in}}%
\pgfpathlineto{\pgfqpoint{3.285773in}{0.482191in}}%
\pgfpathlineto{\pgfqpoint{3.288042in}{0.511321in}}%
\pgfpathlineto{\pgfqpoint{3.290312in}{0.472824in}}%
\pgfpathlineto{\pgfqpoint{3.292581in}{0.468596in}}%
\pgfpathlineto{\pgfqpoint{3.294851in}{0.439234in}}%
\pgfpathlineto{\pgfqpoint{3.299389in}{0.507768in}}%
\pgfpathlineto{\pgfqpoint{3.306198in}{0.458998in}}%
\pgfpathlineto{\pgfqpoint{3.310737in}{0.518682in}}%
\pgfpathlineto{\pgfqpoint{3.313006in}{0.482727in}}%
\pgfpathlineto{\pgfqpoint{3.315276in}{0.476617in}}%
\pgfpathlineto{\pgfqpoint{3.317545in}{0.456822in}}%
\pgfpathlineto{\pgfqpoint{3.319815in}{0.460919in}}%
\pgfpathlineto{\pgfqpoint{3.322084in}{0.504811in}}%
\pgfpathlineto{\pgfqpoint{3.324354in}{0.482188in}}%
\pgfpathlineto{\pgfqpoint{3.326623in}{0.487428in}}%
\pgfpathlineto{\pgfqpoint{3.328893in}{0.409224in}}%
\pgfpathlineto{\pgfqpoint{3.333432in}{0.508384in}}%
\pgfpathlineto{\pgfqpoint{3.335701in}{0.498680in}}%
\pgfpathlineto{\pgfqpoint{3.337971in}{0.491517in}}%
\pgfpathlineto{\pgfqpoint{3.340240in}{0.430614in}}%
\pgfpathlineto{\pgfqpoint{3.344779in}{0.516289in}}%
\pgfpathlineto{\pgfqpoint{3.347049in}{0.483516in}}%
\pgfpathlineto{\pgfqpoint{3.349318in}{0.494767in}}%
\pgfpathlineto{\pgfqpoint{3.351588in}{0.371488in}}%
\pgfpathlineto{\pgfqpoint{3.353857in}{0.408426in}}%
\pgfpathlineto{\pgfqpoint{3.356127in}{0.490294in}}%
\pgfpathlineto{\pgfqpoint{3.358396in}{0.487315in}}%
\pgfpathlineto{\pgfqpoint{3.360666in}{0.499177in}}%
\pgfpathlineto{\pgfqpoint{3.362935in}{0.460558in}}%
\pgfpathlineto{\pgfqpoint{3.365204in}{0.442713in}}%
\pgfpathlineto{\pgfqpoint{3.367474in}{0.508499in}}%
\pgfpathlineto{\pgfqpoint{3.369743in}{0.500849in}}%
\pgfpathlineto{\pgfqpoint{3.372013in}{0.504952in}}%
\pgfpathlineto{\pgfqpoint{3.374282in}{0.338808in}}%
\pgfpathlineto{\pgfqpoint{3.378821in}{0.502751in}}%
\pgfpathlineto{\pgfqpoint{3.381091in}{0.479099in}}%
\pgfpathlineto{\pgfqpoint{3.383360in}{0.502041in}}%
\pgfpathlineto{\pgfqpoint{3.387899in}{0.421117in}}%
\pgfpathlineto{\pgfqpoint{3.390169in}{0.471374in}}%
\pgfpathlineto{\pgfqpoint{3.394708in}{0.508603in}}%
\pgfpathlineto{\pgfqpoint{3.396977in}{0.478619in}}%
\pgfpathlineto{\pgfqpoint{3.399247in}{0.388884in}}%
\pgfpathlineto{\pgfqpoint{3.401516in}{0.506259in}}%
\pgfpathlineto{\pgfqpoint{3.403786in}{0.496683in}}%
\pgfpathlineto{\pgfqpoint{3.406055in}{0.513241in}}%
\pgfpathlineto{\pgfqpoint{3.408325in}{0.441921in}}%
\pgfpathlineto{\pgfqpoint{3.410594in}{0.434118in}}%
\pgfpathlineto{\pgfqpoint{3.412864in}{0.475509in}}%
\pgfpathlineto{\pgfqpoint{3.415133in}{0.469289in}}%
\pgfpathlineto{\pgfqpoint{3.417403in}{0.502359in}}%
\pgfpathlineto{\pgfqpoint{3.421942in}{0.461380in}}%
\pgfpathlineto{\pgfqpoint{3.424211in}{0.448892in}}%
\pgfpathlineto{\pgfqpoint{3.426481in}{0.494070in}}%
\pgfpathlineto{\pgfqpoint{3.428750in}{0.517050in}}%
\pgfpathlineto{\pgfqpoint{3.431019in}{0.486087in}}%
\pgfpathlineto{\pgfqpoint{3.433289in}{0.434564in}}%
\pgfpathlineto{\pgfqpoint{3.435558in}{0.496704in}}%
\pgfpathlineto{\pgfqpoint{3.437828in}{0.483480in}}%
\pgfpathlineto{\pgfqpoint{3.440097in}{0.513978in}}%
\pgfpathlineto{\pgfqpoint{3.442367in}{0.465439in}}%
\pgfpathlineto{\pgfqpoint{3.444636in}{0.474412in}}%
\pgfpathlineto{\pgfqpoint{3.446906in}{0.401877in}}%
\pgfpathlineto{\pgfqpoint{3.451445in}{0.499652in}}%
\pgfpathlineto{\pgfqpoint{3.453714in}{0.494767in}}%
\pgfpathlineto{\pgfqpoint{3.455984in}{0.481652in}}%
\pgfpathlineto{\pgfqpoint{3.458253in}{0.428160in}}%
\pgfpathlineto{\pgfqpoint{3.460523in}{0.494051in}}%
\pgfpathlineto{\pgfqpoint{3.462792in}{0.523350in}}%
\pgfpathlineto{\pgfqpoint{3.465062in}{0.486785in}}%
\pgfpathlineto{\pgfqpoint{3.467331in}{0.477731in}}%
\pgfpathlineto{\pgfqpoint{3.469601in}{0.471088in}}%
\pgfpathlineto{\pgfqpoint{3.471870in}{0.453984in}}%
\pgfpathlineto{\pgfqpoint{3.474140in}{0.502382in}}%
\pgfpathlineto{\pgfqpoint{3.476409in}{0.474345in}}%
\pgfpathlineto{\pgfqpoint{3.478679in}{0.489112in}}%
\pgfpathlineto{\pgfqpoint{3.480948in}{0.446186in}}%
\pgfpathlineto{\pgfqpoint{3.483218in}{0.435605in}}%
\pgfpathlineto{\pgfqpoint{3.485487in}{0.499432in}}%
\pgfpathlineto{\pgfqpoint{3.487757in}{0.504424in}}%
\pgfpathlineto{\pgfqpoint{3.490026in}{0.496558in}}%
\pgfpathlineto{\pgfqpoint{3.492296in}{0.418359in}}%
\pgfpathlineto{\pgfqpoint{3.494565in}{0.485076in}}%
\pgfpathlineto{\pgfqpoint{3.496834in}{0.523986in}}%
\pgfpathlineto{\pgfqpoint{3.499104in}{0.482265in}}%
\pgfpathlineto{\pgfqpoint{3.501373in}{0.496987in}}%
\pgfpathlineto{\pgfqpoint{3.504543in}{0.315625in}}%
\pgfpathmoveto{\pgfqpoint{3.506624in}{0.315625in}}%
\pgfpathlineto{\pgfqpoint{3.508182in}{0.478124in}}%
\pgfpathlineto{\pgfqpoint{3.510451in}{0.478038in}}%
\pgfpathlineto{\pgfqpoint{3.512721in}{0.497584in}}%
\pgfpathlineto{\pgfqpoint{3.514990in}{0.478467in}}%
\pgfpathlineto{\pgfqpoint{3.517260in}{0.414311in}}%
\pgfpathlineto{\pgfqpoint{3.519529in}{0.504686in}}%
\pgfpathlineto{\pgfqpoint{3.521799in}{0.509606in}}%
\pgfpathlineto{\pgfqpoint{3.524068in}{0.508793in}}%
\pgfpathlineto{\pgfqpoint{3.526338in}{0.401114in}}%
\pgfpathlineto{\pgfqpoint{3.530877in}{0.511880in}}%
\pgfpathlineto{\pgfqpoint{3.533146in}{0.472172in}}%
\pgfpathlineto{\pgfqpoint{3.535416in}{0.501917in}}%
\pgfpathlineto{\pgfqpoint{3.537685in}{0.451826in}}%
\pgfpathlineto{\pgfqpoint{3.539955in}{0.445907in}}%
\pgfpathlineto{\pgfqpoint{3.542224in}{0.437140in}}%
\pgfpathlineto{\pgfqpoint{3.544494in}{0.482726in}}%
\pgfpathlineto{\pgfqpoint{3.546763in}{0.506479in}}%
\pgfpathlineto{\pgfqpoint{3.549033in}{0.493961in}}%
\pgfpathlineto{\pgfqpoint{3.551302in}{0.382296in}}%
\pgfpathlineto{\pgfqpoint{3.553572in}{0.511782in}}%
\pgfpathlineto{\pgfqpoint{3.555841in}{0.507140in}}%
\pgfpathlineto{\pgfqpoint{3.558111in}{0.517809in}}%
\pgfpathlineto{\pgfqpoint{3.560380in}{0.379010in}}%
\pgfpathlineto{\pgfqpoint{3.564919in}{0.482248in}}%
\pgfpathlineto{\pgfqpoint{3.567188in}{0.453862in}}%
\pgfpathlineto{\pgfqpoint{3.569458in}{0.495357in}}%
\pgfpathlineto{\pgfqpoint{3.571727in}{0.476590in}}%
\pgfpathlineto{\pgfqpoint{3.573997in}{0.470767in}}%
\pgfpathlineto{\pgfqpoint{3.576266in}{0.358316in}}%
\pgfpathlineto{\pgfqpoint{3.578536in}{0.491482in}}%
\pgfpathlineto{\pgfqpoint{3.580805in}{0.517587in}}%
\pgfpathlineto{\pgfqpoint{3.583075in}{0.499520in}}%
\pgfpathlineto{\pgfqpoint{3.585344in}{0.394556in}}%
\pgfpathlineto{\pgfqpoint{3.587614in}{0.513063in}}%
\pgfpathlineto{\pgfqpoint{3.589883in}{0.492931in}}%
\pgfpathlineto{\pgfqpoint{3.592153in}{0.519844in}}%
\pgfpathlineto{\pgfqpoint{3.594422in}{0.442455in}}%
\pgfpathlineto{\pgfqpoint{3.596692in}{0.475471in}}%
\pgfpathlineto{\pgfqpoint{3.598961in}{0.396722in}}%
\pgfpathlineto{\pgfqpoint{3.601231in}{0.421792in}}%
\pgfpathlineto{\pgfqpoint{3.603500in}{0.483275in}}%
\pgfpathlineto{\pgfqpoint{3.605770in}{0.492250in}}%
\pgfpathlineto{\pgfqpoint{3.608039in}{0.486409in}}%
\pgfpathlineto{\pgfqpoint{3.610309in}{0.415522in}}%
\pgfpathlineto{\pgfqpoint{3.612578in}{0.501562in}}%
\pgfpathlineto{\pgfqpoint{3.614848in}{0.528547in}}%
\pgfpathlineto{\pgfqpoint{3.619387in}{0.463651in}}%
\pgfpathlineto{\pgfqpoint{3.621656in}{0.499791in}}%
\pgfpathlineto{\pgfqpoint{3.623926in}{0.458262in}}%
\pgfpathlineto{\pgfqpoint{3.626195in}{0.506570in}}%
\pgfpathlineto{\pgfqpoint{3.628465in}{0.457736in}}%
\pgfpathlineto{\pgfqpoint{3.630734in}{0.486076in}}%
\pgfpathlineto{\pgfqpoint{3.633003in}{0.452459in}}%
\pgfpathlineto{\pgfqpoint{3.635273in}{0.343488in}}%
\pgfpathlineto{\pgfqpoint{3.637542in}{0.479405in}}%
\pgfpathlineto{\pgfqpoint{3.639812in}{0.507471in}}%
\pgfpathlineto{\pgfqpoint{3.642081in}{0.499304in}}%
\pgfpathlineto{\pgfqpoint{3.644351in}{0.456799in}}%
\pgfpathlineto{\pgfqpoint{3.646620in}{0.437726in}}%
\pgfpathlineto{\pgfqpoint{3.648890in}{0.498547in}}%
\pgfpathlineto{\pgfqpoint{3.651159in}{0.464715in}}%
\pgfpathlineto{\pgfqpoint{3.653429in}{0.491254in}}%
\pgfpathlineto{\pgfqpoint{3.655698in}{0.464606in}}%
\pgfpathlineto{\pgfqpoint{3.657968in}{0.366257in}}%
\pgfpathlineto{\pgfqpoint{3.660237in}{0.479146in}}%
\pgfpathlineto{\pgfqpoint{3.662507in}{0.507326in}}%
\pgfpathlineto{\pgfqpoint{3.664776in}{0.506412in}}%
\pgfpathlineto{\pgfqpoint{3.667046in}{0.417990in}}%
\pgfpathlineto{\pgfqpoint{3.669315in}{0.489230in}}%
\pgfpathlineto{\pgfqpoint{3.671585in}{0.530513in}}%
\pgfpathlineto{\pgfqpoint{3.673854in}{0.484830in}}%
\pgfpathlineto{\pgfqpoint{3.676124in}{0.503893in}}%
\pgfpathlineto{\pgfqpoint{3.678393in}{0.390258in}}%
\pgfpathlineto{\pgfqpoint{3.682932in}{0.459782in}}%
\pgfpathlineto{\pgfqpoint{3.685202in}{0.462863in}}%
\pgfpathlineto{\pgfqpoint{3.687471in}{0.493916in}}%
\pgfpathlineto{\pgfqpoint{3.689741in}{0.489737in}}%
\pgfpathlineto{\pgfqpoint{3.692010in}{0.415664in}}%
\pgfpathlineto{\pgfqpoint{3.694280in}{0.490409in}}%
\pgfpathlineto{\pgfqpoint{3.696549in}{0.516717in}}%
\pgfpathlineto{\pgfqpoint{3.698818in}{0.517171in}}%
\pgfpathlineto{\pgfqpoint{3.701088in}{0.359063in}}%
\pgfpathlineto{\pgfqpoint{3.703357in}{0.458341in}}%
\pgfpathlineto{\pgfqpoint{3.705627in}{0.519224in}}%
\pgfpathlineto{\pgfqpoint{3.707896in}{0.464987in}}%
\pgfpathlineto{\pgfqpoint{3.710166in}{0.504068in}}%
\pgfpathlineto{\pgfqpoint{3.712435in}{0.445710in}}%
\pgfpathlineto{\pgfqpoint{3.714705in}{0.463030in}}%
\pgfpathlineto{\pgfqpoint{3.716974in}{0.364587in}}%
\pgfpathlineto{\pgfqpoint{3.719244in}{0.465162in}}%
\pgfpathlineto{\pgfqpoint{3.721513in}{0.501447in}}%
\pgfpathlineto{\pgfqpoint{3.723783in}{0.506401in}}%
\pgfpathlineto{\pgfqpoint{3.726052in}{0.420494in}}%
\pgfpathlineto{\pgfqpoint{3.728322in}{0.511348in}}%
\pgfpathlineto{\pgfqpoint{3.730591in}{0.470679in}}%
\pgfpathlineto{\pgfqpoint{3.732861in}{0.507644in}}%
\pgfpathlineto{\pgfqpoint{3.735130in}{0.444959in}}%
\pgfpathlineto{\pgfqpoint{3.737400in}{0.471529in}}%
\pgfpathlineto{\pgfqpoint{3.739669in}{0.416940in}}%
\pgfpathlineto{\pgfqpoint{3.741939in}{0.445054in}}%
\pgfpathlineto{\pgfqpoint{3.744208in}{0.494468in}}%
\pgfpathlineto{\pgfqpoint{3.746478in}{0.508669in}}%
\pgfpathlineto{\pgfqpoint{3.748747in}{0.466578in}}%
\pgfpathlineto{\pgfqpoint{3.751017in}{0.499326in}}%
\pgfpathlineto{\pgfqpoint{3.753286in}{0.466369in}}%
\pgfpathlineto{\pgfqpoint{3.755556in}{0.506874in}}%
\pgfpathlineto{\pgfqpoint{3.757825in}{0.454448in}}%
\pgfpathlineto{\pgfqpoint{3.760095in}{0.478590in}}%
\pgfpathlineto{\pgfqpoint{3.762364in}{0.438104in}}%
\pgfpathlineto{\pgfqpoint{3.764633in}{0.432274in}}%
\pgfpathlineto{\pgfqpoint{3.766903in}{0.492610in}}%
\pgfpathlineto{\pgfqpoint{3.769172in}{0.509457in}}%
\pgfpathlineto{\pgfqpoint{3.771442in}{0.482959in}}%
\pgfpathlineto{\pgfqpoint{3.773711in}{0.488732in}}%
\pgfpathlineto{\pgfqpoint{3.775981in}{0.458894in}}%
\pgfpathlineto{\pgfqpoint{3.778250in}{0.504818in}}%
\pgfpathlineto{\pgfqpoint{3.780520in}{0.458660in}}%
\pgfpathlineto{\pgfqpoint{3.782789in}{0.484270in}}%
\pgfpathlineto{\pgfqpoint{3.785059in}{0.451922in}}%
\pgfpathlineto{\pgfqpoint{3.787328in}{0.400461in}}%
\pgfpathlineto{\pgfqpoint{3.789598in}{0.486834in}}%
\pgfpathlineto{\pgfqpoint{3.791867in}{0.509493in}}%
\pgfpathlineto{\pgfqpoint{3.794137in}{0.495108in}}%
\pgfpathlineto{\pgfqpoint{3.796406in}{0.442281in}}%
\pgfpathlineto{\pgfqpoint{3.798676in}{0.507398in}}%
\pgfpathlineto{\pgfqpoint{3.800945in}{0.534555in}}%
\pgfpathlineto{\pgfqpoint{3.803215in}{0.472202in}}%
\pgfpathlineto{\pgfqpoint{3.805484in}{0.490808in}}%
\pgfpathlineto{\pgfqpoint{3.807754in}{0.456841in}}%
\pgfpathlineto{\pgfqpoint{3.810023in}{0.344315in}}%
\pgfpathlineto{\pgfqpoint{3.812293in}{0.471490in}}%
\pgfpathlineto{\pgfqpoint{3.814562in}{0.463706in}}%
\pgfpathlineto{\pgfqpoint{3.816832in}{0.489463in}}%
\pgfpathlineto{\pgfqpoint{3.819101in}{0.484427in}}%
\pgfpathlineto{\pgfqpoint{3.821371in}{0.395071in}}%
\pgfpathlineto{\pgfqpoint{3.823640in}{0.501808in}}%
\pgfpathlineto{\pgfqpoint{3.825910in}{0.520669in}}%
\pgfpathlineto{\pgfqpoint{3.828179in}{0.504876in}}%
\pgfpathlineto{\pgfqpoint{3.830448in}{0.478366in}}%
\pgfpathlineto{\pgfqpoint{3.832718in}{0.490505in}}%
\pgfpathlineto{\pgfqpoint{3.834987in}{0.529868in}}%
\pgfpathlineto{\pgfqpoint{3.837257in}{0.456384in}}%
\pgfpathlineto{\pgfqpoint{3.839526in}{0.499672in}}%
\pgfpathlineto{\pgfqpoint{3.841796in}{0.414665in}}%
\pgfpathlineto{\pgfqpoint{3.844065in}{0.445801in}}%
\pgfpathlineto{\pgfqpoint{3.846335in}{0.392837in}}%
\pgfpathlineto{\pgfqpoint{3.848604in}{0.470963in}}%
\pgfpathlineto{\pgfqpoint{3.850874in}{0.498783in}}%
\pgfpathlineto{\pgfqpoint{3.853143in}{0.503216in}}%
\pgfpathlineto{\pgfqpoint{3.855413in}{0.437305in}}%
\pgfpathlineto{\pgfqpoint{3.857682in}{0.522679in}}%
\pgfpathlineto{\pgfqpoint{3.859952in}{0.470035in}}%
\pgfpathlineto{\pgfqpoint{3.862221in}{0.504378in}}%
\pgfpathlineto{\pgfqpoint{3.864491in}{0.419374in}}%
\pgfpathlineto{\pgfqpoint{3.866760in}{0.457352in}}%
\pgfpathlineto{\pgfqpoint{3.869030in}{0.318773in}}%
\pgfpathlineto{\pgfqpoint{3.871299in}{0.458808in}}%
\pgfpathlineto{\pgfqpoint{3.873569in}{0.495849in}}%
\pgfpathlineto{\pgfqpoint{3.875838in}{0.506929in}}%
\pgfpathlineto{\pgfqpoint{3.878108in}{0.387476in}}%
\pgfpathlineto{\pgfqpoint{3.880377in}{0.514367in}}%
\pgfpathlineto{\pgfqpoint{3.882647in}{0.471013in}}%
\pgfpathlineto{\pgfqpoint{3.884916in}{0.506253in}}%
\pgfpathlineto{\pgfqpoint{3.887186in}{0.438530in}}%
\pgfpathlineto{\pgfqpoint{3.889455in}{0.466231in}}%
\pgfpathlineto{\pgfqpoint{3.891725in}{0.399940in}}%
\pgfpathlineto{\pgfqpoint{3.896264in}{0.496319in}}%
\pgfpathlineto{\pgfqpoint{3.898533in}{0.508510in}}%
\pgfpathlineto{\pgfqpoint{3.900802in}{0.446456in}}%
\pgfpathlineto{\pgfqpoint{3.903072in}{0.507409in}}%
\pgfpathlineto{\pgfqpoint{3.905341in}{0.468935in}}%
\pgfpathlineto{\pgfqpoint{3.907611in}{0.507182in}}%
\pgfpathlineto{\pgfqpoint{3.909880in}{0.447960in}}%
\pgfpathlineto{\pgfqpoint{3.912150in}{0.473833in}}%
\pgfpathlineto{\pgfqpoint{3.914419in}{0.427498in}}%
\pgfpathlineto{\pgfqpoint{3.916689in}{0.441307in}}%
\pgfpathlineto{\pgfqpoint{3.918958in}{0.493989in}}%
\pgfpathlineto{\pgfqpoint{3.921228in}{0.509661in}}%
\pgfpathlineto{\pgfqpoint{3.923497in}{0.471728in}}%
\pgfpathlineto{\pgfqpoint{3.925767in}{0.497932in}}%
\pgfpathlineto{\pgfqpoint{3.928036in}{0.464015in}}%
\pgfpathlineto{\pgfqpoint{3.930306in}{0.506378in}}%
\pgfpathlineto{\pgfqpoint{3.932575in}{0.454510in}}%
\pgfpathlineto{\pgfqpoint{3.934845in}{0.480199in}}%
\pgfpathlineto{\pgfqpoint{3.937114in}{0.444199in}}%
\pgfpathlineto{\pgfqpoint{3.939384in}{0.422304in}}%
\pgfpathlineto{\pgfqpoint{3.941653in}{0.490086in}}%
\pgfpathlineto{\pgfqpoint{3.943923in}{0.510141in}}%
\pgfpathlineto{\pgfqpoint{3.946192in}{0.487087in}}%
\pgfpathlineto{\pgfqpoint{3.948462in}{0.485948in}}%
\pgfpathlineto{\pgfqpoint{3.950731in}{0.455255in}}%
\pgfpathlineto{\pgfqpoint{3.953001in}{0.503517in}}%
\pgfpathlineto{\pgfqpoint{3.955270in}{0.458451in}}%
\pgfpathlineto{\pgfqpoint{3.957540in}{0.485384in}}%
\pgfpathlineto{\pgfqpoint{3.959809in}{0.455939in}}%
\pgfpathlineto{\pgfqpoint{3.962079in}{0.377190in}}%
\pgfpathlineto{\pgfqpoint{3.964348in}{0.483942in}}%
\pgfpathlineto{\pgfqpoint{3.966617in}{0.509897in}}%
\pgfpathlineto{\pgfqpoint{3.968887in}{0.497986in}}%
\pgfpathlineto{\pgfqpoint{3.973426in}{0.439935in}}%
\pgfpathlineto{\pgfqpoint{3.975695in}{0.499121in}}%
\pgfpathlineto{\pgfqpoint{3.977965in}{0.460299in}}%
\pgfpathlineto{\pgfqpoint{3.980234in}{0.489406in}}%
\pgfpathlineto{\pgfqpoint{3.982504in}{0.464691in}}%
\pgfpathlineto{\pgfqpoint{3.984773in}{0.382321in}}%
\pgfpathlineto{\pgfqpoint{3.987043in}{0.474960in}}%
\pgfpathlineto{\pgfqpoint{3.989312in}{0.508854in}}%
\pgfpathlineto{\pgfqpoint{3.991582in}{0.506107in}}%
\pgfpathlineto{\pgfqpoint{3.993851in}{0.405860in}}%
\pgfpathlineto{\pgfqpoint{3.996121in}{0.496843in}}%
\pgfpathlineto{\pgfqpoint{3.998390in}{0.534115in}}%
\pgfpathlineto{\pgfqpoint{4.000660in}{0.482589in}}%
\pgfpathlineto{\pgfqpoint{4.002929in}{0.502673in}}%
\pgfpathlineto{\pgfqpoint{4.005199in}{0.422420in}}%
\pgfpathlineto{\pgfqpoint{4.007468in}{0.423653in}}%
\pgfpathlineto{\pgfqpoint{4.009738in}{0.457799in}}%
\pgfpathlineto{\pgfqpoint{4.012007in}{0.458709in}}%
\pgfpathlineto{\pgfqpoint{4.014277in}{0.491185in}}%
\pgfpathlineto{\pgfqpoint{4.016546in}{0.490498in}}%
\pgfpathlineto{\pgfqpoint{4.018816in}{0.415000in}}%
\pgfpathlineto{\pgfqpoint{4.021085in}{0.490407in}}%
\pgfpathlineto{\pgfqpoint{4.023355in}{0.520002in}}%
\pgfpathlineto{\pgfqpoint{4.025624in}{0.516555in}}%
\pgfpathlineto{\pgfqpoint{4.027894in}{0.428915in}}%
\pgfpathlineto{\pgfqpoint{4.032432in}{0.524560in}}%
\pgfpathlineto{\pgfqpoint{4.034702in}{0.460749in}}%
\pgfpathlineto{\pgfqpoint{4.036971in}{0.503537in}}%
\pgfpathlineto{\pgfqpoint{4.039241in}{0.437646in}}%
\pgfpathlineto{\pgfqpoint{4.041510in}{0.461639in}}%
\pgfpathlineto{\pgfqpoint{4.043780in}{0.383561in}}%
\pgfpathlineto{\pgfqpoint{4.046049in}{0.462102in}}%
\pgfpathlineto{\pgfqpoint{4.048319in}{0.499280in}}%
\pgfpathlineto{\pgfqpoint{4.050588in}{0.508305in}}%
\pgfpathlineto{\pgfqpoint{4.052858in}{0.399392in}}%
\pgfpathlineto{\pgfqpoint{4.055127in}{0.515286in}}%
\pgfpathlineto{\pgfqpoint{4.057397in}{0.469194in}}%
\pgfpathlineto{\pgfqpoint{4.059666in}{0.506985in}}%
\pgfpathlineto{\pgfqpoint{4.061936in}{0.438642in}}%
\pgfpathlineto{\pgfqpoint{4.064205in}{0.469558in}}%
\pgfpathlineto{\pgfqpoint{4.066475in}{0.419146in}}%
\pgfpathlineto{\pgfqpoint{4.068744in}{0.444281in}}%
\pgfpathlineto{\pgfqpoint{4.071014in}{0.493391in}}%
\pgfpathlineto{\pgfqpoint{4.073283in}{0.510144in}}%
\pgfpathlineto{\pgfqpoint{4.075553in}{0.459577in}}%
\pgfpathlineto{\pgfqpoint{4.077822in}{0.504796in}}%
\pgfpathlineto{\pgfqpoint{4.080092in}{0.466365in}}%
\pgfpathlineto{\pgfqpoint{4.082361in}{0.506734in}}%
\pgfpathlineto{\pgfqpoint{4.084631in}{0.449301in}}%
\pgfpathlineto{\pgfqpoint{4.086900in}{0.476319in}}%
\pgfpathlineto{\pgfqpoint{4.089170in}{0.437679in}}%
\pgfpathlineto{\pgfqpoint{4.091439in}{0.432674in}}%
\pgfpathlineto{\pgfqpoint{4.093709in}{0.491803in}}%
\pgfpathlineto{\pgfqpoint{4.095978in}{0.510843in}}%
\pgfpathlineto{\pgfqpoint{4.098247in}{0.478260in}}%
\pgfpathlineto{\pgfqpoint{4.100517in}{0.495493in}}%
\pgfpathlineto{\pgfqpoint{4.102786in}{0.460452in}}%
\pgfpathlineto{\pgfqpoint{4.105056in}{0.505323in}}%
\pgfpathlineto{\pgfqpoint{4.107325in}{0.454479in}}%
\pgfpathlineto{\pgfqpoint{4.109595in}{0.481927in}}%
\pgfpathlineto{\pgfqpoint{4.111864in}{0.450542in}}%
\pgfpathlineto{\pgfqpoint{4.114134in}{0.405701in}}%
\pgfpathlineto{\pgfqpoint{4.116403in}{0.486699in}}%
\pgfpathlineto{\pgfqpoint{4.118673in}{0.510917in}}%
\pgfpathlineto{\pgfqpoint{4.120942in}{0.491674in}}%
\pgfpathlineto{\pgfqpoint{4.123212in}{0.482461in}}%
\pgfpathlineto{\pgfqpoint{4.125481in}{0.449913in}}%
\pgfpathlineto{\pgfqpoint{4.127751in}{0.501760in}}%
\pgfpathlineto{\pgfqpoint{4.130020in}{0.457766in}}%
\pgfpathlineto{\pgfqpoint{4.132290in}{0.486495in}}%
\pgfpathlineto{\pgfqpoint{4.134559in}{0.460232in}}%
\pgfpathlineto{\pgfqpoint{4.136095in}{0.315625in}}%
\pgfpathmoveto{\pgfqpoint{4.137502in}{0.315625in}}%
\pgfpathlineto{\pgfqpoint{4.139098in}{0.479622in}}%
\pgfpathlineto{\pgfqpoint{4.141368in}{0.510314in}}%
\pgfpathlineto{\pgfqpoint{4.143637in}{0.501334in}}%
\pgfpathlineto{\pgfqpoint{4.148176in}{0.430904in}}%
\pgfpathlineto{\pgfqpoint{4.150446in}{0.497013in}}%
\pgfpathlineto{\pgfqpoint{4.152715in}{0.458941in}}%
\pgfpathlineto{\pgfqpoint{4.154985in}{0.489939in}}%
\pgfpathlineto{\pgfqpoint{4.157254in}{0.467658in}}%
\pgfpathlineto{\pgfqpoint{4.159524in}{0.408316in}}%
\pgfpathlineto{\pgfqpoint{4.161793in}{0.469089in}}%
\pgfpathlineto{\pgfqpoint{4.164062in}{0.508908in}}%
\pgfpathlineto{\pgfqpoint{4.166332in}{0.508641in}}%
\pgfpathlineto{\pgfqpoint{4.168601in}{0.431179in}}%
\pgfpathlineto{\pgfqpoint{4.170871in}{0.494768in}}%
\pgfpathlineto{\pgfqpoint{4.173140in}{0.534389in}}%
\pgfpathlineto{\pgfqpoint{4.175410in}{0.483641in}}%
\pgfpathlineto{\pgfqpoint{4.177679in}{0.505082in}}%
\pgfpathlineto{\pgfqpoint{4.179949in}{0.411894in}}%
\pgfpathlineto{\pgfqpoint{4.184488in}{0.451638in}}%
\pgfpathlineto{\pgfqpoint{4.186757in}{0.455462in}}%
\pgfpathlineto{\pgfqpoint{4.189027in}{0.490491in}}%
\pgfpathlineto{\pgfqpoint{4.191296in}{0.492169in}}%
\pgfpathlineto{\pgfqpoint{4.193566in}{0.431381in}}%
\pgfpathlineto{\pgfqpoint{4.195835in}{0.486219in}}%
\pgfpathlineto{\pgfqpoint{4.198105in}{0.520440in}}%
\pgfpathlineto{\pgfqpoint{4.200374in}{0.519084in}}%
\pgfpathlineto{\pgfqpoint{4.202644in}{0.409307in}}%
\pgfpathlineto{\pgfqpoint{4.207183in}{0.523809in}}%
\pgfpathlineto{\pgfqpoint{4.209452in}{0.459470in}}%
\pgfpathlineto{\pgfqpoint{4.211722in}{0.503949in}}%
\pgfpathlineto{\pgfqpoint{4.213991in}{0.439984in}}%
\pgfpathlineto{\pgfqpoint{4.216261in}{0.464945in}}%
\pgfpathlineto{\pgfqpoint{4.218530in}{0.407394in}}%
\pgfpathlineto{\pgfqpoint{4.223069in}{0.498488in}}%
\pgfpathlineto{\pgfqpoint{4.225339in}{0.509999in}}%
\pgfpathlineto{\pgfqpoint{4.227608in}{0.424976in}}%
\pgfpathlineto{\pgfqpoint{4.229878in}{0.513985in}}%
\pgfpathlineto{\pgfqpoint{4.232147in}{0.467616in}}%
\pgfpathlineto{\pgfqpoint{4.234416in}{0.507075in}}%
\pgfpathlineto{\pgfqpoint{4.236686in}{0.440214in}}%
\pgfpathlineto{\pgfqpoint{4.238955in}{0.471894in}}%
\pgfpathlineto{\pgfqpoint{4.241225in}{0.429399in}}%
\pgfpathlineto{\pgfqpoint{4.243494in}{0.437600in}}%
\pgfpathlineto{\pgfqpoint{4.245764in}{0.491545in}}%
\pgfpathlineto{\pgfqpoint{4.248033in}{0.511251in}}%
\pgfpathlineto{\pgfqpoint{4.250303in}{0.467074in}}%
\pgfpathlineto{\pgfqpoint{4.252572in}{0.502989in}}%
\pgfpathlineto{\pgfqpoint{4.254842in}{0.463987in}}%
\pgfpathlineto{\pgfqpoint{4.257111in}{0.506227in}}%
\pgfpathlineto{\pgfqpoint{4.259381in}{0.449797in}}%
\pgfpathlineto{\pgfqpoint{4.261650in}{0.478049in}}%
\pgfpathlineto{\pgfqpoint{4.263920in}{0.444073in}}%
\pgfpathlineto{\pgfqpoint{4.266189in}{0.423187in}}%
\pgfpathlineto{\pgfqpoint{4.268459in}{0.489463in}}%
\pgfpathlineto{\pgfqpoint{4.270728in}{0.511664in}}%
\pgfpathlineto{\pgfqpoint{4.272998in}{0.483097in}}%
\pgfpathlineto{\pgfqpoint{4.275267in}{0.493199in}}%
\pgfpathlineto{\pgfqpoint{4.277537in}{0.457070in}}%
\pgfpathlineto{\pgfqpoint{4.279806in}{0.504158in}}%
\pgfpathlineto{\pgfqpoint{4.282076in}{0.454287in}}%
\pgfpathlineto{\pgfqpoint{4.284345in}{0.483115in}}%
\pgfpathlineto{\pgfqpoint{4.286615in}{0.454800in}}%
\pgfpathlineto{\pgfqpoint{4.288884in}{0.384748in}}%
\pgfpathlineto{\pgfqpoint{4.291154in}{0.483605in}}%
\pgfpathlineto{\pgfqpoint{4.293423in}{0.511427in}}%
\pgfpathlineto{\pgfqpoint{4.297962in}{0.479367in}}%
\pgfpathlineto{\pgfqpoint{4.300231in}{0.444753in}}%
\pgfpathlineto{\pgfqpoint{4.302501in}{0.500221in}}%
\pgfpathlineto{\pgfqpoint{4.304770in}{0.457002in}}%
\pgfpathlineto{\pgfqpoint{4.307040in}{0.487205in}}%
\pgfpathlineto{\pgfqpoint{4.309309in}{0.463192in}}%
\pgfpathlineto{\pgfqpoint{4.311579in}{0.364332in}}%
\pgfpathlineto{\pgfqpoint{4.313848in}{0.475562in}}%
\pgfpathlineto{\pgfqpoint{4.316118in}{0.510536in}}%
\pgfpathlineto{\pgfqpoint{4.318387in}{0.503868in}}%
\pgfpathlineto{\pgfqpoint{4.322926in}{0.421393in}}%
\pgfpathlineto{\pgfqpoint{4.325196in}{0.495172in}}%
\pgfpathlineto{\pgfqpoint{4.327465in}{0.457638in}}%
\pgfpathlineto{\pgfqpoint{4.329735in}{0.490193in}}%
\pgfpathlineto{\pgfqpoint{4.332004in}{0.469733in}}%
\pgfpathlineto{\pgfqpoint{4.334274in}{0.421695in}}%
\pgfpathlineto{\pgfqpoint{4.336543in}{0.482000in}}%
\pgfpathlineto{\pgfqpoint{4.336543in}{0.482000in}}%
\pgfusepath{stroke}%
\end{pgfscope}%
\begin{pgfscope}%
\pgfpathrectangle{\pgfqpoint{2.520957in}{0.320625in}}{\pgfqpoint{1.815586in}{1.164375in}}%
\pgfusepath{clip}%
\pgfsetrectcap%
\pgfsetroundjoin%
\pgfsetlinewidth{1.154312pt}%
\definecolor{currentstroke}{rgb}{0.432967,0.117855,0.506160}%
\pgfsetstrokecolor{currentstroke}%
\pgfsetdash{}{0pt}%
\pgfpathmoveto{\pgfqpoint{2.520957in}{1.451682in}}%
\pgfpathlineto{\pgfqpoint{2.525496in}{1.442221in}}%
\pgfpathlineto{\pgfqpoint{2.545921in}{1.404093in}}%
\pgfpathlineto{\pgfqpoint{2.609467in}{1.281334in}}%
\pgfpathlineto{\pgfqpoint{2.632161in}{1.241097in}}%
\pgfpathlineto{\pgfqpoint{2.657126in}{1.199057in}}%
\pgfpathlineto{\pgfqpoint{2.686629in}{1.151715in}}%
\pgfpathlineto{\pgfqpoint{2.720671in}{1.099389in}}%
\pgfpathlineto{\pgfqpoint{2.761522in}{1.038797in}}%
\pgfpathlineto{\pgfqpoint{2.813720in}{0.963489in}}%
\pgfpathlineto{\pgfqpoint{2.888613in}{0.857548in}}%
\pgfpathlineto{\pgfqpoint{3.008896in}{0.689530in}}%
\pgfpathlineto{\pgfqpoint{3.088328in}{0.579486in}}%
\pgfpathlineto{\pgfqpoint{3.090597in}{0.577612in}}%
\pgfpathlineto{\pgfqpoint{3.092867in}{0.573169in}}%
\pgfpathlineto{\pgfqpoint{3.106483in}{0.555185in}}%
\pgfpathlineto{\pgfqpoint{3.111022in}{0.547957in}}%
\pgfpathlineto{\pgfqpoint{3.113292in}{0.548224in}}%
\pgfpathlineto{\pgfqpoint{3.117831in}{0.541304in}}%
\pgfpathlineto{\pgfqpoint{3.120100in}{0.535349in}}%
\pgfpathlineto{\pgfqpoint{3.122370in}{0.532933in}}%
\pgfpathlineto{\pgfqpoint{3.124639in}{0.537546in}}%
\pgfpathlineto{\pgfqpoint{3.126909in}{0.527711in}}%
\pgfpathlineto{\pgfqpoint{3.129178in}{0.528402in}}%
\pgfpathlineto{\pgfqpoint{3.131448in}{0.519794in}}%
\pgfpathlineto{\pgfqpoint{3.133717in}{0.516348in}}%
\pgfpathlineto{\pgfqpoint{3.135987in}{0.515899in}}%
\pgfpathlineto{\pgfqpoint{3.138256in}{0.516345in}}%
\pgfpathlineto{\pgfqpoint{3.140526in}{0.519813in}}%
\pgfpathlineto{\pgfqpoint{3.142795in}{0.509216in}}%
\pgfpathlineto{\pgfqpoint{3.145065in}{0.501820in}}%
\pgfpathlineto{\pgfqpoint{3.147334in}{0.521718in}}%
\pgfpathlineto{\pgfqpoint{3.149604in}{0.514224in}}%
\pgfpathlineto{\pgfqpoint{3.151873in}{0.517214in}}%
\pgfpathlineto{\pgfqpoint{3.154143in}{0.488027in}}%
\pgfpathlineto{\pgfqpoint{3.156412in}{0.484884in}}%
\pgfpathlineto{\pgfqpoint{3.158682in}{0.502809in}}%
\pgfpathlineto{\pgfqpoint{3.160951in}{0.488580in}}%
\pgfpathlineto{\pgfqpoint{3.163220in}{0.505338in}}%
\pgfpathlineto{\pgfqpoint{3.165490in}{0.486527in}}%
\pgfpathlineto{\pgfqpoint{3.167759in}{0.476254in}}%
\pgfpathlineto{\pgfqpoint{3.170029in}{0.473160in}}%
\pgfpathlineto{\pgfqpoint{3.172298in}{0.499237in}}%
\pgfpathlineto{\pgfqpoint{3.174568in}{0.515144in}}%
\pgfpathlineto{\pgfqpoint{3.176837in}{0.490731in}}%
\pgfpathlineto{\pgfqpoint{3.179107in}{0.455568in}}%
\pgfpathlineto{\pgfqpoint{3.181376in}{0.514615in}}%
\pgfpathlineto{\pgfqpoint{3.183646in}{0.496336in}}%
\pgfpathlineto{\pgfqpoint{3.185915in}{0.516152in}}%
\pgfpathlineto{\pgfqpoint{3.188185in}{0.453187in}}%
\pgfpathlineto{\pgfqpoint{3.190454in}{0.464919in}}%
\pgfpathlineto{\pgfqpoint{3.192724in}{0.457465in}}%
\pgfpathlineto{\pgfqpoint{3.194993in}{0.460239in}}%
\pgfpathlineto{\pgfqpoint{3.197263in}{0.496375in}}%
\pgfpathlineto{\pgfqpoint{3.199532in}{0.490436in}}%
\pgfpathlineto{\pgfqpoint{3.201802in}{0.473406in}}%
\pgfpathlineto{\pgfqpoint{3.204071in}{0.441000in}}%
\pgfpathlineto{\pgfqpoint{3.206341in}{0.502574in}}%
\pgfpathlineto{\pgfqpoint{3.208610in}{0.523339in}}%
\pgfpathlineto{\pgfqpoint{3.213149in}{0.438026in}}%
\pgfpathlineto{\pgfqpoint{3.215419in}{0.502351in}}%
\pgfpathlineto{\pgfqpoint{3.217688in}{0.472605in}}%
\pgfpathlineto{\pgfqpoint{3.219958in}{0.510389in}}%
\pgfpathlineto{\pgfqpoint{3.222227in}{0.461981in}}%
\pgfpathlineto{\pgfqpoint{3.224497in}{0.480290in}}%
\pgfpathlineto{\pgfqpoint{3.226766in}{0.420547in}}%
\pgfpathlineto{\pgfqpoint{3.229035in}{0.439856in}}%
\pgfpathlineto{\pgfqpoint{3.231305in}{0.494498in}}%
\pgfpathlineto{\pgfqpoint{3.233574in}{0.503511in}}%
\pgfpathlineto{\pgfqpoint{3.235844in}{0.485998in}}%
\pgfpathlineto{\pgfqpoint{3.238113in}{0.449501in}}%
\pgfpathlineto{\pgfqpoint{3.240383in}{0.504132in}}%
\pgfpathlineto{\pgfqpoint{3.242652in}{0.529707in}}%
\pgfpathlineto{\pgfqpoint{3.244922in}{0.476369in}}%
\pgfpathlineto{\pgfqpoint{3.247191in}{0.480916in}}%
\pgfpathlineto{\pgfqpoint{3.249461in}{0.469342in}}%
\pgfpathlineto{\pgfqpoint{3.251730in}{0.420189in}}%
\pgfpathlineto{\pgfqpoint{3.254000in}{0.488362in}}%
\pgfpathlineto{\pgfqpoint{3.256269in}{0.469649in}}%
\pgfpathlineto{\pgfqpoint{3.258539in}{0.488959in}}%
\pgfpathlineto{\pgfqpoint{3.260808in}{0.469911in}}%
\pgfpathlineto{\pgfqpoint{3.263078in}{0.432124in}}%
\pgfpathlineto{\pgfqpoint{3.265347in}{0.502850in}}%
\pgfpathlineto{\pgfqpoint{3.267617in}{0.513937in}}%
\pgfpathlineto{\pgfqpoint{3.269886in}{0.497270in}}%
\pgfpathlineto{\pgfqpoint{3.272156in}{0.470814in}}%
\pgfpathlineto{\pgfqpoint{3.274425in}{0.492574in}}%
\pgfpathlineto{\pgfqpoint{3.276695in}{0.528406in}}%
\pgfpathlineto{\pgfqpoint{3.278964in}{0.466498in}}%
\pgfpathlineto{\pgfqpoint{3.281234in}{0.496929in}}%
\pgfpathlineto{\pgfqpoint{3.283503in}{0.359255in}}%
\pgfpathlineto{\pgfqpoint{3.285773in}{0.416823in}}%
\pgfpathlineto{\pgfqpoint{3.288042in}{0.451753in}}%
\pgfpathlineto{\pgfqpoint{3.292581in}{0.497511in}}%
\pgfpathlineto{\pgfqpoint{3.294851in}{0.491415in}}%
\pgfpathlineto{\pgfqpoint{3.297120in}{0.445591in}}%
\pgfpathlineto{\pgfqpoint{3.299389in}{0.517786in}}%
\pgfpathlineto{\pgfqpoint{3.301659in}{0.516809in}}%
\pgfpathlineto{\pgfqpoint{3.303928in}{0.507735in}}%
\pgfpathlineto{\pgfqpoint{3.308467in}{0.446748in}}%
\pgfpathlineto{\pgfqpoint{3.310737in}{0.509758in}}%
\pgfpathlineto{\pgfqpoint{3.313006in}{0.453679in}}%
\pgfpathlineto{\pgfqpoint{3.315276in}{0.495229in}}%
\pgfpathlineto{\pgfqpoint{3.317545in}{0.457911in}}%
\pgfpathlineto{\pgfqpoint{3.319815in}{0.454445in}}%
\pgfpathlineto{\pgfqpoint{3.322084in}{0.329979in}}%
\pgfpathlineto{\pgfqpoint{3.324354in}{0.488913in}}%
\pgfpathlineto{\pgfqpoint{3.326623in}{0.509998in}}%
\pgfpathlineto{\pgfqpoint{3.328893in}{0.500381in}}%
\pgfpathlineto{\pgfqpoint{3.331162in}{0.463840in}}%
\pgfpathlineto{\pgfqpoint{3.333432in}{0.530234in}}%
\pgfpathlineto{\pgfqpoint{3.335701in}{0.507499in}}%
\pgfpathlineto{\pgfqpoint{3.337971in}{0.516729in}}%
\pgfpathlineto{\pgfqpoint{3.340240in}{0.438897in}}%
\pgfpathlineto{\pgfqpoint{3.342510in}{0.442827in}}%
\pgfpathlineto{\pgfqpoint{3.344779in}{0.468172in}}%
\pgfpathlineto{\pgfqpoint{3.347049in}{0.432034in}}%
\pgfpathlineto{\pgfqpoint{3.349318in}{0.482223in}}%
\pgfpathlineto{\pgfqpoint{3.351588in}{0.480885in}}%
\pgfpathlineto{\pgfqpoint{3.353857in}{0.471543in}}%
\pgfpathlineto{\pgfqpoint{3.356127in}{0.409351in}}%
\pgfpathlineto{\pgfqpoint{3.358396in}{0.506919in}}%
\pgfpathlineto{\pgfqpoint{3.360666in}{0.523956in}}%
\pgfpathlineto{\pgfqpoint{3.362935in}{0.392441in}}%
\pgfpathlineto{\pgfqpoint{3.365204in}{0.458023in}}%
\pgfpathlineto{\pgfqpoint{3.367474in}{0.463748in}}%
\pgfpathlineto{\pgfqpoint{3.369743in}{0.416610in}}%
\pgfpathlineto{\pgfqpoint{3.372013in}{0.477618in}}%
\pgfpathlineto{\pgfqpoint{3.374282in}{0.478860in}}%
\pgfpathlineto{\pgfqpoint{3.376552in}{0.479121in}}%
\pgfpathlineto{\pgfqpoint{3.378821in}{0.454027in}}%
\pgfpathlineto{\pgfqpoint{3.381091in}{0.499462in}}%
\pgfpathlineto{\pgfqpoint{3.383360in}{0.525009in}}%
\pgfpathlineto{\pgfqpoint{3.385630in}{0.441763in}}%
\pgfpathlineto{\pgfqpoint{3.387899in}{0.472120in}}%
\pgfpathlineto{\pgfqpoint{3.390169in}{0.448743in}}%
\pgfpathlineto{\pgfqpoint{3.392438in}{0.399490in}}%
\pgfpathlineto{\pgfqpoint{3.394708in}{0.474438in}}%
\pgfpathlineto{\pgfqpoint{3.399247in}{0.485189in}}%
\pgfpathlineto{\pgfqpoint{3.401516in}{0.466150in}}%
\pgfpathlineto{\pgfqpoint{3.406055in}{0.525700in}}%
\pgfpathlineto{\pgfqpoint{3.408325in}{0.457752in}}%
\pgfpathlineto{\pgfqpoint{3.410594in}{0.482405in}}%
\pgfpathlineto{\pgfqpoint{3.412864in}{0.428799in}}%
\pgfpathlineto{\pgfqpoint{3.415133in}{0.347619in}}%
\pgfpathlineto{\pgfqpoint{3.417403in}{0.469008in}}%
\pgfpathlineto{\pgfqpoint{3.421942in}{0.490235in}}%
\pgfpathlineto{\pgfqpoint{3.424211in}{0.477021in}}%
\pgfpathlineto{\pgfqpoint{3.426481in}{0.483632in}}%
\pgfpathlineto{\pgfqpoint{3.428750in}{0.525020in}}%
\pgfpathlineto{\pgfqpoint{3.431019in}{0.467552in}}%
\pgfpathlineto{\pgfqpoint{3.433289in}{0.490572in}}%
\pgfpathlineto{\pgfqpoint{3.435558in}{0.388460in}}%
\pgfpathlineto{\pgfqpoint{3.437828in}{0.373670in}}%
\pgfpathlineto{\pgfqpoint{3.440097in}{0.461615in}}%
\pgfpathlineto{\pgfqpoint{3.444636in}{0.494398in}}%
\pgfpathlineto{\pgfqpoint{3.446906in}{0.485105in}}%
\pgfpathlineto{\pgfqpoint{3.449175in}{0.470433in}}%
\pgfpathlineto{\pgfqpoint{3.451445in}{0.522919in}}%
\pgfpathlineto{\pgfqpoint{3.453714in}{0.473330in}}%
\pgfpathlineto{\pgfqpoint{3.455984in}{0.496994in}}%
\pgfpathlineto{\pgfqpoint{3.458253in}{0.362127in}}%
\pgfpathlineto{\pgfqpoint{3.460523in}{0.416639in}}%
\pgfpathlineto{\pgfqpoint{3.462792in}{0.451091in}}%
\pgfpathlineto{\pgfqpoint{3.467331in}{0.497603in}}%
\pgfpathlineto{\pgfqpoint{3.469601in}{0.491474in}}%
\pgfpathlineto{\pgfqpoint{3.471870in}{0.449125in}}%
\pgfpathlineto{\pgfqpoint{3.474140in}{0.518840in}}%
\pgfpathlineto{\pgfqpoint{3.476409in}{0.516740in}}%
\pgfpathlineto{\pgfqpoint{3.478679in}{0.507132in}}%
\pgfpathlineto{\pgfqpoint{3.483218in}{0.446064in}}%
\pgfpathlineto{\pgfqpoint{3.485487in}{0.509325in}}%
\pgfpathlineto{\pgfqpoint{3.487757in}{0.452978in}}%
\pgfpathlineto{\pgfqpoint{3.490026in}{0.494758in}}%
\pgfpathlineto{\pgfqpoint{3.492296in}{0.458350in}}%
\pgfpathlineto{\pgfqpoint{3.494565in}{0.454061in}}%
\pgfpathlineto{\pgfqpoint{3.496834in}{0.333362in}}%
\pgfpathlineto{\pgfqpoint{3.499104in}{0.489870in}}%
\pgfpathlineto{\pgfqpoint{3.501373in}{0.510260in}}%
\pgfpathlineto{\pgfqpoint{3.503643in}{0.499857in}}%
\pgfpathlineto{\pgfqpoint{3.505912in}{0.467047in}}%
\pgfpathlineto{\pgfqpoint{3.508182in}{0.531121in}}%
\pgfpathlineto{\pgfqpoint{3.510451in}{0.506868in}}%
\pgfpathlineto{\pgfqpoint{3.512721in}{0.516315in}}%
\pgfpathlineto{\pgfqpoint{3.514990in}{0.440442in}}%
\pgfpathlineto{\pgfqpoint{3.517260in}{0.443037in}}%
\pgfpathlineto{\pgfqpoint{3.519529in}{0.467146in}}%
\pgfpathlineto{\pgfqpoint{3.521799in}{0.431868in}}%
\pgfpathlineto{\pgfqpoint{3.524068in}{0.481926in}}%
\pgfpathlineto{\pgfqpoint{3.526338in}{0.481393in}}%
\pgfpathlineto{\pgfqpoint{3.528607in}{0.471116in}}%
\pgfpathlineto{\pgfqpoint{3.530877in}{0.403431in}}%
\pgfpathlineto{\pgfqpoint{3.533146in}{0.508034in}}%
\pgfpathlineto{\pgfqpoint{3.535416in}{0.524154in}}%
\pgfpathlineto{\pgfqpoint{3.537685in}{0.371489in}}%
\pgfpathlineto{\pgfqpoint{3.539955in}{0.456940in}}%
\pgfpathlineto{\pgfqpoint{3.542224in}{0.464686in}}%
\pgfpathlineto{\pgfqpoint{3.544494in}{0.414895in}}%
\pgfpathlineto{\pgfqpoint{3.546763in}{0.476765in}}%
\pgfpathlineto{\pgfqpoint{3.549033in}{0.478393in}}%
\pgfpathlineto{\pgfqpoint{3.551302in}{0.478761in}}%
\pgfpathlineto{\pgfqpoint{3.553572in}{0.455274in}}%
\pgfpathlineto{\pgfqpoint{3.555841in}{0.500076in}}%
\pgfpathlineto{\pgfqpoint{3.558111in}{0.525133in}}%
\pgfpathlineto{\pgfqpoint{3.560380in}{0.438595in}}%
\pgfpathlineto{\pgfqpoint{3.562649in}{0.471303in}}%
\pgfpathlineto{\pgfqpoint{3.564919in}{0.449804in}}%
\pgfpathlineto{\pgfqpoint{3.567188in}{0.398667in}}%
\pgfpathlineto{\pgfqpoint{3.569458in}{0.473875in}}%
\pgfpathlineto{\pgfqpoint{3.573997in}{0.484780in}}%
\pgfpathlineto{\pgfqpoint{3.576266in}{0.466344in}}%
\pgfpathlineto{\pgfqpoint{3.580805in}{0.525943in}}%
\pgfpathlineto{\pgfqpoint{3.583075in}{0.455734in}}%
\pgfpathlineto{\pgfqpoint{3.585344in}{0.481645in}}%
\pgfpathlineto{\pgfqpoint{3.587614in}{0.431040in}}%
\pgfpathlineto{\pgfqpoint{3.589883in}{0.348146in}}%
\pgfpathlineto{\pgfqpoint{3.592153in}{0.468558in}}%
\pgfpathlineto{\pgfqpoint{3.596692in}{0.489796in}}%
\pgfpathlineto{\pgfqpoint{3.598961in}{0.477104in}}%
\pgfpathlineto{\pgfqpoint{3.601231in}{0.484844in}}%
\pgfpathlineto{\pgfqpoint{3.603500in}{0.525372in}}%
\pgfpathlineto{\pgfqpoint{3.605770in}{0.466275in}}%
\pgfpathlineto{\pgfqpoint{3.608039in}{0.489906in}}%
\pgfpathlineto{\pgfqpoint{3.610309in}{0.394597in}}%
\pgfpathlineto{\pgfqpoint{3.612578in}{0.371231in}}%
\pgfpathlineto{\pgfqpoint{3.614848in}{0.461363in}}%
\pgfpathlineto{\pgfqpoint{3.619387in}{0.493975in}}%
\pgfpathlineto{\pgfqpoint{3.621656in}{0.485081in}}%
\pgfpathlineto{\pgfqpoint{3.623926in}{0.472205in}}%
\pgfpathlineto{\pgfqpoint{3.626195in}{0.523454in}}%
\pgfpathlineto{\pgfqpoint{3.628465in}{0.472485in}}%
\pgfpathlineto{\pgfqpoint{3.630734in}{0.496426in}}%
\pgfpathlineto{\pgfqpoint{3.633003in}{0.346298in}}%
\pgfpathlineto{\pgfqpoint{3.635273in}{0.415042in}}%
\pgfpathlineto{\pgfqpoint{3.637542in}{0.451105in}}%
\pgfpathlineto{\pgfqpoint{3.642081in}{0.497219in}}%
\pgfpathlineto{\pgfqpoint{3.644351in}{0.491420in}}%
\pgfpathlineto{\pgfqpoint{3.646620in}{0.452190in}}%
\pgfpathlineto{\pgfqpoint{3.648890in}{0.519638in}}%
\pgfpathlineto{\pgfqpoint{3.651159in}{0.475870in}}%
\pgfpathlineto{\pgfqpoint{3.653429in}{0.501466in}}%
\pgfpathlineto{\pgfqpoint{3.655698in}{0.413326in}}%
\pgfpathlineto{\pgfqpoint{3.657968in}{0.437892in}}%
\pgfpathlineto{\pgfqpoint{3.660237in}{0.435748in}}%
\pgfpathlineto{\pgfqpoint{3.662507in}{0.473530in}}%
\pgfpathlineto{\pgfqpoint{3.664776in}{0.499492in}}%
\pgfpathlineto{\pgfqpoint{3.667046in}{0.496391in}}%
\pgfpathlineto{\pgfqpoint{3.669315in}{0.408752in}}%
\pgfpathlineto{\pgfqpoint{3.671585in}{0.513785in}}%
\pgfpathlineto{\pgfqpoint{3.673854in}{0.516937in}}%
\pgfpathlineto{\pgfqpoint{3.676124in}{0.515047in}}%
\pgfpathlineto{\pgfqpoint{3.678393in}{0.450445in}}%
\pgfpathlineto{\pgfqpoint{3.680663in}{0.411237in}}%
\pgfpathlineto{\pgfqpoint{3.682932in}{0.503919in}}%
\pgfpathlineto{\pgfqpoint{3.685202in}{0.451400in}}%
\pgfpathlineto{\pgfqpoint{3.687471in}{0.495280in}}%
\pgfpathlineto{\pgfqpoint{3.689741in}{0.463014in}}%
\pgfpathlineto{\pgfqpoint{3.692010in}{0.464139in}}%
\pgfpathlineto{\pgfqpoint{3.694280in}{0.400864in}}%
\pgfpathlineto{\pgfqpoint{3.696549in}{0.485472in}}%
\pgfpathlineto{\pgfqpoint{3.698818in}{0.511983in}}%
\pgfpathlineto{\pgfqpoint{3.701088in}{0.505531in}}%
\pgfpathlineto{\pgfqpoint{3.703357in}{0.432859in}}%
\pgfpathlineto{\pgfqpoint{3.705627in}{0.526067in}}%
\pgfpathlineto{\pgfqpoint{3.707896in}{0.507182in}}%
\pgfpathlineto{\pgfqpoint{3.710166in}{0.521496in}}%
\pgfpathlineto{\pgfqpoint{3.712435in}{0.392429in}}%
\pgfpathlineto{\pgfqpoint{3.714705in}{0.459990in}}%
\pgfpathlineto{\pgfqpoint{3.716974in}{0.455632in}}%
\pgfpathlineto{\pgfqpoint{3.719244in}{0.419366in}}%
\pgfpathlineto{\pgfqpoint{3.721513in}{0.478226in}}%
\pgfpathlineto{\pgfqpoint{3.723783in}{0.482398in}}%
\pgfpathlineto{\pgfqpoint{3.726052in}{0.479103in}}%
\pgfpathlineto{\pgfqpoint{3.728322in}{0.443452in}}%
\pgfpathlineto{\pgfqpoint{3.730591in}{0.503557in}}%
\pgfpathlineto{\pgfqpoint{3.732861in}{0.526360in}}%
\pgfpathlineto{\pgfqpoint{3.735130in}{0.428985in}}%
\pgfpathlineto{\pgfqpoint{3.737400in}{0.471011in}}%
\pgfpathlineto{\pgfqpoint{3.739669in}{0.451855in}}%
\pgfpathlineto{\pgfqpoint{3.741939in}{0.390194in}}%
\pgfpathlineto{\pgfqpoint{3.744208in}{0.471521in}}%
\pgfpathlineto{\pgfqpoint{3.748747in}{0.484651in}}%
\pgfpathlineto{\pgfqpoint{3.751017in}{0.470234in}}%
\pgfpathlineto{\pgfqpoint{3.753286in}{0.493681in}}%
\pgfpathlineto{\pgfqpoint{3.755556in}{0.526097in}}%
\pgfpathlineto{\pgfqpoint{3.757825in}{0.454775in}}%
\pgfpathlineto{\pgfqpoint{3.760095in}{0.481706in}}%
\pgfpathlineto{\pgfqpoint{3.762364in}{0.430473in}}%
\pgfpathlineto{\pgfqpoint{3.764633in}{0.334119in}}%
\pgfpathlineto{\pgfqpoint{3.766903in}{0.467133in}}%
\pgfpathlineto{\pgfqpoint{3.771442in}{0.489766in}}%
\pgfpathlineto{\pgfqpoint{3.773711in}{0.478164in}}%
\pgfpathlineto{\pgfqpoint{3.775981in}{0.485618in}}%
\pgfpathlineto{\pgfqpoint{3.778250in}{0.525851in}}%
\pgfpathlineto{\pgfqpoint{3.780520in}{0.465065in}}%
\pgfpathlineto{\pgfqpoint{3.782789in}{0.489787in}}%
\pgfpathlineto{\pgfqpoint{3.785059in}{0.396639in}}%
\pgfpathlineto{\pgfqpoint{3.787328in}{0.376744in}}%
\pgfpathlineto{\pgfqpoint{3.789598in}{0.459712in}}%
\pgfpathlineto{\pgfqpoint{3.794137in}{0.493755in}}%
\pgfpathlineto{\pgfqpoint{3.796406in}{0.486050in}}%
\pgfpathlineto{\pgfqpoint{3.798676in}{0.472924in}}%
\pgfpathlineto{\pgfqpoint{3.800945in}{0.523923in}}%
\pgfpathlineto{\pgfqpoint{3.803215in}{0.471672in}}%
\pgfpathlineto{\pgfqpoint{3.805484in}{0.496256in}}%
\pgfpathlineto{\pgfqpoint{3.807754in}{0.337098in}}%
\pgfpathlineto{\pgfqpoint{3.810023in}{0.416268in}}%
\pgfpathlineto{\pgfqpoint{3.814562in}{0.476253in}}%
\pgfpathlineto{\pgfqpoint{3.816832in}{0.496960in}}%
\pgfpathlineto{\pgfqpoint{3.819101in}{0.492079in}}%
\pgfpathlineto{\pgfqpoint{3.821371in}{0.453513in}}%
\pgfpathlineto{\pgfqpoint{3.823640in}{0.520252in}}%
\pgfpathlineto{\pgfqpoint{3.825910in}{0.475264in}}%
\pgfpathlineto{\pgfqpoint{3.828179in}{0.501278in}}%
\pgfpathlineto{\pgfqpoint{3.830448in}{0.411024in}}%
\pgfpathlineto{\pgfqpoint{3.832718in}{0.438172in}}%
\pgfpathlineto{\pgfqpoint{3.834987in}{0.433870in}}%
\pgfpathlineto{\pgfqpoint{3.837257in}{0.473097in}}%
\pgfpathlineto{\pgfqpoint{3.839526in}{0.499194in}}%
\pgfpathlineto{\pgfqpoint{3.841796in}{0.496894in}}%
\pgfpathlineto{\pgfqpoint{3.844065in}{0.412482in}}%
\pgfpathlineto{\pgfqpoint{3.846335in}{0.514535in}}%
\pgfpathlineto{\pgfqpoint{3.848604in}{0.517585in}}%
\pgfpathlineto{\pgfqpoint{3.850874in}{0.514773in}}%
\pgfpathlineto{\pgfqpoint{3.853143in}{0.455462in}}%
\pgfpathlineto{\pgfqpoint{3.855413in}{0.415356in}}%
\pgfpathlineto{\pgfqpoint{3.857682in}{0.504546in}}%
\pgfpathlineto{\pgfqpoint{3.859952in}{0.450224in}}%
\pgfpathlineto{\pgfqpoint{3.862221in}{0.494854in}}%
\pgfpathlineto{\pgfqpoint{3.864491in}{0.462476in}}%
\pgfpathlineto{\pgfqpoint{3.866760in}{0.463965in}}%
\pgfpathlineto{\pgfqpoint{3.869030in}{0.404813in}}%
\pgfpathlineto{\pgfqpoint{3.871299in}{0.485565in}}%
\pgfpathlineto{\pgfqpoint{3.873569in}{0.511941in}}%
\pgfpathlineto{\pgfqpoint{3.875838in}{0.505964in}}%
\pgfpathlineto{\pgfqpoint{3.878108in}{0.438717in}}%
\pgfpathlineto{\pgfqpoint{3.880377in}{0.527316in}}%
\pgfpathlineto{\pgfqpoint{3.882647in}{0.507602in}}%
\pgfpathlineto{\pgfqpoint{3.884916in}{0.521377in}}%
\pgfpathlineto{\pgfqpoint{3.887186in}{0.405772in}}%
\pgfpathlineto{\pgfqpoint{3.889455in}{0.459367in}}%
\pgfpathlineto{\pgfqpoint{3.891725in}{0.456182in}}%
\pgfpathlineto{\pgfqpoint{3.893994in}{0.417892in}}%
\pgfpathlineto{\pgfqpoint{3.896264in}{0.477405in}}%
\pgfpathlineto{\pgfqpoint{3.898533in}{0.482139in}}%
\pgfpathlineto{\pgfqpoint{3.900802in}{0.478815in}}%
\pgfpathlineto{\pgfqpoint{3.903072in}{0.444607in}}%
\pgfpathlineto{\pgfqpoint{3.905341in}{0.504295in}}%
\pgfpathlineto{\pgfqpoint{3.907611in}{0.526528in}}%
\pgfpathlineto{\pgfqpoint{3.909880in}{0.423784in}}%
\pgfpathlineto{\pgfqpoint{3.912150in}{0.470238in}}%
\pgfpathlineto{\pgfqpoint{3.914419in}{0.452837in}}%
\pgfpathlineto{\pgfqpoint{3.916689in}{0.388505in}}%
\pgfpathlineto{\pgfqpoint{3.918958in}{0.470788in}}%
\pgfpathlineto{\pgfqpoint{3.923497in}{0.484268in}}%
\pgfpathlineto{\pgfqpoint{3.925767in}{0.470686in}}%
\pgfpathlineto{\pgfqpoint{3.928036in}{0.494565in}}%
\pgfpathlineto{\pgfqpoint{3.930306in}{0.526340in}}%
\pgfpathlineto{\pgfqpoint{3.932575in}{0.452571in}}%
\pgfpathlineto{\pgfqpoint{3.934845in}{0.480993in}}%
\pgfpathlineto{\pgfqpoint{3.937114in}{0.432439in}}%
\pgfpathlineto{\pgfqpoint{3.939384in}{0.333238in}}%
\pgfpathlineto{\pgfqpoint{3.941653in}{0.466570in}}%
\pgfpathlineto{\pgfqpoint{3.946192in}{0.489350in}}%
\pgfpathlineto{\pgfqpoint{3.948462in}{0.478346in}}%
\pgfpathlineto{\pgfqpoint{3.950731in}{0.486777in}}%
\pgfpathlineto{\pgfqpoint{3.953001in}{0.526197in}}%
\pgfpathlineto{\pgfqpoint{3.955270in}{0.463620in}}%
\pgfpathlineto{\pgfqpoint{3.957540in}{0.489139in}}%
\pgfpathlineto{\pgfqpoint{3.959809in}{0.401534in}}%
\pgfpathlineto{\pgfqpoint{3.962079in}{0.375117in}}%
\pgfpathlineto{\pgfqpoint{3.964348in}{0.459337in}}%
\pgfpathlineto{\pgfqpoint{3.968887in}{0.493339in}}%
\pgfpathlineto{\pgfqpoint{3.971156in}{0.486106in}}%
\pgfpathlineto{\pgfqpoint{3.973426in}{0.474577in}}%
\pgfpathlineto{\pgfqpoint{3.975695in}{0.524431in}}%
\pgfpathlineto{\pgfqpoint{3.977965in}{0.470722in}}%
\pgfpathlineto{\pgfqpoint{3.980234in}{0.495692in}}%
\pgfpathlineto{\pgfqpoint{3.982362in}{0.315625in}}%
\pgfpathmoveto{\pgfqpoint{3.982749in}{0.315625in}}%
\pgfpathlineto{\pgfqpoint{3.984773in}{0.414883in}}%
\pgfpathlineto{\pgfqpoint{3.989312in}{0.476156in}}%
\pgfpathlineto{\pgfqpoint{3.991582in}{0.496575in}}%
\pgfpathlineto{\pgfqpoint{3.993851in}{0.492075in}}%
\pgfpathlineto{\pgfqpoint{3.996121in}{0.456262in}}%
\pgfpathlineto{\pgfqpoint{3.998390in}{0.521002in}}%
\pgfpathlineto{\pgfqpoint{4.000660in}{0.474662in}}%
\pgfpathlineto{\pgfqpoint{4.002929in}{0.500809in}}%
\pgfpathlineto{\pgfqpoint{4.005199in}{0.406671in}}%
\pgfpathlineto{\pgfqpoint{4.007468in}{0.436932in}}%
\pgfpathlineto{\pgfqpoint{4.009738in}{0.434200in}}%
\pgfpathlineto{\pgfqpoint{4.012007in}{0.473122in}}%
\pgfpathlineto{\pgfqpoint{4.014277in}{0.498865in}}%
\pgfpathlineto{\pgfqpoint{4.016546in}{0.496875in}}%
\pgfpathlineto{\pgfqpoint{4.018816in}{0.419739in}}%
\pgfpathlineto{\pgfqpoint{4.021085in}{0.515869in}}%
\pgfpathlineto{\pgfqpoint{4.023355in}{0.476233in}}%
\pgfpathlineto{\pgfqpoint{4.025624in}{0.504599in}}%
\pgfpathlineto{\pgfqpoint{4.027894in}{0.432379in}}%
\pgfpathlineto{\pgfqpoint{4.030163in}{0.452369in}}%
\pgfpathlineto{\pgfqpoint{4.032432in}{0.407689in}}%
\pgfpathlineto{\pgfqpoint{4.034702in}{0.468607in}}%
\pgfpathlineto{\pgfqpoint{4.036971in}{0.500129in}}%
\pgfpathlineto{\pgfqpoint{4.039241in}{0.500614in}}%
\pgfpathlineto{\pgfqpoint{4.041510in}{0.384944in}}%
\pgfpathlineto{\pgfqpoint{4.043780in}{0.507979in}}%
\pgfpathlineto{\pgfqpoint{4.046049in}{0.516914in}}%
\pgfpathlineto{\pgfqpoint{4.048319in}{0.520909in}}%
\pgfpathlineto{\pgfqpoint{4.050588in}{0.409977in}}%
\pgfpathlineto{\pgfqpoint{4.052858in}{0.359055in}}%
\pgfpathlineto{\pgfqpoint{4.055127in}{0.498230in}}%
\pgfpathlineto{\pgfqpoint{4.057397in}{0.446405in}}%
\pgfpathlineto{\pgfqpoint{4.059666in}{0.494037in}}%
\pgfpathlineto{\pgfqpoint{4.061936in}{0.465756in}}%
\pgfpathlineto{\pgfqpoint{4.064205in}{0.471851in}}%
\pgfpathlineto{\pgfqpoint{4.066475in}{0.435493in}}%
\pgfpathlineto{\pgfqpoint{4.068744in}{0.479811in}}%
\pgfpathlineto{\pgfqpoint{4.071014in}{0.512744in}}%
\pgfpathlineto{\pgfqpoint{4.073283in}{0.510206in}}%
\pgfpathlineto{\pgfqpoint{4.075553in}{0.351242in}}%
\pgfpathlineto{\pgfqpoint{4.077822in}{0.521097in}}%
\pgfpathlineto{\pgfqpoint{4.080092in}{0.506737in}}%
\pgfpathlineto{\pgfqpoint{4.082361in}{0.525245in}}%
\pgfpathlineto{\pgfqpoint{4.084631in}{0.365748in}}%
\pgfpathlineto{\pgfqpoint{4.086900in}{0.471695in}}%
\pgfpathlineto{\pgfqpoint{4.089170in}{0.441052in}}%
\pgfpathlineto{\pgfqpoint{4.091439in}{0.395196in}}%
\pgfpathlineto{\pgfqpoint{4.093709in}{0.472368in}}%
\pgfpathlineto{\pgfqpoint{4.095978in}{0.482398in}}%
\pgfpathlineto{\pgfqpoint{4.098247in}{0.485182in}}%
\pgfpathlineto{\pgfqpoint{4.100517in}{0.463890in}}%
\pgfpathlineto{\pgfqpoint{4.105056in}{0.527819in}}%
\pgfpathlineto{\pgfqpoint{4.107325in}{0.446813in}}%
\pgfpathlineto{\pgfqpoint{4.109595in}{0.480809in}}%
\pgfpathlineto{\pgfqpoint{4.111864in}{0.435717in}}%
\pgfpathlineto{\pgfqpoint{4.113529in}{0.315625in}}%
\pgfpathmoveto{\pgfqpoint{4.114650in}{0.315625in}}%
\pgfpathlineto{\pgfqpoint{4.116403in}{0.463889in}}%
\pgfpathlineto{\pgfqpoint{4.120942in}{0.488959in}}%
\pgfpathlineto{\pgfqpoint{4.123212in}{0.481085in}}%
\pgfpathlineto{\pgfqpoint{4.125481in}{0.486299in}}%
\pgfpathlineto{\pgfqpoint{4.127751in}{0.526338in}}%
\pgfpathlineto{\pgfqpoint{4.130020in}{0.462812in}}%
\pgfpathlineto{\pgfqpoint{4.132290in}{0.489115in}}%
\pgfpathlineto{\pgfqpoint{4.134559in}{0.401088in}}%
\pgfpathlineto{\pgfqpoint{4.136829in}{0.380196in}}%
\pgfpathlineto{\pgfqpoint{4.139098in}{0.457661in}}%
\pgfpathlineto{\pgfqpoint{4.143637in}{0.493249in}}%
\pgfpathlineto{\pgfqpoint{4.145907in}{0.486959in}}%
\pgfpathlineto{\pgfqpoint{4.148176in}{0.475523in}}%
\pgfpathlineto{\pgfqpoint{4.150446in}{0.524991in}}%
\pgfpathlineto{\pgfqpoint{4.152715in}{0.469750in}}%
\pgfpathlineto{\pgfqpoint{4.154985in}{0.495562in}}%
\pgfpathlineto{\pgfqpoint{4.156788in}{0.315625in}}%
\pgfpathmoveto{\pgfqpoint{4.157971in}{0.315625in}}%
\pgfpathlineto{\pgfqpoint{4.159524in}{0.416478in}}%
\pgfpathlineto{\pgfqpoint{4.164062in}{0.475582in}}%
\pgfpathlineto{\pgfqpoint{4.166332in}{0.496289in}}%
\pgfpathlineto{\pgfqpoint{4.168601in}{0.492865in}}%
\pgfpathlineto{\pgfqpoint{4.170871in}{0.457272in}}%
\pgfpathlineto{\pgfqpoint{4.173140in}{0.521581in}}%
\pgfpathlineto{\pgfqpoint{4.175410in}{0.473990in}}%
\pgfpathlineto{\pgfqpoint{4.177679in}{0.500641in}}%
\pgfpathlineto{\pgfqpoint{4.179949in}{0.404243in}}%
\pgfpathlineto{\pgfqpoint{4.182218in}{0.437385in}}%
\pgfpathlineto{\pgfqpoint{4.184488in}{0.431996in}}%
\pgfpathlineto{\pgfqpoint{4.186757in}{0.472641in}}%
\pgfpathlineto{\pgfqpoint{4.189027in}{0.498568in}}%
\pgfpathlineto{\pgfqpoint{4.191296in}{0.497440in}}%
\pgfpathlineto{\pgfqpoint{4.193566in}{0.422491in}}%
\pgfpathlineto{\pgfqpoint{4.195835in}{0.516560in}}%
\pgfpathlineto{\pgfqpoint{4.198105in}{0.475742in}}%
\pgfpathlineto{\pgfqpoint{4.200374in}{0.504429in}}%
\pgfpathlineto{\pgfqpoint{4.202644in}{0.430774in}}%
\pgfpathlineto{\pgfqpoint{4.204913in}{0.452367in}}%
\pgfpathlineto{\pgfqpoint{4.207183in}{0.404575in}}%
\pgfpathlineto{\pgfqpoint{4.209452in}{0.468147in}}%
\pgfpathlineto{\pgfqpoint{4.211722in}{0.499808in}}%
\pgfpathlineto{\pgfqpoint{4.213991in}{0.501063in}}%
\pgfpathlineto{\pgfqpoint{4.216261in}{0.377582in}}%
\pgfpathlineto{\pgfqpoint{4.218530in}{0.508871in}}%
\pgfpathlineto{\pgfqpoint{4.220800in}{0.517637in}}%
\pgfpathlineto{\pgfqpoint{4.223069in}{0.520792in}}%
\pgfpathlineto{\pgfqpoint{4.227608in}{0.338697in}}%
\pgfpathlineto{\pgfqpoint{4.229878in}{0.498972in}}%
\pgfpathlineto{\pgfqpoint{4.232147in}{0.445277in}}%
\pgfpathlineto{\pgfqpoint{4.234416in}{0.493590in}}%
\pgfpathlineto{\pgfqpoint{4.236686in}{0.465189in}}%
\pgfpathlineto{\pgfqpoint{4.238955in}{0.471577in}}%
\pgfpathlineto{\pgfqpoint{4.241225in}{0.437048in}}%
\pgfpathlineto{\pgfqpoint{4.245764in}{0.512733in}}%
\pgfpathlineto{\pgfqpoint{4.248033in}{0.510660in}}%
\pgfpathlineto{\pgfqpoint{4.249264in}{0.315625in}}%
\pgfpathmoveto{\pgfqpoint{4.251309in}{0.315625in}}%
\pgfpathlineto{\pgfqpoint{4.252572in}{0.522520in}}%
\pgfpathlineto{\pgfqpoint{4.254842in}{0.507305in}}%
\pgfpathlineto{\pgfqpoint{4.257111in}{0.525262in}}%
\pgfpathlineto{\pgfqpoint{4.259381in}{0.321675in}}%
\pgfpathlineto{\pgfqpoint{4.261650in}{0.471137in}}%
\pgfpathlineto{\pgfqpoint{4.263920in}{0.442006in}}%
\pgfpathlineto{\pgfqpoint{4.266189in}{0.393194in}}%
\pgfpathlineto{\pgfqpoint{4.268459in}{0.471499in}}%
\pgfpathlineto{\pgfqpoint{4.270728in}{0.482119in}}%
\pgfpathlineto{\pgfqpoint{4.272998in}{0.484862in}}%
\pgfpathlineto{\pgfqpoint{4.275267in}{0.464593in}}%
\pgfpathlineto{\pgfqpoint{4.279806in}{0.528070in}}%
\pgfpathlineto{\pgfqpoint{4.282076in}{0.443820in}}%
\pgfpathlineto{\pgfqpoint{4.284345in}{0.480159in}}%
\pgfpathlineto{\pgfqpoint{4.286615in}{0.437301in}}%
\pgfpathlineto{\pgfqpoint{4.288123in}{0.315625in}}%
\pgfpathmoveto{\pgfqpoint{4.289551in}{0.315625in}}%
\pgfpathlineto{\pgfqpoint{4.291154in}{0.463129in}}%
\pgfpathlineto{\pgfqpoint{4.295693in}{0.488571in}}%
\pgfpathlineto{\pgfqpoint{4.297962in}{0.481434in}}%
\pgfpathlineto{\pgfqpoint{4.300231in}{0.487386in}}%
\pgfpathlineto{\pgfqpoint{4.302501in}{0.526686in}}%
\pgfpathlineto{\pgfqpoint{4.304770in}{0.461236in}}%
\pgfpathlineto{\pgfqpoint{4.307040in}{0.488513in}}%
\pgfpathlineto{\pgfqpoint{4.309309in}{0.405174in}}%
\pgfpathlineto{\pgfqpoint{4.311579in}{0.379453in}}%
\pgfpathlineto{\pgfqpoint{4.313848in}{0.457107in}}%
\pgfpathlineto{\pgfqpoint{4.318387in}{0.492849in}}%
\pgfpathlineto{\pgfqpoint{4.320657in}{0.487122in}}%
\pgfpathlineto{\pgfqpoint{4.322926in}{0.477044in}}%
\pgfpathlineto{\pgfqpoint{4.325196in}{0.525478in}}%
\pgfpathlineto{\pgfqpoint{4.327465in}{0.468677in}}%
\pgfpathlineto{\pgfqpoint{4.329735in}{0.495019in}}%
\pgfpathlineto{\pgfqpoint{4.331712in}{0.315625in}}%
\pgfpathmoveto{\pgfqpoint{4.332481in}{0.315625in}}%
\pgfpathlineto{\pgfqpoint{4.334274in}{0.415408in}}%
\pgfpathlineto{\pgfqpoint{4.336543in}{0.468858in}}%
\pgfpathlineto{\pgfqpoint{4.336543in}{0.468858in}}%
\pgfusepath{stroke}%
\end{pgfscope}%
\begin{pgfscope}%
\pgfpathrectangle{\pgfqpoint{2.520957in}{0.320625in}}{\pgfqpoint{1.815586in}{1.164375in}}%
\pgfusepath{clip}%
\pgfsetrectcap%
\pgfsetroundjoin%
\pgfsetlinewidth{1.154312pt}%
\definecolor{currentstroke}{rgb}{0.569172,0.167454,0.504105}%
\pgfsetstrokecolor{currentstroke}%
\pgfsetdash{}{0pt}%
\pgfpathmoveto{\pgfqpoint{2.520957in}{1.451682in}}%
\pgfpathlineto{\pgfqpoint{2.525496in}{1.442123in}}%
\pgfpathlineto{\pgfqpoint{2.543652in}{1.407752in}}%
\pgfpathlineto{\pgfqpoint{2.564077in}{1.366117in}}%
\pgfpathlineto{\pgfqpoint{2.602658in}{1.287429in}}%
\pgfpathlineto{\pgfqpoint{2.627622in}{1.239304in}}%
\pgfpathlineto{\pgfqpoint{2.652587in}{1.193577in}}%
\pgfpathlineto{\pgfqpoint{2.679821in}{1.145961in}}%
\pgfpathlineto{\pgfqpoint{2.711593in}{1.092654in}}%
\pgfpathlineto{\pgfqpoint{2.750175in}{1.030161in}}%
\pgfpathlineto{\pgfqpoint{2.800103in}{0.951544in}}%
\pgfpathlineto{\pgfqpoint{2.868188in}{0.846554in}}%
\pgfpathlineto{\pgfqpoint{2.972584in}{0.687795in}}%
\pgfpathlineto{\pgfqpoint{3.045207in}{0.578201in}}%
\pgfpathlineto{\pgfqpoint{3.049746in}{0.572637in}}%
\pgfpathlineto{\pgfqpoint{3.052016in}{0.567906in}}%
\pgfpathlineto{\pgfqpoint{3.054285in}{0.566202in}}%
\pgfpathlineto{\pgfqpoint{3.056555in}{0.562502in}}%
\pgfpathlineto{\pgfqpoint{3.058824in}{0.562362in}}%
\pgfpathlineto{\pgfqpoint{3.061094in}{0.554169in}}%
\pgfpathlineto{\pgfqpoint{3.063363in}{0.551486in}}%
\pgfpathlineto{\pgfqpoint{3.067902in}{0.543875in}}%
\pgfpathlineto{\pgfqpoint{3.079250in}{0.530128in}}%
\pgfpathlineto{\pgfqpoint{3.081519in}{0.540288in}}%
\pgfpathlineto{\pgfqpoint{3.083789in}{0.521117in}}%
\pgfpathlineto{\pgfqpoint{3.086058in}{0.521581in}}%
\pgfpathlineto{\pgfqpoint{3.088328in}{0.513245in}}%
\pgfpathlineto{\pgfqpoint{3.090597in}{0.509759in}}%
\pgfpathlineto{\pgfqpoint{3.092867in}{0.508017in}}%
\pgfpathlineto{\pgfqpoint{3.095136in}{0.508107in}}%
\pgfpathlineto{\pgfqpoint{3.097405in}{0.510922in}}%
\pgfpathlineto{\pgfqpoint{3.099675in}{0.507285in}}%
\pgfpathlineto{\pgfqpoint{3.101944in}{0.500479in}}%
\pgfpathlineto{\pgfqpoint{3.104214in}{0.528897in}}%
\pgfpathlineto{\pgfqpoint{3.106483in}{0.493481in}}%
\pgfpathlineto{\pgfqpoint{3.108753in}{0.504321in}}%
\pgfpathlineto{\pgfqpoint{3.111022in}{0.478954in}}%
\pgfpathlineto{\pgfqpoint{3.113292in}{0.476799in}}%
\pgfpathlineto{\pgfqpoint{3.115561in}{0.477421in}}%
\pgfpathlineto{\pgfqpoint{3.117831in}{0.486854in}}%
\pgfpathlineto{\pgfqpoint{3.120100in}{0.500309in}}%
\pgfpathlineto{\pgfqpoint{3.122370in}{0.497572in}}%
\pgfpathlineto{\pgfqpoint{3.124639in}{0.474077in}}%
\pgfpathlineto{\pgfqpoint{3.126909in}{0.523484in}}%
\pgfpathlineto{\pgfqpoint{3.129178in}{0.479435in}}%
\pgfpathlineto{\pgfqpoint{3.131448in}{0.501775in}}%
\pgfpathlineto{\pgfqpoint{3.133717in}{0.446888in}}%
\pgfpathlineto{\pgfqpoint{3.135987in}{0.453562in}}%
\pgfpathlineto{\pgfqpoint{3.138256in}{0.449580in}}%
\pgfpathlineto{\pgfqpoint{3.142795in}{0.498869in}}%
\pgfpathlineto{\pgfqpoint{3.145065in}{0.498333in}}%
\pgfpathlineto{\pgfqpoint{3.147334in}{0.442616in}}%
\pgfpathlineto{\pgfqpoint{3.149604in}{0.518178in}}%
\pgfpathlineto{\pgfqpoint{3.151873in}{0.476402in}}%
\pgfpathlineto{\pgfqpoint{3.154143in}{0.504089in}}%
\pgfpathlineto{\pgfqpoint{3.156412in}{0.434454in}}%
\pgfpathlineto{\pgfqpoint{3.158682in}{0.452985in}}%
\pgfpathlineto{\pgfqpoint{3.160951in}{0.420222in}}%
\pgfpathlineto{\pgfqpoint{3.163220in}{0.469328in}}%
\pgfpathlineto{\pgfqpoint{3.165490in}{0.499526in}}%
\pgfpathlineto{\pgfqpoint{3.167759in}{0.501097in}}%
\pgfpathlineto{\pgfqpoint{3.170029in}{0.390437in}}%
\pgfpathlineto{\pgfqpoint{3.172298in}{0.512279in}}%
\pgfpathlineto{\pgfqpoint{3.174568in}{0.475487in}}%
\pgfpathlineto{\pgfqpoint{3.176837in}{0.506452in}}%
\pgfpathlineto{\pgfqpoint{3.179107in}{0.442789in}}%
\pgfpathlineto{\pgfqpoint{3.181376in}{0.461978in}}%
\pgfpathlineto{\pgfqpoint{3.183646in}{0.373118in}}%
\pgfpathlineto{\pgfqpoint{3.185915in}{0.463022in}}%
\pgfpathlineto{\pgfqpoint{3.188185in}{0.499703in}}%
\pgfpathlineto{\pgfqpoint{3.190454in}{0.503627in}}%
\pgfpathlineto{\pgfqpoint{3.192724in}{0.435788in}}%
\pgfpathlineto{\pgfqpoint{3.194993in}{0.502330in}}%
\pgfpathlineto{\pgfqpoint{3.197263in}{0.516691in}}%
\pgfpathlineto{\pgfqpoint{3.199532in}{0.524817in}}%
\pgfpathlineto{\pgfqpoint{3.201802in}{0.347075in}}%
\pgfpathlineto{\pgfqpoint{3.206341in}{0.493333in}}%
\pgfpathlineto{\pgfqpoint{3.208610in}{0.439903in}}%
\pgfpathlineto{\pgfqpoint{3.210880in}{0.491633in}}%
\pgfpathlineto{\pgfqpoint{3.213149in}{0.466823in}}%
\pgfpathlineto{\pgfqpoint{3.215419in}{0.476602in}}%
\pgfpathlineto{\pgfqpoint{3.217688in}{0.450933in}}%
\pgfpathlineto{\pgfqpoint{3.219958in}{0.474259in}}%
\pgfpathlineto{\pgfqpoint{3.222227in}{0.512693in}}%
\pgfpathlineto{\pgfqpoint{3.224497in}{0.513284in}}%
\pgfpathlineto{\pgfqpoint{3.226766in}{0.431925in}}%
\pgfpathlineto{\pgfqpoint{3.229035in}{0.517313in}}%
\pgfpathlineto{\pgfqpoint{3.231305in}{0.506053in}}%
\pgfpathlineto{\pgfqpoint{3.233574in}{0.527430in}}%
\pgfpathlineto{\pgfqpoint{3.235844in}{0.406063in}}%
\pgfpathlineto{\pgfqpoint{3.238113in}{0.478289in}}%
\pgfpathlineto{\pgfqpoint{3.240383in}{0.427361in}}%
\pgfpathlineto{\pgfqpoint{3.242652in}{0.353843in}}%
\pgfpathlineto{\pgfqpoint{3.244922in}{0.466615in}}%
\pgfpathlineto{\pgfqpoint{3.247191in}{0.481780in}}%
\pgfpathlineto{\pgfqpoint{3.249461in}{0.488557in}}%
\pgfpathlineto{\pgfqpoint{3.251730in}{0.474019in}}%
\pgfpathlineto{\pgfqpoint{3.254000in}{0.494699in}}%
\pgfpathlineto{\pgfqpoint{3.256269in}{0.528493in}}%
\pgfpathlineto{\pgfqpoint{3.258539in}{0.453497in}}%
\pgfpathlineto{\pgfqpoint{3.260808in}{0.485959in}}%
\pgfpathlineto{\pgfqpoint{3.265347in}{0.370223in}}%
\pgfpathlineto{\pgfqpoint{3.267617in}{0.457108in}}%
\pgfpathlineto{\pgfqpoint{3.272156in}{0.491043in}}%
\pgfpathlineto{\pgfqpoint{3.274425in}{0.486806in}}%
\pgfpathlineto{\pgfqpoint{3.276695in}{0.480817in}}%
\pgfpathlineto{\pgfqpoint{3.278964in}{0.526202in}}%
\pgfpathlineto{\pgfqpoint{3.281234in}{0.465748in}}%
\pgfpathlineto{\pgfqpoint{3.283503in}{0.492717in}}%
\pgfpathlineto{\pgfqpoint{3.285773in}{0.368039in}}%
\pgfpathlineto{\pgfqpoint{3.292581in}{0.476105in}}%
\pgfpathlineto{\pgfqpoint{3.294851in}{0.494703in}}%
\pgfpathlineto{\pgfqpoint{3.297120in}{0.491310in}}%
\pgfpathlineto{\pgfqpoint{3.299389in}{0.468445in}}%
\pgfpathlineto{\pgfqpoint{3.301659in}{0.524230in}}%
\pgfpathlineto{\pgfqpoint{3.303928in}{0.471134in}}%
\pgfpathlineto{\pgfqpoint{3.306198in}{0.498089in}}%
\pgfpathlineto{\pgfqpoint{3.308467in}{0.371141in}}%
\pgfpathlineto{\pgfqpoint{3.310737in}{0.428942in}}%
\pgfpathlineto{\pgfqpoint{3.313006in}{0.437900in}}%
\pgfpathlineto{\pgfqpoint{3.315276in}{0.473617in}}%
\pgfpathlineto{\pgfqpoint{3.317545in}{0.497115in}}%
\pgfpathlineto{\pgfqpoint{3.319815in}{0.495999in}}%
\pgfpathlineto{\pgfqpoint{3.322084in}{0.446333in}}%
\pgfpathlineto{\pgfqpoint{3.324354in}{0.520160in}}%
\pgfpathlineto{\pgfqpoint{3.326623in}{0.474316in}}%
\pgfpathlineto{\pgfqpoint{3.328893in}{0.502310in}}%
\pgfpathlineto{\pgfqpoint{3.331162in}{0.416159in}}%
\pgfpathlineto{\pgfqpoint{3.333432in}{0.444841in}}%
\pgfpathlineto{\pgfqpoint{3.335701in}{0.419112in}}%
\pgfpathlineto{\pgfqpoint{3.337971in}{0.470309in}}%
\pgfpathlineto{\pgfqpoint{3.340240in}{0.498863in}}%
\pgfpathlineto{\pgfqpoint{3.342510in}{0.499691in}}%
\pgfpathlineto{\pgfqpoint{3.344779in}{0.394384in}}%
\pgfpathlineto{\pgfqpoint{3.347049in}{0.514927in}}%
\pgfpathlineto{\pgfqpoint{3.349318in}{0.475330in}}%
\pgfpathlineto{\pgfqpoint{3.351588in}{0.505410in}}%
\pgfpathlineto{\pgfqpoint{3.353857in}{0.435891in}}%
\pgfpathlineto{\pgfqpoint{3.356127in}{0.457166in}}%
\pgfpathlineto{\pgfqpoint{3.358396in}{0.379565in}}%
\pgfpathlineto{\pgfqpoint{3.360666in}{0.465412in}}%
\pgfpathlineto{\pgfqpoint{3.362935in}{0.499629in}}%
\pgfpathlineto{\pgfqpoint{3.365204in}{0.502636in}}%
\pgfpathlineto{\pgfqpoint{3.367474in}{0.412263in}}%
\pgfpathlineto{\pgfqpoint{3.369743in}{0.508451in}}%
\pgfpathlineto{\pgfqpoint{3.372013in}{0.474367in}}%
\pgfpathlineto{\pgfqpoint{3.374282in}{0.507335in}}%
\pgfpathlineto{\pgfqpoint{3.376552in}{0.447949in}}%
\pgfpathlineto{\pgfqpoint{3.378821in}{0.467088in}}%
\pgfpathlineto{\pgfqpoint{3.381091in}{0.360270in}}%
\pgfpathlineto{\pgfqpoint{3.383360in}{0.458422in}}%
\pgfpathlineto{\pgfqpoint{3.385630in}{0.499255in}}%
\pgfpathlineto{\pgfqpoint{3.387899in}{0.504835in}}%
\pgfpathlineto{\pgfqpoint{3.390169in}{0.453902in}}%
\pgfpathlineto{\pgfqpoint{3.392438in}{0.496051in}}%
\pgfpathlineto{\pgfqpoint{3.394708in}{0.515195in}}%
\pgfpathlineto{\pgfqpoint{3.396977in}{0.527141in}}%
\pgfpathlineto{\pgfqpoint{3.399247in}{0.416523in}}%
\pgfpathlineto{\pgfqpoint{3.401516in}{0.445939in}}%
\pgfpathlineto{\pgfqpoint{3.403786in}{0.488148in}}%
\pgfpathlineto{\pgfqpoint{3.406055in}{0.434601in}}%
\pgfpathlineto{\pgfqpoint{3.408325in}{0.489865in}}%
\pgfpathlineto{\pgfqpoint{3.410594in}{0.467855in}}%
\pgfpathlineto{\pgfqpoint{3.412864in}{0.480008in}}%
\pgfpathlineto{\pgfqpoint{3.415133in}{0.458438in}}%
\pgfpathlineto{\pgfqpoint{3.417403in}{0.468795in}}%
\pgfpathlineto{\pgfqpoint{3.419672in}{0.512191in}}%
\pgfpathlineto{\pgfqpoint{3.421942in}{0.514686in}}%
\pgfpathlineto{\pgfqpoint{3.424211in}{0.457658in}}%
\pgfpathlineto{\pgfqpoint{3.426481in}{0.511751in}}%
\pgfpathlineto{\pgfqpoint{3.428750in}{0.504171in}}%
\pgfpathlineto{\pgfqpoint{3.431019in}{0.528585in}}%
\pgfpathlineto{\pgfqpoint{3.433289in}{0.428713in}}%
\pgfpathlineto{\pgfqpoint{3.435558in}{0.483408in}}%
\pgfpathlineto{\pgfqpoint{3.437828in}{0.409534in}}%
\pgfpathlineto{\pgfqpoint{3.439962in}{0.315625in}}%
\pgfpathmoveto{\pgfqpoint{3.440186in}{0.315625in}}%
\pgfpathlineto{\pgfqpoint{3.442367in}{0.462223in}}%
\pgfpathlineto{\pgfqpoint{3.444636in}{0.481333in}}%
\pgfpathlineto{\pgfqpoint{3.446906in}{0.491367in}}%
\pgfpathlineto{\pgfqpoint{3.449175in}{0.480160in}}%
\pgfpathlineto{\pgfqpoint{3.451445in}{0.489416in}}%
\pgfpathlineto{\pgfqpoint{3.453714in}{0.528296in}}%
\pgfpathlineto{\pgfqpoint{3.455984in}{0.460201in}}%
\pgfpathlineto{\pgfqpoint{3.458253in}{0.490479in}}%
\pgfpathlineto{\pgfqpoint{3.460523in}{0.400288in}}%
\pgfpathlineto{\pgfqpoint{3.462792in}{0.399144in}}%
\pgfpathlineto{\pgfqpoint{3.465062in}{0.451009in}}%
\pgfpathlineto{\pgfqpoint{3.467331in}{0.475154in}}%
\pgfpathlineto{\pgfqpoint{3.469601in}{0.493068in}}%
\pgfpathlineto{\pgfqpoint{3.471870in}{0.490819in}}%
\pgfpathlineto{\pgfqpoint{3.474140in}{0.472755in}}%
\pgfpathlineto{\pgfqpoint{3.476409in}{0.525090in}}%
\pgfpathlineto{\pgfqpoint{3.478679in}{0.469253in}}%
\pgfpathlineto{\pgfqpoint{3.480948in}{0.496285in}}%
\pgfpathlineto{\pgfqpoint{3.483218in}{0.325077in}}%
\pgfpathlineto{\pgfqpoint{3.485487in}{0.422638in}}%
\pgfpathlineto{\pgfqpoint{3.487757in}{0.441646in}}%
\pgfpathlineto{\pgfqpoint{3.490026in}{0.474497in}}%
\pgfpathlineto{\pgfqpoint{3.492296in}{0.496309in}}%
\pgfpathlineto{\pgfqpoint{3.494565in}{0.494780in}}%
\pgfpathlineto{\pgfqpoint{3.496834in}{0.456343in}}%
\pgfpathlineto{\pgfqpoint{3.499104in}{0.522141in}}%
\pgfpathlineto{\pgfqpoint{3.501373in}{0.473087in}}%
\pgfpathlineto{\pgfqpoint{3.503643in}{0.500897in}}%
\pgfpathlineto{\pgfqpoint{3.505912in}{0.404093in}}%
\pgfpathlineto{\pgfqpoint{3.508182in}{0.440041in}}%
\pgfpathlineto{\pgfqpoint{3.510451in}{0.425466in}}%
\pgfpathlineto{\pgfqpoint{3.512721in}{0.471294in}}%
\pgfpathlineto{\pgfqpoint{3.514990in}{0.498173in}}%
\pgfpathlineto{\pgfqpoint{3.517260in}{0.498802in}}%
\pgfpathlineto{\pgfqpoint{3.519529in}{0.421682in}}%
\pgfpathlineto{\pgfqpoint{3.521799in}{0.517205in}}%
\pgfpathlineto{\pgfqpoint{3.524068in}{0.474941in}}%
\pgfpathlineto{\pgfqpoint{3.526338in}{0.504376in}}%
\pgfpathlineto{\pgfqpoint{3.528607in}{0.429427in}}%
\pgfpathlineto{\pgfqpoint{3.530877in}{0.453199in}}%
\pgfpathlineto{\pgfqpoint{3.533146in}{0.396129in}}%
\pgfpathlineto{\pgfqpoint{3.535416in}{0.467079in}}%
\pgfpathlineto{\pgfqpoint{3.537685in}{0.499321in}}%
\pgfpathlineto{\pgfqpoint{3.539955in}{0.501928in}}%
\pgfpathlineto{\pgfqpoint{3.542224in}{0.371147in}}%
\pgfpathlineto{\pgfqpoint{3.544494in}{0.511264in}}%
\pgfpathlineto{\pgfqpoint{3.546763in}{0.474752in}}%
\pgfpathlineto{\pgfqpoint{3.549033in}{0.506729in}}%
\pgfpathlineto{\pgfqpoint{3.551302in}{0.443671in}}%
\pgfpathlineto{\pgfqpoint{3.553572in}{0.463753in}}%
\pgfpathlineto{\pgfqpoint{3.555454in}{0.315625in}}%
\pgfpathmoveto{\pgfqpoint{3.556234in}{0.315625in}}%
\pgfpathlineto{\pgfqpoint{3.558111in}{0.460932in}}%
\pgfpathlineto{\pgfqpoint{3.560380in}{0.499371in}}%
\pgfpathlineto{\pgfqpoint{3.562649in}{0.504345in}}%
\pgfpathlineto{\pgfqpoint{3.564919in}{0.442583in}}%
\pgfpathlineto{\pgfqpoint{3.567188in}{0.500626in}}%
\pgfpathlineto{\pgfqpoint{3.569458in}{0.516609in}}%
\pgfpathlineto{\pgfqpoint{3.571727in}{0.525789in}}%
\pgfpathlineto{\pgfqpoint{3.573997in}{0.368116in}}%
\pgfpathlineto{\pgfqpoint{3.578536in}{0.491896in}}%
\pgfpathlineto{\pgfqpoint{3.580805in}{0.437343in}}%
\pgfpathlineto{\pgfqpoint{3.583075in}{0.490720in}}%
\pgfpathlineto{\pgfqpoint{3.585344in}{0.466919in}}%
\pgfpathlineto{\pgfqpoint{3.587614in}{0.477767in}}%
\pgfpathlineto{\pgfqpoint{3.589883in}{0.454488in}}%
\pgfpathlineto{\pgfqpoint{3.592153in}{0.472372in}}%
\pgfpathlineto{\pgfqpoint{3.594422in}{0.512570in}}%
\pgfpathlineto{\pgfqpoint{3.596692in}{0.514137in}}%
\pgfpathlineto{\pgfqpoint{3.598961in}{0.441287in}}%
\pgfpathlineto{\pgfqpoint{3.601231in}{0.516172in}}%
\pgfpathlineto{\pgfqpoint{3.603500in}{0.505800in}}%
\pgfpathlineto{\pgfqpoint{3.605770in}{0.528020in}}%
\pgfpathlineto{\pgfqpoint{3.608039in}{0.412190in}}%
\pgfpathlineto{\pgfqpoint{3.610309in}{0.480074in}}%
\pgfpathlineto{\pgfqpoint{3.612578in}{0.422250in}}%
\pgfpathlineto{\pgfqpoint{3.614848in}{0.318436in}}%
\pgfpathlineto{\pgfqpoint{3.617117in}{0.464480in}}%
\pgfpathlineto{\pgfqpoint{3.619387in}{0.481475in}}%
\pgfpathlineto{\pgfqpoint{3.621656in}{0.489499in}}%
\pgfpathlineto{\pgfqpoint{3.623926in}{0.476828in}}%
\pgfpathlineto{\pgfqpoint{3.626195in}{0.493326in}}%
\pgfpathlineto{\pgfqpoint{3.628465in}{0.528660in}}%
\pgfpathlineto{\pgfqpoint{3.630734in}{0.455180in}}%
\pgfpathlineto{\pgfqpoint{3.633003in}{0.487475in}}%
\pgfpathlineto{\pgfqpoint{3.635273in}{0.416176in}}%
\pgfpathlineto{\pgfqpoint{3.637542in}{0.383716in}}%
\pgfpathlineto{\pgfqpoint{3.639812in}{0.454435in}}%
\pgfpathlineto{\pgfqpoint{3.642081in}{0.475690in}}%
\pgfpathlineto{\pgfqpoint{3.644351in}{0.491650in}}%
\pgfpathlineto{\pgfqpoint{3.646620in}{0.488603in}}%
\pgfpathlineto{\pgfqpoint{3.648890in}{0.478680in}}%
\pgfpathlineto{\pgfqpoint{3.651159in}{0.526095in}}%
\pgfpathlineto{\pgfqpoint{3.653429in}{0.466559in}}%
\pgfpathlineto{\pgfqpoint{3.655698in}{0.493873in}}%
\pgfpathlineto{\pgfqpoint{3.657968in}{0.346168in}}%
\pgfpathlineto{\pgfqpoint{3.660237in}{0.412868in}}%
\pgfpathlineto{\pgfqpoint{3.664776in}{0.475415in}}%
\pgfpathlineto{\pgfqpoint{3.667046in}{0.495153in}}%
\pgfpathlineto{\pgfqpoint{3.669315in}{0.492809in}}%
\pgfpathlineto{\pgfqpoint{3.671585in}{0.465358in}}%
\pgfpathlineto{\pgfqpoint{3.673854in}{0.523865in}}%
\pgfpathlineto{\pgfqpoint{3.676124in}{0.471523in}}%
\pgfpathlineto{\pgfqpoint{3.678393in}{0.498975in}}%
\pgfpathlineto{\pgfqpoint{3.680663in}{0.383371in}}%
\pgfpathlineto{\pgfqpoint{3.682932in}{0.433025in}}%
\pgfpathlineto{\pgfqpoint{3.685202in}{0.433419in}}%
\pgfpathlineto{\pgfqpoint{3.687471in}{0.472720in}}%
\pgfpathlineto{\pgfqpoint{3.689741in}{0.497371in}}%
\pgfpathlineto{\pgfqpoint{3.692010in}{0.497191in}}%
\pgfpathlineto{\pgfqpoint{3.694280in}{0.440680in}}%
\pgfpathlineto{\pgfqpoint{3.696549in}{0.519524in}}%
\pgfpathlineto{\pgfqpoint{3.698818in}{0.474356in}}%
\pgfpathlineto{\pgfqpoint{3.701088in}{0.502950in}}%
\pgfpathlineto{\pgfqpoint{3.703357in}{0.420191in}}%
\pgfpathlineto{\pgfqpoint{3.705627in}{0.447713in}}%
\pgfpathlineto{\pgfqpoint{3.707896in}{0.411910in}}%
\pgfpathlineto{\pgfqpoint{3.710166in}{0.469158in}}%
\pgfpathlineto{\pgfqpoint{3.712435in}{0.498926in}}%
\pgfpathlineto{\pgfqpoint{3.714705in}{0.500640in}}%
\pgfpathlineto{\pgfqpoint{3.716974in}{0.369468in}}%
\pgfpathlineto{\pgfqpoint{3.719244in}{0.514112in}}%
\pgfpathlineto{\pgfqpoint{3.721513in}{0.475058in}}%
\pgfpathlineto{\pgfqpoint{3.723783in}{0.505817in}}%
\pgfpathlineto{\pgfqpoint{3.726052in}{0.438039in}}%
\pgfpathlineto{\pgfqpoint{3.728322in}{0.459306in}}%
\pgfpathlineto{\pgfqpoint{3.730591in}{0.358417in}}%
\pgfpathlineto{\pgfqpoint{3.732861in}{0.463919in}}%
\pgfpathlineto{\pgfqpoint{3.735130in}{0.499474in}}%
\pgfpathlineto{\pgfqpoint{3.737400in}{0.503376in}}%
\pgfpathlineto{\pgfqpoint{3.739669in}{0.423256in}}%
\pgfpathlineto{\pgfqpoint{3.741939in}{0.507425in}}%
\pgfpathlineto{\pgfqpoint{3.744208in}{0.473776in}}%
\pgfpathlineto{\pgfqpoint{3.746478in}{0.507505in}}%
\pgfpathlineto{\pgfqpoint{3.748747in}{0.449222in}}%
\pgfpathlineto{\pgfqpoint{3.751017in}{0.468730in}}%
\pgfpathlineto{\pgfqpoint{3.753286in}{0.380172in}}%
\pgfpathlineto{\pgfqpoint{3.755556in}{0.456427in}}%
\pgfpathlineto{\pgfqpoint{3.757825in}{0.498842in}}%
\pgfpathlineto{\pgfqpoint{3.760095in}{0.505392in}}%
\pgfpathlineto{\pgfqpoint{3.762364in}{0.458354in}}%
\pgfpathlineto{\pgfqpoint{3.764633in}{0.494151in}}%
\pgfpathlineto{\pgfqpoint{3.766903in}{0.514987in}}%
\pgfpathlineto{\pgfqpoint{3.769172in}{0.527892in}}%
\pgfpathlineto{\pgfqpoint{3.771442in}{0.424337in}}%
\pgfpathlineto{\pgfqpoint{3.773711in}{0.450773in}}%
\pgfpathlineto{\pgfqpoint{3.775981in}{0.486653in}}%
\pgfpathlineto{\pgfqpoint{3.778250in}{0.431846in}}%
\pgfpathlineto{\pgfqpoint{3.780520in}{0.488956in}}%
\pgfpathlineto{\pgfqpoint{3.782789in}{0.467865in}}%
\pgfpathlineto{\pgfqpoint{3.785059in}{0.480942in}}%
\pgfpathlineto{\pgfqpoint{3.787328in}{0.461078in}}%
\pgfpathlineto{\pgfqpoint{3.789598in}{0.466734in}}%
\pgfpathlineto{\pgfqpoint{3.791867in}{0.511942in}}%
\pgfpathlineto{\pgfqpoint{3.794137in}{0.515331in}}%
\pgfpathlineto{\pgfqpoint{3.796406in}{0.462688in}}%
\pgfpathlineto{\pgfqpoint{3.798676in}{0.510449in}}%
\pgfpathlineto{\pgfqpoint{3.800945in}{0.503764in}}%
\pgfpathlineto{\pgfqpoint{3.803215in}{0.528984in}}%
\pgfpathlineto{\pgfqpoint{3.805484in}{0.431849in}}%
\pgfpathlineto{\pgfqpoint{3.807754in}{0.484822in}}%
\pgfpathlineto{\pgfqpoint{3.810023in}{0.402531in}}%
\pgfpathlineto{\pgfqpoint{3.812293in}{0.347982in}}%
\pgfpathlineto{\pgfqpoint{3.814562in}{0.460034in}}%
\pgfpathlineto{\pgfqpoint{3.816832in}{0.480975in}}%
\pgfpathlineto{\pgfqpoint{3.819101in}{0.492109in}}%
\pgfpathlineto{\pgfqpoint{3.821371in}{0.482285in}}%
\pgfpathlineto{\pgfqpoint{3.823640in}{0.487857in}}%
\pgfpathlineto{\pgfqpoint{3.825910in}{0.528316in}}%
\pgfpathlineto{\pgfqpoint{3.828179in}{0.461331in}}%
\pgfpathlineto{\pgfqpoint{3.830448in}{0.491678in}}%
\pgfpathlineto{\pgfqpoint{3.832718in}{0.391874in}}%
\pgfpathlineto{\pgfqpoint{3.834987in}{0.406078in}}%
\pgfpathlineto{\pgfqpoint{3.837257in}{0.448174in}}%
\pgfpathlineto{\pgfqpoint{3.839526in}{0.474555in}}%
\pgfpathlineto{\pgfqpoint{3.841796in}{0.493519in}}%
\pgfpathlineto{\pgfqpoint{3.844065in}{0.492223in}}%
\pgfpathlineto{\pgfqpoint{3.846335in}{0.470212in}}%
\pgfpathlineto{\pgfqpoint{3.848604in}{0.524828in}}%
\pgfpathlineto{\pgfqpoint{3.850874in}{0.469769in}}%
\pgfpathlineto{\pgfqpoint{3.853143in}{0.497192in}}%
\pgfpathlineto{\pgfqpoint{3.855413in}{0.353734in}}%
\pgfpathlineto{\pgfqpoint{3.857682in}{0.426936in}}%
\pgfpathlineto{\pgfqpoint{3.859952in}{0.437963in}}%
\pgfpathlineto{\pgfqpoint{3.862221in}{0.473749in}}%
\pgfpathlineto{\pgfqpoint{3.864491in}{0.496618in}}%
\pgfpathlineto{\pgfqpoint{3.866760in}{0.495954in}}%
\pgfpathlineto{\pgfqpoint{3.869030in}{0.452310in}}%
\pgfpathlineto{\pgfqpoint{3.871299in}{0.521599in}}%
\pgfpathlineto{\pgfqpoint{3.873569in}{0.473264in}}%
\pgfpathlineto{\pgfqpoint{3.875838in}{0.501575in}}%
\pgfpathlineto{\pgfqpoint{3.878108in}{0.409320in}}%
\pgfpathlineto{\pgfqpoint{3.880377in}{0.442986in}}%
\pgfpathlineto{\pgfqpoint{3.882647in}{0.419928in}}%
\pgfpathlineto{\pgfqpoint{3.884916in}{0.470325in}}%
\pgfpathlineto{\pgfqpoint{3.887186in}{0.498301in}}%
\pgfpathlineto{\pgfqpoint{3.889455in}{0.499735in}}%
\pgfpathlineto{\pgfqpoint{3.891725in}{0.411407in}}%
\pgfpathlineto{\pgfqpoint{3.893994in}{0.516498in}}%
\pgfpathlineto{\pgfqpoint{3.896264in}{0.474812in}}%
\pgfpathlineto{\pgfqpoint{3.898533in}{0.504836in}}%
\pgfpathlineto{\pgfqpoint{3.900802in}{0.431947in}}%
\pgfpathlineto{\pgfqpoint{3.903072in}{0.455378in}}%
\pgfpathlineto{\pgfqpoint{3.905341in}{0.384283in}}%
\pgfpathlineto{\pgfqpoint{3.907611in}{0.465828in}}%
\pgfpathlineto{\pgfqpoint{3.909880in}{0.499259in}}%
\pgfpathlineto{\pgfqpoint{3.912150in}{0.502664in}}%
\pgfpathlineto{\pgfqpoint{3.914419in}{0.394413in}}%
\pgfpathlineto{\pgfqpoint{3.916689in}{0.510373in}}%
\pgfpathlineto{\pgfqpoint{3.918958in}{0.474330in}}%
\pgfpathlineto{\pgfqpoint{3.921228in}{0.506974in}}%
\pgfpathlineto{\pgfqpoint{3.923497in}{0.445141in}}%
\pgfpathlineto{\pgfqpoint{3.925767in}{0.465422in}}%
\pgfpathlineto{\pgfqpoint{3.928036in}{0.346924in}}%
\pgfpathlineto{\pgfqpoint{3.930306in}{0.459278in}}%
\pgfpathlineto{\pgfqpoint{3.932575in}{0.499087in}}%
\pgfpathlineto{\pgfqpoint{3.934845in}{0.504909in}}%
\pgfpathlineto{\pgfqpoint{3.937114in}{0.447982in}}%
\pgfpathlineto{\pgfqpoint{3.939384in}{0.499054in}}%
\pgfpathlineto{\pgfqpoint{3.941653in}{0.516469in}}%
\pgfpathlineto{\pgfqpoint{3.943923in}{0.526577in}}%
\pgfpathlineto{\pgfqpoint{3.946192in}{0.389317in}}%
\pgfpathlineto{\pgfqpoint{3.950731in}{0.490600in}}%
\pgfpathlineto{\pgfqpoint{3.953001in}{0.435176in}}%
\pgfpathlineto{\pgfqpoint{3.955270in}{0.489950in}}%
\pgfpathlineto{\pgfqpoint{3.957540in}{0.467014in}}%
\pgfpathlineto{\pgfqpoint{3.959809in}{0.478726in}}%
\pgfpathlineto{\pgfqpoint{3.962079in}{0.457207in}}%
\pgfpathlineto{\pgfqpoint{3.964348in}{0.470677in}}%
\pgfpathlineto{\pgfqpoint{3.966617in}{0.512426in}}%
\pgfpathlineto{\pgfqpoint{3.968887in}{0.514793in}}%
\pgfpathlineto{\pgfqpoint{3.971156in}{0.448150in}}%
\pgfpathlineto{\pgfqpoint{3.973426in}{0.515066in}}%
\pgfpathlineto{\pgfqpoint{3.975695in}{0.505509in}}%
\pgfpathlineto{\pgfqpoint{3.977965in}{0.528477in}}%
\pgfpathlineto{\pgfqpoint{3.980234in}{0.417103in}}%
\pgfpathlineto{\pgfqpoint{3.982504in}{0.481539in}}%
\pgfpathlineto{\pgfqpoint{3.984773in}{0.417395in}}%
\pgfpathlineto{\pgfqpoint{3.985886in}{0.315625in}}%
\pgfpathmoveto{\pgfqpoint{3.987993in}{0.315625in}}%
\pgfpathlineto{\pgfqpoint{3.989312in}{0.462637in}}%
\pgfpathlineto{\pgfqpoint{3.991582in}{0.481205in}}%
\pgfpathlineto{\pgfqpoint{3.993851in}{0.490271in}}%
\pgfpathlineto{\pgfqpoint{3.996121in}{0.479005in}}%
\pgfpathlineto{\pgfqpoint{3.998390in}{0.492039in}}%
\pgfpathlineto{\pgfqpoint{4.000660in}{0.528750in}}%
\pgfpathlineto{\pgfqpoint{4.002929in}{0.456607in}}%
\pgfpathlineto{\pgfqpoint{4.005199in}{0.488727in}}%
\pgfpathlineto{\pgfqpoint{4.007468in}{0.410746in}}%
\pgfpathlineto{\pgfqpoint{4.009738in}{0.392536in}}%
\pgfpathlineto{\pgfqpoint{4.012007in}{0.452087in}}%
\pgfpathlineto{\pgfqpoint{4.014277in}{0.475210in}}%
\pgfpathlineto{\pgfqpoint{4.016546in}{0.492147in}}%
\pgfpathlineto{\pgfqpoint{4.018816in}{0.490021in}}%
\pgfpathlineto{\pgfqpoint{4.021085in}{0.476667in}}%
\pgfpathlineto{\pgfqpoint{4.023355in}{0.525942in}}%
\pgfpathlineto{\pgfqpoint{4.025624in}{0.467252in}}%
\pgfpathlineto{\pgfqpoint{4.027894in}{0.494832in}}%
\pgfpathlineto{\pgfqpoint{4.030085in}{0.315625in}}%
\pgfpathmoveto{\pgfqpoint{4.030296in}{0.315625in}}%
\pgfpathlineto{\pgfqpoint{4.032432in}{0.417782in}}%
\pgfpathlineto{\pgfqpoint{4.039241in}{0.495519in}}%
\pgfpathlineto{\pgfqpoint{4.041510in}{0.494000in}}%
\pgfpathlineto{\pgfqpoint{4.043780in}{0.462391in}}%
\pgfpathlineto{\pgfqpoint{4.046049in}{0.523475in}}%
\pgfpathlineto{\pgfqpoint{4.048319in}{0.471852in}}%
\pgfpathlineto{\pgfqpoint{4.050588in}{0.499710in}}%
\pgfpathlineto{\pgfqpoint{4.052858in}{0.391651in}}%
\pgfpathlineto{\pgfqpoint{4.055127in}{0.436233in}}%
\pgfpathlineto{\pgfqpoint{4.057397in}{0.429262in}}%
\pgfpathlineto{\pgfqpoint{4.059666in}{0.471923in}}%
\pgfpathlineto{\pgfqpoint{4.061936in}{0.497574in}}%
\pgfpathlineto{\pgfqpoint{4.064205in}{0.498142in}}%
\pgfpathlineto{\pgfqpoint{4.066475in}{0.434906in}}%
\pgfpathlineto{\pgfqpoint{4.068744in}{0.518915in}}%
\pgfpathlineto{\pgfqpoint{4.071014in}{0.474376in}}%
\pgfpathlineto{\pgfqpoint{4.073283in}{0.503477in}}%
\pgfpathlineto{\pgfqpoint{4.075553in}{0.423392in}}%
\pgfpathlineto{\pgfqpoint{4.077822in}{0.450035in}}%
\pgfpathlineto{\pgfqpoint{4.080092in}{0.404720in}}%
\pgfpathlineto{\pgfqpoint{4.082361in}{0.468129in}}%
\pgfpathlineto{\pgfqpoint{4.084631in}{0.498960in}}%
\pgfpathlineto{\pgfqpoint{4.086900in}{0.501397in}}%
\pgfpathlineto{\pgfqpoint{4.089170in}{0.316231in}}%
\pgfpathlineto{\pgfqpoint{4.091439in}{0.513340in}}%
\pgfpathlineto{\pgfqpoint{4.093709in}{0.474800in}}%
\pgfpathlineto{\pgfqpoint{4.095978in}{0.506143in}}%
\pgfpathlineto{\pgfqpoint{4.098247in}{0.439817in}}%
\pgfpathlineto{\pgfqpoint{4.100517in}{0.461069in}}%
\pgfpathlineto{\pgfqpoint{4.102786in}{0.324698in}}%
\pgfpathlineto{\pgfqpoint{4.105056in}{0.462570in}}%
\pgfpathlineto{\pgfqpoint{4.107325in}{0.499316in}}%
\pgfpathlineto{\pgfqpoint{4.109595in}{0.503964in}}%
\pgfpathlineto{\pgfqpoint{4.111864in}{0.431073in}}%
\pgfpathlineto{\pgfqpoint{4.114134in}{0.506457in}}%
\pgfpathlineto{\pgfqpoint{4.116403in}{0.473227in}}%
\pgfpathlineto{\pgfqpoint{4.118673in}{0.507621in}}%
\pgfpathlineto{\pgfqpoint{4.120942in}{0.450298in}}%
\pgfpathlineto{\pgfqpoint{4.123212in}{0.470101in}}%
\pgfpathlineto{\pgfqpoint{4.125481in}{0.391845in}}%
\pgfpathlineto{\pgfqpoint{4.127751in}{0.454600in}}%
\pgfpathlineto{\pgfqpoint{4.130020in}{0.498450in}}%
\pgfpathlineto{\pgfqpoint{4.132290in}{0.505829in}}%
\pgfpathlineto{\pgfqpoint{4.134559in}{0.461938in}}%
\pgfpathlineto{\pgfqpoint{4.136829in}{0.492338in}}%
\pgfpathlineto{\pgfqpoint{4.139098in}{0.514736in}}%
\pgfpathlineto{\pgfqpoint{4.141368in}{0.528512in}}%
\pgfpathlineto{\pgfqpoint{4.143637in}{0.430356in}}%
\pgfpathlineto{\pgfqpoint{4.145907in}{0.454681in}}%
\pgfpathlineto{\pgfqpoint{4.148176in}{0.485259in}}%
\pgfpathlineto{\pgfqpoint{4.150446in}{0.429275in}}%
\pgfpathlineto{\pgfqpoint{4.152715in}{0.488153in}}%
\pgfpathlineto{\pgfqpoint{4.154985in}{0.467866in}}%
\pgfpathlineto{\pgfqpoint{4.157254in}{0.481730in}}%
\pgfpathlineto{\pgfqpoint{4.159524in}{0.463186in}}%
\pgfpathlineto{\pgfqpoint{4.161793in}{0.464807in}}%
\pgfpathlineto{\pgfqpoint{4.164062in}{0.511681in}}%
\pgfpathlineto{\pgfqpoint{4.166332in}{0.515840in}}%
\pgfpathlineto{\pgfqpoint{4.168601in}{0.466764in}}%
\pgfpathlineto{\pgfqpoint{4.170871in}{0.509175in}}%
\pgfpathlineto{\pgfqpoint{4.173140in}{0.503332in}}%
\pgfpathlineto{\pgfqpoint{4.175410in}{0.529290in}}%
\pgfpathlineto{\pgfqpoint{4.177679in}{0.434530in}}%
\pgfpathlineto{\pgfqpoint{4.179949in}{0.486022in}}%
\pgfpathlineto{\pgfqpoint{4.182218in}{0.395273in}}%
\pgfpathlineto{\pgfqpoint{4.184488in}{0.365627in}}%
\pgfpathlineto{\pgfqpoint{4.186757in}{0.458069in}}%
\pgfpathlineto{\pgfqpoint{4.189027in}{0.480648in}}%
\pgfpathlineto{\pgfqpoint{4.191296in}{0.492737in}}%
\pgfpathlineto{\pgfqpoint{4.193566in}{0.484010in}}%
\pgfpathlineto{\pgfqpoint{4.195835in}{0.486362in}}%
\pgfpathlineto{\pgfqpoint{4.198105in}{0.528283in}}%
\pgfpathlineto{\pgfqpoint{4.200374in}{0.462317in}}%
\pgfpathlineto{\pgfqpoint{4.202644in}{0.492700in}}%
\pgfpathlineto{\pgfqpoint{4.204913in}{0.382749in}}%
\pgfpathlineto{\pgfqpoint{4.213991in}{0.493898in}}%
\pgfpathlineto{\pgfqpoint{4.216261in}{0.493376in}}%
\pgfpathlineto{\pgfqpoint{4.218530in}{0.467742in}}%
\pgfpathlineto{\pgfqpoint{4.220800in}{0.524536in}}%
\pgfpathlineto{\pgfqpoint{4.223069in}{0.470217in}}%
\pgfpathlineto{\pgfqpoint{4.225339in}{0.497967in}}%
\pgfpathlineto{\pgfqpoint{4.227608in}{0.368979in}}%
\pgfpathlineto{\pgfqpoint{4.229878in}{0.430385in}}%
\pgfpathlineto{\pgfqpoint{4.232147in}{0.434506in}}%
\pgfpathlineto{\pgfqpoint{4.234416in}{0.473069in}}%
\pgfpathlineto{\pgfqpoint{4.236686in}{0.496873in}}%
\pgfpathlineto{\pgfqpoint{4.238955in}{0.496922in}}%
\pgfpathlineto{\pgfqpoint{4.241225in}{0.448251in}}%
\pgfpathlineto{\pgfqpoint{4.243494in}{0.521064in}}%
\pgfpathlineto{\pgfqpoint{4.245764in}{0.473409in}}%
\pgfpathlineto{\pgfqpoint{4.248033in}{0.502153in}}%
\pgfpathlineto{\pgfqpoint{4.250303in}{0.413466in}}%
\pgfpathlineto{\pgfqpoint{4.252572in}{0.445426in}}%
\pgfpathlineto{\pgfqpoint{4.254842in}{0.414463in}}%
\pgfpathlineto{\pgfqpoint{4.257111in}{0.469437in}}%
\pgfpathlineto{\pgfqpoint{4.259381in}{0.498397in}}%
\pgfpathlineto{\pgfqpoint{4.261650in}{0.500505in}}%
\pgfpathlineto{\pgfqpoint{4.263920in}{0.399188in}}%
\pgfpathlineto{\pgfqpoint{4.266189in}{0.515814in}}%
\pgfpathlineto{\pgfqpoint{4.268459in}{0.474679in}}%
\pgfpathlineto{\pgfqpoint{4.270728in}{0.505221in}}%
\pgfpathlineto{\pgfqpoint{4.272998in}{0.434065in}}%
\pgfpathlineto{\pgfqpoint{4.275267in}{0.457218in}}%
\pgfpathlineto{\pgfqpoint{4.277537in}{0.370310in}}%
\pgfpathlineto{\pgfqpoint{4.279806in}{0.464673in}}%
\pgfpathlineto{\pgfqpoint{4.282076in}{0.499183in}}%
\pgfpathlineto{\pgfqpoint{4.284345in}{0.503271in}}%
\pgfpathlineto{\pgfqpoint{4.286615in}{0.408167in}}%
\pgfpathlineto{\pgfqpoint{4.288884in}{0.509516in}}%
\pgfpathlineto{\pgfqpoint{4.291154in}{0.473925in}}%
\pgfpathlineto{\pgfqpoint{4.293423in}{0.507165in}}%
\pgfpathlineto{\pgfqpoint{4.295693in}{0.446407in}}%
\pgfpathlineto{\pgfqpoint{4.297962in}{0.466851in}}%
\pgfpathlineto{\pgfqpoint{4.300231in}{0.368554in}}%
\pgfpathlineto{\pgfqpoint{4.302501in}{0.457736in}}%
\pgfpathlineto{\pgfqpoint{4.304770in}{0.498806in}}%
\pgfpathlineto{\pgfqpoint{4.307040in}{0.505368in}}%
\pgfpathlineto{\pgfqpoint{4.309309in}{0.452348in}}%
\pgfpathlineto{\pgfqpoint{4.311579in}{0.497530in}}%
\pgfpathlineto{\pgfqpoint{4.313848in}{0.516285in}}%
\pgfpathlineto{\pgfqpoint{4.316118in}{0.527250in}}%
\pgfpathlineto{\pgfqpoint{4.318387in}{0.402466in}}%
\pgfpathlineto{\pgfqpoint{4.322926in}{0.489370in}}%
\pgfpathlineto{\pgfqpoint{4.325196in}{0.433126in}}%
\pgfpathlineto{\pgfqpoint{4.327465in}{0.489253in}}%
\pgfpathlineto{\pgfqpoint{4.329735in}{0.467090in}}%
\pgfpathlineto{\pgfqpoint{4.332004in}{0.479556in}}%
\pgfpathlineto{\pgfqpoint{4.336543in}{0.414240in}}%
\pgfpathlineto{\pgfqpoint{4.336543in}{0.414240in}}%
\pgfusepath{stroke}%
\end{pgfscope}%
\begin{pgfscope}%
\pgfpathrectangle{\pgfqpoint{2.520957in}{0.320625in}}{\pgfqpoint{1.815586in}{1.164375in}}%
\pgfusepath{clip}%
\pgfsetrectcap%
\pgfsetroundjoin%
\pgfsetlinewidth{1.154312pt}%
\definecolor{currentstroke}{rgb}{0.716387,0.214982,0.475290}%
\pgfsetstrokecolor{currentstroke}%
\pgfsetdash{}{0pt}%
\pgfpathmoveto{\pgfqpoint{2.520957in}{1.451682in}}%
\pgfpathlineto{\pgfqpoint{2.525496in}{1.442028in}}%
\pgfpathlineto{\pgfqpoint{2.543652in}{1.407088in}}%
\pgfpathlineto{\pgfqpoint{2.561807in}{1.369156in}}%
\pgfpathlineto{\pgfqpoint{2.616275in}{1.253910in}}%
\pgfpathlineto{\pgfqpoint{2.643509in}{1.199514in}}%
\pgfpathlineto{\pgfqpoint{2.670743in}{1.147522in}}%
\pgfpathlineto{\pgfqpoint{2.702515in}{1.089363in}}%
\pgfpathlineto{\pgfqpoint{2.738827in}{1.025323in}}%
\pgfpathlineto{\pgfqpoint{2.784217in}{0.947655in}}%
\pgfpathlineto{\pgfqpoint{2.845493in}{0.845206in}}%
\pgfpathlineto{\pgfqpoint{2.934003in}{0.699593in}}%
\pgfpathlineto{\pgfqpoint{2.999818in}{0.592358in}}%
\pgfpathlineto{\pgfqpoint{3.002087in}{0.589451in}}%
\pgfpathlineto{\pgfqpoint{3.006626in}{0.581280in}}%
\pgfpathlineto{\pgfqpoint{3.015704in}{0.566697in}}%
\pgfpathlineto{\pgfqpoint{3.022513in}{0.555919in}}%
\pgfpathlineto{\pgfqpoint{3.024782in}{0.556479in}}%
\pgfpathlineto{\pgfqpoint{3.029321in}{0.545042in}}%
\pgfpathlineto{\pgfqpoint{3.031590in}{0.543607in}}%
\pgfpathlineto{\pgfqpoint{3.033860in}{0.537909in}}%
\pgfpathlineto{\pgfqpoint{3.036129in}{0.541319in}}%
\pgfpathlineto{\pgfqpoint{3.038399in}{0.529967in}}%
\pgfpathlineto{\pgfqpoint{3.040668in}{0.529224in}}%
\pgfpathlineto{\pgfqpoint{3.042938in}{0.522427in}}%
\pgfpathlineto{\pgfqpoint{3.045207in}{0.518438in}}%
\pgfpathlineto{\pgfqpoint{3.047477in}{0.515879in}}%
\pgfpathlineto{\pgfqpoint{3.049746in}{0.519495in}}%
\pgfpathlineto{\pgfqpoint{3.052016in}{0.519363in}}%
\pgfpathlineto{\pgfqpoint{3.054285in}{0.507924in}}%
\pgfpathlineto{\pgfqpoint{3.056555in}{0.513192in}}%
\pgfpathlineto{\pgfqpoint{3.058824in}{0.537839in}}%
\pgfpathlineto{\pgfqpoint{3.061094in}{0.511025in}}%
\pgfpathlineto{\pgfqpoint{3.063363in}{0.507542in}}%
\pgfpathlineto{\pgfqpoint{3.065633in}{0.497432in}}%
\pgfpathlineto{\pgfqpoint{3.067902in}{0.481567in}}%
\pgfpathlineto{\pgfqpoint{3.070172in}{0.497713in}}%
\pgfpathlineto{\pgfqpoint{3.072441in}{0.480131in}}%
\pgfpathlineto{\pgfqpoint{3.074711in}{0.495167in}}%
\pgfpathlineto{\pgfqpoint{3.076980in}{0.486021in}}%
\pgfpathlineto{\pgfqpoint{3.079250in}{0.467959in}}%
\pgfpathlineto{\pgfqpoint{3.081519in}{0.472343in}}%
\pgfpathlineto{\pgfqpoint{3.083789in}{0.514645in}}%
\pgfpathlineto{\pgfqpoint{3.086058in}{0.514404in}}%
\pgfpathlineto{\pgfqpoint{3.088328in}{0.464471in}}%
\pgfpathlineto{\pgfqpoint{3.090597in}{0.444729in}}%
\pgfpathlineto{\pgfqpoint{3.092867in}{0.489071in}}%
\pgfpathlineto{\pgfqpoint{3.095136in}{0.455579in}}%
\pgfpathlineto{\pgfqpoint{3.097405in}{0.488237in}}%
\pgfpathlineto{\pgfqpoint{3.099675in}{0.477135in}}%
\pgfpathlineto{\pgfqpoint{3.101944in}{0.459143in}}%
\pgfpathlineto{\pgfqpoint{3.104214in}{0.429787in}}%
\pgfpathlineto{\pgfqpoint{3.106483in}{0.509272in}}%
\pgfpathlineto{\pgfqpoint{3.108753in}{0.517510in}}%
\pgfpathlineto{\pgfqpoint{3.111022in}{0.419928in}}%
\pgfpathlineto{\pgfqpoint{3.113292in}{0.427695in}}%
\pgfpathlineto{\pgfqpoint{3.115561in}{0.479791in}}%
\pgfpathlineto{\pgfqpoint{3.117831in}{0.445196in}}%
\pgfpathlineto{\pgfqpoint{3.120100in}{0.487281in}}%
\pgfpathlineto{\pgfqpoint{3.122370in}{0.478354in}}%
\pgfpathlineto{\pgfqpoint{3.124639in}{0.464484in}}%
\pgfpathlineto{\pgfqpoint{3.126909in}{0.386179in}}%
\pgfpathlineto{\pgfqpoint{3.129178in}{0.506166in}}%
\pgfpathlineto{\pgfqpoint{3.131448in}{0.520397in}}%
\pgfpathlineto{\pgfqpoint{3.133717in}{0.395501in}}%
\pgfpathlineto{\pgfqpoint{3.135987in}{0.444875in}}%
\pgfpathlineto{\pgfqpoint{3.138256in}{0.471412in}}%
\pgfpathlineto{\pgfqpoint{3.140526in}{0.437487in}}%
\pgfpathlineto{\pgfqpoint{3.142795in}{0.485641in}}%
\pgfpathlineto{\pgfqpoint{3.147334in}{0.472709in}}%
\pgfpathlineto{\pgfqpoint{3.149604in}{0.413114in}}%
\pgfpathlineto{\pgfqpoint{3.151873in}{0.501795in}}%
\pgfpathlineto{\pgfqpoint{3.154143in}{0.522484in}}%
\pgfpathlineto{\pgfqpoint{3.156412in}{0.434835in}}%
\pgfpathlineto{\pgfqpoint{3.158682in}{0.461703in}}%
\pgfpathlineto{\pgfqpoint{3.160951in}{0.460532in}}%
\pgfpathlineto{\pgfqpoint{3.163220in}{0.428001in}}%
\pgfpathlineto{\pgfqpoint{3.165490in}{0.483209in}}%
\pgfpathlineto{\pgfqpoint{3.167759in}{0.481037in}}%
\pgfpathlineto{\pgfqpoint{3.170029in}{0.479935in}}%
\pgfpathlineto{\pgfqpoint{3.172298in}{0.442249in}}%
\pgfpathlineto{\pgfqpoint{3.174568in}{0.496413in}}%
\pgfpathlineto{\pgfqpoint{3.176837in}{0.523713in}}%
\pgfpathlineto{\pgfqpoint{3.179107in}{0.454037in}}%
\pgfpathlineto{\pgfqpoint{3.181376in}{0.474260in}}%
\pgfpathlineto{\pgfqpoint{3.185915in}{0.412145in}}%
\pgfpathlineto{\pgfqpoint{3.188185in}{0.479637in}}%
\pgfpathlineto{\pgfqpoint{3.190454in}{0.481664in}}%
\pgfpathlineto{\pgfqpoint{3.192724in}{0.486072in}}%
\pgfpathlineto{\pgfqpoint{3.194993in}{0.459784in}}%
\pgfpathlineto{\pgfqpoint{3.199532in}{0.523933in}}%
\pgfpathlineto{\pgfqpoint{3.201802in}{0.508137in}}%
\pgfpathlineto{\pgfqpoint{3.204071in}{0.459226in}}%
\pgfpathlineto{\pgfqpoint{3.206341in}{0.509533in}}%
\pgfpathlineto{\pgfqpoint{3.208610in}{0.480212in}}%
\pgfpathlineto{\pgfqpoint{3.210880in}{0.518754in}}%
\pgfpathlineto{\pgfqpoint{3.213149in}{0.448327in}}%
\pgfpathlineto{\pgfqpoint{3.215419in}{0.486474in}}%
\pgfpathlineto{\pgfqpoint{3.219958in}{0.363357in}}%
\pgfpathlineto{\pgfqpoint{3.222227in}{0.462755in}}%
\pgfpathlineto{\pgfqpoint{3.224497in}{0.497238in}}%
\pgfpathlineto{\pgfqpoint{3.226766in}{0.498246in}}%
\pgfpathlineto{\pgfqpoint{3.229035in}{0.459364in}}%
\pgfpathlineto{\pgfqpoint{3.231305in}{0.503605in}}%
\pgfpathlineto{\pgfqpoint{3.233574in}{0.535197in}}%
\pgfpathlineto{\pgfqpoint{3.235844in}{0.495693in}}%
\pgfpathlineto{\pgfqpoint{3.238113in}{0.486605in}}%
\pgfpathlineto{\pgfqpoint{3.240383in}{0.489964in}}%
\pgfpathlineto{\pgfqpoint{3.242652in}{0.406337in}}%
\pgfpathlineto{\pgfqpoint{3.244922in}{0.491355in}}%
\pgfpathlineto{\pgfqpoint{3.247191in}{0.451624in}}%
\pgfpathlineto{\pgfqpoint{3.249461in}{0.486301in}}%
\pgfpathlineto{\pgfqpoint{3.251730in}{0.472215in}}%
\pgfpathlineto{\pgfqpoint{3.254000in}{0.415588in}}%
\pgfpathlineto{\pgfqpoint{3.258539in}{0.514563in}}%
\pgfpathlineto{\pgfqpoint{3.260808in}{0.508935in}}%
\pgfpathlineto{\pgfqpoint{3.263078in}{0.467461in}}%
\pgfpathlineto{\pgfqpoint{3.265347in}{0.389949in}}%
\pgfpathlineto{\pgfqpoint{3.267617in}{0.491285in}}%
\pgfpathlineto{\pgfqpoint{3.269886in}{0.446048in}}%
\pgfpathlineto{\pgfqpoint{3.272156in}{0.486402in}}%
\pgfpathlineto{\pgfqpoint{3.274425in}{0.472367in}}%
\pgfpathlineto{\pgfqpoint{3.276695in}{0.443473in}}%
\pgfpathlineto{\pgfqpoint{3.278964in}{0.433479in}}%
\pgfpathlineto{\pgfqpoint{3.281234in}{0.510907in}}%
\pgfpathlineto{\pgfqpoint{3.283503in}{0.514135in}}%
\pgfpathlineto{\pgfqpoint{3.288042in}{0.366904in}}%
\pgfpathlineto{\pgfqpoint{3.290312in}{0.484787in}}%
\pgfpathlineto{\pgfqpoint{3.292581in}{0.446040in}}%
\pgfpathlineto{\pgfqpoint{3.294851in}{0.487596in}}%
\pgfpathlineto{\pgfqpoint{3.297120in}{0.476052in}}%
\pgfpathlineto{\pgfqpoint{3.299389in}{0.455061in}}%
\pgfpathlineto{\pgfqpoint{3.301659in}{0.410690in}}%
\pgfpathlineto{\pgfqpoint{3.303928in}{0.508577in}}%
\pgfpathlineto{\pgfqpoint{3.306198in}{0.517842in}}%
\pgfpathlineto{\pgfqpoint{3.308467in}{0.376776in}}%
\pgfpathlineto{\pgfqpoint{3.313006in}{0.478180in}}%
\pgfpathlineto{\pgfqpoint{3.315276in}{0.442439in}}%
\pgfpathlineto{\pgfqpoint{3.317545in}{0.487273in}}%
\pgfpathlineto{\pgfqpoint{3.319815in}{0.478280in}}%
\pgfpathlineto{\pgfqpoint{3.322084in}{0.465526in}}%
\pgfpathlineto{\pgfqpoint{3.324230in}{0.315625in}}%
\pgfpathmoveto{\pgfqpoint{3.324453in}{0.315625in}}%
\pgfpathlineto{\pgfqpoint{3.326623in}{0.505005in}}%
\pgfpathlineto{\pgfqpoint{3.328893in}{0.520679in}}%
\pgfpathlineto{\pgfqpoint{3.331162in}{0.405470in}}%
\pgfpathlineto{\pgfqpoint{3.333432in}{0.448158in}}%
\pgfpathlineto{\pgfqpoint{3.335701in}{0.469401in}}%
\pgfpathlineto{\pgfqpoint{3.337971in}{0.436526in}}%
\pgfpathlineto{\pgfqpoint{3.340240in}{0.485775in}}%
\pgfpathlineto{\pgfqpoint{3.344779in}{0.474056in}}%
\pgfpathlineto{\pgfqpoint{3.347049in}{0.416832in}}%
\pgfpathlineto{\pgfqpoint{3.349318in}{0.500597in}}%
\pgfpathlineto{\pgfqpoint{3.351588in}{0.522651in}}%
\pgfpathlineto{\pgfqpoint{3.353857in}{0.440797in}}%
\pgfpathlineto{\pgfqpoint{3.356127in}{0.464478in}}%
\pgfpathlineto{\pgfqpoint{3.358396in}{0.457965in}}%
\pgfpathlineto{\pgfqpoint{3.360666in}{0.426530in}}%
\pgfpathlineto{\pgfqpoint{3.362935in}{0.483132in}}%
\pgfpathlineto{\pgfqpoint{3.365204in}{0.481336in}}%
\pgfpathlineto{\pgfqpoint{3.367474in}{0.481221in}}%
\pgfpathlineto{\pgfqpoint{3.369743in}{0.445112in}}%
\pgfpathlineto{\pgfqpoint{3.372013in}{0.494882in}}%
\pgfpathlineto{\pgfqpoint{3.374282in}{0.523681in}}%
\pgfpathlineto{\pgfqpoint{3.376552in}{0.457519in}}%
\pgfpathlineto{\pgfqpoint{3.378821in}{0.476481in}}%
\pgfpathlineto{\pgfqpoint{3.383360in}{0.408959in}}%
\pgfpathlineto{\pgfqpoint{3.385630in}{0.479313in}}%
\pgfpathlineto{\pgfqpoint{3.387899in}{0.481865in}}%
\pgfpathlineto{\pgfqpoint{3.390169in}{0.487225in}}%
\pgfpathlineto{\pgfqpoint{3.392438in}{0.461801in}}%
\pgfpathlineto{\pgfqpoint{3.394708in}{0.487395in}}%
\pgfpathlineto{\pgfqpoint{3.396977in}{0.523668in}}%
\pgfpathlineto{\pgfqpoint{3.399247in}{0.508647in}}%
\pgfpathlineto{\pgfqpoint{3.401516in}{0.466422in}}%
\pgfpathlineto{\pgfqpoint{3.403786in}{0.506523in}}%
\pgfpathlineto{\pgfqpoint{3.406055in}{0.478819in}}%
\pgfpathlineto{\pgfqpoint{3.408325in}{0.518340in}}%
\pgfpathlineto{\pgfqpoint{3.410594in}{0.450786in}}%
\pgfpathlineto{\pgfqpoint{3.412864in}{0.487582in}}%
\pgfpathlineto{\pgfqpoint{3.417403in}{0.376002in}}%
\pgfpathlineto{\pgfqpoint{3.419672in}{0.461695in}}%
\pgfpathlineto{\pgfqpoint{3.421942in}{0.496971in}}%
\pgfpathlineto{\pgfqpoint{3.424211in}{0.499357in}}%
\pgfpathlineto{\pgfqpoint{3.426481in}{0.463192in}}%
\pgfpathlineto{\pgfqpoint{3.431019in}{0.534792in}}%
\pgfpathlineto{\pgfqpoint{3.433289in}{0.497124in}}%
\pgfpathlineto{\pgfqpoint{3.435558in}{0.490223in}}%
\pgfpathlineto{\pgfqpoint{3.437828in}{0.486221in}}%
\pgfpathlineto{\pgfqpoint{3.440097in}{0.398637in}}%
\pgfpathlineto{\pgfqpoint{3.442367in}{0.490313in}}%
\pgfpathlineto{\pgfqpoint{3.444636in}{0.452132in}}%
\pgfpathlineto{\pgfqpoint{3.446906in}{0.486990in}}%
\pgfpathlineto{\pgfqpoint{3.449175in}{0.473120in}}%
\pgfpathlineto{\pgfqpoint{3.451445in}{0.421529in}}%
\pgfpathlineto{\pgfqpoint{3.455984in}{0.513953in}}%
\pgfpathlineto{\pgfqpoint{3.458253in}{0.510307in}}%
\pgfpathlineto{\pgfqpoint{3.460523in}{0.461788in}}%
\pgfpathlineto{\pgfqpoint{3.462792in}{0.371445in}}%
\pgfpathlineto{\pgfqpoint{3.465062in}{0.490014in}}%
\pgfpathlineto{\pgfqpoint{3.467331in}{0.446286in}}%
\pgfpathlineto{\pgfqpoint{3.469601in}{0.486873in}}%
\pgfpathlineto{\pgfqpoint{3.471870in}{0.473188in}}%
\pgfpathlineto{\pgfqpoint{3.474140in}{0.446812in}}%
\pgfpathlineto{\pgfqpoint{3.476409in}{0.428082in}}%
\pgfpathlineto{\pgfqpoint{3.478679in}{0.510120in}}%
\pgfpathlineto{\pgfqpoint{3.480948in}{0.515147in}}%
\pgfpathlineto{\pgfqpoint{3.483218in}{0.424013in}}%
\pgfpathlineto{\pgfqpoint{3.485487in}{0.389802in}}%
\pgfpathlineto{\pgfqpoint{3.487757in}{0.483137in}}%
\pgfpathlineto{\pgfqpoint{3.490026in}{0.445723in}}%
\pgfpathlineto{\pgfqpoint{3.492296in}{0.487834in}}%
\pgfpathlineto{\pgfqpoint{3.494565in}{0.476766in}}%
\pgfpathlineto{\pgfqpoint{3.496834in}{0.457941in}}%
\pgfpathlineto{\pgfqpoint{3.499104in}{0.399156in}}%
\pgfpathlineto{\pgfqpoint{3.501373in}{0.507619in}}%
\pgfpathlineto{\pgfqpoint{3.503643in}{0.518647in}}%
\pgfpathlineto{\pgfqpoint{3.505583in}{0.315625in}}%
\pgfpathmoveto{\pgfqpoint{3.506435in}{0.315625in}}%
\pgfpathlineto{\pgfqpoint{3.508182in}{0.431158in}}%
\pgfpathlineto{\pgfqpoint{3.510451in}{0.476061in}}%
\pgfpathlineto{\pgfqpoint{3.512721in}{0.441485in}}%
\pgfpathlineto{\pgfqpoint{3.514990in}{0.487190in}}%
\pgfpathlineto{\pgfqpoint{3.517260in}{0.478865in}}%
\pgfpathlineto{\pgfqpoint{3.519529in}{0.467938in}}%
\pgfpathlineto{\pgfqpoint{3.521799in}{0.365824in}}%
\pgfpathlineto{\pgfqpoint{3.524068in}{0.503806in}}%
\pgfpathlineto{\pgfqpoint{3.526338in}{0.521249in}}%
\pgfpathlineto{\pgfqpoint{3.528607in}{0.419838in}}%
\pgfpathlineto{\pgfqpoint{3.530877in}{0.453210in}}%
\pgfpathlineto{\pgfqpoint{3.533146in}{0.466609in}}%
\pgfpathlineto{\pgfqpoint{3.535416in}{0.434668in}}%
\pgfpathlineto{\pgfqpoint{3.537685in}{0.485362in}}%
\pgfpathlineto{\pgfqpoint{3.542224in}{0.476118in}}%
\pgfpathlineto{\pgfqpoint{3.544494in}{0.425782in}}%
\pgfpathlineto{\pgfqpoint{3.546763in}{0.499080in}}%
\pgfpathlineto{\pgfqpoint{3.549033in}{0.522975in}}%
\pgfpathlineto{\pgfqpoint{3.551302in}{0.446803in}}%
\pgfpathlineto{\pgfqpoint{3.553572in}{0.468135in}}%
\pgfpathlineto{\pgfqpoint{3.555841in}{0.454117in}}%
\pgfpathlineto{\pgfqpoint{3.558111in}{0.423157in}}%
\pgfpathlineto{\pgfqpoint{3.560380in}{0.482410in}}%
\pgfpathlineto{\pgfqpoint{3.562649in}{0.481603in}}%
\pgfpathlineto{\pgfqpoint{3.564919in}{0.482982in}}%
\pgfpathlineto{\pgfqpoint{3.567188in}{0.450024in}}%
\pgfpathlineto{\pgfqpoint{3.569458in}{0.492926in}}%
\pgfpathlineto{\pgfqpoint{3.571727in}{0.523729in}}%
\pgfpathlineto{\pgfqpoint{3.573997in}{0.505645in}}%
\pgfpathlineto{\pgfqpoint{3.576266in}{0.440182in}}%
\pgfpathlineto{\pgfqpoint{3.578536in}{0.513506in}}%
\pgfpathlineto{\pgfqpoint{3.580805in}{0.482686in}}%
\pgfpathlineto{\pgfqpoint{3.583075in}{0.518992in}}%
\pgfpathlineto{\pgfqpoint{3.585344in}{0.444926in}}%
\pgfpathlineto{\pgfqpoint{3.587614in}{0.483922in}}%
\pgfpathlineto{\pgfqpoint{3.589883in}{0.415913in}}%
\pgfpathlineto{\pgfqpoint{3.591683in}{0.315625in}}%
\pgfpathmoveto{\pgfqpoint{3.592485in}{0.315625in}}%
\pgfpathlineto{\pgfqpoint{3.594422in}{0.468144in}}%
\pgfpathlineto{\pgfqpoint{3.596692in}{0.497657in}}%
\pgfpathlineto{\pgfqpoint{3.598961in}{0.495194in}}%
\pgfpathlineto{\pgfqpoint{3.601231in}{0.445164in}}%
\pgfpathlineto{\pgfqpoint{3.603500in}{0.506840in}}%
\pgfpathlineto{\pgfqpoint{3.605770in}{0.534828in}}%
\pgfpathlineto{\pgfqpoint{3.608039in}{0.492114in}}%
\pgfpathlineto{\pgfqpoint{3.610309in}{0.478331in}}%
\pgfpathlineto{\pgfqpoint{3.612578in}{0.495039in}}%
\pgfpathlineto{\pgfqpoint{3.614848in}{0.421653in}}%
\pgfpathlineto{\pgfqpoint{3.617117in}{0.494009in}}%
\pgfpathlineto{\pgfqpoint{3.619387in}{0.452451in}}%
\pgfpathlineto{\pgfqpoint{3.621656in}{0.485524in}}%
\pgfpathlineto{\pgfqpoint{3.623926in}{0.469604in}}%
\pgfpathlineto{\pgfqpoint{3.626195in}{0.392197in}}%
\pgfpathlineto{\pgfqpoint{3.628465in}{0.473195in}}%
\pgfpathlineto{\pgfqpoint{3.630734in}{0.514927in}}%
\pgfpathlineto{\pgfqpoint{3.633003in}{0.505284in}}%
\pgfpathlineto{\pgfqpoint{3.635273in}{0.475320in}}%
\pgfpathlineto{\pgfqpoint{3.637542in}{0.414688in}}%
\pgfpathlineto{\pgfqpoint{3.639812in}{0.494152in}}%
\pgfpathlineto{\pgfqpoint{3.642081in}{0.447546in}}%
\pgfpathlineto{\pgfqpoint{3.644351in}{0.486223in}}%
\pgfpathlineto{\pgfqpoint{3.646620in}{0.470359in}}%
\pgfpathlineto{\pgfqpoint{3.648890in}{0.433535in}}%
\pgfpathlineto{\pgfqpoint{3.651159in}{0.447259in}}%
\pgfpathlineto{\pgfqpoint{3.653429in}{0.511753in}}%
\pgfpathlineto{\pgfqpoint{3.655698in}{0.511418in}}%
\pgfpathlineto{\pgfqpoint{3.657968in}{0.449576in}}%
\pgfpathlineto{\pgfqpoint{3.660237in}{0.340093in}}%
\pgfpathlineto{\pgfqpoint{3.662507in}{0.488258in}}%
\pgfpathlineto{\pgfqpoint{3.664776in}{0.448437in}}%
\pgfpathlineto{\pgfqpoint{3.667046in}{0.487912in}}%
\pgfpathlineto{\pgfqpoint{3.669315in}{0.474441in}}%
\pgfpathlineto{\pgfqpoint{3.671585in}{0.447406in}}%
\pgfpathlineto{\pgfqpoint{3.673854in}{0.432894in}}%
\pgfpathlineto{\pgfqpoint{3.676124in}{0.509753in}}%
\pgfpathlineto{\pgfqpoint{3.678393in}{0.515652in}}%
\pgfpathlineto{\pgfqpoint{3.680663in}{0.414925in}}%
\pgfpathlineto{\pgfqpoint{3.682932in}{0.398986in}}%
\pgfpathlineto{\pgfqpoint{3.685202in}{0.482340in}}%
\pgfpathlineto{\pgfqpoint{3.687471in}{0.445948in}}%
\pgfpathlineto{\pgfqpoint{3.689741in}{0.488250in}}%
\pgfpathlineto{\pgfqpoint{3.692010in}{0.477087in}}%
\pgfpathlineto{\pgfqpoint{3.694280in}{0.459667in}}%
\pgfpathlineto{\pgfqpoint{3.696549in}{0.390573in}}%
\pgfpathlineto{\pgfqpoint{3.698818in}{0.506628in}}%
\pgfpathlineto{\pgfqpoint{3.701088in}{0.519019in}}%
\pgfpathlineto{\pgfqpoint{3.703357in}{0.359315in}}%
\pgfpathlineto{\pgfqpoint{3.705627in}{0.436006in}}%
\pgfpathlineto{\pgfqpoint{3.707896in}{0.474542in}}%
\pgfpathlineto{\pgfqpoint{3.710166in}{0.441564in}}%
\pgfpathlineto{\pgfqpoint{3.712435in}{0.487525in}}%
\pgfpathlineto{\pgfqpoint{3.716974in}{0.469388in}}%
\pgfpathlineto{\pgfqpoint{3.719244in}{0.376977in}}%
\pgfpathlineto{\pgfqpoint{3.721513in}{0.502761in}}%
\pgfpathlineto{\pgfqpoint{3.723783in}{0.521490in}}%
\pgfpathlineto{\pgfqpoint{3.726052in}{0.428682in}}%
\pgfpathlineto{\pgfqpoint{3.728322in}{0.456419in}}%
\pgfpathlineto{\pgfqpoint{3.730591in}{0.464582in}}%
\pgfpathlineto{\pgfqpoint{3.732861in}{0.434019in}}%
\pgfpathlineto{\pgfqpoint{3.735130in}{0.485453in}}%
\pgfpathlineto{\pgfqpoint{3.739669in}{0.477456in}}%
\pgfpathlineto{\pgfqpoint{3.741939in}{0.429603in}}%
\pgfpathlineto{\pgfqpoint{3.744208in}{0.497752in}}%
\pgfpathlineto{\pgfqpoint{3.746478in}{0.523034in}}%
\pgfpathlineto{\pgfqpoint{3.748747in}{0.451163in}}%
\pgfpathlineto{\pgfqpoint{3.751017in}{0.470583in}}%
\pgfpathlineto{\pgfqpoint{3.753286in}{0.451114in}}%
\pgfpathlineto{\pgfqpoint{3.755556in}{0.421403in}}%
\pgfpathlineto{\pgfqpoint{3.757825in}{0.482278in}}%
\pgfpathlineto{\pgfqpoint{3.760095in}{0.481876in}}%
\pgfpathlineto{\pgfqpoint{3.762364in}{0.484185in}}%
\pgfpathlineto{\pgfqpoint{3.764633in}{0.452371in}}%
\pgfpathlineto{\pgfqpoint{3.769172in}{0.523581in}}%
\pgfpathlineto{\pgfqpoint{3.771442in}{0.506279in}}%
\pgfpathlineto{\pgfqpoint{3.773711in}{0.450885in}}%
\pgfpathlineto{\pgfqpoint{3.775981in}{0.510785in}}%
\pgfpathlineto{\pgfqpoint{3.778250in}{0.481640in}}%
\pgfpathlineto{\pgfqpoint{3.780520in}{0.518761in}}%
\pgfpathlineto{\pgfqpoint{3.782789in}{0.447862in}}%
\pgfpathlineto{\pgfqpoint{3.785059in}{0.485096in}}%
\pgfpathlineto{\pgfqpoint{3.787328in}{0.420201in}}%
\pgfpathlineto{\pgfqpoint{3.788975in}{0.315625in}}%
\pgfpathmoveto{\pgfqpoint{3.790067in}{0.315625in}}%
\pgfpathlineto{\pgfqpoint{3.791867in}{0.467465in}}%
\pgfpathlineto{\pgfqpoint{3.794137in}{0.497462in}}%
\pgfpathlineto{\pgfqpoint{3.796406in}{0.496381in}}%
\pgfpathlineto{\pgfqpoint{3.798676in}{0.450339in}}%
\pgfpathlineto{\pgfqpoint{3.800945in}{0.504711in}}%
\pgfpathlineto{\pgfqpoint{3.803215in}{0.534556in}}%
\pgfpathlineto{\pgfqpoint{3.805484in}{0.493877in}}%
\pgfpathlineto{\pgfqpoint{3.807754in}{0.482684in}}%
\pgfpathlineto{\pgfqpoint{3.810023in}{0.491705in}}%
\pgfpathlineto{\pgfqpoint{3.812293in}{0.417418in}}%
\pgfpathlineto{\pgfqpoint{3.814562in}{0.493138in}}%
\pgfpathlineto{\pgfqpoint{3.816832in}{0.453175in}}%
\pgfpathlineto{\pgfqpoint{3.819101in}{0.486305in}}%
\pgfpathlineto{\pgfqpoint{3.821371in}{0.470553in}}%
\pgfpathlineto{\pgfqpoint{3.823640in}{0.401205in}}%
\pgfpathlineto{\pgfqpoint{3.825910in}{0.471635in}}%
\pgfpathlineto{\pgfqpoint{3.828179in}{0.514385in}}%
\pgfpathlineto{\pgfqpoint{3.830448in}{0.506800in}}%
\pgfpathlineto{\pgfqpoint{3.832718in}{0.470608in}}%
\pgfpathlineto{\pgfqpoint{3.834987in}{0.406564in}}%
\pgfpathlineto{\pgfqpoint{3.837257in}{0.493063in}}%
\pgfpathlineto{\pgfqpoint{3.839526in}{0.448058in}}%
\pgfpathlineto{\pgfqpoint{3.841796in}{0.486823in}}%
\pgfpathlineto{\pgfqpoint{3.844065in}{0.471241in}}%
\pgfpathlineto{\pgfqpoint{3.846335in}{0.437424in}}%
\pgfpathlineto{\pgfqpoint{3.848604in}{0.443853in}}%
\pgfpathlineto{\pgfqpoint{3.850874in}{0.511057in}}%
\pgfpathlineto{\pgfqpoint{3.853143in}{0.512542in}}%
\pgfpathlineto{\pgfqpoint{3.855413in}{0.441460in}}%
\pgfpathlineto{\pgfqpoint{3.857199in}{0.315625in}}%
\pgfpathmoveto{\pgfqpoint{3.858059in}{0.315625in}}%
\pgfpathlineto{\pgfqpoint{3.859952in}{0.486841in}}%
\pgfpathlineto{\pgfqpoint{3.862221in}{0.448442in}}%
\pgfpathlineto{\pgfqpoint{3.864491in}{0.488314in}}%
\pgfpathlineto{\pgfqpoint{3.866760in}{0.475229in}}%
\pgfpathlineto{\pgfqpoint{3.871299in}{0.427016in}}%
\pgfpathlineto{\pgfqpoint{3.873569in}{0.508910in}}%
\pgfpathlineto{\pgfqpoint{3.875838in}{0.516566in}}%
\pgfpathlineto{\pgfqpoint{3.878108in}{0.395180in}}%
\pgfpathlineto{\pgfqpoint{3.880377in}{0.411745in}}%
\pgfpathlineto{\pgfqpoint{3.882647in}{0.480522in}}%
\pgfpathlineto{\pgfqpoint{3.884916in}{0.445416in}}%
\pgfpathlineto{\pgfqpoint{3.887186in}{0.488385in}}%
\pgfpathlineto{\pgfqpoint{3.889455in}{0.477755in}}%
\pgfpathlineto{\pgfqpoint{3.891725in}{0.462321in}}%
\pgfpathlineto{\pgfqpoint{3.893994in}{0.369497in}}%
\pgfpathlineto{\pgfqpoint{3.896264in}{0.505579in}}%
\pgfpathlineto{\pgfqpoint{3.898533in}{0.519699in}}%
\pgfpathlineto{\pgfqpoint{3.900802in}{0.393840in}}%
\pgfpathlineto{\pgfqpoint{3.903072in}{0.442326in}}%
\pgfpathlineto{\pgfqpoint{3.905341in}{0.472173in}}%
\pgfpathlineto{\pgfqpoint{3.907611in}{0.440322in}}%
\pgfpathlineto{\pgfqpoint{3.909880in}{0.487306in}}%
\pgfpathlineto{\pgfqpoint{3.914419in}{0.471632in}}%
\pgfpathlineto{\pgfqpoint{3.916689in}{0.396343in}}%
\pgfpathlineto{\pgfqpoint{3.918958in}{0.501441in}}%
\pgfpathlineto{\pgfqpoint{3.921228in}{0.521934in}}%
\pgfpathlineto{\pgfqpoint{3.923497in}{0.436947in}}%
\pgfpathlineto{\pgfqpoint{3.925767in}{0.460673in}}%
\pgfpathlineto{\pgfqpoint{3.928036in}{0.461400in}}%
\pgfpathlineto{\pgfqpoint{3.930306in}{0.431712in}}%
\pgfpathlineto{\pgfqpoint{3.932575in}{0.484939in}}%
\pgfpathlineto{\pgfqpoint{3.934845in}{0.481310in}}%
\pgfpathlineto{\pgfqpoint{3.937114in}{0.479366in}}%
\pgfpathlineto{\pgfqpoint{3.939384in}{0.436353in}}%
\pgfpathlineto{\pgfqpoint{3.941653in}{0.496065in}}%
\pgfpathlineto{\pgfqpoint{3.943923in}{0.523219in}}%
\pgfpathlineto{\pgfqpoint{3.946192in}{0.502595in}}%
\pgfpathlineto{\pgfqpoint{3.948462in}{0.403716in}}%
\pgfpathlineto{\pgfqpoint{3.950731in}{0.517131in}}%
\pgfpathlineto{\pgfqpoint{3.953001in}{0.484663in}}%
\pgfpathlineto{\pgfqpoint{3.955270in}{0.518839in}}%
\pgfpathlineto{\pgfqpoint{3.957540in}{0.440433in}}%
\pgfpathlineto{\pgfqpoint{3.959809in}{0.480919in}}%
\pgfpathlineto{\pgfqpoint{3.962079in}{0.399172in}}%
\pgfpathlineto{\pgfqpoint{3.964348in}{0.376189in}}%
\pgfpathlineto{\pgfqpoint{3.966617in}{0.472900in}}%
\pgfpathlineto{\pgfqpoint{3.968887in}{0.497880in}}%
\pgfpathlineto{\pgfqpoint{3.971156in}{0.491628in}}%
\pgfpathlineto{\pgfqpoint{3.973426in}{0.422083in}}%
\pgfpathlineto{\pgfqpoint{3.975695in}{0.509661in}}%
\pgfpathlineto{\pgfqpoint{3.977965in}{0.534137in}}%
\pgfpathlineto{\pgfqpoint{3.980234in}{0.487636in}}%
\pgfpathlineto{\pgfqpoint{3.982504in}{0.467720in}}%
\pgfpathlineto{\pgfqpoint{3.984773in}{0.499690in}}%
\pgfpathlineto{\pgfqpoint{3.987043in}{0.432350in}}%
\pgfpathlineto{\pgfqpoint{3.989312in}{0.496496in}}%
\pgfpathlineto{\pgfqpoint{3.991582in}{0.452905in}}%
\pgfpathlineto{\pgfqpoint{3.993851in}{0.484424in}}%
\pgfpathlineto{\pgfqpoint{3.996121in}{0.466458in}}%
\pgfpathlineto{\pgfqpoint{3.998390in}{0.332574in}}%
\pgfpathlineto{\pgfqpoint{4.000660in}{0.478819in}}%
\pgfpathlineto{\pgfqpoint{4.002929in}{0.515085in}}%
\pgfpathlineto{\pgfqpoint{4.005199in}{0.501009in}}%
\pgfpathlineto{\pgfqpoint{4.007468in}{0.482195in}}%
\pgfpathlineto{\pgfqpoint{4.009738in}{0.429755in}}%
\pgfpathlineto{\pgfqpoint{4.012007in}{0.496844in}}%
\pgfpathlineto{\pgfqpoint{4.014277in}{0.448608in}}%
\pgfpathlineto{\pgfqpoint{4.016546in}{0.485711in}}%
\pgfpathlineto{\pgfqpoint{4.018816in}{0.467919in}}%
\pgfpathlineto{\pgfqpoint{4.021085in}{0.420060in}}%
\pgfpathlineto{\pgfqpoint{4.023355in}{0.457706in}}%
\pgfpathlineto{\pgfqpoint{4.025624in}{0.512396in}}%
\pgfpathlineto{\pgfqpoint{4.027894in}{0.508251in}}%
\pgfpathlineto{\pgfqpoint{4.030163in}{0.460582in}}%
\pgfpathlineto{\pgfqpoint{4.032432in}{0.396835in}}%
\pgfpathlineto{\pgfqpoint{4.034702in}{0.491504in}}%
\pgfpathlineto{\pgfqpoint{4.036971in}{0.450359in}}%
\pgfpathlineto{\pgfqpoint{4.039241in}{0.487871in}}%
\pgfpathlineto{\pgfqpoint{4.041510in}{0.472460in}}%
\pgfpathlineto{\pgfqpoint{4.043780in}{0.437693in}}%
\pgfpathlineto{\pgfqpoint{4.046049in}{0.447563in}}%
\pgfpathlineto{\pgfqpoint{4.048319in}{0.510713in}}%
\pgfpathlineto{\pgfqpoint{4.050588in}{0.513080in}}%
\pgfpathlineto{\pgfqpoint{4.052858in}{0.435616in}}%
\pgfpathlineto{\pgfqpoint{4.055127in}{0.338079in}}%
\pgfpathlineto{\pgfqpoint{4.057397in}{0.486200in}}%
\pgfpathlineto{\pgfqpoint{4.059666in}{0.448856in}}%
\pgfpathlineto{\pgfqpoint{4.061936in}{0.488818in}}%
\pgfpathlineto{\pgfqpoint{4.064205in}{0.475574in}}%
\pgfpathlineto{\pgfqpoint{4.066475in}{0.452559in}}%
\pgfpathlineto{\pgfqpoint{4.068744in}{0.423166in}}%
\pgfpathlineto{\pgfqpoint{4.071014in}{0.508012in}}%
\pgfpathlineto{\pgfqpoint{4.073283in}{0.517023in}}%
\pgfpathlineto{\pgfqpoint{4.075553in}{0.365020in}}%
\pgfpathlineto{\pgfqpoint{4.080092in}{0.479246in}}%
\pgfpathlineto{\pgfqpoint{4.082361in}{0.445746in}}%
\pgfpathlineto{\pgfqpoint{4.084631in}{0.488835in}}%
\pgfpathlineto{\pgfqpoint{4.086900in}{0.478272in}}%
\pgfpathlineto{\pgfqpoint{4.089170in}{0.463861in}}%
\pgfpathlineto{\pgfqpoint{4.091439in}{0.357493in}}%
\pgfpathlineto{\pgfqpoint{4.093709in}{0.504644in}}%
\pgfpathlineto{\pgfqpoint{4.095978in}{0.520022in}}%
\pgfpathlineto{\pgfqpoint{4.098247in}{0.409569in}}%
\pgfpathlineto{\pgfqpoint{4.100517in}{0.446227in}}%
\pgfpathlineto{\pgfqpoint{4.102786in}{0.470493in}}%
\pgfpathlineto{\pgfqpoint{4.105056in}{0.440065in}}%
\pgfpathlineto{\pgfqpoint{4.107325in}{0.487504in}}%
\pgfpathlineto{\pgfqpoint{4.111864in}{0.473057in}}%
\pgfpathlineto{\pgfqpoint{4.114134in}{0.403295in}}%
\pgfpathlineto{\pgfqpoint{4.116403in}{0.500267in}}%
\pgfpathlineto{\pgfqpoint{4.118673in}{0.522089in}}%
\pgfpathlineto{\pgfqpoint{4.120942in}{0.442700in}}%
\pgfpathlineto{\pgfqpoint{4.123212in}{0.463483in}}%
\pgfpathlineto{\pgfqpoint{4.125481in}{0.458976in}}%
\pgfpathlineto{\pgfqpoint{4.127751in}{0.430687in}}%
\pgfpathlineto{\pgfqpoint{4.130020in}{0.484938in}}%
\pgfpathlineto{\pgfqpoint{4.132290in}{0.481643in}}%
\pgfpathlineto{\pgfqpoint{4.134559in}{0.480648in}}%
\pgfpathlineto{\pgfqpoint{4.136829in}{0.439450in}}%
\pgfpathlineto{\pgfqpoint{4.139098in}{0.494599in}}%
\pgfpathlineto{\pgfqpoint{4.141368in}{0.523185in}}%
\pgfpathlineto{\pgfqpoint{4.143637in}{0.503400in}}%
\pgfpathlineto{\pgfqpoint{4.145907in}{0.425690in}}%
\pgfpathlineto{\pgfqpoint{4.148176in}{0.514649in}}%
\pgfpathlineto{\pgfqpoint{4.150446in}{0.483904in}}%
\pgfpathlineto{\pgfqpoint{4.152715in}{0.518773in}}%
\pgfpathlineto{\pgfqpoint{4.154985in}{0.443973in}}%
\pgfpathlineto{\pgfqpoint{4.157254in}{0.482177in}}%
\pgfpathlineto{\pgfqpoint{4.159524in}{0.405227in}}%
\pgfpathlineto{\pgfqpoint{4.161793in}{0.367283in}}%
\pgfpathlineto{\pgfqpoint{4.164062in}{0.472465in}}%
\pgfpathlineto{\pgfqpoint{4.166332in}{0.497746in}}%
\pgfpathlineto{\pgfqpoint{4.168601in}{0.492927in}}%
\pgfpathlineto{\pgfqpoint{4.170871in}{0.430700in}}%
\pgfpathlineto{\pgfqpoint{4.173140in}{0.507763in}}%
\pgfpathlineto{\pgfqpoint{4.175410in}{0.534000in}}%
\pgfpathlineto{\pgfqpoint{4.177679in}{0.489844in}}%
\pgfpathlineto{\pgfqpoint{4.179949in}{0.473263in}}%
\pgfpathlineto{\pgfqpoint{4.182218in}{0.496682in}}%
\pgfpathlineto{\pgfqpoint{4.184488in}{0.429743in}}%
\pgfpathlineto{\pgfqpoint{4.186757in}{0.495756in}}%
\pgfpathlineto{\pgfqpoint{4.189027in}{0.453812in}}%
\pgfpathlineto{\pgfqpoint{4.191296in}{0.485287in}}%
\pgfpathlineto{\pgfqpoint{4.193566in}{0.467479in}}%
\pgfpathlineto{\pgfqpoint{4.195835in}{0.359653in}}%
\pgfpathlineto{\pgfqpoint{4.198105in}{0.477637in}}%
\pgfpathlineto{\pgfqpoint{4.200374in}{0.514609in}}%
\pgfpathlineto{\pgfqpoint{4.202644in}{0.502729in}}%
\pgfpathlineto{\pgfqpoint{4.204913in}{0.478163in}}%
\pgfpathlineto{\pgfqpoint{4.207183in}{0.424896in}}%
\pgfpathlineto{\pgfqpoint{4.209452in}{0.495898in}}%
\pgfpathlineto{\pgfqpoint{4.211722in}{0.449345in}}%
\pgfpathlineto{\pgfqpoint{4.213991in}{0.486417in}}%
\pgfpathlineto{\pgfqpoint{4.216261in}{0.468880in}}%
\pgfpathlineto{\pgfqpoint{4.218530in}{0.425024in}}%
\pgfpathlineto{\pgfqpoint{4.220800in}{0.455310in}}%
\pgfpathlineto{\pgfqpoint{4.223069in}{0.511784in}}%
\pgfpathlineto{\pgfqpoint{4.225339in}{0.509523in}}%
\pgfpathlineto{\pgfqpoint{4.227608in}{0.454318in}}%
\pgfpathlineto{\pgfqpoint{4.229878in}{0.381870in}}%
\pgfpathlineto{\pgfqpoint{4.232147in}{0.490268in}}%
\pgfpathlineto{\pgfqpoint{4.234416in}{0.450616in}}%
\pgfpathlineto{\pgfqpoint{4.236686in}{0.488407in}}%
\pgfpathlineto{\pgfqpoint{4.238955in}{0.473333in}}%
\pgfpathlineto{\pgfqpoint{4.241225in}{0.441565in}}%
\pgfpathlineto{\pgfqpoint{4.243494in}{0.443769in}}%
\pgfpathlineto{\pgfqpoint{4.245764in}{0.509973in}}%
\pgfpathlineto{\pgfqpoint{4.248033in}{0.514129in}}%
\pgfpathlineto{\pgfqpoint{4.250303in}{0.424238in}}%
\pgfpathlineto{\pgfqpoint{4.252572in}{0.375142in}}%
\pgfpathlineto{\pgfqpoint{4.254842in}{0.484616in}}%
\pgfpathlineto{\pgfqpoint{4.257111in}{0.448640in}}%
\pgfpathlineto{\pgfqpoint{4.259381in}{0.489126in}}%
\pgfpathlineto{\pgfqpoint{4.261650in}{0.476331in}}%
\pgfpathlineto{\pgfqpoint{4.263920in}{0.455590in}}%
\pgfpathlineto{\pgfqpoint{4.266189in}{0.415133in}}%
\pgfpathlineto{\pgfqpoint{4.268459in}{0.507090in}}%
\pgfpathlineto{\pgfqpoint{4.270728in}{0.517831in}}%
\pgfpathlineto{\pgfqpoint{4.272719in}{0.315625in}}%
\pgfpathmoveto{\pgfqpoint{4.273455in}{0.315625in}}%
\pgfpathlineto{\pgfqpoint{4.275267in}{0.427542in}}%
\pgfpathlineto{\pgfqpoint{4.277537in}{0.477198in}}%
\pgfpathlineto{\pgfqpoint{4.279806in}{0.444940in}}%
\pgfpathlineto{\pgfqpoint{4.282076in}{0.488845in}}%
\pgfpathlineto{\pgfqpoint{4.284345in}{0.478894in}}%
\pgfpathlineto{\pgfqpoint{4.286615in}{0.466368in}}%
\pgfpathlineto{\pgfqpoint{4.288275in}{0.315625in}}%
\pgfpathmoveto{\pgfqpoint{4.289400in}{0.315625in}}%
\pgfpathlineto{\pgfqpoint{4.291154in}{0.503487in}}%
\pgfpathlineto{\pgfqpoint{4.293423in}{0.520595in}}%
\pgfpathlineto{\pgfqpoint{4.295693in}{0.422572in}}%
\pgfpathlineto{\pgfqpoint{4.297962in}{0.451433in}}%
\pgfpathlineto{\pgfqpoint{4.300231in}{0.467792in}}%
\pgfpathlineto{\pgfqpoint{4.302501in}{0.438447in}}%
\pgfpathlineto{\pgfqpoint{4.304770in}{0.487159in}}%
\pgfpathlineto{\pgfqpoint{4.309309in}{0.475167in}}%
\pgfpathlineto{\pgfqpoint{4.311579in}{0.414991in}}%
\pgfpathlineto{\pgfqpoint{4.313848in}{0.498797in}}%
\pgfpathlineto{\pgfqpoint{4.316118in}{0.522412in}}%
\pgfpathlineto{\pgfqpoint{4.318387in}{0.498925in}}%
\pgfpathlineto{\pgfqpoint{4.320657in}{0.361876in}}%
\pgfpathlineto{\pgfqpoint{4.322926in}{0.520419in}}%
\pgfpathlineto{\pgfqpoint{4.325196in}{0.486144in}}%
\pgfpathlineto{\pgfqpoint{4.327465in}{0.518305in}}%
\pgfpathlineto{\pgfqpoint{4.329735in}{0.434456in}}%
\pgfpathlineto{\pgfqpoint{4.332004in}{0.477434in}}%
\pgfpathlineto{\pgfqpoint{4.334274in}{0.455523in}}%
\pgfpathlineto{\pgfqpoint{4.336543in}{0.418154in}}%
\pgfpathlineto{\pgfqpoint{4.336543in}{0.418154in}}%
\pgfusepath{stroke}%
\end{pgfscope}%
\begin{pgfscope}%
\pgfpathrectangle{\pgfqpoint{2.520957in}{0.320625in}}{\pgfqpoint{1.815586in}{1.164375in}}%
\pgfusepath{clip}%
\pgfsetrectcap%
\pgfsetroundjoin%
\pgfsetlinewidth{1.154312pt}%
\definecolor{currentstroke}{rgb}{0.852126,0.276106,0.418573}%
\pgfsetstrokecolor{currentstroke}%
\pgfsetdash{}{0pt}%
\pgfpathmoveto{\pgfqpoint{2.520957in}{1.451682in}}%
\pgfpathlineto{\pgfqpoint{2.525496in}{1.441936in}}%
\pgfpathlineto{\pgfqpoint{2.541382in}{1.411125in}}%
\pgfpathlineto{\pgfqpoint{2.557269in}{1.377640in}}%
\pgfpathlineto{\pgfqpoint{2.648048in}{1.181564in}}%
\pgfpathlineto{\pgfqpoint{2.679821in}{1.117459in}}%
\pgfpathlineto{\pgfqpoint{2.713863in}{1.051317in}}%
\pgfpathlineto{\pgfqpoint{2.752444in}{0.978809in}}%
\pgfpathlineto{\pgfqpoint{2.800103in}{0.891720in}}%
\pgfpathlineto{\pgfqpoint{2.861379in}{0.782176in}}%
\pgfpathlineto{\pgfqpoint{2.947620in}{0.630379in}}%
\pgfpathlineto{\pgfqpoint{2.988470in}{0.559138in}}%
\pgfpathlineto{\pgfqpoint{2.990740in}{0.557622in}}%
\pgfpathlineto{\pgfqpoint{2.993009in}{0.553688in}}%
\pgfpathlineto{\pgfqpoint{2.995279in}{0.547424in}}%
\pgfpathlineto{\pgfqpoint{2.997548in}{0.543156in}}%
\pgfpathlineto{\pgfqpoint{2.999818in}{0.541243in}}%
\pgfpathlineto{\pgfqpoint{3.002087in}{0.535594in}}%
\pgfpathlineto{\pgfqpoint{3.004357in}{0.533799in}}%
\pgfpathlineto{\pgfqpoint{3.011165in}{0.520072in}}%
\pgfpathlineto{\pgfqpoint{3.013435in}{0.527792in}}%
\pgfpathlineto{\pgfqpoint{3.015704in}{0.526942in}}%
\pgfpathlineto{\pgfqpoint{3.017974in}{0.508205in}}%
\pgfpathlineto{\pgfqpoint{3.020243in}{0.503596in}}%
\pgfpathlineto{\pgfqpoint{3.022513in}{0.508563in}}%
\pgfpathlineto{\pgfqpoint{3.024782in}{0.497854in}}%
\pgfpathlineto{\pgfqpoint{3.027052in}{0.504941in}}%
\pgfpathlineto{\pgfqpoint{3.033860in}{0.481308in}}%
\pgfpathlineto{\pgfqpoint{3.036129in}{0.512181in}}%
\pgfpathlineto{\pgfqpoint{3.038399in}{0.519286in}}%
\pgfpathlineto{\pgfqpoint{3.040668in}{0.468213in}}%
\pgfpathlineto{\pgfqpoint{3.042938in}{0.466284in}}%
\pgfpathlineto{\pgfqpoint{3.045207in}{0.486577in}}%
\pgfpathlineto{\pgfqpoint{3.047477in}{0.466327in}}%
\pgfpathlineto{\pgfqpoint{3.049746in}{0.492340in}}%
\pgfpathlineto{\pgfqpoint{3.052016in}{0.482743in}}%
\pgfpathlineto{\pgfqpoint{3.054285in}{0.470510in}}%
\pgfpathlineto{\pgfqpoint{3.056555in}{0.441048in}}%
\pgfpathlineto{\pgfqpoint{3.058824in}{0.505596in}}%
\pgfpathlineto{\pgfqpoint{3.061094in}{0.520166in}}%
\pgfpathlineto{\pgfqpoint{3.063363in}{0.435005in}}%
\pgfpathlineto{\pgfqpoint{3.065633in}{0.451664in}}%
\pgfpathlineto{\pgfqpoint{3.067902in}{0.473021in}}%
\pgfpathlineto{\pgfqpoint{3.070172in}{0.446664in}}%
\pgfpathlineto{\pgfqpoint{3.072441in}{0.488389in}}%
\pgfpathlineto{\pgfqpoint{3.076980in}{0.473153in}}%
\pgfpathlineto{\pgfqpoint{3.079250in}{0.413493in}}%
\pgfpathlineto{\pgfqpoint{3.081519in}{0.500780in}}%
\pgfpathlineto{\pgfqpoint{3.083789in}{0.521926in}}%
\pgfpathlineto{\pgfqpoint{3.086058in}{0.442701in}}%
\pgfpathlineto{\pgfqpoint{3.088328in}{0.462543in}}%
\pgfpathlineto{\pgfqpoint{3.090597in}{0.460847in}}%
\pgfpathlineto{\pgfqpoint{3.092867in}{0.433835in}}%
\pgfpathlineto{\pgfqpoint{3.095136in}{0.485584in}}%
\pgfpathlineto{\pgfqpoint{3.097405in}{0.481666in}}%
\pgfpathlineto{\pgfqpoint{3.099675in}{0.479882in}}%
\pgfpathlineto{\pgfqpoint{3.101944in}{0.436270in}}%
\pgfpathlineto{\pgfqpoint{3.104214in}{0.495322in}}%
\pgfpathlineto{\pgfqpoint{3.106483in}{0.523033in}}%
\pgfpathlineto{\pgfqpoint{3.108753in}{0.502597in}}%
\pgfpathlineto{\pgfqpoint{3.111022in}{0.416939in}}%
\pgfpathlineto{\pgfqpoint{3.113292in}{0.515531in}}%
\pgfpathlineto{\pgfqpoint{3.115561in}{0.484429in}}%
\pgfpathlineto{\pgfqpoint{3.117831in}{0.518697in}}%
\pgfpathlineto{\pgfqpoint{3.120100in}{0.442859in}}%
\pgfpathlineto{\pgfqpoint{3.122370in}{0.481267in}}%
\pgfpathlineto{\pgfqpoint{3.124639in}{0.399781in}}%
\pgfpathlineto{\pgfqpoint{3.126909in}{0.380752in}}%
\pgfpathlineto{\pgfqpoint{3.129178in}{0.473925in}}%
\pgfpathlineto{\pgfqpoint{3.131448in}{0.497783in}}%
\pgfpathlineto{\pgfqpoint{3.133717in}{0.491703in}}%
\pgfpathlineto{\pgfqpoint{3.135987in}{0.421344in}}%
\pgfpathlineto{\pgfqpoint{3.138256in}{0.508635in}}%
\pgfpathlineto{\pgfqpoint{3.140526in}{0.533711in}}%
\pgfpathlineto{\pgfqpoint{3.142795in}{0.488248in}}%
\pgfpathlineto{\pgfqpoint{3.145065in}{0.469385in}}%
\pgfpathlineto{\pgfqpoint{3.147334in}{0.498292in}}%
\pgfpathlineto{\pgfqpoint{3.149604in}{0.433381in}}%
\pgfpathlineto{\pgfqpoint{3.151873in}{0.496685in}}%
\pgfpathlineto{\pgfqpoint{3.154143in}{0.453994in}}%
\pgfpathlineto{\pgfqpoint{3.156412in}{0.484816in}}%
\pgfpathlineto{\pgfqpoint{3.158682in}{0.466087in}}%
\pgfpathlineto{\pgfqpoint{3.160910in}{0.315625in}}%
\pgfpathmoveto{\pgfqpoint{3.160989in}{0.315625in}}%
\pgfpathlineto{\pgfqpoint{3.163220in}{0.479840in}}%
\pgfpathlineto{\pgfqpoint{3.165490in}{0.514610in}}%
\pgfpathlineto{\pgfqpoint{3.167759in}{0.500819in}}%
\pgfpathlineto{\pgfqpoint{3.170029in}{0.481004in}}%
\pgfpathlineto{\pgfqpoint{3.172298in}{0.430721in}}%
\pgfpathlineto{\pgfqpoint{3.174568in}{0.497079in}}%
\pgfpathlineto{\pgfqpoint{3.176837in}{0.449817in}}%
\pgfpathlineto{\pgfqpoint{3.179107in}{0.486133in}}%
\pgfpathlineto{\pgfqpoint{3.181376in}{0.467624in}}%
\pgfpathlineto{\pgfqpoint{3.183646in}{0.417525in}}%
\pgfpathlineto{\pgfqpoint{3.188185in}{0.512031in}}%
\pgfpathlineto{\pgfqpoint{3.190454in}{0.507878in}}%
\pgfpathlineto{\pgfqpoint{3.194993in}{0.399920in}}%
\pgfpathlineto{\pgfqpoint{3.197263in}{0.491902in}}%
\pgfpathlineto{\pgfqpoint{3.199532in}{0.451578in}}%
\pgfpathlineto{\pgfqpoint{3.201802in}{0.488325in}}%
\pgfpathlineto{\pgfqpoint{3.204071in}{0.472165in}}%
\pgfpathlineto{\pgfqpoint{3.206341in}{0.435638in}}%
\pgfpathlineto{\pgfqpoint{3.208610in}{0.450718in}}%
\pgfpathlineto{\pgfqpoint{3.210880in}{0.510431in}}%
\pgfpathlineto{\pgfqpoint{3.213149in}{0.512604in}}%
\pgfpathlineto{\pgfqpoint{3.215419in}{0.434837in}}%
\pgfpathlineto{\pgfqpoint{3.217632in}{0.315625in}}%
\pgfpathmoveto{\pgfqpoint{3.217727in}{0.315625in}}%
\pgfpathlineto{\pgfqpoint{3.219958in}{0.486806in}}%
\pgfpathlineto{\pgfqpoint{3.222227in}{0.450265in}}%
\pgfpathlineto{\pgfqpoint{3.224497in}{0.489380in}}%
\pgfpathlineto{\pgfqpoint{3.226766in}{0.475319in}}%
\pgfpathlineto{\pgfqpoint{3.231305in}{0.429570in}}%
\pgfpathlineto{\pgfqpoint{3.233574in}{0.507866in}}%
\pgfpathlineto{\pgfqpoint{3.235844in}{0.516493in}}%
\pgfpathlineto{\pgfqpoint{3.238113in}{0.366260in}}%
\pgfpathlineto{\pgfqpoint{3.240383in}{0.414578in}}%
\pgfpathlineto{\pgfqpoint{3.242652in}{0.480159in}}%
\pgfpathlineto{\pgfqpoint{3.244922in}{0.447512in}}%
\pgfpathlineto{\pgfqpoint{3.247191in}{0.489591in}}%
\pgfpathlineto{\pgfqpoint{3.249461in}{0.478075in}}%
\pgfpathlineto{\pgfqpoint{3.251730in}{0.462316in}}%
\pgfpathlineto{\pgfqpoint{3.254000in}{0.384618in}}%
\pgfpathlineto{\pgfqpoint{3.256269in}{0.504683in}}%
\pgfpathlineto{\pgfqpoint{3.258539in}{0.519484in}}%
\pgfpathlineto{\pgfqpoint{3.260808in}{0.407654in}}%
\pgfpathlineto{\pgfqpoint{3.265347in}{0.471860in}}%
\pgfpathlineto{\pgfqpoint{3.267617in}{0.442468in}}%
\pgfpathlineto{\pgfqpoint{3.269886in}{0.488584in}}%
\pgfpathlineto{\pgfqpoint{3.274425in}{0.471646in}}%
\pgfpathlineto{\pgfqpoint{3.276695in}{0.385409in}}%
\pgfpathlineto{\pgfqpoint{3.278964in}{0.500568in}}%
\pgfpathlineto{\pgfqpoint{3.281234in}{0.521591in}}%
\pgfpathlineto{\pgfqpoint{3.283503in}{0.495633in}}%
\pgfpathlineto{\pgfqpoint{3.285773in}{0.417824in}}%
\pgfpathlineto{\pgfqpoint{3.288042in}{0.522686in}}%
\pgfpathlineto{\pgfqpoint{3.290312in}{0.486963in}}%
\pgfpathlineto{\pgfqpoint{3.292581in}{0.517612in}}%
\pgfpathlineto{\pgfqpoint{3.294851in}{0.428358in}}%
\pgfpathlineto{\pgfqpoint{3.297120in}{0.474254in}}%
\pgfpathlineto{\pgfqpoint{3.298807in}{0.315625in}}%
\pgfpathmoveto{\pgfqpoint{3.300193in}{0.315625in}}%
\pgfpathlineto{\pgfqpoint{3.301659in}{0.415626in}}%
\pgfpathlineto{\pgfqpoint{3.303928in}{0.479996in}}%
\pgfpathlineto{\pgfqpoint{3.306198in}{0.497723in}}%
\pgfpathlineto{\pgfqpoint{3.308467in}{0.483550in}}%
\pgfpathlineto{\pgfqpoint{3.310737in}{0.382699in}}%
\pgfpathlineto{\pgfqpoint{3.313006in}{0.513830in}}%
\pgfpathlineto{\pgfqpoint{3.315276in}{0.532021in}}%
\pgfpathlineto{\pgfqpoint{3.317545in}{0.476096in}}%
\pgfpathlineto{\pgfqpoint{3.319815in}{0.435383in}}%
\pgfpathlineto{\pgfqpoint{3.322084in}{0.507305in}}%
\pgfpathlineto{\pgfqpoint{3.324354in}{0.445624in}}%
\pgfpathlineto{\pgfqpoint{3.326623in}{0.500622in}}%
\pgfpathlineto{\pgfqpoint{3.328893in}{0.452595in}}%
\pgfpathlineto{\pgfqpoint{3.331162in}{0.481390in}}%
\pgfpathlineto{\pgfqpoint{3.333432in}{0.458628in}}%
\pgfpathlineto{\pgfqpoint{3.335701in}{0.399925in}}%
\pgfpathlineto{\pgfqpoint{3.337971in}{0.487292in}}%
\pgfpathlineto{\pgfqpoint{3.340240in}{0.514775in}}%
\pgfpathlineto{\pgfqpoint{3.342510in}{0.490607in}}%
\pgfpathlineto{\pgfqpoint{3.344779in}{0.493294in}}%
\pgfpathlineto{\pgfqpoint{3.347049in}{0.447774in}}%
\pgfpathlineto{\pgfqpoint{3.349318in}{0.501496in}}%
\pgfpathlineto{\pgfqpoint{3.351588in}{0.449300in}}%
\pgfpathlineto{\pgfqpoint{3.353857in}{0.483691in}}%
\pgfpathlineto{\pgfqpoint{3.356127in}{0.461585in}}%
\pgfpathlineto{\pgfqpoint{3.358396in}{0.360925in}}%
\pgfpathlineto{\pgfqpoint{3.360666in}{0.472556in}}%
\pgfpathlineto{\pgfqpoint{3.362935in}{0.513038in}}%
\pgfpathlineto{\pgfqpoint{3.365204in}{0.500402in}}%
\pgfpathlineto{\pgfqpoint{3.367474in}{0.477244in}}%
\pgfpathlineto{\pgfqpoint{3.369743in}{0.433688in}}%
\pgfpathlineto{\pgfqpoint{3.372013in}{0.497282in}}%
\pgfpathlineto{\pgfqpoint{3.374282in}{0.452748in}}%
\pgfpathlineto{\pgfqpoint{3.376552in}{0.486671in}}%
\pgfpathlineto{\pgfqpoint{3.378821in}{0.467065in}}%
\pgfpathlineto{\pgfqpoint{3.381091in}{0.404950in}}%
\pgfpathlineto{\pgfqpoint{3.383360in}{0.467191in}}%
\pgfpathlineto{\pgfqpoint{3.385630in}{0.511976in}}%
\pgfpathlineto{\pgfqpoint{3.387899in}{0.506453in}}%
\pgfpathlineto{\pgfqpoint{3.392438in}{0.409570in}}%
\pgfpathlineto{\pgfqpoint{3.394708in}{0.493244in}}%
\pgfpathlineto{\pgfqpoint{3.396977in}{0.453091in}}%
\pgfpathlineto{\pgfqpoint{3.399247in}{0.488709in}}%
\pgfpathlineto{\pgfqpoint{3.401516in}{0.471169in}}%
\pgfpathlineto{\pgfqpoint{3.403786in}{0.431046in}}%
\pgfpathlineto{\pgfqpoint{3.406055in}{0.455263in}}%
\pgfpathlineto{\pgfqpoint{3.408325in}{0.510132in}}%
\pgfpathlineto{\pgfqpoint{3.410594in}{0.511596in}}%
\pgfpathlineto{\pgfqpoint{3.412864in}{0.435874in}}%
\pgfpathlineto{\pgfqpoint{3.414905in}{0.315625in}}%
\pgfpathmoveto{\pgfqpoint{3.415297in}{0.315625in}}%
\pgfpathlineto{\pgfqpoint{3.417403in}{0.487965in}}%
\pgfpathlineto{\pgfqpoint{3.419672in}{0.452278in}}%
\pgfpathlineto{\pgfqpoint{3.421942in}{0.490065in}}%
\pgfpathlineto{\pgfqpoint{3.424211in}{0.474808in}}%
\pgfpathlineto{\pgfqpoint{3.426481in}{0.447682in}}%
\pgfpathlineto{\pgfqpoint{3.428750in}{0.437950in}}%
\pgfpathlineto{\pgfqpoint{3.431019in}{0.507787in}}%
\pgfpathlineto{\pgfqpoint{3.433289in}{0.515685in}}%
\pgfpathlineto{\pgfqpoint{3.435558in}{0.372547in}}%
\pgfpathlineto{\pgfqpoint{3.437828in}{0.407415in}}%
\pgfpathlineto{\pgfqpoint{3.440097in}{0.481477in}}%
\pgfpathlineto{\pgfqpoint{3.442367in}{0.449703in}}%
\pgfpathlineto{\pgfqpoint{3.444636in}{0.490477in}}%
\pgfpathlineto{\pgfqpoint{3.446906in}{0.477759in}}%
\pgfpathlineto{\pgfqpoint{3.449175in}{0.460230in}}%
\pgfpathlineto{\pgfqpoint{3.451445in}{0.403949in}}%
\pgfpathlineto{\pgfqpoint{3.453714in}{0.504700in}}%
\pgfpathlineto{\pgfqpoint{3.455984in}{0.518839in}}%
\pgfpathlineto{\pgfqpoint{3.458253in}{0.405822in}}%
\pgfpathlineto{\pgfqpoint{3.462792in}{0.473299in}}%
\pgfpathlineto{\pgfqpoint{3.465062in}{0.445030in}}%
\pgfpathlineto{\pgfqpoint{3.467331in}{0.489735in}}%
\pgfpathlineto{\pgfqpoint{3.471870in}{0.470185in}}%
\pgfpathlineto{\pgfqpoint{3.474140in}{0.352554in}}%
\pgfpathlineto{\pgfqpoint{3.476409in}{0.500716in}}%
\pgfpathlineto{\pgfqpoint{3.478679in}{0.521086in}}%
\pgfpathlineto{\pgfqpoint{3.480948in}{0.494257in}}%
\pgfpathlineto{\pgfqpoint{3.483218in}{0.420472in}}%
\pgfpathlineto{\pgfqpoint{3.485487in}{0.522274in}}%
\pgfpathlineto{\pgfqpoint{3.487757in}{0.487202in}}%
\pgfpathlineto{\pgfqpoint{3.490026in}{0.517260in}}%
\pgfpathlineto{\pgfqpoint{3.492296in}{0.428880in}}%
\pgfpathlineto{\pgfqpoint{3.494565in}{0.473267in}}%
\pgfpathlineto{\pgfqpoint{3.496834in}{0.316305in}}%
\pgfpathlineto{\pgfqpoint{3.499104in}{0.420842in}}%
\pgfpathlineto{\pgfqpoint{3.501373in}{0.481735in}}%
\pgfpathlineto{\pgfqpoint{3.503643in}{0.497518in}}%
\pgfpathlineto{\pgfqpoint{3.505912in}{0.482259in}}%
\pgfpathlineto{\pgfqpoint{3.508182in}{0.400453in}}%
\pgfpathlineto{\pgfqpoint{3.510451in}{0.513440in}}%
\pgfpathlineto{\pgfqpoint{3.512721in}{0.531359in}}%
\pgfpathlineto{\pgfqpoint{3.514990in}{0.475398in}}%
\pgfpathlineto{\pgfqpoint{3.517260in}{0.433242in}}%
\pgfpathlineto{\pgfqpoint{3.519529in}{0.507016in}}%
\pgfpathlineto{\pgfqpoint{3.521799in}{0.447703in}}%
\pgfpathlineto{\pgfqpoint{3.524068in}{0.501223in}}%
\pgfpathlineto{\pgfqpoint{3.526338in}{0.453460in}}%
\pgfpathlineto{\pgfqpoint{3.528607in}{0.481260in}}%
\pgfpathlineto{\pgfqpoint{3.530877in}{0.457017in}}%
\pgfpathlineto{\pgfqpoint{3.533146in}{0.407653in}}%
\pgfpathlineto{\pgfqpoint{3.535416in}{0.488695in}}%
\pgfpathlineto{\pgfqpoint{3.537685in}{0.514230in}}%
\pgfpathlineto{\pgfqpoint{3.539955in}{0.489384in}}%
\pgfpathlineto{\pgfqpoint{3.542224in}{0.493224in}}%
\pgfpathlineto{\pgfqpoint{3.544494in}{0.449526in}}%
\pgfpathlineto{\pgfqpoint{3.546763in}{0.502032in}}%
\pgfpathlineto{\pgfqpoint{3.549033in}{0.450339in}}%
\pgfpathlineto{\pgfqpoint{3.551302in}{0.483768in}}%
\pgfpathlineto{\pgfqpoint{3.553572in}{0.460548in}}%
\pgfpathlineto{\pgfqpoint{3.555841in}{0.336459in}}%
\pgfpathlineto{\pgfqpoint{3.558111in}{0.474615in}}%
\pgfpathlineto{\pgfqpoint{3.560380in}{0.512663in}}%
\pgfpathlineto{\pgfqpoint{3.562649in}{0.499536in}}%
\pgfpathlineto{\pgfqpoint{3.564919in}{0.477226in}}%
\pgfpathlineto{\pgfqpoint{3.567188in}{0.435998in}}%
\pgfpathlineto{\pgfqpoint{3.569458in}{0.497850in}}%
\pgfpathlineto{\pgfqpoint{3.571727in}{0.453806in}}%
\pgfpathlineto{\pgfqpoint{3.573997in}{0.486871in}}%
\pgfpathlineto{\pgfqpoint{3.576266in}{0.466339in}}%
\pgfpathlineto{\pgfqpoint{3.578536in}{0.398541in}}%
\pgfpathlineto{\pgfqpoint{3.580805in}{0.469403in}}%
\pgfpathlineto{\pgfqpoint{3.583075in}{0.511671in}}%
\pgfpathlineto{\pgfqpoint{3.585344in}{0.505785in}}%
\pgfpathlineto{\pgfqpoint{3.589883in}{0.413341in}}%
\pgfpathlineto{\pgfqpoint{3.592153in}{0.493819in}}%
\pgfpathlineto{\pgfqpoint{3.594422in}{0.454212in}}%
\pgfpathlineto{\pgfqpoint{3.596692in}{0.489032in}}%
\pgfpathlineto{\pgfqpoint{3.598961in}{0.470696in}}%
\pgfpathlineto{\pgfqpoint{3.601231in}{0.428105in}}%
\pgfpathlineto{\pgfqpoint{3.603500in}{0.458040in}}%
\pgfpathlineto{\pgfqpoint{3.605770in}{0.509902in}}%
\pgfpathlineto{\pgfqpoint{3.608039in}{0.511073in}}%
\pgfpathlineto{\pgfqpoint{3.610309in}{0.435674in}}%
\pgfpathlineto{\pgfqpoint{3.612578in}{0.333588in}}%
\pgfpathlineto{\pgfqpoint{3.614848in}{0.488553in}}%
\pgfpathlineto{\pgfqpoint{3.617117in}{0.453465in}}%
\pgfpathlineto{\pgfqpoint{3.619387in}{0.490500in}}%
\pgfpathlineto{\pgfqpoint{3.621656in}{0.474523in}}%
\pgfpathlineto{\pgfqpoint{3.623926in}{0.445973in}}%
\pgfpathlineto{\pgfqpoint{3.626195in}{0.441831in}}%
\pgfpathlineto{\pgfqpoint{3.628465in}{0.507617in}}%
\pgfpathlineto{\pgfqpoint{3.630734in}{0.515268in}}%
\pgfpathlineto{\pgfqpoint{3.633003in}{0.371268in}}%
\pgfpathlineto{\pgfqpoint{3.635273in}{0.403772in}}%
\pgfpathlineto{\pgfqpoint{3.637542in}{0.482080in}}%
\pgfpathlineto{\pgfqpoint{3.639812in}{0.450984in}}%
\pgfpathlineto{\pgfqpoint{3.642081in}{0.491025in}}%
\pgfpathlineto{\pgfqpoint{3.644351in}{0.477618in}}%
\pgfpathlineto{\pgfqpoint{3.646620in}{0.459154in}}%
\pgfpathlineto{\pgfqpoint{3.648890in}{0.411675in}}%
\pgfpathlineto{\pgfqpoint{3.651159in}{0.504584in}}%
\pgfpathlineto{\pgfqpoint{3.653429in}{0.518502in}}%
\pgfpathlineto{\pgfqpoint{3.655698in}{0.483846in}}%
\pgfpathlineto{\pgfqpoint{3.657968in}{0.463375in}}%
\pgfpathlineto{\pgfqpoint{3.660237in}{0.528100in}}%
\pgfpathlineto{\pgfqpoint{3.662507in}{0.487768in}}%
\pgfpathlineto{\pgfqpoint{3.664776in}{0.514725in}}%
\pgfpathlineto{\pgfqpoint{3.667046in}{0.399530in}}%
\pgfpathlineto{\pgfqpoint{3.669315in}{0.464076in}}%
\pgfpathlineto{\pgfqpoint{3.671585in}{0.405581in}}%
\pgfpathlineto{\pgfqpoint{3.673854in}{0.436487in}}%
\pgfpathlineto{\pgfqpoint{3.676124in}{0.485774in}}%
\pgfpathlineto{\pgfqpoint{3.678393in}{0.496637in}}%
\pgfpathlineto{\pgfqpoint{3.682932in}{0.450893in}}%
\pgfpathlineto{\pgfqpoint{3.685202in}{0.517264in}}%
\pgfpathlineto{\pgfqpoint{3.687471in}{0.528432in}}%
\pgfpathlineto{\pgfqpoint{3.692010in}{0.378492in}}%
\pgfpathlineto{\pgfqpoint{3.694280in}{0.514504in}}%
\pgfpathlineto{\pgfqpoint{3.696549in}{0.454114in}}%
\pgfpathlineto{\pgfqpoint{3.698818in}{0.503674in}}%
\pgfpathlineto{\pgfqpoint{3.701088in}{0.450070in}}%
\pgfpathlineto{\pgfqpoint{3.703357in}{0.476503in}}%
\pgfpathlineto{\pgfqpoint{3.705627in}{0.446243in}}%
\pgfpathlineto{\pgfqpoint{3.707896in}{0.433994in}}%
\pgfpathlineto{\pgfqpoint{3.710166in}{0.493575in}}%
\pgfpathlineto{\pgfqpoint{3.712435in}{0.513599in}}%
\pgfpathlineto{\pgfqpoint{3.714705in}{0.474718in}}%
\pgfpathlineto{\pgfqpoint{3.716974in}{0.502922in}}%
\pgfpathlineto{\pgfqpoint{3.719244in}{0.458990in}}%
\pgfpathlineto{\pgfqpoint{3.721513in}{0.505365in}}%
\pgfpathlineto{\pgfqpoint{3.723783in}{0.447661in}}%
\pgfpathlineto{\pgfqpoint{3.726052in}{0.480003in}}%
\pgfpathlineto{\pgfqpoint{3.728322in}{0.452016in}}%
\pgfpathlineto{\pgfqpoint{3.730591in}{0.392918in}}%
\pgfpathlineto{\pgfqpoint{3.732861in}{0.482840in}}%
\pgfpathlineto{\pgfqpoint{3.735130in}{0.512874in}}%
\pgfpathlineto{\pgfqpoint{3.737400in}{0.489470in}}%
\pgfpathlineto{\pgfqpoint{3.739669in}{0.490135in}}%
\pgfpathlineto{\pgfqpoint{3.741939in}{0.451420in}}%
\pgfpathlineto{\pgfqpoint{3.744208in}{0.502155in}}%
\pgfpathlineto{\pgfqpoint{3.746478in}{0.452952in}}%
\pgfpathlineto{\pgfqpoint{3.748747in}{0.483858in}}%
\pgfpathlineto{\pgfqpoint{3.751017in}{0.459182in}}%
\pgfpathlineto{\pgfqpoint{3.753286in}{0.331688in}}%
\pgfpathlineto{\pgfqpoint{3.755556in}{0.479699in}}%
\pgfpathlineto{\pgfqpoint{3.757825in}{0.512364in}}%
\pgfpathlineto{\pgfqpoint{3.760095in}{0.497627in}}%
\pgfpathlineto{\pgfqpoint{3.762364in}{0.478855in}}%
\pgfpathlineto{\pgfqpoint{3.764633in}{0.440114in}}%
\pgfpathlineto{\pgfqpoint{3.766903in}{0.499003in}}%
\pgfpathlineto{\pgfqpoint{3.769172in}{0.454897in}}%
\pgfpathlineto{\pgfqpoint{3.771442in}{0.486908in}}%
\pgfpathlineto{\pgfqpoint{3.773711in}{0.464826in}}%
\pgfpathlineto{\pgfqpoint{3.775981in}{0.386076in}}%
\pgfpathlineto{\pgfqpoint{3.778250in}{0.472192in}}%
\pgfpathlineto{\pgfqpoint{3.780520in}{0.511240in}}%
\pgfpathlineto{\pgfqpoint{3.782789in}{0.504561in}}%
\pgfpathlineto{\pgfqpoint{3.787328in}{0.419481in}}%
\pgfpathlineto{\pgfqpoint{3.789598in}{0.494813in}}%
\pgfpathlineto{\pgfqpoint{3.791867in}{0.455785in}}%
\pgfpathlineto{\pgfqpoint{3.794137in}{0.489353in}}%
\pgfpathlineto{\pgfqpoint{3.796406in}{0.469772in}}%
\pgfpathlineto{\pgfqpoint{3.798676in}{0.422125in}}%
\pgfpathlineto{\pgfqpoint{3.803215in}{0.509675in}}%
\pgfpathlineto{\pgfqpoint{3.805484in}{0.510078in}}%
\pgfpathlineto{\pgfqpoint{3.810023in}{0.363838in}}%
\pgfpathlineto{\pgfqpoint{3.812293in}{0.489677in}}%
\pgfpathlineto{\pgfqpoint{3.814562in}{0.455149in}}%
\pgfpathlineto{\pgfqpoint{3.816832in}{0.491009in}}%
\pgfpathlineto{\pgfqpoint{3.819101in}{0.473882in}}%
\pgfpathlineto{\pgfqpoint{3.821371in}{0.442568in}}%
\pgfpathlineto{\pgfqpoint{3.823640in}{0.447647in}}%
\pgfpathlineto{\pgfqpoint{3.825910in}{0.507493in}}%
\pgfpathlineto{\pgfqpoint{3.828179in}{0.514479in}}%
\pgfpathlineto{\pgfqpoint{3.830448in}{0.376516in}}%
\pgfpathlineto{\pgfqpoint{3.832718in}{0.394831in}}%
\pgfpathlineto{\pgfqpoint{3.834987in}{0.483288in}}%
\pgfpathlineto{\pgfqpoint{3.837257in}{0.452904in}}%
\pgfpathlineto{\pgfqpoint{3.839526in}{0.491753in}}%
\pgfpathlineto{\pgfqpoint{3.841796in}{0.477240in}}%
\pgfpathlineto{\pgfqpoint{3.844065in}{0.456925in}}%
\pgfpathlineto{\pgfqpoint{3.846335in}{0.422475in}}%
\pgfpathlineto{\pgfqpoint{3.848604in}{0.504583in}}%
\pgfpathlineto{\pgfqpoint{3.850874in}{0.517875in}}%
\pgfpathlineto{\pgfqpoint{3.853143in}{0.482030in}}%
\pgfpathlineto{\pgfqpoint{3.855413in}{0.463888in}}%
\pgfpathlineto{\pgfqpoint{3.857682in}{0.527664in}}%
\pgfpathlineto{\pgfqpoint{3.859952in}{0.487925in}}%
\pgfpathlineto{\pgfqpoint{3.862221in}{0.514248in}}%
\pgfpathlineto{\pgfqpoint{3.864491in}{0.399557in}}%
\pgfpathlineto{\pgfqpoint{3.866760in}{0.462575in}}%
\pgfpathlineto{\pgfqpoint{3.869030in}{0.412449in}}%
\pgfpathlineto{\pgfqpoint{3.871299in}{0.439662in}}%
\pgfpathlineto{\pgfqpoint{3.873569in}{0.487113in}}%
\pgfpathlineto{\pgfqpoint{3.875838in}{0.496259in}}%
\pgfpathlineto{\pgfqpoint{3.878108in}{0.469304in}}%
\pgfpathlineto{\pgfqpoint{3.880377in}{0.454849in}}%
\pgfpathlineto{\pgfqpoint{3.882647in}{0.516794in}}%
\pgfpathlineto{\pgfqpoint{3.884916in}{0.527708in}}%
\pgfpathlineto{\pgfqpoint{3.889455in}{0.383945in}}%
\pgfpathlineto{\pgfqpoint{3.891725in}{0.514219in}}%
\pgfpathlineto{\pgfqpoint{3.893994in}{0.455724in}}%
\pgfpathlineto{\pgfqpoint{3.896264in}{0.504007in}}%
\pgfpathlineto{\pgfqpoint{3.898533in}{0.450679in}}%
\pgfpathlineto{\pgfqpoint{3.900802in}{0.476081in}}%
\pgfpathlineto{\pgfqpoint{3.903072in}{0.443581in}}%
\pgfpathlineto{\pgfqpoint{3.905341in}{0.437592in}}%
\pgfpathlineto{\pgfqpoint{3.907611in}{0.494596in}}%
\pgfpathlineto{\pgfqpoint{3.909880in}{0.512933in}}%
\pgfpathlineto{\pgfqpoint{3.912150in}{0.472987in}}%
\pgfpathlineto{\pgfqpoint{3.914419in}{0.502844in}}%
\pgfpathlineto{\pgfqpoint{3.916689in}{0.460245in}}%
\pgfpathlineto{\pgfqpoint{3.918958in}{0.505731in}}%
\pgfpathlineto{\pgfqpoint{3.921228in}{0.448464in}}%
\pgfpathlineto{\pgfqpoint{3.923497in}{0.479831in}}%
\pgfpathlineto{\pgfqpoint{3.925767in}{0.450319in}}%
\pgfpathlineto{\pgfqpoint{3.928036in}{0.400993in}}%
\pgfpathlineto{\pgfqpoint{3.930306in}{0.484310in}}%
\pgfpathlineto{\pgfqpoint{3.932575in}{0.512400in}}%
\pgfpathlineto{\pgfqpoint{3.934845in}{0.488351in}}%
\pgfpathlineto{\pgfqpoint{3.937114in}{0.490117in}}%
\pgfpathlineto{\pgfqpoint{3.939384in}{0.452876in}}%
\pgfpathlineto{\pgfqpoint{3.941653in}{0.502642in}}%
\pgfpathlineto{\pgfqpoint{3.943923in}{0.453800in}}%
\pgfpathlineto{\pgfqpoint{3.946192in}{0.483839in}}%
\pgfpathlineto{\pgfqpoint{3.948462in}{0.457995in}}%
\pgfpathlineto{\pgfqpoint{3.950731in}{0.357269in}}%
\pgfpathlineto{\pgfqpoint{3.953001in}{0.481238in}}%
\pgfpathlineto{\pgfqpoint{3.955270in}{0.511974in}}%
\pgfpathlineto{\pgfqpoint{3.959809in}{0.478811in}}%
\pgfpathlineto{\pgfqpoint{3.962079in}{0.441927in}}%
\pgfpathlineto{\pgfqpoint{3.964348in}{0.499503in}}%
\pgfpathlineto{\pgfqpoint{3.966617in}{0.455811in}}%
\pgfpathlineto{\pgfqpoint{3.968887in}{0.487029in}}%
\pgfpathlineto{\pgfqpoint{3.971156in}{0.464023in}}%
\pgfpathlineto{\pgfqpoint{3.973426in}{0.376227in}}%
\pgfpathlineto{\pgfqpoint{3.975695in}{0.473985in}}%
\pgfpathlineto{\pgfqpoint{3.977965in}{0.510935in}}%
\pgfpathlineto{\pgfqpoint{3.980234in}{0.503928in}}%
\pgfpathlineto{\pgfqpoint{3.984773in}{0.422244in}}%
\pgfpathlineto{\pgfqpoint{3.987043in}{0.495330in}}%
\pgfpathlineto{\pgfqpoint{3.989312in}{0.456756in}}%
\pgfpathlineto{\pgfqpoint{3.991582in}{0.489594in}}%
\pgfpathlineto{\pgfqpoint{3.993851in}{0.469242in}}%
\pgfpathlineto{\pgfqpoint{3.996121in}{0.418636in}}%
\pgfpathlineto{\pgfqpoint{4.000660in}{0.509439in}}%
\pgfpathlineto{\pgfqpoint{4.002929in}{0.509584in}}%
\pgfpathlineto{\pgfqpoint{4.007468in}{0.372764in}}%
\pgfpathlineto{\pgfqpoint{4.009738in}{0.490209in}}%
\pgfpathlineto{\pgfqpoint{4.012007in}{0.456186in}}%
\pgfpathlineto{\pgfqpoint{4.014277in}{0.491360in}}%
\pgfpathlineto{\pgfqpoint{4.016546in}{0.473555in}}%
\pgfpathlineto{\pgfqpoint{4.018816in}{0.440711in}}%
\pgfpathlineto{\pgfqpoint{4.021085in}{0.450530in}}%
\pgfpathlineto{\pgfqpoint{4.023355in}{0.507316in}}%
\pgfpathlineto{\pgfqpoint{4.025624in}{0.514085in}}%
\pgfpathlineto{\pgfqpoint{4.027894in}{0.465083in}}%
\pgfpathlineto{\pgfqpoint{4.030163in}{0.483435in}}%
\pgfpathlineto{\pgfqpoint{4.032432in}{0.532286in}}%
\pgfpathlineto{\pgfqpoint{4.034702in}{0.486527in}}%
\pgfpathlineto{\pgfqpoint{4.036971in}{0.510248in}}%
\pgfpathlineto{\pgfqpoint{4.039241in}{0.330093in}}%
\pgfpathlineto{\pgfqpoint{4.041510in}{0.450104in}}%
\pgfpathlineto{\pgfqpoint{4.043780in}{0.436598in}}%
\pgfpathlineto{\pgfqpoint{4.046049in}{0.448850in}}%
\pgfpathlineto{\pgfqpoint{4.048319in}{0.489313in}}%
\pgfpathlineto{\pgfqpoint{4.050588in}{0.494510in}}%
\pgfpathlineto{\pgfqpoint{4.052858in}{0.452960in}}%
\pgfpathlineto{\pgfqpoint{4.055127in}{0.475148in}}%
\pgfpathlineto{\pgfqpoint{4.057397in}{0.519484in}}%
\pgfpathlineto{\pgfqpoint{4.059666in}{0.523447in}}%
\pgfpathlineto{\pgfqpoint{4.061936in}{0.390119in}}%
\pgfpathlineto{\pgfqpoint{4.066475in}{0.520482in}}%
\pgfpathlineto{\pgfqpoint{4.068744in}{0.458446in}}%
\pgfpathlineto{\pgfqpoint{4.071014in}{0.504418in}}%
\pgfpathlineto{\pgfqpoint{4.073283in}{0.444828in}}%
\pgfpathlineto{\pgfqpoint{4.075553in}{0.469805in}}%
\pgfpathlineto{\pgfqpoint{4.077822in}{0.426261in}}%
\pgfpathlineto{\pgfqpoint{4.080092in}{0.451339in}}%
\pgfpathlineto{\pgfqpoint{4.082361in}{0.497371in}}%
\pgfpathlineto{\pgfqpoint{4.084631in}{0.511469in}}%
\pgfpathlineto{\pgfqpoint{4.086900in}{0.448537in}}%
\pgfpathlineto{\pgfqpoint{4.089170in}{0.510707in}}%
\pgfpathlineto{\pgfqpoint{4.091439in}{0.465627in}}%
\pgfpathlineto{\pgfqpoint{4.093709in}{0.506998in}}%
\pgfpathlineto{\pgfqpoint{4.095978in}{0.443244in}}%
\pgfpathlineto{\pgfqpoint{4.098247in}{0.474667in}}%
\pgfpathlineto{\pgfqpoint{4.100517in}{0.437572in}}%
\pgfpathlineto{\pgfqpoint{4.102786in}{0.429489in}}%
\pgfpathlineto{\pgfqpoint{4.105056in}{0.489683in}}%
\pgfpathlineto{\pgfqpoint{4.107325in}{0.511817in}}%
\pgfpathlineto{\pgfqpoint{4.109595in}{0.474025in}}%
\pgfpathlineto{\pgfqpoint{4.111864in}{0.500152in}}%
\pgfpathlineto{\pgfqpoint{4.114134in}{0.461713in}}%
\pgfpathlineto{\pgfqpoint{4.116403in}{0.505642in}}%
\pgfpathlineto{\pgfqpoint{4.118673in}{0.450842in}}%
\pgfpathlineto{\pgfqpoint{4.120942in}{0.479458in}}%
\pgfpathlineto{\pgfqpoint{4.125481in}{0.415897in}}%
\pgfpathlineto{\pgfqpoint{4.127751in}{0.488073in}}%
\pgfpathlineto{\pgfqpoint{4.130020in}{0.511842in}}%
\pgfpathlineto{\pgfqpoint{4.132290in}{0.485703in}}%
\pgfpathlineto{\pgfqpoint{4.134559in}{0.491382in}}%
\pgfpathlineto{\pgfqpoint{4.136829in}{0.455404in}}%
\pgfpathlineto{\pgfqpoint{4.139098in}{0.503611in}}%
\pgfpathlineto{\pgfqpoint{4.141368in}{0.454490in}}%
\pgfpathlineto{\pgfqpoint{4.143637in}{0.483532in}}%
\pgfpathlineto{\pgfqpoint{4.145907in}{0.455711in}}%
\pgfpathlineto{\pgfqpoint{4.148176in}{0.379080in}}%
\pgfpathlineto{\pgfqpoint{4.150446in}{0.483137in}}%
\pgfpathlineto{\pgfqpoint{4.152715in}{0.511406in}}%
\pgfpathlineto{\pgfqpoint{4.157254in}{0.479239in}}%
\pgfpathlineto{\pgfqpoint{4.159524in}{0.444934in}}%
\pgfpathlineto{\pgfqpoint{4.161793in}{0.500338in}}%
\pgfpathlineto{\pgfqpoint{4.164062in}{0.457018in}}%
\pgfpathlineto{\pgfqpoint{4.166332in}{0.487014in}}%
\pgfpathlineto{\pgfqpoint{4.168601in}{0.462526in}}%
\pgfpathlineto{\pgfqpoint{4.170871in}{0.349667in}}%
\pgfpathlineto{\pgfqpoint{4.173140in}{0.476833in}}%
\pgfpathlineto{\pgfqpoint{4.175410in}{0.510561in}}%
\pgfpathlineto{\pgfqpoint{4.177679in}{0.502698in}}%
\pgfpathlineto{\pgfqpoint{4.182218in}{0.427079in}}%
\pgfpathlineto{\pgfqpoint{4.184488in}{0.496282in}}%
\pgfpathlineto{\pgfqpoint{4.186757in}{0.458051in}}%
\pgfpathlineto{\pgfqpoint{4.189027in}{0.489776in}}%
\pgfpathlineto{\pgfqpoint{4.191296in}{0.468178in}}%
\pgfpathlineto{\pgfqpoint{4.193566in}{0.411507in}}%
\pgfpathlineto{\pgfqpoint{4.195835in}{0.467899in}}%
\pgfpathlineto{\pgfqpoint{4.198105in}{0.509178in}}%
\pgfpathlineto{\pgfqpoint{4.200374in}{0.508624in}}%
\pgfpathlineto{\pgfqpoint{4.202644in}{0.438366in}}%
\pgfpathlineto{\pgfqpoint{4.204913in}{0.386547in}}%
\pgfpathlineto{\pgfqpoint{4.207183in}{0.491229in}}%
\pgfpathlineto{\pgfqpoint{4.209452in}{0.457663in}}%
\pgfpathlineto{\pgfqpoint{4.211722in}{0.491747in}}%
\pgfpathlineto{\pgfqpoint{4.213991in}{0.472847in}}%
\pgfpathlineto{\pgfqpoint{4.216261in}{0.436938in}}%
\pgfpathlineto{\pgfqpoint{4.218530in}{0.455062in}}%
\pgfpathlineto{\pgfqpoint{4.220800in}{0.507177in}}%
\pgfpathlineto{\pgfqpoint{4.223069in}{0.513321in}}%
\pgfpathlineto{\pgfqpoint{4.225339in}{0.462316in}}%
\pgfpathlineto{\pgfqpoint{4.227608in}{0.483481in}}%
\pgfpathlineto{\pgfqpoint{4.229878in}{0.531806in}}%
\pgfpathlineto{\pgfqpoint{4.232147in}{0.486618in}}%
\pgfpathlineto{\pgfqpoint{4.234416in}{0.509640in}}%
\pgfpathlineto{\pgfqpoint{4.236686in}{0.334741in}}%
\pgfpathlineto{\pgfqpoint{4.238955in}{0.447761in}}%
\pgfpathlineto{\pgfqpoint{4.241225in}{0.440398in}}%
\pgfpathlineto{\pgfqpoint{4.243494in}{0.451098in}}%
\pgfpathlineto{\pgfqpoint{4.245764in}{0.490231in}}%
\pgfpathlineto{\pgfqpoint{4.248033in}{0.493935in}}%
\pgfpathlineto{\pgfqpoint{4.250303in}{0.449951in}}%
\pgfpathlineto{\pgfqpoint{4.252572in}{0.477229in}}%
\pgfpathlineto{\pgfqpoint{4.254842in}{0.518936in}}%
\pgfpathlineto{\pgfqpoint{4.257111in}{0.522650in}}%
\pgfpathlineto{\pgfqpoint{4.259381in}{0.386983in}}%
\pgfpathlineto{\pgfqpoint{4.263920in}{0.520185in}}%
\pgfpathlineto{\pgfqpoint{4.266189in}{0.459786in}}%
\pgfpathlineto{\pgfqpoint{4.268459in}{0.504520in}}%
\pgfpathlineto{\pgfqpoint{4.270728in}{0.445132in}}%
\pgfpathlineto{\pgfqpoint{4.272998in}{0.469032in}}%
\pgfpathlineto{\pgfqpoint{4.275267in}{0.421162in}}%
\pgfpathlineto{\pgfqpoint{4.277537in}{0.453666in}}%
\pgfpathlineto{\pgfqpoint{4.279806in}{0.498061in}}%
\pgfpathlineto{\pgfqpoint{4.282076in}{0.510663in}}%
\pgfpathlineto{\pgfqpoint{4.284345in}{0.445482in}}%
\pgfpathlineto{\pgfqpoint{4.286615in}{0.510607in}}%
\pgfpathlineto{\pgfqpoint{4.288884in}{0.466594in}}%
\pgfpathlineto{\pgfqpoint{4.291154in}{0.507146in}}%
\pgfpathlineto{\pgfqpoint{4.293423in}{0.443790in}}%
\pgfpathlineto{\pgfqpoint{4.295693in}{0.474215in}}%
\pgfpathlineto{\pgfqpoint{4.297962in}{0.434635in}}%
\pgfpathlineto{\pgfqpoint{4.300231in}{0.433035in}}%
\pgfpathlineto{\pgfqpoint{4.302501in}{0.490750in}}%
\pgfpathlineto{\pgfqpoint{4.304770in}{0.511234in}}%
\pgfpathlineto{\pgfqpoint{4.307040in}{0.472460in}}%
\pgfpathlineto{\pgfqpoint{4.309309in}{0.500115in}}%
\pgfpathlineto{\pgfqpoint{4.311579in}{0.462765in}}%
\pgfpathlineto{\pgfqpoint{4.313848in}{0.505939in}}%
\pgfpathlineto{\pgfqpoint{4.316118in}{0.451481in}}%
\pgfpathlineto{\pgfqpoint{4.318387in}{0.479206in}}%
\pgfpathlineto{\pgfqpoint{4.322926in}{0.420432in}}%
\pgfpathlineto{\pgfqpoint{4.325196in}{0.489191in}}%
\pgfpathlineto{\pgfqpoint{4.327465in}{0.511359in}}%
\pgfpathlineto{\pgfqpoint{4.329735in}{0.484626in}}%
\pgfpathlineto{\pgfqpoint{4.332004in}{0.491273in}}%
\pgfpathlineto{\pgfqpoint{4.334274in}{0.457220in}}%
\pgfpathlineto{\pgfqpoint{4.336543in}{0.504357in}}%
\pgfpathlineto{\pgfqpoint{4.336543in}{0.504357in}}%
\pgfusepath{stroke}%
\end{pgfscope}%
\begin{pgfscope}%
\pgfpathrectangle{\pgfqpoint{2.520957in}{0.320625in}}{\pgfqpoint{1.815586in}{1.164375in}}%
\pgfusepath{clip}%
\pgfsetrectcap%
\pgfsetroundjoin%
\pgfsetlinewidth{1.154312pt}%
\definecolor{currentstroke}{rgb}{0.950210,0.390820,0.362468}%
\pgfsetstrokecolor{currentstroke}%
\pgfsetdash{}{0pt}%
\pgfpathmoveto{\pgfqpoint{2.520957in}{1.451682in}}%
\pgfpathlineto{\pgfqpoint{2.525496in}{1.441846in}}%
\pgfpathlineto{\pgfqpoint{2.541382in}{1.410652in}}%
\pgfpathlineto{\pgfqpoint{2.557269in}{1.376637in}}%
\pgfpathlineto{\pgfqpoint{2.591311in}{1.300023in}}%
\pgfpathlineto{\pgfqpoint{2.632161in}{1.209198in}}%
\pgfpathlineto{\pgfqpoint{2.668473in}{1.131271in}}%
\pgfpathlineto{\pgfqpoint{2.707054in}{1.051205in}}%
\pgfpathlineto{\pgfqpoint{2.750175in}{0.964445in}}%
\pgfpathlineto{\pgfqpoint{2.800103in}{0.866675in}}%
\pgfpathlineto{\pgfqpoint{2.859110in}{0.753666in}}%
\pgfpathlineto{\pgfqpoint{2.936272in}{0.608423in}}%
\pgfpathlineto{\pgfqpoint{2.956698in}{0.570455in}}%
\pgfpathlineto{\pgfqpoint{2.963506in}{0.557731in}}%
\pgfpathlineto{\pgfqpoint{2.965775in}{0.554270in}}%
\pgfpathlineto{\pgfqpoint{2.972584in}{0.541158in}}%
\pgfpathlineto{\pgfqpoint{2.974853in}{0.542040in}}%
\pgfpathlineto{\pgfqpoint{2.977123in}{0.538351in}}%
\pgfpathlineto{\pgfqpoint{2.979392in}{0.528565in}}%
\pgfpathlineto{\pgfqpoint{2.981662in}{0.523820in}}%
\pgfpathlineto{\pgfqpoint{2.983931in}{0.524336in}}%
\pgfpathlineto{\pgfqpoint{2.986201in}{0.516386in}}%
\pgfpathlineto{\pgfqpoint{2.988470in}{0.517517in}}%
\pgfpathlineto{\pgfqpoint{2.993009in}{0.503762in}}%
\pgfpathlineto{\pgfqpoint{2.995279in}{0.500240in}}%
\pgfpathlineto{\pgfqpoint{2.997548in}{0.517661in}}%
\pgfpathlineto{\pgfqpoint{2.999818in}{0.519654in}}%
\pgfpathlineto{\pgfqpoint{3.002087in}{0.486666in}}%
\pgfpathlineto{\pgfqpoint{3.004357in}{0.481741in}}%
\pgfpathlineto{\pgfqpoint{3.006626in}{0.496199in}}%
\pgfpathlineto{\pgfqpoint{3.008896in}{0.479487in}}%
\pgfpathlineto{\pgfqpoint{3.011165in}{0.496744in}}%
\pgfpathlineto{\pgfqpoint{3.013435in}{0.485875in}}%
\pgfpathlineto{\pgfqpoint{3.017974in}{0.460695in}}%
\pgfpathlineto{\pgfqpoint{3.020243in}{0.508178in}}%
\pgfpathlineto{\pgfqpoint{3.022513in}{0.517999in}}%
\pgfpathlineto{\pgfqpoint{3.024782in}{0.443847in}}%
\pgfpathlineto{\pgfqpoint{3.027052in}{0.447912in}}%
\pgfpathlineto{\pgfqpoint{3.029321in}{0.480792in}}%
\pgfpathlineto{\pgfqpoint{3.031590in}{0.455599in}}%
\pgfpathlineto{\pgfqpoint{3.033860in}{0.491165in}}%
\pgfpathlineto{\pgfqpoint{3.038399in}{0.466886in}}%
\pgfpathlineto{\pgfqpoint{3.040668in}{0.416746in}}%
\pgfpathlineto{\pgfqpoint{3.042938in}{0.503554in}}%
\pgfpathlineto{\pgfqpoint{3.045207in}{0.519982in}}%
\pgfpathlineto{\pgfqpoint{3.047477in}{0.429618in}}%
\pgfpathlineto{\pgfqpoint{3.052016in}{0.470071in}}%
\pgfpathlineto{\pgfqpoint{3.054285in}{0.443863in}}%
\pgfpathlineto{\pgfqpoint{3.056555in}{0.489073in}}%
\pgfpathlineto{\pgfqpoint{3.061094in}{0.473654in}}%
\pgfpathlineto{\pgfqpoint{3.063363in}{0.401183in}}%
\pgfpathlineto{\pgfqpoint{3.065633in}{0.499080in}}%
\pgfpathlineto{\pgfqpoint{3.067902in}{0.521719in}}%
\pgfpathlineto{\pgfqpoint{3.070172in}{0.497010in}}%
\pgfpathlineto{\pgfqpoint{3.072441in}{0.384598in}}%
\pgfpathlineto{\pgfqpoint{3.074711in}{0.519897in}}%
\pgfpathlineto{\pgfqpoint{3.076980in}{0.486620in}}%
\pgfpathlineto{\pgfqpoint{3.079250in}{0.517792in}}%
\pgfpathlineto{\pgfqpoint{3.081519in}{0.435094in}}%
\pgfpathlineto{\pgfqpoint{3.083789in}{0.475844in}}%
\pgfpathlineto{\pgfqpoint{3.086058in}{0.342935in}}%
\pgfpathlineto{\pgfqpoint{3.090597in}{0.480290in}}%
\pgfpathlineto{\pgfqpoint{3.092867in}{0.497630in}}%
\pgfpathlineto{\pgfqpoint{3.095136in}{0.485087in}}%
\pgfpathlineto{\pgfqpoint{3.097405in}{0.353233in}}%
\pgfpathlineto{\pgfqpoint{3.099675in}{0.511908in}}%
\pgfpathlineto{\pgfqpoint{3.101944in}{0.531941in}}%
\pgfpathlineto{\pgfqpoint{3.104214in}{0.479805in}}%
\pgfpathlineto{\pgfqpoint{3.106483in}{0.447196in}}%
\pgfpathlineto{\pgfqpoint{3.108753in}{0.504377in}}%
\pgfpathlineto{\pgfqpoint{3.111022in}{0.445209in}}%
\pgfpathlineto{\pgfqpoint{3.113292in}{0.500249in}}%
\pgfpathlineto{\pgfqpoint{3.115561in}{0.454257in}}%
\pgfpathlineto{\pgfqpoint{3.117831in}{0.482433in}}%
\pgfpathlineto{\pgfqpoint{3.120100in}{0.459336in}}%
\pgfpathlineto{\pgfqpoint{3.122370in}{0.396280in}}%
\pgfpathlineto{\pgfqpoint{3.124639in}{0.487162in}}%
\pgfpathlineto{\pgfqpoint{3.126909in}{0.514182in}}%
\pgfpathlineto{\pgfqpoint{3.129178in}{0.492459in}}%
\pgfpathlineto{\pgfqpoint{3.131448in}{0.490081in}}%
\pgfpathlineto{\pgfqpoint{3.133717in}{0.446288in}}%
\pgfpathlineto{\pgfqpoint{3.135987in}{0.501032in}}%
\pgfpathlineto{\pgfqpoint{3.138256in}{0.450920in}}%
\pgfpathlineto{\pgfqpoint{3.140526in}{0.484603in}}%
\pgfpathlineto{\pgfqpoint{3.142795in}{0.462225in}}%
\pgfpathlineto{\pgfqpoint{3.145065in}{0.368801in}}%
\pgfpathlineto{\pgfqpoint{3.147334in}{0.472390in}}%
\pgfpathlineto{\pgfqpoint{3.149604in}{0.512415in}}%
\pgfpathlineto{\pgfqpoint{3.151873in}{0.501489in}}%
\pgfpathlineto{\pgfqpoint{3.154143in}{0.473341in}}%
\pgfpathlineto{\pgfqpoint{3.156412in}{0.431051in}}%
\pgfpathlineto{\pgfqpoint{3.158682in}{0.496770in}}%
\pgfpathlineto{\pgfqpoint{3.160951in}{0.453974in}}%
\pgfpathlineto{\pgfqpoint{3.163220in}{0.487467in}}%
\pgfpathlineto{\pgfqpoint{3.165490in}{0.467591in}}%
\pgfpathlineto{\pgfqpoint{3.167759in}{0.408013in}}%
\pgfpathlineto{\pgfqpoint{3.170029in}{0.466954in}}%
\pgfpathlineto{\pgfqpoint{3.172298in}{0.511331in}}%
\pgfpathlineto{\pgfqpoint{3.174568in}{0.507182in}}%
\pgfpathlineto{\pgfqpoint{3.179107in}{0.404427in}}%
\pgfpathlineto{\pgfqpoint{3.181376in}{0.492688in}}%
\pgfpathlineto{\pgfqpoint{3.183646in}{0.454105in}}%
\pgfpathlineto{\pgfqpoint{3.185915in}{0.489405in}}%
\pgfpathlineto{\pgfqpoint{3.188185in}{0.471597in}}%
\pgfpathlineto{\pgfqpoint{3.190454in}{0.432629in}}%
\pgfpathlineto{\pgfqpoint{3.192724in}{0.455079in}}%
\pgfpathlineto{\pgfqpoint{3.194993in}{0.509470in}}%
\pgfpathlineto{\pgfqpoint{3.197263in}{0.512022in}}%
\pgfpathlineto{\pgfqpoint{3.199532in}{0.428062in}}%
\pgfpathlineto{\pgfqpoint{3.201128in}{0.315625in}}%
\pgfpathmoveto{\pgfqpoint{3.202294in}{0.315625in}}%
\pgfpathlineto{\pgfqpoint{3.204071in}{0.487382in}}%
\pgfpathlineto{\pgfqpoint{3.206341in}{0.453142in}}%
\pgfpathlineto{\pgfqpoint{3.208610in}{0.490686in}}%
\pgfpathlineto{\pgfqpoint{3.210880in}{0.475148in}}%
\pgfpathlineto{\pgfqpoint{3.213149in}{0.448635in}}%
\pgfpathlineto{\pgfqpoint{3.215419in}{0.437994in}}%
\pgfpathlineto{\pgfqpoint{3.217688in}{0.507115in}}%
\pgfpathlineto{\pgfqpoint{3.219958in}{0.515879in}}%
\pgfpathlineto{\pgfqpoint{3.222227in}{0.332381in}}%
\pgfpathlineto{\pgfqpoint{3.226766in}{0.480885in}}%
\pgfpathlineto{\pgfqpoint{3.229035in}{0.450502in}}%
\pgfpathlineto{\pgfqpoint{3.231305in}{0.491059in}}%
\pgfpathlineto{\pgfqpoint{3.233574in}{0.478031in}}%
\pgfpathlineto{\pgfqpoint{3.235844in}{0.460807in}}%
\pgfpathlineto{\pgfqpoint{3.238113in}{0.405050in}}%
\pgfpathlineto{\pgfqpoint{3.240383in}{0.504036in}}%
\pgfpathlineto{\pgfqpoint{3.242652in}{0.518851in}}%
\pgfpathlineto{\pgfqpoint{3.247191in}{0.459281in}}%
\pgfpathlineto{\pgfqpoint{3.249461in}{0.527209in}}%
\pgfpathlineto{\pgfqpoint{3.251730in}{0.487834in}}%
\pgfpathlineto{\pgfqpoint{3.254000in}{0.515066in}}%
\pgfpathlineto{\pgfqpoint{3.256269in}{0.405519in}}%
\pgfpathlineto{\pgfqpoint{3.258539in}{0.465160in}}%
\pgfpathlineto{\pgfqpoint{3.260808in}{0.402018in}}%
\pgfpathlineto{\pgfqpoint{3.263078in}{0.435603in}}%
\pgfpathlineto{\pgfqpoint{3.265347in}{0.485655in}}%
\pgfpathlineto{\pgfqpoint{3.267617in}{0.496700in}}%
\pgfpathlineto{\pgfqpoint{3.272156in}{0.448271in}}%
\pgfpathlineto{\pgfqpoint{3.274425in}{0.516754in}}%
\pgfpathlineto{\pgfqpoint{3.276695in}{0.528650in}}%
\pgfpathlineto{\pgfqpoint{3.278964in}{0.456847in}}%
\pgfpathlineto{\pgfqpoint{3.281234in}{0.351960in}}%
\pgfpathlineto{\pgfqpoint{3.283503in}{0.513690in}}%
\pgfpathlineto{\pgfqpoint{3.285773in}{0.454050in}}%
\pgfpathlineto{\pgfqpoint{3.288042in}{0.503608in}}%
\pgfpathlineto{\pgfqpoint{3.290312in}{0.450655in}}%
\pgfpathlineto{\pgfqpoint{3.292581in}{0.476901in}}%
\pgfpathlineto{\pgfqpoint{3.294851in}{0.446729in}}%
\pgfpathlineto{\pgfqpoint{3.297120in}{0.433316in}}%
\pgfpathlineto{\pgfqpoint{3.299389in}{0.493513in}}%
\pgfpathlineto{\pgfqpoint{3.301659in}{0.513493in}}%
\pgfpathlineto{\pgfqpoint{3.303928in}{0.475665in}}%
\pgfpathlineto{\pgfqpoint{3.306198in}{0.502207in}}%
\pgfpathlineto{\pgfqpoint{3.308467in}{0.458794in}}%
\pgfpathlineto{\pgfqpoint{3.310737in}{0.505291in}}%
\pgfpathlineto{\pgfqpoint{3.313006in}{0.448155in}}%
\pgfpathlineto{\pgfqpoint{3.315276in}{0.480243in}}%
\pgfpathlineto{\pgfqpoint{3.317545in}{0.452162in}}%
\pgfpathlineto{\pgfqpoint{3.319815in}{0.392416in}}%
\pgfpathlineto{\pgfqpoint{3.322084in}{0.482891in}}%
\pgfpathlineto{\pgfqpoint{3.324354in}{0.512733in}}%
\pgfpathlineto{\pgfqpoint{3.326623in}{0.489737in}}%
\pgfpathlineto{\pgfqpoint{3.328893in}{0.489556in}}%
\pgfpathlineto{\pgfqpoint{3.331162in}{0.451330in}}%
\pgfpathlineto{\pgfqpoint{3.333432in}{0.502135in}}%
\pgfpathlineto{\pgfqpoint{3.335701in}{0.453297in}}%
\pgfpathlineto{\pgfqpoint{3.337971in}{0.483991in}}%
\pgfpathlineto{\pgfqpoint{3.340240in}{0.459099in}}%
\pgfpathlineto{\pgfqpoint{3.342510in}{0.335099in}}%
\pgfpathlineto{\pgfqpoint{3.344779in}{0.479958in}}%
\pgfpathlineto{\pgfqpoint{3.347049in}{0.512221in}}%
\pgfpathlineto{\pgfqpoint{3.349318in}{0.497573in}}%
\pgfpathlineto{\pgfqpoint{3.351588in}{0.478519in}}%
\pgfpathlineto{\pgfqpoint{3.353857in}{0.440398in}}%
\pgfpathlineto{\pgfqpoint{3.356127in}{0.499097in}}%
\pgfpathlineto{\pgfqpoint{3.358396in}{0.455218in}}%
\pgfpathlineto{\pgfqpoint{3.360666in}{0.486977in}}%
\pgfpathlineto{\pgfqpoint{3.362935in}{0.464594in}}%
\pgfpathlineto{\pgfqpoint{3.365204in}{0.383177in}}%
\pgfpathlineto{\pgfqpoint{3.367474in}{0.472822in}}%
\pgfpathlineto{\pgfqpoint{3.369743in}{0.511126in}}%
\pgfpathlineto{\pgfqpoint{3.372013in}{0.504297in}}%
\pgfpathlineto{\pgfqpoint{3.376552in}{0.420840in}}%
\pgfpathlineto{\pgfqpoint{3.378821in}{0.495070in}}%
\pgfpathlineto{\pgfqpoint{3.381091in}{0.456169in}}%
\pgfpathlineto{\pgfqpoint{3.383360in}{0.489409in}}%
\pgfpathlineto{\pgfqpoint{3.385630in}{0.469451in}}%
\pgfpathlineto{\pgfqpoint{3.387899in}{0.420031in}}%
\pgfpathlineto{\pgfqpoint{3.392438in}{0.509624in}}%
\pgfpathlineto{\pgfqpoint{3.394708in}{0.509686in}}%
\pgfpathlineto{\pgfqpoint{3.399247in}{0.372336in}}%
\pgfpathlineto{\pgfqpoint{3.401516in}{0.490161in}}%
\pgfpathlineto{\pgfqpoint{3.403786in}{0.455695in}}%
\pgfpathlineto{\pgfqpoint{3.406055in}{0.491111in}}%
\pgfpathlineto{\pgfqpoint{3.408325in}{0.473524in}}%
\pgfpathlineto{\pgfqpoint{3.410594in}{0.440719in}}%
\pgfpathlineto{\pgfqpoint{3.412864in}{0.450114in}}%
\pgfpathlineto{\pgfqpoint{3.415133in}{0.507548in}}%
\pgfpathlineto{\pgfqpoint{3.417403in}{0.514021in}}%
\pgfpathlineto{\pgfqpoint{3.419672in}{0.383896in}}%
\pgfpathlineto{\pgfqpoint{3.421942in}{0.387383in}}%
\pgfpathlineto{\pgfqpoint{3.424211in}{0.484086in}}%
\pgfpathlineto{\pgfqpoint{3.426481in}{0.453735in}}%
\pgfpathlineto{\pgfqpoint{3.428750in}{0.491975in}}%
\pgfpathlineto{\pgfqpoint{3.431019in}{0.476889in}}%
\pgfpathlineto{\pgfqpoint{3.433289in}{0.455210in}}%
\pgfpathlineto{\pgfqpoint{3.435558in}{0.428069in}}%
\pgfpathlineto{\pgfqpoint{3.437828in}{0.504799in}}%
\pgfpathlineto{\pgfqpoint{3.440097in}{0.517405in}}%
\pgfpathlineto{\pgfqpoint{3.442367in}{0.480326in}}%
\pgfpathlineto{\pgfqpoint{3.444636in}{0.466058in}}%
\pgfpathlineto{\pgfqpoint{3.446906in}{0.527940in}}%
\pgfpathlineto{\pgfqpoint{3.449175in}{0.487927in}}%
\pgfpathlineto{\pgfqpoint{3.451445in}{0.513766in}}%
\pgfpathlineto{\pgfqpoint{3.453714in}{0.394895in}}%
\pgfpathlineto{\pgfqpoint{3.455984in}{0.460939in}}%
\pgfpathlineto{\pgfqpoint{3.458253in}{0.417664in}}%
\pgfpathlineto{\pgfqpoint{3.460523in}{0.441832in}}%
\pgfpathlineto{\pgfqpoint{3.462792in}{0.487877in}}%
\pgfpathlineto{\pgfqpoint{3.465062in}{0.495947in}}%
\pgfpathlineto{\pgfqpoint{3.467331in}{0.467043in}}%
\pgfpathlineto{\pgfqpoint{3.469601in}{0.458995in}}%
\pgfpathlineto{\pgfqpoint{3.471870in}{0.516965in}}%
\pgfpathlineto{\pgfqpoint{3.474140in}{0.526988in}}%
\pgfpathlineto{\pgfqpoint{3.476409in}{0.448552in}}%
\pgfpathlineto{\pgfqpoint{3.478679in}{0.402008in}}%
\pgfpathlineto{\pgfqpoint{3.480948in}{0.514957in}}%
\pgfpathlineto{\pgfqpoint{3.483218in}{0.456845in}}%
\pgfpathlineto{\pgfqpoint{3.485487in}{0.504263in}}%
\pgfpathlineto{\pgfqpoint{3.487757in}{0.450248in}}%
\pgfpathlineto{\pgfqpoint{3.490026in}{0.475091in}}%
\pgfpathlineto{\pgfqpoint{3.492296in}{0.440329in}}%
\pgfpathlineto{\pgfqpoint{3.494565in}{0.441263in}}%
\pgfpathlineto{\pgfqpoint{3.496834in}{0.495494in}}%
\pgfpathlineto{\pgfqpoint{3.499104in}{0.512467in}}%
\pgfpathlineto{\pgfqpoint{3.501373in}{0.469296in}}%
\pgfpathlineto{\pgfqpoint{3.503643in}{0.504124in}}%
\pgfpathlineto{\pgfqpoint{3.505912in}{0.461793in}}%
\pgfpathlineto{\pgfqpoint{3.508182in}{0.506187in}}%
\pgfpathlineto{\pgfqpoint{3.510451in}{0.448104in}}%
\pgfpathlineto{\pgfqpoint{3.512721in}{0.478951in}}%
\pgfpathlineto{\pgfqpoint{3.517260in}{0.410341in}}%
\pgfpathlineto{\pgfqpoint{3.519529in}{0.485994in}}%
\pgfpathlineto{\pgfqpoint{3.521799in}{0.512120in}}%
\pgfpathlineto{\pgfqpoint{3.524068in}{0.485596in}}%
\pgfpathlineto{\pgfqpoint{3.526338in}{0.492115in}}%
\pgfpathlineto{\pgfqpoint{3.528607in}{0.455456in}}%
\pgfpathlineto{\pgfqpoint{3.530877in}{0.503544in}}%
\pgfpathlineto{\pgfqpoint{3.533146in}{0.453750in}}%
\pgfpathlineto{\pgfqpoint{3.535416in}{0.483040in}}%
\pgfpathlineto{\pgfqpoint{3.537685in}{0.455581in}}%
\pgfpathlineto{\pgfqpoint{3.539955in}{0.382850in}}%
\pgfpathlineto{\pgfqpoint{3.542224in}{0.483482in}}%
\pgfpathlineto{\pgfqpoint{3.544494in}{0.511806in}}%
\pgfpathlineto{\pgfqpoint{3.546763in}{0.494332in}}%
\pgfpathlineto{\pgfqpoint{3.549033in}{0.481736in}}%
\pgfpathlineto{\pgfqpoint{3.551302in}{0.446209in}}%
\pgfpathlineto{\pgfqpoint{3.553572in}{0.500702in}}%
\pgfpathlineto{\pgfqpoint{3.555841in}{0.456095in}}%
\pgfpathlineto{\pgfqpoint{3.558111in}{0.486381in}}%
\pgfpathlineto{\pgfqpoint{3.560380in}{0.461903in}}%
\pgfpathlineto{\pgfqpoint{3.562649in}{0.335998in}}%
\pgfpathlineto{\pgfqpoint{3.564919in}{0.477196in}}%
\pgfpathlineto{\pgfqpoint{3.567188in}{0.510945in}}%
\pgfpathlineto{\pgfqpoint{3.569458in}{0.501749in}}%
\pgfpathlineto{\pgfqpoint{3.573997in}{0.430440in}}%
\pgfpathlineto{\pgfqpoint{3.576266in}{0.496908in}}%
\pgfpathlineto{\pgfqpoint{3.578536in}{0.457458in}}%
\pgfpathlineto{\pgfqpoint{3.580805in}{0.489154in}}%
\pgfpathlineto{\pgfqpoint{3.583075in}{0.467394in}}%
\pgfpathlineto{\pgfqpoint{3.585344in}{0.406173in}}%
\pgfpathlineto{\pgfqpoint{3.587614in}{0.469205in}}%
\pgfpathlineto{\pgfqpoint{3.589883in}{0.509672in}}%
\pgfpathlineto{\pgfqpoint{3.592153in}{0.507651in}}%
\pgfpathlineto{\pgfqpoint{3.594422in}{0.445450in}}%
\pgfpathlineto{\pgfqpoint{3.596692in}{0.398527in}}%
\pgfpathlineto{\pgfqpoint{3.598961in}{0.492271in}}%
\pgfpathlineto{\pgfqpoint{3.601231in}{0.457419in}}%
\pgfpathlineto{\pgfqpoint{3.603500in}{0.491207in}}%
\pgfpathlineto{\pgfqpoint{3.605770in}{0.471977in}}%
\pgfpathlineto{\pgfqpoint{3.608039in}{0.432890in}}%
\pgfpathlineto{\pgfqpoint{3.610309in}{0.457729in}}%
\pgfpathlineto{\pgfqpoint{3.612578in}{0.507836in}}%
\pgfpathlineto{\pgfqpoint{3.614848in}{0.512399in}}%
\pgfpathlineto{\pgfqpoint{3.617117in}{0.456600in}}%
\pgfpathlineto{\pgfqpoint{3.619387in}{0.487528in}}%
\pgfpathlineto{\pgfqpoint{3.621656in}{0.533042in}}%
\pgfpathlineto{\pgfqpoint{3.623926in}{0.485760in}}%
\pgfpathlineto{\pgfqpoint{3.626195in}{0.508451in}}%
\pgfpathlineto{\pgfqpoint{3.628465in}{0.377621in}}%
\pgfpathlineto{\pgfqpoint{3.630734in}{0.444197in}}%
\pgfpathlineto{\pgfqpoint{3.633003in}{0.443663in}}%
\pgfpathlineto{\pgfqpoint{3.635273in}{0.452132in}}%
\pgfpathlineto{\pgfqpoint{3.637542in}{0.490101in}}%
\pgfpathlineto{\pgfqpoint{3.639812in}{0.493473in}}%
\pgfpathlineto{\pgfqpoint{3.642081in}{0.444250in}}%
\pgfpathlineto{\pgfqpoint{3.646620in}{0.519842in}}%
\pgfpathlineto{\pgfqpoint{3.648890in}{0.521311in}}%
\pgfpathlineto{\pgfqpoint{3.651159in}{0.330517in}}%
\pgfpathlineto{\pgfqpoint{3.653429in}{0.460445in}}%
\pgfpathlineto{\pgfqpoint{3.655698in}{0.522169in}}%
\pgfpathlineto{\pgfqpoint{3.657968in}{0.459420in}}%
\pgfpathlineto{\pgfqpoint{3.660237in}{0.504234in}}%
\pgfpathlineto{\pgfqpoint{3.662507in}{0.442273in}}%
\pgfpathlineto{\pgfqpoint{3.664776in}{0.466808in}}%
\pgfpathlineto{\pgfqpoint{3.667046in}{0.414309in}}%
\pgfpathlineto{\pgfqpoint{3.671585in}{0.498291in}}%
\pgfpathlineto{\pgfqpoint{3.673854in}{0.510277in}}%
\pgfpathlineto{\pgfqpoint{3.676124in}{0.432146in}}%
\pgfpathlineto{\pgfqpoint{3.678393in}{0.513129in}}%
\pgfpathlineto{\pgfqpoint{3.680663in}{0.467411in}}%
\pgfpathlineto{\pgfqpoint{3.682932in}{0.507089in}}%
\pgfpathlineto{\pgfqpoint{3.685202in}{0.440822in}}%
\pgfpathlineto{\pgfqpoint{3.687471in}{0.472086in}}%
\pgfpathlineto{\pgfqpoint{3.689741in}{0.429218in}}%
\pgfpathlineto{\pgfqpoint{3.692010in}{0.438067in}}%
\pgfpathlineto{\pgfqpoint{3.694280in}{0.491779in}}%
\pgfpathlineto{\pgfqpoint{3.696549in}{0.511030in}}%
\pgfpathlineto{\pgfqpoint{3.698818in}{0.465609in}}%
\pgfpathlineto{\pgfqpoint{3.701088in}{0.503474in}}%
\pgfpathlineto{\pgfqpoint{3.703357in}{0.464734in}}%
\pgfpathlineto{\pgfqpoint{3.705627in}{0.506407in}}%
\pgfpathlineto{\pgfqpoint{3.707896in}{0.449391in}}%
\pgfpathlineto{\pgfqpoint{3.710166in}{0.477139in}}%
\pgfpathlineto{\pgfqpoint{3.712435in}{0.440908in}}%
\pgfpathlineto{\pgfqpoint{3.714705in}{0.428818in}}%
\pgfpathlineto{\pgfqpoint{3.716974in}{0.490844in}}%
\pgfpathlineto{\pgfqpoint{3.719244in}{0.511244in}}%
\pgfpathlineto{\pgfqpoint{3.721513in}{0.479126in}}%
\pgfpathlineto{\pgfqpoint{3.723783in}{0.495724in}}%
\pgfpathlineto{\pgfqpoint{3.726052in}{0.460006in}}%
\pgfpathlineto{\pgfqpoint{3.728322in}{0.505219in}}%
\pgfpathlineto{\pgfqpoint{3.730591in}{0.453755in}}%
\pgfpathlineto{\pgfqpoint{3.732861in}{0.481549in}}%
\pgfpathlineto{\pgfqpoint{3.735130in}{0.450195in}}%
\pgfpathlineto{\pgfqpoint{3.737400in}{0.407588in}}%
\pgfpathlineto{\pgfqpoint{3.739669in}{0.486968in}}%
\pgfpathlineto{\pgfqpoint{3.741939in}{0.511093in}}%
\pgfpathlineto{\pgfqpoint{3.744208in}{0.489994in}}%
\pgfpathlineto{\pgfqpoint{3.746478in}{0.485181in}}%
\pgfpathlineto{\pgfqpoint{3.748747in}{0.452196in}}%
\pgfpathlineto{\pgfqpoint{3.751017in}{0.502533in}}%
\pgfpathlineto{\pgfqpoint{3.753286in}{0.457012in}}%
\pgfpathlineto{\pgfqpoint{3.755556in}{0.485408in}}%
\pgfpathlineto{\pgfqpoint{3.757825in}{0.457954in}}%
\pgfpathlineto{\pgfqpoint{3.760095in}{0.356150in}}%
\pgfpathlineto{\pgfqpoint{3.762364in}{0.482033in}}%
\pgfpathlineto{\pgfqpoint{3.764633in}{0.510571in}}%
\pgfpathlineto{\pgfqpoint{3.766903in}{0.498296in}}%
\pgfpathlineto{\pgfqpoint{3.771442in}{0.439704in}}%
\pgfpathlineto{\pgfqpoint{3.773711in}{0.499063in}}%
\pgfpathlineto{\pgfqpoint{3.775981in}{0.458820in}}%
\pgfpathlineto{\pgfqpoint{3.778250in}{0.488611in}}%
\pgfpathlineto{\pgfqpoint{3.780520in}{0.464382in}}%
\pgfpathlineto{\pgfqpoint{3.782789in}{0.377954in}}%
\pgfpathlineto{\pgfqpoint{3.785059in}{0.475107in}}%
\pgfpathlineto{\pgfqpoint{3.787328in}{0.509576in}}%
\pgfpathlineto{\pgfqpoint{3.789598in}{0.504920in}}%
\pgfpathlineto{\pgfqpoint{3.791867in}{0.452856in}}%
\pgfpathlineto{\pgfqpoint{3.794137in}{0.417293in}}%
\pgfpathlineto{\pgfqpoint{3.796406in}{0.494736in}}%
\pgfpathlineto{\pgfqpoint{3.798676in}{0.459318in}}%
\pgfpathlineto{\pgfqpoint{3.800945in}{0.491108in}}%
\pgfpathlineto{\pgfqpoint{3.803215in}{0.469722in}}%
\pgfpathlineto{\pgfqpoint{3.805484in}{0.420346in}}%
\pgfpathlineto{\pgfqpoint{3.810023in}{0.508042in}}%
\pgfpathlineto{\pgfqpoint{3.812293in}{0.510222in}}%
\pgfpathlineto{\pgfqpoint{3.814562in}{0.444698in}}%
\pgfpathlineto{\pgfqpoint{3.819101in}{0.533118in}}%
\pgfpathlineto{\pgfqpoint{3.821371in}{0.485035in}}%
\pgfpathlineto{\pgfqpoint{3.823640in}{0.506353in}}%
\pgfpathlineto{\pgfqpoint{3.825910in}{0.397572in}}%
\pgfpathlineto{\pgfqpoint{3.828179in}{0.436419in}}%
\pgfpathlineto{\pgfqpoint{3.830448in}{0.450883in}}%
\pgfpathlineto{\pgfqpoint{3.832718in}{0.455858in}}%
\pgfpathlineto{\pgfqpoint{3.834987in}{0.491084in}}%
\pgfpathlineto{\pgfqpoint{3.837257in}{0.492032in}}%
\pgfpathlineto{\pgfqpoint{3.839526in}{0.432735in}}%
\pgfpathlineto{\pgfqpoint{3.841796in}{0.485784in}}%
\pgfpathlineto{\pgfqpoint{3.844065in}{0.519565in}}%
\pgfpathlineto{\pgfqpoint{3.846335in}{0.518960in}}%
\pgfpathlineto{\pgfqpoint{3.848604in}{0.397804in}}%
\pgfpathlineto{\pgfqpoint{3.853143in}{0.523060in}}%
\pgfpathlineto{\pgfqpoint{3.855413in}{0.460859in}}%
\pgfpathlineto{\pgfqpoint{3.857682in}{0.503955in}}%
\pgfpathlineto{\pgfqpoint{3.859952in}{0.440220in}}%
\pgfpathlineto{\pgfqpoint{3.862221in}{0.463701in}}%
\pgfpathlineto{\pgfqpoint{3.864491in}{0.395484in}}%
\pgfpathlineto{\pgfqpoint{3.866760in}{0.460559in}}%
\pgfpathlineto{\pgfqpoint{3.869030in}{0.499163in}}%
\pgfpathlineto{\pgfqpoint{3.871299in}{0.508663in}}%
\pgfpathlineto{\pgfqpoint{3.873569in}{0.411292in}}%
\pgfpathlineto{\pgfqpoint{3.875838in}{0.514531in}}%
\pgfpathlineto{\pgfqpoint{3.878108in}{0.469034in}}%
\pgfpathlineto{\pgfqpoint{3.880377in}{0.507028in}}%
\pgfpathlineto{\pgfqpoint{3.882647in}{0.439219in}}%
\pgfpathlineto{\pgfqpoint{3.884916in}{0.469732in}}%
\pgfpathlineto{\pgfqpoint{3.887186in}{0.418916in}}%
\pgfpathlineto{\pgfqpoint{3.889455in}{0.444657in}}%
\pgfpathlineto{\pgfqpoint{3.891725in}{0.493570in}}%
\pgfpathlineto{\pgfqpoint{3.893994in}{0.509924in}}%
\pgfpathlineto{\pgfqpoint{3.896264in}{0.457738in}}%
\pgfpathlineto{\pgfqpoint{3.898533in}{0.505324in}}%
\pgfpathlineto{\pgfqpoint{3.900802in}{0.467006in}}%
\pgfpathlineto{\pgfqpoint{3.903072in}{0.506828in}}%
\pgfpathlineto{\pgfqpoint{3.905341in}{0.448689in}}%
\pgfpathlineto{\pgfqpoint{3.907611in}{0.475256in}}%
\pgfpathlineto{\pgfqpoint{3.909880in}{0.433780in}}%
\pgfpathlineto{\pgfqpoint{3.912150in}{0.437134in}}%
\pgfpathlineto{\pgfqpoint{3.914419in}{0.492971in}}%
\pgfpathlineto{\pgfqpoint{3.916689in}{0.510365in}}%
\pgfpathlineto{\pgfqpoint{3.918958in}{0.473713in}}%
\pgfpathlineto{\pgfqpoint{3.921228in}{0.497963in}}%
\pgfpathlineto{\pgfqpoint{3.923497in}{0.463002in}}%
\pgfpathlineto{\pgfqpoint{3.925767in}{0.506132in}}%
\pgfpathlineto{\pgfqpoint{3.928036in}{0.453635in}}%
\pgfpathlineto{\pgfqpoint{3.930306in}{0.480113in}}%
\pgfpathlineto{\pgfqpoint{3.932575in}{0.445103in}}%
\pgfpathlineto{\pgfqpoint{3.934845in}{0.420826in}}%
\pgfpathlineto{\pgfqpoint{3.937114in}{0.489627in}}%
\pgfpathlineto{\pgfqpoint{3.939384in}{0.510473in}}%
\pgfpathlineto{\pgfqpoint{3.941653in}{0.486052in}}%
\pgfpathlineto{\pgfqpoint{3.943923in}{0.488010in}}%
\pgfpathlineto{\pgfqpoint{3.946192in}{0.456292in}}%
\pgfpathlineto{\pgfqpoint{3.948462in}{0.503918in}}%
\pgfpathlineto{\pgfqpoint{3.950731in}{0.457379in}}%
\pgfpathlineto{\pgfqpoint{3.953001in}{0.484367in}}%
\pgfpathlineto{\pgfqpoint{3.955270in}{0.454189in}}%
\pgfpathlineto{\pgfqpoint{3.957540in}{0.390769in}}%
\pgfpathlineto{\pgfqpoint{3.959809in}{0.485290in}}%
\pgfpathlineto{\pgfqpoint{3.962079in}{0.510191in}}%
\pgfpathlineto{\pgfqpoint{3.964348in}{0.495263in}}%
\pgfpathlineto{\pgfqpoint{3.966617in}{0.475514in}}%
\pgfpathlineto{\pgfqpoint{3.968887in}{0.445710in}}%
\pgfpathlineto{\pgfqpoint{3.971156in}{0.500676in}}%
\pgfpathlineto{\pgfqpoint{3.973426in}{0.459621in}}%
\pgfpathlineto{\pgfqpoint{3.975695in}{0.487940in}}%
\pgfpathlineto{\pgfqpoint{3.977965in}{0.461558in}}%
\pgfpathlineto{\pgfqpoint{3.980234in}{0.316488in}}%
\pgfpathlineto{\pgfqpoint{3.982504in}{0.479166in}}%
\pgfpathlineto{\pgfqpoint{3.984773in}{0.509435in}}%
\pgfpathlineto{\pgfqpoint{3.987043in}{0.502530in}}%
\pgfpathlineto{\pgfqpoint{3.989312in}{0.356745in}}%
\pgfpathlineto{\pgfqpoint{3.991582in}{0.502129in}}%
\pgfpathlineto{\pgfqpoint{3.993851in}{0.534874in}}%
\pgfpathlineto{\pgfqpoint{3.996121in}{0.478619in}}%
\pgfpathlineto{\pgfqpoint{3.998390in}{0.497896in}}%
\pgfpathlineto{\pgfqpoint{4.000660in}{0.441890in}}%
\pgfpathlineto{\pgfqpoint{4.002929in}{0.401959in}}%
\pgfpathlineto{\pgfqpoint{4.005199in}{0.464748in}}%
\pgfpathlineto{\pgfqpoint{4.007468in}{0.461306in}}%
\pgfpathlineto{\pgfqpoint{4.009738in}{0.490411in}}%
\pgfpathlineto{\pgfqpoint{4.012007in}{0.487673in}}%
\pgfpathlineto{\pgfqpoint{4.014277in}{0.319864in}}%
\pgfpathlineto{\pgfqpoint{4.016546in}{0.497478in}}%
\pgfpathlineto{\pgfqpoint{4.018816in}{0.520744in}}%
\pgfpathlineto{\pgfqpoint{4.021085in}{0.510161in}}%
\pgfpathlineto{\pgfqpoint{4.023355in}{0.465441in}}%
\pgfpathlineto{\pgfqpoint{4.025624in}{0.485169in}}%
\pgfpathlineto{\pgfqpoint{4.027894in}{0.528479in}}%
\pgfpathlineto{\pgfqpoint{4.030163in}{0.458310in}}%
\pgfpathlineto{\pgfqpoint{4.032432in}{0.501196in}}%
\pgfpathlineto{\pgfqpoint{4.034702in}{0.423925in}}%
\pgfpathlineto{\pgfqpoint{4.036971in}{0.451306in}}%
\pgfpathlineto{\pgfqpoint{4.039241in}{0.359856in}}%
\pgfpathlineto{\pgfqpoint{4.041510in}{0.468857in}}%
\pgfpathlineto{\pgfqpoint{4.043780in}{0.499103in}}%
\pgfpathlineto{\pgfqpoint{4.046049in}{0.504687in}}%
\pgfpathlineto{\pgfqpoint{4.048319in}{0.421449in}}%
\pgfpathlineto{\pgfqpoint{4.050588in}{0.521295in}}%
\pgfpathlineto{\pgfqpoint{4.052858in}{0.470147in}}%
\pgfpathlineto{\pgfqpoint{4.055127in}{0.505043in}}%
\pgfpathlineto{\pgfqpoint{4.057397in}{0.423749in}}%
\pgfpathlineto{\pgfqpoint{4.059666in}{0.459545in}}%
\pgfpathlineto{\pgfqpoint{4.061936in}{0.356241in}}%
\pgfpathlineto{\pgfqpoint{4.064205in}{0.457302in}}%
\pgfpathlineto{\pgfqpoint{4.066475in}{0.495705in}}%
\pgfpathlineto{\pgfqpoint{4.068744in}{0.507281in}}%
\pgfpathlineto{\pgfqpoint{4.071014in}{0.401297in}}%
\pgfpathlineto{\pgfqpoint{4.073283in}{0.513699in}}%
\pgfpathlineto{\pgfqpoint{4.075553in}{0.470887in}}%
\pgfpathlineto{\pgfqpoint{4.077822in}{0.506317in}}%
\pgfpathlineto{\pgfqpoint{4.080092in}{0.439038in}}%
\pgfpathlineto{\pgfqpoint{4.082361in}{0.466364in}}%
\pgfpathlineto{\pgfqpoint{4.084631in}{0.399359in}}%
\pgfpathlineto{\pgfqpoint{4.089170in}{0.496441in}}%
\pgfpathlineto{\pgfqpoint{4.091439in}{0.508260in}}%
\pgfpathlineto{\pgfqpoint{4.093709in}{0.443680in}}%
\pgfpathlineto{\pgfqpoint{4.095978in}{0.507976in}}%
\pgfpathlineto{\pgfqpoint{4.098247in}{0.469468in}}%
\pgfpathlineto{\pgfqpoint{4.100517in}{0.507172in}}%
\pgfpathlineto{\pgfqpoint{4.102786in}{0.447067in}}%
\pgfpathlineto{\pgfqpoint{4.105056in}{0.472557in}}%
\pgfpathlineto{\pgfqpoint{4.107325in}{0.422085in}}%
\pgfpathlineto{\pgfqpoint{4.109595in}{0.444918in}}%
\pgfpathlineto{\pgfqpoint{4.111864in}{0.494947in}}%
\pgfpathlineto{\pgfqpoint{4.114134in}{0.509093in}}%
\pgfpathlineto{\pgfqpoint{4.116403in}{0.466245in}}%
\pgfpathlineto{\pgfqpoint{4.118673in}{0.500344in}}%
\pgfpathlineto{\pgfqpoint{4.120942in}{0.466188in}}%
\pgfpathlineto{\pgfqpoint{4.123212in}{0.506912in}}%
\pgfpathlineto{\pgfqpoint{4.125481in}{0.453282in}}%
\pgfpathlineto{\pgfqpoint{4.127751in}{0.478052in}}%
\pgfpathlineto{\pgfqpoint{4.130020in}{0.437287in}}%
\pgfpathlineto{\pgfqpoint{4.132290in}{0.432958in}}%
\pgfpathlineto{\pgfqpoint{4.134559in}{0.492557in}}%
\pgfpathlineto{\pgfqpoint{4.136829in}{0.509540in}}%
\pgfpathlineto{\pgfqpoint{4.139098in}{0.480627in}}%
\pgfpathlineto{\pgfqpoint{4.141368in}{0.491208in}}%
\pgfpathlineto{\pgfqpoint{4.143637in}{0.460628in}}%
\pgfpathlineto{\pgfqpoint{4.145907in}{0.505426in}}%
\pgfpathlineto{\pgfqpoint{4.148176in}{0.457585in}}%
\pgfpathlineto{\pgfqpoint{4.150446in}{0.482833in}}%
\pgfpathlineto{\pgfqpoint{4.154985in}{0.412757in}}%
\pgfpathlineto{\pgfqpoint{4.157254in}{0.488805in}}%
\pgfpathlineto{\pgfqpoint{4.159524in}{0.509564in}}%
\pgfpathlineto{\pgfqpoint{4.161793in}{0.491196in}}%
\pgfpathlineto{\pgfqpoint{4.164062in}{0.479619in}}%
\pgfpathlineto{\pgfqpoint{4.166332in}{0.451843in}}%
\pgfpathlineto{\pgfqpoint{4.168601in}{0.502521in}}%
\pgfpathlineto{\pgfqpoint{4.170871in}{0.460393in}}%
\pgfpathlineto{\pgfqpoint{4.173140in}{0.486892in}}%
\pgfpathlineto{\pgfqpoint{4.175410in}{0.457466in}}%
\pgfpathlineto{\pgfqpoint{4.177679in}{0.363678in}}%
\pgfpathlineto{\pgfqpoint{4.179949in}{0.483550in}}%
\pgfpathlineto{\pgfqpoint{4.182218in}{0.509110in}}%
\pgfpathlineto{\pgfqpoint{4.184488in}{0.499354in}}%
\pgfpathlineto{\pgfqpoint{4.186757in}{0.411984in}}%
\pgfpathlineto{\pgfqpoint{4.189027in}{0.503580in}}%
\pgfpathlineto{\pgfqpoint{4.191296in}{0.534240in}}%
\pgfpathlineto{\pgfqpoint{4.193566in}{0.476781in}}%
\pgfpathlineto{\pgfqpoint{4.195835in}{0.494737in}}%
\pgfpathlineto{\pgfqpoint{4.198105in}{0.448290in}}%
\pgfpathlineto{\pgfqpoint{4.200374in}{0.376499in}}%
\pgfpathlineto{\pgfqpoint{4.202644in}{0.469449in}}%
\pgfpathlineto{\pgfqpoint{4.204913in}{0.463617in}}%
\pgfpathlineto{\pgfqpoint{4.207183in}{0.490387in}}%
\pgfpathlineto{\pgfqpoint{4.209452in}{0.485313in}}%
\pgfpathlineto{\pgfqpoint{4.211722in}{0.372663in}}%
\pgfpathlineto{\pgfqpoint{4.213991in}{0.500311in}}%
\pgfpathlineto{\pgfqpoint{4.216261in}{0.519996in}}%
\pgfpathlineto{\pgfqpoint{4.218530in}{0.506830in}}%
\pgfpathlineto{\pgfqpoint{4.220800in}{0.471832in}}%
\pgfpathlineto{\pgfqpoint{4.223069in}{0.487715in}}%
\pgfpathlineto{\pgfqpoint{4.225339in}{0.528808in}}%
\pgfpathlineto{\pgfqpoint{4.227608in}{0.458546in}}%
\pgfpathlineto{\pgfqpoint{4.229878in}{0.500087in}}%
\pgfpathlineto{\pgfqpoint{4.232147in}{0.418057in}}%
\pgfpathlineto{\pgfqpoint{4.234416in}{0.445995in}}%
\pgfpathlineto{\pgfqpoint{4.236686in}{0.397608in}}%
\pgfpathlineto{\pgfqpoint{4.238955in}{0.471622in}}%
\pgfpathlineto{\pgfqpoint{4.241225in}{0.499100in}}%
\pgfpathlineto{\pgfqpoint{4.243494in}{0.502313in}}%
\pgfpathlineto{\pgfqpoint{4.245764in}{0.437653in}}%
\pgfpathlineto{\pgfqpoint{4.248033in}{0.522257in}}%
\pgfpathlineto{\pgfqpoint{4.250303in}{0.470725in}}%
\pgfpathlineto{\pgfqpoint{4.252572in}{0.504179in}}%
\pgfpathlineto{\pgfqpoint{4.254842in}{0.418752in}}%
\pgfpathlineto{\pgfqpoint{4.257111in}{0.455618in}}%
\pgfpathlineto{\pgfqpoint{4.259381in}{0.317755in}}%
\pgfpathlineto{\pgfqpoint{4.261650in}{0.461081in}}%
\pgfpathlineto{\pgfqpoint{4.263920in}{0.496383in}}%
\pgfpathlineto{\pgfqpoint{4.266189in}{0.505582in}}%
\pgfpathlineto{\pgfqpoint{4.268459in}{0.341476in}}%
\pgfpathlineto{\pgfqpoint{4.270728in}{0.515022in}}%
\pgfpathlineto{\pgfqpoint{4.272998in}{0.472012in}}%
\pgfpathlineto{\pgfqpoint{4.275267in}{0.505889in}}%
\pgfpathlineto{\pgfqpoint{4.277537in}{0.436445in}}%
\pgfpathlineto{\pgfqpoint{4.279806in}{0.463204in}}%
\pgfpathlineto{\pgfqpoint{4.282076in}{0.371339in}}%
\pgfpathlineto{\pgfqpoint{4.284345in}{0.457818in}}%
\pgfpathlineto{\pgfqpoint{4.286615in}{0.497410in}}%
\pgfpathlineto{\pgfqpoint{4.288884in}{0.506840in}}%
\pgfpathlineto{\pgfqpoint{4.291154in}{0.430054in}}%
\pgfpathlineto{\pgfqpoint{4.293423in}{0.509581in}}%
\pgfpathlineto{\pgfqpoint{4.295693in}{0.471077in}}%
\pgfpathlineto{\pgfqpoint{4.297962in}{0.507136in}}%
\pgfpathlineto{\pgfqpoint{4.300231in}{0.445580in}}%
\pgfpathlineto{\pgfqpoint{4.302501in}{0.470062in}}%
\pgfpathlineto{\pgfqpoint{4.304770in}{0.409108in}}%
\pgfpathlineto{\pgfqpoint{4.309309in}{0.496399in}}%
\pgfpathlineto{\pgfqpoint{4.311579in}{0.507983in}}%
\pgfpathlineto{\pgfqpoint{4.313848in}{0.458935in}}%
\pgfpathlineto{\pgfqpoint{4.316118in}{0.502327in}}%
\pgfpathlineto{\pgfqpoint{4.318387in}{0.468395in}}%
\pgfpathlineto{\pgfqpoint{4.320657in}{0.507320in}}%
\pgfpathlineto{\pgfqpoint{4.322926in}{0.452563in}}%
\pgfpathlineto{\pgfqpoint{4.325196in}{0.476100in}}%
\pgfpathlineto{\pgfqpoint{4.327465in}{0.429342in}}%
\pgfpathlineto{\pgfqpoint{4.329735in}{0.440374in}}%
\pgfpathlineto{\pgfqpoint{4.332004in}{0.494556in}}%
\pgfpathlineto{\pgfqpoint{4.334274in}{0.350113in}}%
\pgfpathlineto{\pgfqpoint{4.336543in}{0.448059in}}%
\pgfpathlineto{\pgfqpoint{4.336543in}{0.448059in}}%
\pgfusepath{stroke}%
\end{pgfscope}%
\begin{pgfscope}%
\pgfpathrectangle{\pgfqpoint{2.520957in}{0.320625in}}{\pgfqpoint{1.815586in}{1.164375in}}%
\pgfusepath{clip}%
\pgfsetrectcap%
\pgfsetroundjoin%
\pgfsetlinewidth{1.154312pt}%
\definecolor{currentstroke}{rgb}{0.989363,0.557873,0.391671}%
\pgfsetstrokecolor{currentstroke}%
\pgfsetdash{}{0pt}%
\pgfpathmoveto{\pgfqpoint{2.520957in}{1.451682in}}%
\pgfpathlineto{\pgfqpoint{2.525496in}{1.441759in}}%
\pgfpathlineto{\pgfqpoint{2.541382in}{1.410222in}}%
\pgfpathlineto{\pgfqpoint{2.557269in}{1.375789in}}%
\pgfpathlineto{\pgfqpoint{2.589041in}{1.303288in}}%
\pgfpathlineto{\pgfqpoint{2.634431in}{1.200481in}}%
\pgfpathlineto{\pgfqpoint{2.675282in}{1.110870in}}%
\pgfpathlineto{\pgfqpoint{2.718402in}{1.018960in}}%
\pgfpathlineto{\pgfqpoint{2.768330in}{0.915178in}}%
\pgfpathlineto{\pgfqpoint{2.827337in}{0.795150in}}%
\pgfpathlineto{\pgfqpoint{2.899960in}{0.650081in}}%
\pgfpathlineto{\pgfqpoint{2.940811in}{0.569486in}}%
\pgfpathlineto{\pgfqpoint{2.943081in}{0.566081in}}%
\pgfpathlineto{\pgfqpoint{2.947620in}{0.556734in}}%
\pgfpathlineto{\pgfqpoint{2.952159in}{0.547196in}}%
\pgfpathlineto{\pgfqpoint{2.954428in}{0.545284in}}%
\pgfpathlineto{\pgfqpoint{2.956698in}{0.541835in}}%
\pgfpathlineto{\pgfqpoint{2.961236in}{0.529498in}}%
\pgfpathlineto{\pgfqpoint{2.963506in}{0.528486in}}%
\pgfpathlineto{\pgfqpoint{2.965775in}{0.537313in}}%
\pgfpathlineto{\pgfqpoint{2.968045in}{0.519480in}}%
\pgfpathlineto{\pgfqpoint{2.970314in}{0.518408in}}%
\pgfpathlineto{\pgfqpoint{2.972584in}{0.507289in}}%
\pgfpathlineto{\pgfqpoint{2.974853in}{0.502624in}}%
\pgfpathlineto{\pgfqpoint{2.977123in}{0.507051in}}%
\pgfpathlineto{\pgfqpoint{2.979392in}{0.501973in}}%
\pgfpathlineto{\pgfqpoint{2.981662in}{0.506672in}}%
\pgfpathlineto{\pgfqpoint{2.986201in}{0.483177in}}%
\pgfpathlineto{\pgfqpoint{2.988470in}{0.511970in}}%
\pgfpathlineto{\pgfqpoint{2.990740in}{0.512400in}}%
\pgfpathlineto{\pgfqpoint{2.993009in}{0.505584in}}%
\pgfpathlineto{\pgfqpoint{2.995279in}{0.470253in}}%
\pgfpathlineto{\pgfqpoint{2.997548in}{0.479110in}}%
\pgfpathlineto{\pgfqpoint{2.999818in}{0.518211in}}%
\pgfpathlineto{\pgfqpoint{3.002087in}{0.478691in}}%
\pgfpathlineto{\pgfqpoint{3.004357in}{0.500495in}}%
\pgfpathlineto{\pgfqpoint{3.006626in}{0.452462in}}%
\pgfpathlineto{\pgfqpoint{3.008896in}{0.444152in}}%
\pgfpathlineto{\pgfqpoint{3.015704in}{0.503696in}}%
\pgfpathlineto{\pgfqpoint{3.017974in}{0.486670in}}%
\pgfpathlineto{\pgfqpoint{3.020243in}{0.449618in}}%
\pgfpathlineto{\pgfqpoint{3.022513in}{0.517285in}}%
\pgfpathlineto{\pgfqpoint{3.024782in}{0.507951in}}%
\pgfpathlineto{\pgfqpoint{3.027052in}{0.510812in}}%
\pgfpathlineto{\pgfqpoint{3.029321in}{0.428779in}}%
\pgfpathlineto{\pgfqpoint{3.031590in}{0.397208in}}%
\pgfpathlineto{\pgfqpoint{3.033860in}{0.494980in}}%
\pgfpathlineto{\pgfqpoint{3.036129in}{0.461979in}}%
\pgfpathlineto{\pgfqpoint{3.038399in}{0.496749in}}%
\pgfpathlineto{\pgfqpoint{3.040668in}{0.470430in}}%
\pgfpathlineto{\pgfqpoint{3.042938in}{0.453733in}}%
\pgfpathlineto{\pgfqpoint{3.045207in}{0.428720in}}%
\pgfpathlineto{\pgfqpoint{3.047477in}{0.496193in}}%
\pgfpathlineto{\pgfqpoint{3.049746in}{0.513767in}}%
\pgfpathlineto{\pgfqpoint{3.052016in}{0.489080in}}%
\pgfpathlineto{\pgfqpoint{3.054285in}{0.445053in}}%
\pgfpathlineto{\pgfqpoint{3.056555in}{0.521954in}}%
\pgfpathlineto{\pgfqpoint{3.058824in}{0.496554in}}%
\pgfpathlineto{\pgfqpoint{3.061094in}{0.515578in}}%
\pgfpathlineto{\pgfqpoint{3.063363in}{0.394805in}}%
\pgfpathlineto{\pgfqpoint{3.065633in}{0.454496in}}%
\pgfpathlineto{\pgfqpoint{3.070172in}{0.447307in}}%
\pgfpathlineto{\pgfqpoint{3.072441in}{0.490715in}}%
\pgfpathlineto{\pgfqpoint{3.074711in}{0.488879in}}%
\pgfpathlineto{\pgfqpoint{3.076980in}{0.468883in}}%
\pgfpathlineto{\pgfqpoint{3.079250in}{0.429189in}}%
\pgfpathlineto{\pgfqpoint{3.081519in}{0.508167in}}%
\pgfpathlineto{\pgfqpoint{3.083789in}{0.524071in}}%
\pgfpathlineto{\pgfqpoint{3.086058in}{0.478160in}}%
\pgfpathlineto{\pgfqpoint{3.088328in}{0.379685in}}%
\pgfpathlineto{\pgfqpoint{3.090597in}{0.515044in}}%
\pgfpathlineto{\pgfqpoint{3.092867in}{0.472911in}}%
\pgfpathlineto{\pgfqpoint{3.095136in}{0.510870in}}%
\pgfpathlineto{\pgfqpoint{3.097405in}{0.448709in}}%
\pgfpathlineto{\pgfqpoint{3.099675in}{0.475088in}}%
\pgfpathlineto{\pgfqpoint{3.101944in}{0.408707in}}%
\pgfpathlineto{\pgfqpoint{3.104214in}{0.436895in}}%
\pgfpathlineto{\pgfqpoint{3.106483in}{0.491362in}}%
\pgfpathlineto{\pgfqpoint{3.108753in}{0.505037in}}%
\pgfpathlineto{\pgfqpoint{3.111022in}{0.477786in}}%
\pgfpathlineto{\pgfqpoint{3.113292in}{0.470753in}}%
\pgfpathlineto{\pgfqpoint{3.115561in}{0.514552in}}%
\pgfpathlineto{\pgfqpoint{3.117831in}{0.531248in}}%
\pgfpathlineto{\pgfqpoint{3.120100in}{0.454386in}}%
\pgfpathlineto{\pgfqpoint{3.122370in}{0.459804in}}%
\pgfpathlineto{\pgfqpoint{3.124639in}{0.489377in}}%
\pgfpathlineto{\pgfqpoint{3.126909in}{0.426802in}}%
\pgfpathlineto{\pgfqpoint{3.129178in}{0.489518in}}%
\pgfpathlineto{\pgfqpoint{3.131448in}{0.463382in}}%
\pgfpathlineto{\pgfqpoint{3.133717in}{0.482867in}}%
\pgfpathlineto{\pgfqpoint{3.135987in}{0.466680in}}%
\pgfpathlineto{\pgfqpoint{3.138256in}{0.445711in}}%
\pgfpathlineto{\pgfqpoint{3.140526in}{0.504797in}}%
\pgfpathlineto{\pgfqpoint{3.142795in}{0.517382in}}%
\pgfpathlineto{\pgfqpoint{3.145065in}{0.480680in}}%
\pgfpathlineto{\pgfqpoint{3.147334in}{0.501066in}}%
\pgfpathlineto{\pgfqpoint{3.149604in}{0.507556in}}%
\pgfpathlineto{\pgfqpoint{3.151873in}{0.532429in}}%
\pgfpathlineto{\pgfqpoint{3.154143in}{0.438282in}}%
\pgfpathlineto{\pgfqpoint{3.156412in}{0.488110in}}%
\pgfpathlineto{\pgfqpoint{3.158682in}{0.420114in}}%
\pgfpathlineto{\pgfqpoint{3.160951in}{0.384867in}}%
\pgfpathlineto{\pgfqpoint{3.163220in}{0.455853in}}%
\pgfpathlineto{\pgfqpoint{3.167759in}{0.491465in}}%
\pgfpathlineto{\pgfqpoint{3.170029in}{0.488456in}}%
\pgfpathlineto{\pgfqpoint{3.172298in}{0.473639in}}%
\pgfpathlineto{\pgfqpoint{3.174568in}{0.524132in}}%
\pgfpathlineto{\pgfqpoint{3.176837in}{0.469803in}}%
\pgfpathlineto{\pgfqpoint{3.179107in}{0.494330in}}%
\pgfpathlineto{\pgfqpoint{3.181114in}{0.315625in}}%
\pgfpathmoveto{\pgfqpoint{3.181836in}{0.315625in}}%
\pgfpathlineto{\pgfqpoint{3.183646in}{0.407805in}}%
\pgfpathlineto{\pgfqpoint{3.185915in}{0.452032in}}%
\pgfpathlineto{\pgfqpoint{3.188185in}{0.477284in}}%
\pgfpathlineto{\pgfqpoint{3.190454in}{0.495869in}}%
\pgfpathlineto{\pgfqpoint{3.192724in}{0.490121in}}%
\pgfpathlineto{\pgfqpoint{3.194993in}{0.464913in}}%
\pgfpathlineto{\pgfqpoint{3.197263in}{0.522898in}}%
\pgfpathlineto{\pgfqpoint{3.199532in}{0.473094in}}%
\pgfpathlineto{\pgfqpoint{3.201802in}{0.498880in}}%
\pgfpathlineto{\pgfqpoint{3.204071in}{0.384829in}}%
\pgfpathlineto{\pgfqpoint{3.206341in}{0.428719in}}%
\pgfpathlineto{\pgfqpoint{3.208610in}{0.441337in}}%
\pgfpathlineto{\pgfqpoint{3.210880in}{0.474429in}}%
\pgfpathlineto{\pgfqpoint{3.213149in}{0.497721in}}%
\pgfpathlineto{\pgfqpoint{3.215419in}{0.494693in}}%
\pgfpathlineto{\pgfqpoint{3.217688in}{0.442664in}}%
\pgfpathlineto{\pgfqpoint{3.219958in}{0.518740in}}%
\pgfpathlineto{\pgfqpoint{3.222227in}{0.475609in}}%
\pgfpathlineto{\pgfqpoint{3.224497in}{0.502531in}}%
\pgfpathlineto{\pgfqpoint{3.226766in}{0.419450in}}%
\pgfpathlineto{\pgfqpoint{3.229035in}{0.443334in}}%
\pgfpathlineto{\pgfqpoint{3.231305in}{0.426828in}}%
\pgfpathlineto{\pgfqpoint{3.233574in}{0.471842in}}%
\pgfpathlineto{\pgfqpoint{3.235844in}{0.499453in}}%
\pgfpathlineto{\pgfqpoint{3.238113in}{0.497986in}}%
\pgfpathlineto{\pgfqpoint{3.240383in}{0.396244in}}%
\pgfpathlineto{\pgfqpoint{3.242652in}{0.513469in}}%
\pgfpathlineto{\pgfqpoint{3.244922in}{0.517598in}}%
\pgfpathlineto{\pgfqpoint{3.247191in}{0.515850in}}%
\pgfpathlineto{\pgfqpoint{3.249461in}{0.451617in}}%
\pgfpathlineto{\pgfqpoint{3.251730in}{0.408552in}}%
\pgfpathlineto{\pgfqpoint{3.254000in}{0.503914in}}%
\pgfpathlineto{\pgfqpoint{3.256269in}{0.449881in}}%
\pgfpathlineto{\pgfqpoint{3.258539in}{0.494786in}}%
\pgfpathlineto{\pgfqpoint{3.260808in}{0.462647in}}%
\pgfpathlineto{\pgfqpoint{3.263078in}{0.464171in}}%
\pgfpathlineto{\pgfqpoint{3.265347in}{0.404992in}}%
\pgfpathlineto{\pgfqpoint{3.267617in}{0.485641in}}%
\pgfpathlineto{\pgfqpoint{3.269886in}{0.511891in}}%
\pgfpathlineto{\pgfqpoint{3.272156in}{0.505627in}}%
\pgfpathlineto{\pgfqpoint{3.274425in}{0.441613in}}%
\pgfpathlineto{\pgfqpoint{3.276695in}{0.527677in}}%
\pgfpathlineto{\pgfqpoint{3.278964in}{0.507510in}}%
\pgfpathlineto{\pgfqpoint{3.281234in}{0.520744in}}%
\pgfpathlineto{\pgfqpoint{3.283503in}{0.412706in}}%
\pgfpathlineto{\pgfqpoint{3.285773in}{0.456955in}}%
\pgfpathlineto{\pgfqpoint{3.288042in}{0.458398in}}%
\pgfpathlineto{\pgfqpoint{3.290312in}{0.421447in}}%
\pgfpathlineto{\pgfqpoint{3.292581in}{0.478553in}}%
\pgfpathlineto{\pgfqpoint{3.294851in}{0.482028in}}%
\pgfpathlineto{\pgfqpoint{3.297120in}{0.476878in}}%
\pgfpathlineto{\pgfqpoint{3.299389in}{0.436049in}}%
\pgfpathlineto{\pgfqpoint{3.301659in}{0.505573in}}%
\pgfpathlineto{\pgfqpoint{3.303928in}{0.525817in}}%
\pgfpathlineto{\pgfqpoint{3.306198in}{0.412065in}}%
\pgfpathlineto{\pgfqpoint{3.308467in}{0.465748in}}%
\pgfpathlineto{\pgfqpoint{3.310737in}{0.457811in}}%
\pgfpathlineto{\pgfqpoint{3.313006in}{0.401667in}}%
\pgfpathlineto{\pgfqpoint{3.315276in}{0.473495in}}%
\pgfpathlineto{\pgfqpoint{3.317545in}{0.478346in}}%
\pgfpathlineto{\pgfqpoint{3.319815in}{0.481808in}}%
\pgfpathlineto{\pgfqpoint{3.322084in}{0.463921in}}%
\pgfpathlineto{\pgfqpoint{3.326623in}{0.525783in}}%
\pgfpathlineto{\pgfqpoint{3.328893in}{0.445691in}}%
\pgfpathlineto{\pgfqpoint{3.331162in}{0.475744in}}%
\pgfpathlineto{\pgfqpoint{3.333432in}{0.443407in}}%
\pgfpathlineto{\pgfqpoint{3.335701in}{0.383574in}}%
\pgfpathlineto{\pgfqpoint{3.337971in}{0.471331in}}%
\pgfpathlineto{\pgfqpoint{3.342510in}{0.486382in}}%
\pgfpathlineto{\pgfqpoint{3.344779in}{0.470519in}}%
\pgfpathlineto{\pgfqpoint{3.347049in}{0.492129in}}%
\pgfpathlineto{\pgfqpoint{3.349318in}{0.526087in}}%
\pgfpathlineto{\pgfqpoint{3.351588in}{0.457953in}}%
\pgfpathlineto{\pgfqpoint{3.353857in}{0.483494in}}%
\pgfpathlineto{\pgfqpoint{3.356127in}{0.426209in}}%
\pgfpathlineto{\pgfqpoint{3.358387in}{0.315625in}}%
\pgfpathmoveto{\pgfqpoint{3.358403in}{0.315625in}}%
\pgfpathlineto{\pgfqpoint{3.360666in}{0.467089in}}%
\pgfpathlineto{\pgfqpoint{3.365204in}{0.490280in}}%
\pgfpathlineto{\pgfqpoint{3.367474in}{0.478240in}}%
\pgfpathlineto{\pgfqpoint{3.369743in}{0.484055in}}%
\pgfpathlineto{\pgfqpoint{3.372013in}{0.525332in}}%
\pgfpathlineto{\pgfqpoint{3.374282in}{0.466475in}}%
\pgfpathlineto{\pgfqpoint{3.376552in}{0.490064in}}%
\pgfpathlineto{\pgfqpoint{3.378821in}{0.394197in}}%
\pgfpathlineto{\pgfqpoint{3.381091in}{0.370631in}}%
\pgfpathlineto{\pgfqpoint{3.383360in}{0.461585in}}%
\pgfpathlineto{\pgfqpoint{3.387899in}{0.493721in}}%
\pgfpathlineto{\pgfqpoint{3.392438in}{0.473604in}}%
\pgfpathlineto{\pgfqpoint{3.394708in}{0.523672in}}%
\pgfpathlineto{\pgfqpoint{3.396977in}{0.471948in}}%
\pgfpathlineto{\pgfqpoint{3.399247in}{0.495523in}}%
\pgfpathlineto{\pgfqpoint{3.401345in}{0.315625in}}%
\pgfpathmoveto{\pgfqpoint{3.401825in}{0.315625in}}%
\pgfpathlineto{\pgfqpoint{3.403786in}{0.408686in}}%
\pgfpathlineto{\pgfqpoint{3.406055in}{0.454058in}}%
\pgfpathlineto{\pgfqpoint{3.410594in}{0.496572in}}%
\pgfpathlineto{\pgfqpoint{3.412864in}{0.489574in}}%
\pgfpathlineto{\pgfqpoint{3.415133in}{0.458507in}}%
\pgfpathlineto{\pgfqpoint{3.417403in}{0.520719in}}%
\pgfpathlineto{\pgfqpoint{3.419672in}{0.475349in}}%
\pgfpathlineto{\pgfqpoint{3.421942in}{0.499994in}}%
\pgfpathlineto{\pgfqpoint{3.424211in}{0.401266in}}%
\pgfpathlineto{\pgfqpoint{3.426481in}{0.430171in}}%
\pgfpathlineto{\pgfqpoint{3.428750in}{0.443627in}}%
\pgfpathlineto{\pgfqpoint{3.431019in}{0.475245in}}%
\pgfpathlineto{\pgfqpoint{3.433289in}{0.498810in}}%
\pgfpathlineto{\pgfqpoint{3.435558in}{0.493905in}}%
\pgfpathlineto{\pgfqpoint{3.437828in}{0.433073in}}%
\pgfpathlineto{\pgfqpoint{3.440097in}{0.516299in}}%
\pgfpathlineto{\pgfqpoint{3.442367in}{0.516634in}}%
\pgfpathlineto{\pgfqpoint{3.444636in}{0.511004in}}%
\pgfpathlineto{\pgfqpoint{3.446906in}{0.464323in}}%
\pgfpathlineto{\pgfqpoint{3.449175in}{0.434644in}}%
\pgfpathlineto{\pgfqpoint{3.451445in}{0.507075in}}%
\pgfpathlineto{\pgfqpoint{3.453714in}{0.453325in}}%
\pgfpathlineto{\pgfqpoint{3.455984in}{0.495280in}}%
\pgfpathlineto{\pgfqpoint{3.458253in}{0.460265in}}%
\pgfpathlineto{\pgfqpoint{3.460523in}{0.457322in}}%
\pgfpathlineto{\pgfqpoint{3.462367in}{0.315625in}}%
\pgfpathmoveto{\pgfqpoint{3.463153in}{0.315625in}}%
\pgfpathlineto{\pgfqpoint{3.465062in}{0.488785in}}%
\pgfpathlineto{\pgfqpoint{3.467331in}{0.510765in}}%
\pgfpathlineto{\pgfqpoint{3.469601in}{0.500910in}}%
\pgfpathlineto{\pgfqpoint{3.471870in}{0.459883in}}%
\pgfpathlineto{\pgfqpoint{3.474140in}{0.529473in}}%
\pgfpathlineto{\pgfqpoint{3.476409in}{0.506762in}}%
\pgfpathlineto{\pgfqpoint{3.478679in}{0.517288in}}%
\pgfpathlineto{\pgfqpoint{3.480948in}{0.431893in}}%
\pgfpathlineto{\pgfqpoint{3.483218in}{0.445532in}}%
\pgfpathlineto{\pgfqpoint{3.485487in}{0.466054in}}%
\pgfpathlineto{\pgfqpoint{3.487757in}{0.432355in}}%
\pgfpathlineto{\pgfqpoint{3.490026in}{0.482450in}}%
\pgfpathlineto{\pgfqpoint{3.492296in}{0.481594in}}%
\pgfpathlineto{\pgfqpoint{3.494565in}{0.471296in}}%
\pgfpathlineto{\pgfqpoint{3.496834in}{0.398948in}}%
\pgfpathlineto{\pgfqpoint{3.499104in}{0.507558in}}%
\pgfpathlineto{\pgfqpoint{3.501373in}{0.523881in}}%
\pgfpathlineto{\pgfqpoint{3.503643in}{0.380250in}}%
\pgfpathlineto{\pgfqpoint{3.505912in}{0.456200in}}%
\pgfpathlineto{\pgfqpoint{3.508182in}{0.465358in}}%
\pgfpathlineto{\pgfqpoint{3.510451in}{0.418473in}}%
\pgfpathlineto{\pgfqpoint{3.512721in}{0.478181in}}%
\pgfpathlineto{\pgfqpoint{3.514990in}{0.478571in}}%
\pgfpathlineto{\pgfqpoint{3.517260in}{0.477797in}}%
\pgfpathlineto{\pgfqpoint{3.519529in}{0.450223in}}%
\pgfpathlineto{\pgfqpoint{3.521799in}{0.500534in}}%
\pgfpathlineto{\pgfqpoint{3.524068in}{0.524548in}}%
\pgfpathlineto{\pgfqpoint{3.526338in}{0.436554in}}%
\pgfpathlineto{\pgfqpoint{3.528607in}{0.468646in}}%
\pgfpathlineto{\pgfqpoint{3.530877in}{0.453307in}}%
\pgfpathlineto{\pgfqpoint{3.533146in}{0.408734in}}%
\pgfpathlineto{\pgfqpoint{3.535416in}{0.476531in}}%
\pgfpathlineto{\pgfqpoint{3.539955in}{0.482986in}}%
\pgfpathlineto{\pgfqpoint{3.542224in}{0.459641in}}%
\pgfpathlineto{\pgfqpoint{3.546763in}{0.525351in}}%
\pgfpathlineto{\pgfqpoint{3.549033in}{0.452605in}}%
\pgfpathlineto{\pgfqpoint{3.551302in}{0.477901in}}%
\pgfpathlineto{\pgfqpoint{3.553572in}{0.439715in}}%
\pgfpathlineto{\pgfqpoint{3.555841in}{0.386221in}}%
\pgfpathlineto{\pgfqpoint{3.558111in}{0.473056in}}%
\pgfpathlineto{\pgfqpoint{3.562649in}{0.487560in}}%
\pgfpathlineto{\pgfqpoint{3.564919in}{0.469996in}}%
\pgfpathlineto{\pgfqpoint{3.567188in}{0.488950in}}%
\pgfpathlineto{\pgfqpoint{3.569458in}{0.525180in}}%
\pgfpathlineto{\pgfqpoint{3.571727in}{0.463260in}}%
\pgfpathlineto{\pgfqpoint{3.573997in}{0.485631in}}%
\pgfpathlineto{\pgfqpoint{3.576266in}{0.417924in}}%
\pgfpathlineto{\pgfqpoint{3.578434in}{0.315625in}}%
\pgfpathmoveto{\pgfqpoint{3.578605in}{0.315625in}}%
\pgfpathlineto{\pgfqpoint{3.580805in}{0.468548in}}%
\pgfpathlineto{\pgfqpoint{3.585344in}{0.491588in}}%
\pgfpathlineto{\pgfqpoint{3.587614in}{0.477867in}}%
\pgfpathlineto{\pgfqpoint{3.589883in}{0.480239in}}%
\pgfpathlineto{\pgfqpoint{3.592153in}{0.524177in}}%
\pgfpathlineto{\pgfqpoint{3.594422in}{0.470087in}}%
\pgfpathlineto{\pgfqpoint{3.596692in}{0.492011in}}%
\pgfpathlineto{\pgfqpoint{3.598961in}{0.371754in}}%
\pgfpathlineto{\pgfqpoint{3.601231in}{0.377452in}}%
\pgfpathlineto{\pgfqpoint{3.603500in}{0.462442in}}%
\pgfpathlineto{\pgfqpoint{3.608039in}{0.494995in}}%
\pgfpathlineto{\pgfqpoint{3.610309in}{0.484383in}}%
\pgfpathlineto{\pgfqpoint{3.612578in}{0.468242in}}%
\pgfpathlineto{\pgfqpoint{3.614848in}{0.522033in}}%
\pgfpathlineto{\pgfqpoint{3.617117in}{0.514963in}}%
\pgfpathlineto{\pgfqpoint{3.619387in}{0.496019in}}%
\pgfpathlineto{\pgfqpoint{3.621656in}{0.490561in}}%
\pgfpathlineto{\pgfqpoint{3.623926in}{0.463855in}}%
\pgfpathlineto{\pgfqpoint{3.626195in}{0.514004in}}%
\pgfpathlineto{\pgfqpoint{3.628465in}{0.453486in}}%
\pgfpathlineto{\pgfqpoint{3.630734in}{0.492991in}}%
\pgfpathlineto{\pgfqpoint{3.633003in}{0.450727in}}%
\pgfpathlineto{\pgfqpoint{3.637542in}{0.426045in}}%
\pgfpathlineto{\pgfqpoint{3.639812in}{0.493809in}}%
\pgfpathlineto{\pgfqpoint{3.642081in}{0.506737in}}%
\pgfpathlineto{\pgfqpoint{3.644351in}{0.488888in}}%
\pgfpathlineto{\pgfqpoint{3.646620in}{0.487309in}}%
\pgfpathlineto{\pgfqpoint{3.648890in}{0.534502in}}%
\pgfpathlineto{\pgfqpoint{3.651159in}{0.504054in}}%
\pgfpathlineto{\pgfqpoint{3.653429in}{0.506763in}}%
\pgfpathlineto{\pgfqpoint{3.655698in}{0.467683in}}%
\pgfpathlineto{\pgfqpoint{3.657968in}{0.398601in}}%
\pgfpathlineto{\pgfqpoint{3.660237in}{0.479460in}}%
\pgfpathlineto{\pgfqpoint{3.662507in}{0.443790in}}%
\pgfpathlineto{\pgfqpoint{3.664776in}{0.485860in}}%
\pgfpathlineto{\pgfqpoint{3.667046in}{0.478640in}}%
\pgfpathlineto{\pgfqpoint{3.669315in}{0.454486in}}%
\pgfpathlineto{\pgfqpoint{3.671585in}{0.422021in}}%
\pgfpathlineto{\pgfqpoint{3.673854in}{0.512177in}}%
\pgfpathlineto{\pgfqpoint{3.676124in}{0.518563in}}%
\pgfpathlineto{\pgfqpoint{3.678393in}{0.427020in}}%
\pgfpathlineto{\pgfqpoint{3.680663in}{0.418525in}}%
\pgfpathlineto{\pgfqpoint{3.682932in}{0.479494in}}%
\pgfpathlineto{\pgfqpoint{3.685202in}{0.435492in}}%
\pgfpathlineto{\pgfqpoint{3.687471in}{0.483634in}}%
\pgfpathlineto{\pgfqpoint{3.689741in}{0.476747in}}%
\pgfpathlineto{\pgfqpoint{3.692010in}{0.465397in}}%
\pgfpathlineto{\pgfqpoint{3.694280in}{0.388164in}}%
\pgfpathlineto{\pgfqpoint{3.696549in}{0.506984in}}%
\pgfpathlineto{\pgfqpoint{3.698818in}{0.520894in}}%
\pgfpathlineto{\pgfqpoint{3.701088in}{0.337570in}}%
\pgfpathlineto{\pgfqpoint{3.703357in}{0.443298in}}%
\pgfpathlineto{\pgfqpoint{3.705627in}{0.471343in}}%
\pgfpathlineto{\pgfqpoint{3.707896in}{0.432717in}}%
\pgfpathlineto{\pgfqpoint{3.710166in}{0.483265in}}%
\pgfpathlineto{\pgfqpoint{3.712435in}{0.479169in}}%
\pgfpathlineto{\pgfqpoint{3.714705in}{0.472070in}}%
\pgfpathlineto{\pgfqpoint{3.716974in}{0.418506in}}%
\pgfpathlineto{\pgfqpoint{3.719244in}{0.503792in}}%
\pgfpathlineto{\pgfqpoint{3.721513in}{0.522858in}}%
\pgfpathlineto{\pgfqpoint{3.723783in}{0.419604in}}%
\pgfpathlineto{\pgfqpoint{3.726052in}{0.458524in}}%
\pgfpathlineto{\pgfqpoint{3.728322in}{0.462996in}}%
\pgfpathlineto{\pgfqpoint{3.730591in}{0.425212in}}%
\pgfpathlineto{\pgfqpoint{3.732861in}{0.481410in}}%
\pgfpathlineto{\pgfqpoint{3.737400in}{0.478334in}}%
\pgfpathlineto{\pgfqpoint{3.739669in}{0.442652in}}%
\pgfpathlineto{\pgfqpoint{3.741939in}{0.499092in}}%
\pgfpathlineto{\pgfqpoint{3.744208in}{0.524052in}}%
\pgfpathlineto{\pgfqpoint{3.746478in}{0.444941in}}%
\pgfpathlineto{\pgfqpoint{3.748747in}{0.470355in}}%
\pgfpathlineto{\pgfqpoint{3.751017in}{0.451493in}}%
\pgfpathlineto{\pgfqpoint{3.753286in}{0.413296in}}%
\pgfpathlineto{\pgfqpoint{3.755556in}{0.478830in}}%
\pgfpathlineto{\pgfqpoint{3.757825in}{0.480920in}}%
\pgfpathlineto{\pgfqpoint{3.760095in}{0.483826in}}%
\pgfpathlineto{\pgfqpoint{3.762364in}{0.457462in}}%
\pgfpathlineto{\pgfqpoint{3.766903in}{0.524565in}}%
\pgfpathlineto{\pgfqpoint{3.769172in}{0.458598in}}%
\pgfpathlineto{\pgfqpoint{3.771442in}{0.479734in}}%
\pgfpathlineto{\pgfqpoint{3.775981in}{0.390961in}}%
\pgfpathlineto{\pgfqpoint{3.778250in}{0.475177in}}%
\pgfpathlineto{\pgfqpoint{3.782789in}{0.488611in}}%
\pgfpathlineto{\pgfqpoint{3.785059in}{0.468550in}}%
\pgfpathlineto{\pgfqpoint{3.787328in}{0.486128in}}%
\pgfpathlineto{\pgfqpoint{3.789598in}{0.524230in}}%
\pgfpathlineto{\pgfqpoint{3.791867in}{0.467301in}}%
\pgfpathlineto{\pgfqpoint{3.794137in}{0.487367in}}%
\pgfpathlineto{\pgfqpoint{3.798674in}{0.315625in}}%
\pgfpathmoveto{\pgfqpoint{3.798677in}{0.315625in}}%
\pgfpathlineto{\pgfqpoint{3.800945in}{0.470271in}}%
\pgfpathlineto{\pgfqpoint{3.805484in}{0.492712in}}%
\pgfpathlineto{\pgfqpoint{3.807754in}{0.477150in}}%
\pgfpathlineto{\pgfqpoint{3.810023in}{0.476483in}}%
\pgfpathlineto{\pgfqpoint{3.812293in}{0.522888in}}%
\pgfpathlineto{\pgfqpoint{3.814562in}{0.512728in}}%
\pgfpathlineto{\pgfqpoint{3.816832in}{0.487571in}}%
\pgfpathlineto{\pgfqpoint{3.819101in}{0.496319in}}%
\pgfpathlineto{\pgfqpoint{3.821371in}{0.470538in}}%
\pgfpathlineto{\pgfqpoint{3.823640in}{0.516055in}}%
\pgfpathlineto{\pgfqpoint{3.825910in}{0.453968in}}%
\pgfpathlineto{\pgfqpoint{3.828179in}{0.491511in}}%
\pgfpathlineto{\pgfqpoint{3.830448in}{0.444729in}}%
\pgfpathlineto{\pgfqpoint{3.832718in}{0.421197in}}%
\pgfpathlineto{\pgfqpoint{3.834987in}{0.445202in}}%
\pgfpathlineto{\pgfqpoint{3.837257in}{0.495352in}}%
\pgfpathlineto{\pgfqpoint{3.839526in}{0.504194in}}%
\pgfpathlineto{\pgfqpoint{3.841796in}{0.480285in}}%
\pgfpathlineto{\pgfqpoint{3.844065in}{0.493089in}}%
\pgfpathlineto{\pgfqpoint{3.846335in}{0.534614in}}%
\pgfpathlineto{\pgfqpoint{3.848604in}{0.501706in}}%
\pgfpathlineto{\pgfqpoint{3.850874in}{0.501093in}}%
\pgfpathlineto{\pgfqpoint{3.853143in}{0.474960in}}%
\pgfpathlineto{\pgfqpoint{3.855393in}{0.315625in}}%
\pgfpathmoveto{\pgfqpoint{3.855431in}{0.315625in}}%
\pgfpathlineto{\pgfqpoint{3.857682in}{0.484383in}}%
\pgfpathlineto{\pgfqpoint{3.859952in}{0.448585in}}%
\pgfpathlineto{\pgfqpoint{3.862221in}{0.487128in}}%
\pgfpathlineto{\pgfqpoint{3.864491in}{0.476720in}}%
\pgfpathlineto{\pgfqpoint{3.866760in}{0.443308in}}%
\pgfpathlineto{\pgfqpoint{3.869030in}{0.446676in}}%
\pgfpathlineto{\pgfqpoint{3.871299in}{0.512918in}}%
\pgfpathlineto{\pgfqpoint{3.873569in}{0.515297in}}%
\pgfpathlineto{\pgfqpoint{3.878108in}{0.381059in}}%
\pgfpathlineto{\pgfqpoint{3.880377in}{0.484425in}}%
\pgfpathlineto{\pgfqpoint{3.882647in}{0.441709in}}%
\pgfpathlineto{\pgfqpoint{3.884916in}{0.485902in}}%
\pgfpathlineto{\pgfqpoint{3.887186in}{0.475530in}}%
\pgfpathlineto{\pgfqpoint{3.889455in}{0.458096in}}%
\pgfpathlineto{\pgfqpoint{3.891725in}{0.373441in}}%
\pgfpathlineto{\pgfqpoint{3.893994in}{0.508422in}}%
\pgfpathlineto{\pgfqpoint{3.896264in}{0.518450in}}%
\pgfpathlineto{\pgfqpoint{3.898533in}{0.376573in}}%
\pgfpathlineto{\pgfqpoint{3.903072in}{0.477184in}}%
\pgfpathlineto{\pgfqpoint{3.905341in}{0.440259in}}%
\pgfpathlineto{\pgfqpoint{3.907611in}{0.486186in}}%
\pgfpathlineto{\pgfqpoint{3.909880in}{0.478337in}}%
\pgfpathlineto{\pgfqpoint{3.912150in}{0.466005in}}%
\pgfpathlineto{\pgfqpoint{3.914419in}{0.346934in}}%
\pgfpathlineto{\pgfqpoint{3.916689in}{0.505668in}}%
\pgfpathlineto{\pgfqpoint{3.918958in}{0.520933in}}%
\pgfpathlineto{\pgfqpoint{3.921228in}{0.394631in}}%
\pgfpathlineto{\pgfqpoint{3.923497in}{0.447275in}}%
\pgfpathlineto{\pgfqpoint{3.925767in}{0.469879in}}%
\pgfpathlineto{\pgfqpoint{3.928036in}{0.435002in}}%
\pgfpathlineto{\pgfqpoint{3.930306in}{0.484884in}}%
\pgfpathlineto{\pgfqpoint{3.934845in}{0.473456in}}%
\pgfpathlineto{\pgfqpoint{3.937114in}{0.417780in}}%
\pgfpathlineto{\pgfqpoint{3.939384in}{0.501587in}}%
\pgfpathlineto{\pgfqpoint{3.941653in}{0.522670in}}%
\pgfpathlineto{\pgfqpoint{3.943923in}{0.435038in}}%
\pgfpathlineto{\pgfqpoint{3.946192in}{0.462234in}}%
\pgfpathlineto{\pgfqpoint{3.948462in}{0.460071in}}%
\pgfpathlineto{\pgfqpoint{3.950731in}{0.426941in}}%
\pgfpathlineto{\pgfqpoint{3.953001in}{0.482870in}}%
\pgfpathlineto{\pgfqpoint{3.955270in}{0.480950in}}%
\pgfpathlineto{\pgfqpoint{3.957540in}{0.479863in}}%
\pgfpathlineto{\pgfqpoint{3.959809in}{0.442348in}}%
\pgfpathlineto{\pgfqpoint{3.962079in}{0.496675in}}%
\pgfpathlineto{\pgfqpoint{3.964348in}{0.523720in}}%
\pgfpathlineto{\pgfqpoint{3.966617in}{0.452998in}}%
\pgfpathlineto{\pgfqpoint{3.968887in}{0.473550in}}%
\pgfpathlineto{\pgfqpoint{3.971156in}{0.447029in}}%
\pgfpathlineto{\pgfqpoint{3.973426in}{0.413435in}}%
\pgfpathlineto{\pgfqpoint{3.975695in}{0.479891in}}%
\pgfpathlineto{\pgfqpoint{3.977965in}{0.481589in}}%
\pgfpathlineto{\pgfqpoint{3.980234in}{0.485419in}}%
\pgfpathlineto{\pgfqpoint{3.982504in}{0.458077in}}%
\pgfpathlineto{\pgfqpoint{3.987043in}{0.523958in}}%
\pgfpathlineto{\pgfqpoint{3.989312in}{0.507528in}}%
\pgfpathlineto{\pgfqpoint{3.991582in}{0.453564in}}%
\pgfpathlineto{\pgfqpoint{3.993851in}{0.511201in}}%
\pgfpathlineto{\pgfqpoint{3.996121in}{0.481175in}}%
\pgfpathlineto{\pgfqpoint{3.998390in}{0.518919in}}%
\pgfpathlineto{\pgfqpoint{4.000660in}{0.446723in}}%
\pgfpathlineto{\pgfqpoint{4.002929in}{0.485417in}}%
\pgfpathlineto{\pgfqpoint{4.005199in}{0.422752in}}%
\pgfpathlineto{\pgfqpoint{4.007468in}{0.334485in}}%
\pgfpathlineto{\pgfqpoint{4.009738in}{0.465064in}}%
\pgfpathlineto{\pgfqpoint{4.012007in}{0.497463in}}%
\pgfpathlineto{\pgfqpoint{4.014277in}{0.496741in}}%
\pgfpathlineto{\pgfqpoint{4.016546in}{0.452726in}}%
\pgfpathlineto{\pgfqpoint{4.018816in}{0.505616in}}%
\pgfpathlineto{\pgfqpoint{4.021085in}{0.535106in}}%
\pgfpathlineto{\pgfqpoint{4.023355in}{0.493610in}}%
\pgfpathlineto{\pgfqpoint{4.025624in}{0.481604in}}%
\pgfpathlineto{\pgfqpoint{4.027894in}{0.493459in}}%
\pgfpathlineto{\pgfqpoint{4.030163in}{0.416682in}}%
\pgfpathlineto{\pgfqpoint{4.032432in}{0.493093in}}%
\pgfpathlineto{\pgfqpoint{4.034702in}{0.452016in}}%
\pgfpathlineto{\pgfqpoint{4.036971in}{0.485687in}}%
\pgfpathlineto{\pgfqpoint{4.039241in}{0.470388in}}%
\pgfpathlineto{\pgfqpoint{4.041510in}{0.400009in}}%
\pgfpathlineto{\pgfqpoint{4.043780in}{0.471529in}}%
\pgfpathlineto{\pgfqpoint{4.046049in}{0.514925in}}%
\pgfpathlineto{\pgfqpoint{4.048319in}{0.506048in}}%
\pgfpathlineto{\pgfqpoint{4.050588in}{0.474244in}}%
\pgfpathlineto{\pgfqpoint{4.052858in}{0.411533in}}%
\pgfpathlineto{\pgfqpoint{4.055127in}{0.493727in}}%
\pgfpathlineto{\pgfqpoint{4.057397in}{0.447186in}}%
\pgfpathlineto{\pgfqpoint{4.059666in}{0.486162in}}%
\pgfpathlineto{\pgfqpoint{4.061936in}{0.470530in}}%
\pgfpathlineto{\pgfqpoint{4.064205in}{0.434364in}}%
\pgfpathlineto{\pgfqpoint{4.066475in}{0.446418in}}%
\pgfpathlineto{\pgfqpoint{4.068744in}{0.511775in}}%
\pgfpathlineto{\pgfqpoint{4.071014in}{0.511437in}}%
\pgfpathlineto{\pgfqpoint{4.073283in}{0.450046in}}%
\pgfpathlineto{\pgfqpoint{4.075553in}{0.343646in}}%
\pgfpathlineto{\pgfqpoint{4.077822in}{0.488400in}}%
\pgfpathlineto{\pgfqpoint{4.080092in}{0.448424in}}%
\pgfpathlineto{\pgfqpoint{4.082361in}{0.487837in}}%
\pgfpathlineto{\pgfqpoint{4.084631in}{0.474229in}}%
\pgfpathlineto{\pgfqpoint{4.086900in}{0.446357in}}%
\pgfpathlineto{\pgfqpoint{4.089170in}{0.435039in}}%
\pgfpathlineto{\pgfqpoint{4.091439in}{0.509957in}}%
\pgfpathlineto{\pgfqpoint{4.093709in}{0.515210in}}%
\pgfpathlineto{\pgfqpoint{4.095978in}{0.420395in}}%
\pgfpathlineto{\pgfqpoint{4.098247in}{0.391129in}}%
\pgfpathlineto{\pgfqpoint{4.100517in}{0.483270in}}%
\pgfpathlineto{\pgfqpoint{4.102786in}{0.446614in}}%
\pgfpathlineto{\pgfqpoint{4.105056in}{0.488368in}}%
\pgfpathlineto{\pgfqpoint{4.107325in}{0.476651in}}%
\pgfpathlineto{\pgfqpoint{4.109595in}{0.457598in}}%
\pgfpathlineto{\pgfqpoint{4.111864in}{0.403813in}}%
\pgfpathlineto{\pgfqpoint{4.114134in}{0.507166in}}%
\pgfpathlineto{\pgfqpoint{4.116403in}{0.518318in}}%
\pgfpathlineto{\pgfqpoint{4.118312in}{0.315625in}}%
\pgfpathmoveto{\pgfqpoint{4.119243in}{0.315625in}}%
\pgfpathlineto{\pgfqpoint{4.120942in}{0.429742in}}%
\pgfpathlineto{\pgfqpoint{4.123212in}{0.476584in}}%
\pgfpathlineto{\pgfqpoint{4.125481in}{0.443401in}}%
\pgfpathlineto{\pgfqpoint{4.127751in}{0.488119in}}%
\pgfpathlineto{\pgfqpoint{4.130020in}{0.478841in}}%
\pgfpathlineto{\pgfqpoint{4.132290in}{0.466735in}}%
\pgfpathlineto{\pgfqpoint{4.134398in}{0.315625in}}%
\pgfpathmoveto{\pgfqpoint{4.134691in}{0.315625in}}%
\pgfpathlineto{\pgfqpoint{4.136829in}{0.503827in}}%
\pgfpathlineto{\pgfqpoint{4.139098in}{0.520710in}}%
\pgfpathlineto{\pgfqpoint{4.141368in}{0.418634in}}%
\pgfpathlineto{\pgfqpoint{4.143637in}{0.450649in}}%
\pgfpathlineto{\pgfqpoint{4.145907in}{0.468267in}}%
\pgfpathlineto{\pgfqpoint{4.148176in}{0.437865in}}%
\pgfpathlineto{\pgfqpoint{4.150446in}{0.486748in}}%
\pgfpathlineto{\pgfqpoint{4.154985in}{0.474500in}}%
\pgfpathlineto{\pgfqpoint{4.157254in}{0.413260in}}%
\pgfpathlineto{\pgfqpoint{4.159524in}{0.499620in}}%
\pgfpathlineto{\pgfqpoint{4.161793in}{0.522360in}}%
\pgfpathlineto{\pgfqpoint{4.164062in}{0.444833in}}%
\pgfpathlineto{\pgfqpoint{4.166332in}{0.465147in}}%
\pgfpathlineto{\pgfqpoint{4.168601in}{0.457467in}}%
\pgfpathlineto{\pgfqpoint{4.170871in}{0.429041in}}%
\pgfpathlineto{\pgfqpoint{4.173140in}{0.484454in}}%
\pgfpathlineto{\pgfqpoint{4.175410in}{0.481665in}}%
\pgfpathlineto{\pgfqpoint{4.177679in}{0.481119in}}%
\pgfpathlineto{\pgfqpoint{4.179949in}{0.441228in}}%
\pgfpathlineto{\pgfqpoint{4.182218in}{0.494318in}}%
\pgfpathlineto{\pgfqpoint{4.184488in}{0.523223in}}%
\pgfpathlineto{\pgfqpoint{4.186757in}{0.503577in}}%
\pgfpathlineto{\pgfqpoint{4.189027in}{0.426675in}}%
\pgfpathlineto{\pgfqpoint{4.191296in}{0.514631in}}%
\pgfpathlineto{\pgfqpoint{4.193566in}{0.483894in}}%
\pgfpathlineto{\pgfqpoint{4.195835in}{0.518763in}}%
\pgfpathlineto{\pgfqpoint{4.198105in}{0.443705in}}%
\pgfpathlineto{\pgfqpoint{4.200374in}{0.481983in}}%
\pgfpathlineto{\pgfqpoint{4.202644in}{0.403886in}}%
\pgfpathlineto{\pgfqpoint{4.204913in}{0.371012in}}%
\pgfpathlineto{\pgfqpoint{4.207183in}{0.472893in}}%
\pgfpathlineto{\pgfqpoint{4.209452in}{0.497755in}}%
\pgfpathlineto{\pgfqpoint{4.211722in}{0.492369in}}%
\pgfpathlineto{\pgfqpoint{4.213991in}{0.426361in}}%
\pgfpathlineto{\pgfqpoint{4.216261in}{0.508298in}}%
\pgfpathlineto{\pgfqpoint{4.218530in}{0.533855in}}%
\pgfpathlineto{\pgfqpoint{4.220800in}{0.488805in}}%
\pgfpathlineto{\pgfqpoint{4.223069in}{0.470592in}}%
\pgfpathlineto{\pgfqpoint{4.225339in}{0.497919in}}%
\pgfpathlineto{\pgfqpoint{4.227608in}{0.432380in}}%
\pgfpathlineto{\pgfqpoint{4.229878in}{0.496454in}}%
\pgfpathlineto{\pgfqpoint{4.232147in}{0.453862in}}%
\pgfpathlineto{\pgfqpoint{4.234416in}{0.484846in}}%
\pgfpathlineto{\pgfqpoint{4.236686in}{0.466278in}}%
\pgfpathlineto{\pgfqpoint{4.238955in}{0.320406in}}%
\pgfpathlineto{\pgfqpoint{4.241225in}{0.479620in}}%
\pgfpathlineto{\pgfqpoint{4.243494in}{0.514642in}}%
\pgfpathlineto{\pgfqpoint{4.245764in}{0.500822in}}%
\pgfpathlineto{\pgfqpoint{4.248033in}{0.481185in}}%
\pgfpathlineto{\pgfqpoint{4.250303in}{0.430924in}}%
\pgfpathlineto{\pgfqpoint{4.252572in}{0.497142in}}%
\pgfpathlineto{\pgfqpoint{4.254842in}{0.449738in}}%
\pgfpathlineto{\pgfqpoint{4.257111in}{0.486039in}}%
\pgfpathlineto{\pgfqpoint{4.259381in}{0.467437in}}%
\pgfpathlineto{\pgfqpoint{4.261650in}{0.416268in}}%
\pgfpathlineto{\pgfqpoint{4.266189in}{0.512107in}}%
\pgfpathlineto{\pgfqpoint{4.268459in}{0.507482in}}%
\pgfpathlineto{\pgfqpoint{4.272998in}{0.403617in}}%
\pgfpathlineto{\pgfqpoint{4.275267in}{0.492342in}}%
\pgfpathlineto{\pgfqpoint{4.277537in}{0.451728in}}%
\pgfpathlineto{\pgfqpoint{4.279806in}{0.488217in}}%
\pgfpathlineto{\pgfqpoint{4.282076in}{0.471734in}}%
\pgfpathlineto{\pgfqpoint{4.284345in}{0.433350in}}%
\pgfpathlineto{\pgfqpoint{4.286615in}{0.452879in}}%
\pgfpathlineto{\pgfqpoint{4.288884in}{0.510625in}}%
\pgfpathlineto{\pgfqpoint{4.291154in}{0.511948in}}%
\pgfpathlineto{\pgfqpoint{4.293423in}{0.439007in}}%
\pgfpathlineto{\pgfqpoint{4.295204in}{0.315625in}}%
\pgfpathmoveto{\pgfqpoint{4.296065in}{0.315625in}}%
\pgfpathlineto{\pgfqpoint{4.297962in}{0.487747in}}%
\pgfpathlineto{\pgfqpoint{4.300231in}{0.450834in}}%
\pgfpathlineto{\pgfqpoint{4.302501in}{0.489386in}}%
\pgfpathlineto{\pgfqpoint{4.304770in}{0.474758in}}%
\pgfpathlineto{\pgfqpoint{4.307040in}{0.448206in}}%
\pgfpathlineto{\pgfqpoint{4.309309in}{0.435252in}}%
\pgfpathlineto{\pgfqpoint{4.311579in}{0.508276in}}%
\pgfpathlineto{\pgfqpoint{4.313848in}{0.515711in}}%
\pgfpathlineto{\pgfqpoint{4.316118in}{0.388702in}}%
\pgfpathlineto{\pgfqpoint{4.318387in}{0.404896in}}%
\pgfpathlineto{\pgfqpoint{4.320657in}{0.481791in}}%
\pgfpathlineto{\pgfqpoint{4.322926in}{0.448760in}}%
\pgfpathlineto{\pgfqpoint{4.325196in}{0.489868in}}%
\pgfpathlineto{\pgfqpoint{4.327465in}{0.477483in}}%
\pgfpathlineto{\pgfqpoint{4.329735in}{0.459532in}}%
\pgfpathlineto{\pgfqpoint{4.332004in}{0.452466in}}%
\pgfpathlineto{\pgfqpoint{4.334274in}{0.438123in}}%
\pgfpathlineto{\pgfqpoint{4.336543in}{0.483734in}}%
\pgfpathlineto{\pgfqpoint{4.336543in}{0.483734in}}%
\pgfusepath{stroke}%
\end{pgfscope}%
\begin{pgfscope}%
\pgfpathrectangle{\pgfqpoint{2.520957in}{0.320625in}}{\pgfqpoint{1.815586in}{1.164375in}}%
\pgfusepath{clip}%
\pgfsetrectcap%
\pgfsetroundjoin%
\pgfsetlinewidth{1.154312pt}%
\definecolor{currentstroke}{rgb}{0.997351,0.719089,0.493755}%
\pgfsetstrokecolor{currentstroke}%
\pgfsetdash{}{0pt}%
\pgfpathmoveto{\pgfqpoint{2.520957in}{1.451682in}}%
\pgfpathlineto{\pgfqpoint{2.525496in}{1.441674in}}%
\pgfpathlineto{\pgfqpoint{2.541382in}{1.409832in}}%
\pgfpathlineto{\pgfqpoint{2.557269in}{1.375075in}}%
\pgfpathlineto{\pgfqpoint{2.589041in}{1.301926in}}%
\pgfpathlineto{\pgfqpoint{2.634431in}{1.198168in}}%
\pgfpathlineto{\pgfqpoint{2.675282in}{1.107600in}}%
\pgfpathlineto{\pgfqpoint{2.720671in}{1.009674in}}%
\pgfpathlineto{\pgfqpoint{2.775139in}{0.894916in}}%
\pgfpathlineto{\pgfqpoint{2.840954in}{0.758978in}}%
\pgfpathlineto{\pgfqpoint{2.915847in}{0.606770in}}%
\pgfpathlineto{\pgfqpoint{2.920386in}{0.597610in}}%
\pgfpathlineto{\pgfqpoint{2.940811in}{0.557403in}}%
\pgfpathlineto{\pgfqpoint{2.943081in}{0.554293in}}%
\pgfpathlineto{\pgfqpoint{2.945350in}{0.549289in}}%
\pgfpathlineto{\pgfqpoint{2.947620in}{0.542665in}}%
\pgfpathlineto{\pgfqpoint{2.949889in}{0.540098in}}%
\pgfpathlineto{\pgfqpoint{2.952159in}{0.545433in}}%
\pgfpathlineto{\pgfqpoint{2.954428in}{0.530673in}}%
\pgfpathlineto{\pgfqpoint{2.956698in}{0.529633in}}%
\pgfpathlineto{\pgfqpoint{2.958967in}{0.519774in}}%
\pgfpathlineto{\pgfqpoint{2.961236in}{0.515249in}}%
\pgfpathlineto{\pgfqpoint{2.963506in}{0.513184in}}%
\pgfpathlineto{\pgfqpoint{2.965775in}{0.510191in}}%
\pgfpathlineto{\pgfqpoint{2.968045in}{0.513565in}}%
\pgfpathlineto{\pgfqpoint{2.970314in}{0.506455in}}%
\pgfpathlineto{\pgfqpoint{2.972584in}{0.492967in}}%
\pgfpathlineto{\pgfqpoint{2.974853in}{0.514557in}}%
\pgfpathlineto{\pgfqpoint{2.977123in}{0.518914in}}%
\pgfpathlineto{\pgfqpoint{2.979392in}{0.512938in}}%
\pgfpathlineto{\pgfqpoint{2.981662in}{0.480425in}}%
\pgfpathlineto{\pgfqpoint{2.983931in}{0.480844in}}%
\pgfpathlineto{\pgfqpoint{2.986201in}{0.517839in}}%
\pgfpathlineto{\pgfqpoint{2.988470in}{0.477202in}}%
\pgfpathlineto{\pgfqpoint{2.990740in}{0.502021in}}%
\pgfpathlineto{\pgfqpoint{2.993009in}{0.464112in}}%
\pgfpathlineto{\pgfqpoint{2.995279in}{0.461595in}}%
\pgfpathlineto{\pgfqpoint{2.997548in}{0.447965in}}%
\pgfpathlineto{\pgfqpoint{2.999818in}{0.485009in}}%
\pgfpathlineto{\pgfqpoint{3.002087in}{0.506115in}}%
\pgfpathlineto{\pgfqpoint{3.004357in}{0.497999in}}%
\pgfpathlineto{\pgfqpoint{3.006626in}{0.449318in}}%
\pgfpathlineto{\pgfqpoint{3.008896in}{0.521796in}}%
\pgfpathlineto{\pgfqpoint{3.011165in}{0.511324in}}%
\pgfpathlineto{\pgfqpoint{3.013435in}{0.515475in}}%
\pgfpathlineto{\pgfqpoint{3.015704in}{0.441602in}}%
\pgfpathlineto{\pgfqpoint{3.017974in}{0.415310in}}%
\pgfpathlineto{\pgfqpoint{3.020243in}{0.487226in}}%
\pgfpathlineto{\pgfqpoint{3.022513in}{0.448662in}}%
\pgfpathlineto{\pgfqpoint{3.024782in}{0.490994in}}%
\pgfpathlineto{\pgfqpoint{3.027052in}{0.474137in}}%
\pgfpathlineto{\pgfqpoint{3.029321in}{0.465675in}}%
\pgfpathlineto{\pgfqpoint{3.031590in}{0.383355in}}%
\pgfpathlineto{\pgfqpoint{3.033860in}{0.498139in}}%
\pgfpathlineto{\pgfqpoint{3.036129in}{0.517837in}}%
\pgfpathlineto{\pgfqpoint{3.038399in}{0.443897in}}%
\pgfpathlineto{\pgfqpoint{3.040668in}{0.449767in}}%
\pgfpathlineto{\pgfqpoint{3.042938in}{0.468412in}}%
\pgfpathlineto{\pgfqpoint{3.045207in}{0.451177in}}%
\pgfpathlineto{\pgfqpoint{3.047477in}{0.493635in}}%
\pgfpathlineto{\pgfqpoint{3.052016in}{0.471455in}}%
\pgfpathlineto{\pgfqpoint{3.054285in}{0.364063in}}%
\pgfpathlineto{\pgfqpoint{3.056555in}{0.498001in}}%
\pgfpathlineto{\pgfqpoint{3.058824in}{0.520610in}}%
\pgfpathlineto{\pgfqpoint{3.061094in}{0.493063in}}%
\pgfpathlineto{\pgfqpoint{3.063363in}{0.376493in}}%
\pgfpathlineto{\pgfqpoint{3.065633in}{0.512622in}}%
\pgfpathlineto{\pgfqpoint{3.067902in}{0.484924in}}%
\pgfpathlineto{\pgfqpoint{3.070172in}{0.516157in}}%
\pgfpathlineto{\pgfqpoint{3.072441in}{0.447659in}}%
\pgfpathlineto{\pgfqpoint{3.074711in}{0.475450in}}%
\pgfpathlineto{\pgfqpoint{3.076734in}{0.315625in}}%
\pgfpathmoveto{\pgfqpoint{3.077306in}{0.315625in}}%
\pgfpathlineto{\pgfqpoint{3.079250in}{0.431352in}}%
\pgfpathlineto{\pgfqpoint{3.081519in}{0.487786in}}%
\pgfpathlineto{\pgfqpoint{3.083789in}{0.497596in}}%
\pgfpathlineto{\pgfqpoint{3.086058in}{0.483086in}}%
\pgfpathlineto{\pgfqpoint{3.088328in}{0.412617in}}%
\pgfpathlineto{\pgfqpoint{3.090597in}{0.507383in}}%
\pgfpathlineto{\pgfqpoint{3.092867in}{0.529415in}}%
\pgfpathlineto{\pgfqpoint{3.095136in}{0.480906in}}%
\pgfpathlineto{\pgfqpoint{3.097405in}{0.458927in}}%
\pgfpathlineto{\pgfqpoint{3.099675in}{0.495801in}}%
\pgfpathlineto{\pgfqpoint{3.101944in}{0.449052in}}%
\pgfpathlineto{\pgfqpoint{3.104214in}{0.500495in}}%
\pgfpathlineto{\pgfqpoint{3.106483in}{0.461278in}}%
\pgfpathlineto{\pgfqpoint{3.108753in}{0.483924in}}%
\pgfpathlineto{\pgfqpoint{3.113292in}{0.421816in}}%
\pgfpathlineto{\pgfqpoint{3.115561in}{0.493734in}}%
\pgfpathlineto{\pgfqpoint{3.117831in}{0.511618in}}%
\pgfpathlineto{\pgfqpoint{3.122370in}{0.469847in}}%
\pgfpathlineto{\pgfqpoint{3.124639in}{0.507515in}}%
\pgfpathlineto{\pgfqpoint{3.126909in}{0.533291in}}%
\pgfpathlineto{\pgfqpoint{3.129178in}{0.463146in}}%
\pgfpathlineto{\pgfqpoint{3.131448in}{0.488397in}}%
\pgfpathlineto{\pgfqpoint{3.133717in}{0.449930in}}%
\pgfpathlineto{\pgfqpoint{3.135944in}{0.315625in}}%
\pgfpathmoveto{\pgfqpoint{3.136024in}{0.315625in}}%
\pgfpathlineto{\pgfqpoint{3.138256in}{0.470169in}}%
\pgfpathlineto{\pgfqpoint{3.140526in}{0.468522in}}%
\pgfpathlineto{\pgfqpoint{3.142795in}{0.489815in}}%
\pgfpathlineto{\pgfqpoint{3.145065in}{0.483469in}}%
\pgfpathlineto{\pgfqpoint{3.147334in}{0.442182in}}%
\pgfpathlineto{\pgfqpoint{3.149604in}{0.511362in}}%
\pgfpathlineto{\pgfqpoint{3.151873in}{0.481858in}}%
\pgfpathlineto{\pgfqpoint{3.154143in}{0.497013in}}%
\pgfpathlineto{\pgfqpoint{3.158682in}{0.329121in}}%
\pgfpathlineto{\pgfqpoint{3.160951in}{0.476750in}}%
\pgfpathlineto{\pgfqpoint{3.163220in}{0.486432in}}%
\pgfpathlineto{\pgfqpoint{3.165490in}{0.500367in}}%
\pgfpathlineto{\pgfqpoint{3.170029in}{0.450989in}}%
\pgfpathlineto{\pgfqpoint{3.172298in}{0.514513in}}%
\pgfpathlineto{\pgfqpoint{3.174568in}{0.505202in}}%
\pgfpathlineto{\pgfqpoint{3.176837in}{0.505825in}}%
\pgfpathlineto{\pgfqpoint{3.179107in}{0.422213in}}%
\pgfpathlineto{\pgfqpoint{3.181376in}{0.427219in}}%
\pgfpathlineto{\pgfqpoint{3.183646in}{0.501786in}}%
\pgfpathlineto{\pgfqpoint{3.185915in}{0.470280in}}%
\pgfpathlineto{\pgfqpoint{3.188185in}{0.498776in}}%
\pgfpathlineto{\pgfqpoint{3.192724in}{0.434018in}}%
\pgfpathlineto{\pgfqpoint{3.194993in}{0.458293in}}%
\pgfpathlineto{\pgfqpoint{3.197263in}{0.494716in}}%
\pgfpathlineto{\pgfqpoint{3.199532in}{0.509756in}}%
\pgfpathlineto{\pgfqpoint{3.204071in}{0.446371in}}%
\pgfpathlineto{\pgfqpoint{3.206341in}{0.518370in}}%
\pgfpathlineto{\pgfqpoint{3.208610in}{0.497883in}}%
\pgfpathlineto{\pgfqpoint{3.210880in}{0.512462in}}%
\pgfpathlineto{\pgfqpoint{3.213149in}{0.386121in}}%
\pgfpathlineto{\pgfqpoint{3.217688in}{0.470405in}}%
\pgfpathlineto{\pgfqpoint{3.219958in}{0.460776in}}%
\pgfpathlineto{\pgfqpoint{3.222227in}{0.496501in}}%
\pgfpathlineto{\pgfqpoint{3.224497in}{0.483722in}}%
\pgfpathlineto{\pgfqpoint{3.226766in}{0.457004in}}%
\pgfpathlineto{\pgfqpoint{3.229035in}{0.452742in}}%
\pgfpathlineto{\pgfqpoint{3.231305in}{0.504066in}}%
\pgfpathlineto{\pgfqpoint{3.233574in}{0.519224in}}%
\pgfpathlineto{\pgfqpoint{3.235844in}{0.476206in}}%
\pgfpathlineto{\pgfqpoint{3.238113in}{0.402592in}}%
\pgfpathlineto{\pgfqpoint{3.240383in}{0.513438in}}%
\pgfpathlineto{\pgfqpoint{3.242652in}{0.481144in}}%
\pgfpathlineto{\pgfqpoint{3.244922in}{0.511923in}}%
\pgfpathlineto{\pgfqpoint{3.247191in}{0.448428in}}%
\pgfpathlineto{\pgfqpoint{3.249461in}{0.468236in}}%
\pgfpathlineto{\pgfqpoint{3.251730in}{0.379116in}}%
\pgfpathlineto{\pgfqpoint{3.254000in}{0.452970in}}%
\pgfpathlineto{\pgfqpoint{3.256269in}{0.496522in}}%
\pgfpathlineto{\pgfqpoint{3.258539in}{0.499054in}}%
\pgfpathlineto{\pgfqpoint{3.260808in}{0.469093in}}%
\pgfpathlineto{\pgfqpoint{3.263078in}{0.470763in}}%
\pgfpathlineto{\pgfqpoint{3.265347in}{0.509760in}}%
\pgfpathlineto{\pgfqpoint{3.267617in}{0.526613in}}%
\pgfpathlineto{\pgfqpoint{3.269886in}{0.462654in}}%
\pgfpathlineto{\pgfqpoint{3.272156in}{0.445726in}}%
\pgfpathlineto{\pgfqpoint{3.274425in}{0.493284in}}%
\pgfpathlineto{\pgfqpoint{3.276695in}{0.451183in}}%
\pgfpathlineto{\pgfqpoint{3.278964in}{0.499125in}}%
\pgfpathlineto{\pgfqpoint{3.281234in}{0.465999in}}%
\pgfpathlineto{\pgfqpoint{3.283503in}{0.480607in}}%
\pgfpathlineto{\pgfqpoint{3.285773in}{0.446379in}}%
\pgfpathlineto{\pgfqpoint{3.288042in}{0.455437in}}%
\pgfpathlineto{\pgfqpoint{3.290312in}{0.504864in}}%
\pgfpathlineto{\pgfqpoint{3.292581in}{0.510699in}}%
\pgfpathlineto{\pgfqpoint{3.294851in}{0.475726in}}%
\pgfpathlineto{\pgfqpoint{3.297120in}{0.492780in}}%
\pgfpathlineto{\pgfqpoint{3.299389in}{0.505418in}}%
\pgfpathlineto{\pgfqpoint{3.301659in}{0.529277in}}%
\pgfpathlineto{\pgfqpoint{3.303928in}{0.452325in}}%
\pgfpathlineto{\pgfqpoint{3.306198in}{0.482378in}}%
\pgfpathlineto{\pgfqpoint{3.308467in}{0.443004in}}%
\pgfpathlineto{\pgfqpoint{3.310737in}{0.377016in}}%
\pgfpathlineto{\pgfqpoint{3.313006in}{0.474892in}}%
\pgfpathlineto{\pgfqpoint{3.315276in}{0.476639in}}%
\pgfpathlineto{\pgfqpoint{3.317545in}{0.489589in}}%
\pgfpathlineto{\pgfqpoint{3.319815in}{0.475194in}}%
\pgfpathlineto{\pgfqpoint{3.322084in}{0.472062in}}%
\pgfpathlineto{\pgfqpoint{3.324354in}{0.519585in}}%
\pgfpathlineto{\pgfqpoint{3.326623in}{0.514726in}}%
\pgfpathlineto{\pgfqpoint{3.328893in}{0.482252in}}%
\pgfpathlineto{\pgfqpoint{3.331162in}{0.501713in}}%
\pgfpathlineto{\pgfqpoint{3.333432in}{0.483215in}}%
\pgfpathlineto{\pgfqpoint{3.335701in}{0.522112in}}%
\pgfpathlineto{\pgfqpoint{3.337971in}{0.450945in}}%
\pgfpathlineto{\pgfqpoint{3.340240in}{0.490999in}}%
\pgfpathlineto{\pgfqpoint{3.342510in}{0.426167in}}%
\pgfpathlineto{\pgfqpoint{3.344779in}{0.411829in}}%
\pgfpathlineto{\pgfqpoint{3.349318in}{0.490661in}}%
\pgfpathlineto{\pgfqpoint{3.351588in}{0.500646in}}%
\pgfpathlineto{\pgfqpoint{3.353857in}{0.483072in}}%
\pgfpathlineto{\pgfqpoint{3.356127in}{0.492405in}}%
\pgfpathlineto{\pgfqpoint{3.358396in}{0.533387in}}%
\pgfpathlineto{\pgfqpoint{3.360666in}{0.506264in}}%
\pgfpathlineto{\pgfqpoint{3.362935in}{0.495159in}}%
\pgfpathlineto{\pgfqpoint{3.365204in}{0.489335in}}%
\pgfpathlineto{\pgfqpoint{3.367474in}{0.410257in}}%
\pgfpathlineto{\pgfqpoint{3.369743in}{0.495408in}}%
\pgfpathlineto{\pgfqpoint{3.372013in}{0.447773in}}%
\pgfpathlineto{\pgfqpoint{3.374282in}{0.487331in}}%
\pgfpathlineto{\pgfqpoint{3.376552in}{0.469388in}}%
\pgfpathlineto{\pgfqpoint{3.378821in}{0.439132in}}%
\pgfpathlineto{\pgfqpoint{3.381091in}{0.437790in}}%
\pgfpathlineto{\pgfqpoint{3.383360in}{0.509118in}}%
\pgfpathlineto{\pgfqpoint{3.385630in}{0.512177in}}%
\pgfpathlineto{\pgfqpoint{3.387899in}{0.429045in}}%
\pgfpathlineto{\pgfqpoint{3.389545in}{0.315625in}}%
\pgfpathmoveto{\pgfqpoint{3.390624in}{0.315625in}}%
\pgfpathlineto{\pgfqpoint{3.392438in}{0.486859in}}%
\pgfpathlineto{\pgfqpoint{3.394708in}{0.452162in}}%
\pgfpathlineto{\pgfqpoint{3.396977in}{0.489935in}}%
\pgfpathlineto{\pgfqpoint{3.399247in}{0.475359in}}%
\pgfpathlineto{\pgfqpoint{3.401516in}{0.447320in}}%
\pgfpathlineto{\pgfqpoint{3.403786in}{0.440169in}}%
\pgfpathlineto{\pgfqpoint{3.406055in}{0.508352in}}%
\pgfpathlineto{\pgfqpoint{3.408325in}{0.515675in}}%
\pgfpathlineto{\pgfqpoint{3.410594in}{0.388108in}}%
\pgfpathlineto{\pgfqpoint{3.412864in}{0.404440in}}%
\pgfpathlineto{\pgfqpoint{3.415133in}{0.482009in}}%
\pgfpathlineto{\pgfqpoint{3.417403in}{0.449136in}}%
\pgfpathlineto{\pgfqpoint{3.419672in}{0.490091in}}%
\pgfpathlineto{\pgfqpoint{3.421942in}{0.477366in}}%
\pgfpathlineto{\pgfqpoint{3.424211in}{0.459380in}}%
\pgfpathlineto{\pgfqpoint{3.426481in}{0.405455in}}%
\pgfpathlineto{\pgfqpoint{3.428750in}{0.505191in}}%
\pgfpathlineto{\pgfqpoint{3.431019in}{0.518535in}}%
\pgfpathlineto{\pgfqpoint{3.433289in}{0.394921in}}%
\pgfpathlineto{\pgfqpoint{3.437828in}{0.474630in}}%
\pgfpathlineto{\pgfqpoint{3.440097in}{0.445589in}}%
\pgfpathlineto{\pgfqpoint{3.442367in}{0.489748in}}%
\pgfpathlineto{\pgfqpoint{3.446906in}{0.468438in}}%
\pgfpathlineto{\pgfqpoint{3.448748in}{0.315625in}}%
\pgfpathmoveto{\pgfqpoint{3.449538in}{0.315625in}}%
\pgfpathlineto{\pgfqpoint{3.451445in}{0.501752in}}%
\pgfpathlineto{\pgfqpoint{3.453714in}{0.520700in}}%
\pgfpathlineto{\pgfqpoint{3.455984in}{0.492319in}}%
\pgfpathlineto{\pgfqpoint{3.458253in}{0.436582in}}%
\pgfpathlineto{\pgfqpoint{3.460523in}{0.524113in}}%
\pgfpathlineto{\pgfqpoint{3.462792in}{0.487435in}}%
\pgfpathlineto{\pgfqpoint{3.465062in}{0.516738in}}%
\pgfpathlineto{\pgfqpoint{3.467331in}{0.422404in}}%
\pgfpathlineto{\pgfqpoint{3.469601in}{0.470964in}}%
\pgfpathlineto{\pgfqpoint{3.471870in}{0.361132in}}%
\pgfpathlineto{\pgfqpoint{3.476409in}{0.482843in}}%
\pgfpathlineto{\pgfqpoint{3.478679in}{0.497370in}}%
\pgfpathlineto{\pgfqpoint{3.480948in}{0.479097in}}%
\pgfpathlineto{\pgfqpoint{3.483218in}{0.424052in}}%
\pgfpathlineto{\pgfqpoint{3.485487in}{0.515043in}}%
\pgfpathlineto{\pgfqpoint{3.487757in}{0.530552in}}%
\pgfpathlineto{\pgfqpoint{3.492296in}{0.405594in}}%
\pgfpathlineto{\pgfqpoint{3.494565in}{0.510108in}}%
\pgfpathlineto{\pgfqpoint{3.496834in}{0.450340in}}%
\pgfpathlineto{\pgfqpoint{3.499104in}{0.502336in}}%
\pgfpathlineto{\pgfqpoint{3.501373in}{0.452180in}}%
\pgfpathlineto{\pgfqpoint{3.503643in}{0.479420in}}%
\pgfpathlineto{\pgfqpoint{3.508182in}{0.420839in}}%
\pgfpathlineto{\pgfqpoint{3.510451in}{0.490968in}}%
\pgfpathlineto{\pgfqpoint{3.512721in}{0.514107in}}%
\pgfpathlineto{\pgfqpoint{3.514990in}{0.483445in}}%
\pgfpathlineto{\pgfqpoint{3.517260in}{0.498061in}}%
\pgfpathlineto{\pgfqpoint{3.519529in}{0.454437in}}%
\pgfpathlineto{\pgfqpoint{3.521799in}{0.503757in}}%
\pgfpathlineto{\pgfqpoint{3.524068in}{0.449152in}}%
\pgfpathlineto{\pgfqpoint{3.526338in}{0.482014in}}%
\pgfpathlineto{\pgfqpoint{3.528607in}{0.456639in}}%
\pgfpathlineto{\pgfqpoint{3.530877in}{0.352027in}}%
\pgfpathlineto{\pgfqpoint{3.533146in}{0.479110in}}%
\pgfpathlineto{\pgfqpoint{3.535416in}{0.512875in}}%
\pgfpathlineto{\pgfqpoint{3.537685in}{0.494537in}}%
\pgfpathlineto{\pgfqpoint{3.539955in}{0.484681in}}%
\pgfpathlineto{\pgfqpoint{3.542224in}{0.445469in}}%
\pgfpathlineto{\pgfqpoint{3.544494in}{0.500342in}}%
\pgfpathlineto{\pgfqpoint{3.546763in}{0.453505in}}%
\pgfpathlineto{\pgfqpoint{3.549033in}{0.485260in}}%
\pgfpathlineto{\pgfqpoint{3.551302in}{0.462417in}}%
\pgfpathlineto{\pgfqpoint{3.553572in}{0.345076in}}%
\pgfpathlineto{\pgfqpoint{3.555841in}{0.476001in}}%
\pgfpathlineto{\pgfqpoint{3.558111in}{0.512163in}}%
\pgfpathlineto{\pgfqpoint{3.560380in}{0.501012in}}%
\pgfpathlineto{\pgfqpoint{3.562649in}{0.473042in}}%
\pgfpathlineto{\pgfqpoint{3.564919in}{0.432505in}}%
\pgfpathlineto{\pgfqpoint{3.567188in}{0.497249in}}%
\pgfpathlineto{\pgfqpoint{3.569458in}{0.454896in}}%
\pgfpathlineto{\pgfqpoint{3.571727in}{0.487780in}}%
\pgfpathlineto{\pgfqpoint{3.573997in}{0.466962in}}%
\pgfpathlineto{\pgfqpoint{3.576266in}{0.404475in}}%
\pgfpathlineto{\pgfqpoint{3.578536in}{0.468359in}}%
\pgfpathlineto{\pgfqpoint{3.580805in}{0.510864in}}%
\pgfpathlineto{\pgfqpoint{3.583075in}{0.506695in}}%
\pgfpathlineto{\pgfqpoint{3.587614in}{0.407958in}}%
\pgfpathlineto{\pgfqpoint{3.589883in}{0.493193in}}%
\pgfpathlineto{\pgfqpoint{3.592153in}{0.455450in}}%
\pgfpathlineto{\pgfqpoint{3.594422in}{0.489820in}}%
\pgfpathlineto{\pgfqpoint{3.596692in}{0.471094in}}%
\pgfpathlineto{\pgfqpoint{3.598961in}{0.429166in}}%
\pgfpathlineto{\pgfqpoint{3.601231in}{0.458598in}}%
\pgfpathlineto{\pgfqpoint{3.603500in}{0.509210in}}%
\pgfpathlineto{\pgfqpoint{3.605770in}{0.511327in}}%
\pgfpathlineto{\pgfqpoint{3.608039in}{0.428516in}}%
\pgfpathlineto{\pgfqpoint{3.610219in}{0.315625in}}%
\pgfpathmoveto{\pgfqpoint{3.610368in}{0.315625in}}%
\pgfpathlineto{\pgfqpoint{3.612578in}{0.488297in}}%
\pgfpathlineto{\pgfqpoint{3.614848in}{0.454693in}}%
\pgfpathlineto{\pgfqpoint{3.617117in}{0.491197in}}%
\pgfpathlineto{\pgfqpoint{3.619387in}{0.474609in}}%
\pgfpathlineto{\pgfqpoint{3.621656in}{0.445698in}}%
\pgfpathlineto{\pgfqpoint{3.623926in}{0.444051in}}%
\pgfpathlineto{\pgfqpoint{3.626195in}{0.507027in}}%
\pgfpathlineto{\pgfqpoint{3.628465in}{0.515090in}}%
\pgfpathlineto{\pgfqpoint{3.630734in}{0.348355in}}%
\pgfpathlineto{\pgfqpoint{3.633003in}{0.403506in}}%
\pgfpathlineto{\pgfqpoint{3.635273in}{0.482283in}}%
\pgfpathlineto{\pgfqpoint{3.637542in}{0.452553in}}%
\pgfpathlineto{\pgfqpoint{3.639812in}{0.491810in}}%
\pgfpathlineto{\pgfqpoint{3.642081in}{0.477553in}}%
\pgfpathlineto{\pgfqpoint{3.644351in}{0.458105in}}%
\pgfpathlineto{\pgfqpoint{3.646620in}{0.419580in}}%
\pgfpathlineto{\pgfqpoint{3.648890in}{0.504226in}}%
\pgfpathlineto{\pgfqpoint{3.651159in}{0.518043in}}%
\pgfpathlineto{\pgfqpoint{3.653429in}{0.482798in}}%
\pgfpathlineto{\pgfqpoint{3.655698in}{0.461837in}}%
\pgfpathlineto{\pgfqpoint{3.657968in}{0.527162in}}%
\pgfpathlineto{\pgfqpoint{3.660237in}{0.487971in}}%
\pgfpathlineto{\pgfqpoint{3.662507in}{0.514313in}}%
\pgfpathlineto{\pgfqpoint{3.664776in}{0.401880in}}%
\pgfpathlineto{\pgfqpoint{3.667046in}{0.462643in}}%
\pgfpathlineto{\pgfqpoint{3.669315in}{0.412944in}}%
\pgfpathlineto{\pgfqpoint{3.671585in}{0.440152in}}%
\pgfpathlineto{\pgfqpoint{3.673854in}{0.487395in}}%
\pgfpathlineto{\pgfqpoint{3.676124in}{0.496148in}}%
\pgfpathlineto{\pgfqpoint{3.678393in}{0.468806in}}%
\pgfpathlineto{\pgfqpoint{3.680663in}{0.455982in}}%
\pgfpathlineto{\pgfqpoint{3.682932in}{0.516660in}}%
\pgfpathlineto{\pgfqpoint{3.685202in}{0.527401in}}%
\pgfpathlineto{\pgfqpoint{3.689741in}{0.390428in}}%
\pgfpathlineto{\pgfqpoint{3.692010in}{0.514378in}}%
\pgfpathlineto{\pgfqpoint{3.694280in}{0.456445in}}%
\pgfpathlineto{\pgfqpoint{3.696549in}{0.504165in}}%
\pgfpathlineto{\pgfqpoint{3.698818in}{0.450576in}}%
\pgfpathlineto{\pgfqpoint{3.701088in}{0.475493in}}%
\pgfpathlineto{\pgfqpoint{3.703357in}{0.441350in}}%
\pgfpathlineto{\pgfqpoint{3.705627in}{0.440268in}}%
\pgfpathlineto{\pgfqpoint{3.707896in}{0.495289in}}%
\pgfpathlineto{\pgfqpoint{3.710166in}{0.512525in}}%
\pgfpathlineto{\pgfqpoint{3.712435in}{0.470177in}}%
\pgfpathlineto{\pgfqpoint{3.714705in}{0.503783in}}%
\pgfpathlineto{\pgfqpoint{3.716974in}{0.461566in}}%
\pgfpathlineto{\pgfqpoint{3.719244in}{0.506124in}}%
\pgfpathlineto{\pgfqpoint{3.721513in}{0.448212in}}%
\pgfpathlineto{\pgfqpoint{3.723783in}{0.479035in}}%
\pgfpathlineto{\pgfqpoint{3.728322in}{0.410125in}}%
\pgfpathlineto{\pgfqpoint{3.730591in}{0.485990in}}%
\pgfpathlineto{\pgfqpoint{3.732861in}{0.512093in}}%
\pgfpathlineto{\pgfqpoint{3.735130in}{0.485491in}}%
\pgfpathlineto{\pgfqpoint{3.737400in}{0.492194in}}%
\pgfpathlineto{\pgfqpoint{3.739669in}{0.455642in}}%
\pgfpathlineto{\pgfqpoint{3.741939in}{0.503623in}}%
\pgfpathlineto{\pgfqpoint{3.744208in}{0.453731in}}%
\pgfpathlineto{\pgfqpoint{3.746478in}{0.482911in}}%
\pgfpathlineto{\pgfqpoint{3.748747in}{0.455165in}}%
\pgfpathlineto{\pgfqpoint{3.751017in}{0.386105in}}%
\pgfpathlineto{\pgfqpoint{3.753286in}{0.483885in}}%
\pgfpathlineto{\pgfqpoint{3.755556in}{0.511761in}}%
\pgfpathlineto{\pgfqpoint{3.757825in}{0.493717in}}%
\pgfpathlineto{\pgfqpoint{3.760095in}{0.482474in}}%
\pgfpathlineto{\pgfqpoint{3.762364in}{0.447254in}}%
\pgfpathlineto{\pgfqpoint{3.764633in}{0.501023in}}%
\pgfpathlineto{\pgfqpoint{3.766903in}{0.456109in}}%
\pgfpathlineto{\pgfqpoint{3.769172in}{0.486129in}}%
\pgfpathlineto{\pgfqpoint{3.771442in}{0.461149in}}%
\pgfpathlineto{\pgfqpoint{3.773526in}{0.315625in}}%
\pgfpathmoveto{\pgfqpoint{3.773879in}{0.315625in}}%
\pgfpathlineto{\pgfqpoint{3.775981in}{0.478241in}}%
\pgfpathlineto{\pgfqpoint{3.778250in}{0.510947in}}%
\pgfpathlineto{\pgfqpoint{3.780520in}{0.500836in}}%
\pgfpathlineto{\pgfqpoint{3.785059in}{0.433401in}}%
\pgfpathlineto{\pgfqpoint{3.787328in}{0.497558in}}%
\pgfpathlineto{\pgfqpoint{3.789598in}{0.457648in}}%
\pgfpathlineto{\pgfqpoint{3.791867in}{0.488871in}}%
\pgfpathlineto{\pgfqpoint{3.794137in}{0.466446in}}%
\pgfpathlineto{\pgfqpoint{3.796406in}{0.398575in}}%
\pgfpathlineto{\pgfqpoint{3.798676in}{0.471238in}}%
\pgfpathlineto{\pgfqpoint{3.800945in}{0.509787in}}%
\pgfpathlineto{\pgfqpoint{3.803215in}{0.506573in}}%
\pgfpathlineto{\pgfqpoint{3.805484in}{0.449631in}}%
\pgfpathlineto{\pgfqpoint{3.807754in}{0.407791in}}%
\pgfpathlineto{\pgfqpoint{3.810023in}{0.493358in}}%
\pgfpathlineto{\pgfqpoint{3.812293in}{0.457937in}}%
\pgfpathlineto{\pgfqpoint{3.814562in}{0.490996in}}%
\pgfpathlineto{\pgfqpoint{3.816832in}{0.470949in}}%
\pgfpathlineto{\pgfqpoint{3.819101in}{0.427432in}}%
\pgfpathlineto{\pgfqpoint{3.821371in}{0.461429in}}%
\pgfpathlineto{\pgfqpoint{3.823640in}{0.508143in}}%
\pgfpathlineto{\pgfqpoint{3.825910in}{0.511257in}}%
\pgfpathlineto{\pgfqpoint{3.828179in}{0.449884in}}%
\pgfpathlineto{\pgfqpoint{3.832718in}{0.533513in}}%
\pgfpathlineto{\pgfqpoint{3.834987in}{0.485065in}}%
\pgfpathlineto{\pgfqpoint{3.837257in}{0.507108in}}%
\pgfpathlineto{\pgfqpoint{3.839526in}{0.395035in}}%
\pgfpathlineto{\pgfqpoint{3.841796in}{0.439565in}}%
\pgfpathlineto{\pgfqpoint{3.844065in}{0.447937in}}%
\pgfpathlineto{\pgfqpoint{3.846335in}{0.454094in}}%
\pgfpathlineto{\pgfqpoint{3.848604in}{0.490435in}}%
\pgfpathlineto{\pgfqpoint{3.850874in}{0.492646in}}%
\pgfpathlineto{\pgfqpoint{3.853143in}{0.436439in}}%
\pgfpathlineto{\pgfqpoint{3.855413in}{0.484562in}}%
\pgfpathlineto{\pgfqpoint{3.857682in}{0.520118in}}%
\pgfpathlineto{\pgfqpoint{3.859952in}{0.519548in}}%
\pgfpathlineto{\pgfqpoint{3.862221in}{0.400210in}}%
\pgfpathlineto{\pgfqpoint{3.866760in}{0.523523in}}%
\pgfpathlineto{\pgfqpoint{3.869030in}{0.459761in}}%
\pgfpathlineto{\pgfqpoint{3.871299in}{0.503893in}}%
\pgfpathlineto{\pgfqpoint{3.873569in}{0.439627in}}%
\pgfpathlineto{\pgfqpoint{3.875838in}{0.464090in}}%
\pgfpathlineto{\pgfqpoint{3.878108in}{0.401202in}}%
\pgfpathlineto{\pgfqpoint{3.880377in}{0.459534in}}%
\pgfpathlineto{\pgfqpoint{3.882647in}{0.498779in}}%
\pgfpathlineto{\pgfqpoint{3.884916in}{0.509217in}}%
\pgfpathlineto{\pgfqpoint{3.887186in}{0.410252in}}%
\pgfpathlineto{\pgfqpoint{3.889455in}{0.515166in}}%
\pgfpathlineto{\pgfqpoint{3.891725in}{0.468538in}}%
\pgfpathlineto{\pgfqpoint{3.893994in}{0.506920in}}%
\pgfpathlineto{\pgfqpoint{3.896264in}{0.437979in}}%
\pgfpathlineto{\pgfqpoint{3.898533in}{0.469509in}}%
\pgfpathlineto{\pgfqpoint{3.900802in}{0.420002in}}%
\pgfpathlineto{\pgfqpoint{3.903072in}{0.444112in}}%
\pgfpathlineto{\pgfqpoint{3.905341in}{0.493230in}}%
\pgfpathlineto{\pgfqpoint{3.907611in}{0.510242in}}%
\pgfpathlineto{\pgfqpoint{3.909880in}{0.456069in}}%
\pgfpathlineto{\pgfqpoint{3.912150in}{0.506397in}}%
\pgfpathlineto{\pgfqpoint{3.914419in}{0.466948in}}%
\pgfpathlineto{\pgfqpoint{3.916689in}{0.506763in}}%
\pgfpathlineto{\pgfqpoint{3.918958in}{0.447406in}}%
\pgfpathlineto{\pgfqpoint{3.921228in}{0.474618in}}%
\pgfpathlineto{\pgfqpoint{3.923497in}{0.433141in}}%
\pgfpathlineto{\pgfqpoint{3.925767in}{0.437784in}}%
\pgfpathlineto{\pgfqpoint{3.928036in}{0.492953in}}%
\pgfpathlineto{\pgfqpoint{3.930306in}{0.510564in}}%
\pgfpathlineto{\pgfqpoint{3.932575in}{0.471539in}}%
\pgfpathlineto{\pgfqpoint{3.934845in}{0.499678in}}%
\pgfpathlineto{\pgfqpoint{3.937114in}{0.463590in}}%
\pgfpathlineto{\pgfqpoint{3.939384in}{0.506283in}}%
\pgfpathlineto{\pgfqpoint{3.941653in}{0.452400in}}%
\pgfpathlineto{\pgfqpoint{3.943923in}{0.479202in}}%
\pgfpathlineto{\pgfqpoint{3.946192in}{0.443603in}}%
\pgfpathlineto{\pgfqpoint{3.948462in}{0.423775in}}%
\pgfpathlineto{\pgfqpoint{3.950731in}{0.490081in}}%
\pgfpathlineto{\pgfqpoint{3.953001in}{0.510637in}}%
\pgfpathlineto{\pgfqpoint{3.955270in}{0.483749in}}%
\pgfpathlineto{\pgfqpoint{3.957540in}{0.490670in}}%
\pgfpathlineto{\pgfqpoint{3.959809in}{0.457947in}}%
\pgfpathlineto{\pgfqpoint{3.962079in}{0.504510in}}%
\pgfpathlineto{\pgfqpoint{3.964348in}{0.456350in}}%
\pgfpathlineto{\pgfqpoint{3.966617in}{0.483310in}}%
\pgfpathlineto{\pgfqpoint{3.968887in}{0.452262in}}%
\pgfpathlineto{\pgfqpoint{3.971156in}{0.400474in}}%
\pgfpathlineto{\pgfqpoint{3.973426in}{0.486422in}}%
\pgfpathlineto{\pgfqpoint{3.975695in}{0.510401in}}%
\pgfpathlineto{\pgfqpoint{3.980234in}{0.479655in}}%
\pgfpathlineto{\pgfqpoint{3.982504in}{0.449234in}}%
\pgfpathlineto{\pgfqpoint{3.984773in}{0.501682in}}%
\pgfpathlineto{\pgfqpoint{3.987043in}{0.458920in}}%
\pgfpathlineto{\pgfqpoint{3.989312in}{0.486857in}}%
\pgfpathlineto{\pgfqpoint{3.991582in}{0.459452in}}%
\pgfpathlineto{\pgfqpoint{3.993851in}{0.323994in}}%
\pgfpathlineto{\pgfqpoint{3.996121in}{0.481269in}}%
\pgfpathlineto{\pgfqpoint{3.998390in}{0.509776in}}%
\pgfpathlineto{\pgfqpoint{4.000660in}{0.500248in}}%
\pgfpathlineto{\pgfqpoint{4.005199in}{0.435315in}}%
\pgfpathlineto{\pgfqpoint{4.007468in}{0.498143in}}%
\pgfpathlineto{\pgfqpoint{4.009738in}{0.460305in}}%
\pgfpathlineto{\pgfqpoint{4.012007in}{0.489805in}}%
\pgfpathlineto{\pgfqpoint{4.014277in}{0.465471in}}%
\pgfpathlineto{\pgfqpoint{4.016546in}{0.388694in}}%
\pgfpathlineto{\pgfqpoint{4.018816in}{0.474340in}}%
\pgfpathlineto{\pgfqpoint{4.021085in}{0.508709in}}%
\pgfpathlineto{\pgfqpoint{4.023355in}{0.506154in}}%
\pgfpathlineto{\pgfqpoint{4.025624in}{0.405343in}}%
\pgfpathlineto{\pgfqpoint{4.027894in}{0.496857in}}%
\pgfpathlineto{\pgfqpoint{4.030163in}{0.534001in}}%
\pgfpathlineto{\pgfqpoint{4.032432in}{0.482265in}}%
\pgfpathlineto{\pgfqpoint{4.034702in}{0.501869in}}%
\pgfpathlineto{\pgfqpoint{4.036971in}{0.426275in}}%
\pgfpathlineto{\pgfqpoint{4.039241in}{0.418556in}}%
\pgfpathlineto{\pgfqpoint{4.041510in}{0.460291in}}%
\pgfpathlineto{\pgfqpoint{4.043780in}{0.459994in}}%
\pgfpathlineto{\pgfqpoint{4.046049in}{0.491213in}}%
\pgfpathlineto{\pgfqpoint{4.048319in}{0.489325in}}%
\pgfpathlineto{\pgfqpoint{4.050588in}{0.396377in}}%
\pgfpathlineto{\pgfqpoint{4.052858in}{0.493483in}}%
\pgfpathlineto{\pgfqpoint{4.055127in}{0.519976in}}%
\pgfpathlineto{\pgfqpoint{4.057397in}{0.513862in}}%
\pgfpathlineto{\pgfqpoint{4.059666in}{0.446892in}}%
\pgfpathlineto{\pgfqpoint{4.061936in}{0.478288in}}%
\pgfpathlineto{\pgfqpoint{4.064205in}{0.526138in}}%
\pgfpathlineto{\pgfqpoint{4.066475in}{0.460622in}}%
\pgfpathlineto{\pgfqpoint{4.068744in}{0.502494in}}%
\pgfpathlineto{\pgfqpoint{4.071014in}{0.431750in}}%
\pgfpathlineto{\pgfqpoint{4.073283in}{0.455987in}}%
\pgfpathlineto{\pgfqpoint{4.074603in}{0.315625in}}%
\pgfpathmoveto{\pgfqpoint{4.076462in}{0.315625in}}%
\pgfpathlineto{\pgfqpoint{4.077822in}{0.466823in}}%
\pgfpathlineto{\pgfqpoint{4.080092in}{0.499547in}}%
\pgfpathlineto{\pgfqpoint{4.082361in}{0.505762in}}%
\pgfpathlineto{\pgfqpoint{4.084631in}{0.380154in}}%
\pgfpathlineto{\pgfqpoint{4.086900in}{0.518737in}}%
\pgfpathlineto{\pgfqpoint{4.089170in}{0.470568in}}%
\pgfpathlineto{\pgfqpoint{4.091439in}{0.505972in}}%
\pgfpathlineto{\pgfqpoint{4.093709in}{0.430799in}}%
\pgfpathlineto{\pgfqpoint{4.095978in}{0.463057in}}%
\pgfpathlineto{\pgfqpoint{4.098247in}{0.383395in}}%
\pgfpathlineto{\pgfqpoint{4.100517in}{0.454616in}}%
\pgfpathlineto{\pgfqpoint{4.102786in}{0.495611in}}%
\pgfpathlineto{\pgfqpoint{4.105056in}{0.507834in}}%
\pgfpathlineto{\pgfqpoint{4.107325in}{0.427699in}}%
\pgfpathlineto{\pgfqpoint{4.109595in}{0.510919in}}%
\pgfpathlineto{\pgfqpoint{4.111864in}{0.470428in}}%
\pgfpathlineto{\pgfqpoint{4.114134in}{0.506760in}}%
\pgfpathlineto{\pgfqpoint{4.116403in}{0.443136in}}%
\pgfpathlineto{\pgfqpoint{4.118673in}{0.469154in}}%
\pgfpathlineto{\pgfqpoint{4.120942in}{0.410341in}}%
\pgfpathlineto{\pgfqpoint{4.125481in}{0.496090in}}%
\pgfpathlineto{\pgfqpoint{4.127751in}{0.508586in}}%
\pgfpathlineto{\pgfqpoint{4.130020in}{0.454206in}}%
\pgfpathlineto{\pgfqpoint{4.132290in}{0.505052in}}%
\pgfpathlineto{\pgfqpoint{4.134559in}{0.468468in}}%
\pgfpathlineto{\pgfqpoint{4.136829in}{0.507217in}}%
\pgfpathlineto{\pgfqpoint{4.139098in}{0.449788in}}%
\pgfpathlineto{\pgfqpoint{4.141368in}{0.474702in}}%
\pgfpathlineto{\pgfqpoint{4.143637in}{0.427975in}}%
\pgfpathlineto{\pgfqpoint{4.145907in}{0.441335in}}%
\pgfpathlineto{\pgfqpoint{4.148176in}{0.494303in}}%
\pgfpathlineto{\pgfqpoint{4.150446in}{0.509187in}}%
\pgfpathlineto{\pgfqpoint{4.152715in}{0.471816in}}%
\pgfpathlineto{\pgfqpoint{4.154985in}{0.497249in}}%
\pgfpathlineto{\pgfqpoint{4.157254in}{0.464698in}}%
\pgfpathlineto{\pgfqpoint{4.159524in}{0.506597in}}%
\pgfpathlineto{\pgfqpoint{4.161793in}{0.455073in}}%
\pgfpathlineto{\pgfqpoint{4.164062in}{0.479669in}}%
\pgfpathlineto{\pgfqpoint{4.166332in}{0.440802in}}%
\pgfpathlineto{\pgfqpoint{4.168601in}{0.428487in}}%
\pgfpathlineto{\pgfqpoint{4.170871in}{0.491751in}}%
\pgfpathlineto{\pgfqpoint{4.173140in}{0.509447in}}%
\pgfpathlineto{\pgfqpoint{4.175410in}{0.483957in}}%
\pgfpathlineto{\pgfqpoint{4.177679in}{0.487917in}}%
\pgfpathlineto{\pgfqpoint{4.179949in}{0.458694in}}%
\pgfpathlineto{\pgfqpoint{4.182218in}{0.504798in}}%
\pgfpathlineto{\pgfqpoint{4.184488in}{0.458763in}}%
\pgfpathlineto{\pgfqpoint{4.186757in}{0.484018in}}%
\pgfpathlineto{\pgfqpoint{4.189027in}{0.450759in}}%
\pgfpathlineto{\pgfqpoint{4.191296in}{0.406249in}}%
\pgfpathlineto{\pgfqpoint{4.193566in}{0.487936in}}%
\pgfpathlineto{\pgfqpoint{4.195835in}{0.509326in}}%
\pgfpathlineto{\pgfqpoint{4.198105in}{0.493225in}}%
\pgfpathlineto{\pgfqpoint{4.200374in}{0.450387in}}%
\pgfpathlineto{\pgfqpoint{4.202644in}{0.508360in}}%
\pgfpathlineto{\pgfqpoint{4.204913in}{0.534105in}}%
\pgfpathlineto{\pgfqpoint{4.207183in}{0.469431in}}%
\pgfpathlineto{\pgfqpoint{4.209452in}{0.487490in}}%
\pgfpathlineto{\pgfqpoint{4.211722in}{0.461859in}}%
\pgfpathlineto{\pgfqpoint{4.213908in}{0.315625in}}%
\pgfpathmoveto{\pgfqpoint{4.214068in}{0.315625in}}%
\pgfpathlineto{\pgfqpoint{4.216261in}{0.474830in}}%
\pgfpathlineto{\pgfqpoint{4.218530in}{0.465018in}}%
\pgfpathlineto{\pgfqpoint{4.220800in}{0.488615in}}%
\pgfpathlineto{\pgfqpoint{4.223069in}{0.481495in}}%
\pgfpathlineto{\pgfqpoint{4.225339in}{0.422455in}}%
\pgfpathlineto{\pgfqpoint{4.227608in}{0.504880in}}%
\pgfpathlineto{\pgfqpoint{4.229878in}{0.520059in}}%
\pgfpathlineto{\pgfqpoint{4.232147in}{0.498625in}}%
\pgfpathlineto{\pgfqpoint{4.234416in}{0.488403in}}%
\pgfpathlineto{\pgfqpoint{4.236686in}{0.494972in}}%
\pgfpathlineto{\pgfqpoint{4.238955in}{0.530559in}}%
\pgfpathlineto{\pgfqpoint{4.241225in}{0.453009in}}%
\pgfpathlineto{\pgfqpoint{4.243494in}{0.496783in}}%
\pgfpathlineto{\pgfqpoint{4.245764in}{0.392336in}}%
\pgfpathlineto{\pgfqpoint{4.248033in}{0.433532in}}%
\pgfpathlineto{\pgfqpoint{4.250303in}{0.424255in}}%
\pgfpathlineto{\pgfqpoint{4.252572in}{0.474924in}}%
\pgfpathlineto{\pgfqpoint{4.254842in}{0.497635in}}%
\pgfpathlineto{\pgfqpoint{4.257111in}{0.498392in}}%
\pgfpathlineto{\pgfqpoint{4.259381in}{0.462340in}}%
\pgfpathlineto{\pgfqpoint{4.261650in}{0.525870in}}%
\pgfpathlineto{\pgfqpoint{4.263920in}{0.468779in}}%
\pgfpathlineto{\pgfqpoint{4.266189in}{0.501155in}}%
\pgfpathlineto{\pgfqpoint{4.268459in}{0.392366in}}%
\pgfpathlineto{\pgfqpoint{4.270728in}{0.445150in}}%
\pgfpathlineto{\pgfqpoint{4.272998in}{0.400815in}}%
\pgfpathlineto{\pgfqpoint{4.275267in}{0.466320in}}%
\pgfpathlineto{\pgfqpoint{4.277537in}{0.496022in}}%
\pgfpathlineto{\pgfqpoint{4.279806in}{0.502573in}}%
\pgfpathlineto{\pgfqpoint{4.282076in}{0.425152in}}%
\pgfpathlineto{\pgfqpoint{4.284345in}{0.519814in}}%
\pgfpathlineto{\pgfqpoint{4.286615in}{0.472092in}}%
\pgfpathlineto{\pgfqpoint{4.288884in}{0.503662in}}%
\pgfpathlineto{\pgfqpoint{4.291154in}{0.421993in}}%
\pgfpathlineto{\pgfqpoint{4.293423in}{0.453554in}}%
\pgfpathlineto{\pgfqpoint{4.295693in}{0.368864in}}%
\pgfpathlineto{\pgfqpoint{4.297962in}{0.464748in}}%
\pgfpathlineto{\pgfqpoint{4.300231in}{0.497752in}}%
\pgfpathlineto{\pgfqpoint{4.302501in}{0.504045in}}%
\pgfpathlineto{\pgfqpoint{4.304770in}{0.332497in}}%
\pgfpathlineto{\pgfqpoint{4.307040in}{0.515556in}}%
\pgfpathlineto{\pgfqpoint{4.309309in}{0.472829in}}%
\pgfpathlineto{\pgfqpoint{4.311579in}{0.505768in}}%
\pgfpathlineto{\pgfqpoint{4.313848in}{0.435323in}}%
\pgfpathlineto{\pgfqpoint{4.316118in}{0.461371in}}%
\pgfpathlineto{\pgfqpoint{4.318387in}{0.328470in}}%
\pgfpathlineto{\pgfqpoint{4.320657in}{0.460073in}}%
\pgfpathlineto{\pgfqpoint{4.322926in}{0.497898in}}%
\pgfpathlineto{\pgfqpoint{4.325196in}{0.505637in}}%
\pgfpathlineto{\pgfqpoint{4.327465in}{0.421577in}}%
\pgfpathlineto{\pgfqpoint{4.329735in}{0.509865in}}%
\pgfpathlineto{\pgfqpoint{4.332004in}{0.472197in}}%
\pgfpathlineto{\pgfqpoint{4.334274in}{0.507068in}}%
\pgfpathlineto{\pgfqpoint{4.336543in}{0.457324in}}%
\pgfpathlineto{\pgfqpoint{4.336543in}{0.457324in}}%
\pgfusepath{stroke}%
\end{pgfscope}%
\begin{pgfscope}%
\pgfpathrectangle{\pgfqpoint{2.520957in}{0.320625in}}{\pgfqpoint{1.815586in}{1.164375in}}%
\pgfusepath{clip}%
\pgfsetbuttcap%
\pgfsetroundjoin%
\pgfsetlinewidth{0.803000pt}%
\definecolor{currentstroke}{rgb}{0.000000,0.000000,0.000000}%
\pgfsetstrokecolor{currentstroke}%
\pgfsetstrokeopacity{0.700000}%
\pgfsetdash{{2.960000pt}{1.280000pt}}{0.000000pt}%
\pgfpathmoveto{\pgfqpoint{2.520957in}{1.451682in}}%
\pgfpathlineto{\pgfqpoint{3.132787in}{0.315625in}}%
\pgfpathlineto{\pgfqpoint{3.132787in}{0.315625in}}%
\pgfusepath{stroke}%
\end{pgfscope}%
\begin{pgfscope}%
\pgfsetrectcap%
\pgfsetmiterjoin%
\pgfsetlinewidth{0.803000pt}%
\definecolor{currentstroke}{rgb}{0.000000,0.000000,0.000000}%
\pgfsetstrokecolor{currentstroke}%
\pgfsetdash{}{0pt}%
\pgfpathmoveto{\pgfqpoint{2.520957in}{0.320625in}}%
\pgfpathlineto{\pgfqpoint{2.520957in}{1.485000in}}%
\pgfusepath{stroke}%
\end{pgfscope}%
\begin{pgfscope}%
\pgfsetrectcap%
\pgfsetmiterjoin%
\pgfsetlinewidth{0.803000pt}%
\definecolor{currentstroke}{rgb}{0.000000,0.000000,0.000000}%
\pgfsetstrokecolor{currentstroke}%
\pgfsetdash{}{0pt}%
\pgfpathmoveto{\pgfqpoint{4.336543in}{0.320625in}}%
\pgfpathlineto{\pgfqpoint{4.336543in}{1.485000in}}%
\pgfusepath{stroke}%
\end{pgfscope}%
\begin{pgfscope}%
\pgfsetrectcap%
\pgfsetmiterjoin%
\pgfsetlinewidth{0.803000pt}%
\definecolor{currentstroke}{rgb}{0.000000,0.000000,0.000000}%
\pgfsetstrokecolor{currentstroke}%
\pgfsetdash{}{0pt}%
\pgfpathmoveto{\pgfqpoint{2.520957in}{0.320625in}}%
\pgfpathlineto{\pgfqpoint{4.336543in}{0.320625in}}%
\pgfusepath{stroke}%
\end{pgfscope}%
\begin{pgfscope}%
\pgfsetrectcap%
\pgfsetmiterjoin%
\pgfsetlinewidth{0.803000pt}%
\definecolor{currentstroke}{rgb}{0.000000,0.000000,0.000000}%
\pgfsetstrokecolor{currentstroke}%
\pgfsetdash{}{0pt}%
\pgfpathmoveto{\pgfqpoint{2.520957in}{1.485000in}}%
\pgfpathlineto{\pgfqpoint{4.336543in}{1.485000in}}%
\pgfusepath{stroke}%
\end{pgfscope}%
\begin{pgfscope}%
\definecolor{textcolor}{rgb}{0.000000,0.000000,0.000000}%
\pgfsetstrokecolor{textcolor}%
\pgfsetfillcolor{textcolor}%
\pgftext[x=3.428750in,y=1.568333in,,base]{\color{textcolor}{\rmfamily\fontsize{8.000000}{9.600000}\selectfont\catcode`\^=\active\def^{\ifmmode\sp\else\^{}\fi}\catcode`\%=\active\def%{\%}(b) Covariance conditioning}}%
\end{pgfscope}%
\begin{pgfscope}%
\pgfsetbuttcap%
\pgfsetmiterjoin%
\definecolor{currentfill}{rgb}{1.000000,1.000000,1.000000}%
\pgfsetfillcolor{currentfill}%
\pgfsetfillopacity{0.800000}%
\pgfsetlinewidth{1.003750pt}%
\definecolor{currentstroke}{rgb}{0.800000,0.800000,0.800000}%
\pgfsetstrokecolor{currentstroke}%
\pgfsetstrokeopacity{0.800000}%
\pgfsetdash{}{0pt}%
\pgfpathmoveto{\pgfqpoint{3.841321in}{1.203889in}}%
\pgfpathlineto{\pgfqpoint{4.291821in}{1.203889in}}%
\pgfpathquadraticcurveto{\pgfqpoint{4.304599in}{1.203889in}}{\pgfqpoint{4.304599in}{1.216667in}}%
\pgfpathlineto{\pgfqpoint{4.304599in}{1.440278in}}%
\pgfpathquadraticcurveto{\pgfqpoint{4.304599in}{1.453056in}}{\pgfqpoint{4.291821in}{1.453056in}}%
\pgfpathlineto{\pgfqpoint{3.841321in}{1.453056in}}%
\pgfpathquadraticcurveto{\pgfqpoint{3.828544in}{1.453056in}}{\pgfqpoint{3.828544in}{1.440278in}}%
\pgfpathlineto{\pgfqpoint{3.828544in}{1.216667in}}%
\pgfpathquadraticcurveto{\pgfqpoint{3.828544in}{1.203889in}}{\pgfqpoint{3.841321in}{1.203889in}}%
\pgfpathlineto{\pgfqpoint{3.841321in}{1.203889in}}%
\pgfpathclose%
\pgfusepath{stroke,fill}%
\end{pgfscope}%
\begin{pgfscope}%
\pgfsetrectcap%
\pgfsetroundjoin%
\pgfsetlinewidth{1.154312pt}%
\definecolor{currentstroke}{rgb}{0.146785,0.068738,0.334011}%
\pgfsetstrokecolor{currentstroke}%
\pgfsetdash{}{0pt}%
\pgfpathmoveto{\pgfqpoint{3.844516in}{1.410904in}}%
\pgfpathlineto{\pgfqpoint{3.889238in}{1.410904in}}%
\pgfpathlineto{\pgfqpoint{3.933960in}{1.410904in}}%
\pgfusepath{stroke}%
\end{pgfscope}%
\begin{pgfscope}%
\definecolor{textcolor}{rgb}{0.000000,0.000000,0.000000}%
\pgfsetstrokecolor{textcolor}%
\pgfsetfillcolor{textcolor}%
\pgftext[x=3.985071in,y=1.388543in,left,base]{\color{textcolor}{\rmfamily\fontsize{4.600000}{5.520000}\selectfont\catcode`\^=\active\def^{\ifmmode\sp\else\^{}\fi}\catcode`\%=\active\def%{\%}$r=0$}}%
\end{pgfscope}%
\begin{pgfscope}%
\pgfsetrectcap%
\pgfsetroundjoin%
\pgfsetlinewidth{1.154312pt}%
\definecolor{currentstroke}{rgb}{0.716387,0.214982,0.475290}%
\pgfsetstrokecolor{currentstroke}%
\pgfsetdash{}{0pt}%
\pgfpathmoveto{\pgfqpoint{3.844516in}{1.334237in}}%
\pgfpathlineto{\pgfqpoint{3.889238in}{1.334237in}}%
\pgfpathlineto{\pgfqpoint{3.933960in}{1.334237in}}%
\pgfusepath{stroke}%
\end{pgfscope}%
\begin{pgfscope}%
\definecolor{textcolor}{rgb}{0.000000,0.000000,0.000000}%
\pgfsetstrokecolor{textcolor}%
\pgfsetfillcolor{textcolor}%
\pgftext[x=3.985071in,y=1.311876in,left,base]{\color{textcolor}{\rmfamily\fontsize{4.600000}{5.520000}\selectfont\catcode`\^=\active\def^{\ifmmode\sp\else\^{}\fi}\catcode`\%=\active\def%{\%}$r=0.5$}}%
\end{pgfscope}%
\begin{pgfscope}%
\pgfsetrectcap%
\pgfsetroundjoin%
\pgfsetlinewidth{1.154312pt}%
\definecolor{currentstroke}{rgb}{0.997351,0.719089,0.493755}%
\pgfsetstrokecolor{currentstroke}%
\pgfsetdash{}{0pt}%
\pgfpathmoveto{\pgfqpoint{3.844516in}{1.257571in}}%
\pgfpathlineto{\pgfqpoint{3.889238in}{1.257571in}}%
\pgfpathlineto{\pgfqpoint{3.933960in}{1.257571in}}%
\pgfusepath{stroke}%
\end{pgfscope}%
\begin{pgfscope}%
\definecolor{textcolor}{rgb}{0.000000,0.000000,0.000000}%
\pgfsetstrokecolor{textcolor}%
\pgfsetfillcolor{textcolor}%
\pgftext[x=3.985071in,y=1.235209in,left,base]{\color{textcolor}{\rmfamily\fontsize{4.600000}{5.520000}\selectfont\catcode`\^=\active\def^{\ifmmode\sp\else\^{}\fi}\catcode`\%=\active\def%{\%}$r=1$}}%
\end{pgfscope}%
\begin{pgfscope}%
\pgfsetbuttcap%
\pgfsetmiterjoin%
\definecolor{currentfill}{rgb}{1.000000,1.000000,1.000000}%
\pgfsetfillcolor{currentfill}%
\pgfsetlinewidth{0.000000pt}%
\definecolor{currentstroke}{rgb}{0.000000,0.000000,0.000000}%
\pgfsetstrokecolor{currentstroke}%
\pgfsetstrokeopacity{0.000000}%
\pgfsetdash{}{0pt}%
\pgfpathmoveto{\pgfqpoint{4.554414in}{0.320625in}}%
\pgfpathlineto{\pgfqpoint{6.370000in}{0.320625in}}%
\pgfpathlineto{\pgfqpoint{6.370000in}{1.485000in}}%
\pgfpathlineto{\pgfqpoint{4.554414in}{1.485000in}}%
\pgfpathlineto{\pgfqpoint{4.554414in}{0.320625in}}%
\pgfpathclose%
\pgfusepath{fill}%
\end{pgfscope}%
\begin{pgfscope}%
\pgfpathrectangle{\pgfqpoint{4.554414in}{0.320625in}}{\pgfqpoint{1.815586in}{1.164375in}}%
\pgfusepath{clip}%
\pgfsetrectcap%
\pgfsetroundjoin%
\pgfsetlinewidth{0.803000pt}%
\definecolor{currentstroke}{rgb}{0.909804,0.909804,0.909804}%
\pgfsetstrokecolor{currentstroke}%
\pgfsetdash{}{0pt}%
\pgfpathmoveto{\pgfqpoint{4.554414in}{0.320625in}}%
\pgfpathlineto{\pgfqpoint{4.554414in}{1.485000in}}%
\pgfusepath{stroke}%
\end{pgfscope}%
\begin{pgfscope}%
\pgfsetbuttcap%
\pgfsetroundjoin%
\definecolor{currentfill}{rgb}{0.000000,0.000000,0.000000}%
\pgfsetfillcolor{currentfill}%
\pgfsetlinewidth{0.803000pt}%
\definecolor{currentstroke}{rgb}{0.000000,0.000000,0.000000}%
\pgfsetstrokecolor{currentstroke}%
\pgfsetdash{}{0pt}%
\pgfsys@defobject{currentmarker}{\pgfqpoint{0.000000in}{-0.048611in}}{\pgfqpoint{0.000000in}{0.000000in}}{%
\pgfpathmoveto{\pgfqpoint{0.000000in}{0.000000in}}%
\pgfpathlineto{\pgfqpoint{0.000000in}{-0.048611in}}%
\pgfusepath{stroke,fill}%
}%
\begin{pgfscope}%
\pgfsys@transformshift{4.554414in}{0.320625in}%
\pgfsys@useobject{currentmarker}{}%
\end{pgfscope}%
\end{pgfscope}%
\begin{pgfscope}%
\definecolor{textcolor}{rgb}{0.000000,0.000000,0.000000}%
\pgfsetstrokecolor{textcolor}%
\pgfsetfillcolor{textcolor}%
\pgftext[x=4.554414in,y=0.223403in,,top]{\color{textcolor}{\rmfamily\fontsize{6.000000}{7.200000}\selectfont\catcode`\^=\active\def^{\ifmmode\sp\else\^{}\fi}\catcode`\%=\active\def%{\%}0}}%
\end{pgfscope}%
\begin{pgfscope}%
\pgfpathrectangle{\pgfqpoint{4.554414in}{0.320625in}}{\pgfqpoint{1.815586in}{1.164375in}}%
\pgfusepath{clip}%
\pgfsetrectcap%
\pgfsetroundjoin%
\pgfsetlinewidth{0.803000pt}%
\definecolor{currentstroke}{rgb}{0.909804,0.909804,0.909804}%
\pgfsetstrokecolor{currentstroke}%
\pgfsetdash{}{0pt}%
\pgfpathmoveto{\pgfqpoint{5.008310in}{0.320625in}}%
\pgfpathlineto{\pgfqpoint{5.008310in}{1.485000in}}%
\pgfusepath{stroke}%
\end{pgfscope}%
\begin{pgfscope}%
\pgfsetbuttcap%
\pgfsetroundjoin%
\definecolor{currentfill}{rgb}{0.000000,0.000000,0.000000}%
\pgfsetfillcolor{currentfill}%
\pgfsetlinewidth{0.803000pt}%
\definecolor{currentstroke}{rgb}{0.000000,0.000000,0.000000}%
\pgfsetstrokecolor{currentstroke}%
\pgfsetdash{}{0pt}%
\pgfsys@defobject{currentmarker}{\pgfqpoint{0.000000in}{-0.048611in}}{\pgfqpoint{0.000000in}{0.000000in}}{%
\pgfpathmoveto{\pgfqpoint{0.000000in}{0.000000in}}%
\pgfpathlineto{\pgfqpoint{0.000000in}{-0.048611in}}%
\pgfusepath{stroke,fill}%
}%
\begin{pgfscope}%
\pgfsys@transformshift{5.008310in}{0.320625in}%
\pgfsys@useobject{currentmarker}{}%
\end{pgfscope}%
\end{pgfscope}%
\begin{pgfscope}%
\definecolor{textcolor}{rgb}{0.000000,0.000000,0.000000}%
\pgfsetstrokecolor{textcolor}%
\pgfsetfillcolor{textcolor}%
\pgftext[x=5.008310in,y=0.223403in,,top]{\color{textcolor}{\rmfamily\fontsize{6.000000}{7.200000}\selectfont\catcode`\^=\active\def^{\ifmmode\sp\else\^{}\fi}\catcode`\%=\active\def%{\%}10}}%
\end{pgfscope}%
\begin{pgfscope}%
\pgfpathrectangle{\pgfqpoint{4.554414in}{0.320625in}}{\pgfqpoint{1.815586in}{1.164375in}}%
\pgfusepath{clip}%
\pgfsetrectcap%
\pgfsetroundjoin%
\pgfsetlinewidth{0.803000pt}%
\definecolor{currentstroke}{rgb}{0.909804,0.909804,0.909804}%
\pgfsetstrokecolor{currentstroke}%
\pgfsetdash{}{0pt}%
\pgfpathmoveto{\pgfqpoint{5.462207in}{0.320625in}}%
\pgfpathlineto{\pgfqpoint{5.462207in}{1.485000in}}%
\pgfusepath{stroke}%
\end{pgfscope}%
\begin{pgfscope}%
\pgfsetbuttcap%
\pgfsetroundjoin%
\definecolor{currentfill}{rgb}{0.000000,0.000000,0.000000}%
\pgfsetfillcolor{currentfill}%
\pgfsetlinewidth{0.803000pt}%
\definecolor{currentstroke}{rgb}{0.000000,0.000000,0.000000}%
\pgfsetstrokecolor{currentstroke}%
\pgfsetdash{}{0pt}%
\pgfsys@defobject{currentmarker}{\pgfqpoint{0.000000in}{-0.048611in}}{\pgfqpoint{0.000000in}{0.000000in}}{%
\pgfpathmoveto{\pgfqpoint{0.000000in}{0.000000in}}%
\pgfpathlineto{\pgfqpoint{0.000000in}{-0.048611in}}%
\pgfusepath{stroke,fill}%
}%
\begin{pgfscope}%
\pgfsys@transformshift{5.462207in}{0.320625in}%
\pgfsys@useobject{currentmarker}{}%
\end{pgfscope}%
\end{pgfscope}%
\begin{pgfscope}%
\definecolor{textcolor}{rgb}{0.000000,0.000000,0.000000}%
\pgfsetstrokecolor{textcolor}%
\pgfsetfillcolor{textcolor}%
\pgftext[x=5.462207in,y=0.223403in,,top]{\color{textcolor}{\rmfamily\fontsize{6.000000}{7.200000}\selectfont\catcode`\^=\active\def^{\ifmmode\sp\else\^{}\fi}\catcode`\%=\active\def%{\%}20}}%
\end{pgfscope}%
\begin{pgfscope}%
\pgfpathrectangle{\pgfqpoint{4.554414in}{0.320625in}}{\pgfqpoint{1.815586in}{1.164375in}}%
\pgfusepath{clip}%
\pgfsetrectcap%
\pgfsetroundjoin%
\pgfsetlinewidth{0.803000pt}%
\definecolor{currentstroke}{rgb}{0.909804,0.909804,0.909804}%
\pgfsetstrokecolor{currentstroke}%
\pgfsetdash{}{0pt}%
\pgfpathmoveto{\pgfqpoint{5.916103in}{0.320625in}}%
\pgfpathlineto{\pgfqpoint{5.916103in}{1.485000in}}%
\pgfusepath{stroke}%
\end{pgfscope}%
\begin{pgfscope}%
\pgfsetbuttcap%
\pgfsetroundjoin%
\definecolor{currentfill}{rgb}{0.000000,0.000000,0.000000}%
\pgfsetfillcolor{currentfill}%
\pgfsetlinewidth{0.803000pt}%
\definecolor{currentstroke}{rgb}{0.000000,0.000000,0.000000}%
\pgfsetstrokecolor{currentstroke}%
\pgfsetdash{}{0pt}%
\pgfsys@defobject{currentmarker}{\pgfqpoint{0.000000in}{-0.048611in}}{\pgfqpoint{0.000000in}{0.000000in}}{%
\pgfpathmoveto{\pgfqpoint{0.000000in}{0.000000in}}%
\pgfpathlineto{\pgfqpoint{0.000000in}{-0.048611in}}%
\pgfusepath{stroke,fill}%
}%
\begin{pgfscope}%
\pgfsys@transformshift{5.916103in}{0.320625in}%
\pgfsys@useobject{currentmarker}{}%
\end{pgfscope}%
\end{pgfscope}%
\begin{pgfscope}%
\definecolor{textcolor}{rgb}{0.000000,0.000000,0.000000}%
\pgfsetstrokecolor{textcolor}%
\pgfsetfillcolor{textcolor}%
\pgftext[x=5.916103in,y=0.223403in,,top]{\color{textcolor}{\rmfamily\fontsize{6.000000}{7.200000}\selectfont\catcode`\^=\active\def^{\ifmmode\sp\else\^{}\fi}\catcode`\%=\active\def%{\%}30}}%
\end{pgfscope}%
\begin{pgfscope}%
\pgfpathrectangle{\pgfqpoint{4.554414in}{0.320625in}}{\pgfqpoint{1.815586in}{1.164375in}}%
\pgfusepath{clip}%
\pgfsetrectcap%
\pgfsetroundjoin%
\pgfsetlinewidth{0.803000pt}%
\definecolor{currentstroke}{rgb}{0.909804,0.909804,0.909804}%
\pgfsetstrokecolor{currentstroke}%
\pgfsetdash{}{0pt}%
\pgfpathmoveto{\pgfqpoint{6.370000in}{0.320625in}}%
\pgfpathlineto{\pgfqpoint{6.370000in}{1.485000in}}%
\pgfusepath{stroke}%
\end{pgfscope}%
\begin{pgfscope}%
\pgfsetbuttcap%
\pgfsetroundjoin%
\definecolor{currentfill}{rgb}{0.000000,0.000000,0.000000}%
\pgfsetfillcolor{currentfill}%
\pgfsetlinewidth{0.803000pt}%
\definecolor{currentstroke}{rgb}{0.000000,0.000000,0.000000}%
\pgfsetstrokecolor{currentstroke}%
\pgfsetdash{}{0pt}%
\pgfsys@defobject{currentmarker}{\pgfqpoint{0.000000in}{-0.048611in}}{\pgfqpoint{0.000000in}{0.000000in}}{%
\pgfpathmoveto{\pgfqpoint{0.000000in}{0.000000in}}%
\pgfpathlineto{\pgfqpoint{0.000000in}{-0.048611in}}%
\pgfusepath{stroke,fill}%
}%
\begin{pgfscope}%
\pgfsys@transformshift{6.370000in}{0.320625in}%
\pgfsys@useobject{currentmarker}{}%
\end{pgfscope}%
\end{pgfscope}%
\begin{pgfscope}%
\definecolor{textcolor}{rgb}{0.000000,0.000000,0.000000}%
\pgfsetstrokecolor{textcolor}%
\pgfsetfillcolor{textcolor}%
\pgftext[x=6.370000in,y=0.223403in,,top]{\color{textcolor}{\rmfamily\fontsize{6.000000}{7.200000}\selectfont\catcode`\^=\active\def^{\ifmmode\sp\else\^{}\fi}\catcode`\%=\active\def%{\%}40}}%
\end{pgfscope}%
\begin{pgfscope}%
\definecolor{textcolor}{rgb}{0.000000,0.000000,0.000000}%
\pgfsetstrokecolor{textcolor}%
\pgfsetfillcolor{textcolor}%
\pgftext[x=5.462207in,y=0.084514in,,top]{\color{textcolor}{\rmfamily\fontsize{6.000000}{7.200000}\selectfont\catcode`\^=\active\def^{\ifmmode\sp\else\^{}\fi}\catcode`\%=\active\def%{\%}time}}%
\end{pgfscope}%
\begin{pgfscope}%
\pgfpathrectangle{\pgfqpoint{4.554414in}{0.320625in}}{\pgfqpoint{1.815586in}{1.164375in}}%
\pgfusepath{clip}%
\pgfsetrectcap%
\pgfsetroundjoin%
\pgfsetlinewidth{0.803000pt}%
\definecolor{currentstroke}{rgb}{0.909804,0.909804,0.909804}%
\pgfsetstrokecolor{currentstroke}%
\pgfsetdash{}{0pt}%
\pgfpathmoveto{\pgfqpoint{4.554414in}{0.466172in}}%
\pgfpathlineto{\pgfqpoint{6.370000in}{0.466172in}}%
\pgfusepath{stroke}%
\end{pgfscope}%
\begin{pgfscope}%
\pgfsetbuttcap%
\pgfsetroundjoin%
\definecolor{currentfill}{rgb}{0.000000,0.000000,0.000000}%
\pgfsetfillcolor{currentfill}%
\pgfsetlinewidth{0.803000pt}%
\definecolor{currentstroke}{rgb}{0.000000,0.000000,0.000000}%
\pgfsetstrokecolor{currentstroke}%
\pgfsetdash{}{0pt}%
\pgfsys@defobject{currentmarker}{\pgfqpoint{-0.048611in}{0.000000in}}{\pgfqpoint{-0.000000in}{0.000000in}}{%
\pgfpathmoveto{\pgfqpoint{-0.000000in}{0.000000in}}%
\pgfpathlineto{\pgfqpoint{-0.048611in}{0.000000in}}%
\pgfusepath{stroke,fill}%
}%
\begin{pgfscope}%
\pgfsys@transformshift{4.554414in}{0.466172in}%
\pgfsys@useobject{currentmarker}{}%
\end{pgfscope}%
\end{pgfscope}%
\begin{pgfscope}%
\pgfpathrectangle{\pgfqpoint{4.554414in}{0.320625in}}{\pgfqpoint{1.815586in}{1.164375in}}%
\pgfusepath{clip}%
\pgfsetrectcap%
\pgfsetroundjoin%
\pgfsetlinewidth{0.803000pt}%
\definecolor{currentstroke}{rgb}{0.909804,0.909804,0.909804}%
\pgfsetstrokecolor{currentstroke}%
\pgfsetdash{}{0pt}%
\pgfpathmoveto{\pgfqpoint{4.554414in}{0.660234in}}%
\pgfpathlineto{\pgfqpoint{6.370000in}{0.660234in}}%
\pgfusepath{stroke}%
\end{pgfscope}%
\begin{pgfscope}%
\pgfsetbuttcap%
\pgfsetroundjoin%
\definecolor{currentfill}{rgb}{0.000000,0.000000,0.000000}%
\pgfsetfillcolor{currentfill}%
\pgfsetlinewidth{0.803000pt}%
\definecolor{currentstroke}{rgb}{0.000000,0.000000,0.000000}%
\pgfsetstrokecolor{currentstroke}%
\pgfsetdash{}{0pt}%
\pgfsys@defobject{currentmarker}{\pgfqpoint{-0.048611in}{0.000000in}}{\pgfqpoint{-0.000000in}{0.000000in}}{%
\pgfpathmoveto{\pgfqpoint{-0.000000in}{0.000000in}}%
\pgfpathlineto{\pgfqpoint{-0.048611in}{0.000000in}}%
\pgfusepath{stroke,fill}%
}%
\begin{pgfscope}%
\pgfsys@transformshift{4.554414in}{0.660234in}%
\pgfsys@useobject{currentmarker}{}%
\end{pgfscope}%
\end{pgfscope}%
\begin{pgfscope}%
\pgfpathrectangle{\pgfqpoint{4.554414in}{0.320625in}}{\pgfqpoint{1.815586in}{1.164375in}}%
\pgfusepath{clip}%
\pgfsetrectcap%
\pgfsetroundjoin%
\pgfsetlinewidth{0.803000pt}%
\definecolor{currentstroke}{rgb}{0.909804,0.909804,0.909804}%
\pgfsetstrokecolor{currentstroke}%
\pgfsetdash{}{0pt}%
\pgfpathmoveto{\pgfqpoint{4.554414in}{0.854297in}}%
\pgfpathlineto{\pgfqpoint{6.370000in}{0.854297in}}%
\pgfusepath{stroke}%
\end{pgfscope}%
\begin{pgfscope}%
\pgfsetbuttcap%
\pgfsetroundjoin%
\definecolor{currentfill}{rgb}{0.000000,0.000000,0.000000}%
\pgfsetfillcolor{currentfill}%
\pgfsetlinewidth{0.803000pt}%
\definecolor{currentstroke}{rgb}{0.000000,0.000000,0.000000}%
\pgfsetstrokecolor{currentstroke}%
\pgfsetdash{}{0pt}%
\pgfsys@defobject{currentmarker}{\pgfqpoint{-0.048611in}{0.000000in}}{\pgfqpoint{-0.000000in}{0.000000in}}{%
\pgfpathmoveto{\pgfqpoint{-0.000000in}{0.000000in}}%
\pgfpathlineto{\pgfqpoint{-0.048611in}{0.000000in}}%
\pgfusepath{stroke,fill}%
}%
\begin{pgfscope}%
\pgfsys@transformshift{4.554414in}{0.854297in}%
\pgfsys@useobject{currentmarker}{}%
\end{pgfscope}%
\end{pgfscope}%
\begin{pgfscope}%
\pgfpathrectangle{\pgfqpoint{4.554414in}{0.320625in}}{\pgfqpoint{1.815586in}{1.164375in}}%
\pgfusepath{clip}%
\pgfsetrectcap%
\pgfsetroundjoin%
\pgfsetlinewidth{0.803000pt}%
\definecolor{currentstroke}{rgb}{0.909804,0.909804,0.909804}%
\pgfsetstrokecolor{currentstroke}%
\pgfsetdash{}{0pt}%
\pgfpathmoveto{\pgfqpoint{4.554414in}{1.048359in}}%
\pgfpathlineto{\pgfqpoint{6.370000in}{1.048359in}}%
\pgfusepath{stroke}%
\end{pgfscope}%
\begin{pgfscope}%
\pgfsetbuttcap%
\pgfsetroundjoin%
\definecolor{currentfill}{rgb}{0.000000,0.000000,0.000000}%
\pgfsetfillcolor{currentfill}%
\pgfsetlinewidth{0.803000pt}%
\definecolor{currentstroke}{rgb}{0.000000,0.000000,0.000000}%
\pgfsetstrokecolor{currentstroke}%
\pgfsetdash{}{0pt}%
\pgfsys@defobject{currentmarker}{\pgfqpoint{-0.048611in}{0.000000in}}{\pgfqpoint{-0.000000in}{0.000000in}}{%
\pgfpathmoveto{\pgfqpoint{-0.000000in}{0.000000in}}%
\pgfpathlineto{\pgfqpoint{-0.048611in}{0.000000in}}%
\pgfusepath{stroke,fill}%
}%
\begin{pgfscope}%
\pgfsys@transformshift{4.554414in}{1.048359in}%
\pgfsys@useobject{currentmarker}{}%
\end{pgfscope}%
\end{pgfscope}%
\begin{pgfscope}%
\pgfpathrectangle{\pgfqpoint{4.554414in}{0.320625in}}{\pgfqpoint{1.815586in}{1.164375in}}%
\pgfusepath{clip}%
\pgfsetrectcap%
\pgfsetroundjoin%
\pgfsetlinewidth{0.803000pt}%
\definecolor{currentstroke}{rgb}{0.909804,0.909804,0.909804}%
\pgfsetstrokecolor{currentstroke}%
\pgfsetdash{}{0pt}%
\pgfpathmoveto{\pgfqpoint{4.554414in}{1.242422in}}%
\pgfpathlineto{\pgfqpoint{6.370000in}{1.242422in}}%
\pgfusepath{stroke}%
\end{pgfscope}%
\begin{pgfscope}%
\pgfsetbuttcap%
\pgfsetroundjoin%
\definecolor{currentfill}{rgb}{0.000000,0.000000,0.000000}%
\pgfsetfillcolor{currentfill}%
\pgfsetlinewidth{0.803000pt}%
\definecolor{currentstroke}{rgb}{0.000000,0.000000,0.000000}%
\pgfsetstrokecolor{currentstroke}%
\pgfsetdash{}{0pt}%
\pgfsys@defobject{currentmarker}{\pgfqpoint{-0.048611in}{0.000000in}}{\pgfqpoint{-0.000000in}{0.000000in}}{%
\pgfpathmoveto{\pgfqpoint{-0.000000in}{0.000000in}}%
\pgfpathlineto{\pgfqpoint{-0.048611in}{0.000000in}}%
\pgfusepath{stroke,fill}%
}%
\begin{pgfscope}%
\pgfsys@transformshift{4.554414in}{1.242422in}%
\pgfsys@useobject{currentmarker}{}%
\end{pgfscope}%
\end{pgfscope}%
\begin{pgfscope}%
\pgfpathrectangle{\pgfqpoint{4.554414in}{0.320625in}}{\pgfqpoint{1.815586in}{1.164375in}}%
\pgfusepath{clip}%
\pgfsetrectcap%
\pgfsetroundjoin%
\pgfsetlinewidth{0.803000pt}%
\definecolor{currentstroke}{rgb}{0.909804,0.909804,0.909804}%
\pgfsetstrokecolor{currentstroke}%
\pgfsetdash{}{0pt}%
\pgfpathmoveto{\pgfqpoint{4.554414in}{1.436484in}}%
\pgfpathlineto{\pgfqpoint{6.370000in}{1.436484in}}%
\pgfusepath{stroke}%
\end{pgfscope}%
\begin{pgfscope}%
\pgfsetbuttcap%
\pgfsetroundjoin%
\definecolor{currentfill}{rgb}{0.000000,0.000000,0.000000}%
\pgfsetfillcolor{currentfill}%
\pgfsetlinewidth{0.803000pt}%
\definecolor{currentstroke}{rgb}{0.000000,0.000000,0.000000}%
\pgfsetstrokecolor{currentstroke}%
\pgfsetdash{}{0pt}%
\pgfsys@defobject{currentmarker}{\pgfqpoint{-0.048611in}{0.000000in}}{\pgfqpoint{-0.000000in}{0.000000in}}{%
\pgfpathmoveto{\pgfqpoint{-0.000000in}{0.000000in}}%
\pgfpathlineto{\pgfqpoint{-0.048611in}{0.000000in}}%
\pgfusepath{stroke,fill}%
}%
\begin{pgfscope}%
\pgfsys@transformshift{4.554414in}{1.436484in}%
\pgfsys@useobject{currentmarker}{}%
\end{pgfscope}%
\end{pgfscope}%
\begin{pgfscope}%
\pgfpathrectangle{\pgfqpoint{4.554414in}{0.320625in}}{\pgfqpoint{1.815586in}{1.164375in}}%
\pgfusepath{clip}%
\pgfsetrectcap%
\pgfsetroundjoin%
\pgfsetlinewidth{1.304875pt}%
\definecolor{currentstroke}{rgb}{0.000000,0.000000,0.000000}%
\pgfsetstrokecolor{currentstroke}%
\pgfsetdash{}{0pt}%
\pgfpathmoveto{\pgfqpoint{4.554414in}{1.451682in}}%
\pgfpathlineto{\pgfqpoint{4.561676in}{1.436391in}}%
\pgfpathlineto{\pgfqpoint{4.576201in}{1.406957in}}%
\pgfpathlineto{\pgfqpoint{4.594356in}{1.366851in}}%
\pgfpathlineto{\pgfqpoint{4.703292in}{1.119529in}}%
\pgfpathlineto{\pgfqpoint{4.750497in}{1.017420in}}%
\pgfpathlineto{\pgfqpoint{4.804965in}{0.902493in}}%
\pgfpathlineto{\pgfqpoint{4.870326in}{0.767350in}}%
\pgfpathlineto{\pgfqpoint{4.950211in}{0.604841in}}%
\pgfpathlineto{\pgfqpoint{4.975630in}{0.554375in}}%
\pgfpathlineto{\pgfqpoint{4.979261in}{0.551170in}}%
\pgfpathlineto{\pgfqpoint{4.986523in}{0.534733in}}%
\pgfpathlineto{\pgfqpoint{4.990154in}{0.529633in}}%
\pgfpathlineto{\pgfqpoint{4.993785in}{0.518438in}}%
\pgfpathlineto{\pgfqpoint{4.997417in}{0.509829in}}%
\pgfpathlineto{\pgfqpoint{5.001048in}{0.509498in}}%
\pgfpathlineto{\pgfqpoint{5.004679in}{0.513363in}}%
\pgfpathlineto{\pgfqpoint{5.008310in}{0.514557in}}%
\pgfpathlineto{\pgfqpoint{5.011941in}{0.485537in}}%
\pgfpathlineto{\pgfqpoint{5.015573in}{0.480181in}}%
\pgfpathlineto{\pgfqpoint{5.019204in}{0.524005in}}%
\pgfpathlineto{\pgfqpoint{5.022835in}{0.459285in}}%
\pgfpathlineto{\pgfqpoint{5.026466in}{0.464112in}}%
\pgfpathlineto{\pgfqpoint{5.030097in}{0.492551in}}%
\pgfpathlineto{\pgfqpoint{5.033728in}{0.465642in}}%
\pgfpathlineto{\pgfqpoint{5.037360in}{0.466385in}}%
\pgfpathlineto{\pgfqpoint{5.040991in}{0.507495in}}%
\pgfpathlineto{\pgfqpoint{5.044622in}{0.511324in}}%
\pgfpathlineto{\pgfqpoint{5.048253in}{0.517848in}}%
\pgfpathlineto{\pgfqpoint{5.051884in}{0.478452in}}%
\pgfpathlineto{\pgfqpoint{5.055515in}{0.457535in}}%
\pgfpathlineto{\pgfqpoint{5.059147in}{0.441742in}}%
\pgfpathlineto{\pgfqpoint{5.062778in}{0.465675in}}%
\pgfpathlineto{\pgfqpoint{5.066409in}{0.504912in}}%
\pgfpathlineto{\pgfqpoint{5.070040in}{0.520324in}}%
\pgfpathlineto{\pgfqpoint{5.073671in}{0.418200in}}%
\pgfpathlineto{\pgfqpoint{5.077302in}{0.440317in}}%
\pgfpathlineto{\pgfqpoint{5.080934in}{0.493635in}}%
\pgfpathlineto{\pgfqpoint{5.084565in}{0.446004in}}%
\pgfpathlineto{\pgfqpoint{5.088196in}{0.469780in}}%
\pgfpathlineto{\pgfqpoint{5.091827in}{0.509087in}}%
\pgfpathlineto{\pgfqpoint{5.095458in}{0.511889in}}%
\pgfpathlineto{\pgfqpoint{5.099090in}{0.512622in}}%
\pgfpathlineto{\pgfqpoint{5.102721in}{0.479745in}}%
\pgfpathlineto{\pgfqpoint{5.106352in}{0.462385in}}%
\pgfpathlineto{\pgfqpoint{5.109983in}{0.457365in}}%
\pgfpathlineto{\pgfqpoint{5.113614in}{0.467886in}}%
\pgfpathlineto{\pgfqpoint{5.117245in}{0.497596in}}%
\pgfpathlineto{\pgfqpoint{5.120877in}{0.510076in}}%
\pgfpathlineto{\pgfqpoint{5.124508in}{0.484444in}}%
\pgfpathlineto{\pgfqpoint{5.128139in}{0.342956in}}%
\pgfpathlineto{\pgfqpoint{5.131770in}{0.521671in}}%
\pgfpathlineto{\pgfqpoint{5.139032in}{0.377062in}}%
\pgfpathlineto{\pgfqpoint{5.142664in}{0.491192in}}%
\pgfpathlineto{\pgfqpoint{5.146295in}{0.468233in}}%
\pgfpathlineto{\pgfqpoint{5.149926in}{0.427206in}}%
\pgfpathlineto{\pgfqpoint{5.153557in}{0.490982in}}%
\pgfpathlineto{\pgfqpoint{5.157188in}{0.512916in}}%
\pgfpathlineto{\pgfqpoint{5.160819in}{0.530822in}}%
\pgfpathlineto{\pgfqpoint{5.164451in}{0.431768in}}%
\pgfpathlineto{\pgfqpoint{5.168082in}{0.445244in}}%
\pgfpathlineto{\pgfqpoint{5.171713in}{0.470169in}}%
\pgfpathlineto{\pgfqpoint{5.175344in}{0.440866in}}%
\pgfpathlineto{\pgfqpoint{5.178975in}{0.497626in}}%
\pgfpathlineto{\pgfqpoint{5.182606in}{0.516291in}}%
\pgfpathlineto{\pgfqpoint{5.186238in}{0.456571in}}%
\pgfpathlineto{\pgfqpoint{5.189869in}{0.408678in}}%
\pgfpathlineto{\pgfqpoint{5.193500in}{0.499026in}}%
\pgfpathlineto{\pgfqpoint{5.197131in}{0.474594in}}%
\pgfpathlineto{\pgfqpoint{5.200762in}{0.351811in}}%
\pgfpathlineto{\pgfqpoint{5.204394in}{0.494703in}}%
\pgfpathlineto{\pgfqpoint{5.208025in}{0.505202in}}%
\pgfpathlineto{\pgfqpoint{5.211656in}{0.513827in}}%
\pgfpathlineto{\pgfqpoint{5.215287in}{0.454038in}}%
\pgfpathlineto{\pgfqpoint{5.222549in}{0.482879in}}%
\pgfpathlineto{\pgfqpoint{5.226181in}{0.434018in}}%
\pgfpathlineto{\pgfqpoint{5.229812in}{0.492243in}}%
\pgfpathlineto{\pgfqpoint{5.233443in}{0.514639in}}%
\pgfpathlineto{\pgfqpoint{5.237074in}{0.492315in}}%
\pgfpathlineto{\pgfqpoint{5.240705in}{0.434061in}}%
\pgfpathlineto{\pgfqpoint{5.244336in}{0.512462in}}%
\pgfpathlineto{\pgfqpoint{5.247968in}{0.466702in}}%
\pgfpathlineto{\pgfqpoint{5.251599in}{0.400945in}}%
\pgfpathlineto{\pgfqpoint{5.255230in}{0.491863in}}%
\pgfpathlineto{\pgfqpoint{5.258861in}{0.484371in}}%
\pgfpathlineto{\pgfqpoint{5.262492in}{0.452742in}}%
\pgfpathlineto{\pgfqpoint{5.266123in}{0.463927in}}%
\pgfpathlineto{\pgfqpoint{5.269755in}{0.500935in}}%
\pgfpathlineto{\pgfqpoint{5.273386in}{0.526521in}}%
\pgfpathlineto{\pgfqpoint{5.276770in}{0.315625in}}%
\pgfpathmoveto{\pgfqpoint{5.277394in}{0.315625in}}%
\pgfpathlineto{\pgfqpoint{5.280648in}{0.448428in}}%
\pgfpathlineto{\pgfqpoint{5.284279in}{0.488955in}}%
\pgfpathlineto{\pgfqpoint{5.287910in}{0.316776in}}%
\pgfpathlineto{\pgfqpoint{5.291542in}{0.486728in}}%
\pgfpathlineto{\pgfqpoint{5.295173in}{0.514232in}}%
\pgfpathlineto{\pgfqpoint{5.298804in}{0.509760in}}%
\pgfpathlineto{\pgfqpoint{5.302435in}{0.497120in}}%
\pgfpathlineto{\pgfqpoint{5.306066in}{0.497114in}}%
\pgfpathlineto{\pgfqpoint{5.309698in}{0.465392in}}%
\pgfpathlineto{\pgfqpoint{5.313329in}{0.447629in}}%
\pgfpathlineto{\pgfqpoint{5.316960in}{0.480607in}}%
\pgfpathlineto{\pgfqpoint{5.320591in}{0.490085in}}%
\pgfpathlineto{\pgfqpoint{5.324222in}{0.486426in}}%
\pgfpathlineto{\pgfqpoint{5.327853in}{0.421428in}}%
\pgfpathlineto{\pgfqpoint{5.331485in}{0.481065in}}%
\pgfpathlineto{\pgfqpoint{5.335116in}{0.529277in}}%
\pgfpathlineto{\pgfqpoint{5.338747in}{0.444804in}}%
\pgfpathlineto{\pgfqpoint{5.342378in}{0.396011in}}%
\pgfpathlineto{\pgfqpoint{5.346009in}{0.490783in}}%
\pgfpathlineto{\pgfqpoint{5.349640in}{0.433673in}}%
\pgfpathlineto{\pgfqpoint{5.353272in}{0.475194in}}%
\pgfpathlineto{\pgfqpoint{5.356903in}{0.506878in}}%
\pgfpathlineto{\pgfqpoint{5.360534in}{0.518908in}}%
\pgfpathlineto{\pgfqpoint{5.364165in}{0.524838in}}%
\pgfpathlineto{\pgfqpoint{5.367796in}{0.462773in}}%
\pgfpathlineto{\pgfqpoint{5.371427in}{0.450945in}}%
\pgfpathlineto{\pgfqpoint{5.375059in}{0.466443in}}%
\pgfpathlineto{\pgfqpoint{5.378690in}{0.463867in}}%
\pgfpathlineto{\pgfqpoint{5.382321in}{0.492862in}}%
\pgfpathlineto{\pgfqpoint{5.385952in}{0.503998in}}%
\pgfpathlineto{\pgfqpoint{5.389583in}{0.492405in}}%
\pgfpathlineto{\pgfqpoint{5.393215in}{0.407370in}}%
\pgfpathlineto{\pgfqpoint{5.396846in}{0.521897in}}%
\pgfpathlineto{\pgfqpoint{5.400477in}{0.460120in}}%
\pgfpathlineto{\pgfqpoint{5.404108in}{0.418431in}}%
\pgfpathlineto{\pgfqpoint{5.407739in}{0.487331in}}%
\pgfpathlineto{\pgfqpoint{5.411370in}{0.458074in}}%
\pgfpathlineto{\pgfqpoint{5.415002in}{0.447970in}}%
\pgfpathlineto{\pgfqpoint{5.418633in}{0.487372in}}%
\pgfpathlineto{\pgfqpoint{5.422264in}{0.470566in}}%
\pgfpathlineto{\pgfqpoint{5.425895in}{0.486859in}}%
\pgfpathlineto{\pgfqpoint{5.429526in}{0.463694in}}%
\pgfpathlineto{\pgfqpoint{5.433157in}{0.481374in}}%
\pgfpathlineto{\pgfqpoint{5.436789in}{0.485994in}}%
\pgfpathlineto{\pgfqpoint{5.440420in}{0.463514in}}%
\pgfpathlineto{\pgfqpoint{5.444051in}{0.388108in}}%
\pgfpathlineto{\pgfqpoint{5.447682in}{0.507455in}}%
\pgfpathlineto{\pgfqpoint{5.451313in}{0.402578in}}%
\pgfpathlineto{\pgfqpoint{5.454944in}{0.460956in}}%
\pgfpathlineto{\pgfqpoint{5.458576in}{0.497893in}}%
\pgfpathlineto{\pgfqpoint{5.462207in}{0.505191in}}%
\pgfpathlineto{\pgfqpoint{5.465838in}{0.500196in}}%
\pgfpathlineto{\pgfqpoint{5.469469in}{0.484180in}}%
\pgfpathlineto{\pgfqpoint{5.473100in}{0.459496in}}%
\pgfpathlineto{\pgfqpoint{5.476731in}{0.426944in}}%
\pgfpathlineto{\pgfqpoint{5.483994in}{0.507707in}}%
\pgfpathlineto{\pgfqpoint{5.487625in}{0.523577in}}%
\pgfpathlineto{\pgfqpoint{5.491256in}{0.496650in}}%
\pgfpathlineto{\pgfqpoint{5.494887in}{0.449461in}}%
\pgfpathlineto{\pgfqpoint{5.498519in}{0.516738in}}%
\pgfpathlineto{\pgfqpoint{5.502150in}{0.414765in}}%
\pgfpathlineto{\pgfqpoint{5.505781in}{0.439727in}}%
\pgfpathlineto{\pgfqpoint{5.509412in}{0.490886in}}%
\pgfpathlineto{\pgfqpoint{5.513043in}{0.485683in}}%
\pgfpathlineto{\pgfqpoint{5.516674in}{0.424052in}}%
\pgfpathlineto{\pgfqpoint{5.520306in}{0.477901in}}%
\pgfpathlineto{\pgfqpoint{5.523937in}{0.499670in}}%
\pgfpathlineto{\pgfqpoint{5.527568in}{0.526188in}}%
\pgfpathlineto{\pgfqpoint{5.531199in}{0.456649in}}%
\pgfpathlineto{\pgfqpoint{5.534830in}{0.452180in}}%
\pgfpathlineto{\pgfqpoint{5.538461in}{0.451776in}}%
\pgfpathlineto{\pgfqpoint{5.542093in}{0.433505in}}%
\pgfpathlineto{\pgfqpoint{5.545724in}{0.506664in}}%
\pgfpathlineto{\pgfqpoint{5.549355in}{0.525630in}}%
\pgfpathlineto{\pgfqpoint{5.552986in}{0.454437in}}%
\pgfpathlineto{\pgfqpoint{5.556617in}{0.363135in}}%
\pgfpathlineto{\pgfqpoint{5.560248in}{0.486192in}}%
\pgfpathlineto{\pgfqpoint{5.563880in}{0.452781in}}%
\pgfpathlineto{\pgfqpoint{5.567511in}{0.462481in}}%
\pgfpathlineto{\pgfqpoint{5.571142in}{0.494537in}}%
\pgfpathlineto{\pgfqpoint{5.574773in}{0.470124in}}%
\pgfpathlineto{\pgfqpoint{5.578404in}{0.481531in}}%
\pgfpathlineto{\pgfqpoint{5.582035in}{0.469929in}}%
\pgfpathlineto{\pgfqpoint{5.585667in}{0.479861in}}%
\pgfpathlineto{\pgfqpoint{5.589298in}{0.476001in}}%
\pgfpathlineto{\pgfqpoint{5.592929in}{0.448711in}}%
\pgfpathlineto{\pgfqpoint{5.595336in}{0.315625in}}%
\pgfpathmoveto{\pgfqpoint{5.597508in}{0.315625in}}%
\pgfpathlineto{\pgfqpoint{5.600191in}{0.507100in}}%
\pgfpathlineto{\pgfqpoint{5.603823in}{0.338516in}}%
\pgfpathlineto{\pgfqpoint{5.607454in}{0.466962in}}%
\pgfpathlineto{\pgfqpoint{5.611085in}{0.498104in}}%
\pgfpathlineto{\pgfqpoint{5.614716in}{0.504090in}}%
\pgfpathlineto{\pgfqpoint{5.621978in}{0.489360in}}%
\pgfpathlineto{\pgfqpoint{5.625610in}{0.455450in}}%
\pgfpathlineto{\pgfqpoint{5.629241in}{0.387707in}}%
\pgfpathlineto{\pgfqpoint{5.632872in}{0.475231in}}%
\pgfpathlineto{\pgfqpoint{5.636503in}{0.509806in}}%
\pgfpathlineto{\pgfqpoint{5.640134in}{0.523746in}}%
\pgfpathlineto{\pgfqpoint{5.643686in}{0.315625in}}%
\pgfpathmoveto{\pgfqpoint{5.643894in}{0.315625in}}%
\pgfpathlineto{\pgfqpoint{5.647397in}{0.442217in}}%
\pgfpathlineto{\pgfqpoint{5.651028in}{0.482985in}}%
\pgfpathlineto{\pgfqpoint{5.654659in}{0.400767in}}%
\pgfpathlineto{\pgfqpoint{5.658290in}{0.488329in}}%
\pgfpathlineto{\pgfqpoint{5.661921in}{0.515090in}}%
\pgfpathlineto{\pgfqpoint{5.665552in}{0.473008in}}%
\pgfpathlineto{\pgfqpoint{5.669184in}{0.442825in}}%
\pgfpathlineto{\pgfqpoint{5.672815in}{0.487854in}}%
\pgfpathlineto{\pgfqpoint{5.676446in}{0.475973in}}%
\pgfpathlineto{\pgfqpoint{5.680077in}{0.419580in}}%
\pgfpathlineto{\pgfqpoint{5.687340in}{0.507930in}}%
\pgfpathlineto{\pgfqpoint{5.690971in}{0.534151in}}%
\pgfpathlineto{\pgfqpoint{5.694602in}{0.439344in}}%
\pgfpathlineto{\pgfqpoint{5.698233in}{0.401880in}}%
\pgfpathlineto{\pgfqpoint{5.701864in}{0.489061in}}%
\pgfpathlineto{\pgfqpoint{5.705495in}{0.350594in}}%
\pgfpathlineto{\pgfqpoint{5.709127in}{0.486103in}}%
\pgfpathlineto{\pgfqpoint{5.712758in}{0.510341in}}%
\pgfpathlineto{\pgfqpoint{5.716389in}{0.516660in}}%
\pgfpathlineto{\pgfqpoint{5.720020in}{0.511471in}}%
\pgfpathlineto{\pgfqpoint{5.727282in}{0.459333in}}%
\pgfpathlineto{\pgfqpoint{5.730914in}{0.468337in}}%
\pgfpathlineto{\pgfqpoint{5.734545in}{0.475493in}}%
\pgfpathlineto{\pgfqpoint{5.738176in}{0.481419in}}%
\pgfpathlineto{\pgfqpoint{5.741807in}{0.470847in}}%
\pgfpathlineto{\pgfqpoint{5.745438in}{0.433317in}}%
\pgfpathlineto{\pgfqpoint{5.749069in}{0.354262in}}%
\pgfpathlineto{\pgfqpoint{5.752701in}{0.506124in}}%
\pgfpathlineto{\pgfqpoint{5.756332in}{0.389663in}}%
\pgfpathlineto{\pgfqpoint{5.759963in}{0.471139in}}%
\pgfpathlineto{\pgfqpoint{5.763594in}{0.496431in}}%
\pgfpathlineto{\pgfqpoint{5.767225in}{0.501615in}}%
\pgfpathlineto{\pgfqpoint{5.770856in}{0.492194in}}%
\pgfpathlineto{\pgfqpoint{5.774488in}{0.494311in}}%
\pgfpathlineto{\pgfqpoint{5.778119in}{0.449187in}}%
\pgfpathlineto{\pgfqpoint{5.781750in}{0.371138in}}%
\pgfpathlineto{\pgfqpoint{5.785381in}{0.481306in}}%
\pgfpathlineto{\pgfqpoint{5.789012in}{0.511761in}}%
\pgfpathlineto{\pgfqpoint{5.792644in}{0.523555in}}%
\pgfpathlineto{\pgfqpoint{5.796275in}{0.400580in}}%
\pgfpathlineto{\pgfqpoint{5.799906in}{0.445614in}}%
\pgfpathlineto{\pgfqpoint{5.803537in}{0.475345in}}%
\pgfpathlineto{\pgfqpoint{5.806896in}{0.315625in}}%
\pgfpathmoveto{\pgfqpoint{5.807412in}{0.315625in}}%
\pgfpathlineto{\pgfqpoint{5.810799in}{0.495564in}}%
\pgfpathlineto{\pgfqpoint{5.814431in}{0.518597in}}%
\pgfpathlineto{\pgfqpoint{5.818062in}{0.468783in}}%
\pgfpathlineto{\pgfqpoint{5.821693in}{0.426655in}}%
\pgfpathlineto{\pgfqpoint{5.825324in}{0.488871in}}%
\pgfpathlineto{\pgfqpoint{5.828955in}{0.471255in}}%
\pgfpathlineto{\pgfqpoint{5.832586in}{0.323379in}}%
\pgfpathlineto{\pgfqpoint{5.836218in}{0.474037in}}%
\pgfpathlineto{\pgfqpoint{5.839849in}{0.468625in}}%
\pgfpathlineto{\pgfqpoint{5.843480in}{0.493358in}}%
\pgfpathlineto{\pgfqpoint{5.847111in}{0.454094in}}%
\pgfpathlineto{\pgfqpoint{5.850742in}{0.480810in}}%
\pgfpathlineto{\pgfqpoint{5.854373in}{0.492142in}}%
\pgfpathlineto{\pgfqpoint{5.858005in}{0.472927in}}%
\pgfpathlineto{\pgfqpoint{5.861636in}{0.449884in}}%
\pgfpathlineto{\pgfqpoint{5.865267in}{0.520609in}}%
\pgfpathlineto{\pgfqpoint{5.868898in}{0.486270in}}%
\pgfpathlineto{\pgfqpoint{5.872529in}{0.471027in}}%
\pgfpathlineto{\pgfqpoint{5.876160in}{0.492160in}}%
\pgfpathlineto{\pgfqpoint{5.879792in}{0.454094in}}%
\pgfpathlineto{\pgfqpoint{5.883423in}{0.431555in}}%
\pgfpathlineto{\pgfqpoint{5.887054in}{0.486083in}}%
\pgfpathlineto{\pgfqpoint{5.890685in}{0.519084in}}%
\pgfpathlineto{\pgfqpoint{5.894316in}{0.534858in}}%
\pgfpathlineto{\pgfqpoint{5.897948in}{0.466771in}}%
\pgfpathlineto{\pgfqpoint{5.901579in}{0.346717in}}%
\pgfpathlineto{\pgfqpoint{5.905210in}{0.499819in}}%
\pgfpathlineto{\pgfqpoint{5.908841in}{0.418516in}}%
\pgfpathlineto{\pgfqpoint{5.912472in}{0.477511in}}%
\pgfpathlineto{\pgfqpoint{5.916103in}{0.498779in}}%
\pgfpathlineto{\pgfqpoint{5.919735in}{0.502227in}}%
\pgfpathlineto{\pgfqpoint{5.923366in}{0.491622in}}%
\pgfpathlineto{\pgfqpoint{5.926997in}{0.498161in}}%
\pgfpathlineto{\pgfqpoint{5.930628in}{0.439141in}}%
\pgfpathlineto{\pgfqpoint{5.934259in}{0.420002in}}%
\pgfpathlineto{\pgfqpoint{5.937890in}{0.485432in}}%
\pgfpathlineto{\pgfqpoint{5.941522in}{0.512931in}}%
\pgfpathlineto{\pgfqpoint{5.945153in}{0.522414in}}%
\pgfpathlineto{\pgfqpoint{5.948784in}{0.432104in}}%
\pgfpathlineto{\pgfqpoint{5.952415in}{0.447406in}}%
\pgfpathlineto{\pgfqpoint{5.956046in}{0.466179in}}%
\pgfpathlineto{\pgfqpoint{5.959677in}{0.405314in}}%
\pgfpathlineto{\pgfqpoint{5.963309in}{0.501110in}}%
\pgfpathlineto{\pgfqpoint{5.966940in}{0.521563in}}%
\pgfpathlineto{\pgfqpoint{5.974202in}{0.405685in}}%
\pgfpathlineto{\pgfqpoint{5.977833in}{0.488719in}}%
\pgfpathlineto{\pgfqpoint{5.981465in}{0.465328in}}%
\pgfpathlineto{\pgfqpoint{5.985096in}{0.426439in}}%
\pgfpathlineto{\pgfqpoint{5.988727in}{0.483749in}}%
\pgfpathlineto{\pgfqpoint{5.992358in}{0.469722in}}%
\pgfpathlineto{\pgfqpoint{5.995989in}{0.488759in}}%
\pgfpathlineto{\pgfqpoint{5.999620in}{0.462149in}}%
\pgfpathlineto{\pgfqpoint{6.003252in}{0.480859in}}%
\pgfpathlineto{\pgfqpoint{6.006883in}{0.486422in}}%
\pgfpathlineto{\pgfqpoint{6.010514in}{0.464554in}}%
\pgfpathlineto{\pgfqpoint{6.014145in}{0.390554in}}%
\pgfpathlineto{\pgfqpoint{6.017776in}{0.507378in}}%
\pgfpathlineto{\pgfqpoint{6.021407in}{0.409657in}}%
\pgfpathlineto{\pgfqpoint{6.025039in}{0.459452in}}%
\pgfpathlineto{\pgfqpoint{6.028670in}{0.497892in}}%
\pgfpathlineto{\pgfqpoint{6.032301in}{0.505363in}}%
\pgfpathlineto{\pgfqpoint{6.035932in}{0.500787in}}%
\pgfpathlineto{\pgfqpoint{6.039563in}{0.483005in}}%
\pgfpathlineto{\pgfqpoint{6.043194in}{0.460305in}}%
\pgfpathlineto{\pgfqpoint{6.046826in}{0.432394in}}%
\pgfpathlineto{\pgfqpoint{6.050457in}{0.466803in}}%
\pgfpathlineto{\pgfqpoint{6.054088in}{0.507142in}}%
\pgfpathlineto{\pgfqpoint{6.057719in}{0.523425in}}%
\pgfpathlineto{\pgfqpoint{6.061350in}{0.496857in}}%
\pgfpathlineto{\pgfqpoint{6.064981in}{0.450422in}}%
\pgfpathlineto{\pgfqpoint{6.068613in}{0.516368in}}%
\pgfpathlineto{\pgfqpoint{6.072244in}{0.419410in}}%
\pgfpathlineto{\pgfqpoint{6.075875in}{0.437857in}}%
\pgfpathlineto{\pgfqpoint{6.079506in}{0.491213in}}%
\pgfpathlineto{\pgfqpoint{6.083137in}{0.486387in}}%
\pgfpathlineto{\pgfqpoint{6.086769in}{0.431352in}}%
\pgfpathlineto{\pgfqpoint{6.090400in}{0.476003in}}%
\pgfpathlineto{\pgfqpoint{6.094031in}{0.499116in}}%
\pgfpathlineto{\pgfqpoint{6.097662in}{0.526138in}}%
\pgfpathlineto{\pgfqpoint{6.101293in}{0.454545in}}%
\pgfpathlineto{\pgfqpoint{6.104924in}{0.452592in}}%
\pgfpathlineto{\pgfqpoint{6.108556in}{0.454680in}}%
\pgfpathlineto{\pgfqpoint{6.112187in}{0.429681in}}%
\pgfpathlineto{\pgfqpoint{6.115818in}{0.505762in}}%
\pgfpathlineto{\pgfqpoint{6.119449in}{0.525143in}}%
\pgfpathlineto{\pgfqpoint{6.123080in}{0.455894in}}%
\pgfpathlineto{\pgfqpoint{6.126711in}{0.371081in}}%
\pgfpathlineto{\pgfqpoint{6.130343in}{0.486741in}}%
\pgfpathlineto{\pgfqpoint{6.133974in}{0.454616in}}%
\pgfpathlineto{\pgfqpoint{6.137605in}{0.459299in}}%
\pgfpathlineto{\pgfqpoint{6.141236in}{0.493392in}}%
\pgfpathlineto{\pgfqpoint{6.144867in}{0.470322in}}%
\pgfpathlineto{\pgfqpoint{6.148498in}{0.482455in}}%
\pgfpathlineto{\pgfqpoint{6.152130in}{0.469154in}}%
\pgfpathlineto{\pgfqpoint{6.155761in}{0.480131in}}%
\pgfpathlineto{\pgfqpoint{6.159392in}{0.477687in}}%
\pgfpathlineto{\pgfqpoint{6.163023in}{0.450872in}}%
\pgfpathlineto{\pgfqpoint{6.166152in}{0.315625in}}%
\pgfpathmoveto{\pgfqpoint{6.167024in}{0.315625in}}%
\pgfpathlineto{\pgfqpoint{6.170285in}{0.507217in}}%
\pgfpathlineto{\pgfqpoint{6.173917in}{0.357702in}}%
\pgfpathlineto{\pgfqpoint{6.177548in}{0.466197in}}%
\pgfpathlineto{\pgfqpoint{6.181179in}{0.498201in}}%
\pgfpathlineto{\pgfqpoint{6.184810in}{0.504185in}}%
\pgfpathlineto{\pgfqpoint{6.188441in}{0.497249in}}%
\pgfpathlineto{\pgfqpoint{6.192073in}{0.488726in}}%
\pgfpathlineto{\pgfqpoint{6.195704in}{0.456304in}}%
\pgfpathlineto{\pgfqpoint{6.199335in}{0.396739in}}%
\pgfpathlineto{\pgfqpoint{6.202966in}{0.474378in}}%
\pgfpathlineto{\pgfqpoint{6.206597in}{0.509447in}}%
\pgfpathlineto{\pgfqpoint{6.210228in}{0.523654in}}%
\pgfpathlineto{\pgfqpoint{6.213860in}{0.341769in}}%
\pgfpathlineto{\pgfqpoint{6.217491in}{0.441902in}}%
\pgfpathlineto{\pgfqpoint{6.221122in}{0.483940in}}%
\pgfpathlineto{\pgfqpoint{6.224753in}{0.406249in}}%
\pgfpathlineto{\pgfqpoint{6.228384in}{0.487223in}}%
\pgfpathlineto{\pgfqpoint{6.232015in}{0.514630in}}%
\pgfpathlineto{\pgfqpoint{6.235647in}{0.517043in}}%
\pgfpathlineto{\pgfqpoint{6.239278in}{0.516279in}}%
\pgfpathlineto{\pgfqpoint{6.242909in}{0.487490in}}%
\pgfpathlineto{\pgfqpoint{6.246540in}{0.453698in}}%
\pgfpathlineto{\pgfqpoint{6.250171in}{0.435518in}}%
\pgfpathlineto{\pgfqpoint{6.253802in}{0.475162in}}%
\pgfpathlineto{\pgfqpoint{6.257434in}{0.497850in}}%
\pgfpathlineto{\pgfqpoint{6.261065in}{0.504880in}}%
\pgfpathlineto{\pgfqpoint{6.264696in}{0.482991in}}%
\pgfpathlineto{\pgfqpoint{6.268327in}{0.390885in}}%
\pgfpathlineto{\pgfqpoint{6.271958in}{0.526443in}}%
\pgfpathlineto{\pgfqpoint{6.275590in}{0.433700in}}%
\pgfpathlineto{\pgfqpoint{6.279221in}{0.392336in}}%
\pgfpathlineto{\pgfqpoint{6.282852in}{0.485385in}}%
\pgfpathlineto{\pgfqpoint{6.286483in}{0.452317in}}%
\pgfpathlineto{\pgfqpoint{6.290114in}{0.467548in}}%
\pgfpathlineto{\pgfqpoint{6.293745in}{0.499424in}}%
\pgfpathlineto{\pgfqpoint{6.297377in}{0.468779in}}%
\pgfpathlineto{\pgfqpoint{6.301008in}{0.477129in}}%
\pgfpathlineto{\pgfqpoint{6.304639in}{0.473873in}}%
\pgfpathlineto{\pgfqpoint{6.308270in}{0.477184in}}%
\pgfpathlineto{\pgfqpoint{6.311901in}{0.462246in}}%
\pgfpathlineto{\pgfqpoint{6.315532in}{0.425152in}}%
\pgfpathlineto{\pgfqpoint{6.319164in}{0.380457in}}%
\pgfpathlineto{\pgfqpoint{6.322795in}{0.506022in}}%
\pgfpathlineto{\pgfqpoint{6.326426in}{0.368937in}}%
\pgfpathlineto{\pgfqpoint{6.330057in}{0.470813in}}%
\pgfpathlineto{\pgfqpoint{6.333688in}{0.497752in}}%
\pgfpathlineto{\pgfqpoint{6.337319in}{0.502733in}}%
\pgfpathlineto{\pgfqpoint{6.340951in}{0.493863in}}%
\pgfpathlineto{\pgfqpoint{6.344582in}{0.493173in}}%
\pgfpathlineto{\pgfqpoint{6.348213in}{0.450983in}}%
\pgfpathlineto{\pgfqpoint{6.351844in}{0.328470in}}%
\pgfpathlineto{\pgfqpoint{6.355475in}{0.479919in}}%
\pgfpathlineto{\pgfqpoint{6.359106in}{0.511213in}}%
\pgfpathlineto{\pgfqpoint{6.362738in}{0.523503in}}%
\pgfpathlineto{\pgfqpoint{6.366369in}{0.386703in}}%
\pgfpathlineto{\pgfqpoint{6.370000in}{0.457324in}}%
\pgfpathlineto{\pgfqpoint{6.370000in}{0.457324in}}%
\pgfusepath{stroke}%
\end{pgfscope}%
\begin{pgfscope}%
\pgfpathrectangle{\pgfqpoint{4.554414in}{0.320625in}}{\pgfqpoint{1.815586in}{1.164375in}}%
\pgfusepath{clip}%
\pgfsetrectcap%
\pgfsetroundjoin%
\pgfsetlinewidth{0.602250pt}%
\definecolor{currentstroke}{rgb}{0.141176,0.345098,0.650980}%
\pgfsetstrokecolor{currentstroke}%
\pgfsetstrokeopacity{0.120000}%
\pgfsetdash{}{0pt}%
\pgfpathmoveto{\pgfqpoint{4.554414in}{1.478120in}}%
\pgfpathlineto{\pgfqpoint{4.561676in}{1.457841in}}%
\pgfpathlineto{\pgfqpoint{4.568938in}{1.441891in}}%
\pgfpathlineto{\pgfqpoint{4.576201in}{1.429913in}}%
\pgfpathlineto{\pgfqpoint{4.587094in}{1.415698in}}%
\pgfpathlineto{\pgfqpoint{4.612512in}{1.386734in}}%
\pgfpathlineto{\pgfqpoint{4.637931in}{1.359163in}}%
\pgfpathlineto{\pgfqpoint{4.659718in}{1.337271in}}%
\pgfpathlineto{\pgfqpoint{4.681505in}{1.316920in}}%
\pgfpathlineto{\pgfqpoint{4.706923in}{1.294821in}}%
\pgfpathlineto{\pgfqpoint{4.735972in}{1.271185in}}%
\pgfpathlineto{\pgfqpoint{4.772284in}{1.243254in}}%
\pgfpathlineto{\pgfqpoint{4.819489in}{1.208556in}}%
\pgfpathlineto{\pgfqpoint{4.881219in}{1.164698in}}%
\pgfpathlineto{\pgfqpoint{4.968367in}{1.104251in}}%
\pgfpathlineto{\pgfqpoint{5.106352in}{1.010043in}}%
\pgfpathlineto{\pgfqpoint{5.375059in}{0.828161in}}%
\pgfpathlineto{\pgfqpoint{5.451313in}{0.776772in}}%
\pgfpathlineto{\pgfqpoint{5.469469in}{0.764421in}}%
\pgfpathlineto{\pgfqpoint{5.473100in}{0.762659in}}%
\pgfpathlineto{\pgfqpoint{5.476731in}{0.759540in}}%
\pgfpathlineto{\pgfqpoint{5.487625in}{0.752738in}}%
\pgfpathlineto{\pgfqpoint{5.494887in}{0.747457in}}%
\pgfpathlineto{\pgfqpoint{5.498519in}{0.746088in}}%
\pgfpathlineto{\pgfqpoint{5.502150in}{0.742572in}}%
\pgfpathlineto{\pgfqpoint{5.505781in}{0.740635in}}%
\pgfpathlineto{\pgfqpoint{5.509412in}{0.737453in}}%
\pgfpathlineto{\pgfqpoint{5.513043in}{0.737327in}}%
\pgfpathlineto{\pgfqpoint{5.516674in}{0.732551in}}%
\pgfpathlineto{\pgfqpoint{5.520306in}{0.731264in}}%
\pgfpathlineto{\pgfqpoint{5.523937in}{0.728101in}}%
\pgfpathlineto{\pgfqpoint{5.527568in}{0.728104in}}%
\pgfpathlineto{\pgfqpoint{5.531199in}{0.723400in}}%
\pgfpathlineto{\pgfqpoint{5.534830in}{0.721109in}}%
\pgfpathlineto{\pgfqpoint{5.538461in}{0.720763in}}%
\pgfpathlineto{\pgfqpoint{5.542093in}{0.717182in}}%
\pgfpathlineto{\pgfqpoint{5.545724in}{0.714802in}}%
\pgfpathlineto{\pgfqpoint{5.549355in}{0.710873in}}%
\pgfpathlineto{\pgfqpoint{5.552986in}{0.716141in}}%
\pgfpathlineto{\pgfqpoint{5.556617in}{0.705716in}}%
\pgfpathlineto{\pgfqpoint{5.560248in}{0.707107in}}%
\pgfpathlineto{\pgfqpoint{5.563880in}{0.701483in}}%
\pgfpathlineto{\pgfqpoint{5.567511in}{0.707961in}}%
\pgfpathlineto{\pgfqpoint{5.571142in}{0.696196in}}%
\pgfpathlineto{\pgfqpoint{5.574773in}{0.698464in}}%
\pgfpathlineto{\pgfqpoint{5.578404in}{0.697063in}}%
\pgfpathlineto{\pgfqpoint{5.582035in}{0.697314in}}%
\pgfpathlineto{\pgfqpoint{5.585667in}{0.692191in}}%
\pgfpathlineto{\pgfqpoint{5.589298in}{0.685073in}}%
\pgfpathlineto{\pgfqpoint{5.592929in}{0.698112in}}%
\pgfpathlineto{\pgfqpoint{5.596560in}{0.682499in}}%
\pgfpathlineto{\pgfqpoint{5.600191in}{0.686945in}}%
\pgfpathlineto{\pgfqpoint{5.603823in}{0.673770in}}%
\pgfpathlineto{\pgfqpoint{5.607454in}{0.699080in}}%
\pgfpathlineto{\pgfqpoint{5.611085in}{0.669800in}}%
\pgfpathlineto{\pgfqpoint{5.614716in}{0.680807in}}%
\pgfpathlineto{\pgfqpoint{5.618347in}{0.676135in}}%
\pgfpathlineto{\pgfqpoint{5.621978in}{0.687802in}}%
\pgfpathlineto{\pgfqpoint{5.625610in}{0.666484in}}%
\pgfpathlineto{\pgfqpoint{5.629241in}{0.672347in}}%
\pgfpathlineto{\pgfqpoint{5.632872in}{0.686593in}}%
\pgfpathlineto{\pgfqpoint{5.636503in}{0.669461in}}%
\pgfpathlineto{\pgfqpoint{5.640134in}{0.676835in}}%
\pgfpathlineto{\pgfqpoint{5.643765in}{0.646901in}}%
\pgfpathlineto{\pgfqpoint{5.647397in}{0.691606in}}%
\pgfpathlineto{\pgfqpoint{5.651028in}{0.645880in}}%
\pgfpathlineto{\pgfqpoint{5.654659in}{0.673846in}}%
\pgfpathlineto{\pgfqpoint{5.658290in}{0.645102in}}%
\pgfpathlineto{\pgfqpoint{5.661921in}{0.690437in}}%
\pgfpathlineto{\pgfqpoint{5.665552in}{0.654730in}}%
\pgfpathlineto{\pgfqpoint{5.669184in}{0.661382in}}%
\pgfpathlineto{\pgfqpoint{5.672815in}{0.676277in}}%
\pgfpathlineto{\pgfqpoint{5.676446in}{0.672908in}}%
\pgfpathlineto{\pgfqpoint{5.680077in}{0.658563in}}%
\pgfpathlineto{\pgfqpoint{5.683708in}{0.651194in}}%
\pgfpathlineto{\pgfqpoint{5.687340in}{0.690688in}}%
\pgfpathlineto{\pgfqpoint{5.690971in}{0.633712in}}%
\pgfpathlineto{\pgfqpoint{5.694602in}{0.673400in}}%
\pgfpathlineto{\pgfqpoint{5.698233in}{0.633091in}}%
\pgfpathlineto{\pgfqpoint{5.701864in}{0.688513in}}%
\pgfpathlineto{\pgfqpoint{5.705495in}{0.606391in}}%
\pgfpathlineto{\pgfqpoint{5.709127in}{0.669088in}}%
\pgfpathlineto{\pgfqpoint{5.712758in}{0.655673in}}%
\pgfpathlineto{\pgfqpoint{5.716389in}{0.681377in}}%
\pgfpathlineto{\pgfqpoint{5.720020in}{0.664656in}}%
\pgfpathlineto{\pgfqpoint{5.723651in}{0.637203in}}%
\pgfpathlineto{\pgfqpoint{5.727282in}{0.683902in}}%
\pgfpathlineto{\pgfqpoint{5.730914in}{0.656656in}}%
\pgfpathlineto{\pgfqpoint{5.734545in}{0.662177in}}%
\pgfpathlineto{\pgfqpoint{5.738176in}{0.623780in}}%
\pgfpathlineto{\pgfqpoint{5.741807in}{0.693282in}}%
\pgfpathlineto{\pgfqpoint{5.745438in}{0.593559in}}%
\pgfpathlineto{\pgfqpoint{5.749069in}{0.669582in}}%
\pgfpathlineto{\pgfqpoint{5.752701in}{0.657421in}}%
\pgfpathlineto{\pgfqpoint{5.756332in}{0.683232in}}%
\pgfpathlineto{\pgfqpoint{5.759963in}{0.620352in}}%
\pgfpathlineto{\pgfqpoint{5.763594in}{0.664597in}}%
\pgfpathlineto{\pgfqpoint{5.767225in}{0.672562in}}%
\pgfpathlineto{\pgfqpoint{5.770856in}{0.666902in}}%
\pgfpathlineto{\pgfqpoint{5.774488in}{0.673393in}}%
\pgfpathlineto{\pgfqpoint{5.778119in}{0.566352in}}%
\pgfpathlineto{\pgfqpoint{5.781750in}{0.689011in}}%
\pgfpathlineto{\pgfqpoint{5.785381in}{0.633525in}}%
\pgfpathlineto{\pgfqpoint{5.789012in}{0.665154in}}%
\pgfpathlineto{\pgfqpoint{5.792644in}{0.547541in}}%
\pgfpathlineto{\pgfqpoint{5.796275in}{0.691795in}}%
\pgfpathlineto{\pgfqpoint{5.799906in}{0.645409in}}%
\pgfpathlineto{\pgfqpoint{5.807168in}{0.674609in}}%
\pgfpathlineto{\pgfqpoint{5.810799in}{0.674695in}}%
\pgfpathlineto{\pgfqpoint{5.814431in}{0.639391in}}%
\pgfpathlineto{\pgfqpoint{5.818062in}{0.657633in}}%
\pgfpathlineto{\pgfqpoint{5.821693in}{0.684316in}}%
\pgfpathlineto{\pgfqpoint{5.825324in}{0.638340in}}%
\pgfpathlineto{\pgfqpoint{5.828955in}{0.677156in}}%
\pgfpathlineto{\pgfqpoint{5.832586in}{0.631010in}}%
\pgfpathlineto{\pgfqpoint{5.836218in}{0.689132in}}%
\pgfpathlineto{\pgfqpoint{5.839849in}{0.589713in}}%
\pgfpathlineto{\pgfqpoint{5.843480in}{0.665871in}}%
\pgfpathlineto{\pgfqpoint{5.847111in}{0.621656in}}%
\pgfpathlineto{\pgfqpoint{5.850742in}{0.686039in}}%
\pgfpathlineto{\pgfqpoint{5.854373in}{0.668424in}}%
\pgfpathlineto{\pgfqpoint{5.858005in}{0.635057in}}%
\pgfpathlineto{\pgfqpoint{5.861636in}{0.685508in}}%
\pgfpathlineto{\pgfqpoint{5.865267in}{0.662250in}}%
\pgfpathlineto{\pgfqpoint{5.868898in}{0.649414in}}%
\pgfpathlineto{\pgfqpoint{5.872529in}{0.646500in}}%
\pgfpathlineto{\pgfqpoint{5.876160in}{0.691012in}}%
\pgfpathlineto{\pgfqpoint{5.879792in}{0.547692in}}%
\pgfpathlineto{\pgfqpoint{5.883423in}{0.675325in}}%
\pgfpathlineto{\pgfqpoint{5.887054in}{0.664998in}}%
\pgfpathlineto{\pgfqpoint{5.890685in}{0.684951in}}%
\pgfpathlineto{\pgfqpoint{5.894316in}{0.562957in}}%
\pgfpathlineto{\pgfqpoint{5.897948in}{0.664509in}}%
\pgfpathlineto{\pgfqpoint{5.901579in}{0.653073in}}%
\pgfpathlineto{\pgfqpoint{5.905210in}{0.674292in}}%
\pgfpathlineto{\pgfqpoint{5.908841in}{0.681030in}}%
\pgfpathlineto{\pgfqpoint{5.912472in}{0.552856in}}%
\pgfpathlineto{\pgfqpoint{5.916103in}{0.691989in}}%
\pgfpathlineto{\pgfqpoint{5.919735in}{0.644508in}}%
\pgfpathlineto{\pgfqpoint{5.923366in}{0.654424in}}%
\pgfpathlineto{\pgfqpoint{5.926997in}{0.627285in}}%
\pgfpathlineto{\pgfqpoint{5.930628in}{0.692212in}}%
\pgfpathlineto{\pgfqpoint{5.934259in}{0.650480in}}%
\pgfpathlineto{\pgfqpoint{5.937890in}{0.664716in}}%
\pgfpathlineto{\pgfqpoint{5.941522in}{0.683830in}}%
\pgfpathlineto{\pgfqpoint{5.945153in}{0.676834in}}%
\pgfpathlineto{\pgfqpoint{5.948784in}{0.611457in}}%
\pgfpathlineto{\pgfqpoint{5.952415in}{0.661208in}}%
\pgfpathlineto{\pgfqpoint{5.956046in}{0.671063in}}%
\pgfpathlineto{\pgfqpoint{5.959677in}{0.649351in}}%
\pgfpathlineto{\pgfqpoint{5.963309in}{0.687162in}}%
\pgfpathlineto{\pgfqpoint{5.966940in}{0.655278in}}%
\pgfpathlineto{\pgfqpoint{5.970571in}{0.692592in}}%
\pgfpathlineto{\pgfqpoint{5.974202in}{0.617564in}}%
\pgfpathlineto{\pgfqpoint{5.977833in}{0.656264in}}%
\pgfpathlineto{\pgfqpoint{5.981465in}{0.579634in}}%
\pgfpathlineto{\pgfqpoint{5.985096in}{0.688640in}}%
\pgfpathlineto{\pgfqpoint{5.988727in}{0.678650in}}%
\pgfpathlineto{\pgfqpoint{5.992358in}{0.628112in}}%
\pgfpathlineto{\pgfqpoint{5.995989in}{0.695971in}}%
\pgfpathlineto{\pgfqpoint{5.999620in}{0.664410in}}%
\pgfpathlineto{\pgfqpoint{6.003252in}{0.627600in}}%
\pgfpathlineto{\pgfqpoint{6.010514in}{0.682293in}}%
\pgfpathlineto{\pgfqpoint{6.014145in}{0.512268in}}%
\pgfpathlineto{\pgfqpoint{6.017776in}{0.686440in}}%
\pgfpathlineto{\pgfqpoint{6.021407in}{0.686244in}}%
\pgfpathlineto{\pgfqpoint{6.025039in}{0.687704in}}%
\pgfpathlineto{\pgfqpoint{6.028670in}{0.561675in}}%
\pgfpathlineto{\pgfqpoint{6.032301in}{0.655778in}}%
\pgfpathlineto{\pgfqpoint{6.035932in}{0.594587in}}%
\pgfpathlineto{\pgfqpoint{6.039563in}{0.678859in}}%
\pgfpathlineto{\pgfqpoint{6.043194in}{0.694422in}}%
\pgfpathlineto{\pgfqpoint{6.046826in}{0.636829in}}%
\pgfpathlineto{\pgfqpoint{6.050457in}{0.702163in}}%
\pgfpathlineto{\pgfqpoint{6.054088in}{0.646368in}}%
\pgfpathlineto{\pgfqpoint{6.057719in}{0.634635in}}%
\pgfpathlineto{\pgfqpoint{6.061350in}{0.647073in}}%
\pgfpathlineto{\pgfqpoint{6.064981in}{0.686776in}}%
\pgfpathlineto{\pgfqpoint{6.068613in}{0.660395in}}%
\pgfpathlineto{\pgfqpoint{6.072244in}{0.673381in}}%
\pgfpathlineto{\pgfqpoint{6.075875in}{0.705131in}}%
\pgfpathlineto{\pgfqpoint{6.079506in}{0.677252in}}%
\pgfpathlineto{\pgfqpoint{6.083137in}{0.563715in}}%
\pgfpathlineto{\pgfqpoint{6.086769in}{0.653176in}}%
\pgfpathlineto{\pgfqpoint{6.090400in}{0.639628in}}%
\pgfpathlineto{\pgfqpoint{6.094031in}{0.655949in}}%
\pgfpathlineto{\pgfqpoint{6.097662in}{0.654495in}}%
\pgfpathlineto{\pgfqpoint{6.101293in}{0.523409in}}%
\pgfpathlineto{\pgfqpoint{6.104924in}{0.663396in}}%
\pgfpathlineto{\pgfqpoint{6.108556in}{0.624932in}}%
\pgfpathlineto{\pgfqpoint{6.112187in}{0.689040in}}%
\pgfpathlineto{\pgfqpoint{6.115818in}{0.678014in}}%
\pgfpathlineto{\pgfqpoint{6.119449in}{0.687530in}}%
\pgfpathlineto{\pgfqpoint{6.123080in}{0.656962in}}%
\pgfpathlineto{\pgfqpoint{6.126711in}{0.591986in}}%
\pgfpathlineto{\pgfqpoint{6.130343in}{0.638572in}}%
\pgfpathlineto{\pgfqpoint{6.133974in}{0.663457in}}%
\pgfpathlineto{\pgfqpoint{6.137605in}{0.664977in}}%
\pgfpathlineto{\pgfqpoint{6.141236in}{0.657899in}}%
\pgfpathlineto{\pgfqpoint{6.144867in}{0.711620in}}%
\pgfpathlineto{\pgfqpoint{6.148498in}{0.640905in}}%
\pgfpathlineto{\pgfqpoint{6.152130in}{0.629800in}}%
\pgfpathlineto{\pgfqpoint{6.155761in}{0.613497in}}%
\pgfpathlineto{\pgfqpoint{6.159392in}{0.664260in}}%
\pgfpathlineto{\pgfqpoint{6.163023in}{0.537249in}}%
\pgfpathlineto{\pgfqpoint{6.166654in}{0.691896in}}%
\pgfpathlineto{\pgfqpoint{6.170285in}{0.703059in}}%
\pgfpathlineto{\pgfqpoint{6.173917in}{0.668631in}}%
\pgfpathlineto{\pgfqpoint{6.177548in}{0.592203in}}%
\pgfpathlineto{\pgfqpoint{6.181179in}{0.652036in}}%
\pgfpathlineto{\pgfqpoint{6.184810in}{0.657455in}}%
\pgfpathlineto{\pgfqpoint{6.188441in}{0.618168in}}%
\pgfpathlineto{\pgfqpoint{6.192073in}{0.649612in}}%
\pgfpathlineto{\pgfqpoint{6.195704in}{0.611496in}}%
\pgfpathlineto{\pgfqpoint{6.199335in}{0.663219in}}%
\pgfpathlineto{\pgfqpoint{6.202966in}{0.598524in}}%
\pgfpathlineto{\pgfqpoint{6.206597in}{0.690772in}}%
\pgfpathlineto{\pgfqpoint{6.210228in}{0.694527in}}%
\pgfpathlineto{\pgfqpoint{6.213860in}{0.674183in}}%
\pgfpathlineto{\pgfqpoint{6.217491in}{0.586828in}}%
\pgfpathlineto{\pgfqpoint{6.221122in}{0.655519in}}%
\pgfpathlineto{\pgfqpoint{6.224753in}{0.655776in}}%
\pgfpathlineto{\pgfqpoint{6.228384in}{0.641781in}}%
\pgfpathlineto{\pgfqpoint{6.232015in}{0.653187in}}%
\pgfpathlineto{\pgfqpoint{6.235647in}{0.584023in}}%
\pgfpathlineto{\pgfqpoint{6.239278in}{0.663147in}}%
\pgfpathlineto{\pgfqpoint{6.242909in}{0.617142in}}%
\pgfpathlineto{\pgfqpoint{6.246540in}{0.689696in}}%
\pgfpathlineto{\pgfqpoint{6.250171in}{0.686794in}}%
\pgfpathlineto{\pgfqpoint{6.253802in}{0.679043in}}%
\pgfpathlineto{\pgfqpoint{6.257434in}{0.567116in}}%
\pgfpathlineto{\pgfqpoint{6.261065in}{0.658127in}}%
\pgfpathlineto{\pgfqpoint{6.264696in}{0.649313in}}%
\pgfpathlineto{\pgfqpoint{6.268327in}{0.658372in}}%
\pgfpathlineto{\pgfqpoint{6.271958in}{0.654433in}}%
\pgfpathlineto{\pgfqpoint{6.275590in}{0.423900in}}%
\pgfpathlineto{\pgfqpoint{6.279221in}{0.661821in}}%
\pgfpathlineto{\pgfqpoint{6.282852in}{0.632080in}}%
\pgfpathlineto{\pgfqpoint{6.286483in}{0.688391in}}%
\pgfpathlineto{\pgfqpoint{6.290114in}{0.671652in}}%
\pgfpathlineto{\pgfqpoint{6.293745in}{0.691133in}}%
\pgfpathlineto{\pgfqpoint{6.297377in}{0.655788in}}%
\pgfpathlineto{\pgfqpoint{6.301008in}{0.600198in}}%
\pgfpathlineto{\pgfqpoint{6.304639in}{0.635533in}}%
\pgfpathlineto{\pgfqpoint{6.308270in}{0.664874in}}%
\pgfpathlineto{\pgfqpoint{6.311901in}{0.661337in}}%
\pgfpathlineto{\pgfqpoint{6.315532in}{0.664499in}}%
\pgfpathlineto{\pgfqpoint{6.319164in}{0.712273in}}%
\pgfpathlineto{\pgfqpoint{6.322795in}{0.632376in}}%
\pgfpathlineto{\pgfqpoint{6.326426in}{0.633617in}}%
\pgfpathlineto{\pgfqpoint{6.330057in}{0.608418in}}%
\pgfpathlineto{\pgfqpoint{6.333688in}{0.663493in}}%
\pgfpathlineto{\pgfqpoint{6.337319in}{0.576648in}}%
\pgfpathlineto{\pgfqpoint{6.340951in}{0.692058in}}%
\pgfpathlineto{\pgfqpoint{6.344582in}{0.700144in}}%
\pgfpathlineto{\pgfqpoint{6.348213in}{0.671514in}}%
\pgfpathlineto{\pgfqpoint{6.351844in}{0.586372in}}%
\pgfpathlineto{\pgfqpoint{6.355475in}{0.653574in}}%
\pgfpathlineto{\pgfqpoint{6.359106in}{0.655096in}}%
\pgfpathlineto{\pgfqpoint{6.362738in}{0.631740in}}%
\pgfpathlineto{\pgfqpoint{6.366369in}{0.651046in}}%
\pgfpathlineto{\pgfqpoint{6.370000in}{0.653739in}}%
\pgfpathlineto{\pgfqpoint{6.370000in}{0.653739in}}%
\pgfusepath{stroke}%
\end{pgfscope}%
\begin{pgfscope}%
\pgfpathrectangle{\pgfqpoint{4.554414in}{0.320625in}}{\pgfqpoint{1.815586in}{1.164375in}}%
\pgfusepath{clip}%
\pgfsetrectcap%
\pgfsetroundjoin%
\pgfsetlinewidth{0.602250pt}%
\definecolor{currentstroke}{rgb}{0.141176,0.345098,0.650980}%
\pgfsetstrokecolor{currentstroke}%
\pgfsetstrokeopacity{0.120000}%
\pgfsetdash{}{0pt}%
\pgfpathmoveto{\pgfqpoint{4.554414in}{1.476862in}}%
\pgfpathlineto{\pgfqpoint{4.561676in}{1.456008in}}%
\pgfpathlineto{\pgfqpoint{4.568938in}{1.439380in}}%
\pgfpathlineto{\pgfqpoint{4.576201in}{1.426375in}}%
\pgfpathlineto{\pgfqpoint{4.587094in}{1.410255in}}%
\pgfpathlineto{\pgfqpoint{4.601619in}{1.392569in}}%
\pgfpathlineto{\pgfqpoint{4.619775in}{1.372239in}}%
\pgfpathlineto{\pgfqpoint{4.637931in}{1.353783in}}%
\pgfpathlineto{\pgfqpoint{4.656086in}{1.336910in}}%
\pgfpathlineto{\pgfqpoint{4.674242in}{1.321498in}}%
\pgfpathlineto{\pgfqpoint{4.696029in}{1.304534in}}%
\pgfpathlineto{\pgfqpoint{4.721448in}{1.286301in}}%
\pgfpathlineto{\pgfqpoint{4.754128in}{1.264474in}}%
\pgfpathlineto{\pgfqpoint{4.797702in}{1.236869in}}%
\pgfpathlineto{\pgfqpoint{4.866694in}{1.194646in}}%
\pgfpathlineto{\pgfqpoint{4.990154in}{1.120612in}}%
\pgfpathlineto{\pgfqpoint{5.244336in}{0.969730in}}%
\pgfpathlineto{\pgfqpoint{5.549355in}{0.789302in}}%
\pgfpathlineto{\pgfqpoint{5.552986in}{0.788062in}}%
\pgfpathlineto{\pgfqpoint{5.556617in}{0.784927in}}%
\pgfpathlineto{\pgfqpoint{5.574773in}{0.774844in}}%
\pgfpathlineto{\pgfqpoint{5.578404in}{0.774885in}}%
\pgfpathlineto{\pgfqpoint{5.582035in}{0.770235in}}%
\pgfpathlineto{\pgfqpoint{5.592929in}{0.763576in}}%
\pgfpathlineto{\pgfqpoint{5.596560in}{0.761655in}}%
\pgfpathlineto{\pgfqpoint{5.600191in}{0.761800in}}%
\pgfpathlineto{\pgfqpoint{5.603823in}{0.757932in}}%
\pgfpathlineto{\pgfqpoint{5.607454in}{0.759204in}}%
\pgfpathlineto{\pgfqpoint{5.611085in}{0.752707in}}%
\pgfpathlineto{\pgfqpoint{5.618347in}{0.748461in}}%
\pgfpathlineto{\pgfqpoint{5.621978in}{0.747920in}}%
\pgfpathlineto{\pgfqpoint{5.629241in}{0.742108in}}%
\pgfpathlineto{\pgfqpoint{5.632872in}{0.741229in}}%
\pgfpathlineto{\pgfqpoint{5.636503in}{0.739265in}}%
\pgfpathlineto{\pgfqpoint{5.640134in}{0.735858in}}%
\pgfpathlineto{\pgfqpoint{5.643765in}{0.736724in}}%
\pgfpathlineto{\pgfqpoint{5.647397in}{0.742169in}}%
\pgfpathlineto{\pgfqpoint{5.651028in}{0.729846in}}%
\pgfpathlineto{\pgfqpoint{5.658290in}{0.725561in}}%
\pgfpathlineto{\pgfqpoint{5.661921in}{0.727748in}}%
\pgfpathlineto{\pgfqpoint{5.665552in}{0.720710in}}%
\pgfpathlineto{\pgfqpoint{5.669184in}{0.720118in}}%
\pgfpathlineto{\pgfqpoint{5.672815in}{0.717633in}}%
\pgfpathlineto{\pgfqpoint{5.676446in}{0.719848in}}%
\pgfpathlineto{\pgfqpoint{5.680077in}{0.711940in}}%
\pgfpathlineto{\pgfqpoint{5.683708in}{0.721599in}}%
\pgfpathlineto{\pgfqpoint{5.687340in}{0.725828in}}%
\pgfpathlineto{\pgfqpoint{5.690971in}{0.709740in}}%
\pgfpathlineto{\pgfqpoint{5.694602in}{0.710802in}}%
\pgfpathlineto{\pgfqpoint{5.698233in}{0.701343in}}%
\pgfpathlineto{\pgfqpoint{5.701864in}{0.705380in}}%
\pgfpathlineto{\pgfqpoint{5.705495in}{0.702582in}}%
\pgfpathlineto{\pgfqpoint{5.709127in}{0.707253in}}%
\pgfpathlineto{\pgfqpoint{5.712758in}{0.694655in}}%
\pgfpathlineto{\pgfqpoint{5.716389in}{0.729406in}}%
\pgfpathlineto{\pgfqpoint{5.720020in}{0.699540in}}%
\pgfpathlineto{\pgfqpoint{5.723651in}{0.688122in}}%
\pgfpathlineto{\pgfqpoint{5.727282in}{0.688293in}}%
\pgfpathlineto{\pgfqpoint{5.730914in}{0.698005in}}%
\pgfpathlineto{\pgfqpoint{5.734545in}{0.679762in}}%
\pgfpathlineto{\pgfqpoint{5.738176in}{0.713344in}}%
\pgfpathlineto{\pgfqpoint{5.741807in}{0.724668in}}%
\pgfpathlineto{\pgfqpoint{5.745438in}{0.694951in}}%
\pgfpathlineto{\pgfqpoint{5.749069in}{0.672633in}}%
\pgfpathlineto{\pgfqpoint{5.752701in}{0.684752in}}%
\pgfpathlineto{\pgfqpoint{5.756332in}{0.686833in}}%
\pgfpathlineto{\pgfqpoint{5.759963in}{0.672315in}}%
\pgfpathlineto{\pgfqpoint{5.763594in}{0.681211in}}%
\pgfpathlineto{\pgfqpoint{5.767225in}{0.662594in}}%
\pgfpathlineto{\pgfqpoint{5.770856in}{0.686393in}}%
\pgfpathlineto{\pgfqpoint{5.774488in}{0.663762in}}%
\pgfpathlineto{\pgfqpoint{5.778119in}{0.711577in}}%
\pgfpathlineto{\pgfqpoint{5.781750in}{0.707738in}}%
\pgfpathlineto{\pgfqpoint{5.785381in}{0.701444in}}%
\pgfpathlineto{\pgfqpoint{5.789012in}{0.647461in}}%
\pgfpathlineto{\pgfqpoint{5.792644in}{0.681942in}}%
\pgfpathlineto{\pgfqpoint{5.796275in}{0.661446in}}%
\pgfpathlineto{\pgfqpoint{5.799906in}{0.691491in}}%
\pgfpathlineto{\pgfqpoint{5.803537in}{0.673870in}}%
\pgfpathlineto{\pgfqpoint{5.807168in}{0.640799in}}%
\pgfpathlineto{\pgfqpoint{5.810799in}{0.672084in}}%
\pgfpathlineto{\pgfqpoint{5.814431in}{0.673592in}}%
\pgfpathlineto{\pgfqpoint{5.818062in}{0.707802in}}%
\pgfpathlineto{\pgfqpoint{5.821693in}{0.682824in}}%
\pgfpathlineto{\pgfqpoint{5.825324in}{0.706055in}}%
\pgfpathlineto{\pgfqpoint{5.828955in}{0.643394in}}%
\pgfpathlineto{\pgfqpoint{5.832586in}{0.680465in}}%
\pgfpathlineto{\pgfqpoint{5.836218in}{0.621761in}}%
\pgfpathlineto{\pgfqpoint{5.839849in}{0.708941in}}%
\pgfpathlineto{\pgfqpoint{5.843480in}{0.656613in}}%
\pgfpathlineto{\pgfqpoint{5.850742in}{0.646259in}}%
\pgfpathlineto{\pgfqpoint{5.854373in}{0.688263in}}%
\pgfpathlineto{\pgfqpoint{5.858005in}{0.699408in}}%
\pgfpathlineto{\pgfqpoint{5.861636in}{0.643087in}}%
\pgfpathlineto{\pgfqpoint{5.865267in}{0.730807in}}%
\pgfpathlineto{\pgfqpoint{5.868898in}{0.616598in}}%
\pgfpathlineto{\pgfqpoint{5.872529in}{0.654859in}}%
\pgfpathlineto{\pgfqpoint{5.876160in}{0.622254in}}%
\pgfpathlineto{\pgfqpoint{5.879792in}{0.685971in}}%
\pgfpathlineto{\pgfqpoint{5.883423in}{0.594195in}}%
\pgfpathlineto{\pgfqpoint{5.887054in}{0.714917in}}%
\pgfpathlineto{\pgfqpoint{5.890685in}{0.723302in}}%
\pgfpathlineto{\pgfqpoint{5.894316in}{0.701071in}}%
\pgfpathlineto{\pgfqpoint{5.897948in}{0.640653in}}%
\pgfpathlineto{\pgfqpoint{5.901579in}{0.644866in}}%
\pgfpathlineto{\pgfqpoint{5.905210in}{0.620136in}}%
\pgfpathlineto{\pgfqpoint{5.908841in}{0.687990in}}%
\pgfpathlineto{\pgfqpoint{5.912472in}{0.705916in}}%
\pgfpathlineto{\pgfqpoint{5.916103in}{0.659355in}}%
\pgfpathlineto{\pgfqpoint{5.919735in}{0.701864in}}%
\pgfpathlineto{\pgfqpoint{5.923366in}{0.665761in}}%
\pgfpathlineto{\pgfqpoint{5.926997in}{0.668519in}}%
\pgfpathlineto{\pgfqpoint{5.930628in}{0.675539in}}%
\pgfpathlineto{\pgfqpoint{5.934259in}{0.720017in}}%
\pgfpathlineto{\pgfqpoint{5.937890in}{0.588559in}}%
\pgfpathlineto{\pgfqpoint{5.941522in}{0.652018in}}%
\pgfpathlineto{\pgfqpoint{5.945153in}{0.592923in}}%
\pgfpathlineto{\pgfqpoint{5.948784in}{0.694865in}}%
\pgfpathlineto{\pgfqpoint{5.952415in}{0.688611in}}%
\pgfpathlineto{\pgfqpoint{5.956046in}{0.599547in}}%
\pgfpathlineto{\pgfqpoint{5.959677in}{0.690233in}}%
\pgfpathlineto{\pgfqpoint{5.963309in}{0.680514in}}%
\pgfpathlineto{\pgfqpoint{5.966940in}{0.656417in}}%
\pgfpathlineto{\pgfqpoint{5.974202in}{0.723167in}}%
\pgfpathlineto{\pgfqpoint{5.977833in}{0.643116in}}%
\pgfpathlineto{\pgfqpoint{5.981465in}{0.654405in}}%
\pgfpathlineto{\pgfqpoint{5.985096in}{0.645902in}}%
\pgfpathlineto{\pgfqpoint{5.988727in}{0.697042in}}%
\pgfpathlineto{\pgfqpoint{5.992358in}{0.653780in}}%
\pgfpathlineto{\pgfqpoint{5.995989in}{0.660733in}}%
\pgfpathlineto{\pgfqpoint{5.999620in}{0.670836in}}%
\pgfpathlineto{\pgfqpoint{6.003252in}{0.688953in}}%
\pgfpathlineto{\pgfqpoint{6.006883in}{0.621096in}}%
\pgfpathlineto{\pgfqpoint{6.010514in}{0.701379in}}%
\pgfpathlineto{\pgfqpoint{6.014145in}{0.718225in}}%
\pgfpathlineto{\pgfqpoint{6.017776in}{0.649126in}}%
\pgfpathlineto{\pgfqpoint{6.021407in}{0.685933in}}%
\pgfpathlineto{\pgfqpoint{6.025039in}{0.633752in}}%
\pgfpathlineto{\pgfqpoint{6.028670in}{0.682454in}}%
\pgfpathlineto{\pgfqpoint{6.032301in}{0.665432in}}%
\pgfpathlineto{\pgfqpoint{6.035932in}{0.692858in}}%
\pgfpathlineto{\pgfqpoint{6.039563in}{0.593520in}}%
\pgfpathlineto{\pgfqpoint{6.043194in}{0.721790in}}%
\pgfpathlineto{\pgfqpoint{6.046826in}{0.689897in}}%
\pgfpathlineto{\pgfqpoint{6.050457in}{0.617681in}}%
\pgfpathlineto{\pgfqpoint{6.054088in}{0.664020in}}%
\pgfpathlineto{\pgfqpoint{6.057719in}{0.684758in}}%
\pgfpathlineto{\pgfqpoint{6.061350in}{0.614523in}}%
\pgfpathlineto{\pgfqpoint{6.064981in}{0.706653in}}%
\pgfpathlineto{\pgfqpoint{6.068613in}{0.725239in}}%
\pgfpathlineto{\pgfqpoint{6.075875in}{0.629195in}}%
\pgfpathlineto{\pgfqpoint{6.079506in}{0.667642in}}%
\pgfpathlineto{\pgfqpoint{6.083137in}{0.683496in}}%
\pgfpathlineto{\pgfqpoint{6.086769in}{0.602606in}}%
\pgfpathlineto{\pgfqpoint{6.090400in}{0.667875in}}%
\pgfpathlineto{\pgfqpoint{6.094031in}{0.637325in}}%
\pgfpathlineto{\pgfqpoint{6.097662in}{0.683419in}}%
\pgfpathlineto{\pgfqpoint{6.101293in}{0.615569in}}%
\pgfpathlineto{\pgfqpoint{6.104924in}{0.709592in}}%
\pgfpathlineto{\pgfqpoint{6.108556in}{0.711731in}}%
\pgfpathlineto{\pgfqpoint{6.112187in}{0.695449in}}%
\pgfpathlineto{\pgfqpoint{6.115818in}{0.589126in}}%
\pgfpathlineto{\pgfqpoint{6.119449in}{0.676323in}}%
\pgfpathlineto{\pgfqpoint{6.123080in}{0.665455in}}%
\pgfpathlineto{\pgfqpoint{6.126711in}{0.678167in}}%
\pgfpathlineto{\pgfqpoint{6.130343in}{0.671715in}}%
\pgfpathlineto{\pgfqpoint{6.133974in}{0.542858in}}%
\pgfpathlineto{\pgfqpoint{6.137605in}{0.674915in}}%
\pgfpathlineto{\pgfqpoint{6.141236in}{0.660794in}}%
\pgfpathlineto{\pgfqpoint{6.144867in}{0.708343in}}%
\pgfpathlineto{\pgfqpoint{6.148498in}{0.691019in}}%
\pgfpathlineto{\pgfqpoint{6.152130in}{0.704389in}}%
\pgfpathlineto{\pgfqpoint{6.155761in}{0.609491in}}%
\pgfpathlineto{\pgfqpoint{6.159392in}{0.679568in}}%
\pgfpathlineto{\pgfqpoint{6.163023in}{0.602393in}}%
\pgfpathlineto{\pgfqpoint{6.166654in}{0.702985in}}%
\pgfpathlineto{\pgfqpoint{6.170285in}{0.662016in}}%
\pgfpathlineto{\pgfqpoint{6.173917in}{0.639058in}}%
\pgfpathlineto{\pgfqpoint{6.177548in}{0.655755in}}%
\pgfpathlineto{\pgfqpoint{6.181179in}{0.682538in}}%
\pgfpathlineto{\pgfqpoint{6.184810in}{0.702415in}}%
\pgfpathlineto{\pgfqpoint{6.188441in}{0.541841in}}%
\pgfpathlineto{\pgfqpoint{6.192073in}{0.728659in}}%
\pgfpathlineto{\pgfqpoint{6.195704in}{0.578657in}}%
\pgfpathlineto{\pgfqpoint{6.199335in}{0.652706in}}%
\pgfpathlineto{\pgfqpoint{6.202966in}{0.577285in}}%
\pgfpathlineto{\pgfqpoint{6.206597in}{0.688203in}}%
\pgfpathlineto{\pgfqpoint{6.210228in}{0.605555in}}%
\pgfpathlineto{\pgfqpoint{6.213860in}{0.712449in}}%
\pgfpathlineto{\pgfqpoint{6.217491in}{0.727656in}}%
\pgfpathlineto{\pgfqpoint{6.224753in}{0.648714in}}%
\pgfpathlineto{\pgfqpoint{6.228384in}{0.636205in}}%
\pgfpathlineto{\pgfqpoint{6.232015in}{0.638611in}}%
\pgfpathlineto{\pgfqpoint{6.235647in}{0.682879in}}%
\pgfpathlineto{\pgfqpoint{6.239278in}{0.708759in}}%
\pgfpathlineto{\pgfqpoint{6.242909in}{0.674383in}}%
\pgfpathlineto{\pgfqpoint{6.246540in}{0.704369in}}%
\pgfpathlineto{\pgfqpoint{6.250171in}{0.657534in}}%
\pgfpathlineto{\pgfqpoint{6.253802in}{0.671896in}}%
\pgfpathlineto{\pgfqpoint{6.257434in}{0.663244in}}%
\pgfpathlineto{\pgfqpoint{6.261065in}{0.716875in}}%
\pgfpathlineto{\pgfqpoint{6.264696in}{0.624517in}}%
\pgfpathlineto{\pgfqpoint{6.268327in}{0.649173in}}%
\pgfpathlineto{\pgfqpoint{6.271958in}{0.534547in}}%
\pgfpathlineto{\pgfqpoint{6.275590in}{0.692630in}}%
\pgfpathlineto{\pgfqpoint{6.279221in}{0.695048in}}%
\pgfpathlineto{\pgfqpoint{6.282852in}{0.565729in}}%
\pgfpathlineto{\pgfqpoint{6.286483in}{0.695340in}}%
\pgfpathlineto{\pgfqpoint{6.290114in}{0.675810in}}%
\pgfpathlineto{\pgfqpoint{6.293745in}{0.663611in}}%
\pgfpathlineto{\pgfqpoint{6.297377in}{0.685258in}}%
\pgfpathlineto{\pgfqpoint{6.301008in}{0.722426in}}%
\pgfpathlineto{\pgfqpoint{6.304639in}{0.613912in}}%
\pgfpathlineto{\pgfqpoint{6.308270in}{0.655806in}}%
\pgfpathlineto{\pgfqpoint{6.311901in}{0.630545in}}%
\pgfpathlineto{\pgfqpoint{6.315532in}{0.697357in}}%
\pgfpathlineto{\pgfqpoint{6.319164in}{0.670224in}}%
\pgfpathlineto{\pgfqpoint{6.322795in}{0.650645in}}%
\pgfpathlineto{\pgfqpoint{6.326426in}{0.679462in}}%
\pgfpathlineto{\pgfqpoint{6.330057in}{0.686983in}}%
\pgfpathlineto{\pgfqpoint{6.333688in}{0.640517in}}%
\pgfpathlineto{\pgfqpoint{6.337319in}{0.697884in}}%
\pgfpathlineto{\pgfqpoint{6.340951in}{0.720803in}}%
\pgfpathlineto{\pgfqpoint{6.344582in}{0.576403in}}%
\pgfpathlineto{\pgfqpoint{6.348213in}{0.684987in}}%
\pgfpathlineto{\pgfqpoint{6.351844in}{0.653150in}}%
\pgfpathlineto{\pgfqpoint{6.355475in}{0.686994in}}%
\pgfpathlineto{\pgfqpoint{6.359106in}{0.656317in}}%
\pgfpathlineto{\pgfqpoint{6.362738in}{0.694661in}}%
\pgfpathlineto{\pgfqpoint{6.366369in}{0.670159in}}%
\pgfpathlineto{\pgfqpoint{6.370000in}{0.677715in}}%
\pgfpathlineto{\pgfqpoint{6.370000in}{0.677715in}}%
\pgfusepath{stroke}%
\end{pgfscope}%
\begin{pgfscope}%
\pgfpathrectangle{\pgfqpoint{4.554414in}{0.320625in}}{\pgfqpoint{1.815586in}{1.164375in}}%
\pgfusepath{clip}%
\pgfsetrectcap%
\pgfsetroundjoin%
\pgfsetlinewidth{0.602250pt}%
\definecolor{currentstroke}{rgb}{0.141176,0.345098,0.650980}%
\pgfsetstrokecolor{currentstroke}%
\pgfsetstrokeopacity{0.120000}%
\pgfsetdash{}{0pt}%
\pgfpathmoveto{\pgfqpoint{4.554414in}{1.475541in}}%
\pgfpathlineto{\pgfqpoint{4.558045in}{1.463557in}}%
\pgfpathlineto{\pgfqpoint{4.561676in}{1.453688in}}%
\pgfpathlineto{\pgfqpoint{4.568938in}{1.439182in}}%
\pgfpathlineto{\pgfqpoint{4.576201in}{1.427344in}}%
\pgfpathlineto{\pgfqpoint{4.587094in}{1.413089in}}%
\pgfpathlineto{\pgfqpoint{4.597988in}{1.400334in}}%
\pgfpathlineto{\pgfqpoint{4.623406in}{1.373284in}}%
\pgfpathlineto{\pgfqpoint{4.656086in}{1.340063in}}%
\pgfpathlineto{\pgfqpoint{4.685136in}{1.312492in}}%
\pgfpathlineto{\pgfqpoint{4.717816in}{1.283132in}}%
\pgfpathlineto{\pgfqpoint{4.746866in}{1.258487in}}%
\pgfpathlineto{\pgfqpoint{4.779546in}{1.232151in}}%
\pgfpathlineto{\pgfqpoint{4.819489in}{1.201485in}}%
\pgfpathlineto{\pgfqpoint{4.863063in}{1.169460in}}%
\pgfpathlineto{\pgfqpoint{4.913900in}{1.133531in}}%
\pgfpathlineto{\pgfqpoint{4.968367in}{1.096393in}}%
\pgfpathlineto{\pgfqpoint{5.030097in}{1.055674in}}%
\pgfpathlineto{\pgfqpoint{5.099090in}{1.011543in}}%
\pgfpathlineto{\pgfqpoint{5.178975in}{0.961839in}}%
\pgfpathlineto{\pgfqpoint{5.273386in}{0.904505in}}%
\pgfpathlineto{\pgfqpoint{5.385952in}{0.837535in}}%
\pgfpathlineto{\pgfqpoint{5.494887in}{0.773760in}}%
\pgfpathlineto{\pgfqpoint{5.498519in}{0.772484in}}%
\pgfpathlineto{\pgfqpoint{5.502150in}{0.769548in}}%
\pgfpathlineto{\pgfqpoint{5.513043in}{0.763255in}}%
\pgfpathlineto{\pgfqpoint{5.516674in}{0.762986in}}%
\pgfpathlineto{\pgfqpoint{5.520306in}{0.759296in}}%
\pgfpathlineto{\pgfqpoint{5.527568in}{0.755055in}}%
\pgfpathlineto{\pgfqpoint{5.531199in}{0.753206in}}%
\pgfpathlineto{\pgfqpoint{5.534830in}{0.753044in}}%
\pgfpathlineto{\pgfqpoint{5.542093in}{0.746473in}}%
\pgfpathlineto{\pgfqpoint{5.545724in}{0.744503in}}%
\pgfpathlineto{\pgfqpoint{5.549355in}{0.743882in}}%
\pgfpathlineto{\pgfqpoint{5.556617in}{0.740386in}}%
\pgfpathlineto{\pgfqpoint{5.560248in}{0.736596in}}%
\pgfpathlineto{\pgfqpoint{5.563880in}{0.733803in}}%
\pgfpathlineto{\pgfqpoint{5.567511in}{0.735079in}}%
\pgfpathlineto{\pgfqpoint{5.571142in}{0.729625in}}%
\pgfpathlineto{\pgfqpoint{5.574773in}{0.738855in}}%
\pgfpathlineto{\pgfqpoint{5.578404in}{0.728512in}}%
\pgfpathlineto{\pgfqpoint{5.582035in}{0.723379in}}%
\pgfpathlineto{\pgfqpoint{5.585667in}{0.721647in}}%
\pgfpathlineto{\pgfqpoint{5.589298in}{0.722106in}}%
\pgfpathlineto{\pgfqpoint{5.592929in}{0.718232in}}%
\pgfpathlineto{\pgfqpoint{5.596560in}{0.728330in}}%
\pgfpathlineto{\pgfqpoint{5.600191in}{0.712804in}}%
\pgfpathlineto{\pgfqpoint{5.607454in}{0.709395in}}%
\pgfpathlineto{\pgfqpoint{5.611085in}{0.707711in}}%
\pgfpathlineto{\pgfqpoint{5.614716in}{0.707986in}}%
\pgfpathlineto{\pgfqpoint{5.618347in}{0.711102in}}%
\pgfpathlineto{\pgfqpoint{5.621978in}{0.700426in}}%
\pgfpathlineto{\pgfqpoint{5.625610in}{0.699112in}}%
\pgfpathlineto{\pgfqpoint{5.629241in}{0.699603in}}%
\pgfpathlineto{\pgfqpoint{5.632872in}{0.694057in}}%
\pgfpathlineto{\pgfqpoint{5.636503in}{0.714762in}}%
\pgfpathlineto{\pgfqpoint{5.640134in}{0.691311in}}%
\pgfpathlineto{\pgfqpoint{5.643765in}{0.688287in}}%
\pgfpathlineto{\pgfqpoint{5.647397in}{0.692765in}}%
\pgfpathlineto{\pgfqpoint{5.651028in}{0.687205in}}%
\pgfpathlineto{\pgfqpoint{5.654659in}{0.716429in}}%
\pgfpathlineto{\pgfqpoint{5.658290in}{0.691341in}}%
\pgfpathlineto{\pgfqpoint{5.661921in}{0.677457in}}%
\pgfpathlineto{\pgfqpoint{5.669184in}{0.691105in}}%
\pgfpathlineto{\pgfqpoint{5.672815in}{0.709589in}}%
\pgfpathlineto{\pgfqpoint{5.676446in}{0.696695in}}%
\pgfpathlineto{\pgfqpoint{5.680077in}{0.672787in}}%
\pgfpathlineto{\pgfqpoint{5.683708in}{0.672583in}}%
\pgfpathlineto{\pgfqpoint{5.687340in}{0.696724in}}%
\pgfpathlineto{\pgfqpoint{5.690971in}{0.692388in}}%
\pgfpathlineto{\pgfqpoint{5.694602in}{0.696978in}}%
\pgfpathlineto{\pgfqpoint{5.701864in}{0.655636in}}%
\pgfpathlineto{\pgfqpoint{5.705495in}{0.698210in}}%
\pgfpathlineto{\pgfqpoint{5.709127in}{0.651102in}}%
\pgfpathlineto{\pgfqpoint{5.712758in}{0.721679in}}%
\pgfpathlineto{\pgfqpoint{5.716389in}{0.691873in}}%
\pgfpathlineto{\pgfqpoint{5.720020in}{0.650506in}}%
\pgfpathlineto{\pgfqpoint{5.723651in}{0.654701in}}%
\pgfpathlineto{\pgfqpoint{5.727282in}{0.683725in}}%
\pgfpathlineto{\pgfqpoint{5.730914in}{0.660079in}}%
\pgfpathlineto{\pgfqpoint{5.734545in}{0.714687in}}%
\pgfpathlineto{\pgfqpoint{5.738176in}{0.633936in}}%
\pgfpathlineto{\pgfqpoint{5.741807in}{0.647850in}}%
\pgfpathlineto{\pgfqpoint{5.745438in}{0.649387in}}%
\pgfpathlineto{\pgfqpoint{5.749069in}{0.656030in}}%
\pgfpathlineto{\pgfqpoint{5.752701in}{0.667663in}}%
\pgfpathlineto{\pgfqpoint{5.756332in}{0.692038in}}%
\pgfpathlineto{\pgfqpoint{5.759963in}{0.624278in}}%
\pgfpathlineto{\pgfqpoint{5.763594in}{0.640187in}}%
\pgfpathlineto{\pgfqpoint{5.767225in}{0.663551in}}%
\pgfpathlineto{\pgfqpoint{5.770856in}{0.615126in}}%
\pgfpathlineto{\pgfqpoint{5.774488in}{0.707620in}}%
\pgfpathlineto{\pgfqpoint{5.778119in}{0.645438in}}%
\pgfpathlineto{\pgfqpoint{5.781750in}{0.620891in}}%
\pgfpathlineto{\pgfqpoint{5.785381in}{0.667397in}}%
\pgfpathlineto{\pgfqpoint{5.789012in}{0.656600in}}%
\pgfpathlineto{\pgfqpoint{5.792644in}{0.712470in}}%
\pgfpathlineto{\pgfqpoint{5.796275in}{0.676501in}}%
\pgfpathlineto{\pgfqpoint{5.799906in}{0.608576in}}%
\pgfpathlineto{\pgfqpoint{5.803537in}{0.662468in}}%
\pgfpathlineto{\pgfqpoint{5.807168in}{0.682216in}}%
\pgfpathlineto{\pgfqpoint{5.810799in}{0.705214in}}%
\pgfpathlineto{\pgfqpoint{5.814431in}{0.691061in}}%
\pgfpathlineto{\pgfqpoint{5.818062in}{0.647753in}}%
\pgfpathlineto{\pgfqpoint{5.821693in}{0.645710in}}%
\pgfpathlineto{\pgfqpoint{5.825324in}{0.693326in}}%
\pgfpathlineto{\pgfqpoint{5.828955in}{0.685552in}}%
\pgfpathlineto{\pgfqpoint{5.832586in}{0.692776in}}%
\pgfpathlineto{\pgfqpoint{5.836218in}{0.669318in}}%
\pgfpathlineto{\pgfqpoint{5.839849in}{0.591209in}}%
\pgfpathlineto{\pgfqpoint{5.843480in}{0.696068in}}%
\pgfpathlineto{\pgfqpoint{5.847111in}{0.644743in}}%
\pgfpathlineto{\pgfqpoint{5.850742in}{0.679717in}}%
\pgfpathlineto{\pgfqpoint{5.854373in}{0.678142in}}%
\pgfpathlineto{\pgfqpoint{5.858005in}{0.634869in}}%
\pgfpathlineto{\pgfqpoint{5.861636in}{0.691095in}}%
\pgfpathlineto{\pgfqpoint{5.865267in}{0.672554in}}%
\pgfpathlineto{\pgfqpoint{5.868898in}{0.710095in}}%
\pgfpathlineto{\pgfqpoint{5.872529in}{0.703697in}}%
\pgfpathlineto{\pgfqpoint{5.876160in}{0.664041in}}%
\pgfpathlineto{\pgfqpoint{5.879792in}{0.578078in}}%
\pgfpathlineto{\pgfqpoint{5.883423in}{0.683855in}}%
\pgfpathlineto{\pgfqpoint{5.887054in}{0.640848in}}%
\pgfpathlineto{\pgfqpoint{5.890685in}{0.721038in}}%
\pgfpathlineto{\pgfqpoint{5.894316in}{0.682840in}}%
\pgfpathlineto{\pgfqpoint{5.897948in}{0.547890in}}%
\pgfpathlineto{\pgfqpoint{5.901579in}{0.632469in}}%
\pgfpathlineto{\pgfqpoint{5.905210in}{0.665797in}}%
\pgfpathlineto{\pgfqpoint{5.908841in}{0.551136in}}%
\pgfpathlineto{\pgfqpoint{5.912472in}{0.709466in}}%
\pgfpathlineto{\pgfqpoint{5.919735in}{0.588949in}}%
\pgfpathlineto{\pgfqpoint{5.923366in}{0.664627in}}%
\pgfpathlineto{\pgfqpoint{5.926997in}{0.663163in}}%
\pgfpathlineto{\pgfqpoint{5.930628in}{0.712459in}}%
\pgfpathlineto{\pgfqpoint{5.934259in}{0.678543in}}%
\pgfpathlineto{\pgfqpoint{5.937890in}{0.603531in}}%
\pgfpathlineto{\pgfqpoint{5.941522in}{0.659776in}}%
\pgfpathlineto{\pgfqpoint{5.945153in}{0.684756in}}%
\pgfpathlineto{\pgfqpoint{5.948784in}{0.704186in}}%
\pgfpathlineto{\pgfqpoint{5.952415in}{0.691718in}}%
\pgfpathlineto{\pgfqpoint{5.956046in}{0.650412in}}%
\pgfpathlineto{\pgfqpoint{5.959677in}{0.639474in}}%
\pgfpathlineto{\pgfqpoint{5.963309in}{0.694602in}}%
\pgfpathlineto{\pgfqpoint{5.966940in}{0.683143in}}%
\pgfpathlineto{\pgfqpoint{5.970571in}{0.691669in}}%
\pgfpathlineto{\pgfqpoint{5.974202in}{0.670424in}}%
\pgfpathlineto{\pgfqpoint{5.977833in}{0.537524in}}%
\pgfpathlineto{\pgfqpoint{5.981465in}{0.696360in}}%
\pgfpathlineto{\pgfqpoint{5.985096in}{0.638276in}}%
\pgfpathlineto{\pgfqpoint{5.988727in}{0.677209in}}%
\pgfpathlineto{\pgfqpoint{5.992358in}{0.677669in}}%
\pgfpathlineto{\pgfqpoint{5.995989in}{0.640568in}}%
\pgfpathlineto{\pgfqpoint{5.999620in}{0.690374in}}%
\pgfpathlineto{\pgfqpoint{6.003252in}{0.676642in}}%
\pgfpathlineto{\pgfqpoint{6.006883in}{0.709300in}}%
\pgfpathlineto{\pgfqpoint{6.010514in}{0.704391in}}%
\pgfpathlineto{\pgfqpoint{6.014145in}{0.666274in}}%
\pgfpathlineto{\pgfqpoint{6.017776in}{0.531384in}}%
\pgfpathlineto{\pgfqpoint{6.021407in}{0.683645in}}%
\pgfpathlineto{\pgfqpoint{6.025039in}{0.650182in}}%
\pgfpathlineto{\pgfqpoint{6.028670in}{0.721622in}}%
\pgfpathlineto{\pgfqpoint{6.032301in}{0.686204in}}%
\pgfpathlineto{\pgfqpoint{6.035932in}{0.568285in}}%
\pgfpathlineto{\pgfqpoint{6.039563in}{0.628584in}}%
\pgfpathlineto{\pgfqpoint{6.043194in}{0.666401in}}%
\pgfpathlineto{\pgfqpoint{6.046826in}{0.559155in}}%
\pgfpathlineto{\pgfqpoint{6.050457in}{0.710722in}}%
\pgfpathlineto{\pgfqpoint{6.054088in}{0.651783in}}%
\pgfpathlineto{\pgfqpoint{6.057719in}{0.571606in}}%
\pgfpathlineto{\pgfqpoint{6.061350in}{0.663937in}}%
\pgfpathlineto{\pgfqpoint{6.064981in}{0.667715in}}%
\pgfpathlineto{\pgfqpoint{6.068613in}{0.712644in}}%
\pgfpathlineto{\pgfqpoint{6.072244in}{0.680421in}}%
\pgfpathlineto{\pgfqpoint{6.075875in}{0.612263in}}%
\pgfpathlineto{\pgfqpoint{6.079506in}{0.657897in}}%
\pgfpathlineto{\pgfqpoint{6.083137in}{0.686907in}}%
\pgfpathlineto{\pgfqpoint{6.086769in}{0.703271in}}%
\pgfpathlineto{\pgfqpoint{6.090400in}{0.692233in}}%
\pgfpathlineto{\pgfqpoint{6.094031in}{0.653154in}}%
\pgfpathlineto{\pgfqpoint{6.097662in}{0.634905in}}%
\pgfpathlineto{\pgfqpoint{6.101293in}{0.695589in}}%
\pgfpathlineto{\pgfqpoint{6.104924in}{0.681211in}}%
\pgfpathlineto{\pgfqpoint{6.108556in}{0.690631in}}%
\pgfpathlineto{\pgfqpoint{6.112187in}{0.671344in}}%
\pgfpathlineto{\pgfqpoint{6.115818in}{0.500548in}}%
\pgfpathlineto{\pgfqpoint{6.119449in}{0.696544in}}%
\pgfpathlineto{\pgfqpoint{6.123080in}{0.632908in}}%
\pgfpathlineto{\pgfqpoint{6.126711in}{0.675161in}}%
\pgfpathlineto{\pgfqpoint{6.130343in}{0.677187in}}%
\pgfpathlineto{\pgfqpoint{6.133974in}{0.645043in}}%
\pgfpathlineto{\pgfqpoint{6.137605in}{0.689720in}}%
\pgfpathlineto{\pgfqpoint{6.141236in}{0.620504in}}%
\pgfpathlineto{\pgfqpoint{6.148498in}{0.668574in}}%
\pgfpathlineto{\pgfqpoint{6.152130in}{0.672896in}}%
\pgfpathlineto{\pgfqpoint{6.155761in}{0.670525in}}%
\pgfpathlineto{\pgfqpoint{6.159392in}{0.697941in}}%
\pgfpathlineto{\pgfqpoint{6.163023in}{0.680456in}}%
\pgfpathlineto{\pgfqpoint{6.166654in}{0.699879in}}%
\pgfpathlineto{\pgfqpoint{6.170285in}{0.680408in}}%
\pgfpathlineto{\pgfqpoint{6.173917in}{0.636881in}}%
\pgfpathlineto{\pgfqpoint{6.177548in}{0.675250in}}%
\pgfpathlineto{\pgfqpoint{6.181179in}{0.689712in}}%
\pgfpathlineto{\pgfqpoint{6.184810in}{0.711705in}}%
\pgfpathlineto{\pgfqpoint{6.188441in}{0.709942in}}%
\pgfpathlineto{\pgfqpoint{6.192073in}{0.658316in}}%
\pgfpathlineto{\pgfqpoint{6.195704in}{0.558152in}}%
\pgfpathlineto{\pgfqpoint{6.199335in}{0.667834in}}%
\pgfpathlineto{\pgfqpoint{6.202966in}{0.667898in}}%
\pgfpathlineto{\pgfqpoint{6.206597in}{0.714303in}}%
\pgfpathlineto{\pgfqpoint{6.210228in}{0.682656in}}%
\pgfpathlineto{\pgfqpoint{6.213860in}{0.622070in}}%
\pgfpathlineto{\pgfqpoint{6.221122in}{0.689767in}}%
\pgfpathlineto{\pgfqpoint{6.224753in}{0.701556in}}%
\pgfpathlineto{\pgfqpoint{6.228384in}{0.692617in}}%
\pgfpathlineto{\pgfqpoint{6.232015in}{0.656276in}}%
\pgfpathlineto{\pgfqpoint{6.235647in}{0.628537in}}%
\pgfpathlineto{\pgfqpoint{6.239278in}{0.696767in}}%
\pgfpathlineto{\pgfqpoint{6.242909in}{0.678837in}}%
\pgfpathlineto{\pgfqpoint{6.246540in}{0.689090in}}%
\pgfpathlineto{\pgfqpoint{6.250171in}{0.672249in}}%
\pgfpathlineto{\pgfqpoint{6.253802in}{0.569758in}}%
\pgfpathlineto{\pgfqpoint{6.257434in}{0.696691in}}%
\pgfpathlineto{\pgfqpoint{6.261065in}{0.625236in}}%
\pgfpathlineto{\pgfqpoint{6.264696in}{0.672377in}}%
\pgfpathlineto{\pgfqpoint{6.268327in}{0.676253in}}%
\pgfpathlineto{\pgfqpoint{6.271958in}{0.650052in}}%
\pgfpathlineto{\pgfqpoint{6.275590in}{0.688792in}}%
\pgfpathlineto{\pgfqpoint{6.279221in}{0.626298in}}%
\pgfpathlineto{\pgfqpoint{6.282852in}{0.639018in}}%
\pgfpathlineto{\pgfqpoint{6.286483in}{0.665425in}}%
\pgfpathlineto{\pgfqpoint{6.290114in}{0.674840in}}%
\pgfpathlineto{\pgfqpoint{6.293745in}{0.667503in}}%
\pgfpathlineto{\pgfqpoint{6.297377in}{0.699574in}}%
\pgfpathlineto{\pgfqpoint{6.301008in}{0.678078in}}%
\pgfpathlineto{\pgfqpoint{6.304639in}{0.698759in}}%
\pgfpathlineto{\pgfqpoint{6.308270in}{0.680607in}}%
\pgfpathlineto{\pgfqpoint{6.311901in}{0.641786in}}%
\pgfpathlineto{\pgfqpoint{6.315532in}{0.673560in}}%
\pgfpathlineto{\pgfqpoint{6.322795in}{0.710948in}}%
\pgfpathlineto{\pgfqpoint{6.326426in}{0.711440in}}%
\pgfpathlineto{\pgfqpoint{6.330057in}{0.661919in}}%
\pgfpathlineto{\pgfqpoint{6.333688in}{0.505980in}}%
\pgfpathlineto{\pgfqpoint{6.337319in}{0.666967in}}%
\pgfpathlineto{\pgfqpoint{6.340951in}{0.672324in}}%
\pgfpathlineto{\pgfqpoint{6.344582in}{0.714398in}}%
\pgfpathlineto{\pgfqpoint{6.348213in}{0.684275in}}%
\pgfpathlineto{\pgfqpoint{6.351844in}{0.628048in}}%
\pgfpathlineto{\pgfqpoint{6.355475in}{0.651066in}}%
\pgfpathlineto{\pgfqpoint{6.359106in}{0.691573in}}%
\pgfpathlineto{\pgfqpoint{6.362738in}{0.700476in}}%
\pgfpathlineto{\pgfqpoint{6.366369in}{0.692828in}}%
\pgfpathlineto{\pgfqpoint{6.370000in}{0.629942in}}%
\pgfpathlineto{\pgfqpoint{6.370000in}{0.629942in}}%
\pgfusepath{stroke}%
\end{pgfscope}%
\begin{pgfscope}%
\pgfpathrectangle{\pgfqpoint{4.554414in}{0.320625in}}{\pgfqpoint{1.815586in}{1.164375in}}%
\pgfusepath{clip}%
\pgfsetrectcap%
\pgfsetroundjoin%
\pgfsetlinewidth{0.602250pt}%
\definecolor{currentstroke}{rgb}{0.141176,0.345098,0.650980}%
\pgfsetstrokecolor{currentstroke}%
\pgfsetstrokeopacity{0.120000}%
\pgfsetdash{}{0pt}%
\pgfpathmoveto{\pgfqpoint{4.554414in}{1.471508in}}%
\pgfpathlineto{\pgfqpoint{4.561676in}{1.450401in}}%
\pgfpathlineto{\pgfqpoint{4.568938in}{1.433756in}}%
\pgfpathlineto{\pgfqpoint{4.576201in}{1.419684in}}%
\pgfpathlineto{\pgfqpoint{4.587094in}{1.401735in}}%
\pgfpathlineto{\pgfqpoint{4.605250in}{1.374622in}}%
\pgfpathlineto{\pgfqpoint{4.616144in}{1.359558in}}%
\pgfpathlineto{\pgfqpoint{4.634299in}{1.336636in}}%
\pgfpathlineto{\pgfqpoint{4.656086in}{1.311113in}}%
\pgfpathlineto{\pgfqpoint{4.677873in}{1.287498in}}%
\pgfpathlineto{\pgfqpoint{4.699660in}{1.265581in}}%
\pgfpathlineto{\pgfqpoint{4.725079in}{1.241815in}}%
\pgfpathlineto{\pgfqpoint{4.750497in}{1.219596in}}%
\pgfpathlineto{\pgfqpoint{4.779546in}{1.195632in}}%
\pgfpathlineto{\pgfqpoint{4.815858in}{1.167205in}}%
\pgfpathlineto{\pgfqpoint{4.859432in}{1.134579in}}%
\pgfpathlineto{\pgfqpoint{4.917531in}{1.092611in}}%
\pgfpathlineto{\pgfqpoint{4.993785in}{1.039006in}}%
\pgfpathlineto{\pgfqpoint{5.099090in}{0.966411in}}%
\pgfpathlineto{\pgfqpoint{5.229812in}{0.877642in}}%
\pgfpathlineto{\pgfqpoint{5.349640in}{0.797464in}}%
\pgfpathlineto{\pgfqpoint{5.415002in}{0.754237in}}%
\pgfpathlineto{\pgfqpoint{5.418633in}{0.752972in}}%
\pgfpathlineto{\pgfqpoint{5.422264in}{0.749470in}}%
\pgfpathlineto{\pgfqpoint{5.433157in}{0.742769in}}%
\pgfpathlineto{\pgfqpoint{5.436789in}{0.741090in}}%
\pgfpathlineto{\pgfqpoint{5.440420in}{0.737908in}}%
\pgfpathlineto{\pgfqpoint{5.444051in}{0.736858in}}%
\pgfpathlineto{\pgfqpoint{5.447682in}{0.733041in}}%
\pgfpathlineto{\pgfqpoint{5.451313in}{0.730398in}}%
\pgfpathlineto{\pgfqpoint{5.454944in}{0.728905in}}%
\pgfpathlineto{\pgfqpoint{5.458576in}{0.725676in}}%
\pgfpathlineto{\pgfqpoint{5.465838in}{0.724417in}}%
\pgfpathlineto{\pgfqpoint{5.469469in}{0.718532in}}%
\pgfpathlineto{\pgfqpoint{5.476731in}{0.716352in}}%
\pgfpathlineto{\pgfqpoint{5.480363in}{0.713608in}}%
\pgfpathlineto{\pgfqpoint{5.483994in}{0.709393in}}%
\pgfpathlineto{\pgfqpoint{5.487625in}{0.707673in}}%
\pgfpathlineto{\pgfqpoint{5.491256in}{0.710348in}}%
\pgfpathlineto{\pgfqpoint{5.494887in}{0.706014in}}%
\pgfpathlineto{\pgfqpoint{5.498519in}{0.699843in}}%
\pgfpathlineto{\pgfqpoint{5.502150in}{0.713842in}}%
\pgfpathlineto{\pgfqpoint{5.505781in}{0.694968in}}%
\pgfpathlineto{\pgfqpoint{5.509412in}{0.693587in}}%
\pgfpathlineto{\pgfqpoint{5.513043in}{0.690561in}}%
\pgfpathlineto{\pgfqpoint{5.516674in}{0.692127in}}%
\pgfpathlineto{\pgfqpoint{5.520306in}{0.697049in}}%
\pgfpathlineto{\pgfqpoint{5.523937in}{0.690462in}}%
\pgfpathlineto{\pgfqpoint{5.527568in}{0.691627in}}%
\pgfpathlineto{\pgfqpoint{5.531199in}{0.679282in}}%
\pgfpathlineto{\pgfqpoint{5.534830in}{0.679547in}}%
\pgfpathlineto{\pgfqpoint{5.538461in}{0.687833in}}%
\pgfpathlineto{\pgfqpoint{5.542093in}{0.675766in}}%
\pgfpathlineto{\pgfqpoint{5.545724in}{0.680270in}}%
\pgfpathlineto{\pgfqpoint{5.549355in}{0.707879in}}%
\pgfpathlineto{\pgfqpoint{5.552986in}{0.664568in}}%
\pgfpathlineto{\pgfqpoint{5.556617in}{0.665951in}}%
\pgfpathlineto{\pgfqpoint{5.560248in}{0.661727in}}%
\pgfpathlineto{\pgfqpoint{5.563880in}{0.672859in}}%
\pgfpathlineto{\pgfqpoint{5.567511in}{0.673703in}}%
\pgfpathlineto{\pgfqpoint{5.571142in}{0.653167in}}%
\pgfpathlineto{\pgfqpoint{5.574773in}{0.690620in}}%
\pgfpathlineto{\pgfqpoint{5.578404in}{0.680045in}}%
\pgfpathlineto{\pgfqpoint{5.582035in}{0.654682in}}%
\pgfpathlineto{\pgfqpoint{5.585667in}{0.679625in}}%
\pgfpathlineto{\pgfqpoint{5.589298in}{0.645482in}}%
\pgfpathlineto{\pgfqpoint{5.592929in}{0.663655in}}%
\pgfpathlineto{\pgfqpoint{5.596560in}{0.655122in}}%
\pgfpathlineto{\pgfqpoint{5.600191in}{0.643251in}}%
\pgfpathlineto{\pgfqpoint{5.603823in}{0.651202in}}%
\pgfpathlineto{\pgfqpoint{5.607454in}{0.650591in}}%
\pgfpathlineto{\pgfqpoint{5.611085in}{0.662933in}}%
\pgfpathlineto{\pgfqpoint{5.614716in}{0.641122in}}%
\pgfpathlineto{\pgfqpoint{5.618347in}{0.659937in}}%
\pgfpathlineto{\pgfqpoint{5.621978in}{0.692151in}}%
\pgfpathlineto{\pgfqpoint{5.625610in}{0.635369in}}%
\pgfpathlineto{\pgfqpoint{5.629241in}{0.625375in}}%
\pgfpathlineto{\pgfqpoint{5.632872in}{0.648266in}}%
\pgfpathlineto{\pgfqpoint{5.636503in}{0.632132in}}%
\pgfpathlineto{\pgfqpoint{5.640134in}{0.682918in}}%
\pgfpathlineto{\pgfqpoint{5.643765in}{0.691463in}}%
\pgfpathlineto{\pgfqpoint{5.647397in}{0.648848in}}%
\pgfpathlineto{\pgfqpoint{5.651028in}{0.633215in}}%
\pgfpathlineto{\pgfqpoint{5.654659in}{0.613217in}}%
\pgfpathlineto{\pgfqpoint{5.658290in}{0.669813in}}%
\pgfpathlineto{\pgfqpoint{5.661921in}{0.672399in}}%
\pgfpathlineto{\pgfqpoint{5.665552in}{0.633346in}}%
\pgfpathlineto{\pgfqpoint{5.669184in}{0.680085in}}%
\pgfpathlineto{\pgfqpoint{5.672815in}{0.592760in}}%
\pgfpathlineto{\pgfqpoint{5.676446in}{0.652788in}}%
\pgfpathlineto{\pgfqpoint{5.680077in}{0.656603in}}%
\pgfpathlineto{\pgfqpoint{5.683708in}{0.591523in}}%
\pgfpathlineto{\pgfqpoint{5.687340in}{0.682868in}}%
\pgfpathlineto{\pgfqpoint{5.690971in}{0.702871in}}%
\pgfpathlineto{\pgfqpoint{5.694602in}{0.649454in}}%
\pgfpathlineto{\pgfqpoint{5.698233in}{0.626460in}}%
\pgfpathlineto{\pgfqpoint{5.701864in}{0.585608in}}%
\pgfpathlineto{\pgfqpoint{5.705495in}{0.663299in}}%
\pgfpathlineto{\pgfqpoint{5.709127in}{0.641700in}}%
\pgfpathlineto{\pgfqpoint{5.712758in}{0.660520in}}%
\pgfpathlineto{\pgfqpoint{5.716389in}{0.702408in}}%
\pgfpathlineto{\pgfqpoint{5.720020in}{0.663499in}}%
\pgfpathlineto{\pgfqpoint{5.723651in}{0.559224in}}%
\pgfpathlineto{\pgfqpoint{5.727282in}{0.662434in}}%
\pgfpathlineto{\pgfqpoint{5.730914in}{0.643011in}}%
\pgfpathlineto{\pgfqpoint{5.734545in}{0.650905in}}%
\pgfpathlineto{\pgfqpoint{5.738176in}{0.569163in}}%
\pgfpathlineto{\pgfqpoint{5.741807in}{0.670135in}}%
\pgfpathlineto{\pgfqpoint{5.745438in}{0.642305in}}%
\pgfpathlineto{\pgfqpoint{5.749069in}{0.592652in}}%
\pgfpathlineto{\pgfqpoint{5.752701in}{0.656848in}}%
\pgfpathlineto{\pgfqpoint{5.756332in}{0.559405in}}%
\pgfpathlineto{\pgfqpoint{5.759963in}{0.676131in}}%
\pgfpathlineto{\pgfqpoint{5.763594in}{0.698465in}}%
\pgfpathlineto{\pgfqpoint{5.767225in}{0.539114in}}%
\pgfpathlineto{\pgfqpoint{5.770856in}{0.623097in}}%
\pgfpathlineto{\pgfqpoint{5.774488in}{0.597605in}}%
\pgfpathlineto{\pgfqpoint{5.778119in}{0.652956in}}%
\pgfpathlineto{\pgfqpoint{5.781750in}{0.683690in}}%
\pgfpathlineto{\pgfqpoint{5.785381in}{0.677448in}}%
\pgfpathlineto{\pgfqpoint{5.789012in}{0.668653in}}%
\pgfpathlineto{\pgfqpoint{5.792644in}{0.603837in}}%
\pgfpathlineto{\pgfqpoint{5.799906in}{0.669937in}}%
\pgfpathlineto{\pgfqpoint{5.803537in}{0.652264in}}%
\pgfpathlineto{\pgfqpoint{5.807168in}{0.542233in}}%
\pgfpathlineto{\pgfqpoint{5.810799in}{0.674506in}}%
\pgfpathlineto{\pgfqpoint{5.814431in}{0.628385in}}%
\pgfpathlineto{\pgfqpoint{5.818062in}{0.652774in}}%
\pgfpathlineto{\pgfqpoint{5.821693in}{0.607119in}}%
\pgfpathlineto{\pgfqpoint{5.825324in}{0.655004in}}%
\pgfpathlineto{\pgfqpoint{5.828955in}{0.647655in}}%
\pgfpathlineto{\pgfqpoint{5.832586in}{0.618547in}}%
\pgfpathlineto{\pgfqpoint{5.836218in}{0.663180in}}%
\pgfpathlineto{\pgfqpoint{5.839849in}{0.616583in}}%
\pgfpathlineto{\pgfqpoint{5.843480in}{0.662999in}}%
\pgfpathlineto{\pgfqpoint{5.847111in}{0.694600in}}%
\pgfpathlineto{\pgfqpoint{5.850742in}{0.595013in}}%
\pgfpathlineto{\pgfqpoint{5.854373in}{0.623133in}}%
\pgfpathlineto{\pgfqpoint{5.858005in}{0.634899in}}%
\pgfpathlineto{\pgfqpoint{5.861636in}{0.640691in}}%
\pgfpathlineto{\pgfqpoint{5.865267in}{0.681224in}}%
\pgfpathlineto{\pgfqpoint{5.868898in}{0.679011in}}%
\pgfpathlineto{\pgfqpoint{5.872529in}{0.647438in}}%
\pgfpathlineto{\pgfqpoint{5.876160in}{0.659771in}}%
\pgfpathlineto{\pgfqpoint{5.879792in}{0.644316in}}%
\pgfpathlineto{\pgfqpoint{5.883423in}{0.667603in}}%
\pgfpathlineto{\pgfqpoint{5.887054in}{0.553519in}}%
\pgfpathlineto{\pgfqpoint{5.890685in}{0.662867in}}%
\pgfpathlineto{\pgfqpoint{5.894316in}{0.684405in}}%
\pgfpathlineto{\pgfqpoint{5.897948in}{0.657958in}}%
\pgfpathlineto{\pgfqpoint{5.901579in}{0.640690in}}%
\pgfpathlineto{\pgfqpoint{5.905210in}{0.701451in}}%
\pgfpathlineto{\pgfqpoint{5.908841in}{0.639357in}}%
\pgfpathlineto{\pgfqpoint{5.912472in}{0.624369in}}%
\pgfpathlineto{\pgfqpoint{5.916103in}{0.598940in}}%
\pgfpathlineto{\pgfqpoint{5.919735in}{0.667728in}}%
\pgfpathlineto{\pgfqpoint{5.923366in}{0.680125in}}%
\pgfpathlineto{\pgfqpoint{5.926997in}{0.656428in}}%
\pgfpathlineto{\pgfqpoint{5.930628in}{0.678581in}}%
\pgfpathlineto{\pgfqpoint{5.934259in}{0.548741in}}%
\pgfpathlineto{\pgfqpoint{5.937890in}{0.648819in}}%
\pgfpathlineto{\pgfqpoint{5.941522in}{0.653776in}}%
\pgfpathlineto{\pgfqpoint{5.945153in}{0.566167in}}%
\pgfpathlineto{\pgfqpoint{5.948784in}{0.627396in}}%
\pgfpathlineto{\pgfqpoint{5.952415in}{0.657754in}}%
\pgfpathlineto{\pgfqpoint{5.959677in}{0.643692in}}%
\pgfpathlineto{\pgfqpoint{5.963309in}{0.625412in}}%
\pgfpathlineto{\pgfqpoint{5.966940in}{0.681490in}}%
\pgfpathlineto{\pgfqpoint{5.970571in}{0.640208in}}%
\pgfpathlineto{\pgfqpoint{5.974202in}{0.516708in}}%
\pgfpathlineto{\pgfqpoint{5.977833in}{0.656081in}}%
\pgfpathlineto{\pgfqpoint{5.981465in}{0.556293in}}%
\pgfpathlineto{\pgfqpoint{5.985096in}{0.678241in}}%
\pgfpathlineto{\pgfqpoint{5.988727in}{0.696790in}}%
\pgfpathlineto{\pgfqpoint{5.992358in}{0.610316in}}%
\pgfpathlineto{\pgfqpoint{5.995989in}{0.630140in}}%
\pgfpathlineto{\pgfqpoint{5.999620in}{0.548592in}}%
\pgfpathlineto{\pgfqpoint{6.003252in}{0.661700in}}%
\pgfpathlineto{\pgfqpoint{6.006883in}{0.679172in}}%
\pgfpathlineto{\pgfqpoint{6.010514in}{0.660191in}}%
\pgfpathlineto{\pgfqpoint{6.014145in}{0.675699in}}%
\pgfpathlineto{\pgfqpoint{6.017776in}{0.582978in}}%
\pgfpathlineto{\pgfqpoint{6.021407in}{0.644937in}}%
\pgfpathlineto{\pgfqpoint{6.025039in}{0.667997in}}%
\pgfpathlineto{\pgfqpoint{6.028670in}{0.623782in}}%
\pgfpathlineto{\pgfqpoint{6.032301in}{0.671092in}}%
\pgfpathlineto{\pgfqpoint{6.035932in}{0.706980in}}%
\pgfpathlineto{\pgfqpoint{6.039563in}{0.597562in}}%
\pgfpathlineto{\pgfqpoint{6.043194in}{0.625095in}}%
\pgfpathlineto{\pgfqpoint{6.046826in}{0.579366in}}%
\pgfpathlineto{\pgfqpoint{6.050457in}{0.660534in}}%
\pgfpathlineto{\pgfqpoint{6.054088in}{0.657276in}}%
\pgfpathlineto{\pgfqpoint{6.057719in}{0.623189in}}%
\pgfpathlineto{\pgfqpoint{6.061350in}{0.693116in}}%
\pgfpathlineto{\pgfqpoint{6.064981in}{0.675204in}}%
\pgfpathlineto{\pgfqpoint{6.068613in}{0.626917in}}%
\pgfpathlineto{\pgfqpoint{6.072244in}{0.672391in}}%
\pgfpathlineto{\pgfqpoint{6.075875in}{0.623689in}}%
\pgfpathlineto{\pgfqpoint{6.079506in}{0.654200in}}%
\pgfpathlineto{\pgfqpoint{6.083137in}{0.621744in}}%
\pgfpathlineto{\pgfqpoint{6.086769in}{0.641102in}}%
\pgfpathlineto{\pgfqpoint{6.090400in}{0.640364in}}%
\pgfpathlineto{\pgfqpoint{6.094031in}{0.630892in}}%
\pgfpathlineto{\pgfqpoint{6.097662in}{0.658332in}}%
\pgfpathlineto{\pgfqpoint{6.101293in}{0.614850in}}%
\pgfpathlineto{\pgfqpoint{6.104924in}{0.665130in}}%
\pgfpathlineto{\pgfqpoint{6.108556in}{0.694679in}}%
\pgfpathlineto{\pgfqpoint{6.112187in}{0.609283in}}%
\pgfpathlineto{\pgfqpoint{6.115818in}{0.609349in}}%
\pgfpathlineto{\pgfqpoint{6.119449in}{0.632134in}}%
\pgfpathlineto{\pgfqpoint{6.123080in}{0.634255in}}%
\pgfpathlineto{\pgfqpoint{6.126711in}{0.684321in}}%
\pgfpathlineto{\pgfqpoint{6.130343in}{0.689573in}}%
\pgfpathlineto{\pgfqpoint{6.137605in}{0.621215in}}%
\pgfpathlineto{\pgfqpoint{6.141236in}{0.615496in}}%
\pgfpathlineto{\pgfqpoint{6.144867in}{0.669749in}}%
\pgfpathlineto{\pgfqpoint{6.148498in}{0.669322in}}%
\pgfpathlineto{\pgfqpoint{6.152130in}{0.620195in}}%
\pgfpathlineto{\pgfqpoint{6.155761in}{0.679494in}}%
\pgfpathlineto{\pgfqpoint{6.159392in}{0.591667in}}%
\pgfpathlineto{\pgfqpoint{6.163023in}{0.652553in}}%
\pgfpathlineto{\pgfqpoint{6.166654in}{0.646624in}}%
\pgfpathlineto{\pgfqpoint{6.170285in}{0.608917in}}%
\pgfpathlineto{\pgfqpoint{6.173917in}{0.642594in}}%
\pgfpathlineto{\pgfqpoint{6.177548in}{0.646361in}}%
\pgfpathlineto{\pgfqpoint{6.181179in}{0.658554in}}%
\pgfpathlineto{\pgfqpoint{6.184810in}{0.638150in}}%
\pgfpathlineto{\pgfqpoint{6.188441in}{0.643036in}}%
\pgfpathlineto{\pgfqpoint{6.192073in}{0.688306in}}%
\pgfpathlineto{\pgfqpoint{6.195704in}{0.631411in}}%
\pgfpathlineto{\pgfqpoint{6.199335in}{0.598628in}}%
\pgfpathlineto{\pgfqpoint{6.202966in}{0.652503in}}%
\pgfpathlineto{\pgfqpoint{6.206597in}{0.607907in}}%
\pgfpathlineto{\pgfqpoint{6.210228in}{0.679449in}}%
\pgfpathlineto{\pgfqpoint{6.213860in}{0.690981in}}%
\pgfpathlineto{\pgfqpoint{6.217491in}{0.491474in}}%
\pgfpathlineto{\pgfqpoint{6.221122in}{0.653064in}}%
\pgfpathlineto{\pgfqpoint{6.224753in}{0.664864in}}%
\pgfpathlineto{\pgfqpoint{6.228384in}{0.663779in}}%
\pgfpathlineto{\pgfqpoint{6.232015in}{0.610793in}}%
\pgfpathlineto{\pgfqpoint{6.235647in}{0.652256in}}%
\pgfpathlineto{\pgfqpoint{6.239278in}{0.684221in}}%
\pgfpathlineto{\pgfqpoint{6.242909in}{0.676952in}}%
\pgfpathlineto{\pgfqpoint{6.246540in}{0.474432in}}%
\pgfpathlineto{\pgfqpoint{6.250171in}{0.701877in}}%
\pgfpathlineto{\pgfqpoint{6.253802in}{0.599013in}}%
\pgfpathlineto{\pgfqpoint{6.257434in}{0.626050in}}%
\pgfpathlineto{\pgfqpoint{6.261065in}{0.540561in}}%
\pgfpathlineto{\pgfqpoint{6.264696in}{0.658743in}}%
\pgfpathlineto{\pgfqpoint{6.268327in}{0.685324in}}%
\pgfpathlineto{\pgfqpoint{6.271958in}{0.676356in}}%
\pgfpathlineto{\pgfqpoint{6.275590in}{0.672508in}}%
\pgfpathlineto{\pgfqpoint{6.279221in}{0.578253in}}%
\pgfpathlineto{\pgfqpoint{6.282852in}{0.641803in}}%
\pgfpathlineto{\pgfqpoint{6.286483in}{0.665466in}}%
\pgfpathlineto{\pgfqpoint{6.290114in}{0.642247in}}%
\pgfpathlineto{\pgfqpoint{6.293745in}{0.579707in}}%
\pgfpathlineto{\pgfqpoint{6.297377in}{0.670259in}}%
\pgfpathlineto{\pgfqpoint{6.301008in}{0.637086in}}%
\pgfpathlineto{\pgfqpoint{6.304639in}{0.650510in}}%
\pgfpathlineto{\pgfqpoint{6.308270in}{0.527928in}}%
\pgfpathlineto{\pgfqpoint{6.311901in}{0.664321in}}%
\pgfpathlineto{\pgfqpoint{6.315532in}{0.645827in}}%
\pgfpathlineto{\pgfqpoint{6.319164in}{0.602892in}}%
\pgfpathlineto{\pgfqpoint{6.322795in}{0.661535in}}%
\pgfpathlineto{\pgfqpoint{6.326426in}{0.596637in}}%
\pgfpathlineto{\pgfqpoint{6.330057in}{0.669873in}}%
\pgfpathlineto{\pgfqpoint{6.333688in}{0.696609in}}%
\pgfpathlineto{\pgfqpoint{6.337319in}{0.554691in}}%
\pgfpathlineto{\pgfqpoint{6.340951in}{0.624639in}}%
\pgfpathlineto{\pgfqpoint{6.344582in}{0.619898in}}%
\pgfpathlineto{\pgfqpoint{6.348213in}{0.648171in}}%
\pgfpathlineto{\pgfqpoint{6.351844in}{0.682255in}}%
\pgfpathlineto{\pgfqpoint{6.355475in}{0.677998in}}%
\pgfpathlineto{\pgfqpoint{6.359106in}{0.665741in}}%
\pgfpathlineto{\pgfqpoint{6.362738in}{0.616357in}}%
\pgfpathlineto{\pgfqpoint{6.366369in}{0.607815in}}%
\pgfpathlineto{\pgfqpoint{6.370000in}{0.655674in}}%
\pgfpathlineto{\pgfqpoint{6.370000in}{0.655674in}}%
\pgfusepath{stroke}%
\end{pgfscope}%
\begin{pgfscope}%
\pgfpathrectangle{\pgfqpoint{4.554414in}{0.320625in}}{\pgfqpoint{1.815586in}{1.164375in}}%
\pgfusepath{clip}%
\pgfsetrectcap%
\pgfsetroundjoin%
\pgfsetlinewidth{0.602250pt}%
\definecolor{currentstroke}{rgb}{0.141176,0.345098,0.650980}%
\pgfsetstrokecolor{currentstroke}%
\pgfsetstrokeopacity{0.120000}%
\pgfsetdash{}{0pt}%
\pgfpathmoveto{\pgfqpoint{4.554414in}{1.471756in}}%
\pgfpathlineto{\pgfqpoint{4.561676in}{1.450691in}}%
\pgfpathlineto{\pgfqpoint{4.565307in}{1.442031in}}%
\pgfpathlineto{\pgfqpoint{4.572569in}{1.429568in}}%
\pgfpathlineto{\pgfqpoint{4.583463in}{1.414890in}}%
\pgfpathlineto{\pgfqpoint{4.608881in}{1.382646in}}%
\pgfpathlineto{\pgfqpoint{4.630668in}{1.356914in}}%
\pgfpathlineto{\pgfqpoint{4.648824in}{1.337209in}}%
\pgfpathlineto{\pgfqpoint{4.666980in}{1.319183in}}%
\pgfpathlineto{\pgfqpoint{4.688767in}{1.299317in}}%
\pgfpathlineto{\pgfqpoint{4.714185in}{1.277813in}}%
\pgfpathlineto{\pgfqpoint{4.743235in}{1.254683in}}%
\pgfpathlineto{\pgfqpoint{4.783177in}{1.224415in}}%
\pgfpathlineto{\pgfqpoint{4.837645in}{1.184731in}}%
\pgfpathlineto{\pgfqpoint{4.910269in}{1.133285in}}%
\pgfpathlineto{\pgfqpoint{5.008310in}{1.065247in}}%
\pgfpathlineto{\pgfqpoint{5.128139in}{0.983463in}}%
\pgfpathlineto{\pgfqpoint{5.269755in}{0.888197in}}%
\pgfpathlineto{\pgfqpoint{5.407739in}{0.796492in}}%
\pgfpathlineto{\pgfqpoint{5.444051in}{0.772828in}}%
\pgfpathlineto{\pgfqpoint{5.451313in}{0.767856in}}%
\pgfpathlineto{\pgfqpoint{5.454944in}{0.766280in}}%
\pgfpathlineto{\pgfqpoint{5.458576in}{0.762874in}}%
\pgfpathlineto{\pgfqpoint{5.473100in}{0.753712in}}%
\pgfpathlineto{\pgfqpoint{5.476731in}{0.752278in}}%
\pgfpathlineto{\pgfqpoint{5.480363in}{0.748548in}}%
\pgfpathlineto{\pgfqpoint{5.483994in}{0.746458in}}%
\pgfpathlineto{\pgfqpoint{5.487625in}{0.746314in}}%
\pgfpathlineto{\pgfqpoint{5.491256in}{0.741330in}}%
\pgfpathlineto{\pgfqpoint{5.494887in}{0.739317in}}%
\pgfpathlineto{\pgfqpoint{5.498519in}{0.738124in}}%
\pgfpathlineto{\pgfqpoint{5.502150in}{0.734154in}}%
\pgfpathlineto{\pgfqpoint{5.505781in}{0.733066in}}%
\pgfpathlineto{\pgfqpoint{5.509412in}{0.734206in}}%
\pgfpathlineto{\pgfqpoint{5.513043in}{0.727254in}}%
\pgfpathlineto{\pgfqpoint{5.516674in}{0.725120in}}%
\pgfpathlineto{\pgfqpoint{5.520306in}{0.727684in}}%
\pgfpathlineto{\pgfqpoint{5.523937in}{0.719821in}}%
\pgfpathlineto{\pgfqpoint{5.527568in}{0.718760in}}%
\pgfpathlineto{\pgfqpoint{5.531199in}{0.719423in}}%
\pgfpathlineto{\pgfqpoint{5.534830in}{0.712732in}}%
\pgfpathlineto{\pgfqpoint{5.538461in}{0.713353in}}%
\pgfpathlineto{\pgfqpoint{5.542093in}{0.721166in}}%
\pgfpathlineto{\pgfqpoint{5.545724in}{0.706259in}}%
\pgfpathlineto{\pgfqpoint{5.549355in}{0.704048in}}%
\pgfpathlineto{\pgfqpoint{5.552986in}{0.710536in}}%
\pgfpathlineto{\pgfqpoint{5.556617in}{0.698361in}}%
\pgfpathlineto{\pgfqpoint{5.560248in}{0.700656in}}%
\pgfpathlineto{\pgfqpoint{5.563880in}{0.706766in}}%
\pgfpathlineto{\pgfqpoint{5.567511in}{0.692598in}}%
\pgfpathlineto{\pgfqpoint{5.571142in}{0.693516in}}%
\pgfpathlineto{\pgfqpoint{5.574773in}{0.712954in}}%
\pgfpathlineto{\pgfqpoint{5.578404in}{0.685104in}}%
\pgfpathlineto{\pgfqpoint{5.582035in}{0.684285in}}%
\pgfpathlineto{\pgfqpoint{5.585667in}{0.695613in}}%
\pgfpathlineto{\pgfqpoint{5.589298in}{0.676908in}}%
\pgfpathlineto{\pgfqpoint{5.592929in}{0.688331in}}%
\pgfpathlineto{\pgfqpoint{5.596560in}{0.704349in}}%
\pgfpathlineto{\pgfqpoint{5.600191in}{0.678644in}}%
\pgfpathlineto{\pgfqpoint{5.603823in}{0.671294in}}%
\pgfpathlineto{\pgfqpoint{5.607454in}{0.705959in}}%
\pgfpathlineto{\pgfqpoint{5.611085in}{0.663264in}}%
\pgfpathlineto{\pgfqpoint{5.614716in}{0.669404in}}%
\pgfpathlineto{\pgfqpoint{5.618347in}{0.686039in}}%
\pgfpathlineto{\pgfqpoint{5.621978in}{0.656377in}}%
\pgfpathlineto{\pgfqpoint{5.629241in}{0.707598in}}%
\pgfpathlineto{\pgfqpoint{5.632872in}{0.671451in}}%
\pgfpathlineto{\pgfqpoint{5.636503in}{0.647815in}}%
\pgfpathlineto{\pgfqpoint{5.640134in}{0.697304in}}%
\pgfpathlineto{\pgfqpoint{5.643765in}{0.641297in}}%
\pgfpathlineto{\pgfqpoint{5.651028in}{0.686125in}}%
\pgfpathlineto{\pgfqpoint{5.654659in}{0.650922in}}%
\pgfpathlineto{\pgfqpoint{5.658290in}{0.668018in}}%
\pgfpathlineto{\pgfqpoint{5.661921in}{0.709534in}}%
\pgfpathlineto{\pgfqpoint{5.665552in}{0.659969in}}%
\pgfpathlineto{\pgfqpoint{5.669184in}{0.628248in}}%
\pgfpathlineto{\pgfqpoint{5.672815in}{0.686533in}}%
\pgfpathlineto{\pgfqpoint{5.676446in}{0.619859in}}%
\pgfpathlineto{\pgfqpoint{5.680077in}{0.670363in}}%
\pgfpathlineto{\pgfqpoint{5.683708in}{0.694477in}}%
\pgfpathlineto{\pgfqpoint{5.687340in}{0.669372in}}%
\pgfpathlineto{\pgfqpoint{5.690971in}{0.629517in}}%
\pgfpathlineto{\pgfqpoint{5.694602in}{0.707577in}}%
\pgfpathlineto{\pgfqpoint{5.698233in}{0.637226in}}%
\pgfpathlineto{\pgfqpoint{5.701864in}{0.626766in}}%
\pgfpathlineto{\pgfqpoint{5.705495in}{0.676975in}}%
\pgfpathlineto{\pgfqpoint{5.709127in}{0.603281in}}%
\pgfpathlineto{\pgfqpoint{5.712758in}{0.676503in}}%
\pgfpathlineto{\pgfqpoint{5.716389in}{0.704027in}}%
\pgfpathlineto{\pgfqpoint{5.720020in}{0.682934in}}%
\pgfpathlineto{\pgfqpoint{5.723651in}{0.621128in}}%
\pgfpathlineto{\pgfqpoint{5.727282in}{0.701618in}}%
\pgfpathlineto{\pgfqpoint{5.730914in}{0.603039in}}%
\pgfpathlineto{\pgfqpoint{5.738176in}{0.675430in}}%
\pgfpathlineto{\pgfqpoint{5.741807in}{0.631975in}}%
\pgfpathlineto{\pgfqpoint{5.749069in}{0.709255in}}%
\pgfpathlineto{\pgfqpoint{5.752701in}{0.681642in}}%
\pgfpathlineto{\pgfqpoint{5.756332in}{0.636826in}}%
\pgfpathlineto{\pgfqpoint{5.759963in}{0.691446in}}%
\pgfpathlineto{\pgfqpoint{5.763594in}{0.567710in}}%
\pgfpathlineto{\pgfqpoint{5.767225in}{0.658068in}}%
\pgfpathlineto{\pgfqpoint{5.770856in}{0.684388in}}%
\pgfpathlineto{\pgfqpoint{5.774488in}{0.670383in}}%
\pgfpathlineto{\pgfqpoint{5.778119in}{0.621665in}}%
\pgfpathlineto{\pgfqpoint{5.781750in}{0.708184in}}%
\pgfpathlineto{\pgfqpoint{5.785381in}{0.663396in}}%
\pgfpathlineto{\pgfqpoint{5.789012in}{0.605694in}}%
\pgfpathlineto{\pgfqpoint{5.792644in}{0.678019in}}%
\pgfpathlineto{\pgfqpoint{5.796275in}{0.558581in}}%
\pgfpathlineto{\pgfqpoint{5.799906in}{0.672060in}}%
\pgfpathlineto{\pgfqpoint{5.803537in}{0.697150in}}%
\pgfpathlineto{\pgfqpoint{5.807168in}{0.695314in}}%
\pgfpathlineto{\pgfqpoint{5.810799in}{0.663235in}}%
\pgfpathlineto{\pgfqpoint{5.814431in}{0.702871in}}%
\pgfpathlineto{\pgfqpoint{5.818062in}{0.628885in}}%
\pgfpathlineto{\pgfqpoint{5.821693in}{0.585915in}}%
\pgfpathlineto{\pgfqpoint{5.825324in}{0.668086in}}%
\pgfpathlineto{\pgfqpoint{5.828955in}{0.611598in}}%
\pgfpathlineto{\pgfqpoint{5.832586in}{0.677837in}}%
\pgfpathlineto{\pgfqpoint{5.836218in}{0.705288in}}%
\pgfpathlineto{\pgfqpoint{5.839849in}{0.644358in}}%
\pgfpathlineto{\pgfqpoint{5.843480in}{0.534957in}}%
\pgfpathlineto{\pgfqpoint{5.847111in}{0.668132in}}%
\pgfpathlineto{\pgfqpoint{5.850742in}{0.603969in}}%
\pgfpathlineto{\pgfqpoint{5.854373in}{0.677425in}}%
\pgfpathlineto{\pgfqpoint{5.858005in}{0.702892in}}%
\pgfpathlineto{\pgfqpoint{5.861636in}{0.647648in}}%
\pgfpathlineto{\pgfqpoint{5.865267in}{0.543269in}}%
\pgfpathlineto{\pgfqpoint{5.868898in}{0.671447in}}%
\pgfpathlineto{\pgfqpoint{5.872529in}{0.620824in}}%
\pgfpathlineto{\pgfqpoint{5.876160in}{0.671576in}}%
\pgfpathlineto{\pgfqpoint{5.879792in}{0.700891in}}%
\pgfpathlineto{\pgfqpoint{5.883423in}{0.650920in}}%
\pgfpathlineto{\pgfqpoint{5.887054in}{0.571809in}}%
\pgfpathlineto{\pgfqpoint{5.890685in}{0.673172in}}%
\pgfpathlineto{\pgfqpoint{5.894316in}{0.628758in}}%
\pgfpathlineto{\pgfqpoint{5.901579in}{0.698854in}}%
\pgfpathlineto{\pgfqpoint{5.905210in}{0.653516in}}%
\pgfpathlineto{\pgfqpoint{5.908841in}{0.587227in}}%
\pgfpathlineto{\pgfqpoint{5.912472in}{0.674493in}}%
\pgfpathlineto{\pgfqpoint{5.916103in}{0.635643in}}%
\pgfpathlineto{\pgfqpoint{5.919735in}{0.660921in}}%
\pgfpathlineto{\pgfqpoint{5.923366in}{0.696834in}}%
\pgfpathlineto{\pgfqpoint{5.926997in}{0.655717in}}%
\pgfpathlineto{\pgfqpoint{5.930628in}{0.599143in}}%
\pgfpathlineto{\pgfqpoint{5.934259in}{0.675354in}}%
\pgfpathlineto{\pgfqpoint{5.937890in}{0.641002in}}%
\pgfpathlineto{\pgfqpoint{5.941522in}{0.654743in}}%
\pgfpathlineto{\pgfqpoint{5.945153in}{0.694824in}}%
\pgfpathlineto{\pgfqpoint{5.948784in}{0.657579in}}%
\pgfpathlineto{\pgfqpoint{5.952415in}{0.608471in}}%
\pgfpathlineto{\pgfqpoint{5.956046in}{0.675880in}}%
\pgfpathlineto{\pgfqpoint{5.959677in}{0.645429in}}%
\pgfpathlineto{\pgfqpoint{5.963309in}{0.647822in}}%
\pgfpathlineto{\pgfqpoint{5.966940in}{0.692819in}}%
\pgfpathlineto{\pgfqpoint{5.970571in}{0.700920in}}%
\pgfpathlineto{\pgfqpoint{5.974202in}{0.695246in}}%
\pgfpathlineto{\pgfqpoint{5.977833in}{0.684779in}}%
\pgfpathlineto{\pgfqpoint{5.985096in}{0.583818in}}%
\pgfpathlineto{\pgfqpoint{5.988727in}{0.674926in}}%
\pgfpathlineto{\pgfqpoint{5.992358in}{0.670313in}}%
\pgfpathlineto{\pgfqpoint{5.995989in}{0.512649in}}%
\pgfpathlineto{\pgfqpoint{5.999620in}{0.686328in}}%
\pgfpathlineto{\pgfqpoint{6.003252in}{0.691078in}}%
\pgfpathlineto{\pgfqpoint{6.006883in}{0.692064in}}%
\pgfpathlineto{\pgfqpoint{6.010514in}{0.679891in}}%
\pgfpathlineto{\pgfqpoint{6.014145in}{0.623057in}}%
\pgfpathlineto{\pgfqpoint{6.017776in}{0.632951in}}%
\pgfpathlineto{\pgfqpoint{6.021407in}{0.680166in}}%
\pgfpathlineto{\pgfqpoint{6.025039in}{0.695335in}}%
\pgfpathlineto{\pgfqpoint{6.028670in}{0.688874in}}%
\pgfpathlineto{\pgfqpoint{6.032301in}{0.678614in}}%
\pgfpathlineto{\pgfqpoint{6.035932in}{0.623960in}}%
\pgfpathlineto{\pgfqpoint{6.039563in}{0.623068in}}%
\pgfpathlineto{\pgfqpoint{6.043194in}{0.677188in}}%
\pgfpathlineto{\pgfqpoint{6.046826in}{0.693351in}}%
\pgfpathlineto{\pgfqpoint{6.050457in}{0.687097in}}%
\pgfpathlineto{\pgfqpoint{6.054088in}{0.676371in}}%
\pgfpathlineto{\pgfqpoint{6.057719in}{0.630846in}}%
\pgfpathlineto{\pgfqpoint{6.061350in}{0.612993in}}%
\pgfpathlineto{\pgfqpoint{6.064981in}{0.676734in}}%
\pgfpathlineto{\pgfqpoint{6.068613in}{0.693759in}}%
\pgfpathlineto{\pgfqpoint{6.072244in}{0.688808in}}%
\pgfpathlineto{\pgfqpoint{6.075875in}{0.673997in}}%
\pgfpathlineto{\pgfqpoint{6.083137in}{0.598992in}}%
\pgfpathlineto{\pgfqpoint{6.086769in}{0.675515in}}%
\pgfpathlineto{\pgfqpoint{6.090400in}{0.693542in}}%
\pgfpathlineto{\pgfqpoint{6.094031in}{0.689677in}}%
\pgfpathlineto{\pgfqpoint{6.097662in}{0.671538in}}%
\pgfpathlineto{\pgfqpoint{6.101293in}{0.638696in}}%
\pgfpathlineto{\pgfqpoint{6.104924in}{0.577734in}}%
\pgfpathlineto{\pgfqpoint{6.108556in}{0.674204in}}%
\pgfpathlineto{\pgfqpoint{6.112187in}{0.693332in}}%
\pgfpathlineto{\pgfqpoint{6.115818in}{0.690627in}}%
\pgfpathlineto{\pgfqpoint{6.119449in}{0.668968in}}%
\pgfpathlineto{\pgfqpoint{6.123080in}{0.641677in}}%
\pgfpathlineto{\pgfqpoint{6.126711in}{0.532833in}}%
\pgfpathlineto{\pgfqpoint{6.130343in}{0.672704in}}%
\pgfpathlineto{\pgfqpoint{6.133974in}{0.693023in}}%
\pgfpathlineto{\pgfqpoint{6.137605in}{0.689250in}}%
\pgfpathlineto{\pgfqpoint{6.141236in}{0.635932in}}%
\pgfpathlineto{\pgfqpoint{6.144867in}{0.665421in}}%
\pgfpathlineto{\pgfqpoint{6.148498in}{0.684378in}}%
\pgfpathlineto{\pgfqpoint{6.155761in}{0.646444in}}%
\pgfpathlineto{\pgfqpoint{6.159392in}{0.591603in}}%
\pgfpathlineto{\pgfqpoint{6.163023in}{0.672154in}}%
\pgfpathlineto{\pgfqpoint{6.166654in}{0.704727in}}%
\pgfpathlineto{\pgfqpoint{6.170285in}{0.714263in}}%
\pgfpathlineto{\pgfqpoint{6.173917in}{0.612783in}}%
\pgfpathlineto{\pgfqpoint{6.177548in}{0.636030in}}%
\pgfpathlineto{\pgfqpoint{6.181179in}{0.655376in}}%
\pgfpathlineto{\pgfqpoint{6.184810in}{0.662613in}}%
\pgfpathlineto{\pgfqpoint{6.188441in}{0.662068in}}%
\pgfpathlineto{\pgfqpoint{6.192073in}{0.583783in}}%
\pgfpathlineto{\pgfqpoint{6.195704in}{0.654203in}}%
\pgfpathlineto{\pgfqpoint{6.199335in}{0.647759in}}%
\pgfpathlineto{\pgfqpoint{6.202966in}{0.665565in}}%
\pgfpathlineto{\pgfqpoint{6.206597in}{0.660299in}}%
\pgfpathlineto{\pgfqpoint{6.210228in}{0.658836in}}%
\pgfpathlineto{\pgfqpoint{6.213860in}{0.571099in}}%
\pgfpathlineto{\pgfqpoint{6.217491in}{0.651352in}}%
\pgfpathlineto{\pgfqpoint{6.221122in}{0.649564in}}%
\pgfpathlineto{\pgfqpoint{6.224753in}{0.667786in}}%
\pgfpathlineto{\pgfqpoint{6.228384in}{0.658405in}}%
\pgfpathlineto{\pgfqpoint{6.232015in}{0.662065in}}%
\pgfpathlineto{\pgfqpoint{6.235647in}{0.615703in}}%
\pgfpathlineto{\pgfqpoint{6.239278in}{0.641895in}}%
\pgfpathlineto{\pgfqpoint{6.242909in}{0.648960in}}%
\pgfpathlineto{\pgfqpoint{6.246540in}{0.671077in}}%
\pgfpathlineto{\pgfqpoint{6.250171in}{0.656037in}}%
\pgfpathlineto{\pgfqpoint{6.253802in}{0.663384in}}%
\pgfpathlineto{\pgfqpoint{6.257434in}{0.630861in}}%
\pgfpathlineto{\pgfqpoint{6.261065in}{0.632875in}}%
\pgfpathlineto{\pgfqpoint{6.264696in}{0.648637in}}%
\pgfpathlineto{\pgfqpoint{6.268327in}{0.673738in}}%
\pgfpathlineto{\pgfqpoint{6.271958in}{0.653368in}}%
\pgfpathlineto{\pgfqpoint{6.275590in}{0.664643in}}%
\pgfpathlineto{\pgfqpoint{6.282852in}{0.621878in}}%
\pgfpathlineto{\pgfqpoint{6.290114in}{0.676177in}}%
\pgfpathlineto{\pgfqpoint{6.293745in}{0.650384in}}%
\pgfpathlineto{\pgfqpoint{6.297377in}{0.665543in}}%
\pgfpathlineto{\pgfqpoint{6.301008in}{0.650072in}}%
\pgfpathlineto{\pgfqpoint{6.304639in}{0.608495in}}%
\pgfpathlineto{\pgfqpoint{6.308270in}{0.647316in}}%
\pgfpathlineto{\pgfqpoint{6.311901in}{0.678359in}}%
\pgfpathlineto{\pgfqpoint{6.315532in}{0.647077in}}%
\pgfpathlineto{\pgfqpoint{6.319164in}{0.666204in}}%
\pgfpathlineto{\pgfqpoint{6.322795in}{0.656365in}}%
\pgfpathlineto{\pgfqpoint{6.326426in}{0.590544in}}%
\pgfpathlineto{\pgfqpoint{6.330057in}{0.689629in}}%
\pgfpathlineto{\pgfqpoint{6.333688in}{0.718589in}}%
\pgfpathlineto{\pgfqpoint{6.337319in}{0.554550in}}%
\pgfpathlineto{\pgfqpoint{6.340951in}{0.625734in}}%
\pgfpathlineto{\pgfqpoint{6.344582in}{0.654443in}}%
\pgfpathlineto{\pgfqpoint{6.348213in}{0.646165in}}%
\pgfpathlineto{\pgfqpoint{6.351844in}{0.682033in}}%
\pgfpathlineto{\pgfqpoint{6.355475in}{0.688703in}}%
\pgfpathlineto{\pgfqpoint{6.359106in}{0.648205in}}%
\pgfpathlineto{\pgfqpoint{6.362738in}{0.655395in}}%
\pgfpathlineto{\pgfqpoint{6.366369in}{0.714637in}}%
\pgfpathlineto{\pgfqpoint{6.370000in}{0.637450in}}%
\pgfpathlineto{\pgfqpoint{6.370000in}{0.637450in}}%
\pgfusepath{stroke}%
\end{pgfscope}%
\begin{pgfscope}%
\pgfpathrectangle{\pgfqpoint{4.554414in}{0.320625in}}{\pgfqpoint{1.815586in}{1.164375in}}%
\pgfusepath{clip}%
\pgfsetrectcap%
\pgfsetroundjoin%
\pgfsetlinewidth{1.154312pt}%
\definecolor{currentstroke}{rgb}{0.141176,0.345098,0.650980}%
\pgfsetstrokecolor{currentstroke}%
\pgfsetdash{}{0pt}%
\pgfpathmoveto{\pgfqpoint{4.554414in}{1.474926in}}%
\pgfpathlineto{\pgfqpoint{4.561676in}{1.453928in}}%
\pgfpathlineto{\pgfqpoint{4.568938in}{1.438094in}}%
\pgfpathlineto{\pgfqpoint{4.576201in}{1.425824in}}%
\pgfpathlineto{\pgfqpoint{4.587094in}{1.410682in}}%
\pgfpathlineto{\pgfqpoint{4.601619in}{1.392805in}}%
\pgfpathlineto{\pgfqpoint{4.623406in}{1.367939in}}%
\pgfpathlineto{\pgfqpoint{4.645193in}{1.344912in}}%
\pgfpathlineto{\pgfqpoint{4.666980in}{1.323717in}}%
\pgfpathlineto{\pgfqpoint{4.688767in}{1.304149in}}%
\pgfpathlineto{\pgfqpoint{4.714185in}{1.282924in}}%
\pgfpathlineto{\pgfqpoint{4.743235in}{1.260332in}}%
\pgfpathlineto{\pgfqpoint{4.775915in}{1.236425in}}%
\pgfpathlineto{\pgfqpoint{4.815858in}{1.208628in}}%
\pgfpathlineto{\pgfqpoint{4.866694in}{1.174673in}}%
\pgfpathlineto{\pgfqpoint{4.932056in}{1.132452in}}%
\pgfpathlineto{\pgfqpoint{5.015573in}{1.079938in}}%
\pgfpathlineto{\pgfqpoint{5.120877in}{1.015123in}}%
\pgfpathlineto{\pgfqpoint{5.258861in}{0.931587in}}%
\pgfpathlineto{\pgfqpoint{5.440420in}{0.823014in}}%
\pgfpathlineto{\pgfqpoint{5.502150in}{0.786372in}}%
\pgfpathlineto{\pgfqpoint{5.542093in}{0.762978in}}%
\pgfpathlineto{\pgfqpoint{5.549355in}{0.758851in}}%
\pgfpathlineto{\pgfqpoint{5.552986in}{0.757504in}}%
\pgfpathlineto{\pgfqpoint{5.556617in}{0.754216in}}%
\pgfpathlineto{\pgfqpoint{5.563880in}{0.749989in}}%
\pgfpathlineto{\pgfqpoint{5.567511in}{0.748510in}}%
\pgfpathlineto{\pgfqpoint{5.571142in}{0.745537in}}%
\pgfpathlineto{\pgfqpoint{5.574773in}{0.746122in}}%
\pgfpathlineto{\pgfqpoint{5.578404in}{0.744116in}}%
\pgfpathlineto{\pgfqpoint{5.582035in}{0.739464in}}%
\pgfpathlineto{\pgfqpoint{5.592929in}{0.733468in}}%
\pgfpathlineto{\pgfqpoint{5.596560in}{0.733293in}}%
\pgfpathlineto{\pgfqpoint{5.600191in}{0.730821in}}%
\pgfpathlineto{\pgfqpoint{5.603823in}{0.726982in}}%
\pgfpathlineto{\pgfqpoint{5.607454in}{0.729780in}}%
\pgfpathlineto{\pgfqpoint{5.611085in}{0.722035in}}%
\pgfpathlineto{\pgfqpoint{5.618347in}{0.719565in}}%
\pgfpathlineto{\pgfqpoint{5.621978in}{0.718651in}}%
\pgfpathlineto{\pgfqpoint{5.625610in}{0.714727in}}%
\pgfpathlineto{\pgfqpoint{5.629241in}{0.714796in}}%
\pgfpathlineto{\pgfqpoint{5.632872in}{0.711679in}}%
\pgfpathlineto{\pgfqpoint{5.636503in}{0.711959in}}%
\pgfpathlineto{\pgfqpoint{5.643765in}{0.707378in}}%
\pgfpathlineto{\pgfqpoint{5.647397in}{0.712523in}}%
\pgfpathlineto{\pgfqpoint{5.651028in}{0.701243in}}%
\pgfpathlineto{\pgfqpoint{5.654659in}{0.704779in}}%
\pgfpathlineto{\pgfqpoint{5.658290in}{0.698053in}}%
\pgfpathlineto{\pgfqpoint{5.661921in}{0.705697in}}%
\pgfpathlineto{\pgfqpoint{5.665552in}{0.692252in}}%
\pgfpathlineto{\pgfqpoint{5.669184in}{0.694414in}}%
\pgfpathlineto{\pgfqpoint{5.672815in}{0.698891in}}%
\pgfpathlineto{\pgfqpoint{5.676446in}{0.694356in}}%
\pgfpathlineto{\pgfqpoint{5.680077in}{0.685810in}}%
\pgfpathlineto{\pgfqpoint{5.683708in}{0.694945in}}%
\pgfpathlineto{\pgfqpoint{5.687340in}{0.702324in}}%
\pgfpathlineto{\pgfqpoint{5.690971in}{0.692537in}}%
\pgfpathlineto{\pgfqpoint{5.694602in}{0.697030in}}%
\pgfpathlineto{\pgfqpoint{5.698233in}{0.675084in}}%
\pgfpathlineto{\pgfqpoint{5.705495in}{0.685987in}}%
\pgfpathlineto{\pgfqpoint{5.709127in}{0.678622in}}%
\pgfpathlineto{\pgfqpoint{5.712758in}{0.696212in}}%
\pgfpathlineto{\pgfqpoint{5.716389in}{0.708437in}}%
\pgfpathlineto{\pgfqpoint{5.720020in}{0.679418in}}%
\pgfpathlineto{\pgfqpoint{5.723651in}{0.660342in}}%
\pgfpathlineto{\pgfqpoint{5.727282in}{0.687399in}}%
\pgfpathlineto{\pgfqpoint{5.730914in}{0.671038in}}%
\pgfpathlineto{\pgfqpoint{5.734545in}{0.687158in}}%
\pgfpathlineto{\pgfqpoint{5.738176in}{0.683341in}}%
\pgfpathlineto{\pgfqpoint{5.741807in}{0.696908in}}%
\pgfpathlineto{\pgfqpoint{5.745438in}{0.670194in}}%
\pgfpathlineto{\pgfqpoint{5.749069in}{0.682612in}}%
\pgfpathlineto{\pgfqpoint{5.752701in}{0.672882in}}%
\pgfpathlineto{\pgfqpoint{5.756332in}{0.677552in}}%
\pgfpathlineto{\pgfqpoint{5.759963in}{0.671738in}}%
\pgfpathlineto{\pgfqpoint{5.763594in}{0.675812in}}%
\pgfpathlineto{\pgfqpoint{5.767225in}{0.660191in}}%
\pgfpathlineto{\pgfqpoint{5.774488in}{0.682223in}}%
\pgfpathlineto{\pgfqpoint{5.778119in}{0.680058in}}%
\pgfpathlineto{\pgfqpoint{5.781750in}{0.695280in}}%
\pgfpathlineto{\pgfqpoint{5.785381in}{0.678997in}}%
\pgfpathlineto{\pgfqpoint{5.789012in}{0.656642in}}%
\pgfpathlineto{\pgfqpoint{5.792644in}{0.686183in}}%
\pgfpathlineto{\pgfqpoint{5.796275in}{0.670222in}}%
\pgfpathlineto{\pgfqpoint{5.799906in}{0.670981in}}%
\pgfpathlineto{\pgfqpoint{5.803537in}{0.675767in}}%
\pgfpathlineto{\pgfqpoint{5.807168in}{0.675873in}}%
\pgfpathlineto{\pgfqpoint{5.810799in}{0.683824in}}%
\pgfpathlineto{\pgfqpoint{5.814431in}{0.682466in}}%
\pgfpathlineto{\pgfqpoint{5.818062in}{0.678548in}}%
\pgfpathlineto{\pgfqpoint{5.821693in}{0.666304in}}%
\pgfpathlineto{\pgfqpoint{5.825324in}{0.684996in}}%
\pgfpathlineto{\pgfqpoint{5.828955in}{0.666267in}}%
\pgfpathlineto{\pgfqpoint{5.836218in}{0.683767in}}%
\pgfpathlineto{\pgfqpoint{5.839849in}{0.676384in}}%
\pgfpathlineto{\pgfqpoint{5.843480in}{0.672072in}}%
\pgfpathlineto{\pgfqpoint{5.847111in}{0.669770in}}%
\pgfpathlineto{\pgfqpoint{5.850742in}{0.665936in}}%
\pgfpathlineto{\pgfqpoint{5.858005in}{0.683234in}}%
\pgfpathlineto{\pgfqpoint{5.861636in}{0.672697in}}%
\pgfpathlineto{\pgfqpoint{5.865267in}{0.700682in}}%
\pgfpathlineto{\pgfqpoint{5.868898in}{0.684106in}}%
\pgfpathlineto{\pgfqpoint{5.872529in}{0.674547in}}%
\pgfpathlineto{\pgfqpoint{5.876160in}{0.671073in}}%
\pgfpathlineto{\pgfqpoint{5.879792in}{0.676409in}}%
\pgfpathlineto{\pgfqpoint{5.883423in}{0.667975in}}%
\pgfpathlineto{\pgfqpoint{5.887054in}{0.683486in}}%
\pgfpathlineto{\pgfqpoint{5.890685in}{0.706095in}}%
\pgfpathlineto{\pgfqpoint{5.897948in}{0.654877in}}%
\pgfpathlineto{\pgfqpoint{5.901579in}{0.670424in}}%
\pgfpathlineto{\pgfqpoint{5.905210in}{0.677225in}}%
\pgfpathlineto{\pgfqpoint{5.908841in}{0.666789in}}%
\pgfpathlineto{\pgfqpoint{5.912472in}{0.690714in}}%
\pgfpathlineto{\pgfqpoint{5.916103in}{0.665344in}}%
\pgfpathlineto{\pgfqpoint{5.919735in}{0.675216in}}%
\pgfpathlineto{\pgfqpoint{5.923366in}{0.677860in}}%
\pgfpathlineto{\pgfqpoint{5.926997in}{0.657930in}}%
\pgfpathlineto{\pgfqpoint{5.930628in}{0.690469in}}%
\pgfpathlineto{\pgfqpoint{5.934259in}{0.691584in}}%
\pgfpathlineto{\pgfqpoint{5.937890in}{0.644067in}}%
\pgfpathlineto{\pgfqpoint{5.941522in}{0.664857in}}%
\pgfpathlineto{\pgfqpoint{5.945153in}{0.676100in}}%
\pgfpathlineto{\pgfqpoint{5.948784in}{0.682547in}}%
\pgfpathlineto{\pgfqpoint{5.952415in}{0.675508in}}%
\pgfpathlineto{\pgfqpoint{5.956046in}{0.660694in}}%
\pgfpathlineto{\pgfqpoint{5.959677in}{0.664334in}}%
\pgfpathlineto{\pgfqpoint{5.963309in}{0.678792in}}%
\pgfpathlineto{\pgfqpoint{5.966940in}{0.678730in}}%
\pgfpathlineto{\pgfqpoint{5.970571in}{0.690323in}}%
\pgfpathlineto{\pgfqpoint{5.974202in}{0.695646in}}%
\pgfpathlineto{\pgfqpoint{5.977833in}{0.661469in}}%
\pgfpathlineto{\pgfqpoint{5.981465in}{0.666135in}}%
\pgfpathlineto{\pgfqpoint{5.985096in}{0.667583in}}%
\pgfpathlineto{\pgfqpoint{5.988727in}{0.687280in}}%
\pgfpathlineto{\pgfqpoint{5.992358in}{0.660015in}}%
\pgfpathlineto{\pgfqpoint{5.999620in}{0.675885in}}%
\pgfpathlineto{\pgfqpoint{6.003252in}{0.678136in}}%
\pgfpathlineto{\pgfqpoint{6.006883in}{0.687473in}}%
\pgfpathlineto{\pgfqpoint{6.010514in}{0.691034in}}%
\pgfpathlineto{\pgfqpoint{6.014145in}{0.688655in}}%
\pgfpathlineto{\pgfqpoint{6.017776in}{0.657351in}}%
\pgfpathlineto{\pgfqpoint{6.021407in}{0.680236in}}%
\pgfpathlineto{\pgfqpoint{6.025039in}{0.677459in}}%
\pgfpathlineto{\pgfqpoint{6.028670in}{0.694456in}}%
\pgfpathlineto{\pgfqpoint{6.032301in}{0.673947in}}%
\pgfpathlineto{\pgfqpoint{6.035932in}{0.682129in}}%
\pgfpathlineto{\pgfqpoint{6.039563in}{0.648812in}}%
\pgfpathlineto{\pgfqpoint{6.043194in}{0.696076in}}%
\pgfpathlineto{\pgfqpoint{6.046826in}{0.673221in}}%
\pgfpathlineto{\pgfqpoint{6.050457in}{0.692409in}}%
\pgfpathlineto{\pgfqpoint{6.054088in}{0.661872in}}%
\pgfpathlineto{\pgfqpoint{6.057719in}{0.655184in}}%
\pgfpathlineto{\pgfqpoint{6.061350in}{0.666431in}}%
\pgfpathlineto{\pgfqpoint{6.064981in}{0.687487in}}%
\pgfpathlineto{\pgfqpoint{6.068613in}{0.704060in}}%
\pgfpathlineto{\pgfqpoint{6.072244in}{0.679427in}}%
\pgfpathlineto{\pgfqpoint{6.075875in}{0.676561in}}%
\pgfpathlineto{\pgfqpoint{6.079506in}{0.662668in}}%
\pgfpathlineto{\pgfqpoint{6.083137in}{0.666677in}}%
\pgfpathlineto{\pgfqpoint{6.086769in}{0.676768in}}%
\pgfpathlineto{\pgfqpoint{6.090400in}{0.678016in}}%
\pgfpathlineto{\pgfqpoint{6.094031in}{0.664633in}}%
\pgfpathlineto{\pgfqpoint{6.097662in}{0.666373in}}%
\pgfpathlineto{\pgfqpoint{6.101293in}{0.663907in}}%
\pgfpathlineto{\pgfqpoint{6.104924in}{0.684151in}}%
\pgfpathlineto{\pgfqpoint{6.108556in}{0.692396in}}%
\pgfpathlineto{\pgfqpoint{6.112187in}{0.684526in}}%
\pgfpathlineto{\pgfqpoint{6.115818in}{0.666339in}}%
\pgfpathlineto{\pgfqpoint{6.119449in}{0.680655in}}%
\pgfpathlineto{\pgfqpoint{6.123080in}{0.650274in}}%
\pgfpathlineto{\pgfqpoint{6.126711in}{0.668921in}}%
\pgfpathlineto{\pgfqpoint{6.130343in}{0.675152in}}%
\pgfpathlineto{\pgfqpoint{6.133974in}{0.667696in}}%
\pgfpathlineto{\pgfqpoint{6.137605in}{0.677660in}}%
\pgfpathlineto{\pgfqpoint{6.141236in}{0.645705in}}%
\pgfpathlineto{\pgfqpoint{6.144867in}{0.693798in}}%
\pgfpathlineto{\pgfqpoint{6.148498in}{0.676617in}}%
\pgfpathlineto{\pgfqpoint{6.152130in}{0.678088in}}%
\pgfpathlineto{\pgfqpoint{6.155761in}{0.659561in}}%
\pgfpathlineto{\pgfqpoint{6.159392in}{0.674366in}}%
\pgfpathlineto{\pgfqpoint{6.163023in}{0.660789in}}%
\pgfpathlineto{\pgfqpoint{6.166654in}{0.696115in}}%
\pgfpathlineto{\pgfqpoint{6.170285in}{0.693637in}}%
\pgfpathlineto{\pgfqpoint{6.173917in}{0.647436in}}%
\pgfpathlineto{\pgfqpoint{6.177548in}{0.654016in}}%
\pgfpathlineto{\pgfqpoint{6.184810in}{0.690726in}}%
\pgfpathlineto{\pgfqpoint{6.188441in}{0.679123in}}%
\pgfpathlineto{\pgfqpoint{6.192073in}{0.698718in}}%
\pgfpathlineto{\pgfqpoint{6.195704in}{0.628961in}}%
\pgfpathlineto{\pgfqpoint{6.199335in}{0.654961in}}%
\pgfpathlineto{\pgfqpoint{6.202966in}{0.652673in}}%
\pgfpathlineto{\pgfqpoint{6.206597in}{0.691579in}}%
\pgfpathlineto{\pgfqpoint{6.210228in}{0.677765in}}%
\pgfpathlineto{\pgfqpoint{6.213860in}{0.687616in}}%
\pgfpathlineto{\pgfqpoint{6.217491in}{0.694931in}}%
\pgfpathlineto{\pgfqpoint{6.221122in}{0.675709in}}%
\pgfpathlineto{\pgfqpoint{6.224753in}{0.677175in}}%
\pgfpathlineto{\pgfqpoint{6.228384in}{0.668739in}}%
\pgfpathlineto{\pgfqpoint{6.232015in}{0.650089in}}%
\pgfpathlineto{\pgfqpoint{6.235647in}{0.655448in}}%
\pgfpathlineto{\pgfqpoint{6.239278in}{0.689815in}}%
\pgfpathlineto{\pgfqpoint{6.242909in}{0.668211in}}%
\pgfpathlineto{\pgfqpoint{6.246540in}{0.686961in}}%
\pgfpathlineto{\pgfqpoint{6.250171in}{0.682245in}}%
\pgfpathlineto{\pgfqpoint{6.253802in}{0.661897in}}%
\pgfpathlineto{\pgfqpoint{6.257434in}{0.668073in}}%
\pgfpathlineto{\pgfqpoint{6.261065in}{0.684842in}}%
\pgfpathlineto{\pgfqpoint{6.264696in}{0.655882in}}%
\pgfpathlineto{\pgfqpoint{6.268327in}{0.672226in}}%
\pgfpathlineto{\pgfqpoint{6.271958in}{0.656809in}}%
\pgfpathlineto{\pgfqpoint{6.275590in}{0.677883in}}%
\pgfpathlineto{\pgfqpoint{6.279221in}{0.667139in}}%
\pgfpathlineto{\pgfqpoint{6.282852in}{0.630485in}}%
\pgfpathlineto{\pgfqpoint{6.286483in}{0.679061in}}%
\pgfpathlineto{\pgfqpoint{6.290114in}{0.671091in}}%
\pgfpathlineto{\pgfqpoint{6.293745in}{0.668968in}}%
\pgfpathlineto{\pgfqpoint{6.301008in}{0.691905in}}%
\pgfpathlineto{\pgfqpoint{6.304639in}{0.668383in}}%
\pgfpathlineto{\pgfqpoint{6.308270in}{0.661182in}}%
\pgfpathlineto{\pgfqpoint{6.311901in}{0.661432in}}%
\pgfpathlineto{\pgfqpoint{6.319164in}{0.688843in}}%
\pgfpathlineto{\pgfqpoint{6.322795in}{0.681772in}}%
\pgfpathlineto{\pgfqpoint{6.326426in}{0.682258in}}%
\pgfpathlineto{\pgfqpoint{6.330057in}{0.675550in}}%
\pgfpathlineto{\pgfqpoint{6.333688in}{0.692509in}}%
\pgfpathlineto{\pgfqpoint{6.337319in}{0.668437in}}%
\pgfpathlineto{\pgfqpoint{6.340951in}{0.693635in}}%
\pgfpathlineto{\pgfqpoint{6.344582in}{0.690114in}}%
\pgfpathlineto{\pgfqpoint{6.348213in}{0.672973in}}%
\pgfpathlineto{\pgfqpoint{6.351844in}{0.666145in}}%
\pgfpathlineto{\pgfqpoint{6.355475in}{0.677080in}}%
\pgfpathlineto{\pgfqpoint{6.359106in}{0.669880in}}%
\pgfpathlineto{\pgfqpoint{6.366369in}{0.689707in}}%
\pgfpathlineto{\pgfqpoint{6.370000in}{0.657585in}}%
\pgfpathlineto{\pgfqpoint{6.370000in}{0.657585in}}%
\pgfusepath{stroke}%
\end{pgfscope}%
\begin{pgfscope}%
\pgfpathrectangle{\pgfqpoint{4.554414in}{0.320625in}}{\pgfqpoint{1.815586in}{1.164375in}}%
\pgfusepath{clip}%
\pgfsetrectcap%
\pgfsetroundjoin%
\pgfsetlinewidth{0.602250pt}%
\definecolor{currentstroke}{rgb}{0.819608,0.286275,0.356863}%
\pgfsetstrokecolor{currentstroke}%
\pgfsetstrokeopacity{0.120000}%
\pgfsetdash{}{0pt}%
\pgfpathmoveto{\pgfqpoint{4.554414in}{1.463724in}}%
\pgfpathlineto{\pgfqpoint{4.558045in}{1.450282in}}%
\pgfpathlineto{\pgfqpoint{4.561676in}{1.440454in}}%
\pgfpathlineto{\pgfqpoint{4.576201in}{1.410366in}}%
\pgfpathlineto{\pgfqpoint{4.619775in}{1.324494in}}%
\pgfpathlineto{\pgfqpoint{4.641562in}{1.284767in}}%
\pgfpathlineto{\pgfqpoint{4.663349in}{1.247396in}}%
\pgfpathlineto{\pgfqpoint{4.688767in}{1.206089in}}%
\pgfpathlineto{\pgfqpoint{4.717816in}{1.160978in}}%
\pgfpathlineto{\pgfqpoint{4.757759in}{1.101207in}}%
\pgfpathlineto{\pgfqpoint{4.812227in}{1.022009in}}%
\pgfpathlineto{\pgfqpoint{4.888481in}{0.913290in}}%
\pgfpathlineto{\pgfqpoint{4.986523in}{0.775302in}}%
\pgfpathlineto{\pgfqpoint{4.997417in}{0.760659in}}%
\pgfpathlineto{\pgfqpoint{5.001048in}{0.754849in}}%
\pgfpathlineto{\pgfqpoint{5.011941in}{0.740211in}}%
\pgfpathlineto{\pgfqpoint{5.015573in}{0.738229in}}%
\pgfpathlineto{\pgfqpoint{5.019204in}{0.731616in}}%
\pgfpathlineto{\pgfqpoint{5.022835in}{0.729928in}}%
\pgfpathlineto{\pgfqpoint{5.026466in}{0.719320in}}%
\pgfpathlineto{\pgfqpoint{5.030097in}{0.715198in}}%
\pgfpathlineto{\pgfqpoint{5.033728in}{0.709110in}}%
\pgfpathlineto{\pgfqpoint{5.037360in}{0.709801in}}%
\pgfpathlineto{\pgfqpoint{5.040991in}{0.714372in}}%
\pgfpathlineto{\pgfqpoint{5.044622in}{0.698305in}}%
\pgfpathlineto{\pgfqpoint{5.048253in}{0.716263in}}%
\pgfpathlineto{\pgfqpoint{5.051884in}{0.684018in}}%
\pgfpathlineto{\pgfqpoint{5.055515in}{0.681390in}}%
\pgfpathlineto{\pgfqpoint{5.059147in}{0.673986in}}%
\pgfpathlineto{\pgfqpoint{5.062778in}{0.689118in}}%
\pgfpathlineto{\pgfqpoint{5.066409in}{0.692031in}}%
\pgfpathlineto{\pgfqpoint{5.070040in}{0.663501in}}%
\pgfpathlineto{\pgfqpoint{5.073671in}{0.690383in}}%
\pgfpathlineto{\pgfqpoint{5.077302in}{0.661725in}}%
\pgfpathlineto{\pgfqpoint{5.080934in}{0.665950in}}%
\pgfpathlineto{\pgfqpoint{5.084565in}{0.650611in}}%
\pgfpathlineto{\pgfqpoint{5.088196in}{0.696461in}}%
\pgfpathlineto{\pgfqpoint{5.091827in}{0.649395in}}%
\pgfpathlineto{\pgfqpoint{5.095458in}{0.634121in}}%
\pgfpathlineto{\pgfqpoint{5.099090in}{0.645264in}}%
\pgfpathlineto{\pgfqpoint{5.102721in}{0.660816in}}%
\pgfpathlineto{\pgfqpoint{5.106352in}{0.694156in}}%
\pgfpathlineto{\pgfqpoint{5.109983in}{0.678561in}}%
\pgfpathlineto{\pgfqpoint{5.113614in}{0.688085in}}%
\pgfpathlineto{\pgfqpoint{5.117245in}{0.594885in}}%
\pgfpathlineto{\pgfqpoint{5.120877in}{0.663214in}}%
\pgfpathlineto{\pgfqpoint{5.124508in}{0.648111in}}%
\pgfpathlineto{\pgfqpoint{5.128139in}{0.658806in}}%
\pgfpathlineto{\pgfqpoint{5.131770in}{0.658037in}}%
\pgfpathlineto{\pgfqpoint{5.135401in}{0.599256in}}%
\pgfpathlineto{\pgfqpoint{5.139032in}{0.670563in}}%
\pgfpathlineto{\pgfqpoint{5.142664in}{0.590540in}}%
\pgfpathlineto{\pgfqpoint{5.146295in}{0.691214in}}%
\pgfpathlineto{\pgfqpoint{5.149926in}{0.699967in}}%
\pgfpathlineto{\pgfqpoint{5.153557in}{0.668426in}}%
\pgfpathlineto{\pgfqpoint{5.157188in}{0.628664in}}%
\pgfpathlineto{\pgfqpoint{5.160819in}{0.647887in}}%
\pgfpathlineto{\pgfqpoint{5.164451in}{0.679926in}}%
\pgfpathlineto{\pgfqpoint{5.168082in}{0.617618in}}%
\pgfpathlineto{\pgfqpoint{5.171713in}{0.691087in}}%
\pgfpathlineto{\pgfqpoint{5.175344in}{0.712414in}}%
\pgfpathlineto{\pgfqpoint{5.178975in}{0.674920in}}%
\pgfpathlineto{\pgfqpoint{5.182606in}{0.621343in}}%
\pgfpathlineto{\pgfqpoint{5.186238in}{0.636418in}}%
\pgfpathlineto{\pgfqpoint{5.189869in}{0.657049in}}%
\pgfpathlineto{\pgfqpoint{5.193500in}{0.627550in}}%
\pgfpathlineto{\pgfqpoint{5.197131in}{0.702226in}}%
\pgfpathlineto{\pgfqpoint{5.200762in}{0.707683in}}%
\pgfpathlineto{\pgfqpoint{5.204394in}{0.682830in}}%
\pgfpathlineto{\pgfqpoint{5.208025in}{0.631518in}}%
\pgfpathlineto{\pgfqpoint{5.211656in}{0.614378in}}%
\pgfpathlineto{\pgfqpoint{5.215287in}{0.637311in}}%
\pgfpathlineto{\pgfqpoint{5.218918in}{0.655610in}}%
\pgfpathlineto{\pgfqpoint{5.222549in}{0.697596in}}%
\pgfpathlineto{\pgfqpoint{5.226181in}{0.688307in}}%
\pgfpathlineto{\pgfqpoint{5.229812in}{0.686377in}}%
\pgfpathlineto{\pgfqpoint{5.233443in}{0.546079in}}%
\pgfpathlineto{\pgfqpoint{5.237074in}{0.660441in}}%
\pgfpathlineto{\pgfqpoint{5.240705in}{0.644382in}}%
\pgfpathlineto{\pgfqpoint{5.244336in}{0.650542in}}%
\pgfpathlineto{\pgfqpoint{5.247968in}{0.654714in}}%
\pgfpathlineto{\pgfqpoint{5.251599in}{0.610453in}}%
\pgfpathlineto{\pgfqpoint{5.255230in}{0.668197in}}%
\pgfpathlineto{\pgfqpoint{5.258861in}{0.570959in}}%
\pgfpathlineto{\pgfqpoint{5.262492in}{0.692570in}}%
\pgfpathlineto{\pgfqpoint{5.266123in}{0.704624in}}%
\pgfpathlineto{\pgfqpoint{5.273386in}{0.626322in}}%
\pgfpathlineto{\pgfqpoint{5.277017in}{0.645390in}}%
\pgfpathlineto{\pgfqpoint{5.280648in}{0.679218in}}%
\pgfpathlineto{\pgfqpoint{5.284279in}{0.637441in}}%
\pgfpathlineto{\pgfqpoint{5.287910in}{0.637493in}}%
\pgfpathlineto{\pgfqpoint{5.291542in}{0.664625in}}%
\pgfpathlineto{\pgfqpoint{5.295173in}{0.668900in}}%
\pgfpathlineto{\pgfqpoint{5.298804in}{0.638654in}}%
\pgfpathlineto{\pgfqpoint{5.306066in}{0.707070in}}%
\pgfpathlineto{\pgfqpoint{5.309698in}{0.571107in}}%
\pgfpathlineto{\pgfqpoint{5.313329in}{0.641490in}}%
\pgfpathlineto{\pgfqpoint{5.316960in}{0.526579in}}%
\pgfpathlineto{\pgfqpoint{5.320591in}{0.679019in}}%
\pgfpathlineto{\pgfqpoint{5.324222in}{0.679221in}}%
\pgfpathlineto{\pgfqpoint{5.327853in}{0.586017in}}%
\pgfpathlineto{\pgfqpoint{5.331485in}{0.686896in}}%
\pgfpathlineto{\pgfqpoint{5.335116in}{0.645243in}}%
\pgfpathlineto{\pgfqpoint{5.338747in}{0.662026in}}%
\pgfpathlineto{\pgfqpoint{5.342378in}{0.609381in}}%
\pgfpathlineto{\pgfqpoint{5.346009in}{0.694632in}}%
\pgfpathlineto{\pgfqpoint{5.349640in}{0.693689in}}%
\pgfpathlineto{\pgfqpoint{5.353272in}{0.614384in}}%
\pgfpathlineto{\pgfqpoint{5.356903in}{0.699592in}}%
\pgfpathlineto{\pgfqpoint{5.360534in}{0.653769in}}%
\pgfpathlineto{\pgfqpoint{5.364165in}{0.643180in}}%
\pgfpathlineto{\pgfqpoint{5.367796in}{0.648232in}}%
\pgfpathlineto{\pgfqpoint{5.371427in}{0.693653in}}%
\pgfpathlineto{\pgfqpoint{5.375059in}{0.686429in}}%
\pgfpathlineto{\pgfqpoint{5.378690in}{0.626478in}}%
\pgfpathlineto{\pgfqpoint{5.382321in}{0.712301in}}%
\pgfpathlineto{\pgfqpoint{5.385952in}{0.660345in}}%
\pgfpathlineto{\pgfqpoint{5.389583in}{0.625419in}}%
\pgfpathlineto{\pgfqpoint{5.393215in}{0.647261in}}%
\pgfpathlineto{\pgfqpoint{5.396846in}{0.681505in}}%
\pgfpathlineto{\pgfqpoint{5.400477in}{0.654222in}}%
\pgfpathlineto{\pgfqpoint{5.404108in}{0.617479in}}%
\pgfpathlineto{\pgfqpoint{5.407739in}{0.667693in}}%
\pgfpathlineto{\pgfqpoint{5.411370in}{0.665860in}}%
\pgfpathlineto{\pgfqpoint{5.415002in}{0.635992in}}%
\pgfpathlineto{\pgfqpoint{5.418633in}{0.676320in}}%
\pgfpathlineto{\pgfqpoint{5.422264in}{0.707735in}}%
\pgfpathlineto{\pgfqpoint{5.425895in}{0.519346in}}%
\pgfpathlineto{\pgfqpoint{5.429526in}{0.636349in}}%
\pgfpathlineto{\pgfqpoint{5.433157in}{0.530121in}}%
\pgfpathlineto{\pgfqpoint{5.436789in}{0.676950in}}%
\pgfpathlineto{\pgfqpoint{5.440420in}{0.683733in}}%
\pgfpathlineto{\pgfqpoint{5.444051in}{0.624737in}}%
\pgfpathlineto{\pgfqpoint{5.447682in}{0.687844in}}%
\pgfpathlineto{\pgfqpoint{5.451313in}{0.641189in}}%
\pgfpathlineto{\pgfqpoint{5.454944in}{0.661116in}}%
\pgfpathlineto{\pgfqpoint{5.458576in}{0.605980in}}%
\pgfpathlineto{\pgfqpoint{5.462207in}{0.692530in}}%
\pgfpathlineto{\pgfqpoint{5.465838in}{0.644890in}}%
\pgfpathlineto{\pgfqpoint{5.469469in}{0.614384in}}%
\pgfpathlineto{\pgfqpoint{5.473100in}{0.648289in}}%
\pgfpathlineto{\pgfqpoint{5.476731in}{0.654879in}}%
\pgfpathlineto{\pgfqpoint{5.480363in}{0.692586in}}%
\pgfpathlineto{\pgfqpoint{5.483994in}{0.677563in}}%
\pgfpathlineto{\pgfqpoint{5.487625in}{0.687217in}}%
\pgfpathlineto{\pgfqpoint{5.491256in}{0.503038in}}%
\pgfpathlineto{\pgfqpoint{5.494887in}{0.662661in}}%
\pgfpathlineto{\pgfqpoint{5.498519in}{0.655913in}}%
\pgfpathlineto{\pgfqpoint{5.502150in}{0.653434in}}%
\pgfpathlineto{\pgfqpoint{5.505781in}{0.700254in}}%
\pgfpathlineto{\pgfqpoint{5.509412in}{0.691854in}}%
\pgfpathlineto{\pgfqpoint{5.513043in}{0.696806in}}%
\pgfpathlineto{\pgfqpoint{5.516674in}{0.552191in}}%
\pgfpathlineto{\pgfqpoint{5.520306in}{0.649613in}}%
\pgfpathlineto{\pgfqpoint{5.523937in}{0.588779in}}%
\pgfpathlineto{\pgfqpoint{5.527568in}{0.674238in}}%
\pgfpathlineto{\pgfqpoint{5.531199in}{0.702688in}}%
\pgfpathlineto{\pgfqpoint{5.534830in}{0.678039in}}%
\pgfpathlineto{\pgfqpoint{5.538461in}{0.706891in}}%
\pgfpathlineto{\pgfqpoint{5.542093in}{0.521432in}}%
\pgfpathlineto{\pgfqpoint{5.545724in}{0.635151in}}%
\pgfpathlineto{\pgfqpoint{5.549355in}{0.540591in}}%
\pgfpathlineto{\pgfqpoint{5.552986in}{0.676419in}}%
\pgfpathlineto{\pgfqpoint{5.556617in}{0.689920in}}%
\pgfpathlineto{\pgfqpoint{5.560248in}{0.650253in}}%
\pgfpathlineto{\pgfqpoint{5.563880in}{0.688736in}}%
\pgfpathlineto{\pgfqpoint{5.567511in}{0.637941in}}%
\pgfpathlineto{\pgfqpoint{5.571142in}{0.658910in}}%
\pgfpathlineto{\pgfqpoint{5.574773in}{0.611426in}}%
\pgfpathlineto{\pgfqpoint{5.578404in}{0.690392in}}%
\pgfpathlineto{\pgfqpoint{5.585667in}{0.599972in}}%
\pgfpathlineto{\pgfqpoint{5.589298in}{0.646565in}}%
\pgfpathlineto{\pgfqpoint{5.592929in}{0.651090in}}%
\pgfpathlineto{\pgfqpoint{5.596560in}{0.694996in}}%
\pgfpathlineto{\pgfqpoint{5.600191in}{0.685781in}}%
\pgfpathlineto{\pgfqpoint{5.603823in}{0.685116in}}%
\pgfpathlineto{\pgfqpoint{5.607454in}{0.532803in}}%
\pgfpathlineto{\pgfqpoint{5.611085in}{0.661026in}}%
\pgfpathlineto{\pgfqpoint{5.614716in}{0.657371in}}%
\pgfpathlineto{\pgfqpoint{5.618347in}{0.641402in}}%
\pgfpathlineto{\pgfqpoint{5.621978in}{0.656706in}}%
\pgfpathlineto{\pgfqpoint{5.625610in}{0.619090in}}%
\pgfpathlineto{\pgfqpoint{5.629241in}{0.672031in}}%
\pgfpathlineto{\pgfqpoint{5.632872in}{0.451796in}}%
\pgfpathlineto{\pgfqpoint{5.636503in}{0.689754in}}%
\pgfpathlineto{\pgfqpoint{5.640134in}{0.703858in}}%
\pgfpathlineto{\pgfqpoint{5.643765in}{0.660089in}}%
\pgfpathlineto{\pgfqpoint{5.647397in}{0.633570in}}%
\pgfpathlineto{\pgfqpoint{5.651028in}{0.640836in}}%
\pgfpathlineto{\pgfqpoint{5.654659in}{0.681577in}}%
\pgfpathlineto{\pgfqpoint{5.658290in}{0.642035in}}%
\pgfpathlineto{\pgfqpoint{5.661921in}{0.638718in}}%
\pgfpathlineto{\pgfqpoint{5.665552in}{0.667610in}}%
\pgfpathlineto{\pgfqpoint{5.669184in}{0.668771in}}%
\pgfpathlineto{\pgfqpoint{5.672815in}{0.644512in}}%
\pgfpathlineto{\pgfqpoint{5.676446in}{0.669239in}}%
\pgfpathlineto{\pgfqpoint{5.680077in}{0.705752in}}%
\pgfpathlineto{\pgfqpoint{5.683708in}{0.668859in}}%
\pgfpathlineto{\pgfqpoint{5.687340in}{0.645899in}}%
\pgfpathlineto{\pgfqpoint{5.690971in}{0.682998in}}%
\pgfpathlineto{\pgfqpoint{5.694602in}{0.674125in}}%
\pgfpathlineto{\pgfqpoint{5.698233in}{0.607017in}}%
\pgfpathlineto{\pgfqpoint{5.701864in}{0.672489in}}%
\pgfpathlineto{\pgfqpoint{5.705495in}{0.693987in}}%
\pgfpathlineto{\pgfqpoint{5.709127in}{0.639193in}}%
\pgfpathlineto{\pgfqpoint{5.712758in}{0.678738in}}%
\pgfpathlineto{\pgfqpoint{5.716389in}{0.698089in}}%
\pgfpathlineto{\pgfqpoint{5.720020in}{0.682368in}}%
\pgfpathlineto{\pgfqpoint{5.723651in}{0.538666in}}%
\pgfpathlineto{\pgfqpoint{5.727282in}{0.662856in}}%
\pgfpathlineto{\pgfqpoint{5.730914in}{0.663802in}}%
\pgfpathlineto{\pgfqpoint{5.734545in}{0.621605in}}%
\pgfpathlineto{\pgfqpoint{5.738176in}{0.651305in}}%
\pgfpathlineto{\pgfqpoint{5.741807in}{0.629444in}}%
\pgfpathlineto{\pgfqpoint{5.745438in}{0.668635in}}%
\pgfpathlineto{\pgfqpoint{5.749069in}{0.465503in}}%
\pgfpathlineto{\pgfqpoint{5.752701in}{0.691530in}}%
\pgfpathlineto{\pgfqpoint{5.756332in}{0.707199in}}%
\pgfpathlineto{\pgfqpoint{5.759963in}{0.655187in}}%
\pgfpathlineto{\pgfqpoint{5.763594in}{0.630024in}}%
\pgfpathlineto{\pgfqpoint{5.767225in}{0.639452in}}%
\pgfpathlineto{\pgfqpoint{5.770856in}{0.680362in}}%
\pgfpathlineto{\pgfqpoint{5.778119in}{0.623487in}}%
\pgfpathlineto{\pgfqpoint{5.781750in}{0.670429in}}%
\pgfpathlineto{\pgfqpoint{5.785381in}{0.665663in}}%
\pgfpathlineto{\pgfqpoint{5.789012in}{0.644778in}}%
\pgfpathlineto{\pgfqpoint{5.792644in}{0.669288in}}%
\pgfpathlineto{\pgfqpoint{5.796275in}{0.706081in}}%
\pgfpathlineto{\pgfqpoint{5.799906in}{0.542558in}}%
\pgfpathlineto{\pgfqpoint{5.803537in}{0.639223in}}%
\pgfpathlineto{\pgfqpoint{5.807168in}{0.559014in}}%
\pgfpathlineto{\pgfqpoint{5.810799in}{0.676271in}}%
\pgfpathlineto{\pgfqpoint{5.814431in}{0.683696in}}%
\pgfpathlineto{\pgfqpoint{5.818062in}{0.624083in}}%
\pgfpathlineto{\pgfqpoint{5.821693in}{0.689189in}}%
\pgfpathlineto{\pgfqpoint{5.825324in}{0.636141in}}%
\pgfpathlineto{\pgfqpoint{5.828955in}{0.663239in}}%
\pgfpathlineto{\pgfqpoint{5.832586in}{0.557189in}}%
\pgfpathlineto{\pgfqpoint{5.836218in}{0.690113in}}%
\pgfpathlineto{\pgfqpoint{5.839849in}{0.697204in}}%
\pgfpathlineto{\pgfqpoint{5.843480in}{0.644907in}}%
\pgfpathlineto{\pgfqpoint{5.847111in}{0.701466in}}%
\pgfpathlineto{\pgfqpoint{5.850742in}{0.645748in}}%
\pgfpathlineto{\pgfqpoint{5.854373in}{0.645670in}}%
\pgfpathlineto{\pgfqpoint{5.858005in}{0.638955in}}%
\pgfpathlineto{\pgfqpoint{5.861636in}{0.691634in}}%
\pgfpathlineto{\pgfqpoint{5.865267in}{0.692226in}}%
\pgfpathlineto{\pgfqpoint{5.868898in}{0.550127in}}%
\pgfpathlineto{\pgfqpoint{5.872529in}{0.713697in}}%
\pgfpathlineto{\pgfqpoint{5.876160in}{0.651400in}}%
\pgfpathlineto{\pgfqpoint{5.879792in}{0.629251in}}%
\pgfpathlineto{\pgfqpoint{5.883423in}{0.640796in}}%
\pgfpathlineto{\pgfqpoint{5.887054in}{0.681963in}}%
\pgfpathlineto{\pgfqpoint{5.890685in}{0.666107in}}%
\pgfpathlineto{\pgfqpoint{5.894316in}{0.583409in}}%
\pgfpathlineto{\pgfqpoint{5.897948in}{0.673246in}}%
\pgfpathlineto{\pgfqpoint{5.901579in}{0.662472in}}%
\pgfpathlineto{\pgfqpoint{5.905210in}{0.642745in}}%
\pgfpathlineto{\pgfqpoint{5.908841in}{0.670867in}}%
\pgfpathlineto{\pgfqpoint{5.912472in}{0.706135in}}%
\pgfpathlineto{\pgfqpoint{5.916103in}{0.579679in}}%
\pgfpathlineto{\pgfqpoint{5.919735in}{0.633363in}}%
\pgfpathlineto{\pgfqpoint{5.923366in}{0.560167in}}%
\pgfpathlineto{\pgfqpoint{5.926997in}{0.673836in}}%
\pgfpathlineto{\pgfqpoint{5.930628in}{0.687791in}}%
\pgfpathlineto{\pgfqpoint{5.934259in}{0.645983in}}%
\pgfpathlineto{\pgfqpoint{5.937890in}{0.689544in}}%
\pgfpathlineto{\pgfqpoint{5.941522in}{0.630913in}}%
\pgfpathlineto{\pgfqpoint{5.945153in}{0.662305in}}%
\pgfpathlineto{\pgfqpoint{5.948784in}{0.526498in}}%
\pgfpathlineto{\pgfqpoint{5.952415in}{0.686834in}}%
\pgfpathlineto{\pgfqpoint{5.956046in}{0.650330in}}%
\pgfpathlineto{\pgfqpoint{5.959677in}{0.595820in}}%
\pgfpathlineto{\pgfqpoint{5.963309in}{0.655066in}}%
\pgfpathlineto{\pgfqpoint{5.966940in}{0.645863in}}%
\pgfpathlineto{\pgfqpoint{5.970571in}{0.693473in}}%
\pgfpathlineto{\pgfqpoint{5.974202in}{0.685378in}}%
\pgfpathlineto{\pgfqpoint{5.977833in}{0.684071in}}%
\pgfpathlineto{\pgfqpoint{5.981465in}{0.578993in}}%
\pgfpathlineto{\pgfqpoint{5.985096in}{0.660482in}}%
\pgfpathlineto{\pgfqpoint{5.988727in}{0.664352in}}%
\pgfpathlineto{\pgfqpoint{5.992358in}{0.633669in}}%
\pgfpathlineto{\pgfqpoint{5.995989in}{0.658372in}}%
\pgfpathlineto{\pgfqpoint{5.999620in}{0.625147in}}%
\pgfpathlineto{\pgfqpoint{6.003252in}{0.674667in}}%
\pgfpathlineto{\pgfqpoint{6.006883in}{0.560240in}}%
\pgfpathlineto{\pgfqpoint{6.010514in}{0.687100in}}%
\pgfpathlineto{\pgfqpoint{6.014145in}{0.703168in}}%
\pgfpathlineto{\pgfqpoint{6.017776in}{0.596867in}}%
\pgfpathlineto{\pgfqpoint{6.021407in}{0.669316in}}%
\pgfpathlineto{\pgfqpoint{6.025039in}{0.650899in}}%
\pgfpathlineto{\pgfqpoint{6.028670in}{0.677263in}}%
\pgfpathlineto{\pgfqpoint{6.032301in}{0.605307in}}%
\pgfpathlineto{\pgfqpoint{6.035932in}{0.680375in}}%
\pgfpathlineto{\pgfqpoint{6.039563in}{0.674644in}}%
\pgfpathlineto{\pgfqpoint{6.043194in}{0.658230in}}%
\pgfpathlineto{\pgfqpoint{6.046826in}{0.689888in}}%
\pgfpathlineto{\pgfqpoint{6.050457in}{0.670467in}}%
\pgfpathlineto{\pgfqpoint{6.054088in}{0.686950in}}%
\pgfpathlineto{\pgfqpoint{6.057719in}{0.628294in}}%
\pgfpathlineto{\pgfqpoint{6.061350in}{0.664679in}}%
\pgfpathlineto{\pgfqpoint{6.064981in}{0.560784in}}%
\pgfpathlineto{\pgfqpoint{6.068613in}{0.683654in}}%
\pgfpathlineto{\pgfqpoint{6.072244in}{0.649728in}}%
\pgfpathlineto{\pgfqpoint{6.075875in}{0.565072in}}%
\pgfpathlineto{\pgfqpoint{6.079506in}{0.650994in}}%
\pgfpathlineto{\pgfqpoint{6.083137in}{0.643520in}}%
\pgfpathlineto{\pgfqpoint{6.086769in}{0.695566in}}%
\pgfpathlineto{\pgfqpoint{6.090400in}{0.691451in}}%
\pgfpathlineto{\pgfqpoint{6.094031in}{0.681804in}}%
\pgfpathlineto{\pgfqpoint{6.097662in}{0.578009in}}%
\pgfpathlineto{\pgfqpoint{6.101293in}{0.659138in}}%
\pgfpathlineto{\pgfqpoint{6.104924in}{0.664424in}}%
\pgfpathlineto{\pgfqpoint{6.108556in}{0.615028in}}%
\pgfpathlineto{\pgfqpoint{6.112187in}{0.654176in}}%
\pgfpathlineto{\pgfqpoint{6.115818in}{0.634134in}}%
\pgfpathlineto{\pgfqpoint{6.119449in}{0.672199in}}%
\pgfpathlineto{\pgfqpoint{6.123080in}{0.574399in}}%
\pgfpathlineto{\pgfqpoint{6.126711in}{0.688092in}}%
\pgfpathlineto{\pgfqpoint{6.130343in}{0.706566in}}%
\pgfpathlineto{\pgfqpoint{6.133974in}{0.650443in}}%
\pgfpathlineto{\pgfqpoint{6.137605in}{0.636334in}}%
\pgfpathlineto{\pgfqpoint{6.141236in}{0.632764in}}%
\pgfpathlineto{\pgfqpoint{6.144867in}{0.682001in}}%
\pgfpathlineto{\pgfqpoint{6.148498in}{0.656028in}}%
\pgfpathlineto{\pgfqpoint{6.152130in}{0.624936in}}%
\pgfpathlineto{\pgfqpoint{6.155761in}{0.673237in}}%
\pgfpathlineto{\pgfqpoint{6.159392in}{0.664918in}}%
\pgfpathlineto{\pgfqpoint{6.163023in}{0.649948in}}%
\pgfpathlineto{\pgfqpoint{6.166654in}{0.662804in}}%
\pgfpathlineto{\pgfqpoint{6.170285in}{0.704532in}}%
\pgfpathlineto{\pgfqpoint{6.173917in}{0.678928in}}%
\pgfpathlineto{\pgfqpoint{6.177548in}{0.628499in}}%
\pgfpathlineto{\pgfqpoint{6.181179in}{0.688357in}}%
\pgfpathlineto{\pgfqpoint{6.184810in}{0.670788in}}%
\pgfpathlineto{\pgfqpoint{6.188441in}{0.619243in}}%
\pgfpathlineto{\pgfqpoint{6.192073in}{0.669234in}}%
\pgfpathlineto{\pgfqpoint{6.195704in}{0.695285in}}%
\pgfpathlineto{\pgfqpoint{6.199335in}{0.659446in}}%
\pgfpathlineto{\pgfqpoint{6.202966in}{0.672163in}}%
\pgfpathlineto{\pgfqpoint{6.206597in}{0.703423in}}%
\pgfpathlineto{\pgfqpoint{6.210228in}{0.678755in}}%
\pgfpathlineto{\pgfqpoint{6.213860in}{0.580738in}}%
\pgfpathlineto{\pgfqpoint{6.217491in}{0.660619in}}%
\pgfpathlineto{\pgfqpoint{6.221122in}{0.669868in}}%
\pgfpathlineto{\pgfqpoint{6.224753in}{0.547075in}}%
\pgfpathlineto{\pgfqpoint{6.228384in}{0.646697in}}%
\pgfpathlineto{\pgfqpoint{6.232015in}{0.641999in}}%
\pgfpathlineto{\pgfqpoint{6.235647in}{0.669014in}}%
\pgfpathlineto{\pgfqpoint{6.239278in}{0.574649in}}%
\pgfpathlineto{\pgfqpoint{6.242909in}{0.689606in}}%
\pgfpathlineto{\pgfqpoint{6.246540in}{0.709081in}}%
\pgfpathlineto{\pgfqpoint{6.250171in}{0.643483in}}%
\pgfpathlineto{\pgfqpoint{6.253802in}{0.632761in}}%
\pgfpathlineto{\pgfqpoint{6.257434in}{0.630641in}}%
\pgfpathlineto{\pgfqpoint{6.261065in}{0.680596in}}%
\pgfpathlineto{\pgfqpoint{6.264696in}{0.664886in}}%
\pgfpathlineto{\pgfqpoint{6.268327in}{0.596263in}}%
\pgfpathlineto{\pgfqpoint{6.271958in}{0.675835in}}%
\pgfpathlineto{\pgfqpoint{6.275590in}{0.661386in}}%
\pgfpathlineto{\pgfqpoint{6.279221in}{0.650240in}}%
\pgfpathlineto{\pgfqpoint{6.282852in}{0.661860in}}%
\pgfpathlineto{\pgfqpoint{6.286483in}{0.704112in}}%
\pgfpathlineto{\pgfqpoint{6.290114in}{0.601248in}}%
\pgfpathlineto{\pgfqpoint{6.293745in}{0.635626in}}%
\pgfpathlineto{\pgfqpoint{6.297377in}{0.600541in}}%
\pgfpathlineto{\pgfqpoint{6.301008in}{0.672523in}}%
\pgfpathlineto{\pgfqpoint{6.304639in}{0.687441in}}%
\pgfpathlineto{\pgfqpoint{6.308270in}{0.645697in}}%
\pgfpathlineto{\pgfqpoint{6.311901in}{0.690316in}}%
\pgfpathlineto{\pgfqpoint{6.315532in}{0.623245in}}%
\pgfpathlineto{\pgfqpoint{6.319164in}{0.664002in}}%
\pgfpathlineto{\pgfqpoint{6.322795in}{0.581865in}}%
\pgfpathlineto{\pgfqpoint{6.326426in}{0.683891in}}%
\pgfpathlineto{\pgfqpoint{6.330057in}{0.653405in}}%
\pgfpathlineto{\pgfqpoint{6.333688in}{0.592024in}}%
\pgfpathlineto{\pgfqpoint{6.337319in}{0.661372in}}%
\pgfpathlineto{\pgfqpoint{6.340951in}{0.640473in}}%
\pgfpathlineto{\pgfqpoint{6.344582in}{0.692032in}}%
\pgfpathlineto{\pgfqpoint{6.348213in}{0.681160in}}%
\pgfpathlineto{\pgfqpoint{6.351844in}{0.686082in}}%
\pgfpathlineto{\pgfqpoint{6.355475in}{0.666947in}}%
\pgfpathlineto{\pgfqpoint{6.359106in}{0.509178in}}%
\pgfpathlineto{\pgfqpoint{6.362738in}{0.661196in}}%
\pgfpathlineto{\pgfqpoint{6.366369in}{0.653921in}}%
\pgfpathlineto{\pgfqpoint{6.370000in}{0.601129in}}%
\pgfpathlineto{\pgfqpoint{6.370000in}{0.601129in}}%
\pgfusepath{stroke}%
\end{pgfscope}%
\begin{pgfscope}%
\pgfpathrectangle{\pgfqpoint{4.554414in}{0.320625in}}{\pgfqpoint{1.815586in}{1.164375in}}%
\pgfusepath{clip}%
\pgfsetrectcap%
\pgfsetroundjoin%
\pgfsetlinewidth{0.602250pt}%
\definecolor{currentstroke}{rgb}{0.819608,0.286275,0.356863}%
\pgfsetstrokecolor{currentstroke}%
\pgfsetstrokeopacity{0.120000}%
\pgfsetdash{}{0pt}%
\pgfpathmoveto{\pgfqpoint{4.554414in}{1.463570in}}%
\pgfpathlineto{\pgfqpoint{4.558045in}{1.451221in}}%
\pgfpathlineto{\pgfqpoint{4.561676in}{1.442037in}}%
\pgfpathlineto{\pgfqpoint{4.568938in}{1.427611in}}%
\pgfpathlineto{\pgfqpoint{4.619775in}{1.332228in}}%
\pgfpathlineto{\pgfqpoint{4.645193in}{1.285697in}}%
\pgfpathlineto{\pgfqpoint{4.670611in}{1.241280in}}%
\pgfpathlineto{\pgfqpoint{4.699660in}{1.192634in}}%
\pgfpathlineto{\pgfqpoint{4.735972in}{1.134064in}}%
\pgfpathlineto{\pgfqpoint{4.779546in}{1.065999in}}%
\pgfpathlineto{\pgfqpoint{4.830383in}{0.988689in}}%
\pgfpathlineto{\pgfqpoint{4.892113in}{0.896918in}}%
\pgfpathlineto{\pgfqpoint{4.964736in}{0.791030in}}%
\pgfpathlineto{\pgfqpoint{5.022835in}{0.707381in}}%
\pgfpathlineto{\pgfqpoint{5.026466in}{0.705845in}}%
\pgfpathlineto{\pgfqpoint{5.030097in}{0.712302in}}%
\pgfpathlineto{\pgfqpoint{5.033728in}{0.692283in}}%
\pgfpathlineto{\pgfqpoint{5.037360in}{0.686727in}}%
\pgfpathlineto{\pgfqpoint{5.040991in}{0.693329in}}%
\pgfpathlineto{\pgfqpoint{5.044622in}{0.677355in}}%
\pgfpathlineto{\pgfqpoint{5.048253in}{0.673995in}}%
\pgfpathlineto{\pgfqpoint{5.051884in}{0.688130in}}%
\pgfpathlineto{\pgfqpoint{5.055515in}{0.688319in}}%
\pgfpathlineto{\pgfqpoint{5.059147in}{0.685066in}}%
\pgfpathlineto{\pgfqpoint{5.062778in}{0.661810in}}%
\pgfpathlineto{\pgfqpoint{5.066409in}{0.660408in}}%
\pgfpathlineto{\pgfqpoint{5.070040in}{0.668064in}}%
\pgfpathlineto{\pgfqpoint{5.073671in}{0.646295in}}%
\pgfpathlineto{\pgfqpoint{5.077302in}{0.665082in}}%
\pgfpathlineto{\pgfqpoint{5.080934in}{0.678738in}}%
\pgfpathlineto{\pgfqpoint{5.088196in}{0.625110in}}%
\pgfpathlineto{\pgfqpoint{5.091827in}{0.692497in}}%
\pgfpathlineto{\pgfqpoint{5.095458in}{0.656428in}}%
\pgfpathlineto{\pgfqpoint{5.099090in}{0.629911in}}%
\pgfpathlineto{\pgfqpoint{5.102721in}{0.668173in}}%
\pgfpathlineto{\pgfqpoint{5.106352in}{0.647578in}}%
\pgfpathlineto{\pgfqpoint{5.109983in}{0.593113in}}%
\pgfpathlineto{\pgfqpoint{5.113614in}{0.650600in}}%
\pgfpathlineto{\pgfqpoint{5.117245in}{0.683412in}}%
\pgfpathlineto{\pgfqpoint{5.120877in}{0.703093in}}%
\pgfpathlineto{\pgfqpoint{5.124508in}{0.624510in}}%
\pgfpathlineto{\pgfqpoint{5.128139in}{0.593391in}}%
\pgfpathlineto{\pgfqpoint{5.131770in}{0.675706in}}%
\pgfpathlineto{\pgfqpoint{5.135401in}{0.556658in}}%
\pgfpathlineto{\pgfqpoint{5.139032in}{0.649003in}}%
\pgfpathlineto{\pgfqpoint{5.142664in}{0.679880in}}%
\pgfpathlineto{\pgfqpoint{5.146295in}{0.680898in}}%
\pgfpathlineto{\pgfqpoint{5.153557in}{0.656476in}}%
\pgfpathlineto{\pgfqpoint{5.157188in}{0.650093in}}%
\pgfpathlineto{\pgfqpoint{5.160819in}{0.665696in}}%
\pgfpathlineto{\pgfqpoint{5.164451in}{0.643589in}}%
\pgfpathlineto{\pgfqpoint{5.168082in}{0.651812in}}%
\pgfpathlineto{\pgfqpoint{5.171713in}{0.651426in}}%
\pgfpathlineto{\pgfqpoint{5.175344in}{0.597384in}}%
\pgfpathlineto{\pgfqpoint{5.178975in}{0.662272in}}%
\pgfpathlineto{\pgfqpoint{5.182606in}{0.700819in}}%
\pgfpathlineto{\pgfqpoint{5.186238in}{0.665735in}}%
\pgfpathlineto{\pgfqpoint{5.189869in}{0.644164in}}%
\pgfpathlineto{\pgfqpoint{5.193500in}{0.671475in}}%
\pgfpathlineto{\pgfqpoint{5.197131in}{0.622065in}}%
\pgfpathlineto{\pgfqpoint{5.200762in}{0.612589in}}%
\pgfpathlineto{\pgfqpoint{5.204394in}{0.664607in}}%
\pgfpathlineto{\pgfqpoint{5.208025in}{0.691394in}}%
\pgfpathlineto{\pgfqpoint{5.211656in}{0.703899in}}%
\pgfpathlineto{\pgfqpoint{5.215287in}{0.640448in}}%
\pgfpathlineto{\pgfqpoint{5.218918in}{0.505892in}}%
\pgfpathlineto{\pgfqpoint{5.222549in}{0.682027in}}%
\pgfpathlineto{\pgfqpoint{5.226181in}{0.583688in}}%
\pgfpathlineto{\pgfqpoint{5.229812in}{0.645115in}}%
\pgfpathlineto{\pgfqpoint{5.233443in}{0.669220in}}%
\pgfpathlineto{\pgfqpoint{5.237074in}{0.668405in}}%
\pgfpathlineto{\pgfqpoint{5.240705in}{0.652454in}}%
\pgfpathlineto{\pgfqpoint{5.244336in}{0.671040in}}%
\pgfpathlineto{\pgfqpoint{5.247968in}{0.630617in}}%
\pgfpathlineto{\pgfqpoint{5.251599in}{0.394173in}}%
\pgfpathlineto{\pgfqpoint{5.255230in}{0.655856in}}%
\pgfpathlineto{\pgfqpoint{5.258861in}{0.683728in}}%
\pgfpathlineto{\pgfqpoint{5.262492in}{0.695986in}}%
\pgfpathlineto{\pgfqpoint{5.266123in}{0.656421in}}%
\pgfpathlineto{\pgfqpoint{5.269755in}{0.545334in}}%
\pgfpathlineto{\pgfqpoint{5.273386in}{0.692300in}}%
\pgfpathlineto{\pgfqpoint{5.277017in}{0.553586in}}%
\pgfpathlineto{\pgfqpoint{5.280648in}{0.624839in}}%
\pgfpathlineto{\pgfqpoint{5.284279in}{0.666044in}}%
\pgfpathlineto{\pgfqpoint{5.287910in}{0.657780in}}%
\pgfpathlineto{\pgfqpoint{5.291542in}{0.566420in}}%
\pgfpathlineto{\pgfqpoint{5.295173in}{0.660381in}}%
\pgfpathlineto{\pgfqpoint{5.298804in}{0.674721in}}%
\pgfpathlineto{\pgfqpoint{5.302435in}{0.697991in}}%
\pgfpathlineto{\pgfqpoint{5.306066in}{0.634631in}}%
\pgfpathlineto{\pgfqpoint{5.309698in}{0.628683in}}%
\pgfpathlineto{\pgfqpoint{5.313329in}{0.627955in}}%
\pgfpathlineto{\pgfqpoint{5.316960in}{0.602784in}}%
\pgfpathlineto{\pgfqpoint{5.320591in}{0.679904in}}%
\pgfpathlineto{\pgfqpoint{5.324222in}{0.699456in}}%
\pgfpathlineto{\pgfqpoint{5.331485in}{0.571474in}}%
\pgfpathlineto{\pgfqpoint{5.335116in}{0.661061in}}%
\pgfpathlineto{\pgfqpoint{5.338747in}{0.633277in}}%
\pgfpathlineto{\pgfqpoint{5.342378in}{0.621396in}}%
\pgfpathlineto{\pgfqpoint{5.346009in}{0.659461in}}%
\pgfpathlineto{\pgfqpoint{5.349640in}{0.641155in}}%
\pgfpathlineto{\pgfqpoint{5.353272in}{0.663541in}}%
\pgfpathlineto{\pgfqpoint{5.356903in}{0.634704in}}%
\pgfpathlineto{\pgfqpoint{5.360534in}{0.654559in}}%
\pgfpathlineto{\pgfqpoint{5.364165in}{0.661886in}}%
\pgfpathlineto{\pgfqpoint{5.367796in}{0.649898in}}%
\pgfpathlineto{\pgfqpoint{5.371427in}{0.613692in}}%
\pgfpathlineto{\pgfqpoint{5.375059in}{0.680619in}}%
\pgfpathlineto{\pgfqpoint{5.378690in}{0.595823in}}%
\pgfpathlineto{\pgfqpoint{5.382321in}{0.626157in}}%
\pgfpathlineto{\pgfqpoint{5.385952in}{0.669085in}}%
\pgfpathlineto{\pgfqpoint{5.389583in}{0.683793in}}%
\pgfpathlineto{\pgfqpoint{5.393215in}{0.685861in}}%
\pgfpathlineto{\pgfqpoint{5.396846in}{0.642331in}}%
\pgfpathlineto{\pgfqpoint{5.400477in}{0.631521in}}%
\pgfpathlineto{\pgfqpoint{5.404108in}{0.628584in}}%
\pgfpathlineto{\pgfqpoint{5.407739in}{0.618926in}}%
\pgfpathlineto{\pgfqpoint{5.411370in}{0.678871in}}%
\pgfpathlineto{\pgfqpoint{5.415002in}{0.698133in}}%
\pgfpathlineto{\pgfqpoint{5.418633in}{0.627707in}}%
\pgfpathlineto{\pgfqpoint{5.422264in}{0.525104in}}%
\pgfpathlineto{\pgfqpoint{5.425895in}{0.665814in}}%
\pgfpathlineto{\pgfqpoint{5.429526in}{0.635011in}}%
\pgfpathlineto{\pgfqpoint{5.433157in}{0.619812in}}%
\pgfpathlineto{\pgfqpoint{5.436789in}{0.667903in}}%
\pgfpathlineto{\pgfqpoint{5.440420in}{0.649534in}}%
\pgfpathlineto{\pgfqpoint{5.444051in}{0.658873in}}%
\pgfpathlineto{\pgfqpoint{5.447682in}{0.641793in}}%
\pgfpathlineto{\pgfqpoint{5.451313in}{0.656458in}}%
\pgfpathlineto{\pgfqpoint{5.454944in}{0.660713in}}%
\pgfpathlineto{\pgfqpoint{5.458576in}{0.630129in}}%
\pgfpathlineto{\pgfqpoint{5.462207in}{0.645995in}}%
\pgfpathlineto{\pgfqpoint{5.465838in}{0.699094in}}%
\pgfpathlineto{\pgfqpoint{5.469469in}{0.658638in}}%
\pgfpathlineto{\pgfqpoint{5.473100in}{0.636016in}}%
\pgfpathlineto{\pgfqpoint{5.476731in}{0.668269in}}%
\pgfpathlineto{\pgfqpoint{5.480363in}{0.626990in}}%
\pgfpathlineto{\pgfqpoint{5.483994in}{0.610952in}}%
\pgfpathlineto{\pgfqpoint{5.487625in}{0.662350in}}%
\pgfpathlineto{\pgfqpoint{5.491256in}{0.692961in}}%
\pgfpathlineto{\pgfqpoint{5.494887in}{0.708438in}}%
\pgfpathlineto{\pgfqpoint{5.498519in}{0.639746in}}%
\pgfpathlineto{\pgfqpoint{5.502150in}{0.507253in}}%
\pgfpathlineto{\pgfqpoint{5.505781in}{0.675510in}}%
\pgfpathlineto{\pgfqpoint{5.509412in}{0.580690in}}%
\pgfpathlineto{\pgfqpoint{5.513043in}{0.650502in}}%
\pgfpathlineto{\pgfqpoint{5.516674in}{0.674134in}}%
\pgfpathlineto{\pgfqpoint{5.520306in}{0.679096in}}%
\pgfpathlineto{\pgfqpoint{5.523937in}{0.671452in}}%
\pgfpathlineto{\pgfqpoint{5.527568in}{0.668423in}}%
\pgfpathlineto{\pgfqpoint{5.531199in}{0.620651in}}%
\pgfpathlineto{\pgfqpoint{5.534830in}{0.541934in}}%
\pgfpathlineto{\pgfqpoint{5.538461in}{0.653610in}}%
\pgfpathlineto{\pgfqpoint{5.542093in}{0.686746in}}%
\pgfpathlineto{\pgfqpoint{5.545724in}{0.699633in}}%
\pgfpathlineto{\pgfqpoint{5.549355in}{0.492108in}}%
\pgfpathlineto{\pgfqpoint{5.552986in}{0.612273in}}%
\pgfpathlineto{\pgfqpoint{5.556617in}{0.653291in}}%
\pgfpathlineto{\pgfqpoint{5.560248in}{0.569588in}}%
\pgfpathlineto{\pgfqpoint{5.563880in}{0.664564in}}%
\pgfpathlineto{\pgfqpoint{5.567511in}{0.688607in}}%
\pgfpathlineto{\pgfqpoint{5.571142in}{0.646943in}}%
\pgfpathlineto{\pgfqpoint{5.574773in}{0.627976in}}%
\pgfpathlineto{\pgfqpoint{5.578404in}{0.658628in}}%
\pgfpathlineto{\pgfqpoint{5.582035in}{0.650889in}}%
\pgfpathlineto{\pgfqpoint{5.585667in}{0.613143in}}%
\pgfpathlineto{\pgfqpoint{5.589298in}{0.601936in}}%
\pgfpathlineto{\pgfqpoint{5.592929in}{0.624859in}}%
\pgfpathlineto{\pgfqpoint{5.596560in}{0.676976in}}%
\pgfpathlineto{\pgfqpoint{5.600191in}{0.592424in}}%
\pgfpathlineto{\pgfqpoint{5.603823in}{0.650223in}}%
\pgfpathlineto{\pgfqpoint{5.607454in}{0.673145in}}%
\pgfpathlineto{\pgfqpoint{5.611085in}{0.669175in}}%
\pgfpathlineto{\pgfqpoint{5.614716in}{0.651377in}}%
\pgfpathlineto{\pgfqpoint{5.618347in}{0.673291in}}%
\pgfpathlineto{\pgfqpoint{5.621978in}{0.627131in}}%
\pgfpathlineto{\pgfqpoint{5.625610in}{0.550645in}}%
\pgfpathlineto{\pgfqpoint{5.629241in}{0.658792in}}%
\pgfpathlineto{\pgfqpoint{5.632872in}{0.683867in}}%
\pgfpathlineto{\pgfqpoint{5.636503in}{0.694278in}}%
\pgfpathlineto{\pgfqpoint{5.640134in}{0.651938in}}%
\pgfpathlineto{\pgfqpoint{5.643765in}{0.495546in}}%
\pgfpathlineto{\pgfqpoint{5.647397in}{0.693356in}}%
\pgfpathlineto{\pgfqpoint{5.651028in}{0.501229in}}%
\pgfpathlineto{\pgfqpoint{5.654659in}{0.626869in}}%
\pgfpathlineto{\pgfqpoint{5.658290in}{0.664427in}}%
\pgfpathlineto{\pgfqpoint{5.661921in}{0.653367in}}%
\pgfpathlineto{\pgfqpoint{5.665552in}{0.555102in}}%
\pgfpathlineto{\pgfqpoint{5.669184in}{0.666729in}}%
\pgfpathlineto{\pgfqpoint{5.672815in}{0.677527in}}%
\pgfpathlineto{\pgfqpoint{5.676446in}{0.697977in}}%
\pgfpathlineto{\pgfqpoint{5.680077in}{0.639736in}}%
\pgfpathlineto{\pgfqpoint{5.683708in}{0.627429in}}%
\pgfpathlineto{\pgfqpoint{5.687340in}{0.620082in}}%
\pgfpathlineto{\pgfqpoint{5.690971in}{0.615081in}}%
\pgfpathlineto{\pgfqpoint{5.694602in}{0.682003in}}%
\pgfpathlineto{\pgfqpoint{5.698233in}{0.700434in}}%
\pgfpathlineto{\pgfqpoint{5.701864in}{0.628656in}}%
\pgfpathlineto{\pgfqpoint{5.705495in}{0.544060in}}%
\pgfpathlineto{\pgfqpoint{5.709127in}{0.660811in}}%
\pgfpathlineto{\pgfqpoint{5.712758in}{0.629840in}}%
\pgfpathlineto{\pgfqpoint{5.716389in}{0.630464in}}%
\pgfpathlineto{\pgfqpoint{5.720020in}{0.664237in}}%
\pgfpathlineto{\pgfqpoint{5.723651in}{0.642668in}}%
\pgfpathlineto{\pgfqpoint{5.727282in}{0.660823in}}%
\pgfpathlineto{\pgfqpoint{5.730914in}{0.638521in}}%
\pgfpathlineto{\pgfqpoint{5.734545in}{0.654570in}}%
\pgfpathlineto{\pgfqpoint{5.738176in}{0.658753in}}%
\pgfpathlineto{\pgfqpoint{5.741807in}{0.644250in}}%
\pgfpathlineto{\pgfqpoint{5.745438in}{0.600986in}}%
\pgfpathlineto{\pgfqpoint{5.749069in}{0.681174in}}%
\pgfpathlineto{\pgfqpoint{5.752701in}{0.586535in}}%
\pgfpathlineto{\pgfqpoint{5.759963in}{0.670124in}}%
\pgfpathlineto{\pgfqpoint{5.763594in}{0.683152in}}%
\pgfpathlineto{\pgfqpoint{5.767225in}{0.683717in}}%
\pgfpathlineto{\pgfqpoint{5.770856in}{0.647302in}}%
\pgfpathlineto{\pgfqpoint{5.774488in}{0.631674in}}%
\pgfpathlineto{\pgfqpoint{5.778119in}{0.622197in}}%
\pgfpathlineto{\pgfqpoint{5.781750in}{0.625949in}}%
\pgfpathlineto{\pgfqpoint{5.785381in}{0.680220in}}%
\pgfpathlineto{\pgfqpoint{5.789012in}{0.698711in}}%
\pgfpathlineto{\pgfqpoint{5.792644in}{0.622326in}}%
\pgfpathlineto{\pgfqpoint{5.796275in}{0.559409in}}%
\pgfpathlineto{\pgfqpoint{5.799906in}{0.665460in}}%
\pgfpathlineto{\pgfqpoint{5.803537in}{0.631051in}}%
\pgfpathlineto{\pgfqpoint{5.807168in}{0.629333in}}%
\pgfpathlineto{\pgfqpoint{5.810799in}{0.671588in}}%
\pgfpathlineto{\pgfqpoint{5.814431in}{0.650136in}}%
\pgfpathlineto{\pgfqpoint{5.818062in}{0.655661in}}%
\pgfpathlineto{\pgfqpoint{5.821693in}{0.645176in}}%
\pgfpathlineto{\pgfqpoint{5.825324in}{0.656408in}}%
\pgfpathlineto{\pgfqpoint{5.828955in}{0.657395in}}%
\pgfpathlineto{\pgfqpoint{5.832586in}{0.620723in}}%
\pgfpathlineto{\pgfqpoint{5.836218in}{0.652492in}}%
\pgfpathlineto{\pgfqpoint{5.839849in}{0.700880in}}%
\pgfpathlineto{\pgfqpoint{5.843480in}{0.657068in}}%
\pgfpathlineto{\pgfqpoint{5.847111in}{0.630130in}}%
\pgfpathlineto{\pgfqpoint{5.850742in}{0.668938in}}%
\pgfpathlineto{\pgfqpoint{5.854373in}{0.623403in}}%
\pgfpathlineto{\pgfqpoint{5.858005in}{0.620082in}}%
\pgfpathlineto{\pgfqpoint{5.861636in}{0.665485in}}%
\pgfpathlineto{\pgfqpoint{5.865267in}{0.693634in}}%
\pgfpathlineto{\pgfqpoint{5.868898in}{0.707993in}}%
\pgfpathlineto{\pgfqpoint{5.872529in}{0.631838in}}%
\pgfpathlineto{\pgfqpoint{5.876160in}{0.520714in}}%
\pgfpathlineto{\pgfqpoint{5.879792in}{0.673960in}}%
\pgfpathlineto{\pgfqpoint{5.883423in}{0.592816in}}%
\pgfpathlineto{\pgfqpoint{5.887054in}{0.651878in}}%
\pgfpathlineto{\pgfqpoint{5.890685in}{0.673818in}}%
\pgfpathlineto{\pgfqpoint{5.894316in}{0.677411in}}%
\pgfpathlineto{\pgfqpoint{5.897948in}{0.668163in}}%
\pgfpathlineto{\pgfqpoint{5.901579in}{0.670614in}}%
\pgfpathlineto{\pgfqpoint{5.908841in}{0.569882in}}%
\pgfpathlineto{\pgfqpoint{5.912472in}{0.656089in}}%
\pgfpathlineto{\pgfqpoint{5.916103in}{0.687018in}}%
\pgfpathlineto{\pgfqpoint{5.919735in}{0.699001in}}%
\pgfpathlineto{\pgfqpoint{5.923366in}{0.558814in}}%
\pgfpathlineto{\pgfqpoint{5.926997in}{0.615640in}}%
\pgfpathlineto{\pgfqpoint{5.930628in}{0.650741in}}%
\pgfpathlineto{\pgfqpoint{5.934259in}{0.539171in}}%
\pgfpathlineto{\pgfqpoint{5.937890in}{0.667434in}}%
\pgfpathlineto{\pgfqpoint{5.941522in}{0.690455in}}%
\pgfpathlineto{\pgfqpoint{5.945153in}{0.645819in}}%
\pgfpathlineto{\pgfqpoint{5.948784in}{0.621862in}}%
\pgfpathlineto{\pgfqpoint{5.952415in}{0.659899in}}%
\pgfpathlineto{\pgfqpoint{5.956046in}{0.649522in}}%
\pgfpathlineto{\pgfqpoint{5.959677in}{0.598946in}}%
\pgfpathlineto{\pgfqpoint{5.963309in}{0.617276in}}%
\pgfpathlineto{\pgfqpoint{5.966940in}{0.629370in}}%
\pgfpathlineto{\pgfqpoint{5.970571in}{0.675374in}}%
\pgfpathlineto{\pgfqpoint{5.974202in}{0.602333in}}%
\pgfpathlineto{\pgfqpoint{5.977833in}{0.651454in}}%
\pgfpathlineto{\pgfqpoint{5.981465in}{0.672409in}}%
\pgfpathlineto{\pgfqpoint{5.985096in}{0.666807in}}%
\pgfpathlineto{\pgfqpoint{5.988727in}{0.646342in}}%
\pgfpathlineto{\pgfqpoint{5.992358in}{0.674984in}}%
\pgfpathlineto{\pgfqpoint{5.999620in}{0.575437in}}%
\pgfpathlineto{\pgfqpoint{6.003252in}{0.661045in}}%
\pgfpathlineto{\pgfqpoint{6.006883in}{0.683957in}}%
\pgfpathlineto{\pgfqpoint{6.010514in}{0.692912in}}%
\pgfpathlineto{\pgfqpoint{6.014145in}{0.645716in}}%
\pgfpathlineto{\pgfqpoint{6.017776in}{0.567759in}}%
\pgfpathlineto{\pgfqpoint{6.021407in}{0.693468in}}%
\pgfpathlineto{\pgfqpoint{6.025039in}{0.509034in}}%
\pgfpathlineto{\pgfqpoint{6.028670in}{0.629511in}}%
\pgfpathlineto{\pgfqpoint{6.032301in}{0.663519in}}%
\pgfpathlineto{\pgfqpoint{6.035932in}{0.650411in}}%
\pgfpathlineto{\pgfqpoint{6.039563in}{0.593011in}}%
\pgfpathlineto{\pgfqpoint{6.043194in}{0.671453in}}%
\pgfpathlineto{\pgfqpoint{6.046826in}{0.678759in}}%
\pgfpathlineto{\pgfqpoint{6.050457in}{0.696491in}}%
\pgfpathlineto{\pgfqpoint{6.054088in}{0.644308in}}%
\pgfpathlineto{\pgfqpoint{6.061350in}{0.613102in}}%
\pgfpathlineto{\pgfqpoint{6.064981in}{0.622201in}}%
\pgfpathlineto{\pgfqpoint{6.068613in}{0.683364in}}%
\pgfpathlineto{\pgfqpoint{6.072244in}{0.701282in}}%
\pgfpathlineto{\pgfqpoint{6.075875in}{0.624502in}}%
\pgfpathlineto{\pgfqpoint{6.079506in}{0.489813in}}%
\pgfpathlineto{\pgfqpoint{6.083137in}{0.660314in}}%
\pgfpathlineto{\pgfqpoint{6.086769in}{0.625921in}}%
\pgfpathlineto{\pgfqpoint{6.090400in}{0.637559in}}%
\pgfpathlineto{\pgfqpoint{6.094031in}{0.668211in}}%
\pgfpathlineto{\pgfqpoint{6.097662in}{0.643894in}}%
\pgfpathlineto{\pgfqpoint{6.101293in}{0.658136in}}%
\pgfpathlineto{\pgfqpoint{6.104924in}{0.641608in}}%
\pgfpathlineto{\pgfqpoint{6.108556in}{0.654504in}}%
\pgfpathlineto{\pgfqpoint{6.112187in}{0.655659in}}%
\pgfpathlineto{\pgfqpoint{6.115818in}{0.638159in}}%
\pgfpathlineto{\pgfqpoint{6.119449in}{0.584952in}}%
\pgfpathlineto{\pgfqpoint{6.123080in}{0.681505in}}%
\pgfpathlineto{\pgfqpoint{6.126711in}{0.575601in}}%
\pgfpathlineto{\pgfqpoint{6.130343in}{0.633971in}}%
\pgfpathlineto{\pgfqpoint{6.133974in}{0.670920in}}%
\pgfpathlineto{\pgfqpoint{6.137605in}{0.682401in}}%
\pgfpathlineto{\pgfqpoint{6.141236in}{0.681542in}}%
\pgfpathlineto{\pgfqpoint{6.144867in}{0.651345in}}%
\pgfpathlineto{\pgfqpoint{6.148498in}{0.631661in}}%
\pgfpathlineto{\pgfqpoint{6.152130in}{0.615558in}}%
\pgfpathlineto{\pgfqpoint{6.155761in}{0.631389in}}%
\pgfpathlineto{\pgfqpoint{6.159392in}{0.681266in}}%
\pgfpathlineto{\pgfqpoint{6.163023in}{0.699078in}}%
\pgfpathlineto{\pgfqpoint{6.166654in}{0.616410in}}%
\pgfpathlineto{\pgfqpoint{6.170285in}{0.576172in}}%
\pgfpathlineto{\pgfqpoint{6.173917in}{0.664950in}}%
\pgfpathlineto{\pgfqpoint{6.177548in}{0.626964in}}%
\pgfpathlineto{\pgfqpoint{6.181179in}{0.636260in}}%
\pgfpathlineto{\pgfqpoint{6.184810in}{0.674647in}}%
\pgfpathlineto{\pgfqpoint{6.188441in}{0.650512in}}%
\pgfpathlineto{\pgfqpoint{6.192073in}{0.652490in}}%
\pgfpathlineto{\pgfqpoint{6.195704in}{0.647965in}}%
\pgfpathlineto{\pgfqpoint{6.199335in}{0.656242in}}%
\pgfpathlineto{\pgfqpoint{6.202966in}{0.654003in}}%
\pgfpathlineto{\pgfqpoint{6.206597in}{0.609319in}}%
\pgfpathlineto{\pgfqpoint{6.213860in}{0.702350in}}%
\pgfpathlineto{\pgfqpoint{6.217491in}{0.655341in}}%
\pgfpathlineto{\pgfqpoint{6.221122in}{0.623922in}}%
\pgfpathlineto{\pgfqpoint{6.224753in}{0.669394in}}%
\pgfpathlineto{\pgfqpoint{6.228384in}{0.619712in}}%
\pgfpathlineto{\pgfqpoint{6.232015in}{0.626747in}}%
\pgfpathlineto{\pgfqpoint{6.235647in}{0.668098in}}%
\pgfpathlineto{\pgfqpoint{6.239278in}{0.694088in}}%
\pgfpathlineto{\pgfqpoint{6.242909in}{0.707333in}}%
\pgfpathlineto{\pgfqpoint{6.246540in}{0.622299in}}%
\pgfpathlineto{\pgfqpoint{6.250171in}{0.560473in}}%
\pgfpathlineto{\pgfqpoint{6.253802in}{0.672431in}}%
\pgfpathlineto{\pgfqpoint{6.257434in}{0.601496in}}%
\pgfpathlineto{\pgfqpoint{6.261065in}{0.652970in}}%
\pgfpathlineto{\pgfqpoint{6.264696in}{0.673365in}}%
\pgfpathlineto{\pgfqpoint{6.268327in}{0.675629in}}%
\pgfpathlineto{\pgfqpoint{6.271958in}{0.664790in}}%
\pgfpathlineto{\pgfqpoint{6.275590in}{0.672447in}}%
\pgfpathlineto{\pgfqpoint{6.279221in}{0.616003in}}%
\pgfpathlineto{\pgfqpoint{6.282852in}{0.584921in}}%
\pgfpathlineto{\pgfqpoint{6.286483in}{0.658163in}}%
\pgfpathlineto{\pgfqpoint{6.290114in}{0.687136in}}%
\pgfpathlineto{\pgfqpoint{6.293745in}{0.698187in}}%
\pgfpathlineto{\pgfqpoint{6.297377in}{0.582427in}}%
\pgfpathlineto{\pgfqpoint{6.301008in}{0.618319in}}%
\pgfpathlineto{\pgfqpoint{6.304639in}{0.648211in}}%
\pgfpathlineto{\pgfqpoint{6.308270in}{0.367294in}}%
\pgfpathlineto{\pgfqpoint{6.311901in}{0.669781in}}%
\pgfpathlineto{\pgfqpoint{6.315532in}{0.692009in}}%
\pgfpathlineto{\pgfqpoint{6.319164in}{0.644560in}}%
\pgfpathlineto{\pgfqpoint{6.322795in}{0.615400in}}%
\pgfpathlineto{\pgfqpoint{6.326426in}{0.660928in}}%
\pgfpathlineto{\pgfqpoint{6.330057in}{0.648112in}}%
\pgfpathlineto{\pgfqpoint{6.333688in}{0.579507in}}%
\pgfpathlineto{\pgfqpoint{6.337319in}{0.627894in}}%
\pgfpathlineto{\pgfqpoint{6.340951in}{0.632989in}}%
\pgfpathlineto{\pgfqpoint{6.344582in}{0.673776in}}%
\pgfpathlineto{\pgfqpoint{6.348213in}{0.609808in}}%
\pgfpathlineto{\pgfqpoint{6.351844in}{0.652436in}}%
\pgfpathlineto{\pgfqpoint{6.355475in}{0.671529in}}%
\pgfpathlineto{\pgfqpoint{6.359106in}{0.664309in}}%
\pgfpathlineto{\pgfqpoint{6.362738in}{0.640994in}}%
\pgfpathlineto{\pgfqpoint{6.366369in}{0.676397in}}%
\pgfpathlineto{\pgfqpoint{6.370000in}{0.610020in}}%
\pgfpathlineto{\pgfqpoint{6.370000in}{0.610020in}}%
\pgfusepath{stroke}%
\end{pgfscope}%
\begin{pgfscope}%
\pgfpathrectangle{\pgfqpoint{4.554414in}{0.320625in}}{\pgfqpoint{1.815586in}{1.164375in}}%
\pgfusepath{clip}%
\pgfsetrectcap%
\pgfsetroundjoin%
\pgfsetlinewidth{0.602250pt}%
\definecolor{currentstroke}{rgb}{0.819608,0.286275,0.356863}%
\pgfsetstrokecolor{currentstroke}%
\pgfsetstrokeopacity{0.120000}%
\pgfsetdash{}{0pt}%
\pgfpathmoveto{\pgfqpoint{4.554414in}{1.467372in}}%
\pgfpathlineto{\pgfqpoint{4.561676in}{1.444409in}}%
\pgfpathlineto{\pgfqpoint{4.568938in}{1.428501in}}%
\pgfpathlineto{\pgfqpoint{4.627037in}{1.314472in}}%
\pgfpathlineto{\pgfqpoint{4.648824in}{1.274610in}}%
\pgfpathlineto{\pgfqpoint{4.670611in}{1.236883in}}%
\pgfpathlineto{\pgfqpoint{4.696029in}{1.195099in}}%
\pgfpathlineto{\pgfqpoint{4.725079in}{1.149509in}}%
\pgfpathlineto{\pgfqpoint{4.761390in}{1.094714in}}%
\pgfpathlineto{\pgfqpoint{4.804965in}{1.031082in}}%
\pgfpathlineto{\pgfqpoint{4.855801in}{0.958815in}}%
\pgfpathlineto{\pgfqpoint{4.917531in}{0.873027in}}%
\pgfpathlineto{\pgfqpoint{4.990154in}{0.774016in}}%
\pgfpathlineto{\pgfqpoint{4.997417in}{0.765068in}}%
\pgfpathlineto{\pgfqpoint{5.004679in}{0.754329in}}%
\pgfpathlineto{\pgfqpoint{5.015573in}{0.739752in}}%
\pgfpathlineto{\pgfqpoint{5.022835in}{0.729865in}}%
\pgfpathlineto{\pgfqpoint{5.026466in}{0.727157in}}%
\pgfpathlineto{\pgfqpoint{5.030097in}{0.720203in}}%
\pgfpathlineto{\pgfqpoint{5.033728in}{0.715219in}}%
\pgfpathlineto{\pgfqpoint{5.037360in}{0.711547in}}%
\pgfpathlineto{\pgfqpoint{5.040991in}{0.711518in}}%
\pgfpathlineto{\pgfqpoint{5.044622in}{0.713863in}}%
\pgfpathlineto{\pgfqpoint{5.048253in}{0.696486in}}%
\pgfpathlineto{\pgfqpoint{5.051884in}{0.690890in}}%
\pgfpathlineto{\pgfqpoint{5.055515in}{0.691617in}}%
\pgfpathlineto{\pgfqpoint{5.070040in}{0.666522in}}%
\pgfpathlineto{\pgfqpoint{5.073671in}{0.686914in}}%
\pgfpathlineto{\pgfqpoint{5.077302in}{0.659546in}}%
\pgfpathlineto{\pgfqpoint{5.080934in}{0.652455in}}%
\pgfpathlineto{\pgfqpoint{5.084565in}{0.664901in}}%
\pgfpathlineto{\pgfqpoint{5.088196in}{0.684206in}}%
\pgfpathlineto{\pgfqpoint{5.091827in}{0.698084in}}%
\pgfpathlineto{\pgfqpoint{5.095458in}{0.683798in}}%
\pgfpathlineto{\pgfqpoint{5.099090in}{0.675337in}}%
\pgfpathlineto{\pgfqpoint{5.102721in}{0.671064in}}%
\pgfpathlineto{\pgfqpoint{5.106352in}{0.637419in}}%
\pgfpathlineto{\pgfqpoint{5.109983in}{0.637250in}}%
\pgfpathlineto{\pgfqpoint{5.113614in}{0.634152in}}%
\pgfpathlineto{\pgfqpoint{5.117245in}{0.667790in}}%
\pgfpathlineto{\pgfqpoint{5.120877in}{0.685209in}}%
\pgfpathlineto{\pgfqpoint{5.124508in}{0.691077in}}%
\pgfpathlineto{\pgfqpoint{5.128139in}{0.694957in}}%
\pgfpathlineto{\pgfqpoint{5.131770in}{0.594901in}}%
\pgfpathlineto{\pgfqpoint{5.135401in}{0.599077in}}%
\pgfpathlineto{\pgfqpoint{5.139032in}{0.671762in}}%
\pgfpathlineto{\pgfqpoint{5.142664in}{0.581894in}}%
\pgfpathlineto{\pgfqpoint{5.146295in}{0.642882in}}%
\pgfpathlineto{\pgfqpoint{5.149926in}{0.669704in}}%
\pgfpathlineto{\pgfqpoint{5.153557in}{0.687089in}}%
\pgfpathlineto{\pgfqpoint{5.157188in}{0.696893in}}%
\pgfpathlineto{\pgfqpoint{5.160819in}{0.584140in}}%
\pgfpathlineto{\pgfqpoint{5.164451in}{0.602467in}}%
\pgfpathlineto{\pgfqpoint{5.168082in}{0.653833in}}%
\pgfpathlineto{\pgfqpoint{5.171713in}{0.615965in}}%
\pgfpathlineto{\pgfqpoint{5.175344in}{0.631067in}}%
\pgfpathlineto{\pgfqpoint{5.178975in}{0.650922in}}%
\pgfpathlineto{\pgfqpoint{5.182606in}{0.627836in}}%
\pgfpathlineto{\pgfqpoint{5.186238in}{0.674616in}}%
\pgfpathlineto{\pgfqpoint{5.189869in}{0.565484in}}%
\pgfpathlineto{\pgfqpoint{5.193500in}{0.639949in}}%
\pgfpathlineto{\pgfqpoint{5.197131in}{0.668903in}}%
\pgfpathlineto{\pgfqpoint{5.200762in}{0.683711in}}%
\pgfpathlineto{\pgfqpoint{5.204394in}{0.691991in}}%
\pgfpathlineto{\pgfqpoint{5.208025in}{0.590494in}}%
\pgfpathlineto{\pgfqpoint{5.211656in}{0.611002in}}%
\pgfpathlineto{\pgfqpoint{5.215287in}{0.657252in}}%
\pgfpathlineto{\pgfqpoint{5.218918in}{0.619631in}}%
\pgfpathlineto{\pgfqpoint{5.222549in}{0.624983in}}%
\pgfpathlineto{\pgfqpoint{5.226181in}{0.653730in}}%
\pgfpathlineto{\pgfqpoint{5.229812in}{0.634654in}}%
\pgfpathlineto{\pgfqpoint{5.233443in}{0.674381in}}%
\pgfpathlineto{\pgfqpoint{5.237074in}{0.563689in}}%
\pgfpathlineto{\pgfqpoint{5.240705in}{0.640401in}}%
\pgfpathlineto{\pgfqpoint{5.244336in}{0.670219in}}%
\pgfpathlineto{\pgfqpoint{5.247968in}{0.681627in}}%
\pgfpathlineto{\pgfqpoint{5.251599in}{0.688485in}}%
\pgfpathlineto{\pgfqpoint{5.255230in}{0.600397in}}%
\pgfpathlineto{\pgfqpoint{5.258861in}{0.617733in}}%
\pgfpathlineto{\pgfqpoint{5.262492in}{0.659360in}}%
\pgfpathlineto{\pgfqpoint{5.266123in}{0.619895in}}%
\pgfpathlineto{\pgfqpoint{5.269755in}{0.623927in}}%
\pgfpathlineto{\pgfqpoint{5.273386in}{0.657762in}}%
\pgfpathlineto{\pgfqpoint{5.277017in}{0.640491in}}%
\pgfpathlineto{\pgfqpoint{5.280648in}{0.673565in}}%
\pgfpathlineto{\pgfqpoint{5.284279in}{0.568699in}}%
\pgfpathlineto{\pgfqpoint{5.287910in}{0.641563in}}%
\pgfpathlineto{\pgfqpoint{5.291542in}{0.671647in}}%
\pgfpathlineto{\pgfqpoint{5.295173in}{0.679592in}}%
\pgfpathlineto{\pgfqpoint{5.298804in}{0.683839in}}%
\pgfpathlineto{\pgfqpoint{5.302435in}{0.634163in}}%
\pgfpathlineto{\pgfqpoint{5.306066in}{0.497551in}}%
\pgfpathlineto{\pgfqpoint{5.309698in}{0.690081in}}%
\pgfpathlineto{\pgfqpoint{5.313329in}{0.604934in}}%
\pgfpathlineto{\pgfqpoint{5.316960in}{0.565322in}}%
\pgfpathlineto{\pgfqpoint{5.320591in}{0.658039in}}%
\pgfpathlineto{\pgfqpoint{5.324222in}{0.670219in}}%
\pgfpathlineto{\pgfqpoint{5.327853in}{0.678914in}}%
\pgfpathlineto{\pgfqpoint{5.331485in}{0.669071in}}%
\pgfpathlineto{\pgfqpoint{5.335116in}{0.644427in}}%
\pgfpathlineto{\pgfqpoint{5.338747in}{0.680135in}}%
\pgfpathlineto{\pgfqpoint{5.342378in}{0.633687in}}%
\pgfpathlineto{\pgfqpoint{5.346009in}{0.643397in}}%
\pgfpathlineto{\pgfqpoint{5.349640in}{0.639756in}}%
\pgfpathlineto{\pgfqpoint{5.356903in}{0.669126in}}%
\pgfpathlineto{\pgfqpoint{5.360534in}{0.673126in}}%
\pgfpathlineto{\pgfqpoint{5.364165in}{0.670912in}}%
\pgfpathlineto{\pgfqpoint{5.367796in}{0.659431in}}%
\pgfpathlineto{\pgfqpoint{5.371427in}{0.623409in}}%
\pgfpathlineto{\pgfqpoint{5.375059in}{0.616972in}}%
\pgfpathlineto{\pgfqpoint{5.378690in}{0.600111in}}%
\pgfpathlineto{\pgfqpoint{5.382321in}{0.672416in}}%
\pgfpathlineto{\pgfqpoint{5.385952in}{0.688203in}}%
\pgfpathlineto{\pgfqpoint{5.389583in}{0.645435in}}%
\pgfpathlineto{\pgfqpoint{5.393215in}{0.647434in}}%
\pgfpathlineto{\pgfqpoint{5.396846in}{0.640898in}}%
\pgfpathlineto{\pgfqpoint{5.404108in}{0.666105in}}%
\pgfpathlineto{\pgfqpoint{5.407739in}{0.667328in}}%
\pgfpathlineto{\pgfqpoint{5.411370in}{0.662123in}}%
\pgfpathlineto{\pgfqpoint{5.415002in}{0.661463in}}%
\pgfpathlineto{\pgfqpoint{5.418633in}{0.628795in}}%
\pgfpathlineto{\pgfqpoint{5.422264in}{0.622373in}}%
\pgfpathlineto{\pgfqpoint{5.425895in}{0.599313in}}%
\pgfpathlineto{\pgfqpoint{5.429526in}{0.670509in}}%
\pgfpathlineto{\pgfqpoint{5.433157in}{0.688611in}}%
\pgfpathlineto{\pgfqpoint{5.436789in}{0.647182in}}%
\pgfpathlineto{\pgfqpoint{5.440420in}{0.645544in}}%
\pgfpathlineto{\pgfqpoint{5.444051in}{0.642870in}}%
\pgfpathlineto{\pgfqpoint{5.447682in}{0.653830in}}%
\pgfpathlineto{\pgfqpoint{5.451313in}{0.666876in}}%
\pgfpathlineto{\pgfqpoint{5.454944in}{0.663519in}}%
\pgfpathlineto{\pgfqpoint{5.458576in}{0.654703in}}%
\pgfpathlineto{\pgfqpoint{5.462207in}{0.663859in}}%
\pgfpathlineto{\pgfqpoint{5.465838in}{0.632280in}}%
\pgfpathlineto{\pgfqpoint{5.469469in}{0.624562in}}%
\pgfpathlineto{\pgfqpoint{5.473100in}{0.603188in}}%
\pgfpathlineto{\pgfqpoint{5.476731in}{0.669693in}}%
\pgfpathlineto{\pgfqpoint{5.480363in}{0.689310in}}%
\pgfpathlineto{\pgfqpoint{5.483994in}{0.690431in}}%
\pgfpathlineto{\pgfqpoint{5.487625in}{0.698984in}}%
\pgfpathlineto{\pgfqpoint{5.491256in}{0.605548in}}%
\pgfpathlineto{\pgfqpoint{5.494887in}{0.606000in}}%
\pgfpathlineto{\pgfqpoint{5.498519in}{0.659075in}}%
\pgfpathlineto{\pgfqpoint{5.502150in}{0.513199in}}%
\pgfpathlineto{\pgfqpoint{5.505781in}{0.648915in}}%
\pgfpathlineto{\pgfqpoint{5.513043in}{0.693186in}}%
\pgfpathlineto{\pgfqpoint{5.516674in}{0.708001in}}%
\pgfpathlineto{\pgfqpoint{5.520306in}{0.657325in}}%
\pgfpathlineto{\pgfqpoint{5.523937in}{0.632033in}}%
\pgfpathlineto{\pgfqpoint{5.527568in}{0.676481in}}%
\pgfpathlineto{\pgfqpoint{5.531199in}{0.603308in}}%
\pgfpathlineto{\pgfqpoint{5.534830in}{0.599905in}}%
\pgfpathlineto{\pgfqpoint{5.538461in}{0.652797in}}%
\pgfpathlineto{\pgfqpoint{5.542093in}{0.686314in}}%
\pgfpathlineto{\pgfqpoint{5.545724in}{0.701133in}}%
\pgfpathlineto{\pgfqpoint{5.549355in}{0.623244in}}%
\pgfpathlineto{\pgfqpoint{5.552986in}{0.578913in}}%
\pgfpathlineto{\pgfqpoint{5.556617in}{0.658259in}}%
\pgfpathlineto{\pgfqpoint{5.560248in}{0.644506in}}%
\pgfpathlineto{\pgfqpoint{5.563880in}{0.614752in}}%
\pgfpathlineto{\pgfqpoint{5.567511in}{0.601996in}}%
\pgfpathlineto{\pgfqpoint{5.571142in}{0.579751in}}%
\pgfpathlineto{\pgfqpoint{5.574773in}{0.678686in}}%
\pgfpathlineto{\pgfqpoint{5.578404in}{0.604488in}}%
\pgfpathlineto{\pgfqpoint{5.582035in}{0.594085in}}%
\pgfpathlineto{\pgfqpoint{5.585667in}{0.653531in}}%
\pgfpathlineto{\pgfqpoint{5.589298in}{0.684193in}}%
\pgfpathlineto{\pgfqpoint{5.592929in}{0.698983in}}%
\pgfpathlineto{\pgfqpoint{5.596560in}{0.622201in}}%
\pgfpathlineto{\pgfqpoint{5.600191in}{0.557487in}}%
\pgfpathlineto{\pgfqpoint{5.603823in}{0.660658in}}%
\pgfpathlineto{\pgfqpoint{5.607454in}{0.647110in}}%
\pgfpathlineto{\pgfqpoint{5.611085in}{0.622459in}}%
\pgfpathlineto{\pgfqpoint{5.614716in}{0.587862in}}%
\pgfpathlineto{\pgfqpoint{5.618347in}{0.507662in}}%
\pgfpathlineto{\pgfqpoint{5.621978in}{0.679514in}}%
\pgfpathlineto{\pgfqpoint{5.625610in}{0.608152in}}%
\pgfpathlineto{\pgfqpoint{5.629241in}{0.591888in}}%
\pgfpathlineto{\pgfqpoint{5.632872in}{0.655054in}}%
\pgfpathlineto{\pgfqpoint{5.636503in}{0.682661in}}%
\pgfpathlineto{\pgfqpoint{5.640134in}{0.697719in}}%
\pgfpathlineto{\pgfqpoint{5.643765in}{0.619793in}}%
\pgfpathlineto{\pgfqpoint{5.647397in}{0.517190in}}%
\pgfpathlineto{\pgfqpoint{5.651028in}{0.662736in}}%
\pgfpathlineto{\pgfqpoint{5.654659in}{0.648035in}}%
\pgfpathlineto{\pgfqpoint{5.658290in}{0.623262in}}%
\pgfpathlineto{\pgfqpoint{5.661921in}{0.551596in}}%
\pgfpathlineto{\pgfqpoint{5.665552in}{0.564108in}}%
\pgfpathlineto{\pgfqpoint{5.669184in}{0.680264in}}%
\pgfpathlineto{\pgfqpoint{5.672815in}{0.610013in}}%
\pgfpathlineto{\pgfqpoint{5.676446in}{0.593030in}}%
\pgfpathlineto{\pgfqpoint{5.680077in}{0.657094in}}%
\pgfpathlineto{\pgfqpoint{5.683708in}{0.681412in}}%
\pgfpathlineto{\pgfqpoint{5.687340in}{0.696422in}}%
\pgfpathlineto{\pgfqpoint{5.690971in}{0.677889in}}%
\pgfpathlineto{\pgfqpoint{5.694602in}{0.664762in}}%
\pgfpathlineto{\pgfqpoint{5.698233in}{0.673656in}}%
\pgfpathlineto{\pgfqpoint{5.701864in}{0.627608in}}%
\pgfpathlineto{\pgfqpoint{5.705495in}{0.621700in}}%
\pgfpathlineto{\pgfqpoint{5.709127in}{0.633436in}}%
\pgfpathlineto{\pgfqpoint{5.712758in}{0.667977in}}%
\pgfpathlineto{\pgfqpoint{5.716389in}{0.685456in}}%
\pgfpathlineto{\pgfqpoint{5.720020in}{0.688659in}}%
\pgfpathlineto{\pgfqpoint{5.723651in}{0.689751in}}%
\pgfpathlineto{\pgfqpoint{5.727282in}{0.617239in}}%
\pgfpathlineto{\pgfqpoint{5.730914in}{0.607406in}}%
\pgfpathlineto{\pgfqpoint{5.734545in}{0.668983in}}%
\pgfpathlineto{\pgfqpoint{5.738176in}{0.585386in}}%
\pgfpathlineto{\pgfqpoint{5.741807in}{0.645449in}}%
\pgfpathlineto{\pgfqpoint{5.745438in}{0.671061in}}%
\pgfpathlineto{\pgfqpoint{5.749069in}{0.685778in}}%
\pgfpathlineto{\pgfqpoint{5.752701in}{0.693823in}}%
\pgfpathlineto{\pgfqpoint{5.756332in}{0.604658in}}%
\pgfpathlineto{\pgfqpoint{5.759963in}{0.609419in}}%
\pgfpathlineto{\pgfqpoint{5.763594in}{0.651553in}}%
\pgfpathlineto{\pgfqpoint{5.767225in}{0.606287in}}%
\pgfpathlineto{\pgfqpoint{5.770856in}{0.640214in}}%
\pgfpathlineto{\pgfqpoint{5.774488in}{0.659581in}}%
\pgfpathlineto{\pgfqpoint{5.778119in}{0.634569in}}%
\pgfpathlineto{\pgfqpoint{5.781750in}{0.671985in}}%
\pgfpathlineto{\pgfqpoint{5.785381in}{0.586454in}}%
\pgfpathlineto{\pgfqpoint{5.789012in}{0.643270in}}%
\pgfpathlineto{\pgfqpoint{5.792644in}{0.670178in}}%
\pgfpathlineto{\pgfqpoint{5.796275in}{0.682232in}}%
\pgfpathlineto{\pgfqpoint{5.799906in}{0.688574in}}%
\pgfpathlineto{\pgfqpoint{5.803537in}{0.611491in}}%
\pgfpathlineto{\pgfqpoint{5.807168in}{0.617609in}}%
\pgfpathlineto{\pgfqpoint{5.810799in}{0.654987in}}%
\pgfpathlineto{\pgfqpoint{5.814431in}{0.610361in}}%
\pgfpathlineto{\pgfqpoint{5.821693in}{0.661947in}}%
\pgfpathlineto{\pgfqpoint{5.825324in}{0.639976in}}%
\pgfpathlineto{\pgfqpoint{5.828955in}{0.671540in}}%
\pgfpathlineto{\pgfqpoint{5.832586in}{0.587874in}}%
\pgfpathlineto{\pgfqpoint{5.836218in}{0.643955in}}%
\pgfpathlineto{\pgfqpoint{5.839849in}{0.671456in}}%
\pgfpathlineto{\pgfqpoint{5.843480in}{0.680010in}}%
\pgfpathlineto{\pgfqpoint{5.847111in}{0.684695in}}%
\pgfpathlineto{\pgfqpoint{5.850742in}{0.618568in}}%
\pgfpathlineto{\pgfqpoint{5.854373in}{0.623403in}}%
\pgfpathlineto{\pgfqpoint{5.858005in}{0.656950in}}%
\pgfpathlineto{\pgfqpoint{5.861636in}{0.609891in}}%
\pgfpathlineto{\pgfqpoint{5.865267in}{0.635532in}}%
\pgfpathlineto{\pgfqpoint{5.868898in}{0.665208in}}%
\pgfpathlineto{\pgfqpoint{5.876160in}{0.707765in}}%
\pgfpathlineto{\pgfqpoint{5.879792in}{0.634256in}}%
\pgfpathlineto{\pgfqpoint{5.883423in}{0.592682in}}%
\pgfpathlineto{\pgfqpoint{5.887054in}{0.667684in}}%
\pgfpathlineto{\pgfqpoint{5.894316in}{0.583415in}}%
\pgfpathlineto{\pgfqpoint{5.897948in}{0.639997in}}%
\pgfpathlineto{\pgfqpoint{5.901579in}{0.684191in}}%
\pgfpathlineto{\pgfqpoint{5.905210in}{0.703717in}}%
\pgfpathlineto{\pgfqpoint{5.908841in}{0.678370in}}%
\pgfpathlineto{\pgfqpoint{5.912472in}{0.675902in}}%
\pgfpathlineto{\pgfqpoint{5.916103in}{0.661865in}}%
\pgfpathlineto{\pgfqpoint{5.919735in}{0.625179in}}%
\pgfpathlineto{\pgfqpoint{5.923366in}{0.606907in}}%
\pgfpathlineto{\pgfqpoint{5.926997in}{0.604769in}}%
\pgfpathlineto{\pgfqpoint{5.930628in}{0.674677in}}%
\pgfpathlineto{\pgfqpoint{5.934259in}{0.690585in}}%
\pgfpathlineto{\pgfqpoint{5.937890in}{0.642838in}}%
\pgfpathlineto{\pgfqpoint{5.941522in}{0.645050in}}%
\pgfpathlineto{\pgfqpoint{5.945153in}{0.643440in}}%
\pgfpathlineto{\pgfqpoint{5.948784in}{0.651292in}}%
\pgfpathlineto{\pgfqpoint{5.952415in}{0.660984in}}%
\pgfpathlineto{\pgfqpoint{5.956046in}{0.665460in}}%
\pgfpathlineto{\pgfqpoint{5.959677in}{0.661041in}}%
\pgfpathlineto{\pgfqpoint{5.963309in}{0.663618in}}%
\pgfpathlineto{\pgfqpoint{5.966940in}{0.625573in}}%
\pgfpathlineto{\pgfqpoint{5.970571in}{0.610853in}}%
\pgfpathlineto{\pgfqpoint{5.974202in}{0.610053in}}%
\pgfpathlineto{\pgfqpoint{5.977833in}{0.673922in}}%
\pgfpathlineto{\pgfqpoint{5.981465in}{0.690486in}}%
\pgfpathlineto{\pgfqpoint{5.985096in}{0.645288in}}%
\pgfpathlineto{\pgfqpoint{5.988727in}{0.642260in}}%
\pgfpathlineto{\pgfqpoint{5.992358in}{0.645033in}}%
\pgfpathlineto{\pgfqpoint{5.995989in}{0.653244in}}%
\pgfpathlineto{\pgfqpoint{5.999620in}{0.663804in}}%
\pgfpathlineto{\pgfqpoint{6.003252in}{0.662551in}}%
\pgfpathlineto{\pgfqpoint{6.006883in}{0.654469in}}%
\pgfpathlineto{\pgfqpoint{6.010514in}{0.665469in}}%
\pgfpathlineto{\pgfqpoint{6.014145in}{0.629310in}}%
\pgfpathlineto{\pgfqpoint{6.017776in}{0.614966in}}%
\pgfpathlineto{\pgfqpoint{6.021407in}{0.612074in}}%
\pgfpathlineto{\pgfqpoint{6.025039in}{0.672649in}}%
\pgfpathlineto{\pgfqpoint{6.028670in}{0.690785in}}%
\pgfpathlineto{\pgfqpoint{6.032301in}{0.646523in}}%
\pgfpathlineto{\pgfqpoint{6.035932in}{0.639336in}}%
\pgfpathlineto{\pgfqpoint{6.043194in}{0.654296in}}%
\pgfpathlineto{\pgfqpoint{6.046826in}{0.664607in}}%
\pgfpathlineto{\pgfqpoint{6.050457in}{0.658312in}}%
\pgfpathlineto{\pgfqpoint{6.054088in}{0.645805in}}%
\pgfpathlineto{\pgfqpoint{6.057719in}{0.667691in}}%
\pgfpathlineto{\pgfqpoint{6.061350in}{0.632364in}}%
\pgfpathlineto{\pgfqpoint{6.064981in}{0.616587in}}%
\pgfpathlineto{\pgfqpoint{6.068613in}{0.615869in}}%
\pgfpathlineto{\pgfqpoint{6.072244in}{0.671954in}}%
\pgfpathlineto{\pgfqpoint{6.075875in}{0.691293in}}%
\pgfpathlineto{\pgfqpoint{6.079506in}{0.689324in}}%
\pgfpathlineto{\pgfqpoint{6.083137in}{0.694958in}}%
\pgfpathlineto{\pgfqpoint{6.086769in}{0.625048in}}%
\pgfpathlineto{\pgfqpoint{6.090400in}{0.614825in}}%
\pgfpathlineto{\pgfqpoint{6.094031in}{0.656497in}}%
\pgfpathlineto{\pgfqpoint{6.097662in}{0.520401in}}%
\pgfpathlineto{\pgfqpoint{6.101293in}{0.652634in}}%
\pgfpathlineto{\pgfqpoint{6.104924in}{0.675788in}}%
\pgfpathlineto{\pgfqpoint{6.108556in}{0.692924in}}%
\pgfpathlineto{\pgfqpoint{6.112187in}{0.706508in}}%
\pgfpathlineto{\pgfqpoint{6.115818in}{0.648481in}}%
\pgfpathlineto{\pgfqpoint{6.119449in}{0.615533in}}%
\pgfpathlineto{\pgfqpoint{6.123080in}{0.677294in}}%
\pgfpathlineto{\pgfqpoint{6.126711in}{0.596405in}}%
\pgfpathlineto{\pgfqpoint{6.130343in}{0.610704in}}%
\pgfpathlineto{\pgfqpoint{6.133974in}{0.656837in}}%
\pgfpathlineto{\pgfqpoint{6.137605in}{0.686664in}}%
\pgfpathlineto{\pgfqpoint{6.141236in}{0.701246in}}%
\pgfpathlineto{\pgfqpoint{6.144867in}{0.615147in}}%
\pgfpathlineto{\pgfqpoint{6.148498in}{0.542527in}}%
\pgfpathlineto{\pgfqpoint{6.152130in}{0.659463in}}%
\pgfpathlineto{\pgfqpoint{6.155761in}{0.642329in}}%
\pgfpathlineto{\pgfqpoint{6.163023in}{0.550335in}}%
\pgfpathlineto{\pgfqpoint{6.166654in}{0.423873in}}%
\pgfpathlineto{\pgfqpoint{6.170285in}{0.679573in}}%
\pgfpathlineto{\pgfqpoint{6.173917in}{0.597141in}}%
\pgfpathlineto{\pgfqpoint{6.177548in}{0.606785in}}%
\pgfpathlineto{\pgfqpoint{6.181179in}{0.657442in}}%
\pgfpathlineto{\pgfqpoint{6.184810in}{0.684447in}}%
\pgfpathlineto{\pgfqpoint{6.188441in}{0.698758in}}%
\pgfpathlineto{\pgfqpoint{6.192073in}{0.612532in}}%
\pgfpathlineto{\pgfqpoint{6.195704in}{0.323073in}}%
\pgfpathlineto{\pgfqpoint{6.199335in}{0.661938in}}%
\pgfpathlineto{\pgfqpoint{6.202966in}{0.644869in}}%
\pgfpathlineto{\pgfqpoint{6.206597in}{0.606569in}}%
\pgfpathlineto{\pgfqpoint{6.210228in}{0.334161in}}%
\pgfpathlineto{\pgfqpoint{6.213860in}{0.567923in}}%
\pgfpathlineto{\pgfqpoint{6.217491in}{0.680237in}}%
\pgfpathlineto{\pgfqpoint{6.221122in}{0.600424in}}%
\pgfpathlineto{\pgfqpoint{6.224753in}{0.606153in}}%
\pgfpathlineto{\pgfqpoint{6.228384in}{0.659025in}}%
\pgfpathlineto{\pgfqpoint{6.232015in}{0.682904in}}%
\pgfpathlineto{\pgfqpoint{6.235647in}{0.697172in}}%
\pgfpathlineto{\pgfqpoint{6.239278in}{0.608156in}}%
\pgfpathlineto{\pgfqpoint{6.242909in}{0.539289in}}%
\pgfpathlineto{\pgfqpoint{6.246540in}{0.663979in}}%
\pgfpathlineto{\pgfqpoint{6.250171in}{0.645612in}}%
\pgfpathlineto{\pgfqpoint{6.257434in}{0.565486in}}%
\pgfpathlineto{\pgfqpoint{6.261065in}{0.662448in}}%
\pgfpathlineto{\pgfqpoint{6.264696in}{0.699145in}}%
\pgfpathlineto{\pgfqpoint{6.268327in}{0.659147in}}%
\pgfpathlineto{\pgfqpoint{6.271958in}{0.655668in}}%
\pgfpathlineto{\pgfqpoint{6.275590in}{0.656195in}}%
\pgfpathlineto{\pgfqpoint{6.279221in}{0.639300in}}%
\pgfpathlineto{\pgfqpoint{6.282852in}{0.630359in}}%
\pgfpathlineto{\pgfqpoint{6.286483in}{0.582976in}}%
\pgfpathlineto{\pgfqpoint{6.290114in}{0.656055in}}%
\pgfpathlineto{\pgfqpoint{6.293745in}{0.684052in}}%
\pgfpathlineto{\pgfqpoint{6.297377in}{0.678379in}}%
\pgfpathlineto{\pgfqpoint{6.301008in}{0.697333in}}%
\pgfpathlineto{\pgfqpoint{6.304639in}{0.620041in}}%
\pgfpathlineto{\pgfqpoint{6.308270in}{0.618361in}}%
\pgfpathlineto{\pgfqpoint{6.311901in}{0.648918in}}%
\pgfpathlineto{\pgfqpoint{6.315532in}{0.607462in}}%
\pgfpathlineto{\pgfqpoint{6.319164in}{0.641865in}}%
\pgfpathlineto{\pgfqpoint{6.322795in}{0.661705in}}%
\pgfpathlineto{\pgfqpoint{6.326426in}{0.631130in}}%
\pgfpathlineto{\pgfqpoint{6.330057in}{0.670651in}}%
\pgfpathlineto{\pgfqpoint{6.333688in}{0.600354in}}%
\pgfpathlineto{\pgfqpoint{6.337319in}{0.644390in}}%
\pgfpathlineto{\pgfqpoint{6.340951in}{0.668881in}}%
\pgfpathlineto{\pgfqpoint{6.344582in}{0.682295in}}%
\pgfpathlineto{\pgfqpoint{6.348213in}{0.688291in}}%
\pgfpathlineto{\pgfqpoint{6.351844in}{0.619579in}}%
\pgfpathlineto{\pgfqpoint{6.355475in}{0.617214in}}%
\pgfpathlineto{\pgfqpoint{6.359106in}{0.650476in}}%
\pgfpathlineto{\pgfqpoint{6.362738in}{0.600012in}}%
\pgfpathlineto{\pgfqpoint{6.366369in}{0.597087in}}%
\pgfpathlineto{\pgfqpoint{6.370000in}{0.629003in}}%
\pgfpathlineto{\pgfqpoint{6.370000in}{0.629003in}}%
\pgfusepath{stroke}%
\end{pgfscope}%
\begin{pgfscope}%
\pgfpathrectangle{\pgfqpoint{4.554414in}{0.320625in}}{\pgfqpoint{1.815586in}{1.164375in}}%
\pgfusepath{clip}%
\pgfsetrectcap%
\pgfsetroundjoin%
\pgfsetlinewidth{0.602250pt}%
\definecolor{currentstroke}{rgb}{0.819608,0.286275,0.356863}%
\pgfsetstrokecolor{currentstroke}%
\pgfsetstrokeopacity{0.120000}%
\pgfsetdash{}{0pt}%
\pgfpathmoveto{\pgfqpoint{4.554414in}{1.465642in}}%
\pgfpathlineto{\pgfqpoint{4.561676in}{1.443683in}}%
\pgfpathlineto{\pgfqpoint{4.568938in}{1.427584in}}%
\pgfpathlineto{\pgfqpoint{4.608881in}{1.352019in}}%
\pgfpathlineto{\pgfqpoint{4.637931in}{1.298297in}}%
\pgfpathlineto{\pgfqpoint{4.663349in}{1.253791in}}%
\pgfpathlineto{\pgfqpoint{4.688767in}{1.211548in}}%
\pgfpathlineto{\pgfqpoint{4.717816in}{1.165492in}}%
\pgfpathlineto{\pgfqpoint{4.750497in}{1.115812in}}%
\pgfpathlineto{\pgfqpoint{4.786809in}{1.062569in}}%
\pgfpathlineto{\pgfqpoint{4.830383in}{1.000733in}}%
\pgfpathlineto{\pgfqpoint{4.877588in}{0.935764in}}%
\pgfpathlineto{\pgfqpoint{4.928424in}{0.867723in}}%
\pgfpathlineto{\pgfqpoint{4.986523in}{0.791920in}}%
\pgfpathlineto{\pgfqpoint{5.033728in}{0.731536in}}%
\pgfpathlineto{\pgfqpoint{5.037360in}{0.728270in}}%
\pgfpathlineto{\pgfqpoint{5.044622in}{0.717647in}}%
\pgfpathlineto{\pgfqpoint{5.048253in}{0.714209in}}%
\pgfpathlineto{\pgfqpoint{5.051884in}{0.708446in}}%
\pgfpathlineto{\pgfqpoint{5.055515in}{0.704550in}}%
\pgfpathlineto{\pgfqpoint{5.059147in}{0.701817in}}%
\pgfpathlineto{\pgfqpoint{5.062778in}{0.695371in}}%
\pgfpathlineto{\pgfqpoint{5.066409in}{0.690701in}}%
\pgfpathlineto{\pgfqpoint{5.070040in}{0.695877in}}%
\pgfpathlineto{\pgfqpoint{5.073671in}{0.681031in}}%
\pgfpathlineto{\pgfqpoint{5.077302in}{0.676910in}}%
\pgfpathlineto{\pgfqpoint{5.080934in}{0.674468in}}%
\pgfpathlineto{\pgfqpoint{5.084565in}{0.667792in}}%
\pgfpathlineto{\pgfqpoint{5.091827in}{0.676615in}}%
\pgfpathlineto{\pgfqpoint{5.095458in}{0.655800in}}%
\pgfpathlineto{\pgfqpoint{5.099090in}{0.653022in}}%
\pgfpathlineto{\pgfqpoint{5.102721in}{0.672851in}}%
\pgfpathlineto{\pgfqpoint{5.106352in}{0.643587in}}%
\pgfpathlineto{\pgfqpoint{5.109983in}{0.642834in}}%
\pgfpathlineto{\pgfqpoint{5.113614in}{0.633529in}}%
\pgfpathlineto{\pgfqpoint{5.117245in}{0.638266in}}%
\pgfpathlineto{\pgfqpoint{5.120877in}{0.665747in}}%
\pgfpathlineto{\pgfqpoint{5.124508in}{0.675398in}}%
\pgfpathlineto{\pgfqpoint{5.128139in}{0.613482in}}%
\pgfpathlineto{\pgfqpoint{5.135401in}{0.636917in}}%
\pgfpathlineto{\pgfqpoint{5.139032in}{0.634186in}}%
\pgfpathlineto{\pgfqpoint{5.142664in}{0.628423in}}%
\pgfpathlineto{\pgfqpoint{5.146295in}{0.590392in}}%
\pgfpathlineto{\pgfqpoint{5.149926in}{0.641097in}}%
\pgfpathlineto{\pgfqpoint{5.153557in}{0.629195in}}%
\pgfpathlineto{\pgfqpoint{5.157188in}{0.615384in}}%
\pgfpathlineto{\pgfqpoint{5.160819in}{0.638337in}}%
\pgfpathlineto{\pgfqpoint{5.164451in}{0.624446in}}%
\pgfpathlineto{\pgfqpoint{5.168082in}{0.584522in}}%
\pgfpathlineto{\pgfqpoint{5.171713in}{0.651705in}}%
\pgfpathlineto{\pgfqpoint{5.175344in}{0.664017in}}%
\pgfpathlineto{\pgfqpoint{5.178975in}{0.658419in}}%
\pgfpathlineto{\pgfqpoint{5.182606in}{0.650133in}}%
\pgfpathlineto{\pgfqpoint{5.186238in}{0.600911in}}%
\pgfpathlineto{\pgfqpoint{5.189869in}{0.571831in}}%
\pgfpathlineto{\pgfqpoint{5.193500in}{0.642392in}}%
\pgfpathlineto{\pgfqpoint{5.197131in}{0.629267in}}%
\pgfpathlineto{\pgfqpoint{5.200762in}{0.599686in}}%
\pgfpathlineto{\pgfqpoint{5.204394in}{0.665893in}}%
\pgfpathlineto{\pgfqpoint{5.208025in}{0.653085in}}%
\pgfpathlineto{\pgfqpoint{5.211656in}{0.630177in}}%
\pgfpathlineto{\pgfqpoint{5.215287in}{0.653322in}}%
\pgfpathlineto{\pgfqpoint{5.218918in}{0.534210in}}%
\pgfpathlineto{\pgfqpoint{5.222549in}{0.623563in}}%
\pgfpathlineto{\pgfqpoint{5.226181in}{0.649848in}}%
\pgfpathlineto{\pgfqpoint{5.229812in}{0.640965in}}%
\pgfpathlineto{\pgfqpoint{5.233443in}{0.595453in}}%
\pgfpathlineto{\pgfqpoint{5.237074in}{0.656835in}}%
\pgfpathlineto{\pgfqpoint{5.240705in}{0.537596in}}%
\pgfpathlineto{\pgfqpoint{5.244336in}{0.624841in}}%
\pgfpathlineto{\pgfqpoint{5.247968in}{0.640718in}}%
\pgfpathlineto{\pgfqpoint{5.251599in}{0.618419in}}%
\pgfpathlineto{\pgfqpoint{5.255230in}{0.501050in}}%
\pgfpathlineto{\pgfqpoint{5.258861in}{0.654512in}}%
\pgfpathlineto{\pgfqpoint{5.262492in}{0.578238in}}%
\pgfpathlineto{\pgfqpoint{5.266123in}{0.629059in}}%
\pgfpathlineto{\pgfqpoint{5.269755in}{0.635533in}}%
\pgfpathlineto{\pgfqpoint{5.273386in}{0.585960in}}%
\pgfpathlineto{\pgfqpoint{5.277017in}{0.587253in}}%
\pgfpathlineto{\pgfqpoint{5.280648in}{0.649574in}}%
\pgfpathlineto{\pgfqpoint{5.284279in}{0.601159in}}%
\pgfpathlineto{\pgfqpoint{5.287910in}{0.630494in}}%
\pgfpathlineto{\pgfqpoint{5.291542in}{0.623591in}}%
\pgfpathlineto{\pgfqpoint{5.295173in}{0.540814in}}%
\pgfpathlineto{\pgfqpoint{5.298804in}{0.610445in}}%
\pgfpathlineto{\pgfqpoint{5.302435in}{0.642257in}}%
\pgfpathlineto{\pgfqpoint{5.306066in}{0.616462in}}%
\pgfpathlineto{\pgfqpoint{5.309698in}{0.630488in}}%
\pgfpathlineto{\pgfqpoint{5.313329in}{0.604110in}}%
\pgfpathlineto{\pgfqpoint{5.316960in}{0.608537in}}%
\pgfpathlineto{\pgfqpoint{5.320591in}{0.661802in}}%
\pgfpathlineto{\pgfqpoint{5.324222in}{0.680766in}}%
\pgfpathlineto{\pgfqpoint{5.327853in}{0.606172in}}%
\pgfpathlineto{\pgfqpoint{5.331485in}{0.602888in}}%
\pgfpathlineto{\pgfqpoint{5.335116in}{0.578225in}}%
\pgfpathlineto{\pgfqpoint{5.338747in}{0.632494in}}%
\pgfpathlineto{\pgfqpoint{5.342378in}{0.643648in}}%
\pgfpathlineto{\pgfqpoint{5.346009in}{0.612236in}}%
\pgfpathlineto{\pgfqpoint{5.349640in}{0.627923in}}%
\pgfpathlineto{\pgfqpoint{5.353272in}{0.653586in}}%
\pgfpathlineto{\pgfqpoint{5.356903in}{0.664379in}}%
\pgfpathlineto{\pgfqpoint{5.360534in}{0.640614in}}%
\pgfpathlineto{\pgfqpoint{5.364165in}{0.590960in}}%
\pgfpathlineto{\pgfqpoint{5.367796in}{0.596225in}}%
\pgfpathlineto{\pgfqpoint{5.371427in}{0.645333in}}%
\pgfpathlineto{\pgfqpoint{5.375059in}{0.659510in}}%
\pgfpathlineto{\pgfqpoint{5.378690in}{0.648877in}}%
\pgfpathlineto{\pgfqpoint{5.382321in}{0.648709in}}%
\pgfpathlineto{\pgfqpoint{5.385952in}{0.574119in}}%
\pgfpathlineto{\pgfqpoint{5.389583in}{0.606986in}}%
\pgfpathlineto{\pgfqpoint{5.393215in}{0.645379in}}%
\pgfpathlineto{\pgfqpoint{5.396846in}{0.650843in}}%
\pgfpathlineto{\pgfqpoint{5.400477in}{0.628652in}}%
\pgfpathlineto{\pgfqpoint{5.404108in}{0.653558in}}%
\pgfpathlineto{\pgfqpoint{5.407739in}{0.552884in}}%
\pgfpathlineto{\pgfqpoint{5.411370in}{0.616664in}}%
\pgfpathlineto{\pgfqpoint{5.415002in}{0.646981in}}%
\pgfpathlineto{\pgfqpoint{5.418633in}{0.642489in}}%
\pgfpathlineto{\pgfqpoint{5.422264in}{0.603526in}}%
\pgfpathlineto{\pgfqpoint{5.425895in}{0.656261in}}%
\pgfpathlineto{\pgfqpoint{5.429526in}{0.417123in}}%
\pgfpathlineto{\pgfqpoint{5.433157in}{0.623077in}}%
\pgfpathlineto{\pgfqpoint{5.436789in}{0.645155in}}%
\pgfpathlineto{\pgfqpoint{5.440420in}{0.628597in}}%
\pgfpathlineto{\pgfqpoint{5.444051in}{0.538555in}}%
\pgfpathlineto{\pgfqpoint{5.447682in}{0.656181in}}%
\pgfpathlineto{\pgfqpoint{5.451313in}{0.557195in}}%
\pgfpathlineto{\pgfqpoint{5.454944in}{0.627610in}}%
\pgfpathlineto{\pgfqpoint{5.458576in}{0.640302in}}%
\pgfpathlineto{\pgfqpoint{5.465838in}{0.570914in}}%
\pgfpathlineto{\pgfqpoint{5.469469in}{0.652611in}}%
\pgfpathlineto{\pgfqpoint{5.473100in}{0.590056in}}%
\pgfpathlineto{\pgfqpoint{5.476731in}{0.630279in}}%
\pgfpathlineto{\pgfqpoint{5.480363in}{0.631496in}}%
\pgfpathlineto{\pgfqpoint{5.483994in}{0.531931in}}%
\pgfpathlineto{\pgfqpoint{5.487625in}{0.604118in}}%
\pgfpathlineto{\pgfqpoint{5.491256in}{0.646316in}}%
\pgfpathlineto{\pgfqpoint{5.494887in}{0.609390in}}%
\pgfpathlineto{\pgfqpoint{5.498519in}{0.631202in}}%
\pgfpathlineto{\pgfqpoint{5.502150in}{0.616966in}}%
\pgfpathlineto{\pgfqpoint{5.505781in}{0.589054in}}%
\pgfpathlineto{\pgfqpoint{5.509412in}{0.658905in}}%
\pgfpathlineto{\pgfqpoint{5.513043in}{0.681929in}}%
\pgfpathlineto{\pgfqpoint{5.516674in}{0.587878in}}%
\pgfpathlineto{\pgfqpoint{5.520306in}{0.602324in}}%
\pgfpathlineto{\pgfqpoint{5.523937in}{0.593412in}}%
\pgfpathlineto{\pgfqpoint{5.527568in}{0.629273in}}%
\pgfpathlineto{\pgfqpoint{5.531199in}{0.645128in}}%
\pgfpathlineto{\pgfqpoint{5.534830in}{0.625203in}}%
\pgfpathlineto{\pgfqpoint{5.538461in}{0.613437in}}%
\pgfpathlineto{\pgfqpoint{5.542093in}{0.651205in}}%
\pgfpathlineto{\pgfqpoint{5.545724in}{0.667342in}}%
\pgfpathlineto{\pgfqpoint{5.549355in}{0.636186in}}%
\pgfpathlineto{\pgfqpoint{5.552986in}{0.596546in}}%
\pgfpathlineto{\pgfqpoint{5.556617in}{0.584219in}}%
\pgfpathlineto{\pgfqpoint{5.560248in}{0.643355in}}%
\pgfpathlineto{\pgfqpoint{5.563880in}{0.660806in}}%
\pgfpathlineto{\pgfqpoint{5.571142in}{0.645445in}}%
\pgfpathlineto{\pgfqpoint{5.574773in}{0.583729in}}%
\pgfpathlineto{\pgfqpoint{5.578404in}{0.600048in}}%
\pgfpathlineto{\pgfqpoint{5.582035in}{0.644868in}}%
\pgfpathlineto{\pgfqpoint{5.585667in}{0.653438in}}%
\pgfpathlineto{\pgfqpoint{5.589298in}{0.635555in}}%
\pgfpathlineto{\pgfqpoint{5.592929in}{0.651622in}}%
\pgfpathlineto{\pgfqpoint{5.596560in}{0.570384in}}%
\pgfpathlineto{\pgfqpoint{5.603823in}{0.647431in}}%
\pgfpathlineto{\pgfqpoint{5.607454in}{0.646401in}}%
\pgfpathlineto{\pgfqpoint{5.611085in}{0.614684in}}%
\pgfpathlineto{\pgfqpoint{5.614716in}{0.655500in}}%
\pgfpathlineto{\pgfqpoint{5.618347in}{0.533039in}}%
\pgfpathlineto{\pgfqpoint{5.621978in}{0.620123in}}%
\pgfpathlineto{\pgfqpoint{5.625610in}{0.646971in}}%
\pgfpathlineto{\pgfqpoint{5.629241in}{0.634886in}}%
\pgfpathlineto{\pgfqpoint{5.632872in}{0.572910in}}%
\pgfpathlineto{\pgfqpoint{5.636503in}{0.656756in}}%
\pgfpathlineto{\pgfqpoint{5.640134in}{0.521254in}}%
\pgfpathlineto{\pgfqpoint{5.643765in}{0.625856in}}%
\pgfpathlineto{\pgfqpoint{5.647397in}{0.643775in}}%
\pgfpathlineto{\pgfqpoint{5.651028in}{0.616551in}}%
\pgfpathlineto{\pgfqpoint{5.654659in}{0.542455in}}%
\pgfpathlineto{\pgfqpoint{5.658290in}{0.654890in}}%
\pgfpathlineto{\pgfqpoint{5.661921in}{0.576578in}}%
\pgfpathlineto{\pgfqpoint{5.665552in}{0.629516in}}%
\pgfpathlineto{\pgfqpoint{5.669184in}{0.636915in}}%
\pgfpathlineto{\pgfqpoint{5.672815in}{0.577551in}}%
\pgfpathlineto{\pgfqpoint{5.676446in}{0.647192in}}%
\pgfpathlineto{\pgfqpoint{5.680077in}{0.683748in}}%
\pgfpathlineto{\pgfqpoint{5.683708in}{0.568738in}}%
\pgfpathlineto{\pgfqpoint{5.687340in}{0.585937in}}%
\pgfpathlineto{\pgfqpoint{5.690971in}{0.618473in}}%
\pgfpathlineto{\pgfqpoint{5.694602in}{0.614280in}}%
\pgfpathlineto{\pgfqpoint{5.698233in}{0.647171in}}%
\pgfpathlineto{\pgfqpoint{5.701864in}{0.647941in}}%
\pgfpathlineto{\pgfqpoint{5.705495in}{0.568424in}}%
\pgfpathlineto{\pgfqpoint{5.709127in}{0.638091in}}%
\pgfpathlineto{\pgfqpoint{5.712758in}{0.677255in}}%
\pgfpathlineto{\pgfqpoint{5.716389in}{0.615933in}}%
\pgfpathlineto{\pgfqpoint{5.720020in}{0.602088in}}%
\pgfpathlineto{\pgfqpoint{5.723651in}{0.437935in}}%
\pgfpathlineto{\pgfqpoint{5.727282in}{0.633123in}}%
\pgfpathlineto{\pgfqpoint{5.730914in}{0.664103in}}%
\pgfpathlineto{\pgfqpoint{5.734545in}{0.667236in}}%
\pgfpathlineto{\pgfqpoint{5.741807in}{0.595870in}}%
\pgfpathlineto{\pgfqpoint{5.745438in}{0.562725in}}%
\pgfpathlineto{\pgfqpoint{5.749069in}{0.638889in}}%
\pgfpathlineto{\pgfqpoint{5.752701in}{0.659989in}}%
\pgfpathlineto{\pgfqpoint{5.756332in}{0.655086in}}%
\pgfpathlineto{\pgfqpoint{5.759963in}{0.640999in}}%
\pgfpathlineto{\pgfqpoint{5.763594in}{0.593085in}}%
\pgfpathlineto{\pgfqpoint{5.767225in}{0.589386in}}%
\pgfpathlineto{\pgfqpoint{5.770856in}{0.644145in}}%
\pgfpathlineto{\pgfqpoint{5.774488in}{0.656162in}}%
\pgfpathlineto{\pgfqpoint{5.778119in}{0.642798in}}%
\pgfpathlineto{\pgfqpoint{5.781750in}{0.648787in}}%
\pgfpathlineto{\pgfqpoint{5.785381in}{0.583173in}}%
\pgfpathlineto{\pgfqpoint{5.789012in}{0.605109in}}%
\pgfpathlineto{\pgfqpoint{5.792644in}{0.647109in}}%
\pgfpathlineto{\pgfqpoint{5.796275in}{0.649751in}}%
\pgfpathlineto{\pgfqpoint{5.799906in}{0.624308in}}%
\pgfpathlineto{\pgfqpoint{5.803537in}{0.653976in}}%
\pgfpathlineto{\pgfqpoint{5.807168in}{0.562874in}}%
\pgfpathlineto{\pgfqpoint{5.810799in}{0.615764in}}%
\pgfpathlineto{\pgfqpoint{5.814431in}{0.648209in}}%
\pgfpathlineto{\pgfqpoint{5.818062in}{0.640464in}}%
\pgfpathlineto{\pgfqpoint{5.821693in}{0.593930in}}%
\pgfpathlineto{\pgfqpoint{5.825324in}{0.656653in}}%
\pgfpathlineto{\pgfqpoint{5.828955in}{0.492169in}}%
\pgfpathlineto{\pgfqpoint{5.832586in}{0.623130in}}%
\pgfpathlineto{\pgfqpoint{5.836218in}{0.646714in}}%
\pgfpathlineto{\pgfqpoint{5.839849in}{0.626016in}}%
\pgfpathlineto{\pgfqpoint{5.843480in}{0.459862in}}%
\pgfpathlineto{\pgfqpoint{5.847111in}{0.656487in}}%
\pgfpathlineto{\pgfqpoint{5.850742in}{0.553037in}}%
\pgfpathlineto{\pgfqpoint{5.854373in}{0.628077in}}%
\pgfpathlineto{\pgfqpoint{5.858005in}{0.641806in}}%
\pgfpathlineto{\pgfqpoint{5.861636in}{0.600032in}}%
\pgfpathlineto{\pgfqpoint{5.868898in}{0.682432in}}%
\pgfpathlineto{\pgfqpoint{5.872529in}{0.591521in}}%
\pgfpathlineto{\pgfqpoint{5.876160in}{0.577950in}}%
\pgfpathlineto{\pgfqpoint{5.879792in}{0.625178in}}%
\pgfpathlineto{\pgfqpoint{5.883423in}{0.606824in}}%
\pgfpathlineto{\pgfqpoint{5.887054in}{0.646927in}}%
\pgfpathlineto{\pgfqpoint{5.890685in}{0.652561in}}%
\pgfpathlineto{\pgfqpoint{5.894316in}{0.599226in}}%
\pgfpathlineto{\pgfqpoint{5.897948in}{0.631389in}}%
\pgfpathlineto{\pgfqpoint{5.901579in}{0.678415in}}%
\pgfpathlineto{\pgfqpoint{5.905210in}{0.606004in}}%
\pgfpathlineto{\pgfqpoint{5.908841in}{0.603478in}}%
\pgfpathlineto{\pgfqpoint{5.912472in}{0.542113in}}%
\pgfpathlineto{\pgfqpoint{5.916103in}{0.628341in}}%
\pgfpathlineto{\pgfqpoint{5.919735in}{0.663852in}}%
\pgfpathlineto{\pgfqpoint{5.923366in}{0.669855in}}%
\pgfpathlineto{\pgfqpoint{5.926997in}{0.622724in}}%
\pgfpathlineto{\pgfqpoint{5.930628in}{0.598863in}}%
\pgfpathlineto{\pgfqpoint{5.934259in}{0.525970in}}%
\pgfpathlineto{\pgfqpoint{5.937890in}{0.636024in}}%
\pgfpathlineto{\pgfqpoint{5.941522in}{0.660732in}}%
\pgfpathlineto{\pgfqpoint{5.945153in}{0.658879in}}%
\pgfpathlineto{\pgfqpoint{5.948784in}{0.636234in}}%
\pgfpathlineto{\pgfqpoint{5.952415in}{0.598044in}}%
\pgfpathlineto{\pgfqpoint{5.956046in}{0.575040in}}%
\pgfpathlineto{\pgfqpoint{5.959677in}{0.642277in}}%
\pgfpathlineto{\pgfqpoint{5.963309in}{0.657674in}}%
\pgfpathlineto{\pgfqpoint{5.966940in}{0.647943in}}%
\pgfpathlineto{\pgfqpoint{5.970571in}{0.645553in}}%
\pgfpathlineto{\pgfqpoint{5.974202in}{0.591201in}}%
\pgfpathlineto{\pgfqpoint{5.977833in}{0.597042in}}%
\pgfpathlineto{\pgfqpoint{5.981465in}{0.646301in}}%
\pgfpathlineto{\pgfqpoint{5.985096in}{0.652318in}}%
\pgfpathlineto{\pgfqpoint{5.988727in}{0.631870in}}%
\pgfpathlineto{\pgfqpoint{5.992358in}{0.652019in}}%
\pgfpathlineto{\pgfqpoint{5.995989in}{0.577502in}}%
\pgfpathlineto{\pgfqpoint{6.003252in}{0.648501in}}%
\pgfpathlineto{\pgfqpoint{6.006883in}{0.644514in}}%
\pgfpathlineto{\pgfqpoint{6.010514in}{0.607382in}}%
\pgfpathlineto{\pgfqpoint{6.014145in}{0.655911in}}%
\pgfpathlineto{\pgfqpoint{6.017776in}{0.545580in}}%
\pgfpathlineto{\pgfqpoint{6.021407in}{0.619906in}}%
\pgfpathlineto{\pgfqpoint{6.025039in}{0.648360in}}%
\pgfpathlineto{\pgfqpoint{6.028670in}{0.632625in}}%
\pgfpathlineto{\pgfqpoint{6.032301in}{0.606400in}}%
\pgfpathlineto{\pgfqpoint{6.035932in}{0.676852in}}%
\pgfpathlineto{\pgfqpoint{6.039563in}{0.622917in}}%
\pgfpathlineto{\pgfqpoint{6.043194in}{0.488779in}}%
\pgfpathlineto{\pgfqpoint{6.046826in}{0.637179in}}%
\pgfpathlineto{\pgfqpoint{6.050457in}{0.571517in}}%
\pgfpathlineto{\pgfqpoint{6.054088in}{0.644261in}}%
\pgfpathlineto{\pgfqpoint{6.057719in}{0.660693in}}%
\pgfpathlineto{\pgfqpoint{6.061350in}{0.640487in}}%
\pgfpathlineto{\pgfqpoint{6.064981in}{0.579930in}}%
\pgfpathlineto{\pgfqpoint{6.068613in}{0.679454in}}%
\pgfpathlineto{\pgfqpoint{6.072244in}{0.544731in}}%
\pgfpathlineto{\pgfqpoint{6.075875in}{0.598133in}}%
\pgfpathlineto{\pgfqpoint{6.079506in}{0.597698in}}%
\pgfpathlineto{\pgfqpoint{6.083137in}{0.607213in}}%
\pgfpathlineto{\pgfqpoint{6.086769in}{0.662694in}}%
\pgfpathlineto{\pgfqpoint{6.090400in}{0.677127in}}%
\pgfpathlineto{\pgfqpoint{6.094031in}{0.585763in}}%
\pgfpathlineto{\pgfqpoint{6.097662in}{0.597522in}}%
\pgfpathlineto{\pgfqpoint{6.101293in}{0.570733in}}%
\pgfpathlineto{\pgfqpoint{6.104924in}{0.622439in}}%
\pgfpathlineto{\pgfqpoint{6.108556in}{0.662295in}}%
\pgfpathlineto{\pgfqpoint{6.112187in}{0.670196in}}%
\pgfpathlineto{\pgfqpoint{6.115818in}{0.613405in}}%
\pgfpathlineto{\pgfqpoint{6.119449in}{0.602462in}}%
\pgfpathlineto{\pgfqpoint{6.123080in}{0.513732in}}%
\pgfpathlineto{\pgfqpoint{6.126711in}{0.632074in}}%
\pgfpathlineto{\pgfqpoint{6.130343in}{0.661320in}}%
\pgfpathlineto{\pgfqpoint{6.133974in}{0.663020in}}%
\pgfpathlineto{\pgfqpoint{6.137605in}{0.629709in}}%
\pgfpathlineto{\pgfqpoint{6.141236in}{0.602031in}}%
\pgfpathlineto{\pgfqpoint{6.144867in}{0.546221in}}%
\pgfpathlineto{\pgfqpoint{6.148498in}{0.639420in}}%
\pgfpathlineto{\pgfqpoint{6.152130in}{0.658723in}}%
\pgfpathlineto{\pgfqpoint{6.155761in}{0.652724in}}%
\pgfpathlineto{\pgfqpoint{6.159392in}{0.641149in}}%
\pgfpathlineto{\pgfqpoint{6.163023in}{0.598085in}}%
\pgfpathlineto{\pgfqpoint{6.166654in}{0.584182in}}%
\pgfpathlineto{\pgfqpoint{6.170285in}{0.644784in}}%
\pgfpathlineto{\pgfqpoint{6.173917in}{0.654529in}}%
\pgfpathlineto{\pgfqpoint{6.177548in}{0.638916in}}%
\pgfpathlineto{\pgfqpoint{6.181179in}{0.649179in}}%
\pgfpathlineto{\pgfqpoint{6.184810in}{0.588817in}}%
\pgfpathlineto{\pgfqpoint{6.188441in}{0.603179in}}%
\pgfpathlineto{\pgfqpoint{6.192073in}{0.648156in}}%
\pgfpathlineto{\pgfqpoint{6.195704in}{0.648085in}}%
\pgfpathlineto{\pgfqpoint{6.199335in}{0.618677in}}%
\pgfpathlineto{\pgfqpoint{6.202966in}{0.654441in}}%
\pgfpathlineto{\pgfqpoint{6.206597in}{0.569752in}}%
\pgfpathlineto{\pgfqpoint{6.210228in}{0.615203in}}%
\pgfpathlineto{\pgfqpoint{6.213860in}{0.649381in}}%
\pgfpathlineto{\pgfqpoint{6.217491in}{0.638377in}}%
\pgfpathlineto{\pgfqpoint{6.221122in}{0.580236in}}%
\pgfpathlineto{\pgfqpoint{6.224753in}{0.673450in}}%
\pgfpathlineto{\pgfqpoint{6.228384in}{0.627872in}}%
\pgfpathlineto{\pgfqpoint{6.232015in}{0.548763in}}%
\pgfpathlineto{\pgfqpoint{6.235647in}{0.640816in}}%
\pgfpathlineto{\pgfqpoint{6.239278in}{0.543779in}}%
\pgfpathlineto{\pgfqpoint{6.242909in}{0.642138in}}%
\pgfpathlineto{\pgfqpoint{6.246540in}{0.661722in}}%
\pgfpathlineto{\pgfqpoint{6.250171in}{0.646253in}}%
\pgfpathlineto{\pgfqpoint{6.253802in}{0.517753in}}%
\pgfpathlineto{\pgfqpoint{6.257434in}{0.677959in}}%
\pgfpathlineto{\pgfqpoint{6.261065in}{0.401488in}}%
\pgfpathlineto{\pgfqpoint{6.264696in}{0.595524in}}%
\pgfpathlineto{\pgfqpoint{6.268327in}{0.608560in}}%
\pgfpathlineto{\pgfqpoint{6.271958in}{0.595748in}}%
\pgfpathlineto{\pgfqpoint{6.275590in}{0.660899in}}%
\pgfpathlineto{\pgfqpoint{6.279221in}{0.677489in}}%
\pgfpathlineto{\pgfqpoint{6.282852in}{0.561109in}}%
\pgfpathlineto{\pgfqpoint{6.286483in}{0.596482in}}%
\pgfpathlineto{\pgfqpoint{6.290114in}{0.588814in}}%
\pgfpathlineto{\pgfqpoint{6.293745in}{0.616009in}}%
\pgfpathlineto{\pgfqpoint{6.297377in}{0.661555in}}%
\pgfpathlineto{\pgfqpoint{6.301008in}{0.672230in}}%
\pgfpathlineto{\pgfqpoint{6.304639in}{0.602191in}}%
\pgfpathlineto{\pgfqpoint{6.308270in}{0.603283in}}%
\pgfpathlineto{\pgfqpoint{6.311901in}{0.560920in}}%
\pgfpathlineto{\pgfqpoint{6.315532in}{0.627444in}}%
\pgfpathlineto{\pgfqpoint{6.319164in}{0.661234in}}%
\pgfpathlineto{\pgfqpoint{6.322795in}{0.665910in}}%
\pgfpathlineto{\pgfqpoint{6.326426in}{0.622646in}}%
\pgfpathlineto{\pgfqpoint{6.330057in}{0.604480in}}%
\pgfpathlineto{\pgfqpoint{6.333688in}{0.446159in}}%
\pgfpathlineto{\pgfqpoint{6.337319in}{0.636246in}}%
\pgfpathlineto{\pgfqpoint{6.340951in}{0.659403in}}%
\pgfpathlineto{\pgfqpoint{6.344582in}{0.656711in}}%
\pgfpathlineto{\pgfqpoint{6.348213in}{0.636386in}}%
\pgfpathlineto{\pgfqpoint{6.355475in}{0.566691in}}%
\pgfpathlineto{\pgfqpoint{6.359106in}{0.642768in}}%
\pgfpathlineto{\pgfqpoint{6.362738in}{0.656066in}}%
\pgfpathlineto{\pgfqpoint{6.366369in}{0.650349in}}%
\pgfpathlineto{\pgfqpoint{6.370000in}{0.615933in}}%
\pgfpathlineto{\pgfqpoint{6.370000in}{0.615933in}}%
\pgfusepath{stroke}%
\end{pgfscope}%
\begin{pgfscope}%
\pgfpathrectangle{\pgfqpoint{4.554414in}{0.320625in}}{\pgfqpoint{1.815586in}{1.164375in}}%
\pgfusepath{clip}%
\pgfsetrectcap%
\pgfsetroundjoin%
\pgfsetlinewidth{0.602250pt}%
\definecolor{currentstroke}{rgb}{0.819608,0.286275,0.356863}%
\pgfsetstrokecolor{currentstroke}%
\pgfsetstrokeopacity{0.120000}%
\pgfsetdash{}{0pt}%
\pgfpathmoveto{\pgfqpoint{4.554414in}{1.460401in}}%
\pgfpathlineto{\pgfqpoint{4.558045in}{1.449047in}}%
\pgfpathlineto{\pgfqpoint{4.565307in}{1.431795in}}%
\pgfpathlineto{\pgfqpoint{4.572569in}{1.417378in}}%
\pgfpathlineto{\pgfqpoint{4.597988in}{1.369712in}}%
\pgfpathlineto{\pgfqpoint{4.634299in}{1.300919in}}%
\pgfpathlineto{\pgfqpoint{4.659718in}{1.255526in}}%
\pgfpathlineto{\pgfqpoint{4.685136in}{1.212622in}}%
\pgfpathlineto{\pgfqpoint{4.714185in}{1.166000in}}%
\pgfpathlineto{\pgfqpoint{4.746866in}{1.115754in}}%
\pgfpathlineto{\pgfqpoint{4.786809in}{1.056487in}}%
\pgfpathlineto{\pgfqpoint{4.837645in}{0.983266in}}%
\pgfpathlineto{\pgfqpoint{4.899375in}{0.896391in}}%
\pgfpathlineto{\pgfqpoint{4.979261in}{0.785887in}}%
\pgfpathlineto{\pgfqpoint{5.011941in}{0.742551in}}%
\pgfpathlineto{\pgfqpoint{5.015573in}{0.735997in}}%
\pgfpathlineto{\pgfqpoint{5.019204in}{0.731182in}}%
\pgfpathlineto{\pgfqpoint{5.022835in}{0.728250in}}%
\pgfpathlineto{\pgfqpoint{5.030097in}{0.716234in}}%
\pgfpathlineto{\pgfqpoint{5.037360in}{0.709791in}}%
\pgfpathlineto{\pgfqpoint{5.040991in}{0.711509in}}%
\pgfpathlineto{\pgfqpoint{5.044622in}{0.696450in}}%
\pgfpathlineto{\pgfqpoint{5.048253in}{0.691906in}}%
\pgfpathlineto{\pgfqpoint{5.051884in}{0.697906in}}%
\pgfpathlineto{\pgfqpoint{5.055515in}{0.685232in}}%
\pgfpathlineto{\pgfqpoint{5.059147in}{0.676554in}}%
\pgfpathlineto{\pgfqpoint{5.066409in}{0.685855in}}%
\pgfpathlineto{\pgfqpoint{5.070040in}{0.697481in}}%
\pgfpathlineto{\pgfqpoint{5.073671in}{0.659912in}}%
\pgfpathlineto{\pgfqpoint{5.077302in}{0.653545in}}%
\pgfpathlineto{\pgfqpoint{5.080934in}{0.685363in}}%
\pgfpathlineto{\pgfqpoint{5.084565in}{0.656590in}}%
\pgfpathlineto{\pgfqpoint{5.088196in}{0.636925in}}%
\pgfpathlineto{\pgfqpoint{5.091827in}{0.666193in}}%
\pgfpathlineto{\pgfqpoint{5.099090in}{0.692330in}}%
\pgfpathlineto{\pgfqpoint{5.102721in}{0.640190in}}%
\pgfpathlineto{\pgfqpoint{5.106352in}{0.617922in}}%
\pgfpathlineto{\pgfqpoint{5.109983in}{0.685482in}}%
\pgfpathlineto{\pgfqpoint{5.117245in}{0.599809in}}%
\pgfpathlineto{\pgfqpoint{5.120877in}{0.662903in}}%
\pgfpathlineto{\pgfqpoint{5.128139in}{0.685789in}}%
\pgfpathlineto{\pgfqpoint{5.131770in}{0.636266in}}%
\pgfpathlineto{\pgfqpoint{5.135401in}{0.601471in}}%
\pgfpathlineto{\pgfqpoint{5.139032in}{0.688719in}}%
\pgfpathlineto{\pgfqpoint{5.146295in}{0.598934in}}%
\pgfpathlineto{\pgfqpoint{5.149926in}{0.661859in}}%
\pgfpathlineto{\pgfqpoint{5.157188in}{0.675288in}}%
\pgfpathlineto{\pgfqpoint{5.160819in}{0.626417in}}%
\pgfpathlineto{\pgfqpoint{5.164451in}{0.623106in}}%
\pgfpathlineto{\pgfqpoint{5.168082in}{0.691945in}}%
\pgfpathlineto{\pgfqpoint{5.171713in}{0.652290in}}%
\pgfpathlineto{\pgfqpoint{5.175344in}{0.627332in}}%
\pgfpathlineto{\pgfqpoint{5.178975in}{0.661302in}}%
\pgfpathlineto{\pgfqpoint{5.182606in}{0.659944in}}%
\pgfpathlineto{\pgfqpoint{5.186238in}{0.661455in}}%
\pgfpathlineto{\pgfqpoint{5.189869in}{0.595823in}}%
\pgfpathlineto{\pgfqpoint{5.197131in}{0.694344in}}%
\pgfpathlineto{\pgfqpoint{5.200762in}{0.663796in}}%
\pgfpathlineto{\pgfqpoint{5.204394in}{0.651163in}}%
\pgfpathlineto{\pgfqpoint{5.208025in}{0.661177in}}%
\pgfpathlineto{\pgfqpoint{5.211656in}{0.651306in}}%
\pgfpathlineto{\pgfqpoint{5.215287in}{0.648312in}}%
\pgfpathlineto{\pgfqpoint{5.218918in}{0.527870in}}%
\pgfpathlineto{\pgfqpoint{5.222549in}{0.663760in}}%
\pgfpathlineto{\pgfqpoint{5.226181in}{0.695307in}}%
\pgfpathlineto{\pgfqpoint{5.229812in}{0.675501in}}%
\pgfpathlineto{\pgfqpoint{5.233443in}{0.668876in}}%
\pgfpathlineto{\pgfqpoint{5.237074in}{0.660542in}}%
\pgfpathlineto{\pgfqpoint{5.240705in}{0.644073in}}%
\pgfpathlineto{\pgfqpoint{5.244336in}{0.643123in}}%
\pgfpathlineto{\pgfqpoint{5.247968in}{0.599213in}}%
\pgfpathlineto{\pgfqpoint{5.251599in}{0.669298in}}%
\pgfpathlineto{\pgfqpoint{5.255230in}{0.693686in}}%
\pgfpathlineto{\pgfqpoint{5.258861in}{0.683715in}}%
\pgfpathlineto{\pgfqpoint{5.262492in}{0.679321in}}%
\pgfpathlineto{\pgfqpoint{5.266123in}{0.657316in}}%
\pgfpathlineto{\pgfqpoint{5.269755in}{0.639308in}}%
\pgfpathlineto{\pgfqpoint{5.273386in}{0.649861in}}%
\pgfpathlineto{\pgfqpoint{5.277017in}{0.613664in}}%
\pgfpathlineto{\pgfqpoint{5.280648in}{0.666155in}}%
\pgfpathlineto{\pgfqpoint{5.284279in}{0.688241in}}%
\pgfpathlineto{\pgfqpoint{5.287910in}{0.685106in}}%
\pgfpathlineto{\pgfqpoint{5.291542in}{0.679750in}}%
\pgfpathlineto{\pgfqpoint{5.295173in}{0.650028in}}%
\pgfpathlineto{\pgfqpoint{5.298804in}{0.638495in}}%
\pgfpathlineto{\pgfqpoint{5.302435in}{0.665147in}}%
\pgfpathlineto{\pgfqpoint{5.306066in}{0.615846in}}%
\pgfpathlineto{\pgfqpoint{5.309698in}{0.656482in}}%
\pgfpathlineto{\pgfqpoint{5.313329in}{0.678792in}}%
\pgfpathlineto{\pgfqpoint{5.316960in}{0.677398in}}%
\pgfpathlineto{\pgfqpoint{5.320591in}{0.664714in}}%
\pgfpathlineto{\pgfqpoint{5.324222in}{0.642980in}}%
\pgfpathlineto{\pgfqpoint{5.327853in}{0.645999in}}%
\pgfpathlineto{\pgfqpoint{5.331485in}{0.683870in}}%
\pgfpathlineto{\pgfqpoint{5.335116in}{0.609574in}}%
\pgfpathlineto{\pgfqpoint{5.338747in}{0.642121in}}%
\pgfpathlineto{\pgfqpoint{5.342378in}{0.667374in}}%
\pgfpathlineto{\pgfqpoint{5.346009in}{0.659796in}}%
\pgfpathlineto{\pgfqpoint{5.349640in}{0.609277in}}%
\pgfpathlineto{\pgfqpoint{5.353272in}{0.647622in}}%
\pgfpathlineto{\pgfqpoint{5.356903in}{0.664235in}}%
\pgfpathlineto{\pgfqpoint{5.360534in}{0.701395in}}%
\pgfpathlineto{\pgfqpoint{5.364165in}{0.580031in}}%
\pgfpathlineto{\pgfqpoint{5.371427in}{0.659425in}}%
\pgfpathlineto{\pgfqpoint{5.375059in}{0.633322in}}%
\pgfpathlineto{\pgfqpoint{5.378690in}{0.623442in}}%
\pgfpathlineto{\pgfqpoint{5.382321in}{0.660602in}}%
\pgfpathlineto{\pgfqpoint{5.385952in}{0.632010in}}%
\pgfpathlineto{\pgfqpoint{5.389583in}{0.670932in}}%
\pgfpathlineto{\pgfqpoint{5.393215in}{0.620850in}}%
\pgfpathlineto{\pgfqpoint{5.396846in}{0.650222in}}%
\pgfpathlineto{\pgfqpoint{5.400477in}{0.669059in}}%
\pgfpathlineto{\pgfqpoint{5.404108in}{0.677556in}}%
\pgfpathlineto{\pgfqpoint{5.407739in}{0.675703in}}%
\pgfpathlineto{\pgfqpoint{5.411370in}{0.658963in}}%
\pgfpathlineto{\pgfqpoint{5.415002in}{0.628144in}}%
\pgfpathlineto{\pgfqpoint{5.418633in}{0.616865in}}%
\pgfpathlineto{\pgfqpoint{5.422264in}{0.614790in}}%
\pgfpathlineto{\pgfqpoint{5.425895in}{0.678937in}}%
\pgfpathlineto{\pgfqpoint{5.429526in}{0.696297in}}%
\pgfpathlineto{\pgfqpoint{5.433157in}{0.644453in}}%
\pgfpathlineto{\pgfqpoint{5.436789in}{0.626149in}}%
\pgfpathlineto{\pgfqpoint{5.440420in}{0.657482in}}%
\pgfpathlineto{\pgfqpoint{5.444051in}{0.653694in}}%
\pgfpathlineto{\pgfqpoint{5.447682in}{0.644185in}}%
\pgfpathlineto{\pgfqpoint{5.451313in}{0.624918in}}%
\pgfpathlineto{\pgfqpoint{5.454944in}{0.581268in}}%
\pgfpathlineto{\pgfqpoint{5.458576in}{0.681905in}}%
\pgfpathlineto{\pgfqpoint{5.462207in}{0.606900in}}%
\pgfpathlineto{\pgfqpoint{5.465838in}{0.609152in}}%
\pgfpathlineto{\pgfqpoint{5.469469in}{0.662986in}}%
\pgfpathlineto{\pgfqpoint{5.473100in}{0.686520in}}%
\pgfpathlineto{\pgfqpoint{5.476731in}{0.699131in}}%
\pgfpathlineto{\pgfqpoint{5.480363in}{0.556197in}}%
\pgfpathlineto{\pgfqpoint{5.483994in}{0.601768in}}%
\pgfpathlineto{\pgfqpoint{5.487625in}{0.664819in}}%
\pgfpathlineto{\pgfqpoint{5.491256in}{0.630377in}}%
\pgfpathlineto{\pgfqpoint{5.494887in}{0.619674in}}%
\pgfpathlineto{\pgfqpoint{5.498519in}{0.660285in}}%
\pgfpathlineto{\pgfqpoint{5.502150in}{0.644424in}}%
\pgfpathlineto{\pgfqpoint{5.505781in}{0.672881in}}%
\pgfpathlineto{\pgfqpoint{5.509412in}{0.605727in}}%
\pgfpathlineto{\pgfqpoint{5.513043in}{0.651051in}}%
\pgfpathlineto{\pgfqpoint{5.516674in}{0.675451in}}%
\pgfpathlineto{\pgfqpoint{5.520306in}{0.676328in}}%
\pgfpathlineto{\pgfqpoint{5.523937in}{0.664758in}}%
\pgfpathlineto{\pgfqpoint{5.527568in}{0.623936in}}%
\pgfpathlineto{\pgfqpoint{5.531199in}{0.643507in}}%
\pgfpathlineto{\pgfqpoint{5.534830in}{0.691550in}}%
\pgfpathlineto{\pgfqpoint{5.538461in}{0.592524in}}%
\pgfpathlineto{\pgfqpoint{5.542093in}{0.634098in}}%
\pgfpathlineto{\pgfqpoint{5.545724in}{0.664847in}}%
\pgfpathlineto{\pgfqpoint{5.549355in}{0.656676in}}%
\pgfpathlineto{\pgfqpoint{5.552986in}{0.597330in}}%
\pgfpathlineto{\pgfqpoint{5.556617in}{0.637182in}}%
\pgfpathlineto{\pgfqpoint{5.560248in}{0.666000in}}%
\pgfpathlineto{\pgfqpoint{5.563880in}{0.706240in}}%
\pgfpathlineto{\pgfqpoint{5.567511in}{0.460612in}}%
\pgfpathlineto{\pgfqpoint{5.571142in}{0.601767in}}%
\pgfpathlineto{\pgfqpoint{5.574773in}{0.659557in}}%
\pgfpathlineto{\pgfqpoint{5.578404in}{0.629691in}}%
\pgfpathlineto{\pgfqpoint{5.582035in}{0.625713in}}%
\pgfpathlineto{\pgfqpoint{5.585667in}{0.656202in}}%
\pgfpathlineto{\pgfqpoint{5.589298in}{0.631966in}}%
\pgfpathlineto{\pgfqpoint{5.592929in}{0.672903in}}%
\pgfpathlineto{\pgfqpoint{5.596560in}{0.611190in}}%
\pgfpathlineto{\pgfqpoint{5.600191in}{0.649825in}}%
\pgfpathlineto{\pgfqpoint{5.603823in}{0.671546in}}%
\pgfpathlineto{\pgfqpoint{5.607454in}{0.680158in}}%
\pgfpathlineto{\pgfqpoint{5.611085in}{0.679621in}}%
\pgfpathlineto{\pgfqpoint{5.618347in}{0.628836in}}%
\pgfpathlineto{\pgfqpoint{5.621978in}{0.627198in}}%
\pgfpathlineto{\pgfqpoint{5.625610in}{0.601270in}}%
\pgfpathlineto{\pgfqpoint{5.629241in}{0.675943in}}%
\pgfpathlineto{\pgfqpoint{5.632872in}{0.694095in}}%
\pgfpathlineto{\pgfqpoint{5.636503in}{0.646710in}}%
\pgfpathlineto{\pgfqpoint{5.640134in}{0.633653in}}%
\pgfpathlineto{\pgfqpoint{5.643765in}{0.655328in}}%
\pgfpathlineto{\pgfqpoint{5.647397in}{0.655025in}}%
\pgfpathlineto{\pgfqpoint{5.651028in}{0.651589in}}%
\pgfpathlineto{\pgfqpoint{5.654659in}{0.636065in}}%
\pgfpathlineto{\pgfqpoint{5.658290in}{0.601551in}}%
\pgfpathlineto{\pgfqpoint{5.661921in}{0.680991in}}%
\pgfpathlineto{\pgfqpoint{5.665552in}{0.614546in}}%
\pgfpathlineto{\pgfqpoint{5.669184in}{0.595313in}}%
\pgfpathlineto{\pgfqpoint{5.672815in}{0.660162in}}%
\pgfpathlineto{\pgfqpoint{5.676446in}{0.685883in}}%
\pgfpathlineto{\pgfqpoint{5.680077in}{0.699572in}}%
\pgfpathlineto{\pgfqpoint{5.683708in}{0.587129in}}%
\pgfpathlineto{\pgfqpoint{5.687340in}{0.592518in}}%
\pgfpathlineto{\pgfqpoint{5.690971in}{0.666278in}}%
\pgfpathlineto{\pgfqpoint{5.694602in}{0.635968in}}%
\pgfpathlineto{\pgfqpoint{5.698233in}{0.600491in}}%
\pgfpathlineto{\pgfqpoint{5.701864in}{0.653664in}}%
\pgfpathlineto{\pgfqpoint{5.705495in}{0.642236in}}%
\pgfpathlineto{\pgfqpoint{5.709127in}{0.675368in}}%
\pgfpathlineto{\pgfqpoint{5.712758in}{0.591920in}}%
\pgfpathlineto{\pgfqpoint{5.716389in}{0.648944in}}%
\pgfpathlineto{\pgfqpoint{5.720020in}{0.675693in}}%
\pgfpathlineto{\pgfqpoint{5.723651in}{0.677825in}}%
\pgfpathlineto{\pgfqpoint{5.727282in}{0.669864in}}%
\pgfpathlineto{\pgfqpoint{5.730914in}{0.603691in}}%
\pgfpathlineto{\pgfqpoint{5.734545in}{0.639006in}}%
\pgfpathlineto{\pgfqpoint{5.738176in}{0.692659in}}%
\pgfpathlineto{\pgfqpoint{5.741807in}{0.572276in}}%
\pgfpathlineto{\pgfqpoint{5.745438in}{0.631496in}}%
\pgfpathlineto{\pgfqpoint{5.749069in}{0.666281in}}%
\pgfpathlineto{\pgfqpoint{5.752701in}{0.659704in}}%
\pgfpathlineto{\pgfqpoint{5.756332in}{0.619743in}}%
\pgfpathlineto{\pgfqpoint{5.759963in}{0.623161in}}%
\pgfpathlineto{\pgfqpoint{5.767225in}{0.705781in}}%
\pgfpathlineto{\pgfqpoint{5.770856in}{0.557422in}}%
\pgfpathlineto{\pgfqpoint{5.774488in}{0.594662in}}%
\pgfpathlineto{\pgfqpoint{5.778119in}{0.661535in}}%
\pgfpathlineto{\pgfqpoint{5.785381in}{0.609615in}}%
\pgfpathlineto{\pgfqpoint{5.789012in}{0.648775in}}%
\pgfpathlineto{\pgfqpoint{5.792644in}{0.628629in}}%
\pgfpathlineto{\pgfqpoint{5.796275in}{0.675133in}}%
\pgfpathlineto{\pgfqpoint{5.799906in}{0.601643in}}%
\pgfpathlineto{\pgfqpoint{5.803537in}{0.648337in}}%
\pgfpathlineto{\pgfqpoint{5.807168in}{0.672090in}}%
\pgfpathlineto{\pgfqpoint{5.810799in}{0.681441in}}%
\pgfpathlineto{\pgfqpoint{5.814431in}{0.682037in}}%
\pgfpathlineto{\pgfqpoint{5.818062in}{0.649245in}}%
\pgfpathlineto{\pgfqpoint{5.821693in}{0.629134in}}%
\pgfpathlineto{\pgfqpoint{5.825324in}{0.634096in}}%
\pgfpathlineto{\pgfqpoint{5.828955in}{0.585067in}}%
\pgfpathlineto{\pgfqpoint{5.832586in}{0.673153in}}%
\pgfpathlineto{\pgfqpoint{5.836218in}{0.692124in}}%
\pgfpathlineto{\pgfqpoint{5.839849in}{0.648139in}}%
\pgfpathlineto{\pgfqpoint{5.843480in}{0.639201in}}%
\pgfpathlineto{\pgfqpoint{5.847111in}{0.653154in}}%
\pgfpathlineto{\pgfqpoint{5.850742in}{0.655768in}}%
\pgfpathlineto{\pgfqpoint{5.854373in}{0.656419in}}%
\pgfpathlineto{\pgfqpoint{5.858005in}{0.643165in}}%
\pgfpathlineto{\pgfqpoint{5.861636in}{0.613690in}}%
\pgfpathlineto{\pgfqpoint{5.865267in}{0.679967in}}%
\pgfpathlineto{\pgfqpoint{5.868898in}{0.619544in}}%
\pgfpathlineto{\pgfqpoint{5.872529in}{0.578203in}}%
\pgfpathlineto{\pgfqpoint{5.876160in}{0.657459in}}%
\pgfpathlineto{\pgfqpoint{5.879792in}{0.685237in}}%
\pgfpathlineto{\pgfqpoint{5.883423in}{0.699763in}}%
\pgfpathlineto{\pgfqpoint{5.887054in}{0.601803in}}%
\pgfpathlineto{\pgfqpoint{5.890685in}{0.581529in}}%
\pgfpathlineto{\pgfqpoint{5.894316in}{0.667017in}}%
\pgfpathlineto{\pgfqpoint{5.897948in}{0.639850in}}%
\pgfpathlineto{\pgfqpoint{5.901579in}{0.573220in}}%
\pgfpathlineto{\pgfqpoint{5.905210in}{0.647205in}}%
\pgfpathlineto{\pgfqpoint{5.908841in}{0.683189in}}%
\pgfpathlineto{\pgfqpoint{5.912472in}{0.709678in}}%
\pgfpathlineto{\pgfqpoint{5.916103in}{0.643458in}}%
\pgfpathlineto{\pgfqpoint{5.919735in}{0.601208in}}%
\pgfpathlineto{\pgfqpoint{5.923366in}{0.670464in}}%
\pgfpathlineto{\pgfqpoint{5.926997in}{0.621429in}}%
\pgfpathlineto{\pgfqpoint{5.930628in}{0.613580in}}%
\pgfpathlineto{\pgfqpoint{5.934259in}{0.657118in}}%
\pgfpathlineto{\pgfqpoint{5.937890in}{0.692117in}}%
\pgfpathlineto{\pgfqpoint{5.941522in}{0.710342in}}%
\pgfpathlineto{\pgfqpoint{5.945153in}{0.671695in}}%
\pgfpathlineto{\pgfqpoint{5.948784in}{0.651839in}}%
\pgfpathlineto{\pgfqpoint{5.952415in}{0.677636in}}%
\pgfpathlineto{\pgfqpoint{5.956046in}{0.610541in}}%
\pgfpathlineto{\pgfqpoint{5.959677in}{0.601394in}}%
\pgfpathlineto{\pgfqpoint{5.963309in}{0.657538in}}%
\pgfpathlineto{\pgfqpoint{5.966940in}{0.689787in}}%
\pgfpathlineto{\pgfqpoint{5.970571in}{0.704279in}}%
\pgfpathlineto{\pgfqpoint{5.974202in}{0.602369in}}%
\pgfpathlineto{\pgfqpoint{5.977833in}{0.545356in}}%
\pgfpathlineto{\pgfqpoint{5.981465in}{0.661280in}}%
\pgfpathlineto{\pgfqpoint{5.985096in}{0.632645in}}%
\pgfpathlineto{\pgfqpoint{5.988727in}{0.610049in}}%
\pgfpathlineto{\pgfqpoint{5.992358in}{0.643397in}}%
\pgfpathlineto{\pgfqpoint{5.995989in}{0.627611in}}%
\pgfpathlineto{\pgfqpoint{5.999620in}{0.676796in}}%
\pgfpathlineto{\pgfqpoint{6.003252in}{0.584905in}}%
\pgfpathlineto{\pgfqpoint{6.006883in}{0.646795in}}%
\pgfpathlineto{\pgfqpoint{6.010514in}{0.673115in}}%
\pgfpathlineto{\pgfqpoint{6.014145in}{0.683121in}}%
\pgfpathlineto{\pgfqpoint{6.017776in}{0.685144in}}%
\pgfpathlineto{\pgfqpoint{6.021407in}{0.642720in}}%
\pgfpathlineto{\pgfqpoint{6.025039in}{0.628808in}}%
\pgfpathlineto{\pgfqpoint{6.028670in}{0.640979in}}%
\pgfpathlineto{\pgfqpoint{6.032301in}{0.552257in}}%
\pgfpathlineto{\pgfqpoint{6.035932in}{0.669448in}}%
\pgfpathlineto{\pgfqpoint{6.039563in}{0.689535in}}%
\pgfpathlineto{\pgfqpoint{6.043194in}{0.649502in}}%
\pgfpathlineto{\pgfqpoint{6.046826in}{0.645021in}}%
\pgfpathlineto{\pgfqpoint{6.050457in}{0.650142in}}%
\pgfpathlineto{\pgfqpoint{6.054088in}{0.656442in}}%
\pgfpathlineto{\pgfqpoint{6.057719in}{0.661380in}}%
\pgfpathlineto{\pgfqpoint{6.061350in}{0.649944in}}%
\pgfpathlineto{\pgfqpoint{6.064981in}{0.624706in}}%
\pgfpathlineto{\pgfqpoint{6.068613in}{0.678488in}}%
\pgfpathlineto{\pgfqpoint{6.072244in}{0.624557in}}%
\pgfpathlineto{\pgfqpoint{6.075875in}{0.536525in}}%
\pgfpathlineto{\pgfqpoint{6.079506in}{0.653796in}}%
\pgfpathlineto{\pgfqpoint{6.083137in}{0.684236in}}%
\pgfpathlineto{\pgfqpoint{6.086769in}{0.699714in}}%
\pgfpathlineto{\pgfqpoint{6.090400in}{0.613935in}}%
\pgfpathlineto{\pgfqpoint{6.094031in}{0.561535in}}%
\pgfpathlineto{\pgfqpoint{6.097662in}{0.667529in}}%
\pgfpathlineto{\pgfqpoint{6.101293in}{0.643856in}}%
\pgfpathlineto{\pgfqpoint{6.104924in}{0.473883in}}%
\pgfpathlineto{\pgfqpoint{6.108556in}{0.638137in}}%
\pgfpathlineto{\pgfqpoint{6.112187in}{0.680244in}}%
\pgfpathlineto{\pgfqpoint{6.115818in}{0.708514in}}%
\pgfpathlineto{\pgfqpoint{6.119449in}{0.647802in}}%
\pgfpathlineto{\pgfqpoint{6.123080in}{0.612030in}}%
\pgfpathlineto{\pgfqpoint{6.126711in}{0.670394in}}%
\pgfpathlineto{\pgfqpoint{6.130343in}{0.626917in}}%
\pgfpathlineto{\pgfqpoint{6.133974in}{0.596849in}}%
\pgfpathlineto{\pgfqpoint{6.137605in}{0.651796in}}%
\pgfpathlineto{\pgfqpoint{6.141236in}{0.690131in}}%
\pgfpathlineto{\pgfqpoint{6.144867in}{0.709254in}}%
\pgfpathlineto{\pgfqpoint{6.148498in}{0.673551in}}%
\pgfpathlineto{\pgfqpoint{6.152130in}{0.656153in}}%
\pgfpathlineto{\pgfqpoint{6.155761in}{0.676409in}}%
\pgfpathlineto{\pgfqpoint{6.159392in}{0.616377in}}%
\pgfpathlineto{\pgfqpoint{6.163023in}{0.584929in}}%
\pgfpathlineto{\pgfqpoint{6.166654in}{0.654156in}}%
\pgfpathlineto{\pgfqpoint{6.170285in}{0.688727in}}%
\pgfpathlineto{\pgfqpoint{6.173917in}{0.703832in}}%
\pgfpathlineto{\pgfqpoint{6.177548in}{0.613076in}}%
\pgfpathlineto{\pgfqpoint{6.181179in}{0.484405in}}%
\pgfpathlineto{\pgfqpoint{6.184810in}{0.662316in}}%
\pgfpathlineto{\pgfqpoint{6.188441in}{0.637098in}}%
\pgfpathlineto{\pgfqpoint{6.192073in}{0.585905in}}%
\pgfpathlineto{\pgfqpoint{6.195704in}{0.633674in}}%
\pgfpathlineto{\pgfqpoint{6.199335in}{0.623062in}}%
\pgfpathlineto{\pgfqpoint{6.202966in}{0.678672in}}%
\pgfpathlineto{\pgfqpoint{6.206597in}{0.564512in}}%
\pgfpathlineto{\pgfqpoint{6.210228in}{0.644435in}}%
\pgfpathlineto{\pgfqpoint{6.213860in}{0.672909in}}%
\pgfpathlineto{\pgfqpoint{6.217491in}{0.683937in}}%
\pgfpathlineto{\pgfqpoint{6.221122in}{0.687177in}}%
\pgfpathlineto{\pgfqpoint{6.224753in}{0.636364in}}%
\pgfpathlineto{\pgfqpoint{6.228384in}{0.628223in}}%
\pgfpathlineto{\pgfqpoint{6.232015in}{0.645918in}}%
\pgfpathlineto{\pgfqpoint{6.235647in}{0.361205in}}%
\pgfpathlineto{\pgfqpoint{6.239278in}{0.665883in}}%
\pgfpathlineto{\pgfqpoint{6.242909in}{0.687154in}}%
\pgfpathlineto{\pgfqpoint{6.246540in}{0.650272in}}%
\pgfpathlineto{\pgfqpoint{6.250171in}{0.649518in}}%
\pgfpathlineto{\pgfqpoint{6.253802in}{0.647146in}}%
\pgfpathlineto{\pgfqpoint{6.261065in}{0.664760in}}%
\pgfpathlineto{\pgfqpoint{6.264696in}{0.654708in}}%
\pgfpathlineto{\pgfqpoint{6.268327in}{0.632564in}}%
\pgfpathlineto{\pgfqpoint{6.271958in}{0.676972in}}%
\pgfpathlineto{\pgfqpoint{6.275590in}{0.627957in}}%
\pgfpathlineto{\pgfqpoint{6.279221in}{0.503442in}}%
\pgfpathlineto{\pgfqpoint{6.282852in}{0.650237in}}%
\pgfpathlineto{\pgfqpoint{6.286483in}{0.683270in}}%
\pgfpathlineto{\pgfqpoint{6.290114in}{0.699525in}}%
\pgfpathlineto{\pgfqpoint{6.297377in}{0.667981in}}%
\pgfpathlineto{\pgfqpoint{6.301008in}{0.679328in}}%
\pgfpathlineto{\pgfqpoint{6.304639in}{0.627414in}}%
\pgfpathlineto{\pgfqpoint{6.308270in}{0.600052in}}%
\pgfpathlineto{\pgfqpoint{6.311901in}{0.648987in}}%
\pgfpathlineto{\pgfqpoint{6.315532in}{0.673243in}}%
\pgfpathlineto{\pgfqpoint{6.319164in}{0.686691in}}%
\pgfpathlineto{\pgfqpoint{6.322795in}{0.681256in}}%
\pgfpathlineto{\pgfqpoint{6.326426in}{0.664077in}}%
\pgfpathlineto{\pgfqpoint{6.330057in}{0.678446in}}%
\pgfpathlineto{\pgfqpoint{6.333688in}{0.634853in}}%
\pgfpathlineto{\pgfqpoint{6.337319in}{0.643355in}}%
\pgfpathlineto{\pgfqpoint{6.340951in}{0.646634in}}%
\pgfpathlineto{\pgfqpoint{6.344582in}{0.657729in}}%
\pgfpathlineto{\pgfqpoint{6.348213in}{0.665113in}}%
\pgfpathlineto{\pgfqpoint{6.351844in}{0.664388in}}%
\pgfpathlineto{\pgfqpoint{6.355475in}{0.653645in}}%
\pgfpathlineto{\pgfqpoint{6.359106in}{0.675767in}}%
\pgfpathlineto{\pgfqpoint{6.362738in}{0.616344in}}%
\pgfpathlineto{\pgfqpoint{6.366369in}{0.610467in}}%
\pgfpathlineto{\pgfqpoint{6.370000in}{0.614516in}}%
\pgfpathlineto{\pgfqpoint{6.370000in}{0.614516in}}%
\pgfusepath{stroke}%
\end{pgfscope}%
\begin{pgfscope}%
\pgfpathrectangle{\pgfqpoint{4.554414in}{0.320625in}}{\pgfqpoint{1.815586in}{1.164375in}}%
\pgfusepath{clip}%
\pgfsetrectcap%
\pgfsetroundjoin%
\pgfsetlinewidth{1.154312pt}%
\definecolor{currentstroke}{rgb}{0.819608,0.286275,0.356863}%
\pgfsetstrokecolor{currentstroke}%
\pgfsetdash{}{0pt}%
\pgfpathmoveto{\pgfqpoint{4.554414in}{1.464269in}}%
\pgfpathlineto{\pgfqpoint{4.558045in}{1.452010in}}%
\pgfpathlineto{\pgfqpoint{4.561676in}{1.442160in}}%
\pgfpathlineto{\pgfqpoint{4.568938in}{1.426660in}}%
\pgfpathlineto{\pgfqpoint{4.623406in}{1.322281in}}%
\pgfpathlineto{\pgfqpoint{4.648824in}{1.276078in}}%
\pgfpathlineto{\pgfqpoint{4.674242in}{1.232440in}}%
\pgfpathlineto{\pgfqpoint{4.703292in}{1.185103in}}%
\pgfpathlineto{\pgfqpoint{4.735972in}{1.134172in}}%
\pgfpathlineto{\pgfqpoint{4.775915in}{1.074197in}}%
\pgfpathlineto{\pgfqpoint{4.823120in}{1.005506in}}%
\pgfpathlineto{\pgfqpoint{4.877588in}{0.928319in}}%
\pgfpathlineto{\pgfqpoint{4.939318in}{0.842866in}}%
\pgfpathlineto{\pgfqpoint{4.993785in}{0.769135in}}%
\pgfpathlineto{\pgfqpoint{5.030097in}{0.721985in}}%
\pgfpathlineto{\pgfqpoint{5.033728in}{0.715798in}}%
\pgfpathlineto{\pgfqpoint{5.037360in}{0.713012in}}%
\pgfpathlineto{\pgfqpoint{5.040991in}{0.712568in}}%
\pgfpathlineto{\pgfqpoint{5.044622in}{0.705038in}}%
\pgfpathlineto{\pgfqpoint{5.048253in}{0.703648in}}%
\pgfpathlineto{\pgfqpoint{5.051884in}{0.695742in}}%
\pgfpathlineto{\pgfqpoint{5.055515in}{0.691854in}}%
\pgfpathlineto{\pgfqpoint{5.059147in}{0.687025in}}%
\pgfpathlineto{\pgfqpoint{5.062778in}{0.683936in}}%
\pgfpathlineto{\pgfqpoint{5.066409in}{0.683245in}}%
\pgfpathlineto{\pgfqpoint{5.070040in}{0.683732in}}%
\pgfpathlineto{\pgfqpoint{5.073671in}{0.678529in}}%
\pgfpathlineto{\pgfqpoint{5.077302in}{0.664883in}}%
\pgfpathlineto{\pgfqpoint{5.080934in}{0.674132in}}%
\pgfpathlineto{\pgfqpoint{5.084565in}{0.659620in}}%
\pgfpathlineto{\pgfqpoint{5.088196in}{0.676796in}}%
\pgfpathlineto{\pgfqpoint{5.091827in}{0.683031in}}%
\pgfpathlineto{\pgfqpoint{5.095458in}{0.668394in}}%
\pgfpathlineto{\pgfqpoint{5.099090in}{0.670331in}}%
\pgfpathlineto{\pgfqpoint{5.102721in}{0.665316in}}%
\pgfpathlineto{\pgfqpoint{5.109983in}{0.665814in}}%
\pgfpathlineto{\pgfqpoint{5.113614in}{0.661899in}}%
\pgfpathlineto{\pgfqpoint{5.117245in}{0.659765in}}%
\pgfpathlineto{\pgfqpoint{5.120877in}{0.682672in}}%
\pgfpathlineto{\pgfqpoint{5.124508in}{0.672822in}}%
\pgfpathlineto{\pgfqpoint{5.128139in}{0.674076in}}%
\pgfpathlineto{\pgfqpoint{5.131770in}{0.652988in}}%
\pgfpathlineto{\pgfqpoint{5.135401in}{0.612124in}}%
\pgfpathlineto{\pgfqpoint{5.139032in}{0.670407in}}%
\pgfpathlineto{\pgfqpoint{5.142664in}{0.651434in}}%
\pgfpathlineto{\pgfqpoint{5.146295in}{0.668909in}}%
\pgfpathlineto{\pgfqpoint{5.149926in}{0.677103in}}%
\pgfpathlineto{\pgfqpoint{5.153557in}{0.669003in}}%
\pgfpathlineto{\pgfqpoint{5.157188in}{0.671905in}}%
\pgfpathlineto{\pgfqpoint{5.160819in}{0.645067in}}%
\pgfpathlineto{\pgfqpoint{5.164451in}{0.652226in}}%
\pgfpathlineto{\pgfqpoint{5.168082in}{0.664329in}}%
\pgfpathlineto{\pgfqpoint{5.171713in}{0.665626in}}%
\pgfpathlineto{\pgfqpoint{5.175344in}{0.681328in}}%
\pgfpathlineto{\pgfqpoint{5.178975in}{0.663074in}}%
\pgfpathlineto{\pgfqpoint{5.182606in}{0.672245in}}%
\pgfpathlineto{\pgfqpoint{5.186238in}{0.659019in}}%
\pgfpathlineto{\pgfqpoint{5.189869in}{0.633403in}}%
\pgfpathlineto{\pgfqpoint{5.193500in}{0.651425in}}%
\pgfpathlineto{\pgfqpoint{5.197131in}{0.682356in}}%
\pgfpathlineto{\pgfqpoint{5.200762in}{0.681775in}}%
\pgfpathlineto{\pgfqpoint{5.204394in}{0.676117in}}%
\pgfpathlineto{\pgfqpoint{5.208025in}{0.665562in}}%
\pgfpathlineto{\pgfqpoint{5.211656in}{0.672732in}}%
\pgfpathlineto{\pgfqpoint{5.218918in}{0.625321in}}%
\pgfpathlineto{\pgfqpoint{5.222549in}{0.675355in}}%
\pgfpathlineto{\pgfqpoint{5.226181in}{0.675756in}}%
\pgfpathlineto{\pgfqpoint{5.229812in}{0.666425in}}%
\pgfpathlineto{\pgfqpoint{5.233443in}{0.660429in}}%
\pgfpathlineto{\pgfqpoint{5.237074in}{0.657346in}}%
\pgfpathlineto{\pgfqpoint{5.240705in}{0.641137in}}%
\pgfpathlineto{\pgfqpoint{5.244336in}{0.658191in}}%
\pgfpathlineto{\pgfqpoint{5.247968in}{0.656692in}}%
\pgfpathlineto{\pgfqpoint{5.255230in}{0.667984in}}%
\pgfpathlineto{\pgfqpoint{5.258861in}{0.667345in}}%
\pgfpathlineto{\pgfqpoint{5.262492in}{0.681240in}}%
\pgfpathlineto{\pgfqpoint{5.266123in}{0.675465in}}%
\pgfpathlineto{\pgfqpoint{5.269755in}{0.641186in}}%
\pgfpathlineto{\pgfqpoint{5.273386in}{0.665139in}}%
\pgfpathlineto{\pgfqpoint{5.277017in}{0.627016in}}%
\pgfpathlineto{\pgfqpoint{5.280648in}{0.665627in}}%
\pgfpathlineto{\pgfqpoint{5.284279in}{0.662274in}}%
\pgfpathlineto{\pgfqpoint{5.287910in}{0.660823in}}%
\pgfpathlineto{\pgfqpoint{5.291542in}{0.662857in}}%
\pgfpathlineto{\pgfqpoint{5.295173in}{0.662777in}}%
\pgfpathlineto{\pgfqpoint{5.298804in}{0.663593in}}%
\pgfpathlineto{\pgfqpoint{5.302435in}{0.674695in}}%
\pgfpathlineto{\pgfqpoint{5.306066in}{0.674365in}}%
\pgfpathlineto{\pgfqpoint{5.313329in}{0.650711in}}%
\pgfpathlineto{\pgfqpoint{5.316960in}{0.644966in}}%
\pgfpathlineto{\pgfqpoint{5.320591in}{0.670651in}}%
\pgfpathlineto{\pgfqpoint{5.324222in}{0.681306in}}%
\pgfpathlineto{\pgfqpoint{5.327853in}{0.651559in}}%
\pgfpathlineto{\pgfqpoint{5.331485in}{0.670700in}}%
\pgfpathlineto{\pgfqpoint{5.335116in}{0.642095in}}%
\pgfpathlineto{\pgfqpoint{5.338747in}{0.658614in}}%
\pgfpathlineto{\pgfqpoint{5.342378in}{0.644681in}}%
\pgfpathlineto{\pgfqpoint{5.346009in}{0.669091in}}%
\pgfpathlineto{\pgfqpoint{5.349640in}{0.664007in}}%
\pgfpathlineto{\pgfqpoint{5.353272in}{0.651772in}}%
\pgfpathlineto{\pgfqpoint{5.356903in}{0.676312in}}%
\pgfpathlineto{\pgfqpoint{5.360534in}{0.676443in}}%
\pgfpathlineto{\pgfqpoint{5.364165in}{0.651133in}}%
\pgfpathlineto{\pgfqpoint{5.367796in}{0.644232in}}%
\pgfpathlineto{\pgfqpoint{5.371427in}{0.666160in}}%
\pgfpathlineto{\pgfqpoint{5.375059in}{0.668697in}}%
\pgfpathlineto{\pgfqpoint{5.378690in}{0.627628in}}%
\pgfpathlineto{\pgfqpoint{5.382321in}{0.683954in}}%
\pgfpathlineto{\pgfqpoint{5.385952in}{0.666011in}}%
\pgfpathlineto{\pgfqpoint{5.389583in}{0.662195in}}%
\pgfpathlineto{\pgfqpoint{5.396846in}{0.659520in}}%
\pgfpathlineto{\pgfqpoint{5.400477in}{0.652469in}}%
\pgfpathlineto{\pgfqpoint{5.404108in}{0.658837in}}%
\pgfpathlineto{\pgfqpoint{5.407739in}{0.660468in}}%
\pgfpathlineto{\pgfqpoint{5.411370in}{0.663684in}}%
\pgfpathlineto{\pgfqpoint{5.415002in}{0.670579in}}%
\pgfpathlineto{\pgfqpoint{5.418633in}{0.650449in}}%
\pgfpathlineto{\pgfqpoint{5.422264in}{0.674586in}}%
\pgfpathlineto{\pgfqpoint{5.425895in}{0.658557in}}%
\pgfpathlineto{\pgfqpoint{5.429526in}{0.669577in}}%
\pgfpathlineto{\pgfqpoint{5.433157in}{0.658652in}}%
\pgfpathlineto{\pgfqpoint{5.440420in}{0.661264in}}%
\pgfpathlineto{\pgfqpoint{5.444051in}{0.643852in}}%
\pgfpathlineto{\pgfqpoint{5.447682in}{0.664610in}}%
\pgfpathlineto{\pgfqpoint{5.451313in}{0.648063in}}%
\pgfpathlineto{\pgfqpoint{5.458576in}{0.656909in}}%
\pgfpathlineto{\pgfqpoint{5.462207in}{0.665696in}}%
\pgfpathlineto{\pgfqpoint{5.465838in}{0.667843in}}%
\pgfpathlineto{\pgfqpoint{5.469469in}{0.649891in}}%
\pgfpathlineto{\pgfqpoint{5.473100in}{0.657866in}}%
\pgfpathlineto{\pgfqpoint{5.476731in}{0.675630in}}%
\pgfpathlineto{\pgfqpoint{5.480363in}{0.672822in}}%
\pgfpathlineto{\pgfqpoint{5.483994in}{0.666176in}}%
\pgfpathlineto{\pgfqpoint{5.487625in}{0.679215in}}%
\pgfpathlineto{\pgfqpoint{5.491256in}{0.662475in}}%
\pgfpathlineto{\pgfqpoint{5.494887in}{0.677390in}}%
\pgfpathlineto{\pgfqpoint{5.498519in}{0.652069in}}%
\pgfpathlineto{\pgfqpoint{5.502150in}{0.632273in}}%
\pgfpathlineto{\pgfqpoint{5.505781in}{0.677203in}}%
\pgfpathlineto{\pgfqpoint{5.509412in}{0.668327in}}%
\pgfpathlineto{\pgfqpoint{5.513043in}{0.682707in}}%
\pgfpathlineto{\pgfqpoint{5.516674in}{0.681444in}}%
\pgfpathlineto{\pgfqpoint{5.520306in}{0.664539in}}%
\pgfpathlineto{\pgfqpoint{5.523937in}{0.651354in}}%
\pgfpathlineto{\pgfqpoint{5.527568in}{0.664039in}}%
\pgfpathlineto{\pgfqpoint{5.531199in}{0.671808in}}%
\pgfpathlineto{\pgfqpoint{5.534830in}{0.667321in}}%
\pgfpathlineto{\pgfqpoint{5.538461in}{0.676337in}}%
\pgfpathlineto{\pgfqpoint{5.542093in}{0.669897in}}%
\pgfpathlineto{\pgfqpoint{5.545724in}{0.685267in}}%
\pgfpathlineto{\pgfqpoint{5.549355in}{0.632499in}}%
\pgfpathlineto{\pgfqpoint{5.552986in}{0.644581in}}%
\pgfpathlineto{\pgfqpoint{5.556617in}{0.664366in}}%
\pgfpathlineto{\pgfqpoint{5.560248in}{0.648566in}}%
\pgfpathlineto{\pgfqpoint{5.563880in}{0.683545in}}%
\pgfpathlineto{\pgfqpoint{5.567511in}{0.660142in}}%
\pgfpathlineto{\pgfqpoint{5.571142in}{0.641460in}}%
\pgfpathlineto{\pgfqpoint{5.574773in}{0.653949in}}%
\pgfpathlineto{\pgfqpoint{5.578404in}{0.662142in}}%
\pgfpathlineto{\pgfqpoint{5.582035in}{0.639552in}}%
\pgfpathlineto{\pgfqpoint{5.585667in}{0.645134in}}%
\pgfpathlineto{\pgfqpoint{5.589298in}{0.656931in}}%
\pgfpathlineto{\pgfqpoint{5.592929in}{0.674009in}}%
\pgfpathlineto{\pgfqpoint{5.596560in}{0.669320in}}%
\pgfpathlineto{\pgfqpoint{5.600191in}{0.656162in}}%
\pgfpathlineto{\pgfqpoint{5.603823in}{0.667784in}}%
\pgfpathlineto{\pgfqpoint{5.607454in}{0.662155in}}%
\pgfpathlineto{\pgfqpoint{5.611085in}{0.661691in}}%
\pgfpathlineto{\pgfqpoint{5.614716in}{0.650175in}}%
\pgfpathlineto{\pgfqpoint{5.618347in}{0.645594in}}%
\pgfpathlineto{\pgfqpoint{5.621978in}{0.655042in}}%
\pgfpathlineto{\pgfqpoint{5.625610in}{0.622285in}}%
\pgfpathlineto{\pgfqpoint{5.629241in}{0.660780in}}%
\pgfpathlineto{\pgfqpoint{5.632872in}{0.672279in}}%
\pgfpathlineto{\pgfqpoint{5.636503in}{0.680949in}}%
\pgfpathlineto{\pgfqpoint{5.640134in}{0.683115in}}%
\pgfpathlineto{\pgfqpoint{5.643765in}{0.642236in}}%
\pgfpathlineto{\pgfqpoint{5.647397in}{0.665231in}}%
\pgfpathlineto{\pgfqpoint{5.651028in}{0.644001in}}%
\pgfpathlineto{\pgfqpoint{5.654659in}{0.654676in}}%
\pgfpathlineto{\pgfqpoint{5.658290in}{0.646876in}}%
\pgfpathlineto{\pgfqpoint{5.661921in}{0.654367in}}%
\pgfpathlineto{\pgfqpoint{5.665552in}{0.638514in}}%
\pgfpathlineto{\pgfqpoint{5.669184in}{0.663450in}}%
\pgfpathlineto{\pgfqpoint{5.672815in}{0.654754in}}%
\pgfpathlineto{\pgfqpoint{5.676446in}{0.677762in}}%
\pgfpathlineto{\pgfqpoint{5.680077in}{0.688840in}}%
\pgfpathlineto{\pgfqpoint{5.683708in}{0.657981in}}%
\pgfpathlineto{\pgfqpoint{5.687340in}{0.665088in}}%
\pgfpathlineto{\pgfqpoint{5.690971in}{0.666851in}}%
\pgfpathlineto{\pgfqpoint{5.694602in}{0.665475in}}%
\pgfpathlineto{\pgfqpoint{5.698233in}{0.673324in}}%
\pgfpathlineto{\pgfqpoint{5.701864in}{0.652809in}}%
\pgfpathlineto{\pgfqpoint{5.705495in}{0.662492in}}%
\pgfpathlineto{\pgfqpoint{5.709127in}{0.655934in}}%
\pgfpathlineto{\pgfqpoint{5.716389in}{0.675442in}}%
\pgfpathlineto{\pgfqpoint{5.720020in}{0.674982in}}%
\pgfpathlineto{\pgfqpoint{5.723651in}{0.666715in}}%
\pgfpathlineto{\pgfqpoint{5.727282in}{0.656290in}}%
\pgfpathlineto{\pgfqpoint{5.730914in}{0.648704in}}%
\pgfpathlineto{\pgfqpoint{5.734545in}{0.656551in}}%
\pgfpathlineto{\pgfqpoint{5.738176in}{0.665799in}}%
\pgfpathlineto{\pgfqpoint{5.741807in}{0.631167in}}%
\pgfpathlineto{\pgfqpoint{5.745438in}{0.652642in}}%
\pgfpathlineto{\pgfqpoint{5.749069in}{0.669491in}}%
\pgfpathlineto{\pgfqpoint{5.752701in}{0.677478in}}%
\pgfpathlineto{\pgfqpoint{5.756332in}{0.675939in}}%
\pgfpathlineto{\pgfqpoint{5.759963in}{0.649813in}}%
\pgfpathlineto{\pgfqpoint{5.763594in}{0.660251in}}%
\pgfpathlineto{\pgfqpoint{5.767225in}{0.679061in}}%
\pgfpathlineto{\pgfqpoint{5.770856in}{0.655538in}}%
\pgfpathlineto{\pgfqpoint{5.774488in}{0.648133in}}%
\pgfpathlineto{\pgfqpoint{5.778119in}{0.642313in}}%
\pgfpathlineto{\pgfqpoint{5.781750in}{0.657780in}}%
\pgfpathlineto{\pgfqpoint{5.785381in}{0.655650in}}%
\pgfpathlineto{\pgfqpoint{5.789012in}{0.669579in}}%
\pgfpathlineto{\pgfqpoint{5.792644in}{0.655871in}}%
\pgfpathlineto{\pgfqpoint{5.796275in}{0.682368in}}%
\pgfpathlineto{\pgfqpoint{5.803537in}{0.641239in}}%
\pgfpathlineto{\pgfqpoint{5.807168in}{0.642318in}}%
\pgfpathlineto{\pgfqpoint{5.810799in}{0.668753in}}%
\pgfpathlineto{\pgfqpoint{5.814431in}{0.667734in}}%
\pgfpathlineto{\pgfqpoint{5.818062in}{0.643756in}}%
\pgfpathlineto{\pgfqpoint{5.821693in}{0.663353in}}%
\pgfpathlineto{\pgfqpoint{5.825324in}{0.646999in}}%
\pgfpathlineto{\pgfqpoint{5.828955in}{0.654262in}}%
\pgfpathlineto{\pgfqpoint{5.832586in}{0.643043in}}%
\pgfpathlineto{\pgfqpoint{5.836218in}{0.675569in}}%
\pgfpathlineto{\pgfqpoint{5.839849in}{0.683561in}}%
\pgfpathlineto{\pgfqpoint{5.843480in}{0.656902in}}%
\pgfpathlineto{\pgfqpoint{5.847111in}{0.678788in}}%
\pgfpathlineto{\pgfqpoint{5.850742in}{0.649244in}}%
\pgfpathlineto{\pgfqpoint{5.854373in}{0.639857in}}%
\pgfpathlineto{\pgfqpoint{5.858005in}{0.643226in}}%
\pgfpathlineto{\pgfqpoint{5.861636in}{0.664006in}}%
\pgfpathlineto{\pgfqpoint{5.865267in}{0.679874in}}%
\pgfpathlineto{\pgfqpoint{5.868898in}{0.681825in}}%
\pgfpathlineto{\pgfqpoint{5.872529in}{0.685124in}}%
\pgfpathlineto{\pgfqpoint{5.876160in}{0.677036in}}%
\pgfpathlineto{\pgfqpoint{5.879792in}{0.663740in}}%
\pgfpathlineto{\pgfqpoint{5.883423in}{0.667580in}}%
\pgfpathlineto{\pgfqpoint{5.890685in}{0.656147in}}%
\pgfpathlineto{\pgfqpoint{5.894316in}{0.654155in}}%
\pgfpathlineto{\pgfqpoint{5.897948in}{0.657195in}}%
\pgfpathlineto{\pgfqpoint{5.901579in}{0.670791in}}%
\pgfpathlineto{\pgfqpoint{5.905210in}{0.672767in}}%
\pgfpathlineto{\pgfqpoint{5.908841in}{0.667549in}}%
\pgfpathlineto{\pgfqpoint{5.912472in}{0.691654in}}%
\pgfpathlineto{\pgfqpoint{5.916103in}{0.661615in}}%
\pgfpathlineto{\pgfqpoint{5.919735in}{0.670171in}}%
\pgfpathlineto{\pgfqpoint{5.923366in}{0.651480in}}%
\pgfpathlineto{\pgfqpoint{5.926997in}{0.645001in}}%
\pgfpathlineto{\pgfqpoint{5.930628in}{0.665710in}}%
\pgfpathlineto{\pgfqpoint{5.934259in}{0.662617in}}%
\pgfpathlineto{\pgfqpoint{5.937890in}{0.676309in}}%
\pgfpathlineto{\pgfqpoint{5.941522in}{0.685695in}}%
\pgfpathlineto{\pgfqpoint{5.945153in}{0.659015in}}%
\pgfpathlineto{\pgfqpoint{5.948784in}{0.638816in}}%
\pgfpathlineto{\pgfqpoint{5.952415in}{0.669847in}}%
\pgfpathlineto{\pgfqpoint{5.956046in}{0.646618in}}%
\pgfpathlineto{\pgfqpoint{5.959677in}{0.636595in}}%
\pgfpathlineto{\pgfqpoint{5.963309in}{0.654740in}}%
\pgfpathlineto{\pgfqpoint{5.966940in}{0.662446in}}%
\pgfpathlineto{\pgfqpoint{5.970571in}{0.684177in}}%
\pgfpathlineto{\pgfqpoint{5.974202in}{0.653056in}}%
\pgfpathlineto{\pgfqpoint{5.977833in}{0.663098in}}%
\pgfpathlineto{\pgfqpoint{5.981465in}{0.668982in}}%
\pgfpathlineto{\pgfqpoint{5.985096in}{0.654599in}}%
\pgfpathlineto{\pgfqpoint{5.988727in}{0.645728in}}%
\pgfpathlineto{\pgfqpoint{5.992358in}{0.655054in}}%
\pgfpathlineto{\pgfqpoint{5.995989in}{0.641452in}}%
\pgfpathlineto{\pgfqpoint{5.999620in}{0.653785in}}%
\pgfpathlineto{\pgfqpoint{6.003252in}{0.659110in}}%
\pgfpathlineto{\pgfqpoint{6.006883in}{0.659614in}}%
\pgfpathlineto{\pgfqpoint{6.010514in}{0.677786in}}%
\pgfpathlineto{\pgfqpoint{6.014145in}{0.678999in}}%
\pgfpathlineto{\pgfqpoint{6.017776in}{0.652382in}}%
\pgfpathlineto{\pgfqpoint{6.021407in}{0.667520in}}%
\pgfpathlineto{\pgfqpoint{6.025039in}{0.651090in}}%
\pgfpathlineto{\pgfqpoint{6.028670in}{0.668524in}}%
\pgfpathlineto{\pgfqpoint{6.032301in}{0.639260in}}%
\pgfpathlineto{\pgfqpoint{6.035932in}{0.668347in}}%
\pgfpathlineto{\pgfqpoint{6.039563in}{0.666552in}}%
\pgfpathlineto{\pgfqpoint{6.043194in}{0.655363in}}%
\pgfpathlineto{\pgfqpoint{6.046826in}{0.671527in}}%
\pgfpathlineto{\pgfqpoint{6.050457in}{0.672052in}}%
\pgfpathlineto{\pgfqpoint{6.061350in}{0.646564in}}%
\pgfpathlineto{\pgfqpoint{6.064981in}{0.611999in}}%
\pgfpathlineto{\pgfqpoint{6.068613in}{0.676897in}}%
\pgfpathlineto{\pgfqpoint{6.072244in}{0.673880in}}%
\pgfpathlineto{\pgfqpoint{6.075875in}{0.658554in}}%
\pgfpathlineto{\pgfqpoint{6.079506in}{0.661897in}}%
\pgfpathlineto{\pgfqpoint{6.083137in}{0.674537in}}%
\pgfpathlineto{\pgfqpoint{6.086769in}{0.680958in}}%
\pgfpathlineto{\pgfqpoint{6.090400in}{0.667909in}}%
\pgfpathlineto{\pgfqpoint{6.094031in}{0.660732in}}%
\pgfpathlineto{\pgfqpoint{6.097662in}{0.640358in}}%
\pgfpathlineto{\pgfqpoint{6.101293in}{0.649651in}}%
\pgfpathlineto{\pgfqpoint{6.104924in}{0.654954in}}%
\pgfpathlineto{\pgfqpoint{6.108556in}{0.667475in}}%
\pgfpathlineto{\pgfqpoint{6.112187in}{0.683006in}}%
\pgfpathlineto{\pgfqpoint{6.115818in}{0.677245in}}%
\pgfpathlineto{\pgfqpoint{6.119449in}{0.645891in}}%
\pgfpathlineto{\pgfqpoint{6.123080in}{0.660696in}}%
\pgfpathlineto{\pgfqpoint{6.126711in}{0.662996in}}%
\pgfpathlineto{\pgfqpoint{6.130343in}{0.676179in}}%
\pgfpathlineto{\pgfqpoint{6.133974in}{0.657206in}}%
\pgfpathlineto{\pgfqpoint{6.141236in}{0.682257in}}%
\pgfpathlineto{\pgfqpoint{6.144867in}{0.681670in}}%
\pgfpathlineto{\pgfqpoint{6.148498in}{0.651697in}}%
\pgfpathlineto{\pgfqpoint{6.152130in}{0.649653in}}%
\pgfpathlineto{\pgfqpoint{6.155761in}{0.661818in}}%
\pgfpathlineto{\pgfqpoint{6.159392in}{0.658197in}}%
\pgfpathlineto{\pgfqpoint{6.163023in}{0.667380in}}%
\pgfpathlineto{\pgfqpoint{6.166654in}{0.641253in}}%
\pgfpathlineto{\pgfqpoint{6.170285in}{0.683459in}}%
\pgfpathlineto{\pgfqpoint{6.173917in}{0.679389in}}%
\pgfpathlineto{\pgfqpoint{6.177548in}{0.625839in}}%
\pgfpathlineto{\pgfqpoint{6.181179in}{0.662574in}}%
\pgfpathlineto{\pgfqpoint{6.184810in}{0.669939in}}%
\pgfpathlineto{\pgfqpoint{6.188441in}{0.668487in}}%
\pgfpathlineto{\pgfqpoint{6.192073in}{0.648918in}}%
\pgfpathlineto{\pgfqpoint{6.195704in}{0.666344in}}%
\pgfpathlineto{\pgfqpoint{6.199335in}{0.650734in}}%
\pgfpathlineto{\pgfqpoint{6.202966in}{0.664583in}}%
\pgfpathlineto{\pgfqpoint{6.206597in}{0.670027in}}%
\pgfpathlineto{\pgfqpoint{6.210228in}{0.654889in}}%
\pgfpathlineto{\pgfqpoint{6.213860in}{0.674496in}}%
\pgfpathlineto{\pgfqpoint{6.217491in}{0.669684in}}%
\pgfpathlineto{\pgfqpoint{6.224753in}{0.654565in}}%
\pgfpathlineto{\pgfqpoint{6.228384in}{0.641453in}}%
\pgfpathlineto{\pgfqpoint{6.232015in}{0.655896in}}%
\pgfpathlineto{\pgfqpoint{6.235647in}{0.672943in}}%
\pgfpathlineto{\pgfqpoint{6.239278in}{0.665434in}}%
\pgfpathlineto{\pgfqpoint{6.242909in}{0.686504in}}%
\pgfpathlineto{\pgfqpoint{6.246540in}{0.680712in}}%
\pgfpathlineto{\pgfqpoint{6.250171in}{0.641717in}}%
\pgfpathlineto{\pgfqpoint{6.257434in}{0.652613in}}%
\pgfpathlineto{\pgfqpoint{6.261065in}{0.662947in}}%
\pgfpathlineto{\pgfqpoint{6.264696in}{0.675393in}}%
\pgfpathlineto{\pgfqpoint{6.268327in}{0.652289in}}%
\pgfpathlineto{\pgfqpoint{6.271958in}{0.665446in}}%
\pgfpathlineto{\pgfqpoint{6.275590in}{0.659874in}}%
\pgfpathlineto{\pgfqpoint{6.279221in}{0.652008in}}%
\pgfpathlineto{\pgfqpoint{6.282852in}{0.640735in}}%
\pgfpathlineto{\pgfqpoint{6.286483in}{0.678661in}}%
\pgfpathlineto{\pgfqpoint{6.290114in}{0.676759in}}%
\pgfpathlineto{\pgfqpoint{6.293745in}{0.679563in}}%
\pgfpathlineto{\pgfqpoint{6.297377in}{0.660060in}}%
\pgfpathlineto{\pgfqpoint{6.301008in}{0.678658in}}%
\pgfpathlineto{\pgfqpoint{6.304639in}{0.658592in}}%
\pgfpathlineto{\pgfqpoint{6.308270in}{0.620628in}}%
\pgfpathlineto{\pgfqpoint{6.311901in}{0.667088in}}%
\pgfpathlineto{\pgfqpoint{6.319164in}{0.666237in}}%
\pgfpathlineto{\pgfqpoint{6.322795in}{0.661209in}}%
\pgfpathlineto{\pgfqpoint{6.326426in}{0.663092in}}%
\pgfpathlineto{\pgfqpoint{6.330057in}{0.661737in}}%
\pgfpathlineto{\pgfqpoint{6.333688in}{0.608000in}}%
\pgfpathlineto{\pgfqpoint{6.337319in}{0.645727in}}%
\pgfpathlineto{\pgfqpoint{6.340951in}{0.653655in}}%
\pgfpathlineto{\pgfqpoint{6.344582in}{0.676861in}}%
\pgfpathlineto{\pgfqpoint{6.348213in}{0.670558in}}%
\pgfpathlineto{\pgfqpoint{6.355475in}{0.655308in}}%
\pgfpathlineto{\pgfqpoint{6.359106in}{0.657397in}}%
\pgfpathlineto{\pgfqpoint{6.362738in}{0.645207in}}%
\pgfpathlineto{\pgfqpoint{6.366369in}{0.653686in}}%
\pgfpathlineto{\pgfqpoint{6.370000in}{0.616122in}}%
\pgfpathlineto{\pgfqpoint{6.370000in}{0.616122in}}%
\pgfusepath{stroke}%
\end{pgfscope}%
\begin{pgfscope}%
\pgfpathrectangle{\pgfqpoint{4.554414in}{0.320625in}}{\pgfqpoint{1.815586in}{1.164375in}}%
\pgfusepath{clip}%
\pgfsetrectcap%
\pgfsetroundjoin%
\pgfsetlinewidth{0.602250pt}%
\definecolor{currentstroke}{rgb}{0.000000,0.529412,0.423529}%
\pgfsetstrokecolor{currentstroke}%
\pgfsetstrokeopacity{0.120000}%
\pgfsetdash{}{0pt}%
\pgfpathmoveto{\pgfqpoint{4.554414in}{1.455473in}}%
\pgfpathlineto{\pgfqpoint{4.558045in}{1.444109in}}%
\pgfpathlineto{\pgfqpoint{4.565307in}{1.428595in}}%
\pgfpathlineto{\pgfqpoint{4.583463in}{1.391350in}}%
\pgfpathlineto{\pgfqpoint{4.612512in}{1.327957in}}%
\pgfpathlineto{\pgfqpoint{4.652455in}{1.241432in}}%
\pgfpathlineto{\pgfqpoint{4.685136in}{1.173328in}}%
\pgfpathlineto{\pgfqpoint{4.721448in}{1.100335in}}%
\pgfpathlineto{\pgfqpoint{4.761390in}{1.022669in}}%
\pgfpathlineto{\pgfqpoint{4.804965in}{0.940405in}}%
\pgfpathlineto{\pgfqpoint{4.852170in}{0.853631in}}%
\pgfpathlineto{\pgfqpoint{4.903006in}{0.762549in}}%
\pgfpathlineto{\pgfqpoint{4.921162in}{0.731007in}}%
\pgfpathlineto{\pgfqpoint{4.924793in}{0.726834in}}%
\pgfpathlineto{\pgfqpoint{4.928424in}{0.721760in}}%
\pgfpathlineto{\pgfqpoint{4.935687in}{0.705354in}}%
\pgfpathlineto{\pgfqpoint{4.939318in}{0.698975in}}%
\pgfpathlineto{\pgfqpoint{4.942949in}{0.695769in}}%
\pgfpathlineto{\pgfqpoint{4.946580in}{0.687416in}}%
\pgfpathlineto{\pgfqpoint{4.950211in}{0.685691in}}%
\pgfpathlineto{\pgfqpoint{4.953843in}{0.701380in}}%
\pgfpathlineto{\pgfqpoint{4.957474in}{0.678486in}}%
\pgfpathlineto{\pgfqpoint{4.961105in}{0.661555in}}%
\pgfpathlineto{\pgfqpoint{4.964736in}{0.672892in}}%
\pgfpathlineto{\pgfqpoint{4.968367in}{0.654007in}}%
\pgfpathlineto{\pgfqpoint{4.971998in}{0.659319in}}%
\pgfpathlineto{\pgfqpoint{4.975630in}{0.643442in}}%
\pgfpathlineto{\pgfqpoint{4.979261in}{0.641354in}}%
\pgfpathlineto{\pgfqpoint{4.982892in}{0.641672in}}%
\pgfpathlineto{\pgfqpoint{4.986523in}{0.644009in}}%
\pgfpathlineto{\pgfqpoint{4.990154in}{0.651708in}}%
\pgfpathlineto{\pgfqpoint{4.993785in}{0.633354in}}%
\pgfpathlineto{\pgfqpoint{4.997417in}{0.639712in}}%
\pgfpathlineto{\pgfqpoint{5.001048in}{0.676951in}}%
\pgfpathlineto{\pgfqpoint{5.004679in}{0.634257in}}%
\pgfpathlineto{\pgfqpoint{5.008310in}{0.580980in}}%
\pgfpathlineto{\pgfqpoint{5.011941in}{0.650183in}}%
\pgfpathlineto{\pgfqpoint{5.015573in}{0.572203in}}%
\pgfpathlineto{\pgfqpoint{5.019204in}{0.667489in}}%
\pgfpathlineto{\pgfqpoint{5.022835in}{0.690580in}}%
\pgfpathlineto{\pgfqpoint{5.026466in}{0.565523in}}%
\pgfpathlineto{\pgfqpoint{5.030097in}{0.613246in}}%
\pgfpathlineto{\pgfqpoint{5.033728in}{0.619920in}}%
\pgfpathlineto{\pgfqpoint{5.037360in}{0.630858in}}%
\pgfpathlineto{\pgfqpoint{5.040991in}{0.676178in}}%
\pgfpathlineto{\pgfqpoint{5.044622in}{0.682521in}}%
\pgfpathlineto{\pgfqpoint{5.048253in}{0.639752in}}%
\pgfpathlineto{\pgfqpoint{5.051884in}{0.620369in}}%
\pgfpathlineto{\pgfqpoint{5.055515in}{0.573459in}}%
\pgfpathlineto{\pgfqpoint{5.059147in}{0.658946in}}%
\pgfpathlineto{\pgfqpoint{5.062778in}{0.670221in}}%
\pgfpathlineto{\pgfqpoint{5.066409in}{0.649367in}}%
\pgfpathlineto{\pgfqpoint{5.070040in}{0.668600in}}%
\pgfpathlineto{\pgfqpoint{5.073671in}{0.581632in}}%
\pgfpathlineto{\pgfqpoint{5.077302in}{0.634466in}}%
\pgfpathlineto{\pgfqpoint{5.080934in}{0.663323in}}%
\pgfpathlineto{\pgfqpoint{5.084565in}{0.640428in}}%
\pgfpathlineto{\pgfqpoint{5.088196in}{0.555949in}}%
\pgfpathlineto{\pgfqpoint{5.091827in}{0.669535in}}%
\pgfpathlineto{\pgfqpoint{5.095458in}{0.608544in}}%
\pgfpathlineto{\pgfqpoint{5.099090in}{0.646819in}}%
\pgfpathlineto{\pgfqpoint{5.102721in}{0.635853in}}%
\pgfpathlineto{\pgfqpoint{5.106352in}{0.612631in}}%
\pgfpathlineto{\pgfqpoint{5.109983in}{0.677285in}}%
\pgfpathlineto{\pgfqpoint{5.113614in}{0.693857in}}%
\pgfpathlineto{\pgfqpoint{5.117245in}{0.640717in}}%
\pgfpathlineto{\pgfqpoint{5.120877in}{0.619473in}}%
\pgfpathlineto{\pgfqpoint{5.124508in}{0.489078in}}%
\pgfpathlineto{\pgfqpoint{5.128139in}{0.655551in}}%
\pgfpathlineto{\pgfqpoint{5.131770in}{0.646513in}}%
\pgfpathlineto{\pgfqpoint{5.135401in}{0.622056in}}%
\pgfpathlineto{\pgfqpoint{5.139032in}{0.685274in}}%
\pgfpathlineto{\pgfqpoint{5.142664in}{0.668706in}}%
\pgfpathlineto{\pgfqpoint{5.146295in}{0.631395in}}%
\pgfpathlineto{\pgfqpoint{5.149926in}{0.668702in}}%
\pgfpathlineto{\pgfqpoint{5.153557in}{0.592236in}}%
\pgfpathlineto{\pgfqpoint{5.157188in}{0.646890in}}%
\pgfpathlineto{\pgfqpoint{5.160819in}{0.647370in}}%
\pgfpathlineto{\pgfqpoint{5.164451in}{0.560271in}}%
\pgfpathlineto{\pgfqpoint{5.168082in}{0.613415in}}%
\pgfpathlineto{\pgfqpoint{5.171713in}{0.653547in}}%
\pgfpathlineto{\pgfqpoint{5.175344in}{0.639954in}}%
\pgfpathlineto{\pgfqpoint{5.178975in}{0.638157in}}%
\pgfpathlineto{\pgfqpoint{5.182606in}{0.585174in}}%
\pgfpathlineto{\pgfqpoint{5.186238in}{0.661313in}}%
\pgfpathlineto{\pgfqpoint{5.189869in}{0.637601in}}%
\pgfpathlineto{\pgfqpoint{5.193500in}{0.603049in}}%
\pgfpathlineto{\pgfqpoint{5.197131in}{0.653480in}}%
\pgfpathlineto{\pgfqpoint{5.200762in}{0.607822in}}%
\pgfpathlineto{\pgfqpoint{5.204394in}{0.655172in}}%
\pgfpathlineto{\pgfqpoint{5.208025in}{0.686083in}}%
\pgfpathlineto{\pgfqpoint{5.211656in}{0.613776in}}%
\pgfpathlineto{\pgfqpoint{5.215287in}{0.592563in}}%
\pgfpathlineto{\pgfqpoint{5.218918in}{0.640070in}}%
\pgfpathlineto{\pgfqpoint{5.222549in}{0.597901in}}%
\pgfpathlineto{\pgfqpoint{5.226181in}{0.673172in}}%
\pgfpathlineto{\pgfqpoint{5.229812in}{0.688861in}}%
\pgfpathlineto{\pgfqpoint{5.233443in}{0.606792in}}%
\pgfpathlineto{\pgfqpoint{5.237074in}{0.622728in}}%
\pgfpathlineto{\pgfqpoint{5.240705in}{0.573770in}}%
\pgfpathlineto{\pgfqpoint{5.244336in}{0.649572in}}%
\pgfpathlineto{\pgfqpoint{5.247968in}{0.674235in}}%
\pgfpathlineto{\pgfqpoint{5.255230in}{0.659519in}}%
\pgfpathlineto{\pgfqpoint{5.258861in}{0.610885in}}%
\pgfpathlineto{\pgfqpoint{5.262492in}{0.617673in}}%
\pgfpathlineto{\pgfqpoint{5.266123in}{0.664542in}}%
\pgfpathlineto{\pgfqpoint{5.269755in}{0.658216in}}%
\pgfpathlineto{\pgfqpoint{5.273386in}{0.592493in}}%
\pgfpathlineto{\pgfqpoint{5.277017in}{0.691394in}}%
\pgfpathlineto{\pgfqpoint{5.280648in}{0.636695in}}%
\pgfpathlineto{\pgfqpoint{5.284279in}{0.539421in}}%
\pgfpathlineto{\pgfqpoint{5.287910in}{0.646379in}}%
\pgfpathlineto{\pgfqpoint{5.291542in}{0.617774in}}%
\pgfpathlineto{\pgfqpoint{5.295173in}{0.662786in}}%
\pgfpathlineto{\pgfqpoint{5.298804in}{0.664706in}}%
\pgfpathlineto{\pgfqpoint{5.302435in}{0.546788in}}%
\pgfpathlineto{\pgfqpoint{5.306066in}{0.660256in}}%
\pgfpathlineto{\pgfqpoint{5.309698in}{0.687870in}}%
\pgfpathlineto{\pgfqpoint{5.316960in}{0.609517in}}%
\pgfpathlineto{\pgfqpoint{5.320591in}{0.612292in}}%
\pgfpathlineto{\pgfqpoint{5.324222in}{0.662153in}}%
\pgfpathlineto{\pgfqpoint{5.327853in}{0.672108in}}%
\pgfpathlineto{\pgfqpoint{5.331485in}{0.650529in}}%
\pgfpathlineto{\pgfqpoint{5.335116in}{0.670678in}}%
\pgfpathlineto{\pgfqpoint{5.338747in}{0.428668in}}%
\pgfpathlineto{\pgfqpoint{5.342378in}{0.638791in}}%
\pgfpathlineto{\pgfqpoint{5.346009in}{0.654758in}}%
\pgfpathlineto{\pgfqpoint{5.349640in}{0.629513in}}%
\pgfpathlineto{\pgfqpoint{5.353272in}{0.567866in}}%
\pgfpathlineto{\pgfqpoint{5.356903in}{0.664504in}}%
\pgfpathlineto{\pgfqpoint{5.360534in}{0.623412in}}%
\pgfpathlineto{\pgfqpoint{5.364165in}{0.644519in}}%
\pgfpathlineto{\pgfqpoint{5.367796in}{0.607798in}}%
\pgfpathlineto{\pgfqpoint{5.371427in}{0.631875in}}%
\pgfpathlineto{\pgfqpoint{5.375059in}{0.635427in}}%
\pgfpathlineto{\pgfqpoint{5.378690in}{0.634433in}}%
\pgfpathlineto{\pgfqpoint{5.382321in}{0.650428in}}%
\pgfpathlineto{\pgfqpoint{5.385952in}{0.630102in}}%
\pgfpathlineto{\pgfqpoint{5.389583in}{0.631978in}}%
\pgfpathlineto{\pgfqpoint{5.393215in}{0.676978in}}%
\pgfpathlineto{\pgfqpoint{5.396846in}{0.632694in}}%
\pgfpathlineto{\pgfqpoint{5.400477in}{0.468099in}}%
\pgfpathlineto{\pgfqpoint{5.404108in}{0.651976in}}%
\pgfpathlineto{\pgfqpoint{5.407739in}{0.534904in}}%
\pgfpathlineto{\pgfqpoint{5.411370in}{0.666025in}}%
\pgfpathlineto{\pgfqpoint{5.415002in}{0.689597in}}%
\pgfpathlineto{\pgfqpoint{5.418633in}{0.531126in}}%
\pgfpathlineto{\pgfqpoint{5.422264in}{0.617490in}}%
\pgfpathlineto{\pgfqpoint{5.425895in}{0.622209in}}%
\pgfpathlineto{\pgfqpoint{5.429526in}{0.632470in}}%
\pgfpathlineto{\pgfqpoint{5.433157in}{0.674716in}}%
\pgfpathlineto{\pgfqpoint{5.436789in}{0.679221in}}%
\pgfpathlineto{\pgfqpoint{5.440420in}{0.642961in}}%
\pgfpathlineto{\pgfqpoint{5.444051in}{0.623079in}}%
\pgfpathlineto{\pgfqpoint{5.447682in}{0.572970in}}%
\pgfpathlineto{\pgfqpoint{5.451313in}{0.660437in}}%
\pgfpathlineto{\pgfqpoint{5.454944in}{0.667868in}}%
\pgfpathlineto{\pgfqpoint{5.458576in}{0.618372in}}%
\pgfpathlineto{\pgfqpoint{5.462207in}{0.679286in}}%
\pgfpathlineto{\pgfqpoint{5.465838in}{0.649177in}}%
\pgfpathlineto{\pgfqpoint{5.469469in}{0.590568in}}%
\pgfpathlineto{\pgfqpoint{5.473100in}{0.657111in}}%
\pgfpathlineto{\pgfqpoint{5.476731in}{0.573631in}}%
\pgfpathlineto{\pgfqpoint{5.480363in}{0.659961in}}%
\pgfpathlineto{\pgfqpoint{5.483994in}{0.674908in}}%
\pgfpathlineto{\pgfqpoint{5.487625in}{0.639786in}}%
\pgfpathlineto{\pgfqpoint{5.491256in}{0.639955in}}%
\pgfpathlineto{\pgfqpoint{5.494887in}{0.694252in}}%
\pgfpathlineto{\pgfqpoint{5.498519in}{0.628296in}}%
\pgfpathlineto{\pgfqpoint{5.502150in}{0.618365in}}%
\pgfpathlineto{\pgfqpoint{5.505781in}{0.556027in}}%
\pgfpathlineto{\pgfqpoint{5.509412in}{0.655415in}}%
\pgfpathlineto{\pgfqpoint{5.513043in}{0.677434in}}%
\pgfpathlineto{\pgfqpoint{5.516674in}{0.668686in}}%
\pgfpathlineto{\pgfqpoint{5.520306in}{0.664031in}}%
\pgfpathlineto{\pgfqpoint{5.523937in}{0.584712in}}%
\pgfpathlineto{\pgfqpoint{5.527568in}{0.629111in}}%
\pgfpathlineto{\pgfqpoint{5.531199in}{0.661105in}}%
\pgfpathlineto{\pgfqpoint{5.534830in}{0.654258in}}%
\pgfpathlineto{\pgfqpoint{5.538461in}{0.594556in}}%
\pgfpathlineto{\pgfqpoint{5.542093in}{0.671290in}}%
\pgfpathlineto{\pgfqpoint{5.545724in}{0.592633in}}%
\pgfpathlineto{\pgfqpoint{5.549355in}{0.644696in}}%
\pgfpathlineto{\pgfqpoint{5.552986in}{0.642356in}}%
\pgfpathlineto{\pgfqpoint{5.556617in}{0.507729in}}%
\pgfpathlineto{\pgfqpoint{5.560248in}{0.624026in}}%
\pgfpathlineto{\pgfqpoint{5.563880in}{0.652140in}}%
\pgfpathlineto{\pgfqpoint{5.567511in}{0.641579in}}%
\pgfpathlineto{\pgfqpoint{5.571142in}{0.640606in}}%
\pgfpathlineto{\pgfqpoint{5.574773in}{0.568783in}}%
\pgfpathlineto{\pgfqpoint{5.578404in}{0.662639in}}%
\pgfpathlineto{\pgfqpoint{5.582035in}{0.639383in}}%
\pgfpathlineto{\pgfqpoint{5.585667in}{0.597851in}}%
\pgfpathlineto{\pgfqpoint{5.589298in}{0.655502in}}%
\pgfpathlineto{\pgfqpoint{5.592929in}{0.608540in}}%
\pgfpathlineto{\pgfqpoint{5.596560in}{0.653930in}}%
\pgfpathlineto{\pgfqpoint{5.600191in}{0.685944in}}%
\pgfpathlineto{\pgfqpoint{5.603823in}{0.611500in}}%
\pgfpathlineto{\pgfqpoint{5.607454in}{0.600699in}}%
\pgfpathlineto{\pgfqpoint{5.611085in}{0.642230in}}%
\pgfpathlineto{\pgfqpoint{5.614716in}{0.600545in}}%
\pgfpathlineto{\pgfqpoint{5.618347in}{0.671984in}}%
\pgfpathlineto{\pgfqpoint{5.621978in}{0.686744in}}%
\pgfpathlineto{\pgfqpoint{5.625610in}{0.616916in}}%
\pgfpathlineto{\pgfqpoint{5.629241in}{0.633208in}}%
\pgfpathlineto{\pgfqpoint{5.632872in}{0.670102in}}%
\pgfpathlineto{\pgfqpoint{5.640134in}{0.623724in}}%
\pgfpathlineto{\pgfqpoint{5.643765in}{0.615486in}}%
\pgfpathlineto{\pgfqpoint{5.647397in}{0.667521in}}%
\pgfpathlineto{\pgfqpoint{5.651028in}{0.683906in}}%
\pgfpathlineto{\pgfqpoint{5.654659in}{0.668745in}}%
\pgfpathlineto{\pgfqpoint{5.658290in}{0.677717in}}%
\pgfpathlineto{\pgfqpoint{5.661921in}{0.612183in}}%
\pgfpathlineto{\pgfqpoint{5.665552in}{0.576240in}}%
\pgfpathlineto{\pgfqpoint{5.669184in}{0.641069in}}%
\pgfpathlineto{\pgfqpoint{5.672815in}{0.561497in}}%
\pgfpathlineto{\pgfqpoint{5.676446in}{0.673979in}}%
\pgfpathlineto{\pgfqpoint{5.680077in}{0.693752in}}%
\pgfpathlineto{\pgfqpoint{5.683708in}{0.565112in}}%
\pgfpathlineto{\pgfqpoint{5.687340in}{0.611988in}}%
\pgfpathlineto{\pgfqpoint{5.690971in}{0.584487in}}%
\pgfpathlineto{\pgfqpoint{5.694602in}{0.642347in}}%
\pgfpathlineto{\pgfqpoint{5.698233in}{0.677878in}}%
\pgfpathlineto{\pgfqpoint{5.701864in}{0.678811in}}%
\pgfpathlineto{\pgfqpoint{5.705495in}{0.650866in}}%
\pgfpathlineto{\pgfqpoint{5.709127in}{0.610507in}}%
\pgfpathlineto{\pgfqpoint{5.712758in}{0.608023in}}%
\pgfpathlineto{\pgfqpoint{5.716389in}{0.661288in}}%
\pgfpathlineto{\pgfqpoint{5.720020in}{0.667620in}}%
\pgfpathlineto{\pgfqpoint{5.723651in}{0.639806in}}%
\pgfpathlineto{\pgfqpoint{5.727282in}{0.671078in}}%
\pgfpathlineto{\pgfqpoint{5.730914in}{0.490489in}}%
\pgfpathlineto{\pgfqpoint{5.734545in}{0.639969in}}%
\pgfpathlineto{\pgfqpoint{5.738176in}{0.657631in}}%
\pgfpathlineto{\pgfqpoint{5.745438in}{0.590875in}}%
\pgfpathlineto{\pgfqpoint{5.749069in}{0.664394in}}%
\pgfpathlineto{\pgfqpoint{5.752701in}{0.624645in}}%
\pgfpathlineto{\pgfqpoint{5.756332in}{0.645499in}}%
\pgfpathlineto{\pgfqpoint{5.759963in}{0.610473in}}%
\pgfpathlineto{\pgfqpoint{5.763594in}{0.636556in}}%
\pgfpathlineto{\pgfqpoint{5.767225in}{0.638773in}}%
\pgfpathlineto{\pgfqpoint{5.770856in}{0.632055in}}%
\pgfpathlineto{\pgfqpoint{5.774488in}{0.652350in}}%
\pgfpathlineto{\pgfqpoint{5.778119in}{0.631137in}}%
\pgfpathlineto{\pgfqpoint{5.781750in}{0.630418in}}%
\pgfpathlineto{\pgfqpoint{5.785381in}{0.677749in}}%
\pgfpathlineto{\pgfqpoint{5.789012in}{0.632577in}}%
\pgfpathlineto{\pgfqpoint{5.792644in}{0.536579in}}%
\pgfpathlineto{\pgfqpoint{5.796275in}{0.653843in}}%
\pgfpathlineto{\pgfqpoint{5.799906in}{0.526672in}}%
\pgfpathlineto{\pgfqpoint{5.803537in}{0.665131in}}%
\pgfpathlineto{\pgfqpoint{5.807168in}{0.688680in}}%
\pgfpathlineto{\pgfqpoint{5.810799in}{0.653826in}}%
\pgfpathlineto{\pgfqpoint{5.814431in}{0.592231in}}%
\pgfpathlineto{\pgfqpoint{5.818062in}{0.679465in}}%
\pgfpathlineto{\pgfqpoint{5.821693in}{0.629586in}}%
\pgfpathlineto{\pgfqpoint{5.825324in}{0.631571in}}%
\pgfpathlineto{\pgfqpoint{5.828955in}{0.504750in}}%
\pgfpathlineto{\pgfqpoint{5.832586in}{0.654520in}}%
\pgfpathlineto{\pgfqpoint{5.836218in}{0.685106in}}%
\pgfpathlineto{\pgfqpoint{5.839849in}{0.684631in}}%
\pgfpathlineto{\pgfqpoint{5.843480in}{0.658367in}}%
\pgfpathlineto{\pgfqpoint{5.847111in}{0.624602in}}%
\pgfpathlineto{\pgfqpoint{5.850742in}{0.539035in}}%
\pgfpathlineto{\pgfqpoint{5.854373in}{0.650769in}}%
\pgfpathlineto{\pgfqpoint{5.858005in}{0.585726in}}%
\pgfpathlineto{\pgfqpoint{5.861636in}{0.664437in}}%
\pgfpathlineto{\pgfqpoint{5.865267in}{0.690895in}}%
\pgfpathlineto{\pgfqpoint{5.868898in}{0.588958in}}%
\pgfpathlineto{\pgfqpoint{5.872529in}{0.600487in}}%
\pgfpathlineto{\pgfqpoint{5.876160in}{0.623096in}}%
\pgfpathlineto{\pgfqpoint{5.879792in}{0.622523in}}%
\pgfpathlineto{\pgfqpoint{5.883423in}{0.676993in}}%
\pgfpathlineto{\pgfqpoint{5.887054in}{0.687385in}}%
\pgfpathlineto{\pgfqpoint{5.890685in}{0.628751in}}%
\pgfpathlineto{\pgfqpoint{5.894316in}{0.618613in}}%
\pgfpathlineto{\pgfqpoint{5.897948in}{0.540710in}}%
\pgfpathlineto{\pgfqpoint{5.901579in}{0.655279in}}%
\pgfpathlineto{\pgfqpoint{5.905210in}{0.673921in}}%
\pgfpathlineto{\pgfqpoint{5.908841in}{0.661397in}}%
\pgfpathlineto{\pgfqpoint{5.912472in}{0.664991in}}%
\pgfpathlineto{\pgfqpoint{5.916103in}{0.591872in}}%
\pgfpathlineto{\pgfqpoint{5.923366in}{0.663288in}}%
\pgfpathlineto{\pgfqpoint{5.926997in}{0.651488in}}%
\pgfpathlineto{\pgfqpoint{5.930628in}{0.569179in}}%
\pgfpathlineto{\pgfqpoint{5.934259in}{0.671551in}}%
\pgfpathlineto{\pgfqpoint{5.937890in}{0.593178in}}%
\pgfpathlineto{\pgfqpoint{5.941522in}{0.645509in}}%
\pgfpathlineto{\pgfqpoint{5.945153in}{0.643401in}}%
\pgfpathlineto{\pgfqpoint{5.948784in}{0.562496in}}%
\pgfpathlineto{\pgfqpoint{5.952415in}{0.630131in}}%
\pgfpathlineto{\pgfqpoint{5.956046in}{0.650878in}}%
\pgfpathlineto{\pgfqpoint{5.959677in}{0.643473in}}%
\pgfpathlineto{\pgfqpoint{5.963309in}{0.641595in}}%
\pgfpathlineto{\pgfqpoint{5.966940in}{0.564308in}}%
\pgfpathlineto{\pgfqpoint{5.970571in}{0.664612in}}%
\pgfpathlineto{\pgfqpoint{5.974202in}{0.682642in}}%
\pgfpathlineto{\pgfqpoint{5.977833in}{0.675211in}}%
\pgfpathlineto{\pgfqpoint{5.981465in}{0.672247in}}%
\pgfpathlineto{\pgfqpoint{5.985096in}{0.593920in}}%
\pgfpathlineto{\pgfqpoint{5.988727in}{0.614066in}}%
\pgfpathlineto{\pgfqpoint{5.992358in}{0.650536in}}%
\pgfpathlineto{\pgfqpoint{5.995989in}{0.600167in}}%
\pgfpathlineto{\pgfqpoint{5.999620in}{0.667487in}}%
\pgfpathlineto{\pgfqpoint{6.003252in}{0.698630in}}%
\pgfpathlineto{\pgfqpoint{6.006883in}{0.647849in}}%
\pgfpathlineto{\pgfqpoint{6.010514in}{0.531505in}}%
\pgfpathlineto{\pgfqpoint{6.014145in}{0.653720in}}%
\pgfpathlineto{\pgfqpoint{6.017776in}{0.634569in}}%
\pgfpathlineto{\pgfqpoint{6.021407in}{0.644560in}}%
\pgfpathlineto{\pgfqpoint{6.025039in}{0.556601in}}%
\pgfpathlineto{\pgfqpoint{6.028670in}{0.649080in}}%
\pgfpathlineto{\pgfqpoint{6.035932in}{0.620898in}}%
\pgfpathlineto{\pgfqpoint{6.039563in}{0.650911in}}%
\pgfpathlineto{\pgfqpoint{6.043194in}{0.613890in}}%
\pgfpathlineto{\pgfqpoint{6.046826in}{0.651588in}}%
\pgfpathlineto{\pgfqpoint{6.050457in}{0.683372in}}%
\pgfpathlineto{\pgfqpoint{6.057719in}{0.566986in}}%
\pgfpathlineto{\pgfqpoint{6.061350in}{0.642890in}}%
\pgfpathlineto{\pgfqpoint{6.064981in}{0.571617in}}%
\pgfpathlineto{\pgfqpoint{6.068613in}{0.672366in}}%
\pgfpathlineto{\pgfqpoint{6.072244in}{0.691141in}}%
\pgfpathlineto{\pgfqpoint{6.075875in}{0.573736in}}%
\pgfpathlineto{\pgfqpoint{6.079506in}{0.618117in}}%
\pgfpathlineto{\pgfqpoint{6.083137in}{0.594485in}}%
\pgfpathlineto{\pgfqpoint{6.086769in}{0.643111in}}%
\pgfpathlineto{\pgfqpoint{6.090400in}{0.676258in}}%
\pgfpathlineto{\pgfqpoint{6.094031in}{0.675599in}}%
\pgfpathlineto{\pgfqpoint{6.097662in}{0.652534in}}%
\pgfpathlineto{\pgfqpoint{6.101293in}{0.614324in}}%
\pgfpathlineto{\pgfqpoint{6.104924in}{0.606613in}}%
\pgfpathlineto{\pgfqpoint{6.108556in}{0.662492in}}%
\pgfpathlineto{\pgfqpoint{6.112187in}{0.664964in}}%
\pgfpathlineto{\pgfqpoint{6.115818in}{0.630590in}}%
\pgfpathlineto{\pgfqpoint{6.119449in}{0.671783in}}%
\pgfpathlineto{\pgfqpoint{6.123080in}{0.504670in}}%
\pgfpathlineto{\pgfqpoint{6.126711in}{0.640781in}}%
\pgfpathlineto{\pgfqpoint{6.130343in}{0.658722in}}%
\pgfpathlineto{\pgfqpoint{6.133974in}{0.616767in}}%
\pgfpathlineto{\pgfqpoint{6.137605in}{0.606236in}}%
\pgfpathlineto{\pgfqpoint{6.141236in}{0.663724in}}%
\pgfpathlineto{\pgfqpoint{6.144867in}{0.626692in}}%
\pgfpathlineto{\pgfqpoint{6.148498in}{0.646403in}}%
\pgfpathlineto{\pgfqpoint{6.152130in}{0.611084in}}%
\pgfpathlineto{\pgfqpoint{6.155761in}{0.641379in}}%
\pgfpathlineto{\pgfqpoint{6.159392in}{0.680622in}}%
\pgfpathlineto{\pgfqpoint{6.163023in}{0.687696in}}%
\pgfpathlineto{\pgfqpoint{6.166654in}{0.656851in}}%
\pgfpathlineto{\pgfqpoint{6.170285in}{0.614682in}}%
\pgfpathlineto{\pgfqpoint{6.173917in}{0.582605in}}%
\pgfpathlineto{\pgfqpoint{6.177548in}{0.656104in}}%
\pgfpathlineto{\pgfqpoint{6.181179in}{0.635201in}}%
\pgfpathlineto{\pgfqpoint{6.184810in}{0.647077in}}%
\pgfpathlineto{\pgfqpoint{6.188441in}{0.692036in}}%
\pgfpathlineto{\pgfqpoint{6.192073in}{0.663225in}}%
\pgfpathlineto{\pgfqpoint{6.195704in}{0.602505in}}%
\pgfpathlineto{\pgfqpoint{6.199335in}{0.664414in}}%
\pgfpathlineto{\pgfqpoint{6.202966in}{0.614794in}}%
\pgfpathlineto{\pgfqpoint{6.206597in}{0.647438in}}%
\pgfpathlineto{\pgfqpoint{6.210228in}{0.631264in}}%
\pgfpathlineto{\pgfqpoint{6.213860in}{0.600709in}}%
\pgfpathlineto{\pgfqpoint{6.217491in}{0.626702in}}%
\pgfpathlineto{\pgfqpoint{6.221122in}{0.643253in}}%
\pgfpathlineto{\pgfqpoint{6.224753in}{0.646342in}}%
\pgfpathlineto{\pgfqpoint{6.228384in}{0.631811in}}%
\pgfpathlineto{\pgfqpoint{6.232015in}{0.624810in}}%
\pgfpathlineto{\pgfqpoint{6.235647in}{0.672261in}}%
\pgfpathlineto{\pgfqpoint{6.239278in}{0.634904in}}%
\pgfpathlineto{\pgfqpoint{6.242909in}{0.560832in}}%
\pgfpathlineto{\pgfqpoint{6.246540in}{0.652074in}}%
\pgfpathlineto{\pgfqpoint{6.250171in}{0.576114in}}%
\pgfpathlineto{\pgfqpoint{6.253802in}{0.663795in}}%
\pgfpathlineto{\pgfqpoint{6.257434in}{0.689420in}}%
\pgfpathlineto{\pgfqpoint{6.261065in}{0.585067in}}%
\pgfpathlineto{\pgfqpoint{6.268327in}{0.628187in}}%
\pgfpathlineto{\pgfqpoint{6.271958in}{0.622450in}}%
\pgfpathlineto{\pgfqpoint{6.275590in}{0.675385in}}%
\pgfpathlineto{\pgfqpoint{6.279221in}{0.684932in}}%
\pgfpathlineto{\pgfqpoint{6.282852in}{0.631540in}}%
\pgfpathlineto{\pgfqpoint{6.286483in}{0.622216in}}%
\pgfpathlineto{\pgfqpoint{6.290114in}{0.526257in}}%
\pgfpathlineto{\pgfqpoint{6.293745in}{0.656514in}}%
\pgfpathlineto{\pgfqpoint{6.297377in}{0.671804in}}%
\pgfpathlineto{\pgfqpoint{6.301008in}{0.655709in}}%
\pgfpathlineto{\pgfqpoint{6.304639in}{0.666282in}}%
\pgfpathlineto{\pgfqpoint{6.308270in}{0.595106in}}%
\pgfpathlineto{\pgfqpoint{6.315532in}{0.664476in}}%
\pgfpathlineto{\pgfqpoint{6.319164in}{0.647387in}}%
\pgfpathlineto{\pgfqpoint{6.322795in}{0.635583in}}%
\pgfpathlineto{\pgfqpoint{6.326426in}{0.696326in}}%
\pgfpathlineto{\pgfqpoint{6.330057in}{0.619522in}}%
\pgfpathlineto{\pgfqpoint{6.333688in}{0.589514in}}%
\pgfpathlineto{\pgfqpoint{6.337319in}{0.636277in}}%
\pgfpathlineto{\pgfqpoint{6.340951in}{0.632030in}}%
\pgfpathlineto{\pgfqpoint{6.344582in}{0.662607in}}%
\pgfpathlineto{\pgfqpoint{6.348213in}{0.652621in}}%
\pgfpathlineto{\pgfqpoint{6.351844in}{0.613990in}}%
\pgfpathlineto{\pgfqpoint{6.355475in}{0.667298in}}%
\pgfpathlineto{\pgfqpoint{6.359106in}{0.681228in}}%
\pgfpathlineto{\pgfqpoint{6.362738in}{0.658427in}}%
\pgfpathlineto{\pgfqpoint{6.366369in}{0.597697in}}%
\pgfpathlineto{\pgfqpoint{6.370000in}{0.607857in}}%
\pgfpathlineto{\pgfqpoint{6.370000in}{0.607857in}}%
\pgfusepath{stroke}%
\end{pgfscope}%
\begin{pgfscope}%
\pgfpathrectangle{\pgfqpoint{4.554414in}{0.320625in}}{\pgfqpoint{1.815586in}{1.164375in}}%
\pgfusepath{clip}%
\pgfsetrectcap%
\pgfsetroundjoin%
\pgfsetlinewidth{0.602250pt}%
\definecolor{currentstroke}{rgb}{0.000000,0.529412,0.423529}%
\pgfsetstrokecolor{currentstroke}%
\pgfsetstrokeopacity{0.120000}%
\pgfsetdash{}{0pt}%
\pgfpathmoveto{\pgfqpoint{4.554414in}{1.457142in}}%
\pgfpathlineto{\pgfqpoint{4.558045in}{1.445482in}}%
\pgfpathlineto{\pgfqpoint{4.565307in}{1.430171in}}%
\pgfpathlineto{\pgfqpoint{4.583463in}{1.393753in}}%
\pgfpathlineto{\pgfqpoint{4.605250in}{1.346836in}}%
\pgfpathlineto{\pgfqpoint{4.656086in}{1.236548in}}%
\pgfpathlineto{\pgfqpoint{4.692398in}{1.160821in}}%
\pgfpathlineto{\pgfqpoint{4.732341in}{1.080366in}}%
\pgfpathlineto{\pgfqpoint{4.775915in}{0.995289in}}%
\pgfpathlineto{\pgfqpoint{4.823120in}{0.905569in}}%
\pgfpathlineto{\pgfqpoint{4.877588in}{0.804448in}}%
\pgfpathlineto{\pgfqpoint{4.917531in}{0.731649in}}%
\pgfpathlineto{\pgfqpoint{4.921162in}{0.726757in}}%
\pgfpathlineto{\pgfqpoint{4.924793in}{0.722808in}}%
\pgfpathlineto{\pgfqpoint{4.928424in}{0.712418in}}%
\pgfpathlineto{\pgfqpoint{4.932056in}{0.705137in}}%
\pgfpathlineto{\pgfqpoint{4.935687in}{0.707460in}}%
\pgfpathlineto{\pgfqpoint{4.939318in}{0.692111in}}%
\pgfpathlineto{\pgfqpoint{4.946580in}{0.682288in}}%
\pgfpathlineto{\pgfqpoint{4.950211in}{0.672494in}}%
\pgfpathlineto{\pgfqpoint{4.953843in}{0.680236in}}%
\pgfpathlineto{\pgfqpoint{4.957474in}{0.699374in}}%
\pgfpathlineto{\pgfqpoint{4.961105in}{0.676670in}}%
\pgfpathlineto{\pgfqpoint{4.964736in}{0.658240in}}%
\pgfpathlineto{\pgfqpoint{4.968367in}{0.675601in}}%
\pgfpathlineto{\pgfqpoint{4.971998in}{0.633280in}}%
\pgfpathlineto{\pgfqpoint{4.975630in}{0.647044in}}%
\pgfpathlineto{\pgfqpoint{4.979261in}{0.668615in}}%
\pgfpathlineto{\pgfqpoint{4.982892in}{0.673713in}}%
\pgfpathlineto{\pgfqpoint{4.986523in}{0.661709in}}%
\pgfpathlineto{\pgfqpoint{4.990154in}{0.667673in}}%
\pgfpathlineto{\pgfqpoint{4.993785in}{0.611499in}}%
\pgfpathlineto{\pgfqpoint{4.997417in}{0.617473in}}%
\pgfpathlineto{\pgfqpoint{5.001048in}{0.660388in}}%
\pgfpathlineto{\pgfqpoint{5.004679in}{0.677126in}}%
\pgfpathlineto{\pgfqpoint{5.008310in}{0.674893in}}%
\pgfpathlineto{\pgfqpoint{5.011941in}{0.651884in}}%
\pgfpathlineto{\pgfqpoint{5.015573in}{0.619835in}}%
\pgfpathlineto{\pgfqpoint{5.019204in}{0.549895in}}%
\pgfpathlineto{\pgfqpoint{5.022835in}{0.649585in}}%
\pgfpathlineto{\pgfqpoint{5.026466in}{0.679169in}}%
\pgfpathlineto{\pgfqpoint{5.030097in}{0.686785in}}%
\pgfpathlineto{\pgfqpoint{5.033728in}{0.621768in}}%
\pgfpathlineto{\pgfqpoint{5.037360in}{0.621922in}}%
\pgfpathlineto{\pgfqpoint{5.040991in}{0.612137in}}%
\pgfpathlineto{\pgfqpoint{5.044622in}{0.627534in}}%
\pgfpathlineto{\pgfqpoint{5.048253in}{0.676313in}}%
\pgfpathlineto{\pgfqpoint{5.051884in}{0.692020in}}%
\pgfpathlineto{\pgfqpoint{5.055515in}{0.484903in}}%
\pgfpathlineto{\pgfqpoint{5.059147in}{0.612800in}}%
\pgfpathlineto{\pgfqpoint{5.062778in}{0.641747in}}%
\pgfpathlineto{\pgfqpoint{5.066409in}{0.568870in}}%
\pgfpathlineto{\pgfqpoint{5.070040in}{0.668414in}}%
\pgfpathlineto{\pgfqpoint{5.073671in}{0.690364in}}%
\pgfpathlineto{\pgfqpoint{5.077302in}{0.614958in}}%
\pgfpathlineto{\pgfqpoint{5.080934in}{0.578027in}}%
\pgfpathlineto{\pgfqpoint{5.084565in}{0.655842in}}%
\pgfpathlineto{\pgfqpoint{5.088196in}{0.595429in}}%
\pgfpathlineto{\pgfqpoint{5.091827in}{0.652957in}}%
\pgfpathlineto{\pgfqpoint{5.095458in}{0.683108in}}%
\pgfpathlineto{\pgfqpoint{5.099090in}{0.636988in}}%
\pgfpathlineto{\pgfqpoint{5.102721in}{0.573989in}}%
\pgfpathlineto{\pgfqpoint{5.106352in}{0.658912in}}%
\pgfpathlineto{\pgfqpoint{5.109983in}{0.630175in}}%
\pgfpathlineto{\pgfqpoint{5.113614in}{0.619764in}}%
\pgfpathlineto{\pgfqpoint{5.117245in}{0.669559in}}%
\pgfpathlineto{\pgfqpoint{5.120877in}{0.683695in}}%
\pgfpathlineto{\pgfqpoint{5.124508in}{0.688004in}}%
\pgfpathlineto{\pgfqpoint{5.128139in}{0.649074in}}%
\pgfpathlineto{\pgfqpoint{5.135401in}{0.592561in}}%
\pgfpathlineto{\pgfqpoint{5.139032in}{0.648397in}}%
\pgfpathlineto{\pgfqpoint{5.142664in}{0.663531in}}%
\pgfpathlineto{\pgfqpoint{5.146295in}{0.656171in}}%
\pgfpathlineto{\pgfqpoint{5.149926in}{0.472999in}}%
\pgfpathlineto{\pgfqpoint{5.153557in}{0.654651in}}%
\pgfpathlineto{\pgfqpoint{5.157188in}{0.696222in}}%
\pgfpathlineto{\pgfqpoint{5.160819in}{0.609311in}}%
\pgfpathlineto{\pgfqpoint{5.164451in}{0.619730in}}%
\pgfpathlineto{\pgfqpoint{5.168082in}{0.617038in}}%
\pgfpathlineto{\pgfqpoint{5.171713in}{0.615752in}}%
\pgfpathlineto{\pgfqpoint{5.175344in}{0.677263in}}%
\pgfpathlineto{\pgfqpoint{5.178975in}{0.695149in}}%
\pgfpathlineto{\pgfqpoint{5.182606in}{0.572561in}}%
\pgfpathlineto{\pgfqpoint{5.186238in}{0.596062in}}%
\pgfpathlineto{\pgfqpoint{5.189869in}{0.640996in}}%
\pgfpathlineto{\pgfqpoint{5.193500in}{0.534035in}}%
\pgfpathlineto{\pgfqpoint{5.197131in}{0.668215in}}%
\pgfpathlineto{\pgfqpoint{5.200762in}{0.689624in}}%
\pgfpathlineto{\pgfqpoint{5.204394in}{0.623389in}}%
\pgfpathlineto{\pgfqpoint{5.208025in}{0.511519in}}%
\pgfpathlineto{\pgfqpoint{5.211656in}{0.654499in}}%
\pgfpathlineto{\pgfqpoint{5.215287in}{0.606600in}}%
\pgfpathlineto{\pgfqpoint{5.218918in}{0.649297in}}%
\pgfpathlineto{\pgfqpoint{5.222549in}{0.679586in}}%
\pgfpathlineto{\pgfqpoint{5.229812in}{0.604078in}}%
\pgfpathlineto{\pgfqpoint{5.233443in}{0.655486in}}%
\pgfpathlineto{\pgfqpoint{5.237074in}{0.633563in}}%
\pgfpathlineto{\pgfqpoint{5.240705in}{0.606347in}}%
\pgfpathlineto{\pgfqpoint{5.244336in}{0.661794in}}%
\pgfpathlineto{\pgfqpoint{5.247968in}{0.641634in}}%
\pgfpathlineto{\pgfqpoint{5.251599in}{0.639183in}}%
\pgfpathlineto{\pgfqpoint{5.255230in}{0.647691in}}%
\pgfpathlineto{\pgfqpoint{5.258861in}{0.645148in}}%
\pgfpathlineto{\pgfqpoint{5.262492in}{0.599957in}}%
\pgfpathlineto{\pgfqpoint{5.266123in}{0.626723in}}%
\pgfpathlineto{\pgfqpoint{5.269755in}{0.634498in}}%
\pgfpathlineto{\pgfqpoint{5.273386in}{0.658727in}}%
\pgfpathlineto{\pgfqpoint{5.277017in}{0.630174in}}%
\pgfpathlineto{\pgfqpoint{5.280648in}{0.648773in}}%
\pgfpathlineto{\pgfqpoint{5.284279in}{0.645787in}}%
\pgfpathlineto{\pgfqpoint{5.287910in}{0.568340in}}%
\pgfpathlineto{\pgfqpoint{5.291542in}{0.663305in}}%
\pgfpathlineto{\pgfqpoint{5.295173in}{0.700332in}}%
\pgfpathlineto{\pgfqpoint{5.302435in}{0.543526in}}%
\pgfpathlineto{\pgfqpoint{5.306066in}{0.654817in}}%
\pgfpathlineto{\pgfqpoint{5.309698in}{0.570016in}}%
\pgfpathlineto{\pgfqpoint{5.313329in}{0.658533in}}%
\pgfpathlineto{\pgfqpoint{5.316960in}{0.679754in}}%
\pgfpathlineto{\pgfqpoint{5.320591in}{0.676505in}}%
\pgfpathlineto{\pgfqpoint{5.324222in}{0.646108in}}%
\pgfpathlineto{\pgfqpoint{5.327853in}{0.683696in}}%
\pgfpathlineto{\pgfqpoint{5.331485in}{0.619014in}}%
\pgfpathlineto{\pgfqpoint{5.335116in}{0.489083in}}%
\pgfpathlineto{\pgfqpoint{5.338747in}{0.653569in}}%
\pgfpathlineto{\pgfqpoint{5.342378in}{0.620383in}}%
\pgfpathlineto{\pgfqpoint{5.346009in}{0.640893in}}%
\pgfpathlineto{\pgfqpoint{5.349640in}{0.674493in}}%
\pgfpathlineto{\pgfqpoint{5.353272in}{0.637313in}}%
\pgfpathlineto{\pgfqpoint{5.356903in}{0.622075in}}%
\pgfpathlineto{\pgfqpoint{5.360534in}{0.651303in}}%
\pgfpathlineto{\pgfqpoint{5.364165in}{0.634247in}}%
\pgfpathlineto{\pgfqpoint{5.367796in}{0.594632in}}%
\pgfpathlineto{\pgfqpoint{5.371427in}{0.651966in}}%
\pgfpathlineto{\pgfqpoint{5.375059in}{0.636087in}}%
\pgfpathlineto{\pgfqpoint{5.378690in}{0.647128in}}%
\pgfpathlineto{\pgfqpoint{5.382321in}{0.641677in}}%
\pgfpathlineto{\pgfqpoint{5.385952in}{0.645399in}}%
\pgfpathlineto{\pgfqpoint{5.393215in}{0.590605in}}%
\pgfpathlineto{\pgfqpoint{5.396846in}{0.620406in}}%
\pgfpathlineto{\pgfqpoint{5.400477in}{0.663701in}}%
\pgfpathlineto{\pgfqpoint{5.404108in}{0.620226in}}%
\pgfpathlineto{\pgfqpoint{5.407739in}{0.647294in}}%
\pgfpathlineto{\pgfqpoint{5.411370in}{0.650823in}}%
\pgfpathlineto{\pgfqpoint{5.415002in}{0.618530in}}%
\pgfpathlineto{\pgfqpoint{5.418633in}{0.558018in}}%
\pgfpathlineto{\pgfqpoint{5.422264in}{0.673311in}}%
\pgfpathlineto{\pgfqpoint{5.425895in}{0.564896in}}%
\pgfpathlineto{\pgfqpoint{5.429526in}{0.642776in}}%
\pgfpathlineto{\pgfqpoint{5.433157in}{0.664156in}}%
\pgfpathlineto{\pgfqpoint{5.436789in}{0.653097in}}%
\pgfpathlineto{\pgfqpoint{5.440420in}{0.611742in}}%
\pgfpathlineto{\pgfqpoint{5.444051in}{0.673626in}}%
\pgfpathlineto{\pgfqpoint{5.447682in}{0.581889in}}%
\pgfpathlineto{\pgfqpoint{5.451313in}{0.629332in}}%
\pgfpathlineto{\pgfqpoint{5.454944in}{0.666282in}}%
\pgfpathlineto{\pgfqpoint{5.458576in}{0.667529in}}%
\pgfpathlineto{\pgfqpoint{5.462207in}{0.649971in}}%
\pgfpathlineto{\pgfqpoint{5.465838in}{0.666156in}}%
\pgfpathlineto{\pgfqpoint{5.469469in}{0.616657in}}%
\pgfpathlineto{\pgfqpoint{5.473100in}{0.594952in}}%
\pgfpathlineto{\pgfqpoint{5.476731in}{0.660863in}}%
\pgfpathlineto{\pgfqpoint{5.480363in}{0.674052in}}%
\pgfpathlineto{\pgfqpoint{5.483994in}{0.666755in}}%
\pgfpathlineto{\pgfqpoint{5.487625in}{0.626092in}}%
\pgfpathlineto{\pgfqpoint{5.494887in}{0.669916in}}%
\pgfpathlineto{\pgfqpoint{5.498519in}{0.643001in}}%
\pgfpathlineto{\pgfqpoint{5.502150in}{0.633580in}}%
\pgfpathlineto{\pgfqpoint{5.505781in}{0.576630in}}%
\pgfpathlineto{\pgfqpoint{5.513043in}{0.682961in}}%
\pgfpathlineto{\pgfqpoint{5.516674in}{0.701766in}}%
\pgfpathlineto{\pgfqpoint{5.520306in}{0.634736in}}%
\pgfpathlineto{\pgfqpoint{5.523937in}{0.532826in}}%
\pgfpathlineto{\pgfqpoint{5.527568in}{0.663338in}}%
\pgfpathlineto{\pgfqpoint{5.531199in}{0.611335in}}%
\pgfpathlineto{\pgfqpoint{5.534830in}{0.648789in}}%
\pgfpathlineto{\pgfqpoint{5.538461in}{0.658629in}}%
\pgfpathlineto{\pgfqpoint{5.542093in}{0.643517in}}%
\pgfpathlineto{\pgfqpoint{5.545724in}{0.592236in}}%
\pgfpathlineto{\pgfqpoint{5.549355in}{0.674188in}}%
\pgfpathlineto{\pgfqpoint{5.552986in}{0.539107in}}%
\pgfpathlineto{\pgfqpoint{5.556617in}{0.639112in}}%
\pgfpathlineto{\pgfqpoint{5.560248in}{0.662676in}}%
\pgfpathlineto{\pgfqpoint{5.563880in}{0.661578in}}%
\pgfpathlineto{\pgfqpoint{5.567511in}{0.637999in}}%
\pgfpathlineto{\pgfqpoint{5.571142in}{0.671112in}}%
\pgfpathlineto{\pgfqpoint{5.574773in}{0.589535in}}%
\pgfpathlineto{\pgfqpoint{5.578404in}{0.623801in}}%
\pgfpathlineto{\pgfqpoint{5.582035in}{0.663822in}}%
\pgfpathlineto{\pgfqpoint{5.585667in}{0.673119in}}%
\pgfpathlineto{\pgfqpoint{5.589298in}{0.664189in}}%
\pgfpathlineto{\pgfqpoint{5.592929in}{0.660630in}}%
\pgfpathlineto{\pgfqpoint{5.600191in}{0.578021in}}%
\pgfpathlineto{\pgfqpoint{5.603823in}{0.656274in}}%
\pgfpathlineto{\pgfqpoint{5.607454in}{0.677322in}}%
\pgfpathlineto{\pgfqpoint{5.611085in}{0.679290in}}%
\pgfpathlineto{\pgfqpoint{5.614716in}{0.640017in}}%
\pgfpathlineto{\pgfqpoint{5.618347in}{0.625070in}}%
\pgfpathlineto{\pgfqpoint{5.621978in}{0.589960in}}%
\pgfpathlineto{\pgfqpoint{5.625610in}{0.640257in}}%
\pgfpathlineto{\pgfqpoint{5.629241in}{0.676847in}}%
\pgfpathlineto{\pgfqpoint{5.632872in}{0.688683in}}%
\pgfpathlineto{\pgfqpoint{5.636503in}{0.593147in}}%
\pgfpathlineto{\pgfqpoint{5.640134in}{0.622377in}}%
\pgfpathlineto{\pgfqpoint{5.643765in}{0.633455in}}%
\pgfpathlineto{\pgfqpoint{5.647397in}{0.606749in}}%
\pgfpathlineto{\pgfqpoint{5.651028in}{0.671678in}}%
\pgfpathlineto{\pgfqpoint{5.654659in}{0.691002in}}%
\pgfpathlineto{\pgfqpoint{5.658290in}{0.667202in}}%
\pgfpathlineto{\pgfqpoint{5.661921in}{0.613922in}}%
\pgfpathlineto{\pgfqpoint{5.665552in}{0.685107in}}%
\pgfpathlineto{\pgfqpoint{5.669184in}{0.525275in}}%
\pgfpathlineto{\pgfqpoint{5.672815in}{0.622268in}}%
\pgfpathlineto{\pgfqpoint{5.676446in}{0.654142in}}%
\pgfpathlineto{\pgfqpoint{5.680077in}{0.630003in}}%
\pgfpathlineto{\pgfqpoint{5.683708in}{0.642252in}}%
\pgfpathlineto{\pgfqpoint{5.687340in}{0.685759in}}%
\pgfpathlineto{\pgfqpoint{5.690971in}{0.674187in}}%
\pgfpathlineto{\pgfqpoint{5.694602in}{0.665335in}}%
\pgfpathlineto{\pgfqpoint{5.698233in}{0.666766in}}%
\pgfpathlineto{\pgfqpoint{5.701864in}{0.596918in}}%
\pgfpathlineto{\pgfqpoint{5.705495in}{0.620278in}}%
\pgfpathlineto{\pgfqpoint{5.709127in}{0.662635in}}%
\pgfpathlineto{\pgfqpoint{5.712758in}{0.679332in}}%
\pgfpathlineto{\pgfqpoint{5.716389in}{0.677342in}}%
\pgfpathlineto{\pgfqpoint{5.720020in}{0.653918in}}%
\pgfpathlineto{\pgfqpoint{5.723651in}{0.613245in}}%
\pgfpathlineto{\pgfqpoint{5.727282in}{0.558092in}}%
\pgfpathlineto{\pgfqpoint{5.730914in}{0.651151in}}%
\pgfpathlineto{\pgfqpoint{5.734545in}{0.679917in}}%
\pgfpathlineto{\pgfqpoint{5.738176in}{0.686468in}}%
\pgfpathlineto{\pgfqpoint{5.741807in}{0.626876in}}%
\pgfpathlineto{\pgfqpoint{5.745438in}{0.619998in}}%
\pgfpathlineto{\pgfqpoint{5.749069in}{0.600936in}}%
\pgfpathlineto{\pgfqpoint{5.752701in}{0.631621in}}%
\pgfpathlineto{\pgfqpoint{5.756332in}{0.677949in}}%
\pgfpathlineto{\pgfqpoint{5.759963in}{0.692744in}}%
\pgfpathlineto{\pgfqpoint{5.763594in}{0.527167in}}%
\pgfpathlineto{\pgfqpoint{5.767225in}{0.612309in}}%
\pgfpathlineto{\pgfqpoint{5.770856in}{0.636789in}}%
\pgfpathlineto{\pgfqpoint{5.774488in}{0.584818in}}%
\pgfpathlineto{\pgfqpoint{5.778119in}{0.670818in}}%
\pgfpathlineto{\pgfqpoint{5.781750in}{0.691379in}}%
\pgfpathlineto{\pgfqpoint{5.785381in}{0.610050in}}%
\pgfpathlineto{\pgfqpoint{5.789012in}{0.581512in}}%
\pgfpathlineto{\pgfqpoint{5.792644in}{0.653152in}}%
\pgfpathlineto{\pgfqpoint{5.796275in}{0.584581in}}%
\pgfpathlineto{\pgfqpoint{5.799906in}{0.656681in}}%
\pgfpathlineto{\pgfqpoint{5.803537in}{0.684460in}}%
\pgfpathlineto{\pgfqpoint{5.807168in}{0.634946in}}%
\pgfpathlineto{\pgfqpoint{5.810799in}{0.565307in}}%
\pgfpathlineto{\pgfqpoint{5.814431in}{0.658093in}}%
\pgfpathlineto{\pgfqpoint{5.818062in}{0.626573in}}%
\pgfpathlineto{\pgfqpoint{5.821693in}{0.627667in}}%
\pgfpathlineto{\pgfqpoint{5.825324in}{0.671347in}}%
\pgfpathlineto{\pgfqpoint{5.828955in}{0.643099in}}%
\pgfpathlineto{\pgfqpoint{5.832586in}{0.624905in}}%
\pgfpathlineto{\pgfqpoint{5.836218in}{0.653555in}}%
\pgfpathlineto{\pgfqpoint{5.839849in}{0.642443in}}%
\pgfpathlineto{\pgfqpoint{5.843480in}{0.525088in}}%
\pgfpathlineto{\pgfqpoint{5.847111in}{0.648034in}}%
\pgfpathlineto{\pgfqpoint{5.850742in}{0.680951in}}%
\pgfpathlineto{\pgfqpoint{5.854373in}{0.698101in}}%
\pgfpathlineto{\pgfqpoint{5.858005in}{0.605848in}}%
\pgfpathlineto{\pgfqpoint{5.865267in}{0.628549in}}%
\pgfpathlineto{\pgfqpoint{5.868898in}{0.633709in}}%
\pgfpathlineto{\pgfqpoint{5.872529in}{0.664808in}}%
\pgfpathlineto{\pgfqpoint{5.876160in}{0.672371in}}%
\pgfpathlineto{\pgfqpoint{5.879792in}{0.641584in}}%
\pgfpathlineto{\pgfqpoint{5.883423in}{0.618972in}}%
\pgfpathlineto{\pgfqpoint{5.887054in}{0.696625in}}%
\pgfpathlineto{\pgfqpoint{5.890685in}{0.489226in}}%
\pgfpathlineto{\pgfqpoint{5.894316in}{0.608005in}}%
\pgfpathlineto{\pgfqpoint{5.897948in}{0.639066in}}%
\pgfpathlineto{\pgfqpoint{5.901579in}{0.531942in}}%
\pgfpathlineto{\pgfqpoint{5.905210in}{0.669611in}}%
\pgfpathlineto{\pgfqpoint{5.908841in}{0.691635in}}%
\pgfpathlineto{\pgfqpoint{5.912472in}{0.620007in}}%
\pgfpathlineto{\pgfqpoint{5.916103in}{0.512757in}}%
\pgfpathlineto{\pgfqpoint{5.919735in}{0.651323in}}%
\pgfpathlineto{\pgfqpoint{5.923366in}{0.594455in}}%
\pgfpathlineto{\pgfqpoint{5.926997in}{0.655050in}}%
\pgfpathlineto{\pgfqpoint{5.930628in}{0.681411in}}%
\pgfpathlineto{\pgfqpoint{5.937890in}{0.598159in}}%
\pgfpathlineto{\pgfqpoint{5.941522in}{0.654931in}}%
\pgfpathlineto{\pgfqpoint{5.945153in}{0.630473in}}%
\pgfpathlineto{\pgfqpoint{5.948784in}{0.617630in}}%
\pgfpathlineto{\pgfqpoint{5.952415in}{0.664171in}}%
\pgfpathlineto{\pgfqpoint{5.956046in}{0.640535in}}%
\pgfpathlineto{\pgfqpoint{5.959677in}{0.636595in}}%
\pgfpathlineto{\pgfqpoint{5.963309in}{0.648400in}}%
\pgfpathlineto{\pgfqpoint{5.966940in}{0.643431in}}%
\pgfpathlineto{\pgfqpoint{5.970571in}{0.581804in}}%
\pgfpathlineto{\pgfqpoint{5.974202in}{0.631220in}}%
\pgfpathlineto{\pgfqpoint{5.977833in}{0.633916in}}%
\pgfpathlineto{\pgfqpoint{5.981465in}{0.657035in}}%
\pgfpathlineto{\pgfqpoint{5.985096in}{0.632481in}}%
\pgfpathlineto{\pgfqpoint{5.988727in}{0.648177in}}%
\pgfpathlineto{\pgfqpoint{5.992358in}{0.641636in}}%
\pgfpathlineto{\pgfqpoint{5.995989in}{0.552234in}}%
\pgfpathlineto{\pgfqpoint{6.003252in}{0.669754in}}%
\pgfpathlineto{\pgfqpoint{6.006883in}{0.598546in}}%
\pgfpathlineto{\pgfqpoint{6.010514in}{0.646364in}}%
\pgfpathlineto{\pgfqpoint{6.014145in}{0.660890in}}%
\pgfpathlineto{\pgfqpoint{6.017776in}{0.637775in}}%
\pgfpathlineto{\pgfqpoint{6.021407in}{0.634586in}}%
\pgfpathlineto{\pgfqpoint{6.025039in}{0.693790in}}%
\pgfpathlineto{\pgfqpoint{6.032301in}{0.598525in}}%
\pgfpathlineto{\pgfqpoint{6.035932in}{0.661601in}}%
\pgfpathlineto{\pgfqpoint{6.039563in}{0.570326in}}%
\pgfpathlineto{\pgfqpoint{6.043194in}{0.648158in}}%
\pgfpathlineto{\pgfqpoint{6.046826in}{0.676238in}}%
\pgfpathlineto{\pgfqpoint{6.050457in}{0.684584in}}%
\pgfpathlineto{\pgfqpoint{6.057719in}{0.665557in}}%
\pgfpathlineto{\pgfqpoint{6.061350in}{0.626669in}}%
\pgfpathlineto{\pgfqpoint{6.064981in}{0.605911in}}%
\pgfpathlineto{\pgfqpoint{6.068613in}{0.652122in}}%
\pgfpathlineto{\pgfqpoint{6.072244in}{0.637686in}}%
\pgfpathlineto{\pgfqpoint{6.075875in}{0.587650in}}%
\pgfpathlineto{\pgfqpoint{6.079506in}{0.652945in}}%
\pgfpathlineto{\pgfqpoint{6.083137in}{0.633609in}}%
\pgfpathlineto{\pgfqpoint{6.086769in}{0.645801in}}%
\pgfpathlineto{\pgfqpoint{6.090400in}{0.642790in}}%
\pgfpathlineto{\pgfqpoint{6.094031in}{0.642568in}}%
\pgfpathlineto{\pgfqpoint{6.097662in}{0.597164in}}%
\pgfpathlineto{\pgfqpoint{6.101293in}{0.604058in}}%
\pgfpathlineto{\pgfqpoint{6.104924in}{0.620501in}}%
\pgfpathlineto{\pgfqpoint{6.108556in}{0.661839in}}%
\pgfpathlineto{\pgfqpoint{6.112187in}{0.623678in}}%
\pgfpathlineto{\pgfqpoint{6.115818in}{0.647046in}}%
\pgfpathlineto{\pgfqpoint{6.119449in}{0.647339in}}%
\pgfpathlineto{\pgfqpoint{6.123080in}{0.614655in}}%
\pgfpathlineto{\pgfqpoint{6.126711in}{0.556702in}}%
\pgfpathlineto{\pgfqpoint{6.130343in}{0.672551in}}%
\pgfpathlineto{\pgfqpoint{6.133974in}{0.578171in}}%
\pgfpathlineto{\pgfqpoint{6.137605in}{0.643389in}}%
\pgfpathlineto{\pgfqpoint{6.141236in}{0.662451in}}%
\pgfpathlineto{\pgfqpoint{6.144867in}{0.652146in}}%
\pgfpathlineto{\pgfqpoint{6.148498in}{0.611527in}}%
\pgfpathlineto{\pgfqpoint{6.152130in}{0.673728in}}%
\pgfpathlineto{\pgfqpoint{6.155761in}{0.570511in}}%
\pgfpathlineto{\pgfqpoint{6.159392in}{0.631564in}}%
\pgfpathlineto{\pgfqpoint{6.163023in}{0.665865in}}%
\pgfpathlineto{\pgfqpoint{6.166654in}{0.667513in}}%
\pgfpathlineto{\pgfqpoint{6.170285in}{0.650016in}}%
\pgfpathlineto{\pgfqpoint{6.173917in}{0.666836in}}%
\pgfpathlineto{\pgfqpoint{6.177548in}{0.612778in}}%
\pgfpathlineto{\pgfqpoint{6.181179in}{0.601943in}}%
\pgfpathlineto{\pgfqpoint{6.184810in}{0.661249in}}%
\pgfpathlineto{\pgfqpoint{6.188441in}{0.674486in}}%
\pgfpathlineto{\pgfqpoint{6.192073in}{0.670301in}}%
\pgfpathlineto{\pgfqpoint{6.195704in}{0.651787in}}%
\pgfpathlineto{\pgfqpoint{6.199335in}{0.626183in}}%
\pgfpathlineto{\pgfqpoint{6.202966in}{0.547649in}}%
\pgfpathlineto{\pgfqpoint{6.206597in}{0.649198in}}%
\pgfpathlineto{\pgfqpoint{6.210228in}{0.676176in}}%
\pgfpathlineto{\pgfqpoint{6.213860in}{0.682597in}}%
\pgfpathlineto{\pgfqpoint{6.217491in}{0.600693in}}%
\pgfpathlineto{\pgfqpoint{6.221122in}{0.626744in}}%
\pgfpathlineto{\pgfqpoint{6.224753in}{0.680832in}}%
\pgfpathlineto{\pgfqpoint{6.228384in}{0.620001in}}%
\pgfpathlineto{\pgfqpoint{6.232015in}{0.634665in}}%
\pgfpathlineto{\pgfqpoint{6.235647in}{0.628514in}}%
\pgfpathlineto{\pgfqpoint{6.239278in}{0.568924in}}%
\pgfpathlineto{\pgfqpoint{6.242909in}{0.674608in}}%
\pgfpathlineto{\pgfqpoint{6.246540in}{0.699935in}}%
\pgfpathlineto{\pgfqpoint{6.250171in}{0.660388in}}%
\pgfpathlineto{\pgfqpoint{6.253802in}{0.613581in}}%
\pgfpathlineto{\pgfqpoint{6.257434in}{0.670936in}}%
\pgfpathlineto{\pgfqpoint{6.261065in}{0.553892in}}%
\pgfpathlineto{\pgfqpoint{6.264696in}{0.643568in}}%
\pgfpathlineto{\pgfqpoint{6.268327in}{0.665921in}}%
\pgfpathlineto{\pgfqpoint{6.271958in}{0.664582in}}%
\pgfpathlineto{\pgfqpoint{6.275590in}{0.642772in}}%
\pgfpathlineto{\pgfqpoint{6.279221in}{0.671818in}}%
\pgfpathlineto{\pgfqpoint{6.282852in}{0.577080in}}%
\pgfpathlineto{\pgfqpoint{6.286483in}{0.626446in}}%
\pgfpathlineto{\pgfqpoint{6.290114in}{0.662922in}}%
\pgfpathlineto{\pgfqpoint{6.293745in}{0.672751in}}%
\pgfpathlineto{\pgfqpoint{6.297377in}{0.663577in}}%
\pgfpathlineto{\pgfqpoint{6.301008in}{0.661768in}}%
\pgfpathlineto{\pgfqpoint{6.304639in}{0.612864in}}%
\pgfpathlineto{\pgfqpoint{6.308270in}{0.589928in}}%
\pgfpathlineto{\pgfqpoint{6.311901in}{0.657092in}}%
\pgfpathlineto{\pgfqpoint{6.315532in}{0.678037in}}%
\pgfpathlineto{\pgfqpoint{6.319164in}{0.679648in}}%
\pgfpathlineto{\pgfqpoint{6.322795in}{0.642132in}}%
\pgfpathlineto{\pgfqpoint{6.326426in}{0.622824in}}%
\pgfpathlineto{\pgfqpoint{6.330057in}{0.575737in}}%
\pgfpathlineto{\pgfqpoint{6.333688in}{0.642211in}}%
\pgfpathlineto{\pgfqpoint{6.337319in}{0.677881in}}%
\pgfpathlineto{\pgfqpoint{6.340951in}{0.689075in}}%
\pgfpathlineto{\pgfqpoint{6.344582in}{0.599119in}}%
\pgfpathlineto{\pgfqpoint{6.348213in}{0.621272in}}%
\pgfpathlineto{\pgfqpoint{6.351844in}{0.628889in}}%
\pgfpathlineto{\pgfqpoint{6.355475in}{0.611696in}}%
\pgfpathlineto{\pgfqpoint{6.359106in}{0.673153in}}%
\pgfpathlineto{\pgfqpoint{6.362738in}{0.691640in}}%
\pgfpathlineto{\pgfqpoint{6.366369in}{0.579910in}}%
\pgfpathlineto{\pgfqpoint{6.370000in}{0.628097in}}%
\pgfpathlineto{\pgfqpoint{6.370000in}{0.628097in}}%
\pgfusepath{stroke}%
\end{pgfscope}%
\begin{pgfscope}%
\pgfpathrectangle{\pgfqpoint{4.554414in}{0.320625in}}{\pgfqpoint{1.815586in}{1.164375in}}%
\pgfusepath{clip}%
\pgfsetrectcap%
\pgfsetroundjoin%
\pgfsetlinewidth{0.602250pt}%
\definecolor{currentstroke}{rgb}{0.000000,0.529412,0.423529}%
\pgfsetstrokecolor{currentstroke}%
\pgfsetstrokeopacity{0.120000}%
\pgfsetdash{}{0pt}%
\pgfpathmoveto{\pgfqpoint{4.554414in}{1.456590in}}%
\pgfpathlineto{\pgfqpoint{4.558045in}{1.445008in}}%
\pgfpathlineto{\pgfqpoint{4.565307in}{1.430196in}}%
\pgfpathlineto{\pgfqpoint{4.579832in}{1.401755in}}%
\pgfpathlineto{\pgfqpoint{4.597988in}{1.363540in}}%
\pgfpathlineto{\pgfqpoint{4.670611in}{1.207651in}}%
\pgfpathlineto{\pgfqpoint{4.710554in}{1.125754in}}%
\pgfpathlineto{\pgfqpoint{4.757759in}{1.031804in}}%
\pgfpathlineto{\pgfqpoint{4.812227in}{0.925987in}}%
\pgfpathlineto{\pgfqpoint{4.873957in}{0.808394in}}%
\pgfpathlineto{\pgfqpoint{4.928424in}{0.706885in}}%
\pgfpathlineto{\pgfqpoint{4.932056in}{0.703365in}}%
\pgfpathlineto{\pgfqpoint{4.935687in}{0.711131in}}%
\pgfpathlineto{\pgfqpoint{4.939318in}{0.685900in}}%
\pgfpathlineto{\pgfqpoint{4.942949in}{0.679306in}}%
\pgfpathlineto{\pgfqpoint{4.946580in}{0.678316in}}%
\pgfpathlineto{\pgfqpoint{4.950211in}{0.666362in}}%
\pgfpathlineto{\pgfqpoint{4.953843in}{0.666113in}}%
\pgfpathlineto{\pgfqpoint{4.957474in}{0.673285in}}%
\pgfpathlineto{\pgfqpoint{4.961105in}{0.657118in}}%
\pgfpathlineto{\pgfqpoint{4.964736in}{0.661965in}}%
\pgfpathlineto{\pgfqpoint{4.968367in}{0.642218in}}%
\pgfpathlineto{\pgfqpoint{4.975630in}{0.665058in}}%
\pgfpathlineto{\pgfqpoint{4.979261in}{0.658934in}}%
\pgfpathlineto{\pgfqpoint{4.982892in}{0.644322in}}%
\pgfpathlineto{\pgfqpoint{4.986523in}{0.664529in}}%
\pgfpathlineto{\pgfqpoint{4.993785in}{0.589423in}}%
\pgfpathlineto{\pgfqpoint{4.997417in}{0.638816in}}%
\pgfpathlineto{\pgfqpoint{5.001048in}{0.673551in}}%
\pgfpathlineto{\pgfqpoint{5.004679in}{0.689923in}}%
\pgfpathlineto{\pgfqpoint{5.008310in}{0.611248in}}%
\pgfpathlineto{\pgfqpoint{5.011941in}{0.564050in}}%
\pgfpathlineto{\pgfqpoint{5.015573in}{0.658305in}}%
\pgfpathlineto{\pgfqpoint{5.019204in}{0.633468in}}%
\pgfpathlineto{\pgfqpoint{5.022835in}{0.549384in}}%
\pgfpathlineto{\pgfqpoint{5.026466in}{0.638749in}}%
\pgfpathlineto{\pgfqpoint{5.030097in}{0.673886in}}%
\pgfpathlineto{\pgfqpoint{5.033728in}{0.699618in}}%
\pgfpathlineto{\pgfqpoint{5.040991in}{0.553630in}}%
\pgfpathlineto{\pgfqpoint{5.044622in}{0.661328in}}%
\pgfpathlineto{\pgfqpoint{5.048253in}{0.601447in}}%
\pgfpathlineto{\pgfqpoint{5.051884in}{0.623119in}}%
\pgfpathlineto{\pgfqpoint{5.055515in}{0.658839in}}%
\pgfpathlineto{\pgfqpoint{5.059147in}{0.684485in}}%
\pgfpathlineto{\pgfqpoint{5.062778in}{0.700023in}}%
\pgfpathlineto{\pgfqpoint{5.066409in}{0.642804in}}%
\pgfpathlineto{\pgfqpoint{5.070040in}{0.595428in}}%
\pgfpathlineto{\pgfqpoint{5.073671in}{0.670454in}}%
\pgfpathlineto{\pgfqpoint{5.077302in}{0.555763in}}%
\pgfpathlineto{\pgfqpoint{5.080934in}{0.626280in}}%
\pgfpathlineto{\pgfqpoint{5.084565in}{0.660946in}}%
\pgfpathlineto{\pgfqpoint{5.088196in}{0.679558in}}%
\pgfpathlineto{\pgfqpoint{5.091827in}{0.687657in}}%
\pgfpathlineto{\pgfqpoint{5.095458in}{0.603972in}}%
\pgfpathlineto{\pgfqpoint{5.099090in}{0.608878in}}%
\pgfpathlineto{\pgfqpoint{5.102721in}{0.637399in}}%
\pgfpathlineto{\pgfqpoint{5.106352in}{0.523950in}}%
\pgfpathlineto{\pgfqpoint{5.109983in}{0.656389in}}%
\pgfpathlineto{\pgfqpoint{5.113614in}{0.676339in}}%
\pgfpathlineto{\pgfqpoint{5.117245in}{0.638881in}}%
\pgfpathlineto{\pgfqpoint{5.120877in}{0.639263in}}%
\pgfpathlineto{\pgfqpoint{5.124508in}{0.638559in}}%
\pgfpathlineto{\pgfqpoint{5.128139in}{0.646227in}}%
\pgfpathlineto{\pgfqpoint{5.131770in}{0.649690in}}%
\pgfpathlineto{\pgfqpoint{5.135401in}{0.639159in}}%
\pgfpathlineto{\pgfqpoint{5.139032in}{0.612334in}}%
\pgfpathlineto{\pgfqpoint{5.142664in}{0.670246in}}%
\pgfpathlineto{\pgfqpoint{5.146295in}{0.606062in}}%
\pgfpathlineto{\pgfqpoint{5.149926in}{0.585130in}}%
\pgfpathlineto{\pgfqpoint{5.153557in}{0.650961in}}%
\pgfpathlineto{\pgfqpoint{5.157188in}{0.676644in}}%
\pgfpathlineto{\pgfqpoint{5.160819in}{0.689510in}}%
\pgfpathlineto{\pgfqpoint{5.164451in}{0.540039in}}%
\pgfpathlineto{\pgfqpoint{5.171713in}{0.654250in}}%
\pgfpathlineto{\pgfqpoint{5.175344in}{0.614293in}}%
\pgfpathlineto{\pgfqpoint{5.178975in}{0.625209in}}%
\pgfpathlineto{\pgfqpoint{5.182606in}{0.660219in}}%
\pgfpathlineto{\pgfqpoint{5.186238in}{0.639291in}}%
\pgfpathlineto{\pgfqpoint{5.189869in}{0.656222in}}%
\pgfpathlineto{\pgfqpoint{5.193500in}{0.619010in}}%
\pgfpathlineto{\pgfqpoint{5.197131in}{0.645724in}}%
\pgfpathlineto{\pgfqpoint{5.200762in}{0.663185in}}%
\pgfpathlineto{\pgfqpoint{5.204394in}{0.656147in}}%
\pgfpathlineto{\pgfqpoint{5.208025in}{0.610261in}}%
\pgfpathlineto{\pgfqpoint{5.211656in}{0.661232in}}%
\pgfpathlineto{\pgfqpoint{5.215287in}{0.650114in}}%
\pgfpathlineto{\pgfqpoint{5.218918in}{0.670425in}}%
\pgfpathlineto{\pgfqpoint{5.222549in}{0.629923in}}%
\pgfpathlineto{\pgfqpoint{5.226181in}{0.633025in}}%
\pgfpathlineto{\pgfqpoint{5.229812in}{0.639834in}}%
\pgfpathlineto{\pgfqpoint{5.233443in}{0.608022in}}%
\pgfpathlineto{\pgfqpoint{5.240705in}{0.677847in}}%
\pgfpathlineto{\pgfqpoint{5.244336in}{0.672415in}}%
\pgfpathlineto{\pgfqpoint{5.247968in}{0.684747in}}%
\pgfpathlineto{\pgfqpoint{5.251599in}{0.637042in}}%
\pgfpathlineto{\pgfqpoint{5.255230in}{0.619007in}}%
\pgfpathlineto{\pgfqpoint{5.258861in}{0.618449in}}%
\pgfpathlineto{\pgfqpoint{5.262492in}{0.581232in}}%
\pgfpathlineto{\pgfqpoint{5.266123in}{0.667283in}}%
\pgfpathlineto{\pgfqpoint{5.269755in}{0.685741in}}%
\pgfpathlineto{\pgfqpoint{5.273386in}{0.634203in}}%
\pgfpathlineto{\pgfqpoint{5.277017in}{0.620206in}}%
\pgfpathlineto{\pgfqpoint{5.280648in}{0.646973in}}%
\pgfpathlineto{\pgfqpoint{5.284279in}{0.641392in}}%
\pgfpathlineto{\pgfqpoint{5.287910in}{0.622604in}}%
\pgfpathlineto{\pgfqpoint{5.291542in}{0.599614in}}%
\pgfpathlineto{\pgfqpoint{5.295173in}{0.534458in}}%
\pgfpathlineto{\pgfqpoint{5.298804in}{0.672703in}}%
\pgfpathlineto{\pgfqpoint{5.302435in}{0.572285in}}%
\pgfpathlineto{\pgfqpoint{5.309698in}{0.658580in}}%
\pgfpathlineto{\pgfqpoint{5.313329in}{0.677470in}}%
\pgfpathlineto{\pgfqpoint{5.316960in}{0.685720in}}%
\pgfpathlineto{\pgfqpoint{5.320591in}{0.594153in}}%
\pgfpathlineto{\pgfqpoint{5.324222in}{0.611255in}}%
\pgfpathlineto{\pgfqpoint{5.327853in}{0.643858in}}%
\pgfpathlineto{\pgfqpoint{5.331485in}{0.570546in}}%
\pgfpathlineto{\pgfqpoint{5.335116in}{0.650065in}}%
\pgfpathlineto{\pgfqpoint{5.338747in}{0.673684in}}%
\pgfpathlineto{\pgfqpoint{5.342378in}{0.640435in}}%
\pgfpathlineto{\pgfqpoint{5.346009in}{0.642994in}}%
\pgfpathlineto{\pgfqpoint{5.349640in}{0.636188in}}%
\pgfpathlineto{\pgfqpoint{5.353272in}{0.646927in}}%
\pgfpathlineto{\pgfqpoint{5.356903in}{0.653860in}}%
\pgfpathlineto{\pgfqpoint{5.360534in}{0.641627in}}%
\pgfpathlineto{\pgfqpoint{5.364165in}{0.613259in}}%
\pgfpathlineto{\pgfqpoint{5.367796in}{0.669561in}}%
\pgfpathlineto{\pgfqpoint{5.371427in}{0.612594in}}%
\pgfpathlineto{\pgfqpoint{5.375059in}{0.565341in}}%
\pgfpathlineto{\pgfqpoint{5.378690in}{0.649094in}}%
\pgfpathlineto{\pgfqpoint{5.382321in}{0.675316in}}%
\pgfpathlineto{\pgfqpoint{5.385952in}{0.688951in}}%
\pgfpathlineto{\pgfqpoint{5.389583in}{0.564732in}}%
\pgfpathlineto{\pgfqpoint{5.393215in}{0.594545in}}%
\pgfpathlineto{\pgfqpoint{5.396846in}{0.656847in}}%
\pgfpathlineto{\pgfqpoint{5.400477in}{0.620261in}}%
\pgfpathlineto{\pgfqpoint{5.404108in}{0.614867in}}%
\pgfpathlineto{\pgfqpoint{5.407739in}{0.657433in}}%
\pgfpathlineto{\pgfqpoint{5.415002in}{0.699515in}}%
\pgfpathlineto{\pgfqpoint{5.418633in}{0.593047in}}%
\pgfpathlineto{\pgfqpoint{5.422264in}{0.559878in}}%
\pgfpathlineto{\pgfqpoint{5.425895in}{0.658066in}}%
\pgfpathlineto{\pgfqpoint{5.429526in}{0.566028in}}%
\pgfpathlineto{\pgfqpoint{5.433157in}{0.639500in}}%
\pgfpathlineto{\pgfqpoint{5.436789in}{0.667517in}}%
\pgfpathlineto{\pgfqpoint{5.440420in}{0.685765in}}%
\pgfpathlineto{\pgfqpoint{5.444051in}{0.697318in}}%
\pgfpathlineto{\pgfqpoint{5.447682in}{0.614375in}}%
\pgfpathlineto{\pgfqpoint{5.451313in}{0.463522in}}%
\pgfpathlineto{\pgfqpoint{5.454944in}{0.667744in}}%
\pgfpathlineto{\pgfqpoint{5.458576in}{0.536675in}}%
\pgfpathlineto{\pgfqpoint{5.462207in}{0.636467in}}%
\pgfpathlineto{\pgfqpoint{5.465838in}{0.664366in}}%
\pgfpathlineto{\pgfqpoint{5.469469in}{0.677409in}}%
\pgfpathlineto{\pgfqpoint{5.473100in}{0.681021in}}%
\pgfpathlineto{\pgfqpoint{5.476731in}{0.631271in}}%
\pgfpathlineto{\pgfqpoint{5.480363in}{0.614946in}}%
\pgfpathlineto{\pgfqpoint{5.483994in}{0.623768in}}%
\pgfpathlineto{\pgfqpoint{5.487625in}{0.577820in}}%
\pgfpathlineto{\pgfqpoint{5.491256in}{0.665437in}}%
\pgfpathlineto{\pgfqpoint{5.494887in}{0.683595in}}%
\pgfpathlineto{\pgfqpoint{5.498519in}{0.636739in}}%
\pgfpathlineto{\pgfqpoint{5.502150in}{0.622552in}}%
\pgfpathlineto{\pgfqpoint{5.505781in}{0.646566in}}%
\pgfpathlineto{\pgfqpoint{5.509412in}{0.644519in}}%
\pgfpathlineto{\pgfqpoint{5.513043in}{0.634657in}}%
\pgfpathlineto{\pgfqpoint{5.516674in}{0.610173in}}%
\pgfpathlineto{\pgfqpoint{5.520306in}{0.543293in}}%
\pgfpathlineto{\pgfqpoint{5.523937in}{0.672886in}}%
\pgfpathlineto{\pgfqpoint{5.527568in}{0.585679in}}%
\pgfpathlineto{\pgfqpoint{5.531199in}{0.613143in}}%
\pgfpathlineto{\pgfqpoint{5.534830in}{0.657855in}}%
\pgfpathlineto{\pgfqpoint{5.538461in}{0.676511in}}%
\pgfpathlineto{\pgfqpoint{5.542093in}{0.685236in}}%
\pgfpathlineto{\pgfqpoint{5.545724in}{0.584352in}}%
\pgfpathlineto{\pgfqpoint{5.549355in}{0.612058in}}%
\pgfpathlineto{\pgfqpoint{5.552986in}{0.647945in}}%
\pgfpathlineto{\pgfqpoint{5.556617in}{0.585861in}}%
\pgfpathlineto{\pgfqpoint{5.560248in}{0.645064in}}%
\pgfpathlineto{\pgfqpoint{5.563880in}{0.671774in}}%
\pgfpathlineto{\pgfqpoint{5.567511in}{0.641685in}}%
\pgfpathlineto{\pgfqpoint{5.571142in}{0.645378in}}%
\pgfpathlineto{\pgfqpoint{5.574773in}{0.634205in}}%
\pgfpathlineto{\pgfqpoint{5.578404in}{0.647440in}}%
\pgfpathlineto{\pgfqpoint{5.582035in}{0.656818in}}%
\pgfpathlineto{\pgfqpoint{5.585667in}{0.643009in}}%
\pgfpathlineto{\pgfqpoint{5.589298in}{0.563148in}}%
\pgfpathlineto{\pgfqpoint{5.592929in}{0.676071in}}%
\pgfpathlineto{\pgfqpoint{5.596560in}{0.653354in}}%
\pgfpathlineto{\pgfqpoint{5.600191in}{0.661377in}}%
\pgfpathlineto{\pgfqpoint{5.603823in}{0.642607in}}%
\pgfpathlineto{\pgfqpoint{5.611085in}{0.623898in}}%
\pgfpathlineto{\pgfqpoint{5.614716in}{0.533912in}}%
\pgfpathlineto{\pgfqpoint{5.618347in}{0.659911in}}%
\pgfpathlineto{\pgfqpoint{5.621978in}{0.687398in}}%
\pgfpathlineto{\pgfqpoint{5.625610in}{0.673363in}}%
\pgfpathlineto{\pgfqpoint{5.629241in}{0.676774in}}%
\pgfpathlineto{\pgfqpoint{5.632872in}{0.650372in}}%
\pgfpathlineto{\pgfqpoint{5.640134in}{0.591202in}}%
\pgfpathlineto{\pgfqpoint{5.643765in}{0.615742in}}%
\pgfpathlineto{\pgfqpoint{5.647397in}{0.673568in}}%
\pgfpathlineto{\pgfqpoint{5.651028in}{0.690740in}}%
\pgfpathlineto{\pgfqpoint{5.654659in}{0.627465in}}%
\pgfpathlineto{\pgfqpoint{5.658290in}{0.594136in}}%
\pgfpathlineto{\pgfqpoint{5.661921in}{0.650927in}}%
\pgfpathlineto{\pgfqpoint{5.665552in}{0.636426in}}%
\pgfpathlineto{\pgfqpoint{5.669184in}{0.577541in}}%
\pgfpathlineto{\pgfqpoint{5.672815in}{0.561331in}}%
\pgfpathlineto{\pgfqpoint{5.676446in}{0.584806in}}%
\pgfpathlineto{\pgfqpoint{5.680077in}{0.671793in}}%
\pgfpathlineto{\pgfqpoint{5.683708in}{0.426681in}}%
\pgfpathlineto{\pgfqpoint{5.687340in}{0.630815in}}%
\pgfpathlineto{\pgfqpoint{5.690971in}{0.662255in}}%
\pgfpathlineto{\pgfqpoint{5.694602in}{0.675428in}}%
\pgfpathlineto{\pgfqpoint{5.698233in}{0.679164in}}%
\pgfpathlineto{\pgfqpoint{5.701864in}{0.625955in}}%
\pgfpathlineto{\pgfqpoint{5.705495in}{0.618280in}}%
\pgfpathlineto{\pgfqpoint{5.709127in}{0.632981in}}%
\pgfpathlineto{\pgfqpoint{5.712758in}{0.542149in}}%
\pgfpathlineto{\pgfqpoint{5.716389in}{0.660737in}}%
\pgfpathlineto{\pgfqpoint{5.720020in}{0.681412in}}%
\pgfpathlineto{\pgfqpoint{5.723651in}{0.639080in}}%
\pgfpathlineto{\pgfqpoint{5.727282in}{0.627755in}}%
\pgfpathlineto{\pgfqpoint{5.730914in}{0.645335in}}%
\pgfpathlineto{\pgfqpoint{5.734545in}{0.646041in}}%
\pgfpathlineto{\pgfqpoint{5.738176in}{0.641719in}}%
\pgfpathlineto{\pgfqpoint{5.741807in}{0.616897in}}%
\pgfpathlineto{\pgfqpoint{5.745438in}{0.548222in}}%
\pgfpathlineto{\pgfqpoint{5.749069in}{0.672772in}}%
\pgfpathlineto{\pgfqpoint{5.752701in}{0.597067in}}%
\pgfpathlineto{\pgfqpoint{5.756332in}{0.605315in}}%
\pgfpathlineto{\pgfqpoint{5.759963in}{0.656802in}}%
\pgfpathlineto{\pgfqpoint{5.763594in}{0.675550in}}%
\pgfpathlineto{\pgfqpoint{5.767225in}{0.685038in}}%
\pgfpathlineto{\pgfqpoint{5.770856in}{0.568605in}}%
\pgfpathlineto{\pgfqpoint{5.778119in}{0.651876in}}%
\pgfpathlineto{\pgfqpoint{5.781750in}{0.598087in}}%
\pgfpathlineto{\pgfqpoint{5.785381in}{0.638817in}}%
\pgfpathlineto{\pgfqpoint{5.789012in}{0.669479in}}%
\pgfpathlineto{\pgfqpoint{5.796275in}{0.696381in}}%
\pgfpathlineto{\pgfqpoint{5.799906in}{0.548284in}}%
\pgfpathlineto{\pgfqpoint{5.807168in}{0.651715in}}%
\pgfpathlineto{\pgfqpoint{5.810799in}{0.528112in}}%
\pgfpathlineto{\pgfqpoint{5.814431in}{0.649512in}}%
\pgfpathlineto{\pgfqpoint{5.818062in}{0.673336in}}%
\pgfpathlineto{\pgfqpoint{5.821693in}{0.685321in}}%
\pgfpathlineto{\pgfqpoint{5.825324in}{0.691665in}}%
\pgfpathlineto{\pgfqpoint{5.828955in}{0.507612in}}%
\pgfpathlineto{\pgfqpoint{5.836218in}{0.661933in}}%
\pgfpathlineto{\pgfqpoint{5.839849in}{0.591702in}}%
\pgfpathlineto{\pgfqpoint{5.843480in}{0.642724in}}%
\pgfpathlineto{\pgfqpoint{5.847111in}{0.665467in}}%
\pgfpathlineto{\pgfqpoint{5.850742in}{0.673496in}}%
\pgfpathlineto{\pgfqpoint{5.854373in}{0.672503in}}%
\pgfpathlineto{\pgfqpoint{5.858005in}{0.647276in}}%
\pgfpathlineto{\pgfqpoint{5.861636in}{0.616486in}}%
\pgfpathlineto{\pgfqpoint{5.865267in}{0.602583in}}%
\pgfpathlineto{\pgfqpoint{5.868898in}{0.612459in}}%
\pgfpathlineto{\pgfqpoint{5.872529in}{0.671441in}}%
\pgfpathlineto{\pgfqpoint{5.876160in}{0.688437in}}%
\pgfpathlineto{\pgfqpoint{5.879792in}{0.630711in}}%
\pgfpathlineto{\pgfqpoint{5.883423in}{0.597475in}}%
\pgfpathlineto{\pgfqpoint{5.887054in}{0.651573in}}%
\pgfpathlineto{\pgfqpoint{5.890685in}{0.640730in}}%
\pgfpathlineto{\pgfqpoint{5.894316in}{0.607614in}}%
\pgfpathlineto{\pgfqpoint{5.897948in}{0.486230in}}%
\pgfpathlineto{\pgfqpoint{5.905210in}{0.672738in}}%
\pgfpathlineto{\pgfqpoint{5.908841in}{0.530223in}}%
\pgfpathlineto{\pgfqpoint{5.912472in}{0.627566in}}%
\pgfpathlineto{\pgfqpoint{5.916103in}{0.662161in}}%
\pgfpathlineto{\pgfqpoint{5.919735in}{0.674790in}}%
\pgfpathlineto{\pgfqpoint{5.923366in}{0.678842in}}%
\pgfpathlineto{\pgfqpoint{5.926997in}{0.621373in}}%
\pgfpathlineto{\pgfqpoint{5.930628in}{0.619934in}}%
\pgfpathlineto{\pgfqpoint{5.934259in}{0.638649in}}%
\pgfpathlineto{\pgfqpoint{5.937890in}{0.444298in}}%
\pgfpathlineto{\pgfqpoint{5.941522in}{0.657014in}}%
\pgfpathlineto{\pgfqpoint{5.945153in}{0.679764in}}%
\pgfpathlineto{\pgfqpoint{5.948784in}{0.640742in}}%
\pgfpathlineto{\pgfqpoint{5.952415in}{0.630991in}}%
\pgfpathlineto{\pgfqpoint{5.956046in}{0.644281in}}%
\pgfpathlineto{\pgfqpoint{5.959677in}{0.647149in}}%
\pgfpathlineto{\pgfqpoint{5.963309in}{0.646411in}}%
\pgfpathlineto{\pgfqpoint{5.966940in}{0.620608in}}%
\pgfpathlineto{\pgfqpoint{5.970571in}{0.629323in}}%
\pgfpathlineto{\pgfqpoint{5.974202in}{0.686103in}}%
\pgfpathlineto{\pgfqpoint{5.977833in}{0.653266in}}%
\pgfpathlineto{\pgfqpoint{5.981465in}{0.649924in}}%
\pgfpathlineto{\pgfqpoint{5.985096in}{0.651184in}}%
\pgfpathlineto{\pgfqpoint{5.988727in}{0.630591in}}%
\pgfpathlineto{\pgfqpoint{5.992358in}{0.595431in}}%
\pgfpathlineto{\pgfqpoint{5.995989in}{0.592132in}}%
\pgfpathlineto{\pgfqpoint{5.999620in}{0.670536in}}%
\pgfpathlineto{\pgfqpoint{6.003252in}{0.694020in}}%
\pgfpathlineto{\pgfqpoint{6.006883in}{0.671872in}}%
\pgfpathlineto{\pgfqpoint{6.010514in}{0.666754in}}%
\pgfpathlineto{\pgfqpoint{6.014145in}{0.659453in}}%
\pgfpathlineto{\pgfqpoint{6.017776in}{0.614680in}}%
\pgfpathlineto{\pgfqpoint{6.021407in}{0.497750in}}%
\pgfpathlineto{\pgfqpoint{6.025039in}{0.633877in}}%
\pgfpathlineto{\pgfqpoint{6.028670in}{0.677638in}}%
\pgfpathlineto{\pgfqpoint{6.032301in}{0.693677in}}%
\pgfpathlineto{\pgfqpoint{6.039563in}{0.535123in}}%
\pgfpathlineto{\pgfqpoint{6.043194in}{0.652228in}}%
\pgfpathlineto{\pgfqpoint{6.046826in}{0.628582in}}%
\pgfpathlineto{\pgfqpoint{6.050457in}{0.575395in}}%
\pgfpathlineto{\pgfqpoint{6.054088in}{0.625280in}}%
\pgfpathlineto{\pgfqpoint{6.057719in}{0.614951in}}%
\pgfpathlineto{\pgfqpoint{6.061350in}{0.667423in}}%
\pgfpathlineto{\pgfqpoint{6.064981in}{0.578247in}}%
\pgfpathlineto{\pgfqpoint{6.068613in}{0.638279in}}%
\pgfpathlineto{\pgfqpoint{6.072244in}{0.663652in}}%
\pgfpathlineto{\pgfqpoint{6.075875in}{0.671668in}}%
\pgfpathlineto{\pgfqpoint{6.079506in}{0.670644in}}%
\pgfpathlineto{\pgfqpoint{6.083137in}{0.643685in}}%
\pgfpathlineto{\pgfqpoint{6.086769in}{0.620895in}}%
\pgfpathlineto{\pgfqpoint{6.090400in}{0.616871in}}%
\pgfpathlineto{\pgfqpoint{6.094031in}{0.600929in}}%
\pgfpathlineto{\pgfqpoint{6.097662in}{0.667853in}}%
\pgfpathlineto{\pgfqpoint{6.101293in}{0.686688in}}%
\pgfpathlineto{\pgfqpoint{6.104924in}{0.634174in}}%
\pgfpathlineto{\pgfqpoint{6.108556in}{0.605504in}}%
\pgfpathlineto{\pgfqpoint{6.112187in}{0.651372in}}%
\pgfpathlineto{\pgfqpoint{6.115818in}{0.643176in}}%
\pgfpathlineto{\pgfqpoint{6.119449in}{0.621501in}}%
\pgfpathlineto{\pgfqpoint{6.123080in}{0.533990in}}%
\pgfpathlineto{\pgfqpoint{6.126711in}{0.581250in}}%
\pgfpathlineto{\pgfqpoint{6.130343in}{0.673363in}}%
\pgfpathlineto{\pgfqpoint{6.133974in}{0.564159in}}%
\pgfpathlineto{\pgfqpoint{6.137605in}{0.623265in}}%
\pgfpathlineto{\pgfqpoint{6.141236in}{0.661868in}}%
\pgfpathlineto{\pgfqpoint{6.144867in}{0.674227in}}%
\pgfpathlineto{\pgfqpoint{6.148498in}{0.678876in}}%
\pgfpathlineto{\pgfqpoint{6.152130in}{0.615248in}}%
\pgfpathlineto{\pgfqpoint{6.155761in}{0.621073in}}%
\pgfpathlineto{\pgfqpoint{6.159392in}{0.643991in}}%
\pgfpathlineto{\pgfqpoint{6.163023in}{0.537905in}}%
\pgfpathlineto{\pgfqpoint{6.166654in}{0.652579in}}%
\pgfpathlineto{\pgfqpoint{6.170285in}{0.677855in}}%
\pgfpathlineto{\pgfqpoint{6.173917in}{0.682863in}}%
\pgfpathlineto{\pgfqpoint{6.177548in}{0.690627in}}%
\pgfpathlineto{\pgfqpoint{6.181179in}{0.617671in}}%
\pgfpathlineto{\pgfqpoint{6.184810in}{0.612018in}}%
\pgfpathlineto{\pgfqpoint{6.188441in}{0.642985in}}%
\pgfpathlineto{\pgfqpoint{6.192073in}{0.592469in}}%
\pgfpathlineto{\pgfqpoint{6.195704in}{0.656081in}}%
\pgfpathlineto{\pgfqpoint{6.199335in}{0.676961in}}%
\pgfpathlineto{\pgfqpoint{6.202966in}{0.683237in}}%
\pgfpathlineto{\pgfqpoint{6.206597in}{0.683643in}}%
\pgfpathlineto{\pgfqpoint{6.210228in}{0.623119in}}%
\pgfpathlineto{\pgfqpoint{6.213860in}{0.612760in}}%
\pgfpathlineto{\pgfqpoint{6.217491in}{0.654336in}}%
\pgfpathlineto{\pgfqpoint{6.221122in}{0.614981in}}%
\pgfpathlineto{\pgfqpoint{6.224753in}{0.646426in}}%
\pgfpathlineto{\pgfqpoint{6.228384in}{0.664089in}}%
\pgfpathlineto{\pgfqpoint{6.232015in}{0.667456in}}%
\pgfpathlineto{\pgfqpoint{6.235647in}{0.661529in}}%
\pgfpathlineto{\pgfqpoint{6.239278in}{0.658051in}}%
\pgfpathlineto{\pgfqpoint{6.242909in}{0.614303in}}%
\pgfpathlineto{\pgfqpoint{6.246540in}{0.558168in}}%
\pgfpathlineto{\pgfqpoint{6.250171in}{0.630769in}}%
\pgfpathlineto{\pgfqpoint{6.253802in}{0.675368in}}%
\pgfpathlineto{\pgfqpoint{6.257434in}{0.691292in}}%
\pgfpathlineto{\pgfqpoint{6.264696in}{0.542381in}}%
\pgfpathlineto{\pgfqpoint{6.268327in}{0.654038in}}%
\pgfpathlineto{\pgfqpoint{6.271958in}{0.634451in}}%
\pgfpathlineto{\pgfqpoint{6.275590in}{0.487978in}}%
\pgfpathlineto{\pgfqpoint{6.279221in}{0.615828in}}%
\pgfpathlineto{\pgfqpoint{6.282852in}{0.614638in}}%
\pgfpathlineto{\pgfqpoint{6.286483in}{0.669260in}}%
\pgfpathlineto{\pgfqpoint{6.290114in}{0.560191in}}%
\pgfpathlineto{\pgfqpoint{6.293745in}{0.636399in}}%
\pgfpathlineto{\pgfqpoint{6.297377in}{0.664244in}}%
\pgfpathlineto{\pgfqpoint{6.301008in}{0.671401in}}%
\pgfpathlineto{\pgfqpoint{6.304639in}{0.670476in}}%
\pgfpathlineto{\pgfqpoint{6.308270in}{0.640791in}}%
\pgfpathlineto{\pgfqpoint{6.311901in}{0.623345in}}%
\pgfpathlineto{\pgfqpoint{6.315532in}{0.625101in}}%
\pgfpathlineto{\pgfqpoint{6.319164in}{0.590424in}}%
\pgfpathlineto{\pgfqpoint{6.322795in}{0.664963in}}%
\pgfpathlineto{\pgfqpoint{6.326426in}{0.685287in}}%
\pgfpathlineto{\pgfqpoint{6.330057in}{0.636458in}}%
\pgfpathlineto{\pgfqpoint{6.333688in}{0.610225in}}%
\pgfpathlineto{\pgfqpoint{6.337319in}{0.651107in}}%
\pgfpathlineto{\pgfqpoint{6.340951in}{0.644934in}}%
\pgfpathlineto{\pgfqpoint{6.344582in}{0.629664in}}%
\pgfpathlineto{\pgfqpoint{6.348213in}{0.559966in}}%
\pgfpathlineto{\pgfqpoint{6.351844in}{0.652130in}}%
\pgfpathlineto{\pgfqpoint{6.355475in}{0.693084in}}%
\pgfpathlineto{\pgfqpoint{6.359106in}{0.650001in}}%
\pgfpathlineto{\pgfqpoint{6.362738in}{0.634873in}}%
\pgfpathlineto{\pgfqpoint{6.366369in}{0.656863in}}%
\pgfpathlineto{\pgfqpoint{6.370000in}{0.588508in}}%
\pgfpathlineto{\pgfqpoint{6.370000in}{0.588508in}}%
\pgfusepath{stroke}%
\end{pgfscope}%
\begin{pgfscope}%
\pgfpathrectangle{\pgfqpoint{4.554414in}{0.320625in}}{\pgfqpoint{1.815586in}{1.164375in}}%
\pgfusepath{clip}%
\pgfsetrectcap%
\pgfsetroundjoin%
\pgfsetlinewidth{0.602250pt}%
\definecolor{currentstroke}{rgb}{0.000000,0.529412,0.423529}%
\pgfsetstrokecolor{currentstroke}%
\pgfsetstrokeopacity{0.120000}%
\pgfsetdash{}{0pt}%
\pgfpathmoveto{\pgfqpoint{4.554414in}{1.456535in}}%
\pgfpathlineto{\pgfqpoint{4.558045in}{1.444744in}}%
\pgfpathlineto{\pgfqpoint{4.565307in}{1.429096in}}%
\pgfpathlineto{\pgfqpoint{4.579832in}{1.400124in}}%
\pgfpathlineto{\pgfqpoint{4.597988in}{1.361255in}}%
\pgfpathlineto{\pgfqpoint{4.666980in}{1.210945in}}%
\pgfpathlineto{\pgfqpoint{4.703292in}{1.135482in}}%
\pgfpathlineto{\pgfqpoint{4.743235in}{1.055138in}}%
\pgfpathlineto{\pgfqpoint{4.790440in}{0.962861in}}%
\pgfpathlineto{\pgfqpoint{4.844907in}{0.858999in}}%
\pgfpathlineto{\pgfqpoint{4.906637in}{0.743716in}}%
\pgfpathlineto{\pgfqpoint{4.913900in}{0.731046in}}%
\pgfpathlineto{\pgfqpoint{4.917531in}{0.725899in}}%
\pgfpathlineto{\pgfqpoint{4.921162in}{0.716928in}}%
\pgfpathlineto{\pgfqpoint{4.924793in}{0.710185in}}%
\pgfpathlineto{\pgfqpoint{4.928424in}{0.704810in}}%
\pgfpathlineto{\pgfqpoint{4.935687in}{0.690473in}}%
\pgfpathlineto{\pgfqpoint{4.939318in}{0.687230in}}%
\pgfpathlineto{\pgfqpoint{4.942949in}{0.687926in}}%
\pgfpathlineto{\pgfqpoint{4.946580in}{0.694152in}}%
\pgfpathlineto{\pgfqpoint{4.950211in}{0.663483in}}%
\pgfpathlineto{\pgfqpoint{4.953843in}{0.657627in}}%
\pgfpathlineto{\pgfqpoint{4.957474in}{0.659103in}}%
\pgfpathlineto{\pgfqpoint{4.961105in}{0.644057in}}%
\pgfpathlineto{\pgfqpoint{4.964736in}{0.655185in}}%
\pgfpathlineto{\pgfqpoint{4.968367in}{0.669176in}}%
\pgfpathlineto{\pgfqpoint{4.971998in}{0.672511in}}%
\pgfpathlineto{\pgfqpoint{4.975630in}{0.667698in}}%
\pgfpathlineto{\pgfqpoint{4.979261in}{0.641613in}}%
\pgfpathlineto{\pgfqpoint{4.982892in}{0.622645in}}%
\pgfpathlineto{\pgfqpoint{4.986523in}{0.633403in}}%
\pgfpathlineto{\pgfqpoint{4.990154in}{0.627666in}}%
\pgfpathlineto{\pgfqpoint{4.993785in}{0.638293in}}%
\pgfpathlineto{\pgfqpoint{4.997417in}{0.640486in}}%
\pgfpathlineto{\pgfqpoint{5.001048in}{0.629408in}}%
\pgfpathlineto{\pgfqpoint{5.004679in}{0.602160in}}%
\pgfpathlineto{\pgfqpoint{5.008310in}{0.660861in}}%
\pgfpathlineto{\pgfqpoint{5.011941in}{0.571209in}}%
\pgfpathlineto{\pgfqpoint{5.019204in}{0.647426in}}%
\pgfpathlineto{\pgfqpoint{5.022835in}{0.665255in}}%
\pgfpathlineto{\pgfqpoint{5.026466in}{0.669891in}}%
\pgfpathlineto{\pgfqpoint{5.030097in}{0.613372in}}%
\pgfpathlineto{\pgfqpoint{5.033728in}{0.606613in}}%
\pgfpathlineto{\pgfqpoint{5.037360in}{0.614013in}}%
\pgfpathlineto{\pgfqpoint{5.040991in}{0.577592in}}%
\pgfpathlineto{\pgfqpoint{5.044622in}{0.656069in}}%
\pgfpathlineto{\pgfqpoint{5.048253in}{0.675723in}}%
\pgfpathlineto{\pgfqpoint{5.051884in}{0.619433in}}%
\pgfpathlineto{\pgfqpoint{5.055515in}{0.574839in}}%
\pgfpathlineto{\pgfqpoint{5.059147in}{0.642302in}}%
\pgfpathlineto{\pgfqpoint{5.062778in}{0.625351in}}%
\pgfpathlineto{\pgfqpoint{5.066409in}{0.514299in}}%
\pgfpathlineto{\pgfqpoint{5.070040in}{0.614197in}}%
\pgfpathlineto{\pgfqpoint{5.073671in}{0.613883in}}%
\pgfpathlineto{\pgfqpoint{5.077302in}{0.653601in}}%
\pgfpathlineto{\pgfqpoint{5.080934in}{0.586113in}}%
\pgfpathlineto{\pgfqpoint{5.084565in}{0.631268in}}%
\pgfpathlineto{\pgfqpoint{5.088196in}{0.652287in}}%
\pgfpathlineto{\pgfqpoint{5.091827in}{0.648235in}}%
\pgfpathlineto{\pgfqpoint{5.095458in}{0.631740in}}%
\pgfpathlineto{\pgfqpoint{5.099090in}{0.652148in}}%
\pgfpathlineto{\pgfqpoint{5.102721in}{0.608100in}}%
\pgfpathlineto{\pgfqpoint{5.106352in}{0.434711in}}%
\pgfpathlineto{\pgfqpoint{5.109983in}{0.634408in}}%
\pgfpathlineto{\pgfqpoint{5.113614in}{0.663352in}}%
\pgfpathlineto{\pgfqpoint{5.117245in}{0.676655in}}%
\pgfpathlineto{\pgfqpoint{5.120877in}{0.645596in}}%
\pgfpathlineto{\pgfqpoint{5.124508in}{0.592745in}}%
\pgfpathlineto{\pgfqpoint{5.128139in}{0.671299in}}%
\pgfpathlineto{\pgfqpoint{5.131770in}{0.570870in}}%
\pgfpathlineto{\pgfqpoint{5.135401in}{0.589011in}}%
\pgfpathlineto{\pgfqpoint{5.139032in}{0.644640in}}%
\pgfpathlineto{\pgfqpoint{5.142664in}{0.643751in}}%
\pgfpathlineto{\pgfqpoint{5.146295in}{0.615794in}}%
\pgfpathlineto{\pgfqpoint{5.149926in}{0.592619in}}%
\pgfpathlineto{\pgfqpoint{5.153557in}{0.643397in}}%
\pgfpathlineto{\pgfqpoint{5.157188in}{0.684807in}}%
\pgfpathlineto{\pgfqpoint{5.160819in}{0.554556in}}%
\pgfpathlineto{\pgfqpoint{5.164451in}{0.595168in}}%
\pgfpathlineto{\pgfqpoint{5.168082in}{0.629205in}}%
\pgfpathlineto{\pgfqpoint{5.171713in}{0.559954in}}%
\pgfpathlineto{\pgfqpoint{5.175344in}{0.643887in}}%
\pgfpathlineto{\pgfqpoint{5.178975in}{0.666428in}}%
\pgfpathlineto{\pgfqpoint{5.182606in}{0.624281in}}%
\pgfpathlineto{\pgfqpoint{5.186238in}{0.619333in}}%
\pgfpathlineto{\pgfqpoint{5.189869in}{0.633022in}}%
\pgfpathlineto{\pgfqpoint{5.193500in}{0.629985in}}%
\pgfpathlineto{\pgfqpoint{5.197131in}{0.607790in}}%
\pgfpathlineto{\pgfqpoint{5.200762in}{0.564537in}}%
\pgfpathlineto{\pgfqpoint{5.204394in}{0.535252in}}%
\pgfpathlineto{\pgfqpoint{5.208025in}{0.660571in}}%
\pgfpathlineto{\pgfqpoint{5.211656in}{0.429561in}}%
\pgfpathlineto{\pgfqpoint{5.215287in}{0.621234in}}%
\pgfpathlineto{\pgfqpoint{5.218918in}{0.651472in}}%
\pgfpathlineto{\pgfqpoint{5.222549in}{0.661316in}}%
\pgfpathlineto{\pgfqpoint{5.226181in}{0.659580in}}%
\pgfpathlineto{\pgfqpoint{5.233443in}{0.609249in}}%
\pgfpathlineto{\pgfqpoint{5.237074in}{0.589586in}}%
\pgfpathlineto{\pgfqpoint{5.240705in}{0.610674in}}%
\pgfpathlineto{\pgfqpoint{5.244336in}{0.661373in}}%
\pgfpathlineto{\pgfqpoint{5.247968in}{0.678970in}}%
\pgfpathlineto{\pgfqpoint{5.251599in}{0.603792in}}%
\pgfpathlineto{\pgfqpoint{5.255230in}{0.501949in}}%
\pgfpathlineto{\pgfqpoint{5.258861in}{0.644255in}}%
\pgfpathlineto{\pgfqpoint{5.262492in}{0.614870in}}%
\pgfpathlineto{\pgfqpoint{5.266123in}{0.594655in}}%
\pgfpathlineto{\pgfqpoint{5.269755in}{0.642056in}}%
\pgfpathlineto{\pgfqpoint{5.273386in}{0.625701in}}%
\pgfpathlineto{\pgfqpoint{5.277017in}{0.644154in}}%
\pgfpathlineto{\pgfqpoint{5.280648in}{0.612012in}}%
\pgfpathlineto{\pgfqpoint{5.284279in}{0.634922in}}%
\pgfpathlineto{\pgfqpoint{5.287910in}{0.646834in}}%
\pgfpathlineto{\pgfqpoint{5.291542in}{0.633001in}}%
\pgfpathlineto{\pgfqpoint{5.295173in}{0.597950in}}%
\pgfpathlineto{\pgfqpoint{5.298804in}{0.658984in}}%
\pgfpathlineto{\pgfqpoint{5.302435in}{0.597034in}}%
\pgfpathlineto{\pgfqpoint{5.306066in}{0.582756in}}%
\pgfpathlineto{\pgfqpoint{5.309698in}{0.644789in}}%
\pgfpathlineto{\pgfqpoint{5.313329in}{0.663089in}}%
\pgfpathlineto{\pgfqpoint{5.316960in}{0.669982in}}%
\pgfpathlineto{\pgfqpoint{5.324222in}{0.569700in}}%
\pgfpathlineto{\pgfqpoint{5.327853in}{0.673000in}}%
\pgfpathlineto{\pgfqpoint{5.331485in}{0.434051in}}%
\pgfpathlineto{\pgfqpoint{5.335116in}{0.607709in}}%
\pgfpathlineto{\pgfqpoint{5.338747in}{0.644603in}}%
\pgfpathlineto{\pgfqpoint{5.342378in}{0.634303in}}%
\pgfpathlineto{\pgfqpoint{5.346009in}{0.472987in}}%
\pgfpathlineto{\pgfqpoint{5.349640in}{0.640242in}}%
\pgfpathlineto{\pgfqpoint{5.353272in}{0.654870in}}%
\pgfpathlineto{\pgfqpoint{5.356903in}{0.680095in}}%
\pgfpathlineto{\pgfqpoint{5.360534in}{0.606145in}}%
\pgfpathlineto{\pgfqpoint{5.364165in}{0.605618in}}%
\pgfpathlineto{\pgfqpoint{5.367796in}{0.616052in}}%
\pgfpathlineto{\pgfqpoint{5.371427in}{0.550709in}}%
\pgfpathlineto{\pgfqpoint{5.375059in}{0.654919in}}%
\pgfpathlineto{\pgfqpoint{5.378690in}{0.675035in}}%
\pgfpathlineto{\pgfqpoint{5.382321in}{0.620628in}}%
\pgfpathlineto{\pgfqpoint{5.385952in}{0.593060in}}%
\pgfpathlineto{\pgfqpoint{5.389583in}{0.639186in}}%
\pgfpathlineto{\pgfqpoint{5.393215in}{0.624509in}}%
\pgfpathlineto{\pgfqpoint{5.396846in}{0.530419in}}%
\pgfpathlineto{\pgfqpoint{5.400477in}{0.597075in}}%
\pgfpathlineto{\pgfqpoint{5.404108in}{0.601094in}}%
\pgfpathlineto{\pgfqpoint{5.407739in}{0.655573in}}%
\pgfpathlineto{\pgfqpoint{5.411370in}{0.576174in}}%
\pgfpathlineto{\pgfqpoint{5.415002in}{0.629302in}}%
\pgfpathlineto{\pgfqpoint{5.418633in}{0.651782in}}%
\pgfpathlineto{\pgfqpoint{5.422264in}{0.653754in}}%
\pgfpathlineto{\pgfqpoint{5.425895in}{0.643783in}}%
\pgfpathlineto{\pgfqpoint{5.429526in}{0.648181in}}%
\pgfpathlineto{\pgfqpoint{5.433157in}{0.606599in}}%
\pgfpathlineto{\pgfqpoint{5.436789in}{0.505808in}}%
\pgfpathlineto{\pgfqpoint{5.440420in}{0.629522in}}%
\pgfpathlineto{\pgfqpoint{5.444051in}{0.664360in}}%
\pgfpathlineto{\pgfqpoint{5.447682in}{0.679082in}}%
\pgfpathlineto{\pgfqpoint{5.451313in}{0.559052in}}%
\pgfpathlineto{\pgfqpoint{5.454944in}{0.582142in}}%
\pgfpathlineto{\pgfqpoint{5.458576in}{0.641424in}}%
\pgfpathlineto{\pgfqpoint{5.462207in}{0.593796in}}%
\pgfpathlineto{\pgfqpoint{5.469469in}{0.658441in}}%
\pgfpathlineto{\pgfqpoint{5.473100in}{0.629902in}}%
\pgfpathlineto{\pgfqpoint{5.476731in}{0.629925in}}%
\pgfpathlineto{\pgfqpoint{5.480363in}{0.628259in}}%
\pgfpathlineto{\pgfqpoint{5.483994in}{0.635664in}}%
\pgfpathlineto{\pgfqpoint{5.487625in}{0.634747in}}%
\pgfpathlineto{\pgfqpoint{5.491256in}{0.601245in}}%
\pgfpathlineto{\pgfqpoint{5.494887in}{0.630029in}}%
\pgfpathlineto{\pgfqpoint{5.498519in}{0.679560in}}%
\pgfpathlineto{\pgfqpoint{5.502150in}{0.638628in}}%
\pgfpathlineto{\pgfqpoint{5.505781in}{0.620082in}}%
\pgfpathlineto{\pgfqpoint{5.509412in}{0.646775in}}%
\pgfpathlineto{\pgfqpoint{5.513043in}{0.608576in}}%
\pgfpathlineto{\pgfqpoint{5.516674in}{0.576776in}}%
\pgfpathlineto{\pgfqpoint{5.520306in}{0.635114in}}%
\pgfpathlineto{\pgfqpoint{5.523937in}{0.671895in}}%
\pgfpathlineto{\pgfqpoint{5.527568in}{0.689667in}}%
\pgfpathlineto{\pgfqpoint{5.534830in}{0.595846in}}%
\pgfpathlineto{\pgfqpoint{5.538461in}{0.658408in}}%
\pgfpathlineto{\pgfqpoint{5.542093in}{0.493392in}}%
\pgfpathlineto{\pgfqpoint{5.545724in}{0.621439in}}%
\pgfpathlineto{\pgfqpoint{5.549355in}{0.651838in}}%
\pgfpathlineto{\pgfqpoint{5.552986in}{0.666260in}}%
\pgfpathlineto{\pgfqpoint{5.556617in}{0.668353in}}%
\pgfpathlineto{\pgfqpoint{5.560248in}{0.627984in}}%
\pgfpathlineto{\pgfqpoint{5.563880in}{0.603416in}}%
\pgfpathlineto{\pgfqpoint{5.567511in}{0.589918in}}%
\pgfpathlineto{\pgfqpoint{5.571142in}{0.605198in}}%
\pgfpathlineto{\pgfqpoint{5.574773in}{0.662128in}}%
\pgfpathlineto{\pgfqpoint{5.578404in}{0.679057in}}%
\pgfpathlineto{\pgfqpoint{5.582035in}{0.609015in}}%
\pgfpathlineto{\pgfqpoint{5.585667in}{0.520314in}}%
\pgfpathlineto{\pgfqpoint{5.589298in}{0.642098in}}%
\pgfpathlineto{\pgfqpoint{5.592929in}{0.616904in}}%
\pgfpathlineto{\pgfqpoint{5.596560in}{0.585004in}}%
\pgfpathlineto{\pgfqpoint{5.600191in}{0.633606in}}%
\pgfpathlineto{\pgfqpoint{5.603823in}{0.619034in}}%
\pgfpathlineto{\pgfqpoint{5.607454in}{0.647604in}}%
\pgfpathlineto{\pgfqpoint{5.611085in}{0.605337in}}%
\pgfpathlineto{\pgfqpoint{5.614716in}{0.633355in}}%
\pgfpathlineto{\pgfqpoint{5.618347in}{0.647662in}}%
\pgfpathlineto{\pgfqpoint{5.621978in}{0.641483in}}%
\pgfpathlineto{\pgfqpoint{5.625610in}{0.619778in}}%
\pgfpathlineto{\pgfqpoint{5.629241in}{0.656464in}}%
\pgfpathlineto{\pgfqpoint{5.632872in}{0.598190in}}%
\pgfpathlineto{\pgfqpoint{5.636503in}{0.571753in}}%
\pgfpathlineto{\pgfqpoint{5.640134in}{0.641040in}}%
\pgfpathlineto{\pgfqpoint{5.643765in}{0.664823in}}%
\pgfpathlineto{\pgfqpoint{5.647397in}{0.674930in}}%
\pgfpathlineto{\pgfqpoint{5.651028in}{0.566240in}}%
\pgfpathlineto{\pgfqpoint{5.658290in}{0.634091in}}%
\pgfpathlineto{\pgfqpoint{5.661921in}{0.538273in}}%
\pgfpathlineto{\pgfqpoint{5.665552in}{0.642945in}}%
\pgfpathlineto{\pgfqpoint{5.669184in}{0.668368in}}%
\pgfpathlineto{\pgfqpoint{5.672815in}{0.628500in}}%
\pgfpathlineto{\pgfqpoint{5.676446in}{0.609160in}}%
\pgfpathlineto{\pgfqpoint{5.680077in}{0.638337in}}%
\pgfpathlineto{\pgfqpoint{5.683708in}{0.633584in}}%
\pgfpathlineto{\pgfqpoint{5.687340in}{0.613481in}}%
\pgfpathlineto{\pgfqpoint{5.690971in}{0.533273in}}%
\pgfpathlineto{\pgfqpoint{5.694602in}{0.650292in}}%
\pgfpathlineto{\pgfqpoint{5.698233in}{0.685576in}}%
\pgfpathlineto{\pgfqpoint{5.701864in}{0.628844in}}%
\pgfpathlineto{\pgfqpoint{5.705495in}{0.585462in}}%
\pgfpathlineto{\pgfqpoint{5.709127in}{0.649524in}}%
\pgfpathlineto{\pgfqpoint{5.712758in}{0.593194in}}%
\pgfpathlineto{\pgfqpoint{5.716389in}{0.612239in}}%
\pgfpathlineto{\pgfqpoint{5.720020in}{0.649354in}}%
\pgfpathlineto{\pgfqpoint{5.723651in}{0.673929in}}%
\pgfpathlineto{\pgfqpoint{5.727282in}{0.687262in}}%
\pgfpathlineto{\pgfqpoint{5.730914in}{0.608290in}}%
\pgfpathlineto{\pgfqpoint{5.734545in}{0.493019in}}%
\pgfpathlineto{\pgfqpoint{5.738176in}{0.654171in}}%
\pgfpathlineto{\pgfqpoint{5.741807in}{0.564792in}}%
\pgfpathlineto{\pgfqpoint{5.745438in}{0.630077in}}%
\pgfpathlineto{\pgfqpoint{5.749069in}{0.653712in}}%
\pgfpathlineto{\pgfqpoint{5.752701in}{0.660435in}}%
\pgfpathlineto{\pgfqpoint{5.756332in}{0.655664in}}%
\pgfpathlineto{\pgfqpoint{5.759963in}{0.643989in}}%
\pgfpathlineto{\pgfqpoint{5.763594in}{0.602588in}}%
\pgfpathlineto{\pgfqpoint{5.767225in}{0.524078in}}%
\pgfpathlineto{\pgfqpoint{5.770856in}{0.625140in}}%
\pgfpathlineto{\pgfqpoint{5.774488in}{0.665447in}}%
\pgfpathlineto{\pgfqpoint{5.778119in}{0.680550in}}%
\pgfpathlineto{\pgfqpoint{5.781750in}{0.579439in}}%
\pgfpathlineto{\pgfqpoint{5.785381in}{0.563447in}}%
\pgfpathlineto{\pgfqpoint{5.789012in}{0.640411in}}%
\pgfpathlineto{\pgfqpoint{5.792644in}{0.599397in}}%
\pgfpathlineto{\pgfqpoint{5.796275in}{0.623300in}}%
\pgfpathlineto{\pgfqpoint{5.799906in}{0.653623in}}%
\pgfpathlineto{\pgfqpoint{5.803537in}{0.626849in}}%
\pgfpathlineto{\pgfqpoint{5.807168in}{0.635446in}}%
\pgfpathlineto{\pgfqpoint{5.810799in}{0.623175in}}%
\pgfpathlineto{\pgfqpoint{5.814431in}{0.634780in}}%
\pgfpathlineto{\pgfqpoint{5.818062in}{0.637529in}}%
\pgfpathlineto{\pgfqpoint{5.821693in}{0.618921in}}%
\pgfpathlineto{\pgfqpoint{5.825324in}{0.566642in}}%
\pgfpathlineto{\pgfqpoint{5.828955in}{0.660996in}}%
\pgfpathlineto{\pgfqpoint{5.832586in}{0.577066in}}%
\pgfpathlineto{\pgfqpoint{5.836218in}{0.604617in}}%
\pgfpathlineto{\pgfqpoint{5.839849in}{0.648823in}}%
\pgfpathlineto{\pgfqpoint{5.843480in}{0.663079in}}%
\pgfpathlineto{\pgfqpoint{5.847111in}{0.666467in}}%
\pgfpathlineto{\pgfqpoint{5.850742in}{0.614243in}}%
\pgfpathlineto{\pgfqpoint{5.854373in}{0.611717in}}%
\pgfpathlineto{\pgfqpoint{5.858005in}{0.620880in}}%
\pgfpathlineto{\pgfqpoint{5.861636in}{0.571856in}}%
\pgfpathlineto{\pgfqpoint{5.865267in}{0.653109in}}%
\pgfpathlineto{\pgfqpoint{5.868898in}{0.674780in}}%
\pgfpathlineto{\pgfqpoint{5.872529in}{0.668829in}}%
\pgfpathlineto{\pgfqpoint{5.876160in}{0.661593in}}%
\pgfpathlineto{\pgfqpoint{5.879792in}{0.650789in}}%
\pgfpathlineto{\pgfqpoint{5.883423in}{0.607105in}}%
\pgfpathlineto{\pgfqpoint{5.887054in}{0.578533in}}%
\pgfpathlineto{\pgfqpoint{5.890685in}{0.630835in}}%
\pgfpathlineto{\pgfqpoint{5.894316in}{0.652075in}}%
\pgfpathlineto{\pgfqpoint{5.897948in}{0.660326in}}%
\pgfpathlineto{\pgfqpoint{5.901579in}{0.643539in}}%
\pgfpathlineto{\pgfqpoint{5.905210in}{0.561126in}}%
\pgfpathlineto{\pgfqpoint{5.908841in}{0.676952in}}%
\pgfpathlineto{\pgfqpoint{5.916103in}{0.526462in}}%
\pgfpathlineto{\pgfqpoint{5.919735in}{0.640665in}}%
\pgfpathlineto{\pgfqpoint{5.923366in}{0.622875in}}%
\pgfpathlineto{\pgfqpoint{5.926997in}{0.544926in}}%
\pgfpathlineto{\pgfqpoint{5.930628in}{0.619998in}}%
\pgfpathlineto{\pgfqpoint{5.934259in}{0.606667in}}%
\pgfpathlineto{\pgfqpoint{5.937890in}{0.650708in}}%
\pgfpathlineto{\pgfqpoint{5.941522in}{0.599843in}}%
\pgfpathlineto{\pgfqpoint{5.945153in}{0.631067in}}%
\pgfpathlineto{\pgfqpoint{5.948784in}{0.646517in}}%
\pgfpathlineto{\pgfqpoint{5.952415in}{0.647470in}}%
\pgfpathlineto{\pgfqpoint{5.956046in}{0.633824in}}%
\pgfpathlineto{\pgfqpoint{5.959677in}{0.653635in}}%
\pgfpathlineto{\pgfqpoint{5.963309in}{0.597431in}}%
\pgfpathlineto{\pgfqpoint{5.966940in}{0.562101in}}%
\pgfpathlineto{\pgfqpoint{5.970571in}{0.637399in}}%
\pgfpathlineto{\pgfqpoint{5.974202in}{0.666501in}}%
\pgfpathlineto{\pgfqpoint{5.977833in}{0.678737in}}%
\pgfpathlineto{\pgfqpoint{5.981465in}{0.515018in}}%
\pgfpathlineto{\pgfqpoint{5.985096in}{0.592829in}}%
\pgfpathlineto{\pgfqpoint{5.988727in}{0.634059in}}%
\pgfpathlineto{\pgfqpoint{5.992358in}{0.561000in}}%
\pgfpathlineto{\pgfqpoint{5.995989in}{0.640892in}}%
\pgfpathlineto{\pgfqpoint{5.999620in}{0.665467in}}%
\pgfpathlineto{\pgfqpoint{6.003252in}{0.628142in}}%
\pgfpathlineto{\pgfqpoint{6.006883in}{0.618501in}}%
\pgfpathlineto{\pgfqpoint{6.010514in}{0.634279in}}%
\pgfpathlineto{\pgfqpoint{6.014145in}{0.633572in}}%
\pgfpathlineto{\pgfqpoint{6.017776in}{0.619961in}}%
\pgfpathlineto{\pgfqpoint{6.021407in}{0.557503in}}%
\pgfpathlineto{\pgfqpoint{6.025039in}{0.570812in}}%
\pgfpathlineto{\pgfqpoint{6.028670in}{0.661337in}}%
\pgfpathlineto{\pgfqpoint{6.032301in}{0.506156in}}%
\pgfpathlineto{\pgfqpoint{6.035932in}{0.619953in}}%
\pgfpathlineto{\pgfqpoint{6.039563in}{0.652980in}}%
\pgfpathlineto{\pgfqpoint{6.043194in}{0.658729in}}%
\pgfpathlineto{\pgfqpoint{6.046826in}{0.654599in}}%
\pgfpathlineto{\pgfqpoint{6.057719in}{0.599671in}}%
\pgfpathlineto{\pgfqpoint{6.061350in}{0.609697in}}%
\pgfpathlineto{\pgfqpoint{6.064981in}{0.658883in}}%
\pgfpathlineto{\pgfqpoint{6.068613in}{0.677597in}}%
\pgfpathlineto{\pgfqpoint{6.072244in}{0.661625in}}%
\pgfpathlineto{\pgfqpoint{6.075875in}{0.641242in}}%
\pgfpathlineto{\pgfqpoint{6.079506in}{0.662556in}}%
\pgfpathlineto{\pgfqpoint{6.083137in}{0.601605in}}%
\pgfpathlineto{\pgfqpoint{6.086769in}{0.429826in}}%
\pgfpathlineto{\pgfqpoint{6.090400in}{0.639526in}}%
\pgfpathlineto{\pgfqpoint{6.094031in}{0.650409in}}%
\pgfpathlineto{\pgfqpoint{6.097662in}{0.648872in}}%
\pgfpathlineto{\pgfqpoint{6.101293in}{0.610652in}}%
\pgfpathlineto{\pgfqpoint{6.104924in}{0.607643in}}%
\pgfpathlineto{\pgfqpoint{6.108556in}{0.682569in}}%
\pgfpathlineto{\pgfqpoint{6.112187in}{0.578037in}}%
\pgfpathlineto{\pgfqpoint{6.115818in}{0.559270in}}%
\pgfpathlineto{\pgfqpoint{6.119449in}{0.640134in}}%
\pgfpathlineto{\pgfqpoint{6.123080in}{0.609344in}}%
\pgfpathlineto{\pgfqpoint{6.126711in}{0.613014in}}%
\pgfpathlineto{\pgfqpoint{6.130343in}{0.647035in}}%
\pgfpathlineto{\pgfqpoint{6.133974in}{0.620820in}}%
\pgfpathlineto{\pgfqpoint{6.137605in}{0.640250in}}%
\pgfpathlineto{\pgfqpoint{6.141236in}{0.618522in}}%
\pgfpathlineto{\pgfqpoint{6.144867in}{0.632863in}}%
\pgfpathlineto{\pgfqpoint{6.148498in}{0.637476in}}%
\pgfpathlineto{\pgfqpoint{6.152130in}{0.629637in}}%
\pgfpathlineto{\pgfqpoint{6.155761in}{0.601166in}}%
\pgfpathlineto{\pgfqpoint{6.159392in}{0.659460in}}%
\pgfpathlineto{\pgfqpoint{6.163023in}{0.580204in}}%
\pgfpathlineto{\pgfqpoint{6.166654in}{0.599727in}}%
\pgfpathlineto{\pgfqpoint{6.170285in}{0.646058in}}%
\pgfpathlineto{\pgfqpoint{6.173917in}{0.665533in}}%
\pgfpathlineto{\pgfqpoint{6.177548in}{0.672076in}}%
\pgfpathlineto{\pgfqpoint{6.181179in}{0.601660in}}%
\pgfpathlineto{\pgfqpoint{6.184810in}{0.605839in}}%
\pgfpathlineto{\pgfqpoint{6.188441in}{0.622373in}}%
\pgfpathlineto{\pgfqpoint{6.192073in}{0.549269in}}%
\pgfpathlineto{\pgfqpoint{6.195704in}{0.652112in}}%
\pgfpathlineto{\pgfqpoint{6.199335in}{0.673402in}}%
\pgfpathlineto{\pgfqpoint{6.202966in}{0.623328in}}%
\pgfpathlineto{\pgfqpoint{6.206597in}{0.589109in}}%
\pgfpathlineto{\pgfqpoint{6.210228in}{0.641582in}}%
\pgfpathlineto{\pgfqpoint{6.213860in}{0.628945in}}%
\pgfpathlineto{\pgfqpoint{6.217491in}{0.577200in}}%
\pgfpathlineto{\pgfqpoint{6.221122in}{0.599002in}}%
\pgfpathlineto{\pgfqpoint{6.224753in}{0.610127in}}%
\pgfpathlineto{\pgfqpoint{6.228384in}{0.656350in}}%
\pgfpathlineto{\pgfqpoint{6.232015in}{0.569719in}}%
\pgfpathlineto{\pgfqpoint{6.235647in}{0.629332in}}%
\pgfpathlineto{\pgfqpoint{6.239278in}{0.653710in}}%
\pgfpathlineto{\pgfqpoint{6.242909in}{0.650829in}}%
\pgfpathlineto{\pgfqpoint{6.246540in}{0.619294in}}%
\pgfpathlineto{\pgfqpoint{6.250171in}{0.644863in}}%
\pgfpathlineto{\pgfqpoint{6.253802in}{0.638144in}}%
\pgfpathlineto{\pgfqpoint{6.257434in}{0.657055in}}%
\pgfpathlineto{\pgfqpoint{6.261065in}{0.624661in}}%
\pgfpathlineto{\pgfqpoint{6.264696in}{0.621387in}}%
\pgfpathlineto{\pgfqpoint{6.268327in}{0.613373in}}%
\pgfpathlineto{\pgfqpoint{6.271958in}{0.467344in}}%
\pgfpathlineto{\pgfqpoint{6.275590in}{0.652908in}}%
\pgfpathlineto{\pgfqpoint{6.279221in}{0.680074in}}%
\pgfpathlineto{\pgfqpoint{6.282852in}{0.660677in}}%
\pgfpathlineto{\pgfqpoint{6.286483in}{0.654367in}}%
\pgfpathlineto{\pgfqpoint{6.290114in}{0.649952in}}%
\pgfpathlineto{\pgfqpoint{6.297377in}{0.555691in}}%
\pgfpathlineto{\pgfqpoint{6.301008in}{0.634291in}}%
\pgfpathlineto{\pgfqpoint{6.304639in}{0.668741in}}%
\pgfpathlineto{\pgfqpoint{6.308270in}{0.682723in}}%
\pgfpathlineto{\pgfqpoint{6.311901in}{0.515927in}}%
\pgfpathlineto{\pgfqpoint{6.319164in}{0.632323in}}%
\pgfpathlineto{\pgfqpoint{6.322795in}{0.567528in}}%
\pgfpathlineto{\pgfqpoint{6.326426in}{0.640508in}}%
\pgfpathlineto{\pgfqpoint{6.330057in}{0.662481in}}%
\pgfpathlineto{\pgfqpoint{6.333688in}{0.626246in}}%
\pgfpathlineto{\pgfqpoint{6.337319in}{0.624833in}}%
\pgfpathlineto{\pgfqpoint{6.340951in}{0.630317in}}%
\pgfpathlineto{\pgfqpoint{6.344582in}{0.633193in}}%
\pgfpathlineto{\pgfqpoint{6.348213in}{0.624175in}}%
\pgfpathlineto{\pgfqpoint{6.351844in}{0.595956in}}%
\pgfpathlineto{\pgfqpoint{6.355475in}{0.450141in}}%
\pgfpathlineto{\pgfqpoint{6.359106in}{0.661358in}}%
\pgfpathlineto{\pgfqpoint{6.362738in}{0.534391in}}%
\pgfpathlineto{\pgfqpoint{6.366369in}{0.568638in}}%
\pgfpathlineto{\pgfqpoint{6.370000in}{0.624155in}}%
\pgfpathlineto{\pgfqpoint{6.370000in}{0.624155in}}%
\pgfusepath{stroke}%
\end{pgfscope}%
\begin{pgfscope}%
\pgfpathrectangle{\pgfqpoint{4.554414in}{0.320625in}}{\pgfqpoint{1.815586in}{1.164375in}}%
\pgfusepath{clip}%
\pgfsetrectcap%
\pgfsetroundjoin%
\pgfsetlinewidth{0.602250pt}%
\definecolor{currentstroke}{rgb}{0.000000,0.529412,0.423529}%
\pgfsetstrokecolor{currentstroke}%
\pgfsetstrokeopacity{0.120000}%
\pgfsetdash{}{0pt}%
\pgfpathmoveto{\pgfqpoint{4.554414in}{1.457396in}}%
\pgfpathlineto{\pgfqpoint{4.558045in}{1.445267in}}%
\pgfpathlineto{\pgfqpoint{4.565307in}{1.429207in}}%
\pgfpathlineto{\pgfqpoint{4.583463in}{1.391028in}}%
\pgfpathlineto{\pgfqpoint{4.608881in}{1.334174in}}%
\pgfpathlineto{\pgfqpoint{4.645193in}{1.253376in}}%
\pgfpathlineto{\pgfqpoint{4.674242in}{1.191552in}}%
\pgfpathlineto{\pgfqpoint{4.703292in}{1.132218in}}%
\pgfpathlineto{\pgfqpoint{4.739603in}{1.060767in}}%
\pgfpathlineto{\pgfqpoint{4.783177in}{0.977765in}}%
\pgfpathlineto{\pgfqpoint{4.837645in}{0.876648in}}%
\pgfpathlineto{\pgfqpoint{4.903006in}{0.757660in}}%
\pgfpathlineto{\pgfqpoint{4.924793in}{0.720276in}}%
\pgfpathlineto{\pgfqpoint{4.928424in}{0.712249in}}%
\pgfpathlineto{\pgfqpoint{4.939318in}{0.692494in}}%
\pgfpathlineto{\pgfqpoint{4.942949in}{0.691457in}}%
\pgfpathlineto{\pgfqpoint{4.946580in}{0.694101in}}%
\pgfpathlineto{\pgfqpoint{4.950211in}{0.674568in}}%
\pgfpathlineto{\pgfqpoint{4.953843in}{0.666602in}}%
\pgfpathlineto{\pgfqpoint{4.957474in}{0.668303in}}%
\pgfpathlineto{\pgfqpoint{4.961105in}{0.659138in}}%
\pgfpathlineto{\pgfqpoint{4.964736in}{0.647042in}}%
\pgfpathlineto{\pgfqpoint{4.968367in}{0.644116in}}%
\pgfpathlineto{\pgfqpoint{4.971998in}{0.639895in}}%
\pgfpathlineto{\pgfqpoint{4.975630in}{0.663666in}}%
\pgfpathlineto{\pgfqpoint{4.979261in}{0.622796in}}%
\pgfpathlineto{\pgfqpoint{4.982892in}{0.639288in}}%
\pgfpathlineto{\pgfqpoint{4.986523in}{0.658436in}}%
\pgfpathlineto{\pgfqpoint{4.990154in}{0.658957in}}%
\pgfpathlineto{\pgfqpoint{4.993785in}{0.650314in}}%
\pgfpathlineto{\pgfqpoint{4.997417in}{0.648927in}}%
\pgfpathlineto{\pgfqpoint{5.004679in}{0.592084in}}%
\pgfpathlineto{\pgfqpoint{5.008310in}{0.624937in}}%
\pgfpathlineto{\pgfqpoint{5.011941in}{0.664622in}}%
\pgfpathlineto{\pgfqpoint{5.015573in}{0.681626in}}%
\pgfpathlineto{\pgfqpoint{5.019204in}{0.601084in}}%
\pgfpathlineto{\pgfqpoint{5.022835in}{0.564582in}}%
\pgfpathlineto{\pgfqpoint{5.026466in}{0.649605in}}%
\pgfpathlineto{\pgfqpoint{5.030097in}{0.618883in}}%
\pgfpathlineto{\pgfqpoint{5.033728in}{0.594726in}}%
\pgfpathlineto{\pgfqpoint{5.037360in}{0.646486in}}%
\pgfpathlineto{\pgfqpoint{5.040991in}{0.670739in}}%
\pgfpathlineto{\pgfqpoint{5.044622in}{0.690851in}}%
\pgfpathlineto{\pgfqpoint{5.048253in}{0.564717in}}%
\pgfpathlineto{\pgfqpoint{5.051884in}{0.574458in}}%
\pgfpathlineto{\pgfqpoint{5.055515in}{0.646015in}}%
\pgfpathlineto{\pgfqpoint{5.059147in}{0.483807in}}%
\pgfpathlineto{\pgfqpoint{5.062778in}{0.640476in}}%
\pgfpathlineto{\pgfqpoint{5.066409in}{0.665316in}}%
\pgfpathlineto{\pgfqpoint{5.070040in}{0.676400in}}%
\pgfpathlineto{\pgfqpoint{5.073671in}{0.680030in}}%
\pgfpathlineto{\pgfqpoint{5.077302in}{0.593616in}}%
\pgfpathlineto{\pgfqpoint{5.080934in}{0.600666in}}%
\pgfpathlineto{\pgfqpoint{5.084565in}{0.646807in}}%
\pgfpathlineto{\pgfqpoint{5.088196in}{0.607499in}}%
\pgfpathlineto{\pgfqpoint{5.091827in}{0.638556in}}%
\pgfpathlineto{\pgfqpoint{5.095458in}{0.654451in}}%
\pgfpathlineto{\pgfqpoint{5.099090in}{0.654418in}}%
\pgfpathlineto{\pgfqpoint{5.102721in}{0.642877in}}%
\pgfpathlineto{\pgfqpoint{5.106352in}{0.656280in}}%
\pgfpathlineto{\pgfqpoint{5.109983in}{0.599681in}}%
\pgfpathlineto{\pgfqpoint{5.113614in}{0.554542in}}%
\pgfpathlineto{\pgfqpoint{5.117245in}{0.637735in}}%
\pgfpathlineto{\pgfqpoint{5.120877in}{0.669740in}}%
\pgfpathlineto{\pgfqpoint{5.124508in}{0.683119in}}%
\pgfpathlineto{\pgfqpoint{5.128139in}{0.538430in}}%
\pgfpathlineto{\pgfqpoint{5.131770in}{0.584508in}}%
\pgfpathlineto{\pgfqpoint{5.135401in}{0.641283in}}%
\pgfpathlineto{\pgfqpoint{5.139032in}{0.592788in}}%
\pgfpathlineto{\pgfqpoint{5.142664in}{0.633005in}}%
\pgfpathlineto{\pgfqpoint{5.146295in}{0.660269in}}%
\pgfpathlineto{\pgfqpoint{5.149926in}{0.630791in}}%
\pgfpathlineto{\pgfqpoint{5.153557in}{0.636935in}}%
\pgfpathlineto{\pgfqpoint{5.157188in}{0.627293in}}%
\pgfpathlineto{\pgfqpoint{5.160819in}{0.638097in}}%
\pgfpathlineto{\pgfqpoint{5.164451in}{0.641704in}}%
\pgfpathlineto{\pgfqpoint{5.168082in}{0.625286in}}%
\pgfpathlineto{\pgfqpoint{5.171713in}{0.582515in}}%
\pgfpathlineto{\pgfqpoint{5.175344in}{0.663628in}}%
\pgfpathlineto{\pgfqpoint{5.178975in}{0.589001in}}%
\pgfpathlineto{\pgfqpoint{5.182606in}{0.599248in}}%
\pgfpathlineto{\pgfqpoint{5.186238in}{0.649526in}}%
\pgfpathlineto{\pgfqpoint{5.189869in}{0.667182in}}%
\pgfpathlineto{\pgfqpoint{5.193500in}{0.674059in}}%
\pgfpathlineto{\pgfqpoint{5.197131in}{0.596863in}}%
\pgfpathlineto{\pgfqpoint{5.200762in}{0.610849in}}%
\pgfpathlineto{\pgfqpoint{5.204394in}{0.634541in}}%
\pgfpathlineto{\pgfqpoint{5.208025in}{0.481449in}}%
\pgfpathlineto{\pgfqpoint{5.211656in}{0.648169in}}%
\pgfpathlineto{\pgfqpoint{5.215287in}{0.672714in}}%
\pgfpathlineto{\pgfqpoint{5.222549in}{0.677016in}}%
\pgfpathlineto{\pgfqpoint{5.226181in}{0.634020in}}%
\pgfpathlineto{\pgfqpoint{5.229812in}{0.609805in}}%
\pgfpathlineto{\pgfqpoint{5.233443in}{0.617210in}}%
\pgfpathlineto{\pgfqpoint{5.237074in}{0.616525in}}%
\pgfpathlineto{\pgfqpoint{5.240705in}{0.653786in}}%
\pgfpathlineto{\pgfqpoint{5.244336in}{0.669924in}}%
\pgfpathlineto{\pgfqpoint{5.247968in}{0.665873in}}%
\pgfpathlineto{\pgfqpoint{5.251599in}{0.647392in}}%
\pgfpathlineto{\pgfqpoint{5.255230in}{0.662942in}}%
\pgfpathlineto{\pgfqpoint{5.258861in}{0.617289in}}%
\pgfpathlineto{\pgfqpoint{5.262492in}{0.612856in}}%
\pgfpathlineto{\pgfqpoint{5.266123in}{0.636424in}}%
\pgfpathlineto{\pgfqpoint{5.269755in}{0.637823in}}%
\pgfpathlineto{\pgfqpoint{5.273386in}{0.626513in}}%
\pgfpathlineto{\pgfqpoint{5.277017in}{0.603375in}}%
\pgfpathlineto{\pgfqpoint{5.280648in}{0.551648in}}%
\pgfpathlineto{\pgfqpoint{5.284279in}{0.664531in}}%
\pgfpathlineto{\pgfqpoint{5.287910in}{0.507889in}}%
\pgfpathlineto{\pgfqpoint{5.291542in}{0.619448in}}%
\pgfpathlineto{\pgfqpoint{5.295173in}{0.652252in}}%
\pgfpathlineto{\pgfqpoint{5.298804in}{0.667453in}}%
\pgfpathlineto{\pgfqpoint{5.302435in}{0.670405in}}%
\pgfpathlineto{\pgfqpoint{5.306066in}{0.623998in}}%
\pgfpathlineto{\pgfqpoint{5.309698in}{0.609771in}}%
\pgfpathlineto{\pgfqpoint{5.313329in}{0.612170in}}%
\pgfpathlineto{\pgfqpoint{5.316960in}{0.587049in}}%
\pgfpathlineto{\pgfqpoint{5.320591in}{0.660220in}}%
\pgfpathlineto{\pgfqpoint{5.324222in}{0.679127in}}%
\pgfpathlineto{\pgfqpoint{5.327853in}{0.622992in}}%
\pgfpathlineto{\pgfqpoint{5.331485in}{0.584328in}}%
\pgfpathlineto{\pgfqpoint{5.335116in}{0.644575in}}%
\pgfpathlineto{\pgfqpoint{5.338747in}{0.630096in}}%
\pgfpathlineto{\pgfqpoint{5.342378in}{0.565938in}}%
\pgfpathlineto{\pgfqpoint{5.346009in}{0.599337in}}%
\pgfpathlineto{\pgfqpoint{5.349640in}{0.607701in}}%
\pgfpathlineto{\pgfqpoint{5.353272in}{0.660308in}}%
\pgfpathlineto{\pgfqpoint{5.356903in}{0.564904in}}%
\pgfpathlineto{\pgfqpoint{5.360534in}{0.630752in}}%
\pgfpathlineto{\pgfqpoint{5.364165in}{0.656436in}}%
\pgfpathlineto{\pgfqpoint{5.367796in}{0.658482in}}%
\pgfpathlineto{\pgfqpoint{5.371427in}{0.650705in}}%
\pgfpathlineto{\pgfqpoint{5.375059in}{0.646299in}}%
\pgfpathlineto{\pgfqpoint{5.382321in}{0.584111in}}%
\pgfpathlineto{\pgfqpoint{5.389583in}{0.664125in}}%
\pgfpathlineto{\pgfqpoint{5.393215in}{0.681545in}}%
\pgfpathlineto{\pgfqpoint{5.396846in}{0.602329in}}%
\pgfpathlineto{\pgfqpoint{5.400477in}{0.545352in}}%
\pgfpathlineto{\pgfqpoint{5.404108in}{0.649491in}}%
\pgfpathlineto{\pgfqpoint{5.411370in}{0.588066in}}%
\pgfpathlineto{\pgfqpoint{5.415002in}{0.644356in}}%
\pgfpathlineto{\pgfqpoint{5.418633in}{0.670289in}}%
\pgfpathlineto{\pgfqpoint{5.422264in}{0.691198in}}%
\pgfpathlineto{\pgfqpoint{5.425895in}{0.574543in}}%
\pgfpathlineto{\pgfqpoint{5.429526in}{0.568290in}}%
\pgfpathlineto{\pgfqpoint{5.433157in}{0.646735in}}%
\pgfpathlineto{\pgfqpoint{5.436789in}{0.459729in}}%
\pgfpathlineto{\pgfqpoint{5.440420in}{0.639446in}}%
\pgfpathlineto{\pgfqpoint{5.444051in}{0.664709in}}%
\pgfpathlineto{\pgfqpoint{5.447682in}{0.676725in}}%
\pgfpathlineto{\pgfqpoint{5.451313in}{0.681222in}}%
\pgfpathlineto{\pgfqpoint{5.454944in}{0.583157in}}%
\pgfpathlineto{\pgfqpoint{5.458576in}{0.597724in}}%
\pgfpathlineto{\pgfqpoint{5.462207in}{0.647877in}}%
\pgfpathlineto{\pgfqpoint{5.465838in}{0.605158in}}%
\pgfpathlineto{\pgfqpoint{5.469469in}{0.638214in}}%
\pgfpathlineto{\pgfqpoint{5.473100in}{0.654889in}}%
\pgfpathlineto{\pgfqpoint{5.476731in}{0.655693in}}%
\pgfpathlineto{\pgfqpoint{5.480363in}{0.645136in}}%
\pgfpathlineto{\pgfqpoint{5.483994in}{0.655370in}}%
\pgfpathlineto{\pgfqpoint{5.491256in}{0.545277in}}%
\pgfpathlineto{\pgfqpoint{5.494887in}{0.636504in}}%
\pgfpathlineto{\pgfqpoint{5.498519in}{0.669675in}}%
\pgfpathlineto{\pgfqpoint{5.502150in}{0.683376in}}%
\pgfpathlineto{\pgfqpoint{5.505781in}{0.552237in}}%
\pgfpathlineto{\pgfqpoint{5.509412in}{0.581160in}}%
\pgfpathlineto{\pgfqpoint{5.513043in}{0.641915in}}%
\pgfpathlineto{\pgfqpoint{5.516674in}{0.596004in}}%
\pgfpathlineto{\pgfqpoint{5.520306in}{0.630853in}}%
\pgfpathlineto{\pgfqpoint{5.523937in}{0.658785in}}%
\pgfpathlineto{\pgfqpoint{5.527568in}{0.630416in}}%
\pgfpathlineto{\pgfqpoint{5.531199in}{0.638502in}}%
\pgfpathlineto{\pgfqpoint{5.534830in}{0.625802in}}%
\pgfpathlineto{\pgfqpoint{5.538461in}{0.638071in}}%
\pgfpathlineto{\pgfqpoint{5.542093in}{0.643028in}}%
\pgfpathlineto{\pgfqpoint{5.545724in}{0.628283in}}%
\pgfpathlineto{\pgfqpoint{5.549355in}{0.590021in}}%
\pgfpathlineto{\pgfqpoint{5.552986in}{0.663240in}}%
\pgfpathlineto{\pgfqpoint{5.556617in}{0.591388in}}%
\pgfpathlineto{\pgfqpoint{5.560248in}{0.596228in}}%
\pgfpathlineto{\pgfqpoint{5.563880in}{0.648714in}}%
\pgfpathlineto{\pgfqpoint{5.567511in}{0.667379in}}%
\pgfpathlineto{\pgfqpoint{5.571142in}{0.674961in}}%
\pgfpathlineto{\pgfqpoint{5.574773in}{0.591587in}}%
\pgfpathlineto{\pgfqpoint{5.578404in}{0.609708in}}%
\pgfpathlineto{\pgfqpoint{5.582035in}{0.635721in}}%
\pgfpathlineto{\pgfqpoint{5.585667in}{0.518017in}}%
\pgfpathlineto{\pgfqpoint{5.589298in}{0.646911in}}%
\pgfpathlineto{\pgfqpoint{5.592929in}{0.671887in}}%
\pgfpathlineto{\pgfqpoint{5.600191in}{0.678291in}}%
\pgfpathlineto{\pgfqpoint{5.603823in}{0.631115in}}%
\pgfpathlineto{\pgfqpoint{5.607454in}{0.609073in}}%
\pgfpathlineto{\pgfqpoint{5.611085in}{0.619290in}}%
\pgfpathlineto{\pgfqpoint{5.614716in}{0.614410in}}%
\pgfpathlineto{\pgfqpoint{5.618347in}{0.653488in}}%
\pgfpathlineto{\pgfqpoint{5.621978in}{0.670068in}}%
\pgfpathlineto{\pgfqpoint{5.625610in}{0.666976in}}%
\pgfpathlineto{\pgfqpoint{5.629241in}{0.650375in}}%
\pgfpathlineto{\pgfqpoint{5.632872in}{0.661112in}}%
\pgfpathlineto{\pgfqpoint{5.636503in}{0.617162in}}%
\pgfpathlineto{\pgfqpoint{5.640134in}{0.615436in}}%
\pgfpathlineto{\pgfqpoint{5.643765in}{0.635501in}}%
\pgfpathlineto{\pgfqpoint{5.647397in}{0.638130in}}%
\pgfpathlineto{\pgfqpoint{5.651028in}{0.628983in}}%
\pgfpathlineto{\pgfqpoint{5.654659in}{0.609206in}}%
\pgfpathlineto{\pgfqpoint{5.658290in}{0.565648in}}%
\pgfpathlineto{\pgfqpoint{5.661921in}{0.664539in}}%
\pgfpathlineto{\pgfqpoint{5.665552in}{0.524513in}}%
\pgfpathlineto{\pgfqpoint{5.669184in}{0.618061in}}%
\pgfpathlineto{\pgfqpoint{5.672815in}{0.651850in}}%
\pgfpathlineto{\pgfqpoint{5.676446in}{0.667909in}}%
\pgfpathlineto{\pgfqpoint{5.680077in}{0.671542in}}%
\pgfpathlineto{\pgfqpoint{5.683708in}{0.621235in}}%
\pgfpathlineto{\pgfqpoint{5.687340in}{0.609237in}}%
\pgfpathlineto{\pgfqpoint{5.690971in}{0.614438in}}%
\pgfpathlineto{\pgfqpoint{5.694602in}{0.581779in}}%
\pgfpathlineto{\pgfqpoint{5.698233in}{0.659441in}}%
\pgfpathlineto{\pgfqpoint{5.701864in}{0.678539in}}%
\pgfpathlineto{\pgfqpoint{5.705495in}{0.624056in}}%
\pgfpathlineto{\pgfqpoint{5.709127in}{0.589292in}}%
\pgfpathlineto{\pgfqpoint{5.712758in}{0.644066in}}%
\pgfpathlineto{\pgfqpoint{5.716389in}{0.630873in}}%
\pgfpathlineto{\pgfqpoint{5.720020in}{0.576459in}}%
\pgfpathlineto{\pgfqpoint{5.727282in}{0.604617in}}%
\pgfpathlineto{\pgfqpoint{5.730914in}{0.661049in}}%
\pgfpathlineto{\pgfqpoint{5.734545in}{0.557359in}}%
\pgfpathlineto{\pgfqpoint{5.738176in}{0.629879in}}%
\pgfpathlineto{\pgfqpoint{5.741807in}{0.656427in}}%
\pgfpathlineto{\pgfqpoint{5.745438in}{0.659375in}}%
\pgfpathlineto{\pgfqpoint{5.752701in}{0.644767in}}%
\pgfpathlineto{\pgfqpoint{5.759963in}{0.589016in}}%
\pgfpathlineto{\pgfqpoint{5.763594in}{0.618878in}}%
\pgfpathlineto{\pgfqpoint{5.767225in}{0.663698in}}%
\pgfpathlineto{\pgfqpoint{5.770856in}{0.681447in}}%
\pgfpathlineto{\pgfqpoint{5.778119in}{0.531680in}}%
\pgfpathlineto{\pgfqpoint{5.781750in}{0.649466in}}%
\pgfpathlineto{\pgfqpoint{5.785381in}{0.621560in}}%
\pgfpathlineto{\pgfqpoint{5.789012in}{0.580195in}}%
\pgfpathlineto{\pgfqpoint{5.792644in}{0.641945in}}%
\pgfpathlineto{\pgfqpoint{5.796275in}{0.669723in}}%
\pgfpathlineto{\pgfqpoint{5.799906in}{0.691462in}}%
\pgfpathlineto{\pgfqpoint{5.803537in}{0.583544in}}%
\pgfpathlineto{\pgfqpoint{5.807168in}{0.561089in}}%
\pgfpathlineto{\pgfqpoint{5.810799in}{0.647546in}}%
\pgfpathlineto{\pgfqpoint{5.814431in}{0.505804in}}%
\pgfpathlineto{\pgfqpoint{5.818062in}{0.638220in}}%
\pgfpathlineto{\pgfqpoint{5.821693in}{0.664004in}}%
\pgfpathlineto{\pgfqpoint{5.825324in}{0.676992in}}%
\pgfpathlineto{\pgfqpoint{5.828955in}{0.682388in}}%
\pgfpathlineto{\pgfqpoint{5.832586in}{0.568235in}}%
\pgfpathlineto{\pgfqpoint{5.836218in}{0.594337in}}%
\pgfpathlineto{\pgfqpoint{5.839849in}{0.649024in}}%
\pgfpathlineto{\pgfqpoint{5.843480in}{0.602405in}}%
\pgfpathlineto{\pgfqpoint{5.847111in}{0.637815in}}%
\pgfpathlineto{\pgfqpoint{5.850742in}{0.655342in}}%
\pgfpathlineto{\pgfqpoint{5.854373in}{0.656978in}}%
\pgfpathlineto{\pgfqpoint{5.858005in}{0.647401in}}%
\pgfpathlineto{\pgfqpoint{5.861636in}{0.654326in}}%
\pgfpathlineto{\pgfqpoint{5.865267in}{0.601652in}}%
\pgfpathlineto{\pgfqpoint{5.868898in}{0.531437in}}%
\pgfpathlineto{\pgfqpoint{5.872529in}{0.635094in}}%
\pgfpathlineto{\pgfqpoint{5.876160in}{0.669542in}}%
\pgfpathlineto{\pgfqpoint{5.879792in}{0.683575in}}%
\pgfpathlineto{\pgfqpoint{5.883423in}{0.563280in}}%
\pgfpathlineto{\pgfqpoint{5.887054in}{0.577157in}}%
\pgfpathlineto{\pgfqpoint{5.890685in}{0.642616in}}%
\pgfpathlineto{\pgfqpoint{5.894316in}{0.599275in}}%
\pgfpathlineto{\pgfqpoint{5.901579in}{0.657082in}}%
\pgfpathlineto{\pgfqpoint{5.905210in}{0.629972in}}%
\pgfpathlineto{\pgfqpoint{5.908841in}{0.640141in}}%
\pgfpathlineto{\pgfqpoint{5.912472in}{0.624104in}}%
\pgfpathlineto{\pgfqpoint{5.916103in}{0.638024in}}%
\pgfpathlineto{\pgfqpoint{5.919735in}{0.644417in}}%
\pgfpathlineto{\pgfqpoint{5.923366in}{0.631227in}}%
\pgfpathlineto{\pgfqpoint{5.926997in}{0.596794in}}%
\pgfpathlineto{\pgfqpoint{5.930628in}{0.662774in}}%
\pgfpathlineto{\pgfqpoint{5.934259in}{0.593834in}}%
\pgfpathlineto{\pgfqpoint{5.937890in}{0.592561in}}%
\pgfpathlineto{\pgfqpoint{5.941522in}{0.647795in}}%
\pgfpathlineto{\pgfqpoint{5.945153in}{0.667536in}}%
\pgfpathlineto{\pgfqpoint{5.948784in}{0.675861in}}%
\pgfpathlineto{\pgfqpoint{5.952415in}{0.585091in}}%
\pgfpathlineto{\pgfqpoint{5.956046in}{0.608419in}}%
\pgfpathlineto{\pgfqpoint{5.959677in}{0.636986in}}%
\pgfpathlineto{\pgfqpoint{5.963309in}{0.538576in}}%
\pgfpathlineto{\pgfqpoint{5.966940in}{0.645431in}}%
\pgfpathlineto{\pgfqpoint{5.970571in}{0.670938in}}%
\pgfpathlineto{\pgfqpoint{5.977833in}{0.679542in}}%
\pgfpathlineto{\pgfqpoint{5.981465in}{0.627773in}}%
\pgfpathlineto{\pgfqpoint{5.985096in}{0.608256in}}%
\pgfpathlineto{\pgfqpoint{5.988727in}{0.621463in}}%
\pgfpathlineto{\pgfqpoint{5.992358in}{0.611982in}}%
\pgfpathlineto{\pgfqpoint{5.995989in}{0.653120in}}%
\pgfpathlineto{\pgfqpoint{5.999620in}{0.670170in}}%
\pgfpathlineto{\pgfqpoint{6.003252in}{0.668041in}}%
\pgfpathlineto{\pgfqpoint{6.006883in}{0.653258in}}%
\pgfpathlineto{\pgfqpoint{6.010514in}{0.659049in}}%
\pgfpathlineto{\pgfqpoint{6.014145in}{0.616990in}}%
\pgfpathlineto{\pgfqpoint{6.017776in}{0.618055in}}%
\pgfpathlineto{\pgfqpoint{6.021407in}{0.634470in}}%
\pgfpathlineto{\pgfqpoint{6.025039in}{0.638442in}}%
\pgfpathlineto{\pgfqpoint{6.028670in}{0.631490in}}%
\pgfpathlineto{\pgfqpoint{6.032301in}{0.614558in}}%
\pgfpathlineto{\pgfqpoint{6.035932in}{0.576694in}}%
\pgfpathlineto{\pgfqpoint{6.039563in}{0.664490in}}%
\pgfpathlineto{\pgfqpoint{6.043194in}{0.537136in}}%
\pgfpathlineto{\pgfqpoint{6.046826in}{0.616465in}}%
\pgfpathlineto{\pgfqpoint{6.050457in}{0.651391in}}%
\pgfpathlineto{\pgfqpoint{6.054088in}{0.668336in}}%
\pgfpathlineto{\pgfqpoint{6.057719in}{0.672674in}}%
\pgfpathlineto{\pgfqpoint{6.061350in}{0.618079in}}%
\pgfpathlineto{\pgfqpoint{6.064981in}{0.608635in}}%
\pgfpathlineto{\pgfqpoint{6.068613in}{0.616778in}}%
\pgfpathlineto{\pgfqpoint{6.072244in}{0.575285in}}%
\pgfpathlineto{\pgfqpoint{6.075875in}{0.658531in}}%
\pgfpathlineto{\pgfqpoint{6.079506in}{0.677858in}}%
\pgfpathlineto{\pgfqpoint{6.083137in}{0.625098in}}%
\pgfpathlineto{\pgfqpoint{6.086769in}{0.594033in}}%
\pgfpathlineto{\pgfqpoint{6.090400in}{0.643492in}}%
\pgfpathlineto{\pgfqpoint{6.094031in}{0.631667in}}%
\pgfpathlineto{\pgfqpoint{6.097662in}{0.585455in}}%
\pgfpathlineto{\pgfqpoint{6.101293in}{0.581567in}}%
\pgfpathlineto{\pgfqpoint{6.104924in}{0.601031in}}%
\pgfpathlineto{\pgfqpoint{6.108556in}{0.661803in}}%
\pgfpathlineto{\pgfqpoint{6.112187in}{0.547530in}}%
\pgfpathlineto{\pgfqpoint{6.115818in}{0.628887in}}%
\pgfpathlineto{\pgfqpoint{6.119449in}{0.656379in}}%
\pgfpathlineto{\pgfqpoint{6.123080in}{0.660246in}}%
\pgfpathlineto{\pgfqpoint{6.126711in}{0.654472in}}%
\pgfpathlineto{\pgfqpoint{6.130343in}{0.643061in}}%
\pgfpathlineto{\pgfqpoint{6.133974in}{0.615990in}}%
\pgfpathlineto{\pgfqpoint{6.137605in}{0.593739in}}%
\pgfpathlineto{\pgfqpoint{6.141236in}{0.616151in}}%
\pgfpathlineto{\pgfqpoint{6.144867in}{0.663188in}}%
\pgfpathlineto{\pgfqpoint{6.148498in}{0.681284in}}%
\pgfpathlineto{\pgfqpoint{6.152130in}{0.608424in}}%
\pgfpathlineto{\pgfqpoint{6.155761in}{0.509409in}}%
\pgfpathlineto{\pgfqpoint{6.159392in}{0.649400in}}%
\pgfpathlineto{\pgfqpoint{6.163023in}{0.623106in}}%
\pgfpathlineto{\pgfqpoint{6.166654in}{0.569660in}}%
\pgfpathlineto{\pgfqpoint{6.170285in}{0.639227in}}%
\pgfpathlineto{\pgfqpoint{6.173917in}{0.669052in}}%
\pgfpathlineto{\pgfqpoint{6.177548in}{0.691648in}}%
\pgfpathlineto{\pgfqpoint{6.181179in}{0.591116in}}%
\pgfpathlineto{\pgfqpoint{6.184810in}{0.551769in}}%
\pgfpathlineto{\pgfqpoint{6.188441in}{0.648379in}}%
\pgfpathlineto{\pgfqpoint{6.192073in}{0.528191in}}%
\pgfpathlineto{\pgfqpoint{6.195704in}{0.636826in}}%
\pgfpathlineto{\pgfqpoint{6.199335in}{0.663212in}}%
\pgfpathlineto{\pgfqpoint{6.202966in}{0.677207in}}%
\pgfpathlineto{\pgfqpoint{6.206597in}{0.683512in}}%
\pgfpathlineto{\pgfqpoint{6.210228in}{0.544020in}}%
\pgfpathlineto{\pgfqpoint{6.213860in}{0.590395in}}%
\pgfpathlineto{\pgfqpoint{6.217491in}{0.650190in}}%
\pgfpathlineto{\pgfqpoint{6.221122in}{0.599284in}}%
\pgfpathlineto{\pgfqpoint{6.224753in}{0.637358in}}%
\pgfpathlineto{\pgfqpoint{6.228384in}{0.655762in}}%
\pgfpathlineto{\pgfqpoint{6.232015in}{0.658229in}}%
\pgfpathlineto{\pgfqpoint{6.235647in}{0.649620in}}%
\pgfpathlineto{\pgfqpoint{6.239278in}{0.653170in}}%
\pgfpathlineto{\pgfqpoint{6.242909in}{0.602671in}}%
\pgfpathlineto{\pgfqpoint{6.246540in}{0.508638in}}%
\pgfpathlineto{\pgfqpoint{6.250171in}{0.633537in}}%
\pgfpathlineto{\pgfqpoint{6.253802in}{0.669353in}}%
\pgfpathlineto{\pgfqpoint{6.257434in}{0.683712in}}%
\pgfpathlineto{\pgfqpoint{6.261065in}{0.572223in}}%
\pgfpathlineto{\pgfqpoint{6.264696in}{0.572406in}}%
\pgfpathlineto{\pgfqpoint{6.268327in}{0.643307in}}%
\pgfpathlineto{\pgfqpoint{6.271958in}{0.602425in}}%
\pgfpathlineto{\pgfqpoint{6.275590in}{0.625276in}}%
\pgfpathlineto{\pgfqpoint{6.279221in}{0.655200in}}%
\pgfpathlineto{\pgfqpoint{6.282852in}{0.629443in}}%
\pgfpathlineto{\pgfqpoint{6.286483in}{0.641787in}}%
\pgfpathlineto{\pgfqpoint{6.290114in}{0.622232in}}%
\pgfpathlineto{\pgfqpoint{6.293745in}{0.637942in}}%
\pgfpathlineto{\pgfqpoint{6.297377in}{0.645787in}}%
\pgfpathlineto{\pgfqpoint{6.301008in}{0.634030in}}%
\pgfpathlineto{\pgfqpoint{6.304639in}{0.602865in}}%
\pgfpathlineto{\pgfqpoint{6.308270in}{0.662235in}}%
\pgfpathlineto{\pgfqpoint{6.311901in}{0.596197in}}%
\pgfpathlineto{\pgfqpoint{6.315532in}{0.588247in}}%
\pgfpathlineto{\pgfqpoint{6.319164in}{0.646784in}}%
\pgfpathlineto{\pgfqpoint{6.322795in}{0.667650in}}%
\pgfpathlineto{\pgfqpoint{6.326426in}{0.676739in}}%
\pgfpathlineto{\pgfqpoint{6.330057in}{0.577133in}}%
\pgfpathlineto{\pgfqpoint{6.337319in}{0.638258in}}%
\pgfpathlineto{\pgfqpoint{6.340951in}{0.552769in}}%
\pgfpathlineto{\pgfqpoint{6.344582in}{0.643771in}}%
\pgfpathlineto{\pgfqpoint{6.348213in}{0.669889in}}%
\pgfpathlineto{\pgfqpoint{6.355475in}{0.680755in}}%
\pgfpathlineto{\pgfqpoint{6.359106in}{0.624011in}}%
\pgfpathlineto{\pgfqpoint{6.362738in}{0.607323in}}%
\pgfpathlineto{\pgfqpoint{6.366369in}{0.631960in}}%
\pgfpathlineto{\pgfqpoint{6.370000in}{0.586570in}}%
\pgfpathlineto{\pgfqpoint{6.370000in}{0.586570in}}%
\pgfusepath{stroke}%
\end{pgfscope}%
\begin{pgfscope}%
\pgfpathrectangle{\pgfqpoint{4.554414in}{0.320625in}}{\pgfqpoint{1.815586in}{1.164375in}}%
\pgfusepath{clip}%
\pgfsetrectcap%
\pgfsetroundjoin%
\pgfsetlinewidth{1.154312pt}%
\definecolor{currentstroke}{rgb}{0.000000,0.529412,0.423529}%
\pgfsetstrokecolor{currentstroke}%
\pgfsetdash{}{0pt}%
\pgfpathmoveto{\pgfqpoint{4.554414in}{1.456638in}}%
\pgfpathlineto{\pgfqpoint{4.558045in}{1.444927in}}%
\pgfpathlineto{\pgfqpoint{4.565307in}{1.429462in}}%
\pgfpathlineto{\pgfqpoint{4.583463in}{1.392634in}}%
\pgfpathlineto{\pgfqpoint{4.608881in}{1.337333in}}%
\pgfpathlineto{\pgfqpoint{4.652455in}{1.242404in}}%
\pgfpathlineto{\pgfqpoint{4.685136in}{1.173907in}}%
\pgfpathlineto{\pgfqpoint{4.721448in}{1.100377in}}%
\pgfpathlineto{\pgfqpoint{4.761390in}{1.021928in}}%
\pgfpathlineto{\pgfqpoint{4.808596in}{0.931671in}}%
\pgfpathlineto{\pgfqpoint{4.863063in}{0.830028in}}%
\pgfpathlineto{\pgfqpoint{4.913900in}{0.737255in}}%
\pgfpathlineto{\pgfqpoint{4.932056in}{0.705596in}}%
\pgfpathlineto{\pgfqpoint{4.935687in}{0.703861in}}%
\pgfpathlineto{\pgfqpoint{4.939318in}{0.691862in}}%
\pgfpathlineto{\pgfqpoint{4.942949in}{0.688967in}}%
\pgfpathlineto{\pgfqpoint{4.946580in}{0.688175in}}%
\pgfpathlineto{\pgfqpoint{4.950211in}{0.674013in}}%
\pgfpathlineto{\pgfqpoint{4.953843in}{0.680674in}}%
\pgfpathlineto{\pgfqpoint{4.957474in}{0.680461in}}%
\pgfpathlineto{\pgfqpoint{4.961105in}{0.662327in}}%
\pgfpathlineto{\pgfqpoint{4.964736in}{0.660818in}}%
\pgfpathlineto{\pgfqpoint{4.968367in}{0.661210in}}%
\pgfpathlineto{\pgfqpoint{4.971998in}{0.656161in}}%
\pgfpathlineto{\pgfqpoint{4.975630in}{0.659569in}}%
\pgfpathlineto{\pgfqpoint{4.979261in}{0.652227in}}%
\pgfpathlineto{\pgfqpoint{4.982892in}{0.651595in}}%
\pgfpathlineto{\pgfqpoint{4.986523in}{0.655319in}}%
\pgfpathlineto{\pgfqpoint{4.990154in}{0.652350in}}%
\pgfpathlineto{\pgfqpoint{4.993785in}{0.633267in}}%
\pgfpathlineto{\pgfqpoint{4.997417in}{0.639254in}}%
\pgfpathlineto{\pgfqpoint{5.001048in}{0.662145in}}%
\pgfpathlineto{\pgfqpoint{5.004679in}{0.666450in}}%
\pgfpathlineto{\pgfqpoint{5.008310in}{0.651756in}}%
\pgfpathlineto{\pgfqpoint{5.011941in}{0.646047in}}%
\pgfpathlineto{\pgfqpoint{5.015573in}{0.655122in}}%
\pgfpathlineto{\pgfqpoint{5.019204in}{0.643891in}}%
\pgfpathlineto{\pgfqpoint{5.022835in}{0.664458in}}%
\pgfpathlineto{\pgfqpoint{5.026466in}{0.660290in}}%
\pgfpathlineto{\pgfqpoint{5.033728in}{0.667066in}}%
\pgfpathlineto{\pgfqpoint{5.037360in}{0.630858in}}%
\pgfpathlineto{\pgfqpoint{5.040991in}{0.655030in}}%
\pgfpathlineto{\pgfqpoint{5.044622in}{0.673178in}}%
\pgfpathlineto{\pgfqpoint{5.048253in}{0.658845in}}%
\pgfpathlineto{\pgfqpoint{5.051884in}{0.660252in}}%
\pgfpathlineto{\pgfqpoint{5.055515in}{0.634570in}}%
\pgfpathlineto{\pgfqpoint{5.059147in}{0.658634in}}%
\pgfpathlineto{\pgfqpoint{5.062778in}{0.673113in}}%
\pgfpathlineto{\pgfqpoint{5.066409in}{0.644065in}}%
\pgfpathlineto{\pgfqpoint{5.070040in}{0.661359in}}%
\pgfpathlineto{\pgfqpoint{5.073671in}{0.671406in}}%
\pgfpathlineto{\pgfqpoint{5.077302in}{0.629993in}}%
\pgfpathlineto{\pgfqpoint{5.080934in}{0.634391in}}%
\pgfpathlineto{\pgfqpoint{5.084565in}{0.649600in}}%
\pgfpathlineto{\pgfqpoint{5.088196in}{0.651626in}}%
\pgfpathlineto{\pgfqpoint{5.091827in}{0.666901in}}%
\pgfpathlineto{\pgfqpoint{5.099090in}{0.644666in}}%
\pgfpathlineto{\pgfqpoint{5.102721in}{0.630030in}}%
\pgfpathlineto{\pgfqpoint{5.106352in}{0.639558in}}%
\pgfpathlineto{\pgfqpoint{5.109983in}{0.653715in}}%
\pgfpathlineto{\pgfqpoint{5.113614in}{0.671146in}}%
\pgfpathlineto{\pgfqpoint{5.117245in}{0.659547in}}%
\pgfpathlineto{\pgfqpoint{5.120877in}{0.662729in}}%
\pgfpathlineto{\pgfqpoint{5.124508in}{0.667614in}}%
\pgfpathlineto{\pgfqpoint{5.128139in}{0.653299in}}%
\pgfpathlineto{\pgfqpoint{5.131770in}{0.632443in}}%
\pgfpathlineto{\pgfqpoint{5.135401in}{0.626580in}}%
\pgfpathlineto{\pgfqpoint{5.139032in}{0.657886in}}%
\pgfpathlineto{\pgfqpoint{5.142664in}{0.660254in}}%
\pgfpathlineto{\pgfqpoint{5.146295in}{0.643675in}}%
\pgfpathlineto{\pgfqpoint{5.149926in}{0.638833in}}%
\pgfpathlineto{\pgfqpoint{5.153557in}{0.643236in}}%
\pgfpathlineto{\pgfqpoint{5.157188in}{0.678051in}}%
\pgfpathlineto{\pgfqpoint{5.160819in}{0.660239in}}%
\pgfpathlineto{\pgfqpoint{5.164451in}{0.616213in}}%
\pgfpathlineto{\pgfqpoint{5.168082in}{0.618952in}}%
\pgfpathlineto{\pgfqpoint{5.171713in}{0.636688in}}%
\pgfpathlineto{\pgfqpoint{5.175344in}{0.657415in}}%
\pgfpathlineto{\pgfqpoint{5.178975in}{0.667801in}}%
\pgfpathlineto{\pgfqpoint{5.182606in}{0.631526in}}%
\pgfpathlineto{\pgfqpoint{5.186238in}{0.643082in}}%
\pgfpathlineto{\pgfqpoint{5.189869in}{0.651060in}}%
\pgfpathlineto{\pgfqpoint{5.193500in}{0.644552in}}%
\pgfpathlineto{\pgfqpoint{5.197131in}{0.648174in}}%
\pgfpathlineto{\pgfqpoint{5.200762in}{0.661756in}}%
\pgfpathlineto{\pgfqpoint{5.204394in}{0.641776in}}%
\pgfpathlineto{\pgfqpoint{5.208025in}{0.658112in}}%
\pgfpathlineto{\pgfqpoint{5.211656in}{0.645498in}}%
\pgfpathlineto{\pgfqpoint{5.215287in}{0.647271in}}%
\pgfpathlineto{\pgfqpoint{5.218918in}{0.661215in}}%
\pgfpathlineto{\pgfqpoint{5.222549in}{0.664302in}}%
\pgfpathlineto{\pgfqpoint{5.226181in}{0.654120in}}%
\pgfpathlineto{\pgfqpoint{5.229812in}{0.659041in}}%
\pgfpathlineto{\pgfqpoint{5.233443in}{0.629811in}}%
\pgfpathlineto{\pgfqpoint{5.237074in}{0.627343in}}%
\pgfpathlineto{\pgfqpoint{5.240705in}{0.651043in}}%
\pgfpathlineto{\pgfqpoint{5.244336in}{0.664430in}}%
\pgfpathlineto{\pgfqpoint{5.247968in}{0.673305in}}%
\pgfpathlineto{\pgfqpoint{5.251599in}{0.647269in}}%
\pgfpathlineto{\pgfqpoint{5.255230in}{0.647995in}}%
\pgfpathlineto{\pgfqpoint{5.258861in}{0.632191in}}%
\pgfpathlineto{\pgfqpoint{5.262492in}{0.608789in}}%
\pgfpathlineto{\pgfqpoint{5.266123in}{0.650801in}}%
\pgfpathlineto{\pgfqpoint{5.269755in}{0.661565in}}%
\pgfpathlineto{\pgfqpoint{5.273386in}{0.636974in}}%
\pgfpathlineto{\pgfqpoint{5.277017in}{0.661506in}}%
\pgfpathlineto{\pgfqpoint{5.280648in}{0.635527in}}%
\pgfpathlineto{\pgfqpoint{5.284279in}{0.645143in}}%
\pgfpathlineto{\pgfqpoint{5.287910in}{0.630662in}}%
\pgfpathlineto{\pgfqpoint{5.291542in}{0.638299in}}%
\pgfpathlineto{\pgfqpoint{5.295173in}{0.671598in}}%
\pgfpathlineto{\pgfqpoint{5.298804in}{0.662560in}}%
\pgfpathlineto{\pgfqpoint{5.302435in}{0.637431in}}%
\pgfpathlineto{\pgfqpoint{5.306066in}{0.642135in}}%
\pgfpathlineto{\pgfqpoint{5.309698in}{0.661150in}}%
\pgfpathlineto{\pgfqpoint{5.313329in}{0.660398in}}%
\pgfpathlineto{\pgfqpoint{5.316960in}{0.669020in}}%
\pgfpathlineto{\pgfqpoint{5.320591in}{0.652310in}}%
\pgfpathlineto{\pgfqpoint{5.324222in}{0.656408in}}%
\pgfpathlineto{\pgfqpoint{5.327853in}{0.668107in}}%
\pgfpathlineto{\pgfqpoint{5.331485in}{0.621980in}}%
\pgfpathlineto{\pgfqpoint{5.335116in}{0.648147in}}%
\pgfpathlineto{\pgfqpoint{5.338747in}{0.651718in}}%
\pgfpathlineto{\pgfqpoint{5.342378in}{0.630279in}}%
\pgfpathlineto{\pgfqpoint{5.346009in}{0.637095in}}%
\pgfpathlineto{\pgfqpoint{5.349640in}{0.649396in}}%
\pgfpathlineto{\pgfqpoint{5.353272in}{0.646936in}}%
\pgfpathlineto{\pgfqpoint{5.356903in}{0.658951in}}%
\pgfpathlineto{\pgfqpoint{5.360534in}{0.635770in}}%
\pgfpathlineto{\pgfqpoint{5.364165in}{0.638510in}}%
\pgfpathlineto{\pgfqpoint{5.367796in}{0.647463in}}%
\pgfpathlineto{\pgfqpoint{5.371427in}{0.637275in}}%
\pgfpathlineto{\pgfqpoint{5.375059in}{0.640180in}}%
\pgfpathlineto{\pgfqpoint{5.378690in}{0.653070in}}%
\pgfpathlineto{\pgfqpoint{5.382321in}{0.651272in}}%
\pgfpathlineto{\pgfqpoint{5.385952in}{0.659550in}}%
\pgfpathlineto{\pgfqpoint{5.389583in}{0.640582in}}%
\pgfpathlineto{\pgfqpoint{5.393215in}{0.661173in}}%
\pgfpathlineto{\pgfqpoint{5.396846in}{0.632478in}}%
\pgfpathlineto{\pgfqpoint{5.400477in}{0.633158in}}%
\pgfpathlineto{\pgfqpoint{5.404108in}{0.636260in}}%
\pgfpathlineto{\pgfqpoint{5.407739in}{0.644495in}}%
\pgfpathlineto{\pgfqpoint{5.411370in}{0.657852in}}%
\pgfpathlineto{\pgfqpoint{5.415002in}{0.677451in}}%
\pgfpathlineto{\pgfqpoint{5.418633in}{0.644167in}}%
\pgfpathlineto{\pgfqpoint{5.422264in}{0.667571in}}%
\pgfpathlineto{\pgfqpoint{5.425895in}{0.635595in}}%
\pgfpathlineto{\pgfqpoint{5.429526in}{0.631736in}}%
\pgfpathlineto{\pgfqpoint{5.433157in}{0.656420in}}%
\pgfpathlineto{\pgfqpoint{5.436789in}{0.658425in}}%
\pgfpathlineto{\pgfqpoint{5.440420in}{0.658040in}}%
\pgfpathlineto{\pgfqpoint{5.444051in}{0.675513in}}%
\pgfpathlineto{\pgfqpoint{5.447682in}{0.659317in}}%
\pgfpathlineto{\pgfqpoint{5.451313in}{0.655301in}}%
\pgfpathlineto{\pgfqpoint{5.454944in}{0.656799in}}%
\pgfpathlineto{\pgfqpoint{5.458576in}{0.641082in}}%
\pgfpathlineto{\pgfqpoint{5.462207in}{0.655570in}}%
\pgfpathlineto{\pgfqpoint{5.465838in}{0.652199in}}%
\pgfpathlineto{\pgfqpoint{5.469469in}{0.653848in}}%
\pgfpathlineto{\pgfqpoint{5.473100in}{0.658488in}}%
\pgfpathlineto{\pgfqpoint{5.476731in}{0.644264in}}%
\pgfpathlineto{\pgfqpoint{5.480363in}{0.654105in}}%
\pgfpathlineto{\pgfqpoint{5.483994in}{0.658713in}}%
\pgfpathlineto{\pgfqpoint{5.487625in}{0.625333in}}%
\pgfpathlineto{\pgfqpoint{5.491256in}{0.643729in}}%
\pgfpathlineto{\pgfqpoint{5.494887in}{0.675260in}}%
\pgfpathlineto{\pgfqpoint{5.498519in}{0.660471in}}%
\pgfpathlineto{\pgfqpoint{5.502150in}{0.655237in}}%
\pgfpathlineto{\pgfqpoint{5.505781in}{0.618910in}}%
\pgfpathlineto{\pgfqpoint{5.513043in}{0.664094in}}%
\pgfpathlineto{\pgfqpoint{5.516674in}{0.672221in}}%
\pgfpathlineto{\pgfqpoint{5.520306in}{0.641459in}}%
\pgfpathlineto{\pgfqpoint{5.523937in}{0.658135in}}%
\pgfpathlineto{\pgfqpoint{5.527568in}{0.663001in}}%
\pgfpathlineto{\pgfqpoint{5.531199in}{0.640625in}}%
\pgfpathlineto{\pgfqpoint{5.534830in}{0.645370in}}%
\pgfpathlineto{\pgfqpoint{5.538461in}{0.657552in}}%
\pgfpathlineto{\pgfqpoint{5.542093in}{0.663580in}}%
\pgfpathlineto{\pgfqpoint{5.545724in}{0.611098in}}%
\pgfpathlineto{\pgfqpoint{5.549355in}{0.651002in}}%
\pgfpathlineto{\pgfqpoint{5.552986in}{0.652560in}}%
\pgfpathlineto{\pgfqpoint{5.556617in}{0.639904in}}%
\pgfpathlineto{\pgfqpoint{5.560248in}{0.641482in}}%
\pgfpathlineto{\pgfqpoint{5.563880in}{0.656168in}}%
\pgfpathlineto{\pgfqpoint{5.567511in}{0.646573in}}%
\pgfpathlineto{\pgfqpoint{5.571142in}{0.658699in}}%
\pgfpathlineto{\pgfqpoint{5.574773in}{0.634458in}}%
\pgfpathlineto{\pgfqpoint{5.578404in}{0.657432in}}%
\pgfpathlineto{\pgfqpoint{5.582035in}{0.648078in}}%
\pgfpathlineto{\pgfqpoint{5.585667in}{0.644227in}}%
\pgfpathlineto{\pgfqpoint{5.589298in}{0.649256in}}%
\pgfpathlineto{\pgfqpoint{5.592929in}{0.660618in}}%
\pgfpathlineto{\pgfqpoint{5.596560in}{0.653409in}}%
\pgfpathlineto{\pgfqpoint{5.600191in}{0.667630in}}%
\pgfpathlineto{\pgfqpoint{5.603823in}{0.638155in}}%
\pgfpathlineto{\pgfqpoint{5.607454in}{0.650975in}}%
\pgfpathlineto{\pgfqpoint{5.611085in}{0.651424in}}%
\pgfpathlineto{\pgfqpoint{5.614716in}{0.622581in}}%
\pgfpathlineto{\pgfqpoint{5.618347in}{0.656543in}}%
\pgfpathlineto{\pgfqpoint{5.621978in}{0.673059in}}%
\pgfpathlineto{\pgfqpoint{5.625610in}{0.655020in}}%
\pgfpathlineto{\pgfqpoint{5.629241in}{0.664535in}}%
\pgfpathlineto{\pgfqpoint{5.632872in}{0.667850in}}%
\pgfpathlineto{\pgfqpoint{5.636503in}{0.622796in}}%
\pgfpathlineto{\pgfqpoint{5.640134in}{0.624146in}}%
\pgfpathlineto{\pgfqpoint{5.643765in}{0.641690in}}%
\pgfpathlineto{\pgfqpoint{5.647397in}{0.663121in}}%
\pgfpathlineto{\pgfqpoint{5.651028in}{0.673286in}}%
\pgfpathlineto{\pgfqpoint{5.658290in}{0.655682in}}%
\pgfpathlineto{\pgfqpoint{5.661921in}{0.641814in}}%
\pgfpathlineto{\pgfqpoint{5.665552in}{0.655739in}}%
\pgfpathlineto{\pgfqpoint{5.672815in}{0.627883in}}%
\pgfpathlineto{\pgfqpoint{5.676446in}{0.656684in}}%
\pgfpathlineto{\pgfqpoint{5.680077in}{0.672479in}}%
\pgfpathlineto{\pgfqpoint{5.683708in}{0.623541in}}%
\pgfpathlineto{\pgfqpoint{5.687340in}{0.655024in}}%
\pgfpathlineto{\pgfqpoint{5.690971in}{0.650722in}}%
\pgfpathlineto{\pgfqpoint{5.694602in}{0.657572in}}%
\pgfpathlineto{\pgfqpoint{5.698233in}{0.675709in}}%
\pgfpathlineto{\pgfqpoint{5.701864in}{0.661344in}}%
\pgfpathlineto{\pgfqpoint{5.705495in}{0.629011in}}%
\pgfpathlineto{\pgfqpoint{5.709127in}{0.642209in}}%
\pgfpathlineto{\pgfqpoint{5.716389in}{0.658833in}}%
\pgfpathlineto{\pgfqpoint{5.720020in}{0.662275in}}%
\pgfpathlineto{\pgfqpoint{5.723651in}{0.648071in}}%
\pgfpathlineto{\pgfqpoint{5.727282in}{0.662520in}}%
\pgfpathlineto{\pgfqpoint{5.730914in}{0.643574in}}%
\pgfpathlineto{\pgfqpoint{5.734545in}{0.652386in}}%
\pgfpathlineto{\pgfqpoint{5.738176in}{0.663212in}}%
\pgfpathlineto{\pgfqpoint{5.741807in}{0.632874in}}%
\pgfpathlineto{\pgfqpoint{5.745438in}{0.633250in}}%
\pgfpathlineto{\pgfqpoint{5.749069in}{0.658113in}}%
\pgfpathlineto{\pgfqpoint{5.752701in}{0.640730in}}%
\pgfpathlineto{\pgfqpoint{5.756332in}{0.654536in}}%
\pgfpathlineto{\pgfqpoint{5.759963in}{0.664489in}}%
\pgfpathlineto{\pgfqpoint{5.763594in}{0.646462in}}%
\pgfpathlineto{\pgfqpoint{5.767225in}{0.659763in}}%
\pgfpathlineto{\pgfqpoint{5.770856in}{0.652899in}}%
\pgfpathlineto{\pgfqpoint{5.774488in}{0.642672in}}%
\pgfpathlineto{\pgfqpoint{5.778119in}{0.661066in}}%
\pgfpathlineto{\pgfqpoint{5.781750in}{0.661468in}}%
\pgfpathlineto{\pgfqpoint{5.785381in}{0.648903in}}%
\pgfpathlineto{\pgfqpoint{5.789012in}{0.643470in}}%
\pgfpathlineto{\pgfqpoint{5.792644in}{0.655547in}}%
\pgfpathlineto{\pgfqpoint{5.796275in}{0.670315in}}%
\pgfpathlineto{\pgfqpoint{5.799906in}{0.664021in}}%
\pgfpathlineto{\pgfqpoint{5.803537in}{0.658934in}}%
\pgfpathlineto{\pgfqpoint{5.807168in}{0.660831in}}%
\pgfpathlineto{\pgfqpoint{5.810799in}{0.634449in}}%
\pgfpathlineto{\pgfqpoint{5.814431in}{0.639213in}}%
\pgfpathlineto{\pgfqpoint{5.818062in}{0.661254in}}%
\pgfpathlineto{\pgfqpoint{5.821693in}{0.660537in}}%
\pgfpathlineto{\pgfqpoint{5.825324in}{0.671717in}}%
\pgfpathlineto{\pgfqpoint{5.828955in}{0.657268in}}%
\pgfpathlineto{\pgfqpoint{5.832586in}{0.626691in}}%
\pgfpathlineto{\pgfqpoint{5.836218in}{0.660993in}}%
\pgfpathlineto{\pgfqpoint{5.839849in}{0.659466in}}%
\pgfpathlineto{\pgfqpoint{5.843480in}{0.646140in}}%
\pgfpathlineto{\pgfqpoint{5.850742in}{0.662082in}}%
\pgfpathlineto{\pgfqpoint{5.854373in}{0.673575in}}%
\pgfpathlineto{\pgfqpoint{5.858005in}{0.632549in}}%
\pgfpathlineto{\pgfqpoint{5.861636in}{0.643375in}}%
\pgfpathlineto{\pgfqpoint{5.865267in}{0.661650in}}%
\pgfpathlineto{\pgfqpoint{5.868898in}{0.644928in}}%
\pgfpathlineto{\pgfqpoint{5.872529in}{0.659419in}}%
\pgfpathlineto{\pgfqpoint{5.876160in}{0.671130in}}%
\pgfpathlineto{\pgfqpoint{5.883423in}{0.645609in}}%
\pgfpathlineto{\pgfqpoint{5.887054in}{0.674744in}}%
\pgfpathlineto{\pgfqpoint{5.890685in}{0.631885in}}%
\pgfpathlineto{\pgfqpoint{5.894316in}{0.627133in}}%
\pgfpathlineto{\pgfqpoint{5.897948in}{0.636140in}}%
\pgfpathlineto{\pgfqpoint{5.901579in}{0.642272in}}%
\pgfpathlineto{\pgfqpoint{5.905210in}{0.662366in}}%
\pgfpathlineto{\pgfqpoint{5.908841in}{0.670382in}}%
\pgfpathlineto{\pgfqpoint{5.912472in}{0.639382in}}%
\pgfpathlineto{\pgfqpoint{5.916103in}{0.634667in}}%
\pgfpathlineto{\pgfqpoint{5.919735in}{0.654176in}}%
\pgfpathlineto{\pgfqpoint{5.923366in}{0.655746in}}%
\pgfpathlineto{\pgfqpoint{5.926997in}{0.636922in}}%
\pgfpathlineto{\pgfqpoint{5.930628in}{0.656406in}}%
\pgfpathlineto{\pgfqpoint{5.934259in}{0.645796in}}%
\pgfpathlineto{\pgfqpoint{5.937890in}{0.620833in}}%
\pgfpathlineto{\pgfqpoint{5.941522in}{0.647592in}}%
\pgfpathlineto{\pgfqpoint{5.945153in}{0.659743in}}%
\pgfpathlineto{\pgfqpoint{5.948784in}{0.650562in}}%
\pgfpathlineto{\pgfqpoint{5.952415in}{0.643583in}}%
\pgfpathlineto{\pgfqpoint{5.956046in}{0.639558in}}%
\pgfpathlineto{\pgfqpoint{5.959677in}{0.644569in}}%
\pgfpathlineto{\pgfqpoint{5.963309in}{0.635637in}}%
\pgfpathlineto{\pgfqpoint{5.966940in}{0.628679in}}%
\pgfpathlineto{\pgfqpoint{5.974202in}{0.674678in}}%
\pgfpathlineto{\pgfqpoint{5.977833in}{0.669998in}}%
\pgfpathlineto{\pgfqpoint{5.981465in}{0.652452in}}%
\pgfpathlineto{\pgfqpoint{5.985096in}{0.628087in}}%
\pgfpathlineto{\pgfqpoint{5.988727in}{0.632932in}}%
\pgfpathlineto{\pgfqpoint{5.992358in}{0.630186in}}%
\pgfpathlineto{\pgfqpoint{5.995989in}{0.630444in}}%
\pgfpathlineto{\pgfqpoint{5.999620in}{0.664145in}}%
\pgfpathlineto{\pgfqpoint{6.003252in}{0.682513in}}%
\pgfpathlineto{\pgfqpoint{6.006883in}{0.650847in}}%
\pgfpathlineto{\pgfqpoint{6.010514in}{0.650295in}}%
\pgfpathlineto{\pgfqpoint{6.014145in}{0.650431in}}%
\pgfpathlineto{\pgfqpoint{6.017776in}{0.627114in}}%
\pgfpathlineto{\pgfqpoint{6.021407in}{0.627823in}}%
\pgfpathlineto{\pgfqpoint{6.025039in}{0.662546in}}%
\pgfpathlineto{\pgfqpoint{6.028670in}{0.658958in}}%
\pgfpathlineto{\pgfqpoint{6.032301in}{0.661665in}}%
\pgfpathlineto{\pgfqpoint{6.035932in}{0.634900in}}%
\pgfpathlineto{\pgfqpoint{6.039563in}{0.646388in}}%
\pgfpathlineto{\pgfqpoint{6.043194in}{0.643803in}}%
\pgfpathlineto{\pgfqpoint{6.050457in}{0.667802in}}%
\pgfpathlineto{\pgfqpoint{6.054088in}{0.655518in}}%
\pgfpathlineto{\pgfqpoint{6.057719in}{0.651332in}}%
\pgfpathlineto{\pgfqpoint{6.061350in}{0.643643in}}%
\pgfpathlineto{\pgfqpoint{6.064981in}{0.629001in}}%
\pgfpathlineto{\pgfqpoint{6.068613in}{0.661136in}}%
\pgfpathlineto{\pgfqpoint{6.072244in}{0.667145in}}%
\pgfpathlineto{\pgfqpoint{6.075875in}{0.650144in}}%
\pgfpathlineto{\pgfqpoint{6.079506in}{0.663745in}}%
\pgfpathlineto{\pgfqpoint{6.083137in}{0.627013in}}%
\pgfpathlineto{\pgfqpoint{6.086769in}{0.629186in}}%
\pgfpathlineto{\pgfqpoint{6.090400in}{0.652899in}}%
\pgfpathlineto{\pgfqpoint{6.094031in}{0.652421in}}%
\pgfpathlineto{\pgfqpoint{6.097662in}{0.647954in}}%
\pgfpathlineto{\pgfqpoint{6.101293in}{0.654514in}}%
\pgfpathlineto{\pgfqpoint{6.104924in}{0.617668in}}%
\pgfpathlineto{\pgfqpoint{6.108556in}{0.664870in}}%
\pgfpathlineto{\pgfqpoint{6.112187in}{0.642053in}}%
\pgfpathlineto{\pgfqpoint{6.115818in}{0.634281in}}%
\pgfpathlineto{\pgfqpoint{6.119449in}{0.653599in}}%
\pgfpathlineto{\pgfqpoint{6.123080in}{0.630304in}}%
\pgfpathlineto{\pgfqpoint{6.126711in}{0.631773in}}%
\pgfpathlineto{\pgfqpoint{6.130343in}{0.662496in}}%
\pgfpathlineto{\pgfqpoint{6.133974in}{0.608755in}}%
\pgfpathlineto{\pgfqpoint{6.137605in}{0.628890in}}%
\pgfpathlineto{\pgfqpoint{6.141236in}{0.653506in}}%
\pgfpathlineto{\pgfqpoint{6.144867in}{0.656814in}}%
\pgfpathlineto{\pgfqpoint{6.148498in}{0.664408in}}%
\pgfpathlineto{\pgfqpoint{6.152130in}{0.645049in}}%
\pgfpathlineto{\pgfqpoint{6.155761in}{0.616922in}}%
\pgfpathlineto{\pgfqpoint{6.159392in}{0.659863in}}%
\pgfpathlineto{\pgfqpoint{6.163023in}{0.661000in}}%
\pgfpathlineto{\pgfqpoint{6.166654in}{0.649680in}}%
\pgfpathlineto{\pgfqpoint{6.170285in}{0.655096in}}%
\pgfpathlineto{\pgfqpoint{6.173917in}{0.667699in}}%
\pgfpathlineto{\pgfqpoint{6.177548in}{0.677514in}}%
\pgfpathlineto{\pgfqpoint{6.181179in}{0.615559in}}%
\pgfpathlineto{\pgfqpoint{6.184810in}{0.638328in}}%
\pgfpathlineto{\pgfqpoint{6.188441in}{0.669241in}}%
\pgfpathlineto{\pgfqpoint{6.192073in}{0.648108in}}%
\pgfpathlineto{\pgfqpoint{6.195704in}{0.646168in}}%
\pgfpathlineto{\pgfqpoint{6.199335in}{0.666243in}}%
\pgfpathlineto{\pgfqpoint{6.202966in}{0.662286in}}%
\pgfpathlineto{\pgfqpoint{6.206597in}{0.667997in}}%
\pgfpathlineto{\pgfqpoint{6.210228in}{0.649276in}}%
\pgfpathlineto{\pgfqpoint{6.213860in}{0.651591in}}%
\pgfpathlineto{\pgfqpoint{6.221122in}{0.623571in}}%
\pgfpathlineto{\pgfqpoint{6.224753in}{0.656180in}}%
\pgfpathlineto{\pgfqpoint{6.228384in}{0.651150in}}%
\pgfpathlineto{\pgfqpoint{6.232015in}{0.648130in}}%
\pgfpathlineto{\pgfqpoint{6.235647in}{0.654944in}}%
\pgfpathlineto{\pgfqpoint{6.239278in}{0.646969in}}%
\pgfpathlineto{\pgfqpoint{6.242909in}{0.648056in}}%
\pgfpathlineto{\pgfqpoint{6.246540in}{0.668529in}}%
\pgfpathlineto{\pgfqpoint{6.250171in}{0.641310in}}%
\pgfpathlineto{\pgfqpoint{6.257434in}{0.681825in}}%
\pgfpathlineto{\pgfqpoint{6.261065in}{0.605819in}}%
\pgfpathlineto{\pgfqpoint{6.264696in}{0.619287in}}%
\pgfpathlineto{\pgfqpoint{6.268327in}{0.648238in}}%
\pgfpathlineto{\pgfqpoint{6.271958in}{0.638165in}}%
\pgfpathlineto{\pgfqpoint{6.275590in}{0.652021in}}%
\pgfpathlineto{\pgfqpoint{6.279221in}{0.671253in}}%
\pgfpathlineto{\pgfqpoint{6.282852in}{0.636792in}}%
\pgfpathlineto{\pgfqpoint{6.286483in}{0.649985in}}%
\pgfpathlineto{\pgfqpoint{6.290114in}{0.640123in}}%
\pgfpathlineto{\pgfqpoint{6.293745in}{0.651946in}}%
\pgfpathlineto{\pgfqpoint{6.297377in}{0.658587in}}%
\pgfpathlineto{\pgfqpoint{6.301008in}{0.656366in}}%
\pgfpathlineto{\pgfqpoint{6.304639in}{0.658601in}}%
\pgfpathlineto{\pgfqpoint{6.308270in}{0.657947in}}%
\pgfpathlineto{\pgfqpoint{6.311901in}{0.632234in}}%
\pgfpathlineto{\pgfqpoint{6.315532in}{0.654410in}}%
\pgfpathlineto{\pgfqpoint{6.319164in}{0.654927in}}%
\pgfpathlineto{\pgfqpoint{6.322795in}{0.652230in}}%
\pgfpathlineto{\pgfqpoint{6.326426in}{0.677928in}}%
\pgfpathlineto{\pgfqpoint{6.330057in}{0.636467in}}%
\pgfpathlineto{\pgfqpoint{6.333688in}{0.622454in}}%
\pgfpathlineto{\pgfqpoint{6.337319in}{0.654555in}}%
\pgfpathlineto{\pgfqpoint{6.340951in}{0.659913in}}%
\pgfpathlineto{\pgfqpoint{6.344582in}{0.642387in}}%
\pgfpathlineto{\pgfqpoint{6.351844in}{0.650098in}}%
\pgfpathlineto{\pgfqpoint{6.355475in}{0.672386in}}%
\pgfpathlineto{\pgfqpoint{6.359106in}{0.665448in}}%
\pgfpathlineto{\pgfqpoint{6.362738in}{0.663146in}}%
\pgfpathlineto{\pgfqpoint{6.366369in}{0.630164in}}%
\pgfpathlineto{\pgfqpoint{6.370000in}{0.613527in}}%
\pgfpathlineto{\pgfqpoint{6.370000in}{0.613527in}}%
\pgfusepath{stroke}%
\end{pgfscope}%
\begin{pgfscope}%
\pgfpathrectangle{\pgfqpoint{4.554414in}{0.320625in}}{\pgfqpoint{1.815586in}{1.164375in}}%
\pgfusepath{clip}%
\pgfsetrectcap%
\pgfsetroundjoin%
\pgfsetlinewidth{0.602250pt}%
\definecolor{currentstroke}{rgb}{0.478431,0.243137,0.615686}%
\pgfsetstrokecolor{currentstroke}%
\pgfsetstrokeopacity{0.120000}%
\pgfsetdash{}{0pt}%
\pgfpathmoveto{\pgfqpoint{4.554414in}{1.454313in}}%
\pgfpathlineto{\pgfqpoint{4.558045in}{1.443912in}}%
\pgfpathlineto{\pgfqpoint{4.579832in}{1.399802in}}%
\pgfpathlineto{\pgfqpoint{4.597988in}{1.359656in}}%
\pgfpathlineto{\pgfqpoint{4.681505in}{1.171212in}}%
\pgfpathlineto{\pgfqpoint{4.721448in}{1.084991in}}%
\pgfpathlineto{\pgfqpoint{4.765022in}{0.993674in}}%
\pgfpathlineto{\pgfqpoint{4.812227in}{0.897388in}}%
\pgfpathlineto{\pgfqpoint{4.863063in}{0.796192in}}%
\pgfpathlineto{\pgfqpoint{4.910269in}{0.704312in}}%
\pgfpathlineto{\pgfqpoint{4.913900in}{0.699948in}}%
\pgfpathlineto{\pgfqpoint{4.917531in}{0.699042in}}%
\pgfpathlineto{\pgfqpoint{4.924793in}{0.677796in}}%
\pgfpathlineto{\pgfqpoint{4.928424in}{0.683303in}}%
\pgfpathlineto{\pgfqpoint{4.932056in}{0.666303in}}%
\pgfpathlineto{\pgfqpoint{4.935687in}{0.657600in}}%
\pgfpathlineto{\pgfqpoint{4.939318in}{0.659781in}}%
\pgfpathlineto{\pgfqpoint{4.942949in}{0.661006in}}%
\pgfpathlineto{\pgfqpoint{4.946580in}{0.660842in}}%
\pgfpathlineto{\pgfqpoint{4.950211in}{0.637065in}}%
\pgfpathlineto{\pgfqpoint{4.953843in}{0.635222in}}%
\pgfpathlineto{\pgfqpoint{4.957474in}{0.685082in}}%
\pgfpathlineto{\pgfqpoint{4.961105in}{0.641547in}}%
\pgfpathlineto{\pgfqpoint{4.964736in}{0.618787in}}%
\pgfpathlineto{\pgfqpoint{4.968367in}{0.654262in}}%
\pgfpathlineto{\pgfqpoint{4.971998in}{0.634623in}}%
\pgfpathlineto{\pgfqpoint{4.975630in}{0.591458in}}%
\pgfpathlineto{\pgfqpoint{4.979261in}{0.621257in}}%
\pgfpathlineto{\pgfqpoint{4.982892in}{0.668683in}}%
\pgfpathlineto{\pgfqpoint{4.986523in}{0.693771in}}%
\pgfpathlineto{\pgfqpoint{4.990154in}{0.644360in}}%
\pgfpathlineto{\pgfqpoint{4.993785in}{0.606831in}}%
\pgfpathlineto{\pgfqpoint{4.997417in}{0.662199in}}%
\pgfpathlineto{\pgfqpoint{5.001048in}{0.603086in}}%
\pgfpathlineto{\pgfqpoint{5.004679in}{0.606721in}}%
\pgfpathlineto{\pgfqpoint{5.008310in}{0.653166in}}%
\pgfpathlineto{\pgfqpoint{5.011941in}{0.677470in}}%
\pgfpathlineto{\pgfqpoint{5.015573in}{0.690569in}}%
\pgfpathlineto{\pgfqpoint{5.019204in}{0.640027in}}%
\pgfpathlineto{\pgfqpoint{5.022835in}{0.574914in}}%
\pgfpathlineto{\pgfqpoint{5.026466in}{0.672945in}}%
\pgfpathlineto{\pgfqpoint{5.030097in}{0.537139in}}%
\pgfpathlineto{\pgfqpoint{5.033728in}{0.618646in}}%
\pgfpathlineto{\pgfqpoint{5.037360in}{0.656376in}}%
\pgfpathlineto{\pgfqpoint{5.040991in}{0.665093in}}%
\pgfpathlineto{\pgfqpoint{5.044622in}{0.664295in}}%
\pgfpathlineto{\pgfqpoint{5.048253in}{0.630145in}}%
\pgfpathlineto{\pgfqpoint{5.051884in}{0.625329in}}%
\pgfpathlineto{\pgfqpoint{5.055515in}{0.633729in}}%
\pgfpathlineto{\pgfqpoint{5.059147in}{0.563331in}}%
\pgfpathlineto{\pgfqpoint{5.062778in}{0.656584in}}%
\pgfpathlineto{\pgfqpoint{5.066409in}{0.680753in}}%
\pgfpathlineto{\pgfqpoint{5.070040in}{0.673529in}}%
\pgfpathlineto{\pgfqpoint{5.073671in}{0.668287in}}%
\pgfpathlineto{\pgfqpoint{5.077302in}{0.651132in}}%
\pgfpathlineto{\pgfqpoint{5.080934in}{0.621670in}}%
\pgfpathlineto{\pgfqpoint{5.084565in}{0.615415in}}%
\pgfpathlineto{\pgfqpoint{5.088196in}{0.626877in}}%
\pgfpathlineto{\pgfqpoint{5.091827in}{0.659211in}}%
\pgfpathlineto{\pgfqpoint{5.095458in}{0.675008in}}%
\pgfpathlineto{\pgfqpoint{5.099090in}{0.663891in}}%
\pgfpathlineto{\pgfqpoint{5.102721in}{0.636611in}}%
\pgfpathlineto{\pgfqpoint{5.106352in}{0.670847in}}%
\pgfpathlineto{\pgfqpoint{5.109983in}{0.625432in}}%
\pgfpathlineto{\pgfqpoint{5.113614in}{0.611812in}}%
\pgfpathlineto{\pgfqpoint{5.117245in}{0.642576in}}%
\pgfpathlineto{\pgfqpoint{5.120877in}{0.645534in}}%
\pgfpathlineto{\pgfqpoint{5.124508in}{0.639854in}}%
\pgfpathlineto{\pgfqpoint{5.128139in}{0.610876in}}%
\pgfpathlineto{\pgfqpoint{5.131770in}{0.542397in}}%
\pgfpathlineto{\pgfqpoint{5.135401in}{0.669656in}}%
\pgfpathlineto{\pgfqpoint{5.139032in}{0.554958in}}%
\pgfpathlineto{\pgfqpoint{5.142664in}{0.619142in}}%
\pgfpathlineto{\pgfqpoint{5.146295in}{0.656497in}}%
\pgfpathlineto{\pgfqpoint{5.149926in}{0.671438in}}%
\pgfpathlineto{\pgfqpoint{5.153557in}{0.675707in}}%
\pgfpathlineto{\pgfqpoint{5.157188in}{0.618315in}}%
\pgfpathlineto{\pgfqpoint{5.160819in}{0.615939in}}%
\pgfpathlineto{\pgfqpoint{5.164451in}{0.630798in}}%
\pgfpathlineto{\pgfqpoint{5.168082in}{0.545628in}}%
\pgfpathlineto{\pgfqpoint{5.171713in}{0.657664in}}%
\pgfpathlineto{\pgfqpoint{5.175344in}{0.679356in}}%
\pgfpathlineto{\pgfqpoint{5.178975in}{0.634072in}}%
\pgfpathlineto{\pgfqpoint{5.182606in}{0.612640in}}%
\pgfpathlineto{\pgfqpoint{5.186238in}{0.646132in}}%
\pgfpathlineto{\pgfqpoint{5.189869in}{0.640692in}}%
\pgfpathlineto{\pgfqpoint{5.193500in}{0.622795in}}%
\pgfpathlineto{\pgfqpoint{5.197131in}{0.526391in}}%
\pgfpathlineto{\pgfqpoint{5.200762in}{0.587433in}}%
\pgfpathlineto{\pgfqpoint{5.204394in}{0.669417in}}%
\pgfpathlineto{\pgfqpoint{5.208025in}{0.546473in}}%
\pgfpathlineto{\pgfqpoint{5.211656in}{0.623805in}}%
\pgfpathlineto{\pgfqpoint{5.215287in}{0.660053in}}%
\pgfpathlineto{\pgfqpoint{5.218918in}{0.668122in}}%
\pgfpathlineto{\pgfqpoint{5.222549in}{0.665905in}}%
\pgfpathlineto{\pgfqpoint{5.226181in}{0.429936in}}%
\pgfpathlineto{\pgfqpoint{5.229812in}{0.619652in}}%
\pgfpathlineto{\pgfqpoint{5.233443in}{0.679278in}}%
\pgfpathlineto{\pgfqpoint{5.237074in}{0.571599in}}%
\pgfpathlineto{\pgfqpoint{5.240705in}{0.621269in}}%
\pgfpathlineto{\pgfqpoint{5.244336in}{0.650502in}}%
\pgfpathlineto{\pgfqpoint{5.247968in}{0.637750in}}%
\pgfpathlineto{\pgfqpoint{5.251599in}{0.552594in}}%
\pgfpathlineto{\pgfqpoint{5.255230in}{0.650993in}}%
\pgfpathlineto{\pgfqpoint{5.258861in}{0.662758in}}%
\pgfpathlineto{\pgfqpoint{5.262492in}{0.688688in}}%
\pgfpathlineto{\pgfqpoint{5.266123in}{0.607505in}}%
\pgfpathlineto{\pgfqpoint{5.269755in}{0.611205in}}%
\pgfpathlineto{\pgfqpoint{5.273386in}{0.630101in}}%
\pgfpathlineto{\pgfqpoint{5.277017in}{0.484463in}}%
\pgfpathlineto{\pgfqpoint{5.280648in}{0.656409in}}%
\pgfpathlineto{\pgfqpoint{5.284279in}{0.677109in}}%
\pgfpathlineto{\pgfqpoint{5.287910in}{0.632209in}}%
\pgfpathlineto{\pgfqpoint{5.291542in}{0.625300in}}%
\pgfpathlineto{\pgfqpoint{5.295173in}{0.640887in}}%
\pgfpathlineto{\pgfqpoint{5.298804in}{0.638533in}}%
\pgfpathlineto{\pgfqpoint{5.302435in}{0.623521in}}%
\pgfpathlineto{\pgfqpoint{5.306066in}{0.599786in}}%
\pgfpathlineto{\pgfqpoint{5.309698in}{0.521621in}}%
\pgfpathlineto{\pgfqpoint{5.313329in}{0.668902in}}%
\pgfpathlineto{\pgfqpoint{5.316960in}{0.558871in}}%
\pgfpathlineto{\pgfqpoint{5.320591in}{0.619311in}}%
\pgfpathlineto{\pgfqpoint{5.324222in}{0.656620in}}%
\pgfpathlineto{\pgfqpoint{5.327853in}{0.672556in}}%
\pgfpathlineto{\pgfqpoint{5.331485in}{0.677661in}}%
\pgfpathlineto{\pgfqpoint{5.335116in}{0.615723in}}%
\pgfpathlineto{\pgfqpoint{5.338747in}{0.614006in}}%
\pgfpathlineto{\pgfqpoint{5.342378in}{0.630373in}}%
\pgfpathlineto{\pgfqpoint{5.346009in}{0.541852in}}%
\pgfpathlineto{\pgfqpoint{5.349640in}{0.657786in}}%
\pgfpathlineto{\pgfqpoint{5.353272in}{0.678989in}}%
\pgfpathlineto{\pgfqpoint{5.356903in}{0.634075in}}%
\pgfpathlineto{\pgfqpoint{5.360534in}{0.614776in}}%
\pgfpathlineto{\pgfqpoint{5.364165in}{0.645315in}}%
\pgfpathlineto{\pgfqpoint{5.367796in}{0.640664in}}%
\pgfpathlineto{\pgfqpoint{5.371427in}{0.623933in}}%
\pgfpathlineto{\pgfqpoint{5.375059in}{0.552264in}}%
\pgfpathlineto{\pgfqpoint{5.378690in}{0.578710in}}%
\pgfpathlineto{\pgfqpoint{5.382321in}{0.669446in}}%
\pgfpathlineto{\pgfqpoint{5.385952in}{0.550114in}}%
\pgfpathlineto{\pgfqpoint{5.389583in}{0.622938in}}%
\pgfpathlineto{\pgfqpoint{5.393215in}{0.659514in}}%
\pgfpathlineto{\pgfqpoint{5.396846in}{0.668831in}}%
\pgfpathlineto{\pgfqpoint{5.400477in}{0.667816in}}%
\pgfpathlineto{\pgfqpoint{5.404108in}{0.541998in}}%
\pgfpathlineto{\pgfqpoint{5.411370in}{0.679892in}}%
\pgfpathlineto{\pgfqpoint{5.415002in}{0.567254in}}%
\pgfpathlineto{\pgfqpoint{5.418633in}{0.620114in}}%
\pgfpathlineto{\pgfqpoint{5.422264in}{0.650227in}}%
\pgfpathlineto{\pgfqpoint{5.425895in}{0.638868in}}%
\pgfpathlineto{\pgfqpoint{5.429526in}{0.531316in}}%
\pgfpathlineto{\pgfqpoint{5.433157in}{0.647994in}}%
\pgfpathlineto{\pgfqpoint{5.436789in}{0.661992in}}%
\pgfpathlineto{\pgfqpoint{5.440420in}{0.690012in}}%
\pgfpathlineto{\pgfqpoint{5.444051in}{0.604147in}}%
\pgfpathlineto{\pgfqpoint{5.447682in}{0.609430in}}%
\pgfpathlineto{\pgfqpoint{5.451313in}{0.630237in}}%
\pgfpathlineto{\pgfqpoint{5.454944in}{0.510450in}}%
\pgfpathlineto{\pgfqpoint{5.458576in}{0.655982in}}%
\pgfpathlineto{\pgfqpoint{5.462207in}{0.676288in}}%
\pgfpathlineto{\pgfqpoint{5.465838in}{0.631804in}}%
\pgfpathlineto{\pgfqpoint{5.469469in}{0.627372in}}%
\pgfpathlineto{\pgfqpoint{5.473100in}{0.639848in}}%
\pgfpathlineto{\pgfqpoint{5.476731in}{0.638487in}}%
\pgfpathlineto{\pgfqpoint{5.480363in}{0.625067in}}%
\pgfpathlineto{\pgfqpoint{5.483994in}{0.606900in}}%
\pgfpathlineto{\pgfqpoint{5.487625in}{0.554399in}}%
\pgfpathlineto{\pgfqpoint{5.491256in}{0.668740in}}%
\pgfpathlineto{\pgfqpoint{5.494887in}{0.562489in}}%
\pgfpathlineto{\pgfqpoint{5.498519in}{0.617960in}}%
\pgfpathlineto{\pgfqpoint{5.502150in}{0.655907in}}%
\pgfpathlineto{\pgfqpoint{5.505781in}{0.673193in}}%
\pgfpathlineto{\pgfqpoint{5.509412in}{0.679121in}}%
\pgfpathlineto{\pgfqpoint{5.513043in}{0.611842in}}%
\pgfpathlineto{\pgfqpoint{5.516674in}{0.612363in}}%
\pgfpathlineto{\pgfqpoint{5.520306in}{0.631091in}}%
\pgfpathlineto{\pgfqpoint{5.523937in}{0.526074in}}%
\pgfpathlineto{\pgfqpoint{5.527568in}{0.657157in}}%
\pgfpathlineto{\pgfqpoint{5.531199in}{0.678254in}}%
\pgfpathlineto{\pgfqpoint{5.534830in}{0.634325in}}%
\pgfpathlineto{\pgfqpoint{5.538461in}{0.617946in}}%
\pgfpathlineto{\pgfqpoint{5.542093in}{0.644219in}}%
\pgfpathlineto{\pgfqpoint{5.545724in}{0.640801in}}%
\pgfpathlineto{\pgfqpoint{5.549355in}{0.626253in}}%
\pgfpathlineto{\pgfqpoint{5.552986in}{0.572644in}}%
\pgfpathlineto{\pgfqpoint{5.556617in}{0.564958in}}%
\pgfpathlineto{\pgfqpoint{5.560248in}{0.669494in}}%
\pgfpathlineto{\pgfqpoint{5.563880in}{0.557196in}}%
\pgfpathlineto{\pgfqpoint{5.567511in}{0.621358in}}%
\pgfpathlineto{\pgfqpoint{5.571142in}{0.658849in}}%
\pgfpathlineto{\pgfqpoint{5.574773in}{0.669575in}}%
\pgfpathlineto{\pgfqpoint{5.578404in}{0.671659in}}%
\pgfpathlineto{\pgfqpoint{5.582035in}{0.625039in}}%
\pgfpathlineto{\pgfqpoint{5.585667in}{0.620313in}}%
\pgfpathlineto{\pgfqpoint{5.589298in}{0.631420in}}%
\pgfpathlineto{\pgfqpoint{5.592929in}{0.562695in}}%
\pgfpathlineto{\pgfqpoint{5.596560in}{0.657720in}}%
\pgfpathlineto{\pgfqpoint{5.600191in}{0.680273in}}%
\pgfpathlineto{\pgfqpoint{5.603823in}{0.677771in}}%
\pgfpathlineto{\pgfqpoint{5.607454in}{0.676508in}}%
\pgfpathlineto{\pgfqpoint{5.614716in}{0.616251in}}%
\pgfpathlineto{\pgfqpoint{5.618347in}{0.617123in}}%
\pgfpathlineto{\pgfqpoint{5.621978in}{0.624561in}}%
\pgfpathlineto{\pgfqpoint{5.625610in}{0.658631in}}%
\pgfpathlineto{\pgfqpoint{5.629241in}{0.674287in}}%
\pgfpathlineto{\pgfqpoint{5.632872in}{0.669206in}}%
\pgfpathlineto{\pgfqpoint{5.636503in}{0.650071in}}%
\pgfpathlineto{\pgfqpoint{5.640134in}{0.667329in}}%
\pgfpathlineto{\pgfqpoint{5.643765in}{0.622619in}}%
\pgfpathlineto{\pgfqpoint{5.647397in}{0.620033in}}%
\pgfpathlineto{\pgfqpoint{5.651028in}{0.639855in}}%
\pgfpathlineto{\pgfqpoint{5.654659in}{0.643222in}}%
\pgfpathlineto{\pgfqpoint{5.658290in}{0.637140in}}%
\pgfpathlineto{\pgfqpoint{5.661921in}{0.620616in}}%
\pgfpathlineto{\pgfqpoint{5.665552in}{0.584902in}}%
\pgfpathlineto{\pgfqpoint{5.669184in}{0.668995in}}%
\pgfpathlineto{\pgfqpoint{5.672815in}{0.558889in}}%
\pgfpathlineto{\pgfqpoint{5.676446in}{0.617166in}}%
\pgfpathlineto{\pgfqpoint{5.680077in}{0.654904in}}%
\pgfpathlineto{\pgfqpoint{5.683708in}{0.673480in}}%
\pgfpathlineto{\pgfqpoint{5.687340in}{0.679824in}}%
\pgfpathlineto{\pgfqpoint{5.690971in}{0.609851in}}%
\pgfpathlineto{\pgfqpoint{5.694602in}{0.611061in}}%
\pgfpathlineto{\pgfqpoint{5.698233in}{0.630770in}}%
\pgfpathlineto{\pgfqpoint{5.701864in}{0.516349in}}%
\pgfpathlineto{\pgfqpoint{5.705495in}{0.657156in}}%
\pgfpathlineto{\pgfqpoint{5.709127in}{0.677924in}}%
\pgfpathlineto{\pgfqpoint{5.712758in}{0.634116in}}%
\pgfpathlineto{\pgfqpoint{5.716389in}{0.619789in}}%
\pgfpathlineto{\pgfqpoint{5.720020in}{0.643416in}}%
\pgfpathlineto{\pgfqpoint{5.723651in}{0.640561in}}%
\pgfpathlineto{\pgfqpoint{5.727282in}{0.626610in}}%
\pgfpathlineto{\pgfqpoint{5.730914in}{0.582069in}}%
\pgfpathlineto{\pgfqpoint{5.734545in}{0.548758in}}%
\pgfpathlineto{\pgfqpoint{5.738176in}{0.669421in}}%
\pgfpathlineto{\pgfqpoint{5.741807in}{0.558912in}}%
\pgfpathlineto{\pgfqpoint{5.745438in}{0.620632in}}%
\pgfpathlineto{\pgfqpoint{5.749069in}{0.658311in}}%
\pgfpathlineto{\pgfqpoint{5.752701in}{0.670246in}}%
\pgfpathlineto{\pgfqpoint{5.756332in}{0.673097in}}%
\pgfpathlineto{\pgfqpoint{5.759963in}{0.622787in}}%
\pgfpathlineto{\pgfqpoint{5.763594in}{0.619031in}}%
\pgfpathlineto{\pgfqpoint{5.767225in}{0.631517in}}%
\pgfpathlineto{\pgfqpoint{5.770856in}{0.557175in}}%
\pgfpathlineto{\pgfqpoint{5.774488in}{0.657563in}}%
\pgfpathlineto{\pgfqpoint{5.778119in}{0.679943in}}%
\pgfpathlineto{\pgfqpoint{5.781750in}{0.678752in}}%
\pgfpathlineto{\pgfqpoint{5.785381in}{0.678770in}}%
\pgfpathlineto{\pgfqpoint{5.792644in}{0.614674in}}%
\pgfpathlineto{\pgfqpoint{5.796275in}{0.618048in}}%
\pgfpathlineto{\pgfqpoint{5.799906in}{0.623224in}}%
\pgfpathlineto{\pgfqpoint{5.803537in}{0.658593in}}%
\pgfpathlineto{\pgfqpoint{5.807168in}{0.674448in}}%
\pgfpathlineto{\pgfqpoint{5.810799in}{0.670967in}}%
\pgfpathlineto{\pgfqpoint{5.814431in}{0.654348in}}%
\pgfpathlineto{\pgfqpoint{5.818062in}{0.665562in}}%
\pgfpathlineto{\pgfqpoint{5.821693in}{0.621403in}}%
\pgfpathlineto{\pgfqpoint{5.825324in}{0.622010in}}%
\pgfpathlineto{\pgfqpoint{5.828955in}{0.638838in}}%
\pgfpathlineto{\pgfqpoint{5.832586in}{0.643099in}}%
\pgfpathlineto{\pgfqpoint{5.836218in}{0.638034in}}%
\pgfpathlineto{\pgfqpoint{5.839849in}{0.625088in}}%
\pgfpathlineto{\pgfqpoint{5.843480in}{0.595400in}}%
\pgfpathlineto{\pgfqpoint{5.847111in}{0.668754in}}%
\pgfpathlineto{\pgfqpoint{5.850742in}{0.561339in}}%
\pgfpathlineto{\pgfqpoint{5.854373in}{0.615985in}}%
\pgfpathlineto{\pgfqpoint{5.858005in}{0.654204in}}%
\pgfpathlineto{\pgfqpoint{5.861636in}{0.674095in}}%
\pgfpathlineto{\pgfqpoint{5.865267in}{0.681172in}}%
\pgfpathlineto{\pgfqpoint{5.868898in}{0.605790in}}%
\pgfpathlineto{\pgfqpoint{5.872529in}{0.609233in}}%
\pgfpathlineto{\pgfqpoint{5.876160in}{0.631229in}}%
\pgfpathlineto{\pgfqpoint{5.879792in}{0.487895in}}%
\pgfpathlineto{\pgfqpoint{5.883423in}{0.656616in}}%
\pgfpathlineto{\pgfqpoint{5.887054in}{0.677164in}}%
\pgfpathlineto{\pgfqpoint{5.890685in}{0.634120in}}%
\pgfpathlineto{\pgfqpoint{5.894316in}{0.622416in}}%
\pgfpathlineto{\pgfqpoint{5.897948in}{0.642324in}}%
\pgfpathlineto{\pgfqpoint{5.901579in}{0.640606in}}%
\pgfpathlineto{\pgfqpoint{5.905210in}{0.628469in}}%
\pgfpathlineto{\pgfqpoint{5.908841in}{0.592904in}}%
\pgfpathlineto{\pgfqpoint{5.912472in}{0.511536in}}%
\pgfpathlineto{\pgfqpoint{5.916103in}{0.669356in}}%
\pgfpathlineto{\pgfqpoint{5.919735in}{0.563741in}}%
\pgfpathlineto{\pgfqpoint{5.923366in}{0.619102in}}%
\pgfpathlineto{\pgfqpoint{5.926997in}{0.657612in}}%
\pgfpathlineto{\pgfqpoint{5.930628in}{0.670927in}}%
\pgfpathlineto{\pgfqpoint{5.934259in}{0.674724in}}%
\pgfpathlineto{\pgfqpoint{5.937890in}{0.619246in}}%
\pgfpathlineto{\pgfqpoint{5.941522in}{0.617671in}}%
\pgfpathlineto{\pgfqpoint{5.945153in}{0.632484in}}%
\pgfpathlineto{\pgfqpoint{5.948784in}{0.545003in}}%
\pgfpathlineto{\pgfqpoint{5.952415in}{0.656849in}}%
\pgfpathlineto{\pgfqpoint{5.956046in}{0.679282in}}%
\pgfpathlineto{\pgfqpoint{5.959677in}{0.634545in}}%
\pgfpathlineto{\pgfqpoint{5.963309in}{0.610967in}}%
\pgfpathlineto{\pgfqpoint{5.966940in}{0.646943in}}%
\pgfpathlineto{\pgfqpoint{5.970571in}{0.641325in}}%
\pgfpathlineto{\pgfqpoint{5.974202in}{0.624101in}}%
\pgfpathlineto{\pgfqpoint{5.977833in}{0.482971in}}%
\pgfpathlineto{\pgfqpoint{5.981465in}{0.653204in}}%
\pgfpathlineto{\pgfqpoint{5.985096in}{0.691285in}}%
\pgfpathlineto{\pgfqpoint{5.988727in}{0.642601in}}%
\pgfpathlineto{\pgfqpoint{5.992358in}{0.616653in}}%
\pgfpathlineto{\pgfqpoint{5.995989in}{0.655910in}}%
\pgfpathlineto{\pgfqpoint{5.999620in}{0.615158in}}%
\pgfpathlineto{\pgfqpoint{6.003252in}{0.585770in}}%
\pgfpathlineto{\pgfqpoint{6.006883in}{0.642204in}}%
\pgfpathlineto{\pgfqpoint{6.010514in}{0.678586in}}%
\pgfpathlineto{\pgfqpoint{6.014145in}{0.696849in}}%
\pgfpathlineto{\pgfqpoint{6.017776in}{0.651516in}}%
\pgfpathlineto{\pgfqpoint{6.021407in}{0.618457in}}%
\pgfpathlineto{\pgfqpoint{6.025039in}{0.666432in}}%
\pgfpathlineto{\pgfqpoint{6.028670in}{0.570193in}}%
\pgfpathlineto{\pgfqpoint{6.032301in}{0.618252in}}%
\pgfpathlineto{\pgfqpoint{6.035932in}{0.655748in}}%
\pgfpathlineto{\pgfqpoint{6.039563in}{0.676204in}}%
\pgfpathlineto{\pgfqpoint{6.043194in}{0.684470in}}%
\pgfpathlineto{\pgfqpoint{6.046826in}{0.602236in}}%
\pgfpathlineto{\pgfqpoint{6.050457in}{0.605500in}}%
\pgfpathlineto{\pgfqpoint{6.054088in}{0.629674in}}%
\pgfpathlineto{\pgfqpoint{6.057719in}{0.506582in}}%
\pgfpathlineto{\pgfqpoint{6.061350in}{0.657507in}}%
\pgfpathlineto{\pgfqpoint{6.064981in}{0.676882in}}%
\pgfpathlineto{\pgfqpoint{6.068613in}{0.633740in}}%
\pgfpathlineto{\pgfqpoint{6.072244in}{0.623364in}}%
\pgfpathlineto{\pgfqpoint{6.075875in}{0.641604in}}%
\pgfpathlineto{\pgfqpoint{6.079506in}{0.640581in}}%
\pgfpathlineto{\pgfqpoint{6.083137in}{0.629200in}}%
\pgfpathlineto{\pgfqpoint{6.086769in}{0.599346in}}%
\pgfpathlineto{\pgfqpoint{6.090400in}{0.452670in}}%
\pgfpathlineto{\pgfqpoint{6.094031in}{0.669244in}}%
\pgfpathlineto{\pgfqpoint{6.097662in}{0.564125in}}%
\pgfpathlineto{\pgfqpoint{6.101293in}{0.618497in}}%
\pgfpathlineto{\pgfqpoint{6.104924in}{0.657046in}}%
\pgfpathlineto{\pgfqpoint{6.108556in}{0.671547in}}%
\pgfpathlineto{\pgfqpoint{6.112187in}{0.676013in}}%
\pgfpathlineto{\pgfqpoint{6.115818in}{0.616935in}}%
\pgfpathlineto{\pgfqpoint{6.119449in}{0.616259in}}%
\pgfpathlineto{\pgfqpoint{6.123080in}{0.632359in}}%
\pgfpathlineto{\pgfqpoint{6.126711in}{0.538119in}}%
\pgfpathlineto{\pgfqpoint{6.130343in}{0.656773in}}%
\pgfpathlineto{\pgfqpoint{6.133974in}{0.678937in}}%
\pgfpathlineto{\pgfqpoint{6.137605in}{0.634633in}}%
\pgfpathlineto{\pgfqpoint{6.141236in}{0.613387in}}%
\pgfpathlineto{\pgfqpoint{6.144867in}{0.646114in}}%
\pgfpathlineto{\pgfqpoint{6.148498in}{0.641244in}}%
\pgfpathlineto{\pgfqpoint{6.152130in}{0.625086in}}%
\pgfpathlineto{\pgfqpoint{6.155761in}{0.533202in}}%
\pgfpathlineto{\pgfqpoint{6.159392in}{0.651903in}}%
\pgfpathlineto{\pgfqpoint{6.163023in}{0.691046in}}%
\pgfpathlineto{\pgfqpoint{6.166654in}{0.643545in}}%
\pgfpathlineto{\pgfqpoint{6.170285in}{0.620455in}}%
\pgfpathlineto{\pgfqpoint{6.173917in}{0.655310in}}%
\pgfpathlineto{\pgfqpoint{6.177548in}{0.615171in}}%
\pgfpathlineto{\pgfqpoint{6.181179in}{0.584060in}}%
\pgfpathlineto{\pgfqpoint{6.184810in}{0.640743in}}%
\pgfpathlineto{\pgfqpoint{6.188441in}{0.678820in}}%
\pgfpathlineto{\pgfqpoint{6.192073in}{0.697258in}}%
\pgfpathlineto{\pgfqpoint{6.195704in}{0.653928in}}%
\pgfpathlineto{\pgfqpoint{6.199335in}{0.624215in}}%
\pgfpathlineto{\pgfqpoint{6.202966in}{0.666182in}}%
\pgfpathlineto{\pgfqpoint{6.206597in}{0.571490in}}%
\pgfpathlineto{\pgfqpoint{6.210228in}{0.617249in}}%
\pgfpathlineto{\pgfqpoint{6.213860in}{0.655016in}}%
\pgfpathlineto{\pgfqpoint{6.217491in}{0.676784in}}%
\pgfpathlineto{\pgfqpoint{6.221122in}{0.685719in}}%
\pgfpathlineto{\pgfqpoint{6.224753in}{0.597808in}}%
\pgfpathlineto{\pgfqpoint{6.228384in}{0.603324in}}%
\pgfpathlineto{\pgfqpoint{6.232015in}{0.629891in}}%
\pgfpathlineto{\pgfqpoint{6.235647in}{0.476702in}}%
\pgfpathlineto{\pgfqpoint{6.239278in}{0.657115in}}%
\pgfpathlineto{\pgfqpoint{6.242909in}{0.676130in}}%
\pgfpathlineto{\pgfqpoint{6.246540in}{0.633533in}}%
\pgfpathlineto{\pgfqpoint{6.250171in}{0.625592in}}%
\pgfpathlineto{\pgfqpoint{6.253802in}{0.640526in}}%
\pgfpathlineto{\pgfqpoint{6.257434in}{0.640549in}}%
\pgfpathlineto{\pgfqpoint{6.261065in}{0.630670in}}%
\pgfpathlineto{\pgfqpoint{6.264696in}{0.606580in}}%
\pgfpathlineto{\pgfqpoint{6.268327in}{0.534331in}}%
\pgfpathlineto{\pgfqpoint{6.271958in}{0.669088in}}%
\pgfpathlineto{\pgfqpoint{6.275590in}{0.567571in}}%
\pgfpathlineto{\pgfqpoint{6.279221in}{0.617042in}}%
\pgfpathlineto{\pgfqpoint{6.282852in}{0.656325in}}%
\pgfpathlineto{\pgfqpoint{6.286483in}{0.672181in}}%
\pgfpathlineto{\pgfqpoint{6.290114in}{0.677504in}}%
\pgfpathlineto{\pgfqpoint{6.293745in}{0.613098in}}%
\pgfpathlineto{\pgfqpoint{6.297377in}{0.614706in}}%
\pgfpathlineto{\pgfqpoint{6.301008in}{0.633101in}}%
\pgfpathlineto{\pgfqpoint{6.304639in}{0.519712in}}%
\pgfpathlineto{\pgfqpoint{6.308270in}{0.656096in}}%
\pgfpathlineto{\pgfqpoint{6.311901in}{0.678222in}}%
\pgfpathlineto{\pgfqpoint{6.315532in}{0.634967in}}%
\pgfpathlineto{\pgfqpoint{6.319164in}{0.616717in}}%
\pgfpathlineto{\pgfqpoint{6.322795in}{0.645011in}}%
\pgfpathlineto{\pgfqpoint{6.326426in}{0.641391in}}%
\pgfpathlineto{\pgfqpoint{6.330057in}{0.627395in}}%
\pgfpathlineto{\pgfqpoint{6.333688in}{0.561763in}}%
\pgfpathlineto{\pgfqpoint{6.337319in}{0.578414in}}%
\pgfpathlineto{\pgfqpoint{6.340951in}{0.669648in}}%
\pgfpathlineto{\pgfqpoint{6.344582in}{0.559680in}}%
\pgfpathlineto{\pgfqpoint{6.348213in}{0.621418in}}%
\pgfpathlineto{\pgfqpoint{6.351844in}{0.659466in}}%
\pgfpathlineto{\pgfqpoint{6.355475in}{0.668643in}}%
\pgfpathlineto{\pgfqpoint{6.359106in}{0.667957in}}%
\pgfpathlineto{\pgfqpoint{6.362738in}{0.544709in}}%
\pgfpathlineto{\pgfqpoint{6.366369in}{0.595094in}}%
\pgfpathlineto{\pgfqpoint{6.370000in}{0.672405in}}%
\pgfpathlineto{\pgfqpoint{6.370000in}{0.672405in}}%
\pgfusepath{stroke}%
\end{pgfscope}%
\begin{pgfscope}%
\pgfpathrectangle{\pgfqpoint{4.554414in}{0.320625in}}{\pgfqpoint{1.815586in}{1.164375in}}%
\pgfusepath{clip}%
\pgfsetrectcap%
\pgfsetroundjoin%
\pgfsetlinewidth{0.602250pt}%
\definecolor{currentstroke}{rgb}{0.478431,0.243137,0.615686}%
\pgfsetstrokecolor{currentstroke}%
\pgfsetstrokeopacity{0.120000}%
\pgfsetdash{}{0pt}%
\pgfpathmoveto{\pgfqpoint{4.554414in}{1.452830in}}%
\pgfpathlineto{\pgfqpoint{4.558045in}{1.442989in}}%
\pgfpathlineto{\pgfqpoint{4.579832in}{1.398306in}}%
\pgfpathlineto{\pgfqpoint{4.601619in}{1.349469in}}%
\pgfpathlineto{\pgfqpoint{4.670611in}{1.192608in}}%
\pgfpathlineto{\pgfqpoint{4.710554in}{1.105091in}}%
\pgfpathlineto{\pgfqpoint{4.754128in}{1.012181in}}%
\pgfpathlineto{\pgfqpoint{4.804965in}{0.906405in}}%
\pgfpathlineto{\pgfqpoint{4.863063in}{0.788130in}}%
\pgfpathlineto{\pgfqpoint{4.881219in}{0.751617in}}%
\pgfpathlineto{\pgfqpoint{4.884850in}{0.745380in}}%
\pgfpathlineto{\pgfqpoint{4.892113in}{0.729762in}}%
\pgfpathlineto{\pgfqpoint{4.895744in}{0.723688in}}%
\pgfpathlineto{\pgfqpoint{4.903006in}{0.707937in}}%
\pgfpathlineto{\pgfqpoint{4.906637in}{0.702781in}}%
\pgfpathlineto{\pgfqpoint{4.910269in}{0.699428in}}%
\pgfpathlineto{\pgfqpoint{4.913900in}{0.697335in}}%
\pgfpathlineto{\pgfqpoint{4.917531in}{0.679011in}}%
\pgfpathlineto{\pgfqpoint{4.921162in}{0.674276in}}%
\pgfpathlineto{\pgfqpoint{4.924793in}{0.680723in}}%
\pgfpathlineto{\pgfqpoint{4.928424in}{0.657600in}}%
\pgfpathlineto{\pgfqpoint{4.932056in}{0.660616in}}%
\pgfpathlineto{\pgfqpoint{4.935687in}{0.673488in}}%
\pgfpathlineto{\pgfqpoint{4.939318in}{0.655971in}}%
\pgfpathlineto{\pgfqpoint{4.942949in}{0.628741in}}%
\pgfpathlineto{\pgfqpoint{4.946580in}{0.676159in}}%
\pgfpathlineto{\pgfqpoint{4.950211in}{0.655843in}}%
\pgfpathlineto{\pgfqpoint{4.953843in}{0.643550in}}%
\pgfpathlineto{\pgfqpoint{4.957474in}{0.654816in}}%
\pgfpathlineto{\pgfqpoint{4.961105in}{0.639913in}}%
\pgfpathlineto{\pgfqpoint{4.964736in}{0.608820in}}%
\pgfpathlineto{\pgfqpoint{4.968367in}{0.633564in}}%
\pgfpathlineto{\pgfqpoint{4.971998in}{0.670593in}}%
\pgfpathlineto{\pgfqpoint{4.975630in}{0.690318in}}%
\pgfpathlineto{\pgfqpoint{4.979261in}{0.626488in}}%
\pgfpathlineto{\pgfqpoint{4.982892in}{0.577000in}}%
\pgfpathlineto{\pgfqpoint{4.986523in}{0.670256in}}%
\pgfpathlineto{\pgfqpoint{4.990154in}{0.561092in}}%
\pgfpathlineto{\pgfqpoint{4.993785in}{0.632766in}}%
\pgfpathlineto{\pgfqpoint{4.997417in}{0.667272in}}%
\pgfpathlineto{\pgfqpoint{5.001048in}{0.663672in}}%
\pgfpathlineto{\pgfqpoint{5.004679in}{0.639553in}}%
\pgfpathlineto{\pgfqpoint{5.008310in}{0.654411in}}%
\pgfpathlineto{\pgfqpoint{5.011941in}{0.648351in}}%
\pgfpathlineto{\pgfqpoint{5.015573in}{0.657782in}}%
\pgfpathlineto{\pgfqpoint{5.019204in}{0.638517in}}%
\pgfpathlineto{\pgfqpoint{5.022835in}{0.638169in}}%
\pgfpathlineto{\pgfqpoint{5.026466in}{0.627821in}}%
\pgfpathlineto{\pgfqpoint{5.030097in}{0.572098in}}%
\pgfpathlineto{\pgfqpoint{5.033728in}{0.664925in}}%
\pgfpathlineto{\pgfqpoint{5.037360in}{0.693007in}}%
\pgfpathlineto{\pgfqpoint{5.044622in}{0.613896in}}%
\pgfpathlineto{\pgfqpoint{5.048253in}{0.667100in}}%
\pgfpathlineto{\pgfqpoint{5.051884in}{0.594660in}}%
\pgfpathlineto{\pgfqpoint{5.055515in}{0.620068in}}%
\pgfpathlineto{\pgfqpoint{5.059147in}{0.662156in}}%
\pgfpathlineto{\pgfqpoint{5.062778in}{0.675212in}}%
\pgfpathlineto{\pgfqpoint{5.066409in}{0.675421in}}%
\pgfpathlineto{\pgfqpoint{5.070040in}{0.578736in}}%
\pgfpathlineto{\pgfqpoint{5.077302in}{0.668746in}}%
\pgfpathlineto{\pgfqpoint{5.080934in}{0.617366in}}%
\pgfpathlineto{\pgfqpoint{5.084565in}{0.636095in}}%
\pgfpathlineto{\pgfqpoint{5.088196in}{0.643256in}}%
\pgfpathlineto{\pgfqpoint{5.091827in}{0.611564in}}%
\pgfpathlineto{\pgfqpoint{5.095458in}{0.572384in}}%
\pgfpathlineto{\pgfqpoint{5.099090in}{0.671103in}}%
\pgfpathlineto{\pgfqpoint{5.102721in}{0.609487in}}%
\pgfpathlineto{\pgfqpoint{5.106352in}{0.609470in}}%
\pgfpathlineto{\pgfqpoint{5.109983in}{0.664002in}}%
\pgfpathlineto{\pgfqpoint{5.113614in}{0.666320in}}%
\pgfpathlineto{\pgfqpoint{5.117245in}{0.656942in}}%
\pgfpathlineto{\pgfqpoint{5.120877in}{0.630902in}}%
\pgfpathlineto{\pgfqpoint{5.124508in}{0.645634in}}%
\pgfpathlineto{\pgfqpoint{5.128139in}{0.666800in}}%
\pgfpathlineto{\pgfqpoint{5.131770in}{0.618754in}}%
\pgfpathlineto{\pgfqpoint{5.135401in}{0.639677in}}%
\pgfpathlineto{\pgfqpoint{5.139032in}{0.650362in}}%
\pgfpathlineto{\pgfqpoint{5.142664in}{0.587076in}}%
\pgfpathlineto{\pgfqpoint{5.146295in}{0.648644in}}%
\pgfpathlineto{\pgfqpoint{5.149926in}{0.686505in}}%
\pgfpathlineto{\pgfqpoint{5.157188in}{0.631449in}}%
\pgfpathlineto{\pgfqpoint{5.160819in}{0.664186in}}%
\pgfpathlineto{\pgfqpoint{5.168082in}{0.574445in}}%
\pgfpathlineto{\pgfqpoint{5.171713in}{0.654703in}}%
\pgfpathlineto{\pgfqpoint{5.175344in}{0.673383in}}%
\pgfpathlineto{\pgfqpoint{5.178975in}{0.680279in}}%
\pgfpathlineto{\pgfqpoint{5.182606in}{0.578017in}}%
\pgfpathlineto{\pgfqpoint{5.186238in}{0.617111in}}%
\pgfpathlineto{\pgfqpoint{5.189869in}{0.674399in}}%
\pgfpathlineto{\pgfqpoint{5.193500in}{0.581154in}}%
\pgfpathlineto{\pgfqpoint{5.197131in}{0.634132in}}%
\pgfpathlineto{\pgfqpoint{5.200762in}{0.657137in}}%
\pgfpathlineto{\pgfqpoint{5.204394in}{0.636018in}}%
\pgfpathlineto{\pgfqpoint{5.208025in}{0.609782in}}%
\pgfpathlineto{\pgfqpoint{5.211656in}{0.676579in}}%
\pgfpathlineto{\pgfqpoint{5.215287in}{0.662397in}}%
\pgfpathlineto{\pgfqpoint{5.218918in}{0.660170in}}%
\pgfpathlineto{\pgfqpoint{5.222549in}{0.652797in}}%
\pgfpathlineto{\pgfqpoint{5.226181in}{0.624054in}}%
\pgfpathlineto{\pgfqpoint{5.229812in}{0.527274in}}%
\pgfpathlineto{\pgfqpoint{5.233443in}{0.638425in}}%
\pgfpathlineto{\pgfqpoint{5.237074in}{0.676777in}}%
\pgfpathlineto{\pgfqpoint{5.240705in}{0.693154in}}%
\pgfpathlineto{\pgfqpoint{5.244336in}{0.652058in}}%
\pgfpathlineto{\pgfqpoint{5.247968in}{0.592368in}}%
\pgfpathlineto{\pgfqpoint{5.251599in}{0.675894in}}%
\pgfpathlineto{\pgfqpoint{5.255230in}{0.488274in}}%
\pgfpathlineto{\pgfqpoint{5.258861in}{0.625500in}}%
\pgfpathlineto{\pgfqpoint{5.262492in}{0.658517in}}%
\pgfpathlineto{\pgfqpoint{5.266123in}{0.658371in}}%
\pgfpathlineto{\pgfqpoint{5.269755in}{0.642691in}}%
\pgfpathlineto{\pgfqpoint{5.273386in}{0.656801in}}%
\pgfpathlineto{\pgfqpoint{5.277017in}{0.628311in}}%
\pgfpathlineto{\pgfqpoint{5.280648in}{0.575759in}}%
\pgfpathlineto{\pgfqpoint{5.284279in}{0.641531in}}%
\pgfpathlineto{\pgfqpoint{5.287910in}{0.670549in}}%
\pgfpathlineto{\pgfqpoint{5.291542in}{0.683742in}}%
\pgfpathlineto{\pgfqpoint{5.295173in}{0.635484in}}%
\pgfpathlineto{\pgfqpoint{5.298804in}{0.568434in}}%
\pgfpathlineto{\pgfqpoint{5.302435in}{0.680360in}}%
\pgfpathlineto{\pgfqpoint{5.306066in}{0.550148in}}%
\pgfpathlineto{\pgfqpoint{5.309698in}{0.619236in}}%
\pgfpathlineto{\pgfqpoint{5.313329in}{0.658970in}}%
\pgfpathlineto{\pgfqpoint{5.316960in}{0.645041in}}%
\pgfpathlineto{\pgfqpoint{5.320591in}{0.493455in}}%
\pgfpathlineto{\pgfqpoint{5.324222in}{0.662134in}}%
\pgfpathlineto{\pgfqpoint{5.327853in}{0.665221in}}%
\pgfpathlineto{\pgfqpoint{5.331485in}{0.676518in}}%
\pgfpathlineto{\pgfqpoint{5.335116in}{0.638460in}}%
\pgfpathlineto{\pgfqpoint{5.338747in}{0.625583in}}%
\pgfpathlineto{\pgfqpoint{5.342378in}{0.608591in}}%
\pgfpathlineto{\pgfqpoint{5.346009in}{0.617220in}}%
\pgfpathlineto{\pgfqpoint{5.349640in}{0.671242in}}%
\pgfpathlineto{\pgfqpoint{5.353272in}{0.690465in}}%
\pgfpathlineto{\pgfqpoint{5.356903in}{0.667455in}}%
\pgfpathlineto{\pgfqpoint{5.360534in}{0.637827in}}%
\pgfpathlineto{\pgfqpoint{5.364165in}{0.675712in}}%
\pgfpathlineto{\pgfqpoint{5.367796in}{0.590111in}}%
\pgfpathlineto{\pgfqpoint{5.371427in}{0.600628in}}%
\pgfpathlineto{\pgfqpoint{5.375059in}{0.655460in}}%
\pgfpathlineto{\pgfqpoint{5.378690in}{0.659985in}}%
\pgfpathlineto{\pgfqpoint{5.382321in}{0.650618in}}%
\pgfpathlineto{\pgfqpoint{5.393215in}{0.630703in}}%
\pgfpathlineto{\pgfqpoint{5.396846in}{0.616556in}}%
\pgfpathlineto{\pgfqpoint{5.400477in}{0.663529in}}%
\pgfpathlineto{\pgfqpoint{5.404108in}{0.684473in}}%
\pgfpathlineto{\pgfqpoint{5.407739in}{0.644854in}}%
\pgfpathlineto{\pgfqpoint{5.411370in}{0.500368in}}%
\pgfpathlineto{\pgfqpoint{5.415002in}{0.676737in}}%
\pgfpathlineto{\pgfqpoint{5.418633in}{0.609190in}}%
\pgfpathlineto{\pgfqpoint{5.422264in}{0.598060in}}%
\pgfpathlineto{\pgfqpoint{5.425895in}{0.661787in}}%
\pgfpathlineto{\pgfqpoint{5.429526in}{0.655694in}}%
\pgfpathlineto{\pgfqpoint{5.433157in}{0.626356in}}%
\pgfpathlineto{\pgfqpoint{5.444051in}{0.676808in}}%
\pgfpathlineto{\pgfqpoint{5.447682in}{0.617378in}}%
\pgfpathlineto{\pgfqpoint{5.451313in}{0.631851in}}%
\pgfpathlineto{\pgfqpoint{5.454944in}{0.642120in}}%
\pgfpathlineto{\pgfqpoint{5.458576in}{0.542497in}}%
\pgfpathlineto{\pgfqpoint{5.462207in}{0.658929in}}%
\pgfpathlineto{\pgfqpoint{5.465838in}{0.685927in}}%
\pgfpathlineto{\pgfqpoint{5.469469in}{0.665276in}}%
\pgfpathlineto{\pgfqpoint{5.473100in}{0.639793in}}%
\pgfpathlineto{\pgfqpoint{5.476731in}{0.669293in}}%
\pgfpathlineto{\pgfqpoint{5.480363in}{0.618484in}}%
\pgfpathlineto{\pgfqpoint{5.483994in}{0.544942in}}%
\pgfpathlineto{\pgfqpoint{5.487625in}{0.654170in}}%
\pgfpathlineto{\pgfqpoint{5.491256in}{0.664408in}}%
\pgfpathlineto{\pgfqpoint{5.494887in}{0.661532in}}%
\pgfpathlineto{\pgfqpoint{5.498519in}{0.428580in}}%
\pgfpathlineto{\pgfqpoint{5.502150in}{0.639450in}}%
\pgfpathlineto{\pgfqpoint{5.505781in}{0.685682in}}%
\pgfpathlineto{\pgfqpoint{5.509412in}{0.565603in}}%
\pgfpathlineto{\pgfqpoint{5.513043in}{0.621694in}}%
\pgfpathlineto{\pgfqpoint{5.516674in}{0.648536in}}%
\pgfpathlineto{\pgfqpoint{5.520306in}{0.598702in}}%
\pgfpathlineto{\pgfqpoint{5.523937in}{0.646796in}}%
\pgfpathlineto{\pgfqpoint{5.527568in}{0.680222in}}%
\pgfpathlineto{\pgfqpoint{5.531199in}{0.674847in}}%
\pgfpathlineto{\pgfqpoint{5.534830in}{0.672132in}}%
\pgfpathlineto{\pgfqpoint{5.538461in}{0.655574in}}%
\pgfpathlineto{\pgfqpoint{5.542093in}{0.618824in}}%
\pgfpathlineto{\pgfqpoint{5.545724in}{0.558728in}}%
\pgfpathlineto{\pgfqpoint{5.549355in}{0.642005in}}%
\pgfpathlineto{\pgfqpoint{5.552986in}{0.669968in}}%
\pgfpathlineto{\pgfqpoint{5.556617in}{0.680767in}}%
\pgfpathlineto{\pgfqpoint{5.560248in}{0.547091in}}%
\pgfpathlineto{\pgfqpoint{5.563880in}{0.619543in}}%
\pgfpathlineto{\pgfqpoint{5.567511in}{0.651886in}}%
\pgfpathlineto{\pgfqpoint{5.571142in}{0.569701in}}%
\pgfpathlineto{\pgfqpoint{5.574773in}{0.649152in}}%
\pgfpathlineto{\pgfqpoint{5.578404in}{0.679978in}}%
\pgfpathlineto{\pgfqpoint{5.582035in}{0.673417in}}%
\pgfpathlineto{\pgfqpoint{5.585667in}{0.663758in}}%
\pgfpathlineto{\pgfqpoint{5.589298in}{0.659459in}}%
\pgfpathlineto{\pgfqpoint{5.596560in}{0.592200in}}%
\pgfpathlineto{\pgfqpoint{5.600191in}{0.644217in}}%
\pgfpathlineto{\pgfqpoint{5.603823in}{0.661887in}}%
\pgfpathlineto{\pgfqpoint{5.607454in}{0.666295in}}%
\pgfpathlineto{\pgfqpoint{5.611085in}{0.619643in}}%
\pgfpathlineto{\pgfqpoint{5.614716in}{0.621737in}}%
\pgfpathlineto{\pgfqpoint{5.618347in}{0.689382in}}%
\pgfpathlineto{\pgfqpoint{5.621978in}{0.583961in}}%
\pgfpathlineto{\pgfqpoint{5.625610in}{0.595200in}}%
\pgfpathlineto{\pgfqpoint{5.629241in}{0.653627in}}%
\pgfpathlineto{\pgfqpoint{5.632872in}{0.616161in}}%
\pgfpathlineto{\pgfqpoint{5.636503in}{0.631625in}}%
\pgfpathlineto{\pgfqpoint{5.640134in}{0.669914in}}%
\pgfpathlineto{\pgfqpoint{5.647397in}{0.687230in}}%
\pgfpathlineto{\pgfqpoint{5.651028in}{0.634675in}}%
\pgfpathlineto{\pgfqpoint{5.654659in}{0.616957in}}%
\pgfpathlineto{\pgfqpoint{5.658290in}{0.615570in}}%
\pgfpathlineto{\pgfqpoint{5.661921in}{0.627349in}}%
\pgfpathlineto{\pgfqpoint{5.665552in}{0.664779in}}%
\pgfpathlineto{\pgfqpoint{5.669184in}{0.679303in}}%
\pgfpathlineto{\pgfqpoint{5.672815in}{0.585662in}}%
\pgfpathlineto{\pgfqpoint{5.676446in}{0.611711in}}%
\pgfpathlineto{\pgfqpoint{5.680077in}{0.663379in}}%
\pgfpathlineto{\pgfqpoint{5.683708in}{0.623415in}}%
\pgfpathlineto{\pgfqpoint{5.687340in}{0.616325in}}%
\pgfpathlineto{\pgfqpoint{5.690971in}{0.671130in}}%
\pgfpathlineto{\pgfqpoint{5.694602in}{0.669098in}}%
\pgfpathlineto{\pgfqpoint{5.698233in}{0.665899in}}%
\pgfpathlineto{\pgfqpoint{5.701864in}{0.645525in}}%
\pgfpathlineto{\pgfqpoint{5.705495in}{0.635017in}}%
\pgfpathlineto{\pgfqpoint{5.709127in}{0.629958in}}%
\pgfpathlineto{\pgfqpoint{5.712758in}{0.628707in}}%
\pgfpathlineto{\pgfqpoint{5.716389in}{0.660637in}}%
\pgfpathlineto{\pgfqpoint{5.720020in}{0.675583in}}%
\pgfpathlineto{\pgfqpoint{5.723651in}{0.637289in}}%
\pgfpathlineto{\pgfqpoint{5.727282in}{0.587587in}}%
\pgfpathlineto{\pgfqpoint{5.730914in}{0.682931in}}%
\pgfpathlineto{\pgfqpoint{5.734545in}{0.616043in}}%
\pgfpathlineto{\pgfqpoint{5.738176in}{0.571447in}}%
\pgfpathlineto{\pgfqpoint{5.741807in}{0.660112in}}%
\pgfpathlineto{\pgfqpoint{5.745438in}{0.639853in}}%
\pgfpathlineto{\pgfqpoint{5.749069in}{0.491752in}}%
\pgfpathlineto{\pgfqpoint{5.752701in}{0.655526in}}%
\pgfpathlineto{\pgfqpoint{5.759963in}{0.683778in}}%
\pgfpathlineto{\pgfqpoint{5.763594in}{0.612930in}}%
\pgfpathlineto{\pgfqpoint{5.767225in}{0.625034in}}%
\pgfpathlineto{\pgfqpoint{5.770856in}{0.640674in}}%
\pgfpathlineto{\pgfqpoint{5.774488in}{0.598188in}}%
\pgfpathlineto{\pgfqpoint{5.778119in}{0.659495in}}%
\pgfpathlineto{\pgfqpoint{5.781750in}{0.680969in}}%
\pgfpathlineto{\pgfqpoint{5.785381in}{0.665323in}}%
\pgfpathlineto{\pgfqpoint{5.789012in}{0.632012in}}%
\pgfpathlineto{\pgfqpoint{5.792644in}{0.674884in}}%
\pgfpathlineto{\pgfqpoint{5.796275in}{0.623649in}}%
\pgfpathlineto{\pgfqpoint{5.799906in}{0.548359in}}%
\pgfpathlineto{\pgfqpoint{5.803537in}{0.653212in}}%
\pgfpathlineto{\pgfqpoint{5.807168in}{0.649856in}}%
\pgfpathlineto{\pgfqpoint{5.810799in}{0.621794in}}%
\pgfpathlineto{\pgfqpoint{5.814431in}{0.612660in}}%
\pgfpathlineto{\pgfqpoint{5.818062in}{0.662571in}}%
\pgfpathlineto{\pgfqpoint{5.821693in}{0.694217in}}%
\pgfpathlineto{\pgfqpoint{5.825324in}{0.477266in}}%
\pgfpathlineto{\pgfqpoint{5.828955in}{0.604232in}}%
\pgfpathlineto{\pgfqpoint{5.832586in}{0.646336in}}%
\pgfpathlineto{\pgfqpoint{5.836218in}{0.350297in}}%
\pgfpathlineto{\pgfqpoint{5.839849in}{0.655006in}}%
\pgfpathlineto{\pgfqpoint{5.843480in}{0.679601in}}%
\pgfpathlineto{\pgfqpoint{5.847111in}{0.680548in}}%
\pgfpathlineto{\pgfqpoint{5.850742in}{0.675157in}}%
\pgfpathlineto{\pgfqpoint{5.854373in}{0.657289in}}%
\pgfpathlineto{\pgfqpoint{5.858005in}{0.619287in}}%
\pgfpathlineto{\pgfqpoint{5.861636in}{0.601093in}}%
\pgfpathlineto{\pgfqpoint{5.865267in}{0.642507in}}%
\pgfpathlineto{\pgfqpoint{5.868898in}{0.656224in}}%
\pgfpathlineto{\pgfqpoint{5.872529in}{0.655052in}}%
\pgfpathlineto{\pgfqpoint{5.876160in}{0.623096in}}%
\pgfpathlineto{\pgfqpoint{5.879792in}{0.627813in}}%
\pgfpathlineto{\pgfqpoint{5.883423in}{0.694354in}}%
\pgfpathlineto{\pgfqpoint{5.887054in}{0.601474in}}%
\pgfpathlineto{\pgfqpoint{5.890685in}{0.543120in}}%
\pgfpathlineto{\pgfqpoint{5.894316in}{0.648441in}}%
\pgfpathlineto{\pgfqpoint{5.897948in}{0.594685in}}%
\pgfpathlineto{\pgfqpoint{5.901579in}{0.644334in}}%
\pgfpathlineto{\pgfqpoint{5.905210in}{0.670310in}}%
\pgfpathlineto{\pgfqpoint{5.908841in}{0.638229in}}%
\pgfpathlineto{\pgfqpoint{5.912472in}{0.631251in}}%
\pgfpathlineto{\pgfqpoint{5.916103in}{0.643368in}}%
\pgfpathlineto{\pgfqpoint{5.919735in}{0.643638in}}%
\pgfpathlineto{\pgfqpoint{5.923366in}{0.625317in}}%
\pgfpathlineto{\pgfqpoint{5.926997in}{0.477868in}}%
\pgfpathlineto{\pgfqpoint{5.930628in}{0.602709in}}%
\pgfpathlineto{\pgfqpoint{5.934259in}{0.668596in}}%
\pgfpathlineto{\pgfqpoint{5.937890in}{0.566128in}}%
\pgfpathlineto{\pgfqpoint{5.941522in}{0.638169in}}%
\pgfpathlineto{\pgfqpoint{5.945153in}{0.662776in}}%
\pgfpathlineto{\pgfqpoint{5.948784in}{0.659247in}}%
\pgfpathlineto{\pgfqpoint{5.952415in}{0.640930in}}%
\pgfpathlineto{\pgfqpoint{5.956046in}{0.662979in}}%
\pgfpathlineto{\pgfqpoint{5.959677in}{0.615905in}}%
\pgfpathlineto{\pgfqpoint{5.963309in}{0.557221in}}%
\pgfpathlineto{\pgfqpoint{5.966940in}{0.650400in}}%
\pgfpathlineto{\pgfqpoint{5.970571in}{0.673175in}}%
\pgfpathlineto{\pgfqpoint{5.974202in}{0.681259in}}%
\pgfpathlineto{\pgfqpoint{5.977833in}{0.627673in}}%
\pgfpathlineto{\pgfqpoint{5.981465in}{0.587035in}}%
\pgfpathlineto{\pgfqpoint{5.985096in}{0.682516in}}%
\pgfpathlineto{\pgfqpoint{5.988727in}{0.555709in}}%
\pgfpathlineto{\pgfqpoint{5.992358in}{0.623236in}}%
\pgfpathlineto{\pgfqpoint{5.995989in}{0.650666in}}%
\pgfpathlineto{\pgfqpoint{5.999620in}{0.631083in}}%
\pgfpathlineto{\pgfqpoint{6.003252in}{0.621624in}}%
\pgfpathlineto{\pgfqpoint{6.006883in}{0.672458in}}%
\pgfpathlineto{\pgfqpoint{6.010514in}{0.670577in}}%
\pgfpathlineto{\pgfqpoint{6.014145in}{0.678596in}}%
\pgfpathlineto{\pgfqpoint{6.021407in}{0.616180in}}%
\pgfpathlineto{\pgfqpoint{6.025039in}{0.547072in}}%
\pgfpathlineto{\pgfqpoint{6.028670in}{0.633883in}}%
\pgfpathlineto{\pgfqpoint{6.032301in}{0.676919in}}%
\pgfpathlineto{\pgfqpoint{6.035932in}{0.692475in}}%
\pgfpathlineto{\pgfqpoint{6.039563in}{0.564126in}}%
\pgfpathlineto{\pgfqpoint{6.043194in}{0.585104in}}%
\pgfpathlineto{\pgfqpoint{6.046826in}{0.642371in}}%
\pgfpathlineto{\pgfqpoint{6.050457in}{0.568048in}}%
\pgfpathlineto{\pgfqpoint{6.054088in}{0.654049in}}%
\pgfpathlineto{\pgfqpoint{6.057719in}{0.676789in}}%
\pgfpathlineto{\pgfqpoint{6.061350in}{0.635413in}}%
\pgfpathlineto{\pgfqpoint{6.064981in}{0.619034in}}%
\pgfpathlineto{\pgfqpoint{6.068613in}{0.647129in}}%
\pgfpathlineto{\pgfqpoint{6.072244in}{0.638970in}}%
\pgfpathlineto{\pgfqpoint{6.075875in}{0.589857in}}%
\pgfpathlineto{\pgfqpoint{6.079506in}{0.600175in}}%
\pgfpathlineto{\pgfqpoint{6.083137in}{0.614444in}}%
\pgfpathlineto{\pgfqpoint{6.086769in}{0.664129in}}%
\pgfpathlineto{\pgfqpoint{6.090400in}{0.597677in}}%
\pgfpathlineto{\pgfqpoint{6.094031in}{0.641408in}}%
\pgfpathlineto{\pgfqpoint{6.097662in}{0.659717in}}%
\pgfpathlineto{\pgfqpoint{6.101293in}{0.653831in}}%
\pgfpathlineto{\pgfqpoint{6.104924in}{0.630393in}}%
\pgfpathlineto{\pgfqpoint{6.108556in}{0.666713in}}%
\pgfpathlineto{\pgfqpoint{6.112187in}{0.605035in}}%
\pgfpathlineto{\pgfqpoint{6.115818in}{0.596514in}}%
\pgfpathlineto{\pgfqpoint{6.119449in}{0.655060in}}%
\pgfpathlineto{\pgfqpoint{6.123080in}{0.673720in}}%
\pgfpathlineto{\pgfqpoint{6.126711in}{0.678495in}}%
\pgfpathlineto{\pgfqpoint{6.130343in}{0.621057in}}%
\pgfpathlineto{\pgfqpoint{6.133974in}{0.621313in}}%
\pgfpathlineto{\pgfqpoint{6.137605in}{0.626773in}}%
\pgfpathlineto{\pgfqpoint{6.141236in}{0.598488in}}%
\pgfpathlineto{\pgfqpoint{6.144867in}{0.666323in}}%
\pgfpathlineto{\pgfqpoint{6.148498in}{0.686797in}}%
\pgfpathlineto{\pgfqpoint{6.152130in}{0.618632in}}%
\pgfpathlineto{\pgfqpoint{6.155761in}{0.525056in}}%
\pgfpathlineto{\pgfqpoint{6.159392in}{0.656075in}}%
\pgfpathlineto{\pgfqpoint{6.163023in}{0.623636in}}%
\pgfpathlineto{\pgfqpoint{6.166654in}{0.613894in}}%
\pgfpathlineto{\pgfqpoint{6.170285in}{0.661273in}}%
\pgfpathlineto{\pgfqpoint{6.173917in}{0.679575in}}%
\pgfpathlineto{\pgfqpoint{6.177548in}{0.692814in}}%
\pgfpathlineto{\pgfqpoint{6.181179in}{0.602731in}}%
\pgfpathlineto{\pgfqpoint{6.184810in}{0.611399in}}%
\pgfpathlineto{\pgfqpoint{6.188441in}{0.635782in}}%
\pgfpathlineto{\pgfqpoint{6.192073in}{0.613281in}}%
\pgfpathlineto{\pgfqpoint{6.195704in}{0.659257in}}%
\pgfpathlineto{\pgfqpoint{6.199335in}{0.676399in}}%
\pgfpathlineto{\pgfqpoint{6.202966in}{0.669464in}}%
\pgfpathlineto{\pgfqpoint{6.206597in}{0.639894in}}%
\pgfpathlineto{\pgfqpoint{6.210228in}{0.677464in}}%
\pgfpathlineto{\pgfqpoint{6.213860in}{0.622181in}}%
\pgfpathlineto{\pgfqpoint{6.217491in}{0.586609in}}%
\pgfpathlineto{\pgfqpoint{6.221122in}{0.649882in}}%
\pgfpathlineto{\pgfqpoint{6.224753in}{0.636994in}}%
\pgfpathlineto{\pgfqpoint{6.228384in}{0.499584in}}%
\pgfpathlineto{\pgfqpoint{6.232015in}{0.626741in}}%
\pgfpathlineto{\pgfqpoint{6.235647in}{0.619641in}}%
\pgfpathlineto{\pgfqpoint{6.239278in}{0.658755in}}%
\pgfpathlineto{\pgfqpoint{6.242909in}{0.616544in}}%
\pgfpathlineto{\pgfqpoint{6.246540in}{0.642168in}}%
\pgfpathlineto{\pgfqpoint{6.250171in}{0.651908in}}%
\pgfpathlineto{\pgfqpoint{6.253802in}{0.644300in}}%
\pgfpathlineto{\pgfqpoint{6.257434in}{0.614145in}}%
\pgfpathlineto{\pgfqpoint{6.261065in}{0.669333in}}%
\pgfpathlineto{\pgfqpoint{6.264696in}{0.585379in}}%
\pgfpathlineto{\pgfqpoint{6.268327in}{0.616052in}}%
\pgfpathlineto{\pgfqpoint{6.271958in}{0.658482in}}%
\pgfpathlineto{\pgfqpoint{6.275590in}{0.673554in}}%
\pgfpathlineto{\pgfqpoint{6.279221in}{0.674832in}}%
\pgfpathlineto{\pgfqpoint{6.282852in}{0.636013in}}%
\pgfpathlineto{\pgfqpoint{6.286483in}{0.620032in}}%
\pgfpathlineto{\pgfqpoint{6.290114in}{0.606920in}}%
\pgfpathlineto{\pgfqpoint{6.293745in}{0.619670in}}%
\pgfpathlineto{\pgfqpoint{6.297377in}{0.671079in}}%
\pgfpathlineto{\pgfqpoint{6.301008in}{0.688777in}}%
\pgfpathlineto{\pgfqpoint{6.304639in}{0.601841in}}%
\pgfpathlineto{\pgfqpoint{6.308270in}{0.573993in}}%
\pgfpathlineto{\pgfqpoint{6.311901in}{0.653040in}}%
\pgfpathlineto{\pgfqpoint{6.315532in}{0.609015in}}%
\pgfpathlineto{\pgfqpoint{6.319164in}{0.635563in}}%
\pgfpathlineto{\pgfqpoint{6.322795in}{0.669291in}}%
\pgfpathlineto{\pgfqpoint{6.326426in}{0.639761in}}%
\pgfpathlineto{\pgfqpoint{6.330057in}{0.632757in}}%
\pgfpathlineto{\pgfqpoint{6.333688in}{0.643748in}}%
\pgfpathlineto{\pgfqpoint{6.337319in}{0.644152in}}%
\pgfpathlineto{\pgfqpoint{6.340951in}{0.627745in}}%
\pgfpathlineto{\pgfqpoint{6.344582in}{0.524965in}}%
\pgfpathlineto{\pgfqpoint{6.348213in}{0.661500in}}%
\pgfpathlineto{\pgfqpoint{6.351844in}{0.696492in}}%
\pgfpathlineto{\pgfqpoint{6.359106in}{0.564215in}}%
\pgfpathlineto{\pgfqpoint{6.362738in}{0.658071in}}%
\pgfpathlineto{\pgfqpoint{6.366369in}{0.574281in}}%
\pgfpathlineto{\pgfqpoint{6.370000in}{0.611745in}}%
\pgfpathlineto{\pgfqpoint{6.370000in}{0.611745in}}%
\pgfusepath{stroke}%
\end{pgfscope}%
\begin{pgfscope}%
\pgfpathrectangle{\pgfqpoint{4.554414in}{0.320625in}}{\pgfqpoint{1.815586in}{1.164375in}}%
\pgfusepath{clip}%
\pgfsetrectcap%
\pgfsetroundjoin%
\pgfsetlinewidth{0.602250pt}%
\definecolor{currentstroke}{rgb}{0.478431,0.243137,0.615686}%
\pgfsetstrokecolor{currentstroke}%
\pgfsetstrokeopacity{0.120000}%
\pgfsetdash{}{0pt}%
\pgfpathmoveto{\pgfqpoint{4.554414in}{1.453975in}}%
\pgfpathlineto{\pgfqpoint{4.558045in}{1.443874in}}%
\pgfpathlineto{\pgfqpoint{4.583463in}{1.392225in}}%
\pgfpathlineto{\pgfqpoint{4.605250in}{1.343875in}}%
\pgfpathlineto{\pgfqpoint{4.659718in}{1.222249in}}%
\pgfpathlineto{\pgfqpoint{4.696029in}{1.144263in}}%
\pgfpathlineto{\pgfqpoint{4.735972in}{1.061112in}}%
\pgfpathlineto{\pgfqpoint{4.783177in}{0.965499in}}%
\pgfpathlineto{\pgfqpoint{4.837645in}{0.857778in}}%
\pgfpathlineto{\pgfqpoint{4.888481in}{0.759172in}}%
\pgfpathlineto{\pgfqpoint{4.903006in}{0.732588in}}%
\pgfpathlineto{\pgfqpoint{4.910269in}{0.717334in}}%
\pgfpathlineto{\pgfqpoint{4.913900in}{0.713820in}}%
\pgfpathlineto{\pgfqpoint{4.917531in}{0.705762in}}%
\pgfpathlineto{\pgfqpoint{4.928424in}{0.686002in}}%
\pgfpathlineto{\pgfqpoint{4.932056in}{0.678589in}}%
\pgfpathlineto{\pgfqpoint{4.935687in}{0.669393in}}%
\pgfpathlineto{\pgfqpoint{4.939318in}{0.680567in}}%
\pgfpathlineto{\pgfqpoint{4.942949in}{0.699498in}}%
\pgfpathlineto{\pgfqpoint{4.946580in}{0.668001in}}%
\pgfpathlineto{\pgfqpoint{4.950211in}{0.646371in}}%
\pgfpathlineto{\pgfqpoint{4.953843in}{0.677184in}}%
\pgfpathlineto{\pgfqpoint{4.957474in}{0.637643in}}%
\pgfpathlineto{\pgfqpoint{4.961105in}{0.628805in}}%
\pgfpathlineto{\pgfqpoint{4.964736in}{0.666350in}}%
\pgfpathlineto{\pgfqpoint{4.968367in}{0.679297in}}%
\pgfpathlineto{\pgfqpoint{4.971998in}{0.683215in}}%
\pgfpathlineto{\pgfqpoint{4.975630in}{0.597190in}}%
\pgfpathlineto{\pgfqpoint{4.979261in}{0.629317in}}%
\pgfpathlineto{\pgfqpoint{4.982892in}{0.682565in}}%
\pgfpathlineto{\pgfqpoint{4.986523in}{0.590363in}}%
\pgfpathlineto{\pgfqpoint{4.990154in}{0.637968in}}%
\pgfpathlineto{\pgfqpoint{4.993785in}{0.663850in}}%
\pgfpathlineto{\pgfqpoint{4.997417in}{0.646114in}}%
\pgfpathlineto{\pgfqpoint{5.001048in}{0.593964in}}%
\pgfpathlineto{\pgfqpoint{5.004679in}{0.675564in}}%
\pgfpathlineto{\pgfqpoint{5.008310in}{0.669970in}}%
\pgfpathlineto{\pgfqpoint{5.011941in}{0.677573in}}%
\pgfpathlineto{\pgfqpoint{5.015573in}{0.648848in}}%
\pgfpathlineto{\pgfqpoint{5.019204in}{0.632476in}}%
\pgfpathlineto{\pgfqpoint{5.022835in}{0.611065in}}%
\pgfpathlineto{\pgfqpoint{5.026466in}{0.620278in}}%
\pgfpathlineto{\pgfqpoint{5.030097in}{0.677572in}}%
\pgfpathlineto{\pgfqpoint{5.033728in}{0.697597in}}%
\pgfpathlineto{\pgfqpoint{5.040991in}{0.654883in}}%
\pgfpathlineto{\pgfqpoint{5.044622in}{0.677938in}}%
\pgfpathlineto{\pgfqpoint{5.048253in}{0.606011in}}%
\pgfpathlineto{\pgfqpoint{5.051884in}{0.592780in}}%
\pgfpathlineto{\pgfqpoint{5.055515in}{0.659197in}}%
\pgfpathlineto{\pgfqpoint{5.059147in}{0.672961in}}%
\pgfpathlineto{\pgfqpoint{5.062778in}{0.674483in}}%
\pgfpathlineto{\pgfqpoint{5.066409in}{0.625918in}}%
\pgfpathlineto{\pgfqpoint{5.070040in}{0.636691in}}%
\pgfpathlineto{\pgfqpoint{5.073671in}{0.650368in}}%
\pgfpathlineto{\pgfqpoint{5.077302in}{0.546429in}}%
\pgfpathlineto{\pgfqpoint{5.080934in}{0.662467in}}%
\pgfpathlineto{\pgfqpoint{5.084565in}{0.689445in}}%
\pgfpathlineto{\pgfqpoint{5.088196in}{0.674490in}}%
\pgfpathlineto{\pgfqpoint{5.091827in}{0.656525in}}%
\pgfpathlineto{\pgfqpoint{5.095458in}{0.671010in}}%
\pgfpathlineto{\pgfqpoint{5.099090in}{0.632815in}}%
\pgfpathlineto{\pgfqpoint{5.102721in}{0.587581in}}%
\pgfpathlineto{\pgfqpoint{5.106352in}{0.652347in}}%
\pgfpathlineto{\pgfqpoint{5.109983in}{0.668482in}}%
\pgfpathlineto{\pgfqpoint{5.113614in}{0.673832in}}%
\pgfpathlineto{\pgfqpoint{5.117245in}{0.626518in}}%
\pgfpathlineto{\pgfqpoint{5.120877in}{0.622368in}}%
\pgfpathlineto{\pgfqpoint{5.124508in}{0.693653in}}%
\pgfpathlineto{\pgfqpoint{5.128139in}{0.604631in}}%
\pgfpathlineto{\pgfqpoint{5.131770in}{0.591652in}}%
\pgfpathlineto{\pgfqpoint{5.135401in}{0.662882in}}%
\pgfpathlineto{\pgfqpoint{5.139032in}{0.637069in}}%
\pgfpathlineto{\pgfqpoint{5.142664in}{0.601672in}}%
\pgfpathlineto{\pgfqpoint{5.146295in}{0.663964in}}%
\pgfpathlineto{\pgfqpoint{5.153557in}{0.697694in}}%
\pgfpathlineto{\pgfqpoint{5.157188in}{0.603608in}}%
\pgfpathlineto{\pgfqpoint{5.160819in}{0.619171in}}%
\pgfpathlineto{\pgfqpoint{5.164451in}{0.645508in}}%
\pgfpathlineto{\pgfqpoint{5.168082in}{0.595398in}}%
\pgfpathlineto{\pgfqpoint{5.171713in}{0.665018in}}%
\pgfpathlineto{\pgfqpoint{5.175344in}{0.686047in}}%
\pgfpathlineto{\pgfqpoint{5.178975in}{0.683663in}}%
\pgfpathlineto{\pgfqpoint{5.182606in}{0.672629in}}%
\pgfpathlineto{\pgfqpoint{5.186238in}{0.668984in}}%
\pgfpathlineto{\pgfqpoint{5.189869in}{0.628384in}}%
\pgfpathlineto{\pgfqpoint{5.193500in}{0.616288in}}%
\pgfpathlineto{\pgfqpoint{5.197131in}{0.649051in}}%
\pgfpathlineto{\pgfqpoint{5.200762in}{0.657430in}}%
\pgfpathlineto{\pgfqpoint{5.204394in}{0.652627in}}%
\pgfpathlineto{\pgfqpoint{5.208025in}{0.622043in}}%
\pgfpathlineto{\pgfqpoint{5.211656in}{0.562681in}}%
\pgfpathlineto{\pgfqpoint{5.215287in}{0.677228in}}%
\pgfpathlineto{\pgfqpoint{5.218918in}{0.552944in}}%
\pgfpathlineto{\pgfqpoint{5.222549in}{0.638138in}}%
\pgfpathlineto{\pgfqpoint{5.226181in}{0.664679in}}%
\pgfpathlineto{\pgfqpoint{5.229812in}{0.675080in}}%
\pgfpathlineto{\pgfqpoint{5.233443in}{0.671287in}}%
\pgfpathlineto{\pgfqpoint{5.237074in}{0.656868in}}%
\pgfpathlineto{\pgfqpoint{5.240705in}{0.621540in}}%
\pgfpathlineto{\pgfqpoint{5.244336in}{0.563990in}}%
\pgfpathlineto{\pgfqpoint{5.247968in}{0.639292in}}%
\pgfpathlineto{\pgfqpoint{5.251599in}{0.681099in}}%
\pgfpathlineto{\pgfqpoint{5.255230in}{0.696660in}}%
\pgfpathlineto{\pgfqpoint{5.258861in}{0.592015in}}%
\pgfpathlineto{\pgfqpoint{5.262492in}{0.584799in}}%
\pgfpathlineto{\pgfqpoint{5.266123in}{0.655395in}}%
\pgfpathlineto{\pgfqpoint{5.269755in}{0.608324in}}%
\pgfpathlineto{\pgfqpoint{5.273386in}{0.645955in}}%
\pgfpathlineto{\pgfqpoint{5.277017in}{0.674487in}}%
\pgfpathlineto{\pgfqpoint{5.280648in}{0.643946in}}%
\pgfpathlineto{\pgfqpoint{5.284279in}{0.644040in}}%
\pgfpathlineto{\pgfqpoint{5.287910in}{0.645513in}}%
\pgfpathlineto{\pgfqpoint{5.291542in}{0.650234in}}%
\pgfpathlineto{\pgfqpoint{5.295173in}{0.643149in}}%
\pgfpathlineto{\pgfqpoint{5.302435in}{0.559154in}}%
\pgfpathlineto{\pgfqpoint{5.306066in}{0.677454in}}%
\pgfpathlineto{\pgfqpoint{5.309698in}{0.540013in}}%
\pgfpathlineto{\pgfqpoint{5.313329in}{0.635050in}}%
\pgfpathlineto{\pgfqpoint{5.316960in}{0.668634in}}%
\pgfpathlineto{\pgfqpoint{5.320591in}{0.674731in}}%
\pgfpathlineto{\pgfqpoint{5.324222in}{0.669894in}}%
\pgfpathlineto{\pgfqpoint{5.327853in}{0.653619in}}%
\pgfpathlineto{\pgfqpoint{5.335116in}{0.605805in}}%
\pgfpathlineto{\pgfqpoint{5.338747in}{0.632807in}}%
\pgfpathlineto{\pgfqpoint{5.342378in}{0.676600in}}%
\pgfpathlineto{\pgfqpoint{5.346009in}{0.694003in}}%
\pgfpathlineto{\pgfqpoint{5.349640in}{0.673748in}}%
\pgfpathlineto{\pgfqpoint{5.353272in}{0.644749in}}%
\pgfpathlineto{\pgfqpoint{5.356903in}{0.682989in}}%
\pgfpathlineto{\pgfqpoint{5.360534in}{0.607420in}}%
\pgfpathlineto{\pgfqpoint{5.364165in}{0.579468in}}%
\pgfpathlineto{\pgfqpoint{5.367796in}{0.659016in}}%
\pgfpathlineto{\pgfqpoint{5.371427in}{0.663141in}}%
\pgfpathlineto{\pgfqpoint{5.375059in}{0.648054in}}%
\pgfpathlineto{\pgfqpoint{5.378690in}{0.517187in}}%
\pgfpathlineto{\pgfqpoint{5.382321in}{0.652330in}}%
\pgfpathlineto{\pgfqpoint{5.385952in}{0.700756in}}%
\pgfpathlineto{\pgfqpoint{5.389583in}{0.526773in}}%
\pgfpathlineto{\pgfqpoint{5.393215in}{0.607756in}}%
\pgfpathlineto{\pgfqpoint{5.396846in}{0.647256in}}%
\pgfpathlineto{\pgfqpoint{5.400477in}{0.584658in}}%
\pgfpathlineto{\pgfqpoint{5.404108in}{0.658530in}}%
\pgfpathlineto{\pgfqpoint{5.407739in}{0.682147in}}%
\pgfpathlineto{\pgfqpoint{5.411370in}{0.640291in}}%
\pgfpathlineto{\pgfqpoint{5.415002in}{0.633253in}}%
\pgfpathlineto{\pgfqpoint{5.418633in}{0.650105in}}%
\pgfpathlineto{\pgfqpoint{5.422264in}{0.645086in}}%
\pgfpathlineto{\pgfqpoint{5.425895in}{0.613777in}}%
\pgfpathlineto{\pgfqpoint{5.429526in}{0.435025in}}%
\pgfpathlineto{\pgfqpoint{5.433157in}{0.589383in}}%
\pgfpathlineto{\pgfqpoint{5.436789in}{0.675052in}}%
\pgfpathlineto{\pgfqpoint{5.440420in}{0.561096in}}%
\pgfpathlineto{\pgfqpoint{5.444051in}{0.641923in}}%
\pgfpathlineto{\pgfqpoint{5.447682in}{0.667924in}}%
\pgfpathlineto{\pgfqpoint{5.451313in}{0.673026in}}%
\pgfpathlineto{\pgfqpoint{5.454944in}{0.665662in}}%
\pgfpathlineto{\pgfqpoint{5.458576in}{0.660996in}}%
\pgfpathlineto{\pgfqpoint{5.462207in}{0.623062in}}%
\pgfpathlineto{\pgfqpoint{5.465838in}{0.546934in}}%
\pgfpathlineto{\pgfqpoint{5.469469in}{0.643910in}}%
\pgfpathlineto{\pgfqpoint{5.473100in}{0.680410in}}%
\pgfpathlineto{\pgfqpoint{5.476731in}{0.695189in}}%
\pgfpathlineto{\pgfqpoint{5.480363in}{0.569549in}}%
\pgfpathlineto{\pgfqpoint{5.483994in}{0.600392in}}%
\pgfpathlineto{\pgfqpoint{5.487625in}{0.655754in}}%
\pgfpathlineto{\pgfqpoint{5.491256in}{0.601570in}}%
\pgfpathlineto{\pgfqpoint{5.494887in}{0.649140in}}%
\pgfpathlineto{\pgfqpoint{5.498519in}{0.678183in}}%
\pgfpathlineto{\pgfqpoint{5.502150in}{0.645439in}}%
\pgfpathlineto{\pgfqpoint{5.505781in}{0.637886in}}%
\pgfpathlineto{\pgfqpoint{5.509412in}{0.649487in}}%
\pgfpathlineto{\pgfqpoint{5.513043in}{0.650399in}}%
\pgfpathlineto{\pgfqpoint{5.516674in}{0.638322in}}%
\pgfpathlineto{\pgfqpoint{5.520306in}{0.557942in}}%
\pgfpathlineto{\pgfqpoint{5.523937in}{0.600327in}}%
\pgfpathlineto{\pgfqpoint{5.527568in}{0.677031in}}%
\pgfpathlineto{\pgfqpoint{5.531199in}{0.461994in}}%
\pgfpathlineto{\pgfqpoint{5.534830in}{0.638592in}}%
\pgfpathlineto{\pgfqpoint{5.538461in}{0.670043in}}%
\pgfpathlineto{\pgfqpoint{5.542093in}{0.671384in}}%
\pgfpathlineto{\pgfqpoint{5.545724in}{0.653322in}}%
\pgfpathlineto{\pgfqpoint{5.549355in}{0.646152in}}%
\pgfpathlineto{\pgfqpoint{5.552986in}{0.651091in}}%
\pgfpathlineto{\pgfqpoint{5.556617in}{0.676495in}}%
\pgfpathlineto{\pgfqpoint{5.560248in}{0.634318in}}%
\pgfpathlineto{\pgfqpoint{5.563880in}{0.637964in}}%
\pgfpathlineto{\pgfqpoint{5.567511in}{0.636615in}}%
\pgfpathlineto{\pgfqpoint{5.571142in}{0.556201in}}%
\pgfpathlineto{\pgfqpoint{5.574773in}{0.665977in}}%
\pgfpathlineto{\pgfqpoint{5.578404in}{0.695039in}}%
\pgfpathlineto{\pgfqpoint{5.582035in}{0.675908in}}%
\pgfpathlineto{\pgfqpoint{5.585667in}{0.668176in}}%
\pgfpathlineto{\pgfqpoint{5.589298in}{0.666666in}}%
\pgfpathlineto{\pgfqpoint{5.592929in}{0.615054in}}%
\pgfpathlineto{\pgfqpoint{5.596560in}{0.580725in}}%
\pgfpathlineto{\pgfqpoint{5.600191in}{0.653117in}}%
\pgfpathlineto{\pgfqpoint{5.603823in}{0.684416in}}%
\pgfpathlineto{\pgfqpoint{5.607454in}{0.697234in}}%
\pgfpathlineto{\pgfqpoint{5.611085in}{0.548560in}}%
\pgfpathlineto{\pgfqpoint{5.614716in}{0.605845in}}%
\pgfpathlineto{\pgfqpoint{5.618347in}{0.644025in}}%
\pgfpathlineto{\pgfqpoint{5.621978in}{0.524354in}}%
\pgfpathlineto{\pgfqpoint{5.625610in}{0.663861in}}%
\pgfpathlineto{\pgfqpoint{5.629241in}{0.685126in}}%
\pgfpathlineto{\pgfqpoint{5.632872in}{0.641330in}}%
\pgfpathlineto{\pgfqpoint{5.636503in}{0.623236in}}%
\pgfpathlineto{\pgfqpoint{5.640134in}{0.653447in}}%
\pgfpathlineto{\pgfqpoint{5.643765in}{0.645999in}}%
\pgfpathlineto{\pgfqpoint{5.647397in}{0.608933in}}%
\pgfpathlineto{\pgfqpoint{5.651028in}{0.582026in}}%
\pgfpathlineto{\pgfqpoint{5.654659in}{0.612010in}}%
\pgfpathlineto{\pgfqpoint{5.658290in}{0.673531in}}%
\pgfpathlineto{\pgfqpoint{5.661921in}{0.580135in}}%
\pgfpathlineto{\pgfqpoint{5.665552in}{0.644263in}}%
\pgfpathlineto{\pgfqpoint{5.669184in}{0.668426in}}%
\pgfpathlineto{\pgfqpoint{5.672815in}{0.668639in}}%
\pgfpathlineto{\pgfqpoint{5.676446in}{0.656322in}}%
\pgfpathlineto{\pgfqpoint{5.680077in}{0.665368in}}%
\pgfpathlineto{\pgfqpoint{5.683708in}{0.623745in}}%
\pgfpathlineto{\pgfqpoint{5.687340in}{0.470156in}}%
\pgfpathlineto{\pgfqpoint{5.690971in}{0.649019in}}%
\pgfpathlineto{\pgfqpoint{5.694602in}{0.679912in}}%
\pgfpathlineto{\pgfqpoint{5.698233in}{0.693252in}}%
\pgfpathlineto{\pgfqpoint{5.701864in}{0.456184in}}%
\pgfpathlineto{\pgfqpoint{5.705495in}{0.611923in}}%
\pgfpathlineto{\pgfqpoint{5.709127in}{0.654888in}}%
\pgfpathlineto{\pgfqpoint{5.712758in}{0.587735in}}%
\pgfpathlineto{\pgfqpoint{5.716389in}{0.653515in}}%
\pgfpathlineto{\pgfqpoint{5.720020in}{0.681869in}}%
\pgfpathlineto{\pgfqpoint{5.723651in}{0.687607in}}%
\pgfpathlineto{\pgfqpoint{5.727282in}{0.691725in}}%
\pgfpathlineto{\pgfqpoint{5.730914in}{0.644495in}}%
\pgfpathlineto{\pgfqpoint{5.734545in}{0.622514in}}%
\pgfpathlineto{\pgfqpoint{5.738176in}{0.626767in}}%
\pgfpathlineto{\pgfqpoint{5.741807in}{0.633730in}}%
\pgfpathlineto{\pgfqpoint{5.745438in}{0.667292in}}%
\pgfpathlineto{\pgfqpoint{5.749069in}{0.681335in}}%
\pgfpathlineto{\pgfqpoint{5.752701in}{0.672088in}}%
\pgfpathlineto{\pgfqpoint{5.756332in}{0.636486in}}%
\pgfpathlineto{\pgfqpoint{5.759963in}{0.686381in}}%
\pgfpathlineto{\pgfqpoint{5.763594in}{0.625959in}}%
\pgfpathlineto{\pgfqpoint{5.767225in}{0.589742in}}%
\pgfpathlineto{\pgfqpoint{5.770856in}{0.655482in}}%
\pgfpathlineto{\pgfqpoint{5.774488in}{0.643332in}}%
\pgfpathlineto{\pgfqpoint{5.778119in}{0.552931in}}%
\pgfpathlineto{\pgfqpoint{5.781750in}{0.619119in}}%
\pgfpathlineto{\pgfqpoint{5.785381in}{0.615727in}}%
\pgfpathlineto{\pgfqpoint{5.789012in}{0.668812in}}%
\pgfpathlineto{\pgfqpoint{5.792644in}{0.610489in}}%
\pgfpathlineto{\pgfqpoint{5.796275in}{0.646498in}}%
\pgfpathlineto{\pgfqpoint{5.799906in}{0.663078in}}%
\pgfpathlineto{\pgfqpoint{5.803537in}{0.663774in}}%
\pgfpathlineto{\pgfqpoint{5.807168in}{0.649144in}}%
\pgfpathlineto{\pgfqpoint{5.810799in}{0.670071in}}%
\pgfpathlineto{\pgfqpoint{5.814431in}{0.611960in}}%
\pgfpathlineto{\pgfqpoint{5.818062in}{0.587914in}}%
\pgfpathlineto{\pgfqpoint{5.821693in}{0.655639in}}%
\pgfpathlineto{\pgfqpoint{5.825324in}{0.682389in}}%
\pgfpathlineto{\pgfqpoint{5.828955in}{0.692958in}}%
\pgfpathlineto{\pgfqpoint{5.832586in}{0.584759in}}%
\pgfpathlineto{\pgfqpoint{5.839849in}{0.644570in}}%
\pgfpathlineto{\pgfqpoint{5.843480in}{0.436403in}}%
\pgfpathlineto{\pgfqpoint{5.847111in}{0.664437in}}%
\pgfpathlineto{\pgfqpoint{5.850742in}{0.687061in}}%
\pgfpathlineto{\pgfqpoint{5.854373in}{0.641113in}}%
\pgfpathlineto{\pgfqpoint{5.858005in}{0.613946in}}%
\pgfpathlineto{\pgfqpoint{5.861636in}{0.656256in}}%
\pgfpathlineto{\pgfqpoint{5.865267in}{0.645395in}}%
\pgfpathlineto{\pgfqpoint{5.868898in}{0.596915in}}%
\pgfpathlineto{\pgfqpoint{5.872529in}{0.612915in}}%
\pgfpathlineto{\pgfqpoint{5.876160in}{0.625658in}}%
\pgfpathlineto{\pgfqpoint{5.879792in}{0.671574in}}%
\pgfpathlineto{\pgfqpoint{5.883423in}{0.593062in}}%
\pgfpathlineto{\pgfqpoint{5.887054in}{0.646417in}}%
\pgfpathlineto{\pgfqpoint{5.890685in}{0.668850in}}%
\pgfpathlineto{\pgfqpoint{5.894316in}{0.663940in}}%
\pgfpathlineto{\pgfqpoint{5.897948in}{0.645589in}}%
\pgfpathlineto{\pgfqpoint{5.901579in}{0.668870in}}%
\pgfpathlineto{\pgfqpoint{5.905210in}{0.623910in}}%
\pgfpathlineto{\pgfqpoint{5.908841in}{0.526818in}}%
\pgfpathlineto{\pgfqpoint{5.912472in}{0.653324in}}%
\pgfpathlineto{\pgfqpoint{5.916103in}{0.679101in}}%
\pgfpathlineto{\pgfqpoint{5.919735in}{0.690426in}}%
\pgfpathlineto{\pgfqpoint{5.923366in}{0.650864in}}%
\pgfpathlineto{\pgfqpoint{5.926997in}{0.529035in}}%
\pgfpathlineto{\pgfqpoint{5.930628in}{0.688540in}}%
\pgfpathlineto{\pgfqpoint{5.934259in}{0.552856in}}%
\pgfpathlineto{\pgfqpoint{5.937890in}{0.618266in}}%
\pgfpathlineto{\pgfqpoint{5.941522in}{0.661158in}}%
\pgfpathlineto{\pgfqpoint{5.945153in}{0.653104in}}%
\pgfpathlineto{\pgfqpoint{5.948784in}{0.569735in}}%
\pgfpathlineto{\pgfqpoint{5.952415in}{0.651987in}}%
\pgfpathlineto{\pgfqpoint{5.956046in}{0.670048in}}%
\pgfpathlineto{\pgfqpoint{5.959677in}{0.696167in}}%
\pgfpathlineto{\pgfqpoint{5.963309in}{0.622628in}}%
\pgfpathlineto{\pgfqpoint{5.966940in}{0.621968in}}%
\pgfpathlineto{\pgfqpoint{5.970571in}{0.629614in}}%
\pgfpathlineto{\pgfqpoint{5.974202in}{0.581820in}}%
\pgfpathlineto{\pgfqpoint{5.977833in}{0.672856in}}%
\pgfpathlineto{\pgfqpoint{5.981465in}{0.692773in}}%
\pgfpathlineto{\pgfqpoint{5.985096in}{0.633445in}}%
\pgfpathlineto{\pgfqpoint{5.988727in}{0.592656in}}%
\pgfpathlineto{\pgfqpoint{5.992358in}{0.656466in}}%
\pgfpathlineto{\pgfqpoint{5.995989in}{0.635729in}}%
\pgfpathlineto{\pgfqpoint{5.999620in}{0.579056in}}%
\pgfpathlineto{\pgfqpoint{6.003252in}{0.640125in}}%
\pgfpathlineto{\pgfqpoint{6.006883in}{0.630365in}}%
\pgfpathlineto{\pgfqpoint{6.010514in}{0.665489in}}%
\pgfpathlineto{\pgfqpoint{6.014145in}{0.617195in}}%
\pgfpathlineto{\pgfqpoint{6.017776in}{0.648876in}}%
\pgfpathlineto{\pgfqpoint{6.021407in}{0.663836in}}%
\pgfpathlineto{\pgfqpoint{6.025039in}{0.658531in}}%
\pgfpathlineto{\pgfqpoint{6.028670in}{0.636958in}}%
\pgfpathlineto{\pgfqpoint{6.032301in}{0.672744in}}%
\pgfpathlineto{\pgfqpoint{6.035932in}{0.610587in}}%
\pgfpathlineto{\pgfqpoint{6.039563in}{0.598202in}}%
\pgfpathlineto{\pgfqpoint{6.043194in}{0.658976in}}%
\pgfpathlineto{\pgfqpoint{6.046826in}{0.680973in}}%
\pgfpathlineto{\pgfqpoint{6.050457in}{0.689005in}}%
\pgfpathlineto{\pgfqpoint{6.054088in}{0.607064in}}%
\pgfpathlineto{\pgfqpoint{6.057719in}{0.622207in}}%
\pgfpathlineto{\pgfqpoint{6.061350in}{0.642797in}}%
\pgfpathlineto{\pgfqpoint{6.064981in}{0.546879in}}%
\pgfpathlineto{\pgfqpoint{6.068613in}{0.666401in}}%
\pgfpathlineto{\pgfqpoint{6.072244in}{0.689086in}}%
\pgfpathlineto{\pgfqpoint{6.075875in}{0.639215in}}%
\pgfpathlineto{\pgfqpoint{6.079506in}{0.598263in}}%
\pgfpathlineto{\pgfqpoint{6.083137in}{0.659061in}}%
\pgfpathlineto{\pgfqpoint{6.086769in}{0.643994in}}%
\pgfpathlineto{\pgfqpoint{6.090400in}{0.570811in}}%
\pgfpathlineto{\pgfqpoint{6.094031in}{0.632317in}}%
\pgfpathlineto{\pgfqpoint{6.097662in}{0.678201in}}%
\pgfpathlineto{\pgfqpoint{6.101293in}{0.704868in}}%
\pgfpathlineto{\pgfqpoint{6.108556in}{0.527820in}}%
\pgfpathlineto{\pgfqpoint{6.112187in}{0.661128in}}%
\pgfpathlineto{\pgfqpoint{6.115818in}{0.532198in}}%
\pgfpathlineto{\pgfqpoint{6.119449in}{0.651492in}}%
\pgfpathlineto{\pgfqpoint{6.123080in}{0.677383in}}%
\pgfpathlineto{\pgfqpoint{6.126711in}{0.689911in}}%
\pgfpathlineto{\pgfqpoint{6.130343in}{0.692837in}}%
\pgfpathlineto{\pgfqpoint{6.133974in}{0.619255in}}%
\pgfpathlineto{\pgfqpoint{6.137605in}{0.616999in}}%
\pgfpathlineto{\pgfqpoint{6.141236in}{0.653960in}}%
\pgfpathlineto{\pgfqpoint{6.144867in}{0.631159in}}%
\pgfpathlineto{\pgfqpoint{6.148498in}{0.652208in}}%
\pgfpathlineto{\pgfqpoint{6.152130in}{0.660367in}}%
\pgfpathlineto{\pgfqpoint{6.155761in}{0.654243in}}%
\pgfpathlineto{\pgfqpoint{6.159392in}{0.631652in}}%
\pgfpathlineto{\pgfqpoint{6.163023in}{0.675494in}}%
\pgfpathlineto{\pgfqpoint{6.166654in}{0.585119in}}%
\pgfpathlineto{\pgfqpoint{6.170285in}{0.620515in}}%
\pgfpathlineto{\pgfqpoint{6.173917in}{0.662157in}}%
\pgfpathlineto{\pgfqpoint{6.177548in}{0.681809in}}%
\pgfpathlineto{\pgfqpoint{6.181179in}{0.686879in}}%
\pgfpathlineto{\pgfqpoint{6.184810in}{0.628789in}}%
\pgfpathlineto{\pgfqpoint{6.188441in}{0.621519in}}%
\pgfpathlineto{\pgfqpoint{6.192073in}{0.626245in}}%
\pgfpathlineto{\pgfqpoint{6.195704in}{0.601506in}}%
\pgfpathlineto{\pgfqpoint{6.199335in}{0.674001in}}%
\pgfpathlineto{\pgfqpoint{6.202966in}{0.693074in}}%
\pgfpathlineto{\pgfqpoint{6.210228in}{0.566500in}}%
\pgfpathlineto{\pgfqpoint{6.213860in}{0.658823in}}%
\pgfpathlineto{\pgfqpoint{6.217491in}{0.635901in}}%
\pgfpathlineto{\pgfqpoint{6.221122in}{0.586411in}}%
\pgfpathlineto{\pgfqpoint{6.224753in}{0.648675in}}%
\pgfpathlineto{\pgfqpoint{6.228384in}{0.637116in}}%
\pgfpathlineto{\pgfqpoint{6.232015in}{0.663020in}}%
\pgfpathlineto{\pgfqpoint{6.235647in}{0.623322in}}%
\pgfpathlineto{\pgfqpoint{6.239278in}{0.650075in}}%
\pgfpathlineto{\pgfqpoint{6.242909in}{0.662870in}}%
\pgfpathlineto{\pgfqpoint{6.246540in}{0.651489in}}%
\pgfpathlineto{\pgfqpoint{6.250171in}{0.619647in}}%
\pgfpathlineto{\pgfqpoint{6.253802in}{0.674799in}}%
\pgfpathlineto{\pgfqpoint{6.257434in}{0.609049in}}%
\pgfpathlineto{\pgfqpoint{6.261065in}{0.605966in}}%
\pgfpathlineto{\pgfqpoint{6.264696in}{0.661917in}}%
\pgfpathlineto{\pgfqpoint{6.268327in}{0.679375in}}%
\pgfpathlineto{\pgfqpoint{6.271958in}{0.684758in}}%
\pgfpathlineto{\pgfqpoint{6.275590in}{0.621251in}}%
\pgfpathlineto{\pgfqpoint{6.279221in}{0.627300in}}%
\pgfpathlineto{\pgfqpoint{6.282852in}{0.640932in}}%
\pgfpathlineto{\pgfqpoint{6.286483in}{0.577428in}}%
\pgfpathlineto{\pgfqpoint{6.290114in}{0.667835in}}%
\pgfpathlineto{\pgfqpoint{6.293745in}{0.690413in}}%
\pgfpathlineto{\pgfqpoint{6.297377in}{0.683621in}}%
\pgfpathlineto{\pgfqpoint{6.301008in}{0.674332in}}%
\pgfpathlineto{\pgfqpoint{6.304639in}{0.669046in}}%
\pgfpathlineto{\pgfqpoint{6.311901in}{0.582147in}}%
\pgfpathlineto{\pgfqpoint{6.315532in}{0.650063in}}%
\pgfpathlineto{\pgfqpoint{6.319164in}{0.667574in}}%
\pgfpathlineto{\pgfqpoint{6.322795in}{0.672064in}}%
\pgfpathlineto{\pgfqpoint{6.326426in}{0.645966in}}%
\pgfpathlineto{\pgfqpoint{6.330057in}{0.590282in}}%
\pgfpathlineto{\pgfqpoint{6.333688in}{0.696593in}}%
\pgfpathlineto{\pgfqpoint{6.337319in}{0.605860in}}%
\pgfpathlineto{\pgfqpoint{6.340951in}{0.554632in}}%
\pgfpathlineto{\pgfqpoint{6.344582in}{0.656776in}}%
\pgfpathlineto{\pgfqpoint{6.348213in}{0.628549in}}%
\pgfpathlineto{\pgfqpoint{6.351844in}{0.624048in}}%
\pgfpathlineto{\pgfqpoint{6.355475in}{0.661464in}}%
\pgfpathlineto{\pgfqpoint{6.359106in}{0.636618in}}%
\pgfpathlineto{\pgfqpoint{6.362738in}{0.655712in}}%
\pgfpathlineto{\pgfqpoint{6.366369in}{0.635702in}}%
\pgfpathlineto{\pgfqpoint{6.370000in}{0.592188in}}%
\pgfpathlineto{\pgfqpoint{6.370000in}{0.592188in}}%
\pgfusepath{stroke}%
\end{pgfscope}%
\begin{pgfscope}%
\pgfpathrectangle{\pgfqpoint{4.554414in}{0.320625in}}{\pgfqpoint{1.815586in}{1.164375in}}%
\pgfusepath{clip}%
\pgfsetrectcap%
\pgfsetroundjoin%
\pgfsetlinewidth{0.602250pt}%
\definecolor{currentstroke}{rgb}{0.478431,0.243137,0.615686}%
\pgfsetstrokecolor{currentstroke}%
\pgfsetstrokeopacity{0.120000}%
\pgfsetdash{}{0pt}%
\pgfpathmoveto{\pgfqpoint{4.554414in}{1.453884in}}%
\pgfpathlineto{\pgfqpoint{4.558045in}{1.443573in}}%
\pgfpathlineto{\pgfqpoint{4.583463in}{1.391673in}}%
\pgfpathlineto{\pgfqpoint{4.608881in}{1.334963in}}%
\pgfpathlineto{\pgfqpoint{4.656086in}{1.229353in}}%
\pgfpathlineto{\pgfqpoint{4.692398in}{1.151022in}}%
\pgfpathlineto{\pgfqpoint{4.732341in}{1.067630in}}%
\pgfpathlineto{\pgfqpoint{4.775915in}{0.979243in}}%
\pgfpathlineto{\pgfqpoint{4.826752in}{0.878653in}}%
\pgfpathlineto{\pgfqpoint{4.884850in}{0.766121in}}%
\pgfpathlineto{\pgfqpoint{4.910269in}{0.717740in}}%
\pgfpathlineto{\pgfqpoint{4.913900in}{0.710544in}}%
\pgfpathlineto{\pgfqpoint{4.917531in}{0.707298in}}%
\pgfpathlineto{\pgfqpoint{4.921162in}{0.696875in}}%
\pgfpathlineto{\pgfqpoint{4.924793in}{0.690104in}}%
\pgfpathlineto{\pgfqpoint{4.928424in}{0.687879in}}%
\pgfpathlineto{\pgfqpoint{4.932056in}{0.688093in}}%
\pgfpathlineto{\pgfqpoint{4.935687in}{0.687454in}}%
\pgfpathlineto{\pgfqpoint{4.939318in}{0.664428in}}%
\pgfpathlineto{\pgfqpoint{4.942949in}{0.656177in}}%
\pgfpathlineto{\pgfqpoint{4.946580in}{0.685039in}}%
\pgfpathlineto{\pgfqpoint{4.950211in}{0.641571in}}%
\pgfpathlineto{\pgfqpoint{4.953843in}{0.641728in}}%
\pgfpathlineto{\pgfqpoint{4.957474in}{0.658700in}}%
\pgfpathlineto{\pgfqpoint{4.961105in}{0.650216in}}%
\pgfpathlineto{\pgfqpoint{4.964736in}{0.615875in}}%
\pgfpathlineto{\pgfqpoint{4.968367in}{0.645004in}}%
\pgfpathlineto{\pgfqpoint{4.971998in}{0.661620in}}%
\pgfpathlineto{\pgfqpoint{4.975630in}{0.690167in}}%
\pgfpathlineto{\pgfqpoint{4.979261in}{0.608942in}}%
\pgfpathlineto{\pgfqpoint{4.982892in}{0.615295in}}%
\pgfpathlineto{\pgfqpoint{4.986523in}{0.633703in}}%
\pgfpathlineto{\pgfqpoint{4.990154in}{0.567562in}}%
\pgfpathlineto{\pgfqpoint{4.993785in}{0.656603in}}%
\pgfpathlineto{\pgfqpoint{4.997417in}{0.678020in}}%
\pgfpathlineto{\pgfqpoint{5.001048in}{0.632845in}}%
\pgfpathlineto{\pgfqpoint{5.004679in}{0.621887in}}%
\pgfpathlineto{\pgfqpoint{5.008310in}{0.643213in}}%
\pgfpathlineto{\pgfqpoint{5.011941in}{0.637988in}}%
\pgfpathlineto{\pgfqpoint{5.015573in}{0.613711in}}%
\pgfpathlineto{\pgfqpoint{5.019204in}{0.559228in}}%
\pgfpathlineto{\pgfqpoint{5.022835in}{0.560876in}}%
\pgfpathlineto{\pgfqpoint{5.026466in}{0.668806in}}%
\pgfpathlineto{\pgfqpoint{5.030097in}{0.499150in}}%
\pgfpathlineto{\pgfqpoint{5.033728in}{0.627869in}}%
\pgfpathlineto{\pgfqpoint{5.037360in}{0.659527in}}%
\pgfpathlineto{\pgfqpoint{5.040991in}{0.669876in}}%
\pgfpathlineto{\pgfqpoint{5.044622in}{0.669668in}}%
\pgfpathlineto{\pgfqpoint{5.048253in}{0.637153in}}%
\pgfpathlineto{\pgfqpoint{5.051884in}{0.618369in}}%
\pgfpathlineto{\pgfqpoint{5.059147in}{0.605618in}}%
\pgfpathlineto{\pgfqpoint{5.062778in}{0.666194in}}%
\pgfpathlineto{\pgfqpoint{5.066409in}{0.685181in}}%
\pgfpathlineto{\pgfqpoint{5.073671in}{0.555950in}}%
\pgfpathlineto{\pgfqpoint{5.077302in}{0.652167in}}%
\pgfpathlineto{\pgfqpoint{5.080934in}{0.632463in}}%
\pgfpathlineto{\pgfqpoint{5.084565in}{0.445021in}}%
\pgfpathlineto{\pgfqpoint{5.088196in}{0.628553in}}%
\pgfpathlineto{\pgfqpoint{5.091827in}{0.625945in}}%
\pgfpathlineto{\pgfqpoint{5.095458in}{0.661431in}}%
\pgfpathlineto{\pgfqpoint{5.099090in}{0.592869in}}%
\pgfpathlineto{\pgfqpoint{5.102721in}{0.639209in}}%
\pgfpathlineto{\pgfqpoint{5.106352in}{0.661397in}}%
\pgfpathlineto{\pgfqpoint{5.109983in}{0.656766in}}%
\pgfpathlineto{\pgfqpoint{5.113614in}{0.620048in}}%
\pgfpathlineto{\pgfqpoint{5.117245in}{0.655533in}}%
\pgfpathlineto{\pgfqpoint{5.120877in}{0.646706in}}%
\pgfpathlineto{\pgfqpoint{5.124508in}{0.664955in}}%
\pgfpathlineto{\pgfqpoint{5.128139in}{0.631923in}}%
\pgfpathlineto{\pgfqpoint{5.131770in}{0.629689in}}%
\pgfpathlineto{\pgfqpoint{5.135401in}{0.625353in}}%
\pgfpathlineto{\pgfqpoint{5.139032in}{0.547186in}}%
\pgfpathlineto{\pgfqpoint{5.142664in}{0.657121in}}%
\pgfpathlineto{\pgfqpoint{5.146295in}{0.685567in}}%
\pgfpathlineto{\pgfqpoint{5.149926in}{0.669337in}}%
\pgfpathlineto{\pgfqpoint{5.153557in}{0.667391in}}%
\pgfpathlineto{\pgfqpoint{5.157188in}{0.653921in}}%
\pgfpathlineto{\pgfqpoint{5.160819in}{0.611699in}}%
\pgfpathlineto{\pgfqpoint{5.164451in}{0.445351in}}%
\pgfpathlineto{\pgfqpoint{5.168082in}{0.633927in}}%
\pgfpathlineto{\pgfqpoint{5.171713in}{0.675134in}}%
\pgfpathlineto{\pgfqpoint{5.175344in}{0.690802in}}%
\pgfpathlineto{\pgfqpoint{5.178975in}{0.591732in}}%
\pgfpathlineto{\pgfqpoint{5.182606in}{0.550339in}}%
\pgfpathlineto{\pgfqpoint{5.186238in}{0.646663in}}%
\pgfpathlineto{\pgfqpoint{5.189869in}{0.607588in}}%
\pgfpathlineto{\pgfqpoint{5.193500in}{0.629690in}}%
\pgfpathlineto{\pgfqpoint{5.197131in}{0.656584in}}%
\pgfpathlineto{\pgfqpoint{5.200762in}{0.631190in}}%
\pgfpathlineto{\pgfqpoint{5.204394in}{0.648817in}}%
\pgfpathlineto{\pgfqpoint{5.208025in}{0.623873in}}%
\pgfpathlineto{\pgfqpoint{5.211656in}{0.642143in}}%
\pgfpathlineto{\pgfqpoint{5.215287in}{0.651316in}}%
\pgfpathlineto{\pgfqpoint{5.218918in}{0.644348in}}%
\pgfpathlineto{\pgfqpoint{5.222549in}{0.620966in}}%
\pgfpathlineto{\pgfqpoint{5.226181in}{0.665934in}}%
\pgfpathlineto{\pgfqpoint{5.229812in}{0.600284in}}%
\pgfpathlineto{\pgfqpoint{5.233443in}{0.590845in}}%
\pgfpathlineto{\pgfqpoint{5.237074in}{0.650065in}}%
\pgfpathlineto{\pgfqpoint{5.240705in}{0.673550in}}%
\pgfpathlineto{\pgfqpoint{5.244336in}{0.683709in}}%
\pgfpathlineto{\pgfqpoint{5.247968in}{0.571484in}}%
\pgfpathlineto{\pgfqpoint{5.255230in}{0.641986in}}%
\pgfpathlineto{\pgfqpoint{5.258861in}{0.560477in}}%
\pgfpathlineto{\pgfqpoint{5.262492in}{0.649070in}}%
\pgfpathlineto{\pgfqpoint{5.266123in}{0.674154in}}%
\pgfpathlineto{\pgfqpoint{5.269755in}{0.637100in}}%
\pgfpathlineto{\pgfqpoint{5.273386in}{0.626536in}}%
\pgfpathlineto{\pgfqpoint{5.277017in}{0.642293in}}%
\pgfpathlineto{\pgfqpoint{5.280648in}{0.642665in}}%
\pgfpathlineto{\pgfqpoint{5.284279in}{0.633745in}}%
\pgfpathlineto{\pgfqpoint{5.287910in}{0.583317in}}%
\pgfpathlineto{\pgfqpoint{5.291542in}{0.569960in}}%
\pgfpathlineto{\pgfqpoint{5.295173in}{0.669997in}}%
\pgfpathlineto{\pgfqpoint{5.298804in}{0.562394in}}%
\pgfpathlineto{\pgfqpoint{5.302435in}{0.622071in}}%
\pgfpathlineto{\pgfqpoint{5.306066in}{0.660153in}}%
\pgfpathlineto{\pgfqpoint{5.309698in}{0.668218in}}%
\pgfpathlineto{\pgfqpoint{5.313329in}{0.665114in}}%
\pgfpathlineto{\pgfqpoint{5.316960in}{0.544071in}}%
\pgfpathlineto{\pgfqpoint{5.320591in}{0.624953in}}%
\pgfpathlineto{\pgfqpoint{5.324222in}{0.678383in}}%
\pgfpathlineto{\pgfqpoint{5.327853in}{0.585943in}}%
\pgfpathlineto{\pgfqpoint{5.331485in}{0.624289in}}%
\pgfpathlineto{\pgfqpoint{5.335116in}{0.648715in}}%
\pgfpathlineto{\pgfqpoint{5.338747in}{0.630952in}}%
\pgfpathlineto{\pgfqpoint{5.342378in}{0.605616in}}%
\pgfpathlineto{\pgfqpoint{5.346009in}{0.663237in}}%
\pgfpathlineto{\pgfqpoint{5.349640in}{0.666629in}}%
\pgfpathlineto{\pgfqpoint{5.353272in}{0.683946in}}%
\pgfpathlineto{\pgfqpoint{5.356903in}{0.627639in}}%
\pgfpathlineto{\pgfqpoint{5.360534in}{0.615782in}}%
\pgfpathlineto{\pgfqpoint{5.364165in}{0.616059in}}%
\pgfpathlineto{\pgfqpoint{5.367796in}{0.584768in}}%
\pgfpathlineto{\pgfqpoint{5.371427in}{0.665534in}}%
\pgfpathlineto{\pgfqpoint{5.375059in}{0.684729in}}%
\pgfpathlineto{\pgfqpoint{5.378690in}{0.626992in}}%
\pgfpathlineto{\pgfqpoint{5.382321in}{0.595043in}}%
\pgfpathlineto{\pgfqpoint{5.385952in}{0.647744in}}%
\pgfpathlineto{\pgfqpoint{5.389583in}{0.632153in}}%
\pgfpathlineto{\pgfqpoint{5.393215in}{0.536759in}}%
\pgfpathlineto{\pgfqpoint{5.396846in}{0.603220in}}%
\pgfpathlineto{\pgfqpoint{5.400477in}{0.607293in}}%
\pgfpathlineto{\pgfqpoint{5.404108in}{0.664567in}}%
\pgfpathlineto{\pgfqpoint{5.407739in}{0.575400in}}%
\pgfpathlineto{\pgfqpoint{5.411370in}{0.635910in}}%
\pgfpathlineto{\pgfqpoint{5.415002in}{0.660297in}}%
\pgfpathlineto{\pgfqpoint{5.418633in}{0.664467in}}%
\pgfpathlineto{\pgfqpoint{5.425895in}{0.651319in}}%
\pgfpathlineto{\pgfqpoint{5.429526in}{0.617199in}}%
\pgfpathlineto{\pgfqpoint{5.433157in}{0.576698in}}%
\pgfpathlineto{\pgfqpoint{5.436789in}{0.628089in}}%
\pgfpathlineto{\pgfqpoint{5.440420in}{0.670592in}}%
\pgfpathlineto{\pgfqpoint{5.444051in}{0.687271in}}%
\pgfpathlineto{\pgfqpoint{5.447682in}{0.603779in}}%
\pgfpathlineto{\pgfqpoint{5.451313in}{0.551712in}}%
\pgfpathlineto{\pgfqpoint{5.454944in}{0.652648in}}%
\pgfpathlineto{\pgfqpoint{5.458576in}{0.621134in}}%
\pgfpathlineto{\pgfqpoint{5.462207in}{0.606125in}}%
\pgfpathlineto{\pgfqpoint{5.465838in}{0.651458in}}%
\pgfpathlineto{\pgfqpoint{5.469469in}{0.634552in}}%
\pgfpathlineto{\pgfqpoint{5.473100in}{0.652662in}}%
\pgfpathlineto{\pgfqpoint{5.476731in}{0.618245in}}%
\pgfpathlineto{\pgfqpoint{5.480363in}{0.642834in}}%
\pgfpathlineto{\pgfqpoint{5.483994in}{0.656805in}}%
\pgfpathlineto{\pgfqpoint{5.487625in}{0.644708in}}%
\pgfpathlineto{\pgfqpoint{5.491256in}{0.505652in}}%
\pgfpathlineto{\pgfqpoint{5.494887in}{0.672308in}}%
\pgfpathlineto{\pgfqpoint{5.498519in}{0.649903in}}%
\pgfpathlineto{\pgfqpoint{5.502150in}{0.654274in}}%
\pgfpathlineto{\pgfqpoint{5.505781in}{0.643693in}}%
\pgfpathlineto{\pgfqpoint{5.509412in}{0.628042in}}%
\pgfpathlineto{\pgfqpoint{5.513043in}{0.597819in}}%
\pgfpathlineto{\pgfqpoint{5.516674in}{0.589469in}}%
\pgfpathlineto{\pgfqpoint{5.520306in}{0.668442in}}%
\pgfpathlineto{\pgfqpoint{5.523937in}{0.692204in}}%
\pgfpathlineto{\pgfqpoint{5.527568in}{0.666321in}}%
\pgfpathlineto{\pgfqpoint{5.531199in}{0.654142in}}%
\pgfpathlineto{\pgfqpoint{5.534830in}{0.661394in}}%
\pgfpathlineto{\pgfqpoint{5.538461in}{0.602375in}}%
\pgfpathlineto{\pgfqpoint{5.542093in}{0.584313in}}%
\pgfpathlineto{\pgfqpoint{5.545724in}{0.645892in}}%
\pgfpathlineto{\pgfqpoint{5.549355in}{0.677031in}}%
\pgfpathlineto{\pgfqpoint{5.552986in}{0.690302in}}%
\pgfpathlineto{\pgfqpoint{5.556617in}{0.416690in}}%
\pgfpathlineto{\pgfqpoint{5.560248in}{0.589100in}}%
\pgfpathlineto{\pgfqpoint{5.563880in}{0.640033in}}%
\pgfpathlineto{\pgfqpoint{5.567511in}{0.575438in}}%
\pgfpathlineto{\pgfqpoint{5.571142in}{0.647788in}}%
\pgfpathlineto{\pgfqpoint{5.574773in}{0.669499in}}%
\pgfpathlineto{\pgfqpoint{5.578404in}{0.634337in}}%
\pgfpathlineto{\pgfqpoint{5.582035in}{0.636098in}}%
\pgfpathlineto{\pgfqpoint{5.585667in}{0.635962in}}%
\pgfpathlineto{\pgfqpoint{5.589298in}{0.641885in}}%
\pgfpathlineto{\pgfqpoint{5.592929in}{0.639071in}}%
\pgfpathlineto{\pgfqpoint{5.596560in}{0.621278in}}%
\pgfpathlineto{\pgfqpoint{5.600191in}{0.572379in}}%
\pgfpathlineto{\pgfqpoint{5.603823in}{0.669075in}}%
\pgfpathlineto{\pgfqpoint{5.607454in}{0.576158in}}%
\pgfpathlineto{\pgfqpoint{5.614716in}{0.656308in}}%
\pgfpathlineto{\pgfqpoint{5.618347in}{0.672345in}}%
\pgfpathlineto{\pgfqpoint{5.621978in}{0.677178in}}%
\pgfpathlineto{\pgfqpoint{5.625610in}{0.616843in}}%
\pgfpathlineto{\pgfqpoint{5.629241in}{0.616252in}}%
\pgfpathlineto{\pgfqpoint{5.632872in}{0.630398in}}%
\pgfpathlineto{\pgfqpoint{5.636503in}{0.559347in}}%
\pgfpathlineto{\pgfqpoint{5.640134in}{0.659413in}}%
\pgfpathlineto{\pgfqpoint{5.643765in}{0.680993in}}%
\pgfpathlineto{\pgfqpoint{5.647397in}{0.632549in}}%
\pgfpathlineto{\pgfqpoint{5.651028in}{0.601942in}}%
\pgfpathlineto{\pgfqpoint{5.654659in}{0.649044in}}%
\pgfpathlineto{\pgfqpoint{5.658290in}{0.639009in}}%
\pgfpathlineto{\pgfqpoint{5.661921in}{0.605304in}}%
\pgfpathlineto{\pgfqpoint{5.665552in}{0.583983in}}%
\pgfpathlineto{\pgfqpoint{5.669184in}{0.611352in}}%
\pgfpathlineto{\pgfqpoint{5.672815in}{0.667332in}}%
\pgfpathlineto{\pgfqpoint{5.676446in}{0.534394in}}%
\pgfpathlineto{\pgfqpoint{5.680077in}{0.633088in}}%
\pgfpathlineto{\pgfqpoint{5.683708in}{0.662316in}}%
\pgfpathlineto{\pgfqpoint{5.687340in}{0.663455in}}%
\pgfpathlineto{\pgfqpoint{5.694602in}{0.630919in}}%
\pgfpathlineto{\pgfqpoint{5.698233in}{0.640085in}}%
\pgfpathlineto{\pgfqpoint{5.701864in}{0.672158in}}%
\pgfpathlineto{\pgfqpoint{5.705495in}{0.615840in}}%
\pgfpathlineto{\pgfqpoint{5.712758in}{0.639793in}}%
\pgfpathlineto{\pgfqpoint{5.716389in}{0.607511in}}%
\pgfpathlineto{\pgfqpoint{5.723651in}{0.676874in}}%
\pgfpathlineto{\pgfqpoint{5.727282in}{0.669213in}}%
\pgfpathlineto{\pgfqpoint{5.730914in}{0.675934in}}%
\pgfpathlineto{\pgfqpoint{5.738176in}{0.616042in}}%
\pgfpathlineto{\pgfqpoint{5.741807in}{0.586757in}}%
\pgfpathlineto{\pgfqpoint{5.745438in}{0.617378in}}%
\pgfpathlineto{\pgfqpoint{5.749069in}{0.671419in}}%
\pgfpathlineto{\pgfqpoint{5.752701in}{0.688861in}}%
\pgfpathlineto{\pgfqpoint{5.759963in}{0.533216in}}%
\pgfpathlineto{\pgfqpoint{5.763594in}{0.649142in}}%
\pgfpathlineto{\pgfqpoint{5.770856in}{0.597655in}}%
\pgfpathlineto{\pgfqpoint{5.774488in}{0.639220in}}%
\pgfpathlineto{\pgfqpoint{5.778119in}{0.624340in}}%
\pgfpathlineto{\pgfqpoint{5.781750in}{0.657418in}}%
\pgfpathlineto{\pgfqpoint{5.785381in}{0.607410in}}%
\pgfpathlineto{\pgfqpoint{5.789012in}{0.640448in}}%
\pgfpathlineto{\pgfqpoint{5.792644in}{0.657735in}}%
\pgfpathlineto{\pgfqpoint{5.796275in}{0.655906in}}%
\pgfpathlineto{\pgfqpoint{5.799906in}{0.641602in}}%
\pgfpathlineto{\pgfqpoint{5.803537in}{0.660605in}}%
\pgfpathlineto{\pgfqpoint{5.807168in}{0.611521in}}%
\pgfpathlineto{\pgfqpoint{5.810799in}{0.523890in}}%
\pgfpathlineto{\pgfqpoint{5.814431in}{0.641777in}}%
\pgfpathlineto{\pgfqpoint{5.818062in}{0.672843in}}%
\pgfpathlineto{\pgfqpoint{5.821693in}{0.686694in}}%
\pgfpathlineto{\pgfqpoint{5.825324in}{0.549223in}}%
\pgfpathlineto{\pgfqpoint{5.828955in}{0.593541in}}%
\pgfpathlineto{\pgfqpoint{5.832586in}{0.649109in}}%
\pgfpathlineto{\pgfqpoint{5.836218in}{0.601402in}}%
\pgfpathlineto{\pgfqpoint{5.843480in}{0.665470in}}%
\pgfpathlineto{\pgfqpoint{5.847111in}{0.637860in}}%
\pgfpathlineto{\pgfqpoint{5.850742in}{0.640872in}}%
\pgfpathlineto{\pgfqpoint{5.854373in}{0.633340in}}%
\pgfpathlineto{\pgfqpoint{5.858005in}{0.643889in}}%
\pgfpathlineto{\pgfqpoint{5.861636in}{0.647661in}}%
\pgfpathlineto{\pgfqpoint{5.865267in}{0.623990in}}%
\pgfpathlineto{\pgfqpoint{5.868898in}{0.623191in}}%
\pgfpathlineto{\pgfqpoint{5.872529in}{0.683268in}}%
\pgfpathlineto{\pgfqpoint{5.876160in}{0.649072in}}%
\pgfpathlineto{\pgfqpoint{5.879792in}{0.640152in}}%
\pgfpathlineto{\pgfqpoint{5.883423in}{0.651367in}}%
\pgfpathlineto{\pgfqpoint{5.887054in}{0.623154in}}%
\pgfpathlineto{\pgfqpoint{5.890685in}{0.466225in}}%
\pgfpathlineto{\pgfqpoint{5.894316in}{0.626115in}}%
\pgfpathlineto{\pgfqpoint{5.897948in}{0.675670in}}%
\pgfpathlineto{\pgfqpoint{5.901579in}{0.696192in}}%
\pgfpathlineto{\pgfqpoint{5.905210in}{0.659715in}}%
\pgfpathlineto{\pgfqpoint{5.908841in}{0.635996in}}%
\pgfpathlineto{\pgfqpoint{5.912472in}{0.665607in}}%
\pgfpathlineto{\pgfqpoint{5.916103in}{0.583812in}}%
\pgfpathlineto{\pgfqpoint{5.919735in}{0.612070in}}%
\pgfpathlineto{\pgfqpoint{5.923366in}{0.653826in}}%
\pgfpathlineto{\pgfqpoint{5.926997in}{0.677122in}}%
\pgfpathlineto{\pgfqpoint{5.930628in}{0.686537in}}%
\pgfpathlineto{\pgfqpoint{5.934259in}{0.595589in}}%
\pgfpathlineto{\pgfqpoint{5.937890in}{0.603853in}}%
\pgfpathlineto{\pgfqpoint{5.941522in}{0.629537in}}%
\pgfpathlineto{\pgfqpoint{5.945153in}{0.517532in}}%
\pgfpathlineto{\pgfqpoint{5.948784in}{0.659011in}}%
\pgfpathlineto{\pgfqpoint{5.952415in}{0.678266in}}%
\pgfpathlineto{\pgfqpoint{5.956046in}{0.633008in}}%
\pgfpathlineto{\pgfqpoint{5.959677in}{0.618364in}}%
\pgfpathlineto{\pgfqpoint{5.963309in}{0.643853in}}%
\pgfpathlineto{\pgfqpoint{5.966940in}{0.639268in}}%
\pgfpathlineto{\pgfqpoint{5.970571in}{0.617150in}}%
\pgfpathlineto{\pgfqpoint{5.974202in}{0.552964in}}%
\pgfpathlineto{\pgfqpoint{5.977833in}{0.566548in}}%
\pgfpathlineto{\pgfqpoint{5.981465in}{0.668911in}}%
\pgfpathlineto{\pgfqpoint{5.985096in}{0.460681in}}%
\pgfpathlineto{\pgfqpoint{5.988727in}{0.628104in}}%
\pgfpathlineto{\pgfqpoint{5.992358in}{0.659741in}}%
\pgfpathlineto{\pgfqpoint{5.995989in}{0.669172in}}%
\pgfpathlineto{\pgfqpoint{5.999620in}{0.668291in}}%
\pgfpathlineto{\pgfqpoint{6.003252in}{0.638332in}}%
\pgfpathlineto{\pgfqpoint{6.006883in}{0.619277in}}%
\pgfpathlineto{\pgfqpoint{6.010514in}{0.612250in}}%
\pgfpathlineto{\pgfqpoint{6.014145in}{0.606935in}}%
\pgfpathlineto{\pgfqpoint{6.017776in}{0.666106in}}%
\pgfpathlineto{\pgfqpoint{6.021407in}{0.685212in}}%
\pgfpathlineto{\pgfqpoint{6.025039in}{0.622588in}}%
\pgfpathlineto{\pgfqpoint{6.028670in}{0.548227in}}%
\pgfpathlineto{\pgfqpoint{6.032301in}{0.652649in}}%
\pgfpathlineto{\pgfqpoint{6.035932in}{0.632299in}}%
\pgfpathlineto{\pgfqpoint{6.039563in}{0.419559in}}%
\pgfpathlineto{\pgfqpoint{6.043194in}{0.630841in}}%
\pgfpathlineto{\pgfqpoint{6.046826in}{0.627554in}}%
\pgfpathlineto{\pgfqpoint{6.050457in}{0.660978in}}%
\pgfpathlineto{\pgfqpoint{6.054088in}{0.594714in}}%
\pgfpathlineto{\pgfqpoint{6.057719in}{0.639601in}}%
\pgfpathlineto{\pgfqpoint{6.061350in}{0.661498in}}%
\pgfpathlineto{\pgfqpoint{6.064981in}{0.655723in}}%
\pgfpathlineto{\pgfqpoint{6.068613in}{0.615315in}}%
\pgfpathlineto{\pgfqpoint{6.072244in}{0.657509in}}%
\pgfpathlineto{\pgfqpoint{6.075875in}{0.647467in}}%
\pgfpathlineto{\pgfqpoint{6.079506in}{0.663884in}}%
\pgfpathlineto{\pgfqpoint{6.083137in}{0.633010in}}%
\pgfpathlineto{\pgfqpoint{6.086769in}{0.630126in}}%
\pgfpathlineto{\pgfqpoint{6.090400in}{0.625000in}}%
\pgfpathlineto{\pgfqpoint{6.094031in}{0.534655in}}%
\pgfpathlineto{\pgfqpoint{6.097662in}{0.657728in}}%
\pgfpathlineto{\pgfqpoint{6.101293in}{0.686078in}}%
\pgfpathlineto{\pgfqpoint{6.104924in}{0.668810in}}%
\pgfpathlineto{\pgfqpoint{6.108556in}{0.665542in}}%
\pgfpathlineto{\pgfqpoint{6.112187in}{0.654749in}}%
\pgfpathlineto{\pgfqpoint{6.115818in}{0.612064in}}%
\pgfpathlineto{\pgfqpoint{6.119449in}{0.341198in}}%
\pgfpathlineto{\pgfqpoint{6.123080in}{0.635029in}}%
\pgfpathlineto{\pgfqpoint{6.126711in}{0.674995in}}%
\pgfpathlineto{\pgfqpoint{6.130343in}{0.690605in}}%
\pgfpathlineto{\pgfqpoint{6.133974in}{0.588597in}}%
\pgfpathlineto{\pgfqpoint{6.137605in}{0.558122in}}%
\pgfpathlineto{\pgfqpoint{6.141236in}{0.646824in}}%
\pgfpathlineto{\pgfqpoint{6.144867in}{0.606478in}}%
\pgfpathlineto{\pgfqpoint{6.152130in}{0.657981in}}%
\pgfpathlineto{\pgfqpoint{6.155761in}{0.632246in}}%
\pgfpathlineto{\pgfqpoint{6.159392in}{0.647858in}}%
\pgfpathlineto{\pgfqpoint{6.163023in}{0.625187in}}%
\pgfpathlineto{\pgfqpoint{6.166654in}{0.642382in}}%
\pgfpathlineto{\pgfqpoint{6.170285in}{0.650853in}}%
\pgfpathlineto{\pgfqpoint{6.173917in}{0.642321in}}%
\pgfpathlineto{\pgfqpoint{6.177548in}{0.616480in}}%
\pgfpathlineto{\pgfqpoint{6.181179in}{0.666447in}}%
\pgfpathlineto{\pgfqpoint{6.184810in}{0.599628in}}%
\pgfpathlineto{\pgfqpoint{6.188441in}{0.593232in}}%
\pgfpathlineto{\pgfqpoint{6.192073in}{0.650885in}}%
\pgfpathlineto{\pgfqpoint{6.195704in}{0.673249in}}%
\pgfpathlineto{\pgfqpoint{6.199335in}{0.682812in}}%
\pgfpathlineto{\pgfqpoint{6.202966in}{0.578955in}}%
\pgfpathlineto{\pgfqpoint{6.210228in}{0.641604in}}%
\pgfpathlineto{\pgfqpoint{6.213860in}{0.553309in}}%
\pgfpathlineto{\pgfqpoint{6.217491in}{0.649893in}}%
\pgfpathlineto{\pgfqpoint{6.221122in}{0.674933in}}%
\pgfpathlineto{\pgfqpoint{6.224753in}{0.637082in}}%
\pgfpathlineto{\pgfqpoint{6.228384in}{0.624431in}}%
\pgfpathlineto{\pgfqpoint{6.232015in}{0.643284in}}%
\pgfpathlineto{\pgfqpoint{6.235647in}{0.642607in}}%
\pgfpathlineto{\pgfqpoint{6.239278in}{0.632178in}}%
\pgfpathlineto{\pgfqpoint{6.242909in}{0.572109in}}%
\pgfpathlineto{\pgfqpoint{6.246540in}{0.649103in}}%
\pgfpathlineto{\pgfqpoint{6.250171in}{0.690611in}}%
\pgfpathlineto{\pgfqpoint{6.253802in}{0.644153in}}%
\pgfpathlineto{\pgfqpoint{6.257434in}{0.620254in}}%
\pgfpathlineto{\pgfqpoint{6.261065in}{0.655829in}}%
\pgfpathlineto{\pgfqpoint{6.264696in}{0.614035in}}%
\pgfpathlineto{\pgfqpoint{6.268327in}{0.593348in}}%
\pgfpathlineto{\pgfqpoint{6.271958in}{0.645045in}}%
\pgfpathlineto{\pgfqpoint{6.275590in}{0.680165in}}%
\pgfpathlineto{\pgfqpoint{6.279221in}{0.697713in}}%
\pgfpathlineto{\pgfqpoint{6.282852in}{0.647894in}}%
\pgfpathlineto{\pgfqpoint{6.286483in}{0.607802in}}%
\pgfpathlineto{\pgfqpoint{6.290114in}{0.666630in}}%
\pgfpathlineto{\pgfqpoint{6.293745in}{0.533946in}}%
\pgfpathlineto{\pgfqpoint{6.297377in}{0.626655in}}%
\pgfpathlineto{\pgfqpoint{6.301008in}{0.658827in}}%
\pgfpathlineto{\pgfqpoint{6.304639in}{0.675448in}}%
\pgfpathlineto{\pgfqpoint{6.308270in}{0.679964in}}%
\pgfpathlineto{\pgfqpoint{6.311901in}{0.626351in}}%
\pgfpathlineto{\pgfqpoint{6.315532in}{0.610260in}}%
\pgfpathlineto{\pgfqpoint{6.319164in}{0.613205in}}%
\pgfpathlineto{\pgfqpoint{6.322795in}{0.594417in}}%
\pgfpathlineto{\pgfqpoint{6.326426in}{0.666448in}}%
\pgfpathlineto{\pgfqpoint{6.330057in}{0.684169in}}%
\pgfpathlineto{\pgfqpoint{6.333688in}{0.627002in}}%
\pgfpathlineto{\pgfqpoint{6.337319in}{0.590062in}}%
\pgfpathlineto{\pgfqpoint{6.340951in}{0.648529in}}%
\pgfpathlineto{\pgfqpoint{6.344582in}{0.633956in}}%
\pgfpathlineto{\pgfqpoint{6.348213in}{0.559095in}}%
\pgfpathlineto{\pgfqpoint{6.351844in}{0.603280in}}%
\pgfpathlineto{\pgfqpoint{6.355475in}{0.609072in}}%
\pgfpathlineto{\pgfqpoint{6.359106in}{0.664672in}}%
\pgfpathlineto{\pgfqpoint{6.362738in}{0.575601in}}%
\pgfpathlineto{\pgfqpoint{6.366369in}{0.619393in}}%
\pgfpathlineto{\pgfqpoint{6.370000in}{0.645262in}}%
\pgfpathlineto{\pgfqpoint{6.370000in}{0.645262in}}%
\pgfusepath{stroke}%
\end{pgfscope}%
\begin{pgfscope}%
\pgfpathrectangle{\pgfqpoint{4.554414in}{0.320625in}}{\pgfqpoint{1.815586in}{1.164375in}}%
\pgfusepath{clip}%
\pgfsetrectcap%
\pgfsetroundjoin%
\pgfsetlinewidth{0.602250pt}%
\definecolor{currentstroke}{rgb}{0.478431,0.243137,0.615686}%
\pgfsetstrokecolor{currentstroke}%
\pgfsetstrokeopacity{0.120000}%
\pgfsetdash{}{0pt}%
\pgfpathmoveto{\pgfqpoint{4.554414in}{1.453399in}}%
\pgfpathlineto{\pgfqpoint{4.558045in}{1.443303in}}%
\pgfpathlineto{\pgfqpoint{4.579832in}{1.398849in}}%
\pgfpathlineto{\pgfqpoint{4.597988in}{1.358420in}}%
\pgfpathlineto{\pgfqpoint{4.670611in}{1.193863in}}%
\pgfpathlineto{\pgfqpoint{4.706923in}{1.115319in}}%
\pgfpathlineto{\pgfqpoint{4.746866in}{1.031780in}}%
\pgfpathlineto{\pgfqpoint{4.790440in}{0.943261in}}%
\pgfpathlineto{\pgfqpoint{4.844907in}{0.835280in}}%
\pgfpathlineto{\pgfqpoint{4.899375in}{0.729406in}}%
\pgfpathlineto{\pgfqpoint{4.903006in}{0.724358in}}%
\pgfpathlineto{\pgfqpoint{4.906637in}{0.720938in}}%
\pgfpathlineto{\pgfqpoint{4.910269in}{0.708229in}}%
\pgfpathlineto{\pgfqpoint{4.913900in}{0.701364in}}%
\pgfpathlineto{\pgfqpoint{4.917531in}{0.695965in}}%
\pgfpathlineto{\pgfqpoint{4.921162in}{0.687229in}}%
\pgfpathlineto{\pgfqpoint{4.924793in}{0.687122in}}%
\pgfpathlineto{\pgfqpoint{4.928424in}{0.693796in}}%
\pgfpathlineto{\pgfqpoint{4.932056in}{0.669765in}}%
\pgfpathlineto{\pgfqpoint{4.935687in}{0.659582in}}%
\pgfpathlineto{\pgfqpoint{4.939318in}{0.667676in}}%
\pgfpathlineto{\pgfqpoint{4.942949in}{0.652350in}}%
\pgfpathlineto{\pgfqpoint{4.946580in}{0.644763in}}%
\pgfpathlineto{\pgfqpoint{4.950211in}{0.666069in}}%
\pgfpathlineto{\pgfqpoint{4.953843in}{0.647481in}}%
\pgfpathlineto{\pgfqpoint{4.957474in}{0.644907in}}%
\pgfpathlineto{\pgfqpoint{4.961105in}{0.647019in}}%
\pgfpathlineto{\pgfqpoint{4.964736in}{0.646652in}}%
\pgfpathlineto{\pgfqpoint{4.971998in}{0.604423in}}%
\pgfpathlineto{\pgfqpoint{4.975630in}{0.624398in}}%
\pgfpathlineto{\pgfqpoint{4.979261in}{0.664222in}}%
\pgfpathlineto{\pgfqpoint{4.982892in}{0.609252in}}%
\pgfpathlineto{\pgfqpoint{4.986523in}{0.644746in}}%
\pgfpathlineto{\pgfqpoint{4.990154in}{0.657941in}}%
\pgfpathlineto{\pgfqpoint{4.993785in}{0.637005in}}%
\pgfpathlineto{\pgfqpoint{4.997417in}{0.624827in}}%
\pgfpathlineto{\pgfqpoint{5.001048in}{0.688611in}}%
\pgfpathlineto{\pgfqpoint{5.004679in}{0.647713in}}%
\pgfpathlineto{\pgfqpoint{5.008310in}{0.617077in}}%
\pgfpathlineto{\pgfqpoint{5.011941in}{0.659251in}}%
\pgfpathlineto{\pgfqpoint{5.015573in}{0.599552in}}%
\pgfpathlineto{\pgfqpoint{5.022835in}{0.666807in}}%
\pgfpathlineto{\pgfqpoint{5.026466in}{0.684870in}}%
\pgfpathlineto{\pgfqpoint{5.030097in}{0.688801in}}%
\pgfpathlineto{\pgfqpoint{5.033728in}{0.618494in}}%
\pgfpathlineto{\pgfqpoint{5.037360in}{0.616462in}}%
\pgfpathlineto{\pgfqpoint{5.040991in}{0.640102in}}%
\pgfpathlineto{\pgfqpoint{5.044622in}{0.636845in}}%
\pgfpathlineto{\pgfqpoint{5.048253in}{0.645664in}}%
\pgfpathlineto{\pgfqpoint{5.051884in}{0.632266in}}%
\pgfpathlineto{\pgfqpoint{5.055515in}{0.576728in}}%
\pgfpathlineto{\pgfqpoint{5.059147in}{0.571213in}}%
\pgfpathlineto{\pgfqpoint{5.062778in}{0.667347in}}%
\pgfpathlineto{\pgfqpoint{5.066409in}{0.595206in}}%
\pgfpathlineto{\pgfqpoint{5.070040in}{0.640568in}}%
\pgfpathlineto{\pgfqpoint{5.073671in}{0.656115in}}%
\pgfpathlineto{\pgfqpoint{5.077302in}{0.651759in}}%
\pgfpathlineto{\pgfqpoint{5.080934in}{0.624227in}}%
\pgfpathlineto{\pgfqpoint{5.084565in}{0.669642in}}%
\pgfpathlineto{\pgfqpoint{5.088196in}{0.577992in}}%
\pgfpathlineto{\pgfqpoint{5.095458in}{0.660717in}}%
\pgfpathlineto{\pgfqpoint{5.099090in}{0.671827in}}%
\pgfpathlineto{\pgfqpoint{5.102721in}{0.666630in}}%
\pgfpathlineto{\pgfqpoint{5.106352in}{0.652815in}}%
\pgfpathlineto{\pgfqpoint{5.109983in}{0.616849in}}%
\pgfpathlineto{\pgfqpoint{5.113614in}{0.455698in}}%
\pgfpathlineto{\pgfqpoint{5.117245in}{0.645768in}}%
\pgfpathlineto{\pgfqpoint{5.120877in}{0.675777in}}%
\pgfpathlineto{\pgfqpoint{5.124508in}{0.685403in}}%
\pgfpathlineto{\pgfqpoint{5.128139in}{0.601845in}}%
\pgfpathlineto{\pgfqpoint{5.131770in}{0.617795in}}%
\pgfpathlineto{\pgfqpoint{5.135401in}{0.627752in}}%
\pgfpathlineto{\pgfqpoint{5.139032in}{0.600232in}}%
\pgfpathlineto{\pgfqpoint{5.142664in}{0.668906in}}%
\pgfpathlineto{\pgfqpoint{5.146295in}{0.688624in}}%
\pgfpathlineto{\pgfqpoint{5.149926in}{0.606204in}}%
\pgfpathlineto{\pgfqpoint{5.153557in}{0.578855in}}%
\pgfpathlineto{\pgfqpoint{5.157188in}{0.653025in}}%
\pgfpathlineto{\pgfqpoint{5.160819in}{0.598217in}}%
\pgfpathlineto{\pgfqpoint{5.164451in}{0.646355in}}%
\pgfpathlineto{\pgfqpoint{5.168082in}{0.677324in}}%
\pgfpathlineto{\pgfqpoint{5.171713in}{0.637896in}}%
\pgfpathlineto{\pgfqpoint{5.175344in}{0.604547in}}%
\pgfpathlineto{\pgfqpoint{5.178975in}{0.653517in}}%
\pgfpathlineto{\pgfqpoint{5.182606in}{0.637105in}}%
\pgfpathlineto{\pgfqpoint{5.186238in}{0.520379in}}%
\pgfpathlineto{\pgfqpoint{5.189869in}{0.647670in}}%
\pgfpathlineto{\pgfqpoint{5.193500in}{0.678210in}}%
\pgfpathlineto{\pgfqpoint{5.197131in}{0.696236in}}%
\pgfpathlineto{\pgfqpoint{5.200762in}{0.583660in}}%
\pgfpathlineto{\pgfqpoint{5.208025in}{0.634779in}}%
\pgfpathlineto{\pgfqpoint{5.211656in}{0.619627in}}%
\pgfpathlineto{\pgfqpoint{5.215287in}{0.661062in}}%
\pgfpathlineto{\pgfqpoint{5.218918in}{0.675325in}}%
\pgfpathlineto{\pgfqpoint{5.222549in}{0.660657in}}%
\pgfpathlineto{\pgfqpoint{5.226181in}{0.584245in}}%
\pgfpathlineto{\pgfqpoint{5.229812in}{0.688157in}}%
\pgfpathlineto{\pgfqpoint{5.233443in}{0.604980in}}%
\pgfpathlineto{\pgfqpoint{5.237074in}{0.555127in}}%
\pgfpathlineto{\pgfqpoint{5.240705in}{0.650058in}}%
\pgfpathlineto{\pgfqpoint{5.244336in}{0.610981in}}%
\pgfpathlineto{\pgfqpoint{5.251599in}{0.672912in}}%
\pgfpathlineto{\pgfqpoint{5.255230in}{0.635107in}}%
\pgfpathlineto{\pgfqpoint{5.258861in}{0.622870in}}%
\pgfpathlineto{\pgfqpoint{5.262492in}{0.647574in}}%
\pgfpathlineto{\pgfqpoint{5.266123in}{0.635342in}}%
\pgfpathlineto{\pgfqpoint{5.269755in}{0.488639in}}%
\pgfpathlineto{\pgfqpoint{5.273386in}{0.631833in}}%
\pgfpathlineto{\pgfqpoint{5.277017in}{0.627053in}}%
\pgfpathlineto{\pgfqpoint{5.280648in}{0.655437in}}%
\pgfpathlineto{\pgfqpoint{5.284279in}{0.625535in}}%
\pgfpathlineto{\pgfqpoint{5.287910in}{0.644653in}}%
\pgfpathlineto{\pgfqpoint{5.291542in}{0.645116in}}%
\pgfpathlineto{\pgfqpoint{5.295173in}{0.615924in}}%
\pgfpathlineto{\pgfqpoint{5.298804in}{0.519348in}}%
\pgfpathlineto{\pgfqpoint{5.302435in}{0.671119in}}%
\pgfpathlineto{\pgfqpoint{5.306066in}{0.531722in}}%
\pgfpathlineto{\pgfqpoint{5.309698in}{0.636994in}}%
\pgfpathlineto{\pgfqpoint{5.313329in}{0.662618in}}%
\pgfpathlineto{\pgfqpoint{5.316960in}{0.660089in}}%
\pgfpathlineto{\pgfqpoint{5.320591in}{0.638858in}}%
\pgfpathlineto{\pgfqpoint{5.324222in}{0.666316in}}%
\pgfpathlineto{\pgfqpoint{5.327853in}{0.605771in}}%
\pgfpathlineto{\pgfqpoint{5.331485in}{0.602788in}}%
\pgfpathlineto{\pgfqpoint{5.335116in}{0.658085in}}%
\pgfpathlineto{\pgfqpoint{5.338747in}{0.672812in}}%
\pgfpathlineto{\pgfqpoint{5.342378in}{0.671888in}}%
\pgfpathlineto{\pgfqpoint{5.346009in}{0.642010in}}%
\pgfpathlineto{\pgfqpoint{5.349640in}{0.623882in}}%
\pgfpathlineto{\pgfqpoint{5.353272in}{0.592461in}}%
\pgfpathlineto{\pgfqpoint{5.360534in}{0.673012in}}%
\pgfpathlineto{\pgfqpoint{5.364165in}{0.686900in}}%
\pgfpathlineto{\pgfqpoint{5.367796in}{0.649784in}}%
\pgfpathlineto{\pgfqpoint{5.371427in}{0.490170in}}%
\pgfpathlineto{\pgfqpoint{5.375059in}{0.683388in}}%
\pgfpathlineto{\pgfqpoint{5.378690in}{0.538267in}}%
\pgfpathlineto{\pgfqpoint{5.382321in}{0.623522in}}%
\pgfpathlineto{\pgfqpoint{5.385952in}{0.649413in}}%
\pgfpathlineto{\pgfqpoint{5.389583in}{0.623805in}}%
\pgfpathlineto{\pgfqpoint{5.393215in}{0.640367in}}%
\pgfpathlineto{\pgfqpoint{5.396846in}{0.681993in}}%
\pgfpathlineto{\pgfqpoint{5.400477in}{0.671877in}}%
\pgfpathlineto{\pgfqpoint{5.404108in}{0.667761in}}%
\pgfpathlineto{\pgfqpoint{5.407739in}{0.660783in}}%
\pgfpathlineto{\pgfqpoint{5.411370in}{0.604528in}}%
\pgfpathlineto{\pgfqpoint{5.415002in}{0.600090in}}%
\pgfpathlineto{\pgfqpoint{5.418633in}{0.654617in}}%
\pgfpathlineto{\pgfqpoint{5.422264in}{0.679070in}}%
\pgfpathlineto{\pgfqpoint{5.425895in}{0.684843in}}%
\pgfpathlineto{\pgfqpoint{5.429526in}{0.629183in}}%
\pgfpathlineto{\pgfqpoint{5.433157in}{0.613512in}}%
\pgfpathlineto{\pgfqpoint{5.436789in}{0.591880in}}%
\pgfpathlineto{\pgfqpoint{5.440420in}{0.627178in}}%
\pgfpathlineto{\pgfqpoint{5.444051in}{0.675594in}}%
\pgfpathlineto{\pgfqpoint{5.447682in}{0.691036in}}%
\pgfpathlineto{\pgfqpoint{5.451313in}{0.556109in}}%
\pgfpathlineto{\pgfqpoint{5.458576in}{0.641410in}}%
\pgfpathlineto{\pgfqpoint{5.462207in}{0.506390in}}%
\pgfpathlineto{\pgfqpoint{5.465838in}{0.661141in}}%
\pgfpathlineto{\pgfqpoint{5.469469in}{0.684139in}}%
\pgfpathlineto{\pgfqpoint{5.476731in}{0.568309in}}%
\pgfpathlineto{\pgfqpoint{5.480363in}{0.653309in}}%
\pgfpathlineto{\pgfqpoint{5.483994in}{0.624704in}}%
\pgfpathlineto{\pgfqpoint{5.487625in}{0.618800in}}%
\pgfpathlineto{\pgfqpoint{5.491256in}{0.661804in}}%
\pgfpathlineto{\pgfqpoint{5.494887in}{0.638327in}}%
\pgfpathlineto{\pgfqpoint{5.498519in}{0.637926in}}%
\pgfpathlineto{\pgfqpoint{5.502150in}{0.642871in}}%
\pgfpathlineto{\pgfqpoint{5.505781in}{0.643320in}}%
\pgfpathlineto{\pgfqpoint{5.513043in}{0.589792in}}%
\pgfpathlineto{\pgfqpoint{5.516674in}{0.620191in}}%
\pgfpathlineto{\pgfqpoint{5.520306in}{0.663724in}}%
\pgfpathlineto{\pgfqpoint{5.523937in}{0.606582in}}%
\pgfpathlineto{\pgfqpoint{5.527568in}{0.644231in}}%
\pgfpathlineto{\pgfqpoint{5.531199in}{0.657260in}}%
\pgfpathlineto{\pgfqpoint{5.534830in}{0.637098in}}%
\pgfpathlineto{\pgfqpoint{5.538461in}{0.565003in}}%
\pgfpathlineto{\pgfqpoint{5.542093in}{0.671964in}}%
\pgfpathlineto{\pgfqpoint{5.545724in}{0.568693in}}%
\pgfpathlineto{\pgfqpoint{5.549355in}{0.628415in}}%
\pgfpathlineto{\pgfqpoint{5.552986in}{0.664100in}}%
\pgfpathlineto{\pgfqpoint{5.556617in}{0.664802in}}%
\pgfpathlineto{\pgfqpoint{5.560248in}{0.635616in}}%
\pgfpathlineto{\pgfqpoint{5.563880in}{0.658658in}}%
\pgfpathlineto{\pgfqpoint{5.567511in}{0.651035in}}%
\pgfpathlineto{\pgfqpoint{5.571142in}{0.657911in}}%
\pgfpathlineto{\pgfqpoint{5.574773in}{0.648127in}}%
\pgfpathlineto{\pgfqpoint{5.578404in}{0.628150in}}%
\pgfpathlineto{\pgfqpoint{5.582035in}{0.472518in}}%
\pgfpathlineto{\pgfqpoint{5.585667in}{0.638165in}}%
\pgfpathlineto{\pgfqpoint{5.589298in}{0.681277in}}%
\pgfpathlineto{\pgfqpoint{5.592929in}{0.699187in}}%
\pgfpathlineto{\pgfqpoint{5.596560in}{0.634134in}}%
\pgfpathlineto{\pgfqpoint{5.600191in}{0.405293in}}%
\pgfpathlineto{\pgfqpoint{5.603823in}{0.662930in}}%
\pgfpathlineto{\pgfqpoint{5.607454in}{0.596671in}}%
\pgfpathlineto{\pgfqpoint{5.611085in}{0.644497in}}%
\pgfpathlineto{\pgfqpoint{5.614716in}{0.660852in}}%
\pgfpathlineto{\pgfqpoint{5.618347in}{0.655669in}}%
\pgfpathlineto{\pgfqpoint{5.621978in}{0.630382in}}%
\pgfpathlineto{\pgfqpoint{5.625610in}{0.670088in}}%
\pgfpathlineto{\pgfqpoint{5.629241in}{0.569096in}}%
\pgfpathlineto{\pgfqpoint{5.632872in}{0.623717in}}%
\pgfpathlineto{\pgfqpoint{5.636503in}{0.659942in}}%
\pgfpathlineto{\pgfqpoint{5.640134in}{0.671528in}}%
\pgfpathlineto{\pgfqpoint{5.643765in}{0.666200in}}%
\pgfpathlineto{\pgfqpoint{5.647397in}{0.653539in}}%
\pgfpathlineto{\pgfqpoint{5.651028in}{0.615295in}}%
\pgfpathlineto{\pgfqpoint{5.654659in}{0.516567in}}%
\pgfpathlineto{\pgfqpoint{5.658290in}{0.646486in}}%
\pgfpathlineto{\pgfqpoint{5.661921in}{0.676361in}}%
\pgfpathlineto{\pgfqpoint{5.665552in}{0.685779in}}%
\pgfpathlineto{\pgfqpoint{5.669184in}{0.604295in}}%
\pgfpathlineto{\pgfqpoint{5.672815in}{0.616927in}}%
\pgfpathlineto{\pgfqpoint{5.676446in}{0.624988in}}%
\pgfpathlineto{\pgfqpoint{5.680077in}{0.603380in}}%
\pgfpathlineto{\pgfqpoint{5.683708in}{0.669833in}}%
\pgfpathlineto{\pgfqpoint{5.687340in}{0.689036in}}%
\pgfpathlineto{\pgfqpoint{5.690971in}{0.604195in}}%
\pgfpathlineto{\pgfqpoint{5.694602in}{0.579368in}}%
\pgfpathlineto{\pgfqpoint{5.698233in}{0.651883in}}%
\pgfpathlineto{\pgfqpoint{5.701864in}{0.594828in}}%
\pgfpathlineto{\pgfqpoint{5.705495in}{0.648110in}}%
\pgfpathlineto{\pgfqpoint{5.709127in}{0.677870in}}%
\pgfpathlineto{\pgfqpoint{5.712758in}{0.637227in}}%
\pgfpathlineto{\pgfqpoint{5.716389in}{0.603232in}}%
\pgfpathlineto{\pgfqpoint{5.720020in}{0.653340in}}%
\pgfpathlineto{\pgfqpoint{5.723651in}{0.636146in}}%
\pgfpathlineto{\pgfqpoint{5.727282in}{0.545617in}}%
\pgfpathlineto{\pgfqpoint{5.730914in}{0.648497in}}%
\pgfpathlineto{\pgfqpoint{5.734545in}{0.678948in}}%
\pgfpathlineto{\pgfqpoint{5.738176in}{0.696969in}}%
\pgfpathlineto{\pgfqpoint{5.741807in}{0.584900in}}%
\pgfpathlineto{\pgfqpoint{5.749069in}{0.633264in}}%
\pgfpathlineto{\pgfqpoint{5.752701in}{0.620716in}}%
\pgfpathlineto{\pgfqpoint{5.756332in}{0.661332in}}%
\pgfpathlineto{\pgfqpoint{5.759963in}{0.675193in}}%
\pgfpathlineto{\pgfqpoint{5.763594in}{0.661394in}}%
\pgfpathlineto{\pgfqpoint{5.767225in}{0.588562in}}%
\pgfpathlineto{\pgfqpoint{5.770856in}{0.688530in}}%
\pgfpathlineto{\pgfqpoint{5.774488in}{0.603166in}}%
\pgfpathlineto{\pgfqpoint{5.778119in}{0.555071in}}%
\pgfpathlineto{\pgfqpoint{5.781750in}{0.649073in}}%
\pgfpathlineto{\pgfqpoint{5.785381in}{0.609024in}}%
\pgfpathlineto{\pgfqpoint{5.789012in}{0.643888in}}%
\pgfpathlineto{\pgfqpoint{5.792644in}{0.673324in}}%
\pgfpathlineto{\pgfqpoint{5.796275in}{0.634405in}}%
\pgfpathlineto{\pgfqpoint{5.799906in}{0.622432in}}%
\pgfpathlineto{\pgfqpoint{5.803537in}{0.647327in}}%
\pgfpathlineto{\pgfqpoint{5.807168in}{0.634563in}}%
\pgfpathlineto{\pgfqpoint{5.810799in}{0.525352in}}%
\pgfpathlineto{\pgfqpoint{5.814431in}{0.632265in}}%
\pgfpathlineto{\pgfqpoint{5.818062in}{0.626137in}}%
\pgfpathlineto{\pgfqpoint{5.821693in}{0.655137in}}%
\pgfpathlineto{\pgfqpoint{5.825324in}{0.625899in}}%
\pgfpathlineto{\pgfqpoint{5.828955in}{0.644389in}}%
\pgfpathlineto{\pgfqpoint{5.832586in}{0.644088in}}%
\pgfpathlineto{\pgfqpoint{5.836218in}{0.616244in}}%
\pgfpathlineto{\pgfqpoint{5.839849in}{0.500485in}}%
\pgfpathlineto{\pgfqpoint{5.843480in}{0.670988in}}%
\pgfpathlineto{\pgfqpoint{5.847111in}{0.537879in}}%
\pgfpathlineto{\pgfqpoint{5.850742in}{0.637110in}}%
\pgfpathlineto{\pgfqpoint{5.854373in}{0.662250in}}%
\pgfpathlineto{\pgfqpoint{5.858005in}{0.660427in}}%
\pgfpathlineto{\pgfqpoint{5.861636in}{0.639963in}}%
\pgfpathlineto{\pgfqpoint{5.865267in}{0.666257in}}%
\pgfpathlineto{\pgfqpoint{5.868898in}{0.604463in}}%
\pgfpathlineto{\pgfqpoint{5.872529in}{0.603886in}}%
\pgfpathlineto{\pgfqpoint{5.876160in}{0.657967in}}%
\pgfpathlineto{\pgfqpoint{5.879792in}{0.673155in}}%
\pgfpathlineto{\pgfqpoint{5.883423in}{0.672469in}}%
\pgfpathlineto{\pgfqpoint{5.887054in}{0.641917in}}%
\pgfpathlineto{\pgfqpoint{5.890685in}{0.623096in}}%
\pgfpathlineto{\pgfqpoint{5.894316in}{0.590247in}}%
\pgfpathlineto{\pgfqpoint{5.897948in}{0.635307in}}%
\pgfpathlineto{\pgfqpoint{5.901579in}{0.673370in}}%
\pgfpathlineto{\pgfqpoint{5.905210in}{0.687251in}}%
\pgfpathlineto{\pgfqpoint{5.908841in}{0.538890in}}%
\pgfpathlineto{\pgfqpoint{5.912472in}{0.614534in}}%
\pgfpathlineto{\pgfqpoint{5.916103in}{0.642373in}}%
\pgfpathlineto{\pgfqpoint{5.919735in}{0.541569in}}%
\pgfpathlineto{\pgfqpoint{5.923366in}{0.662026in}}%
\pgfpathlineto{\pgfqpoint{5.926997in}{0.685742in}}%
\pgfpathlineto{\pgfqpoint{5.930628in}{0.675188in}}%
\pgfpathlineto{\pgfqpoint{5.934259in}{0.654594in}}%
\pgfpathlineto{\pgfqpoint{5.937890in}{0.674257in}}%
\pgfpathlineto{\pgfqpoint{5.941522in}{0.605197in}}%
\pgfpathlineto{\pgfqpoint{5.945153in}{0.579295in}}%
\pgfpathlineto{\pgfqpoint{5.948784in}{0.654760in}}%
\pgfpathlineto{\pgfqpoint{5.952415in}{0.653486in}}%
\pgfpathlineto{\pgfqpoint{5.956046in}{0.607134in}}%
\pgfpathlineto{\pgfqpoint{5.959677in}{0.640865in}}%
\pgfpathlineto{\pgfqpoint{5.966940in}{0.688141in}}%
\pgfpathlineto{\pgfqpoint{5.970571in}{0.632371in}}%
\pgfpathlineto{\pgfqpoint{5.974202in}{0.618386in}}%
\pgfpathlineto{\pgfqpoint{5.977833in}{0.590745in}}%
\pgfpathlineto{\pgfqpoint{5.981465in}{0.626789in}}%
\pgfpathlineto{\pgfqpoint{5.985096in}{0.676511in}}%
\pgfpathlineto{\pgfqpoint{5.988727in}{0.692817in}}%
\pgfpathlineto{\pgfqpoint{5.992358in}{0.552027in}}%
\pgfpathlineto{\pgfqpoint{5.999620in}{0.638620in}}%
\pgfpathlineto{\pgfqpoint{6.003252in}{0.463936in}}%
\pgfpathlineto{\pgfqpoint{6.006883in}{0.663112in}}%
\pgfpathlineto{\pgfqpoint{6.010514in}{0.684686in}}%
\pgfpathlineto{\pgfqpoint{6.017776in}{0.563809in}}%
\pgfpathlineto{\pgfqpoint{6.021407in}{0.652771in}}%
\pgfpathlineto{\pgfqpoint{6.025039in}{0.623275in}}%
\pgfpathlineto{\pgfqpoint{6.028670in}{0.622264in}}%
\pgfpathlineto{\pgfqpoint{6.032301in}{0.662619in}}%
\pgfpathlineto{\pgfqpoint{6.035932in}{0.637685in}}%
\pgfpathlineto{\pgfqpoint{6.039563in}{0.637054in}}%
\pgfpathlineto{\pgfqpoint{6.043194in}{0.643140in}}%
\pgfpathlineto{\pgfqpoint{6.046826in}{0.642618in}}%
\pgfpathlineto{\pgfqpoint{6.050457in}{0.612003in}}%
\pgfpathlineto{\pgfqpoint{6.054088in}{0.592886in}}%
\pgfpathlineto{\pgfqpoint{6.057719in}{0.619548in}}%
\pgfpathlineto{\pgfqpoint{6.061350in}{0.663206in}}%
\pgfpathlineto{\pgfqpoint{6.064981in}{0.608392in}}%
\pgfpathlineto{\pgfqpoint{6.068613in}{0.644211in}}%
\pgfpathlineto{\pgfqpoint{6.072244in}{0.656298in}}%
\pgfpathlineto{\pgfqpoint{6.075875in}{0.636796in}}%
\pgfpathlineto{\pgfqpoint{6.079506in}{0.567389in}}%
\pgfpathlineto{\pgfqpoint{6.083137in}{0.671901in}}%
\pgfpathlineto{\pgfqpoint{6.086769in}{0.563080in}}%
\pgfpathlineto{\pgfqpoint{6.090400in}{0.629212in}}%
\pgfpathlineto{\pgfqpoint{6.094031in}{0.663872in}}%
\pgfpathlineto{\pgfqpoint{6.097662in}{0.664969in}}%
\pgfpathlineto{\pgfqpoint{6.101293in}{0.635608in}}%
\pgfpathlineto{\pgfqpoint{6.104924in}{0.658840in}}%
\pgfpathlineto{\pgfqpoint{6.108556in}{0.650768in}}%
\pgfpathlineto{\pgfqpoint{6.112187in}{0.658081in}}%
\pgfpathlineto{\pgfqpoint{6.115818in}{0.648278in}}%
\pgfpathlineto{\pgfqpoint{6.119449in}{0.627107in}}%
\pgfpathlineto{\pgfqpoint{6.123080in}{0.475963in}}%
\pgfpathlineto{\pgfqpoint{6.126711in}{0.638970in}}%
\pgfpathlineto{\pgfqpoint{6.130343in}{0.682051in}}%
\pgfpathlineto{\pgfqpoint{6.133974in}{0.699812in}}%
\pgfpathlineto{\pgfqpoint{6.137605in}{0.635331in}}%
\pgfpathlineto{\pgfqpoint{6.141236in}{0.496618in}}%
\pgfpathlineto{\pgfqpoint{6.144867in}{0.662544in}}%
\pgfpathlineto{\pgfqpoint{6.148498in}{0.598312in}}%
\pgfpathlineto{\pgfqpoint{6.152130in}{0.644474in}}%
\pgfpathlineto{\pgfqpoint{6.155761in}{0.660212in}}%
\pgfpathlineto{\pgfqpoint{6.159392in}{0.655955in}}%
\pgfpathlineto{\pgfqpoint{6.163023in}{0.631768in}}%
\pgfpathlineto{\pgfqpoint{6.166654in}{0.670049in}}%
\pgfpathlineto{\pgfqpoint{6.170285in}{0.565104in}}%
\pgfpathlineto{\pgfqpoint{6.173917in}{0.624304in}}%
\pgfpathlineto{\pgfqpoint{6.177548in}{0.659695in}}%
\pgfpathlineto{\pgfqpoint{6.181179in}{0.671858in}}%
\pgfpathlineto{\pgfqpoint{6.184810in}{0.666854in}}%
\pgfpathlineto{\pgfqpoint{6.188441in}{0.653583in}}%
\pgfpathlineto{\pgfqpoint{6.192073in}{0.614121in}}%
\pgfpathlineto{\pgfqpoint{6.195704in}{0.529219in}}%
\pgfpathlineto{\pgfqpoint{6.199335in}{0.646578in}}%
\pgfpathlineto{\pgfqpoint{6.202966in}{0.676762in}}%
\pgfpathlineto{\pgfqpoint{6.206597in}{0.686236in}}%
\pgfpathlineto{\pgfqpoint{6.210228in}{0.604477in}}%
\pgfpathlineto{\pgfqpoint{6.213860in}{0.616064in}}%
\pgfpathlineto{\pgfqpoint{6.217491in}{0.623603in}}%
\pgfpathlineto{\pgfqpoint{6.221122in}{0.604190in}}%
\pgfpathlineto{\pgfqpoint{6.224753in}{0.670312in}}%
\pgfpathlineto{\pgfqpoint{6.228384in}{0.689284in}}%
\pgfpathlineto{\pgfqpoint{6.232015in}{0.603897in}}%
\pgfpathlineto{\pgfqpoint{6.235647in}{0.577986in}}%
\pgfpathlineto{\pgfqpoint{6.239278in}{0.651218in}}%
\pgfpathlineto{\pgfqpoint{6.242909in}{0.593617in}}%
\pgfpathlineto{\pgfqpoint{6.246540in}{0.648812in}}%
\pgfpathlineto{\pgfqpoint{6.250171in}{0.677957in}}%
\pgfpathlineto{\pgfqpoint{6.253802in}{0.636928in}}%
\pgfpathlineto{\pgfqpoint{6.257434in}{0.603638in}}%
\pgfpathlineto{\pgfqpoint{6.261065in}{0.653028in}}%
\pgfpathlineto{\pgfqpoint{6.264696in}{0.635750in}}%
\pgfpathlineto{\pgfqpoint{6.268327in}{0.551463in}}%
\pgfpathlineto{\pgfqpoint{6.271958in}{0.648282in}}%
\pgfpathlineto{\pgfqpoint{6.275590in}{0.637867in}}%
\pgfpathlineto{\pgfqpoint{6.279221in}{0.649144in}}%
\pgfpathlineto{\pgfqpoint{6.282852in}{0.634664in}}%
\pgfpathlineto{\pgfqpoint{6.286483in}{0.646109in}}%
\pgfpathlineto{\pgfqpoint{6.290114in}{0.639757in}}%
\pgfpathlineto{\pgfqpoint{6.293745in}{0.547520in}}%
\pgfpathlineto{\pgfqpoint{6.297377in}{0.660008in}}%
\pgfpathlineto{\pgfqpoint{6.301008in}{0.696901in}}%
\pgfpathlineto{\pgfqpoint{6.304639in}{0.628748in}}%
\pgfpathlineto{\pgfqpoint{6.308270in}{0.472460in}}%
\pgfpathlineto{\pgfqpoint{6.311901in}{0.656294in}}%
\pgfpathlineto{\pgfqpoint{6.315532in}{0.473204in}}%
\pgfpathlineto{\pgfqpoint{6.319164in}{0.649965in}}%
\pgfpathlineto{\pgfqpoint{6.322795in}{0.675114in}}%
\pgfpathlineto{\pgfqpoint{6.326426in}{0.680074in}}%
\pgfpathlineto{\pgfqpoint{6.330057in}{0.669214in}}%
\pgfpathlineto{\pgfqpoint{6.333688in}{0.663415in}}%
\pgfpathlineto{\pgfqpoint{6.337319in}{0.625133in}}%
\pgfpathlineto{\pgfqpoint{6.340951in}{0.610562in}}%
\pgfpathlineto{\pgfqpoint{6.344582in}{0.648251in}}%
\pgfpathlineto{\pgfqpoint{6.348213in}{0.639502in}}%
\pgfpathlineto{\pgfqpoint{6.351844in}{0.539412in}}%
\pgfpathlineto{\pgfqpoint{6.355475in}{0.629046in}}%
\pgfpathlineto{\pgfqpoint{6.359106in}{0.623042in}}%
\pgfpathlineto{\pgfqpoint{6.362738in}{0.655288in}}%
\pgfpathlineto{\pgfqpoint{6.366369in}{0.627048in}}%
\pgfpathlineto{\pgfqpoint{6.370000in}{0.586135in}}%
\pgfpathlineto{\pgfqpoint{6.370000in}{0.586135in}}%
\pgfusepath{stroke}%
\end{pgfscope}%
\begin{pgfscope}%
\pgfpathrectangle{\pgfqpoint{4.554414in}{0.320625in}}{\pgfqpoint{1.815586in}{1.164375in}}%
\pgfusepath{clip}%
\pgfsetrectcap%
\pgfsetroundjoin%
\pgfsetlinewidth{1.154312pt}%
\definecolor{currentstroke}{rgb}{0.478431,0.243137,0.615686}%
\pgfsetstrokecolor{currentstroke}%
\pgfsetdash{}{0pt}%
\pgfpathmoveto{\pgfqpoint{4.554414in}{1.453686in}}%
\pgfpathlineto{\pgfqpoint{4.558045in}{1.443533in}}%
\pgfpathlineto{\pgfqpoint{4.579832in}{1.399294in}}%
\pgfpathlineto{\pgfqpoint{4.597988in}{1.359136in}}%
\pgfpathlineto{\pgfqpoint{4.670611in}{1.195712in}}%
\pgfpathlineto{\pgfqpoint{4.706923in}{1.117447in}}%
\pgfpathlineto{\pgfqpoint{4.746866in}{1.033914in}}%
\pgfpathlineto{\pgfqpoint{4.794071in}{0.937891in}}%
\pgfpathlineto{\pgfqpoint{4.848539in}{0.829808in}}%
\pgfpathlineto{\pgfqpoint{4.899375in}{0.731119in}}%
\pgfpathlineto{\pgfqpoint{4.906637in}{0.718663in}}%
\pgfpathlineto{\pgfqpoint{4.910269in}{0.710617in}}%
\pgfpathlineto{\pgfqpoint{4.913900in}{0.705599in}}%
\pgfpathlineto{\pgfqpoint{4.917531in}{0.699484in}}%
\pgfpathlineto{\pgfqpoint{4.921162in}{0.691015in}}%
\pgfpathlineto{\pgfqpoint{4.924793in}{0.686387in}}%
\pgfpathlineto{\pgfqpoint{4.928424in}{0.684610in}}%
\pgfpathlineto{\pgfqpoint{4.932056in}{0.674976in}}%
\pgfpathlineto{\pgfqpoint{4.935687in}{0.672412in}}%
\pgfpathlineto{\pgfqpoint{4.939318in}{0.667516in}}%
\pgfpathlineto{\pgfqpoint{4.942949in}{0.673129in}}%
\pgfpathlineto{\pgfqpoint{4.946580in}{0.670988in}}%
\pgfpathlineto{\pgfqpoint{4.950211in}{0.652073in}}%
\pgfpathlineto{\pgfqpoint{4.953843in}{0.655294in}}%
\pgfpathlineto{\pgfqpoint{4.957474in}{0.663291in}}%
\pgfpathlineto{\pgfqpoint{4.961105in}{0.642684in}}%
\pgfpathlineto{\pgfqpoint{4.964736in}{0.643036in}}%
\pgfpathlineto{\pgfqpoint{4.968367in}{0.656444in}}%
\pgfpathlineto{\pgfqpoint{4.975630in}{0.671613in}}%
\pgfpathlineto{\pgfqpoint{4.979261in}{0.639684in}}%
\pgfpathlineto{\pgfqpoint{4.982892in}{0.658505in}}%
\pgfpathlineto{\pgfqpoint{4.986523in}{0.668268in}}%
\pgfpathlineto{\pgfqpoint{4.990154in}{0.637952in}}%
\pgfpathlineto{\pgfqpoint{4.993785in}{0.647343in}}%
\pgfpathlineto{\pgfqpoint{4.997417in}{0.662312in}}%
\pgfpathlineto{\pgfqpoint{5.001048in}{0.661872in}}%
\pgfpathlineto{\pgfqpoint{5.004679in}{0.651077in}}%
\pgfpathlineto{\pgfqpoint{5.008310in}{0.653520in}}%
\pgfpathlineto{\pgfqpoint{5.011941in}{0.665483in}}%
\pgfpathlineto{\pgfqpoint{5.015573in}{0.663572in}}%
\pgfpathlineto{\pgfqpoint{5.019204in}{0.631565in}}%
\pgfpathlineto{\pgfqpoint{5.022835in}{0.639174in}}%
\pgfpathlineto{\pgfqpoint{5.026466in}{0.667065in}}%
\pgfpathlineto{\pgfqpoint{5.030097in}{0.664683in}}%
\pgfpathlineto{\pgfqpoint{5.033728in}{0.669143in}}%
\pgfpathlineto{\pgfqpoint{5.037360in}{0.672139in}}%
\pgfpathlineto{\pgfqpoint{5.040991in}{0.659020in}}%
\pgfpathlineto{\pgfqpoint{5.044622in}{0.662375in}}%
\pgfpathlineto{\pgfqpoint{5.048253in}{0.645935in}}%
\pgfpathlineto{\pgfqpoint{5.051884in}{0.618144in}}%
\pgfpathlineto{\pgfqpoint{5.062778in}{0.668978in}}%
\pgfpathlineto{\pgfqpoint{5.066409in}{0.670692in}}%
\pgfpathlineto{\pgfqpoint{5.070040in}{0.647968in}}%
\pgfpathlineto{\pgfqpoint{5.077302in}{0.652732in}}%
\pgfpathlineto{\pgfqpoint{5.080934in}{0.639303in}}%
\pgfpathlineto{\pgfqpoint{5.084565in}{0.664077in}}%
\pgfpathlineto{\pgfqpoint{5.088196in}{0.648477in}}%
\pgfpathlineto{\pgfqpoint{5.091827in}{0.643457in}}%
\pgfpathlineto{\pgfqpoint{5.095458in}{0.663292in}}%
\pgfpathlineto{\pgfqpoint{5.099090in}{0.659857in}}%
\pgfpathlineto{\pgfqpoint{5.102721in}{0.642656in}}%
\pgfpathlineto{\pgfqpoint{5.106352in}{0.656522in}}%
\pgfpathlineto{\pgfqpoint{5.113614in}{0.652662in}}%
\pgfpathlineto{\pgfqpoint{5.117245in}{0.647965in}}%
\pgfpathlineto{\pgfqpoint{5.120877in}{0.652895in}}%
\pgfpathlineto{\pgfqpoint{5.124508in}{0.675498in}}%
\pgfpathlineto{\pgfqpoint{5.128139in}{0.639361in}}%
\pgfpathlineto{\pgfqpoint{5.131770in}{0.613727in}}%
\pgfpathlineto{\pgfqpoint{5.135401in}{0.652597in}}%
\pgfpathlineto{\pgfqpoint{5.139032in}{0.626913in}}%
\pgfpathlineto{\pgfqpoint{5.142664in}{0.646514in}}%
\pgfpathlineto{\pgfqpoint{5.146295in}{0.674377in}}%
\pgfpathlineto{\pgfqpoint{5.149926in}{0.673746in}}%
\pgfpathlineto{\pgfqpoint{5.153557in}{0.675500in}}%
\pgfpathlineto{\pgfqpoint{5.157188in}{0.639971in}}%
\pgfpathlineto{\pgfqpoint{5.164451in}{0.633602in}}%
\pgfpathlineto{\pgfqpoint{5.171713in}{0.661397in}}%
\pgfpathlineto{\pgfqpoint{5.175344in}{0.678832in}}%
\pgfpathlineto{\pgfqpoint{5.178975in}{0.666351in}}%
\pgfpathlineto{\pgfqpoint{5.182606in}{0.643546in}}%
\pgfpathlineto{\pgfqpoint{5.189869in}{0.650882in}}%
\pgfpathlineto{\pgfqpoint{5.193500in}{0.648746in}}%
\pgfpathlineto{\pgfqpoint{5.197131in}{0.668040in}}%
\pgfpathlineto{\pgfqpoint{5.200762in}{0.641434in}}%
\pgfpathlineto{\pgfqpoint{5.204394in}{0.651059in}}%
\pgfpathlineto{\pgfqpoint{5.208025in}{0.619858in}}%
\pgfpathlineto{\pgfqpoint{5.211656in}{0.648964in}}%
\pgfpathlineto{\pgfqpoint{5.215287in}{0.664197in}}%
\pgfpathlineto{\pgfqpoint{5.218918in}{0.660126in}}%
\pgfpathlineto{\pgfqpoint{5.222549in}{0.652942in}}%
\pgfpathlineto{\pgfqpoint{5.226181in}{0.647655in}}%
\pgfpathlineto{\pgfqpoint{5.229812in}{0.664047in}}%
\pgfpathlineto{\pgfqpoint{5.237074in}{0.653798in}}%
\pgfpathlineto{\pgfqpoint{5.240705in}{0.669016in}}%
\pgfpathlineto{\pgfqpoint{5.244336in}{0.657837in}}%
\pgfpathlineto{\pgfqpoint{5.247968in}{0.630021in}}%
\pgfpathlineto{\pgfqpoint{5.251599in}{0.666430in}}%
\pgfpathlineto{\pgfqpoint{5.255230in}{0.667333in}}%
\pgfpathlineto{\pgfqpoint{5.258861in}{0.635385in}}%
\pgfpathlineto{\pgfqpoint{5.262492in}{0.663884in}}%
\pgfpathlineto{\pgfqpoint{5.266123in}{0.655716in}}%
\pgfpathlineto{\pgfqpoint{5.269755in}{0.625275in}}%
\pgfpathlineto{\pgfqpoint{5.273386in}{0.641535in}}%
\pgfpathlineto{\pgfqpoint{5.277017in}{0.648173in}}%
\pgfpathlineto{\pgfqpoint{5.280648in}{0.646001in}}%
\pgfpathlineto{\pgfqpoint{5.284279in}{0.653195in}}%
\pgfpathlineto{\pgfqpoint{5.287910in}{0.648732in}}%
\pgfpathlineto{\pgfqpoint{5.291542in}{0.657377in}}%
\pgfpathlineto{\pgfqpoint{5.295173in}{0.648494in}}%
\pgfpathlineto{\pgfqpoint{5.298804in}{0.609442in}}%
\pgfpathlineto{\pgfqpoint{5.302435in}{0.658579in}}%
\pgfpathlineto{\pgfqpoint{5.306066in}{0.651637in}}%
\pgfpathlineto{\pgfqpoint{5.309698in}{0.640288in}}%
\pgfpathlineto{\pgfqpoint{5.313329in}{0.660790in}}%
\pgfpathlineto{\pgfqpoint{5.316960in}{0.649340in}}%
\pgfpathlineto{\pgfqpoint{5.320591in}{0.647122in}}%
\pgfpathlineto{\pgfqpoint{5.324222in}{0.667974in}}%
\pgfpathlineto{\pgfqpoint{5.327853in}{0.654983in}}%
\pgfpathlineto{\pgfqpoint{5.331485in}{0.659958in}}%
\pgfpathlineto{\pgfqpoint{5.335116in}{0.641284in}}%
\pgfpathlineto{\pgfqpoint{5.338747in}{0.646796in}}%
\pgfpathlineto{\pgfqpoint{5.346009in}{0.666319in}}%
\pgfpathlineto{\pgfqpoint{5.349640in}{0.664105in}}%
\pgfpathlineto{\pgfqpoint{5.353272in}{0.675328in}}%
\pgfpathlineto{\pgfqpoint{5.360534in}{0.645609in}}%
\pgfpathlineto{\pgfqpoint{5.364165in}{0.664992in}}%
\pgfpathlineto{\pgfqpoint{5.367796in}{0.641049in}}%
\pgfpathlineto{\pgfqpoint{5.371427in}{0.647030in}}%
\pgfpathlineto{\pgfqpoint{5.375059in}{0.668953in}}%
\pgfpathlineto{\pgfqpoint{5.378690in}{0.630508in}}%
\pgfpathlineto{\pgfqpoint{5.385952in}{0.671311in}}%
\pgfpathlineto{\pgfqpoint{5.389583in}{0.625324in}}%
\pgfpathlineto{\pgfqpoint{5.393215in}{0.637352in}}%
\pgfpathlineto{\pgfqpoint{5.396846in}{0.660424in}}%
\pgfpathlineto{\pgfqpoint{5.400477in}{0.657773in}}%
\pgfpathlineto{\pgfqpoint{5.404108in}{0.666536in}}%
\pgfpathlineto{\pgfqpoint{5.407739in}{0.657874in}}%
\pgfpathlineto{\pgfqpoint{5.411370in}{0.651587in}}%
\pgfpathlineto{\pgfqpoint{5.415002in}{0.652953in}}%
\pgfpathlineto{\pgfqpoint{5.418633in}{0.648349in}}%
\pgfpathlineto{\pgfqpoint{5.422264in}{0.658021in}}%
\pgfpathlineto{\pgfqpoint{5.425895in}{0.661933in}}%
\pgfpathlineto{\pgfqpoint{5.429526in}{0.629580in}}%
\pgfpathlineto{\pgfqpoint{5.433157in}{0.624617in}}%
\pgfpathlineto{\pgfqpoint{5.436789in}{0.654477in}}%
\pgfpathlineto{\pgfqpoint{5.440420in}{0.666950in}}%
\pgfpathlineto{\pgfqpoint{5.444051in}{0.671080in}}%
\pgfpathlineto{\pgfqpoint{5.447682in}{0.664228in}}%
\pgfpathlineto{\pgfqpoint{5.451313in}{0.644316in}}%
\pgfpathlineto{\pgfqpoint{5.454944in}{0.645334in}}%
\pgfpathlineto{\pgfqpoint{5.458576in}{0.644974in}}%
\pgfpathlineto{\pgfqpoint{5.462207in}{0.651677in}}%
\pgfpathlineto{\pgfqpoint{5.465838in}{0.661670in}}%
\pgfpathlineto{\pgfqpoint{5.469469in}{0.661647in}}%
\pgfpathlineto{\pgfqpoint{5.473100in}{0.657036in}}%
\pgfpathlineto{\pgfqpoint{5.476731in}{0.668198in}}%
\pgfpathlineto{\pgfqpoint{5.480363in}{0.634832in}}%
\pgfpathlineto{\pgfqpoint{5.483994in}{0.629760in}}%
\pgfpathlineto{\pgfqpoint{5.491256in}{0.654761in}}%
\pgfpathlineto{\pgfqpoint{5.494887in}{0.654402in}}%
\pgfpathlineto{\pgfqpoint{5.498519in}{0.652341in}}%
\pgfpathlineto{\pgfqpoint{5.502150in}{0.648574in}}%
\pgfpathlineto{\pgfqpoint{5.505781in}{0.665587in}}%
\pgfpathlineto{\pgfqpoint{5.509412in}{0.652128in}}%
\pgfpathlineto{\pgfqpoint{5.513043in}{0.625794in}}%
\pgfpathlineto{\pgfqpoint{5.516674in}{0.630423in}}%
\pgfpathlineto{\pgfqpoint{5.520306in}{0.649252in}}%
\pgfpathlineto{\pgfqpoint{5.523937in}{0.661168in}}%
\pgfpathlineto{\pgfqpoint{5.527568in}{0.668689in}}%
\pgfpathlineto{\pgfqpoint{5.531199in}{0.663970in}}%
\pgfpathlineto{\pgfqpoint{5.534830in}{0.654449in}}%
\pgfpathlineto{\pgfqpoint{5.538461in}{0.646487in}}%
\pgfpathlineto{\pgfqpoint{5.542093in}{0.655936in}}%
\pgfpathlineto{\pgfqpoint{5.545724in}{0.636814in}}%
\pgfpathlineto{\pgfqpoint{5.552986in}{0.669108in}}%
\pgfpathlineto{\pgfqpoint{5.556617in}{0.664308in}}%
\pgfpathlineto{\pgfqpoint{5.560248in}{0.642878in}}%
\pgfpathlineto{\pgfqpoint{5.563880in}{0.638840in}}%
\pgfpathlineto{\pgfqpoint{5.567511in}{0.638952in}}%
\pgfpathlineto{\pgfqpoint{5.571142in}{0.644828in}}%
\pgfpathlineto{\pgfqpoint{5.574773in}{0.662537in}}%
\pgfpathlineto{\pgfqpoint{5.578404in}{0.674838in}}%
\pgfpathlineto{\pgfqpoint{5.582035in}{0.657920in}}%
\pgfpathlineto{\pgfqpoint{5.585667in}{0.652459in}}%
\pgfpathlineto{\pgfqpoint{5.589298in}{0.663028in}}%
\pgfpathlineto{\pgfqpoint{5.592929in}{0.667447in}}%
\pgfpathlineto{\pgfqpoint{5.596560in}{0.633371in}}%
\pgfpathlineto{\pgfqpoint{5.603823in}{0.673057in}}%
\pgfpathlineto{\pgfqpoint{5.607454in}{0.673434in}}%
\pgfpathlineto{\pgfqpoint{5.611085in}{0.631179in}}%
\pgfpathlineto{\pgfqpoint{5.614716in}{0.643129in}}%
\pgfpathlineto{\pgfqpoint{5.618347in}{0.667808in}}%
\pgfpathlineto{\pgfqpoint{5.621978in}{0.647170in}}%
\pgfpathlineto{\pgfqpoint{5.625610in}{0.654913in}}%
\pgfpathlineto{\pgfqpoint{5.629241in}{0.664338in}}%
\pgfpathlineto{\pgfqpoint{5.632872in}{0.645472in}}%
\pgfpathlineto{\pgfqpoint{5.636503in}{0.641364in}}%
\pgfpathlineto{\pgfqpoint{5.640134in}{0.665367in}}%
\pgfpathlineto{\pgfqpoint{5.643765in}{0.667733in}}%
\pgfpathlineto{\pgfqpoint{5.647397in}{0.659521in}}%
\pgfpathlineto{\pgfqpoint{5.651028in}{0.623707in}}%
\pgfpathlineto{\pgfqpoint{5.654659in}{0.631271in}}%
\pgfpathlineto{\pgfqpoint{5.658290in}{0.650961in}}%
\pgfpathlineto{\pgfqpoint{5.661921in}{0.646524in}}%
\pgfpathlineto{\pgfqpoint{5.665552in}{0.660757in}}%
\pgfpathlineto{\pgfqpoint{5.669184in}{0.662740in}}%
\pgfpathlineto{\pgfqpoint{5.672815in}{0.649860in}}%
\pgfpathlineto{\pgfqpoint{5.676446in}{0.631029in}}%
\pgfpathlineto{\pgfqpoint{5.680077in}{0.653056in}}%
\pgfpathlineto{\pgfqpoint{5.683708in}{0.659839in}}%
\pgfpathlineto{\pgfqpoint{5.687340in}{0.669460in}}%
\pgfpathlineto{\pgfqpoint{5.690971in}{0.649327in}}%
\pgfpathlineto{\pgfqpoint{5.694602in}{0.657706in}}%
\pgfpathlineto{\pgfqpoint{5.698233in}{0.668509in}}%
\pgfpathlineto{\pgfqpoint{5.701864in}{0.643914in}}%
\pgfpathlineto{\pgfqpoint{5.705495in}{0.640408in}}%
\pgfpathlineto{\pgfqpoint{5.709127in}{0.663592in}}%
\pgfpathlineto{\pgfqpoint{5.712758in}{0.631194in}}%
\pgfpathlineto{\pgfqpoint{5.716389in}{0.641362in}}%
\pgfpathlineto{\pgfqpoint{5.720020in}{0.665519in}}%
\pgfpathlineto{\pgfqpoint{5.723651in}{0.667073in}}%
\pgfpathlineto{\pgfqpoint{5.727282in}{0.664860in}}%
\pgfpathlineto{\pgfqpoint{5.730914in}{0.664473in}}%
\pgfpathlineto{\pgfqpoint{5.734545in}{0.650746in}}%
\pgfpathlineto{\pgfqpoint{5.738176in}{0.669071in}}%
\pgfpathlineto{\pgfqpoint{5.741807in}{0.632571in}}%
\pgfpathlineto{\pgfqpoint{5.745438in}{0.642438in}}%
\pgfpathlineto{\pgfqpoint{5.749069in}{0.662674in}}%
\pgfpathlineto{\pgfqpoint{5.752701in}{0.670677in}}%
\pgfpathlineto{\pgfqpoint{5.756332in}{0.660059in}}%
\pgfpathlineto{\pgfqpoint{5.759963in}{0.671956in}}%
\pgfpathlineto{\pgfqpoint{5.763594in}{0.641868in}}%
\pgfpathlineto{\pgfqpoint{5.767225in}{0.618199in}}%
\pgfpathlineto{\pgfqpoint{5.770856in}{0.660579in}}%
\pgfpathlineto{\pgfqpoint{5.774488in}{0.638910in}}%
\pgfpathlineto{\pgfqpoint{5.778119in}{0.653940in}}%
\pgfpathlineto{\pgfqpoint{5.781750in}{0.666367in}}%
\pgfpathlineto{\pgfqpoint{5.785381in}{0.655394in}}%
\pgfpathlineto{\pgfqpoint{5.789012in}{0.650048in}}%
\pgfpathlineto{\pgfqpoint{5.792644in}{0.660078in}}%
\pgfpathlineto{\pgfqpoint{5.796275in}{0.640287in}}%
\pgfpathlineto{\pgfqpoint{5.799906in}{0.639862in}}%
\pgfpathlineto{\pgfqpoint{5.803537in}{0.657465in}}%
\pgfpathlineto{\pgfqpoint{5.807168in}{0.653077in}}%
\pgfpathlineto{\pgfqpoint{5.810799in}{0.652255in}}%
\pgfpathlineto{\pgfqpoint{5.814431in}{0.636792in}}%
\pgfpathlineto{\pgfqpoint{5.818062in}{0.657802in}}%
\pgfpathlineto{\pgfqpoint{5.821693in}{0.675411in}}%
\pgfpathlineto{\pgfqpoint{5.825324in}{0.651003in}}%
\pgfpathlineto{\pgfqpoint{5.828955in}{0.662891in}}%
\pgfpathlineto{\pgfqpoint{5.836218in}{0.617377in}}%
\pgfpathlineto{\pgfqpoint{5.839849in}{0.637988in}}%
\pgfpathlineto{\pgfqpoint{5.843480in}{0.662248in}}%
\pgfpathlineto{\pgfqpoint{5.847111in}{0.662959in}}%
\pgfpathlineto{\pgfqpoint{5.850742in}{0.665252in}}%
\pgfpathlineto{\pgfqpoint{5.854373in}{0.647828in}}%
\pgfpathlineto{\pgfqpoint{5.858005in}{0.645418in}}%
\pgfpathlineto{\pgfqpoint{5.861636in}{0.654182in}}%
\pgfpathlineto{\pgfqpoint{5.865267in}{0.660803in}}%
\pgfpathlineto{\pgfqpoint{5.868898in}{0.630076in}}%
\pgfpathlineto{\pgfqpoint{5.872529in}{0.655687in}}%
\pgfpathlineto{\pgfqpoint{5.876160in}{0.642001in}}%
\pgfpathlineto{\pgfqpoint{5.883423in}{0.671107in}}%
\pgfpathlineto{\pgfqpoint{5.887054in}{0.652140in}}%
\pgfpathlineto{\pgfqpoint{5.890685in}{0.640613in}}%
\pgfpathlineto{\pgfqpoint{5.894316in}{0.642588in}}%
\pgfpathlineto{\pgfqpoint{5.897948in}{0.651845in}}%
\pgfpathlineto{\pgfqpoint{5.901579in}{0.674297in}}%
\pgfpathlineto{\pgfqpoint{5.905210in}{0.666062in}}%
\pgfpathlineto{\pgfqpoint{5.908841in}{0.619235in}}%
\pgfpathlineto{\pgfqpoint{5.912472in}{0.644581in}}%
\pgfpathlineto{\pgfqpoint{5.916103in}{0.659784in}}%
\pgfpathlineto{\pgfqpoint{5.919735in}{0.659203in}}%
\pgfpathlineto{\pgfqpoint{5.923366in}{0.648019in}}%
\pgfpathlineto{\pgfqpoint{5.926997in}{0.665663in}}%
\pgfpathlineto{\pgfqpoint{5.930628in}{0.676993in}}%
\pgfpathlineto{\pgfqpoint{5.934259in}{0.657027in}}%
\pgfpathlineto{\pgfqpoint{5.937890in}{0.643921in}}%
\pgfpathlineto{\pgfqpoint{5.941522in}{0.639110in}}%
\pgfpathlineto{\pgfqpoint{5.945153in}{0.642271in}}%
\pgfpathlineto{\pgfqpoint{5.948784in}{0.647148in}}%
\pgfpathlineto{\pgfqpoint{5.952415in}{0.660263in}}%
\pgfpathlineto{\pgfqpoint{5.956046in}{0.662460in}}%
\pgfpathlineto{\pgfqpoint{5.959677in}{0.665622in}}%
\pgfpathlineto{\pgfqpoint{5.963309in}{0.640824in}}%
\pgfpathlineto{\pgfqpoint{5.966940in}{0.662057in}}%
\pgfpathlineto{\pgfqpoint{5.970571in}{0.648651in}}%
\pgfpathlineto{\pgfqpoint{5.974202in}{0.649888in}}%
\pgfpathlineto{\pgfqpoint{5.977833in}{0.641780in}}%
\pgfpathlineto{\pgfqpoint{5.981465in}{0.667759in}}%
\pgfpathlineto{\pgfqpoint{5.985096in}{0.674176in}}%
\pgfpathlineto{\pgfqpoint{5.988727in}{0.661828in}}%
\pgfpathlineto{\pgfqpoint{5.992358in}{0.642137in}}%
\pgfpathlineto{\pgfqpoint{5.995989in}{0.651690in}}%
\pgfpathlineto{\pgfqpoint{5.999620in}{0.643069in}}%
\pgfpathlineto{\pgfqpoint{6.003252in}{0.624748in}}%
\pgfpathlineto{\pgfqpoint{6.006883in}{0.654134in}}%
\pgfpathlineto{\pgfqpoint{6.010514in}{0.671662in}}%
\pgfpathlineto{\pgfqpoint{6.014145in}{0.671426in}}%
\pgfpathlineto{\pgfqpoint{6.017776in}{0.650316in}}%
\pgfpathlineto{\pgfqpoint{6.021407in}{0.661940in}}%
\pgfpathlineto{\pgfqpoint{6.025039in}{0.646533in}}%
\pgfpathlineto{\pgfqpoint{6.028670in}{0.621659in}}%
\pgfpathlineto{\pgfqpoint{6.032301in}{0.663978in}}%
\pgfpathlineto{\pgfqpoint{6.035932in}{0.664529in}}%
\pgfpathlineto{\pgfqpoint{6.039563in}{0.645879in}}%
\pgfpathlineto{\pgfqpoint{6.043194in}{0.659468in}}%
\pgfpathlineto{\pgfqpoint{6.046826in}{0.654526in}}%
\pgfpathlineto{\pgfqpoint{6.050457in}{0.660824in}}%
\pgfpathlineto{\pgfqpoint{6.054088in}{0.629191in}}%
\pgfpathlineto{\pgfqpoint{6.057719in}{0.648609in}}%
\pgfpathlineto{\pgfqpoint{6.061350in}{0.654655in}}%
\pgfpathlineto{\pgfqpoint{6.068613in}{0.647618in}}%
\pgfpathlineto{\pgfqpoint{6.072244in}{0.664696in}}%
\pgfpathlineto{\pgfqpoint{6.075875in}{0.637399in}}%
\pgfpathlineto{\pgfqpoint{6.079506in}{0.637571in}}%
\pgfpathlineto{\pgfqpoint{6.083137in}{0.651500in}}%
\pgfpathlineto{\pgfqpoint{6.086769in}{0.640618in}}%
\pgfpathlineto{\pgfqpoint{6.090400in}{0.610988in}}%
\pgfpathlineto{\pgfqpoint{6.094031in}{0.652106in}}%
\pgfpathlineto{\pgfqpoint{6.097662in}{0.662135in}}%
\pgfpathlineto{\pgfqpoint{6.101293in}{0.680226in}}%
\pgfpathlineto{\pgfqpoint{6.104924in}{0.653800in}}%
\pgfpathlineto{\pgfqpoint{6.108556in}{0.660223in}}%
\pgfpathlineto{\pgfqpoint{6.112187in}{0.659827in}}%
\pgfpathlineto{\pgfqpoint{6.115818in}{0.622839in}}%
\pgfpathlineto{\pgfqpoint{6.119449in}{0.638405in}}%
\pgfpathlineto{\pgfqpoint{6.123080in}{0.659029in}}%
\pgfpathlineto{\pgfqpoint{6.126711in}{0.672267in}}%
\pgfpathlineto{\pgfqpoint{6.130343in}{0.679952in}}%
\pgfpathlineto{\pgfqpoint{6.133974in}{0.673324in}}%
\pgfpathlineto{\pgfqpoint{6.137605in}{0.625110in}}%
\pgfpathlineto{\pgfqpoint{6.141236in}{0.633913in}}%
\pgfpathlineto{\pgfqpoint{6.144867in}{0.651429in}}%
\pgfpathlineto{\pgfqpoint{6.148498in}{0.659884in}}%
\pgfpathlineto{\pgfqpoint{6.152130in}{0.647295in}}%
\pgfpathlineto{\pgfqpoint{6.155761in}{0.641142in}}%
\pgfpathlineto{\pgfqpoint{6.159392in}{0.650323in}}%
\pgfpathlineto{\pgfqpoint{6.163023in}{0.667335in}}%
\pgfpathlineto{\pgfqpoint{6.166654in}{0.646569in}}%
\pgfpathlineto{\pgfqpoint{6.170285in}{0.640984in}}%
\pgfpathlineto{\pgfqpoint{6.173917in}{0.660236in}}%
\pgfpathlineto{\pgfqpoint{6.177548in}{0.671895in}}%
\pgfpathlineto{\pgfqpoint{6.181179in}{0.666441in}}%
\pgfpathlineto{\pgfqpoint{6.184810in}{0.642406in}}%
\pgfpathlineto{\pgfqpoint{6.188441in}{0.653657in}}%
\pgfpathlineto{\pgfqpoint{6.192073in}{0.666888in}}%
\pgfpathlineto{\pgfqpoint{6.195704in}{0.653396in}}%
\pgfpathlineto{\pgfqpoint{6.199335in}{0.669331in}}%
\pgfpathlineto{\pgfqpoint{6.202966in}{0.674499in}}%
\pgfpathlineto{\pgfqpoint{6.206597in}{0.656421in}}%
\pgfpathlineto{\pgfqpoint{6.213860in}{0.641026in}}%
\pgfpathlineto{\pgfqpoint{6.217491in}{0.651650in}}%
\pgfpathlineto{\pgfqpoint{6.221122in}{0.664329in}}%
\pgfpathlineto{\pgfqpoint{6.224753in}{0.648814in}}%
\pgfpathlineto{\pgfqpoint{6.228384in}{0.658266in}}%
\pgfpathlineto{\pgfqpoint{6.232015in}{0.641941in}}%
\pgfpathlineto{\pgfqpoint{6.235647in}{0.620888in}}%
\pgfpathlineto{\pgfqpoint{6.239278in}{0.651657in}}%
\pgfpathlineto{\pgfqpoint{6.242909in}{0.652373in}}%
\pgfpathlineto{\pgfqpoint{6.246540in}{0.645942in}}%
\pgfpathlineto{\pgfqpoint{6.250171in}{0.668943in}}%
\pgfpathlineto{\pgfqpoint{6.253802in}{0.653633in}}%
\pgfpathlineto{\pgfqpoint{6.257434in}{0.621916in}}%
\pgfpathlineto{\pgfqpoint{6.261065in}{0.652007in}}%
\pgfpathlineto{\pgfqpoint{6.264696in}{0.636418in}}%
\pgfpathlineto{\pgfqpoint{6.268327in}{0.646884in}}%
\pgfpathlineto{\pgfqpoint{6.271958in}{0.666395in}}%
\pgfpathlineto{\pgfqpoint{6.275590in}{0.660117in}}%
\pgfpathlineto{\pgfqpoint{6.279221in}{0.672267in}}%
\pgfpathlineto{\pgfqpoint{6.282852in}{0.644782in}}%
\pgfpathlineto{\pgfqpoint{6.286483in}{0.645844in}}%
\pgfpathlineto{\pgfqpoint{6.290114in}{0.662307in}}%
\pgfpathlineto{\pgfqpoint{6.293745in}{0.657770in}}%
\pgfpathlineto{\pgfqpoint{6.297377in}{0.664129in}}%
\pgfpathlineto{\pgfqpoint{6.301008in}{0.679938in}}%
\pgfpathlineto{\pgfqpoint{6.304639in}{0.654816in}}%
\pgfpathlineto{\pgfqpoint{6.308270in}{0.653103in}}%
\pgfpathlineto{\pgfqpoint{6.311901in}{0.656121in}}%
\pgfpathlineto{\pgfqpoint{6.315532in}{0.628330in}}%
\pgfpathlineto{\pgfqpoint{6.322795in}{0.663528in}}%
\pgfpathlineto{\pgfqpoint{6.326426in}{0.661071in}}%
\pgfpathlineto{\pgfqpoint{6.330057in}{0.660912in}}%
\pgfpathlineto{\pgfqpoint{6.333688in}{0.668642in}}%
\pgfpathlineto{\pgfqpoint{6.337319in}{0.621283in}}%
\pgfpathlineto{\pgfqpoint{6.340951in}{0.645224in}}%
\pgfpathlineto{\pgfqpoint{6.344582in}{0.637655in}}%
\pgfpathlineto{\pgfqpoint{6.348213in}{0.638995in}}%
\pgfpathlineto{\pgfqpoint{6.351844in}{0.666722in}}%
\pgfpathlineto{\pgfqpoint{6.355475in}{0.650426in}}%
\pgfpathlineto{\pgfqpoint{6.359106in}{0.650731in}}%
\pgfpathlineto{\pgfqpoint{6.362738in}{0.645816in}}%
\pgfpathlineto{\pgfqpoint{6.366369in}{0.619560in}}%
\pgfpathlineto{\pgfqpoint{6.370000in}{0.645140in}}%
\pgfpathlineto{\pgfqpoint{6.370000in}{0.645140in}}%
\pgfusepath{stroke}%
\end{pgfscope}%
\begin{pgfscope}%
\pgfsetrectcap%
\pgfsetmiterjoin%
\pgfsetlinewidth{0.803000pt}%
\definecolor{currentstroke}{rgb}{0.000000,0.000000,0.000000}%
\pgfsetstrokecolor{currentstroke}%
\pgfsetdash{}{0pt}%
\pgfpathmoveto{\pgfqpoint{4.554414in}{0.320625in}}%
\pgfpathlineto{\pgfqpoint{4.554414in}{1.485000in}}%
\pgfusepath{stroke}%
\end{pgfscope}%
\begin{pgfscope}%
\pgfsetrectcap%
\pgfsetmiterjoin%
\pgfsetlinewidth{0.803000pt}%
\definecolor{currentstroke}{rgb}{0.000000,0.000000,0.000000}%
\pgfsetstrokecolor{currentstroke}%
\pgfsetdash{}{0pt}%
\pgfpathmoveto{\pgfqpoint{6.370000in}{0.320625in}}%
\pgfpathlineto{\pgfqpoint{6.370000in}{1.485000in}}%
\pgfusepath{stroke}%
\end{pgfscope}%
\begin{pgfscope}%
\pgfsetrectcap%
\pgfsetmiterjoin%
\pgfsetlinewidth{0.803000pt}%
\definecolor{currentstroke}{rgb}{0.000000,0.000000,0.000000}%
\pgfsetstrokecolor{currentstroke}%
\pgfsetdash{}{0pt}%
\pgfpathmoveto{\pgfqpoint{4.554414in}{0.320625in}}%
\pgfpathlineto{\pgfqpoint{6.370000in}{0.320625in}}%
\pgfusepath{stroke}%
\end{pgfscope}%
\begin{pgfscope}%
\pgfsetrectcap%
\pgfsetmiterjoin%
\pgfsetlinewidth{0.803000pt}%
\definecolor{currentstroke}{rgb}{0.000000,0.000000,0.000000}%
\pgfsetstrokecolor{currentstroke}%
\pgfsetdash{}{0pt}%
\pgfpathmoveto{\pgfqpoint{4.554414in}{1.485000in}}%
\pgfpathlineto{\pgfqpoint{6.370000in}{1.485000in}}%
\pgfusepath{stroke}%
\end{pgfscope}%
\begin{pgfscope}%
\definecolor{textcolor}{rgb}{0.000000,0.000000,0.000000}%
\pgfsetstrokecolor{textcolor}%
\pgfsetfillcolor{textcolor}%
\pgftext[x=5.462207in,y=1.568333in,,base]{\color{textcolor}{\rmfamily\fontsize{8.000000}{9.600000}\selectfont\catcode`\^=\active\def^{\ifmmode\sp\else\^{}\fi}\catcode`\%=\active\def%{\%}(c) Finite width}}%
\end{pgfscope}%
\begin{pgfscope}%
\pgfsetbuttcap%
\pgfsetmiterjoin%
\definecolor{currentfill}{rgb}{1.000000,1.000000,1.000000}%
\pgfsetfillcolor{currentfill}%
\pgfsetfillopacity{0.800000}%
\pgfsetlinewidth{1.003750pt}%
\definecolor{currentstroke}{rgb}{0.800000,0.800000,0.800000}%
\pgfsetstrokecolor{currentstroke}%
\pgfsetstrokeopacity{0.800000}%
\pgfsetdash{}{0pt}%
\pgfpathmoveto{\pgfqpoint{5.761155in}{1.065358in}}%
\pgfpathlineto{\pgfqpoint{6.327222in}{1.065358in}}%
\pgfpathquadraticcurveto{\pgfqpoint{6.339444in}{1.065358in}}{\pgfqpoint{6.339444in}{1.077580in}}%
\pgfpathlineto{\pgfqpoint{6.339444in}{1.442222in}}%
\pgfpathquadraticcurveto{\pgfqpoint{6.339444in}{1.454444in}}{\pgfqpoint{6.327222in}{1.454444in}}%
\pgfpathlineto{\pgfqpoint{5.761155in}{1.454444in}}%
\pgfpathquadraticcurveto{\pgfqpoint{5.748933in}{1.454444in}}{\pgfqpoint{5.748933in}{1.442222in}}%
\pgfpathlineto{\pgfqpoint{5.748933in}{1.077580in}}%
\pgfpathquadraticcurveto{\pgfqpoint{5.748933in}{1.065358in}}{\pgfqpoint{5.761155in}{1.065358in}}%
\pgfpathlineto{\pgfqpoint{5.761155in}{1.065358in}}%
\pgfpathclose%
\pgfusepath{stroke,fill}%
\end{pgfscope}%
\begin{pgfscope}%
\pgfsetrectcap%
\pgfsetroundjoin%
\pgfsetlinewidth{1.304875pt}%
\definecolor{currentstroke}{rgb}{0.000000,0.000000,0.000000}%
\pgfsetstrokecolor{currentstroke}%
\pgfsetdash{}{0pt}%
\pgfpathmoveto{\pgfqpoint{5.764211in}{1.412330in}}%
\pgfpathlineto{\pgfqpoint{5.803933in}{1.412330in}}%
\pgfpathlineto{\pgfqpoint{5.843655in}{1.412330in}}%
\pgfusepath{stroke}%
\end{pgfscope}%
\begin{pgfscope}%
\definecolor{textcolor}{rgb}{0.000000,0.000000,0.000000}%
\pgfsetstrokecolor{textcolor}%
\pgfsetfillcolor{textcolor}%
\pgftext[x=5.892544in,y=1.390942in,left,base]{\color{textcolor}{\rmfamily\fontsize{4.400000}{5.280000}\selectfont\catcode`\^=\active\def^{\ifmmode\sp\else\^{}\fi}\catcode`\%=\active\def%{\%}mean field}}%
\end{pgfscope}%
\begin{pgfscope}%
\pgfsetrectcap%
\pgfsetroundjoin%
\pgfsetlinewidth{1.154312pt}%
\definecolor{currentstroke}{rgb}{0.141176,0.345098,0.650980}%
\pgfsetstrokecolor{currentstroke}%
\pgfsetdash{}{0pt}%
\pgfpathmoveto{\pgfqpoint{5.764211in}{1.338424in}}%
\pgfpathlineto{\pgfqpoint{5.803933in}{1.338424in}}%
\pgfpathlineto{\pgfqpoint{5.843655in}{1.338424in}}%
\pgfusepath{stroke}%
\end{pgfscope}%
\begin{pgfscope}%
\definecolor{textcolor}{rgb}{0.000000,0.000000,0.000000}%
\pgfsetstrokecolor{textcolor}%
\pgfsetfillcolor{textcolor}%
\pgftext[x=5.892544in,y=1.317035in,left,base]{\color{textcolor}{\rmfamily\fontsize{4.400000}{5.280000}\selectfont\catcode`\^=\active\def^{\ifmmode\sp\else\^{}\fi}\catcode`\%=\active\def%{\%}$n=32$}}%
\end{pgfscope}%
\begin{pgfscope}%
\pgfsetrectcap%
\pgfsetroundjoin%
\pgfsetlinewidth{1.154312pt}%
\definecolor{currentstroke}{rgb}{0.819608,0.286275,0.356863}%
\pgfsetstrokecolor{currentstroke}%
\pgfsetdash{}{0pt}%
\pgfpathmoveto{\pgfqpoint{5.764211in}{1.264518in}}%
\pgfpathlineto{\pgfqpoint{5.803933in}{1.264518in}}%
\pgfpathlineto{\pgfqpoint{5.843655in}{1.264518in}}%
\pgfusepath{stroke}%
\end{pgfscope}%
\begin{pgfscope}%
\definecolor{textcolor}{rgb}{0.000000,0.000000,0.000000}%
\pgfsetstrokecolor{textcolor}%
\pgfsetfillcolor{textcolor}%
\pgftext[x=5.892544in,y=1.243129in,left,base]{\color{textcolor}{\rmfamily\fontsize{4.400000}{5.280000}\selectfont\catcode`\^=\active\def^{\ifmmode\sp\else\^{}\fi}\catcode`\%=\active\def%{\%}$n=128$}}%
\end{pgfscope}%
\begin{pgfscope}%
\pgfsetrectcap%
\pgfsetroundjoin%
\pgfsetlinewidth{1.154312pt}%
\definecolor{currentstroke}{rgb}{0.000000,0.529412,0.423529}%
\pgfsetstrokecolor{currentstroke}%
\pgfsetdash{}{0pt}%
\pgfpathmoveto{\pgfqpoint{5.764211in}{1.190612in}}%
\pgfpathlineto{\pgfqpoint{5.803933in}{1.190612in}}%
\pgfpathlineto{\pgfqpoint{5.843655in}{1.190612in}}%
\pgfusepath{stroke}%
\end{pgfscope}%
\begin{pgfscope}%
\definecolor{textcolor}{rgb}{0.000000,0.000000,0.000000}%
\pgfsetstrokecolor{textcolor}%
\pgfsetfillcolor{textcolor}%
\pgftext[x=5.892544in,y=1.169223in,left,base]{\color{textcolor}{\rmfamily\fontsize{4.400000}{5.280000}\selectfont\catcode`\^=\active\def^{\ifmmode\sp\else\^{}\fi}\catcode`\%=\active\def%{\%}$n=512$}}%
\end{pgfscope}%
\begin{pgfscope}%
\pgfsetrectcap%
\pgfsetroundjoin%
\pgfsetlinewidth{1.154312pt}%
\definecolor{currentstroke}{rgb}{0.478431,0.243137,0.615686}%
\pgfsetstrokecolor{currentstroke}%
\pgfsetdash{}{0pt}%
\pgfpathmoveto{\pgfqpoint{5.764211in}{1.116706in}}%
\pgfpathlineto{\pgfqpoint{5.803933in}{1.116706in}}%
\pgfpathlineto{\pgfqpoint{5.843655in}{1.116706in}}%
\pgfusepath{stroke}%
\end{pgfscope}%
\begin{pgfscope}%
\definecolor{textcolor}{rgb}{0.000000,0.000000,0.000000}%
\pgfsetstrokecolor{textcolor}%
\pgfsetfillcolor{textcolor}%
\pgftext[x=5.892544in,y=1.095317in,left,base]{\color{textcolor}{\rmfamily\fontsize{4.400000}{5.280000}\selectfont\catcode`\^=\active\def^{\ifmmode\sp\else\^{}\fi}\catcode`\%=\active\def%{\%}$n=2048$}}%
\end{pgfscope}%
\end{pgfpicture}%
\makeatother%
\endgroup%
}
  \captionsetup{font=footnotesize,skip=2pt}
  \caption{Smoothed-Lasso counterpart of Figure~\ref{fig:bures-convergence-triptych}, using the same dimensions $(d_u,d_v)=(16,9)$, teacher,
    initialization scales, and sampling seeds.
    \textbf{(a)} White covariance $\Sigma_0=\alpha\Id_m$ for $\alpha=0.1,0.5,1,5$.
    Dashed curves are the support-gap envelopes $\Delta_0 e^{-4\alpha\kappa t}$; dotted curves use the internal coarse determinant bound of Remark~\ref{rem:coarse-rate}.
    \textbf{(b)} One input-output diagonal pair is scaled by $r\in\{0,1/8,\ldots,1\}$.
    Here $g(\Sigma_0)=2r$; the dashed curve shows the white-reference envelope.
    \textbf{(c)} The population covariance flow from $\Id_m$ and finite-width flows at widths $n=32,128,512,2048$, with five independent trials each.}
  \label{fig:smooth-lasso-triptych}
\end{figure}


\end{document}

%%% Local Variables:
%%% mode: LaTeX
%%% TeX-master: t
%%% End:
